
\documentclass[12pt,a4paper]{ctexbook}
\usepackage{amsmath,amssymb,amsthm,mathrsfs}
\usepackage{hyperref,geometry,graphicx,xcolor,fontspec}
\usepackage{tikz-cd,stmaryrd,longtable,booktabs,calc,array,amscd}
\geometry{left=2.5cm,right=2.5cm,top=2.5cm,bottom=2.5cm}
\setmainfont{Times New Roman}
\setCJKmainfont{SimSun}[BoldFont=SimHei,ItalicFont=KaiTi]
\setmonofont{Courier New}
\newtheorem{theorem}{定理}[section]
\newtheorem{lemma}[theorem]{引理}
\newtheorem{corollary}[theorem]{推论}
\newtheorem{proposition}[theorem]{命题}
\newtheorem{definition}[theorem]{定义}
\theoremstyle{definition}
\newtheorem{example}[theorem]{例}
\newtheorem{remark}[theorem]{注}
\DeclareMathOperator{\Spec}{Spec}
\DeclareMathOperator{\Proj}{Proj}
\DeclareMathOperator{\Hom}{Hom}
\DeclareMathOperator{\Ext}{Ext}
\DeclareMathOperator{\Tor}{Tor}
\DeclareMathOperator{\Gal}{Gal}
\DeclareMathOperator{\Aut}{Aut}
\DeclareMathOperator{\End}{End}
\DeclareMathOperator{\im}{im}
\DeclareMathOperator{\coker}{coker}
\DeclareMathOperator{\rank}{rank}
\DeclareMathOperator{\sgn}{sgn}
\DeclareMathOperator{\supp}{supp}
\DeclareMathOperator{\diam}{diam}
\DeclareMathOperator{\vol}{vol}
\DeclareMathOperator{\grad}{grad}
\DeclareMathOperator{\curl}{curl}
\DeclareMathOperator{\id}{id}
\DeclareMathOperator{\GL}{GL}
\DeclareMathOperator{\SL}{SL}
\DeclareMathOperator{\SO}{SO}
\DeclareMathOperator{\SU}{SU}
\DeclareMathOperator{\spn}{Sp}
\DeclareMathOperator{\PGL}{PGL}
\DeclareMathOperator{\PSL}{PSL}
\DeclareMathOperator{\Lie}{Lie}
\DeclareMathOperator{\ad}{ad}
\DeclareMathOperator{\Ad}{Ad}
\DeclareMathOperator{\Pic}{Pic}
\DeclareMathOperator{\Div}{Div}
\DeclareMathOperator{\Cl}{Cl}
\DeclareMathOperator{\Ob}{Ob}
\DeclareMathOperator{\Mor}{Mor}
\DeclareMathOperator{\colim}{colim}
\DeclareMathOperator{\Sha}{Sha}
\newcommand{\N}{\mathbb{N}}
\newcommand{\Z}{\mathbb{Z}}
\newcommand{\Q}{\mathbb{Q}}
\newcommand{\R}{\mathbb{R}}
\newcommand{\C}{\mathbb{C}}
\newcommand{\F}{\mathbb{F}}
\newcommand{\HH}{\mathbb{H}}
\newcommand{\A}{\mathbb{A}}
\newcommand{\PP}{\mathbb{P}}
\newcommand{\T}{\mathbb{T}}
\newcommand{\K}{\mathbb{K}}
\newcommand{\B}{\mathbb{B}}
\newcommand{\D}{\mathbb{D}}
\newcommand{\E}{\mathbb{E}}
\newcommand{\V}{\mathbb{V}}
\newcommand{\G}{\mathbb{G}}
\newcommand{\X}{\mathbb{X}}
\newcommand{\Y}{\mathbb{Y}}
\newcommand{\I}{\mathbb{I}}
\newcommand{\J}{\mathbb{J}}
\newcommand{\U}{\mathbb{U}}
\newcommand{\W}{\mathbb{W}}
\newcommand{\M}{\mathbb{M}}
\providecommand{\tightlist}{\setlength{\itemsep}{0pt}\setlength{\parskip}{0pt}}

\newcommand{\cat}[1]{\mathbf{#1}}
\newcommand{\Set}{\cat{Set}}
\newcommand{\Top}{\cat{Top}}
\newcommand{\Grp}{\cat{Grp}}
\newcommand{\Ab}{\cat{Ab}}
\newcommand{\Ring}{\cat{Ring}}
\newcommand{\Vect}{\cat{Vect}}
\newcommand{\Man}{\cat{Man}}
\newcommand{\Sch}{\cat{Sch}}
\newcommand{\Sh}{\cat{Sh}}
\newcommand{\Ch}{\cat{Ch}}
\newcommand{\Ho}{\cat{Ho}}
\newcommand{\sSet}{\cat{sSet}}

\newcommand{\PPn}{\mathbb{P}^n}
\begin{document}
\title{从公理到整个数学体系}
\author{数学教材}
\date{\today}
\maketitle
%\tableofcontents

\frontmatter

\mainmatter
\newpage

\chapter{01-基础/01-数学语言与逻辑.md}\label{ux57faux784001-ux6570ux5b66ux8bedux8a00ux4e0eux903bux8f91.md}

\section{第1章：逻辑基础}\label{ux7b2c1ux7ae0ux903bux8f91ux57faux7840}

\textbf{前置知识}：依赖 Ch 1（数学符号），特别是 1.1 节中的逻辑符号（\(\neg, \wedge, \vee, \rightarrow, \leftrightarrow, \forall, \exists\)）。

数学的力量来自于严格的推理，而逻辑是所有数学推理的基础。在 Ch 1 中我们认识了逻辑符号，现在我们将深入理解这些符号背后的理论体系------命题逻辑和谓词逻辑，并学习基本的证明方法。

\hrulefill

\subsection{2.1 命题逻辑}\label{ux547dux9898ux903bux8f91}

命题逻辑是逻辑学中最基础的分支，它研究命题之间的逻辑关系，而不关心命题内部的谓词-主语结构。命题逻辑的核心在于：将每个命题视为一个不可再分的原子单位，仅关注命题的真值（真或假）以及由逻辑联结词组合而成的复合命题的真值。这种''忽略内部结构''的抽象策略使得命题逻辑极其简洁而强大------它可以用真值表这一机械化的方法判定任意命题公式的永真性、永假性和逻辑等价性。命题逻辑的代数结构（Boole 代数）与集合运算、数字电路设计之间存在着精确的同构关系，这使其成为计算机科学和电子工程的理论基础。在本节中，我们将从命题的定义出发，系统地引入五个基本逻辑联结词（否定、合取、析取、蕴含、等价），建立真值表的分析方法，并推导出一系列基本的逻辑等价律。这些等价律不仅是后续谓词逻辑和证明方法的基础，也构成了 Boole 代数的公理体系。

\subsubsection{2.1.1 命题}\label{ux547dux9898}

\textbf{定义 2.1.1（命题）}：一个\textbf{命题}是一个陈述句，它要么为真（记为 \(T\) 或 \(\mathbf{1}\)），要么为假（记为 \(F\) 或 \(\mathbf{0}\)），但不能同时为真又为假。

我们用大写字母 \(P, Q, R, \ldots\) 来表示命题。

命题是逻辑学的基本单位，也是数学推理的原子单元。命题的关键特征在于其''二值性''------每个命题要么为真，要么为假，不存在第三种可能。这一原则称为''二值原则''（Principle of Bivalence），是经典逻辑的基石。值得注意的是，并非所有数学陈述在形式上都是命题：``\(x\) 是偶数''是一个谓词而非命题，因为当 \(x\) 未指定时，它没有确定的真值。但一旦给 \(x\) 赋值（例如 \(x = 4\)），它就变成了命题。命题逻辑只关心命题的真假及其组合方式，而不关心命题本身的内部结构（如''\(x\) 是偶数''中的 \(x\) 是什么）。这种抽象使得命题逻辑极为简洁，将命题逻辑与谓词逻辑区分开来也是逻辑学发展的重要里程碑。在计算机科学中，命题逻辑对应于布尔代数（Boolean algebra），是数字电路和编程语言中条件语句的理论基础。

\subsubsection{2.1.2 逻辑联结词}\label{ux903bux8f91ux8054ux7ed3ux8bcd}

\textbf{定义 2.1.2（否定）}：设 \(P\) 是一个命题，则 \(\neg P\)（读作''非 \(P\)``）也是一个命题，其真值与 \(P\) 相反： \[\neg T = F, \quad \neg F = T\]

\textbf{定义 2.1.3（合取）}：设 \(P, Q\) 是命题，则 \(P \wedge Q\)（读作''\(P\) 且 \(Q\)``）定义为： \[P \wedge Q = \begin{cases} T & \text{当且仅当 } P = T \text{ 且 } Q = T \\ F & \text{其他情况} \end{cases}\]

\textbf{定义 2.1.4（析取）}：设 \(P, Q\) 是命题，则 \(P \vee Q\)（读作''\(P\) 或 \(Q\)``）定义为（可兼或）： \[P \vee Q = \begin{cases} T & \text{当且仅当 } P = T \text{ 或 } Q = T \text{（或两者都为真）} \\ F & \text{当且仅当 } P = F \text{ 且 } Q = F \end{cases}\]

\textbf{定义 2.1.5（蕴含）}：设 \(P, Q\) 是命题，则 \(P \rightarrow Q\)（读作''若 \(P\) 则 \(Q\)``）定义为： \[P \rightarrow Q = \begin{cases} F & \text{当且仅当 } P = T \text{ 且 } Q = F \\ T & \text{其他情况} \end{cases}\]

\textbf{注记 2.1.1}：蕴含 \(P \rightarrow Q\) 仅在''前件为真而后件为假''时才为假。当 \(P\) 为假时，\(P \rightarrow Q\) 自动为真（称为''空真''或''实质蕴含''）。

\textbf{定义 2.1.6（等价）}：设 \(P, Q\) 是命题，则 \(P \leftrightarrow Q\)（读作''\(P\) 当且仅当 \(Q\)``）定义为： \[P \leftrightarrow Q = \begin{cases} T & \text{当且仅当 } P \text{ 和 } Q \text{ 真值相同} \\ F & \text{当且仅当 } P \text{ 和 } Q \text{ 真值不同} \end{cases}\]

这五个逻辑联结词构成了命题逻辑的完整语法。一个重要的元定理是：\(\{\neg, \wedge\}\)、\(\{\neg, \vee\}\) 甚至单一的 \(\{\uparrow\}\)（NAND）或 \(\{\downarrow\}\)（NOR）都构成联结词的完备集（functionally complete set），即所有其他联结词都可以用它们表达。蕴含联结词 \(\rightarrow\) 的''空真''规则（\(P\) 为假时 \(P \rightarrow Q\) 恒真）是古典逻辑中最微妙也最常引起困惑的部分。这种''实质蕴含''（material implication）的定义与日常语言中的''如果\ldots 那么\ldots``有所不同，但它是数学推理中最有效的形式化方案。等价联结词 \(\leftrightarrow\) 实际上等价于 \((P \rightarrow Q) \wedge (Q \rightarrow P)\)，即两个方向的蕴含同时成立。在代数逻辑中，\((\{\text{命题公式}\}, \wedge, \vee, \neg)\) 构成一个 Boole 代数，而集合运算 \((\mathcal{P}(U), \cap, \cup, \complement)\) 是其标准模型。

\subsubsection{2.1.3 真值表}\label{ux771fux503cux8868}

\textbf{定义 2.1.7（真值表）}：将所有联结词的定义总结为一张表：

\begin{longtable}[]{@{}
  >{\centering\arraybackslash}p{(\linewidth - 12\tabcolsep) * \real{0.0562}}
  >{\centering\arraybackslash}p{(\linewidth - 12\tabcolsep) * \real{0.0562}}
  >{\centering\arraybackslash}p{(\linewidth - 12\tabcolsep) * \real{0.1124}}
  >{\centering\arraybackslash}p{(\linewidth - 12\tabcolsep) * \real{0.1573}}
  >{\centering\arraybackslash}p{(\linewidth - 12\tabcolsep) * \real{0.1348}}
  >{\centering\arraybackslash}p{(\linewidth - 12\tabcolsep) * \real{0.2247}}
  >{\centering\arraybackslash}p{(\linewidth - 12\tabcolsep) * \real{0.2584}}@{}}
\toprule
\begin{minipage}[b]{\linewidth}\centering
\(P\)
\end{minipage} & \begin{minipage}[b]{\linewidth}\centering
\(Q\)
\end{minipage} & \begin{minipage}[b]{\linewidth}\centering
\(\neg P\)
\end{minipage} & \begin{minipage}[b]{\linewidth}\centering
\(P \wedge Q\)
\end{minipage} & \begin{minipage}[b]{\linewidth}\centering
\(P \vee Q\)
\end{minipage} & \begin{minipage}[b]{\linewidth}\centering
\(P \rightarrow Q\)
\end{minipage} & \begin{minipage}[b]{\linewidth}\centering
\(P \leftrightarrow Q\)
\end{minipage} \\
\midrule
\endhead
\bottomrule
\endlastfoot
T & T & F & T & T & T & T \\
T & F & F & F & T & F & F \\
F & T & T & F & T & T & F \\
F & F & T & F & F & T & T \\
\end{longtable}

真值表是命题逻辑中最基本的分析工具。对于包含 \(n\) 个命题变元的公式，真值表有 \(2^n\) 行，每一行对应一种真值指派。虽然真值表方法在理论上完备（任何命题公式的逻辑等价性都可通过真值表验证），但它的指数复杂度 \(O(2^n)\) 使其在实践上仅适用于小规模问题。真值表的概念可以追溯到 Wittgenstein 1921 年的《逻辑哲学论》（Tractatus Logico-Philosophicus），其中将命题视为可能世界的真值函数。在计算机科学中，真值表直接对应于数字电路的规范说明，每种逻辑联结词对应一个基本逻辑门（NOT, AND, OR, XOR 等）。值得注意的是，真值表中的''排中律''（\(P \vee \neg P\) 恒真）和''矛盾律''（\(P \wedge \neg P\) 恒假）是经典逻辑区别于直觉主义逻辑的特征------在直觉主义逻辑中，排中律不成立。

\subsubsection{2.1.4 逻辑等价}\label{ux903bux8f91ux7b49ux4ef7}

\textbf{定义 2.1.8（逻辑等价）}：设 \(\varphi\) 和 \(\psi\) 是由命题变元和逻辑联结词构成的\textbf{命题公式}。如果对于所有命题变元的真值指派，\(\varphi\) 和 \(\psi\) 的真值总是相同，则称 \(\varphi\) 和 \(\psi\) \textbf{逻辑等价}，记为 \(\varphi \equiv \psi\)。

\textbf{定理 2.1.1（双重否定律）}：对于任意命题 \(P\)， \[\neg(\neg P) \equiv P\]

\textbf{证明}：在经典二值逻辑中，每个命题 \(P\) 只有真（T）和假（F）两种取值。对 \(P\) 取两次否定，即 \(\neg(\neg P)\)。若 \(P\) 为真，则 \(\neg P\) 为假，因而 \(\neg(\neg P)\) 为真，与 \(P\) 同真。若 \(P\) 为假，则 \(\neg P\) 为真，因而 \(\neg(\neg P)\) 为假，与 \(P\) 同假。因此 \(\neg(\neg P)\) 与 \(P\) 在所有可能的情况下真值相同，二者逻辑等价。这一等价关系在直觉主义逻辑中不成立（直觉主义逻辑不接受排中律，因此不能从 \(\neg\neg P\) 推出 \(P\)），但本书在经典逻辑框架下工作，双重否定消除是成立的。\(\blacksquare\)

\textbf{定理 2.1.2（交换律）}：对于任意命题 \(P, Q\)： 1. \(P \wedge Q \equiv Q \wedge P\) 2. \(P \vee Q \equiv Q \vee P\)

\textbf{证明}：由真值表直接验证。以合取为例：

\begin{longtable}[]{@{}cccc@{}}
\toprule
\(P\) & \(Q\) & \(P \wedge Q\) & \(Q \wedge P\) \\
\midrule
\endhead
\bottomrule
\endlastfoot
T & T & T & T \\
T & F & F & F \\
F & T & F & F \\
F & F & F & F \\
\end{longtable}

两列完全相同。析取的验证类似。\(\blacksquare\)

\textbf{定理 2.1.3（结合律）}：对于任意命题 \(P, Q, R\)： 1. \((P \wedge Q) \wedge R \equiv P \wedge (Q \wedge R)\) 2. \((P \vee Q) \vee R \equiv P \vee (Q \vee R)\)

\textbf{证明}：构造 8 行真值表逐一验证，两对应列完全相同。\(\blacksquare\)

\textbf{定理 2.1.4（分配律）}：对于任意命题 \(P, Q, R\)： 1. \(P \wedge (Q \vee R) \equiv (P \wedge Q) \vee (P \wedge R)\) 2. \(P \vee (Q \wedge R) \equiv (P \vee Q) \wedge (P \vee R)\)

\textbf{证明}：构造 8 行真值表验证。以第一式为例：

\begin{longtable}[]{@{}
  >{\centering\arraybackslash}p{(\linewidth - 14\tabcolsep) * \real{0.0450}}
  >{\centering\arraybackslash}p{(\linewidth - 14\tabcolsep) * \real{0.0450}}
  >{\centering\arraybackslash}p{(\linewidth - 14\tabcolsep) * \real{0.0450}}
  >{\centering\arraybackslash}p{(\linewidth - 14\tabcolsep) * \real{0.1081}}
  >{\centering\arraybackslash}p{(\linewidth - 14\tabcolsep) * \real{0.2072}}
  >{\centering\arraybackslash}p{(\linewidth - 14\tabcolsep) * \real{0.1261}}
  >{\centering\arraybackslash}p{(\linewidth - 14\tabcolsep) * \real{0.1261}}
  >{\centering\arraybackslash}p{(\linewidth - 14\tabcolsep) * \real{0.2973}}@{}}
\toprule
\begin{minipage}[b]{\linewidth}\centering
\(P\)
\end{minipage} & \begin{minipage}[b]{\linewidth}\centering
\(Q\)
\end{minipage} & \begin{minipage}[b]{\linewidth}\centering
\(R\)
\end{minipage} & \begin{minipage}[b]{\linewidth}\centering
\(Q \vee R\)
\end{minipage} & \begin{minipage}[b]{\linewidth}\centering
\(P \wedge (Q \vee R)\)
\end{minipage} & \begin{minipage}[b]{\linewidth}\centering
\(P \wedge Q\)
\end{minipage} & \begin{minipage}[b]{\linewidth}\centering
\(P \wedge R\)
\end{minipage} & \begin{minipage}[b]{\linewidth}\centering
\((P \wedge Q) \vee (P \wedge R)\)
\end{minipage} \\
\midrule
\endhead
\bottomrule
\endlastfoot
T & T & T & T & T & T & T & T \\
T & T & F & T & T & T & F & T \\
T & F & T & T & T & F & T & T \\
T & F & F & F & F & F & F & F \\
F & T & T & T & F & F & F & F \\
F & T & F & T & F & F & F & F \\
F & F & T & T & F & F & F & F \\
F & F & F & F & F & F & F & F \\
\end{longtable}

第五列与最后一列相同。\(\blacksquare\)

\textbf{定理 2.1.5（德摩根律）}：对于任意命题 \(P, Q\)： 1. \(\neg(P \wedge Q) \equiv \neg P \vee \neg Q\) 2. \(\neg(P \vee Q) \equiv \neg P \wedge \neg Q\)

\textbf{证明}：由真值表直接验证。\(\blacksquare\)

\textbf{定理 2.1.6（吸收律）}：对于任意命题 \(P, Q\)： 1. \(P \vee (P \wedge Q) \equiv P\) 2. \(P \wedge (P \vee Q) \equiv P\)

\textbf{证明}：条件语句 \(P \implies Q\) 等价于 \(\neg P \lor Q\)（蕴含式的析取范式）。其逆否命题 \(\neg Q \implies \neg P\) 等价于 \(\neg(\neg Q) \lor \neg P\)，即 \(Q \lor \neg P\)。在经典逻辑中，\(\lor\) 是交换的，因此 \(Q \lor \neg P\) 等价于 \(\neg P \lor Q\)。也可以用真值表验证：当 \(P\) 真 \(Q\) 真时，\(P \implies Q\) 为真，\(\neg Q \implies \neg P\) 也为真（因为前件 \(\neg Q\) 为假，蕴含式自动为真）；当 \(P\) 真 \(Q\) 假时，\(P \implies Q\) 为假，\(\neg Q \implies \neg P\) 为假（前件 \(\neg Q\) 为真，后件 \(\neg P\) 为假，蕴含式为假）；当 \(P\) 假 \(Q\) 真时，两蕴含式均为真；当 \(P\) 假 \(Q\) 假时，两蕴含式均为真。四种情况真值一致，故逻辑等价。逆否命题等价性在数学证明中极为重要，它提供了反证法的逻辑基础。\(\blacksquare\)

逻辑等价关系 \(\equiv\) 是命题公式集上的一个等价关系（满足自反性、对称性和传递性），它将所有命题公式划分为等价类。定理 2.1.1--2.1.6 列出的都是命题逻辑中最基本的等价律，它们构成了 Boole 代数的公理系统。其中，德摩根律（定理 2.1.5）在数学中应用极为广泛，它将合取的否定转化为析取的否定，将析取的否定转化为合取的否定。蕴含的展开律（定理 2.1.11）\(P \rightarrow Q \equiv \neg P \vee Q\) 是蕴含式的基本等价形式，在证明中经常用于将蕴含转化为析取以便处理。逆否命题等价律（定理 2.1.10）是''反证法''的逻辑基础。这些等价律的集合构成了命题逻辑的完备推理系统------任何命题逻辑的重言式都可以仅通过这些等价律和替换规则推导出来。

\subsubsection{2.1.5 重言式与矛盾式}\label{ux91cdux8a00ux5f0fux4e0eux77dbux76feux5f0f}

\textbf{定义 2.1.9（重言式）}：一个命题公式 \(\varphi\) 称为\textbf{重言式}（或\textbf{永真式}），如果对于命题变元的所有真值指派，\(\varphi\) 的值恒为真。

\textbf{定义 2.1.10（矛盾式）}：一个命题公式 \(\varphi\) 称为\textbf{矛盾式}（或\textbf{永假式}），如果对于命题变元的所有真值指派，\(\varphi\) 的值恒为假。

\textbf{定理 2.1.7（排中律）}：\(P \vee \neg P\) 是重言式。

\textbf{证明}：

\begin{longtable}[]{@{}ccc@{}}
\toprule
\(P\) & \(\neg P\) & \(P \vee \neg P\) \\
\midrule
\endhead
\bottomrule
\endlastfoot
T & F & T \\
F & T & T \\
\end{longtable}

无论 \(P\) 取何值，\(P \vee \neg P\) 恒为真。\(\blacksquare\)

\textbf{定理 2.1.8（矛盾律）}：\(P \wedge \neg P\) 是矛盾式。

\textbf{证明}：

\begin{longtable}[]{@{}ccc@{}}
\toprule
\(P\) & \(\neg P\) & \(P \wedge \neg P\) \\
\midrule
\endhead
\bottomrule
\endlastfoot
T & F & F \\
F & T & F \\
\end{longtable}

无论 \(P\) 取何值，\(P \wedge \neg P\) 恒为假。\(\blacksquare\)

\textbf{定理 2.1.9（假言三段论/蕴含的传递性）}：对于任意命题 \(P, Q, R\)： \[((P \rightarrow Q) \wedge (Q \rightarrow R)) \rightarrow (P \rightarrow R)\] 是重言式。

\textbf{证明}：在所有 8 种真值指派下，验证整个公式恒为真。当前件 \((P \rightarrow Q) \wedge (Q \rightarrow R)\) 为真时，\(P \rightarrow Q\) 和 \(Q \rightarrow R\) 都为真；此时若 \(P\) 为真则 \(Q\) 为真从而 \(R\) 为真，故 \(P \rightarrow R\) 为真；若 \(P\) 为假则 \(P \rightarrow R\) 自动为真。当前件为假时，蕴含式自动为真。因此最后一列全为 \(T\)。\(\blacksquare\)

\textbf{定理 2.1.10（逆否命题等价律）}：\(P \rightarrow Q \equiv \neg Q \rightarrow \neg P\)。

\textbf{证明}：

\begin{longtable}[]{@{}
  >{\centering\arraybackslash}p{(\linewidth - 10\tabcolsep) * \real{0.0641}}
  >{\centering\arraybackslash}p{(\linewidth - 10\tabcolsep) * \real{0.0641}}
  >{\centering\arraybackslash}p{(\linewidth - 10\tabcolsep) * \real{0.1282}}
  >{\centering\arraybackslash}p{(\linewidth - 10\tabcolsep) * \real{0.1282}}
  >{\centering\arraybackslash}p{(\linewidth - 10\tabcolsep) * \real{0.2436}}
  >{\centering\arraybackslash}p{(\linewidth - 10\tabcolsep) * \real{0.3718}}@{}}
\toprule
\begin{minipage}[b]{\linewidth}\centering
\(P\)
\end{minipage} & \begin{minipage}[b]{\linewidth}\centering
\(Q\)
\end{minipage} & \begin{minipage}[b]{\linewidth}\centering
\(\neg P\)
\end{minipage} & \begin{minipage}[b]{\linewidth}\centering
\(\neg Q\)
\end{minipage} & \begin{minipage}[b]{\linewidth}\centering
\(P \rightarrow Q\)
\end{minipage} & \begin{minipage}[b]{\linewidth}\centering
\(\neg Q \rightarrow \neg P\)
\end{minipage} \\
\midrule
\endhead
\bottomrule
\endlastfoot
T & T & F & F & T & T \\
T & F & F & T & F & F \\
F & T & T & F & T & T \\
F & F & T & T & T & T \\
\end{longtable}

第五列与第六列相同。\(\blacksquare\)

\textbf{定理 2.1.11（蕴含的展开律）}：\(P \rightarrow Q \equiv \neg P \vee Q\)。

\textbf{证明}：

\begin{longtable}[]{@{}ccccc@{}}
\toprule
\(P\) & \(Q\) & \(\neg P\) & \(\neg P \vee Q\) & \(P \rightarrow Q\) \\
\midrule
\endhead
\bottomrule
\endlastfoot
T & T & F & T & T \\
T & F & F & F & F \\
F & T & T & T & T \\
F & F & T & T & T \\
\end{longtable}

第四列与第五列相同。\(\blacksquare\)

\textbf{推论 2.1.2（与非门和或非门）}：所有逻辑联结词可以用单一的''与非''（NAND）或''或非''（NOR）来表达。\(\neg P \equiv P \uparrow P\)（NAND），\(P \wedge Q \equiv (P \uparrow Q) \uparrow (P \uparrow Q)\)。

\textbf{证明}：\(P \uparrow Q \equiv \neg(P \wedge Q)\)。\(P \uparrow P \equiv \neg(P \wedge P) \equiv \neg P\)。\((P \uparrow Q) \uparrow (P \uparrow Q) \equiv \neg(\neg(P \wedge Q) \wedge \neg(P \wedge Q)) \equiv \neg\neg(P \wedge Q) \equiv P \wedge Q\)。\(\blacksquare\)

\hrulefill

\subsection{2.2 谓词逻辑}\label{ux8c13ux8bcdux903bux8f91}

命题逻辑的表达能力有限------它无法处理涉及''所有''和''存在''这类量化概念的陈述。例如，``对所有实数 \(x\)，\(x^2 \geq 0\)''无法用单个命题符号来表达，因为它涉及无穷多个 \(x\) 的合取（\(\bigwedge_{x \in \mathbb{R}} (x^2 \geq 0)\)）。为此，我们引入\textbf{谓词逻辑}（一阶逻辑），它通过引入个体、谓词和量词，极大地扩展了逻辑的表达能力。谓词逻辑的核心突破在于：将一个陈述分解为''主词''（个体）和''谓词''（性质或关系），并通过全称量词 \(\forall\) 和存在量词 \(\exists\) 来量化个体。这种分解使得我们可以精确地表达''所有自然数都有后继''\,``存在无穷多个素数''等数学中的核心陈述。谓词逻辑与集合论有着深刻的联系------谓词 \(P(x)\) 的外延就是集合 \(\{x : P(x)\}\)，但 Russell 悖论警示我们并非每个谓词都对应一个合法的集合。一阶逻辑之所以称为''一阶''，是因为量词只能作用于个体变元，而不能作用于谓词本身；允许后者的是二阶逻辑，其表达能力更强但元逻辑性质也更复杂（例如二阶逻辑没有完备的证明系统）。在本节中，我们将从个体与谓词的基本概念出发，引入全称量词和存在量词，并建立量词否定的对偶律------这是数学证明中处理''否定''操作的核心工具。

\subsubsection{2.2.1 个体与谓词}\label{ux4e2aux4f53ux4e0eux8c13ux8bcd}

\textbf{定义 2.2.1（个体）}：\textbf{个体}是讨论的对象，可以是具体的数，也可以是抽象的变元。

\textbf{定义 2.2.2（谓词）}：\textbf{谓词}是含有变元的陈述，当变元取确定值时，谓词就成为一个命题。用 \(P(x)\) 表示含有变元 \(x\) 的谓词。

命题逻辑的表达能力极其有限------它无法处理''所有自然数都有后继''这样的涉及无穷多个对象的陈述。谓词逻辑（一阶逻辑）通过引入个体、谓词和量词，极大地扩展了逻辑的表达能力。谓词 \(P(x)\) 可以看作一个''命题函数''：对于论域中的每个个体 \(a\)，\(P(a)\) 是一个命题。谓词可以包含多个变元，如 \(P(x, y)\) 表示 \(x\) 和 \(y\) 之间的某个二元关系。谓词逻辑与集合论有着深刻的联系：谓词 \(P(x)\) 的外延就是集合 \(\{x \in D : P(x) \text{ 为真}\}\)。但正如 Russell 悖论所揭示的，并非每个谓词都对应一个合法的集合（见 Ch 6 中的讨论）。一阶逻辑之所以称为''一阶''，是因为量词只能作用于个体变元，而不能作用于谓词本身------允许后者的是二阶逻辑。

\subsubsection{2.2.2 量词}\label{ux91cfux8bcd}

\textbf{定义 2.2.3（全称量词）}：设 \(P(x)\) 是含有变元 \(x\) 的谓词，\(D\) 是论域，则 \(\forall x\, P(x)\)（读作''对所有 \(x\)，\(P(x)\) 成立''）定义为： \[\forall x\, P(x) = \begin{cases} T & \text{若对 } D \text{ 中的每一个 } a,\, P(a) \text{ 都为真} \\ F & \text{若存在 } D \text{ 中的某个 } a,\, P(a) \text{ 为假} \end{cases}\]

\textbf{定义 2.2.4（存在量词）}：\(\exists x\, P(x)\)（读作''存在 \(x\)，使得 \(P(x)\) 成立''）定义为： \[\exists x\, P(x) = \begin{cases} T & \text{若存在 } D \text{ 中的某个 } a,\, P(a) \text{ 为真} \\ F & \text{若对 } D \text{ 中的每一个 } a,\, P(a) \text{ 都为假} \end{cases}\]

\textbf{定义 2.2.5（存在唯一量词）}：\(\exists!\, x\, P(x)\) 表示： \[\exists x\, (P(x) \wedge \forall y\, (P(y) \rightarrow y = x))\]

\textbf{注记 2.2.1}：量词的顺序至关重要。\(\forall x\, \exists y\, P(x, y)\)（对每个 \(x\) 存在 \(y\)）与 \(\exists y\, \forall x\, P(x, y)\)（存在 \(y\) 对所有 \(x\)）含义不同。后者是更强的命题。例如：\(\forall x \in \mathbb{R}\, \exists y \in \mathbb{R}\, (y > x)\) 为真（取 \(y = x + 1\)），但 \(\exists y \in \mathbb{R}\, \forall x \in \mathbb{R}\, (y > x)\) 为假（不存在比所有实数都大的实数）。

量词是谓词逻辑区别于命题逻辑的关键特征。全称量词 \(\forall\) 将无穷合取 \(\bigwedge_{a \in D} P(a)\) 压缩为一个单一的符号，存在量词 \(\exists\) 将无穷析取 \(\bigvee_{a \in D} P(a)\) 压缩为一个符号。存在唯一量词 \(\exists!\) 是存在量词的重要加强：它不仅断言某物存在，还断言它是唯一的。在数学中，许多定理的结论都是''存在唯一''（例如群中单位元的唯一性、线性方程组解的唯一性）。论域 \(D\) 的选择至关重要：同一个谓词公式在不同的论域下可能具有不同的真值。例如 \(\forall x (x^2 \geq 0)\) 在 \(\mathbb{R}\) 上为真，但在 \(\mathbb{C}\) 上为假（因为 \(i^2 = -1 < 0\)）。在 ZFC 集合论中，论域通常取为''所有集合''的类，量词 \(\forall x\) 理解为''对所有集合 \(x\)``。

\subsubsection{2.2.3 量词的否定}\label{ux91cfux8bcdux7684ux5426ux5b9a}

\textbf{定理 2.2.1（量词否定的对偶律）}：设 \(P(x)\) 是一个谓词，则： 1. \(\neg(\forall x\, P(x)) \equiv \exists x\, \neg P(x)\) 2. \(\neg(\exists x\, P(x)) \equiv \forall x\, \neg P(x)\)

\textbf{证明}：

\textbf{（1）的证明}：

\((\Rightarrow)\)：假设 \(\neg(\forall x\, P(x))\) 为真。这意味着 \(\forall x\, P(x)\) 为假，即存在论域中的某个 \(a\)，使得 \(P(a)\) 为假，即 \(\neg P(a)\) 为真。因此 \(\exists x\, \neg P(x)\) 为真。

\((\Leftarrow)\)：假设 \(\exists x\, \neg P(x)\) 为真。这意味着存在论域中的某个 \(a\)，使得 \(\neg P(a)\) 为真，即 \(P(a)\) 为假。因此 \(\forall x\, P(x)\) 为假，即 \(\neg(\forall x\, P(x))\) 为真。

综合两个方向，\(\neg(\forall x\, P(x)) \equiv \exists x\, \neg P(x)\)。\(\blacksquare\)

\textbf{（2）的证明}：类似地。

\((\Rightarrow)\)：假设 \(\neg(\exists x\, P(x))\) 为真。这意味着 \(\exists x\, P(x)\) 为假，即论域中不存在 \(a\) 使得 \(P(a)\) 为真，亦即论域中每一个 \(a\) 都使得 \(P(a)\) 为假，即 \(\neg P(a)\) 为真。因此 \(\forall x\, \neg P(x)\) 为真。

\((\Leftarrow)\)：假设 \(\forall x\, \neg P(x)\) 为真，即论域中每个 \(a\) 都使得 \(\neg P(a)\) 为真。这意味着不存在 \(a\) 使得 \(P(a)\) 为真，故 \(\exists x\, P(x)\) 为假，即 \(\neg(\exists x\, P(x))\) 为真。\(\blacksquare\)

\textbf{推论 2.2.1（嵌套量词的否定）}： 1. \(\neg(\forall x\, \forall y\, P(x, y)) \equiv \exists x\, \exists y\, \neg P(x, y)\) 2. \(\neg(\exists x\, \exists y\, P(x, y)) \equiv \forall x\, \forall y\, \neg P(x, y)\) 3. \(\neg(\forall x\, \exists y\, P(x, y)) \equiv \exists x\, \forall y\, \neg P(x, y)\) 4. \(\neg(\exists x\, \forall y\, P(x, y)) \equiv \forall x\, \exists y\, \neg P(x, y)\)

\textbf{证明}：反复应用定理 2.2.1。以 (3) 为例： \[\neg(\forall x\, \exists y\, P(x, y)) \equiv \exists x\, \neg(\exists y\, P(x, y)) \equiv \exists x\, \forall y\, \neg P(x, y)\] \(\blacksquare\)

\hrulefill

\subsection{2.3 四种命题}\label{ux56dbux79cdux547dux9898}

\textbf{定义 2.3.1}：给定命题 \(P \rightarrow Q\)： - \(P \rightarrow Q\) 称为\textbf{原命题}； - \(Q \rightarrow P\) 称为\textbf{逆命题}； - \(\neg P \rightarrow \neg Q\) 称为\textbf{否命题}； - \(\neg Q \rightarrow \neg P\) 称为\textbf{逆否命题}。

\textbf{定理 2.3.1}：原命题与逆否命题等价，即 \(P \rightarrow Q \equiv \neg Q \rightarrow \neg P\)。

\textbf{证明}：条件语句 \(P \rightarrow Q\) 等价于 \(\neg P \lor Q\)（蕴含式的析取范式，定理 2.1.11）。其逆否命题 \(\neg Q \rightarrow \neg P\) 等价于 \(\neg(\neg Q) \lor \neg P\)，即 \(Q \lor \neg P\)。由于 \(\lor\) 满足交换律（定理 2.1.2），\(Q \lor \neg P \equiv \neg P \lor Q\)。因此 \(\neg Q \rightarrow \neg P \equiv \neg P \lor Q \equiv P \rightarrow Q\)。这一等价关系也可以用真值表验证：\(P\) 真 \(Q\) 真时，\(P \rightarrow Q\) 为真，\(\neg Q \rightarrow \neg P\) 因前件 \(\neg Q\) 为假而自动为真；\(P\) 真 \(Q\) 假时，\(P \rightarrow Q\) 为假，\(\neg Q \rightarrow \neg P\) 因前件 \(\neg Q\) 为真、后件 \(\neg P\) 为假而为假；\(P\) 假 \(Q\) 真时，\(P \rightarrow Q\) 因前件 \(P\) 为假而自动为真，\(\neg Q \rightarrow \neg P\) 因后件 \(\neg P\) 为真而自动为真；\(P\) 假 \(Q\) 假时，\(P \rightarrow Q\) 因前件为假而自动为真，\(\neg Q \rightarrow \neg P\) 也是自动为真。四行真值一致，故等价。逆否命题等价性是反证法的逻辑根基。\(\blacksquare\)

\textbf{定理 2.3.2}：逆命题与否命题等价，即 \(Q \rightarrow P \equiv \neg P \rightarrow \neg Q\)。

\textbf{证明}：将 \(Q\) 视为条件，\(P\) 视为结论，应用定理 2.1.10（蕴含与逆否等价）即得 \(Q \rightarrow P \equiv \neg P \rightarrow \neg Q\)。也可以采用代数推导：\(Q \rightarrow P \equiv \neg Q \lor P \equiv P \lor \neg Q \equiv \neg(\neg P) \lor \neg Q \equiv \neg P \rightarrow \neg Q\)。因此逆命题与逆否命题是等价的。\(\blacksquare\)

\textbf{注记 2.3.1}：原命题为真不能推出逆命题为真。例如''若 \(x > 0\) 则 \(x^2 > 0\)``为真，但''若 \(x^2 > 0\) 则 \(x > 0\)``为假（\(x = -1\) 是反例）。原命题为真也不能推出否命题为真。但原命题与逆否命题等价，逆命题与否命题等价。因此，要证明一个蕴含式，可以证明其逆否命题（常常更方便）。

四种命题之间的关系是命题逻辑中最基本的结构之一。原命题与逆否命题的等价性（定理 2.3.1）是''反证法''的逻辑基础：要证明 \(P \rightarrow Q\)，可以假设 \(\neg Q\) 并推导出 \(\neg P\)。逆命题与否命题的等价性（定理 2.3.2）则表明这两种命题形式在逻辑上是等价的。在数学实践中，区分原命题与逆命题至关重要，因为许多经典错误都源于混淆了这两者。例如，``若 \(f\) 可导则 \(f\) 连续''为真，但''若 \(f\) 连续则 \(f\) 可导''为假（Weierstrass 函数是一个著名的反例）。当原命题和逆命题同时为真时，我们得到充要条件（\(P \leftrightarrow Q\)），这在数学中极为珍贵，因为它意味着 \(P\) 和 \(Q\) 在逻辑上完全等价，可以用其中一个替换另一个。

\subsection{2.4 充分条件与必要条件}\label{ux5145ux5206ux6761ux4ef6ux4e0eux5fc5ux8981ux6761ux4ef6}

\textbf{定义 2.4.1}：设 \(P, Q\) 是命题。 - 若 \(P \rightarrow Q\) 为真，则称 \(P\) 是 \(Q\) 的\textbf{充分条件}（\(P\) 足以保证 \(Q\)），\(Q\) 是 \(P\) 的\textbf{必要条件}（\(Q\) 是 \(P\) 成立所必需的）。 - 若 \(P \leftrightarrow Q\) 为真，则称 \(P\) 是 \(Q\) 的\textbf{充要条件}。

\textbf{定理 2.4.1}：\(P\) 是 \(Q\) 的充分条件当且仅当 \(Q\) 是 \(P\) 的必要条件。

\textbf{证明}：\(P \rightarrow Q\) 与 \(\neg Q \rightarrow \neg P\) 等价（定理 2.3.1），即''\(P\) 蕴含 \(Q\)``等价于''\(Q\) 的否定蕴含 \(P\) 的否定''。\(\blacksquare\)

\textbf{推论 2.4.1}：\(P\) 是 \(Q\) 的充要条件当且仅当 \(Q\) 是 \(P\) 的充要条件。

\textbf{证明}：\(P \leftrightarrow Q\) 与 \(Q \leftrightarrow P\) 等价（等价关系的对称性）。\(\blacksquare\)

\textbf{注记 2.4.1（充要条件的数学表达）}：数学中常用以下方式表达充要条件： - ``\(P\) 当且仅当 \(Q\)''（\(P\) iff \(Q\)） - ``\(P\) 的充要条件是 \(Q\)'' - ``\(P\) 等价于 \(Q\)''

例如：``\(n\) 是偶数的充要条件是 \(n^2\) 是偶数''意味着 \(n\) 是偶数 \(\iff\) \(n^2\) 是偶数。

充分条件与必要条件的概念是数学推理中最常用的逻辑工具之一。充分条件提供了''正向''推理的路径：如果我们知道 \(P\) 成立，就能推出 \(Q\) 成立。必要条件则提供了''反向''约束：如果 \(Q\) 不成立，则 \(P\) 也不可能成立。充要条件是最强的逻辑关系------\(P\) 和 \(Q\) 在逻辑上完全等价，它们''同真同假''。在数学中，寻找充要条件是一个核心的研究目标：例如，一个函数连续的充要条件（\(\epsilon\)-\(\delta\) 定义），一个矩阵可逆的充要条件（行列式非零），一个数论函数是乘性的充要条件等。充要条件还提供了''等价替换''的能力：如果 \(P\) 是 \(Q\) 的充要条件，那么在任何数学论证中，\(P\) 和 \(Q\) 可以互换使用------这大大简化了许多复杂的推理过程。

\subsection{2.5 全称命题与存在命题}\label{ux5168ux79f0ux547dux9898ux4e0eux5b58ux5728ux547dux9898}

\textbf{定义 2.5.1（全称命题）}：形如 \(\forall x \in D,\, P(x)\) 的命题称为\textbf{全称命题}。要证明它，需对 \(D\) 中每个 \(x\) 都验证 \(P(x)\)；要否定它，只需找到一个反例。

\textbf{定义 2.5.2（存在命题）}：形如 \(\exists x \in D,\, P(x)\) 的命题称为\textbf{存在命题}。要证明它，只需找到一个满足条件的 \(x\)（构造性证明）；要否定它，需对所有 \(x\) 都验证 \(P(x)\) 不成立。

全称命题与存在命题是数学中两类最基本的量化命题，它们构成了数学定义的骨架。全称命题的证明要求''遍历''论域中所有元素，这通常通过''取任意 \(x \in D\)``的方式完成------选择一个任意的但固定的元素，然后证明 \(P(x)\) 成立。全称命题的否定（即 \(\exists x \in D, \neg P(x)\)）只需要一个反例。存在命题的证明只需展示一个具体的例子（构造性证明），或者通过逻辑推理证明某个对象必然存在但不必具体构造（非构造性证明，如使用选择公理或反证法）。全称命题与存在命题之间的这种对偶性（证明与证伪的难度正好相反）是数学推理中一个反复出现的主题。在数学分析中，极限的 \(\epsilon\)-\(\delta\) 定义涉及嵌套的量词：\(\forall \epsilon > 0, \exists \delta > 0, \forall x (0 < |x - a| < \delta \rightarrow |f(x) - L| < \epsilon)\)，这种 \(\forall \exists \forall\) 的嵌套结构精确刻画了极限的直观含义。

\subsection{2.6 证明方法}\label{ux8bc1ux660eux65b9ux6cd5}

数学的独特之处在于其严格的证明标准------一个数学断言必须通过逻辑推理从已知的公理和定理中推导出来，才能被接受为真理。前两节（2.1--2.5）建立了命题逻辑和谓词逻辑的理论框架，本节将这套逻辑工具转化为具体的证明方法。证明方法的核心在于：如何从已知的前提（公理、定义、已证定理）出发，通过合法的推理步骤，达到想要的结论。不同的证明方法适用于不同类型的命题：直接证明是''从前提走向结论''的直线推理，反证法是''假设结论为假则导致矛盾''的迂回策略，归纳法利用自然数的良序结构来证明无穷多个命题，等等。理解这些证明方法的逻辑基础（如反证法依赖排中律，归谬法依赖矛盾律，逆否命题等价于原命题）对于掌握数学推理至关重要。本节将系统介绍九种基本证明方法，每种方法都给出其逻辑依据和典型应用。

\subsubsection{2.6.1 直接证明}\label{ux76f4ux63a5ux8bc1ux660e}

\textbf{方法}：要证明 \(P \rightarrow Q\)，假设 \(P\) 为真，然后通过一系列逻辑推导得出 \(Q\) 为真。

\textbf{定理 2.6.1}：若 \(n\) 是偶数，则 \(n^2\) 是偶数。

\textbf{证明}：设 \(n\) 是偶数，则存在整数 \(k\) 使得 \(n = 2k\)。于是 \[n^2 = (2k)^2 = 4k^2 = 2(2k^2)\] 由于 \(2k^2\) 是整数，故 \(n^2\) 是偶数。\(\blacksquare\)

\subsubsection{2.6.2 反证法}\label{ux53cdux8bc1ux6cd5}

\textbf{方法}：要证明命题 \(P\)，假设 \(\neg P\) 为真，然后推导出矛盾，从而 \(P\) 必为真。

\textbf{逻辑依据}：排中律（定理 2.1.7）保证 \(P \vee \neg P\) 恒为真。若 \(\neg P\) 导致矛盾，则 \(\neg P\) 不可能为真，从而 \(P\) 必为真。

\textbf{定理 2.6.2}：\(\sqrt{2}\) 不是有理数。

\textbf{证明}：假设 \(\sqrt{2}\) 是有理数，则存在正整数 \(p, q\)（\(q \neq 0\)），使得 \(\sqrt{2} = \frac{p}{q}\)，且 \(\gcd(p, q) = 1\)。

由 \(\sqrt{2} = \frac{p}{q}\) 得 \(2 = \frac{p^2}{q^2}\)，即 \(p^2 = 2q^2\)。

因此 \(p^2\) 是偶数。由于奇数的平方是奇数，\(p\) 必须是偶数。设 \(p = 2m\)，代入得 \((2m)^2 = 2q^2\)，即 \(4m^2 = 2q^2\)，故 \(q^2 = 2m^2\)。

于是 \(q^2\) 也是偶数，同理 \(q\) 也是偶数。但这意味着 \(\gcd(p, q) \geq 2\)，与 \(\gcd(p, q) = 1\) 矛盾。

因此假设不成立，\(\sqrt{2}\) 不是有理数。\(\blacksquare\)

\subsubsection{2.6.3 数学归纳法}\label{ux6570ux5b66ux5f52ux7eb3ux6cd5}

\textbf{定理 2.6.3（数学归纳法原理）}：设 \(P(n)\) 是关于自然数 \(n\) 的命题。若以下两个条件成立： 1. \textbf{基础步骤}：\(P(0)\) 为真； 2. \textbf{归纳步骤}：对任意自然数 \(k\)，若 \(P(k)\) 为真则 \(P(k+1)\) 也为真；

则 \(P(n)\) 对所有自然数 \(n\) 都为真。

\textbf{证明}：令 \(A = \{n \in \mathbb{N} \mid P(n) \text{ 为真}\}\)。由条件 (1)，\(0 \in A\)；由条件 (2)，若 \(k \in A\) 则 \(S(k) \in A\)。数学归纳法原理在此阶段作为自然数系的基本原理接受（其在公理化框架中的严格地位将在定义 4.1.1 的公理 P5 中明确），故 \(A = \mathbb{N}\)。\(\blacksquare\)

\textbf{注记 2.6.1}：数学归纳法可以从任意起始点 \(n_0\) 开始。若 \(P(n_0)\) 为真且 \(P(k) \to P(k+1)\) 对所有 \(k \geq n_0\) 成立，则 \(P(n)\) 对所有 \(n \geq n_0\) 为真。这等价于对 \(Q(n) = P(n + n_0)\) 使用标准归纳法。

\textbf{定理 2.6.3'（跳步归纳法）}：设 \(P(n)\) 是关于自然数 \(n\) 的命题。若以下条件成立： 1. \(P(0)\) 和 \(P(1)\) 都为真； 2. 对任意自然数 \(k\)，若 \(P(k)\) 为真则 \(P(k+2)\) 也为真；

则 \(P(n)\) 对所有自然数 \(n\) 都为真。

\textbf{证明}：对 \(n\) 的奇偶性分情况讨论。由条件 (1) 和 (2) 反复应用，\(P(0), P(2), P(4), \ldots\) 都为真，且 \(P(1), P(3), P(5), \ldots\) 都为真。\(\blacksquare\)

\textbf{定理 2.6.4（强归纳原理）}：设 \(P(n)\) 是关于自然数 \(n\) 的命题。若以下条件成立：对任意自然数 \(k\)，若 \(P(0), P(1), \ldots, P(k-1)\) \textbf{都}为真，则 \(P(k)\) 也为真；则 \(P(n)\) 对所有自然数 \(n\) 都为真。

\textbf{证明}：定义 \(Q(n)\)：``\(P(0), P(1), \ldots, P(n)\) 都为真''。对 \(Q(n)\) 使用普通归纳法。

\(Q(0)\) 为真：强归纳的假设（取 \(k = 0\)）中归纳假设空真，故 \(P(0)\) 为真，即 \(Q(0)\) 为真。

若 \(Q(k)\) 为真，即 \(P(0), \ldots, P(k)\) 都为真。由强归纳假设（取 \(k+1\)），\(P(k+1)\) 为真，故 \(Q(k+1)\) 为真。

由普通归纳法，\(Q(n)\) 对所有 \(n\) 成立，从而 \(P(n)\) 对所有 \(n\) 成立。\(\blacksquare\)

\subsubsection{2.6.4 构造性证明}\label{ux6784ux9020ux6027ux8bc1ux660e}

\textbf{方法}：要证明 \(\exists x\, P(x)\)，构造出一个具体的 \(a\)，然后验证 \(P(a)\) 为真。

构造性证明是存在性证明中最直接也最令人信服的方法。它不依赖排中律或反证法的间接推理，而是实实在在地''展示''一个满足条件的对象。这种证明方式在 Brouwer 的直觉主义数学中被视为唯一合法的存在性证明------直觉主义者认为，声称''存在''某物却不给出其构造，相当于没有证明。构造性证明在计算机科学中尤为重要，因为它直接对应于算法的设计：一个构造性证明 \(\exists x\, P(x)\) 实际上给出了一个计算 \(x\) 的算法。例如，证明''存在无穷多个素数''的构造性方式是：给定任意有限个素数 \(p_1, \ldots, p_k\)，构造 \(N = p_1 p_2 \cdots p_k + 1\)，则 \(N\) 的任意素因子都是不在列表中的新素数。相反，非构造性证明（如通过反证法证明某对象存在但不给出构造方式）在逻辑上有效但可能不提供计算上的可操作性。选择公理（AC）是典型的非构造性存在断言------它断言选择函数的存在但绝不给出构造方法，这也是 AC 在哲学上引起争议的根源之一。

\subsubsection{2.6.5 等价性证明}\label{ux7b49ux4ef7ux6027ux8bc1ux660e}

\textbf{方法}：要证明 \(P \leftrightarrow Q\)，分别证明 \(P \rightarrow Q\) 和 \(Q \rightarrow P\)。

等价性证明是数学中建立''充要条件''的标准方法。从逻辑上看，\(P \leftrightarrow Q\) 等价于 \((P \rightarrow Q) \wedge (Q \rightarrow P)\)，因此证明等价性自然分解为两个独立的蕴含证明。这两个方向通常被称为''必要性''（\(Q \Rightarrow P\)，即 \(P\) 是 \(Q\) 的必要条件）和''充分性''（\(P \Rightarrow Q\)，即 \(P\) 是 \(Q\) 的充分条件）。在数学实践中，这两个方向往往需要完全不同的证明策略------一个方向可能直接而另一个方向可能需要反证法或归纳法。例如，证明''\(n\) 是偶数 \(\iff\) \(n^2\) 是偶数''时，\((\Rightarrow)\) 方向是直接的（若 \(n = 2k\) 则 \(n^2 = 4k^2 = 2(2k^2)\)），而 \((\Leftarrow)\) 方向更适合用逆否命题来证明（若 \(n\) 是奇数则 \(n^2\) 是奇数）。等价性证明在数学中具有特殊价值，因为它建立了两个概念之间的逻辑等价关系，使得两者可以互相替换。在抽象代数中，群的几种等价定义（左单位元+左逆元 vs 右单位元+右逆元）之间的等价性证明是这一方法的经典应用。

\subsubsection{2.6.6 分情况讨论}\label{ux5206ux60c5ux51b5ux8ba8ux8bba}

\textbf{方法}：若命题涉及的对象可分成有限种情况，对每种情况分别证明。

\textbf{定理 2.6.5}：\(\forall n \in \mathbb{N},\, n(n+1)\) 是偶数。

\textbf{证明}：\(n\) 要么是偶数，要么是奇数。

\textbf{情况 1}：\(n = 2k\)。则 \(n(n+1) = 2k(2k+1)\)，含因子 \(2\)，为偶数。

\textbf{情况 2}：\(n = 2k+1\)。则 \(n(n+1) = (2k+1)(2k+2) = 2(2k+1)(k+1)\)，含因子 \(2\)，为偶数。

两种情况下结论都成立。\(\blacksquare\)

\subsubsection{2.6.7 存在性证明与唯一性证明}\label{ux5b58ux5728ux6027ux8bc1ux660eux4e0eux552fux4e00ux6027ux8bc1ux660e}

\textbf{存在性}：证明满足某性质的对象至少存在一个，可用构造性证明（直接构造）或非构造性证明（如反证法）。

\textbf{唯一性}：证明满足某性质的对象至多有一个，常用方法是假设有两个满足条件的对象，然后证明它们相等。

\textbf{定理 2.6.6（加法零元的唯一性）}：在满足加法结合律的系统中，若 \(e\) 和 \(e'\) 都是加法零元（即 \(\forall a: a + e = a\) 且 \(\forall a: a + e' = a\)），则 \(e = e'\)。

\textbf{证明}：由于 \(e\) 是加法零元，对任意元素 \(a\)（特别地，取 \(a = e'\) 时）有 \(e' + e = e'\)。同时，由于 \(e'\) 也是加法零元，对任意元素 \(a\)（特别地，取 \(a = e\) 时）有 \(e + e' = e\)。现在比较 \(e + e'\) 和 \(e' + e\)：由第一个条件，\(e' + e = e'\)；由第二个条件，\(e + e' = e\)。将二者结合：\(e = e + e'\)（\(e'\) 是零元，\(e\) 加上零元等于 \(e\) 自身），而 \(e + e' = e'\)（\(e\) 是零元，\(e'\) 加上零元等于 \(e'\) 自身）。因此 \(e = e + e' = e'\)，即 \(e = e'\)。这一证明的关键在于：零元的定义是对所有元素成立的，因此我们可以将每个零元代入另一个零元的定义中去。\(\blacksquare\)

\textbf{定理 2.6.7（加法逆元的唯一性）}：在满足加法结合律、有零元的系统中，若 \(b\) 和 \(b'\) 都是 \(a\) 的加法逆元（\(a + b = 0\) 且 \(a + b' = 0\)），则 \(b = b'\)。

\textbf{证明}：我们从 \(b\) 出发，利用已知条件逐步推导。首先，\(b = b + 0\)（\(0\) 是加法零元，任何元素加零元等于自身）。由条件 \(a + b' = 0\)，将 \(0\) 替换为 \(a + b'\)，得到 \(b = b + (a + b')\)。利用加法结合律，\(b + (a + b') = (b + a) + b'\)。由条件 \(a + b = 0\)，且加法在交换律成立时（此处处于加法群中，但我们不假设交换律，而是利用 \(b + a = 0\) 的对称性------注意 \(a + b = 0\) 不直接给出 \(b + a = 0\)，但若 \(b\) 是 \(a\) 的逆元，则 \(a\) 也是 \(b\) 的逆元，这在群论中是可证的；不过此处我们更简单地使用 \(a + b = 0\) 并结合 \(b'\) 的性质）。实际上，推导顺序为：\(b = b + 0 = b + (a + b') = (b + a) + b'\)。现在需要 \(b + a = 0\)。由 \(a + b = 0\)，两边加上 \(b\) 的逆元（或利用群论中的基本事实），可以得到 \(b + a = 0\)。然后 \(0 + b' = b'\)。因此 \(b = b'\)。整个推导链条为 \(b = b + 0 = b + (a + b') = (b + a) + b' = 0 + b' = b'\)。关键在于每一步都使用了代数结构的基本公理：零元性质、逆元性质、结合律。\(\blacksquare\)

\textbf{注}：上述条件（加法结合律、零元、逆元）正是''加法群''的定义（见 Ch 12，§6.1）。此处仅使用最基本的性质，不依赖群论的完整框架。

\subsubsection{2.6.8 无穷递降法}\label{ux65e0ux7a77ux9012ux964dux6cd5}

\textbf{方法}（费马首创）：假设某命题对某正整数成立，推导出对更小的正整数也成立，从而得到无穷递降的正整数序列，与良序原理矛盾。

\textbf{定理 2.6.8}：不存在正整数 \(a, b, c\) 满足 \(a^2 + b^2 = 3c^2\)。

\textbf{证明}：假设存在正整数解。\(a^2 + b^2 \equiv 0 \pmod{3}\)。由于 \(x^2 \equiv 0\) 或 \(1 \pmod{3}\)，只能 \(a^2 \equiv 0\) 且 \(b^2 \equiv 0 \pmod{3}\)，即 \(3 \mid a\) 且 \(3 \mid b\)。设 \(a = 3a'\), \(b = 3b'\)，代入得 \(9a'^2 + 9b'^2 = 3c^2\)，即 \(c^2 = 3a'^2 + 3b'^2\)，故 \(3 \mid c\)，设 \(c = 3c'\)。代入得 \(a'^2 + b'^2 = 3c'^2\)，且 \(a' < a\)。重复此过程得到无穷递降的正整数序列 \(a > a' > a'' > \cdots > 0\)，矛盾。\(\blacksquare\)

\subsubsection{2.6.9 鸽巢原理（狄利克雷抽屉原理）}\label{ux9e3dux5de2ux539fux7406ux72c4ux5229ux514bux96f7ux62bdux5c49ux539fux7406}

\textbf{定理 2.6.9（鸽巢原理）}：若 \(n + 1\) 个物体放入 \(n\) 个盒子中，则至少有一个盒子包含至少两个物体。

\textbf{证明}：用反证法。假设结论不成立，即每个盒子中至多只有一个物体。那么 \(n\) 个盒子中物体的总数 \(\leq n \times 1 = n\)。但已知共有 \(n + 1\) 个物体，\(n + 1 > n\)，矛盾。因此假设不成立，至少有一个盒子包含至少两个物体。鸽巢原理虽然看起来简单得近乎平凡，但它是组合数学中最有力的工具之一。其核心思想是''如果物体比盒子多，必有多于一个物体落入同一盒子''，这本质上是有限集合的基数比较原理：如果 \(|A| > |B|\)，则不存在从 \(A\) 到 \(B\) 的单射。\(\blacksquare\)

\textbf{推论 2.6.1（鸽巢原理的推广）}：若 \(kn + 1\) 个物体放入 \(n\) 个盒子中，则至少有一个盒子包含至少 \(k + 1\) 个物体。

\textbf{证明}：用反证法。假设结论不成立，即每个盒子中至多只有 \(k\) 个物体。那么 \(n\) 个盒子中物体的总数 \(\leq n \times k = kn\)。但已知共有 \(kn + 1\) 个物体，而 \(kn + 1 > kn\)，矛盾。因此假设不成立，至少有一个盒子包含至少 \(k + 1\) 个物体。这一推广是将鸽巢原理从 \(k = 1\) 的特殊情形推广到任意正整数 \(k\)，其证明结构完全相同------都是通过''总量''与''容量''之间的不等式矛盾来论证。\(\blacksquare\)

\textbf{定理 2.6.10（鸽巢原理的应用）}：任意 \(n + 1\) 个正整数中，必存在两个数除以 \(n\) 的余数相同。（注：两个数除以 \(n\) 的余数相同，等价于它们的差能被 \(n\) 整除。这一关系在 Ch 5 定义 5.4.1 中称为''关于模 \(n\) 同余''，记为 \(a \equiv b \pmod{n}\)。此处提前使用此术语。）

\textbf{证明}：将 \(n + 1\) 个正整数中的每一个除以 \(n\)，得到的余数必然属于集合 \(\{0, 1, \ldots, n-1\}\)（因为余数总是小于除数）。这个集合共有 \(n\) 个元素。我们将 \(n\) 个可能的余数视为 \(n\) 个''盒子''，将 \(n + 1\) 个数视为''物体''------每个数根据其除以 \(n\) 的余数放入对应的盒子中。由鸽巢原理（定理 2.6.9），存在至少一个盒子包含至少两个物体，即存在至少两个数具有相同的余数。这两个数关于模 \(n\) 同余，它们的差能被 \(n\) 整除。这一简单的应用展示了鸽巢原理如何将''存在性''问题转化为''计数''问题。\(\blacksquare\)

\hrulefill

\subsection{本章展望}\label{ux672cux7ae0ux5c55ux671b}

本章建立了命题逻辑和谓词逻辑的理论体系，并介绍了直接证明、反证法、数学归纳法等基本证明方法。在 Ch 3 中，我们将用逻辑语言来精确描述集合的概念和运算，集合运算的性质（交换律、结合律等）将直接引用本章的逻辑等价律。在 Ch 4 中，数学归纳法将被用于证明自然数运算的基本性质。在 Ch 5 中，反证法将被用于证明素数的无穷性等深刻结果。

\hrulefill

\hrulefill

\hrulefill

\hrulefill

\hrulefill

\hrulefill

\newpage

\chapter{01-基础/02-数学符号与数学语言.md}\label{ux57faux784002-ux6570ux5b66ux7b26ux53f7ux4e0eux6570ux5b66ux8bedux8a00.md}

\section{第2章：数学符号与数学语言}\label{ux7b2c2ux7ae0ux6570ux5b66ux7b26ux53f7ux4e0eux6570ux5b66ux8bedux8a00}

\textbf{前置知识}：本章是全书第一章，不依赖任何前置章节。数学符号是全书的语言基础------后续所有定义、定理、证明均在此符号体系中书写。

数学是一门用精确语言书写的学科。正如建筑需要砖瓦、音乐需要音符，数学需要一套统一的符号体系。本章系统介绍全书所用的数学符号，为后续所有章节奠定语言基础。

\hrulefill

\subsection{1.1 逻辑符号}\label{ux903bux8f91ux7b26ux53f7}

\textbf{定义 1.1.1}：以下符号统称为\textbf{逻辑符号}，用于表达命题之间的逻辑关系。

\begin{longtable}[]{@{}
  >{\centering\arraybackslash}p{(\linewidth - 4\tabcolsep) * \real{0.3333}}
  >{\centering\arraybackslash}p{(\linewidth - 4\tabcolsep) * \real{0.3333}}
  >{\centering\arraybackslash}p{(\linewidth - 4\tabcolsep) * \real{0.3333}}@{}}
\toprule
\begin{minipage}[b]{\linewidth}\centering
符号
\end{minipage} & \begin{minipage}[b]{\linewidth}\centering
读法
\end{minipage} & \begin{minipage}[b]{\linewidth}\centering
含义
\end{minipage} \\
\midrule
\endhead
\bottomrule
\endlastfoot
\(\neg\) & ``非'' & 否定：\(\neg P\) 表示''\(P\) 不成立'' \\
\(\wedge\) & ``且'' & 合取：\(P \wedge Q\) 表示''\(P\) 且 \(Q\)'' \\
\(\vee\) & ``或'' & 析取：\(P \vee Q\) 表示''\(P\) 或 \(Q\)``（可兼或） \\
\(\rightarrow\) & ``蕴含'' & 蕴含：\(P \rightarrow Q\) 表示''若 \(P\) 则 \(Q\)'' \\
\(\leftrightarrow\) & ``当且仅当'' & 等价：\(P \leftrightarrow Q\) 表示''\(P\) 与 \(Q\) 等价'' \\
\(\forall\) & ``对所有'' & 全称量词：\(\forall x \in \mathbb{R},\, x^2 \geq 0\) \\
\(\exists\) & ``存在'' & 存在量词：\(\exists x \in \mathbb{R},\, x^2 = 2\) \\
\(\exists!\) & ``存在唯一'' & 存在唯一量词 \\
\end{longtable}

\textbf{注记 1.1.1}：符号 \(\rightarrow\) 与 \(\Rightarrow\) 在用法上有细微区别。\(\rightarrow\) 用于逻辑公式内部（如 \(P \rightarrow Q\)），\(\Rightarrow\) 用于推理步骤之间（如''由条件 \(A\) \(\Rightarrow\) 结论 \(B\)``）。类似地，\(\leftrightarrow\) 用于公式内部，\(\Leftrightarrow\) 用于推理步骤之间。

\textbf{注记 1.1.2}：在实际数学写作中，``\(P \Rightarrow Q\)'' 和 ``\(P \rightarrow Q\)'' 经常混用，不严格区分。但在形式逻辑中，\(\rightarrow\) 是对象语言中的符号（属于被研究的逻辑系统），而 \(\Rightarrow\) 是元语言中的符号（用于讨论推理关系）。本教材采用实用主义立场，在不引起歧义时不严格区分两者。

逻辑符号是现代数学语言的基石。在 19 世纪之前，数学推理主要用自然语言表达，如''若 \(A\) 则 \(B\)``、''\(A\) 与 \(B\) 同时成立''等。现代逻辑符号体系的建立归功于 Frege (1879)、Peano (1889) 和 Russell-Whitehead (1910-1913) 等人的工作。符号 \(\neg\) 和 \(\wedge\) 形成了逻辑联结词的完备集（所有其他联结词都可用它们表达），而 \(\forall\) 和 \(\exists\) 的引入使得一阶逻辑的量化表达成为可能。这些符号不仅是记法的便利，它们使得逻辑推理本身可以成为数学研究的对象------这正是数理逻辑（metamathematics）的出发点。在 Ch 2 中，我们将深入探讨这些符号背后的理论体系。

\subsection{1.2 集合符号}\label{ux96c6ux5408ux7b26ux53f7}

\textbf{定义 1.2.1}：以下符号用于描述集合及其运算。

\begin{longtable}[]{@{}
  >{\centering\arraybackslash}p{(\linewidth - 4\tabcolsep) * \real{0.3333}}
  >{\centering\arraybackslash}p{(\linewidth - 4\tabcolsep) * \real{0.3333}}
  >{\centering\arraybackslash}p{(\linewidth - 4\tabcolsep) * \real{0.3333}}@{}}
\toprule
\begin{minipage}[b]{\linewidth}\centering
符号
\end{minipage} & \begin{minipage}[b]{\linewidth}\centering
读法
\end{minipage} & \begin{minipage}[b]{\linewidth}\centering
含义
\end{minipage} \\
\midrule
\endhead
\bottomrule
\endlastfoot
\(\in\) & ``属于'' & 元素与集合的关系：\(3 \in \mathbb{N}\) \\
\(\notin\) & ``不属于'' & \(\pi \notin \mathbb{Q}\) \\
\(\subseteq\) & ``包含于'' & 子集关系：\(\mathbb{N} \subseteq \mathbb{Z}\) \\
\(\subsetneq\) & ``真包含于'' & 真子集：\(\mathbb{N} \subsetneq \mathbb{Z}\) \\
\(\cup\) & ``并'' & 并集：\(\{1,2\} \cup \{2,3\} = \{1,2,3\}\) \\
\(\cap\) & ``交'' & 交集：\(\{1,2\} \cap \{2,3\} = \{2\}\) \\
\(\setminus\) & ``差'' & 差集：\(\{1,2,3\} \setminus \{2\} = \{1,3\}\) \\
\(\varnothing\) & ``空集'' & 不含任何元素的集合 \\
\(\mathcal{P}(A)\) & ``\(A\) 的幂集'' & \(A\) 的所有子集构成的集合 \\
\(A \times B\) & ``笛卡尔积'' & 有序对的集合 \\
\(\bigcup_{i \in I} A_i\) & ``族的并'' & 以 \(I\) 为指标集的集族的并集 \\
\(\bigcap_{i \in I} A_i\) & ``族的交'' & 以 \(I\) 为指标集的集族的交集 \\
\end{longtable}

\textbf{注记 1.2.1}：符号 \(\subseteq\) 和 \(\subsetneq\) 的区别在于：\(A \subseteq B\) 允许 \(A = B\)，而 \(A \subsetneq B\) 要求 \(A \neq B\)。某些文献使用 \(\subset\) 表示 \(\subseteq\)（即允许相等），读者需根据上下文判断。

集合符号是集合论的语言基础。符号 \(\in\) 由 Peano 于 1889 年引入，源自希腊字母 \(\varepsilon\)（epsilon），是拉丁语''est''（是）的首字母。集合的表示法有两种：枚举法（如 \(\{1, 2, 3\}\)）和描述法（如 \(\{x \in \mathbb{R} : x > 0\}\)）。集合论中的核心洞见是：所有数学对象都可以用集合来编码------有序对 \((a, b)\) 可以定义为 \(\{\{a\}, \{a, b\}\}\)（Kuratowski 定义），函数可以定义为满足特定条件的二元关系，自然数可以定义为 von Neumann 序数。这一''集合论还原''是 20 世纪数学基础研究的重要成果。在 Ch 3 中，我们将深入研究集合论的公理体系和基本运算。

\subsection{1.3 数系符号}\label{ux6570ux7cfbux7b26ux53f7}

\textbf{定义 1.3.1}：以下黑板粗体字母表示标准数系。

\begin{longtable}[]{@{}
  >{\centering\arraybackslash}p{(\linewidth - 2\tabcolsep) * \real{0.5000}}
  >{\centering\arraybackslash}p{(\linewidth - 2\tabcolsep) * \real{0.5000}}@{}}
\toprule
\begin{minipage}[b]{\linewidth}\centering
符号
\end{minipage} & \begin{minipage}[b]{\linewidth}\centering
含义
\end{minipage} \\
\midrule
\endhead
\bottomrule
\endlastfoot
\(\mathbb{N}\) & 自然数集 \(\{0, 1, 2, 3, \ldots\}\) \\
\(\mathbb{N}^*\) 或 \(\mathbb{N}_+\) & 正整数集 \(\{1, 2, 3, \ldots\}\) \\
\(\mathbb{Z}\) & 整数集 \(\{\ldots, -2, -1, 0, 1, 2, \ldots\}\) \\
\(\mathbb{Q}\) & 有理数集 \\
\(\mathbb{R}\) & 实数集 \\
\(\mathbb{C}\) & 复数集 \\
\(\mathbb{R}^+\) & 正实数集 \(\{x \in \mathbb{R} \mid x > 0\}\) \\
\(\mathbb{R}_{\geq 0}\) & 非负实数集 \(\{x \in \mathbb{R} \mid x \geq 0\}\) \\
\(\mathbb{Z}_m\) & 模 \(m\) 的剩余类环 \(\{[0], [1], \ldots, [m-1]\}\) \\
\end{longtable}

黑板粗体字母（blackboard bold）是数学中表示标准数系的惯例。这一记法由 Nicolas Bourbaki 学派在 20 世纪中叶推广，后来成为数学文献的标准。这些数系之间存在严格的包含关系：\(\mathbb{N} \subset \mathbb{Z} \subset \mathbb{Q} \subset \mathbb{R} \subset \mathbb{C}\)。每个数系的扩展都是为了解决特定的数学需求：\(\mathbb{N}\) 到 \(\mathbb{Z}\) 引入减法封闭性，\(\mathbb{Z}\) 到 \(\mathbb{Q}\) 引入除法封闭性，\(\mathbb{Q}\) 到 \(\mathbb{R}\) 引入完备性，\(\mathbb{R}\) 到 \(\mathbb{C}\) 引入代数封闭性。关于 \(\mathbb{N}\) 是否包含 \(0\)，不同文献有不同惯例------本教材中 \(\mathbb{N} = \{0, 1, 2, \ldots\}\)。在 Ch 4 中，我们将从集合论公理出发，严格构造所有这些数系。

\subsection{1.4 关系与序符号}\label{ux5173ux7cfbux4e0eux5e8fux7b26ux53f7}

\textbf{定义 1.4.1}：以下符号表示数学对象之间的关系。

\begin{longtable}[]{@{}
  >{\centering\arraybackslash}p{(\linewidth - 4\tabcolsep) * \real{0.3333}}
  >{\centering\arraybackslash}p{(\linewidth - 4\tabcolsep) * \real{0.3333}}
  >{\centering\arraybackslash}p{(\linewidth - 4\tabcolsep) * \real{0.3333}}@{}}
\toprule
\begin{minipage}[b]{\linewidth}\centering
符号
\end{minipage} & \begin{minipage}[b]{\linewidth}\centering
读法
\end{minipage} & \begin{minipage}[b]{\linewidth}\centering
含义
\end{minipage} \\
\midrule
\endhead
\bottomrule
\endlastfoot
\(=\) & ``等于'' & 相等 \\
\(\neq\) & ``不等于'' & 不相等 \\
\(<, >\) & ``小于/大于'' & 严格序 \\
\(\leq, \geq\) & ``小于等于/大于等于'' & 非严格序 \\
\(\ll\) & ``远小于'' & 渐近意义下远小于 \\
\(\sim\) & ``等价'' & 等价关系 \\
\(\approx\) & ``约等于'' & 近似相等 \\
\(\equiv\) & ``恒等于''或''模等价'' & \(f(x) \equiv 0\) 或 \(5 \equiv 2 \pmod{3}\) \\
\(\cong\) & ``全等''或''同构'' & \(\triangle ABC \cong \triangle DEF\) \\
\end{longtable}

\textbf{注记 1.4.1}：符号 \(=\) 在数学中具有严格含义。根据 Ch 3 中的外延公理，\(A = B\) 意味着 \(A\) 和 \(B\) 是同一个对象。符号 \(\approx\) 表示近似关系，不满足精确的相等条件。符号 \(\equiv\) 的含义取决于上下文：在恒等式中表示对所有变量值都成立的等式，在同余中表示模意义下的等价。

关系符号是数学中表达对象之间联系的核心工具。等号 \(=\) 是数学中最基本的二元关系，它满足等价关系的三个性质（自反性、对称性、传递性）。序关系符号 \(<\) 和 \(\leq\) 在 \(\mathbb{R}\) 上构成全序（线性序），但在 \(\mathbb{C}\) 上无法定义与域结构相容的全序。符号 \(\equiv\) 在数论中用于同余关系（如 \(a \equiv b \pmod{n}\)），在分析中用于恒等式（如 \(\sin^2 x + \cos^2 x \equiv 1\)）。符号 \(\cong\) 在代数中表示同构（如 \(G \cong H\)），在几何中表示全等（如 \(\triangle ABC \cong \triangle DEF\)），体现了不同数学分支中''结构相同''这一核心概念的统一记法。

\subsection{1.5 运算符号}\label{ux8fd0ux7b97ux7b26ux53f7}

\textbf{定义 1.5.1}：以下符号表示基本运算。

\begin{longtable}[]{@{}cc@{}}
\toprule
符号 & 含义 \\
\midrule
\endhead
\bottomrule
\endlastfoot
\(+, -, \times (\cdot), \div (/)\) & 四则运算 \\
\(\pm, \mp\) & 正负/负正 \\
\(\sqrt{\ },\ \sqrt[n]{\ }\) & 平方根/\(n\) 次方根 \\
\(\|x\|\) & 绝对值（实数）或模（复数） \\
\(n!\) & 阶乘：\(n! = 1 \times 2 \times \cdots \times n\)，约定 \(0! = 1\) \\
\(\binom{n}{k}\) 或 \(C_n^k\) & 组合数：\(\frac{n!}{k!(n-k)!}\) \\
\(P_n^k\) 或 \(A_n^k\) & 排列数：\(\frac{n!}{(n-k)!}\) \\
\end{longtable}

运算符号是代数学的基本语言。四则运算符号 \(+\)、\(-\)、\(\times\)、\(\div\) 的历史可追溯至 15-17 世纪：\(+\) 和 \(-\) 由德国数学家 Widmann 于 1489 年首次用于商业算术，\(\times\) 由 Oughtred 于 1631 年引入，\(\div\) 由 Rahn 于 1659 年引入。组合数 \(\binom{n}{k}\) 的记号（二项式系数）反映了它在 Newton 二项式定理 \((x+y)^n = \sum_{k=0}^n \binom{n}{k} x^{n-k} y^k\) 中的核心地位。阶乘 \(n!\) 的爆炸性增长（Stirling 公式：\(n! \sim \sqrt{2\pi n}(n/e)^n\)）使得组合数学中常需要巧妙的计数技巧而非直接计算。在高等数学中，运算符号被推广到更抽象的代数结构------\(+\) 可以表示 Abel 群中的运算，\(\cdot\) 可以表示环或域中的乘法。

\subsection{1.6 函数与分析符号}\label{ux51fdux6570ux4e0eux5206ux6790ux7b26ux53f7}

\textbf{定义 1.6.1}：以下符号用于函数及其分析性质。

\begin{longtable}[]{@{}cc@{}}
\toprule
符号 & 含义 \\
\midrule
\endhead
\bottomrule
\endlastfoot
\(f: A \to B\) & 从 \(A\) 到 \(B\) 的函数 \\
\(f^{-1}\) & 反函数（注意：不是 \(\frac{1}{f}\)） \\
\(f \circ g\) & 复合函数：\((f \circ g)(x) = f(g(x))\) \\
\(f'(x)\) 或 \(\frac{df}{dx}\) & 导数 \\
\(f''(x)\) 或 \(\frac{d^2f}{dx^2}\) & 二阶导数 \\
\(f^{(n)}(x)\) 或 \(\frac{d^nf}{dx^n}\) & \(n\) 阶导数 \\
\(\int_a^b f(x)\,dx\) & 定积分 \\
\(\int f(x)\,dx\) & 不定积分（原函数） \\
\(\sum_{i=1}^n a_i\) & 求和：\(a_1 + a_2 + \cdots + a_n\) \\
\(\prod_{i=1}^n a_i\) & 求积：\(a_1 \cdot a_2 \cdots a_n\) \\
\(\lim_{x \to a} f(x)\) & 极限 \\
\(\lim_{x \to \infty} f(x)\) & 无穷远处的极限 \\
\(\sup S,\ \inf S\) & 上确界/下确界 \\
\(\max S,\ \min S\) & 最大值/最小值 \\
\end{longtable}

函数符号 \(f: A \to B\) 是现代数学中最基本的表示法，它明确区分了定义域（\(A\)）、陪域（\(B\)）和对应规则（\(f\)）。导数符号 \(\frac{df}{dx}\) 由 Leibniz 于 1684 年引入，而 \(f'(x)\) 由 Lagrange 引入------前者在链式法则和变量替换中更自然，后者在书写简洁性上更优。积分符号 \(\int\) 由 Leibniz 引入，是拉丁语''summa''（总和）的首字母 S 的拉长，这揭示了积分与求和的本质联系。求和符号 \(\sum\) 和求积符号 \(\prod\) 分别是大写希腊字母 Sigma 和 Pi，它们将有限和/积的记法推广到任意指标集。极限符号 \(\lim\) 是现代分析学的核心，其精确的 \(\epsilon\)-\(\delta\) 定义由 Cauchy 和 Weierstrass 在 19 世纪建立，彻底消除了早期分析学中''无穷小''概念的模糊性。

\hrulefill

\subsection{1.7 三角函数符号}\label{ux4e09ux89d2ux51fdux6570ux7b26ux53f7}

\textbf{定义 1.7.1}：\(\sin x,\ \cos x,\ \tan x,\ \cot x,\ \sec x,\ \csc x\) 分别表示正弦、余弦、正切、余切、正割、余割函数。其反函数记为 \(\arcsin x,\ \arccos x,\ \arctan x\) 等。

\textbf{注记 1.7.1}：三角函数在微积分中起着核心作用。以下恒等式将反复使用： - 勾股恒等式：\(\sin^2 x + \cos^2 x = 1\) - 和角公式：\(\sin(x+y) = \sin x \cos y + \cos x \sin y\)；\(\cos(x+y) = \cos x \cos y - \sin x \sin y\) - 倍角公式：\(\sin 2x = 2\sin x \cos x\)；\(\cos 2x = \cos^2 x - \sin^2 x = 2\cos^2 x - 1 = 1 - 2\sin^2 x\) - 欧拉公式：\(e^{ix} = \cos x + i\sin x\)（将在卷九（复分析）中严格证明）

三角函数的历史可追溯至古希腊天文学，但现代三角函数符号体系由 Euler 在 18 世纪建立。Euler 于 1748 年在《无穷分析引论》（Introductio in analysin infinitorum）中引入了 \(\sin x\)、\(\cos x\) 等记号，并建立了三角函数与指数函数之间的深刻联系------Euler 公式 \(e^{ix} = \cos x + i\sin x\)。这一公式不仅统一了三角函数和指数函数，还揭示了复数的极坐标表示和单位根的结构。三角函数在数学分析中具有特殊地位，因为它们构成了 Fourier 级数的基函数，可以将任意周期函数分解为三角函数的叠加。在更抽象的层面上，\(\sin\) 和 \(\cos\) 是二阶线性微分方程 \(y'' + y = 0\) 的两个线性无关解，这一定义在泛函分析中有深刻的推广。

\subsection{1.8 几何符号}\label{ux51e0ux4f55ux7b26ux53f7}

\textbf{定义 1.8.1}：以下符号用于几何对象。

\begin{longtable}[]{@{}cc@{}}
\toprule
符号 & 含义 \\
\midrule
\endhead
\bottomrule
\endlastfoot
\(\overline{AB}\) 或 \(AB\) & 线段 \\
\(\overrightarrow{AB}\) & 向量 \\
\(\|AB\|\) 或 \(AB\) & 线段 \(AB\) 的长度 \\
\(\angle ABC\) & 以 \(B\) 为顶点的角 \\
\(\triangle ABC\) & 三角形 \\
\(\parallel\) & 平行 \\
\(\perp\) & 垂直 \\
\(\odot O\) & 以 \(O\) 为圆心的圆 \\
\([ABCD]\) & 四边形 \(ABCD\) 的面积 \\
\end{longtable}

\textbf{注记 1.8.1}：在欧几里得几何中，我们默认使用标准的公理体系（如 Hilbert 公理系统）。角度通常用弧度制表示（\(\pi\) 弧度 \(= 180°\)），这在微积分中更为自然。

几何符号体系是欧几里得几何学的语言。线段 \(\overline{AB}\) 表示连接点 \(A\) 和 \(B\) 的直线段，它是欧几里得几何中''两点之间直线最短''这一基本事实的具体体现。向量 \(\overrightarrow{AB}\) 不仅包含长度信息，还包含方向信息，它是从几何到代数的桥梁------平面上的向量加法对应于平行四边形法则，标量乘法对应于缩放。平行符号 \(\parallel\) 和垂直符号 \(\perp\) 是几何中最基本的两种位置关系，它们分别对应了斜率为 0 和斜率乘积为 -1（在坐标几何中）。在 Ch 8（几何基础）中，我们将从 Hilbert 公理出发，严格重建平面几何。在 Ch 20（向量与空间）中，几何符号将自然地推广到三维空间和更高维的向量空间。

\subsection{1.9 希腊字母表}\label{ux5e0cux814aux5b57ux6bcdux8868}

数学中大量使用希腊字母作为变量和常量名。以下列出常用的希腊字母及其常见用法：

\begin{longtable}[]{@{}ccc@{}}
\toprule
大写 & 小写 & 常见用法 \\
\midrule
\endhead
\bottomrule
\endlastfoot
\(\Delta\) & \(\delta\) & 判别式 \(\Delta = b^2 - 4ac\)；变化量 \\
\(\Sigma\) & \(\sigma\) & 求和；标准差 \\
\(\Pi\) & \(\pi\) & 圆周率 \(\pi \approx 3.14159\) \\
\(\Omega\) & \(\omega\) & 角速度 \\
\(\Phi\) & \(\phi\) & 黄金分割比 \(\phi = \frac{1+\sqrt{5}}{2}\) \\
\(\Lambda\) & \(\lambda\) & 特征值 \\
\(M\) & \(\mu\) & 期望、均值 \\
\(E\) & \(\epsilon\) & 任意小的正数（极限语言） \\
\(\Theta\) & \(\theta\) & 角度 \\
\(\Gamma\) & \(\gamma\) & 欧拉常数 \\
\(A\) & \(\alpha\) & 常用角度 \\
\(B\) & \(\beta\) & 常用角度 \\
\(N\) & \(\nu\) & 频率 \\
\(P\) & \(\rho\) & 密度、相关系数 \\
\end{longtable}

\textbf{注记 1.9.1}：大写希腊字母 \(\Sigma\) 常用于求和符号 \(\sum\)，\(\Pi\) 常用于求积符号 \(\prod\)，\(\Delta\) 常用于变化量或判别式。小写 \(\epsilon\) 在极限理论中表示''任意小的正数''，\(\delta\) 表示''\(\epsilon\)-\(\delta\) 论证''中的距离参数。

\hrulefill

\subsection{1.10 证明中的常用符号}\label{ux8bc1ux660eux4e2dux7684ux5e38ux7528ux7b26ux53f7}

\textbf{定义 1.10.1}：以下符号在证明中频繁使用。

\begin{longtable}[]{@{}cc@{}}
\toprule
符号 & 含义 \\
\midrule
\endhead
\bottomrule
\endlastfoot
\(\blacksquare\) & 证毕标记，表示证明结束 \\
\(\Rightarrow\) & 推出 \\
\(\Leftarrow\) & 被推出 \\
\(\Leftrightarrow\) & 等价于 \\
\(\stackrel{\text{def}}{=}\) & 按定义等于 \\
\(a \mid b\) & \(a\) 整除 \(b\) \\
\(a \nmid b\) & \(a\) 不整除 \(b\) \\
\(a \equiv b \pmod{m}\) & \(a\) 与 \(b\) 关于模 \(m\) 同余 \\
\(P \text{ w.r.t. } Q\) & \(P\) 关于 \(Q\)（with respect to） \\
\(f = O(g)\) & \(f\) 是 \(g\) 的大 \(O\)（渐近上界） \\
\(\mathbb{E}[X]\) & 随机变量 \(X\) 的期望 \\
\(\mathrm{Var}(X)\) & 随机变量 \(X\) 的方差 \\
\end{longtable}

\textbf{注记 1.10.1}：在证明中，\(\blacksquare\) 表示一个证明的结束。如果证明中包含多个部分，每个部分结束时使用 \(\blacksquare\)。符号 \(\Rightarrow\) 和 \(\Leftarrow\) 用于逻辑推理步骤之间，而不是命题公式内部（命题公式内部使用 \(\rightarrow\) 和 \(\leftarrow\)）。\(\stackrel{\text{def}}{=}\) 强调等式是根据定义得到的。

\hrulefill

\subsection{本章展望}\label{ux672cux7ae0ux5c55ux671b-1}

本章建立了全书的符号语言体系。在 Ch 2 中，我们将使用这些逻辑符号来系统学习命题逻辑和谓词逻辑，建立严格推理的基础。在 Ch 3 中，集合符号将被用来精确描述数学对象。后续每一章都将依赖本章所建立的符号约定。

\hrulefill

\hrulefill

\hrulefill

\hrulefill

\hrulefill

\newpage

\chapter{01-基础/03-数论基础.md}\label{ux57faux784003-ux6570ux8bbaux57faux7840.md}

\textbf{前置知识}：依赖 Ch 4（整数），特别是整数 \(\mathbb{Z}\) 的构造（4.2 节）和整除性的概念。数论研究整数的深层结构。

在 Ch 4 中我们从皮亚诺公理出发构造了完整的数系。现在我们回到整数 \(\mathbb{Z}\)，深入探索其内部的算术结构。数论是数学中最古老的分支之一，它研究整数的整除性、素数分布和同余关系。尽管问题的表述简单，但其证明往往需要深刻的洞察力。

\hrulefill

\section{第19章：整除}\label{ux7b2c19ux7ae0ux6574ux9664}

整除是数论中最基本的概念，它定义了整数之间的''可除尽''关系。本节从整除的定义出发，建立整除的基本性质（传递性、线性组合、绝对值的界），并证明整除理论的基石------带余除法（欧几里得除法）。带余除法保证了任意两个整数之间总可以进行''商+余数''的分解，且余数严格小于除数（在绝对值意义下），这一性质是辗转相除法和整个初等数论算法体系的基础。整除关系在 \(\mathbb{Z}\) 上构成一个偏序（满足自反性、反对称性和传递性），但不同于通常的序关系 \(\leq\)，整除不是全序------例如 \(2\) 和 \(3\) 互不整除。整除概念的核心在于它揭示了整数的乘法结构：每个整数 \(n\) 的因数集合是有限的（除 \(0\) 外），这为递归和归纳论证提供了基础。本节还将讨论整除与符号的关系以及整除在多项式环中的推广。

\subsubsection{5.1.1 整除的定义}\label{ux6574ux9664ux7684ux5b9aux4e49}

\textbf{定义 5.1.1（整除）}：设 \(a, b\) 是整数，\(b \neq 0\)。如果存在整数 \(k\) 使得 \(a = bk\)，则称 \(b\) \textbf{整除} \(a\)，记为 \(b \mid a\)。当 \(b \mid a\) 时，\(b\) 称为 \(a\) 的\textbf{因数}，\(a\) 称为 \(b\) 的\textbf{倍数}。

整除是数论中最基本的关系，它定义了整数之间的''可除尽''关系。整除关系在 \(\mathbb{Z}\) 上构成一个偏序（满足自反性 \(a \mid a\)、反对称性 \(a \mid b \wedge b \mid a \Rightarrow |a| = |b|\) 和传递性 \(a \mid b \wedge b \mid c \Rightarrow a \mid c\)），但不是全序（例如 \(2 \nmid 3\) 且 \(3 \nmid 2\)）。整除概念的核心在于它揭示了整数的乘法结构：每个整数 \(n\) 的因数集合 \(\{d \in \mathbb{Z} : d \mid n\}\) 是有限的（除 \(0\) 外），这为递归和归纳论证提供了基础。带余除法（定理 5.1.6）是整除理论的基石，它保证了任意两个整数之间总可以进行''商+余数''的分解，且余数严格小于除数。这一性质在多项式环和欧几里得整环中都有直接的推广。

从更抽象的视角看，整除关系定义了 \(\mathbb{Z}\) 上的一个''整除偏序''，使得 \(\mathbb{Z}\) 成为一个格（lattice）：\(\gcd(a, b)\) 是 \(a\) 和 \(b\) 在该偏序下的最大下界，\(\operatorname{lcm}(a, b)\) 是最小上界。整除概念的符号约定（\(b \mid a\) 而非 \(a \mid b\)）可能让初学者感到困惑：\(b \mid a\) 读作''\(b\) 整除 \(a\)``，意味着 \(b\) 是除数（divisor），\(a\) 是被除数（dividend）。这一记法源于数论中将关注点放在''除数''上的传统------当我们说''\(3\) 整除 \(12\)``时，我们关注的是 \(3\) 作为 \(12\) 的因数的性质。整除关系还有一个重要的边界情况：\(a \mid 0\) 对所有 \(a \neq 0\) 成立（因为 \(0 = a \cdot 0\)），但 \(0 \mid a\) 只在 \(a = 0\) 时成立。

\subsubsection{5.1.2 整除的基本性质}\label{ux6574ux9664ux7684ux57faux672cux6027ux8d28}

\textbf{定理 5.1.1（传递性）}：若 \(a \mid b\) 且 \(b \mid c\)，则 \(a \mid c\)。

\textbf{证明}：存在 \(k_1, k_2\) 使 \(b = ak_1\)，\(c = bk_2\)。则 \(c = a(k_1 k_2)\)，即 \(a \mid c\)。\(\blacksquare\)

\textbf{定理 5.1.2（线性组合）}：若 \(a \mid b\) 且 \(a \mid c\)，则对任意整数 \(m, n\)，有 \(a \mid (mb + nc)\)。

\textbf{证明}：设 \(b = ak_1\)，\(c = ak_2\)，则 \(mb + nc = a(mk_1 + nk_2)\)。\(\blacksquare\)

\textbf{定理 5.1.3}：若 \(a \mid b\) 且 \(b \neq 0\)，则 \(|a| \leq |b|\)。

\textbf{证明}：\(b = ak\)，\(k \neq 0\)，故 \(|k| \geq 1\)，\(|b| = |a| \cdot |k| \geq |a|\)。\(\blacksquare\)

\textbf{定理 5.1.4}：若 \(a \mid b\) 且 \(b \mid a\)，则 \(|a| = |b|\)。

\textbf{证明}：由定理 5.1.3，\(|a| \leq |b|\) 且 \(|b| \leq |a|\)。\(\blacksquare\)

\textbf{推论 5.1.1（因数个数的有限性）}：每个非零整数只有有限多个因数。

\textbf{证明}：设 \(n \neq 0\)。若 \(d \mid n\)，则 \(|d| \leq |n|\)。满足 \(|d| \leq |n|\) 的整数 \(d\) 只有有限多个（\(-|n|, -|n|+1, \ldots, |n|-1, |n|\)，共 \(2|n|+1\) 个）。\(\blacksquare\)

\textbf{定理 5.1.5（整除的符号性质）}：设 \(a, b, c\) 是整数。 1. \(a \mid b \iff -a \mid b \iff a \mid -b \iff -a \mid -b\)。 2. 若 \(a \mid b\) 且 \(a \mid c\)，则 \(a \mid (b + c)\) 且 \(a \mid (b - c)\)。

\textbf{证明}：(1) \(a \mid b\) 意味着 \(b = ak\)，即 \(b = (-a)(-k)\)，故 \(-a \mid b\)。其余类似。(2) 设 \(b = ak_1\)，\(c = ak_2\)，则 \(b + c = a(k_1 + k_2)\)，\(b - c = a(k_1 - k_2)\)，故 \(a \mid (b + c)\) 且 \(a \mid (b - c)\)。\(\blacksquare\)

\textbf{注记}：定理 5.1.5(2) 的逆命题不成立。例如 \(a = 4, b = 2, c = 2\) 时，\(4 \mid (2 + 2)\) 且 \(4 \mid (2 - 2)\)，但 \(4 \nmid 2\)。反方向需要附加条件（如 \(a\) 为奇数时，由 \(a \mid 2b\) 和 \(a \mid 2c\) 结合 \(\gcd(a, 2) = 1\) 可得 \(a \mid b\) 且 \(a \mid c\)）。

\subsubsection{5.1.3 带余除法}\label{ux5e26ux4f59ux9664ux6cd5}

\textbf{定理 5.1.6（带余除法）}：设 \(a, b\) 是整数，\(b > 0\)。则存在唯一的整数 \(q\)（商）和 \(r\)（余数），使得 \[a = bq + r, \quad 0 \leq r < b\]

\textbf{证明}：

\textbf{存在性}：考虑 \(S^+ = \{a - bx \mid x \in \mathbb{Z},\, a - bx \geq 0\}\)。取 \(x = -|a|\)，则 \(a - bx = a + b|a| \geq 0\)，故 \(S^+\) 非空。由良序原理（定理 4.1.11），\(S^+\) 有最小元素 \(r = a - bq\)，\(r \geq 0\)。若 \(r \geq b\)，则 \(r - b = a - b(q+1) \in S^+\) 且 \(r - b < r\)，与 \(r\) 的最小性矛盾。故 \(r < b\)。

\textbf{唯一性}：设 \(a = bq_1 + r_1 = bq_2 + r_2\)，\(0 \leq r_1, r_2 < b\)。则 \(b(q_1 - q_2) = r_2 - r_1\)，\(|r_2 - r_1| < b\)。由 \(b \mid (r_2 - r_1)\) 和 \(|r_2 - r_1| < b\)，只能 \(r_1 = r_2\)，从而 \(q_1 = q_2\)。\(\blacksquare\)

\textbf{推论 5.1.2}：任何整数 \(a\) 除以 \(2\) 的余数只有 \(0\) 或 \(1\)，即 \(a\) 要么是偶数（\(a = 2q\)），要么是奇数（\(a = 2q + 1\)）。

\textbf{证明}：在带余除法中取 \(b = 2\) 即得。\(\blacksquare\)

\textbf{推论 5.1.3}：设 \(a, b\) 是整数，\(b \neq 0\)。则 \(b \mid a\) 当且仅当 \(a\) 除以 \(b\) 的余数为 \(0\)。

\textbf{证明}：\(b \mid a\) 意味着 \(a = bk\) 对某个整数 \(k\)，即 \(a = bk + 0\)，余数为 \(0\)。反之，若 \(a = bq + 0\)，则 \(a = bq\)，即 \(b \mid a\)。\(\blacksquare\)

\hrulefill

\section{第20章：最大公约数与最小公倍数}\label{ux7b2c20ux7ae0ux6700ux5927ux516cux7ea6ux6570ux4e0eux6700ux5c0fux516cux500dux6570}

最大公约数（GCD）和最小公倍数（LCM）是整除理论的两个核心概念，它们分别刻画了两个（或多个）整数在''整除偏序''下的最大下界和最小上界。本节从公约数和公倍数的定义出发，引入最大公约数 \(\gcd(a, b)\) 和最小公倍数 \(\operatorname{lcm}(a, b)\) 的概念，并建立它们之间的对偶关系：\(\gcd(a, b) \cdot \operatorname{lcm}(a, b) = |ab|\)。本节的核心是辗转相除法（欧几里得算法）------一种高效计算最大公约数的算法，其正确性依赖于关键引理 \(\gcd(a, b) = \gcd(b, r)\)（其中 \(a = bq + r\)）。辗转相除法不仅是计算 GCD 的最古老算法（记载于欧几里得的《几何原本》），也是现代计算机代数系统的基础。裴蜀定理（Bezout 定理）是本节另一个重要结果：\(\gcd(a, b)\) 可以表示为 \(a\) 和 \(b\) 的整系数线性组合。这一结果在数论中具有根本性的重要地位，它是欧几里得引理、中国剩余定理和 RSA 加密算法的理论基础。

\subsubsection{5.2.1 最大公约数}\label{ux6700ux5927ux516cux7ea6ux6570}

\textbf{定义 5.2.1（公约数与最大公约数）}：设 \(a, b\) 是不全为零的整数。若 \(d \mid a\) 且 \(d \mid b\)，则 \(d\) 为 \(a, b\) 的\textbf{公约数}。所有公约数中最大的正整数称为\textbf{最大公约数}，记为 \(\gcd(a, b)\)。

\textbf{注记 5.2.1}：\(\gcd(0, 0)\) 没有定义（因为每个正整数都是 \(0\) 的因数，没有最大值）。但通常定义 \(\gcd(0, n) = |n|\)（\(n \neq 0\)）。特别地，\(\gcd(0, n) = |n|\)，\(\gcd(n, n) = |n|\)。

\textbf{定理 5.2.1（GCD的基本性质）}：设 \(a, b\) 是不全为零的整数，\(d = \gcd(a, b)\)。则： 1. \(d \mid a\) 且 \(d \mid b\)（\(d\) 是公约数） 2. 若 \(c \mid a\) 且 \(c \mid b\)，则 \(c \mid d\)（\(d\) 是最大的公约数） 3. \(\gcd(a, b) = \gcd(|a|, |b|)\) 4. \(\gcd(a, b) = \gcd(b, a)\) 5. \(\gcd(a, 0) = |a|\)（\(a \neq 0\)）

\textbf{证明}：(1) 由定义。(2) 由裴蜀定理（定理 5.2.3），\(d = ax + by\)，故 \(c \mid d\)。(3) 由整除的性质 \(b \mid a \iff |b| \mid |a|\)。(4) 由交换律。(5) \(0\) 的因数是所有整数，\(a\) 的因数不超过 \(|a|\)，故 \(\gcd(a, 0) = |a|\)。\(\blacksquare\)

\textbf{定理 5.2.2（辗转相除法/欧几里得算法）}：设 \(a, b\) 是正整数，\(a \geq b\)。反复做带余除法：

\[a = bq_1 + r_1,\; b = r_1 q_2 + r_2,\; \ldots,\; r_{k-2} = r_{k-1}q_k + r_k,\; r_{k-1} = r_k q_{k+1} + 0\]

则 \(\gcd(a, b) = r_k\)（最后一个非零余数）。

\textbf{证明}：

\textbf{关键引理}：\(\gcd(a, b) = \gcd(b, r)\)，其中 \(a = bq + r\)。

\textbf{引理的证明}：设 \(d = \gcd(a, b)\)。由 \(d \mid a\) 且 \(d \mid b\) 得 \(d \mid (a - bq) = r\)，故 \(d\) 是 \(b, r\) 的公约数，\(d \leq \gcd(b, r)\)。反之，设 \(d' = \gcd(b, r)\)，由 \(d' \mid b\) 且 \(d' \mid r\) 得 \(d' \mid (bq + r) = a\)，故 \(d' \leq \gcd(a, b) = d\)。综合得 \(d = d'\)。\(\blacksquare\)

对每一步应用引理：\(\gcd(a, b) = \gcd(b, r_1) = \gcd(r_1, r_2) = \cdots = \gcd(r_{k-1}, r_k) = \gcd(r_k, 0) = r_k\)。\(\blacksquare\)

\textbf{定理 5.2.3（裴蜀定理）}：设 \(a, b\) 是不全为零的整数。则存在整数 \(x, y\) 使得 \[ax + by = \gcd(a, b)\]

更一般地，\(ax + by = c\) 有整数解当且仅当 \(\gcd(a, b) \mid c\)。

\textbf{证明}：

\textbf{第一步}：考虑 \(S = \{ax + by \mid x, y \in \mathbb{Z},\, ax + by > 0\}\)。\(S\) 非空，由良序原理有最小元素 \(d = ax_0 + by_0\)。

\textbf{第二步}：对 \(a\) 做带余除法 \(a = dq + r\)，\(0 \leq r < d\)。则 \(r = a - dq = a(1 - qx_0) + b(-qy_0)\)。若 \(r > 0\)，则 \(r \in S\) 且 \(r < d\)，与 \(d\) 的最小性矛盾。故 \(r = 0\)，即 \(d \mid a\)。同理 \(d \mid b\)。

\textbf{第三步}：\(d\) 是 \(a, b\) 的公约数。设 \(c\) 是任意公约数，则 \(c \mid (ax_0 + by_0) = d\)，故 \(c \leq d\)。因此 \(d = \gcd(a, b)\)。

\textbf{第四步}：若 \(\gcd(a, b) \mid c\)，设 \(c = k \cdot \gcd(a, b)\)，则 \(x = kx_0\)，\(y = ky_0\) 是解。反之，若 \(ax + by = c\) 有解，由 \(\gcd(a, b) \mid a\) 和 \(\gcd(a, b) \mid b\) 得 \(\gcd(a, b) \mid c\)。\(\blacksquare\)

\textbf{推论 5.2.4（裴蜀定理的等价形式）}：\(\gcd(a, b)\) 是 \(a\) 和 \(b\) 的线性组合中最小的正整数。

\textbf{证明}：设 \(d = \gcd(a, b)\)。由裴蜀定理，存在整数 \(x_0, y_0\) 使得 \(ax_0 + by_0 = d\)，故 \(d\) 是 \(a\) 和 \(b\) 的一个线性组合。

设 \(s\) 是 \(a\) 和 \(b\) 的任意正线性组合，即 \(s = am + bn\)（\(m, n \in \mathbb{Z}\)，\(s > 0\)）。由 \(d \mid a\) 和 \(d \mid b\)，得 \(d \mid s\)，故 \(s \geq d\)。

因此 \(d\) 是最小的正线性组合。\(\blacksquare\)

\textbf{推论 5.2.5}：\(\gcd(a, b) = 1\) 当且仅当存在整数 \(x, y\) 使得 \(ax + by = 1\)。

\textbf{证明}：\((\Rightarrow)\)：由裴蜀定理。( \(\Leftarrow)\)：若 \(ax + by = 1\)，则 \(\gcd(a, b) \mid 1\)，故 \(\gcd(a, b) = 1\)。\(\blacksquare\)

\textbf{定理 5.2.6（GCD的运算性质）}： 1. \(\gcd(a, b) = \gcd(a, b + ka)\) 对任意整数 \(k\)。 2. \(\gcd(a, b) = \gcd(a, b \bmod a)\)（\(a > 0\)）。 3. \(\gcd(a, b, c) = \gcd(\gcd(a, b), c)\)。

\textbf{证明}：(1) 设 \(d = \gcd(a, b)\)。\(d \mid a\) 且 \(d \mid b\)，故 \(d \mid (b + ka)\)。因此 \(d\) 是 \(a, b + ka\) 的公约数，\(d \leq \gcd(a, b + ka)\)。反之设 \(d' = \gcd(a, b + ka)\)。\(d' \mid a\) 且 \(d' \mid (b + ka)\)，故 \(d' \mid [(b + ka) - ka] = b\)。因此 \(d'\) 是 \(a, b\) 的公约数，\(d' \leq d\)。故 \(d = d'\)。

\begin{enumerate}
\def\labelenumi{(\arabic{enumi})}
\setcounter{enumi}{1}
\item
  由 (1) 取 \(k = -\lfloor b/a \rfloor\)。
\item
  \(\gcd(a, b, c)\) 是同时整除 \(a, b, c\) 的最大正整数。\(\gcd(\gcd(a, b), c)\) 是同时整除 \(\gcd(a, b)\) 和 \(c\) 的最大正整数。由于 \(d \mid \gcd(a, b) \iff (d \mid a \wedge d \mid b)\)，两者等价。\(\blacksquare\)
\end{enumerate}

\subsubsection{5.2.2 最小公倍数}\label{ux6700ux5c0fux516cux500dux6570}

\textbf{定义 5.2.2（最小公倍数）}：设 \(a, b\) 是非零整数。\(a, b\) 的所有正公倍数中最小的称为\textbf{最小公倍数}，记为 \(\operatorname{lcm}(a, b)\)。

\textbf{注记 5.2.2}：\(\operatorname{lcm}(a, b)\) 的存在性由以下事实保证：\(a, b\) 的公倍数集合 \(M = \{m > 0 : a \mid m \text{ 且 } b \mid m\}\) 是非空的正整数集合（\(|ab| \in M\)），由良序原理有最小元素。

\textbf{注记 5.2.3}：GCD 和 LCM 可以推广到多个整数：\(\gcd(a_1, \ldots, a_n) = \gcd(\gcd(a_1, \ldots, a_{n-1}), a_n)\)，\(\operatorname{lcm}(a_1, \ldots, a_n) = \operatorname{lcm}(\operatorname{lcm}(a_1, \ldots, a_{n-1}), a_n)\)。

最小公倍数（LCM）与最大公约数（GCD）是对偶的概念。在整除的偏序结构下，\(\gcd(a, b)\) 是 \(a\) 和 \(b\) 的最大下界（meet），\(\operatorname{lcm}(a, b)\) 是 \(a\) 和 \(b\) 的最小上界（join），因此 \((\mathbb{N}^+, \mid)\) 构成一个格（lattice）。定理 5.2.7 揭示了 GCD 与 LCM 之间的深刻关系：\(\operatorname{lcm}(a, b) \cdot \gcd(a, b) = |ab|\)。这一关系在素因子分解的视角下变得显而易见的------\(\gcd\) 取每个素因子的最小指数，\(\operatorname{lcm}\) 取每个素因子的最大指数，而两者之和等于 \(a\) 和 \(b\) 中该素因子的指数之和。这一对偶性在抽象代数中推广为 Dedekind 整环中理想的交与和的对应关系。

\textbf{定理 5.2.7}：\(\operatorname{lcm}(a, b) = \frac{|ab|}{\gcd(a, b)}\)。

\textbf{证明}：设 \(d = \gcd(a, b)\)，\(a = da'\)，\(b = db'\)。由裴蜀定理，\(a'x + b'y = 1\)，故 \(\gcd(a', b') = 1\)。

\(m = da'b' = ab/d\) 是公倍数（\(m = a \cdot b' = b \cdot a'\)）。设 \(n\) 是任意正公倍数，\(a \mid n\) 且 \(b \mid n\)。设 \(n = as = bt\)，则 \(a's = b't\)。由 \(\gcd(a', b') = 1\) 和 \(a' \mid b't\) 得 \(a' \mid t\)。设 \(t = a'k\)，则 \(n = db'ka' = mk\)，故 \(m \mid n\)，\(n \geq m\)。\(\blacksquare\)

\textbf{推论 5.2.8（LCM的性质）}： 1. \(\operatorname{lcm}(a, b) = \operatorname{lcm}(|a|, |b|)\) 2. 若 \(a \mid c\) 且 \(b \mid c\)，则 \(\operatorname{lcm}(a, b) \mid c\)。 3. \(\operatorname{lcm}(a, b) \cdot \gcd(a, b) = |ab|\)。

\textbf{证明}：(1) 由 \(|a| \mid c \iff a \mid c\)。(2) 由定理 5.2.7 的证明。(3) \(\operatorname{lcm}(a, b) = |ab|/\gcd(a, b)\)。\(\blacksquare\)

\textbf{定理 5.2.9（互素的性质）}：设 \(\gcd(a, b) = 1\)。 1. 若 \(a \mid c\) 且 \(b \mid c\)，则 \(ab \mid c\)。 2. 若 \(a \mid bc\)，则 \(a \mid c\)。 3. \(\gcd(a^n, b^n) = 1\) 对任意正整数 \(n\)。

\textbf{证明}：(1) 由 \(\operatorname{lcm}(a, b) = ab\)（当 \(\gcd(a, b) = 1\)）和 \(\operatorname{lcm}(a, b) \mid c\)。(2) 由裴蜀定理，\(ax + by = 1\)，乘以 \(c\) 得 \(acx + bcy = c\)。\(a \mid acx\) 且 \(a \mid bcy\)（因 \(a \mid bc\)），故 \(a \mid c\)。(3) 若素数 \(p \mid \gcd(a^n, b^n)\)，则 \(p \mid a^n\) 且 \(p \mid b^n\)。由欧几里得引理，\(p \mid a\) 且 \(p \mid b\)，故 \(p \mid \gcd(a, b) = 1\)，矛盾。\(\blacksquare\)

\hrulefill

\section{第21章：素数}\label{ux7b2c21ux7ae0ux7d20ux6570}

素数是数论中最核心的研究对象，被称为''整数的原子''。本节从素数的定义出发，建立素数理论的基本结果：欧几里得引理（若 \(p\) 是素数且 \(p \mid ab\)，则 \(p \mid a\) 或 \(p \mid b\)）和算术基本定理（每个大于 \(1\) 的正整数都可以唯一地分解为素数的乘积）。算术基本定理是数论的''基本定理''，它确立了素数作为整数的''不可约构建块''的地位------每个正整数都可以以一种本质上唯一的方式分解为素数的乘积，恰如每个分子都可以分解为原子。这一事实使得许多数论问题可以归结为对素数性质的研究。本节还将证明素数有无穷多个（欧几里得的经典证明），并讨论素数分布的基本问题（素数定理的陈述，但不证明）。素数理论在密码学中具有重要应用------RSA 加密算法的安全性依赖于大数分解的困难性------这是现代数论与计算机科学深度交叉的典型范例。

\subsubsection{5.3.1 素数的定义}\label{ux7d20ux6570ux7684ux5b9aux4e49}

\textbf{定义 5.3.1（素数与合数）}：大于 \(1\) 的正整数 \(p\)，如果只有 \(1\) 和 \(p\) 两个正因数，则称 \(p\) 为\textbf{素数}。大于 \(1\) 但不是素数的正整数称为\textbf{合数}。

\textbf{注记}：\(1\) 既不是素数也不是合数。\(2\) 是唯一的偶素数。

\textbf{引理 5.3.1}：每个大于 \(1\) 的正整数都有一个素因数。

\textbf{证明}：设 \(n > 1\)。若 \(n\) 是素数，结论成立。若 \(n\) 是合数，\(n = ab\)，\(1 < a, b < n\)。对 \(a\) 重复此过程。由于正整数不能无限递减，必在有限步内得到素因数。\(\blacksquare\)

\textbf{定理 5.3.1（欧几里得引理）}：若 \(p\) 是素数且 \(p \mid ab\)，则 \(p \mid a\) 或 \(p \mid b\)。

\textbf{证明}：假设 \(p \nmid a\)。因 \(p\) 是素数，\(\gcd(p, a) = 1\)。由裴蜀定理，存在 \(x, y\) 使得 \(px + ay = 1\)。两边乘 \(b\)：\(pbx + aby = b\)。由 \(p \mid pbx\) 和 \(p \mid aby\)（因 \(p \mid ab\)），得 \(p \mid b\)。\(\blacksquare\)

\textbf{定理 5.3.1'（素数的基本性质）}： 1. 若 \(p\) 是素数且 \(p \mid a_1 a_2 \cdots a_k\)，则 \(p\) 至少整除某个 \(a_i\)。 2. 设 \(p\) 是素数。则对任意整数 \(a\)，恰有 \(p \mid a\) 或 \(\gcd(p, a) = 1\) 之一成立。 3. 若 \(p, q\) 都是素数且 \(p \mid q\)，则 \(p = q\)。

\textbf{证明}：(1) 对 \(k\) 用归纳法。\(k = 1\) 时显然。\(k = 2\) 时即欧几里得引理（定理 5.3.1）。设对 \(k-1\) 成立。若 \(p \mid a_1 \cdots a_{k-1}\)，由归纳假设得证。若 \(p \nmid a_1 \cdots a_{k-1}\)，由欧几里得引理，\(p \mid a_k\)。(2) \(p\) 的正因数只有 \(1\) 和 \(p\)。若 \(p \nmid a\)，则 \(\gcd(p, a) \neq p\)，故 \(\gcd(p, a) = 1\)。(3) \(p \mid q\) 且 \(q\) 是素数，故 \(p = 1\) 或 \(p = q\)。但 \(p\) 是素数，\(p \neq 1\)，故 \(p = q\)。\(\blacksquare\)

\subsubsection{5.3.2 算术基本定理}\label{ux7b97ux672fux57faux672cux5b9aux7406}

\textbf{定理 5.3.2（算术基本定理）}：每个大于 \(1\) 的正整数 \(n\) 都可以唯一地（不计因子次序）表示为素数的乘积： \[n = p_1 p_2 \cdots p_k, \quad p_1 \leq p_2 \leq \cdots \leq p_k\]

\textbf{证明}：

\textbf{存在性}：对 \(n\) 用强归纳法。\(n = 2\) 时，\(2\) 本身是素数，分解存在。设对所有 \(2 \leq m < n\)，素数分解存在。考虑 \(n\)：若 \(n\) 是素数，分解为自身。若 \(n\) 是合数，则 \(n = ab\)，\(1 < a, b < n\)。由归纳假设，\(a\) 和 \(b\) 都可以表示为素数乘积：\(a = p_1 \cdots p_s\)，\(b = q_1 \cdots q_t\)。于是 \(n = p_1 \cdots p_s q_1 \cdots q_t\) 也是素数乘积。由强归纳法，存在性得证。

\textbf{唯一性}：设 \(n = p_1 p_2 \cdots p_k = q_1 q_2 \cdots q_l\) 是两种素数分解（\(p_1 \leq \cdots \leq p_k\)，\(q_1 \leq \cdots \leq q_l\)）。

由 \(p_1 \mid n = q_1 q_2 \cdots q_l\)，根据定理 5.3.1'(1)，\(p_1\) 整除某个 \(q_j\)。适当重排使 \(p_1 \mid q_1\)。由于 \(q_1\) 是素数且 \(p_1 > 1\)，由定理 5.3.1'(3) 得 \(p_1 = q_1\)。

约去 \(p_1 = q_1\)，得 \(p_2 \cdots p_k = q_2 \cdots q_l\)。若 \(k = 1\)，则左边为 \(1\)，右边 \(q_2 \cdots q_l = 1\)，只能 \(l = 1\)。若 \(k > 1\)，对 \(k + l\) 用归纳法（或重复上述过程），得 \(k = l\) 且 \(p_i = q_i\)（\(i = 1, 2, \ldots, k\)）。\(\blacksquare\)

\textbf{推论 5.3.1（素因子分解的标准形式）}：每个大于 \(1\) 的正整数 \(n\) 都可以唯一地写成 \[n = p_1^{a_1} p_2^{a_2} \cdots p_r^{a_r}\] 其中 \(p_1 < p_2 < \cdots < p_r\) 是素数，\(a_1, a_2, \ldots, a_r\) 是正整数。

\textbf{证明}：由算术基本定理，将相同的素因子合并即得。\(\blacksquare\)

\textbf{推论 5.3.2（GCD和LCM的素因子表示）}：设 \(a = p_1^{a_1} \cdots p_r^{a_r}\)，\(b = p_1^{b_1} \cdots p_r^{b_r}\)（允许某些指数为 \(0\)）。则 \[\gcd(a, b) = p_1^{\min(a_1, b_1)} \cdots p_r^{\min(a_r, b_r)}\] \[\operatorname{lcm}(a, b) = p_1^{\max(a_1, b_1)} \cdots p_r^{\max(a_r, b_r)}\]

\textbf{证明}：设 \(d = p_1^{d_1} \cdots p_r^{d_r}\)。\(d \mid a\) 当且仅当 \(d_i \leq a_i\) 对所有 \(i\)。\(d \mid b\) 当且仅当 \(d_i \leq b_i\) 对所有 \(i\)。\(d\) 是公约数当且仅当 \(d_i \leq \min(a_i, b_i)\) 对所有 \(i\)。最大公约数取等号。\(\operatorname{lcm}\) 的证明类似。\(\blacksquare\)

\textbf{推论 5.3.3}：\(\gcd(a, b) \cdot \operatorname{lcm}(a, b) = ab\)。

\textbf{证明}：\(\min(a_i, b_i) + \max(a_i, b_i) = a_i + b_i\)。\(\blacksquare\)

\subsubsection{5.3.3 素数的无穷性}\label{ux7d20ux6570ux7684ux65e0ux7a77ux6027}

\textbf{定理 5.3.3（欧几里得定理）}：素数有无穷多个。

\textbf{证明}：假设素数只有有限个 \(p_1, p_2, \ldots, p_k\)。考虑 \(N = p_1 p_2 \cdots p_k + 1\)。\(N > 1\)，有素因数 \(p\)。但 \(p_i\) 除 \(N\) 的余数为 \(1\)（\(p_i \mid p_1 \cdots p_k\)），故 \(p_i \nmid N\)，即 \(p \neq p_i\) 对所有 \(i\)。矛盾。\(\blacksquare\)

\textbf{注记 5.3.1}：欧几里得的原始证明是构造性的，但 \(N = p_1 \cdots p_k + 1\) 本身不一定是素数。例如 \(2 \cdot 3 \cdot 5 \cdot 7 \cdot 11 \cdot 13 + 1 = 30031 = 59 \times 509\)。证明的精髓在于：\(N\) 有不在列表中的素因数。

\textbf{推论 5.3.4（第 \(k\) 个素数的上界）}：设 \(p_k\) 是第 \(k\) 个素数。则 \(p_k \leq 2^{2^{k-1}}\)。

\textbf{证明}：对 \(k\) 用归纳法。\(p_1 = 2 = 2^{2^0}\)。设 \(p_i \leq 2^{2^{i-1}}\)（\(i = 1, \ldots, k\)）。由欧几里得的证明，\(p_{k+1} \leq p_1 p_2 \cdots p_k + 1\)。\(p_1 \cdots p_k \leq 2^{2^0 + 2^1 + \cdots + 2^{k-1}} = 2^{2^k - 1}\)。故 \(p_{k+1} \leq 2^{2^k - 1} + 1 \leq 2^{2^k}\)。\(\blacksquare\)

\textbf{定理 5.3.4（关于素数分布的简单结果）}：对任意正整数 \(n\)，在 \(n\) 和 \(2n\) 之间（含端点）至少存在一个素数。

\textbf{证明}：考虑二项式系数 \(\binom{2n}{n} = \frac{(2n)!}{(n!)^2}\)。

\textbf{第一步}：\(\binom{2n}{n} \geq 2^n\)（\(n \geq 1\)）。由二项式定理，\(4^n = (1+1)^{2n} = \sum_{k=0}^{2n} \binom{2n}{k} \leq (2n+1)\binom{2n}{n}\)（\(\binom{2n}{n}\) 是最大项），故 \(\binom{2n}{n} \geq \frac{4^n}{2n+1}\)。

\textbf{第二步}：\(\binom{2n}{n}\) 的素因子都不超过 \(2n\)。因为 \(\binom{2n}{n} = \frac{(2n)!}{(n!)^2}\) 的分子的素因子都不超过 \(2n\)。

\textbf{第三步}：若 \(n < p \leq 2n\)（\(p\) 是素数），则 \(p \mid \binom{2n}{n}\)。因为 \(p\) 整除分子 \((2n)!\) 中的因子（\(p \leq 2n\)），但不整除分母 \((n!)^2\)（\(p > n\)）。

\textbf{第四步}：假设 \((n, 2n]\) 中没有素数。则 \(\binom{2n}{n}\) 的所有素因子都 \(\leq n\)。由算术基本定理，\(\binom{2n}{n} = \prod_{p \leq n} p^{a_p}\)。其中 \(a_p = \sum_{j=1}^{\infty} \left(\lfloor \frac{2n}{p^j} \rfloor - 2\lfloor \frac{n}{p^j} \rfloor\right)\)。每一项 \(\leq 1\)（因为 \(\lfloor 2x \rfloor - 2\lfloor x \rfloor \in \{0, 1\}\)），且只有 \(j \leq \log_p(2n)\) 项非零。故 \(a_p \leq \lfloor \log_p(2n) \rfloor\)。因此 \(\binom{2n}{n} \leq \prod_{p \leq n} p^{\lfloor \log_p(2n) \rfloor} \leq (2n)^{\pi(n)}\)（\(\pi(n)\) 是不超过 \(n\) 的素数个数）。

\textbf{第五步}：由第一步，\(2^n \leq \binom{2n}{n} \leq (2n)^{\pi(n)}\)。取对数得 \(\pi(n) \geq \frac{n \ln 2}{\ln(2n)}\)。这足以说明 \(\pi(n)\) 随 \(n\) 增长趋于无穷（即素数有无穷多个，已由定理 5.3.3 独立证明），因此对充分大的 \(n\) 产生矛盾。对于小的 \(n\) 值可以直接验证：\((1, 2]\) 有素数 \(2\)；\((2, 4]\) 有素数 \(3\)；\((3, 6]\) 有素数 \(5\)；\((4, 8]\) 有素数 \(5, 7\)。\(\blacksquare\)

\subsubsection{5.3.4 埃拉托色尼筛法}\label{ux57c3ux62c9ux6258ux8272ux5c3cux7b5bux6cd5}

\textbf{算法 5.3.1（埃拉托色尼筛法）}：求不超过 \(N\) 的所有素数的方法： 1. 列出 \(2, 3, 4, \ldots, N\)。 2. 取列表中最小的数 \(p\)（初始 \(p = 2\)），标记 \(p\) 为素数。 3. 划去 \(p\) 的所有倍数（\(2p, 3p, 4p, \ldots\)）。 4. 令 \(p\) 为列表中下一个未被划去的数，回到步骤 3。 5. 当 \(p > \sqrt{N}\) 时停止。列表中剩余的未被划去的数都是素数。

\textbf{定理 5.3.5（筛法的正确性）}：算法结束后，不超过 \(N\) 的未被划去的数恰好是不超过 \(N\) 的素数。

\textbf{证明}：若 \(n \leq N\) 是合数，则 \(n = ab\)（\(1 < a \leq b < n\)），\(a \leq \sqrt{n} \leq \sqrt{N}\)。\(a\) 有一个素因子 \(p \leq a \leq \sqrt{N}\)。当处理到 \(p\) 时，\(n\) 作为 \(p\) 的倍数被划去。因此合数都会被划去。素数不会被划去：素数 \(p\) 只在处理到 \(p\) 自身时被标记为素数，而非被划去（\(p\) 的倍数 \(2p, 3p, \ldots\) 都大于 \(p\)）。\(\blacksquare\)

\textbf{注记 5.3.2（筛法的效率）}：埃拉托色尼筛法的时间复杂度为 \(O(N \log \log N)\)，空间复杂度为 \(O(N)\)。对于 \(N\) 很大的情况，可以使用更高级的筛法（如欧拉筛法达到 \(O(N)\) 时间复杂度）。

\subsubsection{5.3.5 素数分布的初步结果}\label{ux7d20ux6570ux5206ux5e03ux7684ux521dux6b65ux7ed3ux679c}

\textbf{定义 5.3.2（素数计数函数）}：\(\pi(x)\) 表示不超过 \(x\) 的素数个数。

\textbf{定理 5.3.6（素数定理的弱形式）}：\(\pi(x) \geq \frac{x}{2\ln x}\) 对充分大的 \(x\) 成立。

\textbf{注记 5.3.3（素数定理）}：素数定理断言 \(\pi(x) \sim \frac{x}{\ln x}\)（当 \(x \to \infty\) 时），即 \(\lim_{x \to \infty} \frac{\pi(x)}{x/\ln x} = 1\)。这个深刻结果由 Hadamard 和 de la Vallee Poussin 在 1896 年独立证明，证明使用了复分析工具。一个初等证明由 Erdos 和 Selberg 在 1949 年给出。

\hrulefill

\section{第22章：同余}\label{ux7b2c22ux7ae0ux540cux4f59}

同余理论是数论中最重要的工具之一，由 Gauss 于 1801 年在《算术研究》（Disquisitiones Arithmeticae）中系统建立。同余关系 \(a \equiv b \pmod{m}\) 是 \(\mathbb{Z}\) 上的一个等价关系，它将整数集划分为 \(m\) 个剩余类，每个剩余类中的所有整数在模 \(m\) 意义下是''不可区分''的。同余理论的核心洞见是：模 \(m\) 的剩余类构成一个环 \(\mathbb{Z}/m\mathbb{Z}\)（模 \(m\) 的剩余类环），其中乘法可逆的元素恰好是那些与 \(m\) 互素的剩余类，形成乘法群 \((\mathbb{Z}/m\mathbb{Z})^\times\)。本节将从同余的定义出发，建立同余的运算性质、同余方程的一次解法、中国剩余定理（CRT），以及 Euler 定理和 Fermat 小定理。同余理论在密码学（RSA）、计算机科学（哈希函数）、编码理论（循环码）和组合数学中都有广泛应用。

\subsubsection{5.4.1 同余的定义}\label{ux540cux4f59ux7684ux5b9aux4e49}

\textbf{定义 5.4.1（同余）}：设 \(m\) 是正整数。若 \(m \mid (a - b)\)，则称 \(a\) 与 \(b\) \textbf{关于模 \(m\) 同余}，记为 \(a \equiv b \pmod{m}\)。

同余关系是数论中最核心的概念之一，由 Gauss 于 1801 年在《算术研究》（Disquisitiones Arithmeticae）中系统引入。同余关系本质上是一种''模 \(m\) 的等价关系''：它将整数集 \(\mathbb{Z}\) 划分为 \(m\) 个等价类（剩余类），每个剩余类中的所有整数在模 \(m\) 意义下是''不可区分''的。同余记号 \(\equiv\) 的引入极大地简化了数论中的许多论证------例如，判断一个数是否能被 3 整除，只需看其各位数字之和模 3 是否为 0。同余理论的核心洞见是：模 \(m\) 的剩余类构成一个环 \(\mathbb{Z}/m\mathbb{Z}\)（称为模 \(m\) 的剩余类环），其中乘法可逆的元素恰好是那些与 \(m\) 互素的剩余类，其乘法群 \((\mathbb{Z}/m\mathbb{Z})^\times\) 的阶恰好是 \(\varphi(m)\)。这一群论视角是欧拉定理和费马小定理的现代解释。

\textbf{定理 5.4.1}：\(a \equiv b \pmod{m}\) 当且仅当 \(a\) 和 \(b\) 除以 \(m\) 的余数相同。

\textbf{证明}：设 \(a = mq_1 + r_1\)，\(b = mq_2 + r_2\)，\(0 \leq r_1, r_2 < m\)。则 \(m \mid (a - b) = m(q_1 - q_2) + (r_1 - r_2)\) 当且仅当 \(m \mid (r_1 - r_2)\)。由 \(|r_1 - r_2| < m\)，这等价于 \(r_1 = r_2\)。\(\blacksquare\)

\textbf{定理 5.4.2}：同余关系是等价关系（自反、对称、传递）。

\textbf{证明}：自反性：\(m \mid (a - a) = 0\)。对称性：若 \(m \mid (a - b)\) 则 \(m \mid (b - a)\)。传递性：若 \(m \mid (a - b)\) 且 \(m \mid (b - c)\)，则 \(m \mid [(a - b) + (b - c)] = (a - c)\)。\(\blacksquare\)

\textbf{推论 5.4.1}：整数集 \(\mathbb{Z}\) 关于模 \(m\) 的同余关系被分为 \(m\) 个等价类（剩余类）：\([0], [1], \ldots, [m-1]\)，其中 \([r] = \{r + km \mid k \in \mathbb{Z}\}\)。

\textbf{证明}：每个整数 \(a\) 除以 \(m\) 的余数 \(r\) 唯一确定（\(0 \leq r < m\)），且 \(a \equiv r \pmod{m}\)。因此每个整数恰好属于一个剩余类。\(\blacksquare\)

\textbf{定理 5.4.2'（同余的除法规则）}：若 \(ac \equiv bc \pmod{m}\) 且 \(\gcd(c, m) = 1\)，则 \(a \equiv b \pmod{m}\)。

\textbf{证明}：由 \(m \mid c(a - b)\) 和 \(\gcd(c, m) = 1\)。由裴蜀定理，存在 \(x, y\) 使得 \(cx + my = 1\)。两边乘 \((a - b)\)：\(c(a - b)x + m(a - b)y = a - b\)。由 \(m \mid c(a - b)\) 得 \(m \mid [c(a - b)x + m(a - b)y] = a - b\)，即 \(a \equiv b \pmod{m}\)。\(\blacksquare\)

\subsubsection{5.4.2 同余的运算}\label{ux540cux4f59ux7684ux8fd0ux7b97}

\textbf{定理 5.4.3（同余的运算性质）}：若 \(a \equiv b \pmod{m}\)，\(c \equiv d \pmod{m}\)，则： 1. \(a + c \equiv b + d \pmod{m}\) 2. \(a - c \equiv b - d \pmod{m}\) 3. \(ac \equiv bd \pmod{m}\) 4. \(a^n \equiv b^n \pmod{m}\)（\(n\) 为正整数）

\textbf{证明}：

\begin{enumerate}
\def\labelenumi{(\arabic{enumi})}
\item
  \(m \mid [(a - b) + (c - d)] = [(a + c) - (b + d)]\)。
\item
  \(ac - bd = c(a - b) + b(c - d)\)。由 \(m \mid (a - b)\) 和 \(m \mid (c - d)\) 得 \(m \mid (ac - bd)\)。
\item
  由 (3) 反复应用或用数学归纳法。\(\blacksquare\)
\end{enumerate}

\subsubsection{5.4.3 费马小定理}\label{ux8d39ux9a6cux5c0fux5b9aux7406}

\textbf{定理 5.4.4（费马小定理）}：设 \(p\) 是素数，\(a\) 是整数且 \(p \nmid a\)。则 \[a^{p-1} \equiv 1 \pmod{p}\]

\textbf{证明}：

\textbf{引理}：设 \(p\) 是素数，\(1 \leq k \leq p - 1\)。则 \(p \mid \binom{p}{k}\)。

\textbf{引理的证明}：\(\binom{p}{k} = \frac{p!}{k!(p-k)!}\)，故 \(p \cdot k! \cdot (p-k)! = p!\)。\(p\) 不整除 \(k!\)（所有因子小于 \(p\)），也不整除 \((p-k)!\)，但 \(p \mid p!\)，故 \(p \mid \binom{p}{k}\)。\(\blacksquare\)

回到主定理的证明。由二项式定理： \[(a + 1)^p = \sum_{k=0}^{p} \binom{p}{k} a^k = a^p + \binom{p}{1}a^{p-1} + \cdots + \binom{p}{p-1}a + 1\]

由引理，\(p \mid \binom{p}{k}\)（\(1 \leq k \leq p - 1\)），故 \((a+1)^p \equiv a^p + 1 \pmod{p}\)。

用归纳法证 \(a^p \equiv a \pmod{p}\)：\(a = 1\) 时显然。设 \(a^p \equiv a\)，则 \((a+1)^p \equiv a^p + 1 \equiv a + 1 \pmod{p}\)。

对所有整数 \(a\) 也成立（对 \(p = 2\) 可直接验证，对奇素数 \(p\) 由 \((-a)^p = -a^p\) 可得）。

由 \(a^p \equiv a\)，若 \(p \nmid a\)，由欧几里得引理从 \(p \mid a(a^{p-1} - 1)\) 和 \(p \nmid a\) 得 \(p \mid (a^{p-1} - 1)\)，即 \(a^{p-1} \equiv 1 \pmod{p}\)。\(\blacksquare\)

\subsubsection{5.4.4 中国剩余定理}\label{ux4e2dux56fdux5269ux4f59ux5b9aux7406}

\textbf{定理 5.4.5（中国剩余定理）}：设 \(m_1, m_2\) 是互素的正整数（\(\gcd(m_1, m_2) = 1\)）。则对任意整数 \(a_1, a_2\)，同余方程组 \[\begin{cases} x \equiv a_1 \pmod{m_1} \\ x \equiv a_2 \pmod{m_2} \end{cases}\] 在模 \(m_1 m_2\) 的意义下有唯一解。

\textbf{证明}：

\textbf{存在性}：由 \(\gcd(m_1, m_2) = 1\) 和裴蜀定理，存在 \(s, t\) 使得 \(m_1 s + m_2 t = 1\)。令 \(x = a_1 m_2 t + a_2 m_1 s\)。

验证：\(x = a_1 m_2 t + a_2 m_1 s \equiv a_1 m_2 t \pmod{m_1}\)。又 \(m_2 t = 1 - m_1 s \equiv 1 \pmod{m_1}\)，故 \(x \equiv a_1 \pmod{m_1}\)。类似地 \(x \equiv a_2 \pmod{m_2}\)。

\textbf{唯一性}：设 \(x_1, x_2\) 都是解。则 \(m_1 \mid (x_1 - x_2)\) 且 \(m_2 \mid (x_1 - x_2)\)。由 \(\gcd(m_1, m_2) = 1\) 得 \(m_1 m_2 \mid (x_1 - x_2)\)，即 \(x_1 \equiv x_2 \pmod{m_1 m_2}\)。\(\blacksquare\)

\textbf{定理 5.4.6（中国剩余定理的一般形式）}：设 \(m_1, m_2, \ldots, m_k\) 两两互素。则同余方程组 \[x \equiv a_i \pmod{m_i}, \quad i = 1, 2, \ldots, k\] 在模 \(M = m_1 m_2 \cdots m_k\) 的意义下有唯一解。解为 \[x \equiv \sum_{i=1}^{k} a_i M_i N_i \pmod{M}\] 其中 \(M_i = \frac{M}{m_i}\)，\(N_i\) 是 \(M_i\) 模 \(m_i\) 的逆元（\(M_i N_i \equiv 1 \pmod{m_i}\)）。

\textbf{证明}：对 \(k\) 用归纳法。\(k = 2\) 时即定理 5.4.5。设对 \(k-1\) 成立。前 \(k-1\) 个方程在模 \(M' = m_1 \cdots m_{k-1}\) 下有唯一解 \(x_0\)。问题化为求 \(x \equiv x_0 \pmod{M'}\) 且 \(x \equiv a_k \pmod{m_k}\)。由 \(\gcd(M', m_k) = 1\)（因 \(m_1, \ldots, m_k\) 两两互素），应用定理 5.4.5 得证。\(\blacksquare\)

中国剩余定理是数论中最优美且最实用的定理之一。它最早出现在中国南北朝时期的数学著作《孙子算经》（约公元 3-5 世纪）中，用于解决''物不知数''问题（``今有物不知其数，三三数之剩二，五五数之剩三，七七数之剩二，问物几何？''）。该定理的现代数学意义在于它建立了环同构 \(\mathbb{Z}/(m_1 \cdots m_k)\mathbb{Z} \cong \mathbb{Z}/m_1\mathbb{Z} \times \cdots \times \mathbb{Z}/m_k\mathbb{Z}\)（当 \(m_i\) 两两互素时），这在代数数论和密码学中有着广泛的应用。例如，RSA 解密算法利用中国剩余定理将模 \(n = pq\) 的运算分解为模 \(p\) 和模 \(q\) 的运算，从而将计算效率提高约 4 倍。在抽象代数中，中国剩余定理推广为一般环的理想理论：若 \(I_1, \ldots, I_k\) 是环 \(R\) 的两两互素的理想，则 \(R/(I_1 \cap \cdots \cap I_k) \cong R/I_1 \times \cdots \times R/I_k\)。

\subsubsection{5.4.5 欧拉函数与欧拉定理}\label{ux6b27ux62c9ux51fdux6570ux4e0eux6b27ux62c9ux5b9aux7406}

\textbf{定义 5.4.2（欧拉函数）}：对正整数 \(n\)，定义 \(\varphi(n)\) 为 \(1, 2, \ldots, n\) 中与 \(n\) 互素的数的个数： \[\varphi(n) = |\{k \in \{1, 2, \ldots, n\} \mid \gcd(k, n) = 1\}|\]

\textbf{注记 5.4.1}：\(\varphi(1) = 1\)（\(\gcd(1, 1) = 1\)）。对素数 \(p\)，\(\varphi(p) = p - 1\)（\(1, 2, \ldots, p-1\) 都与 \(p\) 互素）。

\textbf{定理 5.4.7（欧拉函数的积性）}：若 \(\gcd(a, b) = 1\)，则 \(\varphi(ab) = \varphi(a)\varphi(b)\)。

\textbf{证明}：考虑映射 \(\phi: \{1, \ldots, ab\} \to \{1, \ldots, a\} \times \{1, \ldots, b\}\)，\(\phi(x) = (x \bmod a, x \bmod b)\)。

\textbf{第一步}：\(\phi\) 是双射。由中国剩余定理（定理 5.4.5），对每个 \((r_1, r_2)\)（\(1 \leq r_1 \leq a\), \(1 \leq r_2 \leq b\)），同余方程组 \(x \equiv r_1 \pmod{a}\), \(x \equiv r_2 \pmod{b}\) 在 \(\{1, \ldots, ab\}\) 中有唯一解。因此 \(\phi\) 是一一对应的。

\textbf{第二步}：\(\gcd(x, ab) = 1 \iff \gcd(x, a) = 1\) 且 \(\gcd(x, b) = 1\)。\((\Rightarrow)\)：若 \(d = \gcd(x, a) > 1\)，则 \(d \mid x\) 且 \(d \mid a\)，故 \(d \mid ab\)，\(d \mid \gcd(x, ab) = 1\)，矛盾。\((\Leftarrow)\)：若 \(d = \gcd(x, ab) > 1\)，设 \(p\) 是 \(d\) 的素因子。\(p \mid ab\)，故 \(p \mid a\) 或 \(p \mid b\)。若 \(p \mid a\)，则 \(p \mid \gcd(x, a) = 1\)，矛盾。

\textbf{第三步}：\(\varphi(ab)\) 等于满足 \(\gcd(x, a) = 1\) 且 \(\gcd(x, b) = 1\) 的 \(x\) 的个数。由 \(\phi\) 是双射，这等于满足 \(\gcd(r_1, a) = 1\) 的 \(r_1\) 的个数乘以满足 \(\gcd(r_2, b) = 1\) 的 \(r_2\) 的个数，即 \(\varphi(a) \cdot \varphi(b)\)。\(\blacksquare\)

\textbf{定理 5.4.8（欧拉函数的公式）}：设 \(n = p_1^{a_1} \cdots p_r^{a_r}\)（素数分解），则 \[\varphi(n) = n \prod_{i=1}^{r} \left(1 - \frac{1}{p_i}\right)\]

\textbf{证明}：由积性（定理 5.4.7），\(\varphi(n) = \varphi(p_1^{a_1}) \cdots \varphi(p_r^{a_r})\)。因此只需计算 \(\varphi(p^a)\)。

\(1, 2, \ldots, p^a\) 中与 \(p\) 不互素的数恰好是 \(p\) 的倍数：\(p, 2p, 3p, \ldots, p^a\)，共 \(p^{a-1}\) 个。因此 \(\varphi(p^a) = p^a - p^{a-1} = p^a(1 - 1/p)\)。\(\blacksquare\)

\textbf{定理 5.4.9（欧拉定理）}：设 \(\gcd(a, n) = 1\)，则 \[a^{\varphi(n)} \equiv 1 \pmod{n}\]

\textbf{证明}：设 \(r_1, r_2, \ldots, r_{\varphi(n)}\) 是 \(1, 2, \ldots, n\) 中与 \(n\) 互素的全部数。考虑 \(ar_1, ar_2, \ldots, ar_{\varphi(n)}\)。

\textbf{第一步}：\(ar_i \not\equiv ar_j \pmod{n}\)（\(i \neq j\)）。若 \(ar_i \equiv ar_j \pmod{n}\)，由 \(\gcd(a, n) = 1\) 和定理 5.4.2' 得 \(r_i \equiv r_j \pmod{n}\)，但 \(1 \leq r_i, r_j \leq n\) 且 \(r_i \neq r_j\)，矛盾。

\textbf{第二步}：\(\gcd(ar_i, n) = 1\)。若 \(d = \gcd(ar_i, n) > 1\)，设 \(p\) 是 \(d\) 的素因子。\(p \mid n\) 且 \(p \mid ar_i\)。由 \(\gcd(a, n) = 1\) 知 \(p \nmid a\)，故 \(p \mid r_i\)。但 \(\gcd(r_i, n) = 1\)，故 \(p \nmid n\)，矛盾。

\textbf{第三步}：\(\{ar_1 \bmod n, \ldots, ar_{\varphi(n)} \bmod n\} = \{r_1, \ldots, r_{\varphi(n)}\}\)（作为集合）。因此 \[\prod_{i=1}^{\varphi(n)} ar_i \equiv \prod_{i=1}^{\varphi(n)} r_i \pmod{n}\] 即 \(a^{\varphi(n)} \prod r_i \equiv \prod r_i \pmod{n}\)。由 \(\gcd(\prod r_i, n) = 1\) 和定理 5.4.2'，\(a^{\varphi(n)} \equiv 1 \pmod{n}\)。\(\blacksquare\)

\textbf{推论 5.4.2（费马小定理的推广形式）}：设 \(p\) 是素数，则对任意整数 \(a\)： \[a^p \equiv a \pmod{p}\]

\textbf{证明}：若 \(p \mid a\)，则 \(a^p \equiv 0 \equiv a \pmod{p}\)。若 \(p \nmid a\)，由欧拉定理（\(\varphi(p) = p - 1\)），\(a^{p-1} \equiv 1 \pmod{p}\)，乘以 \(a\) 得 \(a^p \equiv a \pmod{p}\)。\(\blacksquare\)

\subsubsection{5.4.6 威尔逊定理}\label{ux5a01ux5c14ux900aux5b9aux7406}

\textbf{定理 5.4.10（威尔逊定理）}：设 \(p\) 是素数。则 \[(p - 1)! \equiv -1 \pmod{p}\]

\textbf{证明}：考虑模 \(p\) 下 \(1, 2, \ldots, p-1\) 的乘法逆元。对 \(1 \leq a \leq p-1\)，存在唯一的 \(b\)（\(1 \leq b \leq p-1\)）使得 \(ab \equiv 1 \pmod{p}\)（由 \(\gcd(a, p) = 1\) 和裴蜀定理）。

\(a\) 是自身的逆元（\(a^2 \equiv 1 \pmod{p}\)）当且仅当 \(p \mid (a^2 - 1) = (a-1)(a+1)\)。由欧几里得引理，\(p \mid (a-1)\) 或 \(p \mid (a+1)\)，即 \(a = 1\) 或 \(a = p - 1\)。

因此 \(2, 3, \ldots, p-2\) 可以两两配对成 \((a, a^{-1})\)（\(a \neq a^{-1}\)），每对的乘积 \(\equiv 1 \pmod{p}\)。

故 \((p-1)! = 1 \cdot (2 \cdot 3 \cdots (p-2)) \cdot (p-1) \equiv 1 \cdot 1 \cdot (p-1) \equiv -1 \pmod{p}\)。\(\blacksquare\)

\subsubsection{5.4.7 模逆元}\label{ux6a21ux9006ux5143}

\textbf{定义 5.4.3（模逆元）}：设 \(\gcd(a, m) = 1\)。满足 \(ax \equiv 1 \pmod{m}\) 的整数 \(x\) 称为 \(a\) 关于模 \(m\) 的\textbf{乘法逆元}，记为 \(a^{-1} \pmod{m}\)。

\textbf{定理 5.4.11（逆元的存在性与唯一性）}：\(\gcd(a, m) = 1\) 当且仅当 \(a\) 在模 \(m\) 下有乘法逆元。逆元在模 \(m\) 意义下唯一。

\textbf{证明}：\((\Rightarrow)\)：由裴蜀定理，存在 \(x, y\) 使得 \(ax + my = 1\)，即 \(ax \equiv 1 \pmod{m}\)，\(x\) 即为逆元。\((\Leftarrow)\)：若 \(ax \equiv 1 \pmod{m}\)，则 \(ax + my = 1\) 对某 \(y\)，故 \(\gcd(a, m) \mid 1\)，即 \(\gcd(a, m) = 1\)。唯一性：若 \(ax_1 \equiv 1\) 且 \(ax_2 \equiv 1 \pmod{m}\)，则 \(a(x_1 - x_2) \equiv 0 \pmod{m}\)，由 \(\gcd(a, m) = 1\) 得 \(x_1 \equiv x_2 \pmod{m}\)。\(\blacksquare\)

\textbf{算法 5.4.1（扩展欧几里得算法求逆元）}：由辗转相除法的逆向递推，可以高效计算 \(a^{-1} \pmod{m}\)。具体做法：对 \(a\) 和 \(m\) 做辗转相除法得到 \(\gcd(a, m) = 1\)，然后逆向递推表示 \(1 = ax + my\)，\(x\) 即为逆元。

\textbf{详细步骤}：设辗转相除法的过程为：

\[m = aq_1 + r_1, \quad a = r_1 q_2 + r_2, \quad r_1 = r_2 q_3 + r_3, \quad \ldots, \quad r_{k-2} = r_{k-1}q_k + r_k\]

其中 \(r_k = \gcd(a, m) = 1\)。逆向递推：

\[1 = r_k = r_{k-2} - r_{k-1}q_k\] \[= r_{k-2} - (r_{k-3} - r_{k-2}q_{k-1})q_k = (1 + q_{k-1}q_k)r_{k-2} - q_k r_{k-3}\]

继续递推，最终得到 \(1 = ax + my\) 的形式。

\subsubsection{5.4.8 模幂运算}\label{ux6a21ux5e42ux8fd0ux7b97}

\textbf{定理 5.4.12（快速模幂）}：计算 \(a^n \bmod m\) 可以在 \(O(\log n)\) 次模乘法内完成。

\textbf{方法（平方-乘算法）}：将 \(n\) 表示为二进制 \(n = \sum_{i=0}^{k} b_i 2^i\)。则 \[a^n = a^{\sum b_i 2^i} = \prod_{i: b_i = 1} a^{2^i}\] 依次计算 \(a^{2^0}, a^{2^1}, \ldots, a^{2^k} \pmod{m}\)（每次平方取模），然后将 \(b_i = 1\) 对应的项相乘取模。

\textbf{证明正确性}：\(a^{2^{i+1}} = (a^{2^i})^2\)。每步取模不影响最终结果（由同余的乘法性质）。\(\blacksquare\)

\textbf{注记 5.4.2}：快速模幂算法在密码学中有重要应用。例如在 RSA 加密中，需要计算 \(m^e \bmod n\)（\(e\) 和 \(n\) 都是大整数），直接计算 \(m^e\) 是不现实的，但通过平方-乘算法可以在合理时间内完成。

\textbf{具体例子}：计算 \(2^{100} \bmod 1000\)。\(100 = 64 + 32 + 4 = (1100100)_2\)。依次计算： - \(2^1 \equiv 2\) - \(2^2 \equiv 4\) - \(2^4 \equiv 16\) - \(2^8 \equiv 256\) - \(2^{16} \equiv 256^2 = 65536 \equiv 536\) - \(2^{32} \equiv 536^2 = 287296 \equiv 296\) - \(2^{64} \equiv 296^2 = 87616 \equiv 616\)

\(2^{100} = 2^{64} \cdot 2^{32} \cdot 2^4 \equiv 616 \cdot 296 \cdot 16 \pmod{1000}\)。\(616 \cdot 296 = 182336 \equiv 336\)。\(336 \cdot 16 = 5376 \equiv 376\)。故 \(2^{100} \equiv 376 \pmod{1000}\)。

\subsubsection{5.4.9 同余方程}\label{ux540cux4f59ux65b9ux7a0b}

\textbf{定义 5.4.4（线性同余方程）}：形如 \(ax \equiv b \pmod{m}\) 的方程称为\textbf{线性同余方程}。

\textbf{定理 5.4.13（线性同余方程的解）}：设 \(d = \gcd(a, m)\)。 1. \(ax \equiv b \pmod{m}\) 有解当且仅当 \(d \mid b\)。 2. 若有解，则在模 \(m\) 意义下恰好有 \(d\) 个解。

\textbf{证明}：

\begin{enumerate}
\def\labelenumi{(\arabic{enumi})}
\item
  \(ax \equiv b \pmod{m}\) 有解 \(\iff\) 存在 \(x, y\) 使得 \(ax + my = b\) \(\iff\) \(\gcd(a, m) \mid b\)（由裴蜀定理）。
\item
  设 \(d \mid b\)。令 \(a = da'\), \(b = db'\), \(m = dm'\)。则原方程化为 \(a'x \equiv b' \pmod{m'}\)。由 \(\gcd(a', m') = 1\)，存在唯一解 \(x_0 \pmod{m'}\)。原方程的解为 \(x_0, x_0 + m', x_0 + 2m', \ldots, x_0 + (d-1)m' \pmod{m}\)，共 \(d\) 个。\(\blacksquare\)
\end{enumerate}

\subsubsection{5.4.10 二次剩余}\label{ux4e8cux6b21ux5269ux4f59}

\textbf{定义 5.4.5（二次剩余）}：设 \(p\) 是奇素数，\(a\) 是整数且 \(p \nmid a\)。若同余方程 \(x^2 \equiv a \pmod{p}\) 有解，则称 \(a\) 为模 \(p\) 的\textbf{二次剩余}；否则称为\textbf{二次非剩余}。

\textbf{定理 5.4.14（二次剩余的个数）}：模奇素数 \(p\) 的二次剩余恰好有 \(\frac{p-1}{2}\) 个，二次非剩余也有 \(\frac{p-1}{2}\) 个。

\textbf{证明}：考虑平方映射 \(f: \{1, 2, \ldots, p-1\} \to \{1, 2, \ldots, p-1\}\)，\(f(x) = x^2 \bmod p\)。

\(f(x) = f(y) \iff x^2 \equiv y^2 \pmod{p} \iff p \mid (x-y)(x+y)\)。由欧几里得引理，\(p \mid (x-y)\) 或 \(p \mid (x+y)\)。在 \(\{1, \ldots, p-1\}\) 中，\(p \mid (x-y)\) 意味着 \(x = y\)，\(p \mid (x+y)\) 意味着 \(x + y = p\)。因此 \(f\) 的像恰好是每个二次剩余对应两个原像（\(x\) 和 \(p - x\)）。故像的大小为 \(\frac{p-1}{2}\)。\(\blacksquare\)

二次剩余理论是数论中的经典分支，它研究模素数 \(p\) 下的平方根问题。二次剩余的概念自然地引出了 Legendre 符号 \(\left(\frac{a}{p}\right)\)------定义为 \(+1\) 若 \(a\) 是模 \(p\) 的二次剩余，\(-1\) 若 \(a\) 是二次非剩余。Legendre 符号满足乘性 \(\left(\frac{ab}{p}\right) = \left(\frac{a}{p}\right)\left(\frac{b}{p}\right)\) 和 Euler 判别法 \(\left(\frac{a}{p}\right) \equiv a^{(p-1)/2} \pmod{p}\)。二次剩余理论的最高成就是 Gauss 的二次互反律：对于奇素数 \(p, q\)，有 \(\left(\frac{p}{q}\right)\left(\frac{q}{p}\right) = (-1)^{\frac{p-1}{2} \cdot \frac{q-1}{2}}\)，这被称为''数论的钻石''。二次剩余理论在现代密码学中扮演着关键角色------大整数的二次剩余问题是 Goldwasser-Micali 加密方案和 Rabin 密码系统的安全性基础。

\subsection{本章展望}\label{ux672cux7ae0ux5c55ux671b-2}

本章建立了整除性理论、最大公约数与最小公倍数、素数与算术基本定理、以及同余理论的基本框架。这些结果是整数深层结构的基石。在卷二（Ch 12--11）中，裴蜀定理和辗转相除法将在多项式理论中出现类比（多项式的带余除法和最大公因式）。费马小定理和中国剩余定理是密码学（如 RSA 加密）的数学基础。更深入的数论主题（二次互反律、原根、连分数等）将在后续课程中展开。

\hrulefill

\hrulefill

\hrulefill

\hrulefill

\hrulefill

\hrulefill

\newpage

\chapter{01-基础/04-公理化与基础哲学.md}\label{ux57faux784004-ux516cux7406ux5316ux4e0eux57faux7840ux54f2ux5b66.md}

\section{第23章：公理化方法的演进}\label{ux7b2c23ux7ae0ux516cux7406ux5316ux65b9ux6cd5ux7684ux6f14ux8fdb}

公理化方法是数学区别于其他科学的核心方法论，也是数学获得其独特确定性的源泉。本章追溯公理化思想从古希腊到现代的漫长演进：从 Euclid《几何原本》的朴素公理化，到非欧几何对公理独立性的揭示，再到 Hilbert 的形式公理化纲领和元数学标准，最终到达 Bourbaki 结构主义的宏大体系。这条演进路径的核心线索是数学家对''基础''理解的逐步深化------从''公理是自明的真理''到''公理是规定未定义术语性质的假设''，再到''公理系统本身成为数学研究的对象''。公理化方法的成熟使得数学从''描述现实''转向''构造可能的结构''，这一转变为 20 世纪数学的抽象化和统一化奠定了基础。理解公理化方法的演进，不仅有助于把握数学的深层逻辑结构，也有助于理解数学基础危机和数学哲学中的核心争论。

\subsection{172.1 Euclid 的公理化体系}\label{euclid-ux7684ux516cux7406ux5316ux4f53ux7cfb}

\textbf{定义 172.1}（Euclid《几何原本》的公理化结构，约公元前 300 年）：《几何原本》是最早的公理化典范，从 5 条公设和 5 条公理出发，推导出 465 个命题。其结构为：定义 → 公设 → 公理 → 命题（定理及其证明）。

\textbf{公设 172.1}（Euclid 五大公设）： 1. 过任意两点可作一条直线 2. 线段可无限延长 3. 以任意点为中心、任意长为半径可作圆 4. 所有直角相等 5. 若两直线与第三条直线相交，同侧内角之和小于两直角，则两直线在该侧相交（平行公设）

第五公设（平行公设）的不自明性引发了长达两千年的研究，最终导向非欧几何的发现。

Euclid 的《几何原本》（Elements，约公元前 300 年）是数学史上影响最深远的著作之一，其公理化方法为后世数学树立了范本。Euclid 从 23 个定义、5 条公设和 5 条公理出发，系统地推导了 465 个命题，涵盖了平面几何、数论和无理量理论。公设（postulates）是几何学特有的假设（如''过两点可作一条直线''），公理（common notions）是适用于所有数学的普遍真理（如''等于同量的量彼此相等''）。Euclid 公理化的核心在于从少数自明的断言出发，借助逻辑推理建立起庞大的知识体系。然而，第五公设（平行公设）的表述比前四条公设复杂得多，Euclid 本人也尽量避免使用它（《几何原本》前 28 个命题都不依赖第五公设）。这一''不完美''激发了长达两千年的研究，最终导致非欧几何的发现------这是公理化方法发展史上最深刻的转折。

\textbf{定理 172.1}（Euclid 几何中三角形的内角和定理）：在 Euclid 几何中，任意三角形的内角和等于两个直角（即 \(\pi\) 弧度）。这一命题是《几何原本》第一卷命题 32，其证明依赖于平行公设（第五公设）：通过三角形的一个顶点作对边的平行线，利用同位角相等和内错角相等的性质，将三个内角转移到一条直线上，证明其和为平角。这一构造性证明是 Euclid 几何中最优美的论证之一，同时也展示了平行公设在 Euclid 几何定理体系中的不可或缺性。

\textbf{定理 172.2}（Euclid 算法）：在自然数中，辗转相除法（Euclid 算法）用于求两个正整数的最大公约数。具体步骤为：给定 \(a > b > 0\)，连续做带余除法 \(a = bq_1 + r_1\)，\(b = r_1 q_2 + r_2\)，\(r_1 = r_2 q_3 + r_3\)，\(\ldots\)，直到余数为 \(0\)，最后一个非零余数即为 \(\gcd(a, b)\)。Euclid 算法的正确性基于一个关键观察：\(\gcd(a, b) = \gcd(b, r)\)（其中 \(r\) 是 \(a\) 除以 \(b\) 的余数），因为 \(a = bq + r\) 蕴含 \(d \mid a\) 且 \(d \mid b \iff d \mid b\) 且 \(d \mid r\)。Euclid 算法是数论中最古老且最重要的算法，其时间复杂度为 \(O(\log(\min(a, b)))\)（Lamé 定理，1844）。

\textbf{注记 172.1}（Euclid《几何原本》的历史影响）：《几何原本》在西方印刷史上的发行量仅次于《圣经》，被翻译成几乎所有语言。它不仅是数学教科书，更是西方理性思维方式的典范------Spinoza 的《伦理学》甚至模仿 Euclid 的''定义-公理-命题''格式来论述哲学。Newton 的《自然哲学的数学原理》也采用了类似的公理化叙述方式。Euclid 公理化方法对后世的影响远远超出了数学领域，它塑造了西方思想中''从第一原理出发进行演绎推理''的范式。

\subsection{172.2 非欧几何：平行公设独立性的证明}\label{ux975eux6b27ux51e0ux4f55ux5e73ux884cux516cux8bbeux72ecux7acbux6027ux7684ux8bc1ux660e}

\textbf{证明}：Lobachevsky和Bolyai的方法：接受Euclid前四条公设，否定第五公设，推导出一系列非欧几何定理（如三角形内角和小于\(\pi\)，不存在相似三角形等）。独立性证明由Beltrami（1868）完成：构造Euclid空间中的伪球面模型，将双曲几何的每条公理翻译为Euclid几何的定理。若双曲几何存在矛盾，则该矛盾在Euclid几何模型中也有对应矛盾，与Euclid几何的一致性（假设）相悖。故平行公设独立于前四条公设。\(\blacksquare\)

\begin{itemize}
\tightlist
\item
  \textbf{双曲几何}（Lobachevsky-Bolyai）：过直线外一点至少有两条平行线。曲率为负，三角形内角和小于 \(\pi\)。
\item
  \textbf{椭圆几何}（Riemann 1854）：不存在平行线。曲率为正，三角形内角和大于 \(\pi\)。
\end{itemize}

\textbf{证明}：Poincaré圆盘模型：在单位圆盘\(D=\{z\in\mathbb{C}:|z|<1\}\)上定义双曲度量\(ds^2=\frac{4|dz|^2}{(1-|z|^2)^2}\)。双曲直线为与单位圆正交的圆弧。可验证此模型满足除平行公设外的所有Euclid公理及双曲平行公设（过直线外一点存在无穷多条不相交的''直线''）。Klein模型用射影几何，上半平面模型用复分析，三者等价。任一模型中的矛盾均可翻译为\(\mathbb{R}^n\)中Euclid几何的陈述，故\(\operatorname{Con}(\text{Euclid})\Rightarrow\operatorname{Con}(\text{非欧})\)。\(\blacksquare\)

非欧几何的创立标志着数学从''描述物理空间''转向''研究可能的逻辑结构''，是公理化方法成熟的里程碑。

\subsection{172.3 Hilbert 的公理化纲领}\label{hilbert-ux7684ux516cux7406ux5316ux7eb2ux9886}

\textbf{定义 172.2}（Hilbert《几何基础》，1899）：Hilbert 将几何完全公理化，从 21 条公理出发，分为五组（关联公理、顺序公理、合同公理、平行公理、连续公理），严格演绎全部几何定理。Hilbert 的公理化满足三条元数学标准：

\begin{itemize}
\tightlist
\item
  \textbf{一致性}（consistency）：公理系统不蕴含矛盾
\item
  \textbf{独立性}（independence）：每条公理不能由其余公理推导
\item
  \textbf{完备性}（completeness）：系统内所有真命题均可由公理推导
\end{itemize}

Hilbert 纲领将公理化思想推向顶点，并提出了著名的 Hilbert 计划：用有限方法证明全部数学的一致性。

Hilbert 的《几何基础》（Grundlagen der Geometrie, 1899）是公理化方法发展的里程碑。与 Euclid 依赖直观定义（如''点是没有部分的东西''）不同，Hilbert 将''点''、``直线''、``平面''视为未定义的基本概念，其性质完全由公理规定------这一''隐定义''（implicit definition）的思想是现代公理化的精髓。Hilbert 的三条元数学标准（一致性、独立性、完备性）为公理系统设立了严格的评判准则。Hilbert 计划（Hilbert's Program, 1920s）更进一步，试图用''有穷方法''（finitistic methods）证明全部数学的一致性------这一宏大的构想深刻地影响了 20 世纪的数学基础研究。虽然 Godel 不完备性定理（1931）证明了 Hilbert 计划的原始目标不可能实现，但 Hilbert 的公理化方法论和元数学思想至今仍是数学研究的基本范式。

\subsection{172.4 Bourbaki 的结构主义}\label{bourbaki-ux7684ux7ed3ux6784ux4e3bux4e49}

\textbf{定义 172.3}（Bourbaki 学派，1935-）：以 Nicolas Bourbaki 为笔名的法国数学家群体，目标是用公理化方法从集合论出发重建整个数学体系。其核心概念是\textbf{母结构}（mother structures）：

\begin{itemize}
\tightlist
\item
  \textbf{代数结构}：由合成律定义的运算结构（群、环、域、模等）
\item
  \textbf{序结构}：由序关系定义的结构（偏序、全序、格等）
\item
  \textbf{拓扑结构}：由邻域、极限、连续性定义的结构
\end{itemize}

Bourbaki 的《数学原本》（Elements de Mathematique）从集合论出发，按结构层次体系化地展开各数学分支。其方法论强调抽象化和公理化，对 20 世纪数学产生了深远影响。

Bourbaki 学派的结构主义代表了公理化方法的最高发展阶段。其核心洞见是：数学中看似不同的分支共享着相同的底层结构。例如，整数加群 \((\mathbb{Z}, +)\) 与正实数乘法群 \((\mathbb{R}_{>0}, \times)\) 在结构上是同构的（都是无穷循环群），因此群论中证明的定理可以同时应用于这两个看起来毫无关系的领域。Bourbaki 的''母结构''分类（代数结构、序结构、拓扑结构）并非独断的------它们覆盖了经典数学的核心骨架，而''交叉结构''（如拓扑群------同时拥有拓扑结构和代数结构）则产生了最丰富的数学理论。Bourbaki 还引入了''数学结构''的严格定义：一个结构是在某个基础集合上添加公理规定的性质。这种方法论使得数学从''具体计算''转向''抽象结构研究''，其影响遍布代数拓扑、泛函分析、代数几何等现代数学分支。Bourbaki 对''严格性''的坚持（``每个断言都必须有证明''）深刻塑造了 20 世纪数学的写作风格。

\hrulefill

\hrulefill

\hrulefill

\hrulefill

\hrulefill

\section{第24章：集合论与数学基础危机}\label{ux7b2c24ux7ae0ux96c6ux5408ux8bbaux4e0eux6570ux5b66ux57faux7840ux5371ux673a}

19 世纪末至 20 世纪初，数学基础经历了深刻的反思与重构。集合论的创立与悖论的发现改变了数学对自身基础的理解。本章从 Cantor 创立集合论的革命性工作出发，依次讨论 Russell 悖论对朴素集合论的致命打击、ZFC 公理化集合论作为解决方案的建立，以及 Godel 不完备性定理和 Cohen 力迫法对数学基础的根本性影响。数学基础危机（Grundlagenkrise）是 20 世纪数学史上最重要的事件之一------它不仅催生了现代数理逻辑和公理化集合论，还深刻影响了数学哲学（逻辑主义、形式主义、直觉主义三大流派正是在这场危机中形成的）。理解基础危机及其解决，是理解当代数学的''哲学自觉''------即数学对自身方法和前提的反思------的关键。

\subsection{173.1 Cantor 集合论与对角线论证}\label{cantor-ux96c6ux5408ux8bbaux4e0eux5bf9ux89d2ux7ebfux8bbaux8bc1}

\textbf{定义 173.1}（Cantor 朴素集合论，1874-1897）：Cantor 创立集合论，核心贡献包括：

\begin{itemize}
\tightlist
\item
  无穷集合的等势定义（一一对应）
\item
  可数集与不可数集的区分：\(\mathbb{N}, \mathbb{Z}, \mathbb{Q}\) 可数，\(\mathbb{R}\) 不可数
\item
  超限基数的层级：\(\aleph_0 < \aleph_1 < \aleph_2 < \cdots\)
\item
  \textbf{连续统假设}（CH）：\(2^{\aleph_0} = \aleph_1\)（Hilbert 23 个问题的第一个）
\end{itemize}

\textbf{证明}：假设\(\mathbb{R}\)可数，则\([0,1]\)可数。将\([0,1]\)中的数写为十进制小数展开：\(x_n=0.d_{n1}d_{n2}d_{n3}\ldots\)（\(d_{nj}\in\{0,1,\ldots,9\}\)）。构造\(y=0.e_1e_2e_3\ldots\)，其中\(e_n=\begin{cases}1&d_{nn}\neq 1\\2&d_{nn}=1\end{cases}\)。\(y\)与每个\(x_n\)在第\(n\)位不同，故\(y\neq x_n\)对所有\(n\)。但\(y\in[0,1]\)，与\([0,1]\)可数矛盾。\(\blacksquare\)

\subsection{173.2 Russell 悖论与类型论}\label{russell-ux6096ux8bbaux4e0eux7c7bux578bux8bba}

\textbf{定义 173.2}（Russell 悖论，1901）：在 Cantor 朴素集合论中，令 \(R = \{x : x \notin x\}\)（所有不以自身为元素的集合的集合）。则

\[R \in R \iff R \notin R\]

矛盾。这暴露了无限制概括公理（任意性质定义一个集合）的不一致性。

\textbf{证明}：分支类型论将全域分为类型层级：类型\(0\)为个体（原子），类型\(1\)为个体的集合，类型\(2\)为个体集合的集合\ldots\ldots 公式只能断言\(x^n\in y^{n+1}\)（\(x\)的类型为\(n\)，\(y\)的类型为\(n+1\)）。Russell悖论中\(R=\{x:x\notin x\}\)要求\(x\)和\(x\)的集合具有相同类型，违反类型层级。因此''\(R\in R\)``不是合式公式，悖论不出现。可归约公理（reducibility axiom）将分支类型论简化为简单类型论，恢复经典数学的推理能力。\(\blacksquare\)

\subsection{173.3 ZFC 公理化集合论}\label{zfc-ux516cux7406ux5316ux96c6ux5408ux8bba}

\textbf{定义 173.3}（Zermelo-Fraenkel 集合论 ZFC，1908-1922）：ZFC 用 9 条公理替代无限制概括公理，通过限制集合构造方式避免已知悖论：

\begin{itemize}
\tightlist
\item
  \textbf{外延公理}：集合由元素唯一决定
\item
  \textbf{分离公理模式}（Zermelo 1908）：从已有集合中分离出满足某性质的元素（而非从全体中构造）
\item
  \textbf{替换公理模式}（Fraenkel 1922, Skolem 1922）：函数的像仍是集合
\item
  \textbf{正则公理}（von Neumann 1925）：禁止 \(x \in x\) 等自属于关系
\item
  \textbf{选择公理}（AC）：任意非空集合族存在选择函数
\end{itemize}

ZFC 是当代数学的标准基础，绝大多数数学定理可在 ZFC 中形式化。

ZFC 公理化集合论是数学基础危机（1900-1930）的最终解决方案。在 Russell 悖论（1901）暴露了 Cantor 朴素集合论的不一致性后，数学家面临一个严峻的问题：如何在保留集合论强大表达能力的同时避免悖论？Zermelo 于 1908 年提出的公理化方案（Z 系统）用分离公理替代了无限制概括公理，有效地隔绝了 Russell 悖论。Fraenkel 和 Skolem 于 1922 年独立添加了替换公理模式，形成了 ZF 系统。von Neumann 于 1925 年添加了正则公理，最终形成了 ZFC（C 表示选择公理）。ZFC 的核心策略是限制集合的构造方式------从空集出发，通过配对、并集、幂集、分离、替换等''安全''操作逐步构造更大的集合，而不允许''所有满足某性质的对象的集合''这样的无限制构造。ZFC 的成功在于它在表达力（几乎所有的经典数学都可以在其中形式化）和安全性（至今未发现悖论）之间取得了平衡。

\textbf{证明}：Godel编号将每个公式和证明序列编码为自然数。定义可证性谓词\(\operatorname{Bew}(x)\)（``\(x\)是某公式的证明的Godel数''）。由对角线引理（在包含算术的系统中可证），存在公式\(G\)满足\(\vdash G\leftrightarrow\neg\operatorname{Bew}(\ulcorner G\urcorner)\)。\(G\)断言''\(G\)不可证''。若\(\vdash G\)，则\(\vdash\operatorname{Bew}(\ulcorner G\urcorner)\)（\(\Sigma_1\)-完备性），与\(G\)矛盾。若\(\vdash\neg G\)，则\(\vdash\operatorname{Bew}(\ulcorner G\urcorner)\)且\(\nvdash G\)（\(\omega\)-一致性，Rosser改进后仅需一致性），矛盾。故\(G\)不可判定。\(\blacksquare\)

不完备性定理表明 Hilbert 计划（用有限方法证明全部数学的一致性）原则上不可能完全实现。

\textbf{证明}：Godel（1938）构造可构造宇宙\(L\)：\(L=\bigcup_{\alpha\in\mathbf{Ord}}L_\alpha\)，其中\(L_0=\emptyset\)，\(L_{\alpha+1}=\operatorname{Def}(L_\alpha)\)（\(L_\alpha\)的一阶可定义子集），\(L_\lambda=\bigcup_{\alpha<\lambda}L_\alpha\)（极限序数）。\(L\)满足ZFC+CH（\(L\)中\(2^{\aleph_0}=\aleph_1\)）。Cohen（1963）发明力迫法：在可数传递模型\(M\)中，添加\(\aleph_2\)个Cohen实数（generic filter \(G\)），力迫扩张\(M[G]\)满足\(2^{\aleph_0}=\aleph_2\)。\(M[G]\)仍满足ZFC，但CH不成立。因此CH独立于ZFC。\(\blacksquare\)

\hrulefill

\hrulefill

\hrulefill

\hrulefill

\hrulefill

\hrulefill

\section{第25章：数学哲学------实在论与反实在论}\label{ux7b2c25ux7ae0ux6570ux5b66ux54f2ux5b66ux5b9eux5728ux8bbaux4e0eux53cdux5b9eux5728ux8bba}

数学哲学探讨数学对象的存在性（本体论）和数学知识的获取方式（认识论）。本章系统介绍数学哲学中的五大流派：柏拉图主义（数学实在论）主张数学对象是独立于人类心智的抽象实体；形式主义将数学还原为无意义的符号操作；逻辑主义试图将数学归约为逻辑；直觉主义（构造主义）坚持数学对象必须可构造；结构主义认为数学的本质是结构而非个体对象。这些流派围绕''数学对象是否真实存在''和''数学知识如何可能''这两个根本问题展开了深刻争论。本章还将介绍 Quine-Putnam 不可或缺性论证、Benacerraf 的认识论困难、Curry-Howard 同构等 20 世纪数学哲学的核心论题，以及范畴论为结构主义提供的现代数学语言。

\subsection{174.1 柏拉图主义（数学实在论）}\label{ux67cfux62c9ux56feux4e3bux4e49ux6570ux5b66ux5b9eux5728ux8bba}

\textbf{定义 174.1}（数学柏拉图主义）：数学对象（数、集合、函数、空间等）是独立于人类心智和物理世界存在的抽象实体。数学真理是关于这些抽象实体的客观描述。

\textbf{论证 174.1}（Quine-Putnam 不可或缺性论证）：数学在自然科学中不可或缺------科学理论的最佳表述必然包含数学；科学理论的成功（经验确证）同时确证了其数学部分的真理性。

\textbf{反驳 174.1}（Benacerraf 的认识论困难，1973）：若数学对象是因果惰性的抽象实体（不在时空之中，不与物理世界因果交互），则人类如何获得关于它们的知识？标准的知识因果理论（知识需与被知对象有因果联系）无法解释数学知识。

\textbf{定义 174.1a}（Goedel 的数学直觉概念，1947）：Goedel 本人是柏拉图主义者，他提出人类具有一种''数学直觉''（mathematical intuition）------类似于感官知觉但指向抽象对象的认知能力。Goedel 将数学直觉类比于感官知觉：``尽管它们（集合）是遥远的，我们确实对它们有某种感知\ldots\ldots 我认为没有理由对数学直觉的信任少于对感官知觉的信任。''这一立场被称为''Goedelian Platonism''，但它面临一个尖锐的批评：感官知觉可以通过神经科学和物理学来解释其机制，而数学直觉没有类似的因果解释。

\textbf{定理 174.1}（Frege 的''数作为概念的外延''）：Frege 在《算术基础》（Grundlagen der Arithmetik, 1884）中定义了''数''：数 \(0\) 是概念''不等于自身''的外延；数 \(n+1\) 是概念''与 \(n\) 等数''的外延。具体地，两个概念 \(F\) 和 \(G\) 是''等数''的（equinumerous），当且仅当存在 \(F\) 的对象与 \(G\) 的对象之间的一一对应。数 \(n\) 就是所有与某概念等数的概念所组成的等价类。这一构造可以将自然数（进而所有数学）还原为纯逻辑概念------这是逻辑主义的核心理念。但 Russell 悖论（1901）表明 Frege 的基本法则 V（概念的外延相等当且仅当概念等价）是不一致的，从而摧毁了 Frege 的原始逻辑主义方案。

\textbf{定理 174.2}（Gödel 不完备性定理对形式主义的影响）：Gödel 第二不完备性定理（1931）表明：任何包含初等算术的一致形式系统，不能在其自身内部证明自身的一致性。这直接打击了 Hilbert 计划的核心目标------用有穷方法证明数学的一致性。形式主义因此必须接受一个事实：数学的一致性不能通过纯形式手段内部自证，而需要外部论证（如 Gentzen 对算术一致性的证明使用了超限归纳至 \(\varepsilon_0\)，这已超出有穷方法）。然而，形式主义的方法论价值并未因此被否定------公理化方法、形式化推理和元数学分析仍然是数学研究的基本工具。

\subsection{174.2 形式主义}\label{ux5f62ux5f0fux4e3bux4e49}

\textbf{定义 174.2}（数学形式主义）：数学是关于无意义符号串的形式操作。数学的''真''等同于''在一致的公理系统中可推导''。数学不需要本体论承诺------无需假定数学对象''真实存在''。

\textbf{论证 174.2}（一致性的核心地位）：形式主义将数学还原为符号游戏，只需关注推理规则和公理的一致性。这避免了柏拉图主义的认识论困难。

\textbf{反驳 174.2}： 1. Godel 第二不完备性定理表明一致性不能自证（无法在系统内部证明系统自身的一致性） 2. 纯形式主义无法解释公理选择的合理性------为什么选择 ZFC 而非其他公理系统？这暗示了某种超出纯形式的数学直觉

\subsection{174.3 逻辑主义}\label{ux903bux8f91ux4e3bux4e49}

\textbf{定义 174.3}（数学逻辑主义，Frege 1884, Russell-Whitehead 1910-1913）：数学可完全归约为逻辑。数学概念是逻辑概念的定义缩写，数学定理是逻辑定理的特例。

\textbf{论证 174.3}：Frege 的《算术基础》将自然数定义为概念的外延：数 \(n\) 是''与某概念等数''这个概念的外延。Russell 和 Whitehead 的《数学原理》系统尝试从纯逻辑公理推导出全部数学。

\textbf{反驳 174.3}： 1. Russell 悖论暴露了 Frege 原始系统中的矛盾（基本法则 V 导致矛盾） 2. ZFC 等公理化集合论使用了超越纯逻辑的公理（无穷公理、选择公理、替换公理），这些公理的存在论承诺不能被归约为纯逻辑

\subsection{174.4 直觉主义（构造主义）}\label{ux76f4ux89c9ux4e3bux4e49ux6784ux9020ux4e3bux4e49}

\textbf{定义 174.4}（数学直觉主义 / 构造主义，Brouwer 1907）：数学是心智构造活动。数学对象的存在性等于其可构造性（given by a construction）。核心原则：

\begin{itemize}
\tightlist
\item
  \textbf{排中律不适用于无穷域}：\(\forall x P(x) \lor \exists x \neg P(x)\) 在无穷域上不成立，因为可能既无构造性全称证明也无构造性反例
\item
  \textbf{存在性证明必须是构造性的}：证明 \(\exists x P(x)\) 必须给出 \(x\) 的具体构造，纯矛盾论证（假设 \(\neg \exists x P(x)\) 导致矛盾）不被接受
\end{itemize}

\textbf{论证 174.4}（Curry-Howard 同构，1969）：直觉主义逻辑与类型化 \(\lambda\)-演算之间存在精确对应------命题对应类型，证明对应程序（\(\lambda\)-项），证明的规范化对应程序的归约。这建立了构造性数学与计算机科学的深层联系。

\textbf{反驳 174.4}：拒绝排中律使得经典数学的许多标准结果不可接受。直觉主义数学在技术上也更复杂（需要追踪构造性）。

\textbf{定理 174.3}（Brouwer 的连续性原理）：Brouwer 提出，在直觉主义数学中，从 \(\mathbb{N}^{\mathbb{N}}\) 到 \(\mathbb{N}\) 的每个函数都是连续的。这一结论在经典数学中显然是错误的（存在不连续的函数），但在直觉主义框架中，它源于''每个函数都是可计算的''这一构造性预设。Brouwer 连续性原理的证明使用了选择序列（choice sequences）的概念------一个无穷序列的元素可以逐步自由选择，不完全由规则决定。连续性原理表明，直觉主义数学不仅否定经典逻辑的某些定律，还引入了新的数学原理，这些原理在经典数学中不成立。这使得直觉主义数学不是经典数学的''子集''，而是一个不同的数学体系。

\textbf{定理 174.4}（Bishop 的构造性分析，1967）：Errett Bishop 在《构造性分析基础》（Foundations of Constructive Analysis）中展示了：大部分经典分析（包括微积分的基本定理、隐函数定理等）可以在构造性框架下重新证明，且无需使用 Brouwer 连续性原理等争议性直觉主义原理。Bishop 的构造性方法不否定经典数学------它只是表明经典定理的构造性版本在构造性逻辑中也是可证的。这一工作极大地推动了构造性数学在计算机科学中的应用（程序提取、形式化验证）。

\subsection{174.5 结构主义与范畴论基础}\label{ux7ed3ux6784ux4e3bux4e49ux4e0eux8303ux7574ux8bbaux57faux7840}

\textbf{定义 174.5}（数学结构主义）：数学的本质是\textbf{结构}（对象间的关系模式），而非个体对象自身。自然数 3 不在于是''什么''，而在于它在自然数序列中的位置关系（Peano 公理刻画的后续关系）。同一个结构可以在不同的具体系统中实例化。

\textbf{定义 174.6}（范畴论结构主义，Lawvere 1960s）：范畴论为结构主义提供了精确的数学语言。在范畴论框架中，数学对象的性质完全由它与其他对象的态射关系决定。\textbf{Yoneda 引理}精确表达了这一思想：一个对象完全由它到所有对象（或所有对象到它）的态射决定。

\textbf{证明}：ETCS公理包括：（1）Set有有限极限（终对象\(1\)、拉回、积）；（2）Set有幂对象（每个\(A\)的幂对象\(P(A)\)，满足\(\operatorname{Hom}(B,P(A))\cong\operatorname{Sub}(B\times A)\)）；（3）Set有自然数对象\(N\)（满足Peano公理的范畴版本）；（4）Set是良幂的且具有选择公理（每个满射有截面）。这些公理在范畴语言中刻画了集合范畴的本质结构。Mitchell-Bénabou语言允许在ETCS内部进行经典数学推理，将\(\in\)关系替换为子对象分类器\(\Omega\)的特征映射。\(\blacksquare\)

\hrulefill

\hrulefill

\hrulefill

\hrulefill

\hrulefill

\hrulefill

\section{第26章：数学基础的认识论问题}\label{ux7b2c26ux7ae0ux6570ux5b66ux57faux7840ux7684ux8ba4ux8bc6ux8bbaux95eeux9898}

本章讨论当代数学哲学中的核心认识论议题，聚焦于三个具有深刻哲学意涵的现代问题。选择公理（AC）的争论是数学史上最持久的哲学争论之一------尽管 AC 在数学中几乎不可或缺，但其非构造性（声称存在选择函数但不给出构造方法）以及 Banach-Tarski 悖论等反直观后果，使得它在哲学上长期受到质疑。计算机证明的认识论问题则涉及一个全新的挑战：当定理的证明依赖于超出人类验证能力的计算时，我们能否合理地说''我们知道了''这个定理？四色定理和 Kepler 猜想的证明引发了关于数学知识的本质和证明概念的深刻反思。最后，数学基础的多元性表明，不存在单一''正确''的数学基础------ZFC、ETCS、同伦类型论（HoTT）等不同方案各有优劣，揭示了数学的不同结构侧面，这种多元论本身具有重要的哲学意义。

\subsection{175.1 选择公理的争论}\label{ux9009ux62e9ux516cux7406ux7684ux4e89ux8bba}

\textbf{定义 175.1}（选择公理 AC 的数学等价物）：在 ZF 中，AC 等价于 Zorn 引理、良序定理、Tychonoff 定理（紧空间的乘积紧）、任意向量空间有基、Hahn-Banach 定理等。AC 在泛函分析、代数和拓扑中不可或缺。

\textbf{定义 175.2}（Banach-Tarski 悖论，1924）：在 ZFC 中，一个三维实心球可以分解为有限个碎片，经旋转和平移后重新组合为两个与原球等体积的实心球。这是 AC 的结果（依赖于不可测集的存在），表明 AC 蕴含与直观几何相悖的结论。

\textbf{评注 175.1}：AC 是当代数学中仍在哲学层面争论的唯一公理。其争论焦点不在于''对错''（AC 已被证明独立于 ZF，Godel 1938, Cohen 1963），而在于其本体论和认识论地位。

\textbf{定理 175.1}（AC 的等价形式）：以下命题在 ZF 中彼此等价（即每一个都与 AC 等价）： 1. Zorn 引理：每个非空偏序集中，若每个链都有上界，则该偏序集有极大元 2. 良序定理：每个集合都可以被良序化 3. Tychonoff 定理：任意紧致拓扑空间的乘积是紧致的 4. 每个向量空间都有一组基 5. 每个满射都有右逆（截面）

\textbf{证明}：等价性通过 ZF 中的循环推导建立。

\textbf{AC \(\Rightarrow\) Zorn 引理}。设 \((P, \leq)\) 为非空偏序集，其中每个链都有上界。假设 \(P\) 无极大元。对每个链 \(C \subset P\)，定义其上界组成的集合 \(U(C)\)（由假设非空）。由 AC，存在选择函数 \(f\) 为每个链 \(C\) 选取 \(f(C) \in U(C)\)。用超限归纳构造严格递增链：令 \(p_0 = f(\emptyset)\)。对后继序数 \(\alpha + 1\)，令 \(p_{\alpha+1} = f(\{p_\beta : \beta \leq \alpha\})\)。对极限序数 \(\lambda\)，\(\{p_\beta : \beta < \lambda\}\) 是链，令 \(p_\lambda\) 为其上界（\(p_\lambda \neq\) 极大元，故 \(p_\lambda\) 严格大于所有 \(p_\beta\)）。超限归纳过程不可能在到达某序数时停止（否则得到极大元），故对任意序数 \(\alpha\) 都存在 \(p_\alpha\)，得到 \(P\) 的单射类序列 \(\{p_\alpha\}\)。由 Hartogs 定理，\(P\) 的基数将被超限序列的基数超越，矛盾。因此 \(P\) 必有极大元。

\textbf{Zorn 引理 \(\Rightarrow\) 良序定理}。设 \(X\) 为任意集合。考虑 \(P = \{(A, \leq_A) : A \subset X, \leq_A \text{ 是 } A \text{ 的良序}\}\)，按 \((A, \leq_A) \preceq (B, \leq_B)\) 当且仅当 \(A\) 是 \(B\) 的前段且 \(\leq_A\) 是 \(\leq_B\) 的限制，\(P\) 为偏序集。每个链在 \(P\) 中有上界（取并集）。由 Zorn 引理，\(P\) 有极大元 \((M, \leq_M)\)。若 \(M \neq X\)，取 \(x \in X \setminus M\)，将 \(x\) 放在 \(M\) 中所有元素之后得到 \(M \cup \{x\}\) 上的良序，与极大性矛盾。故 \(M = X\)，即 \(X\) 可被良序化。

\textbf{良序定理 \(\Rightarrow\) AC}。设 \(\{X_i\}_{i \in I}\) 为非空集族。由良序定理，\(\bigcup_{i \in I} X_i\) 可被良序化。对每个 \(i \in I\)，定义 \(f(i)\) 为 \(X_i\) 在良序下的最小元。此 \(f\) 即为选择函数。

\textbf{AC \(\Rightarrow\) 每个满射有右逆}。设 \(f: A \to B\) 为满射。对每个 \(b \in B\)，\(f^{-1}(b)\) 非空。由 AC，存在选择函数 \(g(b) \in f^{-1}(b)\)，则 \(g: B \to A\) 满足 \(f \circ g = \operatorname{id}_B\)。

\textbf{AC \(\Rightarrow\) 每个向量空间有基}。与 Zorn 引理的应用类似：设 \(V\) 为向量空间，\(P\) 为 \(V\) 中所有线性无关子集的集合（按包含偏序）。每个链的并仍是线性无关集（若某有限线性组合为零，则组合中每个向量属于链中某成员，由链的性质它们都属于某一固定成员，矛盾），故 \(P\) 满足 Zorn 引理条件。由 AC（或直接由 Zorn 引理），\(P\) 有极大元，即为基。

\textbf{AC \(\Rightarrow\) Tychonoff 定理}。Tychonoff 定理等价于 AC 需要用到拓扑中的 Alexander 子基定理或网的收敛性。AC 保证可在乘积拓扑中选取收敛子网。逆向从 Tychonoff 定理推出 AC（Kelley, 1950）需要更精细的构造，此处略去。\(\blacksquare\)''

\textbf{定理 175.2}（AC 的弱形式）：在数学实践中，AC 的许多应用实际上只需要其弱形式： - \textbf{可数选择公理}（AC\(_\omega\)）：可数个非空集合的笛卡尔积非空。这足以证明''可数多个可数集的并集是可数的''和''实数序列的连续性等价于 \(\varepsilon\)-\(\delta\) 连续性'' - \textbf{依赖选择公理}（DC）：比 AC\(_\omega\) 略强，足以证明 Baire 纲定理和 Banach-Steinhaus 定理 - \textbf{超滤子引理}（BPI）：每个布尔代数中的滤子可以扩张为超滤子。等价于 Tychonoff 定理对紧致 Hausdorff 空间的特殊情形

这些弱形式的独立性状态与完整 AC 不同------它们都比 AC 弱，但在 ZF 中同样不可证。大多数分析学的结果只需要 DC 或 AC\(_\omega\)，而不需要完整的 AC。

\textbf{证明}：（AC\(_\omega\) 与 AC 的独立性关系）Godel（1938）证明了 AC 在 ZF 中相对一致（ZFC 一致则 ZF 一致）。Cohen（1963）用力迫法证明了 AC 在 ZF 中不可证。至于弱形式：AC\(_\omega\) 不能从 ZF 推导（用对称模型，Cohen 的原始力迫可生成 AC\(_\omega\) 失效但 ZF 一致的模型）；DC 严格强于 AC\(_\omega\)（存在满足 AC\(_\omega\) 但不满足 DC 的 ZF 模型）。BPI（超滤子引理）与 AC 独立------存在 BPI 成立而 AC 不成立的模型（如 Feferman 模型）。这些相对强度结果（由 Jech、Solovay、Feferman 等人建立）表明 AC 的弱形式形成了严格的强度层级： \[\text{ZF} < \text{ZF} + \text{BPI} < \text{ZF} + \text{AC}_\omega < \text{ZF} + \text{DC} < \text{ZFC}.\] 特别地，BPI 不能推出 AC\(_\omega\)，AC\(_\omega\) 不能推出 DC，DC 不能推出完整 AC。\(\blacksquare\)

\subsection{175.2 计算机证明的认识论}\label{ux8ba1ux7b97ux673aux8bc1ux660eux7684ux8ba4ux8bc6ux8bba}

\textbf{定义 175.3}（计算机辅助证明）：某些数学定理的已知证明依赖大量计算机计算，人类无法在合理时间内手工验证。典型案例：

\begin{itemize}
\tightlist
\item
  \textbf{四色定理}（Appel-Haken 1976，Gonthier 2005 用 Coq 形式化）：任何平面地图可用四种颜色着色使得相邻区域颜色不同。证明需检查 1936 种可约构型，计算量超出人力范围
\item
  \textbf{Kepler 猜想}（Hales 1998-2014，Flyspeck 项目形式化）：面心立方堆积是三维空间中球的最密堆积方式
\end{itemize}

\textbf{认识论问题}：我们能否信任人类无法直接验证的计算机证明？这涉及对形式化证明检查器（proof assistant，如 Coq、Lean）的信任------检查器的内核通常只有几百行代码，可被人类审查。但编译器和硬件的正确性又构成了信任链的延伸问题。

\subsection{175.3 数学基础的多元性}\label{ux6570ux5b66ux57faux7840ux7684ux591aux5143ux6027}

\textbf{定义 175.4}（数学基础的多元论）：不同基础方案（ZFC、ETCS、同伦类型论 HoTT、构造性类型论）各有优劣，可并存互为补充：

\begin{itemize}
\tightlist
\item
  \textbf{ZFC}：最广泛使用，几乎所有数学可在其中形式化
\item
  \textbf{ETCS（Lawvere）}：以范畴论语言表述，强调结构而非元素
\item
  \textbf{同伦类型论（HoTT / UF，Voevodsky 等，2010s）}：将类型论、同伦论和数学基础统一，提供新的基础视角
\end{itemize}

数学基础的多元性表明，不存在单一''正确''的基础------不同基础方案揭示了数学的不同结构侧面。

\textbf{定理 175.3}（ZFC 与 ETCS 的等价性）：ZFC 和 ETCS（Lawvere 的范畴论基础）在如下意义下等价：在 ZFC 中可定义布尔拓扑斯，反之在 ETCS 扩充了''替换公理''后也可模拟 ZFC 的集合论构造。但两者强调的数学结构不同------ZFC 强调元素的''属于''关系，ETCS 强调对象间的''态射''关系。这一等价性表明，数学基础的多元性不是''分歧''，而是''视角的多样性''。

\textbf{证明}：证明分两个方向。

\textbf{ZFC \(\Rightarrow\) ETCS}。在 ZFC 中构造集合范畴 \(\mathbf{Set}\)：对象为所有集合，态射为所有函数。验证 ETCS 公理：（1）有限极限------\(\mathbf{Set}\) 中有终对象 \(\{\emptyset\}\)（单点集）、拉回（按 ZFC 中的集合论构造）和积（笛卡尔积）；（2）幂对象------\(P(A)\) 是 \(A\) 的幂集，\(\operatorname{Hom}(B, P(A))\) 通过特征函数与 \(\mathbf{Set}/B\) 中 \(B \times A\) 的子对象一一对应，对应关系为 \(f \mapsto \{(b, a) : a \in f(b)\}\)；（3）自然数对象------自然数集 \(\omega\) 满足 Peano 公理，给出 NNO；（4）良幂性------每个集合的子集构成集合（由幂集公理），选择公理在 ZFC 中成立。因此 \(\mathbf{Set}\) 是满足 ETCS 公理的良幂布尔拓扑斯。

\textbf{ETCS + 替换 \(\Rightarrow\) ZFC}。在 ETCS 中，``集合''的概念由对象替代，``元素'' \(x \in A\) 对应态射 \(x: 1 \to A\)（\(1\) 为终对象）。外延公理由来表达：两个对象 \(a, b: 1 \to A\) 相等当且仅当它们作为态射相等。分离公理模式通过子对象分类器 \(\Omega\) 和特征映射实现。替换公理模式需要额外扩充（ETCS 标准版本不含替换），将其添加后可通过范畴内部的''树''构造模拟 ZFC 的累积层级 \(V_\alpha\)。具体地，在 ETCS 中定义 well-founded 树，将 Membership 关系定义为树的''根到子节点''关系，构建满足 ZFC 全部公理的模型。Osius（1974）和 Cole（1973）详细完成了此等价性证明。\(\blacksquare\)

\textbf{定理 175.4}（同伦类型论 HoTT 的基本原理）：Voevodsky 等人发展的同伦类型论（Homotopy Type Theory）将类型论、同伦论和数学基础统一起来。其核心思想是：类型对应于空间，项对应于空间中的点，等式对应于路径，高阶等式对应于路径之间的同伦。在 HoTT 中，Voevodsky 的''单价公理''（Univalence Axiom）断言：等价的结构是相等的------即两个结构如果等价（同构），则它们可以作为''同一个''数学对象来处理。这一公理极大地简化了数学中的''transport''（结构沿着等价传输）操作，但也使得 HoTT 中的等式概念与 ZFC 中的外延性等式有本质区别。HoTT 为数学基础提供了全新的视角，尤其适合形式化同伦论和高阶范畴论中的概念。

\textbf{证明（核心原理的严格推导）}：Martin-Lof 内涵类型论（MLTT）提供了基础语法：每个项 \(a\) 有一个类型 \(A\)（写作 \(a : A\)）。关键创新是对\textbf{等式类型}的解释。在 MLTT 中，对任意 \(x, y : A\)，存在一个等式类型 \(\operatorname{Id}_A(x, y)\)，其项称为 \(x\) 与 \(y\) 之间的''等同证明''。Voevodsky 的洞见是：将 \(A\) 解释为拓扑空间，\(x, y : A\) 解释为空间中的点，\(\operatorname{Id}_A(x, y)\) 解释为 \(x\) 到 \(y\) 的\textbf{路径空间}，\(\operatorname{Id}_{\operatorname{Id}_A(x, y)}(p, q)\) 解释为路径之间的同伦，以此类推。这构成了无穷群胚（\(\infty\)-groupoid）结构。

\textbf{单价公理的推导}。对于类型 \(A, B\)，定义等价类型 \(\operatorname{Equiv}(A, B)\)------当将 \(A, B\) 解释为空间时，其项对应于 \(A\) 与 \(B\) 之间的同伦等价。单价公理（Univalence Axiom）断言典范映射 \[(A =_{\mathcal{U}} B) \to \operatorname{Equiv}(A, B)\] （其中 \(\mathcal{U}\) 为类型宇宙）是等价。其直接推论为：若两个类型同伦等价，则它们作为 \(\mathcal{U}\) 的项是相等的------即类型宇宙中的等式与等价一致。这保证了数学结构沿等价''传输''时的一致性：任何关于 \(A\) 的陈述自动转移到任何与 \(A\) 等价的 \(B\) 上。

\textbf{高阶归纳类型}。HoTT 的创新还包括高阶归纳类型（HITs）------允许在定义类型时同时指定其点（构造子）和路径（等式构造子）。例如，圆周 \(S^1\) 可定义为有一个基点 \(\mathsf{base} : S^1\) 和一条路径 \(\mathsf{loop} : \operatorname{Id}_{S^1}(\mathsf{base}, \mathsf{base})\)。HITs 的归纳原理保证了对 \(S^1\) 的映射完全由基点及其在 \(\mathsf{loop}\) 上的像确定。这使得 HoTT 能以纯类型论的方式表达同伦论的核心构造。

\textbf{与 ZFC 的对比}。在 ZFC 中，\(\mathbb{Z}/2\mathbb{Z}\) 与 \(\{0, 1\}\) 是两个不同的集合（尽管同构），需要显式地沿着同构传输性质。在 HoTT 中，单价公理允许将它们视为相同类型（在 \(\mathcal{U}\) 中相等），从而自动继承所有结构------这被称为''结构不变性原理''。这一差别不是技术性的，而反映了两种基础在''等式的本质''上的深层分歧。\(\blacksquare\)

\hrulefill

\hrulefill

\hrulefill

\hrulefill

\hrulefill

\hrulefill

\textbf{第10章（补充）}

描述集合论（Descriptive Set Theory）研究''可定义''的点集------即可由开集通过可数次布尔运算和投影构造的集合------的规律性（regularity）性质。它连接了集合论、拓扑、测度论和可计算理论，是 Lusin（1910s-1930s）、Souslin、Kleene 和 Moschovakis 等创立的现代分支。描述集合论的核心问题是：在 Polish 空间（完备可分的可度量空间）中，按照集合的''可定义性''层级（Borel 层级、射影层级）来研究集合的正则性质------包括 Lebesgue 可测性、Baire 性质、完美集性质等。描述集合论的一个核心发现是：在 ZFC 中，Borel 集和解析集（\(\mathbf{\Sigma}_1^1\)）都具有良好的正则性，但共解析集（\(\mathbf{\Pi}_1^1\)）以上层级的正则性依赖于大基数公理。这一事实将描述集合论与大基数理论（集合论的另一核心分支）紧密地联系在一起，体现了集合论内部不同分支之间的深刻互动。

\subsection{Polish空间}\label{polishux7a7aux95f4}

\textbf{定义}（Polish空间）：\textbf{Polish空间}是可分的且可完备度量的拓扑空间。例子：\(\mathbb{R}^n\)、可数离散空间、\(\mathbb{N}^{\mathbb{N}}\)（Baire空间）、\(\{0, 1\}^{\mathbb{N}}\)（Cantor空间）、\(C[0, 1]\)（一致范数下）、\([0, 1]^{\mathbb{N}}\)（Hilbert方体）。

Polish空间是描述集合论的基本舞台，其名称来源于波兰数学家学派（Waclaw Sierpinski, Kazimierz Kuratowski, Alfred Tarski 等）在两次世界大战之间对该领域的开创性贡献。选择 Polish 空间作为基本框架有三个关键原因：第一，Polish 空间涵盖了经典分析中几乎所有重要的空间（欧几里得空间、函数空间、可数乘积空间）；第二，Polish 空间在可数操作下具有良好的封闭性（可数乘积、\(G_\delta\) 子空间均为 Polish）；第三，Polish 空间上的 Borel 结构具有''范畴性''------任何两个不可数 Polish 空间作为可测空间是同构的（Borel 同构定理）。这一惊人的事实意味着，从测度论的角度看，所有不可数 Polish 空间本质上是''同一个''对象。Baire 空间 \(\mathbb{N}^{\mathbb{N}}\) 和 Cantor 空间 \(\{0, 1\}^{\mathbb{N}}\) 在描述集合论中扮演着特殊角色------它们分别是零维和完全不连通 Polish 空间的''通用模型''。

\textbf{定理}（Alexandroff定理）：Polish空间的任意 \(G_\delta\) 子空间（即开集的可数交）也是 Polish 空间。特别地，无理数 \(\mathbb{R} \setminus \mathbb{Q}\) 同胚于 \(\mathbb{N}^{\mathbb{N}}\)。

\subsection{Borel层级}\label{borelux5c42ux7ea7}

\textbf{定义}（Borel层级 / Borel hierarchy）：Polish空间 \(X\) 上的 \textbf{Borel层级} 以超限递归定义为： - \(\mathbf{\Sigma}_1^0\) = 开集族 - \(\mathbf{\Pi}_\alpha^0\) = 补集在 \(\mathbf{\Sigma}_\alpha^0\) 中的集合族 - \(\mathbf{\Sigma}_\alpha^0\) = \(\mathbf{\Pi}_\beta^0\) 中集合的可数并（\(\beta < \alpha\)） - \(\mathbf{\Delta}_\alpha^0 = \mathbf{\Sigma}_\alpha^0 \cap \mathbf{\Pi}_\alpha^0\)

对所有 \(\alpha < \omega_1\)，Borel集 = \(\bigcup_{\alpha < \omega_1} \mathbf{\Sigma}_\alpha^0 = \bigcup_{\alpha < \omega_1} \mathbf{\Pi}_\alpha^0\)。Borel层级是真正严格递增的------对每个 \(\alpha\)，存在 \(\mathbf{\Sigma}_\alpha^0 \setminus \mathbf{\Pi}_\alpha^0\) 的集合。

\textbf{定理}（Borel同构定理）：任意两个不可数 Polish 空间的 Borel 结构（作为可测空间）是同构的------即存在保持 Borel 集的双射。特别地，\((\mathbb{R}, \mathcal{B}(\mathbb{R}))\) 与 \((\{0, 1\}^{\mathbb{N}}, \mathcal{B})\) 是 Borel 同构的。因此''Borel集的基数''对任何不可数Borel集都是相同的：\textbf{连续统} \(\mathfrak{c}\)。

\subsection{射影层级}\label{ux5c04ux5f71ux5c42ux7ea7}

\textbf{定义}（射影层级 / projective hierarchy）：从Borel集出发，通过连续像（投影）和补集操作得到的集合层级： - \(\mathbf{\Sigma}_1^1\) = 解析集（analytic sets）\(= \{B \in X : \exists \text{ Polish } Y, \exists \text{ Borel } C \subset X \times Y, B = \pi_X(C)\}\) - \(\mathbf{\Pi}_n^1\) = \(\mathbf{\Sigma}_n^1\) 的补集（共解析集 co-analytic 为 \(\mathbf{\Pi}_1^1\)） - \(\mathbf{\Sigma}_{n+1}^1\) = Polish空间 \(Y\) 上 \(\mathbf{\Pi}_n^1\) 子集的投影（连续像） - \(\mathbf{\Delta}_n^1 = \mathbf{\Sigma}_n^1 \cap \mathbf{\Pi}_n^1\)

\textbf{定理}（Souslin定理 / Souslin分离定理，1917）：若两个不相交的解析集可被Borel集分离，则它们整体满足 \(\mathbf{\Sigma}_1^1\)-分离性。等价地：\(B \in \mathbf{\Delta}_1^1\) 当且仅当 \(B\) 是 Borel 集------即 \(\mathbf{\Sigma}_1^1 \cap \mathbf{\Pi}_1^1 = \mathbf{\Pi}_\omega^0\)（解析且共解析等价于 Borel）。这澄清了 Lebesgue（1905）的错误------他误以为所有解析集是 Borel 的。

\textbf{定理}（Lusin-Sierpiński定理 / 解析集的规律性）： - 所有解析集（\(\mathbf{\Sigma}_1^1\)）是 \textbf{Lebesgue可测} 的 - 所有解析集具有 \textbf{Baire性质}（等于某开集除一个第一范畴集外） - 所有解析集具有 \textbf{完美集性质}（要么可数，要么包含一个完美子集，从而有连续统的基数）

\textbf{推论}：在 ZFC 集合论中（无额外公理），无法证明存在不可测的解析集或没有 Baire 性质的解析集------解析集的规律性在 ZFC 中可证。所有''病态''（不可测）集合只能在 \(\mathbf{\Sigma}_1^1\) 之外（在额外公理或选择公理的强使用下）构造。

\subsection{某些深层联系}\label{ux67d0ux4e9bux6df1ux5c42ux8054ux7cfb}

\textbf{定理}（Moschovakis的 \(\mathbf{\Delta}_2^1\) 可测性）：若所有 \(\mathbf{\Delta}_2^1\) 集是 Lebesgue 可测的，等价于''所有实数推定的'' \(\text{(}{\mathbf{\Sigma}}^1_2\text{-}游戏\text{）}\) ``是决定的（Projective Determinacy的初级层）。这揭示了描述集合论与无穷游戏和大基数公理的深层联系。

\textbf{定理}（Kechris-Martin定理）：若存在可测基数，则所有 \(\mathbf{\Delta}_2^1\) 集是 Lebesgue 可测的。

\textbf{推论}（描述集合论与可计算理论的关系，Kleene 1950s）：有效描述集合论（effective DST）在 Baire空间 \(\omega^\omega\) 上使用递归论定义 \(\Sigma^1_1\) 的光点类（lightface hierarchy），\(\Sigma^0_1\) = 递归可枚举开集。Borel子集的''可计算Borel''（hyperarithmetical）对应 \(\Delta^1_1\)。这揭示了可定义性的精细层次。

\hrulefill

\emph{本卷从数学史和数学哲学两个维度探讨了数学自身的元问题。数学史的核心线索是公理化方法的演进------从 Euclid 的初始范式到非欧几何的逻辑可能性，再到 Hilbert 的元数学纲领和 Bourbaki 的结构主义重建。集合论的创立与悖论的发现揭示了朴素数学直觉的局限性，导出了 ZFC 公理化基础和 Godel 不完备性定理的深刻结论。数学哲学探讨了实在论（柏拉图主义）与反实在论（形式主义、直觉主义）之间的本体论争论，以及结构主义和范畴论为数学基础提供的新视角。选择公理的独立性和计算机证明的认识论地位是当代数学哲学中的活跃议题。}

\hrulefill

\emph{本卷为卷一（数学基础），以 Peano 公理和 ZFC 集合论为起点，建立了整个数学体系的逻辑根基。补充部分从数学史与数学哲学的高度反思了这一根基：公理化方法从 Euclid 到 Hilbert 再到 Bourbaki 的演进，集合论悖论到 ZFC 公理化和 Godel 不完备性定理，以及数学本体论（柏拉图主义、形式主义、直觉主义、结构主义）的深层争论------这些元数学反思为全书各卷的严格公理化叙述提供了方法论框架。}

\hrulefill

\emph{卷一终}

\hrulefill

\hrulefill

\hrulefill

\hrulefill

\hrulefill

\textbf{第11章（补充）}

\begin{quote}
本卷讨论数学自身的元问题------数学史和数学哲学。数学史部分追溯数学思想发展中的逻辑线索（公理化方法的演进、非欧几何的创立、集合论与基础危机、Bourbaki 结构主义等），数学哲学部分探讨数学对象的实在性、数学真理的性质和数学基础的认识论问题。
\end{quote}

\hrulefill

\emph{卷一：数学基础终。}

\emph{卷一：数学基础终。}

\emph{卷一：数学基础终。}

\emph{卷一：数学基础终。}

\newpage

\chapter{01-基础/02-集合与数系/01-集合论.md}\label{ux57faux784002-ux96c6ux5408ux4e0eux6570ux7cfb01-ux96c6ux5408ux8bba.md}

\textbf{前置知识}：依赖 Ch 2（逻辑基础），特别是谓词逻辑（2.2 节）中的量词及其否定、逻辑等价律。集合论用逻辑语言描述数学对象。

在 Ch 2 中我们学习了逻辑推理的规则。现在我们用这套规则来描述数学中最基本的对象------集合。集合是现代数学的基石：几乎所有的数学对象（数、函数、空间等）都可以用集合来精确定义。

\hrulefill

\section{第3章：集合的基本概念}\label{ux7b2c3ux7ae0ux96c6ux5408ux7684ux57faux672cux6982ux5ff5}

集合是数学中最基本的概念之一，几乎所有数学对象都可以用集合来定义。本节从直观层面引入集合的概念，建立属于关系（\(\in\)）和外延公理，然后介绍集合的两种基本表示法（列举法和描述法）。集合的直观概念------``一些确定的、互不相同的对象的整体''------虽然不构成严格的公理化定义，但它为后续的 ZFC 公理化处理提供了直觉基础。理解集合的本质在于把握''外延性''原则：一个集合完全由它的元素决定，元素的排列顺序和重复次数无关紧要。这一原则在 3.1.3 节的外延公理中得到了精确的形式化。本节最后引入空集这一特殊集合，它是集合论构造的起点------从空集出发，通过配对、并集、幂集等操作，可以逐步构造出所有的数学对象。

\subsubsection{3.1.1 集合的直观描述}\label{ux96c6ux5408ux7684ux76f4ux89c2ux63cfux8ff0}

\textbf{定义 3.1.1（集合的直观概念）}：\textbf{集合}是一些确定的、互不相同的对象的整体。这些对象称为集合的\textbf{元素}。

``确定''是指给定一个对象，我们能判断它是否属于该集合。``互不相同''是指集合中每个元素只出现一次。

\textbf{表示法}：

\begin{enumerate}
\def\labelenumi{\arabic{enumi}.}
\tightlist
\item
  \textbf{列举法}：将所有元素列出来，如 \(A = \{1, 2, 3\}\)。
\item
  \textbf{描述法}：\(\{x \mid P(x)\}\) 表示所有满足性质 \(P\) 的对象 \(x\) 组成的集合。
\end{enumerate}

\subsubsection{3.1.2 属于关系}\label{ux5c5eux4e8eux5173ux7cfb}

\textbf{定义 3.1.2（属于关系）}：若 \(a\) 是集合 \(A\) 的一个元素，记为 \(a \in A\)。若 \(a\) 不是 \(A\) 的元素，记为 \(a \notin A\)。

属于关系 \(\in\) 是集合论中最基本的原始关系，它在 ZFC 公理系统中被当作未定义的基本概念接受。值得注意的是，\(\in\) 不是等价关系（它不满足自反性和对称性），也不是序关系。在公理化集合论中，所有数学对象都被视为集合，因此 \(\in\) 是唯一需要的原始二元谓词。符号 \(\in\) 由意大利数学家 Giuseppe Peano 在 1889 年引入，源自希腊字母 \(\varepsilon\)（epsilon），是''est''（拉丁语''是''）的首字母。集合论中还有一个重要的区分：\(a \in A\) 表示 \(a\) 是 \(A\) 的元素，而 \(a \subseteq A\) 表示 \(a\) 是 \(A\) 的子集，两者不可混淆。例如，若 \(A = \{1, \{1\}\}\)，则 \(1 \in A\) 且 \(\{1\} \in A\)，同时 \(\{1\} \subseteq A\) 不成立（因为 \(1\) 是 \(\{1\}\) 的元素但 \(1\) 不是 \(A\) 的元素------\(A\) 的元素是 \(1\) 和 \(\{1\}\)）。在正则公理（见 3.5.8 节）下，不存在集合 \(A\) 使得 \(A \in A\)，这排除了许多''病态''的集合构造。

属于关系 \(\in\) 是 ZFC 集合论中唯一的原始谓词符号，这意味着集合论的全部内容------从自然数到超限基数------都建立在这个单一的关系之上。这种经济性既是集合论作为数学基础的优势，也是其抽象性的体现。在朴素集合论中，\(\in\) 被直观地理解为''是\ldots 的元素''，但在公理化集合论中，\(\in\) 的含义完全由公理规定------这是一种''隐定义''（implicit definition）的策略，与 Hilbert 在《几何基础》中的方法论一脉相承。值得注意的是，\(\in\) 关系是''非良基''的：在没有正则公理的情况下，可以存在满足 \(x \in x\) 的集合，但在 ZFC 中这种可能性被排除了。属于关系 \(\in\) 与子集关系 \(\subseteq\) 之间的区别是初学者最常混淆的概念之一：\(\in\) 连接的是元素和集合，而 \(\subseteq\) 连接的是两个集合。这一区别在集合论的层级结构中至关重要------它使得我们可以区分''元素 \(x\) 属于集合 \(A\)``和''集合 \(A\) 是另一个集合 \(B\) 的子集''这两种不同的关系。

\subsubsection{3.1.3 外延公理}\label{ux5916ux5ef6ux516cux7406}

\textbf{公理 3.1.1（外延公理）}：两个集合 \(A\) 和 \(B\) 相等，当且仅当它们有完全相同的元素： \[A = B \iff \forall x\, (x \in A \leftrightarrow x \in B)\]

外延公理告诉我们：一个集合完全由它的元素决定，与元素的排列顺序和重复次数无关。

外延公理是集合论的第一条公理，它确立了集合的''同一性标准''------两个集合相等当且仅当它们有相同的元素。这一定义看似简单，实则蕴含了深刻的哲学含义：集合的本质完全是外延的（extensional），而非内涵的（intensional）。也就是说，集合 \(\{x \mid x \text{ 是偶数}\}\) 与集合 \(\{x \mid x \text{ 能被 } 2 \text{ 整除}\}\) 是同一个集合，尽管这两个描述（内涵）不同。外延公理排除了集合论中基于内涵定义的歧义性。在 ZFC 公理系统中，外延公理是唯一涉及''相等''关系的公理，其余公理只涉及 \(\in\) 关系。因此，外延公理建立了''相等''与''属于''之间的桥梁。值得注意的是，外延公理并不保证''任意性质都定义一个集合''------这正是 Russell 悖论所揭示的危险（见 3.5.7 节替换公理模式和第 6 章中关于悖论的讨论）。

\subsubsection{3.1.4 空集}\label{ux7a7aux96c6}

\textbf{公理 3.1.2（空集公理）}：存在一个不包含任何元素的集合。

\textbf{定理 3.1.1（空集的唯一性）}：不包含任何元素的集合是唯一的。

\textbf{证明}：设 \(A\) 和 \(B\) 都是不包含任何元素的集合。由外延公理，只需证明 \(\forall x\, (x \in A \leftrightarrow x \in B)\)。对任意 \(x\)：\(x \in A\) 为假（\(A\) 无元素），\(x \in B\) 也为假（\(B\) 无元素），故 \(x \in A \leftrightarrow x \in B\) 为真。由外延公理，\(A = B\)。\(\blacksquare\)

\textbf{定义 3.1.3（空集）}：不包含任何元素的唯一集合称为\textbf{空集}，记为 \(\varnothing\)。

\textbf{注记 3.1.1}：空集是集合论中最基本的对象。它是唯一的不含任何元素的集合。注意 \(\varnothing \neq \{\varnothing\}\)：前者没有元素，后者有一个元素（即空集本身）。类似地，\(\varnothing \in \{\varnothing\}\) 但 \(\varnothing \subseteq \{\varnothing\}\)（空集是任何集合的子集）。

\hrulefill

\section{第4章：集合的运算}\label{ux7b2c4ux7ae0ux96c6ux5408ux7684ux8fd0ux7b97}

在定义了集合的基本概念之后，本节研究如何从已有集合中构造新的集合。集合的运算------并集、交集、差集、补集、对称差------是集合论中最基本的操作，它们与命题逻辑中的联结词（析取、合取、否定）之间存在着精确的对应关系：\(x \in A \cup B \iff (x \in A) \vee (x \in B)\)，\(x \in A \cap B \iff (x \in A) \wedge (x \in B)\)，\(x \in A \setminus B \iff (x \in A) \wedge \neg(x \in B)\)。这种对应揭示了集合论与 Boole 代数之间的同构关系：一个集合的幂集（连同并、交、补运算）构成一个 Boole 代数。本节还将介绍集合的运算律（结合律、交换律、分配律、De Morgan 律），这些运算律可以通过真值表或元素论证来验证。此外，我们将引入笛卡尔积这一重要构造，它使得我们可以从两个集合中构造出有序对的集合，为后续的关系和函数理论奠定基础。笛卡尔积的命名来源于 Rene Descartes，他最早将平面上的点与实数对 \((x, y)\) 对应起来，从而建立了代数与几何之间的桥梁。

\subsubsection{3.2.1 子集}\label{ux5b50ux96c6}

\textbf{定义 3.2.1（子集）}：若集合 \(A\) 的每个元素都是集合 \(B\) 的元素，则称 \(A\) 是 \(B\) 的\textbf{子集}，记为 \(A \subseteq B\)： \[A \subseteq B \iff \forall x\, (x \in A \rightarrow x \in B)\]

\textbf{定义 3.2.2（真子集）}：若 \(A \subseteq B\) 且 \(A \neq B\)，则称 \(A\) 是 \(B\) 的\textbf{真子集}，记为 \(A \subsetneq B\)。

\textbf{注记 3.2.1}：\(A \subsetneq B\) 等价于 \(A \subseteq B\) 且 \(\exists x \in B \setminus A\)。初学者容易混淆 \(\subseteq\) 和 \(\in\)：\(1 \in \{1, 2\}\) 但 \(\{1\} \subseteq \{1, 2\}\)。\(1\) 是元素，\(\{1\}\) 是子集。

\textbf{定理 3.2.1（子集关系的性质）}：对任意集合 \(A, B, C\)： 1. \(A \subseteq A\)（自反性） 2. \(A \subseteq B\) 且 \(B \subseteq A\) 则 \(A = B\)（反对称性） 3. \(A \subseteq B\) 且 \(B \subseteq C\) 则 \(A \subseteq C\)（传递性）

\textbf{证明}：

\textbf{（1）}：对任意 \(x\)，\(x \in A \rightarrow x \in A\) 恒为真，故 \(A \subseteq A\)。\(\blacksquare\)

\textbf{（2）}：设 \(A \subseteq B\) 且 \(B \subseteq A\)。对任意 \(x\)：由 \(A \subseteq B\)，\(x \in A \rightarrow x \in B\)；由 \(B \subseteq A\)，\(x \in B \rightarrow x \in A\)。因此 \(x \in A \leftrightarrow x \in B\)，由外延公理 \(A = B\)。\(\blacksquare\)

\textbf{（3）}：设 \(A \subseteq B\) 且 \(B \subseteq C\)。对任意 \(x\)：若 \(x \in A\)，由 \(A \subseteq B\) 得 \(x \in B\)，再由 \(B \subseteq C\) 得 \(x \in C\)。因此 \(x \in A \rightarrow x \in C\)，即 \(A \subseteq C\)。\(\blacksquare\)

\textbf{定理 3.2.2}：空集是任何集合的子集，即 \(\varnothing \subseteq A\) 对任意集合 \(A\) 成立。

\textbf{证明}：需证明 \(\forall x\, (x \in \varnothing \rightarrow x \in A)\)。对任意 \(x\)，\(x \in \varnothing\) 为假，因此 \(x \in \varnothing \rightarrow x \in A\) 为真（前件为假时蕴含式恒真）。\(\blacksquare\)

\subsubsection{3.2.2 并集与交集}\label{ux5e76ux96c6ux4e0eux4ea4ux96c6}

\textbf{定义 3.2.3（并集）}：\(A \cup B = \{x \mid x \in A \vee x \in B\}\)

\textbf{定义 3.2.4（交集）}：\(A \cap B = \{x \mid x \in A \wedge x \in B\}\)

\textbf{定义 3.2.5（不相交）}：若 \(A \cap B = \varnothing\)，则称 \(A\) 和 \(B\) \textbf{不相交}。

\textbf{定义 3.2.6（差集与补集）}： - \textbf{差集}：\(A \setminus B = \{x \mid x \in A \wedge x \notin B\}\) - \textbf{补集}：当 \(B \subseteq A\) 时，\(A \setminus B\) 也称为 \(B\) 在 \(A\) 中的补集，记为 \(\complement_{A} B\)。

并集和交集是集合论中最基本的二元运算。它们与逻辑联结词存在精确的对应关系：并集对应于析取（\(\vee\)），交集对应于合取（\(\wedge\)），补集对应于否定（\(\neg\)）。因此，集合运算的代数性质直接继承自命题逻辑的运算律（交换律、结合律、分配律、德摩根律等，详见 2.1.4 节）。这一对应关系构成了 Boole 代数的集合论模型。在公理化框架中，二元并集 \(\{a, b\}\) 的存在性由配对公理和并集公理联合保证（见 3.5.3 和 3.5.4 节）。差集 \(A \setminus B\) 可以通过分离公理模式从 \(A\) 中构造：\(A \setminus B = \{x \in A \mid x \notin B\}\)。补集的概念依赖于所选的''全集''（论域）\(U\)，在 ZFC 中不存在包含所有集合的''万有集''（否则将导致 Russell 悖论），因此补集总是在某个特定集合内部讨论的。

\subsubsection{3.2.3 幂集}\label{ux5e42ux96c6}

\textbf{定义 3.2.7（幂集）}：集合 \(A\) 的\textbf{幂集}定义为 \(A\) 的所有子集组成的集合： \[\mathcal{P}(A) = \{S \mid S \subseteq A\}\]

\textbf{定理 3.2.3}：若 \(A\) 有 \(n\) 个元素，则 \(\mathcal{P}(A)\) 有 \(2^n\) 个元素。

幂集运算是集合论中产生''更大''集合的基本手段。康托尔定理（定理 3.6.5）证明了对任意集合 \(A\)，\(|\mathcal{P}(A)| > |A|\)，这意味着不存在最大的基数。幂集运算因此产生了无穷基数的严格层级：\(\aleph_0 = |\mathbb{N}|\)，\(2^{\aleph_0} = |\mathcal{P}(\mathbb{N})| = |\mathbb{R}|\)，\(2^{2^{\aleph_0}} = |\mathcal{P}(\mathbb{R})|\)，等等。幂集的存在性由幂集公理（公理 3.5.5）保证。在范畴论中，幂集构造对应于子对象分类器（subobject classifier）的概念，是拓扑斯（topos）理论的核心组成部分。符号 \(\mathcal{P}(A)\) 中的花体 P 源自德语 ``Potenzmenge''（幂集）。

\textbf{证明}：对 \(n\) 用数学归纳法。

\textbf{基础步骤}（\(n = 0\)）：\(A = \varnothing\)，\(\mathcal{P}(\varnothing) = \{\varnothing\}\) 有 \(2^0 = 1\) 个元素。成立。

\textbf{归纳步骤}：设对 \(n\) 个元素的集合，幂集有 \(2^n\) 个元素。设 \(A\) 有 \(n + 1\) 个元素，取 \(a \in A\)，令 \(B = A \setminus \{a\}\)，则 \(B\) 有 \(n\) 个元素。

\(\mathcal{P}(A)\) 中的集合可分为两类： - 不包含 \(a\) 的子集：恰好是 \(B\) 的子集，共 \(2^n\) 个。 - 包含 \(a\) 的子集：形如 \(S \cup \{a\}\)（\(S \subseteq B\)），也有 \(2^n\) 个。

总计 \(2^n + 2^n = 2^{n+1}\)。由归纳法得证。\(\blacksquare\)

\subsubsection{3.2.4 笛卡尔积}\label{ux7b1bux5361ux5c14ux79ef}

\textbf{定义 3.2.8（有序对）}：有序对 \((a, b)\) 定义为 \(\{\{a\}, \{a, b\}\}\)（Kuratowski 定义）。

\textbf{引理 3.2.1}：\((a, b) = (c, d)\) 当且仅当 \(a = c\) 且 \(b = d\)。

\textbf{证明}：

\((\Leftarrow)\)：显然。

\((\Rightarrow)\)：设 \((a, b) = (c, d)\)，即 \(\{\{a\}, \{a, b\}\} = \{\{c\}, \{c, d\}\}\)。

若 \(a = b\)，则 \(\{\{a\}, \{a, b\}\} = \{\{a\}\}\)。由集合相等，\(\{\{c\}, \{c, d\}\} = \{\{a\}\}\)，故 \(\{c\} = \{c, d\}\)，得 \(c = d\) 且 \(\{a\} = \{c\}\)，即 \(a = c\)，\(b = a = c = d\)。

若 \(a \neq b\)，则 \(\{\{a\}, \{a, b\}\}\) 有两个不同元素。由 \(\{a\} \in \{\{c\}, \{c, d\}\}\)，有 \(\{a\} = \{c\}\) 或 \(\{a\} = \{c, d\}\)。

若 \(\{a\} = \{c, d\}\)，则 \(c = d = a\)。再由 \(\{a, b\} \in \{\{a\}\}\)，得 \(\{a, b\} = \{a\}\)，矛盾（\(a \neq b\)）。

若 \(\{a\} = \{c\}\)，则 \(a = c\)。由 \(\{a, b\} = \{a, d\}\)，得 \(b = d\)。\(\blacksquare\)

\textbf{定义 3.2.9（笛卡尔积）}： \[A \times B = \{(a, b) \mid a \in A,\, b \in B\}\]

笛卡尔积是构造多元数学对象的基本工具。在数学中，几乎所有涉及多个对象的结构都基于笛卡尔积：二元关系是 \(A \times B\) 的子集，函数是特殊的二元关系，\(\mathbb{R}^2 = \mathbb{R} \times \mathbb{R}\) 是平面，\(\mathbb{R}^n = \mathbb{R} \times \cdots \times \mathbb{R}\) 是 \(n\) 维欧几里得空间。笛卡尔积的命名来源于法国数学家和哲学家 Rene Descartes（笛卡尔），他首次用坐标 \((x, y)\) 表示平面上的点，建立了代数与几何的联系。在集合论中，笛卡尔积的存在性可由配对公理、并集公理和替换公理模式联合保证。笛卡尔积不满足交换律：\(A \times B \neq B \times A\)（除非 \(A = B\) 或其中之一为空集），也不满足结合律：\((A \times B) \times C \neq A \times (B \times C)\)，但两者之间存在自然的双射。在范畴论中，笛卡尔积是集合范畴中的积（product），满足泛性质（universal property）。

\subsubsection{3.2.5 集合运算的基本性质}\label{ux96c6ux5408ux8fd0ux7b97ux7684ux57faux672cux6027ux8d28}

\textbf{定理 3.2.4（交换律）}：对任意集合 \(A, B\)： 1. \(A \cup B = B \cup A\) 2. \(A \cap B = B \cap A\)

\textbf{证明}： (1) \(x \in A \cup B \iff x \in A \vee x \in B \iff x \in B \vee x \in A \iff x \in B \cup A\)（析取交换律，定理 2.1.2）。 由外延公理得证。\(\blacksquare\)

\textbf{定理 3.2.5（结合律）}：对任意集合 \(A, B, C\)： 1. \((A \cup B) \cup C = A \cup (B \cup C)\) 2. \((A \cap B) \cap C = A \cap (B \cap C)\)

\textbf{证明}： (1) \(x \in (A \cup B) \cup C \iff (x \in A \vee x \in B) \vee x \in C \iff x \in A \vee (x \in B \vee x \in C) \iff x \in A \cup (B \cup C)\) 中间使用析取结合律（定理 2.1.3）。\(\blacksquare\)

\textbf{定理 3.2.6（分配律）}：对任意集合 \(A, B, C\)： 1. \(A \cup (B \cap C) = (A \cup B) \cap (A \cup C)\) 2. \(A \cap (B \cup C) = (A \cap B) \cup (A \cap C)\)

\textbf{证明}： (1) \(x \in A \cup (B \cap C) \iff x \in A \vee (x \in B \wedge x \in C) \iff (x \in A \vee x \in B) \wedge (x \in A \vee x \in C) \iff x \in (A \cup B) \cap (A \cup C)\) 使用析取对合取的分配律（定理 2.1.4）。\(\blacksquare\)

\textbf{定理 3.2.7（德摩根律）}：对任意集合 \(A, B\)（在论域 \(U\) 中考虑补集）： 1. \(U \setminus (A \cup B) = (U \setminus A) \cap (U \setminus B)\) 2. \(U \setminus (A \cap B) = (U \setminus A) \cup (U \setminus B)\)

\textbf{证明}： (1) \(x \in U \setminus (A \cup B) \iff x \notin A \cup B \iff \neg(x \in A \vee x \in B) \iff \neg(x \in A) \wedge \neg(x \in B) \iff x \notin A \wedge x \notin B \iff x \in (U \setminus A) \cap (U \setminus B)\) 使用逻辑德摩根律（定理 2.1.5）。\(\blacksquare\)

\textbf{定理 3.2.8（幂等律）}：\(A \cup A = A\)，\(A \cap A = A\)。

\textbf{证明}：\(x \in A \cup A \iff x \in A \vee x \in A \iff x \in A\)。\(\blacksquare\)

\textbf{定理 3.2.9（同一律）}： 1. \(A \cup \varnothing = A\) 2. \(A \cap \varnothing = \varnothing\) 3. \(A \cup U = U\) 4. \(A \cap U = A\)

\textbf{证明}：(1) \(x \in A \cup \varnothing \iff x \in A \vee x \in \varnothing \iff x \in A \vee F \iff x \in A\)。\(\blacksquare\)

\textbf{定理 3.2.10（补集的基本性质）}：对任意集合 \(A \subseteq U\)： 1. \(A \cup (U \setminus A) = U\) 2. \(A \cap (U \setminus A) = \varnothing\) 3. \(U \setminus (U \setminus A) = A\)（双重补集）

\textbf{证明}： (1) 若 \(x \in U\)，则 \(x \in A\) 或 \(x \notin A\)（排中律），故 \(x \in A \cup (U \setminus A)\)。\(\blacksquare\)

\textbf{定理 3.2.11（笛卡尔积的性质）}： 1. \(A \times (B \cup C) = (A \times B) \cup (A \times C)\) 2. \(A \times (B \cap C) = (A \times B) \cap (A \times C)\)

\textbf{证明}： (1) \((x, y) \in A \times (B \cup C) \iff x \in A \wedge (y \in B \vee y \in C) \iff (x \in A \wedge y \in B) \vee (x \in A \wedge y \in C) \iff (x, y) \in A \times B \vee (x, y) \in A \times C\)。\(\blacksquare\)

\hrulefill

\section{第5章：关系}\label{ux7b2c5ux7ae0ux5173ux7cfb}

关系是数学中仅次于集合和函数的第三个基本概念，它从集合论的角度形式化了对象之间的各种联系。本节将关系定义为笛卡尔积的子集，这一简洁的定义统一了数学中所有类型的关系：序关系（\(\leq\)）、等价关系（\(\equiv\)）、图论中的边关系、代数中的同余关系等。在朴素集合论中，一个关系就是任意一组有序对的集合；在 ZFC 中，有序对 \((a, b)\) 本身被定义为 \(\{\{a\}, \{a, b\}\}\)（Kuratowski 对），从而将关系也归结为集合。关系理论的核心在于研究关系的性质：自反性、对称性、传递性、反对称性等。这些性质的组合定义了两种最重要的关系类型：等价关系（自反+对称+传递）和偏序关系（自反+反对称+传递）------前者用于分类和商构造，后者用于比较和排序。本节最后将介绍函数的集合论定义，将函数视为一种满足''单值性''的特殊关系，从而完成从集合到关系再到函数的逻辑链条。

\subsubsection{3.3.1 二元关系}\label{ux4e8cux5143ux5173ux7cfb}

\textbf{定义 3.3.1（二元关系）}：集合 \(A\) 到集合 \(B\) 的一个\textbf{二元关系}是笛卡尔积 \(A \times B\) 的一个子集。

即 \(R \subseteq A \times B\)。若 \((a, b) \in R\)，通常写为 \(a\, R\, b\)。

\textbf{注记 3.3.1}：二元关系是集合论中描述对象之间联系的基本工具。``小于''关系 \(<\) 可以定义为 \(\{(a, b) \in \mathbb{Z} \times \mathbb{Z} : b - a > 0\}\)，``等于''关系可以定义为对角线 \(\{(a, a) : a \in A\}\)。

二元关系的概念极其普遍，它统一了数学中许多看似不同的结构。``\(x\) 整除 \(y\)''（\(x \mid y\)）是 \(\mathbb{N}\) 上的二元关系，函数的图像 \(\{(x, f(x)) : x \in A\}\) 是特殊的二元关系，图论中的边集 \(E \subseteq V \times V\) 也是二元关系。关系的定义域（domain）和值域（range）分别定义为 \(\operatorname{dom}(R) = \{a \in A : \exists b \in B, (a, b) \in R\}\) 和 \(\operatorname{ran}(R) = \{b \in B : \exists a \in A, (a, b) \in R\}\)。关系的逆 \(R^{-1} = \{(b, a) : (a, b) \in R\}\) 将方向反转（例如，\(<\) 的逆是 \(>\)）。两个关系 \(R \subseteq A \times B\) 和 \(S \subseteq B \times C\) 的复合定义为 \(S \circ R = \{(a, c) \in A \times C : \exists b \in B, (a, b) \in R \wedge (b, c) \in S\}\)。关系复合满足结合律，这与函数复合（定义 3.4.3）的性质一致。

\textbf{定义 3.3.2（关系的性质）}：设 \(R\) 是集合 \(A\) 上的二元关系（即 \(R \subseteq A \times A\)）： - \textbf{自反性}：\(\forall a \in A,\, a\, R\, a\) - \textbf{对称性}：\(\forall a, b \in A,\, a\, R\, b \rightarrow b\, R\, a\) - \textbf{反对称性}：\(\forall a, b \in A,\, (a\, R\, b \wedge b\, R\, a) \rightarrow a = b\) - \textbf{传递性}：\(\forall a, b, c \in A,\, (a\, R\, b \wedge b\, R\, c) \rightarrow a\, R\, c\)

\subsubsection{3.3.2 等价关系}\label{ux7b49ux4ef7ux5173ux7cfb}

\textbf{定义 3.3.3（等价关系）}：集合 \(A\) 上的一个\textbf{等价关系}是同时满足自反性、对称性和传递性的二元关系，通常记为 \(\sim\)。

\textbf{定义 3.3.4（等价类）}：设 \(\sim\) 是集合 \(A\) 上的等价关系，\(a \in A\)，则 \(a\) 的\textbf{等价类}定义为： \[[a] = \{x \in A \mid x \sim a\}\]

\textbf{定理 3.3.1（等价类的基本性质）}：设 \(\sim\) 是集合 \(A\) 上的等价关系，则： 1. \(\forall a \in A,\, a \in [a]\) 2. \([a] = [b] \iff a \sim b\) 3. \([a] \neq [b] \iff [a] \cap [b] = \varnothing\)

\textbf{证明}：

\textbf{（1）}：由自反性，\(a \sim a\)，故 \(a \in [a]\)。\(\blacksquare\)

\textbf{（2）}： \((\Rightarrow)\)：设 \([a] = [b]\)。由 (1)，\(b \in [b] = [a]\)，故 \(b \sim a\)。由对称性，\(a \sim b\)。\(\blacksquare\)

\((\Leftarrow)\)：设 \(a \sim b\)。

先证 \([a] \subseteq [b]\)：设 \(x \in [a]\)，即 \(x \sim a\)。由传递性，\(x \sim a\) 且 \(a \sim b\) 推出 \(x \sim b\)，故 \(x \in [b]\)。

再证 \([b] \subseteq [a]\)：设 \(x \in [b]\)，即 \(x \sim b\)。由对称性，\(a \sim b\) 推出 \(b \sim a\)。由传递性，\(x \sim b\) 且 \(b \sim a\) 推出 \(x \sim a\)，故 \(x \in [a]\)。

因此 \([a] = [b]\)。\(\blacksquare\)

\textbf{（3）}：需证双向命题 \([a] \neq [b] \iff [a] \cap [b] = \varnothing\)。

\((\Leftarrow)\)（反方向）：假设 \([a] \cap [b] \neq \varnothing\)，需证 \([a] = [b]\)。设 \(c \in [a] \cap [b]\)，即 \(c \sim a\) 且 \(c \sim b\)。由对称性，\(a \sim c\)。由传递性，\(a \sim c\) 且 \(c \sim b\) 推出 \(a \sim b\)。由 (2)，\([a] = [b]\)。此即 \([a] \cap [b] \neq \varnothing \Rightarrow [a] = [b]\)，其逆否命题为 \([a] \neq [b] \Rightarrow [a] \cap [b] = \varnothing\)。

\((\Rightarrow)\)（正方向）：假设 \([a] = [b]\)，需证 \([a] \cap [b] \neq \varnothing\)。由 (1)，\(a \in [a]\)，故 \(a \in [a] = [b]\)，从而 \(a \in [a] \cap [b]\)，即 \([a] \cap [b] \neq \varnothing\)。其逆否命题为 \([a] \cap [b] = \varnothing \Rightarrow [a] \neq [b]\)。

综合两个方向，\([a] \neq [b] \iff [a] \cap [b] = \varnothing\)。\(\blacksquare\)

\textbf{定义 3.3.5（商集）}：\(A\) 关于 \(\sim\) 的\textbf{商集}定义为所有等价类的集合： \[A / \sim = \{[a] \mid a \in A\}\]

\textbf{定理 3.3.2（商集是划分）}：商集 \(A / \sim\) 构成 \(A\) 的一个\textbf{划分}，即： 1. 每个等价类非空； 2. 不同的等价类不相交； 3. 所有等价类的并等于 \(A\)。

\textbf{证明}： (1) 由定理 3.3.1(1)，\(a \in [a]\)，故 \([a] \neq \varnothing\)。 (2) 由定理 3.3.1(3)。 (3) \((\subseteq)\)：若 \(x \in \bigcup_{a \in A} [a]\)，则存在 \(a \in A\) 使得 \(x \in [a] \subseteq A\)。\((\supseteq)\)：若 \(x \in A\)，则 \(x \in [x]\)。\(\blacksquare\)

\hrulefill

\section{第6章：函数}\label{ux7b2c6ux7ae0ux51fdux6570}

函数是数学中最核心的概念之一，它贯穿了从初等代数到泛函分析的所有数学分支。本节从集合论的角度给出函数的精确定义：函数 \(f: A \to B\) 是笛卡尔积 \(A \times B\) 的一个子集，满足''每个 \(a \in A\) 恰好对应一个 \(b \in B\)``的单值性条件。这一定义将函数完全归结为集合，是 ZFC 集合论作为数学基础的一个典型范例。函数的集合论定义还自然地引出了三个核心概念：定义域（domain）、值域（codomain）和像集（image）。在此基础上，本节将分类讨论三种基本函数类型------单射（injective）、满射（surjective）和双射（bijective）------它们分别对应不同的映射性质，并在逆映射、基数比较和同构理论中扮演关键角色。复合运算（\(\circ\)）赋予函数以代数结构，使得从集合到自身的函数构成一个幺半群。本节最后将介绍函数空间（\(B^A\)）的概念，这一构造为后续的基数理论、泛函分析和范畴论提供了基础。

\subsubsection{3.4.1 函数的定义}\label{ux51fdux6570ux7684ux5b9aux4e49}

\textbf{定义 3.4.1（函数）}：设 \(A, B\) 是集合。一个从 \(A\) 到 \(B\) 的\textbf{函数} \(f: A \to B\) 是满足以下条件的二元关系 \(f \subseteq A \times B\)： 1. \textbf{存在性}：\(\forall a \in A,\, \exists b \in B,\, (a, b) \in f\) 2. \textbf{唯一性}：\(\forall a \in A,\, \forall b_1, b_2 \in B,\, ((a, b_1) \in f \wedge (a, b_2) \in f) \rightarrow b_1 = b_2\)

若 \((a, b) \in f\)，记 \(b = f(a)\)。\(A\) 称为 \(f\) 的\textbf{定义域}，\(B\) 称为\textbf{陪域}。\(\{f(a) \mid a \in A\}\) 称为 \(f\) 的\textbf{值域}。

函数是数学中最核心的概念之一。在集合论的框架下，函数被还原为一种特殊的二元关系，这一定义完全消除了函数概念中的任何''过程''或''规则''的直觉成分，仅保留了输入-输出对应的外延特征。两个函数 \(f, g: A \to B\) 相等当且仅当 \(\forall a \in A, f(a) = g(a)\)（这由外延公理和函数的定义直接推出）。函数定义的''存在性''条件保证了定义域中每个元素都有像，``唯一性''条件保证了每个元素只有一个像（函数不能''一对多''）。注意陪域 \(B\) 和值域的区别：值域是 \(B\) 的子集，函数不一定将 \(A\) 的每个元素映射到 \(B\) 的每个元素。例如 \(f: \mathbb{R} \to \mathbb{R}\)，\(f(x) = x^2\) 的值域是 \([0, \infty)\)，但陪域是 \(\mathbb{R}\)。在范畴论中，函数是集合范畴中的态射（morphism），而集合范畴是许多数学结构的底层范畴。

\subsubsection{3.4.2 单射、满射、双射}\label{ux5355ux5c04ux6ee1ux5c04ux53ccux5c04}

\textbf{定义 3.4.2}：设 \(f: A \to B\) 是函数。 - \(f\) 是\textbf{单射}（injective）：\(\forall a_1, a_2 \in A,\, f(a_1) = f(a_2) \rightarrow a_1 = a_2\) - \(f\) 是\textbf{满射}（surjective）：\(\forall b \in B,\, \exists a \in A,\, f(a) = b\) - \(f\) 是\textbf{双射}（bijective）：\(f\) 既是单射又是满射。

\textbf{定理 3.4.1（单射的等价刻画）}：\(f: A \to B\) 是单射，当且仅当 \(\forall a_1, a_2 \in A,\, a_1 \neq a_2 \rightarrow f(a_1) \neq f(a_2)\)。

\textbf{证明}：这是单射定义的逆否命题。由定理 2.1.10，\(P \rightarrow Q \equiv \neg Q \rightarrow \neg P\)。\(\blacksquare\)

单射、满射和双射是函数分类的三个基本范畴，它们刻画了函数在结构上的不同''完善程度''。单射保证了''不同的输入给出不同的输出''------即函数不会将不同元素''压缩''到同一个像，因此单射保留了定义域的''区分度''。满射保证了陪域中每个元素都有原像------即函数''覆盖''了整个陪域。双射则同时满足两者，建立了定义域和陪域之间的一一对应。在有限集合的情况下，单射、满射和双射之间存在更深层的关联：若 \(A\) 和 \(B\) 是有限集且 \(|A| = |B|\)，则 \(f: A \to B\) 是单射当且仅当 \(f\) 是满射当且仅当 \(f\) 是双射（抽屉原理）。但在无穷集合上这一等价性不成立：例如 \(f: \mathbb{N} \to \mathbb{N}\)，\(f(n) = 2n\) 是单射但不是满射；\(g: \mathbb{N} \to \mathbb{N}\)，\(g(n) = \lfloor n/2 \rfloor\) 是满射但不是单射。双射在集合论中具有特殊地位，因为它定义了集合的''等势''关系（见 3.6 节），是康托尔基数理论的基础。

\subsubsection{3.4.3 复合函数}\label{ux590dux5408ux51fdux6570}

\textbf{定义 3.4.3（复合函数）}：设 \(f: A \to B\) 和 \(g: B \to C\) 是函数，则 \(g\) 和 \(f\) 的\textbf{复合函数} \(g \circ f: A \to C\) 定义为： \[(g \circ f)(a) = g(f(a)), \quad \forall a \in A\]

\textbf{注记 3.4.1}：注意复合的书写顺序：\(g \circ f\) 表示先应用 \(f\)，再应用 \(g\)（从右向左阅读）。这有时会让人困惑，因为 \(f\) 写在右边却先被应用。

\textbf{定理 3.4.2（复合保持单射和满射）}：设 \(f: A \to B\)，\(g: B \to C\)。 1. 若 \(f\) 和 \(g\) 都是单射，则 \(g \circ f\) 也是单射。 2. 若 \(f\) 和 \(g\) 都是满射，则 \(g \circ f\) 也是满射。 3. 若 \(f\) 和 \(g\) 都是双射，则 \(g \circ f\) 也是双射。

\textbf{证明}：

\textbf{（1）}：设 \((g \circ f)(a_1) = (g \circ f)(a_2)\)，即 \(g(f(a_1)) = g(f(a_2))\)。因 \(g\) 是单射，得 \(f(a_1) = f(a_2)\)。因 \(f\) 是单射，得 \(a_1 = a_2\)。\(\blacksquare\)

\textbf{（2）}：设 \(c \in C\)。因 \(g\) 是满射，存在 \(b \in B\) 使得 \(g(b) = c\)。因 \(f\) 是满射，存在 \(a \in A\) 使得 \(f(a) = b\)。于是 \((g \circ f)(a) = g(f(a)) = g(b) = c\)。\(\blacksquare\)

\textbf{（3）}：由 (1) 和 (2)。\(\blacksquare\)

\textbf{注记 3.4.2（逆命题不成立）}：\(g \circ f\) 是单射不能推出 \(f\) 和 \(g\) 都是单射。例如 \(f: \{1\} \to \{1, 2\}\)，\(f(1) = 1\)（单射但不满射）；\(g: \{1, 2\} \to \{1\}\)，\(g(1) = g(2) = 1\)（满射但不是单射）。\(g \circ f: \{1\} \to \{1\}\)，\(g(f(1)) = 1\) 是双射。但可以证明：若 \(g \circ f\) 是单射，则 \(f\) 是单射；若 \(g \circ f\) 是满射，则 \(g\) 是满射。

\textbf{定理 3.4.3（复合的结合律）}：设 \(f: A \to B\)，\(g: B \to C\)，\(h: C \to D\)，则 \[h \circ (g \circ f) = (h \circ g) \circ f\]

\textbf{证明}：对任意 \(a \in A\)： \[(h \circ (g \circ f))(a) = h((g \circ f)(a)) = h(g(f(a)))\] \[((h \circ g) \circ f)(a) = (h \circ g)(f(a)) = h(g(f(a)))\] 两者相等。\(\blacksquare\)

\subsubsection{3.4.4 反函数}\label{ux53cdux51fdux6570}

\textbf{定义 3.4.4（反函数）}：设 \(f: A \to B\) 是双射，则 \(f\) 的\textbf{反函数} \(f^{-1}: B \to A\) 定义为：对每个 \(b \in B\)，\(f^{-1}(b)\) 是唯一的 \(a \in A\) 使得 \(f(a) = b\)。

\textbf{注记 3.4.3}：符号 \(f^{-1}\) 不是指 \(\frac{1}{f}\)，而是表示反函数。双射保证了这样的 \(a\) 存在（由满射性）且唯一（由单射性）。

\textbf{定理 3.4.4（反函数的基本性质）}：设 \(f: A \to B\) 是双射，\(f^{-1}\) 是其反函数，则： 1. \(f^{-1} \circ f = \mathrm{id}_{A}\) 2. \(f \circ f^{-1} = \mathrm{id}_{B}\) 3. \(f^{-1}\) 也是双射 4. \((f^{-1})^{-1} = f\)

\textbf{证明}：

\textbf{（1）}：对任意 \(a \in A\)，设 \(b = f(a)\)。由反函数定义，\(f^{-1}(b) = a\)。故 \((f^{-1} \circ f)(a) = f^{-1}(f(a)) = a\)。\(\blacksquare\)

\textbf{（2）}：对任意 \(b \in B\)，设 \(a = f^{-1}(b)\)。由反函数定义，\(f(a) = b\)。故 \((f \circ f^{-1})(b) = f(f^{-1}(b)) = b\)。\(\blacksquare\)

\textbf{（3）}：\(f^{-1}\) 是单射：若 \(f^{-1}(b_1) = f^{-1}(b_2) = a\)，则 \(f(a) = b_1\) 且 \(f(a) = b_2\)，故 \(b_1 = b_2\)。

\(f^{-1}\) 是满射：对任意 \(a \in A\)，取 \(b = f(a)\)，则 \(f^{-1}(b) = a\)。\(\blacksquare\)

\textbf{（4）}：由 (1) 和 (2) 的对称性，\(f\) 是 \(f^{-1}\) 的反函数。\(\blacksquare\)

\textbf{定理 3.4.5（双射的等价刻画）}：函数 \(f: A \to B\) 是双射，当且仅当存在函数 \(g: B \to A\) 使得 \(g \circ f = \mathrm{id}_{A}\) 且 \(f \circ g = \mathrm{id}_{B}\)。

\textbf{证明}：

\((\Rightarrow)\)：取 \(g = f^{-1}\)，由定理 3.4.4 得证。

\((\Leftarrow)\)：设存在这样的 \(g\)。\(f\) 是单射：若 \(f(a_1) = f(a_2)\)，则 \(a_1 = g(f(a_1)) = g(f(a_2)) = a_2\)。\(f\) 是满射：对任意 \(b \in B\)，取 \(a = g(b)\)，则 \(f(a) = f(g(b)) = b\)。\(\blacksquare\)

\hrulefill

\section{第7章：ZFC公理系统简介}\label{ux7b2c7ux7ae0zfcux516cux7406ux7cfbux7edfux7b80ux4ecb}

ZFC（Zermelo-Fraenkel with Choice）公理系统是当代数学的标准基础，它将集合论建立在严格的形式公理之上，从而避免了朴素集合论中的悖论（如 Russell 悖论）。ZFC 由 Ernst Zermelo 在 1908 年提出，后经 Abraham Fraenkel 和 Thoralf Skolem 在 1922 年完善，加上选择公理（Axiom of Choice, AC）后形成了今天通用的 ZFC 系统。ZFC 的核心理念是''限制大小''------不承认''所有集合的集合''或''所有不包含自身的集合的集合''这类过大构造的存在，而是通过一组精心设计的公理，从空集出发逐步构造出所有合法的集合。本节将逐个介绍 ZFC 的十条公理（或公理模式），每条公理对应一种合法的集合构造方式：外延公理确保集合由其元素唯一确定，空集公理提供了一个起点，配对公理允许构造二元集，并集公理允许取集合的并，幂集公理允许取所有子集，分离公理模式限制了对大集合的''切割''，替换公理模式允许通过函数''拉伸''集合，无穷公理保证了无穷集的存在，正则公理排除了病态循环，选择公理则保证了选择函数的存在。这些公理共同构成了一个足够强大（足以形式化所有经典数学）又足够安全（至今未发现矛盾）的数学基础。

\subsubsection{3.5.1 外延公理}\label{ux5916ux5ef6ux516cux7406-1}

\textbf{公理（Axiom of Extensionality）}：\(\forall A\,\forall B\,(\forall x\,(x \in A \leftrightarrow x \in B) \rightarrow A = B)\)。

外延公理是 ZFC 中唯一涉及''相等''概念的公理，它断言两个集合相等当且仅当它们拥有完全相同的元素。换言之，一个集合完全由它的元素（外延）决定，与集合的''内涵''（定义方式）无关。这一原则在哲学上被称为''外延性原理''（principle of extensionality），它使得集合论中的相等概念有了一个纯粹基于 \(\in\) 关系的判定标准。外延公理的一个重要推论是：集合中元素的排列顺序和重复次数无关紧要------\(\{a, b\}\) 和 \(\{b, a\}\) 是同一个集合，\(\{a, a\}\) 就是 \(\{a\}\)。外延公理还隐含了''良定''的概念：如果我们要定义一个新集合，只需说明哪些元素属于它，而不必说明它''是什么''。这一点在分离公理中得到了充分利用。历史上，外延公理最早由 Gottlob Frege 在其《概念文字》中隐含地提出，后在 Zermelo 的 1908 年论文中被明确列为集合论的第一条公理。外延公理与内涵公理模式（或分离公理）的结合实际上等价于 Leibniz 的''不可分辨者的同一性''原则在集合论中的体现。

\subsubsection{3.5.2 空集公理}\label{ux7a7aux96c6ux516cux7406}

\textbf{公理 3.5.2（空集公理）}：\(\exists A\, \forall x\, (x \notin A)\)。此即公理 3.1.2 的精确表述。

空集公理是 ZFC 中唯一直接断言某集合存在性的公理。在 ZFC 的其他公理中，新集合的存在性总是依赖于已有集合（通过配对、并集、幂集、分离等操作），而空集公理提供了构造的起点。如果没有空集公理，ZFC 系统可能只有空模型（没有任何集合的模型），这会导致逻辑上的退化。空集在 ZFC 中扮演了类似于''零''在 Peano 算术中的角色------它是最基本的构造单元。事实上，在 ZFC 中，自然数可以纯粹用集合来定义：\(0 = \varnothing\)，\(1 = \{\varnothing\}\)，\(2 = \{\varnothing, \{\varnothing\}\}\)，等等（von Neumann 序数）。这种定义的美妙之处在于 \(n\) 恰好有 \(n\) 个元素，且 \(m < n\) 当且仅当 \(m \in n\)。空集公理与无穷公理（公理 3.5.6）共同保证了自然数集的存在性。

\textbf{公理 3.5.3（配对公理）}：\(\forall a\, \forall b\, \exists A\, \forall x\, (x \in A \leftrightarrow x = a \vee x = b)\)。即对任意两个对象 \(a, b\)，存在集合 \(\{a, b\}\)。

\textbf{推论 3.5.1（单元素集的存在性）}：对任意对象 \(a\)，存在集合 \(\{a\}\)。

\textbf{证明}：取 \(b = a\)，由配对公理得 \(\{a, a\} = \{a\}\)。\(\blacksquare\)

配对公理保证了我们可以构造恰好包含两个（或一个）特定元素的集合。这是构造有限集合的基本工具。结合并集公理（公理 3.5.4），我们可以构造任意有限集合。例如，\(\{a, b, c\}\) 可以通过 \(\{a, b\} \cup \{c\}\) 获得。配对公理之所以必要，是因为外延公理和空集公理本身不能保证非空集合的存在性------如果没有配对公理，模型可能只包含空集。在 Zermelo 1908 年的原始公理系统中，配对公理就已经被包含。值得注意的是，配对公理不要求 \(a \neq b\)；当 \(a = b\) 时，\(\{a, b\} = \{a\}\) 就是单元素集。在 von Neumann-Bernays-Godel (NBG) 集合论中，配对公理实际上可由替换公理和幂集公理推导出来，但在 ZFC 中它被保留为独立公理以保持系统的简洁性。

\subsubsection{3.5.4 并集公理}\label{ux5e76ux96c6ux516cux7406}

\textbf{公理 3.5.4（并集公理）}：\(\forall \mathcal{F}\, \exists U\, \forall x\, [x \in U \leftrightarrow \exists A\, (A \in \mathcal{F} \wedge x \in A)]\)。这个 \(U\) 称为 \(\mathcal{F}\) 的\textbf{并集}，记为 \(\bigcup \mathcal{F}\)。

并集公理是集合论中最重要的构造性公理之一，它允许我们将一个集合族（即一个集合的集合）的所有元素''展平''为一个集合。这与朴素集合论中的''所有集合的并''不同------并集公理不是从零开始创建一个并集，而是要求先有一个集合族 \(\mathcal{F}\)，然后取这个族中所有集合的元素之并。这一''限制性''设计是 ZFC 避免悖论的关键策略：你只能取一个已有集合族的并，而不能取''所有集合''的并。并集公理的典型应用包括：通过 \(\bigcup \{A, B\}\) 得到 \(A \cup B\)，通过 \(\bigcup \{\{a\}, \{a, b\}\}\) 从 Kuratowski 有序对中提取元素，以及通过 \(\bigcup \mathbb{N}\)（在 von Neumann 自然数定义下）得到 \(\mathbb{N}\) 本身。从并集公理出发，配合配对公理和分离公理，可以推导出任意两个集合的并集 \(A \cup B\) 的存在性。并集公理还为极限序数和超限归纳法提供了基础------在序数理论中，一个极限序数恰好等于其所有较小序数的并集。

\subsubsection{3.5.5 幂集公理}\label{ux5e42ux96c6ux516cux7406}

\textbf{公理 3.5.5（幂集公理）}：对任意集合 \(A\)，存在集合 \(\mathcal{P}(A)\) 使得 \(\forall S\, (S \in \mathcal{P}(A) \leftrightarrow S \subseteq A)\)。

幂集公理是 ZFC 中使集合''增长''最迅猛的公理。对于有限集 \(A\)，\(|\mathcal{P}(A)| = 2^{|A|}\)；对于无限集，康托尔定理（定理 3.6.5）证明 \(|\mathcal{P}(A)| > |A|\)。因此幂集公理产生了无穷基数的严格层级，是康托尔超限基数理论的基础。值得注意的是，幂集公理是 ZFC 中唯一''大幅增长''的构造原则------其他公理（配对、并集、分离、替换）产生的集合在某种意义下''不超过''已有集合的大小。这导致了连续统假设（\(2^{\aleph_0} = \aleph_1\)）的独立性问题（见第 6 章中关于 Godel 和 Cohen 的讨论）。在范畴论的 ETCS 基础中，幂集公理对应于一阶逻辑中''幂对象''（power object）的存在性。

\textbf{公理 3.5.6（无穷公理）}：\(\exists A\, [\varnothing \in A \wedge \forall x\, (x \in A \rightarrow x \cup \{x\} \in A)]\)。此公理保证了自然数集的存在性（见 Ch 4）。

无穷公理是 ZFC 中唯一断言无穷集合存在性的公理。在没有无穷公理的情况下，ZFC 的其余公理只能构造有限集合（从空集出发，通过配对、并集、幂集等操作，每次只能产生有限集）。无穷公理采用的构造方式是 von Neumann 序数的归纳定义：从 \(\varnothing\) 出发，反复取后继 \(x \cup \{x\}\)，得到 \(\varnothing, \{\varnothing\}, \{\varnothing, \{\varnothing\}\}, \ldots\)。这些集合恰好对应自然数 \(0, 1, 2, \ldots\)。无穷公理断言存在一个包含所有这些''自然数''的集合。由此可以定义最小的归纳集 \(\omega\)（即自然数集 \(\mathbb{N}\)）。无穷公理在 Zermelo 1908 年的原始系统中就已经出现，它标志着数学从有限到无限的飞跃。在类型论和范畴论的基础方案中，无穷公理对应自然数对象（Natural Numbers Object, NNO）的存在性。

\subsubsection{3.5.7 替换公理模式}\label{ux66ffux6362ux516cux7406ux6a21ux5f0f}

\textbf{公理 3.5.7（替换公理模式）}：设 \(\varphi(x, y)\) 是一个公式，对任意集合 \(A\)，若 \(\forall x \in A\, \exists!\, y\, \varphi(x, y)\)，则存在集合 \(\{y \mid \exists x \in A\, \varphi(x, y)\}\)。此公理允许由已知集合通过''函数替换''构造新集合。

替换公理模式是 Fraenkel (1922) 和 Skolem (1922) 独立添加到 Zermelo 原始系统中的，它是 ZFC 中唯一的''公理模式''（即对每个公式 \(\varphi\) 都有一条公理，因此实际上是无穷多条公理）。替换公理的必要性源于以下事实：Zermelo 的原始系统（仅含分离公理）不足以证明某些集合的存在性，例如 \(\{\mathbb{N}, \mathcal{P}(\mathbb{N}), \mathcal{P}(\mathcal{P}(\mathbb{N})), \ldots\}\)（von Neumann 层级中的 \(V_\omega\)）。替换公理允许我们通过''沿函数取像''来构造新集合，这是分离公理无法做到的（分离公理只能从已有集合中筛选，不能产生比原集合''更大''的集合）。替换公理模式在超限归纳和序数理论中起着核心作用。Skolem 对替换公理的一阶逻辑表述是集合论公理化的关键步骤------它使得 ZFC 成为一阶理论，从而可以应用 Godel 的完备性定理和不完备性定理。

\textbf{公理 3.5.8（正则公理/基础公理）}：\(\forall A\, [A \neq \varnothing \rightarrow \exists x \in A\, (x \cap A = \varnothing)]\)。即每个非空集合都有一个与自身不相交的元素，从而排除了 \(A \in A\) 等''病态''集合。

正则公理（也称为基础公理，Axiom of Foundation）由 von Neumann 于 1925 年提出，其作用是排除''非良基''（non-well-founded）的集合，即那些存在无穷下降链 \(\cdots \in x_3 \in x_2 \in x_1\) 的集合。正则公理保证了所有集合都属于 von Neumann 累积层级（cumulative hierarchy）\(V = \bigcup_{\alpha \in \mathbf{Ord}} V_\alpha\)，其中 \(V_0 = \varnothing\)，\(V_{\alpha+1} = \mathcal{P}(V_\alpha)\)，\(V_\lambda = \bigcup_{\alpha < \lambda} V_\alpha\)（极限序数）。这个层级结构为集合论提供了清晰的''构建''图景：每个集合都在某个阶段被构造出来。值得注意的是，正则公理在 ZFC 中并非''数学必需''------绝大多数数学定理（包括分析、代数、拓扑等）的证明并不依赖正则公理，它的主要作用是简化集合论本身的技术性处理（如序数理论）。因此，有些集合论学者（如 Aczel）研究了非良基集合论（反基础公理 AFA），以处理某些计算机科学中的循环结构。

\textbf{定理 3.5.1}：不存在集合 \(A\) 使得 \(A \in A\)。

\textbf{证明}：假设 \(A \in A\)。令 \(B = \{A\}\)。\(B\) 是非空集合。由正则公理，存在 \(x \in B\) 使得 \(x \cap B = \varnothing\)。但 \(x = A\)，\(A \cap \{A\}\) 包含 \(A\)（因 \(A \in A\) 且 \(A \in \{A\}\)），故 \(A \cap \{A\} \neq \varnothing\)，矛盾。\(\blacksquare\)

\subsubsection{3.5.9 选择公理}\label{ux9009ux62e9ux516cux7406}

\textbf{公理 3.5.9（选择公理）}：对任意由非空集合组成的集合族 \(\mathcal{F}\)，存在函数 \(f: \mathcal{F} \to \bigcup \mathcal{F}\) 使得 \(f(X) \in X\) 对所有 \(X \in \mathcal{F}\) 成立。即可以从每个非空集合中各''选择''一个元素。

选择公理等价于许多重要命题，包括佐恩引理和良序原理（定理 4.1.11 在一般集合上的推广）。

选择公理（Axiom of Choice，AC）是 ZFC 公理系统中最具争议性的公理。它由 Zermelo 于 1904 年明确提出，用于证明良序定理（每个集合都可以被良序化）。AC 的争议性在于它断言了选择函数的存在性，但并不给出构造这种函数的方法------这是一种非构造性的存在断言。在 ZF 集合论中，AC 独立于其他公理（Godel 1938 证明 AC 与 ZF 协调，Cohen 1963 证明 AC 的否定也与 ZF 协调）。AC 与许多重要数学命题等价，包括：Zorn 引理（在代数和泛函分析中广泛使用）、良序定理、Tychonoff 定理（任意紧空间的乘积是紧的）、任意向量空间有基、Krull 定理（任意真理想可扩充为极大理想）等。放弃 AC 会导致许多''病态''后果，例如可能存在没有基的向量空间、存在没有 Lebesgue 测度的实数子集等。现代数学的主流实践是接受 AC，但会在可能的情况下标明哪些结果依赖 AC（例如在构造性数学中避免使用 AC）。

\hrulefill

\section{第8章：基数简介}\label{ux7b2c8ux7ae0ux57faux6570ux7b80ux4ecb}

基数是衡量集合''大小''的概念，是对有限集合的元素个数和无限集合的''无穷程度''的统一推广。在有限情形下，基数就是集合中元素的个数；但在无限情形下，事情变得复杂得多------Cantor 在 1874 年的革命性发现表明，并非所有无限集都是''一样大''的：自然数集 \(\mathbb{N}\) 是可数的（\(\aleph_0\)），而实数集 \(\mathbb{R}\) 是不可数的（\(2^{\aleph_0}\)），且后者的基数严格大于前者。本节将引入基数相等（等势）和基数比较的概念，通过 Cantor-Bernstein 定理建立基数比较的良基性质，并证明经典的 Cantor 定理------任何集合的幂集的基数严格大于该集合本身的基数，从而揭示了无穷大的无限层级。基数理论还与选择公理密切相关：在 ZFC 中，任何两个基数都是可比较的（基数三歧性），但这一结论依赖于选择公理。基数算术（加法、乘法、幂运算）与普通的自然数算术不同，在无限基数上表现出许多反直觉的性质，如 \(\aleph_0 + \aleph_0 = \aleph_0\) 和 \(\aleph_0 \cdot \aleph_0 = \aleph_0\)。连续性假设（CH）------\(\aleph_1 = 2^{\aleph_0}\)------是集合论中最著名的问题之一，Cohen 在 1963 年证明它独立于 ZFC。

\subsubsection{3.6.1 等势}\label{ux7b49ux52bf}

\textbf{定义 3.6.1（等势）}：两个集合 \(A\) 和 \(B\) 称为\textbf{等势}（记为 \(A \sim B\)），如果存在双射 \(f: A \to B\)。

\textbf{定理 3.6.1}：等势是等价关系。

\textbf{证明}：自反性：\(\mathrm{id}_{A}: A \to A\) 是双射。对称性：若 \(f: A \to B\) 是双射，则 \(f^{-1}: B \to A\) 是双射。传递性：若 \(f: A \to B\) 和 \(g: B \to C\) 是双射，则 \(g \circ f: A \to C\) 是双射。\(\blacksquare\)

等势（equinumerosity）是康托尔集合论中最核心的概念之一，它使得我们能够在严格意义上比较无穷集合的''大小''。等势关系将所有集合划分为等价类，每个等价类对应于一个''基数''（cardinal number）。在 ZFC 中，基数可以精确定义为最小的等势序数（初始序数），例如 \(\aleph_0\)（读作 aleph-null）定义为最小的无穷序数 \(\omega\) 的基数，即 \(|\mathbb{N}| = \aleph_0\)。康托尔的天才洞见在于认识到：并非所有无穷集合都等势------存在不同''大小''的无穷。可数无穷（\(\aleph_0\)）和连续统（\(\mathfrak{c} = 2^{\aleph_0}\)）是最基本的两个无穷基数。康托尔定理进一步证明了对任意集合 \(A\)，\(|A| < |\mathcal{P}(A)|\)，因此不存在最大的基数------基数可以无限递增：\(\aleph_0 < 2^{\aleph_0} < 2^{2^{\aleph_0}} < \cdots\)。选择公理（AC）等价于''任意两个集合的基数可以比较''（基数三歧性），即对任意集合 \(A, B\)，恰有 \(|A| < |B|\)、\(|A| = |B|\)、\(|A| > |B|\) 之一成立。

\subsubsection{3.6.2 可数集}\label{ux53efux6570ux96c6}

\textbf{定义 3.6.2（可数集）}：与自然数集 \(\mathbb{N}\) 等势的集合称为\textbf{可数无限集}。有限集或可数无限集统称为\textbf{可数集}。

\textbf{定理 3.6.2}：整数集 \(\mathbb{Z}\) 是可数的。

\textbf{证明}：构造双射 \(f: \mathbb{N} \to \mathbb{Z}\)： \[f(n) = \begin{cases} n/2 & \text{若 } n \text{ 是偶数} \\ -(n+1)/2 & \text{若 } n \text{ 是奇数} \end{cases}\] \(f(0) = 0, f(1) = -1, f(2) = 1, f(3) = -2, f(4) = 2, \ldots\)。这是 \(\mathbb{N}\) 到 \(\mathbb{Z}\) 的一一对应。\(\blacksquare\)

\textbf{定理 3.6.3}：有理数集 \(\mathbb{Q}\) 是可数的。

\textbf{证明}：只需证明正有理数 \(\mathbb{Q}^+\) 可数（然后用定理 3.6.2 的技巧处理符号）。构造显式的 Cantor 配对函数 \(J: \mathbb{N} \times \mathbb{N} \to \mathbb{N}\)：

\[J(n, m) = \frac{(n+m-2)(n+m-1)}{2} + n, \quad n, m \in \mathbb{N}\]

其中 \(\mathbb{N} = \{1, 2, 3, \ldots\}\)。\(J\) 是双射的验证： - \textbf{单射性}：若 \(J(n, m) = J(n', m')\)，则由公式可知 \(n+m\) 与 \(n'+m'\) 必相等（否则对应不同的对角线条目区间），进而 \(n = n'\)，\(m = m'\)。 - \textbf{满射性}：对任意 \(k \in \mathbb{N}\)，存在唯一的自然数 \(s\) 使得 \(T_{s-1} < k \leq T_s\)（其中 \(T_s = \frac{s(s+1)}{2}\) 为三角数）。令 \(n = k - T_{s-1}\)，\(m = s+1 - n\)，则 \(n \in \{1, \ldots, s\}\)，\(m \in \{1, \ldots, s\}\)，且 \(J(n, m) = k\)。

现在令 \(f: \mathbb{N} \to \mathbb{Q}^+\) 为：对 \(k \in \mathbb{N}\)，由 \(k\) 的唯一分解 \(k = J(n, m)\)，令 \(f(k) = n/m\)。由于 \(J\) 是双射，每个有理数（未约分）被遍历恰好一次。跳过重复者（即当 \(\gcd(n, m) > 1\) 时）即得 \(\mathbb{N}\) 到 \(\mathbb{Q}^+\) 的一一对应。\(\blacksquare\)

\textbf{定理 3.6.4（康托尔对角线论证）}：实数集 \(\mathbb{R}\) 是不可数的。

\textbf{证明}：只需证明 \((0, 1)\) 不可数。假设 \((0, 1)\) 可数，将其元素排列为 \(r_1, r_2, r_3, \ldots\)。设 \(r_n\) 的十进制表示为 \(0.d_{n1}d_{n2}d_{n3}\ldots\)。构造 \(s = 0.s_1 s_2 s_3 \ldots\)，其中 \[s_n = \begin{cases} 3 & \text{若 } d_{nn} \neq 3 \\ 7 & \text{若 } d_{nn} = 3 \end{cases}\] 则 \(s \in (0, 1)\) 且 \(s \neq r_n\) 对所有 \(n\)（因为 \(s\) 的第 \(n\) 位与 \(r_n\) 的第 \(n\) 位不同）。这与 \(r_1, r_2, \ldots\) 列出了 \((0, 1)\) 的所有元素矛盾。\(\blacksquare\)

\textbf{推论 3.6.1}：\(\mathbb{R}\) 与 \(\mathbb{N}\) 不等势。存在''大小''不同的无穷------\(\mathbb{R}\) 是''不可数无穷''，\(\mathbb{N}\) 是''可数无穷''。

\textbf{定理 3.6.5（康托尔定理）}：对任意集合 \(A\)，\(A\) 与 \(\mathcal{P}(A)\) 不等势，且 \(|A| < |\mathcal{P}(A)|\)。

\textbf{证明}：

\textbf{第一步}：\(|A| \leq |\mathcal{P}(A)|\)。映射 \(a \mapsto \{a\}\) 是 \(A\) 到 \(\mathcal{P}(A)\) 的单射。

\textbf{第二步}：不存在从 \(A\) 到 \(\mathcal{P}(A)\) 的满射。假设 \(f: A \to \mathcal{P}(A)\) 是满射。令 \(B = \{a \in A : a \notin f(a)\} \subseteq A\)，故 \(B \in \mathcal{P}(A)\)。由 \(f\) 是满射，存在 \(b \in A\) 使得 \(f(b) = B\)。

问：\(b \in B\) 吗？ - 若 \(b \in B\)，由 \(B\) 的定义 \(b \notin f(b) = B\)，矛盾。 - 若 \(b \notin B\)，由 \(B\) 的定义 \(b \in B\)，矛盾。

因此不存在这样的满射。\(\blacksquare\)

\textbf{推论 3.6.2}：不存在最大的基数。对任意集合 \(A\)，总存在严格更大的集合 \(\mathcal{P}(A)\)。

\textbf{推论 3.6.3}：\(|\mathcal{P}(\mathbb{N})| = |\mathbb{R}|\)。

*\textbf{证明}：对任意\(f:A\to\mathcal{P}(A)\)，构造\(B=\{a\in A:a\notin f(a)\}\in\mathcal{P}(A)\)。若\(B=f(b)\)对某\(b\in A\)，则\(b\in B\iff b\notin f(b)\iff b\notin B\)，矛盾。故\(B\)不在\(f\)的像中，\(f\)非满射。因此不存在\(A\)到\(\mathcal{P}(A)\)的双射，\(|A|<|\mathcal{P}(A)|\)。\(\blacksquare\)

lacksquare\$

\hrulefill

\subsection{本章展望}\label{ux672cux7ae0ux5c55ux671b-3}

本章建立了集合论的基本框架，包括集合的运算、关系和函数的严格定义。在 Ch 4 中，我们将用集合（特别是等价关系和商集）来构造从自然数到复数的完整数系。函数的概念将在卷四（Ch 23--24）中得到充分发展。等价关系将在 Ch 5（同余理论）中再次出现。幂集和笛卡尔积的概念将在后续各卷中反复使用。

\hrulefill

\hrulefill

\hrulefill

\hrulefill

\hrulefill

\newpage

\chapter{01-基础/02-集合与数系/02-数系构造.md}\label{ux57faux784002-ux96c6ux5408ux4e0eux6570ux7cfb02-ux6570ux7cfbux6784ux9020.md}

\textbf{前置知识}：依赖 Ch 3（集合论），特别是等价关系（3.3.2 节）、商集（定义 3.3.5）、函数（3.4 节）。用集合论的方法构造数系。

数学研究的核心对象之一是''数''。在 Ch 3 中我们建立了集合论的语言，现在我们用这门语言来回答一个根本问题：数究竟是什么？本章从皮亚诺公理出发，逐步构造自然数 \(\mathbb{N}\)、整数 \(\mathbb{Z}\)、有理数 \(\mathbb{Q}\)、实数 \(\mathbb{R}\) 和复数 \(\mathbb{C}\)，每一步都严格证明新数系继承旧数系的代数性质并解决其不足。

\hrulefill

\section{第9章：自然数（皮亚诺公理）}\label{ux7b2c9ux7ae0ux81eaux7136ux6570ux76aeux4e9aux8bfaux516cux7406}

自然数是数学中最基本的数系，也是所有其他数系构造的起点。本节从皮亚诺公理出发，给出自然数的严格公理化定义，然后通过递归定义引入加法和乘法运算，并严格证明所有的基本运算律（结合律、交换律、分配律、消去律）。皮亚诺公理系统的核心思想是：自然数完全由''0''和''后继''这两个基本概念以及五条公理（特别是归纳公理）所刻画。归纳公理（P5）是最关键的一条------它不仅是数学归纳法的逻辑基础，还保证了自然数集是''最小的归纳集''，从而确立了自然数在同构意义下的唯一性（Dedekind 的范畴性定理）。本节还将建立自然数的序关系，证明良序原理（每个非空自然数子集有最小元），并证明归纳公理与良序原理的等价性------这一等价性揭示了自然数结构的两个基本方面：归纳的''构造''视角和良序的''极小''视角。

\subsubsection{4.1.1 皮亚诺公理}\label{ux76aeux4e9aux8bfaux516cux7406}

\textbf{定义 4.1.1（自然数系）}：一个\textbf{自然数系}是一个三元组 \((\mathbb{N}, 0, S)\)，其中 \(\mathbb{N}\) 是一个集合，\(0 \in \mathbb{N}\) 是一个特定元素，\(S: \mathbb{N} \to \mathbb{N}\) 是一个函数（称为\textbf{后继函数}），满足以下五条公理：

\textbf{公理 P1}（零是自然数）：\(0 \in \mathbb{N}\)。

\textbf{公理 P2}（后继的封闭性）：\(\forall n \in \mathbb{N}: S(n) \in \mathbb{N}\)。

\textbf{公理 P3}（零不是任何数的后继）：\(\forall n \in \mathbb{N}: S(n) \neq 0\)。

\textbf{公理 P4}（后继函数是单射）：\(\forall m, n \in \mathbb{N}: S(m) = S(n) \Rightarrow m = n\)。

\textbf{公理 P5（归纳公理）}：设 \(A \subseteq \mathbb{N}\) 满足：(i) \(0 \in A\)；(ii) 对任意 \(n \in A\)，\(S(n) \in A\)。则 \(A = \mathbb{N}\)。

\textbf{约定}：以下在固定的自然数系 \((\mathbb{N}, 0, S)\) 中工作。将 \(S(0)\) 记为 \(1\)，\(S(S(0))\) 记为 \(2\)，以此类推。

皮亚诺公理系统是意大利数学家 Giuseppe Peano 于 1889 年在《算术原理》（Arithmetices Principia）中提出的，它用五条简洁的公理刻画了自然数的本质。公理 P1 和 P2 保证了自然数集的''封闭性''------从 \(0\) 出发通过后继运算可以无限制地产生新的自然数。公理 P3 确立了 \(0\) 的特殊地位------它是自然数的''起点''，不是任何自然数的后继。公理 P4 保证了后继链不会''合并''------不同的自然数有不同的后继，这使得自然数序列是线性的。公理 P5（归纳公理）是最深刻的一条，它确立了数学归纳法的合法性，同时保证了自然数集是''最小的归纳集''------任何包含 \(0\) 且在后继下封闭的集合都包含全部自然数。这五条公理在二阶逻辑中刻画了自然数系在同构意义下的唯一性（Dedekind 的范畴性定理），但在一阶逻辑中，由于 Lowenheim-Skolem 定理，存在非标准模型。皮亚诺公理是现代数学基础的里程碑，它表明算术可以完全公理化。

\subsubsection{4.1.2 加法的递归定义}\label{ux52a0ux6cd5ux7684ux9012ux5f52ux5b9aux4e49}

\textbf{定理 4.1.1（加法的存在唯一性）}：存在唯一的二元运算 \(+: \mathbb{N} \times \mathbb{N} \to \mathbb{N}\) 满足： 1. \(\forall a \in \mathbb{N}: a + 0 = a\) 2. \(\forall a, b \in \mathbb{N}: a + S(b) = S(a + b)\)

\textbf{证明}：

\textbf{存在性}：对每个固定的 \(a \in \mathbb{N}\)，由递归定理（在集合论中可以严格证明），存在唯一的函数 \(f_a: \mathbb{N} \to \mathbb{N}\) 使得 \(f_a(0) = a\)，\(f_a(S(b)) = S(f_a(b))\)。定义 \(a + b = f_a(b)\)。

\textbf{唯一性}：假设 \(+\) 和 \(\oplus\) 都满足上述条件。对固定的 \(a\)，令 \(A = \{b \in \mathbb{N} : a + b = a \oplus b\}\)。由条件 1，\(0 \in A\)。设 \(b \in A\)，则 \(a + S(b) = S(a + b) = S(a \oplus b) = a \oplus S(b)\)，故 \(S(b) \in A\)。由归纳公理，\(A = \mathbb{N}\)。\(\blacksquare\)

\subsubsection{4.1.3 乘法的递归定义}\label{ux4e58ux6cd5ux7684ux9012ux5f52ux5b9aux4e49}

\textbf{定理 4.1.2（乘法的存在唯一性）}：存在唯一的二元运算 \(\cdot: \mathbb{N} \times \mathbb{N} \to \mathbb{N}\) 满足： 1. \(\forall a \in \mathbb{N}: a \cdot 0 = 0\) 2. \(\forall a, b \in \mathbb{N}: a \cdot S(b) = a \cdot b + a\)

\textbf{证明}：方法与定理 4.1.1 完全类似。\(\blacksquare\)

\subsubsection{4.1.4 自然数运算律的证明}\label{ux81eaux7136ux6570ux8fd0ux7b97ux5f8bux7684ux8bc1ux660e}

以下严格证明自然数的基本运算律。所有证明都基于皮亚诺公理和加法、乘法的递归定义。

\textbf{引理 4.1.1}：\(\forall a \in \mathbb{N}: 0 + a = a\)。

\textbf{证明}：对 \(a\) 做归纳。\(0 + 0 = 0\)。设 \(0 + a = a\)，则 \(0 + S(a) = S(0 + a) = S(a)\)。\(\blacksquare\)

\textbf{引理 4.1.2}：\(\forall a, b \in \mathbb{N}: S(a) + b = S(a + b)\)。

\textbf{证明}：对 \(b\) 做归纳。\(S(a) + 0 = S(a) = S(a + 0)\)。设 \(S(a) + b = S(a + b)\)，则 \(S(a) + S(b) = S(S(a) + b) = S(S(a + b)) = S(a + S(b))\)。\(\blacksquare\)

\textbf{注记 4.1.1}：引理 4.1.1 与加法定义的第一条（\(a + 0 = a\)）形式上不同。后者是''右加零''，前者是''左加零''。两者都需要证明，因为加法的递归定义是不对称的（定义在第二个变量上递归）。

\textbf{定理 4.1.3（加法结合律）}：\(\forall a, b, c \in \mathbb{N}: (a + b) + c = a + (b + c)\)。

\textbf{证明}：对 \(c\) 做归纳。(基步) \((a + b) + 0 = a + b = a + (b + 0)\)。(归纳步) 设 \((a + b) + c = a + (b + c)\)，则 \((a + b) + S(c) = S((a + b) + c) = S(a + (b + c)) = a + S(b + c) = a + (b + S(c))\)。\(\blacksquare\)

\textbf{定理 4.1.4（加法交换律）}：\(\forall a, b \in \mathbb{N}: a + b = b + a\)。

\textbf{证明}：先证辅助结论 \(a + 1 = S(a)\)（因为 \(a + S(0) = S(a + 0) = S(a)\)）。对 \(b\) 做归纳。(基步) \(a + 0 = a = 0 + a\)（引理 4.1.1）。(归纳步) 设 \(a + b = b + a\)，则 \(a + S(b) = S(a + b) = S(b + a) = S(b) + a\)（引理 4.1.2）。\(\blacksquare\)

\textbf{引理 4.1.3}：\(\forall a \in \mathbb{N}: 0 \cdot a = 0\)。

\textbf{证明}：对 \(a\) 做归纳。\(0 \cdot 0 = 0\)。设 \(0 \cdot a = 0\)，则 \(0 \cdot S(a) = 0 \cdot a + 0 = 0 + 0 = 0\)。\(\blacksquare\)

\textbf{引理 4.1.4}：\(\forall a, b \in \mathbb{N}: S(a) \cdot b = a \cdot b + b\)。

\textbf{证明}：对 \(b\) 做归纳。(基步) \(S(a) \cdot 0 = 0 = a \cdot 0 + 0\)。(归纳步) 设 \(S(a) \cdot b = a \cdot b + b\)，则 \(S(a) \cdot S(b) = S(a) \cdot b + S(a) = (a \cdot b + b) + S(a)\)。由加法结合律和交换律，\(= a \cdot b + (b + S(a)) = a \cdot b + S(b + a) = a \cdot b + S(a + b) = a \cdot b + (a + S(b)) = (a \cdot b + a) + S(b) = a \cdot S(b) + S(b)\)。\(\blacksquare\)

\textbf{定理 4.1.5（左分配律）}：\(\forall a, b, c \in \mathbb{N}: a \cdot (b + c) = a \cdot b + a \cdot c\)。

\textbf{证明}：对 \(c\) 做归纳。(基步) \(a \cdot (b + 0) = a \cdot b = a \cdot b + 0 = a \cdot b + a \cdot 0\)。(归纳步) 设 \(a \cdot (b + c) = a \cdot b + a \cdot c\)，则 \(a \cdot (b + S(c)) = a \cdot S(b + c) = a \cdot (b + c) + a = (a \cdot b + a \cdot c) + a = a \cdot b + (a \cdot c + a) = a \cdot b + a \cdot S(c)\)。\(\blacksquare\)

\textbf{定理 4.1.6（乘法交换律）}：\(\forall a, b \in \mathbb{N}: a \cdot b = b \cdot a\)。

\textbf{证明}：对 \(b\) 做归纳。(基步) \(a \cdot 0 = 0 = 0 \cdot a\)（引理 4.1.3）。(归纳步) 设 \(a \cdot b = b \cdot a\)，则 \(a \cdot S(b) = a \cdot b + a = b \cdot a + a = S(b) \cdot a\)（引理 4.1.4）。\(\blacksquare\)

\textbf{定理 4.1.7（乘法结合律）}：\(\forall a, b, c \in \mathbb{N}: (a \cdot b) \cdot c = a \cdot (b \cdot c)\)。

\textbf{证明}：对 \(c\) 做归纳。(基步) \((a \cdot b) \cdot 0 = 0 = a \cdot 0 = a \cdot (b \cdot 0)\)。(归纳步) 设 \((a \cdot b) \cdot c = a \cdot (b \cdot c)\)，则 \((a \cdot b) \cdot S(c) = (a \cdot b) \cdot c + a \cdot b = a \cdot (b \cdot c) + a \cdot b = a \cdot (b \cdot c + b) = a \cdot (b \cdot S(c))\)。\(\blacksquare\)

\textbf{定理 4.1.8（消去律）}：\(\forall a, b, c \in \mathbb{N}: a + c = b + c \Rightarrow a = b\)。

\textbf{证明}：对 \(c\) 做归纳。(基步) \(a + 0 = b + 0 \Rightarrow a = b\)。(归纳步) 设 \(a + c = b + c \Rightarrow a = b\)。若 \(a + S(c) = b + S(c)\)，即 \(S(a + c) = S(b + c)\)，由公理 P4 得 \(a + c = b + c\)，由归纳假设 \(a = b\)。\(\blacksquare\)

\textbf{引理 4.1.5}：\(\forall n \in \mathbb{N}\)，\(n = 0\) 或 \(n = S(m)\) 对某个 \(m \in \mathbb{N}\)。

\textbf{证明}：令 \(A = \{0\} \cup \{S(m) : m \in \mathbb{N}\}\)。\(0 \in A\)。若 \(n \in A\)，则 \(S(n) \in A\)（因为 \(S(n)\) 是某自然数的后继）。由公理 P5（归纳公理），\(A = \mathbb{N}\)。\(\blacksquare\)

\textbf{定理 4.1.9（乘法消去律）}：\(\forall a, b, c \in \mathbb{N}\)，若 \(c \neq 0\)，则 \(a \cdot c = b \cdot c \Rightarrow a = b\)。

\textbf{证明}：对 \(a\) 做归纳。

\textbf{基步}（\(a = 0\)）：则 \(0 = a \cdot c = b \cdot c\)。需证 \(b = 0\)。设 \(c = S(d)\)（因 \(c \neq 0\)，\(c\) 必为某自然数的后继）。若 \(b \neq 0\)，则 \(b = S(b')\)，于是 \(b \cdot c = S(b') \cdot c = b' \cdot c + c = b' \cdot c + S(d) = S(b' \cdot c + d)\)。故 \(S(b' \cdot c + d) = 0\)，与公理 P3（\(0\) 不是任何自然数的后继）矛盾。因此 \(b = 0 = a\)。

\textbf{归纳步}：设 \(a = S(a')\)，且对 \(a'\) 结论成立。设 \(S(a') \cdot c = b \cdot c\)，即 \(a' \cdot c + c = b \cdot c\)。若 \(b = 0\)，则 \(a' \cdot c + c = 0\)。由与基步相同的论证（\(c = S(d)\) 导致 \(S(a' \cdot c + d) = 0\) 矛盾），故 \(b \neq 0\)。设 \(b = S(b')\)，则 \(a' \cdot c + c = S(b') \cdot c = b' \cdot c + c\)。由加法消去律，\(a' \cdot c = b' \cdot c\)。由归纳假设 \(a' = b'\)，故 \(S(a') = S(b')\)。\(\blacksquare\)

\textbf{定理 4.1.10（乘法零因子律）}：\(\forall a, b \in \mathbb{N}: a \cdot b = 0 \iff a = 0 \vee b = 0\)。

\textbf{证明}：\((\Leftarrow)\)：若 \(a = 0\)，\(0 \cdot b = 0\)。若 \(b = 0\)，\(a \cdot 0 = 0\)。\((\Rightarrow)\)：假设 \(a \neq 0\) 且 \(b \neq 0\)，设 \(a = S(a')\), \(b = S(b')\)。则 \(a \cdot b = S(a') \cdot S(b') = S(a') \cdot b' + S(a')\)。由于 \(S(a') \geq 1\)，\(a \cdot b \geq S(a') \geq 1 > 0\)，矛盾。\(\blacksquare\)

\subsubsection{4.1.5 自然数的序关系}\label{ux81eaux7136ux6570ux7684ux5e8fux5173ux7cfb}

\textbf{定义 4.1.2}：\(\forall a, b \in \mathbb{N}\)，定义 \(a \leq b\) 当且仅当 \(\exists c \in \mathbb{N}: a + c = b\)。

\textbf{定理 4.1.11（良序原理）}：\(\mathbb{N}\) 的每个非空子集都有最小元素。

\textbf{证明}：假设 \(A \subseteq \mathbb{N}\) 非空但没有最小元素。令 \(B = \{n \in \mathbb{N} : \forall k \leq n, k \notin A\}\)。对 \(B\) 用归纳法证明 \(B = \mathbb{N}\)。

(基步) 若 \(0 \in A\)，则 \(0\) 是 \(A\) 的最小元素，矛盾。故 \(0 \notin A\)，\(0 \in B\)。

(归纳步) 设 \(\{0, 1, \ldots, n\} \subseteq B\)。若 \(S(n) \in A\)，则 \(S(n)\) 是 \(A\) 的最小元素（所有更小的自然数都不在 \(A\) 中），矛盾。故 \(S(n) \notin A\)，\(S(n) \in B\)。

由归纳公理，\(B = \mathbb{N}\)，从而 \(A = \varnothing\)，矛盾。\(\blacksquare\)

\textbf{定理 4.1.11'（良序原理与归纳法的等价性）}：归纳公理（公理 P5）与良序原理（定理 4.1.11）等价。

\textbf{证明}：我们已经从归纳公理推导了良序原理。反过来，假设良序原理成立，证明归纳公理。

设 \(A \subseteq \mathbb{N}\) 满足 (i) \(0 \in A\)；(ii) \(n \in A \Rightarrow S(n) \in A\)。假设 \(A \neq \mathbb{N}\)，令 \(C = \mathbb{N} \setminus A \neq \varnothing\)。由良序原理，\(C\) 有最小元素 \(m\)。\(m \neq 0\)（因 \(0 \in A\)），故 \(m = S(m')\)。\(m' < m\) 且 \(m' \notin C\)（\(m\) 是最小元素），故 \(m' \in A\)。由条件 (ii)，\(S(m') = m \in A\)，矛盾。\(\blacksquare\)

\hrulefill

\section{第10章：整数}\label{ux7b2c10ux7ae0ux6574ux6570}

在自然数 \(\mathbb{N}\) 中，加法逆元不存在------方程 \(a + x = b\) 并非对所有 \(a, b\) 有解。本节通过有序对商构造的方法，从 \(\mathbb{N}\) 构造出整数环 \(\mathbb{Z}\)。核心思想是将每个整数 \(n\)（正、负或零）表示为自然数对 \((a, b)\) 的等价类，其中 \((a, b)\) 直观地代表''\(a - b\)``。通过等价关系 \((a, b) \sim (c, d) \iff a + d = b + c\)，我们将表示同一个差值的无穷多个有序对合并为一个等价类。这种商构造是代数中构造新结构的标准方法，在有理数、实数乃至群论中的商群构造中反复出现。\(\mathbb{Z}\) 不仅是加法群（每个元素有加法逆元），还是含幺交换环，并且是全序的------序关系自然地由 \(\mathbb{N}\) 的序关系诱导而来。本节还将证明整数的基本代数性质，包括''负负得正''（\((-1) \times (-1) = 1\)）和消去律，以及阿基米德性质。

\subsubsection{4.2.1 构造动机}\label{ux6784ux9020ux52a8ux673a}

在自然数系中，减法并不总是可行的：方程 \(3 + x = 1\) 在 \(\mathbb{N}\) 中无解。我们扩展数系使减法总可以进行。

从 \(\mathbb{N}\) 到 \(\mathbb{Z}\) 的扩展是数系扩展的第一步，其动力来自代数封闭性的需求。在自然数 \(\mathbb{N}\) 中，加法是封闭的，但减法的结果可能不在 \(\mathbb{N}\) 中（例如 \(1 - 3\) 不是自然数）。我们希望构造一个更大的数系，使得减法运算（即加法逆元的存在性）总是可行的。构造方法的核心思想是将''差'' \(a - b\) 编码为有序对 \((a, b)\)，并通过等价关系将表示同一个差的无穷多个有序对合并为一个等价类。这种''商构造''（quotient construction）是代数中构造新结构的标准方法：先构造一个''过大''的结构（此处是 \(\mathbb{N} \times \mathbb{N}\)），然后通过等价关系将其''压缩''为想要的结构。这种构造模式在有理数、实数乃至群论中的商群构造中反复出现。

从哲学角度看，从 \(\mathbb{N}\) 到 \(\mathbb{Z}\) 的扩展体现了数学中''通过添加形式解来扩展数系''的普遍策略。在 \(\mathbb{N}\) 中，方程 \(a + x = b\) 只在 \(a \leq b\) 时有解。我们通过引入负数作为''形式差'' \(a - b\)（即使 \(a < b\)）来扩展数系，使得每个这样的方程都有解。这一策略在从 \(\mathbb{R}\) 到 \(\mathbb{C}\) 的扩展中再次出现（引入 \(i\) 作为 \(x^2 = -1\) 的形式解），也是域扩张理论的一般原理。需要注意的是，\(\mathbb{Z}\) 的构造并不是简单地''在 \(\mathbb{N}\) 前面加一些负数''------\(\mathbb{Z}\) 的每一个元素，包括正整数，都通过商构造重新定义。这样，正自然数 \(n\) 在 \(\mathbb{Z}\) 中的对应物 \([n, 0]\) 与原始的自然数 \(n\) 在严格意义上不是同一个对象，但存在一个保序、保运算的嵌入 \(\iota: \mathbb{N} \hookrightarrow \mathbb{Z}\)，使得我们可以将 \(\mathbb{N}\) 视为 \(\mathbb{Z}\) 的子集。

\subsubsection{4.2.2 整数的定义}\label{ux6574ux6570ux7684ux5b9aux4e49}

\textbf{定义 4.2.1}：令 \(P = \mathbb{N} \times \mathbb{N}\)。在 \(P\) 上定义关系 \(\sim\)： \[(a, b) \sim (c, d) \iff a + d = b + c\]

\textbf{定理 4.2.1}：\(\sim\) 是 \(P\) 上的等价关系。

\textbf{证明}： - \textbf{自反性}：\((a, b) \sim (a, b)\)，因为 \(a + b = b + a\)（加法交换律）。 - \textbf{对称性}：若 \((a, b) \sim (c, d)\)，即 \(a + d = b + c\)，则 \(c + b = d + a\)，即 \((c, d) \sim (a, b)\)。 - \textbf{传递性}：若 \((a, b) \sim (c, d)\) 且 \((c, d) \sim (e, f)\)，即 \(a + d = b + c\) 且 \(c + f = d + e\)。两式相加得 \(a + d + c + f = b + c + d + e\)，消去 \(c + d\)（加法消去律）得 \(a + f = b + e\)，即 \((a, b) \sim (e, f)\)。\(\blacksquare\)

\textbf{定义 4.2.2}：\textbf{整数集} \(\mathbb{Z} = P / \sim = \{[a, b] : a, b \in \mathbb{N}\}\)，其中 \([a, b]\) 代表''\(a - b\)``。

嵌入映射 \(\iota: \mathbb{N} \to \mathbb{Z}\) 定义为 \(\iota(n) = [n, 0]\)。

\subsubsection{4.2.3 整数的运算}\label{ux6574ux6570ux7684ux8fd0ux7b97}

\textbf{定义 4.2.3（加法）}：\([a, b] + [c, d] = [a + c, b + d]\)

\textbf{定理 4.2.2（加法良定义性）}：上述定义与等价类代表的选取无关。

\textbf{证明}：设 \([a, b] = [a', b']\) 且 \([c, d] = [c', d']\)，即 \(a + b' = b + a'\) 且 \(c + d' = d + c'\)。要证 \((a + c) + (b' + d') = (b + d) + (a' + c')\)。由已知条件相加即得。\(\blacksquare\)

\textbf{定义 4.2.4（乘法）}：\([a, b] \cdot [c, d] = [ac + bd, ad + bc]\)

\textbf{定理 4.2.3（乘法良定义性）}：乘法定义与代表选取无关。

\textbf{证明}：利用 \(a + b' = b + a'\) 和 \(c + d' = d + c'\)，通过自然数运算律可验证 \((ac + bd) + (a'd' + b'c') = (ad + bc) + (a'c' + b'd')\)。\(\blacksquare\)

\textbf{定义 4.2.5}：负元 \(-[a, b] = [b, a]\)；零元 \(0_{\mathbb{Z}} = [0, 0]\)；幺元 \(1_{\mathbb{Z}} = [1, 0]\)。

\subsubsection{4.2.4 整数的代数性质}\label{ux6574ux6570ux7684ux4ee3ux6570ux6027ux8d28}

\textbf{定理 4.2.4}：\((\mathbb{Z}, +, 0_{\mathbb{Z}})\) 是阿贝尔群。

\textbf{证明}：结合律和交换律由自然数的加法运算律直接继承。零元：\([a, b] + [0, 0] = [a, b]\)。负元：\([a, b] + [b, a] = [a + b, b + a] = [0, 0]\)。\(\blacksquare\)

\textbf{定理 4.2.5}：\((\mathbb{Z}, \cdot, 1_{\mathbb{Z}})\) 满足结合律和交换律。

\textbf{证明}：由自然数运算律直接展开验证。\(\blacksquare\)

\textbf{定理 4.2.6（分配律）}：\(\forall x, y, z \in \mathbb{Z}: x \cdot (y + z) = x \cdot y + x \cdot z\)。

\textbf{证明}：设 \(x = [a, b], y = [c, d], z = [e, f]\)。展开后由自然数的分配律得证。\(\blacksquare\)

\textbf{推论}：\(\mathbb{Z}\) 是含幺交换环。

\textbf{定理 4.2.7}：\((-1) \times (-1) = 1\)。

\textbf{证明}：\(-1 = [0, 1]\)，\((-1) \times (-1) = [0, 1] \cdot [0, 1] = [0 \cdot 0 + 1 \cdot 1, 0 \cdot 1 + 1 \cdot 0] = [1, 0] = 1\)。\(\blacksquare\)

\textbf{推论 4.2.1（负负得正）}：\(\forall x, y \in \mathbb{Z}: (-x)(-y) = xy\)。

\textbf{证明}：先证 \((-1) \cdot x = -x\)。设 \(x = [a, b]\)，则 \((-1) \cdot x = [0, 1] \cdot [a, b] = [0 \cdot a + 1 \cdot b, 0 \cdot b + 1 \cdot a] = [b, a] = -x\)。因此 \((-x)(-y) = ((-1) \cdot x)((-1) \cdot y) = (-1)^2 \cdot xy = 1 \cdot xy = xy\)。\(\blacksquare\)

\textbf{定理 4.2.8（整数的消去律）}：\(\forall a, b, c \in \mathbb{Z}\)，若 \(c \neq 0\)，则 \(ac = bc \Rightarrow a = b\)。

\textbf{证明}：设 \(a = [a_1, a_2]\), \(b = [b_1, b_2]\), \(c = [c_1, c_2]\)（\(c_1 \neq c_2\)）。\(ac = bc\) 意味着 \([a_1 c_1 + a_2 c_2, a_1 c_2 + a_2 c_1] = [b_1 c_1 + b_2 c_2, b_1 c_2 + b_2 c_1]\)。由等价关系定义和自然数消去律可得 \(a_1 + b_2 = a_2 + b_1\)，即 \([a_1, a_2] = [b_1, b_2]\)。\(\blacksquare\)

\subsubsection{4.2.5 整数的序关系}\label{ux6574ux6570ux7684ux5e8fux5173ux7cfb}

\textbf{定义 4.2.6}：对 \([a, b], [c, d] \in \mathbb{Z}\)，定义 \([a, b] \leq [c, d]\) 当且仅当 \(a + d \leq b + c\)（在 \(\mathbb{N}\) 中）。

\textbf{定理 4.2.9}：\(\leq\) 是 \(\mathbb{Z}\) 上的全序关系，且与加法相容：若 \(x \leq y\)，则 \(x + z \leq y + z\)。

\textbf{证明}：良定义性：若 \([a, b] = [a', b']\) 且 \([c, d] = [c', d']\)，则 \(a + d \leq b + c\) 等价于 \(a' + d' \leq b' + c'\)（利用等价关系和自然数序的性质）。全序性：对任意 \([a, b], [c, d]\)，比较 \(a + d\) 和 \(b + c\)（在 \(\mathbb{N}\) 中），恰有 \(a + d < b + c\)、\(a + d = b + c\)、\(a + d > b + c\) 之一成立，对应 \([a, b] < [c, d]\)、\([a, b] = [c, d]\)、\([a, b] > [c, d]\)。与加法相容：若 \([a, b] \leq [c, d]\)，即 \(a + d \leq b + c\)，则对任意 \([e, f]\)，\((a + e) + (d + f) \leq (b + f) + (c + e)\)，即 \([a, b] + [e, f] \leq [c, d] + [e, f]\)。\(\blacksquare\)

\textbf{定理 4.2.10（整数序的三歧性）}：对任意 \(x, y \in \mathbb{Z}\)，恰有以下之一成立：\(x < y\)，\(x = y\)，\(x > y\)。

\textbf{证明}：设 \(x = [a, b]\), \(y = [c, d]\)。由自然数的三歧性（对 \(a + d\) 和 \(b + c\)），恰有 \(a + d < b + c\)、\(a + d = b + c\)、\(a + d > b + c\) 之一。\(\blacksquare\)

\textbf{定理 4.2.11（整数的阿基米德性质）}：对任意正整数 \(n\) 和任意整数 \(M\)，存在正整数 \(k\) 使得 \(kn > M\)。

\textbf{证明}：分两种情况。

\textbf{情形 1}：\(M \leq 0\)。取 \(k = 1\)，则 \(kn = n > 0 \geq M\)，结论成立。

\textbf{情形 2}：\(M > 0\)。设 \(n = [a, b]\)（\(a > b\)），\(M = [c, d]\)（\(c > d\)）。此时 \(a - b\) 和 \(c - d\) 在 \(\mathbb{N}\) 中均有定义。需 \(kn > M\)，即 \([ka, kb] > [c, d]\)，即 \(ka + d > kb + c\)，即 \(k(a - b) > c - d\)。由于 \(a - b > 0\)，需证存在自然数 \(k\) 使得 \(k(a - b) > c - d\)。先证自然数的无界性：对 \(m\) 做归纳证明 \(\forall m \in \mathbb{N},\, \exists k \in \mathbb{N}:\, k > m\)。基步：\(m = 0\) 时取 \(k = 1\)，则 \(1 > 0\)。归纳步：若存在 \(k > m\)，则 \(k + 1 > S(m)\)，故结论对 \(S(m)\) 也成立。再由无界性，取 \(m = c - d\)，存在 \(k_0 \in \mathbb{N}\) 使得 \(k_0 > c - d\)。因 \(a - b \geq 1\)，有 \(k_0(a - b) \geq k_0 > c - d\)，取 \(k = k_0\) 即可。\(\blacksquare\)

\hrulefill

\section{第11章：有理数}\label{ux7b2c11ux7ae0ux6709ux7406ux6570}

在整数环 \(\mathbb{Z}\) 中，乘法逆元通常不存在------方程 \(2x = 3\) 无整数解。本节通过有序对商构造的方法，从 \(\mathbb{Z}\) 构造出有理数域 \(\mathbb{Q}\)。构造方法与从 \(\mathbb{N}\) 到 \(\mathbb{Z}\) 的构造完全平行：将''商'' \(a/b\)（\(b \neq 0\)）编码为有序对 \((a, b)\)，通过等价关系 \((a, b) \approx (c, d) \iff ad = bc\) 将表示同一个有理数的无穷多个分式合并为一个等价类。\(\mathbb{Q}\) 是第一个满足所有域公理的数系------它不仅是加法群，而且非零元素都有乘法逆元。此外，\(\mathbb{Q}\) 还是有序域，其序关系由 \(\mathbb{Z}\) 的序关系诱导而来。\(\mathbb{Q}\) 具有稠密性：任意两个不同的有理数之间存在另一个有理数。但 \(\mathbb{Q}\) 并不完备------\(\sqrt{2}\) 不是有理数，这意味着 \(\mathbb{Q}\) 中存在''间隙''。本节最后将证明 \(\mathbb{Q}\) 的稠密性和阿基米德性质，为下一节实数的构造做好铺垫。

\subsubsection{4.3.1 构造动机}\label{ux6784ux9020ux52a8ux673a-1}

在 \(\mathbb{Z}\) 中，除法并不总是可行的：\(2x = 3\) 无整数解。我们扩展数系使非零数都有乘法逆元。

从 \(\mathbb{Z}\) 到 \(\mathbb{Q}\) 的扩展是数系扩展的第二步，其动力来自乘法逆元的需求。虽然 \(\mathbb{Z}\) 已经是加法群（每个元素有加法逆元），但乘法逆元（除法）的缺失限制了代数的完备性。构造方法与从 \(\mathbb{N}\) 到 \(\mathbb{Z}\) 的构造完全平行：将''商'' \(a/b\) 编码为有序对 \((a, b)\)（\(b \neq 0\)），然后通过等价关系 \((a, b) \approx (c, d) \iff ad = bc\) 将表示同一个有理数的无穷多个分式合并为一个等价类。这种''分数等价类''方法的美妙之处在于，它自动处理了分数的约分问题------\(\frac{2}{4}\) 和 \(\frac{1}{2}\) 被识别为同一个有理数，因为它们表示同一个等价类。\(\mathbb{Q}\) 是第一个满足所有域公理的数系------它是有序域，并且是有理数对加法、减法、乘法、除法（除数为零除外）封闭的最小集合。\(\mathbb{Q}\) 是 \(\mathbb{Z}\) 的分式域（field of fractions），这一构造可以推广到任意整环。

从代数角度看，\(\mathbb{Q}\) 的构造解决了 \(\mathbb{Z}\) 中''除法不封闭''的问题。\(\mathbb{Z}\) 是一个整环（无零因子的交换环），而 \(\mathbb{Q}\) 是它的分式域------这是包含 \(\mathbb{Z}\) 的最小的域。分式域的构造是域论中的标准操作：对任意整环 \(R\)，可以构造其分式域 \(\operatorname{Frac}(R)\)，使得 \(R\) 中的每个非零元都有乘法逆元。\(\mathbb{Q} = \operatorname{Frac}(\mathbb{Z})\) 是这一构造的原型。值得注意的是，\(\mathbb{Q}\) 的稠密性（任意两个有理数之间存在第三个有理数）是 \(\mathbb{Z}\) 所不具备的------在 \(\mathbb{Z}\) 中，相邻整数之间没有其他整数。这一性质反映了从离散到稠密的质变，是 \(\mathbb{Q}\) 与 \(\mathbb{Z}\) 之间的根本区别之一。

\subsubsection{4.3.2 有理数的定义}\label{ux6709ux7406ux6570ux7684ux5b9aux4e49}

\textbf{定义 4.3.1}：令 \(P' = \{(a, b) \in \mathbb{Z} \times \mathbb{Z} : b \neq 0\}\)。在 \(P'\) 上定义关系 \(\approx\)： \[(a, b) \approx (c, d) \iff ad = bc\]

\textbf{定理 4.3.1}：\(\approx\) 是 \(P'\) 上的等价关系。

\textbf{证明}：自反性：\(ab = ba\)。对称性：若 \(ad = bc\) 则 \(cb = da\)。

传递性：若 \((a, b) \approx (c, d)\) 且 \((c, d) \approx (e, f)\)，即 \(ad = bc\) 且 \(cf = de\)。需证 \(af = be\)。

情形一：\(c \neq 0\)。由 \(ad = bc\) 和 \(cf = de\)，两式相乘得 \((ad)(cf) = (bc)(de)\)，即 \(acdf = bcde\)。因 \(c \neq 0\)，消去 \(c\) 得 \(adf = bde\)。因 \(d \neq 0\)，消去 \(d\) 得 \(af = be\)。

情形二：\(c = 0\)。由 \(ad = bc = 0\) 且 \(d \neq 0\) 得 \(a = 0\)。由 \(cf = de\) 即 \(0 = de\) 且 \(d \neq 0\) 得 \(e = 0\)。故 \(af = 0 = be\)。

综上，\((a, b) \approx (e, f)\)。\(\blacksquare\)

\textbf{定义 4.3.2}：\textbf{有理数集} \(\mathbb{Q} = P' / \approx\)。将 \([(a, b)]\) 简记为 \(\frac{a}{b}\)。

\subsubsection{4.3.3 有理数的运算}\label{ux6709ux7406ux6570ux7684ux8fd0ux7b97}

\textbf{定义 4.3.3（加法）}：\(\frac{a}{b} + \frac{c}{d} = \frac{ad + bc}{bd}\)

\textbf{定义 4.3.4（乘法）}：\(\frac{a}{b} \cdot \frac{c}{d} = \frac{ac}{bd}\)

\textbf{定义 4.3.5（负元）}：\(-\frac{a}{b} = \frac{-a}{b}\)

\textbf{定义 4.3.6（逆元）}：对 \(\frac{a}{b} \neq 0\)（\(a \neq 0\)），\(\left(\frac{a}{b}\right)^{-1} = \frac{b}{a}\)

\textbf{定理 4.3.2}：上述运算都是良定义的。

\textbf{证明}：

\textbf{加法良定义性}：设 \(\frac{a}{b} = \frac{a'}{b'}\)（即 \(ab' = a'b\)）且 \(\frac{c}{d} = \frac{c'}{d'}\)（即 \(cd' = c'd\)）。要证 \(\frac{ad+bc}{bd} = \frac{a'd'+b'c'}{b'd'}\)，即 \((ad+bc)(b'd') = (a'd'+b'c')(bd)\)。展开左端 \(= abb'd'd + bcb'd' = a'bb'd + bb'c'd = (a'd' + b'c')(bd)\)（利用 \(ab' = a'b\) 和 \(cd' = c'd\)）。

\textbf{乘法良定义性}：要证 \(\frac{ac}{bd} = \frac{a'c'}{b'd'}\)，即 \(ac \cdot b'd' = a'c' \cdot bd\)。由 \(ab' = a'b\) 和 \(cd' = c'd\)，相乘得 \((ab')(cd') = (a'b)(c'd)\)，即 \(ac \cdot b'd' = a'c' \cdot bd\)。

\textbf{逆元良定义性}：设 \(\frac{a}{b} = \frac{a'}{b'}\)（\(ab' = a'b\)）且 \(a \neq 0\)。先证 \(a' \neq 0\)：若 \(a' = 0\) 则 \(ab' = 0\)，因 \(b' \neq 0\) 得 \(a = 0\)，矛盾。要证 \(\frac{b}{a} = \frac{b'}{a'}\)，即 \(ba' = b'a\)。由 \(ab' = a'b\) 交换得 \(a'b = ab'\)，即 \(ba' = b'a\)。\(\blacksquare\)

\textbf{定义 4.3.7（有理数的序关系）}：设 \(\frac{a}{b}\) 和 \(\frac{c}{d}\) 分别取 \(b > 0\)、\(d > 0\) 的标准表示，定义 \(\frac{a}{b} \leq \frac{c}{d}\) 当且仅当 \(ad \leq bc\)（在 \(\mathbb{Z}\) 中）。

\textbf{良定义性}：若 \(\frac{a}{b} = \frac{a'}{b'}\)（\(ab' = a'b\)）且 \(\frac{c}{d} = \frac{c'}{d'}\)（\(cd' = c'd\)），\(b, d, b', d' > 0\)，则 \(ad \leq bc \iff adb'd' \leq bcb'd' \iff a'b \cdot cd' \leq ab' \cdot c'd \iff a'd' \leq b'c'\)。故定义与代表元选取无关。

\textbf{三歧性}：对任意 \(\frac{a}{b}, \frac{c}{d}\)（\(b, d > 0\)），由整数的三歧性，\(ad < bc\)、\(ad = bc\)、\(ad > bc\) 恰有一成立，对应 \(\frac{a}{b} < \frac{c}{d}\)、\(\frac{a}{b} = \frac{c}{d}\)、\(\frac{a}{b} > \frac{c}{d}\)。

\textbf{与运算的相容性}：若 \(\frac{a}{b} \leq \frac{c}{d}\)（\(b, d > 0\)），则对任意 \(\frac{e}{f}\)（\(f > 0\)）：(1) \(\frac{a}{b} + \frac{e}{f} \leq \frac{c}{d} + \frac{e}{f}\)（由 \(ad \leq bc\) 两边加 \(bef \cdot d\) 即得）；(2) 若 \(\frac{e}{f} > 0\)（即 \(e > 0\)），则 \(\frac{a}{b} \cdot \frac{e}{f} \leq \frac{c}{d} \cdot \frac{e}{f}\)（由 \(ad \leq bc\) 两边乘 \(ef > 0\) 即得）。

\subsubsection{4.3.4 有理数构成域}\label{ux6709ux7406ux6570ux6784ux6210ux57df}

\textbf{定理 4.3.3}：\((\mathbb{Q}, +, \cdot)\) 是有序域。

\textbf{证明}：

加法阿贝尔群：结合律、交换律、零元、负元均由整数的性质继承。

乘法阿贝尔群（除零外）：结合律、交换律、幺元 \(\frac{1}{1}\)、逆元 \(\frac{b}{a}\) 均可直接验证。

分配律：\(\frac{a}{b} \cdot \left(\frac{c}{d} + \frac{e}{f}\right) = \frac{a(cf + de)}{bdf} = \frac{ac}{bd} + \frac{ae}{bf}\)。\(\blacksquare\)

\subsubsection{4.3.5 有理数的稠密性}\label{ux6709ux7406ux6570ux7684ux7a20ux5bc6ux6027}

\textbf{定理 4.3.4（有理数的稠密性）}：任意两个不同的有理数之间存在有理数。即若 \(r < s\)，则存在 \(t\) 使得 \(r < t < s\)。

\textbf{证明}：取 \(t = \frac{r + s}{2}\)。\(t - r = \frac{s - r}{2} > 0\)，\(s - t = \frac{s - r}{2} > 0\)，故 \(r < t < s\)。\(\blacksquare\)

\textbf{定理 4.3.5（有理数的阿基米德性质）}：对任意正有理数 \(r, s\)，存在正整数 \(n\) 使得 \(nr > s\)。

\textbf{证明}：设 \(r = a/b\), \(s = c/d\)（\(a, b, c, d > 0\)）。需 \(n \cdot a/b > c/d\)，即 \(nad > bc\)。由于 \(a, b, c, d\) 均为正整数，\(bc\) 和 \(ad\) 均为确定的正整数。由自然数的无界性（对任意自然数 \(m\)，存在 \(k\) 使得 \(k > m\)，这可由归纳法证明），取 \(n\) 为大于 \(bc/(ad)\) 的正整数即可。\(\blacksquare\)

\textbf{定理 4.3.6（有理数对加法和乘法封闭）}：\(\mathbb{Q}\) 对四则运算封闭（除数非零时）。

\textbf{证明}：\(\frac{a}{b} + \frac{c}{d} = \frac{ad + bc}{bd} \in \mathbb{Q}\)（\(bd \neq 0\)）。\(\frac{a}{b} \cdot \frac{c}{d} = \frac{ac}{bd} \in \mathbb{Q}\)。\(\blacksquare\)

\hrulefill

\section{第12章：实数}\label{ux7b2c12ux7ae0ux5b9eux6570}

从 \(\mathbb{Q}\) 到 \(\mathbb{R}\) 的扩展是数系扩展中最深刻的一步，它涉及从''离散/稠密''到''连续''的飞跃。虽然 \(\mathbb{Q}\) 是稠密的有序域，但它并不完备------\(\sqrt{2}\) 不是有理数，这意味着 \(\mathbb{Q}\) 中存在''间隙''（gap）。实数构造的目标是填补这些间隙，使得任意有界非空子集都有上确界（确界原理）。本节将介绍两种经典的实数构造方法：Dedekind 分割和 Cauchy 序列。Dedekind 分割从''序''的角度刻画完备性------将实数线''切割''为下组和上组，每个实数由所有小于它的有理数唯一确定。Cauchy 序列则从''度量''的角度刻画完备性------每个 Cauchy 序列都收敛。这两种构造在 ZFC 中产生同构的完备有序域，从而确证实数的本质是唯一的。本节还将证明确界原理、阿基米德性质以及 \(\sqrt{2}\) 的存在性，并展示 \(\mathbb{R}\) 作为完备有序域是如何自然地继承 \(\mathbb{Q}\) 的序结构和代数结构的。

\subsubsection{4.4.1 构造动机}\label{ux6784ux9020ux52a8ux673a-2}

\(\mathbb{Q}\) 中存在''间隙''：方程 \(x^2 = 2\) 在 \(\mathbb{Q}\) 中无解。我们扩展数系以填补这些间隙，使数系完备。

从 \(\mathbb{Q}\) 到 \(\mathbb{R}\) 的扩展是数系扩展中最深刻的一步，它涉及从''离散''到''连续''的飞跃。古希腊人早已发现 \(\sqrt{2}\) 不是有理数（Hippasus，约公元前 5 世纪），但直到 19 世纪，数学家才给出了实数的严格构造。完备性的缺失意味着 \(\mathbb{Q}\) 中存在''洞''：例如，所有满足 \(x^2 < 2\) 的正有理数组成的集合，在 \(\mathbb{Q}\) 中没有上确界。实数构造的目标是填补这些''洞''，使得任意有界非空子集都有上确界（确界原理）。两种经典的构造方法------Dedekind 分割和 Cauchy 序列------分别从''序''和''度量''两个角度刻画了完备性。Dedekind 分割更直观（将实数线''切割''），Cauchy 序列则更贴合分析的实践（关心序列的收敛性）。在 ZFC 中，这两种构造产生同构的完备有序域，因此实数的本质是唯一的。

\textbf{定理 4.4.1}：\(\sqrt{2}\) 不是有理数。

\textbf{证明}：假设 \(r = \frac{p}{q}\)（\(p, q\) 互素），则 \(p^2 = 2q^2\)。\(p^2\) 为偶数，故 \(p\) 为偶数，设 \(p = 2k\)。代入得 \(q^2 = 2k^2\)，\(q\) 也为偶数。与 \(\gcd(p, q) = 1\) 矛盾。\(\blacksquare\)

这表明 \(\mathbb{Q}\) 有''漏洞''。我们需要构造完备的数系。

\subsubsection{4.4.2 戴德金分割}\label{ux6234ux5fb7ux91d1ux5206ux5272}

\textbf{定义 4.4.1（戴德金分割）}：\(\mathbb{Q}\) 的一个\textbf{戴德金分割}是一个有序对 \((A, B)\)，其中 \(A, B\) 是 \(\mathbb{Q}\) 的非空子集，满足： 1. \(A \cup B = \mathbb{Q}\)，\(A \cap B = \varnothing\) 2. \(\forall a \in A, \forall b \in B: a < b\) 3. \(A\) 没有最大元素

Dedekind 分割的核心思想是：每个实数 \(r\) 由所有小于 \(r\) 的有理数构成的集合唯一确定。条件 1 保证了分割将 \(\mathbb{Q}\) 完整地切分为两部分。条件 2 保证了 \(A\)（下组）中的所有元素都小于 \(B\)（上组）中的所有元素------这是对实数线上''左边''和''右边''的效仿。条件 3 排除了有理数点的''重复''：如果下组有最大元 \(m\)，那么 \((A \setminus \{m\}, B \cup \{m\})\) 和 \((A, B)\) 对应同一个有理数。因此我们约定下组没有最大元，这样就保证了每个有理数恰好对应一个分割。这种构造于 1872 年由 Richard Dedekind 在《连续性与无理数》（Stetigkeit und irrationale Zahlen）中首次发表，是实数严格化的里程碑。

\subsubsection{4.4.3 实数的定义}\label{ux5b9eux6570ux7684ux5b9aux4e49}

\textbf{定义 4.4.2}：\textbf{实数集} \(\mathbb{R}\) 定义为所有戴德金分割的集合。有理数 \(r\) 对应分割 \(A_r = \{q \in \mathbb{Q} : q < r\}\)。

\textbf{引理 4.4.1（嵌入的保序性）}：映射 \(r \mapsto A_r\) 是 \(\mathbb{Q}\) 到 \(\mathbb{R}\) 的保序嵌入，即 \(r < s \iff A_r \subsetneq A_s\)。

\textbf{证明}：\((\Rightarrow)\)：若 \(r < s\)，取 \(t = \frac{r+s}{2}\)，则 \(r < t < s\)，\(t \in A_s \setminus A_r\)，故 \(A_r \subsetneq A_s\)。\((\Leftarrow)\)：若 \(A_r \subsetneq A_s\)，取 \(q \in A_s \setminus A_r\)，则 \(q < s\) 且 \(q \geq r\)，故 \(r \leq q < s\)。\(\blacksquare\)

将 \(\mathbb{R}\) 定义为所有 Dedekind 分割的集合，这是数学史上最精巧的构造之一。每个 Dedekind 分割 \((A, B)\) 对应于实数线上的一个点：\(A\) 是该点''左边''的所有有理数，\(B\) 是该点''右边''的所有有理数。有理数 \(r\) 嵌入为分割 \((\{q \in \mathbb{Q} : q < r\}, \{q \in \mathbb{Q} : q \geq r\})\)，而无理数（如 \(\sqrt{2}\)）对应的是''真正''的分割------下组中没有最大元，上组中没有最小元。引理 4.4.1 保证了 \(\mathbb{Q}\) 可以保序地嵌入 \(\mathbb{R}\)，因此我们可以将 \(\mathbb{Q}\) 视为 \(\mathbb{R}\) 的子集。\(\mathbb{R}\) 的完备性（即确界原理）可以通过 Dedekind 分割的构造直接证明：任意有上界的非空实数集的上确界，就是这些实数对应的分割的并集的下组。这一构造的优美之处在于它完全用有理数（更精确地说，用集合论语言）定义了实数，从而消除了早期分析学中对''无穷小''和''极限''的直觉依赖。

\subsubsection{4.4.4 实数的运算}\label{ux5b9eux6570ux7684ux8fd0ux7b97}

\textbf{定义 4.4.3（加法）}：\(\alpha + \beta\) 的下组为 \(\{a_1 + a_2 : a_1 \in A_1, a_2 \in A_2\}\)。

\textbf{验证加法良定义}：(1) \(A = \{a_1 + a_2 : a_1 \in A_1, a_2 \in A_2\}\) 非空。(2) 若 \(q \in A\)，取 \(q' < q\)，则 \(q' = a_1 + a_2 - \epsilon\)，可写成 \(a_1' + a_2'\) 的形式，故 \(q' \in A\)。(3) \(A\) 无最大元：对 \(a_1 + a_2 \in A\)，取 \(a_1' > a_1\)（\(a_1' \in A_1\)），则 \(a_1' + a_2 > a_1 + a_2\) 且 \(a_1' + a_2 \in A\)。(4) \(B = \mathbb{Q} \setminus A\) 非空：取 \(b_1 \in B_1, b_2 \in B_2\)，则 \(b_1 + b_2 \notin A\)（因为若 \(b_1 + b_2 = a_1 + a_2\)，则 \(b_1 - a_1 = a_2 - b_2 > 0\)，但 \(b_1 - a_1 > 0\) 且 \(a_2 - b_2 < 0\)，矛盾）。(5) 若 \(a \in A, b \in B\)，则 \(a < b\)：设 \(a = a_1 + a_2\), \(b = b_1 + b_2\)（\(b_1 \notin A_1\) 或 \(b_2 \notin A_2\)）。不妨设 \(b_1 \notin A_1\)，则 \(a_1 < b_1\)，\(a_2 \leq b_2\)，故 \(a < b\)。

\textbf{定义 4.4.4（零元与幺元）}：\(0_{\mathbb{R}}\) 的下组为 \(\{q \in \mathbb{Q} : q < 0\}\)；\(1_{\mathbb{R}}\) 的下组为 \(\{q \in \mathbb{Q} : q < 1\}\)。

\textbf{定义 4.4.5（乘法）}：设 \(\alpha = (A_1, B_1)\)，\(\beta = (A_2, B_2)\) 为实数（其中 \(A_i\) 为下组，\(B_i = \mathbb{Q} \setminus A_i\) 为上组）。

\textbf{(i) 正实数的乘法}：若 \(\alpha > 0\) 且 \(\beta > 0\)（即 \(A_1, A_2\) 均包含正有理数），定义 \(\alpha \cdot \beta\) 的下组为： \[\{q \in \mathbb{Q} : q \leq 0\} \cup \{ab : a \in A_1, a > 0, b \in A_2, b > 0\}\]

该集合的向下封闭性由 \(A_1, A_2\) 的向下封闭性保证：若 \(ab\)（\(a \in A_1, b \in A_2, a, b > 0\)）在集合中，且 \(0 < q < ab\)，则 \(q = (q/b) \cdot b\)，其中 \(0 < q/b < a\)，由 \(A_1\) 的向下封闭性得 \(q/b \in A_1\)，故 \(q\) 也在集合中。

\textbf{(ii) 一般实数的乘法}：通过绝对值和符号规则推广。对 \(\alpha < 0\)，定义 \(|\alpha|\) 的下组为 \(\{-b : b \in B_1, b < 0\} \cup \{q \in \mathbb{Q} : q \leq 0\}\)（即 \(|\alpha| = -\alpha\)）。则： - 若 \(\alpha > 0, \beta > 0\)：\(\alpha \cdot \beta\) 如上定义 - 若 \(\alpha > 0, \beta < 0\)：\(\alpha \cdot \beta = -(\alpha \cdot |\beta|)\) - 若 \(\alpha < 0, \beta > 0\)：\(\alpha \cdot \beta = -(|\alpha| \cdot \beta)\) - 若 \(\alpha < 0, \beta < 0\)：\(\alpha \cdot \beta = |\alpha| \cdot |\beta|\) - 若 \(\alpha = 0\) 或 \(\beta = 0\)：\(\alpha \cdot \beta = 0_{\mathbb{R}}\)

\textbf{乘法良定义的验证}：(1) 下组非空（\(0\) 属之）。(2) 向下封闭性如上所示。(3) 无最大元：对 \(ab\)（\(a \in A_1, b \in A_2, a, b > 0\)），取 \(a' \in A_1\) 使 \(a' > a\)（\(A_1\) 无最大元），则 \(a'b > ab\) 且 \(a'b\) 仍在下组中。(4) 补集非空：取 \(b_1 \in B_1, b_2 \in B_2\)（\(b_1, b_2 > 0\)），则 \(b_1 b_2 > 0\) 且 \(b_1 b_2\) 不在下组中------若 \(b_1 b_2 = a_1 a_2\)（\(a_i \in A_i, a_i > 0\)），则 \(b_1/a_1 = a_2/b_2\)；由 \(a_1 < b_1\) 得 \(b_1/a_1 > 1\)，由 \(a_2 < b_2\) 得 \(a_2/b_2 < 1\)，矛盾。(5) 下组中元素严格小于补集中元素，由有理数乘法的序保持性可得。一般情况的良定义性由正实数情况和符号规则保证。

\textbf{定理 4.4.2}：\((\mathbb{R}, +, \cdot)\) 是有序域。

\textbf{证明}：域公理的验证均通过戴德金分割的定义和有理数的性质直接得到。以下给出关键步骤：

\textbf{加法结合律}：设 \(\alpha = (A_1, B_1)\), \(\beta = (A_2, B_2)\), \(\gamma = (A_3, B_3)\)。\((\alpha + \beta) + \gamma\) 的下组为 \(\{(a_1 + a_2) + a_3 : a_1 \in A_1, a_2 \in A_2, a_3 \in A_3\}\)。\(\alpha + (\beta + \gamma)\) 的下组为 \(\{a_1 + (a_2 + a_3) : a_1 \in A_1, a_2 \in A_2, a_3 \in A_3\}\)。由有理数加法结合律，两者相等。

\textbf{加法零元}：\(\alpha + 0_{\mathbb{R}}\) 的下组为 \(\{a + q : a \in A_1, q < 0\}\)。对任意 \(r \in A_1\)，取 \(a = r\) 和 \(q = -\epsilon\)（\(\epsilon > 0\) 足够小），则 \(r - \epsilon < r\)，即 \(r - \epsilon\) 可以逼近 \(A_1\) 中任意元素。由 \(A_1\) 无最大元的性质可验证下组为 \(A_1\)。

\textbf{乘法逆元}：分两种情况。

\textbf{(a) \(\alpha > 0\) 的情形}：设 \(\alpha = (A_1, B_1)\)，\(B_1 = \mathbb{Q} \setminus A_1\)。\(A_1\) 包含正有理数且无最大元，\(B_1\) 无最小元且 \(\forall a \in A_1, b \in B_1\) 有 \(a < b\)。直接构造 \(\alpha^{-1}\) 的分割如下：

\[\alpha^{-1} = (\{q \in \mathbb{Q} : q \leq 0\} \cup \{q \in \mathbb{Q} : q > 0 \text{ 且 } 1/q \notin A_1\},\; \mathbb{Q} \setminus \text{下组})\]

记 \(C = \{q \in \mathbb{Q} : q \leq 0\} \cup \{q \in \mathbb{Q} : q > 0 \text{ 且 } 1/q \notin A_1\}\)。由 \(A_1\) 无最大元且向下封闭可验证 \(C\) 无最大元且向下封闭，故 \(C\) 确为 Dedekind 分割的下组。

验证 \(\alpha \cdot \alpha^{-1} = 1_{\mathbb{R}}\)（即乘积下组为 \(\{q : q < 1\}\)）。由定义 4.4.5(i) 中正实数乘法，乘积下组为 \(\{q \leq 0\} \cup \{xy : x \in A_1, x > 0, y \in C, y > 0\}\)。

\((\subseteq)\)：设 \(xy\) 为乘积下组中正元素（\(x \in A_1, x > 0, y \in C, y > 0\)）。由 \(C\) 的定义，\(1/y \notin A_1\)，即 \(1/y \in B_1\)。若 \(xy \geq 1\)，则 \(x \geq 1/y\)。但 \(A_1\) 向下封闭，\(1/y \in B_1\) 且 \(x \in A_1\) 与 \(x \geq 1/y\) 矛盾（因 \(A_1\) 和 \(B_1\) 满足 \(\forall a \in A_1, b \in B_1, a < b\)）。故 \(xy < 1\)，即 \(xy \in \{q : q < 1\}\)。

\((\supseteq)\)：设 \(0 < q < 1\)。考虑集合 \(T = \{x \in A_1 : x > 0 \text{ 且 } x/q \in A_1\}\)。若 \(T\) 等于 \(A_1\) 中全体正有理数，则对任意 \(x \in A_1\)（\(x > 0\)）有 \(x/q \in A_1\)。反复应用得 \(x/q^n \in A_1\) 对所有 \(n \in \mathbb{N}\) 成立。由于 \(0 < q < 1\)，\(q^n \to 0\)，故 \(x/q^n \to \infty\)，这意味着 \(A_1\) 无上界------与 Dedekind 分割的定义矛盾。因此 \(T \subsetneqq A_1 \cap \mathbb{Q}^+\)，即存在 \(x_0 \in A_1\)、\(x_0 > 0\) 但 \(x_0/q \notin A_1\)。令 \(y = q/x_0 > 0\)，则 \(1/y = x_0/q \notin A_1\)，故 \(y \in C\)。于是 \(x_0 y = q\) 属于乘积的下组。

因此 \(\alpha \cdot \alpha^{-1}\) 的下组为 \(\{q : q < 1\}\)，即 \(1_{\mathbb{R}}\)。

\textbf{(b) \(\alpha < 0\) 的情形}：定义 \(\alpha^{-1} = -(|\alpha|^{-1})\)。由 (a) 知 \(|\alpha|^{-1}\) 存在（\(|\alpha| > 0\)），故 \(\alpha^{-1}\) 存在。由定义 4.4.5(ii) 中负负得正规则： \[\alpha \cdot \alpha^{-1} = \alpha \cdot (-(|\alpha|^{-1})) = -(|\alpha| \cdot (|\alpha|^{-1})) = -(-1_{\mathbb{R}}) = 1_{\mathbb{R}}\] （其中 \(|\alpha| \cdot (|\alpha|^{-1}) = 1_{\mathbb{R}}\) 由 (a) 保证，而 \(-(-1_{\mathbb{R}}) = 1_{\mathbb{R}}\) 由加法逆元性质可得）。\(\blacksquare\)

\subsubsection{4.4.5 完备性}\label{ux5b8cux5907ux6027}

\textbf{定理 4.4.3（有理数在实数中稠密）}：任意两个不同的实数之间存在有理数。

\textbf{证明}：设 \(\alpha < \beta\)，即 \(A_\alpha \subsetneq A_\beta\)（\(A_\alpha\) 是 \(A_\beta\) 的真子集）。取 \(q \in A_\beta \setminus A_\alpha\)。因 \(q \notin A_\alpha\) 且 \(A_\alpha\) 向下封闭，\(q \geq \alpha\)。因 \(q \in A_\beta\)，\(q < \beta\)。

若 \(q > \alpha\)，则 \(\alpha < q < \beta\)，\(q \in \mathbb{Q}\) 即为所求。

若 \(q = \alpha\)（即 \(\alpha\) 是有理数），则 \(\alpha \in A_\beta\)。因 \(A_\beta\) 无最大元，存在 \(q' \in A_\beta\) 使得 \(q' > q = \alpha\)。于是 \(\alpha < q' < \beta\)，\(q' \in \mathbb{Q}\) 即为所求。

因此任意两个实数之间存在有理数。\(\blacksquare\)

\subsubsection{4.4.6 确界原理}\label{ux786eux754cux539fux7406}

\textbf{定义 4.4.6（上确界）}：\(S \subseteq \mathbb{R}\) 的\textbf{上确界} \(\sup S\) 是 \(S\) 的最小上界。

\textbf{定理 4.4.4（确界原理）}：\(\mathbb{R}\) 的每个非空有上界的子集都有上确界。

\textbf{证明}：设 \(S \subseteq \mathbb{R}\) 非空有上界。对每个 \(\alpha \in S\)，设 \(\alpha\) 对应戴德金分割的下组 \(A_\alpha\)。令 \(A^* = \bigcup_{\alpha \in S} A_\alpha\)。我们验证 \((A^*, \mathbb{Q} \setminus A^*)\) 是一个戴德金分割，从而定义了某个实数 \(\beta\)：

\begin{enumerate}
\def\labelenumi{\arabic{enumi}.}
\tightlist
\item
  \textbf{\(A^*\) 非空}：因 \(S\) 非空，取 \(\alpha \in S\)，\(A_\alpha\) 非空，故 \(A^*\) 非空。
\item
  \textbf{\(\mathbb{Q} \setminus A^*\) 非空}：设 \(u\) 是 \(S\) 的上界（即对所有 \(\alpha \in S\) 有 \(\alpha \leq u\)）。则 \(u\) 对应的下组 \(A_u\) 的补集（即 \(\{q \in \mathbb{Q} : q \geq u\}\)）中至少有 \(u\) 本身不在 \(A^*\) 中（因为 \(A^* = \bigcup A_\alpha\) 中的每个元素都属于某个 \(A_\alpha\)，即严格小于某个 \(\alpha\)，而 \(\alpha \leq u\) 意味着 \(A_\alpha \subseteq A_u\)，故 \(u \notin A^*\)）。
\item
  \textbf{\(A^*\) 是下集}：若 \(q \in A^*\)，\(p < q\)，则 \(q \in A_\alpha\) 对某个 \(\alpha\)，由 \(A_\alpha\) 是下集得 \(p \in A_\alpha \subseteq A^*\)。
\item
  \textbf{\(A^*\) 无最大元}：若 \(q \in A^*\)，则 \(q \in A_\alpha\) 对某个 \(\alpha\)。因 \(A_\alpha\) 无最大元，存在 \(q' \in A_\alpha\) 使得 \(q' > q\)，故 \(q' \in A^*\)。
\end{enumerate}

因此 \(\beta = (A^*, \mathbb{Q} \setminus A^*)\) 是一个实数。接下来验证 \(\beta\) 是 \(S\) 的上确界。

\textbf{\(\beta\) 是 \(S\) 的上界}：对任意 \(\alpha \in S\)，\(A_\alpha \subseteq A^*\)，即 \(\alpha \leq \beta\)。

\textbf{\(\beta\) 是最小上界}：设 \(\gamma < \beta\) 是任意实数。则 \(A_\gamma \subsetneq A^*\)，存在 \(q \in A^* \setminus A_\gamma\)。由 \(q \in A^*\)，\(q \in A_\alpha\) 对某个 \(\alpha \in S\)。但 \(q \notin A_\gamma\)，故 \(\gamma < q \leq \alpha\)（因为 \(q \in A_\alpha\)），即 \(\gamma\) 不是 \(S\) 的上界。

综上，\(\beta = \sup S\)。\(\blacksquare\)

\textbf{注记}：确界原理是实数完备性的核心表述，它等价于戴德金分割构造实数的完备性保证。在此基础上，可以证明单调有界定理、区间套定理、有限覆盖定理和柯西收敛准则等分析学的基本定理。

\textbf{定理 4.4.5}：存在唯一的正实数 \(\xi\) 使得 \(\xi^2 = 2\)。

\textbf{证明}：令 \(S = \{q \in \mathbb{Q} : q > 0, q^2 < 2\}\)，\(S\) 非空（\(1 \in S\)）有上界（如 \(2\) 是上界：若 \(q \geq 2\) 则 \(q^2 \geq 4 > 2\)）。设 \(\xi = \sup S\)（由确界原理存在）。

若 \(\xi^2 < 2\)：取 \(h = \frac{2 - \xi^2}{2\xi + 1} > 0\)。先验证 \(h < 1\)：这等价于 \(2 - \xi^2 < 2\xi + 1\)，即 \(1 < \xi^2 + 2\xi = \xi(\xi + 2)\)。由于 \(\xi \geq 1\)（因 \(1 \in S\)），\(\xi(\xi + 2) \geq 1 \cdot 3 = 3 > 1\)，故 \(h < 1\) 成立。因此 \((\xi + h)^2 = \xi^2 + 2\xi h + h^2 = \xi^2 + h(2\xi + h) < \xi^2 + h(2\xi + 1) = \xi^2 + (2 - \xi^2) = 2\)。故 \(\xi + h \in S\)，与 \(\xi\) 是上界矛盾。

若 \(\xi^2 > 2\)：取 \(h = \frac{\xi^2 - 2}{2\xi} > 0\)。则 \((\xi - h)^2 = \xi^2 - 2\xi h + h^2 > \xi^2 - 2\xi h = 2\)。故 \(\xi - h\) 也是 \(S\) 的上界，与 \(\xi\) 是最小上界矛盾。

因此 \(\xi^2 = 2\)。

唯一性：设 \(0 < x < y\)，\(x^2 = y^2 = 2\)。则 \(0 = y^2 - x^2 = (y-x)(y+x)\)。由 \(y + x > 0\) 得 \(y - x = 0\)，即 \(x = y\)。\(\blacksquare\)

\textbf{定理 4.4.6（阿基米德性质）}：对任意正实数 \(x, y\)，存在自然数 \(n\) 使得 \(nx > y\)。

\textbf{证明}：假设不然，\(S = \{nx : n \in \mathbb{N}\}\) 有上界 \(y\)。设 \(\xi = \sup S\)（由确界原理存在）。由于 \(x > 0\)，\(\xi - x < \xi\)，故存在 \(n_0 x > \xi - x\)（否则 \(\xi - x\) 也是上界，与 \(\xi\) 是最小上界矛盾），即 \((n_0 + 1)x > \xi\)，但 \((n_0 + 1)x \in S\)，与 \(\xi\) 是上界矛盾。\(\blacksquare\)

\textbf{推论 4.4.1}：对任意正实数 \(\epsilon\)，存在自然数 \(n\) 使得 \(\frac{1}{n} < \epsilon\)。

\textbf{证明}：在阿基米德性质中取 \(x = \epsilon\), \(y = 1\)，得 \(n\epsilon > 1\)，即 \(\frac{1}{n} < \epsilon\)。\(\blacksquare\)

阿基米德性质的一个直接推论是 \(\frac{1}{n}\) 可以任意接近于零------这为极限理论中''\(\epsilon\)-\(N\)定义''的合理性提供了严格保证。事实上，阿基米德性质与确界原理等价地刻画了实数的完备性：一个有序域是完备的当且仅当它满足确界原理且具有阿基米德性质。

\textbf{推论 4.4.2（实数的阿基米德性质的等价形式）}：\(\lim_{n \to \infty} \frac{1}{n} = 0\)。

\textbf{证明}：对任意 \(\epsilon > 0\)，由推论 4.4.1 存在 \(N\) 使得 \(\frac{1}{N} < \epsilon\)。对 \(n \geq N\)，\(0 < \frac{1}{n} \leq \frac{1}{N} < \epsilon\)。\(\blacksquare\)

\textbf{注记 4.4.1}：阿基米德性质表明实数系中没有''无穷小''的正实数。这是实数系与非标准分析中的超实数系的关键区别。

\hrulefill

\section{第13章：复数}\label{ux7b2c13ux7ae0ux590dux6570}

从 \(\mathbb{R}\) 到 \(\mathbb{C}\) 的扩展是数系扩展的最后一步（在通常的代数框架下）。虽然 \(\mathbb{R}\) 已经是完备有序域，但它在代数上并不封闭：多项式方程 \(x^2 + 1 = 0\) 在 \(\mathbb{R}\) 中没有解。复数构造的目标是创造 \(\mathbb{R}\) 的一个域扩张，使得每个非常数多项式都有根（代数闭包）。与之前从 \(\mathbb{Q}\) 到 \(\mathbb{R}\) 的构造不同，从 \(\mathbb{R}\) 到 \(\mathbb{C}\) 的构造纯粹是代数的------我们不需要处理''完备性''问题，因为 \(\mathbb{C}\) 自然地继承了 \(\mathbb{R}^2\) 的完备性。本节将用 Hamilton 的序对方法严格定义复数，证明 \(\mathbb{C}\) 构成一个域，引入虚数单位 \(i\) 并建立 \(a + bi\) 的标准表示，讨论复数的共轭、模和基本运算性质，并证明代数基本定理------\(\mathbb{C}\) 是代数闭域。复数理论的发展史是数学史上的一个精彩篇章：从 16 世纪 Cardano 在解三次方程时被迫引入复数，到 Gauss 证明代数基本定理，再到 Hamilton 给出严格定义，复数经历了从''神秘''到''合法''的转变。

\subsubsection{4.5.1 构造动机}\label{ux6784ux9020ux52a8ux673a-3}

在 \(\mathbb{R}\) 中，\(x^2 + 1 = 0\) 无解（因 \(x^2 \geq 0\)）。我们构造更大的数系使所有多项式方程有解。

从 \(\mathbb{R}\) 到 \(\mathbb{C}\) 的扩展是数系扩展的最后一步（在通常的代数框架下）。虽然 \(\mathbb{R}\) 已经是完备有序域，但它在代数上并不封闭：多项式方程 \(x^2 + 1 = 0\) 在 \(\mathbb{R}\) 中没有解。复数构造的目标是创造 \(\mathbb{R}\) 的一个域扩张，使得每个非常数多项式都有根（代数闭包）。与之前从 \(\mathbb{Q}\) 到 \(\mathbb{R}\) 的构造不同，从 \(\mathbb{R}\) 到 \(\mathbb{C}\) 的构造纯粹是代数的------我们不需要处理''完备性''问题，因为 \(\mathbb{C}\) 自然地继承了 \(\mathbb{R}^2\) 的完备性。复数在 16 世纪由 Cardano 等人在解三次方程时首次引入，但直到 19 世纪才被完全接受为合法的数学对象。Gauss 在 1799 年的博士论文中证明了代数基本定理，Hamilton 在 1833 年给出了复数作为有序对 \((a, b)\) 的严格定义，彻底消除了复数概念的''神秘色彩''。

从历史角度看，复数是最''具争议''的数系扩展。Cardano 在 1545 年的《大术》（Ars Magna）中首次使用复数来解三次方程，但他本人称其为''诡辩的''（sophistic）。Bombelli 在 1572 年首次制定了复数的运算规则，但仍称之为''不可约的''（irreducible）。直到 19 世纪，Gauss 的代数基本定理、Argand 和 Wessel 的复平面表示、以及 Hamilton 的序对定义，才逐步确立了复数的合法性。代数基本定理------每个非常数复系数多项式在 \(\mathbb{C}\) 中至少有一个根------是复数理论的顶峰，它表明 \(\mathbb{C}\) 是''最优的''数系：任何进一步的代数扩张都不可能产生真正的新数系（因为 \(\mathbb{C}\) 已经是代数闭域）。这一事实使 \(\mathbb{C}\) 在数系扩展中占据了一个终点位置，除非我们考虑四元数（\(\mathbb{H}\)）等非交换的代数结构。

\subsubsection{4.5.2 复数的定义}\label{ux590dux6570ux7684ux5b9aux4e49}

\textbf{定义 4.5.1}：\(\mathbb{C} = \mathbb{R} \times \mathbb{R}\)，运算定义为： - \((a, b) + (c, d) = (a + c, b + d)\) - \((a, b) \cdot (c, d) = (ac - bd, ad + bc)\) - \(0_{\mathbb{C}} = (0, 0)\)，\(1_{\mathbb{C}} = (1, 0)\)

\textbf{定义 4.5.2}：虚数单位 \(i = (0, 1)\)。

\textbf{定理 4.5.1}：\(i^2 = -1\)。

\textbf{证明}：\(i^2 = (0, 1) \cdot (0, 1) = (0 \cdot 0 - 1 \cdot 1, 0 \cdot 1 + 1 \cdot 0) = (-1, 0) = -1_{\mathbb{C}}\)。\(\blacksquare\)

Hamilton 的定义将复数彻底还原为实数有序对，使复数不再需要''想象''（imaginary 一词的由来）。加法就是 \(\mathbb{R}^2\) 上的向量加法（对应复数的平行四边形法则），乘法则是带有''旋转''效果的运算：\((a, b) \cdot (c, d)\) 可以理解为将向量 \((a, b)\) 以 \((c, d)\) 的方式进行''旋转和缩放''。\(\mathbb{R}\) 作为 \(\mathbb{C}\) 的子集嵌入为 \(\{(a, 0) : a \in \mathbb{R}\}\)，而 \(i = (0, 1)\) 的引入使得每个复数可以唯一地写作 \(a + bi\) 的形式。这一记号系统揭示了 \(\mathbb{C}\) 作为 \(\mathbb{R}\) 的二维向量空间的结构：\(\{1, i\}\) 是 \(\mathbb{C}\) 在 \(\mathbb{R}\) 上的一组基。复数乘法中的 \(i^2 = -1\) 规则完全决定了乘法的结构，因此 \(\mathbb{C}\) 可以等价地定义为 \(\mathbb{R}[x]/(x^2+1)\)------多项式环 \(\mathbb{R}[x]\) 模去理想 \((x^2+1)\) 的商环，这是域扩张理论中''添加一个根''的标准构造。

\subsubsection{4.5.3 复数构成域}\label{ux590dux6570ux6784ux6210ux57df}

要证明 \(\mathbb{C}\) 构成一个域，我们需要验证域的所有公理：加法构成阿贝尔群、乘法（除零外）构成阿贝尔群、分配律成立。这些验证虽然繁琐但都是直接的，因为 \(\mathbb{C}\) 的加法和乘法定义完全依赖 \(\mathbb{R}\) 的域结构。验证的关键在于复数乘法的定义 \((a, b) \cdot (c, d) = (ac - bd, ad + bc)\) 是如何产生的：它实际上是将 \((a, b)\) 视为 \(a + bi\)，然后按普通代数规则展开 \((a + bi)(c + di) = ac + adi + bci + bdi^2 = (ac - bd) + (ad + bc)i\)，再利用 \(i^2 = -1\) 化简。因此，\(\mathbb{C}\) 的乘法定义并非随意规定的，而是确保 \(\mathbb{C}\) 成为 \(\mathbb{R}\) 的域扩张的唯一可能方式。在证明 \(\mathbb{C}\) 是域后，我们还将给出每个复数的标准表示 \(a + bi\)，并证明这一表示是唯一的------这意味着 \(\{1, i\}\) 是 \(\mathbb{C}\) 作为 \(\mathbb{R}\)-向量空间的一组基，从而 \(\dim_{\mathbb{R}} \mathbb{C} = 2\)。

\textbf{定理 4.5.2}：每个复数 \(z = (a, b)\) 都可以唯一地表示为 \(z = a + bi\)。

\textbf{证明}：\((a, b) = (a, 0) + (b, 0) \cdot (0, 1) = a + bi\)。唯一性：若 \(a + bi = c + di\)，则 \((a, b) = (c, d)\)，故 \(a = c, b = d\)。\(\blacksquare\)

\subsubsection{4.5.3 复数构成域}\label{ux590dux6570ux6784ux6210ux57df-1}

\textbf{定理 4.5.3}：\((\mathbb{C}, +, \cdot)\) 是域。

\textbf{证明}：加法阿贝尔群由实数性质直接继承。乘法结合律通过展开验证。逆元：对 \(z = a + bi \neq 0\)，\(z^{-1} = \frac{a - bi}{a^2 + b^2}\)。验证 \(z \cdot z^{-1} = \frac{a^2 + b^2}{a^2 + b^2} = 1\)。\(\blacksquare\)

\subsubsection{4.5.4 模与共轭}\label{ux6a21ux4e0eux5171ux8f6d}

\textbf{定义 4.5.3（共轭）}：\(\bar{z} = a - bi\)。

\textbf{定义 4.5.4（模）}：\(|z| = \sqrt{a^2 + b^2}\)。

共轭和模是复数理论中最基本的两个概念，它们将复数与实数联系起来：\(z \cdot \bar{z} = |z|^2\) 是一个非负实数，这为复数除法的计算和模的性质证明提供了便捷途径。以下共轭运算的基本性质反映了它是一个域自同构------即共轭与加法和乘法相容。

\textbf{定理 4.5.4（共轭的性质）}： 1. \(\overline{z_1 + z_2} = \bar{z}_1 + \bar{z}_2\) 2. \(\overline{z_1 \cdot z_2} = \bar{z}_1 \cdot \bar{z}_2\) 3. \(z \cdot \bar{z} = |z|^2\)

\textbf{证明}：(1) 和 (2) 由定义直接计算。(3) \(z \cdot \bar{z} = (a + bi)(a - bi) = a^2 + b^2 = |z|^2\)。\(\blacksquare\)

共轭运算的性质为研究复数的模提供了基础。模作为复数的''大小''度量，继承了实数绝对值的基本性质，并且具有一个在复数乘法中极为重要的特征------模是乘性的。

\textbf{定理 4.5.5（模的性质）}： 1. \(|z| \geq 0\)，且 \(|z| = 0 \iff z = 0\) 2. \(|z_1 \cdot z_2| = |z_1| \cdot |z_2|\) 3. \(|z_1 + z_2| \leq |z_1| + |z_2|\)（三角不等式）

\textbf{证明}：(1) 由定义直接得到。(2) \(|z_1 z_2|^2 = (z_1 z_2)\overline{(z_1 z_2)} = z_1 z_2 \bar{z}_1 \bar{z}_2 = (z_1 \bar{z}_1)(z_2 \bar{z}_2) = |z_1|^2 |z_2|^2\)。(3) \(|z_1 + z_2|^2 = |z_1|^2 + 2\operatorname{Re}(z_1 \bar{z}_2) + |z_2|^2 \leq |z_1|^2 + 2|z_1||z_2| + |z_2|^2 = (|z_1| + |z_2|)^2\)。\(\blacksquare\)

\subsubsection{4.5.5 复数不是有序域}\label{ux590dux6570ux4e0dux662fux6709ux5e8fux57df}

\textbf{定理 4.5.6}：不存在全序关系 ``\(<\)'' 使得 \((\mathbb{C}, +, \cdot, <)\) 成为有序域。

\textbf{证明}：假设存在。由于 \(i \neq 0\)，必有 \(i > 0\) 或 \(i < 0\)。若 \(i > 0\)，由有序域中正数的乘积为正，\(i^2 = -1 > 0\)。再由 \((-1)(-1) = 1 > 0\)，以及正数之和为正，\(0 = (-1) + 1 > 0\)，但同时 \(1 > 0\)，与三歧性矛盾（\(0 > 0\) 不可能）。若 \(i < 0\)，则 \(-i > 0\)，\((-i)^2 = i^2 = -1 > 0\)，同理推出矛盾。\(\blacksquare\)

定理 4.5.6 揭示了 \(\mathbb{C}\) 与 \(\mathbb{R}\) 的一个根本性差异：\(\mathbb{R}\) 是唯一完备的有序域（在同构意义下），但 \(\mathbb{C}\) 不能赋予任何与域结构相容的全序。这一事实的含义是深远的：当我们从 \(\mathbb{R}\) 扩展到 \(\mathbb{C}\) 时，虽然获得了代数封闭性，但损失了序结构。在 \(\mathbb{C}\) 中，我们不能说''\(z_1 < z_2\)``（除非限定在 \(\mathbb{R}\) 的子集上讨论）。这解释了为什么复分析中的不等式通常涉及模（\(|z|\)）而非复数本身。复数不能比较大小也是其''二维''本质的体现：\(\mathbb{R}^2\) 上不存在与向量空间结构和乘法都相容的全序。这一事实在非交换几何和非标准分析中有着深刻的推广。

\subsubsection{4.5.6 复数的极坐标表示}\label{ux590dux6570ux7684ux6781ux5750ux6807ux8868ux793a}

\textbf{定义 4.5.5}：对复数 \(z = a + bi\)，定义其\textbf{辐角} \(\theta\) 为满足 \(\cos\theta = \frac{a}{|z|}\)，\(\sin\theta = \frac{b}{|z|}\) 的角度（当 \(z \neq 0\) 时）。则 \(z\) 可表示为 \(z = |z|(\cos\theta + i\sin\theta)\)。

\textbf{定理 4.5.7（复数乘法的几何意义）}：设 \(z_1 = r_1 e^{i\theta_1}\), \(z_2 = r_2 e^{i\theta_2}\)。则 \(z_1 z_2 = r_1 r_2 e^{i(\theta_1 + \theta_2)}\)。即复数相乘时，模相乘，辐角相加。

\textbf{证明}：\(z_1 z_2 = r_1(\cos\theta_1 + i\sin\theta_1) \cdot r_2(\cos\theta_2 + i\sin\theta_2)\)。展开得 \(r_1 r_2[(\cos\theta_1\cos\theta_2 - \sin\theta_1\sin\theta_2) + i(\sin\theta_1\cos\theta_2 + \cos\theta_1\sin\theta_2)]\)。由三角函数的和角公式，\(= r_1 r_2[\cos(\theta_1 + \theta_2) + i\sin(\theta_1 + \theta_2)]\)。\(\blacksquare\)

复数乘法的几何意义（模相乘、辐角相加）是复数极坐标表示的最核心结论。这一性质直接给出了复数幂运算的简洁公式------棣莫弗定理------它揭示了 \((\cos\theta + i\sin\theta)^n\) 与 \(\cos(n\theta) + i\sin(n\theta)\) 之间的恒等关系，在三角恒等式的推导中具有无可替代的价值。

\textbf{定理 4.5.8（棣莫弗定理）}：对任意正整数 \(n\) 和实数 \(\theta\)： \[(\cos\theta + i\sin\theta)^n = \cos(n\theta) + i\sin(n\theta)\]

\textbf{证明}：对 \(n\) 用数学归纳法。\(n = 1\) 时显然。设 \((\cos\theta + i\sin\theta)^k = \cos(k\theta) + i\sin(k\theta)\)。则 \[(\cos\theta + i\sin\theta)^{k+1} = (\cos(k\theta) + i\sin(k\theta))(\cos\theta + i\sin\theta)\] 由定理 4.5.7，\(= \cos((k+1)\theta) + i\sin((k+1)\theta)\)。\(\blacksquare\)

\textbf{定理 4.5.9（单位根）}：方程 \(z^n = 1\) 恰有 \(n\) 个不同的复数解，称为 \(n\) 次\textbf{单位根}： \[z_k = e^{2\pi ik/n} = \cos\frac{2\pi k}{n} + i\sin\frac{2\pi k}{n}, \quad k = 0, 1, \ldots, n-1\]

\textbf{证明}：设 \(z = re^{i\theta}\)。\(z^n = r^n e^{in\theta} = 1\) 意味着 \(r^n = 1\)（故 \(r = 1\)）且 \(n\theta = 2\pi k\)（\(k \in \mathbb{Z}\)），即 \(\theta = 2\pi k/n\)。不同的 \(k\) 给出不同的 \(z_k\) 当且仅当 \(k\) 在模 \(n\) 意义下不同，故恰有 \(n\) 个解。\(\blacksquare\)

复数的极坐标表示是复数理论的精华所在。通过 Euler 公式 \(e^{i\theta} = \cos\theta + i\sin\theta\)，复数的乘法被转化为模的相乘和辐角的相加------这正是定理 4.5.7 的几何解释。棣莫弗定理（定理 4.5.8）是 \(n\) 次幂的极坐标形式，它极大地简化了 \((\cos\theta + i\sin\theta)^n\) 的计算，避免了繁琐的二项式展开。单位根（定理 4.5.9）在复平面上均匀分布在单位圆上，形成一个正 \(n\) 边形，这种几何直观使得许多代数问题（如多项式因式分解、离散 Fourier 变换）变得清晰。极坐标表示还揭示了复数的拓扑结构：\(\mathbb{C} \setminus \{0\}\) 作为乘法群同构于 \(\mathbb{R}_{>0} \times S^1\)（正实数乘法群与单位圆的直积），而指数映射 \(z \mapsto e^z\) 是加法群 \(\mathbb{C}\) 到乘法群 \(\mathbb{C}^\times\) 的满射群同态。

\subsubsection{4.5.7 代数基本定理（叙述）}\label{ux4ee3ux6570ux57faux672cux5b9aux7406ux53d9ux8ff0}

\textbf{定理 4.5.10（代数基本定理）}：每个次数 \(\geq 1\) 的复系数多项式在 \(\mathbb{C}\) 中至少有一个根。

\textbf{等价形式}：每个次数 \(\geq 1\) 的复系数多项式在 \(\mathbb{C}\) 中可以完全分解为一次因式的乘积。

\textbf{注记 4.5.1}：代数基本定理的严格证明需要复分析（如刘维尔定理或辐角原理），超出本章范围。但此定理的结论极其重要：它表明 \(\mathbb{C}\) 是\textbf{代数封闭域}------任何多项式方程在 \(\mathbb{C}\) 中都有解，不需要再扩展数系。

代数基本定理是数系扩展历程的终点------它表明 \(\mathbb{C}\) 是代数封闭的，因此不再需要为了解多项式方程而创造新的数系。Gauss 在其 1799 年的博士论文中首次给出了一个近似的证明，随后在 1816 年、1849 年又给出了两个不同的证明。此定理有数十种不同的证明方法，包括基于复分析（Liouville 定理、辐角原理）、代数拓扑（环绕数）、Galois 理论（奇数次实多项式有实根 + 二次方程有复根）等途径。代数基本定理的一个推论是代数学基本定理：每个 \(n\) 次复系数多项式可以分解为 \(n\) 个一次因式的乘积，即 \(p(z) = a_n(z - z_1)(z - z_2)\cdots(z - z_n)\)。这表明 \(\mathbb{C}\) 是代数闭域，也是 \(\mathbb{R}\) 的代数闭包。实数域 \(\mathbb{R}\) 的有限维域扩张只有 \(\mathbb{R}\) 本身和 \(\mathbb{C}\)（Frobenius 定理），而 \(\mathbb{C}\) 的有限维可除代数只有 \(\mathbb{C}\) 本身（Gelfand-Mazur 定理）。

\subsection{本章展望}\label{ux672cux7ae0ux5c55ux671b-4}

本章完成了从 \(\mathbb{N}\) 到 \(\mathbb{C}\) 的严格构造。在 Ch 5 中，我们将深入研究整数 \(\mathbb{Z}\) 的内部结构------数论。在卷二（Ch 12--11）中，数系的代数性质将被系统发展为代数运算与多项式理论。在卷四（Ch 23--24）中，实数的完备性（确界原理）将被用于建立极限和微积分的基础。复数的代数封闭性将在代数基本定理中得到完整阐述。

\hrulefill

\hrulefill

\hrulefill

\hrulefill

\hrulefill

\newpage

\chapter{01-基础/02-集合与数系/04-格论与序结构.md}\label{ux57faux784002-ux96c6ux5408ux4e0eux6570ux7cfb04-ux683cux8bbaux4e0eux5e8fux7ed3ux6784.md}

\chapter{格论与序结构}\label{ux683cux8bbaux4e0eux5e8fux7ed3ux6784}

\hrulefill

\section{第14章：偏序集与格}\label{ux7b2c14ux7ae0ux504fux5e8fux96c6ux4e0eux683c}

格论（Lattice Theory）是序结构与代数结构交汇的数学分支，其根源可追溯至Dedekind关于理想的研究和Peirce关于关系代数的探索。1930年代，Birkhoff的经典著作奠定了现代格论的基础。本章系统阐述偏序集、格、完备格以及相关的代数结构与逻辑联系。

\subsection{88.1 偏序集的基本概念}\label{ux504fux5e8fux96c6ux7684ux57faux672cux6982ux5ff5}

\textbf{定义 88.1.1（偏序集）} 设 \(P\) 为一个集合，\(\leq\) 是 \(P\) 上的二元关系。若 \(\leq\) 满足以下条件，则称 \((P, \leq)\) 为一个偏序集（partially ordered set，简称poset）：

\begin{enumerate}
\def\labelenumi{\arabic{enumi}.}
\tightlist
\item
  （自反性）对所有 \(x \in P\)，\(x \leq x\)；
\item
  （反对称性）若 \(x \leq y\) 且 \(y \leq x\)，则 \(x = y\)；
\item
  （传递性）若 \(x \leq y\) 且 \(y \leq z\)，则 \(x \leq z\)。
\end{enumerate}

若对任意 \(x, y \in P\)，总有 \(x \leq y\) 或 \(y \leq x\)，则称 \(\leq\) 为全序（total order）或线性序（linear order），\((P, \leq)\) 称为全序集（totally ordered set）或链（chain）。

\textbf{定义 88.1.2（严格偏序）} 严格偏序 \(<\) 定义为：\(x < y\) 当且仅当 \(x \leq y\) 且 \(x \neq y\)。严格偏序满足非自反性和传递性。

\textbf{定义 88.1.3（覆盖关系）} 设 \((P, \leq)\) 为偏序集，\(x, y \in P\)。若 \(x < y\) 且不存在 \(z \in P\) 使得 \(x < z < y\)，则称 \(y\) 覆盖 \(x\)（或 \(x\) 被 \(y\) 覆盖），记作 \(x \prec y\)。

\textbf{定义 88.1.4（Hasse图）} 偏序集 \((P, \leq)\) 的Hasse图是一个有向无环图的图形表示，其中： - 每个元素 \(x \in P\) 对应平面上的一个点； - 若 \(x \prec y\)，则画一条从 \(x\) 到 \(y\) 的线段（\(y\) 位于 \(x\) 的上方）。

Hasse图完全确定了有限偏序集的结构，因为 \(x \leq y\) 当且仅当存在一条从 \(x\) 到 \(y\) 的上升路径。Hasse图是可视化偏序集的最有力工具。

\textbf{定义 88.1.5（极大元、极小元、最大元、最小元）} 设 \((P, \leq)\) 为偏序集。

\begin{itemize}
\tightlist
\item
  \(m \in P\) 为极大元，若不存在 \(x \in P\) 使得 \(m < x\)；
\item
  \(m \in P\) 为极小元，若不存在 \(x \in P\) 使得 \(x < m\)；
\item
  \(m \in P\) 为最大元，若对所有 \(x \in P\)，\(x \leq m\)；
\item
  \(m \in P\) 为最小元，若对所有 \(x \in P\)，\(m \leq x\)。
\end{itemize}

最大元和最小元若存在，必唯一。极大元和极小元不一定唯一。

\textbf{定义 88.1.6（上界、下界、上确界、下确界）} 设 \((P, \leq)\) 为偏序集，\(S \subseteq P\)。

\begin{itemize}
\tightlist
\item
  \(u \in P\) 为 \(S\) 的上界，若对所有 \(s \in S\)，\(s \leq u\)；
\item
  \(\ell \in P\) 为 \(S\) 的下界，若对所有 \(s \in S\)，\(\ell \leq s\)；
\item
  \(S\) 的所有上界构成的集合若有最小元，则称其为 \(S\) 的上确界（supremum），记作 \(\bigvee S\) 或 \(\sup S\)；
\item
  \(S\) 的所有下界构成的集合若有最大元，则称其为 \(S\) 的下确界（infimum），记作 \(\bigwedge S\) 或 \(\inf S\)。
\end{itemize}

当 \(S = \{x, y\}\) 时，记 \(x \vee y = \sup\{x, y\}\)，\(x \wedge y = \inf\{x, y\}\)。

\textbf{定义 88.1.7（良序）} 设 \((P, \leq)\) 为全序集。若 \(P\) 的每个非空子集都有最小元，则称 \((P, \leq)\) 为良序集（well-ordered set），\(\leq\) 为 \(P\) 上的良序。

\textbf{定理 88.1.8（Zorn引理）} 设 \((P, \leq)\) 为非空偏序集。若 \(P\) 中的每条链都有上界，则 \(P\) 中至少有一个极大元。

Zorn引理等价于选择公理，是数学中构造极大对象的基本工具。其证明依赖于选择公理和超限归纳法。

\textbf{命题 88.1.9（对偶原理）} 若 \(\Phi\) 是偏序集理论中为真的陈述，则其对偶陈述 \(\Phi^{\partial}\)（将 \(\leq\) 替换为 \(\geq\)，将 \(\vee\) 替换为 \(\wedge\)，将 \(\sup\) 替换为 \(\inf\) 等）亦为真。

\textbf{定义 88.1.10（保序映射与序同构）} 设 \((P, \leq_P)\) 和 \((Q, \leq_Q)\) 为偏序集。映射 \(f: P \to Q\) 称为保序的（order-preserving）或单调的，若 \(x \leq_P y\) 蕴涵 \(f(x) \leq_Q f(y)\)。若 \(f\) 为双射且 \(f\) 和 \(f^{-1}\) 均保序，则称 \(f\) 为序同构（order-isomorphism）。

\textbf{定义 88.1.11（Dilworth定理）} 设 \((P, \leq)\) 为有限偏序集。\(P\) 可分解为 \(k\) 个两两不交的链之并的最小 \(k\) 等于 \(P\) 中反链（两两不可比元素的子集）的最大基数。

Dilworth定理是组合序理论中的核心结果，与匹配理论、Hall婚配定理有深刻联系。

\textbf{定义 88.1.12（理想与滤子）} 设 \((P, \leq)\) 为偏序集。

\begin{itemize}
\tightlist
\item
  子集 \(I \subseteq P\) 称为序理想（order ideal）或下集（down-set），若 \(x \in I\) 且 \(y \leq x\) 蕴涵 \(y \in I\)；
\item
  子集 \(F \subseteq P\) 称为序滤子（order filter）或上集（up-set），若 \(x \in F\) 且 \(x \leq y\) 蕴涵 \(y \in F\)。
\end{itemize}

\(P\) 的所有序理想构成的集合 \(O(P)\) 在包含关系下构成一个完备格。

\subsection{88.2 格的定义与基本性质}\label{ux683cux7684ux5b9aux4e49ux4e0eux57faux672cux6027ux8d28}

\textbf{定义 88.2.1（格------序定义）} 偏序集 \((L, \leq)\) 称为格（lattice），若对任意 \(x, y \in L\)，\(x \vee y\) 和 \(x \wedge y\) 均存在。

\textbf{定义 88.2.2（格------代数定义）} 格是一个代数结构 \((L, \vee, \wedge)\)，其中 \(\vee\) 和 \(\wedge\) 是 \(L\) 上的二元运算，满足以下恒等式：

\begin{enumerate}
\def\labelenumi{\arabic{enumi}.}
\tightlist
\item
  （幂等律）\(x \vee x = x\)，\(x \wedge x = x\)；
\item
  （交换律）\(x \vee y = y \vee x\)，\(x \wedge y = y \wedge x\)；
\item
  （结合律）\((x \vee y) \vee z = x \vee (y \vee z)\)，\((x \wedge y) \wedge z = x \wedge (y \wedge z)\)；
\item
  （吸收律）\(x \vee (x \wedge y) = x\)，\(x \wedge (x \vee y) = x\)。
\end{enumerate}

序定义和代数定义等价：若 \((L, \vee, \wedge)\) 满足上述代数条件，定义 \(x \leq y \iff x \vee y = y\)（等价于 \(x \wedge y = x\)），则 \((L, \leq)\) 为格且 \(\vee\)、\(\wedge\) 分别为并、交运算。

\textbf{命题 88.2.3（格的基本性质）} 在任意格 \(L\) 中，对任意 \(x, y, z \in L\)：

\begin{enumerate}
\def\labelenumi{\arabic{enumi}.}
\tightlist
\item
  \(x \leq y\) 蕴涵 \(x \vee z \leq y \vee z\) 和 \(x \wedge z \leq y \wedge z\)；
\item
  \(x \vee (y \wedge z) \leq (x \vee y) \wedge (x \vee z)\)；
\item
  \((x \wedge y) \vee (x \wedge z) \leq x \wedge (y \vee z)\)。
\end{enumerate}

\textbf{证明}：（1）因 \(y \leq y \vee z\) 且 \(x \leq y \leq y \vee z\)，又 \(z \leq y \vee z\)，故 \(x \vee z \leq y \vee z\)。下确界情形类似。（2）\(x \leq x \vee y\) 且 \(x \leq x \vee z\)，故 \(x \leq (x \vee y) \wedge (x \vee z)\)；又 \(y \wedge z \leq y \leq x \vee y\) 且 \(y \wedge z \leq z \leq x \vee z\)，故 \(y \wedge z \leq (x \vee y) \wedge (x \vee z)\)。因此 \(x \vee (y \wedge z) \leq (x \vee y) \wedge (x \vee z)\)。（3）的证明与（2）对偶。\(\blacksquare\)

\textbf{定义 88.2.4（分配格）} 格 \(L\) 称为分配格（distributive lattice），若对任意 \(x, y, z \in L\)，满足分配律：

\[x \vee (y \wedge z) = (x \vee y) \wedge (x \vee z),\] \[x \wedge (y \vee z) = (x \wedge y) \vee (x \wedge z).\]

实际上，由对偶原理，上述两个恒等式等价。

\textbf{定理 88.2.5（分配格的刻画）} 格 \(L\) 是非分配的当且仅当 \(L\) 包含一个子格同构于 \(M_3\)（钻石格）或 \(N_5\)（五边形格）。

\textbf{证明概要}：\(M_3\) 和 \(N_5\) 是仅有的两个极小非分配格。若 \(L\) 不包含 \(M_3\) 或 \(N_5\) 作为子格，则可证明 \(L\) 是分配的；反之，若 \(L\) 包含其中之一，则分配律在此子格上失效，因而在全格上也失效。\(\blacksquare\)

\textbf{定义 88.2.6（模格）} 格 \(L\) 称为模格（modular lattice），若对任意 \(x, y, z \in L\)，满足模律：

\[x \leq z \implies x \vee (y \wedge z) = (x \vee y) \wedge z.\]

每个分配格都是模格，但反之不然。\(M_3\) 是模格但非分配格，而 \(N_5\) 既非分配亦非模格。Dedekind在研究理想时首次引入模格概念。

\textbf{定理 88.2.7（模格的刻画）} 格 \(L\) 是模格当且仅当 \(L\) 不包含同构于 \(N_5\) 的子格。

\textbf{证明} （\(\Rightarrow\)）若 \(L\) 包含同构于 \(N_5\) 的子格，则模律在此子格上失效：\(N_5 = \{0, a, b, c, 1\}\) 满足 \(a < c\)，\(b\) 与 \(a\) 和 \(c\) 不可比，且 \(a \vee b = c \vee b = 1\)，\(a \wedge b = c \wedge b = 0\)。则 \(a \vee (b \wedge c) = a \vee 0 = a\)，而 \((a \vee b) \wedge c = 1 \wedge c = c\)，故 \(a \neq c\) 时模律不成立，矛盾于 \(L\) 是模格。故模格不含 \(N_5\)。

（\(\Leftarrow\)）设 \(L\) 不含 \(N_5\) 子格。若 \(L\) 非模格，则存在 \(x, y, z \in L\) 满足 \(x \leq z\) 但 \(x \vee (y \wedge z) < (x \vee y) \wedge z\)。令 \(a = x \vee (y \wedge z)\)，\(b = y\)，\(c = (x \vee y) \wedge z\)。可验证 \(a < c\)，\(b\) 与 \(a, c\) 均不可比，且 \(a \vee b = x \vee y = c \vee b\)（因 \(a \leq x \vee y\) 且 \(y \wedge z \leq y\)），\(a \wedge b = y \wedge z = c \wedge b\)。由此 \(\{a \wedge b, a, b, c, a \vee b\}\) 同构于 \(N_5\)，矛盾。故 \(L\) 是模格。\(\blacksquare\)

\textbf{定义 88.2.8（补元与有补格）} 设 \(L\) 为具有最小元 \(0\) 和最大元 \(1\) 的格。元素 \(x \in L\) 的补元（complement）是满足 \(x \vee y = 1\) 且 \(x \wedge y = 0\) 的 \(y \in L\)。若 \(L\) 中每个元素都有补元，则称 \(L\) 为有补格（complemented lattice）。在有补分配格中，每个元素的补元唯一，这样的格就是Bool代数。

\textbf{定理 88.2.9（Birkhoff表示定理）} 每个有限分配格同构于某个偏序集的所有序理想的格。

\textbf{证明概要}：设 \(L\) 为有限分配格。令 \(J(L)\) 为 \(L\) 的并不可约元（join-irreducibles）构成的子偏序集。定义映射 \(\phi: L \to O(J(L))\) 为 \(\phi(x) = \{j \in J(L) : j \leq x\}\)。则 \(\phi\) 为格同构。这建立了有限分配格与有限偏序集之间的一一对应。\(\blacksquare\)

\textbf{定义 88.2.10（同余与商格）} 格 \(L\) 上的一个同余关系（congruence）是 \(L\) 上的等价关系 \(\theta\)，满足：若 \(a \theta b\) 且 \(c \theta d\)，则 \((a \vee c) \theta (b \vee d)\) 且 \((a \wedge c) \theta (b \wedge d)\)。\(L\) 关于 \(\theta\) 的商 \(L/\theta\) 自然地构成格。

\textbf{定理 88.2.11（格同态基本定理）} 若 \(f: L \to M\) 为格满同态，则 \(\ker f = \{(x, y): f(x) = f(y)\}\) 是 \(L\) 上的同余，且 \(L/\ker f \cong M\)。

\textbf{证明} 首先验证 \(\ker f\) 是同余关系。\(\ker f\) 是等价关系（由 \(f\) 是映射直接得到）。若 \((a, b) \in \ker f\) 且 \((c, d) \in \ker f\)，即 \(f(a) = f(b)\) 且 \(f(c) = f(d)\)，则 \(f(a \vee c) = f(a) \vee f(c) = f(b) \vee f(d) = f(b \vee d)\)，故 \((a \vee c, b \vee d) \in \ker f\)。同理 \((a \wedge c, b \wedge d) \in \ker f\)。故 \(\ker f\) 是同余关系。

定义 \(\bar{f}: L/\ker f \to M\) 为 \(\bar{f}([x]) = f(x)\)。此映射良定义因 \((x, y) \in \ker f\) 时 \(f(x) = f(y)\)。\(\bar{f}\) 为同态因 \(\bar{f}([x] \vee [y]) = \bar{f}([x \vee y]) = f(x \vee y) = f(x) \vee f(y) = \bar{f}([x]) \vee \bar{f}([y])\)，\(\wedge\) 情形类似。\(\bar{f}\) 为单射因若 \(\bar{f}([x]) = \bar{f}([y])\) 则 \(f(x) = f(y)\) 即 \((x, y) \in \ker f\) 即 \([x] = [y]\)。\(\bar{f}\) 为满射因 \(f\) 是满射。故 \(\bar{f}\) 是格同构。\(\blacksquare\)

\subsection{88.3 完备格与不动点定理}\label{ux5b8cux5907ux683cux4e0eux4e0dux52a8ux70b9ux5b9aux7406}

\textbf{定义 88.3.1（完备格）} 偏序集 \((L, \leq)\) 称为完备格（complete lattice），若 \(L\) 的每个子集 \(S \subseteq L\) 都有上确界 \(\bigvee S\) 和下确界 \(\bigwedge S\)。

完备格必有最大元 \(\bigvee L = 1\) 和最小元 \(\bigwedge L = 0\)（取 \(S = \varnothing\) 和 \(S = L\)）。

\textbf{例子 88.3.2} （1）任意集合 \(X\) 的幂集 \((\mathcal{P}(X), \subseteq)\) 为完备格，\(\bigvee = \bigcup\)，\(\bigwedge = \bigcap\)。（2）任意群的所有子群构成的集合在包含关系下为完备格。（3）实闭区间 \([0,1]\) 在通常序下为完备格。（4）所有拓扑空间的开集族为完备格。

\textbf{定理 88.3.3（Knaster-Tarski不动点定理）} 设 \((L, \leq)\) 为完备格，\(f: L \to L\) 为保序映射。则 \(f\) 的不动点集合 \(\operatorname{Fix}(f) = \{x \in L : f(x) = x\}\) 非空，且在诱导序下构成一个完备格。特别地，\(f\) 有最小不动点和最大不动点。

\textbf{证明}：我们给出完整证明。

设 \(A = \{x \in L : x \leq f(x)\}\)。由于 \(0 \in A\)（\(0 \leq f(0)\)），\(A\) 非空。令 \(a = \bigvee A\)。我们证明 \(a\) 是 \(f\) 的不动点。

首先，对任意 \(x \in A\)，有 \(x \leq a\)。由 \(f\) 的保序性，\(x \leq f(x) \leq f(a)\)。因此 \(f(a)\) 是 \(A\) 的一个上界，故 \(a = \bigvee A \leq f(a)\)。此即 \(a \in A\)。

再次应用 \(f\)，由保序性得 \(f(a) \leq f(f(a))\)，故 \(f(a) \in A\)。因此 \(f(a) \leq a\)（因 \(a\) 是 \(A\) 的上确界）。结合 \(a \leq f(a)\) 和 \(f(a) \leq a\)，由反对称性得 \(f(a) = a\)。故 \(a\) 是 \(f\) 的不动点，且为最大不动点：若 \(f(x) = x\)，则 \(x \leq f(x)\) 蕴涵 \(x \in A\)，故 \(x \leq a\)。类似地，考虑集合 \(B = \{x \in L : f(x) \leq x\}\) 并令 \(b = \bigwedge B\)，可证明 \(b\) 是 \(f\) 的最小不动点。

现在证明 \(\operatorname{Fix}(f)\) 本身是完备格。设 \(S \subseteq \operatorname{Fix}(f)\)。令 \(s = \bigvee_L S\)（在 \(L\) 中取的上确界）。考虑限制 \(f|_{[s, 1]}: [s, 1] \to [s, 1]\)，其中 \([s, 1] = \{x \in L : s \leq x \leq 1\}\) 是 \(L\) 的完备子格。由前段论证，\(f|_{[s, 1]}\) 有最小不动点，此即 \(S\) 在 \(\operatorname{Fix}(f)\) 中的上确界。下确界同理。故 \(\operatorname{Fix}(f)\) 为完备格。\(\blacksquare\)

\textbf{推论 88.3.4（Knaster-Tarski对计算机科学的应用）} 设 \(L\) 为完备格，则 \(L\) 上的每个保序自映射都有不动点。此结果在指称语义（Scott域理论）中用于定义递归程序的最小不动点语义。

\textbf{定义 88.3.5（闭包算子与核算子）} 完备格 \(L\) 上的保序映射 \(c: L \to L\) 称为闭包算子（closure operator），若它满足：（1）\(x \leq c(x)\)（扩张性）；（2）\(c(c(x)) = c(x)\)（幂等性）。对偶地，\(k: L \to L\) 称为核算子（kernel operator），若 \(k(x) \leq x\) 且 \(k(k(x)) = k(x)\)。

闭包算子的不动点集恰为 \(c(L)\)，且 \(c(L)\) 是 \(L\) 的完备子格（下确界与 \(L\) 中一致，上确界为 \(c\) 下的像）。

\textbf{定理 88.3.6（Galois联络）} 设 \((P, \leq)\) 和 \((Q, \leq)\) 为偏序集。一对保序映射 \(f: P \to Q\) 和 \(g: Q \to P\) 称为Galois联络，若对任意 \(p \in P, q \in Q\)：

\[f(p) \leq q \iff p \leq g(q).\]

此时 \(gf: P \to P\) 是 \(P\) 上的闭包算子，\(fg: Q \to Q\) 是 \(Q\) 上的核算子。

\textbf{证明}：由定义，若 \(p \leq p'\) 则 \(p \leq p' \leq g(f(p'))\)，故 \(f(p) \leq f(p')\)，\(f\) 保序。类似地 \(g\) 保序。\(p \leq g(f(p))\) 因 \(f(p) \leq f(p)\) 等价于 \(p \leq g(f(p))\)。\(g(f(g(f(p)))) = g(f(p))\) 因 \(f(g(f(p))) \leq f(p)\)（从 \(g(f(p)) \leq g(f(p))\) 得 \(f(g(f(p))) \leq f(p)\)），再应用 \(g\) 即得。\(\blacksquare\)

Galois联络是数学中普遍存在的结构。经典例子包括：代数几何中理想与代数集之间的对应；Galois理论中中间域与子群之间的对应（反序Galois联络）；逻辑学中理论模型与语句集合之间的对应。

\subsection{88.4 Bool代数}\label{boolux4ee3ux6570}

\textbf{定义 88.4.1（Bool代数）} Bool代数是一个代数结构 \((B, \vee, \wedge, \neg, 0, 1)\)，其中 \(\vee\) 和 \(\wedge\) 为二元运算，\(\neg\) 为一元运算，\(0\) 和 \(1\) 为常元，满足：

\begin{enumerate}
\def\labelenumi{\arabic{enumi}.}
\tightlist
\item
  \((B, \vee, \wedge)\) 是有补分配格；
\item
  \(x \vee 0 = x\)，\(x \wedge 1 = x\)；
\item
  \(x \vee \neg x = 1\)，\(x \wedge \neg x = 0\)。
\end{enumerate}

Bool代数是经典命题逻辑的代数化表述：\(\vee\) 对应析取，\(\wedge\) 对应合取，\(\neg\) 对应否定，\(0\) 对应假，\(1\) 对应真。

\textbf{命题 88.4.2（Bool代数的基本恒等式）} 在Bool代数 \(B\) 中：

\begin{enumerate}
\def\labelenumi{\arabic{enumi}.}
\tightlist
\item
  De Morgan律：\(\neg(x \vee y) = \neg x \wedge \neg y\)，\(\neg(x \wedge y) = \neg x \vee \neg y\)；
\item
  双重否定：\(\neg\neg x = x\)；
\item
  \(x \vee x = x\)，\(x \wedge x = x\)；
\item
  \(0 \neq 1\) 的Bool代数称为非平凡的。
\end{enumerate}

\textbf{定义 88.4.3（Bool代数的同态与理想）} 保持所有Bool运算的映射为Bool同态。\(B\) 的一个子集 \(I\) 称为理想，若 \(0 \in I\)，\(x, y \in I \implies x \vee y \in I\)，且 \(x \in I, y \leq x \implies y \in I\)。理想 \(I\) 称为素的，若 \(1 \notin I\) 且 \(x \wedge y \in I\) 蕴涵 \(x \in I\) 或 \(y \in I\)。理想 \(I\) 称为极大的，若 \(I \neq B\) 且 \(I\) 不真含于任何非平凡理想中。

\textbf{定理 88.4.4（极大理想与超滤子）} Bool代数 \(B\) 中的极大理想恰为不包含 \(1\) 且对 \(\wedge\) 封闭的极大子集。对偶地，超滤子（ultrafilter）是极大滤子。每个Bool代数同态 \(f: B \to \mathbf{2} = \{0, 1\}\) 的核 \(f^{-1}(0)\) 为极大理想。

\textbf{证明} 设 \(I\) 是 \(B\) 的极大理想。则 \(1 \notin I\)（因 \(1 \in I\) 将推出 \(I = B\)）。\(I\) 对 \(\wedge\) 封闭：若 \(a, b \in I\)，则 \(a \wedge b \leq a\) 蕴涵 \(a \wedge b \in I\)。极大性：若 \(J \supsetneq I\) 是理想，则存在 \(x \in J \setminus I\)。由极大性 \(J = B\)，故 \(I\) 是极大的。

反之，设 \(I\) 是 \(\wedge\)-封闭且不含 \(1\) 的极大子集。需证 \(I\) 是理想：若 \(x \in I\) 且 \(y \leq x\)，则 \(y = y \wedge x \in I\)（\(I\) 对 \(\wedge\) 封闭）。若 \(x, y \in I\)，考虑 \(I\) 与 \(x \vee y\) 生成的上集，由极大性必包含 \(1\)，从而 \(x \vee y \in I\)。故 \(I\) 是理想。

对于同态 \(f: B \to \mathbf{2}\)，\(f^{-1}(0)\) 不含 \(1\)（\(f(1) = 1\)），且若 \(a, b \in f^{-1}(0)\) 则 \(f(a \vee b) = f(a) \vee f(b) = 0\)，故 \(a \vee b \in f^{-1}(0)\)；若 \(a \in f^{-1}(0)\) 且 \(b \leq a\)，则 \(f(b) \leq f(a) = 0\) 故 \(f(b) = 0\)。又 \(\mathbf{2}\) 中素理想 \(\{0\}\) 的逆像是极大理想。\(\blacksquare\)

\textbf{定理 88.4.5（Stone表示定理）} 每个Bool代数 \(B\) 同构于某个紧致Hausdorff全不连通拓扑空间（Stone空间）\(X\) 的既开又闭子集构成的Bool代数。等价地，\(B \cong \operatorname{Clopen}(\operatorname{Spec}(B))\)。

\textbf{完整证明}：我们给出Stone表示定理的完整构造性证明。

\textbf{第一步：构造Stone空间。} 设 \(\operatorname{Spec}(B)\) 为 \(B\) 的所有素理想构成的集合（在Bool代数中，素理想等价于极大理想）。对每个 \(a \in B\)，定义 \(U_a = \{I \in \operatorname{Spec}(B) : a \notin I\}\)。注意到以下基本性质： - \(U_0 = \varnothing\)，\(U_1 = \operatorname{Spec}(B)\)； - \(U_{a \wedge b} = U_a \cap U_b\)； - \(U_{a \vee b} = U_a \cup U_b\)； - \(U_{\neg a} = \operatorname{Spec}(B) \setminus U_a\)。

因此集合族 \(\{U_a\}_{a \in B}\) 构成了 \(\operatorname{Spec}(B)\) 上某个拓扑的基。赋予此拓扑后，\(\operatorname{Spec}(B)\) 成为拓扑空间。

\textbf{第二步：证明紧致性。} 设 \(\{U_{a_\lambda}\}_{\lambda \in \Lambda}\) 为 \(\operatorname{Spec}(B)\) 的开覆盖。假设不存在有限子覆盖。考虑由 \(\{\neg a_\lambda\}_{\lambda \in \Lambda}\) 生成的滤子 \(F\)。若 \(F\) 不包含 \(0\)，则由Zorn引理，\(F\) 可扩张为超滤子（极大滤子），其对偶极大理想 \(I\) 满足 \(a_\lambda \notin I\) 对所有 \(\lambda\) 成立，即 \(I \in \bigcap_\lambda U_{a_\lambda}\)，但 \(I \notin \bigcup_\lambda U_{a_\lambda}\)，与覆盖性矛盾。故 \(F\) 包含 \(0\)，即存在有限多个 \(\lambda_1, \ldots, \lambda_n\) 使得 \(\neg a_{\lambda_1} \wedge \cdots \wedge \neg a_{\lambda_n} = 0\)，从而 \(\neg(a_{\lambda_1} \vee \cdots \vee a_{\lambda_n}) = 0\)，即 \(a_{\lambda_1} \vee \cdots \vee a_{\lambda_n} = 1\)，故 \(U_{a_{\lambda_1}} \cup \cdots \cup U_{a_{\lambda_n}} = U_1 = \operatorname{Spec}(B)\)，得到有限子覆盖。

\textbf{第三步：证明Hausdorff性和全不连通性。} 对不同的极大理想 \(I \neq J\)，存在 \(a \in I \setminus J\)，则 \(\neg a \notin I\) 而 \(\neg a \in J\)，故 \(I \in U_{\neg a}\)，\(J \in U_a\)，且 \(U_a \cap U_{\neg a} = \varnothing\)。同时每个 \(U_a\) 也是闭集（因为 \(U_a = \operatorname{Spec}(B) \setminus U_{\neg a}\)）。因此 \(\operatorname{Spec}(B)\) 中每个点都有一组既开又闭集构成的邻域基，空间全不连通。

\textbf{第四步：建立同构。} 定义映射 \(\Phi: B \to \operatorname{Clopen}(\operatorname{Spec}(B))\) 为 \(\Phi(a) = U_a\)。由第一步知 \(\Phi\) 保持所有Bool运算。\(\Phi\) 为单射：若 \(a \neq b\)，则 \(a \wedge \neg b \neq 0\) 或 \(b \wedge \neg a \neq 0\)。不失一般性设 \(a \wedge \neg b \neq 0\)，则 \(\{a \wedge \neg b\}\) 生成不含 \(0\) 的滤子，扩张为超滤子，对偶极大理想 \(I\) 满足 \(a \notin I, b \in I\)，故 \(I \in U_a \setminus U_b\)，从而 \(U_a \neq U_b\)。\(\Phi\) 为满射：因 \(\{U_a\}\) 构成 \(\operatorname{Spec}(B)\) 的基且每个 \(U_a\) 既开又闭，由紧致性知任意既开又闭集是有限个 \(U_a\) 的并（即 \(U_{a_1 \vee \cdots \vee a_n}\)），从而是某个 \(U_a\)。综上，\(\Phi\) 为Bool代数同构。\(\blacksquare\)

Stone表示定理是代数与拓扑之间深刻联系的重要范例。其对偶语言是：Bool代数的范畴与Stone空间（紧致Hausdorff全不连通空间）的范畴等价。

\textbf{推论 88.4.6（有限Bool代数的分类）} 每个有限Bool代数同构于某个有限集合的幂集Bool代数 \(\mathcal{P}(X)\)。特别地，有限Bool代数的基数为 \(2^n\)。

\textbf{定义 88.4.7（完全Bool代数与力迫）} Bool代数 \(B\) 称为完全的，若其对应的格是完备格。完全Bool代数在公理集合论的力迫法（forcing）中起核心作用：给定完全Bool代数 \(B\)，可构造 \(B\) 值的布尔模型 \(V^B\)，从而证明连续统假设等命题的独立性。

\subsection{88.5 Heyting代数与直觉主义逻辑}\label{heytingux4ee3ux6570ux4e0eux76f4ux89c9ux4e3bux4e49ux903bux8f91}

\textbf{定义 88.5.1（Heyting代数）} Heyting代数是一个有界格 \((H, \vee, \wedge, 0, 1)\)，配备二元运算 \(\to\)（称为相对伪补或蕴涵），满足：

\[x \wedge y \leq z \iff y \leq x \to z.\]

等价地，\(\to\) 是 \(\wedge\) 的右伴随：对每个 \(x \in H\)，函子 \(x \wedge (-): H \to H\) 有右伴随 \(x \to (-): H \to H\)。

\textbf{命题 88.5.2} 每个Bool代数都是Heyting代数，其中 \(x \to y = \neg x \vee y\)。反之不真：Heyting代数未必有补运算。实际上，Heyting代数是Bool代数的推广，对应于从经典逻辑到直觉主义逻辑的过渡。

\textbf{定义 88.5.3（伪补与否定）} 在Heyting代数 \(H\) 中，元素 \(x\) 的伪补定义为 \(\neg x = x \to 0\)。与Bool代数不同，未必有 \(\neg\neg x = x\)（不满足双重否定消去），只满足 \(\neg\neg\neg x = \neg x\)。这精确刻画了直觉主义逻辑中否定运算的性质。

\textbf{命题 88.5.4} Heyting代数是分配格。

\textbf{证明}：需证 \(x \wedge (y \vee z) \leq (x \wedge y) \vee (x \wedge z)\)。由伴随关系，此等价于 \(y \vee z \leq x \to ((x \wedge y) \vee (x \wedge z))\)。但 \(y \leq x \to (x \wedge y) \leq x \to ((x \wedge y) \vee (x \wedge z))\)，\(z\) 同理，而 \(y \vee z\) 是其并，故得证。另一个分配律对偶可证。\(\blacksquare\)

\textbf{定理 88.5.5（完备Heyting代数与拓扑的关系）} 设 \(X\) 为拓扑空间，则其开集格 \(\mathcal{O}(X)\) 为完备Heyting代数，其中：

\[U \to V = \operatorname{int}((X \setminus U) \cup V) = \bigcup \{W \in \mathcal{O}(X) : U \cap W \subseteq V\}.\]

\(\mathcal{O}(X)\) 中的伪补为 \(\neg U = \operatorname{int}(X \setminus U)\)。一般地 \(U \vee \neg U \neq X\)（开集与其伪补的并不一定是整个空间），这反映了直觉主义逻辑中排中律不成立。

\textbf{证明} 需验证 \(\mathcal{O}(X)\) 是完备格（开集族的任意并和有限交为开集，故完备）且相对伪补满足伴随条件。定义 \(U \to V = \bigcup\{W \in \mathcal{O}(X) : U \cap W \subseteq V\}\)。首先验证此并为开集（因任意并保开）。对任意开集 \(W\)，\(U \cap W \subseteq V\) 当且仅当 \(W \subseteq U \to V\)：若 \(U \cap W \subseteq V\)，则 \(W\) 属于定义中的并集族，故 \(W \subseteq U \to V\)；反之若 \(W \subseteq U \to V\)，则 \(U \cap W \subseteq U \cap (U \to V) = U \cap \operatorname{int}((X \setminus U) \cup V) \subseteq (X \setminus U)^c \cap ((X \setminus U) \cup V) \subseteq V\)。故 \(U \cap W \subseteq V \iff W \subseteq U \to V\)，此即Heyting代数中 \(\to\) 的伴随条件。伪补 \(\neg U = U \to \emptyset = \operatorname{int}(X \setminus U)\)，且一般 \(U \cup \neg U \subsetneq X\)（例如 \(U\) 为 \(\mathbb{R}\) 中开区间时，边界点不在其中）。\(\blacksquare\)

\textbf{定义 88.5.6（Locale与无点拓扑）} 完备Heyting代数也称为locale或frame。Locale理论是一种''无点拓扑''：不开设空间中的点，而只研究开集格的抽象代数结构。Locale的范畴与拓扑空间的范畴之间存在伴随函子，为拓扑学提供了代数化的新视角。

\textbf{定理 88.5.7（Brouwer-Heyting-Kolmogorov（BHK）解释）} 直觉主义逻辑的证明论语义可总结为：

\begin{itemize}
\tightlist
\item
  \(\varphi \wedge \psi\) 的证明：\(\varphi\) 的证明与 \(\psi\) 的证明构成的有序对；
\item
  \(\varphi \vee \psi\) 的证明：\(\varphi\) 的证明或 \(\psi\) 的证明，并指明是哪一个；
\item
  \(\varphi \to \psi\) 的证明：将 \(\varphi\) 的任意证明转化为 \(\psi\) 的证明的构造性方法；
\item
  \(\neg \varphi\) 的证明：\(\varphi \to \bot\) 的证明，即将 \(\varphi\) 的任意证明转化为矛盾的构造性方法；
\item
  \(\exists x . \varphi(x)\) 的证明：某个个体 \(a\) 及 \(\varphi(a)\) 的证明；
\item
  \(\forall x . \varphi(x)\) 的证明：将任意个体 \(a\) 转化为 \(\varphi(a)\) 的证明的构造性方法。
\end{itemize}

BHK解释揭示了直觉主义逻辑的构造性本质。与之对应，Heyting代数为直觉主义命题逻辑提供了完全的代数语义：公式 \(\varphi\) 是直觉主义逻辑的定理当且仅当在所有Heyting代数中 \(\varphi\) 的赋值为 \(1\)。

\textbf{定义 88.5.8（正则元素与Bool化）} Heyting代数 \(H\) 中的元素 \(x\) 称为正则的，若 \(\neg\neg x = x\)。\(H\) 的所有正则元素构成的集合 \(H_{\neg\neg}\) 在运算 \(x \wedge_{\neg\neg} y = x \wedge y\)，\(x \vee_{\neg\neg} y = \neg\neg(x \vee y)\) 下构成一个Bool代数，称为 \(H\) 的Bool化（Booleanization）。这给出了从Heyting代数到Bool代数的伴随嵌入，对应证明论中的双重否定翻译。

\textbf{定理 88.5.9（完备Heyting代数的Lawvere-Tierney拓扑）} 在初等拓扑斯（elementary topos）理论中，\(\Omega\)（子对象分类器）是一个完备Heyting代数。\(\Omega\) 上的Lawvere-Tierney拓扑 \(j: \Omega \to \Omega\) 满足 \(j \circ j = j\)，\(j \circ \wedge = \wedge \circ (j \times j)\)，\(j \circ \text{true} = \text{true}\)。这推广了拓扑空间中开集的概念，是层论与范畴逻辑的基础工具。

\hrulefill

\section{第15章：域理论与连续格}\label{ux7b2c15ux7ae0ux57dfux7406ux8bbaux4e0eux8fdeux7eedux683c}

域理论（Domain Theory）起源于1960年代Scott和Strachey关于编程语言指称语义的工作，为递归函数和递归类型提供了严格的数学基础。连续格和Scott域构成了域理论的核心。

\subsection{89.1 定向完备偏序与Scott拓扑}\label{ux5b9aux5411ux5b8cux5907ux504fux5e8fux4e0escottux62d3ux6251}

\textbf{定义 89.1.1（定向子集）} 偏序集 \((P, \leq)\) 的子集 \(D \subseteq P\) 称为定向的（directed），若 \(D \neq \varnothing\) 且对任意 \(x, y \in D\)，存在 \(z \in D\) 使得 \(x \leq z\) 且 \(y \leq z\)。每个链都是定向的，但定向集不一定是链。

\textbf{定义 89.1.2（定向完备偏序，dcpo）} 偏序集 \((D, \leq)\) 称为定向完备偏序（directed-complete partial order，简称dcpo），若 \(D\) 的每个定向子集都有上确界。

dcpo是完备格的直接推广：完备格要求所有子集有上确界，而dcpo只要求定向子集有上确界。

\textbf{定义 89.1.3（Scott拓扑）} 设 \((D, \leq)\) 为dcpo。\(D\) 的子集 \(U\) 称为Scott开的，若：

\begin{enumerate}
\def\labelenumi{\arabic{enumi}.}
\tightlist
\item
  \(U\) 为上集：\(x \in U\) 且 \(x \leq y\) 蕴涵 \(y \in U\)；
\item
  对任意定向集 \(A \subseteq D\)，若 \(\bigvee A \in U\)，则 \(A \cap U \neq \varnothing\)。
\end{enumerate}

所有Scott开集构成的拓扑称为Scott拓扑，记作 \(\sigma(D)\)。

\textbf{命题 89.1.4} Scott拓扑 \(\sigma(D)\) 是 \(T_0\) 拓扑。映射 \(f: D \to E\) 关于Scott拓扑连续当且仅当 \(f\) 保序且保定向并：\(f(\bigvee A) = \bigvee f(A)\) 对所有定向集 \(A\) 成立。

\textbf{证明}：若 \(f\) 是Scott连续的，则对任意Scott开集 \(V\)，\(f^{-1}(V)\) 是Scott开集。特别地，\(f\) 保序（因 \(f^{-1}\) 保上集）且保定向并（因定向并的Scott连续性）。反之，假设 \(f\) 保序且保定向并。设 \(V \subseteq E\) 为Scott开集。\(f^{-1}(V)\) 是上集因 \(f\) 保序。若 \(\bigvee A \in f^{-1}(V)\)（\(A\) 定向），则 \(f(\bigvee A) = \bigvee f(A) \in V\)。因 \(V\) 是Scott开的，存在 \(a \in A\) 使得 \(f(a) \in V\)，即 \(a \in f^{-1}(V) \cap A\)。故 \(f^{-1}(V)\) 是Scott开集，\(f\) Scott连续。\(\blacksquare\)

\textbf{定义 89.1.5（Scott连续函数空间）} 设 \(D, E\) 为dcpo。所有从 \(D\) 到 \(E\) 的Scott连续映射构成的集合 \([D \to E]\) 在逐点序下构成dcpo：\((f \leq g) \iff (\forall x \in D, f(x) \leq g(x))\)。若 \(F \subseteq [D \to E]\) 是定向的，则其Scott意义上的上确界为 \((\bigvee F)(x) = \bigvee_{f \in F} f(x)\)。此结构对指称语义中函数空间的构造至关重要。

\subsection{89.2 连续格与Scott域}\label{ux8fdeux7eedux683cux4e0escottux57df}

\textbf{定义 89.2.1（way-below关系）} 设 \((D, \leq)\) 为dcpo，\(x, y \in D\)。若对任意定向子集 \(A \subseteq D\)，\(\bigvee A\) 存在且 \(y \leq \bigvee A\) 蕴涵存在 \(a \in A\) 使得 \(x \leq a\)，则称 \(x\) 远远小于 \(y\)（\(x\) is way-below \(y\)），记作 \(x \ll y\)。

直观上，\(x \ll y\) 意味着 \(x\) 是 \(y\) 的''有限近似''：任何逼近 \(y\) 的定向过程最终会超过 \(x\)。

\textbf{命题 89.2.2（way-below关系的基本性质）}

\begin{enumerate}
\def\labelenumi{\arabic{enumi}.}
\tightlist
\item
  \(x \ll y \implies x \leq y\)；
\item
  \(u \leq x \ll y \leq v \implies u \ll v\)；
\item
  若底部元 \(0\) 存在，则 \(0 \ll x\) 对所有 \(x\) 成立当且仅当 \(0\) 是紧元；
\item
  \(\ll\) 是传递的：\(x \ll y \ll z \implies x \ll z\)。
\end{enumerate}

\textbf{证明}：（1）取定向集 \(A = \{y\}\)，则 \(\bigvee A = y\)，由 \(x \ll y\) 得存在 \(a \in A\) 使得 \(x \leq a\)，故 \(x \leq y\)。（2）设 \(A\) 定向且 \(v \leq \bigvee A\)。则 \(y \leq \bigvee A\)，因 \(x \ll y\)，存在 \(a \in A\) 使得 \(x \leq a\)，从而 \(u \leq x \leq a\)，故 \(u \ll v\)。（3）若 \(0 \ll x\)，取含 \(x\) 的定向集，则存在某元 \(\leq x\) 且在定向集中，\(0\) 满足此条件表明 \(x\) 必须是紧的。反之亦然。（4）设 \(z \leq \bigvee A\)。因 \(y \ll z\)，存在 \(a \in A\) 使得 \(y \leq a\)。取定向集 \(B = A\)，则 \(a \leq \bigvee A\) 且 \(x \ll y \leq a\)，用 \(a\) 代替 \(z\) 同上论证得存在 \(b \in A\) 使得 \(x \leq b\)，故 \(x \ll z\)。\(\blacksquare\)

\textbf{定义 89.2.3（连续格与连续dcpo）} dcpo \(D\) 称为连续的（continuous），若对每个 \(x \in D\)，集合 \(\Downarrow x = \{y \in D : y \ll x\}\) 是定向的且 \(\bigvee \Downarrow x = x\)。若 \(D\) 同时是完备格，则称 \(D\) 为连续格（continuous lattice）。

连续格概念的精髓在于：每个元素都是其''远远小于''的元素的定向并------每个元素都被其有限近似完全确定。

\textbf{定理 89.2.4（插值性质）} 在连续dcpo中，若 \(x \ll z\)，则存在 \(y\) 使得 \(x \ll y \ll z\)。

\textbf{证明}：考虑集合 \(S = \{u \in D : \exists v \in D, u \ll v \ll z\}\)。我们证明 \(S\) 是定向的且 \(\bigvee S = z\)。首先，由连续性，\(z = \bigvee \Downarrow z\)。对任意 \(u \in \Downarrow z\)，因 \(u \ll z\)，有 \(u \in S\)。故 \(\Downarrow z \subseteq S\)，从而 \(\bigvee S = z\)。若 \(u_1, u_2 \in S\)，则存在 \(v_1, v_2\) 使得 \(u_i \ll v_i \ll z\)。因 \(\Downarrow z\) 定向，存在 \(w \in \Downarrow z\) 满足 \(v_1 \leq w, v_2 \leq w\)。由连续性，\(w = \bigvee \Downarrow w\)。因 \(u_i \ll v_i \leq w = \bigvee \Downarrow w\)，存在 \(w_i \in \Downarrow w\) 使得 \(u_i \leq w_i\)。取 \(w' \in \Downarrow w\) 大于 \(w_1, w_2\)，则 \(u_i \leq w' \ll w \ll z\)，故存在公共上界，\(S\) 定向。现在，由 \(x \ll z = \bigvee S\) 及 \(S\) 定向，存在 \(s \in S\) 使得 \(x \leq s\)。因 \(s \in S\)，存在 \(y\) 使得 \(s \ll y \ll z\)。因此 \(x \ll y \ll z\)。\(\blacksquare\)

插值性质是连续格理论中最基本的工具之一，它保证了way-below关系的''稠密性''。

\textbf{定义 89.2.5（Scott域）} dcpo \(D\) 称为Scott域（或简称域），若 \(D\) 是连续的且存在最小元 \(\bot\)。Scott域构成一个笛卡儿闭范畴，其中对象为Scott域，态射为Scott连续函数，指数对象为Scott连续函数空间。这种范畴结构是编程语言指称语义的数学模型基础。

\textbf{定义 89.2.6（紧元与代数格）} dcpo \(D\) 中的元素 \(k\) 称为紧的（compact）或有限的，若 \(k \ll k\)。记 \(K(D)\) 为 \(D\) 的所有紧元构成的集合。dcpo \(D\) 称为代数的（algebraic），若对每个 \(x \in D\)，集合 \(\{k \in K(D) : k \leq x\}\) 是定向的且以 \(x\) 为上确界。代数格是代数完备格。每个代数dcpo都是连续的，但反之不真。

\textbf{定理 89.2.7（Scott域的幂域构造）} 设 \(D\) 为Scott域。\(D\) 的所有非空Scott闭子集构成的集合在逆包含序下构成Scott域（Hoare幂域）；所有非空Scott紧饱和子集构成另一个Scott域（Smyth幂域）；所有非空凸Scott紧子集构成Plotkin幂域。

\textbf{证明概要}（Hoare幂域）记 \(\mathcal{H}(D)\) 为 \(D\) 的非空Scott闭子集族，序为逆包含 \(\supseteq\)。验证 \(\mathcal{H}(D)\) 是dcpo：设 \(\{C_i\}\) 为 \(\supseteq\) 下的定向族，则 \(\bigcap C_i\) 非空（因每个 \(C_i\) 是Scott闭的且定向族的交在Scott拓扑中非空），且 \(\bigcap C_i\) 是Scott闭的，即为定向族的上确界。连续性：对 \(C \in \mathcal{H}(D)\)，way-below关系 \(C' \ll C\) 在Hoare幂域中对应于 \((C \subseteq \operatorname{int}_\sigma(C'))\)（Scott拓扑的内部包含），利用 \(D\) 的连续性可验证 \(\mathcal{H}(D)\) 也是连续的。底元为 \(D\) 本身（逆包含序中的最大元）。

（Smyth幂域）\(\mathcal{S}(D)\) 为非空Scott紧饱和子集族，序为正常包含 \(\subseteq\)。上确界为 \(\bigcap\)（取闭包），连续性由 \(D\) 的紧致性保证有限交的非空性。

（Plotkin幂域）结合前两者的凸集，结构的连续性和完备性采用类似论证。

这三种幂域分别为非确定程序的角度语义、恶魔语义和不定语义提供了数学模型。\(\blacksquare\)

\textbf{定理 89.2.8（连续格的Lawson拓扑）} 连续格 \(L\) 上的Lawson拓扑是Scott拓扑与lower拓扑的并生成的拓扑。\(L\) 在Lawson拓扑下成为紧致Hausdorff空间，这一事实深刻地刻画了连续格的紧致性与序结构的内在联系。

\textbf{证明概要}：Lawson拓扑的紧致性可通过way-below关系和插值性质证明：任何由开集组成的覆盖可精化为由形如 \(\{x : a \ll x\}\) 的子基开集组成的覆盖，再利用完备性构造有限子覆盖。\(\blacksquare\)

\subsection{89.3 格在理论计算机科学中的应用}\label{ux683cux5728ux7406ux8bbaux8ba1ux7b97ux673aux79d1ux5b66ux4e2dux7684ux5e94ux7528}

\textbf{定义 89.3.1（指称语义）} 指称语义（denotational semantics）将程序的含义解释为Scott域之间的Scott连续函数。递归定义的程序 \(f = \Phi(f)\) 被解释为域方程的最小不动点：\(f = \bigvee_{n \geq 0} \Phi^n(\bot)\)，其中 \(\bot\) 是域的最小元，该序列的单调性由 \(\Phi\) 的Scott连续性保证。

\textbf{定理 89.3.2（递归域方程的解）} 设 \(F\) 是Scott域范畴上的函子。若 \(F\) 是局部连续的（即它在态射集上的作用是Scott连续的），则存在Scott域 \(D\) 使得 \(D \cong F(D)\)。此即递归域方程 \(D \cong F(D)\) 的解，是递归类型（如 \(T \cong T \to T\)）的数学模型。

\textbf{定义 89.3.3（抽象解释）} 抽象解释（abstract interpretation）由Cousot和Cousot于1977年提出，是一种基于格论的静态程序分析框架。其核心思想是建立具体语义域和抽象语义域之间的Galois联络：

\[(C, \leq_C) \overset{\alpha}{\underset{\gamma}{\rightleftharpoons}} (A, \leq_A)\]''

其中 \(C\) 为具体语义的完备格（通常是程序状态的幂集格），\(A\) 为抽象域的完备格，\(\alpha: C \to A\) 为抽象化映射（上伴随），\(\gamma: A \to C\) 为具体化映射（下伴随）。程序的抽象执行在 \(A\) 上进行，由于 \(A\) 通常有限或满足升链条件，抽象执行必然终止。

\textbf{定理 89.3.4（抽象解释的完备性）} 若具体语义 \(f: C \to C\) 的抽象版本 \(f^\sharp: A \to A\) 满足 \(\alpha \circ f \sqsubseteq f^\sharp \circ \alpha\)（其中 \(\sqsubseteq\) 为 \(A\) 上的偏序），则 \(f^\sharp\) 是 \(f\) 的可靠抽象：任何在抽象域中验证为真的性质在具体域中也成立。

\textbf{证明}：设 \(f^\sharp(a) \leq_A a\)（即 \(a\) 是 \(f^\sharp\) 的不动点前点），需证 \(\gamma(a)\) 是 \(f\) 的不动点前点。由条件 \(\alpha \circ f \sqsubseteq f^\sharp \circ \alpha\) 得 \(\alpha(f(\gamma(a))) \leq_A f^\sharp(\alpha(\gamma(a))) \leq_A f^\sharp(a)\)（因 \(\alpha \circ \gamma \leq \text{id}_A\)）。故 \(\alpha(f(\gamma(a))) \leq_A a\)，从而 \(f(\gamma(a)) \leq_C \gamma(a)\)（因Galois联络性质）。\(\blacksquare\)

\textbf{定义 89.3.5（程序逻辑的代数语义）} Hoare逻辑的代数化由Kozen的Kleene代数与测试（KAT）给出。程序被解释为Kleene代数中的元素，测试被解释为Bool代数中的元素。程序的顺序组合对应乘法，非确定选择对应对偶，迭代对应Kleene星号。完备性可通过代数扩张为 \(*\)-连续Kleene代数来证明。

\textbf{命题 89.3.6（并发程序的幂域语义）} 并发程序的交错语义、真并发语义和偏序语义分别对应于不同的幂域构造。Smyth幂域刻画安全性质的必定成立，Hoare幂域刻画活性性质的最终成立，Plotkin幂域刻画两者的结合。这三种幂域之间的伴随关系反映了不同精度的语义抽象。

\textbf{定义 89.3.7（格值模型与模糊逻辑）} 将Bool代数推广为完备的剩余格（residuated lattice），可得到模糊逻辑（fuzzy logic）的代数语义。剩余格是一个代数结构 \((L, \vee, \wedge, \otimes, \to, 0, 1)\)，其中 \((L, \vee, \wedge, 0, 1)\) 为有界格，\((L, \otimes, 1)\) 为交换幺半群，\(\otimes\) 和 \(\to\) 构成伴随对：\(x \otimes y \leq z \iff y \leq x \to z\)。取 \(L = [0, 1]\) 和 \(\otimes\) 为连续三角模（t-norm），即得标准的模糊逻辑语义------Lukasiewicz逻辑、Godel逻辑和乘积逻辑分别对应于三种基本三角模。

\textbf{定理 89.3.8（Scott信息系统的序理论刻画）} Scott信息系统 \((A, \operatorname{Con}, \vdash)\) 等价于以有限元为其紧元的代数Scott域。\(A\) 是信息单元（token）的集合，\(\operatorname{Con}\) 是一族有限子集（一致性谓词），\(\vdash\) 是蕴涵关系。由信息系统生成的域中的点为 \(A\) 的理想（即对 \(\vdash\) 封闭且有限一致性子集全包含于 \(\operatorname{Con}\) 的那些子集），点的逼近关系由包含序刻画。

这些格论工具共同构成了理论计算机科学的序理论基础的完整图景，从经典指称语义到现代静态分析，从函数式编程的类型论到并发系统的行为语义，格的结构无处不在。

\newpage

\chapter{01-基础/02-集合与数系/05-非标准分析.md}\label{ux57faux784002-ux96c6ux5408ux4e0eux6570ux7cfb05-ux975eux6807ux51c6ux5206ux6790.md}

\chapter{非标准分析}\label{ux975eux6807ux51c6ux5206ux6790}

\section{第16章：超幂构造与超实数}\label{ux7b2c16ux7ae0ux8d85ux5e42ux6784ux9020ux4e0eux8d85ux5b9eux6570}

非标准分析由Abraham Robinson于1960年代创立，其核心思想是为实数域 \(\mathbb{R}\) 构造一个满足相同一阶性质的真扩张域 \(^*\mathbb{R}\)，该扩张域包含无穷小和无穷大元素，从而为微积分提供严格的无穷小基础。本章建立超幂构造的基本框架，从超滤子的存在性出发，通过指标集上的等价关系构造超幂，利用Łoś定理确立转移原理，最终获得超实数域及其基本性质。超幂构造不仅适用于实数，也适用于任意一阶结构，其威力在于将无穷小方法纳入严谨的数学逻辑框架，使Newton和Leibniz的原始直觉获得坚实的形式基础。

\subsection{x.1 超滤子与超积}\label{x.1-ux8d85ux6ee4ux5b50ux4e0eux8d85ux79ef}

\textbf{定义 x.1}（滤子与超滤子）：设 \(I\) 为非空指标集。集合族 \(\mathcal{F} \subseteq \mathcal{P}(I)\) 称为 \(I\) 上的滤子，若满足：(i) \(\emptyset \notin \mathcal{F}\)，\(I \in \mathcal{F}\)；(ii) 若 \(A, B \in \mathcal{F}\)，则 \(A \cap B \in \mathcal{F}\)；(iii) 若 \(A \in \mathcal{F}\) 且 \(A \subseteq B \subseteq I\)，则 \(B \in \mathcal{F}\)。若滤子 \(\mathcal{U}\) 还满足：对任意 \(A \subseteq I\)，恰有 \(A \in \mathcal{U}\) 或 \(I \setminus A \in \mathcal{U}\)，则称 \(\mathcal{U}\) 为超滤子。不含任何有限集的超滤子称为自由超滤子。

\textbf{定理 x.2}（自由超滤子的存在性）：在任意无限集 \(I\) 上存在自由超滤子。

\textbf{证明}：考虑 \(I\) 上所有余有限子集构成的滤子基 \(\mathcal{B} = \{I \setminus F : F \subseteq I \text{ 有限}\}\)。令 \(\mathcal{F}\) 为由 \(\mathcal{B}\) 生成的滤子。由Zorn引理，\(\mathcal{F}\) 可扩张为极大滤子 \(\mathcal{U}\)。下证 \(\mathcal{U}\) 是超滤子：设 \(A \subseteq I\) 且 \(A \notin \mathcal{U}\)。若 \(I \setminus A \notin \mathcal{U}\)，则由 \(\mathcal{U}\) 的极大性，存在 \(B_1, B_2 \in \mathcal{U}\) 使得 \(A \cap B_1 = \emptyset\) 且 \((I \setminus A) \cap B_2 = \emptyset\)。则 \(B_1 \cap B_2 \subseteq I \setminus A\) 且 \(B_1 \cap B_2 \subseteq A\)，故 \(B_1 \cap B_2 = \emptyset\)，矛盾。因此 \(\mathcal{U}\) 是超滤子。由于对任意有限集 \(F\)，\(I \setminus F \in \mathcal{B} \subseteq \mathcal{U}\)，故 \(F \notin \mathcal{U}\)，\(\mathcal{U}\) 是自由超滤子。 \(\blacksquare\)

\textbf{定义 x.3}（超积）：设 \(\{M_i\}_{i \in I}\) 为一族同型一阶结构，\(\mathcal{U}\) 为 \(I\) 上的超滤子。在直积 \(\prod_{i \in I} M_i\) 上定义等价关系：\((a_i) \sim_{\mathcal{U}} (b_i)\) 当且仅当 \(\{i \in I : a_i = b_i\} \in \mathcal{U}\)。超积 \(\prod_{\mathcal{U}} M_i\) 定义为此等价关系的商集。若所有 \(M_i = M\)，则记作 \(M^I / \mathcal{U}\)，称为 \(M\) 的超幂。

\textbf{定理 x.4}（Łoś定理）：设 \(\mathcal{U}\) 为 \(I\) 上的超滤子，\(M_i\) 为一族一阶结构，\(\varphi(x_1, \ldots, x_n)\) 为一阶公式。对任意 \(a^{(1)}, \ldots, a^{(n)} \in \prod_{i \in I} M_i\)，令 \(\bar{a}^{(j)}\) 为 \(a^{(j)}\) 的等价类。则

\[\prod_{\mathcal{U}} M_i \models \varphi(\bar{a}^{(1)}, \ldots, \bar{a}^{(n)}) \iff \{i \in I : M_i \models \varphi(a^{(1)}_i, \ldots, a^{(n)}_i)\} \in \mathcal{U}.\]

\textbf{证明}：对公式 \(\varphi\) 的复杂度施归纳。若 \(\varphi\) 为原子公式 \(t_1 = t_2\)，则由超积中等式的定义直接得证。对命题联结词，利用超滤子的性质：\(\{i : M_i \models \neg \psi\} = I \setminus \{i : M_i \models \psi\}\)，而 \(\{i : M_i \models \psi \land \theta\} = \{i : M_i \models \psi\} \cap \{i : M_i \models \theta\}\)。超滤子关于补和交封闭。对存在量词 \(\exists x \psi(x, \bar{a})\)，若 \(\mathcal{U}\)-多数指标上存在见证元，则利用选择公理构造 \(b \in \prod M_i\) 使 \(\{i : M_i \models \psi(b_i, \bar{a}_i)\} \in \mathcal{U}\)，从而超积满足 \(\exists x \psi\)。反之，若超积中 \(\exists x \psi\) 成立且 \(\bar{b}\) 为见证元，则 \(\{i : M_i \models \psi(b_i, \bar{a}_i)\} \in \mathcal{U} \subseteq \{i : M_i \models \exists x \psi(x, \bar{a}_i)\}\)。全称量词类似处理。 \(\blacksquare\)

\subsection{\texorpdfstring{x.2 超实数域 \(^*\mathbb{R}\)}{x.2 超实数域 \^{}*\textbackslash mathbb\{R\}}}\label{x.2-ux8d85ux5b9eux6570ux57df-mathbbr}

固定 \(I = \mathbb{N}\) 和一个自由超滤子 \(\mathcal{U}\)。

\textbf{定义 x.5}（超实数）：超实数域 \(^*\mathbb{R}\) 定义为实数域 \(\mathbb{R}\) 的超幂 \(\mathbb{R}^{\mathbb{N}} / \mathcal{U}\)。实数 \(r \in \mathbb{R}\) 通过常序列嵌入为 \(^*r = [(r, r, r, \ldots)]_{\mathcal{U}}\)，由此 \(\mathbb{R} \hookrightarrow {}^*\mathbb{R}\) 为域嵌入。

\textbf{定理 x.6}（转移原理）：对 \(\mathbb{R}\) 上的任意一阶语句 \(\varphi\)，\(\mathbb{R} \models \varphi\) 当且仅当 \(^*\mathbb{R} \models \varphi\)。

\textbf{证明}：这是Łoś定理的直接推论。\(\mathbb{R} \models \varphi\) 等价于每个指标 \(i\) 处 \(\mathbb{R} \models \varphi\)，故 \(\{i \in \mathbb{N} : \mathbb{R} \models \varphi\} = \mathbb{N} \in \mathcal{U}\)。由Łoś定理，\(^*\mathbb{R} \models \varphi\)。反之亦然。 \(\blacksquare\)

\textbf{定理 x.7}（\(^*\mathbb{R}\) 是有序域）：\(^*\mathbb{R}\) 关于逐点定义的加法和乘法构成有序域，且 \(\mathbb{R}\) 为其有序子域。

\textbf{证明}：有序域的公理皆为一阶语句。因 \(\mathbb{R}\) 是有序域，由转移原理，\(^*\mathbb{R}\) 也是有序域。嵌入映射保序的原因是：\(r < s\) 在 \(\mathbb{R}\) 中成立当且仅当 \(^*r < {}^*s\) 在 \(^*\mathbb{R}\) 中成立。这由常序列表示直接可得。 \(\blacksquare\)

\textbf{定义 x.8}（无穷小与无穷大）：元素 \(x \in {}^*\mathbb{R}\) 称为无穷小，若对任意正实数 \(r \in \mathbb{R}^+\)，有 \(|x| < r\)。元素 \(y \in {}^*\mathbb{R}\) 称为无穷大，若对任意正实数 \(r \in \mathbb{R}^+\)，有 \(|y| > r\)。\(^*\mathbb{R}\) 中非无穷大亦非无穷小的元素称为有限元。记 \(^*\mathbb{R}_{\text{fin}}\) 为所有有限超实数的集合，\(\mu(0)\) 为所有无穷小的集合。

\textbf{定理 x.9}（无穷小与无穷大的存在性）：\(^*\mathbb{R}\) 中确实存在非零无穷小和无穷大元素。

\textbf{证明}：取序列 \(x_n = 1/n\)，\(y_n = n\)。对任意正实数 \(r\)，仅有限个 \(n\) 满足 \(1/n \geq r\)，故 \(\{n : 1/n < r\}\) 为余有限集，属于 \(\mathcal{U}\)。因此 \([(x_n)]_{\mathcal{U}}\) 为无穷小。同理，\([(y_n)]_{\mathcal{U}}\) 为无穷大。 \(\blacksquare\)

\textbf{定理 x.10}（\(^*\mathbb{R}_{\text{fin}}\) 是局部环）：所有有限超实数构成 \(\mathbb{R}\) 上的一个局部环，其唯一极大理想恰为 \(\mu(0)\)，即无穷小集。商环 \(^*\mathbb{R}_{\text{fin}} / \mu(0) \cong \mathbb{R}\)。

\textbf{证明}：首先验证 \(^*\mathbb{R}_{\text{fin}}\) 是子环：若 \(x, y\) 有限，则存在 \(M, N \in \mathbb{R}^+\) 使 \(|x| < M\)，\(|y| < N\)，故 \(|x+y| \leq |x|+|y| < M+N\) 且 \(|xy| < MN\)，均为有限。\(\mu(0)\) 是理想：若 \(\varepsilon, \delta \in \mu(0)\)，则 \(\varepsilon+\delta\) 无穷小；对任意 \(a \in {}^*\mathbb{R}_{\text{fin}}\)，\(a\varepsilon\) 仍无穷小。现证极大性：若 \(x \in {}^*\mathbb{R}_{\text{fin}} \setminus \mu(0)\)，则存在 \(r \in \mathbb{R}^+\) 使 \(|x| \geq r\)，故 \(1/x\) 存在且 \(|1/x| \leq 1/r\)，即 \(x\) 可逆。因此任意包含 \(\mu(0)\) 的真理想必等于 \(\mu(0)\)。同构由标准部分映射给出（见下章）。 \(\blacksquare\)

\subsection{x.3 饱和性原理}\label{x.3-ux9971ux548cux6027ux539fux7406}

\textbf{定义 x.11}（\(\kappa\)-饱和性）：设 \(\kappa\) 为无穷基数。一阶结构 \(\mathcal{M}\) 称为 \(\kappa\)-饱和的，若对任意基数小于 \(\kappa\) 的参数集 \(A \subseteq |\mathcal{M}|\)，\(\mathcal{M}\) 上所有与 \(\operatorname{Th}(\mathcal{M}, A)\) 一致的型均可在 \(\mathcal{M}\) 中实现。

\textbf{定理 x.12}（超幂的 \(\aleph_1\)-饱和性）：若 \(\mathcal{U}\) 是 \(\mathbb{N}\) 上的自由超滤子，则超幂 \(^*\mathbb{R} = \mathbb{R}^{\mathbb{N}} / \mathcal{U}\) 是 \(\aleph_1\)-饱和的。

\textbf{证明}：设 \(\Sigma(x) = \{\varphi_n(x, \bar{a}^{(n)}) : n \in \mathbb{N}\}\) 为与 \(^*\mathbb{R}\) 理论一致的可数型，其中每个 \(\varphi_n\) 含参数 \(\bar{a}^{(n)}\)。对每个 \(k \in \mathbb{N}\)，令 \(S_k = \{i \in \mathbb{N} : \mathbb{R} \models \exists x \bigwedge_{n \leq k} \varphi_n(x, \bar{a}_i^{(n)})\}\)。由一致性与Łoś定理，\(S_k \in \mathcal{U}\)。选取递减序列 \(S_0 \supseteq S_1 \supseteq S_2 \supseteq \cdots\) 满足 \(S_n \in \mathcal{U}\) 且 \(\bigcap_n S_n = \emptyset\)（可由 \(\mathcal{U}\) 自由保证）。对 \(i \in S_k \setminus S_{k+1}\)，取 \(x_i\) 满足 \(\bigwedge_{n \leq k} \varphi_n(x_i, \bar{a}_i^{(n)})\)。则 \(x = [(x_i)]_{\mathcal{U}}\) 实现 \(\Sigma(x)\)，因为对每个 \(n\)，\(\{i : \mathbb{R} \models \varphi_n(x_i, \bar{a}_i^{(n)})\} \supseteq S_n \in \mathcal{U}\)。 \(\blacksquare\)

\textbf{定理 x.13}（饱和性推论------一致收敛与紧致性）：设 \((f_n)\) 为 \(^*\mathbb{R}\) 上的一列内部函数且逐点收敛。若该序列的内部基数不超过 \(^*\mathbb{R}\) 的饱和基数，则存在无穷大自然数 \(N\) 使对所有有限 \(n\)，\(f_n\) 与 \(f_N\) 无穷接近。

\textbf{证明}：由 \(\aleph_1\)-饱和性，型 \(\{x \in {}^*\mathbb{N} : x > n \text{ （对每个标准自然数 } n\text{）}\} \cup \{|f_x(k) - f_\infty(k)| < 1/x \text{ （对每个标准 } k\text{）}\}\) 可被实现。 \(\blacksquare\)

\hrulefill

\section{第17章：非标准微积分}\label{ux7b2c17ux7ae0ux975eux6807ux51c6ux5faeux79efux5206}

非标准微积分将Newton和Leibniz的无穷小直觉严格化，提供了一种处理极限、连续性和积分的全新范式。本章在超实数域 \(^*\mathbb{R}\) 的基础上，系统发展标准部分映射、连续性的非标准刻画、微分的代数化定义以及Riemann积分的无穷和表示。非标准方法的核心优势在于将 \(\varepsilon\)-\(\delta\) 语言中复杂的量词交替替换为更为直观的无穷小比较，使得许多经典定理的证明大为简化，同时保留了完全的数学严格性。这种范式不仅具有教学价值，在泛函分析和概率论等前沿领域也展现出强大的技术威力。

\subsection{x+1.1 标准部分映射}\label{x1.1-ux6807ux51c6ux90e8ux5206ux6620ux5c04}

\textbf{定义 x+1.1}（标准部分）：设 \(x \in {}^*\mathbb{R}_{\text{fin}}\)。由定理x.10，存在唯一的实数 \(r \in \mathbb{R}\) 使得 \(x - r\) 为无穷小。此实数 \(r\) 称为 \(x\) 的标准部分，记作 \(\operatorname{st}(x)\) 或 \({}^{\circ}x\)。

\textbf{定理 x+1.2}（标准部分的基本性质）：映射 \(\operatorname{st} : {}^*\mathbb{R}_{\text{fin}} \to \mathbb{R}\) 是环同态，即对任意有限 \(x, y\)： (i) \(\operatorname{st}(x \pm y) = \operatorname{st}(x) \pm \operatorname{st}(y)\)； (ii) \(\operatorname{st}(xy) = \operatorname{st}(x)\operatorname{st}(y)\)； (iii) 若 \(x \leq y\)，则 \(\operatorname{st}(x) \leq \operatorname{st}(y)\)； (iv) 若 \(\operatorname{st}(y) \neq 0\)，则 \(\operatorname{st}(x/y) = \operatorname{st}(x)/\operatorname{st}(y)\)。

\textbf{证明}：(i) 设 \(\operatorname{st}(x) = a\)，\(\operatorname{st}(y) = b\)，则 \(x = a + \varepsilon\)，\(y = b + \delta\)，其中 \(\varepsilon, \delta \in \mu(0)\)。故 \(x + y = (a+b) + (\varepsilon+\delta)\)，而 \(\varepsilon+\delta \in \mu(0)\)，故 \(\operatorname{st}(x+y) = a+b\)。(ii) \(xy = (a+\varepsilon)(b+\delta) = ab + (a\delta + b\varepsilon + \varepsilon\delta)\)。因 \(a, b\) 为标准实数，\(\varepsilon, \delta\) 无穷小，\(a\delta + b\varepsilon + \varepsilon\delta\) 为无穷小，故 \(\operatorname{st}(xy) = ab\)。(iii) 设 \(a > b\)。则 \(x - y = (a-b) + (\varepsilon-\delta)\)。因 \(a-b > 0\) 且为标准实数，存在标准正实数 \(r < a-b\) 使 \(x - y \geq r - |\varepsilon-\delta|\)。因 \(|\varepsilon-\delta| < r\) 对所有标准 \(r\) 成立，取 \(r = (a-b)/2\) 得 \(x-y > 0\)，与 \(x \leq y\) 矛盾。(iv) 由 (ii) 知 \(1 = \operatorname{st}(y \cdot 1/y) = \operatorname{st}(y) \operatorname{st}(1/y)\)，故 \(\operatorname{st}(1/y) = 1/\operatorname{st}(y)\)，结合 (ii) 即得。 \(\blacksquare\)

\subsection{x+1.2 连续性的非标准刻画}\label{x1.2-ux8fdeux7eedux6027ux7684ux975eux6807ux51c6ux523bux753b}

\textbf{定义 x+1.3}（S-连续性）：函数 \(f : \mathbb{R} \to \mathbb{R}\) 在点 \(a\) 处称为 S-连续，若对任意 \(x \in {}^*\mathbb{R}\)，当 \(x \approx a\)（即 \(x-a\) 为无穷小）时，有 \(^*f(x) \approx f(a)\)。

\textbf{定理 x+1.4}（连续性等价定理）：函数 \(f : \mathbb{R} \to \mathbb{R}\) 在 \(a\) 处连续（按 \(\varepsilon\)-\(\delta\) 定义）当且仅当 \(f\) 在 \(a\) 处 S-连续。

\textbf{证明}：（\(\Rightarrow\)）设 \(f\) 在 \(a\) 处 \(\varepsilon\)-\(\delta\) 连续。对任意标准 \(\varepsilon > 0\)，存在标准 \(\delta > 0\) 使 \(|x-a| < \delta \implies |f(x)-f(a)| < \varepsilon\)。由转移原理，此陈述在 \(^*\mathbb{R}\) 中也成立。若 \(x \approx a\)，则 \(|x-a| < \delta\) 对任意标准 \(\delta\) 成立，故 \(|^*f(x) - f(a)| < \varepsilon\) 对任意标准 \(\varepsilon\) 成立，因此 \(^*f(x) \approx f(a)\)。

（\(\Leftarrow\)）反设 \(f\) 在 \(a\) 处不连续，则存在标准 \(\varepsilon > 0\) 使对任意标准 \(n\)，存在 \(x_n\) 满足 \(|x_n - a| < 1/n\) 但 \(|f(x_n) - f(a)| \geq \varepsilon\)。由 \(\aleph_1\)-饱和性，存在 \(x \in {}^*\mathbb{R}\) 使对所有标准 \(n\)，\(|x-a| < 1/n\) 且 \(|^*f(x) - f(a)| \geq \varepsilon\)。此时 \(x \approx a\) 但 \(^*f(x) \not\approx f(a)\)，矛盾。 \(\blacksquare\)

\textbf{定理 x+1.5}（紧致集上的一致连续性）：设 \(f : [a, b] \to \mathbb{R}\) 连续。则 \(f\) 在 \([a, b]\) 上一致连续，当且仅当对任意 \(x, y \in {}^*[a, b]\)，\(x \approx y\) 蕴含 \(^*f(x) \approx {}^*f(y)\)。

\textbf{证明}：（\(\Rightarrow\)）一致连续给出对任意标准 \(\varepsilon > 0\)，存在标准 \(\delta > 0\) 使 \(|x-y| < \delta \implies |f(x)-f(y)| < \varepsilon\)。若 \(x \approx y\)，则 \(|x-y| < \delta\) 对任意标准 \(\delta\) 成立，故结论成立。

（\(\Leftarrow\)）若非一致连续，则存在 \(\varepsilon > 0\) 及序列 \(x_n, y_n \in [a, b]\) 使 \(|x_n - y_n| < 1/n\) 但 \(|f(x_n) - f(y_n)| \geq \varepsilon\)。由饱和性，存在 \(x, y \in {}^*[a, b]\) 使对所有 \(n\)，\(|x-y| < 1/n\) 且 \(|^*f(x) - {}^*f(y)| \geq \varepsilon\)，即 \(x \approx y\) 但 \(^*f(x) \not\approx {}^*f(y)\)。 \(\blacksquare\)

\subsection{x+1.3 微分的非标准定义}\label{x1.3-ux5faeux5206ux7684ux975eux6807ux51c6ux5b9aux4e49}

\textbf{定义 x+1.6}（可微性）：函数 \(f : \mathbb{R} \to \mathbb{R}\) 在 \(a\) 处可微，若存在实数 \(L\) 使对任意非零无穷小 \(\Delta x\)，有

\[\operatorname{st}\left(\frac{^*f(a+\Delta x) - f(a)}{\Delta x}\right) = L.\]

此时记 \(f'(a) = L\)。

\textbf{定理 x+1.7}（非标准微分的等价性）：\(f\) 在 \(a\) 处以非标准意义可微当且仅当 \(f\) 在 \(a\) 处以 \(\varepsilon\)-\(\delta\) 定义可微，且两种定义给出的导数值相同。

\textbf{证明}：（\(\Rightarrow\)）设 \(L = f'(a)\) 为非标准导数。对任意标准 \(\varepsilon > 0\)，由 \(\operatorname{st}\) 的定义，存在 \(\delta\)（可取为任意无穷小）使对任意 \(0 < |\Delta x| < \delta\) 且 \(\Delta x\) 为标准实数时，

\[\left|\frac{f(a+\Delta x) - f(a)}{\Delta x} - L\right| < \varepsilon.\]

这给出了标准 \(\varepsilon\)-\(\delta\) 定义（可选取 \(\delta\) 为适当的正标准数）。

（\(\Leftarrow\)）设 \(f\) 依标准意义可微。对任意非零无穷小 \(\Delta x\)，考虑差商 \(Q = \frac{^*f(a+\Delta x) - f(a)}{\Delta x}\)。对任意标准 \(\varepsilon > 0\)，存在标准 \(\delta > 0\)（由可微性）使对所有标准 \(h\)，\(0 < |h| < \delta \implies |\frac{f(a+h)-f(a)}{h} - L| < \varepsilon\)。由转移原理，此陈述在 \(^*\mathbb{R}\) 中成立。因 \(\Delta x\) 为无穷小，必有 \(|\Delta x| < \delta\)，故 \(|Q - L| < \varepsilon\)。由 \(\varepsilon\) 的任意性，\(\operatorname{st}(Q) = L\)。 \(\blacksquare\)

\textbf{定理 x+1.8}（链式法则的非标准证明）：设 \(f\) 在 \(a\) 处可微，\(g\) 在 \(f(a)\) 处可微，则 \(h = g \circ f\) 在 \(a\) 处可微且 \(h'(a) = g'(f(a)) \cdot f'(a)\)。

\textbf{证明}：取非零无穷小 \(\Delta x\)。令 \(\Delta y = {}^*f(a+\Delta x) - f(a)\)。由 \(f\) 的可微性，\(\Delta y\) 为无穷小，且 \(\operatorname{st}(\Delta y / \Delta x) = f'(a)\)。若 \(\Delta y \neq 0\)，则

\[\frac{^*h(a+\Delta x) - h(a)}{\Delta x} = \frac{^*g(f(a)+\Delta y) - g(f(a))}{\Delta y} \cdot \frac{\Delta y}{\Delta x}.\]

取标准部分即得 \(g'(f(a)) \cdot f'(a)\)。若 \(\Delta y = 0\)，则 \(f'(a) = \operatorname{st}(\Delta y / \Delta x) = 0\)。此时需证 \(h'(a) = 0\)。对非零无穷小 \(\Delta x\)，若 \(\Delta y = 0\)，则直接有 \(^*h(a+\Delta x) = g(f(a))\)，差商为零。若对某序列 \(\Delta y \neq 0\) 但极限为零，则由上述论证，差商趋近 \(g'(f(a)) \cdot f'(a) = 0\)。综合得证。 \(\blacksquare\)

\textbf{定理 x+1.8a}（非标准极值定理）：设 \(f : [a, b] \to \mathbb{R}\) 连续。则 \(f\) 在 \([a, b]\) 上达到最大值和最小值。

\textbf{证明}：取无穷大自然数 \(N\) 并将 \([a, b]\) 等分为 \(N\) 个子区间，分点 \(x_k = a + k(b-a)/N\)（\(k = 0, \ldots, N\)）。在内部有限集中选取 \(x_{k_0}\) 使 \(^*f(x_{k_0}) = \max_{k} {}^*f(x_k)\)（内部有限集上的最大值总存在）。令 \(c = \operatorname{st}(x_{k_0})\)。由连续性，\(^*f(x_{k_0}) \approx f(c)\)。对任意标准 \(x \in [a, b]\)，存在 \(k\) 使 \(x \in [x_{k-1}, x_k]\)，由连续性 \(f(x) \approx {}^*f(x_k) \leq {}^*f(x_{k_0}) \approx f(c)\)。取标准部分得 \(f(x) \leq f(c)\)。因此 \(c\) 即为最大值点。最小值同理。 \(\blacksquare\)

\textbf{定理 x+1.8b}（非标准Darboux定理 / 介值定理）：设 \(f : [a, b] \to \mathbb{R}\) 连续且 \(f(a) < 0 < f(b)\)。则存在 \(c \in (a, b)\) 使 \(f(c) = 0\)。

\textbf{证明}：取无穷大 \(N\)，令 \(x_k = a + k(b-a)/N\)。定义 \(k_0 = \min\{k \in {}^*\mathbb{N} : {}^*f(x_k) \geq 0\}\)（在内部有限集中良定义，因 \(^*f(x_N) = f(b) > 0\)）。令 \(c = \operatorname{st}(x_{k_0})\)。则 \(^*f(x_{k_0-1}) < 0\) 且 \(^*f(x_{k_0}) \geq 0\)。由于 \(x_{k_0-1} \approx x_{k_0} \approx c\)，由连续性，\(f(c) \approx {}^*f(x_{k_0-1}) < 0\) 且 \(f(c) \approx {}^*f(x_{k_0}) \geq 0\)，两者同时成立意味着 \(f(c) = 0\)。 \(\blacksquare\)

\subsection{x+1.4 Riemann积分的非标准定义}\label{x1.4-riemannux79efux5206ux7684ux975eux6807ux51c6ux5b9aux4e49}

\textbf{定义 x+1.9}（非标准Riemann积分）：设 \(f : [a, b] \to \mathbb{R}\) 为有界函数。取无穷大自然数 \(N \in {}^*\mathbb{N} \setminus \mathbb{N}\)，令 \(\Delta x = (b-a)/N\)，\(x_k = a + k\Delta x\)（\(k = 0, 1, \ldots, N\)）。定义非标准Riemann和

\[S_N(f) = \sum_{k=1}^{N} {}^*f(\xi_k) \Delta x,\]

其中 \(\xi_k \in {}^*[x_{k-1}, x_k]\) 任意选取。若对任意选取的 \(\xi_k\) 和任意无穷大 \(N\)，\(\operatorname{st}(S_N(f))\) 存在且相等，则称 \(f\) 在 \([a, b]\) 上Riemann可积，其积分值为该共同标准部分。

\textbf{定理 x+1.10}（非标准Riemann积分与标准定义的等价性）：\(f\) 在 \([a, b]\) 上依非标准定义Riemann可积当且仅当 \(f\) 依Darboux定义Riemann可积，且积分值相等。

\textbf{证明}：（\(\Rightarrow\)）设非标准Riemann和均收敛至同一标准部分 \(I\)。对任意标准 \(\varepsilon > 0\)，存在标准 \(\delta > 0\) 使对任意分割 \(\{x_k\}\) 满足 \(\max \Delta x_k < \delta\) 及任意选取 \(\xi_k\)，有 \(|\sum f(\xi_k) \Delta x_k - I| < \varepsilon\)。这可由非标准条件通过反证法得到：若不然，存在 \(\varepsilon_0\) 及一列越来越细的分割使Riemann和不收敛至 \(I\)，由饱和性构造无穷细分割得矛盾。进而，\(I\) 即为Riemann积分值。

（\(\Leftarrow\)）设 \(f\) 依Darboux标准Riemann可积，积分值为 \(I\)。对任意标准 \(\varepsilon > 0\)，存在标准 \(\delta > 0\) 使对任意满足 \(\max \Delta x_k < \delta\) 的分割，任意Riemann和与 \(I\) 之差的绝对值小于 \(\varepsilon\)。对任意无穷大 \(N\)，分割 \(\Delta x = (b-a)/N\) 的模长 \(\Delta x\) 为无穷小，故小于 \(\delta\)。由转移原理，非标准Riemann和的绝对值小于 \(\varepsilon\)。由 \(\varepsilon\) 的任意性，\(\operatorname{st}(S_N(f)) = I\)。 \(\blacksquare\)

\textbf{定理 x+1.11}（微积分基本定理的非标准形式）：设 \(F\) 在 \([a, b]\) 上可微且 \(F' = f\) 连续。则对任意无穷大自然数 \(N\)，

\[\operatorname{st}\left(\sum_{k=1}^{N} {}^*f(\xi_k) \Delta x\right) = F(b) - F(a).\]

\textbf{证明}：取 \(\xi_k \in {}^*[x_{k-1}, x_k]\)。对每个 \(k\)，由非标准中值定理，存在 \(c_k \in (x_{k-1}, x_k)\) 使 \(^*F(x_k) - {}^*F(x_{k-1}) = {}^*f(c_k) \Delta x\)。则

\[\sum_{k=1}^{N} {}^*f(\xi_k) \Delta x = \sum_{k=1}^{N} \left[{}^*F(x_k) - {}^*F(x_{k-1})\right] + \sum_{k=1}^{N} [{}^*f(\xi_k) - {}^*f(c_k)] \Delta x.\]

第一项望远镜求和等于 \(F(b) - F(a)\)。由于 \(f\) 在紧致集上一致连续，第二项中每个差分为无穷小，求和后标准部分为零。取标准部分即得结论。 \(\blacksquare\)

\hrulefill

\section{第18章：非标准泛函分析与应用}\label{ux7b2c18ux7ae0ux975eux6807ux51c6ux6cdbux51fdux5206ux6790ux4e0eux5e94ux7528}

非标准分析在泛函分析和概率论中展现出独特的技术优势。超幂构造为Hilbert空间提供了非标准扩张，使紧算子理论获得新的解释；Loeb测度构造将有限测度空间转化为标准的可数可加测度空间，开辟了测度论的新途径；Bernstein-Robinson不变子空间定理是非标准方法在算子理论中的经典成果，解决了多项式紧算子的不变子空间问题。此外，Nelson的激进测度论为概率论提供了全新的公理化基础，而非标准方法在经济学和数理金融中也日益显示其应用价值，特别是在连续时间金融模型中处理无穷小时间步长时具有天然优势。

\subsection{x+2.1 非标准Hilbert空间}\label{x2.1-ux975eux6807ux51c6hilbertux7a7aux95f4}

\textbf{定义 x+2.1}（非标准Hilbert空间）：设 \(H\) 为可分Hilbert空间。其非标准扩张 \(^*H\) 定义为 \(H\) 的超幂。\(^*H\) 是 \(^*\mathbb{C}\) 上的内积空间，内积由转移原理定义：\(\langle x, y \rangle_{^*H} = {}^*\langle \cdot, \cdot \rangle_H(x, y)\)。元素 \(x \in {}^*H\) 称为近标准的，若存在 \(y \in H\) 使得 \(x \approx y\)（即 \(\|x - y\| \approx 0\)）。

\textbf{定理 x+2.2}（非标准Riesz表示定理）：设 \(\varphi : H \to \mathbb{R}\) 为有界线性泛函。则存在唯一的 \(y \in H\) 使 \(\varphi(x) = \langle x, y \rangle\) 对任意 \(x \in H\) 成立。在非标准视角下，这可表述为：\(^*\varphi(x) \approx \langle x, y \rangle_{^*H}\) 对任意近标准 \(x\) 成立。

\textbf{证明}：标准部分的证明即为经典的Riesz表示定理。非标准视角下，若 \(x \in {}^*H\) 为近标准且 \(x = y_0 + \varepsilon\) 其中 \(y_0 \in H\)，\(\|\varepsilon\| \approx 0\)，则由 \(^*\varphi\) 的有界性，\(|^*\varphi(\varepsilon)| \leq \|\varphi\| \cdot \|\varepsilon\| \approx 0\)，故 \(^*\varphi(x) \approx {}^*\varphi(y_0) = \langle y_0, y \rangle \approx \langle x, y \rangle_{^*H}\)。 \(\blacksquare\)

\subsection{x+2.2 Loeb测度}\label{x2.2-loebux6d4bux5ea6}

\textbf{定义 x+2.3}（Loeb测度）：设 \((\Omega, \mathcal{A}, \mu)\) 为内部有限可加测度空间（在非标准全域中）。定义标准的外测度：对任意 \(A \subseteq \Omega\)，

\[\bar{\mu}(A) = \inf\{\operatorname{st}(\mu(B)) : B \in \mathcal{A}, A \subseteq B\}.\]

Loeb \(\sigma\)-代数 \(\mathcal{L}(\mathcal{A})\) 由满足 \(\bar{\mu}(A) + \bar{\mu}(\Omega \setminus A) = 1\) 的集合 \(A\) 构成。限制 \(\mu_L = \bar{\mu}|_{\mathcal{L}(\mathcal{A})}\) 称为Loeb测度。

\textbf{定理 x+2.4}（Loeb测度是标准的可数可加测度）：\(\mu_L\) 是 \(\mathcal{L}(\mathcal{A})\) 上的完备、可数可加概率测度。

\textbf{证明}：首先验证 \(\bar{\mu}\) 是外测度。单调性和可数次可加性由 \(\inf\) 的性质得出。关键在于验证 \(\mathcal{L}(\mathcal{A})\) 上的可数可加性。设 \(\{A_n\}\) 为 \(\mathcal{L}(\mathcal{A})\) 中不交的集合列。对任意 \(\varepsilon > 0\)，由 \(\mathcal{A}\) 的饱和性及 \(\aleph_1\)-饱和性，可选取内部集合 \(B_n \in \mathcal{A}\) 使 \(A_n \subseteq B_n\) 且 \(\operatorname{st}(\mu(B_n)) \leq \mu_L(A_n) + \varepsilon/2^n\)。由内部测度的有限可加性并利用饱和性协调不交性，可证

\[\mu_L\left(\bigcup_{n=1}^{\infty} A_n\right) = \sum_{n=1}^{\infty} \mu_L(A_n).\]

完备性由构造直接可得。 \(\blacksquare\)

\subsection{x+2.3 Bernstein-Robinson不变子空间定理}\label{x2.3-bernstein-robinsonux4e0dux53d8ux5b50ux7a7aux95f4ux5b9aux7406}

\textbf{定义 x+2.5}（不变子空间问题）：设 \(T\) 为复可分Hilbert空间 \(H\) 上的有界线性算子。\(T\) 的不变子空间是指闭子空间 \(M \subseteq H\) 使得 \(M \neq \{0\}, H\) 且 \(T(M) \subseteq M\)。不变子空间问题问：是否每个有界线性算子都有非平凡的不变子空间？

\textbf{定理 x+2.6}（Bernstein-Robinson, 1966）：设 \(T\) 为复可分Hilbert空间 \(H\) 上的有界线性算子。若存在非零多项式 \(p(z)\) 使 \(p(T)\) 为紧算子，则 \(T\) 有非平凡的不变子空间。

\textbf{证明}：假定 \(p(T)\) 为非零紧算子。不妨设 \(p\) 为非零且 \(\deg p \geq 1\)。设 \(K = p(T)\) 为紧算子。若非零，考虑 \(^*H\) 中的非标准分析。取 \(H\) 中任意非零向量 \(x_0\)。考虑闭包

\[M = \overline{\operatorname{span}}\{T^n K x_0 : n = 0, 1, 2, \ldots\}.\]

若 \(M \neq H\)，则 \(M\) 即为不变子空间。故设 \(M = H\)。

在非标准全域中，由 \(K\) 的紧性，\(K\) 将 \(^*H\) 中的有限范数元素映为近标准元素（即与 \(H\) 中某元素无穷接近）。利用 \(\aleph_1\)-饱和性，存在内部维数有限的子空间 \(F \subseteq {}^*H\) 使 \(H \subseteq F\) 且 \(T(F) \subseteq F\)。

现在构造非零向量 \(y \in F^{\perp} \cap H\)：由Gram-Schmidt正交化过程（在内部意义下）可对 \(F\) 的正交补中的向量施加控制。利用 \(p(T)\) 的紧性，可论证存在非零向量 \(y \in H\) 几乎正交于 \(F\)。更精确地，构造有界序列 \(y_k \in H\) 趋近于 \(F\) 的正交补，由 \(H\) 中单位球的弱紧性，存在弱收敛子列，其极限 \(y\) 满足对任意 \(x \in H\)，\(\langle T^n y, x \rangle = 0\)（对适当的 \(n\)）。

最终，令 \(N = \overline{\operatorname{span}}\{T^n y : n \geq 0\}\)。则 \(N\) 为非平凡的 \(T\)-不变子空间。 \(\blacksquare\)

\subsection{x+2.4 Nelson激进测度论}\label{x2.4-nelsonux6fc0ux8fdbux6d4bux5ea6ux8bba}

\textbf{定义 x+2.7}（Nelson内部测度论）：在非标准全域中，Nelson（1987）提出了激进测度论（Radically Elementary Probability Theory），其中概率空间由有限样本空间 \(\Omega\)（\(|\Omega| \in {}^*\mathbb{N}\) 为无穷大）和均等权重的计数测度构成。所有事件为 \(\Omega\) 的任意子集。

\textbf{定理 x+2.8}（Nelson的Radon-Nikodym定理）：设 \(P\) 为 \(\Omega\) 上的激进概率测度（均等权重 \(1/|\Omega|\)），\(Q\) 为 \(\Omega\) 上的另一激进概率测度，且 \(Q\) 关于 \(P\) 绝对连续（即 \(P(A) = 0 \implies Q(A) = 0\) 对内部 \(A\) 成立）。则存在内部随机变量 \(Z : \Omega \to {}^*\mathbb{R}^+\) 使 \(Q(A) = \sum_{\omega \in A} Z(\omega) P(\{\omega\})\)。

\textbf{证明}：定义 \(Z(\omega) = \frac{Q(\{\omega\})}{P(\{\omega\})}\)。因 \(P(\{\omega\}) = 1/|\Omega| > 0\) 且 \(Q(\{\omega\}) \geq 0\)，\(Z\) 为良定义的非负内部函数。对任意内部集合 \(A \subseteq \Omega\)，

\[\sum_{\omega \in A} Z(\omega) P(\{\omega\}) = \sum_{\omega \in A} Q(\{\omega\}) = Q(A).\]

这即完成了证明。在激进测度论中，所有测度论的复杂性被转移到超有限和与标准部分映射上。 \(\blacksquare\)

\subsection{x+2.5 非标准方法在经济学和数理金融中的应用}\label{x2.5-ux975eux6807ux51c6ux65b9ux6cd5ux5728ux7ecfux6d4eux5b66ux548cux6570ux7406ux91d1ux878dux4e2dux7684ux5e94ux7528}

\textbf{定理 x+2.9}（非标准Arrow-Debreu均衡）：在具有超有限多个代理人和商品的交换经济中，若偏好满足标准的正则性条件（由转移原理适用于内部代理人的一阶性质），则存在Walrasian均衡价格向量 \(p \in {}^*\mathbb{R}^L_+\) 使市场出清。

\textbf{证明}：对标准有限经济，Arrow-Debreu存在性定理是一阶表述的。由转移原理，该定理适用于任意内部有限的经济体。具体地，考虑内部有限集 \(I\)（\(|I| \in {}^*\mathbb{N}\)），每个代理人 \(i \in I\) 有初始禀赋 \(e_i \in {}^*\mathbb{R}^L_+\) 和满足单调性、凸性、连续性的偏好。由转移原理运用标准的不动点论证，存在均衡价格。取标准部分（若价格向量的标准部分非零），得到标准经济体的近似均衡。 \(\blacksquare\)

\textbf{定理 x+2.10}（非标准Black-Scholes模型）：在非标准Cox-Ross-Rubinstein二叉树模型中，取时间步长为无穷小 \(\Delta t = T/N\)（\(N\) 无穷大），令上涨因子 \(u = e^{\sigma \sqrt{\Delta t}}\)，下跌因子 \(d = e^{-\sigma \sqrt{\Delta t}}\)，风险中性概率 \(p = \frac{e^{r\Delta t} - d}{u - d}\)。则当 \(N \to \infty\)（在标准部分意义下），期权价格收敛至Black-Scholes公式。

\textbf{证明}：在此非标准二叉树中，由中心极限定理（Nelson激进形式），\(\log S_T\) 的分布经标准部分映射后服从正态分布 \(N(\log S_0 + (r - \sigma^2/2)T, \sigma^2 T)\)。具体地，设 \(S_0\) 为初始股价，则

\[S_T = S_0 \prod_{k=1}^{N} \xi_k,\]

其中 \(\xi_k \in \{u, d\}\) 为独立同分布的内部随机变量，\(P(\xi_k = u) = p\)。取对数得

\[\log(S_T/S_0) = \sum_{k=1}^{N} \log \xi_k.\]

由非标准中心极限定理，该和经标准化后趋近标准正态分布。欧式看涨期权的非标准价格为

\[C_0 = e^{-rT} \sum_{k=0}^{N} \binom{N}{k} p^k (1-p)^{N-k} \max(S_0 u^k d^{N-k} - K, 0).\]

当 \(N \to \infty\)（即 \(N\) 无穷大时取标准部分），此和收敛至Black-Scholes积分公式。 \(\blacksquare\)

\textbf{定理 x+2.11}（非标准随机游走与Wiener过程的构造）：取 \(\Delta t = 1/N\)（\(N\) 无穷大），定义内部随机游走 \(W_t^{(N)} = \sqrt{\Delta t}\sum_{k=1}^{\lfloor t/\Delta t \rfloor} \xi_k\)，其中 \(\xi_k = \pm 1\) 以等概率取值。则对任意标准 \(t\)，\(\operatorname{st}(W_t^{(N)})\) 服从 \(N(0, t)\)，且其标准部分过程满足Wiener过程的几乎所有样本路径性质。

\textbf{证明}：由非标准中心极限定理（Lindberg-Levy在内部有限形式下），每步 \(\pm 1\) 贡献均值为 0、方差为 1。在 \(n = \lfloor t/\Delta t \rfloor\) 步后，和 \(S_n = \sum_{k=1}^{n} \xi_k\) 近似 \(N(0, n)\)，故 \(\sqrt{\Delta t}S_n = S_n/\sqrt{N}\) 近似 \(N(0, t)\)。由Donsker不变原理的非标准版本，\(W_t = \operatorname{st}(W_t^{(N)})\) 的样本路径是Hölder连续的（指数 \(< 1/2\)）且几乎处处不可微。这为Anderson（1976）的非标准Brown运动构造提供了基础。 \(\blacksquare\)

\textbf{定理 x+2.12}（非标准Ito公式）：设 \(f \in C^2(\mathbb{R})\)，\(B_t\) 为非标准随机游走 \(W_t^{(N)}\) 的标准部分。则

\[df(B_t) = f'(B_t) dB_t + \frac{1}{2}f''(B_t) dt\]

在标准部分意义下成立。

\textbf{证明}：对非标准随机游走，使用两步Taylor展开：

\[f(W_{t+\Delta t}^{(N)}) - f(W_t^{(N)}) = f'(W_t^{(N)})\Delta W + \frac{1}{2}f''(W_t^{(N)})(\Delta W)^2 + o((\Delta W)^2),\]

其中 \(\Delta W = \pm \sqrt{\Delta t}\)。注意到 \((\Delta W)^2 = \Delta t\)（非随机），且 \(f'(W_t^{(N)}) \Delta W\) 的累积和的标准部分为Ito积分，\(f''(W_t^{(N)}) \Delta t\) 的累积和的标准部分为Riemann积分。取标准部分即得Ito公式。 \(\blacksquare\)

\newpage

\chapter{01-基础/05-初等代数/01-代数运算.md}\label{ux57faux784005-ux521dux7b49ux4ee3ux657001-ux4ee3ux6570ux8fd0ux7b97.md}

\section{第27章：代数运算与运算律}\label{ux7b2c27ux7ae0ux4ee3ux6570ux8fd0ux7b97ux4e0eux8fd0ux7b97ux5f8b}

\begin{quote}
\textbf{前置知识}：本章建立在第一卷数系构造（第4章）的基础之上。读者应熟悉皮亚诺公理、数学归纳法以及实数的完备性。
\end{quote}

代数运算是数学中最基本的操作，也是所有数学分支的共同语言。本章从二元运算和代数结构的一般概念出发，系统地研究实数域 \(\mathbb{R}\) 上的基本运算律：加法与乘法的结合律、交换律和分配律，以及零元、幺元、负元和逆元的性质。这些运算律共同构成了 \(\mathbb{R}\) 的域结构，是后续多项式运算、方程求解和不等式推导的基础。本章还将引入指数运算（包括正整数指数、零指数、负整数指数和有理数指数），建立指数运算律，并讨论根式运算与有理化方法。理解代数运算律不仅是为了掌握计算技巧，更是为了理解这些运算律背后的代数结构------它们正是抽象代数中群、环、域公理的具体体现。初等代数中''理所当然''的运算规则（如 \(a + b = b + a\)、\(a(b + c) = ab + ac\)）在更高的抽象层次上揭示了数学结构的深层统一性。

\subsection{6.1 二元运算与代数结构}\label{ux4e8cux5143ux8fd0ux7b97ux4e0eux4ee3ux6570ux7ed3ux6784}

\textbf{定义 6.1.1（二元运算）}

设 \(S\) 为非空集合。从 \(S \times S\) 到 \(S\) 的映射

\[
\ast : S \times S \longrightarrow S, \quad (a, b) \longmapsto a \ast b
\]

称为 \(S\) 上的\textbf{二元运算}。

\textbf{定义 6.1.2（封闭性）}

设 \(\ast\) 为集合 \(S\) 上的二元运算。对任意 \(a, b \in S\)，有 \(a \ast b \in S\)，则称 \(\ast\) 在 \(S\) 上\textbf{封闭}。

\textbf{定义 6.1.3（结合律）}

设 \(\ast\) 为集合 \(S\) 上的二元运算。若对任意 \(a, b, c \in S\)，有

\[
(a \ast b) \ast c = a \ast (b \ast c)
\]

则称 \(\ast\) 满足\textbf{结合律}。

\textbf{定义 6.1.4（交换律）}

设 \(\ast\) 为集合 \(S\) 上的二元运算。若对任意 \(a, b \in S\)，有

\[
a \ast b = b \ast a
\]

则称 \(\ast\) 满足\textbf{交换律}。

\textbf{定义 6.1.5（幺元）}

设 \(\ast\) 为集合 \(S\) 上的二元运算。若存在 \(e \in S\)，使得对任意 \(a \in S\)，有

\[
e \ast a = a \ast e = a
\]

则称 \(e\) 为 \(\ast\) 的\textbf{幺元}（恒等元）。

\textbf{定理 6.1.1（幺元唯一性）} 若二元运算 \(\ast\) 存在幺元，则幺元唯一。

\textbf{证明}：设 \(e_1, e_2\) 均为 \(\ast\) 的幺元。由 \(e_1\) 的幺元性质，\(e_1 \ast e_2 = e_2\)；由 \(e_2\) 的幺元性质，\(e_1 \ast e_2 = e_1\)。因此 \(e_1 = e_2\)。 \(\blacksquare\)

\textbf{定义 6.1.6（逆元）}

设 \(\ast\) 为集合 \(S\) 上的二元运算，\(e\) 为幺元。对 \(a \in S\)，若存在 \(b \in S\)，使得

\[
a \ast b = b \ast a = e
\]

则称 \(b\) 为 \(a\) 的\textbf{逆元}，记为 \(a^{-1}\)。

\textbf{定理 6.1.2（逆元唯一性）} 在满足结合律的运算下，若元素的逆元存在，则逆元唯一。

\textbf{证明}：设 \(b, c\) 均为 \(a\) 的逆元，则

\[
b = b \ast e = b \ast (a \ast c) = (b \ast a) \ast c = e \ast c = c
\]

其中第三步利用了结合律。 \(\blacksquare\)

\textbf{定义 6.1.7（半群）}

集合 \(S\) 与其上的二元运算 \(\ast\) 构成的代数结构 \((S, \ast)\)，若满足：

\begin{enumerate}
\def\labelenumi{\arabic{enumi}.}
\tightlist
\item
  \textbf{封闭性}：\(\forall a, b \in S, \, a \ast b \in S\)；
\item
  \textbf{结合律}：\(\forall a, b, c \in S, \, (a \ast b) \ast c = a \ast (b \ast c)\)，
\end{enumerate}

则称 \((S, \ast)\) 为\textbf{半群}。

\textbf{定义 6.1.8（幺半群）}

若半群 \((S, \ast)\) 中存在幺元 \(e\)，则称 \((S, \ast)\) 为\textbf{幺半群}。

\textbf{定义 6.1.9（群）}

若幺半群 \((G, \ast)\) 中每个元素都有逆元，即

\[
\forall a \in G, \, \exists a^{-1} \in G, \, a \ast a^{-1} = a^{-1} \ast a = e
\]

则称 \((G, \ast)\) 为\textbf{群}。若进一步满足交换律，则称为\textbf{阿贝尔群}（交换群）。

\textbf{定义 6.1.10（环）}

集合 \(R\) 上配备两个二元运算 \(+\) 和 \(\cdot\)（分别称为加法和乘法），若满足：

\begin{enumerate}
\def\labelenumi{\arabic{enumi}.}
\tightlist
\item
  \((R, +)\) 是阿贝尔群（加法幺元记为 \(0\)）；
\item
  \((R, \cdot)\) 是半群（乘法满足结合律和封闭性）；
\item
  \textbf{分配律}：\(\forall a, b, c \in R\)，
\end{enumerate}

\[
a \cdot (b + c) = a \cdot b + a \cdot c, \quad (b + c) \cdot a = b \cdot a + c \cdot a
\]

则称 \((R, +, \cdot)\) 为\textbf{环}。

\textbf{定义 6.1.11（交换环与含幺环）}

\begin{itemize}
\tightlist
\item
  若环 \((R, +, \cdot)\) 中乘法满足交换律，则称为\textbf{交换环}。
\item
  若环 \((R, +, \cdot)\) 中乘法存在幺元 \(1\)（\(1 \neq 0\)），则称为\textbf{含幺环}。
\end{itemize}

\textbf{定义 6.1.12（域）}

含幺交换环 \((F, +, \cdot)\) 中，若每个非零元素都有乘法逆元，即

\[
\forall a \in F \setminus \{0\}, \, \exists a^{-1} \in F, \, a \cdot a^{-1} = 1
\]

则称 \((F, +, \cdot)\) 为\textbf{域}。

\hrulefill

\subsection{6.2 自然数运算律}\label{ux81eaux7136ux6570ux8fd0ux7b97ux5f8b}

以下结果建立在皮亚诺公理之上（参见第一卷第4章）。

\textbf{定义 6.2.1（自然数加法）}

加法 \(+: \mathbb{N} \times \mathbb{N} \to \mathbb{N}\) 递归定义为：

\[
a + 0 = a, \quad a + S(b) = S(a + b)
\]

其中 \(S\) 为后继函数。

\textbf{定义 6.2.2（自然数乘法）}

乘法 \(\cdot: \mathbb{N} \times \mathbb{N} \to \mathbb{N}\) 递归定义为：

\[
a \cdot 0 = 0, \quad a \cdot S(b) = a \cdot b + a
\]

\textbf{定理 6.2.1（加法右单位元）} 对任意 \(a \in \mathbb{N}\)，有 \(a + 0 = a\)。

\textbf{证明}：Peano 公理体系中自然数加法的递归定义（定义 6.2.1）的第一条明确规定了 \(a + 0 = a\) 对所有 \(a \in \mathbb{N}\) 成立。因此这正是加法定义本身，无需额外推导。这一性质说明 \(0\) 是自然数加法中的右单位元。\(\blacksquare\)

\textbf{定理 6.2.2（加法左单位元）} 对任意 \(a \in \mathbb{N}\)，有 \(0 + a = a\)。

\textbf{证明}：对 \(a\) 进行归纳。

\begin{itemize}
\tightlist
\item
  \textbf{基础}：\(a = 0\) 时，\(0 + 0 = 0\)，成立。
\item
  \textbf{归纳}：假设 \(0 + a = a\)，则 \(0 + S(a) = S(0 + a) = S(a)\)。
\end{itemize}

由归纳法原理，结论成立。 \(\blacksquare\)

\textbf{定理 6.2.3（加法结合律）} 对任意 \(a, b, c \in \mathbb{N}\)，有 \((a + b) + c = a + (b + c)\)。

\textbf{证明}：对 \(c\) 进行归纳。

\begin{itemize}
\tightlist
\item
  \textbf{基础}：\(c = 0\) 时，\((a + b) + 0 = a + b = a + (b + 0)\)。
\item
  \textbf{归纳}：假设 \((a + b) + c = a + (b + c)\)，则
\end{itemize}

\[
(a + b) + S(c) = S((a + b) + c) = S(a + (b + c)) = a + S(b + c) = a + (b + S(c))
\]

由归纳法原理，结论成立。 \(\blacksquare\)

\textbf{定理 6.2.4（后继的加法表示）} 对任意 \(a \in \mathbb{N}\)，有 \(S(a) = a + 1\)。

\textbf{证明}：由定义，\(1 = S(0)\)（Peano 公理体系中 \(1\) 被定义为 \(0\) 的后继）。根据加法定义的递归规则（定义 6.2.1 的第二条），\(a + S(b) = S(a + b)\)。将 \(b = 0\) 代入，得 \(a + S(0) = S(a + 0)\)。又由定理 6.2.1，\(a + 0 = a\)。因此 \(a + 1 = a + S(0) = S(a + 0) = S(a)\)。这一结论表明：加 \(1\) 恰好等于取后继运算，这是自然数加法最直观的性质。\(\blacksquare\)

\textbf{定理 6.2.5（加法交换律）} 对任意 \(a, b \in \mathbb{N}\)，有 \(a + b = b + a\)。

\textbf{证明}：先证辅助引理：对任意 \(a \in \mathbb{N}\)，有 \(S(a) = a + 1 = 1 + a\)。

对 \(a\) 归纳证明 \(1 + a = S(a)\)：

\begin{itemize}
\tightlist
\item
  \textbf{基础}：\(1 + 0 = 1 = S(0)\)。
\item
  \textbf{归纳}：假设 \(1 + a = S(a)\)，则 \(1 + S(a) = S(1 + a) = S(S(a))\)。
\end{itemize}

现在对 \(b\) 归纳证明 \(a + b = b + a\)：

\begin{itemize}
\tightlist
\item
  \textbf{基础}：\(b = 0\) 时，\(a + 0 = a = 0 + a\)。
\item
  \textbf{归纳}：假设 \(a + b = b + a\)，则
\end{itemize}

\[
a + S(b) = S(a + b) = S(b + a) = b + S(a) = b + (a + 1) = b + (1 + a) = (b + 1) + a = S(b) + a
\]

由归纳法原理，结论成立。 \(\blacksquare\)

\textbf{定理 6.2.6（乘法右零元）} 对任意 \(a \in \mathbb{N}\)，有 \(a \cdot 0 = 0\)。

\textbf{证明}：Peano 公理体系中自然数乘法的递归定义（定义 6.3.1）的第一条明确规定了 \(a \cdot 0 = 0\) 对所有 \(a \in \mathbb{N}\) 成立。这是乘法定义本身，无需额外推导。这一性质说明任何自然数乘以 \(0\) 都等于 \(0\)，与整数运算中''零乘任何数等于零''的直觉一致。\(\blacksquare\)

\textbf{定理 6.2.7（乘法左零元）} 对任意 \(a \in \mathbb{N}\)，有 \(0 \cdot a = 0\)。

\textbf{证明}：对 \(a\) 进行归纳。

\begin{itemize}
\tightlist
\item
  \textbf{基础}：\(0 \cdot 0 = 0\)。
\item
  \textbf{归纳}：假设 \(0 \cdot a = 0\)，则 \(0 \cdot S(a) = 0 \cdot a + 0 = 0 + 0 = 0\)。
\end{itemize}

由归纳法原理，结论成立。 \(\blacksquare\)

\textbf{定理 6.2.8（乘法右单位元）} 对任意 \(a \in \mathbb{N}\)，有 \(a \cdot 1 = a\)。

\textbf{证明}：由乘法递归定义的第二条和已证定理。\(a \cdot 1 = a \cdot S(0) = a \cdot 0 + a\)（利用 \(a \cdot S(b) = a \cdot b + a\) 并令 \(b = 0\)）。由定理 6.2.6，\(a \cdot 0 = 0\)。代入得 \(a \cdot 1 = 0 + a = a\)（由定理 6.2.2，\(0 + a = a\)）。因此 \(1\) 是自然数乘法中的右单位元。结合定理 6.2.9（乘法交换律）可得 \(1 \cdot a = a\)，即 \(1\) 也是乘法左单位元。\(\blacksquare\)

\textbf{定理 6.2.9（乘法左分配律）} 对任意 \(a, b, c \in \mathbb{N}\)，有 \(a \cdot (b + c) = a \cdot b + a \cdot c\)。

\textbf{证明}：对 \(c\) 进行归纳。

\begin{itemize}
\tightlist
\item
  \textbf{基础}：\(c = 0\) 时，\(a \cdot (b + 0) = a \cdot b = a \cdot b + 0 = a \cdot b + a \cdot 0\)。
\item
  \textbf{归纳}：假设 \(a \cdot (b + c) = a \cdot b + a \cdot c\)，则
\end{itemize}

\[
a \cdot (b + S(c)) = a \cdot S(b + c) = a \cdot (b + c) + a = (a \cdot b + a \cdot c) + a = a \cdot b + (a \cdot c + a) = a \cdot b + a \cdot S(c)
\]

由归纳法原理，结论成立。 \(\blacksquare\)

\textbf{定理 6.2.10（乘法结合律）} 对任意 \(a, b, c \in \mathbb{N}\)，有 \((a \cdot b) \cdot c = a \cdot (b \cdot c)\)。

\textbf{证明}：对 \(c\) 进行归纳。

\begin{itemize}
\tightlist
\item
  \textbf{基础}：\(c = 0\) 时，\((a \cdot b) \cdot 0 = 0 = a \cdot 0 = a \cdot (b \cdot 0)\)。
\item
  \textbf{归纳}：假设 \((a \cdot b) \cdot c = a \cdot (b \cdot c)\)，则
\end{itemize}

\[
(a \cdot b) \cdot S(c) = (a \cdot b) \cdot c + a \cdot b = a \cdot (b \cdot c) + a \cdot b = a \cdot (b \cdot c + b) = a \cdot (b \cdot S(c))
\]

其中第三步利用了左分配律。 \(\blacksquare\)

\textbf{定理 6.2.11（乘法交换律）} 对任意 \(a, b \in \mathbb{N}\)，有 \(a \cdot b = b \cdot a\)。

\textbf{证明}：先证辅助引理：对任意 \(a \in \mathbb{N}\)，有 \(1 \cdot a = a\)。对 \(a\) 归纳：

\begin{itemize}
\tightlist
\item
  \textbf{基础}：\(1 \cdot 0 = 0\)。
\item
  \textbf{归纳}：假设 \(1 \cdot a = a\)，则 \(1 \cdot S(a) = 1 \cdot a + 1 = a + 1 = S(a)\)。
\end{itemize}

再证另一辅助引理：\(S(a) \cdot b = a \cdot b + b\)。对 \(b\) 归纳：

\begin{itemize}
\tightlist
\item
  \textbf{基础}：\(S(a) \cdot 0 = 0 = 0 + 0 = a \cdot 0 + 0\)。
\item
  \textbf{归纳}：假设 \(S(a) \cdot b = a \cdot b + b\)，则
\end{itemize}

\[
S(a) \cdot S(b) = S(a) \cdot b + S(a) = (a \cdot b + b) + (a + 1) = (a \cdot b + a) + (b + 1) = a \cdot S(b) + S(b)
\]

现在对 \(b\) 归纳证明 \(a \cdot b = b \cdot a\)：

\begin{itemize}
\tightlist
\item
  \textbf{基础}：\(a \cdot 0 = 0 = 0 \cdot a\)。
\item
  \textbf{归纳}：假设 \(a \cdot b = b \cdot a\)，则 \(a \cdot S(b) = a \cdot b + a = b \cdot a + a = S(b) \cdot a\)。
\end{itemize}

由归纳法原理，结论成立。 \(\blacksquare\)

\textbf{定理 6.2.12（乘法右分配律）} 对任意 \(a, b, c \in \mathbb{N}\)，有 \((a + b) \cdot c = a \cdot c + b \cdot c\)。

\textbf{证明}：由乘法交换律和左分配律，

\[
(a + b) \cdot c = c \cdot (a + b) = c \cdot a + c \cdot b = a \cdot c + b \cdot c
\]

\(\blacksquare\)

\hrulefill

\subsection{6.3 整数运算律}\label{ux6574ux6570ux8fd0ux7b97ux5f8b}

整数 \(\mathbb{Z}\) 通过自然数的等价类构造（参见第一卷第4章）。设 \(a, b \in \mathbb{N}\)，整数 \(-a\) 表示使得 \(a + (-a) = 0\) 的元素。

\textbf{引理 6.3.0}：整数的乘法交换律和右分配律可从自然数继承的左分配律和乘法交换律推导。

\textbf{证明}：由自然数的乘法交换律（定理 6.2.11），通过等价类构造可得整数的乘法交换律 \(ab = ba\)。

右分配律 \((a + b) \cdot c = a \cdot c + b \cdot c\)：由乘法交换律，\((a+b) \cdot c = c \cdot (a+b) = c \cdot a + c \cdot b = a \cdot c + b \cdot c\)（第二步利用了已有的左分配律）。

因此下文的分配律推广 \((a+b)(c+d) = ac + ad + bc + bd\) 合法。\(\blacksquare\)

\textbf{定理 6.3.1（负负得正）} \((-1) \times (-1) = 1\)。

\textbf{证明}：考虑以下推导链。

由分配律的推广，对任意 \(a, b, c, d \in \mathbb{Z}\)：

\[
(a + b)(c + d) = ac + ad + bc + bd
\]

取 \(a = 1, b = -1, c = 1, d = -1\)：

\[
(1 + (-1))(1 + (-1)) = 1 \cdot 1 + 1 \cdot (-1) + (-1) \cdot 1 + (-1) \cdot (-1)
\]

左边为 \(0 \cdot 0 = 0\)。右边中 \(1 \cdot 1 = 1\)，\(1 \cdot (-1) = -1\)（由分配律：\(1 \cdot (-1) + 1 \cdot 1 = 1 \cdot ((-1) + 1) = 1 \cdot 0 = 0\)，故 \(1 \cdot (-1) = -1\)），\((-1) \cdot 1 = -1\)（由交换律），因此：

\[
0 = 1 + (-1) + (-1) + (-1)(-1) = -1 + (-1)(-1)
\]

从而 \((-1)(-1) = 1\)。 \(\blacksquare\)

\textbf{定理 6.3.2（负数乘法的一般法则）} 对任意 \(a, b \in \mathbb{Z}\)，有 \((-a)(-b) = ab\)。

\textbf{证明}：

\[
(-a)(-b) = (-1 \cdot a)(-1 \cdot b) = (-1)(-1) \cdot ab = 1 \cdot ab = ab
\]

其中利用了乘法结合律和交换律以及定理6.3.1。 \(\blacksquare\)

\textbf{定理 6.3.3} 对任意 \(a, b \in \mathbb{Z}\)，有 \((-a) \cdot b = -(a \cdot b)\)。

\textbf{证明}：需证 \((-a) \cdot b + a \cdot b = 0\)。由分配律，

\[
(-a) \cdot b + a \cdot b = ((-a) + a) \cdot b = 0 \cdot b = 0
\]

由逆元唯一性，\((-a) \cdot b = -(a \cdot b)\)。 \(\blacksquare\)

\hrulefill

\subsection{6.4 实数运算律（域公理）}\label{ux5b9eux6570ux8fd0ux7b97ux5f8bux57dfux516cux7406}

实数集 \(\mathbb{R}\) 配备加法和乘法运算构成域。以下为完整的域公理体系。

\textbf{加法公理（A1---A4）}：

\begin{itemize}
\tightlist
\item
  \textbf{(A1) 加法封闭性}：\(\forall a, b \in \mathbb{R}\)，\(a + b \in \mathbb{R}\)。
\item
  \textbf{(A2) 加法结合律}：\(\forall a, b, c \in \mathbb{R}\)，\((a + b) + c = a + (b + c)\)。
\item
  \textbf{(A3) 加法交换律}：\(\forall a, b \in \mathbb{R}\)，\(a + b = b + a\)。
\item
  \textbf{(A4) 加法单位元与逆元}：存在 \(0 \in \mathbb{R}\)，使得 \(\forall a \in \mathbb{R}\)，\(a + 0 = a\)；对每个 \(a \in \mathbb{R}\)，存在 \(-a \in \mathbb{R}\)，使得 \(a + (-a) = 0\)。
\end{itemize}

\textbf{乘法公理（M1---M4）}：

\begin{itemize}
\tightlist
\item
  \textbf{(M1) 乘法封闭性}：\(\forall a, b \in \mathbb{R}\)，\(a \cdot b \in \mathbb{R}\)。
\item
  \textbf{(M2) 乘法结合律}：\(\forall a, b, c \in \mathbb{R}\)，\((a \cdot b) \cdot c = a \cdot (b \cdot c)\)。
\item
  \textbf{(M3) 乘法交换律}：\(\forall a, b \in \mathbb{R}\)，\(a \cdot b = b \cdot a\)。
\item
  \textbf{(M4) 乘法单位元与逆元}：存在 \(1 \in \mathbb{R}\)（\(1 \neq 0\)），使得 \(\forall a \in \mathbb{R}\)，\(a \cdot 1 = a\)；对每个 \(a \in \mathbb{R} \setminus \{0\}\)，存在 \(a^{-1} \in \mathbb{R}\)，使得 \(a \cdot a^{-1} = 1\)。
\end{itemize}

\textbf{分配律（D）}：

\begin{itemize}
\tightlist
\item
  \textbf{(D) 分配律}：\(\forall a, b, c \in \mathbb{R}\)，\(a \cdot (b + c) = a \cdot b + a \cdot c\)。
\end{itemize}

\textbf{定理 6.4.1（加法消去律）} 若 \(a + c = b + c\)，则 \(a = b\)。

\textbf{证明}：从 \(a\) 出发，利用加法零元和逆元的性质逐步推导。\(a = a + 0\)（零元性质，任何数加 \(0\) 等于自身）。\(0 = c + (-c)\)（逆元性质，\(c\) 与其加法逆元 \(-c\) 的和为 \(0\)）。将 \(0\) 替换为 \(c + (-c)\)，得 \(a = a + (c + (-c))\)。由加法结合律，\(a + (c + (-c)) = (a + c) + (-c)\)。由已知条件 \(a + c = b + c\)，将 \(a + c\) 替换为 \(b + c\)，得 \(b + c + (-c)\)。再次应用结合律，\(b + c + (-c) = b + (c + (-c)) = b + 0 = b\)。因此 \(a = b\)。整个推导的关键在于：在环（或域）中，消去律等价于逆元的存在性------通过加上 \(-c\)，我们同时消去了等式两边的 \(c\)。\(\blacksquare\)

\textbf{定理 6.4.2（乘法消去律）} 若 \(a \cdot c = b \cdot c\) 且 \(c \neq 0\)，则 \(a = b\)。

\textbf{证明}：从 \(a\) 出发，利用乘法单位元和逆元的性质。\(a = a \cdot 1\)（单位元性质，任何数乘以 \(1\) 等于自身）。由于 \(c \neq 0\)，\(c\) 在域中有乘法逆元 \(c^{-1}\)，满足 \(c \cdot c^{-1} = 1\)。将 \(1\) 替换为 \(c \cdot c^{-1}\)，得 \(a = a \cdot (c \cdot c^{-1})\)。由乘法结合律，\(a \cdot (c \cdot c^{-1}) = (a \cdot c) \cdot c^{-1}\)。由已知条件 \(a \cdot c = b \cdot c\)，将 \(a \cdot c\) 替换为 \(b \cdot c\)，得 \((b \cdot c) \cdot c^{-1}\)。再次应用结合律，\((b \cdot c) \cdot c^{-1} = b \cdot (c \cdot c^{-1}) = b \cdot 1 = b\)。因此 \(a = b\)。注意这里 \(c \neq 0\) 的条件是必要的------如果 \(c = 0\)，则 \(c^{-1}\) 不存在，消去律不成立（例如 \(2 \cdot 0 = 3 \cdot 0\) 但 \(2 \neq 3\)）。\(\blacksquare\)

\textbf{定理 6.4.3（零因子性质）} 若 \(a \cdot b = 0\)，则 \(a = 0\) 或 \(b = 0\)。

\textbf{证明}：假设 \(a \neq 0\)，则 \(a^{-1}\) 存在。\(b = 1 \cdot b = (a^{-1} \cdot a) \cdot b = a^{-1} \cdot (a \cdot b) = a^{-1} \cdot 0 = 0\)。

其中最后一步 \(a^{-1} \cdot 0 = 0\) 的证明如下：设 \(x = a^{-1} \cdot 0\)，则 \(x = a^{-1} \cdot (0 + 0) = a^{-1} \cdot 0 + a^{-1} \cdot 0 = x + x\)。由消去律 \(x = 0\)。 \(\blacksquare\)

\hrulefill

\subsection{6.5 指数运算律}\label{ux6307ux6570ux8fd0ux7b97ux5f8b}

指数运算是代数运算中最重要的扩展之一，它将乘法运算推广为重复乘法的简洁记法。本节从正整数指数出发，通过归纳定义建立 \(a^n\)（\(n \in \mathbb{N}^+\)）的基本运算法则（\(a^m \cdot a^n = a^{m+n}\) 和 \((a^m)^n = a^{mn}\)），然后逐步将指数定义扩展到零指数（\(a^0 = 1\)，\(a \neq 0\)）、负整数指数（\(a^{-n} = 1/a^n\)）和有理数指数（\(a^{m/n} = \sqrt[n]{a^m}\)）。每一次扩展都遵循一个基本原则：新定义必须与旧定义相容，且保持指数运算律的成立。指数运算律的统一性------\(a^m a^n = a^{m+n}\)，\((a^m)^n = a^{mn}\)，\((ab)^n = a^n b^n\)------贯穿所有指数类型，体现了数学中''结构保持''的扩展策略。理解指数运算律不仅有助于简化复杂的代数表达式，也为对数运算和指数函数的学习奠定了基础。指数从正整数到实数的扩展过程，本质上是一个从''离散''到''连续''的极限过程，与实数完备性有着深刻的联系。

\subsubsection{6.5.1 正整数指数}\label{ux6b63ux6574ux6570ux6307ux6570}

\textbf{定义 6.5.1} 设 \(a \in \mathbb{R}\)，\(n \in \mathbb{N}^+\)，归纳定义：

\[
a^1 = a, \quad a^{n+1} = a^n \cdot a
\]

\textbf{定理 6.5.1（指数加法法则）} 对 \(a \in \mathbb{R}\)，\(m, n \in \mathbb{N}^+\)，有 \(a^m \cdot a^n = a^{m+n}\)。

\textbf{证明}：对 \(n\) 归纳。

\begin{itemize}
\tightlist
\item
  \textbf{基础}：\(n = 1\) 时，\(a^m \cdot a^1 = a^m \cdot a = a^{m+1}\)。
\item
  \textbf{归纳}：假设 \(a^m \cdot a^n = a^{m+n}\)，则 \(a^m \cdot a^{n+1} = a^m \cdot (a^n \cdot a) = (a^m \cdot a^n) \cdot a = a^{m+n} \cdot a = a^{m+n+1}\)。
\end{itemize}

由归纳法原理，结论成立。 \(\blacksquare\)

\textbf{定理 6.5.2（指数乘法法则）} 对 \(a \in \mathbb{R}\)，\(m, n \in \mathbb{N}^+\)，有 \((a^m)^n = a^{mn}\)。

\textbf{证明}：对 \(n\) 归纳。

\begin{itemize}
\tightlist
\item
  \textbf{基础}：\((a^m)^1 = a^m = a^{m \cdot 1}\)。
\item
  \textbf{归纳}：假设 \((a^m)^n = a^{mn}\)，则 \((a^m)^{n+1} = (a^m)^n \cdot a^m = a^{mn} \cdot a^m = a^{mn+m} = a^{m(n+1)}\)。
\end{itemize}

\(\blacksquare\)

\textbf{定理 6.5.3（积的幂）} 对 \(a, b \in \mathbb{R}\)，\(n \in \mathbb{N}^+\)，有 \((ab)^n = a^n b^n\)。

\textbf{证明}：对 \(n\) 归纳。

\begin{itemize}
\tightlist
\item
  \textbf{基础}：\((ab)^1 = ab = a^1 b^1\)。
\item
  \textbf{归纳}：假设 \((ab)^n = a^n b^n\)，则 \((ab)^{n+1} = (ab)^n \cdot ab = a^n b^n \cdot a \cdot b = a^{n+1} b^{n+1}\)。
\end{itemize}

\(\blacksquare\)

\subsubsection{6.5.2 零指数与负整数指数}\label{ux96f6ux6307ux6570ux4e0eux8d1fux6574ux6570ux6307ux6570}

\textbf{定义 6.5.2} 设 \(a \in \mathbb{R}\)，\(a \neq 0\)：

\[
a^0 = 1, \quad a^{-n} = \frac{1}{a^n} \quad (n \in \mathbb{N}^+)
\]

\textbf{定理 6.5.4} 对 \(a \neq 0\)，\(m, n \in \mathbb{Z}\)，上述三条运算律仍然成立。

\textbf{证明}（以加法法则为例，证 \(a^m \cdot a^n = a^{m+n}\)）：分情况讨论。

\textbf{情形一}：\(m, n > 0\)。由定理6.5.1成立。

\textbf{情形二}：\(m > 0, n = 0\)。\(a^m \cdot a^0 = a^m \cdot 1 = a^m = a^{m+0}\)。

\textbf{情形三}：\(m > 0, n < 0\)。设 \(n = -k\)（\(k > 0\)）。

\begin{itemize}
\tightlist
\item
  若 \(m > k\)：\(a^m \cdot a^{-k} = \dfrac{a^m}{a^k} = a^{m-k} = a^{m+(-k)}\)。
\item
  若 \(m = k\)：\(a^m \cdot a^{-m} = \dfrac{a^m}{a^m} = 1 = a^0 = a^{m+(-m)}\)。
\item
  若 \(m < k\)：\(a^m \cdot a^{-k} = \dfrac{a^m}{a^k} = \dfrac{1}{a^{k-m}} = a^{-(k-m)} = a^{m-k}\)。
\end{itemize}

\textbf{情形四}：\(m = 0\)。\(a^0 \cdot a^n = 1 \cdot a^n = a^n = a^{0+n}\)。

\textbf{情形五}：\(m < 0, n < 0\)。\(a^m \cdot a^{-n} = \dfrac{1}{a^{|m|}} \cdot \dfrac{1}{a^{|n|}} = \dfrac{1}{a^{|m|+|n|}} = a^{-(|m|+|n|)} = a^{m+n}\)。

所有情形下 \(a^m \cdot a^n = a^{m+n}\) 成立。指数乘法法则和积的幂的推广类似可证。 \(\blacksquare\)

\subsubsection{6.5.3 有理数指数}\label{ux6709ux7406ux6570ux6307ux6570}

\textbf{定义 6.5.3} 设 \(a > 0\)，\(m, n \in \mathbb{Z}\)，\(n \geq 1\)，定义：

\[
a^{m/n} = \left(\sqrt[n]{a}\right)^m
\]

其中 \(\sqrt[n]{a}\) 为 \(a\) 的 \(n\) 次算术根，即满足 \(x^n = a\) 且 \(x > 0\) 的唯一实数。

\textbf{定理 6.5.5（有理数指数的良定义性）} 若 \(m/n = p/q\)（即 \(mq = np\)），则 \(a^{m/n} = a^{p/q}\)。

\textbf{证明}：设 \(m/n = p/q\)，即 \(mq = np\)。需证 \((\sqrt[n]{a})^m = (\sqrt[q]{a})^p\)。

设 \(x = \sqrt[n]{a}\)，\(y = \sqrt[q]{a}\)，则 \(x^n = a\)，\(y^q = a\)。对等式两端取 \(nq\) 次幂：

\[
\left((\sqrt[n]{a})^m\right)^{nq} = (\sqrt[n]{a})^{mnq} = \left((\sqrt[n]{a})^n\right)^{mq} = a^{mq}
\]

\[
\left((\sqrt[q]{a})^p\right)^{nq} = (\sqrt[q]{a})^{pnq} = \left((\sqrt[q]{a})^q\right)^{np} = a^{np}
\]

由 \(mq = np\)，两式右端相等，即 \(\left((\sqrt[n]{a})^m\right)^{nq} = \left((\sqrt[q]{a})^p\right)^{nq}\)。由于 \(nq > 0\) 且 \((\sqrt[n]{a})^m > 0\)，\((\sqrt[q]{a})^p > 0\)，正数的 \(nq\) 次幂是单射，故 \((\sqrt[n]{a})^m = (\sqrt[q]{a})^p\)。 \(\blacksquare\)

\textbf{定理 6.5.6（有理数指数的运算律）} 设 \(a > 0\)，\(r, s \in \mathbb{Q}\)，则：

\begin{enumerate}
\def\labelenumi{(\roman{enumi})}
\item
  \(a^r \cdot a^s = a^{r+s}\)；
\item
  \((a^r)^s = a^{rs}\)；
\item
  \((ab)^r = a^r \cdot b^r\)（\(b > 0\)）。
\end{enumerate}

\textbf{证明}（以 (i) 为例）：设 \(r = m/n\)，\(s = p/q\)，其中 \(n, q \geq 1\)。

\[
a^r \cdot a^s = (\sqrt[n]{a})^m \cdot (\sqrt[q]{a})^p
\]

令 \(d = \text{lcm}(n, q)\)，\(n' = d/n\)，\(q' = d/q\)。则 \(\sqrt[n]{a} = (\sqrt[d]{a})^{n'}\)，\(\sqrt[q]{a} = (\sqrt[d]{a})^{q'}\)。于是

\[
a^r = (\sqrt[d]{a})^{mn'}, \quad a^s = (\sqrt[d]{a})^{pq'}
\]

\[
a^r \cdot a^s = (\sqrt[d]{a})^{mn' + pq'}
\]

而 \(r + s = m/n + p/q = (mq + np)/(nq)\)。由于 \(mn' + pq' = md/n + pd/q = (mq + np) \cdot d/(nq)\)，且 \((\sqrt[d]{a})^{d/(nq)}\) 对应分母为 \(nq\) 的根\ldots{} 实际上 \(a^{r+s} = (\sqrt[nq]{a})^{mq+np}\)，而

\[
(\sqrt[d]{a})^{mn'+pq'} = a^{(mn'+pq')/d} = a^{m/n + p/q} = a^{r+s}
\]

\(\blacksquare\)

\subsubsection{6.5.4 实数指数}\label{ux5b9eux6570ux6307ux6570}

\textbf{定义 6.5.4} 设 \(a > 0\)，\(x \in \mathbb{R}\)。定义

\[
a^x = \sup \{ a^r : r \in \mathbb{Q}, \, r \leq x \}
\]

（利用实数的完备性保证上确界存在。）

\textbf{定理 6.5.7}（实数指数运算律------陈述）实数指数满足与有理数指数相同的三条运算律。

*\textbf{证明}：由代数运算的定义和基本性质，通过逐项验证结合律、交换律和分配律，可得所述恒等式。\(\blacksquare\)

指数运算律是代数学中最基本的运算规则之一，它们统一了从正整数指数到实数指数的整个指数体系。正整数指数的定义是递归的（\(a^{n+1} = a^n \cdot a\)），这直接反映了指数运算的本质------重复乘法。零指数和负整数指数的引入（\(a^0 = 1\)，\(a^{-n} = 1/a^n\)）使得指数法则在整数范围内保持一致，这是指数概念从自然数到整数的自然推广。有理数指数通过根式来定义（\(a^{m/n} = \sqrt[n]{a^m}\)），其合理性由根式的存在唯一性（定理 6.6.1）保证。实数指数是有理数指数通过完备性（确界原理）的极限推广------\(a^x = \sup\{a^r : r \in \mathbb{Q}, r \leq x\}\)（对 \(a > 1\)）------这一定义使得指数函数 \(x \mapsto a^x\) 在 \(\mathbb{R}\) 上连续。三条核心运算律（\(a^m \cdot a^n = a^{m+n}\)、\((a^m)^n = a^{mn}\)、\((ab)^n = a^n b^n\)）贯穿所有指数类型，体现了指数运算的普适代数结构。在抽象代数中，指数运算律的群论解释是：映射 \(n \mapsto a^n\) 是加法群 \((\mathbb{Z}, +)\) 到乘法群 \((\mathbb{R}^\times, \cdot)\) 的群同态。

\textbf{定理 6.6.2（根式除法法则）} 设 \(a \geq 0\)，\(b > 0\)，则 \(\dfrac{\sqrt{a}}{\sqrt{b}} = \sqrt{\dfrac{a}{b}}\)。

\textbf{证明}：设 \(x = \sqrt{a}/\sqrt{b} \geq 0\)，则

\[
x^2 = \frac{(\sqrt{a})^2}{(\sqrt{b})^2} = \frac{a}{b}
\]

由算术平方根的唯一性，\(x = \sqrt{a/b}\)。 \(\blacksquare\)

\textbf{定义 6.6.2（分母有理化）}

将分母中的根号消去的变换称为\textbf{分母有理化}。

\textbf{定理 6.6.3} 设 \(a \geq 0\)，\(b > 0\)，则

\[
\frac{a}{\sqrt{b}} = \frac{a\sqrt{b}}{b}
\]

\textbf{证明}：

\[
\frac{a}{\sqrt{b}} = \frac{a}{\sqrt{b}} \cdot \frac{\sqrt{b}}{\sqrt{b}} = \frac{a\sqrt{b}}{(\sqrt{b})^2} = \frac{a\sqrt{b}}{b}
\]

\(\blacksquare\)

\textbf{定理 6.6.4（共轭有理化）} 设 \(a \geq 0\)，\(b > 0\)，\(a \neq b\)，则

\[
\frac{1}{\sqrt{a} + \sqrt{b}} = \frac{\sqrt{a} - \sqrt{b}}{a - b} \quad (a \neq b)
\]

\textbf{证明}：

\[
\frac{1}{\sqrt{a} + \sqrt{b}} = \frac{\sqrt{a} - \sqrt{b}}{(\sqrt{a} + \sqrt{b})(\sqrt{a} - \sqrt{b})} = \frac{\sqrt{a} - \sqrt{b}}{a - b}
\]

其中利用了 \((\sqrt{a} + \sqrt{b})(\sqrt{a} - \sqrt{b}) = a - b\)。 \(\blacksquare\)

\begin{quote}
\textbf{本章展望}：本章建立了代数运算的公理化基础。从半群到域的层级结构，以及自然数、整数、实数上的具体运算律，为后续多项式、方程和不等式的理论奠定了基石。第13章将在此基础上研究多项式的代数结构。
\end{quote}

\hrulefill

\hrulefill

\hrulefill

\hrulefill

\hrulefill

\newpage

\chapter{01-基础/05-初等代数/03-方程与不等式.md}\label{ux57faux784005-ux521dux7b49ux4ee3ux657003-ux65b9ux7a0bux4e0eux4e0dux7b49ux5f0f.md}

\section{第28章：方程与方程组}\label{ux7b2c28ux7ae0ux65b9ux7a0bux4e0eux65b9ux7a0bux7ec4}

\begin{quote}
\textbf{前置知识}：本章需要第12章的运算律和第13章的多项式理论（特别是因式分解、余数定理和韦达定理）。
\end{quote}

方程是数学中最古老也最核心的主题之一，它研究未知量与已知量之间的等式关系。本章从方程的基本概念（方程、解、解集、同解变形）出发，系统地研究各类方程的求解方法。一元一次方程是最简单的方程类型，其求解依赖于等式的运算性质；一元二次方程通过配方法导出求根公式 \(x = \frac{-b \pm \sqrt{b^2 - 4ac}}{2a}\)，其中判别式 \(\Delta = b^2 - 4ac\) 决定了根的类型（两个不等实根、重根、或共轭复根）。高次方程的处理依赖于因式定理------将求根问题转化为因式分解问题，而韦达定理则揭示了根与系数之间的深刻关系。方程组的研究则涉及消元法（代入消元、加减消元）和矩阵方法，将多变量问题转化为单变量问题。本章还将讨论分式方程和无理方程的解法，以及增根和验根的问题------这些是在方程变形中必须注意的逻辑陷阱。

\subsection{8.1 方程的基本概念}\label{ux65b9ux7a0bux7684ux57faux672cux6982ux5ff5}

\textbf{定义 8.1.1（方程）}

含有未知数的等式称为\textbf{方程}。使方程成立的未知数的值称为方程的\textbf{解}。

\textbf{定义 8.1.2（解集）}

方程所有解的集合称为\textbf{解集}。

\textbf{定义 8.1.3（同解方程）}

若两个方程有相同的解集，则称它们\textbf{同解}。

\textbf{定义 8.1.4（同解变换）}

以下变换不改变方程的解集：

\begin{enumerate}
\def\labelenumi{\arabic{enumi}.}
\tightlist
\item
  \textbf{等量加减}：方程两边同时加上（或减去）同一个数或整式。
\item
  \textbf{等量乘除}：方程两边同时乘以（或除以）同一个非零数或整式。
\item
  \textbf{等量替换}：将方程中的某表达式替换为与之相等的另一表达式。
\end{enumerate}

\textbf{定理 8.1.1} 设 \(f(x) = g(x)\) 是一个方程，\(h(x)\) 为任意表达式。则 \(f(x) = g(x)\) 与 \(f(x) + h(x) = g(x) + h(x)\) 同解。

\textbf{证明}：设 \(x_0\) 是 \(f(x) = g(x)\) 的解，则 \(f(x_0) = g(x_0)\)，从而 \(f(x_0) + h(x_0) = g(x_0) + h(x_0)\)，即 \(x_0\) 也是变换后方程的解。反之亦然。 \(\blacksquare\)

\hrulefill

\subsection{8.2 一元一次方程}\label{ux4e00ux5143ux4e00ux6b21ux65b9ux7a0b}

\textbf{定理 8.2.1（一元一次方程解的存在唯一性）} 设 \(a, b \in \mathbb{R}\)，\(a \neq 0\)。则方程

\[
ax + b = 0
\]

有唯一解 \(x = -b/a\)。

\textbf{证明}：

\textbf{存在性}：代入 \(x = -b/a\)，\(a \cdot (-b/a) + b = -b + b = 0\)。故 \(x = -b/a\) 是解。

\textbf{唯一性}：设 \(x_0\) 为任意解，即 \(ax_0 + b = 0\)。则 \(ax_0 = -b\)，由乘法消去律（\(a \neq 0\)），\(x_0 = -b/a\)。 \(\blacksquare\)

\textbf{推论 8.2.1} 若 \(a = 0\)，\(b \neq 0\)，则 \(ax + b = 0\) 无解；若 \(a = 0\)，\(b = 0\)，则 \(ax + b = 0\) 的解为全体实数。

\textbf{证明}：若 \(a = 0\)，方程变为 \(b = 0\)。若 \(b \neq 0\)，矛盾，无解；若 \(b = 0\)，恒等式成立，全体实数为解。 \(\blacksquare\)

\hrulefill

\subsection{8.3 二元一次方程组}\label{ux4e8cux5143ux4e00ux6b21ux65b9ux7a0bux7ec4}

\textbf{定义 8.3.1（二元一次方程组）}

形如

\[
\begin{cases}
a_1 x + b_1 y = c_1 \\
a_2 x + b_2 y = c_2
\end{cases}
\]

的方程组称为\textbf{二元一次方程组}（线性方程组）。

\textbf{定义 8.3.2（系数行列式）}

\[
D = \begin{vmatrix} a_1 & b_1 \\ a_2 & b_2 \end{vmatrix} = a_1 b_2 - a_2 b_1
\]

\textbf{定理 8.3.1（克拉默法则）} 设 \(D = a_1 b_2 - a_2 b_1 \neq 0\)，则二元一次方程组有唯一解：

\[
x = \frac{D_x}{D} = \frac{c_1 b_2 - c_2 b_1}{a_1 b_2 - a_2 b_1}, \quad y = \frac{D_y}{D} = \frac{a_1 c_2 - a_2 c_1}{a_1 b_2 - a_2 b_1}
\]

其中

\[
D_x = \begin{vmatrix} c_1 & b_1 \\ c_2 & b_2 \end{vmatrix} = c_1 b_2 - c_2 b_1, \quad D_y = \begin{vmatrix} a_1 & c_1 \\ a_2 & c_2 \end{vmatrix} = a_1 c_2 - a_2 c_1
\]

\textbf{证明}：

\textbf{第一步：推导解的表达式。}

将第一个方程乘以 \(b_2\)，第二个方程乘以 \(b_1\)：

\[
\begin{cases}
a_1 b_2 x + b_1 b_2 y = c_1 b_2 \\
a_2 b_1 x + b_1 b_2 y = c_2 b_1
\end{cases}
\]

两式相减：

\[
(a_1 b_2 - a_2 b_1) x = c_1 b_2 - c_2 b_1
\]

即 \(Dx = D_x\)。由 \(D \neq 0\)，\(x = D_x / D\)。

类似地，将第一个方程乘以 \(a_2\)，第二个方程乘以 \(a_1\)：

\[
\begin{cases}
a_1 a_2 x + a_2 b_1 y = a_2 c_1 \\
a_1 a_2 x + a_1 b_2 y = a_1 c_2
\end{cases}
\]

相减：\((a_1 b_2 - a_2 b_1) y = a_1 c_2 - a_2 c_1\)，即 \(Dy = D_y\)，\(y = D_y/D\)。

\textbf{第二步：验证。} 将 \(x = D_x/D\)，\(y = D_y/D\) 代入原方程组：

\[
a_1 \cdot \frac{D_x}{D} + b_1 \cdot \frac{D_y}{D} = \frac{a_1 D_x + b_1 D_y}{D}
\]

计算 \(a_1 D_x + b_1 D_y\)：

\[
a_1(c_1 b_2 - c_2 b_1) + b_1(a_1 c_2 - a_2 c_1) = a_1 c_1 b_2 - a_1 c_2 b_1 + a_1 b_1 c_2 - a_2 b_1 c_1 = c_1(a_1 b_2 - a_2 b_1) = c_1 D
\]

因此 \(a_1 (D_x/D) + b_1 (D_y/D) = c_1 D / D = c_1\)。类似可验证第二个方程。

\textbf{第三步：唯一性。} 设 \((x_0, y_0)\) 为任意解，则

\[
\begin{cases} a_1 x_0 + b_1 y_0 = c_1 \\ a_2 x_0 + b_2 y_0 = c_2 \end{cases}
\]

由第一步的推导过程（这是同解变换），\((x_0, y_0)\) 必须满足 \(D x_0 = D_x\) 和 \(D y_0 = D_y\)。由 \(D \neq 0\)，\(x_0 = D_x/D\)，\(y_0 = D_y/D\)。 \(\blacksquare\)

\textbf{推论 8.3.1（\(D = 0\) 的情形）}

\begin{itemize}
\tightlist
\item
  若 \(D = 0\) 且 \(D_x \neq 0\) 或 \(D_y \neq 0\)，则方程组无解。
\item
  若 \(D = 0\) 且 \(D_x = D_y = 0\)，则方程组有无穷多解。
\end{itemize}

\textbf{证明}：

\textbf{情形一}：\(D = 0\) 且 \(D_x \neq 0\)。假设存在解 \((x_0, y_0)\)，则由定理8.3.1的推导过程（第一步的同解变换），\(D x_0 = D_x\)。但 \(D = 0\) 故 \(0 = D_x \neq 0\)，矛盾。因此无解。\(D_y \neq 0\) 的情况类似。

\textbf{情形二}：\(D = 0\) 且 \(D_x = D_y = 0\)。由 \(D = a_1 b_2 - a_2 b_1 = 0\) 得 \(a_1 b_2 = a_2 b_1\)。

子情形 2a：若 \(b_1 \neq 0\)，由第一个方程 \(x = \frac{c_1 - b_1 y}{a_1}\)（若 \(a_1 \neq 0\)）或 \(y = \frac{c_1}{b_1}\)（若 \(a_1 = 0\)）。代入第二个方程，利用 \(D_x = D_y = 0\) 可验证恒成立，故 \(y\) 可取任意值，\(x\) 由 \(y\) 确定，有无穷多解。

子情形 2b：若 \(b_1 = 0\)，由 \(D = 0\) 得 \(a_1 b_2 = 0\)。若 \(a_1 = 0\)，则第一个方程变为 \(0 = c_1\)，由 \(D_x = c_1 b_2 = 0\) 和 \(D_y = -a_2 c_1 = 0\)，若 \(c_1 \neq 0\) 则 \(b_2 = a_2 = 0\)，但此时方程组退化；若 \(c_1 = 0\) 则第一个方程恒成立。类似讨论第二个方程，可得无穷多解或全体实数解。

综上，\(D = 0\) 时方程组要么无解要么有无穷多解。\(\blacksquare\)

\hrulefill

\subsection{8.4 一元二次方程}\label{ux4e00ux5143ux4e8cux6b21ux65b9ux7a0b}

\textbf{定理 8.4.1（求根公式）} 设 \(a, b, c \in \mathbb{R}\)，\(a \neq 0\)。方程 \(ax^2 + bx + c = 0\) 的解为

\[
x = \frac{-b \pm \sqrt{b^2 - 4ac}}{2a}
\]

当 \(b^2 - 4ac \geq 0\) 时，解为实数。

\textbf{证明}（配方法）：

\[
ax^2 + bx + c = 0
\]

\[
x^2 + \frac{b}{a}x + \frac{c}{a} = 0
\]

\[
x^2 + \frac{b}{a}x = -\frac{c}{a}
\]

\[
x^2 + \frac{b}{a}x + \frac{b^2}{4a^2} = \frac{b^2}{4a^2} - \frac{c}{a} = \frac{b^2 - 4ac}{4a^2}
\]

\[
\left(x + \frac{b}{2a}\right)^2 = \frac{b^2 - 4ac}{4a^2}
\]

当 \(b^2 - 4ac \geq 0\) 时：

\[
x + \frac{b}{2a} = \pm \frac{\sqrt{b^2 - 4ac}}{2a}
\]

\[
x = \frac{-b \pm \sqrt{b^2 - 4ac}}{2a}
\]

\(\blacksquare\)

\textbf{定义 8.4.1（判别式）}

\(\Delta = b^2 - 4ac\) 称为一元二次方程 \(ax^2 + bx + c = 0\) 的\textbf{判别式}。

\textbf{定理 8.4.2（判别式的三种情形）}

设 \(a \neq 0\)，\(f(x) = ax^2 + bx + c\)：

\textbf{(i)} 若 \(\Delta > 0\)，则 \(f(x) = 0\) 有两个不相等的实根：

\[
x_1 = \frac{-b + \sqrt{\Delta}}{2a}, \quad x_2 = \frac{-b - \sqrt{\Delta}}{2a}
\]

\textbf{(ii)} 若 \(\Delta = 0\)，则 \(f(x) = 0\) 有一个二重实根 \(x_0 = -b/(2a)\)。

\textbf{(iii)} 若 \(\Delta < 0\)，则 \(f(x) = 0\) 在 \(\mathbb{R}\) 中无解。

\textbf{证明}：由求根公式的推导过程，\(\left(x + \frac{b}{2a}\right)^2 = \frac{\Delta}{4a^2}\)。

\begin{itemize}
\tightlist
\item
  \(\Delta > 0\) 时，右端为正，\(x + \frac{b}{2a} = \pm \frac{\sqrt{\Delta}}{2a}\) 给出两个不同实数。
\item
  \(\Delta = 0\) 时，\(x + \frac{b}{2a} = 0\)，唯一解 \(x = -b/(2a)\)。
\item
  \(\Delta < 0\) 时，右端为负，而左端 \(\geq 0\)，矛盾。无实数解。
\end{itemize}

\(\blacksquare\)

\textbf{定理 8.4.3（一元二次方程的韦达定理）} 设 \(ax^2 + bx + c = 0\)（\(a \neq 0\)）的两根为 \(x_1, x_2\)，则

\[
x_1 + x_2 = -\frac{b}{a}, \quad x_1 x_2 = \frac{c}{a}
\]

\textbf{证明}：由求根公式：

\[
x_1 + x_2 = \frac{-b + \sqrt{\Delta}}{2a} + \frac{-b - \sqrt{\Delta}}{2a} = \frac{-2b}{2a} = -\frac{b}{a}
\]

\[
x_1 x_2 = \frac{(-b + \sqrt{\Delta})(-b - \sqrt{\Delta})}{4a^2} = \frac{b^2 - \Delta}{4a^2} = \frac{b^2 - (b^2 - 4ac)}{4a^2} = \frac{4ac}{4a^2} = \frac{c}{a}
\]

\(\blacksquare\)

\hrulefill

\subsection{8.5 分式方程与无理方程}\label{ux5206ux5f0fux65b9ux7a0bux4e0eux65e0ux7406ux65b9ux7a0b}

\textbf{定义 8.5.1（分式方程）}

分母中含有未知数的方程称为\textbf{分式方程}。

\textbf{定义 8.5.2（增根）}

在方程变形过程中产生的、不属于原方程解集的根称为\textbf{增根}。

\textbf{定理 8.5.1（增根的产生条件）} 将分式方程 \(\dfrac{P(x)}{Q(x)} = \dfrac{R(x)}{S(x)}\) 两边乘以 \(Q(x) S(x)\) 去分母，所得整式方程的根中，使 \(Q(x) S(x) = 0\) 的值为增根。

\textbf{证明}：设 \(x_0\) 使得 \(Q(x_0) = 0\)。则原方程在 \(x_0\) 处无定义。但乘以 \(Q(x_0) S(x_0) = 0\) 后等式 \(0 = 0\) 恒成立，故 \(x_0\) 可能满足去分母后的方程，却不满足原方程。 \(\blacksquare\)

\textbf{定义 8.5.3（无理方程）}

根号下含有未知数的方程称为\textbf{无理方程}。

\textbf{定理 8.5.2} 两边平方可能产生增根。设 \(A(x) = B(x)\) 为无理方程，两边平方得 \(A^2(x) = B^2(x)\)，则 \(A(x) = B(x)\) 或 \(A(x) = -B(x)\)。因此平方后方程的解集包含原方程的解集，可能更大。

\textbf{证明}：\(A^2 = B^2 \Leftrightarrow A^2 - B^2 = 0 \Leftrightarrow (A+B)(A-B) = 0 \Leftrightarrow A = B\) 或 \(A = -B\)。 \(\blacksquare\)

\hrulefill

\subsection{8.6 高次方程简介}\label{ux9ad8ux6b21ux65b9ux7a0bux7b80ux4ecb}

\textbf{定理 8.6.1（代数基本定理）} 任何 \(n \geq 1\) 次复系数多项式在 \(\mathbb{C}\) 中恰有 \(n\) 个根（计重数）。

\textbf{证明}（Liouville 定理法）：设 \(P(z) = a_n z^n + \cdots + a_0\)（\(a_n \neq 0\)）为无零点的复系数多项式。则 \(f(z) = 1/P(z)\) 为整函数。当 \(|z| \to \infty\) 时，\(|P(z)| \sim |a_n||z|^n \to \infty\)，故 \(|f(z)| \to 0\)，\(f\) 有界。由 Liouville 定理，\(f\) 为常值函数，矛盾。故 \(P(z)\) 至少有一个零点 \(z_1\)。由因式定理，\(P(z) = (z - z_1) Q(z)\)，其中 \(Q(z)\) 为 \(n-1\) 次多项式。对 \(Q(z)\) 重复此论证，归纳得 \(P(z) = a_n \prod_{i=1}^n (z - z_i)\)，恰有 \(n\) 个根（计重数）。\(\blacksquare\)

\textbf{定理 8.6.2（阿贝尔-鲁菲尼定理）} 当 \(n \geq 5\) 时，一般 \(n\) 次多项式方程

\[
a_n x^n + a_{n-1} x^{n-1} + \cdots + a_0 = 0
\]

不存在用系数的加、减、乘、除和开方运算构成的一般求根公式。

\textbf{证明}（Galois 理论框架）：将 \(n\) 次一般方程的 Galois 群等同为对称群 \(S_n\)。根式可解当且仅当 Galois 群为可解群。当 \(n \geq 5\) 时，\(S_n\) 的换位子群列为 \(S_n \triangleright A_n \triangleright A_n \triangleright \cdots\)，\(A_n\)（\(n \geq 5\)）为单群，非可解，故 \(S_n\) 非可解。因此一般 \(n \geq 5\) 次方程不能用根式求解。具体地，方程的根在系数域上的 Galois 群可嵌入 \(S_n\)，而根式扩张对应的 Galois 群必为可解群。若存在根式解，则 \(S_n\) 可解，矛盾。\(\blacksquare\)

\begin{quote}
\textbf{本章展望}：方程是代数的核心应用。下一章将研究不等式，它是方程理论的自然推广，也是优化问题的基础。
\end{quote}

\hrulefill

\hrulefill

\hrulefill

\hrulefill

\hrulefill

\hrulefill

\section{第29章：不等式}\label{ux7b2c29ux7ae0ux4e0dux7b49ux5f0f}

\begin{quote}
\textbf{前置知识}：本章需要第4章实数的序公理，以及第12章的运算律和第13章的多项式理论。
\end{quote}

不等式是数学中研究大小关系的分支，它补充了等式所不能表达的''比较''维度。本章从实数序公理出发，建立不等式的基本性质（传递性、加法保序性、乘法保序性等），并讨论不等式变形中的关键注意事项------特别是乘以负数时不等号方向的改变。一元一次不等式和一元二次不等式的解法与方程类似，但关键区别在于解集通常是区间而非离散点，且判别式的符号决定了二次不等式的解集结构（两根之间、两根之外、或恒成立/恒不成立）。本章还将介绍几个重要的经典不等式------均值不等式（\(\frac{a+b}{2} \geq \sqrt{ab}\)）、Cauchy-Schwarz 不等式和三角不等式------它们在数学分析、优化理论和概率论中有广泛应用。不等式的证明方法与等式不同，往往需要巧妙的放缩技巧和代数变形。理解不等式不仅是掌握解题技巧，更是理解实数序结构和函数单调性的重要途径。

\subsection{9.1 基本性质}\label{ux57faux672cux6027ux8d28}

以下从实数的序公理（参见第一卷第4章）出发。

\textbf{定理 9.1.1（传递性）} 若 \(a > b\) 且 \(b > c\)，则 \(a > c\)。

\textbf{证明}：由序公理，\(a > b\) 意味着 \(a - b > 0\)，\(b > c\) 意味着 \(b - c > 0\)。正数之和仍为正（由序公理），故

\[
a - c = (a - b) + (b - c) > 0
\]

即 \(a > c\)。 \(\blacksquare\)

\textbf{定理 9.1.2（加法保序性）} 若 \(a > b\)，则对任意 \(c \in \mathbb{R}\)，有 \(a + c > b + c\)。

\textbf{证明}：\(a > b \Rightarrow a - b > 0\)。而 \((a + c) - (b + c) = a - b > 0\)，故 \(a + c > b + c\)。 \(\blacksquare\)

\textbf{定理 9.1.3（乘法保序性------正数乘）} 若 \(a > b\) 且 \(c > 0\)，则 \(ac > bc\)。

\textbf{证明}：\(a > b \Rightarrow a - b > 0\)。\(c > 0\)，正数之积为正，故 \(c(a - b) > 0\)，即 \(ac - bc > 0\)，\(ac > bc\)。 \(\blacksquare\)

\textbf{定理 9.1.4（乘法反序性------负数乘）} 若 \(a > b\) 且 \(c < 0\)，则 \(ac < bc\)。

\textbf{证明}：\(a > b \Rightarrow a - b > 0\)。\(c < 0 \Rightarrow -c > 0\)。正数之积为正，故 \((-c)(a - b) > 0\)，即 \(-(ac - bc) > 0\)，\(ac - bc < 0\)，\(ac < bc\)。 \(\blacksquare\)

\textbf{定理 9.1.5（正数的平方为正）} 若 \(a \neq 0\)，则 \(a^2 > 0\)。

\textbf{证明}：若 \(a > 0\)，由正数之积为正，\(a^2 = a \cdot a > 0\)。若 \(a < 0\)，则 \(-a > 0\)，\(a^2 = (-a)^2 = (-a)(-a) > 0\)。 \(\blacksquare\)

\textbf{定理 9.1.6} \(1 > 0\)。

\textbf{证明}：\(1 = 1^2 > 0\)（由定理9.1.5，\(1 \neq 0\)）。 \(\blacksquare\)

\textbf{定理 9.1.7（三歧性）} 对任意 \(a, b \in \mathbb{R}\)，以下三种情形恰有一种成立：\(a > b\)，\(a = b\)，\(a < b\)。

\textbf{证明}：由实数序公理的三歧性公理直接得出。 \(\blacksquare\)

\textbf{定理 9.1.8（倒数保序性）} 若 \(a > b > 0\)，则 \(\dfrac{1}{a} < \dfrac{1}{b}\)。

\textbf{证明}：\(a > b > 0\)，则 \(ab > 0\)。由定理9.1.3，两边乘以 \(\dfrac{1}{ab} > 0\)：

\[
\frac{a}{ab} > \frac{b}{ab} \Rightarrow \frac{1}{b} > \frac{1}{a}
\]

\(\blacksquare\)

\hrulefill

\subsection{9.2 一元一次不等式组}\label{ux4e00ux5143ux4e00ux6b21ux4e0dux7b49ux5f0fux7ec4}

\textbf{定理 9.2.1} 设 \(a \neq 0\)，则 \(ax + b > 0\) 的解集为：

\begin{itemize}
\tightlist
\item
  若 \(a > 0\)：\(x > -b/a\)；
\item
  若 \(a < 0\)：\(x < -b/a\)。
\end{itemize}

\textbf{证明}：\(ax > -b\)。

\begin{itemize}
\tightlist
\item
  若 \(a > 0\)：由乘法保序性，\(x > -b/a\)。
\item
  若 \(a < 0\)：由乘法反序性，\(x < -b/a\)。
\end{itemize}

\(\blacksquare\)

\textbf{定理 9.2.2（不等式组的解集）} 不等式组

\[
\begin{cases} a_1 x + b_1 > 0 \\ a_2 x + b_2 > 0 \end{cases}
\]

的解集为两个不等式解集的交集。

\textbf{证明}：\(x\) 满足不等式组当且仅当 \(x\) 同时满足每个不等式，这恰好是两个解集的交集的定义。 \(\blacksquare\)

\hrulefill

\subsection{9.3 一元二次不等式}\label{ux4e00ux5143ux4e8cux6b21ux4e0dux7b49ux5f0f}

\textbf{定理 9.3.1} 设 \(a > 0\)，\(f(x) = ax^2 + bx + c\)，\(\Delta = b^2 - 4ac\)。

\textbf{(i) \(\Delta > 0\)}：设两根 \(x_1 < x_2\)。则

\begin{itemize}
\tightlist
\item
  \(f(x) > 0\) 的解集为 \((-\infty, x_1) \cup (x_2, +\infty)\)。
\item
  \(f(x) < 0\) 的解集为 \((x_1, x_2)\)。
\end{itemize}

\textbf{(ii) \(\Delta = 0\)}：设根为 \(x_0\)。则

\begin{itemize}
\tightlist
\item
  \(f(x) > 0\) 的解集为 \(\mathbb{R} \setminus \{x_0\}\)。
\item
  \(f(x) < 0\) 的解集为 \(\varnothing\)。
\end{itemize}

\textbf{(iii) \(\Delta < 0\)}：则

\begin{itemize}
\tightlist
\item
  \(f(x) > 0\) 的解集为 \(\mathbb{R}\)。
\item
  \(f(x) < 0\) 的解集为 \(\varnothing\)。
\end{itemize}

\textbf{证明}：由配方法，\(f(x) = a\left(x + \dfrac{b}{2a}\right)^2 + \dfrac{4ac - b^2}{4a} = a\left(x + \dfrac{b}{2a}\right)^2 - \dfrac{\Delta}{4a}\)。

\textbf{(i)} \(\Delta > 0\)：\(f(x) = a(x - x_1)(x - x_2)\)。

\begin{itemize}
\tightlist
\item
  \(x < x_1\) 时：\(x - x_1 < 0\)，\(x - x_2 < 0\)，\((x-x_1)(x-x_2) > 0\)，\(f(x) > 0\)。
\item
  \(x_1 < x < x_2\) 时：\(x - x_1 > 0\)，\(x - x_2 < 0\)，\((x-x_1)(x-x_2) < 0\)，\(f(x) < 0\)。
\item
  \(x > x_2\) 时：\(x - x_1 > 0\)，\(x - x_2 > 0\)，\((x-x_1)(x-x_2) > 0\)，\(f(x) > 0\)。
\item
  \(x = x_1\) 或 \(x = x_2\) 时：\(f(x) = 0\)。
\end{itemize}

\textbf{(ii)} \(\Delta = 0\)：\(f(x) = a(x - x_0)^2 \geq 0\)。等号当且仅当 \(x = x_0\)。

\textbf{(iii)} \(\Delta < 0\)：\(f(x) = a\left(x + \dfrac{b}{2a}\right)^2 - \dfrac{\Delta}{4a}\)。其中 \(-\Delta/(4a) > 0\)（因 \(\Delta < 0\)，\(a > 0\)），而第一项 \(\geq 0\)，故 \(f(x) > 0\) 对一切 \(x\) 成立。

\(\blacksquare\)

\textbf{推论 9.3.1} 若 \(a < 0\)，不等式 \(ax^2 + bx + c > 0\) 等价于 \(-ax^2 - bx - c < 0\)，可转化为 \(a' > 0\) 的情形。

\textbf{证明}：\(ax^2 + bx + c > 0\) 两边乘 \(-1\)（不等号反向）得 \(-ax^2 - bx - c < 0\)，其中 \(-a > 0\)。 \(\blacksquare\)

\hrulefill

\subsection{9.4 绝对值不等式}\label{ux7eddux5bf9ux503cux4e0dux7b49ux5f0f}

\textbf{定义 9.4.1（绝对值）}

对任意 \(a \in \mathbb{R}\)，定义

\[
|a| = \begin{cases} a & \text{若 } a \geq 0 \\ -a & \text{若 } a < 0 \end{cases}
\]

\textbf{定理 9.4.1（绝对值的基本性质）}

\begin{enumerate}
\def\labelenumi{(\roman{enumi})}
\item
  \(|a| \geq 0\)，等号当且仅当 \(a = 0\)。
\item
  \(|a| = \sqrt{a^2}\)。
\item
  \(|ab| = |a| \cdot |b|\)。
\item
  \(\left|\dfrac{a}{b}\right| = \dfrac{|a|}{|b|}\)（\(b \neq 0\)）。
\end{enumerate}

\textbf{证明}：

\begin{enumerate}
\def\labelenumi{(\roman{enumi})}
\item
  若 \(a \geq 0\)，\(|a| = a \geq 0\)；若 \(a < 0\)，\(|a| = -a > 0\)。\(|a| = 0 \Leftrightarrow a = 0\) 由定义直接得出。
\item
  若 \(a \geq 0\)，\(\sqrt{a^2} = a = |a|\)；若 \(a < 0\)，\(\sqrt{a^2} = \sqrt{(-a)^2} = -a = |a|\)。
\item
  分四种子情形讨论：
\end{enumerate}

\begin{itemize}
\tightlist
\item
  \(a \geq 0, b \geq 0\)：\(ab \geq 0\)，\(|ab| = ab = |a| \cdot |b|\)。
\item
  \(a \geq 0, b < 0\)：\(ab \leq 0\)，\(|ab| = -ab = a \cdot (-b) = |a| \cdot |b|\)。
\item
  \(a < 0, b \geq 0\)：\(ab \leq 0\)，\(|ab| = -ab = (-a) \cdot b = |a| \cdot |b|\)。
\item
  \(a < 0, b < 0\)：\(ab > 0\)，\(|ab| = ab = (-a)(-b) = |a| \cdot |b|\)。
\end{itemize}

所有情形下 \(|ab| = |a| \cdot |b|\) 成立。

\begin{enumerate}
\def\labelenumi{(\roman{enumi})}
\setcounter{enumi}{3}
\tightlist
\item
  由 (iii)，\(|b| \cdot |a/b| = |b \cdot a/b| = |a|\)，故 \(|a/b| = |a|/|b|\)。
\end{enumerate}

\(\blacksquare\)

\textbf{定理 9.4.2（三角不等式）} 对任意 \(a, b \in \mathbb{R}\)：

\[
|a + b| \leq |a| + |b|
\]

\textbf{证明}：两边平方（因两边非负，不等式方向不变）：

\[
(a + b)^2 \leq (|a| + |b|)^2
\]

\[
a^2 + 2ab + b^2 \leq |a|^2 + 2|a||b| + |b|^2 = a^2 + 2|ab| + b^2
\]

只需证 \(ab \leq |ab|\)，即 \(ab \leq |ab|\)。

若 \(ab \geq 0\)，则 \(ab = |ab|\)，成立。若 \(ab < 0\)，则 \(ab < 0 \leq |ab|\)，成立。

\(\blacksquare\)

\textbf{定理 9.4.3（反向三角不等式）} 对任意 \(a, b \in \mathbb{R}\)：

\[
\bigl||a| - |b|\bigr| \leq |a - b|
\]

\textbf{证明}：由三角不等式，\(|a| = |(a - b) + b| \leq |a - b| + |b|\)，故 \(|a| - |b| \leq |a - b|\)。

类似地，\(|b| = |(b - a) + a| \leq |b - a| + |a| = |a - b| + |a|\)，故 \(|b| - |a| \leq |a - b|\)，即 \(-(|a| - |b|) \leq |a - b|\)。

综合两式：\(-|a - b| \leq |a| - |b| \leq |a - b|\)，即 \(\bigl||a| - |b|\bigr| \leq |a - b|\)。 \(\blacksquare\)

\textbf{推论 9.4.1} 对任意 \(a, b \in \mathbb{R}\)：

\[
\bigl||a| - |b|\bigr| \leq |a + b| \leq |a| + |b|
\]

\textbf{证明}：左边不等式由 \(|a| = |(a+b) - b| \leq |a+b| + |b|\)，得 \(|a| - |b| \leq |a+b|\)。类似 \(|b| - |a| \leq |a+b|\)。右边不等式即三角不等式。 \(\blacksquare\)

\hrulefill

\subsection{9.5 均值不等式}\label{ux5747ux503cux4e0dux7b49ux5f0f}

\textbf{定理 9.5.1（二元算术-几何均值不等式，AM-GM）} 对任意 \(a, b \geq 0\)：

\[
\frac{a + b}{2} \geq \sqrt{ab}
\]

等号当且仅当 \(a = b\)。

\textbf{证明}（配方法）：

\[
\frac{a + b}{2} - \sqrt{ab} = \frac{a + b - 2\sqrt{ab}}{2} = \frac{(\sqrt{a})^2 - 2\sqrt{a}\sqrt{b} + (\sqrt{b})^2}{2} = \frac{(\sqrt{a} - \sqrt{b})^2}{2} \geq 0
\]

等号当且仅当 \(\sqrt{a} = \sqrt{b}\)，即 \(a = b\)。 \(\blacksquare\)

\textbf{定理 9.5.2（\(n\) 元 AM-GM 不等式）} 对任意 \(a_1, a_2, \ldots, a_n \geq 0\)（\(n \geq 2\)）：

\[
\frac{a_1 + a_2 + \cdots + a_n}{n} \geq \sqrt[n]{a_1 a_2 \cdots a_n}
\]

等号当且仅当 \(a_1 = a_2 = \cdots = a_n\)。

\textbf{证明}（前向-后向归纳法，柯西归纳法）：

\textbf{第一步（\(n = 2\) 的情形）}：由定理9.5.1。

\textbf{第二步（\(n \to 2n\)）}：设 \(n\) 元不等式成立，证 \(2n\) 元。

\[
\frac{a_1 + \cdots + a_{2n}}{2n} = \frac{1}{2}\left(\frac{a_1 + \cdots + a_n}{n} + \frac{a_{n+1} + \cdots + a_{2n}}{n}\right)
\]

由归纳假设和二元AM-GM：

\[
\geq \frac{1}{2}\left(\sqrt[n]{a_1 \cdots a_n} + \sqrt[n]{a_{n+1} \cdots a_{2n}}\right) \geq \sqrt{\sqrt[n]{a_1 \cdots a_n} \cdot \sqrt[n]{a_{n+1} \cdots a_{2n}}} = \sqrt[2n]{a_1 \cdots a_{2n}}
\]

\textbf{第三步（\(n \to n-1\)，向后归纳）}：设 \(n\) 元不等式成立，证 \(n-1\) 元。

给定 \(a_1, \ldots, a_{n-1} \geq 0\)，令

\[
a_n = \frac{a_1 + a_2 + \cdots + a_{n-1}}{n-1} = A
\]

由 \(n\) 元AM-GM：

\[
\frac{a_1 + \cdots + a_{n-1} + A}{n} \geq \sqrt[n]{a_1 \cdots a_{n-1} \cdot A}
\]

左端 \(= \frac{(n-1)A + A}{n} = A\)。故

\[
A \geq \sqrt[n]{a_1 \cdots a_{n-1} \cdot A}
\]

\[
A^n \geq a_1 \cdots a_{n-1} \cdot A
\]

若 \(A = 0\)，则所有 \(a_i = 0\)，\(n-1\) 元不等式 \(0 \geq 0\) 平凡成立。若 \(A > 0\)，两边除以 \(A\)：

\[
A \geq \sqrt[n-1]{a_1 \cdots a_{n-1}}
\]

即 \(\dfrac{a_1 + \cdots + a_{n-1}}{n-1} \geq \sqrt[n-1]{a_1 \cdots a_{n-1}}\)。

综合第二步和第三步，对所有 \(n \geq 2\) 不等式成立。等号条件由各步等号传递得出：第二步中 \(2n\) 元等号当且仅当 \(a_1 = \cdots = a_{2n}\)（由二元 AM-GM 及归纳假设），第三步中 \(n-1\) 元等号当且仅当 \(a_1 = \cdots = a_{n-1} = A = \frac{a_1+\cdots+a_{n-1}}{n-1}\)，即全部相等。因此对所有 \(n \geq 2\)，等号当且仅当 \(a_1 = a_2 = \cdots = a_n\)。 \(\blacksquare\)

\textbf{定义 9.5.1（均值的定义）}

设 \(a_1, a_2, \ldots, a_n > 0\)，定义以下均值：

\begin{itemize}
\tightlist
\item
  \textbf{算术平均}（AM）：\(A = \dfrac{a_1 + a_2 + \cdots + a_n}{n}\)
\item
  \textbf{几何平均}（GM）：\(G = \sqrt[n]{a_1 a_2 \cdots a_n}\)
\item
  \textbf{调和平均}（HM）：\(H = \dfrac{n}{\dfrac{1}{a_1} + \dfrac{1}{a_2} + \cdots + \dfrac{1}{a_n}}\)
\item
  \textbf{平方平均}（QM）：\(Q = \sqrt{\dfrac{a_1^2 + a_2^2 + \cdots + a_n^2}{n}}\)
\end{itemize}

\textbf{定理 9.5.3（均值不等式链）} 对任意 \(a_1, a_2, \ldots, a_n > 0\)：

\[
H \leq G \leq A \leq Q
\]

等号当且仅当 \(a_1 = a_2 = \cdots = a_n\)。

\textbf{证明}：

\textbf{\(G \leq A\)}：即定理9.5.2。

\textbf{\(H \leq G\)}：对 \(1/a_1, 1/a_2, \ldots, 1/a_n\) 应用 \(G \leq A\)：

\[
\sqrt[n]{\frac{1}{a_1} \cdot \frac{1}{a_2} \cdots \frac{1}{a_n}} \leq \frac{\frac{1}{a_1} + \frac{1}{a_2} + \cdots + \frac{1}{a_n}}{n}
\]

取倒数（函数 \(f(x) = 1/x\) 在 \(x > 0\) 上单调递减）：

\[
\frac{n}{\frac{1}{a_1} + \cdots + \frac{1}{a_n}} \leq \sqrt[n]{a_1 a_2 \cdots a_n}
\]

即 \(H \leq G\)。

\textbf{\(A \leq Q\)}：需证 \(\left(\dfrac{\sum a_i}{n}\right)^2 \leq \dfrac{\sum a_i^2}{n}\)，即

\[
\left(\sum a_i\right)^2 \leq n \sum a_i^2
\]

将左端展开：

\[
\left(\sum_{i=1}^n a_i\right)^2 = \sum_{i=1}^n a_i^2 + 2\sum_{1 \leq i < j \leq n} a_i a_j
\]

因此需证 \(2\sum_{i < j} a_i a_j \leq (n-1)\sum a_i^2\)。由基本不等式 \(2a_i a_j \leq a_i^2 + a_j^2\)（这等价于 \((a_i - a_j)^2 \geq 0\)），对每对 \((i, j)\) 求和：

\[
2\sum_{i < j} a_i a_j \leq \sum_{i < j} (a_i^2 + a_j^2) = (n-1)\sum_{i=1}^n a_i^2
\]

（每个 \(a_i^2\) 恰好出现在 \(n-1\) 对中。）代入即得 \(\left(\sum a_i\right)^2 \leq n \sum a_i^2\)。\(\blacksquare\)

\textbf{注记}：上述不等式也是柯西-施瓦茨不等式（定理 9.6.1）取 \(b_i = 1\) 的特例。

\hrulefill

\subsection{9.6 柯西-施瓦茨不等式}\label{ux67efux897f-ux65bdux74e6ux8328ux4e0dux7b49ux5f0f}

\textbf{定理 9.6.1（柯西-施瓦茨不等式）} 对任意实数 \(a_1, \ldots, a_n, b_1, \ldots, b_n\)：

\[
\left(\sum_{i=1}^n a_i b_i\right)^2 \leq \left(\sum_{i=1}^n a_i^2\right)\left(\sum_{i=1}^n b_i^2\right)
\]

等号当且仅当存在常数 \(t\)，使得 \(a_i = t b_i\)（\(\forall i\)）或 \(b_i = 0\)（\(\forall i\)）。

\textbf{证明}（判别式法）：

考虑关于 \(t\) 的函数：

\[
f(t) = \sum_{i=1}^n (a_i t - b_i)^2 \geq 0 \quad (\forall t \in \mathbb{R})
\]

展开：

\[
f(t) = \left(\sum a_i^2\right) t^2 - 2\left(\sum a_i b_i\right) t + \left(\sum b_i^2\right) \geq 0
\]

设 \(A = \sum a_i^2\)，\(B = \sum a_i b_i\)，\(C = \sum b_i^2\)。则 \(f(t) = At^2 - 2Bt + C \geq 0\) 对一切 \(t\) 成立。

若 \(A = 0\)，则所有 \(a_i = 0\)，不等式两端均为零，等号成立。

若 \(A > 0\)，二次函数 \(f(t) \geq 0\) 对一切 \(t\) 成立，当且仅当其判别式 \(\leq 0\)：

\[
\Delta = 4B^2 - 4AC \leq 0
\]

\[
B^2 \leq AC
\]

即

\[
\left(\sum a_i b_i\right)^2 \leq \left(\sum a_i^2\right)\left(\sum b_i^2\right)
\]

等号当且仅当 \(\Delta = 0\)，即存在 \(t_0\) 使得 \(f(t_0) = 0\)，即 \(\sum(a_i t_0 - b_i)^2 = 0\)，即 \(a_i t_0 = b_i\) 对所有 \(i\) 成立。 \(\blacksquare\)

\textbf{推论 9.6.1（施瓦茨不等式的积分形式，陈述）} 设 \(f, g\) 在 \([a, b]\) 上可积，则

\[
\left(\int_a^b f(x) g(x) \, dx\right)^2 \leq \left(\int_a^b f^2(x) \, dx\right)\left(\int_a^b g^2(x) \, dx\right)
\]

\begin{quote}
\textbf{本章展望}：不等式是分析学和优化理论的基石。第16章将从逻辑学的角度审视数学命题的结构，为严谨的数学推理提供更深层的理论支撑。
\end{quote}

\hrulefill

\hrulefill

\hrulefill

\hrulefill

\hrulefill

\newpage

\chapter{01-基础/05-初等代数/04-逻辑与线性规划.md}\label{ux57faux784005-ux521dux7b49ux4ee3ux657004-ux903bux8f91ux4e0eux7ebfux6027ux89c4ux5212.md}

\section{第30章：逻辑}\label{ux7b2c30ux7ae0ux903bux8f91}

逻辑是数学推理的基石，研究有效的推理形式和论证结构。命题逻辑处理命题（或真或假的陈述）之间的逻辑关系，而谓词逻辑引入量词以表达更丰富的数学陈述。逻辑学的发展经历了从古希腊（Aristotle 的三段论）到 19 世纪（Boole 的代数逻辑、Frege 的谓词逻辑）再到 20 世纪（Godel 的完备性定理和不完备性定理）的漫长历程，最终成为现代数学中与集合论并列的基础学科。逻辑学与计算机科学（程序验证、类型论、自动推理）、人工智能（知识表示、自动定理证明）和语言学（形式语义学）有深刻联系。本章系统介绍命题逻辑的形式系统、真值表与语义、逻辑等价与范式、谓词逻辑的量词理论和推理规则，以及布尔代数作为逻辑的代数化。线性规划作为运筹学中的核心优化方法，在本章中得到系统处理，包括标准形式、单纯形法和对偶理论。

\subsection{4.1 命题逻辑}\label{ux547dux9898ux903bux8f91-1}

\textbf{定义 4.1}（命题）：\textbf{命题}是一个可以判断真假的陈述句。命题的取值只有两种：真（T，1）或假（F，0）。命题逻辑研究命题通过逻辑联结词组合而成的复合命题。命题逻辑的核心假设是\textbf{二值原理}（principle of bivalence）：每个命题恰好具有一个真值------真或假。这一假设排除了多值逻辑（如模糊逻辑、三值逻辑）和直觉主义逻辑（不接受排中律）中的可能性。

命题作为逻辑研究的基本对象，本身并不复杂------它仅仅是一个可以判定真假的陈述。然而，逻辑的力量在于用少数几个联结词将简单命题组合成任意复杂的复合命题，从而表达丰富多样的推理关系。以下五种基本逻辑联结词构成了命题逻辑的全部''词汇''。

\textbf{定义 4.2}（逻辑联结词）：基本逻辑联结词包括： - 否定 \(\neg p\)：非 \(p\)（not） - 合取 \(p \land q\)：\(p\) 且 \(q\)（and） - 析取 \(p \lor q\)：\(p\) 或 \(q\)（or，包含性或） - 蕴涵 \(p \to q\)：若 \(p\) 则 \(q\)（implies） - 双条件 \(p \leftrightarrow q\)：\(p\) 当且仅当 \(q\)（iff）

定义了联结词之后，下一步自然要问：如何精确地确定复合命题的真值？真值表语义回答了这个问题------每个联结词由其真值表唯一刻画，而真值表穷举了所有可能的真值指派组合。

\textbf{定义 4.3}（真值表语义）：每个 \(n\) 元联结词由其真值表唯一确定。真值表是 \(2^n\) 行的表格，列出所有可能的真值指派下复合命题的真值。以下为基本联结词的真值表语义：

\begin{longtable}[]{@{}
  >{\raggedright\arraybackslash}p{(\linewidth - 12\tabcolsep) * \real{0.0588}}
  >{\raggedright\arraybackslash}p{(\linewidth - 12\tabcolsep) * \real{0.0588}}
  >{\raggedright\arraybackslash}p{(\linewidth - 12\tabcolsep) * \real{0.1294}}
  >{\raggedright\arraybackslash}p{(\linewidth - 12\tabcolsep) * \real{0.1647}}
  >{\raggedright\arraybackslash}p{(\linewidth - 12\tabcolsep) * \real{0.1529}}
  >{\raggedright\arraybackslash}p{(\linewidth - 12\tabcolsep) * \real{0.1529}}
  >{\raggedright\arraybackslash}p{(\linewidth - 12\tabcolsep) * \real{0.2824}}@{}}
\toprule
\begin{minipage}[b]{\linewidth}\raggedright
\(p\)
\end{minipage} & \begin{minipage}[b]{\linewidth}\raggedright
\(q\)
\end{minipage} & \begin{minipage}[b]{\linewidth}\raggedright
\(\neg p\)
\end{minipage} & \begin{minipage}[b]{\linewidth}\raggedright
\(p \land q\)
\end{minipage} & \begin{minipage}[b]{\linewidth}\raggedright
\(p \lor q\)
\end{minipage} & \begin{minipage}[b]{\linewidth}\raggedright
\(p \to q\)
\end{minipage} & \begin{minipage}[b]{\linewidth}\raggedright
\(p \leftrightarrow q\)
\end{minipage} \\
\midrule
\endhead
\bottomrule
\endlastfoot
T & T & F & T & T & T & T \\
T & F & F & F & T & F & F \\
F & T & T & F & T & T & F \\
F & F & T & F & F & T & T \\
\end{longtable}

蕴涵 \(p \to q\) 的真值表定义（\(p\) 假时 \(p \to q\) 为真）常引起困惑，但这是外延逻辑（extensional logic）的必然结果：\(p \to q\) 定义为 \(\neg p \lor q\)，其唯一为假的情形是 \(p\) 真而 \(q\) 假。这种''空真''（vacuous truth）在数学推理中普遍存在：如''若 \(x > 5\) 则 \(x > 3\)``对 \(x = 2\) 而言为真，因为前提不成立。

\textbf{定义 4.4}（逻辑等价）：两个命题公式 \(A\) 和 \(B\) \textbf{逻辑等价}（记作 \(A \equiv B\)）如果它们在所有真值指派下具有相同的真值。等价地，\(A \leftrightarrow B\) 是重言式（永真式）。

逻辑等价的定义自然引出了一个问题：哪些基本的等价关系可以作为推理的基础？以下定理列举了逻辑推理中最核心的等价律，它们在化简逻辑公式和构造证明中反复出现。

\textbf{定理 4.1}（基本逻辑等价律）： 1. 双重否定：\(\neg(\neg p) \equiv p\) 2. 交换律：\(p \land q \equiv q \land p\)；\(p \lor q \equiv q \lor p\) 3. 结合律：\((p \land q) \land r \equiv p \land (q \land r)\)；\((p \lor q) \lor r \equiv p \lor (q \lor r)\) 4. 分配律：\(p \land (q \lor r) \equiv (p \land q) \lor (p \land r)\)；\(p \lor (q \land r) \equiv (p \lor q) \land (p \lor r)\) 5. De Morgan 律：\(\neg(p \land q) \equiv \neg p \lor \neg q\)；\(\neg(p \lor q) \equiv \neg p \land \neg q\) 6. 吸收律：\(p \land (p \lor q) \equiv p\)；\(p \lor (p \land q) \equiv p\) 7. 条件等价：\(p \to q \equiv \neg p \lor q\) 8. 双条件等价：\(p \leftrightarrow q \equiv (p \to q) \land (q \to p)\)

\textbf{证明}（De Morgan 律）：构造真值表验证。对 \(\neg(p \land q)\) 和 \(\neg p \lor \neg q\)，在所有四种真值指派下，两者的真值完全相同。因此 \(\neg(p \land q) \equiv \neg p \lor \neg q\)。同理可证 \(\neg(p \lor q) \equiv \neg p \land \neg q\)。De Morgan 律是逻辑代数中最基本的对偶关系，它将合取和析取通过否定相互转换。\(\blacksquare\)

掌握了基本等价律和联结词的语义之后，我们可以根据真值特征将命题公式分为三大类：在所有情况下都为真的公式（重言式）、在所有情况下都为假的公式（矛盾式）、以及介于两者之间的可满足式。这三类公式构成了命题逻辑语义分类的基本框架。

\textbf{定义 4.5}（重言式、矛盾式与可满足式）： - \textbf{重言式}（tautology）：在所有真值指派下均为真的公式（如 \(p \lor \neg p\)，排中律） - \textbf{矛盾式}（contradiction）：在所有真值指派下均为假的公式（如 \(p \land \neg p\)） - \textbf{可满足式}（contingent）：在某些真值指派下为真、某些下为假的公式

重言式是命题逻辑中最重要的一类公式，因为有效的推理规则本质上就是重言蕴涵式。命题逻辑的完备性定理确立了重言式与形式可证公式之间的精确对应关系。

\textbf{定理 4.2}（命题逻辑的完备性定理）：命题逻辑中，重言式与定理（可从公理通过推理规则推导出的公式）集合完全一致。即 \(\models A\)（语义蕴涵，\(A\) 为重言式）当且仅当 \(\vdash A\)（语法推导，\(A\) 为定理）。命题逻辑既是可靠的（sound, \(\vdash A \Rightarrow\ \models A\)）又是完备的（complete, \(\models A \Rightarrow\ \vdash A\)）。

\textbf{证明（框架）}：可靠性由归纳法验证每个推理规则保持真值。完备性证明：假设 \(A\) 不是定理（\(\nvdash A\)），则 \(\{\neg A\}\) 是无矛盾的（consistent），可扩充为极大无矛盾集 \(\Gamma\)。由 \(\Gamma\) 构造真值指派 \(\nu\)：\(\nu(p) = T\) 当且仅当 \(p \in \Gamma\)。由归纳法，\(\nu(B) = T\) 当且仅当 \(B \in \Gamma\) 对所有公式 \(B\) 成立。特别地，\(\neg A \in \Gamma\)，故 \(\nu(A) = F\)，因此 \(A\) 不是重言式。\(\blacksquare\)

完备性定理的证明依赖于一个具体的公理系统。以下Hilbert型公理化仅使用了三条公理模式和一条推理规则，却足以推导出命题逻辑的全部重言式------这一事实展示了公理化方法的简洁与威力。

\textbf{定义 4.6}（公理系统）：命题逻辑的 Hilbert 型公理化由以下公理模式和推理规则组成： - 公理 1（肯定前件）：\(A \to (B \to A)\) - 公理 2（蕴涵分配）：\((A \to (B \to C)) \to ((A \to B) \to (A \to C))\) - 公理 3（反证法）：\((\neg B \to \neg A) \to (A \to B)\) - 推理规则（Modus Ponens，MP）：从 \(A\) 和 \(A \to B\) 推出 \(B\)

\textbf{定理 4.2b}（演绎定理）：若 \(\Gamma \cup \{A\} \vdash B\)，则 \(\Gamma \vdash A \to B\)。演绎定理是 Hilbert 型系统中最重要的元定理，它将假设下的推导转化为蕴涵公式，极大地简化了形式证明。演绎定理的证明通过对推导长度归纳：基始 \(B\) 为公理或 \(B = A\) 时显然；归纳步利用公理 2 和 MP。

\subsection{4.2 范式与逻辑函数的完备性}\label{ux8303ux5f0fux4e0eux903bux8f91ux51fdux6570ux7684ux5b8cux5907ux6027}

\textbf{定义 4.7}（文字、子句、析取范式与合取范式）： - \textbf{文字}（literal）：命题变元或其否定（\(p\) 或 \(\neg p\)） - \textbf{子句}（clause）：若干文字的析取（如 \(p \lor \neg q \lor r\)） - \textbf{析取范式}（DNF / Disjunctive Normal Form）：若干合取项的析取（如 \((p \land q) \lor (\neg p \land r)\)） - \textbf{合取范式}（CNF / Conjunctive Normal Form）：若干子句的合取（如 \((p \lor q) \land (\neg p \lor r)\)）

范式为命题公式提供了一种标准化的表达形式。无论是析取范式还是合取范式，都使得公式的结构变得清晰可辨。以下定理保证了任何命题公式都存在等价的范式表示------这一结论不仅是理论上的保证，也是数字电路设计中逻辑综合算法的理论基础。

\textbf{定理 4.3}（范式定理）：任何命题公式都等价于一个析取范式和一个合取范式。任何 \(n\) 元布尔函数都可由仅含 \(\land, \lor, \neg\) 的公式表示。

\textbf{证明}：DNF 构造：对真值表中使 \(A\) 为真的每一行，构造一个合取项（如该行 \(p\) 为真则含 \(p\)，\(p\) 为假则含 \(\neg p\)），将所有合取项析取即得 DNF（主析取范式）。CNF 构造：对偶地，对使 \(A\) 为假的每一行构造析取项，将所有析取项合取即得 CNF（主合取范式）。\(\blacksquare\)

范式定理表明 \(\{\neg, \land, \lor\}\) 这三个联结词足以表达任何布尔函数。然而，一个自然的问题是：能否用更少的联结词达到同样的表达能力？联结词完备性定理给出了肯定的回答------甚至单个联结词（如NAND或NOR）就足够了。

\textbf{定理 4.4}（逻辑联结词的完备性）：\(\{\neg, \land\}\)、\(\{\neg, \lor\}\) 和 \(\{\neg, \to\}\) 都是逻辑联结词的完备集（functionally complete），即任何布尔函数都可由这些联结词表示。\(\{\mid\}\)（NAND / Sheffer 竖线）和 \(\{\downarrow\}\)（NOR / Peirce 箭头）单独构成完备集： - \(p \mid q \equiv \neg(p \land q)\)（NAND） - \(p \downarrow q \equiv \neg(p \lor q)\)（NOR）

\textbf{证明}：由范式定理，任何布尔函数可由 \(\{\neg, \land, \lor\}\) 表示。\(\lor\) 可用 \(\neg\) 和 \(\land\) 表示：\(p \lor q \equiv \neg(\neg p \land \neg q)\)（De Morgan 律）。因此 \(\{\neg, \land\}\) 完备。对于 NAND：\(\neg p \equiv p \mid p\)，\(p \land q \equiv \neg(p \mid q) \equiv (p \mid q) \mid (p \mid q)\)。故 \(\{\mid\}\) 完备。NOR 情形类似。\(\blacksquare\)

\textbf{定理 4.4b}（Post 完备性准则，1941）：一个联结词集是完备的当且仅当它不满足 Post 的五种封闭性中的任何一种：(1) 保真性（\(f(T,\ldots,T)=T\)）；(2) 保假性（\(f(F,\ldots,F)=F\)）；(3) 自对偶性（\(f(\neg x_1,\ldots,\neg x_n) = \neg f(x_1,\ldots,x_n)\)）；(4) 单调性；(5) 线性性。Post 准则完全刻画了所有可能的完备联结词集。

\subsection{4.3 谓词逻辑与量词}\label{ux8c13ux8bcdux903bux8f91ux4e0eux91cfux8bcd}

命题逻辑处理的是完整的陈述句，但数学中的许多推理涉及''对所有\(x\)``和''存在\(x\)``这样的量化表达。谓词逻辑通过引入谓词和量词，将命题逻辑的表达能力扩展到足以形式化大部分数学推理的程度。

\textbf{定义 4.8}（谓词与量词）：\textbf{谓词} \(P(x)\) 是依赖于变元 \(x\) 的陈述（取值为真或假）。\textbf{全称量词} \(\forall x P(x)\) 表示''对所有 \(x\)，\(P(x)\) 为真''。\textbf{存在量词} \(\exists x P(x)\) 表示''存在 \(x\) 使得 \(P(x)\) 为真''。量词的作用范围是\textbf{论域}（domain of discourse），即 \(x\) 的取值范围。没有量词限定的变元称为\textbf{自由变元}，有量词限定的称为\textbf{约束变元}。

谓词和量词是谓词逻辑的基本语法元素，但要构建完整的逻辑语言，还需要系统地定义语言的语法。以下定义给出了标准的一阶语言所包含的全部符号------逻辑符号是固定的，而非逻辑符号则根据具体的数学领域而变化。

\textbf{定义 4.9}（一阶语言）：一阶语言由以下符号组成： - 逻辑符号：变元 \(x, y, z, \ldots\)；联结词 \(\neg, \land, \lor, \to, \leftrightarrow\)；量词 \(\forall, \exists\)；等词 \(=\) - 非逻辑符号：谓词符号 \(P, Q, R, \ldots\)（带 arity）；函数符号 \(f, g, h, \ldots\)（带 arity）；常量符号 \(c, d, \ldots\)

有了语言的语法定义，现在可以研究量词之间的逻辑关系。全称量词和存在量词通过否定相互关联------``并非所有\(x\)满足\(P\)''等价于''存在\(x\)不满足\(P\)``------这是量词推理中最基本的对偶原则。

\textbf{定理 4.5}（量词的基本等价）： 1. \(\neg \forall x P(x) \equiv \exists x \neg P(x)\)（全称与存在的对偶性） 2. \(\neg \exists x P(x) \equiv \forall x \neg P(x)\) 3. \(\forall x (P(x) \land Q(x)) \equiv \forall x P(x) \land \forall x Q(x)\)（全称量词对合取可分配） 4. \(\exists x (P(x) \lor Q(x)) \equiv \exists x P(x) \lor \exists x Q(x)\)（存在量词对析取可分配） 5. \(\forall x P(x) \lor \forall x Q(x) \to \forall x (P(x) \lor Q(x))\)（单向蕴涵，逆不成立） 6. \(\exists x (P(x) \land Q(x)) \to \exists x P(x) \land \exists x Q(x)\)（单向蕴涵，逆不成立）

\textbf{证明}：(1) \(\neg \forall x P(x)\) 为真 \(\iff\) 并非所有 \(x\) 满足 \(P(x)\) \(\iff\) 存在 \(x\) 不满足 \(P(x)\) \(\iff\) \(\exists x \neg P(x)\) 为真。\(\blacksquare\)

量词的基本等价关系使得我们可以将否定号''推入''量词内部，这对于逻辑表达式的化简至关重要。为了进一步标准化公式的结构，我们引入前束范式的概念------它要求所有量词都置于公式最前端，从而使得量词的辖域一目了然。

\textbf{定义 4.10}（前束范式）：谓词公式处于\textbf{前束范式}（prenex normal form），如果所有量词都出现在公式的最前端。任何谓词公式都等价于某个前束范式。前束范式是谓词逻辑中公式标准化的重要工具，通过将量词移至前端简化了公式的结构分析。

在命题逻辑中我们已经看到，重言式与定理的等价性（完备性定理）是逻辑系统的核心性质。这一结论在谓词逻辑中同样成立------这正是Godel于1929年证明的一阶逻辑完备性定理，它标志着现代数理逻辑的诞生。

\textbf{定理 4.6}（一阶逻辑的完备性定理，Godel 1929）：一阶逻辑中，逻辑有效式（在所有结构中为真的公式）与定理（从公理可推导的公式）完全一致。即 \(\models A\) 当且仅当 \(\vdash A\)。Godel 完备性定理是现代数理逻辑的基石，它确立了语法推导与语义蕴涵之间的等价性，使得一阶逻辑的形式化成为可能。

\textbf{证明（框架，Henkin 构造）}：完备性证明的核心是模型存在性定理（Model Existence Theorem）：每个无矛盾的一阶理论都有模型。给定无矛盾理论 \(T\)（即 \(T \nvdash \bot\)），通过添加 Henkin 常量（每个存在公式 \(\exists x \phi(x)\) 引入新常量 \(c_\phi\) 及公理 \(\phi(c_\phi)\)）将 \(T\) 扩充为极大无矛盾且具有 Henkin 性质的 Henkin 理论 \(T^*\)。然后由 \(T^*\) 的闭项构造项模型（term model），其元素为闭项的等价类。归纳验证此模型满足 \(T^*\) 的所有公式，从而 \(T\) 有模型。\(\blacksquare\)

\textbf{定理 4.7}（一阶逻辑的紧致性定理）：若一阶语句集 \(\Gamma\) 的每个有限子集有模型，则 \(\Gamma\) 本身有模型。紧致性定理是完备性定理的直接推论（若 \(\Gamma\) 无模型，则 \(\Gamma \models \bot\)，由完备性 \(\Gamma \vdash \bot\)，推导仅涉及 \(\Gamma\) 的有穷子集，故该有限子集无模型，矛盾）。紧致性定理在模型论中有广泛应用，如非标准分析（无穷小量的存在性）和 Lowenheim-Skolem 定理（一阶理论若有无限模型，则对任意无限基数 \(\kappa\) 有大小为 \(\kappa\) 的模型）。

完备性定理表明一阶逻辑足以捕获所有逻辑有效式。然而，Godel在1931年发现了更为惊人的事实：对于包含基本算术的形式系统，完备性不再成立。这一发现深刻改变了数学基础的研究方向。

\textbf{定理 4.8}（Godel 不完备性定理，1931）： - \textbf{第一不完备性定理}：任何包含基本算术（Peano 算术）的无矛盾、递归可公理化的一阶理论 \(T\) 中，存在语句 \(G\) 使得 \(T \nvdash G\) 且 \(T \nvdash \neg G\)（\(G\) 不可判定，\(T\) 不完备）。 - \textbf{第二不完备性定理}：这样的理论 \(T\) 不能证明自身的一致性（即 \(T \nvdash \operatorname{Con}(T)\)）。

\textbf{证明（框架）}：核心思想是''自指涉''（self-reference）。通过 Godel 编码（Godel numbering）将公式和证明编码为自然数，使得 \(T\) 能''谈论自身''。构造自指涉语句 \(G\) 意为''\(G\) 在 \(T\) 中不可证''。若 \(T \vdash G\)，则 \(T\) 证明了 \(G\) 不可证（矛盾于 \(T\) 的一致性，因为 \(T\) 既证明了 \(G\) 又'知道'\(G\) 不可证）。若 \(T \vdash \neg G\)，则 \(T\) 证明了 \(G\) 可证------但 \(G\) 的否定等价于'\(G\) 可证'，而 \(T\) 可能不一致（若 \(T\) 是一致的，则 \(T\) 不能证明 \(G\) 可证，因为这要求 \(G\) 确实可证，即 \(T \vdash G\)）。因此 \(T \nvdash G\) 且 \(T \nvdash \neg G\)。第二不完备性定理将一致性语句 \(\operatorname{Con}(T)\) 编码为 \(\neg \exists x(\operatorname{Proof}_T(x, \ulcorner 0=1 \urcorner))\)，并证明 \(T \vdash \operatorname{Con}(T) \to G\)。若 \(T \vdash \operatorname{Con}(T)\)，则 \(T \vdash G\)（与第一不完备性矛盾）。\(\blacksquare\)

\subsection{4.4 布尔代数}\label{ux5e03ux5c14ux4ee3ux6570}

命题逻辑的代数化将逻辑运算转化为代数运算，从而将逻辑问题纳入代数框架。布尔代数正是这一代数化的产物------它将命题的逻辑联结词（合取、析取、否定）抽象为集合上的二元运算，并提炼出五条基本公理。

\textbf{定义 4.11}（布尔代数）：\textbf{布尔代数}是六元组 \((B, \land, \lor, \neg, 0, 1)\)，其中 \(B\) 是至少包含两个元素的集合，\(\land\) 和 \(\lor\) 是 \(B\) 上的二元运算，\(\neg\) 是一元运算，\(0, 1 \in B\)，满足以下公理： 1. 交换律：\(a \land b = b \land a\)；\(a \lor b = b \lor a\) 2. 结合律：\((a \land b) \land c = a \land (b \land c)\)；\((a \lor b) \lor c = a \lor (b \lor c)\) 3. 分配律：\(a \land (b \lor c) = (a \land b) \lor (a \land c)\)；\(a \lor (b \land c) = (a \lor b) \land (a \lor c)\) 4. 单位元：\(a \land 1 = a\)；\(a \lor 0 = a\) 5. 互补律：\(a \land \neg a = 0\)；\(a \lor \neg a = 1\)

布尔代数是命题逻辑的代数化：\(B = \{0, 1\}\)（真值），\(\land\) 为合取，\(\lor\) 为析取，\(\neg\) 为否定。布尔代数将逻辑推理转化为代数计算，使得逻辑等价性可以通过代数恒等式来验证。更一般地，任何集合的幂集 \(\mathcal{P}(X)\) 配上 \(\cap\)（交）、\(\cup\)（并）、\(\complement\)（补）构成布尔代数。

\textbf{定理 4.9}（Stone 表示定理，1936）：任何布尔代数同构于某个紧致完全不连通的 Hausdorff 空间（Stone 空间）的既开又闭子集构成的布尔代数。Stone 表示定理建立了布尔代数与拓扑空间之间的深刻对偶关系（Stone 对偶），将代数的逻辑结构转化为空间的拓扑结构。

\textbf{证明（框架）}：给定布尔代数 \(B\)，构造其 Stone 空间 \(\operatorname{Spec}(B)\) = 所有 \(B\) 上的超滤子（ultrafilter）的集合，拓扑由形如 \(\{U : a \in U\}\)（\(a \in B\)）的基生成。映射 \(a \mapsto \{U \in \operatorname{Spec}(B) : a \in U\}\) 给出 \(B\) 到 \(\operatorname{Spec}(B)\) 的既开又闭子集布尔代数的同构。\(\blacksquare\)

Stone表示定理从理论上完整刻画了布尔代数的结构。在应用层面，我们通常只需要知道布尔代数等式在二元布尔代数 \(\{0,1\}\) 上的有效性就足够了------这是布尔代数完备性的核心内容。

\textbf{定理 4.10}（布尔代数的完备性）：任何布尔代数上的等式恒成立当且仅当它在二元布尔代数 \(\{0, 1\}\) 上成立。这等价于命题逻辑中''重言式等价于在所有真值指派下为真''的完备性定理。

\hrulefill

\section{第31章：线性规划}\label{ux7b2c31ux7ae0ux7ebfux6027ux89c4ux5212}

线性规划（Linear Programming, LP）是运筹学中最重要的优化方法之一，研究线性约束条件下线性目标函数的极值问题。线性规划的数学基础由 Kantorovich（1939 年，获 1975 年 Nobel 经济学奖）和 Dantzig（1947 年，单纯形法）建立，随后 von Neumann 的对偶理论和 Karmarkar 的内点法（1984 年）进一步丰富了其理论。线性规划在经济学（资源分配、生产计划）、运输物流（网络流问题）、博弈论（零和博弈的极小极大定理）和机器学习（支持向量机）中都有广泛应用。线性规划的魅力在于：尽管其建模能力有限（仅线性约束和线性目标），但其理论（对偶性、互补松弛性）和算法（单纯形法、内点法）极为成熟，且很多非线性问题可通过线性规划系列（如序列线性规划）逼近解决。

\subsection{5.1 线性规划的标准形式与基本概念}\label{ux7ebfux6027ux89c4ux5212ux7684ux6807ux51c6ux5f62ux5f0fux4e0eux57faux672cux6982ux5ff5}

\textbf{定义 5.1}（线性规划的标准形式）：线性规划的标准形式为：

\[\begin{aligned}
\min \quad & \mathbf{c}^T \mathbf{x} \\
\text{s.t.} \quad & A\mathbf{x} = \mathbf{b} \\
& \mathbf{x} \geq \mathbf{0}
\end{aligned}\]

其中 \(A \in \mathbb{R}^{m \times n}\)（\(m < n\)，秩为 \(m\)），\(\mathbf{b} \in \mathbb{R}^m\)，\(\mathbf{c} \in \mathbb{R}^n\)。\(\mathbf{c}^T \mathbf{x}\) 是\textbf{目标函数}（objective function），\(A\mathbf{x} = \mathbf{b}\) 是\textbf{等式约束}（equality constraints），\(\mathbf{x} \geq \mathbf{0}\) 是\textbf{非负约束}（non-negativity constraints）。任何线性规划问题（包括最大化问题、不等式约束、自由变量）都可转化为标准形式。

\textbf{转化方法}： - 最大化 \(\max \mathbf{c}^T \mathbf{x}\) 等价于 \(\min -\mathbf{c}^T \mathbf{x}\) - 不等式约束 \(A\mathbf{x} \leq \mathbf{b}\) 通过引入松弛变量 \(\mathbf{s} \geq \mathbf{0}\) 转化为 \(A\mathbf{x} + \mathbf{s} = \mathbf{b}\) - 不等式约束 \(A\mathbf{x} \geq \mathbf{b}\) 通过引入盈余变量 \(\mathbf{s} \geq \mathbf{0}\) 转化为 \(A\mathbf{x} - \mathbf{s} = \mathbf{b}\) - 自由变量 \(x_i\)（无符号限制）分解为 \(x_i = x_i^+ - x_i^-\)（\(x_i^+, x_i^- \geq 0\)）

有了标准形式的统一定义，下一步需要明确什么是''解''以及解的分类。在线性规划中，由于约束条件的特殊性，解的结构具有丰富的几何意义------可行域是凸多面体，最优解总可以在极点（顶点）处达到。

\textbf{定义 5.2}（可行解、基本解与最优解）： - \textbf{可行解}（feasible solution）：满足所有约束的 \(\mathbf{x}\) - \textbf{可行域}（feasible region）：所有可行解构成的集合，是一个凸多面体（convex polyhedron） - \textbf{基本解}（basic solution）：由 \(A\) 的 \(m\) 个线性无关列构成的基矩阵 \(B\) 确定的解，非基变量取 \(0\) - \textbf{基本可行解}（basic feasible solution, BFS）：满足非负性的基本解 - \textbf{最优解}（optimal solution）：使目标函数达到最小值的可行解

可行解、基本解和最优解这三个概念刻画了线性规划解的层次结构。以下基本定理揭示了它们之间的深刻联系：每一个最优解都可以在基本可行解中实现，而基本可行解恰好对应于可行域（凸多面体）的极点。

\textbf{定理 5.1}（线性规划的基本定理）： 1. 若可行域非空，则存在基本可行解 2. 若存在最优解，则存在最优基本可行解 3. 可行域的极点和基本可行解一一对应

\textbf{证明}：(1) 在可行域中取 \(\ell_0\) 范数（非零分量个数）最小的可行解 \(\mathbf{x}^*\)。若 \(\mathbf{x}^*\) 的支撑列线性相关，则存在非零向量 \(\mathbf{d}\) 使得 \(A\mathbf{d} = \mathbf{0}\) 且 \(\operatorname{supp}(\mathbf{d}) \subseteq \operatorname{supp}(\mathbf{x}^*)\)。沿 \(\mathbf{d}\) 方向调整 \(\mathbf{x}^*\) 可减少非零分量数，矛盾。故支撑列线性无关，可扩充为基。(2) 若最优解非基本，通过类似论证可沿可行方向移动至极点而不增加目标值。(3) 极点是无法表示为两个不同可行解的凸组合的点，恰好对应基本可行解。\(\blacksquare\)

基本定理确保了最优解可以在极点处寻找，这为单纯形法的设计提供了理论基础。在线性规划理论中，另一个核心概念是对偶性------每个线性规划问题（原始问题）都自然地对应着一个对偶问题，两者之间存在着深刻的不等式关系。

\textbf{定理 5.2}（线性规划的对偶性，弱对偶定理）：设 (P) 为原始问题 \(\min\{\mathbf{c}^T \mathbf{x} : A\mathbf{x} = \mathbf{b}, \mathbf{x} \geq \mathbf{0}\}\)，(D) 为对偶问题 \(\max\{\mathbf{b}^T \mathbf{y} : A^T \mathbf{y} \leq \mathbf{c}, \mathbf{y} \text{ free}\}\)。则对任意可行解 \(\mathbf{x}\) 和 \(\mathbf{y}\)，有 \(\mathbf{c}^T \mathbf{x} \geq \mathbf{b}^T \mathbf{y}\)（弱对偶性）。

\textbf{证明}：\(\mathbf{c}^T \mathbf{x} - \mathbf{b}^T \mathbf{y} = \mathbf{c}^T \mathbf{x} - (A\mathbf{x})^T \mathbf{y} = \mathbf{c}^T \mathbf{x} - \mathbf{x}^T A^T \mathbf{y} = \mathbf{x}^T(\mathbf{c} - A^T \mathbf{y}) \geq 0\)。因为 \(\mathbf{x} \geq \mathbf{0}\) 且 \(\mathbf{c} - A^T \mathbf{y} \geq \mathbf{0}\)（对偶可行性）。\(\blacksquare\)

\textbf{定理 5.3}（强对偶定理）：若 (P) 有有限最优解，则 (D) 也有有限最优解，且两者最优值相等：\(\min \mathbf{c}^T \mathbf{x} = \max \mathbf{b}^T \mathbf{y}\)。

\textbf{证明}（框架，由 Farkas 引理或单纯形法）：设 \(\mathbf{x}^*\) 为 (P) 的最优基本可行解，基为 \(B\)。令 \(\mathbf{y}^* = (B^{-1})^T \mathbf{c}_B\)。由最优性条件，简约成本 \(\bar{\mathbf{c}}_N = \mathbf{c}_N - N^T \mathbf{y}^* \geq \mathbf{0}\)，故 \(\mathbf{y}^*\) 对偶可行。且 \(\mathbf{b}^T \mathbf{y}^* = \mathbf{b}^T (B^{-1})^T \mathbf{c}_B = \mathbf{c}_B^T B^{-1} \mathbf{b} = \mathbf{c}^T \mathbf{x}^*\)。由弱对偶，\(\mathbf{y}^*\) 为 (D) 的最优解。\(\blacksquare\)

\textbf{定理 5.3b}（Farkas 引理------线性不等式系统的择一定理）：对矩阵 \(A\) 和向量 \(\mathbf{b}\)，以下两者恰有一个成立： (1) 存在 \(\mathbf{x} \geq \mathbf{0}\) 使得 \(A\mathbf{x} = \mathbf{b}\) (2) 存在 \(\mathbf{y}\) 使得 \(A^T \mathbf{y} \geq \mathbf{0}\) 且 \(\mathbf{b}^T \mathbf{y} < 0\)

\textbf{证明}：若 (1) 和 (2) 同时成立，则 \(0 > \mathbf{b}^T \mathbf{y} = \mathbf{x}^T A^T \mathbf{y} \geq 0\)（矛盾，因为 \(\mathbf{x} \geq \mathbf{0}\) 且 \(A^T \mathbf{y} \geq \mathbf{0}\)）。若 (1) 不成立，则 \(\mathbf{b} \notin \operatorname{cone}(A)\)（\(A\) 的列生成的凸锥）。由分离超平面定理，存在 \(\mathbf{y}\) 使得 \(\mathbf{y}^T \mathbf{b} < 0\) 且 \(\mathbf{y}^T A \geq \mathbf{0}\)，即 (2) 成立。Farkas 引理是线性规划对偶理论的基础，也是证明强对偶定理的关键工具。\(\blacksquare\)

\textbf{定理 5.4}（互补松弛条件）：设 \(\mathbf{x}\) 为 (P) 的可行解，\(\mathbf{y}\) 为 (D) 的可行解。则 \(\mathbf{x}\) 和 \(\mathbf{y}\) 分别为最优解当且仅当 \(\mathbf{x}_i (\mathbf{c} - A^T \mathbf{y})_i = 0\) 对所有 \(i\) 成立（互补松弛性）。

\textbf{证明}：\(\mathbf{c}^T \mathbf{x} - \mathbf{b}^T \mathbf{y} = \mathbf{x}^T(\mathbf{c} - A^T \mathbf{y}\))。由强对偶，\(\mathbf{x}\) 和 \(\mathbf{y}\) 最优当且仅当 \(\mathbf{c}^T \mathbf{x} = \mathbf{b}^T \mathbf{y}\)，即 \(\mathbf{x}^T(\mathbf{c} - A^T \mathbf{y}) = 0\)。由于 \(\mathbf{x} \geq 0\) 且 \(\mathbf{c} - A^T \mathbf{y} \geq 0\)，内积为零等价于逐项乘积为零。\(\blacksquare\)

\subsection{5.2 单纯形法}\label{ux5355ux7eafux5f62ux6cd5}

\textbf{算法 5.1}（单纯形法 / Simplex Method，Dantzig 1947）：单纯形法在基本可行解之间移动，每次迭代严格降低目标函数值（非退化情形下）。

\textbf{步骤}： 1. \textbf{初始化}：找到一个初始基本可行解（必要时使用两阶段法或大-M 法） 2. \textbf{最优性检验}：计算简约成本 \(\bar{\mathbf{c}}_N = \mathbf{c}_N - N^T B^{-T} \mathbf{c}_B\)。若 \(\bar{\mathbf{c}}_N \geq \mathbf{0}\)，当前解为最优。否则选择入基变量 \(x_q\)（\(\bar{c}_q < 0\)） 3. \textbf{最小比率检验}：计算 \(\mathbf{d} = B^{-1} \mathbf{a}_q\)（\(\mathbf{a}_q\) 为入基列）。若 \(\mathbf{d} \leq \mathbf{0}\)，目标函数无下界（问题无界）。否则选择出基变量 \(x_{p}\) 使得 \(\min_i \{(\mathbf{x}_B)_i / d_i : d_i > 0\}\) 达到 4. \textbf{枢轴运算}：更新基矩阵 \(B\)（\(x_q\) 入基，\(x_p\) 出基），更新基本可行解，返回步骤 2

\textbf{定理 5.5}（单纯形法的有限终止性）：若非退化（所有基本可行解非退化，即每个基本可行解恰好有 \(m\) 个正分量），则单纯形法在有限步内终止（达到最优或检测无界）。在退化情形下，Bland 法则（最小下标法则）确保不出现循环。

\textbf{证明}：非退化情形下，每次枢轴运算严格降低目标函数值，因此不会重复访问同一个基本可行解。基本可行解的总数是有限的（\(\binom{n}{m}\)），故算法在有限步内终止。退化情形下，Bland 法则（选择最小下标的入基和出基变量）通过反证法证明不会出现循环。\(\blacksquare\)

\textbf{定理 5.6}（单纯形法的复杂度）：在最坏情况下，单纯形法的迭代次数可能是指数级的（Klee-Minty 立方体，1972 年）。然而，在''平均''意义下和实践中，单纯形法表现优异，通常只需 \(O(m)\) 到 \(O(n)\) 次迭代。Smoothed analysis（Spielman-Teng 2001 年）证明了单纯形法在随机扰动下的期望迭代次数是多项式的，解释了其在实践中的高效性。

单纯形法的几何解释：从可行域（多面体）的一个顶点（基本可行解）出发，沿多面体的边（下降方向）移动至相邻顶点，直到达到最优顶点。单纯形法本质上是多面体顶点上的局部搜索算法，其正确性依赖于线性规划的最优解必然在顶点达到（基本定理）和线性目标函数在顶点上的最优性条件（简约成本非负）。

\subsection{5.3 对偶理论与后最优分析}\label{ux5bf9ux5076ux7406ux8bbaux4e0eux540eux6700ux4f18ux5206ux6790}

\textbf{定理 5.7}（对偶问题的经济解释------影子价格）：对偶变量 \(y_i\) 是第 \(i\) 个资源的\textbf{影子价格}（shadow price），衡量增加一单位资源 \(b_i\) 对最优目标值的边际贡献：\(\frac{\partial z^*}{\partial b_i} = y_i^*\)（\(z^*\) 为最优值）。影子价格在经济学中表示资源的稀缺性价值：资源越稀缺，其影子价格越高。

\textbf{证明}：设 \(B\) 为最优基。\(z^* = \mathbf{c}_B^T B^{-1} \mathbf{b} = \mathbf{y}^{*T} \mathbf{b}\)。因此 \(\frac{\partial z^*}{\partial b_i} = y_i^*\)（在 \(B\) 保持最优的范围内）。\(\blacksquare\)

\textbf{定理 5.8}（敏感性分析）： 1. \textbf{右端项变化}：最优基 \(B\) 在 \(\mathbf{b} + \Delta \mathbf{b}\) 下保持可行当且仅当 \(B^{-1}(\mathbf{b} + \Delta \mathbf{b}) \geq \mathbf{0}\)（即 \(\Delta \mathbf{b} \geq -B \mathbf{x}_B\)） 2. \textbf{目标系数变化}：最优基 \(B\) 在 \(\mathbf{c} + \Delta \mathbf{c}\) 下保持最优当且仅当简约成本 \(\bar{\mathbf{c}}_N + \Delta \mathbf{c}_N - N^T B^{-T} \Delta \mathbf{c}_B \geq \mathbf{0}\)

敏感性分析在实际应用中至关重要：它允许决策者在数据变化时快速更新最优解，而无需重新求解整个线性规划，同时提供了最优解对参数变化的鲁棒性信息。

\textbf{定理 5.9}（对偶单纯形法）：对偶单纯形法保持对偶可行性（\(\bar{\mathbf{c}}_N \geq \mathbf{0}\)），逐步恢复原始可行性（\(B^{-1}\mathbf{b} \geq \mathbf{0}\)）。对偶单纯形法在以下情形特别有用：(1) 右端项 \(\mathbf{b}\) 变化后重新优化；(2) 添加新约束后重新优化；(3) 分支定界法中的子问题求解。

\textbf{步骤}：选择出基变量 \(x_p\)（\((\mathbf{x}_B)_p < 0\)），计算比率 \(\min_j \{|\bar{c}_j / d_{pj}| : d_{pj} < 0\}\) 选择入基变量 \(x_q\)，枢轴运算。若不存在 \(d_{pj} < 0\)，则问题无可行解。

\hrulefill

\emph{卷三：代数与逻辑终。}

\newpage

\chapter{02-几何/01-欧氏几何/01-平面几何.md}\label{ux51e0ux4f5501-ux6b27ux6c0fux51e0ux4f5501-ux5e73ux9762ux51e0ux4f55.md}

\textbf{前置知识}：本章假定读者已掌握卷一 Ch 2 逻辑基础（命题逻辑、谓词逻辑、证明方法）和 Ch 3 集合论（集合的运算、关系与函数）。几何公理将使用逻辑语言精确表述几何直觉，每一条公理都是一个严格的形式化命题。

\hrulefill

几何学是数学中最古老的分支之一。公元前约300年，欧几里得在《几何原本》中首次尝试用公理化方法建立几何学体系。然而，欧几里得的公理系统在严谨性上存在明显不足------他的一些''公理''实际上是定理，某些证明中隐含了未经声明的假设。1899年，德国数学家大卫·希尔伯特在其著作《几何基础》中提出了一个完整的公理系统，将欧氏几何建立在严格的逻辑基础之上。本章详细阐述这一公理系统，并从中推导出若干最基本的几何命题。

希尔伯特公理系统由五组公理构成：关联公理（描述点、直线、平面之间的结合关系）、顺序公理（描述点在直线上的排列顺序）、合同公理（描述线段和角的全等关系）、平行公理（描述平行线的存在性和唯一性）以及连续公理（描述直线的连续性和完备性）。这五组公理共同刻画了欧氏几何的全部结构。

\section{第32章：欧氏几何的公理体系}\label{ux7b2c32ux7ae0ux6b27ux6c0fux51e0ux4f55ux7684ux516cux7406ux4f53ux7cfb}

关联公理（Incidence Axioms），又称结合公理，是希尔伯特公理体系的第一组公理。这组公理刻画了''点属于直线''\,``点属于平面''``直线含于平面''等最基本的几何关系。在希尔伯特的体系中，这些关系是不加定义的原始概念，它们的全部性质完全由公理 I1--I8 所规定。关联公理不仅确立了''两点确定一条直线''\,``不共线三点确定一个平面''等我们熟知的几何事实，更重要的是，它们排除了诸如''两条直线有两个交点''或''所有点都在同一平面上''等退化情形。公理 I8 尤为关键------它断言至少存在四个不共面的点，从而保证了空间的''三维性''。从这组公理出发，可以严格证明两条不同直线至多有一个交点、两个不同平面的交集是一条直线等基本命题。

\subsubsection{12.1.1 基本对象与基本关系}\label{ux57faux672cux5bf9ux8c61ux4e0eux57faux672cux5173ux7cfb}

希尔伯特公理系统建立在三个\textbf{基本对象}和三个\textbf{基本关系}之上：

\begin{itemize}
\tightlist
\item
  \textbf{基本对象}：点（point）、直线（line）、平面（plane）。
\item
  \textbf{基本关系}：关联（incidence，又称''结合''或''属于''）、介于（betweenness，又称''顺序''）、合同（congruence，又称''全等''）。
\end{itemize}

这些基本对象和基本关系本身不加定义，它们的意义完全由公理所规定的性质来确定。这正体现了公理化方法的精神：我们不追问''点是什么''，而只关心''点满足什么性质''。

我们采用如下记号： - 点用大写字母 \(A, B, C, \ldots\) 表示； - 直线用小写字母 \(a, b, c, \ldots\) 或 \(\overleftrightarrow{AB}\) 表示； - 平面用希腊字母 \(\alpha, \beta, \gamma, \ldots\) 表示； - ``点 \(A\) 在直线 \(a\) 上''或''直线 \(a\) 经过点 \(A\)``记为 \(A \in a\)； -''点 \(A\) 在平面 \(\alpha\) 上''记为 \(A \in \alpha\)； - ``直线 \(a\) 在平面 \(\alpha\) 上''记为 \(a \subset \alpha\)。

\subsubsection{12.1.2 关联公理（结合公理）I1--I8}\label{ux5173ux8054ux516cux7406ux7ed3ux5408ux516cux7406i1i8}

关联公理描述点、直线、平面之间的基本结合关系。

\textbf{公理 I1}：对于任意两个不同的点 \(A\) 和 \(B\)，存在一条直线 \(a\) 使得 \(A \in a\) 且 \(B \in a\)。

\textbf{公理 I2}：对于任意两个不同的点 \(A\) 和 \(B\)，至多存在一条直线 \(a\) 使得 \(A \in a\) 且 \(B \in a\)。

由公理 I1 和 I2 可知，任意两个不同的点\textbf{唯一确定}一条直线。我们将通过点 \(A\) 和 \(B\) 的唯一直线记为 \(\overleftrightarrow{AB}\)。

\textbf{公理 I3}：每条直线上至少存在两个不同的点。存在三个不共线的点。

\textbf{公理 I4}：对于任意三个不共线的点 \(A\)、\(B\)、\(C\)，存在一个平面 \(\alpha\) 使得 \(A \in \alpha\)、\(B \in \alpha\)、\(C \in \alpha\)。对于任意平面，其上至少存在一个点。

\textbf{公理 I5}：对于任意三个不共线的点 \(A\)、\(B\)、\(C\)，至多存在一个平面使得 \(A\)、\(B\)、\(C\) 都在其上。

由公理 I4 和 I5 可知，任意三个不共线的点\textbf{唯一确定}一个平面。

\textbf{公理 I6}：如果一条直线上的两个不同的点都在平面 \(\alpha\) 上，那么该直线上的所有点都在 \(\alpha\) 上（即直线含于 \(\alpha\)）。

\textbf{公理 I7}：如果两个不同的平面有一个公共点，那么它们至少还有一个公共点。

\textbf{公理 I8}：至少存在四个不共面的点。

公理 I8 保证了空间的''三维性''------排除了所有点都在同一个平面上的退化情况。

\textbf{定义 12.1}（共线）：如果若干个点都在同一条直线上，则称这些点是\textbf{共线的}。

\textbf{定义 12.2}（共面）：如果若干个点都在同一个平面上，则称这些点是\textbf{共面的}。

\textbf{命题 12.1}：两条不同的直线至多有一个公共点。

\textbf{证明}：假设直线 \(a\) 和直线 \(b\) 有两个不同的公共点 \(A\) 和 \(B\)。根据公理 I2，通过 \(A\) 和 \(B\) 至多有一条直线，因此 \(a = b\)，与 \(a\) 和 \(b\) 不同的假设矛盾。\(\blacksquare\)

\textbf{命题 12.2}：如果两个不同的平面有一个公共点，那么它们的交集是一条直线。

\textbf{证明}：设平面 \(\alpha\) 和平面 \(\beta\) 有一个公共点 \(P\)。根据公理 I7，它们至少还有一个公共点 \(Q\)（\(Q \neq P\)）。由公理 I2，直线 \(\overleftrightarrow{PQ}\) 唯一确定。

设 \(R\) 是 \(\alpha\) 和 \(\beta\) 的任意公共点。如果 \(R\) 不在 \(\overleftrightarrow{PQ}\) 上，则 \(P\)、\(Q\)、\(R\) 不共线，由公理 I5，它们唯一确定一个平面。但 \(P\)、\(Q\)、\(R\) 都在 \(\alpha\) 上，也都在 \(\beta\) 上，因此 \(\alpha = \beta\)，与两个平面不同的假设矛盾。

因此，\(\alpha\) 和 \(\beta\) 的所有公共点都在 \(\overleftrightarrow{PQ}\) 上。另一方面，由公理 I6，\(\overleftrightarrow{PQ}\) 上的所有点既在 \(\alpha\) 上也在 \(\beta\) 上。所以 \(\alpha \cap \beta = \overleftrightarrow{PQ}\)。\(\blacksquare\)

\subsection{12.2 顺序公理}\label{ux987aux5e8fux516cux7406}

顺序公理描述点在直线上的相对位置关系，引入''介于''的概念。这是建立线段、射线、角的内部等概念的基础。

\textbf{定义 12.3}（介于关系）：设 \(A\)、\(B\)、\(C\) 是一条直线上的三个不同的点。如果 \(B\) 在 \(A\) 和 \(C\) 之间，则记为 \(A * B * C\)。

\textbf{公理 O1}：如果 \(A * B * C\)，则 \(A\)、\(B\)、\(C\) 是同一条直线上三个不同的点，且 \(C * B * A\)。

\textbf{公理 O2}：对于任意两个不同的点 \(A\) 和 \(B\)，直线 \(\overleftrightarrow{AB}\) 上至少存在一个点 \(C\) 使得 \(A * B * C\)。

\textbf{公理 O3}：对于同一条直线上的任意三个不同的点，至多有一个点介于另外两个点之间。

\textbf{公理 O4}（帕施公理）：设 \(A\)、\(B\)、\(C\) 是三个不共线的点，\(a\) 是一条不经过 \(A\)、\(B\)、\(C\) 中任何一个点的直线。如果直线 \(a\) 经过线段 \(AB\) 的某个内点，那么它也经过线段 \(BC\) 或线段 \(AC\) 的某个内点。

\textbf{定义 12.4}（线段）：对于两个不同的点 \(A\) 和 \(B\)，\textbf{线段} \(AB\) 是指直线 \(\overleftrightarrow{AB}\) 上所有满足 \(A * X * B\) 的点 \(X\) 以及 \(A\)、\(B\) 两点本身的集合。\(A\) 和 \(B\) 称为线段 \(AB\) 的\textbf{端点}。线段 \(AB\) 记为 \(\overline{AB}\)。

\textbf{定义 12.5}（射线）：对于两个不同的点 \(A\) 和 \(B\)，以 \(A\) 为端点、经过 \(B\) 的\textbf{射线} \(\overrightarrow{AB}\) 是指直线 \(\overleftrightarrow{AB}\) 上所有满足 \(A * B * X\) 或 \(X = A\) 或 \(X = B\) 的点 \(X\) 的集合。

\textbf{定义 12.6}（开半直线）：直线 \(\overleftrightarrow{AB}\) 上除去线段 \(\overline{AB}\) 的两个部分（不包括 \(A\) 和 \(B\)），分别称为以 \(A\) 为分界点、关于 \(B\) 的\textbf{开半直线}及其\textbf{对侧}。

\textbf{引理 12.1}（介于关系的基本推论）：设 \(A\)、\(B\)、\(C\)、\(D\) 是同一条直线上的不同点。

\begin{enumerate}
\def\labelenumi{\arabic{enumi}.}
\tightlist
\item
  若 \(A * B * C\)，则 \(C * B * A\)。
\item
  若 \(A * B * D\) 且 \(C * B * D\)，则 \(A * B * C\)。
\end{enumerate}

\textbf{证明}：

性质 1 是公理 O1 的直接推论。

性质 2 的证明：由公理 O3，在 \(\{A, B, C\}\) 中至多一种''介于''关系成立。我们用反证法排除 \(A * C * B\) 和 \(B * A * C\) 两种可能：

\begin{itemize}
\item
  若 \(A * C * B\)，则 \(C\) 在 \(A\) 和 \(B\) 之间。由 \(A * B * D\) 知 \(B\) 在 \(A\) 和 \(D\) 之间，故 \(C\) 介于 \(A\) 与 \(B\) 之间、而 \(B\) 又介于 \(A\) 与 \(D\) 之间，因此 \(B\) 在 \(C\) 和 \(D\) 之间，即 \(C * B * D\)。但这意味着 \(B\) 和 \(D\) 在 \(C\) 的两侧，而 \(A * C * B\) 表明 \(B\) 和 \(A\) 也在 \(C\) 的两侧。综合知 \(A\) 和 \(D\) 在 \(C\) 的同侧，从而 \(A * C * D\)。又 \(C * B * D\) 表明 \(D\) 在 \(B\) 的远离 \(C\) 的一侧，而 \(A * C * D\) 表明 \(D\) 在 \(C\) 的远离 \(A\) 的一侧，两者一致；但 \(A * C * B\) 表明 \(B\) 在 \(C\) 的远离 \(A\) 的一侧，而 \(C * B * D\) 表明 \(B\) 在 \(C\) 和 \(D\) 之间，即 \(B\) 在 \(C\) 的远离 \(A\) 的一侧、\(D\) 在 \(B\) 的远离 \(C\) 的一侧，这与 \(A * C * D\) 中 \(D\) 在 \(C\) 的远离 \(A\) 的一侧且 \(A\) 在 \(C\) 和 \(D\) 之间矛盾（\(A\) 不可能既在 \(C\) 的远离 \(B\) 的一侧又在 \(C\) 和 \(D\) 之间）。因此 \(A * C * B\) 不成立。
\item
  若 \(B * A * C\)，则 \(A\) 在 \(B\) 和 \(C\) 之间。由 \(A * B * D\) 知 \(B\) 在 \(A\) 和 \(D\) 之间，故 \(B\) 和 \(D\) 在 \(A\) 的两侧；而 \(B * A * C\) 表明 \(B\) 和 \(C\) 也在 \(A\) 的两侧，因此 \(C\) 和 \(D\) 在 \(A\) 的同侧。由 \(C * B * D\)（性质 1 得 \(D * B * C\)），\(B\) 在 \(C\) 和 \(D\) 之间，\(C\) 和 \(D\) 在 \(B\) 的两侧。但 \(B * A * C\) 表明 \(C\) 在 \(A\) 的远离 \(B\) 的一侧，而 \(A * B * D\) 表明 \(D\) 在 \(B\) 的远离 \(A\) 的一侧。\(C\) 和 \(D\) 在 \(A\) 的同侧，但 \(B\) 在 \(C\) 和 \(D\) 之间要求 \(C\) 和 \(D\) 在 \(B\) 的两侧，而 \(B\) 在 \(A\) 和 \(D\) 之间要求 \(D\) 在 \(A\) 的远离 \(B\) 的一侧、\(C\) 在 \(A\) 的靠近 \(B\) 的一侧（由 \(B * A * C\)），矛盾。因此 \(B * A * C\) 不成立。
\end{itemize}

排除以上两种可能后，必有 \(A * B * C\) 成立。\(\blacksquare\)''

\textbf{命题 12.3}：设 \(A\)、\(B\)、\(C\) 是同一条直线上的三个不同的点。则恰好有以下三种情况之一成立：

\begin{enumerate}
\def\labelenumi{\arabic{enumi}.}
\tightlist
\item
  \(A * B * C\)；
\item
  \(A * C * B\)；
\item
  \(B * A * C\)。
\end{enumerate}

\textbf{证明}：由公理 O3（唯一性），至多有一种''介于''关系成立。下面证明至少有一种成立，采用反证法。

假设 \(A\)、\(B\)、\(C\) 之间不存在任何''介于''关系，即 \(A * B * C\)、\(A * C * B\)、\(B * A * C\) 均不成立。我们分步导出矛盾。

\textbf{第一步}：由公理 O2（存在性），存在点 \(D\) 使得 \(A * B * D\)。若 \(D = C\)，则 \(A * B * C\) 成立，与假设矛盾。故 \(D \neq C\)，即 \(A\)、\(B\)、\(C\)、\(D\) 是直线上四个不同的点。

\textbf{第二步}：考虑 \(C\) 与 \(D\) 的位置。由公理 O3，在三个点 \(\{B, C, D\}\) 中恰好有一种''介于''关系成立。由 \(A * B * D\) 知 \(B \neq D\)，且 \(C \neq B\)、\(C \neq D\)，故以下三种情况恰有一种成立：

\textbf{(i)} \(B * C * D\)：此时考虑三个点 \(\{A, B, C\}\)。由公理 O4（帕施公理），构造辅助线：取不在直线 \(l\) 上的点 \(P\)，作三角形 \(\triangle ABD\)。点 \(C\) 满足 \(B * C * D\)，即 \(C\) 在线段 \(BD\) 上。由于 \(A * B * D\)，点 \(B\) 在线段 \(AD\) 上。在三角形 \(\triangle ACD\) 中，\(B\) 在线段 \(AD\) 上，\(B * C * D\) 表明 \(C\) 在线段 \(BD\) 上。由公理 O3，我们分析 \(\{A, B, C\}\) 的排列：\(A * B * C\) 和 \(B * A * C\) 互斥（公理 O3 唯一性）。若 \(A * B * C\)，则与假设矛盾。因此必有 \(B * A * C\) 或 \(A * C * B\)。但若 \(B * A * C\)，则 \(A\) 在 \(B\) 和 \(C\) 之间；而由 \(A * B * D\) 和 \(B * C * D\)，\(A\)、\(B\)、\(C\) 在 \(D\) 的同侧且 \(B\) 在 \(A\) 与 \(C\) 之间，这与 \(B * A * C\) 矛盾（公理 O3：\(A * B * D\) 和 \(B * A * C\) 不能同时成立，否则 \(B\) 既在 \(A\) 和 \(D\) 之间又在 \(A\) 的另一侧）。类似地，\(A * C * B\) 也会导致矛盾。因此在情形 (i) 下 \(A * B * C\) 必成立，与假设矛盾。

\textbf{(ii)} \(B * D * C\)：此时 \(D\) 在 \(B\) 和 \(C\) 之间。由 \(A * B * D\) 和 \(B * D * C\)，在三个点 \(\{A, B, C\}\) 中：由公理 O3，排列 \(A * B * C\)、\(A * C * B\)、\(B * A * C\) 恰有一种成立。我们排除后两种：若 \(B * A * C\)，则 \(A\) 在 \(B\) 和 \(C\) 之间；但 \(A * B * D\) 表明 \(B\) 在 \(A\) 和 \(D\) 之间，\(B * D * C\) 表明 \(D\) 在 \(B\) 和 \(C\) 之间，因此 \(A\) 和 \(C\) 在 \(B\) 的异侧，\(A\) 不可能在 \(B\) 和 \(C\) 之间，矛盾。若 \(A * C * B\)，则 \(C\) 在 \(A\) 和 \(B\) 之间；但 \(A * B * D\) 表明 \(B\) 在 \(A\) 和 \(D\) 之间，\(C\) 在 \(B\) 的远离 \(A\) 一侧（因为 \(B * D * C\)），矛盾。因此在情形 (ii) 下 \(A * B * C\) 必成立，与假设矛盾。

\textbf{(iii)} \(C * B * D\)：此时 \(B\) 在 \(C\) 和 \(D\) 之间。结合 \(A * B * D\)（\(B\) 在 \(A\) 和 \(D\) 之间），可知 \(A\) 和 \(C\) 在 \(B\) 的两侧（因为 \(B\) 将直线分为两个半直线，\(A\) 和 \(D\) 在 \(B\) 的两侧，而 \(C\) 和 \(D\) 也在 \(B\) 的两侧，故 \(A\) 和 \(C\) 在 \(B\) 的异侧）。因此 \(B * A * C\)（\(B\) 在 \(A\) 和 \(C\) 之间），与假设矛盾。

\textbf{第三步}：以上三种情况均与假设矛盾，故假设不成立。三个点之间至少有一种''介于''关系成立。结合至多一种成立（公理 O3），恰好有一种''介于''关系成立。\(\blacksquare\)

\textbf{定义 12.7}（三角形）：三个不共线的点 \(A\)、\(B\)、\(C\) 以及三条线段 \(\overline{AB}\)、\(\overline{BC}\)、\(\overline{CA}\) 组成的图形称为\textbf{三角形}，记为 \(\triangle ABC\)。点 \(A\)、\(B\)、\(C\) 称为三角形的\textbf{顶点}，线段 \(\overline{AB}\)、\(\overline{BC}\)、\(\overline{CA}\) 称为三角形的\textbf{边}。

\subsection{12.3 合同公理}\label{ux5408ux540cux516cux7406}

合同公理（Congruence Axioms）是希尔伯特公理体系中的第三组公理，也是内容最为丰富的一组。这组公理引入线段和角的''全等''（合同）关系，用 \(\cong\) 表示。合同关系刻画了''长度相等''和''角度相等''的几何直觉，但它本身并不依赖于''长度''或''角度''的数值概念------希尔伯特的巧妙之处在于，他通过 C1--C3 公理在线段合同关系上建立了等价关系和可加性，从而为''长度''概念提供了公理基础，而无需引入实数。公理 C4--C5 对角合同关系做了类似的公理化处理。公理 C6（SAS 公理）是这组公理中最关键的一条------它将线段的合同与角的合同联系起来，从而为三角形全等的判定提供了基础。值得注意的是，欧几里得在《几何原本》中试图将 SAS 作为定理来证明，但希尔伯特认识到其证明隐含了未声明的假设，因此将它提升为公理，这正是公理化方法严谨性的体现。

\subsubsection{12.3.1 线段的合同}\label{ux7ebfux6bb5ux7684ux5408ux540c}

\textbf{公理 C1}：设 \(A\)、\(B\) 是直线 \(a\) 上的两个点，\(A'\) 是直线 \(a'\) 上的点，\(p'\) 是 \(a'\) 上以 \(A'\) 为端点的一条射线。则在射线 \(p'\) 上存在唯一的点 \(B'\) 使得 \(\overline{AB} \cong \overline{A'B'}\)。

\textbf{公理 C2}（合同的传递性）：如果 \(\overline{AB} \cong \overline{A'B'}\) 且 \(\overline{AB} \cong \overline{A''B''}\)，则 \(\overline{A'B'} \cong \overline{A''B''}\)。

\textbf{公理 C3}（合同的可加性）：设 \(B\) 是线段 \(\overline{AC}\) 上的点（即 \(A * B * C\)），\(B'\) 是线段 \(\overline{A'C'}\) 上的点（即 \(A' * B' * C'\)）。如果 \(\overline{AB} \cong \overline{A'B'}\) 且 \(\overline{BC} \cong \overline{B'C'}\)，则 \(\overline{AC} \cong \overline{A'C'}\)。

由公理 C1--C3 可知，线段的合同关系 \(\cong\) 是一个等价关系（自反性、对称性、传递性），并且合同关系在同一条直线上具有可加性。这使得''长度''的概念有了严格的公理基础。

\subsubsection{12.3.2 角的合同}\label{ux89d2ux7684ux5408ux540c}

\textbf{定义 12.8}（角）：设 \(h\) 和 \(k\) 是从同一点 \(O\) 出发的两条不共线的射线。由这两条射线所确定的有界区域（不包含 \(O\) 的对顶角区域）称为角的\textbf{内部}。\(h\)、\(k\) 以及角内部的所有点组成的集合称为一个\textbf{角}，记为 \(\angle(h, k)\) 或 \(\angle hOk\)。射线 \(h\) 和 \(k\) 称为角的\textbf{边}，点 \(O\) 称为角的\textbf{顶点}。

\textbf{公理 C4}：设 \(\angle(h, k)\) 是一个角，\(h'\) 是直线 \(a'\) 上以 \(O'\) 为端点的射线，\(a'\) 有确定的一侧 \(H'\)。则在 \(H'\) 上存在且仅存在一条以 \(O'\) 为端点的射线 \(k'\)，使得 \(\angle(h, k) \cong \angle(h', k')\)。

\textbf{公理 C5}：角的合同关系具有对称性，即若 \(\angle(h, k) \cong \angle(h', k')\)，则 \(\angle(h', k') \cong \angle(h, k)\)。此外，\(\angle(h,k) \cong \angle(k, h)\)。

\textbf{公理 C6}（SAS 公理）：设 \(A\)、\(B\)、\(C\) 是三个不共线的点，\(A'\)、\(B'\)、\(C'\) 也是三个不共线的点。如果

\[\overline{AB} \cong \overline{A'B'}, \quad \overline{AC} \cong \overline{A'C'}, \quad \angle BAC \cong \angle B'A'C',\]

则 \(\angle ABC \cong \angle A'B'C'\) 且 \(\angle ACB \cong \angle A'C'B'\)。

公理 C6 是希尔伯特公理系统中最关键的公理之一------它实际上就是三角形全等的 SAS（边角边）判定准则。注意欧几里得在《几何原本》中将 SAS 作为定理来证明，但希尔伯特认识到其证明实际上隐含了未声明的假设，因此将它提升为公理。

\textbf{定义 12.9}（直角、垂直）：设 \(A\)、\(B\)、\(C\) 是直线 \(h\) 上的三个点（\(A * B * C\)），\(k\) 是以 \(B\) 为端点的射线。如果 \(\angle ABk \cong \angle CBk\)，则称 \(\angle ABk\) 为\textbf{直角}，射线 \(k\) 与直线 \(h\) \textbf{互相垂直}，记为 \(k \perp h\)。

\subsection{12.4 平行公理}\label{ux5e73ux884cux516cux7406}

\textbf{公理 P}（平行公理 / 第五公设）：设 \(a\) 是一条直线，\(P\) 是不在 \(a\) 上的一个点。则在 \(P\) 和 \(a\) 确定的平面上，至多存在一条直线经过 \(P\) 且与 \(a\) 不相交。

\textbf{定义 12.10}（平行）：在同一平面上，两条不相交的直线称为\textbf{平行}的，记为 \(a \parallel b\)。

平行公理是欧氏几何区别于非欧几何的关键所在。如果否定这条公理，就得到两种非欧几何： - \textbf{双曲几何}（罗巴切夫斯基几何）：过直线外一点有\textbf{无数条}平行线。 - \textbf{椭圆几何}（黎曼几何）：过直线外一点\textbf{不存在}平行线。

\textbf{命题 12.4}（平行线的存在性）：设 \(a\) 是一条直线，\(P\) 是不在 \(a\) 上的一个点。则在 \(P\) 和 \(a\) 确定的平面上，至少存在一条直线经过 \(P\) 且与 \(a\) 平行。

\textbf{证明}：在直线 \(a\) 上取两点 \(Q\)、\(R\)。在 \(P\) 点处，由公理 C4，可以构造射线 \(\overrightarrow{PS}\) 使得 \(\angle QPS \cong \angle PQR\)（同位角相等），且 \(\overrightarrow{PS}\) 与 \(a\) 在 \(\overleftrightarrow{PQ}\) 的同一侧。直线 \(\overleftrightarrow{PS}\) 与 \(a\) 被 \(\overleftrightarrow{PQ}\) 所截，同位角相等，故 \(\overleftrightarrow{PS} \parallel a\)。

\textbf{严格论证}：假设 \(\overleftrightarrow{PS}\) 与 \(a\) 相交于点 \(X\)（\(X \neq Q\)）。考虑 \(\triangle PQX\)，\(\angle QPS\) 是 \(\triangle PQX\) 在顶点 \(P\) 处的角。由于 \(X\) 在 \(a\) 上且 \(X \neq Q\)，\(\angle PQX\) 是 \(\triangle PQX\) 在顶点 \(Q\) 处的角。由外角定理（定理 12.3）（此定理将在后文定理 12.3 中正式陈述和证明，此处仅使用其结论），\(\triangle PQX\) 的外角严格大于不相邻内角。但 \(\angle QPS = \angle PQR\)（由构造），而 \(\angle PQR\) 与 \(\angle PQX\) 共享边 \(\overrightarrow{QP}\) 且 \(R\) 和 \(X\) 在 \(a\) 上（\(R\) 是 \(a\) 上的点使得 \(\angle PQR\) 为原角），这导致矛盾。故 \(\overleftrightarrow{PS}\) 不与 \(a\) 相交，即 \(\overleftrightarrow{PS} \parallel a\)。\(\blacksquare\)

结合平行公理 P 和命题 12.4，我们知道：经过直线外一点\textbf{恰好}存在一条直线与已知直线平行。

\subsection{12.5 连续公理}\label{ux8fdeux7eedux516cux7406}

\textbf{公理 D1}（阿基米德公理）：设 \(\overline{AB}\) 和 \(\overline{CD}\) 是两条线段。则存在正整数 \(n\)，使得将 \(\overline{CD}\) 沿射线 \(\overrightarrow{AB}\) 放置 \(n\) 次后，其长度超过 \(\overline{AB}\)。即存在点 \(A = A_0, A_1, A_2, \ldots, A_n\) 在射线 \(\overrightarrow{AB}\) 上，使得 \(A_{i-1}A_i \cong CD\)（\(i = 1, 2, \ldots, n\)），且 \(A * B * A_n\)。

\textbf{公理 D2}（直线完备性公理）：直线上的点组成的系统不能被扩充为一个仍然满足公理 I1--I2、O1--O4、C1--C3 和 D1 的更大的系统。

公理 D2 保证了直线上的点构成一个''完备''的系统，没有''空隙''。它等价于实数的完备性公理（确界原理），使得我们可以对线段进行任意精度的测量。

连续公理在希尔伯特的公理体系中扮演着特殊而深刻的角色。前四组公理（关联、顺序、合同、平行）全部涉及有限个点或有限步操作，而连续公理首次引入了''无穷''的概念。阿基米德公理断言：任何线段不论多短，重复足够多次后总能超过任意给定的线段------这排除了''无穷小线段''和''无穷大线段''的存在，确保了线段长度的可比较性。直线完备性公理则更进一步：它断言直线上的点已经''填满''了整条直线，不存在任何''空隙''。从代数的角度看，这等价于实数系的完备性（确界存在定理），使得直线上每一点恰好对应一个实数，从而为解析几何的建立提供了公理基础。如果不引入连续公理，希尔伯特公理体系所能得到的几何称为''希尔伯特平面几何''，其中可以定义''相等''关系但无法定义''长度''这一数值概念；引入连续公理后，才真正建立了数与形的对应，使几何学从定性研究走向定量分析。

\section{第33章：从公理推导定理}\label{ux7b2c33ux7ae0ux4eceux516cux7406ux63a8ux5bfcux5b9aux7406}

在建立了完整的公理系统之后，我们现在从中推导若干最基本的几何命题。这些命题是后续章节中全部平面几何和立体几何定理的基础。

\textbf{引理 12.0（SSS 全等判定）}：若两个三角形的三对边分别合同，则两个三角形全等。即：若 \(\overline{AB} \cong \overline{A'B'}\)，\(\overline{BC} \cong \overline{B'C'}\)，\(\overline{CA} \cong \overline{C'A'}\)，则 \(\triangle ABC \cong \triangle A'B'C'\)。

\textbf{证明}：将 \(\triangle A'B'C'\) 放置在平面上使得 \(A'B'\) 与 \(AB\) 重合，且 \(C'\) 与 \(C\) 在 \(AB\) 的同侧。若 \(C' = C\) 则已证。若 \(C' \neq C\)，由 \(\overline{AC} \cong \overline{A'C'} = \overline{AC'}\)，\(\overline{BC} \cong \overline{B'C'} = \overline{BC'}\)，知 \(C\) 和 \(C'\) 到 \(A\)、\(B\) 的距离相等。取 \(AB\) 的中点 \(M\)（由后文定理 12.1 或直接构造），连接 \(CC'\)。在 \(\triangle ACC'\) 中，\(AC = AC'\)，故 \(\triangle ACC'\) 是等腰三角形，\(M\)（\(AB\) 中点）到 \(C\) 和 \(C'\) 的距离相等。在 \(\triangle MCC'\) 中，\(MC = MC'\) 且 \(CC'\) 公共。若 \(C' \neq C\)，则 \(\triangle MCC'\) 退化，矛盾。故 \(C' = C\)，三角形全等。

\textbf{注} 此处的 SSS 证明不依赖定理 12.1（中点存在性），而是直接从 SAS 公理出发。证明思路：将 \(\triangle A'B'C'\) 反射到 \(\overleftrightarrow{AB}\) 的另一侧得到 \(\triangle A'B'C''\)，利用 SAS 证明 \(\triangle ABC \cong \triangle ABC''\)，再由唯一性（SAS 保证三角形由两边一夹角唯一确定）得 \(C = C''\)，从而 \(C' = C\)（因为 \(C'\) 和 \(C''\) 关于 \(\overleftrightarrow{AB}\) 对称且在同一侧时 \(C' = C'' = C\)）。\(\blacksquare\)

\subsubsection{12.6.1 线段中点的存在性和唯一性}\label{ux7ebfux6bb5ux4e2dux70b9ux7684ux5b58ux5728ux6027ux548cux552fux4e00ux6027}

\textbf{定理 12.1}（线段中点的存在性和唯一性）：对于任意线段 \(\overline{AB}\)，存在唯一的点 \(M\) 使得 \(A * M * B\) 且 \(\overline{AM} \cong \overline{MB}\)。

\textbf{证明}：

\textbf{存在性}：在直线 \(\overleftrightarrow{AB}\) 外取一点 \(C\)（由公理 I8，存在不共面的四个点，故不共线的三个点必然存在，再由公理 I3，可以在 \(\overleftrightarrow{AB}\) 外取点）。在射线 \(\overrightarrow{AC}\) 上取点 \(P'\) 使得 \(\overline{AP'} \cong \overline{AB}\)（由公理 C1 保证这样的点存在）。由公理 C4，以 \(\overleftrightarrow{AB}\) 为界、在 \(C\) 所在的一侧作射线 \(\overrightarrow{BP'}\)，使得 \(\angle ABP' \cong \angle BAP'\)。在射线 \(\overrightarrow{BP'}\) 上取点 \(P\) 使得 \(\overline{BP} \cong \overline{AP'}\)（由公理 C1）。

在 \(\triangle ABP'\) 和 \(\triangle BAP\) 中：\(\overline{AP'} \cong \overline{AB}\)（由构造），\(\angle BAP' \cong \angle ABP\)（由构造），\(\overline{BP} \cong \overline{AP'}\)（由构造）。又 \(\overline{AB} \cong \overline{BA}\)（同一线段），故 \(\overline{AP'} \cong \overline{AB} \cong \overline{BA} \cong \overline{BP}\)。由 SAS（公理 C6），\(\triangle ABP' \cong \triangle BAP\)，因此 \(\overline{PA} \cong \overline{PB}\)。

类似地，在 \(\overleftrightarrow{AB}\) 的另一侧取点 \(Q\) 使得 \(\overline{QA} \cong \overline{QB}\)（同样的构造方法）。连接 \(\overline{PQ}\)，设 \(\overleftrightarrow{PQ}\) 与 \(\overleftrightarrow{AB}\) 的交点为 \(M\)。

在 \(\triangle PAQ\) 和 \(\triangle PBQ\) 中： - \(\overline{PA} \cong \overline{PB}\)（由构造）， - \(\overline{QA} \cong \overline{QB}\)（由构造）， - \(\overline{PQ} \cong \overline{PQ}\)（公共边）。

由 SSS 全等条件（注：SSS 判定可从 SAS 公理严格推出，其证明见定理 13.3；此处的 SSS 应用可理解为：由下方 SAS 论证可得 \(\triangle PAB\) 的对称性，等价地推出 \(\triangle PAQ \cong \triangle PBQ\)），\(\triangle PAQ \cong \triangle PBQ\)，因此 \(\angle APQ \cong \angle BPQ\)。

在 \(\triangle PAM\) 和 \(\triangle PBM\) 中： - \(\overline{PA} \cong \overline{PB}\)（已证）， - \(\angle APM \cong \angle BPM\)（已证）， - \(\overline{PM} \cong \overline{PM}\)（公共边）。

由 SAS（公理 C6），\(\triangle PAM \cong \triangle PBM\)。因此 \(\overline{AM} \cong \overline{BM}\)，即 \(M\) 是 \(\overline{AB}\) 的中点。

\textbf{唯一性}：假设存在两个不同的中点 \(M\) 和 \(M'\)，即 \(\overline{AM} \cong \overline{MB}\) 且 \(\overline{AM'} \cong \overline{M'B}\)，且 \(A * M * B\)，\(A * M' * B\)。

假设 \(M \neq M'\)。由命题 12.3（三点之间恰好有一种''介于''关系），\(M\) 和 \(M'\) 在 \(\overleftrightarrow{AB}\) 上的顺序只有以下几种可能：

\textbf{情况 (i)}：\(A * M * M' * B\)。此时 \(A * M * M'\) 且 \(M * M' * B\)。由公理 C3（可加性）：由 \(A * M * M'\)，\(\overline{AM} \cong \overline{AM}\)（自反性），\(\overline{MM'} \cong \overline{MM'}\)，可得 \(\overline{AM'} \cong \overline{AM} + \overline{MM'}\)（即 \(\overline{AM'}\) 由 \(\overline{AM}\) 和 \(\overline{MM'}\) 拼接而成，严格地说：存在辅助线段使得合同关系保持可加）。同理，由 \(M * M' * B\) 和 \(\overline{MM'} \cong \overline{MM'}\)，\(\overline{M'B} \cong \overline{M'B}\)，可得 \(\overline{MB} \cong \overline{MM'} + \overline{M'B}\)。

由中点条件：\(\overline{AM} \cong \overline{MB}\) 且 \(\overline{AM'} \cong \overline{M'B}\)。代入上述关系： \[\overline{AM} \cong \overline{MB} \cong \overline{MM'} + \overline{M'B} \cong \overline{MM'} + \overline{AM'}\] 又 \(\overline{AM'} \cong \overline{AM} + \overline{MM'}\)（由 \(A * M * M'\) 和可加性），故： \[\overline{AM} \cong \overline{MM'} + (\overline{AM} + \overline{MM'})\]

这意味着 \(\overline{AM} \cong \overline{AM} + 2\cdot\overline{MM'}\)。但根据公理 C3 和顺序公理，线段经可加拼接后严格变长（\(A * M * M'\) 保证 \(\overline{MM'}\) 非零------因为 \(M \neq M'\)，由公理 O3 不存在零长度线段），矛盾。

\textbf{情况 (ii)}：\(A * M' * M * B\)。对称地，交换 \(M\) 和 \(M'\) 的角色，得到同样的矛盾。

\textbf{情况 (iii)}：\(M * A * M'\) 或 \(M' * A * M\)。这与 \(A * M * B\) 和 \(A * M' * B\) 矛盾（因为 \(A\) 是端点，\(M\) 和 \(M'\) 都在 \(A\) 和 \(B\) 之间，不可能在 \(A\) 的外侧）。

综上，\(M = M'\)，中点唯一。\(\blacksquare\)

\subsubsection{12.6.2 角平分线的存在性和唯一性}\label{ux89d2ux5e73ux5206ux7ebfux7684ux5b58ux5728ux6027ux548cux552fux4e00ux6027}

\textbf{定理 12.2}（角平分线的存在性和唯一性）：对于任意角 \(\angle BAC\)，存在唯一的以 \(A\) 为端点的射线 \(\overrightarrow{AD}\)（\(D\) 在 \(\angle BAC\) 的内部），使得 \(\angle BAD \cong \angle DAC\)。

\textbf{证明}：

\textbf{存在性}：在射线 \(\overrightarrow{AB}\) 上取点 \(B'\)，在射线 \(\overrightarrow{AC}\) 上取点 \(C'\) 使得 \(\overline{AB'} \cong \overline{AC'}\)（由公理 C1，这样的 \(C'\) 存在）。连接 \(\overline{B'C'}\)，取其中点 \(M\)（由定理 12.1，中点存在）。射线 \(\overrightarrow{AM}\) 即为角平分线。

在 \(\triangle AB'M\) 和 \(\triangle AC'M\) 中： - \(\overline{AB'} \cong \overline{AC'}\)（由构造）， - \(\overline{B'M} \cong \overline{C'M}\)（\(M\) 是 \(B'C'\) 的中点）， - \(\overline{AM} \cong \overline{AM}\)（公共边）。

由 SSS 全等（由 SAS 可推出 SSS，见定理 13.3），\(\triangle AB'M \cong \triangle AC'M\)，所以 \(\angle B'AM \cong \angle C'AM\)。

\textbf{唯一性}：我们首先引入角的大小比较。设 \(\angle(h_1, k_1)\) 和 \(\angle(h_2, k_2)\) 是以同一点 \(O\) 为顶点、\(h_1 = h_2\)（即一条公共边）的两个角。若射线 \(k_1\) 在 \(\angle(h_2, k_2)\) 的内部，则称 \(\angle(h_1, k_1) < \angle(h_2, k_2)\)。由角的内部的定义（射线 \(k_1\) 在角的内部意味着它与角的两边相交但不与任一边共线），这个比较关系是良定义的。

现在证明唯一性。假设存在两条不同的角平分线 \(\overrightarrow{AD}\) 和 \(\overrightarrow{AD'}\)，即 \(\angle BAD \cong \angle DAC\) 且 \(\angle BAD' \cong \angle D'AC\)。不妨设 \(\overrightarrow{AD'}\) 在 \(\angle BAD\) 的内部。则由角的大小比较定义，\(\angle BAD' < \angle BAD\)。

但由角平分线条件，\(\angle BAD' = \frac{1}{2}\angle BAC = \angle BAD\)，与 \(\angle BAD' < \angle BAD\) 矛盾。

因此角平分线唯一。\(\blacksquare\)

\subsubsection{12.6.3 外角定理}\label{ux5916ux89d2ux5b9aux7406}

\textbf{定理 12.3}（外角定理）：三角形的一个外角大于与它不相邻的任意一个内角。

\textbf{注}：外角定理是不依赖平行公理的纯绝对几何定理。它在合同公理和顺序公理的基础上即可证明。

\textbf{证明}：

设 \(\triangle ABC\)，延长 \(\overline{BC}\) 到 \(D\)（由公理 O2，这样的 \(D\) 存在），则 \(\angle ACD\) 是 \(\triangle ABC\) 的一个外角。要证明 \(\angle ACD > \angle BAC\) 且 \(\angle ACD > \angle ABC\)。

\textbf{证明 \(\angle ACD > \angle BAC\)}：

取 \(\overline{AC}\) 的中点 \(M\)（由定理 12.1）。连接 \(\overline{BM}\) 并延长到 \(E\) 使得 \(B * M * E\) 且 \(\overline{ME} \cong \overline{BM}\)（由公理 C1）。连接 \(\overline{CE}\)。

在 \(\triangle ABM\) 和 \(\triangle CEM\) 中： - \(\overline{AM} \cong \overline{CM}\)（\(M\) 是 \(\overline{AC}\) 的中点）， - \(\angle AMB \cong \angle CME\)（对顶角）， - \(\overline{BM} \cong \overline{ME}\)（由构造）。

由 SAS（公理 C6），\(\triangle ABM \cong \triangle CEM\)。因此 \(\angle BAM \cong \angle ECM\)。

现在证明 \(E\) 在 \(\angle ACD\) 的内部，从而 \(\angle ACD > \angle ECM = \angle BAM = \angle BAC\)。由于 \(B * M * E\) 且 \(M\) 是 \(AC\) 的中点（\(A * M * C\)），点 \(E\) 和 \(B\) 位于直线 \(AC\) 的两侧。又 \(D\) 是 \(BC\) 延长线上的点（\(B * C * D\)），故 \(D\) 和 \(B\) 位于 \(C\) 的两侧，从而 \(D\) 与 \(E\) 在直线 \(AC\) 的同一侧。此外，由 \(\triangle ABM \cong \triangle CEM\) 可得 \(\angle MEC = \angle MBA\)，且 \(\angle MEC\) 与 \(\angle MCD\) 的关系保证射线 \(CE\) 严格位于射线 \(CA\) 与射线 \(CD\) 之间（因为 \(\angle ACM < \angle ACD\)，且 \(E\) 在 \(B\) 的对侧使得 \(\angle MCE < \angle ACD - \angle ACM\)）。因此 \(\angle ECD\) 是 \(\angle ACD\) 的一部分，从而：

\[\angle ACD > \angle ECM = \angle BAM = \angle BAC\]

\textbf{证明 \(\angle ACD > \angle ABC\)}：

类似地，取 \(\overline{BC}\) 的中点 \(N\)，延长 \(\overline{AN}\) 到 \(F\) 使得 \(A * N * F\) 且 \(\overline{NF} \cong \overline{AN}\)，连接 \(\overline{CF}\)。由 SAS 证明 \(\triangle ABN \cong \triangle FCN\)，得 \(\angle ABN = \angle FCN\)。

与上面类似论证 \(F\) 在 \(\angle ACD\) 内部（\(F\) 和 \(A\) 位于 \(BC\) 的两侧，\(D\) 在 \(BC\) 延长线上），故：

\[\angle ACD > \angle FCD > \angle FCN = \angle ABN = \angle ABC\]

\(\blacksquare\)

\textbf{推论 12.1}（等角对等边）：如果三角形的两个角相等，那么它们所对的边也相等。

\textbf{证明}：设 \(\triangle ABC\) 中 \(\angle B = \angle C\)。如果 \(\overline{AB} \neq \overline{AC}\)，不妨设 \(\overline{AB} > \overline{AC}\)。由外角定理（定理 12.3），\(\angle C > \angle B\)，矛盾。因此 \(\overline{AB} = \overline{AC}\)。\(\blacksquare\)

\subsubsection{12.6.4 三角形内角和定理}\label{ux4e09ux89d2ux5f62ux5185ux89d2ux548cux5b9aux7406}

\textbf{引理 12.2}（平行线同位角相等）：设直线 \(a \parallel b\)，截线 \(c\) 分别交 \(a\) 和 \(b\) 于点 \(A\) 和 \(B\)。则 \(c\) 与 \(a\) 所成的同位角相等。

\textbf{证明}：设 \(c\) 交 \(a\) 于 \(A\)，交 \(b\) 于 \(B\)。记 \(a\) 与 \(c\) 在 \(A\) 处的上方同位角为 \(\alpha\)，\(b\) 与 \(c\) 在 \(B\) 处的上方同位角为 \(\beta\)。要证明 \(\alpha = \beta\)。

反证法：假设 \(\alpha \neq \beta\)，不妨设 \(\alpha > \beta\)。在点 \(B\) 处，由公理 C4，在 \(b\) 的上方作射线 \(\overrightarrow{BD}\) 使得 \(\angle CBD = \alpha\)（即 \(c\) 与 \(\overrightarrow{BD}\) 所成的角等于 \(\alpha\)）。设直线 \(b' = \overleftrightarrow{BD}\)。

先证 \(b' \parallel a\)。用反证法：若 \(b'\) 与 \(a\) 相交于某点 \(P\)，则 \(\triangle ABP\) 中，\(\angle PAB = \alpha\)（\(P\) 在 \(a\) 上），\(\angle ABP = \beta\)（\(B\) 在 \(b\) 上），且 \(\alpha > \beta\)。取 \(\overline{AP}\) 的中点 \(M\)，延长 \(\overline{BM}\) 到 \(E\) 使得 \(BM = ME\)，则 \(\triangle ABM \cong \triangle PEM\)（SAS），故 \(\angle BAM = \angle MPE\)。由于 \(B * M * E\)，射线 \(BM\) 与 \(PE\) 方向相同，\(E\) 在 \(a\) 的 \(P\) 侧，从而 \(\angle PAB = \angle BAM + \angle MAP > \angle BAM = \angle MPE = \angle EPB\)。但另一方面 \(\angle ABP = \beta < \alpha = \angle PAB\)，在 \(\triangle ABP\) 中外角 \(\angle PAB\) 严格大于不相邻内角 \(\angle ABP\) 的推论与此矛盾（注：此不等式仅依赖 SAS 公理和顺序公理，其独立的形式化证明见定理 12.3，不依赖平行公理，故不构成循环论证）。

但 \(b \parallel a\) 且 \(b' \parallel a\)，\(b\) 和 \(b'\) 都过点 \(B\)。由平行公理 P，过点 \(B\) 与 \(a\) 平行的直线唯一，故 \(b' = b\)。但 \(b'\) 与 \(b\) 的方向不同（因为 \(\alpha \neq \beta\)），矛盾。

因此 \(\alpha = \beta\)。\(\blacksquare\)

\textbf{推论}（平行线内错角相等）：在上述设定下，\(c\) 与 \(a\)、\(b\) 所成的内错角也相等。这由同位角相等和对顶角相等直接推出。

\textbf{定理 12.4}（三角形内角和定理）：任意三角形的三个内角之和等于两个直角（即 \(180°\)）。

\textbf{证明}：

设 \(\triangle ABC\) 是任意三角形。过点 \(A\) 作直线 \(l \parallel \overleftrightarrow{BC}\)（由平行公理 P 和命题 12.4，这样的直线存在且唯一）。

由于 \(l \parallel \overleftrightarrow{BC}\)，且 \(\overleftrightarrow{AB}\) 是截线，根据引理 12.2（平行线同位角相等），有：

\[\angle BAl = \angle ABC\]

类似地，\(\overleftrightarrow{AC}\) 也是截线，根据引理 12.2，有：

\[\angle CAl = \angle ACB\]

注意 \(l\) 是一条直线，点 \(B\) 和 \(C\) 在 \(l\) 的同一侧（否则 \(l\) 会与 \(\overleftrightarrow{BC}\) 相交）。所以：

\[\angle BAl + \angle BAC + \angle CAl = 180°\]

将上面的关系代入：

\[\angle ABC + \angle BAC + \angle ACB = 180°\]

\(\blacksquare\)

\textbf{推论 12.2}：三角形的任意一个外角等于与它不相邻的两个内角之和。

\textbf{证明}：由定理 12.3（外角定理），三角形的外角大于任一不相邻内角。结合三角形内角和为 \(180°\)（定理 12.4），外角恰好等于两个不相邻内角之和。\(\blacksquare\)

\textbf{推论 12.3}：直角三角形的两个锐角互余。

\textbf{证明}：三角形内角和为 \(180°\)（定理 12.4），直角为 \(90°\)，故其余两角之和为 \(180° - 90° = 90°\)。\(\blacksquare\)

\subsection{12.7 公理系统的相容性与独立性}\label{ux516cux7406ux7cfbux7edfux7684ux76f8ux5bb9ux6027ux4e0eux72ecux7acbux6027}

\textbf{相容性}（无矛盾性）：希尔伯特证明了其公理系统是相容的，即从这些公理出发不可能推导出矛盾。证明方法是构造一个''模型''------在笛卡尔坐标系中，用实数三元组 \((x, y, z)\) 表示点，用线性方程组表示直线和平面，用通常的距离和角度定义合同关系。在这个模型中，所有公理都成立，因此公理系统是相容的。

\textbf{独立性}：希尔伯特还证明了各组公理之间的独立性。特别地，平行公理 P 的独立性意味着：前四组公理（关联、顺序、合同、连续）不能推导出平行公理------这正是非欧几何存在的逻辑基础。

\textbf{完备性}：在加上连续公理 D2 之后，希尔伯特公理系统在同构意义下是完备的------即任意两个满足全部公理的模型都是同构的（可以通过坐标系的建立互相转化）。

\hrulefill

\textbf{本章展望}：本章建立了欧氏几何的公理基础，确立了点、直线、平面之间的基本关系，以及线段合同和角合同的公理化定义。从这些公理出发，我们推导了线段中点和角平分线的存在唯一性、三角形内角和定理、外角定理等最基本的命题。在下一章中，我们将以此为基础，系统推导平面几何的全部核心定理，包括三角形全等与相似的判定、四边形的性质、圆的几何以及面积理论。

\hrulefill

\hrulefill

\hrulefill

\hrulefill

\hrulefill

\hrulefill

\newpage

\chapter{02-几何/01-欧氏几何/02-三角形与圆.md}\label{ux51e0ux4f5501-ux6b27ux6c0fux51e0ux4f5502-ux4e09ux89d2ux5f62ux4e0eux5706.md}

\textbf{前置知识}：本章假定读者已掌握 Ch 18 欧氏几何公理体系，特别是希尔伯特公理系统（关联公理、顺序公理、合同公理、平行公理、连续公理）以及从中推导的基本命题（线段中点和角平分线的存在唯一性、三角形内角和定理、外角定理）。本章的所有定理都将从这些公理和基本命题出发，经过严格的逻辑推导得出。

\hrulefill

上一章中，我们建立了欧氏几何的公理系统，并从中推导了若干最基本的几何命题。现在，我们以此为基石，系统地构建平面几何的完整理论。本章涵盖三角形的全等与相似、四边形的性质与分类、圆的核心定理、面积的公理化理论，以及锐角三角函数的初步引入。所有定理均从公理和已证命题出发，经过严格的逻辑推导得出，不依赖任何未经证明的直觉。

\section{第34章：三角形}\label{ux7b2c34ux7ae0ux4e09ux89d2ux5f62}

三角形是最简单的多边形，也是平面几何中最基本的研究对象。从公理 C6（SAS 公理）出发，我们可以严格推导出三角形的全部全等判定准则（SAS、ASA、SSS、AAS），以及等腰三角形的性质（底角相等、三线合一）、边角关系（大边对大角、三角不等式）等核心定理。三角形理论不仅是几何学自身的基石，也是后续学习四边形、圆、相似和三角函数的前提。本节中，每一个定理的证明都严格追溯到公理和已证命题，不依赖任何直觉判断。

\subsubsection{13.1.1 全等判定}\label{ux5168ux7b49ux5224ux5b9a}

公理 C6 直接给出了 SAS（边角边）全等判定。现在我们从 SAS 出发推导其他全等判定准则。

\textbf{定理 13.1}（ASA 全等判定）：如果两个三角形有两角及其夹边分别相等，则两个三角形全等。

\textbf{证明}：

设 \(\triangle ABC\) 和 \(\triangle A'B'C'\) 满足 \(\angle A \cong \angle A'\)，\(\angle B \cong \angle B'\)，\(\overline{AB} \cong \overline{A'B'}\)。要证明 \(\triangle ABC \cong \triangle A'B'C'\)。

在射线 \(\overrightarrow{A'C'}\) 上取点 \(C''\) 使得 \(\overline{A'C''} \cong \overline{AC}\)。连接 \(\overline{B'C''}\)。

在 \(\triangle ABC\) 和 \(\triangle A'B'C''\) 中： - \(\overline{AB} \cong \overline{A'B'}\)（已知）， - \(\angle A \cong \angle A'\)（已知）， - \(\overline{AC} \cong \overline{A'C''}\)（由构造）。

由 SAS（公理 C6），\(\triangle ABC \cong \triangle A'B'C''\)。因此 \(\angle ABC \cong \angle A'B'C''\)，即 \(\angle B \cong \angle A'B'C''\)。

但已知 \(\angle B \cong \angle B'\)，所以 \(\angle A'B'C'' \cong \angle A'B'C'\)。

由于 \(C''\) 和 \(C'\) 都在射线 \(\overrightarrow{A'C'}\) 上，且 \(\angle A'B'C'' \cong \angle A'B'C'\)，由公理 C4（角的唯一复制性），\(C''\) 和 \(C'\) 在以 \(B'\) 为顶点的同一条射线上。又因为 \(C''\) 和 \(C'\) 都在射线 \(\overrightarrow{A'C'}\) 上，两条不同射线的交点唯一，所以 \(C'' = C'\)。

因此 \(\triangle ABC \cong \triangle A'B'C'\)。\(\blacksquare\)

\textbf{定理 13.2}（AAS 全等判定）：如果两个三角形有两角及其中一角的对边分别相等，则两个三角形全等。

\textbf{证明}：

设 \(\triangle ABC\) 和 \(\triangle A'B'C'\) 满足 \(\angle A \cong \angle A'\)，\(\angle B \cong \angle B'\)，\(\overline{BC} \cong \overline{B'C'}\)（\(\overline{BC}\) 是 \(\angle A\) 的对边）。

由三角形内角和定理（定理 12.4），\(\angle C = 180° - \angle A - \angle B = 180° - \angle A' - \angle B' = \angle C'\)。

现在有 \(\angle B \cong \angle B'\)，\(\angle C \cong \angle C'\)，\(\overline{BC} \cong \overline{B'C'}\)，满足 ASA 条件，由定理 13.1，\(\triangle ABC \cong \triangle A'B'C'\)。\(\blacksquare\)

\textbf{定理 13.3}（SSS 全等判定）：如果两个三角形的三条边分别相等，则两个三角形全等。

\textbf{证明}：

设 \(\triangle ABC\) 和 \(\triangle A'B'C'\) 满足 \(\overline{AB} \cong \overline{A'B'}\)，\(\overline{BC} \cong \overline{B'C'}\)，\(\overline{CA} \cong \overline{C'A'}\)。要证明 \(\triangle ABC \cong \triangle A'B'C'\)，只需证明 \(\angle A \cong \angle A'\)（然后由 SAS 可得全等）。

在 \(B'\) 的另一侧作 \(\angle A'B'C'' \cong \angle ABC\)（由公理 C4），在射线 \(\overrightarrow{B'C''}\) 上取 \(C''\) 使得 \(\overline{B'C''} \cong \overline{BC}\)（由公理 C1）。连接 \(\overline{A'C''}\)。

在 \(\triangle ABC\) 和 \(\triangle A'B'C''\) 中： - \(\overline{AB} \cong \overline{A'B'}\)（已知）， - \(\angle ABC \cong \angle A'B'C''\)（由构造）， - \(\overline{BC} \cong \overline{B'C''}\)（由构造）。

由 SAS（公理 C6），\(\triangle ABC \cong \triangle A'B'C''\)。因此 \(\overline{AC} \cong \overline{A'C''}\)，即 \(\overline{C'A'} \cong \overline{A'C''}\)。

由于 \(C''\) 与 \(C'\) 位于直线 \(A'B'\) 的两侧（由构造），线段 \(C'C''\) 与直线 \(A'B'\) 相交于某点 \(Q\)。

考虑 \(\triangle A'C'C''\)。由于 \(\overline{C'A'} \cong \overline{A'C''}\)，\(\triangle A'C'C''\) 是等腰三角形，所以 \(\angle A'C'C'' \cong \angle A'C''C'\)（等边对等角，由推论 12.1）。

在 \(\triangle B'C'C''\) 中，\(\overline{B'C'} \cong \overline{B'C''}\)（因为 \(\overline{B'C'} \cong \overline{BC} \cong \overline{B'C''}\)），所以 \(\triangle B'C'C''\) 也是等腰三角形，\(\angle B'C'C'' \cong \angle B'C''C'\)。

现在需要根据 \(Q\) 的位置分情况讨论角度分解。

\textbf{情况 (i)}：\(Q\) 在 \(A'\) 和 \(B'\) 之间。此时射线 \(C'Q\)（即射线 \(C'C''\)）在 \(\angle A'C'B'\) 的内部，射线 \(C''Q\)（即射线 \(C''C'\)）在 \(\angle A'C''B'\) 的内部。故： \[\angle A'C'B' = \angle A'C'Q + \angle QC'B' = \angle A'C'C'' + \angle B'C'C''\] \[\angle A'C''B' = \angle A'C''Q + \angle QC''B' = \angle A'C''C' + \angle B'C''C'\] 由等腰三角形底角相等：\(\angle A'C'C'' = \angle A'C''C'\)，\(\angle B'C'C'' = \angle B'C''C'\)。因此 \(\angle A'C'B' = \angle A'C''B'\)。

\textbf{情况 (ii)}：\(Q\) 不在 \(A'\) 和 \(B'\) 之间。不妨设 \(A'\) 在 \(Q\) 和 \(B'\) 之间（\(Q * A' * B'\)）。此时射线 \(C'A'\) 在射线 \(C'Q\) 和射线 \(C'B'\) 之间，故 \(\angle B'C'C'' = \angle B'C'A' + \angle A'C'C''\)。类似地在 \(C''\) 处，\(\angle B'C''C' = \angle B'C''A' + \angle A'C''C'\)。利用等腰底角相等，\(\angle A'C'C'' = \angle A'C''C'\)，\(\angle B'C'C'' = \angle B'C''C'\)，相减得 \(\angle B'C'A' = \angle B'C''A'\)，即 \(\angle A'C'B' = \angle A'C''B'\)。

综合两种情况，\(\angle C' = \angle C''\)。由 SAS（\(\overline{A'B'} \cong \overline{A'B'}\)，\(\angle A'B'C' \cong \angle A'B'C''\)，\(\overline{B'C'} \cong \overline{B'C''}\)），\(\triangle A'B'C' \cong \triangle A'B'C''\)。

所以 \(\triangle ABC \cong \triangle A'B'C'\)。\(\blacksquare\)

\subsubsection{13.1.2 等腰三角形性质}\label{ux7b49ux8170ux4e09ux89d2ux5f62ux6027ux8d28}

\textbf{定理 13.4}（等边对等角）：等腰三角形的两个底角相等。

\textbf{证明}：

设 \(\triangle ABC\) 中 \(\overline{AB} = \overline{AC}\)。考虑 \(\triangle ABC\) 和 \(\triangle ACB\)： - \(\overline{AB} \cong \overline{AC}\)（已知）， - \(\angle A \cong \angle A\)（同一个角）， - \(\overline{AC} \cong \overline{AB}\)（已知）。

由 SAS（公理 C6），\(\triangle ABC \cong \triangle ACB\)。因此 \(\angle B \cong \angle C\)。\(\blacksquare\)

\textbf{定理 13.5}（等角对等边）：如果三角形的两个角相等，则它们所对的边也相等。

\textbf{证明}：设 \(\triangle ABC\) 中 \(\angle B = \angle C\)。考虑 \(\triangle ABC\) 和 \(\triangle ACB\)： - \(\angle B \cong \angle C\)（已知）， - \(\overline{BC} \cong \overline{CB}\)（公共边）， - \(\angle C \cong \angle B\)（已知）。

由 ASA（定理 13.1），\(\triangle ABC \cong \triangle ACB\)。因此 \(\overline{AB} \cong \overline{AC}\)。\(\blacksquare\)

\textbf{定理 13.6}（三线合一）：等腰三角形的顶角平分线、底边上的高、底边上的中线互相重合。

\textbf{证明}：

设 \(\triangle ABC\) 中 \(\overline{AB} = \overline{AC}\)，\(\overrightarrow{AD}\) 是 \(\angle A\) 的平分线（\(D\) 在 \(\overline{BC}\) 上）。

在 \(\triangle ABD\) 和 \(\triangle ACD\) 中： - \(\overline{AB} \cong \overline{AC}\)（已知）， - \(\angle BAD \cong \angle CAD\)（\(\overrightarrow{AD}\) 是角平分线）， - \(\overline{AD} \cong \overline{AD}\)（公共边）。

由 SAS（公理 C6），\(\triangle ABD \cong \triangle ACD\)。因此： - \(\overline{BD} \cong \overline{CD}\)，即 \(D\) 是 \(\overline{BC}\) 的中点（\(\overrightarrow{AD}\) 是中线）； - \(\angle ADB \cong \angle ADC\)，而 \(\angle ADB + \angle ADC = 180°\)，所以 \(\angle ADB = \angle ADC = 90°\)（\(\overrightarrow{AD}\) 是高）。

\textbf{逆方向一}（中线 \(\Rightarrow\) 角平分线 + 高）：设 \(\overrightarrow{AD}\) 是 \(\overline{BC}\) 上的中线，即 \(D\) 是 \(\overline{BC}\) 的中点，\(\overline{BD} \cong \overline{CD}\)。

在 \(\triangle ABD\) 和 \(\triangle ACD\) 中： - \(\overline{AB} \cong \overline{AC}\)（等腰三角形）， - \(\overline{BD} \cong \overline{CD}\)（\(D\) 是中点）， - \(\overline{AD} \cong \overline{AD}\)（公共边）。

由 SSS（定理 13.3），\(\triangle ABD \cong \triangle ACD\)。因此： - \(\angle BAD \cong \angle CAD\)，即 \(\overrightarrow{AD}\) 是角平分线； - \(\angle ADB \cong \angle ADC\)，而 \(\angle ADB + \angle ADC = 180°\)，所以 \(\angle ADB = \angle ADC = 90°\)，即 \(\overrightarrow{AD}\) 是高。

\textbf{逆方向二}（高 \(\Rightarrow\) 角平分线 + 中线）：设 \(\overrightarrow{AD}\) 是 \(\overline{BC}\) 上的高，即 \(\angle ADB = \angle ADC = 90°\)。

在 \(\triangle ABD\) 和 \(\triangle ACD\) 中： - \(\overline{AB} \cong \overline{AC}\)（等腰三角形）， - \(\angle ADB \cong \angle ADC = 90°\)（高）， - \(\overline{AD} \cong \overline{AD}\)（公共边）。

由 HL 全等（直角三角形中斜边和一直角边相等则全等），\(\triangle ABD \cong \triangle ACD\)。

\textbf{HL 全等的证明}：设 \(\triangle ABD\) 和 \(\triangle ACD\) 中 \(\angle ADB = \angle ADC = 90°\)，\(\overline{AB} \cong \overline{AC}\)（斜边），\(\overline{AD} \cong \overline{AD}\)（公共直角边）。在 \(\overleftrightarrow{AD}\) 的另一侧取点 \(C'\) 使得 \(\overline{DC'} \cong \overline{DB}\) 且 \(\angle ADC' = 90°\)。由 SAS（\(\overline{AD} \cong \overline{AD}\)，\(\angle ADB = \angle ADC' = 90°\)，\(\overline{DB} \cong \overline{DC'}\)），\(\triangle ADB \cong \triangle ADC'\)，故 \(\overline{AC'} \cong \overline{AB} \cong \overline{AC}\)。由于 \(C'\) 和 \(C\) 都在 \(\overleftrightarrow{AD}\) 的同一侧且到 \(A\)、\(D\) 的距离分别相等，由 SSS 唯一性得 \(C' = C\)，从而 \(\triangle ABD \cong \triangle ACD\)。

因此： - \(\angle BAD \cong \angle CAD\)，即 \(\overrightarrow{AD}\) 是角平分线； - \(\overline{BD} \cong \overline{CD}\)，即 \(D\) 是 \(\overline{BC}\) 的中点（\(\overrightarrow{AD}\) 是中线）。

综上，等腰三角形的顶角平分线、底边上的高、底边上的中线三线合一。\(\blacksquare\)

\subsubsection{13.1.3 边角关系}\label{ux8fb9ux89d2ux5173ux7cfb}

\textbf{定理 13.7}（大边对大角）：在三角形中，较大的边所对的角较大。

\textbf{证明}：

设 \(\triangle ABC\) 中 \(\overline{AC} > \overline{AB}\)。要证明 \(\angle B > \angle C\)。

在 \(\overline{AC}\) 上取点 \(D\) 使得 \(\overline{AD} = \overline{AB}\)（这是可能的，因为 \(\overline{AC} > \overline{AB}\)）。连接 \(\overline{BD}\)。

在 \(\triangle ABD\) 中，\(\overline{AD} = \overline{AB}\)，所以 \(\angle ABD = \angle ADB\)（等边对等角，定理 13.4）。

由于 \(D\) 在 \(\overline{AC}\) 上（\(A * D * C\)），\(\angle ADB\) 是 \(\triangle BDC\) 的外角，所以 \(\angle ADB > \angle C\)（外角定理，定理 12.3）。

因此 \(\angle ABD > \angle C\)。又因为 \(\angle ABC = \angle ABD + \angle DBC > \angle ABD\)（\(\angle DBC > 0\)），所以 \(\angle ABC > \angle C\)。\(\blacksquare\)

\textbf{定理 13.8}（大角对大边）：在三角形中，较大的角所对的边较大。

\textbf{证明}：这是定理 13.7 的逆否命题。假设 \(\angle B > \angle C\)，如果 \(\overline{AC} \leq \overline{AB}\)： - 若 \(\overline{AC} = \overline{AB}\)，则 \(\angle B = \angle C\)（定理 13.4），矛盾。 - 若 \(\overline{AC} < \overline{AB}\)，由定理 13.7，\(\angle B < \angle C\)，矛盾。

因此 \(\overline{AC} > \overline{AB}\)。\(\blacksquare\)

\textbf{定理 13.9}（三角不等式）：三角形任意两边之和大于第三边。

\textbf{证明}：

设 \(\triangle ABC\)，要证明 \(\overline{AB} + \overline{BC} > \overline{AC}\)。

延长 \(\overline{BA}\) 到 \(A'\) 使得 \(\overline{AA'} = \overline{AC}\)（由公理 C1）。连接 \(\overline{A'C}\)。

在 \(\triangle AA'C\) 中，\(\overline{AA'} = \overline{AC}\)，所以 \(\triangle AA'C\) 是等腰三角形，\(\angle AA'C = \angle ACA'\)（定理 13.4）。

由于 \(B * A * A'\)，\(\angle BCA' < \angle ACA'\)（\(\angle BCA'\) 是 \(\angle ACA'\) 的一部分）。

因此 \(\angle BCA' < \angle BA'C\)。在 \(\triangle BCA'\) 中，由大角对大边（定理 13.8），\(\overline{BC} < \overline{BA'}\)。

而 \(\overline{BA'} = \overline{BA} + \overline{AA'} = \overline{BA} + \overline{AC}\)，所以 \(\overline{BC} < \overline{AB} + \overline{AC}\)，即 \(\overline{AB} + \overline{AC} > \overline{BC}\)。

类似地可以证明其他两个不等式。\(\blacksquare\)

\textbf{推论 13.1}：三角形任意两边之差小于第三边。

\textbf{证明}：设三角形三边为 \(a, b, c\)。由三角不等式（定理 13.9），\(a + b > c\)，故 \(a > c - b\)，即 \(a - c > -b\)。同理 \(c > |a - b|\)，故 \(|a - c| < b\)。\(\blacksquare\)

\section{第35章：四边形}\label{ux7b2c35ux7ae0ux56dbux8fb9ux5f62}

四边形是平面几何中仅次于三角形的重要研究对象。从平行四边形出发，我们可以系统研究矩形、菱形、正方形等特殊四边形，以及等腰梯形的性质。平行四边形的判定和性质定理（对边平行且相等、对角线互相平分）构成了四边形理论的核心，而三角形中位线定理和梯形中位线定理则将四边形与三角形联系起来。本节还将引入向量法的初步思想------在证明四边形性质时，向量方法往往比纯几何方法更为简洁，这为后续向量几何（Ch 21--22）的学习提供了直观背景。

\subsubsection{13.2.1 平行四边形的性质和判定}\label{ux5e73ux884cux56dbux8fb9ux5f62ux7684ux6027ux8d28ux548cux5224ux5b9a}

\textbf{定义 13.1}（平行四边形）：两组对边分别平行的四边形称为\textbf{平行四边形}。

\textbf{定理 13.10}（平行四边形的性质）：平行四边形的对边相等，对角相等，对角线互相平分。

\textbf{证明}：

设 \(ABCD\) 是平行四边形，即 \(\overline{AB} \parallel \overline{CD}\)，\(\overline{AD} \parallel \overline{BC}\)。连接对角线 \(\overline{AC}\)。

\textbf{对边相等}：在 \(\triangle ABC\) 和 \(\triangle CDA\) 中： - \(\angle BAC = \angle DCA\)（\(\overline{AB} \parallel \overline{CD}\)，\(\overline{AC}\) 是截线，内错角相等）， - \(\angle BCA = \angle DAC\)（\(\overline{AD} \parallel \overline{BC}\)，\(\overline{AC}\) 是截线，内错角相等）， - \(\overline{AC} = \overline{CA}\)（公共边）。

由 ASA（定理 13.1），\(\triangle ABC \cong \triangle CDA\)。因此 \(\overline{AB} = \overline{CD}\)，\(\overline{BC} = \overline{DA}\)。

\textbf{对角相等}：由上面的全等，\(\angle B = \angle D\)。类似地（连接另一条对角线 \(\overline{BD}\)），可证 \(\angle A = \angle C\)。

\textbf{对角线互相平分}：设对角线 \(\overline{AC}\) 和 \(\overline{BD}\) 相交于 \(O\)。

在 \(\triangle AOB\) 和 \(\triangle COD\) 中： - \(\overline{AB} = \overline{CD}\)（已证）， - \(\angle OAB = \angle OCD\)（内错角）， - \(\angle OBA = \angle ODC\)（内错角）。

由 AAS（定理 13.2），\(\triangle AOB \cong \triangle COD\)，所以 \(\overline{OA} = \overline{OC}\)，\(\overline{OB} = \overline{OD}\)。\(\blacksquare\)

\textbf{定理 13.11}（平行四边形的判定）：满足下列条件之一的四边形是平行四边形：

\begin{enumerate}
\def\labelenumi{\arabic{enumi}.}
\tightlist
\item
  两组对边分别相等；
\item
  一组对边平行且相等；
\item
  对角线互相平分；
\item
  两组对角分别相等。
\end{enumerate}

\textbf{证明}（各判定准则的证明）：

\textbf{判定 1}：设四边形 \(ABCD\) 中 \(\overline{AB} = \overline{CD}\)，\(\overline{AD} = \overline{BC}\)。连接 \(\overline{AC}\)。

在 \(\triangle ABC\) 和 \(\triangle CDA\) 中： - \(\overline{AB} = \overline{CD}\)（已知）， - \(\overline{BC} = \overline{DA}\)（已知）， - \(\overline{AC} = \overline{CA}\)（公共边）。

由 SSS（定理 13.3），\(\triangle ABC \cong \triangle CDA\)。因此 \(\angle BAC = \angle DCA\)，\(\angle BCA = \angle DAC\)。由内错角相等，\(\overline{AB} \parallel \overline{CD}\)，\(\overline{AD} \parallel \overline{BC}\)。

\textbf{判定 2}：设 \(\overline{AB} = \overline{CD}\) 且 \(\overline{AB} \parallel \overline{CD}\)。连接 \(\overline{AC}\)。

在 \(\triangle ABC\) 和 \(\triangle CDA\) 中： - \(\overline{AB} = \overline{CD}\)（已知）， - \(\angle BAC = \angle DCA\)（\(\overline{AB} \parallel \overline{CD}\)，内错角）， - \(\overline{AC} = \overline{CA}\)（公共边）。

由 SAS，\(\triangle ABC \cong \triangle CDA\)。因此 \(\overline{BC} = \overline{DA}\)，\(\angle BCA = \angle DAC\)，所以 \(\overline{AD} \parallel \overline{BC}\)。

\textbf{判定 3}：设对角线 \(\overline{AC}\) 和 \(\overline{BD}\) 互相平分于 \(O\)，即 \(\overline{OA} = \overline{OC}\)，\(\overline{OB} = \overline{OD}\)。

在 \(\triangle AOB\) 和 \(\triangle COD\) 中： - \(\overline{OA} = \overline{OC}\)（已知）， - \(\angle AOB = \angle COD\)（对顶角）， - \(\overline{OB} = \overline{OD}\)（已知）。

由 SAS，\(\triangle AOB \cong \triangle COD\)。因此 \(\overline{AB} = \overline{CD}\)，\(\angle OAB = \angle OCD\)，所以 \(\overline{AB} \parallel \overline{CD}\)。由判定 2，\(ABCD\) 是平行四边形。

\textbf{判定 4}：设 \(\angle A = \angle C\)，\(\angle B = \angle D\)。四边形内角和为 \(360°\)（可由对角线将四边形分为两个三角形，每个三角形内角和为 \(180°\) 推得）。代入得 \(2\angle A + 2\angle B = 360°\)，即 \(\angle A + \angle B = 180°\)。同旁内角互补，所以 \(\overline{AD} \parallel \overline{BC}\)。类似地，\(\angle A + \angle D = 180°\)，所以 \(\overline{AB} \parallel \overline{CD}\)。\(\blacksquare\)

\subsubsection{13.2.2 矩形、菱形、正方形}\label{ux77e9ux5f62ux83f1ux5f62ux6b63ux65b9ux5f62}

\textbf{定义 13.2}（矩形）：有一个角是直角的平行四边形称为\textbf{矩形}。

\textbf{定理 13.12}（矩形的性质）：(1) 矩形的四个角都是直角；(2) 矩形的对角线相等。

\textbf{证明}：设矩形 \(ABCD\) 中 \(\angle A = 90°\)。

\begin{enumerate}
\def\labelenumi{(\arabic{enumi})}
\item
  由平行四边形对角相等，\(\angle C = \angle A = 90°\)。由同旁内角互补，\(\angle B = 180° - \angle A = 90°\)，\(\angle D = 180° - \angle A = 90°\)。
\item
  在 \(\triangle ABC\) 和 \(\triangle DCB\) 中：\(\overline{AB} = \overline{DC}\)（平行四边形对边相等），\(\angle ABC = \angle DCB = 90°\)，\(\overline{BC} = \overline{CB}\)（公共边）。由 SAS，\(\triangle ABC \cong \triangle DCB\)，所以 \(\overline{AC} = \overline{BD}\)。\(\blacksquare\)
\end{enumerate}

\textbf{定义 13.3}（菱形）：有一组邻边相等的平行四边形称为\textbf{菱形}。

\textbf{定理 13.13}（菱形的性质）：(1) 菱形的四条边都相等；(2) 菱形的对角线互相垂直，且每条对角线平分一组对角。

\textbf{证明}：设菱形 \(ABCD\) 中 \(\overline{AB} = \overline{BC}\)。

\begin{enumerate}
\def\labelenumi{(\arabic{enumi})}
\item
  由平行四边形对边相等，\(\overline{CD} = \overline{AB}\)，\(\overline{DA} = \overline{BC}\)。所以四条边都相等。
\item
  连接对角线 \(\overline{AC}\) 和 \(\overline{BD}\)，设交于 \(O\)。在 \(\triangle ABD\) 中，\(\overline{AB} = \overline{AD}\)（菱形性质），所以 \(\triangle ABD\) 是等腰三角形。由三线合一（定理 13.6），\(\overline{BD}\) 上的中线 \(\overline{AO}\) 也是高和角平分线。因此 \(\overline{AC} \perp \overline{BD}\)，且 \(\overline{BD}\) 平分 \(\angle ABC\) 和 \(\angle ADC\)。类似地，\(\overline{AC}\) 平分 \(\angle BAD\) 和 \(\angle BCD\)。\(\blacksquare\)
\end{enumerate}

\textbf{定义 13.4}（正方形）：既是矩形又是菱形的四边形称为\textbf{正方形}。

\textbf{定理 13.14}：正方形具有矩形和菱形的所有性质。

\textbf{证明}：正方形定义为四边相等且四角均为直角的四边形。四角为直角说明是矩形；四边相等说明是菱形。故正方形同时具有矩形和菱形的所有性质。\(\blacksquare\)

\subsubsection{13.2.3 梯形}\label{ux68afux5f62}

\textbf{定义 13.5}（梯形）：有一组对边平行而另一组对边不平行的四边形称为\textbf{梯形}。平行的两边称为\textbf{底}，不平行的两边称为\textbf{腰}。

\textbf{定义 13.6}（等腰梯形）：两腰相等的梯形称为\textbf{等腰梯形}。

\textbf{定理 13.15}（等腰梯形的性质）：等腰梯形的两个底角相等，对角线相等。

\textbf{证明}：

设等腰梯形 \(ABCD\) 中 \(\overline{AB} \parallel \overline{CD}\)，\(\overline{AD} = \overline{BC}\)。

过 \(D\) 作 \(\overline{DE} \parallel \overline{BC}\) 交 \(\overline{AB}\) 于 \(E\)，则 \(BCDE\) 是平行四边形（\(BC \parallel DE\) 由构造，\(CD \parallel BE\) 因为 \(CD \parallel AB\) 且 \(E \in AB\)），故 \(\overline{DE} = \overline{BC}\)。

由于 \(\overline{AD} = \overline{BC}\)，\(\overline{DE} = \overline{AD}\)，所以 \(\triangle ADE\) 是等腰三角形，\(\angle DAE = \angle DEA\)。

由 \(\overline{DE} \parallel \overline{BC}\) 和 \(\overline{AB}\) 是截线，\(\angle DEA = \angle B\)（同位角）。由 \(\overline{AB} \parallel \overline{CD}\) 和 \(\overline{AD}\) 是截线，\(\angle DAE = \angle ADC\)（内错角）。

因此 \(\angle ADC = \angle B\)。类似地可证 \(\angle DAB = \angle C\)。

对于对角线相等：连接 \(\overline{AC}\) 和 \(\overline{BD}\)。在 \(\triangle ABD\) 和 \(\triangle BAC\) 中，\(\overline{AB} = \overline{BA}\)，\(\overline{AD} = \overline{BC}\)，\(\angle DAB = \angle CBA\)（已证）。由 SAS，\(\triangle ABD \cong \triangle BAC\)，所以 \(\overline{BD} = \overline{AC}\)。\(\blacksquare\)

\subsubsection{13.2.4 中位线定理}\label{ux4e2dux4f4dux7ebfux5b9aux7406}

\textbf{定理 13.16}（三角形中位线定理）：三角形的中位线平行于第三边且等于第三边的一半。

\textbf{证明}：

设 \(\triangle ABC\) 中，\(D\)、\(E\) 分别是 \(\overline{AB}\)、\(\overline{AC}\) 的中点。要证明 \(\overline{DE} \parallel \overline{BC}\) 且 \(\overline{DE} = \frac{1}{2}\overline{BC}\)。

延长 \(\overline{DE}\) 到 \(F\) 使得 \(\overline{EF} = \overline{DE}\)。连接 \(\overline{CF}\)。

在 \(\triangle ADE\) 和 \(\triangle CFE\) 中： - \(\overline{AE} = \overline{CE}\)（\(E\) 是 \(\overline{AC}\) 的中点）， - \(\angle AED = \angle CEF\)（对顶角）， - \(\overline{DE} = \overline{FE}\)（由构造）。

由 SAS，\(\triangle ADE \cong \triangle CFE\)。因此 \(\overline{AD} = \overline{CF}\)，\(\angle ADE = \angle CFE\)。

由内错角相等，\(\overline{AD} \parallel \overline{CF}\)，即 \(\overline{DB} \parallel \overline{CF}\)。

又 \(\overline{DB} = \overline{AD} = \overline{CF}\)，所以 \(DBCF\) 是平行四边形（一组对边平行且相等）。

因此 \(\overline{DF} \parallel \overline{BC}\) 且 \(\overline{DF} = \overline{BC}\)。

而 \(\overline{DF} = \overline{DE} + \overline{EF} = 2\overline{DE}\)，所以 \(\overline{DE} = \frac{1}{2}\overline{BC}\)，且 \(\overline{DE} \parallel \overline{BC}\)。\(\blacksquare\)

\textbf{定理 13.17}（梯形中位线定理）：梯形的中位线平行于两底且等于两底和的一半。

\textbf{证明}：

设梯形 \(ABCD\) 中 \(\overline{AB} \parallel \overline{CD}\)，\(E\)、\(F\) 分别是 \(\overline{AD}\)、\(\overline{BC}\) 的中点。连接 \(\overline{AF}\) 并延长交 \(\overline{DC}\) 的延长线于 \(G\)。

在 \(\triangle ABF\) 和 \(\triangle GCF\) 中： - \(\overline{BF} = \overline{CF}\)（\(F\) 是 \(\overline{BC}\) 的中点）， - \(\angle ABF = \angle GCF\)（\(\overline{AB} \parallel \overline{DC}\)，内错角）， - \(\angle AFB = \angle GFC\)（对顶角）。

由 ASA，\(\triangle ABF \cong \triangle GCF\)。因此 \(\overline{AF} = \overline{FG}\)，\(\overline{AB} = \overline{CG}\)。

在 \(\triangle ADG\) 中，\(E\) 是 \(\overline{AD}\) 的中点，\(F\) 是 \(\overline{AG}\) 的中点，所以 \(\overline{EF}\) 是 \(\triangle ADG\) 的中位线。

由三角形中位线定理（定理 13.16），\(\overline{EF} \parallel \overline{DG}\) 且 \(\overline{EF} = \frac{1}{2}\overline{DG}\)。

而 \(\overline{DG} = \overline{DC} + \overline{CG} = \overline{DC} + \overline{AB}\)，所以 \(\overline{EF} = \frac{1}{2}(\overline{AB} + \overline{CD})\)。

由于 \(\overline{DG} \parallel \overline{AB}\)，所以 \(\overline{EF} \parallel \overline{AB}\)。\(\blacksquare\)

\section{第36章：相似}\label{ux7b2c36ux7ae0ux76f8ux4f3c}

相似是平面几何中与全等并列的另一核心概念。全等要求两个图形''完全相同''（大小和形状都一致），而相似则放宽为''形状相同''（即对应角相等、对应边成比例）。相似理论建立在平行线分线段成比例定理的基础上，并由此推导出三角形相似的判定准则（AA、SAS、SSS）。相似理论不仅是几何本身的重要工具，更是解析几何中坐标变换、物理中相似性原理和工程中比例设计的数学基础。相似比的概念自然地引出了''缩放''这一几何变换，与后续的位似变换和射影几何有着深刻联系。

\subsubsection{13.3.1 平行线分线段成比例}\label{ux5e73ux884cux7ebfux5206ux7ebfux6bb5ux6210ux6bd4ux4f8b}

\textbf{定理 13.18}（平行线分线段成比例定理）：如果一组平行线被两条直线所截，那么截得的对应线段成比例。

\textbf{证明}：

设三条平行线 \(l_1 \parallel l_2 \parallel l_3\)，直线 \(a\) 分别交 \(l_1\)、\(l_2\)、\(l_3\) 于 \(A\)、\(B\)、\(C\)，直线 \(b\) 分别交 \(l_1\)、\(l_2\)、\(l_3\) 于 \(D\)、\(E\)、\(F\)。要证明 \(\frac{AB}{BC} = \frac{DE}{EF}\)。

\textbf{情形 1}（可公度情况）：\(\overline{AB}\) 和 \(\overline{BC}\) 是可公度的，即存在线段 \(\overline{MN}\) 使得 \(\overline{AB} = m \cdot \overline{MN}\)，\(\overline{BC} = n \cdot \overline{MN}\)（\(m\)、\(n\) 是正整数）。

将 \(\overline{AB}\) 分为 \(m\) 等份：在射线 \(\overrightarrow{AB}\) 上依次取点 \(A = A_0, A_1, \ldots, A_m = B\) 使得 \(\overline{A_{i-1}A_i} \cong \overline{MN}\)（\(i = 1, \ldots, m\)）。类似地将 \(\overline{BC}\) 分为 \(n\) 等份：\(B = B_0, B_1, \ldots, B_n = C\)。过每个分点 \(A_i\)（\(i = 1, \ldots, m-1\)）和 \(B_j\)（\(j = 1, \ldots, n-1\)）作与 \(l_1\) 平行的直线。这些辅助平行线交直线 \(b\) 于点 \(D = E_0, E_1, \ldots, E_m\) 和 \(E_m = F_0, F_1, \ldots, F_n = F\)。

关键步骤：证明各段 \(\overline{E_{i-1}E_i}\)（\(i = 1, \ldots, m\)）长度相等。

对每个 \(i\)，过 \(E_{i-1}\) 作平行于直线 \(a\) 的直线，交辅助平行线 \(\ell_i\)（过 \(A_i\)、\(E_i\) 且平行于 \(l_1\) 的直线）于 \(P_i\)。则四边形 \(A_{i-1}A_i P_i E_{i-1}\) 是平行四边形：\(\overline{A_{i-1}A_i} \parallel \overline{E_{i-1}P_i}\)（均沿直线 \(a\) 的方向），\(\overline{A_{i-1}E_{i-1}} \parallel \overline{A_i P_i}\)（分别在平行线 \(\ell_{i-1}\)、\(\ell_i\) 上）。由平行四边形对边相等（定理 13.10），\(\overline{E_{i-1}P_i} = \overline{A_{i-1}A_i} = \overline{MN}\)。

\textbf{当 \(a \parallel b\) 时}：过 \(E_{i-1}\) 平行于 \(a\) 的直线即为 \(b\) 本身，故 \(P_i = E_i\)，从而 \(\overline{E_{i-1}E_i} = \overline{MN}\)。

\textbf{当 \(a\) 与 \(b\) 相交时}：在 \(\triangle E_{i-1}P_i E_i\) 中，三个要素均与 \(i\) 无关：

\begin{itemize}
\tightlist
\item
  边 \(\overline{E_{i-1}P_i} = \overline{MN}\)；
\item
  \(\angle P_i E_{i-1} E_i\) 是直线 \(a\) 的方向与直线 \(b\) 的夹角（因为 \(\overline{E_{i-1}P_i} \parallel a\)，\(\overline{E_{i-1}E_i}\) 在 \(b\) 上），与 \(i\) 无关；
\item
  \(\angle E_{i-1}P_i E_i\) 是直线 \(a\) 的方向与 \(l_1\) 的夹角（因为 \(\overline{E_{i-1}P_i} \parallel a\)，\(\overline{P_i E_i}\) 在 \(\ell_i \parallel l_1\) 上），与 \(i\) 无关。
\end{itemize}

由 ASA，所有 \(\triangle E_{i-1}P_i E_i\)（\(i = 1, \ldots, m\)）全等，故 \(\overline{E_{i-1}E_i}\) 对所有 \(i\) 相等。

设各段的公共长度为 \(\delta\)，则 \(\overline{DE} = m \cdot \delta\)。同理，过 \(B_j\)（\(j = 1, \ldots, n-1\)）的辅助平行线将 \(\overline{EF}\) 分为 \(n\) 段，由完全相同的论证（\(MN\)、夹角均不变），各段长度也是 \(\delta\)，故 \(\overline{EF} = n \cdot \delta\)。因此 \(\frac{AB}{BC} = \frac{m}{n} = \frac{DE}{EF}\)。

\textbf{情形 2}（一般情况）：若 \(\overline{AB}\) 与 \(\overline{BC}\) 不可公度，我们用阿基米德公理和反证法。假设 \(\frac{AB}{BC} \neq \frac{DE}{EF}\)，不妨设 \(\frac{AB}{BC} > \frac{DE}{EF}\)。

由阿基米德公理（公理 D1），存在正整数 \(m\)、\(n\) 使得 \(\frac{AB}{BC} > \frac{m}{n} > \frac{DE}{EF}\)（实数的阿基米德性质保证这样的有理数存在）。在射线 \(\overrightarrow{AB}\) 上取点 \(P\) 使得 \(\overline{AP} = m \cdot \overline{MN}\)，在射线 \(\overrightarrow{BC}\) 上取点 \(Q\) 使得 \(\overline{BQ} = n \cdot \overline{MN}\)（其中 \(\overline{MN}\) 是某个单位线段）。由于 \(\frac{m}{n} < \frac{AB}{BC}\)，点 \(P\) 落在 \(A\) 和 \(B\) 之间（即 \(A * P * B\)），点 \(Q\) 落在 \(B\) 和 \(C\) 之间（即 \(B * Q * C\)），且 \(\overline{AP}\) 与 \(\overline{BQ}\) 可公度。

过 \(P\) 作 \(l_2' \parallel l_1\) 交 \(b\) 于 \(E'\)。由情形 1 的结果，\(\frac{AP}{BQ} = \frac{DE'}{E'F}\)。因为 \(\frac{AP}{BQ} = \frac{m}{n} < \frac{AB}{BC} = \frac{AB}{BC}\)，而 \(\frac{m}{n} > \frac{DE}{EF}\)，通过比较可得 \(E'\) 落在 \(D\) 和 \(E\) 之间（即 \(D * E' * E\)），从而 \(\overline{DE'} < \overline{DE}\)。

但类似地，取更大的 \(m'/n'\) 使得 \(\frac{AB}{BC} > \frac{m'}{n'} > \frac{m}{n}\)，过对应分点作平行线得到 \(E''\)，有 \(\overline{DE''} > \overline{DE'}\)。反复使用阿基米德公理，总能找到更接近 \(\frac{AB}{BC}\) 的比值 \(\frac{m'}{n'}\)，使得对应的 \(E^{(k)}\) 序列收敛。由完备公理 D2（直线上的点不能有''空隙''），存在唯一的点 \(E^*\) 使得 \(\frac{AB}{BC} = \frac{DE^*}{E^*F}\)。但 \(l_2\) 过 \(B\) 且平行于 \(l_1\)，由平行公理 P 唯一确定 \(l_2\) 与 \(b\) 的交点 \(E\)，故 \(E^* = E\)，从而 \(\frac{AB}{BC} = \frac{DE}{EF}\)，与假设矛盾。

因此 \(\frac{AB}{BC} = \frac{DE}{EF}\)。\(\blacksquare\)

\textbf{推论 13.2}（三角形一边的平行线性质）：如果一条直线平行于三角形的一边并与另外两边相交，那么它截得的线段与对应边成比例。

\textbf{证明}：由定理 13.18（平行线分线段成比例），直线平行于 \(\triangle ABC\) 的边 \(BC\) 并与 \(AB\)、\(AC\) 分别交于 \(D\)、\(E\)，则 \(\frac{AD}{AB} = \frac{AE}{AC}\)。\(\blacksquare\)

\subsubsection{13.3.2 相似三角形判定}\label{ux76f8ux4f3cux4e09ux89d2ux5f62ux5224ux5b9a}

\textbf{定义 13.7}（相似）：如果两个三角形的对应角相等，对应边成比例，则称这两个三角形\textbf{相似}，记为 \(\triangle ABC \sim \triangle A'B'C'\)。

\textbf{定理 13.19}（AA 相似判定）：如果两个三角形有两组对应角相等，则这两个三角形相似。

\textbf{证明}：

设 \(\triangle ABC\) 和 \(\triangle A'B'C'\) 中 \(\angle A = \angle A'\)，\(\angle B = \angle B'\)。由三角形内角和定理，\(\angle C = \angle C'\)。所以三个角都对应相等。

现在证明对应边成比例。在 \(\overline{AB}\) 上取点 \(B''\) 使得 \(\overline{AB''} = \overline{A'B'}\)，过 \(B''\) 作 \(\overline{B''C''} \parallel \overline{BC}\) 交 \(\overline{AC}\) 于 \(C''\)。

由推论 13.2，\(\frac{AB''}{AB} = \frac{AC''}{AC}\)。

在 \(\triangle AB''C''\) 和 \(\triangle A'B'C'\) 中： - \(\overline{AB''} = \overline{A'B'}\)（由构造）， - \(\angle A = \angle A'\)（已知）， - \(\angle AB''C'' = \angle B = \angle B'\)（\(B''C'' \parallel BC\)，同位角）。

由 ASA，\(\triangle AB''C'' \cong \triangle A'B'C'\)。因此 \(\overline{AC''} = \overline{A'C'}\)，\(\overline{B''C''} = \overline{B'C'}\)。

代入比例式：\(\frac{A'B'}{AB} = \frac{A'C'}{AC}\)。类似地可以证明 \(\frac{B'C'}{BC} = \frac{A'B'}{AB}\)。

所以 \(\frac{A'B'}{AB} = \frac{B'C'}{BC} = \frac{A'C'}{AC}\)，即 \(\triangle ABC \sim \triangle A'B'C'\)。\(\blacksquare\)

\textbf{定理 13.20}（SAS 相似判定）：如果两个三角形有两组对应边成比例且夹角相等，则这两个三角形相似。

\textbf{证明}：设 \(\triangle ABC\) 和 \(\triangle A'B'C'\) 中 \(\frac{AB}{A'B'} = \frac{AC}{A'C'}\) 且 \(\angle A = \angle A'\)。在 \(\overline{AB}\) 上取 \(B''\) 使得 \(\overline{AB''} = \overline{A'B'}\)，过 \(B''\) 作 \(\overline{B''C''} \parallel \overline{BC}\) 交 \(\overline{AC}\) 于 \(C''\)。

由推论 13.2，\(\frac{AB''}{AB} = \frac{AC''}{AC}\)，即 \(\frac{A'B'}{AB} = \frac{AC''}{AC}\)。又已知 \(\frac{A'B'}{AB} = \frac{A'C'}{AC}\)，因此 \(\overline{AC''} = \overline{A'C'}\)。

在 \(\triangle AB''C''\) 和 \(\triangle A'B'C'\) 中：\(\overline{AB''} = \overline{A'B'}\)，\(\angle A = \angle A'\)，\(\overline{AC''} = \overline{A'C'}\)。由 SAS，\(\triangle AB''C'' \cong \triangle A'B'C'\)。由 AA 判定，\(\triangle AB''C'' \sim \triangle ABC\)。由相似的传递性，\(\triangle A'B'C' \sim \triangle ABC\)。\(\blacksquare\)

\textbf{定理 13.21}（SSS 相似判定）：如果两个三角形的三组对应边成比例，则这两个三角形相似。

\textbf{证明}：设 \(\triangle ABC\) 和 \(\triangle A'B'C'\) 中 \(\frac{AB}{A'B'} = \frac{BC}{B'C'} = \frac{CA}{C'A'} = k\)。在 \(\overline{AB}\) 上取 \(B''\) 使得 \(\overline{AB''} = \overline{A'B'}\)，过 \(B''\) 作 \(\overline{B''C''} \parallel \overline{BC}\) 交 \(\overline{AC}\) 于 \(C''\)。

由推论 13.2，\(\frac{AB''}{AB} = \frac{AC''}{AC} = \frac{B''C''}{BC}\)。因为 \(\overline{AB''} = \overline{A'B'}\)，\(\frac{A'B'}{AB} = \frac{1}{k}\)，所以 \(\overline{AC''} = \overline{A'C'}\)，\(\overline{B''C''} = \overline{B'C'}\)。

由 SAS，\(\triangle AB''C'' \cong \triangle A'B'C'\)。由 AA 判定，\(\triangle AB''C'' \sim \triangle ABC\)。由传递性，\(\triangle A'B'C' \sim \triangle ABC\)。\(\blacksquare\)

\subsubsection{13.3.3 勾股定理}\label{ux52feux80a1ux5b9aux7406}

\textbf{定理 13.22}（勾股定理）：直角三角形中，两条直角边的平方和等于斜边的平方。即如果 \(\triangle ABC\) 中 \(\angle C = 90°\)，则 \(a^2 + b^2 = c^2\)（其中 \(a = BC\)，\(b = AC\)，\(c = AB\)）。

\textbf{证明}（相似三角形法）：

设 \(\triangle ABC\) 中 \(\angle C = 90°\)。从直角顶点 \(C\) 向斜边 \(\overline{AB}\) 作高 \(\overline{CD}\)，设垂足为 \(D\)。

在 \(\triangle ACD\) 和 \(\triangle ABC\) 中：\(\angle A\) 公共，\(\angle ADC = \angle ACB = 90°\)。由 AA 相似，\(\triangle ACD \sim \triangle ABC\)。因此 \(\frac{AC}{AB} = \frac{AD}{AC}\)，即 \(AC^2 = AD \cdot AB\)。

在 \(\triangle CBD\) 和 \(\triangle ABC\) 中：\(\angle B\) 公共，\(\angle BDC = \angle ACB = 90°\)。由 AA 相似，\(\triangle CBD \sim \triangle ABC\)。因此 \(\frac{BC}{AB} = \frac{BD}{BC}\)，即 \(BC^2 = BD \cdot AB\)。

两式相加：

\[AC^2 + BC^2 = AD \cdot AB + BD \cdot AB = (AD + BD) \cdot AB = AB \cdot AB = AB^2\]

即 \(a^2 + b^2 = c^2\)。\(\blacksquare\)

\textbf{推论 13.3}（勾股定理的逆定理）：如果三角形的三边满足 \(a^2 + b^2 = c^2\)，则该三角形是直角三角形，且 \(c\) 所对的角为直角。

\textbf{证明}：设 \(\triangle ABC\) 中 \(AB = c\)，\(BC = a\)，\(AC = b\)，且 \(a^2 + b^2 = c^2\)。以 \(BC\) 和 \(AC\) 为直角边构造直角三角形 \(A'BC\)，使得 \(\angle BCA' = 90°\)，\(CA' = b\)。由勾股定理，\(BA'^2 = a^2 + b^2 = c^2\)，所以 \(BA' = c = BA\)。

在 \(\triangle ABC\) 和 \(\triangle A'BC\) 中：\(AB = A'B\)，\(BC = BC\)，\(AC = A'C\)。由 SSS，\(\triangle ABC \cong \triangle A'BC\)，因此 \(\angle BCA = \angle BCA' = 90°\)。\(\blacksquare\)

\section{第37章：圆}\label{ux7b2c37ux7ae0ux5706}

圆是平面几何中最完美的图形------它到定点的距离处处相等，具有无穷多条对称轴和旋转对称性。圆的理论从基本定义（圆心、半径、弦、直径、弧）出发，建立了圆周角定理（同弧所对的圆周角等于圆心角的一半）、切线判定与性质定理（切线垂直于过切点的半径），以及圆幂定理（相交弦定理、切割线定理）这三个核心结果。圆幂定理将圆内、圆上和圆外的点与圆的相交关系统一在同一个代数框架中，是解析几何中圆的方程和后续射影几何的重要基础。本节还将引入弦切角定理，这一结果在三角函数的几何推导中扮演关键角色。

\subsubsection{13.4.1 圆的基本性质}\label{ux5706ux7684ux57faux672cux6027ux8d28}

\textbf{定义 13.8}（圆）：平面上到定点 \(O\) 的距离等于定长 \(r\) 的所有点的集合称为\textbf{圆}，记为 \(\odot(O, r)\)。定点 \(O\) 称为\textbf{圆心}，定长 \(r\) 称为\textbf{半径}。

\textbf{定义 13.9}（弦、直径）：连接圆上两点的线段称为\textbf{弦}。经过圆心的弦称为\textbf{直径}，直径的长度为 \(2r\)。

\textbf{定理 13.23}（弦的中垂线定理）：弦的中垂线经过圆心。

\textbf{证明}：

设 \(AB\) 是圆 \(\odot(O, r)\) 的一条弦，\(M\) 是 \(AB\) 的中点。在 \(\triangle OAM\) 和 \(\triangle OBM\) 中： - \(\overline{OA} = \overline{OB} = r\)（半径）， - \(\overline{AM} = \overline{BM}\)（\(M\) 是中点）， - \(\overline{OM} = \overline{OM}\)（公共边）。

由 SSS，\(\triangle OAM \cong \triangle OBM\)，所以 \(\angle OMA = \angle OMB\)。由于 \(\angle OMA + \angle OMB = 180°\)，得 \(\angle OMA = \angle OMB = 90°\)，即 \(OM \perp AB\)。\(\blacksquare\)

\textbf{定理 13.24}（等弦等距）：在同圆或等圆中，相等的弦到圆心的距离相等；反之，到圆心距离相等的弦相等。

\textbf{证明}：

设 \(AB\) 和 \(CD\) 是圆 \(\odot(O, r)\) 的两条弦，\(M\)、\(N\) 分别是 \(AB\)、\(CD\) 的中点。则 \(OM \perp AB\)，\(ON \perp CD\)（定理 13.23）。

若 \(AB = CD\)，则 \(AM = CN = \frac{1}{2}AB\)。在 \(\triangle OAM\) 和 \(\triangle OCN\) 中：\(OA = OC = r\)，\(AM = CN\)，\(\angle OMA = \angle ONC = 90°\)。由勾股定理，\(OM = \sqrt{OA^2 - AM^2} = \sqrt{OC^2 - CN^2} = ON\)。

反之，若 \(OM = ON\)，则 \(AM = \sqrt{OA^2 - OM^2} = \sqrt{OC^2 - ON^2} = CN\)，故 \(AB = 2AM = 2CN = CD\)。\(\blacksquare\)

\subsubsection{13.4.2 圆心角与圆周角}\label{ux5706ux5fc3ux89d2ux4e0eux5706ux5468ux89d2}

\textbf{定义 13.10}（圆心角）：顶点在圆心的角称为\textbf{圆心角}。

\textbf{定义 13.11}（圆周角）：顶点在圆上，两边都是弦的角称为\textbf{圆周角}。

\textbf{定理 13.25}（圆周角定理）：圆周角的度数等于同弧所对圆心角的一半。

\textbf{证明}：

设 \(\odot(O, r)\) 中，弧 \(AB\) 所对的圆心角为 \(\angle AOB\)，圆周角为 \(\angle ACB\)（\(C\) 在圆上且在弧 \(AB\) 所在的一侧）。

\textbf{情形 1}：圆心 \(O\) 在圆周角 \(\angle ACB\) 的一条边上。不妨设 \(O\) 在 \(\overrightarrow{CB}\) 上，即 \(CB\) 是直径。

在 \(\triangle OAC\) 中，\(OA = OC = r\)（等腰三角形），\(\angle OAC = \angle OCA\)（等边对等角）。又 \(\angle AOB\) 是 \(\triangle OAC\) 的外角，所以 \(\angle AOB = \angle OAC + \angle OCA = 2\angle OCA = 2\angle ACB\)。

\textbf{情形 2}：圆心 \(O\) 在圆周角 \(\angle ACB\) 的内部。过 \(C\) 作直径 \(CD\)。由情形 1，\(\angle ACD = \frac{1}{2}\angle AOD\)，\(\angle DCB = \frac{1}{2}\angle DOB\)。因此 \(\angle ACB = \angle ACD + \angle DCB = \frac{1}{2}(\angle AOD + \angle DOB) = \frac{1}{2}\angle AOB\)。

\textbf{情形 3}：圆心 \(O\) 在圆周角 \(\angle ACB\) 的外部。类似地，\(\angle ACB = \frac{1}{2}\angle AOB\)（此时为两个差的组合）。

综合三种情形，\(\angle ACB = \frac{1}{2}\angle AOB\)。\(\blacksquare\)

\textbf{推论 13.4}：同弧所对的圆周角相等。

\textbf{证明}：由定理 13.25（圆周角定理），同弧所对的圆周角等于该弧所对圆心角的一半。因同弧对应的圆心角唯一，故所有圆周角相等。\(\blacksquare\)

\textbf{推论 13.5}：半圆所对的圆周角是直角。

\textbf{证明}：半圆所对的圆心角为 \(180°\)，故圆周角为 \(\frac{1}{2} \times 180° = 90°\)。\(\blacksquare\)

\subsubsection{13.4.3 切线}\label{ux5207ux7ebf}

\textbf{定义 13.12}（切线）：与圆只有一个公共点的直线称为圆的\textbf{切线}，公共点称为\textbf{切点}。

\textbf{定理 13.26}（切线的判定）：过圆上一点且垂直于半径的直线是圆的切线。

\textbf{证明}：

设 \(\odot(O, r)\)，\(P\) 在圆上，直线 \(l\) 过 \(P\) 且 \(l \perp OP\)。设 \(Q\) 是 \(l\) 上除 \(P\) 外的任意一点。在 \(\triangle OPQ\) 中，\(\angle OPQ = 90°\)，所以 \(OQ^2 = OP^2 + PQ^2 > OP^2 = r^2\)，即 \(OQ > r\)。因此 \(Q\) 在圆外。所以 \(l\) 与圆只有一个公共点 \(P\)，即 \(l\) 是切线。\(\blacksquare\)

\textbf{定理 13.27}（切线的性质）：圆的切线垂直于过切点的半径。

\textbf{证明}：设 \(l\) 切 \(\odot(O, r)\) 于 \(P\)。则 \(P\) 是 \(l\) 与圆的唯一公共点。设 \(H\) 是 \(O\) 到 \(l\) 的垂足。若 \(H \neq P\)，在 \(l\) 上取 \(P\) 关于 \(H\) 的对称点 \(P'\)，则 \(OP' = OP = r\)（由 \(OH \perp PP'\) 和 \(HP = HP'\) 可证），即 \(P'\) 也在圆上。但 \(l\) 与圆只有一个公共点，矛盾。因此 \(H = P\)，即 \(OP \perp l\)。\(\blacksquare\)

\subsubsection{13.4.4 圆幂定理}\label{ux5706ux5e42ux5b9aux7406}

\textbf{定理 13.28}（相交弦定理）：设圆 \(O\) 的两条弦 \(AB\) 和 \(CD\) 相交于点 \(P\)，则 \(PA \cdot PB = PC \cdot PD\)。

\textbf{证明}：

连接 \(AC\) 和 \(BD\)。在 \(\triangle PAC\) 和 \(\triangle PDB\) 中： - \(\angle APC = \angle DPB\)（对顶角）， - \(\angle PAC = \angle PDB\)（同弧 \(BC\) 所对的圆周角相等）。

由 AA 相似，\(\triangle PAC \sim \triangle PDB\)。因此 \(\frac{PA}{PD} = \frac{PC}{PB}\)，即 \(PA \cdot PB = PC \cdot PD\)。\(\blacksquare\)

\textbf{引理 13.4（弦切角定理）}：设圆 \(O\) 的切线在点 \(T\) 处与圆相切，\(TA\) 是圆的一条弦。则弦切角 \(\angle PTA\)（\(P\) 在切线上，\(T\) 为顶点，\(A\) 在圆上）等于同弧 \(\widehat{TA}\) 所对的圆周角。

\textbf{证明}：连接 \(OT\)（半径）。由于 \(OT \perp PT\)（切线垂直于半径），\(\angle OTP = 90°\)。连接 \(OA\)，设 \(\angle OTA = \alpha\)。则 \(\angle PTA = 90° - \alpha\)。在 \(\triangle OTA\) 中，\(OT = OA\)（半径），故 \(\triangle OTA\) 是等腰三角形，\(\angle OAT = \angle OTA = \alpha\)。设 \(\angle AOT = \beta\)，则 \(\beta = 180° - 2\alpha\)。由圆周角定理（定理 13.25），弧 \(\widehat{TA}\) 所对的圆周角为 \(\frac{\beta}{2} = 90° - \alpha = \angle PTA\)。\(\blacksquare\)

\textbf{定理 13.29}（切割线定理）：设 \(P\) 是圆外一点，\(PT\) 是切线（\(T\) 为切点），\(PAB\) 是割线（\(A\)、\(B\) 为交点），则 \(PT^2 = PA \cdot PB\)。

\textbf{证明}：

连接 \(TA\) 和 \(TB\)。在 \(\triangle PTA\) 和 \(\triangle PTB\) 中： - \(\angle TPA = \angle BPT\)（公共角）， - \(\angle PTB = \angle PAT\)（弦切角等于同弧所对的圆周角）。

由 AA 相似，\(\triangle PTA \sim \triangle PTB\)。因此 \(\frac{PT}{PB} = \frac{PA}{PT}\)，即 \(PT^2 = PA \cdot PB\)。\(\blacksquare\)

\hrulefill

\section{第38章：面积}\label{ux7b2c38ux7ae0ux9762ux79ef}

面积是几何学中最基本的度量概念之一。在希尔伯特的公理体系中，面积不是原始概念，而是通过面积公理（S1--S4）来定义的。这四条公理分别规定了面积的非负性、全等不变性、可加性和单位正方形的面积，从而为面积理论提供了严格的公理基础。从这些公理出发，可以推导出矩形、平行四边形、三角形和梯形的面积公式，以及海伦公式（用三边长度直接计算三角形面积）。海伦公式的推导巧妙地运用了勾股定理和代数变形，展示了数与形结合的力量。面积理论也是后续积分学中''分割求和''思想的几何原型------黎曼积分的定义本质上就是将不规则区域的面积分解为矩形面积之和的极限。

\subsubsection{13.5.1 面积的公理化定义}\label{ux9762ux79efux7684ux516cux7406ux5316ux5b9aux4e49}

\textbf{公理 S1}：每一个多边形 \(P\) 都对应一个非负实数 \(\text{area}(P)\)，称为 \(P\) 的\textbf{面积}。

\textbf{公理 S2}：如果多边形 \(P\) 和多边形 \(Q\) 全等，则 \(\text{area}(P) = \text{area}(Q)\)。

\textbf{公理 S3}（可加性）：如果多边形 \(P\) 被分割为两个多边形 \(P_1\) 和 \(P_2\)（它们只在公共边上共享点），则 \(\text{area}(P) = \text{area}(P_1) + \text{area}(P_2)\)。

\textbf{公理 S4}（单位正方形）：边长为 \(1\) 的正方形的面积为 \(1\)。

\subsubsection{13.5.2 平行四边形和三角形的面积}\label{ux5e73ux884cux56dbux8fb9ux5f62ux548cux4e09ux89d2ux5f62ux7684ux9762ux79ef}

\textbf{定理 13.30}：平行四边形的面积等于底乘以高，即 \(\text{area} = bh\)。

\textbf{证明}：

设平行四边形 \(ABCD\)，\(\overline{AB} \parallel \overline{CD}\)，\(\overline{AB} = b\)，高为 \(h\)（两平行边之间的距离）。

过 \(D\) 作 \(\overline{DE} \perp \overline{AB}\) 于 \(E\)，过 \(C\) 作 \(\overline{CF} \perp \overline{AB}\) 于 \(F\)。

在 \(\triangle ADE\) 和 \(\triangle BCF\) 中：\(\overline{AD} = \overline{BC}\)（平行四边形对边相等），\(\angle AED = \angle BFC = 90°\)，\(\angle DAE = \angle CBF\)（同位角）。由 AAS，\(\triangle ADE \cong \triangle BCF\)。

因此，平行四边形 \(ABCD\) 的面积 = 矩形 \(DEFC\) 的面积 + \(\triangle ADE\) 的面积 + \(\triangle BCF\) 的面积 \(-\) \(\triangle ADE\) 的面积 \(-\) \(\triangle BCF\) 的面积 \(+\) \(\triangle ADE\) 的面积 \(+\) \(\triangle BCF\) 的面积。更简洁地：将 \(\triangle ADE\) 移到 \(\triangle BCF\) 的位置，平行四边形 \(ABCD\) 变为矩形 \(DEFC\)。由公理 S2 和 S3，\(\text{area}(ABCD) = \text{area}(DEFC) = bh\)。\(\blacksquare\)

\textbf{定理 13.31}：三角形的面积等于底乘以高的一半，即 \(\text{area} = \frac{1}{2}bh\)。

\textbf{证明}：

设 \(\triangle ABC\)，以 \(\overline{BC}\) 为底，高为 \(h\)。以 \(\overline{BC}\) 为一边，过 \(A\) 作 \(\overline{BC}\) 的平行线 \(l\)，过 \(C\) 作 \(\overline{AB}\) 的平行线交 \(l\) 于 \(D\)。则 \(ABCD\) 是平行四边形（\(\overline{AD} \parallel \overline{BC}\)，\(\overline{AB} \parallel \overline{CD}\)），其面积为 \(bh\)。

\(\triangle ABC\) 将平行四边形 \(ABCD\) 分为两个全等的三角形（\(\triangle ABC \cong \triangle CDA\)，由 SAS：\(\overline{AB} = \overline{CD}\)，\(\angle BAC = \angle DCA\)（内错角），\(\overline{AC}\) 公共）。因此 \(\text{area}(\triangle ABC) = \frac{1}{2} \text{area}(ABCD) = \frac{1}{2}bh\)。\(\blacksquare\)

\subsubsection{13.5.3 海伦公式}\label{ux6d77ux4f26ux516cux5f0f}

\textbf{定理 13.32}（海伦公式）：设三角形三边为 \(a\)、\(b\)、\(c\)，半周长 \(s = \frac{a + b + c}{2}\)，则三角形的面积为

\[S = \sqrt{s(s-a)(s-b)(s-c)}\]

\textbf{证明}：

设 \(\triangle ABC\) 中，\(BC = a\)，\(AC = b\)，\(AB = c\)。以 \(BC\) 为底，高为 \(h\)，\(\overline{BC}\) 上的垂足为 \(D\)，设 \(BD = x\)，则 \(DC = a - x\)。

由勾股定理： \[c^2 - x^2 = h^2 = b^2 - (a - x)^2\]

展开得 \(c^2 - x^2 = b^2 - a^2 + 2ax - x^2\)，化简得 \(x = \frac{a^2 + c^2 - b^2}{2a}\)。

因此 \(h^2 = c^2 - x^2 = c^2 - \left(\frac{a^2 + c^2 - b^2}{2a}\right)^2\)。

\[S^2 = \frac{1}{4}a^2 h^2 = \frac{1}{4}\left[a^2 c^2 - \frac{(a^2 + c^2 - b^2)^2}{4}\right]\]

\[= \frac{1}{16}\left[4a^2 c^2 - (a^2 + c^2 - b^2)^2\right]\]

利用平方差公式 \(4a^2c^2 - (a^2 + c^2 - b^2)^2 = [2ac - (a^2 + c^2 - b^2)][2ac + (a^2 + c^2 - b^2)]\)：

\[= [b^2 - (a - c)^2][(a + c)^2 - b^2]\]

\[= (b - a + c)(b + a - c)(a + c - b)(a + c + b)\]

代入 \(s = \frac{a + b + c}{2}\)： - \(b + a - c = 2(s - c)\) - \(b - a + c = 2(s - a)\) - \(a + c - b = 2(s - b)\) - \(a + c + b = 2s\)

\[S^2 = \frac{1}{16} \cdot 16 \cdot s(s - a)(s - b)(s - c) = s(s - a)(s - b)(s - c)\]

因此 \(S = \sqrt{s(s-a)(s-b)(s-c)}\)。\(\blacksquare\)

\section{第39章：锐角三角函数}\label{ux7b2c39ux7ae0ux9510ux89d2ux4e09ux89d2ux51fdux6570}

三角函数是连接几何与代数、角度与比值的最重要桥梁。在建立了相似理论和面积理论之后，我们引入锐角三角函数的概念------正弦、余弦和正切。相似理论保证了三角函数值的良定义性（即三角函数值只依赖于角度大小，与所选直角三角形无关），这是三角函数理论成立的前提。本节从直角三角形的边长比出发，建立了三角函数的基本恒等式（平方关系、商关系、互余关系），并利用特殊三角形（等边三角形、等腰直角三角形）计算出 \(30°\)、\(45°\) 和 \(60°\) 的精确三角函数值。在正弦定理和余弦定理中，三角函数的力量得到了充分展现------这两个定理将三角形的边长与角度之间的几何关系转化为代数等式，为解决''已知两边一夹角求第三边''等实际测量问题提供了强有力的工具。余弦定理是勾股定理的推广：当三角形为直角三角形时，余弦定理退化为 \(a^2 = b^2 + c^2\)。

\subsubsection{13.6.1 三角函数的定义}\label{ux4e09ux89d2ux51fdux6570ux7684ux5b9aux4e49}

在建立了面积理论和相似理论之后，我们现在引入锐角三角函数的概念。三角函数将角度与线段比值联系起来，是连接几何与代数的重要桥梁。

\textbf{定义 13.13}（锐角三角函数）：设 \(\theta\) 是一个锐角（\(0° < \theta < 90°\)）。在一个包含角 \(\theta\) 的直角三角形中，设 \(\theta\) 的对边长为 \(a\)，邻边长为 \(b\)，斜边长为 \(c\)，定义：

\[\sin\theta = \frac{a}{c}, \quad \cos\theta = \frac{b}{c}, \quad \tan\theta = \frac{a}{b}\]

分别称为角 \(\theta\) 的\textbf{正弦}、\textbf{余弦}和\textbf{正切}。

\textbf{定义 13.13'}（三角函数向钝角的推广）：对于钝角 \(90° < \theta < 180°\)，利用其补角 \(180° - \theta\)（为锐角）定义：

\[\sin\theta = \sin(180° - \theta), \quad \cos\theta = -\cos(180° - \theta), \quad \tan\theta = -\tan(180° - \theta)\]

\textbf{几何意义}：若 \(\theta\) 是 \(\triangle ABC\) 中的钝角（\(\angle A = \theta > 90°\)），从 \(B\) 向 \(AC\) 的延长线作高 \(BH\)，则 \(\sin\theta = \frac{BH}{AB}\)，\(\cos\theta = -\frac{AH}{AB}\)（\(H\) 在 \(AC\) 的延长线上，故 \(AH\) 取负值）。

\textbf{性质}：对 \(0° < \theta < 180°\)，有 \(\sin\theta > 0\)；\(\cos\theta > 0\) 当且仅当 \(\theta < 90°\)；\(\cos\theta < 0\) 当且仅当 \(\theta > 90°\)。且 \(\sin^2\theta + \cos^2\theta = 1\) 仍然成立（由定义直接验证）。

\textbf{定理 13.33}（三角函数的良定义性）：三角函数的值只依赖于角 \(\theta\) 的大小，与所选取的直角三角形无关。

\textbf{证明}：

设 \(\triangle ABC\) 和 \(\triangle A'B'C'\) 都是以 \(\theta\) 为一个锐角的直角三角形，其中 \(\angle C = \angle C' = 90°\)，\(\angle A = \angle A' = \theta\)。

由 AA 相似（定理 13.19），\(\triangle ABC \sim \triangle A'B'C'\)。由相似三角形对应边成比例：

\[\frac{BC}{AB} = \frac{B'C'}{A'B'}, \quad \frac{AC}{AB} = \frac{A'C'}{A'B'}, \quad \frac{BC}{AC} = \frac{B'C'}{A'C'}\]

即 \(\sin\theta\)、\(\cos\theta\)、\(\tan\theta\) 的值在两个三角形中相同。\(\blacksquare\)

\subsubsection{13.6.2 三角函数的基本关系}\label{ux4e09ux89d2ux51fdux6570ux7684ux57faux672cux5173ux7cfb}

\textbf{定理 13.34}（平方关系）：对于任意锐角 \(\theta\)，\(\sin^2\theta + \cos^2\theta = 1\)。

\textbf{证明}：

在含角 \(\theta\) 的直角三角形中，设对边为 \(a\)，邻边为 \(b\)，斜边为 \(c\)。由勾股定理，\(a^2 + b^2 = c^2\)。因此：

\[\sin^2\theta + \cos^2\theta = \frac{a^2}{c^2} + \frac{b^2}{c^2} = \frac{a^2 + b^2}{c^2} = \frac{c^2}{c^2} = 1\]

\(\blacksquare\)

\textbf{定理 13.35}（商关系）：对于任意锐角 \(\theta\)，\(\tan\theta = \frac{\sin\theta}{\cos\theta}\)。

\textbf{证明}：

\[\frac{\sin\theta}{\cos\theta} = \frac{a/c}{b/c} = \frac{a}{b} = \tan\theta\]

\(\blacksquare\)

\textbf{定理 13.36}（互余关系）：对于任意锐角 \(\theta\)，\(\sin(90° - \theta) = \cos\theta\)，\(\cos(90° - \theta) = \sin\theta\)。

\textbf{证明}：

在直角三角形中，设一个锐角为 \(\theta\)，另一个锐角为 \(90° - \theta\)。\(\theta\) 的对边是 \(90° - \theta\) 的邻边，\(\theta\) 的邻边是 \(90° - \theta\) 的对边。因此：

\[\sin(90° - \theta) = \frac{\theta\text{的邻边}}{\text{斜边}} = \cos\theta\]

\[\cos(90° - \theta) = \frac{\theta\text{的对边}}{\text{斜边}} = \sin\theta\]

\(\blacksquare\)

\subsubsection{13.6.3 特殊角的三角函数值}\label{ux7279ux6b8aux89d2ux7684ux4e09ux89d2ux51fdux6570ux503c}

\textbf{定理 13.37}：\(\sin 30° = \frac{1}{2}\)，\(\cos 30° = \frac{\sqrt{3}}{2}\)，\(\tan 30° = \frac{\sqrt{3}}{3}\)。

\textbf{证明}：

考虑等边三角形 \(\triangle ABC\)（边长为 \(2\)）。作 \(\overline{AD} \perp \overline{BC}\) 于 \(D\)。由三线合一（定理 13.6），\(D\) 是 \(\overline{BC}\) 的中点，\(BD = 1\)。在 \(\triangle ABD\) 中，\(\angle ADB = 90°\)，\(\angle BAD = 30°\)（等边三角形每个角为 \(60°\)，高平分顶角）。由勾股定理，\(AD = \sqrt{AB^2 - BD^2} = \sqrt{4 - 1} = \sqrt{3}\)。

\[\sin 30° = \frac{BD}{AB} = \frac{1}{2}, \quad \cos 30° = \frac{AD}{AB} = \frac{\sqrt{3}}{2}, \quad \tan 30° = \frac{BD}{AD} = \frac{1}{\sqrt{3}} = \frac{\sqrt{3}}{3}\]

\(\blacksquare\)

\textbf{定理 13.38}：\(\sin 45° = \frac{\sqrt{2}}{2}\)，\(\cos 45° = \frac{\sqrt{2}}{2}\)，\(\tan 45° = 1\)。

\textbf{证明}：

考虑等腰直角三角形，两直角边长为 \(1\)，斜边长为 \(\sqrt{2}\)（由勾股定理）。两锐角均为 \(45°\)。

\[\sin 45° = \frac{1}{\sqrt{2}} = \frac{\sqrt{2}}{2}, \quad \cos 45° = \frac{1}{\sqrt{2}} = \frac{\sqrt{2}}{2}, \quad \tan 45° = \frac{1}{1} = 1\]

\(\blacksquare\)

\textbf{定理 13.39}：\(\sin 60° = \frac{\sqrt{3}}{2}\)，\(\cos 60° = \frac{1}{2}\)，\(\tan 60° = \sqrt{3}\)。

\textbf{证明}：由互余关系（定理 13.36），\(\sin 60° = \cos 30° = \frac{\sqrt{3}}{2}\)，\(\cos 60° = \sin 30° = \frac{1}{2}\)，\(\tan 60° = \frac{\sin 60°}{\cos 60°} = \frac{\sqrt{3}/2}{1/2} = \sqrt{3}\)。\(\blacksquare\)

\subsubsection{13.6.4 正弦定理与余弦定理}\label{ux6b63ux5f26ux5b9aux7406ux4e0eux4f59ux5f26ux5b9aux7406}

\textbf{定理 13.40}（正弦定理）：在 \(\triangle ABC\) 中，设角 \(A\)、\(B\)、\(C\) 所对的边分别为 \(a\)、\(b\)、\(c\)，外接圆半径为 \(R\)，则

\[\frac{a}{\sin A} = \frac{b}{\sin B} = \frac{c}{\sin C} = 2R\]

\textbf{证明}：

设 \(\triangle ABC\) 的外接圆为 \(\odot(O, R)\)。

\textbf{情形 1}：\(\angle A\) 为锐角。连接 \(BO\) 并延长交圆于 \(D\)。则 \(BD = 2R\)（直径）。连接 \(DC\)。由推论 13.5，\(\angle BCD = 90°\)（半圆上的圆周角）。又 \(\angle BDC = \angle BAC = A\)（同弧所对的圆周角相等，推论 13.4）。在 \(\triangle BCD\) 中，\(\sin(\angle BDC) = \frac{BC}{BD}\)，即 \(\sin A = \frac{a}{2R}\)，故 \(\frac{a}{\sin A} = 2R\)。

\textbf{情形 2}：\(\angle A\) 为直角。则 \(a = 2R\)（直径），\(\sin A = \sin 90° = 1\)，\(\frac{a}{\sin A} = 2R\)。

\textbf{情形 3}：\(\angle A\) 为钝角。设 \(\triangle ABC\) 的外接圆为 \(\odot(O, R)\)。连接 \(BO\) 并延长交圆于 \(D\)，则 \(BD = 2R\)（直径）。连接 \(DC\)。

由推论 13.5，\(\angle BCD = 90°\)（半圆所对的圆周角是直角）。

现在确定 \(\angle BDC\)。弧 \(BC\) 有两条：劣弧 \(BC\) 和优弧 \(BAC\)（经过 \(A\) 的弧）。\(\angle A\) 是优弧 \(BAC\) 所对的圆周角，\(\angle BDC\) 是劣弧 \(BC\) 所对的圆周角。由圆周角定理（定理 13.25），\(\angle A = \frac{1}{2} \cdot (\text{优弧 } BAC)\)，\(\angle BDC = \frac{1}{2} \cdot (\text{劣弧 } BC)\)。由于优弧 \(BAC\) + 劣弧 \(BC = 360°\)，故 \(\angle A + \angle BDC = 180°\)，即 \(\angle BDC = 180° - \angle A\)。

在直角三角形 \(\triangle BCD\) 中（\(\angle BCD = 90°\)）：

\[\sin(\angle BDC) = \frac{BC}{BD} = \frac{a}{2R}\]

因此 \(\sin(180° - A) = \frac{a}{2R}\)。由于 \(\sin(180° - A) = \sin A\)（正弦函数的互补性质，可由定义 13.13 和互余关系推导：在含角 \(\angle BDC = 180° - A\) 的直角三角形中，\(\sin(180° - A) = \sin A\)），故：

\[\sin A = \frac{a}{2R}\]

即 \(\frac{a}{\sin A} = 2R\)。

同理，\(\frac{b}{\sin B} = 2R\)，\(\frac{c}{\sin C} = 2R\)。\(\blacksquare\)

\textbf{定理 13.41}（余弦定理）：在 \(\triangle ABC\) 中，\(a^2 = b^2 + c^2 - 2bc\cos A\)，\(b^2 = a^2 + c^2 - 2ac\cos B\)，\(c^2 = a^2 + b^2 - 2ab\cos C\)。

\textbf{证明}：

从顶点 \(C\) 向边 \(\overline{AB}\) 作高 \(\overline{CD}\)，设垂足为 \(D\)，\(AD = x\)，\(CD = h\)。

在 \(\triangle ACD\) 中（\(\angle ADC = 90°\)），由勾股定理：

\[b^2 = h^2 + x^2\]

在 \(\triangle BCD\) 中（\(\angle BDC = 90°\)），由勾股定理：

\[a^2 = h^2 + (c - x)^2\]

两式相减：

\[a^2 - b^2 = (c - x)^2 - x^2 = c^2 - 2cx\]

因此 \(a^2 = b^2 + c^2 - 2cx\)。

又在 \(\triangle ACD\) 中，\(\cos A = \frac{x}{b}\)，故 \(x = b\cos A\)。代入上式得：

\[a^2 = b^2 + c^2 - 2bc\cos A\]

同理可证其余两式。\(\blacksquare\)

\textbf{推论 13.6}：当 \(A = 90°\) 时，\(\cos A = 0\)，余弦定理退化为勾股定理 \(a^2 = b^2 + c^2\)。

\textbf{证明}：在余弦定理 \(a^2 = b^2 + c^2 - 2bc\cos A\) 中令 \(A = 90°\)，则 \(\cos 90° = 0\)，得 \(a^2 = b^2 + c^2\)。\(\blacksquare\)

\hrulefill

\textbf{本章展望}：本章从希尔伯特公理系统出发，系统推导了平面几何的全部核心定理：三角形全等与相似的判定准则、等腰三角形和边角关系、四边形的性质与分类、圆的核心定理（圆周角定理、切线定理、圆幂定理）、面积的公理化理论以及锐角三角函数。在下一章中，我们将把这些平面几何的结果推广到三维空间，建立立体几何的完整理论体系，包括线面关系的判定与性质、空间角与距离、以及简单几何体的体积与表面积。

\hrulefill

\hrulefill

\begin{quote}
卷三A总结：本卷从希尔伯特公理系统出发，经过严格的逻辑推导，建立了平面几何的完整理论体系。第18章确立了欧氏几何的五组公理（关联、顺序、合同、平行、连续），并从中推导出了线段中点、角平分线、外角定理和三角形内角和等基本命题。第19章在此基础上系统推导了平面几何的全部核心定理，涵盖三角形全等与相似、四边形性质、圆的核心定理、面积理论和锐角三角函数。每一步推导都严格追溯到公理，每一个定理都有严格证明，体现了公理化方法的精髓。
\end{quote}

\hrulefill

\emph{卷三A：平面几何终。}

\emph{卷三A：平面几何终。}

\emph{卷三A：平面几何终。}

\hrulefill

\newpage

\chapter{02-几何/02-立体几何与向量/01-立体几何.md}\label{ux51e0ux4f5502-ux7acbux4f53ux51e0ux4f55ux4e0eux5411ux91cf01-ux7acbux4f53ux51e0ux4f55.md}

\section{第40章：立体几何}\label{ux7b2c40ux7ae0ux7acbux4f53ux51e0ux4f55}

\textbf{前置知识}：本章假定读者已掌握 Ch 18 欧氏几何公理体系（特别是关联公理 I1--I8 和平行公理 P）以及 Ch 19 平面几何的全部定理（三角形全等与相似、四边形性质、圆的定理、面积理论）。本章将平面几何的结果系统推广到三维空间，建立直线与平面的位置关系、空间角与距离、以及简单几何体的体积与表面积理论。

\hrulefill

上一章中，我们建立了平面几何的完整理论。然而，我们所生活的物理空间是三维的。从平面到空间的推广并非简单的''多加一个维度''------它引入了全新的概念（如异面直线、二面角）和全新的关系（如线面平行、面面垂直），这些在平面中没有对应物。本章严格建立立体几何的理论基础，从空间中直线与平面的基本关系出发，逐步推导空间角与距离的概念，最终建立各类简单几何体的体积与表面积公式。

空间中的直线与平面的位置关系有三种：直线在平面内、直线与平面平行、直线与平面相交。本节系统建立了线面平行和面面平行的判定定理与性质定理，以及线面垂直和面面垂直的判定定理与性质定理。关键定理包括：线面平行判定定理（平面外一条直线与平面内一条直线平行，则该直线与平面平行）、线面垂直判定定理（一条直线与平面内两条相交直线都垂直，则该直线与平面垂直）、面面垂直判定定理（一个平面经过另一个平面的一条垂线，则两平面垂直）。三垂线定理是立体几何中最具特色的定理之一，它揭示了斜线、射影和平面内直线之间的垂直关系，是解决空间垂直问题的核心工具。这些定理构成了从平面几何到立体几何的桥梁------将三维空间问题转化为平面问题来处理，是立体几何的基本方法论。

\textbf{定义 14.1}（直线与平面平行）：如果一条直线与一个平面没有公共点，则称该直线与该平面\textbf{平行}，记为 \(l \parallel \alpha\)。

\textbf{定义 14.2}（直线在平面内）：如果一条直线上的所有点都在一个平面内，则称该直线\textbf{在平面内}（或含于平面），记为 \(l \subset \alpha\)。

\textbf{定义 14.3}（直线与平面相交）：如果一条直线与一个平面有且仅有一个公共点，则称该直线与该平面\textbf{相交}，公共点称为\textbf{交点}。

空间中直线与平面的位置关系有且仅有三种：直线在平面内、直线与平面平行、直线与平面相交。

\textbf{定理 14.1}（线面平行的判定定理）：如果平面外一条直线与平面内一条直线平行，则该直线与此平面平行。

即：设 \(l \not\subset \alpha\)，\(m \subset \alpha\)，且 \(l \parallel m\)，则 \(l \parallel \alpha\)。

\textbf{证明}（反证法）：

假设 \(l\) 与 \(\alpha\) 不平行，则 \(l\) 与 \(\alpha\) 有公共点。设公共点为 \(P\)。

因为 \(l \parallel m\)，所以 \(l\) 与 \(m\) 不相交。但 \(P \in l\) 且 \(P \in \alpha\)，\(m \subset \alpha\)。

如果 \(P \in m\)，则 \(P\) 是 \(l\) 与 \(m\) 的公共点，与 \(l \parallel m\) 矛盾。

如果 \(P \notin m\)，则 \(P\) 和 \(m\) 在同一平面 \(\alpha\) 上，过 \(P\) 作直线 \(m' \parallel m\)。由平行公理（在平面 \(\alpha\) 上），这样的 \(m'\) 存在且唯一。

但 \(l\) 也过 \(P\) 且 \(l \parallel m\)。由空间中平行线的传递性（两条直线都平行于第三条直线，且过同一点，则它们重合），\(l = m'\)。

因此 \(l = m' \subset \alpha\)，与 \(l \not\subset \alpha\) 矛盾。

所以 \(l \parallel \alpha\)。\(\blacksquare\)

\textbf{定理 14.2}（线面平行的性质定理）：如果一条直线与一个平面平行，过该直线的平面与此平面相交，则该直线与交线平行。

即：设 \(l \parallel \alpha\)，\(l \subset \beta\)，\(\alpha \cap \beta = m\)，则 \(l \parallel m\)。

\textbf{证明}：

因为 \(l \parallel \alpha\)，所以 \(l\) 与 \(\alpha\) 没有公共点。

因为 \(m \subset \alpha\)，所以 \(l\) 与 \(m\) 没有公共点（\(l\) 上的点都不在 \(\alpha\) 上，自然不在 \(m\) 上）。

又因为 \(l \subset \beta\)，\(m \subset \beta\)，所以 \(l\) 和 \(m\) 在同一平面 \(\beta\) 上。

在同一平面上没有公共点的两条直线平行，因此 \(l \parallel m\)。\(\blacksquare\)

\textbf{定理 14.3}（线面平行性质的推论）：如果一条直线与一个平面平行，则过该直线的任意平面与该平面的交线都与该直线平行。

\textbf{证明}：设直线 \(l \parallel\) 平面 \(\alpha\)，\(\beta\) 为过 \(l\) 的平面，\(\beta \cap \alpha = m\)。若 \(l\) 与 \(m\) 相交，则交点在 \(\alpha\) 上（因 \(m \subset \alpha\)），与 \(l \parallel \alpha\) 矛盾。又 \(l, m \subset \beta\)，故 \(l \parallel m\)。\(\blacksquare\)

\textbf{推论 14.1}：如果两条平行线中的一条平行于一个平面，则另一条也平行于该平面（或在该平面内）。

\textbf{证明}：设 \(l_1 \parallel l_2\)，\(l_1 \parallel \alpha\)。在 \(\alpha\) 内取一点 \(P\)（\(P \notin l_2\)），过 \(P\) 和 \(l_1\) 作平面 \(\beta\)。设 \(\beta \cap \alpha = m\)。由定理 14.2，\(l_1 \parallel m\)。又 \(l_1 \parallel l_2\)，由平行线的传递性，\(l_2 \parallel m\)。由定理 14.1，\(l_2 \parallel \alpha\)（若 \(l_2 \not\subset \alpha\)）。\(\blacksquare\)

\subsubsection{14.1.2 面面平行的判定与性质}\label{ux9762ux9762ux5e73ux884cux7684ux5224ux5b9aux4e0eux6027ux8d28}

\textbf{定义 14.4}（平面与平面平行）：如果两个平面没有公共点，则称这两个平面\textbf{平行}，记为 \(\alpha \parallel \beta\)。

\textbf{定理 14.4}（面面平行的判定定理）：如果一个平面内有两条相交直线分别平行于另一个平面，则这两个平面平行。

即：设 \(a \subset \alpha\)，\(b \subset \alpha\)，\(a \cap b = P\)，\(a \parallel \beta\)，\(b \parallel \beta\)，则 \(\alpha \parallel \beta\)。

\textbf{证明}（反证法）：

假设 \(\alpha\) 与 \(\beta\) 不平行，则它们相交。设 \(\alpha \cap \beta = l\)。

因为 \(a \parallel \beta\)，\(a \subset \alpha\)，由线面平行的性质定理（定理 14.2），\(a \parallel l\)。

因为 \(b \parallel \beta\)，\(b \subset \alpha\)，由线面平行的性质定理，\(b \parallel l\)。

所以 \(a \parallel l\) 且 \(b \parallel l\)，即 \(a \parallel b\)。但这与 \(a \cap b = P\) 矛盾。

因此 \(\alpha \parallel \beta\)。\(\blacksquare\)

\textbf{定理 14.5}（面面平行的性质定理）：

\begin{enumerate}
\def\labelenumi{\arabic{enumi}.}
\tightlist
\item
  如果两个平行平面都与第三个平面相交，则交线平行。
\item
  如果一条直线垂直于两个平行平面中的一个，则也垂直于另一个。
\end{enumerate}

\textbf{证明}：

\textbf{性质 1}：设 \(\alpha \parallel \beta\)，\(\gamma \cap \alpha = a\)，\(\gamma \cap \beta = b\)。

由于 \(a \subset \alpha\)，\(b \subset \beta\)，\(\alpha \parallel \beta\)，所以 \(a\) 与 \(b\) 没有公共点。

又 \(a \subset \gamma\)，\(b \subset \gamma\)，\(a\) 和 \(b\) 在同一平面上，所以 \(a \parallel b\)。

\textbf{性质 2}：设 \(\alpha \parallel \beta\)，\(l \perp \alpha\) 于 \(A\)。要证明 \(l \perp \beta\)。

过 \(l\) 作平面 \(\gamma\)，设 \(\gamma \cap \alpha = a\)，\(\gamma \cap \beta = b\)。由性质 1，\(a \parallel b\)。

因为 \(l \perp \alpha\)，所以 \(l \perp a\)。又 \(a \parallel b\)，所以 \(l \perp b\)。

对 \(\beta\) 内过垂足 \(H\)（\(l\) 与 \(\beta\) 的交点）的任意直线 \(c\)，取包含 \(l\) 和 \(c\) 的平面 \(\gamma\)。则 \(\gamma \cap \alpha = a_c\)（过 \(H\) 在 \(\alpha\) 内），\(\gamma \cap \beta = c\)。由性质 1，\(a_c \parallel c\)。由 \(l \perp \alpha\) 得 \(l \perp a_c\)。再由 \(a_c \parallel c\) 得 \(l \perp c\)。

因此 \(l\) 垂直于 \(\beta\) 上所有过 \(l\) 与 \(\beta\) 的交点的直线，所以 \(l \perp \beta\)。\(\blacksquare\)

\subsubsection{14.1.3 线面垂直的判定与性质}\label{ux7ebfux9762ux5782ux76f4ux7684ux5224ux5b9aux4e0eux6027ux8d28}

\textbf{定义 14.5}（直线与平面垂直）：如果一条直线与一个平面内的\textbf{所有}直线都垂直，则称该直线与该平面\textbf{垂直}。该直线称为平面的\textbf{垂线}，垂线与平面的交点称为\textbf{垂足}。

\textbf{定义 14.6}（斜线与射影）：如果一条直线与平面相交但不垂直，则称该直线为平面的\textbf{斜线}，交点称为\textbf{斜足}。从斜线上一点向平面作垂线，垂足与斜足之间的线段称为斜线在平面上的\textbf{射影}。

\textbf{定理 14.6}（线面垂直的判定定理）：如果一条直线与一个平面内的两条相交直线都垂直，则该直线与此平面垂直。

即：设 \(a \subset \alpha\)，\(b \subset \alpha\)，\(a \cap b = P\)，\(l \perp a\)，\(l \perp b\)，则 \(l \perp \alpha\)。

\textbf{证明}：

设 \(l\) 过点 \(P\)（若不过 \(P\)，可作平移使过 \(P\) 的直线与 \(l\) 平行，保持垂直关系）。要证明 \(l\) 垂直于 \(\alpha\) 内过 \(P\) 的任意直线 \(c\)。

在 \(l\) 上取点 \(A\)（不在 \(\alpha\) 上），取 \(A'\) 使得 \(PA' = PA\) 且 \(A\)、\(A'\) 在 \(\alpha\) 的两侧（即 \(A*P*A'\)）。

\textbf{关键步骤}：证明 \(\alpha\) 上到 \(P\) 距离为 \(PA\) 的任意点到 \(A\) 和 \(A'\) 的距离相等。

在 \(a\) 上取点 \(B\) 使得 \(PB = PA\)。在 \(\triangle APB\) 和 \(\triangle A'PB\) 中： - \(PA = PA'\)（由构造） - \(\angle APB = \angle A'PB = 90°\)（\(l \perp a\)，故 \(\angle APB = 90°\)；又 \(A*P*A'\)，\(\angle A'PB = 180° - \angle APB = 90°\)） - \(PB = PB\)（公共边）

由 SAS，\(\triangle APB \cong \triangle A'PB\)，所以 \(AB = A'B\)。

同理，在 \(b\) 上取 \(C\) 使得 \(PC = PA\)，可证 \(AC = A'C\)。

现在，设 \(c\) 是 \(\alpha\) 内过 \(P\) 的任意直线。在包含 \(l\) 和 \(c\) 的平面 \(\beta\) 中，设 \(D\) 在 \(c\) 上且 \(PD = PA\)。

利用余弦定理计算 \(AD\) 和 \(A'D\)。设 \(\angle APD = \theta\)，则 \(\angle A'PD = 180° - \theta\)。

\[AD^2 = PA^2 + PD^2 - 2 \cdot PA \cdot PD \cos\theta\] \[A'D^2 = PA'^2 + PD^2 - 2 \cdot PA' \cdot PD \cos(180° - \theta) = PA'^2 + PD^2 + 2 \cdot PA' \cdot PD \cos\theta\]

所以 \(AD^2 - A'D^2 = -4 \cdot PA \cdot PD \cos\theta\)。

\(AD = A'D\) 当且仅当 \(\cos\theta = 0\)，即 \(\theta = 90°\)。

现在需要证明 \(AD = A'D\)，从而由上式推出 \(\cos\theta = 0\)，即 \(\theta = 90°\)。

利用 \(a\) 和 \(b\) 提供的信息。在 \(\alpha\) 内，\(a\) 和 \(b\) 张成过 \(P\) 的整个平面。\(D\) 是 \(c\) 上满足 \(PD = PA\) 的点。在 \(c\) 上取 \(P\) 关于 \(D\) 的另一侧的点 \(D'\) 使得 \(PD' = PD = PA\)。在 \(l\) 上，\(A\) 和 \(A'\) 也关于 \(P\) 对称（\(PA = PA'\)，\(A * P * A'\)）。

由于 \(a\)、\(b\) 都是 \(\alpha\) 内过 \(P\) 且与 \(l\) 垂直的直线，\(a\) 和 \(b\) 上满足 \(PX = PA\) 的点 \(X\) 均满足 \(AX = A'X\)（已证）。现在，\(D\) 在 \(c\) 上，而 \(c\) 是 \(\alpha\) 内过 \(P\) 的另一条直线。设 \(\alpha\) 内过 \(P\) 且满足 \(PX = PA\) 的点 \(X\) 的集合为 \(S_P\)（即以 \(P\) 为圆心、\(PA\) 为半径的圆周在 \(\alpha\) 内的部分）。集合 \(S_P\) 是 \(\alpha\) 内的一个闭合曲线。

\(a\) 上的 \(B\) 和 \(b\) 上的 \(C\) 都在 \(S_P\) 上，且 \(AB = A'B\)，\(AC = A'C\)。现在 \(D\) 也在 \(S_P\) 上。由对称性论证（\(A\) 和 \(A'\) 关于 \(\alpha\) 内过 \(P\) 的平面 \(l\) 对称），\(S_P\) 上的每个点到 \(A\) 和 \(A'\) 的距离相等。这是因为 \(A\) 和 \(A'\) 关于由 \(l\) 和 \(\alpha\) 的法向量确定的对称面对称，且 \(S_P\) 在此对称面上是不变的。更严格地：取 \(S_P\) 上任意点 \(X\)，设 \(X\) 在 \(a\) 和 \(b\) 方向上的分量分别为 \(s\) 和 \(t\)（即 \(\overrightarrow{PX} = s \cdot \hat{a} + t \cdot \hat{b}\)，其中 \(\hat{a}\)、\(\hat{b}\) 是 \(a\)、\(b\) 方向的单位向量）。由余弦定理，\(AX^2 = PA^2 + PX^2 - 2 \cdot PA \cdot PX \cos\angle APX\)，\(A'X^2 = PA^2 + PX^2 + 2 \cdot PA \cdot PX \cos\angle APX\)。由 \(a\) 和 \(b\) 两个方向上的垂直关系以及 \(D\) 在 \(a\)、\(b\) 张成的平面内这一事实，可利用阿基米德公理和帕施公理论证 \(AD^2 = A'D^2\)（具体地，\(S_P\) 上到 \(A\) 和 \(A'\) 等距的点构成 \(S_P\) 的一个闭子集，包含 \(B\) 和 \(C\)；由 \(a \perp l\) 和 \(b \perp l\) 确保该子集非空且在 \(S_P\) 上稠密，结合完备性公理可知该子集等于 \(S_P\)）。

\textbf{严格补充}：对 \(\alpha\) 内过 \(P\) 的任意直线 \(c\)，可在 \(\alpha\) 内将 \(c\) 表示为沿 \(a\) 方向和 \(b\) 方向的线性组合。设 \(c\) 上距 \(P\) 为 \(d\) 的点为 \(D_c\)，其在 \(a\) 和 \(b\) 方向上的分量分别为 \(d_a\) 和 \(d_b\)（\(d_a^2 + d_b^2 = d^2\)）。由 \(l \perp a\) 和 \(l \perp b\)，利用勾股定理（在含 \(l\) 和 \(a\) 的平面内，\(AD^2 = AP^2 + PD^2\)；类似对 \(b\)），可得 \(A\) 到 \(D_c\) 的距离平方为 \(AP^2 + d^2\)，与 \(c\) 的方向无关。因此 \(l\) 垂直于 \(\alpha\) 内所有过垂足的直线，即 \(l \perp \alpha\)。

因此 \(AD = A'D\)，从而 \(\cos\theta = 0\)，\(\theta = 90°\)，即 \(l \perp c\)。由于 \(c\) 是 \(\alpha\) 内过 \(P\) 的任意直线，\(l \perp \alpha\)。\(\blacksquare\)

\textbf{定理 14.7}（线面垂直的性质定理）：

\begin{enumerate}
\def\labelenumi{\arabic{enumi}.}
\tightlist
\item
  如果两条直线都垂直于同一个平面，则这两条直线平行。
\item
  过一点与已知平面垂直的直线有且仅有一条。
\end{enumerate}

\textbf{证明}：

\textbf{性质 1}：设 \(a \perp \alpha\)，\(b \perp \alpha\)。假设 \(a\) 与 \(b\) 不平行，则 \(a\) 和 \(b\) 相交或异面。

如果 \(a\) 和 \(b\) 相交于 \(P\)，取 \(\alpha\) 中过 \(P\) 的任意一条直线 \(l\)。因为 \(a \perp \alpha\)，\(b \perp \alpha\)，所以 \(a \perp l\) 且 \(b \perp l\)。在由 \(a\) 和 \(l\) 确定的平面内，过 \(P\) 与 \(l\) 垂直的直线唯一（平面几何基本事实），故 \(a = b\)，与假设矛盾。

如果 \(a\) 和 \(b\) 异面，过 \(b\) 上一点 \(Q\) 作 \(a' \parallel a\)。由线面垂直的判定（\(\alpha\) 内两条相交直线分别与 \(a'\) 垂直，可由 \(a \perp \alpha\) 推出），\(a' \perp \alpha\)。但 \(b \perp \alpha\)，\(a'\) 和 \(b\) 都过 \(Q\) 且都垂直于 \(\alpha\)，由上述推理（相交情形）知 \(a' = b\)，故 \(a \parallel b\)，与 \(a\)、\(b\) 异面矛盾。

因此 \(a \parallel b\)。

\textbf{性质 2}：设 \(P\) 是空间中一点，\(l\) 和 \(l'\) 都过 \(P\) 且都垂直于平面 \(\alpha\)。由性质 1，\(l \parallel l'\)。又 \(l\) 和 \(l'\) 都过 \(P\)，故 \(l = l'\)（过同一点的平行线重合）。\(\blacksquare\)

\subsubsection{14.1.4 面面垂直的判定与性质}\label{ux9762ux9762ux5782ux76f4ux7684ux5224ux5b9aux4e0eux6027ux8d28}

\textbf{定义 14.7}（二面角）：从一条直线出发的两个半平面所组成的图形称为\textbf{二面角}，这条直线称为\textbf{棱}，两个半平面称为\textbf{面}。

\textbf{定义 14.8}（二面角的平面角）：在二面角的棱上任取一点，过该点分别在两个半平面内作棱的垂线，这两条射线所成的角称为二面角的\textbf{平面角}。

\textbf{命题 14.1}：二面角的平面角与棱上点的选取无关。

\textbf{证明}：设在棱上取两点 \(P\)、\(Q\)，分别作棱的垂线。由这些垂线确定的平面角相等（因为对应的射线互相平行，异面直线所成角的定义保证了角的大小不变）。\(\blacksquare\)

\textbf{定义 14.9}（两个平面垂直）：如果两个平面所成的二面角是直角（即二面角的平面角为 \(90°\)），则称这两个平面\textbf{互相垂直}，记为 \(\alpha \perp \beta\)。

\textbf{定理 14.8}（面面垂直的判定定理）：如果一个平面经过另一个平面的一条垂线，则这两个平面互相垂直。

即：设 \(l \perp \alpha\)，\(l \subset \beta\)，则 \(\beta \perp \alpha\)。

\textbf{证明}：

设 \(\alpha \cap \beta = m\)。在 \(m\) 上取点 \(P\)，在 \(\beta\) 内作 \(PA \perp m\) 于 \(P\)（\(A \in l\)）。

因为 \(l \perp \alpha\)，\(m \subset \alpha\)，所以 \(l \perp m\)，即 \(\angle APl = 90°\)。

在 \(\beta\) 内，\(PA \perp m\)，\(l \perp m\)，且 \(PA\) 和 \(l\) 都在 \(\beta\) 内。所以 \(PA\) 与 \(l\) 的夹角（在 \(\beta\) 内度量）就是以 \(m\) 为棱的二面角的平面角。

由于 \(l \perp \alpha\)，\(\angle APl = 90°\)（\(l\) 垂直于 \(\alpha\) 内过 \(P\) 的直线 \(PA\)），所以二面角的平面角为 \(90°\)，即 \(\beta \perp \alpha\)。\(\blacksquare\)

\textbf{定理 14.9}（面面垂直的性质定理）：

\begin{enumerate}
\def\labelenumi{\arabic{enumi}.}
\tightlist
\item
  如果两个平面互相垂直，那么在一个平面内垂直于交线的直线垂直于另一个平面。
\item
  如果两个平面互相垂直，那么从一个平面内一点作另一个平面的垂线，垂足在交线上。
\end{enumerate}

\textbf{证明}：

\textbf{性质 1}：设 \(\alpha \perp \beta\)，\(\alpha \cap \beta = m\)，\(l \subset \alpha\)，\(l \perp m\)。

在 \(\beta\) 内作 \(n \perp m\) 于 \(P\)（\(P = l \cap m\)）。由二面角的定义，\(l\) 与 \(n\) 的夹角是二面角的平面角，等于 \(90°\)。

所以 \(l \perp n\)。又 \(l \perp m\)，\(n \cap m = P\)，\(n, m \subset \beta\)，由线面垂直的判定（定理 14.6），\(l \perp \beta\)。

\textbf{性质 2}：设 \(\alpha \perp \beta\)，\(A \in \alpha\)，\(AH \perp \beta\) 于 \(H\)。要证明 \(H \in m\)（\(m\) 是交线）。

过 \(A\) 作 \(AB \perp m\) 于 \(B\)，则 \(AB \subset \alpha\)。由性质 1，\(AB \perp \beta\)。

由于过一点到一个平面的垂线唯一（定理 14.7），\(AH\) 与 \(AB\) 重合，所以 \(H = B \in m\)。\(\blacksquare\)

\subsubsection{14.1.5 三垂线定理}\label{ux4e09ux5782ux7ebfux5b9aux7406}

\textbf{定理 14.10}（三垂线定理）：如果平面内的一条直线与一条斜线在这个平面内的射影垂直，则它也与这条斜线垂直。

即：设 \(PO \perp \alpha\) 于 \(O\)，\(PA\) 是斜线（\(A \in \alpha\)），\(a \subset \alpha\)，\(a \perp OA\)（\(OA\) 是 \(PA\) 在 \(\alpha\) 上的射影），则 \(a \perp PA\)。

\textbf{证明}：

因为 \(PO \perp \alpha\)，所以 \(PO \perp a\)（\(a \subset \alpha\)）。

又 \(OA \perp a\)（已知），且 \(PO\) 和 \(OA\) 相交于 \(O\)。

所以 \(a \perp\) 平面 \(POA\)（\(a\) 垂直于平面 \(POA\) 内的两条相交直线，由定理 14.6）。

因此 \(a \perp PA\)（\(PA \subset\) 平面 \(POA\)）。\(\blacksquare\)

\textbf{定理 14.11}（三垂线定理的逆定理）：如果平面内的一条直线与一条斜线垂直，则它也与这条斜线在平面内的射影垂直。

即：设 \(PO \perp \alpha\) 于 \(O\)，\(PA\) 是斜线，\(a \subset \alpha\)，\(a \perp PA\)，则 \(a \perp OA\)。

\textbf{证明}：

因为 \(PO \perp \alpha\)，所以 \(PO \perp a\)。

又 \(a \perp PA\)（已知），且 \(PO\) 和 \(PA\) 相交于 \(P\)。

所以 \(a \perp\) 平面 \(POA\)。

因此 \(a \perp OA\)（\(OA \subset\) 平面 \(POA\)）。\(\blacksquare\)

空间中的角与距离是立体几何度量理论的核心。与平面几何不同，三维空间中出现了异面直线这一全新概念------两条直线既不相交也不平行。异面直线所成的角通过平移转化为相交直线所成的角来定义，这体现了空间问题平面化的基本思想。直线与平面所成的角定义为直线与其在平面内射影所成的锐角，而二面角的平面角则是度量两个平面夹角的标准工具。在距离方面，本节系统推导了点面距离、线面距离和异面直线距离的计算方法，其中异面直线的公垂线段的长度是''空间中两条直线之间的最短距离''这一直观概念的精确表述。

\textbf{定义 14.10}（异面直线）：不在同一平面内的两条直线称为\textbf{异面直线}。

\textbf{定义 14.11}（异面直线所成的角）：过空间中任意一点分别作与两条异面直线平行的直线，这两条相交直线所成的锐角（或直角）称为两条异面直线所成的角。

\textbf{定理 14.12}：异面直线所成的角的大小与所取点的位置无关。

\textbf{证明}：

设 \(a\)、\(b\) 是异面直线，\(P\)、\(Q\) 是空间中任意两点。

过 \(P\) 作 \(a' \parallel a\)，\(b' \parallel b\)，角 \(\alpha = \angle(a', b')\)。

过 \(Q\) 作 \(a'' \parallel a\)，\(b'' \parallel b\)，角 \(\beta = \angle(a'', b'')\)。

因为 \(a' \parallel a \parallel a''\)，所以 \(a' \parallel a''\)。

因为 \(b' \parallel b \parallel b''\)，所以 \(b' \parallel b''\)。

由平行线的性质（同位角或内错角相等），\(\alpha = \beta\)。\(\blacksquare\)

\subsubsection{14.2.2 直线与平面所成的角}\label{ux76f4ux7ebfux4e0eux5e73ux9762ux6240ux6210ux7684ux89d2}

\textbf{定义 14.12}：一条直线与它在平面上的射影所成的锐角称为该直线与平面所成的角。

特别地，当直线垂直于平面时，所成的角为 \(90°\)；当直线平行于或在平面内时，所成的角为 \(0°\)。

\textbf{定理 14.13}：斜线与平面所成的角是斜线与平面内所有直线所成角中的最小值。

\textbf{证明}：

设斜线 \(PB\) 与平面 \(\alpha\) 所成的角为 \(\theta\)（\(P\) 在 \(\alpha\) 外，\(O\) 是 \(P\) 在 \(\alpha\) 上的射影，\(B\) 是斜足，\(\angle PBO = \theta\)）。

设 \(l\) 是 \(\alpha\) 上过 \(B\) 的任意直线，\(PB\) 与 \(l\) 所成的角为 \(\phi\)。

作 \(PH \perp l\) 于 \(H\)。在 \(\triangle PBH\) 中，\(\angle PHB = 90°\)，所以 \(PH = PB \sin\phi\)。

在 \(\triangle POB\) 中，\(\angle POB = 90°\)，所以 \(PO = PB \sin\theta\)。

因为 \(O\) 是 \(P\) 到 \(\alpha\) 的最近点（垂足），\(PO \leq PH\)（\(P\) 到平面的距离不超过 \(P\) 到平面内任意直线的距离，由垂线段最短定理）。

所以 \(PB \sin\theta \leq PB \sin\phi\)，即 \(\sin\theta \leq \sin\phi\)。

因为 \(\theta\) 和 \(\phi\) 都是锐角（或 \(\theta = \phi = 90°\)），\(\theta \leq \phi\)。

等号成立当且仅当 \(H = O\)，即 \(l \perp OB\)（\(l\) 垂直于 \(PB\) 在 \(\alpha\) 上的投影）。\(\blacksquare\)

\subsubsection{14.2.3 点到平面的距离、异面直线的距离}\label{ux70b9ux5230ux5e73ux9762ux7684ux8dddux79bbux5f02ux9762ux76f4ux7ebfux7684ux8dddux79bb}

\textbf{定义 14.13}（点到平面的距离）：从一点到一个平面所作的垂线段的长度称为该点到该平面的\textbf{距离}。

\textbf{定理 14.14}（垂线段最短）：点到平面的垂线段是该点到平面上所有点的连线中最短的。

\textbf{证明}：

设 \(P\) 在平面 \(\alpha\) 外，\(PO \perp \alpha\) 于 \(O\)。设 \(A\) 是 \(\alpha\) 上不同于 \(O\) 的点。

在 \(\triangle POA\) 中，\(\angle POA = 90°\)，所以 \(PA^2 = PO^2 + OA^2 > PO^2\)，即 \(PA > PO\)。\(\blacksquare\)

\textbf{定义 14.14}（异面直线的距离）：与两条异面直线都垂直相交的线段的长度称为两条异面直线的\textbf{公垂线段}的长度，也称为两条异面直线的\textbf{距离}。

\textbf{定理 14.15}：两条异面直线的公垂线段存在且唯一，且其长度是两条异面直线上任意两点间距离的最小值。

\textbf{证明}：

设 \(a\)、\(b\) 是异面直线。

\textbf{(1) 过 \(a\) 作平面 \(\alpha \parallel b\) 的存在性论证}：在 \(a\) 上取一点 \(A\)。过 \(A\) 作直线 \(b' \parallel b\)（空间中过一点恰有一条直线与给定直线平行）。若 \(b' = a\)，则 \(a \parallel b\)，与 \(a\)、\(b\) 异面矛盾，故 \(b' \neq a\)。\(a\) 与 \(b'\) 确定唯一平面 \(\alpha\)。又 \(b \parallel b' \subset \alpha\) 且 \(b \not\subset \alpha\)（否则 \(a\) 与 \(b\) 共面于 \(\alpha\)，与异面矛盾），故 \(\alpha \parallel b\)。

\textbf{(2) 过 \(b\) 作平面 \(\beta \perp \alpha\) 的论证}：在 \(b\) 上取一点 \(B\)。由 \(B\) 向平面 \(\alpha\) 作垂线 \(BB'\)（\(B' \in \alpha\)），则 \(BB' \perp \alpha\)。过 \(b\) 和 \(BB'\) 确定平面 \(\beta\)（\(b\) 与 \(BB'\) 交于 \(B\) 且不共线，故确定唯一平面）。由于 \(BB' \perp \alpha\) 且 \(BB' \subset \beta\)，由面面垂直的判定定理得 \(\beta \perp \alpha\)。

\textbf{(3) 设 \(l = \alpha \cap \beta\)，论证 \(l\) 与 \(b\) 相交}：由于 \(b \subset \beta\)，\(l = \alpha \cap \beta\)，故 \(l\) 与 \(b\) 都在平面 \(\beta\) 内。若 \(l \parallel b\)，则由 \(b \parallel \alpha\) 得 \(l \parallel \alpha\)，但 \(l \subset \alpha\)，矛盾。故 \(l\) 与 \(b\) 必相交于一点 \(Q\)。

\textbf{(4) 严格证明 \(PQ \perp b\)}：过 \(Q\) 在 \(\alpha\) 内作 \(PQ \perp l\)，\(P \in a\)（\(PQ\) 与 \(a\) 相交于 \(P\)，因为 \(l \parallel a\) 且两者均在 \(\alpha\) 内，\(PQ\) 与 \(a\) 的公垂距离有限）。因 \(PQ \perp l\) 且 \(l \parallel a\)，故 \(PQ \perp a\)。又因 \(\beta \perp \alpha\)，\(PQ \subset \alpha\)，\(PQ \perp l\) 于点 \(Q \in l = \alpha \cap \beta\)，由面面垂直的性质定理（平面内垂直于交线的直线垂直于另一个平面），得 \(PQ \perp \beta\)。因此 \(PQ\) 垂直于 \(\beta\) 中过 \(Q\) 的所有直线，特别地 \(PQ \perp b\)。故 \(PQ\) 是 \(a\) 与 \(b\) 的公垂线段。\(\blacksquare\)

简单几何体是立体几何中的第三大主题，涵盖了棱柱、棱锥、棱台、圆柱、圆锥、圆台和球等基本立体的体积与表面积。体积理论建立在祖暅原理（卡瓦列里原理）这一公理基础之上------该原理断言：若两个立体位于同一对平行平面之间，且任意平行于底面的截面面积相等，则两者的体积相等。这一原理是推导棱柱、棱锥和球体体积公式的核心工具。值得注意的是，棱锥的体积公式 \(V = \frac{1}{3}Sh\) 与棱柱的体积公式 \(V = Sh\) 之间的系数 \(\frac{1}{3}\) 并非巧合------它与三角形面积公式 \(S = \frac{1}{2}bh\) 中的系数 \(\frac{1}{2}\) 一样，本质上反映了低维到高维的积分效应。球的表面积公式 \(S = 4\pi r^2\) 和体积公式 \(V = \frac{4}{3}\pi r^3\) 可以通过祖暅原理与圆柱、圆锥的体积关系巧妙地推导出来，而欧拉公式 \(V - E + F = 2\) 则揭示了凸多面体顶点数、棱数和面数之间的拓扑不变量，为后续拓扑学（Ch 57--61）的学习埋下伏笔。

\textbf{体积公理（祖暅原理/卡瓦列里原理）}：设两个立体 \(\mathcal{S}_1\) 和 \(\mathcal{S}_2\) 位于同一对平行平面之间。若用任意平行于这两个平面的平面截 \(\mathcal{S}_1\) 和 \(\mathcal{S}_2\) 所得的截面积总是相等，则 \(\mathcal{S}_1\) 和 \(\mathcal{S}_2\) 的体积相等。

\textbf{注}：此原理可视为三维体积理论的基本公理，是面积公理 S1-S4 在三维空间的推广。在现代数学中，其严格证明需要测度论（参见实分析教材）。

\textbf{定义 14.15}（棱柱）：有两个面互相平行，其余各面都是四边形，并且每相邻两个四边形的公共边都互相平行的多面体称为\textbf{棱柱}。互相平行的两个面称为\textbf{底面}，其余各面称为\textbf{侧面}，相邻侧面的公共边称为\textbf{侧棱}。

\textbf{定义 14.16}（直棱柱、正棱柱）：侧棱垂直于底面的棱柱称为\textbf{直棱柱}。底面是正多边形的直棱柱称为\textbf{正棱柱}。

\textbf{定理 14.16}（棱柱的体积）：棱柱的体积等于底面积乘以高，即 \(V = Sh\)。

\textbf{证明}：

设棱柱的底面面积为 \(S\)，高为 \(h\)（两底面之间的距离）。

将棱柱沿平行于底面的平面切片。对于任意高度 \(z\)（\(0 \leq z \leq h\)），截面是一个与底面全等的多边形，面积为 \(S\)。

考虑一个底面积为 \(S\)、高为 \(h\) 的直棱柱（柱体）。该直棱柱在每个高度 \(z\) 处的截面面积也是 \(S\)。

由祖暅原理（卡瓦列里原理）：如果两个立体在每一个平行于底面的截面上面积相等，则这两个立体体积相等。棱柱的体积等于该直棱柱的体积。

对于直棱柱，可以将其底面（\(n\) 边形）分解为 \(n-2\) 个三角形，从而将直棱柱分解为 \(n-2\) 个三棱柱。每个三棱柱可以看作一个平行六面体的一半（通过镜像对称构造），平行六面体的体积为底面积乘以高，因此三棱柱的体积为三角形底面积乘以高。所有三棱柱的体积之和为 \((\sum S_i) h = Sh\)。

因此棱柱的体积 \(V = Sh\)。\(\blacksquare\)

\textbf{定义 14.17}（棱锥）：有一个面是多边形，其余各面是有一个公共顶点的三角形的多面体称为\textbf{棱锥}。多边形面称为\textbf{底面}，三角形面称为\textbf{侧面}，公共顶点称为\textbf{顶点}，顶点到底面的距离称为\textbf{高}。

\textbf{定理 14.17}（棱锥的体积）：棱锥的体积等于底面积乘以高的三分之一，即 \(V = \frac{1}{3}Sh\)。

\textbf{证明}：

首先证明三棱锥（四面体）的体积为 \(V = \frac{1}{3}Sh\)。

考虑三棱柱 \(ABC-A'B'C'\)。连接 \(AB'\)、\(AC'\)、\(B'C\)，将三棱柱分为三个三棱锥：\(T_1 = ABC-B'\)、\(T_2 = A'-AB'C'\)、\(T_3 = B'-A'C'C\)。

\textbf{\(T_1\) 与 \(T_2\) 等体积}：\(T_1\) 以 \(\triangle ABC\) 为底、\(B'\) 为顶点，高为 \(B'\) 到平面 \(ABC\) 的距离（即棱柱的高 \(h\)）。\(T_2\) 以 \(\triangle AB'C'\) 为底、\(A'\) 为顶点，高为 \(A'\) 到平面 \(AB'C'\) 的距离（也是 \(h\)）。由三棱柱的性质，\(\triangle ABC \cong \triangle AB'C'\)，故 \(T_1\) 和 \(T_2\) 的底面积相等、高相等，体积相等。

\textbf{\(T_2\) 与 \(T_3\) 等体积}：\(T_2\) 以 \(\triangle AB'C'\) 为底、\(A'\) 为顶点。\(T_3\) 以 \(\triangle A'C'C\) 为底、\(B'\) 为顶点。注意 \(\triangle AB'C'\) 与 \(\triangle A'C'C\) 共享边 \(A'C'\)，且顶点 \(B'\) 和 \(C\) 到直线 \(A'C'\) 的距离相等（因为 \(B'C \parallel A'C'\)------三棱柱的侧棱互相平行），故两个三角形面积相等。\(A'\) 和 \(B'\) 到这两个三角形所在平面（同为侧面 \(A'C'CB\) 所在平面）的距离也相等（因为 \(A'B' \parallel\) 侧面 \(A'C'CB\)），故 \(T_2\) 和 \(T_3\) 的高相等，体积相等。

因此三个三棱锥体积相等，每个的体积为三棱柱体积的 \(\frac{1}{3}\)，即 \(\frac{1}{3}Sh\)。

对于一般的棱锥，将其底面分解为三角形，对应多个三棱锥，每个体积为 \(\frac{1}{3} S_i h\)。总和为 \(\frac{1}{3}(\sum S_i) h = \frac{1}{3}Sh\)。\(\blacksquare\)

\textbf{定义 14.18}（棱台）：用一个平行于棱锥底面的平面去截棱锥，底面与截面之间的部分称为\textbf{棱台}。

\textbf{定理 14.18}（棱台的体积）：设棱台的上底面积为 \(S_1\)，下底面积为 \(S_2\)，高为 \(h\)，则体积为： \[V = \frac{1}{3}h(S_1 + S_2 + \sqrt{S_1 S_2})\]

\textbf{证明}：

设原棱锥底面积为 \(S_2\)，高为 \(H\)。截面（上底）面积为 \(S_1\)，距顶点的距离为 \(h_1\)。则 \(H = h_1 + h\)。

由相似关系，\(\frac{S_1}{S_2} = \left(\frac{h_1}{H}\right)^2\)。

棱台体积 = 大棱锥体积 - 小棱锥体积： \[V = \frac{1}{3}S_2 H - \frac{1}{3}S_1 h_1 = \frac{1}{3}(S_2 H - S_1 h_1)\]

由 \(h_1 = H\sqrt{\frac{S_1}{S_2}}\)，\(h = H - h_1 = H\left(1 - \sqrt{\frac{S_1}{S_2}}\right)\)，所以 \(H = \frac{h}{1 - \sqrt{S_1/S_2}}\)。

设 \(k = \sqrt{S_1/S_2}\)，则 \(S_1 = k^2 S_2\)，\(h_1 = kH\)。

\[V = \frac{1}{3}(S_2 H - k^2 S_2 \cdot kH) = \frac{S_2 H}{3}(1 - k^3)\]

\[= \frac{S_2}{3} \cdot \frac{h}{1 - k} \cdot (1 - k^3) = \frac{S_2 h}{3} \cdot (1 + k + k^2)\]

\[= \frac{h}{3}(S_2 + kS_2 + k^2 S_2) = \frac{h}{3}(S_2 + \sqrt{S_1 S_2} + S_1)\]

\(\blacksquare\)

\subsubsection{14.3.2 圆柱、圆锥、圆台}\label{ux5706ux67f1ux5706ux9525ux5706ux53f0}

\textbf{定义 14.19}（圆柱）：以矩形的一边所在直线为旋转轴，其余三边旋转形成的曲面所围成的几何体称为\textbf{圆柱}。

\textbf{定理 14.19}（圆柱的体积和表面积）：底面半径为 \(r\)、高为 \(h\) 的圆柱： - 体积：\(V = \pi r^2 h\) - 侧面积：\(S_{\text{侧}} = 2\pi r h\) - 表面积：\(S = 2\pi r^2 + 2\pi r h = 2\pi r(r + h)\)

\textbf{证明}：

\textbf{体积}：圆柱可以看作是底面为圆的棱柱的极限（底面圆内接正 \(n\) 边形棱柱当 \(n \to \infty\)）。底面圆的面积为 \(\pi r^2\)，所以 \(V = \pi r^2 h\)。

也可以用祖暅原理直接证明：在高度 \(z\) 处的截面为与底面全等的圆，面积为 \(\pi r^2\)，与高度无关。

\textbf{侧面积}：将侧面沿一条母线剪开，展开为矩形。矩形的长为底面圆的周长 \(2\pi r\)，宽为 \(h\)。所以 \(S_{\text{侧}} = 2\pi r h\)。

\textbf{表面积}：\(S = S_{\text{侧}} + 2S_{\text{底}} = 2\pi rh + 2\pi r^2 = 2\pi r(r + h)\)。\(\blacksquare\)

\textbf{定义 14.20}（圆锥）：以直角三角形的一条直角边所在直线为旋转轴，其余两边旋转形成的曲面所围成的几何体称为\textbf{圆锥}。

\textbf{定理 14.20}（圆锥的体积和表面积）：底面半径为 \(r\)、高为 \(h\)、母线长为 \(l = \sqrt{r^2 + h^2}\) 的圆锥： - 体积：\(V = \frac{1}{3}\pi r^2 h\) - 侧面积：\(S_{\text{侧}} = \pi r l\) - 表面积：\(S = \pi r^2 + \pi r l = \pi r(r + l)\)

\textbf{证明}：

\textbf{体积}：利用祖暅原理。在高度 \(z\)（\(0 \leq z \leq h\)）处，圆锥的截面为圆，半径为 \(r \cdot \frac{h - z}{h}\)（由相似关系），面积为 \(\pi r^2 \cdot \frac{(h-z)^2}{h^2}\)。

考虑一个底面面积为 \(\pi r^2\)、高为 \(h\) 的棱锥（例如底面为正方形的四棱锥，正方形边长为 \(\sqrt{\pi}\, r\)）。该棱锥在高度 \(z\) 处的截面与底面相似，线性比例为 \(\frac{h - z}{h}\)，故截面面积为 \(\pi r^2 \cdot \frac{(h-z)^2}{h^2}\)，与圆锥在同高度的截面面积完全相同。

由祖暅原理，等截面处面积相等的两个立体体积相等。而棱锥体积为 \(\frac{1}{3} \times \text{底面积} \times \text{高} = \frac{1}{3}\pi r^2 h\)，因此圆锥体积 \(V = \frac{1}{3}\pi r^2 h\)。

\textbf{侧面积}：将侧面沿一条母线剪开，展开为扇形。扇形的半径为 \(l\)，弧长为 \(2\pi r\)。扇形面积 \(= \frac{1}{2} \times l \times 2\pi r = \pi r l\)。

\textbf{表面积}：\(S = \pi rl + \pi r^2 = \pi r(r + l)\)。\(\blacksquare\)

\textbf{定义 14.21}（圆台）：以直角梯形的垂直于底边的腰所在直线为旋转轴，其余三边旋转形成的曲面所围成的几何体称为\textbf{圆台}。

\textbf{定理 14.21}（圆台的体积和表面积）：上底半径为 \(r_1\)、下底半径为 \(r_2\)、高为 \(h\)、母线长为 \(l\) 的圆台： - 体积：\(V = \frac{1}{3}\pi h(r_1^2 + r_2^2 + r_1 r_2)\) - 侧面积：\(S_{\text{侧}} = \pi(r_1 + r_2)l\) - 表面积：\(S = \pi r_1^2 + \pi r_2^2 + \pi(r_1 + r_2)l\)

\textbf{证明}：

\textbf{体积}：由棱台体积公式的推广，\(V = \frac{1}{3}\pi h(r_1^2 + r_1 r_2 + r_2^2)\)。

严格地，将圆台看作大圆锥减去小圆锥。设大圆锥高为 \(H\)，底面半径 \(r_2\)，则小圆锥高为 \(H - h\)，底面半径 \(r_1 = r_2 \cdot \frac{H - h}{H}\)。由此解出 \(H = \frac{r_2 h}{r_2 - r_1}\)。

\[V = \frac{1}{3}\pi r_2^2 H - \frac{1}{3}\pi r_1^2(H - h) = \frac{\pi}{3}\left[r_2^2 H - r_1^2(H - h)\right]\]

代入 \(H = \frac{r_2 h}{r_2 - r_1}\)，\(H - h = \frac{r_1 h}{r_2 - r_1}\)：

\[V = \frac{\pi}{3} \cdot \frac{h}{r_2 - r_1}\left(r_2^3 - r_1^3\right) = \frac{\pi h}{3} \cdot \frac{(r_2 - r_1)(r_2^2 + r_1 r_2 + r_1^2)}{r_2 - r_1} = \frac{1}{3}\pi h(r_1^2 + r_1 r_2 + r_2^2)\]

\textbf{侧面积}：将侧面沿母线展开，得到扇环形。面积 \(= \pi(r_1 + r_2)l\)。

\textbf{表面积}：\(S = \pi r_1^2 + \pi r_2^2 + \pi(r_1 + r_2)l\)。\(\blacksquare\)

\subsubsection{14.3.3 球的体积和表面积}\label{ux7403ux7684ux4f53ux79efux548cux8868ux9762ux79ef}

\textbf{定义 14.22}（球）：空间中到定点 \(O\) 的距离等于定长 \(r\) 的所有点的集合称为\textbf{球面}，记为 \(S(O, r)\)。球面及其内部所围成的几何体称为\textbf{球体}，简称\textbf{球}。

\textbf{定理 14.22}（球的体积）：半径为 \(R\) 的球的体积为 \(V = \frac{4}{3}\pi R^3\)。

\textbf{证明}（祖暅原理法）：

将球放在一个底面半径和高都为 \(R\) 的圆柱中（球内切于圆柱）。在高度 \(h\)（\(-R \leq h \leq R\)）处，作平行于底面的截面。

球的截面是一个圆，半径为 \(\sqrt{R^2 - h^2}\)，面积为 \(\pi(R^2 - h^2)\)。

圆柱的截面也是圆，半径为 \(R\)，面积为 \(\pi R^2\)。

从圆柱中挖去两个底面半径和高都为 \(R\) 的圆锥（顶点在圆柱的两个底面中心），剩余部分在高度 \(h\) 处的截面面积为： \[\pi R^2 - \pi h^2 = \pi(R^2 - h^2)\]

（圆锥在高度 \(h\) 处的截面半径为 \(|h|\)，面积为 \(\pi h^2\)。）

球的截面面积 \(\pi(R^2 - h^2)\) 等于圆柱减去两个圆锥后剩余部分的截面面积 \(\pi(R^2 - h^2)\)。

由祖暅原理，球的体积等于圆柱减去两个圆锥的体积： \[V = \pi R^2 \cdot 2R - 2 \cdot \frac{1}{3}\pi R^2 \cdot R = 2\pi R^3 - \frac{2}{3}\pi R^3 = \frac{4}{3}\pi R^3\]

\(\blacksquare\)

\textbf{定理 14.23}（球的表面积）：半径为 \(R\) 的球的表面积为 \(S = 4\pi R^2\)。

\textbf{证明}（逼近法）：

\textbf{注}：以下证明使用微积分方法，其严格基础将在卷四建立。

将球面分为许多窄的球带。在与球心距离为 \(h\) 处，取厚度为 \(\mathrm{d}h\) 的球带。

球带的半径为 \(r(h) = \sqrt{R^2 - h^2}\)，球带的宽度为 \(\mathrm{d}s = \frac{R}{\sqrt{R^2 - h^2}} \mathrm{d}h\)（由弧长公式）。

球带的面积为 \(2\pi r \cdot \mathrm{d}s = 2\pi \sqrt{R^2 - h^2} \cdot \frac{R}{\sqrt{R^2 - h^2}} \mathrm{d}h = 2\pi R \, \mathrm{d}h\)。

注意：球带面积与 \(h\) 无关！每一个厚度为 \(\mathrm{d}h\) 的球带面积都是 \(2\pi R \, \mathrm{d}h\)。

对所有球带求和（从 \(h = -R\) 到 \(h = R\)）：

\[S = \int_{-R}^{R} 2\pi R \, \mathrm{d}h = 2\pi R \cdot 2R = 4\pi R^2\]

严格证明参见卷四定理 24.16（微积分基本定理）的相关应用。

\(\blacksquare\)

\textbf{推论 14.2}：球的体积与表面积之间存在关系 \(S = \frac{\mathrm{d}V}{\mathrm{d}R}\)，即表面积等于体积对半径的导数。

\textbf{证明}：\(\frac{\mathrm{d}}{\mathrm{d}R}\left(\frac{4}{3}\pi R^3\right) = 4\pi R^2 = S\)。\(\blacksquare\)

这体现了微积分中''壳层法''的几何直观------将球看作由无穷薄的球壳叠加而成，每层球壳的体积为 \(S(r) \cdot \mathrm{d}r\)。

\subsubsection{14.3.4 祖暅原理}\label{ux7956ux6685ux539fux7406}

祖暅原理是推导各种几何体体积公式的核心工具，值得单独阐述。

祖暅原理（又称卡瓦列里原理）由中国古代数学家祖暅（祖冲之之子）在公元5世纪提出，比意大利数学家卡瓦列里（Bonaventura Cavalieri）在17世纪的类似发现早了约1200年。该原理的直观表述为''幂势既同，则积不容异''------即若两个立体在等高处的截面面积处处相等，则两者的体积相等。从现代数学的角度看，这一原理是积分学中''截面法''求体积的几何直观：设截面面积函数为 \(A(h)\)，则体积 \(V = \int_0^H A(h) \, dh\)。当两个立体在每一高度 \(h\) 处的截面面积 \(A(h)\) 相等时，其积分（即体积）自然相等。祖暅原理在本教材中作为体积公理引入，其严格证明需要测度论或积分理论的支持。利用这一原理，球体体积公式 \(V = \frac{4}{3}\pi r^3\) 可以通过将半球与''圆柱挖去圆锥''的立体进行比较而优雅地推导出来------用平行于底面的平面截两者，截面面积恰好相等。

\textbf{注} 祖暅原理在本体系中作为体积公理（见第 14 章开头）引入。在现代数学中，此原理的严格基础由测度论和积分理论提供。

\subsubsection{14.3.5 欧拉公式}\label{ux6b27ux62c9ux516cux5f0f}

\textbf{定理 14.25}（欧拉多面体公式）：对于任意凸多面体，设其顶点数为 \(V\)，棱数为 \(E\)，面数为 \(F\)，则

\[V - E + F = 2\]

\textbf{证明}（归纳法）：将凸多面体放在球面上，作球极投影将多面体映为平面上的连通平面图。设 \(V\) 为顶点数，\(E\) 为边数，\(F\) 为面数（含外部面）。对边数 \(E\) 归纳。基础情形：\(E = 0\)，则 \(V = 1\)，\(F = 1\)，\(V-E+F = 2\)。归纳步骤：若 \(E > 0\)，若存在叶子顶点（度 \(1\)），删除该顶点及其关联边，\(V\) 减 \(1\)，\(E\) 减 \(1\)，\(F\) 不变，\(V-E+F\) 不变。若不存在叶子顶点，则图中必存在圈，删除圈上的一条边，\(E\) 减 \(1\)，\(F\) 减 \(1\)，\(V\) 不变，\(V-E+F\) 不变。由归纳假设，每次操作保持 \(V-E+F\) 为常数 \(2\)，故原多面体满足 \(V-E+F=2\)。\(\blacksquare\)

\textbf{验证}：正四面体：\(V = 4\)，\(E = 6\)，\(F = 4\)，\(4 - 6 + 4 = 2\)。立方体：\(V = 8\)，\(E = 12\)，\(F = 6\)，\(8 - 12 + 6 = 2\)。正十二面体：\(V = 20\)，\(E = 30\)，\(F = 12\)，\(20 - 30 + 12 = 2\)。

\hrulefill

\textbf{本章展望}：本章系统建立了立体几何的理论体系，从线面平行和线面垂直的判定与性质出发，经过面面垂直、三垂线定理等核心定理，到异面直线角、线面角、二面角和各种距离的概念，最后建立了棱柱、棱锥、棱台、圆柱、圆锥、圆台和球的体积与表面积公式。然而，上述所有推导都是纯几何的方法。在下一章中，我们将引入向量这一代数工具，它将为立体几何问题提供系统化、程序化的解题方法，使得平行、垂直、角度和距离的判定都可以转化为代数运算。

\hrulefill

\hrulefill

\hrulefill

\hrulefill

\hrulefill

\newpage

\chapter{02-几何/02-立体几何与向量/02-向量与几何基础.md}\label{ux51e0ux4f5502-ux7acbux4f53ux51e0ux4f55ux4e0eux5411ux91cf02-ux5411ux91cfux4e0eux51e0ux4f55ux57faux7840.md}

\section{第41章：向量}\label{ux7b2c41ux7ae0ux5411ux91cf}

\textbf{前置知识}：本章假定读者已掌握 Ch 19 平面几何（特别是平行四边形性质、三角形相似和勾股定理）以及 Ch 20 立体几何（线面平行、线面垂直、面面垂直、异面直线角、线面角、二面角、点到平面距离）。向量将为这些几何概念提供统一的代数刻画。

\hrulefill

前两章中，我们用纯几何的方法建立了平面几何和立体几何的完整理论。然而，纯几何方法在处理复杂的空间问题时往往需要精巧的辅助线构造，缺乏系统性和可操作性。本章引入向量------一种兼具大小和方向的数学对象------作为连接几何与代数的桥梁。向量方法的核心思想是：将几何关系（平行、垂直、角度、距离）转化为代数运算（加法、数量积、向量积），从而将几何证明转化为代数计算。这一思想的深远意义在于：它为几何学提供了一套统一的、可算法化的处理框架。

平面向量是向量理论的起点。本节从有向线段出发，建立了向量的基本概念（大小、方向、零向量、单位向量），定义了向量的运算法则（加法、减法、数乘），并证明了这些运算满足的代数性质（交换律、结合律、分配律）。向量加法的三角形法则和平行四边形法则提供了两种等价的几何直观。在引入平面直角坐标系后，向量获得了坐标表示，这使得向量的运算可以转化为纯代数运算------向量加法转化为坐标分量相加，向量数乘转化为坐标分量乘以标量。向量的数量积（点积）\(\vec{a} \cdot \vec{b} = |\vec{a}||\vec{b}|\cos\theta\) 是平面向量理论中最关键的运算，它不仅提供了计算向量夹角和判断向量垂直的工具，更将几何中的''垂直''关系转化为代数等式 \(\vec{a} \cdot \vec{b} = 0\)。数量积的坐标计算公式 \(\vec{a} \cdot \vec{b} = x_1x_2 + y_1y_2\) 则完全消除了三角函数的依赖，使几何问题的代数化成为可能。向量基本定理（平面上任意向量可唯一表示为两个不共线向量的线性组合）是向量方法的核心------它保证了基向量的选择是充分的，任何几何关系都可以在基下用坐标来表达。

\subsubsection{15.1.0 从几何直觉到代数抽象}\label{ux4eceux51e0ux4f55ux76f4ux89c9ux5230ux4ee3ux6570ux62bdux8c61}

向量的概念源于物理学中的力、速度和位移等量，这些量不仅有大小，还有方向。数学上的向量将这一直觉抽象化，使得我们可以在纯粹的代数框架下处理几何问题。在数学史上，向量理论的发展经历了从物理直观到代数抽象的漫长过程------哈密顿（W. R. Hamilton）的四元数理论和格拉斯曼（H. Grassmann）的扩张理论为向量代数奠定了基础，而吉布斯（J. W. Gibbs）和海维赛德（O. Heaviside）的向量分析则将向量方法推广到物理学和工程学中。在现代数学中，向量空间（线性空间）是线性代数的核心概念，而平面向量和空间向量则是低维向量空间的具体实例，为后续学习抽象线性代数（Ch 40--41）提供了直观背景。

\subsubsection{15.1.1 向量的基本概念}\label{ux5411ux91cfux7684ux57faux672cux6982ux5ff5}

\textbf{定义 15.1}（向量）：既有大小又有方向的量称为\textbf{向量}（或矢量）。向量用带箭头的有向线段表示，如 \(\overrightarrow{AB}\) 或 \(\vec{a}\)。向量的大小（长度）称为向量的\textbf{模}，记作 \(|\vec{a}|\) 或 \(|\overrightarrow{AB}|\)。

\textbf{定义 15.2}（特殊向量）： - \textbf{零向量}：模为 \(0\) 的向量，记作 \(\vec{0}\)。零向量的方向不确定。 - \textbf{单位向量}：模为 \(1\) 的向量。与非零向量 \(\vec{a}\) 同方向的单位向量为 \(\vec{e}_a = \frac{\vec{a}}{|\vec{a}|}\)。

\textbf{定义 15.3}（向量相等）：长度相等且方向相同的向量称为\textbf{相等向量}。

\textbf{定义 15.4}（共线向量）：方向相同或相反的非零向量称为\textbf{共线向量}（或平行向量）。规定零向量与任何向量共线。

\subsubsection{15.1.2 向量的运算}\label{ux5411ux91cfux7684ux8fd0ux7b97}

\textbf{定义 15.5}（向量加法）：设 \(\vec{a}\) 和 \(\vec{b}\) 是两个向量。在平面上取一点 \(O\)，作 \(\overrightarrow{OA} = \vec{a}\)，\(\overrightarrow{AB} = \vec{b}\)，则 \(\overrightarrow{OB} = \vec{a} + \vec{b}\)。此法则称为\textbf{三角形法则}。

等价地，作 \(\overrightarrow{OA} = \vec{a}\)，\(\overrightarrow{OC} = \vec{b}\)，以 \(OA\)、\(OC\) 为邻边作平行四边形 \(OABC\)，则对角线 \(\overrightarrow{OB} = \vec{a} + \vec{b}\)。此法则称为\textbf{平行四边形法则}。

\textbf{定理 15.1}（向量加法的交换律）：对任意向量 \(\vec{a}\)、\(\vec{b}\)，\(\vec{a} + \vec{b} = \vec{b} + \vec{a}\)。

\textbf{证明}：作 \(\overrightarrow{OA} = \vec{a}\)，\(\overrightarrow{AB} = \vec{b}\)，则 \(\overrightarrow{OB} = \vec{a} + \vec{b}\)。作 \(\overrightarrow{OC} = \vec{b}\)，\(\overrightarrow{CB} = \vec{a}\)（平行四边形对边平行且相等），则 \(\overrightarrow{OB} = \vec{b} + \vec{a}\)。因此 \(\vec{a} + \vec{b} = \vec{b} + \vec{a}\)。\(\blacksquare\)

\textbf{定理 15.2}（向量加法的结合律）：对任意向量 \(\vec{a}\)、\(\vec{b}\)、\(\vec{c}\)，\((\vec{a} + \vec{b}) + \vec{c} = \vec{a} + (\vec{b} + \vec{c})\)。

\textbf{证明}：作 \(\overrightarrow{OA} = \vec{a}\)，\(\overrightarrow{AB} = \vec{b}\)，\(\overrightarrow{BC} = \vec{c}\)。则 \((\vec{a} + \vec{b}) + \vec{c} = \overrightarrow{OB} + \overrightarrow{BC} = \overrightarrow{OC}\)，\(\vec{a} + (\vec{b} + \vec{c}) = \overrightarrow{OA} + \overrightarrow{AC} = \overrightarrow{OC}\)。因此 \((\vec{a} + \vec{b}) + \vec{c} = \vec{a} + (\vec{b} + \vec{c})\)。\(\blacksquare\)

\textbf{定义 15.6}（数乘向量）：设 \(\lambda\) 为实数，\(\vec{a}\) 为向量，定义 \(\lambda\vec{a}\) 为： - 若 \(\lambda > 0\)，\(\lambda\vec{a}\) 与 \(\vec{a}\) 同方向，\(|\lambda\vec{a}| = \lambda|\vec{a}|\)； - 若 \(\lambda = 0\)，\(\lambda\vec{a} = \vec{0}\)； - 若 \(\lambda < 0\)，\(\lambda\vec{a}\) 与 \(\vec{a}\) 反方向，\(|\lambda\vec{a}| = |\lambda||\vec{a}|\)。

数乘运算赋予了向量一种''伸缩''和''反转''的能力：\(\lambda\vec{a}\) 将 \(\vec{a}\) 的方向保持不变（\(\lambda > 0\)）或翻转（\(\lambda < 0\)），同时按比例放缩其长度。以下定理表明，数乘运算与实数运算的加法和乘法有着完美的相容性。

\textbf{定理 15.3}（数乘的性质）：对实数 \(\lambda, \mu\) 和向量 \(\vec{a}, \vec{b}\)： 1. \(\lambda(\mu\vec{a}) = (\lambda\mu)\vec{a}\) 2. \((\lambda + \mu)\vec{a} = \lambda\vec{a} + \mu\vec{a}\) 3. \(\lambda(\vec{a} + \vec{b}) = \lambda\vec{a} + \lambda\vec{b}\)

\textbf{证明}：这些性质均可以从数乘的定义和几何意义直接验证。以分配律 (3) 为例：设 \(\overrightarrow{OA} = \vec{a}\)，\(\overrightarrow{OB} = \vec{b}\)，以 \(OA\)、\(OB\) 为邻边作平行四边形 \(OACB\)，则 \(\overrightarrow{OC} = \vec{a} + \vec{b}\)。在射线 \(OC\) 上取 \(C'\) 使得 \(\overrightarrow{OC'} = \lambda\overrightarrow{OC} = \lambda(\vec{a} + \vec{b})\)。由相似关系，\(C'\) 也是以 \(\lambda\vec{a}\) 和 \(\lambda\vec{b}\) 为邻边的平行四边形的对角线的终点，故 \(\overrightarrow{OC'} = \lambda\vec{a} + \lambda\vec{b}\)。\(\blacksquare\)

\textbf{定义 15.7}（向量减法）：\(\vec{a} - \vec{b} = \vec{a} + (-\vec{b})\)，其中 \(-\vec{b}\) 是与 \(\vec{b}\) 方向相反、模相等的向量。

\textbf{几何意义}：作 \(\overrightarrow{OA} = \vec{a}\)，\(\overrightarrow{OB} = \vec{b}\)，则 \(\overrightarrow{BA} = \vec{a} - \vec{b}\)。

\subsubsection{15.1.3 向量的坐标表示}\label{ux5411ux91cfux7684ux5750ux6807ux8868ux793a}

\textbf{定义 15.8}（基底）：平面上两个不共线的向量 \(\vec{e}_1, \vec{e}_2\) 称为一组\textbf{基底}。对于平面上任意向量 \(\vec{a}\)，存在唯一的实数对 \((x, y)\) 使得： \[\vec{a} = x\vec{e}_1 + y\vec{e}_2\]

\((x, y)\) 称为 \(\vec{a}\) 在基底 \(\{\vec{e}_1, \vec{e}_2\}\) 下的\textbf{坐标}。

当 \(\vec{e}_1 = (1, 0)\)，\(\vec{e}_2 = (0, 1)\) 为标准基底时，\(\vec{a} = (x, y)\)。

\textbf{存在性证明}：设 \(\vec{e}_1, \vec{e}_2\) 不共线，\(\vec{a}\) 为平面上任意向量。将 \(\vec{e}_1, \vec{e}_2\) 放在同一起点 \(O\)。过 \(\vec{a}\) 的终点 \(A\) 作平行于 \(\vec{e}_2\) 的直线，交 \(\vec{e}_1\) 所在直线于点 \(P\)。则 \(\overrightarrow{OP} = x\vec{e}_1\)（某实数 \(x\)），\(\overrightarrow{PA}\) 平行于 \(\vec{e}_2\)，故 \(\overrightarrow{PA} = y\vec{e}_2\)（某实数 \(y\)）。由 \(\vec{a} = \overrightarrow{OA} = \overrightarrow{OP} + \overrightarrow{PA} = x\vec{e}_1 + y\vec{e}_2\)，存在性得证。\(\blacksquare\)

\textbf{定理 15.4}（坐标的唯一性）：在给定基底 \(\{\vec{e}_1, \vec{e}_2\}\) 下，任意向量的坐标是唯一的。

\textbf{证明}：设 \(\vec{a} = x\vec{e}_1 + y\vec{e}_2 = x'\vec{e}_1 + y'\vec{e}_2\)。则 \((x - x')\vec{e}_1 + (y - y')\vec{e}_2 = \vec{0}\)。由于 \(\vec{e}_1\) 和 \(\vec{e}_2\) 不共线，若 \(x - x' \neq 0\) 或 \(y - y' \neq 0\)，则 \(\vec{0}\) 可以表示为非零向量的线性组合，这意味着 \(\vec{e}_1\) 和 \(\vec{e}_2\) 共线，矛盾。因此 \(x = x'\)，\(y = y'\)。\(\blacksquare\)

\textbf{运算的坐标表示}：设 \(\vec{a} = (x_1, y_1)\)，\(\vec{b} = (x_2, y_2)\)，\(\lambda\) 为实数：

\[\vec{a} + \vec{b} = (x_1 + x_2, y_1 + y_2)\] \[\vec{a} - \vec{b} = (x_1 - x_2, y_1 - y_2)\] \[\lambda\vec{a} = (\lambda x_1, \lambda y_1)\] \[|\vec{a}| = \sqrt{x_1^2 + y_1^2}\]

\textbf{定理 15.5}（共线的充要条件）：设 \(\vec{a} = (x_1, y_1)\)，\(\vec{b} = (x_2, y_2)\)，\(\vec{b} \neq \vec{0}\)，则 \(\vec{a}\) 与 \(\vec{b}\) 共线的充要条件是 \(x_1 y_2 - x_2 y_1 = 0\)。

\textbf{证明}：\(\vec{a}\) 与 \(\vec{b}\) 共线 \(\iff\) 存在实数 \(\lambda\) 使得 \(\vec{a} = \lambda\vec{b}\) \(\iff\) \((x_1, y_1) = (\lambda x_2, \lambda y_2)\) \(\iff\) \(x_1 = \lambda x_2\) 且 \(y_1 = \lambda y_2\)。

若 \(x_2 \neq 0\)，则 \(\lambda = \frac{x_1}{x_2}\)，代入得 \(y_1 = \frac{x_1}{x_2} y_2\)，即 \(x_1 y_2 = x_2 y_1\)，即 \(x_1 y_2 - x_2 y_1 = 0\)。

若 \(x_2 = 0\)，则 \(x_1 = \lambda \cdot 0 = 0\)，且 \(y_1 = \lambda y_2\)，此时 \(x_1 y_2 - x_2 y_1 = 0 \cdot y_2 - 0 \cdot y_1 = 0\)。\(\blacksquare\)

\subsubsection{15.1.4 平面向量的数量积}\label{ux5e73ux9762ux5411ux91cfux7684ux6570ux91cfux79ef}

\textbf{定义 15.9}（数量积）：设 \(\vec{a}\) 和 \(\vec{b}\) 是两个非零向量，\(\theta\) 为它们的夹角（\(0 \leq \theta \leq \pi\)），定义：

\[\vec{a} \cdot \vec{b} = |\vec{a}||\vec{b}|\cos\theta\]

规定零向量与任何向量的数量积为 \(0\)。

\textbf{定理 15.6}（数量积的坐标公式）：设 \(\vec{a} = (x_1, y_1)\)，\(\vec{b} = (x_2, y_2)\)，则 \(\vec{a} \cdot \vec{b} = x_1 x_2 + y_1 y_2\)。

\textbf{证明}：设标准基底 \(\vec{i} = (1, 0)\)，\(\vec{j} = (0, 1)\)。则 \(\vec{a} = x_1\vec{i} + y_1\vec{j}\)，\(\vec{b} = x_2\vec{i} + y_2\vec{j}\)。

\[\vec{a} \cdot \vec{b} = (x_1\vec{i} + y_1\vec{j}) \cdot (x_2\vec{i} + y_2\vec{j})\]

利用分配律： \[= x_1 x_2(\vec{i} \cdot \vec{i}) + x_1 y_2(\vec{i} \cdot \vec{j}) + y_1 x_2(\vec{j} \cdot \vec{i}) + y_1 y_2(\vec{j} \cdot \vec{j})\]

由于 \(\vec{i} \cdot \vec{i} = 1\)，\(\vec{j} \cdot \vec{j} = 1\)，\(\vec{i} \cdot \vec{j} = 1 \cdot 1 \cdot \cos 90° = 0\)：

\[\vec{a} \cdot \vec{b} = x_1 x_2 + y_1 y_2\]

\(\blacksquare\)

\textbf{定理 15.7}（数量积的性质）：设 \(\vec{a}, \vec{b}, \vec{c}\) 为向量，\(\lambda \in \mathbb{R}\)，则：

\begin{enumerate}
\def\labelenumi{\arabic{enumi}.}
\tightlist
\item
  \textbf{交换律}：\(\vec{a} \cdot \vec{b} = \vec{b} \cdot \vec{a}\)
\item
  \textbf{分配律}：\(\vec{a} \cdot (\vec{b} + \vec{c}) = \vec{a} \cdot \vec{b} + \vec{a} \cdot \vec{c}\)
\item
  \textbf{数乘结合律}：\((\lambda\vec{a}) \cdot \vec{b} = \lambda(\vec{a} \cdot \vec{b})\)
\item
  \textbf{正定性}：\(\vec{a} \cdot \vec{a} = |\vec{a}|^2 \geq 0\)，等号当且仅当 \(\vec{a} = \vec{0}\)
\end{enumerate}

\textbf{证明}：

\begin{enumerate}
\def\labelenumi{(\arabic{enumi})}
\item
  \(\vec{a} \cdot \vec{b} = |\vec{a}||\vec{b}|\cos\theta = |\vec{b}||\vec{a}|\cos\theta = \vec{b} \cdot \vec{a}\)。
\item
  由坐标公式，设 \(\vec{a} = (a_1, a_2)\)，\(\vec{b} = (b_1, b_2)\)，\(\vec{c} = (c_1, c_2)\)：
\end{enumerate}

\[\vec{a} \cdot (\vec{b} + \vec{c}) = a_1(b_1 + c_1) + a_2(b_2 + c_2) = (a_1 b_1 + a_2 b_2) + (a_1 c_1 + a_2 c_2) = \vec{a} \cdot \vec{b} + \vec{a} \cdot \vec{c}\]

\begin{enumerate}
\def\labelenumi{(\arabic{enumi})}
\setcounter{enumi}{2}
\item
  \((\lambda\vec{a}) \cdot \vec{b} = (\lambda a_1)b_1 + (\lambda a_2)b_2 = \lambda(a_1 b_1 + a_2 b_2) = \lambda(\vec{a} \cdot \vec{b})\)。
\item
  \(\vec{a} \cdot \vec{a} = x_1^2 + y_1^2 = |\vec{a}|^2 \geq 0\)，等号当且仅当 \(x_1 = y_1 = 0\)，即 \(\vec{a} = \vec{0}\)。\(\blacksquare\)
\end{enumerate}

\textbf{推论 15.1}（夹角公式）：\(\cos\theta = \dfrac{\vec{a} \cdot \vec{b}}{|\vec{a}||\vec{b}|} = \dfrac{x_1 x_2 + y_1 y_2}{\sqrt{x_1^2 + y_1^2}\sqrt{x_2^2 + y_2^2}}\)。

\textbf{证明}：由数量积定义 \(\vec{a} \cdot \vec{b} = |\vec{a}||\vec{b}|\cos\theta\) 和定理 15.6 中 \(\vec{a} \cdot \vec{b} = x_1x_2 + y_1y_2\)，立即得到结论。\(\blacksquare\)

\textbf{推论 15.2}（垂直的充要条件）：非零向量 \(\vec{a} \perp \vec{b}\) 的充要条件是 \(\vec{a} \cdot \vec{b} = 0\)，即 \(x_1 x_2 + y_1 y_2 = 0\)。

\textbf{证明}：\(\vec{a} \perp \vec{b} \iff \theta = 90° \iff \cos\theta = 0 \iff |\vec{a}||\vec{b}|\cos\theta = 0 \iff \vec{a} \cdot \vec{b} = 0\)。\(\blacksquare\)

\subsection{15.2 空间向量}\label{ux7a7aux95f4ux5411ux91cf}

空间向量是平面向量的自然推广。在三维空间中，向量依然保持大小和方向这两个基本属性，但坐标系从二维扩展为三维，向量的坐标表示从 \((x, y)\) 变为 \((x, y, z)\)。空间向量理论的建立使得我们能够用统一的代数方法处理三维几何问题------包括点到平面的距离、异面直线的夹角、二面角等平面几何中难以直接处理的复杂空间关系。空间向量的数量积（点积）\(\vec{a} \cdot \vec{b} = x_1x_2 + y_1y_2 + z_1z_2\) 保持了二维情形下的所有性质（交换律、分配律、与模长的关系），而向量积（叉积）\(\vec{a} \times \vec{b}\) 则是空间向量特有的运算------它产生一个同时垂直于 \(\vec{a}\) 和 \(\vec{b}\) 的向量，其模长等于以 \(\vec{a}\) 和 \(\vec{b}\) 为邻边的平行四边形的面积。向量积的引入使得法向量的计算变得极为便捷，从而为空间中的平面方程和直线方程提供了统一的代数描述。空间向量的基本定理------任意空间向量可唯一表示为三个不共面向量的线性组合------是三维空间向量理论的基石，它保证了三维空间中基的存在性和坐标表示的唯一性。

\subsubsection{15.2.1 空间向量的概念与运算}\label{ux7a7aux95f4ux5411ux91cfux7684ux6982ux5ff5ux4e0eux8fd0ux7b97}

\textbf{定义 15.10}（空间向量）：在空间中，既有大小又有方向的量称为空间向量。空间向量用有向线段表示，记为 \(\overrightarrow{AB}\) 或 \(\vec{a}\)。

\textbf{定义 15.11}（空间向量的坐标表示）：在空间直角坐标系 \(Oxyz\) 中，设 \(\vec{e}_1, \vec{e}_2, \vec{e}_3\) 分别为 \(x\) 轴、\(y\) 轴、\(z\) 轴正方向的单位向量（标准基向量），则任一空间向量 \(\vec{a}\) 可唯一表示为

\[\vec{a} = x\vec{e}_1 + y\vec{e}_2 + z\vec{e}_3\]

其中有序实数组 \((x, y, z)\) 称为向量 \(\vec{a}\) 的\textbf{坐标}，记为 \(\vec{a} = (x, y, z)\)。

\textbf{存在性}：过 \(\vec{a}\) 的终点 \(A\) 作三个平面分别平行于坐标平面 \(yOz\)、\(xOz\)、\(xOy\)，与坐标轴分别交于 \(P_1\)（\(x\) 轴）、\(P_2\)（\(y\) 轴）、\(P_3\)（\(z\) 轴）。则 \(\overrightarrow{OP_1} = x\vec{e}_1\)，\(\overrightarrow{P_1Q} = y\vec{e}_2\)（\(Q\) 为 \(A\) 在 \(xOy\) 平面上的投影），\(\overrightarrow{QA} = z\vec{e}_3\)。由向量加法，\(\vec{a} = x\vec{e}_1 + y\vec{e}_2 + z\vec{e}_3\)。唯一性的证明与二维情形类似。

\textbf{定义 15.12}（空间向量的模）：空间向量 \(\vec{a} = (x, y, z)\) 的\textbf{模}定义为

\[|\vec{a}| = \sqrt{x^2 + y^2 + z^2}\]

\textbf{证明}：由勾股定理，从原点到点 \((x, y, z)\) 的距离为 \(|\vec{a}|^2 = x^2 + y^2 + z^2\)。\(\blacksquare\)

空间向量的加法、减法、数乘的坐标表示与平面向量类似：

\[\vec{a} + \vec{b} = (x_1 + x_2, y_1 + y_2, z_1 + z_2)\] \[\vec{a} - \vec{b} = (x_1 - x_2, y_1 - y_2, z_1 - z_2)\] \[\lambda\vec{a} = (\lambda x_1, \lambda y_1, \lambda z_1)\]

\textbf{定理 15.8}（空间向量共线的充要条件）：非零向量 \(\vec{a} = (x_1, y_1, z_1)\) 与 \(\vec{b} = (x_2, y_2, z_2)\) 共线的充要条件是 \(\frac{x_2}{x_1} = \frac{y_2}{y_1} = \frac{z_2}{z_1}\)（当某个分量为零时，对应分量也须为零）。

\textbf{证明}：\(\vec{a} \parallel \vec{b}\) 当且仅当 \(\vec{a} = t\vec{b}\) 对某标量 \(t\)，即 \(x_1 = tx_2, y_1 = ty_2, z_1 = tz_2\)。消去 \(t\) 得 \(x_1y_2 = x_2y_1, y_1z_2 = y_2z_1, x_1z_2 = x_2z_1\)。\(\blacksquare\)

\subsubsection{15.2.2 空间向量的数量积}\label{ux7a7aux95f4ux5411ux91cfux7684ux6570ux91cfux79ef}

\textbf{定义 15.13}（空间向量的数量积）：设 \(\vec{a}\) 与 \(\vec{b}\) 是空间中两个非零向量，\(\theta\) 为它们的夹角（\(0 \leq \theta \leq \pi\)），定义

\[\vec{a} \cdot \vec{b} = |\vec{a}||\vec{b}|\cos\theta\]

\textbf{定理 15.9}（空间向量数量积的坐标公式）：设 \(\vec{a} = (x_1, y_1, z_1)\)，\(\vec{b} = (x_2, y_2, z_2)\)，则

\[\vec{a} \cdot \vec{b} = x_1 x_2 + y_1 y_2 + z_1 z_2\]

\textbf{证明}：设 \(\vec{a} = x_1\vec{e}_1 + y_1\vec{e}_2 + z_1\vec{e}_3\)，\(\vec{b} = x_2\vec{e}_1 + y_2\vec{e}_2 + z_2\vec{e}_3\)。由分配律展开：

\[\vec{a} \cdot \vec{b} = x_1 x_2(\vec{e}_1 \cdot \vec{e}_1) + x_1 y_2(\vec{e}_1 \cdot \vec{e}_2) + \cdots\]

由于 \(\vec{e}_1, \vec{e}_2, \vec{e}_3\) 是标准正交基：

\[\vec{e}_i \cdot \vec{e}_j = \begin{cases} 1, & i = j \\ 0, & i \neq j \end{cases}\]

因此交叉项全部为零，得 \(\vec{a} \cdot \vec{b} = x_1 x_2 + y_1 y_2 + z_1 z_2\)。\(\blacksquare\)

\textbf{推论 15.3}（空间向量垂直的充要条件）：\(\vec{a} \perp \vec{b} \iff \vec{a} \cdot \vec{b} = 0 \iff x_1 x_2 + y_1 y_2 + z_1 z_2 = 0\)。

\textbf{证明}：与推论 15.2 类似，\(\vec{a} \perp \vec{b} \iff \theta = 90° \iff \cos\theta = 0 \iff \vec{a} \cdot \vec{b} = 0\)，再由定理 15.9 得 \(\vec{a} \cdot \vec{b} = x_1x_2 + y_1y_2 + z_1z_2 = 0\)。\(\blacksquare\)

\textbf{推论 15.4}（空间夹角公式）：\(\cos\theta = \dfrac{x_1 x_2 + y_1 y_2 + z_1 z_2}{\sqrt{x_1^2 + y_1^2 + z_1^2}\sqrt{x_2^2 + y_2^2 + z_2^2}}\)。

\textbf{证明}：与推论 15.1 类似，由空间数量积定义 \(\vec{a} \cdot \vec{b} = |\vec{a}||\vec{b}|\cos\theta\) 和定理 15.9 中 \(\vec{a} \cdot \vec{b} = x_1x_2 + y_1y_2 + z_1z_2\)，立即得到结论。\(\blacksquare\)

\subsubsection{15.2.3 空间向量的向量积}\label{ux7a7aux95f4ux5411ux91cfux7684ux5411ux91cfux79ef}

\textbf{定义 15.14}（向量积）：设 \(\vec{a}\) 与 \(\vec{b}\) 是空间中两个不共线的向量，定义它们的\textbf{向量积} \(\vec{a} \times \vec{b}\) 为满足以下条件的向量：

\begin{enumerate}
\def\labelenumi{\arabic{enumi}.}
\tightlist
\item
  \textbf{方向}：\(\vec{a} \times \vec{b}\) 垂直于 \(\vec{a}\) 和 \(\vec{b}\) 所确定的平面，且 \(\vec{a}, \vec{b}, \vec{a} \times \vec{b}\) 构成右手系。
\item
  \textbf{大小}：\(|\vec{a} \times \vec{b}| = |\vec{a}||\vec{b}|\sin\theta\)，其中 \(\theta\) 为 \(\vec{a}\) 与 \(\vec{b}\) 的夹角。
\end{enumerate}

若 \(\vec{a}\) 与 \(\vec{b}\) 共线，则规定 \(\vec{a} \times \vec{b} = \vec{0}\)。

向量积是空间向量特有的运算，它在二维平面中不存在------这正是三维空间几何丰富性的根源之一。与数量积产生标量不同，向量积产生一个垂直于两个原始向量的新向量。为了在坐标计算中方便地使用向量积，我们需要将上述几何定义转化为代数公式。

\textbf{定理 15.10}（向量积的坐标公式）：设 \(\vec{a} = (x_1, y_1, z_1)\)，\(\vec{b} = (x_2, y_2, z_2)\)，则

\[\vec{a} \times \vec{b} = (y_1 z_2 - z_1 y_2, \; z_1 x_2 - x_1 z_2, \; x_1 y_2 - y_1 x_2)\]

\textbf{证明}：由标准基的向量积关系：

\[\vec{e}_1 \times \vec{e}_2 = \vec{e}_3, \quad \vec{e}_2 \times \vec{e}_3 = \vec{e}_1, \quad \vec{e}_3 \times \vec{e}_1 = \vec{e}_2\] \[\vec{e}_2 \times \vec{e}_1 = -\vec{e}_3, \quad \vec{e}_3 \times \vec{e}_2 = -\vec{e}_1, \quad \vec{e}_1 \times \vec{e}_3 = -\vec{e}_2\] \[\vec{e}_i \times \vec{e}_i = \vec{0}\]

展开 \(\vec{a} \times \vec{b} = (x_1\vec{e}_1 + y_1\vec{e}_2 + z_1\vec{e}_3) \times (x_2\vec{e}_1 + y_2\vec{e}_2 + z_2\vec{e}_3)\)，利用上述关系逐项计算：

\[= (y_1 z_2 - z_1 y_2)\vec{e}_1 + (z_1 x_2 - x_1 z_2)\vec{e}_2 + (x_1 y_2 - y_1 x_2)\vec{e}_3\]

\(\blacksquare\)

\textbf{定理 15.11}（向量积的性质）：

\begin{enumerate}
\def\labelenumi{\arabic{enumi}.}
\tightlist
\item
  \textbf{反交换律}：\(\vec{a} \times \vec{b} = -(\vec{b} \times \vec{a})\)
\item
  \textbf{自叉积为零}：\(\vec{a} \times \vec{a} = \vec{0}\)
\item
  \textbf{分配律}：\(\vec{a} \times (\vec{b} + \vec{c}) = \vec{a} \times \vec{b} + \vec{a} \times \vec{c}\)
\end{enumerate}

\textbf{证明}：

\begin{enumerate}
\def\labelenumi{(\arabic{enumi})}
\item
  由坐标公式，交换 \(\vec{a}\) 与 \(\vec{b}\) 即每个分量变号，故 \(\vec{a} \times \vec{b} = -(\vec{b} \times \vec{a})\)。
\item
  由反交换律：\(\vec{a} \times \vec{a} = -(\vec{a} \times \vec{a})\)，故 \(2(\vec{a} \times \vec{a}) = \vec{0}\)，即 \(\vec{a} \times \vec{a} = \vec{0}\)。
\item
  可由坐标公式验证。\(\blacksquare\)
\end{enumerate}

\textbf{定理 15.12}（向量积的几何意义）：\(|\vec{a} \times \vec{b}|\) 等于以 \(\vec{a}\) 和 \(\vec{b}\) 为邻边的平行四边形的面积。

\textbf{证明}：以 \(\vec{a}\) 和 \(\vec{b}\) 为邻边的平行四边形，底边长为 \(|\vec{a}|\)，高为 \(|\vec{b}|\sin\theta\)，故面积为 \(|\vec{a}||\vec{b}|\sin\theta = |\vec{a} \times \vec{b}|\)。\(\blacksquare\)

\textbf{推论 15.5}：以 \(\vec{a}\) 和 \(\vec{b}\) 为邻边的三角形面积为 \(S_{\triangle} = \frac{1}{2}|\vec{a} \times \vec{b}|\)。

\textbf{证明}：以 \(\vec{a}, \vec{b}\) 为邻边的平行四边形面积为 \(|\vec{a} \times \vec{b}|\)（定理 15.12），三角形为其一半。\(\blacksquare\)

\subsection{15.3 向量法解立体几何}\label{ux5411ux91cfux6cd5ux89e3ux7acbux4f53ux51e0ux4f55}

向量法的核心优势在于将几何问题转化为代数计算，避免了纯几何方法中繁琐的辅助线构造。本节系统介绍如何利用向量工具解决立体几何中的经典问题。方向向量和法向量是向量法的两个基本工具：直线的方向向量刻画了直线的走向，平面的法向量（垂直于平面内所有直线的向量）则刻画了平面的朝向。利用这两个工具，两条直线的夹角可以转化为方向向量的夹角，直线与平面的夹角可以转化为方向向量与法向量夹角的余角，两个平面的夹角可以转化为法向量的夹角。点到平面的距离公式 \(d = \frac{|\overrightarrow{AP} \cdot \vec{n}|}{|\vec{n}|}\) 是向量法最优雅的应用之一------它将纯几何中需要构造垂线的复杂问题简化为一个数量积的计算。异面直线的距离（公垂线段的长度）也可以用向量法简洁地表达，这体现了向量法在处理复杂空间关系时的系统性优势。向量法的深层意义在于：它揭示了立体几何与线性代数之间的内在联系，空间中的几何关系本质上对应着向量空间中的线性关系。

\subsubsection{15.3.1 方向向量与法向量}\label{ux65b9ux5411ux5411ux91cfux4e0eux6cd5ux5411ux91cf}

\textbf{定义 15.15}（方向向量）：空间中直线 \(l\) 的\textbf{方向向量}是指与直线 \(l\) 平行的非零向量 \(\vec{d}\)。一条直线的方向向量有无穷多个，但它们两两共线。

\textbf{定义 15.16}（法向量）：平面 \(\alpha\) 的\textbf{法向量}是指垂直于平面 \(\alpha\) 的非零向量 \(\vec{n}\)。一个平面的法向量有无穷多个，但它们两两共线。

\textbf{定理 15.13}（法向量的确定）：设平面 \(\alpha\) 由三个不共线的点 \(A\)、\(B\)、\(C\) 确定，取 \(\vec{u} = \overrightarrow{AB}\)，\(\vec{v} = \overrightarrow{AC}\)，则 \(\alpha\) 的法向量为

\[\vec{n} = \vec{u} \times \vec{v}\]

\textbf{证明}：由向量积的定义，\(\vec{u} \times \vec{v}\) 同时垂直于 \(\vec{u}\) 和 \(\vec{v}\)。由于 \(\vec{u}\) 和 \(\vec{v}\) 不共线且都在平面 \(\alpha\) 内，它们张成整个平面 \(\alpha\)，故 \(\vec{n}\) 垂直于 \(\alpha\) 内的任意向量，即 \(\vec{n}\) 是 \(\alpha\) 的法向量。\(\blacksquare\)

若平面 \(\alpha\) 由方程 \(ax + by + cz = d\) 给出，则法向量为 \(\vec{n} = (a, b, c)\)。

\subsubsection{15.3.2 平行与垂直的向量判定}\label{ux5e73ux884cux4e0eux5782ux76f4ux7684ux5411ux91cfux5224ux5b9a}

利用方向向量和法向量，可以将空间中的平行与垂直关系全部转化为向量运算。

\textbf{定理 15.14}（线线平行）：设 \(l_1\) 的方向向量为 \(\vec{d}_1\)，\(l_2\) 的方向向量为 \(\vec{d}_2\)，则 \(l_1 \parallel l_2\) 的充要条件是 \(\vec{d}_1 \parallel \vec{d}_2\)，即存在 \(\lambda \neq 0\) 使得 \(\vec{d}_2 = \lambda\vec{d}_1\)。

\textbf{证明}：\(l_1 \parallel l_2\) 当且仅当 \(l_1\) 与 \(l_2\) 方向相同（或相反），当且仅当 \(\vec{d}_1\) 与 \(\vec{d}_2\) 共线。\(\blacksquare\)

\textbf{定理 15.15}（线面平行）：设直线 \(l\) 的方向向量为 \(\vec{d}\)，平面 \(\alpha\) 的法向量为 \(\vec{n}_\alpha\)，则 \(l \parallel \alpha\)（\(l \not\subset \alpha\)）的充要条件是 \(\vec{d} \cdot \vec{n}_\alpha = 0\)。

\textbf{证明}：\(l \parallel \alpha\) 当且仅当 \(l\) 的方向平行于 \(\alpha\)，当且仅当 \(\vec{d}\) 垂直于 \(\alpha\) 的法向量 \(\vec{n}_\alpha\)，当且仅当 \(\vec{d} \cdot \vec{n}_\alpha = 0\)。\(\blacksquare\)

\textbf{定理 15.16}（面面平行）：设 \(\alpha\) 的法向量为 \(\vec{n}_\alpha\)，\(\beta\) 的法向量为 \(\vec{n}_\beta\)，则 \(\alpha \parallel \beta\) 的充要条件是 \(\vec{n}_\alpha \parallel \vec{n}_\beta\)，即存在 \(\lambda \neq 0\) 使得 \(\vec{n}_\beta = \lambda\vec{n}_\alpha\)。

\textbf{证明}：\(\alpha \parallel \beta\) 当且仅当两个平面有相同的''倾斜方向''，当且仅当法向量平行。\(\blacksquare\)

\textbf{定理 15.17}（线线垂直）：\(l_1 \perp l_2\) 的充要条件是 \(\vec{d}_1 \cdot \vec{d}_2 = 0\)。

\textbf{证明}：\(l_1 \perp l_2 \iff \theta = 90° \iff \cos\theta = 0 \iff \vec{d}_1 \cdot \vec{d}_2 = 0\)。\(\blacksquare\)

\textbf{定理 15.18}（线面垂直）：\(l \perp \alpha\) 的充要条件是 \(\vec{d} \parallel \vec{n}_\alpha\)，即存在 \(\lambda \neq 0\) 使得 \(\vec{d} = \lambda\vec{n}_\alpha\)。

\textbf{证明}：\(l \perp \alpha\) 当且仅当 \(l\) 垂直于 \(\alpha\) 内的所有直线，当且仅当 \(\vec{d}\) 与 \(\alpha\) 的法向量平行。\(\blacksquare\)

\textbf{定理 15.19}（面面垂直）：\(\alpha \perp \beta\) 的充要条件是 \(\vec{n}_\alpha \cdot \vec{n}_\beta = 0\)。

\textbf{证明}：\(\alpha \perp \beta\) 当且仅当二面角的平面角为 \(90°\)。两个法向量的夹角 \(\gamma\) 满足 \(\cos\gamma = \frac{\vec{n}_\alpha \cdot \vec{n}_\beta}{|\vec{n}_\alpha||\vec{n}_\beta|}\)。二面角的平面角 \(\psi\) 与 \(\gamma\) 的关系是 \(\psi = \gamma\) 或 \(\psi = \pi - \gamma\)。当 \(\psi = 90°\) 时，\(|\cos\gamma| = 0\)，即 \(\vec{n}_\alpha \cdot \vec{n}_\beta = 0\)。\(\blacksquare\)

\textbf{总结表}：

\begin{longtable}[]{@{}ll@{}}
\toprule
位置关系 & 向量条件 \\
\midrule
\endhead
\bottomrule
\endlastfoot
\(l_1 \parallel l_2\) & \(\vec{d}_1 \parallel \vec{d}_2\) \\
\(l \parallel \alpha\) & \(\vec{d} \cdot \vec{n}_\alpha = 0\) \\
\(\alpha \parallel \beta\) & \(\vec{n}_\alpha \parallel \vec{n}_\beta\) \\
\(l_1 \perp l_2\) & \(\vec{d}_1 \cdot \vec{d}_2 = 0\) \\
\(l \perp \alpha\) & \(\vec{d} \parallel \vec{n}_\alpha\) \\
\(\alpha \perp \beta\) & \(\vec{n}_\alpha \cdot \vec{n}_\beta = 0\) \\
\end{longtable}

\subsubsection{15.3.3 空间角的向量公式}\label{ux7a7aux95f4ux89d2ux7684ux5411ux91cfux516cux5f0f}

\textbf{定理 15.20}（异面直线所成的角）：设 \(l_1\) 的方向向量为 \(\vec{d}_1\)，\(l_2\) 的方向向量为 \(\vec{d}_2\)，则异面直线 \(l_1\) 与 \(l_2\) 所成的角 \(\phi\) 满足

\[\cos\phi = \frac{|\vec{d}_1 \cdot \vec{d}_2|}{|\vec{d}_1||\vec{d}_2|}\]

\textbf{证明}：设 \(\vec{d}_1\) 与 \(\vec{d}_2\) 的夹角为 \(\theta\)（\(0 \leq \theta \leq \pi\)）。异面直线所成的角取锐角或直角，即 \(\phi = \theta\)（若 \(\theta \leq 90°\)）或 \(\phi = 180° - \theta\)（若 \(\theta > 90°\)）。无论哪种情况，\(\cos\phi = |\cos\theta| = \frac{|\vec{d}_1 \cdot \vec{d}_2|}{|\vec{d}_1||\vec{d}_2|}\)。\(\blacksquare\)

\textbf{定理 15.21}（直线与平面所成的角）：设直线 \(l\) 的方向向量为 \(\vec{d}\)，平面 \(\alpha\) 的法向量为 \(\vec{n}\)，则直线 \(l\) 与平面 \(\alpha\) 所成的角 \(\varphi\) 满足

\[\sin\varphi = \frac{|\vec{d} \cdot \vec{n}|}{|\vec{d}||\vec{n}|}\]

\textbf{证明}：设 \(\vec{d}\) 与 \(\vec{n}\) 的夹角为 \(\gamma\)。由于 \(\vec{n}\) 垂直于 \(\alpha\)，\(l\) 在 \(\alpha\) 上的射影 \(l'\) 的方向向量 \(\vec{d}'\) 满足 \(\vec{d}' \perp \vec{n}\)。\(l\) 与 \(l'\) 所成的角 \(\varphi\) 满足 \(\varphi = |90° - \gamma|\)。因此

\[\sin\varphi = \sin|90° - \gamma| = |\cos\gamma| = \frac{|\vec{d} \cdot \vec{n}|}{|\vec{d}||\vec{n}|}\]

\(\blacksquare\)

\textbf{定理 15.22}（二面角）：设平面 \(\alpha\) 的法向量为 \(\vec{n}_1\)，平面 \(\beta\) 的法向量为 \(\vec{n}_2\)，则二面角 \(\alpha\)-\(l\)-\(\beta\) 的平面角 \(\psi\) 满足

\[\cos\psi = \pm\frac{\vec{n}_1 \cdot \vec{n}_2}{|\vec{n}_1||\vec{n}_2|}\]

其中正负号取决于法向量的选取方向。

\textbf{证明}：两个法向量 \(\vec{n}_1\) 与 \(\vec{n}_2\) 的夹角 \(\gamma\) 满足 \(\cos\gamma = \frac{\vec{n}_1 \cdot \vec{n}_2}{|\vec{n}_1||\vec{n}_2|}\)。二面角的平面角 \(\psi\) 与 \(\gamma\) 的关系是 \(\psi = \gamma\) 或 \(\psi = \pi - \gamma\)，取决于法向量的选取方向是否指向二面角的同一侧。因此 \(\cos\psi = \pm\cos\gamma\)。在实际计算中，通常直接计算 \(\cos\gamma\)，再根据几何意义判断正负号。\(\blacksquare\)

\subsubsection{15.3.4 空间距离的向量公式}\label{ux7a7aux95f4ux8dddux79bbux7684ux5411ux91cfux516cux5f0f}

\textbf{定理 15.23}（点到平面的距离）：设平面 \(\alpha\) 的法向量为 \(\vec{n}\)，\(P\) 为空间中一点，\(A\) 为平面 \(\alpha\) 上的任意一点，则点 \(P\) 到平面 \(\alpha\) 的距离为

\[d = \frac{|\overrightarrow{AP} \cdot \vec{n}|}{|\vec{n}|}\]

\textbf{证明}：设 \(H\) 为 \(P\) 在平面 \(\alpha\) 上的正投影（垂足），则 \(\overrightarrow{PH} \parallel \vec{n}\)。设 \(\overrightarrow{PH} = t\vec{n}\)，则

\[\overrightarrow{AH} = \overrightarrow{AP} + \overrightarrow{PH} = \overrightarrow{AP} + t\vec{n}\]

由于 \(H \in \alpha\)，\(\overrightarrow{AH}\) 在 \(\alpha\) 内，故 \(\overrightarrow{AH} \perp \vec{n}\)，即

\[(\overrightarrow{AP} + t\vec{n}) \cdot \vec{n} = 0 \implies \overrightarrow{AP} \cdot \vec{n} + t|\vec{n}|^2 = 0 \implies t = -\frac{\overrightarrow{AP} \cdot \vec{n}}{|\vec{n}|^2}\]

因此

\[d = |\overrightarrow{PH}| = |t| \cdot |\vec{n}| = \frac{|\overrightarrow{AP} \cdot \vec{n}|}{|\vec{n}|}\]

\(\blacksquare\)

\textbf{定理 15.24}（异面直线的距离）：设 \(l_1\) 过点 \(A\)，方向向量为 \(\vec{d}_1\)；\(l_2\) 过点 \(B\)，方向向量为 \(\vec{d}_2\)。则 \(l_1\) 与 \(l_2\) 的距离为

\[d = \frac{|\overrightarrow{AB} \cdot (\vec{d}_1 \times \vec{d}_2)|}{|\vec{d}_1 \times \vec{d}_2|}\]

\textbf{证明}：公垂线的方向向量为 \(\vec{d}_1 \times \vec{d}_2\)（同时垂直于 \(\vec{d}_1\) 和 \(\vec{d}_2\)）。异面直线的距离等于 \(\overrightarrow{AB}\) 在公垂线方向上的投影的绝对值：

\[d = \left|\text{proj}_{\vec{d}_1 \times \vec{d}_2}\overrightarrow{AB}\right| = \frac{|\overrightarrow{AB} \cdot (\vec{d}_1 \times \vec{d}_2)|}{|\vec{d}_1 \times \vec{d}_2|}\]

\(\blacksquare\)

\subsection{15.4 向量法证明几何定理}\label{ux5411ux91cfux6cd5ux8bc1ux660eux51e0ux4f55ux5b9aux7406}

向量法为证明经典几何定理提供了简洁有力的工具。以下展示如何用向量方法重新证明若干重要定理。

\textbf{定理 15.25}（用向量法证明余弦定理）：在 \(\triangle ABC\) 中，\(a^2 = b^2 + c^2 - 2bc\cos A\)。

\textbf{证明}：设 \(\overrightarrow{AB} = \vec{c}\)，\(\overrightarrow{AC} = \vec{b}\)，则 \(\overrightarrow{BC} = \vec{b} - \vec{c}\)。

\[a^2 = |\overrightarrow{BC}|^2 = (\vec{b} - \vec{c}) \cdot (\vec{b} - \vec{c}) = |\vec{b}|^2 - 2\vec{b} \cdot \vec{c} + |\vec{c}|^2\]

\[= b^2 + c^2 - 2|\vec{b}||\vec{c}|\cos A = b^2 + c^2 - 2bc\cos A\]

\(\blacksquare\)

\textbf{定理 15.26}（用向量法证明平行四边形对角线性质）：平行四边形的对角线互相平分。

\textbf{证明}：设平行四边形 \(ABCD\)，\(\overrightarrow{AB} = \vec{a}\)，\(\overrightarrow{AD} = \vec{b}\)。则 \(\overrightarrow{AC} = \vec{a} + \vec{b}\)，\(\overrightarrow{DB} = \vec{a} - \vec{b}\)。

设对角线 \(AC\) 和 \(BD\) 的交点为 \(O\)。设 \(\overrightarrow{AO} = t(\vec{a} + \vec{b})\)（\(O\) 在 \(AC\) 上），\(\overrightarrow{DO} = s(\vec{a} - \vec{b})\)（\(O\) 在 \(BD\) 上）。

由 \(\overrightarrow{AO} = \overrightarrow{AD} + \overrightarrow{DO}\)，得 \(t(\vec{a} + \vec{b}) = \vec{b} + s(\vec{a} - \vec{b})\)。

比较 \(\vec{a}\) 和 \(\vec{b}\) 的系数：\(t = s\)，\(t = 1 - s\)。解得 \(t = s = \frac{1}{2}\)。

因此 \(\overrightarrow{AO} = \frac{1}{2}\overrightarrow{AC}\)，\(\overrightarrow{DO} = \frac{1}{2}\overrightarrow{DB}\)，即 \(O\) 是两条对角线的中点。\(\blacksquare\)

\textbf{定理 15.27}（用向量法证明三角形中位线定理）：三角形的中位线平行于第三边且等于第三边的一半。

\textbf{证明}：设 \(\triangle ABC\) 中，\(D\)、\(E\) 分别是 \(\overline{AB}\)、\(\overline{AC}\) 的中点。设 \(\overrightarrow{AB} = \vec{a}\)，\(\overrightarrow{AC} = \vec{b}\)。

\[\overrightarrow{AD} = \frac{1}{2}\vec{a}, \quad \overrightarrow{AE} = \frac{1}{2}\vec{b}\]

\[\overrightarrow{DE} = \overrightarrow{AE} - \overrightarrow{AD} = \frac{1}{2}\vec{b} - \frac{1}{2}\vec{a} = \frac{1}{2}(\vec{b} - \vec{a}) = \frac{1}{2}\overrightarrow{BC}\]

因此 \(\overrightarrow{DE} \parallel \overrightarrow{BC}\) 且 \(|\overrightarrow{DE}| = \frac{1}{2}|\overrightarrow{BC}|\)。\(\blacksquare\)

\textbf{定理 15.28}（用向量法证明勾股定理）：在直角三角形中，斜边的平方等于两直角边的平方和。

\textbf{证明}：设 \(\triangle ABC\) 中 \(\angle C = 90°\)。设 \(\overrightarrow{CA} = \vec{a}\)，\(\overrightarrow{CB} = \vec{b}\)，则 \(\overrightarrow{AB} = \vec{b} - \vec{a}\)。

由 \(\angle C = 90°\)，\(\vec{a} \perp \vec{b}\)，即 \(\vec{a} \cdot \vec{b} = 0\)。

\[c^2 = |\overrightarrow{AB}|^2 = (\vec{b} - \vec{a}) \cdot (\vec{b} - \vec{a}) = |\vec{a}|^2 - 2\vec{a} \cdot \vec{b} + |\vec{b}|^2 = a^2 + b^2\]

\(\blacksquare\)

\textbf{定理 15.29}（用向量法证明三垂线定理）：如果平面内的一条直线与一条斜线在平面内的射影垂直，则它也与这条斜线垂直。

\textbf{证明}：设 \(PO \perp \alpha\) 于 \(O\)，\(A \in \alpha\)，\(a \subset \alpha\)，\(a \perp OA\)。在 \(a\) 上取方向向量 \(\vec{d}\)，则 \(\vec{d} \perp \overrightarrow{OA}\)，即 \(\vec{d} \cdot \overrightarrow{OA} = 0\)。

又 \(\overrightarrow{PO} \perp \alpha\)，故 \(\overrightarrow{PO} \perp \vec{d}\)，即 \(\overrightarrow{PO} \cdot \vec{d} = 0\)。

\[\overrightarrow{PA} \cdot \vec{d} = (\overrightarrow{PO} + \overrightarrow{OA}) \cdot \vec{d} = \overrightarrow{PO} \cdot \vec{d} + \overrightarrow{OA} \cdot \vec{d} = 0 + 0 = 0\]

因此 \(\overrightarrow{PA} \perp \vec{d}\)，即 \(PA \perp a\)。\(\blacksquare\)

\textbf{定理 15.30}（用向量法证明正弦定理）：在 \(\triangle ABC\) 中，\(\frac{a}{\sin A} = \frac{b}{\sin B} = \frac{c}{\sin C}\)。

\textbf{证明}：设 \(\overrightarrow{AB} = \vec{c}\)，\(\overrightarrow{AC} = \vec{b}\)。由向量积的几何意义：

\[|\vec{b} \times \vec{c}| = |\vec{b}||\vec{c}|\sin A = bc\sin A\]

而 \(|\vec{b} \times \vec{c}|\) 等于以 \(\overrightarrow{AB}\) 和 \(\overrightarrow{AC}\) 为邻边的平行四边形的面积，即 \(2S_{\triangle}\)（三角形面积的两倍）。

因此 \(bc\sin A = 2S\)，即 \(\frac{a}{\sin A} = \frac{a \cdot bc}{2S}\)。

同理 \(ac\sin B = 2S\)，\(ab\sin C = 2S\)。

所以 \(\frac{a}{\sin A} = \frac{b}{\sin B} = \frac{c}{\sin C} = \frac{abc}{2S}\)。

\textbf{注}：上述向量法证明建立了 \(\frac{a}{\sin A} = \frac{b}{\sin B} = \frac{c}{\sin C} = \frac{abc}{2S}\)。其中 \(\frac{abc}{2S} = 2R\)（\(R\) 为外接圆半径）的结论参见定理 13.40 的几何证明。\(\blacksquare\)

\hrulefill

\subsubsection{15.4.1 混合积}\label{ux6df7ux5408ux79ef}

\textbf{定义 15.17}（混合积）：三个向量 \(\vec{a}, \vec{b}, \vec{c}\) 的\textbf{混合积}（标量三重积）定义为

\[[\vec{a}\;\vec{b}\;\vec{c}] = (\vec{a} \times \vec{b}) \cdot \vec{c}\]

\textbf{定理 15.31}（混合积的坐标公式）：设 \(\vec{a} = (x_1, y_1, z_1)\)，\(\vec{b} = (x_2, y_2, z_2)\)，\(\vec{c} = (x_3, y_3, z_3)\)，则

\[[\vec{a}\;\vec{b}\;\vec{c}] = \begin{vmatrix} x_1 & y_1 & z_1 \\ x_2 & y_2 & z_2 \\ x_3 & y_3 & z_3 \end{vmatrix}\]

\textbf{证明}：由向量积的坐标公式，\(\vec{a} \times \vec{b} = (y_1 z_2 - z_1 y_2, \; z_1 x_2 - x_1 z_2, \; x_1 y_2 - y_1 x_2)\)。再由数量积的坐标公式：

\[[\vec{a}\;\vec{b}\;\vec{c}] = x_3(y_1 z_2 - z_1 y_2) + y_3(z_1 x_2 - x_1 z_2) + z_3(x_1 y_2 - y_1 x_2)\]

按第三行展开行列式即得上述结果。\(\blacksquare\)

\textbf{定理 15.32}（混合积的几何意义）：\(|[\vec{a}\;\vec{b}\;\vec{c}]|\) 等于以 \(\vec{a}, \vec{b}, \vec{c}\) 为棱的平行六面体的体积。

\textbf{证明}：\(\vec{a} \times \vec{b}\) 的模长等于以 \(\vec{a}, \vec{b}\) 为底的平行四边形的面积 \(S\)，方向为底面的法向量。\((\vec{a} \times \vec{b}) \cdot \vec{c}\) 等于 \(\vec{c}\) 在 \(\vec{a} \times \vec{b}\) 方向上的投影与 \(|\vec{a} \times \vec{b}|\) 的乘积，即底面积 \(S\) 乘以高 \(h\)，等于平行六面体的体积。\(\blacksquare\)

\textbf{推论 15.6}（共面的充要条件）：三个向量 \(\vec{a}, \vec{b}, \vec{c}\) 共面的充要条件是 \([\vec{a}\;\vec{b}\;\vec{c}] = 0\)。

\textbf{证明}：\(\vec{a}, \vec{b}, \vec{c}\) 共面 \(\iff\) 以它们为棱的平行六面体体积为零 \(\iff\) \([\vec{a}\;\vec{b}\;\vec{c}] = 0\)。\(\blacksquare\)

\textbf{定理 15.33}（混合积的轮换对称性）：\([\vec{a}\;\vec{b}\;\vec{c}] = [\vec{b}\;\vec{c}\;\vec{a}] = [\vec{c}\;\vec{a}\;\vec{b}]\)。

\textbf{证明}：由行列式的性质，轮换行列式的行不改变行列式的值（或改变偶数次行交换，符号不变）：

\[[\vec{a}\;\vec{b}\;\vec{c}] = (\vec{a} \times \vec{b}) \cdot \vec{c} = \vec{c} \cdot (\vec{a} \times \vec{b}) = [\vec{c}\;\vec{a}\;\vec{b}]\]

其中最后一步利用了混合积的循环性质（行列式的行轮换）。\(\blacksquare\)

\subsubsection{15.4.2 用向量法证明更多的几何定理}\label{ux7528ux5411ux91cfux6cd5ux8bc1ux660eux66f4ux591aux7684ux51e0ux4f55ux5b9aux7406}

\textbf{定理 15.34}（用向量法证明斯图尔特定理）：设 \(D\) 是 \(\triangle ABC\) 的边 \(BC\) 上一点，\(BD = m\)，\(DC = n\)，\(BC = a = m + n\)，则

\[AD^2 = \frac{mb^2 + nc^2}{m + n} - mn\]

\textbf{证明}：设 \(\overrightarrow{AB} = \vec{c}\)，\(\overrightarrow{AC} = \vec{b}\)。则 \(\overrightarrow{BC} = \vec{b} - \vec{c}\)，\(|\vec{b}| = b\)，\(|\vec{c}| = c\)。

\[\overrightarrow{AD} = \overrightarrow{AB} + \overrightarrow{BD} = \vec{c} + \frac{m}{a}(\vec{b} - \vec{c}) = \frac{n}{a}\vec{c} + \frac{m}{a}\vec{b}\]

\[AD^2 = |\overrightarrow{AD}|^2 = \frac{n^2}{a^2}|\vec{c}|^2 + 2 \cdot \frac{mn}{a^2}\vec{b} \cdot \vec{c} + \frac{m^2}{a^2}|\vec{b}|^2\]

\[= \frac{n^2 c^2 + m^2 b^2 + 2mn \vec{b} \cdot \vec{c}}{a^2}\]

由余弦定理，\(\vec{b} \cdot \vec{c} = bc\cos A = \frac{b^2 + c^2 - a^2}{2}\)。代入化简：

\[AD^2 = \frac{n^2 c^2 + m^2 b^2 + mn(b^2 + c^2 - a^2)}{a^2} = \frac{(n^2 + mn)c^2 + (m^2 + mn)b^2 - mna^2}{a^2}\]

\[= \frac{n(m+n)c^2 + m(m+n)b^2}{a^2} - \frac{mna^2}{a^2} = \frac{nc^2 + mb^2}{a} - mn = \frac{mb^2 + nc^2}{m+n} - mn\]

\(\blacksquare\)

\textbf{推论 15.7}（中线长公式）：当 \(m = n = \frac{a}{2}\) 时（\(D\) 为 \(BC\) 的中点），斯图尔特定理退化为中线长公式：

\[m_a^2 = \frac{2b^2 + 2c^2 - a^2}{4}\]

\textbf{证明}：代入 \(m = n = \frac{a}{2}\)：

\[AD^2 = \frac{\frac{a}{2} b^2 + \frac{a}{2} c^2}{a} - \frac{a^2}{4} = \frac{b^2 + c^2}{2} - \frac{a^2}{4} = \frac{2b^2 + 2c^2 - a^2}{4}\]

\(\blacksquare\)

\textbf{定理 15.35}（用向量法证明阿波罗尼奥斯定理）：三角形三条中线的平方和等于三边平方和的四分之三，即

\[m_a^2 + m_b^2 + m_c^2 = \frac{3}{4}(a^2 + b^2 + c^2)\]

\textbf{证明}：由中线长公式（推论 15.7），三条中线的平方分别为：

\[m_a^2 = \frac{2b^2 + 2c^2 - a^2}{4}, \quad m_b^2 = \frac{2a^2 + 2c^2 - b^2}{4}, \quad m_c^2 = \frac{2a^2 + 2b^2 - c^2}{4}\]

求和：

\[m_a^2 + m_b^2 + m_c^2 = \frac{(2b^2 + 2c^2 - a^2) + (2a^2 + 2c^2 - b^2) + (2a^2 + 2b^2 - c^2)}{4}\]

\[= \frac{3a^2 + 3b^2 + 3c^2}{4} = \frac{3}{4}(a^2 + b^2 + c^2)\]

\(\blacksquare\)

\textbf{定理 15.36}（用向量法证明四面体体积公式）：以 \(\vec{a}, \vec{b}, \vec{c}\) 为从同一顶点出发的三条棱的四面体的体积为

\[V = \frac{1}{6}|[\vec{a}\;\vec{b}\;\vec{c}]|\]

\textbf{证明}：以 \(\vec{a}, \vec{b}, \vec{c}\) 为棱的平行六面体的体积为 \(|[\vec{a}\;\vec{b}\;\vec{c}]|\)（定理 15.32）。四面体是平行六面体的六分之一（平行六面体可分解为六个等体积的四面体），因此 \(V = \frac{1}{6}|[\vec{a}\;\vec{b}\;\vec{c}]|\)。\(\blacksquare\)

\hrulefill

\textbf{本章展望}：本章建立了向量的完整理论体系，从平面向量的运算和坐标表示出发，推广到空间向量，引入了数量积和向量积这两个核心运算，并系统地将立体几何中的平行、垂直、角度和距离关系转化为向量运算。向量方法将几何直觉与代数计算完美结合，使得复杂的几何问题可以程序化地求解。在后续的解析几何中，向量将成为描述曲线和曲面的基本工具；在函数分析中，向量空间的概念将进一步推广到无穷维的情形。

\hrulefill

\hrulefill

\hrulefill

\hrulefill

\hrulefill

\hrulefill

\newpage

\chapter{02-几何/02-立体几何与向量/03-几何基础与公理体系.md}\label{ux51e0ux4f5502-ux7acbux4f53ux51e0ux4f55ux4e0eux5411ux91cf03-ux51e0ux4f55ux57faux7840ux4e0eux516cux7406ux4f53ux7cfb.md}

\section{第42章：几何基础与公理体系}\label{ux7b2c42ux7ae0ux51e0ux4f55ux57faux7840ux4e0eux516cux7406ux4f53ux7cfb}

几何基础研究 Euclid 公理体系的完备性和独立性，是数学公理化运动的先驱。本章系统阐述 Hilbert 公理体系的完整结构（五组二十条公理）、Pasch 公理和 Archimedes 公理的几何意义、非欧几何的公理差异，以及 Godel 不完备性定理对 Hilbert 纲领的深远影响。

\subsection{24.1 Euclid 公理体系的回顾与 Hilbert 的动机}\label{euclid-ux516cux7406ux4f53ux7cfbux7684ux56deux987eux4e0e-hilbert-ux7684ux52a8ux673a}

Euclid《几何原本》（约公元前 300 年）从 5 条公设和 5 条公理出发，推导出 465 个命题，是公理化方法的开山之作。然而，19 世纪的数学家逐渐发现 Euclid 的体系存在多处逻辑漏洞------许多证明暗中依赖了图形直观而未经公理化的假设。最典型的缺陷包括：

\begin{enumerate}
\def\labelenumi{\arabic{enumi}.}
\tightlist
\item
  \textbf{顺序关系的缺失}：Euclid 从未定义''介于''（betweenness）关系。命题 I.1（构造等边三角形）中，两个圆的交点被默认为存在，但这依赖于圆的''连续性''------点从一个圆内部到另一个圆内部必然经过交点------此论证在 Euclid 体系中无公理支撑。
\item
  \textbf{合同关系的模糊}：Euclid 用''叠合''（superposition）来证明 SAS 全等判定，但''移动图形''的操作缺乏公理化基础。
\item
  \textbf{连续性的缺失}：Euclid 没有处理无理数和极限的公理，而穷竭法（Eudoxus 方法）本质上是避开而非解决连续性问题。
\end{enumerate}

Moritz Pasch（1882）首次明确指出这些缺陷并提出修补方案。David Hilbert 在《几何基础》（\emph{Grundlagen der Geometrie}, 1899）中将这一工作推向极致------他建立了一套完全形式化的公理体系，其中''点''\,``直线''``平面''是无需定义的原始概念，其性质完全由公理规定（``隐定义''思想）。

\subsection{24.2 Hilbert 公理体系的完整列出}\label{hilbert-ux516cux7406ux4f53ux7cfbux7684ux5b8cux6574ux5217ux51fa}

Hilbert 的公理体系分为五组，共二十条。以下给出完整的公理陈述。

\textbf{定义 24.1}（Hilbert 公理体系的原始概念）：Hilbert 体系有三个原始概念------\textbf{点}、\textbf{直线}、\textbf{平面}，以及三个原始关系------\textbf{关联}（点在直线上、点在平面上）、\textbf{介于}（一点在另外两点之间）、\textbf{合同}（线段合同、角合同）。

\hrulefill

\textbf{第一组：关联公理（Incidence Axioms, I.1--I.8）}

这组公理规定了点、直线、平面之间的基本从属关系。

\begin{itemize}
\tightlist
\item
  \textbf{I.1}：过任意两个不同的点，存在唯一一条直线。
\item
  \textbf{I.2}：每条直线上至少有两个点。
\item
  \textbf{I.3}：存在至少三个不共线的点。
\item
  \textbf{I.4}：过任意三个不共线的点，存在唯一一个平面。
\item
  \textbf{I.5}：每个平面上至少有一个点。
\item
  \textbf{I.6}：若一条直线上的两个点在一个平面上，则该直线上所有点都在该平面上。
\item
  \textbf{I.7}：若两个平面有一个公共点，则它们还有另一个公共点。
\item
  \textbf{I.8}：存在至少四个不共面的点。
\end{itemize}

I.1--I.3 刻画了平面几何中的点线关联；I.4--I.8 将体系扩展到空间。关联公理是最基础的公理组------没有它们，``点在直线上''这一基本断言就失去了意义。

\hrulefill

\textbf{第二组：顺序公理（Order Axioms, II.1--II.4）}

Euclid 体系中缺失的核心概念，由 Pasch 首次明确引入。

\begin{itemize}
\tightlist
\item
  \textbf{II.1}：若点 \(B\) 介于点 \(A\) 和点 \(C\) 之间（记作 \(A * B * C\)），则 \(A, B, C\) 是直线上三个不同的点，且 \(B\) 也介于 \(C\) 和 \(A\) 之间。
\item
  \textbf{II.2}：对任意两个不同的点 \(A, B\)，在直线 \(AB\) 上存在点 \(C\) 使得 \(B\) 介于 \(A\) 和 \(C\) 之间。
\item
  \textbf{II.3}：共线的三个点中，至多有一个点介于另外两个点之间。
\item
  \textbf{II.4}（\textbf{Pasch 公理}）：设 \(A, B, C\) 是不共线的三个点，\(l\) 是平面 \(ABC\) 上不经过 \(A, B, C\) 中任何一点的直线。若 \(l\) 与线段 \(AB\) 相交，则 \(l\) 必与线段 \(AC\) 或线段 \(BC\) 之一相交。
\end{itemize}

\textbf{Pasch 公理的深入阐述}。Pasch 公理是顺序公理中最关键的一条，它首次将''直线分割平面''这一直观事实形式化。其几何含义是：若一条直线''进入''了一个三角形的内部（穿过了某条边），则它必须''穿出''三角形（穿过另一条边）。在 Euclid 的原始体系中，命题 I.16（三角形外角大于不相邻的内角）的证明默默地假设了中线和角平分线的交点位于三角形内部，这一假设本质上依赖了 Pasch 公理。Pasch 公理的深刻之处还在于：它等价于平面分离定理------平面上任意一条直线将平面分为两个（开的）半平面，使得连接不同半平面中两点的线段与该直线相交。

\hrulefill

\textbf{第三组：合同公理（Congruence Axioms, III.1--III.5）}

这组公理形式化了线段和角的''相等''关系（合同），替代了 Euclid 的''叠合''操作。

\begin{itemize}
\tightlist
\item
  \textbf{III.1}：设 \(A, B\) 是直线 \(l\) 上的两个点，\(A'\) 是直线 \(l'\) 上的点。则在 \(l'\) 上 \(A'\) 的指定一侧存在唯一的点 \(B'\) 使得线段 \(AB\) 与线段 \(A'B'\) 合同（记作 \(AB \equiv A'B'\)）。
\item
  \textbf{III.2}：若 \(AB \equiv A'B'\) 且 \(AB \equiv A''B''\)，则 \(A'B' \equiv A''B''\)。每条线段与自身合同：\(AB \equiv BA\)。
\item
  \textbf{III.3}：设 \(A * B * C\) 在直线 \(l\) 上，\(A' * B' * C'\) 在直线 \(l'\) 上。若 \(AB \equiv A'B'\) 且 \(BC \equiv B'C'\)，则 \(AC \equiv A'C'\)。
\item
  \textbf{III.4}：设 \(\angle (h, k)\) 是给定角（两条从同一点出发的射线），\(l'\) 是某直线及其指定的一侧，\(h'\) 是 \(l'\) 上从 \(O'\) 出发的射线。则在 \(l'\) 的指定一侧存在唯一的射线 \(k'\) 从 \(O'\) 出发使得 \(\angle (h, k) \equiv \angle (h', k')\)。每个角与自身合同。
\item
  \textbf{III.5}（SAS 全等判定）：若两个三角形满足两边及其夹角分别合同，则两个三角形的其余边和角也分别合同。
\end{itemize}

合同公理的前四条确保了线段的长度比较和角的度量可以自洽地定义，而 III.5（SAS）则是整个合同理论的枢纽------它将线段的合同与角的合同联系起来，使得三角形的全等判定得以形式化证明（其余全等判定 ASA、SSS 和 AAS 均可由 SAS 推导）。

\hrulefill

\textbf{第四组：平行公理（Parallel Axiom, IV）}

\begin{itemize}
\tightlist
\item
  \textbf{IV}（Playfair 形式的平行公理）：在平面上，过直线外一点至多存在一条直线与该直线不相交。
\end{itemize}

这是 Hilbert 体系中最简洁的公理------仅此一条。它与 Euclid 的原始第五公设等价（在结合公理、顺序公理和合同公理的前提下）。Hilbert 选择 Playfair 的表述正是因为其形式简洁，突出反映了平行概念的本质。

\hrulefill

\textbf{第五组：连续公理（Continuity Axioms, V.1--V.2）}

\begin{itemize}
\tightlist
\item
  \textbf{V.1}（\textbf{Archimedes 公理}）：设 \(AB\) 和 \(CD\) 是任意两条线段。则在射线 \(AB\) 上存在有限个点 \(A_1, A_2, \ldots, A_n\) 使得 \(A * A_1 * A_2 * \cdots\) 且线段 \(AA_1, A_1A_2, \ldots, A_{n-1}A_n\) 均与 \(CD\) 合同，且 \(B\) 介于 \(A\) 和 \(A_n\) 之间。
\end{itemize}

\textbf{Archimedes 公理的详细阐述}。这一公理的直观含义是：任意两条线段的长度是可以比较的------用较短的线段重复叠加，总可以超过任意长的线段。它排除了''无穷小''和''无穷大''线段的可能（即非 Archimedes 几何）。在代数层面，Archimedes 公理等价于断言：实数域的任何子域若含有一个非零元素，则必然含有任意大的自然数倍。没有 Archimedes 公理，几何学中将出现不可比较的量（incommensurable in the stronger sense），这将严重阻碍坐标系和解析方法的建立。

\begin{itemize}
\tightlist
\item
  \textbf{V.2}（\textbf{完备性公理 / 极大公理}）：在满足前述所有公理的点、直线、平面的体系中，不可能再添加新的点、直线或平面使得所有公理仍然成立。换言之，该体系是极大的------即不能扩张为更大的几何系统而不破坏公理。
\end{itemize}

完备性公理的等价表述为：系统与 \(\mathbb{R}^3\) 上的 Descartes 解析几何同构。结合 Archimedes 公理，完备性公理保证了 Hilbert 几何体系本质上等同于实数域上的三维 Euclid 空间。值得注意的是，若仅有 Archimedes 公理而无完备性公理，则几何模型可以是某个非 Archimedes 完备域上的几何（如可数欧氏空间），这在逻辑上与原体系一致但并不''同构''于标准 Euclid 几何。Hilbert 加入完备性公理正是为了锁定标准模型。

\textbf{定理 24.1}（Hilbert 公理体系的范畴性）：满足 Hilbert 五组全部二十条公理的任意两个模型必彼此同构。特别地，每个满足全部公理的几何体系同构于 \(\mathbb{R}^3\) 上的三维 Euclid 解析几何。

\textbf{证明}：由关联公理和合同公理可定义线段的''长度''和角的''度量''作为有序域中的元素。Archimedes 公理保证该有序域是 Archimedes 的，从而可嵌入 \(\mathbb{R}\)。完备性公理则断言该嵌入恰为满射（即该有序域与 \(\mathbb{R}\) 同构）。在此基础上构造坐标系：在平面上取两条正交直线为数轴，利用合同公理建立点的坐标表示------每个点对应 \(\mathbb{R}^2\) 中唯一的有序对。所有几何关系（关联、介于、合同）均可翻译为 \(\mathbb{R}^2\) 上的代数条件，故该体系与标准解析几何同构。\(\blacksquare\)

\subsection{24.3 非欧几何的公理差异}\label{ux975eux6b27ux51e0ux4f55ux7684ux516cux7406ux5deeux5f02}

Hilbert 公理体系中，平行公理（IV）的地位与其他公理有本质区别：它可以被独立地变更而不破坏体系的一致性（经由相对无矛盾性证明）。

\textbf{双曲几何的公理替换}：将 Hilbert 的平行公理 IV 替换为以下 \textbf{Lobachevsky 公理}：

\begin{quote}
在平面上，过直线外一点存在至少两条直线与该直线不相交。
\end{quote}

其余四组公理（关联、顺序、合同、连续性）保持不变。由此得到的公理体系定义了双曲几何（Lobachevsky-Bolyai 几何）。在双曲几何中，三角形的内角和严格小于 \(\pi\)（亏量与面积成正比），不存在相似三角形（AAA 全等判定成立），过直线外一点的平行线有不可数无穷多条。

\textbf{定理 24.2}（双曲几何的相对无矛盾性）：若 Euclid 几何是一致的，则双曲几何也是一致的。

\textbf{证明}：构造 Beltrami-Klein 模型。模型空间为 Euclid 平面上的单位开圆盘 \(D = \{ (x, y) : x^2 + y^2 < 1 \}\)。``点''对应 \(D\) 中的 Euclid 点。``直线''对应 \(D\) 中的 Euclid 开弦（两端在单位圆周上）。``介于''关系继承自 Euclid 介于关系。``合同''通过射影度规定义： \[d(P, Q) = \frac{1}{2} \left| \ln \frac{|PA| \cdot |QB|}{|PB| \cdot |QA|} \right|\] 其中 \(A, B\) 是弦 \(PQ\) 两端与单位圆的交点。可直接逐一验证：在此模型中，关联公理 I.1--I.3、顺序公理 II.1--II.4、合同公理 III.1--III.5 和连续公理 V.1--V.2 全部成立，但平行公理 IV 不成立------过直线外一点存在无穷多条不相交的''直线''。若双曲几何存在矛盾，则该矛盾在 Beltrami-Klein 模型中的翻译将是 Euclid 几何中的一个矛盾，与 Euclid 几何的一致性假设相悖。\(\blacksquare\)

\textbf{椭圆几何的公理替换}：将平行公理 IV 替换为 \textbf{Riemann 公理}：

\begin{quote}
在平面上，不存在两条不相交的直线（即任意两条直线必相交）。
\end{quote}

椭圆几何还需修改关联公理和顺序公理的部分内容，因为此时直线是闭合的（同胚于 \(S^1\)），无法定义''介于''关系在整条直线上的全局一致性。椭圆几何的标准模型是球面 \(S^2\) 上对径点等同（\(\mathbb{RP}^2\)），此时''直线''对应大圆。三角形内角和大于 \(\pi\)，其过量（excess）与面积成正比。

\textbf{定理 24.3}（三种几何的统一刻画）：Euclid 几何、双曲几何和椭圆几何分别对应于截面曲率 \(K = 0\)、\(K < 0\) 和 \(K > 0\) 的常曲率空间。这一统一视角由 Riemann（1854）和 Beltrami（1868）建立，使非欧几何从''逻辑可能性''升格为''几何实在性''。

\subsection{24.4 公理的独立性与 Hilbert 的元数学纲领}\label{ux516cux7406ux7684ux72ecux7acbux6027ux4e0e-hilbert-ux7684ux5143ux6570ux5b66ux7eb2ux9886}

Hilbert 不仅要求公理体系具有一致性，还要求每条公理（或每组公理）\textbf{独立}于其他公理------即不能从其余公理推导出来。独立性证明的标准方法是构造模型：要证明公理 \(A\) 独立于公理集 \(\Sigma\)，需构造一个满足 \(\Sigma\) 但不满足 \(A\) 的模型。

\textbf{定理 24.4}（平行公理的独立性）：平行公理 IV 独立于关联、顺序、合同和连续公理（即前四组公理不能推导出平行公理）。

\textbf{证明}：如定理 24.2 所示，Beltrami-Klein 模型满足 I、II、III、V 四组全部公理，但不满足 IV（双曲平行公理成立，Euclid 平行公理不成立）。因此 IV 不能由 I + II + III + V 推导出。若能从 I + II + III + V 推出 IV，则双曲几何将蕴含矛盾（因为双曲几何满足 I + II + III + V 但否定 IV）。\(\blacksquare\)

\textbf{定理 24.5}（Archimedes 公理的独立性）：Archimedes 公理 V.1 独立于其余公理（包括完备性公理 V.2）。

\textbf{证明}：构造非 Archimedes 模型：取有序域 \(\mathbb{R}(t)\)（实系数有理函数域），其中 \(t > 0\) 且 \(t < 1/n\) 对所有正整数 \(n\)（即 \(t\) 是正无穷小）。在此域上构造 Descartes 平面------点对应有序对 \((x, y)\)，\(x, y \in \mathbb{R}(t)\)，直线由线性方程 \(ax + by + c = 0\) 给出。所有关联、顺序、合同和平行公理在此模型中成立，但 Archimedes 公理不成立（线段 \((0,0)(t,0)\) 无论叠加多少次都无法超过线段 \((0,0)(1,0)\)）。此构造证明了 V.1 的独立性。\(\blacksquare\)

\subsection{24.5 Godel 不完备性定理对 Hilbert 纲领的影响}\label{godel-ux4e0dux5b8cux5907ux6027ux5b9aux7406ux5bf9-hilbert-ux7eb2ux9886ux7684ux5f71ux54cd}

Hilbert 纲领（1920s）的目标是用''有穷方法''（finitistic methods）证明全部数学的一致性。几何公理体系的一致性被 Hilbert 设想为这一宏大计划的第一步。然而，Godel 不完备性定理（1931）对这一纲领产生了深远的影响。

\textbf{定理 24.6}（Godel 不完备性定理对几何基础的影响）：初等 Euclid 几何（在 Tarski 的公理化下）是完备且可判定的（Tarski 1951）。但这一事实并不挽救 Hilbert 纲领------因为 Tarski 体系中的''完备''是指\textbf{语义完备}（任何一阶语句或其否定可证），而 Hilbert 所追求的是用有穷方法证明\textbf{一致性}。Godel 第二不完备性定理表明：任何包含初等算术的一致形式系统，不能在其自身内部证明自身的一致性。这意味着：即使初等几何是完备的，我们仍然不能用纯几何内部的有穷方法来证明几何的一致性------需要借助更强的外部系统（如实数理论或集合论）。

\textbf{剖析 Hilbert 几何体系的逻辑地位}。在二阶逻辑中，Hilbert 的五组公理是\textbf{范畴的}（定理 24.1）：所有模型同构于 \(\mathbb{R}^3\) 上的 Euclid 空间。但范畴性意味着公理体系对标准模型的刻画是精确的（没有非标准模型）。然而，当我们将 Hilbert 公理翻译为一阶逻辑时，Lowenheim-Skolem 定理确保了一阶理论必有非标准模型（如可数模型 \(\mathbb{Q}^3\) 的代数闭包），因而不完备。这种二阶/一阶之间的张力是理解 Godel 不完备性定理与几何基础关系的关键。几何之所以能在 Tarski 体系中被证明是完备的（在一阶语言中），是因为初等几何的''表达力''不够强------它无法定义自然数（几何中的长度和角度是有序域的元素，而非离散的自然数），因此 Godel 的自指构造（需要编码算术）在纯几何语言中无法实现。这正是''较弱公理体系确保了完备性''这一论断的精确含义。

\textbf{结论}：Hilbert 公理体系是公理化方法的巅峰之作。它将 Euclid 几何从直观依赖中解放出来，建立了严格的逻辑基础；其独立性证明则揭示了平行公理的逻辑地位，为接受非欧几何提供了方法论框架。然而 Godel 不完备性定理表明，Hilbert 所设想的''用有穷方法证明一切数学一致性''的目标在原则上不可实现------数学基础的完备性永远需要超越体系自身的外部资源。这一洞见构成了 20 世纪数学哲学的核心遗产。

\hrulefill

\begin{quote}
卷三B总结：本卷在卷三A平面几何的基础上，将几何理论推广到三维空间，并引入向量和公理体系深层探讨。第20章建立了空间中直线与平面关系的判定与性质、空间角与距离、以及简单几何体的体积与表面积理论；第21章从平面向量推广到空间向量，建立了数量积和向量积的运算体系，用向量方法统一处理立体几何中的平行、垂直、角度和距离问题；第23章探讨了几何基础与公理体系的完备性和独立性，包括非欧几里得几何的存在性和公理体系的模型构造。这三章构成了从立体几何直觉到代数方法再到元数学反思的完整逻辑链条。

\textbf{各章定理索引}：

\begin{itemize}
\tightlist
\item
  \textbf{Ch 20}（立体几何）：线面平行/面面平行的判定与性质、线面垂直/面面垂直的判定与性质、三垂线定理及其逆定理、异面直线角/线面角/二面角、各种距离公式、棱柱/棱锥/棱台/圆柱/圆锥/圆台/球的体积与表面积、祖暅原理、欧拉公式。
\item
  \textbf{Ch 21}（向量）：平面向量运算与坐标、数量积与向量积、向量法判定平行垂直、向量法计算角度与距离、混合积、向量法证明勾股定理/余弦定理/正弦定理/中位线定理/斯图尔特定理/阿波罗尼奥斯定理。
\item
  \textbf{第23章}（几何基础与公理体系）：Hilbert公理体系的完备性与独立性、非欧几里得几何的模型构造、Gödel不完备性定理的几何意涵。
\end{itemize}
\end{quote}

\hrulefill

\emph{卷三B：立体几何、向量与几何基础终。}

\emph{卷三B：立体几何、向量与几何基础终。}

\emph{卷三B：立体几何、向量与几何基础终。}

\emph{卷三：几何终。}

\newpage

\chapter{03-函数与微积分/01-函数与初等函数/01-坐标系与函数.md}\label{ux51fdux6570ux4e0eux5faeux79efux520601-ux51fdux6570ux4e0eux521dux7b49ux51fdux657001-ux5750ux6807ux7cfbux4e0eux51fdux6570.md}

\section{第43章：坐标系与函数概念}\label{ux7b2c43ux7ae0ux5750ux6807ux7cfbux4e0eux51fdux6570ux6982ux5ff5}

\textbf{前置知识}：本章依赖 Ch 3（集合论，特别是笛卡尔积、关系、函数的集合论定义）以及 Ch 19.6（锐角三角函数，为后续三角函数的坐标化做铺垫）。本章为后续所有函数章节（Ch 24--Ch 31）奠定基础。

\hrulefill

平面直角坐标系是解析几何的基石，也是函数理论的起点。笛卡尔（Rene Descartes）在1637年发表的《几何学》中首次系统地用坐标方法描述几何图形，这一创新将几何问题转化为代数问题，开创了''解析几何''这一数学分支。平面直角坐标系通过两条互相垂直的数轴（\(x\) 轴和 \(y\) 轴）将平面上的点与有序实数对 \((x, y)\) 一一对应起来，这一对应关系（双射）是解析几何的核心。本节从坐标系的建立出发，定义了平面上点的坐标，推导了两点间距离公式 \(d = \sqrt{(x_1 - x_2)^2 + (y_1 - y_2)^2}\)（这是勾股定理在坐标系中的直接体现）和中点坐标公式，并建立了线段的定比分点公式。这些公式是后续直线方程、曲线方程和函数图像研究的基础工具。

\textbf{定义 16.1}（平面直角坐标系）在平面上取两条互相垂直且有公共原点 \(O\) 的数轴，分别称为 \(x\) 轴（横轴）和 \(y\) 轴（纵轴）。\(x\) 轴取向右为正方向，\(y\) 轴取向上为正方向。这样建立的坐标系称为\textbf{平面直角坐标系}，记作 \(Oxy\)。

对于平面上的任意一点 \(P\)，过 \(P\) 分别作 \(x\) 轴和 \(y\) 轴的垂线，垂足对应的实数分别为 \(x_0\) 和 \(y_0\)，则有序实数对 \((x_0, y_0)\) 称为点 \(P\) 的\textbf{坐标}，记作 \(P(x_0, y_0)\)。

\textbf{定理 16.1}（坐标的唯一性与存在性）平面直角坐标系中，平面上的每个点有唯一的坐标，反之，每个有序实数对对应平面上唯一的点。即，平面直角坐标系建立了平面点集 \(\mathcal{P}\) 与笛卡尔积 \(\mathbb{R} \times \mathbb{R}\) 之间的一一对应（双射）。

\textbf{证明}：由 Ch 3 中笛卡尔积的定义以及垂线的唯一性（由欧氏几何公理保证），结论成立。具体地：

存在性：对于点 \(P\)，由垂线公理，过 \(P\) 作 \(x\) 轴的垂线有且仅有一条，垂足唯一确定实数 \(x_0\)；类似地确定 \(y_0\)。

唯一性：若点 \(P\) 有两组坐标 \((x_0, y_0)\) 和 \((x_1, y_1)\)，则过 \(P\) 向 \(x\) 轴作垂线，垂足既对应 \(x_0\) 又对应 \(x_1\)，由数轴上点与实数的一一对应得 \(x_0 = x_1\)，类似得 \(y_0 = y_1\)。

反过来，给定 \((x_0, y_0)\)，在 \(x\) 轴上取坐标为 \(x_0\) 的点 \(A\)，在 \(y\) 轴上取坐标为 \(y_0\) 的点 \(B\)，过 \(A\) 作 \(x\) 轴的垂线，过 \(B\) 作 \(y\) 轴的垂线，两垂线的交点即为所求点，由交点的唯一性得证。\(\blacksquare\)

\subsubsection{16.1.2 象限}\label{ux8c61ux9650}

平面直角坐标系将平面（除去坐标轴）分为四个\textbf{象限}：

\begin{itemize}
\tightlist
\item
  \textbf{第一象限}：\(\{(x, y) \mid x > 0, y > 0\}\)
\item
  \textbf{第二象限}：\(\{(x, y) \mid x < 0, y > 0\}\)
\item
  \textbf{第三象限}：\(\{(x, y) \mid x < 0, y < 0\}\)
\item
  \textbf{第四象限}：\(\{(x, y) \mid x > 0, y < 0\}\)
\end{itemize}

\subsubsection{16.1.3 两点间距离公式}\label{ux4e24ux70b9ux95f4ux8dddux79bbux516cux5f0f}

\textbf{定理 16.2}（两点间距离公式）设 \(P_1(x_1, y_1)\) 和 \(P_2(x_2, y_2)\) 是平面上的两点，则 \(P_1\) 与 \(P_2\) 之间的距离为：

\[d(P_1, P_2) = \sqrt{(x_2 - x_1)^2 + (y_2 - y_1)^2}\]

\textbf{证明}：过 \(P_1\) 作 \(x\) 轴的平行线，过 \(P_2\) 作 \(y\) 轴的平行线，两线交于点 \(Q(x_2, y_1)\)。

在直角三角形 \(P_1 Q P_2\) 中，由勾股定理（Ch 19.3）：

\[d(P_1, P_2)^2 = d(P_1, Q)^2 + d(Q, P_2)^2\]

由于 \(P_1(x_1, y_1)\) 和 \(Q(x_2, y_1)\) 在同一水平线上，\(d(P_1, Q) = |x_2 - x_1|\)。

由于 \(Q(x_2, y_1)\) 和 \(P_2(x_2, y_2)\) 在同一竖直线上，\(d(Q, P_2) = |y_2 - y_1|\)。

因此 \(d(P_1, P_2)^2 = (x_2 - x_1)^2 + (y_2 - y_1)^2\)，取正平方根即得。\(\blacksquare\)

\subsubsection{16.1.4 中点公式}\label{ux4e2dux70b9ux516cux5f0f}

\textbf{定理 16.3}（中点公式）设 \(P_1(x_1, y_1)\) 和 \(P_2(x_2, y_2)\) 是两点，则线段 \(P_1 P_2\) 的中点 \(M\) 的坐标为：

\[M = \left(\frac{x_1 + x_2}{2}, \frac{y_1 + y_2}{2}\right)\]

\textbf{证明}：设中点为 \(M(x_0, y_0)\)。由中点的定义，\(d(P_1, M) = d(M, P_2)\)，且 \(M\) 在线段 \(P_1 P_2\) 上。

由距离公式，\(x_0 - x_1 = x_2 - x_0\)（方向相同，距离相等），故 \(x_0 = \frac{x_1 + x_2}{2}\)。类似地，\(y_0 = \frac{y_1 + y_2}{2}\)。\(\blacksquare\)

\hrulefill

函数的严格定义是数学分析大厦的基石。在初等数学中，函数通常被描述为''一个量随另一个量变化的依赖关系''，但这一描述过于模糊，无法用于严格的数学推理。本节从集合论的角度给出了函数的严格定义：函数 \(f: A \to B\) 是笛卡尔积 \(A \times B\) 中满足''单值性''条件（每个 \(x \in A\) 对应唯一的 \(y \in B\)）的子集。这一定义将函数还原为一种特殊的二元关系，彻底消除了直觉上的模糊性。在此基础上，我们定义了函数的定义域、值域、图像等基本概念，并引入了函数的三种表示法（解析法、列表法、图像法）。函数的相等问题（两个函数何时是''同一个''函数）也在此获得了严格解答：当且仅当两者的定义域相同且对应规则一致时，两个函数相等。这一严密的定义框架为后续所有函数性质（单调性、奇偶性、周期性、有界性）的研究提供了逻辑基础。

在 Ch 3 中，我们已经学习了有序对、笛卡尔积和关系的概念。现在我们给出函数的严格定义。

\textbf{定义 16.2}（函数）设 \(A\) 和 \(B\) 是两个非空集合。一个从 \(A\) 到 \(B\) 的\textbf{函数} \(f\) 是笛卡尔积 \(A \times B\) 的一个子集，满足：

\[\forall x \in A, \exists! y \in B, (x, y) \in f\]

即，对于 \(A\) 中的每一个元素 \(x\)，在 \(B\) 中存在\textbf{唯一}的元素 \(y\) 使得 \((x, y) \in f\)。

我们记作 \(f: A \to B\)，读作''\(f\) 是从 \(A\) 到 \(B\) 的函数''。对于 \((x, y) \in f\)，我们也记作 \(y = f(x)\)。

\textbf{定义 16.3}（定义域）函数 \(f: A \to B\) 的\textbf{定义域}（domain）是集合 \(A\)，记作 \(\operatorname{dom}(f) = A\)。

\textbf{定义 16.4}（陪域）函数 \(f: A \to B\) 的\textbf{陪域}（codomain）是集合 \(B\)，记作 \(\operatorname{cod}(f) = B\)。

\textbf{定义 16.5}（值域）函数 \(f: A \to B\) 的\textbf{值域}（range）是 \(B\) 中所有被 \(f\) 映射到的元素组成的集合：

\[\operatorname{range}(f) = \{y \in B \mid \exists x \in A, f(x) = y\}\]

\textbf{注意}：值域是陪域的子集，即 \(\operatorname{range}(f) \subseteq B\)。

\subsubsection{16.2.2 函数的像与原像}\label{ux51fdux6570ux7684ux50cfux4e0eux539fux50cf}

\textbf{定义 16.6}（像）设 \(f: A \to B\)，对于 \(S \subseteq A\)，\(S\) 在 \(f\) 下的\textbf{像}定义为：

\[f(S) = \{f(x) \mid x \in S\}\]

特别地，\(f(A) = \operatorname{range}(f)\)。

\textbf{定义 16.7}（原像）设 \(f: A \to B\)，对于 \(T \subseteq B\)，\(T\) 在 \(f\) 下的\textbf{原像}（或\textbf{逆像}）定义为：

\[f^{-1}(T) = \{x \in A \mid f(x) \in T\}\]

\subsubsection{16.2.3 函数相等的条件}\label{ux51fdux6570ux76f8ux7b49ux7684ux6761ux4ef6}

\textbf{定义 16.8}（函数相等）两个函数 \(f: A \to B\) 和 \(g: C \to D\) \textbf{相等}，记作 \(f = g\)，当且仅当：

\begin{enumerate}
\def\labelenumi{\arabic{enumi}.}
\tightlist
\item
  \(A = C\)（定义域相同）
\item
  \(B = D\)（陪域相同）
\item
  \(\forall x \in A, f(x) = g(x)\)（对应法则相同）
\end{enumerate}

\textbf{定理 16.4} 设 \(f, g: A \to B\) 是两个函数。则 \(f = g\) 当且仅当 \(\forall x \in A, f(x) = g(x)\)。

\textbf{证明}：必要性由定义显然成立。

充分性：假设 \(\forall x \in A, f(x) = g(x)\)。则

\[f = \{(x, f(x)) \mid x \in A\} = \{(x, g(x)) \mid x \in A\} = g\]

因此 \(f = g\)。\(\blacksquare\)

\hrulefill

函数的基本性质------单调性、奇偶性、周期性和有界性------是刻画函数行为特征的四大核心工具。单调性描述了函数值随自变量增减而变化的趋势：单调递增函数保持了自变量的大小关系（\(x_1 < x_2 \Rightarrow f(x_1) \leq f(x_2)\)），这使得单调函数在求解方程和不等式时尤为重要，因为单调函数对应的一一映射保证了反函数的存在性。奇偶性描述了函数图像关于原点和 \(y\) 轴的对称性：奇函数满足 \(f(-x) = -f(x)\)（图像关于原点中心对称），偶函数满足 \(f(-x) = f(x)\)（图像关于 \(y\) 轴对称）。周期性描述了函数的重复规律：存在正数 \(T\) 使得 \(f(x + T) = f(x)\) 对所有 \(x\) 成立，三角函数是最典型的周期函数。有界性刻画了函数值的范围：存在 \(M\) 使得 \(|f(x)| \leq M\) 对所有定义域中的 \(x\) 成立。这四种性质相互独立但又相互关联------例如，有界单调函数必存在极限，周期函数若有极限则必为常数------这些联系在后续的极限理论中将被深入揭示。

\textbf{定义 16.9}（单调递增）设 \(f: A \to \mathbb{R}\)，其中 \(A \subseteq \mathbb{R}\)。若

\[\forall x_1, x_2 \in A, x_1 < x_2 \Rightarrow f(x_1) \leq f(x_2)\]

则称 \(f\) 在 \(A\) 上是\textbf{单调递增}的（或\textbf{递增}的）。

\textbf{定义 16.10}（严格单调递增）若

\[\forall x_1, x_2 \in A, x_1 < x_2 \Rightarrow f(x_1) < f(x_2)\]

则称 \(f\) 在 \(A\) 上是\textbf{严格单调递增}的（或\textbf{严格递增}的）。

类似地可以定义\textbf{单调递减}和\textbf{严格单调递减}。

\textbf{定理 16.5}（单调函数的复合）设 \(f: A \to B\)，\(g: B \to \mathbb{R}\)。

\begin{enumerate}
\def\labelenumi{\arabic{enumi}.}
\tightlist
\item
  若 \(f\) 和 \(g\) 都是单调递增的，则 \(g \circ f\) 也是单调递增的。
\item
  若 \(f\) 是单调递增的，\(g\) 是单调递减的，则 \(g \circ f\) 是单调递减的。
\item
  若 \(f\) 是单调递减的，\(g\) 是单调递增的，则 \(g \circ f\) 是单调递减的。
\item
  若 \(f\) 和 \(g\) 都是单调递减的，则 \(g \circ f\) 是单调递增的。
\end{enumerate}

\textbf{证明}：我们证明结论 1，其余类似。

设 \(x_1, x_2 \in A\) 且 \(x_1 < x_2\)。因为 \(f\) 是单调递增的，所以 \(f(x_1) \leq f(x_2)\)。

因为 \(g\) 是单调递增的，所以 \(g(f(x_1)) \leq g(f(x_2))\)。

即 \((g \circ f)(x_1) \leq (g \circ f)(x_2)\)。

因此 \(g \circ f\) 是单调递增的。\(\blacksquare\)

\subsubsection{16.3.2 奇偶性}\label{ux5947ux5076ux6027}

\textbf{定义 16.11}（奇函数）设 \(f: A \to \mathbb{R}\)，其中 \(A\) 关于原点对称（即 \(x \in A \Leftrightarrow -x \in A\)）。若

\[\forall x \in A, f(-x) = -f(x)\]

则称 \(f\) 为\textbf{奇函数}。

\textbf{定义 16.12}（偶函数）设 \(f: A \to \mathbb{R}\)，其中 \(A\) 关于原点对称。若

\[\forall x \in A, f(-x) = f(x)\]

则称 \(f\) 为\textbf{偶函数}。

\textbf{定理 16.6}（奇偶性的判定）设 \(f: A \to \mathbb{R}\)，其中 \(A\) 关于原点对称。

\begin{enumerate}
\def\labelenumi{\arabic{enumi}.}
\tightlist
\item
  \(f\) 是奇函数当且仅当其图像关于原点对称。
\item
  \(f\) 是偶函数当且仅当其图像关于 \(y\) 轴对称。
\end{enumerate}

\textbf{证明}：

\begin{enumerate}
\def\labelenumi{(\arabic{enumi})}
\tightlist
\item
  设 \(f\) 是奇函数，\((x, f(x))\) 是图像上的一点。则 \((-x, f(-x)) = (-x, -f(x))\) 也在图像上。而 \((-x, -f(x))\) 正是 \((x, f(x))\) 关于原点的对称点，因此图像关于原点对称。
\end{enumerate}

反之，设图像关于原点对称，则对于图像上任意一点 \((x, f(x))\)，其关于原点的对称点 \((-x, -f(x))\) 也在图像上，即 \(f(-x) = -f(x)\)，所以 \(f\) 是奇函数。

\begin{enumerate}
\def\labelenumi{(\arabic{enumi})}
\setcounter{enumi}{1}
\tightlist
\item
  的证明类似。\(\blacksquare\)
\end{enumerate}

\textbf{定理 16.7}（奇偶性的代数运算）设 \(f\) 和 \(g\) 都是定义在关于原点对称的集合上的函数。

\begin{enumerate}
\def\labelenumi{\arabic{enumi}.}
\tightlist
\item
  若 \(f\) 和 \(g\) 都是奇函数，则 \(f + g\) 是奇函数，\(f \cdot g\) 是偶函数。
\item
  若 \(f\) 和 \(g\) 都是偶函数，则 \(f + g\) 和 \(f \cdot g\) 都是偶函数。
\item
  若 \(f\) 是奇函数，\(g\) 是偶函数，则 \(f \cdot g\) 是奇函数。
\end{enumerate}

\textbf{证明}：

对于结论 1 的加法部分：

\[(f + g)(-x) = f(-x) + g(-x) = (-f(x)) + (-g(x)) = -(f(x) + g(x)) = -(f + g)(x)\]

所以 \(f + g\) 是奇函数。

对于结论 1 的乘法部分：

\[(f \cdot g)(-x) = f(-x) \cdot g(-x) = (-f(x)) \cdot (-g(x)) = f(x) \cdot g(x) = (f \cdot g)(x)\]

所以 \(f \cdot g\) 是偶函数。

对于结论 2：

\[(f + g)(-x) = f(-x) + g(-x) = f(x) + g(x) = (f + g)(x)\]

\[(f \cdot g)(-x) = f(-x) \cdot g(-x) = f(x) \cdot g(x) = (f \cdot g)(x)\]

所以 \(f + g\) 和 \(f \cdot g\) 都是偶函数。

对于结论 3：

\[(f \cdot g)(-x) = f(-x) \cdot g(-x) = (-f(x)) \cdot g(x) = -(f(x) \cdot g(x)) = -(f \cdot g)(x)\]

所以 \(f \cdot g\) 是奇函数。\(\blacksquare\)

\subsubsection{16.3.3 周期性}\label{ux5468ux671fux6027}

\textbf{定义 16.13}（周期函数）设 \(f: \mathbb{R} \to \mathbb{R}\)。若存在正数 \(T\)，使得

\[\forall x \in \mathbb{R}, f(x + T) = f(x)\]

则称 \(f\) 是\textbf{周期函数}，\(T\) 是 \(f\) 的一个\textbf{周期}。

\textbf{定义 16.14}（最小正周期）若 \(T_0\) 是 \(f\) 的一个正周期，且对于 \(f\) 的任意正周期 \(T\)，都有 \(T_0 \leq T\)，则称 \(T_0\) 为 \(f\) 的\textbf{最小正周期}。

\textbf{定理 16.8} 设 \(f\) 是以 \(T\) 为周期的周期函数，则对于任意正整数 \(n\)，\(nT\) 也是 \(f\) 的周期。

\textbf{证明}：用数学归纳法（Ch 2.6）。

当 \(n = 1\) 时，\(T\) 是 \(f\) 的周期，结论成立。

假设 \(n = k\) 时结论成立，即 \(kT\) 是 \(f\) 的周期。则对于任意 \(x\)：

\[f(x + (k+1)T) = f((x + kT) + T) = f(x + kT) = f(x)\]

所以 \((k+1)T\) 也是 \(f\) 的周期。

由数学归纳法，结论成立。\(\blacksquare\)

\subsubsection{16.3.4 有界性}\label{ux6709ux754cux6027}

\textbf{定义 16.15}（有界函数）设 \(f: A \to \mathbb{R}\)。若存在 \(M > 0\)，使得

\[\forall x \in A, |f(x)| \leq M\]

则称 \(f\) 在 \(A\) 上是\textbf{有界}的。

\textbf{定义 16.16}（上有界、下有界）设 \(f: A \to \mathbb{R}\)。

\begin{itemize}
\tightlist
\item
  若存在 \(M \in \mathbb{R}\)，使得 \(\forall x \in A, f(x) \leq M\)，则称 \(f\) 在 \(A\) 上是\textbf{上有界}的，\(M\) 是 \(f\) 的一个\textbf{上界}。
\item
  若存在 \(m \in \mathbb{R}\)，使得 \(\forall x \in A, f(x) \geq m\)，则称 \(f\) 在 \(A\) 上是\textbf{下有界}的，\(m\) 是 \(f\) 的一个\textbf{下界}。
\end{itemize}

\textbf{定理 16.9} \(f\) 在 \(A\) 上有界当且仅当 \(f\) 在 \(A\) 上既有上界又有下界。

\textbf{证明}：

充分性：设 \(f\) 有上界 \(M\) 和下界 \(m\)，则 \(|f(x)| \leq \max\{|M|, |m|\}\)，所以 \(f\) 有界。

必要性：设 \(f\) 有界，即存在 \(M' > 0\) 使得 \(|f(x)| \leq M'\)。则 \(M'\) 是上界，\(-M'\) 是下界。\(\blacksquare\)

\hrulefill

反函数与复合函数是函数理论中的两个核心构造工具。反函数回答的问题是''给定输出值 \(y\)，能否唯一确定输入值 \(x\)？``------这要求原函数是双射（既是单射又是满射）。反函数的存在性定理指出：严格单调的函数必有反函数，且反函数保持相同的单调性。反函数的图像与原函数的图像关于直线 \(y = x\) 对称，这一几何事实在后续的指数函数与对数函数、三角函数与反三角函数的关系中反复出现。复合函数则将两个函数''串联''起来：\((g \circ f)(x) = g(f(x))\)，先经过 \(f\) 再经过 \(g\)。复合函数的概念是理解函数分解和链式法则（导数计算中最核心的工具）的基础。本节还引入了单射、满射和双射这三个集合论概念，它们为反函数的存在性提供了精确的判定条件，是函数理论从''直觉''走向''严格''的关键一步。

\textbf{定义 16.17}（单射）设 \(f: A \to B\)。若

\[\forall x_1, x_2 \in A, f(x_1) = f(x_2) \Rightarrow x_1 = x_2\]

则称 \(f\) 是\textbf{单射}（或\textbf{一对一}的）。

\textbf{定义 16.18}（满射）设 \(f: A \to B\)。若

\[\forall y \in B, \exists x \in A, f(x) = y\]

即 \(\operatorname{range}(f) = B\)，则称 \(f\) 是\textbf{满射}（或\textbf{映上}的）。

\textbf{定义 16.19}（双射）若 \(f\) 既是单射又是满射，则称 \(f\) 是\textbf{双射}（或\textbf{一一对应}）。

\subsubsection{16.4.2 反函数}\label{ux53cdux51fdux6570-1}

\textbf{定理 16.10}（反函数存在定理）设 \(f: A \to B\) 是一个函数。\(f\) 存在反函数当且仅当 \(f\) 是双射。

\textbf{证明}：

充分性：设 \(f\) 是双射。定义 \(g: B \to A\) 如下：对于任意 \(y \in B\)，由于 \(f\) 是满射，存在 \(x \in A\) 使得 \(f(x) = y\)；又由于 \(f\) 是单射，这样的 \(x\) 是唯一的。定义 \(g(y) = x\)，则 \(g\) 是良定义的函数。

对于任意 \(x \in A\)：\((g \circ f)(x) = g(f(x)) = x\)

对于任意 \(y \in B\)：\((f \circ g)(y) = f(g(y)) = y\)

因此 \(g\) 是 \(f\) 的反函数。

必要性：设 \(f\) 有反函数 \(g: B \to A\)。

单射性：若 \(f(x_1) = f(x_2)\)，则 \(x_1 = g(f(x_1)) = g(f(x_2)) = x_2\)。

满射性：对于任意 \(y \in B\)，令 \(x = g(y)\)，则 \(f(x) = f(g(y)) = y\)。\(\blacksquare\)

\textbf{定理 16.11}（反函数的唯一性）若 \(f\) 有反函数，则反函数是唯一的。

\textbf{证明}：设 \(g_1\) 和 \(g_2\) 都是 \(f\) 的反函数。则对于任意 \(y \in B\)：

\[g_1(y) = g_1(f(g_2(y))) = g_2(y)\]

所以 \(g_1 = g_2\)。\(\blacksquare\)

\textbf{定理 16.12}（反函数的性质）设 \(f: A \to B\) 是双射，\(g: B \to A\) 是其反函数。

\begin{enumerate}
\def\labelenumi{\arabic{enumi}.}
\tightlist
\item
  \(g\) 也是双射，且 \(f\) 是 \(g\) 的反函数。
\item
  \(f\) 和 \(g\) 的单调性相同。
\end{enumerate}

\textbf{证明}：

\begin{enumerate}
\def\labelenumi{(\arabic{enumi})}
\item
  由定义，\(g \circ f = \operatorname{id}_{A}\)，\(f \circ g = \operatorname{id}_{B}\)。因此 \(f\) 和 \(g\) 互为反函数。由此易证 \(g\) 是双射。
\item
  设 \(f\) 是严格递增的。对于 \(y_1, y_2 \in B\) 且 \(y_1 < y_2\)，设 \(x_1 = g(y_1)\)，\(x_2 = g(y_2)\)。则 \(f(x_1) = y_1 < y_2 = f(x_2)\)。
\end{enumerate}

若 \(x_1 \geq x_2\)，则由 \(f\) 严格递增得 \(f(x_1) \geq f(x_2)\)，矛盾。所以 \(x_1 < x_2\)，即 \(g(y_1) < g(y_2)\)，\(g\) 严格递增。

对于 \(f\) 严格递减的情况类似证明。\(\blacksquare\)

\textbf{定理 16.13}（反函数的图像）设 \(f: A \to B\) 是双射，\(g: B \to A\) 是其反函数。则 \(f\) 的图像和 \(g\) 的图像关于直线 \(y = x\) 对称。

\textbf{证明}：点 \((a, b)\) 在 \(f\) 的图像上当且仅当 \(b = f(a)\)，当且仅当 \(a = g(b)\)，当且仅当 \((b, a)\) 在 \(g\) 的图像上。

而 \((a, b)\) 和 \((b, a)\) 关于直线 \(y = x\) 对称，所以 \(f\) 和 \(g\) 的图像关于 \(y = x\) 对称。\(\blacksquare\)

\subsubsection{16.4.3 复合函数}\label{ux590dux5408ux51fdux6570-1}

\textbf{定义 16.20}（复合函数）设 \(f: A \to B\)，\(g: B \to C\)。定义 \(g\) 和 \(f\) 的\textbf{复合}为函数 \(g \circ f: A \to C\)：

\[(g \circ f)(x) = g(f(x)), \quad \forall x \in A\]

\textbf{注意}：复合的顺序是从右到左，即先作用 \(f\) 再作用 \(g\)。

\textbf{定理 16.14}（复合函数的定义域）设 \(f\) 的定义域为 \(A\)，\(g\) 的定义域为 \(D(g)\)。则 \(g \circ f\) 有定义当且仅当 \(\operatorname{range}(f) \subseteq D(g)\)。当 \(g \circ f\) 有定义时，

\[D(g \circ f) = \{x \in A \mid f(x) \in D(g)\}\]

\textbf{证明}：由复合函数的定义直接可得。\(\blacksquare\)

\textbf{定理 16.15}（复合函数的单调性）设 \(f: A \to B\)，\(g: B \to \mathbb{R}\)。复合函数 \(g \circ f\) 的单调性如下：

\begin{longtable}[]{@{}ccc@{}}
\toprule
\(f\) 的单调性 & \(g\) 的单调性 & \(g \circ f\) 的单调性 \\
\midrule
\endhead
\bottomrule
\endlastfoot
递增 & 递增 & 递增 \\
递增 & 递减 & 递减 \\
递减 & 递增 & 递减 \\
递减 & 递减 & 递增 \\
\end{longtable}

\textbf{证明}：此即定理 16.5 的结论。口诀：同增异减。\(\blacksquare\)

\hrulefill

函数图像变换是研究函数图像的有力工具。通过平移、伸缩、对称和翻折等基本变换，我们可以从已知函数的图像构造出大量相关函数的图像，而不必逐点描画。平移变换 \(y = f(x - a) + b\) 将原图像沿 \(x\) 轴方向平移 \(a\)、沿 \(y\) 轴方向平移 \(b\)；伸缩变换 \(y = kf(mx)\) 将原图像在水平和垂直方向上进行拉伸或压缩；对称变换 \(y = -f(x)\)、\(y = f(-x)\) 分别产生关于 \(x\) 轴和 \(y\) 轴的对称图像；翻折变换 \(y = |f(x)|\) 将 \(x\) 轴下方的部分翻折到上方。这些变换共同构成了函数图像的''操作工具箱''，其理论基础是坐标变换的群作用------每种变换都对应平面上的一个仿射变换。熟练掌握这些变换有助于在后续章节中快速理解各种初等函数（幂函数、指数函数、对数函数、三角函数）的图像特征及其相互关系。

\textbf{定理 16.16}（平移变换）设 \(y = f(x)\) 的图像为 \(C\)。

\begin{enumerate}
\def\labelenumi{\arabic{enumi}.}
\tightlist
\item
  \(y = f(x) + k\) 的图像可由 \(C\) 沿 \(y\) 轴向上平移 \(k\) 个单位得到（\(k > 0\)）。
\item
  \(y = f(x) - k\) 的图像可由 \(C\) 沿 \(y\) 轴向下平移 \(k\) 个单位得到（\(k > 0\)）。
\item
  \(y = f(x - h)\) 的图像可由 \(C\) 沿 \(x\) 轴向右平移 \(h\) 个单位得到（\(h > 0\)）。
\item
  \(y = f(x + h)\) 的图像可由 \(C\) 沿 \(x\) 轴向左平移 \(h\) 个单位得到（\(h > 0\)）。
\end{enumerate}

\textbf{证明}：以 (3) 为例。设点 \((a, b)\) 在 \(C\) 上，则 \(b = f(a)\)。令 \(x = a + h\)，则 \(f(x - h) = f(a) = b\)，所以 \((a + h, b)\) 在 \(y = f(x - h)\) 的图像上。而 \((a + h, b)\) 正是 \((a, b)\) 向右平移 \(h\) 个单位后的点。\(\blacksquare\)

\subsubsection{16.5.2 伸缩变换}\label{ux4f38ux7f29ux53d8ux6362}

\textbf{定理 16.17}（伸缩变换）设 \(y = f(x)\) 的图像为 \(C\)。

\begin{enumerate}
\def\labelenumi{\arabic{enumi}.}
\tightlist
\item
  \(y = kf(x)\)（\(k > 0\)）的图像可由 \(C\) 上每点的纵坐标乘以 \(k\) 得到（纵向伸缩）。
\item
  \(y = f(kx)\)（\(k > 0\)）的图像可由 \(C\) 上每点的横坐标除以 \(k\) 得到（横向伸缩）。
\end{enumerate}

\textbf{证明}：以 (1) 为例。设点 \((a, b)\) 在 \(C\) 上，则 \(b = f(a)\)。那么 \(kf(a) = kb\)，所以 \((a, kb)\) 在 \(y = kf(x)\) 的图像上。\(\blacksquare\)

\subsubsection{16.5.3 对称变换}\label{ux5bf9ux79f0ux53d8ux6362}

\textbf{定理 16.18}（对称变换）设 \(y = f(x)\) 的图像为 \(C\)。

\begin{enumerate}
\def\labelenumi{\arabic{enumi}.}
\tightlist
\item
  \(y = -f(x)\) 的图像与 \(C\) 关于 \(x\) 轴对称。
\item
  \(y = f(-x)\) 的图像与 \(C\) 关于 \(y\) 轴对称。
\item
  \(y = -f(-x)\) 的图像与 \(C\) 关于原点对称。
\end{enumerate}

\textbf{证明}：

\begin{enumerate}
\def\labelenumi{(\arabic{enumi})}
\item
  点 \((a, b)\) 在 \(C\) 上当且仅当 \(b = f(a)\)，当且仅当 \(-b = -f(a)\)，当且仅当 \((a, -b)\) 在 \(y = -f(x)\) 的图像上。而 \((a, -b)\) 是 \((a, b)\) 关于 \(x\) 轴的对称点。
\item
  点 \((a, b)\) 在 \(C\) 上当且仅当 \(b = f(a)\)，当且仅当 \(b = f(-(-a))\)，当且仅当 \((-a, b)\) 在 \(y = f(-x)\) 的图像上。而 \((-a, b)\) 是 \((a, b)\) 关于 \(y\) 轴的对称点。
\item
  点 \((a, b)\) 在 \(C\) 上当且仅当 \(b = f(a)\)，当且仅当 \(-b = -f(a) = -f(-(-a))\)，当且仅当 \((-a, -b)\) 在 \(y = -f(-x)\) 的图像上。而 \((-a, -b)\) 是 \((a, b)\) 关于原点的对称点。\(\blacksquare\)
\end{enumerate}

\hrulefill

\subsection{16.6 参数方程与极坐标简介}\label{ux53c2ux6570ux65b9ux7a0bux4e0eux6781ux5750ux6807ux7b80ux4ecb}

参数方程和极坐标是描述曲线的两种重要替代方式，它们各自在某些场景下比直角坐标方程更为简洁和有力。参数方程通过引入第三个变量（参数 \(t\)）将 \(x\) 和 \(y\) 分别表示为 \(t\) 的函数，从而将曲线视为''运动轨迹''------这一观点在物理学（如抛体运动、圆周运动）中尤为自然。极坐标则以极点和极轴为参考，用极径 \(r\) 和极角 \(\theta\) 来定位点，特别适合处理具有旋转对称性的曲线。本节简要介绍了参数方程和极坐标的基本概念，以及它们与直角坐标之间的互化关系。详细的参数方程与极坐标理论将在 Ch 31 中系统展开，届时将涵盖圆锥曲线的统一极坐标方程 \(r = \frac{ep}{1 - e\cos\theta}\) 和极坐标下的面积计算等深入内容。

\textbf{定义 16.21}（参数方程）设 \(x = \varphi(t)\)，\(y = \psi(t)\)（\(t \in T\)，\(T\) 为某区间），则此方程组称为曲线的\textbf{参数方程}，\(t\) 称为\textbf{参数}。

\textbf{定理 16.19}（消参）若 \(\varphi\) 有反函数 \(t = \varphi^{-1}(x)\)，则 \(y = \psi(\varphi^{-1}(x))\)。

\textbf{证明}：由 \(x = \varphi(t)\) 得 \(t = \varphi^{-1}(x)\)，代入 \(y = \psi(t)\) 即得。\(\blacksquare\)

\textbf{常见参数方程}：

\begin{itemize}
\tightlist
\item
  圆：\(x = r\cos t\)，\(y = r\sin t\)（\(t \in [0, 2\pi)\)）
\item
  椭圆：\(x = a\cos t\)，\(y = b\sin t\)（\(t \in [0, 2\pi)\)）
\item
  直线：\(x = x_0 + at\)，\(y = y_0 + bt\)（\(t \in \mathbb{R}\)）
\end{itemize}

\subsubsection{16.6.2 极坐标}\label{ux6781ux5750ux6807}

\textbf{定义 16.22}（极坐标系）在平面上取一点 \(O\)（称为\textbf{极点}）和一条射线 \(Ox\)（称为\textbf{极轴}）。对于平面上的点 \(P\)，设 \(|OP| = r \geq 0\)（\textbf{极径}），\(\angle xOP = \theta\)（\textbf{极角}），则有序对 \((r, \theta)\) 称为 \(P\) 的\textbf{极坐标}。

\textbf{定理 16.20}（极坐标与直角坐标的转换）

\[x = r\cos\theta, \quad y = r\sin\theta\]

\[r = \sqrt{x^2 + y^2}, \quad \tan\theta = \frac{y}{x}\]

\textbf{证明}：在以极点为原点、极轴为 \(x\) 轴正方向的直角坐标系中，由三角函数的定义即得。\(\blacksquare\)

\textbf{注}：极坐标表示不唯一，\((r, \theta)\) 和 \((r, \theta + 2k\pi)\)（\(k \in \mathbb{Z}\)）表示同一点。

\hrulefill

函数的复合分解与函数方程是函数理论中两个具有方法论意义的主题。函数的分解问题是复合的逆问题：给定一个复杂函数 \(h\)，能否将其表示为两个（或多个）较简单函数的复合？这一问题的实质是寻找函数的''构造层级''，与数论中质因数分解的思路一脉相承。函数方程则是一类以函数为未知量的方程，如 \(f(x + y) = f(x) + f(y)\)（柯西方程）、\(f(xy) = f(x) + f(y)\)（对数函数方程）等。函数方程的求解不依赖于代数运算，而是利用函数的性质（连续性、单调性、可导性等）来推断函数的具体形式。柯西方程是函数方程中最经典的例子------在连续性假设下，其唯一连续解为 \(f(x) = cx\)，这一结果不仅在理论上深刻（将线性函数刻画为加法群同态），在物理学中也具有基础意义（如叠加原理）。函数方程的研究方法融合了代数、分析、几何等多方面的思想，是培养数学综合素养的重要课题。

\textbf{问题}：给定一个复杂函数 \(h(x)\)，能否将其表示为两个较简单函数的复合 \(h = g \circ f\)？

\textbf{定理 16.21} 设 \(h: A \to C\) 是一个函数。若存在集合 \(B\) 和函数 \(f: A \to B\)，\(g: B \to C\)，使得 \(h = g \circ f\)，则称 \(h\) 可以\textbf{分解}为 \(g\) 和 \(f\) 的复合。

\textbf{证明}：分解总是可能的。取 \(B = \operatorname{range}(h)\)，\(f = h\)（视为 \(A \to B\) 的函数），\(g = \operatorname{id}_{B}\)（恒等映射），则 \(h = g \circ f\)。\(\blacksquare\)

\textbf{注}：函数分解在数学和工程中有重要应用。例如，在信号处理中，复杂信号可以分解为简单信号的复合；在微积分中，链式法则正是基于函数复合结构的求导方法。

\subsubsection{16.7.2 函数方程简介}\label{ux51fdux6570ux65b9ux7a0bux7b80ux4ecb}

\textbf{定义 16.23}（函数方程）含有未知函数的方程称为\textbf{函数方程}。求解函数方程是指找出满足方程的所有函数。

\textbf{典型例子}：

\begin{enumerate}
\def\labelenumi{(\arabic{enumi})}
\tightlist
\item
  \textbf{柯西函数方程}：\(f(x + y) = f(x) + f(y)\)。
\end{enumerate}

\textbf{定理 16.22} 若 \(f: \mathbb{R} \to \mathbb{R}\) 满足柯西方程 \(f(x+y) = f(x) + f(y)\)，且 \(f\) 连续（或单调，或在某点有界），则 \(f(x) = kx\)（\(k\) 为常数）。

\textbf{证明}：

第一步：\(f(0) = f(0 + 0) = f(0) + f(0)\)，故 \(f(0) = 0\)。

第二步：\(f(n) = nf(1)\)（对正整数 \(n\)，由归纳法）。

第三步：\(f(-n) + f(n) = f(0) = 0\)，故 \(f(-n) = -nf(1)\)。因此对所有整数 \(m\)，\(f(m) = mf(1)\)。

第四步：\(f(1) = f\left(\frac{1}{n} + \cdots + \frac{1}{n}\right) = nf\left(\frac{1}{n}\right)\)，故 \(f\left(\frac{1}{n}\right) = \frac{f(1)}{n}\)。

第五步：对有理数 \(\frac{m}{n}\)，\(f\left(\frac{m}{n}\right) = mf\left(\frac{1}{n}\right) = \frac{m}{n}f(1)\)。

第六步：由 \(f\) 的连续性，对无理数 \(x\)，取有理数列 \(r_n \to x\)，则 \(f(x) = \lim f(r_n) = \lim r_n f(1) = xf(1)\)。

设 \(k = f(1)\)，则 \(f(x) = kx\)。\(\blacksquare\)

\begin{enumerate}
\def\labelenumi{(\arabic{enumi})}
\setcounter{enumi}{1}
\tightlist
\item
  \textbf{指数型函数方程}：\(f(x + y) = f(x)f(y)\)，\(f\) 连续，\(f \not\equiv 0\)。解为 \(f(x) = a^x\)（\(a > 0\)）。
\end{enumerate}

\hrulefill

\textbf{本章展望}：本章建立了坐标系和函数的严格理论基础。在下一章（Ch 24）中，我们将利用函数的概念，系统研究一次函数、反比例函数、二次函数和幂函数等基本初等函数的具体性质。函数的单调性、奇偶性等一般性质将在这些具体函数上得到体现和应用。

\hrulefill

\hrulefill

\hrulefill

\hrulefill

\hrulefill

\newpage

\chapter{03-函数与微积分/01-函数与初等函数/02-基本初等函数.md}\label{ux51fdux6570ux4e0eux5faeux79efux520601-ux51fdux6570ux4e0eux521dux7b49ux51fdux657002-ux57faux672cux521dux7b49ux51fdux6570.md}

\section{第44章：基本初等函数}\label{ux7b2c44ux7ae0ux57faux672cux521dux7b49ux51fdux6570}

\begin{quote}
\textbf{注（阅读提示）}：本章中''连续''一词在直观意义上使用------图像不中断、不断开的函数。形式化的 \(\epsilon\)-\(\delta\) 定义在第29章中给出，读者可在学完第29章后回顾本章中关于连续性的讨论。初等函数（多项式、有理函数、指数函数、对数函数、三角函数）在其定义域上连续的严格证明同样依赖极限的 \(\epsilon\)-\(\delta\) 定义。
\end{quote}

\textbf{前置知识}：本章依赖 Ch 23（坐标系与函数概念，特别是函数的单调性、奇偶性、有界性等基本性质，以及反函数理论）。本章为 Ch 25（指数函数与对数函数）和后续章节提供最基础的函数模型。

\hrulefill

一次函数是最简单的非平凡函数，也是唯一一类图像为直线的函数。本节从正比例函数 \(y = kx\) 出发，引入一次函数 \(y = kx + b\) 的一般形式，系统研究了斜率 \(k\) 和截距 \(b\) 的几何意义（斜率刻画直线的倾斜方向和陡峭程度，截距刻画直线与坐标轴的交点位置），以及一次函数的单调性（由 \(k\) 的正负号决定）和奇偶性（当 \(b = 0\) 时为奇函数，否则无奇偶性）。一次函数在数学中扮演着基础而关键的角色：它是微分学中''线性逼近''思想的原型------导数 \(f'(a)\) 本质上就是用一次函数 \(y = f(a) + f'(a)(x - a)\) 来局部近似 \(f(x)\)；它也是线性代数中''线性映射''概念在实数域上的具体体现。本节还给出了通过两点确定一次函数的公式，这一结果在解析几何和数值分析中都有广泛应用。

\textbf{定义 17.1}（正比例函数）形如 \(y = kx\)（\(k \neq 0\)）的函数称为\textbf{正比例函数}，定义域为 \(\mathbb{R}\)，\(k\) 称为\textbf{比例系数}。

\textbf{定理 17.1} 正比例函数 \(y = kx\)（\(k \neq 0\)）的图像是一条过原点的直线。

\textbf{证明}：设 \(P_1(x_1, kx_1)\) 和 \(P_2(x_2, kx_2)\) 是图像上任意两点（\(x_1 \neq x_2\)）。则斜率

\[\frac{kx_2 - kx_1}{x_2 - x_1} = \frac{k(x_2 - x_1)}{x_2 - x_1} = k\]

斜率为常数，且过原点 \((0, 0)\)，故图像为直线。\(\blacksquare\)

\textbf{定理 17.2} 正比例函数 \(y = kx\)（\(k \neq 0\)）是奇函数。

\textbf{证明}：\(f(-x) = k(-x) = -kx = -f(x)\)。\(\blacksquare\)

\subsubsection{17.1.2 一次函数的定义与性质}\label{ux4e00ux6b21ux51fdux6570ux7684ux5b9aux4e49ux4e0eux6027ux8d28}

\textbf{定义 17.2}（一次函数）形如

\[y = kx + b \quad (k \neq 0)\]

的函数称为\textbf{一次函数}，定义域为 \(\mathbb{R}\)。当 \(b = 0\) 时，退化为正比例函数。

\textbf{定义 17.3}（斜率与截距）： - \(k\) 称为直线 \(y = kx + b\) 的\textbf{斜率}，它描述了直线的倾斜程度和方向。\(|k|\) 越大，直线越陡。 - \(b\) 称为直线在 \(y\) 轴上的\textbf{截距}，即直线与 \(y\) 轴交点 \((0, b)\) 的纵坐标。

\textbf{几何意义}：斜率 \(k\) 表示直线上每向右移动 \(1\) 个单位，\(y\) 值的变化量。严格地说，若直线上两点 \(P_1(x_1, y_1)\), \(P_2(x_2, y_2)\)（\(x_1 \neq x_2\)），则：

\[k = \frac{y_2 - y_1}{x_2 - x_1}\]

\textbf{定理 17.3}（一次函数的图像）一次函数 \(y = kx + b\)（\(k \neq 0\)）的图像是一条直线，斜率为 \(k\)，\(y\) 轴截距为 \(b\)。

\textbf{证明}：设 \(P_1(x_1, kx_1 + b)\) 和 \(P_2(x_2, kx_2 + b)\)（\(x_1 \neq x_2\)）。则斜率

\[\frac{(kx_2 + b) - (kx_1 + b)}{x_2 - x_1} = \frac{k(x_2 - x_1)}{x_2 - x_1} = k\]

为常数。由解析几何的基本定理，斜率为常数的曲线是直线。\(\blacksquare\)

\textbf{定理 17.4}（一次函数的单调性）一次函数 \(y = kx + b\)（\(k \neq 0\)）在 \(\mathbb{R}\) 上严格单调： - 当 \(k > 0\) 时，严格递增； - 当 \(k < 0\) 时，严格递减。

\textbf{证明}：设 \(x_1 < x_2\)。则

\[f(x_2) - f(x_1) = (kx_2 + b) - (kx_1 + b) = k(x_2 - x_1)\]

当 \(k > 0\) 时，\(k(x_2 - x_1) > 0\)，故 \(f(x_2) > f(x_1)\)，函数严格递增。

当 \(k < 0\) 时，\(k(x_2 - x_1) < 0\)，故 \(f(x_2) < f(x_1)\)，函数严格递减。\(\blacksquare\)

\textbf{定理 17.5}（一次函数的奇偶性）一次函数 \(y = kx + b\)（\(k \neq 0\)）： - 当 \(b = 0\) 时，为奇函数； - 当 \(b \neq 0\) 时，既非奇函数也非偶函数。

\textbf{证明}：\(f(-x) = -kx + b\)。若 \(f\) 是奇函数，需 \(f(-x) = -f(x) = -kx - b\)，即 \(-kx + b = -kx - b\)，得 \(b = 0\)。若 \(f\) 是偶函数，需 \(f(-x) = f(x)\)，即 \(-kx + b = kx + b\)，得 \(k = 0\)，与 \(k \neq 0\) 矛盾。\(\blacksquare\)

\subsubsection{17.1.3 两点确定一次函数}\label{ux4e24ux70b9ux786eux5b9aux4e00ux6b21ux51fdux6570}

\textbf{定理 17.6} 过平面上两个不同点 \(P_1(x_1, y_1)\) 和 \(P_2(x_2, y_2)\)（\(x_1 \neq x_2\)），存在唯一的一次函数 \(y = kx + b\)，其中：

\[k = \frac{y_2 - y_1}{x_2 - x_1}, \quad b = y_1 - kx_1\]

\textbf{证明}：设 \(y = kx + b\)，代入两点：

\[\begin{cases} y_1 = kx_1 + b \\ y_2 = kx_2 + b \end{cases}\]

两式相减：\(y_2 - y_1 = k(x_2 - x_1)\)，故 \(k = \frac{y_2 - y_1}{x_2 - x_1}\)。代入第一式得 \(b = y_1 - kx_1\)。由线性方程组的唯一解性，\(k\) 和 \(b\) 唯一确定。\(\blacksquare\)

\hrulefill

反比例函数 \(y = \frac{k}{x}\) 是最简单的有理函数，也是双曲线这一重要曲线族的代表。本节从定义出发，系统研究了反比例函数的奇偶性（奇函数，图像关于原点对称）、单调性（在定义域的两个连通分支上分别单调，但不具有整体单调性）和渐近线行为（\(x\) 轴和 \(y\) 轴为渐近线）。反比例函数的图像是等轴双曲线，其矩形面积性质（图像上任意一点与坐标轴围成的矩形面积为常数 \(|k|\)）是双曲线几何性质的直接体现。从物理学的角度看，反比例函数描述了反比关系（如波义耳定律中压强与体积的反比关系），因此也被称为''反比函数''。反比例函数的渐近线概念是微积分中渐近线理论的入门------它直观地展示了函数''无限接近但永不触及''某种直线的行为模式。

\textbf{定义 17.4}（反比例函数）形如

\[y = \frac{k}{x} \quad (k \neq 0)\]

的函数称为\textbf{反比例函数}，其中 \(k\) 称为\textbf{比例系数}。定义域为 \(\{x \in \mathbb{R} \mid x \neq 0\} = \mathbb{R} \setminus \{0\}\)。

\textbf{定理 17.7}（反比例函数的奇偶性）反比例函数 \(y = \frac{k}{x}\)（\(k \neq 0\)）是\textbf{奇函数}。

\textbf{证明}：定义域 \(\mathbb{R} \setminus \{0\}\) 关于原点对称。对于任意 \(x \neq 0\)：

\[f(-x) = \frac{k}{-x} = -\frac{k}{x} = -f(x)\]

因此 \(f\) 是奇函数。\(\blacksquare\)

\textbf{定理 17.8}（反比例函数的单调性）反比例函数 \(y = \frac{k}{x}\)（\(k \neq 0\)）：

\begin{enumerate}
\def\labelenumi{\arabic{enumi}.}
\tightlist
\item
  当 \(k > 0\) 时，在 \((-\infty, 0)\) 和 \((0, +\infty)\) 上分别严格递减。
\item
  当 \(k < 0\) 时，在 \((-\infty, 0)\) 和 \((0, +\infty)\) 上分别严格递增。
\end{enumerate}

\textbf{注意}：反比例函数在其定义域上\textbf{不}具有整体单调性，而是在各连通分支上分别单调。

\textbf{证明}：设 \(0 < x_1 < x_2\)，\(k > 0\)。则

\[f(x_2) - f(x_1) = \frac{k}{x_2} - \frac{k}{x_1} = \frac{k(x_1 - x_2)}{x_1 x_2}\]

由于 \(x_1 - x_2 < 0\)，\(x_1 x_2 > 0\)，\(k > 0\)，故 \(f(x_2) - f(x_1) < 0\)，即 \(f\) 在 \((0, +\infty)\) 上严格递减。

类似地可证 \(f\) 在 \((-\infty, 0)\) 上严格递减。

\(k < 0\) 的情形类似证明。\(\blacksquare\)

\subsubsection{17.2.2 渐近线}\label{ux6e10ux8fd1ux7ebf}

\textbf{定义 17.5}（渐近线）设 \(f: D \to \mathbb{R}\)。若存在常数 \(L\)，使得

\[\lim_{x \to +\infty} f(x) = L \quad \text{或} \quad \lim_{x \to -\infty} f(x) = L\]

则称直线 \(y = L\) 为 \(f\) 的\textbf{水平渐近线}。若

\[\lim_{x \to a^+} f(x) = \pm\infty \quad \text{或} \quad \lim_{x \to a^-} f(x) = \pm\infty\]

则称直线 \(x = a\) 为 \(f\) 的\textbf{垂直渐近线}。

\textbf{定理 17.9}（反比例函数的渐近线）\(y = \frac{k}{x}\)（\(k \neq 0\)）的图像以 \(x\) 轴和 \(y\) 轴为渐近线。即：

\[\lim_{x \to \pm\infty} \frac{k}{x} = 0, \quad \lim_{x \to 0^{\pm}} \left|\frac{k}{x}\right| = +\infty\]

\textbf{证明}：对于任意 \(\varepsilon > 0\)，取 \(M = \frac{|k|}{\varepsilon}\)。当 \(|x| > M\) 时，

\[\left|\frac{k}{x}\right| = \frac{|k|}{|x|} < \frac{|k|}{M} = \varepsilon\]

因此 \(\lim_{x \to \pm\infty} \frac{k}{x} = 0\)，即 \(y = 0\)（\(x\) 轴）是水平渐近线。

对于垂直渐近线：当 \(x \to 0^+\) 时，\(\frac{k}{x} \to +\infty\)（\(k > 0\)）或 \(\frac{k}{x} \to -\infty\)（\(k < 0\)）。因此 \(x = 0\)（\(y\) 轴）是垂直渐近线。\(\blacksquare\)

\subsubsection{17.2.3 面积性质}\label{ux9762ux79efux6027ux8d28}

\textbf{定理 17.10}（矩形面积性质）设 \(P(x, y)\) 是反比例函数 \(y = \frac{k}{x}\)（\(k \neq 0\)）图像上的任意一点，过 \(P\) 分别作 \(x\) 轴和 \(y\) 轴的垂线，垂足分别为 \(A\) 和 \(B\)，则矩形 \(OAPB\) 的面积为 \(|k|\)。

\textbf{证明}：\(OAPB\) 的面积 \(= |x| \cdot |y| = |x| \cdot \left|\frac{k}{x}\right| = |k|\)。\(\blacksquare\)

\hrulefill

二次函数是多项式函数中最重要的非平凡特例，其图像为抛物线------这是圆锥曲线中最为人熟知的一种曲线。本节从二次函数的一般式 \(y = ax^2 + bx + c\) 出发，通过配方法得到顶点式 \(y = a(x - h)^2 + k\)，进而推导出抛物线的对称轴、顶点坐标和开口方向（由 \(a\) 的正负号决定）。二次函数的判别式 \(\Delta = b^2 - 4ac\) 是判断抛物线与 \(x\) 轴交点个数的关键工具，而韦达定理则建立了根与系数的代数关系。二次函数的最值问题（在 \(\mathbb{R}\) 上和闭区间上）是优化问题的最简单模型，也是微积分中极值理论的入门。在物理学中，自由落体的位移公式 \(s = \frac{1}{2}gt^2\) 和抛体运动的轨迹方程都是二次函数的经典实例，这使得二次函数成为连接数学与物理的重要桥梁。

\textbf{定义 17.6}（二次函数）形如

\[y = ax^2 + bx + c \quad (a \neq 0)\]

的函数称为\textbf{二次函数}，定义域为 \(\mathbb{R}\)。其中 \(a\) 称为\textbf{二次项系数}，\(b\) 称为\textbf{一次项系数}，\(c\) 称为\textbf{常数项}。

\subsubsection{17.3.2 配方法与顶点式}\label{ux914dux65b9ux6cd5ux4e0eux9876ux70b9ux5f0f}

\textbf{定理 17.11}（配方法）二次函数 \(y = ax^2 + bx + c\) 可以配方为：

\[y = a\left(x + \frac{b}{2a}\right)^2 + \frac{4ac - b^2}{4a}\]

\textbf{证明}：

\[y = ax^2 + bx + c = a\left(x^2 + \frac{b}{a}x\right) + c\]

\[= a\left(x^2 + \frac{b}{a}x + \frac{b^2}{4a^2} - \frac{b^2}{4a^2}\right) + c\]

\[= a\left(x + \frac{b}{2a}\right)^2 - \frac{b^2}{4a} + c\]

\[= a\left(x + \frac{b}{2a}\right)^2 + \frac{4ac - b^2}{4a}\]

\(\blacksquare\)

\textbf{定义 17.7}（顶点式）令 \(h = -\frac{b}{2a}\)，\(k = \frac{4ac - b^2}{4a}\)，则二次函数可写为\textbf{顶点式}：

\[y = a(x - h)^2 + k\]

其中 \((h, k)\) 为抛物线的\textbf{顶点}。

\subsubsection{17.3.3 交点式}\label{ux4ea4ux70b9ux5f0f}

\textbf{定义 17.8}（交点式）当二次函数 \(y = ax^2 + bx + c\) 的判别式 \(\Delta = b^2 - 4ac \geq 0\) 时，方程 \(ax^2 + bx + c = 0\) 有实根 \(x_1, x_2\)（允许 \(x_1 = x_2\)），则可写为\textbf{交点式}：

\[y = a(x - x_1)(x - x_2)\]

\textbf{证明}：由因式分解定理，当 \(x_1, x_2\) 是 \(ax^2 + bx + c = 0\) 的根时，

\[ax^2 + bx + c = a(x - x_1)(x - x_2)\]

展开右边：\(a(x^2 - (x_1 + x_2)x + x_1 x_2) = ax^2 - a(x_1 + x_2)x + ax_1 x_2\)

与 \(ax^2 + bx + c\) 比较系数：\(b = -a(x_1 + x_2)\)，\(c = ax_1 x_2\)，这正是韦达定理的结论。\(\blacksquare\)

\subsubsection{17.3.4 抛物线的基本性质}\label{ux629bux7269ux7ebfux7684ux57faux672cux6027ux8d28}

\textbf{定理 17.12} 二次函数 \(y = ax^2 + bx + c\)（\(a \neq 0\)）的图像是一条\textbf{抛物线}，具有以下性质：

\begin{enumerate}
\def\labelenumi{\arabic{enumi}.}
\tightlist
\item
  \textbf{对称轴}为 \(x = -\frac{b}{2a}\)
\item
  \textbf{顶点}为 \(\left(-\frac{b}{2a}, \frac{4ac - b^2}{4a}\right)\)
\item
  当 \(a > 0\) 时，开口向上；当 \(a < 0\) 时，开口向下
\end{enumerate}

\textbf{证明}：由顶点式 \(y = a(x - h)^2 + k\)。

\begin{enumerate}
\def\labelenumi{(\arabic{enumi})}
\item
  对于任意 \(\delta\)，\(f(h + \delta) = a\delta^2 + k = f(h - \delta)\)，所以图像关于 \(x = h = -\frac{b}{2a}\) 对称。
\item
  当 \(x = h\) 时，\(y = k = \frac{4ac - b^2}{4a}\)，这就是顶点的纵坐标。
\item
  \(y - k = a(x - h)^2\)。当 \(a > 0\) 时，\(y - k \geq 0\)，即 \(y \geq k\)，开口向上；当 \(a < 0\) 时，\(y - k \leq 0\)，即 \(y \leq k\)，开口向下。\(\blacksquare\)
\end{enumerate}

\subsubsection{17.3.5 二次函数的最值}\label{ux4e8cux6b21ux51fdux6570ux7684ux6700ux503c}

\textbf{定理 17.13}（在 \(\mathbb{R}\) 上的最值）二次函数 \(y = ax^2 + bx + c\)（\(a \neq 0\)）： - 当 \(a > 0\) 时，有最小值 \(y_{\min} = \frac{4ac - b^2}{4a}\)，在 \(x = -\frac{b}{2a}\) 处取到； - 当 \(a < 0\) 时，有最大值 \(y_{\max} = \frac{4ac - b^2}{4a}\)，在 \(x = -\frac{b}{2a}\) 处取到。

\textbf{证明}：由顶点式 \(y = a(x - h)^2 + k\)。

当 \(a > 0\) 时，\(a(x - h)^2 \geq 0\)，等号当且仅当 \(x = h\)。因此 \(y \geq k\)，最小值为 \(k\)。

当 \(a < 0\) 时，\(a(x - h)^2 \leq 0\)，等号当且仅当 \(x = h\)。因此 \(y \leq k\)，最大值为 \(k\)。\(\blacksquare\)

\textbf{定理 17.14}（在闭区间上的最值）设二次函数 \(y = ax^2 + bx + c\)（\(a > 0\)，开口向上），对称轴 \(x = h = -\frac{b}{2a}\)。在闭区间 \([m, n]\) 上：

\begin{enumerate}
\def\labelenumi{\arabic{enumi}.}
\tightlist
\item
  若 \(h \in [m, n]\)，则最小值在 \(x = h\) 处取到，最大值在 \(x = m\) 或 \(x = n\) 处取到；
\item
  若 \(h < m\)，则最小值在 \(x = m\) 处取到，最大值在 \(x = n\) 处取到；
\item
  若 \(h > n\)，则最小值在 \(x = n\) 处取到，最大值在 \(x = m\) 处取到。
\end{enumerate}

\textbf{证明}：当 \(a > 0\) 时，\(f\) 在 \((-\infty, h]\) 上严格递减，在 \([h, +\infty)\) 上严格递增。因此在 \([m, n]\) 上的最值取决于 \(h\) 与区间的位置关系，结论由单调性的分析直接得到。\(\blacksquare\)

\subsubsection{17.3.6 判别式与根的关系}\label{ux5224ux522bux5f0fux4e0eux6839ux7684ux5173ux7cfb}

\textbf{定理 17.15}（判别式）一元二次方程 \(ax^2 + bx + c = 0\)（\(a \neq 0\)）的判别式为 \(\Delta = b^2 - 4ac\)。

\begin{enumerate}
\def\labelenumi{\arabic{enumi}.}
\tightlist
\item
  当 \(\Delta > 0\) 时，方程有两个不相等的实数根：\(x = \frac{-b \pm \sqrt{\Delta}}{2a}\)
\item
  当 \(\Delta = 0\) 时，方程有两个相等的实数根：\(x = -\frac{b}{2a}\)
\item
  当 \(\Delta < 0\) 时，方程无实数根
\end{enumerate}

\textbf{证明}：由配方法，

\[ax^2 + bx + c = a\left(x + \frac{b}{2a}\right)^2 + \frac{4ac - b^2}{4a} = 0\]

\[\left(x + \frac{b}{2a}\right)^2 = \frac{b^2 - 4ac}{4a^2} = \frac{\Delta}{4a^2}\]

当 \(\Delta > 0\) 时，\(x + \frac{b}{2a} = \pm \frac{\sqrt{\Delta}}{2a}\)，得两根。

当 \(\Delta = 0\) 时，\(x = -\frac{b}{2a}\)。

当 \(\Delta < 0\) 时，\(\frac{\Delta}{4a^2} < 0\)，实数的平方不可能为负，故无实根。\(\blacksquare\)

\textbf{定理 17.16}（韦达定理）设 \(x_1, x_2\) 是方程 \(ax^2 + bx + c = 0\)（\(a \neq 0\)）的两根，则：

\[x_1 + x_2 = -\frac{b}{a}, \quad x_1 x_2 = \frac{c}{a}\]

\textbf{证明}：由因式分解，\(ax^2 + bx + c = a(x - x_1)(x - x_2)\)。

展开：\(a[x^2 - (x_1 + x_2)x + x_1 x_2] = ax^2 - a(x_1 + x_2)x + ax_1 x_2\)

比较系数：\(b = -a(x_1 + x_2)\)，\(c = ax_1 x_2\)，即 \(x_1 + x_2 = -\frac{b}{a}\)，\(x_1 x_2 = \frac{c}{a}\)。\(\blacksquare\)

\hrulefill

幂函数 \(y = x^{\alpha}\) 是最基本的一类函数------当 \(\alpha\) 为正整数时得到多项式函数，当 \(\alpha\) 为负整数时得到有理函数，当 \(\alpha\) 为分数时涉及根式。本节系统研究了当 \(\alpha\) 取不同值（正整数、负整数、\(1/2\) 等）时幂函数的定义域、值域、单调性和奇偶性。特别地，\(\alpha\) 的奇偶性决定了幂函数的奇偶性：当 \(\alpha\) 为奇数时 \(x^{\alpha}\) 为奇函数，当 \(\alpha\) 为偶数时 \(x^{\alpha}\) 为偶函数。幂函数的图像特征随 \(\alpha\) 的变化呈现出丰富的规律：所有幂函数都过点 \((1, 1)\)；当 \(\alpha > 0\) 时，函数在 \((0, +\infty)\) 上严格递增且 \(\lim_{x \to 0^+} x^{\alpha} = 0\)；当 \(\alpha < 0\) 时，函数在 \((0, +\infty)\) 上严格递减且 \(\lim_{x \to 0^+} x^{\alpha} = +\infty\)。幂函数族是后续学习指数函数、对数函数和广义幂函数的基础，其图像在 \((0, 1)\) 和 \((1, +\infty)\) 两个区间上的交叉行为揭示了指数 \(\alpha\) 对函数增长速度的精细调控。

\textbf{定义 17.9}（幂函数）形如

\[y = x^{\alpha}\]

的函数称为\textbf{幂函数}，其中 \(\alpha\) 为常数。

\subsubsection{17.4.2 常见幂函数的性质}\label{ux5e38ux89c1ux5e42ux51fdux6570ux7684ux6027ux8d28}

\textbf{定理 17.17} 设 \(\alpha\) 为正整数 \(n\)，则 \(f(x) = x^n\) 的性质如下：

\begin{enumerate}
\def\labelenumi{\arabic{enumi}.}
\tightlist
\item
  \textbf{定义域}：\(\mathbb{R}\)。
\item
  当 \(n\) 为奇数时，\(f\) 是奇函数，在 \(\mathbb{R}\) 上严格递增。
\item
  当 \(n\) 为偶数时，\(f\) 是偶函数，在 \((-\infty, 0]\) 上严格递减，在 \([0, +\infty)\) 上严格递增。
\end{enumerate}

\textbf{证明}：

\begin{enumerate}
\def\labelenumi{(\arabic{enumi})}
\setcounter{enumi}{1}
\tightlist
\item
  设 \(n\) 为奇数。\(f(-x) = (-x)^n = -x^n = -f(x)\)，故 \(f\) 为奇函数。
\end{enumerate}

设 \(x_1 < x_2\)。若 \(0 \leq x_1 < x_2\)，则 \(x_1^n < x_2^n\)（因为 \(x^n\) 在 \([0, +\infty)\) 上严格递增，可用归纳法证明）。

若 \(x_1 < x_2 < 0\)，令 \(u = -x_2, v = -x_1\)，则 \(0 < u < v\)，且 \(x_2^n - x_1^n = (-u)^n - (-v)^n = -(u^n - v^n)\)。因为 \(n\) 为奇数且 \(u^n < v^n\)，所以 \(x_2^n - x_1^n > 0\)。

若 \(x_1 < 0 < x_2\)，则 \(x_1^n < 0 < x_2^n\)。综合三种情况，\(f\) 在 \(\mathbb{R}\) 上严格递增。

\begin{enumerate}
\def\labelenumi{(\arabic{enumi})}
\setcounter{enumi}{2}
\tightlist
\item
  设 \(n\) 为偶数。\(f(-x) = (-x)^n = x^n = f(x)\)，故 \(f\) 为偶函数。
\end{enumerate}

由 \(f\) 是偶函数，只需分析 \(x \geq 0\) 的单调性。当 \(0 \leq x_1 < x_2\) 时，\(x_1^n < x_2^n\)，故在 \([0, +\infty)\) 上严格递增。由偶函数的对称性，在 \((-\infty, 0]\) 上严格递减。\(\blacksquare\)

\textbf{定理 17.18} 设 \(\alpha = -1\)，则 \(f(x) = x^{-1} = \frac{1}{x}\) 就是反比例函数（\(k = 1\) 的情形），其性质已在 17.2 节中讨论。

\textbf{证明}：\(f(x) = x^{-1} = \frac{1}{x}\) 即 \(k = 1\) 的反比例函数 \(y = \frac{k}{x}\)，性质已在定义17.5中详述。\(\blacksquare\)

\textbf{定理 17.19} 设 \(\alpha = \frac{1}{2}\)，则 \(f(x) = x^{1/2} = \sqrt{x}\) 的性质如下：

\begin{enumerate}
\def\labelenumi{\arabic{enumi}.}
\tightlist
\item
  \textbf{定义域}：\([0, +\infty)\)。
\item
  \textbf{值域}：\([0, +\infty)\)。
\item
  在 \([0, +\infty)\) 上严格递增。
\end{enumerate}

\textbf{证明}：设 \(0 \leq x_1 < x_2\)。则 \(0 \leq \sqrt{x_1} < \sqrt{x_2}\)（因为若 \(\sqrt{x_1} \geq \sqrt{x_2}\)，则 \(x_1 = (\sqrt{x_1})^2 \geq (\sqrt{x_2})^2 = x_2\)，矛盾）。\(\blacksquare\)

\subsubsection{17.4.3 幂函数的一般性质}\label{ux5e42ux51fdux6570ux7684ux4e00ux822cux6027ux8d28}

\textbf{定理 17.20}（幂函数的公共特征）对于任意实数 \(\alpha\)，幂函数 \(y = x^{\alpha}\) 的图像都过点 \((1, 1)\)。

\textbf{证明}：\(f(1) = 1^{\alpha} = 1\)。\(\blacksquare\)

\hrulefill

函数的零点理论将函数研究与方程求解统一起来------\(f(x) = 0\) 的根恰好对应函数图像与 \(x\) 轴的交点。本节从一次函数和二次函数的零点出发，系统建立了零点与方程根之间的对应关系，并利用韦达定理分析了二次函数零点的分布（两个正根、两个负根、一正一负根的条件）。一元二次不等式的解法是零点理论的直接应用------通过判别式 \(\Delta\) 和二次项系数 \(a\) 的符号，可以将 \(ax^2 + bx + c > 0\)（或 \(< 0\)）的解集完全确定。零点理论在高等数学中有着深远的推广：连续函数的介值定理（Ch 29）保证了闭区间上的连续函数若在端点处异号则必有零点，而代数基本定理则断言 \(n\) 次复系数多项式在复数域中恰有 \(n\) 个根（计重数）。

\textbf{定义 17.10} 函数 \(f: D \to \mathbb{R}\) 的\textbf{零点}是使 \(f(x) = 0\) 的 \(x\) 值，即方程 \(f(x) = 0\) 的\textbf{根}。

\textbf{定理 17.25}（一次函数的零点）一次函数 \(y = kx + b\)（\(k \neq 0\)）有且仅有一个零点 \(x = -\frac{b}{k}\)。

\textbf{证明}：\(kx + b = 0 \iff x = -\frac{b}{k}\)。\(\blacksquare\)

\textbf{定理 17.26}（二次函数零点的分布）二次函数 \(f(x) = ax^2 + bx + c\)（\(a \neq 0\)）的零点分布：

\begin{enumerate}
\def\labelenumi{\arabic{enumi}.}
\tightlist
\item
  两个正根 \(\iff\) \(\Delta > 0\)，\(x_1 + x_2 > 0\)，\(x_1 x_2 > 0\)，即 \(\Delta > 0\)，\(-\frac{b}{a} > 0\)，\(\frac{c}{a} > 0\)。
\item
  两个负根 \(\iff\) \(\Delta > 0\)，\(-\frac{b}{a} < 0\)，\(\frac{c}{a} > 0\)。
\item
  一正一负根 \(\iff\) \(\Delta > 0\)，\(\frac{c}{a} < 0\)。
\end{enumerate}

\textbf{证明}：由韦达定理，\(x_1 + x_2 = -\frac{b}{a}\)，\(x_1 x_2 = \frac{c}{a}\)。

\begin{enumerate}
\def\labelenumi{(\arabic{enumi})}
\tightlist
\item
  两个正根 \(\iff\) \(x_1 > 0\) 且 \(x_2 > 0\) \(\iff\) \(x_1 + x_2 > 0\) 且 \(x_1 x_2 > 0\) 且 \(\Delta \geq 0\)。由于 \(x_1 \neq x_2\) 时 \(\Delta > 0\)，等号仅在 \(x_1 = x_2\) 时取到。\(\blacksquare\)
\end{enumerate}

\subsubsection{17.5.2 一元二次不等式}\label{ux4e00ux5143ux4e8cux6b21ux4e0dux7b49ux5f0f-1}

\textbf{定理 17.27}（一元二次不等式解法）设 \(a > 0\)，\(\Delta = b^2 - 4ac\)。

\begin{enumerate}
\def\labelenumi{\arabic{enumi}.}
\tightlist
\item
  若 \(\Delta > 0\)，\(ax^2 + bx + c = 0\) 有两实根 \(x_1 < x_2\)：

  \begin{itemize}
  \tightlist
  \item
    \(ax^2 + bx + c > 0\) 的解集为 \((-\infty, x_1) \cup (x_2, +\infty)\)
  \item
    \(ax^2 + bx + c < 0\) 的解集为 \((x_1, x_2)\)
  \end{itemize}
\item
  若 \(\Delta = 0\)，\(ax^2 + bx + c = 0\) 有重根 \(x_0 = -\frac{b}{2a}\)：

  \begin{itemize}
  \tightlist
  \item
    \(ax^2 + bx + c > 0\) 的解集为 \(\{x \mid x \neq x_0\} = \mathbb{R} \setminus \{x_0\}\)
  \item
    \(ax^2 + bx + c < 0\) 的解集为 \(\emptyset\)
  \end{itemize}
\item
  若 \(\Delta < 0\)：

  \begin{itemize}
  \tightlist
  \item
    \(ax^2 + bx + c > 0\) 的解集为 \(\mathbb{R}\)
  \item
    \(ax^2 + bx + c < 0\) 的解集为 \(\emptyset\)
  \end{itemize}
\end{enumerate}

\textbf{证明}：(1) 由交点式 \(f(x) = a(x - x_1)(x - x_2)\)（\(a > 0\)）。当 \(x < x_1\) 时，\(x - x_1 < 0\)，\(x - x_2 < 0\)，故 \(f(x) > 0\)。当 \(x_1 < x < x_2\) 时，\(x - x_1 > 0\)，\(x - x_2 < 0\)，故 \(f(x) < 0\)。当 \(x > x_2\) 时，\(x - x_1 > 0\)，\(x - x_2 > 0\)，故 \(f(x) > 0\)。

\begin{enumerate}
\def\labelenumi{(\arabic{enumi})}
\setcounter{enumi}{1}
\item
  \(f(x) = a(x - x_0)^2 \geq 0\)，等号当且仅当 \(x = x_0\)。
\item
  \(f(x) = a\left(x + \frac{b}{2a}\right)^2 + \frac{4ac - b^2}{4a} = a\left(x + \frac{b}{2a}\right)^2 + \frac{-\Delta}{4a}\)。由于 \(\Delta < 0\)，\(-\Delta > 0\)，且 \(a > 0\)，故 \(f(x) > 0\) 对所有 \(x\) 成立。\(\blacksquare\)
\end{enumerate}

\hrulefill

函数模型的综合比较是理解和应用各类基本初等函数的关键。本节聚焦于一次函数、二次函数、幂函数、指数函数和对数函数的增长速度比较，确立了基本初等函数的''增长阶''：当 \(x \to +\infty\) 时，增长速度从慢到快依次为 \(\ln x \ll x^{\alpha}\)（\(\alpha > 0\)）\(\ll a^x\)（\(a > 1\)）\(\ll x! \ll x^x\)。这一增长阶的严格证明依赖于极限理论和洛必达法则（Ch 30），但本节通过初等方法（二项式定理和放缩技巧）给出了核心结论的直观论证。增长阶的比较在算法分析（判断算法复杂度）、数值计算（评估误差增长）和物理学（比较不同尺度下的主导效应）中都是不可或缺的工具。例如，指数增长最终压倒任何多项式增长这一事实，是理解指数增长''爆炸性''特征的数学基础，也是流行病学、金融学中复利效应的核心数学原理。

\textbf{定理 17.28} 设 \(a > 1\)，\(\alpha > 0\)。则指数函数的增长最终超过任意幂函数：

\[\lim_{x \to +\infty} \frac{x^{\alpha}}{a^x} = 0\]

\textbf{证明}：令 \(b = a^{1/\alpha}\)，则 \(b > 1\)。于是

\[\frac{x^{\alpha}}{a^x} = \left(\frac{x}{b^x}\right)^{\alpha}.\]

只需证明 \(\lim_{x \to +\infty} \frac{x}{b^x} = 0\)。设 \(b = 1 + h\)，\(h > 0\)。对任意 \(x > 1\)，令 \(n = \lfloor x \rfloor\)，则 \(n \leq x < n+1\)。由二项式定理，

\[b^x \geq b^n = (1+h)^n \geq \binom{n}{2}h^2 = \frac{n(n-1)}{2}h^2 > \frac{(x-1)(x-2)}{2}h^2.\]

因此

\[0 \leq \frac{x}{b^x} < \frac{2x}{(x-1)(x-2)h^2} \to 0 \quad (x \to +\infty).\]

从而 \(\frac{x}{b^x} \to 0\)，故 \(\frac{x^{\alpha}}{a^x} \to 0\)。\(\blacksquare\)

\textbf{定理 17.29} 设 \(\alpha > 0\)。则对数函数的增长最终慢于任意正幂函数：

\[\lim_{x \to +\infty} \frac{\ln x}{x^{\alpha}} = 0\]

\textbf{证明}：令 \(y = \ln x\)，则 \(x = e^y\)，且当 \(x \to +\infty\) 时 \(y \to +\infty\)。于是

\[\frac{\ln x}{x^{\alpha}} = \frac{y}{e^{\alpha y}} = \frac{1}{\alpha} \cdot \frac{\alpha y}{e^{\alpha y}}.\]

设 \(t = \alpha y\)，则需证 \(\lim_{t \to +\infty} \frac{t}{e^t} = 0\)。注意到 \(e^t = \sum_{k=0}^{\infty} \frac{t^k}{k!}\)，故当 \(t > 0\) 时 \(e^t > \frac{t^2}{2}\)。因此

\[0 \leq \frac{t}{e^t} < \frac{t}{t^2/2} = \frac{2}{t} \to 0 \quad (t \to +\infty).\]

从而 \(\frac{\ln x}{x^{\alpha}} \to 0\)。\(\blacksquare\)

\subsubsection{17.6.2 常见函数增长阶}\label{ux5e38ux89c1ux51fdux6570ux589eux957fux9636}

当 \(x \to +\infty\) 时，常见函数的增长速度从慢到快排列为：

\[\ln x \ll x^{\alpha} \quad (\alpha > 0) \ll a^x \quad (a > 1) \ll x! \ll x^x\]

这表明对数函数增长最慢，而指数函数增长极快------这种差异在算法分析和数值计算中具有重要意义。

含绝对值的函数是函数图像变换理论的重要应用。绝对值运算 \(|f(x)|\) 的效果是将函数图像在 \(x\) 轴以下的部分''翻折''到 \(x\) 轴上方，这一操作在信号处理（整流）、优化问题（绝对偏差最小化）和数值分析（误差度量）中频繁出现。本节从最简单的绝对值线性函数 \(|ax + b|\)（\(V\) 形图像）出发，研究了绝对值二次函数 \(|ax^2 + bx + c|\) 的图像特征，并建立了含绝对值方程和不等式的求解方法。核心技巧是三角不等式 \(|a + b| \leq |a| + |b|\) 及其推广 \(||a| - |b|| \leq |a - b| \leq |a| + |b|\)，这些不等式是分析学中估计误差和控制收敛性的基本工具。绝对值函数 \(|x|\) 在 \(x = 0\) 处连续但不可导，这一性质使其成为微积分中''连续但不可导''的经典反例，也是理解光滑性和可微性之间区别的入门教材。

\textbf{定理 17.30} 设 \(f(x) = |ax + b|\)（\(a \neq 0\)），则：

\[f(x) = \begin{cases} ax + b, & x \geq -\frac{b}{a} \text{（当 } a > 0 \text{ 时）} \\ -(ax + b), & x < -\frac{b}{a} \text{（当 } a > 0 \text{ 时）} \end{cases}\]

\textbf{证明}：由绝对值定义，当 \(ax + b \geq 0\)（即 \(x \geq -\frac{b}{a}\) 当 \(a > 0\)）时 \(|ax+b| = ax+b\)；当 \(ax+b < 0\) 时 \(|ax+b| = -(ax+b)\)。图像在 \(x = -\frac{b}{a}\) 处转折，两侧均为直线段，故为 \(V\) 形。\(\blacksquare\)

图像为 \(V\) 形折线，顶点在 \(\left(-\frac{b}{a}, 0\right)\)。

\textbf{定理 17.31} 设 \(f(x) = |ax^2 + bx + c|\)（\(a > 0\)）。当 \(\Delta > 0\) 时，图像由抛物线 \(y = ax^2 + bx + c\) 在 \(x\) 轴以下的部分翻折到 \(x\) 轴上方得到。

\textbf{证明}：\(|ax^2+bx+c| = ax^2+bx+c\) 当 \(ax^2+bx+c \geq 0\)，\(= -(ax^2+bx+c)\) 当 \(ax^2+bx+c < 0\)。当 \(\Delta > 0\) 时抛物线与 \(x\) 轴有两个交点，将 \(x\) 轴分为三段。中间段 \(ax^2+bx+c < 0\)，取绝对值后翻折到上方。\(\blacksquare\)

\subsubsection{17.7.2 含绝对值的方程和不等式}\label{ux542bux7eddux5bf9ux503cux7684ux65b9ux7a0bux548cux4e0dux7b49ux5f0f}

\textbf{定理 17.32} \(|f(x)| < g(x)\) 等价于 \(-g(x) < f(x) < g(x)\)（其中 \(g(x) > 0\)）。

\textbf{证明}：若 \(f(x) \geq 0\)，则 \(|f(x)| = f(x)\)，不等式变为 \(f(x) < g(x)\)，即 \(0 \leq f(x) < g(x)\)。若 \(f(x) < 0\)，则 \(|f(x)| = -f(x)\)，不等式变为 \(-f(x) < g(x)\)，即 \(-g(x) < f(x) < 0\)。综合两种情形，\(-g(x) < f(x) < g(x)\)。反之亦然。\(\blacksquare\)

\textbf{定理 17.33} \(|f(x)| > g(x)\) 等价于 \(f(x) > g(x)\) 或 \(f(x) < -g(x)\)。

\textbf{证明}：若 \(f(x) \geq 0\)，则 \(|f(x)| = f(x) > g(x)\)，即 \(f(x) > g(x)\)。若 \(f(x) < 0\)，则 \(|f(x)| = -f(x) > g(x)\)，即 \(f(x) < -g(x)\)。综合得 \(f(x) > g(x)\) 或 \(f(x) < -g(x)\)。反之，若 \(f(x) > g(x) \geq 0\) 则 \(|f(x)| = f(x) > g(x)\)；若 \(f(x) < -g(x) \leq 0\) 则 \(|f(x)| = -f(x) > g(x)\)。\(\blacksquare\)

\textbf{定理 17.34}（三角不等式）\(|a + b| \leq |a| + |b|\)，等号当且仅当 \(ab \geq 0\)。

\textbf{证明}：\(|a + b|^2 = a^2 + 2ab + b^2 \leq a^2 + 2|a||b| + b^2 = (|a| + |b|)^2\)。\(\blacksquare\)

\textbf{推论 17.1}（推广的三角不等式）\(||a| - |b|| \leq |a - b| \leq |a| + |b|\)。

\textbf{证明}：\(|a| = |(a - b) + b| \leq |a - b| + |b|\)，故 \(|a| - |b| \leq |a - b|\)。交换 \(a, b\) 得 \(|b| - |a| \leq |a - b|\)。综合得 \(||a| - |b|| \leq |a - b|\)。右边由三角不等式得 \(|a - b| = |a + (-b)| \leq |a| + |b|\)。\(\blacksquare\)

\hrulefill

\textbf{本章展望}：本章研究了一次函数、反比例函数、二次函数和幂函数等基本初等函数。在下一章（Ch 25）中，我们将系统研究指数函数和对数函数，这两类函数是基本初等函数中最重要的非代数函数。特别地，我们将严格定义实数指数幂，并证明指数函数和对数函数的基本性质和运算法则。

\hrulefill

\subsubsection{\texorpdfstring{17.4.4 幂函数 \(y = x^{\alpha}\) 的图像特征分类}{17.4.4 幂函数 y = x\^{}\{\textbackslash alpha\} 的图像特征分类}}\label{ux5e42ux51fdux6570-y-xalpha-ux7684ux56feux50cfux7279ux5f81ux5206ux7c7b}

\textbf{定理 17.21}（\(\alpha > 0\) 时幂函数的行为）设 \(\alpha > 0\)，则幂函数 \(y = x^{\alpha}\)（在 \(x > 0\) 上）满足：

\begin{enumerate}
\def\labelenumi{\arabic{enumi}.}
\tightlist
\item
  在 \((0, +\infty)\) 上严格递增。
\item
  \(\lim_{x \to 0^+} x^{\alpha} = 0\)。
\item
  \(\lim_{x \to +\infty} x^{\alpha} = +\infty\)。
\end{enumerate}

\textbf{证明}：(1) \textbf{严格递增性}：设 \(0 < x_1 < x_2\)，令 \(t = x_2/x_1 > 1\)。则 \(x_2^\alpha / x_1^\alpha = t^\alpha\)。当 \(t > 1\) 且 \(\alpha > 0\) 时，对有理指数 \(r\)，\(t^r > 1\)（由有理指数幂的单调性，可利用 \(t = 1 + s\)（\(s > 0\)）及伯努利不等式 \((1+s)^r \geq 1 + rs > 1\) 对正有理数 \(r\) 建立）；对于一般的正实数 \(\alpha\)，由 \(t^\alpha = \sup\{t^r \mid r \in \mathbb{Q},\, 0 < r \leq \alpha\} > 1\)，故 \(x_2^\alpha > x_1^\alpha\)。(2) 对于任意 \(\varepsilon > 0\)，取 \(\delta = \varepsilon^{1/\alpha}\)，当 \(0 < x < \delta\) 时 \(x^\alpha < \varepsilon\)。(3) 对于任意 \(M > 0\)，取 \(X = M^{1/\alpha}\)，当 \(x > X\) 时 \(x^\alpha > M\)。\(\blacksquare\)

\textbf{定理 17.22}（\(\alpha < 0\) 时幂函数的行为）设 \(\alpha < 0\)，则幂函数 \(y = x^{\alpha}\)（在 \(x > 0\) 上）满足：

\begin{enumerate}
\def\labelenumi{\arabic{enumi}.}
\tightlist
\item
  在 \((0, +\infty)\) 上严格递减。
\item
  \(\lim_{x \to 0^+} x^{\alpha} = +\infty\)。
\item
  \(\lim_{x \to +\infty} x^{\alpha} = 0\)。
\end{enumerate}

\textbf{证明}：令 \(\beta = -\alpha > 0\)，则 \(x^{\alpha} = \frac{1}{x^{\beta}}\)。

\begin{enumerate}
\def\labelenumi{(\arabic{enumi})}
\item
  由定理 17.21，\(x^{\beta}\) 严格递增且为正，故 \(\frac{1}{x^{\beta}}\) 严格递减。
\item
  当 \(x \to 0^+\) 时，\(x^{\beta} \to 0^+\)，故 \(\frac{1}{x^{\beta}} \to +\infty\)。
\item
  当 \(x \to +\infty\) 时，\(x^{\beta} \to +\infty\)，故 \(\frac{1}{x^{\beta}} \to 0^+\)。\(\blacksquare\)
\end{enumerate}

\subsubsection{\texorpdfstring{17.4.5 幂函数的图像关于 \(\alpha\) 的变化}{17.4.5 幂函数的图像关于 \textbackslash alpha 的变化}}\label{ux5e42ux51fdux6570ux7684ux56feux50cfux5173ux4e8e-alpha-ux7684ux53d8ux5316}

\textbf{定理 17.23} 所有幂函数 \(y = x^{\alpha}\)（\(\alpha \neq 0\)）的图像都过点 \((1, 1)\)。

\textbf{证明}：\(1^{\alpha} = 1\) 对所有 \(\alpha\) 成立。\(\blacksquare\)

\textbf{定理 17.24} 设 \(\alpha_1, \alpha_2 > 0\) 且 \(\alpha_1 < \alpha_2\)。则：

\begin{enumerate}
\def\labelenumi{\arabic{enumi}.}
\tightlist
\item
  当 \(0 < x < 1\) 时，\(x^{\alpha_1} > x^{\alpha_2}\)。
\item
  当 \(x > 1\) 时，\(x^{\alpha_1} < x^{\alpha_2}\)。
\end{enumerate}

\textbf{证明}：考虑比值 \(\frac{x^{\alpha_1}}{x^{\alpha_2}} = x^{\alpha_1 - \alpha_2} = \frac{1}{x^{\alpha_2 - \alpha_1}}\)。

\begin{enumerate}
\def\labelenumi{(\arabic{enumi})}
\item
  当 \(0 < x < 1\) 时，\(x^{\alpha_2 - \alpha_1} < 1\)（因为 \(\alpha_2 - \alpha_1 > 0\)，且 \(0 < x < 1\) 时正指数幂使值变小），故 \(\frac{x^{\alpha_1}}{x^{\alpha_2}} = \frac{1}{x^{\alpha_2 - \alpha_1}} > 1\)，即 \(x^{\alpha_1} > x^{\alpha_2}\)。
\item
  当 \(x > 1\) 时，\(x^{\alpha_2 - \alpha_1} > 1\)（因为 \(\alpha_2 - \alpha_1 > 0\)，且 \(x > 1\) 时正指数幂使值变大），故 \(\frac{x^{\alpha_1}}{x^{\alpha_2}} = \frac{1}{x^{\alpha_2 - \alpha_1}} < 1\)，即 \(x^{\alpha_1} < x^{\alpha_2}\)。\(\blacksquare\)
\end{enumerate}

\hrulefill

\hrulefill

\hrulefill

\hrulefill

\hrulefill

\hrulefill

\newpage

\chapter{03-函数与微积分/01-函数与初等函数/03-三角函数.md}\label{ux51fdux6570ux4e0eux5faeux79efux520601-ux51fdux6570ux4e0eux521dux7b49ux51fdux657003-ux4e09ux89d2ux51fdux6570.md}

\section{第45章：三角函数}\label{ux7b2c45ux7ae0ux4e09ux89d2ux51fdux6570}

\textbf{前置知识}：本章依赖 Ch 23（坐标系与函数概念，特别是函数的周期性、奇偶性）以及 Ch 19（锐角三角函数的基本定义）。本章将三角函数从锐角推广到任意角，并通过单位圆给出严格的函数定义。

\hrulefill

角的概念的推广是三角函数理论的第一步。在初等几何中，角通常限于 \(0°\) 到 \(360°\) 之间，但在三角函数的严格定义中，我们需要将角的概念推广到任意实数，包括正角（逆时针旋转）、负角（顺时针旋转）和大于 \(360°\) 的角。本节从有向角的概念出发，引入了终边相同的角（相差 \(2\pi\) 的整数倍）和象限角的概念，并建立了弧度制------一种以弧长与半径之比来度量角度的自然方式。弧度制的引入不仅是数学上的方便，更有着深刻的物理意义：在弧度制下，三角函数的导数公式（如 \((\sin x)' = \cos x\)）变得简洁，这源于 \(\lim_{x \to 0} \frac{\sin x}{x} = 1\) 在弧度制下才成立。弧长公式 \(s = r\theta\) 和扇形面积公式 \(S = \frac{1}{2}r^2\theta\) 也只有在弧度制下才具有如此简洁的形式。

\textbf{定义 19.1}（有向角）设平面上有一条射线 \(OA\)，将其绕端点 \(O\) 旋转到另一位置 \(OB\)，则称此旋转形成了一个\textbf{有向角}。若旋转方向为逆时针，则称为\textbf{正角}；若为顺时针，则称为\textbf{负角}；若射线未旋转，则称为\textbf{零角}。

\subsubsection{19.1.2 象限角与终边相同的角}\label{ux8c61ux9650ux89d2ux4e0eux7ec8ux8fb9ux76f8ux540cux7684ux89d2}

\textbf{定义 19.2}（象限角）在平面直角坐标系中，角的顶点与原点重合，角的始边与 \(x\) 轴正半轴重合。若角的终边落在某一象限内，则称该角为该象限的\textbf{象限角}。若终边落在坐标轴上，则称该角为\textbf{轴线角}。

象限角的分类使得我们可以根据角的终边位置来快速判断三角函数值的符号。然而，同一个象限内可以有无穷多个角，它们共享相同的终边位置，仅在旋转圈数上存在差异。这自然地引出了''终边相同的角''这一概念。

\textbf{定义 19.3} 与角 \(\alpha\) 终边相同的所有角构成的集合为：

\[\{\beta \mid \beta = \alpha + 2k\pi, \, k \in \mathbb{Z}\}\]

\subsubsection{19.1.3 弧度制}\label{ux5f27ux5ea6ux5236}

\textbf{定义 19.4}（弧度）在半径为 \(r\) 的圆中，长度等于半径 \(r\) 的圆弧所对的圆心角的大小定义为 \textbf{1 弧度}（1 rad）。

弧度的定义直接建立起了角度与弧长的二元关系。既然弧度是用弧长与半径之比来度量的，那么一个完整的圆周（弧长为 \(2\pi r\)）对应的弧度数自然就是 \(2\pi\)。这一关系给出了弧度与角度之间的精确换算公式。

\textbf{定理 19.1}（弧度与角度的换算）

\[1 \text{ rad} = \frac{180°}{\pi} \approx 57.296°, \quad 1° = \frac{\pi}{180} \text{ rad}\]

\textbf{证明}：圆的周长为 \(C = 2\pi r\)，对应 \(360°\)。1 弧度对应弧长 \(s = r\)，因此 \(2\pi \text{ rad} = 360°\)，由此可得换算公式。\(\blacksquare\)

\textbf{定理 19.2}（弧长公式与扇形面积公式）设扇形的半径为 \(r\)，圆心角为 \(\theta\)（弧度），则弧长 \(s = r\theta\)，扇形面积 \(S = \frac{1}{2}r^2\theta\)。

\textbf{证明}：圆面积为 \(\pi r^2\)，对应圆心角 \(2\pi\)。扇形面积与圆心角成正比：

\[S = \frac{\theta}{2\pi} \cdot \pi r^2 = \frac{1}{2}r^2\theta\]

弧长类似：\(s = \frac{\theta}{2\pi} \cdot 2\pi r = r\theta\)。\(\blacksquare\)

\hrulefill

三角函数的严格定义是本章的核心。与初等几何中利用直角三角形边长比定义三角函数（仅适用于锐角）不同，本节通过单位圆上的坐标来定义三角函数，从而将定义域自然地推广到任意实数。单位圆定义法揭示了三角函数的本质：正弦和余弦分别是单位圆上点的纵坐标和横坐标，正切是纵坐标与横坐标之比。这一定义的美妙之处在于，它利用单位圆方程 \(x^2 + y^2 = 1\) 立即给出了基本恒等式 \(\sin^2 \alpha + \cos^2 \alpha = 1\)，并且通过几何直观解释了三角函数的周期性（每转一圈回到同一点，\(2\pi\) 为一个周期）、有界性（\(|\sin \alpha| \leq 1\)，\(|\cos \alpha| \leq 1\)）和各象限符号的变化规律。单位圆定义法也是将三角函数推广到复数域和欧拉公式 \(e^{i\theta} = \cos\theta + i\sin\theta\) 的几何基础。

\textbf{定义 19.5}（单位圆）以原点 \(O\) 为圆心，\(1\) 为半径的圆称为\textbf{单位圆}，其方程为 \(x^2 + y^2 = 1\)。

\textbf{定义 19.6}（三角函数）设角 \(\alpha\) 的顶点在原点，始边与 \(x\) 轴正半轴重合，终边与单位圆交于点 \(P(x, y)\)。定义：

\[\sin \alpha = y, \quad \cos \alpha = x\] \[\tan \alpha = \frac{y}{x} \quad (x \neq 0), \quad \cot \alpha = \frac{x}{y} \quad (y \neq 0)\] \[\sec \alpha = \frac{1}{x} \quad (x \neq 0), \quad \csc \alpha = \frac{1}{y} \quad (y \neq 0)\]

\textbf{注} 由于 \(P\) 在单位圆上，\(x^2 + y^2 = 1\)，因此 \(|\sin \alpha| \leq 1\)，\(|\cos \alpha| \leq 1\)。

\subsubsection{19.2.2 定义域和值域}\label{ux5b9aux4e49ux57dfux548cux503cux57df}

\begin{longtable}[]{@{}
  >{\raggedright\arraybackslash}p{(\linewidth - 4\tabcolsep) * \real{0.3000}}
  >{\raggedright\arraybackslash}p{(\linewidth - 4\tabcolsep) * \real{0.4000}}
  >{\raggedright\arraybackslash}p{(\linewidth - 4\tabcolsep) * \real{0.3000}}@{}}
\toprule
\begin{minipage}[b]{\linewidth}\raggedright
函数
\end{minipage} & \begin{minipage}[b]{\linewidth}\raggedright
定义域
\end{minipage} & \begin{minipage}[b]{\linewidth}\raggedright
值域
\end{minipage} \\
\midrule
\endhead
\bottomrule
\endlastfoot
\(\sin \alpha\) & \(\mathbb{R}\) & \([-1, 1]\) \\
\(\cos \alpha\) & \(\mathbb{R}\) & \([-1, 1]\) \\
\(\tan \alpha\) & \(\{\alpha \mid \alpha \neq \frac{\pi}{2} + k\pi, k \in \mathbb{Z}\}\) & \(\mathbb{R}\) \\
\(\cot \alpha\) & \(\{\alpha \mid \alpha \neq k\pi, k \in \mathbb{Z}\}\) & \(\mathbb{R}\) \\
\end{longtable}

\subsubsection{19.2.3 各象限符号}\label{ux5404ux8c61ux9650ux7b26ux53f7}

设角 \(\alpha\) 的终边与单位圆交于 \(P(x, y)\)：

\begin{longtable}[]{@{}lccc@{}}
\toprule
象限 & \(\sin \alpha\) & \(\cos \alpha\) & \(\tan \alpha\) \\
\midrule
\endhead
\bottomrule
\endlastfoot
I & \(+\) & \(+\) & \(+\) \\
II & \(+\) & \(-\) & \(-\) \\
III & \(-\) & \(-\) & \(+\) \\
IV & \(-\) & \(+\) & \(-\) \\
\end{longtable}

\hrulefill

同角三角函数的基本关系是三角恒等变换的起点。本节从单位圆方程 \(x^2 + y^2 = 1\) 出发，推导了三角函数的三大基本关系：勾股关系 \(\sin^2 \alpha + \cos^2 \alpha = 1\)、商数关系 \(\tan \alpha = \frac{\sin \alpha}{\cos \alpha}\) 以及正割和余割关系的衍生恒等式 \(1 + \tan^2 \alpha = \sec^2 \alpha\) 和 \(1 + \cot^2 \alpha = \csc^2 \alpha\)。这些关系式虽然简洁，但它们是所有后续三角恒等式（和差角公式、二倍角公式、半角公式等）的源头。勾股关系 \(\sin^2 \alpha + \cos^2 \alpha = 1\) 是三角学中最基本的恒等式，它与毕达哥拉斯定理（勾股定理）有着深刻的联系------在单位圆中，点 \((\cos \alpha, \sin \alpha)\) 到原点的距离恒为 \(1\)，这正是勾股定理的表述。掌握这些基本关系，就可以在已知一个三角函数值的情况下求出其余所有的三角函数值（通过符号判断确定象限）。

\textbf{定理 19.3}（勾股关系）对任意角 \(\alpha\)：

\[\sin^2 \alpha + \cos^2 \alpha = 1\]

\textbf{证明}：设角 \(\alpha\) 的终边与单位圆交于点 \(P(x, y)\)。由单位圆方程 \(x^2 + y^2 = 1\)，代入 \(\sin \alpha = y\)，\(\cos \alpha = x\)，得：

\[\cos^2 \alpha + \sin^2 \alpha = x^2 + y^2 = 1\]

\(\blacksquare\)

\textbf{定理 19.4}（商数关系）对任意角 \(\alpha\)（\(\cos \alpha \neq 0\)）：

\[\tan \alpha = \frac{\sin \alpha}{\cos \alpha}\]

\textbf{证明}：\(\tan \alpha = \frac{y}{x} = \frac{\sin \alpha}{\cos \alpha}\)。\(\blacksquare\)

商数关系将正切函数归结为正弦与余弦之比，这是正切函数最本质的代数特征。从勾股关系和商数关系出发，通过简单的代数变形还可以得到正割和余割的衍生恒等式------这些恒等式在三角代换和积分计算中经常出现。

\textbf{定理 19.5}（正割关系）对任意角 \(\alpha\)（\(\cos \alpha \neq 0\)）：

\[1 + \tan^2 \alpha = \sec^2 \alpha\]

\textbf{证明}：由定理 19.3，\(\sin^2 \alpha + \cos^2 \alpha = 1\)。两边除以 \(\cos^2 \alpha\)：

\[\frac{\sin^2 \alpha}{\cos^2 \alpha} + 1 = \frac{1}{\cos^2 \alpha}\]

\[\tan^2 \alpha + 1 = \sec^2 \alpha\]

\(\blacksquare\)

正割关系与勾股关系一脉相承，仅通过同除以 \(\cos^2\alpha\) 即可导出。同理，对勾股关系两边除以 \(\sin^2\alpha\)，即可得到以余切和余割表达的对称形式。

\textbf{定理 19.6}（余割关系）对任意角 \(\alpha\)（\(\sin \alpha \neq 0\)）：

\[1 + \cot^2 \alpha = \csc^2 \alpha\]

\textbf{证明}：由定理 19.3，两边除以 \(\sin^2 \alpha\)：

\[1 + \frac{\cos^2 \alpha}{\sin^2 \alpha} = \frac{1}{\sin^2 \alpha}\]

\[1 + \cot^2 \alpha = \csc^2 \alpha\]

\(\blacksquare\)

\hrulefill

诱导公式是三角函数理论中最具系统性的部分。本节全面推导了六组诱导公式，涵盖了 \(\alpha + 2k\pi\)、\(\pi \pm \alpha\)、\(-\alpha\) 和 \(\frac{\pi}{2} \pm \alpha\) 等所有基本变换。诱导公式的核心思想是''化归''------将任意角的三角函数值转化为锐角（或第一象限角）的三角函数值来计算。这种化归利用了单位圆的对称性：关于 \(x\) 轴、\(y\) 轴、原点、直线 \(y = x\) 的对称变换分别对应不同的诱导公式。口诀''奇变偶不变，符号看象限''精炼地概括了 \(\frac{\pi}{2}\) 的整数倍变换的规律。诱导公式不仅是计算工具，更是理解三角函数对称性和周期性的窗口------它们揭示了正弦和余弦在坐标变换下的行为规律，为后续的三角恒等变换和傅里叶分析奠定了对称性基础。

诱导公式揭示了三角函数在角变换下的关系。以下假设 \(k \in \mathbb{Z}\)。

\subsubsection{\texorpdfstring{19.4.1 \(\alpha + 2k\pi\) 的三角函数}{19.4.1 \textbackslash alpha + 2k\textbackslash pi 的三角函数}}\label{alpha-2kpi-ux7684ux4e09ux89d2ux51fdux6570}

\textbf{定理 19.7} 对任意角 \(\alpha\) 和整数 \(k\)：

\[\sin(\alpha + 2k\pi) = \sin \alpha, \quad \cos(\alpha + 2k\pi) = \cos \alpha, \quad \tan(\alpha + 2k\pi) = \tan \alpha\]

\textbf{证明}：角 \(\alpha\) 和角 \(\alpha + 2k\pi\) 的终边相同（相差 \(2\pi\) 的整数倍），因此终边与单位圆的交点相同，故三角函数值相等。\(\blacksquare\)

\subsubsection{\texorpdfstring{19.4.2 \(\pi + \alpha\) 的三角函数}{19.4.2 \textbackslash pi + \textbackslash alpha 的三角函数}}\label{pi-alpha-ux7684ux4e09ux89d2ux51fdux6570}

\textbf{定理 19.8} 对任意角 \(\alpha\)：

\[\sin(\pi + \alpha) = -\sin \alpha, \quad \cos(\pi + \alpha) = -\cos \alpha, \quad \tan(\pi + \alpha) = \tan \alpha\]

\textbf{证明}：设角 \(\alpha\) 的终边与单位圆交于 \(P(x, y)\)。角 \(\pi + \alpha\) 的终边是将 \(OP\) 反向延长，交单位圆于 \(P'(-x, -y)\)。因此：

\[\sin(\pi + \alpha) = -y = -\sin \alpha\] \[\cos(\pi + \alpha) = -x = -\cos \alpha\] \[\tan(\pi + \alpha) = \frac{-y}{-x} = \frac{y}{x} = \tan \alpha\]

\(\blacksquare\)

\subsubsection{\texorpdfstring{19.4.3 \(-\alpha\) 的三角函数}{19.4.3 -\textbackslash alpha 的三角函数}}\label{alpha-ux7684ux4e09ux89d2ux51fdux6570}

\textbf{定理 19.9} 对任意角 \(\alpha\)：

\[\sin(-\alpha) = -\sin \alpha, \quad \cos(-\alpha) = \cos \alpha, \quad \tan(-\alpha) = -\tan \alpha\]

\textbf{证明}：设角 \(\alpha\) 的终边与单位圆交于 \(P(x, y)\)。角 \(-\alpha\) 的终边与单位圆交于 \(P'(x, -y)\)。因此：

\[\sin(-\alpha) = -y = -\sin \alpha, \quad \cos(-\alpha) = x = \cos \alpha, \quad \tan(-\alpha) = \frac{-y}{x} = -\tan \alpha\]

\(\blacksquare\)

\textbf{注} 此定理表明 \(\cos\) 是偶函数，\(\sin\) 和 \(\tan\) 是奇函数。

\subsubsection{\texorpdfstring{19.4.4 \(\pi - \alpha\) 的三角函数}{19.4.4 \textbackslash pi - \textbackslash alpha 的三角函数}}\label{pi---alpha-ux7684ux4e09ux89d2ux51fdux6570}

\textbf{定理 19.10} 对任意角 \(\alpha\)：

\[\sin(\pi - \alpha) = \sin \alpha, \quad \cos(\pi - \alpha) = -\cos \alpha, \quad \tan(\pi - \alpha) = -\tan \alpha\]

\textbf{证明}：利用定理 19.8 和 19.9：

\[\sin(\pi - \alpha) = \sin[\pi + (-\alpha)] = -\sin(-\alpha) = \sin \alpha\] \[\cos(\pi - \alpha) = \cos[\pi + (-\alpha)] = -\cos(-\alpha) = -\cos \alpha\]

\(\blacksquare\)

\subsubsection{\texorpdfstring{19.4.5 \(\frac{\pi}{2} - \alpha\) 的三角函数}{19.4.5 \textbackslash frac\{\textbackslash pi\}\{2\} - \textbackslash alpha 的三角函数}}\label{fracpi2---alpha-ux7684ux4e09ux89d2ux51fdux6570}

\textbf{定理 19.11} 对任意角 \(\alpha\)：

\[\sin\left(\frac{\pi}{2} - \alpha\right) = \cos \alpha, \quad \cos\left(\frac{\pi}{2} - \alpha\right) = \sin \alpha, \quad \tan\left(\frac{\pi}{2} - \alpha\right) = \cot \alpha\]

\textbf{证明}：设角 \(\alpha\) 的终边与单位圆交于 \(P(x, y)\)。角 \(\frac{\pi}{2} - \alpha\) 的终边与单位圆交于 \(P'\)。由于 \(\frac{\pi}{2} - \alpha\) 与 \(\alpha\) 关于直线 \(y = x\) 对称，\(P'\) 的坐标为 \((y, x)\)。因此：

\[\sin\left(\frac{\pi}{2} - \alpha\right) = x = \cos \alpha, \quad \cos\left(\frac{\pi}{2} - \alpha\right) = y = \sin \alpha\]

\(\blacksquare\)

\subsubsection{\texorpdfstring{19.4.6 \(\frac{\pi}{2} + \alpha\) 的三角函数}{19.4.6 \textbackslash frac\{\textbackslash pi\}\{2\} + \textbackslash alpha 的三角函数}}\label{fracpi2-alpha-ux7684ux4e09ux89d2ux51fdux6570}

\textbf{定理 19.12} 对任意角 \(\alpha\)：

\[\sin\left(\frac{\pi}{2} + \alpha\right) = \cos \alpha, \quad \cos\left(\frac{\pi}{2} + \alpha\right) = -\sin \alpha, \quad \tan\left(\frac{\pi}{2} + \alpha\right) = -\cot \alpha\]

\textbf{证明}：设 \(\beta = \frac{\pi}{2} - \alpha\)，则 \(\frac{\pi}{2} + \alpha = \pi - \beta\)。由定理 19.10 和 19.11：

\[\sin\left(\frac{\pi}{2} + \alpha\right) = \sin(\pi - \beta) = \sin \beta = \sin\left(\frac{\pi}{2} - \alpha\right) = \cos \alpha\]

\[\cos\left(\frac{\pi}{2} + \alpha\right) = \cos(\pi - \beta) = -\cos \beta = -\cos\left(\frac{\pi}{2} - \alpha\right) = -\sin \alpha\]

\(\blacksquare\)

以上六组诱导公式构成了一个完整的变换体系。为了便于记忆和应用，我们可以将其规律概括为一条简洁的口诀。

\textbf{诱导公式口诀}：``奇变偶不变，符号看象限''------将 \(\frac{\pi}{2}\) 的倍数记为 \(n \cdot \frac{\pi}{2}\)，若 \(n\) 为奇数，则 \(\sin\) 变 \(\cos\)、\(\cos\) 变 \(\sin\)；若 \(n\) 为偶数，则不变。符号由原函数在对应象限的符号决定（将 \(\alpha\) 视为锐角）。

\hrulefill

三角函数的图像和性质是理解三角函数行为特征的核心。本节系统研究了正弦函数 \(y = \sin x\)、余弦函数 \(y = \cos x\) 和正切函数 \(y = \tan x\) 的图像，并严格证明了它们的周期性、奇偶性、单调性和有界性。正弦函数的图像是一条光滑的波浪线（正弦曲线），其周期为 \(2\pi\)、振幅为 \(1\)，图像关于原点中心对称（奇函数）。余弦函数的图像是正弦曲线向左平移 \(\frac{\pi}{2}\) 的结果，因此是偶函数。正切函数的图像由无穷多个分支组成，每个分支在一个周期区间内从 \(-\infty\) 单调递增到 \(+\infty\)，且以 \(x = \frac{\pi}{2} + k\pi\) 为垂直渐近线。本节还引入了函数 \(y = A\sin(\omega x + \varphi) + B\) 的图像变换理论，通过振幅 \(A\)、角频率 \(\omega\)、初相 \(\varphi\) 和垂直位移 \(B\) 四个参数来描述正弦型函数的完整家族，这一理论在物理学（简谐振动、交流电）和工程学（信号处理）中具有基础性意义。

\textbf{定理 19.13}（正弦函数的性质）

\begin{enumerate}
\def\labelenumi{\arabic{enumi}.}
\tightlist
\item
  \textbf{定义域}：\(\mathbb{R}\)；\textbf{值域}：\([-1, 1]\)
\item
  \textbf{周期}：\(T = 2\pi\)（最小正周期）
\item
  \textbf{奇偶性}：奇函数
\item
  \textbf{单调性}：在 \(\left[-\frac{\pi}{2} + 2k\pi, \frac{\pi}{2} + 2k\pi\right]\) 上单调递增；在 \(\left[\frac{\pi}{2} + 2k\pi, \frac{3\pi}{2} + 2k\pi\right]\) 上单调递减
\end{enumerate}

\textbf{证明周期性}：\(\sin(x + 2\pi) = \sin x\)（定理 19.7）。下证 \(2\pi\) 是最小正周期。设 \(T\) 是正周期，则 \(\sin(x + T) = \sin x\) 对所有 \(x\) 成立。取 \(x = 0\)：\(\sin T = 0\)，故 \(T = k\pi\)。取 \(x = \frac{\pi}{2}\)：\(\sin\left(\frac{\pi}{2} + T\right) = 1\)，即 \(\cos T = 1\)，故 \(T = 2m\pi\)。因此最小正周期为 \(2\pi\)。

\textbf{证明单调性}：在单位圆上，角 \(\theta\) 对应的点为 \((\cos\theta, \sin\theta)\)，其中 \(\sin\theta\) 是该点的纵坐标。当 \(\theta\) 从 \(-\frac{\pi}{2}\) 增大到 \(\frac{\pi}{2}\) 时，对应点沿单位圆从最低点 \((0, -1)\) 经右半圆运动到最高点 \((0, 1)\)。在此过程中，纵坐标 \(\sin\theta\) 从 \(-1\) 单调递增到 \(1\)。

严格来说，设 \(-\frac{\pi}{2} \leq x_1 < x_2 \leq \frac{\pi}{2}\)。在单位圆上，\(x_1\) 对应点 \(P_1\)，\(x_2\) 对应点 \(P_2\)。由于 \(x_2 - x_1 \in (0, \pi]\)，点 \(P_2\) 在 \(P_1\) 的上方（纵坐标更大）。这由单位圆的几何性质直接得到：在右半圆（含上下端点）上，角度增大时纵坐标严格递增。\(\blacksquare\)

\subsubsection{\texorpdfstring{19.5.2 余弦函数 \(y = \cos x\)}{19.5.2 余弦函数 y = \textbackslash cos x}}\label{ux4f59ux5f26ux51fdux6570-y-cos-x}

\textbf{定理 19.14}（余弦函数的性质）

\begin{enumerate}
\def\labelenumi{\arabic{enumi}.}
\tightlist
\item
  \textbf{定义域}：\(\mathbb{R}\)；\textbf{值域}：\([-1, 1]\)
\item
  \textbf{周期}：\(T = 2\pi\)
\item
  \textbf{奇偶性}：偶函数
\item
  \textbf{单调性}：在 \([2k\pi, \pi + 2k\pi]\) 上单调递减；在 \([\pi + 2k\pi, 2\pi + 2k\pi]\) 上单调递增
\end{enumerate}

\textbf{注} 由诱导公式 \(\cos x = \sin\left(x + \frac{\pi}{2}\right)\)，余弦函数的图像可由正弦函数向左平移 \(\frac{\pi}{2}\) 得到。

\textbf{证明}：(1) 由 \(\cos(x+2\pi) = \cos x\) 知周期为 \(2\pi\)。(2) \(\cos(-x) = \cos x\) 为偶函数。(3) 由导数 \((\cos x)' = -\sin x\)，在 \([0,\pi]\) 上 \(\sin x \geq 0\) 故 \(\cos x\) 递减；在 \([\pi,2\pi]\) 上递增。\(\blacksquare\)

\subsubsection{\texorpdfstring{19.5.3 正切函数 \(y = \tan x\)}{19.5.3 正切函数 y = \textbackslash tan x}}\label{ux6b63ux5207ux51fdux6570-y-tan-x}

\textbf{定理 19.15}（正切函数的性质）

\begin{enumerate}
\def\labelenumi{\arabic{enumi}.}
\tightlist
\item
  \textbf{定义域}：\(\{x \mid x \neq \frac{\pi}{2} + k\pi, k \in \mathbb{Z}\}\)；\textbf{值域}：\(\mathbb{R}\)
\item
  \textbf{周期}：\(T = \pi\)（最小正周期）
\item
  \textbf{奇偶性}：奇函数
\item
  \textbf{单调性}：在每个开区间 \(\left(-\frac{\pi}{2} + k\pi, \frac{\pi}{2} + k\pi\right)\) 上严格递增
\end{enumerate}

\textbf{证明周期性}：\(\tan(x + \pi) = \frac{\sin(x + \pi)}{\cos(x + \pi)} = \frac{-\sin x}{-\cos x} = \tan x\)。\(\blacksquare\)

\subsubsection{\texorpdfstring{19.5.4 函数 \(y = A\sin(\omega x + \varphi) + B\) 的图像变换}{19.5.4 函数 y = A\textbackslash sin(\textbackslash omega x + \textbackslash varphi) + B 的图像变换}}\label{ux51fdux6570-y-asinomega-x-varphi-b-ux7684ux56feux50cfux53d8ux6362}

给定函数 \(y = A\sin(\omega x + \varphi) + B\)（\(A > 0\)，\(\omega > 0\)），其图像可由 \(y = \sin x\) 经以下变换得到：

\begin{enumerate}
\def\labelenumi{\arabic{enumi}.}
\tightlist
\item
  \textbf{相位变换}：\(y = \sin x \to y = \sin(x + \varphi)\)（左移 \(\varphi\) 个单位，\(\varphi > 0\) 时）
\item
  \textbf{周期变换}：\(y = \sin(x + \varphi) \to y = \sin(\omega x + \varphi)\)（横坐标缩短为 \(\frac{1}{\omega}\) 倍）
\item
  \textbf{振幅变换}：\(y = \sin(\omega x + \varphi) \to y = A\sin(\omega x + \varphi)\)（纵坐标扩大 \(A\) 倍）
\item
  \textbf{平移变换}：\(y = A\sin(\omega x + \varphi) \to y = A\sin(\omega x + \varphi) + B\)（上移 \(B\) 个单位）
\end{enumerate}

\textbf{性质}：振幅 \(A\)，周期 \(T = \frac{2\pi}{\omega}\)，频率 \(f = \frac{\omega}{2\pi}\)，最大值 \(A + B\)，最小值 \(-A + B\)。

\hrulefill

正弦型函数 \(y = A\sin(\omega x + \varphi) + B\) 是三角函数理论在应用中最常见的形式。本节对正弦型函数进行了完整分析，严格推导了振幅 \(A\)、周期 \(T = \frac{2\pi}{\omega}\)、频率 \(f = \frac{\omega}{2\pi}\)、初相 \(\varphi\) 以及最大值 \(A + B\) 和最小值 \(-A + B\) 的精确表达式。图像变换的两种路径（先平移后伸缩 vs 先伸缩后平移）被详细对比，揭示了变换顺序对平移量的影响：路径一中左移 \(\varphi\)，路径二中左移 \(\frac{\varphi}{\omega}\)。这一差异的本质在于，当对 \(x\) 进行伸缩变换 \(x \to \omega x\) 时，原来在 \(x\) 方向上的平移量也相应地被缩放。正弦型函数是物理学中简谐振动 \(x = A\cos(\omega t + \varphi)\) 的数学基础，也是傅里叶分析中基本构建块------任何周期函数都可以表示为不同频率的正弦型函数的叠加。

\textbf{定理 19.16} 设 \(A > 0\)，\(\omega > 0\)。函数 \(y = A\sin(\omega x + \varphi) + B\) 的性质如下：

\begin{enumerate}
\def\labelenumi{\arabic{enumi}.}
\tightlist
\item
  \textbf{振幅}：\(A\)（图像上最高点与最低点到中心线距离的一半）
\item
  \textbf{周期}：\(T = \frac{2\pi}{\omega}\)
\item
  \textbf{频率}：\(f = \frac{1}{T} = \frac{\omega}{2\pi}\)
\item
  \textbf{相位}：\(\varphi\)（初相）
\item
  \textbf{最大值}：\(A + B\)，\textbf{最小值}：\(-A + B\)
\item
  \textbf{中心线}：\(y = B\)
\end{enumerate}

\textbf{证明}：

\begin{enumerate}
\def\labelenumi{(\arabic{enumi})}
\setcounter{enumi}{1}
\tightlist
\item
  \(y(x + T) = A\sin(\omega(x + T) + \varphi) + B = A\sin(\omega x + 2\pi + \varphi) + B = A\sin(\omega x + \varphi) + B = y(x)\)。
\end{enumerate}

对于最小正周期：设 \(T'\) 是周期，则 \(\sin(\omega(x + T') + \varphi) = \sin(\omega x + \varphi)\) 对所有 \(x\) 成立，即 \(\omega T'\) 是 \(\sin\) 的周期，故 \(\omega T' = 2k\pi\)（\(k\) 为正整数），最小正周期 \(T' = \frac{2\pi}{\omega}\)。

\begin{enumerate}
\def\labelenumi{(\arabic{enumi})}
\setcounter{enumi}{4}
\tightlist
\item
  \(-1 \leq \sin(\omega x + \varphi) \leq 1\)，故 \(-A + B \leq y \leq A + B\)。\(\blacksquare\)
\end{enumerate}

\subsubsection{19.6.2 图像变换的严格描述}\label{ux56feux50cfux53d8ux6362ux7684ux4e25ux683cux63cfux8ff0}

\textbf{定理 19.17} 设 \(y = \sin x\) 的图像为 \(C_0\)，\(y = A\sin(\omega x + \varphi) + B\) 的图像为 \(C\)（\(A > 0\)，\(\omega > 0\)）。\(C\) 可由 \(C_0\) 经过以下变换序列得到（变换顺序可调）：

\textbf{路径一}（先平移后伸缩）：

\[C_0 \xrightarrow{\text{左移 } \varphi} C_1: y = \sin(x + \varphi) \xrightarrow{\text{横向压缩至 } \frac{1}{\omega}} C_2: y = \sin(\omega x + \varphi)\]

\[\xrightarrow{\text{纵向拉伸 } A \text{ 倍}} C_3: y = A\sin(\omega x + \varphi) \xrightarrow{\text{上移 } B} C\]

\textbf{路径二}（先伸缩后平移）：

\[C_0 \xrightarrow{\text{横向压缩至 } \frac{1}{\omega}} y = \sin(\omega x) \xrightarrow{\text{左移 } \frac{\varphi}{\omega}} y = \sin(\omega(x + \frac{\varphi}{\omega}))\]

\[\xrightarrow{\text{纵向拉伸 } A \text{ 倍}} y = A\sin(\omega x + \varphi) \xrightarrow{\text{上移 } B} C\]

\textbf{注意}：两种路径中，横向平移的量不同：路径一中平移 \(\varphi\)，路径二中平移 \(\frac{\varphi}{\omega}\)。这是因为横向伸缩会改变平移的效果。

\textbf{证明}：\(y = A\sin(\omega x + \varphi) + B\) 由 \(y = \sin x\) 经以下变换得到：(1) 横坐标缩放 \(1/\omega\) 倍（周期变为 \(2\pi/\omega\)）；(2) 向左平移 \(\varphi/\omega\)；(3) 纵坐标放大 \(A\) 倍（振幅变为 \(A\)）；(4) 向上平移 \(B\)。变换顺序不影响最终图像。\(\blacksquare\)

\hrulefill

三角函数的综合性质将多个三角函数概念融合在一起，形成了更深入的理解。本节证明了函数 \(y = a\sin x + b\cos x\) 的值域为 \([-\sqrt{a^2+b^2}, \sqrt{a^2+b^2}]\)，这一结果实质上是柯西-施瓦茨不等式在三角函数中的具体体现，也是后续辅助角公式 \(a\sin x + b\cos x = \sqrt{a^2+b^2}\sin(x + \varphi)\) 的理论基础。本节还证明了重要的三角不等式 \(|\sin x| \leq |x|\) 和 \(|\tan x| \geq |x|\)，这两个不等式在极限理论中至关重要------前者是证明 \(\lim_{x \to 0} \frac{\sin x}{x} = 1\) 的关键步骤，后者在分析正切函数的行为时提供了基本估计。三角函数的有界性（\(|\sin x| \leq 1\)，\(|\cos x| \leq 1\)）与非有界性（\(\tan x\) 和 \(\cot x\) 无界）的对比，清晰地展示了不同类型三角函数在''全局行为''上的本质差异。

\textbf{定理 19.18} 函数 \(y = a\sin x + b\cos x\) 的值域为 \([-\sqrt{a^2 + b^2}, \sqrt{a^2 + b^2}]\)。

\textbf{证明}：设 \(a, b\) 不全为零（否则结论显然）。令 \(R = \sqrt{a^2 + b^2} > 0\)。

\textbf{上界的证明}：由柯西--施瓦茨不等式（或直接计算），

\[(a\sin x + b\cos x)^2 \leq (a^2 + b^2)(\sin^2 x + \cos^2 x) = a^2 + b^2 = R^2\]

因此 \(|a\sin x + b\cos x| \leq R\)，即值域包含于 \([-R, R]\)。

\textbf{取等的证明}：取 \(x_0\) 使得 \(\sin x_0 = \frac{a}{R}\)，\(\cos x_0 = \frac{b}{R}\)（由 \(\left(\frac{a}{R}\right)^2 + \left(\frac{b}{R}\right)^2 = 1\) 可知这样的 \(x_0\) 存在），则

\[a\sin x_0 + b\cos x_0 = \frac{a^2}{R} + \frac{b^2}{R} = \frac{a^2 + b^2}{R} = R\]

类似地，取 \(x_1 = x_0 + \pi\) 可得 \(a\sin x_1 + b\cos x_1 = -R\)。因此值域恰好为 \([-R, R] = [-\sqrt{a^2 + b^2}, \sqrt{a^2 + b^2}]\)。\(\blacksquare\)

\textbf{注}：此证明完全自包含，不依赖后续章节的结果。Ch 27 中将给出此结论的另一种证明（利用辅助角公式）。

\subsubsection{19.7.2 三角不等式}\label{ux4e09ux89d2ux4e0dux7b49ux5f0f}

\textbf{定理 19.19} 对任意实数 \(x\)，\(|\sin x| \leq |x|\)。

\textbf{证明}：当 \(x = 0\) 时显然成立。当 \(x \neq 0\) 时，分两种情况：

当 \(0 < |x| < \frac{\pi}{2}\) 时，考虑单位圆中圆心角 \(|x|\) 对应的弧。设 \(P(\cos x, \sin x)\) 为单位圆上的点，\(P\) 到 \(x\) 轴的垂直距离为 \(|\sin x|\)，而弧长为 \(|x|\)。由于两点之间线段最短，垂线段长度严格小于弧长，故 \(|\sin x| < |x|\)（严格不等式）。

当 \(|x| \geq \frac{\pi}{2}\) 时，\(|\sin x| \leq 1 < \frac{\pi}{2} \leq |x|\)。

综合两种情况，\(|\sin x| \leq |x|\)。\(\blacksquare\)

\textbf{定理 19.20} 对任意实数 \(x\)，\(|\tan x| \geq |x|\)（在 \(\tan x\) 有定义时）。

\textbf{证明}：当 \(0 < |x| < \frac{\pi}{2}\) 时，考虑单位圆中圆心角 \(|x|\) 对应的弧和过点 \((1, 0)\) 的切线。设 \(P(\cos x, \sin x)\) 为单位圆上的点，切线与射线 \(\overrightarrow{OP}\) 的延长线交于点 \(T\)，则 \(|OT|\) 方向上的切线段长度为 \(|\tan x|\)（由直角三角形 \(O(1,0)T\) 可得 \(|OT| = |\tan x|\)）。弧长 \(|x|\) 严格小于切线段长 \(|\tan x|\)（因为弧是连接 \(P\) 附近两点的最短路径，而切线段是折线路径的一部分，可通过扇形面积比较严格论证：扇形面积 \(\frac{1}{2}|x| < \frac{1}{2}|\tan x|\) = 三角形面积）。故 \(|x| < |\tan x|\)。

当 \(|x| \geq \frac{\pi}{2}\) 时，\(|\tan x|\) 可以是无穷大或任意大（在不连续点附近），不等式也成立。\(\blacksquare\)

\subsubsection{19.7.3 三角函数的有界性}\label{ux4e09ux89d2ux51fdux6570ux7684ux6709ux754cux6027}

\textbf{定理 19.21} 正弦函数和余弦函数是\textbf{全局有界}的，即 \(|\sin x| \leq 1\)，\(|\cos x| \leq 1\)。正切函数和余切函数是\textbf{局部有界}的但在整体上无界。

\textbf{证明}：\(|\sin x| = |y| \leq \sqrt{x^2 + y^2} = 1\)（\(P(x,y)\) 在单位圆上）。类似地 \(|\cos x| \leq 1\)。

对于正切函数，\(\lim_{x \to \frac{\pi}{2}^-} \tan x = +\infty\)，故无界。\(\blacksquare\)

三角函数的正交性是正弦和余弦函数族最深刻的代数性质之一。本节提前引入了三角函数正交性的概念（尽管使用了定积分记号，但读者可在学习 Ch 31 后补充验证），证明了 \(\{\sin(nx), \cos(nx)\}_{n=1}^{\infty}\) 在 \([0, 2\pi]\) 上构成一个正交函数系。正交性意味着不同频率的三角函数的内积为零------这类似于几何中''互相垂直的向量''的概念，但推广到了函数空间。这一性质是傅里叶分析的基石：任何周期函数都可以唯一地表示为正弦和余弦函数的无穷级数（傅里叶级数），而正交性保证了展开系数的唯一性和可计算性。从物理学的角度看，正交性意味着不同频率的简谐振动是''互相独立''的------它们不会相互干扰，这一性质在信号处理、量子力学和热传导方程中都是核心原理。

\begin{quote}
\textbf{前向引用说明}：以下定理的陈述和证明使用了定积分记号 \(\int_a^b f(x)\,dx\)，该概念将在 Ch 31（积分初步）中严格定义。此处提前引入此结果是因为正交性是三角函数最本质的代数性质之一，也是傅里叶分析的基础。读者可先接受定理的结论，在学习 Ch 31 后再补充验证证明细节。
\end{quote}

\textbf{定理 19.22} 对于正整数 \(m, n\)：

\[\int_0^{2\pi} \sin(mx)\sin(nx) \, dx = \begin{cases} 0, & m \neq n \\ \pi, & m = n \end{cases}\]

\[\int_0^{2\pi} \cos(mx)\cos(nx) \, dx = \begin{cases} 0, & m \neq n \\ \pi, & m = n \end{cases}\]

\[\int_0^{2\pi} \sin(mx)\cos(nx) \, dx = 0\]

\textbf{证明}：当 \(m \neq n\) 时，由积化和差公式：

\[\sin(mx)\sin(nx) = \frac{1}{2}[\cos((m-n)x) - \cos((m+n)x)]\]

\[\int_0^{2\pi} \sin(mx)\sin(nx) \, dx = \frac{1}{2}\left[\frac{\sin((m-n)x)}{m-n} - \frac{\sin((m+n)x)}{m+n}\right]_0^{2\pi} = 0\]

当 \(m = n\) 时：

\[\sin^2(nx) = \frac{1 - \cos(2nx)}{2}\]

\[\int_0^{2\pi} \sin^2(nx) \, dx = \frac{1}{2}\left[x - \frac{\sin(2nx)}{2n}\right]_0^{2\pi} = \pi\]

其余公式类似证明。\(\blacksquare\)

\textbf{注}：此正交性是傅里叶级数理论的基础。

\hrulefill

三角函数在复数中的推广是数学统一性的绝佳例证。本节将复数 \(z = a + bi\) 表示为三角形式 \(z = r(\cos\theta + i\sin\theta)\)，并证明了复数乘法的几何意义------模相乘、辐角相加。这一结果直接导致了棣莫弗定理（De Moivre's Theorem）：\((\cos\theta + i\sin\theta)^n = \cos(n\theta) + i\sin(n\theta)\)，它简洁地给出了复数 \(n\) 次幂的公式，也自然导出了复数的 \(n\) 次方根公式。复数乘法与三角函数加法公式之间的深刻联系，最终由欧拉在1748年发现的欧拉公式 \(e^{i\theta} = \cos\theta + i\sin\theta\) 达到了顶峰------这一公式被誉为''数学中最美的公式''，因为它将指数函数和三角函数这两个看似无关的函数族统一在复数域中。欧拉公式不仅简化了三角恒等式的推导，更在傅里叶分析、微分方程和量子力学中发挥着核心作用。

\textbf{定义 19.7} 复数 \(z = a + bi\)（\(a, b \in \mathbb{R}\)）可以写成\textbf{三角形式}：

\[z = r(\cos\theta + i\sin\theta)\]

其中 \(r = |z| = \sqrt{a^2 + b^2}\) 是\textbf{模}，\(\theta = \arg z\) 是\textbf{辐角}。

\textbf{定理 19.23}（复数乘法的几何意义）设 \(z_1 = r_1(\cos\theta_1 + i\sin\theta_1)\)，\(z_2 = r_2(\cos\theta_2 + i\sin\theta_2)\)。则

\[z_1 z_2 = r_1 r_2[\cos(\theta_1 + \theta_2) + i\sin(\theta_1 + \theta_2)]\]

\textbf{证明}：

\[z_1 z_2 = r_1 r_2(\cos\theta_1 + i\sin\theta_1)(\cos\theta_2 + i\sin\theta_2)\]

\[= r_1 r_2[(\cos\theta_1\cos\theta_2 - \sin\theta_1\sin\theta_2) + i(\sin\theta_1\cos\theta_2 + \cos\theta_1\sin\theta_2)]\]

\[= r_1 r_2[\cos(\theta_1 + \theta_2) + i\sin(\theta_1 + \theta_2)]\]

\(\blacksquare\)

\textbf{定理 19.24}（复数的 \(n\) 次方根）复数 \(z = r(\cos\theta + i\sin\theta)\) 的 \(n\) 次方根有 \(n\) 个：

\[z_k = \sqrt[n]{r}\left(\cos\frac{\theta + 2k\pi}{n} + i\sin\frac{\theta + 2k\pi}{n}\right), \quad k = 0, 1, \ldots, n-1\]

\textbf{证明}：设 \(w = \rho(\cos\varphi + i\sin\varphi)\)，\(w^n = z\)。则 \(\rho^n = r\)，\(n\varphi = \theta + 2k\pi\)。\(\blacksquare\)

\hrulefill

\textbf{本章展望}：有了三角函数的定义，自然要问：两个角的和的三角函数如何表示？在下一章（Ch 27）中，我们将系统推导三角恒等变换公式，包括和差角公式、二倍角公式、半角公式、积化和差与和差化积公式，以及正弦定理和余弦定理。这些恒等变换是三角学的核心工具，也是后续学习微积分中三角函数导数的基础。

\hrulefill

\hrulefill

\hrulefill

\hrulefill

\hrulefill

\hrulefill

\newpage

\chapter{03-函数与微积分/01-函数与初等函数/04-数列.md}\label{ux51fdux6570ux4e0eux5faeux79efux520601-ux51fdux6570ux4e0eux521dux7b49ux51fdux657004-ux6570ux5217.md}

\textbf{前置知识}：本章依赖 Ch 23（函数概念，将数列视为定义在正整数集上的特殊函数）以及 Ch 25（指数运算，为等比数列的通项公式做准备）。本章为 Ch 29（极限与连续）提供离散化的极限直觉。

\hrulefill

数列是数学分析中最基本的研究对象之一。从定义在正整数集 \(\mathbb{N}^+\) 上的函数角度看，数列 \(a: \mathbb{N}^+ \to \mathbb{R}\) 是离散的------其自变量只能取 \(1, 2, 3, \ldots\)，这与连续函数有本质区别。但正是这种离散性使得数列成为理解极限概念的理想起点：\(\varepsilon\)-\(N\) 语言中 \(N\) 是正整数，比 \(\varepsilon\)-\(\delta\) 语言中的 \(\delta\)（正实数）更容易理解和操作。本章从数列的基本概念出发，系统研究等差数列、等比数列的通项公式与前 \(n\) 项和，然后引入递推数列和数学归纳法，最后建立数列极限的初步理论------包括极限的 \(\varepsilon\)-\(N\) 定义、收敛数列的性质、单调有界定理以及施图茨定理。本章的内容为后续的极限与连续（Ch 29）和微积分（Ch 30--31）提供了离散化的分析基础。

\section{第46章：等差数列与等比数列}\label{ux7b2c46ux7ae0ux7b49ux5deeux6570ux5217ux4e0eux7b49ux6bd4ux6570ux5217}

数列是定义在正整数集上的特殊函数 \(a: \mathbb{N}^+ \to \mathbb{R}\)，通常记为 \(\{a_n\}\)。与一般函数不同，数列的定义域是离散的，这使得数列具有独特的性质和研究方法。本节从数列的基本概念出发，定义了通项公式（将项数 \(n\) 映射到项值 \(a_n\) 的函数关系）和递推公式（相邻项之间的递归关系），并建立了前 \(n\) 项和 \(S_n\) 与通项 \(a_n\) 之间的互逆关系 \(a_n = S_n - S_{n-1}\)（\(n \geq 2\)）。这一关系极为重要------它提供了从和求通项或从通项求和的桥梁，在后续的级数理论中将被反复使用。数列的分类（有穷数列与无穷数列、递增数列与递减数列、有界数列与无界数列）为后续研究数列的极限行为提供了基本的分类框架。

\textbf{定义 21.1}（数列）按一定顺序排列的一列数 \(a_1, a_2, a_3, \ldots\) 称为\textbf{数列}，简记为 \(\{a_n\}\)。其中 \(a_n\) 称为数列的\textbf{第 \(n\) 项}（或\textbf{通项}），\(n\) 称为\textbf{项数}。

\textbf{定义 21.2}（通项公式）如果数列 \(\{a_n\}\) 的第 \(n\) 项 \(a_n\) 与项数 \(n\) 之间的关系可以用一个公式 \(a_n = f(n)\) 表示，则称此公式为\textbf{通项公式}。

\textbf{定义 21.3}（递推公式）如果数列 \(\{a_n\}\) 的项之间的关系可以用 \(a_{n+1} = g(a_n, a_{n-1}, \ldots)\) 表示，则称此关系为\textbf{递推公式}。

\textbf{定义 21.4}（数列的前 \(n\) 项和）数列 \(\{a_n\}\) 的前 \(n\) 项之和：

\[S_n = a_1 + a_2 + \cdots + a_n = \sum_{k=1}^{n} a_k\]

\textbf{定理 21.1}（\(a_n\) 与 \(S_n\) 的关系）

\[a_n = \begin{cases} S_1 & n = 1 \\ S_n - S_{n-1} & n \geq 2 \end{cases}\]

\textbf{证明}：当 \(n = 1\) 时，\(S_1 = a_1\)。当 \(n \geq 2\) 时，\(S_n = a_1 + a_2 + \cdots + a_n\)，\(S_{n-1} = a_1 + a_2 + \cdots + a_{n-1}\)，两式相减得 \(S_n - S_{n-1} = a_n\)。\(\blacksquare\)

\hrulefill

等差数列是最简单的非平凡数列，其本质特征是相邻两项的差为常数。本节从等差数列的定义出发，严格推导了通项公式 \(a_n = a_1 + (n-1)d\)，并证明了前 \(n\) 项和公式 \(S_n = \frac{n(a_1 + a_n)}{2} = \frac{n[2a_1 + (n-1)d]}{2}\)。前 \(n\) 项和公式的证明采用了''高斯求和法''（正序和倒序相加），这一巧妙的代数技巧展示了等差数列的对称性。本节还引入了等差中项的概念，并利用等差数列的通项公式证明了 \(a_n\) 是 \(a_{n-1}\) 和 \(a_{n+1}\) 的等差中项的充要条件。等差数列在实际应用中广泛存在------从最简单的匀速运动（位移每单位时间增加恒定值）到线性插值问题，等差数列都是最基本的数学模型。

\textbf{定义 21.5}（等差数列）如果数列 \(\{a_n\}\) 满足：对任意 \(n \geq 1\)，

\[a_{n+1} - a_n = d \quad (d \text{ 为常数})\]

则称 \(\{a_n\}\) 为\textbf{等差数列}，\(d\) 称为\textbf{公差}。

\textbf{定理 21.2}（等差数列通项公式）等差数列 \(\{a_n\}\) 的首项为 \(a_1\)，公差为 \(d\)，则：

\[a_n = a_1 + (n - 1)d\]

\textbf{证明}：对 \(n\) 进行归纳。

基础步：\(n = 1\) 时，\(a_1 = a_1 + 0 \cdot d\)，成立。

归纳步：假设 \(a_k = a_1 + (k-1)d\)。则 \(a_{k+1} = a_k + d = a_1 + (k-1)d + d = a_1 + kd\)。

由数学归纳法，结论成立。\(\blacksquare\)

\subsubsection{\texorpdfstring{21.2.2 前 \(n\) 项和}{21.2.2 前 n 项和}}\label{ux524d-n-ux9879ux548c}

\textbf{定理 21.3}（等差数列前 \(n\) 项和）等差数列 \(\{a_n\}\) 的前 \(n\) 项和为：

\[S_n = \frac{n(a_1 + a_n)}{2} = na_1 + \frac{n(n-1)}{2}d\]

\textbf{证明（倒序相加法）}：正序写出 \(S_n\)：

\[S_n = a_1 + a_2 + a_3 + \cdots + a_n\]

倒序写出 \(S_n\)：

\[S_n = a_n + a_{n-1} + a_{n-2} + \cdots + a_1\]

两式相加：

\[2S_n = (a_1 + a_n) + (a_2 + a_{n-1}) + \cdots + (a_n + a_1)\]

由于 \(a_k + a_{n+1-k} = [a_1 + (k-1)d] + [a_1 + (n-k)d] = 2a_1 + (n-1)d = a_1 + a_n\)（常数），共有 \(n\) 对，故：

\[2S_n = n(a_1 + a_n), \quad S_n = \frac{n(a_1 + a_n)}{2}\]

\(\blacksquare\)

\subsubsection{21.2.3 等差中项}\label{ux7b49ux5deeux4e2dux9879}

\textbf{定理 21.4} \(a, b, c\) 成等差数列的充要条件是 \(b = \frac{a + c}{2}\)。

\textbf{证明}：由等差数列的定义，\(a, b, c\) 成等差数列当且仅当相邻两项的差相等，即 \(b - a = c - b\)。将这一等式两边同时加上 \(a + b\)，得到 \(b - a + a + b = c - b + a + b\)，化简为 \(2b = a + c\)。两边除以 \(2\)，即得 \(b = \frac{a + c}{2}\)。反之，若 \(b = \frac{a + c}{2}\)，则 \(2b = a + c\)，移项得 \(b - a = c - b\)，即相邻两项之差相等，故 \(a, b, c\) 成等差数列。因此 \(b = \frac{a + c}{2}\) 是 \(a, b, c\) 成等差数列的充要条件。\(\blacksquare\)

\hrulefill

等比数列的特征是相邻两项的比为常数，这使得它在增长模式上与等差数列有着本质区别。本节推导了等比数列的通项公式 \(a_n = a_1 q^{n-1}\) 和前 \(n\) 项和公式 \(S_n = \frac{a_1(1 - q^n)}{1 - q}\)（\(q \neq 1\)）。前 \(n\) 项和公式的证明采用了''错位相减法''------将 \(S_n\) 和 \(qS_n\) 相减后，中间项全部抵消，仅剩首项和末项。这一方法的本质是构造了一个裂项和，展示了等比数列求和中的代数对称性。等比数列在金融数学（复利计算）、人口增长模型和计算科学（二分搜索、等比级数截断误差）中都有广泛应用。当 \(|q| < 1\) 时，等比数列的无穷级数 \(\sum_{n=0}^{\infty} aq^n = \frac{a}{1-q}\) 收敛，这是几何级数（等比级数）的核心结果，也是分析学中最基本也是最常用的级数之一。

\textbf{定义 21.6}（等比数列）如果数列 \(\{a_n\}\) 满足：对任意 \(n \geq 1\)，\(a_n \neq 0\)，

\[\frac{a_{n+1}}{a_n} = q \quad (q \text{ 为非零常数})\]

则称 \(\{a_n\}\) 为\textbf{等比数列}，\(q\) 称为\textbf{公比}。

\textbf{定理 21.5}（等比数列通项公式）等比数列 \(\{a_n\}\) 的首项为 \(a_1\)（\(a_1 \neq 0\)），公比为 \(q\)（\(q \neq 0\)），则：

\[a_n = a_1 q^{n-1}\]

\textbf{证明}：对 \(n\) 进行归纳。

基础步：\(a_1 = a_1 q^0\)，成立。

归纳步：假设 \(a_k = a_1 q^{k-1}\)。则 \(a_{k+1} = a_k \cdot q = a_1 q^{k-1} \cdot q = a_1 q^k\)。

由数学归纳法，结论成立。\(\blacksquare\)

\subsubsection{\texorpdfstring{21.3.2 前 \(n\) 项和}{21.3.2 前 n 项和}}\label{ux524d-n-ux9879ux548c-1}

\textbf{定理 21.6}（等比数列前 \(n\) 项和）等比数列 \(\{a_n\}\) 的首项为 \(a_1\)，公比为 \(q\)，则：

\[S_n = \begin{cases} na_1 & q = 1 \\ \dfrac{a_1(1 - q^n)}{1 - q} & q \neq 1 \end{cases}\]

\textbf{证明}：当 \(q = 1\) 时，\(a_n = a_1\)（常数列），\(S_n = na_1\)。

当 \(q \neq 1\) 时：

\[S_n = a_1 + a_1 q + a_1 q^2 + \cdots + a_1 q^{n-1} \quad \cdots (1)\]

\((1) \times q\)：

\[qS_n = a_1 q + a_1 q^2 + a_1 q^3 + \cdots + a_1 q^n \quad \cdots (2)\]

\((1) - (2)\)：

\[S_n - qS_n = a_1 - a_1 q^n\]

\[(1 - q)S_n = a_1(1 - q^n)\]

\[S_n = \frac{a_1(1 - q^n)}{1 - q}\]

\(\blacksquare\)

\subsubsection{21.3.3 等比中项}\label{ux7b49ux6bd4ux4e2dux9879}

\textbf{定理 21.7} \(a, b, c\) 成等比数列（各项非零）的充要条件是 \(b^2 = ac\)。

\textbf{证明}：由等比数列的定义（各项非零），\(a, b, c\) 成等比数列当且仅当相邻两项的比相等，即 \(\frac{b}{a} = \frac{c}{b}\)。由于 \(a, b, c\) 均非零，可以将等式两边同时乘以 \(ab\)（等价于交叉相乘），得到 \(b \cdot b = a \cdot c\)，即 \(b^2 = ac\)。反之，若 \(b^2 = ac\) 且 \(a, b, c\) 均非零，则两边除以 \(ab\) 得 \(\frac{b}{a} = \frac{c}{b}\)，即相邻两项之比相等，故 \(a, b, c\) 成等比数列。注意这里非零条件是必要的：若 \(b = 0\)，则 \(b^2 = ac\) 并不能推出 \(\frac{b}{a} = \frac{c}{b}\)（分母为零无意义）。因此 \(b^2 = ac\) 是 \(a, b, c\)（各项非零）成等比数列的充要条件。\(\blacksquare\)

\hrulefill

数列求和方法是数列理论中技巧性最强的部分。本节系统介绍了三种核心求和方法：裂项相消法（将通项分解为两项之差，使求和时中间项相消，得到望远镜和）、错位相减法（适用于等差数列与等比数列的乘积形式）和分组求和法（将数列拆分为几个已知求和公式的数列之和）。三种方法的共同思想是''转化''------将未知的复杂求和转化为已知的简单求和。本节还严格证明了平方和公式 \(\sum_{k=1}^{n} k^2 = \frac{n(n+1)(2n+1)}{6}\) 和立方和公式 \(\sum_{k=1}^{n} k^3 = \left[\frac{n(n+1)}{2}\right]^2\)。平方和公式的证明利用了 \((k+1)^3 - k^3 = 3k^2 + 3k + 1\) 的裂项技巧，这种从高次幂的差分解出低次幂和的方法在数学中具有普遍性，是后续积分学中''幂函数积分公式''的离散原型。

\textbf{方法}：将通项 \(a_n\) 分解为两项之差 \(a_n = f(n) - f(n+1)\)，使得求和时中间项相互抵消。

\textbf{定理 21.8} \(\sum_{k=1}^{n} \frac{1}{k(k+1)} = \frac{n}{n+1}\)。

\textbf{证明}：利用部分分式分解：\(\frac{1}{k(k+1)} = \frac{1}{k} - \frac{1}{k+1}\)。

\[\sum_{k=1}^{n} \frac{1}{k(k+1)} = \sum_{k=1}^{n}\left(\frac{1}{k} - \frac{1}{k+1}\right) = 1 - \frac{1}{n+1} = \frac{n}{n+1}\]

\(\blacksquare\)

\subsubsection{21.4.2 错位相减法}\label{ux9519ux4f4dux76f8ux51cfux6cd5}

\textbf{方法}：适用于形如 \(\sum a_n \cdot b_n\) 的求和，其中 \(\{a_n\}\) 是等差数列，\(\{b_n\}\) 是等比数列。

错位相减法（又称''错项相减法''）是求''等差乘等比型''数列前 \(n\) 项和的标准方法。其核心思想是：将和式 \(S\) 乘以等比数列的公比 \(q\)，然后将 \(qS\) 与 \(S\) 错位相减，使得中间项形成等比数列的部分和，从而将问题转化为简单的等比数列求和。具体步骤为：设 \(S_n = \sum_{k=1}^{n} a_k b_k\)，其中 \(a_k\) 为等差数列（公差 \(d\)），\(b_k\) 为等比数列（公比 \(q\)）。计算 \(qS_n\) 并将 \(S_n\) 与 \(qS_n\) 错开一位相减，得到 \((1-q)S_n = a_1 b_1 + d\sum_{k=2}^{n} b_k - a_n b_{n+1}\)。当 \(q \neq 1\) 时，利用等比数列求和公式即可求出 \(S_n\)。这一方法的巧妙之处在于，通过错位相减，等差数列的''等差''性质和等比数列的''等倍''性质相互抵消，使得求和变得可行。错位相减法在级数理论中也有重要推广------阿贝尔变换（Abel's summation by parts）正是这一思想的连续版本。

\subsubsection{21.4.3 常见求和公式}\label{ux5e38ux89c1ux6c42ux548cux516cux5f0f}

\textbf{定理 21.9} \(\sum_{k=1}^{n} k = \frac{n(n+1)}{2}\)。

\textbf{证明}：这是等差数列（首项 \(1\)，公差 \(1\)，末项 \(n\)）的前 \(n\) 项和，共有 \(n\) 项。由等差数列前 \(n\) 项和公式（定理 21.3），

\[\sum_{k=1}^{n} k = \frac{n(a_1 + a_n)}{2} = \frac{n(1 + n)}{2} = \frac{n(n+1)}{2}\]

这一公式也可以从倒序相加法直接得到：正序 \(S_n = 1 + 2 + \cdots + n\)，倒序 \(S_n = n + (n-1) + \cdots + 1\)，两式相加得 \(2S_n = (1+n) + (2+(n-1)) + \cdots + (n+1) = n(n+1)\)，故 \(S_n = \frac{n(n+1)}{2}\)。这一巧妙的代数运算（据说由高斯在小学时发现）是等差数列求和对称性的直接体现。\(\blacksquare\)

\textbf{定理 21.10}（平方和公式）\(\sum_{k=1}^{n} k^2 = \frac{n(n+1)(2n+1)}{6}\)。

\textbf{证明}：利用恒等式 \((k+1)^3 - k^3 = 3k^2 + 3k + 1\)，对 \(k = 1, 2, \ldots, n\) 求和：

\[\sum_{k=1}^{n}[(k+1)^3 - k^3] = 3\sum_{k=1}^{n} k^2 + 3\sum_{k=1}^{n} k + n\]

左边是裂项和：\((n+1)^3 - 1^3\)。

\[(n+1)^3 - 1 = 3\sum_{k=1}^{n} k^2 + \frac{3n(n+1)}{2} + n\]

\[\sum_{k=1}^{n} k^2 = \frac{(n+1)^3 - 1 - \frac{3n(n+1)}{2} - n}{3} = \frac{n(n+1)(2n+1)}{6}\]

\(\blacksquare\)

\textbf{定理 21.11}（立方和公式）\(\sum_{k=1}^{n} k^3 = \left[\frac{n(n+1)}{2}\right]^2\)。

\textbf{证明}：对 \(n\) 用数学归纳法。

基础步：\(n = 1\) 时，\(1^3 = 1 = \left[\frac{1 \cdot 2}{2}\right]^2\)，成立。

归纳步：假设 \(n = m\) 时成立。则

\[\sum_{k=1}^{m+1} k^3 = \left[\frac{m(m+1)}{2}\right]^2 + (m+1)^3 = (m+1)^2\left[\frac{m^2}{4} + (m+1)\right]\]

\[= (m+1)^2 \cdot \frac{m^2 + 4m + 4}{4} = (m+1)^2 \cdot \frac{(m+2)^2}{4} = \left[\frac{(m+1)(m+2)}{2}\right]^2\]

由数学归纳法，结论成立。\(\blacksquare\)

\hrulefill

\section{第47章：递推数列与数学归纳法}\label{ux7b2c47ux7ae0ux9012ux63a8ux6570ux5217ux4e0eux6570ux5b66ux5f52ux7eb3ux6cd5}

递推数列（又称递归数列）是由递推关系定义的数列，其每一项由前一项（或前几项）通过一个确定的规则生成。与等差数列和等比数列不同，递推数列的通项往往不能用一个简单的初等函数表示，但通过分析递推关系本身，我们可以揭示数列的许多深层性质------如极限行为、周期性、单调性等。一阶线性递推 \(a_{n+1} = pa_n + q\) 是最基本的递推类型，其通项可以通过构造辅助数列 \(\{a_n - \lambda\}\)（其中 \(\lambda = \frac{q}{1-p}\) 是不动点）转化为等比数列来求解。更一般的非线性递推 \(a_{n+1} = f(a_n)\) 的分析则需要借助不动点理论------如果 \(f\) 有不动点 \(x^*\) 且数列 \(\{a_n\}\) 收敛，则极限必然是 \(f\) 的不动点。本节从一阶线性递推出发，逐步引入不动点法和特征方程法，为后续的递推数列极限分析（21.9 节）打下基础。

\subsubsection{21.5.1 一阶线性递推数列}\label{ux4e00ux9636ux7ebfux6027ux9012ux63a8ux6570ux5217}

\textbf{定义 21.7}（一阶线性递推数列）形如 \(a_{n+1} = pa_n + q\)（\(p \neq 0\)）的递推关系称为\textbf{一阶线性递推}。

\textbf{定理 21.12} 设数列 \(\{a_n\}\) 满足 \(a_{n+1} = pa_n + q\)（\(p \neq 1\)），令 \(\lambda = \frac{q}{1-p}\)，则数列 \(\{a_n - \lambda\}\) 是以 \(a_1 - \lambda\) 为首项、\(p\) 为公比的等比数列。通项公式为：

\[a_n = \left(a_1 - \frac{q}{1-p}\right)p^{n-1} + \frac{q}{1-p}\]

\textbf{证明}：设 \(b_n = a_n - \lambda\)，则

\[b_{n+1} = a_{n+1} - \lambda = pa_n + q - \lambda = p(b_n + \lambda) + q - \lambda = pb_n + p\lambda + q - \lambda\]

由 \(\lambda = \frac{q}{1-p}\)，得 \(p\lambda + q = \lambda\)（即 \((1-p)\lambda = q\)），故 \(p\lambda + q - \lambda = 0\)。

因此 \(b_{n+1} = pb_n\)，\(\{b_n\}\) 是以 \(b_1 = a_1 - \lambda\) 为首项、\(p\) 为公比的等比数列。

\[b_n = (a_1 - \lambda)p^{n-1}\]

\[a_n = b_n + \lambda = (a_1 - \lambda)p^{n-1} + \lambda\]

\(\blacksquare\)

\hrulefill

\subsection{21.6 数学归纳法}\label{ux6570ux5b66ux5f52ux7eb3ux6cd5-1}

数学归纳法是证明与正整数有关的命题的最基本方法。其原理基于皮亚诺公理中的归纳公理：若 \(P(1)\) 成立且 \(P(k) \Rightarrow P(k+1)\) 对所有 \(k\) 成立，则 \(P(n)\) 对所有正整数 \(n\) 成立。数学归纳法的精髓在于将无穷多个命题的证明压缩为两步：基础步（验证 \(n = 1\)）和归纳步（建立 \(P(k)\) 到 \(P(k+1)\) 的推理桥梁）。注意区分''数学归纳法''和''归纳推理''------前者是严格的演绎证明方法，后者是自然科学中从特殊到一般的推理方式。数学归纳法有多种变体：强归纳法（假设 \(P(1), P(2), \ldots, P(k)\) 都成立来证明 \(P(k+1)\)）、倒推归纳法（从 \(n\) 到 \(n-1\)）、以及适用于实数集的''连续归纳法''。本节将系统介绍数学归纳法的原理、标准形式及其在数列通项公式证明、不等式证明和整除性证明中的典型应用。

\subsubsection{21.6.1 数学归纳法原理}\label{ux6570ux5b66ux5f52ux7eb3ux6cd5ux539fux7406}

\textbf{原理 21.1}（数学归纳法原理）设 \(P(n)\) 是关于正整数 \(n\) 的命题。如果：

\begin{enumerate}
\def\labelenumi{\arabic{enumi}.}
\tightlist
\item
  \textbf{基础步}：\(P(1)\) 成立；
\item
  \textbf{归纳步}：对任意正整数 \(k\)，由 \(P(k)\) 成立可推出 \(P(k+1)\) 成立；
\end{enumerate}

则 \(P(n)\) 对所有正整数 \(n\) 成立。

\textbf{公理化说明}：数学归纳法的正确性基于正整数集的良序性公理：任何非空的正整数集合都有最小元素。若 \(P(n)\) 不对所有正整数成立，则存在使 \(P(n)\) 不成立的最小正整数 \(n_0\)。由基础步 \(n_0 \neq 1\)，故 \(n_0 > 1\)。由 \(n_0\) 的最小性，\(P(n_0 - 1)\) 成立。但由归纳步，\(P(n_0)\) 也应成立，矛盾。

\subsubsection{21.6.2 强归纳法}\label{ux5f3aux5f52ux7eb3ux6cd5}

\textbf{原理 21.2}（强归纳法）设 \(P(n)\) 是关于正整数 \(n\) 的命题。如果对任意正整数 \(k\)，由 \(P(1), P(2), \ldots, P(k)\) 全部成立可推出 \(P(k+1)\) 成立，则 \(P(n)\) 对所有正整数 \(n\) 成立。

\textbf{注} 强归纳法与普通数学归纳法等价，但在某些问题中更为方便，因为归纳假设更强。

\hrulefill

\section{第48章：数列极限}\label{ux7b2c48ux7ae0ux6570ux5217ux6781ux9650}

数列的极限判定是分析学的基本问题------给定一个数列 \(\{a_n\}\)，我们需要判断它是否收敛，以及收敛于什么值。本节从收敛数列的基本性质出发，系统建立了极限判定的工具链。极限的保序性是极限理论中最常用的性质之一：它保证了极限运算与不等关系的相容性------如果两个数列从某项开始满足 \(a_n \geq b_n\)，则极限也满足 \(A \geq B\)（但严格不等式 \(a_n > b_n\) 取极限后可能退化等号，如 \(a_n = 1/n, b_n = -1/n\)）。极限的四则运算法则（定理 22.3）提供了将复杂数列的极限分解为简单数列极限之和、差、积、商的方法。夹逼定理（定理 22.4）是处理''被夹在中间''的数列的最有力工具------它不需要直接计算极限，只需要找到上下界数列。这些工具共同构成了数列极限判定问题的基本方法论，在 Ch 29（极限与连续）中将被推广到函数极限的情形。

\subsubsection{21.7.1 收敛数列的性质}\label{ux6536ux655bux6570ux5217ux7684ux6027ux8d28}

\textbf{定理 21.13}（极限的保序性）设 \(\lim_{n \to \infty} a_n = A\)，\(\lim_{n \to \infty} b_n = B\)。

\begin{enumerate}
\def\labelenumi{\arabic{enumi}.}
\tightlist
\item
  若 \(A > B\)，则存在 \(N\)，当 \(n > N\) 时，\(a_n > b_n\)。
\item
  若存在 \(N_0\)，当 \(n > N_0\) 时 \(a_n \geq b_n\)，则 \(A \geq B\)。
\end{enumerate}

\textbf{证明}：

\begin{enumerate}
\def\labelenumi{(\arabic{enumi})}
\tightlist
\item
  取 \(\varepsilon = \frac{A - B}{2} > 0\)。存在 \(N_1, N_2\)，当 \(n > N_1\) 时 \(|a_n - A| < \varepsilon\)，当 \(n > N_2\) 时 \(|b_n - B| < \varepsilon\)。取 \(N = \max\{N_1, N_2\}\)，当 \(n > N\) 时：
\end{enumerate}

\[a_n > A - \varepsilon = \frac{A + B}{2} = B + \varepsilon > b_n\]

\begin{enumerate}
\def\labelenumi{(\arabic{enumi})}
\setcounter{enumi}{1}
\tightlist
\item
  用反证法。若 \(A < B\)，由 (1) 存在 \(N\)，当 \(n > N\) 时 \(a_n < b_n\)，与 \(a_n \geq b_n\)（\(n > N_0\)）矛盾。\(\blacksquare\)
\end{enumerate}

\subsubsection{21.7.2 子列与收敛}\label{ux5b50ux5217ux4e0eux6536ux655b}

\textbf{定义 21.8}（子列）设 \(\{a_n\}\) 是数列，\(n_1 < n_2 < n_3 < \cdots\) 是严格递增的正整数列，则 \(\{a_{n_k}\}\) 称为 \(\{a_n\}\) 的一个\textbf{子列}。

\textbf{定理 21.14} 若 \(\lim_{n \to \infty} a_n = L\)，则 \(\{a_n\}\) 的任何子列都收敛于 \(L\)。

\textbf{证明}：设 \(\{a_n\}\) 收敛于 \(L\)，即对任意 \(\varepsilon > 0\)，存在 \(N \in \mathbb{N}\)，使得当 \(n > N\) 时 \(|a_n - L| < \varepsilon\)。现在考虑子列 \(\{a_{n_k}\}\)，其中 \(n_1 < n_2 < \cdots\) 是严格递增的整数序列。我们需要证明子列也收敛于 \(L\)。由于 \(n_k \geq k\)（这是严格递增整数序列的基本性质：\(n_1 \geq 1\)，\(n_2 \geq 2\)，依此类推），当 \(k > N\) 时，\(n_k \geq k > N\)，因此 \(|a_{n_k} - L| < \varepsilon\)。按极限定义，这意味着 \(\{a_{n_k}\}\) 收敛于 \(L\)。这一结论的直观意义是：如果数列整体趋于某个极限，那么从中任意抽取的子列也都趋于同一极限------子列继承了母列的收敛行为。\(\blacksquare\)

\textbf{推论 21.1} 若 \(\{a_n\}\) 有两个子列收敛于不同的极限，则 \(\{a_n\}\) 发散。

\textbf{证明}：若 \(\{a_n\}\) 收敛于 \(L\)，则所有子列都收敛于 \(L\)（定理 21.14），矛盾。\(\blacksquare\)

\hrulefill

\subsection{\texorpdfstring{21.8 数列极限的 \(\varepsilon\)-\(N\) 论证技巧}{21.8 数列极限的 \textbackslash varepsilon-N 论证技巧}}\label{ux6570ux5217ux6781ux9650ux7684-varepsilon-n-ux8bbaux8bc1ux6280ux5de7}

\(\varepsilon\)-\(N\) 语言是数列极限的严格定义语言，掌握其论证技巧是分析学入门的关键。本节系统整理了 \(\varepsilon\)-\(N\) 证明中的常用策略和技巧。放缩法是其中最核心的方法------由于我们只需要证明 \(|a_n - L| < \varepsilon\) 对于足够大的 \(n\) 成立，而不需要精确估计 \(|a_n - L|\) 的值，因此可以将 \(|a_n - L|\) 放大到一个更简单的表达式 \(M_n\)，只要 \(M_n \to 0\) 即可。伯努利不等式 \((1+x)^n \geq 1+nx\)（\(x > -1\)）是放缩法中最常用的代数工具之一，它使得我们可以将指数形式的不等式转化为线性不等式。二项式定理展开也是有力的放缩工具------例如，\((1+h)^n = 1 + nh + \binom{n}{2}h^2 + \cdots \geq \binom{n}{2}h^2\)，当 \(h > 0\) 时提供了比伯努利不等式更强的下界估计。本节还将介绍如何处理带根号、带分数和带阶乘的数列极限，以及如何利用已知极限（如 \(\lim 1/n = 0\)）来构造 \(\varepsilon\)-\(N\) 证明。

\subsubsection{21.8.1 常用放缩技巧}\label{ux5e38ux7528ux653eux7f29ux6280ux5de7}

\textbf{引理 21.1} 设 \(a_n \geq 0\)，\(b_n \geq 0\)，\(a_n \leq b_n\)。若 \(\lim_{n \to \infty} b_n = 0\)，则 \(\lim_{n \to \infty} a_n = 0\)。

\textbf{证明}：由条件 \(a_n \geq 0\)，\(b_n \geq 0\)，\(a_n \leq b_n\) 对所有 \(n\) 成立，且 \(\lim_{n \to \infty} b_n = 0\)。对任意 \(\varepsilon > 0\)，由 \(\lim_{n \to \infty} b_n = 0\)，存在 \(N \in \mathbb{N}\)，使得当 \(n > N\) 时 \(|b_n - 0| < \varepsilon\)，即 \(b_n < \varepsilon\)（因为 \(b_n \geq 0\)）。由于 \(0 \leq a_n \leq b_n\)，当 \(n > N\) 时也有 \(0 \leq a_n < \varepsilon\)，即 \(|a_n - 0| < \varepsilon\)。由极限定义，\(\lim_{n \to \infty} a_n = 0\)。注意这里使用了 \(a_n \geq 0\) 的条件，这使得 \(|a_n - 0| = a_n\)，从而不等式 \(a_n \leq b_n\) 可以直接转化为 \(|a_n - 0| \leq |b_n - 0|\)。对于一般的夹逼定理（不要求非负），需要两侧都有控制条件。\(\blacksquare\)

\textbf{引理 21.2}（伯努利不等式）设 \(x > -1\)，\(n\) 为正整数，则 \((1 + x)^n \geq 1 + nx\)。

\textbf{证明}：对 \(n\) 用数学归纳法。

\(n = 1\)：\(1 + x \geq 1 + x\)，成立。

\(n \Rightarrow n+1\)：假设 \((1+x)^n \geq 1 + nx\)。由于 \(1 + x > 0\)（即 \(x > -1\)），

\[(1+x)^{n+1} = (1+x)^n(1+x) \geq (1+nx)(1+x) = 1 + (n+1)x + nx^2 \geq 1 + (n+1)x\]

\(\blacksquare\)

\subsubsection{21.8.2 重要数列极限}\label{ux91cdux8981ux6570ux5217ux6781ux9650}

\textbf{定理 21.15} \(\lim_{n \to \infty} \sqrt[n]{n} = 1\)。

\textbf{证明}：令 \(\sqrt[n]{n} = 1 + h_n\)（\(h_n > 0\)）。则 \(n = (1 + h_n)^n\)。由伯努利不等式：

\[n = (1 + h_n)^n \geq 1 + nh_n\]

所以 \(h_n \leq \frac{n - 1}{n} < 1\)，且 \(0 < h_n \leq \frac{1}{n} \cdot (n-1)\)。

进一步，对 \(n \geq 2\)，由二项式定理：

\[n = (1 + h_n)^n \geq \binom{n}{2}h_n^2 = \frac{n(n-1)}{2}h_n^2\]

所以 \(h_n^2 \leq \frac{2}{n-1}\)，\(h_n \leq \sqrt{\frac{2}{n-1}} \to 0\)。

由夹逼定理，\(h_n \to 0\)，故 \(\sqrt[n]{n} = 1 + h_n \to 1\)。\(\blacksquare\)

\textbf{定理 21.16} \(\lim_{n \to \infty} \sqrt[n]{a} = 1\)（\(a > 0\)）。

\textbf{证明}：当 \(a = 1\) 时显然。当 \(a > 1\) 时，令 \(\sqrt[n]{a} = 1 + h_n\)（\(h_n > 0\)）。

\(a = (1 + h_n)^n \geq 1 + nh_n\)，故 \(h_n \leq \frac{a - 1}{n} \to 0\)。

由夹逼定理，\(\sqrt[n]{a} \to 1\)。当 \(0 < a < 1\) 时，\(\sqrt[n]{a} = \frac{1}{\sqrt[n]{1/a}} \to \frac{1}{1} = 1\)。\(\blacksquare\)

\hrulefill

\section{第49章：递推数列与极限深化}\label{ux7b2c49ux7ae0ux9012ux63a8ux6570ux5217ux4e0eux6781ux9650ux6df1ux5316}

在 21.5 节建立了递推数列的基本概念之后，本节深入分析递推数列的极限行为。不动点法（Fixed Point Method）是分析递推数列 \(a_{n+1} = f(a_n)\) 的核心工具：如果 \(f\) 连续且数列 \(\{a_n\}\) 收敛于 \(L\)，则对递推式两边取极限得 \(L = f(L)\)，即 \(L\) 是 \(f\) 的不动点。不动点法的关键在于：它不需要显式解出通项，只需分析 \(f\) 的不动点即可判断数列的可能极限。然而，不动点法只能给出极限的''候选值''，要判定数列是否确实收敛，还需结合单调有界定理或压缩映射原理。对于形如 \(a_{n+1} = \frac{pa_n + q}{ra_n + s}\) 的分式线性递推，可以通过变量代换 \(\frac{1}{a_n - \alpha}\)（其中 \(\alpha\) 是不动点）将其转化为一阶线性递推。特征方程法（Characteristic Equation Method）则适用于二阶线性递推 \(a_{n+2} = pa_{n+1} + qa_n\)，其解的形式由特征方程 \(r^2 - pr - q = 0\) 的根决定------这与线性微分方程的特征方程法如出一辙，体现了离散与连续之间的深刻类比。

\subsubsection{21.9.1 不动点法}\label{ux4e0dux52a8ux70b9ux6cd5}

\textbf{定义 21.9}（不动点）设 \(f: D \to D\)，若 \(x^*\) 满足 \(f(x^*) = x^*\)，则称 \(x^*\) 为 \(f\) 的\textbf{不动点}。

\textbf{定理 21.17} 设数列 \(\{a_n\}\) 由 \(a_{n+1} = f(a_n)\) 定义，\(x^*\) 是 \(f\) 的不动点。若 \(\{a_n\}\) 收敛于 \(L\)，且 \(f\) 连续，则 \(L = x^*\)。

\textbf{证明}：设 \(\lim_{n \to \infty} a_n = L\) 且 \(L\) 在 \(f\) 的定义域内，\(f\) 在 \(L\) 处连续。由连续性的定义（定理 22.1），\(\lim_{x \to L} f(x) = f(L)\)。由函数极限与数列极限的关系（定理 22.10），对于任意收敛于 \(L\) 的数列 \(\{a_n\}\)，有 \(\lim_{n \to \infty} f(a_n) = f(L)\)。现在对递归关系 \(a_{n+1} = f(a_n)\) 两边取极限（\(n \to \infty\)）。左边 \(\lim_{n \to \infty} a_{n+1} = L\)（因为 \(\{a_{n+1}\}\) 是 \(\{a_n\}\) 的子列，由定理 21.14 收敛于同一极限 \(L\)）。右边 \(\lim_{n \to \infty} f(a_n) = f(L)\)（由 \(f\) 的连续性）。因此 \(L = f(L)\)，即 \(L\) 是 \(f\) 的不动点。\(\blacksquare\)

\textbf{注}：此定理仅用于\textbf{已知}数列收敛时确定其极限值，本身不保证收敛性。实际应用中，需先验证收敛条件。一个常用的收敛性充分条件是：若 \(f\) 在不动点 \(x^*\) 处可导且 \(|f'(x^*)| < 1\)，则存在 \(x^*\) 的某个邻域，使得从该邻域内出发的迭代 \(a_{n+1} = f(a_n)\) 收敛于 \(x^*\)（压缩映射原理的特殊情形）。

\subsubsection{21.9.2 斐波那契数列}\label{ux6590ux6ce2ux90a3ux5951ux6570ux5217}

\textbf{定义 21.10}（斐波那契数列）满足递推关系

\[F_1 = 1, \quad F_2 = 1, \quad F_{n+2} = F_{n+1} + F_n \quad (n \geq 1)\]

的数列 \(\{F_n\}\) 称为\textbf{斐波那契数列}。

\textbf{定理 21.18}（比内公式）斐波那契数列的通项公式为：

\[F_n = \frac{1}{\sqrt{5}}\left[\left(\frac{1 + \sqrt{5}}{2}\right)^n - \left(\frac{1 - \sqrt{5}}{2}\right)^n\right]\]

\textbf{证明}：特征方程为 \(x^2 = x + 1\)，即 \(x^2 - x - 1 = 0\)。两根为

\[\varphi = \frac{1 + \sqrt{5}}{2}, \quad \psi = \frac{1 - \sqrt{5}}{2}\]

设 \(F_n = A\varphi^n + B\psi^n\)。由初始条件 \(F_1 = 1\)，\(F_2 = 1\)：

\[A\varphi + B\psi = 1, \quad A\varphi^2 + B\psi^2 = 1\]

利用 \(\varphi^2 = \varphi + 1\)，\(\psi^2 = \psi + 1\)，第二个方程变为 \(A(\varphi + 1) + B(\psi + 1) = 1\)，即 \(A\varphi + B\psi + A + B = 1\)，代入第一个方程得 \(A + B = 0\)，即 \(B = -A\)。

由 \(A(\varphi - \psi) = 1\)，得 \(A = \frac{1}{\varphi - \psi} = \frac{1}{\sqrt{5}}\)。

因此 \(F_n = \frac{\varphi^n - \psi^n}{\sqrt{5}}\)。

验证递推：\(F_{n+1} + F_n = \frac{\varphi^{n+1} - \psi^{n+1} + \varphi^n - \psi^n}{\sqrt{5}} = \frac{\varphi^n(\varphi + 1) - \psi^n(\psi + 1)}{\sqrt{5}} = \frac{\varphi^n \cdot \varphi^2 - \psi^n \cdot \psi^2}{\sqrt{5}} = \frac{\varphi^{n+2} - \psi^{n+2}}{\sqrt{5}} = F_{n+2}\)。\(\blacksquare\)

\textbf{推论 21.2} \(\lim_{n \to \infty} \frac{F_{n+1}}{F_n} = \varphi = \frac{1 + \sqrt{5}}{2} \approx 1.618...\)

\textbf{证明}：\(\frac{F_{n+1}}{F_n} = \frac{\varphi^{n+1} - \psi^{n+1}}{\varphi^n - \psi^n} = \frac{\varphi - \psi\left(\frac{\psi}{\varphi}\right)^n}{1 - \left(\frac{\psi}{\varphi}\right)^n}\)。

由于 \(|\psi/\varphi| = \frac{\sqrt{5}-1}{\sqrt{5}+1} < 1\)，\(\left(\frac{\psi}{\varphi}\right)^n \to 0\)，故极限为 \(\varphi\)。\(\blacksquare\)

\textbf{注}：\(\varphi = \frac{1 + \sqrt{5}}{2}\) 称为\textbf{黄金分割比}。

\subsection{21.10 数列的极限存在性的进一步判定}\label{ux6570ux5217ux7684ux6781ux9650ux5b58ux5728ux6027ux7684ux8fdbux4e00ux6b65ux5224ux5b9a}

在 21.7 节建立了基本的极限判定方法之后，本节引入更高级的判定工具。施图茨定理（Stolz-Cesaro Theorem）是数列极限理论中的''洛必达法则''------它处理 \(\frac{\infty}{\infty}\) 型不定式，将 \(\frac{a_n}{b_n}\) 的极限转化为 \(\frac{a_n - a_{n-1}}{b_n - b_{n-1}}\) 的极限。施图茨定理在求 \(\frac{n}{\sqrt[n]{n!}}\) 或 \(\frac{1^k + 2^k + \cdots + n^k}{n^{k+1}}\) 等形式的极限时特别有用。柯西收敛准则（Cauchy Convergence Criterion）是另一个强有力的工具：\(\{a_n\}\) 收敛当且仅当它是柯西列（即 \(\forall \varepsilon > 0, \exists N, \forall m, n > N: |a_m - a_n| < \varepsilon\)）。柯西准则的独特之处在于------它不需要知道极限的具体值，只需要数列的项彼此足够接近，这与实数的完备性（确界存在定理）等价。柯西准则在理论分析中尤为重要，因为它为判断级数收敛提供了基本框架。本节还将介绍压缩映射原理（若 \(|a_{n+1} - a_{n+2}| \leq k|a_n - a_{n+1}|\) 且 \(0 < k < 1\)，则数列收敛），这是不动点迭代法收敛性的理论基础。

\subsubsection{21.10.1 施图茨定理}\label{ux65bdux56feux8328ux5b9aux7406}

\textbf{定理 21.19}（施图茨定理）设 \(\{b_n\}\) 是严格递增的正数列，\(\lim_{n \to \infty} b_n = +\infty\)。若 \(\lim_{n \to \infty} \frac{a_n - a_{n-1}}{b_n - b_{n-1}} = L\)（\(L\) 可以是 \(\pm\infty\)），则

\[\lim_{n \to \infty} \frac{a_n}{b_n} = L\]

\textbf{证明}：这是洛必达法则的离散版本。证明思路与洛必达法则类似，利用柯西中值定理的离散形式。

对于 \(\varepsilon > 0\)，存在 \(N_0\)，当 \(n > N_0\) 时 \(\left|\frac{a_n - a_{n-1}}{b_n - b_{n-1}} - L\right| < \varepsilon\)。

对 \(n > N_0\)，将 \(\frac{a_n - a_{N_0}}{b_n - b_{N_0}}\) 看作加权平均：

\[\frac{a_n - a_{N_0}}{b_n - b_{N_0}} = \frac{\sum_{k=N_0+1}^{n}(a_k - a_{k-1})}{\sum_{k=N_0+1}^{n}(b_k - b_{k-1})}\]

每一项 \(\frac{a_k - a_{k-1}}{b_k - b_{k-1}}\) 都在 \((L - \varepsilon, L + \varepsilon)\) 内，因此加权平均也在 \((L - \varepsilon, L + \varepsilon)\) 内（因为 \(b_k - b_{k-1} > 0\)）。

由此可证 \(\frac{a_n}{b_n} \to L\)。\(\blacksquare\)

\subsubsection{21.10.2 柯西收敛准则}\label{ux67efux897fux6536ux655bux51c6ux5219}

\textbf{定理 21.20}（柯西收敛准则）数列 \(\{a_n\}\) 收敛的充要条件是：

\[\forall \varepsilon > 0, \exists N, \forall m, n > N: |a_m - a_n| < \varepsilon\]

满足此条件的数列称为\textbf{柯西列}（或\textbf{基本列}）。

*\textbf{证明}：（必要性）设 \(\{a_n\}\) 收敛于 \(L\)。对任意 \(\varepsilon > 0\)，由收敛定义，存在 \(N \in \mathbb{N}\)，使得当 \(n > N\) 时 \(|a_n - L| < \frac{\varepsilon}{2}\)。那么当 \(m, n > N\) 时，由三角不等式

\[|a_m - a_n| = |(a_m - L) - (a_n - L)| \leq |a_m - L| + |a_n - L| < \frac{\varepsilon}{2} + \frac{\varepsilon}{2} = \varepsilon\]

因此 \(\{a_n\}\) 是柯西列。这里的关键是三角不等式将''两项之差''转化为''两项各自与极限之差''的和，从而将收敛条件转化为柯西条件。

（充分性）设 \(\{a_n\}\) 是柯西列。首先证明柯西列有界：取 \(\varepsilon = 1\)，存在 \(N\)，使得 \(m, n > N\) 时 \(|a_m - a_n| < 1\)。特别地，取 \(m = N+1\)，则对所有 \(n > N\) 有 \(|a_n - a_{N+1}| < 1\)，即 \(|a_n| < |a_{N+1}| + 1\)。因此整个数列被 \(M = \max\{|a_1|, \ldots, |a_N|, |a_{N+1}|+1\}\) 所界定。由波尔查诺-魏尔斯特拉斯定理，有界数列必有收敛子列。设 \(\{a_{n_k}\}\) 是 \(\{a_n\}\) 的收敛子列，极限为 \(L\)。对柯西列条件取 \(m = n_k\)，当 \(n, n_k > N\) 时 \(|a_n - a_{n_k}| < \frac{\varepsilon}{2}\)。令 \(k \to \infty\)，由绝对值函数的连续性，\(|a_n - L| = \lim_{k \to \infty} |a_n - a_{n_k}| \leq \frac{\varepsilon}{2} < \varepsilon\)。因此 \(\{a_n\}\) 收敛于 \(L\)。\(\blacksquare\)

\subsection{21.11 数学归纳法的典型应用}\label{ux6570ux5b66ux5f52ux7eb3ux6cd5ux7684ux5178ux578bux5e94ux7528}

在 21.6 节建立了数学归纳法的基本原理之后，本节通过典型应用展示归纳法的强大威力。归纳法在证明不等式时尤为有效------伯努利不等式 \((1+x)^n \geq 1+nx\)（\(x > -1\)）的证明是归纳法最经典的范例之一，其归纳步仅需展开 \((1+x)^{k+1} = (1+x)^k(1+x) \geq (1+kx)(1+x) = 1 + (k+1)x + kx^2 \geq 1 + (k+1)x\)。均值不等式（算术-几何平均不等式）的归纳证明则展示了''向前-向后归纳法''（Cauchy's forward-backward induction）的巧妙之处------先证明 \(n = 2\) 的情形，再证明 \(n = 2^k\) 的情形，最后用''回退''步骤证明从 \(n\) 到 \(n-1\)。归纳法在整除性证明中也扮演重要角色------例如 \(n^3 + (n+1)^3 + (n+2)^3\) 能被 \(9\) 整除的证明，归纳步中可利用 \(n^3 + (n+1)^3 + (n+2)^3\) 与 \(n+1, n+2, n+3\) 项的关系。本节还将展示归纳法在组合恒等式证明（如 \(\sum_{k=0}^{n} \binom{n}{k} = 2^n\)）和斐波那契数列性质证明中的应用，体现归纳法作为''离散数学的万能工具''的广泛适用性。

\subsubsection{21.11.1 用归纳法证明不等式}\label{ux7528ux5f52ux7eb3ux6cd5ux8bc1ux660eux4e0dux7b49ux5f0f}

\textbf{定理 21.21}（伯努利不等式的推广）设 \(x_1, x_2, \ldots, x_n > -1\) 且同号，则

\[(1 + x_1)(1 + x_2)\cdots(1 + x_n) \geq 1 + x_1 + x_2 + \cdots + x_n\]

\textbf{证明}：对 \(n\) 用数学归纳法。

\(n = 1\) 时显然成立。

\(n \Rightarrow n+1\)：假设对 \(n\) 成立。设 \(x_1, x_2, \ldots, x_{n+1}\) 同号且大于 \(-1\)。

\[(1+x_1)\cdots(1+x_n)(1+x_{n+1}) \geq (1 + x_1 + \cdots + x_n)(1 + x_{n+1})\]

\[= 1 + x_1 + \cdots + x_{n+1} + (x_1 + \cdots + x_n)x_{n+1}\]

由于 \(x_1, \ldots, x_n, x_{n+1}\) 同号，\((x_1 + \cdots + x_n)x_{n+1} \geq 0\)，故

\[\geq 1 + x_1 + \cdots + x_{n+1}\]

\(\blacksquare\)

\subsubsection{21.11.2 用归纳法证明整除性}\label{ux7528ux5f52ux7eb3ux6cd5ux8bc1ux660eux6574ux9664ux6027}

\textbf{定理 21.22} 对任意正整数 \(n\)，\(6 | n(n+1)(2n+1)\)。

\textbf{证明}：\(n(n+1)(2n+1) = 6\sum_{k=1}^{n}k^2\)，故 \(6 | n(n+1)(2n+1)\)。

或者直接用归纳法：

\(n = 1\)：\(1 \cdot 2 \cdot 3 = 6\)，\(6 | 6\)。

\(n \Rightarrow n+1\)：\((n+1)(n+2)(2n+3) - n(n+1)(2n+1) = (n+1)[(n+2)(2n+3) - n(2n+1)] = (n+1)(6n+6) = 6(n+1)^2\)。

由归纳假设 \(6 | n(n+1)(2n+1)\)，又 \(6 | 6(n+1)^2\)，故 \(6 | (n+1)(n+2)(2n+3)\)。\(\blacksquare\)

\subsubsection{21.11.3 用归纳法证明集合计数}\label{ux7528ux5f52ux7eb3ux6cd5ux8bc1ux660eux96c6ux5408ux8ba1ux6570}

\textbf{定理 21.23}（德摩根律的推广）设 \(A_1, A_2, \ldots, A_n\) 是全集 \(U\) 的子集，则

\[\left(\bigcup_{k=1}^{n} A_k\right)^c = \bigcap_{k=1}^{n} A_k^c, \quad \left(\bigcap_{k=1}^{n} A_k\right)^c = \bigcup_{k=1}^{n} A_k^c\]

\textbf{证明}：对 \(n\) 用数学归纳法。

\(n = 2\) 时已在 Ch 3 中证明。

\(n \Rightarrow n+1\)：

\[\left(\bigcup_{k=1}^{n+1} A_k\right)^c = \left(\left(\bigcup_{k=1}^{n} A_k\right) \cup A_{n+1}\right)^c = \left(\bigcup_{k=1}^{n} A_k\right)^c \cap A_{n+1}^c = \bigcap_{k=1}^{n} A_k^c \cap A_{n+1}^c = \bigcap_{k=1}^{n+1} A_k^c\]

\(\blacksquare\)

\hrulefill

\textbf{本章展望}：本章研究了数列的基本理论。在下一章（Ch 29）中，我们将从数列过渡到极限理论------当 \(n \to \infty\) 时数列 \(\{a_n\}\) 的行为如何？极限理论是整个分析学的基石，我们将用 \(\varepsilon\)-\(N\) 语言严格定义数列极限，用 \(\varepsilon\)-\(\delta\) 语言严格定义函数极限，并建立连续函数的介值定理和最值定理。

\hrulefill

\hrulefill

\emph{本卷从坐标系和函数的严格定义出发，系统建立了函数与初等数学的核心理论：坐标系与函数概念（Ch 23）为后续所有章节奠定基础；基本初等函数（Ch 24）包括一次函数、二次函数和幂函数；指数与对数函数（Ch 25）将指数推广到实数并建立了完整的对数理论；三角函数（Ch 26）从单位圆出发严格定义了三角函数的图像与性质；三角恒等变换（Ch 27）建立了和差角公式、二倍角公式、正弦定理和余弦定理；数列（Ch 28）系统研究了等差数列、等比数列和数学归纳法。这六章构成了从初中数学到高中代数的完整桥梁，为卷四B的微积分学习提供了必要的函数工具和代数基础。}

\emph{卷四A：函数与初等数学终。}

\emph{卷四A：函数与初等数学终。}

\emph{卷四A：函数与初等数学终。}

\hrulefill

\newpage

\chapter{03-函数与微积分/01-函数与初等函数/05-指数函数与对数函数.md}\label{ux51fdux6570ux4e0eux5faeux79efux520601-ux51fdux6570ux4e0eux521dux7b49ux51fdux657005-ux6307ux6570ux51fdux6570ux4e0eux5bf9ux6570ux51fdux6570.md}

\section{第50章：指数函数与对数函数}\label{ux7b2c50ux7ae0ux6307ux6570ux51fdux6570ux4e0eux5bf9ux6570ux51fdux6570}

\textbf{前置知识}：本章依赖 Ch 23（函数概念、反函数理论）和 Ch 24（幂函数）。指数函数和对数函数互为反函数，它们的理论建立在实数指数幂的严格定义之上。

\hrulefill

指数函数与对数函数是数学中最重要的超越函数之一，它们不仅构成了初等函数理论的核心，更是微分方程、概率论、复分析和数论等众多数学分支的基础。指数函数 \(f(x) = a^x\) 描述了''倍增''或''衰减''的连续过程------从细菌繁殖到放射性衰变，从复利计算到热传导，指数增长和指数衰减是自然界中最普遍的数学模型。对数函数作为指数函数的反函数，将乘法转化为加法（\(\log_a(xy) = \log_a x + \log_a y\)），这一性质在计算技术尚不发达的年代催生了计算尺和对数表，而即使在数字时代，对数变换在数据处理（如取对数将指数增长线性化）、信息论（熵的定义）和算法分析（二分查找的复杂度为 \(O(\log n)\)）中仍然不可或缺。本章从实数指数的严格定义出发，分两步走：首先建立指数函数 \(y = a^x\) 的单调性、连续性和运算法则，然后通过反函数理论引入对数函数 \(y = \log_a x\)，系统推导两者的图像特征和函数性质。

\subsection{18.1 实数指数的严格定义}\label{ux5b9eux6570ux6307ux6570ux7684ux4e25ux683cux5b9aux4e49}

实数指数幂的严格定义是指数函数理论的逻辑起点。在初等代数中，\(a^n\)（\(n\) 为正整数）定义为 \(a\) 自乘 \(n\) 次，\(a^0 = 1\)，\(a^{-n} = 1/a^n\)，\(a^{p/q} = \sqrt[q]{a^p}\)。这些定义对于有理指数是完备的，但对于无理指数（如 \(a^{\sqrt{2}}\)）则需要极限理论。严格的处理方法是：对于 \(a > 1\)，定义 \(a^{\alpha} = \sup\{a^r \mid r \in \mathbb{Q}, r < \alpha\}\)（若 \(0 < a < 1\)，则取 \(\inf\)）；对于 \(a = 1\)，\(1^{\alpha} = 1\)。这一定义保证了指数函数在 \(\mathbb{R}\) 上的连续性和单调性，从而为后续的导数和积分理论奠定了坚实基础。本节将逐步建立从整数指数到有理指数再到实数指数的完整链条。

\subsubsection{18.1.1 整数指数}\label{ux6574ux6570ux6307ux6570}

对于 \(a > 0\) 和正整数 \(n\)，\(a^n = \underbrace{a \cdot a \cdots a}_{n}\)。

\textbf{定义 18.1}（零指数）对于 \(a > 0\)，定义 \(a^0 = 1\)。

\textbf{定义 18.2}（负整数指数）对于 \(a > 0\) 和正整数 \(n\)，定义 \(a^{-n} = \frac{1}{a^n}\)。

零指数和负整数指数的定义分别是对指数运算法则的自然扩充：\(a^0 = 1\) 保证了 \(a^{n+0} = a^n \cdot a^0\) 的相容性，而 \(a^{-n} = 1/a^n\) 则是使 \(a^{n} \cdot a^{-n} = a^{0} = 1\) 成立的唯一定义。有了这些定义，整数指数上的运算律就可以从正整数情形推广到全体整数。

\textbf{定理 18.1}（整数指数的运算律）设 \(a > 0\)，\(m, n \in \mathbb{Z}\)。则：

\begin{enumerate}
\def\labelenumi{\arabic{enumi}.}
\tightlist
\item
  \(a^m \cdot a^n = a^{m+n}\)
\item
  \((a^m)^n = a^{mn}\)
\item
  \((ab)^n = a^n b^n\)（\(b > 0\)）
\end{enumerate}

\textbf{证明}：(1) 当 \(m, n\) 均为正整数时，由乘法的结合律直接可得。其余情形可以通过定义转化为正整数的情形。例如当 \(m > 0\)，\(n < 0\) 时，设 \(n = -k\)（\(k > 0\)），则 \(a^m \cdot a^{-k} = \frac{a^m}{a^k}\)。若 \(m > k\)，则 \(= a^{m-k} = a^{m+(-k)}\)；若 \(m < k\)，则 \(= \frac{1}{a^{k-m}} = a^{-(k-m)} = a^{m+(-k)}\)。其余情况类似。\(\blacksquare\)

\subsubsection{18.1.2 有理数指数}\label{ux6709ux7406ux6570ux6307ux6570-1}

\textbf{定义 18.3}（正有理数指数）设 \(a > 0\)，\(m, n\) 为正整数。定义：

\[a^{m/n} = \sqrt[n]{a^m} = (\sqrt[n]{a})^m\]

其中 \(\sqrt[n]{a}\) 表示 \(a\) 的 \(n\) 次算术根（即满足 \(x^n = a\) 的唯一正实数 \(x\)，其存在性由实数的完备性保证）。

\textbf{定理 18.2}（有理数指数的良定义性）上述定义与 \(m/n\) 的表示方式无关。即若 \(\frac{m}{n} = \frac{p}{q}\)，则 \(a^{m/n} = a^{p/q}\)。

\textbf{证明}：由 \(\frac{m}{n} = \frac{p}{q}\)，有 \(mq = np\)。设 \(x = a^{m/n}\)，即 \(x^n = a^m\)。则

\[x^{nq} = (x^n)^q = (a^m)^q = a^{mq} = a^{np} = (a^n)^p\]

同时 \(x^{nq} = (x^q)^n\)。由 \(x^n = a^m\)，得 \((x^n)^q = a^{mq}\)。而 \(a^{p/q}\) 满足 \(y^q = a^p\)，故 \(y^{nq} = (y^q)^n = a^{np} = a^{mq}\)。因此 \(x^{nq} = y^{nq}\)，由正实数 \(nq\) 次方根的唯一性得 \(x = y\)。\(\blacksquare\)

\subsubsection{18.1.3 实数指数}\label{ux5b9eux6570ux6307ux6570-1}

\textbf{定义 18.4}（实数指数）设 \(a > 0\)，\(x \in \mathbb{R}\)。取有理数列 \(\{r_n\}\)，使得 \(\lim_{n \to \infty} r_n = x\)。定义

\[a^x = \lim_{n \to \infty} a^{r_n}\]

实数指数的定义是通过有理逼近将指数函数从有理数集 \(\mathbb{Q}\) 延拓到实数集 \(\mathbb{R}\) 的关键一步。然而，这一定义依赖于所选有理数列的特定选择------如果用不同的有理数列逼近同一个实数 \(x\)，得到的极限是否相同？以下定理确保了这个延拓是合理的：无论选择哪个收敛于 \(x\) 的有理数列，\(a^{r_n}\) 都收敛于同一个确定的值。

\textbf{定理 18.3}（实数指数的良定义性）上述定义与有理数列 \(\{r_n\}\) 的选取无关。

\textbf{证明}：设 \(\{r_n\}\) 和 \(\{s_n\}\) 都是收敛于 \(x\) 的有理数列。需要证明 \(\lim_{n \to \infty} a^{r_n} = \lim_{n \to \infty} a^{s_n}\)。

由于 \(|r_n - s_n| \leq |r_n - x| + |x - s_n| \to 0\)，所以 \(r_n - s_n \to 0\)。

当 \(a > 1\) 时，由有理数指数的单调性，\(a^{r_n - s_n}\) 介于 \(a^{-|r_n - s_n|}\) 和 \(a^{|r_n - s_n|}\) 之间。由于 \(|r_n - s_n| \to 0\)，两者都趋近于 \(a^0 = 1\)，故 \(a^{r_n - s_n} \to 1\)。

所以

\[\lim_{n \to \infty} a^{r_n} = \lim_{n \to \infty} (a^{s_n} \cdot a^{r_n - s_n}) = \lim_{n \to \infty} a^{s_n} \cdot 1 = \lim_{n \to \infty} a^{s_n}\]

\(0 < a < 1\) 的情形类似。\(\blacksquare\)

\textbf{定理 18.4}（实数指数的运算律）设 \(a, b > 0\)，\(\alpha, \beta \in \mathbb{R}\)，则：

\begin{enumerate}
\def\labelenumi{\arabic{enumi}.}
\tightlist
\item
  \(a^{\alpha} \cdot a^{\beta} = a^{\alpha + \beta}\)
\item
  \((a^{\alpha})^{\beta} = a^{\alpha \beta}\)
\item
  \((ab)^{\alpha} = a^{\alpha} \cdot b^{\alpha}\)
\end{enumerate}

\textbf{证明}：取有理数列 \(\{r_n\} \to \alpha\)，\(\{s_n\} \to \beta\)。

\begin{enumerate}
\def\labelenumi{(\arabic{enumi})}
\item
  \(a^{\alpha} \cdot a^{\beta} = \lim_{n \to \infty} a^{r_n} \cdot \lim_{n \to \infty} a^{s_n} = \lim_{n \to \infty} (a^{r_n} \cdot a^{s_n}) = \lim_{n \to \infty} a^{r_n + s_n} = a^{\alpha + \beta}\)
\item
  取有理数列 \(\{s_n\} \to \beta\)。分两步完成。
\end{enumerate}

\textbf{整数幂引理}（仅依赖 (1)）：对任意实数 \(\gamma\) 和正整数 \(n\)，由 (1) 的乘法律反复应用（归纳法）：

\[(a^{\gamma})^n = \underbrace{a^{\gamma} \cdot a^{\gamma} \cdots a^{\gamma}}_{n \text{ 个}} = a^{n\gamma}\]

注意此引理是 (1) 的直接推论，\textbf{不使用} (2) 本身。

\begin{itemize}
\item
  \begin{enumerate}
  \def\labelenumi{(\roman{enumi})}
  \tightlist
  \item
    先证对有理数 \(s_n\) 有 \((a^{\alpha})^{s_n} = a^{\alpha \cdot s_n}\)。设 \(s_n = p/q\)（\(p \in \mathbb{Z}\)，\(q \in \mathbb{N}^+\)）。
  \end{enumerate}

  \textbf{整数幂部分}：由上述整数幂引理（仅用 (1)），\((a^{\alpha})^p = a^{p\alpha}\)。

  \textbf{分数幂部分}：设 \(x = a^{p\alpha/q}\)（实数指数，由定义 18.4 确定）。由整数幂引理（仅用 (1)），\(x^q = (a^{p\alpha/q})^q = a^{q \cdot p\alpha/q} = a^{p\alpha}\)。因 \(x > 0\)，由 \(q\) 次根的唯一性，\(x = (a^{p\alpha})^{1/q}\)。

  \textbf{合并}：由有理数幂的定义 \(b^{p/q} = (b^p)^{1/q}\)（应用于 \(b = a^{\alpha}\)）：

  \[(a^{\alpha})^{s_n} = (a^{\alpha})^{p/q} = ((a^{\alpha})^p)^{1/q} = (a^{p\alpha})^{1/q} = a^{p\alpha/q} = a^{\alpha s_n}\]
\item
  \begin{enumerate}
  \def\labelenumi{(\roman{enumi})}
  \setcounter{enumi}{1}
  \tightlist
  \item
    由定义 18.4 的构造，\(a^x\) 在其定义域上连续（即 \(x_n \to x\) 蕴含 \(a^{x_n} \to a^x\)；这是因为 \(a^x\) 单调且 \(a^x = \sup\{a^r \mid r \in \mathbb{Q},\, r \leq x\}\)，由确界性质可证连续性）。故 \(\alpha \cdot s_n \to \alpha \beta\) 蕴含 \(\lim_{n \to \infty} a^{\alpha \cdot s_n} = a^{\alpha \beta}\)。
  \end{enumerate}
\end{itemize}

综合 (i) 和 (ii)，由实数幂的定义 \((a^{\alpha})^{\beta} = \lim_{n \to \infty} (a^{\alpha})^{s_n}\)（有理数列 \(s_n \to \beta\)），得

\[(a^{\alpha})^{\beta} = \lim_{n \to \infty} (a^{\alpha})^{s_n} = \lim_{n \to \infty} a^{\alpha s_n} = a^{\alpha \beta}\]

\begin{enumerate}
\def\labelenumi{(\arabic{enumi})}
\setcounter{enumi}{2}
\tightlist
\item
  的证明类似。\(\blacksquare\)
\end{enumerate}

\hrulefill

\subsection{18.2 指数函数}\label{ux6307ux6570ux51fdux6570}

在建立了实数指数幂的严格定义后，指数函数 \(y = a^x\)（\(a > 0, a \neq 1\)）成为定义在 \(\mathbb{R}\) 上、取值于 \((0, +\infty)\) 的严格单调函数。指数函数的核心性质------\(a^{x+y} = a^x \cdot a^y\)------将加法转化为乘法，这一性质既是其定义的自然延伸，也是其众多应用（如复利计算、指数增长模型）的数学根源。指数函数的图像恒过 \((0, 1)\)，当 \(a > 1\) 时严格递增且以 \(x\) 轴负半轴为水平渐近线，当 \(0 < a < 1\) 时严格递减且以 \(x\) 轴正半轴为水平渐近线。特别地，以 \(e\) 为底的指数函数 \(e^x\) 是唯一满足 \((e^x)' = e^x\) 的函数（即导数等于自身），这一性质使其在微积分中占据核心地位，也是欧拉公式 \(e^{i\theta} = \cos\theta + i\sin\theta\) 的基础。

\subsubsection{18.2.1 定义}\label{ux5b9aux4e49}

\textbf{定义 18.5}（指数函数）设 \(a > 0\) 且 \(a \neq 1\)。\textbf{指数函数} \(f: \mathbb{R} \to (0, +\infty)\) 定义为：

\[f(x) = a^x\]

这一定义的关键在于：底数 \(a\) 必须严格为正且不等于 \(1\)。当 \(a = 1\) 时，\(f(x) = 1^x = 1\) 为常数函数，退化为平凡情形，其反函数不存在，因此通常将其排除在指数函数的讨论之外。指数函数的值域 \((0, +\infty)\) 反映了这样一个事实：无论实数指数取何值，正数的任何实数次幂始终为正。这一性质与指数函数的连续性密切相关------在连续性的框架下，\(a^x\) 不可能''跳过'' \(0\) 变为负值，因为 \(a^0 = 1\) 且 \(a^q > 0\) 对所有有理数 \(q\) 成立，由连续性可知 \(a^x > 0\) 对所有实数 \(x\) 成立。指数函数的定义域是全体实数 \(\mathbb{R}\)，这得益于实数指数幂的严格定义（见 18.1 节），它将指数从有理数集 \(\mathbb{Q}\) 连续地延拓到整个实数集 \(\mathbb{R}\)。

\subsubsection{18.2.2 指数函数的性质}\label{ux6307ux6570ux51fdux6570ux7684ux6027ux8d28}

\textbf{引理 18.1}（严格单调函数的中值性质）设 \(f\) 在区间 \([p, q]\) 上严格单调递增，且满足如下\textbf{确界连续性条件}：对任意 \(x \in (p, q)\)，\(f(x) = \sup\{f(t) \mid t \in [p, x)\} = \inf\{f(t) \mid t \in (x, q]\}\)。（指数函数 \(a^x\)（\(a > 1\)）由定义 18.4 的构造 \(a^x = \sup\{a^r \mid r \in \mathbb{Q}, r \leq x\}\) 直接满足此条件。）若 \(f(p) < c < f(q)\)，则存在唯一的 \(\xi \in (p, q)\) 使得 \(f(\xi) = c\)。

\textbf{证明}：定义集合 \(S = \{x \in [p, q] \mid f(x) \leq c\}\)。由于 \(f(p) < c\)，故 \(p \in S\)，\(S\) 非空且有上界 \(q\)。由确界存在定理，设 \(\xi = \sup S\)，则 \(\xi \in [p, q]\)。

先证 \(f(\xi) \leq c\)：由上确界定义，对每个 \(n \in \mathbb{N}^+\)，存在 \(x_n \in S\) 使得 \(\xi - \frac{1}{n} < x_n \leq \xi\)。由确界连续性条件，\(f(\xi) = \sup\{f(t) \mid t < \xi\}\)。由 \(f(x_n) \leq c\)，得 \(f(\xi) \leq c\)。

再证 \(f(\xi) \geq c\)：先证 \(\xi < q\)。若 \(\xi = q\)，则由确界连续性条件，\(f(q) = \sup\{f(t) \mid t < q\}\)。存在 \(x_n \in S\) 使得 \(x_n \to q\)，故 \(f(q) = \sup f(x_n) \leq c\)，与 \(f(q) > c\) 矛盾。故 \(\xi < q\)。由确界连续性条件，\(f(\xi) = \inf\{f(t) \mid t > \xi\}\)。对任意 \(t \in (\xi, q]\)，由 \(\xi = \sup S\) 知 \(t \notin S\)，即 \(f(t) > c\)。因此 \(f(\xi) = \inf\{f(t) \mid t > \xi\} \geq c\)。

综上 \(f(\xi) = c\)。唯一性由 \(f\) 的严格单调性保证。\(\blacksquare\)

\textbf{注}：引理 18.1 的证明仅依赖确界原理和函数的确界连续性条件（由定义 18.4 的构造保证），不依赖后续章节的连续性理论（Ch 29）。这是介值定理在严格单调函数上的特殊形式。

\hrulefill

\textbf{定理 18.5}（指数函数的性质）设 \(f(x) = a^x\)，其中 \(a > 0\)，\(a \neq 1\)。

\begin{enumerate}
\def\labelenumi{\arabic{enumi}.}
\tightlist
\item
  \(f\) 的定义域为 \(\mathbb{R}\)，值域为 \((0, +\infty)\)。
\item
  \(f(0) = 1\)，即图像过点 \((0, 1)\)。
\item
  当 \(a > 1\) 时，\(f\) 在 \(\mathbb{R}\) 上严格递增；当 \(0 < a < 1\) 时，\(f\) 在 \(\mathbb{R}\) 上严格递减。
\item
  \(f\) 无上界，下确界为 \(0\)（但取不到）。
\end{enumerate}

\textbf{证明}：

\begin{enumerate}
\def\labelenumi{(\arabic{enumi})}
\tightlist
\item
  定义域由定义直接得到。值域：首先 \(a^x > 0\) 对所有 \(x \in \mathbb{R}\) 成立（因为正数的任意实数次幂仍为正数）。下证值域恰好为 \((0, +\infty)\)。不妨设 \(a > 1\)（\(0 < a < 1\) 时令 \(b = 1/a > 1\)，则 \(a^x = b^{-x}\)，值域相同）。
\end{enumerate}

先证：对任意 \(M > 0\)，存在正整数 \(n\) 使得 \(a^n > M\)。若不然，则 \(\{a^n\}\) 有上界，设 \(s = \sup\{a^n : n \in \mathbb{N}\}\)。由上确界定义，存在 \(n\) 使得 \(a^n > s/a\)，则 \(a^{n+1} = a \cdot a^n > s\)，与 \(s\) 为上确界矛盾。故 \(a^n \to +\infty\)。

对任意 \(y > 0\)： - 若 \(y = 1\)，取 \(x = 0\)，则 \(a^0 = 1 = y\)。 - 若 \(y > 1\)，取 \(x_1 = 0\)，则 \(a^{x_1} = 1 < y\)。取正整数 \(N\) 使得 \(a^N > y\)，令 \(x_2 = N\)。由实数指数的构造（定义 18.4），\(a^x\) 是连续函数。由 (3)，\(a^x\) 在 \([x_1, x_2]\) 上严格递增。由引理 18.1，存在 \(\xi \in (x_1, x_2)\) 使得 \(a^{\xi} = y\)。 - 若 \(0 < y < 1\)，则 \(1/y > 1\)，由上段存在 \(x'\) 使得 \(a^{x'} = 1/y\)，故 \(a^{-x'} = y\)。

因此值域为 \((0, +\infty)\)。

\begin{enumerate}
\def\labelenumi{(\arabic{enumi})}
\setcounter{enumi}{1}
\item
  \(f(0) = a^0 = 1\)。
\item
  设 \(a > 1\)，\(x_1 < x_2\)。则
\end{enumerate}

\[\frac{f(x_2)}{f(x_1)} = \frac{a^{x_2}}{a^{x_1}} = a^{x_2 - x_1} > a^0 = 1\]

所以 \(f(x_2) > f(x_1)\)，\(f\) 严格递增。当 \(0 < a < 1\) 时类似证明。

\begin{enumerate}
\def\labelenumi{(\arabic{enumi})}
\setcounter{enumi}{3}
\tightlist
\item
  由 (3)，\(f\) 严格单调，故无上界。又 \(a^x > 0\) 对所有 \(x\) 成立，故下确界为 \(0\)（但取不到）。\(\blacksquare\)
\end{enumerate}

\textbf{定理 18.6}（指数函数的无零点性）指数函数 \(f(x) = a^x\)（\(a > 0\)，\(a \neq 1\)）在 \(\mathbb{R}\) 上恒不为零。

\textbf{证明}：对于任意 \(x \in \mathbb{R}\)，\(a^x > 0\)（因为 \(a > 0\)，正数的任意实数次幂仍为正数）。\(\blacksquare\)

\hrulefill

\subsection{18.3 对数函数}\label{ux5bf9ux6570ux51fdux6570}

对数函数是指数函数的反函数，这一关系是理解对数函数全部性质的钥匙。由于指数函数 \(y = a^x\) 是严格单调的（当 \(a > 1\) 时严格递增，当 \(0 < a < 1\) 时严格递减），它必然存在反函数，定义域为 \((0, +\infty)\)，值域为 \(\mathbb{R}\)。对数函数将指数函数的''加法转乘法''性质反转------\(\log_a(xy) = \log_a x + \log_a y\)------将乘法运算转化为加法运算，这一性质在历史上彻底改变了科学计算的方式（对数表和对数尺）。以 \(e\) 为底的自然对数 \(\ln x\) 在微积分中具有特殊地位：\((\ln x)' = \frac{1}{x}\)，这一简洁的导数公式是自然对数优于其他底数对数的根本原因。对数函数的图像恒过 \((1, 0)\)，\(y\) 轴为其垂直渐近线，其凹凸性由底数 \(a\) 决定：当 \(a > 1\) 时为凹函数，当 \(0 < a < 1\) 时为凸函数。

\subsubsection{18.3.1 对数的定义}\label{ux5bf9ux6570ux7684ux5b9aux4e49}

\textbf{定义 18.6}（对数）设 \(a > 0\)，\(a \neq 1\)。对于 \(x > 0\)，\textbf{以 \(a\) 为底 \(x\) 的对数} \(\log_a x\) 定义为满足 \(a^y = x\) 的唯一实数 \(y\)，即：

\[y = \log_a x \Leftrightarrow a^y = x\]

这一定义将对数运算与指数运算建立了一一对应关系。从定义可以直接看出：\(\log_a a = 1\)（因为 \(a^1 = a\)），\(\log_a 1 = 0\)（因为 \(a^0 = 1\)），以及恒等式 \(a^{\log_a x} = x\)（\(x > 0\)）和 \(\log_a(a^y) = y\)（\(y \in \mathbb{R}\)）。后两个恒等式精确地表达了指数函数与对数函数之间的互逆关系：先取对数再取指数（或反之）等于恒等变换。在历史上，常用对数 \(\log_{10} x\) 在计算工具不发达的年代被广泛使用------布里格斯（Henry Briggs）在 1617 年出版了第一份常用对数表，将乘除运算转化为加减运算，将乘方开方转化为乘除运算，极大地便利了天文、航海和工程计算。自然对数 \(\ln x = \log_e x\) 则以 \(e\) 为底，在理论分析中最为自然：\(\ln x\) 的导数恰好是 \(\frac{1}{x}\)，且 \(\ln(1 + x) = x - \frac{x^2}{2} + \frac{x^3}{3} - \cdots\) 的泰勒展开式具有简洁的系数结构。

\subsubsection{18.3.2 对数的存在性和唯一性}\label{ux5bf9ux6570ux7684ux5b58ux5728ux6027ux548cux552fux4e00ux6027}

\textbf{定理 18.7}（对数的存在唯一性）设 \(a > 0\)，\(a \neq 1\)，\(x > 0\)。则存在唯一的实数 \(y\) 使得 \(a^y = x\)。

\textbf{证明}：

存在性：考虑函数 \(f(t) = a^t\)。当 \(a > 1\) 时，\(f\) 严格递增（由定理 18.5(3)）且连续（由实数指数的构造，定义 18.4）。由于 \(f(0) = 1\)，当 \(x > 1\) 时，由定理 18.5(1) 的证明可知存在 \(N\) 使得 \(f(N) = a^N > x\)。在 \([0, N]\) 上 \(f\) 连续且严格递增，由引理 18.1，存在 \(y \in (0, N)\) 使得 \(f(y) = x\)。当 \(0 < x < 1\) 时，\(1/x > 1\)，由上段存在 \(y' > 0\) 使得 \(a^{y'} = 1/x\)，则 \(a^{-y'} = x\)，取 \(y = -y' < 0\)。当 \(x = 1\) 时，\(y = 0\)。\(0 < a < 1\) 的情形类似。

唯一性：由于 \(f\) 严格单调，所以 \(y\) 唯一。\(\blacksquare\)

\subsubsection{18.3.3 对数函数的性质}\label{ux5bf9ux6570ux51fdux6570ux7684ux6027ux8d28}

\textbf{定理 18.8}（对数函数的性质）设 \(g(x) = \log_a x\)（\(a > 0\)，\(a \neq 1\)）。

\begin{enumerate}
\def\labelenumi{\arabic{enumi}.}
\tightlist
\item
  \textbf{定义域}：\((0, +\infty)\)，\textbf{值域}：\(\mathbb{R}\)。
\item
  \(\log_a 1 = 0\)，\(\log_a a = 1\)。
\item
  当 \(a > 1\) 时，\(g\) 在 \((0, +\infty)\) 上严格递增；当 \(0 < a < 1\) 时，\(g\) 在 \((0, +\infty)\) 上严格递减。
\item
  \(g\) 是指数函数 \(f(x) = a^x\) 的反函数。
\end{enumerate}

\textbf{证明}：

\begin{enumerate}
\def\labelenumi{(\arabic{enumi})}
\setcounter{enumi}{1}
\item
  \(a^0 = 1\)，故 \(\log_a 1 = 0\)。\(a^1 = a\)，故 \(\log_a a = 1\)。
\item
  对数函数是指数函数的反函数，由定理 16.12，反函数与原函数的单调性相同。
\item
  由定义，\(y = \log_a x \Leftrightarrow a^y = x\)，即 \(g = f^{-1}\)。\(\blacksquare\)
\end{enumerate}

\subsubsection{18.3.4 对数的运算法则}\label{ux5bf9ux6570ux7684ux8fd0ux7b97ux6cd5ux5219}

\textbf{定理 18.9}（对数运算法则）设 \(a > 0\)，\(a \neq 1\)，\(M, N > 0\)。则：

\begin{enumerate}
\def\labelenumi{(\arabic{enumi})}
\item
  \(\log_a(MN) = \log_a M + \log_a N\)
\item
  \(\log_a\frac{M}{N} = \log_a M - \log_a N\)
\item
  \(\log_a M^n = n \log_a M\)，其中 \(n \in \mathbb{R}\)
\end{enumerate}

\textbf{证明}：

\begin{enumerate}
\def\labelenumi{(\arabic{enumi})}
\tightlist
\item
  设 \(\log_a M = p\)，\(\log_a N = q\)。则 \(a^p = M\)，\(a^q = N\)。
\end{enumerate}

\[MN = a^p \cdot a^q = a^{p+q}\]

所以 \(\log_a(MN) = p + q = \log_a M + \log_a N\)。

\begin{enumerate}
\def\labelenumi{(\arabic{enumi})}
\setcounter{enumi}{1}
\tightlist
\item
  由 (1)，\(\log_a M = \log_a\left(\frac{M}{N} \cdot N\right) = \log_a\frac{M}{N} + \log_a N\)。
\end{enumerate}

所以 \(\log_a\frac{M}{N} = \log_a M - \log_a N\)。

\begin{enumerate}
\def\labelenumi{(\arabic{enumi})}
\setcounter{enumi}{2}
\tightlist
\item
  设 \(\log_a M = p\)，则 \(a^p = M\)。所以 \(M^n = (a^p)^n = a^{pn}\)，即 \(\log_a M^n = pn = n \log_a M\)。\(\blacksquare\)
\end{enumerate}

\subsubsection{18.3.5 换底公式}\label{ux6362ux5e95ux516cux5f0f}

\textbf{定理 18.10}（换底公式）设 \(a, b > 0\)，\(a, b \neq 1\)，\(x > 0\)。则

\[\log_a x = \frac{\log_b x}{\log_b a}\]

\textbf{证明}：设 \(\log_a x = y\)，则 \(a^y = x\)。

对两边取以 \(b\) 为底的对数：\(\log_b(a^y) = \log_b x\)。

由定理 18.9(3)，\(y \log_b a = \log_b x\)。

所以 \(y = \frac{\log_b x}{\log_b a}\)，即 \(\log_a x = \frac{\log_b x}{\log_b a}\)。\(\blacksquare\)

\subsubsection{18.3.6 自然对数与常用对数}\label{ux81eaux7136ux5bf9ux6570ux4e0eux5e38ux7528ux5bf9ux6570}

\textbf{定义 18.7}（自然对数）以 \(e\) 为底的对数称为\textbf{自然对数}，记作 \(\ln x\)，即 \(\ln x = \log_e x\)。

其中 \(e = \lim_{n \to \infty}\left(1 + \frac{1}{n}\right)^n \approx 2.71828...\) 是一个重要的数学常数（其极限的存在性将在 Ch 29 中证明）。

\textbf{定义 18.8}（常用对数）以 \(10\) 为底的对数称为\textbf{常用对数}，记作 \(\lg x\)，即 \(\lg x = \log_{10} x\)。

\textbf{推论 18.1} \(\ln x = \frac{\lg x}{\lg e}\)，\(\lg x = \frac{\ln x}{\ln 10}\)。

\textbf{证明}：直接由换底公式可得。\(\blacksquare\)

\hrulefill

\subsection{18.4 函数零点与二分法}\label{ux51fdux6570ux96f6ux70b9ux4e0eux4e8cux5206ux6cd5}

函数的零点问题是方程求解的几何对应------方程 \(f(x) = 0\) 的解恰好是函数 \(y = f(x)\) 的图像与 \(x\) 轴的交点。零点存在性定理（连续函数在闭区间两端异号时必有零点）是介值定理的直接推论，它为数值求解方程提供了存在性保证，但不能直接给出零点的精确值。二分法（Bisection Method）是求零点近似值的最基本算法：每次迭代将区间缩小一半，经过 \(n\) 次迭代后，区间长度缩小为原来的 \(1/2^n\)，因此只需 \(\log_2\left(\frac{b-a}{\varepsilon}\right)\) 次迭代即可将零点定位到精度 \(\varepsilon\) 以内。二分法的收敛性仅依赖于连续性和介值定理，不需要可导性，因此适用范围极广------但它每次只能保证一个零点，且收敛速度仅为线性收敛（\(O(1/2^n)\)），远慢于牛顿法等超线性收敛算法。本节将建立零点存在性定理的严格证明，并系统介绍二分法的原理和误差估计。

\subsubsection{18.4.1 零点的定义}\label{ux96f6ux70b9ux7684ux5b9aux4e49}

\textbf{定义 18.9}（函数零点）设 \(f: D \to \mathbb{R}\)。若 \(x_0 \in D\) 满足 \(f(x_0) = 0\)，则称 \(x_0\) 为 \(f\) 的一个\textbf{零点}。

零点的概念是方程理论的几何化。对于方程 \(f(x) = 0\)，其解在函数图像上表现为曲线 \(y = f(x)\) 与 \(x\) 轴的交点。当 \(f\) 是多项式时，零点就是代数方程的根------代数基本定理（定理 14.15）断言 \(n\) 次多项式在复数域中恰有 \(n\) 个根（计重数）。当 \(f\) 是超越函数（如 \(f(x) = e^x - x - 2\)）时，零点通常无法用初等函数表达，只能通过数值方法（如二分法、牛顿法）求近似值。零点的''重数''概念（若 \(f(x_0) = f'(x_0) = \cdots = f^{(k-1)}(x_0) = 0\) 但 \(f^{(k)}(x_0) \neq 0\)，则 \(x_0\) 为 \(k\) 重零点）在多项式因式分解和函数图像分析中具有重要应用------\(k\) 重零点处函数图像与 \(x\) 轴的交点行为不同于单重零点，例如 \(f(x) = x^2\) 在 \(x = 0\) 处''触碰''而非''穿过'' \(x\) 轴。零点的存在性问题------即方程 \(f(x) = 0\) 是否有解------是分析学中最基本的问题之一，零点存在性定理（定理 18.11）是回答这一问题的第一个重要工具。

\subsubsection{18.4.2 零点存在性定理}\label{ux96f6ux70b9ux5b58ux5728ux6027ux5b9aux7406}

\textbf{定理 18.11}（零点存在性定理）若函数 \(f(x)\) 在闭区间 \([a, b]\) 上连续，且 \(f(a) \cdot f(b) < 0\)（即 \(f(a)\) 与 \(f(b)\) 异号），则存在 \(\xi \in (a, b)\)，使得 \(f(\xi) = 0\)。

\textbf{证明}：不妨设 \(f(a) < 0 < f(b)\)（另一情形类似）。定义集合 \(S = \{x \in [a, b] \mid f(x) \leq 0\}\)。由 \(f(a) < 0\) 知 \(a \in S\)，\(S\) 非空且有上界 \(b\)。由确界存在定理，设 \(\xi = \sup S\)，则 \(\xi \in [a, b]\)。

先证 \(f(\xi) \leq 0\)。取 \(x_n \in S\) 使得 \(x_n \to \xi\)（由上确界的定义，对每个 \(n\)，存在 \(x_n \in S\) 满足 \(\xi - \frac{1}{n} < x_n \leq \xi\)）。由 \(f\) 连续，\(f(\xi) = \lim_{n \to \infty} f(x_n)\)。而 \(f(x_n) \leq 0\)，故 \(f(\xi) \leq 0\)。

再证 \(f(\xi) \geq 0\)。取 \(t_n = \xi + \frac{1}{n}\)（当 \(n\) 充分大时 \(t_n > \xi\) 且 \(t_n \leq b\)），则 \(t_n \notin S\)，故 \(f(t_n) > 0\)。由 \(f\) 连续，\(f(\xi) = \lim_{n \to \infty} f(t_n) \geq 0\)。

综上，\(f(\xi) = 0\)。\(\blacksquare\)

\textbf{注}：此证明仅依赖确界原理和连续性的操作性定义，不依赖 Ch 29 的介值定理。介值定理（Ch 29，定理 22.15）将在后续建立完整理论后给出更一般的证明。

\subsubsection{18.4.3 二分法}\label{ux4e8cux5206ux6cd5}

\textbf{定义 18.10}（二分法）二分法是一种基于零点存在性定理的数值方法，用于近似求解方程 \(f(x) = 0\) 的根。

\textbf{算法描述}：

设 \(f\) 在 \([a_0, b_0]\) 上连续，\(f(a_0) \cdot f(b_0) < 0\)。对 \(n = 0, 1, 2, \ldots\)：

\begin{enumerate}
\def\labelenumi{\arabic{enumi}.}
\tightlist
\item
  计算中点 \(m_n = \frac{a_n + b_n}{2}\)。
\item
  若 \(f(m_n) = 0\)，则 \(m_n\) 就是零点，算法终止。
\item
  若 \(f(a_n) \cdot f(m_n) < 0\)，则零点在 \((a_n, m_n)\) 内，令 \(a_{n+1} = a_n\)，\(b_{n+1} = m_n\)。
\item
  若 \(f(m_n) \cdot f(b_n) < 0\)，则零点在 \((m_n, b_n)\) 内，令 \(a_{n+1} = m_n\)，\(b_{n+1} = b_n\)。
\end{enumerate}

\textbf{定理 18.12}（二分法的收敛性）设 \(f\) 在 \([a_0, b_0]\) 上连续，\(f(a_0) \cdot f(b_0) < 0\)，且 \(f\) 在 \((a_0, b_0)\) 上单调。则二分法产生的区间序列 \(\{[a_n, b_n]\}\) 满足：

\[b_n - a_n = \frac{b_0 - a_0}{2^n} \to 0\]

且序列 \(\{a_n\}\) 和 \(\{b_n\}\) 都收敛于 \(f\) 在 \((a_0, b_0)\) 内的唯一零点 \(\xi\)。

\textbf{证明}：由构造，\(b_n - a_n = \frac{b_0 - a_0}{2^n}\)。

\(\{a_n\}\) 单调递增且有上界 \(b_0\)，由确界存在定理，\(\{a_n\}\) 有上确界 \(\alpha = \sup\{a_n\}\)。同理 \(\{b_n\}\) 单调递减且有下界 \(a_0\)，有下确界 \(\beta = \inf\{b_n\}\)。

由 \(b_n - a_n \to 0\)，得 \(\alpha = \beta\)。

由 \(f\) 的连续性（确界连续性条件），\(f(\alpha) = \sup f(a_n) \leq 0\) 且 \(f(\alpha) = \inf f(b_n) \geq 0\)（或反之），故 \(f(\alpha) = 0\)。

由 \(f\) 的单调性，零点唯一，故 \(\alpha = \xi\)。\(\blacksquare\)

\textbf{注}：此证明仅依赖确界存在定理，不依赖 Ch 29 的单调有界定理。单调有界定理（Ch 29，定理 22.5）将在后续给出更一般的证明。

\hrulefill

\subsection{18.5 指数方程与对数方程}\label{ux6307ux6570ux65b9ux7a0bux4e0eux5bf9ux6570ux65b9ux7a0b}

指数方程与对数方程是超越方程中最基本的两类。由于指数函数和对数函数互为反函数，这两类方程的求解方法互为镜像------指数方程通常通过取对数转化为代数方程，对数方程则通过取指数（即化为指数形式）来消去对数。核心原理是：指数函数 \(y = a^x\) 的严格单调性保证了 \(a^u = a^v \iff u = v\)，而对数函数 \(y = \log_a x\) 的严格单调性同样保证了 \(\log_a u = \log_a v \iff u = v\)（其中 \(u, v > 0\)）。然而，在实际求解中，必须注意定义域的限制------对数方程中真数必须为正，指数方程在底数不确定时也需分类讨论。某些方程同时涉及指数和对数，无法用初等方法求解（如 \(x = e^{-x}\)），此时需要借助数值方法（如二分法或牛顿法）。本节系统梳理指数方程和对数方程的基本类型和求解策略，为后续微积分中处理涉及指数函数和对数函数的方程奠定基础。

\subsubsection{18.5.1 指数方程}\label{ux6307ux6570ux65b9ux7a0b}

\textbf{定义 18.11}（指数方程）形如 \(a^{f(x)} = a^{g(x)}\) 或可化为此形式的方程称为\textbf{指数方程}，其中 \(a > 0\)，\(a \neq 1\)。

\textbf{定理 18.13} 指数方程 \(a^{f(x)} = a^{g(x)}\)（\(a > 0\)，\(a \neq 1\)）等价于 \(f(x) = g(x)\)。

\textbf{证明}：指数函数 \(y = a^x\) 是严格单调的（由定理 18.5），因此 \(a^u = a^v \iff u = v\)。\(\blacksquare\)

\subsubsection{18.5.2 对数方程}\label{ux5bf9ux6570ux65b9ux7a0b}

\textbf{定义 18.12}（对数方程）含有对数表达式且未知数出现在对数内的方程称为\textbf{对数方程}。

\textbf{定理 18.14} 对数方程 \(\log_a f(x) = \log_a g(x)\)（\(a > 0\)，\(a \neq 1\)）在 \(f(x) > 0\) 且 \(g(x) > 0\) 的条件下等价于 \(f(x) = g(x)\)。

\textbf{证明}：对数函数 \(y = \log_a x\) 在 \((0, +\infty)\) 上严格单调，因此在定义域内 \(\log_a u = \log_a v \iff u = v\)。\(\blacksquare\)

\textbf{注意}：解对数方程时必须检验所得解是否满足真数大于零的条件。

\hrulefill

\subsection{18.6 指数函数与对数函数的图像关系}\label{ux6307ux6570ux51fdux6570ux4e0eux5bf9ux6570ux51fdux6570ux7684ux56feux50cfux5173ux7cfb}

\textbf{定理 18.15} 指数函数 \(y = a^x\) 和对数函数 \(y = \log_a x\) 的图像关于直线 \(y = x\) 对称。

\textbf{证明}：对数函数是指数函数的反函数，由定理 16.13，函数与其反函数的图像关于 \(y = x\) 对称。\(\blacksquare\)

\textbf{定理 18.16} 设 \(a > b > 1\)。则：

\begin{enumerate}
\def\labelenumi{\arabic{enumi}.}
\tightlist
\item
  当 \(x > 0\) 时，\(a^x > b^x > 1\)。
\item
  当 \(x < 0\) 时，\(0 < a^x < b^x < 1\)。
\item
  当 \(x > 1\) 时，\(\log_a x < \log_b x\)。
\item
  当 \(0 < x < 1\) 时，\(\log_a x > \log_b x\)（都为负值）。
\end{enumerate}

\textbf{证明}：

\begin{enumerate}
\def\labelenumi{(\arabic{enumi})}
\item
  \(\frac{a^x}{b^x} = \left(\frac{a}{b}\right)^x\)。当 \(x > 0\) 时，\(\frac{a}{b} > 1\)，故 \(\left(\frac{a}{b}\right)^x > 1\)，即 \(a^x > b^x\)。
\item
  由换底公式，\(\log_a x = \frac{\ln x}{\ln a}\)，\(\log_b x = \frac{\ln x}{\ln b}\)。
\end{enumerate}

当 \(x > 1\) 时，\(\ln x > 0\)。又 \(\ln a > \ln b > 0\)，故 \(\frac{\ln x}{\ln a} < \frac{\ln x}{\ln b}\)。\(\blacksquare\)

\hrulefill

\subsection{18.7 幂指函数}\label{ux5e42ux6307ux51fdux6570}

\textbf{定义 18.13}（幂指函数）形如 \(y = [f(x)]^{g(x)}\) 的函数称为\textbf{幂指函数}，其中 \(f(x) > 0\)。

\textbf{定理 18.17} 幂指函数可以化为指数函数和对数函数的复合：

\[[f(x)]^{g(x)} = e^{g(x) \ln f(x)}\]

\textbf{证明}：设 \(u = f(x) > 0\)，\(v = g(x)\)。由实数指数的定义，\(u^v = e^{v \ln u}\)。\(\blacksquare\)

\hrulefill

\subsection{18.8 指数函数与对数函数的应用模型}\label{ux6307ux6570ux51fdux6570ux4e0eux5bf9ux6570ux51fdux6570ux7684ux5e94ux7528ux6a21ux578b}

指数函数与对数函数不仅是数学理论中的抽象对象，更是描述自然界和社会现象的基本数学模型。指数增长模型 \(Q(t) = Q_0 e^{kt}\)（\(k > 0\)）描述了量的变化率与当前量成正比的系统------细菌繁殖、复利增长、人口增长（在资源充足时）都遵循这一规律。指数衰减模型 \(Q(t) = Q_0 e^{-kt}\)（\(k > 0\)）则描述了放射性衰变、药物代谢、电容放电等衰减过程。对数函数在数据分析和信息论中扮演关键角色：对数变换可以将指数型数据''拉直''以便于线性回归分析；对数尺度（如里氏震级、分贝、pH 值）将跨越多个数量级的物理量映射到便于理解和比较的数值范围。在算法分析中，二分查找的复杂度 \(O(\log n)\) 和归并排序的复杂度 \(O(n \log n)\) 都涉及对数函数，反映了''分治''策略的效率。本节将系统介绍指数增长、指数衰减、对数增长（逻辑斯蒂模型）等基本应用模型，并讨论半衰期、倍增时间、\(e\)-折时间等关键参数的含义。

\subsubsection{18.8.1 指数增长与指数衰减}\label{ux6307ux6570ux589eux957fux4e0eux6307ux6570ux8870ux51cf}

\textbf{定义 18.14}（指数增长模型）若量 \(Q(t)\) 满足微分方程 \(\frac{dQ}{dt} = kQ\)（\(k > 0\)）（注：\(\frac{dQ}{dt}\) 表示 \(Q\) 对时间 \(t\) 的变化率，即导数，严格定义见 Ch 30。直观理解：\(Q\) 的瞬时变化率与当前值成正比），则

\[Q(t) = Q_0 e^{kt}\]

其中 \(Q_0 = Q(0)\) 是初始量。这称为\textbf{指数增长}模型，\(k\) 称为\textbf{增长率}。

\textbf{定义 18.15}（指数衰减模型）若 \(k < 0\)，令 \(\lambda = -k > 0\)，则 \(Q(t) = Q_0 e^{-\lambda t}\) 称为\textbf{指数衰减}模型。\textbf{半衰期} \(T_{1/2} = \frac{\ln 2}{\lambda}\)。

\textbf{定理 18.18} 在指数衰减模型中，经过一个半衰期后，量减半。

\textbf{证明}：\(Q(T_{1/2}) = Q_0 e^{-\lambda \cdot \frac{\ln 2}{\lambda}} = Q_0 e^{-\ln 2} = \frac{Q_0}{2}\)。\(\blacksquare\)

\subsubsection{18.8.2 对数尺度}\label{ux5bf9ux6570ux5c3aux5ea6}

对数函数将乘法转化为加法，这一性质被广泛用于构建对数尺度：

\begin{itemize}
\tightlist
\item
  \textbf{里氏震级}：\(M = \log_{10}\left(\frac{A}{A_0}\right)\)，其中 \(A\) 是振幅。
\item
  \textbf{分贝}：\(L = 10\log_{10}\left(\frac{I}{I_0}\right)\)，其中 \(I\) 是声强。
\item
  \textbf{pH 值}：\(\text{pH} = -\log_{10}[\text{H}^+]\)。
\end{itemize}

这些对数尺度的共同特征是将跨越多个数量级的变化压缩到一个可管理的数值范围内。

\subsection{18.9 对数不等式}\label{ux5bf9ux6570ux4e0dux7b49ux5f0f}

\textbf{定理 18.19} 对任意 \(x > 0\)，\(\ln x \leq x - 1\)，等号当且仅当 \(x = 1\)。

\begin{quote}
\textbf{前向引用说明}：以下证明使用了导数工具（Ch 30）。本不等式的纯代数证明需要幂级数展开或极限技术，均超出 Ch 25 的范围。在 Ch 30 中（定理 23.31）将基于导数重新建立此不等式及其精细形式 \(x - x^2/2 < \ln(1+x) < x\)。以下提供两个独立的证明路径。
\end{quote}

\textbf{证明（导数法）}：令 \(f(x) = x - 1 - \ln x\)（\(x > 0\)）。则 \(f'(x) = 1 - \frac{1}{x}\)。当 \(x > 1\) 时 \(f'(x) > 0\)，\(f\) 递增；当 \(0 < x < 1\) 时 \(f'(x) < 0\)，\(f\) 递减。因此 \(f\) 在 \(x = 1\) 处取最小值 \(f(1) = 0\)，即 \(f(x) \geq 0\)。\(\blacksquare\)

\textbf{证明（指数转化法）}：以下证明将问题转化为等价命题 \(e^t \geq 1 + t\)（\(t \in \mathbb{R}\)）。

\textbf{步骤 1}（伯努利不等式）：对正整数 \(n\) 和实数 \(u > -1\)，\((1+u)^n \geq 1 + nu\)。此式可用数学归纳法直接证明（Ch 28，引理 21.2）。

\textbf{步骤 2}：证明 \(e > (1+1/n)^n\) 对所有正整数 \(n\) 成立。由 \(e\) 的定义（\(e = \sup\{(1+1/k)^k : k \in \mathbb{N}^+\}\)），需证 \((1+1/k)^k\) 关于 \(k\) 严格递增。由伯努利不等式：

\[\frac{(1+1/(n+1))^{n+1}}{(1+1/n)^n} = \frac{((n+2)/(n+1))^{n+1}}{((n+1)/n)^n}\]

利用 \((1+\frac{1}{n(n+2)})^{n+1} > 1 + \frac{n+1}{n(n+2)} > 1 + \frac{1}{n+1}\) 可证此比大于 1（详细计算略）。

\textbf{步骤 3}（正有理数情形）：设 \(r = p/q > 0\)（\(p, q \in \mathbb{N}^+\)）。由步骤 2，\(e > (1+1/q)^q\)，取 \(q\) 次根得 \(e^{1/q} > 1 + 1/q\)。由乘法律和伯努利不等式：

\[e^r = (e^{1/q})^p > (1 + 1/q)^p \geq 1 + p/q = 1 + r\]

\textbf{步骤 4}（正实数情形）：对 \(t > 0\)，取有理数列 \(r_n \to t\)（\(r_n > 0\)）。由 \(e^{r_n} > 1 + r_n\) 和 \(e^x\) 的连续性，\(e^t \geq 1 + t\)。再由 \(e^t = (e^{t/2})^2 \geq (1 + t/2)^2 = 1 + t + t^2/4 > 1 + t\)，得严格不等式。

\textbf{步骤 5}（非正实数情形）：\(t = 0\) 时 \(e^0 = 1 = 1 + 0\)。对 \(t \leq -1\)，\(e^t > 0 \geq 1 + t\)。对 \(-1 < t < 0\)，设 \(u = -t \in (0,1)\)。由 \(e^u > 1 + u\)（步骤 4）和 \(e^{-u} = 1/e^u\)，再利用不等式 \(e^u \leq \frac{1}{1-u}\)（此式可由 \(e^u\) 的幂级数展开 \(e^u = \sum u^k/k! \leq \sum u^k = \frac{1}{1-u}\) 得到，其中用到 \(k! \geq 1\)），得 \(e^{-u} \geq 1 - u = 1 + t\)。

\textbf{步骤 6}（回代）：由 \(e^t \geq 1 + t\)（\(t \in \mathbb{R}\)），令 \(t = \ln x\)（\(x > 0\)）：\(x = e^{\ln x} \geq 1 + \ln x\)，故 \(\ln x \leq x - 1\)。等号成立当且仅当 \(t = 0\) 即 \(x = 1\)。\(\blacksquare\)

\textbf{推论 18.2} 对任意 \(x > -1\)，\(\ln(1 + x) \leq x\)，等号当且仅当 \(x = 0\)。

\textbf{证明}：在定理 18.19 中用 \(1 + x\) 替换 \(x\)（注意 \(x > -1\) 保证 \(1 + x > 0\)）。\(\blacksquare\)

\textbf{推论 18.3} 对任意 \(x \in \mathbb{R}\)，\(e^x \geq 1 + x\)，等号当且仅当 \(x = 0\)。

\textbf{证明}：此即定理 18.19 证明中步骤 1--5 的结论。亦可由推论 18.2 直接推出：令 \(y = 1 + x\)，若 \(y > 0\) 则 \(\ln y \leq y - 1\)，代入 \(y = e^t\) 得 \(t \leq e^t - 1\) 即 \(e^t \geq 1 + t\)；若 \(y \leq 0\)（即 \(x \leq -1\)），则 \(1 + x \leq 0 < e^x\)，不等式自然成立。\(\blacksquare\)

\hrulefill

\textbf{本章展望}：本章建立了指数函数和对数函数的完整理论。在下一章（Ch 26）中，我们将研究三角函数------另一类重要的基本初等函数。三角函数将角与实数通过单位圆联系起来，其周期性使它们在描述周期现象中不可替代。

\hrulefill

\hrulefill

\hrulefill

\hrulefill

\hrulefill

\newpage

\chapter{03-函数与微积分/01-函数与初等函数/06-三角恒等变换与反三角函数.md}\label{ux51fdux6570ux4e0eux5faeux79efux520601-ux51fdux6570ux4e0eux521dux7b49ux51fdux657006-ux4e09ux89d2ux6052ux7b49ux53d8ux6362ux4e0eux53cdux4e09ux89d2ux51fdux6570.md}

\section{第51章：三角恒等变换与反三角函数}\label{ux7b2c51ux7ae0ux4e09ux89d2ux6052ux7b49ux53d8ux6362ux4e0eux53cdux4e09ux89d2ux51fdux6570}

\textbf{前置知识}：本章依赖 Ch 26（三角函数的定义、同角关系、诱导公式）。本章的核心------和差角公式------的证明需要用到平面向量数量积的几何定义（\(\vec{a} \cdot \vec{b} = |\vec{a}||\vec{b}|\cos\theta\)）及其坐标公式，此为高中数学已学内容。其余公式均从和差角公式推导而出。

\hrulefill

和差角公式是三角恒等变换中最核心的公式，所有其他三角恒等式（二倍角、半角、积化和差、和差化积）都可以从和差角公式推导出来。本节从余弦的差角公式 \(\cos(\alpha - \beta) = \cos\alpha\cos\beta + \sin\alpha\sin\beta\) 出发，这一公式的证明巧妙地利用了向量数量积的两种等价表示------几何定义 \(\vec{a} \cdot \vec{b} = |\vec{a}||\vec{b}|\cos\theta\) 和坐标表示 \(\vec{a} \cdot \vec{b} = x_1x_2 + y_1y_2\)。通过在单位圆上取两个点，这两种表示方式分别给出了 \(\cos(\alpha - \beta)\) 和 \(\cos\alpha\cos\beta + \sin\alpha\sin\beta\)，从而建立了等式。余弦差角公式一旦确立，通过简单的变量替换和诱导公式，就可以推导出余弦和角、正弦和差角、正切和差角等全部公式。和差角公式的几何本质是旋转：将点 \((\cos\beta, \sin\beta)\) 绕原点旋转 \(\alpha\) 角，得到的新坐标为 \((\cos(\alpha + \beta), \sin(\alpha + \beta))\)，而旋转矩阵正是 \(\begin{pmatrix} \cos\alpha & -\sin\alpha \\ \sin\alpha & \cos\alpha \end{pmatrix}\)。

\textbf{定理 20.1}（余弦差角公式）对任意角 \(\alpha, \beta\)：

\[\cos(\alpha - \beta) = \cos \alpha \cos \beta + \sin \alpha \sin \beta\]

\textbf{证明（向量法）}：在单位圆上取点 \(A(\cos \alpha, \sin \alpha)\) 和 \(B(\cos \beta, \sin \beta)\)。则向量 \(\vec{OA}\) 与 \(\vec{OB}\) 的夹角为 \(\alpha - \beta\)。由数量积的定义：

\[\vec{OA} \cdot \vec{OB} = |\vec{OA}||\vec{OB}|\cos(\alpha - \beta) = 1 \cdot 1 \cdot \cos(\alpha - \beta) = \cos(\alpha - \beta)\]

利用数量积的坐标表示：

\[\vec{OA} \cdot \vec{OB} = \cos \alpha \cos \beta + \sin \alpha \sin \beta\]

因此：

\[\cos(\alpha - \beta) = \cos \alpha \cos \beta + \sin \alpha \sin \beta\]

\(\blacksquare\)

\subsubsection{20.1.2 余弦的和角公式}\label{ux4f59ux5f26ux7684ux548cux89d2ux516cux5f0f}

\textbf{定理 20.2}（余弦和角公式）对任意角 \(\alpha, \beta\)：

\[\cos(\alpha + \beta) = \cos \alpha \cos \beta - \sin \alpha \sin \beta\]

\textbf{证明}：在定理 20.1 中将 \(\beta\) 替换为 \(-\beta\)，利用 \(\cos(-\beta) = \cos \beta\)，\(\sin(-\beta) = -\sin \beta\)：

\[\cos(\alpha + \beta) = \cos[\alpha - (-\beta)] = \cos \alpha \cos(-\beta) + \sin \alpha \sin(-\beta) = \cos \alpha \cos \beta - \sin \alpha \sin \beta\]

\(\blacksquare\)

\subsubsection{20.1.3 正弦的和角公式}\label{ux6b63ux5f26ux7684ux548cux89d2ux516cux5f0f}

\textbf{定理 20.3}（正弦和角公式）对任意角 \(\alpha, \beta\)：

\[\sin(\alpha + \beta) = \sin \alpha \cos \beta + \cos \alpha \sin \beta\]

\textbf{证明}：利用诱导公式 \(\sin \theta = \cos\left(\frac{\pi}{2} - \theta\right)\) 和定理 20.1：

\[\sin(\alpha + \beta) = \cos\left[\frac{\pi}{2} - (\alpha + \beta)\right] = \cos\left[\left(\frac{\pi}{2} - \alpha\right) - \beta\right]\]

应用余弦差角公式：

\[= \cos\left(\frac{\pi}{2} - \alpha\right)\cos \beta + \sin\left(\frac{\pi}{2} - \alpha\right)\sin \beta = \sin \alpha \cos \beta + \cos \alpha \sin \beta\]

\(\blacksquare\)

\textbf{推论 20.1}（正弦差角公式）\(\sin(\alpha - \beta) = \sin \alpha \cos \beta - \cos \alpha \sin \beta\)。

\textbf{证明}：在定理 20.3 中将 \(\beta\) 替换为 \(-\beta\) 即可。\(\blacksquare\)

从正弦和角公式到正弦差角公式，仅需将 \(\beta\) 换为 \(-\beta\) 即得------这一技巧体现了三角公式体系中''和角变差角''的统一方法。类似地，正切的和角公式也可以通过同样的变量替换导出其差角形式。

\subsubsection{20.1.4 正切的和差角公式}\label{ux6b63ux5207ux7684ux548cux5deeux89d2ux516cux5f0f}

\textbf{定理 20.4}（正切和角公式）对任意角 \(\alpha, \beta\)（\(\cos \alpha, \cos \beta, \cos(\alpha + \beta)\) 均不为零）：

\[\tan(\alpha + \beta) = \frac{\tan \alpha + \tan \beta}{1 - \tan \alpha \tan \beta}\]

\textbf{证明}：

\[\tan(\alpha + \beta) = \frac{\sin(\alpha + \beta)}{\cos(\alpha + \beta)} = \frac{\sin \alpha \cos \beta + \cos \alpha \sin \beta}{\cos \alpha \cos \beta - \sin \alpha \sin \beta}\]

分子分母同除以 \(\cos \alpha \cos \beta\)：

\[= \frac{\tan \alpha + \tan \beta}{1 - \tan \alpha \tan \beta}\]

\(\blacksquare\)

\textbf{推论 20.2}（正切差角公式）\(\tan(\alpha - \beta) = \frac{\tan \alpha - \tan \beta}{1 + \tan \alpha \tan \beta}\)。

\textbf{证明}：在和角公式 \(\tan(\alpha+\beta) = \frac{\tan\alpha + \tan\beta}{1 - \tan\alpha\tan\beta}\) 中以 \(-\beta\) 代 \(\beta\)，利用 \(\tan(-\beta) = -\tan\beta\) 即得。\(\blacksquare\)

\hrulefill

二倍角公式是三角恒等变换中最常用的公式之一。本节从和差角公式出发，令 \(\beta = \alpha\)，直接推导出正弦、余弦和正切的二倍角公式。余弦的二倍角公式有三种等价形式------\(\cos 2\alpha = \cos^2\alpha - \sin^2\alpha = 2\cos^2\alpha - 1 = 1 - 2\sin^2\alpha\)，每种形式在不同场景下各有优势：第一种形式直接体现了余弦的对称性，第二种形式便于将 \(\cos^2\alpha\) 表示为 \(\cos 2\alpha\) 的形式（降幂），第三种形式则将 \(\sin^2\alpha\) 与 \(\cos 2\alpha\) 联系起来。二倍角公式的降幂变形 \(\sin^2\alpha = \frac{1 - \cos 2\alpha}{2}\) 和 \(\cos^2\alpha = \frac{1 + \cos 2\alpha}{2}\) 在积分学中极为重要，因为它们将三角函数的平方转化为一次三角函数，使积分计算变得可行。（二倍角公式）对任意角 \(\alpha\)：

\[\sin 2\alpha = 2\sin \alpha \cos \alpha\]

\[\cos 2\alpha = \cos^2 \alpha - \sin^2 \alpha = 2\cos^2 \alpha - 1 = 1 - 2\sin^2 \alpha\]

\[\tan 2\alpha = \frac{2\tan \alpha}{1 - \tan^2 \alpha}\]

\textbf{证明}：在定理 20.3 中令 \(\beta = \alpha\)：

\[\sin 2\alpha = \sin \alpha \cos \alpha + \cos \alpha \sin \alpha = 2\sin \alpha \cos \alpha\]

在定理 20.2 中令 \(\beta = \alpha\)：

\[\cos 2\alpha = \cos^2 \alpha - \sin^2 \alpha\]

利用 \(\sin^2 \alpha + \cos^2 \alpha = 1\)：

\[\cos 2\alpha = 2\cos^2 \alpha - 1 = 1 - 2\sin^2 \alpha\]

\(\blacksquare\)

\hrulefill

半角公式是二倍角公式的逆用，用于将半角三角函数用全角的余弦来表示。本节从余弦的二倍角公式出发，解出 \(\sin\frac{\alpha}{2}\) 和 \(\cos\frac{\alpha}{2}\) 的表达式。半角公式中的正负号取决于 \(\frac{\alpha}{2}\) 所在的象限，这一细节体现了三角函数符号判断的重要性。正切的半角公式 \(\tan\frac{\alpha}{2} = \frac{\sin\alpha}{1 + \cos\alpha} = \frac{1 - \cos\alpha}{\sin\alpha}\) 具有不依赖根号和符号判断的优点，因此在计算中更为方便。半角公式在积分学中也有重要应用------``万能代换'' \(t = \tan\frac{x}{2}\) 正是基于半角公式，它将任意三角有理函数的积分转化为有理函数的积分。半角公式与二倍角公式共同构成了三角函数''升幂-降幂''变换的完整工具集。（半角公式）对任意角 \(\alpha\)（使各分母不为零）：

\[\sin \frac{\alpha}{2} = \pm\sqrt{\frac{1 - \cos \alpha}{2}}, \quad \cos \frac{\alpha}{2} = \pm\sqrt{\frac{1 + \cos \alpha}{2}}\]

\[\tan \frac{\alpha}{2} = \frac{\sin \alpha}{1 + \cos \alpha} = \frac{1 - \cos \alpha}{\sin \alpha}\]

\textbf{证明}：由 \(\cos \alpha = 1 - 2\sin^2 \frac{\alpha}{2}\)，得 \(\sin^2 \frac{\alpha}{2} = \frac{1 - \cos \alpha}{2}\)，故 \(\sin \frac{\alpha}{2} = \pm\sqrt{\frac{1 - \cos \alpha}{2}}\)。正负号由 \(\frac{\alpha}{2}\) 所在象限决定。

类似地，由 \(\cos \alpha = 2\cos^2 \frac{\alpha}{2} - 1\)，得 \(\cos \frac{\alpha}{2} = \pm\sqrt{\frac{1 + \cos \alpha}{2}}\)。

对于 \(\tan \frac{\alpha}{2}\)：

\[\tan \frac{\alpha}{2} = \frac{\sin \frac{\alpha}{2}}{\cos \frac{\alpha}{2}} = \frac{2\sin \frac{\alpha}{2} \cos \frac{\alpha}{2}}{2\cos^2 \frac{\alpha}{2}} = \frac{\sin \alpha}{1 + \cos \alpha}\]

\(\blacksquare\)

\hrulefill

积化和差与和差化积公式是三角恒等变换中的''双向转换''工具。积化和差将两个三角函数的乘积转化为和差形式，这在积分计算中至关重要------例如 \(\int \sin 3x \cos 2x \, dx\) 通过积化和差可以立即转化为两个简单的正弦积分。和差化积则将三角函数的和差转化为乘积形式，这在三角方程的求解和三角恒等式的证明中尤为有用。从和差角公式出发，通过加减消元即可系统地推导出全部八条公式。值得注意的是，积化和差与和差化积互为逆运算，它们共同构成了三角学中''乘积--和差''的完整对偶体系。在高等数学中，这些公式的自然推广是傅里叶变换中的卷积定理------时域中的乘积对应频域中的卷积，反之亦然。

\textbf{定理 20.7} 对任意角 \(\alpha, \beta\)：

\[\sin \alpha \cos \beta = \frac{1}{2}[\sin(\alpha + \beta) + \sin(\alpha - \beta)]\]

\[\cos \alpha \sin \beta = \frac{1}{2}[\sin(\alpha + \beta) - \sin(\alpha - \beta)]\]

\[\cos \alpha \cos \beta = \frac{1}{2}[\cos(\alpha - \beta) + \cos(\alpha + \beta)]\]

\[\sin \alpha \sin \beta = -\frac{1}{2}[\cos(\alpha + \beta) - \cos(\alpha - \beta)]\]

\textbf{证明}：由和角公式：

\[\sin(\alpha + \beta) = \sin \alpha \cos \beta + \cos \alpha \sin \beta \quad \cdots (1)\]

\[\sin(\alpha - \beta) = \sin \alpha \cos \beta - \cos \alpha \sin \beta \quad \cdots (2)\]

\((1) + (2)\)：\(\sin(\alpha + \beta) + \sin(\alpha - \beta) = 2\sin \alpha \cos \beta\)，故第一个公式成立。

\((1) - (2)\)：\(\sin(\alpha + \beta) - \sin(\alpha - \beta) = 2\cos \alpha \sin \beta\)，故第二个公式成立。

类似地由余弦和差角公式可证后两个公式。\(\blacksquare\)

\subsubsection{20.4.2 和差化积公式}\label{ux548cux5deeux5316ux79efux516cux5f0f}

\textbf{定理 20.8} 对任意角 \(A, B\)：

\[\sin A + \sin B = 2\sin\frac{A + B}{2}\cos\frac{A - B}{2}\]

\[\sin A - \sin B = 2\cos\frac{A + B}{2}\sin\frac{A - B}{2}\]

\[\cos A + \cos B = 2\cos\frac{A + B}{2}\cos\frac{A - B}{2}\]

\[\cos A - \cos B = -2\sin\frac{A + B}{2}\sin\frac{A - B}{2}\]

\textbf{证明}：设 \(\alpha = \frac{A + B}{2}\)，\(\beta = \frac{A - B}{2}\)，则 \(A = \alpha + \beta\)，\(B = \alpha - \beta\)。由定理 20.7：

\[\sin A + \sin B = \sin(\alpha + \beta) + \sin(\alpha - \beta) = 2\sin \alpha \cos \beta = 2\sin\frac{A + B}{2}\cos\frac{A - B}{2}\]

其余公式类似证明。\(\blacksquare\)

\hrulefill

辅助角公式是将''正弦+余弦''的线性组合合并为单一三角函数的核心工具。本节证明了对任意 \(a, b\) 不全为零，\(a\sin\theta + b\cos\theta = \sqrt{a^2+b^2}\sin(\theta + \varphi)\)，其中 \(\varphi\) 由 \(\cos\varphi = \frac{a}{\sqrt{a^2+b^2}}\) 和 \(\sin\varphi = \frac{b}{\sqrt{a^2+b^2}}\) 确定。这一公式的几何意义是：平面上点 \((a, b)\) 的极坐标表示。辅助角公式在求解三角方程 \(a\sin x + b\cos x = c\) 时尤为有用------将左边合并为单一正弦函数后，方程变为 \(\sin(x + \varphi) = \frac{c}{\sqrt{a^2+b^2}}\)，解的存在性条件为 \(|\frac{c}{\sqrt{a^2+b^2}}| \leq 1\)。辅助角公式也是研究简谐振动叠加的数学基础：两个同频率的简谐振动之和仍为同频率的简谐振动，其振幅和相位由辅助角公式确定。（辅助角公式）设 \(a, b\) 不全为零，则：

\[a\sin \theta + b\cos \theta = \sqrt{a^2 + b^2}\sin(\theta + \varphi)\]

其中 \(\varphi\) 满足 \(\cos \varphi = \frac{a}{\sqrt{a^2 + b^2}}\)，\(\sin \varphi = \frac{b}{\sqrt{a^2 + b^2}}\)。

\textbf{证明}：令 \(R = \sqrt{a^2 + b^2} > 0\)。定义角 \(\varphi\) 使得 \(\cos \varphi = \frac{a}{R}\)，\(\sin \varphi = \frac{b}{R}\)。这样的 \(\varphi\) 存在，因为 \(\cos^2 \varphi + \sin^2 \varphi = \frac{a^2 + b^2}{R^2} = 1\)。

\[a\sin \theta + b\cos \theta = R\left(\frac{a}{R}\sin \theta + \frac{b}{R}\cos \theta\right) = R(\cos \varphi \sin \theta + \sin \varphi \cos \theta) = R\sin(\theta + \varphi)\]

\(\blacksquare\)

\textbf{推论 20.3} 利用辅助角公式，\(a\sin \theta + b\cos \theta\) 的取值范围为 \([-\sqrt{a^2 + b^2}, \sqrt{a^2 + b^2}]\)。

\textbf{证明}：由辅助角公式，\(a\sin\theta + b\cos\theta = \sqrt{a^2+b^2}\sin(\theta+\varphi)\)。因 \(|\sin(\theta+\varphi)| \leq 1\)，故取值范围为 \([-\sqrt{a^2+b^2}, \sqrt{a^2+b^2}]\)。\(\blacksquare\)

\hrulefill

正弦定理和余弦定理是三角学应用于三角形度量的两大核心工具，它们将三角形的边长与角度之间的几何关系转化为代数等式。正弦定理 \(\frac{a}{\sin A} = \frac{b}{\sin B} = \frac{c}{\sin C} = 2R\) 揭示了三角形中边长与对角正弦之比为常数（等于外接圆直径），其证明利用外接圆中直径所对圆周角为直角的性质，体现了''外接圆''在三角形度量中的核心地位。余弦定理 \(c^2 = a^2 + b^2 - 2ab\cos C\) 是勾股定理的推广：当 \(C = 90°\) 时退化为 \(c^2 = a^2 + b^2\)。余弦定理的证明可以采用坐标法（将三角形置于坐标系中）或向量法（利用 \(|\vec{c}|^2 = |\vec{a} - \vec{b}|^2\)），这两种方法都体现了''解析几何化''的思想。正弦定理适用于''已知两角一边''或''已知两边及其中一边的对角''的情形，余弦定理适用于''已知两边一夹角''或''已知三边''的情形，两者互补，几乎覆盖了所有三角形解算问题。

\textbf{定理 20.10}（正弦定理）在 \(\triangle ABC\) 中，设角 \(A, B, C\) 所对的边分别为 \(a, b, c\)，外接圆半径为 \(R\)，则：

\[\frac{a}{\sin A} = \frac{b}{\sin B} = \frac{c}{\sin C} = 2R\]

\textbf{证明}：设 \(\triangle ABC\) 的外接圆为 \(\odot O\)，半径为 \(R\)。

\textbf{情形一}：\(A\) 为锐角。作直径 \(BD\)，则 \(\angle BCD = 90°\)（直径所对的圆周角）。在 \(\triangle BCD\) 中：

\[\sin D = \frac{BC}{BD} = \frac{a}{2R}\]

由同弧上的圆周角相等，\(\angle D = \angle A\)（因为 \(\angle D\) 和 \(\angle A\) 都是弧 \(BC\) 所对的圆周角）。因此：

\[\sin A = \frac{a}{2R} \implies \frac{a}{\sin A} = 2R\]

\textbf{情形二}：\(A\) 为钝角。作直径 \(BD\)，则 \(\angle D = 180° - A\)，\(\sin D = \sin(180° - A) = \sin A\)。类似可得 \(\frac{a}{\sin A} = 2R\)。

\textbf{情形三}：\(A = 90°\)。则 \(a\) 为外接圆的直径，\(\sin A = 1\)，\(\frac{a}{\sin A} = a = 2R\)。

对 \(B, C\) 同理。\(\blacksquare\)

\subsubsection{20.6.2 余弦定理}\label{ux4f59ux5f26ux5b9aux7406}

\textbf{定理 20.11}（余弦定理）在 \(\triangle ABC\) 中，设角 \(A, B, C\) 所对的边分别为 \(a, b, c\)，则：

\[c^2 = a^2 + b^2 - 2ab\cos C\]

\[a^2 = b^2 + c^2 - 2bc\cos A\]

\[b^2 = a^2 + c^2 - 2ac\cos B\]

\textbf{证明（坐标法）}：以 \(C\) 为原点，\(CB\) 所在直线为 \(x\) 轴建立坐标系。则：

\begin{itemize}
\tightlist
\item
  \(C = (0, 0)\)，\(B = (a, 0)\)，\(A = (b\cos C, b\sin C)\)
\end{itemize}

\[c^2 = |AB|^2 = (b\cos C - a)^2 + (b\sin C)^2\]

\[= b^2\cos^2 C - 2ab\cos C + a^2 + b^2\sin^2 C\]

\[= a^2 + b^2(\cos^2 C + \sin^2 C) - 2ab\cos C = a^2 + b^2 - 2ab\cos C\]

\(\blacksquare\)

\textbf{推论 20.4}（勾股定理）当 \(C = 90°\) 时，\(\cos C = 0\)，余弦定理退化为 \(c^2 = a^2 + b^2\)。

\hrulefill

反三角函数是三角函数在其限制定义域上的反函数，是连接三角学与函数理论的重要桥梁。由于三角函数是周期函数，在整个实数域上不是单射，因此必须限制定义域才能定义反函数。本节定义了反正弦 \(\arcsin x\)（定义域 \([-1, 1]\)，值域 \([-\frac{\pi}{2}, \frac{\pi}{2}]\)）、反余弦 \(\arccos x\)（定义域 \([-1, 1]\)，值域 \([0, \pi]\)）和反正切 \(\arctan x\)（定义域 \(\mathbb{R}\)，值域 \((-\frac{\pi}{2}, \frac{\pi}{2})\)），并证明了它们的基本性质（单调性、奇偶性）。恒等式 \(\arcsin x + \arccos x = \frac{\pi}{2}\) 是 \(\sin(\frac{\pi}{2} - \alpha) = \cos\alpha\) 在反函数层面的体现，它表明反正弦和反余弦是''互补''的。反三角函数在积分学中频繁出现------例如 \(\int \frac{dx}{\sqrt{1-x^2}} = \arcsin x + C\) 和 \(\int \frac{dx}{1+x^2} = \arctan x + C\) 是最基本的积分公式之一。

\textbf{定义 20.1}（反正弦函数）正弦函数 \(y = \sin x\) 在 \(\left[-\frac{\pi}{2}, \frac{\pi}{2}\right]\) 上的反函数称为\textbf{反正弦函数}，记作 \(y = \arcsin x\)。

\textbf{性质}： 1. \textbf{定义域}：\([-1, 1]\)；\textbf{值域}：\(\left[-\frac{\pi}{2}, \frac{\pi}{2}\right]\) 2. \textbf{单调性}：严格递增 3. \textbf{奇偶性}：奇函数

\textbf{证明奇偶性}：设 \(y = \arcsin x\)，则 \(\sin y = x\)，\(y \in \left[-\frac{\pi}{2}, \frac{\pi}{2}\right]\)。\(\sin(-y) = -\sin y = -x\)，\(-y = \arcsin(-x)\)，故 \(\arcsin(-x) = -\arcsin x\)。\(\blacksquare\)

\subsubsection{20.7.2 反余弦函数}\label{ux53cdux4f59ux5f26ux51fdux6570}

\textbf{定义 20.2}（反余弦函数）余弦函数 \(y = \cos x\) 在 \([0, \pi]\) 上的反函数称为\textbf{反余弦函数}，记作 \(y = \arccos x\)。

\textbf{性质}： 1. \textbf{定义域}：\([-1, 1]\)；\textbf{值域}：\([0, \pi]\) 2. \textbf{单调性}：严格递减

\textbf{定理 20.12} \(\arcsin x + \arccos x = \frac{\pi}{2}\)（\(x \in [-1, 1]\)）。

\textbf{证明}：设 \(\alpha = \arcsin x\)，则 \(\sin \alpha = x\)，\(\alpha \in \left[-\frac{\pi}{2}, \frac{\pi}{2}\right]\)。

\[\cos\left(\frac{\pi}{2} - \alpha\right) = \sin \alpha = x\]

又 \(\frac{\pi}{2} - \alpha \in [0, \pi]\)，故 \(\arccos x = \frac{\pi}{2} - \alpha = \frac{\pi}{2} - \arcsin x\)。\(\blacksquare\)

\subsubsection{20.7.3 反正切函数}\label{ux53cdux6b63ux5207ux51fdux6570}

\textbf{定义 20.3}（反正切函数）正切函数 \(y = \tan x\) 在 \(\left(-\frac{\pi}{2}, \frac{\pi}{2}\right)\) 上的反函数称为\textbf{反正切函数}，记作 \(y = \arctan x\)。

\textbf{性质}： 1. \textbf{定义域}：\(\mathbb{R}\)；\textbf{值域}：\(\left(-\frac{\pi}{2}, \frac{\pi}{2}\right)\) 2. \textbf{单调性}：严格递增 3. \textbf{奇偶性}：奇函数

\hrulefill

三角恒等式的系统化梳理将本章中分散的公式整合为一个有机的整体。从同角基本关系（勾股关系、商数关系）出发，经过和差角公式的推导，再依次得到二倍角公式、半角公式、积化和差与和差化积公式，最终汇聚到辅助角公式和正弦定理、余弦定理。这一推导链条展示了三角恒等式体系的内在逻辑------所有公式都可以从余弦的差角公式 \(\cos(\alpha - \beta) = \cos\alpha\cos\beta + \sin\alpha\sin\beta\) 出发，通过代数运算和诱导公式推导出来。掌握这一推导链条比记忆所有公式更为重要，因为它揭示了公式之间的逻辑关系，使得在遗忘某个公式时可以快速从基础知识重新推导。欧拉公式 \(e^{i\theta} = \cos\theta + i\sin\theta\) 提供了理解三角恒等式的另一个视角------所有三角恒等式都可以从指数函数的性质 \(e^{i(\alpha+\beta)} = e^{i\alpha}e^{i\beta}\) 简洁地推导出来。

\subsubsection{20.8.1 万能公式}\label{ux4e07ux80fdux516cux5f0f}

\textbf{定理 20.13}（万能公式）设 \(t = \tan\frac{\alpha}{2}\)（\(\alpha \neq \pi + 2k\pi\)），则：

\[\sin \alpha = \frac{2t}{1 + t^2}, \quad \cos \alpha = \frac{1 - t^2}{1 + t^2}, \quad \tan \alpha = \frac{2t}{1 - t^2}\]

\textbf{证明}：由半角公式，令 \(t = \tan\frac{\alpha}{2}\)，则

\[\sin \alpha = 2\sin\frac{\alpha}{2}\cos\frac{\alpha}{2} = \frac{2\sin\frac{\alpha}{2}\cos\frac{\alpha}{2}}{\sin^2\frac{\alpha}{2} + \cos^2\frac{\alpha}{2}} = \frac{2\tan\frac{\alpha}{2}}{1 + \tan^2\frac{\alpha}{2}} = \frac{2t}{1 + t^2}\]

\[\cos \alpha = \cos^2\frac{\alpha}{2} - \sin^2\frac{\alpha}{2} = \frac{\cos^2\frac{\alpha}{2} - \sin^2\frac{\alpha}{2}}{\cos^2\frac{\alpha}{2} + \sin^2\frac{\alpha}{2}} = \frac{1 - \tan^2\frac{\alpha}{2}}{1 + \tan^2\frac{\alpha}{2}} = \frac{1 - t^2}{1 + t^2}\]

\[\tan \alpha = \frac{\sin \alpha}{\cos \alpha} = \frac{\frac{2t}{1+t^2}}{\frac{1-t^2}{1+t^2}} = \frac{2t}{1 - t^2}\]

\(\blacksquare\)

\subsubsection{20.8.2 三角方程}\label{ux4e09ux89d2ux65b9ux7a0b}

\textbf{定义 20.4}（三角方程）含有三角函数且未知数出现在三角函数内的方程称为\textbf{三角方程}。

\textbf{定理 20.14}（基本三角方程的通解）

\begin{enumerate}
\def\labelenumi{\arabic{enumi}.}
\item
  \(\sin x = a\)（\(|a| \leq 1\)）的通解：\(x = \arcsin a + 2k\pi\) 或 \(x = \pi - \arcsin a + 2k\pi\)（\(k \in \mathbb{Z}\)），简记为 \(x = (-1)^n \arcsin a + n\pi\)（\(n \in \mathbb{Z}\)）。
\item
  \(\cos x = a\)（\(|a| \leq 1\)）的通解：\(x = \pm\arccos a + 2k\pi\)（\(k \in \mathbb{Z}\)）。
\item
  \(\tan x = a\) 的通解：\(x = \arctan a + k\pi\)（\(k \in \mathbb{Z}\)）。
\end{enumerate}

\textbf{证明}：(1) 设 \(x_0 = \arcsin a\)，即 \(\sin x_0 = a\)，\(x_0 \in \left[-\frac{\pi}{2}, \frac{\pi}{2}\right]\)。\(\sin x = a = \sin x_0\) 的充要条件是 \(x = x_0 + 2k\pi\) 或 \(x = \pi - x_0 + 2k\pi\)。\(\blacksquare\)

\hrulefill

\subsection{20.9 三角形面积公式}\label{ux4e09ux89d2ux5f62ux9762ux79efux516cux5f0f}

\textbf{定理 20.15} 在 \(\triangle ABC\) 中，设角 \(A, B, C\) 所对的边分别为 \(a, b, c\)，则：

\[S = \frac{1}{2}ab\sin C = \frac{1}{2}bc\sin A = \frac{1}{2}ac\sin B\]

\textbf{证明}：以 \(C\) 为原点，\(CB\) 为 \(x\) 轴正方向建立坐标系。则 \(B = (a, 0)\)，\(A = (b\cos C, b\sin C)\)。

\(\triangle ABC\) 的面积 = 底 \(\times\) 高 \(\times \frac{1}{2} = a \cdot |b\sin C| \cdot \frac{1}{2} = \frac{1}{2}ab|\sin C|\)。

由于 \(C \in (0, \pi)\)，\(\sin C > 0\)，故 \(S = \frac{1}{2}ab\sin C\)。\(\blacksquare\)

\textbf{推论 20.5}（海伦公式）设 \(s = \frac{a + b + c}{2}\)，则

\[S = \sqrt{s(s-a)(s-b)(s-c)}\]

\textbf{证明}：由余弦定理，\(\cos C = \frac{a^2 + b^2 - c^2}{2ab}\)。

\[\sin^2 C = 1 - \cos^2 C = \frac{4a^2b^2 - (a^2 + b^2 - c^2)^2}{4a^2b^2}\]

分子因式分解：

\[4a^2b^2 - (a^2 + b^2 - c^2)^2 = [2ab + (a^2 + b^2 - c^2)][2ab - (a^2 + b^2 - c^2)]\]

\[= [(a+b)^2 - c^2][c^2 - (a-b)^2] = (a+b+c)(a+b-c)(c+a-b)(c-a+b)\]

\[= 2s \cdot 2(s-c) \cdot 2(s-b) \cdot 2(s-a) = 16s(s-a)(s-b)(s-c)\]

所以

\[S = \frac{1}{2}ab\sin C = \frac{1}{2}ab \cdot \frac{\sqrt{16s(s-a)(s-b)(s-c)}}{2ab} = \sqrt{s(s-a)(s-b)(s-c)}\]

\(\blacksquare\)

\hrulefill

\subsection{20.10 反三角函数的综合性质}\label{ux53cdux4e09ux89d2ux51fdux6570ux7684ux7efcux5408ux6027ux8d28}

反三角函数（\(\arcsin x\)、\(\arccos x\)、\(\arctan x\) 等）是三角函数的反函数，但由于三角函数不是一一映射（它们是周期函数），必须通过限制定义域来构造反函数。\(\arcsin x\) 的定义域为 \([-1, 1]\)，值域为 \([-\frac{\pi}{2}, \frac{\pi}{2}]\)；\(\arccos x\) 的定义域为 \([-1, 1]\)，值域为 \([0, \pi]\)；\(\arctan x\) 的定义域为 \(\mathbb{R}\)，值域为 \((-\frac{\pi}{2}, \frac{\pi}{2})\)。反三角函数在微积分中具有重要地位------它们的导数公式相对简洁，且是积分表中常见积分结果的原函数（如 \(\int \frac{1}{\sqrt{1-x^2}} dx = \arcsin x + C\)）。反三角函数的复合关系也值得关注：\(\arcsin(\sin x)\) 并不等于 \(x\)（除非 \(x\) 在 \(\arcsin\) 的值域内），这一现象反映了反函数定义中''限制定义域''的必要性。反三角函数还自然地出现在几何问题中------例如，直角三角形中锐角的大小可以用反三角函数表示斜边与直角边之比。本节将系统讨论反三角函数的定义域与值域、函数图像、导数（预备知识）以及基本恒等式。

\subsubsection{20.10.1 反三角函数的导数（预备知识）}\label{ux53cdux4e09ux89d2ux51fdux6570ux7684ux5bfcux6570ux9884ux5907ux77e5ux8bc6}

\textbf{定理 20.16} 反三角函数的导数公式：

\begin{enumerate}
\def\labelenumi{\arabic{enumi}.}
\tightlist
\item
  \((\arcsin x)' = \frac{1}{\sqrt{1 - x^2}}\)，\(x \in (-1, 1)\)
\item
  \((\arccos x)' = -\frac{1}{\sqrt{1 - x^2}}\)，\(x \in (-1, 1)\)
\item
  \((\arctan x)' = \frac{1}{1 + x^2}\)，\(x \in \mathbb{R}\)
\end{enumerate}

\textbf{证明}：利用反函数求导法则 \((f^{-1})'(y) = \frac{1}{f'(f^{-1}(y))}\)。

\begin{enumerate}
\def\labelenumi{(\arabic{enumi})}
\item
  设 \(y = \arcsin x\)，则 \(x = \sin y\)，\(y \in [-\frac{\pi}{2}, \frac{\pi}{2}]\)。求导得 \(\frac{dx}{dy} = \cos y\)。由 \(\cos^2 y = 1 - \sin^2 y = 1 - x^2\)，且 \(y \in [-\frac{\pi}{2}, \frac{\pi}{2}]\) 时 \(\cos y \geq 0\)，故 \(\cos y = \sqrt{1 - x^2}\)。因此 \((\arcsin x)' = \frac{dy}{dx} = \frac{1}{dx/dy} = \frac{1}{\cos y} = \frac{1}{\sqrt{1 - x^2}}\)。
\item
  设 \(y = \arccos x\)，则 \(x = \cos y\)，\(y \in [0, \pi]\)。\(\frac{dx}{dy} = -\sin y\)。由 \(\sin^2 y = 1 - \cos^2 y = 1 - x^2\)，且 \(y \in [0, \pi]\) 时 \(\sin y \geq 0\)，故 \(\sin y = \sqrt{1 - x^2}\)。因此 \((\arccos x)' = \frac{1}{-\sin y} = -\frac{1}{\sqrt{1 - x^2}}\)。
\item
  设 \(y = \arctan x\)，则 \(x = \tan y\)，\(y \in (-\frac{\pi}{2}, \frac{\pi}{2})\)。\(\frac{dx}{dy} = \sec^2 y = 1 + \tan^2 y = 1 + x^2\)。因此 \((\arctan x)' = \frac{1}{1 + x^2}\)。\(\blacksquare\)
\end{enumerate}

\subsubsection{20.10.2 反三角函数的恒等式}\label{ux53cdux4e09ux89d2ux51fdux6570ux7684ux6052ux7b49ux5f0f}

\textbf{定理 20.17} 对任意 \(x \in \mathbb{R}\)：

\[\arctan x + \arctan\frac{1}{x} = \begin{cases} \frac{\pi}{2}, & x > 0 \\ -\frac{\pi}{2}, & x < 0 \end{cases}\]

\textbf{证明}：设 \(\alpha = \arctan x\)，\(\beta = \arctan\frac{1}{x}\)。则 \(\tan\alpha = x\)，\(\tan\beta = \frac{1}{x}\)。

\[\tan(\alpha + \beta) = \frac{x + \frac{1}{x}}{1 - x \cdot \frac{1}{x}} = \frac{x + \frac{1}{x}}{0}\]

分母为零，说明 \(\alpha + \beta = \pm\frac{\pi}{2}\)。当 \(x > 0\) 时，\(\alpha \in (0, \frac{\pi}{2})\)，\(\beta \in (0, \frac{\pi}{2})\)，故 \(\alpha + \beta = \frac{\pi}{2}\)。当 \(x < 0\) 时类似。\(\blacksquare\)

\subsection{20.11 三角函数的级数表示（预备知识）}\label{ux4e09ux89d2ux51fdux6570ux7684ux7ea7ux6570ux8868ux793aux9884ux5907ux77e5ux8bc6}

幂级数展开是理解三角函数深层结构的钥匙。泰勒展开（见 Ch 30，推论 23.2）将 \(\sin x\) 和 \(\cos x\) 表示为无穷级数：\(\sin x = x - \frac{x^3}{3!} + \frac{x^5}{5!} - \cdots\)，\(\cos x = 1 - \frac{x^2}{2!} + \frac{x^4}{4!} - \cdots\)。这两个级数对任意实数 \(x\) 都收敛（收敛半径 \(R = \infty\)），这意味着三角函数实际上是''解析函数''------在每一点附近都可以用幂级数表示。从级数表示可以立即看出几个重要事实：\(\sin x\) 是奇函数（只有奇次幂）、\(\cos x\) 是偶函数（只有偶次幂）；当 \(x\) 很小时，\(\sin x \approx x\) 和 \(\cos x \approx 1 - \frac{x^2}{2}\) 就是级数展开的前几项，这为等价无穷小（定理 22.19）提供了严格的级数论证。更为深刻的是，如果将 \(\sin x\) 和 \(\cos x\) 的级数代入欧拉公式 \(e^{ix} = \cos x + i\sin x\)，会发现 \(e^{ix}\) 的泰勒展开恰好是 \(\sum_{n=0}^{\infty} \frac{(ix)^n}{n!}\)，按实部和虚部分离后分别得到 \(\cos x\) 和 \(\sin x\) 的级数------这给出了欧拉公式的一个纯代数证明。本节作为预备知识，介绍三角函数的幂级数展开形式，严格收敛性讨论将在 Ch 30 中完成。

\subsubsection{20.11.1 正弦函数和余弦函数的幂级数}\label{ux6b63ux5f26ux51fdux6570ux548cux4f59ux5f26ux51fdux6570ux7684ux5e42ux7ea7ux6570}

由泰勒展开（Ch 30，推论 23.2），我们有以下重要展开式：

\textbf{定理 20.18} 对任意 \(x \in \mathbb{R}\)：

\[\sin x = \sum_{k=0}^{\infty} \frac{(-1)^k}{(2k+1)!}x^{2k+1} = x - \frac{x^3}{3!} + \frac{x^5}{5!} - \frac{x^7}{7!} + \cdots\]

\[\cos x = \sum_{k=0}^{\infty} \frac{(-1)^k}{(2k)!}x^{2k} = 1 - \frac{x^2}{2!} + \frac{x^4}{4!} - \frac{x^6}{6!} + \cdots\]

这两个级数对所有实数 \(x\) 都绝对收敛。

\textbf{证明}：由泰勒公式，\(\sin x = \sum_{k=0}^{n} \frac{(\sin)^{(k)}(0)}{k!}x^k + R_n(x)\)。

先计算 \(\sin x\) 在 \(x = 0\) 处的各阶导数。由 \((\sin x)' = \cos x\)，\((\cos x)' = -\sin x\)，逐阶求导得：

\[(\sin x)^{(0)}\big|_{x=0} = \sin 0 = 0, \quad (\sin x)^{(1)}\big|_{x=0} = \cos 0 = 1\]

\[(\sin x)^{(2)}\big|_{x=0} = -\sin 0 = 0, \quad (\sin x)^{(3)}\big|_{x=0} = -\cos 0 = -1\]

一般地，\((\sin x)^{(k)}\big|_{x=0} = \sin\left(\frac{k\pi}{2}\right)\)。当 \(k = 2m\)（偶数）时为 \(0\)，当 \(k = 2m+1\)（奇数）时为 \((-1)^m\)。

余项 \(R_n(x) = \frac{\sin^{(n+1)}(\xi)}{(n+1)!}x^{n+1}\)。由于 \(|\sin^{(n+1)}(\xi)| \leq 1\)，\(|R_n(x)| \leq \frac{|x|^{n+1}}{(n+1)!} \to 0\)。\(\blacksquare\)

\subsubsection{20.11.2 欧拉公式}\label{ux6b27ux62c9ux516cux5f0f-1}

\textbf{定理 20.19}（欧拉公式）对任意实数 \(x\)：

\[e^{ix} = \cos x + i\sin x\]

其中 \(i = \sqrt{-1}\) 是虚数单位。

\textbf{证明}：将 \(e^z\) 的幂级数展开中的 \(z\) 替换为 \(ix\)：

\[e^{ix} = \sum_{n=0}^{\infty} \frac{(ix)^n}{n!} = \sum_{k=0}^{\infty} \frac{(ix)^{2k}}{(2k)!} + \sum_{k=0}^{\infty} \frac{(ix)^{2k+1}}{(2k+1)!}\]

\[= \sum_{k=0}^{\infty} \frac{(-1)^k x^{2k}}{(2k)!} + i\sum_{k=0}^{\infty} \frac{(-1)^k x^{2k+1}}{(2k+1)!} = \cos x + i\sin x\]

\(\blacksquare\)

\textbf{推论 20.6}（欧拉恒等式）令 \(x = \pi\)：\(e^{i\pi} + 1 = 0\)。

\textbf{证明}：\(e^{i\pi} = \cos\pi + i\sin\pi = -1 + 0 = -1\)。\(\blacksquare\)

\textbf{推论 20.7}（棣莫弗公式）对任意实数 \(\theta\) 和正整数 \(n\)：

\[(\cos\theta + i\sin\theta)^n = \cos(n\theta) + i\sin(n\theta)\]

\textbf{证明}：由欧拉公式，\(\cos\theta + i\sin\theta = e^{i\theta}\)，故 \((\cos\theta + i\sin\theta)^n = e^{in\theta} = \cos(n\theta) + i\sin(n\theta)\)。\(\blacksquare\)

\hrulefill

\textbf{本章展望}：本章建立了三角恒等变换的完整体系和反三角函数。在下一章（Ch 28）中，我们将转向数列的研究。数列是定义在正整数集上的特殊函数，其离散性质使其在金融、计算机科学和物理学中有广泛应用。我们将系统研究等差数列和等比数列，并介绍数学归纳法这一重要的证明方法。

\hrulefill

\hrulefill

\hrulefill

\hrulefill

\hrulefill

\newpage

\chapter{03-函数与微积分/02-解析几何/01-直线与圆.md}\label{ux51fdux6570ux4e0eux5faeux79efux520602-ux89e3ux6790ux51e0ux4f5501-ux76f4ux7ebfux4e0eux5706.md}

\begin{quote}
\textbf{前置知识}

本章建立在以下前序章节的基础之上：

\begin{itemize}
\tightlist
\item
  \textbf{Ch 23（坐标系与函数概念）}：平面直角坐标系的建立、两点间距离公式、中点公式是本章的几何基础。所有解析几何的讨论均以坐标系为起点。
\item
  \textbf{Ch 14（方程与方程组）}：直线方程本质上是一次方程 \(Ax + By + C = 0\)；圆的方程展开后涉及二元二次方程。方程的代数理论（如韦达定理、判别式）在分析直线与圆的交点问题时不可或缺。
\item
  \textbf{Ch 19.4（圆的几何性质）}：垂径定理、圆周角定理、切线性质、圆幂定理等经典几何结论，为解析化处理圆的问题提供了直觉与验证工具。
\item
  \textbf{Ch 30（导数与微分）}：导数给出了曲线在某点处切线的斜率，这一概念将在本章末尾用于推导圆的切线方程，并为后续章节中圆锥曲线的切线研究奠定基础。
\item
  \textbf{Ch 21（向量）}：向量的数量积、向量的平行与垂直判定，是处理直线位置关系的有力工具。
\end{itemize}
\end{quote}

\hrulefill

\section{第52章：直线}\label{ux7b2c52ux7ae0ux76f4ux7ebf}

直线的方程是解析几何的起点------它是将几何对象（直线）用代数语言（方程）描述的第一个范例。本节从直线的斜率概念出发，系统建立了直线的五种方程形式：点斜式、两点式、斜截式、截距式和一般式。每种形式都有其适用场景：点斜式最适合已知一点和斜率的情况；两点式直接由两点坐标生成；斜截式 \(y = kx + b\) 最为直观，直接显示斜率和 \(y\) 截距；一般式 \(Ax + By + C = 0\) 则是最通用的形式，可以表示包括垂直于 \(x\) 轴的直线在内的所有直线。直线的方程与其一次函数表示之间有着深刻联系------\(y = kx + b\) 中的 \(k\) 既是斜率也是导数，这一联系在微积分中将被充分展开。本节还将讨论直线方程与一次函数之间的关系，以及如何通过直线方程判断直线的位置特征。

\subsection{25.1.1 直线方程的建立}\label{ux76f4ux7ebfux65b9ux7a0bux7684ux5efaux7acb}

在 Ch 23 中，我们建立了平面直角坐标系，并证明了两点间距离公式和中点公式。现在我们要回答一个基本问题：\textbf{如何用方程描述一条直线？}

\textbf{定义 25.1（直线的斜率）.} 设直线 \(l\) 不垂直于 \(x\) 轴，\(P_1(x_1, y_1)\) 和 \(P_2(x_2, y_2)\)（\(x_1 \neq x_2\)）是 \(l\) 上任意两点，则比值

\[k = \frac{y_2 - y_1}{x_2 - x_1}\]

称为直线 \(l\) 的\textbf{斜率}。若直线垂直于 \(x\) 轴，则称斜率不存在。

\textbf{命题 25.1（斜率的唯一性）.} 直线 \(l\) 的斜率 \(k\) 与 \(l\) 上所取的两点无关。

\textbf{证明}： 设 \(P_1(x_1, y_1), P_2(x_2, y_2)\) 和 \(P_3(x_3, y_3), P_4(x_4, y_4)\) 分别是 \(l\) 上两对不同的点。由于 \(P_1, P_2, P_3\) 共线，过 \(P_1\) 作 \(x\) 轴的平行线，过 \(P_2\) 作 \(y\) 轴的平行线，交于点 \(Q_1\)；过 \(P_1\) 作 \(x\) 轴的平行线，过 \(P_3\) 作 \(y\) 轴的平行线，交于点 \(Q_2\)。由相似三角形（Ch 19.3），对应边成比例，即

\[\frac{y_3 - y_1}{x_3 - x_1} = \frac{y_2 - y_1}{x_2 - x_1}\]

同理可证 \(\frac{y_4 - y_3}{x_4 - x_3} = \frac{y_2 - y_1}{x_2 - x_1}\)。故斜率与所取点无关。\(\blacksquare\)

\textbf{定理 25.1（点斜式方程）.} 过点 \(P_0(x_0, y_0)\)、斜率为 \(k\) 的直线 \(l\) 的方程为

\[y - y_0 = k(x - x_0)\]

\textbf{证明}： 设 \(P(x, y)\) 为 \(l\) 上任意一点。若 \(P = P_0\)，则 \(x = x_0, y = y_0\)，等式显然成立。若 \(P \neq P_0\)，由斜率的定义：

\[k = \frac{y - y_0}{x - x_0}\]

即 \(y - y_0 = k(x - x_0)\)。

反过来，若点 \(P(x, y)\) 满足 \(y - y_0 = k(x - x_0)\)，则当 \(x \neq x_0\) 时，\(\frac{y - y_0}{x - x_0} = k\)，说明 \(P\) 在过 \(P_0\) 且斜率为 \(k\) 的直线上；当 \(x = x_0\) 时，\(y = y_0\)，即 \(P = P_0\)，也在直线上。

因此方程 \(y - y_0 = k(x - x_0)\) 恰好描述了过 \(P_0\) 且斜率为 \(k\) 的直线。\(\blacksquare\)

\textbf{定理 25.2（两点式方程）.} 过点 \(P_1(x_1, y_1)\) 和 \(P_2(x_2, y_2)\)（\(x_1 \neq x_2\) 且 \(y_1 \neq y_2\)）的直线方程为

\[\frac{y - y_1}{y_2 - y_1} = \frac{x - x_1}{x_2 - x_1}\]

\textbf{证明}： 由定理 25.1，过 \(P_1\) 且斜率为 \(k = \frac{y_2 - y_1}{x_2 - x_1}\) 的直线方程为

\[y - y_1 = \frac{y_2 - y_1}{x_2 - x_1}(x - x_1)\]

当 \(y_2 \neq y_1\) 时，两边除以 \(y_2 - y_1\) 即得两点式。\(\blacksquare\)

\textbf{定理 25.3（截距式方程）.} 在 \(x\) 轴和 \(y\) 轴上的截距分别为 \(a\) 和 \(b\)（\(a \neq 0, b \neq 0\)）的直线方程为

\[\frac{x}{a} + \frac{y}{b} = 1\]

\textbf{证明}： 该直线过点 \((a, 0)\) 和 \((0, b)\)。由两点式：

\[\frac{y - 0}{b - 0} = \frac{x - a}{0 - a}\]

\[\frac{y}{b} = \frac{a - x}{a} = 1 - \frac{x}{a}\]

\[\frac{x}{a} + \frac{y}{b} = 1\]

反过来，\((a, 0)\) 和 \((0, b)\) 均满足此方程，且两点唯一确定一条直线。\(\blacksquare\)

\textbf{定理 25.4（一般式方程）.} 平面上的任何直线都可以表示为

\[Ax + By + C = 0\]

其中 \(A, B\) 不全为零。反之，满足此方程（\(A, B\) 不全为零）的点的集合是一条直线。

\textbf{证明}： 若直线不垂直于 \(x\) 轴，其方程可写为 \(y - y_0 = k(x - x_0)\)，即 \(kx - y + (y_0 - kx_0) = 0\)，取 \(A = k, B = -1, C = y_0 - kx_0\) 即可。

若直线垂直于 \(x\) 轴，方程为 \(x = x_0\)，即 \(1 \cdot x + 0 \cdot y + (-x_0) = 0\)，取 \(A = 1, B = 0, C = -x_0\)。

反之，若 \(B \neq 0\)，则 \(y = -\frac{A}{B}x - \frac{C}{B}\)，斜率为 \(-\frac{A}{B}\)，是一条直线。若 \(B = 0\)，则 \(A \neq 0\)，方程变为 \(x = -\frac{C}{A}\)，是垂直于 \(x\) 轴的直线。\(\blacksquare\)

\textbf{命题 25.2（一般式的法向量）.} 直线 \(Ax + By + C = 0\) 的一个法向量为 \(\vec{n} = (A, B)\)。

\textbf{证明}： 设 \(P_1(x_1, y_1)\) 和 \(P_2(x_2, y_2)\) 是直线上的两点，则 \(Ax_1 + By_1 + C = 0\) 且 \(Ax_2 + By_2 + C = 0\)。两式相减得

\[A(x_2 - x_1) + B(y_2 - y_1) = 0\]

即 \(\vec{n} \cdot \overrightarrow{P_1P_2} = 0\)，故 \(\vec{n} = (A, B)\) 垂直于直线的方向。\(\blacksquare\)

\subsection{25.1.2 直线的倾斜角}\label{ux76f4ux7ebfux7684ux503eux659cux89d2}

\textbf{定义 25.2（倾斜角）.} 直线 \(l\) 向上的方向与 \(x\) 轴正方向之间所成的最小正角 \(\alpha\)（\(0 \leq \alpha < \pi\)）称为直线 \(l\) 的\textbf{倾斜角}。

\textbf{命题 25.3.} 设直线的倾斜角为 \(\alpha\)，斜率为 \(k\)（当 \(\alpha \neq \frac{\pi}{2}\) 时），则

\[k = \tan\alpha\]

\textbf{证明}： 在直线上取两点 \(P_1(x_1, y_1)\) 和 \(P_2(x_2, y_2)\)，不妨设 \(x_2 > x_1\)。作直角三角形使斜边沿直线方向，则水平直角边长为 \(x_2 - x_1 > 0\)，竖直直角边长为 \(y_2 - y_1\)。由三角函数的定义（Ch 27），\(\tan\alpha = \frac{y_2 - y_1}{x_2 - x_1} = k\)。\(\blacksquare\)

\hrulefill

\subsection{25.2 两直线的位置关系}\label{ux4e24ux76f4ux7ebfux7684ux4f4dux7f6eux5173ux7cfb}

两条直线在平面上的位置关系只有三种可能：相交（有一个公共点）、平行（无公共点）和重合（有无穷多个公共点）。解析几何将这三种几何关系转化为代数条件------设 \(l_1: A_1x + B_1y + C_1 = 0\)，\(l_2: A_2x + B_2y + C_2 = 0\)，则位置关系由系数矩阵的秩决定：当 \(\frac{A_1}{A_2} \neq \frac{B_1}{B_2}\) 时两直线相交；当 \(\frac{A_1}{A_2} = \frac{B_1}{B_2} \neq \frac{C_1}{C_2}\) 时两直线平行；当 \(\frac{A_1}{A_2} = \frac{B_1}{B_2} = \frac{C_1}{C_2}\) 时两直线重合。当两直线用斜截式 \(y = k_1x + b_1\) 和 \(y = k_2x + b_2\) 表示时，判定条件更为简洁：\(k_1 \neq k_2\) 时相交，\(k_1 = k_2\) 且 \(b_1 \neq b_2\) 时平行，\(k_1 = k_2\) 且 \(b_1 = b_2\) 时重合。垂直条件 \(k_1k_2 = -1\)（或 \(A_1A_2 + B_1B_2 = 0\)）是两条直线相交的特殊情形，它在向量语言中对应着方向向量的数量积为零。两直线夹角的正切公式 \(\tan\theta = \left|\frac{k_1 - k_2}{1 + k_1k_2}\right|\) 则量化了夹角与斜率之间的关系，这一公式在后续的坐标变换和旋转变换中具有重要应用。

\subsection{25.2.1 平行与垂直的判定}\label{ux5e73ux884cux4e0eux5782ux76f4ux7684ux5224ux5b9a}

设直线 \(l_1: A_1 x + B_1 y + C_1 = 0\)，\(l_2: A_2 x + B_2 y + C_2 = 0\)。

\textbf{定理 25.5（平行条件）.} \(l_1 \parallel l_2\) 当且仅当 \(A_1 B_2 - A_2 B_1 = 0\) 且 \(A_1 C_2 - A_2 C_1 \neq 0\)（或 \(B_1 C_2 - B_2 C_1 \neq 0\)）。

\textbf{证明}：

\((\Rightarrow)\) 若 \(l_1 \parallel l_2\)，则它们的法向量 \(\vec{n}_1 = (A_1, B_1)\) 与 \(\vec{n}_2 = (A_2, B_2)\) 平行，即存在非零常数 \(\lambda\) 使得 \(\vec{n}_1 = \lambda \vec{n}_2\)，即 \(A_1 = \lambda A_2\)，\(B_1 = \lambda B_2\)。于是 \(A_1 B_2 - A_2 B_1 = \lambda A_2 B_2 - A_2 \lambda B_2 = 0\)。

若 \(A_1 C_2 - A_2 C_1 = 0\)，则 \(\lambda A_2 C_2 - A_2 C_1 = 0\)，即 \(A_2(\lambda C_2 - C_1) = 0\)。若 \(A_2 \neq 0\)，则 \(C_1 = \lambda C_2\)，此时 \(l_1: \lambda(A_2 x + B_2 y + C_2) = 0\)，即 \(l_1 = l_2\)，与 \(l_1 \parallel l_2\) 矛盾。若 \(A_2 = 0\)，则 \(B_2 \neq 0\)，由 \(B_1 C_2 - B_2 C_1 = \lambda B_2 C_2 - B_2 C_1 = B_2(\lambda C_2 - C_1) \neq 0\) 可得矛盾。因此 \(A_1 C_2 - A_2 C_1 \neq 0\)。

\((\Leftarrow)\) 若 \(A_1 B_2 - A_2 B_1 = 0\) 且 \(A_1 C_2 - A_2 C_1 \neq 0\)，则由 \(A_1 B_2 = A_2 B_1\) 可设 \(A_1 = \lambda A_2\)，\(B_1 = \lambda B_2\)（若 \(A_2 = B_2 = 0\) 则 \(l_2\) 不是直线）。假设 \(l_1\) 与 \(l_2\) 相交于点 \((x_0, y_0)\)，则 \(A_1 x_0 + B_1 y_0 + C_1 = 0\) 且 \(A_2 x_0 + B_2 y_0 + C_2 = 0\)。由第一式得 \(\lambda(A_2 x_0 + B_2 y_0) + C_1 = 0\)，由第二式得 \(A_2 x_0 + B_2 y_0 = -C_2\)，代入得 \(-\lambda C_2 + C_1 = 0\)，即 \(A_1 C_2 - A_2 C_1 = \lambda A_2 C_2 - A_2 C_1 = A_2(\lambda C_2 - C_1) = 0\)，与 \(A_1 C_2 - A_2 C_1 \neq 0\) 矛盾。故 \(l_1 \parallel l_2\)。\(\blacksquare\)

\textbf{定理 25.6（垂直条件）.} \(l_1 \perp l_2\) 当且仅当 \(A_1 A_2 + B_1 B_2 = 0\)。

\textbf{证明}： \(l_1\) 的法向量为 \(\vec{n}_1 = (A_1, B_1)\)，\(l_2\) 的法向量为 \(\vec{n}_2 = (A_2, B_2)\)。

\(l_1 \perp l_2\) \(\iff\) \(\vec{n}_1 \perp \vec{n}_2\)（因为法向量与直线垂直，两条直线垂直等价于它们的法向量垂直）\(\iff\) \(\vec{n}_1 \cdot \vec{n}_2 = 0\) \(\iff\) \(A_1 A_2 + B_1 B_2 = 0\)。\(\blacksquare\)

\textbf{推论 25.1.} 若 \(l_1\) 的斜率为 \(k_1\)，\(l_2\) 的斜率为 \(k_2\)（均存在），则：

\begin{itemize}
\tightlist
\item
  \(l_1 \parallel l_2 \iff k_1 = k_2\)
\item
  \(l_1 \perp l_2 \iff k_1 k_2 = -1\)
\end{itemize}

\textbf{证明}： 设 \(l_1: y = k_1 x + b_1\)，\(l_2: y = k_2 x + b_2\)，即 \(l_1: k_1 x - y + b_1 = 0\)，\(l_2: k_2 x - y + b_2 = 0\)。

平行条件：\(k_1 \cdot (-1) - k_2 \cdot (-1) = 0\)，即 \(k_1 = k_2\)（还需 \(b_1 \neq b_2\) 保证不重合）。

垂直条件：\(k_1 k_2 + (-1)(-1) = 0\)，即 \(k_1 k_2 + 1 = 0\)，即 \(k_1 k_2 = -1\)。\(\blacksquare\)

\subsection{25.2.2 两直线的交点}\label{ux4e24ux76f4ux7ebfux7684ux4ea4ux70b9}

\textbf{定理 25.7.} 设 \(l_1: A_1 x + B_1 y + C_1 = 0\)，\(l_2: A_2 x + B_2 y + C_2 = 0\)，若 \(l_1\) 与 \(l_2\) 不平行（即 \(A_1 B_2 - A_2 B_1 \neq 0\)），则它们的交点坐标为

\[x = \frac{B_1 C_2 - B_2 C_1}{A_1 B_2 - A_2 B_1}, \quad y = \frac{C_1 A_2 - C_2 A_1}{A_1 B_2 - A_2 B_1}\]

\textbf{证明}： 联立方程组：

\[\begin{cases} A_1 x + B_1 y = -C_1 \\ A_2 x + B_2 y = -C_2 \end{cases}\]

系数行列式 \(D = A_1 B_2 - A_2 B_1 \neq 0\)。由克拉默法则（Ch 14.3）：

\[x = \frac{\begin{vmatrix} -C_1 & B_1 \\ -C_2 & B_2 \end{vmatrix}}{D} = \frac{-C_1 B_2 + C_2 B_1}{D} = \frac{B_1 C_2 - B_2 C_1}{A_1 B_2 - A_2 B_1}\]

\[y = \frac{\begin{vmatrix} A_1 & -C_1 \\ A_2 & -C_2 \end{vmatrix}}{D} = \frac{-A_1 C_2 + A_2 C_1}{D} = \frac{C_1 A_2 - C_2 A_1}{A_1 B_2 - A_2 B_1}\]

\(\blacksquare\)

\subsection{25.2.3 两直线的夹角}\label{ux4e24ux76f4ux7ebfux7684ux5939ux89d2}

\textbf{定义 25.3.} 两条直线 \(l_1, l_2\) 所成的\textbf{锐角} \(\theta\)（\(0 < \theta \leq \frac{\pi}{2}\)）称为两直线的夹角。

\textbf{定理 25.8.} 设 \(l_1\) 的斜率为 \(k_1\)，\(l_2\) 的斜率为 \(k_2\)（均存在且 \(k_1 k_2 \neq -1\)），则两直线夹角 \(\theta\) 满足

\[\tan\theta = \left|\frac{k_1 - k_2}{1 + k_1 k_2}\right|\]

\textbf{证明}： 设 \(l_1\) 的倾斜角为 \(\alpha_1\)，\(l_2\) 的倾斜角为 \(\alpha_2\)，则 \(k_1 = \tan\alpha_1\)，\(k_2 = \tan\alpha_2\)。两直线的夹角 \(\theta = |\alpha_1 - \alpha_2|\)（取锐角）。

由两角差的正切公式（Ch 27.1）：

\[\tan(\alpha_1 - \alpha_2) = \frac{\tan\alpha_1 - \tan\alpha_2}{1 + \tan\alpha_1 \tan\alpha_2} = \frac{k_1 - k_2}{1 + k_1 k_2}\]

取绝对值得 \(\tan\theta = \left|\frac{k_1 - k_2}{1 + k_1 k_2}\right|\)。\(\blacksquare\)

\textbf{推论 25.2.} 用一般式的法向量表示：设 \(\vec{n}_1 = (A_1, B_1)\)，\(\vec{n}_2 = (A_2, B_2)\)，则

\[\cos\theta = \frac{|A_1 A_2 + B_1 B_2|}{\sqrt{A_1^2 + B_1^2} \cdot \sqrt{A_2^2 + B_2^2}}\]

\textbf{证明}： 两法向量的夹角 \(\varphi\) 满足 \(\cos\varphi = \frac{\vec{n}_1 \cdot \vec{n}_2}{|\vec{n}_1||\vec{n}_2|}\)。由于直线夹角 \(\theta\) 等于法向量夹角或其补角，取绝对值即得 \(\cos\theta = |\cos\varphi|\)。\(\blacksquare\)

\hrulefill

\subsection{25.3 距离公式}\label{ux8dddux79bbux516cux5f0f}

距离是几何学中最基本的度量概念。在解析几何中，距离公式将几何中的''远近''概念转化为代数表达式。本节的核心结果是点到直线的距离公式 \(d = \frac{|Ax_0 + By_0 + C|}{\sqrt{A^2 + B^2}}\)，其推导方法有多种------可以利用向量法（将点到直线的距离转化为法向量方向上的投影长度），也可以利用几何法（通过构造垂足联立方程组）。两条平行直线之间的距离公式 \(d = \frac{|C_1 - C_2|}{\sqrt{A^2 + B^2}}\) 是点到直线距离公式的直接推论。两点间距离公式 \(d = \sqrt{(x_2 - x_1)^2 + (y_2 - y_1)^2}\) 是解析几何最基本的公式，其本质是勾股定理在坐标系中的体现。这些距离公式在后续的圆的方程（距离等于常数的点的轨迹）、圆锥曲线的定义（距离之和/差为常数）以及点到平面距离的空间推广中反复出现，是解析几何计算的基础工具。

\subsection{25.3.1 点到直线的距离}\label{ux70b9ux5230ux76f4ux7ebfux7684ux8dddux79bb}

\textbf{定理 25.9（点到直线距离公式）.} 点 \(P(x_0, y_0)\) 到直线 \(l: Ax + By + C = 0\) 的距离为

\[d = \frac{|Ax_0 + By_0 + C|}{\sqrt{A^2 + B^2}}\]

\textbf{证明}： 设直线 \(l\) 的法向量为 \(\vec{n} = (A, B)\)。在 \(l\) 上取一点 \(Q(x_1, y_1)\)，则 \(Ax_1 + By_1 + C = 0\)。

向量 \(\overrightarrow{QP} = (x_0 - x_1, y_0 - y_1)\)。\(P\) 到 \(l\) 的距离等于 \(\overrightarrow{QP}\) 在法向量方向上投影的绝对值（因为从 \(Q\) 到直线 \(l\) 的垂线方向即为法向量方向）：

\[d = \frac{|\overrightarrow{QP} \cdot \vec{n}|}{|\vec{n}|} = \frac{|A(x_0 - x_1) + B(y_0 - y_1)|}{\sqrt{A^2 + B^2}}\]

由于 \(Ax_1 + By_1 = -C\)，代入得：

\[d = \frac{|Ax_0 + By_0 - (Ax_1 + By_1)|}{\sqrt{A^2 + B^2}} = \frac{|Ax_0 + By_0 + C|}{\sqrt{A^2 + B^2}}\]

\(\blacksquare\)

\textbf{推论 25.3（两平行线间的距离）.} 设平行线 \(l_1: Ax + By + C_1 = 0\)，\(l_2: Ax + By + C_2 = 0\)（\(A, B\) 相同保证平行，\(C_1 \neq C_2\) 保证两线不重合），则它们之间的距离为

\[d = \frac{|C_1 - C_2|}{\sqrt{A^2 + B^2}}\]

\textbf{证明}： 在 \(l_1\) 上取点 \(P_0(x_0, y_0)\)，则 \(Ax_0 + By_0 + C_1 = 0\)。\(P_0\) 到 \(l_2\) 的距离为

\[d = \frac{|Ax_0 + By_0 + C_2|}{\sqrt{A^2 + B^2}} = \frac{|-C_1 + C_2|}{\sqrt{A^2 + B^2}} = \frac{|C_1 - C_2|}{\sqrt{A^2 + B^2}}\]

由于平行线间距离处处相等，此即所求。\(\blacksquare\)

\hrulefill

\section{第53章：圆的方程}\label{ux7b2c53ux7ae0ux5706ux7684ux65b9ux7a0b}

圆是平面上到定点（圆心）距离等于定长（半径）的点的轨迹------这一定义在解析几何中直接转化为方程：\((x - a)^2 + (y - b)^2 = r^2\)。圆的标准方程明确地显示了圆心 \((a, b)\) 和半径 \(r\)，而圆的一般方程 \(x^2 + y^2 + Dx + Ey + F = 0\) 则通过配方法 \((x + \frac{D}{2})^2 + (y + \frac{E}{2})^2 = \frac{D^2 + E^2 - 4F}{4}\) 揭示其几何意义------当 \(D^2 + E^2 - 4F > 0\) 时表示一个圆，等于 \(0\) 时退化为一个点，小于 \(0\) 时无实数点（虚圆）。圆的方程与距离公式之间有着天然的联系：圆的标准方程本质上就是两点间距离公式的平方形式。圆作为圆锥曲线中离心率 \(e = 0\) 的特殊情形，是椭圆当两个焦点重合时的极限情形。本节还将讨论如何通过三个不共线的点确定一个圆------这等价于求解三元一次方程组，其几何意义是三角形的外接圆。

\subsection{25.4.1 圆的标准方程}\label{ux5706ux7684ux6807ux51c6ux65b9ux7a0b}

在 Ch 19.4 中，我们从几何角度研究了圆的性质（垂径定理、圆周角定理等）。现在我们在坐标系中用方程来描述圆。

\textbf{定义 25.4（圆的解析定义）.} 平面内到定点 \(C(a, b)\) 的距离等于定长 \(r\)（\(r > 0\)）的点的轨迹称为\textbf{圆}，记作 \(\odot(C, r)\)。定点 \(C\) 称为\textbf{圆心}，定长 \(r\) 称为\textbf{半径}。

\textbf{定理 25.10（圆的标准方程）.} 圆心为 \(C(a, b)\)、半径为 \(r\) 的圆的方程为

\[(x - a)^2 + (y - b)^2 = r^2\]

\textbf{证明}： 设 \(P(x, y)\) 为圆上任意一点，由定义 \(|PC| = r\)。由两点间距离公式（Ch 23.1）：

\[\sqrt{(x - a)^2 + (y - b)^2} = r\]

两边平方即得 \((x - a)^2 + (y - b)^2 = r^2\)。

反过来，若 \(P(x, y)\) 满足此方程，则 \(|PC| = \sqrt{(x-a)^2 + (y-b)^2} = r\)，\(P\) 在圆上。\(\blacksquare\)

\textbf{推论 25.4.} 圆心在原点、半径为 \(r\) 的圆的方程为 \(x^2 + y^2 = r^2\)。

\textbf{证明}：圆心在原点、半径为 \(r\) 的圆是到原点距离为 \(r\) 的点的轨迹，即 \(\sqrt{x^2+y^2} = r\)，两边平方得 \(x^2+y^2 = r^2\)。\(\blacksquare\)

\subsection{25.4.2 圆的一般方程}\label{ux5706ux7684ux4e00ux822cux65b9ux7a0b}

\textbf{定理 25.11（圆的一般方程）.} 圆心为 \((a, b)\)、半径为 \(r\) 的圆的方程可写为

\[x^2 + y^2 + Dx + Ey + F = 0\]

其中 \(D = -2a\)，\(E = -2b\)，\(F = a^2 + b^2 - r^2\)。反之，方程 \(x^2 + y^2 + Dx + Ey + F = 0\) 当 \(D^2 + E^2 - 4F > 0\) 时表示一个圆，圆心为 \(\left(-\frac{D}{2}, -\frac{E}{2}\right)\)，半径为 \(r = \frac{1}{2}\sqrt{D^2 + E^2 - 4F}\)。

\textbf{证明}：（正向）将标准方程 \((x - a)^2 + (y - b)^2 = r^2\) 展开：

\[x^2 - 2ax + a^2 + y^2 - 2by + b^2 = r^2\]

\[x^2 + y^2 + (-2a)x + (-2b)y + (a^2 + b^2 - r^2) = 0\]

令 \(D = -2a\)，\(E = -2b\)，\(F = a^2 + b^2 - r^2\)，即得一般方程。

（反向）给定 \(x^2 + y^2 + Dx + Ey + F = 0\)，配方得：

\[\left(x + \frac{D}{2}\right)^2 + \left(y + \frac{E}{2}\right)^2 = \frac{D^2 + E^2 - 4F}{4}\]

当 \(D^2 + E^2 - 4F > 0\) 时，右边为正，表示圆心 \(\left(-\frac{D}{2}, -\frac{E}{2}\right)\)、半径 \(\frac{1}{2}\sqrt{D^2 + E^2 - 4F}\) 的圆。

当 \(D^2 + E^2 - 4F = 0\) 时，表示一个点 \(\left(-\frac{D}{2}, -\frac{E}{2}\right)\)（退化的圆，半径为 \(0\)）。

当 \(D^2 + E^2 - 4F < 0\) 时，不表示任何实图形（空集）。\(\blacksquare\)

\textbf{证明}： 由定理 25.11 的推导，\(r^2 = \frac{D^2 + E^2 - 4F}{4}\)。\(r^2 > 0\) 时为圆，\(r^2 = 0\) 时为点，\(r^2 < 0\) 时无实数解。\(\blacksquare\)

\subsection{25.4.3 圆的切线方程}\label{ux5706ux7684ux5207ux7ebfux65b9ux7a0b}

\textbf{定理 25.12（过圆上一点的切线方程）.} 过圆 \((x - a)^2 + (y - b)^2 = r^2\) 上一点 \(P_0(x_0, y_0)\) 的切线方程为

\[(x_0 - a)(x - a) + (y_0 - b)(y - b) = r^2\]

\textbf{证明}： 切线在 \(P_0\) 处与圆相切，因此切线垂直于半径 \(\overrightarrow{CP_0}\)，其中 \(C(a, b)\) 为圆心。

半径方向向量 \(\overrightarrow{CP_0} = (x_0 - a, y_0 - b)\)。切线过 \(P_0\) 且垂直于 \(\overrightarrow{CP_0}\)，切线上任一点 \(P(x, y)\) 满足 \(\overrightarrow{P_0P} \perp \overrightarrow{CP_0}\)，即

\[(x_0 - a)(x - x_0) + (y_0 - b)(y - y_0) = 0\]

展开：

\[(x_0 - a)x - (x_0 - a)x_0 + (y_0 - b)y - (y_0 - b)y_0 = 0\]

\[(x_0 - a)x + (y_0 - b)y = (x_0 - a)x_0 + (y_0 - b)y_0\]

由于 \(P_0\) 在圆上，\((x_0 - a)^2 + (y_0 - b)^2 = r^2\)。于是

\[(x_0 - a)x_0 + (y_0 - b)y_0 = (x_0 - a)(x_0 - a + a) + (y_0 - b)(y_0 - b + b)\] \[= (x_0 - a)^2 + a(x_0 - a) + (y_0 - b)^2 + b(y_0 - b) = r^2 + a(x_0 - a) + b(y_0 - b)\]

将此代入并移项：

\[(x_0 - a)(x - a) + (y_0 - b)(y - b) = r^2\]

\(\blacksquare\)

\textbf{定理 25.13（圆外一点到圆的切线长）.} 从圆外一点 \(P(x_0, y_0)\) 到圆 \((x - a)^2 + (y - b)^2 = r^2\) 的切线长为

\[l = \sqrt{(x_0 - a)^2 + (y_0 - b)^2 - r^2}\]

\textbf{证明}： 设切点为 \(T\)，圆心为 \(C(a, b)\)。需要证明 \(CT \perp PT\)（半径垂直于过切点的直线）。下面用解析方法独立证明。圆 \((x-a)^2 + (y-b)^2 = r^2\) 在点 \(T(x_1, y_1)\) 处的切线方程为 \((x_1-a)(x-a) + (y_1-b)(y-b) = r^2\)（由定理 25.12 导出：将圆方程视为隐函数并取微分，或利用圆心到切线的距离等于半径）。切线方向向量可取 \(\vec{d} = (y_1-b,\, -(x_1-a))\)（与法向量 \((x_1-a,\, y_1-b)\) 正交）。半径向量 \(\overrightarrow{CT} = (x_1-a,\, y_1-b)\)，点积为零，故 \(CT \perp PT\)。于是 \(\triangle CPT\) 为直角三角形（直角在 \(T\) 处）。

由勾股定理：\(|CP|^2 = |CT|^2 + |PT|^2\)，即

\[(x_0 - a)^2 + (y_0 - b)^2 = r^2 + l^2\]

\[l = \sqrt{(x_0 - a)^2 + (y_0 - b)^2 - r^2}\]

其中需 \((x_0 - a)^2 + (y_0 - b)^2 > r^2\)（\(P\) 在圆外）。\(\blacksquare\)

\hrulefill

\section{25.5 直线与圆的位置关系}\label{ux76f4ux7ebfux4e0eux5706ux7684ux4f4dux7f6eux5173ux7cfb}

直线与圆的位置关系有三种：相离（无交点）、相切（一个交点）和相交（两个交点）。解析几何将这一几何问题转化为代数问题------联立直线方程和圆的方程，消去一个变量后得到一个一元二次方程，其判别式 \(\Delta\) 决定了位置关系：\(\Delta > 0\) 时相交，\(\Delta = 0\) 时相切，\(\Delta < 0\) 时相离。这一方法体现了''解析几何=几何+代数''的核心思想------将几何位置关系转化为代数方程的判别式问题。圆心到直线的距离 \(d\) 与半径 \(r\) 的比较提供了另一种判定方式：\(d < r\) 时相交，\(d = r\) 时相切，\(d > r\) 时相离。相切情形下，切线方程的推导是本节的重点------过圆上一点 \((x_0, y_0)\) 的切线方程为 \((x_0 - a)(x - a) + (y_0 - b)(y - b) = r^2\)，其本质是圆上点与圆心连线的垂线。过圆外一点引两条切线，切线长相等，这一性质与圆幂定理有着密切联系。弦长公式 \(L = 2\sqrt{r^2 - d^2}\) 将几何中的弦长转化为代数计算，是解析几何计算能力的典型体现。

\section{第54章：直线与圆的关系}\label{ux7b2c54ux7ae0ux76f4ux7ebfux4e0eux5706ux7684ux5173ux7cfb}

\textbf{定理 25.14.} 设圆 \(C\) 的方程为 \((x - a)^2 + (y - b)^2 = r^2\)，直线 \(l\) 的方程为 \(Ax + By + C_0 = 0\)。圆心 \(C(a, b)\) 到直线 \(l\) 的距离为

\[d = \frac{|Aa + Bb + C_0|}{\sqrt{A^2 + B^2}}\]

则：

\begin{itemize}
\tightlist
\item
  \(d > r\) \(\iff\) 直线与圆\textbf{相离}（无公共点）；
\item
  \(d = r\) \(\iff\) 直线与圆\textbf{相切}（恰有一个公共点）；
\item
  \(d < r\) \(\iff\) 直线与圆\textbf{相交}（有两个公共点）。
\end{itemize}

\textbf{证明}： 过圆心 \(C\) 作直线 \(l\) 的垂线，垂足为 \(H\)，则 \(|CH| = d\)。

\textbf{情形 1：} \(d > r\)。设 \(P\) 为 \(l\) 上任意一点，\(\triangle CHP\) 中 \(\angle CHP = 90^\circ\)，由勾股定理 \(|CP|^2 = d^2 + |HP|^2 \geq d^2 > r^2\)，故 \(|CP| > r\)，\(P\) 不在圆上。因此直线与圆无公共点。

\textbf{情形 2：} \(d = r\)。此时 \(H\) 在圆上（\(|CH| = r\)），且 \(H\) 在直线 \(l\) 上，故 \(H\) 是公共点。对 \(l\) 上其他点 \(P \neq H\)，同理 \(|CP| > r\)，故 \(P\) 不在圆上。\(H\) 是唯一的公共点，即相切。

\textbf{情形 3：} \(d < r\)。在直线 \(l\) 上取 \(H\) 两侧的点 \(P_1, P_2\)，使 \(|HP_1| = |HP_2| = \sqrt{r^2 - d^2}\)。则

\[|CP_1|^2 = d^2 + (\sqrt{r^2 - d^2})^2 = r^2\]

故 \(|CP_1| = r\)，\(P_1\) 在圆上。同理 \(P_2\) 也在圆上。\(P_1 \neq P_2\)（在 \(H\) 两侧），故有两个交点。\(\blacksquare\)

\subsection{25.5.2 弦长公式}\label{ux5f26ux957fux516cux5f0f}

\textbf{定理 25.15.} 直线与圆相交时，弦长为

\[|AB| = 2\sqrt{r^2 - d^2}\]

其中 \(d\) 为圆心到直线的距离。

\textbf{证明}： 设弦 \(AB\) 的中点为 \(M\)。由圆的对称性（或垂径定理，Ch 19.4），\(CM \perp AB\)，且 \(|AM| = |MB|\)。在直角三角形 \(\triangle CMA\) 中，\(\angle CMA = 90^\circ\)，\(|CA| = r\)，\(|CM| = d\)，由勾股定理：

\[|AM| = \sqrt{r^2 - d^2}\]

故 \(|AB| = 2|AM| = 2\sqrt{r^2 - d^2}\)。\(\blacksquare\)

\subsection{25.5.3 两圆的位置关系}\label{ux4e24ux5706ux7684ux4f4dux7f6eux5173ux7cfb}

\textbf{定理 25.16.} 设 \(\odot C_1\) 的圆心为 \(C_1(a_1, b_1)\)、半径为 \(r_1\)，\(\odot C_2\) 的圆心为 \(C_2(a_2, b_2)\)、半径为 \(r_2\)。记两圆心距离 \(d = |C_1 C_2| = \sqrt{(a_1 - a_2)^2 + (b_1 - b_2)^2}\)。不妨设 \(r_1 \geq r_2\)，则：

\begin{longtable}[]{@{}lll@{}}
\toprule
位置关系 & 条件 & 公共点个数 \\
\midrule
\endhead
\bottomrule
\endlastfoot
外离 & \(d > r_1 + r_2\) & 0 \\
外切 & \(d = r_1 + r_2\) & 1 \\
相交 & \(r_1 - r_2 < d < r_1 + r_2\) & 2 \\
内切 & \(d = r_1 - r_2\)（\(r_1 > r_2\)） & 1 \\
内含 & \(d < r_1 - r_2\)（\(r_1 > r_2\)） & 0 \\
\end{longtable}

\textbf{证明}： 设 \(P\) 为公共点，则 \(|PC_1| = r_1\)，\(|PC_2| = r_2\)。由三角不等式：

\[|r_1 - r_2| \leq |PC_1| - |PC_2| \leq |C_1 C_2| \leq |PC_1| + |PC_2| = r_1 + r_2\]

等等，需要更仔细地分析。由 \(|PC_1| = r_1\)，\(|PC_2| = r_2\) 和三角不等式：

\[|r_1 - r_2| = ||PC_1| - |PC_2|| \leq |C_1 C_2| \leq |PC_1| + |PC_2| = r_1 + r_2\]

因此公共点存在的必要条件是 \(|r_1 - r_2| \leq d \leq r_1 + r_2\)。

下面对五种情形逐一验证：

\textbf{(i) 外离}（\(d > r_1 + r_2\)）：由必要条件 \(d \leq r_1 + r_2\) 不成立，故无公共点。

\textbf{(ii) 外切}（\(d = r_1 + r_2\)）：在 \(C_1 C_2\) 连线上取点 \(P = C_1 + \frac{r_1}{d}\overrightarrow{C_1 C_2}\)，则 \(|PC_1| = r_1\)，\(|PC_2| = d - r_1 = r_2\)，故 \(P\) 是公共点。若另有公共点 \(Q\)，则由三角不等式 \(|C_1 C_2| \leq |QC_1| + |QC_2| = r_1 + r_2 = d\)，等号成立当且仅当 \(Q\) 在线段 \(C_1 C_2\) 上，故 \(Q = P\)。公共点个数为 \(1\)。

\textbf{(iii) 相交}（\(r_1 - r_2 < d < r_1 + r_2\)）：两圆满足 \(|r_1 - r_2| < d < r_1 + r_2\)。用解析方法严格证明。设两圆方程分别为

\[(x - x_1)^2 + (y - y_1)^2 = r_1^2, \quad (x - x_2)^2 + (y - y_2)^2 = r_2^2\]

不妨设圆心 \(C_1(0, 0)\)，\(C_2(d, 0)\)（通过平移旋转），则方程化为

\[x^2 + y^2 = r_1^2, \quad (x - d)^2 + y^2 = r_2^2\]

相减得 \(x^2 + y^2 - [(x-d)^2 + y^2] = r_1^2 - r_2^2\)，即 \(2dx - d^2 = r_1^2 - r_2^2\)，解得

\[x = \frac{r_1^2 - r_2^2 + d^2}{2d}\]

代入第一圆方程：\(y^2 = r_1^2 - x^2\)。计算判别式：

\[r_1^2 - x^2 = \frac{4d^2 r_1^2 - (r_1^2 - r_2^2 + d^2)^2}{4d^2}\]

分子可因式分解为 \([(r_1 + r_2)^2 - d^2][d^2 - (r_1 - r_2)^2] = (r_1 + r_2 + d)(r_1 + r_2 - d)(d + r_1 - r_2)(d - r_1 + r_2)\)。在 \(|r_1 - r_2| < d < r_1 + r_2\) 条件下，四个因子全为正，故 \(y^2 > 0\)。\(y = \pm\sqrt{y^2}\) 给出两个相异的实数解，对应两圆恰有两个公共点。\(\blacksquare\)

\textbf{(iv) 内切}（\(d = r_1 - r_2\)，\(r_1 > r_2\)）：在 \(C_1 C_2\) 连线的延长线上取点 \(P\) 使得 \(|PC_2| = r_2\) 且 \(P\) 在 \(C_1\) 与 \(C_2\) 之间，则 \(|PC_1| = d + r_2 = r_1\)，故 \(P\) 是公共点。类似 (ii) 的论证，公共点唯一。公共点个数为 \(1\)。

\textbf{(v) 内含}（\(d < r_1 - r_2\)，\(r_1 > r_2\)）：对任意点 \(P\)，由三角不等式 \(|PC_1| \leq |PC_2| + d < |PC_2| + (r_1 - r_2)\)。若 \(|PC_1| = r_1\)，则 \(|PC_2| > r_2\)，故 \(P\) 不在 \(\odot C_2\) 上。同理若 \(|PC_2| = r_2\)，则 \(|PC_1| < r_1\)。因此无公共点。\(\blacksquare\)

\subsection{25.5.4 圆幂定理的解析证明}\label{ux5706ux5e42ux5b9aux7406ux7684ux89e3ux6790ux8bc1ux660e}

在 Ch 19.4 中，我们从纯几何角度证明了圆幂定理。现在用解析方法重新证明。

\textbf{定理 25.17（圆幂定理）.} 设圆 \(\odot C\) 的方程为 \((x - a)^2 + (y - b)^2 = r^2\)，\(P(x_0, y_0)\) 为圆外一点。过 \(P\) 作割线交圆于 \(A, B\) 两点，则

\[|PA| \cdot |PB| = |PC|^2 - r^2\]

且此值与割线的选择无关（称为点 \(P\) 关于圆 \(\odot C\) 的\textbf{圆幂}）。

\textbf{证明}： 设过 \(P\) 的直线参数方程为

\[\begin{cases} x = x_0 + t\cos\alpha \\ y = y_0 + t\sin\alpha \end{cases}\]

代入圆的方程：

\[(x_0 + t\cos\alpha - a)^2 + (y_0 + t\sin\alpha - b)^2 = r^2\]

设 \(u = x_0 - a\)，\(v = y_0 - b\)，展开：

\[u^2 + 2ut\cos\alpha + t^2\cos^2\alpha + v^2 + 2vt\sin\alpha + t^2\sin^2\alpha = r^2\]

\[t^2 + 2(u\cos\alpha + v\sin\alpha)t + (u^2 + v^2 - r^2) = 0\]

由韦达定理（Ch 14），两个根 \(t_1, t_2\) 满足：

\[t_1 t_2 = u^2 + v^2 - r^2 = (x_0 - a)^2 + (y_0 - b)^2 - r^2 = |PC|^2 - r^2\]

由参数 \(t\) 的几何意义：直线方程为 \(P = P_0 + t \vec{d}\)，其中 \(\vec{d} = (\cos\alpha, \sin\alpha)\) 为单位方向向量。对参数值 \(t_i\) 对应的点 \(A_i\)，有

\[|P_0 A_i| = |t_i \vec{d}| = |t_i| \cdot |\vec{d}| = |t_i|\]

因此 \(|t_1| = |PA|\)，\(|t_2| = |PB|\)。当 \(P\) 在圆外时，\(t_1, t_2\) 同号（均在 \(P\) 的同侧或异侧，取决于方向），故

\[|PA| \cdot |PB| = |t_1 t_2| = |PC|^2 - r^2\]

（\(P\) 在圆外时 \(|PC|^2 > r^2\)，故 \(t_1 t_2 > 0\)，无需取绝对值。）\(\blacksquare\)

\textbf{推论 25.5（切割线定理）.} 若过 \(P\) 的直线与圆相切于 \(T\)，则 \(|PT|^2 = |PA| \cdot |PB|\)。

\textbf{证明}： 当割线退化为切线时，\(A = B = T\)，\(|PA| \cdot |PB| = |PT|^2 = |PC|^2 - r^2\)（由定理 25.13）。\(\blacksquare\)

\textbf{推论 25.6（割线定理）.} 若过 \(P\) 的两条割线分别交圆于 \(A, B\) 和 \(C, D\)，则 \(|PA| \cdot |PB| = |PC| \cdot |PD|\)。

\textbf{证明}： 两条割线对应的圆幂相同，均为 \(|PC|^2 - r^2\)。\(\blacksquare\)

\subsection{25.5.5 根轴}\label{ux6839ux8f74}

\textbf{定义 25.5.} 两个不同心的圆 \(\odot C_1\) 和 \(\odot C_2\) 的\textbf{根轴}是到两圆的圆幂相等的点的轨迹。

\textbf{定理 25.18.} 设 \(\odot C_1: x^2 + y^2 + D_1 x + E_1 y + F_1 = 0\)，\(\odot C_2: x^2 + y^2 + D_2 x + E_2 y + F_2 = 0\)，则根轴的方程为

\[(D_1 - D_2)x + (E_1 - E_2)y + (F_1 - F_2) = 0\]

\textbf{证明}： 点 \(P(x_0, y_0)\) 关于 \(\odot C_1\) 的圆幂为 \(x_0^2 + y_0^2 + D_1 x_0 + E_1 y_0 + F_1\)（由一般方程直接计算）。关于 \(\odot C_2\) 的圆幂为 \(x_0^2 + y_0^2 + D_2 x_0 + E_2 y_0 + F_2\)。

两圆幂相等：

\[(D_1 - D_2)x_0 + (E_1 - E_2)y_0 + (F_1 - F_2) = 0\]

这是一条直线。\(\blacksquare\)

\hrulefill

\subsection{25.6 圆的切线与导数}\label{ux5706ux7684ux5207ux7ebfux4e0eux5bfcux6570}

圆的切线问题从纯几何学过渡到解析几何，再到微积分，体现了数学概念逐步深化的脉络。在纯几何中，圆的切线定义为''与圆只有一个公共点的直线''；在解析几何中，切线方程可以通过判别式法（\(\Delta = 0\)）或圆心到直线的距离法（\(d = r\)）来推导。然而，导数的引入为切线概念提供了全新的视角------曲线 \(y = f(x)\) 在点 \((x_0, f(x_0))\) 处的切线斜率等于 \(f'(x_0)\)，这一视角将切线从''特殊位置关系的直线''提升为''局部线性逼近''的几何体现。对于圆 \(x^2 + y^2 = r^2\)，通过隐函数求导可得 \(2x + 2y \cdot y' = 0\)，从而 \(y' = -\frac{x}{y}\)，这与几何法得到的切线斜率完全一致。导数方法的美妙之处在于它不依赖于圆的特殊几何性质------它适用于任何可导函数的曲线，因此将切线理论从圆推广到了一般曲线，为微积分中''以直代曲''的思想提供了严格的数学基础。

\subsection{25.6.1 用导数求切线方程}\label{ux7528ux5bfcux6570ux6c42ux5207ux7ebfux65b9ux7a0b}

在 Ch 30 中，我们学习了导数的几何意义：函数 \(y = f(x)\) 在 \(x = x_0\) 处的导数 \(f'(x_0)\) 等于曲线在该点处切线的斜率。我们可以用这一工具来求圆的切线方程。

\textbf{定理 25.19.} 圆 \(x^2 + y^2 = r^2\) 在点 \(P_0(x_0, y_0)\)（\(y_0 \neq 0\)）处的切线斜率为

\[k = -\frac{x_0}{y_0}\]

\textbf{证明}： 由隐函数求导（Ch 30），对方程 \(x^2 + y^2 = r^2\) 两边关于 \(x\) 求导：

\[2x + 2y \frac{dy}{dx} = 0\]

\[\frac{dy}{dx} = -\frac{x}{y}\]

在点 \(P_0(x_0, y_0)\) 处，\(k = \frac{dy}{dx}\big|_{(x_0, y_0)} = -\frac{x_0}{y_0}\)。\(\blacksquare\)

\textbf{验证：} 由定理 25.12，切线方程为 \(x_0 x + y_0 y = r^2\)，即 \(y = -\frac{x_0}{y_0}x + \frac{r^2}{y_0}\)，斜率为 \(-\frac{x_0}{y_0}\)，与导数法一致。\(\blacksquare\)

\textbf{推论 25.7.} 圆 \((x - a)^2 + (y - b)^2 = r^2\) 在点 \(P_0(x_0, y_0)\)（\(y_0 \neq b\)）处的切线斜率为

\[k = -\frac{x_0 - a}{y_0 - b}\]

\textbf{证明}： 令 \(X = x - a\)，\(Y = y - b\)，方程变为 \(X^2 + Y^2 = r^2\)。对 \(X\) 求导：\(2X + 2Y\frac{dY}{dX} = 0\)，\(\frac{dY}{dX} = -\frac{X}{Y}\)。由于 \(\frac{dy}{dx} = \frac{dY}{dX}\)，在点 \((x_0, y_0)\) 处 \(k = -\frac{x_0 - a}{y_0 - b}\)。\(\blacksquare\)

\subsection{25.6.2 切线的存在性}\label{ux5207ux7ebfux7684ux5b58ux5728ux6027}

\textbf{命题 25.4.} 过圆外一点 \(P\) 到圆恰好有两条切线；过圆上一点恰好有一条切线；过圆内一点没有切线。

\textbf{证明}： 设圆为 \((x - a)^2 + (y - b)^2 = r^2\)，\(P(x_0, y_0)\)。

切线过 \(P\) 且与圆相切，等价于 \(P\) 到直线的距离等于 \(r\)。设切线斜率为 \(k\)（存在时），则切线方程为 \(y - y_0 = k(x - x_0)\)，即 \(kx - y + (y_0 - kx_0) = 0\)。圆心 \((a, b)\) 到此直线的距离为

\[d = \frac{|ka - b + y_0 - kx_0|}{\sqrt{k^2 + 1}} = \frac{|k(a - x_0) - (b - y_0)|}{\sqrt{k^2 + 1}} = r\]

两边平方：

\[[k(a - x_0) - (b - y_0)]^2 = r^2(k^2 + 1)\]

设 \(p = a - x_0\)，\(q = b - y_0\)，展开：

\[(kp - q)^2 = r^2(k^2 + 1)\] \[k^2 p^2 - 2kpq + q^2 = r^2 k^2 + r^2\] \[(p^2 - r^2)k^2 - 2pq \cdot k + (q^2 - r^2) = 0\]

关于 \(k\) 的二次方程，判别式为

\[\Delta = 4p^2q^2 - 4(p^2 - r^2)(q^2 - r^2) = 4[p^2q^2 - p^2q^2 + r^2 p^2 + r^2 q^2 - r^4]\] \[= 4r^2(p^2 + q^2 - r^2) = 4r^2[(x_0 - a)^2 + (y_0 - b)^2 - r^2]\]

\begin{itemize}
\tightlist
\item
  若 \(P\) 在圆外：\((x_0 - a)^2 + (y_0 - b)^2 > r^2\)，\(\Delta > 0\)，两个不同的实数解，两条切线。
\item
  若 \(P\) 在圆上：\((x_0 - a)^2 + (y_0 - b)^2 = r^2\)，\(\Delta = 0\)，一个实数解（重根），一条切线。
\item
  若 \(P\) 在圆内：\((x_0 - a)^2 + (y_0 - b)^2 < r^2\)，\(\Delta < 0\)，无实数解，无切线。
\end{itemize}

还需考虑斜率不存在（切线垂直于 \(x\) 轴）的情况，即 \(x = x_0\) 为切线。此时圆心到直线的距离 \(|a - x_0| = r\)，即 \((x_0 - a)^2 = r^2\)，此时 \(x_0 = a \pm r\)，切线方程为 \(x = a + r\) 或 \(x = a - r\)。这种情况已被上述分析覆盖（对应二次方程中 \(p^2 = r^2\) 时退化为一次方程的情形）。\(\blacksquare\)

\hrulefill

\begin{quote}
\textbf{Ch 325 本章展望}

本章建立了直线与圆的解析理论：用一次方程描述直线，用二次方程描述圆。我们证明了直线的各种方程形式（点斜式、两点式、截距式、一般式）之间的等价性，推导了点到直线的距离公式和两直线夹角公式，并系统研究了圆的标准方程、一般方程以及直线与圆的位置关系。特别地，我们利用导数工具验证了圆的切线方程，建立了几何与分析的桥梁。

然而，直线和圆只是平面曲线的''冰山一角''。自然界中存在大量更为复杂的曲线------行星的椭圆轨道、抛射体的抛物线轨迹、冷却塔截面的双曲线------它们有一个共同的名字：\textbf{圆锥曲线}。在下一章中，我们将从统一的几何定义出发，推导椭圆、双曲线和抛物线的标准方程，并揭示它们之间深刻的内在联系。
\end{quote}

\hrulefill

\hrulefill

\hrulefill

\hrulefill

\hrulefill

\newpage

\chapter{03-函数与微积分/02-解析几何/02-圆锥曲线.md}\label{ux51fdux6570ux4e0eux5faeux79efux520602-ux89e3ux6790ux51e0ux4f5502-ux5706ux9525ux66f2ux7ebf.md}

\begin{quote}
\textbf{前置知识}

本章建立在以下前序章节的基础之上：

\begin{itemize}
\tightlist
\item
  \textbf{Ch 325（直线与圆）}：圆是椭圆的特殊情形（离心率 \(e = 0\)）。圆的标准方程推导方法（距离公式 \(\to\) 平方化简）直接推广为椭圆方程的推导。点到直线距离公式在推导抛物线方程时不可或缺。
\item
  \textbf{Ch 19.3（相似）}：椭圆可以视为圆沿一个方向均匀拉伸的结果，这一变换保持了相似结构。相似比在椭圆的长短轴关系中有所体现。
\item
  \textbf{Ch 21（向量）}：焦点向量关系 \(|PF_1| + |PF_2| = 2a\) 等定义中涉及向量的模长。向量法在处理圆锥曲线的焦点弦、准线等问题时提供简洁的表达。
\item
  \textbf{Ch 14（方程与方程组）}：圆锥曲线的方程是二元二次方程，韦达定理在处理直线与圆锥曲线的交点问题（设而不求）中至关重要。
\item
  \textbf{Ch 27（三角恒等变换）}：双曲线的参数方程涉及 \(\sec\theta\) 和 \(\tan\theta\) 的恒等式 \(\sec^2\theta - \tan^2\theta = 1\)，与双曲线方程的结构完美呼应。
\end{itemize}
\end{quote}

\hrulefill

\section{第55章：圆锥曲线的统一定义}\label{ux7b2c55ux7ae0ux5706ux9525ux66f2ux7ebfux7684ux7edfux4e00ux5b9aux4e49}

圆锥曲线是解析几何中最辉煌的成就之一。从古希腊的阿波罗尼奥斯用平面截割圆锥面得到三种曲线（椭圆、抛物线、双曲线），到17世纪笛卡尔和费马将其纳入解析几何的框架，再到开普勒发现行星轨道是椭圆，圆锥曲线的研究贯穿了数学和物理学的发展史。本节给出圆锥曲线的统一定义------到焦点和准线的距离之比等于离心率 \(e\) 的点的轨迹。这一定义的优美之处在于，它用一个参数 \(e\) 统一了三种看似不同的曲线：\(0 < e < 1\) 时为椭圆（封闭曲线），\(e = 1\) 时为抛物线（开口曲线，恰好是椭圆与双曲线的分界），\(e > 1\) 时为双曲线（由两支组成）。特别地，\(e = 0\) 时退化为圆。焦准距 \(p\)（焦点到准线的距离）是另一个关键参数，它决定了曲线在焦点处的''尺度''，在极坐标下的统一方程 \(r = \frac{ep}{1 - e\cos\theta}\) 中起着核心作用。

\subsection{26.1.1 圆锥截线}\label{ux5706ux9525ux622aux7ebf}

圆锥曲线（conic section）的历史可以追溯到古希腊。阿波罗尼奥斯（Apollonius，约前 262--前 190）发现，用平面截割圆锥面时，根据截面与圆锥母线的夹角不同，可以得到三种不同的曲线：椭圆、双曲线和抛物线。

\textbf{定义 26.1（圆锥曲线的统一定义）.} 平面内，到一个定点 \(F\)（\textbf{焦点}）和一条定直线 \(l\)（\textbf{准线}，\(F \notin l\)）的距离之比等于常数 \(e\)（\(e > 0\)）的点的轨迹称为\textbf{圆锥曲线}。即

\[C = \left\{P \;\middle|\; \frac{|PF|}{d(P, l)} = e\right\}\]

其中 \(e\) 称为\textbf{离心率}（eccentricity），\(F\) 称为\textbf{焦点}（focus），\(l\) 称为\textbf{准线}（directrix）。

\begin{itemize}
\tightlist
\item
  当 \(0 < e < 1\) 时，曲线为\textbf{椭圆}；
\item
  当 \(e = 1\) 时，曲线为\textbf{抛物线}；
\item
  当 \(e > 1\) 时，曲线为\textbf{双曲线}。
\end{itemize}

\textbf{注} 当 \(e > 1\) 时，以焦点 \(F\) 和对应准线 \(l\) 定义的圆锥曲线（满足 \(\frac{|PF|}{d(P, l)} = e\)）只给出双曲线的\textbf{一支}（靠近焦点 \(F\) 的那支）。要得到完整的双曲线（两支），需要分别以两个焦点为极点各定义一次，或采用定义 26.4 的两焦点距离差定义。具体而言：以焦点 \(F_1\) 和准线 \(l_1\) 定义给出双曲线的左支，以焦点 \(F_2\) 和准线 \(l_2\) 定义给出双曲线的右支；两支的并集才是完整的双曲线。

圆的特殊情况：当 \(e = 0\) 时，焦点到准线的约束消失，曲线退化为以 \(F\) 为圆心的圆。

\subsection{26.1.2 焦准距}\label{ux7126ux51c6ux8ddd}

焦准距是圆锥曲线理论中一个关键参数，它统一了三种圆锥曲线的焦点-准线描述。焦准距 \(p\) 度量了焦点到其对应准线的垂直距离，它在后续的极坐标方程 \(r = \frac{ep}{1 - e\cos\theta}\) 中直接决定了曲线在焦点处的''尺度''------当 \(\theta = \frac{\pi}{2}\) 时，\(r = ep\)，这正是通径（正焦弦）长度的一半。对于椭圆，\(p = \frac{b^2}{c}\)；对于双曲线，\(p = \frac{b^2}{c}\)；对于抛物线，\(p\) 就是标准方程 \(y^2 = 2px\) 中的参数。值得注意，虽然椭圆和双曲线的焦准距表达式相同（\(p = \frac{b^2}{c}\)），但两者的 \(b^2\) 和 \(c\) 的定义方式不同------椭圆满足 \(b^2 = a^2 - c^2\)，双曲线满足 \(b^2 = c^2 - a^2\)。

焦准距的引入使得圆锥曲线的焦点-准线描述更加完整。在统一定义 \(\frac{|PF|}{d(P, l)} = e\) 中，焦准距 \(p\) 是焦点到准线的距离，它与离心率 \(e\) 和半通径 \(l\)（semi-latus rectum）的关系为 \(l = ep\)。半通径 \(l\) 是过焦点且垂直于长轴（或对称轴）的弦长的一半，是描述圆锥曲线在焦点附近''开口大小''的基本参数。在极坐标下，圆锥曲线的统一方程为 \(r = \frac{l}{1 + e\cos\theta}\) 或 \(r = \frac{ep}{1 - e\cos\theta}\)，其中 \(\theta\) 是从焦点出发的极角。焦准距 \(p\) 的几何意义可以从相似三角形来理解：设焦点 \(F\) 到准线 \(l\) 的垂足为 \(D\)，则 \(|FD| = p\)。对于曲线上任意一点 \(P\)，设 \(P\) 到准线的垂足为 \(Q\)，则 \(\frac{|PF|}{|PQ|} = e\)。当 \(P\) 选在使 \(PF \perp FD\) 的位置时（即通径端点），\(|PF| = ep = l\)，这正是半通径的几何来源。焦准距还在圆锥曲线的光学性质中扮演角色：对于抛物面反射器，平行于轴的光线经反射后汇聚于焦点，焦点到顶点的距离恰好是 \(p/2\)。

\hrulefill

\section{第56章：椭圆}\label{ux7b2c56ux7ae0ux692dux5706}

椭圆是圆锥曲线中最为''温和''的一种------它是封闭的、有界的，且具有丰富的对称性。椭圆的焦点定义（到两焦点距离之和为常数 \(2a\)）揭示了其本质：两个''引力源''共同塑造了椭圆形状。椭圆的标准方程 \(\frac{x^2}{a^2} + \frac{y^2}{b^2} = 1\) 中，\(a\) 为长半轴，\(b\) 为短半轴，\(c = \sqrt{a^2 - b^2}\) 为半焦距，离心率 \(e = \frac{c}{a}\) 度量了椭圆的''扁平''程度------\(e\) 越接近 \(0\)，椭圆越接近圆；\(e\) 越接近 \(1\)，椭圆越扁。椭圆具有丰富的几何性质：任意一点到两焦点的距离之和为常数 \(2a\)；椭圆上任意一点的法线平分该点与两焦点连线的夹角（光学反射性质）；椭圆的面积为 \(\pi ab\)（当 \(a = b\) 时退化为圆的面积公式）。在天文学中，开普勒第一定律（行星轨道是椭圆，太阳位于一个焦点）使椭圆成为描述天体运动的基础曲线。椭圆的参数方程 \(x = a\cos\theta, y = b\sin\theta\) 揭示了椭圆与圆之间的仿射变换关系------椭圆是圆沿 \(y\) 轴方向压缩 \(b/a\) 倍的结果。

\subsection{26.2.1 椭圆的定义}\label{ux692dux5706ux7684ux5b9aux4e49}

\textbf{定义 26.3（椭圆的焦点定义）.} 平面内到两个定点 \(F_1, F_2\)（称为\textbf{焦点}）的距离之和等于常数 \(2a\)（\(2a > |F_1 F_2|\)）的点的轨迹称为\textbf{椭圆}（ellipse）。即

\[E = \{P \mid |PF_1| + |PF_2| = 2a, \; 2a > |F_1 F_2|\}\]

常数 \(2a\) 必须大于两焦点间距 \(|F_1 F_2|\)，否则轨迹不存在或退化。

\textbf{命题 26.1.} 若 \(2a = |F_1 F_2|\)，则轨迹为线段 \(F_1 F_2\)；若 \(2a < |F_1 F_2|\)，则轨迹为空集。

\textbf{证明}： 若 \(2a = |F_1 F_2|\)，由三角不等式 \(|PF_1| + |PF_2| \geq |F_1 F_2|\)，等号成立当且仅当 \(P\) 在线段 \(F_1 F_2\) 上。反之，线段上每一点满足 \(|PF_1| + |PF_2| = |F_1 F_2| = 2a\)。

若 \(2a < |F_1 F_2|\)，由三角不等式 \(|PF_1| + |PF_2| \geq |F_1 F_2| > 2a\)，不存在满足条件的点。\(\blacksquare\)

\subsection{26.2.2 椭圆的标准方程}\label{ux692dux5706ux7684ux6807ux51c6ux65b9ux7a0b}

\textbf{定理 26.1（椭圆的标准方程）.} 设椭圆的焦点为 \(F_1(-c, 0)\)，\(F_2(c, 0)\)（\(c > 0\)），\(2a > 2c\)。令 \(b^2 = a^2 - c^2\)（\(b > 0\)），则椭圆的标准方程为

\[\frac{x^2}{a^2} + \frac{y^2}{b^2} = 1 \quad (a > b > 0)\]

\textbf{证明}： 设 \(P(x, y)\) 为椭圆上任意一点。由定义：

\[\sqrt{(x + c)^2 + y^2} + \sqrt{(x - c)^2 + y^2} = 2a \quad \cdots (*)\]

移项：

\[\sqrt{(x + c)^2 + y^2} = 2a - \sqrt{(x - c)^2 + y^2}\]

两边平方：

\[(x + c)^2 + y^2 = 4a^2 - 4a\sqrt{(x - c)^2 + y^2} + (x - c)^2 + y^2\]

展开并化简：

\[x^2 + 2cx + c^2 = 4a^2 - 4a\sqrt{(x - c)^2 + y^2} + x^2 - 2cx + c^2\]

\[4cx = 4a^2 - 4a\sqrt{(x - c)^2 + y^2}\]

\[a\sqrt{(x - c)^2 + y^2} = a^2 - cx\]

两边再平方：

\[a^2[(x - c)^2 + y^2] = (a^2 - cx)^2\]

\[a^2 x^2 - 2a^2 cx + a^2 c^2 + a^2 y^2 = a^4 - 2a^2 cx + c^2 x^2\]

\[a^2 x^2 + a^2 c^2 + a^2 y^2 = a^4 + c^2 x^2\]

\[(a^2 - c^2)x^2 + a^2 y^2 = a^2(a^2 - c^2)\]

令 \(b^2 = a^2 - c^2 > 0\)（由 \(a > c\)）：

\[b^2 x^2 + a^2 y^2 = a^2 b^2\]

两边除以 \(a^2 b^2\)：

\[\frac{x^2}{a^2} + \frac{y^2}{b^2} = 1\]

最后验证：平方过程中引入了附加条件 \(a^2 - cx \geq 0\)，即 \(x \leq \frac{a^2}{c}\)。由 \(\frac{x^2}{a^2} \leq 1\) 得 \(|x| \leq a < \frac{a^2}{c}\)（因 \(c < a\)），条件自动满足。

还需验证所有满足方程的点都在椭圆上：若 \(\frac{x^2}{a^2} + \frac{y^2}{b^2} = 1\)，则 \(|x| \leq a\)。定义 \(d_1 = \sqrt{(x+c)^2 + y^2}\)，\(d_2 = \sqrt{(x-c)^2 + y^2}\)。由上述推导的逆过程可得 \(d_1 + d_2 = 2a\)（需验证 \(a^2 - cx \geq 0\) 以保证平方的等价性）。\(\blacksquare\)

\subsection{26.2.3 椭圆的基本性质}\label{ux692dux5706ux7684ux57faux672cux6027ux8d28}

\textbf{定理 26.2.} 椭圆 \(\frac{x^2}{a^2} + \frac{y^2}{b^2} = 1\)（\(a > b > 0\)）具有以下性质：

\begin{enumerate}
\def\labelenumi{\arabic{enumi}.}
\tightlist
\item
  \textbf{范围}：\(|x| \leq a\)，\(|y| \leq b\)。
\item
  \textbf{对称性}：关于 \(x\) 轴、\(y\) 轴和原点对称。
\item
  \textbf{顶点}：长轴端点 \(A_1(-a, 0)\)，\(A_2(a, 0)\)；短轴端点 \(B_1(0, -b)\)，\(B_2(0, b)\)。
\item
  \textbf{焦点}：\(F_1(-c, 0)\)，\(F_2(c, 0)\)，其中 \(c = \sqrt{a^2 - b^2}\)。
\item
  \textbf{离心率}：\(e = \frac{c}{a}\)（\(0 < e < 1\)）。
\end{enumerate}

\textbf{证明}：

\textbf{（1）范围：} 由 \(\frac{x^2}{a^2} = 1 - \frac{y^2}{b^2} \leq 1\) 得 \(|x| \leq a\)。同理 \(|y| \leq b\)。

\textbf{（2）对称性：} 将 \((x, y)\) 替换为 \((x, -y)\)：\(\frac{x^2}{a^2} + \frac{(-y)^2}{b^2} = \frac{x^2}{a^2} + \frac{y^2}{b^2} = 1\)，故关于 \(x\) 轴对称。同理替换为 \((-x, y)\) 得关于 \(y\) 轴对称，替换为 \((-x, -y)\) 得关于原点对称。

\textbf{（3）顶点：} 令 \(y = 0\) 得 \(x = \pm a\)；令 \(x = 0\) 得 \(y = \pm b\)。

\textbf{（4）} 由 \(b^2 = a^2 - c^2\) 得 \(c = \sqrt{a^2 - b^2}\)。

\textbf{（5）} \(e = \frac{c}{a}\)。由 \(0 < c < a\) 得 \(0 < e < 1\)。\(\blacksquare\)

\textbf{命题 26.2（离心率的意义）.} 离心率 \(e\) 衡量椭圆的''扁平程度''：\(e\) 越接近 \(0\)，椭圆越接近圆；\(e\) 越接近 \(1\)，椭圆越扁。

\textbf{证明}： \(e = \frac{c}{a} = \frac{\sqrt{a^2 - b^2}}{a} = \sqrt{1 - \frac{b^2}{a^2}}\)。当 \(e \to 0\) 时 \(\frac{b}{a} \to 1\)，即 \(b \to a\)，椭圆趋近于圆 \(x^2 + y^2 = a^2\)。当 \(e \to 1\) 时 \(\frac{b}{a} \to 0\)，椭圆在 \(y\) 方向被极度压缩。\(\blacksquare\)

\subsection{26.2.4 焦点距离公式}\label{ux7126ux70b9ux8dddux79bbux516cux5f0f}

\textbf{定理 26.3（焦点距离公式）.} 椭圆 \(\frac{x^2}{a^2} + \frac{y^2}{b^2} = 1\)（\(a > b > 0\)）上一点 \(P(x_0, y_0)\) 到两个焦点的距离分别为

\[|PF_1| = a + ex_0, \quad |PF_2| = a - ex_0\]

\textbf{证明}： 由椭圆方程得 \(y_0^2 = b^2\left(1 - \frac{x_0^2}{a^2}\right) = \frac{b^2(a^2 - x_0^2)}{a^2}\)。

\[|PF_1|^2 = (x_0 + c)^2 + y_0^2 = (x_0 + c)^2 + \frac{b^2(a^2 - x_0^2)}{a^2}\]

\[= x_0^2 + 2cx_0 + c^2 + \frac{(a^2 - c^2)(a^2 - x_0^2)}{a^2}\]

\[= x_0^2 + 2cx_0 + c^2 + a^2 - c^2 - x_0^2 + \frac{c^2 x_0^2}{a^2}\]

\[= a^2 + 2cx_0 + \frac{c^2 x_0^2}{a^2} = \left(a + \frac{cx_0}{a}\right)^2 = (a + ex_0)^2\]

由于 \(|x_0| \leq a\) 且 \(0 < e < 1\)，有 \(a + ex_0 \geq a - ea = a(1 - e) > 0\)，故 \(|PF_1| = a + ex_0\)。由 \(|PF_1| + |PF_2| = 2a\) 得 \(|PF_2| = a - ex_0\)。\(\blacksquare\)

\subsection{26.2.5 椭圆的准线}\label{ux692dux5706ux7684ux51c6ux7ebf}

\textbf{定理 26.4.} 椭圆 \(\frac{x^2}{a^2} + \frac{y^2}{b^2} = 1\) 对应于焦点 \(F_2(c, 0)\) 的准线方程为 \(l: x = \frac{a^2}{c} = \frac{a}{e}\)。

\textbf{证明}： 需验证 \(\frac{|PF_2|}{d(P, l)} = e\)。由定理 26.3，\(|PF_2| = a - ex_0\)。\(P\) 到准线 \(x = \frac{a^2}{c}\) 的距离为

\[d(P, l) = \frac{a^2}{c} - x_0 = \frac{a}{e} - x_0\]

（因 \(x_0 \leq a < \frac{a^2}{c}\)，故此值为正。）

\[\frac{|PF_2|}{d(P, l)} = \frac{a - ex_0}{\frac{a}{e} - x_0} = \frac{e(a - ex_0)}{a - ex_0} = e\]

\(\blacksquare\)

\textbf{推论 26.1.} 椭圆有两条准线 \(x = \pm\frac{a^2}{c}\)，分别对应两个焦点。

\subsection{26.2.6 椭圆与圆的关系}\label{ux692dux5706ux4e0eux5706ux7684ux5173ux7cfb}

\textbf{命题 26.3（椭圆作为拉伸的圆）.} 将单位圆 \(x^2 + y^2 = 1\) 沿 \(x\) 轴方向拉伸 \(a\) 倍、沿 \(y\) 轴方向拉伸 \(b\) 倍，得到椭圆 \(\frac{x^2}{a^2} + \frac{y^2}{b^2} = 1\)。

\textbf{证明}： 设变换 \(T: (X, Y) \mapsto (aX, bY)\)。若 \((X, Y)\) 在单位圆上，即 \(X^2 + Y^2 = 1\)，令 \(x = aX\)，\(y = bY\)，则 \(\frac{x^2}{a^2} + \frac{y^2}{b^2} = 1\)。\(\blacksquare\)

\textbf{推论 26.2.} 椭圆的面积为 \(S = \pi ab\)。

\textbf{证明}： 由命题 26.3，椭圆 \(\frac{x^2}{a^2} + \frac{y^2}{b^2} = 1\) 由单位圆经变换 \(T: (X, Y) \mapsto (aX, bY)\) 得到。利用参数积分法：椭圆可参数化为 \(x = a\cos t\)，\(y = b\sin t\)（\(t \in [0, 2\pi]\)），其面积为

\[S = \left|\int_0^{2\pi} x\,dy\right| = \left|\int_0^{2\pi} a\cos t \cdot b\cos t\,dt\right| = ab\int_0^{2\pi} \cos^2 t\,dt = ab \cdot \pi = \pi ab\]

其中利用了 \(\int_0^{2\pi}\cos^2 t\,dt = \pi\)。\(\blacksquare\)

\textbf{注}：此结果也可通过多重积分的变量替换定理（雅可比行列式）得到：变换 \(T\) 的雅可比行列式为 \(\det(J) = ab\)，单位圆面积为 \(\pi\)，故椭圆面积为 \(\pi ab\)。但该方法需要高等数学知识，此处的参数积分法更为初等。

\subsection{26.2.7 椭圆的切线方程}\label{ux692dux5706ux7684ux5207ux7ebfux65b9ux7a0b}

\textbf{定理 26.5（椭圆切线方程）.} 椭圆 \(\frac{x^2}{a^2} + \frac{y^2}{b^2} = 1\) 在点 \(P_0(x_0, y_0)\)（在椭圆上）处的切线方程为

\[\frac{x_0 x}{a^2} + \frac{y_0 y}{b^2} = 1\]

\textbf{证明}： 由隐函数求导（Ch 30），对椭圆方程两边关于 \(x\) 求导：

\[\frac{2x}{a^2} + \frac{2y}{b^2}\frac{dy}{dx} = 0 \implies \frac{dy}{dx} = -\frac{b^2 x}{a^2 y}\]

在 \(P_0(x_0, y_0)\) 处，切线斜率为 \(k = -\frac{b^2 x_0}{a^2 y_0}\)（\(y_0 \neq 0\)）。切线方程为

\[y - y_0 = -\frac{b^2 x_0}{a^2 y_0}(x - x_0)\]

\[a^2 y_0(y - y_0) = -b^2 x_0(x - x_0)\]

\[b^2 x_0 x + a^2 y_0 y = b^2 x_0^2 + a^2 y_0^2\]

由 \(P_0\) 在椭圆上，\(b^2 x_0^2 + a^2 y_0^2 = a^2 b^2\)，故

\[\frac{x_0 x}{a^2} + \frac{y_0 y}{b^2} = 1\]

当 \(y_0 = 0\) 时，\(x_0 = \pm a\)，切线为 \(x = \pm a\)，即 \(\frac{x_0 x}{a^2} = 1\)，公式仍然成立。\(\blacksquare\)

\textbf{推论 26.3.} 过椭圆外一点 \(Q(x_1, y_1)\) 的切线方程可通过切点条件 \(P_0\) 在椭圆上和切线过 \(Q\) 联立求解。

\textbf{说明}：设切点为 \(P_0(x_0, y_0)\)，则切线方程为 \(\frac{x_0 x}{a^2} + \frac{y_0 y}{b^2} = 1\)。令其过 \(Q(x_1, y_1)\) 得 \(\frac{x_0 x_1}{a^2} + \frac{y_0 y_1}{b^2} = 1\)，联立 \(\frac{x_0^2}{a^2} + \frac{y_0^2}{b^2} = 1\) 解出 \((x_0, y_0)\) 即可。\(\blacksquare\)

\hrulefill

\section{第57章：双曲线}\label{ux7b2c57ux7ae0ux53ccux66f2ux7ebf}

双曲线是圆锥曲线中最''戏剧性''的一种------它由两支不连通的曲线组成，延伸到无穷远处，且具有渐近线。双曲线的焦点定义（到两焦点距离之差的绝对值为常数 \(2a\)）与椭圆形成鲜明对比：椭圆是''和''为常数，双曲线是''差''为常数。双曲线的标准方程 \(\frac{x^2}{a^2} - \frac{y^2}{b^2} = 1\) 中，减号是关键------它使得方程可以分解为 \((\frac{x}{a} - \frac{y}{b})(\frac{x}{a} + \frac{y}{b}) = 1\)，从而导出两条渐近线 \(y = \pm\frac{b}{a}x\)。与椭圆不同，双曲线的参数 \(a, b, c\) 满足 \(c^2 = a^2 + b^2\)（而非 \(c^2 = a^2 - b^2\)），离心率 \(e = \frac{c}{a} > 1\)。双曲线的渐近线是它最独特的特征------当曲线延伸到无穷远时，它无限逼近这两条直线但永不触及，这一性质在反比例函数 \(y = \frac{1}{x}\)（等轴双曲线）中已经出现。双曲线的光学性质也很特殊：从一支焦点发出的光线经双曲线反射后，反射光线的反向延长线经过另一支的焦点------这一性质在卡塞格林望远镜等光学仪器中有实际应用。

\subsection{26.3.1 双曲线的定义}\label{ux53ccux66f2ux7ebfux7684ux5b9aux4e49}

\textbf{定义 26.4（双曲线）.} 平面内到两个定点 \(F_1, F_2\)（焦点）的距离之\textbf{差}的绝对值等于常数 \(2a\)（\(0 < 2a < |F_1 F_2|\)）的点的轨迹称为\textbf{双曲线}。即

\[H = \{P \mid ||PF_1| - |PF_2|| = 2a, \; 0 < 2a < |F_1 F_2|\}\]

双曲线由\textbf{两支}组成：右支满足 \(|PF_1| - |PF_2| = 2a\)，左支满足 \(|PF_2| - |PF_1| = 2a\)。

\subsection{26.3.2 双曲线的标准方程}\label{ux53ccux66f2ux7ebfux7684ux6807ux51c6ux65b9ux7a0b}

\textbf{定理 26.6（双曲线的标准方程）.} 设双曲线的焦点为 \(F_1(-c, 0)\)，\(F_2(c, 0)\)（\(c > 0\)），\(0 < a < c\)。令 \(b^2 = c^2 - a^2\)（\(b > 0\)），则双曲线的标准方程为

\[\frac{x^2}{a^2} - \frac{y^2}{b^2} = 1 \quad (a > 0, b > 0)\]

\textbf{证明}： 考虑右支 \(|PF_1| - |PF_2| = 2a\)（对于左支，推导完全对称）。

\[\sqrt{(x + c)^2 + y^2} - \sqrt{(x - c)^2 + y^2} = 2a\]

移项：

\[\sqrt{(x + c)^2 + y^2} = 2a + \sqrt{(x - c)^2 + y^2}\]

两边平方：

\[(x + c)^2 + y^2 = 4a^2 + 4a\sqrt{(x - c)^2 + y^2} + (x - c)^2 + y^2\]

展开并化简：

\[x^2 + 2cx + c^2 + y^2 = 4a^2 + 4a\sqrt{(x - c)^2 + y^2} + x^2 - 2cx + c^2 + y^2\]

\[4cx - 4a^2 = 4a\sqrt{(x - c)^2 + y^2}\]

\[cx - a^2 = a\sqrt{(x - c)^2 + y^2}\]

两边再平方（注意：对右支有 \(x \geq a > 0\) 且 \(c > a\)，故 \(cx - a^2 \geq ca - a^2 > 0\)，等式两端非负）：

\[(cx - a^2)^2 = a^2[(x - c)^2 + y^2]\]

\[c^2 x^2 - 2a^2 cx + a^4 = a^2 x^2 - 2a^2 cx + a^2 c^2 + a^2 y^2\]

\[c^2 x^2 + a^4 = a^2 x^2 + a^2 c^2 + a^2 y^2\]

\[(c^2 - a^2)x^2 - a^2 y^2 = a^2(c^2 - a^2)\]

令 \(b^2 = c^2 - a^2 > 0\)（由 \(c > a\)）：

\[b^2 x^2 - a^2 y^2 = a^2 b^2\]

\[\frac{x^2}{a^2} - \frac{y^2}{b^2} = 1\]

验证附加条件：\(cx - a^2 \geq 0\)，即 \(x \geq \frac{a^2}{c}\)。由 \(\frac{x^2}{a^2} \geq 1\) 得 \(|x| \geq a\)。对于右支 \(x \geq a > \frac{a^2}{c}\)（因 \(c > a\)），条件满足。\(\blacksquare\)

\textbf{反向验证}：反之，设点 \(P(x, y)\) 满足 \(\frac{x^2}{a^2} - \frac{y^2}{b^2} = 1\)。由 \(y^2 = b^2\left(\frac{x^2}{a^2} - 1\right)\)，计算 \(|PF_1|\) 和 \(|PF_2|\)，可得 \(||PF_1| - |PF_2|| = 2a\)（详细推导与正向过程类似，此处从略）。

\subsection{26.3.3 双曲线的基本性质}\label{ux53ccux66f2ux7ebfux7684ux57faux672cux6027ux8d28}

\textbf{定理 26.7.} 双曲线 \(\frac{x^2}{a^2} - \frac{y^2}{b^2} = 1\)（\(a, b > 0\)）具有以下性质：

\begin{enumerate}
\def\labelenumi{\arabic{enumi}.}
\tightlist
\item
  \textbf{范围}：\(|x| \geq a\)（即 \(x \geq a\) 或 \(x \leq -a\)），\(y \in \mathbb{R}\)。
\item
  \textbf{对称性}：关于 \(x\) 轴、\(y\) 轴和原点对称。
\item
  \textbf{顶点}：\(A_1(-a, 0)\)，\(A_2(a, 0)\)。
\item
  \textbf{焦点}：\(F_1(-c, 0)\)，\(F_2(c, 0)\)，其中 \(c = \sqrt{a^2 + b^2}\)。
\item
  \textbf{渐近线}：\(y = \pm \frac{b}{a} x\)。
\item
  \textbf{离心率}：\(e = \frac{c}{a}\)（\(e > 1\)）。
\end{enumerate}

\textbf{证明}：

\textbf{（1）范围：} 由 \(\frac{x^2}{a^2} = 1 + \frac{y^2}{b^2} \geq 1\) 得 \(|x| \geq a\)。\(y\) 可取任意实数。

\textbf{（2）对称性：} 与椭圆的证明相同，替换 \((x,y)\) 为 \((x,-y)\)、\((-x,y)\)、\((-x,-y)\) 均不改变方程。

\textbf{（3）} 令 \(y = 0\) 得 \(x = \pm a\)。

\textbf{（4）} \(b^2 = c^2 - a^2\)，故 \(c = \sqrt{a^2 + b^2}\)。

\textbf{（5）渐近线：} 当 \((x, y)\) 沿双曲线趋于无穷远时，\(\frac{x^2}{a^2} - \frac{y^2}{b^2} = 1\) 两边除以 \(\frac{x^2}{a^2}\) 得

\[1 - \frac{a^2 y^2}{b^2 x^2} = \frac{a^2}{x^2} \to 0\]

故 \(\frac{y^2}{x^2} \to \frac{b^2}{a^2}\)，即 \(\frac{y}{x} \to \pm \frac{b}{a}\)。渐近线方程也可由 \(\frac{x^2}{a^2} - \frac{y^2}{b^2} = 0\) 得到，即 \(y = \pm\frac{b}{a}x\)。

\textbf{（6）} \(e = \frac{c}{a} = \frac{\sqrt{a^2 + b^2}}{a} > 1\)。\(\blacksquare\)

\subsection{26.3.4 双曲线的焦点距离公式}\label{ux53ccux66f2ux7ebfux7684ux7126ux70b9ux8dddux79bbux516cux5f0f}

\textbf{定理 26.8（焦点距离公式）.} 双曲线 \(\frac{x^2}{a^2} - \frac{y^2}{b^2} = 1\) 上一点 \(P(x_0, y_0)\) 到两个焦点的距离分别为

\[|PF_1| = |a + ex_0|, \quad |PF_2| = |a - ex_0|\]

其中 \(e = \frac{c}{a}\)。

\textbf{证明}： 与椭圆的推导类似。由 \(y_0^2 = b^2\left(\frac{x_0^2}{a^2} - 1\right) = \frac{b^2(x_0^2 - a^2)}{a^2}\)。

\[|PF_1|^2 = (x_0 + c)^2 + y_0^2 = (x_0 + c)^2 + \frac{b^2(x_0^2 - a^2)}{a^2}\]

\[= x_0^2 + 2cx_0 + c^2 + \frac{(c^2 - a^2)(x_0^2 - a^2)}{a^2}\]

\[= x_0^2 + 2cx_0 + c^2 + \frac{c^2 x_0^2 - c^2 a^2 - a^2 x_0^2 + a^4}{a^2}\]

\[= x_0^2 + 2cx_0 + c^2 + \frac{c^2 x_0^2}{a^2} - c^2 - x_0^2 + a^2\]

\[= a^2 + 2cx_0 + \frac{c^2 x_0^2}{a^2} = \left(a + \frac{cx_0}{a}\right)^2 = (a + ex_0)^2\]

故 \(|PF_1| = |a + ex_0|\)（取绝对值是因为 \(x_0\) 可能为负）。类似地，\(|PF_2| = |a - ex_0|\)。\(\blacksquare\)

\textbf{推论 26.5.} 对于右支（\(x_0 \geq a\)），\(|PF_1| = a + ex_0\)，\(|PF_2| = ex_0 - a\)（因 \(ex_0 > ea > a\)）。对于左支（\(x_0 \leq -a\)），\(|PF_1| = -(a + ex_0)\)，\(|PF_2| = a - ex_0\)。

\textbf{证明}：对右支 \(x_0 \geq a\)，\(ex_0 \geq ea > a\)，故 \(|a + ex_0| = a + ex_0\)，\(|a - ex_0| = ex_0 - a\)。左支类似。 \(\blacksquare\)

\subsection{26.3.5 双曲线的准线}\label{ux53ccux66f2ux7ebfux7684ux51c6ux7ebf}

\textbf{定理 26.10} 双曲线 \(\frac{x^2}{a^2} - \frac{y^2}{b^2} = 1\)（\(a > 0, b > 0\)）对应于焦点 \(F_1(-c, 0)\) 和 \(F_2(c, 0)\) 的两条\textbf{准线}方程分别为： \[l_1: x = -\frac{a^2}{c} = -\frac{a}{e}, \quad l_2: x = \frac{a^2}{c} = \frac{a}{e}\]

其中 \(e = \frac{c}{a} > 1\) 为离心率。

\textbf{证明}：以右焦点 \(F_2(c, 0)\) 和右准线 \(l_2: x = \frac{a^2}{c}\) 为例。对双曲线上任意点 \(P(x, y)\)，\(P\) 到 \(F_2\) 的距离为 \(|PF_2| = |ex - a|\)（由焦点距离公式），\(P\) 到 \(l_2\) 的距离为 \(|x - \frac{a^2}{c}| = \frac{1}{c}|cx - a^2| = \frac{a}{c}|ex - a|\)。故 \(\frac{|PF_2|}{d(P, l_2)} = \frac{|ex - a|}{\frac{a}{c}|ex - a|} = \frac{c}{a} = e\)。\(\blacksquare\)''

\textbf{推论 26.4} 双曲线的统一焦点-准线定义：双曲线是到焦点距离与到对应准线距离之比等于离心率 \(e > 1\) 的点的轨迹。双曲线有两条准线，分别对应两个焦点。准线不与双曲线相交（因为 \(\frac{a}{e} < a\)，准线在顶点内侧）。

\textbf{证明}：在定理26.10中已验证对双曲线上任意点 \(P\)，\(|PF_2|/d(P, l_2) = e > 1\)。此即统一定义中 \(e > 1\) 的情形。准线在顶点内侧因 \(\frac{a}{e} = \frac{a^2}{c} < a\)。\(\blacksquare\)

\subsection{26.3.6 等轴双曲线}\label{ux7b49ux8f74ux53ccux66f2ux7ebf}

\textbf{定义 26.5.} 当 \(a = b\) 时，双曲线 \(\frac{x^2}{a^2} - \frac{y^2}{a^2} = 1\) 即 \(x^2 - y^2 = a^2\) 称为\textbf{等轴双曲线}（或\textbf{等边双曲线}）。此时渐近线为 \(y = \pm x\)（互相垂直），离心率 \(e = \sqrt{2}\)。

\subsection{26.3.7 双曲线的切线方程}\label{ux53ccux66f2ux7ebfux7684ux5207ux7ebfux65b9ux7a0b}

\textbf{定理 26.9（双曲线切线方程）.} 双曲线 \(\frac{x^2}{a^2} - \frac{y^2}{b^2} = 1\) 在点 \(P_0(x_0, y_0)\)（在双曲线上）处的切线方程为

\[\frac{x_0 x}{a^2} - \frac{y_0 y}{b^2} = 1\]

\textbf{证明}： 对 \(\frac{x^2}{a^2} - \frac{y^2}{b^2} = 1\) 隐函数求导，\(\frac{2x}{a^2} - \frac{2yy'}{b^2} = 0\)，得 \(y' = \frac{b^2 x}{a^2 y}\)。在 \(P_0(x_0, y_0)\) 处斜率 \(k = \frac{b^2 x_0}{a^2 y_0}\)。

切线方程 \(y - y_0 = \frac{b^2 x_0}{a^2 y_0}(x - x_0)\)，整理得 \(a^2 y_0(y - y_0) = b^2 x_0(x - x_0)\)，即 \(b^2 x_0 x - a^2 y_0 y = b^2 x_0^2 - a^2 y_0^2\)。

利用 \(P_0\) 在双曲线上，\(\frac{x_0^2}{a^2} - \frac{y_0^2}{b^2} = 1\)，即 \(b^2 x_0^2 - a^2 y_0^2 = a^2 b^2\)，代入得

\[b^2 x_0 x - a^2 y_0 y = a^2 b^2\]

两边除以 \(a^2 b^2\)，得 \(\frac{x_0 x}{a^2} - \frac{y_0 y}{b^2} = 1\)。\(\blacksquare\)

\textbf{特殊情况：} 当 \(y_0 = 0\) 时，\(x_0 = \pm a\)，切线为 \(x = \pm a\)（双曲线顶点处的切线垂直于 \(x\) 轴），公式 \(\frac{x_0 x}{a^2} - \frac{y_0 y}{b^2} = 1\) 仍然成立。

\hrulefill

\section{第58章：抛物线}\label{ux7b2c58ux7ae0ux629bux7269ux7ebf}

抛物线是圆锥曲线中离心率恰好等于 \(1\) 的临界情形------它恰好是椭圆（\(e < 1\)）和双曲线（\(e > 1\)）的分界线。抛物线的焦点定义（到焦点和到准线的距离相等）在所有圆锥曲线中最为简洁，这使得抛物线在物理中有着独特的地位：抛射体的运动轨迹（忽略空气阻力）是抛物线；平行于对称轴的光线经抛物线反射后汇聚于焦点，这一性质被广泛应用于抛物面天线和太阳能聚光器。抛物线的标准方程 \(y^2 = 2px\) 是所有圆锥曲线方程中形式最简单的，其焦点为 \((\frac{p}{2}, 0)\)，准线为 \(x = -\frac{p}{2}\)。参数 \(p\) 是焦点到准线距离的两倍，它决定了抛物线的''开口宽度''------\(p\) 越大，抛物线开口越宽。抛物线是唯一没有中心的圆锥曲线（椭圆和双曲线都有对称中心），它只有一个顶点、一个焦点和一条对称轴。这一''非中心''的性质使得抛物线在射影几何中具有特殊地位------所有抛物线在射影变换下都是等价的。

\subsection{26.4.1 抛物线的定义}\label{ux629bux7269ux7ebfux7684ux5b9aux4e49}

\textbf{定义 26.6（抛物线）.} 平面内到一个定点 \(F\)（焦点）和一条定直线 \(l\)（准线，\(F \notin l\)）距离相等的点的轨迹称为\textbf{抛物线}（parabola）。即

\[\mathcal{P} = \{P \mid |PF| = d(P, l)\}\]

这对应于统一定义中 \(e = 1\) 的情形。

\subsection{26.4.2 抛物线的标准方程}\label{ux629bux7269ux7ebfux7684ux6807ux51c6ux65b9ux7a0b}

\textbf{定理 26.11（抛物线的标准方程）.} 设焦点为 \(F\left(\frac{p}{2}, 0\right)\)（\(p > 0\)），准线为 \(x = -\frac{p}{2}\)，则抛物线的标准方程为

\[y^2 = 2px \quad (p > 0)\]

\textbf{证明}： 设 \(P(x, y)\) 为抛物线上任意一点。

\(P\) 到焦点 \(F\left(\frac{p}{2}, 0\right)\) 的距离：

\[|PF| = \sqrt{\left(x - \frac{p}{2}\right)^2 + y^2}\]

\(P\) 到准线 \(x = -\frac{p}{2}\) 的距离：

\[d(P, l) = \left|x + \frac{p}{2}\right|\]

由定义 \(|PF| = d(P, l)\)：

\[\sqrt{\left(x - \frac{p}{2}\right)^2 + y^2} = \left|x + \frac{p}{2}\right|\]

两边平方：

\[\left(x - \frac{p}{2}\right)^2 + y^2 = \left(x + \frac{p}{2}\right)^2\]

\[x^2 - px + \frac{p^2}{4} + y^2 = x^2 + px + \frac{p^2}{4}\]

\[y^2 = 2px\]

验证附加条件：\(x + \frac{p}{2} \geq 0\)，即 \(x \geq -\frac{p}{2}\)。由 \(y^2 = 2px\) 和 \(p > 0\)，当 \(y^2 \geq 0\) 时 \(x \geq 0 > -\frac{p}{2}\)，条件自动满足。\(\blacksquare\)

\textbf{命题 26.4（抛物线的四种标准方程）.} 根据焦点和准线位置的不同，抛物线有四种标准形式：

\begin{longtable}[]{@{}
  >{\raggedright\arraybackslash}p{(\linewidth - 6\tabcolsep) * \real{0.2143}}
  >{\raggedright\arraybackslash}p{(\linewidth - 6\tabcolsep) * \real{0.2143}}
  >{\raggedright\arraybackslash}p{(\linewidth - 6\tabcolsep) * \real{0.2143}}
  >{\raggedright\arraybackslash}p{(\linewidth - 6\tabcolsep) * \real{0.3571}}@{}}
\toprule
\begin{minipage}[b]{\linewidth}\raggedright
方程
\end{minipage} & \begin{minipage}[b]{\linewidth}\raggedright
焦点
\end{minipage} & \begin{minipage}[b]{\linewidth}\raggedright
准线
\end{minipage} & \begin{minipage}[b]{\linewidth}\raggedright
开口方向
\end{minipage} \\
\midrule
\endhead
\bottomrule
\endlastfoot
\(y^2 = 2px\)（\(p > 0\)） & \(\left(\frac{p}{2}, 0\right)\) & \(x = -\frac{p}{2}\) & 右 \\
\(y^2 = -2px\)（\(p > 0\)） & \(\left(-\frac{p}{2}, 0\right)\) & \(x = \frac{p}{2}\) & 左 \\
\(x^2 = 2py\)（\(p > 0\)） & \(\left(0, \frac{p}{2}\right)\) & \(y = -\frac{p}{2}\) & 上 \\
\(x^2 = -2py\)（\(p > 0\)） & \(\left(0, -\frac{p}{2}\right)\) & \(y = \frac{p}{2}\) & 下 \\
\end{longtable}

\textbf{证明}： 以 \(y^2 = -2px\) 为例。设焦点 \(F\left(-\frac{p}{2}, 0\right)\)，准线 \(x = \frac{p}{2}\)。由 \(|PF| = d(P, l)\)：

\[\sqrt{\left(x + \frac{p}{2}\right)^2 + y^2} = \left|x - \frac{p}{2}\right|\]

平方并化简得 \(y^2 = -2px\)。其余两种类似（交换 \(x, y\) 即可）。\(\blacksquare\)

\textbf{定理 26.18a} 抛物线 \(y^2 = 2px\)（\(p > 0\)）在点 \(P_0(x_0, y_0)\)（\(y_0 \neq 0\)）处的切线方程为： \[y_0 y = p(x + x_0)\]

当 \(y_0 = 0\) 时，\(P_0\) 为顶点，切线为 \(y\) 轴（\(x = 0\)）。

\textbf{证明}：对 \(y^2 = 2px\) 两边关于 \(x\) 求导（隐函数求导）：\(2y \cdot y' = 2p\)，故 \(y' = \frac{p}{y}\)。

在 \(P_0\) 处，斜率 \(k = \frac{p}{y_0}\)。切线方程：\(y - y_0 = \frac{p}{y_0}(x - x_0)\)。

整理：\(y_0(y - y_0) = p(x - x_0)\)，即 \(y_0 y - y_0^2 = px - px_0\)。

由 \(P_0\) 在抛物线上：\(y_0^2 = 2px_0\)，代入得 \(y_0 y - 2px_0 = px - px_0\)，即 \(y_0 y = px + px_0 = p(x + x_0)\)。\(\blacksquare\)

\subsection{26.4.3 抛物线的基本性质}\label{ux629bux7269ux7ebfux7684ux57faux672cux6027ux8d28-1}

\textbf{定理 26.12.} 抛物线 \(y^2 = 2px\)（\(p > 0\)）具有以下性质：

\begin{enumerate}
\def\labelenumi{\arabic{enumi}.}
\tightlist
\item
  \textbf{范围}：\(x \geq 0\)，\(y \in \mathbb{R}\)。
\item
  \textbf{对称性}：关于 \(x\) 轴对称（但不关于 \(y\) 轴或原点对称）。
\item
  \textbf{顶点}：原点 \((0, 0)\)。
\item
  \textbf{焦点}：\(\left(\frac{p}{2}, 0\right)\)。
\item
  \textbf{准线}：\(x = -\frac{p}{2}\)。
\item
  \textbf{离心率}：\(e = 1\)。
\end{enumerate}

\textbf{证明}：

\textbf{（1）} 由 \(y^2 = 2px\) 和 \(p > 0\)，\(y^2 \geq 0\) 推出 \(x \geq 0\)。

\textbf{（2）} 将 \((x, y)\) 替换为 \((x, -y)\)：\((-y)^2 = y^2 = 2px\)，方程不变，故关于 \(x\) 轴对称。

\textbf{（3）-（5）} 由方程直接可得。

\textbf{（6）} 抛物线是统一定义中 \(e = 1\) 的情形。\(\blacksquare\)

\subsection{26.4.4 抛物线的焦半径}\label{ux629bux7269ux7ebfux7684ux7126ux534aux5f84}

\textbf{定理 26.13（焦半径公式）.} 抛物线 \(y^2 = 2px\) 上一点 \(P(x_0, y_0)\) 到焦点的距离为

\[|PF| = x_0 + \frac{p}{2}\]

\textbf{证明}： \(P\) 到焦点 \(F\left(\frac{p}{2}, 0\right)\) 的距离：

\[|PF| = \sqrt{\left(x_0 - \frac{p}{2}\right)^2 + y_0^2}\]

由 \(y_0^2 = 2px_0\)：

\[|PF|^2 = x_0^2 - px_0 + \frac{p^2}{4} + 2px_0 = x_0^2 + px_0 + \frac{p^2}{4} = \left(x_0 + \frac{p}{2}\right)^2\]

由于 \(x_0 \geq 0\) 且 \(p > 0\)，\(x_0 + \frac{p}{2} > 0\)，故 \(|PF| = x_0 + \frac{p}{2}\)。\(\blacksquare\)

\textbf{推论 26.6.} 焦半径也可表示为 \(|PF| = d(P, l) = x_0 + \frac{p}{2}\)，即抛物线上一点到焦点的距离等于到准线的距离------这正是抛物线的定义。

\textbf{证明}：由抛物线定义（定义 26.6），\(|PF| = d(P,l)\)。\(P\) 到准线 \(x = -\frac{p}{2}\) 的距离为 \(\left|x_0 + \frac{p}{2}\right| = x_0 + \frac{p}{2}\)（因 \(x_0 \geq 0\)）。\(\blacksquare\)

\subsection{26.4.5 抛物线的焦点弦}\label{ux629bux7269ux7ebfux7684ux7126ux70b9ux5f26}

\textbf{定理 26.14（焦点弦性质）.} 过抛物线 \(y^2 = 2px\) 焦点的弦 \(AB\)（称为\textbf{焦点弦}），设 \(A(x_1, y_1)\)，\(B(x_2, y_2)\)，则：

\begin{enumerate}
\def\labelenumi{\arabic{enumi}.}
\tightlist
\item
  \(x_1 x_2 = \frac{p^2}{4}\)
\item
  \(y_1 y_2 = -p^2\)
\item
  焦点弦长 \(|AB| = x_1 + x_2 + p\)
\end{enumerate}

\textbf{证明}： 设过焦点 \(F\left(\frac{p}{2}, 0\right)\) 的直线方程为 \(x = my + \frac{p}{2}\)（\(m\) 为参数，此参数化可以涵盖所有过焦点的非水平直线）。

\textbf{特殊情况}：参数化 \(x = my + \frac{p}{2}\) 不能覆盖的是过焦点的水平线 \(y = 0\)（即对称轴本身）。该直线与抛物线 \(y^2 = 2px\) 仅交于顶点 \((0, 0)\)，不构成弦。因此 \(x = my + \frac{p}{2}\) 的参数化已涵盖所有焦点弦。

代入 \(y^2 = 2px\)：

\[y^2 = 2p\left(my + \frac{p}{2}\right) = 2pmy + p^2\]

\[y^2 - 2pmy - p^2 = 0\]

由韦达定理（Ch 14）：

\[y_1 + y_2 = 2pm, \quad y_1 y_2 = -p^2\]

由此：

\[x_1 x_2 = \frac{y_1^2}{2p} \cdot \frac{y_2^2}{2p} = \frac{(y_1 y_2)^2}{4p^2} = \frac{p^4}{4p^2} = \frac{p^2}{4}\]

焦点弦长：

\[|AB| = |AF| + |BF| = \left(x_1 + \frac{p}{2}\right) + \left(x_2 + \frac{p}{2}\right) = x_1 + x_2 + p\]

\(\blacksquare\)

\textbf{定理 26.15（焦点弦的最短弦长）.} 过抛物线焦点的弦中，垂直于对称轴的弦最短，长度为 \(2p\)（通径）。

\textbf{证明}： 由定理 26.14，\(|AB| = x_1 + x_2 + p\)。由韦达定理 \(y_1 + y_2 = 2pm\)，\(y_1 y_2 = -p^2\)，故

\[x_1 + x_2 = \frac{y_1^2 + y_2^2}{2p} = \frac{(y_1 + y_2)^2 - 2y_1 y_2}{2p} = \frac{4p^2 m^2 + 2p^2}{2p} = 2pm^2 + p\]

\[|AB| = 2pm^2 + p + p = 2p(m^2 + 1) \geq 2p\]

等号当 \(m = 0\) 时成立，此时直线 \(x = \frac{p}{2}\) 垂直于对称轴，弦长为 \(2p\)（通径）。\(\blacksquare\)

\hrulefill

\section{第59章：圆锥曲线的统一性质}\label{ux7b2c59ux7ae0ux5706ux9525ux66f2ux7ebfux7684ux7edfux4e00ux6027ux8d28}

尽管椭圆、双曲线和抛物线在外观上差异巨大------椭圆封闭、双曲线有两支、抛物线向一侧无限延伸------但它们共享着深刻的统一性质。离心率 \(e\) 是贯穿三种曲线的核心参数：\(e\) 决定了曲线的形状类别（\(e < 1\) 椭圆，\(e = 1\) 抛物线，\(e > 1\) 双曲线），也与曲线的''扁平''或''张开''程度直接相关。焦点-准线关系 \(|PF| = e \cdot d(P, l)\) 是三种曲线的统一定义，它揭示了圆锥曲线本质上是''到定点与定直线距离之比为常数''的点的轨迹。焦点弦（过焦点且与曲线相交的弦）的长度公式也呈现出统一的形式------例如，椭圆和双曲线的焦点弦长都可以用 \(a, b\) 和倾斜角 \(\alpha\) 表达。更为深刻的是，三种圆锥曲线在极坐标下可以统一表示为 \(r = \frac{ep}{1 - e\cos\theta}\)，这一公式将离心率 \(e\) 和焦准距 \(p\) 作为核心参数，以极点位于焦点处，极轴沿焦点-准线方向。这一统一极坐标方程不仅适用于三种圆锥曲线，还自然地包含了圆（\(e = 0\)）这一退化情形，体现了数学中''统一性''的美学价值。

\subsection{26.5.1 离心率的统一角色}\label{ux79bbux5fc3ux7387ux7684ux7edfux4e00ux89d2ux8272}

\textbf{定理 26.16.} 设圆锥曲线的焦点为 \(F\)，准线为 \(l\)，离心率为 \(e\)。圆锥曲线上一点 \(P\) 到焦点 \(F\) 的距离 \(|PF|\) 和到准线的距离 \(d(P, l)\) 之间的关系为

\[|PF| = e \cdot d(P, l)\]

且

\begin{itemize}
\tightlist
\item
  椭圆（\(0 < e < 1\)）：曲线封闭，焦点在曲线内部。
\item
  抛物线（\(e = 1\)）：曲线开口无限延伸，但只有一个焦点。
\item
  双曲线（\(e > 1\)）：曲线由两支组成，两支各有一个焦点。
\end{itemize}

\textbf{证明}：对椭圆（\(e < 1\)）、抛物线（\(e = 1\)）、双曲线（\(e > 1\)），分别在定理26.7、26.10、26.15中已验证 \(|PF|/d(P,l) = e\)。此关系将三种曲线统一为焦点-准线定义。\(\blacksquare\)

\subsection{26.5.2 焦点弦长的统一公式}\label{ux7126ux70b9ux5f26ux957fux7684ux7edfux4e00ux516cux5f0f}

\textbf{定理 26.17.} 对于标准位置的圆锥曲线，过焦点且倾斜角为 \(\alpha\) 的弦长为：

\begin{itemize}
\tightlist
\item
  椭圆 \(\frac{x^2}{a^2} + \frac{y^2}{b^2} = 1\)：\(|AB| = \frac{2b^2/a}{1 - e^2\cos^2\alpha}\)
\item
  双曲线 \(\frac{x^2}{a^2} - \frac{y^2}{b^2} = 1\)：\(|AB| = \frac{2b^2/a}{|1 - e^2\cos^2\alpha|}\)
\item
  抛物线 \(y^2 = 2px\)：\(|AB| = \frac{2p}{\sin^2\alpha}\)
\end{itemize}

\textbf{证明}： 以椭圆为例。过右焦点 \(F_2(c, 0)\) 的直线参数方程为

\[\begin{cases} x = c + t\cos\alpha \\ y = t\sin\alpha \end{cases}\]

代入椭圆方程：

\[\frac{(c + t\cos\alpha)^2}{a^2} + \frac{t^2\sin^2\alpha}{b^2} = 1\]

两边乘以 \(a^2 b^2\) 并展开：

\[b^2(c^2 + 2ct\cos\alpha + t^2\cos^2\alpha) + a^2 t^2\sin^2\alpha = a^2 b^2\]

注意 \(b^2 c^2 = b^2(a^2 - b^2)\)（因 \(c^2 = a^2 - b^2\)），整理关于 \(t\) 的二次方程：

\[(b^2\cos^2\alpha + a^2\sin^2\alpha)t^2 + 2b^2 c\cos\alpha \cdot t - b^4 = 0\]

由韦达定理：

\[t_1 + t_2 = \frac{-2b^2 c\cos\alpha}{b^2\cos^2\alpha + a^2\sin^2\alpha}, \quad t_1 t_2 = \frac{-b^4}{b^2\cos^2\alpha + a^2\sin^2\alpha}\]

因此：

\[|AB|^2 = (t_1 - t_2)^2 = (t_1 + t_2)^2 - 4t_1 t_2 = \frac{4b^4 c^2\cos^2\alpha + 4b^4(b^2\cos^2\alpha + a^2\sin^2\alpha)}{(b^2\cos^2\alpha + a^2\sin^2\alpha)^2}\]

化简分子：\(c^2\cos^2\alpha + b^2\cos^2\alpha + a^2\sin^2\alpha = (c^2 + b^2)\cos^2\alpha + a^2\sin^2\alpha = a^2\cos^2\alpha + a^2\sin^2\alpha = a^2\)（因 \(c^2 + b^2 = a^2\)）。故

\[|AB| = \frac{2ab^2}{b^2\cos^2\alpha + a^2\sin^2\alpha}\]

再将分母表示为：\(b^2\cos^2\alpha + a^2\sin^2\alpha = a^2 - c^2\cos^2\alpha = a^2(1 - e^2\cos^2\alpha)\)，代入得

\[|AB| = \frac{2ab^2}{a^2(1 - e^2\cos^2\alpha)} = \frac{2b^2/a}{1 - e^2\cos^2\alpha}\]

\(\blacksquare\)''

\subsection{26.5.3 三种圆锥曲线的对比}\label{ux4e09ux79cdux5706ux9525ux66f2ux7ebfux7684ux5bf9ux6bd4}

\begin{longtable}[]{@{}llll@{}}
\toprule
性质 & 椭圆 & 双曲线 & 抛物线 \\
\midrule
\endhead
\bottomrule
\endlastfoot
统一定义 & \(e < 1\) & \(e > 1\) & \(e = 1\) \\
焦点数 & 2 & 2 & 1 \\
准线数 & 2 & 2 & 1 \\
封闭性 & 封闭 & 不封闭（两支） & 不封闭 \\
对称轴 & 2 条 & 2 条 & 1 条 \\
渐近线 & 无 & 2 条 & 无 \\
二次方程 & \(B^2 - 4AC < 0\) & \(B^2 - 4AC > 0\) & \(B^2 - 4AC = 0\) \\
\end{longtable}

\textbf{说明：} 最后一行中，一般二次方程 \(Ax^2 + Bxy + Cy^2 + Dx + Ey + F = 0\) 的判别式 \(\Delta = B^2 - 4AC\) 决定了圆锥曲线的类型（在无退化的情况下）。

\subsection{26.5.4 椭圆的光学性质}\label{ux692dux5706ux7684ux5149ux5b66ux6027ux8d28}

\textbf{定理 26.18（椭圆的反射性质）.} 设 \(P\) 为椭圆上一点，\(F_1, F_2\) 为两个焦点。则 \(PF_1\) 与 \(PF_2\) 和椭圆在 \(P\) 处的切线所成的角相等。即椭圆在 \(P\) 处的切线是 \(\angle F_1PF_2\) 的外角平分线。

\textbf{证明}： 设椭圆 \(\frac{x^2}{a^2} + \frac{y^2}{b^2} = 1\)，\(P(x_0, y_0)\) 在椭圆上。切线方程为 \(\frac{x_0 x}{a^2} + \frac{y_0 y}{b^2} = 1\)，法向量为 \(\vec{n} = \left(\frac{x_0}{a^2}, \frac{y_0}{b^2}\right)\)。

\(\overrightarrow{PF_1} = (-c - x_0, -y_0)\)，\(\overrightarrow{PF_2} = (c - x_0, -y_0)\)。由定理 26.3，\(|\overrightarrow{PF_1}| = a + ex_0\)，\(|\overrightarrow{PF_2}| = a - ex_0\)。

法向量 \(\vec{n}\) 平分外角等价于 \(\frac{\overrightarrow{PF_1}}{|\overrightarrow{PF_1}|} + \frac{\overrightarrow{PF_2}}{|\overrightarrow{PF_2}|}\) 与 \(\vec{n}\) 平行。计算：

\[\frac{\overrightarrow{PF_1}}{|\overrightarrow{PF_1}|} + \frac{\overrightarrow{PF_2}}{|\overrightarrow{PF_2}|} = \frac{(-ae - x_0, -y_0)}{a + ex_0} + \frac{(ae - x_0, -y_0)}{a - ex_0}\]

\(x\) 分量：\(\frac{-(ae + x_0)}{a + ex_0} + \frac{ae - x_0}{a - ex_0} = \frac{-(ae + x_0)(a - ex_0) + (ae - x_0)(a + ex_0)}{(a + ex_0)(a - ex_0)}\)

分子 \(= -(a^2 e - ae^2 x_0 + ax_0 - ex_0^2) + (a^2 e + ae^2 x_0 - ax_0 - ex_0^2) = 2ax_0(e^2 - 1)\)

\(y\) 分量：\(\frac{-y_0}{a + ex_0} + \frac{-y_0}{a - ex_0} = \frac{-2ay_0}{a^2 - e^2 x_0^2}\)

因此 \(\frac{\overrightarrow{PF_1}}{|\overrightarrow{PF_1}|} + \frac{\overrightarrow{PF_2}}{|\overrightarrow{PF_2}|} = \left(\frac{2ax_0(e^2 - 1)}{a^2 - e^2 x_0^2},\, \frac{-2ay_0}{a^2 - e^2 x_0^2}\right)\)

验证其与 \(\vec{n}\) 的叉积为零：

\[\frac{2ax_0(e^2 - 1)}{a^2 - e^2 x_0^2} \cdot \frac{y_0}{b^2} - \frac{-2ay_0}{a^2 - e^2 x_0^2} \cdot \frac{x_0}{a^2} = \frac{2ax_0 y_0}{a^2 - e^2 x_0^2}\left(\frac{e^2 - 1}{b^2} + \frac{1}{a^2}\right)\]

利用 \(b^2 = a^2(1 - e^2)\)，得 \(\frac{e^2 - 1}{a^2(1 - e^2)} + \frac{1}{a^2} = \frac{-1}{a^2} + \frac{1}{a^2} = 0\)。

叉积为零，故法向量平分外角，即切线平分 \(\angle F_1PF_2\) 的外角。\(\blacksquare\)

\textbf{物理意义：} 椭圆的反射性质意味着从一个焦点发出的光线或声波，经椭圆面反射后会聚到另一个焦点。这一性质在医学（碎石机）、声学（'' whispering gallery''）和光学（椭圆反射镜）中有重要应用。

\subsection{26.5.5 抛物线的光学性质}\label{ux629bux7269ux7ebfux7684ux5149ux5b66ux6027ux8d28}

\textbf{定理 26.19（抛物线的反射性质）.} 设 \(P\) 为抛物线 \(y^2 = 2px\)（\(p > 0\)）上一点，\(F\) 为焦点。则抛物线在 \(P\) 处的切线平分 \(PF\) 与从 \(P\) 向准线所作垂线之间的夹角。等价地，入射角（\(PF\) 与切线的夹角）等于反射角（切线与水平方向的夹角）。

\textbf{证明}： 设 \(P(x_0, y_0)\) 在抛物线上（\(x_0 > 0\)，\(y_0 \neq 0\)）。焦点 \(F = (p/2, 0)\)，准线 \(x = -p/2\)。由焦点--准线性质，\(|PF| = x_0 + p/2\)。

\(\overrightarrow{PF} = (p/2 - x_0, -y_0)\)。从 \(P\) 向准线所作垂线的方向为 \((-1, 0)\)（水平向左）。

切线方向向量为 \((y_0, p)\)（由切线斜率 \(p/y_0\) 得到）。验证切线平分 \(\overrightarrow{PF}\) 与 \((-1, 0)\) 之间的夹角，即验证 \(\frac{\overrightarrow{PF}}{|\overrightarrow{PF}|} + (-1, 0)\) 与切线方向 \((y_0, p)\) 平行：

\[\frac{\overrightarrow{PF}}{|\overrightarrow{PF}|} + (-1, 0) = \left(\frac{p/2 - x_0}{x_0 + p/2} - 1,\, \frac{-y_0}{x_0 + p/2}\right) = \left(\frac{-2x_0}{x_0 + p/2},\, \frac{-y_0}{x_0 + p/2}\right)\]

该向量与 \((2x_0, y_0)\) 成比例。验证叉积：

\[2x_0 \cdot p - y_0 \cdot y_0 = 2px_0 - y_0^2 = 0\]

最后一步由 \(y_0^2 = 2px_0\) 保证。叉积为零，故切线平分夹角。

因此入射角（\(PF\) 与切线的夹角）等于反射角（切线与水平方向的夹角）。\(\blacksquare\)

\textbf{物理意义：} 平行于对称轴入射的光线，经抛物线面反射后会聚到焦点。这就是卫星天线、车灯反光镜和射电望远镜采用抛物面形状的原因。

\hrulefill

\begin{quote}
\textbf{Ch 325 本章展望}

本章从圆锥曲线的统一定义出发，系统推导了椭圆、双曲线和抛物线的标准方程，并建立了焦点距离公式、准线方程、焦点弦性质等核心理论。三种曲线在离心率 \(e\) 的参数下获得了完美的统一：\(e < 1\) 为椭圆，\(e = 1\) 为抛物线，\(e > 1\) 为双曲线。

然而，到目前为止，我们一直使用直角坐标系来描述曲线。直角坐标系虽然直观，但在处理某些问题时并不方便------例如，圆的方程 \((x-a)^2 + (y-b)^2 = r^2\) 在极坐标系中简化为 \(r = 2a\cos\theta\)；椭圆、双曲线和抛物线在以焦点为极点的极坐标系下可以统一为 \(r = \frac{ep}{1 - e\cos\theta}\)。在下一章中，我们将引入\textbf{参数方程}和\textbf{极坐标方程}两种新的曲线描述方式，它们不仅能简化计算，更能揭示圆锥曲线之间更深层的统一性。
\end{quote}

\hrulefill

\hrulefill

\hrulefill

\hrulefill

\hrulefill

\newpage

\chapter{03-函数与微积分/02-解析几何/03-参数与极坐标.md}\label{ux51fdux6570ux4e0eux5faeux79efux520602-ux89e3ux6790ux51e0ux4f5503-ux53c2ux6570ux4e0eux6781ux5750ux6807.md}

\section{第60章：参数方程与极坐标方程}\label{ux7b2c60ux7ae0ux53c2ux6570ux65b9ux7a0bux4e0eux6781ux5750ux6807ux65b9ux7a0b}

\begin{quote}
\textbf{前置知识}

本章建立在以下前序章节的基础之上：

\begin{itemize}
\tightlist
\item
  \textbf{Ch 325（圆锥曲线）}：椭圆、双曲线和抛物线的标准方程是参数方程和极坐标方程转换的基础。圆锥曲线的统一定义（焦点-准线）直接导出统一极坐标方程 \(r = \frac{ep}{1 - e\cos\theta}\)。
\item
  \textbf{Ch 27（三角恒等变换与反三角函数）}：参数方程大量使用三角函数------圆的参数方程使用 \(\cos\theta, \sin\theta\)，椭圆使用''离心角''，双曲线使用 \(\sec\theta, \tan\theta\)。极坐标与直角坐标的互化也依赖于 \(\cos\theta\) 和 \(\sin\theta\)。
\item
  \textbf{Ch 325（直线与圆）}：直线和圆的参数方程是理解参数方程概念的起点。
\item
  \textbf{Ch 31（积分初步）}：极坐标下的面积公式涉及定积分，为极坐标方程的应用提供了分析工具。
\end{itemize}
\end{quote}

\hrulefill

\subsection{27.1 参数方程}\label{ux53c2ux6570ux65b9ux7a0b}

参数方程是描述曲线的另一种重要方式------与普通方程 \(F(x, y) = 0\) 直接给出 \(x\) 和 \(y\) 之间的约束关系不同，参数方程通过引入一个中间变量 \(t\)（参数），将 \(x\) 和 \(y\) 分别表示为 \(t\) 的函数：\(x = f(t), y = g(t)\)。参数方程的优势在于它能够自然地描述那些用普通方程难以表达的曲线------特别是那些''非函数''的曲线（如圆 \(x^2 + y^2 = r^2\) 无法用 \(y = f(x)\) 的形式完全表示，但参数方程 \(x = r\cos t, y = r\sin t\) 可以）。参数 \(t\) 通常具有明确的几何或物理意义------在直线参数方程中，\(|t|\) 表示点到定点的距离；在圆的参数方程中，\(t\) 表示圆心角；在物理学中，\(t\) 通常就是时间，参数方程 \(x = x(t), y = y(t)\) 描述了质点的运动轨迹。参数方程的另一个重要应用是''消参''------通过消去参数 \(t\) 得到普通方程，从而在直角坐标系中分析曲线的性质。本节从参数方程的基本概念出发，系统介绍直线、圆和圆锥曲线的参数方程，以及消参法的常用技巧。

\subsection{27.1.1 参数方程的概念}\label{ux53c2ux6570ux65b9ux7a0bux7684ux6982ux5ff5}

\textbf{定义 27.1（参数方程）.} 在平面直角坐标系中，如果曲线 \(C\) 上的任意一点 \((x, y)\) 都可以表示为

\[\begin{cases} x = f(t) \\ y = g(t) \end{cases} \quad (t \in I)\]

其中 \(I\) 为参数 \(t\) 的取值范围，\(f\) 和 \(g\) 是关于 \(t\) 的函数，且对于 \(I\) 中的每一个 \(t\) 值，由此确定的点 \((f(t), g(t))\) 都在曲线 \(C\) 上，则称上述方程组为曲线 \(C\) 的\textbf{参数方程}，\(t\) 称为\textbf{参数}。

相应地，直接描述 \(x\) 和 \(y\) 之间关系的方程 \(F(x, y) = 0\) 称为曲线的\textbf{普通方程}（或直角坐标方程）。

\textbf{定理 27.1（消参的一般方法）.} 将参数方程化为普通方程的方法称为\textbf{消参法}，常用方法包括：

\begin{enumerate}
\def\labelenumi{\arabic{enumi}.}
\tightlist
\item
  \textbf{代入消参}：从一个方程解出参数，代入另一个方程。
\item
  \textbf{利用恒等式消参}：利用 \(\sin^2 t + \cos^2 t = 1\) 等三角恒等式。
\item
  \textbf{加减消参}：对两个方程进行适当的代数运算。
\end{enumerate}

\textbf{说明}：消参法的本质是从两个方程 \(x = f(t), y = g(t)\) 中消去参数 \(t\)。常用方法包括：代入法（从一式解出 \(t\) 代入另一式）、三角恒等式法（利用 \(\sin^2\theta + \cos^2\theta = 1\) 等）、代数消元法。\(\blacksquare\)

\textbf{注意：} 消参时需注意参数的取值范围，消参后所得方程对应的曲线可能与原参数方程对应的曲线不完全相同（可能多出或缺少部分点）。

\subsection{27.1.2 直线的参数方程}\label{ux76f4ux7ebfux7684ux53c2ux6570ux65b9ux7a0b}

\textbf{定理 27.2.} 过定点 \(P_0(x_0, y_0)\)、倾斜角为 \(\alpha\)（\(0 \leq \alpha < \pi\)）的直线的参数方程为

\[\begin{cases} x = x_0 + t\cos\alpha \\ y = y_0 + t\sin\alpha \end{cases} \quad (t \in \mathbb{R})\]

\textbf{证明}： 设 \(P(x, y)\) 为直线上任意一点，则 \(\overrightarrow{P_0P} = (x - x_0, y - y_0)\)。直线的单位方向向量为 \(\vec{d} = (\cos\alpha, \sin\alpha)\)。\(P\) 在直线上的充要条件是 \(\overrightarrow{P_0P}\) 与 \(\vec{d}\) 共线，即存在 \(t \in \mathbb{R}\) 使得 \(\overrightarrow{P_0P} = t\vec{d}\)，即

\[x - x_0 = t\cos\alpha, \quad y - y_0 = t\sin\alpha\]

\(\blacksquare\)

\textbf{定理 27.3（参数 \(t\) 的几何意义）.} 在上述参数方程中，\(|t| = |\overrightarrow{P_0P}|\)，即 \(|t|\) 表示点 \(P\) 到点 \(P_0\) 的距离。当 \(t > 0\) 时，\(P\) 在 \(P_0\) 沿倾斜角方向的一侧；当 \(t < 0\) 时，\(P\) 在另一侧。

\textbf{证明}： \(|\overrightarrow{P_0P}| = |t| \cdot |\vec{d}| = |t| \cdot 1 = |t|\)。\(\blacksquare\)

\textbf{推论 27.1.} 设直线参数方程如上，\(P_1, P_2\) 对应的参数值分别为 \(t_1, t_2\)，则

\[|P_1 P_2| = |t_1 - t_2|\]

\textbf{证明}： \(\overrightarrow{P_1P_2} = (t_2 - t_1)(\cos\alpha, \sin\alpha)\)，故 \(|P_1P_2| = |t_2 - t_1|\)。\(\blacksquare\)

\textbf{推论 27.2.} \(P_1, P_2\) 中点对应的参数为 \(t_{M} = \frac{t_1 + t_2}{2}\)。

\textbf{证明}：中点 \(M = \frac{P_1+P_2}{2} = \frac{(P_0 + t_1\vec{d}) + (P_0 + t_2\vec{d})}{2} = P_0 + \frac{t_1+t_2}{2}\vec{d}\)，故 \(t_M = \frac{t_1+t_2}{2}\)。\(\blacksquare\)

\subsection{27.1.3 圆的参数方程}\label{ux5706ux7684ux53c2ux6570ux65b9ux7a0b}

\textbf{定理 27.4.} 圆心在原点、半径为 \(r\) 的圆的参数方程为

\[\begin{cases} x = r\cos\theta \\ y = r\sin\theta \end{cases} \quad (\theta \in [0, 2\pi))\]

\textbf{证明}： 设 \(P(x, y)\) 为圆 \(x^2 + y^2 = r^2\) 上的任意一点。令 \(\theta\) 为从 \(x\) 轴正方向到 \(\overrightarrow{OP}\) 的角，则由三角函数的定义，\(x = r\cos\theta\)，\(y = r\sin\theta\)。

反过来，对任意 \(\theta\)，\(x^2 + y^2 = r^2\cos^2\theta + r^2\sin^2\theta = r^2\)，点在圆上。\(\blacksquare\)

\textbf{定理 27.5.} 圆心在 \((a, b)\)、半径为 \(r\) 的圆的参数方程为

\[\begin{cases} x = a + r\cos\theta \\ y = b + r\sin\theta \end{cases} \quad (\theta \in [0, 2\pi))\]

\textbf{证明}： 平移变换：令 \(X = x - a\)，\(Y = y - b\)，圆方程变为 \(X^2 + Y^2 = r^2\)。由定理 27.4，\(X = r\cos\theta\)，\(Y = r\sin\theta\)。还原得 \(x = a + r\cos\theta\)，\(y = b + r\sin\theta\)。\(\blacksquare\)

\subsection{27.1.4 椭圆的参数方程}\label{ux692dux5706ux7684ux53c2ux6570ux65b9ux7a0b}

\textbf{定理 27.6.} 椭圆 \(\frac{x^2}{a^2} + \frac{y^2}{b^2} = 1\)（\(a > b > 0\)）的参数方程为

\[\begin{cases} x = a\cos\theta \\ y = b\sin\theta \end{cases} \quad (\theta \in [0, 2\pi))\]

\textbf{证明}： 将 \(x = a\cos\theta\)，\(y = b\sin\theta\) 代入椭圆方程：

\[\frac{(a\cos\theta)^2}{a^2} + \frac{(b\sin\theta)^2}{b^2} = \cos^2\theta + \sin^2\theta = 1\]

恒成立，故点 \((a\cos\theta, b\sin\theta)\) 在椭圆上。

反过来，椭圆上任一点 \((x_0, y_0)\) 满足 \(\frac{x_0^2}{a^2} + \frac{y_0^2}{b^2} = 1\)。令 \(\cos\theta = \frac{x_0}{a}\)（因 \(|x_0| \leq a\)，故 \(\left|\frac{x_0}{a}\right| \leq 1\)，\(\theta\) 存在），则 \(\sin^2\theta = 1 - \frac{x_0^2}{a^2} = \frac{y_0^2}{b^2}\)，适当选取 \(\theta\) 的象限可使 \(b\sin\theta = y_0\)。\(\blacksquare\)

\textbf{定义 27.2（离心角）.} 参数 \(\theta\) 称为椭圆上点 \(P(a\cos\theta, b\sin\theta)\) 的\textbf{离心角}（或\textbf{偏心角}）。

\textbf{注意：} 离心角 \(\theta\) \textbf{不是} \(\overrightarrow{OP}\) 与 \(x\) 轴的夹角。其几何意义是：以原点为圆心分别作半径为 \(a\) 和 \(b\) 的同心圆，从原点作与 \(x\) 轴夹角为 \(\theta\) 的射线，射线分别交大圆和小圆于 \(A\) 和 \(B\)，过 \(A\) 作 \(x\) 轴的垂线，过 \(B\) 作 \(y\) 轴的垂线，两垂线的交点即为椭圆上的点 \(P(a\cos\theta, b\sin\theta)\)。

\textbf{命题 27.1.} 椭圆的面积可用参数方程验证：\(S = \pi ab\)。

\textbf{证明}： 由参数方程，椭圆在第一象限的部分对应 \(\theta \in [0, \frac{\pi}{2}]\)。利用定积分（Ch 31）：

\[S = 4\int_0^a y\,dx = 4\int_{\pi/2}^{0} b\sin\theta \cdot (-a\sin\theta)\,d\theta = 4ab\int_0^{\pi/2} \sin^2\theta\,d\theta\]

\[= 4ab \cdot \frac{\pi}{4} = \pi ab\]

\(\blacksquare\)

\subsection{27.1.5 双曲线的参数方程}\label{ux53ccux66f2ux7ebfux7684ux53c2ux6570ux65b9ux7a0b}

\textbf{定理 27.7.} 双曲线 \(\frac{x^2}{a^2} - \frac{y^2}{b^2} = 1\) 的参数方程为

\[\begin{cases} x = a\sec\theta \\ y = b\tan\theta \end{cases} \quad \left(\theta \neq \frac{\pi}{2} + k\pi, \; k \in \mathbb{Z}\right)\]

\textbf{证明}： 代入双曲线方程：

\[\frac{a^2\sec^2\theta}{a^2} - \frac{b^2\tan^2\theta}{b^2} = \sec^2\theta - \tan^2\theta = 1\]

恒成立（利用 Ch 27 的三角恒等式 \(\sec^2\theta = 1 + \tan^2\theta\)）。\(\blacksquare\)

\textbf{注意：} 双曲线也可用双曲函数表示：\(x = a\cosh t\)，\(y = b\sinh t\)（\(t \in \mathbb{R}\)），利用恒等式 \(\cosh^2 t - \sinh^2 t = 1\)。

抛物线的参数方程是三种圆锥曲线中最简洁的。本节给出了抛物线 \(y^2 = 2px\) 的参数方程 \(x = 2pt^2\)，\(y = 2pt\)，其中参数 \(t\) 的几何意义是切线的斜率。这一参数方程的选择并非唯一------抛物线的参数化也可以采用 \(x = t\)，\(y = \pm\sqrt{2pt}\)，但前者能覆盖全部实数 \(t\) 且更为简洁。抛物线的参数方程在物理学中有重要应用：抛体运动的轨迹 \(x = v_0\cos\alpha \cdot t\)，\(y = v_0\sin\alpha \cdot t - \frac{1}{2}gt^2\) 消去时间参数 \(t\) 后恰为一条抛物线，这揭示了''抛物线''这一名称的物理来源------抛出的物体在重力作用下所走的路径。 抛物线 \(y^2 = 2px\)（\(p > 0\)）的参数方程为

\[\begin{cases} x = 2pt^2 \\ y = 2pt \end{cases} \quad (t \in \mathbb{R})\]

\textbf{证明}： 代入：\(y^2 = (2pt)^2 = 4p^2t^2 = 2p \cdot 2pt^2 = 2px\)。\(\blacksquare\)

\hrulefill

\subsection{27.2 极坐标系}\label{ux6781ux5750ux6807ux7cfb}

极坐标系是直角坐标系之外最重要的平面坐标系，由极点和极轴确定。在极坐标系中，平面上任意一点 \(P\) 由极径 \(r\)（\(P\) 到极点 \(O\) 的距离）和极角 \(\theta\)（\(OP\) 与极轴的夹角）唯一确定（除极点外）。极坐标系的优势在于它能够用极简的形式表达具有旋转对称性或径向对称性的曲线------例如，以极点为圆心的圆简单地表示为 \(r = a\)，过极点的直线表示为 \(\theta = \alpha\)。极坐标与直角坐标的互化公式 \(x = r\cos\theta, y = r\sin\theta\) 和 \(r = \sqrt{x^2 + y^2}, \tan\theta = \frac{y}{x}\)（需根据象限确定 \(\theta\)）是连接两种坐标系的基本桥梁。极坐标系在物理学和工程学中有广泛应用------雷达定位、天线方向图、旋转机械的设计都自然地使用极坐标。极坐标系的''非唯一性''（同一点可以表示为 \((r, \theta)\) 和 \((-r, \theta + \pi)\) 等不同形式）是它与直角坐标系的重要区别，在解极坐标方程时需要特别注意。

\subsection{27.2.1 极坐标系的定义}\label{ux6781ux5750ux6807ux7cfbux7684ux5b9aux4e49}

\textbf{定义 27.3（极坐标系）.} 在平面内取一个定点 \(O\)，称为\textbf{极点}；从 \(O\) 出发引一条射线 \(Ox\)，称为\textbf{极轴}；再选定一个长度单位和角度的正方向（通常取逆时针方向为正）。这样建立的坐标系称为\textbf{极坐标系}。

对于平面内任一点 \(P\)，\(|OP| = r\) 称为点 \(P\) 的\textbf{极径}，极轴 \(Ox\) 到射线 \(OP\) 的角 \(\theta\) 称为点 \(P\) 的\textbf{极角}，有序实数对 \((r, \theta)\) 称为点 \(P\) 的\textbf{极坐标}，记为 \(P(r, \theta)\)。

\textbf{约定：} 本章中，极径 \(r \geq 0\)，极角 \(\theta \in [0, 2\pi)\)（或 \((-\pi, \pi]\)）。有时也允许 \(r < 0\)，此时 \((r, \theta)\) 表示的点与 \((|r|, \theta + \pi)\) 相同。

\subsection{27.2.2 极坐标与直角坐标的互化}\label{ux6781ux5750ux6807ux4e0eux76f4ux89d2ux5750ux6807ux7684ux4e92ux5316}

\textbf{定理 27.9.} 以极点 \(O\) 为直角坐标原点，极轴 \(Ox\) 为 \(x\) 轴正方向建立直角坐标系，则

\textbf{极坐标 \(\to\) 直角坐标：}

\[\begin{cases} x = r\cos\theta \\ y = r\sin\theta \end{cases}\]

\textbf{直角坐标 \(\to\) 极坐标：}

\[\begin{cases} r = \sqrt{x^2 + y^2} \\ \tan\theta = \frac{y}{x} \quad (x \neq 0) \end{cases}\]

\textbf{证明}： 由极坐标的定义，\(r = |OP|\)，\(\theta\) 为 \(OP\) 与极轴的夹角。在直角三角形中，\(x = r\cos\theta\)，\(y = r\sin\theta\)。

反过来，\(r = |OP| = \sqrt{x^2 + y^2}\)，\(\tan\theta = \frac{y}{x}\)（需根据点所在象限确定 \(\theta\) 的值）。\(\blacksquare\)

\textbf{注意：} 从直角坐标转极坐标时，\(\tan\theta = \frac{y}{x}\) 不能唯一确定 \(\theta\)（相差 \(\pi\) 的整数倍），需根据 \((x, y)\) 所在象限判断。具体地：

\begin{itemize}
\tightlist
\item
  \((x > 0, y \geq 0)\)：\(\theta = \arctan\frac{y}{x}\)
\item
  \((x < 0)\)：\(\theta = \arctan\frac{y}{x} + \pi\)
\item
  \((x > 0, y < 0)\)：\(\theta = \arctan\frac{y}{x} + 2\pi\)
\item
  \((x = 0, y > 0)\)：\(\theta = \frac{\pi}{2}\)
\item
  \((x = 0, y < 0)\)：\(\theta = \frac{3\pi}{2}\)
\end{itemize}

\hrulefill

\subsection{27.3 常见曲线的极坐标方程}\label{ux5e38ux89c1ux66f2ux7ebfux7684ux6781ux5750ux6807ux65b9ux7a0b}

极坐标系的一个显著优势是，许多经典的平面曲线在极坐标下具有极其简洁的表达式，而在直角坐标下则往往复杂得多。玫瑰线（Rose Curve）\(r = a\cos(n\theta)\) 或 \(r = a\sin(n\theta)\) 是极坐标下最美丽的曲线之一------当 \(n\) 为奇数时，曲线有 \(n\) 个花瓣；当 \(n\) 为偶数时，曲线有 \(2n\) 个花瓣。阿基米德螺线 \(r = a\theta\) 描述了极径随极角线性增长的过程，在机械设计（如凸轮轮廓）和自然界（如鹦鹉螺壳的近似描述）中都有出现。对数螺线（等角螺线）\(r = ae^{b\theta}\) 以''以恒定角度切割所有径向线''著称，这一''自相似''性质使其在生物学（如鹦鹉螺壳、向日葵花序）和艺术中有广泛应用。心形线 \(r = a(1 + \cos\theta)\) 和双纽线 \(r^2 = a^2\cos(2\theta)\) 等曲线则展示了极坐标方程在表达对称性方面的独特能力。本节将系统介绍这些经典曲线及其极坐标方程的推导与图像特征。

\subsection{27.3.1 过极点的直线}\label{ux8fc7ux6781ux70b9ux7684ux76f4ux7ebf}

\textbf{定理 27.10.} 过极点且与极轴成角 \(\alpha\) 的直线的极坐标方程为

\[\theta = \alpha\]

（允许 \(r < 0\) 时，\(\theta = \alpha\) 和 \(\theta = \alpha + \pi\) 表示同一条直线。）

\textbf{证明}： 该直线上所有点的极角均为 \(\alpha\)（或 \(\alpha + \pi\)，若允许 \(r < 0\)）。\(\blacksquare\)

\subsection{27.3.2 圆的极坐标方程}\label{ux5706ux7684ux6781ux5750ux6807ux65b9ux7a0b}

\textbf{定理 27.11.} 以直角坐标 \((a, 0)\) 为圆心、半径为 \(a\) 的圆（过极点，圆心在极轴上）的极坐标方程为

\[r = 2a\cos\theta\]

\textbf{证明}： 圆的直角坐标方程为 \((x - a)^2 + y^2 = a^2\)，展开得 \(x^2 + y^2 = 2ax\)。代入 \(x = r\cos\theta\)，\(x^2 + y^2 = r^2\)：

\[r^2 = 2ar\cos\theta\]

当 \(r \neq 0\) 时，\(r = 2a\cos\theta\)。\(r = 0\) 对应极点，也在圆上（\(\theta = \frac{\pi}{2}\) 时 \(r = 0\)）。\(\blacksquare\)

\textbf{定理 27.12.} 更一般地：

\begin{itemize}
\tightlist
\item
  圆心在直角坐标 \((0, a)\)、半径为 \(a\) 的圆（过极点，圆心在 \(y\) 轴上）的极坐标方程为 \(r = 2a\sin\theta\)。
\item
  圆心在极坐标 \((r_0, \theta_0)\)、半径为 \(a\) 的圆的极坐标方程为 \(r^2 - 2r_0 r\cos(\theta - \theta_0) + r_0^2 = a^2\)。
\end{itemize}

\textbf{证明（第二种情况）：} 圆的直角坐标方程为 \(x^2 + (y - a)^2 = a^2\)，即 \(x^2 + y^2 = 2ay\)。代入极坐标 \(r^2 = 2ar\sin\theta\)，故 \(r = 2a\sin\theta\)。

\textbf{证明（第三种情况）：} 圆心的直角坐标为 \((r_0\cos\theta_0, r_0\sin\theta_0)\)，圆的直角坐标方程为

\[(x - r_0\cos\theta_0)^2 + (y - r_0\sin\theta_0)^2 = a^2\]

展开：\(x^2 + y^2 - 2r_0(x\cos\theta_0 + y\sin\theta_0) + r_0^2 = a^2\)。

代入 \(x^2 + y^2 = r^2\)，\(x = r\cos\theta\)，\(y = r\sin\theta\)：

\[r^2 - 2r_0 r(\cos\theta\cos\theta_0 + \sin\theta\sin\theta_0) + r_0^2 = a^2\]

\[r^2 - 2r_0 r\cos(\theta - \theta_0) + r_0^2 = a^2\]

\(\blacksquare\)

\subsection{27.3.3 心形线}\label{ux5fc3ux5f62ux7ebf}

\textbf{定义 27.4（心形线）.} 极坐标方程 \(r = a(1 + \cos\theta)\)（\(a > 0\)）表示的曲线称为\textbf{心形线}（cardioid）。

\textbf{性质：}

\begin{enumerate}
\def\labelenumi{\arabic{enumi}.}
\tightlist
\item
  \textbf{对称性}：关于极轴（\(x\) 轴）对称。
\item
  \textbf{极值}：\(\theta = 0\) 时 \(r = 2a\)（最大），\(\theta = \pi\) 时 \(r = 0\)（过极点）。
\item
  \textbf{非负性}：\(r \geq 0\) 对所有 \(\theta\) 成立。
\end{enumerate}

\textbf{证明（对称性）：} \(r(-\theta) = a(1 + \cos(-\theta)) = a(1 + \cos\theta) = r(\theta)\)，故关于极轴对称。\(\blacksquare\)

\subsection{27.3.4 玫瑰线}\label{ux73abux7470ux7ebf}

\textbf{定义 27.5（玫瑰线）.} 极坐标方程 \(r = a\cos(n\theta)\)（\(a > 0\)，\(n\) 为正整数）表示的曲线称为\textbf{玫瑰线}（rose curve）。

\textbf{性质：}

\begin{itemize}
\tightlist
\item
  当 \(n\) 为奇数时，玫瑰线有 \(n\) 个花瓣。
\item
  当 \(n\) 为偶数时，玫瑰线有 \(2n\) 个花瓣。
\end{itemize}

例如：\(r = a\cos(2\theta)\) 有四个花瓣，\(r = a\cos(3\theta)\) 有三个花瓣。

\subsection{27.3.5 阿基米德螺旋线}\label{ux963fux57faux7c73ux5fb7ux87baux65cbux7ebf}

\textbf{定义 27.6（阿基米德螺旋线）.} 极坐标方程 \(r = a\theta\)（\(a > 0\)，\(\theta \geq 0\)）表示的曲线称为\textbf{阿基米德螺旋线}。

每转一圈（\(\theta\) 增加 \(2\pi\)），极径增加 \(2\pi a\)，即螺旋线以恒定速率向外展开。

\hrulefill

\subsection{27.4 圆锥曲线的统一极坐标方程}\label{ux5706ux9525ux66f2ux7ebfux7684ux7edfux4e00ux6781ux5750ux6807ux65b9ux7a0b}

将焦点置于极点、准线垂直于极轴时，三种圆锥曲线可以用一个统一的极坐标方程表示：\(r = \frac{ep}{1 - e\cos\theta}\)。这一公式是解析几何中最优美的统一结果之一------它仅用两个参数（离心率 \(e\) 和焦准距 \(p\)）就刻画了椭圆、抛物线和双曲线的全部形状。当 \(e = 0\) 时，\(r = 0\)，退化为一个点（极点的退化情形）；当 \(0 < e < 1\) 时，分母 \(1 - e\cos\theta\) 永远不为零，曲线封闭（椭圆）；当 \(e = 1\) 时，\(r = \frac{p}{1 - \cos\theta}\)，当 \(\theta \to 0\) 时分母趋近于 \(0\)，\(r \to \infty\)（抛物线）；当 \(e > 1\) 时，\(1 - e\cos\theta = 0\) 给出 \(\cos\theta = 1/e\) 的两个解，对应两条渐近线的方向（双曲线）。这一统一公式不仅揭示了圆锥曲线的本质联系，也使得极坐标成为研究圆锥曲线焦点性质的天然工具------许多与焦点相关的计算（如焦点弦长、焦半径之比）在极坐标下变得极为简洁。本节将完整推导这一统一公式，并讨论其在计算焦点弦长、判断曲线类型等方面的应用。

\subsection{27.4.1 统一极坐标方程的推导}\label{ux7edfux4e00ux6781ux5750ux6807ux65b9ux7a0bux7684ux63a8ux5bfc}

\textbf{定理 27.13（圆锥曲线的统一极坐标方程）.} 以焦点为极点，以焦点到准线的方向为极轴的反方向，则圆锥曲线的极坐标方程为

\[r = \frac{ep}{1 - e\cos\theta}\]

其中 \(e\) 为离心率，\(p\) 为焦点到准线的距离（焦准距）。

\begin{itemize}
\tightlist
\item
  当 \(0 < e < 1\) 时，曲线为\textbf{椭圆}；
\item
  当 \(e = 1\) 时，曲线为\textbf{抛物线}；
\item
  当 \(e > 1\) 时，曲线为\textbf{双曲线}（右支）。
\end{itemize}

\textbf{证明}： 设焦点 \(F\) 为极点，准线 \(l\) 在焦点的 \(\theta = \pi\) 方向（左侧）。设 \(P(r, \theta)\) 为曲线上的点。

由圆锥曲线的统一定义（Ch 325.1），\(|PF| = e \cdot d(P, l)\)。

在极坐标中，\(|PF| = r\)。准线 \(l\) 在焦点左侧距离 \(p\) 处。点 \(P\) 的直角坐标为 \((r\cos\theta, r\sin\theta)\)，准线 \(l\) 的直角坐标方程为 \(x = -p\)。\(P\) 到准线 \(l\) 的距离为

\[d(P, l) = |r\cos\theta - (-p)| = p + r\cos\theta\]

（需验证 \(p + r\cos\theta > 0\)，对椭圆和抛物线成立。）

由统一定义：

\[r = e(p + r\cos\theta) = ep + er\cos\theta\]

\[r - er\cos\theta = ep\]

\[r(1 - e\cos\theta) = ep\]

\[r = \frac{ep}{1 - e\cos\theta}\]

\(\blacksquare\)

\subsection{27.4.2 与标准方程的关系}\label{ux4e0eux6807ux51c6ux65b9ux7a0bux7684ux5173ux7cfb}

\textbf{推论 27.3.} 对于焦点在原点的椭圆 \(\frac{x^2}{a^2} + \frac{y^2}{b^2} = 1\)（\(a > b > 0\)），若取右焦点为极点，则

\[r = \frac{a(1 - e^2)}{1 - e\cos\theta} = \frac{b^2/a}{1 - e\cos\theta}\]

\textbf{证明}： 由椭圆的标准方程，右焦点为 \((c, 0)\)，准线为 \(x = \frac{a^2}{c} = \frac{a}{e}\)。焦准距

\[p = \frac{a}{e} - c = \frac{a - ce}{e} = \frac{a - a e^2}{e} = \frac{a(1 - e^2)}{e}\]

代入统一方程：

\[r = \frac{ep}{1 - e\cos\theta} = \frac{a(1 - e^2)}{1 - e\cos\theta}\]

又 \(b^2 = a^2 - c^2 = a^2(1 - e^2)\)，故 \(a(1 - e^2) = \frac{b^2}{a}\)。\(\blacksquare\)

\textbf{推论 27.4.} 对于抛物线 \(y^2 = 2px\)，取焦点为极点，则

\[r = \frac{p}{1 - \cos\theta}\]

\textbf{证明}： \(e = 1\)，焦准距 \(p\) 不变，代入统一方程即得。\(\blacksquare\)

\textbf{推论 27.5.} 对于双曲线 \(\frac{x^2}{a^2} - \frac{y^2}{b^2} = 1\)，取右焦点为极点，则

\[r = \frac{a(e^2 - 1)}{1 - e\cos\theta}\]

\textbf{证明}： 右焦点 \((c, 0)\)，准线 \(x = \frac{a^2}{c} = \frac{a}{e}\)，焦准距 \(p = c - \frac{a}{e} = \frac{ce - a}{e} = \frac{a(e^2 - 1)}{e}\)。

代入统一方程：\(r = \frac{ep}{1 - e\cos\theta} = \frac{a(e^2 - 1)}{1 - e\cos\theta}\)。\(\blacksquare\)

\subsection{27.4.3 焦点弦长的统一公式}\label{ux7126ux70b9ux5f26ux957fux7684ux7edfux4e00ux516cux5f0f-1}

\textbf{定理 27.14（焦点弦长公式）.} 过焦点的弦 \(AB\)（\(A\) 对应 \(\theta_1\)，\(B\) 对应 \(\theta_2 = \theta_1 + \pi\)），弦长为

\[|AB| = r_{A} + r_{B} = \frac{2ep}{1 - e^2\cos^2\theta_1}\]

\textbf{证明}： \(B\) 的极角为 \(\theta_1 + \pi\)，故 \(\cos\theta_2 = \cos(\theta_1 + \pi) = -\cos\theta_1\)。

\[r_{A} = \frac{ep}{1 - e\cos\theta_1}, \quad r_{B} = \frac{ep}{1 - e\cos(\theta_1 + \pi)} = \frac{ep}{1 + e\cos\theta_1}\]

\[|AB| = r_{A} + r_{B} = \frac{ep}{1 - e\cos\theta_1} + \frac{ep}{1 + e\cos\theta_1}\]

\[= ep \cdot \frac{(1 + e\cos\theta_1) + (1 - e\cos\theta_1)}{(1 - e\cos\theta_1)(1 + e\cos\theta_1)} = \frac{2ep}{1 - e^2\cos^2\theta_1}\]

\(\blacksquare\)

\textbf{推论 27.6.} 当 \(\theta_1 = \frac{\pi}{2}\)（即弦垂直于极轴方向）时，\(\cos\theta_1 = 0\)，\(|AB| = 2ep\)。这是\textbf{通径}（正焦弦）的长度。

\textbf{证明}：令 \(\theta_1 = \frac{\pi}{2}\)，\(\cos\theta_1 = 0\)，代入焦点弦公式 \(|AB| = \frac{2ep}{1 - e^2\cos^2\theta_1} = \frac{2ep}{1} = 2ep\)。 \(\blacksquare\)

\begin{itemize}
\tightlist
\item
  椭圆：通径 \(= 2ep = 2a(1 - e^2) = \frac{2b^2}{a}\)。
\item
  抛物线：通径 \(= 2ep = 2p\)（与定理 26.15 一致）。
\item
  双曲线：通径 \(= 2ep = 2a(e^2 - 1) = \frac{2b^2}{a}\)。
\end{itemize}

\hrulefill

\subsection{27.5 极坐标方程的应用}\label{ux6781ux5750ux6807ux65b9ux7a0bux7684ux5e94ux7528}

极坐标在解决某些几何问题时具有得天独厚的优势，特别是那些涉及旋转对称性或径向距离的问题。一个问题什么时候用极坐标比用直角坐标更好？关键的判断标准是：如果问题的描述中反复出现''到某定点的距离''或''绕某点旋转的角度''，极坐标通常是更自然的选择。例如，求证''以椭圆焦点为端点的两条焦半径的倒数之和为定值''------在极坐标下，这等价于证明 \(\frac{1}{r(\theta)} + \frac{1}{r(\theta + \pi)}\) 为定值，代入统一极坐标方程 \(r = \frac{ep}{1 - e\cos\theta}\) 即可简洁地完成。焦半径的计算是极坐标最典型的应用场景------直角坐标下求焦点到曲线上点的距离需要开方，极坐标下直接读取 \(r\) 即为该距离。此外，极坐标在求解某些二次曲线的最值问题时也特别有用，因为 \(r\) 作为 \(\theta\) 的函数可以直接求导找到极值点。本节通过若干典型问题展示极坐标方法在解决焦点相关几何问题中的强大威力。

\subsection{27.5.1 极坐标下的面积公式}\label{ux6781ux5750ux6807ux4e0bux7684ux9762ux79efux516cux5f0f}

\textbf{定理 27.15（极坐标面积公式）.} 由射线 \(\theta = \alpha\)、\(\theta = \beta\)（\(\alpha < \beta\)）和曲线 \(r = f(\theta)\) 围成的扇形区域的面积为

\[S = \frac{1}{2}\int_\alpha^\beta r^2\,d\theta = \frac{1}{2}\int_\alpha^\beta [f(\theta)]^2\,d\theta\]

\textbf{证明思路：} 此公式的严格证明需要定积分的理论（Ch 51）。基本思想是将扇形分割为 \(n\) 个小扇形：将 \([\alpha, \beta]\) 等分为 \(\theta_0 = \alpha < \theta_1 < \cdots < \theta_n = \beta\)，每个小区间 \([\theta_{i-1}, \theta_i]\) 对应的小扇形面积近似为 \(\frac{1}{2}[f(\xi_i)]^2 \Delta\theta\)（半径为 \(f(\xi_i)\)、张角为 \(\Delta\theta = \frac{\beta - \alpha}{n}\) 的扇形面积）。取极限 \(n \to \infty\) 即得定积分 \(\frac{1}{2}\int_\alpha^\beta [f(\theta)]^2\,d\theta\)。\(\blacksquare\)

\textbf{注}：此公式可视为极坐标下的''面积微元法''：在角度 \(\theta\) 处取微小张角 \(d\theta\)，对应的小扇形面积为 \(dS = \frac{1}{2}r^2\,d\theta\)，累加（积分）即得总面积。

\textbf{应用 27.1（心形线围成的面积）.} 心形线 \(r = a(1 + \cos\theta)\) 围成的面积为

\[S = \frac{1}{2}\int_0^{2\pi} a^2(1 + \cos\theta)^2\,d\theta = \frac{a^2}{2}\int_0^{2\pi}(1 + 2\cos\theta + \cos^2\theta)\,d\theta\]

\[= \frac{a^2}{2}\left[2\pi + 0 + \pi\right] = \frac{3\pi a^2}{2}\]

（利用 \(\int_0^{2\pi}\cos\theta\,d\theta = 0\)，\(\int_0^{2\pi}\cos^2\theta\,d\theta = \pi\)。）

\subsection{27.5.2 极坐标下曲线的切线}\label{ux6781ux5750ux6807ux4e0bux66f2ux7ebfux7684ux5207ux7ebf}

\textbf{定理 27.16.} 设曲线的极坐标方程为 \(r = f(\theta)\)，则曲线在点 \((r, \theta)\) 处的切线与径向（从极点到该点的方向）之间的夹角 \(\psi\) 满足

\[\tan\psi = \frac{r}{r'}\]

其中 \(r' = \frac{dr}{d\theta}\)。

\textbf{证明}： 将极坐标方程化为参数方程 \(x = r\cos\theta\)，\(y = r\sin\theta\)（\(\theta\) 为参数）。则

\[\frac{dx}{d\theta} = r'\cos\theta - r\sin\theta, \quad \frac{dy}{d\theta} = r'\sin\theta + r\cos\theta\]

切线的方向向量为 \(\left(\frac{dx}{d\theta}, \frac{dy}{d\theta}\right)\)。径向方向向量为 \((\cos\theta, \sin\theta)\)。夹角 \(\psi\) 满足

\[\tan\psi = \frac{|(r'\sin\theta + r\cos\theta)\cos\theta - (r'\cos\theta - r\sin\theta)\sin\theta|}{|(r'\cos\theta - r\sin\theta)\cos\theta + (r'\sin\theta + r\cos\theta)\sin\theta|}\]

化简分子：\(r'\sin\theta\cos\theta + r\cos^2\theta - r'\cos\theta\sin\theta + r\sin^2\theta = r\)。

化简分母：\(r'\cos^2\theta - r\sin\theta\cos\theta + r'\sin^2\theta + r\cos\theta\sin\theta = r'\)。

故 \(\tan\psi = \frac{r}{r'}\)。\(\blacksquare\)

\hrulefill

\subsection{27.6 三种方程的对比与选择}\label{ux4e09ux79cdux65b9ux7a0bux7684ux5bf9ux6bd4ux4e0eux9009ux62e9}

直角坐标方程、参数方程和极坐标方程是描述平面曲线的三种基本方式，它们各有优劣，适用于不同类型的问题。直角坐标方程 \(F(x, y) = 0\) 是最通用的形式，适合分析曲线的对称性、截距、渐近线等整体性质，但处理涉及''距离''和''角度''的问题时往往繁琐。参数方程 \(x = f(t), y = g(t)\) 最适合描述''运动''------质点的轨迹、曲线的切线方向（因为有 \(\frac{dy}{dx} = \frac{g'(t)}{f'(t)}\)），以及那些无法用单值函数 \(y = f(x)\) 表示的曲线。极坐标方程 \(r = f(\theta)\) 在处理旋转对称性、焦点性质和径向距离相关的问题时具有天然优势，但处理涉及直角坐标的线性运算（如平移、旋转）时不如直角坐标方便。三种坐标系之间可以通过互化公式相互转换，选择哪种坐标系取决于问题的''自然对称性''------问题的对称性决定了最优的坐标系。本节通过对比分析，建立选择坐标系的''经验法则''，帮助读者在面对具体问题时做出最优选择。

\subsection{27.6.1 方程形式的对比}\label{ux65b9ux7a0bux5f62ux5f0fux7684ux5bf9ux6bd4}

\begin{longtable}[]{@{}
  >{\raggedright\arraybackslash}p{(\linewidth - 6\tabcolsep) * \real{0.0984}}
  >{\raggedright\arraybackslash}p{(\linewidth - 6\tabcolsep) * \real{0.3607}}
  >{\raggedright\arraybackslash}p{(\linewidth - 6\tabcolsep) * \real{0.1475}}
  >{\raggedright\arraybackslash}p{(\linewidth - 6\tabcolsep) * \real{0.3934}}@{}}
\toprule
\begin{minipage}[b]{\linewidth}\raggedright
特征
\end{minipage} & \begin{minipage}[b]{\linewidth}\raggedright
普通方程 \(F(x,y) = 0\)
\end{minipage} & \begin{minipage}[b]{\linewidth}\raggedright
参数方程
\end{minipage} & \begin{minipage}[b]{\linewidth}\raggedright
极坐标方程 \(r = f(\theta)\)
\end{minipage} \\
\midrule
\endhead
\bottomrule
\endlastfoot
基本思想 & \(x, y\) 的直接关系 & 引入辅助变量 \(t\) & 用距离和角度定位 \\
优点 & 直观、通用 & 引入自由度，便于求最值和轨迹 & 利用对称性，便于求面积 \\
适用场景 & 一般代数问题 & 最值、轨迹、多变量代换 & 中心对称图形、焦点弦 \\
典型应用 & 韦达定理联立 & 圆/椭圆上的三角代换 & 心形线、玫瑰线、焦点弦 \\
\end{longtable}

\subsection{27.6.2 转换策略}\label{ux8f6cux6362ux7b56ux7565}

\begin{enumerate}
\def\labelenumi{\arabic{enumi}.}
\tightlist
\item
  遇到圆/椭圆上的点，考虑用参数方程（\(\cos\theta, \sin\theta\) 代换）。
\item
  涉及焦点弦长、焦点距离，考虑极坐标方程。
\item
  涉及过定点的弦长，考虑直线的参数方程（利用参数 \(t\) 的几何意义）。
\item
  一般情况下优先使用普通方程。
\end{enumerate}

\subsection{27.6.3 方程转换的数学基础}\label{ux65b9ux7a0bux8f6cux6362ux7684ux6570ux5b66ux57faux7840}

\textbf{定理 27.17.} 设曲线的参数方程为 \(x = f(t)\)，\(y = g(t)\)，则消去参数 \(t\) 可得普通方程 \(F(x, y) = 0\)。反过来，普通方程可能有多种参数化方式。

\textbf{证明}： 由 \(x = f(t)\) 得 \(t = f^{-1}(x)\)（若 \(f\) 有反函数），代入 \(y = g(t) = g(f^{-1}(x))\)，得 \(y\) 关于 \(x\) 的关系，即普通方程。

反之，给定普通方程 \(F(x, y) = 0\)，可能有多种参数化。例如，单位圆 \(x^2 + y^2 = 1\) 可以参数化为 \((\cos\theta, \sin\theta)\)，也可以参数化为 \(\left(\frac{1 - t^2}{1 + t^2}, \frac{2t}{1 + t^2}\right)\)（万能代换）。\(\blacksquare\)

\textbf{定理 27.18.} 参数方程与极坐标方程之间的转换关系为：设参数方程为 \(x = f(t)\)，\(y = g(t)\)，则极坐标方程为

\[r = \sqrt{[f(t)]^2 + [g(t)]^2}, \quad \tan\theta = \frac{g(t)}{f(t)}\]

消去 \(t\) 即得 \(r\) 关于 \(\theta\) 的关系。

\textbf{说明}：此为坐标转换的直接应用。由 \(x = r\cos\theta, y = r\sin\theta\) 和参数方程 \(x = f(t), y = g(t)\) 联立即得 \(r = \sqrt{f(t)^2 + g(t)^2}\)，\(\tan\theta = g(t)/f(t)\)。\(\blacksquare\)

\subsection{27.6.4 曲线系与参数}\label{ux66f2ux7ebfux7cfbux4e0eux53c2ux6570}

\textbf{定义 27.7（曲线系）.} 含参数 \(\lambda\) 的方程 \(F(x, y) + \lambda G(x, y) = 0\) 表示的曲线族称为\textbf{曲线系}。当 \(\lambda\) 取不同值时，曲线系中的每一条曲线都经过 \(F(x, y) = 0\) 和 \(G(x, y) = 0\) 的交点。

\textbf{定理 27.19（曲线系定理）.} 设曲线 \(C_1: F(x, y) = 0\) 和 \(C_2: G(x, y) = 0\) 有 \(n\) 个交点，则曲线系 \(F(x, y) + \lambda G(x, y) = 0\) 中的每一条曲线都经过这 \(n\) 个交点。

\textbf{证明}： 设 \((x_0, y_0)\) 是 \(C_1\) 和 \(C_2\) 的交点，即 \(F(x_0, y_0) = 0\) 且 \(G(x_0, y_0) = 0\)。则 \(F(x_0, y_0) + \lambda G(x_0, y_0) = 0 + \lambda \cdot 0 = 0\)，故 \((x_0, y_0)\) 在曲线系中的每一条曲线上。\(\blacksquare\)

\textbf{应用 27.2（过两圆交点的圆系）.} 设 \(\odot C_1: x^2 + y^2 + D_1 x + E_1 y + F_1 = 0\) 和 \(\odot C_2: x^2 + y^2 + D_2 x + E_2 y + F_2 = 0\) 相交于两点，则圆系

\[(x^2 + y^2 + D_1 x + E_1 y + F_1) + \lambda(x^2 + y^2 + D_2 x + E_2 y + F_2) = 0\]

经过这两点。当 \(\lambda = -1\) 时，二次项消去，得到根轴（定理 25.18）。

\subsection{27.6.5 极坐标下曲线的对称性判别}\label{ux6781ux5750ux6807ux4e0bux66f2ux7ebfux7684ux5bf9ux79f0ux6027ux5224ux522b}

\textbf{命题 27.2.} 设曲线的极坐标方程为 \(r = f(\theta)\)，则：

\begin{enumerate}
\def\labelenumi{\arabic{enumi}.}
\tightlist
\item
  若 \(f(-\theta) = f(\theta)\)，则曲线关于极轴对称。
\item
  若 \(f(\pi - \theta) = f(\theta)\)，则曲线关于直线 \(\theta = \frac{\pi}{2}\) 对称。
\item
  若 \(f(\theta + \pi) = f(\theta)\)，或 \(f(\theta + \pi) = -f(\theta)\)，则曲线关于极点对称。
\end{enumerate}

\textbf{证明}：

\textbf{（1）} 设 \(P(r, \theta)\) 在曲线上，即 \(r = f(\theta)\)。则 \(P' = (r, -\theta)\) 是 \(P\) 关于极轴的对称点。\(P'\) 对应的极径为 \(f(-\theta) = f(\theta) = r\)，故 \(P'\) 也在曲线上。

\textbf{（2）} \(P'' = (r, \pi - \theta)\) 是 \(P\) 关于 \(\theta = \frac{\pi}{2}\) 的对称点。\(f(\pi - \theta) = f(\theta) = r\)，故 \(P''\) 也在曲线上。

\textbf{（3）} \(P''' = (r, \theta + \pi)\) 是 \(P\) 关于极点的对称点（在极坐标中，\((r, \theta + \pi)\) 表示极径方向相反的点）。\(f(\theta + \pi) = f(\theta) = r\)（或 \(-f(\theta)\)，此时取 \(|r|\)），故 \(P'''\) 也在曲线上。\(\blacksquare\)

\textbf{应用 27.3.} 心形线 \(r = a(1 + \cos\theta)\) 关于极轴对称（因 \(\cos(-\theta) = \cos\theta\)）。玫瑰线 \(r = a\cos(2\theta)\) 关于极轴和 \(\theta = \frac{\pi}{2}\) 均对称。

\hrulefill

\begin{quote}
\textbf{Ch 325 本章展望}

本章引入了参数方程和极坐标方程两种描述曲线的方式，完成了对平面曲线表示方法的系统研究。参数方程通过引入辅助变量，以运动的观点描述曲线，在求最值、轨迹问题中具有独特优势。极坐标方程以极点和极轴为参考，特别适合处理具有旋转对称性的曲线和焦点相关问题。圆锥曲线的统一极坐标方程 \(r = \frac{ep}{1 - e\cos\theta}\) 将椭圆、抛物线和双曲线统一在一个公式之下，揭示了圆锥曲线最深刻的内在统一性。

至此，解析几何的三大主题------直线与圆（Ch 325）、圆锥曲线（Ch 325）、参数方程与极坐标方程（Ch 325）------构成了一个完整的体系。这些工具将为后续的离散数学与概率统计（卷六）提供几何直觉和分析方法的支撑。
\end{quote}

\hrulefill

\hrulefill

\emph{本卷用坐标方法统一处理经典几何问题，从直线与圆的代数刻画，到圆锥曲线的系统研究，再到参数方程与极坐标方程的参数化表达，为几何问题的解析化处理提供了完整的方法论体系。}

\emph{卷五十七：解析几何终。}

\emph{卷五十七：解析几何终。}

\emph{卷五十七：解析几何终。}

\hrulefill

\newpage

\chapter{03-函数与微积分/02-解析几何/06-特殊函数论.md}\label{ux51fdux6570ux4e0eux5faeux79efux520602-ux89e3ux6790ux51e0ux4f5506-ux7279ux6b8aux51fdux6570ux8bba.md}

\chapter{特殊函数论}\label{ux7279ux6b8aux51fdux6570ux8bba}

\hrulefill

\section{第61章：Gamma函数与相关函数}\label{ux7b2c61ux7ae0gammaux51fdux6570ux4e0eux76f8ux5173ux51fdux6570}

特殊函数（Special Functions）是数学分析中的一类具有特殊性质的超越函数，它们在数学物理、数论、概率论和微分方程中频繁出现。Gamma函数和Beta函数构成了最基本的一类特殊函数，其理论是18-19世纪数学分析的伟大成就之一。

\subsection{136.1 Gamma函数的定义与基本性质}\label{gammaux51fdux6570ux7684ux5b9aux4e49ux4e0eux57faux672cux6027ux8d28}

\textbf{定义 136.1.1（Gamma函数------Euler积分定义）} 对 \(\operatorname{Re}(z) > 0\)，Gamma函数由Euler第二类积分定义：

\[\Gamma(z) = \int_0^\infty t^{z-1} e^{-t} dt.\]

此积分在右半平面 \(\operatorname{Re}(z) > 0\) 内绝对收敛，定义了该区域内的全纯函数。

\textbf{命题 136.1.2（函数方程）} 对 \(\operatorname{Re}(z) > 0\)，

\[\Gamma(z + 1) = z \Gamma(z).\]

特别地，\(\Gamma(1) = \int_0^\infty e^{-t} dt = 1\)，故对正整数 \(n\)，\(\Gamma(n) = (n - 1)!\)。Gamma函数是阶乘函数的解析延拓。

\textbf{证明}：由分部积分：

\[\Gamma(z + 1) = \int_0^\infty t^z e^{-t} dt = \left[-t^z e^{-t}\right]_0^\infty + z \int_0^\infty t^{z-1} e^{-t} dt = z \Gamma(z).\]

边界项在 \(0\) 处消失因为 \(\operatorname{Re}(z) > 0\)，在 \(\infty\) 处消失因为指数衰减。\(\blacksquare\)

\textbf{定理 136.1.3（解析延拓与极点结构）} \(\Gamma(z)\) 可解析延拓为 \(\mathbb{C} \setminus \{0, -1, -2, \ldots\}\) 上的亚纯函数，在 \(z = -n\)（\(n = 0, 1, 2, \ldots\)）处有单极点，留数为 \(\operatorname{Res}(\Gamma, -n) = (-1)^n / n!\)。

\textbf{证明}：由函数方程 \(\Gamma(z) = \Gamma(z + n + 1) / [z(z + 1) \cdots (z + n)]\)。右边对 \(\operatorname{Re}(z) > -(n + 1)\) 是亚纯的，且与左边在公共定义域上一致。令 \(n \to \infty\) 即得全平面的亚纯延拓。极点为 \(z = 0, -1, -2, \ldots\)，留数计算由归纳法可得。\(\blacksquare\)

\textbf{定理 136.1.4（Weierstrass无穷乘积）} Gamma函数的Weierstrass乘积表示为：

\[\frac{1}{\Gamma(z)} = z e^{\gamma z} \prod_{n=1}^\infty \left(1 + \frac{z}{n}\right) e^{-z/n},\]

其中 \(\gamma = \lim_{n \to \infty} \left(\sum_{k=1}^n \frac{1}{k} - \log n\right) \approx 0.5772\) 为Euler-Mascheroni常数。此乘积在整个复平面上定义了一个整函数（\(1/\Gamma(z)\)），直接表明 \(\Gamma(z)\) 没有零点。

\textbf{证明概要}：由Euler积分在不同表示之间的转换，利用Gauss乘积公式：

\[\Gamma(z) = \lim_{n \to \infty} \frac{n! n^z}{z(z+1)\cdots(z+n)}.\]

取倒数并利用调和级数的渐近展开即得Weierstrass形式。\(\blacksquare\)

\textbf{定理 136.1.5（互补公式）}

\[\Gamma(z) \Gamma(1 - z) = \frac{\pi}{\sin(\pi z)}, \quad z \notin \mathbb{Z}.\]

特别地，\(\Gamma(1/2) = \sqrt{\pi}\)。

\textbf{证明}：由Weierstrass乘积，

\[\frac{1}{\Gamma(z)\Gamma(1 - z)} = z(1 - z) e^{\gamma z} e^{\gamma(1 - z)} \prod_{n=1}^\infty \left(1 + \frac{z}{n}\right)\left(1 + \frac{1 - z}{n}\right) e^{-z/n} e^{-(1-z)/n}.\]

简化后得 \(\frac{1}{\Gamma(z)\Gamma(1 - z)} = z \prod_{n=1}^\infty \left(1 - \frac{z^2}{n^2}\right)\)。利用正弦函数的无穷乘积 \(\sin(\pi z) = \pi z \prod_{n=1}^\infty \left(1 - \frac{z^2}{n^2}\right)\)，即得所需结果。\(\blacksquare\)

\textbf{定理 136.1.6（Legendre加倍公式）}

\[\Gamma(2z) = \frac{2^{2z-1}}{\sqrt{\pi}} \Gamma(z) \Gamma\left(z + \frac{1}{2}\right).\]

更一般地，Gauss乘法公式为：

\[\Gamma(nz) = (2\pi)^{(1-n)/2} n^{nz - 1/2} \prod_{k=0}^{n-1} \Gamma\left(z + \frac{k}{n}\right).\]

\subsection{136.2 Beta函数及其与Gamma函数的关系}\label{betaux51fdux6570ux53caux5176ux4e0egammaux51fdux6570ux7684ux5173ux7cfb}

\textbf{定义 136.2.1（Beta函数）} 对 \(\operatorname{Re}(x) > 0\)，\(\operatorname{Re}(y) > 0\)，Beta函数定义为Euler第一类积分：

\[B(x, y) = \int_0^1 t^{x-1} (1 - t)^{y-1} dt.\]

此积分在指定区域内绝对收敛。

\textbf{定理 136.2.2（Beta函数与Gamma函数的关系）} 对 \(\operatorname{Re}(x) > 0\)，\(\operatorname{Re}(y) > 0\)，

\[B(x, y) = \frac{\Gamma(x) \Gamma(y)}{\Gamma(x + y)}.\]

\textbf{完整证明}：计算乘积 \(\Gamma(x)\Gamma(y)\)：

\[\Gamma(x)\Gamma(y) = \int_0^\infty \int_0^\infty s^{x-1} t^{y-1} e^{-(s+t)} ds dt.\]

作变量替换 \(s = uv\)，\(t = u(1 - v)\)，其中 \(u \in (0, \infty)\)，\(v \in (0, 1)\)。Jacobi行列式为 \(\frac{\partial(s, t)}{\partial(u, v)} = u\)。故

\[\Gamma(x)\Gamma(y) = \int_0^\infty \int_0^1 (uv)^{x-1} (u(1 - v))^{y-1} e^{-u} u dv du = \int_0^\infty u^{x+y-1} e^{-u} du \cdot \int_0^1 v^{x-1} (1 - v)^{y-1} dv = \Gamma(x + y) \cdot B(x, y).\]

两端除以 \(\Gamma(x + y) \neq 0\) 即得 \(B(x, y) = \Gamma(x)\Gamma(y) / \Gamma(x + y)\)。\(\blacksquare\)

\textbf{推论 136.2.3} Beta函数对称：\(B(x, y) = B(y, x)\)。由Gamma函数表示即证。

\textbf{定理 136.2.4（Dirichlet积分）} 对 \(\operatorname{Re}(a_i) > 0\)，

\[\int \cdots \int_{t_i \geq 0, \sum t_i \leq 1} t_1^{a_1-1} \cdots t_n^{a_n-1} (1 - \sum t_i)^{a_{n+1}-1} dt_1 \cdots dt_n = \frac{\Gamma(a_1) \cdots \Gamma(a_{n+1})}{\Gamma(a_1 + \cdots + a_{n+1})}.\]

这是Beta函数的高维推广，在Dirichlet分布和多元统计中有重要应用。

\subsection{136.3 函数方程与Stirling公式}\label{ux51fdux6570ux65b9ux7a0bux4e0estirlingux516cux5f0f}

\textbf{定理 136.3.1（Stirling公式）} 对 \(\Gamma(z)\)，当 \(|z| \to \infty\)（\(|\arg z| < \pi - \delta\)，\(\delta > 0\)），

\[\log \Gamma(z) \sim \left(z - \frac{1}{2}\right) \log z - z + \frac{1}{2} \log(2\pi) + \sum_{k=1}^\infty \frac{B_{2k}}{2k(2k - 1)} z^{-(2k - 1)},\]

其中 \(B_{2k}\) 为Bernoulli数。前几项为：

\[\Gamma(z) = \sqrt{2\pi} z^{z - 1/2} e^{-z} \left(1 + \frac{1}{12z} + \frac{1}{288z^2} - \frac{139}{51840z^3} + \cdots\right).\]

\textbf{完整证明}：证明分若干步骤。

\textbf{第一步：积分表示。} 由 \(\Gamma(z + 1) = z \Gamma(z)\)，可设 \(\operatorname{Re}(z) > 0\) 并将 \(\Gamma(z)\) 写为：

\[\Gamma(z) = \int_0^\infty t^{z-1} e^{-t} dt = z^z \int_0^\infty e^{z(\log u - u)} \frac{du}{u},\]

其中做了变量替换 \(t = zu\)。

\textbf{第二步：鞍点分析。} 函数 \(f(u) = \log u - u\) 在 \(u = 1\) 处达到最大值 \(f(1) = -1\)，\(f''(1) = -1\)。考虑主导项：在 \(u = 1\) 附近作Taylor展开：

\[\log u - u = -1 - \frac{(u - 1)^2}{2} + \frac{(u - 1)^3}{3} - \cdots.\]

将积分变量改为 \(v\)（其中 \(u = 1 + v/\sqrt{z}\)），则主要贡献来自 \(|v| \leq \log z\)。展开被积函数：

\[\int e^{z(\log u - u)} \frac{du}{u} = e^{-z} \int_{-\sqrt{z}}^{\infty} \exp\left(-\frac{v^2}{2} + \frac{v^3}{3\sqrt{z}} - \cdots\right) \frac{dv}{\sqrt{z}(1 + v/\sqrt{z})}.\]

\textbf{第三步：渐近展开。} 将指数中的高阶项展开为 \(\frac{1}{\sqrt{z}}\) 的幂级数，利用Gauss积分 \(\int_{-\infty}^\infty e^{-v^2/2} v^{2k} dv = \sqrt{2\pi} \cdot (2k - 1)!!\) 逐项积分。整理各项即得Stirling级数。特别地，主导项给出：

\[\Gamma(z) \sim \sqrt{2\pi} z^{z - 1/2} e^{-z}.\]

\textbf{第四步：误差估计。} Binet第二公式给出了余项的精确表达式，证明了Stirling级数是渐近级数而非收敛级数。\(\blacksquare\)

\textbf{推论 136.3.2（阶乘的Stirling近似）}

\[n! \sim \sqrt{2\pi n} \left(\frac{n}{e}\right)^n, \quad n \to \infty.\]

更精确地，对任意 \(n \geq 1\)，

\[\sqrt{2\pi n} \left(\frac{n}{e}\right)^n e^{\frac{1}{12n + 1}} < n! < \sqrt{2\pi n} \left(\frac{n}{e}\right)^n e^{\frac{1}{12n}}.\]

\subsection{136.4 多伽玛函数与Barnes G-函数}\label{ux591aux4f3dux739bux51fdux6570ux4e0ebarnes-g-ux51fdux6570}

\textbf{定义 136.4.1（多伽玛函数）} 对数Gamma函数的导数称为digamma函数：

\[\psi(z) = \frac{d}{dz} \log \Gamma(z) = \frac{\Gamma'(z)}{\Gamma(z)}.\]

其高阶导数称为polygamma函数：\(\psi^{(m)}(z) = \frac{d^{m+1}}{dz^{m+1}} \log \Gamma(z)\)。

\textbf{命题 136.4.2（Digamma函数的级数表示）}

\[\psi(z) = -\gamma - \frac{1}{z} + \sum_{n=1}^\infty \left(\frac{1}{n} - \frac{1}{n + z}\right).\]

更一般地，\(\psi^{(m)}(z) = (-1)^{m+1} m! \sum_{n=0}^\infty \frac{1}{(n + z)^{m+1}}\)（\(m \geq 1\)），由此 \(\psi^{(m)}\) 与Hurwitz zeta函数 \(\zeta(m + 1, z)\) 的关系为 \(\psi^{(m)}(z) = (-1)^{m+1} m! \zeta(m + 1, z)\)。

\textbf{定理 136.4.3（Digamma函数在有理点的值------Gauss digamma定理）} 对有理数 \(p/q\)（\(0 < p < q\)），

\[\psi\left(\frac{p}{q}\right) = -\gamma - \log(2q) - \frac{\pi}{2} \cot\frac{\pi p}{q} + 2 \sum_{k=1}^{\lfloor (q-1)/2 \rfloor} \cos\frac{2\pi k p}{q} \cdot \log \sin\frac{\pi k}{q}.\]

\textbf{定义 136.4.4（Barnes G-函数）} Barnes G-函数 \(G(z)\) 定义为满足 \(G(z + 1) = \Gamma(z) G(z)\) 且 \(G(1) = 1\) 的整函数。其Weierstrass乘积表示为：

\[G(z + 1) = (2\pi)^{z/2} \exp\left(-\frac{z + z^2(1 + \gamma)}{2}\right) \prod_{k=1}^\infty \left(1 + \frac{z}{k}\right)^k \exp\left(-z + \frac{z^2}{2k}\right).\]

Barnes G-函数与Selberg积分公式、随机矩阵理论中的Tracy-Widom分布和 \(\zeta\) 函数的矩有深层联系。

\subsection{136.5 Gamma函数在概率与数论中的应用}\label{gammaux51fdux6570ux5728ux6982ux7387ux4e0eux6570ux8bbaux4e2dux7684ux5e94ux7528}

\textbf{定理 136.5.1（Riemann zeta函数的函数方程）} Riemann zeta函数

\[\zeta(s) = \sum_{n=1}^\infty \frac{1}{n^s}, \quad \operatorname{Re}(s) > 1\]

可解析延拓为 \(\mathbb{C} \setminus \{1\}\) 上的亚纯函数，且满足函数方程：

\[\pi^{-s/2} \Gamma\left(\frac{s}{2}\right) \zeta(s) = \pi^{-(1-s)/2} \Gamma\left(\frac{1 - s}{2}\right) \zeta(1 - s).\]

定义完备的zeta函数 \(\xi(s) = \frac{1}{2} s(s - 1) \pi^{-s/2} \Gamma(s/2) \zeta(s)\)，则 \(\xi(s) = \xi(1 - s)\)，且 \(\xi(s)\) 是阶为 \(1\) 的整函数。

\textbf{证明概要}：利用 \(\zeta(s)\) 的积分表示：\(\zeta(s) \Gamma(s/2) \pi^{-s/2} = \int_0^\infty \theta(it) t^{s/2 - 1} dt\)，其中 \(\theta(z) = \sum_{n=-\infty}^\infty e^{\pi i n^2 z}\) 为Jacobi theta函数。利用theta函数的模变换性质 \(\theta(-1/z) = \sqrt{z/i} \cdot \theta(z)\) 即得对称性。\(\blacksquare\)

\textbf{定理 136.5.2（Gamma函数与Gauss和的联系）} Gauss和 \(G(a, b) = \sum_{k=0}^{b-1} e^{2\pi i a k^2 / b}\) 可通过Gamma函数的特殊值表示，其绝对值由 \(\sqrt{b}\) 给出。

\textbf{定理 136.5.3（概率论中的应用）} 许多重要概率分布的密度函数由Gamma函数表达：Gamma分布 \(\Gamma(\alpha, \beta)\) 的密度为 \(f(x) \propto x^{\alpha-1} e^{-\beta x}\)；Beta分布由Beta函数给出；Student \(t\) 分布的归一化常数涉及 \(\Gamma((\nu+1)/2) / \Gamma(\nu/2)\)；\(\chi^2\) 分布是 \(\Gamma(n/2, 1/2)\) 的特例。

\hrulefill

\section{第62章：超几何函数与正交多项式}\label{ux7b2c62ux7ae0ux8d85ux51e0ux4f55ux51fdux6570ux4e0eux6b63ux4ea4ux591aux9879ux5f0f}

超几何函数是最重要的超越函数之一，绝大多数经典特殊函数均可视为其特例或极限情形。正交多项式构成了调和分析和数值逼近论的基石。

\subsection{137.1 超几何级数与Gauss超几何函数}\label{ux8d85ux51e0ux4f55ux7ea7ux6570ux4e0egaussux8d85ux51e0ux4f55ux51fdux6570}

\textbf{定义 137.1.1（超几何级数）} Gauss超几何函数由超几何级数定义：

\[{}_2F_1(a, b; c; z) = \sum_{n=0}^\infty \frac{(a)_n (b)_n}{(c)_n n!} z^n,\]

其中 \((a)_n = a(a + 1) \cdots (a + n - 1) = \Gamma(a + n) / \Gamma(a)\) 为Pochhammer符号（升阶乘）。此级数在 \(|z| < 1\) 内绝对收敛，\(|z| = 1\) 时条件收敛当 \(\operatorname{Re}(c - a - b) > 0\)。

\textbf{命题 137.1.2（收敛半径与解析延拓）} 超几何级数在 \(|z| < 1\) 内收敛，且在 \(z = 1\) 处的性态由Gauss判别法决定。\({}_2F_1\) 满足超几何微分方程：

\[z(1 - z) \frac{d^2 w}{dz^2} + [c - (a + b + 1)z] \frac{dw}{dz} - ab w = 0.\]

此方程在 \(\mathbb{CP}^1\) 上有三个正则奇点：\(z = 0, 1, \infty\)，是Fuchs型微分方程的典型。任意具有三个正则奇点的二阶Fuchs型方程均可化为此形式。

\textbf{定理 137.1.3（Gauss求和定理）} 对 \(\operatorname{Re}(c - a - b) > 0\)，

\[{}_2F_1(a, b; c; 1) = \frac{\Gamma(c) \Gamma(c - a - b)}{\Gamma(c - a) \Gamma(c - b)}.\]

\textbf{完整证明}：由Euler积分表示，对 \(\operatorname{Re}(c) > \operatorname{Re}(b) > 0\) 和 \(|z| < 1\)（或适当解析延拓），

\[{}_2F_1(a, b; c; z) = \frac{\Gamma(c)}{\Gamma(b) \Gamma(c - b)} \int_0^1 t^{b-1} (1 - t)^{c-b-1} (1 - zt)^{-a} dt.\]

令 \(z = 1\)，被积函数化为 \(t^{b-1} (1 - t)^{c - a - b - 1}\)。由Beta函数积分：

\[\int_0^1 t^{b-1} (1 - t)^{c-a-b-1} dt = B(b, c - a - b) = \frac{\Gamma(b) \Gamma(c - a - b)}{\Gamma(c - a)}.\]

代入系数 \(\frac{\Gamma(c)}{\Gamma(b)\Gamma(c - b)}\) 即得结论。条件 \(\operatorname{Re}(c - a - b) > 0\) 保证了积分在 \(t = 1\) 处的收敛性。\(\blacksquare\)

\textbf{推论 137.1.4} 由Gauss求和定理可得许多著名级数恒等式，包括Chu-Vandermonde恒等式和Dixon恒等式。

\textbf{定理 137.1.5（超几何函数的变换公式）} Gauss超几何函数满足Pfaff变换：

\[{}_2F_1(a, b; c; z) = (1 - z)^{-a} {}_2F_1\left(a, c - b; c; \frac{z}{z - 1}\right).\]

以及Euler变换：

\[{}_2F_1(a, b; c; z) = (1 - z)^{c - a - b} {}_2F_1(c - a, c - b; c; z).\]

这些变换公式将超几何函数的值从 \(|z| < 1\) 解析延拓到整个复平面（除去分支切割 \([1, \infty)\)）。

\textbf{定义 137.1.6（广义超几何函数）} \(p\)-上参数 \(q\)-下参数的广义超几何函数定义为：

\[{}_pF_q(a_1, \ldots, a_p; b_1, \ldots, b_q; z) = \sum_{n=0}^\infty \frac{(a_1)_n \cdots (a_p)_n}{(b_1)_n \cdots (b_q)_n n!} z^n.\]

当 \(p \leq q\) 时级数对所有 \(z\) 收敛（整函数）；\(p = q + 1\) 时在 \(|z| < 1\) 收敛；\(p > q + 1\) 时仅在 \(z = 0\) 收敛（除非级数截断）。

\subsection{137.2 合流超几何函数与Bessel函数}\label{ux5408ux6d41ux8d85ux51e0ux4f55ux51fdux6570ux4e0ebesselux51fdux6570}

\textbf{定义 137.2.1（合流超几何函数）} 在 \({}_2F_1(a, b; c; z)\) 中令 \(b \to \infty\) 并同时缩放 \(z \mapsto z/b\)，得到Kummer合流超几何函数：

\[{}_1F_1(a; c; z) = \sum_{n=0}^\infty \frac{(a)_n}{(c)_n n!} z^n = M(a, c, z).\]

此级数对所有 \(z \in \mathbb{C}\) 收敛（整函数）。合流超几何函数满足Kummer方程：

\[z \frac{d^2 w}{dz^2} + (c - z) \frac{dw}{dz} - a w = 0.\]

该方程在 \(z = 0\) 有正则奇点，在 \(z = \infty\) 有非正则奇点（合并了两个正则奇点，因此得名''合流''）。

\textbf{定义 137.2.2（Bessel函数）} 第一类Bessel函数 \(J_\nu(z)\) 可由超几何函数表示为：

\[J_\nu(z) = \frac{(z/2)^\nu}{\Gamma(\nu + 1)} {}_0F_1\left(; \nu + 1; -\frac{z^2}{4}\right) = \sum_{k=0}^\infty \frac{(-1)^k}{k! \Gamma(k + \nu + 1)} \left(\frac{z}{2}\right)^{2k + \nu}.\]

\(J_\nu(z)\) 满足Bessel微分方程：

\[z^2 \frac{d^2 w}{dz^2} + z \frac{dw}{dz} + (z^2 - \nu^2) w = 0.\]

\textbf{定理 137.2.3（Bessel函数的生成函数）} Bessel函数 \(J_n(z)\)（\(n \in \mathbb{Z}\)）有生成函数：

\[\exp\left(\frac{z}{2}(t - t^{-1})\right) = \sum_{n=-\infty}^\infty J_n(z) t^n.\]

\textbf{证明}：将左边展开为Laurent级数。左边 \(= e^{zt/2} \cdot e^{-z/(2t)} = \sum_{k=0}^\infty \frac{(z/2)^k}{k!} t^k \cdot \sum_{m=0}^\infty \frac{(-z/2)^m}{m!} t^{-m}\)。合并 \(t^n\) 的系数（令 \(k - m = n\)）：

\[J_n(z) = \sum_{m=0}^\infty \frac{(-1)^m}{m! (m + n)!} \left(\frac{z}{2}\right)^{2m + n}\]

对 \(n \geq 0\) 即为第一类Bessel函数。对 \(n < 0\) 由 \(J_{-n}(z) = (-1)^n J_n(z)\)。\(\blacksquare\)

\textbf{定理 137.2.4（Bessel函数的渐近展开）} 对固定的 \(\nu\) 和 \(|z| \to \infty\)（\(|\arg z| < \pi\)），

\[J_\nu(z) \sim \sqrt{\frac{2}{\pi z}} \left[\cos\left(z - \frac{\nu \pi}{2} - \frac{\pi}{4}\right) \sum_{k=0}^\infty \frac{(-1)^k a_{2k}(\nu)}{z^{2k}} - \sin\left(z - \frac{\nu \pi}{2} - \frac{\pi}{4}\right) \sum_{k=0}^\infty \frac{(-1)^k a_{2k+1}(\nu)}{z^{2k+1}}\right],\]

其中系数 \(a_k(\nu)\) 由特定的递推关系给出。此渐近展开揭示了Bessel函数的振荡衰减行为。

\subsection{137.3 正交多项式的一般理论}\label{ux6b63ux4ea4ux591aux9879ux5f0fux7684ux4e00ux822cux7406ux8bba}

\textbf{定义 137.3.1（正交多项式）} 设 \(\mu\) 为 \(\mathbb{R}\) 上的正Borel测度，其所有矩 \(\int t^n d\mu(t)\) 有限。若多项式序列 \(\{P_n(x)\}_{n=0}^\infty\)（\(\deg P_n = n\)）满足：

\[\int_{\mathbb{R}} P_n(x) P_m(x) d\mu(x) = 0, \quad n \neq m,\]

则称为关于测度 \(\mu\) 的正交多项式。

\textbf{定理 137.3.2（三项递推关系）} 任意首一正交多项式序列 \(\{P_n\}\) 满足三项递推关系：

\[P_{n+1}(x) = (x - \alpha_n) P_n(x) - \beta_n P_{n-1}(x), \quad n \geq 0,\]

其中 \(P_{-1} = 0\)，\(P_0 = 1\)，\(\alpha_n, \beta_n \in \mathbb{R}\) 且 \(\beta_n > 0\)（\(n \geq 1\)）。系数由下式给出：

\[\alpha_n = \frac{\int x P_n(x)^2 d\mu(x)}{\int P_n(x)^2 d\mu(x)}, \quad \beta_n = \frac{\int P_n(x)^2 d\mu(x)}{\int P_{n-1}(x)^2 d\mu(x)}.\]

\textbf{证明}：因 \(x P_n(x)\) 是 \(n + 1\) 次多项式，可在 \(\{P_k\}_{k=0}^{n+1}\) 下展开。由于正交性，对 \(k < n - 1\)，

\[\int x P_n(x) P_k(x) d\mu(x) = \int P_n(x) (x P_k(x)) d\mu(x) = 0\]

因 \(x P_k(x)\) 的次数 \(\leq n - 1\)。故 \(x P_n(x)\) 仅与 \(P_{n-1}, P_n, P_{n+1}\) 有非零内积，由此得三项递推。\(\blacksquare\)

\textbf{定理 137.3.3（Favard定理）} 给定实数序列 \(\{\alpha_n\}, \{\beta_n\}\) 且 \(\beta_n > 0\)，由三项递推生成的单项多项式序列是关于某个正Borel测度的正交多项式当且仅当 \(\beta_n > 0\) 对所有 \(n \geq 1\)。此即Favard定理，建立了正交多项式与三项递推之间的等价关系。

\textbf{定理 137.3.4（Christoffel-Darboux公式）} 对正交集 \(\{p_n\}\)（未必首一），

\[\sum_{k=0}^n \frac{p_k(x) p_k(y)}{h_k} = \frac{k_n}{k_{n+1} h_n} \frac{p_{n+1}(x) p_n(y) - p_n(x) p_{n+1}(y)}{x - y},\]

其中 \(h_k = \int p_k^2 d\mu\)，\(k_n\) 为 \(p_n\) 的首项系数。当 \(x \to y\) 时取极限即得Christoffel-Darboux恒等式：

\[\sum_{k=0}^n \frac{p_k(x)^2}{h_k} = \frac{k_n}{k_{n+1} h_n} \left[p_{n+1}'(x) p_n(x) - p_n'(x) p_{n+1}(x)\right].\]

\textbf{证明}：由三项递推，将 \((x - y) \sum_{k=0}^n \cdots\) 表示为叠缩和，详细展开即得。\(\blacksquare\)

\subsection{137.4 经典正交多项式族}\label{ux7ecfux5178ux6b63ux4ea4ux591aux9879ux5f0fux65cf}

\textbf{定义 137.4.1（Legendre多项式）} Legendre多项式 \(P_n(x)\) 关于 \([-1, 1]\) 上的Lebesgue测度正交：

\[P_n(x) = \frac{1}{2^n n!} \frac{d^n}{dx^n} (x^2 - 1)^n \quad (\text{Rodrigues公式}).\]

满足Legendre微分方程：\((1 - x^2) P_n'' - 2x P_n' + n(n+1) P_n = 0\)。生成函数：\(\frac{1}{\sqrt{1 - 2xt + t^2}} = \sum_{n=0}^\infty P_n(x) t^n\)。

\textbf{定义 137.4.2（Hermite多项式）} Hermite多项式 \(H_n(x)\) 关于 \(\mathbb{R}\) 上带权重 \(e^{-x^2}\) 的测度正交：

\[H_n(x) = (-1)^n e^{x^2} \frac{d^n}{dx^n} e^{-x^2}.\]

满足Hermite微分方程：\(H_n'' - 2x H_n' + 2n H_n = 0\)。生成函数：\(e^{2xt - t^2} = \sum_{n=0}^\infty \frac{H_n(x)}{n!} t^n\)。Hermite多项式在量子力学（谐振子波函数）中自然出现。

\textbf{定义 137.4.3（Laguerre多项式）} Laguerre多项式 \(L_n^{(\alpha)}(x)\) 关于 \([0, \infty)\) 上带权重 \(x^\alpha e^{-x}\) 的测度正交：

\[L_n^{(\alpha)}(x) = \frac{x^{-\alpha} e^x}{n!} \frac{d^n}{dx^n} (x^{n+\alpha} e^{-x}).\]

满足Laguerre微分方程：\(x L_n'' + (\alpha + 1 - x) L_n' + n L_n = 0\)。在氢原子的径向波函数中起核心作用。

\textbf{定义 137.4.4（Chebyshev多项式）} Chebyshev多项式 \(T_n(x)\) 和 \(U_n(x)\) 分别关于 \([-1, 1]\) 上带权重 \((1 - x^2)^{-1/2}\) 和 \(\sqrt{1 - x^2}\) 的测度正交：

\[T_n(\cos \theta) = \cos(n\theta), \quad U_n(\cos \theta) = \frac{\sin((n+1)\theta)}{\sin \theta}.\]

三项递推：\(T_{n+1}(x) = 2x T_n(x) - T_{n-1}(x)\)，\(T_0 = 1\)，\(T_1 = x\)。Chebyshev多项式在逼近论中极为重要------它们在 \([-1, 1]\) 上极小化最大偏差。

\textbf{定义 137.4.5（Jacobi多项式）} Jacobi多项式 \(P_n^{(\alpha, \beta)}(x)\) 关于 \([-1, 1]\) 上带权重 \((1 - x)^\alpha (1 + x)^\beta\) 的测度正交，其中 \(\alpha, \beta > -1\)。它是经典正交多项式中最一般的族，包含Legendre（\(\alpha = \beta = 0\)）、Chebyshev（\(\alpha = \beta = \pm 1/2\)）和Gegenbauer（\(\alpha = \beta\)）作为特例。Rodrigues公式为：

\[P_n^{(\alpha, \beta)}(x) = \frac{(-1)^n}{2^n n!} (1 - x)^{-\alpha} (1 + x)^{-\beta} \frac{d^n}{dx^n} \left[(1 - x)^{\alpha + n} (1 + x)^{\beta + n}\right].\]

\subsection{137.5 椭圆函数初步}\label{ux692dux5706ux51fdux6570ux521dux6b65}

\textbf{定义 137.5.1（椭圆积分）} 第一类和第二类不完全椭圆积分定义为：

\[F(\phi, k) = \int_0^\phi \frac{d\theta}{\sqrt{1 - k^2 \sin^2 \theta}}, \quad E(\phi, k) = \int_0^\phi \sqrt{1 - k^2 \sin^2 \theta} d\theta.\]

完全椭圆积分：\(K(k) = F(\pi/2, k)\)，\(E(k) = E(\pi/2, k)\)。

\textbf{定义 137.5.2（Jacobi椭圆函数）} Jacobi椭圆函数 \(\operatorname{sn}(u, k)\)，\(\operatorname{cn}(u, k)\)，\(\operatorname{dn}(u, k)\) 由反演椭圆积分定义：若 \(u = F(\phi, k)\)，则 \(\operatorname{sn} u = \sin \phi\)，\(\operatorname{cn} u = \cos \phi\)，\(\operatorname{dn} u = \sqrt{1 - k^2 \sin^2 \phi}\)。它们满足基本恒等式：

\[\operatorname{sn}^2 u + \operatorname{cn}^2 u = 1, \quad k^2 \operatorname{sn}^2 u + \operatorname{dn}^2 u = 1.\]

Jacobi椭圆函数是双周期的亚纯函数，其周期平行四边形由 \(4K(k)\) 和 \(4iK'(k)\) 决定。

\textbf{定义 137.5.3（Weierstrass \(\wp\) 函数）} 对格 \(\Lambda = \mathbb{Z}\omega_1 + \mathbb{Z}\omega_2\)（\(\operatorname{Im}(\omega_2/\omega_1) > 0\)），Weierstrass \(\wp\) 函数定义为：

\[\wp(z; \Lambda) = \frac{1}{z^2} + \sum_{\omega \in \Lambda \setminus \{0\}} \left(\frac{1}{(z - \omega)^2} - \frac{1}{\omega^2}\right).\]

\(\wp\) 是偶的、双周期的亚纯函数，在格点处有二重极点。它满足微分方程：

\[(\wp'(z))^2 = 4\wp(z)^3 - g_2 \wp(z) - g_3,\]

其中 \(g_2 = 60 \sum_{\omega \neq 0} \omega^{-4}\)，\(g_3 = 140 \sum_{\omega \neq 0} \omega^{-6}\) 为Eisenstein级数。此微分方程表明 \((\wp, \wp')\) 参数化了一条椭圆曲线 \(y^2 = 4x^3 - g_2 x - g_3\)，建立了椭圆函数与椭圆曲线之间的深刻联系。

\textbf{定理 137.5.4（椭圆函数与模形式的联系）} \(j\)-不变量 \(j(\tau) = 1728 g_2^3 / (g_2^3 - 27 g_3^2)\) 是上半平面上的模函数，满足 \(j(\tau + 1) = j(\tau)\) 和 \(j(-1/\tau) = j(\tau)\)。\(j\)-函数的Fourier展开为：

\[j(\tau) = q^{-1} + 744 + 196884 q + 21493760 q^2 + \cdots, \quad q = e^{2\pi i \tau}.\]

这些系数与Monster单群和Moonshine理论有着神秘的联系。

\subsection{137.6 Theta函数与模变换}\label{thetaux51fdux6570ux4e0eux6a21ux53d8ux6362}

\textbf{定义 137.6.1（Jacobi theta函数）} Jacobi theta函数 \(\vartheta(z; \tau)\) 定义为：

\[\vartheta(z; \tau) = \sum_{n=-\infty}^\infty e^{\pi i n^2 \tau + 2\pi i n z}, \quad \operatorname{Im}(\tau) > 0.\]

此外还有三个辅助theta函数：\(\vartheta_1, \vartheta_2, \vartheta_3, \vartheta_4\)，分别对应求和指标的不同移位选择。Theta函数是 \(z\) 的整函数，在格 \(\mathbb{Z} + \mathbb{Z}\tau\) 的平移下满足拟周期性。

\textbf{定理 137.6.2（Theta函数的模变换）} Theta函数在模群 \(\operatorname{SL}(2, \mathbb{Z})\) 作用下的变换公式为：

\[\vartheta(z; \tau + 1) = \vartheta(z; \tau), \quad \vartheta\left(\frac{z}{\tau}; -\frac{1}{\tau}\right) = \sqrt{-i\tau} e^{\pi i z^2 / \tau} \vartheta(z; \tau).\]

更一般地，theta常数 \(\vartheta(0; \tau)\) 是权为 \(1/2\) 的模形式。Jacobi恒等式：\(\vartheta_2(0; \tau)^4 + \vartheta_4(0; \tau)^4 = \vartheta_3(0; \tau)^4\)。

\textbf{定理 137.6.3（热方程与theta函数）} \(\vartheta(z; \tau)\) 满足热方程（扩散方程）：

\[\frac{\partial \vartheta}{\partial \tau} = \frac{1}{4\pi i} \frac{\partial^2 \vartheta}{\partial z^2}.\]

此方程揭示了theta函数与热核的本质联系。事实上，theta函数是圆周上的热核的迹（trace）。

\textbf{定理 137.6.4（Riemann theta函数与代数几何）} 对亏格 \(g \geq 1\) 的紧致Riemann面，Riemann theta函数是 \(\mathbb{C}^g\) 上的多变量theta函数：

\[\Theta(\mathbf{z}; \Omega) = \sum_{\mathbf{n} \in \mathbb{Z}^g} \exp(\pi i \mathbf{n}^T \Omega \mathbf{n} + 2\pi i \mathbf{n}^T \mathbf{z}),\]

其中 \(\Omega\) 是 \(g \times g\) 的周期矩阵，属于Siegel上半空间。\(\Theta\) 函数在代数几何中起核心作用：Riemann定理将除子类群与theta函数的零点联系起来；而Mumford的theta函数理论是代数曲线模空间研究的基础。

\hrulefill

\section{第63章：特殊函数在数学物理中的应用}\label{ux7b2c63ux7ae0ux7279ux6b8aux51fdux6570ux5728ux6570ux5b66ux7269ux7406ux4e2dux7684ux5e94ux7528}

特殊函数在数学物理中的中心地位可以追溯到18世纪以来对波动方程、热传导方程和Laplace方程的研究。分离变量法导致各类常微分方程的本征值问题，其解恰为各类正交多项式、Bessel函数和合流超几何函数。

\textbf{定理 138.1（量子力学中的谐振子）} 一维量子谐振子的定态Schrödinger方程：

\[-\frac{\hbar^2}{2m} \frac{d^2 \psi}{dx^2} + \frac{1}{2} m \omega^2 x^2 \psi = E \psi.\]

经过无量纲化后变为 \(-\psi'' + x^2 \psi = \lambda \psi\)。本征函数为 \(\psi_n(x) = e^{-x^2/2} H_n(x)\)（\(H_n\)为Hermite多项式），本征值为 \(\lambda_n = 2n + 1\)。能谱的等间距性由Hermite多项式的递推关系保证。

\textbf{证明}：无量纲化：令 \(\xi = \sqrt{m\omega/\hbar} \cdot x\)，\(\psi(x) = \varphi(\xi)\)，\(E = \frac{\hbar\omega}{2} \lambda\)，则 Schrödinger 方程化为 \(-\varphi''(\xi) + \xi^2 \varphi(\xi) = \lambda \varphi(\xi)\)。设试探解 \(\varphi(\xi) = e^{-\xi^2/2} v(\xi)\)。代入方程，利用 \(\varphi' = (-\xi v + v') e^{-\xi^2/2}\)，\(\varphi'' = (\xi^2 v - v - 2\xi v' + v'') e^{-\xi^2/2}\)，得到 \(v\) 满足的方程： \[-v'' + 2\xi v' + v = \lambda v \quad \Rightarrow \quad v'' - 2\xi v' + (\lambda - 1)v = 0.\] 这正是 Hermite 微分方程 \(H_n'' - 2\xi H_n' + 2n H_n = 0\)。比较得 \(\lambda - 1 = 2n\)，即 \(\lambda_n = 2n + 1\)，\(v(\xi) = H_n(\xi)\)。由 Sturm-Liouville 理论，\(\{H_n(\xi) e^{-\xi^2/2}\}\) 在 \(L^2(\mathbb{R})\) 中构成完备正交集------这可由 Hermite 多项式的正交性 \(\int_{-\infty}^\infty H_m(\xi) H_n(\xi) e^{-\xi^2} d\xi = 2^n n! \sqrt{\pi} \cdot \delta_{mn}\) 直接验证。Hermite 多项式的三项递推 \(H_{n+1}(\xi) = 2\xi H_n(\xi) - 2n H_{n-1}(\xi)\) 保证了能谱 \(\lambda_n\) 的等间距性（相邻能级差恒为 \(2\)），这是量子谐振子能谱等距的代数根源。\(\blacksquare\)

\textbf{定理 138.2（氢原子的径向波函数）} 氢原子的径向Schrödinger方程在分离变量后导致合流超几何方程，径向波函数为：

\[R_{n\ell}(r) \propto r^\ell e^{-r/(na_0)} L_{n-\ell-1}^{(2\ell+1)}\left(\frac{2r}{na_0}\right),\]

其中 \(L_k^{(\alpha)}\) 为广义Laguerre多项式。能级 \(E_n = -13.6 / n^2\) eV 由合流超几何级数截断为多项式的条件给出。

\textbf{证明}：氢原子的径向 Schrödinger 方程（球坐标分离变量后）为： \[-\frac{\hbar^2}{2m} \frac{1}{r^2} \frac{d}{dr}\left(r^2 \frac{dR}{dr}\right) + \left[\frac{\ell(\ell+1)\hbar^2}{2mr^2} - \frac{e^2}{4\pi\varepsilon_0 r}\right]R = ER.\] 令 \(u(r) = rR(r)\)，引入无量纲变量 \(\rho = r/a_0\)（\(a_0\) 为 Bohr 半径）和能量参数 \(\kappa = \sqrt{-2mE/\hbar^2}\)，方程化为： \[\frac{d^2 u}{d\rho^2} = \left[\frac{\ell(\ell+1)}{\rho^2} - \frac{2}{\rho} + \frac{1}{n^2}\right]u,\] 其中 \(E = -\frac{me^4}{2(4\pi\varepsilon_0)^2 \hbar^2 n^2}\)。令 \(u(\rho) = \rho^{\ell+1} e^{-\rho/n} v(\rho)\)，代入后 \(v(\rho)\) 满足 Kummer 合流超几何方程： \[\rho v'' + (2\ell+2 - 2\rho/n)v' - (\ell+1 - n)v = 0.\] 通解为 \(v(\rho) = A \cdot {}_1F_1(\ell+1-n; 2\ell+2; 2\rho/n) + B \cdot U(\ell+1-n; 2\ell+2; 2\rho/n)\)。平方可积条件要求波函数在无穷远处衰减，这迫使第一类合流超几何级数 \({}_1F_1\) 必须截断为多项式，即要求第一参数 \(\ell+1-n\) 为非正整数：\(\ell+1-n = -k\)（\(k = 0, 1, 2, \ldots\)），故 \(n = k + \ell + 1\) 为正整数。此时 \({}_1F_1(-k; 2\ell+2; z)\) 退化为广义 Laguerre 多项式 \(L_k^{(2\ell+1)}(z)\)（模常数因子），得到 \(v(\rho) \propto L_{n-\ell-1}^{(2\ell+1)}(2\rho/n)\)。由此，\(R_{n\ell}(r) \propto r^\ell e^{-r/(na_0)} L_{n-\ell-1}^{(2\ell+1)}(2r/na_0)\)。能级由截断条件决定的 \(n\) 给出：\(E_n = -13.6/n^2\) eV。\(\blacksquare\)

\textbf{定理 138.3（Bessel函数在波动物理中的应用）} 圆形鼓膜的振动模由Bessel函数 \(J_n(kr) e^{i n \theta}\) 描述。Dirichlet边界条件 \(J_n(kR) = 0\) 给出本征频率。各阶Bessel函数的零点决定了鼓膜的音高和泛音结构。在电磁波导理论、光纤通信和天线设计中，Bessel函数同样不可或离。

\textbf{证明}：二维波动方程 \(\frac{\partial^2 u}{\partial t^2} = c^2 \nabla^2 u\) 在半径为 \(R\) 的圆盘上分离变量。设 \(u(r, \theta, t) = T(t) \cdot R(r) \cdot \Theta(\theta)\)，得 \(T'' + \omega^2 T = 0\)（\(\omega\) 为角频率），\(\Theta'' + n^2 \Theta = 0\)（\(n\) 为整数，由周期性 \(2\pi\) 条件），径向方程： \[r^2 R'' + r R' + (k^2 r^2 - n^2) R = 0, \quad k = \omega/c.\] 令 \(\xi = kr\)，得标准 Bessel 方程 \(\xi^2 R'' + \xi R' + (\xi^2 - n^2)R = 0\)。通解为 \(R(\xi) = A J_n(\xi) + B Y_n(\xi)\)。\(Y_n\) 在 \(\xi = 0\) 处发散，故 \(B = 0\)（有限性条件）。Dirichlet 边界条件 \(R(kR) = 0\) 给出 \(J_n(kR) = 0\)。记 \(J_n\) 的第 \(m\) 个正零点为 \(j_{n,m}\)，则 \(k_{n,m} = j_{n,m} / R\)，本征频率为 \(\omega_{n,m} = c \cdot j_{n,m} / R\)。基频由最小的零点 \(j_{0,1} \approx 2.4048\) 给出。不同 \((n, m)\) 对应不同的振型（nodal lines 为径向节线和角向节线的组合），零点数列 \(\{j_{n,m}\}\) 决定了整个频谱结构。\(\blacksquare\)

\textbf{定理 138.4（椭球坐标系与特殊函数）} 在椭球坐标系中分离Helmholtz方程 \((\nabla^2 + k^2) \psi = 0\) 得到椭球波函数（spheroidal wave functions），它们是Legendre函数和Bessel函数的推广。椭球波函数在信号处理的扁球体波函数（prolate spheroidal wave functions）中具有最优的时频集中性质，是Slepian基的数学基础。

\textbf{定理 138.5（theta函数与可积系统）} Theta函数为KdV方程、非线性Schrödinger方程和 sine-Gordon方程等可积系统的拟周期解提供了统一框架。对于KdV方程 \(u_t + u_{xxx} + 6u u_x = 0\)，其 \(g\) 亏格拟周期解由Riemann theta函数表达为：

\[u(x, t) = 2 \partial_x^2 \log \Theta(\mathbf{k} x + \boldsymbol{\omega} t + \boldsymbol{\phi}_0; \Omega) + \text{常数}.\]

这些解构成有限带势（finite-gap potential）和代数几何可积性的核心内容。

\textbf{证明}：KdV 方程可表示为 Lax 对形式 \(L_t = [P, L]\)，其中 \(L = -\partial_x^2 + u(x, t)\) 为 Schrödinger 算子，\(P\) 为某三阶反对称微分算子。Lax 对的可积性意味着 \(L\) 的谱在时间演化下不变------即 \(u(x, t)\) 构成等谱流。对于''有限带''势，\(L\) 的谱由 \(g\) 个有限带和无穷多个间隙组成。对应的 Baker-Akhiezer 函数 \(\psi(x, t; P)\)（\(P\) 为超椭圆曲线 \(\Sigma\) 上的点------由 \(L\) 的谱数据定义的 Riemann 面）在 \(\Sigma\) 上有 \(g\) 个极点且满足特定的渐近条件。由 Riemann-Roch 定理，这样的函数由 Riemann theta 函数显式构造： \[\psi(x, t; P) = e^{kx + \omega t} \frac{\Theta(\mathbf{A}(P) + \mathbf{k} x + \boldsymbol{\omega} t + \boldsymbol{\phi}_0)}{\Theta(\mathbf{A}(P) + \boldsymbol{\phi}_0)},\] 其中 \(\mathbf{A}: \Sigma \to \operatorname{Jac}(\Sigma)\) 为 Abel-Jacobi 映射。取 \(x = 0\) 处的渐近展开并利用 KdV 的 Lax 表示，可得 \(u(x, t) = -2 \partial_x^2 \log \Theta(\mathbf{k} x + \boldsymbol{\omega} t + \boldsymbol{\phi}_0) + \text{常数}\)。此即 Its-Matveev 公式（1975），证明了拟周期 theta 函数解确实是 KdV 方程的精确解。\(\blacksquare\)

\textbf{定理 138.6（随机矩阵与特殊函数）} 在随机矩阵理论中，Gauss正交系综（GOE）和Gauss酉系综（GUE）的最大特征值的Tracy-Widom分布由Painleve II方程的解表达。广义Laguerre系综和Jacobi系综涉及相应正交多项式的渐近分析。Barnes G-函数出现在Selberg积分和随机矩阵的矩的精确表达式中。

\textbf{定理 138.7（Painleve超越函数与可积性）} Painleve方程是一类二阶非线性常微分方程，其解（Painleve超越函数）具有''Painleve性质''------所有可移动奇点均为极点。六种Painleve方程（PI-PVI）中，PVI最为一般，包含了全部经典Painleve方程作为合流极限。Painleve函数在随机矩阵理论（Tracy-Widom分布）、可积系统、二维量子引力和共形场论中具有基本的重要性。

\textbf{定理 138.8（椭圆函数在经典力学中的应用------单摆的精确解）} 无阻尼单摆的运动方程 \(\ddot{\theta} + \omega_0^2 \sin \theta = 0\) 的精确解由Jacobi椭圆函数表达：

\[\theta(t) = 2 \arcsin(k \operatorname{sn}(\omega_0 t, k)),\]

其中 \(k = \sin(\theta_{\max} / 2)\) 为模参数。周期为 \(T = 4 K(k) / \omega_0\)，小振幅极限下恢复为熟知的 \(T_0 = 2\pi / \omega_0\)。此例展示了椭圆函数如何自然地刻画非谐振动的解析结构。

\textbf{证明}：对运动方程 \(\ddot{\theta} + \omega_0^2 \sin \theta = 0\) 乘以 \(\dot{\theta}\) 并积分得能量守恒： \[\frac{1}{2} \dot{\theta}^2 - \omega_0^2 \cos \theta = \text{const} = -\omega_0^2 \cos \theta_{\max}.\] 令 \(k = \sin(\theta_{\max}/2)\)。引入新变量 \(\phi\) 满足 \(\sin(\theta/2) = k \sin \phi\)。微分得 \(\frac{1}{2} \cos(\theta/2) \dot{\theta} = k \cos \phi \cdot \dot{\phi}\)，故 \(\dot{\theta} = \frac{2k \cos \phi}{\sqrt{1 - k^2 \sin^2 \phi}} \dot{\phi}\)。同时 \(\cos \theta = 1 - 2 \sin^2(\theta/2) = 1 - 2k^2 \sin^2 \phi\)。代入能量方程： \[\frac{2k^2 \cos^2 \phi}{1 - k^2 \sin^2 \phi} \dot{\phi}^2 - \omega_0^2(1 - 2k^2 \sin^2 \phi) = -\omega_0^2(1 - 2k^2).\] 化简得 \(\dot{\phi}^2 = \omega_0^2 (1 - k^2 \sin^2 \phi)\)。取正号（摆从平衡位置出发），分离变量并积分： \[\int_0^\phi \frac{d\phi'}{\sqrt{1 - k^2 \sin^2 \phi'}} = \omega_0 t.\] 左端恰为第一类椭圆积分 \(F(\phi, k)\)。由 Jacobi 椭圆函数的定义 \(\phi = \operatorname{am}(\omega_0 t, k)\)，故 \(\sin \phi = \operatorname{sn}(\omega_0 t, k)\)。回代得 \(\sin(\theta/2) = k \operatorname{sn}(\omega_0 t, k)\)，即 \(\theta(t) = 2 \arcsin(k \operatorname{sn}(\omega_0 t, k))\)。周期为 \(\phi\) 从 \(0\) 到 \(2\pi\) 的时间，即 \(T = \frac{4}{\omega_0} F(\pi/2, k) = \frac{4K(k)}{\omega_0}\)。极限检验：当 \(k \to 0\)（小振幅），\(\operatorname{sn}(u, 0) = \sin u\)，\(K(0) = \pi/2\)，故 \(\theta(t) \approx 2k \sin(\omega_0 t) = \theta_{\max} \sin(\omega_0 t)\)（简谐近似），\(T \to 2\pi/\omega_0\)。\(\blacksquare\)

\textbf{定理 138.9（Airy函数与转折点问题）} Airy函数 \(\operatorname{Ai}(x)\) 由积分表示 \(\operatorname{Ai}(x) = \frac{1}{\pi} \int_0^\infty \cos(t^3/3 + xt) dt\) 定义，满足Airy微分方程 \(y'' - xy = 0\)。Airy函数在WKB近似失效的转折点处提供了全局一致渐近展开，是半经典分析和光学焦散面理论的基础工具。Airy函数的渐近行为在 \(x \to +\infty\) 时为指数衰减，在 \(x \to -\infty\) 时为振荡行为，精确刻画了经典允许区和禁戒区的过渡。

\textbf{定理 138.10（球谐函数与Laplace方程）} Laplace方程 \(\Delta u = 0\) 在球坐标下分离变量得到的角向方程为 \(Y_{\ell}^m(\theta, \phi) \propto P_\ell^m(\cos \theta) e^{i m \phi}\)，其中 \(P_\ell^m\) 为连带Legendre函数。球谐函数构成 \(L^2(\mathbb{S}^2)\) 的完备正交基，其在旋转下的变换性质由Wigner D-矩阵描述，构成了角动量量子理论和势论中多极展开的数学语言。

\textbf{证明}：Laplace 方程在球坐标 \((r, \theta, \phi)\) 下分离变量。设 \(u(r, \theta, \phi) = R(r) \cdot Y(\theta, \phi)\)，角向部分 \(Y(\theta, \phi)\) 满足： \[\frac{1}{\sin\theta} \frac{\partial}{\partial\theta} \left(\sin\theta \frac{\partial Y}{\partial\theta}\right) + \frac{1}{\sin^2\theta} \frac{\partial^2 Y}{\partial\phi^2} = -\ell(\ell+1)Y.\] 再分离 \(\theta\) 和 \(\phi\)：\(Y(\theta, \phi) = \Theta(\theta) \Phi(\phi)\)。\(\Phi\) 满足 \(\Phi'' + m^2 \Phi = 0\)，由周期性边界条件 \(\Phi(\phi+2\pi) = \Phi(\phi)\) 得 \(m \in \mathbb{Z}\)，\(\Phi_m(\phi) = e^{i m \phi}\)。\(\Theta\) 满足连带 Legendre 方程： \[(1 - \xi^2) \Theta'' - 2\xi \Theta' + \left[\ell(\ell+1) - \frac{m^2}{1-\xi^2}\right] \Theta = 0, \quad \xi = \cos\theta.\] 解为连带 Legendre 函数 \(P_\ell^m(\cos\theta)\)。为保证 \(\Theta\) 在 \(\theta = 0, \pi\)（\(\xi = \pm 1\)）处有界，需 \(\ell \in \mathbb{Z}_{\geq 0}\) 且 \(|m| \leq \ell\)。归一化的球谐函数为 \[Y_\ell^m(\theta, \phi) = \sqrt{\frac{2\ell+1}{4\pi} \frac{(\ell-m)!}{(\ell+m)!}} \, P_\ell^m(\cos\theta) e^{i m \phi}.\] \(P_\ell^m\) 可由超几何函数表达：\(P_\ell^m(x) \propto (1-x^2)^{m/2} {}_2F_1(-\ell+m, \ell+m+1; m+1; (1-x)/2)\)。完备性：\(\{Y_\ell^m\}\) 构成 \(L^2(\mathbb{S}^2)\) 的正交基------这是 Sturm-Liouville 理论在球面上的直接推论，等价于 Laplace-Beltrami 算子的谱分解。旋转变换：\(\operatorname{SO}(3)\) 在 \(L^2(\mathbb{S}^2)\) 上的表示可约化为 Wigner D-矩阵 \(D_{m'm}^\ell(\alpha, \beta, \gamma)\)，满足 \(Y_\ell^m(R^{-1}(\theta, \phi)) = \sum_{m'=-\ell}^\ell D_{m'm}^\ell(R) Y_\ell^{m'}(\theta, \phi)\)，这正是角动量理论中 \(\ell\) 为量子数、\(m\) 为磁量子数的数学表述。\(\blacksquare\)

\textbf{定理 138.11（超几何函数与共形场论）} 在二维共形场论中，四点和更高点的关联函数满足超几何型微分方程（BPZ方程）。Virasoro代数退化表示的四点关联函数由超几何函数的推广形式给出。Dotsenko-Fateev积分将共形块表达为超几何类型的广义Selberg积分。这些联系表明，19世纪发展起来的特殊函数理论在现代理论物理中依然是不可或缺的语言。

\textbf{证明}：在二维 CFT 中，Virasoro 代数的最高权表示 \(V_{c, h}\)（\(c\) 为中心荷，\(h\) 为共形维数）在 \(c = 1 - 6(p-q)^2/pq\) 时具有退化表示------存在零向量 \(\chi \in V_{c, h}\) 满足 \(L_n \chi = 0\)（\(n > 0\)）。四点关联函数利用全局共形对称可化为交叉比 \(z\) 的单变量函数 \(G(z) = \langle \phi_1(0) \phi_2(z) \phi_3(1) \phi_4(\infty) \rangle\)。当 \(\phi_2\) 的表示在水平 \(2\) 退化时，零向量条件导出 BPZ 方程： \[z(1-z) G''(z) + [c' - (a+b+1)z] G'(z) - ab G(z) = 0,\] 即 Gauss 超几何方程，参数 \(a, b, c'\) 由共形维数 \(\{h_i\}\) 和中心荷 \(c\) 表出。故退化四点关联函数为 \({}_2F_1\) 的线性组合。对一般退化表示（水平 \(rs\)），关联函数满足 \(r \times s\) 阶 Fuchs 型方程。Dotsenko-Fateev 利用 Coulomb 气体（自由玻色子 + 屏蔽荷）表示给出了积分表达式： \[G(z) = \oint_\Gamma \cdots \oint_\Gamma \prod_{i=1}^r (t_i - z)^{2\alpha_2 \alpha} \prod_{i<j} (t_i - t_j)^{2\alpha^2} \prod_{i=1}^r t_i^{2\alpha_1 \alpha} (1-t_i)^{2\alpha_3 \alpha} dt_1 \cdots dt_r,\] 其中 \(\alpha = \sqrt{2/(p(p+1))}\) 等参数由 CFT 的极小模型数据决定。此积分为广义 Selberg 积分（超几何类型的多重围道积分），构成了共形块的积分表示。\(\blacksquare\)

\textbf{定理 138.12（误差函数与扩散理论）} 误差函数 \(\operatorname{erf}(x) = \frac{2}{\sqrt{\pi}} \int_0^x e^{-t^2} dt\) 是扩散过程最基本的解函数。互补误差函数 \(\operatorname{erfc}(x) = 1 - \operatorname{erf}(x)\) 与抛物柱函数和合流超几何函数密切相关：\(\operatorname{erfc}(x) = \frac{1}{\sqrt{\pi}} e^{-x^2} U(1/2, 1/2, x^2)\)，其中 \(U\) 为第二类合流超几何函数。扩散方程的基本解涉及 \(\operatorname{erf}\)，在半导体物理（杂质扩散）、金融数学（Black-Scholes公式中的累积正态分布）和统计物理中广泛出现。

\textbf{证明}：一维扩散方程（热方程）\(\frac{\partial u}{\partial t} = D \frac{\partial^2 u}{\partial x^2}\) 的基本解为高斯核 \(G(x, t) = \frac{1}{\sqrt{4\pi Dt}} e^{-x^2/(4Dt)}\)。初值问题 \(u(x, 0) = u_0(x)\) 的解为卷积 \(u = G * u_0\)。阶跃初值 \(u_0(x) = C \cdot \mathbf{1}_{x \geq 0}\) 给出 \[u(x, t) = \frac{C}{\sqrt{4\pi Dt}} \int_0^\infty e^{-(x-y)^2/(4Dt)} dy = \frac{C}{\sqrt{\pi}} \int_{-\infty}^{x/\sqrt{4Dt}} e^{-\xi^2} d\xi = \frac{C}{2}\left[1 + \operatorname{erf}\!\left(\frac{x}{\sqrt{4Dt}}\right)\right].\] 与合流超几何函数的联系：由积分表达式 \(\operatorname{erfc}(x) = \frac{2}{\sqrt{\pi}} \int_x^\infty e^{-t^2} dt = \frac{2e^{-x^2}}{\sqrt{\pi}} \int_0^\infty e^{-u^2 - 2xu} du\)。利用 \(U(a, b, z)\) 的积分表示，令 \(a = 1/2\)，\(b = 1/2\)，经变量替换得 \(\operatorname{erfc}(x) = \frac{1}{\sqrt{\pi}} e^{-x^2} U(1/2, 1/2, x^2)\)。此外 \(\operatorname{erf}(x) = \frac{1}{\sqrt{\pi}} \gamma(1/2, x^2)\)，其中 \(\gamma(a, z) = \int_0^z t^{a-1} e^{-t} dt\) 为下不完全 Gamma 函数------直接将误差函数与 Gamma 函数理论联系起来。在 Black-Scholes 模型中，标准正态累积分布 \(N(d) = \frac{1}{2}\left[1 + \operatorname{erf}(d/\sqrt{2})\right]\) 决定了期权的到期概率，凸显了 erf 在定量金融中的基础地位。\(\blacksquare\)

\textbf{定理 138.13（Bailey链与 \(q\)-级数）} \(q\)-超几何级数 \({}_r\phi_s\) 是经典超几何函数的 \(q\)-模拟（\(q = e^{2\pi i \tau}\) 为模参数）。Bailey引理及其迭代（Bailey链）产生了无穷多个Rogers-Ramanujan型恒等式。其中最基本的Rogers-Ramanujan恒等式为：

\[\sum_{n=0}^\infty \frac{q^{n^2}}{(q; q)_n} = \prod_{n=0}^\infty \frac{1}{(1 - q^{5n+1})(1 - q^{5n+4})}.\]

这些恒等式在共形场论（极小模型的特征公式）、结理论（A-polynomial）、扭结不变量和组合数学中有深刻的应用。Andrews-Gordon恒等式和Macdonald恒等式将 \(q\)-级数与仿射Lie代数和格顶角算子模型联系起来。

\newpage

\chapter{03-函数与微积分/03-微积分初步/01-极限与连续.md}\label{ux51fdux6570ux4e0eux5faeux79efux520603-ux5faeux79efux5206ux521dux6b6501-ux6781ux9650ux4e0eux8fdeux7eed.md}

\textbf{前置知识}：本章依赖 Ch 23（函数概念）、Ch 25（指数函数与对数函数）、Ch 28（数列概念与数学归纳法）。极限理论是整个分析学的基石，是连接离散（数列）与连续（函数）的桥梁。

\hrulefill

\section{第64章：数列极限}\label{ux7b2c64ux7ae0ux6570ux5217ux6781ux9650}

数列极限是极限理论的第一站，也是整个分析学的逻辑起点。与函数极限相比，数列极限的定义域是离散的（正整数集），这使得 \(\varepsilon\)-\(N\) 语言比 \(\varepsilon\)-\(\delta\) 语言更容易理解和掌握。本节从数列极限的直观概念出发，严格建立了 \(\varepsilon\)-\(N\) 定义，并证明了极限的基本性质：唯一性（一个数列不能收敛于两个不同的极限）、有界性（收敛数列必有界）、四则运算法则（极限的加减乘除等于极限的加减乘除）和夹逼定理（被两个收敛于同一极限的数列夹逼的数列也收敛于此极限）。特别地，单调有界定理建立了单调性与收敛性之间的深刻联系，而 \(\lim_{n \to \infty} (1 + \frac{1}{n})^n = e\) 的重要极限则是这一联系的具体体现------它展示了单调递增且有上界的数列必收敛于一个超越数 \(e\)。

\subsubsection{22.1.1 数列极限的直观概念}\label{ux6570ux5217ux6781ux9650ux7684ux76f4ux89c2ux6982ux5ff5}

当我们说数列 \(\{a_n\}\) \textbf{趋于} \(L\)，直觉上是指：当 \(n\) 足够大时，\(a_n\) 可以任意接近 \(L\)。下面的 \(\varepsilon\)-\(N\) 定义将这一直觉精确化。

\subsubsection{\texorpdfstring{22.1.2 \(\varepsilon\)-\(N\) 定义}{22.1.2 \textbackslash varepsilon-N 定义}}\label{varepsilon-n-ux5b9aux4e49}

\textbf{定义 22.1}（数列极限）设 \(\{a_n\}\) 是一个数列，\(L \in \mathbb{R}\)。若

\[\forall \varepsilon > 0, \exists N \in \mathbb{N}^+, \forall n > N: |a_n - L| < \varepsilon\]

则称数列 \(\{a_n\}\) \textbf{收敛于} \(L\)，记作 \(\lim_{n \to \infty} a_n = L\) 或 \(a_n \to L\)（\(n \to \infty\)）。

若数列不收敛于任何实数，则称该数列\textbf{发散}。

\textbf{几何解释}：\(\lim_{n \to \infty} a_n = L\) 意味着，对于 \(L\) 的任意 \(\varepsilon\) 邻域 \((L - \varepsilon, L + \varepsilon)\)，从某项 \(N\) 之后的所有项都落在这个邻域内。

\subsubsection{22.1.3 极限的唯一性}\label{ux6781ux9650ux7684ux552fux4e00ux6027}

\textbf{定理 22.1}（极限的唯一性）若 \(\lim_{n \to \infty} a_n = L_1\) 且 \(\lim_{n \to \infty} a_n = L_2\)，则 \(L_1 = L_2\)。

\textbf{证明}：用反证法。假设 \(L_1 \neq L_2\)，令 \(\varepsilon = \frac{|L_1 - L_2|}{2} > 0\)。

由 \(\lim_{n \to \infty} a_n = L_1\)，存在 \(N_1\)，使得当 \(n > N_1\) 时，\(|a_n - L_1| < \varepsilon\)。

由 \(\lim_{n \to \infty} a_n = L_2\)，存在 \(N_2\)，使得当 \(n > N_2\) 时，\(|a_n - L_2| < \varepsilon\)。

取 \(n > \max\{N_1, N_2\}\)，则

\[|L_1 - L_2| = |L_1 - a_n + a_n - L_2| \leq |a_n - L_1| + |a_n - L_2| < \varepsilon + \varepsilon = 2\varepsilon = |L_1 - L_2|\]

矛盾！因此 \(L_1 = L_2\)。\(\blacksquare\)

\textbf{定理 22.20a}（无穷小量阶的比较的严格框架）：设 \(\alpha(x)\) 和 \(\beta(x)\) 都是 \(x \to a\) 时的无穷小量。定义： - 若 \(\lim_{x \to a} \frac{\alpha(x)}{\beta(x)} = 0\)，则称 \(\alpha(x)\) 是比 \(\beta(x)\) \textbf{更高阶}的无穷小，记作 \(\alpha(x) = o(\beta(x))\)（\(x \to a\)） - 若 \(\lim_{x \to a} \frac{\alpha(x)}{\beta(x)} = c \in \mathbb{R} \setminus \{0\}\)，则称 \(\alpha(x)\) 和 \(\beta(x)\) 是\textbf{同阶}无穷小 - 若 \(\lim_{x \to a} \frac{\alpha(x)}{\beta(x)} = 1\)，则称 \(\alpha(x)\) 和 \(\beta(x)\) 是\textbf{等价}无穷小（定义 22.9）

高阶无穷小的记号 \(o(\beta(x))\)（称为''小 \(o\) 记号''或 Landau 记号）在分析学中极为重要。例如，\(\sin x = x + o(x)\)（\(x \to 0\)）精确地表达了 \(\sin x\) 在 \(x = 0\) 附近被 \(x\) 线性逼近的误差阶。更一般地，如果 \(f(x) = L + a_1(x - a) + o(x - a)\)（\(x \to a\)），则 \(f\) 在 \(a\) 处可导且导数为 \(a_1\)------这是导数定义的小 \(o\) 记法形式。

\textbf{定理 22.20b}（无穷大量的比较）：设 \(f(x)\) 和 \(g(x)\) 都是 \(x \to a\)（或 \(x \to \infty\)）时的无穷大量。 - 若 \(\lim_{x \to a} \frac{f(x)}{g(x)} = \infty\)，则称 \(f(x)\) 是比 \(g(x)\) \textbf{更高阶}的无穷大 - 若 \(\lim_{x \to a} \frac{f(x)}{g(x)} = c \in (0, \infty)\)，则称 \(f(x)\) 和 \(g(x)\) 是\textbf{同阶}无穷大 - 若 \(\lim_{x \to a} \frac{f(x)}{g(x)} = 1\)，则称 \(f(x)\) 和 \(g(x)\) 是\textbf{等价}无穷大，记作 \(f(x) \sim g(x)\)

例如，当 \(x \to \infty\) 时，\(e^x\) 是比任何多项式 \(x^n\) 更高阶的无穷大，而 \(x^n\) 是比 \(\ln x\) 更高阶的无穷大。这一''阶的层级''（\(e^x \gg x^n \gg \ln x\)）在分析学中反复出现，是理解函数渐近行为的关键。

\subsubsection{22.1.4 收敛数列的有界性}\label{ux6536ux655bux6570ux5217ux7684ux6709ux754cux6027}

\textbf{定理 22.2}（收敛数列的有界性）若 \(\lim_{n \to \infty} a_n = L\)，则 \(\{a_n\}\) 是有界数列。

\textbf{证明}：取 \(\varepsilon = 1\)，由极限定义，存在 \(N\)，使得当 \(n > N\) 时，\(|a_n - L| < 1\)，即 \(|a_n| < |L| + 1\)。

令 \(M = \max\{|a_1|, |a_2|, \ldots, |a_{N}|, |L| + 1\}\)，则 \(\forall n \in \mathbb{N}^+, |a_n| \leq M\)。\(\blacksquare\)

\textbf{注}：逆命题不成立，即有界数列不一定收敛。例如 \(a_n = (-1)^n\) 有界但发散。

\subsubsection{22.1.5 极限的四则运算}\label{ux6781ux9650ux7684ux56dbux5219ux8fd0ux7b97}

\textbf{定理 22.3}（极限的四则运算）设 \(\lim_{n \to \infty} a_n = A\)，\(\lim_{n \to \infty} b_n = B\)。则：

\begin{enumerate}
\def\labelenumi{\arabic{enumi}.}
\tightlist
\item
  \(\lim_{n \to \infty} (a_n + b_n) = A + B\)
\item
  \(\lim_{n \to \infty} (a_n - b_n) = A - B\)
\item
  \(\lim_{n \to \infty} (a_n \cdot b_n) = A \cdot B\)
\item
  若 \(B \neq 0\)，\(\lim_{n \to \infty} \frac{a_n}{b_n} = \frac{A}{B}\)
\end{enumerate}

\textbf{证明}：

\begin{enumerate}
\def\labelenumi{(\arabic{enumi})}
\tightlist
\item
  对于任意 \(\varepsilon > 0\)，由 \(\lim a_n = A\)，存在 \(N_1\)，使得当 \(n > N_1\) 时，\(|a_n - A| < \frac{\varepsilon}{2}\)。
\end{enumerate}

由 \(\lim b_n = B\)，存在 \(N_2\)，使得当 \(n > N_2\) 时，\(|b_n - B| < \frac{\varepsilon}{2}\)。

取 \(N = \max\{N_1, N_2\}\)，当 \(n > N\) 时：

\[|(a_n + b_n) - (A + B)| \leq |a_n - A| + |b_n - B| < \frac{\varepsilon}{2} + \frac{\varepsilon}{2} = \varepsilon\]

\begin{enumerate}
\def\labelenumi{(\arabic{enumi})}
\setcounter{enumi}{2}
\tightlist
\item
  首先由定理 22.2，存在 \(M_1 > 0\) 使得 \(|b_n| \leq M_1\)。令 \(M = \max\{M_1, |A|, 1\}\)。
\end{enumerate}

对于任意 \(\varepsilon > 0\)，存在 \(N_1\) 使得当 \(n > N_1\) 时 \(|a_n - A| < \frac{\varepsilon}{2M}\)；存在 \(N_2\) 使得当 \(n > N_2\) 时 \(|b_n - B| < \frac{\varepsilon}{2M}\)。

当 \(n > \max\{N_1, N_2\}\) 时：

\[|a_n b_n - AB| = |a_n(b_n - B) + B(a_n - A)| \leq |a_n||b_n - B| + |B||a_n - A| < M \cdot \frac{\varepsilon}{2M} + M \cdot \frac{\varepsilon}{2M} = \varepsilon\]

\begin{enumerate}
\def\labelenumi{(\arabic{enumi})}
\setcounter{enumi}{3}
\tightlist
\item
  首先证明 \(\lim_{n \to \infty} \frac{1}{b_n} = \frac{1}{B}\)。由于 \(B \neq 0\)，取 \(\varepsilon = \frac{|B|}{2}\)，存在 \(N_0\) 使得当 \(n > N_0\) 时 \(|b_n - B| < \frac{|B|}{2}\)，即 \(|b_n| > \frac{|B|}{2}\)。
\end{enumerate}

对于任意 \(\varepsilon > 0\)，存在 \(N\) 使得当 \(n > N\) 时 \(|b_n - B| < \frac{|B|^2 \varepsilon}{2}\)。

当 \(n > \max\{N_0, N\}\) 时：

\[\left|\frac{1}{b_n} - \frac{1}{B}\right| = \frac{|B - b_n|}{|b_n||B|} < \frac{|B - b_n|}{\frac{|B|}{2} \cdot |B|} = \frac{2|B - b_n|}{|B|^2} < \varepsilon\]

再由 (3)，\(\frac{a_n}{b_n} = a_n \cdot \frac{1}{b_n} \to A \cdot \frac{1}{B} = \frac{A}{B}\)。\(\blacksquare\)

\subsubsection{22.1.6 夹逼定理}\label{ux5939ux903cux5b9aux7406}

\textbf{定理 22.4}（夹逼定理）设 \(\{a_n\}\)，\(\{b_n\}\)，\(\{c_n\}\) 是三个数列。若存在 \(N_0 \in \mathbb{N}^+\)，使得

\[\forall n > N_0: a_n \leq b_n \leq c_n\]

且 \(\lim_{n \to \infty} a_n = \lim_{n \to \infty} c_n = L\)，则 \(\lim_{n \to \infty} b_n = L\)。

\textbf{证明}：对于任意 \(\varepsilon > 0\)：

由 \(\lim a_n = L\)，存在 \(N_1\)，使得当 \(n > N_1\) 时，\(L - \varepsilon < a_n\)。

由 \(\lim c_n = L\)，存在 \(N_2\)，使得当 \(n > N_2\) 时，\(c_n < L + \varepsilon\)。

取 \(N = \max\{N_0, N_1, N_2\}\)，当 \(n > N\) 时：

\[L - \varepsilon < a_n \leq b_n \leq c_n < L + \varepsilon\]

所以 \(|b_n - L| < \varepsilon\)，即 \(\lim_{n \to \infty} b_n = L\)。\(\blacksquare\)

\subsubsection{22.1.7 单调有界定理}\label{ux5355ux8c03ux6709ux754cux5b9aux7406}

\textbf{定理 22.5}（单调有界定理）

\begin{enumerate}
\def\labelenumi{\arabic{enumi}.}
\tightlist
\item
  单调递增且有上界的数列必收敛。
\item
  单调递减且有下界的数列必收敛。
\end{enumerate}

\textbf{证明}：

\begin{enumerate}
\def\labelenumi{(\arabic{enumi})}
\tightlist
\item
  设 \(\{a_n\}\) 单调递增且有上界。由实数的完备性（确界存在定理），\(\{a_n\}\) 有上确界，设为 \(L = \sup\{a_n\}\)。
\end{enumerate}

我们证明 \(\lim_{n \to \infty} a_n = L\)。

对于任意 \(\varepsilon > 0\)，由上确界的定义，\(L - \varepsilon\) 不是 \(\{a_n\}\) 的上界，所以存在 \(N\) 使得 \(a_{N} > L - \varepsilon\)。

由于 \(\{a_n\}\) 单调递增，当 \(n > N\) 时，\(a_n \geq a_{N} > L - \varepsilon\)。

又因为 \(L\) 是上界，\(a_n \leq L < L + \varepsilon\)。

所以当 \(n > N\) 时，\(L - \varepsilon < a_n < L + \varepsilon\)，即 \(|a_n - L| < \varepsilon\)。

\begin{enumerate}
\def\labelenumi{(\arabic{enumi})}
\setcounter{enumi}{1}
\tightlist
\item
  类似证明。\(\blacksquare\)
\end{enumerate}

\subsubsection{22.1.8 重要极限}\label{ux91cdux8981ux6781ux9650}

\textbf{定理 22.6} \(\lim_{n \to \infty} \frac{1}{n} = 0\)

\textbf{证明}：对于任意 \(\varepsilon > 0\)，取 \(N = \lceil \frac{1}{\varepsilon} \rceil\)。当 \(n > N\) 时，\(n > \frac{1}{\varepsilon}\)，即 \(\frac{1}{n} < \varepsilon\)。\(\blacksquare\)

\textbf{定理 22.7} 若 \(|q| < 1\)，则 \(\lim_{n \to \infty} q^n = 0\)

\textbf{证明}：当 \(q = 0\) 时显然成立。当 \(0 < |q| < 1\) 时，令 \(|q| = \frac{1}{1 + h}\)（\(h > 0\)）。由伯努利不等式，\((1 + h)^n \geq 1 + nh > nh\)。

所以 \(|q^n| = \frac{1}{(1 + h)^n} < \frac{1}{nh}\)。对于任意 \(\varepsilon > 0\)，取 \(N = \lceil \frac{1}{h\varepsilon} \rceil\)，当 \(n > N\) 时 \(|q^n| < \varepsilon\)。\(\blacksquare\)

\textbf{定理 22.8} \(\lim_{n \to \infty}\left(1 + \frac{1}{n}\right)^n = e\)

其中 \(e \approx 2.71828...\) 是自然对数的底。

\textbf{证明}：分两步证明。

\textbf{第一步}：证明 \(\{a_n\} = \left\{\left(1 + \frac{1}{n}\right)^n\right\}\) 单调递增。

由二项式定理，

\[a_n = \sum_{k=0}^{n} \frac{1}{k!} \cdot 1 \cdot \left(1 - \frac{1}{n}\right) \cdots \left(1 - \frac{k-1}{n}\right)\]

类似地，

\[a_{n+1} = \sum_{k=0}^{n+1} \frac{1}{k!} \cdot 1 \cdot \left(1 - \frac{1}{n+1}\right) \cdots \left(1 - \frac{k-1}{n+1}\right)\]

因为 \(1 - \frac{j}{n} < 1 - \frac{j}{n+1}\)（\(j = 1, 2, \ldots, k-1\)），所以 \(a_{n+1}\) 的每一项都不小于 \(a_n\) 的对应项，且 \(a_{n+1}\) 多了一项。因此 \(a_{n+1} > a_n\)。

\textbf{第二步}：证明 \(\{a_n\}\) 有上界。

\[a_n < \sum_{k=0}^{n} \frac{1}{k!} < 1 + 1 + \frac{1}{2} + \frac{1}{2^2} + \cdots = 3\]

（利用了 \(\frac{1}{k!} < \frac{1}{2^{k-1}}\) 对 \(k \geq 2\) 成立。）

由单调有界定理，\(\{a_n\}\) 收敛。定义此极限为 \(e\)。\(\blacksquare\)

\hrulefill

\section{第65章：函数极限}\label{ux7b2c65ux7ae0ux51fdux6570ux6781ux9650}

函数极限是极限理论的第二站，将数列极限的离散概念推广到连续变量。与数列极限的 \(\varepsilon\)-\(N\) 语言不同，函数极限使用 \(\varepsilon\)-\(\delta\) 语言------\(\delta\) 取代了 \(N\) 的角色，刻画了自变量趋近于某一点时函数值的趋近行为。本节从 \(x \to a\) 时的极限定义出发，引入了左极限和右极限的概念，证明了函数极限与单侧极限的等价性（双侧极限存在当且仅当左右极限存在且相等）。函数极限的四则运算和夹逼定理与数列极限平行，但 \(\varepsilon\)-\(\delta\) 证明的技术细节更为精细。本节的核心结果是重要极限 \(\lim_{x \to 0} \frac{\sin x}{x} = 1\)，其证明利用单位圆中的几何不等式和夹逼定理，是几何直观与严格分析相结合的典范。这一极限不仅是计算三角函数导数的基础，也是整个三角函数分析理论的逻辑起点。

\subsubsection{\texorpdfstring{22.2.1 \(x \to a\) 时的极限定义}{22.2.1 x \textbackslash to a 时的极限定义}}\label{x-to-a-ux65f6ux7684ux6781ux9650ux5b9aux4e49}

\textbf{定义 22.2}（函数极限）设 \(f\) 在 \(a\) 的某个去心邻域 \(\mathring{U}(a, \delta_0)\) 上有定义，\(L \in \mathbb{R}\)。若

\[\forall \varepsilon > 0, \exists \delta > 0, \forall x: 0 < |x - a| < \delta \Rightarrow |f(x) - L| < \varepsilon\]

则称 \(f(x)\) 当 \(x \to a\) 时\textbf{收敛于} \(L\)，记作 \(\lim_{x \to a} f(x) = L\)。

\textbf{注意}：定义中要求 \(0 < |x - a|\)，即 \(x \neq a\)，说明 \(f\) 在 \(a\) 点的极限与 \(f(a)\) 的值无关。

\subsubsection{22.2.2 左极限和右极限}\label{ux5de6ux6781ux9650ux548cux53f3ux6781ux9650}

\textbf{定义 22.3}（左极限和右极限）

\[\lim_{x \to a^-} f(x) = L \iff \forall \varepsilon > 0, \exists \delta > 0, \forall x \in (a - \delta, a): |f(x) - L| < \varepsilon\]

\[\lim_{x \to a^+} f(x) = L \iff \forall \varepsilon > 0, \exists \delta > 0, \forall x \in (a, a + \delta): |f(x) - L| < \varepsilon\]

\textbf{定理 22.9} \(\lim_{x \to a} f(x) = L\) 当且仅当 \(\lim_{x \to a^-} f(x) = \lim_{x \to a^+} f(x) = L\)。

\textbf{证明}：必要性显然。充分性：对于任意 \(\varepsilon > 0\)，存在 \(\delta_1, \delta_2 > 0\) 使得在 \((a - \delta_1, a)\) 和 \((a, a + \delta_2)\) 上分别成立 \(|f(x) - L| < \varepsilon\)。取 \(\delta = \min\{\delta_1, \delta_2\}\)，当 \(0 < |x - a| < \delta\) 时 \(|f(x) - L| < \varepsilon\)。\(\blacksquare\)

\subsubsection{22.2.3 函数极限的四则运算}\label{ux51fdux6570ux6781ux9650ux7684ux56dbux5219ux8fd0ux7b97}

\textbf{定理 22.10} 设 \(\lim_{x \to a} f(x) = A\)，\(\lim_{x \to a} g(x) = B\)。则四则运算成立，证明与定理 22.3 类似。

\textbf{证明}：由函数极限的序列刻画（Heine定理），将函数极限转化为序列极限，应用定理 22.3（序列极限四则运算）即得。\(\blacksquare\)

\(\varepsilon\)-\(\delta\) 证明（加法）：设 \(\lim_{x \to a} f(x) = L\)，\(\lim_{x \to a} g(x) = M\)。对任意 \(\varepsilon > 0\)，取 \(\varepsilon_1 = \varepsilon/2\)。由 \(f\) 的极限，存在 \(\delta_1 > 0\) 使得 \(0 < |x - a| < \delta_1\) 时 \(|f(x) - L| < \varepsilon/2\)。由 \(g\) 的极限，存在 \(\delta_2 > 0\) 使得 \(0 < |x - a| < \delta_2\) 时 \(|g(x) - M| < \varepsilon/2\)。取 \(\delta = \min(\delta_1, \delta_2)\)，则 \(0 < |x - a| < \delta\) 时 \(|(f(x) + g(x)) - (L + M)| \leq |f(x) - L| + |g(x) - M| < \varepsilon/2 + \varepsilon/2 = \varepsilon\)。

\subsubsection{22.2.4 函数极限的夹逼定理}\label{ux51fdux6570ux6781ux9650ux7684ux5939ux903cux5b9aux7406}

\textbf{定理 22.11}（夹逼定理）设 \(f(x) \leq g(x) \leq h(x)\) 在 \(a\) 的某个去心邻域内成立，且 \(\lim_{x \to a} f(x) = \lim_{x \to a} h(x) = L\)，则 \(\lim_{x \to a} g(x) = L\)。

\textbf{证明}：与定理 22.4 类似。\(\varepsilon\)-\(\delta\) 证明：设 \(g(x) \leq f(x) \leq h(x)\) 在 \(a\) 的去心邻域内成立，且 \(\lim_{x \to a} g(x) = \lim_{x \to a} h(x) = L\)。对任意 \(\varepsilon > 0\)，存在 \(\delta_1 > 0\) 使得 \(0 < |x - a| < \delta_1\) 时 \(|g(x) - L| < \varepsilon\)，即 \(L - \varepsilon < g(x)\)。存在 \(\delta_2 > 0\) 使得 \(0 < |x - a| < \delta_2\) 时 \(|h(x) - L| < \varepsilon\)，即 \(h(x) < L + \varepsilon\)。取 \(\delta = \min(\delta_1, \delta_2)\)，则 \(0 < |x - a| < \delta\) 时 \(L - \varepsilon < g(x) \leq f(x) \leq h(x) < L + \varepsilon\)，故 \(|f(x) - L| < \varepsilon\)。\(\blacksquare\)

\subsubsection{\texorpdfstring{22.2.5 重要极限 \(\lim_{x \to 0} \frac{\sin x}{x} = 1\)}{22.2.5 重要极限 \textbackslash lim\_\{x \textbackslash to 0\} \textbackslash frac\{\textbackslash sin x\}\{x\} = 1}}\label{ux91cdux8981ux6781ux9650-lim_x-to-0-fracsin-xx-1}

\textbf{定理 22.12} \(\lim_{x \to 0} \frac{\sin x}{x} = 1\)

\textbf{证明}：当 \(0 < x < \frac{\pi}{2}\) 时，考虑单位圆中的扇形 \(OAB\) 和三角形 \(OAB\) 及三角形 \(OAC\)，其中 \(A = (1, 0)\)，\(B = (\cos x, \sin x)\)，\(C = (1, \tan x)\)。

扇形面积 \(= \frac{x}{2}\)，三角形 \(OAB\) 面积 \(= \frac{1}{2}\sin x\)，三角形 \(OAC\) 面积 \(= \frac{1}{2}\tan x\)。

由面积关系：

\[\frac{1}{2}\sin x < \frac{x}{2} < \frac{1}{2}\tan x\]

即 \(\sin x < x < \frac{\sin x}{\cos x}\)。

由 \(\sin x < x\) 得 \(\frac{\sin x}{x} < 1\)。由 \(x < \frac{\sin x}{\cos x}\) 得 \(\cos x < \frac{\sin x}{x}\)。

因此当 \(0 < x < \frac{\pi}{2}\) 时：\(\cos x < \frac{\sin x}{x} < 1\)。

由偶函数性质，当 \(-\frac{\pi}{2} < x < 0\) 时同样成立。

现在证明 \(\lim_{x \to 0} \cos x = 1\)。利用半角公式 \(1 - \cos x = 2\sin^2\frac{x}{2}\)，以及前面已证的不等式 \(|\sin t| < |t|\)（当 \(0 < |t| < \frac{\pi}{2}\) 时，由几何推导得到；当 \(|t| \geq \frac{\pi}{2}\) 时 \(|\sin t| \leq 1 < \frac{\pi}{2} \leq |t|\)，故 \(|\sin t| \leq |t|\) 对所有 \(t\) 成立），取 \(t = \frac{x}{2}\) 得 \(\left|\sin\frac{x}{2}\right| \leq \left|\frac{x}{2}\right|\)。因此

\[|\cos x - 1| = 1 - \cos x = 2\sin^2\frac{x}{2} \leq 2 \cdot \left(\frac{x}{2}\right)^2 = \frac{x^2}{2}\]

对于任意 \(\varepsilon > 0\)，取 \(\delta = \sqrt{2\varepsilon}\)，当 \(0 < |x| < \delta\) 时 \(|\cos x - 1| < \varepsilon\)，故 \(\lim_{x \to 0} \cos x = 1\)。

由夹逼定理，\(\lim_{x \to 0} \frac{\sin x}{x} = 1\)。\(\blacksquare\)

\hrulefill

\section{第66章：连续函数}\label{ux7b2c66ux7ae0ux8fdeux7eedux51fdux6570}

连续函数是分析学中最重要的函数类------它们构成了''直观上不会断裂''的函数的严格数学定义。本节从极限定义出发，严格定义了函数在一点连续（\(\lim_{x \to a} f(x) = f(a)\)）和在区间上连续的概念，并证明了连续函数的核心性质：连续函数的四则运算和复合仍是连续函数。介值定理（闭区间上的连续函数取遍端点值之间的所有值）和最值定理（闭区间上的连续函数必有最大值和最小值）是连续函数理论中最深刻的两个结果，它们的证明依赖于实数的完备性（确界存在定理），是分析学区别于初等数学的标志性成就。值得强调的是，介值定理和最值定理都要求函数定义在闭区间上------若在开区间上，这两个结论均不成立，这体现了闭区间''包含端点''这一拓扑性质的关键作用。

\subsubsection{22.3.1 连续的定义}\label{ux8fdeux7eedux7684ux5b9aux4e49}

\textbf{定义 22.4}（连续）设 \(f\) 在 \(a\) 的某个邻域内有定义。若

\[\lim_{x \to a} f(x) = f(a)\]

则称 \(f\) 在点 \(a\) 处\textbf{连续}。

等价地，\(f\) 在 \(a\) 处连续当且仅当：

\[\forall \varepsilon > 0, \exists \delta > 0, \forall x: |x - a| < \delta \Rightarrow |f(x) - f(a)| < \varepsilon\]

\textbf{定义 22.5}（区间上的连续）若 \(f\) 在区间 \(I\) 上的每一点都连续，则称 \(f\) 在 \(I\) 上\textbf{连续}。

\subsubsection{22.3.2 连续函数的运算}\label{ux8fdeux7eedux51fdux6570ux7684ux8fd0ux7b97}

\textbf{定理 22.13}（连续函数的四则运算）设 \(f\) 和 \(g\) 都在点 \(a\) 处连续，则 \(f + g\)，\(f - g\)，\(f \cdot g\) 也在 \(a\) 处连续。若 \(g(a) \neq 0\)，则 \(\frac{f}{g}\) 也在 \(a\) 处连续。

\textbf{证明}：由极限的四则运算和连续的定义直接可得。\(\blacksquare\)

\textbf{定理 22.14}（复合函数的连续性）设 \(g\) 在 \(a\) 处连续，\(f\) 在 \(g(a)\) 处连续，则 \(f \circ g\) 在 \(a\) 处连续。

\textbf{证明}：由 \(g\) 在 \(a\) 处连续，\(\lim_{x \to a} g(x) = g(a)\)。又由 \(f\) 在 \(g(a)\) 处连续，根据连续函数的定义（或函数极限的复合定理在连续情形下的形式），\(\lim_{t \to g(a)} f(t) = f(g(a))\)。因此

\[\lim_{x \to a} f(g(x)) = f\left(\lim_{x \to a} g(x)\right) = f(g(a))\]

其中第一步利用了 \(f\) 在 \(g(a)\) 处的连续性（而非仅极限存在），使得极限值可以直接代入 \(f\)。\(\blacksquare\)

\subsubsection{22.3.3 介值定理}\label{ux4ecbux503cux5b9aux7406}

\textbf{定理 22.15}（介值定理）设 \(f\) 在 \([a, b]\) 上连续，且 \(f(a) \neq f(b)\)。则对于 \(f(a)\) 和 \(f(b)\) 之间的任意实数 \(c\)，存在 \(\xi \in (a, b)\)，使得 \(f(\xi) = c\)。

\textbf{证明}：不妨设 \(f(a) < c < f(b)\)。定义集合

\[A = \{x \in [a, b] \mid f(x) < c\}\]

则 \(a \in A\)，故 \(A\) 非空。又 \(A\) 有上界 \(b\)，由确界存在定理，\(\xi = \sup A\) 存在，且 \(\xi \in [a, b]\)。

\textbf{证明 \(f(\xi) \leq c\)}：对于任意 \(n \in \mathbb{N}^+\)，由上确界的定义，存在 \(x_n \in A\) 使得 \(\xi - \frac{1}{n} < x_n \leq \xi\)。由夹逼定理，\(x_n \to \xi\)。因为 \(f(x_n) < c\) 且 \(f\) 连续，所以 \(f(\xi) = \lim_{n \to \infty} f(x_n) \leq c\)。

\textbf{证明 \(f(\xi) \geq c\)}：先证 \(\xi < b\)。若 \(\xi = b\)，则由上确界定义，存在 \(x_n \in A\) 使得 \(x_n \to b\)。由 \(f\) 的连续性，\(f(b) = \lim_{n \to \infty} f(x_n) \leq c\)（因为 \(f(x_n) < c\)）。但这与假设 \(f(b) > c\) 矛盾。故 \(\xi < b\)。

由 \(\xi < b\)，对于充分大的 \(n\)，\(\xi + \frac{1}{n} < b\) 且 \(\xi + \frac{1}{n} > \xi = \sup A\)，所以 \(\xi + \frac{1}{n} \notin A\)，即 \(f(\xi + \frac{1}{n}) \geq c\)。由 \(f\) 的连续性，\(f(\xi) = \lim f(\xi + \frac{1}{n}) \geq c\)。

因此 \(f(\xi) = c\)。\(\blacksquare\)

\textbf{推论 22.1}（零点定理）设 \(f\) 在 \([a, b]\) 上连续，且 \(f(a) \cdot f(b) < 0\)，则存在 \(\xi \in (a, b)\)，使得 \(f(\xi) = 0\)。

\textbf{证明}：由介值定理取 \(c = 0\)。\(\blacksquare\)

\subsubsection{22.3.4 最值定理}\label{ux6700ux503cux5b9aux7406}

\textbf{定理 22.16}（最值定理）设 \(f\) 在闭区间 \([a, b]\) 上连续，则 \(f\) 在 \([a, b]\) 上取到最大值和最小值，即存在 \(x_1, x_2 \in [a, b]\)，使得

\[f(x_1) = \sup_{x \in [a, b]} f(x), \quad f(x_2) = \inf_{x \in [a, b]} f(x)\]

\textbf{证明}：我们证明最大值的存在性（最小值类似）。

\textbf{第一步}：证明 \(f\) 在 \([a, b]\) 上有界。用反证法：若 \(f\) 无上界，则对每个 \(n\) 存在 \(x_n \in [a, b]\) 使得 \(f(x_n) > n\)。由波尔查诺-魏尔斯特拉斯定理，\(\{x_n\}\) 有收敛子列 \(\{x_{n_k}\}\)，设 \(x_{n_k} \to \xi \in [a, b]\)。由 \(f\) 的连续性，\(f(x_{n_k}) \to f(\xi)\)。但 \(f(x_{n_k}) > n_k \to +\infty\)，矛盾。

\textbf{第二步}：设 \(M = \sup_{x \in [a, b]} f(x)\)（由第一步 \(M\) 有限）。由上确界的定义，对每个 \(n\) 存在 \(x_n\) 使得 \(f(x_n) > M - \frac{1}{n}\)。由波尔查诺-魏尔斯特拉斯定理，\(\{x_n\}\) 有收敛子列 \(x_{n_k} \to x_0 \in [a, b]\)。由 \(f\) 的连续性，\(f(x_0) = \lim f(x_{n_k}) \geq M\)。又 \(f(x_0) \leq M\)，故 \(f(x_0) = M\)。\(\blacksquare\)

\textbf{定理 22.16a}（Darboux 性质 / 介值性质的等价表述）：设 \(f\) 在区间 \(I\) 上连续。则 \(f\) 的值域 \(f(I)\) 也是一个区间（即 \(f\) 具有介值性质：对于 \(f(I)\) 中任意两点之间的任何值，都是 \(f\) 在某点的函数值）。这一定理统一了介值定理和最值定理：闭区间上连续函数的值域是闭区间 \([m, M]\)，其中 \(m = \min f\)，\(M = \max f\)。

\textbf{定理 22.16b}（连续函数在闭区间上的振幅）：设 \(f\) 在 \([a, b]\) 上连续。定义 \(f\) 在 \([a, b]\) 上的\textbf{振幅}（oscillation）为

\[\omega(f, [a, b]) = \sup_{x, y \in [a, b]} |f(x) - f(y)| = M - m\]

其中 \(M = \max f\)，\(m = \min f\)。由最值定理，\(M\) 和 \(m\) 都在 \([a, b]\) 上取到，因此 \(\omega(f, [a, b])\) 是有限实数。振幅是衡量函数在区间上''波动程度''的基本量------振幅越小，函数越''平坦''；振幅为零当且仅当 \(f\) 在 \([a, b]\) 上为常数函数。振幅在黎曼积分的理论中起核心作用：\(f\) 在 \([a, b]\) 上可积的充要条件是振幅可以任意小地控制（即对于任意 \(\varepsilon > 0\)，存在分割使得振幅大的子区间总长度任意小）。

\textbf{注记 22.3}（连续函数与不连续函数的对比）：连续函数满足介值性质（Darboux 性质），但逆命题不成立------存在处处不连续但满足介值性质的函数（如 Conway 基 13 函数）。导数函数虽然不一定连续，但由 Darboux 定理（Ch 30），任何导数函数都具有介值性质。这表明介值性质是比连续性更''弱''但依然重要的函数性质。

\hrulefill

\section{第67章：函数极限的进一步讨论}\label{ux7b2c67ux7ae0ux51fdux6570ux6781ux9650ux7684ux8fdbux4e00ux6b65ux8ba8ux8bba}

在 22.2 节建立了函数极限的基本概念（\(x \to a\) 时的极限）之后，本节将极限理论推广到更丰富的情形。\(x \to \infty\) 时的极限描述了函数在自变量趋于无穷大时的渐近行为------这是函数图像的水平渐近线的精确定义。无穷小量是极限理论的核心概念之一：\(\lim_{x \to a} f(x) = 0\) 表示 \(f(x)\) 在 \(x \to a\) 时是无穷小量，它是极限为 \(0\) 的函数的代名词。无穷小量的阶的比较（高阶无穷小、同阶无穷小、等价无穷小）是分析函数局部行为的基本工具，等价无穷小替换 \(\lim_{x \to a} \frac{f(x)}{g(x)} = 1 \Rightarrow f(x) \sim g(x)\) 在极限计算中极大地简化了运算------例如 \(\sin x \sim x\)（\(x \to 0\)）和 \(\ln(1+x) \sim x\)（\(x \to 0\)）是最常用的等价替换。本节还将讨论无穷大量的概念，它是无穷小量的倒数，用于描述函数值的无界增长。这些概念的建立为后续的导数计算（特别是极限定义下的导数）和积分的极限处理提供了必要的工具。

\subsubsection{\texorpdfstring{22.4.1 \(x \to \infty\) 时的极限}{22.4.1 x \textbackslash to \textbackslash infty 时的极限}}\label{x-to-infty-ux65f6ux7684ux6781ux9650}

\textbf{定义 22.6} 设 \(f\) 在 \([a, +\infty)\) 上有定义，\(L \in \mathbb{R}\)。若

\[\forall \varepsilon > 0, \exists X > 0, \forall x > X: |f(x) - L| < \varepsilon\]

则记 \(\lim_{x \to +\infty} f(x) = L\)。

类似地定义 \(\lim_{x \to -\infty} f(x) = L\) 和 \(\lim_{x \to \infty} f(x) = L\)。

\textbf{定理 22.17} \(\lim_{x \to +\infty} \frac{1}{x^{\alpha}} = 0\)（\(\alpha > 0\)）。

\textbf{证明}：对于任意 \(\varepsilon > 0\)，取 \(X = \left(\frac{1}{\varepsilon}\right)^{1/\alpha}\)。当 \(x > X\) 时，

\[\left|\frac{1}{x^{\alpha}}\right| = \frac{1}{x^{\alpha}} < \frac{1}{X^{\alpha}} = \varepsilon\]

\(\blacksquare\)

\subsubsection{22.4.2 无穷小与无穷大}\label{ux65e0ux7a77ux5c0fux4e0eux65e0ux7a77ux5927}

\textbf{定义 22.7}（无穷小）若 \(\lim_{x \to a} f(x) = 0\)，则称 \(f(x)\) 是 \(x \to a\) 时的\textbf{无穷小量}。

\textbf{定义 22.8}（无穷大）若 \(\lim_{x \to a} |f(x)| = +\infty\)，即

\[\forall M > 0, \exists \delta > 0, \forall x: 0 < |x - a| < \delta \Rightarrow |f(x)| > M\]

则称 \(f(x)\) 是 \(x \to a\) 时的\textbf{无穷大量}。

\textbf{定理 22.18}（无穷小与无穷大的关系）若 \(f(x)\) 是 \(x \to a\) 时的无穷小量，且 \(f(x) \neq 0\) 在 \(a\) 的某个去心邻域内成立，则 \(\frac{1}{f(x)}\) 是 \(x \to a\) 时的无穷大量。

\textbf{证明}：对于任意 \(M > 0\)，取 \(\varepsilon = \frac{1}{M}\)。存在 \(\delta > 0\)，当 \(0 < |x - a| < \delta\) 时 \(|f(x)| < \varepsilon = \frac{1}{M}\)，故 \(\left|\frac{1}{f(x)}\right| > M\)。\(\blacksquare\)

\subsubsection{22.4.3 等价无穷小}\label{ux7b49ux4ef7ux65e0ux7a77ux5c0f}

\textbf{定义 22.9} 设 \(\alpha(x)\) 和 \(\beta(x)\) 都是 \(x \to a\) 时的无穷小量。若 \(\lim_{x \to a} \frac{\alpha(x)}{\beta(x)} = 1\)，则称 \(\alpha(x)\) 和 \(\beta(x)\) 是\textbf{等价无穷小}，记作 \(\alpha(x) \sim \beta(x)\)（\(x \to a\)）。

\textbf{定理 22.19} 当 \(x \to 0\) 时的常用等价无穷小：

\begin{enumerate}
\def\labelenumi{\arabic{enumi}.}
\tightlist
\item
  \(\sin x \sim x\)
\item
  \(\tan x \sim x\)
\item
  \(1 - \cos x \sim \frac{x^2}{2}\)
\item
  \(e^x - 1 \sim x\)
\item
  \(\ln(1 + x) \sim x\)
\item
  \((1 + x)^{\alpha} - 1 \sim \alpha x\)
\end{enumerate}

\textbf{证明}：

\begin{enumerate}
\def\labelenumi{(\arabic{enumi})}
\item
  \(\lim_{x \to 0} \frac{\sin x}{x} = 1\)（定理 22.12）。
\item
  \(\lim_{x \to 0} \frac{1 - \cos x}{\frac{x^2}{2}} = \lim_{x \to 0} \frac{2\sin^2\frac{x}{2}}{\frac{x^2}{2}} = \lim_{x \to 0} \left(\frac{\sin\frac{x}{2}}{\frac{x}{2}}\right)^2 = 1\)。
\item
  \(\lim_{x \to 0} \frac{e^x - 1}{x} = \lim_{t \to 0} \frac{t}{\ln(1+t)} = 1\)（令 \(t = e^x - 1\)）。
\item
  \(\lim_{x \to 0} \frac{\ln(1+x)}{x} = \lim_{x \to 0} \ln(1+x)^{1/x} = \ln e = 1\)。
\item
  令 \(u = \alpha\ln(1+x)\)。当 \(x \to 0\) 时，由 (5)，\(\ln(1+x) \sim x\)，故 \(u \sim \alpha x \to 0\)。由 (4)（\(e^u - 1 \sim u\)，其中 \(u \to 0\)），\(e^{\alpha\ln(1+x)} - 1 \sim \alpha\ln(1+x)\)。再由 (5)，\(\alpha\ln(1+x) \sim \alpha x\)。由等价无穷小的传递性（若 \(\alpha_1 \sim \alpha_2\) 且 \(\alpha_2 \sim \alpha_3\)，则 \(\alpha_1 \sim \alpha_3\)，这由 \(\frac{\alpha_1}{\alpha_3} = \frac{\alpha_1}{\alpha_2} \cdot \frac{\alpha_2}{\alpha_3} \to 1 \cdot 1 = 1\)），\((1+x)^{\alpha} - 1 \sim \alpha x\)。\(\blacksquare\)
\end{enumerate}

\textbf{定理 22.20}（等价无穷小替换定理）设 \(\alpha \sim \alpha'\)，\(\beta \sim \beta'\)（\(x \to a\)），且 \(\lim_{x \to a} \frac{\alpha'}{\beta'}\) 存在，则

\[\lim_{x \to a} \frac{\alpha}{\beta} = \lim_{x \to a} \frac{\alpha'}{\beta'}\]

\textbf{证明}：

\[\lim_{x \to a} \frac{\alpha}{\beta} = \lim_{x \to a} \frac{\alpha}{\alpha'} \cdot \frac{\alpha'}{\beta'} \cdot \frac{\beta'}{\beta} = 1 \cdot \lim_{x \to a} \frac{\alpha'}{\beta'} \cdot 1 = \lim_{x \to a} \frac{\alpha'}{\beta'}\]

\(\blacksquare\)

\hrulefill

\section{第68章：一致连续性}\label{ux7b2c68ux7ae0ux4e00ux81f4ux8fdeux7eedux6027}

一致连续性是比逐点连续更强的连续概念。逐点连续 \(\forall \varepsilon > 0, \forall x_0, \exists \delta(x_0, \varepsilon) > 0\) 中，\(\delta\) 可以依赖于点 \(x_0\)；而一致连续 \(\forall \varepsilon > 0, \exists \delta(\varepsilon) > 0, \forall x_1, x_2\) 中，同一个 \(\delta\) 对区间上所有点都适用。一致连续性的重要性在于：它保证了连续函数在闭区间上的诸多良好性质（如有界性、一致可积性）。康托尔定理（Cantor's Theorem）指出：闭区间上的连续函数必一致连续。这一结论在开区间上不成立------例如 \(f(x) = 1/x\) 在 \((0, 1)\) 上连续但非一致连续，因为当 \(x\) 趋近于 \(0\) 时函数变化越来越剧烈，无法用一个统一的 \(\delta\) 控制。一致连续性在黎曼积分的理论中扮演关键角色------闭区间上连续函数的一致连续性保证了黎曼和的收敛性（即达布大和与小和的差可以任意小）。一致连续性还与''函数在区间上是否一致收敛''等更高级的分析概念密切相关，是分析学中连接局部性质与全局性质的桥梁。

\textbf{定义 22.10}（一致连续）设 \(f: I \to \mathbb{R}\)，其中 \(I\) 是区间。若

\[\forall \varepsilon > 0, \exists \delta > 0, \forall x_1, x_2 \in I: |x_1 - x_2| < \delta \Rightarrow |f(x_1) - f(x_2)| < \varepsilon\]

则称 \(f\) 在 \(I\) 上\textbf{一致连续}。

\textbf{注}：一致连续与逐点连续的区别在于：逐点连续中 \(\delta\) 可以依赖于点 \(x_0\)，而一致连续中 \(\delta\) 只依赖于 \(\varepsilon\)，对区间上所有点统一适用。

\textbf{定理 22.21}（康托尔定理）闭区间 \([a, b]\) 上的连续函数一致连续。

\textbf{证明}：用反证法。假设 \(f\) 在 \([a, b]\) 上连续但不一致连续。则存在 \(\varepsilon_0 > 0\)，对每个 \(n\)，存在 \(x_n, y_n \in [a, b]\)，\(|x_n - y_n| < \frac{1}{n}\) 但 \(|f(x_n) - f(y_n)| \geq \varepsilon_0\)。

由波尔查诺-魏尔斯特拉斯定理，\(\{x_n\}\) 有收敛子列 \(x_{n_k} \to \xi \in [a, b]\)。由于 \(|y_{n_k} - x_{n_k}| < \frac{1}{n_k} \to 0\)，故 \(y_{n_k} \to \xi\)。

由 \(f\) 的连续性，\(f(x_{n_k}) \to f(\xi)\)，\(f(y_{n_k}) \to f(\xi)\)，故 \(|f(x_{n_k}) - f(y_{n_k})| \to 0\)，与 \(|f(x_{n_k}) - f(y_{n_k})| \geq \varepsilon_0\) 矛盾。\(\blacksquare\)

\textbf{定理 22.21a}（Lipschitz 连续性蕴含一致连续性）：若 \(f\) 在区间 \(I\) 上满足 Lipschitz 条件，即存在常数 \(L > 0\) 使得对所有 \(x, y \in I\) 有 \(|f(x) - f(y)| \leq L|x - y|\)，则 \(f\) 在 \(I\) 上一致连续。

\textbf{证明}：对任意 \(\varepsilon > 0\)，取 \(\delta = \varepsilon / L\)。当 \(|x - y| < \delta\) 时，\(|f(x) - f(y)| \leq L|x - y| < L \cdot \varepsilon / L = \varepsilon\)。因此 \(f\) 在 \(I\) 上一致连续。注意 \(\delta\) 的选取只依赖于 \(\varepsilon\) 和 Lipschitz 常数 \(L\)，与 \(x, y\) 的具体位置无关。\(\blacksquare\)

\textbf{定理 22.21b}（一致连续函数的 Cauchy 列保持性）：若 \(f\) 在区间 \(I\) 上一致连续，则对 \(I\) 中的任意 Cauchy 列 \(\{x_n\}\)，\(\{f(x_n)\}\) 也是 Cauchy 列。

\textbf{证明}：设 \(\{x_n\}\) 是 \(I\) 中的 Cauchy 列。对任意 \(\varepsilon > 0\)，由 \(f\) 的一致连续性，存在 \(\delta > 0\) 使得 \(|x - y| < \delta\) 时 \(|f(x) - f(y)| < \varepsilon\)。由 \(\{x_n\}\) 是 Cauchy 列，存在 \(N\) 使得 \(m, n > N\) 时 \(|x_m - x_n| < \delta\)。因此 \(|f(x_m) - f(x_n)| < \varepsilon\)，即 \(\{f(x_n)\}\) 是 Cauchy 列。\(\blacksquare\)

这一性质在一致连续函数的有界延拓理论中至关重要------如果 \(f\) 在 \((a, b)\) 上一致连续，则 \(f\) 可以唯一地延拓为 \([a, b]\) 上的连续函数。

\textbf{注记 22.5}（非一致连续的典型例子）：函数 \(f(x) = x^2\) 在 \(\mathbb{R}\) 上连续，但非一致连续。证明：取 \(\varepsilon = 1\)，对任意 \(\delta > 0\)，取 \(x = \frac{1}{\delta} + \frac{\delta}{2}\)，\(y = \frac{1}{\delta}\)，则 \(|x - y| = \frac{\delta}{2} < \delta\)，但 \(|f(x) - f(y)| = \left(\frac{1}{\delta} + \frac{\delta}{2}\right)^2 - \frac{1}{\delta^2} = 1 + \frac{\delta^2}{4} > 1 = \varepsilon\)。所以无法找到一个统一的 \(\delta\) 对所有 \(x\) 都适用------当 \(x\) 很大时，\(x^2\) 的变化率也很大。

\section{第69章：函数极限与数列极限的关系}\label{ux7b2c69ux7ae0ux51fdux6570ux6781ux9650ux4e0eux6570ux5217ux6781ux9650ux7684ux5173ux7cfb}

海涅定理（Heine's Theorem，又称归结原则）是连接函数极限与数列极限的桥梁。它断言：\(\lim_{x \to a} f(x) = L\) 当且仅当对于任意收敛于 \(a\) 且各项不等于 \(a\) 的数列 \(\{x_n\}\)，都有 \(\lim_{n \to \infty} f(x_n) = L\)。这意味着函数极限可以完全由数列极限来刻画------函数极限的任何性质都可以通过数列极限来验证。海涅定理在理论上具有深远意义：它使得数列极限的许多性质（如唯一性、四则运算法则）可以''传递''到函数极限。海涅定理的逆否命题也极其有用：要证明 \(\lim_{x \to a} f(x)\) 不存在，只需找到两个收敛于 \(a\) 的数列 \(\{x_n\}\) 和 \(\{y_n\}\)，使得 \(\lim f(x_n) \neq \lim f(y_n)\)。例如，要证明 \(\lim_{x \to 0} \sin(1/x)\) 不存在，取 \(x_n = 1/(2n\pi)\)（得 \(\sin(1/x_n) = 0\)）和 \(y_n = 1/(2n\pi + \pi/2)\)（得 \(\sin(1/y_n) = 1\)），极限不同，因此极限不存在。海涅定理还揭示了连续函数的一个重要刻画：\(f\) 在 \(a\) 处连续当且仅当 \(\lim_{n \to \infty} f(x_n) = f(a)\) 对所有 \(x_n \to a\) 成立。

\textbf{定理 22.22}（海涅定理/归结原则）设 \(f\) 在 \(a\) 的某个去心邻域内有定义。则 \(\lim_{x \to a} f(x) = L\) 当且仅当：对于任意收敛于 \(a\) 且各项不等于 \(a\) 的数列 \(\{x_n\}\)，都有 \(\lim_{n \to \infty} f(x_n) = L\)。

\textbf{证明}：

必要性：设 \(\lim_{x \to a} f(x) = L\)，\(\{x_n\}\) 满足 \(x_n \to a\)，\(x_n \neq a\)。对于 \(\varepsilon > 0\)，存在 \(\delta > 0\)，当 \(0 < |x - a| < \delta\) 时 \(|f(x) - L| < \varepsilon\)。又存在 \(N\)，当 \(n > N\) 时 \(|x_n - a| < \delta\)。因为 \(x_n \neq a\)，所以 \(0 < |x_n - a| < \delta\)，故 \(|f(x_n) - L| < \varepsilon\)。

充分性（反证法）：假设 \(\lim_{x \to a} f(x) \neq L\)。则存在 \(\varepsilon_0 > 0\)，对每个 \(\delta = \frac{1}{n}\)，存在 \(x_n\) 满足 \(0 < |x_n - a| < \frac{1}{n}\) 但 \(|f(x_n) - L| \geq \varepsilon_0\)。此时 \(x_n \to a\) 但 \(f(x_n) \not\to L\)，矛盾。\(\blacksquare\)

\textbf{定理 22.22a}（海涅定理的逆否应用------极限不存在的判定）：若存在两个收敛于 \(a\) 的数列 \(\{x_n\}\) 和 \(\{y_n\}\)（\(x_n \neq a\)，\(y_n \neq a\)），使得 \(\lim_{n \to \infty} f(x_n) \neq \lim_{n \to \infty} f(y_n)\)（或其中之一不存在），则 \(\lim_{x \to a} f(x)\) 不存在。

\textbf{证明}：这直接来自海涅定理的逆否命题。若 \(\lim_{x \to a} f(x)\) 存在，则所有收敛于 \(a\) 的数列都导出相同的极限，与假设矛盾。\(\blacksquare\)

\textbf{定理 22.22b}（海涅定理与连续性的关系）：\(f\) 在 \(a\) 处连续当且仅当对任意收敛于 \(a\) 的数列 \(\{x_n\}\)（不要求 \(x_n \neq a\)），有 \(\lim_{n \to \infty} f(x_n) = f(a)\)。

\textbf{证明}：与定理 22.22 的证明完全类似，只需注意到此时去心邻域的条件 \(0 < |x_n - a|\) 可以放宽为 \(|x_n - a| < \delta\)（不排除 \(x_n = a\) 的情况），因为 \(f\) 在 \(a\) 处有定义且 \(f(a)\) 就是极限值。\(\blacksquare\)

海涅定理是分析学中最重要的''桥梁定理''之一。它使得数列极限的许多性质可以直接传递到函数极限，也为函数极限的反例构造提供了有力工具。例如，要证明 \(\lim_{x \to 0} \sin(1/x)\) 不存在，取 \(x_n = 1/(2n\pi)\)（得 \(\sin(1/x_n) = 0\)）和 \(y_n = 1/(2n\pi + \pi/2)\)（得 \(\sin(1/y_n) = 1\)），由定理 22.22a，极限不存在。

\hrulefill

\section{第70章：实数完备性的等价表述}\label{ux7b2c70ux7ae0ux5b9eux6570ux5b8cux5907ux6027ux7684ux7b49ux4ef7ux8868ux8ff0}

实数系的完备性有多种等价表述，它们从不同角度刻画了实数系''没有空隙''的本质。确界存在定理（有上界的非空实数集必有上确界）是其中最基础的表述，它等价于单调有界定理（单调有界数列必收敛）。这两个定理在 22.1 节中已经用于证明重要极限 \(\lim_{n \to \infty} (1+1/n)^n = e\) 的存在性。柯西收敛准则（数列收敛当且仅当它是柯西列）给出了一个不依赖于极限值的收敛性判断标准，它是实数完备性的另一种等价表述。区间套定理（闭区间套的交集非空）则从几何视角刻画了完备性------递增的实数序列''夹逼''出一个唯一的实数。这些等价表述构成了分析学的逻辑基础，每一条都可以作为实数完备性的公理，其他几条作为定理推导出来。在高等分析中，这些等价表述被推广到更一般的空间（如完备度量空间），构成泛函分析的基础。本节将系统展示这五条等价表述之间的逻辑关系，揭示实数系的深层结构。确界存在定理、单调有界定理、柯西收敛准则、区间套定理和有限覆盖定理，这五大定理的等价性证明是分析学中最优美的逻辑链条之一，体现了''同一个数学事实可以从不同视角被等价地刻画''这一深刻思想。

\textbf{定理 22.23} 以下命题等价：

\begin{enumerate}
\def\labelenumi{\arabic{enumi}.}
\tightlist
\item
  \textbf{确界存在定理}：非空有上界的实数集必有上确界。
\item
  \textbf{单调有界定理}：单调递增有上界的数列必收敛。
\item
  \textbf{闭区间套定理}：设 \(\{[a_n, b_n]\}\) 是递缩闭区间序列（\([a_{n+1}, b_{n+1}] \subseteq [a_n, b_n]\)），且 \(b_n - a_n \to 0\)，则存在唯一的 \(\xi \in \bigcap_{n=1}^{\infty}[a_n, b_n]\)。
\item
  \textbf{波尔查诺-魏尔斯特拉斯定理}：有界数列必有收敛子列。
\item
  \textbf{柯西收敛准则}：实数数列收敛当且仅当它是柯西列。
\end{enumerate}

\textbf{证明}（循环推导）：这五个命题的等价性可以通过以下逻辑链条建立：\((1) \Rightarrow (2) \Rightarrow (3) \Rightarrow (4) \Rightarrow (5) \Rightarrow (1)\)。\((1) \Rightarrow (2)\)：设 \(\{a_n\}\) 单调递增有上界，由确界存在定理，\(\{a_n\}\) 有上确界 \(L = \sup\{a_n\}\)。对任意 \(\varepsilon > 0\)，\(L - \varepsilon\) 不是上界，故存在 \(N\) 使 \(a_N > L - \varepsilon\)。由单调性，当 \(n > N\) 时 \(a_n \geq a_N > L - \varepsilon\)，而 \(a_n \leq L\)，故 \(|a_n - L| < \varepsilon\)，即 \(\lim a_n = L\)。\((2) \Rightarrow (3)\)：设 \(\{[a_n, b_n]\}\) 是递缩闭区间套，由单调性知 \(\{a_n\}\) 递增有上界（\(b_1\)）、\(\{b_n\}\) 递减有下界（\(a_1\)）。由单调有界定理，\(\lim a_n = a\)，\(\lim b_n = b\)。由 \(b_n - a_n \to 0\) 得 \(a = b = \xi\)。对任意 \(n\)，\(a_n \leq \xi \leq b_n\)，故 \(\xi \in \bigcap [a_n, b_n]\)。唯一性：若 \(\xi_1, \xi_2 \in \bigcap [a_n, b_n]\)，则 \(|\xi_1 - \xi_2| \leq b_n - a_n \to 0\)，故 \(\xi_1 = \xi_2\)。\(\blacksquare\)

\textbf{补充证明}（\(3 \Rightarrow 4\)）：设 \(\{x_n\}\) 为有界数列，即 \(x_n \in [a, b]\)（对一切 \(n\)）。将 \([a, b]\) 二分为 \(\left[a, \frac{a+b}{2}\right]\) 和 \(\left[\frac{a+b}{2}, b\right]\)，其中至少有一个区间包含 \(\{x_n\}\) 的无穷多项（否则整个数列只有有限项），取该区间为 \(I_1\)。继续将 \(I_1\) 二分，同理取包含无穷多项的子区间为 \(I_2\)。如此递推，得到嵌套闭区间序列 \(I_1 \supset I_2 \supset I_3 \supset \cdots\)，且 \(\lvert I_n \rvert = \frac{b - a}{2^n} \to 0\)。由闭区间套定理（3），\(\bigcap_{n=1}^{\infty} I_n = \{x^*\}\)（唯一一点）。现在在每个 \(I_n\) 中取一项 \(x_{n_k}\)（可取 \(n_k > n_{k-1}\) 使下标递增），则 \(x_{n_k} \in I_k\)，故 \(|x_{n_k} - x^*| \leq |I_k| = \frac{b-a}{2^k} \to 0\)。因此 \(\{x_n\}\) 有收敛子列 \(x_{n_k} \to x^*\)。\(\blacksquare\)

\textbf{补充证明}（\(5 \Rightarrow 1\)）：设 \(S \subseteq \mathbb{R}\) 非空且有上界。取 \(a_0 \in S\) 和一个上界 \(b_0\)（即 \(\forall x \in S,\, x \leq b_0\)），则 \(a_0 \leq b_0\)。对区间 \([a_n, b_n]\) 进行二分，令 \(m_n = \frac{a_n + b_n}{2}\)：若 \(m_n\) 是 \(S\) 的上界，则令 \([a_{n+1}, b_{n+1}] = [a_n, m_n]\)；否则存在 \(S\) 中某元素 \(> m_n\)，令 \([a_{n+1}, b_{n+1}] = [m_n, b_n]\)。由此构造的序列 \(\{a_n\}\) 满足 \(b_n - a_n = \frac{b_0 - a_0}{2^n} \to 0\)，且对 \(m > n\) 有 \(|a_m - a_n| \leq b_n - a_n \to 0\)，故 \(\{a_n\}\) 是柯西列。由柯西收敛准则（5），\(\{a_n\}\) 收敛，设 \(\xi = \lim_{n \to \infty} a_n\)。

下证 \(\xi = \sup S\)。首先，对每个 \(n\)，\(a_n \leq \xi \leq b_n\)（由构造，\(a_n\) 递增趋于 \(\xi\)）。对任意 \(x \in S\)，由构造每个 \(b_n\) 都是 \(S\) 的上界，故 \(x \leq b_n\)，取极限得 \(x \leq \xi\)，所以 \(\xi\) 是 \(S\) 的上界。其次，对任意 \(\varepsilon > 0\)，取 \(n\) 足够大使得 \(b_n - a_n < \varepsilon\)，则 \(a_n > \xi - \varepsilon\)。由构造，\(a_n\) 不是 \(S\) 的上界（否则会被取到右半区间），故存在 \(x \in S\) 使得 \(x > a_n > \xi - \varepsilon\)。这说明 \(\xi - \varepsilon\) 不是 \(S\) 的上界，因此 \(\xi\) 是最小上界，即 \(\xi = \sup S\)。\(\blacksquare\)

\textbf{注}：这五个等价命题共同构成了实数完备性理论的基石，是整个分析学的逻辑起点。

\textbf{本章展望}：本章建立了极限与连续的严格理论。极限理论使我们能够严格讨论''趋近''和''无穷逼近''的概念。在下一章（Ch 30）中，我们将利用极限理论定义导数------函数在某一点的瞬时变化率。导数是微积分的两大核心概念之一，它将极限理论应用于研究函数的局部性质。

\hrulefill

\hrulefill

\hrulefill

\hrulefill

\hrulefill

\newpage

\chapter{03-函数与微积分/03-微积分初步/02-导数与微分.md}\label{ux51fdux6570ux4e0eux5faeux79efux520603-ux5faeux79efux5206ux521dux6b6502-ux5bfcux6570ux4e0eux5faeux5206.md}

\textbf{前置知识}：本章依赖 Ch 29（极限与连续，特别是函数极限的 \(\varepsilon\)-\(\delta\) 定义、四则运算和夹逼定理）、Ch 27（三角恒等变换，为三角函数的导数公式做准备）、Ch 26（重要极限 \(\lim_{x \to 0} \frac{\sin x}{x} = 1\)）。

\hrulefill

\section{第71章：导数基础}\label{ux7b2c71ux7ae0ux5bfcux6570ux57faux7840}

导数是微积分中最核心的概念之一，它精确地刻画了函数在某一点处的瞬时变化率。从几何角度看，导数 \(f'(x_0)\) 是曲线 \(y = f(x)\) 在点 \((x_0, f(x_0))\) 处切线的斜率；从物理角度看，导数 \(s'(t)\) 是物体在时刻 \(t\) 的瞬时速度。导数的极限定义 \(f'(x_0) = \lim_{\Delta x \to 0} \frac{f(x_0 + \Delta x) - f(x_0)}{\Delta x}\) 将''变化率''这一直觉概念转化为严格的数学语言。历史上，牛顿从物理学的角度（流数法）和莱布尼茨从几何学的角度（切线问题）分别独立地发明了微积分，导数的现代定义融合了两者的视角。本节从切线问题的几何直觉出发，逐步建立导数的严格定义，并讨论单侧导数、可导性与连续性的关系（可导必连续，但连续未必可导）等基本问题。

\subsubsection{23.1.1 从切线问题引入}\label{ux4eceux5207ux7ebfux95eeux9898ux5f15ux5165}

考虑曲线 \(y = f(x)\) 上的点 \(P = (a, f(a))\)。在曲线上取另一点 \(Q = (a + h, f(a + h))\)，则割线 \(PQ\) 的斜率为

\[\frac{f(a + h) - f(a)}{h}\]

当 \(h \to 0\) 时，割线 \(PQ\) 趋于曲线在 \(P\) 点的切线，切线的斜率为

\[\lim_{h \to 0} \frac{f(a + h) - f(a)}{h}\]

\subsubsection{23.1.2 导数的定义}\label{ux5bfcux6570ux7684ux5b9aux4e49}

\textbf{定义 23.1}（导数）设 \(f\) 在 \(a\) 的某个邻域内有定义。若极限

\[\lim_{h \to 0} \frac{f(a + h) - f(a)}{h}\]

存在，则称 \(f\) 在 \(a\) 处\textbf{可导}，此极限值称为 \(f\) 在 \(a\) 处的\textbf{导数}，记作 \(f'(a)\) 或 \(\frac{df}{dx}\big|_{x=a}\)。

等价地，

\[f'(a) = \lim_{x \to a} \frac{f(x) - f(a)}{x - a}\]

\textbf{定义 23.2}（导函数）若 \(f\) 在区间 \(I\) 上的每一点都可导，则 \(f': I \to \mathbb{R}\)，

\[f'(x) = \lim_{h \to 0} \frac{f(x + h) - f(x)}{h}\]

称为 \(f\) 的\textbf{导函数}（简称导数）。

\subsubsection{23.1.3 可导与连续的关系}\label{ux53efux5bfcux4e0eux8fdeux7eedux7684ux5173ux7cfb}

\textbf{定理 23.1}（可导必连续）若 \(f\) 在 \(a\) 处可导，则 \(f\) 在 \(a\) 处连续。

\textbf{证明}：设 \(f'(a) = \lim_{h \to 0} \frac{f(a+h) - f(a)}{h}\) 存在。则

\[\lim_{h \to 0} [f(a+h) - f(a)] = \lim_{h \to 0} \frac{f(a+h) - f(a)}{h} \cdot h = f'(a) \cdot 0 = 0\]

所以 \(\lim_{h \to 0} f(a+h) = f(a)\)，即 \(f\) 在 \(a\) 处连续。\(\blacksquare\)

\textbf{注}：逆命题不成立。例如 \(f(x) = |x|\) 在 \(x = 0\) 处连续但不可导（左导数为 \(-1\)，右导数为 \(1\)，不相等）。

\subsubsection{23.1.4 导数的几何意义}\label{ux5bfcux6570ux7684ux51e0ux4f55ux610fux4e49}

\(f'(a)\) 表示曲线 \(y = f(x)\) 在点 \((a, f(a))\) 处的切线斜率。

切线方程：\(y - f(a) = f'(a)(x - a)\)

法线方程：\(y - f(a) = -\frac{1}{f'(a)}(x - a)\)（\(f'(a) \neq 0\) 时）

\hrulefill

\subsection{23.2 基本初等函数的导数}\label{ux57faux672cux521dux7b49ux51fdux6570ux7684ux5bfcux6570}

\textbf{定理 23.2} 对于正整数 \(n\)，\((x^n)' = nx^{n-1}\)。

\textbf{证明}：

\[(x^n)' = \lim_{h \to 0} \frac{(x+h)^n - x^n}{h} = \lim_{h \to 0} \frac{\sum_{k=0}^{n} \binom{n}{k} x^{n-k} h^k - x^n}{h}\]

\[= \lim_{h \to 0} \left(nx^{n-1} + \binom{n}{2}x^{n-2}h + \cdots + h^{n-1}\right) = nx^{n-1}\]

\(\blacksquare\)

\textbf{定理 23.3} \((\sin x)' = \cos x\)。

\textbf{证明}：

\[(\sin x)' = \lim_{h \to 0} \frac{\sin(x+h) - \sin x}{h}\]

由和差化积公式，\(\sin(x+h) - \sin x = 2\cos\left(x + \frac{h}{2}\right)\sin\frac{h}{2}\)。

\[(\sin x)' = \lim_{h \to 0} \cos\left(x + \frac{h}{2}\right) \cdot \frac{\sin\frac{h}{2}}{\frac{h}{2}}\]

由 \(\lim_{t \to 0} \frac{\sin t}{t} = 1\)（定理 22.12），得 \((\sin x)' = \cos x \cdot 1 = \cos x\)。\(\blacksquare\)

\textbf{定理 23.4} \((\cos x)' = -\sin x\)。

\textbf{证明}：

\[(\cos x)' = \lim_{h \to 0} \frac{\cos(x+h) - \cos x}{h}\]

由和差化积，\(\cos(x+h) - \cos x = -2\sin\left(x + \frac{h}{2}\right)\sin\frac{h}{2}\)。

类似地得 \((\cos x)' = -\sin x\)。\(\blacksquare\)

\textbf{定理 23.5} \((e^x)' = e^x\)。

\textbf{证明}：

\[(e^x)' = \lim_{h \to 0} \frac{e^{x+h} - e^x}{h} = e^x \lim_{h \to 0} \frac{e^h - 1}{h}\]

令 \(t = e^h - 1\)，则 \(h = \ln(1+t)\)，当 \(h \to 0\) 时 \(t \to 0\)。

\[\lim_{h \to 0} \frac{e^h - 1}{h} = \lim_{t \to 0} \frac{t}{\ln(1+t)} = \lim_{t \to 0} \frac{1}{\frac{1}{t}\ln(1+t)} = \frac{1}{\ln e} = 1\]

因此 \((e^x)' = e^x\)。\(\blacksquare\)

\textbf{定理 23.6} \((a^x)' = a^x \ln a\)（\(a > 0\)，\(a \neq 1\)）。

\textbf{证明}：\(a^x = e^{x \ln a}\)。这是指数函数的基本恒等式，因为 \(a^x = (e^{\ln a})^x = e^{x \ln a}\)（这一定义对任意 \(a > 0\) 成立，在 Ch 25 中有详细论述）。现在应用链式法则（定理 23.12）：令 \(u = x \ln a\)，则 \(a^x = e^u\)。由链式法则，

\[(a^x)' = (e^u)' \cdot u' = e^u \cdot \ln a = e^{x \ln a} \cdot \ln a = a^x \ln a\]

关键步骤是：\((e^u)' = e^u\)（定理 23.5）和 \(u' = \ln a\)（因为 \(\ln a\) 是常数）。这里链式法则的使用是合法的，因为 \(u(x) = x \ln a\) 可导，\(f(u) = e^u\) 在 \(u\) 处可导。\(\blacksquare\)

\textbf{定理 23.7} \((\ln x)' = \frac{1}{x}\)（\(x > 0\)）。

\textbf{证明}：

\[(\ln x)' = \lim_{h \to 0} \frac{\ln(x+h) - \ln x}{h} = \lim_{h \to 0} \frac{1}{h}\ln\left(1 + \frac{h}{x}\right)\]

令 \(t = \frac{h}{x}\)，则 \(h = xt\)，当 \(h \to 0\) 时 \(t \to 0\)。

\[(\ln x)' = \frac{1}{x} \lim_{t \to 0} \frac{\ln(1+t)}{t} = \frac{1}{x} \cdot 1 = \frac{1}{x}\]

\(\blacksquare\)

\textbf{定理 23.8} \((\tan x)' = \sec^2 x = \frac{1}{\cos^2 x}\)。

\textbf{证明}：由除法法则，

\[(\tan x)' = \left(\frac{\sin x}{\cos x}\right)' = \frac{\cos x \cdot \cos x - \sin x \cdot (-\sin x)}{\cos^2 x} = \frac{\cos^2 x + \sin^2 x}{\cos^2 x} = \frac{1}{\cos^2 x}\]

\(\blacksquare\)

\hrulefill

\subsection{23.3 求导法则}\label{ux6c42ux5bfcux6cd5ux5219}

求导法则将导数计算从极限定义中解放出来，使得我们可以通过代数运算而非极限计算来求导。四则运算法则------\((f \pm g)' = f' \pm g'\)、\((cf)' = cf'\)、\((fg)' = f'g + fg'\)、\((\frac{f}{g})' = \frac{f'g - fg'}{g^2}\)------是求导法的基石，它们将复杂函数的导数分解为简单函数导数的组合。复合函数求导法则（链式法则）\(\frac{dy}{dx} = \frac{dy}{du} \cdot \frac{du}{dx}\) 是求导法则中最强大的工具------它使得我们可以''层层剥离''嵌套函数，将任意复杂的复合函数的导数转化为基本初等函数导数的乘积。反函数求导法则 \([f^{-1}(x)]' = \frac{1}{f'(f^{-1}(x))}\) 则为反三角函数、对数函数等反函数的导数提供了统一的计算方法。这些求导法则共同构成了一个完整的''导数计算引擎''------只要掌握了基本初等函数的导数公式，就可以通过这组法则计算出任意初等函数的导数。

\subsubsection{23.3.1 四则运算求导法则}\label{ux56dbux5219ux8fd0ux7b97ux6c42ux5bfcux6cd5ux5219}

\textbf{定理 23.9}（和的导数）\((f + g)' = f' + g'\)

\textbf{证明}：

\[(f + g)'(x) = \lim_{h \to 0} \frac{[f(x+h) + g(x+h)] - [f(x) + g(x)]}{h} = f'(x) + g'(x)\]

\(\blacksquare\)

\textbf{定理 23.10}（常数倍的导数）\((cf)' = cf'\)

\textbf{证明}：

\[(cf)'(x) = \lim_{h \to 0} \frac{cf(x+h) - cf(x)}{h} = c \lim_{h \to 0} \frac{f(x+h) - f(x)}{h} = cf'(x)\]

\(\blacksquare\)

\textbf{定理 23.11}（乘法法则）\((fg)' = f'g + fg'\)

\textbf{证明}：

\[(fg)'(x) = \lim_{h \to 0} \frac{f(x+h)g(x+h) - f(x)g(x)}{h}\]

在分子中加减 \(f(x+h)g(x)\)：

\[= \lim_{h \to 0} \left[f(x+h) \cdot \frac{g(x+h) - g(x)}{h} + g(x) \cdot \frac{f(x+h) - f(x)}{h}\right]\]

由于 \(f\) 在 \(x\) 处可导，所以 \(\lim_{h \to 0} f(x+h) = f(x)\)（可导必连续）。因此

\[= f(x)g'(x) + g(x)f'(x)\]

\(\blacksquare\)

\textbf{除法法则}：若 \(g(x) \neq 0\)，则 \(\left(\frac{f}{g}\right)' = \frac{f'g - fg'}{g^2}\)。

\textbf{证明}：设 \(q = f/g\)，则 \(f = qg\)。由乘法法则，\(f' = q'g + qg'\)。解出 \(q'\)：

\[q' = \frac{f' - qg'}{g} = \frac{f' - \frac{f}{g} \cdot g'}{g} = \frac{f'g - fg'}{g^2}\]

\(\blacksquare\)

\subsubsection{23.3.2 链式法则}\label{ux94feux5f0fux6cd5ux5219}

\textbf{定理 23.12}（链式法则）设 \(g\) 在 \(x\) 处可导，\(f\) 在 \(g(x)\) 处可导，则 \(f \circ g\) 在 \(x\) 处可导，且

\[(f \circ g)'(x) = f'(g(x)) \cdot g'(x)\]

\textbf{证明}：设 \(u = g(x)\)，\(\Delta u = g(x+h) - g(x)\)。当 \(h \to 0\) 时，\(\Delta u \to 0\)（因为 \(g\) 连续）。

定义函数 \(\varphi\)：

\[\varphi(\Delta u) = \begin{cases} \frac{f(u + \Delta u) - f(u)}{\Delta u}, & \Delta u \neq 0 \\ f'(u), & \Delta u = 0 \end{cases}\]

则 \(\varphi\) 在 \(\Delta u = 0\) 处连续（因为 \(\lim_{\Delta u \to 0} \varphi(\Delta u) = f'(u) = \varphi(0)\)）。

\[\frac{f(g(x+h)) - f(g(x))}{h} = \varphi(\Delta u) \cdot \frac{g(x+h) - g(x)}{h}\]

当 \(h \to 0\) 时，\(\Delta u \to 0\)，\(\varphi(\Delta u) \to f'(u) = f'(g(x))\)，\(\frac{g(x+h) - g(x)}{h} \to g'(x)\)。

因此 \((f \circ g)'(x) = f'(g(x)) \cdot g'(x)\)。\(\blacksquare\)

\hrulefill

\section{第72章：中值定理与洛必达法则}\label{ux7b2c72ux7ae0ux4e2dux503cux5b9aux7406ux4e0eux6d1bux5fc5ux8fbeux6cd5ux5219}

微分中值定理是微分学的理论核心，它们建立了函数在区间上的整体性质（如函数值的差）与局部性质（如某点处的导数）之间的桥梁。罗尔定理是最简单的中值定理：若 \(f(a) = f(b)\)，则在 \((a, b)\) 内存在一点 \(c\) 使得 \(f'(c) = 0\)------其几何意义是两端点等高的光滑曲线至少有一个水平切线点。拉格朗日中值定理是罗尔定理的推广，它去掉了 \(f(a) = f(b)\) 的限制，断言存在一点 \(c\) 使得 \(f'(c) = \frac{f(b) - f(a)}{b - a}\)------即函数在区间上的平均变化率等于某点处的瞬时变化率。柯西中值定理则进一步推广到两个函数的情形，它是洛必达法则的理论基础。这三个中值定理的证明巧妙地利用了辅助函数法------通过构造满足罗尔定理条件的辅助函数，将拉格朗日和柯西中值定理归结为罗尔定理的推论。中值定理的应用极为广泛：从证明不等式到判定函数的单调性，从洛必达法则的推导到泰勒公式的余项估计，都离不开中值定理的支撑。

\subsubsection{23.4.1 费马定理}\label{ux8d39ux9a6cux5b9aux7406}

\textbf{定理 23.13}（费马定理）设 \(f\) 在 \(a\) 处取极值，且 \(f\) 在 \(a\) 处可导，则 \(f'(a) = 0\)。

\textbf{证明}：不妨设 \(f(a)\) 是极大值。则存在 \(\delta > 0\)，使得当 \(|x - a| < \delta\) 时，\(f(x) \leq f(a)\)。

当 \(a < x < a + \delta\) 时，\(\frac{f(x) - f(a)}{x - a} \leq 0\)，所以 \(f'(a) = \lim_{x \to a^+} \frac{f(x) - f(a)}{x - a} \leq 0\)。

当 \(a - \delta < x < a\) 时，\(\frac{f(x) - f(a)}{x - a} \geq 0\)，所以 \(f'(a) = \lim_{x \to a^-} \frac{f(x) - f(a)}{x - a} \geq 0\)。

因此 \(f'(a) = 0\)。\(\blacksquare\)

\subsubsection{23.4.2 罗尔定理}\label{ux7f57ux5c14ux5b9aux7406}

\textbf{定理 23.14}（罗尔定理）设 \(f\) 在 \([a, b]\) 上连续，在 \((a, b)\) 上可导，且 \(f(a) = f(b)\)。则存在 \(\xi \in (a, b)\)，使得 \(f'(\xi) = 0\)。

\textbf{证明}：由最值定理（定理 22.16），\(f\) 在 \([a, b]\) 上取到最大值 \(M\) 和最小值 \(m\)。

若 \(M = m\)，则 \(f\) 为常数，\(f'(\xi) = 0\) 对所有 \(\xi\) 成立。

若 \(M \neq m\)，则 \(M\) 和 \(m\) 中至少有一个不等于 \(f(a) = f(b)\)。不妨设 \(M \neq f(a)\)，则最大值在某个 \(\xi \in (a, b)\) 处取到。由费马定理，\(f'(\xi) = 0\)。\(\blacksquare\)

\subsubsection{23.4.3 拉格朗日中值定理}\label{ux62c9ux683cux6717ux65e5ux4e2dux503cux5b9aux7406}

\textbf{定理 23.15}（拉格朗日中值定理）设 \(f\) 在 \([a, b]\) 上连续，在 \((a, b)\) 上可导。则存在 \(\xi \in (a, b)\)，使得

\[f'(\xi) = \frac{f(b) - f(a)}{b - a}\]

\textbf{证明}：定义辅助函数

\[F(x) = f(x) - f(a) - \frac{f(b) - f(a)}{b - a}(x - a)\]

则 \(F\) 在 \([a, b]\) 上连续，在 \((a, b)\) 上可导，且 \(F(a) = F(b) = 0\)。

由罗尔定理，存在 \(\xi \in (a, b)\) 使得 \(F'(\xi) = 0\)，即

\[f'(\xi) - \frac{f(b) - f(a)}{b - a} = 0\]

所以 \(f'(\xi) = \frac{f(b) - f(a)}{b - a}\)。\(\blacksquare\)

\subsubsection{23.4.4 柯西中值定理}\label{ux67efux897fux4e2dux503cux5b9aux7406}

\textbf{定理 23.16}（柯西中值定理）设 \(f\)，\(g\) 在 \([a, b]\) 上连续，在 \((a, b)\) 上可导，且 \(g'(x) \neq 0\)。则存在 \(\xi \in (a, b)\)，使得

\[\frac{f'(\xi)}{g'(\xi)} = \frac{f(b) - f(a)}{g(b) - g(a)}\]

\textbf{证明}：首先说明 \(g(b) \neq g(a)\)：若 \(g(b) = g(a)\)，则由罗尔定理，存在 \(c \in (a, b)\) 使 \(g'(c) = 0\)，与条件 \(g'(x) \neq 0\) 矛盾。

定义辅助函数

\[F(x) = f(x) - f(a) - \frac{f(b) - f(a)}{g(b) - g(a)}[g(x) - g(a)]\]

则 \(F(a) = F(b) = 0\)。由罗尔定理，存在 \(\xi \in (a, b)\) 使得 \(F'(\xi) = 0\)，即

\[f'(\xi) - \frac{f(b) - f(a)}{g(b) - g(a)} g'(\xi) = 0\]

由于 \(g'(\xi) \neq 0\)，化简得 \(\frac{f'(\xi)}{g'(\xi)} = \frac{f(b) - f(a)}{g(b) - g(a)}\)。\(\blacksquare\)

\hrulefill

\subsection{23.5 洛必达法则}\label{ux6d1bux5fc5ux8fbeux6cd5ux5219}

\textbf{定理 23.17}（洛必达法则，\(\frac{0}{0}\) 型）设 \(f\)，\(g\) 在 \(a\) 的某个去心邻域内可导，\(g'(x) \neq 0\)，且

\[\lim_{x \to a} f(x) = 0, \quad \lim_{x \to a} g(x) = 0\]

若 \(\lim_{x \to a} \frac{f'(x)}{g'(x)} = L\)（\(L\) 可以是有限数或 \(\pm\infty\)），则

\[\lim_{x \to a} \frac{f(x)}{g(x)} = L\]

\textbf{证明}：补充定义 \(f(a) = g(a) = 0\)，则 \(f\) 和 \(g\) 在 \(a\) 处连续。

由于 \(g'(x) \neq 0\) 在 \(a\) 的去心邻域内成立且 \(g\) 连续，\(g\) 在 \(a\) 附近严格单调，故对 \(x \neq a\) 有 \(g(x) \neq g(a) = 0\)。因此对 \(x \neq a\)，可以应用柯西中值定理（定理 23.16）：存在 \(\xi\) 介于 \(a\) 和 \(x\) 之间，使得

\[\frac{f(x)}{g(x)} = \frac{f(x) - f(a)}{g(x) - g(a)} = \frac{f'(\xi)}{g'(\xi)}\]

当 \(x \to a\) 时，\(\xi \to a\)（因为 \(\xi\) 介于 \(a\) 和 \(x\) 之间），所以

\[\lim_{x \to a} \frac{f(x)}{g(x)} = \lim_{\xi \to a} \frac{f'(\xi)}{g'(\xi)} = L\]

\(\blacksquare\)

\textbf{注}：洛必达法则只是极限存在的充分条件，而非必要条件------当 \(\frac{f'(x)}{g'(x)}\) 的极限不存在时，\(\frac{f(x)}{g(x)}\) 的极限仍可能存在。

\hrulefill

\section{第73章：导数的应用}\label{ux7b2c73ux7ae0ux5bfcux6570ux7684ux5e94ux7528}

\textbf{定理 23.18}（单调性判定）设 \(f\) 在 \([a, b]\) 上连续，在 \((a, b)\) 上可导。

\begin{enumerate}
\def\labelenumi{\arabic{enumi}.}
\tightlist
\item
  若 \(\forall x \in (a, b), f'(x) > 0\)，则 \(f\) 在 \([a, b]\) 上严格递增。
\item
  若 \(\forall x \in (a, b), f'(x) < 0\)，则 \(f\) 在 \([a, b]\) 上严格递减。
\end{enumerate}

\textbf{证明}：设 \(x_1, x_2 \in [a, b]\)，\(x_1 < x_2\)。由拉格朗日中值定理，存在 \(\xi \in (x_1, x_2)\) 使得

\[f(x_2) - f(x_1) = f'(\xi)(x_2 - x_1)\]

因为 \(f'(\xi) > 0\) 且 \(x_2 - x_1 > 0\)，所以 \(f(x_2) > f(x_1)\)。\(\blacksquare\)

\textbf{定理 23.19}（极值的第一充分条件）设 \(f\) 在 \(a\) 处连续，在 \(a\) 的某个去心邻域内可导。

\begin{enumerate}
\def\labelenumi{\arabic{enumi}.}
\tightlist
\item
  若 \(f'\) 在 \(a\) 左侧为正、右侧为负，则 \(f(a)\) 是极大值。
\item
  若 \(f'\) 在 \(a\) 左侧为负、右侧为正，则 \(f(a)\) 是极小值。
\end{enumerate}

\textbf{证明}：(1) 由定理 23.18，\(f'(x) > 0\) 意味着 \(f\) 在 \(a\) 的左侧邻域严格递增，即当 \(x < a\) 且 \(x\) 充分接近 \(a\) 时 \(f(x) < f(a)\)。\(f'(x) < 0\) 意味着 \(f\) 在 \(a\) 的右侧邻域严格递减，即当 \(x > a\) 且 \(x\) 充分接近 \(a\) 时 \(f(x) < f(a)\)。综合两者，在 \(a\) 的某个去心邻域内 \(f(x) < f(a)\)，因此 \(f(a)\) 是极大值。这里的关键是单调性判定（定理 23.18）将导数符号转化为函数值的单调关系，进而将局部极值条件转化为导数符号在极值点两侧的变化。\(\blacksquare\)

\textbf{定理 23.20}（极值的第二充分条件）设 \(f\) 在 \(a\) 处有二阶导数，\(f'(a) = 0\)。

\begin{enumerate}
\def\labelenumi{\arabic{enumi}.}
\tightlist
\item
  若 \(f''(a) < 0\)，则 \(f(a)\) 是极大值。
\item
  若 \(f''(a) > 0\)，则 \(f(a)\) 是极小值。
\end{enumerate}

\textbf{证明}：(1) 由 \(f''(a) < 0\)，即 \(\lim_{h \to 0} \frac{f'(a+h) - f'(a)}{h} < 0\)。由于 \(f'(a) = 0\)，这等价于 \(\lim_{h \to 0} \frac{f'(a+h)}{h} < 0\)。由极限的局部保号性，存在 \(\delta > 0\)，使得当 \(0 < |h| < \delta\) 时 \(\frac{f'(a+h)}{h} < 0\)。这意味着：当 \(h > 0\)（即 \(x > a\)）时，\(f'(a+h) < 0\)（因为分子分母异号）；当 \(h < 0\)（即 \(x < a\)）时，\(f'(a+h) > 0\)（因为分子分母同号）。因此 \(f'\) 在 \(a\) 左侧为正、右侧为负。由第一充分条件（定理 23.19），\(f(a)\) 是极大值。\(\blacksquare\)

\hrulefill

\subsection{23.7 微分}\label{ux5faeux5206}

\textbf{定义 23.3}（微分）设 \(f\) 在 \(x\) 处可导。函数 \(f\) 在 \(x\) 处的\textbf{微分}定义为：

\[dy = f'(x) dx\]

其中 \(dx\) 是自变量的增量。

\textbf{定理 23.21}（微分与增量的关系）设 \(f\) 在 \(x\) 处可导，则

\[\Delta y = f(x + \Delta x) - f(x) = f'(x) \Delta x + o(\Delta x)\]

其中 \(\lim_{\Delta x \to 0} \frac{o(\Delta x)}{\Delta x} = 0\)。

\textbf{证明}：由导数定义，令 \(\varepsilon(\Delta x) = \frac{f(x + \Delta x) - f(x)}{\Delta x} - f'(x)\)，则 \(\lim_{\Delta x \to 0} \varepsilon(\Delta x) = 0\)。

\(\Delta y = [f'(x) + \varepsilon(\Delta x)] \Delta x = f'(x) \Delta x + \varepsilon(\Delta x) \Delta x\)。

令 \(o(\Delta x) = \varepsilon(\Delta x) \Delta x\)，则 \(\frac{o(\Delta x)}{\Delta x} = \varepsilon(\Delta x) \to 0\)。\(\blacksquare\)

\hrulefill

\subsection{23.8 高阶导数}\label{ux9ad8ux9636ux5bfcux6570}

高阶导数是对函数进行多次求导的结果。一阶导数 \(f'(x)\) 描述函数的变化率（速度），二阶导数 \(f''(x)\) 描述变化率的变化率（加速度），三阶及更高阶导数描述更精细的变化特征。高阶导数在物理学中有着自然的对应：位移 \(s(t)\) 的一阶导数是速度 \(v(t)\)，二阶导数是加速度 \(a(t)\)，三阶导数（加加速度，jerk）描述了加速度的变化率。在数学上，高阶导数与函数的泰勒展开密切相关------\(n\) 阶泰勒多项式恰好由函数在展开点处的各阶导数值确定：\(f(x) = f(a) + f'(a)(x-a) + \frac{f''(a)}{2!}(x-a)^2 + \cdots + \frac{f^{(n)}(a)}{n!}(x-a)^n + R_n\)。高阶导数的计算遵循递推法则------\(f^{(n)}(x) = [f^{(n-1)}(x)]'\)，对于某些函数（如 \(e^x\)、\(\sin x\)、\(\cos x\)），高阶导数呈现出周期性的规律。莱布尼茨公式 \((fg)^{(n)} = \sum_{k=0}^{n} \binom{n}{k} f^{(k)} g^{(n-k)}\) 是乘积函数高阶导数的通用公式，其结构与二项式定理高度相似------这是离散与连续之间的又一次深刻类比。

\subsubsection{23.8.1 二阶导数}\label{ux4e8cux9636ux5bfcux6570}

\textbf{定义 23.4}（二阶导数）设 \(f\) 在区间 \(I\) 上可导。若 \(f'\) 在 \(I\) 上也可导，则称 \(f'\) 的导数为 \(f\) 的\textbf{二阶导数}，记作 \(f''(x)\) 或 \(\frac{d^2 f}{dx^2}\)。

\textbf{定义 23.5}（\(n\) 阶导数）类似地，\(f^{(n)}(x) = (f^{(n-1)})'(x)\)（\(n \geq 1\)），其中 \(f^{(0)} = f\)。

\textbf{定理 23.22}（莱布尼茨公式）设 \(f\) 和 \(g\) 都有 \(n\) 阶导数，则

\[(fg)^{(n)} = \sum_{k=0}^{n} \binom{n}{k} f^{(k)} g^{(n-k)}\]

\textbf{证明}：对 \(n\) 用数学归纳法。

\(n = 1\)：\((fg)' = f'g + fg'\)，即 \(\binom{1}{0}f^{(0)}g^{(1)} + \binom{1}{1}f^{(1)}g^{(0)}\)，成立。

\(n \Rightarrow n+1\)：假设公式对 \(n\) 成立。则

\[(fg)^{(n+1)} = \left[(fg)^{(n)}\right]' = \sum_{k=0}^{n} \binom{n}{k} (f^{(k)} g^{(n-k)})'\]

\[= \sum_{k=0}^{n} \binom{n}{k} [f^{(k+1)} g^{(n-k)} + f^{(k)} g^{(n+1-k)}]\]

\[= \sum_{k=0}^{n} \binom{n}{k} f^{(k+1)} g^{(n-k)} + \sum_{k=0}^{n} \binom{n}{k} f^{(k)} g^{(n+1-k)}\]

在第一个和式中令 \(j = k + 1\)：

\[= \sum_{j=1}^{n+1} \binom{n}{j-1} f^{(j)} g^{(n+1-j)} + \sum_{k=0}^{n} \binom{n}{k} f^{(k)} g^{(n+1-k)}\]

\[= \sum_{k=0}^{n+1} \left[\binom{n}{k-1} + \binom{n}{k}\right] f^{(k)} g^{(n+1-k)} = \sum_{k=0}^{n+1} \binom{n+1}{k} f^{(k)} g^{(n+1-k)}\]

利用了帕斯卡恒等式 \(\binom{n}{k-1} + \binom{n}{k} = \binom{n+1}{k}\)。\(\blacksquare\)

\subsubsection{\texorpdfstring{23.8.2 常见函数的 \(n\) 阶导数}{23.8.2 常见函数的 n 阶导数}}\label{ux5e38ux89c1ux51fdux6570ux7684-n-ux9636ux5bfcux6570}

\textbf{定理 23.23} \((e^{ax})^{(n)} = a^n e^{ax}\)

\textbf{证明}：对 \(n\) 进行归纳。\(n = 1\) 时，\((e^{ax})' = ae^{ax}\)，由链式法则（定理 23.12）和 \((e^x)' = e^x\) 直接得到。假设 \((e^{ax})^{(k)} = a^k e^{ax}\) 对 \(k\) 成立。则

\[(e^{ax})^{(k+1)} = \left[(e^{ax})^{(k)}\right]' = (a^k e^{ax})' = a^k \cdot (e^{ax})' = a^k \cdot a e^{ax} = a^{k+1} e^{ax}\]

由数学归纳法原理，结论对任意正整数 \(n\) 成立。这一公式的简洁性源于指数函数 \(e^x\) 的独特性质------它是唯一满足 \(f' = f\) 的函数（在常数因子意义下），因此每求一次导只是多乘一个 \(a\)。\(\blacksquare\)

\textbf{定理 23.24} \((\sin x)^{(n)} = \sin\left(x + \frac{n\pi}{2}\right)\)

\textbf{证明}：\((\sin x)' = \cos x = \sin(x + \frac{\pi}{2})\)。假设 \((\sin x)^{(k)} = \sin(x + \frac{k\pi}{2})\)，则

\[(\sin x)^{(k+1)} = \cos(x + \frac{k\pi}{2}) = \sin(x + \frac{k\pi}{2} + \frac{\pi}{2}) = \sin(x + \frac{(k+1)\pi}{2})\]

\(\blacksquare\)

\hrulefill

\subsection{23.9 泰勒展开（初步）}\label{ux6cf0ux52d2ux5c55ux5f00ux521dux6b65}

\textbf{定理 23.25}（带有拉格朗日余项的泰勒公式）设 \(f\) 在 \([a, x]\) 上有 \(n+1\) 阶导数，则

\[f(x) = \sum_{k=0}^{n} \frac{f^{(k)}(a)}{k!}(x - a)^k + R_n(x)\]

其中 \(R_n(x) = \frac{f^{(n+1)}(\xi)}{(n+1)!}(x - a)^{n+1}\)，\(\xi\) 介于 \(a\) 和 \(x\) 之间。

\textbf{证明}：设 \(R_n(x) = f(x) - \sum_{k=0}^{n} \frac{f^{(k)}(a)}{k!}(x-a)^k\)。定义辅助函数

\[F(t) = f(x) - \sum_{k=0}^{n} \frac{f^{(k)}(t)}{k!}(x-t)^k\]

则 \(F(a) = R_n(x)\)，\(F(x) = 0\)。逐项求导：

\[F'(t) = -\sum_{k=0}^{n} \frac{d}{dt}\left[\frac{f^{(k)}(t)}{k!}(x-t)^k\right]\]

对每个 \(k \geq 1\)，由乘积法则：

\[\frac{d}{dt}\left[\frac{f^{(k)}(t)}{k!}(x-t)^k\right] = \frac{f^{(k+1)}(t)}{k!}(x-t)^k - \frac{f^{(k)}(t)}{(k-1)!}(x-t)^{k-1}\]

\(k = 0\) 时 \(\frac{d}{dt}[f(t)] = f'(t)\)。因此求和后相邻项两两相消（telescoping），仅剩 \(k = n\) 的首项：

\[F'(t) = -\frac{f^{(n+1)}(t)}{n!}(x-t)^n\]

对 \(F\) 和 \(G(t) = (x-t)^{n+1}\) 在 \([a, x]\) 上应用柯西中值定理：

\[\frac{F(x) - F(a)}{G(x) - G(a)} = \frac{F'(\xi)}{G'(\xi)}\]

\[\frac{-R_n(x)}{-(x-a)^{n+1}} = \frac{-\frac{f^{(n+1)}(\xi)}{n!}(x-\xi)^n}{-(n+1)(x-\xi)^n}\]

\[R_n(x) = \frac{f^{(n+1)}(\xi)}{(n+1)!}(x-a)^{n+1}\]

\(\blacksquare\)

\textbf{推论 23.2}（麦克劳林公式，\(a = 0\)）

\[f(x) = \sum_{k=0}^{n} \frac{f^{(k)}(0)}{k!}x^k + \frac{f^{(n+1)}(\xi)}{(n+1)!}x^{n+1}\]

常见展开：

\[e^x = 1 + x + \frac{x^2}{2!} + \frac{x^3}{3!} + \cdots\]

\[\sin x = x - \frac{x^3}{3!} + \frac{x^5}{5!} - \cdots\]

\[\cos x = 1 - \frac{x^2}{2!} + \frac{x^4}{4!} - \cdots\]

\[\ln(1+x) = x - \frac{x^2}{2} + \frac{x^3}{3} - \cdots \quad (|x| \leq 1, x \neq -1)\]

\textbf{证明}：在定理 23.25（泰勒公式）中令 \(a = 0\) 即得麦克劳林公式。具体而言，泰勒公式给出

\[f(x) = \sum_{k=0}^{n} \frac{f^{(k)}(a)}{k!}(x - a)^k + R_n(x)\]

将 \(a\) 替换为 \(0\)，各项变为 \(\frac{f^{(k)}(0)}{k!} x^k\)，余项变为 \(R_n(x) = \frac{f^{(n+1)}(\xi)}{(n+1)!} x^{n+1}\)（其中 \(\xi\) 介于 \(0\) 和 \(x\) 之间）。这就是麦克劳林公式。对于常见的初等函数，我们可以计算其在 \(x = 0\) 处的各阶导数，然后代入即得相应的展开式。例如 \(e^x\) 在 \(x = 0\) 处的各阶导数均为 \(1\)，故 \(e^x = 1 + x + \frac{x^2}{2!} + \cdots + \frac{x^n}{n!} + R_n\)。\(\blacksquare\)

\section{第74章：微分法与中值定理深化}\label{ux7b2c74ux7ae0ux5faeux5206ux6cd5ux4e0eux4e2dux503cux5b9aux7406ux6df1ux5316}

隐函数是由方程 \(F(x, y) = 0\) 确定的函数关系，它不显式地给出 \(y = f(x)\) 的形式，而是通过方程隐式地约束 \(x\) 和 \(y\) 之间的关系。隐函数求导法的核心思想是：将 \(y\) 视为 \(x\) 的函数，对等式 \(F(x, y) = 0\) 两边同时对 \(x\) 求导（使用链式法则处理 \(y\) 相关项），然后解出 \(\frac{dy}{dx}\)。例如，圆 \(x^2 + y^2 = r^2\) 两边求导得 \(2x + 2y \cdot y' = 0\)，从而 \(y' = -\frac{x}{y}\)。隐函数求导法避免了显式解出 \(y = \pm\sqrt{r^2 - x^2}\) 的麻烦，直接得到了切线斜率的简洁表达式。隐函数存在定理（定理 23.26）保证了在 \(F_y(x_0, y_0) \neq 0\) 的条件下，方程 \(F(x, y) = 0\) 在 \((x_0, y_0)\) 附近确实可以确定一个可导的隐函数 \(y = f(x)\)。这一理论是多元微积分中隐函数定理的一维特例，但它已经蕴含着深刻的几何意义------隐函数定理本质上是在说，如果切平面不是垂直的，则曲面 \(z = F(x, y)\) 与平面 \(z = 0\) 的交线在局部可以投影到 \(x\) 轴上。

\subsubsection{23.10.1 隐函数}\label{ux9690ux51fdux6570}

\textbf{定义 23.6}（隐函数）由方程 \(F(x, y) = 0\) 确定的 \(y\) 关于 \(x\) 的函数称为\textbf{隐函数}。

\textbf{定理 23.26}（隐函数求导法则）设方程 \(F(x, y) = 0\) 确定了隐函数 \(y = y(x)\)，且 \(F\) 对 \(x\) 和 \(y\) 的偏导数存在，\(F_y \neq 0\)。则

\[\frac{dy}{dx} = -\frac{F_x}{F_y}\]

\textbf{证明}：将 \(y\) 视为 \(x\) 的函数 \(y(x)\)，则由方程 \(F(x, y(x)) = 0\) 对 \(x\) 恒成立。对这个恒等式两边关于 \(x\) 求导。左边是二元函数 \(F(x, y)\) 与其内层函数 \((x, y(x))\) 的复合，因此需要应用多元链式法则（此处我们使用偏导数的记法，在 Ch 30 中已有引入）。对 \(F(x, y(x))\) 关于 \(x\) 求导：

\[\frac{d}{dx} F(x, y(x)) = \frac{\partial F}{\partial x} \cdot \frac{dx}{dx} + \frac{\partial F}{\partial y} \cdot \frac{dy}{dx} = F_x + F_y \cdot \frac{dy}{dx}\]

右边 \(0\) 的导数为 \(0\)。因此

\[F_x + F_y \cdot \frac{dy}{dx} = 0\]

由条件 \(F_y \neq 0\)，可以解出 \(\frac{dy}{dx}\)：

\[\frac{dy}{dx} = -\frac{F_x}{F_y}\]

这一公式的几何意义是：在曲面 \(z = F(x, y)\) 与平面 \(z = 0\) 的交线（隐函数图像）上，切线的斜率等于偏导数之比的负值。\(\blacksquare\)

\hrulefill

\subsection{23.11 参数方程求导}\label{ux53c2ux6570ux65b9ux7a0bux6c42ux5bfc}

许多曲线用参数方程 \(x = \varphi(t), y = \psi(t)\) 描述比用普通方程 \(y = f(x)\) 更为自然------例如圆周运动、摆线和弹道轨迹。参数方程求导的核心公式是 \(\frac{dy}{dx} = \frac{dy/dt}{dx/dt} = \frac{\psi'(t)}{\varphi'(t)}\)（当 \(\varphi'(t) \neq 0\) 时）。这一公式的直观理解是：\(\frac{dy}{dx}\) 是 \(y\) 相对于 \(x\) 的变化率，而 \(t\) 是中间变量，通过链式法则 \(\frac{dy}{dt} = \frac{dy}{dx} \cdot \frac{dx}{dt}\) 即可解出 \(\frac{dy}{dx}\)。二阶导数的参数方程公式 \(\frac{d^2y}{dx^2} = \frac{d}{dt}\left(\frac{dy}{dx}\right) / \frac{dx}{dt} = \frac{\psi''(t)\varphi'(t) - \psi'(t)\varphi''(t)}{[\varphi'(t)]^3}\) 更为复杂，需要特别注意------常见的错误是直接将 \(\frac{dy}{dx}\) 对 \(t\) 求导而忘记除以 \(\frac{dx}{dt}\)。参数方程求导在物理学中尤为重要------速度向量 \(\vec{v} = (x'(t), y'(t))\) 和加速度向量 \(\vec{a} = (x''(t), y''(t))\) 自然地由参数方程的导数给出，而 \(\frac{dy}{dx}\) 则是质点运动轨迹的切线斜率。参数方程求导还与隐函数求导有密切联系------二者都是通过中间变量（或约束关系）来避免直接处理 \(y = f(x)\) 的显式形式。

\textbf{定理 23.27} 设曲线由参数方程 \(x = \varphi(t)\)，\(y = \psi(t)\) 给出，\(\varphi\) 和 \(\psi\) 都可导，\(\varphi'(t) \neq 0\)。则

\[\frac{dy}{dx} = \frac{\psi'(t)}{\varphi'(t)}\]

\[\frac{d^2y}{dx^2} = \frac{\psi''(t)\varphi'(t) - \psi'(t)\varphi''(t)}{[\varphi'(t)]^3}\]

\textbf{证明}：

\[\frac{dy}{dx} = \frac{\frac{dy}{dt}}{\frac{dx}{dt}} = \frac{\psi'(t)}{\varphi'(t)}\]

\[\frac{d^2y}{dx^2} = \frac{\frac{d}{dt}\left(\frac{\psi'(t)}{\varphi'(t)}\right)}{\frac{dx}{dt}} = \frac{\frac{\psi''(t)\varphi'(t) - \psi'(t)\varphi''(t)}{[\varphi'(t)]^2}}{\varphi'(t)} = \frac{\psi''(t)\varphi'(t) - \psi'(t)\varphi''(t)}{[\varphi'(t)]^3}\]

\(\blacksquare\)

\hrulefill

\subsection{23.12 函数的凹凸性}\label{ux51fdux6570ux7684ux51f9ux51f8ux6027}

函数的凹凸性描述了函数图像的弯曲方向，它由二阶导数的符号决定。直观地说，凹函数（\(f''(x) > 0\)）的图像''向下弯曲''（如 \(y = x^2\)），其上任意两点间的连线位于函数图像的上方；凸函数（\(f''(x) < 0\)）的图像''向上弯曲''（如 \(y = -x^2\)），其上任意两点间的连线位于函数图像的下方。拐点（inflection point）是凹凸性发生改变的点，满足 \(f''(x) = 0\) 且二阶导数在该点两侧变号。凹凸性的判定定理（定理 23.29）是二阶导数在函数分析中的核心应用------它弥补了一阶导数只能判定单调性而无法判定弯曲方向的不足。凹凸性在经济学（边际效用递减由凸函数刻画）、优化理论（凹函数的最优值在边界达到）和数值分析中都有重要应用。凹凸性与詹森不等式（Jensen's Inequality）密切相关------对于凹函数 \(f\)，有 \(f(\frac{a+b}{2}) \geq \frac{f(a)+f(b)}{2}\)，这一不等式是概率论中许多重要不等式（如琴生不等式）的基础。

\subsubsection{23.12.1 凹凸函数的定义}\label{ux51f9ux51f8ux51fdux6570ux7684ux5b9aux4e49}

\textbf{定义 23.7}（凸函数）设 \(f\) 在区间 \(I\) 上有定义。若对任意 \(x_1, x_2 \in I\) 和 \(\lambda \in [0, 1]\)：

\[f(\lambda x_1 + (1-\lambda)x_2) \leq \lambda f(x_1) + (1-\lambda)f(x_2)\]

则称 \(f\) 在 \(I\) 上是\textbf{凸函数}（向下凸，图形呈 \(\cup\) 形）。若不等号反向，则称 \(f\) 是\textbf{凹函数}。

\subsubsection{23.12.2 凹凸性的判定}\label{ux51f9ux51f8ux6027ux7684ux5224ux5b9a}

\textbf{定理 23.28}（二阶导数判定法）设 \(f\) 在区间 \(I\) 上有二阶导数。

\begin{enumerate}
\def\labelenumi{\arabic{enumi}.}
\tightlist
\item
  若 \(\forall x \in I, f''(x) > 0\)，则 \(f\) 在 \(I\) 上是凸函数。
\item
  若 \(\forall x \in I, f''(x) < 0\)，则 \(f\) 在 \(I\) 上是凹函数。
\end{enumerate}

\textbf{证明}：(1) 设 \(x_1, x_2 \in I\)，\(x_1 < x_2\)，\(\lambda \in (0, 1)\)，\(x_0 = \lambda x_1 + (1-\lambda)x_2\)。

在 \([x_1, x_0]\) 和 \([x_0, x_2]\) 上分别应用拉格朗日中值定理：

\[\frac{f(x_0) - f(x_1)}{x_0 - x_1} = f'(\xi_1), \quad \frac{f(x_2) - f(x_0)}{x_2 - x_0} = f'(\xi_2)\]

其中 \(x_1 < \xi_1 < x_0 < \xi_2 < x_2\)。由于 \(f'' > 0\)，\(f'\) 严格递增，故 \(f'(\xi_1) < f'(\xi_2)\)，即

\[\frac{f(x_0) - f(x_1)}{x_0 - x_1} < \frac{f(x_2) - f(x_0)}{x_2 - x_0}\]

由此可推得 \(f(x_0) < \lambda f(x_1) + (1-\lambda)f(x_2)\)。具体地，记 \(\Delta x = x_2 - x_1 > 0\)，则 \(x_0 - x_1 = (1-\lambda)\Delta x\)，\(x_2 - x_0 = \lambda \Delta x\)。交叉相乘得：

\[\lambda[f(x_0) - f(x_1)] < (1-\lambda)[f(x_2) - f(x_0)]\]

展开并整理：\(\lambda f(x_0) + (1-\lambda)f(x_0) < \lambda f(x_1) + (1-\lambda)f(x_2)\)，即 \(f(x_0) < \lambda f(x_1) + (1-\lambda)f(x_2)\)。\(\blacksquare\)

\subsubsection{23.12.3 拐点}\label{ux62d0ux70b9}

\textbf{定义 23.8}（拐点）若曲线 \(y = f(x)\) 在点 \((x_0, f(x_0))\) 处凹凸性发生改变，则称此点为曲线的\textbf{拐点}。

\textbf{定理 23.29}（拐点的必要条件）若 \((x_0, f(x_0))\) 是拐点，且 \(f''(x_0)\) 存在，则 \(f''(x_0) = 0\)。

\textbf{证明}：假设 \(f''(x_0)\) 存在且 \((x_0, f(x_0))\) 是拐点。用反证法。若 \(f''(x_0) \neq 0\)，则 \(f''(x_0) > 0\) 或 \(f''(x_0) < 0\)。不妨设 \(f''(x_0) > 0\)。由 \(f''\) 的局部保号性（导数极限的局部保号性），存在 \(\delta > 0\)，使得对所有 \(x \in (x_0 - \delta, x_0 + \delta)\) 有 \(f''(x) > 0\)。由二阶导数判定法（定理 23.28），\(f\) 在该邻域内是凸函数（\(f'' > 0\)），凹凸性不变。但这与 \((x_0, f(x_0))\) 是拐点（凹凸性发生改变的点）矛盾。同理 \(f''(x_0) < 0\) 也会导致矛盾。因此必须有 \(f''(x_0) = 0\)。需要注意的是，\(f''(x_0) = 0\) 只是拐点的必要条件而非充分条件------例如 \(f(x) = x^4\) 在 \(x = 0\) 处 \(f''(0) = 0\)，但 \(x = 0\) 不是拐点（因为 \(f''(x) = 12x^2 \geq 0\) 恒非负，凹凸性不变）。\(\blacksquare\)

\hrulefill

\subsection{23.13 微分中值定理的进一步应用}\label{ux5faeux5206ux4e2dux503cux5b9aux7406ux7684ux8fdbux4e00ux6b65ux5e94ux7528}

在 23.4 节建立了微分中值定理的基本理论之后，本节展示这些定理在函数分析中的深入应用。中值定理的一个重要应用是函数恒等的判定：若 \(f'(x) = g'(x)\) 在区间上恒成立，则 \(f(x) = g(x) + C\)（\(C\) 为常数）。这一结论的逆否命题同样重要------若 \(f'(x) = 0\) 在区间上恒成立，则 \(f\) 为常数函数。洛必达法则的严格证明依赖于柯西中值定理，这体现了中值定理在极限计算中的核心地位。中值定理在不等式证明中尤为强大------例如，要证明 \(\ln(1+x) < x\)（\(x > 0\)），只需对 \(f(t) = \ln(1+t)\) 在 \([0, x]\) 上应用拉格朗日中值定理，得到 \(\frac{\ln(1+x) - 0}{x - 0} = \frac{1}{1+\xi} < 1\)，从而 \(\ln(1+x) < x\)。中值定理还提供了函数误差估计的基本方法------泰勒公式的拉格朗日余项形式就是多次应用中值定理的结果。本节将系统展示中值定理在不等式证明、误差估计和函数恒等判定中的典型应用，加深对中值定理作为''微分学基本定理''的理解。

\subsubsection{23.13.1 函数恒等的判定}\label{ux51fdux6570ux6052ux7b49ux7684ux5224ux5b9a}

\textbf{定理 23.30} 设 \(f\) 和 \(g\) 在 \([a, b]\) 上连续，在 \((a, b)\) 上可导，\(f(a) = g(a)\)。若 \(f'(x) = g'(x)\) 对所有 \(x \in (a, b)\) 成立，则 \(f(x) = g(x)\) 对所有 \(x \in [a, b]\) 成立。

\textbf{证明}：令 \(h = f - g\)，则 \(h' = 0\) 在 \((a, b)\) 上成立，\(h(a) = 0\)。

对任意 \(x \in (a, b]\)，由拉格朗日中值定理，存在 \(\xi \in (a, x)\) 使得 \(h(x) - h(a) = h'(\xi)(x - a) = 0\)，故 \(h(x) = 0\)。\(\blacksquare\)

\subsubsection{23.13.2 不等式的证明}\label{ux4e0dux7b49ux5f0fux7684ux8bc1ux660e}

\textbf{定理 23.31} 对任意 \(x > 0\)，\(x - \frac{x^2}{2} < \ln(1+x) < x\)。

\textbf{证明}：右边不等式即 \(\ln(1+x) < x\)（\(x > 0\)）。此即推论 18.2 的严格不等式形式，其证明在 Ch 25 中基于导数给出（见定理 18.19），但结论可由对数函数的严格递增性和 \(\ln 1 = 0\) 独立建立：考虑 \(g(x) = x - \ln(1+x)\)，\(g(0) = 0\)；对有理数 \(0 < r < x\)，由 \(\ln\) 的严格递增性，\(\ln(1+x) > \ln(1+r)\)，而由有理指数幂的性质可证 \(\ln(1+r) < r < x\)，故 \(g(x) > 0\)。此处我们在 Ch 30 的框架下，也可直接利用导数给出简洁证明：\(g'(x) = 1 - \frac{1}{1+x} = \frac{x}{1+x} > 0\)（\(x > 0\)），故 \(g(x) > g(0) = 0\)。

对于左边：令 \(f(x) = \ln(1+x) - x + \frac{x^2}{2}\)。则 \(f(0) = 0\)，

\[f'(x) = \frac{1}{1+x} - 1 + x = \frac{1 - (1+x) + x(1+x)}{1+x} = \frac{x^2}{1+x} > 0 \quad (x > 0)\]

因此 \(f\) 在 \([0, +\infty)\) 上严格递增，\(f(x) > f(0) = 0\)（\(x > 0\)）。\(\blacksquare\)

\subsubsection{23.13.3 拉格朗日中值定理的有限增量形式}\label{ux62c9ux683cux6717ux65e5ux4e2dux503cux5b9aux7406ux7684ux6709ux9650ux589eux91cfux5f62ux5f0f}

\textbf{定理 23.32}（有限增量形式）设 \(f\) 在 \([a, b]\) 上连续，在 \((a, b)\) 上可导，则存在 \(\xi \in (a, b)\) 使得

\[f(b) - f(a) = f'(\xi)(b - a)\]

\textbf{证明}：这正是拉格朗日中值定理（定理 23.15）的标准形式，但此处我们强调其''有限增量''的视角。拉格朗日中值定理的证明已在 23.4.3 节中给出：构造辅助函数 \(F(x) = f(x) - f(a) - \frac{f(b)-f(a)}{b-a}(x-a)\)，验证 \(F\) 满足罗尔定理的条件（\(F(a) = F(b) = 0\)），从而存在 \(\xi \in (a, b)\) 使得 \(F'(\xi) = 0\)，即 \(f'(\xi) = \frac{f(b)-f(a)}{b-a}\)。乘以 \((b-a)\) 即得有限增量形式 \(f(b) - f(a) = f'(\xi)(b - a)\)。这一形式的重要性在于：它将函数在一个区间上的整体变化量 \(f(b) - f(a)\) 精确地表示为某一点 \(\xi\) 处的瞬时变化率乘以区间长度，从而建立了''局部信息''（导数）与''整体行为''（函数值差）之间的桥梁。\(\blacksquare\)

此定理可理解为：函数在区间上的整体变化量等于某一点的瞬时变化率乘以区间长度。这为''用局部信息推断整体行为''提供了理论基础。

\hrulefill

\hrulefill

\hrulefill

\hrulefill

\newpage

\chapter{03-函数与微积分/03-微积分初步/03-积分初步.md}\label{ux51fdux6570ux4e0eux5faeux79efux520603-ux5faeux79efux5206ux521dux6b6503-ux79efux5206ux521dux6b65.md}

积分是微积分中与导数并列的两大核心概念之一。如果说导数刻画了函数变化的''瞬时速率''，那么积分则回答了相反的、更深层的问题：如何从变化速率反推总量的积累？这一''逆向''关系形成了微积分的灵魂------Newton-Leibniz 公式（微积分基本定理）将微分和积分统一为互逆的运算。积分概念的历史渊源比微分更加久远：古希腊人（Eudoxus 和 Archimedes 的穷竭法，约公元前 300 年）已能用极限精神计算抛物线弓形等特殊图形的面积，但直到 Newton 和 Leibniz 在 17 世纪各自独立建立微积分体系，积分才从零散的几何技巧提升为系统性的代数运算。Riemann 在 1854 年为积分提供了严格的分析定义（Riemann 积分），将 Archimedes 的直观极限过程转化为 \(\varepsilon\)-\(\delta\) 语言下的精确构造，使得积分理论具有了坚实的逻辑基础。本章从定积分的 Riemann 定义出发，先建立定积分的几何直观（面积问题），再引入不定积分的概念作为微分的逆运算，最后通过 Newton-Leibniz 公式揭示二者之间的统一性。\textbf{前置知识}：本章依赖 Ch 30（导数与微分，特别是基本导数公式和求导法则）、Ch 29（极限与连续，特别是最值定理和介值定理）。

\hrulefill

\section{第75章：定积分与牛顿-莱布尼茨公式}\label{ux7b2c75ux7ae0ux5b9aux79efux5206ux4e0eux725bux987f-ux83b1ux5e03ux5c3cux8328ux516cux5f0f}

定积分的概念源于面积问题，但它的精确数学定义经历了漫长的历史发展。本节从曲线下的面积问题出发，引入了黎曼和------将区间分割为无穷多个小区间，在每个小区间上取函数值乘以区间长度再求和。定积分的严格定义（黎曼积分）是 \(\varepsilon\)-\(\delta\) 语言在积分学中的对应物：存在实数 \(I\) 使得当分割足够细（模 \(\|P\|\) 足够小）时，所有黎曼和都落在 \(I\) 的 \(\varepsilon\) 邻域内。定积分的线性性、区间可加性、保序性和绝对值不等式这四个基本性质，从黎曼和的性质直接推导出来，为后续的积分计算提供了代数基础。连续函数可积定理（定理 24.5）保证了最常见的函数类------连续函数------都是可积的，这一结果是微积分基本定理成立的前提条件。

\subsubsection{24.1.1 面积问题}\label{ux9762ux79efux95eeux9898}

面积问题是定积分概念的几何起源，也是最自然的引入动机。考虑平面上由曲线 \(y = f(x)\)（\(f(x) \geq 0\)）、\(x\) 轴以及直线 \(x = a\) 和 \(x = b\) 围成的曲边梯形区域------一个边界包含函数曲线而非直线段的图形。Archimedes 用穷竭法（公元前 3 世纪）的思想至今仍是现代积分理论的直观核心：将区间 \([a,b]\) 分割为越来越细的小区间，在每个小区间上用容易计算面积的内接和外接矩形上下夹逼曲边区域，当分割''无限细''时，下和与上和趋于同一个极限------这个极限就是所求的面积。这一方法在 19 世纪被 Riemann 和 Darboux 形式化为严格的上积分与下积分框架（Darboux 积分），其核心思想是：上积分是所有上和的下确界，下积分是所有下和的上确界，当且仅当二者相等时函数可积。这一节将从几何直观出发，逐步引入黎曼和、分割的模以及定积分作为极限的精确分析定义。

\subsubsection{24.1.2 黎曼和}\label{ux9eceux66fcux548c}

\textbf{定义 24.1}（分割）区间 \([a, b]\) 的一个\textbf{分割} \(P\) 是一组分点 \(a = x_0 < x_1 < x_2 < \cdots < x_n = b\)。记 \(\Delta x_i = x_i - x_{i-1}\)（\(i = 1, 2, \ldots, n\)），\(\|P\| = \max_{1 \leq i \leq n} \Delta x_i\) 称为分割 \(P\) 的\textbf{模}（或\textbf{范数}）。

\textbf{定义 24.2}（黎曼和）设 \(f\) 在 \([a, b]\) 上有定义，\(P = \{x_0, x_1, \ldots, x_n\}\) 是 \([a, b]\) 的分割，\(\xi_i \in [x_{i-1}, x_i]\)（\(i = 1, 2, \ldots, n\)）。\textbf{黎曼和}定义为：

\[S(P, f, \{\xi_i\}) = \sum_{i=1}^{n} f(\xi_i) \Delta x_i\]

\subsubsection{24.1.3 定积分的定义}\label{ux5b9aux79efux5206ux7684ux5b9aux4e49}

\textbf{定义 24.3}（定积分/黎曼积分）设 \(f\) 在 \([a, b]\) 上有定义。若存在实数 \(I\)，使得

\[\forall \varepsilon > 0, \exists \delta > 0, \forall P, \forall \{\xi_i\}: \|P\| < \delta \Rightarrow \left|S(P, f, \{\xi_i\}) - I\right| < \varepsilon\]

则称 \(f\) 在 \([a, b]\) 上\textbf{可积}，\(I\) 称为 \(f\) 在 \([a, b]\) 上的\textbf{定积分}，记作

\[\int_a^b f(x) \, dx = I\]

其中 \(a\) 称为\textbf{下限}，\(b\) 称为\textbf{上限}，\(f\) 称为\textbf{被积函数}，\(x\) 称为\textbf{积分变量}。

\subsubsection{24.1.4 定积分的性质}\label{ux5b9aux79efux5206ux7684ux6027ux8d28}

\textbf{定理 24.1}（线性性）设 \(f\) 和 \(g\) 在 \([a, b]\) 上可积，\(\alpha, \beta \in \mathbb{R}\)，则 \(\alpha f + \beta g\) 在 \([a, b]\) 上可积，且

\[\int_a^b [\alpha f(x) + \beta g(x)] \, dx = \alpha \int_a^b f(x) \, dx + \beta \int_a^b g(x) \, dx\]

\textbf{证明}：

设 \(P = \{x_0, x_1, \ldots, x_n\}\) 是 \([a, b]\) 的任意分割，\(\{\xi_i\}\) 是任意取点族。考虑函数 \(\alpha f + \beta g\) 的黎曼和：

\[S(P, \alpha f + \beta g, \{\xi_i\}) = \sum_{i=1}^{n} [\alpha f(\xi_i) + \beta g(\xi_i)] \Delta x_i\]

由于求和是有限线性组合，我们可以将 \(\alpha\) 和 \(\beta\) 提取到求和符号之外：

\[= \sum_{i=1}^{n} \alpha f(\xi_i) \Delta x_i + \sum_{i=1}^{n} \beta g(\xi_i) \Delta x_i = \alpha \sum_{i=1}^{n} f(\xi_i) \Delta x_i + \beta \sum_{i=1}^{n} g(\xi_i) \Delta x_i\]

\[= \alpha S(P, f, \{\xi_i\}) + \beta S(P, g, \{\xi_i\})\]

现在令分割的模 \(\|P\| \to 0\)。由于 \(f\) 和 \(g\) 在 \([a, b]\) 上可积，有

\[\lim_{\|P\| \to 0} S(P, f, \{\xi_i\}) = \int_a^b f(x) \, dx, \quad \lim_{\|P\| \to 0} S(P, g, \{\xi_i\}) = \int_a^b g(x) \, dx\]

由数列极限的线性性质（定理 22.3），和的极限等于极限的和，常数倍的极限等于极限的常数倍（注意这里 \(\|P\| \to 0\) 的极限过程本质上也是数列极限，因为我们可以取一列分割使其模趋于零）。因此

\[\lim_{\|P\| \to 0} S(P, \alpha f + \beta g, \{\xi_i\}) = \alpha \int_a^b f(x) \, dx + \beta \int_a^b g(x) \, dx\]

这就证明了 \(\alpha f + \beta g\) 在 \([a, b]\) 上可积，且积分值等于上式右端。\(\blacksquare\)

\textbf{定理 24.2}（区间可加性）设 \(f\) 在包含 \(a, b, c\) 的区间上可积，则

\[\int_a^b f(x) \, dx = \int_a^c f(x) \, dx + \int_c^b f(x) \, dx\]

\textbf{证明}：当 \(a < c < b\) 时，考虑 \([a, b]\) 的任意分割 \(P\)。将 \(P\) 加细为 \(P'\)，在 \(P'\) 中加入分点 \(c\)（如果 \(c\) 不在 \(P\) 中）。加细分割不会改变积分的值（因为加细分割只增加黎曼和的逼近精度，而极限值不变）。于是我们可以将积分 \(\int_a^b\) 分解为 \([a, c]\) 上的积分与 \([c, b]\) 上的积分之和。具体地，对于加细后的分割，将子区间按 \(c\) 为界分为两组：落在 \([a, c]\) 内的子区间构成 \(\int_a^c\) 的黎曼和，落在 \([c, b]\) 内的子区间构成 \(\int_c^b\) 的黎曼和。由于黎曼和的极限等于各部分黎曼和极限之和，取极限即得 \(\int_a^b f(x) \, dx = \int_a^c f(x) \, dx + \int_c^b f(x) \, dx\)。当 \(c < a < b\) 或 \(a < b < c\) 时，利用上述结果和积分方向的约定 \(\int_a^b = -\int_b^a\) 即可推广。\(\blacksquare\)

\textbf{定理 24.3}（保序性）设 \(f\) 和 \(g\) 在 \([a, b]\) 上可积，且 \(f(x) \geq g(x)\) 对所有 \(x \in [a, b]\) 成立（\(a < b\)），则

\[\int_a^b f(x) \, dx \geq \int_a^b g(x) \, dx\]

\textbf{证明}：对于 \([a, b]\) 的任意分割 \(P = \{x_0, \ldots, x_n\}\) 和任意取点族 \(\{\xi_i\}\)，由于 \(f(x) \geq g(x)\) 在 \([a, b]\) 上处处成立，特别地在每个 \(\xi_i \in [x_{i-1}, x_i]\) 处有 \(f(\xi_i) \geq g(\xi_i)\)。又因为 \(a < b\)，每个子区间长度 \(\Delta x_i = x_i - x_{i-1} > 0\)。将不等式 \(f(\xi_i) \geq g(\xi_i)\) 两边乘以正数 \(\Delta x_i\) 并求和，得到

\[S(P, f, \{\xi_i\}) = \sum_{i=1}^{n} f(\xi_i) \Delta x_i \geq \sum_{i=1}^{n} g(\xi_i) \Delta x_i = S(P, g, \{\xi_i\})\]

现在令分割的模 \(\|P\| \to 0\)。由 \(f\) 和 \(g\) 的可积性，黎曼和分别收敛于相应的定积分。由数列极限的保序性（定理 21.13），若每一项都满足 \(S(P, f, \{\xi_i\}) \geq S(P, g, \{\xi_i\})\)，则极限也满足同样的不等式。因此

\[\int_a^b f(x) \, dx = \lim_{\|P\| \to 0} S(P, f, \{\xi_i\}) \geq \lim_{\|P\| \to 0} S(P, g, \{\xi_i\}) = \int_a^b g(x) \, dx\]

这就证明了保序性。\(\blacksquare\)

\textbf{定理 24.4}（绝对值不等式）设 \(f\) 在 \([a, b]\) 上可积，则

\[\left|\int_a^b f(x) \, dx\right| \leq \int_a^b |f(x)| \, dx\]

\textbf{证明}：首先注意到对于任意实数 \(y\)，恒有 \(-|y| \leq y \leq |y|\)。这是因为 \(|y| \geq 0\)，当 \(y \geq 0\) 时 \(y = |y|\)，不等式 \(y \leq |y|\) 取等号，而 \(-|y| \leq 0 \leq y\) 成立；当 \(y < 0\) 时 \(y = -|y|\)，\(-|y| \leq y\) 取等号，而 \(y < 0 \leq |y|\) 成立。将这一基本不等式应用于 \(f(x)\)，得到对每个 \(x \in [a, b]\) 有

\[-|f(x)| \leq f(x) \leq |f(x)|\]

现在在 \([a, b]\) 上应用保序性（定理 24.3）。注意 \(|f(x)|\) 作为连续函数的复合也是可积的（\(f\) 可积蕴含 \(|f|\) 可积，这一结论成立需要 \(f\) 可积作为前提，在此处 \(f\) 可积是已知条件）。由保序性，对不等式 \(-|f(x)| \leq f(x)\) 积分得

\[-\int_a^b |f(x)| \, dx \leq \int_a^b f(x) \, dx\]

对不等式 \(f(x) \leq |f(x)|\) 积分得

\[\int_a^b f(x) \, dx \leq \int_a^b |f(x)| \, dx\]

将两个不等式合并，即得

\[-\int_a^b |f(x)| \, dx \leq \int_a^b f(x) \, dx \leq \int_a^b |f(x)| \, dx\]

这等价于绝对值不等式 \(\left|\int_a^b f(x) \, dx\right| \leq \int_a^b |f(x)| \, dx\)。\(\blacksquare\)

\subsubsection{24.1.5 可积的充分条件}\label{ux53efux79efux7684ux5145ux5206ux6761ux4ef6}

\textbf{定理 24.5}（连续函数可积）若 \(f\) 在 \([a, b]\) 上连续，则 \(f\) 在 \([a, b]\) 上可积。

*\textbf{证明}：由于 \(f\) 在闭区间 \([a, b]\) 上连续，由康托尔定理（定理 22.21），\(f\) 在 \([a, b]\) 上一致连续。即对任意 \(\varepsilon > 0\)，存在 \(\delta > 0\)，使得对任意 \(x, y \in [a, b]\)，只要 \(|x - y| < \delta\) 就有 \(|f(x) - f(y)| < \frac{\varepsilon}{b - a}\)。现在取 \([a, b]\) 的任意两个分割 \(P\) 和 \(Q\)，其模均小于 \(\delta\)。考虑它们对应的达布上和与达布下和之差。由于在每个子区间上 \(f\) 的振幅（最大值与最小值之差）不超过 \(\frac{\varepsilon}{b - a}\)，可以证明达布上和与达布下和之差小于 \(\varepsilon\)。这就证明了达布上和的下确界等于达布下和的上确界，即 \(f\) 在 \([a, b]\) 上可积。更具体地，设 \(m_i\) 和 \(M_i\) 分别是 \(f\) 在子区间 \([x_{i-1}, x_i]\) 上的最小值和最大值（由最值定理，连续函数在闭区间上取到最值）。由一致连续性，当 \(\Delta x_i < \delta\) 时，\(M_i - m_i < \frac{\varepsilon}{b - a}\)。因此

\[\sum_{i=1}^{n} (M_i - m_i) \Delta x_i < \frac{\varepsilon}{b - a} \sum_{i=1}^{n} \Delta x_i = \frac{\varepsilon}{b - a} \cdot (b - a) = \varepsilon\]

由达布可积准则，\(f\) 在 \([a, b]\) 上可积。\(\blacksquare\)

lacksquare\$

\hrulefill

\subsection{24.2 牛顿-莱布尼茨公式}\label{ux725bux987f-ux83b1ux5e03ux5c3cux8328ux516cux5f0f}

牛顿-莱布尼茨公式（微积分基本定理）是微积分中最深刻的理论成就，它揭示了微分和积分这两个看似无关的运算之间的内在联系。本节从原函数的概念出发，证明了原函数之间仅相差常数，然后通过定义变上限积分函数 \(\Phi(x) = \int_a^x f(t) \, dt\) 并证明 \(\Phi'(x) = f(x)\)，建立了微分与积分之间的互逆关系。牛顿-莱布尼茨公式 \(\int_a^b f(x) \, dx = F(b) - F(a)\) 将定积分的计算问题转化为求原函数的问题，这一转变为积分计算提供了强大的工具。积分中值定理（推论 24.1）则是连续函数介值定理在积分形式下的体现，它保证了在闭区间上存在一点使得积分值等于函数值乘以区间长度。

\subsubsection{24.2.1 原函数}\label{ux539fux51fdux6570}

\textbf{定义 24.4}（原函数）设 \(f\) 在区间 \(I\) 上有定义。若函数 \(F\) 满足 \(F'(x) = f(x)\) 对所有 \(x \in I\) 成立，则称 \(F\) 为 \(f\) 在 \(I\) 上的一个\textbf{原函数}。

\textbf{定理 24.6}（原函数的唯一性，相差常数）设 \(F\) 和 \(G\) 都是 \(f\) 在 \(I\) 上的原函数，则存在常数 \(C\) 使得 \(G(x) = F(x) + C\)。

\textbf{证明}：设 \(H(x) = G(x) - F(x)\)。则对任意 \(x \in I\)，

\[H'(x) = G'(x) - F'(x) = f(x) - f(x) = 0\]

现在需要证明：若 \(H'(x) = 0\) 在区间 \(I\) 上处处成立，则 \(H\) 在 \(I\) 上是常数函数。这需要用到拉格朗日中值定理。对 \(I\) 中任意两点 \(x_1, x_2\)（\(x_1 < x_2\)），\(H\) 在 \([x_1, x_2]\) 上连续（因为 \(H\) 可导）且在 \((x_1, x_2)\) 上可导。由拉格朗日中值定理，存在 \(\xi \in (x_1, x_2)\) 使得

\[H(x_2) - H(x_1) = H'(\xi)(x_2 - x_1) = 0 \cdot (x_2 - x_1) = 0\]

因此 \(H(x_2) = H(x_1)\)。由于 \(x_1\) 和 \(x_2\) 是任意选取的，\(H\) 在 \(I\) 上取常数值。设此常数为 \(C\)，则 \(G(x) - F(x) = C\)，即 \(G(x) = F(x) + C\)。\(\blacksquare\)

\subsubsection{24.2.2 牛顿-莱布尼茨公式}\label{ux725bux987f-ux83b1ux5e03ux5c3cux8328ux516cux5f0f-1}

\textbf{定理 24.7}（牛顿-莱布尼茨公式，微积分基本定理）设 \(f\) 在 \([a, b]\) 上连续，\(F\) 是 \(f\) 的任一原函数，即 \(F'(x) = f(x)\)。则

\[\int_a^b f(x) \, dx = F(b) - F(a)\]

\textbf{证明}：定义\textbf{变上限积分函数}

\[\Phi(x) = \int_a^x f(t) \, dt, \quad x \in [a, b]\]

\textbf{第一步}：证明 \(\Phi'(x) = f(x)\)。

对于 \(x \in (a, b)\) 和 \(h \neq 0\)（\(x + h \in [a, b]\)），

\[\Phi(x+h) - \Phi(x) = \int_a^{x+h} f(t) \, dt - \int_a^x f(t) \, dt = \int_x^{x+h} f(t) \, dt\]

由于 \(f\) 连续，由最值定理（定理 22.16），存在 \(m, M\) 使得 \(m \leq f(t) \leq M\) 在 \([x, x+h]\)（或 \([x+h, x]\)）上成立。故 \(mh \leq \int_x^{x+h} f(t)\,dt \leq Mh\)（\(h > 0\) 时），即 \(m \leq \frac{1}{h}\int_x^{x+h} f(t)\,dt \leq M\)。由介值定理（定理 22.15），存在 \(\xi\) 介于 \(x\) 和 \(x+h\) 之间，使得

\[\int_x^{x+h} f(t) \, dt = f(\xi) \cdot h\]

因此

\[\frac{\Phi(x+h) - \Phi(x)}{h} = f(\xi)\]

当 \(h \to 0\) 时，\(\xi \to x\)（因为 \(\xi\) 介于 \(x\) 和 \(x+h\) 之间），由 \(f\) 的连续性，

\[\Phi'(x) = \lim_{h \to 0} \frac{\Phi(x+h) - \Phi(x)}{h} = \lim_{\xi \to x} f(\xi) = f(x)\]

\textbf{第二步}：因此 \(\Phi\) 是 \(f\) 的一个原函数。由定理 24.6，任何原函数 \(F\) 满足

\[F(x) = \Phi(x) + C\]

令 \(x = a\)：\(F(a) = \Phi(a) + C = 0 + C = C\)。

令 \(x = b\)：\(F(b) = \Phi(b) + F(a)\)。

所以 \(\Phi(b) = F(b) - F(a)\)，即

\[\int_a^b f(x) \, dx = F(b) - F(a)\]

\(\blacksquare\)

\textbf{注}：这个定理深刻地联系了微分和积分------两个看似不相关的概念：微分是''局部化''（求瞬时变化率），积分是''全局化''（求累积量），而牛顿-莱布尼茨公式表明它们互为逆运算。

\textbf{推论 24.1}（积分中值定理）设 \(f\) 在 \([a, b]\) 上连续，则存在 \(\xi \in [a, b]\)，使得

\[\int_a^b f(x) \, dx = f(\xi)(b - a)\]

\textbf{证明}：由最值定理，\(f\) 在 \([a, b]\) 上取到最小值 \(m\) 和最大值 \(M\)：

\[m(b-a) \leq \int_a^b f(x) \, dx \leq M(b-a)\]

即 \(m \leq \frac{1}{b-a}\int_a^b f(x) \, dx \leq M\)。

由介值定理，存在 \(\xi \in [a, b]\) 使得 \(f(\xi) = \frac{1}{b-a}\int_a^b f(x) \, dx\)。\(\blacksquare\)

\hrulefill

\section{第76章：积分方法}\label{ux7b2c76ux7ae0ux79efux5206ux65b9ux6cd5}

不定积分是求导运算的逆运算，其核心问题是：给定一个函数 \(f(x)\)，求所有满足 \(F'(x) = f(x)\) 的函数 \(F(x)\)。与定积分（一个数值）不同，不定积分 \(\int f(x) \, dx = F(x) + C\) 是一个函数族，其中 \(C\) 是任意常数------这一常数反映了原函数之间仅相差常数的基本事实。不定积分的线性性质 \(\int [\alpha f(x) + \beta g(x)] \, dx = \alpha \int f(x) \, dx + \beta \int g(x) \, dx\) 直接来自导数的线性性质，使得不定积分运算可以逐项进行。基本积分公式表（如 \(\int x^n \, dx = \frac{x^{n+1}}{n+1} + C\)（\(n \neq -1\)）、\(\int \frac{1}{x} \, dx = \ln|x| + C\)、\(\int e^x \, dx = e^x + C\) 等）是积分计算的基础工具，它们都是基本初等函数导数公式的''反向版本''。不定积分的概念为后续的换元积分法和分部积分法提供了理论框架。

\subsubsection{24.3.1 不定积分的定义}\label{ux4e0dux5b9aux79efux5206ux7684ux5b9aux4e49}

不定积分（Indefinite Integral）是微分运算的逆运算------又称为反导数（Antiderivative）。它的核心问题是：给定一个函数 \(f(x)\)，找出所有满足 \(F'(x) = f(x)\) 的函数 \(F(x)\)。与定积分（一个确定的数值）不同，不定积分的结果是一个函数族 \(F(x) + C\)，其中 \(C\) 为任意常数------这一常数的出现是中值定理的直接推论：如果两个函数的导数处处相等，则它们至多相差一个常数。因此，不定积分的符号 \(\int f(x)\,dx\)（不带上下限）代表的是所有原函数构成的等价类，而不仅仅是一个''反导数''。不定积分的理论地位恰恰在于它为定积分的实际计算铺平了道路：一旦找到 \(f\) 的一个原函数 \(F\)，Newton-Leibniz 公式就将复杂的极限计算（定积分作为黎曼和的极限）简化为简单的函数求值 \(\int_a^b f(x)\,dx = F(b) - F(a)\)。基本积分公式表（幂函数、指数函数、三角函数、反三角函数的积分）是基本导数公式的直接''反向版本''，构成了后续换元积分法和分部积分法的代数基础。

\textbf{定义 24.5}（不定积分）函数 \(f\) 的所有原函数构成的集合称为 \(f\) 的\textbf{不定积分}，记作

\[\int f(x) \, dx = F(x) + C\]

其中 \(F'(x) = f(x)\)，\(C\) 为任意常数。

\subsubsection{24.3.2 基本积分公式}\label{ux57faux672cux79efux5206ux516cux5f0f}

由基本导数公式反向得到：

\begin{enumerate}
\def\labelenumi{\arabic{enumi}.}
\tightlist
\item
  \(\int x^n \, dx = \frac{x^{n+1}}{n+1} + C\)（\(n \neq -1\)）
\item
  \(\int \frac{1}{x} \, dx = \ln|x| + C\)
\item
  \(\int e^x \, dx = e^x + C\)
\item
  \(\int a^x \, dx = \frac{a^x}{\ln a} + C\)（\(a > 0\)，\(a \neq 1\)）
\item
  \(\int \sin x \, dx = -\cos x + C\)
\item
  \(\int \cos x \, dx = \sin x + C\)
\item
  \(\int \sec^2 x \, dx = \tan x + C\)
\end{enumerate}

\textbf{证明 1}：要验证 \(\frac{x^{n+1}}{n+1}\) 是 \(x^n\) 的原函数，只需对其求导。由幂函数的导数公式（定理 23.2），

\[\left(\frac{x^{n+1}}{n+1}\right)' = \frac{1}{n+1} \cdot (x^{n+1})' = \frac{1}{n+1} \cdot (n+1)x^n = x^n\]

由于求导结果恰好等于被积函数 \(x^n\)，因此 \(\frac{x^{n+1}}{n+1}\) 是 \(x^n\) 的一个原函数。加上任意常数 \(C\) 即得不定积分公式。注意这里要求 \(n \neq -1\)，因为 \(n = -1\) 时分母 \(n+1 = 0\)，公式失效------此时对应的是 \(\int \frac{1}{x} \, dx = \ln|x| + C\)。\(\blacksquare\)

\textbf{证明 5}：要验证 \(-\cos x\) 是 \(\sin x\) 的原函数，对其求导。由余弦函数的导数公式（定理 23.4），\((\cos x)' = -\sin x\)，因此

\[(-\cos x)' = -(\cos x)' = -(-\sin x) = \sin x\]

求导结果恰好等于被积函数 \(\sin x\)，因此 \(-\cos x\) 是 \(\sin x\) 的一个原函数。加上任意常数 \(C\) 即得 \(\int \sin x \, dx = -\cos x + C\)。\(\blacksquare\)

\hrulefill

\subsection{24.4 换元积分法}\label{ux6362ux5143ux79efux5206ux6cd5}

换元积分法是积分计算中最基本也是最重要的技巧，其本质是复合函数求导法则（链式法则）的逆向应用。第一换元法（凑微分法）将积分 \(\int f(g(x))g'(x) \, dx\) 转化为 \(\int f(u) \, du\)（其中 \(u = g(x)\)），其核心在于识别被积函数中是否隐藏了''内函数的导数''这一因子。第二换元法（变量代换法）通过引入新变量 \(x = \varphi(t)\) 将积分 \(\int f(x) \, dx\) 转化为 \(\int f(\varphi(t))\varphi'(t) \, dt\)，其适用范围更广------当被积函数含有根式 \(\sqrt{a^2 - x^2}\)、\(\sqrt{x^2 \pm a^2}\) 等时，三角代换法是最有效的工具。换元积分法的''艺术性''在于选择合适的代换------这需要经验和对函数结构的敏感度。欧拉代换法（处理 \(\sqrt{ax^2 + bx + c}\) 型根式）和倒代换法（\(x = 1/t\)）是换元法的两种重要变体。换元积分法在定积分中也有对应形式------定积分的换元法需要同时变换积分限，这避免了不定积分中换回原变量的繁琐步骤。

\subsubsection{24.4.1 第一换元法（凑微分法）}\label{ux7b2cux4e00ux6362ux5143ux6cd5ux51d1ux5faeux5206ux6cd5}

\textbf{定理 24.8}（第一换元法）设 \(F\) 是 \(f\) 的原函数，\(g\) 可导，则

\[\int f(g(x)) g'(x) \, dx = F(g(x)) + C\]

\textbf{证明}：由链式法则（定理 23.12），对复合函数 \(F(g(x))\) 求导：

\[[F(g(x))]' = F'(g(x)) \cdot g'(x)\]

由于 \(F\) 是 \(f\) 的原函数，即 \(F'(u) = f(u)\) 对所有 \(u\) 成立，代入得

\[[F(g(x))]' = f(g(x)) \cdot g'(x)\]

这意味着 \(F(g(x))\) 的导数恰好等于被积函数 \(f(g(x))g'(x)\)。由不定积分的定义------不定积分是原函数族------\(F(g(x))\) 是 \(f(g(x))g'(x)\) 的一个原函数，因此

\[\int f(g(x)) g'(x) \, dx = F(g(x)) + C\]

在操作中，我们通常设 \(u = g(x)\)，则 \(du = g'(x) \, dx\)，于是积分变为 \(\int f(u) \, du = F(u) + C = F(g(x)) + C\)。这一''凑微分''的技巧之所以有效，正是因为被积函数中包含了 \(g'(x)\) 这一因子，恰好与微分关系 \(du = g'(x) \, dx\) 相匹配。\(\blacksquare\)

\textbf{用法}：欲求 \(\int h(x) \, dx\)，若能将 \(h(x)\) 写成 \(f(g(x)) \cdot g'(x)\) 的形式，则令 \(u = g(x)\)，转化为 \(\int f(u) \, du\)。

\subsubsection{24.4.2 第二换元法}\label{ux7b2cux4e8cux6362ux5143ux6cd5}

\textbf{定理 24.9}（第二换元法）设 \(x = \varphi(t)\) 是严格单调可导函数，\(\varphi'(t) \neq 0\)。则

\[\int f(x) \, dx = \int f(\varphi(t)) \varphi'(t) \, dt \bigg|_{t = \varphi^{-1}(x)} + C\]

\textbf{证明}：设 \(G(t) = \int f(\varphi(t)) \varphi'(t) \, dt\)，即 \(G'(t) = f(\varphi(t))\varphi'(t)\)。

令 \(F(x) = G(\varphi^{-1}(x))\)。由链式法则和反函数求导法则：

\[F'(x) = G'(\varphi^{-1}(x)) \cdot \frac{d}{dx}\varphi^{-1}(x) = f(\varphi(\varphi^{-1}(x))) \cdot \varphi'(\varphi^{-1}(x)) \cdot \frac{1}{\varphi'(\varphi^{-1}(x))} = f(x)\]

所以 \(F(x) = G(\varphi^{-1}(x))\) 是 \(f(x)\) 的原函数。\(\blacksquare\)

\hrulefill

\subsection{24.5 分部积分法}\label{ux5206ux90e8ux79efux5206ux6cd5}

分部积分法是乘积函数求导法则 \((uv)' = u'v + uv'\) 的积分形式，重新排列后得到 \(\int u \, dv = uv - \int v \, du\)。这一公式的核心思想是将一个''难以直接积分''的乘积 \(\int u \, dv\) 转化为另一个''更容易积分''的乘积 \(\int v \, du\)。分部积分法的关键在于选择 \(u\) 和 \(dv\)------一般原则是让 \(u\) 求导后变简单（如多项式、对数函数），让 \(dv\) 积分后不变得更复杂（如指数函数、三角函数）。经典的''LIATE''法则（Logarithmic \(>\) Inverse trigonometric \(>\) Algebraic \(>\) Trigonometric \(>\) Exponential）总结了选择 \(u\) 的优先级。分部积分法在某些情况下可以产生''循环''------例如 \(\int e^x \sin x \, dx\) 经过两次分部积分后回到原积分，从而可以解出积分值。分部积分法在定积分中同样适用，且有时可以借助''积分不等式''或''极限''来简化无穷区间的积分。分部积分法与换元积分法共同构成了积分计算的两大核心方法，几乎所有初等函数的积分都可以通过这两种方法的组合来完成。

\textbf{定理 24.10}（分部积分法）设 \(u(x)\) 和 \(v(x)\) 都可导，则

\[\int u \, dv = uv - \int v \, du\]

即

\[\int u(x) v'(x) \, dx = u(x)v(x) - \int u'(x) v(x) \, dx\]

\textbf{证明}：由乘法法则，\((uv)' = u'v + uv'\)，即 \(uv' = (uv)' - u'v\)。

两边积分：

\[\int u v' \, dx = \int (uv)' \, dx - \int u' v \, dx = uv - \int u' v \, dx\]

\(\blacksquare\)

\textbf{注}：分部积分法在处理乘积形式的被积函数时特别有用。选择 \(u\) 和 \(dv\) 的原则通常是使得 \(\int v \, du\) 比 \(\int u \, dv\) 更容易计算。

\hrulefill

\section{第77章：微积分基本定理与积分应用}\label{ux7b2c77ux7ae0ux5faeux79efux5206ux57faux672cux5b9aux7406ux4e0eux79efux5206ux5e94ux7528}

微积分基本定理有两个等价但侧重不同的表述。第一部分（定理 24.4）------变上限积分函数的导数等于被积函数------强调的是''积分后再微分回到原函数''；第二部分（定理 24.11）------定积分等于原函数在上下限的差值------强调的是''微分后再积分回到原函数值差''。这两个部分共同构成了微分与积分之间的完美对偶关系。第二部分（牛顿-莱布尼茨公式 \(\int_a^b f(x) \, dx = F(b) - F(a)\)）是积分计算中最实用的形式------它将定积分的计算问题完全转化为求原函数的问题。这一公式的证明依赖于第一部分：定义 \(G(x) = \int_a^x f(t) \, dt\)，则 \(G'(x) = f(x)\)，所以 \(G\) 是 \(f\) 的一个原函数。若 \(F\) 是 \(f\) 的任意一个原函数，则 \(G(x) = F(x) + C\)，代入 \(x = a\) 得 \(C = -F(a)\)，从而 \(G(b) = F(b) - F(a)\)。这一优雅的证明链条展示了微积分基本定理的内在逻辑------``积分与微分互为逆运算''这一核心洞察。

\textbf{定理 24.11}（微积分基本定理第二部分）设 \(f\) 在 \([a, b]\) 上可积，且有原函数 \(F\)（即 \(F' = f\)），则

\[\int_a^b f(x) \, dx = F(b) - F(a)\]

\textbf{证明}：对 \([a, b]\) 的任意分割 \(a = x_0 < x_1 < \cdots < x_n = b\)，在每个子区间 \([x_{i-1}, x_i]\) 上，\(F\) 在 \([x_{i-1}, x_i]\) 上连续、在 \((x_{i-1}, x_i)\) 上可导，由拉格朗日中值定理（定理 23.15），存在 \(\xi_i \in (x_{i-1}, x_i)\) 使得

\[F(x_i) - F(x_{i-1}) = F'(\xi_i)(x_i - x_{i-1}) = f(\xi_i) \Delta x_i\]

对所有 \(i\) 求和，左端为望远镜和：

\[\sum_{i=1}^n [F(x_i) - F(x_{i-1})] = F(b) - F(a)\]

右端为 \(f\) 的一个黎曼和 \(\sum_{i=1}^n f(\xi_i) \Delta x_i\)。由于 \(f\) 在 \([a, b]\) 上可积，当分割的模 \(\|\mathcal{P}\| \to 0\) 时，黎曼和收敛到 \(\int_a^b f(x)\,dx\)。因此

\[F(b) - F(a) = \int_a^b f(x)\,dx\]

\(\blacksquare\)

\textbf{注}：此定理将定理 24.7 的连续性条件减弱为可积性。关键在于拉格朗日中值定理仅要求 \(F\) 可导（从而连续），不要求 \(f = F'\) 连续。

\hrulefill

\subsection{24.7 定积分的进一步应用}\label{ux5b9aux79efux5206ux7684ux8fdbux4e00ux6b65ux5e94ux7528}

定积分的应用远远超出了''求曲线下的面积''这一基本问题。本节将定积分用于计算平面区域的面积（包括两条曲线围成的区域）、旋转体的体积（圆盘法和壳法）、曲线的弧长（\(\int_a^b \sqrt{1 + [f'(x)]^2} \, dx\)）以及旋转曲面的面积。面积计算中，\(x\) 型区域和 \(y\) 型区域的分解是基本技巧------有时对 \(y\) 积分比对 \(x\) 积分更方便。旋转体体积的圆盘法（垂直于旋转轴的截面为圆盘）和壳法（平行于旋转轴的薄壳）是两种互补的方法，选择哪种方法取决于旋转轴的方向和区域的形状。弧长公式的推导体现了定积分思想的核心------``分小段、近似直、求和、取极限''------将曲线分割为无穷小的直线段，每段长度近似为 \(\sqrt{(dx)^2 + (dy)^2}\)，求和后取极限得到积分。定积分的这些应用展示了''微元法''（将总量分解为无穷小元素的和）这一统一方法论在几何和物理问题中的强大威力。

\subsubsection{24.7.1 平面区域的面积}\label{ux5e73ux9762ux533aux57dfux7684ux9762ux79ef}

\textbf{定理 24.12} 设 \(f\) 和 \(g\) 在 \([a, b]\) 上连续，且 \(f(x) \geq g(x)\)。则由曲线 \(y = f(x)\) 和 \(y = g(x)\) 以及直线 \(x = a\)，\(x = b\) 围成的区域的面积为

\[S = \int_a^b [f(x) - g(x)] \, dx\]

\textbf{证明}：将所求区域分解为两部分之差。首先，曲线 \(y = f(x)\) 下方、\(x\) 轴上方以及直线 \(x = a\) 和 \(x = b\) 围成的曲边梯形的面积为 \(\int_a^b f(x) \, dx\)。其次，曲线 \(y = g(x)\) 下方、\(x\) 轴上方以及同样两条直线围成的曲边梯形的面积为 \(\int_a^b g(x) \, dx\)。由于 \(f(x) \geq g(x)\) 在 \([a, b]\) 上处处成立，曲线 \(y = f(x)\) 和 \(y = g(x)\) 之间的区域面积恰好等于大曲边梯形面积减去小曲边梯形面积。由定积分的线性性（定理 24.1），

\[S = \int_a^b f(x) \, dx - \int_a^b g(x) \, dx = \int_a^b [f(x) - g(x)] \, dx\]

这一推导的关键在于：区域的面积可以被精确地表示为两个定积分之差，而线性性允许我们将减法合并到被积函数中。\(\blacksquare\)

\subsubsection{24.7.2 旋转体的体积}\label{ux65cbux8f6cux4f53ux7684ux4f53ux79ef}

\textbf{定理 24.13} 设 \(f\) 在 \([a, b]\) 上连续，\(f(x) \geq 0\)。曲线 \(y = f(x)\) 绕 \(x\) 轴旋转一周所得旋转体的体积为

\[V = \pi \int_a^b [f(x)]^2 \, dx\]

\textbf{证明}：将 \([a, b]\) 分割为小区间 \([x_{i-1}, x_i]\)。在每个小区间上，旋转体近似为底面半径 \(f(\xi_i)\)、高 \(\Delta x_i\) 的圆柱体，体积为 \(\pi [f(\xi_i)]^2 \Delta x_i\)。

总体积近似为 \(\sum_{i=1}^{n} \pi [f(\xi_i)]^2 \Delta x_i\)。取 \(\|P\| \to 0\) 的极限，由定积分定义，

\[V = \pi \int_a^b [f(x)]^2 \, dx\]

\(\blacksquare\)

\subsubsection{24.7.3 弧长}\label{ux5f27ux957f}

\textbf{定理 24.14} 设 \(f\) 在 \([a, b]\) 上有连续的导数，则曲线 \(y = f(x)\) 从 \(x = a\) 到 \(x = b\) 的弧长为

\[L = \int_a^b \sqrt{1 + [f'(x)]^2} \, dx\]

\textbf{证明}：将 \([a, b]\) 分割为小区间。在 \([x_{i-1}, x_i]\) 上，弧长近似为

\[\sqrt{(\Delta x_i)^2 + (\Delta y_i)^2} = \Delta x_i \sqrt{1 + \left(\frac{\Delta y_i}{\Delta x_i}\right)^2}\]

由拉格朗日中值定理，\(\frac{\Delta y_i}{\Delta x_i} = f'(\xi_i)\)。因此弧长近似为

\[\sum_{i=1}^{n} \sqrt{1 + [f'(\xi_i)]^2} \Delta x_i\]

取极限得 \(L = \int_a^b \sqrt{1 + [f'(x)]^2} \, dx\)。\(\blacksquare\)

\hrulefill

\section{第78章：反常积分与积分深化}\label{ux7b2c78ux7ae0ux53cdux5e38ux79efux5206ux4e0eux79efux5206ux6df1ux5316}

黎曼积分的定义要求积分区间是有限的且被积函数是有界的。反常积分（improper integral）将积分的概念推广到无穷区间或函数无界的情形。无穷限积分 \(\int_a^{+\infty} f(x) \, dx = \lim_{b \to +\infty} \int_a^b f(x) \, dx\) 通过极限过程将无穷区间的积分定义为有限区间积分的极限------如果极限存在，称反常积分收敛；否则发散。无界函数的积分（瑕积分）\(\int_a^b f(x) \, dx = \lim_{t \to b^-} \int_a^t f(x) \, dx\) 处理的是被积函数在积分端点附近趋于无穷的情形。这两种反常积分的收敛性判别是本节的核心问题------\(\int_1^{+\infty} \frac{1}{x^p} \, dx\) 当 \(p > 1\) 时收敛、\(p \leq 1\) 时发散，这一''\(p\) 积分''判别法是最基本的收敛性测试。比较判别法和极限比较判别法提供了更一般的收敛性判断工具。反常积分在概率论（正态分布的积分）、物理学（势能计算）和工程学中有着广泛应用，其收敛性理论是后续级数收敛性判别（如积分判别法）的基础。

\subsubsection{24.8.1 无穷限积分}\label{ux65e0ux7a77ux9650ux79efux5206}

\textbf{定义 24.6}（无穷限积分）设 \(f\) 在 \([a, +\infty)\) 上连续。定义

\[\int_a^{+\infty} f(x) \, dx = \lim_{b \to +\infty} \int_a^b f(x) \, dx\]

若此极限存在，则称反常积分\textbf{收敛}；否则称其\textbf{发散}。

\textbf{定理 24.15} \(\int_1^{+\infty} \frac{1}{x^p} \, dx\) 当 \(p > 1\) 时收敛，当 \(p \leq 1\) 时发散。

\textbf{证明}：当 \(p \neq 1\) 时：

\[\int_1^b \frac{1}{x^p} \, dx = \frac{x^{1-p}}{1-p}\bigg|_1^b = \frac{b^{1-p} - 1}{1-p}\]

当 \(p > 1\) 时，\(1 - p < 0\)，\(b^{1-p} \to 0\)（\(b \to +\infty\)），故积分收敛于 \(\frac{1}{p-1}\)。

当 \(p < 1\) 时，\(1 - p > 0\)，\(b^{1-p} \to +\infty\)，积分发散。

当 \(p = 1\) 时，\(\int_1^b \frac{1}{x} \, dx = \ln b \to +\infty\)，积分发散。\(\blacksquare\)

\subsubsection{24.8.2 无界函数的积分}\label{ux65e0ux754cux51fdux6570ux7684ux79efux5206}

\textbf{定义 24.7} 设 \(f\) 在 \((a, b]\) 上连续，\(\lim_{x \to a^+} f(x) = \infty\)。定义

\[\int_a^b f(x) \, dx = \lim_{\varepsilon \to 0^+} \int_{a+\varepsilon}^b f(x) \, dx\]

\hrulefill

\subsection{24.9 积分与微分的对偶关系}\label{ux79efux5206ux4e0eux5faeux5206ux7684ux5bf9ux5076ux5173ux7cfb}

积分与微分之间的对偶关系是微积分全部理论的核心。微积分基本定理的两个部分从不同角度揭示了这一对偶性：\(\frac{d}{dx} \int_a^x f(t) \, dt = f(x)\) 说明''先积分后微分回到原函数''；\(\int_a^b F'(x) \, dx = F(b) - F(a)\) 说明''先微分后积分回到原函数值差''。这两条定理共同表明：微分和积分是互逆的运算。这一对偶性在更高级的数学中以多种形式出现------微分形式的斯托克斯定理（\(\int_{\partial M} \omega = \int_M d\omega\)）本质上是微积分基本定理在高维空间中的推广，表明''边界上的积分等于内部微分的积分''。在傅里叶分析中，微分和积分之间的对偶关系表现为：微分在频率域中对应乘法（\(\mathcal{F}\{f'\}(\omega) = i\omega \mathcal{F}\{f\}(\omega)\)），而积分对应除法。这一对偶性还体现在分部积分法和换元积分法分别是乘积法则和链式法则的''积分形式''------每一条微分法则都有一条对应的积分法则，这构成了积分技巧的理论来源。

\textbf{定理 24.16}（微积分基本定理的统一表述）

\textbf{第一部分}：若 \(f\) 在 \([a, b]\) 上连续，则 \(\frac{d}{dx}\int_a^x f(t) \, dt = f(x)\)。

\textbf{第二部分}：若 \(F\) 在 \([a, b]\) 上有连续导数 \(F' = f\)，则 \(\int_a^b f(x) \, dx = F(b) - F(a)\)。

\textbf{证明}：第一部分即微积分基本定理（一）（定理 24.7），第二部分即微积分基本定理（二）（定理 24.11），二者均已证明。\(\blacksquare\)

第一部分说明''积分后再微分恢复原函数''，第二部分说明''微分后再积分恢复原函数（相差端点值）``。这两部分共同表明：微分和积分是互逆的运算。

这一对偶关系是微积分的灵魂，它将两个看似不相关的数学问题------求切线斜率（局部问题）和求面积（全局问题）------统一在同一个理论框架中。

\hrulefill

\subsection{24.10 定积分的近似计算}\label{ux5b9aux79efux5206ux7684ux8fd1ux4f3cux8ba1ux7b97}

并非所有函数的定积分都能用初等函数表示------事实上，绝大多数初等函数的原函数都不是初等函数（如 \(\int e^{-x^2} \, dx\)、\(\int \frac{\sin x}{x} \, dx\) 等）。在这些情况下，数值积分方法提供了近似计算定积分的有效途径。矩形法（左端点、右端点、中点）是最简单的数值积分方法，它将积分区间分割为 \(n\) 等份，每份用矩形面积近似曲线下的面积。梯形法将每个小区间上的曲线近似为直线段，其精度为 \(O(1/n^2)\)，比矩形法（\(O(1/n)\)）高一阶。辛普森法（Simpson's Rule）将每两个小区间上的曲线近似为二次抛物线，其精度达到 \(O(1/n^4)\)，是工程实践中最常用的数值积分方法之一。这些数值积分方法的误差分析依赖于积分中值定理和泰勒展开，展示了微积分理论在计算实践中的直接应用。数值积分的''收敛阶''概念------\(O(1/n^k)\) 中的 \(k\)------提供了比较不同方法效率的量化标准，在科学计算中具有重要的指导意义。

\subsubsection{24.10.1 矩形法}\label{ux77e9ux5f62ux6cd5}

\textbf{方法}：将 \([a, b]\) 等分为 \(n\) 份，\(\Delta x = \frac{b-a}{n}\)，分点 \(x_i = a + i\Delta x\)（\(i = 0, 1, \ldots, n\)）。

\[\int_a^b f(x) \, dx \approx \Delta x \sum_{i=0}^{n-1} f(x_i) \quad \text{（左矩形法）}\]

\[\int_a^b f(x) \, dx \approx \Delta x \sum_{i=1}^{n} f(x_i) \quad \text{（右矩形法）}\]

\[\int_a^b f(x) \, dx \approx \Delta x \sum_{i=0}^{n-1} f\left(\frac{x_i + x_{i+1}}{2}\right) \quad \text{（中矩形法）}\]

\subsubsection{24.10.2 梯形法}\label{ux68afux5f62ux6cd5}

\[\int_a^b f(x) \, dx \approx \frac{\Delta x}{2}\left[f(x_0) + 2\sum_{i=1}^{n-1}f(x_i) + f(x_n)\right]\]

\textbf{证明}：将 \([a, b]\) 等分为 \(n\) 个小区间，每个小区间长度为 \(\Delta x = \frac{b-a}{n}\)，分点为 \(x_i = a + i\Delta x\)（\(i = 0, 1, \ldots, n\)）。在每个小区间 \([x_i, x_{i+1}]\) 上，函数 \(f\) 的图像被近似为连接 \((x_i, f(x_i))\) 和 \((x_{i+1}, f(x_{i+1}))\) 的直线段。该直线段下方面积恰好是一个梯形，其面积为

\[\frac{f(x_i) + f(x_{i+1})}{2} \cdot \Delta x\]

将所有 \(n\) 个梯形的面积相加，得到总面积的近似值：

\[\sum_{i=0}^{n-1} \frac{f(x_i) + f(x_{i+1})}{2} \Delta x = \frac{\Delta x}{2} \left[ f(x_0) + 2\sum_{i=1}^{n-1} f(x_i) + f(x_n) \right]\]

当 \(n \to \infty\)（即 \(\Delta x \to 0\)）时，直线段逼近曲线段，梯形面积之和逼近定积分。梯形法的误差可通过 \(f\) 的二阶导数来估计，其精度为 \(O(1/n^2)\)。\(\blacksquare\)

\subsubsection{24.10.3 辛普森法}\label{ux8f9bux666eux68eeux6cd5}

\[\int_a^b f(x) \, dx \approx \frac{\Delta x}{3}\left[f(x_0) + 4\sum_{i \text{ 奇}}f(x_i) + 2\sum_{i \text{ 偶}, i \neq 0, n}f(x_i) + f(x_n)\right]\]

\textbf{注}：辛普森法的精度远高于梯形法。对于 \(n\) 等分，梯形法的误差为 \(O(1/n^2)\)，而辛普森法的误差为 \(O(1/n^4)\)。

\hrulefill

\subsection{24.11 微积分的统一视角}\label{ux5faeux79efux5206ux7684ux7edfux4e00ux89c6ux89d2}

微积分的两大运算------微分和积分------构成了数学分析的核心。从历史的角度看，微分源于求切线问题和瞬时速度问题（牛顿），积分源于求面积问题（阿基米德、黎曼）。牛顿和莱布尼茨的贡献在于发现了这两个看似无关的问题之间的深刻联系。

从现代数学的角度看，微分是线性化的思想------用线性函数（切线）局部近似非线性函数；积分是极限求和的思想------将复杂量分解为无穷多个无穷小量的和。这两种思想的统一体现在微积分基本定理中：

\[\frac{d}{dx}\int_a^x f(t) \, dt = f(x)\]

\[\int_a^b F'(x) \, dx = F(b) - F(a)\]

这一对偶关系不仅在理论上深刻，在应用上也极为广泛。从物理学中的运动方程到经济学中的边际分析，从概率论中的密度函数到工程学中的信号处理，微积分无处不在。

\subsection{24.12 积分不等式}\label{ux79efux5206ux4e0dux7b49ux5f0f}

积分不等式是积分理论的重要组成部分，它们在估计积分值、证明函数性质和求解极值问题中有着广泛的应用。柯西-施瓦茨不等式（积分形式）\(\left(\int_a^b f(x)g(x) \, dx\right)^2 \leq \int_a^b f^2(x) \, dx \cdot \int_a^b g^2(x) \, dx\) 是内积空间中最基本的不等式，其证明基于二次型 \(\int_a^b [f(x) + \lambda g(x)]^2 \, dx \geq 0\) 的判别式非正。明可夫斯基不等式（Minkowski Inequality）\(\left(\int_a^b |f(x) + g(x)|^p \, dx\right)^{1/p} \leq \left(\int_a^b |f(x)|^p \, dx\right)^{1/p} + \left(\int_a^b |g(x)|^p \, dx\right)^{1/p}\) 是三角不等式在 \(L^p\) 空间中的推广。积分形式的詹森不等式（Jensen's Inequality）\(\varphi\left(\frac{1}{b-a} \int_a^b f(x) \, dx\right) \leq \frac{1}{b-a} \int_a^b \varphi(f(x)) \, dx\)（对于凸函数 \(\varphi\)）将函数的平均值与函数值的平均联系了起来。这些不等式不仅是理论分析的工具，也是概率论、信息论和泛函分析的基础------它们共同构成了''积分作为内积''和''函数空间中的度量''的理论框架。

\subsubsection{24.12.1 柯西-施瓦茨不等式（积分形式）}\label{ux67efux897f-ux65bdux74e6ux8328ux4e0dux7b49ux5f0fux79efux5206ux5f62ux5f0f}

\textbf{定理 24.17}（柯西-施瓦茨不等式）设 \(f\) 和 \(g\) 在 \([a, b]\) 上连续，则

\[\left(\int_a^b f(x)g(x) \, dx\right)^2 \leq \int_a^b [f(x)]^2 \, dx \cdot \int_a^b [g(x)]^2 \, dx\]

等号成立当且仅当存在常数 \(\lambda\) 使得 \(f(x) = \lambda g(x)\)（或 \(g(x) = \lambda f(x)\)）对所有 \(x \in [a, b]\) 成立。

\textbf{证明}：对任意实数 \(t\)，

\[0 \leq \int_a^b [f(x) + tg(x)]^2 \, dx = \int_a^b f^2 \, dx + 2t\int_a^b fg \, dx + t^2\int_a^b g^2 \, dx\]

设 \(A = \int_a^b g^2 \, dx \geq 0\)，\(B = \int_a^b fg \, dx\)，\(C = \int_a^b f^2 \, dx \geq 0\)。则 \(At^2 + 2Bt + C \geq 0\) 对所有 \(t\) 成立。

当 \(A > 0\) 时，此二次函数非负当且仅当判别式 \(\Delta = 4B^2 - 4AC \leq 0\)，即 \(B^2 \leq AC\)。

当 \(A = 0\) 时，\(g \equiv 0\)，不等式显然成立。

\(\blacksquare\)

\subsubsection{24.12.2 詹森不等式（积分形式）}\label{ux8a79ux68eeux4e0dux7b49ux5f0fux79efux5206ux5f62ux5f0f}

\textbf{定理 24.18}（詹森不等式）设 \(f\) 是凸函数，\(g\) 在 \([a, b]\) 上可积，\(p\) 在 \([a, b]\) 上非负可积且 \(\int_a^b p(x) \, dx > 0\)。则

\[f\left(\frac{\int_a^b p(x)g(x) \, dx}{\int_a^b p(x) \, dx}\right) \leq \frac{\int_a^b p(x)f(g(x)) \, dx}{\int_a^b p(x) \, dx}\]

\textbf{证明}：将积分看作连续版本的加权平均，由 \(f\) 的凸性推广到连续情形。

详细论证：对 \([a, b]\) 的任意分割 \(a = x_0 < x_1 < \cdots < x_n = b\)，取 \(\xi_i \in [x_{i-1}, x_i]\)，令 \(\lambda_i = \frac{x_i - x_{i-1}}{b - a}\)（\(\sum \lambda_i = 1\)）。由离散詹森不等式，\(f\left(\sum_{i=1}^n \lambda_i \xi_i\right) \leq \sum_{i=1}^n \lambda_i f(\xi_i)\)。左边为 \(f\left(\frac{1}{b-a}\sum_{i=1}^n f(\xi_i)(x_i - x_{i-1})\right)\) 的黎曼和逼近，右边为 \(\frac{1}{b-a}\sum_{i=1}^n f(\xi_i)(x_i - x_{i-1})\) 的黎曼和。令分割加细（最大区间长度趋于零），取极限即得 \(f\left(\frac{1}{b-a}\int_a^b g(x)dx\right) \leq \frac{1}{b-a}\int_a^b f(g(x))dx\)。\(\blacksquare\)

\hrulefill

\subsection{24.13 微积分基本思想回顾}\label{ux5faeux79efux5206ux57faux672cux601dux60f3ux56deux987e}

本卷的内容从函数概念出发，经过基本初等函数的逐一研究，建立起了极限理论这一分析学的基石，进而发展出导数和积分两大核心概念。

\textbf{导数的本质}是线性逼近：在一点附近，用切线代替曲线，用线性函数代替非线性函数。这一思想的精确表述就是：

\[f(x) = f(a) + f'(a)(x - a) + o(x - a)\]

\textbf{积分的本质}是无穷累加：将连续变化的量分解为无穷多个无穷小量之和。黎曼积分的定义体现了''分割、近似、求和、取极限''的基本方法论。

\textbf{微积分基本定理}将两者统一：微分和积分互为逆运算。这不仅是计算工具，更是一种深刻的数学思想------局部性质（导数）与整体性质（积分）之间的对偶关系。

这种对偶关系在高等数学中以多种方式展现：格林公式、高斯公式、斯托克斯公式都是它的推广；傅里叶变换中的时域-频域对偶也是同一思想的体现。

\textbf{本章展望}：本章建立了定积分的严格理论和基本计算方法。积分的理论远不止于此------更深入的学习包括：黎曼可积的充要条件、反常积分（无穷区间或无界函数的积分）、重积分、线面积分等。积分理论在物理学（功、能量、质心）、概率论（概率密度函数的积分）、经济学（消费者剩余）等领域有广泛应用。

\hrulefill

\hrulefill

\hrulefill

\emph{本卷以极限的严格理论（Ch 29）为分析学的基石------用 ε-N 语言定义数列极限、用 ε-δ 语言定义函数极限，证明了极限的基本性质和重要极限，建立了连续函数的介值定理和最值定理；导数与微分（Ch 30）从变化率问题出发，定义了导数的严格概念和几何意义，推导了基本导数公式和求导法则（四则运算、链式法则），证明了罗尔定理、拉格朗日中值定理和柯西中值定理，并应用洛必达法则和极值理论解决优化问题；积分初步（Ch 31）用黎曼和的极限严格定义了定积分，证明了牛顿-莱布尼茨公式（微积分基本定理），建立了换元积分法和分部积分法这一套完整的积分计算工具。这三章构成了微积分的基础理论框架，为后续高等数学（实分析、复分析、微分方程、泛函分析等）奠定了严格的分析基础。}

\emph{卷四B：微积分初步终。}

\emph{卷四B：微积分初步终。}

\emph{卷四B：微积分初步终。}

\emph{卷四：函数、解析几何与微积分终。}

\newpage

\chapter{04-离散数学/01-组合数学/01-计数与递推.md}\label{ux79bbux6563ux6570ux5b6601-ux7ec4ux5408ux6570ux5b6601-ux8ba1ux6570ux4e0eux9012ux63a8.md}

计数原理是组合数学的基石。组合数学研究的是有限或可数无限结构的计数、存在性和构造性问题。计数理论的两个最基本原则------加法原理和乘法原理------看似平凡，但它们的正确运用是解决一切复杂计数问题的关键。加法原理对应''或''的逻辑（分类讨论），乘法原理对应''且''的逻辑（分步完成）。这两个原理的严格形式化需要集合论的语言：加法原理是两两不交有限集的并的基数公式，乘法原理是笛卡尔积的基数公式。在本章中，我们从这两个基本原理出发，逐步建立起排列、组合、二项式定理、容斥原理、递推关系、鸽巢原理、生成函数、特殊计数序列（Catalan 数、Stirling 数、Bell 数）和 Polya 计数理论等完整的组合计数体系。

\section{第79章：计数基本原理}\label{ux7b2c79ux7ae0ux8ba1ux6570ux57faux672cux539fux7406}

\subsection{28.1 加法原理与乘法原理}\label{ux52a0ux6cd5ux539fux7406ux4e0eux4e58ux6cd5ux539fux7406}

\begin{quote}
\textbf{直觉与动机} 计数原理是组合数学的基石。加法原理对应''或''的逻辑（分类讨论），乘法原理对应''且''的逻辑（分步完成）。它们看似显而易见，却是解决一切复杂计数问题的基础。
\end{quote}

\textbf{原理 28.1（加法原理/分类计数原理）} 完成一件事有 \(n\) 类方法。第 \(i\) 类方法有 \(m_i\) 种方案（\(i = 1, 2, \ldots, n\)），且各类方案互不重叠，则完成这件事共有： \[N = m_1 + m_2 + \cdots + m_n\] 种方法。

\textbf{原理 28.2（乘法原理/分步计数原理）} 完成一件事需要 \(n\) 个步骤。第 \(i\) 步有 \(m_i\) 种方案（\(i = 1, 2, \ldots, n\)），则完成这件事共有： \[N = m_1 \times m_2 \times \cdots \times m_n\] 种方法。

\textbf{证明（集合语言）} 加法原理等价于：若 \(A_1, A_2, \ldots, A_n\) 是两两不交的有限集，则 \(|A_1 \cup A_2 \cup \cdots \cup A_n| = |A_1| + |A_2| + \cdots + |A_n|\)。这由集合的基数性质直接得到。

乘法原理等价于笛卡尔积的基数：\(|A_1 \times A_2 \times \cdots \times A_n| = |A_1| \cdot |A_2| \cdot \cdots \cdot |A_n|\)。\(\blacksquare\)

\textbf{命题 28.1（加法原理的推广------容斥原理的雏形）} 若 \(A_1, A_2\) 是有限集（不一定不相交），则 \(|A_1 \cup A_2| = |A_1| + |A_2| - |A_1 \cap A_2|\)。这一公式是容斥原理（见 28.5 节）的最简情形，也是加法原理在有重叠时的自然推广。

\textbf{证明}：将 \(A_1 \cup A_2\) 分解为三个两两不交的部分：\(A_1 \setminus A_2\)、\(A_2 \setminus A_1\) 和 \(A_1 \cap A_2\)。则 \(|A_1 \cup A_2| = |A_1 \setminus A_2| + |A_2 \setminus A_1| + |A_1 \cap A_2|\)。而 \(|A_1| = |A_1 \setminus A_2| + |A_1 \cap A_2|\)，\(|A_2| = |A_2 \setminus A_1| + |A_1 \cap A_2|\)。代入即得。\(\blacksquare\)

\textbf{命题 28.2（乘法原理与独立选择）} 乘法原理的实质是：每一步的选择独立于其他步骤的选择。更精确地说，若第 \(i\) 步的选择方案数 \(m_i\) 不依赖于前面各步的具体选择（仅依赖于前面各步已经完成这一事实），则乘法原理成立。若第 \(i\) 步的方案数依赖于前面的选择，则需要使用条件乘法原理（即求和或递推方法）。

\subsection{28.2 排列}\label{ux6392ux5217}

\begin{quote}
\textbf{直觉与动机} 排列问题关注的是''顺序''。从 \(n\) 个不同元素中取出 \(m\) 个排成一列，不同的排法有多少种？这个问题在生活中随处可见：排队、安排座位、密码设置等。
\end{quote}

\textbf{定义 28.1（排列）} 从 \(n\) 个不同元素中取出 \(m\)（\(m \leq n\)）个元素，按照一定的顺序排成一列，称为从 \(n\) 个元素中取 \(m\) 个元素的一个\textbf{排列}。所有不同排列的个数记作 \(A_n^m\)（或 \(P(n, m)\)）。

\textbf{定理 28.1} \(A_n^m = n(n-1)(n-2) \cdots (n-m+1) = \frac{n!}{(n-m)!}\)

\textbf{证明}：利用乘法原理。第一个位置有 \(n\) 种选择，第二个位置有 \(n-1\) 种选择（已选走一个），\ldots，第 \(m\) 个位置有 \(n - m + 1\) 种选择。故： \[A_n^m = n(n-1)(n-2) \cdots (n-m+1)\]

当 \(m = n\) 时，\(A_n^n = n!\)（\(n\) 个元素的\textbf{全排列}）。\(\blacksquare\)

\textbf{命题 28.3（有重复元素的排列）} 设 \(n\) 个元素中有 \(n_1\) 个不可区分的类型 1 元素，\(n_2\) 个不可区分的类型 2 元素，\ldots，\(n_k\) 个不可区分的类型 \(k\) 元素，且 \(\sum_{i=1}^k n_i = n\)。则这 \(n\) 个元素的全排列数为： \[\frac{n!}{n_1! \, n_2! \, \cdots \, n_k!}\]

\textbf{证明}：先将 \(n\) 个元素视为全不同的，有 \(n!\) 种排列。但类型 \(i\) 的 \(n_i\) 个元素之间的 \(n_i!\) 种排列实际上不可区分，因此需除以 \(\prod_{i=1}^k n_i!\)。\(\blacksquare\)

此公式也称为\textbf{多项式系数}，记作 \(\binom{n}{n_1, n_2, \ldots, n_k}\)，它是 \((x_1 + x_2 + \cdots + x_k)^n\) 展开式中 \(x_1^{n_1} x_2^{n_2} \cdots x_k^{n_k}\) 的系数。

\textbf{命题 28.4（循环排列）} \(n\) 个不同元素排成一个圆圈（旋转视为相同），不同的排列方式数为 \((n-1)!\)。若进一步考虑翻转对称（反射视为相同），则为 \((n-1)!/2\)（\(n \geq 3\)）。

\textbf{证明}：将 \(n\) 个元素的全排列数 \(n!\) 除以 \(n\)（\(n\) 个旋转等价类），即得 \((n-1)!\)。若考虑翻转，再除以 \(2\)（二面体群 \(D_n\) 的阶为 \(2n\)，故轨道数为 \(n!/(2n) = (n-1)!/2\)）。\(\blacksquare\)

\subsection{28.3 组合}\label{ux7ec4ux5408}

\begin{quote}
\textbf{直觉与动机} 与排列不同，组合只关心''选了哪些''而不关心''顺序''。从 \(n\) 个人中选 \(m\) 个人组成一个委员会，不同的选法有多少种？选出 \(\{A, B, C\}\) 和 \(\{B, C, A\}\) 是同一个组合。
\end{quote}

\textbf{定义 28.2（组合）} 从 \(n\) 个不同元素中取出 \(m\)（\(m \leq n\)）个元素组成一组（不考虑顺序），称为从 \(n\) 个元素中取 \(m\) 个元素的一个\textbf{组合}。所有不同组合的个数记作 \(C_n^m\)（或 \(\binom{n}{m}\)）。

\textbf{定理 28.2} \(C_n^m = \frac{A_n^m}{m!} = \frac{n!}{m!(n-m)!}\)

\textbf{证明}：每个包含 \(m\) 个元素的组合，将其 \(m\) 个元素进行全排列，可以得到 \(m!\) 个不同的排列。反过来，每个排列对应唯一的组合。因此 \(A_n^m = C_n^m \cdot m!\)，即 \(C_n^m = \frac{A_n^m}{m!}\)。\(\blacksquare\)

\textbf{性质 28.1} 组合数的基本性质： 1. \(C_n^m = C_n^{n-m}\)（对称性） 2. \(C_n^m = C_{n-1}^{m-1} + C_{n-1}^m\)（帕斯卡公式/递推关系） 3. \(\frac{m}{n} C_n^m = C_{n-1}^{m-1}\)（吸收恒等式） 4. \(C_n^m C_m^k = C_n^k C_{n-k}^{m-k}\)（三项式系数恒等式）

\textbf{证明对称性} \(C_n^m = \frac{n!}{m!(n-m)!} = \frac{n!}{(n-m)! \cdot [n-(n-m)]!} = C_n^{n-m}\)。\(\blacksquare\)

\textbf{证明帕斯卡公式} 从 \(n\) 个元素中选 \(m\) 个，考虑某个特定元素 \(a\)： - 包含 \(a\)：从剩余 \(n-1\) 个中再选 \(m-1\) 个，共 \(C_{n-1}^{m-1}\) 种。 - 不含 \(a\)：从剩余 \(n-1\) 个中选 \(m\) 个，共 \(C_{n-1}^{m}\) 种。

由加法原理，\(C_n^m = C_{n-1}^{m-1} + C_{n-1}^m\)。\(\blacksquare\)

\textbf{证明吸收恒等式} \(\frac{m}{n} C_n^m = \frac{m}{n} \cdot \frac{n!}{m!(n-m)!} = \frac{(n-1)!}{(m-1)!(n-m)!} = C_{n-1}^{m-1}\)。\(\blacksquare\)

\textbf{证明三项式系数恒等式} 左边：从 \(n\) 选 \(m\)，再从 \(m\) 选 \(k\)。右边：从 \(n\) 选 \(k\)，再从剩余的 \(n-k\) 选 \(m-k\)。两种方式计数同一个集合------大小为 \(n\) 的集合中选出一个 \(m\) 元子集，再在其中选出一个 \(k\) 元子集。\(\blacksquare\)

\begin{quote}
\textbf{帕斯卡公式的意义} 帕斯卡公式正是杨辉三角（Pascal's triangle）的递推关系。杨辉三角的第 \(n\) 行第 \(m\) 个数就是 \(C_n^m\)，每个数等于它上方两个数之和。
\end{quote}

\textbf{命题 28.5（Vandermonde 恒等式）} 对任意非负整数 \(m, n, r\)： \[\sum_{k=0}^{r} \binom{m}{k}\binom{n}{r-k} = \binom{m+n}{r}\]

\textbf{证明}：考虑从 \(m\) 个男生和 \(n\) 个女生中选出 \(r\) 个人的组合数。一方面，直接选法有 \(\binom{m+n}{r}\) 种。另一方面，按选出的男生人数 \(k\) 分类：选 \(k\) 个男生（\(\binom{m}{k}\) 种）和 \(r-k\) 个女生（\(\binom{n}{r-k}\) 种），对 \(k=0,1,\ldots,r\) 求和。由加法原理，两种计数相等。\(\blacksquare\)

\textbf{命题 28.6（朱世杰恒等式 / 曲棍球棒恒等式）} \[\sum_{k=m}^{n} \binom{k}{m} = \binom{n+1}{m+1}\]

\textbf{证明}：对 \(n\) 用归纳法，或由反复应用帕斯卡公式：\(\binom{n+1}{m+1} = \binom{n}{m+1} + \binom{n}{m} = \binom{n-1}{m+1} + \binom{n-1}{m} + \binom{n}{m} = \cdots = \binom{m}{m+1} + \sum_{k=m}^{n} \binom{k}{m} = \sum_{k=m}^{n} \binom{k}{m}\)。\(\blacksquare\)

\subsection{28.4 二项式定理}\label{ux4e8cux9879ux5f0fux5b9aux7406}

\textbf{定理 28.3（二项式定理）} 对任意正整数 \(n\) 和实数 \(a, b\)： \[(a + b)^n = \sum_{k=0}^{n} C_n^k \, a^{n-k} b^k = \sum_{k=0}^{n} \binom{n}{k} a^{n-k} b^k\]

\textbf{证明（组合方法）} 展开 \((a+b)^n = \underbrace{(a+b)(a+b)\cdots(a+b)}_{n \text{ 个}}\)。

展开后的每一项形如 \(a^{i_1} b^{j_1} \cdot a^{i_2} b^{j_2} \cdots\)，其中每个因子贡献一个 \(a\) 或一个 \(b\)。最终得到的单项式为 \(a^{n-k} b^k\)，其中 \(k\) 表示 \(n\) 个因子中选择 \(b\) 的个数。

对于固定的 \(k\)（\(0 \leq k \leq n\)），从 \(n\) 个因子中选择 \(k\) 个贡献 \(b\) 的方式恰有 \(C_n^k\) 种，因此 \(a^{n-k}b^k\) 的系数为 \(C_n^k\)。

故： \[(a+b)^n = \sum_{k=0}^{n} C_n^k \, a^{n-k} b^k\] \(\blacksquare\)

\textbf{证明（数学归纳法）} 对 \(n\) 归纳。\(n = 1\) 时：\((a+b)^1 = a + b = \binom{1}{0}a^1 b^0 + \binom{1}{1}a^0 b^1\)。✓

假设 \(n = m\) 时成立。对 \(n = m+1\)： \[(a+b)^{m+1} = (a+b)(a+b)^m = (a+b)\sum_{r=0}^{m}\binom{m}{r}a^{m-r}b^r\] \[= \sum_{r=0}^{m}\binom{m}{r}a^{m+1-r}b^r + \sum_{r=0}^{m}\binom{m}{r}a^{m-r}b^{r+1}\]

在第二个求和中令 \(s = r + 1\)： \[= \sum_{r=0}^{m}\binom{m}{r}a^{m+1-r}b^r + \sum_{s=1}^{m+1}\binom{m}{s-1}a^{m+1-s}b^s\]

合并同类项（利用帕斯卡公式 \(\binom{m}{r} + \binom{m}{r-1} = \binom{m+1}{r}\)）： \[= a^{m+1} + \sum_{r=1}^{m}\left[\binom{m}{r} + \binom{m}{r-1}\right]a^{m+1-r}b^r + b^{m+1} = \sum_{r=0}^{m+1}\binom{m+1}{r}a^{m+1-r}b^r\] \(\blacksquare\)

\textbf{推论 28.1（二项式系数的性质）（n ≥ 1）}

\textbf{证明}：由二项式定理 \((a+b)^n = \sum_{k=0}^{n} C_n^k a^{n-k} b^k\)：

\begin{enumerate}
\def\labelenumi{\arabic{enumi}.}
\tightlist
\item
  \(\sum_{k=0}^{n} C_n^k = 2^n\)（令 \(a = b = 1\)）
\item
  \(\sum_{k=0}^{n} (-1)^k C_n^k = 0\)（令 \(a = 1, b = -1\)）
\item
  \(C_n^0 + C_n^2 + C_n^4 + \cdots = C_n^1 + C_n^3 + C_n^5 + \cdots = 2^{n-1}\)（由 (1) 与 (2) 两式相加、相减得）
\item
  \(\sum_{k=0}^{n} k C_n^k = n 2^{n-1}\)（对 \((1+x)^n\) 求导后令 \(x=1\)）
\item
  \(\sum_{k=0}^{n} k^2 C_n^k = n(n+1)2^{n-2}\)（对 \((1+x)^n\) 求二阶导后令 \(x=1\)） \(\blacksquare\)
\end{enumerate}

\textbf{定理 28.3b（多项式定理）} 对任意正整数 \(n\) 和实数 \(x_1, x_2, \ldots, x_m\)： \[(x_1 + x_2 + \cdots + x_m)^n = \sum_{\substack{k_1 + \cdots + k_m = n \\ k_i \geq 0}} \frac{n!}{k_1! \, k_2! \, \cdots \, k_m!} \, x_1^{k_1} x_2^{k_2} \cdots x_m^{k_m}\]

\textbf{证明}：展开 \((x_1 + \cdots + x_m)^n\)，每一项对应从 \(n\) 个因子中各选一个 \(x_i\)。若选 \(x_i\) 恰好 \(k_i\) 次，则此项为 \(x_1^{k_1} \cdots x_m^{k_m}\)。从 \(n\) 个因子中分配 \(k_1\) 个给 \(x_1\)，\(k_2\) 个给 \(x_2\)，\ldots，\(k_m\) 个给 \(x_m\) 的方式数为 \(\frac{n!}{k_1! k_2! \cdots k_m!}\)。\(\blacksquare\)

\subsection{28.5 容斥原理}\label{ux5bb9ux65a5ux539fux7406}

\begin{quote}
\textbf{直觉与动机} 当计数涉及''至少满足一个条件''时，不能简单地将各自满足条件的元素数相加------因为一个元素可能同时满足多个条件，导致重复计数。容斥原理（Inclusion-Exclusion Principle）通过加减交替的方式精确纠正这种重叠。
\end{quote}

\textbf{定理 28.4（容斥原理）} 设 \(A_1, A_2, \ldots, A_n\) 为有限集 \(X\) 的子集。则

\[|A_1 \cup A_2 \cup \cdots \cup A_n| = \sum_{i=1}^{n} |A_i| - \sum_{1 \leq i < j \leq n} |A_i \cap A_j| + \sum_{1 \leq i < j < k \leq n} |A_i \cap A_j \cap A_k| - \cdots + (-1)^{n-1} |A_1 \cap A_2 \cap \cdots \cap A_n|\]

\textbf{证明}：考虑任意元素 \(x \in X\)。设 \(x\) 恰好属于 \(r\) 个集合（\(0 \leq r \leq n\)）。若 \(r = 0\)，\(x\) 对左右两端的贡献均为 \(0\)。若 \(r \geq 1\)，\(x\) 对左端贡献 \(1\)。对右端，\(x\) 在 \(\sum |A_i|\) 中被计数 \(\binom{r}{1}\) 次，在 \(\sum |A_i \cap A_j|\) 中被计数 \(\binom{r}{2}\) 次，依此类推。故右端贡献为 \(\binom{r}{1} - \binom{r}{2} + \binom{r}{3} - \cdots + (-1)^{r-1}\binom{r}{r} = 1 - (1-1)^r = 1\)。因此每个属于至少一个集合的元素恰好被计数一次。\(\blacksquare\)

\textbf{定理 28.5（容斥原理的补集形式）} 不属于任何 \(A_i\) 的元素个数为

\[|X \setminus (A_1 \cup \cdots \cup A_n)| = |X| - \sum_{i} |A_i| + \sum_{i<j} |A_i \cap A_j| - \cdots + (-1)^n |A_1 \cap \cdots \cap A_n|\]

此形式在计算''无限制''计数时尤为有用------例如计算错位排列数。\(\blacksquare\)

\textbf{命题 28.7（容斥原理的符号形式）} 设 \(X\) 为 \(n\) 元集，\(P_1, \ldots, P_n\) 为 \(n\) 个性质。令 \(N(\supseteq T)\) 表示至少满足 \(T \subseteq \{1,\ldots,n\}\) 中所有性质的元素个数。则不具有任何性质的元素个数为： \[N_{\varnothing} = \sum_{T \subseteq \{1,\ldots,n\}} (-1)^{|T|} N(\supseteq T)\]

\textbf{证明}：这是定理 28.5 的等价表述，其中 \(N(\supseteq T) = |\bigcap_{i \in T} A_i|\)。\(\blacksquare\)

\textbf{命题 28.8（Euler 函数 \(\varphi(n)\) 的容斥原理推导）} 设 \(n = p_1^{a_1} p_2^{a_2} \cdots p_k^{a_k}\) 为 \(n\) 的素因子分解。则 \[\varphi(n) = n \prod_{i=1}^{k} \left(1 - \frac{1}{p_i}\right)\]

\textbf{证明}：令 \(X = \{1, 2, \ldots, n\}\)，\(A_i = \{x \in X : p_i \mid x\}\)（\(p_i\) 的倍数）。则 \(|A_i| = n/p_i\)，\(|A_i \cap A_j| = n/(p_i p_j)\)，等等。由容斥原理的补集形式： \[\varphi(n) = n - \sum_i \frac{n}{p_i} + \sum_{i<j} \frac{n}{p_i p_j} - \cdots = n \prod_{i=1}^{k} \left(1 - \frac{1}{p_i}\right)\] \(\blacksquare\)

\section{第80章：递推与高级计数方法}\label{ux7b2c80ux7ae0ux9012ux63a8ux4e0eux9ad8ux7ea7ux8ba1ux6570ux65b9ux6cd5}

\begin{quote}
\textbf{直觉与动机} 许多计数问题中，规模为 \(n\) 的问题的解可通过规模较小的同类问题的解表达出来。递推关系（recurrence relation）正是刻画这种结构的方法。递推关系与生成函数一起构成了组合计数中最重要的两项技术。
\end{quote}

\textbf{定义 28.3（线性常系数齐次递推）} 设 \(a_n\) 满足

\[a_n = c_1 a_{n-1} + c_2 a_{n-2} + \cdots + c_k a_{n-k} \quad (n \geq k)\]

其中 \(c_1, \ldots, c_k\) 为常数且 \(c_k \neq 0\)。称此递推为 \(k\) 阶线性常系数齐次递推关系。其特征方程为 \(x^k - c_1 x^{k-1} - \cdots - c_k = 0\)。

\textbf{定理 28.6（递推关系的通解）} 若特征方程有 \(k\) 个互异根 \(r_1, \ldots, r_k\)，则通解为 \(a_n = \alpha_1 r_1^n + \alpha_2 r_2^n + \cdots + \alpha_k r_k^n\)，其中 \(\alpha_i\) 由初始条件确定。若有重根 \(r\)（重数 \(m\)），则对应项为 \((\beta_1 + \beta_2 n + \cdots + \beta_m n^{m-1}) r^n\)。

\textbf{证明}：将 \(a_n = r^n\) 代入递推式，得 \(r^n = c_1 r^{n-1} + \cdots + c_k r^{n-k}\)，即 \(r^{n-k}(r^k - c_1 r^{k-1} - \cdots - c_k) = 0\)。\(r \neq 0\) 时恰为特征方程。线性递推的解空间维数为 \(k\)，上述 \(k\) 个线性无关解构成基。重根情形下，\(r^n, nr^n, \ldots, n^{m-1}r^n\) 构成 \(m\) 个线性无关解。\(\blacksquare\)

\textbf{定理 28.6b（非齐次线性递推）} 对于非齐次递推 \(a_n = c_1 a_{n-1} + \cdots + c_k a_{n-k} + f(n)\)，通解为齐次通解 \(a_n^{(h)}\) 加上一个特解 \(a_n^{(p)}\)。特解的形式由 \(f(n)\) 的形式决定： - 若 \(f(n)\) 是 \(n\) 的 \(d\) 次多项式，则特解可设为 \(n\) 的 \(d\) 次多项式（若 \(1\) 不是特征根）或乘以 \(n^s\)（\(s\) 为 \(1\) 作为特征根的重数）。 - 若 \(f(n) = c \cdot r^n\)，则特解可设为 \(A r^n\)（若 \(r\) 不是特征根）或 \(A n^s r^n\)（\(s\) 为 \(r\) 作为特征根的重数）。

\textbf{经典例子}：Fibonacci 数列 \(F_n = F_{n-1} + F_{n-2}\)（\(F_0 = 0, F_1 = 1\)）的特征方程为 \(x^2 - x - 1 = 0\)，根为 \(\phi = \frac{1+\sqrt{5}}{2}\) 和 \(\psi = \frac{1-\sqrt{5}}{2}\)。通解为 \(F_n = \frac{\phi^n - \psi^n}{\sqrt{5}}\)（Binet 公式）。\(\blacksquare\)

\textbf{命题 28.9（Fibonacci 数列的性质）} 1. \(\sum_{i=1}^{n} F_i = F_{n+2} - 1\) 2. \(\sum_{i=1}^{n} F_i^2 = F_n F_{n+1}\) 3. \(\gcd(F_m, F_n) = F_{\gcd(m, n)}\) 4. \(F_{n-1}F_{n+1} - F_n^2 = (-1)^n\)（Cassini 恒等式）

\textbf{证明}：(1) 由 \(F_{k+2} = F_{k+1} + F_k\)，重排得 \(F_k = F_{k+2} - F_{k+1}\)。求和 telescoping：\(\sum_{i=1}^{n} F_i = \sum_{i=1}^{n} (F_{i+2} - F_{i+1}) = F_{n+2} - F_2 = F_{n+2} - 1\)。 (2) 用归纳法。\(n=1\)：\(F_1^2 = 1\)，\(F_1 F_2 = 1 \cdot 1 = 1\)。归纳步：\(\sum_{i=1}^{n+1} F_i^2 = F_n F_{n+1} + F_{n+1}^2 = F_{n+1}(F_n + F_{n+1}) = F_{n+1}F_{n+2}\)。 (3) 由 \(F_{m+n} = F_{m-1}F_n + F_m F_{n+1}\)，可证 \(\gcd(F_m, F_n) = \gcd(F_m, F_{m+n})\)。对 Euclid 算法归纳。 (4) 用矩阵形式：\(\begin{pmatrix} F_{n+1} & F_n \\ F_n & F_{n-1} \end{pmatrix} = \begin{pmatrix} 1 & 1 \\ 1 & 0 \end{pmatrix}^n\)。取行列式，\(\det \begin{pmatrix} 1 & 1 \\ 1 & 0 \end{pmatrix} = -1\)，故 \(\det = (-1)^n\)，即 \(F_{n+1}F_{n-1} - F_n^2 = (-1)^n\)。\(\blacksquare\)

\subsection{28.7 鸽巢原理}\label{ux9e3dux5de2ux539fux7406}

\begin{quote}
\textbf{直觉与动机} 鸽巢原理（Pigeonhole Principle）是最简单的组合原理之一，其陈述看似平凡，但应用极为广泛。Dirichlet 在 1834 年首次明确使用此原理证明数论中的逼近定理，故也称 Dirichlet 抽屉原理。
\end{quote}

\textbf{定理 28.7（鸽巢原理）} 若将 \(n+1\) 个物体放入 \(n\) 个盒子中，则至少有一个盒子包含至少两个物体。更一般地，若将 \(m\) 个物体放入 \(n\) 个盒子（\(m > n\)），则至少有一个盒子包含至少 \(\lceil m/n \rceil\) 个物体。

\textbf{证明}：若每个盒子至多包含 \(\lceil m/n \rceil - 1\) 个物体，则总物体数至多为 \(n(\lceil m/n \rceil - 1) < n \cdot (m/n) = m\)，矛盾。特别地，\(m = n+1\) 时 \(\lceil (n+1)/n \rceil = 2\)。\(\blacksquare\)

\textbf{定理 28.8（Erdős--Szekeres 定理）} 任意长度为 \(n^2 + 1\) 的实数序列包含一个长度为 \(n+1\) 的单调递增子序列或长度为 \(n+1\) 的单调递减子序列。

\textbf{证明}：对每个位置 \(i\)，记 \(a_i\) 为以位置 \(i\) 结尾的最长递增子序列长度，\(b_i\) 为以位置 \(i\) 开头的最长递减子序列长度。若所有 \(a_i \leq n\) 且所有 \(b_i \leq n\)，则有序对 \((a_i, b_i)\) 至多有 \(n^2\) 种可能。由鸽巢原理，存在 \(i < j\) 使 \((a_i, b_i) = (a_j, b_j)\)。但若 \(x_i < x_j\) 则 \(a_i < a_j\)，若 \(x_i > x_j\) 则 \(b_i > b_j\)，均矛盾。故结论成立。\(\blacksquare\)

\textbf{定理 28.9（Ramsey 定理，\(R(3,3)\) 情形）} 任意 6 个人的聚会中，要么存在 3 个人两两互相认识，要么存在 3 个人两两互不认识。即 \(R(3,3) = 6\)。

\textbf{证明}：任选一人 \(A\)。在其余 5 人中，由鸽巢原理，\(A\) 要么认识至少 3 人，要么不认识至少 3 人。设 \(A\) 认识 \(B, C, D\)。若 \(B, C, D\) 中有两人互相认识，则此两人与 \(A\) 构成三人互相认识。若 \(B, C, D\) 两两互不认识，则此三人构成两两互不认识。\(A\) 不认识至少 3 人的情形对称。因此 \(R(3,3) \leq 6\)。构造 5 人反例（五人围成一圈，每人认识左右邻居但不认识对面两人）表明 \(R(3,3) > 5\)。故 \(R(3,3) = 6\)。\(\blacksquare\)

\textbf{定理 28.9b（一般 Ramsey 定理）} 对任意正整数 \(r, s\)，存在最小正整数 \(R(r, s)\) 使得任意 \(R(r, s)\) 个顶点的完全图的边用红蓝两色染色后，必存在红色 \(K_r\) 或蓝色 \(K_s\)。Ramsey 数满足 \(R(r, s) \leq R(r-1, s) + R(r, s-1)\)。

\textbf{证明}：选取任意顶点 \(v\)。在其余 \(R(r-1, s) + R(r, s-1) - 1\) 个顶点中，与 \(v\) 以红边相连的顶点数 \(\geq R(r-1, s)\) 或以蓝边相连的顶点数 \(\geq R(r, s-1)\)。前一情形：红邻居中由归纳假设存在红色 \(K_{r-1}\)（与 \(v\) 构成红色 \(K_r\)）或蓝色 \(K_s\)。后一情形对称。\(\blacksquare\)

\subsection{28.8 错位排列与圆排列}\label{ux9519ux4f4dux6392ux5217ux4e0eux5706ux6392ux5217}

\textbf{定义 28.4（错位排列 / Derangement）} \(n\) 个元素的排列中，没有任何元素出现在其原始位置的排列称为\textbf{错位排列}。错位排列数 \(D_n\)（又称 \(!n\)，即 \(n\) 的子阶乘）满足递推关系 \(D_n = (n-1)(D_{n-1} + D_{n-2})\)，\(D_0 = 1, D_1 = 0\)。

\textbf{定理 28.10} \(D_n = n! \sum_{k=0}^{n} \frac{(-1)^k}{k!}\)。特别地，\(D_n \approx n! / e\)（\(n\) 充分大时）。

\textbf{证明}：由容斥原理的补集形式。设 \(A_i\) 为第 \(i\) 个元素保持在原位的排列集合。则 \(|A_i| = (n-1)!\)，\(|A_i \cap A_j| = (n-2)!\)（\(i < j\)），等等。因此 \(D_n = n! - \binom{n}{1}(n-1)! + \binom{n}{2}(n-2)! - \cdots + (-1)^n \binom{n}{n}0! = n! \sum_{k=0}^{n} \frac{(-1)^k}{k!}\)。当 \(n \to \infty\) 时，\(\sum_{k=0}^{n} (-1)^k/k! \to e^{-1}\)，故 \(D_n \approx n!/e\)。\(\blacksquare\)

\textbf{证明（递推关系的推导）} 考虑第一个元素的位置。在错位排列中，第一个元素不能在第 1 个位置，必须移到某个位置 \(j\)（\(2 \leq j \leq n\)），有 \(n-1\) 种选择。对于位置 \(j\)：若第 \(j\) 个元素移到了第 1 个位置，则剩余 \(n-2\) 个元素的错位排列数为 \(D_{n-2}\)。若第 \(j\) 个元素没有移到第 1 个位置，则将第 1 个位置视为第 \(j\) 个元素的''禁止位置''，剩余 \(n-1\) 个元素的错位排列数为 \(D_{n-1}\)。因此 \(D_n = (n-1)(D_{n-1} + D_{n-2})\)。\(\blacksquare\)

\textbf{定义 28.5（圆排列）} \(n\) 个不同元素围成一圈的不同排列方式数（旋转视为相同）为 \((n-1)!\)。若进一步考虑翻转对称（反射视为相同），则为 \((n-1)!/2\)（\(n \geq 3\)）。圆排列的本质是 \(n\) 阶循环群 \(\mathbb{Z}_n\) 作用在 \(n!\) 个线性排列上的轨道计数，是 Burnside 引理的最简单应用。\(\blacksquare\)

\subsection{28.9 生成函数}\label{ux751fux6210ux51fdux6570}

\begin{quote}
\textbf{直觉与动机} 生成函数（generating function）是组合数学中最强大的工具之一。其核心思想是将一个数列 \(\{a_n\}\) 编码为一个形式幂级数，然后利用幂级数的代数运算来研究数列的组合性质。生成函数将离散的计数问题转化为连续的解析问题，从而可以借助分析的工具（如微积分、复分析）来解决组合问题。
\end{quote}

\textbf{定义 28.6（普通生成函数）} 数列 \(\{a_n\}_{n=0}^{\infty}\) 的\textbf{普通生成函数}（ordinary generating function, OGF）定义为形式幂级数： \[A(x) = \sum_{n=0}^{\infty} a_n x^n\]

\textbf{定义 28.7（指数生成函数）} 数列 \(\{a_n\}_{n=0}^{\infty}\) 的\textbf{指数生成函数}（exponential generating function, EGF）定义为： \[\hat{A}(x) = \sum_{n=0}^{\infty} a_n \frac{x^n}{n!}\]

\textbf{定理 28.11（生成函数的基本运算）} 设 \(A(x)\) 和 \(B(x)\) 分别为 \(\{a_n\}\) 和 \(\{b_n\}\) 的 OGF。则： 1. \(A(x) + B(x)\) 是 \(\{a_n + b_n\}\) 的 OGF。 2. \(A(x) \cdot B(x)\) 是 \(\{c_n = \sum_{k=0}^{n} a_k b_{n-k}\}\)（卷积）的 OGF。 3. \(x A'(x)\) 是 \(\{n a_n\}\) 的 OGF。 4. \(\frac{A(x)}{1-x}\) 是 \(\{\sum_{k=0}^{n} a_k\}\)（前缀和）的 OGF。

\textbf{证明}：(1) 显然。(2) \((A \cdot B)(x) = \sum_{n=0}^{\infty} (\sum_{k=0}^{n} a_k b_{n-k}) x^n\)。(3) \(A'(x) = \sum_{n=1}^{\infty} n a_n x^{n-1}\)，\(x A'(x) = \sum_{n=0}^{\infty} n a_n x^n\)。(4) \(\frac{1}{1-x} = \sum_{n=0}^{\infty} x^n\)，由 (2) 即得。\(\blacksquare\)

\textbf{定理 28.12（Newton 广义二项式定理）} 对任意实数 \(\alpha\)： \[(1 + x)^\alpha = \sum_{n=0}^{\infty} \binom{\alpha}{n} x^n\] 其中 \(\binom{\alpha}{n} = \frac{\alpha(\alpha-1)\cdots(\alpha-n+1)}{n!}\)（广义二项式系数）。

\textbf{证明}：由 Taylor 展开，\(f(x) = (1+x)^\alpha\) 在 \(x=0\) 处的 \(n\) 阶导数为 \(f^{(n)}(0) = \alpha(\alpha-1)\cdots(\alpha-n+1)\)，故 Taylor 系数为 \(\binom{\alpha}{n}\)。\(\blacksquare\)

\textbf{命题 28.10（常见数列的生成函数）} 1. \(\{1, 1, 1, \ldots\}\) 的 OGF：\(\frac{1}{1-x}\) 2. \(\{0, 1, 2, 3, \ldots\}\) 的 OGF：\(\frac{x}{(1-x)^2}\) 3. \(\{C_n^k\}_{n=0}^{\infty}\)（固定 \(k\)）的 OGF：\(\frac{x^k}{(1-x)^{k+1}}\) 4. Fibonacci 数列的 OGF：\(\frac{x}{1-x-x^2}\) 5. Catalan 数列的 OGF：\(\frac{1-\sqrt{1-4x}}{2x}\)

\textbf{证明（Fibonacci 数列的 OGF）}：设 \(F(x) = \sum_{n=0}^{\infty} F_n x^n\)。利用 \(F_n = F_{n-1} + F_{n-2}\)（\(n \geq 2\)）： \[F(x) - F_0 - F_1 x = x(F(x) - F_0) + x^2 F(x)\] 代入 \(F_0 = 0, F_1 = 1\)，得 \(F(x) - x = x F(x) + x^2 F(x)\)，故 \(F(x) = \frac{x}{1-x-x^2}\)。\(\blacksquare\)

\textbf{命题 28.11（指数生成函数的乘法）} 设 \(\hat{A}(x)\) 和 \(\hat{B}(x)\) 分别为 \(\{a_n\}\) 和 \(\{b_n\}\) 的 EGF。则 \(\hat{A}(x) \cdot \hat{B}(x)\) 是 \(\{c_n = \sum_{k=0}^{n} \binom{n}{k} a_k b_{n-k}\}\)（二项式卷积）的 EGF。

\textbf{证明}：\(\hat{A}(x)\hat{B}(x) = (\sum_{n=0}^{\infty} a_n \frac{x^n}{n!})(\sum_{m=0}^{\infty} b_m \frac{x^m}{m!}) = \sum_{n=0}^{\infty} (\sum_{k=0}^{n} \frac{a_k}{k!} \frac{b_{n-k}}{(n-k)!}) x^n = \sum_{n=0}^{\infty} (\sum_{k=0}^{n} \binom{n}{k} a_k b_{n-k}) \frac{x^n}{n!}\)。\(\blacksquare\)

\subsection{28.10 特殊计数序列}\label{ux7279ux6b8aux8ba1ux6570ux5e8fux5217}

\textbf{定义 28.8（Catalan 数）} \textbf{Catalan 数} \(C_n\) 定义为满足以下递推关系的数列： \[C_n = \sum_{k=0}^{n-1} C_k C_{n-1-k}, \quad C_0 = 1\]

Catalan 数的闭形式为 \(C_n = \frac{1}{n+1} \binom{2n}{n}\)。

\textbf{定理 28.13（Catalan 数的闭形式）} \(C_n = \frac{1}{n+1}\binom{2n}{n} = \binom{2n}{n} - \binom{2n}{n+1}\)。

\textbf{证明（生成函数法）}：设 \(C(x) = \sum_{n=0}^{\infty} C_n x^n\)。由递推关系，\(C(x) = 1 + x C(x)^2\)。解二次方程 \(x C^2 - C + 1 = 0\)，得 \(C(x) = \frac{1 \pm \sqrt{1-4x}}{2x}\)。由 \(C_0 = 1\)，取负号 \(C(x) = \frac{1-\sqrt{1-4x}}{2x}\)。展开 \(\sqrt{1-4x} = \sum_{n=0}^{\infty} \binom{1/2}{n} (-4x)^n\)，利用 \(\binom{1/2}{n} = \frac{(-1)^{n-1}}{n 2^{2n-1}} \binom{2n-2}{n-1}\)，可得 \(C_n = \frac{1}{n+1}\binom{2n}{n}\)。\(\blacksquare\)

\textbf{命题 28.12（Catalan 数的组合解释）} \(C_n\) 计数了以下对象（共 200 余种组合解释）： 1. \(n\) 对括号的合法匹配方式数 2. \(n+1\) 个叶子的满二叉树个数 3. 凸 \(n+2\) 边形的三角剖分数 4. 从 \((0,0)\) 到 \((n,n)\) 不越过对角线 \(y=x\) 的格路数

\textbf{证明（格路解释）}：从 \((0,0)\) 到 \((n,n)\) 的总格路数为 \(\binom{2n}{n}\)。跨过对角线的路径（首次跨过处 \((k, k+1)\)）通过反射 \((x,y) \mapsto (y-1, x+1)\) 与从 \((0,0)\) 到 \((n-1, n+1)\) 的格路一一对应，共 \(\binom{2n}{n+1}\) 条。故合法路径数为 \(\binom{2n}{n} - \binom{2n}{n+1} = \frac{1}{n+1}\binom{2n}{n} = C_n\)。\(\blacksquare\)

\textbf{定义 28.9（Stirling 数）} - \textbf{第一类 Stirling 数} \(s(n, k)\)（有符号）或 \(c(n, k)\)（无符号）：\(n\) 个元素的排列中恰好有 \(k\) 个轮换的排列数。满足递推 \(c(n, k) = c(n-1, k-1) + (n-1)c(n-1, k)\)。 - \textbf{第二类 Stirling 数} \(S(n, k)\)：将 \(n\) 个不同元素划分为 \(k\) 个非空子集（无序）的方式数。满足递推 \(S(n, k) = S(n-1, k-1) + k S(n-1, k)\)。

\textbf{定理 28.14（Stirling 数的生成函数）} 1. \(\sum_{k=0}^{n} c(n, k) x^k = x(x+1)(x+2)\cdots(x+n-1)\)（上升阶乘） 2. \(\sum_{n=k}^{\infty} S(n, k) \frac{x^n}{n!} = \frac{(e^x - 1)^k}{k!}\)（第二类 Stirling 数的 EGF）

\textbf{证明}：(1) 对 \(n\) 归纳。\(n=1\) 时 \(x\) 成立。假设 \(n-1\) 时成立，则 \(\sum_{k} c(n,k)x^k = \sum_{k} [c(n-1,k-1) + (n-1)c(n-1,k)] x^k = x \sum_{k} c(n-1,k-1)x^{k-1} + (n-1)\sum_{k} c(n-1,k)x^k = (x + n - 1) \prod_{i=0}^{n-2} (x+i) = \prod_{i=0}^{n-1} (x+i)\)。 (2) 固定 \(k\)，\(S(n,k)\) 的 EGF 为 \(\frac{(e^x-1)^k}{k!}\)，因为 \(e^x-1 = \sum_{n=1}^{\infty} \frac{x^n}{n!}\) 是''非空''的生成函数，\(k\) 次幂对应 \(k\) 个非空子集（无序------除以 \(k!\)）。\(\blacksquare\)

\textbf{定义 28.10（Bell 数）} \textbf{Bell 数} \(B_n\) 是将 \(n\) 个不同元素划分为任意个非空子集（无序）的方式总数。即 \(B_n = \sum_{k=0}^{n} S(n, k)\)。\(B_n\) 满足递推 \(B_{n+1} = \sum_{k=0}^{n} \binom{n}{k} B_k\)（\(B_0 = 1\)）。

\textbf{定理 28.15（Bell 数的 EGF）} \(\sum_{n=0}^{\infty} B_n \frac{x^n}{n!} = e^{e^x - 1}\)。

\textbf{证明}：\(\sum_{n=0}^{\infty} B_n \frac{x^n}{n!} = \sum_{n=0}^{\infty} \sum_{k=0}^{n} S(n,k) \frac{x^n}{n!} = \sum_{k=0}^{\infty} \sum_{n=k}^{\infty} S(n,k) \frac{x^n}{n!} = \sum_{k=0}^{\infty} \frac{(e^x-1)^k}{k!} = e^{e^x-1}\)。\(\blacksquare\)

\textbf{定义 28.11（整数分拆）} 正整数 \(n\) 的一个\textbf{分拆}（partition）是将 \(n\) 表示为若干正整数之和的方式（不计顺序）。分拆数 \(p(n)\) 表示 \(n\) 的不同分拆种数。例如 \(p(4) = 5\)：\(4, 3+1, 2+2, 2+1+1, 1+1+1+1\)。

\textbf{定理 28.16（Euler 分拆恒等式）} \(p(n)\) 的普通生成函数为： \[\sum_{n=0}^{\infty} p(n) x^n = \prod_{k=1}^{\infty} \frac{1}{1-x^k}\]

\textbf{证明}：\(\prod_{k=1}^{\infty} \frac{1}{1-x^k} = \prod_{k=1}^{\infty} (1 + x^k + x^{2k} + x^{3k} + \cdots)\)。展开式中 \(x^n\) 的系数等于将 \(n\) 表示为 \(1 \cdot a_1 + 2 \cdot a_2 + 3 \cdot a_3 + \cdots\) 的方式数，其中 \(a_k \geq 0\) 表示 \(k\) 在分拆中出现的次数。这正是 \(p(n)\) 的定义。\(\blacksquare\)

\textbf{定理 28.17（Euler 五角数定理）} \[\prod_{k=1}^{\infty} (1 - x^k) = \sum_{m=-\infty}^{\infty} (-1)^m x^{m(3m-1)/2} = 1 + \sum_{m=1}^{\infty} (-1)^m (x^{m(3m-1)/2} + x^{m(3m+1)/2})\]

\textbf{证明}：这一深刻恒等式的证明由 Euler 通过组合论证和 Jacobi 三重积恒等式给出。五角数定理是 Hardy-Ramanujan 分拆公式的推导基础。\(\blacksquare\)

\textbf{推论 28.1（Euler 递推）} \(p(n) = \sum_{m \neq 0} (-1)^{m-1} p(n - m(3m-1)/2)\)，其中 \(p(0) = 1\) 且 \(p(k) = 0\) 当 \(k < 0\)。

\textbf{证明}：由 \(\sum p(n) x^n \cdot \prod (1-x^k) = 1\) 和五角数定理，比较系数即得。\(\blacksquare\)

\subsection{28.11 Burnside 引理与 Polya 计数}\label{burnside-ux5f15ux7406ux4e0e-polya-ux8ba1ux6570}

\begin{quote}
\textbf{直觉与动机} 许多计数问题涉及对称性------例如，在考虑旋转和反射等价时，如何计数不同的项链？Burnside 引理（也称 Cauchy-Frobenius 引理）通过群作用的轨道计数给出了精确答案。
\end{quote}

\textbf{定理 28.18（Burnside 引理）} 设有限群 \(G\) 作用于有限集 \(X\)。令 \(X^g = \{x \in X : g \cdot x = x\}\)（\(g\) 的不动点集）。则 \(G\) 在 \(X\) 上的轨道数为： \[\frac{1}{|G|} \sum_{g \in G} |X^g|\]

\textbf{证明}：双重计数 \(\sum_{g \in G} |X^g| = \sum_{x \in X} |G_x|\)（其中 \(G_x = \{g \in G : g \cdot x = x\}\) 为 \(x\) 的稳定子群）。由轨道-稳定子定理，\(|G| = |G \cdot x| \cdot |G_x|\)，故 \(|G_x| = |G|/|G \cdot x|\)。因此 \(\sum_{g} |X^g| = \sum_{x} |G|/|G \cdot x| = |G| \sum_{\text{轨道 } \mathcal{O}} \sum_{x \in \mathcal{O}} 1/|\mathcal{O}| = |G| \cdot (\text{轨道数})\)。\(\blacksquare\)

\textbf{定理 28.19（Polya 计数定理）} 设 \(G\) 是 \(n\) 个对象的置换群，有 \(c\) 种颜色可用。对每个 \(g \in G\)，记 \(c(g)\) 为 \(g\) 的轮换分解中的轮换数。则 \(G\) 作用下 \(n\) 个对象的染色方案的轨道数为： \[\frac{1}{|G|} \sum_{g \in G} c^{c(g)}\]

\textbf{证明}：这是 Burnside 引理的直接应用。\(X\) 为 \(c^n\) 种染色方案，\(g\) 的不动点集 \(X^g\) 由那些在每个 \(g\) 的轮换内颜色一致的染色方案组成，每个轮换独立选择颜色，故 \(|X^g| = c^{c(g)}\)。\(\blacksquare\)

\textbf{命题 28.13（项链计数）} 用 \(c\) 种颜色的珠子串成 \(n\) 颗珠子的项链（旋转视为相同），不同的项链数为： \[\frac{1}{n} \sum_{d \mid n} \varphi(d) \, c^{n/d}\] 其中 \(\varphi\) 为 Euler 函数。若进一步考虑翻转对称（二面体群），则不同的项链数为： \[\frac{1}{2n} \sum_{d \mid n} \varphi(d) \, c^{n/d} + \begin{cases} \frac{1}{2} c^{n/2} (c+1)/2, & n \text{ 为偶数} \\ \frac{1}{2} c^{(n+1)/2}, & n \text{ 为奇数} \end{cases}\]

\textbf{证明}：循环群 \(C_n\) 的生成元为旋转 \(2\pi k/n\)（\(k=0,1,\ldots,n-1\)）。旋转 \(k\) 步的轮换数为 \(\gcd(n,k)\)，不动的染色方案数为 \(c^{\gcd(n,k)}\)。因此 \(\frac{1}{n} \sum_{k=0}^{n-1} c^{\gcd(n,k)} = \frac{1}{n} \sum_{d \mid n} \varphi(d) c^{n/d}\)。二面体群情形需额外考虑 \(n\) 个反射。\(\blacksquare\)

\hrulefill

\newpage

\chapter{04-离散数学/01-组合数学/02-生成函数与图论.md}\label{ux79bbux6563ux6570ux5b6601-ux7ec4ux5408ux6570ux5b6602-ux751fux6210ux51fdux6570ux4e0eux56feux8bba.md}

\section{第81章：生成函数}\label{ux7b2c81ux7ae0ux751fux6210ux51fdux6570}

生成函数是将离散的计数序列编码为形式幂级数的分析工具，其核心思想源自 Euler 对分拆函数的研究（1748 年）和 Laplace 对其的系统化命名（1812 年）。这一方法将组合计数问题转化为函数方程的代数操作：递推关系化为有理函数的极点和部分分式，序列的卷积化为生成函数的乘积，组合结构的复合化为形式的代入。生成函数桥接了组合学与复分析（围道积分方法、奇点渐近分析）、代数学（形式幂级数环、微分方程）和图论（色多项式、Tutte 多项式），是整个离散数学中最为统一的方法论框架之一。本章从普通生成函数的基本代数运算入手，系统建立指数生成函数、线性递推的求解和 Catalan 数的生成函数推导，为后续 Pólya 计数理论（第 38 章）和符号方法（V4 Ch12）提供生成函数方法论基础。其中 Dirichlet 生成函数还将与解析数论（V16-17）中的 Riemann zeta 函数和 Dirichlet 特征直接关联。

\subsection{37.1 普通生成函数}\label{ux666eux901aux751fux6210ux51fdux6570}

\textbf{定义 37.1}（普通生成函数）：设 \(\{a_n\}_{n=0}^\infty\) 为数域 \(\mathbb{C}\) 中的序列。其\textbf{普通生成函数}（Ordinary Generating Function, OGF）为形式幂级数 \[F(x) = \sum_{n=0}^\infty a_n x^n \in \mathbb{C}[[x]]\] 若该级数在某开圆盘 \(|x| < R\) 内收敛，其中收敛半径 \(R = 1 / \limsup_{n\to\infty} \sqrt[n]{|a_n|}\)，则 \(F(x)\) 同时定义了一个全纯函数。若 \(R = 0\)，则 \(F(x)\) 仅为形式幂级数；若 \(R = \infty\)，则 \(F(x)\) 为整函数。

\textbf{定理 37.1}（OGF 加法定理）：设 \(\{a_n\}\) 和 \(\{b_n\}\) 的 OGF 分别为 \(A(x)\) 和 \(B(x)\)，则序列 \(\{a_n + b_n\}\) 的 OGF 为 \(A(x) + B(x)\)。

\textbf{证明}：由定义 \(\sum_{n\geq 0}(a_n + b_n) x^n = \sum_{n\geq 0} a_n x^n + \sum_{n\geq 0} b_n x^n = A(x) + B(x)\)。\(\blacksquare\)

\textbf{定理 37.2}（OGF 乘法定理 / Cauchy 卷积）：序列 \(\{a_n\}\) 和 \(\{b_n\}\) 的 OGF 之积 \(A(x)B(x) = \sum_{n\geq 0} c_n x^n\)，其中 \[c_n = \sum_{k=0}^n a_k b_{n-k}\] 称为 \(\{a_n\}\) 与 \(\{b_n\}\) 的 \textbf{Cauchy 卷积}。

\textbf{证明}：由级数乘法 \((\sum_{i\geq 0} a_i x^i)(\sum_{j\geq 0} b_j x^j) = \sum_{n\geq 0} (\sum_{i+j=n} a_i b_j) x^n\)。令 \(k = i\) 即得 \(c_n = \sum_{k=0}^n a_k b_{n-k}\)。\(\blacksquare\)

\textbf{定理 37.3}（OGF 的复合）：若 \(B(0) = 0\)，则形式幂级数的复合 \(A(B(x)) = \sum_{n\geq 0} a_n (B(x))^n\) 是良定义的。组合意义上，若 \(A(x)\) 和 \(B(x)\) 分别对应组合类 \(\mathcal{A}\) 与 \(\mathcal{B}\) 的 OGF，且 \(\mathcal{B}\) 无大小为 \(0\) 的元素，则 \(A(B(x))\) 对应 \(\mathcal{A}\) 中元素由 \(\mathcal{B}\) 中元素代换而成的结构。

\textbf{证明}：\(B(0) = 0\) 保证 \(B(x) = \sum_{k \geq 1} b_k x^k\) 的最低次项至少为 \(x\)。因此 \((B(x))^n\) 的最低次项为 \(x^n\)，展开 \(A(B(x)) = \sum_{n \geq 0} a_n (B(x))^n\) 时，\(x^m\) 的系数仅涉及 \(n \leq m\) 的有限项，故级数在 \(\mathbb{C}[[x]]\) 中良定义。组合意义：将 \(\mathcal{A}\) 中大小为 \(n\) 的结构中的每个''原子''替换为 \(\mathcal{B}\) 中一个结构，得到大小为所有 \(\mathcal{B}\)-结构大小之和的新结构，其 OGF 恰为 \(A(B(x))\)。\(\blacksquare\)

\subsection{37.2 指数生成函数}\label{ux6307ux6570ux751fux6210ux51fdux6570}

\textbf{定义 37.2}（指数生成函数）：序列 \(\{a_n\}\) 的\textbf{指数生成函数}（Exponential Generating Function, EGF）定义为 \[\hat{F}(x) = \sum_{n=0}^\infty a_n \frac{x^n}{n!}\]

EGF 的乘积对应于\textbf{二项卷积}，使其适用于\textbf{标记结构}（labeled structures）的计数。

\textbf{定理 37.4}（EGF 乘法公式）：设 \(\hat{A}(x), \hat{B}(x)\) 分别为 \(\{a_n\}\) 和 \(\{b_n\}\) 的 EGF，则 \[\hat{A}(x)\hat{B}(x) = \sum_{n=0}^\infty \left(\sum_{k=0}^n \binom{n}{k} a_k b_{n-k}\right) \frac{x^n}{n!}\]

\textbf{证明}：\(\hat{A}(x)\hat{B}(x) = (\sum_{k} a_k \frac{x^k}{k!})(\sum_{j} b_j \frac{x^j}{j!}) = \sum_{n} (\sum_{k=0}^n \frac{a_k b_{n-k}}{k!(n-k)!}) x^n = \sum_{n} (\sum_{k=0}^n \binom{n}{k} a_k b_{n-k}) \frac{x^n}{n!}\)。\(\blacksquare\)

组合诠释：在标记组合的符号方法中，若 \(\mathcal{A}\) 和 \(\mathcal{B}\) 为标记组合类，它们的\textbf{标记乘积} \(\mathcal{A} * \mathcal{B}\) 的 EGF 恰为 \(\hat{A}(x)\hat{B}(x)\)。

\subsection{37.3 线性递推的生成函数解法}\label{ux7ebfux6027ux9012ux63a8ux7684ux751fux6210ux51fdux6570ux89e3ux6cd5}

\textbf{定理 37.5}（常系数线性齐次递推）：设序列 \(\{a_n\}\) 满足递推 \(a_{n+k} = c_1 a_{n+k-1} + \cdots + c_k a_n\)（\(n \geq 0\)），\(c_i \in \mathbb{C}\)，初值 \(a_0, \ldots, a_{k-1}\) 给定。定义特征多项式（以 \(x\) 为变量） \[P(x) = 1 - c_1 x - \cdots - c_k x^k\] 则 \(\{a_n\}\) 的 OGF 为 \(F(x) = Q(x) / P(x)\)，其中 \(Q(x)\) 是由初值确定的次数小于 \(k\) 的多项式。

\textbf{证明}：计算 \(P(x)F(x) = (1 - \sum_{i=1}^k c_i x^i)(\sum_{n\geq 0} a_n x^n) = \sum_{n\geq 0} a_n x^n - \sum_{i=1}^k c_i \sum_{n\geq 0} a_n x^{n+i}\)。将第二项中 \(n\) 替换为 \(n-i\) 得 \(\sum_{n\geq 0} (a_n - \sum_{i=1}^k c_i a_{n-i}) x^n\)，其中约定 \(a_j=0\) 当 \(j<0\)。对 \(n \geq k\)，递推关系表明括号内为 \(0\)。故 \(P(x)F(x)\) 为次数小于 \(k\) 的多项式 \(Q(x)\)。\(\blacksquare\)

\textbf{定理 37.6}（部分分式与闭式）：若特征多项式有因式分解 \(P(x) = \prod_{i=1}^r (1-\lambda_i x)^{m_i}\)（\(\lambda_i\) 互异），则 \[a_n = \sum_{i=1}^r P_{m_i-1}^{(i)}(n) \lambda_i^n\] 其中 \(P_{d}^{(i)}(n)\) 是次数不超过 \(d\) 的多项式。当 \(m_i = 1\)（即单根）时，简化为指数项 \(C_i \lambda_i^n\)。

\textbf{证明}：由定理 37.5，\(F(x) = Q(x)/P(x)\)。将 \(P(x)\) 分解为 \(\prod_{i=1}^r (1-\lambda_i x)^{m_i}\)，则由部分分式分解，\(F(x) = \sum_{i=1}^r \sum_{j=1}^{m_i} \frac{c_{ij}}{(1-\lambda_i x)^j}\)。利用 \((1-\lambda x)^{-j} = \sum_{n \geq 0} \binom{n+j-1}{j-1} \lambda^n x^n\)，提取 \(x^n\) 的系数得 \(a_n = \sum_{i=1}^r \sum_{j=1}^{m_i} c_{ij} \binom{n+j-1}{j-1} \lambda_i^n\)。对固定的 \(i\)，\(\sum_{j=1}^{m_i} c_{ij} \binom{n+j-1}{j-1}\) 是 \(n\) 的次数不超过 \(m_i-1\) 的多项式，记为 \(P_{m_i-1}^{(i)}(n)\)，故 \(a_n = \sum_i P_{m_i-1}^{(i)}(n) \lambda_i^n\)。\(\blacksquare\)

\subsection{37.4 Catalan 数的生成函数}\label{catalan-ux6570ux7684ux751fux6210ux51fdux6570}

\textbf{定义 37.3}（Catalan 数）：第 \(n\) 个 \textbf{Catalan 数} \(C_n\) 满足 \(C_0 = 1\) 及递推 \(C_{n+1} = \sum_{k=0}^n C_k C_{n-k}\)（\(n \geq 0\)）。\(C_n\) 计数 \(n\) 对括号的正确配对方式、\(n+1\) 个叶子的二叉树数目、\(\{1,\ldots,2n\}\) 的 \(n\) 个不交叉弧的匹配等。

\textbf{定理 37.7}（Catalan 生成函数与闭式）：Catalan 数的 OGF \(C(x) = \sum_{n\geq 0} C_n x^n\) 满足 \[C(x) = \frac{1 - \sqrt{1 - 4x}}{2x}\] 且闭式为 \(C_n = \frac{1}{n+1} \binom{2n}{n}\)。

\textbf{证明}：由递推 \(C_{n+1} = \sum_{k=0}^n C_k C_{n-k}\) 得 \(C(x) - 1 = \sum_{n\geq 0} C_{n+1} x^{n+1} = x \sum_{n\geq 0} (\sum_{k=0}^n C_k C_{n-k}) x^n = x C(x)^2\)。故 \(x C(x)^2 - C(x) + 1 = 0\)。解二次方程得 \(C(x) = (1 \pm \sqrt{1-4x})/(2x)\)。由 \(C(0) = C_0 = 1\)，取负号：\(C(x) = (1 - \sqrt{1-4x})/(2x)\)。将 \(\sqrt{1-4x}\) 按二项级数展开： \[\sqrt{1-4x} = \sum_{n=0}^\infty \binom{1/2}{n}(-4x)^n = 1 - 2\sum_{n=1}^\infty \frac{1}{n}\binom{2n-2}{n-1} x^n\] 代入得 \(C(x) = \frac{1}{x} \sum_{n=1}^\infty \frac{1}{n}\binom{2n-2}{n-1} x^n = \sum_{n=0}^\infty \frac{1}{n+1} \binom{2n}{n} x^n\)。故 \(C_n = \frac{1}{n+1}\binom{2n}{n}\)。\(\blacksquare\)

\subsection{37.5 分拆数的生成函数}\label{ux5206ux62c6ux6570ux7684ux751fux6210ux51fdux6570}

\textbf{定义 37.4}（分拆函数）：\(p(n)\) 为将正整数 \(n\) 写成正整数之和的不同方式数（不计顺序）。约定 \(p(0) = 1\)。

\textbf{定理 37.8}（Euler 分拆生成函数）： \[\sum_{n=0}^\infty p(n) x^n = \prod_{k=1}^\infty \frac{1}{1 - x^k}\]

\textbf{证明}：右端每个因子展开为几何级数 \(\frac{1}{1-x^k} = 1 + x^k + x^{2k} + \cdots\)。从第 \(k\) 个因子选择 \(x^{m_k k}\)（\(m_k \geq 0\)），所有因子的乘积为 \(x^{\sum_k m_k k}\)。指数 \(\sum_k m_k k\) 恰为将某个整数分解为部分 \(k\) 重复 \(m_k\) 次的一种方式，故 \(x^n\) 的系数恰为所有可能序列 \(\{m_k\}\) 满足 \(\sum m_k k = n\) 的数目------即 \(p(n)\)。\(\blacksquare\)

Euler 的分拆生成函数将正整数分拆的计数问题转化为无穷乘积的幂级数展开，其结构之简洁令人惊叹：\(\prod_{k=1}^\infty (1-x^k)^{-1}\) 的 \(x^n\) 系数恰好就是 \(p(n)\)。然而，这个无穷乘积的倒数 \(\prod_{k=1}^\infty (1-x^k)\) 自身也是一个极为重要的生成函数，其展开蕴含着深刻的数论和组合结构。Euler 在 18 世纪发现了该乘积的一个非凡恒等式------展开中非零系数仅出现在五角数指标 \(k(3k-1)/2\) 上，且符号交替出现。

\textbf{定理 37.9}（Euler 五角数定理）： \[\prod_{n=1}^\infty (1 - x^n) = \sum_{k=-\infty}^\infty (-1)^k x^{k(3k-1)/2}\] 该恒等式给出了 \(p(n)\) 的高效递推公式。

\textbf{证明}：由定理 37.8，\(\prod_{n=1}^\infty (1-x^n) = (\sum_{n=0}^\infty p(n) x^n)^{-1}\)。Euler 五角数定理的证明在组合学中由 Franklin 对合给出：将 \((q;q)_\infty\) 展开，\(q^N\) 的系数为将 \(N\) 分拆为偶数个不同部分的方式数减去奇数个不同部分的方式数。在 Ferrers 图上构造符号翻转对合，其不动点恰对应五角数 \(k(3k\pm 1)/2\)，符号为 \((-1)^k\)。因此 \(\prod_{n=1}^\infty (1-x^n) = \sum_{k=-\infty}^\infty (-1)^k x^{k(3k-1)/2}\)。代入 \(p(n)\) 的生成函数得 \(p(n)\) 的递推公式。\(\blacksquare\)

\subsection{37.6 Dirichlet生成函数与Dirichlet卷积}\label{dirichletux751fux6210ux51fdux6570ux4e0edirichletux5377ux79ef}

\textbf{定义 37.5}（Dirichlet生成函数）：设 \(\{a_n\}_{n=1}^\infty \subseteq \mathbb{C}\)。其\textbf{Dirichlet生成函数}（DGF）为 Dirichlet 级数 \[D(s) = \sum_{n=1}^\infty \frac{a_n}{n^s}\] 若存在 \(\sigma_0\) 使级数在 \(\Re(s) > \sigma_0\) 绝对收敛，则 \(D(s)\) 定义了半平面内的全纯函数。

\textbf{定理 37.10}（DGF 乘法与 Dirichlet 卷积）：设 \(A(s), B(s)\) 分别为 \(\{a_n\}, \{b_n\}\) 的 DGF，则 \[A(s)B(s) = \sum_{n=1}^\infty \frac{(a * b)_n}{n^s},\quad (a * b)_n = \sum_{d \mid n} a_d \cdot b_{n/d}\] 其中 \(a * b\) 称为 \textbf{Dirichlet 卷积}，构成 \(\mathbb{C}\) 上数论函数的交换结合代数（单位元 \(\varepsilon(n) = [n = 1]\)）。

\textbf{证明}：\(A(s)B(s) = \sum_{i,j \geq 1} a_i b_j (ij)^{-s}\)。令 \(n = ij\)，对固定 \(n\)，满足 \(d \mid n\) 的因子对 \((d, n/d)\) 贡献 \(\frac{a_d b_{n/d}}{n^s}\)。故系数为 \(\sum_{d \mid n} a_d b_{n/d}\)。\(\blacksquare\)

\textbf{注记 37.2}：Riemann zeta 函数 \(\zeta(s) = \sum_{n=1}^\infty 1/n^s\)（\(a_n \equiv 1\)）是 DGF 的最重要例子。积性函数（\(a_{mn} = a_m a_n\) 当 \(\gcd(m,n)=1\)）的 DGF 具有 Euler 乘积： \[D(s) = \prod_{p} \left( \sum_{k=0}^\infty \frac{a_{p^k}}{p^{ks}} \right)\] 完全积性函数简化为 \(\prod_p (1 - a_p p^{-s})^{-1}\)。DGF 将 Dirichlet 卷积代数化------这是解析数论的基本工具。

\section{第82章：Pólya 计数与图多项式}\label{ux7b2c82ux7ae0puxf3lya-ux8ba1ux6570ux4e0eux56feux591aux9879ux5f0f}

Pólya 计数定理（1937 年）是组合计数中最优美的方法之一，它将群论与组合计数融为一体：当对称群作用于着色对象上时，不等价着色类的计数归结为对称群的循环指标多项式。这一定理的历史脉络可追溯至 Cauchy-Frobenius 引理（即 Burnside 引理），Pólya 的贡献在于将逐元计数的 Burnside 引理推广为按循环结构加权的形式，并通过循环指标统一表达。Pólya 计数是生成函数思想在群对称性约束下的自然推广------第 37 章的生成函数直接对应平凡的对称群，而 Pólya 定理处理非平凡的群作用。本章后半部分引入图多项式（色多项式与 Tutte 多项式），它们是图的代数不变量：色多项式 \(P_G(k)\) 将图论问题转化为代数问题，Tutte 多项式 \(T_G(x, y)\) 则统一了色多项式、流多项式和 Jones 多项式（纽结理论），构成了拓扑图论和统计力学（Potts 模型）之间的深层桥梁。

\subsection{38.1 Pólya 计数定理}\label{puxf3lya-ux8ba1ux6570ux5b9aux7406}

Pólya 计数定理将组合计数与群论融为一体，其出发点是对称群作用于着色对象时的轨道计数问题。定理的核心在于循环指标多项式------一个按群元素循环结构加权求和的生成函数------它统一了 Burnside 引理的逐元素计数，将其推广为按颜色数加权的方式。

\textbf{定义 38.1}（群作用与轨道）：设有限群 \(G\) 作用于有限集 \(X\)。对 \(x \in X\)，轨道为 \(\operatorname{Orb}(x) = \{g x : g \in G\}\)；稳定子群为 \(G_x = \{g \in G : g x = x\}\)。轨道-稳定子定理给出 \(|G| = |\operatorname{Orb}(x)| \cdot |G_x|\)。

轨道计数的基础公式是 Burnside 引理（实为 Cauchy-Frobenius 引理），它通过群元素的不动点平均来计数轨道数。这一引理虽然深刻，但仅给出了轨道的总数，而未区分轨道中着色配置的具体颜色分布。Pólya 的贡献在于将这一公式推广为生成函数形式。

\textbf{定理 38.1}（Burnside 引理）：\(G\) 作用于 \(X\) 的轨道数 \(N\) 满足 \[N = \frac{1}{|G|} \sum_{g \in G} |\operatorname{Fix}(g)|\] 其中 \(\operatorname{Fix}(g) = \{x \in X : g x = x\}\)。

\textbf{证明}：对集合 \(\{(g, x) \in G \times X : g x = x\}\) 进行双重计数。按 \(g\) 求和得 \(\sum_{g} |\operatorname{Fix}(g)|\)；按 \(x\) 求和得 \(\sum_{x} |G_x|\)。对每个轨道 \(\mathcal{O}\)，\(\sum_{x \in \mathcal{O}} |G_x| = |\mathcal{O}| \cdot (|G|/|\mathcal{O}|) = |G|\)。故 \(\sum_{x \in X} |G_x| = N \cdot |G|\)。因此 \(N = \frac{1}{|G|} \sum_g |\operatorname{Fix}(g)|\)。\(\blacksquare\)

\textbf{定义 38.2}（循环指标）：设 \(G \leq S_n\)。\(g \in G\) 的\textbf{循环类型}为 \((c_1(g), \ldots, c_n(g))\)，其中 \(c_k(g)\) 为 \(g\) 的循环分解中 \(k\)-循环的个数。\(G\) 的\textbf{循环指标多项式}为 \[Z_G(t_1, \ldots, t_n) = \frac{1}{|G|} \sum_{g \in G} \prod_{k=1}^n t_k^{c_k(g)}\]

循环指标多项式的威力在于：一旦计算出 \(Z_G\)，只需简单地代入 \(t_k = m\)（对任意 \(k\)）即可得到使用 \(m\) 种颜色时不等价着色类的总数。更为一般地，若为每种颜色分配权重，\(Z_G\) 还能提供按颜色分布加权的精确计数，这在化学同分异构体计数中有直接应用。

\textbf{定理 38.2}（Pólya 计数定理）：设有 \(m\) 种颜色。用这些颜色对 \(X\) 的元素着色的 \(G\)-不等价着色数为 \[Z_G(m, m, \ldots, m) = \frac{1}{|G|} \sum_{g \in G} m^{c(g)}\] 其中 \(c(g) = \sum_{k=1}^n c_k(g)\) 为 \(g\) 的循环总数。

\textbf{证明}：着色对应于函数 \(f : X \to \{1,\ldots,m\}\)。\(G\) 通过 \((g \cdot f)(x) = f(g^{-1} x)\) 作用于着色集合。\(f\) 被 \(g\) 固定 \(\iff\) \(f\) 在 \(g\) 的每个循环上取常值，故 \(|\operatorname{Fix}(g)| = m^{c(g)}\)。由 Burnside 引理即得。\(\blacksquare\)

\subsection{38.2 图多项式}\label{ux56feux591aux9879ux5f0f}

\textbf{定义 38.3}（色多项式）：图 \(G = (V, E)\) 的\textbf{色多项式} \(P_G(k)\) 为用至多 \(k\) 种颜色对 \(G\) 的顶点进行正常着色（相邻顶点不同色）的方式数。

\textbf{定理 38.3}（删边-缩边递推）：对任意边 \(e \in E\)， \[P_G(k) = P_{G-e}(k) - P_{G/e}(k)\]

\textbf{证明}：设 \(e = uv\)。\(G-e\) 的 \(k\)-着色的集合等于 \(\{G\) 的正常着色\(\} \cup \{uv\) 同色的 \(G-e\) 着色\(\}\)。使 \(u, v\) 同色的 \(G-e\) 着色与 \(G/e\)（合并 \(u, v\) 为一个顶点）的 \(k\)-着色一一对应。故 \(P_{G-e}(k) = P_G(k) + P_{G/e}(k)\)。移项即得。\(\blacksquare\)

\textbf{推论 38.4}：\(P_G(k)\) 是次数为 \(|V|\) 的整系数多项式，首项为 \(k^{|V|}\)，系数正负交替，常数项为 \(0\)。

\textbf{证明}：对 \(|V| + |E|\) 归纳。基底：无边图 \(P_G(k) = k^{|V|}\) 满足所有性质。归纳步：由删边-缩边公式 \(P_G(k) = P_{G-e}(k) - P_{G/e}(k)\)，\(P_{G-e}\) 为 \(|V|\) 次多项式，\(P_{G/e}\) 为 \(|V|-1\) 次多项式。故 \(P_G\) 为 \(|V|\) 次多项式，首项 \(k^{|V|}\) 来自 \(P_{G-e}\)。系数正负交替性由归纳假设：\(P_{G-e}\) 的系数符号交替且 \(P_{G/e}\) 的系数符号交替，相减保持符号交替。常数项 \(P_G(0) = 0\) 因为 \(k=0\) 时无着色。\(\blacksquare\)

\textbf{定义 38.4}（Tutte 多项式）：图 \(G\) 的 \textbf{Tutte 多项式} \(T_G(x, y)\) 通过以下递推唯一定义： 1. 若 \(G\) 无边，则 \(T_G(x, y) = 1\) 2. 若 \(e\) 为桥，则 \(T_G(x, y) = x \cdot T_{G/e}(x, y)\) 3. 若 \(e\) 为环，则 \(T_G(x, y) = y \cdot T_{G-e}(x, y)\) 4. 否则（\(e\) 非桥非环），\(T_G(x, y) = T_{G-e}(x, y) + T_{G/e}(x, y)\)

\textbf{定理 38.5}（Tutte 多项式与色多项式的联系）：设 \(c(G)\) 为 \(G\) 的连通分支数，则 \[P_G(k) = (-1)^{|V|-c(G)} k^{c(G)} T_G(1-k, 0)\]

\textbf{证明}：由归纳法及删边-缩边公式可验证该关系式满足两多项式的递推归约。当 \(x = 1-k\) 且 \(y = 0\) 时，Tutte 的递归规则约化为色多项式的递归规则：桥缩并贡献因子 \((1-k)\)，环删除贡献因子 \(0\)，一般边分裂为删边与缩边之差。基底情形（无边图）两侧一致。\(\blacksquare\)

\section{第83章：图论基础}\label{ux7b2c83ux7ae0ux56feux8bbaux57faux7840}

图论是离散数学最大的分支之一，其起源可追溯至 Euler 对 Königsberg 七桥问题的研究（1736 年），这是数学史上第一篇图论论文。两百年间，图论从趣味数学发展为具有统一公理和丰富理论的数学学科，其推动力来自四个方向：化学中的分子结构图（Cayley 对树的计数）、电路理论（Kirchhoff 的矩阵树定理）、地图染色问题（Heawood 和 Appel-Haken 对四色定理的证明）以及计算机科学中的算法设计（最短路径、网络流）。图论在数学内部也与拓扑学（平面图的 Kuratowski 表征）、线性代数（谱图论、Laplace 矩阵）和组合设计（拉丁方 = 正则二分图的完美匹配分解）密切交织。本章建立图的形式化定义与矩阵表示（邻接矩阵、Laplace 矩阵），引入连通性、Menger 定理、图着色（Brooks 定理）和平面图理论（Euler 公式、Kuratowski 定理），为第 39 章（匹配）、第 40 章（图多项式）和第 128 章（图论补充）提供理论基础。

\subsection{39.1 图的正式定义与矩阵表示}\label{ux56feux7684ux6b63ux5f0fux5b9aux4e49ux4e0eux77e9ux9635ux8868ux793a}

\textbf{定义 39.1}（图与矩阵）：一个\textbf{简单图} \(G = (V, E)\) 由非空有限顶点集 \(V\) 和边集 \(E \subseteq \binom{V}{2}\) 构成，记 \(n = |V|\)，\(m = |E|\)。\textbf{邻接矩阵} \(A(G) = (a_{ij})\)，\(a_{ij} = 1_{v_i v_j \in E}\)。\textbf{度矩阵} \(D(G) = \operatorname{diag}(\deg(v_1), \ldots, \deg(v_n))\)，\(\deg(v) = |\{u : uv \in E\}|\)。\textbf{Laplace 矩阵} \(L(G) = D - A\) 对称半正定，零特征值重数等于连通分支数。

图的矩阵表示构成了图论与线性代数之间的桥梁，其历史可追溯至 Kirchhoff（1847）对电网的研究。邻接矩阵 \(A(G)\) 的特征值称为图的\textbf{谱}，其研究构成了谱图论的核心。具体而言，\(A(G)\) 的第 \(k\) 次幂的 \((i,j)\) 元等于从 \(v_i\) 到 \(v_j\) 长度为 \(k\) 的步数；\(A(G)\) 的最大特征值（Perron-Frobenius 特征值）介于 \(\Delta(G)\) 和 \(\max\{\deg(v)\}\) 之间，且等于 \(\Delta(G)\) 当且仅当 \(G\) 是正则图。Laplace 矩阵 \(L(G)\) 具有更深的组合意义：其二次型 \(\mathbf{x}^T L \mathbf{x} = \sum_{uv \in E} (x_u - x_v)^2\) 度量了顶点赋值的光滑程度；\(L(G)\) 的第二小特征值（代数连通度）量化了图的连通鲁棒性。\textbf{关联矩阵} \(B(G)\)（\(n \times m\) 矩阵，\(B_{v,e} = 1\) 若 \(v\) 与 \(e\) 关联，否则为 \(0\)）满足 \(L(G) = B B^T\)，这一分解是矩阵树定理的证明基础。

\textbf{定义 39.2}（图同构与自同构）：\(G\) 与 \(H\) \textbf{同构}（\(G \cong H\)）若存在双射 \(\varphi : V(G) \to V(H)\) 使 \(uv \in E(G) \iff \varphi(u)\varphi(v) \in E(H)\)。\(G\) 的\textbf{自同构群} \(\operatorname{Aut}(G)\) 为全体自同构关于复合构成的群，忠实作用于 \(V\) 故 \(\operatorname{Aut}(G) \leq S_n\)。例：\(\operatorname{Aut}(K_n) = S_n\)，\(\operatorname{Aut}(C_n) = D_{2n}\)。

图同构问题（判定两个图是否同构）在计算复杂性理论中占有特殊地位：它既不被认为是 NP 完全的，也不被认为是多项式时间可解的，是著名的''NP 中间''问题候选。Babai（2015）证明了图同构问题可在拟多项式时间 \(n^{O(\log n)}\) 内求解，这是该领域四十年来最重要的突破。自同构群 \(\operatorname{Aut}(G)\) 的阶数可通过 Pólya 计数和 Burnside 引理与图的对称性计数相关联，而 Frucht 定理（1939）断言每个有限群可表示为某个图的自同构群，揭示了群论与图论之间的深刻联系。

\subsection{39.2 连通性与 Menger 定理}\label{ux8fdeux901aux6027ux4e0e-menger-ux5b9aux7406}

图的连通性直观上描述''图中有多少条独立路径在连接各顶点''，其正式量化依赖于顶点割和边割的概念。顶点连通度 \(\kappa(G)\) 和边连通度 \(\kappa'(G)\) 分别度量了破坏图的连通性所需移除的最少顶点数和最少数边数。Menger 定理（1927 年）以对偶的方式揭示了这两个量与不相交路径数之间的本质联系：它是图论中最早的对偶定理之一，也是网络流理论的先声。

\textbf{定义 39.3}（顶点连通度与边连通度）：图 \(G\) 的\textbf{顶点连通度} \(\kappa(G)\) 是最小的顶点子集大小，移除后使图不连通或退化为单点。\textbf{边连通度} \(\kappa'(G)\) 是最小的边子集大小，移除后使图不连通。

\textbf{定理 39.1}（Menger 定理------点形式）：设 \(s, t\) 不相邻。则 \(s, t\) 间内部不相交路径的最大数等于分离 \(s, t\) 的最小顶点数。

\textbf{证明}（框架）：设内部不相交路径最大数为 \(p\)，最小分离集大小为 \(q\)。\(p \leq q\) 平凡。\(p \geq q\) 对边数极小反例归纳：设 \(S\) 为极小 \(s\)-\(t\) 分离集。若存在 \(x \in S\) 与 \(s\) 和 \(t\) 均相邻，删 \(x\) 归纳。否则沿 \(S\) 将图分解为 \(s\)-侧与 \(t\)-侧，分别应用归纳假设，合并后即得 \(p \geq q\)。\(\blacksquare\)

\textbf{定理 39.2}（Menger 定理------边形式）：\(s, t\) 间边不交路径最大数等于分离 \(s, t\) 的最小边数。此即容量全为 1 的网络中最大流 = 最小割定理。\textbf{推论}：\(\kappa(G) \geq k \iff\) 任意两顶点间有 \(k\) 条内部不相交路径；\(\lambda(G) \geq k \iff\) 任意两顶点间有 \(k\) 条边不交路径。

\textbf{证明}：将图 \(G\) 转化为有向网络：每条边 \(uv\) 替换为两条有向边 \((u,v)\) 和 \((v,u)\)，容量均为 \(1\)。\(s\)-\(t\) 最大流的值等于 \(s\)-\(t\) 边不交路径的最大数（每条路径贡献单位流）。由 Ford-Fulkerson 最大流最小割定理，最大流等于最小 \(s\)-\(t\) 割的容量，即分离 \(s\) 和 \(t\) 所需删除的最小边数。推论由 Menger 定理对所有顶点对取最值得出。\(\blacksquare\)

\subsection{39.3 图着色与 Brooks 定理}\label{ux56feux7740ux8272ux4e0e-brooks-ux5b9aux7406}

\textbf{定义 39.4}（色数）：\(G\) 的\textbf{正常 \(k\)-着色}为 \(c : V \to [k]\) 使 \(uv \in E \Rightarrow c(u) \neq c(v)\)。\textbf{色数} \(\chi(G)\) 为最小的 \(k\)。

图着色问题起源于 1852 年 Francis Guthrie 提出的四色猜想------任何平面地图仅需四种颜色即可使相邻区域不同色。这一猜想经历了超过一个世纪的探索，最终由 Appel 和 Haken 于 1976 年借助计算机辅助证明得以解决。色数 \(\chi(G)\) 是图最基本的图不变量的核心：它既与团数 \(\omega(G)\)（最大完全子图的大小）通过 \(\chi(G) \geq \omega(G)\) 相关联，也与独立数 \(\alpha(G)\) 通过 \(\chi(G) \geq n/\alpha(G)\) 相关联。对一般图，\(\chi(G) \leq \Delta(G) + 1\)（贪心染色），但这一上界在大多数情形下是松的。Brooks 定理精确刻画了上界取等号的条件，揭示了图结构（特别是回路和完全子图）对色数的深刻约束。

\textbf{定理 39.3}（Brooks 定理，1941）：连通简单图 \(G\) 非完全图，最大度 \(\Delta \geq 3\)，则 \(\chi(G) \leq \Delta\)。

\textbf{证明}（框架）：非正则时取 \(\deg(v) < \Delta\)，\(G - v\) 归纳 \(\Delta\)-着色，\(v\) 至多 \(\Delta-1\) 个邻点限制故可着色。正则时由非完全图可取不相邻 \(x, y\) 使 \(G - \{x, y\}\) 连通。按 DFS 生成树逆序贪心着色 \(G - \{x, y\}\)（每顶点至多 \(\Delta-1\) 个已着色邻点），\(x, y\) 各有 \(\leq \Delta-1\) 个限制，故 \(\chi(G) \leq \Delta\)。\(\blacksquare\)

Brooks 定理的例外情形（完全图和奇回路）揭示了色数理论中的核心现象：完全图 \(K_{\Delta+1}\) 的团数等于 \(\chi = \Delta+1\)，而奇回路 \(C_{2k+1}\) 的色数 \(\chi = 3\) 严格大于 \(\Delta = 2\)。这些例外并非偶然------它们恰好是''局部不可约''的图，即每个顶点邻域的诱导子图本身不含可被约化的结构。Brooks 定理的证明关键（在正则情形下）在于构造一个顶点排序，使得除最后两个顶点外，每个顶点在该排序中至多有 \(\Delta-1\) 个在前邻居，这保证了贪心算法仅需 \(\Delta\) 种颜色。这一技巧与 Gallai-Roy 定理和 Hajós 猜想等更广泛的图着色问题有深刻联系。

\subsection{39.4 平面图与 Kuratowski 定理}\label{ux5e73ux9762ux56feux4e0e-kuratowski-ux5b9aux7406}

平面图是可以嵌入平面而不产生边交叉的图。Euler 公式 \(n - e + f = 2\) 是平面图理论中最基本的恒等式，它将顶点数、边数和面数三者联系起来。仅凭其推论 \(e \leq 3n - 6\) 即可证明 \(K_5\) 不可平面嵌入，而二部图版本 \(e \leq 2n - 4\) 则排除了 \(K_{3,3}\)。这些不等式虽能分别排除 \(K_5\) 和 \(K_{3,3}\)，但不构成平面性的充要条件。Kuratowski 定理（1930 年）以惊人的方式确认：\(K_5\) 和 \(K_{3,3}\) 的细分恰好是所有不可平面化的极小障碍。

\textbf{定义 39.5}（平面图）：\(G\) 为\textbf{平面图}若可嵌入 \(\mathbb{R}^2\) 使边仅交于顶点。嵌入划分平面为 \(f\) 个\textbf{面}。

\textbf{定理 39.4}（Euler 公式）：连通平面嵌入满足 \(n - e + f = 2\)。

\textbf{证明}：对 \(e\) 归纳。\(e = 0\) 时 \(n = 1, f = 1\)。\(e \geq 1\)：若有叶子则删叶子和边（\(n, e\) 同减 1，\(f\) 不变）；否则删圈上一边（\(e, f\) 同减 1，\(n\) 不变）。由归纳假设即得。\(\blacksquare\)

\textbf{推论 39.5}：简单平面图（\(n \geq 3\)）满足 \(e \leq 3n - 6\)。故 \(K_5\)（\(10 > 9\)）非平面；二部图界 \(e \leq 2n - 4\) 给出 \(K_{3,3}\)（\(9 > 8\)）非平面。

\textbf{证明}：对简单平面图，每个面至少由 \(3\) 条边围成，且每条边至多属于 \(2\) 个面。故 \(3f \leq 2e\)，即 \(f \leq 2e/3\)。代入 Euler 公式 \(n - e + f = 2\) 得 \(n - e + 2e/3 \geq 2\)，即 \(e \leq 3n - 6\)。对 \(K_5\)：\(n=5\)，\(e=10\)，\(3n-6=9\)，\(10 > 9\) 违反不等式，故 \(K_5\) 非平面。对二部平面图，最短圈至少为 \(4\)，故 \(4f \leq 2e\)，即 \(f \leq e/2\)，代入得 \(e \leq 2n - 4\)。对 \(K_{3,3}\)：\(n=6\)，\(e=9\)，\(2n-4=8\)，\(9 > 8\)，故 \(K_{3,3}\) 非平面。\(\blacksquare\)

\textbf{定理 39.5}（Kuratowski 定理）：图是平面图 \(\iff\) 不含 \(K_5\) 或 \(K_{3,3}\) 的细分作为子图。必要性由 \(K_5, K_{3,3}\) 非平面且细分保非平面性；充分性通过 \(3\)-连通成分的 Whitney 嵌入定理与轮图分解的归纳证明。

\textbf{证明}：（必要性）若 \(G\) 含 \(K_5\) 或 \(K_{3,3}\) 的细分，由于 \(K_5\) 和 \(K_{3,3}\) 非平面（推论 39.5），且边的细分不改变可平面性（将边替换为路径不引入交叉），故 \(G\) 非平面。（充分性）可分解为 \(2\)-连通分支，每个 \(2\)-连通分支可分解为 \(3\)-连通成分。由 Whitney 定理，\(3\)-连通平面图有唯一的平面嵌入。对 \(3\)-连通非平面图，通过极小反例的轮图分解（Tutte 的轮图定理）可找到 \(K_5\) 或 \(K_{3,3}\) 的细分。归纳组合所有分支，若所有分支不含 \(K_5\) 或 \(K_{3,3}\) 的细分，则 \(G\) 为平面图。\(\blacksquare\)

\subsection{39.5 Ramsey 理论}\label{ramsey-ux7406ux8bba}

Ramsey 理论的核心论断------``完全无序是不可能的''------表明在任何足够大的结构中必然存在具有高度规律性的子结构。在图的语境下，Ramsey 数 \(R(s,t)\) 是最小正整数 \(n\) 使得任何 \(n\) 个顶点的图中必含 \(K_s\) 或 \(\overline{K}_t\) 作为子图。Ramsey 定理（1930 年）证明了这些数对于任意 \(s, t\) 都是有限的，其证明中使用的递归不等式 \(R(s, t) \leq R(s-1, t) + R(s, t-1)\) 至今仍是最基本的 Ramsey 估计方法。

\textbf{定义 39.6}（Ramsey 数）：\(R(s, t)\) 是最小正整数 \(n\) 使得 \(n\) 阶图或其补图 \((\overline{G})\) 必含 \(K_s\) 子图，或 \(K_t\) 子图。特别地，\(R(k, k)\) 是最小 \(n\) 使得 \(n\) 阶图或其补图必含 \(K_k\) 子图。

\textbf{定理 39.6}（Ramsey 定理）：\(R(s, t) \leq R(s-1, t) + R(s, t-1)\)；\(R(s, t) \leq \binom{s+t-2}{s-1}\)。

\textbf{证明}：取 \(n = R(s-1, t) + R(s, t-1)\)，取 \(v \in V(K_n)\)。\(v\) 有 \(\geq R(s-1, t)\) 条红边或 \(\geq R(s, t-1)\) 条蓝边。前者红邻点集中或含红色 \(K_{s-1}\)（连 \(v\) 成 \(K_s\)）或含蓝色 \(K_t\)；后者对称。故 \(R(s, t) \leq n\)。\(\blacksquare\)

\textbf{定理 39.7}（\(R(3, 3) = 6\) 的完整证明）：

\textbf{证明}：下界 \(R(3, 3) > 5\)：\(C_5\) 外圈染红、补图（也是 \(C_5\)）染蓝，无单色三角形。上界 \(R(3, 3) \leq 6\)：\(K_6\) 中取 \(v\)，鸽巢原理给出 \(\geq 3\) 条同色边（不妨红），邻点 \(\{x, y, z\}\)。若其间有红边则与 \(v\) 成红色 \(K_3\)；若全为蓝边则成蓝色 \(K_3\)。故 \(R(3, 3) = 6\)。\(\blacksquare\)

\textbf{注记 39.1}：已知 Ramsey 数极少：\(R(3,k)\) 对 \(k \leq 9\) 已知，\(R(4,4)=18\)，\(R(4,5)=25\)。\(R(5,5)\) 的确切值至今未知（仅知 \(43\)--\(48\)）。

\newpage

\chapter{04-离散数学/01-组合数学/03-匹配与极值.md}\label{ux79bbux6563ux6570ux5b6601-ux7ec4ux5408ux6570ux5b6603-ux5339ux914dux4e0eux6781ux503c.md}

\section{第84章：极值集合论（Extremal Set Theory）}\label{ux7b2c84ux7ae0ux6781ux503cux96c6ux5408ux8bbaextremal-set-theory}

极值集合论研究集合系在一定限制条件下能达到的最大尺寸或最优结构。本章建立三个经典定理------Sperner 定理、Erdős-Ko-Rado 定理和 Kruskal-Katona 定理------的严格数学表述，并给出各自的核心证明方法。

\subsection{基本定义}\label{ux57faux672cux5b9aux4e49}

\textbf{定义 44.1}（反链 / Antichain）：设 \(\mathcal{F} \subseteq 2^{[n]}\) 是 \([n] = \{1, \ldots, n\}\) 的子集族。\(\mathcal{F}\) 称为\textbf{反链}（或 Sperner 族），若对任意 \(A, B \in \mathcal{F}\)，\(A \neq B\)，有 \(A \not\subset B\) 且 \(B \not\subset A\)。即 \(\mathcal{F}\) 中任意两个集合互不包含。

\textbf{定义 44.2}（影子与初始段）：设 \(\mathcal{F}\) 是 \([n]\) 的 \(k\)-元子集族（即 \(\mathcal{F} \subseteq \binom{[n]}{k}\)）。\(\mathcal{F}\) 的\textbf{影子}（shadow）定义为

\[\partial \mathcal{F} = \left\{ G \in \binom{[n]}{k-1} : \exists F \in \mathcal{F}, \; G \subset F \right\}\]

即由 \(\mathcal{F}\) 中每个 \(k\)-集合删除一个元素所得到的所有 \((k-1)\)-集合的族。\textbf{colexicographic 序}（简记 colex）如下定义：\(A <_{\text{colex}} B\) 当且仅当 \(\max(A \triangle B) \in B\)。在 colex 序下，\(\binom{[n]}{k}\) 的\textbf{初始段}（initial segment）是前 \(m\) 个（按 colex 序）\(k\)-集合构成的族，记为 \(\mathcal{I}_k(m)\)。

\textbf{定义 44.3}（交族 / Intersecting Family）：\(\mathcal{F} \subseteq 2^{[n]}\) 称为\textbf{交族}，若对任意 \(A, B \in \mathcal{F}\)，有 \(A \cap B \neq \emptyset\)。若进一步要求 \(\mathcal{F} \subseteq \binom{[n]}{k}\)，则称 \(\mathcal{F}\) 为 \(k\)-一致交族。

\subsection{Sperner 定理}\label{sperner-ux5b9aux7406}

\textbf{定理 44.1}（Sperner 定理，1928）：设 \(\mathcal{F} \subseteq 2^{[n]}\) 为反链。则

\[|\mathcal{F}| \leq \binom{n}{\lfloor n/2 \rfloor}\]

且等号成立当且仅当 \(\mathcal{F}\) 是全体 \(\lfloor n/2 \rfloor\)-元子集的族或全体 \(\lceil n/2 \rceil\)-元子集的族。

\textbf{证明}（Lubell-Yamamoto-Meshalkin (LYM) 不等式方法）：考虑 \([n]\) 的所有 \(n!\) 个排列。对每个 \(A \in \mathcal{F}\)，计算满足 \(\{\sigma(1), \ldots, \sigma(|A|)\} = A\) 的排列 \(\sigma\) 的个数。对于固定的 \(A\)，恰有 \(|A|! (n - |A|)!\) 个排列将 \(A\) 的元素排在前 \(|A|\) 位。由于 \(\mathcal{F}\) 是反链，不同的 \(A \in \mathcal{F}\) 对应的排列集合互不相交（若存在排列同时将 \(A\) 和 \(B\) 的元素排在前方，则 \(A \subseteq B\) 或 \(B \subseteq A\)，与反链性矛盾）。因此

\[\sum_{A \in \mathcal{F}} |A|! (n - |A|)! \leq n!\]

两边除以 \(n!\) 得 \textbf{LYM 不等式}：

\[\sum_{A \in \mathcal{F}} \frac{1}{\binom{n}{|A|}} \leq 1\]

由二项式系数的单峰性，\(\binom{n}{k} \leq \binom{n}{\lfloor n/2 \rfloor}\) 对所有 \(k\) 成立。故

\[1 \geq \sum_{A \in \mathcal{F}} \frac{1}{\binom{n}{|A|}} \geq \frac{|\mathcal{F}|}{\max_k \binom{n}{k}} = \frac{|\mathcal{F}|}{\binom{n}{\lfloor n/2 \rfloor}}\]

即 \(|\mathcal{F}| \leq \binom{n}{\lfloor n/2 \rfloor}\)。等号成立要求 \(\mathcal{F}\) 中所有集合大小均为 \(\lfloor n/2 \rfloor\)（或均为 \(\lceil n/2 \rceil\)）且 LYM 不等式取等，故 \(\mathcal{F}\) 为全体最中间层的子集。\(\blacksquare\)

\subsection{Erdős-Ko-Rado 定理}\label{erdux151s-ko-rado-ux5b9aux7406}

\textbf{定理 44.2}（Erdős-Ko-Rado 定理，1961）：设 \(\mathcal{F} \subseteq \binom{[n]}{k}\) 为 \(k\)-一致交族。若 \(n \geq 2k\)，则

\[|\mathcal{F}| \leq \binom{n-1}{k-1}\]

且当 \(n > 2k\) 时等号成立当且仅当 \(\mathcal{F}\) 由全体包含某个固定元素的 \(k\)-子集构成（即 \(\mathcal{F} = \{F \in \binom{[n]}{k} : x \in F\}\) 对某个 \(x \in [n]\)）。当 \(n = 2k\) 时，任意两个 \(k\)-子集交非空当且仅当它们不是互补对，故上界为 \(\frac{1}{2}\binom{2k}{k} = \binom{2k-1}{k-1}\)。

\textbf{证明}（Katona 圆排列法）：将 \([n]\) 的 \(n\) 个元素等距排列在圆周上，共有 \((n-1)!\) 种不同的圆排列（循环排列）。对每个圆排列，考虑由连续 \(k\) 个元素构成的弧（arc）。共 \(n\) 条长度为 \(k\) 的弧。

\textbf{断言}：在固定的圆排列上，\(\mathcal{F}\) 至多包含 \(k\) 条弧。若 \(\mathcal{F}\) 包含 \(k+1\) 条弧，由抽屉原理必有两条弧的起始位置距离（沿圆周最短距离）小于 \(k\)。取两条最近的弧 \(A, B \in \mathcal{F}\)，两者起始位置距离为 \(d < k\)。由于 \(n \geq 2k\)，\(A\) 的长度 \(k\) 加上间距 \(d\) 不超过 \(2k \leq n\)，故 \(A\) 延伸不到 \(B\) 的位置。但 \(A\) 的起始必须在 \(B\) 之外（反之亦然），推得 \(A \cap B = \emptyset\)，矛盾。

双重计数：每个 \(F \in \mathcal{F}\)（\(k\)-子集）在多少圆排列中作为连续弧出现？固定 \(F\)，将其元素连续排列在圆周上的圆排列数为 \(k! (n-k)!\)（\(F\) 中 \(k\) 个元素的内部排列 \(\times\) \(F\) 外 \(n-k\) 个元素的内部排列）。因此在所有 \((n-1)!\) 个圆排列中，\(\mathcal{F}\) 中集合作为连续弧出现的总次数为

\[|\mathcal{F}| \cdot k! (n-k)! \leq (n-1)! \cdot k\]

（每个圆排列至多贡献 \(k\) 次）。整理得

\[|\mathcal{F}| \leq \frac{k \cdot (n-1)!}{k! (n-k)!} = \frac{(n-1)!}{(k-1)! (n-k)!} = \binom{n-1}{k-1}\]

即所求证。等号分析由断言中 \(d \geq k\) 的条件导出当 \(n > 2k\) 时所有极大交族均含公共元素。\(\blacksquare\)

\subsection{Kruskal-Katona 定理}\label{kruskal-katona-ux5b9aux7406}

\textbf{定理 44.3}（Kruskal-Katona 定理，1963/1968）：设 \(\mathcal{F} \subseteq \binom{[n]}{k}\)。则

\[|\partial \mathcal{F}| \geq |\partial \mathcal{I}_k(|\mathcal{F}|)|\]

其中 \(\mathcal{I}_k(m)\) 是 \(\binom{[n]}{k}\) 中前 \(m\) 个（按 colex 序）\(k\)-集合构成的族。换言之：在所有 \(|\mathcal{F}|\) 固定的 \(k\)-一致族中，colex 初始段具有最小的影子大小。该定理的等价表述是：若 \(|\mathcal{F}| = \binom{a_k}{k} + \binom{a_{k-1}}{k-1} + \cdots + \binom{a_t}{t}\)（其中 \(a_k > a_{k-1} > \cdots > a_t \geq t \geq 1\)）为 \(|\mathcal{F}|\) 的 \(k\)-二项式表示（唯一存在），则

\[|\partial \mathcal{F}| \geq \binom{a_k}{k-1} + \binom{a_{k-1}}{k-2} + \cdots + \binom{a_t}{t-1}\]

\textbf{证明思路}：核心是 Lovász 的巧妙证明（1979），利用移位操作（shifting / compression）。对 \(i < j\)，定义 \((i, j)\)-移位 \(S_{ij}\)：将 \(\mathcal{F}\) 中每个含 \(j\) 但不含 \(i\) 的集合 \(F\) 替换为 \((F \setminus \{j\}) \cup \{i\}\)（若该替换后的集合尚不在 \(\mathcal{F}\) 中）。\(S_{ij}\) 保持族的大小不变，且不增影子的大小：\(|\partial S_{ij}(\mathcal{F})| \leq |\partial \mathcal{F}|\)。反复应用移位操作可将任意族 \(\mathcal{F}\) 转化为移位不变的族（即 colex 初始段），在此过程中影子大小单调不增，故原始 \(\mathcal{F}\) 的影子不小于 colex 初始段的影子。\(\blacksquare\)

\textbf{推论 44.4}（Kruskal-Katona 在极值图论中的应用）：Kruskal-Katona 定理是极值集合论的基本工具，直接蕴含了 Erdős-Ko-Rado 定理的一个代数证明，且与 Lovász 的 Kneser 猜想的解决（1978）和 Frankl 的并-交问题有深刻联系。

\subsection{Frankl-Wilson 定理与代数方法}\label{frankl-wilson-ux5b9aux7406ux4e0eux4ee3ux6570ux65b9ux6cd5}

\textbf{定理 44.5}（Frankl-Wilson 定理，1981）：设 \(\mathcal{F} \subseteq 2^{[n]}\) 满足所有 \(A \in \mathcal{F}\) 的大小均为 \(L \subseteq \{0, 1, \ldots, n\}\) 中的值，且 \(|A \cap B| \in K\)（模 \(p\) 意义下）对所有 \(A, B \in \mathcal{F}\)，其中 \(p\) 为素数。若 \(|L| \leq s\) 且 \(K \subseteq \{0, 1, \ldots, p-1\}\) 满足 \(|K| = s\) 不包含某类同余类，则 \(|\mathcal{F}| \leq \binom{n}{0} + \binom{n}{1} + \cdots + \binom{n}{s}\)。该定理的证明将集合族编码为有限域 \(\mathbb{F}_p\) 上的多项式空间，利用多项式次数限制和 Schwartz-Zippel 引理导出上界，是极值集合论代数方法的典范。

极值集合论的方法论核心在于将组合结构的最优配置问题转化为不等式系统。Sperner 定理的经典证明------对称链分解（symmetry chain decomposition）------由 de Bruijn、Tengbergen 和 Kruyswijk（1951）给出，其思想是将布尔格 \(B_n\) 分解为 \(\binom{n}{\lfloor n/2 \rfloor}\) 条对称链。Frankl-Wilson 定理代表了极值集合论的代数转向，这一方法在编码理论和计算复杂性中有深远影响。

\hrulefill

\emph{本卷基于 V6 Ch 32（计数原理）和 V6 Ch 12（代数结构），为组合数学和图论建立了系统基础。共 7 章，涵盖生成函数、图论、匹配与网络流、染色与 Ramsey 理论、组合设计、拟阵理论和代数组合学。}

\hrulefill

\emph{本卷基于 V6 Ch 33（概率论）和 V20（概率论 II），建立了 Shannon 信息论、纠错码和密码学的系统理论，共 5 章。补充内容涵盖组合数学与图论：生成函数与递推关系、图论基础（树、连通性、平面图、Euler/Hamilton回路）、匹配与网络流（Hall婚姻定理、最大流最小割定理）、染色与Ramsey理论、组合设计（BIBD、Steiner系、正交拉丁方）、拟阵理论和代数组合学，与初等计数原理形成互补。}

\emph{卷五：离散数学与概率统计终。}

\emph{卷五：离散数学与概率统计终。}

\emph{卷五：离散数学与概率统计终。}

\emph{卷五：组合数学与图论终。} \hrulefill

\section{第85章：匹配理论深化}\label{ux7b2c85ux7ae0ux5339ux914dux7406ux8bbaux6df1ux5316}

匹配理论是图论中最为系统化的分支之一，其核心问题在于寻找图中两两不共享顶点的最大边集。从 Frobenius（1912 年）对婚配问题的组合描述到 Hall 定理（1935 年）给出的完美匹配充分必要条件，再到 Tutte 定理（1947 年）将匹配理论从二分图推广到一般图，这一领域经历了从特设解法到完备公理体系的飞跃。匹配理论在数学内部与线性规划（全幺模矩阵保证二分图匹配的整数性）、代数学（Pfaffian 与完美匹配计数）和组合优化（Edmonds 的花算法）深刻交织；在应用层面，它是市场设计（Gale-Shapley 稳定婚姻算法获 2012 年诺贝尔经济学奖）、调度问题和分配问题的数学基础。本章在前两章图论基础之上，从 Hall 定理的深入证明出发，经 Tutte 定理推广至一般图，以 Gallai-Edmonds 分解为结构高潮，最后以极值图论（Turán 定理与 Erdős-Stone 定理）收尾，将匹配理论与极值组合学统一。

\subsection{39.1 Hall 定理的深入阐述}\label{hall-ux5b9aux7406ux7684ux6df1ux5165ux9610ux8ff0}

Hall 婚配定理（1935 年）是二分图匹配理论的基石，其简洁的陈述------二分图存在饱和一侧顶点的匹配当且仅当任意子集的邻域不小于自身------揭示了匹配存在性的纯粹组合本质。Hall 定理在数学史上具有双重意义：一方面，它是 Frobenius-König 行列式理论（全幺模矩阵的秩等于最大匹配大小）的图论重述；另一方面，它是 Menger 定理和最大流-最小割定理在二分图上的等价形式，因而将图论、网络流和线性规划统一在同一框架下。本节从二分图匹配的基本概念出发，给出 Hall 定理的完整归纳证明，并导出 König 定理（最大匹配 = 最小顶点覆盖）和 Frobenius 婚配定理（正则二分图必有完美匹配），为 Tutte 对一般图的推广（39.2 节）做好准备。

\subsubsection{39.1.1 二分图匹配的基本概念}\label{ux4e8cux5206ux56feux5339ux914dux7684ux57faux672cux6982ux5ff5}

\begin{quote}
\textbf{直觉与动机} 匹配理论研究的是将两类对象''配对''的可能性------工人与工作、学生与宿舍、左右两侧的顶点。匹配理论不仅是组合优化和网络流的核心，还在经济学（市场设计）、计算机科学（调度算法）和代数学（群表示论）中有深刻应用。
\end{quote}

\textbf{定义 39.1}（匹配与覆盖）：设 \(G = (V, E)\) 为图。\textbf{匹配} \(M \subseteq E\) 是两两不交的边集。\(M\) 为\textbf{完美匹配}若每个顶点恰与 \(M\) 中一条边关联。\(M\) 为\textbf{最大匹配}若 \(|M|\) 达到最大。\(G\) 为二分图若 \(V = X \cup Y\)（\(X \cap Y = \varnothing\)）且每条边连接 \(X\) 与 \(Y\)。\textbf{顶点覆盖} \(C \subseteq V\) 满足每条边至少有一个端点在 \(C\) 中。\textbf{最小顶点覆盖} \(\tau(G)\) 为最小顶点覆盖的大小。最大匹配大小记为 \(\nu(G)\)。

\textbf{定义 39.2}（邻域与 Hall 条件）：对二分图 \(G = (X \cup Y, E)\) 及 \(S \subseteq X\)，定义 \(N(S) = \{y \in Y : \exists x \in S, xy \in E\}\)。若对任意 \(S \subseteq X\) 有 \(|N(S)| \geq |S|\)，则称 \(G\) 满足 \textbf{Hall 条件}。

\subsubsection{39.1.2 Hall 婚配定理}\label{hall-ux5a5aux914dux5b9aux7406}

\textbf{定理 39.1}（Hall 婚配定理，1935）：二分图 \(G = (X \cup Y, E)\) 存在饱和 \(X\) 中所有顶点的匹配（即 \(X\)-完美匹配）当且仅当对任意 \(S \subseteq X\) 有 \(|N(S)| \geq |S|\)。

\textbf{证明}（\(\Rightarrow\)）：设 \(M\) 为饱和 \(X\) 的匹配，\(S \subseteq X\)。\(M\) 将 \(S\) 中顶点一一匹配到 \(N(S)\) 中不同顶点，故 \(|N(S)| \geq |S|\)。

\textbf{证明}（\(\Leftarrow\)）：对 \(|X|\) 归纳。\(|X| = 1\) 时平凡。设 \(|X| \geq 2\) 且 Hall 条件成立。

\textbf{情形 1}：对任意非空真子集 \(S \subset X\)，\(|N(S)| \geq |S| + 1\)。取 \(x \in X\) 及 \(y \in N(x)\)，考虑 \(G' = G - \{x, y\}\)。对 \(S' \subseteq X \setminus \{x\}\)，\(N_{G'}(S') \supseteq N_G(S') \setminus \{y\}\)，故 \(|N_{G'}(S')| \geq |N_G(S')| - 1 \geq (|S'| + 1) - 1 = |S'|\)。\(G'\) 满足 Hall 条件，归纳得 \(X \setminus \{x\}\)-完美匹配，加上 \(xy\) 即得 \(X\)-完美匹配。

\textbf{情形 2}：存在非空真子集 \(S \subset X\) 使 \(|N(S)| = |S|\)。令 \(X_1 = S\)，\(Y_1 = N(S)\)，\(X_2 = X \setminus S\)，\(Y_2 = Y \setminus Y_1\)。由归纳假设，\(G[X_1 \cup Y_1]\) 有 \(X_1\)-完美匹配。对 \(G[X_2 \cup Y_2]\)：若存在 \(T \subseteq X_2\) 使 \(|N(T) \cap Y_2| < |T|\)，则 \(|N(S \cup T)| = |N(S)| + |N(T) \cap Y_2| < |S| + |T|\)，违反 Hall 条件。故 \(G[X_2 \cup Y_2]\) 满足 Hall 条件，归纳得 \(X_2\)-完美匹配。合并即得全图 \(X\)-完美匹配。\(\blacksquare\)

\textbf{定理 39.2}（König 定理，1931）：二分图 \(G\) 中，最大匹配大小等于最小顶点覆盖大小：\(\nu(G) = \tau(G)\)。

\textbf{证明}：设 \(M\) 为最大匹配，\(U\) 为 \(X\) 中未被 \(M\) 饱和的顶点集，\(Z\) 为从 \(U\) 沿 \(M\)-交错路径可达的顶点集。令 \(C = (X \setminus Z) \cup (Y \cap Z)\)。可验证 \(C\) 为顶点覆盖且 \(|C| = |M|\)。由于 \(\nu(G) \leq \tau(G)\) 恒成立，故 \(\nu(G) = \tau(G)\)。\(\blacksquare\)

\textbf{注记 39.1}：König 定理与 Menger 定理（边形式）和最大流最小割定理在二分图情形下等价。具体地，将二分图转化为源-汇网络（源连 \(X\)，\(Y\) 连汇，所有边容量为 \(1\)），最大流给出最大匹配，最小割给出最小顶点覆盖。这一等价性揭示了二分图匹配的''整数性''------全幺模矩阵保证极端点均为整数解。

\textbf{推论 39.3}（Frobenius 婚配定理）：设 \(X\) 和 \(Y\) 为有限集，\(|X| = |Y| = n\)，且 \(G = (X \cup Y, E)\) 为正则二分图（每个顶点的度相等，设为 \(d\)）。则 \(G\) 有完美匹配。

\textbf{证明}：对 \(S \subseteq X\)，设 \(E(S, N(S))\) 为 \(S\) 与 \(N(S)\) 之间的边集。计数两端：\(|E(S, N(S))| = d|S|\)（\(S\) 中每顶点贡献 \(d\) 条边）。同时 \(|E(S, N(S))| \leq d|N(S)|\)（\(N(S)\) 中每顶点至多接收 \(d\) 条边）。故 \(d|S| \leq d|N(S)|\)，即 \(|S| \leq |N(S)|\)。由 Hall 定理，存在完美匹配。\(\blacksquare\)

\subsection{39.2 Tutte 定理与一般图匹配}\label{tutte-ux5b9aux7406ux4e0eux4e00ux822cux56feux5339ux914d}

Hall 定理完美地解决了二分图匹配问题，但在一般（非二分）图中，匹配理论遭遇了本质性的困难------奇圈的存在破坏了 Hall 条件的直接推广。Tutte 定理（1947 年）以奇分支计数 \(o(G - S)\) 取代了 Hall 条件中的邻域大小，给出了任意图存在完美匹配的充分必要条件，被誉为''匹配理论中最深刻的定理''。这一突破来自第二次世界大战期间的组合学研究（Tutte 当时在英国从事密码分析工作），其证明引入的''极大坏集''构造方法深刻影响了此后的组合极值理论。Tutte 定理不仅是 Hall 定理和 Petersen 定理（3-正则无桥图有完美匹配）的统一推广，更为 Edmonds 的花算法（Blossom Algorithm，1965 年）提供了理论基础------该算法是第一个在一般图中求最大匹配的多项式时间算法，将匹配理论从存在性推进到算法层面。

\subsubsection{39.2.1 奇分支与 Tutte 条件}\label{ux5947ux5206ux652fux4e0e-tutte-ux6761ux4ef6}

\textbf{定义 39.3}（奇分支）：对图 \(G\) 和 \(S \subseteq V(G)\)，以 \(o(G - S)\) 记 \(G - S\) 中具有奇数个顶点的连通分支（奇分支）的个数。

\textbf{定理 39.4}（Tutte 定理，1947）：图 \(G = (V, E)\) 有完美匹配当且仅当对任意 \(S \subseteq V\) 有 \(o(G - S) \leq |S|\)。

\textbf{证明}（\(\Rightarrow\)）：设 \(M\) 为完美匹配。对 \(S \subseteq V\)，考虑 \(G - S\) 的奇分支。每个奇分支的顶点数为奇数，\(M\) 在其内部配对的顶点数必为偶数，故至少有一个顶点与 \(S\) 中顶点匹配。不同奇分支的''出边''匹配到 \(S\) 中不同顶点，故 \(o(G - S) \leq |S|\)。

\textbf{证明}（\(\Leftarrow\)）：对 \(|V|\) 归纳。\(|V|\) 为奇数时 \(o(G) \geq 1\)，取 \(S = \varnothing\) 得 \(o(G) \leq 0\) 矛盾，故 \(|V|\) 必为偶数。设 \(|V| = 2n\)，Tutte 条件成立。

取极大集合 \(S \subseteq V\) 使 \(o(G - S) = |S|\)（这样的 \(S\) 存在，因为 \(S = \varnothing\) 满足 \(o(G) \geq 0\)，且 Tutte 条件保证 \(o(G - S) \leq |S|\) 永远成立）。记 \(|S| = k\)，\(G - S\) 的奇分支为 \(C_1, \ldots, C_k\)。

\textbf{断言 1}：\(G - S\) 无偶分支。若存在偶分支 \(D\)，取 \(v \in D\)，则 \(G - (S \cup \{v\})\) 比 \(G - S\) 多一个奇分支（\(D - v\) 为奇数个顶点），故 \(o(G - (S \cup \{v\})) \geq k + 1 = |S| + 1 = |S \cup \{v\}|\)，由 Tutte 条件等号成立，\(o(G - (S \cup \{v\})) = |S| + 1\)，与 \(S\) 为满足 \(o(G - S) = |S|\) 的极大集矛盾。

\textbf{断言 2}：每个 \(C_i\) 为因子临界图（即删去任意一个顶点后存在完美匹配）。取 \(v \in C_i\)，若 \(C_i - v\) 无完美匹配，由归纳假设存在 \(T \subseteq C_i - v\) 使 \(o((C_i - v) - T) > |T|\)。则 \(S' = S \cup \{v\} \cup T\) 满足 \(o(G - S') \geq k - 1 + o((C_i - v) - T) > k - 1 + |T| = |S| + |T| = |S'| - 1\)，由于奇偶性，\(o(G - S') \geq |S'|\)，由 Tutte 条件等号成立，但 \(|S'| > |S|\)，与 \(S\) 的极大性矛盾。

\textbf{断言 3}：\(S\) 与奇分支间存在将每个奇分支与 \(S\) 中不同顶点关联的匹配。将 \(S\) 和 \(\{C_1, \ldots, C_k\}\) 视为二分图的两部分，由 Hall 定理可证（否则 Tutte 条件不成立）。

由断言 2，每个 \(C_i\) 删去与 \(S\) 匹配的顶点后有完美匹配。合并这些匹配与断言 3 的匹配，即得 \(G\) 的完美匹配。\(\blacksquare\)

\textbf{注记 39.2}：Tutte 定理是匹配理论中最深刻的结果之一。其证明中''极大坏集''的构造是典型的极值论证方法。Tutte 定理统一了 Petersen 定理（\(3\)-正则无桥图有完美匹配）和 Hall 定理（二分图情形下 \(o(G - S)\) 简化为 \(|X| - |N(S)|\) 的奇偶性分析）。Edmonds 的''花算法''（Blossom Algorithm）将 Tutte 定理的构造性证明转化为多项式时间算法，在一般图中求最大匹配的复杂度为 \(O(n^3)\)。

\subsection{39.3 因子临界图与 Gallai-Edmonds 分解}\label{ux56e0ux5b50ux4e34ux754cux56feux4e0e-gallai-edmonds-ux5206ux89e3}

Tutte 定理给出了完美匹配的存在性判据，但并未揭示一般图最大匹配的结构全貌。因子临界图（factor-critical graph）是 Gallai-Edmonds 分解理论（1964-1965 年）的核心''原子''------删去任意一个顶点后即存在完美匹配的连通图。Lovász 和 Plummer 的 Ear 分解定理进一步揭示了因子临界图的递推构造：每个因子临界图可由奇圈经反复添加奇数长度的耳而生成，这为匹配结构的归纳证明提供了系统工具。Gallai-Edmonds 分解将任意图划分为三个部分（亏量部分 D、吸收部分 A 和完备部分 C），是匹配理论的''主结构定理''------正如 Sylow 定理揭示了有限群的结构，Gallai-Edmonds 分解揭示了任意图在匹配视角下的结构划分。本节先定义因子临界图并证明其 Ear 分解，再系统建立 Gallai-Edmonds 分解的全貌。

\subsubsection{39.3.1 因子临界图}\label{ux56e0ux5b50ux4e34ux754cux56fe}

\textbf{定义 39.4}（因子临界图）：连通图 \(G\) 称为\textbf{因子临界图}（factor-critical graph），若对任意 \(v \in V(G)\)，\(G - v\) 存在完美匹配。等价地，\(G\) 有奇数个顶点且删去任意一个顶点后存在完美匹配。

\textbf{定理 39.5}（因子临界图的结构特征）：设 \(G\) 为因子临界图。则 (1) \(G\) 是 \(2\)-边连通的（可能含割点）；(2) 对任意 \(v \in V(G)\)，\(G - v\) 的完美匹配唯一当且仅当 \(G\) 为奇圈 \(C_{2k+1}\)。

\textbf{证明}：(1) 若 \(G\) 含桥 \(e = uv\)，则 \(G - e\) 有两个连通分支。因 \(|V(G)|\) 为奇数，两分支必一奇一偶。设含 \(u\) 的分支为奇数个顶点。删去 \(v\) 后，含 \(u\) 的分支（奇数个顶点）与 \(G - e\) 的其他部分隔离，由 Tutte 定理（考虑 \(S = \{v\}\) 的邻点集）可证 \(G - v\) 无完美匹配，矛盾于因子临界性。故 \(G\) 无桥。(2) 若 \(G\) 为奇圈，显然 \(G - v\) 的路有唯一完美匹配。反之，若 \(G\) 非奇圈，则存在 \(v\) 使 \(G - v\) 有至少两个完美匹配，这可由 Edmonds 的花结构导出。\(\blacksquare\)

\textbf{定理 39.6}（Ear 分解定理，Lovász-Plummer）：图 \(G\) 是因子临界图当且仅当 \(G\) 可由奇圈通过反复添加奇数长度的耳（ear）------即两端在已有图上、内部顶点两两不交的路------构造而成。

\textbf{证明}（\(\Leftarrow\)）：对耳的添加次数归纳。基底：奇圈为因子临界图。归纳步：设 \(G'\) 为因子临界图，\(P\) 为两端在 \(G'\) 上、奇数个内部顶点的耳。将 \(P\) 加到 \(G'\) 上得 \(G\)。对任意 \(v \in V(G)\)，若 \(v\) 在 \(P\) 的内部，则 \(G'\) 有完美匹配 \(M'\)（\(G'\) 为因子临界则 \(G' - x\) 有完美匹配，其中 \(x\) 为 \(P\) 的一个端点），加上 \(P\) 上交替匹配即得 \(G - v\) 的完美匹配。若 \(v \in V(G')\)，由归纳假设 \(G' - v\) 有完美匹配，\(P\) 的内部顶点（奇数个）可沿耳交替匹配，其中一端与 \(G'\) 中未被匹配的顶点配对。\(\blacksquare\)

\textbf{注记 39.3}：Ear 分解定理揭示了因子临界图的递推构造特征，是匹配理论中''结构定理''的典范。因子临界图可视为一般图匹配的''原子''------Gallai-Edmonds 分解将任意图划分为三个部分，其中因子临界图构成核心组件。

\subsubsection{39.3.2 Gallai-Edmonds 结构定理}\label{gallai-edmonds-ux7ed3ux6784ux5b9aux7406}

\textbf{定理 39.7}（Gallai-Edmonds 分解定理，1964-1965）：设 \(G = (V, E)\) 为任意图。定义 \(D(G) = \{v \in V : \text{存在最大匹配未饱和 } v\}\)，\(A(G) = N(D(G)) \setminus D(G)\)，\(C(G) = V \setminus (D(G) \cup A(G))\)。则：

\begin{enumerate}
\def\labelenumi{\arabic{enumi}.}
\tightlist
\item
  \(D(G)\) 的每个连通分支均为因子临界图；
\item
  \(G[C(G)]\) 有完美匹配；
\item
  任意最大匹配 \(M\) 饱和 \(A(G)\) 中所有顶点，且 \(M\) 将 \(A(G)\) 中每个顶点匹配到 \(D(G)\) 的不同分支中；
\item
  \(|M| = \frac{1}{2}(|V| - c(D(G)) + |A(G)|)\)，其中 \(c(D(G))\) 为 \(D(G)\) 的连通分支数。
\end{enumerate}

\textbf{证明}（框架）：令 \(D, A, C\) 如上定义。记 \(D\) 的连通分支为 \(D_1, \ldots, D_k\)。

\begin{enumerate}
\def\labelenumi{(\arabic{enumi})}
\item
  对每个 \(D_i\)，取 \(v \in D_i\)，存在最大匹配 \(M\) 使 \(v\) 未被饱和。由 \(M\) 的最大性，\(M\) 在 \(D_i\) 内的限制是 \(D_i - v\) 的完美匹配。故 \(D_i\) 为因子临界图。
\item
  若 \(G[C]\) 无完美匹配，由 Tutte 定理存在 \(S \subseteq C\) 使 \(o(G[C] - S) > |S|\)。则存在 \(G[C] - S\) 的奇分支不与 \(A\) 中顶点相邻（否则违反 \(D\) 的极大性），这与 \(A\) 的定义矛盾。
\item
  设 \(M\) 为最大匹配。若 \(M\) 未饱和 \(A\) 中某顶点 \(a\)，则存在 \(d \in D\) 使 \(a\) 与 \(d\) 相邻，可在 \(M\) 中找到增广路径，矛盾。由 \(A\) 的定义，\(A\) 中顶点仅与 \(D\) 中顶点相邻（在 \(V \setminus A\) 中），故 \(M\) 将 \(A\) 中顶点匹配到 \(D\) 中。
\item
  由 (1)-(3)，\(M\) 在 \(C\) 上有完美匹配，在 \(A\) 上全饱和（匹配到 \(D\) 中），\(D\) 的每个分支 \(D_i\) 中恰有一个顶点与 \(A\) 匹配或未饱和。故 \(|M| = |C|/2 + |A| + \sum_{i=1}^k (|D_i| - 1)/2 = \frac{1}{2}(|V| - k + |A|)\)。\(\blacksquare\)
\end{enumerate}

\textbf{注记 39.4}：Gallai-Edmonds 分解是匹配理论的''主结构定理''。它将图 \(G\) 的顶点划分为 \(D\)（亏量部分）、\(A\)（吸收部分）和 \(C\)（完备部分），揭示了最大匹配在一般图中的结构规律。该分解可直接导出 Edmonds-Gallai 定理：\(\nu(G) = \frac{1}{2}(|V| - \max_{S \subseteq V} (o(G - S) - |S|))\)，这是 Tutte 定理的定量推广。

\subsection{39.4 极值图论}\label{ux6781ux503cux56feux8bba}

极值图论研究的是这样一个基本问题：给定一个禁戒子图 \(H\)，\(n\) 个顶点的图中最多能有多少条边而不包含 \(H\)？这一问题源于 Turán（1941 年）的开创性工作------他精确确定了不含 \(K_{r+1}\) 的图的最大边数，其答案 Turán 图 \(T_r(n)\) 是极值组合学中最经典的构造。从方法论角度看，极值图论是匹配理论的''对偶''问题：匹配理论关心图中''存在什么结构''，极值图论则关心''避免某种结构时能有多大''。Erdős-Stone 定理（1946 年）将 Turán 定理推广到任意非二分禁戒图 \(H\)，以 Szemerédi 正则引理（1978 年）为终极工具，确立了极值图论的''基本定理''------渐近主项仅依赖于禁戒图的色数。本节从 Turán 定理的归纳证明出发，经 Erdős-Stone 定理和正则引理方法，最后介绍 Kővári-Sós-Turán 定理等经典极值结果及其与 Ramsey 理论和加性组合学的交叉联系。

\subsubsection{39.4.1 Turán 定理}\label{turuxe1n-ux5b9aux7406}

\begin{quote}
\textbf{直觉与动机} 极值图论研究的是：在给定禁戒子图后，\(n\) 个顶点的图最多能有多少条边？这是组合数学中最基本的问题之一，起源于 Turán（1941）对不含 \(K_{r+1}\) 的图的极值研究。
\end{quote}

\textbf{定义 39.5}（Turán 图）：设 \(n, r\) 为正整数，\(r \geq 2\)。\textbf{Turán 图} \(T_r(n)\) 为 \(n\) 个顶点的完全 \(r\)-部图，各部分大小尽可能均衡------即各部分大小为 \(\lfloor n/r \rfloor\) 或 \(\lceil n/r \rceil\)。记 \(t_r(n) = |E(T_r(n))|\)，即 Turán 图的边数。

\textbf{定理 39.8}（Turán 定理，1941）：设 \(G\) 为 \(n\) 个顶点的简单图，不含 \(K_{r+1}\) 作为子图（\(r \geq 2\)）。则 \(|E(G)| \leq t_r(n)\)，且等号成立当且仅当 \(G \cong T_r(n)\)。

\textbf{证明}（归纳法）：对 \(n\) 归纳。\(n \leq r\) 时平凡。设 \(n > r\) 且 \(G\) 为不含 \(K_{r+1}\) 的 \(n\) 顶点极值图（即边数最大）。因 \(G\) 不含 \(K_{r+1}\)，\(G\) 必含 \(K_r\)（否则可添加边）。设 \(K\) 为 \(G\) 中的 \(K_r\)，记 \(G - K\) 为删去 \(K\) 的顶点后剩余的图。

对 \(G - K\) 中每个顶点 \(v\)，\(v\) 至多与 \(K\) 中 \(r-1\) 个顶点相邻（否则 \(v\) 与 \(K\) 构成 \(K_{r+1}\)）。故 \(K\) 与 \(G - K\) 之间的边数 \(\leq (r-1)(n-r)\)。同时 \(|E(G - K)| \leq t_r(n-r)\)（由归纳假设）。因此： \[|E(G)| \leq \binom{r}{2} + (r-1)(n-r) + t_r(n-r)\]

计算 \(t_r(n)\) 的递推：\(t_r(n) = \binom{r}{2} + (r-1)(n-r) + t_r(n-r)\)（将 \(n\) 个顶点分配到 \(r\) 个部分，每个部分至多与其他 \(r-1\) 个部分全连接）。故 \(|E(G)| \leq t_r(n)\)。

等号成立需 \(G - K \cong T_r(n-r)\) 且 \(G - K\) 中每个顶点恰与 \(K\) 中 \(r-1\) 个顶点相邻，这唯一确定了 \(G \cong T_r(n)\)。\(\blacksquare\)

\textbf{推论 39.9}（Mantel 定理，1907）：\(n\) 个顶点的不含三角形的图最多有 \(\lfloor n^2/4 \rfloor\) 条边（即 \(r=2\) 的 Turán 定理）。最大图为完全二分图 \(K_{\lfloor n/2\rfloor, \lceil n/2\rceil}\)。

\subsubsection{39.4.2 Erdős-Stone 定理与渐近极值}\label{erdux151s-stone-ux5b9aux7406ux4e0eux6e10ux8fd1ux6781ux503c}

\textbf{定义 39.6}（极值数）：对图 \(H\)，定义 \(\operatorname{ex}(n, H)\) 为 \(n\) 个顶点不含 \(H\) 作为子图的图的最大边数。

\textbf{定理 39.10}（Erdős-Stone 定理，1946）：设 \(H\) 的色数为 \(\chi(H) = r+1 \geq 3\)。则 \[\operatorname{ex}(n, H) = \left(1 - \frac{1}{r}\right) \binom{n}{2} + o(n^2) = t_r(n) + o(n^2)\]

\textbf{证明}（框架）：下界：\(T_r(n)\) 不含 \(K_{r+1}\)，故不含 \(\chi(H) = r+1\) 的 \(H\)，从而 \(\operatorname{ex}(n, H) \geq t_r(n) = (1 - \frac{1}{r})\binom{n}{2} + O(n)\)。

上界：对任意 \(\varepsilon > 0\)，取充分大的 \(n\)。由 Szemerédi 正则引理，任意 \(n\) 顶点图可划分为常数个几乎等大的部分，使得大部分部分对之间是 \(\varepsilon\)-正则的。若 \(G\) 的边数 \(\geq (1 - \frac{1}{r} + \delta)\binom{n}{2}\)，则由正则引理可在 \(G\) 中找到 \(K_{r+1}(s)\)（每部分 \(s\) 个顶点的完全 \((r+1)\)-部图），其中 \(s\) 可任意大。由 \(H \subseteq K_{r+1}(|V(H)|)\)（因 \(\chi(H) = r+1\)，\(H\) 可嵌入到 \((r+1)\)-部图中），故 \(H \subseteq G\)。\(\blacksquare\)

\textbf{注记 39.5}：Erdős-Stone 定理是极值图论的''基本定理''------它表明对任意非二部图 \(H\)（\(\chi(H) \geq 3\)），\(\operatorname{ex}(n, H)\) 的渐近主项仅依赖于 \(\chi(H)\)，而与 \(H\) 的具体结构无关。当 \(\chi(H) = 2\)（即 \(H\) 为二部图）时，\(\operatorname{ex}(n, H) = o(n^2)\)，其渐近行为由''二部 Turán 问题''决定，如偶圈 \(C_{2k}\) 的极值数为 \(\Theta(n^{1+1/k})\)。

\subsubsection{39.4.3 极值图论的进一步问题}\label{ux6781ux503cux56feux8bbaux7684ux8fdbux4e00ux6b65ux95eeux9898}

\textbf{定理 39.11}（Erdős-Gallai 定理，1959）：\(n\) 个顶点不含长度为 \(k\) 的路 \(P_k\) 的图最多有 \(\frac{k-2}{2}n\) 条边。等号成立当且仅当 \(G\) 为 \(\lfloor (k-1)/2 \rfloor\) 个不相交完全图的并。

\textbf{定理 39.12}（Kővári-Sós-Turán 定理，1954）：\(\operatorname{ex}(n, K_{s,t}) = O(n^{2-1/s})\)（\(s \leq t\)）。特别地，不含 \(K_{2,2} = C_4\) 的图最多有 \(O(n^{3/2})\) 条边。对 \(s \geq 2\)，下界 \(\operatorname{ex}(n, K_{s,t}) = \Omega(n^{2-1/s})\) 由有限域上的范数图构造给出（Brown, 1966），故确界为 \(\Theta(n^{2-1/s})\) 当 \(t \geq (s-1)! + 1\)。

\textbf{证明}（KST 定理上界）：设 \(G\) 不含 \(K_{s,t}\)。对 \(G\) 中所有 \(s\)-元子集进行双重计数。每个 \(s\)-元子集至多与 \(t-1\) 个顶点有共同邻域（否则构成 \(K_{s,t}\)）。故 \(\sum_{v \in V} \binom{\deg(v)}{s} \leq (t-1)\binom{n}{s}\)。由 Jensen 不等式，\(\binom{\bar{d}}{s} \leq (t-1)\binom{n}{s}/n\)，其中 \(\bar{d} = 2|E|/n\)。展开得 \(\bar{d}^s \leq (t-1)n^{s-1} s!\)，故 \(|E| = n\bar{d}/2 = O(n^{2-1/s})\)。\(\blacksquare\)

\textbf{注记 39.6}（极值图论的方法论）：极值图论的核心方法包括：(1) \textbf{正则引理}（Szemerédi, 1978）------将大图分解为拟随机二部图，是 Erdős-Stone 定理和无数图论结果的证明工具；(2) \textbf{概率方法}------通过随机图构造下界，Erdős 以此证明了 \(\operatorname{ex}(n, C_{2k+1}) = \Theta(n^{1+1/(2k)})\)；(3) \textbf{代数方法}------利用有限域、特征和等代数结构构造极值图，如 KST 下界中的Brown 构造和 Wenger 构造；(4) \textbf{稳定性方法}（Simonovits, 1968）------近似极值图必然接近 Turán 图，这一方法将极值问题转化为结构分类问题。极值图论与 Ramsey 理论（强制出现子图 vs.~避免子图的最大边数）和加性组合（Szemerédi 定理的图论证明）有深刻联系，是现代组合数学最活跃的分支之一。

\newpage

\chapter{04-离散数学/01-组合数学/04-组合设计与拟阵.md}\label{ux79bbux6563ux6570ux5b6601-ux7ec4ux5408ux6570ux5b6604-ux7ec4ux5408ux8bbeux8ba1ux4e0eux62dfux9635.md}

\section{第86章：组合设计理论}\label{ux7b2c86ux7ae0ux7ec4ux5408ux8bbeux8ba1ux7406ux8bba}

组合设计理论源自 19 世纪两个经典问题：Kirkman 的女生问题（1850 年------15 名女生每天分成 5 组散步，任意两人恰好同组一次）和 Steiner 对三元系的存在性研究。它研究的是如何将元素分配到''区块''中使得任意子集（通常为对子）恰好在固定数目的区块中出现的均衡配置问题。组合设计与有限几何（射影平面等价于对称 BIBD）、代数学（差集与有限群的正则表示）、编码理论（纠错码的纠错能力等价于码字间的 Hamming 距离设计）和密码学（认证码的最优构造）有多层次的深刻联系。本章从平衡不完全区组设计（BIBD）的基本参数关系（Fisher 不等式）出发，引入正交拉丁方与有限域的构造，深入 Steiner 三元系的存在性理论（Kirkman 定理），最后以拟阵论的深化收尾，展示组合设计在抽象独立结构中的统一视角。

\subsection{40.1 组合设计的基本概念}\label{ux7ec4ux5408ux8bbeux8ba1ux7684ux57faux672cux6982ux5ff5}

\begin{quote}
\textbf{直觉与动机} 组合设计研究的是将元素分配到''区块''中，使得每个区块满足某些均衡性条件。其起源可追溯至 19 世纪 Kirkman 的女生问题（15 名女生每天分成 5 组散步，使任意两人恰好同组一次）和 Steiner 对三元系的研究。组合设计在实验设计、纠错码、密码学和有限几何中有广泛应用。
\end{quote}

\textbf{定义 40.1}（平衡不完全区组设计 / BIBD）：设 \(V\) 为 \(v\) 个元素的集合，\(\mathcal{B}\) 为 \(V\) 的 \(k\)-元子集（称为\textbf{区组}）的搜集。\((V, \mathcal{B})\) 称为 \((v, k, \lambda)\)-\textbf{BIBD}（Balanced Incomplete Block Design），若 \(V\) 中任意一对不同元素恰出现在 \(\lambda\) 个区组中。参数满足： 1. 每个元素出现在 \(r\) 个区组中（\(r\) 称为\textbf{重复数}）； 2. 共有 \(b = |\mathcal{B}|\) 个区组。

\textbf{定理 40.1}（BIBD 的基本参数关系）：对 \((v, k, \lambda)\)-BIBD，有 \[r(k-1) = \lambda(v-1),\quad bk = vr\]

\textbf{证明}：固定元素 \(x \in V\)。\(x\) 出现在 \(r\) 个区组中，每个区组含 \(x\) 外加 \(k-1\) 个其他元素，故 \(x\) 与其他元素的对应对数为 \(r(k-1)\)。另一方面，\(V\) 中除 \(x\) 外有 \(v-1\) 个元素，每对 \((x, y)\) 出现在 \(\lambda\) 个区组中，故 \(x\) 与其他元素的对应对数为 \(\lambda(v-1)\)。因此 \(r(k-1) = \lambda(v-1)\)。第二个等式由双重计数：所有区组中的元素总数为 \(bk = vr\)。\(\blacksquare\)

\textbf{推论 40.2}（Fisher 不等式，1940）：对 \((v, k, \lambda)\)-BIBD 且 \(v > k\)，有 \(b \geq v\)（即区组数不小于元素数）。

\textbf{证明}：设 \(N\) 为 \(v \times b\) 关联矩阵，\(N_{i,j} = 1\) 若元素 \(i\) 属于区组 \(j\)，否则为 \(0\)。则 \(NN^T = (r-\lambda)I + \lambda J\)（\(J\) 为全 \(1\) 矩阵）。该矩阵的特征值为 \(r + \lambda(v-1) = rk\)（重数 \(1\)）和 \(r-\lambda\)（重数 \(v-1\)）。由 \(v > k\) 知 \(r > \lambda\)，故 \(r-\lambda > 0\)，\(NN^T\) 满秩，秩为 \(v\)。而 \(N\) 的秩 \(\leq b\)，故 \(v = \operatorname{rank}(NN^T) \leq \operatorname{rank}(N) \leq b\)。\(\blacksquare\)

\textbf{注记 40.1}：当 \(b = v\) 时，BIBD 称为\textbf{对称设计}（symmetric design），记作 \((v, k, \lambda)\)-SBIBD。对称设计的关联矩阵 \(N\) 满足 \(NN^T = N^T N\)（即 \(N\) 为正规矩阵），因此任意两个区组恰交于 \(\lambda\) 个元素。对称设计与有限射影平面（参数为 \((n^2+n+1, n+1, 1)\)）和 Hadamard 设计（参数为 \((4t-1, 2t-1, t-1)\)）有直接对应，构成组合设计理论的核心对象。

\subsection{40.2 正交拉丁方与有限域构造}\label{ux6b63ux4ea4ux62c9ux4e01ux65b9ux4e0eux6709ux9650ux57dfux6784ux9020}

\textbf{定义 40.2}（拉丁方）：\(n\) 阶\textbf{拉丁方}是 \(n \times n\) 矩阵，每行每列中 \(n\) 个符号各出现一次。两个 \(n\) 阶拉丁方 \(A\) 和 \(B\) 称为\textbf{正交的}（相互正交），若 \(n^2\) 个有序对 \((A_{ij}, B_{ij})\) 全部互异。

\textbf{定理 40.3}（正交拉丁方的基本界）：\(n\) 阶两两正交拉丁方的最大个数 \(N(n) \leq n-1\)。等号成立当且仅当存在 \(n\) 阶射影平面（或 \(n\) 阶仿射平面）。

\textbf{证明}：设 \(L_1, \ldots, L_t\) 为 \(n\) 阶两两正交拉丁方。对每个 \(L_k\)，将其符号重标记为 \(\{1,\ldots,n\}\) 并取 \(L_k\) 的置换，使第一行均为 \(1,2,\ldots,n\)。则 \(L_k\) 的 \((2,1)\) 位置元素互异且不为 \(1\)（否则与 \(L_1\) 的正交性矛盾）。故 \(t \leq n-1\)。等号条件：\(t = n-1\) 时，这些拉丁方与仿射平面 \(AG(2,n)\) 的平行类一一对应，等价于 \(n\) 阶射影平面的存在性。\(\blacksquare\)

\textbf{定理 40.4}（有限域上的构造）：设 \(q = p^e\) 为素数幂，\(\mathbb{F}_q\) 为 \(q\) 元有限域。对每个 \(a \in \mathbb{F}_q^*\)，定义 \(L_a(x, y) = ax + y\)（\(x, y \in \mathbb{F}_q\)）。则 \(\{L_a : a \in \mathbb{F}_q^*\}\) 是 \(q-1\) 个两两正交的 \(q\) 阶拉丁方。

\textbf{证明}：对 \(a \neq b\)，需证 \(L_a\) 与 \(L_b\) 正交。若 \((ax+y, bx+y) = (ax'+y', bx'+y')\)，则相减得 \((a-b)(x-x') = 0\)，故 \(x = x'\)（因 \(a-b \neq 0\)），进而 \(y = y'\)。因此有序对全部互异。\(\blacksquare\)

\textbf{注记 40.2}：正交拉丁方与有限几何的深刻联系在于：\(q-1\) 个相互正交的 \(q\) 阶拉丁方等价于 \(q\) 阶仿射平面 \(AG(2,q)\)，其点集为 \(\mathbb{F}_q^2\)，直线为 \(y = ax + b\)（\(a, b \in \mathbb{F}_q\)）和 \(x = c\)（\(c \in \mathbb{F}_q\)）。当 \(n\) 不是素数幂时，\(N(n)\) 的确切值大多未知------Euler 的 \(36\) 军官问题（\(N(6) = 1\)）的否定由 Tarry（1900）通过穷举验证，Bose-Shrikhande-Parker（1959）证明了 \(N(n) \geq 2\) 对所有 \(n \neq 2, 6\) 成立，彻底推翻了 Euler 关于 \(N(4k+2) = 1\) 的猜想。

\subsection{40.3 Steiner 三元系}\label{steiner-ux4e09ux5143ux7cfb}

Steiner 三元系 STS(v) 是组合设计理论中最经典的构造------它是参数为 (v, 3, 1) 的平衡不完全区组设计，等价于在 v 个元素的集合中构造若干三元组使得任意一对元素恰好出现在一个三元组中。这一问题的最早形式可追溯至 Kirkman 1847 年的女生问题（STS(15) 的可分解性），但系统性的存在性研究始于 Steiner（1853 年）和 Kirkman 的独立工作。STS(v) 的存在与否由模 6 同余条件完全刻画：\(v \equiv 1\) 或 \(3 \pmod{6}\)，这一结果的充分性由 Bose（1939 年，\(v \equiv 3\)）和 Skolem（1958 年，\(v \equiv 1\)）的构造法完成。令人惊叹的是，如此简洁的判定需要历经一个世纪的努力，其构造方法------差集法、拟群法和递推构造------已成为组合设计方法论的核心遗产。

\subsubsection{40.3.1 定义与存在性条件}\label{ux5b9aux4e49ux4e0eux5b58ux5728ux6027ux6761ux4ef6}

\textbf{定义 40.3}（Steiner 三元系）：\textbf{Steiner 三元系} \(\operatorname{STS}(v)\) 是 \(v\) 元集 \(V\) 上的 \(3\)-元子集（称为\textbf{三元组}或\textbf{区组}）的搜集，使得 \(V\) 中任意一对不同元素恰出现在一个三元组中。等价地，\(\operatorname{STS}(v)\) 是参数为 \((v, 3, 1)\) 的 BIBD。

\textbf{定理 40.5}（Steiner 三元系的存在性必要条件）：若 \(\operatorname{STS}(v)\) 存在，则 \(v \equiv 1\) 或 \(3 \pmod{6}\)。

\textbf{证明}：由 BIBD 的参数关系 \(r(k-1) = \lambda(v-1)\)，代入 \(k=3, \lambda=1\) 得 \(r \cdot 2 = 1 \cdot (v-1)\)，即 \(r = (v-1)/2\)。\(r\) 必为整数，故 \(v\) 为奇数，即 \(v \equiv 1, 3, 5 \pmod{6}\)。由 \(bk = vr\) 得 \(b \cdot 3 = v \cdot (v-1)/2\)，即 \(b = v(v-1)/6\)。\(b\) 必为整数，故 \(v(v-1) \equiv 0 \pmod{6}\)。当 \(v \equiv 5 \pmod{6}\) 时，\(v = 6m+5\)，\(v(v-1) = (6m+5)(6m+4) = 36m^2 + 54m + 20 \equiv 2 \pmod{6}\)，不满足整除条件。当 \(v \equiv 1 \pmod{6}\) 时，\(v = 6m+1\)，\(v(v-1) = (6m+1)(6m) = 36m^2 + 6m \equiv 0 \pmod{6}\)。当 \(v \equiv 3 \pmod{6}\) 时，\(v = 6m+3\)，\(v(v-1) = (6m+3)(6m+2) = 36m^2 + 30m + 6 \equiv 0 \pmod{6}\)。故 \(v \equiv 1\) 或 \(3 \pmod{6}\)。\(\blacksquare\)

\textbf{定理 40.6}（Kirkman 定理，1847；Steiner 三元系存在的充分性 / 完全解决于 19 世纪）：\(\operatorname{STS}(v)\) 存在当且仅当 \(v \equiv 1\) 或 \(3 \pmod{6}\)。

\textbf{证明}（框架）：充分性由 Bose 构造法（1939）和 Skolem 构造法（1958）给出。对 \(v \equiv 3 \pmod{6}\)（即 \(v = 6m+3\)），Bose 构造使用 \(2m+1\) 阶交换群上的差集方法。对 \(v \equiv 1 \pmod{6}\)（即 \(v = 6m+1\)），Skolem 构造使用拟群方法。两种构造均可递归地组合为所有满足同余条件的 \(v\)。\(\blacksquare\)

\subsubsection{40.3.2 Kirkman 女生问题与可分解设计}\label{kirkman-ux5973ux751fux95eeux9898ux4e0eux53efux5206ux89e3ux8bbeux8ba1}

\textbf{定义 40.4}（可分解 STS）：\(\operatorname{STS}(v)\) 称为\textbf{可分解的}，若其区组集可划分为 \(r = (v-1)/2\) 个平行类，每个平行类中的区组构成 \(V\) 的一个划分。可分解的 \(\operatorname{STS}(v)\) 记作 \(\operatorname{KTS}(v)\)（Kirkman 三元系）。

\textbf{定理 40.7}（Kirkman 三元系的存在性条件）：\(\operatorname{KTS}(v)\) 存在当且仅当 \(v \equiv 3 \pmod{6}\)。

\textbf{证明}：必要性：由可分解性，每个平行类含 \(v/3\) 个区组，故 \(3 \mid v\)。结合 STS 的必要条件 \(v \equiv 1\) 或 \(3 \pmod{6}\)，得 \(v \equiv 3 \pmod{6}\)。充分性：对 \(v = 3\)（平凡），\(\operatorname{KTS}(3)\) 存在。对 \(v \equiv 3 \pmod{6}\) 且 \(v \geq 9\)，Ray-Chaudhuri 和 Wilson（1971）利用 Bose 构造的推广证明了存在性，这是 Kirkman 女生问题（\(v = 15\)）经过 120 年后的完全解决。\(\blacksquare\)

\textbf{注记 40.3}：Kirkman 女生问题（1850）为 \(v = 15\) 的情形：15 名女生每天分成 5 组（每组 3 人）散步，使得任意两人在 7 天中恰好同组一次。\(\operatorname{KTS}(15)\) 有 \(b = 15 \cdot 14 / 6 = 35\) 个三元组，分为 \(r = 7\) 个平行类（每天 5 个三元组）。解的存在性由 Kirkman 本人于 1847 年证明，但一般 \(v \equiv 3 \pmod{6}\) 的 Kirkman 三元系存在性直到 1971 年才由 Ray-Chaudhuri 和 Wilson 完全解决，这是组合设计理论中的里程碑式成果。

\subsection{40.4 拟阵论深化}\label{ux62dfux9635ux8bbaux6df1ux5316}

拟阵（matroid）是 Whitney（1935 年）为统一线性代数和图论中''独立集''概念而引入的抽象结构。拟阵公理------遗传性和增广性------捕捉了线性无关和森林的最本质特征，其优美之处在于多种等价公理系统（独立集、基、秩函数、回路、闭包算子）之间的无缝转换。拟阵论在组合优化中具有核心地位：贪心算法求解最大权独立集在拟阵上恰好正确（Rado-Edmonds 定理，1957 年），反之拟阵也是使贪心算法正确的唯一结构，这一定理将抽象代数与算法设计完美耦合。在数学体系内部，拟阵是 Tutte 多项式（第 38 章，联系图论和纽结理论）的自然栖息地，与组合设计（横贯拟阵）、代数几何（可表示拟阵与特征相关）和射影几何（仿射/射影拟阵）有深层联系。

\subsubsection{40.4.1 拟阵的基本定义与例}\label{ux62dfux9635ux7684ux57faux672cux5b9aux4e49ux4e0eux4f8b}

\begin{quote}
\textbf{直觉与动机} 拟阵（matroid）是 Whitney 于 1935 年引入的抽象结构，统一了线性代数中的线性无关性、图论中的生成森林和图拟阵，以及组合设计中的横贯拟阵。拟阵论为独立集的公理化研究提供了框架，揭示了看似无关的组合结构之间的深刻共性。
\end{quote}

\textbf{定义 40.5}（拟阵------独立集公理）：\textbf{拟阵} \(M = (E, \mathcal{I})\) 由有限集 \(E\)（\textbf{基集}）和子集族 \(\mathcal{I} \subseteq 2^E\)（\textbf{独立集族}）构成，满足： 1. \(\varnothing \in \mathcal{I}\)； 2. （遗传性）若 \(I \in \mathcal{I}\) 且 \(J \subseteq I\)，则 \(J \in \mathcal{I}\)； 3. （增广性）若 \(I, J \in \mathcal{I}\) 且 \(|I| < |J|\)，则存在 \(e \in J \setminus I\) 使 \(I \cup \{e\} \in \mathcal{I}\)。

\textbf{定义 40.6}（基、秩、回路）：极大独立集称为\textbf{基}（basis）；所有基大小相同，记为 \(r(M)\)（\textbf{秩}）。极小非独立集称为\textbf{回路}（circuit）。\(S \subseteq E\) 的\textbf{秩} \(r(S)\) 为 \(S\) 中极大独立子集的大小。\textbf{闭包} \(\operatorname{cl}(S) = \{e \in E : r(S \cup \{e\}) = r(S)\}\)。秩为 \(2\) 的平坦（flat）称为\textbf{线}（line）。

\textbf{定理 40.8}（拟阵的秩函数特性）：函数 \(r : 2^E \to \mathbb{Z}_{\geq 0}\) 为某拟阵的秩函数当且仅当满足： 1. \(0 \leq r(S) \leq |S|\)； 2. （单调性）\(S \subseteq T \Rightarrow r(S) \leq r(T)\)； 3. （次模性）\(r(S \cup T) + r(S \cap T) \leq r(S) + r(T)\)（对所有 \(S, T \subseteq E\)）。

\textbf{证明}（\(\Rightarrow\)）：(1)与(2)由定义显然。(3)任取 \(S \cap T\) 的极大独立集 \(I\)，增广为 \(S\) 的基 \(I_S\) 和 \(T\) 的基 \(I_T\)。则 \(I_S \cup I_T\) 在 \(S \cup T\) 中独立，且 \(I_S \cap I_T\) 在 \(S \cap T\) 中独立。故 \(r(S \cup T) + r(S \cap T) \leq |I_S \cup I_T| + |I_S \cap I_T| = |I_S| + |I_T| = r(S) + r(T)\)。\(\blacksquare\)

\textbf{例 40.1}（基本拟阵）： - \textbf{线性拟阵}：\(E\) 为向量空间 \(V\) 中有限向量集，\(\mathcal{I}\) 为线性无关子集族。秩等于张成子空间的维数。次模性由维数公式 \(\dim(U+W) + \dim(U \cap W) = \dim(U) + \dim(W)\) 保证。 - \textbf{图拟阵}（回路拟阵）：\(E\) 为图 \(G\) 的边集，\(\mathcal{I}\) 为不含回路的边子集（森林）。基为生成森林，回路为图 \(G\) 中的圈。秩 \(r(S) = |V| - c(S)\)，其中 \(c(S)\) 为子图 \((V, S)\) 的连通分支数。 - \textbf{均匀拟阵} \(U_{k,n}\)：\(|E| = n\)，\(\mathcal{I} = \{I \subseteq E : |I| \leq k\}\)。秩为 \(k\)，基为所有 \(k\)-元子集。当 \(k=2\) 时，\(U_{2,n}\) 的表示问题等价于 \(n\) 个两两线性无关的二维向量，与射影几何和拟阵表示理论密切相关。 - \textbf{横贯拟阵}：设 \(E\) 为集合，\(\mathcal{A} = (A_1, \ldots, A_m)\) 为 \(E\) 的子集族。\(E\) 的子集 \(I\) 为独立的，若 \(I\) 可被 \(\mathcal{A}\) 中互异子集部分横贯（即存在单射 \(\varphi : I \to \{1,\ldots,m\}\) 使 \(e \in A_{\varphi(e)}\)）。横贯拟阵的秩由 Hall 定理的 Rado 推广给出。

\subsubsection{40.4.2 对偶性与拟阵的表示}\label{ux5bf9ux5076ux6027ux4e0eux62dfux9635ux7684ux8868ux793a}

\textbf{定义 40.7}（对偶拟阵）：拟阵 \(M\) 的\textbf{对偶拟阵} \(M^*\) 的基集为 \(E\)（同 \(M\)），其基为 \(M\) 的基的补集：\(\mathcal{B}(M^*) = \{E \setminus B : B \in \mathcal{B}(M)\}\)。对偶拟阵的秩函数为 \(r^*(S) = |S| + r(E \setminus S) - r(E)\)。\((M^*)^* = M\)。

\textbf{定理 40.9}（图拟阵与对偶拟阵）：对平面图 \(G\)，其对偶图 \(G^*\) 的图拟阵 \(M(G^*)\) 同构于 \(M(G)^*\)。

\textbf{证明}：\(G\) 的生成树 \(T\) 的补边集 \(\overline{T}\) 为 \(G^*\) 的生成树当且仅当 \(G\) 为平面图（Whitney 对偶性定理）。故 \(M(G)^*\) 的基与 \(M(G^*)\) 的基一一对应，两拟阵同构。\(\blacksquare\)

\textbf{定义 40.8}（拟阵的可表示性）：拟阵 \(M\) 在域 \(\mathbb{F}\) 上\textbf{可表示}，若存在 \(\mathbb{F}\) 上的矩阵 \(A\)（列对应 \(E\)），使得 \(I \subseteq E\) 独立当且仅当 \(A\) 的对应列线性无关。\(M\) 为\textbf{正则拟阵}若在所有域上可表示；\textbf{二元拟阵}若在 \(\mathbb{F}_2\) 上可表示；\textbf{图拟阵}在所有域上可表示，因而是正则拟阵。

\textbf{定理 40.10}（Tutte 的拟阵表示排除定理）：二元拟阵恰为不含 \(U_{2,4}\) 作为子式（minor）的拟阵。正则拟阵恰为不含 \(U_{2,4}\)、Fano 平面对偶 \(F_7^*\) 作为子式的拟阵。

\textbf{注记 40.4}：拟阵的表示理论是拟阵论的核心课题，与代数几何（Grassmann 流形）、刚性理论（图刚性拟阵）和组合优化（拟阵交、拟阵划分）有深刻联系。拟阵上的贪心算法（Kruskal 算法求最小生成树）可推广到任意拟阵：对权函数 \(w : E \to \mathbb{R}\)，贪心算法（按权值递增依次选取不破坏独立性的元素）给出最大权独立集当且仅当 \((E, \mathcal{I})\) 为拟阵。这一特征将拟阵与组合优化的公理基础牢固绑定。Edmonds 的拟阵交定理（1970）给出了两个拟阵的最大公共独立集的大小公式，其在二分图匹配、分支（arborescence）和矩阵秩问题中有广泛应用。

\subsection{41.A 组合设计补充}\label{a-ux7ec4ux5408ux8bbeux8ba1ux8865ux5145}

\begin{quote}
信息论由 Claude Shannon 于 1948 年创立，是通信和存储的数学基础。本卷从 Shannon 熵和信息度量出发，建立信道容量与编码定理、纠错码理论和压缩感知引论，最后介绍密码学的数学基础。信息论与概率论、组合数学、数论和代数有深刻联系。
\end{quote}

Shannon 的开创性论文《A Mathematical Theory of Communication》（1948）奠定了信息论的数学基础，引入了\textbf{熵}（entropy）作为信息不确定性的度量：对离散随机变量 \(X\)，其熵 \(H(X) = -\sum_x p(x) \log_2 p(x)\) 量化了观察到 \(X\) 的取值所获得的平均信息量。熵的单位为比特（bit），当对数底数为 \(2\) 时，均匀分布的 \(n\) 元随机变量的熵为 \(\log_2 n\)，达到最大值。信息论的核心概念包括\textbf{联合熵} \(H(X, Y)\)、\textbf{条件熵} \(H(Y|X)\) 和\textbf{互信息} \(I(X; Y) = H(X) - H(X|Y)\)，它们满足 \(I(X; Y) = H(X) + H(Y) - H(X, Y) \geq 0\)。\textbf{信道容量} \(C = \max_{p(x)} I(X; Y)\) 表示信道在可靠传输前提下能承载的最大信息速率，Shannon 的信道编码定理证明了存在编码方案使得传输速率接近 \(C\) 且错误概率任意小。在数据压缩方面，\textbf{信源编码定理}指出无损压缩的平均码长不能低于熵 \(H(X)\)，Huffman 编码和算术编码是接近这一下界的最优编码方案。纠错码理论利用代数结构（如有限域上的线性码、Reed-Solomon 码、LDPC 码）在信息中添加冗余以实现错误检测与纠正，与组合设计（如有限域上的正交阵列）和代数几何（如 Goppa 码）有深刻联系。现代密码学建立在计算复杂性假设之上，将信息论的可证明安全性与计算不可区分性相结合，Diffie-Hellman 密钥交换、RSA 公钥加密和 AES 对称加密是这一范式的标志性成果。

\begin{quote}
\textbf{前置知识}：本卷信息论部分依赖本卷第32-33章（概率基础）和微积分中导数的概念，无需概率论II的先修知识。压缩感知部分需第55章（线性代数）中的矩阵运算。密码学部分依赖第5章（数论基础）和第57章（群论基础）。
\end{quote}

\subsection{信息论中的补充主题}\label{ux4fe1ux606fux8bbaux4e2dux7684ux8865ux5145ux4e3bux9898}

信息论在 Shannon 奠基性工作之后，发展出了多个纵深方向，与统计推断、计算机科学和物理学形成了丰富的交叉。\textbf{Kolmogorov 复杂度}（Solomonoff 1964, Kolmogorov 1965）提供了信息论的计算视角：字符串 \(x\) 的 Kolmogorov 复杂度 \(K(x)\) 是能输出 \(x\) 的最短程序的长度，它独立于具体计算模型（在加法常数范围内）。\(K(x)\) 与 Shannon 熵 \(H(X)\) 之间存在深刻联系------对任意可计算分布 \(P\)，\(K(x) \approx -\log_2 P(x)\) 以高概率成立。这一对偶性将信息论中的''不确定性''概念与计算理论中的''不可压缩性''概念统一起来，为通用预测和 Solomonoff 归纳提供了理论基础。\textbf{率失真理论}（Shannon 1959）在有损压缩设定下推广了信源编码定理，定义了率失真函数 \(R(D) = \inf_{p(\hat{x}|x): \mathbb{E}[d(X,\hat{X})] \leq D} I(X; \hat{X})\)，其中 \(D\) 为可容忍的失真水平。对 Gauss 信源，\(R(D) = \frac{1}{2}\log_2(\sigma^2/D)\)（\(0 \leq D \leq \sigma^2\)），揭示了信号压缩中失真与码率之间的最优权衡。\textbf{网络信息论}（Cover-Thomas）将点对点通信模型推广到多用户场景，包括多址接入信道（MAC）、广播信道（BC）和中继信道，其容量区域的计算涉及叠加编码、速率分割和联合典型编码等高级技术。\textbf{量子信息论}（Schumacher 1995）将 Shannon 熵推广为 von Neumann 熵 \(S(\rho) = -\operatorname{Tr}(\rho \log_2 \rho)\)，建立了量子信源编码定理和量子信道容量的 Holevo 界，为量子计算和量子通信提供了信息论基础。信息论的方法和工具------特别是熵、互信息和典型性------已成为现代数学和工程中通用的分析语言，从机器学习中的信息瓶颈理论到统计力学中的最大熵原理，无一不体现其深远影响。

\hrulefill

\section{第87章：组合设计}\label{ux7b2c87ux7ae0ux7ec4ux5408ux8bbeux8ba1}

组合设计理论研究将元素安排成满足特定平衡条件的子集族。起源于统计学中的实验设计（Fisher, 1920s），在编码理论、密码学和计算机科学中有广泛应用。

\subsection{131.1 区组设计}\label{ux533aux7ec4ux8bbeux8ba1}

\textbf{定义 131.1}（平衡不完全区组设计 / BIBD）：一个 \((v, k, \lambda)\)-BIBD 是 \(v\) 元集 \(X\) 的一族 \(k\) 元子集（称为\textbf{区组}），使得 \(X\) 中任意两个不同元素恰出现在 \(\lambda\) 个区组中。

\textbf{命题 131.1}（BIBD 的必要条件）：若 \((v, k, \lambda)\)-BIBD 存在，参数满足： - \(r = \frac{\lambda(v-1)}{k-1}\) 是整数（每个元素出现的次数） - \(b = \frac{vr}{k}\) 是整数（区组总数） - Fisher 不等式：\(b \geq v\)（即 \(r \geq k\)）

\textbf{定义 131.2}（射影平面）：阶数为 \(n\) 的射影平面是 \((n^2+n+1, n+1, 1)\)-BIBD。等价于存在 \(n^2+n+1\) 个点和同样多条线，每线 \(n+1\) 点，每点 \(n+1\) 线，任意两点恰在一条线上，任意两线恰交于一点。已知当 \(n\) 为素数幂时存在，\(n=6\) 时不存在（Bruck-Ryser-Chowla 定理），\(n=10\) 时不存在（计算机搜索，1989）。

\textbf{定理 131.2}（Bruck-Ryser-Chowla 定理，1949）：若阶数为 \(n\) 的射影平面存在，且 \(n \equiv 1 \text{ 或 } 2 \pmod{4}\)，则 \(n\) 可表为两个整数的平方和。

\textbf{证明}：设 \(n\) 阶射影平面存在，其关联矩阵 \(A\)（\(v \times v\)，\(v = n^2+n+1\)）满足 \(A A^T = n I + J\)（\(J\) 为全 \(1\) 矩阵）。考虑有理二次型 \(Q(x) = x_1^2 + \cdots + x_v^2\)。由 Lagrange 四平方和定理，\(Q\) 在 \(\mathbb{Q}^v\) 上等价于 \(y_1^2 + \cdots + y_v^2\)。方程 \(A A^T = n I + J\) 在有理数域上给出 \(n\) 的表示条件。具体地，若 \(n \equiv 1, 2 \pmod{4}\)，由 Hasse-Minkowski 局部-整体原理，二次型 \(x_1^2 + \cdots + x_v^2 - n y^2\) 的各项不变量迫使 \(n\) 为两个平方和。其核心是：\(A A^T\) 的 Gram 矩阵结构蕴含 \(n\) 在 \(\mathbb{Q}\) 上可表为平方和，由 Legendre 三平方和定理的推广得 \(n = a^2 + b^2\)。\(\blacksquare\)

\subsection{131.2 正交拉丁方}\label{ux6b63ux4ea4ux62c9ux4e01ux65b9}

\textbf{定义 131.3}（拉丁方）：\(n\) 阶\textbf{拉丁方}是 \(n \times n\) 矩阵，每行和每列都是 \(\{1, \ldots, n\}\) 的排列。两个拉丁方称为\textbf{正交}的，如果叠加后每对 \((i,j)\) 恰好出现一次。

\textbf{定理 131.3}（正交拉丁方与射影平面的关系）：\(n-1\) 个两两正交的 \(n\) 阶拉丁方（MOLS）的存在性等价于 \(n\) 阶射影平面的存在性。

\textbf{证明}：设 \(n\) 阶射影平面存在，其点集为 \(\{P_0,\ldots,P_{n^2+n}\}\)，线集为 \(\{L_0,\ldots,L_{n^2+n}\}\)。固定线 \(L_\infty\) 上的 \(n+1\) 个点，将剩余 \(n^2\) 个点按通过 \(L_\infty\) 上两点的线族分组，得到 \(n-1\) 个平行类（每个对应 \(L_\infty\) 上一个非选定点）。每个平行类将 \(n^2\) 个点划分为 \(n\) 条线，每条线 \(n\) 个点。对每个平行类构造一个 \(n \times n\) 拉丁方：行对应一条参考线的 \(n\) 个点，列对应另一条参考线的 \(n\) 个点，格子 \((i,j)\) 填入通过行点 \(i\) 和列点 \(j\) 的线在平行类中的编号。这 \(n-1\) 个拉丁方两两正交。反之，给定 \(n-1\) 个 MOLS，可构造仿射平面并通过添加无穷远线得到射影平面。\(\blacksquare\)

\textbf{定理 131.4}（Euler 猜想与反例）：Euler 在 1782 年猜想不存在两个正交的 6 阶拉丁方（36 军官问题）。1901 年 Tarry 验证为真。Euler 进一步猜想对 \(n \equiv 2 \pmod{4}\) 均无正交拉丁方。1959 年 Bose-Shrikhande-Parker 推翻此猜想，构造了 10 阶正交拉丁方，并证明对 \(n \neq 2, 6\) 均存在。

\textbf{证明}（要点概述）：(1) \(n=2\) 时，两个 \(2\) 阶拉丁方只有两种可能，经枚举可知它们不可能正交。(2) \(n=6\) 时，若存在两个正交的 \(6\) 阶拉丁方，则它们可转化为 \(6\) 阶射影平面的等价问题。由 Bruck-Ryser-Chowla 定理（定理 131.2），\(n=6\) 时 \(n \equiv 2 \pmod{4}\)，\(6\) 不能表示为两个整数的平方和，故 \(6\) 阶射影平面不存在，从而不存在两个正交的 \(6\) 阶拉丁方。Tarry 在 1901 年通过穷举搜索验证了这一结论。(3) \(n=10\) 时，Bose、Shrikhande 和 Parker 在 1959 年利用有限域和差集构造出了两个正交的 \(10\) 阶拉丁方，从而推翻了 Euler 对 \(n \equiv 2 \pmod{4}\) 的一般性猜想。最终结果：对 \(n \neq 2, 6\)，总存在一对正交的 \(n\) 阶拉丁方。\(\blacksquare\)

\subsection{131.3 Steiner 系与 t-设计}\label{steiner-ux7cfbux4e0e-t-ux8bbeux8ba1}

\textbf{定义 131.4}（\(t\)-设计）：一个 \(t\)-\((v, k, \lambda)\) 设计是 \(v\) 元集的一族 \(k\) 元子集，使得任意 \(t\) 个不同元素恰出现在 \(\lambda\) 个区组中。BIBD 是 \(t=2\) 的特例。

\textbf{定义 131.5}（Steiner 系）：\(S(t, k, v)\) 是 \(t\)-\((v, k, 1)\) 设计。即任意 \(t\) 个元素恰出现在 1 个区组中。

\textbf{定理 131.5}（大设计猜想 / 已解决）：对 \(t \geq 6\)，仅存在平凡的 \(t\)-设计（即所有 \(k\) 元子集均取）。

\textbf{证明}：该结果由 Magliveras-Leavitt (1983) 以及 Teirlinck (1987) 等人的工作最终确立。核心思路：对 \(t \geq 6\)，利用 Kramer-Mesner 方法，将 \(t\)-设计的存在性转化为求解整数方程组 \(M_{t,k}^v \cdot \mathbf{x} = \lambda \mathbf{1}\)，其中 \(M_{t,k}^v\) 为关联矩阵。当 \(t \geq 6\) 时，该方程组的解空间维数由 Ray-Chaudhuri-Wilson 定理控制：对任意 \(t\)-\((v, k, \lambda)\) 设计，区组数 \(b\) 满足 \(b \geq \binom{v}{\lfloor t/2 \rfloor}\)。结合整数规划约束，当 \(t \geq 6\) 且 \(k < v\) 时，仅有平凡解（所有 \(k\) 元子集均取）满足条件。这一结论的最终严格证明依赖于 Wilson 的基础构作法（Wilson's fundamental construction）与有限单群分类定理的交互。\(\blacksquare\)

\subsection{131.4 有限域与组合设计}\label{ux6709ux9650ux57dfux4e0eux7ec4ux5408ux8bbeux8ba1}

\textbf{定理 131.6}（有限域上的射影空间）：\(\mathbb{F}_q\) 上的 \(d\) 维射影空间 \(\operatorname{PG}(d, q)\) 中，点数为 \(\frac{q^{d+1}-1}{q-1}\)，每个 \(k\) 维子空间的点数为 \(\frac{q^{k+1}-1}{q-1}\)。\(\operatorname{PG}(2, q)\) 是 \(q\) 阶射影平面。

\textbf{证明}：\(\operatorname{PG}(d, q)\) 定义为 \(\mathbb{F}_q^{d+1}\) 中所有一维子空间的集合。\(\mathbb{F}_q^{d+1}\) 共有 \(q^{d+1}\) 个向量，减去零向量余 \(q^{d+1}-1\) 个非零向量。每个一维子空间含 \(q-1\) 个非零向量，故点数为 \((q^{d+1}-1)/(q-1)\)。\(k\) 维子空间对应 \(\mathbb{F}_q^{d+1}\) 的 \((k+1)\) 维子空间，非零向量数为 \(q^{k+1}-1\)，故其包含的点数为 \((q^{k+1}-1)/(q-1)\)。\(\operatorname{PG}(2, q)\) 的点数为 \((q^3-1)/(q-1) = q^2+q+1\)，线数为 \((q^2+q+1)\)，每线含 \(q+1\) 个点，每点处有 \(q+1\) 条线，任意两点在唯一一条线上，故为 \(q\) 阶射影平面。\(\blacksquare\)

\textbf{定理 131.7}（差集与 Singer 定理）：\(\operatorname{PG}(d, q)\) 的自同构群在点上有一个正则循环子群（Singer 循环），对应一个\textbf{差集}------这是构造对称设计的标准方法。

\textbf{证明}：考虑 \(\mathbb{F}_{q^{d+1}}\) 的乘法群 \(\mathbb{F}_{q^{d+1}}^\times\)，它是 \(q^{d+1}-1\) 阶循环群。\(\operatorname{PG}(d, q)\) 的点可等同于 \(\mathbb{F}_{q^{d+1}}^\times / \mathbb{F}_q^\times\)，即商群 \(G\)，其阶为 \(v = (q^{d+1}-1)/(q-1)\)。乘法群 \(\mathbb{F}_{q^{d+1}}^\times\) 中一个生成元 \(\alpha\) 诱导 \(G\) 上的一个 \(v\) 阶循环置换 \(\sigma\)（Singer 循环）。取 \(\sigma\) 下一条线的轨道：设 \(D = \{i \in \mathbb{Z}_v : \alpha^i \in L\}\) 为某条线 \(L\) 的坐标集，则 \(|D| = (q+1)/(q-1) = q+1\)（当 \(d=2\)）或一般地 \(|D| = (q^{k+1}-1)/(q-1)\)。由 \(\sigma\) 的正则性，\(D\) 是 \(\mathbb{Z}_v\) 中的差集，即每个非零元素恰有 \(\lambda\) 种方式表为 \(D\) 中两元素之差，这给出了对称 \((v, k, \lambda)\) 设计的构造。\(\blacksquare\)

\textbf{定义 131.6}（Hadamard 矩阵）：\(n\) 阶 Hadamard 矩阵 \(H\) 是 \(n \times n\) 的 \(\pm 1\) 矩阵，满足 \(H H^T = n I_n\)。Hadamard 矩阵存在当且仅当 \(n = 1, 2\) 或 \(n \equiv 0 \pmod{4}\)（Hadamard 猜想，未完全解决）。

\textbf{定理 131.8}（Hadamard 设计与 Hadamard 矩阵的等价性）：\((4m-1, 2m-1, m-1)\) 对称 BIBD 的存在性等价于 \(4m\) 阶 Hadamard 矩阵的存在性。

\textbf{证明}：设 \(H\) 为 \(4m\) 阶 Hadamard 矩阵，标准化使得第一行和第一列全为 \(1\)。删除第一行和第一列，将剩余 \((4m-1) \times (4m-1)\) 子矩阵中的 \(-1\) 替换为 \(0\)，\(1\) 替换为 \(1\)，得到关联矩阵 \(A\)。\(H H^T = 4m I\) 蕴含 \(A\) 的每行恰有 \(2m-1\) 个 \(1\)（因原行有 \(2m\) 个 \(1\) 和 \(2m\) 个 \(-1\)，标准化后除去第一列的 \(1\)），且任意两行的内积为 \(m-1\)（因 \(H\) 中两行内积为 \(0\)，标准化后两行均有 \(2m\) 个 \(1\) 和 \(2m\) 个 \(-1\)，重合位置为 \(m\) 个同号，扣除第一列后余 \(m-1\)）。故 \(A\) 是 \((4m-1, 2m-1, m-1)\) 对称 BIBD 的关联矩阵。反之，从该设计的关联矩阵 \(A\) 可重构 \(H\)：补全 \(1\) 的行和列，将 \(0\) 换为 \(-1\)。\(\blacksquare\)

\hrulefill

\subsection{组合设计的前沿与未解问题}\label{ux7ec4ux5408ux8bbeux8ba1ux7684ux524dux6cbfux4e0eux672aux89e3ux95eeux9898}

组合设计理论在 20 世纪后半叶取得了巨大进展，但仍有大量核心问题未解决。\textbf{Hadamard 猜想}（1893）断言对任意正整数 \(m\)，存在 \(4m\) 阶 Hadamard 矩阵。目前已知最小未解决的阶数为 \(668\)（即 \(m=167\)）。通过 Paley 构造（利用有限域上的二次特征）和 Williamson 构造（利用四元数代数），Hadamard 矩阵的构造问题与有限域理论、差集和编码理论有深刻联系。\textbf{大设计猜想}（所有 \(t \geq 6\) 的非平凡 \(t\)-设计均不存在）已被 Teirlinck（1987）以构造性方式否定------他构造了 \(t \geq 6\) 的非平凡单纯 \(t\)-设计，但区组数极大，参数的非平凡性仅体现在组合意义上。\textbf{有限射影平面}的存在性仅对 \(n\) 为素数幂已知，\(n = 12\) 是第一个既非素数幂又未被完全排除的阶数。Bruck-Ryser-Chowla 定理排除了 \(n \equiv 1, 2 \pmod{4}\) 且不可表为两个平方和的阶数，但 \(n = 12\)（\(12 \equiv 0 \pmod{4}\)）不受此限制，至今未知是否存在 12 阶射影平面。\textbf{MOLS 猜想}（Euler 的 36 军官问题）已完全解决：对 \(n \neq 2, 6\)，总存在一对正交拉丁方。但 \(N(n)\)（\(n\) 阶两两正交拉丁方的最大个数）的精确值仅对少数 \(n\) 已知，\(N(n) \leq n-1\)，等号成立当且仅当 \(n\) 阶射影平面存在。组合设计理论还包含\textbf{有限几何}（如广义四边形、极空间）和\textbf{关联结构}（如 \(t\)-设计、平衡不完全区组设计、可分组设计）的丰富理论，后者在统计实验设计、纠错码和密码学中有广泛应用。组合设计的现代研究大量使用代数方法（有限域、群论、表示论）和计算工具，出现了许多新的构造技术，如 Wilson 的基础构作法（1972-1975）和 Lam 的大规模计算机搜索（1989 年证明不存在 10 阶射影平面）。

\hrulefill

\subsection{加法组合学}\label{ux52a0ux6cd5ux7ec4ux5408ux5b66}

加法组合学（Additive Combinatorics）研究加法群中子集的加性结构，是现代组合数学的核心分支。其基本问题可追溯至 Lagrange（1770）和 Cauchy（1813），但系统的理论框架建立于 20 世纪后半叶，以 Szemerédi、Freiman、Ruzsa 和 Gowers 等人的工作为里程碑。

\subsection{和集与差集}\label{ux548cux96c6ux4e0eux5deeux96c6}

\textbf{定义 1}（和集与差集）：设 \((G, +)\) 为加法群，\(A, B \subseteq G\) 为非空子集。定义

\[A + B = \{a + b : a \in A,\; b \in B\}, \qquad A - B = \{a - b : a \in A,\; b \in B\}.\]

记 \(kA = A + A + \cdots + A\)（\(k\) 次）。若 \(A = B\)，则 \(A + A\) 和 \(A - A\) 分别称为\textbf{和集}与\textbf{差集}。

\textbf{命题 1}（平凡界）：设 \(A, B \subseteq G\) 为有限子集，则

\[\max(|A|, |B|) \leq |A + B| \leq |A||B|,\]

且 \(|A + B| = |A||B|\) 当且仅当所有和 \(a + b\) 两两不同。

\emph{证明}：下界取 \(b_0 \in B\)，则 \(|A + B| \geq |A + b_0| = |A|\)；上界为乘积集的基数显然。等号条件：若存在 \((a_1, b_1) \neq (a_2, b_2)\) 使得 \(a_1 + b_1 = a_2 + b_2\)，则 \(|A + B| < |A||B|\)；反之若所有和不同，则 \(|A + B| = |A||B|\)。\(\blacksquare\)

\textbf{定义 2}（加倍常数）：子集 \(A \subseteq G\) 的\textbf{加倍常数}为 \(\sigma[A] = |A + A| / |A|\)。称 \(A\) 具有\textbf{小加倍}若 \(\sigma[A]\) 有界。

\subsection{Ruzsa 三角不等式}\label{ruzsa-ux4e09ux89d2ux4e0dux7b49ux5f0f}

\textbf{定理 1}（Ruzsa 三角不等式，1994）：设 \(A, B, C \subseteq G\) 为有限非空子集，则

\[|A||B - C| \leq |A - B||A - C|.\]

\emph{证明}：对每个 \(d \in B - C\)，固定一个表示 \(d = b_d - c_d\)（其中 \(b_d \in B, c_d \in C\)）。构造映射

\[\varphi : A \times (B - C) \to (A - B) \times (A - C), \quad \varphi(a, d) = (a - b_d, a - c_d).\]

验证 \(\varphi\) 是单射：若 \((a - b_d, a - c_d) = (a' - b_{d'}, a' - c_{d'})\)，则相减得 \((b_d - c_d) - (b_{d'} - c_{d'}) = 0\)，故 \(d = d'\)，进而 \(b_d = b_{d'}, c_d = c_{d'}\)；由 \(a - b_d = a' - b_d\) 知 \(a = a'\)。因此

\[|A||B - C| = |A \times (B - C)| \leq |(A - B) \times (A - C)| = |A - B||A - C|.\]

\(\blacksquare\)

\textbf{推论 1}：取 \(B = C = A\) 得 \(|A||A - A| \leq |A - A|^2\)（平凡）。取 \(B = A, C = -A\) 得 \(|A||A + A| \leq |A - A|^2\)。这表明 \(|A + A|\) 和 \(|A - A|\) 通过 Ruzsa 三角不等式等价控制：

\[\frac{|A - A|}{|A|} \leq \left(\frac{|A + A|}{|A|}\right)^2.\]

\subsection{Plünnecke 不等式}\label{pluxfcnnecke-ux4e0dux7b49ux5f0f}

\textbf{定理 2}（Plünnecke 不等式，1970；Ruzsa 简化，1989）：设 \(A, B \subseteq G\) 为有限子集，满足 \(|A + B| = \alpha|A|\)。则对任意非负整数 \(h, k\)，

\[|hB - kB| \leq \alpha^{h+k} |A|.\]

特别地，取 \(B = A\) 且 \(|A + A| = K|A|\)，则 \(|kA| \leq K^k |A|\)。

Plünnecke 不等式是加法组合学中迭代估计的基石。其核心思想是：若 \(B\) 对 \(A\) 的加法扩张受控（即 \(|A+B|\) 作为 \(|A|\) 的倍数有界），则 \(B\) 的任意多次加法与减法的组合扩张也受控。这一结果在直觉上并非显然------单独的 \(|A+B|\) 控制并不意味着 \(|2B-2B|\) 同样受控，因为 \(B\) 自身可能具有复杂的结构。Plünnecke 的原创证明（1970）使用了复杂的图论方法，构造了分层有向图（Plünnecke 图），通过 Menger 定理和 Hall 婚配定理的结合导出指数界。Ruzsa（1989）给出了一个更简洁的证明，利用 Plünnecke 图的边不相交路径性质，避免了原证明中的大量组合计算。该不等式在 Freiman 定理的证明中扮演关键角色------它允许将 \(|A+A|\) 的有界性转化为 \(|kA|\) 的多项式界，从而将加法结构问题转化为几何覆盖问题。

\textbf{证明}：构造 Plünnecke 图：顶点集为 \(\bigcup_{i=0}^{h+k} (A + iB)\)，从 \(x \in A + iB\) 到 \(y \in A + (i+1)B\) 连有向边当且仅当 \(y - x \in B\)。由条件 \(|A+B| = \alpha |A|\)，零层到一层的边密度受控。选取 \(A\) 中元 \(a_0\)，考虑从 \(a_0\) 出发到第 \(h\) 层的路径集。由 Menger 定理，存在 \(|A|\) 条边不相交路径从第 \(0\) 层到第 \(h\) 层。利用 Hall 婚配定理，可在每层选取子集使得前向边构成双射，从而第 \(h\) 层的顶点数 \(\leq \alpha^h |A|\)。对称地考虑反向路径，第 \(k\) 层的顶点数 \(\leq \alpha^k |A|\)。\(hB - kB\) 的基数不超过第 \(h+k\) 层顶点数，故 \(|hB - kB| \leq \alpha^{h+k} |A|\)。\(\blacksquare\)

\subsection{Freiman 定理}\label{freiman-ux5b9aux7406}

\textbf{定理 3}（Freiman 定理，1973；Ruzsa 改进，1994）：存在绝对常数 \(C > 0\) 和 \(d \geq 1\) 使得：若 \(A \subseteq \mathbb{Z}\) 是有限子集且 \(|A + A| \leq K|A|\)，则 \(A\) 包含在一个维数不超过 \(d = \lfloor K - 1 \rfloor\) 的广义算术级数

\[P = \left\{a_0 + \sum_{i=1}^d n_i a_i : 0 \leq n_i < L_i\right\}\]

中，且 \(|P| \leq e^{CK^{C}} |A|\)。

\textbf{证明}（Ruzsa 三步法）：

\textbf{第 1 步：Bogolyubov 型引理。} 设 \(|A+A| \leq K|A|\)。由 Plünnecke 不等式，\(|2A-2A| \leq K^4 |A|\)。通过 Fourier 分析：定义 \(\hat{1}_A(\xi) = \sum_{a \in A} e^{2\pi i a \xi}\)，大频谱 \(\operatorname{Spec}_\varepsilon(A) = \{\xi \in \mathbb{R}/\mathbb{Z} : |\hat{1}_A(\xi)| \geq \varepsilon |A|\}\) 满足 \(|\operatorname{Spec}_\varepsilon(A)| \leq K/\varepsilon^2\)。取 \(\varepsilon = 1/(2K)\)，则 \(2A-2A\) 包含 Bohr 集 \(B(\operatorname{Spec}_\varepsilon(A), \tfrac{1}{4}) = \{x : \|x \xi\| \leq \tfrac{1}{4}, \forall \xi \in \operatorname{Spec}_\varepsilon(A)\}\)，其大小为 \(|B| \geq (8K)^{-O(K^2)} |A|\)。

\textbf{第 2 步：Bohr 集到广义算术级数。} 将 \(A\) 嵌入循环群 \(\mathbb{Z}_N\)（\(N\) 充分大），其上 \(d\) 维 Bohr 集可嵌入 \(\mathbb{R}^d\) 中，由 Minkowski 第二定理，该 Bohr 集被一个维数 \(\leq d\) 的广义算术级数 \(P\) 覆盖，且 \(|P| \leq C^d |B|\)。这里 \(d = |\operatorname{Spec}_\varepsilon(A)|\)。

\textbf{第 3 步：提升回 \(\mathbb{Z}\)。} 在 \(\mathbb{Z}_N\) 中得到的广义算术级数 \(P\) 覆盖了 \(A\) 的副本，提升回 \(\mathbb{Z}\) 后 \(A\) 包含在对应的广义算术级数中，维数 \(d \leq C K^{O(1)}\)，大小 \(|P| \leq \exp(K^{O(1)}) |A|\)。目前最佳界：\(d \leq K^{1+o(1)}\)（Sanders 2012），\(|P| \leq \exp(K^{3+o(1)}) |A|\)。\(\blacksquare\)

\subsection{Szemerédi 定理}\label{szemeruxe9di-ux5b9aux7406}

\textbf{定理 4}（Szemerédi 定理，1975）：对任意整数 \(k \geq 3\) 和 \(\delta > 0\)，存在 \(N(k, \delta) \in \mathbb{N}\) 使得：若 \(N \geq N(k, \delta)\) 且 \(A \subseteq \{1, \ldots, N\}\) 满足 \(|A| \geq \delta N\)，则 \(A\) 包含长度为 \(k\) 的等差数列。

\emph{Roth 定理（\(k = 3\)）的 Fourier 分析证明概述}：设 \(A \subseteq \mathbb{Z}_N\)（\(N\) 为素数，避免边界效应），\(|A| = \delta N\)。定义 \(A\) 的平衡函数 \(f = 1_A - \delta\)。三元等差数列的计数可表为 Fourier 系数的卷积：

\[\Lambda(A) = \sum_{x, d} 1_A(x) 1_A(x+d) 1_A(x+2d) = N^{-1} \sum_{\xi} \hat{1}_A(\xi)^2 \hat{1}_A(-2\xi).\]

若 \(A\) 不含三元等差数列，则 \(\Lambda(A) = \Lambda(\delta N, \delta N, \delta N)\)（纯随机情形），导致 \(\sum_{\xi \neq 0} |\hat{f}(\xi)|^3 \geq c \delta^3 N^2\)。由 \(L^2\) 范数约束 \(\sum_{\xi} |\hat{f}(\xi)|^2 = \delta(1-\delta)N\)，知存在一个非零 Fourier 系数满足 \(|\hat{f}(\xi)| \geq c' \delta^2 N\)，这表明 \(A\) 在某个长为 \(N/|\xi|\) 的算术级数上有密度增量 \(\delta + c \delta^2\)。迭代此过程（密度增量论证）得 \(N\) 的上界与 \(\delta\) 的关系 \(N \leq \exp(\exp(C/\delta))\)，证毕。\(\blacksquare\)

\emph{一般 \(k\) 的 Furstenberg 遍历论证明概述（1977）}：设 \(X = \{0, 1\}^{\mathbb{Z}}\) 为 \(\mathbb{Z}\) 在移位 \(\sigma\) 下的 Bernoulli 系统的闭子移位，点 \(\omega \in X\) 编码为 \(A\) 的特征函数序列。定义遍历平均

\[\frac{1}{N} \sum_{n=1}^N \int f \cdot \sigma^n f \cdot \cdots \cdot \sigma^{(k-1)n} f \, d\mu\]

（其中 \(f = 1_A\)）。由多重遍历定理（Furstenberg 结构理论），该平均的极限 \(> 0\) 当且仅当存在具有 \(k\) 项等差数列的结构。通过将动力系统分解为 Kronecker 因子和弱混合扩展，逐层归纳：若系统无 \(k\)-AP，则通过与紧交换群的因子映射导出矛盾。定性形式下 \(N(k, \delta)\) 的存在性得证。Gowers (2001) 给出了定量界 \(N(k, \delta) \leq \exp(\exp(\delta^{-c_k}))\)，用到高阶 Fourier 分析和 Gowers 范数。\(\blacksquare\)

\subsection{Green-Tao 定理}\label{green-tao-ux5b9aux7406}

\textbf{定理 5}（Green-Tao 定理，2008）：素数集合包含任意长的等差数列。即对任意 \(k \geq 3\)，存在无穷多长度为 \(k\) 的素数等差数列。

\textbf{证明}：采用 Green-Tao 三步法。

\textbf{第 1 步：W-技巧。} 令 \(W = \prod_{p \leq w} p\)（\(w = \log \log N\) 充分大慢增）。将素数限制在与 \(W\) 互素的剩余类中，消除小素数同余障碍。设 \(\widetilde{\mathbb{P}} = \{n \leq N/W : Wn + 1 \text{ 为素数}\}\)，则 \(|\widetilde{\mathbb{P}}| \sim \frac{W}{\phi(W)} \frac{N/W}{\log N}\)，由 Dirichlet 特征和及 Siegel-Walfisz 定理，密度 \(\delta = |\widetilde{\mathbb{P}}|/(N/W) \sim 1/\log N\)。

\textbf{第 2 步：伪随机测度构造。} 构造 \(\mathbb{Z}_N\) 上的函数 \(\nu(n) = \frac{\phi(W)}{W} \log R \cdot \Lambda_R(Wn+1)^2\)（截断 Selberg 筛权，\(R = N^{1/(k+1)}\)）。验证：(a) \(\mathbb{E} \nu = 1 + o(1)\)；(b) \(\nu(n) \geq c_k \cdot 1_{\widetilde{\mathbb{P}}}(n)\)（主导性质）；(c) \(\|\nu-1\|_{U^{k-1}} = o(1)\)（Gowers 范数伪随机性），由 Goldston-Yıldırım 筛法估计保证。

\textbf{第 3 步：转移原理。} 设 \(f(n) = c_k \cdot 1_{\widetilde{\mathbb{P}}}(n)\)，则 \(0 \leq f \leq \nu\) 且 \(\mathbb{E} f \geq \delta\)。将 \(f\) 正交分解为 \(f = f_U + f_{U^\perp}\)（Gowers 反均匀性分解），其中 \(f_U = \mathbb{E}(f \mid \mathcal{B})\)（\(\mathcal{B}\) 为 \(U^{k-1}\)-范数控制因子）。由 Szemerédi 定理（有界情形），\(f_U\) 包含 \(k\)-AP；由 \(\nu\) 的伪随机性，\(f_{U^\perp}\) 的贡献可忽略。故 \(\widetilde{\mathbb{P}}\) 包含 \(k\)-AP，从而素数包含任意长的等差数列。\(\blacksquare\)

\subsection{Bourgain-Katz-Tao Sum-Product 定理}\label{bourgain-katz-tao-sum-product-ux5b9aux7406}

\textbf{定理 6}（Bourgain-Katz-Tao，2004）：存在绝对常数 \(\varepsilon > 0\) 使得：对任意素数 \(p\) 和 \(A \subseteq \mathbb{F}_p\) 满足 \(p^\delta < |A| < p^{1-\delta}\)（\(\delta > 0\)），有

\[\max(|A + A|, |A \cdot A|) \geq |A|^{1+\varepsilon}.\]

该定理表明 \(\mathbb{F}_p\) 中任何子集不可能同时在小加法和小乘法下具有小扩张------要么和集大，要么积集大。定理的关键推论是：\(\mathbb{F}_p\) 中的指数和 \(\sum_{a \in A} e^{2\pi i a x / p}\) 具有非平凡的界。证明融合了加法组合学（Balog-Szemerédi-Gowers）和代数几何（Lefschetz 原理、Katz-Tao 多项式方法），常数 \(\varepsilon\) 依赖于 \(\delta\)。

Sum-Product 现象是加法组合学中最深刻的发现之一，其根源可追溯至 Erdős 和 Szemerédi 对整数中加法与乘法相互作用的研究。在整数环 \(\mathbb{Z}\) 中，Erdős-Szemerédi 猜想（1983）断言 \(\max(|A+A|, |A \cdot A|) \geq |A|^{2-o(1)}\)，目前最佳界约为 \(|A|^{4/3+o(1)}\)。在有限域 \(\mathbb{F}_p\) 中，Bourgain-Katz-Tao 定理（2004）首次证明了当 \(|A|\) 在 \(p^\delta\) 和 \(p^{1-\delta}\) 之间时存在非平凡的 sum-product 扩张。该定理的深远影响体现在多个领域：在极值组合中，它给出了有限域上点-线关联问题的改进界；在理论计算机科学中，它被用于构造显式随机性提取器（extractors）和伪随机生成器；在解析数论中，它改进了 Kloosterman 和等指数和的估计。这一结果后来被 Bourgain-Glibichuk-Konyagin（2006）推广到任意有限域，并被 Rudnev-Stevens（2022）将常数 \(\varepsilon\) 改进至至少 \(1/13\)。

\textbf{证明}：反设 \(\max(|A+A|, |A \cdot A|) < |A|^{1+\varepsilon}\)。由 Balog-Szemerédi-Gowers 定理，存在大子集 \(A' \subseteq A\) 使得 \(|A'+A'|, |A' \cdot A'| \leq |A|^{1+O(\varepsilon)}\)。考虑多项式 \(P(x,y) = f(x) + g(y)\)（其中 \(f\) 为适当选取的多项式），其零点集 \(Z = \{(x,y) \in A' \times A' : P(x,y) = 0\}\) 的大小由 Schwartz-Zippel 引理控制在 \(O(|A'|)\)。另一方面，由和-积展开的缓慢增长，\(Z\) 的大小至少为 \(|A'|^{1+\delta}\)（\(\delta > 0\)）。当 \(\varepsilon\) 充分小时，\(|A'|^{1+\delta} > O(|A'|)\)，矛盾。\(\varepsilon\) 的绝对常数存在性由 Bourgain 的多尺度归纳论证保证，目前最佳 \(\varepsilon \geq 1/13\)（Rudnev-Stevens 2022）。\(\blacksquare\)

\subsection{Balog-Szemerédi-Gowers 定理}\label{balog-szemeruxe9di-gowers-ux5b9aux7406}

\textbf{定理 7}（Balog-Szemerédi-Gowers，1994/1998）：设 \(A, B \subseteq G\)（加法群）为有限子集且 \(|A| = |B| = n\)。若和集 \(A + B\) 满足 \(|A + B| \leq K n\)，则存在子集 \(A' \subseteq A, B' \subseteq B\) 满足 \(|A'| \geq n/K, |B'| \geq n/K\) 且 \(|A' + B'| \leq K^{O(1)} n\)。更一般地，若加法能量 \(E(A, B) = |\{(a_1, b_1, a_2, b_2) : a_1 + b_1 = a_2 + b_2\}| \geq n^3 / K\)，则存在稠密子集 \(A', B'\) 使 \(|A' + B'| \leq K^{O(1)} n\)。

\textbf{证明}：假设 \(E(A,B) \geq n^3/K\)。定义图 \(G = (A \times B, E)\)，其中 \((a,b) \in E\) 当且仅当 \(a+b \in S\)（\(S\) 为 \(A+B\) 中高重数和的集合）。由加法能量假设，\(|E| \geq n^2/(2K)\)。由 Balog-Szemerédi 图论引理，存在 \(A' \subseteq A\)，\(B' \subseteq B\)，\(|A'|, |B'| \geq n/(8K^2)\)，使得 \(A' \times B'\) 几乎完全包含在 \(E\) 中。对任意 \(a \in A'\)，\(b, b' \in B'\)，有 \((a, b), (a, b') \in E\)，故 \(a+b, a+b' \in S\)。由 Gowers 路径论证：\(a+b' = (a+b) - (b) + (b')\)，因此 \(|A'+B'| \leq |A'+B| \cdot |B| \leq K^4 n\)。更精细的界可得 \(|A'+B'| \leq K^8 n\)。若初始条件为 \(|A+B| \leq Kn\)，则 \(E(A,B) \geq n^3/K\)，应用上述论证即得结论。\(\blacksquare\)

\hrulefill

\subsection{组合物种理论}\label{ux7ec4ux5408ux7269ux79cdux7406ux8bba}

组合物种理论（Theory of Combinatorial Species）由 André Joyal 于 1981 年引入，为计数组合学提供了范畴论统一框架。其核心思想是：组合结构不仅是计数的对象，而是带有对称群作用的结构族，生成函数等计数工具天然地从这些结构的函子性质中涌现。

\subsection{范畴论定义}\label{ux8303ux7574ux8bbaux5b9aux4e49}

\textbf{定义 1}（组合物种）：设 \(\mathbb{B}\) 为有限集与双射构成的小范畴。一个\textbf{组合物种}是一个函子

\[F : \mathbb{B} \to \mathbf{Set},\]

满足：

\begin{enumerate}
\def\labelenumi{\arabic{enumi}.}
\tightlist
\item
  对每个有限集 \(U\)，\(F[U]\) 是 \(U\) 上的 \(F\)-结构集合；
\item
  对每个双射 \(\sigma : U \to V\)，\(F[\sigma] : F[U] \to F[V]\) 是沿 \(\sigma\) 的结构运输映射；
\item
  函子性：\(F[\operatorname{id}_U] = \operatorname{id}_{F[U]}\)，且 \(F[\tau \circ \sigma] = F[\tau] \circ F[\sigma]\)。
\end{enumerate}

\textbf{命题 1}（对称群表示）：对基数为 \(n\) 的集合 \([n] = \{1, \ldots, n\}\)，\(F[n]\) 带有对称群 \(\mathfrak{S}_n\) 的作用：

\[\rho_n : \mathfrak{S}_n \to \operatorname{Aut}(F[n]), \quad \rho_n(\sigma) = F[\sigma],\]

使得 \(|F[n]| = f_n\) 是 \(n\) 元标记结构数。\(\{f_n\}\) 的生成函数为物种的指数生成函数（见下文）。

\subsection{基本物种}\label{ux57faux672cux7269ux79cd}

以下是 6 个基本组合物种，每个均以函子形式给出：

\textbf{定义 2}（空物种 \(0\)）：\(0[U] = \varnothing\)（空集），\(0[\sigma]\) 为唯一可能映射。无任何结构。

\textbf{定义 3}（空集物种 \(1\)）：\(1[\varnothing] = \{\star\}\)，当 \(U \neq \varnothing\) 时 \(1[U] = \varnothing\)。\(1[\sigma]\) 为唯一可能映射。唯一的空集上结构。

\textbf{定义 4}（单元素物种 \(X\)）：\(X[U] = \{u_0\}\) 若 \(|U| = 1\)（且 \(u_0\) 为 \(U\) 中唯一元素），否则 \(X[U] = \varnothing\)。\(X[\sigma]\) 将单元素结构沿 \(\sigma\) 运输。

\textbf{定义 5}（集合物种 \(E\)）：\(E[U] = \{U\}\)（单一元素，即集合 \(U\) 本身），\(E[\sigma](U) = \sigma(U) = V\)。所有元素无区分的集合，其结构就是集合本身。

\textbf{定义 6}（置换物种 \(S\)）：\(S[U] = \{\text{$U$ 上的所有置换}\}\)。对双射 \(\sigma : U \to V\)，\(S[\sigma](\pi) = \sigma \circ \pi \circ \sigma^{-1}\)（共轭运输）。

\textbf{定义 7}（循环置换物种 \(C\)）：\(C[U] = \{\text{$U$ 上的循环置换}\}\)。运输同 \(S\) 的限制。注意 \(C \neq S\)：\(|S[n]| = n!\)，\(|C[n]| = (n-1)!\)（\(n \geq 1\)）。

\subsection{物种运算}\label{ux7269ux79cdux8fd0ux7b97}

\textbf{定义 8}（和 \(F + G\)）：\((F + G)[U] = F[U] \sqcup G[U]\)（不相交并）。函子作用逐分量：\((F + G)[\sigma](s) = F[\sigma](s)\) 若 \(s \in F[U]\)，否则 \(G[\sigma](s)\)。

\textbf{定义 9}（积 \(F \cdot G\)）： \[(F \cdot G)[U] = \bigsqcup_{U_1 \sqcup U_2 = U} F[U_1] \times G[U_2].\]

即 \(U\) 的一个 \(F \cdot G\)-结构是将 \(U\) 划分为两个不相交子集 \(U_1, U_2\)，在 \(U_1\) 上放置 \(F\)-结构，在 \(U_2\) 上放置 \(G\)-结构。函子作用：\((\sigma : U \to V)\) 将划分 \(\{U_1, U_2\}\) 映射到 \(\{\sigma(U_1), \sigma(U_2)\}\)，并沿限制双射运输各分量上的结构。

\textbf{定义 10}（复合 \(F \circ G\)）： \[(F \circ G)[U] = \bigsqcup_{\pi \in \operatorname{Par}(U)} \left( F[\pi] \times \prod_{B \in \pi} G[B] \right),\]

其中 \(\operatorname{Par}(U)\) 是 \(U\) 的所有集合划分。\(U\) 的一个 \((F \circ G)\)-结构是：先将 \(U\) 划分为若干块，每块上放置 \(G\)-结构，再在块集上放置 \(F\)-结构。这对应了 \(F\) 在每个 \(G\)-结构块上的复合。

\textbf{定义 11}（导数 \(F'\)）：\((F')[U] = F[U \sqcup \{\star_U\}]\)，其中 \(\star_U \notin U\) 是一个区分的新元素。对双射 \(\sigma : U \to V\)，\(F'[\sigma] = F[\sigma \sqcup \operatorname{id}_{\{\star\}}]\)。导数捕捉了''指定一个独特点''的操作------一个 \(F'\)-结构等价于在 \(U\) 加上一个标记元素后施加 \(F\)。

\textbf{命题 2}（运算与指数生成函数的对应）：设物种 \(F\) 的指数生成函数为

\[F(x) = \sum_{n \geq 0} f_n \frac{x^n}{n!}, \quad f_n = |F[n]|.\]

则以下对应关系成立：

\begin{itemize}
\tightlist
\item
  \((F + G)(x) = F(x) + G(x)\)；
\item
  \((F \cdot G)(x) = F(x) \cdot G(x)\)；
\item
  \((F \circ G)(x) = F(G(x))\)（复合）；
\item
  \(F'(x) = \frac{d}{dx} F(x)\)（导数）。
\end{itemize}

\textbf{证明}：对 \(F \cdot G\)，由定义，\((F \cdot G)[n]\) 的结构数为将 \([n]\) 划分为两子集 \(U_1\)（大小 \(k\)）和 \(U_2\)（大小 \(n-k\)）并在两子集上分别放置 \(F\)-和 \(G\)-结构。\(U_1\) 有 \(\binom{n}{k}\) 种选择，故 \(f \cdot g_n = \sum_{k=0}^n \binom{n}{k} f_k g_{n-k}\)。两边乘以 \(x^n/n!\) 并求和：\(\sum_{n} f \cdot g_n \frac{x^n}{n!} = \sum_{n} \sum_{k=0}^n \frac{f_k}{k!} \frac{g_{n-k}}{(n-k)!} x^n = (\sum f_k \frac{x^k}{k!})(\sum g_j \frac{x^j}{j!})\)。对 \(F \circ G\)，\((F \circ G)[n]\) 的结构数为 \(\sum_{\pi \in \operatorname{Par}(n)} f_{|\pi|} \prod_{B \in \pi} g_{|B|}\)，由 Faà di Bruno 公式，这恰为 \(F(G(x))\) 的 Taylor 系数。对 \(F'\)，\(f'_n = |F'[n]| = |F[n \cup \{\star\}]| = f_{n+1}\)，故 \(F'(x) = \sum f_{n+1} \frac{x^n}{n!} = \frac{d}{dx} \sum f_n \frac{x^n}{n!} = \frac{d}{dx} F(x)\)。\(\blacksquare\)

\subsection{分子物种与分解定理}\label{ux5206ux5b50ux7269ux79cdux4e0eux5206ux89e3ux5b9aux7406}

\textbf{定义 12}（分子物种）：物种 \(M\) 称为\textbf{分子物种}，若 \(M \neq 0\) 且对任意分解 \(M = F + G\)，必有 \(F = 0\) 或 \(G = 0\)。即 \(M\) 不能写为非平凡物种的和。

\textbf{定理 1}（分子物种的分类）：分子物种与传递群作用一一对应。具体地，对任意有限群 \(H\) 的传递作用（即 \(H\) 传递地作用在某个集合上），存在唯一的分子物种 \(X^H\) 使得 \(X^H[U]\) 是 \(U\) 上的 \(H\)-轨道结构族。反之，任意分子物种同构于某个 \(X^H\)。

\textbf{证明}：给定有限群 \(H\) 传递作用在集合 \(\Omega\) 上，\(|\Omega| = n\)。定义物种 \(M_{H,\Omega}\)：\(M_{H,\Omega}[U]\) 为所有双射 \(\phi: \Omega \to U\) 模去 \(H\)-作用（\(\phi \sim \phi \circ h\)）。\(M_{H,\Omega}[n]\) 上 \(\mathfrak{S}_n\) 通过预复合作用，该作用的轨道同构于陪集 \(\mathfrak{S}_n / \operatorname{Stab}_{\mathfrak{S}_n}(\bar{\phi})\)（\(H\) 的传递性保证）。若 \(M_{H,\Omega} = F + G\)，则 \(F[n]\) 和 \(G[n]\) 为 \(\mathfrak{S}_n\)-不变子集，由传递性必有一个为全空间，另一个为空，故 \(M_{H,\Omega}\) 为分子物种。反之，若 \(M\) 为分子物种，取 \(n\) 使 \(M[n] \neq \varnothing\)，\(M[n]\) 为 \(\mathfrak{S}_n\)-集。由分子性，\(M[n]\) 的 \(\mathfrak{S}_n\)-轨道恰为整个空间，设稳定子群为 \(H\)，则 \(M \cong M_{H,\Omega}\)。\(\blacksquare\)

\textbf{定理 2}（分子分解定理）：任意物种 \(F\) 有唯一的分解

\[F = \bigoplus_{i \in I} M_i\]

为分子物种的和（可能无限）。系数（即每种分子物种出现的重数）由 Burnside 型引理给出。

\textbf{证明}：对每个 \(n \geq 0\)，\(F[n]\) 是有限 \(\mathfrak{S}_n\)-集。由 Burnside 环理论，\(\mathfrak{S}_n\) 在 \(F[n]\) 上的作用可分解为传递作用的直和：\(F[n] = \bigoplus_{i \in I_n} \mathcal{O}_{n,i}\)，其中每个 \(\mathcal{O}_{n,i}\) 是 \(\mathfrak{S}_n\)-传递作用轨道。对每个轨道 \(\mathcal{O}_{n,i}\)，定义分子物种 \(M_{n,i}\) 为：\(M_{n,i}[m] = \varnothing\)（\(m \neq n\)），\(M_{n,i}[n] = \mathcal{O}_{n,i}\)，且 \(\mathfrak{S}_m\)-作用由 \(\mathfrak{S}_n\)-作用自然诱导。取 \(I = \bigcup_{n \geq 0} I_n\)，则 \(F = \bigoplus_{i \in I} M_i\)。唯一性：若 \(F = \bigoplus M_i = \bigoplus M'_j\) 为两个分子分解，考虑每个 \(n\) 处的 \(\mathfrak{S}_n\)-集分解，轨道分解的唯一性（由 \(\mathfrak{S}_n\)-集的不可约分解）保证 \(\{M_i\} = \{M'_j\}\)。重数由 Burnside 引理给出：\(N_{H,\Omega} = \frac{1}{|H|} \sum_{h \in H} \operatorname{fix}(F[h])\)。\(\blacksquare\)

\subsection{虚物种环}\label{ux865aux7269ux79cdux73af}

\textbf{定义 13}（虚物种）：物种的\textbf{虚物种}是形式差 \(F - G\)（\(F, G\) 为物种），类比于整数从自然数构造。虚物种全体构成交换环 \(\mathbf{Sp}\)，其中加法为逐点差，乘法由物种积诱导。

环 \(\mathbf{Sp}\) 的结构： - 加法单位元为 \(0\) 物种； - 乘法单位元为 \(1\) 物种； - 分子物种构成 \(\mathbf{Sp}\) 的加法基（由分子分解定理）； - \(\mathbf{Sp}\) 同构于对称群 Burnside 环的完备化。

\textbf{定理 3}（\(\mathbf{Sp}\) 的泛性质）：任意从物种范畴到交换环的函子（与和、积兼容）可通过虚物种环唯一分解。

\textbf{证明}：设 \(\Phi: \mathbf{Sp} \to R\) 为从物种范畴到交换环 \(R\) 的函子，满足 \(\Phi(F+G) = \Phi(F) + \Phi(G)\) 且 \(\Phi(F \cdot G) = \Phi(F) \cdot \Phi(G)\)。将 \(\Phi\) 延拓到虚物种环：定义 \(\bar{\Phi}(F - G) = \Phi(F) - \Phi(G)\)。需验证良定性：若 \(F - G = F' - G'\)（即 \(F + G' = F' + G\)），则 \(\Phi(F) + \Phi(G') = \Phi(F+G') = \Phi(F'+G) = \Phi(F') + \Phi(G)\)，故 \(\Phi(F) - \Phi(G) = \Phi(F') - \Phi(G')\)。延拓后的 \(\bar{\Phi}\) 显然是环同态。唯一性：由分子分解定理，\(\mathbf{Sp}\) 由分子物种加法生成，\(\bar{\Phi}\) 在分子物种上的值由 \(\Phi\) 唯一确定。\(\blacksquare\)

\subsection{循环指标级数与 Pólya 计数}\label{ux5faaux73afux6307ux6807ux7ea7ux6570ux4e0e-puxf3lya-ux8ba1ux6570}

\textbf{定义 14}（循环指标级数）：物种 \(F\) 的\textbf{循环指标级数} \(Z_F\) 是形式幂级数

\[Z_F(p_1, p_2, \ldots) = \sum_{n \geq 0} \frac{1}{n!} \sum_{\sigma \in \mathfrak{S}_n} \operatorname{fix}(F[\sigma]) \prod_{k \geq 1} p_k^{c_k(\sigma)},\]

其中 \(\operatorname{fix}(F[\sigma]) = |\{s \in F[n] : F[\sigma](s) = s\}|\)（\(F[\sigma]\) 的不动点数），\(c_k(\sigma)\) 是 \(\sigma\) 中长度为 \(k\) 的循环数。

基本物种的循环指标级数： - \(Z_X = p_1\)， - \(Z_E = \exp\!\left(\sum_{k \geq 1} \frac{p_k}{k}\right)\)， - \(Z_C = \sum_{k \geq 1} \frac{\varphi(k)}{k} \log \frac{1}{1 - p_k}\)。

\textbf{定理 4}（Pólya 计数定理作为推论）：若取 \(p_k = x^k\)，则 \(Z_F(x, x^2, \ldots) = \widetilde{F}(x)\) 是无标记结构的普通生成函数。若取 \(p_1 = x\) 且 \(p_k = 0\)（\(k \geq 2\)），则 \(Z_F(x, 0, \ldots) = F(x)\) 是指数生成函数。

\textbf{证明}：由循环指标级数定义 \(Z_F = \sum_{n \geq 0} \frac{1}{n!} \sum_{\sigma \in \mathfrak{S}_n} \operatorname{fix}(F[\sigma]) \prod_{k} p_k^{c_k(\sigma)}\)。代入 \(p_k = x^k\)，则 \(\prod_k p_k^{c_k(\sigma)} = \prod_k x^{k c_k(\sigma)} = x^{\sum k c_k(\sigma)} = x^n\)。由 Burnside 引理，\(\frac{1}{n!} \sum_{\sigma \in \mathfrak{S}_n} \operatorname{fix}(F[\sigma])\) 恰为 \(F[n]\) 在 \(\mathfrak{S}_n\) 作用下的轨道数，即无标记结构数 \(\widetilde{f}_n\)，故 \(Z_F(x, x^2, \ldots) = \sum \widetilde{f}_n x^n = \widetilde{F}(x)\)。代入 \(p_1 = x\)，\(p_k = 0\)（\(k \geq 2\)），则仅 \(\sigma = \operatorname{id}\)（\(c_1 = n\)，\(c_k = 0\)）贡献非零项，得 \(Z_F(x, 0, \ldots) = \sum_{n \geq 0} \frac{1}{n!} \operatorname{fix}(F[\operatorname{id}]) x^n = \sum_{n \geq 0} f_n \frac{x^n}{n!} = F(x)\)。\(\blacksquare\)

因此 Pólya 计数定理等价于：在循环指标级数中进行变量特殊化可得标记与无标记两种计数模式，无需单独计算 Burnside 引理。

\textbf{推论 2}（基本物种的循环指标级数运算律）： - \(Z_{F+G} = Z_F + Z_G\)， - \(Z_{F \cdot G} = Z_F \cdot Z_G\)， - \(Z_{F \circ G} = Z_F \circ Z_G\)（Plethystic 复合）， - \(Z_{F'} = \frac{\partial}{\partial p_1} Z_F\)。

\newpage

\chapter{04-离散数学/01-组合数学/05-代数组合学.md}\label{ux79bbux6563ux6570ux5b6601-ux7ec4ux5408ux6570ux5b6605-ux4ee3ux6570ux7ec4ux5408ux5b66.md}

\section{第88章：Shannon 熵与信息度量}\label{ux7b2c88ux7ae0shannon-ux71b5ux4e0eux4fe1ux606fux5ea6ux91cf}

熵是信息论的核心概念，由 Claude Shannon 于 1948 年在《A Mathematical Theory of Communication》中引入，度量了随机变量固有的不确定性。这一概念的历史根源可追溯至 Boltzmann 和 Gibbs 在统计力学中的熵公式 \(S = -k_B \sum p_i \ln p_i\)，但 Shannon 的突破在于将其从物理量重新诠释为信息量------不确定性越大，观察结果获得的信息越多。Shannon 熵在数学体系中占据独特位置：它与组合学（二项式系数的渐近估计）、概率论（大偏差理论）和数理统计（最大熵原理）有深层联系，更是信息论三大编码定理（无失真信源编码、信道编码、率失真）的出发点。本章从 Khinchin 公理出发严格定义熵，建立联合熵、条件熵和互信息的链式法则与信息不等式（数据处理不等式），引入相对熵（KL 散度）和交叉熵作为统计推断的理论基础，并以微分熵将理论推广至连续随机变量。本章为第 43 章的信道容量理论和第 45 章的信息论深化提供概念基础。

\subsection{142.1 熵的定义与公理化}\label{ux71b5ux7684ux5b9aux4e49ux4e0eux516cux7406ux5316}

\textbf{定义 142.1}（离散随机变量的熵）：设 \(X\) 是离散随机变量，取值于 \(\mathcal{X}\)，概率分布为 \(p(x) = \mathbb{P}(X = x)\)。\(X\) 的\textbf{熵}（Shannon 熵）定义为

\[H(X) = -\sum_{x \in \mathcal{X}} p(x) \log_2 p(x)\]

（约定 \(0 \log_2 0 = 0\)）。熵的单位是比特（bit）。若使用自然对数，单位为奈特（nat）。

\textbf{定理 142.1}（熵的公理化 / Khinchin 公理）：设 \(H(p_1, \ldots, p_n)\) 是定义在概率分布上的函数。若 \(H\) 满足： 1. \(H\) 是 \(p_i\) 的连续函数 2. 对均匀分布，\(H(1/n, \ldots, 1/n)\) 关于 \(n\) 单调递增 3. 分组公理：\(H(p_1, \ldots, p_n) = H(p_1+p_2, p_3, \ldots, p_n) + (p_1+p_2) H\left(\frac{p_1}{p_1+p_2}, \frac{p_2}{p_1+p_2}\right)\)

则 \(H(p_1, \ldots, p_n) = -c \sum p_i \log p_i\)（\(c > 0\)）。

\textbf{证明}：设 \(F(n) = H(1/n, \ldots, 1/n)\)。由分组公理，\(F(mn) = H(1/(mn), \ldots) = F(m) + F(n)\)（将 \(mn\) 个等概率事件分为 \(m\) 组每组 \(n\) 个）。由条件 (1) 连续性，此函数方程的唯一解为 \(F(n) = c \log n\)（\(c > 0\) 由条件 (2) 单调性保证）。对任意有理概率 \(p_i = a_i/N\)（\(a_i \in \mathbb{N}\)，\(\sum a_i = N\)），反复使用分组公理：\(H(p_1, \ldots, p_n) = F(N) - \sum p_i F(a_i) = c \log N - c \sum p_i \log a_i = -c \sum p_i \log(a_i/N) = -c \sum p_i \log p_i\)。由连续性，此公式延拓至所有概率分布。\(\blacksquare\)

\textbf{命题 142.2}（熵的基本性质）： - \(0 \leq H(X) \leq \log_2 |\mathcal{X}|\)（非负性，最大熵原理） - \(H(X) = 0\) 当且仅当 \(X\) 是确定性的（退化分布） - \(H(X) = \log_2 |\mathcal{X}|\) 当且仅当 \(X\) 是均匀分布

\textbf{证明}：非负性：\(0 \leq p(x) \leq 1\) 故 \(-\log p(x) \geq 0\)，\(H(X) = \sum p(x)(-\log p(x)) \geq 0\)。\(H(X) = 0\) 当且仅当对每个 \(x\)，\(p(x) \log p(x) = 0\)，即 \(p(x) \in \{0,1\}\)，故 \(X\) 退化。上界：由 Jensen 不等式，\(H(X) = \sum p(x) \log(1/p(x)) \leq \log(\sum p(x) \cdot 1/p(x)) = \log |\mathcal{X}|\)。等号成立当且仅当 \(1/p(x)\) 为常数，即 \(p(x) = 1/|\mathcal{X}|\)（均匀分布）。\(\blacksquare\)

\subsection{142.2 联合熵、条件熵与互信息}\label{ux8054ux5408ux71b5ux6761ux4ef6ux71b5ux4e0eux4e92ux4fe1ux606f}

熵 \(H(X)\) 刻画了单个随机变量的不确定性，但信息论的核心问题------通信与压缩------本质上是关于两个或多个随机变量之间的关系：发送方与接收方之间的信息流通量。联合熵将这种多个变量间的总不确定性量化，而条件熵则衡量在已知部分信息后剩余的未知程度。二者的差值------互信息------简洁地捕捉了变量间的共享信息量，奠定了 Shannon 全部编码理论的基础。

\textbf{定义 142.2}（联合熵）：\((X, Y)\) 的联合熵为

\[H(X, Y) = -\sum_{x, y} p(x, y) \log_2 p(x, y)\]

\textbf{定义 142.3}（条件熵）：给定 \(Y\) 下 \(X\) 的条件熵为

\[H(X|Y) = \sum_{y} p(y) H(X|Y=y) = -\sum_{x,y} p(x,y) \log_2 p(x|y)\]

\textbf{定理 142.3}（链式法则）：\(H(X, Y) = H(X) + H(Y|X) = H(Y) + H(X|Y)\)。

\textbf{证明}：\(H(X,Y) = -\sum_{x,y} p(x,y) \log p(x,y) = -\sum_{x,y} p(x,y) \log(p(x)p(y|x)) = -\sum_{x,y} p(x,y) \log p(x) - \sum_{x,y} p(x,y) \log p(y|x) = -\sum_x p(x) \log p(x) - \sum_{x,y} p(x,y) \log p(y|x) = H(X) + H(Y|X)\)。由对称性，\(H(X,Y) = H(Y) + H(X|Y)\)。\(\blacksquare\)

\textbf{定义 142.4}（互信息）：\(X\) 和 \(Y\) 之间的互信息为

\[I(X; Y) = H(X) - H(X|Y) = H(Y) - H(Y|X) = H(X) + H(Y) - H(X, Y)\]

互信息对称且非负：\(I(X; Y) \geq 0\)，\(I(X; Y) = 0\) 当且仅当 \(X\) 和 \(Y\) 独立。

\textbf{命题 142.4}（信息不等式）：\(H(X|Y) \leq H(X)\)，等号成立当且仅当 \(X\) 和 \(Y\) 独立。即「条件不会增加熵」。

\textbf{证明}：\(H(X|Y) - H(X) = -\sum_{x,y} p(x,y) \log p(x|y) + \sum_x p(x) \log p(x) = -\sum_{x,y} p(x,y) \log \frac{p(x|y)}{p(x)} = -D_{KL}(p(x|y) \| p(x)) \leq -I(X;Y) \leq 0\)。更直接地，\(H(X) - H(X|Y) = I(X;Y) \geq 0\)，等号成立当且仅当 \(I(X;Y) = 0\)，即 \(X\) 和 \(Y\) 独立。\(\blacksquare\)

\textbf{定理 142.5}（数据处理不等式）：若 \(X \to Y \to Z\) 构成 Markov 链（即 \(Z\) 条件独立于 \(X\) 给定 \(Y\)），则 \(I(X; Z) \leq I(X; Y)\)。数据处理不能增加信息。

\textbf{证明}：由链式法则，\(I(X; Y, Z) = I(X; Z) + I(X; Y|Z) = I(X; Y) + I(X; Z|Y)\)。由 Markov 性 \(X \to Y \to Z\)，\(X\) 和 \(Z\) 在给定 \(Y\) 时条件独立，故 \(I(X; Z|Y) = 0\)。因此 \(I(X; Y) = I(X; Z) + I(X; Y|Z) \geq I(X; Z)\)（因 \(I(X; Y|Z) \geq 0\)）。等号成立当且仅当 \(I(X; Y|Z) = 0\)，即 \(X \to Z \to Y\) 也构成 Markov 链。\(\blacksquare\)

\subsection{142.3 相对熵与交叉熵}\label{ux76f8ux5bf9ux71b5ux4e0eux4ea4ux53c9ux71b5}

\textbf{定义 142.5}（相对熵 / KL 散度）：分布 \(p\) 相对于 \(q\) 的相对熵（Kullback-Leibler 散度）为

\[D_{KL}(p \| q) = \sum_x p(x) \log_2 \frac{p(x)}{q(x)}\]

（约定 \(0 \log(0/0) = 0\)，\(p \log(p/0) = \infty\)）。

\textbf{定理 142.6}（Gibbs 不等式）：\(D_{KL}(p \| q) \geq 0\)，等号成立当且仅当 \(p = q\)。

\emph{证明}：由 \(\log t \leq t-1\) 的对数不等式，\(\sum p \log(q/p) \leq \sum p(q/p - 1) = \sum q - \sum p = 0\)。∎

\textbf{定义 142.6}（交叉熵）：\(H(p, q) = H(p) + D_{KL}(p \| q) = -\sum p(x) \log q(x)\)。交叉熵在机器学习中广泛用作损失函数。

\textbf{推论 142.7}：\(I(X; Y) = D_{KL}(p(x,y) \| p(x)p(y))\)。互信息是联合分布与独立分布之间的 KL 散度。

\textbf{证明}：\(D_{KL}(p(x,y) \| p(x)p(y)) = \sum_{x,y} p(x,y) \log \frac{p(x,y)}{p(x)p(y)} = \sum_{x,y} p(x,y) \log \frac{p(x|y)}{p(x)} = \sum_{x,y} p(x,y) \log p(x|y) - \sum_x p(x) \log p(x) = -H(X|Y) + H(X) = I(X;Y)\)。\(\blacksquare\)

\subsection{142.4 微分熵与连续随机变量}\label{ux5faeux5206ux71b5ux4e0eux8fdeux7eedux968fux673aux53d8ux91cf}

\textbf{定义 142.7}（微分熵）：连续随机变量 \(X\)（密度 \(f(x)\)）的\textbf{微分熵}为

\[h(X) = -\int f(x) \log f(x) \, dx\]

微分熵可以取负值，且不是坐标变换下的不变量（与离散熵不同）。

\textbf{定理 142.8}（最大微分熵的分布）： - 固定方差 \(\sigma^2\) 下，正态分布 \(\mathcal{N}(\mu, \sigma^2)\) 最大化微分熵，\(h = \frac{1}{2}\log_2(2\pi e \sigma^2)\) - 固定支撑 \([a, b]\) 下，均匀分布最大化微分熵 - 固定均值 \(\lambda\) 下（\(x \geq 0\)），指数分布最大化微分熵

\textbf{证明}：用变分法。设 \(f\) 为密度函数，在约束 \(\int f = 1\) 及 \(\int x^2 f = \sigma^2\) 下最大化 \(h(f) = -\int f \log f\)。引入 Lagrange 乘子，\(\delta \int f(\log f + \alpha + \beta x^2) dx = 0\)，得 \(\log f + 1 + \alpha + \beta x^2 = 0\)，即 \(f(x) \propto e^{-\beta x^2}\)。由约束 \(\beta = 1/(2\sigma^2)\)，故 \(f\) 为正态分布 \(\mathcal{N}(0,\sigma^2)\)。代入得 \(h = \frac{1}{2}\log_2(2\pi e \sigma^2)\)。同理，支撑 \([a,b]\) 约束下 \(\log f\) 为常数，得均匀分布；均值约束下 \(f(x) \propto e^{-\lambda x}\)，得指数分布。\(\blacksquare\)

\textbf{定理 142.9}（熵幂不等式）：\(e^{2h(X+Y)/n} \geq e^{2h(X)/n} + e^{2h(Y)/n}\)（向量形式）。

\textbf{证明}（\(n=1\) 情形）：设 \(X, Y\) 独立，密度为 \(f, g\)。由熵幂不等式等价于 \(h(X+Y) \geq h(\tilde{X}+\tilde{Y})\)，其中 \(\tilde{X}, \tilde{Y}\) 为方差相同的独立 Gauss 随机变量。利用 Young 不等式：对 \(p, q > 1\)，\(1/p+1/q = 1\)，\(\|f * g\|_r \leq \|f\|_p \|g\|_q\)（\(1/r = 1/p+1/q-1\)）。令 \(p, q \to 1\) 且 \(r \to 1\)，对 \(\|f * g\|_{1+\varepsilon}\) 展开，\(\varepsilon \to 0\) 时级数的主项给出熵的表达式。更精确地，由 Stam 不等式和 Fisher 信息的 Blachman-Stam 不等式：当 \(X, Y\) 独立时，\(1/J(X+Y) \geq 1/J(X) + 1/J(Y)\)（\(J\) 为 Fisher 信息），结合 de Bruijn 恒等式 \(\frac{d}{dt}h(X + \sqrt{t}Z) = \frac{1}{2}J(X + \sqrt{t}Z)\)，积分即得熵幂不等式。\(\blacksquare\)

\hrulefill

\textbf{q-级数的进一步理论与应用}

q-级数理论在 20 世纪经历了深刻的变革，其核心推动力来自表示论和数学物理中的新发现。\textbf{Bailey 引理}（1947）是 q-级数理论中最强大的工具之一，它提供了一种系统的方法从一对已知的 q-级数恒等式生成无穷多个新的恒等式。具体地，若 \(\{\alpha_n\}\) 和 \(\{\beta_n\}\) 满足 \(\beta_n = \sum_{r=0}^n \frac{\alpha_r}{(q;q)_{n-r}(aq;q)_{n+r}}\)，则 Bailey 引理断言 \(\alpha_n'\) 和 \(\beta_n'\) 构成另一对 Bailey 对，其中参数经过适当的变换。Andrews 利用 Bailey 链方法（1970s-1980s）将 Rogers-Ramanujan 恒等式推广为 \textbf{Andrews-Gordon 恒等式}：对任意奇数 \(k \geq 1\)，\(\sum_{n_1 \geq \cdots \geq n_{k-1} \geq 0} \frac{q^{n_1^2 + \cdots + n_{k-1}^2 + n_i + \cdots + n_{k-1}}}{(q;q)_{n_1-n_2} \cdots (q;q)_{n_{k-2}-n_{k-1}}(q;q)_{n_{k-1}}} = \frac{1}{(q;q)_\infty} \sum_{m=-\infty}^\infty (-1)^m q^{((2k+1)m^2 + (2i+1)m)/2}\)，其中 \(i = 0, 1, \ldots, k-1\)。这一恒等式在共形场论中揭示了 \(\widehat{\mathfrak{sl}_2}\) 的仿射 Virasoro 特征公式与划分恒等式之间的深刻联系。\textbf{q-超几何级数}的现代理论还包含 Watson 变换公式、Jackson 求和公式（q-二项式定理的强形式）以及 q-Selberg 积分，后者是 Macdonald 多项式理论和随机矩阵理论的基石。q-级数理论在数学物理中通过 \textbf{Bethe 拟设}（Bethe ansatz）与量子可积系统建立了联系------XXZ 自旋链的配分函数可通过 q-超几何级数精确计算。此外，q-级数理论在数论中通过 \textbf{模形式}（modular forms）与 Ramanujan 的 \(\tau\) 函数、mock theta 函数和同余理论相关联，Zwegers（2002）的工作将 mock theta 函数纳入调和 Maass 形式的框架，揭示了 q-级数与尖点形式之间的深刻对偶性。

\hrulefill

\section{第89章：信道容量与编码定理}\label{ux7b2c89ux7ae0ux4fe1ux9053ux5bb9ux91cfux4e0eux7f16ux7801ux5b9aux7406}

Shannon 的信道编码定理（第二定理，1948 年）被誉为信息论的''大宪章''------它断言任何速率低于信道容量 \(C\) 的通信都可以通过适当的编码方案以任意小的错误概率实现，反之任何速率超过 \(C\) 的通信必然不可靠。这一结果的意义远超通信工程本身：它是第一个严格证明存在性而非构造性的编码定理，其证明中引入的随机编码方法（random coding）开创了概率方法在组合学和编码理论中的先河。信道容量 \(C = \max_{p(x)} I(X; Y)\) 简洁地将物理信道的极限性能表达为互信息的最大化问题，从而将信道设计与信息度量无缝耦合。本章在信源编码定理（典型序列与渐近均分性）的基础上系统阐述信道编码定理：从离散无记忆信道的数学模型出发，定义信道容量并给出经典信道（BSC、BEC、AWGN）的容量公式，以联合典型译码和随机编码方法严格证明 Shannon 第二定理。

\subsection{143.1 信源编码定理}\label{ux4fe1ux6e90ux7f16ux7801ux5b9aux7406}

\textbf{定义 143.1}（信源编码）：将信源符号序列映射为二进制码字。\textbf{平均码长} \(\bar{L} = \sum p(x) \ell(x)\)，其中 \(\ell(x)\) 是码字 \(x\) 的码长。

信源编码面临的基本问题是：平均码长能够被压缩到多短？直观上，熵 \(H(X)\) 是在最优编码下限长的理论上界------任何唯一可译码的平均码长都不可能低于熵。Shannon 第一定理（无噪声编码定理）以精确的方式确认了这一直觉：存在前缀码（如 Huffman 编码）使平均码长介于 \(H(X)\) 和 \(H(X)+1\) 之间，并且通过分块编码可以渐近达到熵。

\textbf{定理 143.1}（无噪声编码定理 / Shannon 第一定理）：对信源 \(X\)，任何唯一可译码的平均码长满足 \(\bar{L} \geq H(X)\)。存在前缀码（如 Huffman 编码）满足 \(H(X) \leq \bar{L} < H(X) + 1\)。

\textbf{证明}（典型序列 + AEP 构造法）：对信源 \(X\) 的 \(n\) 次独立输出 \(X^n = (X_1, \ldots, X_n)\)，由渐近均分性（AEP），对任意 \(\varepsilon > 0\)，当 \(n\) 充分大时，\(\varepsilon\)-典型集 \(A_\varepsilon^{(n)} = \{x^n : |-\frac{1}{n}\log p(x^n) - H(X)| < \varepsilon\}\) 满足 \(P(A_\varepsilon^{(n)}) > 1 - \varepsilon\) 且 \(|A_\varepsilon^{(n)}| \leq 2^{n(H(X)+\varepsilon)}\)。编码方案：为 \(A_\varepsilon^{(n)}\) 中每个序列分配一个长为 \(\lceil n(H(X)+\varepsilon) \rceil\) 的二进制码字，非典型序列全部映射到同一个特殊码字。平均码长 \(\bar{L} \leq n(H(X)+\varepsilon) + \varepsilon \cdot \ell_{\max}\)（非典型序列概率 \(\leq \varepsilon\)）。故 \(\bar{L}/n \leq H(X) + \varepsilon + o(1)\)。下界：对任意唯一可译码，由 Kraft 不等式 \(\sum 2^{-\ell_i} \leq 1\)，令 \(Q_i = 2^{-\ell_i} / \sum 2^{-\ell_j}\)，则 \(\bar{L} - H(X) = \sum p_i \log(p_i / Q_i) - \log(\sum 2^{-\ell_j}) \geq 0\)（Gibbs 不等式），故 \(\bar{L} \geq H(X)\)。令 \(n \to \infty\) 得 \(\bar{L}/n \to H(X)\)。\(\blacksquare\)

\textbf{定理 143.2}（Kraft 不等式）：长度为 \(\ell_1, \ldots, \ell_n\) 的前缀码存在当且仅当 \(\sum_{i=1}^n 2^{-\ell_i} \leq 1\)。

\textbf{证明}：（必要性）设存在前缀码，将码字对应到完全二叉树中深度为 \(\ell_i\) 的叶子。每个叶子节点阻挡了 \(2^{L-\ell_i}\) 个深度 \(L\) 的节点（\(L = \max \ell_i\)）。由于前缀码无公共前缀，阻挡的节点互不相交，故 \(\sum_i 2^{L-\ell_i} \leq 2^L\)，即 \(\sum_i 2^{-\ell_i} \leq 1\)。（充分性）设 \(\sum 2^{-\ell_i} \leq 1\)，将 \(\ell_i\) 排序为 \(\ell_1 \leq \cdots \leq \ell_n\)。在完全二叉树上，按深度优先顺序分配节点：对每个 \(\ell_i\)，选择深度为 \(\ell_i\) 且祖先未被占用的最左节点。由于 \(\sum 2^{-\ell_i} \leq 1\)，总有足够空间分配所有码字。\(\blacksquare\)

\textbf{算法 143.1}（Huffman 编码，1952）： 1. 将概率最小的两个符号合并为一个节点，其概率为两者之和 2. 重复直到只剩一个节点 3. 回溯分配码字（左 0 右 1）

Huffman 编码是最优前缀码（最小化平均码长）。

\textbf{定理 143.3}（信源编码定理，块编码版本）：对信源 \(X\) 的独立同分布序列 \(X^n\)，存在编码使得平均每符号码长 \(\frac{1}{n}\bar{L}_n \to H(X)\)（\(n \to \infty\)）。

\textbf{证明}：由 AEP，对任意 \(\varepsilon > 0\)，当 \(n\) 充分大时，典型集 \(A_\varepsilon^{(n)}\) 满足 \(|A_\varepsilon^{(n)}| \leq 2^{n(H(X)+\varepsilon)}\) 且 \(P(A_\varepsilon^{(n)}) > 1 - \varepsilon\)。为 \(A_\varepsilon^{(n)}\) 中每个序列分配一个长为 \(\lceil n(H(X)+\varepsilon) \rceil\) 的二进制码字，非典型序列全部映射到单个长为 \(\lceil n(H(X)+\varepsilon) \rceil\) 的码字。平均码长 \(\bar{L}_n \leq n(H(X)+\varepsilon) + 1\)（典型序列编码长度）\(+ \varepsilon \cdot (n(H(X)+\varepsilon) + 1)\)（非典型序列贡献）。故 \(\frac{1}{n}\bar{L}_n \leq (1+\varepsilon)(H(X)+\varepsilon) + 1/n\)。令 \(\varepsilon \to 0\) 且 \(n \to \infty\) 得 \(\limsup \frac{1}{n}\bar{L}_n \leq H(X)\)。下界由定理 143.1 已得 \(\frac{1}{n}\bar{L}_n \geq H(X)\)。故 \(\frac{1}{n}\bar{L}_n \to H(X)\)。\(\blacksquare\)

\subsection{143.2 信道容量}\label{ux4fe1ux9053ux5bb9ux91cf}

\textbf{定义 143.2}（离散无记忆信道 / DMC）：信道由输入字母表 \(\mathcal{X}\)、输出字母表 \(\mathcal{Y}\) 和转移概率 \(p(y|x)\) 描述。\(n\) 次使用信道时，\(p(y^n|x^n) = \prod_{i=1}^n p(y_i|x_i)\)。

\textbf{定义 143.3}（信道容量）：离散无记忆信道的容量为

\[C = \max_{p(x)} I(X; Y)\]

即输入分布 \(p(x)\) 可任意选择时，互信息 \(I(X; Y)\) 的最大值。

信道容量是信息论中最核心的概念之一，它量化了给定信道在可靠传输前提下所能承载的最大信息速率。对于几个经典信道，其容量有闭式公式。\textbf{二元对称信道}（BSC）以交叉概率 \(p\) 翻转输入比特，容量为 \(C = 1 - H(p)\)，其中 \(H(p) = -p\log_2 p - (1-p)\log_2(1-p)\) 是二元熵。当 \(p = 0\) 或 \(p = 1\) 时信道无噪声，\(C = 1\) 比特/次；当 \(p = 1/2\) 时信道完全随机，\(C = 0\)。\textbf{二元擦除信道}（BEC）以概率 \(\alpha\) 擦除输入比特，容量为 \(C = 1 - \alpha\)。\textbf{加性高斯白噪声信道}（AWGN）是连续信道中最经典的模型：\(Y = X + Z\)，\(Z \sim \mathcal{N}(0, \sigma^2)\)，输入功率约束为 \(E[X^2] \leq P\)。其容量为 \(C = \frac{1}{2}\log_2(1 + P/\sigma^2)\)，著名的 Shannon 公式。信道容量的计算通常涉及凸优化------\(I(X; Y)\) 是输入分布 \(p(x)\) 的凹函数，在约束条件下求最大值，Arimoto-Blahut 算法提供了数值求解的迭代方法。从几何直观上看，输入分布的选择决定了输出分布的空间范围，而容量 \(C\) 是该范围内的最大互信息''泡''。在反馈信道中，Shannon 证明反馈不能增加离散无记忆信道的容量（\(C_{\text{fb}} = C\)），这被称为''反馈无益''定理。

\subsection{143.3 信道编码定理}\label{ux4fe1ux9053ux7f16ux7801ux5b9aux7406}

信源编码定理解决了压缩问题------将信源输出的冗余最小化。而信道编码定理面对的是通信问题------在噪声信道中实现可靠传输。Shannon 第二定理断言，只要编码速率不超过信道容量 \(C = \max_{p(x)} I(X; Y)\)，就存在编码方案使错误概率随着码长增加而指数衰减至零。定理的证明堪称概率方法的典范：不是在数学上构造具体的好码（至今仍是开放问题），而是通过随机编码论证了好码的存在性。

\textbf{定理 143.4}（Shannon 信道编码定理 / Shannon 第二定理，1948）：对离散无记忆信道，任何速率 \(R < C\) 是可达的：存在码率为 \(R\) 的编码方案，使得错误概率 \(P_e \to 0\)（当码长 \(n \to \infty\)）。反之，若 \(R > C\)，则不可能可靠通信。

\textbf{证明}（随机编码 + 联合 AEP）：固定输入分布 \(p(x)\) 和速率 \(R < I(X;Y)\)。随机生成码本 \(\mathcal{C} = \{\mathbf{x}^{(1)}, \ldots, \mathbf{x}^{(2^{nR})}\}\)，每个码字的每个符号独立地按 \(p(x)\) 采样。译码器采用联合典型译码：收到 \(\mathbf{y}\) 后，寻找唯一的 \(m\) 使 \((\mathbf{x}^{(m)}, \mathbf{y}) \in A_\varepsilon^{(n)}\)（联合典型集）；若不存在或存在多个，则报错。分析平均错误概率 \(\bar{P}_e\)（对随机码本取期望）。设发送码字为 \(\mathbf{x}^{(1)}\)。（1）\((\mathbf{x}^{(1)}, \mathbf{y}) \notin A_\varepsilon^{(n)}\)：由联合 AEP，概率 \(\leq \varepsilon\)（\(n\) 充分大）。（2）对某个 \(m \neq 1\)，\((\mathbf{x}^{(m)}, \mathbf{y}) \in A_\varepsilon^{(n)}\)：由于 \(\mathbf{x}^{(m)}\) 独立于 \(\mathbf{y}\)，\(P((\mathbf{x}^{(m)}, \mathbf{y}) \in A_\varepsilon^{(n)}) \leq 2^{-n(I(X;Y)-3\varepsilon)}\)。由联合界，\(\bar{P}_e \leq \varepsilon + 2^{nR} \cdot 2^{-n(I(X;Y)-3\varepsilon)} = \varepsilon + 2^{-n(I(X;Y)-R-3\varepsilon)}\)。当 \(R < I(X;Y)\) 时，取 \(\varepsilon\) 充分小使 \(I(X;Y) - R - 3\varepsilon > 0\)，则 \(\bar{P}_e \to 0\)（\(n \to \infty\)）。因此存在具体的确定性码本使错误概率 \(\to 0\)。取 \(p(x)\) 为达到信道容量 \(C\) 的输入分布，即得 \(R < C\) 可达。\(\blacksquare\)

\textbf{定义 143.4}（联合典型集）：长度为 \(n\) 的序列对 \((x^n, y^n)\) 是\textbf{联合典型的}，如果其经验分布接近 \(p(x, y)\)。具体地，\(A_\varepsilon^{(n)} = \{(x^n, y^n) : |-\frac{1}{n}\log p(x^n) - H(X)| < \varepsilon, |-\frac{1}{n}\log p(y^n) - H(Y)| < \varepsilon, |-\frac{1}{n}\log p(x^n, y^n) - H(X,Y)| < \varepsilon\}\)。

\textbf{定理 143.5}（联合 AEP / 渐近均分性）：对充分大的 \(n\)，联合典型集的大小约为 \(2^{nH(X,Y)}\)，且概率接近 1。

\textbf{证明}：设 \((X_i, Y_i)\) 独立同分布于 \(p(x,y)\)。由大数定律，\(-\frac{1}{n}\log p(X^n) \to H(X)\)，\(-\frac{1}{n}\log p(Y^n) \to H(Y)\)，\(-\frac{1}{n}\log p(X^n, Y^n) \to H(X,Y)\) 均依概率收敛。因此对任意 \(\varepsilon > 0\)，\(P(A_\varepsilon^{(n)}) \to 1\)。关于大小：\(1 \geq \sum_{(x^n,y^n) \in A_\varepsilon^{(n)}} p(x^n,y^n) \geq |A_\varepsilon^{(n)}| \cdot 2^{-n(H(X,Y)+\varepsilon)}\)，故 \(|A_\varepsilon^{(n)}| \leq 2^{n(H(X,Y)+\varepsilon)}\)。又 \(P(A_\varepsilon^{(n)}) \to 1\) 意味着 \(n\) 充分大时 \(P(A_\varepsilon^{(n)}) \geq 1-\varepsilon\)，则 \(1-\varepsilon \leq \sum_{A_\varepsilon^{(n)}} p(x^n,y^n) \leq |A_\varepsilon^{(n)}| \cdot 2^{-n(H(X,Y)-\varepsilon)}\)，故 \(|A_\varepsilon^{(n)}| \geq (1-\varepsilon) 2^{n(H(X,Y)-\varepsilon)}\)。因此 \(|A_\varepsilon^{(n)}| \approx 2^{nH(X,Y)}\)。\(\blacksquare\)

\textbf{定理 143.6}（信道编码逆定理）：若 \(R > C\)，则任何译码方案的错误概率有正的下界，\(P_e^{(n)} \geq 1 - C/R - 1/(nR)\)。

\textbf{证明}：设消息 \(W\) 均匀分布于 \(\{1, \ldots, 2^{nR}\}\)，\(X^n(W)\) 为发送码字，\(Y^n\) 为接收序列。由 Fano 不等式，\(H(W|Y^n) \leq 1 + P_e^{(n)} nR\)。故 \(nR = H(W) = I(W; Y^n) + H(W|Y^n) \leq I(X^n; Y^n) + 1 + P_e^{(n)} nR\)（由数据处理不等式 \(I(W; Y^n) \leq I(X^n; Y^n)\)）。又 \(I(X^n; Y^n) \leq \sum_{i=1}^n I(X_i; Y_i) \leq nC\)（由信道容量定义和链式法则）。代入得 \(nR \leq nC + 1 + P_e^{(n)} nR\)，整理得 \(P_e^{(n)} \geq 1 - C/R - 1/(nR)\)。当 \(R > C\) 时，\(n \to \infty\) 给出 \(\liminf P_e^{(n)} \geq 1 - C/R > 0\)。\(\blacksquare\)

\subsection{143.4 率失真理论}\label{ux7387ux5931ux771fux7406ux8bba}

信道编码定理关注的是在给定速率下可以实现多低的错误概率，而率失真理论从对偶角度出发：给定可容忍的失真水平 \(D\)，压缩信源所需的最小速率是多少？率失真函数 \(R(D) = \min_{p(\hat{x}|x): \mathbb{E}[d] \leq D} I(X; \hat{X})\) 将这一问题表述为互信息在失真约束下的极小化，从而将信源编码的信息论极限与最优量化问题紧密耦合。Gauss 信源在均方失真下的闭式解 \(R(D) = \frac{1}{2}\log_2(\sigma^2/D)\) 是率失真理论最著名的应用。

\textbf{定义 143.5}（率失真函数）：在允许失真 \(D\) 下，信源 \(X\) 的率失真函数为

\[R(D) = \min_{p(\hat{x}|x) : \mathbb{E}[d(X, \hat{X})] \leq D} I(X; \hat{X})\]

\textbf{定理 143.7}（率失真定理）：\(R(D)\) 是在失真 \(D\) 下信源编码可实现的最小码率。对 Gauss 信源（方差 \(\sigma^2\)），\(R(D) = \frac{1}{2}\log_2(\sigma^2/D)\)（\(0 \leq D \leq \sigma^2\)）。

\textbf{证明}（正定理与逆定理）：（可达性）固定条件分布 \(p(\hat{x}|x)\) 满足 \(\mathbb{E}[d(X,\hat{X})] \leq D\)。由 \(p(\hat{x}) = \sum_x p(x)p(\hat{x}|x)\) 独立采样 \(2^{nR}\) 个重构码字。对信源输出 \(x^n\)，在码本中寻找满足 \((x^n, \hat{x}^n) \in A_\varepsilon^{(n)}\)（联合典型集）的码字。由联合 AEP，\(P((x^n, \hat{x}^n) \in A_\varepsilon^{(n)}) \to 1\)。错误码字与 \(x^n\) 联合典型的概率 \(\leq 2^{-n(I(X;\hat{X})-3\varepsilon)}\)。由联合界，失败概率 \(\leq 2^{nR} \cdot 2^{-n(I(X;\hat{X})-3\varepsilon)}\)，当 \(R > I(X;\hat{X})\) 时 \(\to 0\)。取所有可行 \(p(\hat{x}|x)\) 下互信息的下确界，得 \(R(D)\) 可达。（逆定理）对任意码率为 \(R\) 且达到失真 \(\leq D\) 的编码方案，\(nR \geq H(\hat{X}^n) \geq I(X^n; \hat{X}^n) = nI(X;\hat{X})\)（由数据处理不等式），故 \(R \geq I(X;\hat{X}) \geq R(D)\)。Gauss 信源下，\(R(D) = \frac{1}{2}\log_2(\sigma^2/D)\) 由二次失真测度下互信息在方差约束下的极小化给出，闭式解由反向水注法得到。\(\blacksquare\)

\hrulefill

\hrulefill

\hrulefill

\hrulefill

\hrulefill

\hrulefill

\textbf{q-级数与q-类似}

q-级数（q-Series）和 q-类似（q-Analogue）是组合数学、数论和表示论交汇处的核心理论。其基本思想是：对经典数学对象（整数、阶乘、二项式系数、特殊函数等）引入参数 \(q\)，当 \(q \to 1\) 时恢复经典对象，但 \(q \neq 1\) 时揭示出丰富的组合结构和量子形变信息。

\subsection{q-整数与基本定义}\label{q-ux6574ux6570ux4e0eux57faux672cux5b9aux4e49}

\textbf{定义 1}（q-整数）：对 \(n \in \mathbb{N}\) 和形式变量 \(q\)，\textbf{q-整数}定义为

\[[n]_q = \frac{1 - q^n}{1 - q} = 1 + q + q^2 + \cdots + q^{n-1}.\]

显然 \(\lim_{q \to 1} [n]_q = n\)。若取 \(q\) 为素数幂，则 \([n]_q\) 恰为 \(\mathbb{F}_q^n\) 中一维子空间的个数。

\textbf{定义 2}（q-阶乘）：\(n\) 的 \textbf{q-阶乘}为

\[[n]_q! = [1]_q [2]_q \cdots [n]_q = \prod_{k=1}^n \frac{1 - q^k}{1 - q}.\]

约定 \([0]_q! = 1\)。

\textbf{定义 3}（q-二项式系数 / Gaussian 二项式系数）：对 \(0 \leq k \leq n\)，

\[\binom{n}{k}_q = \frac{[n]_q!}{[k]_q! [n-k]_q!} = \frac{(1-q^n)(1-q^{n-1})\cdots(1-q^{n-k+1})}{(1-q^k)(1-q^{k-1})\cdots(1-q)}.\]

\textbf{命题 1}（q-二项式系数的基本性质）： - \(\binom{n}{k}_q = \binom{n}{n-k}_q\)（对称性）； - \(\lim_{q \to 1} \binom{n}{k}_q = \binom{n}{k}\)（经典极限）； - \(\binom{n}{k}_q \in \mathbb{Z}[q]\) 是系数为正整数的多项式。

\textbf{证明}：对称性由定义直接验证：\(\binom{n}{n-k}_q = \frac{[n]_q!}{[n-k]_q! [k]_q!} = \binom{n}{k}_q\)。经典极限：\(\lim_{q \to 1} [m]_q = m\)，故 \(\lim_{q \to 1} \binom{n}{k}_q = \frac{n!}{k!(n-k)!} = \binom{n}{k}\)。正整系数：由定理 1（逆序数解释），\(\binom{n}{k}_q = \sum_w q^{\operatorname{inv}(w)}\)，每个 \(w\) 贡献 \(q\) 的非负整数次幂，系数为满足给定逆序数的二元字个数，故为非负整数。\(\blacksquare\)

\subsection{q-二项式系数的组合解释}\label{q-ux4e8cux9879ux5f0fux7cfbux6570ux7684ux7ec4ux5408ux89e3ux91ca}

\textbf{定理 1}（逆序数解释）：对二元字 \(w \in \{0, 1\}^n\) 恰含 \(k\) 个 \(1\)，定义其\textbf{逆序数} \(\operatorname{inv}(w) = \#\{(i, j) : i < j, w_i = 1, w_j = 0\}\)（即每个 \(1\) 右边的 \(0\) 的个数之和）。则

\[\binom{n}{k}_q = \sum_{w \in \{0,1\}^n, \; |w|_1 = k} q^{\operatorname{inv}(w)}.\]

\textbf{证明}：用归纳法。设 \(a_{n,k}(q) = \sum_{w} q^{\operatorname{inv}(w)}\) 为上述生成函数。分类讨论 \(w\) 的末位： - 若末位为 \(0\)，删除末位不影响已有 \(1\) 的逆序数，\(w\) 对应长度为 \(n-1\)、含 \(k\) 个 \(1\) 的字，逆序数相同。 - 若末位为 \(1\)，删除末位后，剩余的 \(n-k-1\) 个 \(0\) 原先均在此 \(1\) 左边（即每个 \(0\) 对逆序数的贡献为 \(1\)），因此该 \(1\) 贡献 \(n-k\) 个逆序对。\(w\) 对应长度为 \(n-1\)、含 \(k-1\) 个 \(1\) 的字，额外因子 \(q^{n-k}\)。

得递推式 \(a_{n,k}(q) = a_{n-1,k}(q) + q^{n-k} a_{n-1,k-1}(q)\)，这正是 q-二项式系数的标准递推式 \(\binom{n}{k}_q = \binom{n-1}{k}_q + q^{n-k} \binom{n-1}{k-1}_q\)（可由定义直接验证）。边界条件 \(\binom{n}{0}_q = \binom{n}{n}_q = 1\) 与 \(a_{n,0} = a_{n,n} = 1\) 一致，故 \(a_{n,k} = \binom{n}{k}_q\)。\(\blacksquare\)

\textbf{定理 2}（有限域上子空间计数）：\(\mathbb{F}_q\) 上 \(n\) 维向量空间 \(\mathbb{F}_q^n\) 中 \(k\) 维子空间的个数为 \(\binom{n}{k}_q\)。

\textbf{证明}：计算有序基的选取方式再除以子空间内基的选取方式。\(\mathbb{F}_q^n\) 中 \(k\) 个线性无关向量的有序选取数为

\[(q^n - 1)(q^n - q)\cdots(q^n - q^{k-1}),\]

因为第 \(i\) 个向量不能落在前 \(i-1\) 个向量张成的 \(q^{i-1}\) 元子空间中。给定一个 \(k\) 维子空间 \(W\)，\(W\) 中有序基的选取数为

\[(q^k - 1)(q^k - q)\cdots(q^k - q^{k-1}).\]

每个 \(k\) 维子空间对应同样多的有序基，因此子空间个数为

\[\frac{(q^n - 1)(q^n - q)\cdots(q^n - q^{k-1})}{(q^k - 1)(q^k - q)\cdots(q^k - q^{k-1})} = \frac{\prod_{i=0}^{k-1} (q^{n-i} - 1)}{\prod_{i=0}^{k-1} (q^{k-i} - 1)} \cdot \frac{\prod_{i=0}^{k-1} q^{k-1-i}}{\prod_{i=0}^{k-1} q^{k-1-i}} = \binom{n}{k}_q.\]

\(\blacksquare\)

\textbf{推论 1}：\(\mathbb{F}_q^n\) 中 \(k\) 维子空间的个数也是 \(k\) 维商空间的个数，亦是 Grassmann 流形 \(\operatorname{Gr}(k, n)\) 在 \(\mathbb{F}_q\) 上的点数。

\textbf{证明}：对偶性：\(\mathbb{F}_q^n\) 中 \(k\) 维子空间与 \((n-k)\) 维商空间一一对应（取零化子），故 \(k\) 维商空间个数为 \(\binom{n}{n-k}_q = \binom{n}{k}_q\)。Grassmann 流形 \(\operatorname{Gr}(k, n)\) 的参数化：\(\operatorname{Gr}(k, n)\) 的参数化将 \(\mathbb{F}_q\) 上 \(k\) 维子空间一一对应到 \(\mathbb{F}_q\)-有理点，故 \(\mathbb{F}_q\)-点数为 \(\binom{n}{k}_q\)。\(\blacksquare\)

\subsection{q-二项式定理}\label{q-ux4e8cux9879ux5f0fux5b9aux7406}

\textbf{定理 3}（q-二项式定理）：

\[\prod_{i=0}^{n-1} (1 + q^i x) = \sum_{k=0}^n q^{\binom{k}{2}} \binom{n}{k}_q x^k.\]

\emph{证明（组合证明）}：展开左式 \(\prod_{i=0}^{n-1} (1 + q^i x)\)。每次乘法选择 \(q^i x\) 项或 \(1\) 项。若选了 \(k\) 个 \(q^i x\) 项（对应脚标集合 \(S = \{i_1 < i_2 < \cdots < i_k\} \subseteq \{0, \ldots, n-1\}\)），则贡献为 \(q^{i_1 + i_2 + \cdots + i_k} x^k\)。所有这样的 \(S\) 的 \(q\) 指数之和为 \(\sum_{S} q^{\sum_{j=1}^k i_j}\)。

将 \(S\) 对应到二元字 \(w\)：\(w_j = 1\) 当且仅当 \(j \in S\)，则 \(\sum_{j=1}^k i_j = \binom{k}{2} + \operatorname{inv}(w)\)（因为 \(i_1 + \cdots + i_k = 0 + 1 + \cdots + (k-1) + \operatorname{inv}(w)\)，将 \(S\) 的最小元右移 \(0\) 位、次小元右移 \(1\) 位等）。由定理 1，

\[\sum_{S} q^{\sum i_j} = q^{\binom{k}{2}} \sum_{w} q^{\operatorname{inv}(w)} = q^{\binom{k}{2}} \binom{n}{k}_q.\]

\(\blacksquare\)

\textbf{推论 2}（\(q \to 1\) 极限）：\(\lim_{q \to 1} \prod_{i=0}^{n-1} (1 + q^i x) = (1 + x)^n = \sum_{k=0}^n \binom{n}{k} x^k\)，恢复经典二项式定理。

\textbf{推论 3}（\(x = -1\) 的特例）：

\[\prod_{i=0}^{n-1} (1 - q^i) = \sum_{k=0}^n (-1)^k q^{\binom{k}{2}} \binom{n}{k}_q.\]

\subsection{Euler 五边形数定理}\label{euler-ux4e94ux8fb9ux5f62ux6570ux5b9aux7406}

无穷乘积 \((q; q)_\infty = \prod_{k=1}^\infty (1 - q^k)\) 是组合数学中最精妙的生成函数之一------一方面，其展开系数给出了偶数部分不同分拆与奇数部分不同分拆的计数差；另一方面，Euler 在 1750 年惊人地发现该乘积展开中的非零系数仅出现在五角数指标上。五角数 \(m(3m-1)/2\) 的命名源自其几何意义：\(m\) 个点可以排列成五边形，计数为 \(1, 5, 12, 22, \ldots\)，而广义五角数进一步包含了 \(m(3m+1)/2\)。

\textbf{定义 4}（q-Pochhammer 符号）：对 \(n \in \mathbb{N} \cup \{\infty\}\)，

\[(a; q)_n = \prod_{k=0}^{n-1} (1 - a q^k), \qquad (a; q)_\infty = \prod_{k=0}^\infty (1 - a q^k) \quad (|q| < 1).\]

\textbf{定理 4}（Euler 五边形数定理，1750）：

\[(q; q)_\infty = \sum_{m=-\infty}^\infty (-1)^m q^{m(3m-1)/2} = 1 + \sum_{m=1}^\infty (-1)^m \left(q^{m(3m-1)/2} + q^{m(3m+1)/2}\right).\]

指数 \(m(3m-1)/2\) 和 \(m(3m+1)/2\)（\(m \geq 1\)）正是广义五边形数：\(1, 2, 5, 7, 12, 15, 22, 26, \ldots\)

\textbf{证明}（Franklin 对合，1881）：将 \((q; q)_\infty = \prod_{n \geq 1} (1 - q^n)\) 展开，\(q^N\) 的系数为将 \(N\) 分拆为偶数个不同部分的方式数减去分拆为奇数个不同部分的方式数。对 \(N\) 的任意不同部分分拆 \(\lambda = (\lambda_1 > \lambda_2 > \cdots > \lambda_t)\)，绘制 Ferrers 图。设 \(s\) 为最小部分（最底部行的方块数），\(r\) 为从右上角沿对角线向左下方向延伸的连续方块数（即满足 \(\lambda_1, \lambda_2, \ldots\) 中连续递减的行数）。定义 Franklin 对合 \(\iota\)：（a）若 \(s \leq r\)：将底部行移至最右列的右端，\(t\) 变为 \(t-1\)。（b）若 \(s > r\)：将最右列移至底部行的下方，\(t\) 变为 \(t+1\)。该对合是符号翻转的（改变 \(t\) 的奇偶性），不动点仅在 Ferrers 图无法移动时发生------这意味着 \(s = r\) 或 \(s = r+1\) 且底部行与最右列相连。这两类形状恰对应指数 \(m(3m-1)/2\)（\(s=r=m\)）和 \(m(3m+1)/2\)（\(s=m+1, r=m\)），每类各有一个不动点。因此所有非不动项相互抵消，仅剩不动点项贡献 \((-1)^m\)，定理得证。\(\blacksquare\)

\textbf{推论 4}（五边形数定理的划分论形式）：设 \(p(n)\) 为 \(n\) 的划分个数，则

\[p(n) = p(n-1) + p(n-2) - p(n-5) - p(n-7) + p(n-12) + p(n-15) - \cdots,\]

其中下标取五边形数，符号为 \(+ + - - + + - - \cdots\)。

\subsection{Rogers-Ramanujan 恒等式}\label{rogers-ramanujan-ux6052ux7b49ux5f0f}

\textbf{定理 5}（第一 Rogers-Ramanujan 恒等式，1894/1919）：

\[\sum_{n=0}^\infty \frac{q^{n^2}}{(q; q)_n} = \frac{1}{(q; q^5)_\infty (q^4; q^5)_\infty}.\]

\textbf{证明}（Rogers 的 \(q\)-差分方程法）：设 \(H_k(q) = \sum_{n=0}^\infty \frac{q^{n^2+kn}}{(q; q)_n}\)。由级数重排，\(H_k(q) - H_{k+1}(q) = q^{k+1} H_{k+2}(q)\)，且 \(H_k(q) = H_{k+1}(q) + q^{k+1} H_{k+2}(q)\)。将 \(G(q) = H_0(q) = \sum_{n=0}^\infty q^{n^2}/(q; q)_n\) 代入此递推，解得 \(G(q) = \frac{1}{(q; q)_\infty} \sum_{m=-\infty}^\infty (-1)^m q^{m(5m-1)/2}\)。由 Jacobi 三重积恒等式，\(\sum_{m=-\infty}^\infty (-1)^m q^{m(5m-1)/2} = (q; q)_\infty \prod_{k=1}^\infty \frac{1}{(1-q^{5k-4})(1-q^{5k-1})} = \frac{(q; q)_\infty}{(q; q^5)_\infty (q^4; q^5)_\infty}\)，故 \(G(q) = \frac{1}{(q; q^5)_\infty (q^4; q^5)_\infty}\)。\(\blacksquare\)

\textbf{定理 6}（第二 Rogers-Ramanujan 恒等式）：

\[\sum_{n=0}^\infty \frac{q^{n^2 + n}}{(q; q)_n} = \frac{1}{(q^2; q^5)_\infty (q^3; q^5)_\infty}.\]

\textbf{证明}：设 \(H_k(q) = \sum_{n=0}^\infty \frac{q^{n^2+kn}}{(q; q)_n}\)，则 \(H_1(q) = \sum_{n=0}^\infty \frac{q^{n^2+n}}{(q; q)_n}\) 为第二恒等式的左边。由 \(q\)-差分方程 \(H_k(q) = H_{k+1}(q) + q^{k+1} H_{k+2}(q)\)，取 \(k=1\) 得 \(H_1(q) = H_2(q) + q^2 H_3(q)\)。逐步迭代该递推关系，利用 \(H_k(q) \to 1\)（\(k \to \infty\)），得 \(H_1(q) = \frac{1}{(q; q)_\infty} \sum_{m=-\infty}^\infty (-1)^m q^{m(5m+3)/2}\)。由 Jacobi 三重积恒等式，\(\sum_{m=-\infty}^\infty (-1)^m q^{m(5m+3)/2} = (q; q)_\infty \prod_{k=1}^\infty \frac{1}{(1-q^{5k-3})(1-q^{5k-2})} = \frac{(q; q)_\infty}{(q^2; q^5)_\infty (q^3; q^5)_\infty}\)。故 \(H_1(q) = \frac{1}{(q^2; q^5)_\infty (q^3; q^5)_\infty}\)。\(\blacksquare\)

\emph{组合解释}：第一等式左边 \(\sum q^{n^2} / (q; q)_n\) 的 \(q^N\) 系数是 \(N\) 的分拆中相邻部分之差 \(\geq 2\) 的分拆数。右边无限乘积 \(\prod_{k \geq 1} (1 - q^{5k-4})^{-1} (1 - q^{5k-1})^{-1}\) 的 \(q^N\) 系数是 \(N\) 的分拆中所有部分模 \(5\) 余 \(1\) 或 \(4\) 的分拆数。两者相等。第二等式类似：左边对应相邻部分之差 \(\geq 2\) 且最小部分 \(\geq 2\) 的分拆，右边对应所有部分模 \(5\) 余 \(2\) 或 \(3\) 的分拆。

这两个恒等式将受限分拆（模 \(5\) 同余条件）与无限和（\(q\)-超几何级数特例）联系起来，是模形式理论和表示论（\(\widehat{\mathfrak{sl}_2}\) 的 Virasoro 特征公式）的基础。其证明通常用到 \(q\)-差分方程（Rogers）或 Bailey 引理（Andrews-Gordon 推广）。

\subsection{Heine 变换公式}\label{heine-ux53d8ux6362ux516cux5f0f}

\textbf{定义 5}（基本超几何级数）：\(q\)-超几何级数 \(_r\phi_s\) 定义为

\[_r\phi_s\left(\begin{matrix} a_1, \ldots, a_r \\ b_1, \ldots, b_s \end{matrix}; q, z\right) = \sum_{n=0}^\infty \frac{(a_1; q)_n \cdots (a_r; q)_n}{(q; q)_n (b_1; q)_n \cdots (b_s; q)_n} \left[(-1)^n q^{\binom{n}{2}}\right]^{1+s-r} z^n.\]

\textbf{定理 7}（Heine 变换公式，1847）：

\[_2\phi_1\left(\begin{matrix} a, b \\ c \end{matrix}; q, z\right) = \frac{(b; q)_\infty (az; q)_\infty}{(c; q)_\infty (z; q)_\infty} \,_2\phi_1\left(\begin{matrix} c/b, z \\ az \end{matrix}; q, b\right).\]

\textbf{证明}：设左边 \(L(a,b,c;z) = \sum_{n=0}^\infty \frac{(a;q)_n (b;q)_n}{(q;q)_n (c;q)_n} z^n\)。利用 \(q\)-差分算子 \(D_q f(z) = \frac{f(z)-f(qz)}{(1-q)z}\)，验证 \(L\) 满足 \(q\)-差分方程：\((1-q)(c - abqz) L(a,b,c;qz) = (c - (a+b)qz + abqz^2) L(a,b,c;z) - (c - qz)(1 - z) L(a,b,c;q^{-1}z)\)。右边 \(R(a,b,c;z) = \frac{(b;q)_\infty (az;q)_\infty}{(c;q)_\infty (z;q)_\infty} \,_2\phi_1(c/b,z;az;q,b)\) 满足同样的 \(q\)-差分方程。比较 \(z=0\) 处两边均为 \(1\)，由 \(q\)-差分方程解的唯一性，\(L = R\) 对所有 \(|q|,|z| < 1\) 成立。在 \(q \to 1\) 时，此变换还原为 Euler 变换 \(_2F_1(a,b;c;z) = (1-z)^{-a} \,_2F_1(a,c-b;c;z/(z-1))\)。\(\blacksquare\)

这是经典 Gauss 超几何函数恒等式 \(_2F_1(a,b;c;z)\) 的 \(q\)-类似。当 \(q \to 1\) 时，Heine 变换还原为 Euler-Pfaff 变换 \(_2F_1(a,b;c;z) = (1-z)^{-a} \,_2F_1(a,c-b;c;z/(z-1))\)。

Heine 变换是 \(q\)-级数理论的基本工具，由它可以推出 Rogers-Ramanujan 恒等式、Bailey 链方法以及大量划分恒等式。证明基于 \(q\)-差分方程和归纳法，核心是验证两边满足相同的 \(q\)-递推关系并比较 \(z = 0\) 处的初始条件。

\newpage

\chapter{04-离散数学/01-组合数学/06-图论.md}\label{ux79bbux6563ux6570ux5b6601-ux7ec4ux5408ux6570ux5b6606-ux56feux8bba.md}

\chapter{图论}\label{ux56feux8bba}

\hrulefill

\section{第90章：图论基本概念与树}\label{ux7b2c90ux7ae0ux56feux8bbaux57faux672cux6982ux5ff5ux4e0eux6811}

图论是离散数学的核心分支，研究由顶点和边构成的离散结构的性质与算法。本章建立图论的基本概念框架，并系统讨论树的理论。

\subsection{图的基本定义}\label{ux56feux7684ux57faux672cux5b9aux4e49}

\textbf{定义 53.1（简单图）} 一个\textbf{简单图} \(G\) 是一个有序对 \((V(G), E(G))\)，其中 \(V(G)\) 是一个非空有限集，称为\textbf{顶点集}；\(E(G)\) 是 \(V(G)\) 的二元无序子集构成的集合，称为\textbf{边集}。每条边 \(\{u, v\} \in E(G)\) 连接两个不同的顶点 \(u \neq v\)，且任意两个顶点之间至多有一条边。

若允许 \(u = v\)，则称对应边为\textbf{自环}（loop）。若允许两个顶点间有多条边，则称对应图为\textbf{多重图}（multigraph）。允许多重边但不允许自环的图称为\textbf{一般图}。

\textbf{定义 53.2（有向图）} 一个\textbf{有向图}（directed graph）\(D\) 是一个有序对 \((V(D), A(D))\)，其中 \(A(D) \subseteq V(D) \times V(D)\) 是\textbf{弧集}。每条弧 \((u, v) \in A(D)\) 是从 \(u\) 指向 \(v\) 的有向边。

\textbf{定义 53.3（度序列）} 设 \(G\) 是简单图，\(V(G) = \{v_1, \ldots, v_n\}\)。顶点 \(v\) 的\textbf{度} \(\deg(v)\) 是与 \(v\) 关联的边数。序列 \((\deg(v_1), \ldots, \deg(v_n))\)（通常按非增排列）称为 \(G\) 的\textbf{度序列}。

\textbf{定理 53.4（握手引理）} 对于任意图 \(G\)， \[\sum_{v \in V(G)} \deg(v) = 2|E(G)|.\] 特别地，奇度顶点的个数为偶数。

\textbf{证明} 每条边贡献两个度计数，故总和为 \(2|E|\)。

\textbf{定理 53.5（Havel-Hakimi定理）} 设 \(d_1 \geq d_2 \geq \cdots \geq d_n \geq 0\) 是整数序列。则存在简单图以该序列为度序列当且仅当 \(\sum_{i=1}^{n} d_i\) 为偶数且序列 \[(d_2-1, d_3-1, \ldots, d_{d_1+1}-1, d_{d_1+2}, \ldots, d_n)\] 经重排为非增序列后也是某个图的度序列。

\textbf{证明} 必要性：若存在图 \(G\) 以 \((d_1, \ldots, d_n)\) 为度序列，设 \(v_1\) 是度最大的顶点，\(\deg(v_1) = d_1\)。将 \(v_1\) 与 \(v_2, \ldots, v_{d_1+1}\) 相连后删除 \(v_1\)，剩余图的度序列恰为所述序列（经适当的对应重排）。充分性可通过归纳法证明：由归纳假设构造出 \(n-1\) 个顶点的图，再添加最大度顶点并与适当顶点相连即得原序列的图。算法的计算复杂度为 \(O(n^2 \log n)\)。\(\blacksquare\)

\textbf{定义 53.6（图的同构）} 两个图 \(G\) 和 \(H\) 称为\textbf{同构}的，记为 \(G \cong H\)，若存在双射 \(\varphi: V(G) \to V(H)\) 使得对任意 \(u, v \in V(G)\)，\(\{u, v\} \in E(G)\) 当且仅当 \(\{\varphi(u), \varphi(v)\} \in E(H)\)。

\textbf{定义 53.7（完全图与补图）} \(n\) 个顶点的\textbf{完全图} \(K_n\) 是所有可能的 \(\binom{n}{2}\) 条边均存在的简单图。图 \(G\) 的\textbf{补图} \(\overline{G}\) 满足 \(V(\overline{G}) = V(G)\) 且 \(\{u, v\} \in E(\overline{G}) \iff \{u, v\} \notin E(G)\)。

\subsection{图的连通性}\label{ux56feux7684ux8fdeux901aux6027}

\textbf{定义 53.8（连通图）} 图 \(G\) 是\textbf{连通的}，若对任意 \(u, v \in V(G)\) 存在连接 \(u\) 和 \(v\) 的路径。极大连通子图称为 \(G\) 的\textbf{连通分支}。

\textbf{定义 53.9（割点与桥）} 顶点 \(v \in V(G)\) 称为\textbf{割点}（cut-vertex），若 \(G - v\) 的连通分支数大于 \(G\) 的连通分支数。边 \(e \in E(G)\) 称为\textbf{桥}（bridge），若 \(G - e\) 的连通分支数大于 \(G\) 的连通分支数。

\textbf{定理 53.10（割点特征）} 在连通图 \(G\) 中，\(v\) 是割点当且仅当存在顶点 \(x, y \neq v\) 使得 \(v\) 在 \(x\) 到 \(y\) 的每条路上。

\textbf{证明}（\(\Rightarrow\)）设 \(v\) 是割点，则 \(G - v\) 不连通。取 \(G - v\) 的两个不同连通分支中的顶点 \(x\) 和 \(y\)。在 \(G\) 中，\(x\) 到 \(y\) 的任一条路径必经过 \(v\)（否则该路径在 \(G - v\) 中存在，与 \(x\) 和 \(y\) 在不同分支中矛盾）。故 \(v\) 在 \(x\) 到 \(y\) 的每条路上。

（\(\Leftarrow\)）设存在 \(x, y \neq v\) 使得 \(v\) 在 \(x\) 到 \(y\) 的每条路上。则在 \(G - v\) 中 \(x\) 和 \(y\) 之间没有路径，故 \(x\) 和 \(y\) 属于 \(G - v\) 的不同连通分支。因 \(G\) 是连通的而 \(G - v\) 不是，故 \(v\) 是割点。\(\blacksquare\)

\textbf{定义 53.11（\(k\)-连通性）} 图 \(G\) 是\textbf{\(k\)-连通的}（\(k\)-connected），若 \(|V(G)| > k\) 且删除任意 \(k-1\) 个顶点后图仍连通。\(G\) 的\textbf{连通度} \(\kappa(G)\) 是使得 \(G\) 为 \(k\)-连通的最大 \(k\)。\(G\) 是\textbf{\(k\)-边连通的}，若删除任意 \(k-1\) 条边后仍连通。\(G\) 的\textbf{边连通度} \(\kappa'(G)\) 类似定义。

\textbf{定理 53.12（Menger定理 - 顶点形式）} 设 \(G\) 是图，\(A, B \subseteq V(G)\)。则分离 \(A\) 和 \(B\) 所需的最少顶点数等于两两内部不相交的 \(A\)-\(B\) 路径的最大条数。

\textbf{证明} 对边的条数进行归纳。设 \(k\) 是分离 \(A\) 和 \(B\) 所需的最少顶点数。当 \(k = 0\) 时显然。假设定理对边数较小的图成立。若 \(k \geq 1\)，设 \(S\) 是极小分离集，\(|S| = k\)。任取 \(s \in S\)，将 \(G\) 收缩边 \(\{s, x\}\)（其中 \(x\) 在 \(S\) 中有邻），对新图应用归纳假设。通过仔细构造，可证原始图中存在 \(k\) 条内部不相交的路径。\(\blacksquare\)

\textbf{证明}（流方法）：对给定的图 \(G\) 和顶点集 \(A, B\)，构造有向网络：每个顶点 \(v\) 拆分为 \(v_{\text{in}}\) 和 \(v_{\text{out}}\)，其间连容量为 \(1\) 的弧；原图的每条无向边 \(uv\) 替换为两个方向弧 \(u_{\text{out}} \to v_{\text{in}}\) 和 \(v_{\text{out}} \to u_{\text{in}}\)（容量无穷大）。源点为所有 \(a \in A\) 的 \(a_{\text{out}}\)，汇点为所有 \(b \in B\) 的 \(b_{\text{in}}\)。此网络中整数最大流的值等于两两内部不相交的 \(A\)-\(B\) 路径的最大条数。由最大流最小割定理（Ford-Fulkerson），最大流等于最小割容量。最小割对应分离 \(A\) 和 \(B\) 的最小顶点集（容量为1的弧被割断对应删去该顶点）。此即Menger定理的顶点形式。边形式的证明类似------将边替换为容量为1的弧。\(\blacksquare\)

\textbf{定理 53.13（Menger定理 - 边形式）} 图 \(G\) 中连接两个顶点 \(u\) 和 \(v\) 的边不相交路径的最大条数等于分离 \(u\) 和 \(v\) 所需的最少边数。

\textbf{证明} 构造线图 \(L(G)\) 并将边形式归约为顶点形式：\(G\) 中边不相交的 \(u\)-\(v\) 路径对应于 \(L(G)\) 中（修改后的）顶点不相交路径。\(\blacksquare\)

\textbf{定理 53.14（Menger定理 - 全局形式）} 图 \(G\) 是 \(k\)-连通的当且仅当任意两个不同的顶点之间存在至少 \(k\) 条内部不相交的路径。

\textbf{证明}（\(\Rightarrow\)）设 \(G\) 是 \(k\)-连通的。任取两个不同的顶点 \(u\) 和 \(v\)。若 \(u\) 和 \(v\) 不相邻，令 \(A = N(u)\)（\(u\) 的邻居集），\(B = N(v)\)。由于 \(G\) 是 \(k\)-连通的，分离 \(u\) 和 \(v\) 至少需要 \(k\) 个顶点。由 Menger 顶点形式定理（定理 53.12），存在至少 \(k\) 条内部不相交的 \(u\)-\(v\) 路径。若 \(u\) 和 \(v\) 相邻，在 \(G\) 中删去边 \(uv\) 得 \(G'\)，\(G'\) 是 \((k-1)\)-连通的（因删除 \(u\) 和 \(v\) 后至多再删 \(k-2\) 个顶点即不连通）。由归纳，\(G'\) 中 \(u\) 和 \(v\) 间存在至少 \(k-1\) 条内部不相交的路径，加上直接边 \(uv\) 共得 \(k\) 条。

（\(\Leftarrow\)）设任意两个不同顶点间存在至少 \(k\) 条内部不相交的路径。若删除 \(k-1\) 个顶点后图不连通，则存在两个顶点 \(u\) 和 \(v\) 在 \(G\) 中但不在同一剩余连通分支中。然而 Menger 定理保证删除 \(k-1\) 个顶点无法分离 \(u\) 和 \(v\)（因存在 \(k\) 条内部不相交的 \(u\)-\(v\) 路径），矛盾。故 \(G\) 是 \(k\)-连通的。\(\blacksquare\)

\subsection{树与森林}\label{ux6811ux4e0eux68eeux6797}

\textbf{定义 53.15（树与森林）} 不含圈的连通图称为\textbf{树}（tree）。不含圈的图（的每个连通分支是树）称为\textbf{森林}（forest）。

\textbf{定理 53.16（树的基本等价条件）} 设 \(T\) 是 \(n\) 个顶点的图。以下条件等价：（i）\(T\) 是树；（ii）\(T\) 连通且恰有 \(n-1\) 条边；（iii）\(T\) 无圈且恰有 \(n-1\) 条边；（iv）\(T\) 中任意两点间恰有一条路径。

\textbf{证明} 采用循环推理：(i)\(\Rightarrow\)(ii)\(\Rightarrow\)(iii)\(\Rightarrow\)(iv)\(\Rightarrow\)(i)。

(i)\(\Rightarrow\)(ii)：设 \(T\) 是树（连通且无圈）。对 \(n\) 归纳证明边数为 \(n-1\)。\(n=1\) 时边数为 \(0=n-1\)。设 \(n>1\)，取 \(T\) 的一条最长路径的端点 \(v\)（必为叶子，\(\deg(v)=1\)）。删除 \(v\) 得 \(T'\)，\(T'\) 仍无圈且连通，故为树。由归纳假设 \(|E(T')|=(n-1)-1\)，故 \(|E(T)|=|E(T')|+1=n-1\)。连通性由定义直接得到。

(ii)\(\Rightarrow\)(iii)：设 \(T\) 连通且 \(|E(T)|=n-1\)。需证 \(T\) 无圈。反设 \(T\) 含圈 \(C\)。由 \(T\) 连通，可逐次删去圈上的边（每次删边保持连通性），直至得到连通无圈子图 \(T'\)。\(T'\) 是树，故 \(|E(T')|=|V(T')|-1=n-1\)。但删去了至少一条边，故 \(|E(T')|<|E(T)|=n-1\)，矛盾。故 \(T\) 无圈。

(iii)\(\Rightarrow\)(iv)：设 \(T\) 无圈且 \(|E(T)|=n-1\)。首先证 \(T\) 连通：设 \(T\) 有 \(c\) 个连通分支。每分支是无圈连通图（树），第 \(j\) 分支的边数为 \(n_j-1\)（其中 \(n_j\) 为顶点数）。总和：\(|E|=\sum(n_j-1)=n-c\)。由已知 \(|E|=n-1\)，故 \(c=1\)，即 \(T\) 连通。在树中，任意两点间存在路径（由连通性）。若有两条不同路径，则将形成圈（从起点沿第一条路径到终点再沿第二条路径反回），与无圈矛盾。故路径唯一。

(iv)\(\Rightarrow\)(i)：设任意两点间恰有一条路径。由连通性得 \(T\) 连通。若有圈，则圈上任意两点间存在两条不同的路径（沿圈的两个方向），与唯一性矛盾。故 \(T\) 无圈，即为树。\(\blacksquare\)

\textbf{定理 53.17（Cayley公式）} \(n\) 个标号顶点的不同树的个数为 \(n^{n-2}\)。

\textbf{证明（Prufer编码）} 构造双射 \(\varphi: \mathcal{T}_n \to \{1, 2, \ldots, n\}^{n-2}\)，其中 \(\mathcal{T}_n\) 是 \([n] = \{1, \ldots, n\}\) 上全体标号树。

编码算法：给定树 \(T\)，重复以下操作 \(n-2\) 次------找出当前树中标号最小的叶子，记录其唯一邻居的标号，删除该叶子。得到长度为 \(n-2\) 的序列 \((a_1, \ldots, a_{n-2})\)，其中每个 \(a_i \in [n]\)。

解码算法：给定序列 \((a_1, \ldots, a_{n-2})\)，令 \(L\) 为不在该序列中出现的标号集合（恰有 \(n - (n-2) = 2\) 个）。反复进行：取 \(L\) 中最小元 \(b\)，连接 \(b\) 与 \(a_1\)，从 \(L\) 中删除 \(b\)，从序列中删除 \(a_1\)；若 \(a_1\) 不再出现在剩余序列中，将其加入 \(L\)。最后连接 \(L\) 中剩余的两点。该过程恰重构出原始树。由于序列每个位置有 \(n\) 种选择且两过程互逆，故树的总数为 \(n^{n-2}\)。\(\blacksquare\)

\textbf{推论 53.18} \(n\) 个标号顶点的森林（假定为有根森林，指定 \(k\) 个连通分支的根）的个数为 \(k n^{n-k-1}\)。

\subsection{最小生成树}\label{ux6700ux5c0fux751fux6210ux6811}

\textbf{定义 53.19（最小生成树）} 给定连通图 \(G\) 和权函数 \(w: E(G) \to \mathbb{R}\)，\textbf{最小生成树}（Minimum Spanning Tree, MST）是 \(G\) 的一棵生成树 \(T\) 使得 \(w(T) = \sum_{e \in E(T)} w(e)\) 最小。

\textbf{算法 53.20（Kruskal算法）} （1）初始化 \(T = \emptyset\)；（2）将边按权值非递减排序；（3）依序遍历每条边，若加入该边不构成圈，则将其加入 \(T\)；（4）当 \(|T| = |V(G)| - 1\) 时停止，输出 \(T\)。

\textbf{算法 53.21（Prim算法）} （1）任选顶点 \(r\) 为根，初始化树 \(T = (\{r\}, \emptyset)\)；（2）在所有一端在 \(T\) 中、一端不在 \(T\) 中的边中选取权最小的边，将其及其关联的新顶点加入 \(T\)；（3）重复至 \(T\) 包含所有顶点。

\textbf{定理 53.22（拟阵贪心框架下的最优性）} Kruskal算法和Prim算法均产生最小生成树。

\textbf{证明（拟阵框架）} 图 \(G\) 的边集 \(E\) 的所有无圈子集构成一个\textbf{拟阵} \(\mathcal{M}(G) = (E, \mathcal{I})\)，其中 \(\mathcal{I}\) 是全体森林。拟阵的基恰为生成树（假设 \(G\) 连通）。MST问题等价于在拟阵 \(\mathcal{M}(G)\) 中求最小权基。拟阵的贪婪算法（按权递增序选择不破坏独立性的元素）保证得到最小权基------设贪婪算法输出 \(B = \{e_1, \ldots, e_{n-1}\}\)（按选择顺序），若存在更小权的基 \(B^*\)，考虑最小下标 \(i\) 使得 \(e_i \notin B^*\)，则由拟阵交换性质可构造出矛盾。Kruskal算法正是拟阵贪婪算法的直接实现，Prim算法通过维护切的最优边保证拟阵最优性条件。\(\blacksquare\)

\textbf{定理 53.23（割性质与圈性质）} （割性质）对任意割 \(C\)，\(C\) 中权最小的边属于某个MST。（圈性质）对任意圈，圈上权最大的边不属于任何MST（若最大权唯一）。

\textbf{证明}：（割性质）设 \(C\) 为图 \(G\) 的一个割，\(e\) 为 \(C\) 中权最小的边。假设存在不含 \(e\) 的MST \(T\)。在 \(T\) 中加入 \(e\) 形成唯一圈，该圈必包含 \(C\) 中另一条边 \(f \neq e\)（因为 \(e\) 是连接割两部分的边）。由 \(e\) 的最小性，\(w(e) \leq w(f)\)。在 \(T\) 中用 \(e\) 替换 \(f\) 得到新的生成树 \(T'\)，且 \(w(T') \leq w(T)\)。故存在包含 \(e\) 的MST。

（圈性质）设 \(C\) 为一个圈，\(e\) 为 \(C\) 上权最大的边。假设MST \(T\) 包含 \(e\)。删去 \(e\) 将 \(T\) 分为两个连通分支。圈 \(C\) 必包含连接这两个分支的另一条边 \(f\)。由 \(e\) 的最大性，\(w(f) \leq w(e)\)。用 \(f\) 替换 \(e\) 得 \(T'\) 且 \(w(T') \leq w(T)\)。若 \(w(f) < w(e)\)，则 \(T'\) 是更优的MST，矛盾。因此若 \(e\) 是严格最大权，它不可能属于任何MST。\(\blacksquare\)

\subsection{图的矩阵表示}\label{ux56feux7684ux77e9ux9635ux8868ux793a}

\textbf{定义 53.24（邻接矩阵）} 设 \(G\) 是 \(n\) 个顶点的简单图。\(G\) 的\textbf{邻接矩阵} \(A = A(G) \in \mathbb{R}^{n \times n}\) 定义为 \[A_{ij} = \begin{cases} 1, & \{v_i, v_j\} \in E(G) \\ 0, & \text{否则} \end{cases}\]

\textbf{定义 53.25（度矩阵与Laplace矩阵）} 设 \(D = \operatorname{diag}(\deg(v_1), \ldots, \deg(v_n))\) 为\textbf{度矩阵}。\(G\) 的\textbf{Laplace矩阵}（Laplacian matrix）定义为 \(L = D - A\)。等价地， \[L_{ij} = \begin{cases} \deg(v_i), & i = j \\ -1, & i \neq j, \{v_i, v_j\} \in E(G) \\ 0, & \text{否则} \end{cases}\]

\textbf{命题 53.26（Laplace矩阵的基本性质）} （1）\(L\) 是对称半正定矩阵；（2）\(L\) 的最小特征值为 \(0\)，对应特征向量为全1向量 \(\mathbf{1} = (1, 1, \ldots, 1)^T\)；（3）\(0\) 作为 \(L\) 的特征值的重数等于 \(G\) 的连通分支数。

\textbf{定义 53.27（关联矩阵）} 对 \(G\) 的每条边任意指定方向，定义\textbf{关联矩阵} \(B \in \mathbb{R}^{n \times m}\)（\(m = |E|\)）： \[B_{ve} = \begin{cases} 1, & \text{若 } v \text{ 是边 } e \text{ 的尾} \\ -1, & \text{若 } v \text{ 是边 } e \text{ 的头} \\ 0, & \text{否则} \end{cases}\] 则有 \(L = BB^T\)，与方向的选择无关。

\textbf{定理 53.28（Kirchhoff矩阵树定理）} 设 \(G\) 是连通图，\(L\) 是其Laplace矩阵。则 \(G\) 的生成树个数等于 \(L\) 的任意一个 \((n-1) \times (n-1)\) 主子式的行列式值。即对任意 \(i \in \{1, \ldots, n\}\)， \[\tau(G) = \det(L_i),\] 其中 \(L_i\) 是 \(L\) 删去第 \(i\) 行和第 \(i\) 列得到的矩阵。

\textbf{证明（概要）} 利用Cauchy-Binet公式。\(L = BB^T\)，其中 \(B\) 是 \((n-1) \times m\) 的简约关联矩阵（删去某一行）。Cauchy-Binet公式给出： \[\det(BB^T) = \sum_{S \subseteq [m], |S|=n-1} \det(B_S)^2,\] 其中 \(B_S\) 是 \(B\) 关于列子集 \(S\) 的子矩阵。\(\det(B_S) = \pm 1\) 当且仅当 \(S\) 对应的边集构成一棵生成树，否则 \(\det(B_S) = 0\)。故 \(\det(L_i)\) 恰为生成树的总数。该定理为Cayley公式提供了独立证明：\(K_n\) 的Laplace矩阵为 \(L = nI - J\)，删去一行一列后主子式值为 \(n^{n-2}\)。\(\blacksquare\)

图论是离散数学中历史最悠久的分支之一，其源头可追溯至 Euler 对 Königsberg 七桥问题的研究（1736 年）------这是数学史上被公认为第一篇图论论文的工作。从 19 世纪 Cayley 用树对化学分子的计数、Kirchhoff 用生成树分析电网，到 20 世纪 Tutte 的因子理论和 Erdős-Rényi 的随机图模型，图论经历了从具体问题驱动到统一理论体系的蜕变。图论的核心力量在于它将组合结构可视化，并通过递归（删边-缩边递推）、对偶（平面图与对偶图）、代数（邻接矩阵和 Laplace 矩阵的谱）和概率（随机图方法）四重方法深入分析。本章在组合学基础（V6 Ch 32 计数原理）之上建立图的形式化定义，涵盖树与连通性、平面图与 Euler 公式、Euler 回路与 Hamilton 回路，并引入有向图（竞赛图的 Landau 定理）和谱图论导引，为第 49 章的匹配与网络流和第 50 章的染色与 Ramsey 理论奠定图论基础。

\subsection{128.1 图的基本概念}\label{ux56feux7684ux57faux672cux6982ux5ff5}

\textbf{定义 128.1}（图）：一个\textbf{图} \(G = (V, E)\) 由顶点集 \(V\)（有限集）和边集 \(E \subseteq V \times V\) 组成。若无向图，边 \(\{u, v\}\) 是无序对；若有向图，边 \((u, v)\) 是有序对（弧）。

\textbf{定义 128.2}（基本概念）： - \textbf{度}：顶点 \(v\) 的度 \(\deg(v)\) 是与 \(v\) 关联的边数。在有向图中，出度 \(\deg^+(v)\) 和入度 \(\deg^-(v)\) 分别计数出边和入边。 - \textbf{握手引理}：\(\sum_{v \in V} \deg(v) = 2|E|\)。（每条边贡献两个度） - \textbf{路}：长度为 \(k\) 的路是顶点序列 \(v_0 v_1 \cdots v_k\)，其中 \(\{v_i, v_{i+1}\} \in E\)，且顶点互异。 - \textbf{回路}：起点和终点相同的路（\(v_0 = v_k\)，\(k \geq 3\)）。 - \textbf{连通}：图 \(G\) 是连通的，如果任意两个顶点之间存在路。 - \textbf{连通分支}：图的最大连通子图。

\textbf{定义 128.3}（特殊图类）： - \textbf{完全图} \(K_n\)：\(n\) 个顶点，每对顶点之间都有边（\(\binom{n}{2}\) 条边） - \textbf{二分图}：顶点集可划分为 \(V = A \cup B\)，所有边都在 \(A\) 和 \(B\) 之间 - \textbf{完全二分图} \(K_{m,n}\)：\(|A| = m, |B| = n\)，\(A\) 和 \(B\) 之间所有可能的边 - \textbf{路图} \(P_n\)：\(n\) 个顶点排成一条路 - \textbf{回路图} \(C_n\)：\(n\) 个顶点排成一个圈

\textbf{命题 128.1}（二分图的无奇圈特征）：图 \(G\) 是二分图当且仅当 \(G\) 不含奇长度回路。

\emph{证明}：若 \(G\) 是二分图，路在 \(A\) 和 \(B\) 之间交替，故回路长度必为偶数。反之，固定根顶点 \(v_0\)，对每个顶点 \(v\)，令 \(d(v)\) 为从 \(v_0\) 到 \(v\) 的最短路长度，按 \(d(v)\) 的奇偶性划分 \(V\)。无奇回路保证此划分是良定义的二分划分。∎

\subsection{128.2 树与生成树}\label{ux6811ux4e0eux751fux6210ux6811}

\textbf{定义 128.4}（树）：\textbf{树}是连通无回路的图。等价刻画：\(n\) 个顶点的树恰有 \(n-1\) 条边，且任意两点间存在唯一的路。

\textbf{定理 128.2}（Cayley 公式）：\(n\) 个标号顶点的不同树的个数为 \(n^{n-2}\)。

\textbf{证明}（Prufer 编码）：每个标号树唯一对应一个长度为 \(n-2\) 的序列，每个位置可从 \(1\) 到 \(n\) 取值。编码算法：反复删除度数为 \(1\) 的最小标号顶点，记录其邻居。解码算法：给定序列 \(a_1,\ldots,a_{n-2}\)，设 \(L\) 为 \(1,\ldots,n\) 中不在序列中出现的标号集合。对 \(i=1,\ldots,n-2\)，连接 \(a_i\) 与 \(L\) 中最小元素，删除 \(L\) 中该元素，若 \(a_i\) 不再出现于后序序列中则加入 \(L\)。最终 \(L\) 中剩余两元素相连。编码与解码互逆，故双射得证。因此树的总数为 \(n^{n-2}\)。\(\blacksquare\)

Cayley 公式通过 Prufer 编码这一漂亮的双射给出了标号树的精确计数。然而，当图不是树而是一般连通图时，计数所有生成树（即包含所有顶点的树子图）就变得更为复杂。Kirchhoff 的矩阵树定理（1847 年）为这一问题提供了代数解法：生成树个数等于 Laplace 矩阵的任意 \((n-1)\) 阶主子式，将组合计数问题转化为矩阵行列式的计算。

\textbf{定义 128.5}（生成树）：连通图 \(G\) 的\textbf{生成树}是 \(G\) 的包含所有顶点的树子图。

\textbf{定理 128.3}（Kirchhoff 矩阵树定理）：连通图 \(G\) 的生成树个数等于其 Laplace 矩阵 \(L = D - A\)（\(D\) 为度对角矩阵，\(A\) 为邻接矩阵）的任意一个 \((n-1)\) 阶主子式。

\textbf{证明}：对边数 \(m\) 归纳。\(m=n-1\) 时 \(G\) 为树，\(L\) 的 \(n-1\) 阶主子式为 \(1\)。归纳步：设结论对 \(m-1\) 边成立。对图 \(G\) 中边 \(e\)，删除 \(e\) 得 \(G-e\)，收缩 \(e\) 得 \(G/e\)。由矩阵运算性质 \(\det L_{G\setminus e}=\det L_G-\det L_{G/e}\)（沿 \(e\) 展开）。由归纳假设，\(\det L_{G\setminus e}\) 为 \(G\setminus e\) 生成树数，\(\det L_{G/e}\) 为 \(G/e\) 生成树数（即含 \(e\) 的生成树数）。二者和为 \(G\) 的生成树总数，即 \(\det L_G\)。\(\blacksquare\)

\textbf{定义 128.6}（最小生成树）：带权图 \(G\) 的\textbf{最小生成树}（MST）是权和最小的生成树。Kruskal 算法（贪心选择最小权边，不形成回路）和 Prim 算法均可找到 MST。

\textbf{定理 128.4}（Kruskal 算法的正确性）：对任意连通带权图，Kruskal 算法输出最小生成树。

\textbf{证明}：使用\textbf{割性质}（cut property）：对任意割 \((S, V\setminus S)\)，连接割两端的最小权边必属于某个最小生成树。证明：设 \(e\) 为连接 \(S\) 和 \(V\setminus S\) 的最小权边。若某 MST \(T\) 不含 \(e\)，则 \(T \cup \{e\}\) 含唯一回路，该回路必包含另一条割边 \(e'\)。由割性质 \(w(e) \leq w(e')\)，故 \(T' = T \cup \{e\} \setminus \{e'\}\) 的权不超过 \(T\)，\(T'\) 仍为生成树。因此 \(e\) 属于某 MST。Kruskal 算法每次选择不形成回路的最小权边，该边连接了两个不同连通分支，即为连接这两个分支的割中的最小权边，故属于某 MST。归纳可得算法输出的边集构成最小生成树。\(\blacksquare\)

\subsection{128.3 连通性与 Menger 定理}\label{ux8fdeux901aux6027ux4e0e-menger-ux5b9aux7406-1}

\textbf{定义 128.7}（连通度）： - \textbf{点连通度} \(\kappa(G)\)：使图不连通或退化为单点所需删除的最少顶点数 - \textbf{边连通度} \(\lambda(G)\)：使图不连通所需删除的最少边数 - \(G\) 是 \(k\)-连通的，如果 \(\kappa(G) \geq k\)

\textbf{命题 128.5}（Whitney 不等式）：\(\kappa(G) \leq \lambda(G) \leq \delta(G)\)，其中 \(\delta(G)\) 是最小度。

\textbf{证明}：\(\kappa(G) \leq \lambda(G)\)：设 \(\lambda(G) = k\)，存在 \(k\) 条边 \(\{e_1,\ldots,e_k\}\) 删除后图不连通。对每条边 \(e_i = \{u_i, v_i\}\)，若 \(u_i\) 和 \(v_i\) 均非孤立点，删除 \(u_i\)（或 \(v_i\)）即可断开该边。至多删除 \(k\) 个顶点即可断开所有 \(k\) 条边，故 \(\kappa(G) \leq k = \lambda(G)\)。\(\lambda(G) \leq \delta(G)\)：取最小度顶点 \(v\)（\(\deg(v) = \delta(G)\)），删除 \(v\) 的所有关联边即可断开 \(v\) 与图的其余部分，故 \(\lambda(G) \leq \delta(G)\)。\(\blacksquare\)

\textbf{定理 128.6}（Menger 定理，点版本）：图 \(G\) 中两个不相邻顶点 \(u, v\) 之间的内部不相交的 \(u\)-\(v\) 路的最大条数等于分离 \(u\) 和 \(v\) 所需删除的最少顶点数。

\textbf{证明}：将图的每个顶点 \(x\)（\(x\neq u,v\)）拆分为 \(x_{in},x_{out}\)，边容量设为 1（入边接 \(x_{in}\)，出边接 \(x_{out}\)，\(x_{in}\to x_{out}\) 容量 1），原始无向边替换为双向有向边均容量 \(\infty\)。最大流 \(=\) 最小割（容量为整数）。每条容量 1 的内部弧限制了路径不能共享内部顶点，故最大流 \(=\) 不相交路径数，最小割 \(=\) 分离顶点数。由最大流最小割定理得证。Menger 给出了不依赖网络流的自然证明。\(\blacksquare\)

\textbf{定理 128.7}（Menger 定理，边版本）：图 \(G\) 中两个顶点 \(u, v\) 之间的边不相交的 \(u\)-\(v\) 路的最大条数等于分离 \(u\) 和 \(v\) 所需删除的最少边数。

\textbf{证明}：构建辅助有向图：将 \(G\) 的每条无向边替换为两条方向相反的有向边，容量均为 \(1\)。\(u\) 为源点，\(v\) 为汇点。由最大流最小割定理，最大 \(u\)-\(v\) 流的值等于最小 \(u\)-\(v\) 割的容量。由于每条边容量为 \(1\)，最大流的值等于边不相交的 \(u\)-\(v\) 路的最大条数（每条路贡献流量 \(1\)）。最小割的容量等于分离 \(u\) 和 \(v\) 所需删除的最少边数。因此最大边不相交路径数等于最小边割数。\(\blacksquare\)

\textbf{推论 128.8}（全局 Menger 定理）：图 \(G\) 是 \(k\)-连通的当且仅当任意两个顶点之间存在至少 \(k\) 条内部不相交的路。

\subsection{128.4 Euler 回路与 Hamilton 回路}\label{euler-ux56deux8defux4e0e-hamilton-ux56deux8def}

\textbf{定理 128.9}（Euler 定理，1736）：连通图 \(G\) 有 Euler 回路（经过每条边恰好一次的回路）当且仅当每个顶点的度数为偶数。有 Euler 路（经过每条边恰好一次的路，起终点可不同）当且仅当恰好有两个奇数度顶点。

\textbf{证明}：（\(\Rightarrow\)）若存在 Euler 回路，沿回路遍历时每次进入某顶点就沿另一条边离开，故度数为偶数。（\(\Leftarrow\)）所有顶点度数为偶数时，从任意顶点出发沿边遍历，不重复使用边，直到局部无可用边停止。停止点必为出发点（否则奇数度矛盾）。若仍有未访问边，由连通性存在与已访问子图接触的未访问边，在其上继续过程，最终合并所有子回路得 Euler 回路。Euler 路：添加连接两奇度顶点的虚边使全图偶度，找 Euler 回路后移除虚边。\(\blacksquare\)

\emph{证明}：必要性：每进入一个顶点就离开一次，故度数为偶数。充分性：使用 Fleury 算法或 Hierholzer 算法（贪心合并回路）。∎

Euler 回路问题关注的是''边''的遍历------恰经过每条边一次------其判定条件（所有顶点的度数为偶数）简洁而优美。Hamilton 回路则关注''顶点''的遍历------恰经过每个顶点一次------然而与之形成鲜明对比的是，判定一般图是否含有 Hamilton 回路是一个 NP 完全问题。尽管如此，在满足一定最小度条件的稠密图中，存在保证 Hamilton 回路成立的充分条件。

\textbf{定义 128.8}（Hamilton 回路）：Hamilton 回路是经过每个顶点恰好一次的回路。判定 Hamilton 回路的存在性是 NP 完全的。

\textbf{定理 128.10}（Dirac 定理，1952）：若 \(G\) 是 \(n \geq 3\) 个顶点的图，且 \(\delta(G) \geq n/2\)，则 \(G\) 有 Hamilton 回路。

\textbf{证明}：考虑 \(G\) 中最长路 \(P=v_1\cdots v_k\)。若 \(k=n\) 且 \(v_1\) 与 \(v_n\) 相邻则得 Hamilton 回路。否则由于 \(\delta(G)\ge n/2\)，\(v_1\) 的 \(\ge n/2\) 个邻居必在 \(P\) 上。考虑 \(v_k\) 的邻居：设 \(N(v_1)=\{v_{i_1},\ldots\}\)。若 \(v_{i-1}\in N(v_k)\) 对某 \(i\)，则 \(v_1\cdots v_{i-1}v_k v_{k-1}\cdots v_i v_1\) 构成回路。由抽屉原理和 \(\delta\ge n/2\) 保证该 \(i\) 存在。得到回路后若未覆盖全部顶点，由连通度可将回路扩展。因此最长路必为 Hamilton 路且端点相邻，证得 Hamilton 回路。\(\blacksquare\)

\textbf{定理 128.11}（Ore 定理，1960）：若对 \(G\) 中任意不相邻的顶点 \(u, v\)，有 \(\deg(u) + \deg(v) \geq n\)，则 \(G\) 有 Hamilton 回路。

\textbf{证明}：反证法。假设 \(G\) 满足 Ore 条件但无 Hamilton 回路。在 \(G\) 中添加尽可能多的边而不产生 Hamilton 回路，得极大非 Hamilton 图 \(G'\)。\(G'\) 仍有 Ore 条件（添加边只会增加度数）。\(G'\) 中必存在不相邻顶点 \(u, v\)（否则 \(G'\) 为完全图，必有 Hamilton 回路）。添加边 \(\{u,v\}\) 到 \(G'\) 产生 Hamilton 回路 \(u = v_1 \to v_2 \to \cdots \to v_n = v \to u\)。对每个 \(i\)（\(2 \leq i \leq n\)），若 \(\{u, v_i\} \in E(G')\)，则 \(\{v, v_{i-1}\} \notin E(G')\)（否则 \(u \to v_i \to \cdots \to v_n = v \to v_{i-1} \to \cdots \to v_1 = u\) 构成 \(G'\) 中的 Hamilton 回路）。因此 \(\deg_{G'}(v) \leq (n-1) - \deg_{G'}(u)\)，即 \(\deg(u) + \deg(v) \leq n-1\)，与 Ore 条件矛盾。故 \(G\) 必有 Hamilton 回路。\(\blacksquare\)

\subsection{128.5 平面图}\label{ux5e73ux9762ux56fe}

\textbf{定义 128.9}（平面图）：图 \(G\) 是\textbf{平面图}，如果它可以画在平面上且边不交叉（仅在端点相交）。平面图的画法将平面划分为\textbf{面}（face）。

\textbf{定理 128.12}（Euler 公式，1758）：对连通平面图 \(G\)，有

\[n - m + f = 2\]

其中 \(n = |V|\)，\(m = |E|\)，\(f\) 为面数（包括外部无界面）。

\emph{证明}：对边数 \(m\) 归纳。奠基：\(m = 0\) 时连通图仅有 \(n = 1\)，\(f = 1\)，\(1 - 0 + 1 = 2\) 成立。 归纳假设 \(m \geq 1\)。考虑两种情况：（1）若 \(G\) 是树，则 \(m = n - 1\)，\(f = 1\)（仅外部面），\(n - (n-1) + 1 = 2\)。（2）若 \(G\) 含回路，取回路上的任意边 \(e\)。\(G \setminus e\) 仍连通且面数减少 1（\(e\) 分隔的两个面合并为一个），边数减 1，顶点数不变。由归纳假设 \(n - (m-1) + (f-1) = 2\)，即 \(n - m + f = 2\)。∎

\textbf{推论 128.13}：对 \(n \geq 3\) 的简单平面图，\(m \leq 3n - 6\)。若图中无三角形，则 \(m \leq 2n - 4\)。

\emph{证明}：每个面至少由 3 条边围成，而每条边至多属于 2 个面，故 \(3f \leq 2m\)。代入 Euler 公式：\(2 = n - m + f \leq n - m + 2m/3 = n - m/3\)，得 \(m \leq 3n - 6\)。∎

\textbf{推论 128.14}：\(K_5\) 和 \(K_{3,3}\) 不是平面图。

\emph{证明}：\(K_5\)：\(n = 5, m = 10\)，\(10 > 3 \cdot 5 - 6 = 9\)。\(K_{3,3}\)：\(n = 6, m = 9\)，无三角形，\(9 > 2 \cdot 6 - 4 = 8\)。∎

\textbf{定理 128.15}（Kuratowski 定理，1930）：图 \(G\) 是平面图当且仅当 \(G\) 不含 \(K_5\) 或 \(K_{3,3}\) 的细分（即通过插入度数为 2 的顶点得到的子图）。

\textbf{证明}：（\(\Rightarrow\)）\(K_5\) 和 \(K_{3,3}\) 均非平面。\(K_5\)：5 顶点 10 边 \textgreater{} \(3n-6=9\) 违反平面图边数界。\(K_{3,3}\)：二分图无奇回路，若平面由 Euler 公式 \(F\le 2E/4=4\)（最小面度 \(\ge4\)），面积和 \(\sum\deg(F)=2E=18\)，矛盾。（\(\Leftarrow\)）充分性由 Wagner 定理（图论基本定理）：无 \(K_5\) 或 \(K_{3,3}\) 细分的图必为平面图。证明通过极小反例的构造分析：取顶点数最少的非平面 \(K_{5}/K_{3,3}\) 细分自由图，证明其 2-连通性，分析 2-点割和块分解，最终导出矛盾或归约到 \(K_5\)/\(K_{3,3}\) 细分。\(\blacksquare\)

\emph{充分性简述}（若不含 \(K_5\) 或 \(K_{3,3}\) 细分则平面）：由推论 128.14 知 \(K_5\) 和 \(K_{3,3}\) 本身非平面，其细分亦然。充分性证明较复杂：对顶点数归纳，利用 3-连通图的平面嵌入性质（Tutte 关于凸嵌入的定理），最终归结为不含此类子图的图必可平面嵌入。详细证明超出本教材范围。∎

\subsection{128.6 有向图与竞赛图}\label{ux6709ux5411ux56feux4e0eux7adeux8d5bux56fe}

有向图将边的无序对替换为有序弧，从而引入了方向的维度。竞赛图是有向图中最基本且最具组合结构的一类------每对顶点之间恰有一条弧，可以理解为完全图的定向。Landau 在 1953 年给出了刻画竞赛图得分序列（即各顶点的出度序列）的充要条件，其简洁优雅的形式和基于 Hall 条件的证明使其成为有向图理论中的经典结果。

\textbf{定义 128.10}（有向图与竞赛图）：\textbf{有向图} \(D = (V, A)\) 由顶点集 \(V\) 和弧集 \(A \subseteq V \times V\) 组成（弧是有序对）。\textbf{竞赛图}（tournament）是每对顶点之间恰有一条弧的有向图，即完全图的定向。竞赛图 \(T\) 的\textbf{得分序列}为各顶点的出度 \(s_1 \leq s_2 \leq \cdots \leq s_n\)。

\textbf{定理 128.16}（Landau 定理，1953）：非负整数序列 \((s_1, \ldots, s_n)\)（\(s_1 \leq \cdots \leq s_n\)）是某竞赛图的得分序列当且仅当对所有 \(k = 1,\ldots,n\)，\(\sum_{i=1}^k s_i \geq \binom{k}{2}\)，且 \(k=n\) 时等号成立。该定理的证明使用网络流模型：构造二分图，左侧为顶点对，右侧为顶点，流的存在性由 Hall 条件保证。关于竞赛图的基本结构性质，Rédei 定理（1934）断言每个竞赛图都有 Hamilton 路（有向路经过所有顶点），证明通过归纳法：删除任意顶点后归纳构造 Hamilton 路，再将删除的顶点插入适当位置。Moon 定理（1968）进一步推广：每个强连通竞赛图是顶点泛圈的（vertex-pancyclic），即每个顶点包含在所有长度 \(k = 3,\ldots,n\) 的有向回路中。竞赛图也是社会选择理论（Arrow 不可能定理的组合版本）和排序算法分析中的重要模型。

\textbf{注记 128.1}（图论与代数的联系）：图的邻接矩阵 \(A(G)\) 的谱（特征值集合）蕴含了图的深层结构信息。谱图论（spectral graph theory）通过 \(A(G)\) 和 \(L(G)\) 的特征值研究图的扩张性、等周不等式和随机游走性质。对 \(d\)-正则图，\(A(G)\) 的最大特征值为 \(d\)，第二大特征值 \(\lambda_2\) 控制图的扩张性质：\(\lambda_2\) 越小，图的边扩张越好。Alon-Boppana 定理给出 \(\lambda_2 \geq 2\sqrt{d-1} - o(1)\) 的下界，而 Ramanujan 图恰好达到此界，是扩张图和编码理论中的理想构造。图的 Laplace 矩阵 \(L(G)\) 的谱与图上的随机游走、有效电阻和生成树计数（Kirchhoff 矩阵树定理）有直接联系：\(L(G)\) 的非零特征值乘积等于 \(n\) 乘以生成树个数。

\hrulefill

\section{第91章：图的染色与平面性}\label{ux7b2c91ux7ae0ux56feux7684ux67d3ux8272ux4e0eux5e73ux9762ux6027}

\subsection{点染色与色数}\label{ux70b9ux67d3ux8272ux4e0eux8272ux6570}

\textbf{定义 54.1（点染色）} 图 \(G\) 的\textbf{正常 \(k\)-点染色}是一个函数 \(c: V(G) \to \{1, 2, \ldots, k\}\)，使得相邻顶点获不同颜色。\(G\) 的\textbf{色数} \(\chi(G)\) 是使 \(G\) 可正常染色的最小颜色数。

\textbf{定义 54.2（团数与退化性）} \(G\) 的\textbf{团数} \(\omega(G)\) 是 \(G\) 中最大完全子图的顶点数。显然 \(\chi(G) \geq \omega(G)\)。\(G\) 的\textbf{退化性} \(\delta^*(G)\) 是每个子图的最小度的最大值。

\textbf{定理 54.3（Brooks定理）} 设 \(G\) 是连通图，不是完全图也不是奇圈。则 \(\chi(G) \leq \Delta(G)\)，其中 \(\Delta(G)\) 是 \(G\) 的最大度。

\textbf{证明} 对 \(|V(G)|\) 归纳。设 \(n = |V(G)|\)，\(\Delta = \Delta(G)\)。若 \(G\) 的连通度 \(\kappa(G) \leq 1\)，则存在割点 \(v\) 将 \(G\) 分离为分支，对每个分支及 \(v\) 应用归纳假设。故可设 \(G\) 是 \(2\)-连通的，从而 \(\delta(G) \geq 2\) 且没有割点。

关键情形：\(G\) 是 \(3\)-连通图。固定顶点 \(x\)，考虑从 \(x\) 出发的广度优先搜索树 \(T\)。设 \(v_n = x\)，令 \(v_1, v_2\) 是 \(x\) 的两个邻居。由于 \(G\) 不是完全图，可选择 \(v_1, v_2\) 非相邻。按BFS序逆序染色：对 \(i = n-1, \ldots, 1\)，每个顶点 \(v_i\) 在图中有一个编号更大的邻点（其父节点），故在染色时 \(v_i\) 至多有 \(\Delta - 1\) 个已染色的邻居，可用不超过 \(\Delta\) 种颜色完成染色。\(v_n = x\) 由于 \(v_1\) 和 \(v_2\) 获相同颜色（因它们不相邻且与 \(x\) 的剩余 \(\Delta-2\) 个邻居共占用至多 \(\Delta-2\) 种颜色），故 \(x\) 可用一种颜色完成染色。\(\Delta = 2\) 或 \(G\) 不是 \(3\)-连通的情形需单独处理，可利用归纳假设或直接验证。\(\blacksquare\)

\textbf{定理 54.4（四色定理）} 每个平面图是 \(4\)-可染色的，即 \(\chi(G) \leq 4\)。

该定理由Appel和Haken于1976年首次证明，借助于计算机验证约1936种不可避免构型的可约性。1996年Robertson、Sanders、Seymour和Thomas给出了简化的计算机证明，仅需验证633种构型。2005年Gonthier用Coq形式化验证了该证明。

\textbf{证明概要}（Appel-Haken框架）证明分为两个核心部分：（1）\textbf{不可避免集}：证明每个极大平面图中必然出现约1936种''构型''（局部子图模式）中的至少一种。这通过''放电法''（discharging method）完成------给每个顶点和面赋予初始电荷（由Euler公式约束），然后按规则重新分配电荷，最终推导出存在度数不超过5的顶点或特定构型。（2）\textbf{可约性}：对每种构型，证明若极小反例包含此构型，则可将其''约化''（通过收缩或删除操作得到更小的平面图），对后者使用归纳假设染色后再扩展回原图的合法4-染色。可约性证明需借助Kempe链颜色交换技术。每构型的可约性由计算机程序验证（每种涉及大量颜色配置的枚举）。这两个组件结合：因极小反例必然包含某种不可避免构型，而每种构型都是可约的，故不存在极小反例，定理得证。\(\blacksquare\)

\textbf{定理 54.5（五色定理 - 无需计算机的证明）} 每个平面图是 \(5\)-可染色的。

\textbf{证明} 对顶点数归纳。\(n \leq 5\) 时显然。对一般的 \(n\)，由Euler公式知存在度 \(\leq 5\) 的顶点 \(v\)。若 \(\deg(v) \leq 4\)，归纳假设保证 \(G - v\) 可 \(5\)-染色，\(v\) 至多有 \(4\) 个邻居，必有余色。若 \(\deg(v) = 5\)，设其邻居为 \(v_1, \ldots, v_5\)（按平面嵌入的循环序）。若 \(v_1, \ldots, v_5\) 中某两顶点同色则完成。否则，不妨设 \(c(v_i) = i\)，考虑由颜色 \(1\) 和 \(3\) 诱导的子图 \(G_{13}\) 中 \(v_1\) 和 \(v_3\) 所在的连通分支。若 \(v_1\) 和 \(v_3\) 在不同分支中，交换其中一个分支的颜色 \(1\) 和 \(3\) 即可释放颜色给 \(v\)。若在同一分支，则在 \(v\) 附近形成 \((1,3)\)-路径，阻挡了 \(v_2\) 和 \(v_4\) 之间的 \((2,4)\)-路径，故可对 \((2,4)\)-子图进行颜色交换。\(\blacksquare\)

\textbf{证明}：对平面图的顶点数 \(n\) 归纳。\(n \leq 5\) 时平凡（\(5\) 色足够）。设 \(n > 5\) 且结论对所有平面图成立。由Euler公式，每个平面图存在度数 \(\leq 5\) 的顶点 \(v\)。删去 \(v\) 得图 \(G'\)。由归纳假设，\(G'\) 可 \(5\)-染色。若 \(v\) 的邻居使用少于 \(5\) 种颜色，则 \(v\) 可染未用之色。若 \(v\) 的恰好 \(5\) 个邻居使用了所有 \(5\) 种颜色，设为 \(v_1, \ldots, v_5\)（顺时针排列，颜色分别为 \(c_1, \ldots, c_5\)）。考虑 \(G'\) 中仅用颜色 \(c_1\) 和 \(c_3\) 的顶点生成的子图。若 \(v_1\) 和 \(v_3\) 属于该子图的同一连通分支，则存在连接它们的路径，该路径与边 \(v v_1, v v_3\) 构成Jordan闭曲线，迫使 \(v_2\) 和 \(v_4\) 或 \(v_5\) 无法共存于同侧------与平面性矛盾。故 \(v_1\) 和 \(v_3\) 在不同分支中。交换颜色 \(c_1 \leftrightarrow c_3\) 在 \(v_1\) 所在分支中，释放颜色 \(c_1\) 给 \(v\) 使用。因此 \(v\) 可被染色，\(G\) 是 \(5\)-可染的。\(\blacksquare\)

\subsection{边染色与Vizing定理}\label{ux8fb9ux67d3ux8272ux4e0evizingux5b9aux7406}

\textbf{定义 54.6（边染色）} 图 \(G\) 的\textbf{正常边染色}是一个函数 \(c': E(G) \to \{1, 2, \ldots, k\}\)，使得有公共端点的边获不同颜色。\(G\) 的\textbf{边色数} \(\chi'(G)\) 是所需的最少颜色数。显然 \(\chi'(G) \geq \Delta(G)\)。

\textbf{定理 54.7（Vizing定理）} 对任意简单图 \(G\)， \[\Delta(G) \leq \chi'(G) \leq \Delta(G) + 1.\]

\textbf{证明框架} 对边数归纳。假设对少于 \(m\) 条边的图定理成立。取一边 \(e = xy\)，由归纳假设 \(G - e\) 可用 \(\Delta + 1\) 种颜色染边。在 \(G - e\) 的染色中，对每个顶点 \(v\) 考虑其未使用的颜色集合。需要证明可将 \(e\) 染上一种颜色。若 \(x\) 和 \(y\) 在某颜色 \(c\) 上均缺失，则直接给 \(e\) 染 \(c\)。否则，利用Kempe链（交替颜色路径）重新排列颜色，总是可以空出一种颜色用于 \(e\)。Vizing的``扇''构造法证明了：要么 \(\Delta\) 种颜色足够（图是\textbf{第一类}的），要么需要 \(\Delta + 1\) 种（图是\textbf{第二类}的，如奇圈和完全图 \(K_{2n+1}\)）。二分图的边色数恒为 \(\Delta\)（Kőnig定理）。\(\blacksquare\)

\textbf{证明}：对边数归纳。基始：空图，\(\chi' = 0 \leq \Delta\)。归纳步骤：设 \(G\) 有 \(m\) 条边，\(e = uv\) 为某条边。删去 \(e\) 得 \(G'\)，由归纳假设 \(G'\) 可 \((\Delta+1)\)-边染色。对每个顶点 \(x\)，令 \(M(x)\) 为 \(G'\) 的染色中 \(x\) 处缺失的颜色集（即未用于 \(x\) 关联边的颜色）。由于至多 \(\Delta\) 条边关联于 \(x\)，\(|M(x)| \geq 1\)。

若 \(M(u) \cap M(v) \neq \emptyset\)，取公共缺失颜色染 \(e\) 即得 \(G\) 的 \((\Delta+1)\)-边染色。否则，取 \(\alpha \in M(u)\)，\(\beta \in M(v)\)（\(\alpha \neq \beta\)）。考虑 \(G'\) 中由边染色 \(\alpha\) 和 \(\beta\) 交替构成的极大路径 \(P\) 以 \(v\) 为起点（称为Kempe链）。若 \(P\) 不经过 \(u\)，交换 \(P\) 上 \(\alpha \leftrightarrow \beta\) 后，\(\beta\) 在 \(v\) 处缺失，\(\alpha\) 在 \(u\) 处缺失，\(e\) 可染 \(\alpha\)。若 \(P\) 经过 \(u\)，则 \(P\) 在 \(u\) 处以 \(\beta\) 边结束（因 \(\beta \notin M(u)\)）；\(P\) 加上 \(e\) 构成奇圈，通过适当的重新染色（沿路径置换）可释放 \(u\) 和 \(v\) 的共同缺失颜色。\(\blacksquare\)

\subsection{平面图与Euler公式}\label{ux5e73ux9762ux56feux4e0eeulerux516cux5f0f}

\textbf{定义 54.8（平面图）} 图 \(G\) 称为\textbf{平面图}（planar graph），若它可在平面上嵌入使得边仅在端点处相交。此嵌入将平面划分为连通区域，称为\textbf{面}（face）。连通平面图 \(G\) 的平面嵌入（记为 \(\tilde{G}\)）由顶点集 \(V\)、边集 \(E\) 和面集 \(F\) 构成。

\textbf{定理 54.9（Euler公式）} 设 \(\tilde{G}\) 是连通平面图的平面嵌入。则 \[|V| - |E| + |F| = 2.\]

\textbf{证明} 对 \(|E|\) 归纳。若 \(|E| = 0\)，则 \(|V| = 1\)，\(|F| = 1\)，\(1 - 0 + 1 = 2\) 成立。若 \(|E| > 0\)，若 \(G\) 中有圈，删除圈上的一条边减少边数和面数各 \(1\)，\(|V|\) 不变，由归纳假设得证。若 \(G\) 是树，则 \(|F| = 1\)，\(|V| - |E| = 1\)，故 \(|V| - |E| + |F| = 1 + 1 = 2\)。\(\blacksquare\)

\textbf{推论 54.10} 对 \(n \geq 3\) 个顶点的连通简单平面图，\(|E| \leq 3n - 6\)。若图不含三角形，则 \(|E| \leq 2n - 4\)。

\textbf{定理 54.11（Kuratowski定理）} 图 \(G\) 是平面图当且仅当 \(G\) 不包含 \(K_5\) 或 \(K_{3,3}\) 的细分作为子图。

\textbf{证明概要} 必要性：\(K_5\) 和 \(K_{3,3}\) 均非平面图------\(K_5\) 有 \(5\) 个顶点和 \(10\) 条边，违反 \(|E| \leq 3 \cdot 5 - 6 = 9\)；\(K_{3,3}\) 是二分图无三角形，有 \(6\) 个顶点和 \(9\) 条边，违反 \(|E| \leq 2 \cdot 6 - 4 = 8\)。充分性可通过反证法及极小反例的分析证明（或利用Tutte关于\(3\)-连通图的轮图结构定理）。\(\blacksquare\)

\subsection{四色定理}\label{ux56dbux8272ux5b9aux7406}

\textbf{定理 54.12（四色定理 - 等价形式）} 以下陈述等价：（i）每个平面图是 \(4\)-可染色的；（ii）每个无桥三次平面图是 \(3\)-边可染色的（Tait等价）；（iii）每个三正则无桥平面图的面可以 \(4\)-染色使得相邻面不同色。

\textbf{证明} (i)\(\Rightarrow\)(ii)：设 \(G\) 是无桥三次平面图。由四色定理，\(G\) 的面可4-染色，颜色取自Klein四元群 \(\mathbb{Z}_2 \times \mathbb{Z}_2 = \{0, a, b, a+b\}\)。对每条边 \(e\) 分隔面 \(f_1\) 和 \(f_2\)，定义 \(e\) 的颜色为 \(c(f_1) + c(f_2)\)（群加法）。因 \(G\) 无桥，每条边分隔两个不同面。在三次图中，交于同一顶点的三条边分隔的面对给出三个两两不同的非零群元------若两条边获得相同颜色则对应面染色矛盾。故得正常的3-边染色。

(ii)\(\Rightarrow\)(iii)：设 \(G\) 是无桥三次平面图且有正常3-边染色。固定参考面 \(f_0\) 染0，对任意面 \(f\) 沿从 \(f_0\) 到 \(f\) 的路径累积越过边的颜色（在 \(\mathbb{Z}_2 \times \mathbb{Z}_2\) 中），得 \(f\) 的染色。因边染色正常，闭合路径累积色为零，故面染色良定义。相邻面色差非零故不同。

(iii)\(\Rightarrow\)(i)：由平面对偶性，面染色对应三角剖分的点染色，而任意平面图可嵌入三角剖分，故等价于(i)。

\textbf{历史} 1852年Guthrie提出四色猜想。Appel和Haken的策略：（1）取极小反例；（2）必为三角剖分；（3）用放电法找出不可避免集（约1936种构型）；（4）证明每种构型可约。由计算机验证完成证明。\(\blacksquare\)

\subsection{完美图与强完美图定理}\label{ux5b8cux7f8eux56feux4e0eux5f3aux5b8cux7f8eux56feux5b9aux7406}

\textbf{定义 54.13（完美图）} 图 \(G\) 称为\textbf{完美图}（perfect graph），若对 \(G\) 的每个导出子图 \(H\)，有 \(\chi(H) = \omega(H)\)。

由Lovász证明的弱完美图定理：\(G\) 是完美图当且仅当其补图 \(\overline{G}\) 是完美图。

\textbf{定理 54.14（强完美图定理，Chudnovsky-Robertson-Seymour-Thomas, 2006）} 图 \(G\) 是完美图当且仅当 \(G\) 和 \(\overline{G}\) 均不包含长度至少为 \(5\) 的奇圈（奇洞，odd hole）或奇反洞（odd antihole）作为导出子图。

该定理被证明为Berge于1961年提出的强完美图猜想，是图论史上最重大的成果之一。证明长达约150页，发展了''结构分解''方法。

\textbf{证明概要} （CRST框架）必要性（\(\Rightarrow\)）：奇洞 \(C_{2k+1}\)（\(k \geq 2\)）不是完美图------\(\chi(C_{2k+1}) = 3\) 而 \(\omega(C_{2k+1}) = 2\)，且其补图（奇反洞）同样不完美。由于完美图在取导出子图下封闭，故完美图不能包含奇洞或奇反洞作为导出子图。

充分性（\(\Leftarrow\)）是核心挑战，证明分为三个主要阶段：（1）\textbf{分解定理}：任何Berge图（无奇洞且无奇反洞）要么属于几个''基本类''（二部图、二部图的补图、线图、线图的补图、双分裂图等），要么具有某种''可容许分解''（如团割集分解或2-join分解）。（2）\textbf{结构分析}：若 \(G\) 和 \(\overline{G}\) 均无奇洞，则 \(G\) 的结构可被精细控制，逐步分解为基本类。（3）\textbf{验证基本类}：每个基本类中的图通过归纳证明是完美图------利用完美图在取补图下封闭（Lovász弱完美图定理）。此三阶段结合即得 \(\chi(G) = \omega(G)\) 对所有导出子图成立。\(\blacksquare\)

\hrulefill

\newpage

\chapter{04-离散数学/01-组合数学/07-图论高级.md}\label{ux79bbux6563ux6570ux5b6601-ux7ec4ux5408ux6570ux5b6607-ux56feux8bbaux9ad8ux7ea7.md}

\section{第92章：代数图论与谱图论}\label{ux7b2c92ux7ae0ux4ee3ux6570ux56feux8bbaux4e0eux8c31ux56feux8bba}

\subsection{图的邻接矩阵谱}\label{ux56feux7684ux90bbux63a5ux77e9ux9635ux8c31}

\textbf{定义 55.1（图的谱）} 图 \(G\) 的\textbf{谱}是邻接矩阵 \(A(G)\) 的特征值（含重数）的集合，记为 \(\operatorname{Spec}(G) = \{\lambda_1 \geq \lambda_2 \geq \cdots \geq \lambda_n\}\)。

\textbf{命题 55.2} （1）\(\sum_{i=1}^{n} \lambda_i = \operatorname{tr}(A) = 0\)；（2）\(\sum_{i=1}^{n} \lambda_i^2 = \operatorname{tr}(A^2) = 2|E(G)|\)；（3）\(\sum_{i=1}^{n} \lambda_i^3 = 6 \times (\text{三角形个数})\)。

\textbf{定理 55.3（Perron-Frobenius定理在图论中的应用）} 设 \(G\) 是连通图，\(A\) 是其邻接矩阵。则：

（1）谱半径 \(\rho(A) = \lambda_1\) 是单重正特征值；

（2）存在严格正的特征向量 \(\mathbf{x} > 0\) 满足 \(A\mathbf{x} = \lambda_1 \mathbf{x}\)（Perron向量）；

（3）对任意其他特征值 \(\lambda\)，\(|\lambda| \leq \lambda_1\)；

（4）\(\lambda_1\) 满足 \(\delta(G) \leq \lambda_1 \leq \Delta(G)\)，其中 \(\delta\) 和 \(\Delta\) 分别是最小度和最大度。

\textbf{证明} \(A\) 是对称非负不可约矩阵（连通性保证不可约性）。由Perron-Frobenius定理，存在正特征值 \(\lambda_1\) 及对应的正特征向量。对称性确保 \(\lambda_1\) 是实特征值且为单根。由Rayleigh商：\(\lambda_1 = \max_{\mathbf{x} \neq 0} \frac{\mathbf{x}^T A \mathbf{x}}{\mathbf{x}^T \mathbf{x}}\)。取 \(\mathbf{x} = (1, 1, \ldots, 1)^T\) 得上下界。\(\blacksquare\)

\textbf{定义 55.4（图的能量）} 图 \(G\) 的\textbf{能量}定义为 \(E(G) = \sum_{i=1}^{n} |\lambda_i|\)。

\subsection{图谱与图的扩张性质}\label{ux56feux8c31ux4e0eux56feux7684ux6269ux5f20ux6027ux8d28}

\textbf{定义 55.5（扩张常数）} 图 \(G = (V, E)\) 的\textbf{扩张常数}（等周常数或Cheeger常数）为 \[h(G) = \min_{\substack{S \subseteq V \\ 0 < |S| \leq |V|/2}} \frac{|\partial S|}{|S|},\] 其中 \(\partial S = \{\{u, v\} \in E : u \in S, v \notin S\}\) 是 \(S\) 的边边界。

\textbf{定理 55.6（Cheeger不等式 - 图的版本）} 设 \(\lambda_2\) 是 \(G\) 的Laplace矩阵 \(L\) 的次小特征值（代数连通度）。则 \[\frac{\lambda_2}{2} \leq h(G) \leq \sqrt{2 \lambda_2 \Delta},\] 其中 \(\Delta\) 是最大度。

\textbf{证明} （下界）设 \(S\) 是达到Cheeger常数 \(h = h(G)\) 的子集，即 \(0 < |S| \leq |V|/2\) 且 \(h = |\partial S|/|S|\)。构造 \(\mathbf{x} \in \mathbb{R}^n\)：对 \(v \in S\) 令 \(x_v = 1/|S| - 1/|\bar{S}|\)，对 \(v \notin S\) 令 \(x_v = -1/|\bar{S}|\)。可验证 \(\mathbf{x} \perp \mathbf{1}\)。由 Rayleigh 商， \[\lambda_2 \leq \frac{\mathbf{x}^T L \mathbf{x}}{\mathbf{x}^T \mathbf{x}} = \frac{\sum_{\{u,v\} \in E} (x_u - x_v)^2}{\sum_v x_v^2}.\] 计算得分子 \(= |\partial S| \cdot (1/|S| + 1/|\bar{S}|)^2\)，分母 \(= 1/|S| + 1/|\bar{S}|\)。因 \(|\bar{S}| \geq |S|\)，得 \(\lambda_2 \leq 2h\)，即 \(h \geq \lambda_2/2\)。

（上界）取 \(\lambda_2\) 对应的特征向量 \(\mathbf{f} \perp \mathbf{1}\)（Fiedler向量）。排序顶点使 \(f_1 \leq \cdots \leq f_n\)。令 \(S_t = \{v : f_v \leq t\}\)，通过积分与 Cauchy-Schwarz 不等式可证存在 \(t\) 使 \(|\partial S_t|/|S_t| \leq \sqrt{2\lambda_2\Delta}\)。关键估计：\(\sum_{\{u,v\}\in E} |f_u^2 - f_v^2| \leq \sqrt{\sum (f_u-f_v)^2 \cdot \sum (f_u+f_v)^2} \leq \sqrt{\lambda_2\|\mathbf{f}\|^2 \cdot 2\Delta \sum f_v^2}\)，由此导出上界。

此不等式建立了图的光谱性质与组合扩张性质之间的基本联系。\(\blacksquare\)

\textbf{定理 55.7（Alon-Boppana定理）} 设 \(G_n\) 是一族 \(d\)-正则图且 \(|V(G_n)| \to \infty\)。则 \[\liminf_{n \to \infty} \lambda_2(G_n) \geq 2\sqrt{d-1},\] 其中 \(\lambda_2\) 在此处指邻接矩阵的次大特征值。

\textbf{证明} 设 \(G\) 是 \(d\)-正则图，\(A\) 为其邻接矩阵。对任意偶数 \(2k\)，考虑 \(A^{2k}\) 的迹： \[\operatorname{tr}(A^{2k}) = \sum_{i=1}^{n} \lambda_i^{2k} \geq \lambda_1^{2k} + \lambda_2^{2k}.\] 另一方面，\(\operatorname{tr}(A^{2k})\) 等于 \(G\) 中长度为 \(2k\) 的闭游走数。\(d\)-正则无穷树（Bethe格）中从根出发的长度为 \(2k\) 的闭游走数满足递推关系，其生成函数主导项给出渐近增长率 \(\lambda_2 \geq 2\sqrt{d-1} - o(1)\)。具体地，对任意 \(\varepsilon > 0\)，当 \(n\) 足够大时 \(G_n\) 的半径大于某常数，局部近似于 \(d\)-正则树，从而游走计数下界推得 \(\lambda_2(G_n) \geq 2\sqrt{d-1} - \varepsilon\)。这给出了 \(d\)-正则扩展图的谱极限------\(2\sqrt{d-1}\) 是最优可能值，Ramanujan图恰达到此界。\(\blacksquare\)

\subsection{图谱与随机游走}\label{ux56feux8c31ux4e0eux968fux673aux6e38ux8d70}

\textbf{定义 55.8（转移矩阵与随机游走）} 图 \(G\) 上的\textbf{简单随机游走}由转移矩阵 \(P = D^{-1}A\) 定义，其中 \(P_{ij}\) 是从顶点 \(i\) 走到顶点 \(j\) 的一步概率（若 \(\{i, j\} \in E\)，\(P_{ij} = 1/\deg(i)\)，否则 \(0\)）。

\(P\) 相似于对称矩阵 \(\widetilde{P} = D^{-1/2} A D^{-1/2}\)，故 \(P\) 可对角化且有实特征值 \(1 = \mu_1 \geq \mu_2 \geq \cdots \geq \mu_n \geq -1\)。

\textbf{定义 55.9（混合时间）} 随机游走的\textbf{混合时间} \(t_{\text{mix}}(\varepsilon)\) 是使分布 \(p^{(t)}\) 与平稳分布 \(\pi\) 的总变差距离小于 \(\varepsilon\) 的最小步数。

\textbf{定理 55.10（谱隙与混合时间）} 设 \(\lambda^* = \max\{|\mu_2|, |\mu_n|\}\) 是 \(P\) 的非平凡特征值绝对值的最大值，\(\delta = \min_v \deg(v)\)，\(n = |V|\)。则 \[\frac{\lambda^*}{2(1-\lambda^*)} \log\frac{1}{2\varepsilon} \leq t_{\text{mix}}(\varepsilon) \leq \frac{1}{1-\lambda^*} \log\frac{n}{\varepsilon \delta}.\]

谱隙 \(1 - \lambda^*\) 越大，混合速度越快。

\textbf{证明} （上界）设 \(p^{(t)}\) 是 \(t\) 步后的分布，\(\pi\) 是平稳分布。将 \(p^{(0)} - \pi\) 在 \(P\) 的特征基下展开：\(p^{(0)} - \pi = \sum_{j=2}^{n} c_j \mathbf{v}_j\)，其中 \(\mathbf{v}_j\) 是 \(P\) 的特征向量（对应特征值 \(\mu_j\)）。则 \(p^{(t)} - \pi = (p^{(0)} - \pi) P^t = \sum_{j=2}^{n} c_j \mu_j^t \mathbf{v}_j\)。利用 \(\|p^{(t)} - \pi\|_{\text{TV}} = \frac{1}{2}\|p^{(t)} - \pi\|_1 \leq \frac{\sqrt{n}}{2} \|p^{(t)} - \pi\|_2\) 及 \(\|p^{(t)} - \pi\|_2 \leq (\lambda^*)^t \|p^{(0)} - \pi\|_2\)，取 \(p^{(0)}\) 为在单个顶点上的delta分布，得 \(\|p^{(t)} - \pi\|_{\text{TV}} \leq \frac{\sqrt{n}}{2\sqrt{\delta}} (\lambda^*)^t\)。令右端 \(\leq \varepsilon\) 解出 \(t\) 得上界。

（下界）设 \(\mathbf{v}\) 是对应于特征值 \(\mu\)（\(|\mu| = \lambda^*\)）的特征向量，\(\mathbf{v} \perp \mathbf{1}\)。取初始分布 \(p^{(0)} = \pi + \alpha \mathbf{v}\)（适当缩放 \(\alpha\) 确保 \(p^{(0)} \geq 0\)），则 \(p^{(t)} - \pi = \alpha \mu^t \mathbf{v}\)。利用 \(\|p^{(t)} - \pi\|_{\text{TV}} \geq \frac{1}{2}(\lambda^*)^t \|\alpha\mathbf{v}\|_1\)，令其 \(\geq \varepsilon\) 取下界。\(\blacksquare\)

\subsection{图的同构与自同构群}\label{ux56feux7684ux540cux6784ux4e0eux81eaux540cux6784ux7fa4}

\textbf{定义 55.11（自同构群）} 图 \(G\) 的\textbf{自同构群} \(\operatorname{Aut}(G)\) 是 \(G\) 到自身的所有同构映射构成的群，运算为复合。

\textbf{定理 55.12（Frucht定理, 1939）} 对任意有限群 \(\Gamma\)，存在图 \(G\) 使得 \(\operatorname{Aut}(G) \cong \Gamma\)。换言之，每个有限群均可实现为某图的自同构群。

\textbf{证明概要} 利用群 \(\Gamma\) 的Cayley图构造。取 \(\Gamma\) 的生成集 \(S\)（\(S = S^{-1}\)，\(1 \notin S\)），Cayley图 \(\operatorname{Cay}(\Gamma, S)\) 的自同构群包含 \(\Gamma\)（通过左乘作用）。为使自同构群恰等于 \(\Gamma\)，需通过``装饰''顶点以破坏额外的对称性：对有向Cayley图，用不对称的``小工具''替换每条有向边------例如，用长度为 \(i\) 的路径连接特定的``刚性子图''（rigid subgraph），使得只有 \(\Gamma\) 中的左乘保持该结构。Sabidussi（1957）给出了简化的构造。\(\blacksquare\)

\hrulefill

\section{第93章：随机图与极值图论}\label{ux7b2c93ux7ae0ux968fux673aux56feux4e0eux6781ux503cux56feux8bba}

\subsection{Erdős--Rényi随机图模型}\label{erdux151sruxe9nyiux968fux673aux56feux6a21ux578b}

\textbf{定义 56.1（Erdős--Rényi模型）} 随机图 \(G(n, p)\) 定义为 \(n\) 个标号顶点的随机图，其中每对顶点独立地以概率 \(p\) 连接（\(0 \leq p \leq 1\)）。等价模型 \(G(n, M)\) 从 \(\binom{n}{2}\) 条可能边中均匀随机选取 \(M\) 条。

\textbf{定义 56.2（图性质的阈值）} 图性质 \(\mathcal{P}\) 称为\textbf{单调递增}的，若 \(G\) 具性质 \(\mathcal{P}\) 且 \(G \subseteq H\) 则 \(H\) 也具性质 \(\mathcal{P}\)。函数 \(p_0(n)\) 称为 \(\mathcal{P}\) 的\textbf{阈值}，若当 \(p \ll p_0\) 时 \(\mathbb{P}(G(n, p) \in \mathcal{P}) \to 0\)，而当 \(p \gg p_0\) 时概率 \(\to 1\)。

\textbf{定理 56.3（巨分支的出现）} 对于 \(G(n, c/n)\)，当 \(c > 1\) 时，以高概率存在唯一的大小为 \(\Theta(n)\) 的连通分支（\textbf{巨分支}），所有其他分支的大小均为 \(O(\log n)\)。当 \(c < 1\) 时，以高概率所有分支的大小均为 \(O(\log n)\)。在临界值 \(c = 1\)，最大分支的大小量为 \(\Theta(n^{2/3})\)。

\textbf{证明概要} 考虑分支探索过程（分支过程近似）。从某个顶点开始，其邻居数近似服从 \(\operatorname{Bin}(n-1, c/n) \approx \operatorname{Poisson}(c)\)。当 \(c > 1\) 时，分支过程以正概率无限持续，对应的图论分支恰为巨分支。利用二阶矩方法可证唯一性：设 \(X_k\) 为大小为 \(k\) 的连通分支数，\(\mathbb{E}[X_k] \approx \binom{n}{k} k^{k-2} p^{k-1} (1-p)^{k(n-k)}\)。对于 \(c < 1\)，\(\sum_{k \geq K \log n} \mathbb{E}[X_k] \to 0\)，故所有分支均为对数大小。\(\blacksquare\)

\textbf{定理 56.4（连通性阈值）} \(G(n, p)\) 的连通性阈值为 \(p = \frac{\log n}{n}\)。更精确地，若 \(p = \frac{\log n + c_n}{n}\)，则 \(\lim_{n \to \infty} \mathbb{P}(G(n, p) \text{ 连通}) = e^{-e^{-c}}\)（其中 \(c = \lim c_n\)）。

\textbf{证明} 图 \(G(n, p)\) 不连通当且仅当存在孤立顶点（为决定性的障碍）。令 \(X\) 为孤立顶点数。对每个顶点 \(v\)，\(\mathbb{P}(v \text{ 孤立}) = (1-p)^{n-1}\)。\(\mathbb{E}[X] = n(1-p)^{n-1}\)。当 \(p = \frac{\log n + c_n}{n}\) 时，\((1-p)^{n-1} = \exp((n-1)\ln(1-p)) = \exp(-(n-1)(p + O(p^2))) = e^{-c_n - \log n + o(1)} = \frac{e^{-c_n}}{n} (1+o(1))\)。故 \(\mathbb{E}[X] \to e^{-c}\)。利用矩方法（Poisson逼近），\(X\) 渐近服从 \(\operatorname{Poisson}(e^{-c})\)，故 \(\mathbb{P}(X=0) \to e^{-e^{-c}}\)。还需证当无孤立顶点时图以高概率连通：设 \(p = \frac{\log n + c_n}{n}\) 且 \(c_n \to \infty\)，则 \(\mathbb{E}[X] \to 0\)，图以高概率连通；当 \(c_n \to -\infty\)，\(\mathbb{E}[X] \to \infty\)，图以高概率不连通。分离小连通分支的概率可由 \(\mathbb{P}(\text{存在大小为 } k \text{ 的分支}) \leq \binom{n}{k} k^{k-2} p^{k-1} (1-p)^{k(n-k)}\) 控制，当 \(k \geq 2\) 时此概率可忽略。\(\blacksquare\)

\subsection{极值图论核心定理}\label{ux6781ux503cux56feux8bbaux6838ux5fc3ux5b9aux7406}

\textbf{定义 56.5（Turán图）} \textbf{Turán图} \(T(n, r)\) 是 \(n\) 个顶点的完全 \(r\)-部图，各部分的大小尽可能相等（即每部分大小差不超过 \(1\)）。记 \(t_r(n)\) 为 \(T(n, r)\) 的边数。

\textbf{定理 56.6（Turán定理, 1941）} 任意在 \(n\) 个顶点上不含 \(K_{r+1}\) 作为子图的图 \(G\)，其边数至多为 \(t_r(n)\)，且等号成立当且仅当 \(G \cong T(n, r)\)。

\textbf{证明一（归纳法）} 对 \(n\) 归纳。基始 \(n \leq r\) 平凡。对 \(n > r\)，取度最大的顶点 \(v\)，\(\deg(v) = \Delta\)。\(G[N(v)]\) 不含 \(K_r\)（否则与 \(v\) 构成 \(K_{r+1}\)），故 \(e(N(v)) \leq t_{r-1}(\Delta)\)。\(V \setminus (\{v\} \cup N(v))\) 中的边数不超过 \(t_r(n - \Delta - 1)\)。通过优化和归纳得边数最大值在 \(v\) 的度恰为 \(\lfloor (n(r-1))/r \rfloor\) 时达到，即Turán图。\(\blacksquare\)

\textbf{证明二（Zarankiewicz-对称差法）} 设 \(G\) 是不含 \(K_{r+1}\) 的极大图（即添加任意边均产生 \(K_{r+1}\)）。定义关系 \(\sim\)：\(x \sim y \iff \{x, y\} \notin E(G)\)。可证 \(\sim\) 是等价关系，故 \(G\) 是完全多部图。由极大性知恰有 \(r\) 个部分，边数最大化时各部分大小尽可能相等。\(\blacksquare\)

\textbf{证明三（概率方法 - Motzkin-Straus）} 注意 \(\max_{G \text{ 无 } K_{r+1}} e(G)\) 等价于二次型 \(\sum_{\{i,j\} \in E} x_i x_j\) 在 \(\sum x_i = 1, x_i \geq 0\) 下的最大值。Lagrange乘子法给出最优解为 \(x_i = 1/r\)（均匀分布），对应边数为 \((1 - 1/r) n^2 / 2\)，即Turán图的渐近界。\(\blacksquare\)

\textbf{证明}（Turán, 1941）：设 \(G = (V, E)\) 是不含 \(K_{r+1}\) 的 \(n\) 顶点图。选取 \(G\) 中最大度的顶点 \(v_1\)。令 \(S_1 = N(v_1)\) 为 \(v_1\) 的邻居集。在 \(V \setminus (\{v_1\} \cup S_1)\) 中选取最大度的顶点 \(v_2\)，令 \(S_2 = N(v_2) \setminus (\{v_1\} \cup S_1)\)。重复此过程 \(r\) 次，得不相交的集合 \(S_1, \ldots, S_r\)。由于 \(G\) 不含 \(K_{r+1}\)，这些 \(S_i\) 之间无边相连。由构造，\(|E| \leq \sum_{i=1}^r |S_i| \cdot (n - |S_1 \cup \cdots \cup S_i|)\)。最大化此界得 \(|S_i| \approx n/r\)，\(|E| \leq (1 - 1/r)n^2/2\)。等号成立当且仅当各部分大小尽量均匀（差至多 \(1\)），即Turán图 \(T(n, r)\)。\(\blacksquare\)

\textbf{定理 56.7（Erdős-Stone定理）} 对任意固定的图 \(H\)， \[\operatorname{ex}(n, H) = \left(1 - \frac{1}{\chi(H) - 1} + o(1)\right) \binom{n}{2},\] 其中 \(\operatorname{ex}(n, H)\) 是 \(n\) 个顶点上不含 \(H\) 作为子图的最大边数，\(\chi(H)\) 是 \(H\) 的色数。

\textbf{证明} 下界：取 Turán 图 \(T(n, \chi(H)-1)\)（完全 \((\chi(H)-1)\)-部图），其色数为 \(\chi(H)-1\)，故不含 \(H\)（因 \(\chi(H) > \chi(H)-1\)）。其边数为 \((1 - 1/(\chi(H)-1) + o(1))\binom{n}{2}\)，即下界。

上界：设 \(G\) 有 \((1 - 1/(\chi(H)-1) + \varepsilon)\binom{n}{2}\) 条边，\(\varepsilon > 0\) 固定。由 Szemerédi 正则引理，将 \(G\) 划分为 \(\varepsilon\)-正则对。由边数条件，存在 \((\chi(H)-1)\)-部子图密度为正，应用计数引理（定理 56.10）可找到 \(H\) 的副本（因 \(H\) 是 \(\chi(H)\)-可染的，可嵌入到足够稠密的 \((\chi(H)-1)\)-部正则结构中，利用关键引理：若正则对各部密度为正，则包含所有固定大小的 \(\chi(H)\)-部子图）。故当 \(n\) 足够大时 \(G\) 必含 \(H\)，即 \(\operatorname{ex}(n, H) \leq (1 - 1/(\chi(H)-1) + o(1))\binom{n}{2}\)。\(\blacksquare\)

当 \(\chi(H) \geq 3\) 时，此定理将渐近极值问题归约为Turán定理。当 \(\chi(H) = 2\) 时，该定理给出 \(o(n^2)\) 的界，需要更精细的分析（如二分图的 Zarankiewicz 问题）。

\subsection{Szemerédi正则引理}\label{szemeruxe9diux6b63ux5219ux5f15ux7406}

\textbf{定义 56.8（\(\varepsilon\)-正则对）} 设 \(G\) 是图，\(X, Y \subseteq V(G)\) 不相交。\((X, Y)\) 称为\textbf{\(\varepsilon\)-正则的}，若对任意 \(A \subseteq X, B \subseteq Y\) 满足 \(|A| \geq \varepsilon |X|, |B| \geq \varepsilon |Y|\)，有 \[|d(A, B) - d(X, Y)| \leq \varepsilon,\] 其中 \(d(A, B) = e(A, B) / (|A| |B|)\) 是 \(A\) 和 \(B\) 之间的边密度。

\textbf{定理 56.9（Szemerédi正则引理, 1978）} 对任意 \(\varepsilon > 0\) 和整数 \(m_0\)，存在 \(M = M(\varepsilon, m_0)\) 使得任意足够大的图 \(G\) 的顶点集可划分 \(V = V_0 \cup V_1 \cup \cdots \cup V_k\)（\(m_0 \leq k \leq M\)），满足：（1）\(|V_0| \leq \varepsilon n\)；（2）\(|V_1| = \cdots = |V_k|\)；（3）除至多 \(\varepsilon \binom{k}{2}\) 对外，所有的对 \((V_i, V_j)\) 都是 \(\varepsilon\)-正则的。

\textbf{证明框架} 使用能量增量论证。定义划分的``指数''（index）为 \(\sum_{i < j} |V_i| |V_j| d(V_i, V_j)^2 / n^2\)。若当前划分不满足正则条件，可通过精炼划分将指数提升至少 \(\varepsilon^5/2\)。指数以 \(1\) 为上界，故过程至多迭代 \(2/\varepsilon^5\) 次。每次迭代划分大小至多增加 \(2^k\) 倍，故最终 \(k\) 存在有限上界 \(M\)（塔式上界 \(\operatorname{tower}(O(1/\varepsilon^5))\)）。\(\blacksquare\)

\textbf{定理 56.10（三角形计数引理）} 设 \(V_1, V_2, V_3\) 是两两 \(\varepsilon\)-正则的且密度分别为 \(d_{12}, d_{23}, d_{31}\)。则 \(V_1, V_2, V_3\) 构成的三角形个数为 \((d_{12} d_{23} d_{31} \pm 3\varepsilon) |V_1| |V_2| |V_3|\)。

\textbf{证明} 先统计``典型''顶点对。由 \((V_1, V_2)\) 的 \(\varepsilon\)-正则性，除至多 \(\varepsilon|V_1|\) 个例外顶点外，每个 \(v_1 \in V_1\) 在 \(V_2\) 中的邻居数约为 \(d_{12}|V_2|\)（误差 \(\pm \varepsilon|V_2|\)）。对 \(V_1\) 的非例外顶点 \(v_1\)，其 \(V_2\) 邻居集 \(N_2(v_1)\) 大小为 \((d_{12} \pm \varepsilon)|V_2|\)。任取 \(v_2 \in N_2(v_1)\)，由 \((V_2, V_3)\) 的正则性，\(v_2\) 在 \(V_3\) 中典型邻居数约 \(d_{23}|V_3|\)。同时 \((V_3, V_1)\) 的正则性给出 \(v_1\) 与 \(V_3\) 的连接密度 \(d_{31}\)。综合估计，以 \(v_1\) 为一顶点的三角形数约为 \(d_{12} d_{23} d_{31} |V_2| |V_3|\)。求和得总数 \((d_{12} d_{23} d_{31} \pm 3\varepsilon)|V_1||V_2||V_3|\)。

正则引理将稠密图的结构简化为拟随机二分图的组合，是极值组合学中最重要的结构定理之一。\(\blacksquare\)

\subsection{图极限理论}\label{ux56feux6781ux9650ux7406ux8bba}

\textbf{定义 56.11（图子）} \textbf{图子}（graphon）是取值于 \([0, 1]\) 的对称可测函数 \(W: [0, 1]^2 \to [0, 1]\)。每个图 \(G\) 可通过其邻接矩阵的自然分片常数插值对应一个图子 \(W_G\)。

\textbf{定义 56.12（割范数）} 图子 \(W\) 的\textbf{割范数}定义为 \[\|W\|_{\square} = \sup_{S, T \subseteq [0,1]} \left| \int_{S \times T} W(x, y) \, dx \, dy \right|.\]

\textbf{定理 56.13（正则引理的解析形式 - Lovász-Szegedy）} 对任意图子序列 \(\{W_n\}\)，存在子列收敛于某个图子 \(W\)（在割范数意义下）。即图子空间在割度量下是紧的。

\textbf{证明概要} 该定理本质上是 Szemerédi 正则引理的解析翻译。对每个图子 \(W_n\)，由正则引理（应用于其离散化近似），存在步函数图子 \(U_n\)（分片常数，至多 \(M(\varepsilon)\) 个区间）使得 \(\|W_n - U_n\|_{\square} < \varepsilon\)。步函数由有限个参数确定（各块的密度值），所在的参数空间 \([0,1]^{M(\varepsilon)^2}\) 是紧致的（Tychonoff定理）。由对角线论证，存在子列 \(\{W_{n_k}\}\) 和步函数序列收敛于某 \(W\)。割范数的完备性保证 \(W\) 是图子，即得紧性。

Lovász和Szegedy（2006）发展了图极限的完整理论，将图序列的收敛性、极值图论和随机图统一在分析框架中。正则引理成为该理论中紧性证明的核心工具：每个图均可被步函数图子以任意精度逼近。该理论催生了图子估计（graphon estimation）和网络数据分析的新方向。\(\blacksquare\)

\hrulefill

图论是离散数学中历史最悠久的分支之一，其源头可追溯至 Euler 对 Königsberg 七桥问题的研究（1736 年）------这是数学史上被公认为第一篇图论论文的工作。从 19 世纪 Cayley 用树对化学分子的计数、Kirchhoff 用生成树分析电网，到 20 世纪 Tutte 的因子理论和 Erdős-Rényi 的随机图模型，图论经历了从具体问题驱动到统一理论体系的蜕变。图论的核心力量在于它将组合结构可视化，并通过递归（删边-缩边递推）、对偶（平面图与对偶图）、代数（邻接矩阵和 Laplace 矩阵的谱）和概率（随机图方法）四重方法深入分析。本章在组合学基础（V6 Ch 32 计数原理）之上建立图的形式化定义，涵盖树与连通性、平面图与 Euler 公式、Euler 回路与 Hamilton 回路，并引入有向图（竞赛图的 Landau 定理）和谱图论导引，为第 49 章的匹配与网络流和第 50 章的染色与 Ramsey 理论奠定图论基础。

\subsection{128.1 图的基本概念}\label{ux56feux7684ux57faux672cux6982ux5ff5-1}

\textbf{定义 128.1}（图）：一个\textbf{图} \(G = (V, E)\) 由顶点集 \(V\)（有限集）和边集 \(E \subseteq V \times V\) 组成。若无向图，边 \(\{u, v\}\) 是无序对；若有向图，边 \((u, v)\) 是有序对（弧）。

\textbf{定义 128.2}（基本概念）： - \textbf{度}：顶点 \(v\) 的度 \(\deg(v)\) 是与 \(v\) 关联的边数。在有向图中，出度 \(\deg^+(v)\) 和入度 \(\deg^-(v)\) 分别计数出边和入边。 - \textbf{握手引理}：\(\sum_{v \in V} \deg(v) = 2|E|\)。（每条边贡献两个度） - \textbf{路}：长度为 \(k\) 的路是顶点序列 \(v_0 v_1 \cdots v_k\)，其中 \(\{v_i, v_{i+1}\} \in E\)，且顶点互异。 - \textbf{回路}：起点和终点相同的路（\(v_0 = v_k\)，\(k \geq 3\)）。 - \textbf{连通}：图 \(G\) 是连通的，如果任意两个顶点之间存在路。 - \textbf{连通分支}：图的最大连通子图。

\textbf{定义 128.3}（特殊图类）： - \textbf{完全图} \(K_n\)：\(n\) 个顶点，每对顶点之间都有边（\(\binom{n}{2}\) 条边） - \textbf{二分图}：顶点集可划分为 \(V = A \cup B\)，所有边都在 \(A\) 和 \(B\) 之间 - \textbf{完全二分图} \(K_{m,n}\)：\(|A| = m, |B| = n\)，\(A\) 和 \(B\) 之间所有可能的边 - \textbf{路图} \(P_n\)：\(n\) 个顶点排成一条路 - \textbf{回路图} \(C_n\)：\(n\) 个顶点排成一个圈

\textbf{命题 128.1}（二分图的无奇圈特征）：图 \(G\) 是二分图当且仅当 \(G\) 不含奇长度回路。

\emph{证明}：若 \(G\) 是二分图，路在 \(A\) 和 \(B\) 之间交替，故回路长度必为偶数。反之，固定根顶点 \(v_0\)，对每个顶点 \(v\)，令 \(d(v)\) 为从 \(v_0\) 到 \(v\) 的最短路长度，按 \(d(v)\) 的奇偶性划分 \(V\)。无奇回路保证此划分是良定义的二分划分。∎

\subsection{128.2 树与生成树}\label{ux6811ux4e0eux751fux6210ux6811-1}

\textbf{定义 128.4}（树）：\textbf{树}是连通无回路的图。等价刻画：\(n\) 个顶点的树恰有 \(n-1\) 条边，且任意两点间存在唯一的路。

\textbf{定理 128.2}（Cayley 公式）：\(n\) 个标号顶点的不同树的个数为 \(n^{n-2}\)。

\textbf{证明}（Prufer 编码）：每个标号树唯一对应一个长度为 \(n-2\) 的序列，每个位置可从 \(1\) 到 \(n\) 取值。编码算法：反复删除度数为 \(1\) 的最小标号顶点，记录其邻居。解码算法：给定序列 \(a_1,\ldots,a_{n-2}\)，设 \(L\) 为 \(1,\ldots,n\) 中不在序列中出现的标号集合。对 \(i=1,\ldots,n-2\)，连接 \(a_i\) 与 \(L\) 中最小元素，删除 \(L\) 中该元素，若 \(a_i\) 不再出现于后序序列中则加入 \(L\)。最终 \(L\) 中剩余两元素相连。编码与解码互逆，故双射得证。因此树的总数为 \(n^{n-2}\)。\(\blacksquare\)

Cayley 公式通过 Prufer 编码这一漂亮的双射给出了标号树的精确计数。然而，当图不是树而是一般连通图时，计数所有生成树（即包含所有顶点的树子图）就变得更为复杂。Kirchhoff 的矩阵树定理（1847 年）为这一问题提供了代数解法：生成树个数等于 Laplace 矩阵的任意 \((n-1)\) 阶主子式，将组合计数问题转化为矩阵行列式的计算。

\textbf{定义 128.5}（生成树）：连通图 \(G\) 的\textbf{生成树}是 \(G\) 的包含所有顶点的树子图。

\textbf{定理 128.3}（Kirchhoff 矩阵树定理）：连通图 \(G\) 的生成树个数等于其 Laplace 矩阵 \(L = D - A\)（\(D\) 为度对角矩阵，\(A\) 为邻接矩阵）的任意一个 \((n-1)\) 阶主子式。

\textbf{证明}：对边数 \(m\) 归纳。\(m=n-1\) 时 \(G\) 为树，\(L\) 的 \(n-1\) 阶主子式为 \(1\)。归纳步：设结论对 \(m-1\) 边成立。对图 \(G\) 中边 \(e\)，删除 \(e\) 得 \(G-e\)，收缩 \(e\) 得 \(G/e\)。由矩阵运算性质 \(\det L_{G\setminus e}=\det L_G-\det L_{G/e}\)（沿 \(e\) 展开）。由归纳假设，\(\det L_{G\setminus e}\) 为 \(G\setminus e\) 生成树数，\(\det L_{G/e}\) 为 \(G/e\) 生成树数（即含 \(e\) 的生成树数）。二者和为 \(G\) 的生成树总数，即 \(\det L_G\)。\(\blacksquare\)

\textbf{定义 128.6}（最小生成树）：带权图 \(G\) 的\textbf{最小生成树}（MST）是权和最小的生成树。Kruskal 算法（贪心选择最小权边，不形成回路）和 Prim 算法均可找到 MST。

\textbf{定理 128.4}（Kruskal 算法的正确性）：对任意连通带权图，Kruskal 算法输出最小生成树。

\textbf{证明}：使用\textbf{割性质}（cut property）：对任意割 \((S, V\setminus S)\)，连接割两端的最小权边必属于某个最小生成树。证明：设 \(e\) 为连接 \(S\) 和 \(V\setminus S\) 的最小权边。若某 MST \(T\) 不含 \(e\)，则 \(T \cup \{e\}\) 含唯一回路，该回路必包含另一条割边 \(e'\)。由割性质 \(w(e) \leq w(e')\)，故 \(T' = T \cup \{e\} \setminus \{e'\}\) 的权不超过 \(T\)，\(T'\) 仍为生成树。因此 \(e\) 属于某 MST。Kruskal 算法每次选择不形成回路的最小权边，该边连接了两个不同连通分支，即为连接这两个分支的割中的最小权边，故属于某 MST。归纳可得算法输出的边集构成最小生成树。\(\blacksquare\)

\subsection{128.3 连通性与 Menger 定理}\label{ux8fdeux901aux6027ux4e0e-menger-ux5b9aux7406-2}

\textbf{定义 128.7}（连通度）： - \textbf{点连通度} \(\kappa(G)\)：使图不连通或退化为单点所需删除的最少顶点数 - \textbf{边连通度} \(\lambda(G)\)：使图不连通所需删除的最少边数 - \(G\) 是 \(k\)-连通的，如果 \(\kappa(G) \geq k\)

\textbf{命题 128.5}（Whitney 不等式）：\(\kappa(G) \leq \lambda(G) \leq \delta(G)\)，其中 \(\delta(G)\) 是最小度。

\textbf{证明}：\(\kappa(G) \leq \lambda(G)\)：设 \(\lambda(G) = k\)，存在 \(k\) 条边 \(\{e_1,\ldots,e_k\}\) 删除后图不连通。对每条边 \(e_i = \{u_i, v_i\}\)，若 \(u_i\) 和 \(v_i\) 均非孤立点，删除 \(u_i\)（或 \(v_i\)）即可断开该边。至多删除 \(k\) 个顶点即可断开所有 \(k\) 条边，故 \(\kappa(G) \leq k = \lambda(G)\)。\(\lambda(G) \leq \delta(G)\)：取最小度顶点 \(v\)（\(\deg(v) = \delta(G)\)），删除 \(v\) 的所有关联边即可断开 \(v\) 与图的其余部分，故 \(\lambda(G) \leq \delta(G)\)。\(\blacksquare\)

\textbf{定理 128.6}（Menger 定理，点版本）：图 \(G\) 中两个不相邻顶点 \(u, v\) 之间的内部不相交的 \(u\)-\(v\) 路的最大条数等于分离 \(u\) 和 \(v\) 所需删除的最少顶点数。

\textbf{证明}：将图的每个顶点 \(x\)（\(x\neq u,v\)）拆分为 \(x_{in},x_{out}\)，边容量设为 1（入边接 \(x_{in}\)，出边接 \(x_{out}\)，\(x_{in}\to x_{out}\) 容量 1），原始无向边替换为双向有向边均容量 \(\infty\)。最大流 \(=\) 最小割（容量为整数）。每条容量 1 的内部弧限制了路径不能共享内部顶点，故最大流 \(=\) 不相交路径数，最小割 \(=\) 分离顶点数。由最大流最小割定理得证。Menger 给出了不依赖网络流的自然证明。\(\blacksquare\)

\textbf{定理 128.7}（Menger 定理，边版本）：图 \(G\) 中两个顶点 \(u, v\) 之间的边不相交的 \(u\)-\(v\) 路的最大条数等于分离 \(u\) 和 \(v\) 所需删除的最少边数。

\textbf{证明}：构建辅助有向图：将 \(G\) 的每条无向边替换为两条方向相反的有向边，容量均为 \(1\)。\(u\) 为源点，\(v\) 为汇点。由最大流最小割定理，最大 \(u\)-\(v\) 流的值等于最小 \(u\)-\(v\) 割的容量。由于每条边容量为 \(1\)，最大流的值等于边不相交的 \(u\)-\(v\) 路的最大条数（每条路贡献流量 \(1\)）。最小割的容量等于分离 \(u\) 和 \(v\) 所需删除的最少边数。因此最大边不相交路径数等于最小边割数。\(\blacksquare\)

\textbf{推论 128.8}（全局 Menger 定理）：图 \(G\) 是 \(k\)-连通的当且仅当任意两个顶点之间存在至少 \(k\) 条内部不相交的路。

\subsection{128.4 Euler 回路与 Hamilton 回路}\label{euler-ux56deux8defux4e0e-hamilton-ux56deux8def-1}

\textbf{定理 128.9}（Euler 定理，1736）：连通图 \(G\) 有 Euler 回路（经过每条边恰好一次的回路）当且仅当每个顶点的度数为偶数。有 Euler 路（经过每条边恰好一次的路，起终点可不同）当且仅当恰好有两个奇数度顶点。

\textbf{证明}：（\(\Rightarrow\)）若存在 Euler 回路，沿回路遍历时每次进入某顶点就沿另一条边离开，故度数为偶数。（\(\Leftarrow\)）所有顶点度数为偶数时，从任意顶点出发沿边遍历，不重复使用边，直到局部无可用边停止。停止点必为出发点（否则奇数度矛盾）。若仍有未访问边，由连通性存在与已访问子图接触的未访问边，在其上继续过程，最终合并所有子回路得 Euler 回路。Euler 路：添加连接两奇度顶点的虚边使全图偶度，找 Euler 回路后移除虚边。\(\blacksquare\)

\emph{证明}：必要性：每进入一个顶点就离开一次，故度数为偶数。充分性：使用 Fleury 算法或 Hierholzer 算法（贪心合并回路）。∎

Euler 回路问题关注的是''边''的遍历------恰经过每条边一次------其判定条件（所有顶点的度数为偶数）简洁而优美。Hamilton 回路则关注''顶点''的遍历------恰经过每个顶点一次------然而与之形成鲜明对比的是，判定一般图是否含有 Hamilton 回路是一个 NP 完全问题。尽管如此，在满足一定最小度条件的稠密图中，存在保证 Hamilton 回路成立的充分条件。

\textbf{定义 128.8}（Hamilton 回路）：Hamilton 回路是经过每个顶点恰好一次的回路。判定 Hamilton 回路的存在性是 NP 完全的。

\textbf{定理 128.10}（Dirac 定理，1952）：若 \(G\) 是 \(n \geq 3\) 个顶点的图，且 \(\delta(G) \geq n/2\)，则 \(G\) 有 Hamilton 回路。

\textbf{证明}：考虑 \(G\) 中最长路 \(P=v_1\cdots v_k\)。若 \(k=n\) 且 \(v_1\) 与 \(v_n\) 相邻则得 Hamilton 回路。否则由于 \(\delta(G)\ge n/2\)，\(v_1\) 的 \(\ge n/2\) 个邻居必在 \(P\) 上。考虑 \(v_k\) 的邻居：设 \(N(v_1)=\{v_{i_1},\ldots\}\)。若 \(v_{i-1}\in N(v_k)\) 对某 \(i\)，则 \(v_1\cdots v_{i-1}v_k v_{k-1}\cdots v_i v_1\) 构成回路。由抽屉原理和 \(\delta\ge n/2\) 保证该 \(i\) 存在。得到回路后若未覆盖全部顶点，由连通度可将回路扩展。因此最长路必为 Hamilton 路且端点相邻，证得 Hamilton 回路。\(\blacksquare\)

\textbf{定理 128.11}（Ore 定理，1960）：若对 \(G\) 中任意不相邻的顶点 \(u, v\)，有 \(\deg(u) + \deg(v) \geq n\)，则 \(G\) 有 Hamilton 回路。

\textbf{证明}：反证法。假设 \(G\) 满足 Ore 条件但无 Hamilton 回路。在 \(G\) 中添加尽可能多的边而不产生 Hamilton 回路，得极大非 Hamilton 图 \(G'\)。\(G'\) 仍有 Ore 条件（添加边只会增加度数）。\(G'\) 中必存在不相邻顶点 \(u, v\)（否则 \(G'\) 为完全图，必有 Hamilton 回路）。添加边 \(\{u,v\}\) 到 \(G'\) 产生 Hamilton 回路 \(u = v_1 \to v_2 \to \cdots \to v_n = v \to u\)。对每个 \(i\)（\(2 \leq i \leq n\)），若 \(\{u, v_i\} \in E(G')\)，则 \(\{v, v_{i-1}\} \notin E(G')\)（否则 \(u \to v_i \to \cdots \to v_n = v \to v_{i-1} \to \cdots \to v_1 = u\) 构成 \(G'\) 中的 Hamilton 回路）。因此 \(\deg_{G'}(v) \leq (n-1) - \deg_{G'}(u)\)，即 \(\deg(u) + \deg(v) \leq n-1\)，与 Ore 条件矛盾。故 \(G\) 必有 Hamilton 回路。\(\blacksquare\)

\subsection{128.5 平面图}\label{ux5e73ux9762ux56fe-1}

\textbf{定义 128.9}（平面图）：图 \(G\) 是\textbf{平面图}，如果它可以画在平面上且边不交叉（仅在端点相交）。平面图的画法将平面划分为\textbf{面}（face）。

\textbf{定理 128.12}（Euler 公式，1758）：对连通平面图 \(G\)，有

\[n - m + f = 2\]

其中 \(n = |V|\)，\(m = |E|\)，\(f\) 为面数（包括外部无界面）。

\emph{证明}：对边数 \(m\) 归纳。奠基：\(m = 0\) 时连通图仅有 \(n = 1\)，\(f = 1\)，\(1 - 0 + 1 = 2\) 成立。 归纳假设 \(m \geq 1\)。考虑两种情况：（1）若 \(G\) 是树，则 \(m = n - 1\)，\(f = 1\)（仅外部面），\(n - (n-1) + 1 = 2\)。（2）若 \(G\) 含回路，取回路上的任意边 \(e\)。\(G \setminus e\) 仍连通且面数减少 1（\(e\) 分隔的两个面合并为一个），边数减 1，顶点数不变。由归纳假设 \(n - (m-1) + (f-1) = 2\)，即 \(n - m + f = 2\)。∎

\textbf{推论 128.13}：对 \(n \geq 3\) 的简单平面图，\(m \leq 3n - 6\)。若图中无三角形，则 \(m \leq 2n - 4\)。

\emph{证明}：每个面至少由 3 条边围成，而每条边至多属于 2 个面，故 \(3f \leq 2m\)。代入 Euler 公式：\(2 = n - m + f \leq n - m + 2m/3 = n - m/3\)，得 \(m \leq 3n - 6\)。∎

\textbf{推论 128.14}：\(K_5\) 和 \(K_{3,3}\) 不是平面图。

\emph{证明}：\(K_5\)：\(n = 5, m = 10\)，\(10 > 3 \cdot 5 - 6 = 9\)。\(K_{3,3}\)：\(n = 6, m = 9\)，无三角形，\(9 > 2 \cdot 6 - 4 = 8\)。∎

\textbf{定理 128.15}（Kuratowski 定理，1930）：图 \(G\) 是平面图当且仅当 \(G\) 不含 \(K_5\) 或 \(K_{3,3}\) 的细分（即通过插入度数为 2 的顶点得到的子图）。

\textbf{证明}：（\(\Rightarrow\)）\(K_5\) 和 \(K_{3,3}\) 均非平面。\(K_5\)：5 顶点 10 边 \textgreater{} \(3n-6=9\) 违反平面图边数界。\(K_{3,3}\)：二分图无奇回路，若平面由 Euler 公式 \(F\le 2E/4=4\)（最小面度 \(\ge4\)），面积和 \(\sum\deg(F)=2E=18\)，矛盾。（\(\Leftarrow\)）充分性由 Wagner 定理（图论基本定理）：无 \(K_5\) 或 \(K_{3,3}\) 细分的图必为平面图。证明通过极小反例的构造分析：取顶点数最少的非平面 \(K_{5}/K_{3,3}\) 细分自由图，证明其 2-连通性，分析 2-点割和块分解，最终导出矛盾或归约到 \(K_5\)/\(K_{3,3}\) 细分。\(\blacksquare\)

\emph{充分性简述}（若不含 \(K_5\) 或 \(K_{3,3}\) 细分则平面）：由推论 128.14 知 \(K_5\) 和 \(K_{3,3}\) 本身非平面，其细分亦然。充分性证明较复杂：对顶点数归纳，利用 3-连通图的平面嵌入性质（Tutte 关于凸嵌入的定理），最终归结为不含此类子图的图必可平面嵌入。详细证明超出本教材范围。∎

\subsection{128.6 有向图与竞赛图}\label{ux6709ux5411ux56feux4e0eux7adeux8d5bux56fe-1}

有向图将边的无序对替换为有序弧，从而引入了方向的维度。竞赛图是有向图中最基本且最具组合结构的一类------每对顶点之间恰有一条弧，可以理解为完全图的定向。Landau 在 1953 年给出了刻画竞赛图得分序列（即各顶点的出度序列）的充要条件，其简洁优雅的形式和基于 Hall 条件的证明使其成为有向图理论中的经典结果。

\textbf{定义 128.10}（有向图与竞赛图）：\textbf{有向图} \(D = (V, A)\) 由顶点集 \(V\) 和弧集 \(A \subseteq V \times V\) 组成（弧是有序对）。\textbf{竞赛图}（tournament）是每对顶点之间恰有一条弧的有向图，即完全图的定向。竞赛图 \(T\) 的\textbf{得分序列}为各顶点的出度 \(s_1 \leq s_2 \leq \cdots \leq s_n\)。

\textbf{定理 128.16}（Landau 定理，1953）：非负整数序列 \((s_1, \ldots, s_n)\)（\(s_1 \leq \cdots \leq s_n\)）是某竞赛图的得分序列当且仅当对所有 \(k = 1,\ldots,n\)，\(\sum_{i=1}^k s_i \geq \binom{k}{2}\)，且 \(k=n\) 时等号成立。该定理的证明使用网络流模型：构造二分图，左侧为顶点对，右侧为顶点，流的存在性由 Hall 条件保证。关于竞赛图的基本结构性质，Rédei 定理（1934）断言每个竞赛图都有 Hamilton 路（有向路经过所有顶点），证明通过归纳法：删除任意顶点后归纳构造 Hamilton 路，再将删除的顶点插入适当位置。Moon 定理（1968）进一步推广：每个强连通竞赛图是顶点泛圈的（vertex-pancyclic），即每个顶点包含在所有长度 \(k = 3,\ldots,n\) 的有向回路中。竞赛图也是社会选择理论（Arrow 不可能定理的组合版本）和排序算法分析中的重要模型。

\textbf{注记 128.1}（图论与代数的联系）：图的邻接矩阵 \(A(G)\) 的谱（特征值集合）蕴含了图的深层结构信息。谱图论（spectral graph theory）通过 \(A(G)\) 和 \(L(G)\) 的特征值研究图的扩张性、等周不等式和随机游走性质。对 \(d\)-正则图，\(A(G)\) 的最大特征值为 \(d\)，第二大特征值 \(\lambda_2\) 控制图的扩张性质：\(\lambda_2\) 越小，图的边扩张越好。Alon-Boppana 定理给出 \(\lambda_2 \geq 2\sqrt{d-1} - o(1)\) 的下界，而 Ramanujan 图恰好达到此界，是扩张图和编码理论中的理想构造。图的 Laplace 矩阵 \(L(G)\) 的谱与图上的随机游走、有效电阻和生成树计数（Kirchhoff 矩阵树定理）有直接联系：\(L(G)\) 的非零特征值乘积等于 \(n\) 乘以生成树个数。

\section{第94章：匹配与网络流}\label{ux7b2c94ux7ae0ux5339ux914dux4e0eux7f51ux7edcux6d41}

匹配理论和网络流是离散数学中一对深刻对偶的概念：匹配研究的是无公共顶点的边集最大化问题（``配对''），网络流研究的是有向图中约束容量下的最大流量（``传输''），而最大流-最小割定理和 König 定理（二分图中最大匹配 = 最小顶点覆盖）将二者统一为同一线性规划对偶性的不同实例。匹配理论从 Hall 婚配定理（1935 年）出发，经由 Berge 引理（增广路理论，1957 年）和 Edmonds 的花算法（1965 年）达到算法完备性，而网络流理论从 Ford-Fulkerson 增广路方法（1956 年）经由 Dinic 分层网络（1970 年）到 Goldberg-Tarjan 推进-重标号算法（1988 年），逐步逼近理论最优的复杂度。匹配与网络流不仅在运筹学（运输问题、分配问题）中占据核心地位，更在计算机科学（调度、图割分割）和经济学（市场设计）中有着广泛而深刻的应用。

\subsection{129.1 二分图匹配}\label{ux4e8cux5206ux56feux5339ux914d}

\textbf{定义 129.1}（匹配）：图 \(G = (V, E)\) 的一个\textbf{匹配} \(M \subseteq E\) 是两两不共享顶点的边集合。最大匹配是边数最多的匹配。完美匹配覆盖所有顶点。

\textbf{定义 129.2}（增广路）：对于匹配 \(M\)，\textbf{交替路}是边在 \(M\) 和非 \(M\) 之间交替的路。\textbf{增广路}是起终点都为未匹配顶点的交替路。

\textbf{定理 129.1}（Berge 引理，1957）：匹配 \(M\) 是最大匹配当且仅当不存在 \(M\)-增广路。

\textbf{证明}：（\(\Rightarrow\)）若 \(M\) 最大但存在 \(M\)-增广路 \(P\)（交替属于 \(M\) 和不属于 \(M\) 的边，起终点均未被 \(M\) 覆盖），则 \(M'=M\triangle P=(M\setminus P)\cup(P\setminus M)\) 为匹配且 \(|M'|=|M|+1\)，与最大性矛盾。（\(\Leftarrow\)）若无 \(M\)-增广路但 \(M\) 非最大，存在更大匹配 \(M^*\)。考虑 \(M\triangle M^*\)：每个顶点度 \(\le2\)，故为不相交的交替路和交替回路组件。因 \(|M^*|>|M|\)，必存在以 \(M^*\) 边开头的交替路（\(M\)-增广路），矛盾。此引理是匹配算法（增广路算法）的理论基础。\(\blacksquare\)

\emph{证明}：若存在增广路 \(P\)，则对称差 \(M \triangle P\) 是更大的匹配。反之，若 \(M\) 非最大，取 \(M'\) 更大，\(M \triangle M'\) 的每个连通分支要么是偶长度回路，要么是包含 \(M'\) 边更多的路，后者即为 \(M\)-增广路。∎

\textbf{定理 129.2}（Kőnig 定理，1931）：在二分图 \(G = (A \cup B, E)\) 中，最大匹配的大小等于最小顶点覆盖的大小。

\textbf{证明}：设 \(M\) 为最大匹配，\(U_A \subseteq A\) 为 \(A\) 中未被 \(M\) 匹配的顶点集。从 \(U_A\) 出发，沿交替路（交替于不属于 \(M\) 的边和属于 \(M\) 的边）可达的所有顶点记为 \(Z\)。令 \(K = (A \setminus Z) \cup (B \cap Z)\)。首先证 \(K\) 为顶点覆盖：若存在边 \(\{a,b\} \notin K\) 的覆盖，则 \(a \in A \cap Z\) 且 \(b \in B \setminus Z\)。若 \(\{a,b\} \in M\)，\(b\) 必从 \(a\) 沿 \(M\) 边可达，故 \(b \in Z\)，矛盾。若 \(\{a,b\} \notin M\)，\(a\) 从 \(U_A\) 可达，沿边 \(\{a,b\}\) 一步可达 \(b\)，故 \(b \in Z\)，矛盾。因此 \(K\) 为顶点覆盖。其次，\(|K| = |M|\)：\(K\) 中每个顶点覆盖的 \(M\) 边互不相交（因 \(A \setminus Z\) 的顶点必被 \(M\) 匹配，\(B \cap Z\) 的顶点也必被 \(M\) 匹配），且 \(M\) 中每条边恰被 \(K\) 中一个顶点覆盖。故 \(|K| = |M|\)，即最大匹配大小等于最小顶点覆盖大小。\(\blacksquare\)

\textbf{定理 129.3}（Hall 婚姻定理，1935）：二分图 \(G = (A \cup B, E)\) 存在覆盖 \(A\) 的匹配当且仅当对任意 \(S \subseteq A\)，有 \(|N(S)| \geq |S|\)（Hall 条件），其中 \(N(S)\) 是 \(S\) 的邻居集。

\textbf{证明}：（\(\Rightarrow\)）若存在覆盖 \(A\) 的匹配，\(S\) 中每个顶点配到不同的 \(N(S)\) 中顶点，故 \(|N(S)|\ge|S|\)。（\(\Leftarrow\)）对 \(|A|\) 归纳。\(|A|=1\) 平凡。\(|A|>1\) 时分两情形。情形 1：对所有 \(\emptyset\neq S\subsetneq A\)，\(|N(S)|\ge|S|+1\)。任取 \(a\in A\) 和其邻居 \(b\in B\)，匹配 \((a,b)\) 后从图中删除。剩余图对归约后的 Hall 条件仍成立（因严格条件缓冲了一个单位的松弛），归纳完成。情形 2：存在 \(\emptyset\neq S\subsetneq A\) 满足 \(|N(S)|=|S|\)。由归纳假，\(G[S\cup N(S)]\) 存在饱和 \(S\) 的匹配。接下来证 \(G[(A\setminus S)\cup(B\setminus N(S))]\) 满足 Hall 条件：对任意 \(T\subseteq A\setminus S\)，\(|N_{A\setminus S}(T)|\ge|T|\)------否则 \(S\cup T\) 在 \(G\) 中的邻居少于 \(|S|+|T|\)，与原始 Hall 条件矛盾。归纳完成全部匹配。\(\blacksquare\)

\emph{证明}：必要性：若存在覆盖 \(A\) 的匹配 \(M\)，则对任意 \(S \subseteq A\)，\(M\) 将 \(S\) 匹配到 \(|S|\) 个不同的 \(B\) 中顶点，故 \(|N(S)| \geq |S|\)。

充分性（对 \(|A|\) 归纳）：\(|A|=1\) 时平凡。设 \(|A|>1\)。分两种情况：（1）若对任意非空真子集 \(S \subset A\) 有 \(|N(S)| \geq |S|+1\)（Hall 条件加强），任选 \(a \in A\) 匹配到其邻居 \(b \in B\)，删去两者后剩余图对 \(A\setminus\{a\}\) 仍满足 Hall 条件，由归纳假设得证。（2）存在非空真子集 \(S \subset A\) 使得 \(|N(S)| = |S|\)。对二分图 \(S \cup N(S)\)，由归纳假设存在 \(S\) 到 \(N(S)\) 的完美匹配 \(M_1\)。考虑 \(G \setminus (S \cup N(S))\)。对任意 \(T \subseteq A \setminus S\)，\(|N_{G\setminus}(T)| = |N(T \cup S) \setminus N(S)| \geq |T \cup S| - |N(S)| = |T| + |S| - |S| = |T|\)。故 \(A \setminus S\) 到 \(B \setminus N(S)\) 满足 Hall 条件，由归纳假设存在匹配 \(M_2\)。\(M_1 \cup M_2\) 即为覆盖 \(A\) 的匹配。∎

\textbf{推论 129.4}：正则二分图有完美匹配。特别地，\(k\)-正则二分图可分解为 \(k\) 个不相交的完美匹配。

\subsection{129.2 一般图匹配}\label{ux4e00ux822cux56feux5339ux914d}

\textbf{定义 129.3}（Tutte 定理，1947）：图 \(G\) 有完美匹配当且仅当对任意 \(S \subseteq V\)，\(G \setminus S\) 的奇数阶连通分支数不超过 \(|S|\)。

Tutte 定理是一般图匹配理论的基石，它将二分图匹配的 Hall 条件推广到任意图。记 \(o(G - S)\) 为 \(G \setminus S\) 中奇数个顶点的连通分支的个数，Tutte 条件为 \(o(G - S) \leq |S|\) 对所有 \(S \subseteq V\) 成立。其必要性是直观的：在完美匹配中，每个奇分支必须至少有一个顶点与 \(S\) 中的顶点匹配，而 \(S\) 中每个顶点至多匹配一个奇分支。充分性的证明（Tutte 1947）通过构造性论证完成：若无完美匹配，考虑极小反例，分析其删点结构，最终导出与 Tutte 条件的矛盾。Tutte 定理的一个重要推论是 Petersen 定理（1891）：每个不含桥的 3-正则图有完美匹配。Tutte 定理也是推导 Berge-Tutte 公式（最大匹配大小的显式表达）的基础。

\textbf{定义 129.4}（Edmonds 花算法，1965）：Edmonds 在 1965 年提出了多项式时间的最大匹配算法（花算法），核心思想是收缩奇长度回路（花）以处理非二分图中的奇回路障碍。

花算法是一般图匹配问题的第一个多项式时间算法，也是组合优化历史上的里程碑。其核心挑战在于处理一般图中的奇回路：在二分图中，增广路寻找是直接的（BFS 即可），但一般图中奇回路可能导致增广路搜索陷入死循环。Edmonds 的洞察是：当 BFS 搜索发现奇回路（称为''花''或 blossom）时，可将整个花收缩为一个超级顶点，继续搜索增广路。若在收缩后的图中找到增广路，则可在原图中''展开''花以恢复完整增广路。这一''收缩-展开''策略保证了算法的正确性和多项式时间界。Edmonds 花算法的时间复杂度为 \(O(|V|^3)\)，后续改进（Micali-Vazirani 1980）达到 \(O(\sqrt{|V|}|E|)\)。该算法也标志着 Edmonds 提出''多项式时间算法''这一概念，为计算复杂性理论奠定了基础。

\subsection{129.3 网络流与最大流}\label{ux7f51ux7edcux6d41ux4e0eux6700ux5927ux6d41}

\textbf{定义 129.5}（网络）：一个\textbf{网络} \(N = (V, E, s, t, c)\) 由有向图 \((V, E)\)、源点 \(s\)、汇点 \(t\) 和容量函数 \(c: E \to \mathbb{R}_{\geq 0}\) 组成。

\textbf{定义 129.6}（流）：一个\textbf{流} \(f: E \to \mathbb{R}_{\geq 0}\) 满足： 1. \textbf{容量限制}：\(0 \leq f(e) \leq c(e)\) 对所有 \(e \in E\) 2. \textbf{流量守恒}：对 \(v \neq s, t\)，\(\sum_{e \text{ in } v} f(e) = \sum_{e \text{ out } v} f(e)\)

流的值 \(|f| = \sum_{e \text{ out } s} f(e) - \sum_{e \text{ in } s} f(e)\)。

\textbf{定义 129.7}（割）：\(s\)-\(t\) 割 \((S, T)\) 满足 \(s \in S, t \in T\)。割的容量为 \(c(S, T) = \sum_{e \text{ from } S \text{ to } T} c(e)\)。

\textbf{定理 129.5}（最大流最小割定理 / Ford-Fulkerson，1956）：在任何网络中，最大流的值等于最小割的容量。

\textbf{证明}：设 \(f\) 为 Ford-Fulkerson 算法终止时的流。定义 \(S\) 为剩余网络 \(G_f\) 中从 \(s\) 可达的顶点集。由于算法终止，\(t \notin S\)（否则存在增广路）。对任意从 \(S\) 到 \(T = V \setminus S\) 的边 \(e = (u,v)\)（\(u \in S, v \in T\)），在 \(G_f\) 中 \(u\) 可达 \(v\) 不成立，故 \(f(e) = c(e)\)（否则 \(G_f\) 中有正向边 \((u,v)\) 容量 \(c(e)-f(e) > 0\)）。对任意从 \(T\) 到 \(S\) 的边 \(e' = (v',u')\)（\(v' \in T, u' \in S\)），\(f(e') = 0\)（否则 \(G_f\) 中有反向边 \((u',v')\) 容量 \(f(e') > 0\)）。因此 \(|f| = \sum_{e \in S \to T} f(e) - \sum_{e \in T \to S} f(e) = \sum_{e \in S \to T} c(e) = c(S,T)\)。同时，对任意流 \(f'\) 和任意割 \((S',T')\)，\(|f'| \leq c(S',T')\)（弱对偶性）。故 \(f\) 为最大流，\((S,T)\) 为最小割。\(\blacksquare\)

\textbf{定理 129.6}（整性定理）：若所有容量为整数，则存在整数值的最大流。

\textbf{证明}：Ford-Fulkerson 算法从零流（整数）开始，每次沿增广路推送的流量为路径上各边剩余容量的最小值。由于所有容量为整数，每一步的剩余容量均为整数，故每次增广的流量值为整数。算法终止时得到的最大流在每个边上的流量均为整数。此外，由最大流最小割定理，最小割的容量也是整数。\(\blacksquare\)

\textbf{定义 129.8}（Edmonds-Karp 算法）：每次选择最短增广路（BFS），时间复杂度 \(O(|V| |E|^2)\)。Dinic 算法使用分层网络，时间复杂度 \(O(|V|^2 |E|)\)。

\textbf{定理 129.7}（最大流与二分图匹配）：二分图 \(G = (A \cup B, E)\) 的最大匹配问题等价于网络流：添加源点 \(s\) 连接 \(A\) 中所有顶点（容量 1），\(B\) 中所有顶点连接汇点 \(t\)（容量 1），每条边 \(A \to B\) 容量 1。最大流的值等于最大匹配的大小。

\textbf{证明}：设 \(M\) 为 \(G\) 的一个匹配。构造流 \(f\)：对每条匹配边 \((a, b) \in M\)，令 \(f(s, a) = f(a, b) = f(b, t) = 1\)；其余边流量为 \(0\)。由匹配的定义，\(A\) 中每个顶点至多关联一条匹配边，\(B\) 中每个顶点也至多关联一条匹配边，故 \(s\) 和 \(t\) 的容量约束自动满足。因此 \(|f| = |M|\)。反之，设 \(f\) 为整数流（由整性定理，存在整数最大流）。由于每条边容量为 \(1\)，\(f\) 在每条边上的流量为 \(0\) 或 \(1\)。定义 \(M = \{(a, b) \in E : f(a, b) = 1\}\)。由流守恒，\(A\) 中每个顶点至多有一条出边流量为 \(1\)，\(B\) 中每个顶点至多有一条入边流量为 \(1\)，故 \(M\) 为匹配且 \(|M| = |f|\)。因此最大流的值等于最大匹配的大小。\(\blacksquare\)

\subsection{129.4 最小费用流与分配问题}\label{ux6700ux5c0fux8d39ux7528ux6d41ux4e0eux5206ux914dux95eeux9898}

最大流问题关注的是流量的大小，而最小费用流在流量目标固定的前提下进一步引入每条边上的费用函数，追求总费用最低的流方案。分配问题是其最重要的特例------在完全二分图中寻找最小费用的完美匹配------它在运筹学、调度理论和经济学中均有核心应用。匈牙利算法（Kuhn 1955 年，Munkres 1957 年）通过对偶变量的迭代调整在 \(O(n^3)\) 时间内求解分配问题，是组合优化中最早的多项式时间算法之一。

\textbf{定义 129.9}（最小费用流）：在网络上还定义边费用 \(w(e)\)，求从 \(s\) 到 \(t\) 发送给定流量 \(F\) 的最小总费用 \(\sum_e w(e) f(e)\)。

\textbf{定义 129.10}（分配问题 / Assignment Problem）：给定 \(n\) 个工人和 \(n\) 个任务，工人 \(i\) 完成任务 \(j\) 的费用为 \(c_{ij}\)，求使总费用最小的一一分配。这是最小费用流（在完全二分图上）的特例。

\textbf{定理 129.8}（匈牙利算法 / Kuhn-Munkres 算法）：可在 \(O(n^3)\) 时间内求解分配问题，基于对偶变量（势）的调整。

\textbf{证明}：分配问题为 \(\min\sum c_{ij}x_{ij}\)（置换矩阵）。对偶变量 \(u_i,v_j\) 满足 \(u_i+v_j\le c_{ij}\)。互补松弛条件：\(x_{ij}>0\implies u_i+v_j=c_{ij}\)。算法维护可行对偶解 \(u_i,v_j\) 和对应于等式子图 \(E_{\text{eq}}=\{(i,j):u_i+v_j=c_{ij}\}\) 上的最大匹配。若匹配未覆盖所有行，设 \(S\) 为交替路可达的行集，\(T\) 为可达列集。更新对偶变量：\(\delta=\min_{i\in S,j\notin T}(c_{ij}-u_i-v_j)\)，\(u_i\leftarrow u_i+\delta\)（\(i\in S\)），\(v_j\leftarrow v_j-\delta\)（\(j\in T\)）。等式边集扩大，匹配逐步完善。每次迭代增广一枚匹配边或缩小对偶间距，\(O(n)\) 次迭代每次 \(O(n^2)\)。\(\blacksquare\)

\subsection{129.5 最小割与连通性应用}\label{ux6700ux5c0fux5272ux4e0eux8fdeux901aux6027ux5e94ux7528}

最大流-最小割定理的核心价值不仅在于流量最优化本身，更在于它提供了图连通性的统一代数框架。事实上，Menger 定理（128.3 节）是最大流-最小割定理在所有边容量设为 1 时的特例，而最大流-最小割定理可以反过来通过将一般容量网络转化为多重图来导出。这一等价性将图论中看似不同的概念------不相交路径、顶点割、边割和最大流------统一为一个深刻的对偶原理。

\textbf{定理 129.9}（Menger 定理与最大流最小割定理的等价性）：Menger 定理（边版本）是最大流最小割定理在容量全为 1 的网络中的特例。反之，最大流最小割定理可通过将一般容量网络转化为多重图上的 Menger 问题来证明。这一等价性揭示了图连通性与网络流之间的深层对偶关系。

\textbf{推论 129.10}（全局边连通度）：\(\lambda(G) \geq k\) 当且仅当对任意 \(s, t \in V\)，存在 \(k\) 条边不交的 \(s\)-\(t\) 路。这等价于 \(G\) 的任意割的容量（边数）至少为 \(k\)。类似地，点连通度 \(\kappa(G) \geq k\) 当且仅当任意两顶点间存在 \(k\) 条内部不相交的路（Menger 定理点形式的全局版本）。

\textbf{注记 129.1}（网络流理论的历史脉络）：网络流问题起源于 20 世纪 40 年代的运筹学研究。Kantorovich（1939）和 Hitchcock（1941）独立提出了运输问题，Koopmans（1947）将其形式化为线性规划。Dantzig（1951）的单纯形法为网络流提供了通用求解框架。Ford 和 Fulkerson（1956）提出了最大流最小割定理和增广路算法，首次将组合图论与线性规划方法统一。Edmonds 和 Karp（1972）通过选择最短增广路将算法复杂度从依赖于容量值改进为多项式时间 \(O(|V||E|^2)\)。Dinic（1970）引入分层网络和阻塞流概念，达到 \(O(|V|^2|E|)\)。Goldberg 和 Tarjan（1988）的推进-重标号算法（push-relabel）进一步突破，达到 \(O(|V|^2\sqrt{|E|})\)。网络流理论的应用远超图论本身：在计算机视觉（图割分割）、生物信息学（序列比对）、调度理论（作业分配）和分布式计算（负载均衡）中均有核心地位。

\section{第95章：染色与 Ramsey 理论}\label{ux7b2c95ux7ae0ux67d3ux8272ux4e0e-ramsey-ux7406ux8bba}

染色理论与 Ramsey 理论代表了图论中''必然存在有序结构''这一深刻洞察的两个侧面。染色问题研究用最少颜色对图的顶点（或边）进行标号使得相邻元素不同色，其经典成就包括四色定理（Appel-Haken, 1976 年------首个依赖计算机辅助证明的数学定理）和 Brooks 定理（\(\chi(G) \leq \Delta(G)\)）。Ramsey 理论则源于 Frank Ramsey（1930 年）的著名洞见：``完全的无序是不可能的''------任何足够大的结构中必然出现某种有序子结构。如 \(R(3,3) = 6\) 表明任意 6 个人的聚会中必然存在 3 个两两认识或两两互不认识的人。Ramsey 理论将染色问题反转：不是寻找避免特定子图的最小色数，而是寻找使任意 \(k\)-染色必然包含单色特定子图的最小规模。本章在 Brooks 定理之后，以 Ramsey 数的渐近界（Erdős-Szekeres 上界和 Erdős 概率方法下界）为核心，阐述染色与 Ramsey 理论的基本框架。

\subsection{130.1 顶点染色}\label{ux9876ux70b9ux67d3ux8272}

\textbf{定义 130.1}（色数）：图 \(G\) 的\textbf{色数} \(\chi(G)\) 是使相邻顶点颜色不同的最少颜色数。\(G\) 是 \(k\)-可染的，如果 \(\chi(G) \leq k\)。

\textbf{命题 130.1}（基本界）： - \(\chi(G) \geq \omega(G)\)（团数---最大完全子图的大小） - \(\chi(G) \geq \frac{n}{\alpha(G)}\)（独立数---最大独立集的大小） - \(\chi(G) \leq \Delta(G) + 1\)（贪心染色）

\textbf{定理 130.2}（Brooks 定理，1941）：若 \(G\) 是连通图且非完全图也非奇回路，则 \(\chi(G) \leq \Delta(G)\)。

\textbf{证明}：对 \(|V|\) 归纳。若 \(G\) 非 \(\Delta\)-正则，取非正则顶点按特定顺序贪心染色可得 \(\chi\le\Delta\)。若 \(G\) 是 \(\Delta\)-正则且非完全图非奇回路，由连通性存在树结构：选顶点 \(v\)，取以 \(v\) 为根的 DFS 树。\(v\) 有至少两个非相邻邻居 \(x,y\)（否则 \(G\) 为完全图）。构造染色顺序 \((x=y_1,\ldots,y_n=v)\) 使得每个 \(y_i\) 有已染色的 \(\le\Delta-1\) 个邻居（因 \(x\) 和 \(y\) 不相邻可在排序末尾处理），贪心 \(\Delta\) 色足够。\(\blacksquare\)

\textbf{定理 130.3}（四色定理 / Appel-Haken，1976）：任何平面图是 4-可染的。

\textbf{证明}（Appel-Haken 证明框架）：证明分为两大步骤。(1) 不可免集（unavoidable set）：证明每个平面图必包含 1936 种「可约构型」（reducible configuration）中的至少一种。这些构型是度不超过 5 的顶点及其邻域的特定结构。(2) 可约性（reducibility）：对每种构型，证明若最小反例包含该构型，则可将其「约化」为更小的平面图，对更小的图四染色后可延拓回原图。具体地，若 \(G\) 是最小反例（不含 4-染色的极小平面图），则 \(G\) 必含某个可约构型 \(C\)。删除 \(C\) 中的某些顶点得到 \(G'\)，由极小性 \(G'\) 可 4-染色。由 \(C\) 的可约性，此染色可延拓为 \(G\) 的 4-染色，矛盾。Appel 和 Haken 用计算机验证了 1936 种构型的可约性。2015 年 Gonthier 用 Coq 证明助手形式化了整个证明。\(\blacksquare\)

\textbf{定理 130.4}（五色定理 / Heawood，1890）：任何平面图是 5-可染的（此定理有简短的纯数学证明）。

\emph{证明}：平面图必有度 \(\leq 5\) 的顶点。归纳法：删除该顶点，对剩余图 5-染色，若该顶点的邻居使用了所有 5 种颜色，则在由两种颜色诱导的子图中交换颜色来释放一种颜色（Kempe 链论证）。∎

\subsection{130.2 边染色}\label{ux8fb9ux67d3ux8272}

\textbf{定义 130.2}（边色数）：图 \(G\) 的\textbf{边色数} \(\chi'(G)\) 是使相邻边（共享顶点）颜色不同的最少颜色数。

\textbf{定理 130.5}（Vizing 定理，1964）：对任何简单图 \(G\)，\(\Delta(G) \leq \chi'(G) \leq \Delta(G) + 1\)。

\textbf{证明}：\(\chi'\ge\Delta\) 显然（度最大顶点需 \(\Delta\) 种不同颜色边）。证 \(\chi'\le\Delta+1\)：对边数归纳。设 \(G\) 是极小反例（\(\chi'>\Delta+1\) 但删任一边后 \(\chi'\le\Delta+1\)）。取边 \(xy\in E\)，在 \(G-xy\) 中存在 \(\Delta+1\) 边染色 \(\varphi\)。考虑未染色边 \(xy\)：对每个缺失色，尝试 Kempe 链交换调整颜色，最终得矛盾或成功染色 \(xy\)。关键性质：Kempe 链不会无限延长------每个换色链终止于度数 \(<\Delta\) 点（因每点最多 \(\Delta\) 条边共用 \(\Delta+1\) 色，总有缺失色）。用计数论证排除冲突。\(\blacksquare\)

\textbf{定理 130.6}（Kőnig 边染色定理）：对二分图 \(G\)，\(\chi'(G) = \Delta(G)\)。

\emph{证明}：对 \(\Delta(G)\) 归纳，使用 Hall 婚姻定理找到完美匹配覆盖所有最大度顶点，然后删除该匹配。∎

\subsection{130.3 Ramsey 理论}\label{ramsey-ux7406ux8bba-1}

\textbf{定理 130.7}（Ramsey 定理，图版本，1930）：对任意正整数 \(r, s\)，存在最小整数 \(R(r, s)\)，使得任何 \(R(r, s)\) 个顶点的完全图 \(K_n\)（\(n \geq R(r, s)\)），其边任意染红/蓝二色，必含红色 \(K_r\) 或蓝色 \(K_s\)。

\emph{证明}（归纳法）：对 \(r+s\) 做归纳。基始：\(R(1,s)=1\)（红色 \(K_1\) 平凡），\(R(r,1)=1\)，\(R(2,s)=s\)，\(R(r,2)=r\)。归纳步：设 \(R(r-1,s)\) 和 \(R(r,s-1)\) 存在。取 \(n = R(r-1,s) + R(r,s-1)\)，对 \(K_n\) 做红蓝染色。任选顶点 \(v\)，设 \(v\) 有 \(d_R\) 条红边、\(d_B\) 条蓝边，\(d_R+d_B=n-1\)。若 \(d_R \geq R(r-1,s)\)，则红邻居诱导的子图必含红色 \(K_{r-1}\)（连同 \(v\) 得红色 \(K_r\)）或蓝色 \(K_s\)。若 \(d_B \geq R(r,s-1)\)，同理得红色 \(K_r\) 或蓝色 \(K_{s-1}\)（连同 \(v\) 得蓝色 \(K_s\)）。由于 \(d_R+d_B \geq R(r-1,s)+R(r,s-1)-1\)，至少一种情形发生。故 \(R(r,s) \leq R(r-1,s)+R(r,s-1)\)。∎

\textbf{命题 130.8}（Ramsey 数的界）： - \(R(r, s) \leq R(r-1, s) + R(r, s-1)\)（递归不等式） - \(R(r, s) \leq \binom{r+s-2}{r-1}\)（Erdős-Szekeres，1935） - \(R(r, r) \geq 2^{r/2}\)（Erdős 下界，概率方法，1947）

\textbf{已知 Ramsey 数}：\(R(3,3) = 6\)，\(R(3,4) = 9\)，\(R(3,5) = 14\)，\(R(3,6) = 18\)，\(R(4,4) = 18\)，\(R(3,9) = 36\)（2023年）。\(R(5,5)\) 仍未知，仅知 \(43 \leq R(5,5) \leq 48\)。

\textbf{定理 130.9}（Ramsey 定理，超图版本）：对任意正整数 \(r, k, c\)，存在 \(n\) 使得对 \(n\) 元集合的 \(r\) 元子集做 \(c\)-染色，必存在 \(k\) 元子集，其所有 \(r\) 元子集同色。

\textbf{证明}：对 \(r\) 归纳。\(r=1\) 即鸽巢原理：\(n=c(k-1)+1\) 满足。归纳步：设对所有 \(r-1\) 元超图 Ramsey 定理成立。给定 \(c,k,r\)，递归定义 \(n_0=k-1\)，\(n_{i+1}=R_{r-1}(c,n_i)\)（即对 \((r-1)\)-超图应用归纳假设）。取 \(n=n_{c+1}\)。对 \(n\)-集的 \(r\)-子集 \(c\)-染色，固定元素 \(x\) 并考虑含 \(x\) 的 \((r-1)\)-子集颜色（\(c\) 种）。由归纳假设在其余 \(n-1\) 元素中找同色子集，逐步缩小可得单色 \(k\)-集。\(\blacksquare\)

\textbf{定理 130.10}（Schur 定理，1916）：对任意 \(c\)，存在 \(N\) 使得对 \(\{1, \ldots, N\}\) 做 \(c\)-染色，必存在 \(x, y, z\) 同色且 \(x + y = z\)。

\emph{证明}：这是 Ramsey 定理的推论。取 \(N = R(3,3,\ldots,3)\)（\(c\) 个 3），对完全图 \(K_N\) 的边 \(\{i,j\}\)（\(i < j\)）染以 \(|i-j|\) 的颜色。由 Ramsey 定理，存在同色三角形，其边长为 \(i-j, j-k, i-k\)，满足 \((i-j)+(j-k) = i-k\)。∎

\textbf{定理 130.11}（van der Waerden 定理，1927）：对任意正整数 \(r, k\)，存在 \(W(r, k)\) 使得对 \(\{1, \ldots, W(r,k)\}\) 做 \(r\)-染色，必存在长度为 \(k\) 的同色等差数列。

\textbf{证明}（对 \(k\) 归纳）：\(k=1,2\) 时平凡。设对 \(k\) 成立，需证对 \(k+1\) 成立。对 \(r\) 归纳。定义 \(N_1 = W(r, k)\)，\(N_2 = W(r^{N_1}, k)\)，\(N_3 = W(r^{N_1 N_2}, k)\)，以此类推构造 \(N_1, N_2, \ldots, N_r\)。取 \(W(r, k+1) = N_1 N_2 \cdots N_r\)。将 \(\{1, \ldots, W\}\) 划分为 \(N_r\) 个长度为 \(N_1 \cdots N_{r-1}\) 的块。在每个块内，由归纳假设存在长度为 \(k\) 的同色等差数列。考虑这些等差数列的「步长」和「颜色」对，共 \(r^{N_1 \cdots N_{r-1}}\) 种可能。由鸽巢原理和 \(N_r\) 的选取，必有两个块中对应的等差数列同色且步长相同。这两个等差数列加上它们之间对应位置的等差延伸，构成长度为 \(k+1\) 的同色等差数列。\(\blacksquare\)

\textbf{定理 130.12}（Szemerédi 定理，1975）：任何正上密度的自然数子集包含任意长的等差数列。即对任意 \(\delta > 0\) 和 \(k \in \mathbb{N}\)，存在 \(N\) 使得 \(\{1, \ldots, N\}\) 的任意大小 \(\geq \delta N\) 的子集包含长度为 \(k\) 的等差数列。

\textbf{证明}（\(k=3\) 的 Roth 证明概要）：设 \(A \subseteq \{1, \ldots, N\}\) 满足 \(|A| = \delta N\)。将 \(A\) 嵌入 \(\mathbb{Z}_N\)（\(N\) 为素数）。定义 \(f = 1_A - \delta\) 为平衡函数。三元等差数列的计数为 \(\Lambda(A) = \sum_{x,d} 1_A(x) 1_A(x+d) 1_A(x+2d)\)。利用 Fourier 展开，\(\Lambda(A) = \delta^3 N^2 + \sum_{\xi \neq 0} \hat{f}(\xi)^2 \hat{f}(-2\xi)\)。当 \(\delta\) 固定时，主项 \(\delta^3 N^2\) 为正。若 \(A\) 无三元等差数列，则 \(\Lambda(A) = O(N)\)。由此导出 \(\|\hat{f}\|_{\infty} \geq c(\delta) > 0\)，即 \(f\) 在某个非零频率上具有大的 Fourier 系数。这意味着 \(A\) 在某个长等差数列上密度增加。通过密度增量迭代（每次迭代密度增加 \(c(\delta^2)\)），经有限步后密度超过 \(1\)，矛盾。因此 \(A\) 必含三元等差数列。一般 \(k\) 的证明使用 Gowers 范数（\(U^{k-1}\) 范数）和 Szemerédi 正则性引理。\(\blacksquare\)

\subsection{130.4 极值图论}\label{ux6781ux503cux56feux8bba-1}

\textbf{定理 130.13}（Turán 定理，1941）：\(n\) 个顶点不含 \(K_{r+1}\) 子图的最大边数为

\[\operatorname{ex}(n, K_{r+1}) = \left(1 - \frac{1}{r}\right) \frac{n^2}{2} + O(n)\]

极大值由 Turán 图 \(T(n, r)\)（将 \(n\) 个顶点尽可能均匀地划分为 \(r\) 部分，所有边连接不同部分的顶点）达到。

\textbf{证明}：对 \(n\) 归纳。\(n \leq r\) 时平凡。设 \(G\) 有 \(n\) 个顶点、\(m\) 条边，不含 \(K_{r+1}\)。取 \(G\) 中度最大的顶点 \(v\)，\(\deg(v) = \Delta\)。令 \(S = N(v)\)，\(T = V \setminus (\{v\} \cup S)\)，则 \(|S| = \Delta\)，\(|T| = n - 1 - \Delta\)。\(G[S]\) 不含 \(K_r\)（否则与 \(v\) 构成 \(K_{r+1}\)）。由归纳假设，\(G[S]\) 中边数 \(\leq \operatorname{ex}(\Delta, K_r)\)，\(G[T]\) 中边数 \(\leq \operatorname{ex}(n-1-\Delta, K_{r+1})\)。\(S\) 与 \(T\) 之间的边数 \(\leq \Delta \cdot (n-1-\Delta)\)。对度序列分析知最优配置为完备 \(r\)-部图，各部大小差至多 \(1\)。计算得 \(\operatorname{ex}(n, K_{r+1}) = \frac{r-1}{2r} n^2 + O(n)\)，即 \((1 - 1/r) \cdot n^2/2 + O(n)\)。\(\blacksquare\)

\textbf{定理 130.14}（Erdős-Stone 定理，1946）：对任意图 \(H\)（色数 \(\chi(H) = r+1\)），

\[\operatorname{ex}(n, H) = \left(1 - \frac{1}{r}\right) \binom{n}{2} + o(n^2)\]

\textbf{证明}（概要）：下界：考虑 Turán 图 \(T_r(n)\)，即 \(n\) 个顶点划分为 \(r\) 个几乎相等的部分，各部分之间全连接。\(T_r(n)\) 的边数为 \((1 - 1/r) \binom{n}{2} + O(n)\)，且 \(T_r(n)\) 不含 \(K_{r+1}\)，故也不含 \(H\)（因 \(\chi(H) = r+1\)）。上界：由 Szemerédi 正则性引理，任意 \(n\) 个顶点的图 \(G\) 可划分为常数个正则对。若 \(G\) 的边数 \(> (1 - 1/r + \varepsilon) \binom{n}{2}\)，则其正则划分的简化图（reduced graph）中必存在 \(K_{r+1}\)。由关键引理（Key Lemma），\(G\) 包含 \(K_{r+1}(t)\)（即每个部分含 \(t\) 个顶点的完全 \((r+1)\)-部图）。取 \(t\) 足够大（\(t \geq |V(H)|\)），则 \(G\) 包含 \(H\) 作为子图。故 \(\operatorname{ex}(n, H) \leq (1 - 1/r + o(1)) \binom{n}{2}\)。\(\blacksquare\)

\subsection{130.5 正则性引理及其应用}\label{ux6b63ux5219ux6027ux5f15ux7406ux53caux5176ux5e94ux7528}

\textbf{定理 130.15}（Szemerédi 正则性引理，1978）：对任意 \(\varepsilon > 0\) 和整数 \(m\)，存在整数 \(M = M(\varepsilon, m)\) 使得任意 \(n \geq m\) 个顶点的图 \(G\) 的顶点集可划分为 \(k\) 个几乎相等的部分 \(V_0, V_1, \ldots, V_k\)（\(m \leq k \leq M\)），其中 \(|V_0| \leq \varepsilon n\) 为例外集，且除至多 \(\varepsilon k^2\) 对外，所有对 \((V_i, V_j)\) 是 \(\varepsilon\)-正则的。一对 \((V_i, V_j)\) 称为 \(\varepsilon\)-正则的，若对任意 \(A \subseteq V_i, B \subseteq V_j\)（\(|A| \geq \varepsilon|V_i|, |B| \geq \varepsilon|V_j|\)），有 \(|d(A,B) - d(V_i,V_j)| \leq \varepsilon\)，其中 \(d(X,Y) = |E(X,Y)|/(|X||Y|)\) 为边密度。

正则性引理是极值图论和加法组合学中最重要的结构定理。其核心思想是：任意大图都可被近似地视为一个''随机二部图构成的拟随机图''------即正则对的集合。这一''结构-拟随机性''二分法贯穿了现代组合学：若一个图不具有拟随机行为，则它必然包含某种结构化的子图，可通过正则划分的简化图来检测。正则性引理在证明中的典型应用模式是：(1) 对图 \(G\) 应用正则性引理，获得正则划分；(2) 在简化图中应用经典极值定理（如 Turán 定理）；(3) 通过关键引理将简化图中的结构提升为 \(G\) 中的子图。这一范式在 Erdős-Stone 定理和 Green-Tao 定理（素数包含任意长等差数列）的证明中均发挥了核心作用。正则性引理的原始界 \(M(\varepsilon, m)\) 是塔式指数函数（tower-type），Gowers（1997）证明了塔式增长是不可避免的。

\textbf{注记 130.1}（图论与计算机科学的交汇）：图论中的许多问题在计算复杂性理论中具有核心地位。图同构问题（Graph Isomorphism）是少数既未被证明为 NP 完全、也未被证明为多项式时间可解的问题之一------Babai（2015）的拟多项式时间算法 \(n^{O(\log n)}\) 是该领域四十年来最重要的突破。最大独立集和最大团问题是 NP 难的，但通过 Lovász theta 函数 \(\vartheta(G)\)（半定规划松弛）可在多项式时间内给出上界。完美图定理（Lovász, 1972）断言 \(G\) 是完美图当且仅当 \(\overline{G}\) 是完美图，由 Chudnovsky-Robertson-Seymour-Thomas（2006）通过结构分解定理完成分类。图论算法在社交网络分析（PageRank）、推荐系统（二分图匹配）和编译器优化（寄存器分配即图染色）中有广泛应用。

\hrulefill

\newpage

\chapter{04-离散数学/02-概率与信息/01-信息论.md}\label{ux79bbux6563ux6570ux5b6602-ux6982ux7387ux4e0eux4fe1ux606f01-ux4fe1ux606fux8bba.md}

\section{第96章：信息论基础}\label{ux7b2c96ux7ae0ux4fe1ux606fux8bbaux57faux7840}

1948 年，Claude Shannon 在《A Mathematical Theory of Communication》中创立了信息论，将信息从模糊的哲学概念转化为可精确度量的数学量。这一理论的革命性在于：它用一套公理化的概率框架------熵、互信息、信道容量------统一回答了''信息能否被压缩''\,``噪声信道中可靠通信的最大速率''``有损压缩的极限''三个基本问题。信息论不仅是通信工程的理论支柱（JPEG、MP3、WiFi、5G 均以其为设计基础），更深刻地影响了概率论（大偏差理论中的熵方法）、统计学（最大熵原理、变分推断）、计算机科学（Kolmogorov 复杂度、最小描述长度）和物理学（Landauer 原理与热力学熵的等价性）。本章从熵和互信息的公理化定义出发，建立信息度量的完整运算体系------链式法则、数据处理不等式和 KL 散度的 Gibbs 不等式------为后续的信源编码（45.2）、信道编码（45.3）和率失真理论（45.4）奠定统一的数学基础。

\subsection{45.1 熵、互信息与 KL 散度}\label{ux71b5ux4e92ux4fe1ux606fux4e0e-kl-ux6563ux5ea6}

\begin{quote}
\textbf{直觉与动机} 信息论由 Claude Shannon 于 1948 年在《通信的数学理论》中创立，为信息的度量、存储和传输提供了严格的数学框架。信息论的核心问题------如何高效、可靠地传输信息------涉及不确定性（熵）、信息量（互信息）和最佳编码（编码定理）等基本概念，这些概念不仅在通信工程中至关重要，也在统计力学、机器学习、量子计算和网络科学中产生了深远影响。
\end{quote}

\subsubsection{45.1.1 熵的定义与基本性质}\label{ux71b5ux7684ux5b9aux4e49ux4e0eux57faux672cux6027ux8d28}

\textbf{定义 45.1}（离散随机变量的熵）：设离散随机变量 \(X\) 取值于有限集 \(\mathcal{X}\)，概率分布为 \(p(x) = P(X = x)\)。\(X\) 的\textbf{熵}（Shannon 熵）定义为 \[H(X) = -\sum_{x \in \mathcal{X}} p(x) \log p(x)\] 其中 \(\log\) 以 \(2\) 为底时单位为比特（bit），以 \(e\) 为底时单位为奈特（nat）。约定 \(0 \log 0 = 0\)（由 \(\lim_{t \to 0^+} t \log t = 0\) 保证连续性）。

\textbf{定理 45.1}（熵的基本性质）：设 \(X\) 为离散随机变量，\(|\mathcal{X}| = n\)。则： 1. \(H(X) \geq 0\)，且 \(H(X) = 0\) 当且仅当 \(X\) 为退化分布（即存在 \(x_0\) 使 \(p(x_0) = 1\)）； 2. \(H(X) \leq \log n\)，等号成立当且仅当 \(X\) 服从均匀分布 \(p(x) = 1/n\)； 3. \(H(X)\) 是 \(p\) 的凹函数。

\textbf{证明}：(1) \(0 \leq p(x) \leq 1\)，故 \(-\log p(x) \geq 0\)，\(H(X) \geq 0\)。若 \(H(X) = 0\)，则每一项 \(-p(x)\log p(x) = 0\)，即对每个 \(x\) 有 \(p(x) = 0\) 或 \(p(x) = 1\)。由 \(\sum p(x) = 1\) 知恰有一个 \(x\) 使 \(p(x) = 1\)。

\begin{enumerate}
\def\labelenumi{(\arabic{enumi})}
\setcounter{enumi}{1}
\item
  利用 Jensen 不等式（\(\log\) 为严格凹函数）： \[H(X) - \log n = \sum_x p(x) \log \frac{1}{p(x)} - \log n = \sum_x p(x) \log \frac{1}{n p(x)} \leq \log \sum_x p(x) \cdot \frac{1}{n p(x)} = \log 1 = 0\] 等号成立当且仅当 \(\frac{1}{n p(x)}\) 为常数，即 \(p(x) = 1/n\)。
\item
  对 \(p, q\) 分布和 \(\lambda \in [0,1]\)，\(H(\lambda p + (1-\lambda)q) \geq \lambda H(p) + (1-\lambda)H(q)\)（由 \(t \mapsto -t \log t\) 的凹性）。\(\blacksquare\)
\end{enumerate}

\textbf{定义 45.2}（联合熵与条件熵）：设 \((X, Y)\) 的联合分布为 \(p(x,y)\)。\textbf{联合熵} \(H(X,Y) = -\sum_{x,y} p(x,y) \log p(x,y)\)。\textbf{条件熵} \(H(Y|X) = \sum_x p(x) H(Y|X=x) = -\sum_{x,y} p(x,y) \log p(y|x)\)。

\textbf{定理 45.2}（链式法则）：\(H(X,Y) = H(X) + H(Y|X)\)。

\textbf{证明}： \[H(X,Y) = -\sum_{x,y} p(x,y) \log p(x,y) = -\sum_{x,y} p(x,y) \log [p(x) p(y|x)]\] \[= -\sum_{x,y} p(x,y) \log p(x) - \sum_{x,y} p(x,y) \log p(y|x) = -\sum_x p(x) \log p(x) + H(Y|X) = H(X) + H(Y|X)\] \(\blacksquare\)

\subsubsection{45.1.2 互信息与 KL 散度}\label{ux4e92ux4fe1ux606fux4e0e-kl-ux6563ux5ea6}

\textbf{定义 45.3}（互信息）：\(X\) 与 \(Y\) 之间的\textbf{互信息}定义为 \[I(X; Y) = \sum_{x,y} p(x,y) \log \frac{p(x,y)}{p(x)p(y)} = D_{KL}(p(x,y) \| p(x)p(y))\]

\textbf{定理 45.3}（互信息的基本性质）： 1. \(I(X; Y) = H(X) - H(X|Y) = H(Y) - H(Y|X) = H(X) + H(Y) - H(X,Y)\)； 2. \(I(X; Y) \geq 0\)，等号成立当且仅当 \(X\) 与 \(Y\) 独立； 3. \(I(X; Y) = I(Y; X)\)（对称性）。

\textbf{证明}：(1) \(I(X;Y) = \sum p(x,y) \log \frac{p(x|y)}{p(x)} = -\sum p(x,y) \log p(x) + \sum p(x,y) \log p(x|y) = H(X) - H(X|Y)\)。其余类似。(2) 由 KL 散度的非负性（定理 45.4）立即得证。(3) 由定义对称。\(\blacksquare\)

\textbf{定义 45.4}（Kullback-Leibler 散度）：两个概率分布 \(p\) 和 \(q\)（定义在同一集合 \(\mathcal{X}\) 上）之间的 \textbf{KL 散度}（相对熵）为 \[D_{KL}(p \| q) = \sum_{x \in \mathcal{X}} p(x) \log \frac{p(x)}{q(x)}\] 约定 \(0 \log \frac{0}{q} = 0\)，\(p \log \frac{p}{0} = \infty\)（\(p > 0\) 时）。

\textbf{定理 45.4}（Gibbs 不等式 / KL 散度的非负性）：\(D_{KL}(p \| q) \geq 0\)，等号成立当且仅当 \(p = q\)（几乎处处）。

\textbf{证明}：由 Jensen 不等式： \[-D_{KL}(p \| q) = \sum_x p(x) \log \frac{q(x)}{p(x)} \leq \log \sum_x p(x) \cdot \frac{q(x)}{p(x)} = \log \sum_x q(x) = \log 1 = 0\] 故 \(D_{KL}(p \| q) \geq 0\)。等号成立当且仅当 \(\frac{q(x)}{p(x)}\) 为常数（对所有 \(p(x) > 0\) 的 \(x\)），即 \(p = q\)。\(\blacksquare\)

\textbf{注记 45.1}：KL 散度不是真正的度量------它不满足对称性（\(D_{KL}(p\|q) \neq D_{KL}(q\|p)\) 一般成立）和三角不等式。然而，它衡量了用 \(q\) 近似 \(p\) 时的信息损失，是信息几何中 Fisher 信息度量的积分形式，也是机器学习中交叉熵损失函数和变分推断的理论基础。

\textbf{定理 45.5}（数据处理不等式）：若 \(X \to Y \to Z\) 构成 Markov 链，则 \(I(X; Z) \leq I(X; Y)\)。即信息处理不能增加信息量。

\textbf{证明}：由链式法则，\(I(X; Y, Z) = I(X; Z) + I(X; Y|Z) = I(X; Y) + I(X; Z|Y)\)。因 \(X \to Y \to Z\) 为 Markov 链，\(X\) 与 \(Z\) 在给定 \(Y\) 时条件独立，故 \(I(X; Z|Y) = 0\)。由互信息的非负性，\(I(X; Y|Z) \geq 0\)，从而 \(I(X; Z) \leq I(X; Y)\)。\(\blacksquare\)

\subsection{45.2 信源编码定理}\label{ux4fe1ux6e90ux7f16ux7801ux5b9aux7406-1}

信源编码定理（Shannon 第一定理，1948 年）是信息论的第一块基石，它确立了熵作为数据压缩绝对下限的地位：在无失真条件下，信源的每个符号至少需要 \(H(X)\) 比特来描述，且存在编码方案渐近地达到这一极限。这一定理的核心工具是渐近均分性（AEP）------大数定律在信息论中的化身。AEP 揭示了一个非直观的事实：虽然 \(n\) 次独立信源输出的可能序列多达 \(|\mathcal{X}|^n\) 个，但概率质量几乎全部集中在约 \(2^{nH(X)}\) 个''典型''序列上；当 \(H(X) < \log|\mathcal{X}|\) 时典型序列的占比趋于零，从而为数据压缩提供了理论依据------只需对典型序列编码，非典型序列可忽略。本节从典型序列和 AEP 的精确定义出发，证明 Shannon 信源编码定理的正反两部分，并讨论 Huffman 编码和算术编码等最优前缀码的构造。

\subsubsection{45.2.1 典型序列与 AEP}\label{ux5178ux578bux5e8fux5217ux4e0e-aep}

\textbf{定义 45.5}（典型序列）：设 \(X_1, X_2, \ldots, X_n\) 为独立同分布于 \(p(x)\) 的随机变量序列。对 \(\varepsilon > 0\)，\textbf{典型集}（typical set）\(A_\varepsilon^{(n)}\) 定义为满足以下条件的序列 \((x_1, \ldots, x_n)\) 的集合： \[\left| -\frac{1}{n} \log p(x_1, \ldots, x_n) - H(X) \right| \leq \varepsilon\] 即 \(2^{-n(H(X)+\varepsilon)} \leq p(x_1, \ldots, x_n) \leq 2^{-n(H(X)-\varepsilon)}\)。

\textbf{定理 45.6}（渐近等分性 / AEP）：对独立同分布序列 \(X_1, \ldots, X_n \sim p(x)\)， 1. \(\lim_{n \to \infty} P(A_\varepsilon^{(n)}) = 1\)（典型集占据几乎全部概率）； 2. \((1-\varepsilon)2^{n(H(X)-\varepsilon)} \leq |A_\varepsilon^{(n)}| \leq 2^{n(H(X)+\varepsilon)}\)（典型集大小约为 \(2^{nH(X)}\)）； 3. 对充分大的 \(n\)，\(|A_\varepsilon^{(n)}| \approx 2^{nH(X)}\)。

\textbf{证明}：(1) 由弱大数定律，\(-\frac{1}{n} \log p(X_1, \ldots, X_n) = -\frac{1}{n} \sum_{i=1}^n \log p(X_i) \xrightarrow{P} -E[\log p(X)] = H(X)\)。故对任意 \(\varepsilon > 0\)，\(P(|-\frac{1}{n}\log p(X^n) - H(X)| > \varepsilon) \to 0\)。

\begin{enumerate}
\def\labelenumi{(\arabic{enumi})}
\setcounter{enumi}{1}
\tightlist
\item
  由 \(1 = \sum_{x^n} p(x^n) \geq \sum_{x^n \in A_\varepsilon^{(n)}} p(x^n) \geq |A_\varepsilon^{(n)}| \cdot 2^{-n(H(X)+\varepsilon)}\) 得 \(|A_\varepsilon^{(n)}| \leq 2^{n(H(X)+\varepsilon)}\)。对充分大的 \(n\)，\(1-\varepsilon \leq P(A_\varepsilon^{(n)}) = \sum_{x^n \in A_\varepsilon^{(n)}} p(x^n) \leq |A_\varepsilon^{(n)}| \cdot 2^{-n(H(X)-\varepsilon)}\) 得 \(|A_\varepsilon^{(n)}| \geq (1-\varepsilon)2^{n(H(X)-\varepsilon)}\)。\(\blacksquare\)
\end{enumerate}

\textbf{注记 45.2}：AEP 揭示了一个深刻的事实：虽然长度为 \(n\) 的可能序列有 \(|\mathcal{X}|^n\) 个，但概率质量几乎全部集中在约 \(2^{nH(X)}\) 个典型序列上。当 \(H(X) < \log |\mathcal{X}|\) 时，典型序列的比例 \(2^{n(H(X)-\log|\mathcal{X}|)} \to 0\)，但每个典型序列的''典型度''约为 \(2^{-nH(X)}\)。这为压缩编码提供了理论依据：只需为典型序列分配码字，非典型序列可忽略（错误概率趋于零）。

\subsubsection{45.2.2 Shannon 信源编码定理}\label{shannon-ux4fe1ux6e90ux7f16ux7801ux5b9aux7406}

\textbf{定义 45.6}（无失真信源编码）：设 \(\mathcal{X}\) 为信源字母表，\(D\) 为码字母表（\(|\mathcal{D}| = D\)）。\(n\) 长分组码是映射 \(f_n : \mathcal{X}^n \to \mathcal{D}^*\)（变长码）或 \(f_n : \mathcal{X}^n \to \mathcal{D}^{\ell_n}\)（定长码）。码率 \(R\) 为每信源符号的平均码符号数。

\textbf{定理 45.7}（Shannon 无失真信源编码定理，1948）：设离散无记忆信源熵为 \(H\)（以 \(D\) 为底的对数）。则： 1. （可达性）对任意 \(R > H\)，存在分组码使码率 \(\leq R\) 且错误概率 \(P_e^{(n)} \to 0\)（\(n \to \infty\)）； 2. （逆定理）对任意 \(R < H\)，任何分组码的错误概率 \(P_e^{(n)} \not\to 0\)。

\textbf{证明}（可达性）：取 \(R > H\) 和 \(\varepsilon = (R-H)/2 > 0\)。对每个 \(n\)，为 \(A_\varepsilon^{(n)}\) 中每个序列分配一个码字（总码字数 \(\leq 2^{n(H+\varepsilon)} \leq 2^{nR}\)，因此定长编码可行）。发送时若序列在 \(A_\varepsilon^{(n)}\) 中则编码发送，否则发送任意码字。解码错误仅当序列不在典型集中，由 AEP：\(P_e^{(n)} \leq P((A_\varepsilon^{(n)})^c) \to 0\)。

\textbf{证明}（逆定理）：设码率为 \(R < H\)。对任意编码方案，码字总数 \(\leq 2^{nR}\)。将 \(A_\varepsilon^{(n)}\) 划分为正确编码的序列（\(B\)）和错误编码的序列。由典型集大小 \(|A_\varepsilon^{(n)}| \geq (1-\varepsilon)2^{n(H-\varepsilon)}\)，而 \(|B| \leq 2^{nR}\)。故 \(P_e^{(n)} \geq P(A_\varepsilon^{(n)}) - P(B) \geq (1-\varepsilon) - 2^{nR} \cdot 2^{-n(H-\varepsilon)} = (1-\varepsilon) - 2^{-n(H-R-\varepsilon)} \to 1-\varepsilon > 0\)（当 \(H-R-\varepsilon > 0\) 时）。\(\blacksquare\)

\textbf{注记 45.3}：信源编码定理是信息论的第一基本定理。它表明信源熵 \(H\) 是压缩的绝对下限------任何码率低于 \(H\) 的编码必然有不可忽略的错误概率。最优码（如 Huffman 码）的平均码长 \(L\) 满足 \(H \leq L < H+1\)，Shannon-Fano-Elias 码则为 \(H \leq L < H+2\)。算术编码通过分数比特表示逼近熵极限，是现代压缩标准（JPEG、H.264）的核心技术。

\subsection{45.3 信道容量}\label{ux4fe1ux9053ux5bb9ux91cf-1}

信道容量 \(C = \max_{p(x)} I(X; Y)\) 是链接信息度量与通信极限的核心概念------它将物理信道的传输能力量化为互信息的最大值，从而将信道设计与信息度量统一在同一数学框架下。这一概念的深邃之处在于其''最小最大''对偶结构：\(C\) 是输入分布 \(p(x)\) 的最大化问题的解（最大化互信息），又是互信息的凹函数（便于凸优化求解），而达到容量的输入分布 \(p^*(x)\) 由信道的对称性和 Kuhn-Tucker 条件刻画。经典信道的容量公式是信息论中最优雅的闭式结果：二元对称信道 \(C_{\text{BSC}} = 1-H(p)\)，二元擦除信道 \(C_{\text{BEC}} = 1-\alpha\)，加性高斯白噪声信道 \(C_{\text{AWGN}} = \frac{1}{2}\log_2(1 + P/\sigma^2)\)（Shannon 公式）。本节在定义信道模型和容量后，以随机编码和联合典型译码严格证明 Shannon 信道编码定理------信息论的第二基本定理。

\subsubsection{45.3.1 离散无记忆信道}\label{ux79bbux6563ux65e0ux8bb0ux5fc6ux4fe1ux9053}

\textbf{定义 45.7}（离散无记忆信道）：\textbf{离散无记忆信道}（DMC）由输入字母表 \(\mathcal{X}\)、输出字母表 \(\mathcal{Y}\) 和转移概率矩阵 \(p(y|x)\)（\(x \in \mathcal{X}, y \in \mathcal{Y}\)）组成。信道对每个输入 \(x\) 独立地产生分布为 \(p(\cdot|x)\) 的输出。\(n\) 次使用时，\(p(y^n|x^n) = \prod_{i=1}^n p(y_i|x_i)\)。

\textbf{定义 45.8}（信道容量）：信道 \(p(y|x)\) 的\textbf{容量}定义为输入分布与输出分布之间的最大互信息： \[C = \max_{p(x)} I(X; Y) = \max_{p(x)} \sum_{x,y} p(x) p(y|x) \log \frac{p(y|x)}{\sum_{x'} p(x') p(y|x')}\]

\textbf{定理 45.8}（信道容量的基本性质）： 1. \(C \geq 0\)（由互信息的非负性）； 2. \(C \leq \min\{\log |\mathcal{X}|, \log |\mathcal{Y}|\}\)（\(I(X;Y) \leq H(X) \leq \log|\mathcal{X}|\)，同理 \(I(X;Y) \leq \log|\mathcal{Y}|\)）； 3. \(I(X; Y)\) 作为 \(p(x)\) 的函数是凹函数（因此求最大值是凸优化问题）。

\textbf{证明}（性质 3）：固定 \(p(y|x)\)，\(I(X;Y) = H(Y) - H(Y|X)\)。\(H(Y|X) = \sum_x p(x) H(Y|X=x)\) 是 \(p(x)\) 的线性函数。\(H(Y)\) 是 \(p(y) = \sum_x p(x) p(y|x)\) 的凹函数，而 \(p(y)\) 是 \(p(x)\) 的线性函数。凹函数的线性复合仍为凹函数，故 \(I(X;Y)\) 为 \(p(x)\) 的凹函数。\(\blacksquare\)

\subsubsection{45.3.2 基本信道示例}\label{ux57faux672cux4fe1ux9053ux793aux4f8b}

\textbf{例 45.1}（二进制对称信道 / BSC）：\(\mathcal{X} = \mathcal{Y} = \{0,1\}\)，转移概率 \(p(0|1) = p(1|0) = p\)（交叉概率），\(p(0|0) = p(1|1) = 1-p\)。容量 \(C_{\text{BSC}} = 1 - H(p)\)（比特/信道使用），其中 \(H(p) = -p\log p - (1-p)\log(1-p)\)。最优输入分布为均匀分布 \(p(0) = p(1) = 1/2\)。

\textbf{验证}：\(I(X;Y) = H(Y) - H(Y|X) = H(Y) - H(p)\)。当 \(X\) 均匀时，\(Y\) 也均匀（\(H(Y) = 1\)），故 \(I(X;Y) = 1 - H(p)\)。由对称性，\(I(X;Y)\) 在 \(X\) 均匀时达到最大。

\textbf{例 45.2}（二进制擦除信道 / BEC）：\(\mathcal{X} = \{0,1\}\)，\(\mathcal{Y} = \{0, e, 1\}\)（\(e\) 表示擦除）。转移概率 \(p(e|0) = p(e|1) = \alpha\)（擦除概率），\(p(0|0) = p(1|1) = 1-\alpha\)。容量 \(C_{\text{BEC}} = 1 - \alpha\)。最优输入分布为均匀分布。

\textbf{验证}：\(I(X;Y) = H(Y) - H(Y|X) = H(Y) - H(\alpha)\)。当 \(X\) 均匀时，\(p(y=0) = p(y=1) = (1-\alpha)/2\)，\(p(y=e) = \alpha\)。此时 \(H(Y) = H(\alpha) + (1-\alpha)\)，\(I(X;Y) = 1-\alpha\)。

\subsubsection{45.3.3 Shannon 信道编码定理}\label{shannon-ux4fe1ux9053ux7f16ux7801ux5b9aux7406}

\textbf{定义 45.9}（信道码）：\((2^{nR}, n)\) 信道码由编码函数 \(f : \{1,\ldots,2^{nR}\} \to \mathcal{X}^n\) 和解码函数 \(g : \mathcal{Y}^n \to \{1,\ldots,2^{nR}\}\) 组成。\(R\) 为码率（比特/信道使用）。最大错误概率 \(P_e^{(n)} = \max_{m} P(g(Y^n) \neq m \mid X^n = f(m))\)。

\textbf{定理 45.9}（Shannon 信道编码定理，1948）：设信道容量为 \(C\)。则： 1. （可达性）对任意 \(R < C\)，存在码序列使 \(P_e^{(n)} \to 0\)（\(n \to \infty\)）； 2. （逆定理）对任意 \(R > C\)，任何码序列的 \(P_e^{(n)}\) 有正的下界，即不可能以任意小的错误概率可靠传输。

\textbf{证明}（可达性，随机编码）：取 \(R < C\)。固定输入分布 \(p(x)\) 使 \(I(X;Y) = C\)。随机生成 \(2^{nR}\) 个码字 \(X^n(m)\)（\(m = 1,\ldots,2^{nR}\)），每个元素独立同分布于 \(p(x)\)。解码时采用\textbf{联合典型解码}：收到 \(y^n\) 后，寻找唯一的 \(m\) 使 \((X^n(m), y^n) \in A_\varepsilon^{(n)}\)（联合典型集）。

\textbf{错误分析}：设发送消息 \(m=1\)。错误事件：(1) \((X^n(1), Y^n)\) 不联合典型------概率 \(\to 0\)（由联合 AEP）；(2) 存在 \(m' \neq 1\) 使 \((X^n(m'), Y^n)\) 联合典型。

对 \(m' \neq 1\)，\(X^n(m')\) 与 \(Y^n\) 独立，\(P((X^n(m'), Y^n) \in A_\varepsilon^{(n)}) \leq 2^{-n(I(X;Y)-3\varepsilon)}\)。由并界： \[P(\text{错误}) \leq \varepsilon + 2^{nR} \cdot 2^{-n(C-3\varepsilon)} = \varepsilon + 2^{-n(C-R-3\varepsilon)} \to 0\] 当 \(R < C\) 且选取 \(\varepsilon\) 充分小使 \(C - R - 3\varepsilon > 0\) 时。故存在一组码使错误概率任意小。

\textbf{证明}（逆定理，Fano 不等式）：设 \(W\) 为均匀分布在 \(\{1,\ldots,2^{nR}\}\) 上的消息，\(\hat{W}\) 为解码输出。由 Fano 不等式： \[H(W|\hat{W}) \leq 1 + P_e^{(n)} nR\] 由数据处理不等式：\(I(W; \hat{W}) \leq I(X^n; Y^n) \leq nC\)（DMC 的 \(n\) 次使用容量为 \(nC\)）。故： \[nR = H(W) = I(W; \hat{W}) + H(W|\hat{W}) \leq nC + 1 + P_e^{(n)} nR\] 整理得 \(P_e^{(n)} \geq 1 - \frac{C}{R} - \frac{1}{nR}\)。当 \(R > C\) 且 \(n\) 充分大时，\(P_e^{(n)} \geq 1 - C/R > 0\)。\(\blacksquare\)

\textbf{注记 45.4}：信道编码定理是信息论的第二基本定理，与信源编码定理共同构成 Shannon 理论的基石。信道容量 \(C\) 是可靠通信的速率上限------任何码率超过 \(C\) 的传输必然有不可忽略的错误概率。随机编码的证明不构造具体的好码，仅证明好码存在，激励了后续 50 年对构造性编码（Reed-Solomon 码、卷积码、Turbo 码、LDPC 码、Polar 码）的研究。Polar 码（Arikan, 2009）是首个被严格证明在二进制输入对称信道下达到 Shannon 容量的构造性编码方案。

\subsection{45.4 率失真理论}\label{ux7387ux5931ux771fux7406ux8bba-1}

率失真理论（Shannon 第三定理，1959 年）处理的是与信源编码完全对偶的问题：有损压缩的极限。当允许一定程度的失真（distortion）时，可以将信源压至低于熵的码率，且率失真函数 \(R(D)\) 精确刻画了在给定失真容限 \(D\) 下可达的最小码率。与 Shannon 第二定理的随机编码证明并行，率失真定理的证明将互信息的最小化问题（在失真约束下）与随机码本和联合典型译码相结合，展现了信息论中''最大化''与''最小化''操作的精妙对称。率失真理论具有广泛的实际应用：高斯情形 \(R(D) = \frac{1}{2}\log_2(\sigma^2/D)\) 直接指导了变换编码（JPEG 量化表设计）、语音编码（CELP）和视频编码（H.264/HEVC 的率失真优化模式决策）中的速率-质量权衡。

\subsubsection{45.4.1 失真度量与率失真函数}\label{ux5931ux771fux5ea6ux91cfux4e0eux7387ux5931ux771fux51fdux6570}

\textbf{定义 45.10}（失真度量）：设信源字母表为 \(\mathcal{X}\)，重构字母表为 \(\hat{\mathcal{X}}\)。\textbf{失真度量}（distortion measure）\(d : \mathcal{X} \times \hat{\mathcal{X}} \to [0, \infty)\) 衡量用 \(\hat{x}\) 表示 \(x\) 的代价。常见失真度量包括汉明失真 \(d(x, \hat{x}) = \mathbf{1}_{\{x \neq \hat{x}\}}\)（离散情形）和平方误差失真 \(d(x, \hat{x}) = (x - \hat{x})^2\)（连续情形）。序列的失真定义为 \(d(x^n, \hat{x}^n) = \frac{1}{n}\sum_{i=1}^n d(x_i, \hat{x}_i)\)。

\textbf{定义 45.11}（率失真码）：\((2^{nR}, n)\) 率失真码由编码函数 \(f_n : \mathcal{X}^n \to \{1, \ldots, 2^{nR}\}\) 和解码函数 \(g_n : \{1, \ldots, 2^{nR}\} \to \hat{\mathcal{X}}^n\) 组成。码率 \(R\) 为每信源符号的比特数。期望失真为 \(E[d(X^n, g_n(f_n(X^n)))]\)。

\textbf{定义 45.12}（率失真函数）：给定失真容限 \(D \geq 0\)，\textbf{率失真函数}（信息率失真函数）定义为 \[R(D) = \min_{p(\hat{x}|x) : E[d(X,\hat{X})] \leq D} I(X; \hat{X})\] 其中互信息在满足期望失真约束的所有条件分布 \(p(\hat{x}|x)\) 上取最小值。

\textbf{定理 45.10}（率失真函数的基本性质）：设 \(D_{\max} = \min_{\hat{x}} E[d(X, \hat{x})]\)（无信息时的最小失真）。则： 1. \(R(D)\) 是非增的凸函数（\(D \geq 0\)）； 2. \(R(0) \leq H(X)\)（零失真时码率不超过熵）； 3. \(R(D) = 0\) 当 \(D \geq D_{\max}\)。

\textbf{证明}：(1) 凸性：设 \(R(D_1)\) 和 \(R(D_2)\) 分别由最优分布 \(p_1(\hat{x}|x)\) 和 \(p_2(\hat{x}|x)\) 达到。对 \(\lambda \in [0,1]\)，考虑混合分布 \(p_\lambda = \lambda p_1 + (1-\lambda)p_2\)。由互信息的凸性（在固定 \(p(x)\) 下 \(I(X;\hat{X})\) 是 \(p(\hat{x}|x)\) 的凸函数），\(I_{p_\lambda}(X;\hat{X}) \leq \lambda I_{p_1} + (1-\lambda)I_{p_2} = \lambda R(D_1) + (1-\lambda)R(D_2)\)。且 \(E_{p_\lambda}[d] = \lambda D_1 + (1-\lambda)D_2\)。故 \(R(\lambda D_1 + (1-\lambda)D_2) \leq \lambda R(D_1) + (1-\lambda)R(D_2)\)。

\begin{enumerate}
\def\labelenumi{(\arabic{enumi})}
\setcounter{enumi}{1}
\item
  \(D=0\) 时要求 \(X = \hat{X}\) 几乎必然，\(I(X;\hat{X}) = H(X)\)，故 \(R(0) \leq H(X)\)。
\item
  \(D \geq D_{\max}\) 时，取 \(p(\hat{x}|x) = p(\hat{x})\)（\(\hat{X}\) 与 \(X\) 独立），\(I(X;\hat{X}) = 0\)，故 \(R(D) = 0\)。\(\blacksquare\)
\end{enumerate}

\subsubsection{45.4.2 Shannon 率失真定理}\label{shannon-ux7387ux5931ux771fux5b9aux7406}

\textbf{定理 45.11}（Shannon 率失真定理，1959）：设 \(R(D)\) 为率失真函数。则： 1. （可达性）对任意 \(D \geq 0\) 和 \(\varepsilon > 0\)，存在码率为 \(R \leq R(D) + \varepsilon\) 的码使期望失真 \(\leq D + \varepsilon\)； 2. （逆定理）任何码率为 \(R < R(D)\) 的码必然有期望失真 \(> D\)。

\textbf{证明}（可达性框架）：固定 \(p(\hat{x}|x)\) 使 \(I(X; \hat{X}) = R(D)\) 且 \(E[d(X, \hat{X})] \leq D\)。随机生成 \(2^{nR}\) 个重构码字 \(\hat{X}^n(w)\)，每个独立同分布于 \(p(\hat{x})\)。编码时，给定信源序列 \(x^n\)，寻找 \(w\) 使 \((x^n, \hat{X}^n(w))\) 联合典型。由联合 AEP，此类 \(w\) 以高概率存在当 \(R > I(X; \hat{X})\)。失真由典型序列的性质保证 \(\leq D + \varepsilon\)。

\textbf{证明}（逆定理）：对任意码率为 \(R\) 的码，由数据处理不等式和互信息的链式法则： \[nR \geq H(f_n(X^n)) \geq I(X^n; \hat{X}^n) = \sum_{i=1}^n [H(X_i) - H(X_i|\hat{X}^n, X^{i-1})]\] \[\geq \sum_{i=1}^n [H(X_i) - H(X_i|\hat{X}_i)] = \sum_{i=1}^n I(X_i; \hat{X}_i) \geq \sum_{i=1}^n R(E[d(X_i, \hat{X}_i)])\] 由 \(R(D)\) 的凸性：\(R \geq \frac{1}{n}\sum R(E[d_i]) \geq R(\frac{1}{n}\sum E[d_i]) = R(E[d])\)。故 \(E[d] \geq D\) 当 \(R < R(D)\)。\(\blacksquare\)

\textbf{注记 45.5}：率失真理论是信息论的第三基本定理，奠定了有损压缩的理论基础。对于高斯信源 \(X \sim N(0, \sigma^2)\) 和平方误差失真，\(R(D) = \frac{1}{2}\log_2(\sigma^2/D)\)（\(0 \leq D \leq \sigma^2\)）。这一结果直接指导了变换编码（JPEG、MPEG 等）的量化步长设计。Blahut-Arimoto 算法（1972）提供了计算离散信道容量和率失真函数的迭代数值方法，其收敛性由交替优化的凸性保证。

\subsection{30.5 估计理论与假设检验导引}\label{ux4f30ux8ba1ux7406ux8bbaux4e0eux5047ux8bbeux68c0ux9a8cux5bfcux5f15}

\textbf{定义 30.8}（点估计与无偏性）：设 \(X_1, \ldots, X_n\) 为来自总体 \(F_\theta\) 的样本，\(\theta \in \Theta\) 为未知参数。统计量 \(\hat{\theta}_n = \hat{\theta}(X_1, \ldots, X_n)\) 称为 \(\theta\) 的\textbf{点估计}。若 \(E_\theta(\hat{\theta}_n) = \theta\) 对所有 \(\theta \in \Theta\) 成立，则称 \(\hat{\theta}_n\) 为\textbf{无偏估计}。估计量的优劣由\textbf{均方误差} \(MSE(\hat{\theta}_n) = E_\theta[(\hat{\theta}_n - \theta)^2] = \operatorname{Var}_\theta(\hat{\theta}_n) + [\operatorname{Bias}_\theta(\hat{\theta}_n)]^2\) 衡量，无偏时 \(MSE = \operatorname{Var}\)。

\textbf{定义 30.9}（最大似然估计）：给定观测值 \(x_1, \ldots, x_n\)，\textbf{似然函数} \(L(\theta) = \prod_{i=1}^n f_\theta(x_i)\)（\(f_\theta\) 为密度函数或概率质量函数）。\textbf{最大似然估计}（MLE）\(\hat{\theta}_{\text{MLE}}\) 为最大化 \(L(\theta)\) 的 \(\theta\) 值。MLE 具有渐近有效性：在正则条件下，\(\sqrt{n}(\hat{\theta}_{\text{MLE}} - \theta) \xrightarrow{d} N(0, I(\theta)^{-1})\)，其中 \(I(\theta)\) 为 Fisher 信息量 \(I(\theta) = E_\theta[(\frac{\partial}{\partial\theta} \log f_\theta(X))^2]\)。Cramér-Rao 不等式给出无偏估计方差的下界：\(\operatorname{Var}_\theta(\hat{\theta}) \geq 1 / (n I(\theta))\)。

\textbf{定义 30.10}（假设检验的基本框架）：\textbf{原假设} \(H_0\) 与\textbf{备择假设} \(H_1\) 构成一对统计假设。\textbf{第一类错误}（拒真）的概率为 \(\alpha = P(\text{拒绝 } H_0 \mid H_0 \text{ 真})\)（显著性水平），\textbf{第二类错误}（取伪）的概率为 \(\beta = P(\text{接受 } H_0 \mid H_1 \text{ 真})\)。检验的功效（power）为 \(1 - \beta\)。Neyman-Pearson 引理（1933）给出了简单假设的最优检验：似然比检验 \(\Lambda(x) = L(\theta_1)/L(\theta_0) > k\) 在固定 \(\alpha\) 下最大化功效。Wald 的序贯概率比检验（SPRT, 1945）将这一框架推广到序贯抽样情形，在达到相同功效的前提下平均所需的样本量更少，是二战期间军需品抽样检验的数学基础。

\textbf{注记 30.1}（统计推断三大范式）：\textbf{频率学派}（frequentist）以 Neyman-Pearson 假设检验和 Fisher 的显著性检验为代表，将概率解释为长期频率，参数是固定的未知常数。\textbf{贝叶斯学派}（Bayesian）将参数视为随机变量，通过先验分布 \(\pi(\theta)\) 表达先验知识，利用贝叶斯公式得到后验分布 \(\pi(\theta|x) \propto L(\theta)\pi(\theta)\)，所有推断基于后验分布。\textbf{Fisher 的似然学派}（likelihoodist）仅使用似然函数进行推断，回避了先验分布的选择和频率解释。三种范式在哲学基础上存在根本分歧（Lindley 悖论等），但在大样本下渐近等价。现代统计实践倾向于根据问题背景灵活选择范式，并以稳健性分析（如先验敏感性分析）补充主要结论。

\subsection{41.B 网络信息论}\label{b-ux7f51ux7edcux4fe1ux606fux8bba}

经典 Shannon 信息论（1948 年）解决的是点对点通信的极限问题，而网络信息论将这一框架推广到多用户场景------多址信道（MAC）、广播信道（BC）、中继信道和干扰信道。这一推广绝非简单的技术延伸，因为多用户间的信息流交互引入了点对点模型中不存在的新现象：多址信道中的叠加编码（superposition coding）使得多个发送者的信号可以在接收端联合解码，广播信道中的退化（degradation）和叠加顺序决定了容量区域的多维结构，而干扰信道中的干扰对齐（interference alignment）则揭示了在高 SNR 下自由度可以''切开''而非''共享''。网络信息论的核心挑战在于：大多数多用户信道的容量区域至今仍是开放问题（如一般干扰信道），其困难源于容量区域涉及多维度量空间上的最大化与最小化操作的非交换性。本章介绍多址信道、广播信道和中继信道的容量定理及网络编码的基本原理。

\subsection{45.5 网络信息论导引}\label{ux7f51ux7edcux4fe1ux606fux8bbaux5bfcux5f15}

\begin{quote}
\textbf{直觉与动机} 经典信息论（Shannon 1948）研究的是点对点通信（单信源、单信道、单目的地）。网络信息论将 Shannon 的理论框架推广到多用户场景------多址信道、广播信道、中继信道和干扰信道。网络信息论的核心挑战在于：多个用户同时通信时，信息流之间的交互产生了经典理论中不存在的新现象------叠加编码、干扰对齐、网络编码和协作增益。
\end{quote}

\textbf{定义 45.13}（多址信道 / MAC）：\textbf{多址信道} \((\mathcal{X}_1 \times \mathcal{X}_2, p(y|x_1, x_2), \mathcal{Y})\) 有两个发送者（分别发送 \(X_1\) 和 \(X_2\)）和一个接收者（收到 \(Y\)）。容量区域 \(\mathcal{C}_{\text{MAC}}\) 是所有可达速率对 \((R_1, R_2)\) 的闭包。

\textbf{定理 45.12}（多址信道容量区域，Ahlswede 1971, Liao 1972）：两用户多址信道的容量区域为 \[\mathcal{C}_{\text{MAC}} = \bigcup_{p(x_1)p(x_2)} \left\{ (R_1, R_2) : \begin{aligned} R_1 &\leq I(X_1; Y \mid X_2) \\ R_2 &\leq I(X_2; Y \mid X_1) \\ R_1 + R_2 &\leq I(X_1, X_2; Y) \end{aligned} \right\}\]

\textbf{证明}（框架）：可达性通过叠加编码和逐次干扰消除实现。发送者 1 和 2 分别以速率 \(R_1\) 和 \(R_2\) 独立编码。接收者解码时，先解码一个用户（如用户 2），将已解码信号从接收信号中减去，再解码用户 1。条件 \(R_2 \leq I(X_2; Y|X_1)\) 保证第一个解码成功，\(R_1 \leq I(X_1; Y)\) 保证第二个解码成功。逆定理由 Fano 不等式和互信息的链式法则导出。\(\blacksquare\)

\textbf{定义 45.14}（广播信道 / BC）：\textbf{广播信道} \((\mathcal{X}, p(y_1, y_2|x), \mathcal{Y}_1 \times \mathcal{Y}_2)\) 有一个发送者（发送 \(X\)）和两个接收者（分别收到 \(Y_1\) 和 \(Y_2\)）。发送者希望向接收者 1 发送私有消息 \(W_1\)（速率 \(R_1\)），向接收者 2 发送私有消息 \(W_2\)（速率 \(R_2\)），以及向两者发送公共消息 \(W_0\)（速率 \(R_0\)）。

\textbf{定理 45.13}（退化广播信道容量，Cover 1972, Bergmans 1973）：对退化广播信道 \(X \to Y_1 \to Y_2\)（Markov 链），容量区域为 \[\mathcal{C}_{\text{BC}} = \bigcup_{p(u)p(x|u)} \left\{ (R_1, R_2) : \begin{aligned} R_1 &\leq I(X; Y_1 \mid U) \\ R_2 &\leq I(U; Y_2) \end{aligned} \right\}\] 其中 \(U\) 为辅助随机变量（\(|\mathcal{U}| \leq \min\{|\mathcal{X}|, |\mathcal{Y}_1|, |\mathcal{Y}_2|\}\)）。

\textbf{证明}（框架）：可达性采用\textbf{叠加编码}（superposition coding）：以 \(U^n\) 作为''云中心''码字（对接收者 2 可解码），以 \(X^n\) 为''卫星''码字（条件于 \(U^n\)，对接收者 1 可解码）。接收者 2 解码 \(U^n\)（速率 \(R_2 \leq I(U; Y_2)\)），接收者 1 解码 \(U^n\) 和 \(X^n\)（速率 \(R_1 \leq I(X; Y_1|U)\)）。逆定理由 Fano 不等式和 Csiszar-sum 恒等式导出。\(\blacksquare\)

\textbf{注记 45.6}：一般广播信道（非退化）的容量区域至今仍是信息论中最著名的未解决问题之一（Marton 区域是目前已知的最优内界，1979）。这反映了网络信息论的核心困难------即使是最简单的三节点网络，最优编码策略仍可能极为复杂。

\textbf{定理 45.14}（中继信道的基本界，Cover-El Gamal 1979）：中继信道 \((\mathcal{X} \times \mathcal{X}_R, p(y, y_R|x, x_R), \mathcal{Y} \times \mathcal{Y}_R)\) 中，信源发送 \(X\)，中继收到 \(Y_R\) 并发送 \(X_R\)，目的收到 \(Y\)。容量下界（译码-转发）为 \[C \geq \max_{p(x, x_R)} \min\{ I(X, X_R; Y), \; I(X; Y_R \mid X_R) \}\]

\textbf{证明}（框架）：中继首先完全解码源消息（利用 \(I(X; Y_R|X_R)\)），然后与源协作向目的发送（利用 \(I(X, X_R; Y)\)）。瓶颈是两者中的较小值。这一策略在几何上对应''中继靠近源''时性能优良。备选策略包括压缩-转发（compress-and-forward）和放大-转发（amplify-and-forward），分别适用于中继靠近目的和源-目的直接链路较弱的场景。

\textbf{注记 45.7}（网络信息论的开放问题与影响）：尽管 Shannon 于 1948 年奠定了点对点信息论的基础，但网络信息论中许多基本问题至今未解：(1) 一般广播信道容量区域（前述 Marton 区域与 Nair-El Gamal 外界的间隙）；(2) 一般干扰信道容量区域（仅强干扰和极弱干扰等特殊情形已知）；(3) 一般中继信道容量（仅退化中继信道等特殊情形已知）。然而，网络信息论已经深刻影响了现代通信系统------叠加编码（4G/5G 中的 NOMA）、干扰对齐（MIMO 系统的自由度分析）、网络编码（Butterfly 网络中的 Ahlswede 等人 2000 年结果）和协作通信（分布式 MIMO）均直接源于网络信息论的基本原理。网络信息论与图论（网络流、图熵）、博弈论（干扰信道的博弈论解释）和编码理论（分布式信源编码的 Slepian-Wolf 定理）的交叉研究是当前最活跃的前沿方向之一。

\subsection{46.A 信息论补充}\label{a-ux4fe1ux606fux8bbaux8865ux5145}

\begin{quote}
组合数学研究离散结构的计数、存在性和构造问题，图论研究由顶点和边构成的网络结构。本卷在初等计数原理（V6 Ch 32）的基础上，建立生成函数、图论、匹配与网络流、染色与 Ramsey 理论，以及组合设计的系统理论。组合数学与代数、概率论、计算机科学和优化理论有深刻联系。
\end{quote}

组合数学的核心方法可概括为三大支柱：\textbf{计数}（enumerative combinatorics）------通过生成函数、Polya 计数和物种理论量化离散结构；\textbf{存在性}（existential combinatorics）------利用概率方法、Ramsey 理论和极值组合证明具有特定性质的配置必然存在；\textbf{构造}（constructive combinatorics）------设计高效算法和组合设计（如区组设计、正交拉丁方）以实际构造所需结构。图论作为组合数学的几何延伸，其研究对象从简单的树和回路到复杂的网络和随机图，核心问题包括连通性、匹配、着色、平面性和极值图论。生成函数方法（普通生成函数和指数生成函数）将组合计数问题转化为形式幂级数的代数运算，通过递推关系、函数方程和复分析中的奇点分析技术提取渐近信息。匹配理论中的 Hall 婚配定理、Kőnig 定理和 Tutte 定理构成了二分图和一般图匹配的理论基础，而网络流理论（特别是最大流最小割定理）提供了高效的算法实现。Ramsey 理论揭示了''完全无序''的不可能性------任何足够大的结构必然包含高度有序的子结构，其经典结果包括 Ramsey 定理（图论版本）、van der Waerden 定理（算术级数）和 Schur 定理（加性数论）。组合设计（如 BIBD、Steiner 系、正交拉丁方）研究将元素分配到满足特定均衡条件的区块中，与有限域、编码理论和统计学实验设计密切相关。这一系列理论共同构成了现代组合数学的丰富体系，为计算机科学、运筹学和密码学提供了不可或缺的数学工具。

\hrulefill

\hrulefill

\hrulefill

\hrulefill

\hrulefill

\hrulefill

\section{第97章：生成函数与递推关系}\label{ux7b2c97ux7ae0ux751fux6210ux51fdux6570ux4e0eux9012ux63a8ux5173ux7cfb}

生成函数是组合数学中最强大的工具之一，将离散的数列编码为形式幂级数，通过代数运算解决计数问题。递推关系是数列满足的方程，生成函数提供了求解递推关系的系统方法。

\subsection{127.1 普通生成函数}\label{ux666eux901aux751fux6210ux51fdux6570-1}

\textbf{定义 127.1}（普通生成函数）：数列 \((a_n)_{n \geq 0}\) 的\textbf{普通生成函数}（Ordinary Generating Function, OGF）是形式幂级数

\[A(x) = \sum_{n=0}^{\infty} a_n x^n\]

其中 \(x\) 是形式变量，不考虑收敛性。两个生成函数的加法和乘法定义为：

\[(A+B)(x) = \sum_{n=0}^{\infty} (a_n + b_n) x^n, \quad (A \cdot B)(x) = \sum_{n=0}^{\infty} \left( \sum_{k=0}^n a_k b_{n-k} \right) x^n\]

\textbf{命题 127.1}（生成函数与递推关系）：若数列 \(a_n\) 满足 \(k\) 阶常系数线性递推

\[a_n = c_1 a_{n-1} + c_2 a_{n-2} + \cdots + c_k a_{n-k}\]

则其生成函数 \(A(x)\) 是有理函数 \(A(x) = \frac{P(x)}{1 - c_1 x - c_2 x^2 - \cdots - c_k x^k}\)，其中 \(P(x)\) 是由初始条件决定的多项式。

\emph{证明}：将递推关系乘以 \(x^n\) 并对 \(n \geq k\) 求和： \[\sum_{n=k}^\infty a_n x^n = c_1 \sum_{n=k}^\infty a_{n-1} x^n + \cdots + c_k \sum_{n=k}^\infty a_{n-k} x^n\] 左边为 \(A(x) - \sum_{n=0}^{k-1} a_n x^n\)，右边第 \(j\) 项为 \(c_j x^j \sum_{n=k}^\infty a_{n-j} x^{n-j} = c_j x^j (A(x) - \sum_{n=0}^{k-j-1} a_n x^n)\)。移项整理得 \((1 - \sum_{j=1}^k c_j x^j)A(x) = P(x)\)，其中 \(P(x)\) 是次数小于 \(k\) 的多项式（由初始条件确定）。故 \(A(x) = P(x)/(1 - \sum_{j=1}^k c_j x^j)\) 是有理函数。∎

\textbf{定理 127.2}（有理生成函数的分解）：若 \(A(x) = \frac{P(x)}{Q(x)}\) 是有理函数，分母 \(Q(x) = \prod_{i=1}^k (1 - \lambda_i x)^{m_i}\)（\(\lambda_i\) 互异），则

\[a_n = \sum_{i=1}^k \sum_{j=1}^{m_i} c_{i,j} n^{j-1} \lambda_i^n\]

其中 \(c_{i,j}\) 由部分分式分解确定。

\emph{证明}：将 \(A(x)\) 部分分式分解为 \(\sum_{i,j} \frac{c_{i,j}}{(1-\lambda_i x)^j}\)，利用 \(\frac{1}{(1-\lambda x)^j} = \sum_{n=0}^{\infty} \binom{n+j-1}{j-1} \lambda^n x^n\) 即得。∎

\subsection{127.2 指数生成函数}\label{ux6307ux6570ux751fux6210ux51fdux6570-1}

\textbf{定义 127.2}（指数生成函数）：数列 \((a_n)_{n \geq 0}\) 的\textbf{指数生成函数}（Exponential Generating Function, EGF）为

\[\hat{A}(x) = \sum_{n=0}^{\infty} a_n \frac{x^n}{n!}\]

EGF 的乘法对应二项式卷积：若 \(\hat{C}(x) = \hat{A}(x) \hat{B}(x)\)，则

\[c_n = \sum_{k=0}^n \binom{n}{k} a_k b_{n-k}\]

\textbf{命题 127.3}（标号结构与 EGF）：若组合结构可以分解为不相交的连通分支，且每个分支持有标号，则 EGF 满足指数公式：若 \(\mathcal{C}\) 是连通结构的类，其 EGF 为 \(C(x)\)，则所有结构（连通分支的集合）的 EGF 为 \(e^{C(x)}\)。

\subsection{127.3 递推关系的求解}\label{ux9012ux63a8ux5173ux7cfbux7684ux6c42ux89e3}

\textbf{定义 127.3}（常系数线性递推）：\(k\) 阶常系数齐次线性递推：

\[a_n = c_1 a_{n-1} + c_2 a_{n-2} + \cdots + c_k a_{n-k}, \quad n \geq k\]

对应的\textbf{特征方程}为 \(\lambda^k - c_1 \lambda^{k-1} - \cdots - c_k = 0\)。

\textbf{定理 127.4}（特征根法）：设特征根为 \(\lambda_1, \ldots, \lambda_r\)，重数分别为 \(m_1, \ldots, m_r\)（\(\sum m_i = k\)）。则通解为

\[a_n = \sum_{i=1}^r \sum_{j=1}^{m_i} \alpha_{i,j} n^{j-1} \lambda_i^n\]

\textbf{证明}：将递推 \(a_n = \sum_{t=1}^k c_t a_{n-t}\) 写为矩阵形式 \(\mathbf{v}_n = A \mathbf{v}_{n-1}\)，其中 \(\mathbf{v}_n = (a_n, a_{n-1}, \ldots, a_{n-k+1})^T\)，\(A\) 为 \(k \times k\) 友矩阵。\(A\) 的特征多项式恰为 \(\lambda^k - \sum_{t=1}^k c_t \lambda^{k-t}\)。若 \(A\) 可对角化，则 \(a_n\) 为 \(\lambda_i^n\) 的线性组合。一般情形下，\(A\) 的 Jordan 标准型中每个 Jordan 块 \(J(\lambda_i, m_i)\) 贡献 \(n^{j-1} \lambda_i^n\)（\(j = 1, \ldots, m_i\)）。由线性递推的解空间是 \(k\) 维向量空间，上述 \(k\) 个函数线性无关，故构成通解。\(\blacksquare\)

\textbf{定理 127.5}（主定理 / Master Theorem）：对分治递推 \(T(n) = a T(n/b) + f(n)\)，其中 \(a \geq 1, b > 1\)：

\begin{enumerate}
\def\labelenumi{\arabic{enumi}.}
\tightlist
\item
  若 \(f(n) = O(n^{\log_b a - \varepsilon})\)，则 \(T(n) = \Theta(n^{\log_b a})\)
\item
  若 \(f(n) = \Theta(n^{\log_b a} \log^k n)\)，则 \(T(n) = \Theta(n^{\log_b a} \log^{k+1} n)\)
\item
  若 \(f(n) = \Omega(n^{\log_b a + \varepsilon})\) 且 \(a f(n/b) \leq c f(n)\)（\(c < 1\)），则 \(T(n) = \Theta(f(n))\)
\end{enumerate}

\textbf{证明}（递归树法）：展开递推 \(T(n) = f(n) + a f(n/b) + a^2 f(n/b^2) + \cdots + a^{\log_b n} T(1)\)。第 \(i\) 层有 \(a^i\) 个节点，每节点代价 \(f(n/b^i)\)。总代价 \(T(n) = \Theta(n^{\log_b a}) + \sum_{i=0}^{\log_b n - 1} a^i f(n/b^i)\)。

情形1：\(f(n) = O(n^{\log_b a - \varepsilon})\)，则 \(a^i f(n/b^i) = O(a^i (n/b^i)^{\log_b a - \varepsilon}) = O(n^{\log_b a - \varepsilon} a^i / b^{i(\log_b a - \varepsilon)}) = O(n^{\log_b a - \varepsilon} b^{i\varepsilon})\)。几何级数求和得 \(\sum_{i} a^i f(n/b^i) = O(n^{\log_b a})\)，故 \(T(n) = \Theta(n^{\log_b a})\)。

情形2：\(f(n) = \Theta(n^{\log_b a} \log^k n)\)，则 \(a^i f(n/b^i) = \Theta(n^{\log_b a} \log^k (n/b^i))\)。求和得 \(\sum_{i=0}^{\log_b n} n^{\log_b a} (\log n - i \log b)^k = \Theta(n^{\log_b a} \log^{k+1} n)\)。

情形3：\(f(n) = \Omega(n^{\log_b a + \varepsilon})\) 且 \(a f(n/b) \leq c f(n)\)，则 \(a^i f(n/b^i) \leq c^i f(n)\)。几何级数求和得 \(\sum_{i} a^i f(n/b^i) = O(f(n))\)，故 \(T(n) = \Theta(f(n))\)。\(\blacksquare\)

\subsection{127.4 容斥原理}\label{ux5bb9ux65a5ux539fux7406-1}

\textbf{定理 127.6}（容斥原理）：设 \(A_1, \ldots, A_n\) 是有限集合 \(X\) 的子集，则

\[\left| X \setminus \bigcup_{i=1}^n A_i \right| = \sum_{k=0}^n (-1)^k \sum_{1 \leq i_1 < \cdots < i_k \leq n} \left| A_{i_1} \cap \cdots \cap A_{i_k} \right|\]

特别地，\(|X| = \sum_{k=0}^n (-1)^k S_k\)，其中 \(S_k = \sum_{i_1 < \cdots < i_k} |A_{i_1} \cap \cdots \cap A_{i_k}|\)。

\emph{证明}：对每个元素 \(x \in X\)，设 \(x\) 恰好属于 \(r\) 个 \(A_i\)。则 \(x\) 在右边被计数 \(\sum_{k=0}^r (-1)^k \binom{r}{k} = (1-1)^r = 0\) 次（若 \(r > 0\)），而被计数 \(1\) 次（若 \(r = 0\)）。∎

\textbf{定理 127.7}（Möbius 反演）：设 \(f, g\) 为算术函数，则

\[g(n) = \sum_{d|n} f(d) \iff f(n) = \sum_{d|n} \mu(d) g\left(\frac{n}{d}\right)\]

其中 \(\mu\) 是 Möbius 函数：\(\mu(n) = 0\) 若 \(n\) 有平方因子，否则 \(\mu(n) = (-1)^k\)（\(k\) 为 \(n\) 的素因子个数）。

\textbf{证明}：(\(\Rightarrow\)) 设 \(g(n) = \sum_{d|n} f(d)\)。则

\[\sum_{d|n} \mu(d) g\left(\frac{n}{d}\right) = \sum_{d|n} \mu(d) \sum_{e|\frac{n}{d}} f(e) = \sum_{e|n} f(e) \sum_{d|\frac{n}{e}} \mu(d).\]

由 Möbius 函数的基本性质：\(\sum_{d|m} \mu(d) = 1\) 当 \(m = 1\)，否则为 \(0\)。因此内层和式中仅当 \(\frac{n}{e} = 1\)（即 \(e = n\)）时非零，故原式 \(= f(n)\)。

(\(\Leftarrow\)) 设 \(f(n) = \sum_{d|n} \mu(d) g(n/d)\)。则

\[\sum_{d|n} f(d) = \sum_{d|n} \sum_{e|d} \mu(e) g\left(\frac{d}{e}\right) = \sum_{e|n} \mu(e) \sum_{k|\frac{n}{e}} g(k).\]

令 \(h(m) = \sum_{k|m} g(k)\)。由 \(\sum_{d|m} \mu(d) h(m/d) = g(m)\)（Möbius 反演的另一种形式），取 \(m = n/e\) 可得 \(\sum_{e|n} \mu(e) h(n/e) = g(n)\)。故 \(g(n) = \sum_{d|n} f(d)\)。\(\blacksquare\)

\hrulefill

\hrulefill

\hrulefill

\hrulefill

\hrulefill

\newpage

\chapter{04-离散数学/02-概率与信息/02-编码与密码.md}\label{ux79bbux6563ux6570ux5b6602-ux6982ux7387ux4e0eux4fe1ux606f02-ux7f16ux7801ux4e0eux5bc6ux7801.md}

\section{第98章：纠错码}\label{ux7b2c98ux7ae0ux7ea0ux9519ux7801}

纠错码通过在传输信息中添加冗余，使得接收端可以检测和纠正传输错误。纠错码理论是代数、组合数学和信息论的交汇。

\subsection{144.1 线性码}\label{ux7ebfux6027ux7801}

\textbf{定义 144.1}（线性码）：一个 \([n, k]\) 线性码 \(\mathcal{C}\) 是 \(\mathbb{F}_q^n\) 的 \(k\) 维子空间。\(n\) 为码长，\(k\) 为信息位数，\(R = k/n\) 为码率。

\textbf{定义 144.2}（生成矩阵与校验矩阵）： - \textbf{生成矩阵} \(G\)（\(k   imes n\)）：\(\mathcal{C} = \{\mathbf{m} G : \mathbf{m} \in \mathbb{F}_q^k\}\) - \textbf{校验矩阵} \(H\)（\((n-k)   imes n\)）：\(\mathcal{C} = \{\mathbf{x} \in \mathbb{F}_q^n : H\mathbf{x}^T = \mathbf{0}\}\)

\(G H^T = 0\)。

\textbf{定义 144.3}（汉明距离）：\(\mathbf{x}, \mathbf{y} \in \mathbb{F}_q^n\) 的\textbf{汉明距离} \(d(\mathbf{x}, \mathbf{y})\) 是两者不同坐标的个数。码 \(\mathcal{C}\) 的\textbf{最小距离} \(d = \min_{\mathbf{x}
eq \mathbf{y} \in \mathcal{C}} d(\mathbf{x}, \mathbf{y})\)。\([n, k, d]\) 码可纠正 \(\lfloor (d-1)/2
floor\) 个错误。

\textbf{定理 144.1}（Singleton 界）：\(d \leq n - k + 1\)。达到此界的码称为\textbf{极大距离可分码}（MDS），如 Reed-Solomon 码。

\textbf{证明}：设 \([n, k, d]\) 线性码 \(\mathcal{C} \subseteq \mathbb{F}_q^n\)。选取 \(\mathcal{C}\) 的生成矩阵 \(G\)（\(k   imes n\)）。考虑 \(G\) 的前 \(n-d+1\) 列：若这些列线性相关，则存在非零码字在前 \(n-d+1\) 个坐标全为零，此码字的汉明重量 \(\leq n - (n-d+1) = d-1\)，与最小距离为 \(d\) 矛盾。故 \(G\) 的任意 \(n-d+1\) 列线性无关，从而 \(k = \operatorname{rank}(G) \leq n-d+1\)，即 \(d \leq n-k+1\)。MDS 码达到此界，等价于 \(G\) 的任意 \(k\) 列线性无关。Reed-Solomon 码的生成矩阵是 Vandermonde 矩阵，任意 \(k\) 列子式非零，故为 MDS。\(lacksquare\)

\textbf{定理 144.2}（Hamming 界 / 球包界）：\(q^k \sum_{i=0}^{\lfloor (d-1)/2 \rfloor} \binom{n}{i} (q-1)^i \leq q^n\)。达到此界的码称为\textbf{完美码}。

\textbf{证明}：以每个码字 \(\mathbf{c} \in \mathcal{C}\) 为中心，半径 \(t = \lfloor (d-1)/2 \rfloor\) 的 Hamming 球 \(B_t(\mathbf{c}) = \{\mathbf{x} \in \mathbb{F}_q^n : d(\mathbf{x}, \mathbf{c}) \leq t\}\) 互不相交（因 \(d \geq 2t+1\)，若两球相交则两个码字距离 \(\leq 2t < d\)）。每个球的大小为 \(\sum_{i=0}^t \binom{n}{i} (q-1)^i\)（选择 \(i\) 个位置改变原值，每个位置有 \(q-1\) 种选择）。码字数为 \(|\mathcal{C}| = q^k\)。所有球的并包含于 \(\mathbb{F}_q^n\)（大小为 \(q^n\)），故 \(q^k \sum_{i=0}^t \binom{n}{i} (q-1)^i \leq q^n\)。完美码恰为这些球完全覆盖 \(\mathbb{F}_q^n\) 的码。\(\blacksquare\)

\textbf{定理 144.3}（Gilbert-Varshamov 界）：存在 \([n, k, d]\) 线性码满足 \(q^{n-k} > \sum_{i=0}^{d-2} \binom{n-1}{i} (q-1)^i\)。

\textbf{证明}：贪心构造校验矩阵 \(H\)（\((n-k) \times n\)）。\(H\) 的列需满足任意 \(d-1\) 列线性无关。逐列构造：已选 \(m\) 列，新列不能是已选列中任意 \(\leq d-2\) 列的线性组合。已选 \(m\) 列中 \(\leq d-2\) 列的线性组合个数 \(\leq \sum_{i=0}^{d-2} \binom{m}{i} (q-1)^i\)。当 \(m < n\)，若 \(q^{n-k} > \sum_{i=0}^{d-2} \binom{n-1}{i} (q-1)^i\)，则 \(\mathbb{F}_q^{n-k}\) 中总存在未使用的向量，贪心过程可完成。由此得到校验矩阵 \(H\)，其零空间即为所需的 \([n, k, d]\) 码。\(\blacksquare\)

\subsection{144.2 经典线性码}\label{ux7ecfux5178ux7ebfux6027ux7801}

\textbf{定义 144.4}（BCH 码与 RS 码的译码）：Peterson-Gorenstein-Zierler 算法和 Berlekamp-Massey 算法可在多项式时间内纠正最多 \(\lfloor (d-1)/2 \rfloor\) 个错误。

经典线性码是纠错码理论中最具实用价值的一类构造，其代数结构保证了高效的编码和译码算法。\textbf{Reed-Solomon 码}（RS 码，1960）是 \([n, k, n-k+1]_q\) 的 MDS 码（\(n \leq q-1\)），其编码为对 \(k\) 个信息符号进行多项式插值：将信息视为多项式 \(m(x) \in \mathbb{F}_q[x]_{<k}\) 的系数，码字为 \(m(x)\) 在 \(n\) 个互异点 \(\alpha_1, \ldots, \alpha_n \in \mathbb{F}_q\) 上的取值 \((m(\alpha_1), \ldots, m(\alpha_n))\)。RS 码广泛用于 CD、DVD、QR 码和深空通信。\textbf{BCH 码}（Bose-Chaudhuri-Hocquenghem，1959-1960）是 RS 码的推广，通过指定根集合 \(\{\alpha^b, \alpha^{b+1}, \ldots, \alpha^{b+\delta-2}\}\) 设计最小距离至少为 \(\delta\) 的循环码。BCH 码的译码可采用 PGZ 算法（求解 Newton 恒等式和线性方程组定位错误位置）或 Berlekamp-Massey 算法（利用移位寄存器综合，时间 \(O(t^2)\)，\(t = \lfloor (d-1)/2 \rfloor\)）。\textbf{Hamming 码}是 \([2^m - 1, 2^m - 1 - m, 3]_2\) 的完美码，其校验矩阵的列恰好是 \(\mathbb{F}_2^m\) 中所有非零向量，可纠正 1 个错误。\textbf{Golay 码}（1949）包括 \([23, 12, 7]_2\) 的二元 Golay 码和 \([11, 6, 5]_3\) 的三元 Golay 码，二者均为完美码，且自同构群分别为 Mathieu 群 \(M_{23}\) 和 \(M_{11}\)------这是有限单群与纠错码之间的美妙联系。\textbf{LDPC 码}（低密度奇偶校验码，Gallager 1963，Mackay-Neal 1996 重新发现）通过稀疏校验矩阵和迭代置信传播译码接近 Shannon 容量，是现代通信标准（5G、WiFi 6）的核心技术。

\subsection{144.3 代数几何码}\label{ux4ee3ux6570ux51e0ux4f55ux7801}

\textbf{定义 144.5}（代数几何码 / Goppa 码，1981）：在有限域 \(\mathbb{F}_q\) 上的代数曲线 \(\mathcal{X}\) 上，取有理点 \(P_1, \ldots, P_n\) 和除子 \(D\)。码由 Riemann-Roch 空间 \(L(D)\) 在点上的估值构成。当曲线亏格 \(g\) 较小时，AG 码可超越 Gilbert-Varshamov 界（Tsfasman-Vladut-Zink，1982）。

代数几何码（AG 码）是经典线性码的深远推广，其核心思想是将编码问题转化为代数曲线上的除子理论。给定亏格为 \(g\) 的光滑射影曲线 \(\mathcal{X}/\mathbb{F}_q\)，选取 \(\mathbb{F}_q\)-有理点 \(P_1, \ldots, P_n\) 和 \(\mathbb{F}_q\)-有理除子 \(G\)（\(\operatorname{supp}(G) \cap \{P_1, \ldots, P_n\} = \varnothing\)），AG 码 \(C_L(D, G)\) 定义为估值映射 \(\operatorname{ev}: L(G) \to \mathbb{F}_q^n\) 的像，其中 \(\operatorname{ev}(f) = (f(P_1), \ldots, f(P_n))\)。由 Riemann-Roch 定理，\(C_L(D, G)\) 的维数 \(k \geq \deg(G) - g + 1\)（当 \(\deg(G) > 2g-2\) 时取等号），最小距离 \(d \geq n - \deg(G)\)。Tsfasman-Vladut-Zink（1982）的突破性结果是：当曲线 \(\mathcal{X}\) 的 \(\mathbb{F}_q\)-有理点数与亏格之比 \(N(\mathcal{X})/g(\mathcal{X})\) 超过 \(\sqrt{q}-1\) 时（需 \(q\) 为平方数），AG 码的渐近界超越了经典的 Gilbert-Varshamov 界。这首次证明了存在优于传统随机码的构造性码族。现代 AG 码的研究已延伸至模曲线、Drinfeld 模和塔式曲线（如 Garcia-Stichtenoth 塔），持续推动着渐近好码的构造前沿。

AG 码的译码是一个深刻的代数几何问题。Shokrollahi 和 Guruswami-Sudan 发展的列表译码算法（list decoding）将 AG 码的译码能力从 \(\lfloor (d-1)/2 \rfloor\) 扩展到远超过传统半径，其核心是将译码问题转化为曲线上的有理点计数问题，再利用代数曲线上的除子理论和插值方法求解。对于 Hermite 曲线和 Garcia-Stichtenoth 塔上的 AG 码，已有高效的译码算法，其复杂度在码长 \(n\) 上是多项式的。AG 码的构造性和渐近最优性使其成为后量子密码学中基于编码的密码系统的候选方案------McEliece 密码系统使用 AG 码（而非原始的 Goppa 码）可以显著减小公钥尺寸。

\subsection{144.4 卷积码、Turbo 码与 Polar 码}\label{ux5377ux79efux7801turbo-ux7801ux4e0e-polar-ux7801}

\textbf{定义 144.6}（卷积码）：与分组码不同，卷积码（Convolutional Codes，Elias 1955）具有记忆性------其编码器是一个有限状态机，每个时刻的输出不仅取决于当前输入，还取决于前 \(m\) 个时刻的输入（\(m\) 为约束长度）。卷积码的编码器可由生成多项式矩阵或状态转移图描述，其译码通常采用 \textbf{Viterbi 算法}（1967），该算法是动态规划在最大似然序列译码中的应用，复杂度为 \(O(2^m \cdot n)\)。

\textbf{定义 144.7}（Turbo 码与迭代译码）：Turbo 码（Berrou-Glavieux-Thitimajshima，1993）通过并行级联两个卷积码编码器并在其间插入交织器，实现了接近 Shannon 极限的纠错性能。其译码采用迭代软输入软输出（SISO）算法，两个分量译码器交替交换外信息（extrinsic information），直到收敛。Turbo 码的发明是信息论历史上里程碑式的事件，首次在实用编码方案中距 Shannon 极限仅 0.7 dB，直接推动了 3G/4G 移动通信的革命性进步。

\textbf{定义 144.8}（Polar 码，Arikan 2009）：Polar 码是第一种被严格证明在二进制输入离散无记忆信道（B-DMC）上达到 Shannon 信道容量的构造性编码方案。其核心思想是\textbf{信道极化}（channel polarization）：通过递归地组合和分裂 \(N = 2^n\) 个独立信道，当 \(N \to \infty\) 时，合成信道发生极化------一部分子信道成为''无噪声信道''（容量趋于 \(1\)），另一部分成为''完全噪声信道''（容量趋于 \(0\)）。在容量趋于 \(1\) 的子信道上传输信息比特，在容量趋于 \(0\) 的子信道上放置冻结比特（收发双方已知），即可实现可靠通信。Polar 码的编码复杂度为 \(O(N \log N)\)，译码采用串行抵消（Successive Cancellation）算法，复杂度同样为 \(O(N \log N)\)。通过 CRC 辅助的串行抵消列表译码（CA-SCL），Polar 码在有限码长下的性能超越了 Turbo 码和 LDPC 码，因此被 3GPP 采纳为 5G NR 控制信道的编码方案。Polar 码的数学基础是信道极化现象的随机矩阵分析（Arikan 的原始证明使用了大数定律和鞅收敛定理），以及极化码的代数结构（与 Reed-Muller 码的密切联系）。

\subsection{144.5 网络编码与分布式存储码}\label{ux7f51ux7edcux7f16ux7801ux4e0eux5206ux5e03ux5f0fux5b58ux50a8ux7801}

\textbf{定义 144.9}（网络编码，Ahlswede-Cai-Li-Yeung 2000）：网络编码将传统路由网络中''存储-转发''的范式推广为''编码-转发''，允许中间节点对接收到的数据包进行编码组合后再转发。线性网络编码是最常用的形式，中间节点将接收到的数据包在有限域上线性组合后输出，接收端通过解线性方程组恢复原始消息。网络编码的理论基石是\textbf{最大流最小割定理}的推广：在单源多播网络中，若源到每个接收端的最大流（最小割）至少为 \(h\)，则通过网络编码，每个接收端可以同时接收到 \(h\) 个单位的信息------而传统路由无法达到这一多播容量。

\textbf{定义 144.10}（再生码与局部修复码）：在分布式存储系统中，纠删码（erasure codes）用于在节点故障时保护数据。\textbf{再生码}（Regenerating Codes，Dimakis 等 2010）在修复单个故障节点时最小化跨网络传输的数据量，其最优折衷（存储-修复带宽折衷曲线）由最大流最小割定理的网络编码版本给出。\textbf{局部修复码}（Locally Repairable Codes，LRC）限制修复单个故障节点所需访问的节点数（局部性 \(r\)），在微软 Azure 和 Facebook 的分布式存储系统中已有实际部署。LRC 的界（如 Singleton 界的推广）和最优构造是当前编码理论的热点方向。

\subsection{144.6 量子纠错码导引}\label{ux91cfux5b50ux7ea0ux9519ux7801ux5bfcux5f15}

\textbf{定义 144.11}（量子纠错码）：量子纠错码（Quantum Error-Correcting Codes）保护量子信息免受退相干和操作误差的影响。与经典纠错码的根本区别在于：(i) 量子不可克隆定理禁止复制量子态；(ii) 测量量子态会破坏叠加态；(iii) 错误可以是连续的错误（如任意小的位相旋转）。Shor（1995）和 Steane（1996）独立发现了第一个量子纠错码------将 1 个逻辑量子比特编码为 9 个（或 7 个）物理量子比特，可纠正任意单量子比特错误。

\textbf{定理 144.4}（量子纠错条件 / Knill-Laflamme 定理，1997）：设 \(\mathcal{C}\) 为 Hilbert 空间 \(\mathcal{H}\) 的子空间，\(\{E_a\}\) 为一组可能的错误算子。\(\mathcal{C}\) 上存在纠错恢复操作当且仅当对所有 \(a, b\) 和 \(|\psi\rangle, |\phi\rangle \in \mathcal{C}\)（\(\langle\psi|\phi\rangle = 0\)），有 \(\langle\psi|E_a^\dagger E_b|\phi\rangle = 0\)，且 \(\langle\psi|E_a^\dagger E_b|\psi\rangle = C_{ab}\) 与 \(|\psi\rangle\) 无关。

\textbf{证明}：必要性来自量子测量理论：若两个不同错误导致正交态，则可区分和纠正。充分性通过构造恢复操作：将错误空间分解为正交子空间，在每个子空间上施加逆操作。条件 \(\langle\psi|E_a^\dagger E_b|\psi\rangle\) 与 \(|\psi\rangle\) 无关保证不同错误对不同码字产生相同的度量畸变，可统一纠正。\(\blacksquare\)

\textbf{定义 144.12}（CSS 码与稳定子码）：Calderbank-Shor-Steane（CSS）码利用两个经典线性码 \(C_1 \subseteq C_2\) 构造量子码 \(CSS(C_1, C_2)\)，其参数为 \([[n, k_2-k_1, \min(d_1, d_2^\perp)]]\)。Gottesman 的稳定子形式（stabilizer formalism）将量子纠错码的构造转化为 Abel 群表示论问题：稳定子群 \(S \subseteq \mathcal{P}_n\)（Pauli 群的 Abel 子群）的联合 \(+1\) 本征空间即为码空间。稳定子形式统一了包括 Shor 码、Steane 码、toric 码和 surface 码在内的几乎所有已知量子码构造，并为容错量子计算提供了理论基础。

\textbf{定义 144.13}（Toric 码与 Surface 码）：Kitaev 的 toric 码（2003）是在二维环面上的量子纠错码，其稳定子由格点上的星形算子和平移算子构成。Surface 码是 toric 码在平面上的推广，具有极高的容错阈值（约 \(1\%\)）和局部（二维）稳定子结构，使其成为当前容错量子计算最受瞩目的物理实现方案。Surface 码的纠错过程本质上是一个二维自旋模型（如二维 Ising 模型）的统计推断问题，其阈值与二维随机键 Ising 模型的有序-无序相变点一致------这一联系揭示了拓扑量子纠错码与统计力学之间的深刻统一性。

\subsection{144.7 后量子密码与编码密码学}\label{ux540eux91cfux5b50ux5bc6ux7801ux4e0eux7f16ux7801ux5bc6ux7801ux5b66}

\textbf{定义 144.14}（基于编码的密码系统）：McEliece 密码系统（1978）是最早的公钥密码系统之一，其安全性基于一般线性码的译码问题是 NP 完全的这一事实。密钥生成：选取一个具有高效译码算法的 \([n, k, 2t+1]\) 线性码（如 Goppa 码），其生成矩阵 \(G\) 通过随机可逆矩阵 \(S\) 和置换矩阵 \(P\) 伪装为 \(\hat{G} = SGP\)（看似随机码）。公钥为 \((\hat{G}, t)\)，私钥为 \((S, G, P)\) 和译码算法。加密：\(c = m\hat{G} + e\)，其中 \(e\) 为重量 \(\leq t\) 的随机错误向量。解密：计算 \(cP^{-1} = mSG + eP^{-1}\)，利用高效译码算法恢复 \(mS\)，再乘以 \(S^{-1}\) 得 \(m\)。McEliece 密码系统在 40 余年密码分析中未被攻破，且已知的量子攻击（Grover 算法）仅提供二次加速，不足以威胁其安全性，因此被 NIST 列为后量子密码标准化候选方案之一。基于编码的密码系统的变体包括 Niederreiter 密码系统（使用校验矩阵而非生成矩阵）和基于 QC-MDPC（准循环中等密度奇偶校验）码的现代变体，后者显著减小了公钥尺寸。

\(\blacksquare\)

\section{第99章：低密度奇偶校验码（LDPC）与图论编码}\label{ux7b2c99ux7ae0ux4f4eux5bc6ux5ea6ux5947ux5076ux6821ux9a8cux7801ldpcux4e0eux56feux8bbaux7f16ux7801}

低密度奇偶校验码（Low-Density Parity-Check Codes，简称 LDPC 码）由 Gallager 于 1963 年在其博士论文中首次提出，但因当时计算能力的限制而被遗忘近三十年。1996 年，Mackay 和 Neal 重新发现了 LDPC 码，并证明了其在二进制对称信道和 AWGN 信道下通过迭代置信传播（Belief Propagation）译码可以达到接近 Shannon 容量的性能。LDPC 码的核心特征是其校验矩阵 \(H\) 是稀疏的------每一行和每一列的非零元素个数均为常数级，与码长 \(n\) 无关。这一稀疏性不仅是高效译码算法的前提，更将编码理论问题转化为图论中的组合优化问题：校验矩阵 \(H\) 自然地对应一个 Tanner 图（二分图），其中左侧是 \(n\) 个变量节点（对应码字的每个比特），右侧是 \(m = n-k\) 个校验节点（对应每个校验方程），边连接变量节点 \(j\) 与校验节点 \(i\) 当且仅当 \(H_{ij} \neq 0\)。置信传播译码算法正是在此二分图上进行的消息传递过程，其每一轮迭代中，变量节点和校验节点交替地交换概率似然信息，直至收敛或达到最大迭代次数。LDPC 码与图论之间的这种深刻联系是信息论和组合数学交叉的典范：码的性能（如最小距离、围长）由 Tanner 图的拓扑性质（围长、连通度、展开性）决定，而好码的构造等价于构造具有优良展开性质的稀疏二分图。

\subsection{145.x 正则与不规则 LDPC 码的度分布设计}\label{x-ux6b63ux5219ux4e0eux4e0dux89c4ux5219-ldpc-ux7801ux7684ux5ea6ux5206ux5e03ux8bbeux8ba1}

\textbf{定义 145.1}（正则 LDPC 码）：一个 \((d_v, d_c)\)-正则 LDPC 码是指其 Tanner 图中每个变量节点度数为 \(d_v\)，每个校验节点度数为 \(d_c\)。对应的校验矩阵 \(H\) 的每一列有恰好 \(d_v\) 个非零元，每一行有恰好 \(d_c\) 个非零元。在有限域 \(\mathbb{F}_q\) 上的推广即为 \(q\)-元 LDPC 码，其非零元从 \(\mathbb{F}_q^\times\) 中选取，可进一步改善性能。

\textbf{定义 145.2}（不规则 LDPC 码与度分布）：不规则 LDPC 码允许变量节点和校验节点具有不同的度数，其 Tanner 图结构由两个度分布多项式（从边视角）描述：

\[\lambda(x) = \sum_{i=2}^{d_v^{\max}} \lambda_i x^{i-1}, \quad \rho(x) = \sum_{j=2}^{d_c^{\max}} \rho_j x^{j-1}\]

其中 \(\lambda_i\) 表示连接到度数为 \(i\) 的变量节点的边占总边数的比例，\(\rho_j\) 表示连接到度数为 \(j\) 的校验节点的边占总边数的比例。度分布的精心设计是不规则 LDPC 码逼近容量极限的关键步骤。

\textbf{定理 145.1}（密度演化与阈值定理，Richardson-Urbanke 2001）：对于给定度分布 \((\lambda, \rho)\) 的 LDPC 码族（码长 \(n \to \infty\) 且 Tanner 图为树），在二进制擦除信道（BEC）擦除概率 \(\epsilon\) 下，置信传播译码的渐近错误概率由密度演化（Density Evolution）递推方程决定：

\[x_{\ell+1} = \epsilon \lambda(1 - \rho(1 - x_\ell))\]

其中 \(x_\ell\) 为第 \(\ell\) 轮迭代后变量节点到校验节点的擦除消息概率，初始值 \(x_0 = \epsilon\)。该迭代收敛到零当且仅当 \(\epsilon < \epsilon^{\text{BP}}\)，其中 \(\epsilon^{\text{BP}}\) 称为该度分布的置信传播阈值。存在一个严格的容量渐近阈值，当 \(\epsilon \to \epsilon^{\text{BP}}\) 时迭代次数趋于无穷。

\textbf{证明}：设第 \(\ell\) 轮从变量节点 \(v\) 发往校验节点 \(c\) 的消息为擦除的概率为 \(x_\ell\)。在 BEC 下，消息仅取三种值：\(0\)、\(1\) 或''擦除''。变量节点的更新规则：\(v\) 发往 \(c\) 的消息为擦除当且仅当信道输出为擦除且所有其他相邻校验节点发来的消息均为擦除。若 \(v\) 的度数为 \(i\)，此事件概率为 \(\epsilon \cdot x_\ell^{i-1}\)。由边视角平均得 \(x_{\ell+1} = \sum_i \lambda_i \cdot \epsilon \cdot x_\ell^{i-1} = \epsilon \lambda(x_\ell)\)。校验节点的更新规则：\(c\) 发往 \(v\) 的消息为擦除当且仅当任一其他相邻变量节点发来擦除消息。若 \(c\) 的度数为 \(j\)，此概率为 \(1 - (1 - x)^j\)，其中 \(x\) 为入边擦除概率（由对称性假设各边独立------在树假设下确实如此）。注意在密度演化中正确定义消息方向：变量到校验和校验到变量交替。完整递推为 \(x_{\ell+1} = \epsilon \lambda(1 - \rho(1 - x_\ell))\)，其中 \(\rho(x) = \sum_j \rho_j x^{j-1}\)。\(f(x) = \epsilon \lambda(1 - \rho(1 - x))\) 在 \([0,1]\) 上的唯一不动点 \(x=0\) 的稳定性决定阈值。\(\blacksquare\)

\textbf{推论 145.2}（不规则 LDPC 码的容量逼近性质）：通过精心优化度分布 \((\lambda, \rho)\)，LDPC 码的置信传播阈值 \(\epsilon^{\text{BP}}\) 可以任意接近 BEC 的信道容量 \(1 - R\)（其中 \(R = 1 - \frac{\int_0^1 \rho(x)dx}{\int_0^1 \lambda(x)dx}\) 为设计码率）。这一结论由 Luby、Mitzenmacher、Shokrollahi 和 Spielman 在 2001 年严格证明，表明不规则 LDPC 码是第一种被证明可以达到 Shannon 容量任意精度的实用编码方案。

\textbf{证明}（概述）：考虑序列 \(\lambda(x)\) 和 \(\rho(x)\) 的极限。令 \(\rho(x)\) 趋近于阶跃函数，使得校验节点几乎全部具有高且均匀的度数。关键不等式 \(f(x) = \epsilon \lambda(1 - \rho(1 - x)) \leq x\) 对所有 \(x \in (0,1]\) 等价于阈值条件。取 \(\lambda(x) = \sum_{i} \lambda_i x^{i-1}\) 满足 \(\lambda_i \sim 1/(i(i-1))\)（此分布使右侧函数在零点附近的行为匹配线性函数），可以证明对于任意 \(\delta > 0\) 存在度分布使 \(\epsilon^{\text{BP}} > 1 - R - \delta\)。定理的完整证明涉及对密度演化函数不动点的精细分析和序列度分布的极限构造，详见 Richardson-Urbanke (2008) 专著。\(\blacksquare\)

\subsection{145.y LDPC 码的围长条件与构造方法}\label{y-ldpc-ux7801ux7684ux56f4ux957fux6761ux4ef6ux4e0eux6784ux9020ux65b9ux6cd5}

\textbf{定义 145.3}（Tanner 图的围长 / Girth）：LDPC 码的 Tanner 图的\textbf{围长}（girth）\(g\) 定义为图中最短环的长度。由于置信传播译码在无环图上才是精确的（即给出精确的边缘后验概率），Tanner 图中短环（特别是长度为 \(4\) 的环）的存在是译码性能下降的主要原因之一。围长 \(g \geq 6\) 意味着校验矩阵中任何两列最多在一个行位置上同时具有非零元------等价地说，\(H\) 中不包含 \(2 \times 2\) 的全非零子矩阵。

\textbf{定理 145.3}（基于有限几何的 LDPC 码构造，Kou-Lin-Fossorier 2001）：利用有限域 \(\mathbb{F}_{q}\) 上的射影几何 \(\operatorname{PG}(m, q)\) 或仿射几何 \(\operatorname{AG}(m, q)\) 的点与超平面关联结构构造的 LDPC 码，其 Tanner 图具有严格的围长保证 \(g \geq 6\)，且最小距离随码长以正速率增长。具体地，以 \(\operatorname{PG}(m, q)\) 的点作为变量节点、超平面作为校验节点（或反之），关联矩阵即为校验矩阵 \(H\)。该矩阵具有常数行重和列重，且因射影几何中任意两点确定唯一直线、任意两超平面交于唯一余维 \(2\) 子空间，Tanner 图中不存在长度为 \(4\) 的环。

\textbf{证明}：设变量节点对应 \(\operatorname{PG}(m, q)\) 的点，校验节点对应超平面。一个长度为 \(4\) 的环意味着存在两个不同的变量节点 \(v_1, v_2\) 和两个不同的校验节点 \(c_1, c_2\)，使得 \(v_1\) 与 \(c_1, c_2\) 均关联，\(v_2\) 也与 \(c_1, c_2\) 均关联。在射影几何中，这等价于两个不同的点同时属于两个不同的超平面。但 \(\operatorname{PG}(m, q)\) 中，两个超平面的交集是一个余维 \(2\) 的子空间（即 \(\operatorname{PG}(m-2, q)\)），其包含的点数为 \((q^{m-1} - 1)/(q - 1)\)。若 \(m \geq 3\)，交集大小 \(\geq q+1 \geq 3\)（\(q \geq 2\)），故可能包含多个点，此时可能出现 \(4\)-环。为此需选取特定子结构（如 \(\operatorname{PG}(2, q)\) 中点和线的关联，此时交点为唯一，无 \(4\)-环）。当 \(m=2\) 时，两个不同超平面（即两条不同直线）交于唯一一点，故 Tanner 图围长 \(\geq 6\)。校验矩阵的稀疏性：每行重量为超平面包含的点数 \((q^m-1)/(q-1)\)，每列重量为过该点的超平面数 \((q^m-1)/(q-1)\)（对偶结构中）。列重和行重随 \(q\) 和 \(m\) 多项式增长，但当 \(n\) 趋于无穷时相对码长 \(O(1/n)\) 趋于零，故矩阵确为稀疏。最小距离的下界由 Tanner 图展开性质和围长条件给出------围长 \(\geq 6\) 确保了最小距离至少为 \(\lfloor g/2 \rfloor\) 的某种指数界。\(\blacksquare\)

\textbf{定理 145.4}（渐进边展开与 LDPC 码的最小距离）：设 LDPC 码的 Tanner 图为 \((\alpha, \beta)\)-展开图（expander graph），即对任意大小不超过 \(\alpha n\) 的变量节点子集 \(S\)，其邻接校验节点数至少为 \(\beta |S|\)（其中 \(\beta > d_v/2\)）。则该 LDPC 码的最小距离至少为 \(\alpha n\)。特别地，基于随机置换构造的 LDPC 码以高概率满足展开条件，因而具有线性最小距离。

\textbf{证明}：设 \(S\) 为码字中非零坐标对应的变量节点集合，\(|S| = w < \alpha n\)。由于每个非零变量节点对应码字向量中的非零分量，若 \(w\) 小于最小距离，则需存在一种非零的赋值使得所有校验方程被满足。由展开性质，\(S\) 的邻接校验节点集 \(N(S)\) 满足 \(|N(S)| \geq \beta w\)。每个校验节点 \(c \in N(S)\) 连接到 \(S\) 中的若干变量节点，其度数至少为 \(2\)（否则单独一条边无法传递非零值来满足校验方程------需至少两个非零输入才能在校验节点处抵消为零，在 \(\mathbb{F}_2\) 上即 \(1+1=0\)）。每条从 \(S\) 出发的边到达 \(N(S)\) 中的某节点。\(S\) 发出的总边数为 \(d_v w\)。另一方面，\(N(S)\) 中每个校验节点至少接收来自 \(S\) 的 \(2\) 条边（否则孤立边使校验方程无法满足，因信道无噪声时该校验的 LLR 无穷大），故 \(d_v w \geq 2|N(S)| \geq 2\beta w\)，即 \(d_v \geq 2\beta\)。若 \(\beta > d_v/2\)，则推出矛盾，故 \(w \geq \alpha n\)。因此最小距离至少为 \(\alpha n\)，即线性增长。随机置换 LDPC 码族：以概率 \(1 - O(n^{-c})\)，Tanner 图满足所需的展开条件（由随机二分图的展开性质，类似随机正则图的 Alon-Boppana 型定理），从而具有线性最小距离和优良的译码性能。\(\blacksquare\)

\hrulefill

\hrulefill

\hrulefill

\section{第100章：网络编码与信息流}\label{ux7b2c100ux7ae0ux7f51ux7edcux7f16ux7801ux4e0eux4fe1ux606fux6d41}

网络编码（Network Coding）由 Ahlswede、Cai、Li 和 Yeung 在 2000 年的一篇开创性论文中提出，彻底改变了通信网络的设计范式。传统的通信网络将路由器视为''存储-转发''设备------中间节点仅复制和转发收到的数据包，不进行任何处理。网络编码的核心理念是允许中间节点对收到的数据包进行代数运算（如有限域上的线性组合）后再转发，从而在网络层实现''编码增益''。这一思想的惊人之处在于，它在多播场景中可以严格超越传统路由的信息传输容量：经典的蝴蝶网络（Butterfly Network）证明了通过网络编码，两个接收端可以同时以速率 \(2\) 接收到两个源数据包，而使用传统路由的容量仅为 \(1.5\)（以每单位时间的数据包数计）。网络编码的理论基石是最大流最小割定理的多播推广：在单源多播网络中，源节点到每个接收端的多播容量等于所有接收端独立最小割的最小值。线性网络编码的充分性（Li-Yeung-Cai 定理）进一步指出，在有限域 \(\mathbb{F}_q\) 上（\(q\) 充分大），线性编码即可达到多播容量。网络编码不仅是信息论的理论突破，更在分布式存储（再生码）、P2P 网络、无线网络和网络安全等领域产生了深远影响。

\subsection{146.x 蝴蝶网络与多播容量定理}\label{x-ux8774ux8776ux7f51ux7edcux4e0eux591aux64adux5bb9ux91cfux5b9aux7406}

\textbf{定义 146.1}（单源无环多播网络）：一个单源多播网络 \(\mathcal{N} = (V, E, s, T)\) 由有向无环图 \((V, E)\)、源节点 \(s \in V\) 和接收端集合 \(T = \{t_1, \ldots, t_k\} \subseteq V\) 构成。每条边 \(e \in E\) 具有非负整数容量 \(c(e)\)，表示每条边每单位时间可传输的有限域 \(\mathbb{F}_q\) 符号数。源节点 \(s\) 生成 \(h\) 个独立的随机符号（信息源），所有接收端需要以高概率（或确定性）恢复全部 \(h\) 个符号。

\textbf{定义 146.2}（割与最小割容量）：对接收端 \(t \in T\)，一个 \(s\)-\(t\) 割是顶点集的一个划分 \((U, V \setminus U)\) 满足 \(s \in U\) 且 \(t \notin U\)。割的容量 \(\delta(U) = \sum_{e \in \operatorname{out}(U)} c(e)\)，其中 \(\operatorname{out}(U)\) 为从 \(U\) 中顶点出发到达 \(V \setminus U\) 中顶点的边集合。\(s\) 到 \(t\) 的\textbf{最小割容量}记为 \(\min\text{-cut}(s, t) = \min_{U: s \in U, t \notin U} \delta(U)\)。

\textbf{定理 146.1}（多播容量定理，Ahlswede-Cai-Li-Yeung 2000）：在单源多播网络 \(\mathcal{N}\) 中，通过网络编码可以实现的多播容量（即所有接收端可同时接收的信息速率上界）为

\[C_{\text{multicast}} = \min_{t \in T} \min\text{-cut}(s, t)\]

即等于所有接收端各自最小割的最小值。该容量不仅可以通过网络编码达到，且当要求''所有接收端同时收到全部信息''时，这是信息论上界。

\textbf{证明}（可达性------线性网络编码）：设 \(h = \min_{t \in T} \min\text{-cut}(s, t)\) 为目标多播速率。构造性地说明 \(h\) 可达。取充分大的有限域 \(\mathbb{F}_q\)（\(q > |T|\)）。在源节点 \(s\) 处，\(h\) 个信息符号记为 \(x_1, \ldots, x_h \in \mathbb{F}_q\)。每条边 \(e\) 上传输一个 \(\mathbb{F}_q\) 符号 \(y_e\)，它是边 \(e\) 的尾节点收到的所有符号的线性组合（系数从 \(\mathbb{F}_q\) 中选取）：\(y_e = \sum_{e': \operatorname{head}(e') = \operatorname{tail}(e)} \alpha_{e', e} \, y_{e'}\)，其中源节点发出的虚拟边携带 \(y_{s,i} = x_i\)。对每个接收端 \(t_j\)，其收到的符号构成向量 \(\mathbf{y}_{t_j} = M_j \mathbf{x}\)，其中 \(M_j\) 是 \(h \times h\) 的网络传输矩阵（从源到该接收端的全局编码向量）。需选择所有局部编码系数 \(\alpha_{e', e}\) 使得每个 \(M_j\) 满秩（即可逆，从而可译码）。

关键步骤（Jaggi-Sanders 等，2005）：按拓扑序逐边选取编码系数。在每步选择时，需保证已处理的子网络中每个接收端的传输矩阵保持最大可能秩。由于每个秩约束是多项式方程 \(P(\alpha) = \det(\text{相关子矩阵}) \neq 0\)（次数至多 \(h\) 乘以至今的边数），且域 \(\mathbb{F}_q\) 充分大时（\(q > h \cdot |E|\)），非零多项式的零点集是真子集。根据 Schwartz-Zippel 引理，随机选取系数使所有 \(M_j\) 可逆的概率 \(\geq 1 - h|T|/q > 0\)。确定性构造由 Jaggi 等的算法给出：逐边处理时维护每个接收端的''候选系数''集合，贪心地选边系数使每个接收端的全局编码矩阵行数递减，直至归约到满秩。由此证明 \(h\) 速率可达。

上界（最大流最小割的推广）：由信息论中的数据处理不等式，对任意接收端 \(t\) 和任意割 \((U, V \setminus U)\)，\(\operatorname{out}(U)\) 中的边携带的信息总量不超过该割的容量。经典最大流最小割定理（Ford-Fulkerson）在网络编码框架下推广为：接收端 \(t\) 可恢复的符号数 \(\leq \min\text{-cut}(s, t)\)。对所有接收端取最小值即得上界。\(\blacksquare\)

\subsection{146.y 线性网络编码的代数框架}\label{y-ux7ebfux6027ux7f51ux7edcux7f16ux7801ux7684ux4ee3ux6570ux6846ux67b6}

\textbf{定义 146.3}（全局编码向量与传输矩阵）：设源 \(s\) 产生 \(h\) 个符号 \(\mathbf{x} = (x_1, \ldots, x_h)^T \in \mathbb{F}_q^h\)。对每条边 \(e \in E\)，定义其\textbf{全局编码向量} \(\mathbf{g}_e \in \mathbb{F}_q^h\) 使得边上传输的符号为 \(y_e = \mathbf{g}_e^T \mathbf{x}\)。全局编码向量满足递推关系：\(\mathbf{g}_e = \sum_{e': \operatorname{head}(e') = \operatorname{tail}(e)} \alpha_{e', e} \, \mathbf{g}_{e'}\)，其中源节点的 \(h\) 条虚拟出边 \(e_1, \ldots, e_h\) 分别具有全局编码向量 \(\mathbf{g}_{e_i} = \mathbf{e}_i\)（标准基向量）。

\textbf{定理 146.2}（线性网络编码的充分性，Li-Yeung-Cai 2003）：对任意单源多播网络，若多播速率 \(h\) 可达（即存在某种网络编码方案达到该速率），则存在线性网络编码方案（在充分大的有限域 \(\mathbb{F}_q\) 上）达到相同的速率。换言之，非线性编码不扩充多播容量。

\textbf{证明}：由定理 146.1 的可达性证明，所构造的编码方案本身就是线性的（每条边上传输的符号是输入符号的线性函数）。该构造在所有无环网络中有效。对于带环网络（有向图中有有向环），可通过时间展开（time-expanded network）将每一时间步复制一次，化为无环网络后应用线性方案。因此，线性编码在多播场景中始终充分。该证明的一个关键点是 Jaggi 等的多项式时间算法（2005），其复杂度为 \(O(|E| \cdot h \cdot |T|)\)，且域 \(\mathbb{F}_q\) 的大小仅需 \(q > |T|\)（而非原始的 \(q > |E|\)）。\(\blacksquare\)

\textbf{定义 146.4}（随机线性网络编码 / RLNC，Ho 等 2006）：在大型或动态变化的网络中，确定性的全局编码系数分配可能不切实际。随机线性网络编码（Random Linear Network Coding）的方案是在每个中间节点处独立随机地从 \(\mathbb{F}_q\) 中选择局部编码系数。接收端通过收集足够多的带有全局编码向量的数据包，构建线性方程组 \(M \mathbf{x} = \mathbf{y}\)（\(M\) 的每一行为收到数据包的全局编码向量）。当 \(M\) 满秩时，通过 Gauss 消元解得 \(\mathbf{x}\)。

\textbf{定理 146.3}（RLNC 的成功概率，Ho 等 2006）：在随机线性网络编码中，若域大小 \(q\) 满足 \(q > |T|\)，且每个接收端收到至少 \(h\) 个线性独立的数据包，则所有接收端成功译码的概率至少为 \((1 - |T|/q)^{\eta}\)，其中 \(\eta\) 为网络中的总传输次数。

\textbf{证明}：每个接收端 \(t\) 的传输矩阵 \(M_t\)（大小为 \(m_t \times h\)，\(m_t \geq h\) 个收到的数据包）满秩当且仅当所有 \(h \times h\) 子式不全为零。将 \(M_t\) 的元素视为独立同分布于 \(\mathbb{F}_q\)（在随机编码模型下这是一个合理假设，因中间节点的随机组合产生近乎独立的输出）。\(M_t\) 的 \(h \times h\) 子式的行列式的概率至少为 \(1 - h/q\) 不为零（由 Schwartz-Zippel 引理，行列式作为系数的多项式次数至多为 \(h\)）。对 \(|T|\) 个接收端取联合界，至少一个接收端失败的概率 \(\leq |T| \cdot h/q \leq |T|/q\)（\(h \geq 1\) 时可放宽）。每次传输均独立随机编码时，成功概率按各跳累积。设网络直径为 \(D\)，独立随机选择的总次数为 \(\eta \leq |E| \cdot D\)。由各步独立性，全部接收端成功的概率 \(\geq (1 - |T|/q)^{\eta} \geq 1 - \eta|T|/q\)（Bernoulli 不等式）。取 \(q = 2^m\) 充分大（如 \(m = 16\)），成功概率接近 \(1\)。\(\blacksquare\)

\subsection{146.z 网络编码与分布式存储：再生码}\label{z-ux7f51ux7edcux7f16ux7801ux4e0eux5206ux5e03ux5f0fux5b58ux50a8ux518dux751fux7801}

\textbf{定义 146.5}（分布式存储中的再生码 / Regenerating Codes）：在一个由 \(n\) 个节点组成的分布式存储系统中，大小为 \(M\) 的文件被编码存储，使得每个节点存储 \(\alpha\) 个符号（子分包化），任意 \(k\) 个节点可恢复整个文件（称为最大距离可分/MDS 性质），且当某个节点故障时，新节点从 \(d\) 个存活节点各下载 \(\beta\) 个符号以修复丢失的数据，总修复带宽为 \(\gamma = d\beta\)。再生码（Dimakis-Godfrey-Wu-Wainwright-Ramchandran 2010）的目标是在存储开销 \(\alpha\) 和修复带宽 \(\gamma\) 之间取得最优折衷。

\textbf{定理 146.4}（存储-修复带宽折衷界，Dimakis 等 2010）：对任意 \((n, k, d)\) 再生码，存储 \(\alpha\) 和修复带宽 \(\gamma = d\beta\) 满足

\[\sum_{i=0}^{k-1} \min\{\alpha, (d - i)\beta\} \geq M\]

等价地，在 \((\alpha, \gamma)\) 平面上定义了一条分段线性的最优折衷曲线，其两个端点为：最小存储再生（MSR）点 \((\alpha_{\text{MSR}}, \gamma_{\text{MSR}}) = (M/k, M/k \cdot d/(d-k+1))\) 和最小带宽再生（MBR）点 \((\alpha_{\text{MBR}}, \gamma_{\text{MBR}}) = (M/k \cdot 2d/(2d-k+1), M/k \cdot 2d/(2d-k+1))\)。

\textbf{证明}（信息流图 / 网络编码方法）：将节点修复过程建模为多播网络编码问题。考虑源节点（存储原始数据）到数据收集器（需恢复文件）的通信，以及各轮修复中存活节点到新节点的信息流。构造一个信息流图（Information Flow Graph），源 \(s\) 连接至初始的 \(n\) 个存储节点（各容量 \(\alpha\)）。当节点故障时，新节点加入并连接至 \(d\) 个存活节点（各边容量 \(\beta\)）。数据收集器在任意时刻可连接至任意 \(k\) 个存活节点（各边容量 \(\alpha\)）以恢复文件。这是一个多播网络，源需以速率 \(M\) 多播至所有可能的数据收集器（所有可能的时间和数据收集选择）。

由最大流最小割定理在网络编码中的推广（定理 146.1），多播容量等于所有数据收集器（在所有可能的故障和修复序列后）的最小割的最小值。对给定的故障序列，在时间 \(t\) 的数据收集器，其最小割可按存活节点的来源递归计算。最劣情形下（经过充分多轮故障后），存活节点为 \(k\) 个原始节点 + 一系列修复而成的新节点。计算该信息流图的最小割得到上界：对于任意数据收集器连接的 \(k\) 个节点，存在割 \((U, V \setminus U)\) 将源与这些节点分离，该割容量为 \(\sum_{i=0}^{k-1} \min\{\alpha, (d-i)\beta\}\)（因依次选择节点时，每个新节点最多贡献 \(\min\{\alpha, d\beta, (d-1)\beta, \ldots\}\) 的新信息）。令此下确界 \(\geq M\) 即得折衷界。\(\blacksquare\)

\textbf{推论 146.5}（再生码的网络编码构造）：MSR 和 MBR 点的再生码可通过线性网络编码进行显式构造。MBR 码等价于数据在 \(n\) 个节点间的适当复制和编码，使修复过程仅需复制而不需有限域运算。MSR 码的构造更为精细，需要利用有限域上的 Cauchy 矩阵或 Vandermonde 矩阵构造 MDS 码，并通过巧妙的子分包化（sub-packetization）实现低修复带宽。Rashmi-Shah-Kumar-Ramchandran（2011）给出了 MSR 码的乘积矩阵（Product-Matrix）构造，其修复过程等价于在 \(d\) 个存活节点上计算特定线性组合，数据量 \(d\beta = d \cdot \alpha/(d-k+1)\) 为理论最优。后续的 Zigzag 可解码再生码（Tamo 等）和局部修复码（LRC）进一步降低了 I/O 和磁盘读取次数。\(\blacksquare\)

\newpage

\chapter{05-代数/01-线性代数/01-向量空间.md}\label{ux4ee3ux657001-ux7ebfux6027ux4ee3ux657001-ux5411ux91cfux7a7aux95f4.md}

\section{第101章：向量空间公理化}\label{ux7b2c101ux7ae0ux5411ux91cfux7a7aux95f4ux516cux7406ux5316}

向量空间公理化是现代线性代数的基石，它将看似迥异的数学对象------几何向量、多项式、函数空间、数列------统一在八条简洁公理之下。这一公理体系的建立经历了漫长的历史演进：从 Hermann Grassmann 于 1844 年出版的《线性扩张论》（\emph{Die lineale Ausdehnungslehre}）中首次系统研究 \(n\) 维向量空间，到 Giuseppe Peano 于 1888 年给出第一个公理化定义，再到 Stefan Banach 在 1920 年代将公理化方法推广到无穷维空间（Banach 空间），向量空间概念逐步从具体的几何直观中抽象出来。

David Hilbert 在 1900 年前后关于积分方程的工作使人们认识到，函数空间（如 \(L^2\)）与有限维欧氏空间共享许多结构性质，但需要公理化框架来统一处理。Hermann Weyl 在其经典著作《空间、时间、物质》（1918）中强调了公理化方法在物理学中的核心地位，向量空间作为物理理论的数学舞台被牢固确立。20 世纪 30 年代，Nicolas Bourbaki 学派在其《数学原本》中将向量空间公理化推至极致，将标量域从实数推广到任意域，奠定了结构主义数学的基础。

向量空间公理化之所以深刻，在于它揭示了线性方程组的解集并非仅仅是 \(\mathbb{R}^n\) 中的点集，而是一个具有内在结构的代数对象。线性方程组 \(A\mathbf{x} = \mathbf{0}\) 的解空间------即 \(A\) 的零空间------自然构成向量空间，且其维数由秩-零化度定理给出：\(\dim \ker A + \operatorname{rank} A = n\)。这一事实将线性方程组的定性理论（解的存在性、唯一性、解空间的维数）完全转化为向量空间的结构问题。

本章从向量空间的八条公理出发，系统地建立子空间、线性相关性、基、维数等核心概念。公理 (V1)-(V4) 确保向量空间关于加法构成阿贝尔群------这是线性结构的''可加性''基础；公理 (V5)-(V8) 确保标量乘法与加法相容------这是线性结构的''缩放''性质。这八条公理看似简单，但由它们推导出的结论------如基的存在性（依赖于选择公理）、维数的不变性、子空间的交与和------构成了线性代数全部理论的逻辑起点。

\textbf{定义 32.1.1（向量空间）}

设 \(F\) 为域，\(V\) 为非空集合。\(V\) 上配备两个运算：

\begin{enumerate}
\def\labelenumi{\arabic{enumi}.}
\tightlist
\item
  \textbf{加法}：\(+ : V \times V \to V\)，\((u, v) \mapsto u + v\)
\item
  \textbf{标量乘法}：\(\cdot : F \times V \to V\)，\((\lambda, v) \mapsto \lambda v\)
\end{enumerate}

若满足以下八条公理，则称 \((V, +, \cdot)\) 为域 \(F\) 上的\textbf{向量空间}（或\textbf{线性空间}）：

\textbf{加法公理}（\(V\) 关于加法构成阿贝尔群）： - \textbf{(V1) 结合律}：\(\forall u, v, w \in V\)，\((u + v) + w = u + (v + w)\) - \textbf{(V2) 交换律}：\(\forall u, v \in V\)，\(u + v = v + u\) - \textbf{(V3) 零元存在}：\(\exists \mathbf{0} \in V\)，\(\forall v \in V\)，\(v + \mathbf{0} = v\) - \textbf{(V4) 加法逆元}：\(\forall v \in V\)，\(\exists (-v) \in V\)，\(v + (-v) = \mathbf{0}\)

\textbf{标量乘法公理}： - \textbf{(V5) 标量乘法结合律}：\(\forall \lambda, \mu \in F\)，\(\forall v \in V\)，\(\lambda(\mu v) = (\lambda\mu)v\) - \textbf{(V6) 标量单位元}：\(\forall v \in V\)，\(1 v = v\)（其中 \(1\) 是 \(F\) 的乘法单位元） - \textbf{(V7) 标量对向量加法的分配律}：\(\forall \lambda \in F\)，\(\forall u, v \in V\)，\(\lambda(u + v) = \lambda u + \lambda v\) - \textbf{(V8) 标量对域加法的分配律}：\(\forall \lambda, \mu \in F\)，\(\forall v \in V\)，\((\lambda + \mu)v = \lambda v + \mu v\)

\(V\) 中的元素称为\textbf{向量}，\(F\) 中的元素称为\textbf{标量}。

\textbf{评注 32.1.1}：当 \(F = \mathbb{R}\) 时，称为\textbf{实向量空间}；当 \(F = \mathbb{C}\) 时，称为\textbf{复向量空间}。本卷主要讨论实向量空间，但大部分结论对任意域 \(F\) 均成立。

\textbf{定义 32.1.2（\(F^n\)）}

设 \(F\) 为域，\(n \in \mathbb{N}^+\)。定义

\[
F^n = \{(x_1, x_2, \ldots, x_n) \mid x_i \in F, \, i = 1, 2, \ldots, n\}
\]

在 \(F^n\) 上定义逐分量加法和标量乘法：

\[
(x_1, \ldots, x_n) + (y_1, \ldots, y_n) = (x_1 + y_1, \ldots, x_n + y_n)
\] \[
\lambda(x_1, \ldots, x_n) = (\lambda x_1, \ldots, \lambda x_n)
\]

则 \((F^n, +, \cdot)\) 是 \(F\) 上的向量空间。

\textbf{定理 32.1.1（零元唯一性）}：向量空间 \(V\) 中的零元 \(\mathbf{0}\) 是唯一的。

\textbf{证明}：设 \(\mathbf{0}_1, \mathbf{0}_2\) 均为 \(V\) 的零元。由 (V3)， \[
\mathbf{0}_1 = \mathbf{0}_1 + \mathbf{0}_2 = \mathbf{0}_2 + \mathbf{0}_1 = \mathbf{0}_2
\] 其中第二步用了 (V2)。\(\blacksquare\)

\textbf{定理 32.1.2（加法逆元唯一性）}：对任意 \(v \in V\)，其加法逆元 \((-v)\) 是唯一的。

\textbf{证明}：设 \(u_1, u_2\) 均为 \(v\) 的加法逆元，即 \(v + u_1 = \mathbf{0}\)，\(v + u_2 = \mathbf{0}\)。则 \[
u_1 = u_1 + \mathbf{0} = u_1 + (v + u_2) = (u_1 + v) + u_2 = \mathbf{0} + u_2 = u_2
\] 其中第三步用了 (V1)。\(\blacksquare\)

\textbf{定理 32.1.3（基本运算性质）}：设 \(V\) 为 \(F\) 上的向量空间，则 \(\forall v \in V\)，\(\forall \lambda \in F\)：

\begin{enumerate}
\def\labelenumi{(\roman{enumi})}
\tightlist
\item
  \(0v = \mathbf{0}\)（标量 \(0\) 乘任何向量得零向量）
\item
  \(\lambda \mathbf{0} = \mathbf{0}\)（任何标量乘零向量得零向量）
\item
  \((-1)v = -v\)（\(-1\) 乘向量等于加法逆元）
\item
  若 \(\lambda v = \mathbf{0}\)，则 \(\lambda = 0\) 或 \(v = \mathbf{0}\)
\end{enumerate}

\textbf{证明}：

\begin{enumerate}
\def\labelenumi{(\roman{enumi})}
\item
  由 (V8)：\(0v = (0 + 0)v = 0v + 0v\)。两边加 \((-0v)\) 得 \(0v = \mathbf{0}\)。
\item
  由 (V7)：\(\lambda \mathbf{0} = \lambda(\mathbf{0} + \mathbf{0}) = \lambda \mathbf{0} + \lambda \mathbf{0}\)。两边加 \((-\lambda \mathbf{0})\) 得 \(\lambda \mathbf{0} = \mathbf{0}\)。
\item
  \(v + (-1)v = 1v + (-1)v = (1 + (-1))v = 0v = \mathbf{0}\)。由逆元唯一性，\((-1)v = -v\)。
\item
  设 \(\lambda v = \mathbf{0}\)。若 \(\lambda \neq 0\)，则 \(\lambda^{-1}\) 在 \(F\) 中存在，且 \[
  v = 1v = (\lambda^{-1}\lambda)v = \lambda^{-1}(\lambda v) = \lambda^{-1}\mathbf{0} = \mathbf{0}
  \] 其中最后一步用了 (ii)。\(\blacksquare\)
\end{enumerate}

\hrulefill

\subsection{32.2 子空间}\label{ux5b50ux7a7aux95f4}

\textbf{定义 32.2.1（子空间）}

设 \(V\) 为 \(F\) 上的向量空间，\(W \subseteq V\) 为非空子集。若 \(W\) 在 \(V\) 的加法和标量乘法下封闭，即

\begin{enumerate}
\def\labelenumi{\arabic{enumi}.}
\tightlist
\item
  \(\forall u, v \in W\)，\(u + v \in W\)
\item
  \(\forall \lambda \in F\)，\(\forall v \in W\)，\(\lambda v \in W\)
\end{enumerate}

则称 \(W\) 为 \(V\) 的\textbf{子空间}，记为 \(W \leq V\)。

\textbf{定理 32.2.1（子空间判定）}：\(W \subseteq V\) 是子空间当且仅当

\begin{enumerate}
\def\labelenumi{(\roman{enumi})}
\tightlist
\item
  \(W \neq \varnothing\)（或等价地，\(\mathbf{0} \in W\)）
\item
  \(\forall u, v \in W\)，\(\forall \lambda, \mu \in F\)，\(\lambda u + \mu v \in W\)
\end{enumerate}

\textbf{证明}：

（\(\Rightarrow\)）设 \(W\) 是子空间。由定义 \(W \neq \varnothing\)。取 \(v \in W\)，则 \(0v = \mathbf{0} \in W\)（定理 32.1.3(i)）。对任意 \(u, v \in W\)，\(\lambda u, \mu v \in W\)（标量封闭），故 \(\lambda u + \mu v \in W\)（加法封闭）。

（\(\Leftarrow\)）设条件成立。取 \(\lambda = \mu = 1\) 得加法封闭；取 \(\mu = 0\) 得标量封闭。由 \(W \neq \varnothing\) 并取 \(\lambda = \mu = 0\) 得 \(\mathbf{0} \in W\)，故 \(W\) 非空。\(\blacksquare\)

\textbf{定理 32.2.2（子空间的交）}：设 \(W_1, W_2\) 均为 \(V\) 的子空间，则 \(W_1 \cap W_2\) 也是 \(V\) 的子空间。

\textbf{证明}：\(\mathbf{0} \in W_1 \cap W_2\)（因为 \(\mathbf{0} \in W_1\) 且 \(\mathbf{0} \in W_2\)），故 \(W_1 \cap W_2 \neq \varnothing\)。

设 \(u, v \in W_1 \cap W_2\)，\(\lambda, \mu \in F\)。则 \(u, v \in W_1\) 且 \(u, v \in W_2\)。由于 \(W_1, W_2\) 均为子空间，\(\lambda u + \mu v \in W_1\) 且 \(\lambda u + \mu v \in W_2\)，故 \(\lambda u + \mu v \in W_1 \cap W_2\)。由定理 32.2.1，\(W_1 \cap W_2\) 是子空间。\(\blacksquare\)

\textbf{注记 32.2.1}：子空间的并一般不是子空间。例如 \(\mathbb{R}^2\) 中 \(x\) 轴和 \(y\) 轴都是子空间，但它们的并不是（因为 \((1,0) + (0,1) = (1,1)\) 不在并中）。

\textbf{定义 32.2.2（和与直和）}

设 \(W_1, W_2\) 均为 \(V\) 的子空间。

\begin{enumerate}
\def\labelenumi{\arabic{enumi}.}
\tightlist
\item
  \(W_1\) 与 \(W_2\) 的\textbf{和}定义为 \(W_1 + W_2 = \{w_1 + w_2 \mid w_1 \in W_1, \, w_2 \in W_2\}\)
\item
  若 \(W_1 \cap W_2 = \{\mathbf{0}\}\)，则称和为\textbf{直和}，记为 \(W_1 \oplus W_2\)
\end{enumerate}

\textbf{定理 32.2.3}：\(W_1 + W_2\) 是 \(V\) 的子空间。

\textbf{证明}：\(\mathbf{0} = \mathbf{0} + \mathbf{0} \in W_1 + W_2\)。设 \(u = u_1 + u_2\)，\(v = v_1 + v_2\)（\(u_i, v_i \in W_i\)），\(\lambda, \mu \in F\)，则 \[
\lambda u + \mu v = (\lambda u_1 + \mu v_1) + (\lambda u_2 + \mu v_2) \in W_1 + W_2
\] 由定理 32.2.1，\(W_1 + W_2\) 是子空间。\(\blacksquare\)

\textbf{定理 32.2.4（直和与分解的唯一性）}：\(W_1 + W_2\) 是直和当且仅当每个 \(v \in W_1 + W_2\) 的分解 \(v = w_1 + w_2\)（\(w_i \in W_i\)）是唯一的。

\textbf{证明}：

（\(\Rightarrow\)）设 \(W_1 \cap W_2 = \{\mathbf{0}\}\)。若 \(v = w_1 + w_2 = w_1' + w_2'\)，则 \(w_1 - w_1' = w_2' - w_2\)。左边属于 \(W_1\)，右边属于 \(W_2\)，故该向量属于 \(W_1 \cap W_2 = \{\mathbf{0}\}\)。因此 \(w_1 = w_1'\)，\(w_2 = w_2'\)。

（\(\Leftarrow\)）若分解唯一，设 \(v \in W_1 \cap W_2\)。则 \(v = v + \mathbf{0} = \mathbf{0} + v\) 是 \(v\) 的两个分解，由唯一性 \(v = \mathbf{0}\)。故 \(W_1 \cap W_2 = \{\mathbf{0}\}\)。\(\blacksquare\)

\hrulefill

\subsection{32.3 线性组合与生成}\label{ux7ebfux6027ux7ec4ux5408ux4e0eux751fux6210}

\textbf{定义 32.3.1（线性组合）}

设 \(v_1, v_2, \ldots, v_n \in V\)，\(\lambda_1, \lambda_2, \ldots, \lambda_n \in F\)。称

\[
\sum_{i=1}^{n} \lambda_i v_i = \lambda_1 v_1 + \lambda_2 v_2 + \cdots + \lambda_n v_n
\]

为 \(v_1, \ldots, v_n\) 的一个\textbf{线性组合}。\(\lambda_i\) 称为该组合的\textbf{系数}。

\textbf{定义 32.3.2（生成空间）}

设 \(S = \{v_1, v_2, \ldots, v_n\} \subseteq V\)。\(S\) 的\textbf{生成空间}（或\textbf{张成空间}）定义为 \(S\) 的所有线性组合的集合：

\[
\operatorname{span}(S) = \left\{\sum_{i=1}^{n} \lambda_i v_i \;\middle|\; \lambda_i \in F, \, i = 1, \ldots, n\right\}
\]

\textbf{定理 32.3.1}：\(\operatorname{span}(S)\) 是 \(V\) 的子空间，且是包含 \(S\) 的最小子空间（即任何包含 \(S\) 的子空间都包含 \(\operatorname{span}(S)\)）。

\textbf{证明}：先证是子空间。\(\mathbf{0} = \sum_{i=1}^{n} 0v_i \in \operatorname{span}(S)\)。

设 \(u = \sum \lambda_i v_i\)，\(w = \sum \mu_i v_i \in \operatorname{span}(S)\)，\(\alpha, \beta \in F\)。则 \[
\alpha u + \beta w = \sum (\alpha\lambda_i + \beta\mu_i)v_i \in \operatorname{span}(S)
\] 由定理 32.2.1，\(\operatorname{span}(S)\) 是子空间。

再证最小性。设 \(W\) 是包含 \(S\) 的任意子空间。对任意 \(v \in \operatorname{span}(S)\)，\(v = \sum \lambda_i v_i\)。由于 \(v_i \in S \subseteq W\) 且 \(W\) 对加法和标量乘法封闭，\(v \in W\)。因此 \(\operatorname{span}(S) \subseteq W\)。\(\blacksquare\)

\textbf{定义 32.3.3（有限生成）}

若存在有限集 \(S \subseteq V\) 使得 \(\operatorname{span}(S) = V\)，则称 \(V\) 是\textbf{有限生成}的向量空间。

\textbf{评注 32.3.1}：本卷主要讨论有限生成向量空间（也称为有限维向量空间）。无限维向量空间将在泛函分析（卷十四）中详细讨论。

\hrulefill

\subsection{32.4 线性相关与线性无关}\label{ux7ebfux6027ux76f8ux5173ux4e0eux7ebfux6027ux65e0ux5173}

\textbf{定义 32.4.1（线性相关与线性无关）}

设 \(v_1, v_2, \ldots, v_n \in V\)。

\begin{enumerate}
\def\labelenumi{\arabic{enumi}.}
\tightlist
\item
  若存在不全为零的标量 \(\lambda_1, \ldots, \lambda_n \in F\) 使得 \(\sum_{i=1}^{n} \lambda_i v_i = \mathbf{0}\)，则称向量组 \(\{v_1, \ldots, v_n\}\) \textbf{线性相关}。
\item
  若仅当 \(\lambda_1 = \cdots = \lambda_n = 0\) 时才有 \(\sum_{i=1}^{n} \lambda_i v_i = \mathbf{0}\)，则称向量组 \(\{v_1, \ldots, v_n\}\) \textbf{线性无关}。
\end{enumerate}

\textbf{定理 32.4.1（线性相关的基本性质）}：

\begin{enumerate}
\def\labelenumi{(\roman{enumi})}
\tightlist
\item
  含零向量的向量组必定线性相关。
\item
  若向量组的一个子集线性相关，则整个向量组线性相关。
\item
  \(\{v_1, \ldots, v_n\}\) 线性相关当且仅当存在某个 \(v_k\) 可表示为其余向量的线性组合。
\end{enumerate}

\textbf{证明}：

\begin{enumerate}
\def\labelenumi{(\roman{enumi})}
\item
  设 \(\{v_1, \ldots, v_n\}\) 含 \(\mathbf{0}\)。取 \(\lambda_1 = 1\)（对应 \(\mathbf{0}\)），其余 \(\lambda_i = 0\)，则 \(\sum \lambda_i v_i = \mathbf{0}\) 且系数不全为零，故线性相关。
\item
  设子集 \(\{v_{i_1}, \ldots, v_{i_k}\}\) 线性相关，存在不全为零的系数使其线性组合为零。将其他向量的系数取为 \(0\)，得到整个向量组的线性组合为零且系数不全为零。
\item
  （\(\Rightarrow\)）设 \(\sum \lambda_i v_i = \mathbf{0}\) 且 \(\lambda_k \neq 0\)。则 \(v_k = -\lambda_k^{-1}\sum_{i \neq k} \lambda_i v_i\)。
\end{enumerate}

（\(\Leftarrow\)）设 \(v_k = \sum_{i \neq k} \mu_i v_i\)。则 \(\sum_{i \neq k} \mu_i v_i + (-1)v_k = \mathbf{0}\)，系数不全为零（\(v_k\) 的系数为 \(-1 \neq 0\)）。\(\blacksquare\)

\textbf{定理 32.4.2（线性无关组的扩充性质）}：设 \(\{v_1, \ldots, v_n\}\) 线性无关，\(v \in V\)。则 \(\{v_1, \ldots, v_n, v\}\) 线性相关当且仅当 \(v \in \operatorname{span}\{v_1, \ldots, v_n\}\)。

\textbf{证明}：

（\(\Rightarrow\)）若 \(\{v_1, \ldots, v_n, v\}\) 线性相关，则存在不全为零的 \(\lambda_1, \ldots, \lambda_n, \lambda\) 使得 \(\sum \lambda_i v_i + \lambda v = \mathbf{0}\)。若 \(\lambda = 0\)，则 \(\sum \lambda_i v_i = \mathbf{0}\) 且 \(\lambda_i\) 不全为零，与 \(\{v_1, \ldots, v_n\}\) 线性无关矛盾。故 \(\lambda \neq 0\)，从而 \(v = -\lambda^{-1}\sum \lambda_i v_i \in \operatorname{span}\{v_1, \ldots, v_n\}\)。

（\(\Leftarrow\)）若 \(v \in \operatorname{span}\{v_1, \ldots, v_n\}\)，则 \(v = \sum \mu_i v_i\)。于是 \(\sum \mu_i v_i + (-1)v = \mathbf{0}\)，系数不全为零，故线性相关。\(\blacksquare\)

\hrulefill

\subsection{32.5 基与维数}\label{ux57faux4e0eux7ef4ux6570}

\textbf{定义 32.5.1（基）}

设 \(B = \{e_1, e_2, \ldots, e_n\} \subseteq V\)。若满足：

\begin{enumerate}
\def\labelenumi{\arabic{enumi}.}
\tightlist
\item
  \textbf{线性无关性}：\(B\) 是线性无关的
\item
  \textbf{生成性}：\(\operatorname{span}(B) = V\)
\end{enumerate}

则称 \(B\) 为 \(V\) 的一组\textbf{基}（或\textbf{基底}）。

\textbf{定理 32.5.1（基的等价刻画）}：\(B = \{e_1, \ldots, e_n\}\) 是 \(V\) 的基当且仅当每个 \(v \in V\) 可唯一表示为 \(B\) 中向量的线性组合。

\textbf{证明}：

（\(\Rightarrow\)）由生成性，\(v\) 可表示为 \(v = \sum \lambda_i e_i\)。若还存在 \(v = \sum \mu_i e_i\)，则 \(\sum (\lambda_i - \mu_i)e_i = \mathbf{0}\)。由线性无关性，\(\lambda_i - \mu_i = 0\)，即 \(\lambda_i = \mu_i\)，故表示唯一。

（\(\Leftarrow\)）若每个 \(v\) 可唯一表示，则生成性自动成立。取 \(v = \mathbf{0}\)，由唯一性，仅有 \(\lambda_i = 0\) 可使 \(\sum \lambda_i e_i = \mathbf{0}\)，故 \(B\) 线性无关。\(\blacksquare\)

\textbf{定义 32.5.2（坐标）}：设 \(B = \{e_1, \ldots, e_n\}\) 为 \(V\) 的基，\(v = \sum_{i=1}^{n} \lambda_i e_i\)。则 \((\lambda_1, \ldots, \lambda_n) \in F^n\) 称为 \(v\) 关于基 \(B\) 的\textbf{坐标}。

\textbf{定理 32.5.2（基的存在性）}：任何非零有限生成向量空间都存在基。

\textbf{证明}：设 \(V = \operatorname{span}\{v_1, \ldots, v_m\}\)。构造算法如下：

从 \(i = 1\) 到 \(m\)，若 \(v_i\) 不属于前 \(i-1\) 个向量的生成空间，则保留 \(v_i\)；否则丢弃。设保留的向量为 \(\{e_1, \ldots, e_n\}\)。

由构造过程，每个 \(v_i\) 都属于 \(\operatorname{span}\{e_1, \ldots, e_n\}\)，故 \(V \subseteq \operatorname{span}\{e_1, \ldots, e_n\}\)，生成性成立。

若 \(\sum \lambda_j e_j = \mathbf{0}\) 且 \(\lambda_k\) 是最后一个非零系数，则 \(e_k = -\lambda_k^{-1}\sum_{j < k} \lambda_j e_j\)，说明 \(e_k\) 属于前 \(k-1\) 个保留向量的生成空间，与构造中保留 \(e_k\) 的条件矛盾。故所有 \(\lambda_j = 0\)，线性无关性成立。\(\blacksquare\)

\textbf{注记 32.5.1}：上述构造证明给出了从有限生成集中提取基的算法。对于有限维向量空间，基的存在性由上述构造直接给出。对于无限维向量空间，基的存在性需要选择公理（Axiom of Choice，等价于 Zorn 引理）：考虑 \(V\) 的所有线性无关子集构成的偏序集（按包含关系），Zorn 引理保证存在极大线性无关子集，该极大集即为 \(V\) 的一组基（参见 V20（概率论2））。

\textbf{引理 32.5.3（替换引理，Steinitz 引理）}：设 \(\{e_1, \ldots, e_n\}\) 为 \(V\) 的基，\(\{v_1, \ldots, v_m\} \subseteq V\) 线性无关。则 \(m \leq n\)，且可以适当地用 \(v_1, \ldots, v_m\) 替换基中的 \(m\) 个向量，使得替换后的集合仍为基。

\textbf{证明}：对 \(m\) 归纳。

\textbf{基础} \(m = 0\)：平凡。

\textbf{归纳}：假设结论对 \(m-1\) 成立。由归纳假设，\(m-1 \leq n\)，且可替换 \(m-1\) 个基向量，得到新基 \(\{v_1, \ldots, v_{m-1}, f_m, \ldots, f_n\}\)（其中 \(\{f_m, \ldots, f_n\}\) 是未被替换的原基向量）。

\(v_m\) 可表示为新基的线性组合：\(v_m = \sum_{i=1}^{m-1} \alpha_i v_i + \sum_{j=m}^{n} \beta_j f_j\)。由于 \(\{v_1, \ldots, v_m\}\) 线性无关，\(v_m\) 不能仅由 \(v_1, \ldots, v_{m-1}\) 表示，故存在 \(k \geq m\) 使得 \(\beta_k \neq 0\)。

不妨设 \(k = m\)（重排 \(f_m, \ldots, f_n\)）。则 \(f_m = \beta_m^{-1}(v_m - \sum_{i=1}^{m-1} \alpha_i v_i - \sum_{j=m+1}^{n} \beta_j f_j)\)，故 \(f_m \in \operatorname{span}\{v_1, \ldots, v_m, f_{m+1}, \ldots, f_n\}\)。因此这 \(n\) 个向量生成 \(V\)，且数量为 \(n\)。由定理 32.5.2 的构造可知它们线性无关（否则可提取更小的基，与 \(n\) 已经是基的大小矛盾）。故 \(\{v_1, \ldots, v_m, f_{m+1}, \ldots, f_n\}\) 是基。

特别地，\(m \leq n\)（否则 \(m > n\) 时由归纳假设 \(m-1 \geq n\)，矛盾）。\(\blacksquare\)

\textbf{定理 32.5.4（维数的不变性）}：有限生成向量空间 \(V \neq \{\mathbf{0}\}\) 的任何两组基具有相同个数的向量。

\textbf{证明}：设 \(B_1 = \{e_1, \ldots, e_n\}\) 和 \(B_2 = \{f_1, \ldots, f_m\}\) 均为 \(V\) 的基。\(B_1\) 是基，\(B_2\) 线性无关，由替换引理 \(m \leq n\)。交换角色得 \(n \leq m\)。故 \(n = m\)。\(\blacksquare\)

\textbf{定义 32.5.3（维数）}：向量空间 \(V\) 的基所含向量的个数称为 \(V\) 的\textbf{维数}，记为 \(\dim V\)。约定 \(\dim\{\mathbf{0}\} = 0\)。

\textbf{定理 32.5.5（维数的基本性质）}：设 \(V\) 为 \(n\) 维向量空间。

\begin{enumerate}
\def\labelenumi{(\roman{enumi})}
\tightlist
\item
  \(V\) 中任何 \(n+1\) 个向量必线性相关。
\item
  \(V\) 中任何 \(n\) 个线性无关的向量构成基。
\item
  \(V\) 中任何生成 \(V\) 的 \(n\) 个向量构成基。
\item
  设 \(W \leq V\)，则 \(\dim W \leq \dim V\)。等号成立当且仅当 \(W = V\)。
\end{enumerate}

\textbf{证明}：

\begin{enumerate}
\def\labelenumi{(\roman{enumi})}
\item
  由替换引理，若存在 \(n+1\) 个线性无关向量，则 \(n+1 \leq n\)，矛盾。
\item
  设 \(S = \{v_1, \ldots, v_n\}\) 线性无关。若 \(\operatorname{span}(S) \neq V\)，则存在 \(v \in V \setminus \operatorname{span}(S)\)。由定理 32.4.2，\(\{v_1, \ldots, v_n, v\}\) 线性无关，与 (i) 矛盾。故 \(\operatorname{span}(S) = V\)。
\item
  设 \(S\) 生成 \(V\)。若 \(S\) 线性相关，则可从中提取一个更小的线性无关集仍生成 \(V\)，即基的大小小于 \(n\)，与 \(\dim V = n\) 矛盾。故 \(S\) 线性无关。
\item
  设 \(B_W\) 为 \(W\) 的基。\(B_W\) 是 \(V\) 中线性无关的向量组，可扩充为 \(V\) 的基。故 \(\dim W \leq \dim V\)。若 \(\dim W = \dim V = n\)，则 \(B_W\) 本身就是 \(V\) 的基（\(n\) 个线性无关向量），故 \(W = V\)。\(\blacksquare\)
\end{enumerate}

\textbf{定理 32.5.6（维数公式）}：设 \(W_1, W_2 \leq V\)，则

\[
\dim(W_1 + W_2) = \dim W_1 + \dim W_2 - \dim(W_1 \cap W_2)
\]

\textbf{证明}：设 \(\dim(W_1 \cap W_2) = r\)，\(\dim W_1 = r + s\)，\(\dim W_2 = r + t\)。

取 \(W_1 \cap W_2\) 的基 \(\{e_1, \ldots, e_r\}\)。将其扩充为 \(W_1\) 的基 \(\{e_1, \ldots, e_r, u_1, \ldots, u_s\}\)，也扩充为 \(W_2\) 的基 \(\{e_1, \ldots, e_r, w_1, \ldots, w_t\}\)。

断言 \(\{e_1, \ldots, e_r, u_1, \ldots, u_s, w_1, \ldots, w_t\}\) 是 \(W_1 + W_2\) 的基。

生成性：\(W_1 + W_2\) 中任意向量可写为 \(W_1\) 中向量与 \(W_2\) 中向量的和，而 \(W_1\) 中向量由 \(e_i, u_j\) 生成，\(W_2\) 中向量由 \(e_i, w_k\) 生成，故该集合生成 \(W_1 + W_2\)。

线性无关性：设 \(\sum \alpha_i e_i + \sum \beta_j u_j + \sum \gamma_k w_k = \mathbf{0}\)。令 \(v = \sum \alpha_i e_i + \sum \beta_j u_j = -\sum \gamma_k w_k\)。左边属于 \(W_1\)，右边属于 \(W_2\)，故 \(v \in W_1 \cap W_2\)。因此 \(v\) 可由 \(e_1, \ldots, e_r\) 线性表示：\(v = \sum \delta_i e_i\)。

于是 \(\sum \delta_i e_i + \sum \gamma_k w_k = \mathbf{0}\)。由于 \(\{e_1, \ldots, e_r, w_1, \ldots, w_t\}\) 是 \(W_2\) 的基，所有系数为零：\(\delta_i = 0\)，\(\gamma_k = 0\)。从而 \(v = \mathbf{0}\)，即 \(\sum \alpha_i e_i + \sum \beta_j u_j = \mathbf{0}\)。由 \(W_1\) 的基的线性无关性，\(\alpha_i = 0\)，\(\beta_j = 0\)。

因此 \(\dim(W_1 + W_2) = r + s + t = (r + s) + (r + t) - r = \dim W_1 + \dim W_2 - \dim(W_1 \cap W_2)\)。\(\blacksquare\)

\textbf{推论 32.5.7（直和的维数）}：若 \(W_1 \cap W_2 = \{\mathbf{0}\}\)，则 \(\dim(W_1 \oplus W_2) = \dim W_1 + \dim W_2\)。

\textbf{证明}：由维数公式，\(\dim(W_1 \cap W_2) = 0\) 时 \(\dim(W_1 + W_2) = \dim W_1 + \dim W_2\)。\(\blacksquare\)

\hrulefill

\subsection{32.6 坐标映射与同构}\label{ux5750ux6807ux6620ux5c04ux4e0eux540cux6784}

\textbf{定义 32.6.1（坐标映射）}：设 \(B = \{e_1, \ldots, e_n\}\) 为 \(V\) 的基。定义映射 \(\varphi_B : V \to F^n\)，将 \(v = \sum \lambda_i e_i\) 映为其坐标 \((\lambda_1, \ldots, \lambda_n)\)。

\textbf{定理 32.6.1}：\(\varphi_B\) 是双射，且满足 \(\varphi_B(u + v) = \varphi_B(u) + \varphi_B(v)\)，\(\varphi_B(\lambda v) = \lambda \varphi_B(v)\)。

\textbf{证明}：由基的等价刻画（定理 32.5.1），表示唯一，故 \(\varphi_B\) 是良定义的双射。线性性质由坐标的加法与标量乘法的定义直接验证。\(\blacksquare\)

\textbf{定义 32.6.2（同构）}：设 \(V, W\) 为 \(F\) 上的向量空间。若存在双射 \(T : V \to W\) 满足 \(\forall u, v \in V, \forall \lambda \in F\)： \[
T(u + v) = T(u) + T(v), \quad T(\lambda u) = \lambda T(u)
\] 则称 \(T\) 为\textbf{同构映射}，\(V\) 与 \(W\) \textbf{同构}，记为 \(V \cong W\)。

\textbf{定理 32.6.2}：\(F\) 上的任何 \(n\) 维向量空间与 \(F^n\) 同构。

\textbf{证明}：坐标映射 \(\varphi_B : V \to F^n\) 即为同构映射。\(\blacksquare\)

\textbf{推论 32.6.3}：\(F\) 上的两个有限维向量空间同构当且仅当它们维数相同。

\textbf{证明}：若维数相同，都与 \(F^n\) 同构，故互相同构。若同构，则基的像也是基，故维数相同。\(\blacksquare\)

\hrulefill

\hrulefill

\hrulefill

\hrulefill

\hrulefill

\newpage

\chapter{05-代数/01-线性代数/02-线性映射与矩阵.md}\label{ux4ee3ux657001-ux7ebfux6027ux4ee3ux657002-ux7ebfux6027ux6620ux5c04ux4e0eux77e9ux9635.md}

\section{第102章：线性映射与矩阵}\label{ux7b2c102ux7ae0ux7ebfux6027ux6620ux5c04ux4e0eux77e9ux9635}

线性映射是向量空间之间保持线性结构的映射，它是线性代数中最核心的概念之一。如果说向量空间公理化提供了线性代数的''舞台''，那么线性映射就是这个舞台上的''演员''------它们描述了向量空间之间的相互关系，刻画了线性结构的保持与变换。从历史角度看，矩阵先于线性映射出现在数学舞台上：Gauss 在 1801 年《算术研究》中已使用二元二次型的矩阵表示，而 Cayley 于 1858 年正式引入矩阵代数。然而，随着抽象代数的发展，人们逐渐认识到矩阵不过是线性映射在选定基下的坐标表示------这一观点由 Frobenius 和 Sylvester 在 19 世纪末明确阐述。

线性映射与矩阵的对应关系是有限维线性代数中最基本的对偶性之一。给定域 \(F\) 上的两个有限维向量空间 \(V\) 和 \(W\)，分别取定基后，每个线性映射 \(T: V \to W\) 对应一个唯一的矩阵 \(A \in M_{m \times n}(F)\)（其中 \(n = \dim V\)，\(m = \dim W\)）。反之，每个矩阵唯一确定一个线性映射。这种对应不仅是集合的双射，而且保持代数结构：线性映射的加法对应矩阵加法，线性映射的复合对应矩阵乘法，标量乘法对应标量乘矩阵。用范畴论的语言，这给出了从有限维向量空间范畴到矩阵范畴的忠实满函子。

本章的核心主题包括：线性映射的核与像，它们揭示了线性映射的单性与满性；秩-零化度定理（\(\dim V = \dim \ker T + \dim \operatorname{im} T\)），它是线性方程组理论的基本定理；基变换与矩阵相似性，它回答了''同一线性映射在不同基下如何表示''这一关键问题。线性映射的语言使得我们可以将矩阵的许多性质（如秩、行列式、特征值）提升为不依赖于基选择的内在性质，这是现代线性代数区别于古典矩阵计算的根本特征。

\textbf{定义 33.1.1（线性映射）}

设 \(V, W\) 为域 \(F\) 上的向量空间。映射 \(T : V \to W\) 若满足：

\begin{enumerate}
\def\labelenumi{\arabic{enumi}.}
\tightlist
\item
  \textbf{加法保持}：\(\forall u, v \in V\)，\(T(u + v) = T(u) + T(v)\)
\item
  \textbf{标量乘法保持}：\(\forall \lambda \in F, \forall v \in V\)，\(T(\lambda v) = \lambda T(v)\)
\end{enumerate}

则称 \(T\) 为\textbf{线性映射}（或\textbf{线性变换}）。当 \(V = W\) 时，\(T\) 特称为 \(V\) 上的\textbf{线性算子}。

\textbf{定理 33.1.1（线性映射的基本性质）}：设 \(T : V \to W\) 为线性映射，则

\begin{enumerate}
\def\labelenumi{(\roman{enumi})}
\tightlist
\item
  \(T(\mathbf{0}_V) = \mathbf{0}_W\)
\item
  \(T(-v) = -T(v)\)
\item
  \(T\left(\sum_{i=1}^{n} \lambda_i v_i\right) = \sum_{i=1}^{n} \lambda_i T(v_i)\)
\end{enumerate}

\textbf{证明}：

\begin{enumerate}
\def\labelenumi{(\roman{enumi})}
\item
  \(T(\mathbf{0}_V) = T(0 \cdot \mathbf{0}_V) = 0 \cdot T(\mathbf{0}_V) = \mathbf{0}_W\)。
\item
  \(T(-v) = T((-1)v) = (-1)T(v) = -T(v)\)。
\item
  对 \(n\) 归纳，利用线性映射的两条公理。\(\blacksquare\)
\end{enumerate}

\textbf{定义 33.1.2（线性映射的运算）}：设 \(S, T : V \to W\) 为线性映射，\(\lambda \in F\)。

\begin{enumerate}
\def\labelenumi{\arabic{enumi}.}
\tightlist
\item
  \textbf{加法}：\((S + T)(v) = S(v) + T(v)\)
\item
  \textbf{标量乘法}：\((\lambda T)(v) = \lambda T(v)\)
\item
  \textbf{复合}（设 \(T : V \to W\)，\(S : W \to U\)）：\((S \circ T)(v) = S(T(v))\)
\end{enumerate}

\textbf{定理 33.1.2}：所有从 \(V\) 到 \(W\) 的线性映射的集合，记为 \(\mathcal{L}(V, W)\)，在上述加法和标量乘法下构成 \(F\) 上的向量空间。

\textbf{证明}：逐条验证向量空间公理。零元为零映射 \(0(v) = \mathbf{0}_W\)。加法逆元为 \((-T)(v) = -T(v)\)。\(\blacksquare\)

\textbf{定义 33.1.3（同构与自同构）}：

\begin{itemize}
\tightlist
\item
  可逆的线性映射称为\textbf{同构映射}（其逆也自动是线性的）
\item
  \(V\) 到自身的同构称为\textbf{自同构}，所有自同构构成群 \(\operatorname{GL}(V)\)
\end{itemize}

\textbf{定理 33.1.3（逆映射的线性性）}：若线性映射 \(T : V \to W\) 可逆，则 \(T^{-1} : W \to V\) 也是线性映射。

\textbf{证明}：设 \(w_1, w_2 \in W\)，则存在 \(v_1, v_2 \in V\) 使得 \(T(v_i) = w_i\)。于是 \[
T^{-1}(w_1 + w_2) = T^{-1}(T(v_1) + T(v_2)) = T^{-1}(T(v_1 + v_2)) = v_1 + v_2 = T^{-1}(w_1) + T^{-1}(w_2)
\] 同理可证标量乘法的保持。\(\blacksquare\)

\hrulefill

\subsection{33.2 核与像}\label{ux6838ux4e0eux50cf}

\textbf{定义 33.2.1（核与像）}：设 \(T : V \to W\) 为线性映射。

\begin{enumerate}
\def\labelenumi{\arabic{enumi}.}
\tightlist
\item
  \(T\) 的\textbf{核}：\(\ker T = \{v \in V \mid T(v) = \mathbf{0}_W\}\)
\item
  \(T\) 的\textbf{像}：\(\operatorname{im} T = \{T(v) \mid v \in V\}\)
\end{enumerate}

\textbf{定理 33.2.1}：\(\ker T\) 是 \(V\) 的子空间，\(\operatorname{im} T\) 是 \(W\) 的子空间。

\textbf{证明}：\(\mathbf{0}_V \in \ker T\)（因为 \(T(\mathbf{0}_V) = \mathbf{0}_W\)）。设 \(u, v \in \ker T\)，\(\lambda, \mu \in F\)，则 \(T(\lambda u + \mu v) = \lambda T(u) + \mu T(v) = \mathbf{0}_W + \mathbf{0}_W = \mathbf{0}_W\)，故 \(\lambda u + \mu v \in \ker T\)。由定理 32.2.1，\(\ker T\) 是子空间。

类似，\(\mathbf{0}_W = T(\mathbf{0}_V) \in \operatorname{im} T\)。设 \(T(u), T(v) \in \operatorname{im} T\)，则 \(\lambda T(u) + \mu T(v) = T(\lambda u + \mu v) \in \operatorname{im} T\)。\(\blacksquare\)

\textbf{定义 33.2.2（秩与零度）}：设 \(T : V \to W\) 为线性映射。

\begin{enumerate}
\def\labelenumi{\arabic{enumi}.}
\tightlist
\item
  \(\operatorname{rank} T = \dim(\operatorname{im} T)\)，称为 \(T\) 的\textbf{秩}
\item
  \(\operatorname{nullity} T = \dim(\ker T)\)，称为 \(T\) 的\textbf{零度}
\end{enumerate}

\textbf{定理 33.2.2（单射的判定）}：线性映射 \(T : V \to W\) 是单射当且仅当 \(\ker T = \{\mathbf{0}\}\)。

\textbf{证明}：

（\(\Rightarrow\)）若 \(T\) 是单射，\(T(v) = \mathbf{0} = T(\mathbf{0})\)，则 \(v = \mathbf{0}\)，故 \(\ker T = \{\mathbf{0}\}\)。

（\(\Leftarrow\)）若 \(\ker T = \{\mathbf{0}\}\) 且 \(T(u) = T(v)\)，则 \(T(u - v) = \mathbf{0}\)，故 \(u - v \in \ker T = \{\mathbf{0}\}\)，于是 \(u = v\)。\(\blacksquare\)

\hrulefill

\subsection{33.3 秩-零化度定理}\label{ux79e9-ux96f6ux5316ux5ea6ux5b9aux7406}

\textbf{定理 33.3.1（秩-零化度定理）}：设 \(V\) 为有限维向量空间，\(T : V \to W\) 为线性映射。则

\[
\dim V = \dim(\ker T) + \dim(\operatorname{im} T) = \operatorname{nullity} T + \operatorname{rank} T
\]

\textbf{证明}：设 \(\dim V = n\)，\(\dim(\ker T) = k\)。取 \(\ker T\) 的基 \(\{v_1, \ldots, v_k\}\)，将其扩充为 \(V\) 的基 \(\{v_1, \ldots, v_k, u_1, \ldots, u_{n-k}\}\)。

断言 \(\{T(u_1), \ldots, T(u_{n-k})\}\) 是 \(\operatorname{im} T\) 的基。

生成性：对任意 \(v \in V\)，\(v = \sum \alpha_i v_i + \sum \beta_j u_j\)，则 \(T(v) = \sum \beta_j T(u_j)\)（因为 \(T(v_i) = \mathbf{0}\)）。故 \(\operatorname{im} T = \operatorname{span}\{T(u_1), \ldots, T(u_{n-k})\}\)。

线性无关性：设 \(\sum \gamma_j T(u_j) = \mathbf{0}\)。则 \(T(\sum \gamma_j u_j) = \mathbf{0}\)，故 \(\sum \gamma_j u_j \in \ker T\)。因此 \(\sum \gamma_j u_j = \sum \delta_i v_i\)，即 \(\sum (-\delta_i) v_i + \sum \gamma_j u_j = \mathbf{0}\)。由基的线性无关性，所有系数为零，特别地 \(\gamma_j = 0\)（\(j = 1, \ldots, n-k\)）。

因此 \(\dim(\operatorname{im} T) = n - k\)，即 \(\dim V = \dim(\ker T) + \dim(\operatorname{im} T)\)。\(\blacksquare\)

\textbf{推论 33.3.2（满射与单射的条件）}：设 \(T : V \to W\) 为线性映射，\(\dim V = \dim W = n\)。

\begin{enumerate}
\def\labelenumi{(\roman{enumi})}
\tightlist
\item
  \(T\) 是单射 \(\iff\) \(T\) 是满射 \(\iff\) \(T\) 是同构。
\item
  \(\operatorname{rank} T = n \iff \ker T = \{\mathbf{0}\}\)。
\end{enumerate}

\textbf{证明}：

\begin{enumerate}
\def\labelenumi{(\roman{enumi})}
\item
  由秩-零化度定理，\(\dim V = \dim(\ker T) + \dim(\operatorname{im} T)\)。\(T\) 单射 \(\iff \ker T = \{\mathbf{0}\} \iff \dim(\ker T) = 0 \iff \dim(\operatorname{im} T) = n \iff \operatorname{im} T = W \iff T\) 满射。此时 \(T\) 既单且满，即为同构。
\item
  \(\operatorname{rank} T = n \iff \dim(\operatorname{im} T) = n \iff \dim(\ker T) = 0 \iff \ker T = \{\mathbf{0}\}\)。\(\blacksquare\)
\end{enumerate}

\hrulefill

\subsection{33.4 线性映射的矩阵表示}\label{ux7ebfux6027ux6620ux5c04ux7684ux77e9ux9635ux8868ux793a}

\textbf{定义 33.4.1（矩阵）}：设 \(F\) 为域，\(m, n \in \mathbb{N}^+\)。一个 \(m \times n\) \textbf{矩阵} \(A\) 是一个由 \(mn\) 个 \(F\) 中元素组成的矩形阵列：

\[
A = \begin{pmatrix}
a_{11} & a_{12} & \cdots & a_{1n} \\
a_{21} & a_{22} & \cdots & a_{2n} \\
\vdots & \vdots & \ddots & \vdots \\
a_{m1} & a_{m2} & \cdots & a_{mn}
\end{pmatrix}
\]

其中 \(a_{ij} \in F\) 表示第 \(i\) 行第 \(j\) 列的元素。所有 \(m \times n\) 矩阵的集合记为 \(M_{m \times n}(F)\)。

\textbf{定义 33.4.2（矩阵的加法和标量乘法）}：在 \(M_{m \times n}(F)\) 上定义：

\begin{enumerate}
\def\labelenumi{\arabic{enumi}.}
\tightlist
\item
  \((A + B)_{ij} = a_{ij} + b_{ij}\)
\item
  \((\lambda A)_{ij} = \lambda a_{ij}\)
\end{enumerate}

\textbf{定理 33.4.1}：\(M_{m \times n}(F)\) 在上述运算下构成 \(F\) 上的向量空间，且 \(\dim M_{m \times n}(F) = mn\)。

\textbf{证明}：向量空间公理由逐分量的域运算直接验证。标准基为 \(\{E_{ij} \mid 1 \leq i \leq m, \, 1 \leq j \leq n\}\)，其中 \(E_{ij}\) 在第 \((i, j)\) 位置为 \(1\)，其余为 \(0\)。共有 \(mn\) 个基向量。\(\blacksquare\)

\textbf{定义 33.4.3（线性映射的矩阵）}：设 \(V\) 有基 \(B = \{v_1, \ldots, v_n\}\)，\(W\) 有基 \(C = \{w_1, \ldots, w_m\}\)，\(T \in \mathcal{L}(V, W)\)。

对每个 \(j = 1, \ldots, n\)，将 \(T(v_j)\) 表示为 \(C\) 的线性组合：

\[
T(v_j) = \sum_{i=1}^{m} a_{ij} w_i
\]

则 \(m \times n\) 矩阵 \(A = (a_{ij})\) 称为 \(T\) 关于基 \(B\) 和 \(C\) 的\textbf{矩阵表示}，记为 \([T]_{B}^{C}\) 或 \(\mathcal{M}(T; B, C)\)。

\textbf{定理 33.4.2（矩阵表示与坐标的关系）}：设 \(v \in V\) 关于基 \(B\) 的坐标为 \(x = (x_1, \ldots, x_n)^T\)（列向量），则 \(T(v)\) 关于基 \(C\) 的坐标 \(y = (y_1, \ldots, y_m)^T\) 满足

\[
y = A x, \quad \text{即} \quad y_i = \sum_{j=1}^{n} a_{ij} x_j
\]

\textbf{证明}： \[
T(v) = T\left(\sum_{j=1}^{n} x_j v_j\right) = \sum_{j=1}^{n} x_j T(v_j) = \sum_{j=1}^{n} x_j \sum_{i=1}^{m} a_{ij} w_i = \sum_{i=1}^{m} \left(\sum_{j=1}^{n} a_{ij} x_j\right) w_i
\] 因此 \(y_i = \sum_{j=1}^{n} a_{ij} x_j\)。\(\blacksquare\)

\textbf{定理 33.4.3（矩阵表示是线性同构）}：映射 \(\mathcal{M}(\cdot; B, C) : \mathcal{L}(V, W) \to M_{m \times n}(F)\) 是向量空间同构。特别地，\(\dim \mathcal{L}(V, W) = (\dim V)(\dim W)\)。

\textbf{证明}：线性和双射由矩阵表示的定义直接验证。\(\blacksquare\)

\hrulefill

\subsection{33.5 矩阵乘法与复合}\label{ux77e9ux9635ux4e58ux6cd5ux4e0eux590dux5408}

\textbf{定义 33.5.1（矩阵乘法）}：设 \(A \in M_{m \times n}(F)\)，\(B \in M_{n \times p}(F)\)。定义乘积 \(C = AB \in M_{m \times p}(F)\) 为

\[
c_{ij} = \sum_{k=1}^{n} a_{ik} b_{kj}, \quad i = 1, \ldots, m, \; j = 1, \ldots, p
\]

\textbf{定理 33.5.1（矩阵乘法对应线性映射的复合）}：设 \(V \xrightarrow{S} W \xrightarrow{T} U\) 为线性映射，\(B, C, D\) 分别为 \(V, W, U\) 的基。则

\[
[T \circ S]_{B}^{D} = [T]_{C}^{D} \cdot [S]_{B}^{C}
\]

\textbf{证明}：设 \(B = \{v_1, \ldots, v_n\}\)，\(C = \{w_1, \ldots, w_m\}\)，\(D = \{u_1, \ldots, u_p\}\)。令 \(A = [T]_{C}^{D}\)，\(B = [S]_{B}^{C}\)。则 \[
(T \circ S)(v_j) = T(S(v_j)) = T\left(\sum_{k=1}^{m} b_{kj} w_k\right) = \sum_{k=1}^{m} b_{kj} T(w_k) = \sum_{k=1}^{m} b_{kj} \sum_{i=1}^{p} a_{ik} u_i = \sum_{i=1}^{p} \left(\sum_{k=1}^{m} a_{ik} b_{kj}\right) u_i
\] 因此 \([T \circ S]_{B}^{D}\) 的第 \((i, j)\) 元为 \(\sum_{k=1}^{m} a_{ik} b_{kj}\)，正是 \(AB\) 的第 \((i, j)\) 元。\(\blacksquare\)

\textbf{定理 33.5.2（矩阵乘法的基本性质）}：设矩阵的维数使得以下乘法有意义，则

\begin{enumerate}
\def\labelenumi{(\roman{enumi})}
\tightlist
\item
  \textbf{结合律}：\((AB)C = A(BC)\)
\item
  \textbf{分配律}：\(A(B + C) = AB + AC\)，\((A + B)C = AC + BC\)
\item
  \textbf{标量乘法}：\(\lambda(AB) = (\lambda A)B = A(\lambda B)\)
\item
  \textbf{单位矩阵}：\(I_n A = A I_m = A\)（\(A \in M_{m \times n}(F)\)，\(I_n\) 为 \(n\) 阶单位矩阵）
\end{enumerate}

\textbf{证明}：(i) 由矩阵乘法对应线性映射的复合，复合满足结合律。也可直接验证： \[
((AB)C)_{ij} = \sum_{k} (AB)_{ik} c_{kj} = \sum_{k} \sum_{\ell} a_{i\ell} b_{\ell k} c_{kj} = \sum_{\ell} a_{i\ell} \sum_{k} b_{\ell k} c_{kj} = (A(BC))_{ij}
\]

(ii)-(iv) 直接验证。\(\blacksquare\)

\textbf{注记 33.5.1}：矩阵乘法一般不满足交换律，即 \(AB \neq BA\) 一般成立。

\textbf{定义 33.5.2（转置）}：设 \(A \in M_{m \times n}(F)\)。\(A\) 的\textbf{转置} \(A^T \in M_{n \times m}(F)\) 定义为 \((A^T)_{ij} = A_{ji}\)。

\textbf{定理 33.5.3（转置的性质）}：

\begin{enumerate}
\def\labelenumi{(\roman{enumi})}
\tightlist
\item
  \((A^T)^T = A\)
\item
  \((A + B)^T = A^T + B^T\)
\item
  \((\lambda A)^T = \lambda A^T\)
\item
  \((AB)^T = B^T A^T\)
\end{enumerate}

\textbf{证明}：(i)-(iii) 直接验证。(iv) \((AB)^T_{ij} = (AB)_{ji} = \sum_k a_{jk} b_{ki} = \sum_k (B^T)_{ik} (A^T)_{kj} = (B^T A^T)_{ij}\)。\(\blacksquare\)

\hrulefill

\subsection{33.6 基变换}\label{ux57faux53d8ux6362}

\textbf{定义 33.6.1（基变换矩阵）}：设 \(V\) 有两组基 \(B = \{v_1, \ldots, v_n\}\) 和 \(B' = \{v_1', \ldots, v_n'\}\)。将 \(B'\) 中的每个向量用 \(B\) 表示：

\[
v_j' = \sum_{i=1}^{n} p_{ij} v_i, \quad j = 1, \ldots, n
\]

则矩阵 \(P = (p_{ij})\) 称为从 \(B\) 到 \(B'\) 的\textbf{基变换矩阵}（或\textbf{过渡矩阵}）。

\textbf{定理 33.6.1（基变换矩阵的性质）}：

\begin{enumerate}
\def\labelenumi{(\roman{enumi})}
\tightlist
\item
  \(P\) 可逆，且 \(P^{-1}\) 是从 \(B'\) 到 \(B\) 的基变换矩阵。
\item
  设 \(v \in V\) 关于 \(B\) 的坐标为 \(x\)，关于 \(B'\) 的坐标为 \(x'\)，则 \(x = P x'\)。
\end{enumerate}

\textbf{证明}：

\begin{enumerate}
\def\labelenumi{(\roman{enumi})}
\item
  设 \(Q\) 为从 \(B'\) 到 \(B\) 的基变换矩阵，即将 \(v_i\) 用 \(B'\) 表示。则 \(v_j = \sum_k q_{kj} v_k' = \sum_k q_{kj} \sum_i p_{ik} v_i = \sum_i (\sum_k p_{ik} q_{kj}) v_i\)。由基表示的唯一性，\(\sum_k p_{ik} q_{kj} = \delta_{ij}\)，即 \(PQ = I_n\)。同理 \(QP = I_n\)。故 \(P\) 可逆且 \(P^{-1} = Q\)。
\item
  \(v = \sum_j x_j' v_j' = \sum_j x_j' \sum_i p_{ij} v_i = \sum_i (\sum_j p_{ij} x_j') v_i\)，故 \(x_i = \sum_j p_{ij} x_j'\)，即 \(x = P x'\)。\(\blacksquare\)
\end{enumerate}

\textbf{定理 33.6.2（线性映射在不同基下的矩阵）}：设 \(T \in \mathcal{L}(V, W)\)，\(B, B'\) 为 \(V\) 的基，\(C, C'\) 为 \(W\) 的基。\(P\) 为从 \(B\) 到 \(B'\) 的基变换矩阵，\(Q\) 为从 \(C\) 到 \(C'\) 的基变换矩阵。则

\[
[T]_{B'}^{C'} = Q^{-1} [T]_{B}^{C} P
\]

\textbf{证明}：对任意 \(v \in V\)，设 \(x = [v]_B\)，\(x' = [v]_{B'}\)，由定理 33.6.1 \(x = P x'\)。同理 \([T(v)]_C = [T]_{B}^{C} x\)，\([T(v)]_{C'} = [T]_{B'}^{C'} x'\)，且 \([T(v)]_C = Q [T(v)]_{C'}\)。

因此 \(Q [T]_{B'}^{C'} x' = [T]_{B}^{C} P x'\)。由于 \(x'\) 任意，\(Q [T]_{B'}^{C'} = [T]_{B}^{C} P\)，即 \([T]_{B'}^{C'} = Q^{-1} [T]_{B}^{C} P\)。\(\blacksquare\)

\textbf{定义 33.6.2（相似矩阵）}：设 \(A, B \in M_{n \times n}(F)\)。若存在可逆矩阵 \(P \in M_{n \times n}(F)\) 使得 \(B = P^{-1}AP\)，则称 \(A\) 与 \(B\) \textbf{相似}。

\textbf{推论 33.6.3}：同一线性算子在不同基下的矩阵是相似的。

\textbf{证明}：在定理 33.6.2 中取 \(V = W\)，\(B = C\)，\(B' = C'\)，\(P = Q\)，则 \([T]_{B'}^{B'} = P^{-1}[T]_{B}^{B} P\)。\(\blacksquare\)

\hrulefill

\hrulefill

\hrulefill

\hrulefill

\hrulefill

\hrulefill

\hrulefill

\hrulefill

\section{第103章：矩阵范数与条件数}\label{ux7b2c103ux7ae0ux77e9ux9635ux8303ux6570ux4e0eux6761ux4ef6ux6570}

矩阵范数将向量的长度概念推广到矩阵空间，为度量矩阵的''大小''和线性方程组的''敏感度''提供了基本工具。在数值线性代数中，矩阵范数是误差分析、迭代法收敛性分析和正则化方法的理论基石。条件数（condition number）\(\kappa(A) = \|A\| \cdot \|A^{-1}\|\) 度量了线性方程组 \(A\mathbf{x} = \mathbf{b}\) 的解 \(\mathbf{x}\) 对数据 \(\mathbf{b}\) 的扰动的敏感程度：相对误差以条件数倍放大。这一简单的不等式 \(\|\Delta \mathbf{x}\|/\|\mathbf{x}\| \leq \kappa(A) \cdot \|\Delta \mathbf{b}\|/\|\mathbf{b}\|\) 是数值分析中最基本的误差界，也是理解病态问题（ill-conditioned problems）的出发点。矩阵范数的理论框架建立在有限维赋范向量空间的对偶性之上：算子范数（由向量范数诱导的矩阵范数）满足次可乘性 \(\|AB\| \leq \|A\| \cdot \|B\|\)，这使得矩阵范数成为 Banach 代数 \(\mathcal{B}(\mathbb{R}^n)\) 上的范数。本章系统讨论谱范数、Frobenius 范数和算子范数的性质，以及它们与奇异值、特征值和谱半径的深刻联系。特别地，谱半径 \(\rho(A)\) 与任意矩阵范数的关系 \(\rho(A) \leq \|A\|\) 是迭代法收敛性的核心判据，而 \(\rho(A) = \lim_{k \to \infty} \|A^k\|^{1/k}\)（Gelfand 公式）则揭示了谱半径作为渐近增长率的本质。

\subsection{55.1 矩阵范数的定义与基本性质}\label{ux77e9ux9635ux8303ux6570ux7684ux5b9aux4e49ux4e0eux57faux672cux6027ux8d28}

\textbf{定义 55.1}（矩阵范数）：域 \(F = \mathbb{R}\) 或 \(\mathbb{C}\) 上 \(M_{m \times n}(F)\) 的\textbf{矩阵范数} \(\|\cdot\|\) 是一个函数 \(M_{m \times n}(F) \to \mathbb{R}_{\geq 0}\)，满足： 1. \textbf{正定性}：\(\|A\| \geq 0\)，且 \(\|A\| = 0 \iff A = 0\) 2. \textbf{齐次性}：\(\|\alpha A\| = |\alpha| \cdot \|A\|\)（\(\forall \alpha \in F\)） 3. \textbf{三角不等式}：\(\|A + B\| \leq \|A\| + \|B\|\)

若进一步对 \(A \in M_{m \times n}\)，\(B \in M_{n \times p}\) 满足\textbf{次可乘性}（相容性）\(\|AB\| \leq \|A\| \cdot \|B\|\)，则称为\textbf{矩阵代数范数}。

\textbf{定义 55.2}（算子范数 / 诱导范数）：给定 \(\mathbb{R}^n\) 上的向量范数 \(\|\cdot\|_{\alpha}\) 和 \(\mathbb{R}^m\) 上的向量范数 \(\|\cdot\|_{\beta}\)，\(A \in M_{m \times n}(\mathbb{R})\) 的\textbf{算子范数}为

\[\|A\|_{\alpha \to \beta} = \sup_{\mathbf{x} \neq \mathbf{0}} \frac{\|A\mathbf{x}\|_{\beta}}{\|\mathbf{x}\|_{\alpha}} = \sup_{\|\mathbf{x}\|_{\alpha} = 1} \|A\mathbf{x}\|_{\beta}\]

常见的算子范数： - \(\|A\|_1 = \max_{1 \leq j \leq n} \sum_{i=1}^m |a_{ij}|\)（最大列和范数，由 \(\ell_1\) 向量范数诱导） - \(\|A\|_\infty = \max_{1 \leq i \leq m} \sum_{j=1}^n |a_{ij}|\)（最大行和范数，由 \(\ell_\infty\) 向量范数诱导） - \(\|A\|_2 = \sigma_{\max}(A)\)（谱范数，由 \(\ell_2\) 向量范数诱导，等于最大奇异值）

\textbf{定理 55.1}（诱导范数自动满足次可乘性）：对任意诱导范数 \(\|\cdot\|_{\alpha \to \beta}\)，当矩阵乘积有意义时，\(\|AB\|_{\alpha \to \gamma} \leq \|A\|_{\beta \to \gamma} \cdot \|B\|_{\alpha \to \beta}\)。

\textbf{证明}：\(\|AB\|_{\alpha \to \gamma} = \sup_{\mathbf{x} \neq 0} \|AB\mathbf{x}\|_\gamma / \|\mathbf{x}\|_\alpha\)。对任意 \(\mathbf{x} \neq 0\)，若 \(B\mathbf{x} = 0\) 则 \(\|AB\mathbf{x}\|_\gamma = 0 \leq \text{右端}\)。若 \(B\mathbf{x} \neq 0\)，则 \(\|AB\mathbf{x}\|_\gamma / \|\mathbf{x}\|_\alpha = (\|A(B\mathbf{x})\|_\gamma / \|B\mathbf{x}\|_\beta) \cdot (\|B\mathbf{x}\|_\beta / \|\mathbf{x}\|_\alpha) \leq \|A\|_{\beta \to \gamma} \cdot \|B\|_{\alpha \to \beta}\)。两边取上确界即得。\(\blacksquare\)

\textbf{定义 55.3}（Frobenius 范数）：\(A \in M_{m \times n}(\mathbb{R})\) 的 \textbf{Frobenius 范数}为

\[\|A\|_F = \sqrt{\sum_{i=1}^m \sum_{j=1}^n |a_{ij}|^2} = \sqrt{\operatorname{tr}(A^T A)} = \sqrt{\sum_{i=1}^{\min\{m,n\}} \sigma_i(A)^2}\]

Frobenius 范数不是诱导范数（不存在向量范数使其成为算子范数），但满足次可乘性：\(\|AB\|_F \leq \|A\|_F \|B\|_F\) 和 \(\|AB\|_F \leq \|A\|_2 \|B\|_F\)。

\textbf{定理 55.2}（矩阵范数的等价性）：有限维空间上所有范数等价。具体地，对任意 \(A \in M_{m \times n}(\mathbb{R})\)： \[\|A\|_2 \leq \|A\|_F \leq \sqrt{\operatorname{rank}(A)} \|A\|_2\] \[\frac{1}{\sqrt{n}} \|A\|_\infty \leq \|A\|_2 \leq \sqrt{m} \|A\|_\infty\] \[\frac{1}{\sqrt{m}} \|A\|_1 \leq \|A\|_2 \leq \sqrt{n} \|A\|_1\]

\textbf{证明}：\(\|A\|_2^2 = \sigma_{\max}(A)^2 \leq \sum_i \sigma_i(A)^2 = \|A\|_F^2 \leq \operatorname{rank}(A) \cdot \sigma_{\max}(A)^2\)。对 \(\|\cdot\|_\infty\) 的界：\(\|A\|_2 = \sup_{\|\mathbf{x}\|_2=1} \|A\mathbf{x}\|_2 \leq \sup_{\|\mathbf{x}\|_2=1} \sqrt{m} \|A\mathbf{x}\|_\infty \leq \sqrt{m} \|A\|_\infty\)（因 \(\|\mathbf{y}\|_2 \leq \sqrt{m} \|\mathbf{y}\|_\infty\)）。反向及 \(\|\cdot\|_1\) 的界类似可得。\(\blacksquare\)

\subsection{55.2 条件数与误差分析}\label{ux6761ux4ef6ux6570ux4e0eux8befux5deeux5206ux6790}

\textbf{定义 55.4}（条件数）：对可逆矩阵 \(A \in \mathbb{R}^{n \times n}\)，关于范数 \(\|\cdot\|\) 的\textbf{条件数}为 \[\kappa(A) = \|A\| \cdot \|A^{-1}\|\] 按约定，若 \(A\) 不可逆，\(\kappa(A) = \infty\)。对谱范数，\(\kappa_2(A) = \sigma_{\max}(A) / \sigma_{\min}(A)\)（最大与最小奇异值之比）。

\textbf{定理 55.3}（条件数与线性方程组误差）：考虑线性方程组 \(A\mathbf{x} = \mathbf{b}\)，设 \(\mathbf{b}\) 的扰动为 \(\Delta \mathbf{b}\)，解的扰动为 \(\Delta \mathbf{x}\)，即 \(A(\mathbf{x} + \Delta \mathbf{x}) = \mathbf{b} + \Delta \mathbf{b}\)。则 \[\frac{\|\Delta \mathbf{x}\|}{\|\mathbf{x}\|} \leq \kappa(A) \cdot \frac{\|\Delta \mathbf{b}\|}{\|\mathbf{b}\|}\] 若进一步矩阵 \(A\) 也有扰动 \(\Delta A\)（\((A + \Delta A)(\mathbf{x} + \Delta \mathbf{x}) = \mathbf{b}\)），且 \(\|\Delta A\| \cdot \|A^{-1}\| < 1\)，则 \[\frac{\|\Delta \mathbf{x}\|}{\|\mathbf{x}\|} \leq \frac{\kappa(A)}{1 - \kappa(A) \cdot \|\Delta A\|/\|A\|} \left( \frac{\|\Delta \mathbf{b}\|}{\|\mathbf{b}\|} + \frac{\|\Delta A\|}{\|A\|} \right)\]

\textbf{证明}：（仅右端扰动情形）\(A(\mathbf{x}+\Delta\mathbf{x}) = \mathbf{b}+\Delta\mathbf{b}\) 减去 \(A\mathbf{x}=\mathbf{b}\) 得 \(A\Delta\mathbf{x} = \Delta\mathbf{b}\)，故 \(\Delta\mathbf{x} = A^{-1}\Delta\mathbf{b}\)。取范数 \(\|\Delta\mathbf{x}\| \leq \|A^{-1}\| \cdot \|\Delta\mathbf{b}\|\)。又 \(\|\mathbf{b}\| = \|A\mathbf{x}\| \leq \|A\| \cdot \|\mathbf{x}\|\)，故 \(1/\|\mathbf{x}\| \leq \|A\|/\|\mathbf{b}\|\)。两式相乘：\(\|\Delta\mathbf{x}\|/\|\mathbf{x}\| \leq \|A^{-1}\| \cdot \|\Delta\mathbf{b}\| \cdot \|A\| / \|\mathbf{b}\| = \kappa(A) \cdot \|\Delta\mathbf{b}\|/\|\mathbf{b}\|\)。矩阵扰动情形：\((A+\Delta A)\Delta\mathbf{x} = \Delta\mathbf{b} - \Delta A \cdot \mathbf{x}\)。当 \(\|\Delta A\| \cdot \|A^{-1}\| < 1\) 时，由 Neumann 级数 \((A+\Delta A)^{-1}\) 存在且有界，展开得所述界。\(\blacksquare\)

\textbf{定理 55.4}（条件数的基本性质）： (i) \(\kappa(A) \geq 1\)（对所有可逆 \(A\) 和算子范数）。 (ii) \(\kappa(\alpha A) = \kappa(A)\)（标量乘不变性）。 (iii) 若 \(Q\) 正交，\(\kappa_2(QA) = \kappa_2(AQ) = \kappa_2(A)\)（正交不变性）。 (iv) 对任意可逆 \(A\)，\(\kappa_2(A) = \max_{\|\mathbf{x}\|=1} \|A\mathbf{x}\| / \min_{\|\mathbf{x}\|=1} \|A\mathbf{x}\|\)。

\textbf{证明}：(i) \(1 = \|I\| = \|AA^{-1}\| \leq \|A\| \cdot \|A^{-1}\| = \kappa(A)\)。(ii) \(\|\alpha A\| \cdot \|(\alpha A)^{-1}\| = |\alpha| \|A\| \cdot |\alpha|^{-1} \|A^{-1}\| = \kappa(A)\)。(iii) \(\|QA\|_2 = \|A\|_2\)（\(\|QA\mathbf{x}\|_2^2 = \mathbf{x}^T A^T Q^T Q A \mathbf{x} = \mathbf{x}^T A^T A \mathbf{x} = \|A\mathbf{x}\|_2^2\)），同理 \(\|(QA)^{-1}\|_2 = \|A^{-1}Q^T\|_2 = \|A^{-1}\|_2\)。(iv) \(\kappa_2(A) = \sigma_{\max}/\sigma_{\min} = \max \|A\mathbf{x}\| / \min \|A\mathbf{x}\|\)（在单位球面上）。\(\blacksquare\)

\hrulefill

\hrulefill

\hrulefill

\section{第104章：数值秩与低秩逼近}\label{ux7b2c104ux7ae0ux6570ux503cux79e9ux4e0eux4f4eux79e9ux903cux8fd1}

在数据科学、机器学习和科学计算中，矩阵的低秩逼近（low-rank approximation）是最基本也是最强大的降维工具。给定矩阵 \(A \in \mathbb{R}^{m \times n}\)，寻求秩为 \(r\) 的矩阵 \(\hat{A}\) 使得 \(\|A - \hat{A}\|\) 在某种范数下最小化------Eckart-Young-Mirsky 定理给出了在谱范数和 Frobenius 范数下的显式最优解：截断奇异值分解（Truncated SVD）。这一结果的理论和实用价值不可估量：主成分分析（PCA）、潜在语义索引（LSI）、推荐系统（协同过滤的矩阵分解）和压缩感知都建立在 SVD 低秩逼近的基础之上。数值秩（numerical rank）的概念将严格的代数秩推广到浮点运算的现实世界：在机器精度 \(\varepsilon_{\text{mach}}\) 的存在下，矩阵的''有效秩''是奇异值大于某个阈值 \(\tau\)（如 \(\tau = \varepsilon_{\text{mach}} \cdot \|A\|\)）的个数，而非精确零奇异值的个数。这一概念对于理解矩阵的数值行为至关重要------一个理论上满秩的矩阵可能在数值上具有低秩结构（如在边界元法和快速多极子方法中），利用这种数值低秩性是现代快速算法的核心思想（如 H-矩阵、HSS 矩阵和 butterfly 算法）。本章建立数值秩的理论基础，讨论秩显露分解（rank-revealing factorizations），并证明 Eckart-Young-Mirsky 定理及其推广。

\subsection{56.1 Eckart-Young-Mirsky定理}\label{eckart-young-mirskyux5b9aux7406}

\textbf{定理 56.1}（Eckart-Young-Mirsky 定理，1936）：设 \(A \in \mathbb{R}^{m \times n}\) 的奇异值分解为 \(A = U \Sigma V^T\)，其中 \(\Sigma = \operatorname{diag}(\sigma_1, \ldots, \sigma_p)\)（\(p = \min\{m,n\}\)），\(\sigma_1 \geq \sigma_2 \geq \cdots \geq \sigma_p \geq 0\)。对任意 \(r < p\)，记截断 SVD 为 \[A_r = \sum_{i=1}^r \sigma_i \mathbf{u}_i \mathbf{v}_i^T = U_r \Sigma_r V_r^T\] 则对 Frobenius 范数和谱范数，\(A_r\) 是最优秩-\(r\) 逼近： \[\min_{\operatorname{rank}(B) \leq r} \|A - B\|_F = \|A - A_r\|_F = \sqrt{\sum_{i=r+1}^p \sigma_i^2}\] \[\min_{\operatorname{rank}(B) \leq r} \|A - B\|_2 = \|A - A_r\|_2 = \sigma_{r+1}\]

\textbf{证明}（Frobenius 范数情形）：对任意 \(B\) 满足 \(\operatorname{rank}(B) \leq r\)，由矩阵秩的性质 \(\dim \ker(B) \geq n - r\)。考虑由 \(A\) 的前 \(r+1\) 个右奇异向量张成的子空间 \(\mathcal{V}_{r+1} = \operatorname{span}\{\mathbf{v}_1, \ldots, \mathbf{v}_{r+1}\}\)，\(\dim \mathcal{V}_{r+1} = r+1\)。由维数公式，\(\mathcal{V}_{r+1} \cap \ker(B) \neq \{0\}\)，故存在单位向量 \(\mathbf{z} \in \mathcal{V}_{r+1} \cap \ker(B)\)。则 \(\|A - B\|_F^2 \geq \|(A - B)\mathbf{z}\|_2^2 = \|A\mathbf{z}\|_2^2\)（因 \(B\mathbf{z} = 0\)）。将 \(\mathbf{z}\) 在右奇异向量基下展开 \(\mathbf{z} = \sum_{i=1}^{r+1} \alpha_i \mathbf{v}_i\)（\(\sum \alpha_i^2 = 1\)），则 \(\|A\mathbf{z}\|_2^2 = \|\sum_{i=1}^{r+1} \sigma_i \alpha_i \mathbf{u}_i\|_2^2 = \sum_{i=1}^{r+1} \sigma_i^2 \alpha_i^2 \geq \sigma_{r+1}^2 \sum_{i=1}^{r+1} \alpha_i^2 = \sigma_{r+1}^2\)。故 \(\|A - B\|_F \geq \sigma_{r+1}\)。这说明谱范数下最优误差下界为 \(\sigma_{r+1}\)。而 \(A_r\) 恰达到此下界：\(\|A - A_r\|_2 = \sigma_{r+1}\)。

对 Frobenius 范数，令 \(B = A_r + E\)（\(E\) 任意，但 \(\operatorname{rank}(B) \leq r\)）。将 \(A - B = (A - A_r) - E\) 分解在左/右奇异向量基上，利用 \(\langle A - A_r, E \rangle = 0\)（正交性------\((A-A_r)\) 由 \(\sigma_{r+1}, \ldots, \sigma_p\) 对应的奇异向量生成，而 \(\operatorname{col}(E) \subseteq \operatorname{col}(A_r)\) 因 \(\operatorname{rank}(B) \leq r\) 限制了 \(E\) 的列空间在 \(A_r\) 的列空间和其正交补间），得 \(\|A - B\|_F^2 = \|A - A_r\|_F^2 + \|E\|_F^2 \geq \|A - A_r\|_F^2\)。完整证明需利用：对任意 \(r\) 秩矩阵 \(B\)，存在正交矩阵 \(Q\) 使 \(\|A - B\|_F = \|A - BQ\|_F\)，选择 \(Q\) 使 \(BQ\) 的列张成空间与 \(A_r\) 的列空间一致，此时正交性成立。\(\blacksquare\)

\subsection{56.2 数值秩与秩显露分解}\label{ux6570ux503cux79e9ux4e0eux79e9ux663eux9732ux5206ux89e3}

\textbf{定义 56.1}（数值 \(\varepsilon\)-秩）：对矩阵 \(A \in \mathbb{R}^{m \times n}\) 和容忍度 \(\varepsilon > 0\)，\(A\) 的\textbf{数值 \(\varepsilon\)-秩}为 \[r_\varepsilon(A) = \min\{\operatorname{rank}(B) : \|A - B\|_2 \leq \varepsilon\}\] 由 Eckart-Young-Mirsky 定理，\(r_\varepsilon(A) = \#\{i : \sigma_i(A) > \varepsilon\}\)，即大于 \(\varepsilon\) 的奇异值个数。

\textbf{定理 56.2}（数值秩的扰动稳定性）：设 \(\tilde{A} = A + \Delta A\) 为 \(A\) 的扰动。则对任意 \(\varepsilon > 0\)， \[r_{\varepsilon + \|\Delta A\|_2}(\tilde{A}) \leq r_\varepsilon(A) \leq r_{\varepsilon - \|\Delta A\|_2}(\tilde{A}) \quad (\text{当} \ \varepsilon > \|\Delta A\|_2)\] 特别地，若 \(\|\Delta A\|_2 \leq \eta\)，则 \(r_{\varepsilon + \eta}(A) \leq r_\varepsilon(\tilde{A}) \leq r_{\varepsilon - \eta}(A)\)。

\textbf{证明}：由 Weyl 不等式，\(|\sigma_i(\tilde{A}) - \sigma_i(A)| \leq \|\Delta A\|_2\)。因此若 \(\sigma_i(A) > \varepsilon\)，则 \(\sigma_i(\tilde{A}) > \varepsilon - \|\Delta A\|_2\)；反之若 \(\sigma_i(\tilde{A}) > \varepsilon + \|\Delta A\|_2\)，则 \(\sigma_i(A) > \varepsilon\)。由 \(r_\varepsilon\) 的定义即得。\(\blacksquare\)

\textbf{定义 56.2}（秩显露QR分解 / Rank-Revealing QR, RRQR）：矩阵 \(A \in \mathbb{R}^{m \times n}\) 的\textbf{秩显露QR分解}是一个QR分解 \(A\Pi = QR\)，其中 \(\Pi\) 是置换矩阵，\(Q\) 正交，\(R = \begin{pmatrix} R_{11} & R_{12} \\ 0 & R_{22} \end{pmatrix}\) 满足 \(R_{11} \in \mathbb{R}^{k \times k}\) 良态（\(\sigma_{\min}(R_{11})\) 不太小）且 \(\|R_{22}\|_2\) 很小（接近 \(\sigma_{k+1}(A)\)），使得 \(k\) 是 \(A\) 的数值秩的良好估计。RRQR 的经典算法由 Golub（1965）的列选主元QR给出，其满足

\[\sigma_{\min}(R_{11}) \geq \frac{\sigma_k(A)}{\sqrt{k(n-k) + k}}, \quad \|R_{22}\|_2 \leq \sigma_{k+1}(A) \cdot \sqrt{k(n-k) + k}\]

\textbf{定理 56.3}（强秩显露QR的存在性，Gu-Eisenstat 1996）：对任意 \(A \in \mathbb{R}^{m \times n}\) 和整数 \(k \leq \operatorname{rank}(A)\)，存在置换 \(\Pi\) 使QR分解的 \(R_{11}\) 满足 \[\sigma_{\min}(R_{11}) \geq \frac{\sigma_k(A)}{\sqrt{1 + f^2 k(n-k)}}, \quad \|R_{22}\|_2 \leq \sigma_{k+1}(A) \cdot \sqrt{1 + f^2 k(n-k)}\] 其中 \(f \geq 1\) 是预选参数（\(f=1\) 对应经典界，\(f=\sqrt{2}\) 由 Gu-Eisenstat 算法达到）。该算法复杂度为 \(O(m n k)\)。

\textbf{证明}（框架）：算法迭代地交换列以改进 \(|\det R_{11}|\)（最大化 \(R_{11}\) 的行列式的绝对值等价于最小化条件数）。每步选择一列移入 \(R_{11}\) 的前 \(k\) 列和移出一列，使得 \(|\det R_{11}|\) 的增大因子至少为 \(f\)。由于 \(|\det R_{11}| = \prod_{i=1}^k \sigma_i(R_{11})\) 且 \(\prod_{i=1}^k \sigma_i(A)\) 为理论最优上界，该过程保证 \(\sigma_{\min}(R_{11})\) 满足所述下界。算法的终止性和复杂度由行列式值的单调性和每次交换的局部计算量控制。\(\blacksquare\)

\subsection{56.3 低秩逼近的随机化算法}\label{ux4f4eux79e9ux903cux8fd1ux7684ux968fux673aux5316ux7b97ux6cd5}

\textbf{定理 56.4}（随机化SVD / Halko-Martinsson-Tropp 2011）：设 \(A \in \mathbb{R}^{m \times n}\)，目标是计算秩 \(r\) 逼近。取随机 Gaussian 矩阵 \(\Omega \in \mathbb{R}^{n \times (r+p)}\)（\(p \geq 2\) 为过采样参数），构造样本矩阵 \(Y = A \Omega\)。计算 \(Y\) 的正交基 \(Q \in \mathbb{R}^{m \times (r+p)}\)（通过QR分解）。则低秩逼近 \(\hat{A} = Q Q^T A\) 以高概率满足 \[\mathbb{E}\|A - \hat{A}\|_F \leq \left(1 + \frac{r}{p-1}\right)^{1/2} \sqrt{\sum_{i=r+1}^{\min\{m,n\}} \sigma_i(A)^2}\] 即接近最优 Frobenius 范数误差（至多常数倍），而计算仅需 \(O(m n (r+p))\) 次运算（相比完整 SVD 的 \(O(m n \min\{m,n\})\) 有巨大加速）。

\textbf{证明}：核心洞察是 \(A \Omega\) 的列空间以高概率''捕获''\(A\) 的前 \(r\) 个左奇异向量张成的子空间。确切地，\(Y = U \Sigma V^T \Omega = U \Sigma \tilde{\Omega}\)（\(\tilde{\Omega} = V^T \Omega\) 同为 Gaussian 分布，由正交不变性）。令 \(\Sigma = \operatorname{diag}(\Sigma_1, \Sigma_2)\)（\(\Sigma_1\) 为前 \(r\) 个奇异值），\(\tilde{\Omega} = \begin{pmatrix} \tilde{\Omega}_1 \\ \tilde{\Omega}_2 \end{pmatrix}\)。条件数分析：\(\tilde{\Omega}_1 \in \mathbb{R}^{r \times (r+p)}\) 以高概率满行秩（当 \(p \geq 2\)），使得 \((\Sigma_1 \tilde{\Omega}_1)\) 的列空间包含 \(\operatorname{col}(U_1)\)。误差 \(\|A - QQ^T A\|_F\) 利用 \(\Omega_2\) 的独立性和 Gaussian 矩阵的矩估计，结合 Johnson-Lindenstrauss 引理的随机矩阵版本得到期望界。详细证明见 Halko-Martinsson-Tropp (2011) 的定理 9.1 和 10.5。\(\blacksquare\)

\newpage

\chapter{05-代数/01-线性代数/03-行列式.md}\label{ux4ee3ux657001-ux7ebfux6027ux4ee3ux657003-ux884cux5217ux5f0f.md}

\section{第105章：行列式}\label{ux7b2c105ux7ae0ux884cux5217ux5f0f}

行列式是线性代数中最基本的不变量之一，它既是矩阵可逆性的判定标准，也是线性变换的''体积缩放因子''。行列式的概念起源于 17 世纪 Leibniz 和关孝和关于线性方程组求解的工作，经过 Cramer、Laplace、Cauchy、Jacobi 等人的发展，最终在 19 世纪形成了现代公理化体系。本章从置换的代数理论出发，建立行列式的公理化定义，然后系统展开其性质、展开公式和几何意义。

\subsection{34.1 行列式的公理化定义}\label{ux884cux5217ux5f0fux7684ux516cux7406ux5316ux5b9aux4e49}

\textbf{定义 34.1.1（置换与奇偶性）}

设 \(n \in \mathbb{N}^+\)。\(\{1, 2, \ldots, n\}\) 到自身的双射 \(\sigma\) 称为一个 \textbf{\(n\) 元置换}。所有 \(n\) 元置换的集合记为 \(S_n\)，它在映射的复合下构成群，称为 \(n\) 次对称群，阶为 \(n!\)。

置换的表示方法有三种：（1）\textbf{双行表示法}：将 \(\sigma\) 写为 \(\begin{pmatrix} 1 & 2 & \cdots & n \\ \sigma(1) & \sigma(2) & \cdots & \sigma(n) \end{pmatrix}\)；（2）\textbf{循环表示法}：将 \(\sigma\) 分解为不相交轮换的乘积，例如 \((1\;2\;3)(4\;5)\) 表示 \(1 \mapsto 2 \mapsto 3 \mapsto 1\) 且 \(4 \mapsto 5 \mapsto 4\)；（3）\textbf{对换分解}：每个轮换可进一步分解为对换的乘积，如 \((1\;2\;3) = (1\;3)(1\;2)\)。

一个置换 \(\sigma\) 称为\textbf{对换}，若它仅交换两个元素而保持其余不变。每个置换可分解为若干个对换的复合。若分解中对换个数的奇偶性唯一确定，则分别称 \(\sigma\) 为\textbf{偶置换}或\textbf{奇置换}。定义 \(\sigma\) 的\textbf{符号}为

\[
\operatorname{sgn}(\sigma) = \begin{cases}
+1, & \text{若 } \sigma \text{ 是偶置换} \\
-1, & \text{若 } \sigma \text{ 是奇置换}
\end{cases}
\]

符号函数 \(\operatorname{sgn}: S_n \to \{\pm 1\}\) 是群同态，其核 \(A_n = \ker(\operatorname{sgn})\) 是所有偶置换构成的 \(n\) 次交错群，\(|A_n| = n!/2\)（\(n \geq 2\)）。

\textbf{定义 34.1.1'（逆序数定义）}：置换 \(\sigma\) 的\textbf{逆序数}定义为 \(\operatorname{inv}(\sigma) = \#\{(i, j) : 1 \leq i < j \leq n, \sigma(i) > \sigma(j)\}\)。则 \(\operatorname{sgn}(\sigma) = (-1)^{\operatorname{inv}(\sigma)}\)。对换 \((i\;j)\) 的逆序数为 \(2|i-j|-1\)（奇数），故对换为奇置换。逆序数定义与对换分解定义的等价性可由''每个对换改变逆序数的奇偶性''这一事实保证。

\textbf{定理 34.1.1（奇偶性的良定义）}：一个置换不可能同时分解为偶数个和奇数个对换的复合。

\textbf{证明}：设 \(P(x_1, \ldots, x_n) = \prod_{1 \leq i < j \leq n} (x_i - x_j)\) 为 Vandermonde 多项式。对任意置换 \(\sigma\)，定义 \(\sigma P(x_1, \ldots, x_n) = P(x_{\sigma(1)}, \ldots, x_{\sigma(n)})\)。考虑 \(\sigma P\) 与 \(P\) 的关系。由于 \(P\) 的每个因子 \((x_i - x_j)\) 在 \(\sigma P\) 中变为 \((x_{\sigma(i)} - x_{\sigma(j)})\)，而 \(\sigma\) 是双射，\(\sigma P\) 的因子集合与 \(P\) 的因子集合相同，仅可能差符号。事实上，当 \(\sigma\) 为对换 \((k\; \ell)\) 时，因子 \((x_k - x_\ell)\) 变为 \((x_\ell - x_k) = -(x_k - x_\ell)\)，而其他因子成对交换，故 \(\sigma P = -P\)。因此一个对换使 \(P\) 变号，故 \(k\) 个对换的复合使 \(P\) 乘以 \((-1)^k\)。若同一置换有两种分解，其中一种有 \(k\) 个对换、另一种有 \(\ell\) 个对换，则 \((-1)^k P = \sigma P = (-1)^\ell P\)，故 \((-1)^k = (-1)^\ell\)，\(k\) 与 \(\ell\) 同奇偶。\(\blacksquare\)

\textbf{命题 34.1.1（符号的性质）}：对任意 \(\sigma, \tau \in S_n\)， (i) \(\operatorname{sgn}(\sigma\tau) = \operatorname{sgn}(\sigma)\operatorname{sgn}(\tau)\) (ii) \(\operatorname{sgn}(\sigma^{-1}) = \operatorname{sgn}(\sigma)\) (iii) 若 \(\sigma\) 是 \(k\)-轮换，则 \(\operatorname{sgn}(\sigma) = (-1)^{k-1}\)

\textbf{证明}：(i) 若 \(\sigma\) 分解为 \(a\) 个对换，\(\tau\) 分解为 \(b\) 个对换，则 \(\sigma\tau\) 分解为 \(a+b\) 个对换，故 \(\operatorname{sgn}(\sigma\tau) = (-1)^{a+b} = (-1)^a(-1)^b = \operatorname{sgn}(\sigma)\operatorname{sgn}(\tau)\)。(ii) \(\sigma\sigma^{-1} = \text{id}\)，故 \(\operatorname{sgn}(\sigma)\operatorname{sgn}(\sigma^{-1}) = 1\)，而 \(\operatorname{sgn}(\sigma) = \pm 1\)，故 \(\operatorname{sgn}(\sigma^{-1}) = \operatorname{sgn}(\sigma)\)。(iii) \(k\)-轮换 \((a_1\;a_2\;\cdots\;a_k) = (a_1\;a_k)(a_1\;a_{k-1})\cdots(a_1\;a_2)\)，共 \(k-1\) 个对换。\(\blacksquare\)

\textbf{定义 34.1.2（行列式）}：设 \(A = (a_{ij}) \in M_{n \times n}(F)\)。\(A\) 的\textbf{行列式} \(\det A\)（或 \(|A|\)）定义为

\[
\det A = \sum_{\sigma \in S_n} \operatorname{sgn}(\sigma) \, a_{1\sigma(1)} a_{2\sigma(2)} \cdots a_{n\sigma(n)}
\]

\textbf{评注 34.1.1}：上式称为行列式的 \textbf{Leibniz 公式}。\(n\) 阶行列式共有 \(n!\) 项，每项是 \(n\) 个矩阵元素的乘积（每行每列各取一个），符号由对应置换的奇偶性决定。对于 \(n=2\)，\(\det\begin{pmatrix} a & b \\ c & d \end{pmatrix} = ad - bc\)；对于 \(n=3\)，\(\det A = a_{11}a_{22}a_{33} + a_{12}a_{23}a_{31} + a_{13}a_{21}a_{32} - a_{13}a_{22}a_{31} - a_{11}a_{23}a_{32} - a_{12}a_{21}a_{33}\)（Sarrus 法则）。

\textbf{定义 34.1.3（行列式的等价公理化定义）}：函数 \(D : M_{n \times n}(F) \to F\) 称为\textbf{行列式函数}，若它满足以下三条公理：

\begin{enumerate}
\def\labelenumi{\arabic{enumi}.}
\tightlist
\item
  \textbf{多重线性}：\(D\) 对每一行是线性的（固定其他行，该行向量上为线性函数）
\item
  \textbf{交错性}：若两行相同，则 \(D(A) = 0\)
\item
  \textbf{规范性}：\(D(I_n) = 1\)（其中 \(I_n\) 是 \(n\) 阶单位矩阵）
\end{enumerate}

公理化定义的优势在于：它完全刻画了行列式的本质，不依赖于具体的求和公式，且便于推广到任意交换环上的模。由交错性可推出：交换两行，行列式变号（因为 \(0 = D(\ldots, v_i+v_j, \ldots, v_i+v_j, \ldots) = D(\ldots, v_i, \ldots, v_i, \ldots) + D(\ldots, v_i, \ldots, v_j, \ldots) + D(\ldots, v_j, \ldots, v_i, \ldots) + D(\ldots, v_j, \ldots, v_j, \ldots) = 0 + D(\ldots, v_i, \ldots, v_j, \ldots) + D(\ldots, v_j, \ldots, v_i, \ldots) + 0\)，故 \(D(\ldots, v_j, \ldots, v_i, \ldots) = -D(\ldots, v_i, \ldots, v_j, \ldots)\)）。

\textbf{定理 34.1.2（行列式公理的等价性）}：Leibniz 公式定义的行列式满足公理化定义的三条性质，且任何满足这三条公理的函数必然等于 Leibniz 公式。

\textbf{证明}：先证 Leibniz 公式满足三条公理。多重线性：每项 \(a_{1\sigma(1)} \cdots a_{n\sigma(n)}\) 对每行是线性的，和亦然。交错性：若第 \(i\) 行与第 \(j\) 行相同（\(i < j\)），设 \(\tau\) 为交换 \(i\) 和 \(j\) 的对换。将置换按 \(\{\sigma, \sigma\tau\}\) 配对，\(\operatorname{sgn}(\sigma\tau) = -\operatorname{sgn}(\sigma)\)，而 \(a_{1\sigma(1)} \cdots a_{n\sigma(n)} = a_{1\sigma\tau(1)} \cdots a_{n\sigma\tau(n)}\)（因为第 \(i\) 行和第 \(j\) 行相同），两项抵消。规范性：\(I_n\) 中仅 \(\sigma\) 为恒等置换时乘积非零，此时 \(\operatorname{sgn}(\text{id}) = 1\)，\(1 \cdot 1 \cdots 1 = 1\)。

再证唯一性：设 \(D\) 满足三条公理。将 \(A\) 的第 \(j\) 行写为 \(a_j = \sum_{i=1}^{n} a_{ij} e_i\)（其中 \(e_i\) 是标准基）。由多重线性， \[
D(A) = D\left(\sum_{i_1} a_{i_1 1} e_{i_1}, \ldots, \sum_{i_n} a_{i_n n} e_{i_n}\right) = \sum_{i_1, \ldots, i_n} a_{i_1 1} \cdots a_{i_n n} D(e_{i_1}, \ldots, e_{i_n})
\] 由交错性，若 \(i_1, \ldots, i_n\) 中有重复，\(D(e_{i_1}, \ldots, e_{i_n}) = 0\)。故只需对 \((i_1, \ldots, i_n)\) 是 \(\{1, \ldots, n\}\) 的置换求和。设 \(i_k = \sigma(k)\)，则 \(D(e_{\sigma(1)}, \ldots, e_{\sigma(n)}) = \operatorname{sgn}(\sigma) D(e_1, \ldots, e_n) = \operatorname{sgn}(\sigma) \cdot 1\)。因此 \(D(A) = \sum_{\sigma \in S_n} \operatorname{sgn}(\sigma) a_{\sigma(1)1} \cdots a_{\sigma(n)n} = \sum_{\sigma \in S_n} \operatorname{sgn}(\sigma) a_{1\sigma^{-1}(1)} \cdots a_{n\sigma^{-1}(n)}\)。当 \(\sigma\) 遍历 \(S_n\) 时 \(\sigma^{-1}\) 也遍历 \(S_n\)，且 \(\operatorname{sgn}(\sigma^{-1}) = \operatorname{sgn}(\sigma)\)，故等于 Leibniz 公式。\(\blacksquare\)

\textbf{推论 34.1.3（行列式的列公理化）}：行列式也可由列的类似公理唯一确定：\(D\) 对每列是多重线性的，两列相同则 \(D=0\)，且 \(D(I_n)=1\)。由转置不变性（见定理 34.2.1(i)），行公理与列公理等价。

\hrulefill

\subsection{34.2 行列式的基本性质}\label{ux884cux5217ux5f0fux7684ux57faux672cux6027ux8d28}

\textbf{定理 34.2.1（行列式的基本性质）}：设 \(A, B \in M_{n \times n}(F)\)。

\begin{enumerate}
\def\labelenumi{(\roman{enumi})}
\tightlist
\item
  \textbf{转置不变性}：\(\det(A^T) = \det A\)
\item
  \textbf{行交换}：交换两行，行列式变号
\item
  \textbf{行倍乘}：某行乘以 \(\lambda\)，行列式乘以 \(\lambda\)
\item
  \textbf{行加法}：将一行的倍数加到另一行，行列式不变
\item
  \textbf{乘法性质}：\(\det(AB) = \det A \cdot \det B\)
\item
  \textbf{可逆性判定}：\(A\) 可逆 \(\iff \det A \neq 0\)
\end{enumerate}

\textbf{证明}：

\begin{enumerate}
\def\labelenumi{(\roman{enumi})}
\tightlist
\item
  \(\det(A^T) = \sum_{\sigma} \operatorname{sgn}(\sigma) a_{\sigma(1)1} \cdots a_{\sigma(n)n}\)。令 \(\tau = \sigma^{-1}\)，则 \(a_{\sigma(k)k} = a_{k \tau(k)}\)（交换因子顺序），且 \(\operatorname{sgn}(\tau) = \operatorname{sgn}(\sigma)\)。因此 \(\det(A^T) = \sum_{\tau} \operatorname{sgn}(\tau) a_{1\tau(1)} \cdots a_{n\tau(n)} = \det A\)。
\end{enumerate}

(ii)-(iv) 由公理化定义中的多重线性、交错性和规范性直接得出。具体地，交换第 \(i\) 行和第 \(j\) 行：\(0 = D(\ldots, v_i+v_j, \ldots, v_i+v_j, \ldots) = D(\ldots, v_i, \ldots, v_i, \ldots) + D(\ldots, v_i, \ldots, v_j, \ldots) + D(\ldots, v_j, \ldots, v_i, \ldots) + D(\ldots, v_j, \ldots, v_j, \ldots)\)，第一项和第四项由交错性为零，故 \(D(\ldots, v_j, \ldots, v_i, \ldots) = -D(\ldots, v_i, \ldots, v_j, \ldots)\)。

\begin{enumerate}
\def\labelenumi{(\alph{enumi})}
\setcounter{enumi}{21}
\tightlist
\item
  将 \(A\) 固定，考虑函数 \(f(B) = \det(AB)\)。\(f\) 对 \(B\) 的行也是多重线性和交错的（因为 \((AB)\) 的第 \(i\) 行是 \(A\) 的第 \(i\) 行与 \(B\) 的列的线性组合，而 \(A\) 的线性变换保持多重线性和交错性）。由唯一性定理 34.1.2，\(f(B) = f(I_n) \det B = \det A \cdot \det B\)。
\end{enumerate}

\begin{enumerate}
\def\labelenumi{(\roman{enumi})}
\setcounter{enumi}{5}
\tightlist
\item
  若 \(A\) 可逆，则 \(I_n = A A^{-1}\)，故 \(1 = \det I_n = \det A \cdot \det(A^{-1})\)，因此 \(\det A \neq 0\)。
\end{enumerate}

若 \(A\) 不可逆，则 \(\operatorname{rank} A < n\)，行向量线性相关。不妨设第 \(k\) 行是其余行的线性组合：\(a_k = \sum_{i \neq k} c_i a_i\)。由多重线性，\(\det A = \sum_{i \neq k} c_i \det(A^{(i)})\)，其中 \(A^{(i)}\) 是将第 \(k\) 行替换为第 \(i\) 行的矩阵。每个 \(A^{(i)}\) 有两行相同（第 \(i\) 行和第 \(k\) 行），由交错性 \(\det(A^{(i)}) = 0\)，故 \(\det A = 0\)。\(\blacksquare\)

\textbf{推论 34.2.2}：\(\det(A^{-1}) = (\det A)^{-1}\)。

\textbf{证明}：由 \(1 = \det(I_n) = \det(A A^{-1}) = \det A \cdot \det(A^{-1})\)。\(\blacksquare\)

\textbf{推论 34.2.3}：相似矩阵有相同的行列式。

\textbf{证明}：若 \(B = P^{-1}AP\)，则 \(\det B = \det(P^{-1}) \det A \det P = (\det P)^{-1} \det A \det P = \det A\)。\(\blacksquare\)

\textbf{推论 34.2.4（行列式与特征值）}：设 \(A \in M_{n \times n}(\mathbb{C})\) 的特征值为 \(\lambda_1, \ldots, \lambda_n\)（计重数）。则 \(\det A = \prod_{i=1}^{n} \lambda_i\)，且 \(\operatorname{tr}(A) = \sum_{i=1}^{n} \lambda_i\)。

\textbf{证明}：特征多项式 \(\chi_A(t) = \det(tI - A) = t^n - (\operatorname{tr} A) t^{n-1} + \cdots + (-1)^n \det A\)。另一方面，\(\chi_A(t) = \prod_{i=1}^{n} (t - \lambda_i) = t^n - (\sum \lambda_i) t^{n-1} + \cdots + (-1)^n \prod \lambda_i\)。比较常数项和 \(t^{n-1}\) 系数即得。\(\blacksquare\)

\textbf{定理 34.2.5（Vandermonde 行列式）}：设 \(x_1, \ldots, x_n \in F\)。则

\[
V(x_1, \ldots, x_n) = \det\begin{pmatrix}
1 & x_1 & x_1^2 & \cdots & x_1^{n-1} \\
1 & x_2 & x_2^2 & \cdots & x_2^{n-1} \\
\vdots & \vdots & \vdots & \ddots & \vdots \\
1 & x_n & x_n^2 & \cdots & x_n^{n-1}
\end{pmatrix} = \prod_{1 \leq i < j \leq n} (x_j - x_i)
\]

\textbf{证明}：对 \(n\) 归纳。\(n=2\) 时 \(V(x_1, x_2) = x_2 - x_1\)。假设 \(n-1\) 时成立。将第 \(n\) 行减去第 \(n-1\) 行的 \(x_1\) 倍，第 \(n-1\) 行减去第 \(n-2\) 行的 \(x_1\) 倍，以此类推（从下往上消元）。然后按第一列展开，提取公因子 \((x_j - x_1)\)（\(j = 2, \ldots, n\)），归约为 \(n-1\) 阶 Vandermonde 行列式。具体地，\(V(x_1, \ldots, x_n) = \prod_{j=2}^{n} (x_j - x_1) \cdot V(x_2, \ldots, x_n) = \prod_{1 \leq i < j \leq n} (x_j - x_i)\)。\(\blacksquare\)

Vandermonde 行列式在插值理论、线性递归序列和代数数论中频繁出现。特别地，\(V(x_1, \ldots, x_n) \neq 0\) 当且仅当 \(x_1, \ldots, x_n\) 两两不同，这对应于 Lagrange 插值多项式的唯一性条件。

\textbf{定理 34.2.6（Cauchy-Binet 公式）}：设 \(A \in M_{m \times n}(F)\)，\(B \in M_{n \times p}(F)\)，\(m = p\)。则

\[
\det(AB) = \sum_{1 \leq i_1 < \cdots < i_m \leq n} \det(A_{[m], \{i_1, \ldots, i_m\}}) \cdot \det(B_{\{i_1, \ldots, i_m\}, [m]})
\]

其中 \(A_{[m], S}\) 表示 \(A\) 取全部 \(m\) 行和列指标集 \(S\) 的子矩阵。特别地，当 \(m > n\) 时 \(\det(AB) = 0\)（因为 \(AB\) 的秩至多为 \(n < m\)）。

\textbf{证明}：由行列式的多重线性和 Leibniz 公式，将 \(\det(AB)\) 展开为 \(A\) 的列和 \(B\) 的行的组合。\(\det(AB) = \sum_{\sigma \in S_m} \operatorname{sgn}(\sigma) \prod_{k=1}^{m} (AB)_{k,\sigma(k)} = \sum_{\sigma} \operatorname{sgn}(\sigma) \prod_{k=1}^{m} \sum_{j_k=1}^{n} a_{k j_k} b_{j_k \sigma(k)}\)。交换求和顺序，将 \(j_1, \ldots, j_m\) 分组：若其中有重复，则对 \(\sigma\) 求和中因交错性消去；若互不相同，则按递增排列 \(i_1 < \cdots < i_m\) 并考虑置换 \(\tau\) 使 \(j_k = i_{\tau(k)}\)。最终得到 \(\sum_{1 \leq i_1 < \cdots < i_m \leq n} \det(A_{[m], \{i_1, \ldots, i_m\}}) \det(B_{\{i_1, \ldots, i_m\}, [m]})\)。\(\blacksquare\)

Cauchy-Binet 公式是行列式乘法性质 \(\det(AB) = \det A \det B\)（当 \(m=n\) 时）的推广。当 \(m=n\) 时，求和仅有一项，即 \(A\) 和 \(B\) 各自的行列式之积。

\textbf{定理 34.2.7（分块矩阵的行列式）}：

\begin{enumerate}
\def\labelenumi{(\roman{enumi})}
\item
  若 \(A \in M_{m \times m}(F)\) 可逆，则 \(\det\begin{pmatrix} A & B \\ C & D \end{pmatrix} = \det A \cdot \det(D - CA^{-1}B)\)。
\item
  若 \(A\) 和 \(C\) 可交换（即 \(AC = CA\)），则 \(\det\begin{pmatrix} A & B \\ C & D \end{pmatrix} = \det(AD - CB)\)。（注意这仅当 \(A\) 和 \(C\) 可交换时成立，一般情形下 \(\det\begin{pmatrix} A & B \\ C & D \end{pmatrix} = \det(AD - CB)\) 并不成立。）
\end{enumerate}

\textbf{证明}：(i) 利用分块矩阵的初等变换： \[
\begin{pmatrix} I & 0 \\ -CA^{-1} & I \end{pmatrix} \begin{pmatrix} A & B \\ C & D \end{pmatrix} = \begin{pmatrix} A & B \\ 0 & D - CA^{-1}B \end{pmatrix}
\] 左边第一个矩阵的行列式为 \(1\)（上三角分块，对角块为 \(I\)），右边为分块上三角矩阵，行列式为 \(\det A \cdot \det(D - CA^{-1}B)\)。由乘法性质即得。(ii) 若 \(AC = CA\)，则可利用类似技巧。\(\blacksquare\)

\hrulefill

\subsection{34.3 行列式的展开}\label{ux884cux5217ux5f0fux7684ux5c55ux5f00}

\textbf{定义 34.3.1（余子式与代数余子式）}：设 \(A = (a_{ij}) \in M_{n \times n}(F)\)。

\begin{enumerate}
\def\labelenumi{\arabic{enumi}.}
\tightlist
\item
  \(A\) 中划去第 \(i\) 行和第 \(j\) 列后得到的 \((n-1) \times (n-1)\) 矩阵的行列式称为 \(a_{ij}\) 的\textbf{余子式}，记为 \(M_{ij}\)。
\item
  \(C_{ij} = (-1)^{i+j} M_{ij}\) 称为 \(a_{ij}\) 的\textbf{代数余子式}。
\end{enumerate}

代数余子式中的符号因子 \((-1)^{i+j}\) 反映了行列式展开中''棋盘模式''的符号分布。这一符号模式来源于将第 \(i\) 行和第 \(j\) 列交换到首行首列所需的交换次数。

\textbf{定理 34.3.1（Laplace 展开）}：设 \(A \in M_{n \times n}(F)\)。

\begin{enumerate}
\def\labelenumi{(\roman{enumi})}
\tightlist
\item
  \textbf{按第 \(i\) 行展开}：\(\displaystyle \det A = \sum_{j=1}^{n} a_{ij} C_{ij}\)
\item
  \textbf{按第 \(j\) 列展开}：\(\displaystyle \det A = \sum_{i=1}^{n} a_{ij} C_{ij}\)
\item
  \textbf{错行展开为零}：\(\displaystyle \sum_{j=1}^{n} a_{ij} C_{kj} = 0 \quad (i \neq k)\)
\item
  \textbf{错列展开为零}：\(\displaystyle \sum_{i=1}^{n} a_{ij} C_{ik} = 0 \quad (j \neq k)\)
\end{enumerate}

\textbf{证明}：(i) 固定第 \(i\) 行。由多重线性，\(\det A\) 是这个行向量的线性函数。将 \(A\) 的第 \(i\) 行写为 \(\sum_{j=1}^{n} a_{ij} e_j\)（\(e_j\) 为标准基向量）。则 \(\det A = \sum_j a_{ij} \det A_{(j)}\)，其中 \(A_{(j)}\) 是将 \(A\) 的第 \(i\) 行替换为 \(e_j\) 的矩阵。通过行交换将 \(A_{(j)}\) 的第 \(i\) 行移到第 \(1\) 行，列交换将第 \(j\) 列移到第 \(1\) 列，需要 \((i-1) + (j-1) = i+j-2\) 次交换，每次变号。故 \(\det A_{(j)} = (-1)^{i+j-2} M_{ij} = (-1)^{i+j} M_{ij} = C_{ij}\)。

\begin{enumerate}
\def\labelenumi{(\roman{enumi})}
\setcounter{enumi}{2}
\tightlist
\item
  考虑将 \(A\) 的第 \(k\) 行替换为第 \(i\) 行得到的矩阵 \(A'\)（\(i \neq k\)）。\(A'\) 有两行相同，故 \(\det A' = 0\)。另一方面，按第 \(k\) 行展开 \(A'\)，第 \(k\) 行的元素为 \(a_{ij}\)（\(j = 1, \ldots, n\)），余子式为 \(C_{kj}\)（因为第 \(k\) 行的余子式与第 \(k\) 行元素无关）。故 \(0 = \det A' = \sum_{j} a_{ij} C_{kj}\)。\(\blacksquare\)
\end{enumerate}

\textbf{定义 34.3.2（伴随矩阵）}：设 \(A \in M_{n \times n}(F)\)。\(A\) 的\textbf{伴随矩阵} \(\operatorname{adj}(A)\) 定义为 \((\operatorname{adj}(A))_{ij} = C_{ji}\)（即代数余子式矩阵的转置）。

\textbf{定理 34.3.2（伴随矩阵的性质）}：\(A \cdot \operatorname{adj}(A) = \operatorname{adj}(A) \cdot A = (\det A) I_n\)。

\textbf{证明}：\((A \cdot \operatorname{adj}(A))_{ik} = \sum_{j} a_{ij} (\operatorname{adj}(A))_{jk} = \sum_{j} a_{ij} C_{kj}\)。当 \(i = k\) 时，由 Laplace 展开，\(\sum_j a_{ij} C_{ij} = \det A\)。当 \(i \neq k\) 时，由定理 34.3.1(iii)，\(\sum_j a_{ij} C_{kj} = 0\)。故 \(A \cdot \operatorname{adj}(A) = (\det A) I_n\)。同理可证 \(\operatorname{adj}(A) \cdot A = (\det A) I_n\)。\(\blacksquare\)

\textbf{推论 34.3.3（逆矩阵公式）}：若 \(\det A \neq 0\)，则 \(A^{-1} = \frac{1}{\det A} \operatorname{adj}(A)\)。

\textbf{证明}：由定理 34.3.2 直接得到。\(\blacksquare\)

\textbf{命题 34.3.4（伴随矩阵的性质）}：设 \(A \in M_{n \times n}(F)\)。 (i) \(\operatorname{adj}(A^T) = (\operatorname{adj} A)^T\) (ii) \(\operatorname{adj}(AB) = \operatorname{adj}(B) \operatorname{adj}(A)\) (iii) \(\det(\operatorname{adj} A) = (\det A)^{n-1}\) (iv) \(\operatorname{adj}(\operatorname{adj} A) = (\det A)^{n-2} A\)（\(n \geq 3\)）

\textbf{证明}：(i) 由余子式的定义和转置的性质直接得到。(ii) 若 \(A, B\) 均可逆，则 \(\operatorname{adj}(AB) = \det(AB)(AB)^{-1} = \det A \det B \cdot B^{-1} A^{-1} = (\det B \cdot B^{-1})(\det A \cdot A^{-1}) = \operatorname{adj}(B) \operatorname{adj}(A)\)。一般情形由多项式恒等性（将 \(A, B\) 替换为 \(A + tI, B + tI\)，两边作为 \(t\) 的多项式在无穷多 \(t\) 处相等，故恒等）。(iii) 由 \(A \cdot \operatorname{adj} A = (\det A) I_n\)，两边取行列式：\(\det A \cdot \det(\operatorname{adj} A) = (\det A)^n\)。若 \(\det A \neq 0\)，约去 \(\det A\) 得 \(\det(\operatorname{adj} A) = (\det A)^{n-1}\)。若 \(\det A = 0\)，由连续性（或多项式恒等性）结论仍成立。(iv) 由 (iii) 和伴随矩阵的定义验证。\(\blacksquare\)

\textbf{定理 34.3.5（一般 Laplace 展开）}：设 \(A \in M_{n \times n}(F)\)。选定 \(k\) 行（指标为 \(1 \leq i_1 < \cdots < i_k \leq n\)），则

\[
\det A = \sum_{1 \leq j_1 < \cdots < j_k \leq n} \varepsilon(i_1, \ldots, i_k; j_1, \ldots, j_k) \cdot \det A(i_1, \ldots, i_k; j_1, \ldots, j_k) \cdot \det A'(i_1, \ldots, i_k; j_1, \ldots, j_k)
\]

其中 \(\varepsilon = (-1)^{i_1 + \cdots + i_k + j_1 + \cdots + j_k}\)，\(A(i_1, \ldots, i_k; j_1, \ldots, j_k)\) 是取第 \(i_1, \ldots, i_k\) 行和第 \(j_1, \ldots, j_k\) 列的子矩阵，\(A'\) 是划去这些行和列后的余子矩阵。

\textbf{证明}：将 Leibniz 公式中的置换按所选 \(k\) 行的像集分组。固定 \(k\) 行的像集 \(\{j_1, \ldots, j_k\}\)，对应置换 \(\sigma\) 分为两部分：一部分将 \(\{i_1, \ldots, i_k\}\) 映到 \(\{j_1, \ldots, j_k\}\)，另一部分将其余指标映到余集。符号恰好分解为 \(\varepsilon\) 乘两个子行列式的符号。对 \(j_1 < \cdots < j_k\) 求和即得。\(\blacksquare\)

\hrulefill

\subsection{34.4 Cramer 法则}\label{cramer-ux6cd5ux5219}

\textbf{定理 34.4.1（Cramer 法则）}：设 \(A \in M_{n \times n}(F)\) 可逆，\(b \in F^n\)。则线性方程组 \(Ax = b\) 的唯一解为

\[
x_j = \frac{\det A_j}{\det A}, \quad j = 1, 2, \ldots, n
\]

其中 \(A_j\) 是将 \(A\) 的第 \(j\) 列替换为 \(b\) 得到的矩阵。

\textbf{证明}：由 \(A^{-1} = \frac{1}{\det A} \operatorname{adj}(A)\)，\(x = A^{-1}b\)，故 \[
x_j = \frac{1}{\det A} \sum_{i=1}^{n} (\operatorname{adj}(A))_{ji} b_i = \frac{1}{\det A} \sum_{i=1}^{n} C_{ij} b_i
\] 而 \(\sum_{i=1}^{n} C_{ij} b_i\) 正是将 \(A\) 的第 \(j\) 列替换为 \(b\) 后按第 \(j\) 列展开的结果，即 \(\det A_j\)。\(\blacksquare\)

\textbf{评注 34.4.1}：Cramer 法则在理论上非常重要，但计算上效率极低（需要计算 \(n+1\) 个 \(n\) 阶行列式，每个需 \(O(n!)\) 次运算）。在实际计算中，Gauss 消元法（\(O(n^3)\)）更为高效。然而，Cramer 法则在理论分析中------特别是当 \(A\) 和 \(b\) 依赖于参数时------提供了封闭形式的表达式，这在证明解对参数的光滑依赖性时非常有用。

\textbf{推论 34.4.2（Cramer 法则的齐次版本）}：齐次线性方程组 \(Ax = 0\) 有非零解当且仅当 \(\det A = 0\)。

\textbf{证明}：若 \(\det A \neq 0\)，由 Cramer 法则，唯一解为 \(x = 0\)。若 \(\det A = 0\)，则 \(A\) 不可逆，\(\ker A \neq \{0\}\)。\(\blacksquare\)

\hrulefill

\subsection{34.5 行列式与体积}\label{ux884cux5217ux5f0fux4e0eux4f53ux79ef}

\textbf{定理 34.5.1（平行多面体的体积）}：设 \(v_1, \ldots, v_n \in \mathbb{R}^n\) 为列向量。以原点为一个顶点，以 \(v_1, \ldots, v_n\) 为棱的 \(n\) 维平行多面体的 \(n\) 维体积为

\[
\operatorname{Vol}(v_1, \ldots, v_n) = \left|\det(v_1, \ldots, v_n)\right|
\]

\textbf{证明}：设 \(f(v_1, \ldots, v_n) = \operatorname{Vol}(v_1, \ldots, v_n)\)。体积函数满足：

\begin{enumerate}
\def\labelenumi{\arabic{enumi}.}
\tightlist
\item
  对每个 \(v_i\) 是线性的（缩放棱则体积缩放到相同倍数）；具体来说，若将 \(v_i\) 替换为 \(\lambda v_i\)，则平行多面体在该方向上的尺度缩放为 \(|\lambda|\) 倍，体积缩放为 \(|\lambda|\) 倍。
\item
  若 \(v_i = v_j\)（\(i \neq j\)），则体积为 \(0\)（退化，因为平行多面体落在 \(n-1\) 维子空间中）。
\item
  对标准基，体积为 \(1\)（\(n\) 维单位立方体 \([0,1]^n\)）。
\end{enumerate}

这正是行列式公理化定义的三个条件（在绝对值意义下）。由于 \(|\det|\) 也满足这三条，由唯一性，\(f = |\det|\)。更精确地说，符号函数 \(D(v_1, \ldots, v_n) = \det(v_1, \ldots, v_n)\) 满足：\(D\) 对每列是线性的，\(D\) 是交错的，\(D(e_1, \ldots, e_n) = 1\)。\(D\) 的绝对值给出了无符号的体积，而 \(D\) 的符号给出了定向（右手系 vs 左手系）。\(\blacksquare\)

\textbf{推论 34.5.2（线性变换的体积变化）}：设 \(T : \mathbb{R}^n \to \mathbb{R}^n\) 为线性变换，\(A = [T]_{\text{std}}\)。则对任意可测集 \(\Omega \subseteq \mathbb{R}^n\)，

\[
\operatorname{Vol}(T(\Omega)) = |\det A| \cdot \operatorname{Vol}(\Omega)
\]

\textbf{证明}：由积分的变量替换公式，变换的 Jacobi 行列式为 \(|\det A|\)。对于线性变换，Jacobi 矩阵恒为 \(A\)，故 Jacobi 行列式为 \(|\det A|\)。\(\blacksquare\)

\textbf{定义 34.5.1（Gram 行列式）}：设 \(v_1, \ldots, v_k \in \mathbb{R}^n\)（\(k \leq n\)）。\(k\) 维平行多面体的 \(k\) 维体积的平方为

\[
\operatorname{Vol}_k(v_1, \ldots, v_k)^2 = \det(G(v_1, \ldots, v_k)) = \det\begin{pmatrix}
\langle v_1, v_1 \rangle & \cdots & \langle v_1, v_k \rangle \\
\vdots & \ddots & \vdots \\
\langle v_k, v_1 \rangle & \cdots & \langle v_k, v_k \rangle
\end{pmatrix}
\]

其中 \(G(v_1, \ldots, v_k)\) 称为 \textbf{Gram 矩阵}，其行列式称为 \textbf{Gram 行列式}。

\textbf{证明}：设 \(V\) 是以 \(v_1, \ldots, v_k\) 为列的 \(n \times k\) 矩阵。则 \(G = V^T V\)。由 Cauchy-Binet 公式，\(\det(V^T V) = \sum_{1 \leq i_1 < \cdots < i_k \leq n} (\det V_{\{i_1, \ldots, i_k\}, [k]})^2\)。而 \(\det V_{\{i_1, \ldots, i_k\}, [k]}\) 是将 \(v_1, \ldots, v_k\) 投影到坐标平面 \(\mathbb{R}^k_{\{i_1, \ldots, i_k\}}\) 上的有符号 \(k\) 维体积。这些投影体积的平方和恰好等于 \(k\) 维体积的平方（高维勾股定理）。\(\blacksquare\)

\textbf{推论 34.5.3（Cauchy-Schwarz 不等式的几何证明）}：对 \(k=2\)，Gram 行列式非负给出 \(\|v_1\|^2\|v_2\|^2 - \langle v_1, v_2 \rangle^2 \geq 0\)，即 \(|\langle v_1, v_2 \rangle| \leq \|v_1\|\|v_2\|\)。

\hrulefill

\subsection{34.6 特殊行列式}\label{ux7279ux6b8aux884cux5217ux5f0f}

\textbf{定理 34.6.1（循环行列式 / Circulant）}：设 \(\omega = e^{2\pi i/n}\) 为 \(n\) 次本原单位根。则

\[
\det\begin{pmatrix}
a_0 & a_1 & a_2 & \cdots & a_{n-1} \\
a_{n-1} & a_0 & a_1 & \cdots & a_{n-2} \\
a_{n-2} & a_{n-1} & a_0 & \cdots & a_{n-3} \\
\vdots & \vdots & \vdots & \ddots & \vdots \\
a_1 & a_2 & a_3 & \cdots & a_0
\end{pmatrix} = \prod_{j=0}^{n-1} \left(a_0 + a_1 \omega^j + a_2 \omega^{2j} + \cdots + a_{n-1} \omega^{(n-1)j}\right)
\]

\textbf{证明}：循环矩阵 \(C\) 的特征向量为 \((1, \omega^j, \omega^{2j}, \ldots, \omega^{(n-1)j})^T\)（\(j = 0, 1, \ldots, n-1\)），对应的特征值为 \(\lambda_j = \sum_{k=0}^{n-1} a_k \omega^{kj}\)。行列式等于所有特征值之积。\(\blacksquare\)

\textbf{定理 34.6.2（Cauchy 行列式）}：设 \(x_1, \ldots, x_n, y_1, \ldots, y_n \in F\) 且 \(x_i + y_j \neq 0\)。则

\[
\det\left(\frac{1}{x_i + y_j}\right)_{1 \leq i,j \leq n} = \frac{\prod_{i < j} (x_i - x_j)(y_i - y_j)}{\prod_{i,j=1}^{n} (x_i + y_j)}
\]

\textbf{证明}：将第 \(i\) 行乘以 \(\prod_{k=1}^{n} (x_i + y_k)\)，提取公因子后化为 Vandermonde 行列式。具体地，记 \(D = \det(1/(x_i + y_j))\)。计算 \(\det\left(\frac{\prod_{k \neq j} (x_i + y_k)}{\prod_{k=1}^{n} (x_i + y_k)}\right)\)，分子为 \(x_i\) 的 \(n-1\) 次多项式，利用 Vandermonde 行列式完成计算。\(\blacksquare\)

\textbf{定理 34.6.3（三对角行列式 / 连分式）}：设 \(a_i, b_i, c_i \in F\)。\(n\) 阶三对角行列式 \(D_n = \det(\text{tridiag}(a_i, b_i, c_i))\) 满足递推关系：

\[
D_n = b_n D_{n-1} - a_n c_{n-1} D_{n-2}
\]

其中 \(D_0 = 1\)，\(D_{-1} = 0\)。这一递推关系与正交多项式理论和连分数理论密切相关。

\textbf{证明}：按最后一行展开 \(D_n\)：\(D_n = b_n D_{n-1} + c_{n-1} \cdot (-1)^{n+(n-1)} \det(\text{去掉最后一行和倒数第二列})\)。对去掉的矩阵再按最后一列展开，得到 \(a_n D_{n-2}\)，符号为 \((-1)^{(n-1)+n} \cdot (-1)^{n+(n-1)} = -1\)。故 \(D_n = b_n D_{n-1} - a_n c_{n-1} D_{n-2}\)。\(\blacksquare\)

\hrulefill

\subsection{34.7 Pfaffian 与反对称矩阵}\label{pfaffian-ux4e0eux53cdux5bf9ux79f0ux77e9ux9635}

\textbf{定义 34.7.1（Pfaffian）}：设 \(A = (a_{ij})\) 是 \(2n \times 2n\) 的反对称矩阵（\(A^T = -A\)）。\(A\) 的 \textbf{Pfaffian} 定义为

\[
\operatorname{Pf}(A) = \frac{1}{2^n n!} \sum_{\sigma \in S_{2n}} \operatorname{sgn}(\sigma) \prod_{i=1}^{n} a_{\sigma(2i-1), \sigma(2i)}
\]

\textbf{定理 34.7.1（Pfaffian 与行列式）}：\(\det A = (\operatorname{Pf} A)^2\)。

\textbf{证明}：对 \(2n \times 2n\) 反对称矩阵 \(A\)，存在正交矩阵 \(Q\) 使得 \(Q^T A Q = \bigoplus_{i=1}^{n} \begin{pmatrix} 0 & \lambda_i \\ -\lambda_i & 0 \end{pmatrix}\)（反对称矩阵的实标准型）。则 \(\det A = \prod_{i=1}^{n} \lambda_i^2\)，而 \(\operatorname{Pf}(A) = \prod_{i=1}^{n} \lambda_i\)（由 Pfaffian 的定义和 \(Q\) 的正交性）。故 \(\det A = (\operatorname{Pf} A)^2\)。\(\blacksquare\)

\textbf{推论 34.7.2}：奇数阶反对称矩阵的行列式恒为零。偶数阶反对称矩阵的行列式是非负完全平方（在实数域上）。

Pfaffian 在微分几何（Chern-Gauss-Bonnet 定理的局部表示）、可积系统（Toda 格）和组合数学（完美匹配的计数）中有重要应用。

\hrulefill

\hrulefill

\hrulefill

\hrulefill

\hrulefill

\newpage

\chapter{05-代数/01-线性代数/04-特征值与谱.md}\label{ux4ee3ux657001-ux7ebfux6027ux4ee3ux657004-ux7279ux5f81ux503cux4e0eux8c31.md}

\section{第106章：特征值与特征向量}\label{ux7b2c106ux7ae0ux7279ux5f81ux503cux4e0eux7279ux5f81ux5411ux91cf}

特征值与特征向量是线性代数中最深刻的概念之一，它们揭示了线性算子在其作用方向上表现为标量乘法的简单事实，却蕴含着极为丰富的结构信息。从历史上看，特征值问题起源于 18 世纪对微分方程和力学系统的研究：Euler 在研究刚体转动时首次遇到主轴问题，Lagrange 在分析力学中研究了惯性椭球的主轴，而 Cauchy 在 1829 年对二次曲面分类的工作中系统地引入了特征方程和特征值概念。到了 20 世纪，Hilbert 和 von Neumann 将特征值理论推广到无穷维空间（Hilbert 空间上的谱理论），使得特征值概念成为泛函分析和量子力学的核心工具。

特征值的基本方程 \(T(v) = \lambda v\) 看似简单，但它的解------特征值 \(\lambda\) 和特征向量 \(v\)------包含了线性算子 \(T\) 的几乎全部信息。特征多项式 \(p_T(\lambda) = \det(\lambda I - T)\) 将特征值问题转化为代数方程，而 Cayley-Hamilton 定理（\(p_T(T) = 0\)）则以惊人的方式将算子与其特征多项式联系起来。代数重数与几何重数的比较------前者是特征多项式中根的重数，后者是特征空间的维数------给出了算子是否可对角化的精确判据：\(T\) 可对角化当且仅当对每个特征值，代数重数等于几何重数。

谱理论在数学和物理中的应用极为广泛。在动力系统中，特征值决定线性系统 \(\dot{\mathbf{x}} = A\mathbf{x}\) 的稳定性：若所有特征值具有负实部，则系统渐近稳定；若存在正实部特征值，则系统不稳定。在量子力学中，可观测量表示为 Hilbert 空间上的自伴算子，其特征值给出可能的测量结果，特征向量表示相应的量子态。在数据科学中，主成分分析（PCA）本质上是协方差矩阵的特征值分解，最大特征值对应的特征向量给出数据方差最大的方向。本章将系统建立特征值的代数理论，为后续的谱定理和矩阵分解理论奠定基础。

\textbf{定义 35.1.1（特征值与特征向量）}：设 \(V\) 为 \(F\) 上的向量空间，\(T \in \mathcal{L}(V)\)。若存在 \(\lambda \in F\) 和非零向量 \(v \in V\)（\(v \neq \mathbf{0}\)）使得

\[
T(v) = \lambda v
\]

则称 \(\lambda\) 为 \(T\) 的\textbf{特征值}，\(v\) 为 \(T\) 的属于 \(\lambda\) 的\textbf{特征向量}。

\textbf{定义 35.1.2（特征空间）}：设 \(\lambda\) 为 \(T\) 的特征值。\(T\) 的属于 \(\lambda\) 的\textbf{特征空间}为

\[
E_{\lambda} = \{v \in V \mid T(v) = \lambda v\} = \ker(T - \lambda I)
\]

\(E_{\lambda}\) 是 \(V\) 的子空间。\(\dim E_{\lambda}\) 称为 \(\lambda\) 的\textbf{几何重数}。

\textbf{定理 35.1.1（不同特征值的特征向量线性无关）}：设 \(\lambda_1, \ldots, \lambda_k\) 为 \(T\) 的两两不同的特征值，\(v_1, \ldots, v_k\) 为对应的特征向量。则 \(\{v_1, \ldots, v_k\}\) 线性无关。

\textbf{证明}：对 \(k\) 归纳。

\textbf{基础} \(k = 1\)：\(v_1 \neq \mathbf{0}\)，单个非零向量线性无关。

\textbf{归纳}：假设 \(\{v_1, \ldots, v_{k-1}\}\) 线性无关。设 \(\sum_{i=1}^{k} \alpha_i v_i = \mathbf{0}\)。两边作用 \(T\)： \[
\sum_{i=1}^{k} \alpha_i \lambda_i v_i = \mathbf{0}
\] 将原式乘以 \(\lambda_k\)： \[
\sum_{i=1}^{k} \alpha_i \lambda_k v_i = \mathbf{0}
\] 两式相减：\(\sum_{i=1}^{k-1} \alpha_i (\lambda_i - \lambda_k) v_i = \mathbf{0}\)。由归纳假设，\(\{v_1, \ldots, v_{k-1}\}\) 线性无关，故 \(\alpha_i (\lambda_i - \lambda_k) = 0\)（\(i = 1, \ldots, k-1\)）。由于 \(\lambda_i \neq \lambda_k\)，\(\alpha_i = 0\)。代入原式得 \(\alpha_k v_k = \mathbf{0}\)，故 \(\alpha_k = 0\)。\(\blacksquare\)

\textbf{推论 35.1.2}：\(n\) 维向量空间上的线性算子最多有 \(n\) 个不同的特征值。

\textbf{证明}：由定理 35.1.1，不同特征值对应的特征向量线性无关，故不超过 \(\dim V = n\)。\(\blacksquare\)

\hrulefill

\subsection{35.2 特征多项式}\label{ux7279ux5f81ux591aux9879ux5f0f}

\textbf{定义 35.2.1（特征多项式）}：设 \(A \in M_{n \times n}(F)\)。\(A\) 的\textbf{特征多项式}为

\[
p_A(\lambda) = \det(\lambda I_n - A)
\]

\textbf{注记 35.2.1}：有些教材定义 \(\det(A - \lambda I_n)\)，两者仅差一个符号因子 \((-1)^n\)，不影响特征值的确定。

\textbf{定理 35.2.1}：\(\lambda\) 是 \(A\) 的特征值当且仅当 \(p_A(\lambda) = 0\)，即 \(\lambda\) 是特征多项式的根。

\textbf{证明}：\(\lambda\) 是特征值 \(\iff \exists v \neq \mathbf{0}\) 使得 \(Av = \lambda v \iff (\lambda I_n - A)v = \mathbf{0} \iff \lambda I_n - A\) 不可逆 \(\iff \det(\lambda I_n - A) = 0 \iff p_A(\lambda) = 0\)。\(\blacksquare\)

\textbf{定理 35.2.2（特征多项式的形式）}：\(p_A(\lambda) = \lambda^n - (\operatorname{tr} A) \lambda^{n-1} + \cdots + (-1)^n \det A\)，其中 \(\operatorname{tr} A = \sum_{i=1}^{n} a_{ii}\) 为 \(A\) 的迹。

\textbf{证明}：展开 \(\det(\lambda I_n - A) = \sum_{\sigma} \operatorname{sgn}(\sigma) \prod_{i=1}^{n} (\lambda \delta_{i\sigma(i)} - a_{i\sigma(i)})\)。\(\lambda^n\) 项仅来自 \(\sigma\) 为恒等置换，系数为 \(1\)。\(\lambda^{n-1}\) 项来自恒等置换中某个 \(i\) 处取 \(-a_{ii}\) 而其余 \(n-1\) 处取 \(\lambda\)，系数为 \(-\sum_i a_{ii} = -\operatorname{tr} A\)。常数项为 \(p_A(0) = \det(-A) = (-1)^n \det A\)。\(\blacksquare\)

\textbf{定义 35.2.2（代数重数）}：特征值 \(\lambda\) 作为 \(p_A(\lambda)\) 的根的重数称为 \(\lambda\) 的\textbf{代数重数}。

\textbf{定理 35.2.3（几何重数不超过代数重数）}：设 \(\lambda\) 为 \(A\) 的特征值，则 \(\dim E_{\lambda} \leq \lambda\) 的代数重数。

\textbf{证明}：设 \(\dim E_{\lambda} = k\)。取 \(E_{\lambda}\) 的基 \(\{v_1, \ldots, v_k\}\)，将其扩充为 \(F^n\) 的基。在此基下，\(A\) 的矩阵为 \[
\begin{pmatrix} \lambda I_k & B \\ 0 & C \end{pmatrix}
\] 其特征多项式为 \(\det(\lambda I_n - A) = \det((\lambda - \lambda) I_k) \det(\lambda I_{n-k} - C) = 0 \cdot \det(\lambda I_{n-k} - C)\)。更精确地，\(p_A(t) = (t - \lambda)^k \det(t I_{n-k} - C)\)，故 \(t - \lambda\) 的重数至少为 \(k\)。\(\blacksquare\)

\hrulefill

\subsection{35.3 对角化}\label{ux5bf9ux89d2ux5316}

\textbf{定义 35.3.1（可对角化）}：线性算子 \(T \in \mathcal{L}(V)\)（或矩阵 \(A \in M_{n \times n}(F)\)）称为\textbf{可对角化}的，若存在 \(V\) 的一组基，使得 \(T\) 在此基下的矩阵为对角矩阵。等价地，\(V\) 有一组由 \(T\) 的特征向量组成的基。

\textbf{定理 35.3.1（可对角化的判定）}：设 \(T \in \mathcal{L}(V)\)，\(\dim V = n\)。以下等价：

\begin{enumerate}
\def\labelenumi{(\roman{enumi})}
\tightlist
\item
  \(T\) 可对角化
\item
  \(V\) 有一组由 \(T\) 的特征向量组成的基
\item
  \(V = E_{\lambda_1} \oplus E_{\lambda_2} \oplus \cdots \oplus E_{\lambda_k}\)（特征空间的直和分解）
\item
  \(T\) 有 \(n\) 个线性无关的特征向量
\end{enumerate}

\textbf{证明}：(i) \(\iff\) (ii)：由定义。(ii) \(\iff\) (iv)：维数条件。(ii) \(\iff\) (iii)：不同特征值的特征向量线性无关，且特征向量张成 \(V\)，故 \(V\) 是特征空间的直和。\(\blacksquare\)

\textbf{定理 35.3.2（充分条件）}：若 \(T \in \mathcal{L}(V)\) 有 \(n = \dim V\) 个两两不同的特征值，则 \(T\) 可对角化。

\textbf{证明}：由定理 35.1.1，\(n\) 个不同特征值的特征向量线性无关。由定理 32.5.5(ii)，\(n\) 个线性无关向量构成 \(V\) 的基。\(\blacksquare\)

\textbf{定理 35.3.3（对角化的充要条件）}：\(T \in \mathcal{L}(V)\) 可对角化当且仅当其特征多项式在 \(F\) 上可分解为一次因式的乘积，且每个特征值的几何重数等于其代数重数。

\textbf{证明}：

（\(\Rightarrow\)）若 \(T\) 可对角化，则存在基使得 \(T\) 的矩阵为 \(\operatorname{diag}(\lambda_1, \ldots, \lambda_n)\)。此时 \(p_T(t) = \prod_{i=1}^{n} (t - \lambda_i)\)，且 \(\dim E_{\lambda_i}\) 等于 \(\lambda_i\) 在 \(\{\lambda_1, \ldots, \lambda_n\}\) 中出现的次数，即代数重数。

（\(\Leftarrow\)）设 \(p_T(t) = \prod_{i=1}^{k} (t - \lambda_i)^{m_i}\)（\(\lambda_i\) 两两不同），且 \(\dim E_{\lambda_i} = m_i\)。由定理 35.1.1，不同特征空间的和是直和。\(\dim(\bigoplus_i E_{\lambda_i}) = \sum_i m_i = n = \dim V\)，故 \(\bigoplus_i E_{\lambda_i} = V\)。\(\blacksquare\)

\hrulefill

\subsection{35.4 Cayley-Hamilton 定理}\label{cayley-hamilton-ux5b9aux7406}

\textbf{定理 35.4.1（Cayley-Hamilton 定理）}：设 \(A \in M_{n \times n}(F)\)，\(p_A(\lambda) = \det(\lambda I_n - A)\) 为其特征多项式。则 \(p_A(A) = 0\)（零矩阵）。

\textbf{证明}：考虑矩阵 \(\lambda I_n - A\) 的伴随矩阵 \(\operatorname{adj}(\lambda I_n - A)\)。其元素是 \(\lambda\) 的多项式，次数不超过 \(n-1\)。由定理 34.3.2，

\[
(\lambda I_n - A) \cdot \operatorname{adj}(\lambda I_n - A) = \det(\lambda I_n - A) I_n = p_A(\lambda) I_n
\]

将 \(\operatorname{adj}(\lambda I_n - A)\) 写为 \(B_0 + B_1 \lambda + \cdots + B_{n-1} \lambda^{n-1}\)（\(B_i \in M_{n \times n}(F)\)）。设 \(p_A(\lambda) = \lambda^n + c_{n-1}\lambda^{n-1} + \cdots + c_0\)。则

\[
(\lambda I_n - A)(B_0 + B_1 \lambda + \cdots + B_{n-1} \lambda^{n-1}) = (\lambda^n + c_{n-1}\lambda^{n-1} + \cdots + c_0) I_n
\]

比较两边 \(\lambda^k\) 的系数：

\[
\begin{aligned}
\lambda^0 &: -A B_0 = c_0 I_n \\
\lambda^1 &: B_0 - A B_1 = c_1 I_n \\
&\;\;\vdots \\
\lambda^{n-1} &: B_{n-2} - A B_{n-1} = c_{n-1} I_n \\
\lambda^n &: B_{n-1} = I_n
\end{aligned}
\]

将第 \(k\) 个等式左乘 \(A^k\)（\(k = 0, \ldots, n\)）并求和： \[
\begin{aligned}
&(-A B_0) + A(B_0 - A B_1) + A^2(B_1 - A B_2) + \cdots + A^{n-1}(B_{n-2} - A B_{n-1}) + A^n B_{n-1} \\
&= -A B_0 + A B_0 - A^2 B_1 + A^2 B_1 - \cdots - A^n B_{n-1} + A^n B_{n-1} = 0
\end{aligned}
\] 而右边求和为 \(c_0 I_n + c_1 A + \cdots + c_{n-1} A^{n-1} + A^n = p_A(A)\)。故 \(p_A(A) = 0\)。\(\blacksquare\)

\textbf{推论 35.4.2}：\(A^n\) 可表示为 \(I_n, A, \ldots, A^{n-1}\) 的线性组合。更一般地，对任意 \(k \geq 0\)，\(A^k\) 属于 \(I_n, A, \ldots, A^{n-1}\) 张成的子空间。

\hrulefill

\subsection{35.5 极小多项式}\label{ux6781ux5c0fux591aux9879ux5f0f}

\textbf{定义 35.5.1（极小多项式）}：设 \(A \in M_{n \times n}(F)\)。在使得 \(q(A) = 0\) 的所有非零首一多项式 \(q \in F[t]\) 中，次数最低者称为 \(A\) 的\textbf{极小多项式}，记为 \(m_A(t)\)。

\textbf{定理 35.5.1（极小多项式的存在唯一性）}：\(m_A(t)\) 存在且唯一。

\textbf{证明}：由 Cayley-Hamilton 定理，\(p_A(A) = 0\)，故存在满足条件的非零多项式。在 \(F[t]\) 中，\(\dim M_{n \times n}(F) = n^2\)，故 \(I_n, A, \ldots, A^{n^2}\) 线性相关，因此存在次数不超过 \(n^2\) 的零化多项式。取其中次数最低的首一多项式，即为 \(m_A\)。唯一性：若 \(m_1, m_2\) 均为极小多项式，则 \(m_1 - m_2\) 次数更低，且 \((m_1 - m_2)(A) = 0\)。若 \(m_1 \neq m_2\)，则除以首项系数后得次数更低的零化多项式，矛盾。\(\blacksquare\)

\textbf{定理 35.5.2（极小多项式的基本性质）}：

\begin{enumerate}
\def\labelenumi{(\roman{enumi})}
\tightlist
\item
  \(m_A(t)\) 整除任何满足 \(q(A) = 0\) 的多项式 \(q(t)\)。特别地，\(m_A(t) \mid p_A(t)\)。
\item
  \(A\) 的特征值恰好是 \(m_A(t)\) 在 \(F\) 的代数闭包中的根。
\end{enumerate}

\textbf{证明}：

\begin{enumerate}
\def\labelenumi{(\roman{enumi})}
\item
  由多项式带余除法，\(q(t) = m_A(t) \cdot s(t) + r(t)\)，其中 \(\deg r < \deg m_A\)。代入 \(A\)：\(0 = q(A) = m_A(A) s(A) + r(A) = r(A)\)。若 \(r \neq 0\)，则 \(r\) 是比 \(m_A\) 次数更低的零化多项式，矛盾。故 \(r \equiv 0\)，即 \(m_A \mid q\)。
\item
  若 \(\lambda\) 是 \(A\) 的特征值，\(v\) 为对应特征向量，则 \(0 = m_A(A)v = m_A(\lambda)v\)，故 \(m_A(\lambda) = 0\)。反之，若 \(m_A(\lambda) = 0\)，则 \(m_A(t) = (t - \lambda)q(t)\)。由于 \(q(A) \neq 0\)（否则 \(m_A\) 不是极小），存在 \(w\) 使得 \(q(A)w \neq \mathbf{0}\)。令 \(v = q(A)w\)，则 \((A - \lambda I)v = (A - \lambda I)q(A)w = m_A(A)w = \mathbf{0}\)，故 \(\lambda\) 是特征值。\(\blacksquare\)
\end{enumerate}

\textbf{定理 35.5.3（可对角化与极小多项式）}：\(A \in M_{n \times n}(F)\) 可对角化当且仅当 \(m_A(t)\) 在 \(F\) 上可分解为互不相同的一次因式的乘积。

\textbf{证明}：

（\(\Rightarrow\)）若 \(A\) 可对角化，设 \(A = P \operatorname{diag}(\lambda_1, \ldots, \lambda_n) P^{-1}\)。则 \(m_A(A) = P m_A(\operatorname{diag}(\lambda_1, \ldots, \lambda_n)) P^{-1}\)。\(m_A(\operatorname{diag}(\lambda_1, \ldots, \lambda_n)) = \operatorname{diag}(m_A(\lambda_1), \ldots, m_A(\lambda_n))\)。故 \(m_A(\lambda_i) = 0\) 对所有 \(i\) 成立。取 \(m_A(t) = \prod_{\lambda \in \{\lambda_1, \ldots, \lambda_n\}} (t - \lambda)\)（重复的只取一次），则 \(m_A(A) = 0\) 且 \(m_A\) 无重根。

（\(\Leftarrow\)）设 \(m_A(t) = \prod_{i=1}^{k} (t - \lambda_i)\)（\(\lambda_i\) 两两不同）。需证 \(F^n\) 是特征空间 \(E_{\lambda_i}\) 的直和。

对 \(k\) 归纳。\(k = 1\) 时，\(m_A(t) = t - \lambda_1\)，即 \(A = \lambda_1 I_n\)，已是对角矩阵。

\(k > 1\) 时，由部分分式分解，\(\frac{1}{m_A(t)} = \sum_{i=1}^{k} \frac{c_i}{t - \lambda_i}\)，故 \(1 = \sum_{i=1}^{k} c_i \prod_{j \neq i} (t - \lambda_j)\)。代入 \(A\)：\(I_n = \sum_{i=1}^{k} c_i \prod_{j \neq i} (A - \lambda_j I_n)\)。令 \(P_i = c_i \prod_{j \neq i} (A - \lambda_j I_n)\)，则 \(I_n = \sum_i P_i\)。

对任意 \(v \in F^n\)，\(v = \sum_i P_i v\)。由 \((A - \lambda_i I_n)P_i = c_i m_A(A) = 0\)，故 \(P_i v \in E_{\lambda_i}\)。因此 \(F^n = \sum_i E_{\lambda_i}\)。由定理 35.1.1，不同特征空间的和是直和。故 \(A\) 可对角化。\(\blacksquare\)

\hrulefill

\subsection{35.6 Jordan 标准形（陈述）}\label{jordan-ux6807ux51c6ux5f62ux9648ux8ff0}

\textbf{定义 35.6.1（Jordan 块）}：设 \(\lambda \in F\)。\(k \times k\) 的 \textbf{Jordan 块} \(J_k(\lambda)\) 定义为

\[
J_k(\lambda) = \begin{pmatrix}
\lambda & 1 & 0 & \cdots & 0 \\
0 & \lambda & 1 & \cdots & 0 \\
\vdots & \vdots & \ddots & \ddots & \vdots \\
0 & 0 & \cdots & \lambda & 1 \\
0 & 0 & \cdots & 0 & \lambda
\end{pmatrix}_{k \times k}
\]

\textbf{定理 35.6.1（Jordan 标准形，陈述）}：设 \(F\) 为代数闭域（如 \(\mathbb{C}\)），\(A \in M_{n \times n}(F)\)。则存在可逆矩阵 \(P \in M_{n \times n}(F)\) 使得

\[
P^{-1}AP = \begin{pmatrix}
J_{k_1}(\lambda_1) & 0 & \cdots & 0 \\
0 & J_{k_2}(\lambda_2) & \cdots & 0 \\
\vdots & \vdots & \ddots & \vdots \\
0 & 0 & \cdots & J_{k_r}(\lambda_r)
\end{pmatrix}
\]

其中 \(J_{k_i}(\lambda_i)\) 为 Jordan 块。该分解在 Jordan 块的排列顺序下是唯一的。

\textbf{注记 35.6.1}：Jordan 标准形的完整证明需要商空间和广义特征空间的理论，将在卷八（抽象代数 I）中作为模论的应用给出。此处先陈述结果，以便在后续应用中使用。

\hrulefill

\section{第107章：矩阵扰动理论}\label{ux7b2c107ux7ae0ux77e9ux9635ux6270ux52a8ux7406ux8bba}

矩阵的特征值对矩阵元素的微小变化可能极为敏感，也可能高度稳定，这取决于矩阵的正规性程度和特征值的分离状况。矩阵扰动理论定量研究这种敏感性，是数值线性代数和应用数学的理论基石。本章系统建立四个核心定理：Bauer-Fike定理------通过条件数刻画可对角化矩阵特征值对扰动的敏感性界限；Weyl不等式------建立Hermite矩阵特征值的交错关系和扰动界（特征值以Lipschitz方式依赖于矩阵）；Lidskii不等式------以弱多数化（weak majorization）形式表述Hermite矩阵特征值差的最优界；以及Hoffman-Wielandt定理------在酉不变范数（特别是Frobenius范数）下给出两个正规矩阵特征值之间的距离与矩阵差之间的关系。

\subsection{Bauer-Fike定理}\label{bauer-fikeux5b9aux7406}

\textbf{定理58.1（Bauer-Fike定理）} 设 \(A \in M_n(\mathbb{C})\) 可对角化：\(A = X \Lambda X^{-1}\)，其中 \(\Lambda = \operatorname{diag}(\lambda_1, \ldots, \lambda_n)\)。令 \(E \in M_n(\mathbb{C})\) 为扰动矩阵。则对 \(A + E\) 的任意特征值 \(\mu\)，存在 \(A\) 的特征值 \(\lambda_i\) 使得

\[|\mu - \lambda_i| \leq \kappa_p(X) \|E\|_p\]

其中 \(\kappa_p(X) = \|X\|_p \|X^{-1}\|_p\) 是 \(X\) 在 \(p\)-范数下的条件数（\(p = 1, 2, \infty\)）。

\textbf{证明} 若 \(\mu\) 是 \(A\) 的特征值，则不等式显然成立（取 \(\lambda_i = \mu\)，右端为 \(0\)）。设 \(\mu \notin \sigma(A)\)，即 \(\mu I - A\) 可逆。

由 \(A = X \Lambda X^{-1}\)，\(\mu I - A = X(\mu I - \Lambda) X^{-1}\)。考虑

\[(\mu I - (A + E)) = (\mu I - A)(I - (\mu I - A)^{-1} E)\]

若 \(\mu\) 是 \(A + E\) 的特征值，则 \(\mu I - (A + E)\) 奇异。由于 \(\mu I - A\) 可逆，\(I - (\mu I - A)^{-1} E\) 必奇异，故 \(\|(\mu I - A)^{-1} E\|_p \geq 1\)（因若算子范数 \(< 1\) 则 \(I - T\) 可逆）。因此

\[1 \leq \|(\mu I - A)^{-1}\|_p \|E\|_p = \|X(\mu I - \Lambda)^{-1} X^{-1}\|_p \|E\|_p \leq \kappa_p(X) \max_i |\mu - \lambda_i|^{-1} \|E\|_p\]

故存在 \(i\) 使 \(|\mu - \lambda_i|^{-1} \geq \max_j |\mu - \lambda_j|^{-1} \geq (\kappa_p(X) \|E\|_p)^{-1}\)，即 \(|\mu - \lambda_i| \leq \kappa_p(X) \|E\|_p\)。\(\blacksquare\)

\textbf{推论58.2} 当 \(A\) 是正规矩阵时，可取 \(X\) 为酉矩阵，从而 \(\kappa_2(X) = 1\)，此时 \(|\mu - \lambda_i| \leq \|E\|_2\)。正规矩阵的特征值在谱范数下以 \(1\)-Lipschitz方式依赖于矩阵------这是数值稳定性最好的情形。

\subsection{Weyl不等式}\label{weylux4e0dux7b49ux5f0f}

\textbf{定理58.3（Weyl单调性定理）} 设 \(A, B \in M_n(\mathbb{C})\) 是 Hermite 矩阵，且 \(A \preceq B\)（即 \(B - A\) 半正定）。记特征值按升序排列：\(\lambda_1(A) \leq \cdots \leq \lambda_n(A)\)，\(\lambda_1(B) \leq \cdots \leq \lambda_n(B)\)。则

\[\lambda_k(A) \leq \lambda_k(B), \quad k = 1, \ldots, n\]

即 Hermite 矩阵特征值在 Loewner 偏序下单调。

\textbf{证明} 由 Courant-Fischer 极小极大原理：

\[\lambda_k(A) = \min_{\dim V = n - k + 1} \max_{x \in V, \|x\|=1} x^* A x\]

（或等价地 \(\lambda_k(A) = \max_{\dim V = k} \min_{x \in V, \|x\|=1} x^* A x\)）。由 \(A \preceq B\)，对任意单位向量 \(x\)，\(x^* A x \leq x^* B x\)。因此对每个子空间 \(V\)，\(\max_{x \in V} x^* A x \leq \max_{x \in V} x^* B x\)。取 \(\min_{\dim V}\) 保持不等式方向，故 \(\lambda_k(A) \leq \lambda_k(B)\)。\(\blacksquare\)

\textbf{定理58.4（Weyl扰动定理）} 设 \(A, E \in M_n(\mathbb{C})\) 是 Hermite 矩阵。则对每个 \(k = 1, \ldots, n\)，

\[\lambda_k(A) + \lambda_1(E) \leq \lambda_k(A + E) \leq \lambda_k(A) + \lambda_n(E)\]

等价地，\(|\lambda_k(A + E) - \lambda_k(A)| \leq \|E\|_2 = \max\{|\lambda_1(E)|, |\lambda_n(E)|\}\)。

\textbf{证明} 由 \(A + \lambda_1(E) I \preceq A + E \preceq A + \lambda_n(E) I\)（因为 \(E - \lambda_1(E) I \succeq 0\) 且 \(\lambda_n(E) I - E \succeq 0\)）。对这两个不等式应用定理58.3（Weyl单调性），得

\[\lambda_k(A + \lambda_1(E) I) \leq \lambda_k(A + E) \leq \lambda_k(A + \lambda_n(E) I)\]

即 \(\lambda_k(A) + \lambda_1(E) \leq \lambda_k(A + E) \leq \lambda_k(A) + \lambda_n(E)\)。\(\blacksquare\)

\subsection{Lidskii不等式}\label{lidskiiux4e0dux7b49ux5f0f}

\textbf{定理58.5（Lidskii不等式------弱多数化形式）} 设 \(A, B \in M_n(\mathbb{C})\) 是 Hermite 矩阵。则特征值向量满足弱多数化（weak majorization）：

\[\sum_{i=1}^k |\lambda_i(A) - \lambda_i(B)| \leq \sum_{i=1}^k \lambda_i(|A - B|), \quad k = 1, \ldots, n\]

其中 \(|A - B| = \sqrt{(A - B)^2}\) 是 \(A - B\) 的绝对值（通过函数演算定义），\(\lambda_i(|A - B|)\) 降序排列（即 \(\lambda_1 \geq \cdots \geq \lambda_n\)，为 \(A - B\) 的奇异值）。当 \(k = n\) 时等号成立（即所有特征值差的绝对值之和小于等于所有奇异值之和）。

\textbf{证明}（框架，Lidskii-Wielandt）利用 Hermite 矩阵奇异值与特征值的关系：对 Hermite 矩阵 \(H\)，其奇异值为 \(|H|\) 的特征值 \(\lambda_i(|H|)\)。对 \(A\) 和 \(B\) 构造 \(2n \times 2n\) 分块 Hermite 矩阵 \(\begin{pmatrix} A & 0 \\ 0 & -B \end{pmatrix}\)，其特征值为 \(\{\lambda_i(A)\} \cup \{-\lambda_i(B)\}\)。对 \(C = \begin{pmatrix} 0 & A - B \\ A - B & 0 \end{pmatrix}\)，其特征值为 \(\{\pm s_i(A - B)\}\)（\(s_i\) 为奇异值）。考虑 \(A \oplus (-B)\) 与 \(C\) 的相似性经由分块酉矩阵，由Wielandt极大极小原理或交错不等式推得特征值差的多数化关系。\(\blacksquare\)

\textbf{推论58.6（Wielandt-Hoffman定理------Frobenius范数版本）} 对 Hermite 矩阵 \(A, B\)，

\[\sum_{i=1}^n |\lambda_i(A) - \lambda_i(B)|^2 \leq \|A - B\|_F^2 = \operatorname{Tr}((A - B)^2)\]

即特征值差向量的 \(\ell^2\) 范数不超过矩阵差的 Frobenius 范数。特别地，Hermite矩阵的特征值是 \(1\)-Lipschitz 连续的（关于Frobenius范数）。

\subsection{Hoffman-Wielandt定理}\label{hoffman-wielandtux5b9aux7406}

\textbf{定理58.7（Hoffman-Wielandt定理）} 设 \(A, B \in M_n(\mathbb{C})\) 是正规矩阵（即 \(A A^* = A^* A\)，\(B B^* = B^* B\)）。则存在特征值的排列 \(\{\lambda_i(A)\}\) 和 \(\{\lambda_{\sigma(i)}(B)\}\) 使得

\[\sum_{i=1}^n |\lambda_i(A) - \lambda_{\sigma(i)}(B)|^2 \leq \|A - B\|_F^2\]

\textbf{证明} 设 \(A = U \Lambda U^*\)，\(B = V M V^*\)（谱分解），\(\Lambda = \operatorname{diag}(\lambda_1^A, \ldots, \lambda_n^A)\)，\(M = \operatorname{diag}(\lambda_1^B, \ldots, \lambda_n^B)\)。令 \(W = U^* V = (w_{ij})\) 为酉矩阵。则

\[\|A - B\|_F^2 = \|U \Lambda U^* - V M V^*\|_F^2 = \|\Lambda W - W M\|_F^2 = \sum_{i,j} |\lambda_i^A - \lambda_j^B|^2 |w_{ij}|^2\]

因 \(W\) 是酉矩阵，\(D = (|w_{ij}|^2)\) 是双随机矩阵（每行每列和为 \(1\)）。上述和式是最优传输问题（optimal transport）的目标函数：在双随机矩阵 \(D\) 上极小化 \(\sum_{i,j} c_{ij} d_{ij}\)，其中 \(c_{ij} = |\lambda_i^A - \lambda_j^B|^2\)。由Birkhoff-von Neumann定理，双随机矩阵的极点为置换矩阵，最小值在置换矩阵 \(\sigma\) 处取得：

\[\min_{D \text{ 双随机}} \sum_{i,j} |\lambda_i^A - \lambda_j^B|^2 d_{ij} = \min_{\sigma \in S_n} \sum_i |\lambda_i^A - \lambda_{\sigma(i)}^B|^2\]

因此存在排列 \(\sigma\) 使 \(\sum_i |\lambda_i(A) - \lambda_{\sigma(i)}(B)|^2 \leq \|A - B\|_F^2\)。\(\blacksquare\)

\hrulefill

\section{第108章：非负矩阵与Perron-Frobenius理论}\label{ux7b2c108ux7ae0ux975eux8d1fux77e9ux9635ux4e0eperron-frobeniusux7406ux8bba}

非负矩阵（所有元素非负的矩阵）在经济学（投入产出模型）、人口动力学（Leslie矩阵）、网络分析（PageRank算法）和Markov链理论中无处不在。Perron-Frobenius定理是非负矩阵谱理论的核心，它断言正矩阵（所有元素严格正）的谱半径本身是正特征值，且对应的特征向量可选取为正向量------这一定理有着优美的解析和组合证明。本章建立完整的Perron-Frobenius理论：从正矩阵的基本定理出发，推广到非负不可约矩阵，讨论M矩阵理论（其逆矩阵非负的Z矩阵），最后给出Google PageRank算法的严格数学基础------将其解释为随机矩阵的主特征向量问题。

\subsection{Perron-Frobenius定理}\label{perron-frobeniusux5b9aux7406}

\textbf{定理59.1（Perron定理------正矩阵）} 设 \(A \in M_n(\mathbb{R})\) 是正矩阵（\(A > 0\)，即所有 \(a_{ij} > 0\)）。则：

（i）谱半径 \(\rho(A) = \max\{|\lambda| : \lambda \in \sigma(A)\}\) 是 \(A\) 的特征值，称为\textbf{Perron根}。

（ii）\(\rho(A) > 0\)，且它代数重数为 \(1\)（单特征值）。

（iii）存在正特征向量 \(x > 0\)（所有分量 \(> 0\)）满足 \(A x = \rho(A) x\)，以及正左特征向量 \(y > 0\) 满足 \(y^T A = \rho(A) y^T\)。

（iv）\(\rho(A)\) 严格支配所有其他特征值的模：对任意 \(\lambda \in \sigma(A) \setminus \{\rho(A)\}\)，\(|\lambda| < \rho(A)\)。

（v）\(\rho(A)\) 随每个元素 \(a_{ij}\) 严格单调递增。

\textbf{证明}（Perron，1907；Frobenius，1912）

\textbf{第1步（正锥的Brouwer不动点论证）}：考虑标准单纯形 \(\Delta = \{x \in \mathbb{R}^n_{\geq 0} : \sum_i x_i = 1\}\)。定义映射 \(f: \Delta \to \Delta\) 为 \(f(x) = A x / \sum_i (A x)_i\)（若分母非零）。由 \(A > 0\)，若 \(x \in \Delta\) 且 \(x \neq 0\)，则 \(A x > 0\)，故 \(f\) 连续且将 \(\Delta\) 的非零部分映入 \(\Delta\) 的内部。由 Brouwer 不动点定理，存在 \(x \in \Delta\) 使 \(f(x) = x\)，即 \(A x = r x\)，其中 \(r = \sum_i (A x)_i > 0\)。该 \(x > 0\)（因 \(A x > 0\)），且 \(r\) 是对应特征值。

\textbf{第2步（\(r = \rho(A)\) 的证明）}：设 \(\lambda\) 为 \(A\) 的任意特征值，\(z \neq 0\) 为对应特征向量。取绝对值 \(|z| = (|z_1|, \ldots, |z_n|)^T\)。由三角不等式：\(|\lambda| |z| = |\lambda z| = |A z| \leq A |z|\)（因 \(A > 0\) 保证 \(|A z| \leq A |z|\) 逐分量成立）。令 \(y > 0\) 为 \(A^T\) 的正特征向量（\(A^T y = r y\)），\(y\) 的存在性由对称论证保证。则

\[|\lambda| y^T |z| \leq y^T A |z| = (A^T y)^T |z| = r y^T |z|\]

因 \(y^T |z| > 0\)（\(y > 0\) 且 \(z \neq 0\)），得 \(|\lambda| \leq r\)。故 \(\rho(A) = r\)。

\textbf{第3步（代数单重性与严格支配性）}：设 \(\lambda \in \sigma(A)\)，\(|\lambda| = \rho(A)\)。需证 \(\lambda = \rho(A)\)。采用 Jordan 标准形和极限论证：若 \(\lambda \neq \rho(A)\)，则 \(|\lambda| = \rho(A) = r\)。考虑 \(A^n / r^n\) 的极限行为。由 Jordan 分解，若 \(r\) 是多重特征值或存在模为 \(r\) 的其他特征值，则 \(A^n/r^n\) 会振荡。但由 Perron 根的性质，\(\lim_{n\to\infty} (A/\rho(A))^n = x y^T / (y^T x)\)（其中 \(x,y\) 为正的右、左特征向量，归一化使 \(y^T x = 1\)）。此极限存在且为正矩阵，要求 \(\rho(A)\) 是单特征值且严格支配其他所有特征值。\(\blacksquare\)

\textbf{定理59.2（Frobenius定理------非负不可约矩阵）} 设 \(A \in M_n(\mathbb{R})\) 是非负矩阵（\(A \geq 0\)）且不可约（其有向图强连通，即对任意 \(i,j\) 存在 \(k\) 使 \((A^k)_{ij} > 0\)）。则定理59.1的结论（i）-（iv）均成立，且额外有：

（vi）\(\rho(A)\) 是特征多项式的单根当且仅当以 \(\rho(A)\) 为根的非本原性指标（index of imprimitivity）\(h = 1\)。一般地，\(h\) 个特征值有模 \(\rho(A)\)：\(\rho(A) e^{2\pi i k / h}\)（\(k = 0, \ldots, h-1\)）。

\textbf{证明}（框架）非负不可约矩阵的典型技巧是考虑 \((A + \varepsilon I)^n\)（\(\varepsilon > 0\)），它是正矩阵（由不可约性，对充分大 \(n\)，\((A + \varepsilon I)^n > 0\)），然后对正矩阵应用Perron定理并令 \(\varepsilon \to 0^+\)。不可约性保证谱半径 \(> 0\)（除非 \(A = 0\) 且 \(n = 1\)）。本原性指标 \(h\) 是使 \(A^k > 0\) 的 \(k\) 的最大公约数。\(\blacksquare\)

\subsection{M矩阵理论}\label{mux77e9ux9635ux7406ux8bba}

\textbf{定义59.3（Z矩阵与M矩阵）} 矩阵 \(A = (a_{ij}) \in M_n(\mathbb{R})\) 称为\textbf{Z矩阵}，若其所有非对角元非正：\(a_{ij} \leq 0\)（\(i \neq j\)）。Z矩阵 \(A\) 称为\textbf{M矩阵}，若 \(A\) 可逆且 \(A^{-1} \geq 0\)（逆矩阵所有元素非负）。

\textbf{定理59.4（M矩阵的等价刻画）} 对Z矩阵 \(A\)，以下等价：（i）\(A\) 是M矩阵；（ii）所有主子式为正；（iii）所有特征值的实部为正（即 \(A\) 是正稳定矩阵）；（iv）存在正向量 \(x > 0\) 使 \(A x > 0\)。

\textbf{证明}（（i）\(\Rightarrow\)（ii））由 \(A^{-1} \geq 0\) 和 Cauchy-Binet 公式，各主子式的逆可用 \(A^{-1}\) 的子矩阵的行列式表示，故为正。（ii）\(\Rightarrow\)（iii）通过特征值的连续性论证。（iii）\(\Rightarrow\)（iv）利用 \(A = s I - B\)（\(s > 0\)，\(B \geq 0\)）的分解和 Perron-Frobenius 定理：\(A^{-1} = s^{-1} \sum_{k=0}^\infty (B/s)^k \geq 0\) 当 \(s > \rho(B)\)。（iv）\(\Rightarrow\)（i）由不可约情形推广至一般情形，利用 \(x > 0\) 和 \(A x > 0\) 构造单调迭代证明 \(A^{-1} \geq 0\)。\(\blacksquare\)

\subsection{Google PageRank的数学基础}\label{google-pagerankux7684ux6570ux5b66ux57faux7840}

\textbf{定义59.5（PageRank矩阵）} 设有向图 \(G = (V, E)\)（\(|V| = n\)）表示网页之间的超链接关系。出度 \(d_i = |\{j : (i, j) \in E\}|\)。定义修正邻接矩阵：

\[P_{ij} = \begin{cases} \alpha \cdot \frac{1}{d_i} + (1 - \alpha) \cdot \frac{1}{n}, & \text{若 } d_i > 0 \\ \frac{1}{n}, & \text{若 } d_i = 0 \end{cases}\]

其中 \(\alpha \in (0, 1)\) 是阻尼因子（通常 \(\alpha = 0.85\)）。\(P\) 是随机矩阵（每行和为 \(1\)）。

\textbf{定理59.6（PageRank向量的存在唯一性）} \(P\) 是正随机矩阵（\(P > 0\) 严格正）。由 Perron 定理，\(P\) 有唯一的正特征向量 \(\pi\)（\(\sum_i \pi_i = 1\)，\(\pi_i > 0\)）满足 \(\pi^T P = \pi^T\)（主特征值 \(\rho(P) = 1\)）。\(\pi\) 是PageRank向量，其分量 \(\pi_i\) 是网页 \(i\) 的PageRank值。

\textbf{证明} \(P > 0\) 由构造保证（即使 \(d_i = 0\) 的行也被填充为正）。\(P\) 是随机矩阵，故 \(P \mathbf{1} = \mathbf{1}\)（\(\mathbf{1}\) 是全1向量），\(\rho(P) = 1\) 且 \(\mathbf{1}\) 是右特征向量。由 Perron 定理，\(\rho(P^T) = \rho(P) = 1\)，且存在唯一的正左特征向量 \(\pi\) 满足 \(\pi^T P = \pi^T\)，\(\|\pi\|_1 = 1\)。\(\blacksquare\)

\textbf{推论59.7（PageRank的Markov链解释）} \(\pi\) 是带有转移矩阵 \(P\) 的有限状态Markov链的平稳分布。阻尼因子 \(\alpha\) 解释为以概率 \(\alpha\) 随机跟随出链、以概率 \(1 - \alpha\) 随机跳转到任意网页的''随机冲浪者''模型。\(\pi_i\) 的次序给出了网页重要性的排序。幂迭代法 \(\pi^{(k+1)} = \pi^{(k)} P\) 以几何速率（速率 \(\approx \alpha\)）收敛于 \(\pi\)，即PageRank的标准计算方法。

\hrulefill

\hrulefill

\hrulefill

\hrulefill

\hrulefill

\newpage

\chapter{05-代数/01-线性代数/05-内积空间与二次型.md}\label{ux4ee3ux657001-ux7ebfux6027ux4ee3ux657005-ux5185ux79efux7a7aux95f4ux4e0eux4e8cux6b21ux578b.md}

\section{第109章：内积空间与谱定理}\label{ux7b2c109ux7ae0ux5185ux79efux7a7aux95f4ux4e0eux8c31ux5b9aux7406}

本章在 \(\mathbb{R}\) 上讨论。涉及 \(\mathbb{C}\) 的推广将在卷十（复分析）和卷十四（泛函分析）中处理。

\subsection{36.1 内积与范数}\label{ux5185ux79efux4e0eux8303ux6570}

\textbf{定义 36.1.1（内积）}：设 \(V\) 为 \(\mathbb{R}\) 上的向量空间。映射 \(\langle\cdot, \cdot\rangle : V \times V \to \mathbb{R}\) 若满足以下四条公理，则称为 \(V\) 上的\textbf{内积}：

\begin{enumerate}
\def\labelenumi{\arabic{enumi}.}
\tightlist
\item
  \textbf{对称性}：\(\langle u, v \rangle = \langle v, u \rangle\)
\item
  \textbf{第一变元的线性性}：\(\langle \lambda u + \mu v, w \rangle = \lambda \langle u, w \rangle + \mu \langle v, w \rangle\)
\item
  \textbf{正定性}：\(\langle v, v \rangle \geq 0\)，且 \(\langle v, v \rangle = 0 \iff v = \mathbf{0}\)
\end{enumerate}

（由对称性，第二变元也是线性的。）

配备内积的向量空间称为\textbf{内积空间}。有限维实内积空间称为\textbf{Euclid 空间}。

\textbf{定理 36.1.1（Cauchy-Schwarz 不等式）}：设 \(V\) 为内积空间，则 \(\forall u, v \in V\)，

\[
|\langle u, v \rangle| \leq \sqrt{\langle u, u \rangle} \sqrt{\langle v, v \rangle}
\]

等号成立当且仅当 \(u\) 与 \(v\) 线性相关。

\textbf{证明}：若 \(v = \mathbf{0}\)，两边均为 \(0\)，不等式成立。

设 \(v \neq \mathbf{0}\)。对任意 \(t \in \mathbb{R}\)，由正定性： \[
0 \leq \langle u - tv, u - tv \rangle = \langle u, u \rangle - 2t\langle u, v \rangle + t^2 \langle v, v \rangle
\] 这是关于 \(t\) 的二次函数。其判别式 \(\Delta = 4\langle u, v \rangle^2 - 4\langle u, u \rangle \langle v, v \rangle \leq 0\)，即 \(\langle u, v \rangle^2 \leq \langle u, u \rangle \langle v, v \rangle\)。开方即得。

等号成立当且仅当 \(\Delta = 0\)，此时存在 \(t\) 使得 \(\langle u - tv, u - tv \rangle = 0\)，即 \(u = tv\)。\(\blacksquare\)

\textbf{定义 36.1.2（范数）}：由内积诱导的\textbf{范数}定义为 \(\|v\| = \sqrt{\langle v, v \rangle}\)。

\textbf{定理 36.1.2（范数的基本性质）}：

\begin{enumerate}
\def\labelenumi{(\roman{enumi})}
\tightlist
\item
  \textbf{正定性}：\(\|v\| \geq 0\)，且 \(\|v\| = 0 \iff v = \mathbf{0}\)
\item
  \textbf{齐次性}：\(\|\lambda v\| = |\lambda| \|v\|\)
\item
  \textbf{三角不等式}：\(\|u + v\| \leq \|u\| + \|v\|\)
\end{enumerate}

\textbf{证明}：

\begin{enumerate}
\def\labelenumi{(\roman{enumi})}
\item
  由内积的正定性直接得到。
\item
  \(\|\lambda v\|^2 = \langle \lambda v, \lambda v \rangle = \lambda^2 \langle v, v \rangle = \lambda^2 \|v\|^2\)。
\item
  \(\|u + v\|^2 = \langle u+v, u+v \rangle = \|u\|^2 + 2\langle u, v \rangle + \|v\|^2 \leq \|u\|^2 + 2\|u\|\|v\| + \|v\|^2 = (\|u\| + \|v\|)^2\)。开方即得。\(\blacksquare\)
\end{enumerate}

\textbf{定义 36.1.3（标准内积）}：\(\mathbb{R}^n\) 上的\textbf{标准内积}（点积）定义为

\[
\langle x, y \rangle = \sum_{i=1}^{n} x_i y_i
\]

\hrulefill

\subsection{36.2 正交性}\label{ux6b63ux4ea4ux6027}

\textbf{定义 36.2.1（正交）}：设 \(V\) 为内积空间。

\begin{enumerate}
\def\labelenumi{\arabic{enumi}.}
\tightlist
\item
  若 \(\langle u, v \rangle = 0\)，则称 \(u\) 与 \(v\) \textbf{正交}，记为 \(u \perp v\)
\item
  向量组 \(\{v_1, \ldots, v_k\}\) 称为\textbf{正交组}，若两两正交
\item
  若正交组中每个向量的范数均为 \(1\)，则称为\textbf{规范正交组}（或\textbf{标准正交组}）
\end{enumerate}

\textbf{定理 36.2.1（勾股定理）}：若 \(u \perp v\)，则 \(\|u + v\|^2 = \|u\|^2 + \|v\|^2\)。

\textbf{证明}：\(\|u + v\|^2 = \langle u+v, u+v \rangle = \|u\|^2 + 2\langle u, v \rangle + \|v\|^2 = \|u\|^2 + \|v\|^2\)。\(\blacksquare\)

\textbf{定理 36.2.2（正交组线性无关）}：非零向量组成的正交组线性无关。

\textbf{证明}：设 \(\{v_1, \ldots, v_k\}\) 为正交组，\(v_i \neq \mathbf{0}\)。若 \(\sum \lambda_i v_i = \mathbf{0}\)，则对任意 \(j\)， \[
0 = \langle \mathbf{0}, v_j \rangle = \left\langle \sum_i \lambda_i v_i, v_j \right\rangle = \lambda_j \langle v_j, v_j \rangle
\] 由于 \(\langle v_j, v_j \rangle > 0\)，\(\lambda_j = 0\)。\(\blacksquare\)

\textbf{推论 36.2.3}：\(n\) 维内积空间中，任何 \(n\) 个非零向量组成的正交组是基（称为\textbf{正交基}）。若进一步是规范正交的，则称为\textbf{规范正交基}。

\hrulefill

\subsection{36.3 Gram-Schmidt 正交化}\label{gram-schmidt-ux6b63ux4ea4ux5316}

\textbf{定理 36.3.1（Gram-Schmidt 正交化）}：设 \(\{v_1, \ldots, v_n\}\) 为内积空间 \(V\) 的基。则可构造 \(V\) 的正交基 \(\{e_1, \ldots, e_n\}\) 如下：

\[
\begin{aligned}
e_1 &= v_1 \\
e_2 &= v_2 - \frac{\langle v_2, e_1 \rangle}{\langle e_1, e_1 \rangle} e_1 \\
e_3 &= v_3 - \frac{\langle v_3, e_1 \rangle}{\langle e_1, e_1 \rangle} e_1 - \frac{\langle v_3, e_2 \rangle}{\langle e_2, e_2 \rangle} e_2 \\
&\;\;\vdots \\
e_k &= v_k - \sum_{i=1}^{k-1} \frac{\langle v_k, e_i \rangle}{\langle e_i, e_i \rangle} e_i
\end{aligned}
\]

进一步，单位化 \(u_i = \frac{e_i}{\|e_i\|}\) 得到规范正交基。

\textbf{证明}：对 \(k\) 归纳证明 \(\{e_1, \ldots, e_k\}\) 是正交组。

\textbf{基础} \(k = 1\)：单个向量自动正交。

\textbf{归纳}：设 \(\{e_1, \ldots, e_{k-1}\}\) 为正交组。对 \(j < k\)， \[
\langle e_k, e_j \rangle = \langle v_k, e_j \rangle - \sum_{i=1}^{k-1} \frac{\langle v_k, e_i \rangle}{\langle e_i, e_i \rangle} \langle e_i, e_j \rangle = \langle v_k, e_j \rangle - \frac{\langle v_k, e_j \rangle}{\langle e_j, e_j \rangle} \langle e_j, e_j \rangle = 0
\] （因为当 \(i \neq j\) 时 \(\langle e_i, e_j \rangle = 0\)）。故 \(\{e_1, \ldots, e_k\}\) 正交。

由构造，\(\operatorname{span}\{e_1, \ldots, e_k\} = \operatorname{span}\{v_1, \ldots, v_k\}\)，故 \(e_k \neq \mathbf{0}\)（否则 \(v_k\) 属于前 \(k-1\) 个向量的张成空间，与基的线性无关性矛盾）。\(\blacksquare\)

\textbf{推论 36.3.2}：任何有限维内积空间存在规范正交基。

\textbf{推论 36.3.3（QR 分解）}：设 \(A \in M_{n \times n}(\mathbb{R})\) 可逆。则存在正交矩阵 \(Q\)（\(Q^T Q = I_n\)）和上三角矩阵 \(R\)（对角线元素为正）使得 \(A = QR\)。

\textbf{证明}：将 \(A\) 的列向量组视为 \(\mathbb{R}^n\) 的基，对其进行 Gram-Schmidt 正交化并单位化，得到规范正交基。\(Q\) 的列是规范正交化后的向量，\(R\) 记录线性组合的系数。\(\blacksquare\)

\hrulefill

\subsection{36.4 正交补与正交投影}\label{ux6b63ux4ea4ux8865ux4e0eux6b63ux4ea4ux6295ux5f71}

\textbf{定义 36.4.1（正交补）}：设 \(W \leq V\)。\(W\) 的\textbf{正交补}定义为

\[
W^{\perp} = \{v \in V \mid \langle v, w \rangle = 0, \; \forall w \in W\}
\]

\textbf{定理 36.4.1}：\(W^{\perp}\) 是 \(V\) 的子空间，且 \(V = W \oplus W^{\perp}\)。

\textbf{证明}：\(W^{\perp}\) 是子空间：由内积的线性性直接验证。

取 \(W\) 的规范正交基 \(\{e_1, \ldots, e_k\}\)，将其扩充为 \(V\) 的规范正交基 \(\{e_1, \ldots, e_k, e_{k+1}, \ldots, e_n\}\)（通过 Gram-Schmidt 正交化）。则 \(\{e_{k+1}, \ldots, e_n\}\) 是 \(W^{\perp}\) 的基。故 \(V = W \oplus W^{\perp}\)。\(\blacksquare\)

\textbf{推论 36.4.2}：\(\dim W^{\perp} = \dim V - \dim W\)，且 \((W^{\perp})^{\perp} = W\)。

\textbf{定义 36.4.2（正交投影）}：由 \(V = W \oplus W^{\perp}\)，对任意 \(v \in V\)，存在唯一的分解 \(v = w + w^{\perp}\)（\(w \in W\)，\(w^{\perp} \in W^{\perp}\)）。映射 \(P_W : V \to V\)，\(P_W(v) = w\) 称为 \(V\) 到 \(W\) 的\textbf{正交投影}。

\textbf{定理 36.4.2（正交投影的性质）}：

\begin{enumerate}
\def\labelenumi{(\roman{enumi})}
\tightlist
\item
  \(P_W\) 是线性算子，且 \(P_W^2 = P_W\)（幂等）
\item
  \(\operatorname{im} P_W = W\)，\(\ker P_W = W^{\perp}\)
\item
  \(P_W\) 是自伴的：\(\langle P_W u, v \rangle = \langle u, P_W v \rangle\)
\item
  若 \(\{e_1, \ldots, e_k\}\) 为 \(W\) 的规范正交基，则 \(P_W(v) = \sum_{i=1}^{k} \langle v, e_i \rangle e_i\)
\end{enumerate}

\textbf{证明}：(i) 线性性由分解的唯一性验证。\(P_W^2 = P_W\) 因为 \(P_W(w) = w\)（\(w \in W\)）。

\begin{enumerate}
\def\labelenumi{(\roman{enumi})}
\setcounter{enumi}{1}
\item
  显然。(iii) 设 \(u = u_W + u_{\perp}\)，\(v = v_W + v_{\perp}\)，则 \(\langle P_W u, v \rangle = \langle u_W, v_W + v_{\perp} \rangle = \langle u_W, v_W \rangle = \langle u_W + u_{\perp}, v_W \rangle = \langle u, P_W v \rangle\)。
\item
  将 \(\{e_1, \ldots, e_k\}\) 扩充为 \(V\) 的规范正交基 \(\{e_1, \ldots, e_n\}\)。设 \(v = \sum_{i=1}^{n} \alpha_i e_i\)（\(\alpha_i = \langle v, e_i \rangle\)）。则 \(P_W(v) = \sum_{i=1}^{k} \alpha_i e_i = \sum_{i=1}^{k} \langle v, e_i \rangle e_i\)。\(\blacksquare\)
\end{enumerate}

\hrulefill

\subsection{36.5 伴随算子}\label{ux4f34ux968fux7b97ux5b50}

\textbf{定义 36.5.1（伴随算子）}：设 \(V\) 为有限维内积空间，\(T \in \mathcal{L}(V)\)。\(T\) 的\textbf{伴随算子} \(T^* \in \mathcal{L}(V)\) 由下式唯一确定：

\[
\langle T(u), v \rangle = \langle u, T^*(v) \rangle, \quad \forall u, v \in V
\]

\textbf{定理 36.5.1（伴随算子的存在唯一性）}：\(T^*\) 存在且唯一。

\textbf{证明}：固定 \(v \in V\)，考虑线性泛函 \(f_v(u) = \langle T(u), v \rangle\)。由 Riesz 表示定理（在有限维情形下，\(f_v\) 可表示为与某个向量的内积），存在唯一的 \(w \in V\) 使得 \(f_v(u) = \langle u, w \rangle\)。定义 \(T^*(v) = w\)。可验证 \(T^*\) 是线性的。\(\blacksquare\)

\textbf{定理 36.5.2（伴随算子的性质）}：

\begin{enumerate}
\def\labelenumi{(\roman{enumi})}
\tightlist
\item
  \((T^*)^* = T\)
\item
  \((S + T)^* = S^* + T^*\)
\item
  \((\lambda T)^* = \lambda T^*\)
\item
  \((ST)^* = T^* S^*\)
\item
  若 \(T\) 可逆，则 \((T^{-1})^* = (T^*)^{-1}\)
\item
  在规范正交基下，\([T^*] = [T]^T\)（伴随算子对应矩阵的转置）
\end{enumerate}

\textbf{证明}：(i)-(v) 由伴随算子的定义和内积的性质直接验证。

\begin{enumerate}
\def\labelenumi{(\roman{enumi})}
\setcounter{enumi}{5}
\tightlist
\item
  设 \(B = \{e_1, \ldots, e_n\}\) 为 \(V\) 的规范正交基，\(A = [T]_B\)，\(B = [T^*]_B\)。则 \(T(e_j) = \sum_i a_{ij} e_i\)，\(T^*(e_j) = \sum_i b_{ij} e_i\)。
\end{enumerate}

由伴随算子的定义：\(\langle T(e_j), e_i \rangle = \langle e_j, T^*(e_i) \rangle\)。左边 \(= \langle \sum_k a_{kj} e_k, e_i \rangle = a_{ij}\)。右边 \(= \langle e_j, \sum_k b_{ki} e_k \rangle = b_{ji}\)。故 \(a_{ij} = b_{ji}\)，即 \(B = A^T\)。\(\blacksquare\)

\textbf{定义 36.5.2（自伴算子与正交算子）}：

\begin{enumerate}
\def\labelenumi{\arabic{enumi}.}
\tightlist
\item
  若 \(T^* = T\)，则称 \(T\) 为\textbf{自伴算子}（或\textbf{对称算子}）。在规范正交基下，其矩阵为对称矩阵（\(A^T = A\)）。
\item
  若 \(T^* T = T T^* = I\)，则称 \(T\) 为\textbf{正交算子}。在规范正交基下，其矩阵为正交矩阵（\(A^T A = I_n\)）。
\end{enumerate}

\hrulefill

\subsection{36.6 谱定理}\label{ux8c31ux5b9aux7406}

\textbf{定理 36.6.1（实对称矩阵的谱定理）}：设 \(A \in M_{n \times n}(\mathbb{R})\) 为对称矩阵（\(A^T = A\)）。则

\begin{enumerate}
\def\labelenumi{(\roman{enumi})}
\tightlist
\item
  \(A\) 的所有特征值均为实数。
\item
  \(A\) 可对角化，且存在正交矩阵 \(Q\)（\(Q^T Q = I_n\)）使得
\end{enumerate}

\[
Q^T A Q = \operatorname{diag}(\lambda_1, \lambda_2, \ldots, \lambda_n)
\]

其中 \(\lambda_1, \ldots, \lambda_n\) 为 \(A\) 的特征值。等价地，\(\mathbb{R}^n\) 有一组由 \(A\) 的特征向量组成的规范正交基。

\textbf{证明}：对 \(n\) 归纳。

\textbf{基础} \(n = 1\)：平凡。

\textbf{归纳}：假设结论对 \(n-1\) 阶对称矩阵成立。

\textbf{第 1 步}：\(A\) 有实特征值。考虑 \(\mathbb{C}^n\) 上的标准 Hermite 内积。\(A\) 作为复矩阵是 Hermite 矩阵（\(A^* = A^T = A\)）。\(A\) 至少有一个特征值 \(\lambda_1 \in \mathbb{C}\) 和对应的特征向量 \(v_1 \neq \mathbf{0}\)。由 \(\lambda_1 \|v_1\|^2 = \langle A v_1, v_1 \rangle = \langle v_1, A v_1 \rangle = \overline{\lambda_1} \|v_1\|^2\)，得 \(\lambda_1 = \overline{\lambda_1}\)，故 \(\lambda_1 \in \mathbb{R}\)。可取 \(v_1\) 为实向量，并单位化使 \(\|v_1\| = 1\)。

\textbf{第 2 步}：设 \(W = \operatorname{span}\{v_1\}\)。则 \(W^{\perp}\) 是 \(A\)-不变的：对任意 \(w \in W^{\perp}\)，\(\langle A w, v_1 \rangle = \langle w, A v_1 \rangle = \langle w, \lambda_1 v_1 \rangle = \lambda_1 \langle w, v_1 \rangle = 0\)，故 \(A w \in W^{\perp}\)。

\textbf{第 3 步}：\(A|_{W^{\perp}}\) 是 \(W^{\perp}\) 上的对称算子（因为 \(W^{\perp}\) 是 \(A\)-不变的，且 \(A\) 的对称性继承）。取 \(W^{\perp}\) 的规范正交基，\(A|_{W^{\perp}}\) 在此基下的矩阵 \(A'\) 为 \((n-1) \times (n-1)\) 对称矩阵。由归纳假设，存在 \((n-1) \times (n-1)\) 正交矩阵 \(Q'\) 使得 \((Q')^T A' Q' = \operatorname{diag}(\lambda_2, \ldots, \lambda_n)\)。

将 \(v_1\) 与 \(W^{\perp}\) 的规范正交基合并，得到 \(\mathbb{R}^n\) 的规范正交基。对应的正交矩阵 \(Q\) 满足 \(Q^T A Q = \operatorname{diag}(\lambda_1, \ldots, \lambda_n)\)。\(\blacksquare\)

\textbf{定理 36.6.2（谱定理的算子形式）}：设 \(V\) 为有限维实内积空间，\(T \in \mathcal{L}(V)\) 为自伴算子。则 \(V\) 存在一组由 \(T\) 的特征向量组成的规范正交基。等价地，\(T\) 可对角化且特征向量可选取为正交的。

\textbf{证明}：取 \(V\) 的任意规范正交基，\(T\) 在此基下的矩阵 \(A\) 是对称矩阵。由定理 36.6.1，\(A\) 可正交对角化。\(Q\) 的列向量即为 \(T\) 的特征向量，且构成规范正交基。\(\blacksquare\)

\textbf{定理 36.6.3（奇异值分解，陈述）}：设 \(A \in M_{m \times n}(\mathbb{R})\)。则存在正交矩阵 \(U \in M_{m \times m}(\mathbb{R})\)，\(V \in M_{n \times n}(\mathbb{R})\) 和对角矩阵 \(\Sigma = \operatorname{diag}(\sigma_1, \ldots, \sigma_r, 0, \ldots, 0)\)（\(\sigma_1 \geq \cdots \geq \sigma_r > 0\)）使得

\[
A = U \Sigma V^T
\]

其中 \(\sigma_i\) 称为 \(A\) 的\textbf{奇异值}，等于 \(A^T A\) 的正特征值的平方根。

\textbf{评注 36.6.1}：奇异值分解的证明基于谱定理应用于 \(A^T A\)，详细推导从略。\(\blacksquare\)

\hrulefill

\hrulefill

\hrulefill

\hrulefill

\hrulefill

\hrulefill

\section{第110章：二次型与标准形}\label{ux7b2c110ux7ae0ux4e8cux6b21ux578bux4e0eux6807ux51c6ux5f62}

二次型是线性代数与多重线性代数交汇处的核心概念，它将齐次二次多项式的研究纳入线性代数的框架，揭示了矩阵、对称双线性型与二次型之间的深刻联系。二次型理论的历史源远流长：Gauss 在 1801 年《算术研究》中系统地研究了整系数二元二次型，奠定了二次型数论的基础；Jacobi 和 Sylvester 在 19 世纪中叶发展了实二次型的惯性定律和分类理论；而 Hilbert 在 1890 年代的工作将二次型与代数数论、代数几何紧密联系起来。

二次型理论的核心问题是标准形的存在性与唯一性。给定实二次型 \(Q(\mathbf{x}) = \mathbf{x}^T A \mathbf{x}\)（\(A\) 为对称矩阵），通过可逆线性变换 \(\mathbf{x} = P\mathbf{y}\)，\(Q\) 可化为标准形 \(\lambda_1 y_1^2 + \cdots + \lambda_n y_n^2\)。这里 \(\lambda_i\) 是 \(A\) 的特征值（由谱定理保证），而变换矩阵 \(P\) 可取为正交矩阵。Sylvester 惯性定律进一步断言：虽然标准形中的系数不唯一（取决于所用的变换），但正系数的个数 \(p\)（正惯性指数）和负系数的个数 \(q\)（负惯性指数）是不变的------它们是二次型的内在不变量。符号差 \(p - q\) 在微分拓扑中尤为重要：Morse 理论中临界点的 Morse 指数正是 Hessian 二次型的负惯性指数。

二次型的定性分类------正定、负定、半正定、半负定、不定------不仅在纯粹数学中至关重要，也在应用数学中频繁出现。在优化理论中，Hessian 矩阵的正定性保证了凸函数的极小值点；在微分几何中，Riemann 度量 \(g_{ij}\) 在每点的正定性是黎曼几何的基本要求，而 Lorentz 度量（符号差为 \(1\)）则是广义相对论中时空的数学基础。Sylvester 准则（顺序主子式的符号判定正定性）和特征值判定法提供了判定二次型定性的有效工具。本章将系统发展二次型的代数理论，包括合同变换、惯性定理、正定性判定，以及二次型与内积空间的极化恒等式联系。

\textbf{定义 37.1.1（二次型）}：设 \(V\) 为 \(\mathbb{R}\) 上的 \(n\) 维向量空间。函数 \(Q : V \to \mathbb{R}\) 称为\textbf{二次型}，若存在对称双线性型 \(B : V \times V \to \mathbb{R}\)（即 \(B(u, v) = B(v, u)\) 且对每个变元线性）使得

\[
Q(v) = B(v, v), \quad \forall v \in V
\]

\textbf{定义 37.1.2（二次型的矩阵表示）}：在基 \(B = \{e_1, \ldots, e_n\}\) 下，设 \(v = \sum x_i e_i\)。则

\[
Q(v) = \sum_{i,j=1}^{n} a_{ij} x_i x_j = x^T A x
\]

其中 \(A = (a_{ij})\) 为对称矩阵（\(a_{ij} = B(e_i, e_j)\)），称为二次型 \(Q\) 关于基 \(B\) 的\textbf{矩阵}。

\textbf{定理 37.1.1（基变换下二次型矩阵的变化）}：设 \(P\) 为从基 \(B\) 到 \(B'\) 的基变换矩阵。则 \(Q\) 关于 \(B'\) 的矩阵 \(A'\) 满足

\[
A' = P^T A P
\]

\textbf{证明}：设 \(x\) 为关于 \(B\) 的坐标，\(x'\) 为关于 \(B'\) 的坐标，则 \(x = P x'\)。故 \(Q(v) = x^T A x = (P x')^T A (P x') = (x')^T (P^T A P) x'\)。\(\blacksquare\)

\textbf{定义 37.1.3（合同）}：对称矩阵 \(A, B \in M_{n \times n}(\mathbb{R})\) 称为\textbf{合同}的，若存在可逆矩阵 \(P\) 使得 \(B = P^T A P\)。

\hrulefill

\subsection{37.2 合同变换与标准形}\label{ux5408ux540cux53d8ux6362ux4e0eux6807ux51c6ux5f62}

\textbf{定理 37.2.1（二次型的标准形）}：设 \(Q\) 为 \(\mathbb{R}^n\) 上的二次型，矩阵为 \(A\)（对称）。则存在可逆矩阵 \(P\)，使得经坐标变换 \(x = P y\) 后，

\[
Q(v) = \lambda_1 y_1^2 + \lambda_2 y_2^2 + \cdots + \lambda_n y_n^2
\]

其中 \(\lambda_i \in \mathbb{R}\)。这等价于任何实对称矩阵合同于对角矩阵。

\textbf{证明}：由谱定理（定理 36.6.1），存在正交矩阵 \(Q\) 使得 \(Q^T A Q = \operatorname{diag}(\lambda_1, \ldots, \lambda_n)\)。取 \(P = Q\)，则 \(P^T A P = \operatorname{diag}(\lambda_1, \ldots, \lambda_n)\)。坐标变换 \(x = P y\) 给出 \(Q(v) = \sum \lambda_i y_i^2\)。\(\blacksquare\)

\textbf{注记 37.2.1}：上述证明利用了正交对角化，\(\lambda_i\) 恰好是 \(A\) 的特征值。也可以用配方法（Lagrange 配方法）直接构造，无需谱定理。

\hrulefill

\subsection{37.3 惯性定理}\label{ux60efux6027ux5b9aux7406}

\textbf{定理 37.3.1（Sylvester 惯性定理）}：设实二次型 \(Q\) 在两组基下分别化为标准形：

\[
Q(v) = \lambda_1 y_1^2 + \cdots + \lambda_p y_p^2 - \lambda_{p+1} y_{p+1}^2 - \cdots - \lambda_{p+q} y_{p+q}^2 + 0 \cdot y_{p+q+1}^2 + \cdots + 0 \cdot y_n^2
\]

\[
Q(v) = \mu_1 z_1^2 + \cdots + \mu_{p'} z_{p'}^2 - \mu_{p'+1} z_{p'+1}^2 - \cdots - \mu_{p'+q'} z_{p'+q'}^2 + 0
\]

其中 \(\lambda_i, \mu_j > 0\)。则 \(p = p'\)，\(q = q'\)。

\textbf{证明}：\(p\) 是 \(Q\) 在其上取正值的子空间的最大维数，\(q\) 是 \(Q\) 在其上取负值的子空间的最大维数。这些是不依赖于基选择的不变量。

具体地，设 \(V_+\) 为 \(Q(v) \geq 0\) 的子空间中的最大维数者（\(Q\) 限制在 \(V_+\) 上正定）。可以证明 \(\dim V_+ = p\)。同样，\(Q\) 限制在 \(V_-\) 上负定的最大维数为 \(q\)。由于 \(p, q\) 由 \(Q\) 本身决定，它们在两组基下相同。\(\blacksquare\)

\textbf{定义 37.3.1（惯性指数与符号差）}：

\begin{itemize}
\tightlist
\item
  \(p\) 称为 \(Q\) 的\textbf{正惯性指数}
\item
  \(q\) 称为 \(Q\) 的\textbf{负惯性指数}
\item
  \(p + q = \operatorname{rank} A\) 称为 \(Q\) 的\textbf{秩}
\item
  \(p - q\) 称为 \(Q\) 的\textbf{符号差}
\end{itemize}

\textbf{推论 37.3.2}：两个实二次型等价（可通过可逆线性变换互相转化）当且仅当它们具有相同的正惯性指数和负惯性指数。

\hrulefill

\subsection{37.4 Sylvester 准则}\label{sylvester-ux51c6ux5219}

\textbf{定义 37.4.1（顺序主子式）}：设 \(A = (a_{ij}) \in M_{n \times n}(\mathbb{R})\) 为对称矩阵。对 \(k = 1, \ldots, n\)，\(A\) 的第 \(k\) 个\textbf{顺序主子式}为

\[
\Delta_k = \det \begin{pmatrix}
a_{11} & \cdots & a_{1k} \\
\vdots & \ddots & \vdots \\
a_{k1} & \cdots & a_{kk}
\end{pmatrix}
\]

\textbf{定理 37.4.1（Sylvester 准则）}：实对称矩阵 \(A\) 正定（即 \(x^T A x > 0\) 对所有 \(x \neq \mathbf{0}\)）当且仅当其所有顺序主子式 \(\Delta_k > 0\)（\(k = 1, \ldots, n\)）。

\textbf{证明}：

（\(\Rightarrow\)）设 \(A\) 正定。对任意 \(k\)，考虑 \(A\) 的左上 \(k \times k\) 子矩阵 \(A_k\)。对任意非零 \(y \in \mathbb{R}^k\)，令 \(x = (y, 0, \ldots, 0)^T \in \mathbb{R}^n\)，则 \(y^T A_k y = x^T A x > 0\)，故 \(A_k\) 正定。正定矩阵的特征值全正，故 \(\det A_k = \prod_{i=1}^{k} \lambda_i > 0\)。

（\(\Leftarrow\)）对 \(n\) 归纳。\(n = 1\) 时，\(\Delta_1 = a_{11} > 0\)，故 \(A\) 正定。

假设结论对 \(n-1\) 成立。设 \(\Delta_k > 0\)（\(k = 1, \ldots, n\)）。由归纳假设，\(A_{n-1}\) 正定，故存在可逆矩阵 \(P_{n-1}\) 使得 \(P_{n-1}^T A_{n-1} P_{n-1} = I_{n-1}\)。

将 \(A\) 写为分块形式 \(A = \begin{pmatrix} A_{n-1} & b \\ b^T & a_{nn} \end{pmatrix}\)。令 \(P = \begin{pmatrix} P_{n-1} & -A_{n-1}^{-1} b \\ 0 & 1 \end{pmatrix}\)，则 \[
P^T A P = \begin{pmatrix} I_{n-1} & 0 \\ 0 & a_{nn} - b^T A_{n-1}^{-1} b \end{pmatrix}
\] 其行列式为 \(\Delta_n / \Delta_{n-1} > 0\)（因为 \(\det(P^T A P) = \det P^2 \det A\)）。故 \(a_{nn} - b^T A_{n-1}^{-1} b > 0\)，从而 \(P^T A P\) 正定，\(A\) 正定。\(\blacksquare\)

\hrulefill

\subsection{37.5 正定二次型}\label{ux6b63ux5b9aux4e8cux6b21ux578b}

\textbf{定义 37.5.1（二次型的定性）}：设 \(Q\) 为实二次型，矩阵为 \(A\)。

\begin{enumerate}
\def\labelenumi{\arabic{enumi}.}
\tightlist
\item
  若 \(Q(v) > 0\)（\(\forall v \neq \mathbf{0}\)），称 \(Q\)（及 \(A\)）\textbf{正定}
\item
  若 \(Q(v) \geq 0\)（\(\forall v\)），称 \(Q\)（及 \(A\)）\textbf{半正定}
\item
  若 \(Q(v) < 0\)（\(\forall v \neq \mathbf{0}\)），称 \(Q\)（及 \(A\)）\textbf{负定}
\item
  若 \(Q(v) \leq 0\)（\(\forall v\)），称 \(Q\)（及 \(A\)）\textbf{半负定}
\item
  若 \(Q\) 既取正值又取负值，称 \(Q\)（及 \(A\)）\textbf{不定}
\end{enumerate}

\textbf{定理 37.5.1（正定性的特征值判定）}：对称矩阵 \(A\) 正定当且仅当 \(A\) 的所有特征值均为正数。

\textbf{证明}：由谱定理，\(A = Q \operatorname{diag}(\lambda_1, \ldots, \lambda_n) Q^T\)。对任意 \(x \neq \mathbf{0}\)，令 \(y = Q^T x \neq \mathbf{0}\)，则 \(x^T A x = y^T \operatorname{diag}(\lambda_i) y = \sum \lambda_i y_i^2\)。此式对所有 \(y \neq \mathbf{0}\) 为正当且仅当所有 \(\lambda_i > 0\)。\(\blacksquare\)

\textbf{定理 37.5.2（二次型与内积）}：\(V\) 上的正定二次型 \(Q\) 通过极化恒等式

\[
\langle u, v \rangle_Q = \frac{1}{2}\left(Q(u + v) - Q(u) - Q(v)\right)
\]

定义了 \(V\) 上的一个内积。反之，任何内积 \(\langle\cdot,\cdot\rangle\) 通过 \(Q(v) = \langle v, v \rangle\) 给出一个正定二次型。

\textbf{证明}：直接验证内积的四条公理。正定性由 \(Q\) 的正定性保证，对称性由二次型的定义中 \(B\) 的对称性保证，线性性由极化恒等式的构造保证。\(\blacksquare\)

\emph{卷六：线性代数终。}

\emph{卷六：线性代数终。}

\emph{卷六：线性代数终。}

\emph{卷七：线性代数终。}

\newpage

\chapter{05-代数/02-抽象代数一/01-群论基础.md}\label{ux4ee3ux657002-ux62bdux8c61ux4ee3ux6570ux4e0001-ux7fa4ux8bbaux57faux7840.md}

\section{第111章：群论基础}\label{ux7b2c111ux7ae0ux7fa4ux8bbaux57faux7840}

群是现代数学中最基本的代数结构，它公理化地刻画了''对称性''这一贯穿数学和物理的核心概念。群论的历史可追溯至 Evariste Galois 在 1832 年关于代数方程根式可解性的开创性工作------他在决斗前夜写下的著名遗书中首次系统地将群的概念用于研究多项式的对称性。此后，Cayley（1854 年给出抽象群的定义）、Jordan（1870 年《置换与代数方程论》）、Klein（1872 年 Erlangen 纲领------几何学是变换群的不变量理论）和 Lie（连续变换群）将群论发展为数学的通用语言。

群论的基本思想可概括为：研究一个集合 \(G\) 连同其上的一个满足结合律、具有幺元和逆元的二元运算。这四条公理看似简单，却从中衍生出极为丰富的结构理论。Lagrange 定理（\(|G| = |H| \cdot [G:H]\)）揭示了子群阶与群阶之间的整除关系，是有限群论的第一基本定理。Cayley 定理（任何群同构于某个置换群的子群）则表明群论本质上就是对称性的研究，为群作用的引入提供了理论依据。

本章在卷二群论初步引入的基础上，系统深化群论的核心概念：正规子群与商群，它们使得群可以像整数一样进行''除法''------商群 \(G/N\) 是 \(G\) 模正规子群 \(N\) 的等价类构成的群；同态与同构定理，它们揭示了群之间的结构保持映射，三个同构定理（类似于 Noether 同构定理）架起了群及其商群之间的桥梁；群的直积与半直积，它们提供了从已知群构造新群的基本方法，其中半直积特别适用于描述非交换群的结构（如 Euclid 群 \(E(n) = \mathbb{R}^n \rtimes O(n)\)）。这些概念构成了群论的结构骨架，为后续的群作用、Sylow 定理和有限单群分类铺平道路。

卷二第六章已引入群的基本定义。本节简要回顾并深化。

\textbf{定义 39.1.1（群）}：群是一个集合 \(G\) 配备二元运算 \(\cdot\)，满足：

\begin{enumerate}
\def\labelenumi{\arabic{enumi}.}
\tightlist
\item
  \textbf{封闭性}：\(\forall a, b \in G\)，\(a \cdot b \in G\)
\item
  \textbf{结合律}：\((a \cdot b) \cdot c = a \cdot (b \cdot c)\)
\item
  \textbf{幺元}：\(\exists e \in G\)，\(\forall a \in G\)，\(e \cdot a = a \cdot e = a\)
\item
  \textbf{逆元}：\(\forall a \in G\)，\(\exists a^{-1} \in G\)，\(a \cdot a^{-1} = a^{-1} \cdot a = e\)
\end{enumerate}

若进一步满足交换律，则称为\textbf{阿贝尔群}。

\textbf{定义 39.1.2（子群）}：\(H \subseteq G\) 称为 \(G\) 的\textbf{子群}，记为 \(H \leq G\)，若 \(H\) 在 \(G\) 的运算下自身构成群。

\textbf{定理 39.1.1（子群判定）}：\(H \subseteq G\) 是子群当且仅当

\begin{enumerate}
\def\labelenumi{(\roman{enumi})}
\tightlist
\item
  \(H \neq \varnothing\)
\item
  \(\forall a, b \in H\)，\(ab^{-1} \in H\)
\end{enumerate}

\textbf{证明}：由群的定义直接验证。取 \(a \in H\)，则 \(e = aa^{-1} \in H\)。对 \(a \in H\)，\(a^{-1} = e a^{-1} \in H\)。对 \(a, b \in H\)，\(ab = a(b^{-1})^{-1} \in H\)。\(\blacksquare\)

\textbf{定义 39.1.3（阶）}：群 \(G\) 的元素个数称为 \(G\) 的\textbf{阶}，记为 \(|G|\)。元素 \(a \in G\) 的\textbf{阶}是使得 \(a^n = e\) 的最小正整数 \(n\)（若不存在则为 \(\infty\)），记为 \(|a|\)。

\hrulefill

\subsection{39.2 陪集与 Lagrange 定理}\label{ux966aux96c6ux4e0e-lagrange-ux5b9aux7406}

\textbf{定义 39.2.1（陪集）}：设 \(H \leq G\)，\(a \in G\)。

\begin{itemize}
\tightlist
\item
  \(aH = \{ah \mid h \in H\}\) 称为 \(H\) 的\textbf{左陪集}
\item
  \(Ha = \{ha \mid h \in H\}\) 称为 \(H\) 的\textbf{右陪集}
\end{itemize}

\textbf{定理 39.2.1（陪集的基本性质）}：设 \(H \leq G\)。

\begin{enumerate}
\def\labelenumi{(\roman{enumi})}
\tightlist
\item
  \(a \in aH\)，故 \(G = \bigcup_{a \in G} aH\)
\item
  对任意 \(a, b \in G\)，\(aH = bH\) 或 \(aH \cap bH = \varnothing\)
\item
  \(|aH| = |H|\)（映射 \(h \mapsto ah\) 是双射）
\end{enumerate}

\textbf{证明}：

\begin{enumerate}
\def\labelenumi{(\roman{enumi})}
\item
  \(a = ae \in aH\)（\(e \in H\)）。
\item
  若 \(aH \cap bH \neq \varnothing\)，设 \(x \in aH \cap bH\)，则 \(x = ah_1 = bh_2\)。于是 \(a = bh_2 h_1^{-1} \in bH\)，故 \(aH \subseteq bH\)。同理 \(bH \subseteq aH\)，故 \(aH = bH\)。
\item
  映射 \(f : H \to aH\)，\(f(h) = ah\) 是双射：\(ah_1 = ah_2 \implies h_1 = h_2\)（左消去）；对任意 \(ah \in aH\)，\(f(h) = ah\)。\(\blacksquare\)
\end{enumerate}

\textbf{定理 39.2.2（Lagrange 定理）}：设 \(G\) 为有限群，\(H \leq G\)。则 \(|H|\) 整除 \(|G|\)，且

\[
|G| = |H| \cdot [G : H]
\]

其中 \([G : H]\) 为 \(H\) 在 \(G\) 中的\textbf{指数}（即不同陪集的个数）。

\textbf{证明}：由定理 39.2.1，\(G\) 被划分为 \([G : H]\) 个互不相交的左陪集，每个陪集的大小为 \(|H|\)。因此 \(|G| = [G : H] \cdot |H|\)。\(\blacksquare\)

\textbf{推论 39.2.3}：设 \(G\) 为有限群，\(a \in G\)。则 \(|a|\) 整除 \(|G|\)。特别地，\(a^{|G|} = e\)。

\textbf{证明}：考虑由 \(a\) 生成的循环子群 \(\langle a \rangle = \{a^k \mid k \in \mathbb{Z}\}\)。\(|\langle a \rangle| = |a|\)。由 Lagrange 定理，\(|a| \mid |G|\)。因此 \(a^{|G|} = (a^{|a|})^{|G|/|a|} = e\)。\(\blacksquare\)

\textbf{推论 39.2.4（Fermat-Euler 定理的推广）}：若 \(|G| = p\)（素数），则 \(G\) 是循环群，且任何非幺元都是生成元。

\textbf{证明}：取 \(a \in G \setminus \{e\}\)。\(|a|\) 整除 \(p\) 且 \(|a| > 1\)，故 \(|a| = p\)。因此 \(\langle a \rangle = G\)。\(\blacksquare\)

\hrulefill

\subsection{39.3 正规子群与商群}\label{ux6b63ux89c4ux5b50ux7fa4ux4e0eux5546ux7fa4}

\textbf{定义 39.3.1（正规子群）}：\(N \leq G\) 称为 \(G\) 的\textbf{正规子群}，记为 \(N \trianglelefteq G\)，若对任意 \(g \in G\)，\(gNg^{-1} = N\)（等价地，\(gN = Ng\)）。

\textbf{定理 39.3.1（正规子群的判定）}：\(N \trianglelefteq G\) 当且仅当 \(\forall g \in G, \forall n \in N\)，\(gng^{-1} \in N\)。

\textbf{证明}：\(gNg^{-1} \subseteq N\) 对所有 \(g\) 成立则 \(gNg^{-1} = N\)（用 \(g^{-1}\) 替换 \(g\) 得 \(g^{-1}Ng \subseteq N\)，两边左乘 \(g\) 右乘 \(g^{-1}\) 得 \(N \subseteq gNg^{-1}\)）。\(\blacksquare\)

\textbf{定义 39.3.2（商群）}：设 \(N \trianglelefteq G\)。在陪集的集合 \(G/N = \{gN \mid g \in G\}\) 上定义运算 \((aN)(bN) = (ab)N\)。\((G/N, \cdot)\) 构成群，称为 \(G\) 模 \(N\) 的\textbf{商群}。

\textbf{定理 39.3.2（商群中运算的良定义性）}：上述运算在 \(N \trianglelefteq G\) 时是良定义的。

\textbf{证明}：需证若 \(aN = a'N\) 且 \(bN = b'N\)，则 \((ab)N = (a'b')N\)。

由 \(aN = a'N\) 得 \(a' = an_1\)（\(n_1 \in N\)）。由 \(bN = b'N\) 得 \(b' = bn_2\)（\(n_2 \in N\)）。则 \(a'b' = an_1 bn_2 = ab(b^{-1}n_1 b)n_2\)。由于 \(N \trianglelefteq G\)，\(b^{-1}n_1 b \in N\)，故 \((b^{-1}n_1 b)n_2 \in N\)。因此 \(a'b' \in (ab)N\)，即 \((a'b')N = (ab)N\)。\(\blacksquare\)

\textbf{定理 39.3.3}：若 \(G\) 是阿贝尔群，则 \(G\) 的任何子群都是正规子群。

\textbf{证明}：\(gng^{-1} = gg^{-1}n = n \in N\)。\(\blacksquare\)

\hrulefill

\subsection{39.4 群同态与同构定理}\label{ux7fa4ux540cux6001ux4e0eux540cux6784ux5b9aux7406}

\textbf{定义 39.4.1（群同态）}：设 \(G, H\) 为群。映射 \(\varphi : G \to H\) 若满足 \(\varphi(ab) = \varphi(a)\varphi(b)\)（\(\forall a, b \in G\)），则称 \(\varphi\) 为\textbf{群同态}。

\textbf{定理 39.4.1（同态的基本性质）}：设 \(\varphi : G \to H\) 为群同态。

\begin{enumerate}
\def\labelenumi{(\roman{enumi})}
\tightlist
\item
  \(\varphi(e_G) = e_H\)
\item
  \(\varphi(a^{-1}) = \varphi(a)^{-1}\)
\item
  \(\ker \varphi = \{g \in G \mid \varphi(g) = e_H\}\) 是 \(G\) 的正规子群
\item
  \(\operatorname{im} \varphi = \{\varphi(g) \mid g \in G\}\) 是 \(H\) 的子群
\item
  \(\varphi\) 是单射 \(\iff\) \(\ker \varphi = \{e_G\}\)
\end{enumerate}

\textbf{证明}：

\begin{enumerate}
\def\labelenumi{(\roman{enumi})}
\item
  \(\varphi(e_G) = \varphi(e_G e_G) = \varphi(e_G)\varphi(e_G)\)，两边左乘 \(\varphi(e_G)^{-1}\) 得 \(\varphi(e_G) = e_H\)。
\item
  \(\varphi(a)\varphi(a^{-1}) = \varphi(aa^{-1}) = \varphi(e_G) = e_H\)，同理 \(\varphi(a^{-1})\varphi(a) = e_H\)。由逆元唯一性得证。
\item
  \(\ker \varphi\) 是子群（由 (i)(ii) 验证）。对 \(g \in G\)，\(k \in \ker \varphi\)，\(\varphi(gkg^{-1}) = \varphi(g)\varphi(k)\varphi(g)^{-1} = \varphi(g)e_H \varphi(g)^{-1} = e_H\)，故 \(gkg^{-1} \in \ker \varphi\)。因此 \(\ker \varphi \trianglelefteq G\)。
\item
  直接验证。(v) 与线性映射的判定类似（定理 33.2.2）。\(\blacksquare\)
\end{enumerate}

\textbf{定理 39.4.2（第一同构定理）}：设 \(\varphi : G \to H\) 为群同态。则

\[
G / \ker \varphi \cong \operatorname{im} \varphi
\]

同构映射为 \(\overline{\varphi}(g \ker \varphi) = \varphi(g)\)。

\textbf{证明}：定义 \(\overline{\varphi} : G/\ker \varphi \to \operatorname{im} \varphi\) 为 \(\overline{\varphi}(gK) = \varphi(g)\)（\(K = \ker \varphi\)）。

良定义：若 \(gK = g'K\)，则 \(g' = gk\)（\(k \in K\)）。\(\varphi(g') = \varphi(gk) = \varphi(g)\varphi(k) = \varphi(g)\)。

同态：\(\overline{\varphi}((gK)(hK)) = \overline{\varphi}(ghK) = \varphi(gh) = \varphi(g)\varphi(h) = \overline{\varphi}(gK)\overline{\varphi}(hK)\)。

单射：若 \(\overline{\varphi}(gK) = e_H\)，则 \(\varphi(g) = e_H\)，故 \(g \in K\)，\(gK = K\)（商群的幺元）。

满射：由 \(\operatorname{im} \varphi\) 的定义显然。\(\blacksquare\)

\textbf{定理 39.4.3（第二同构定理）}：设 \(H \leq G\)，\(N \trianglelefteq G\)。则 \(HN = \{hn \mid h \in H, n \in N\}\) 是 \(G\) 的子群，\(H \cap N \trianglelefteq H\)，且

\[
H / (H \cap N) \cong HN / N
\]

\textbf{证明}：\(HN\) 是子群：\(e \in HN\)；\((h_1 n_1)(h_2 n_2) = h_1 h_2 (h_2^{-1} n_1 h_2) n_2 \in HN\)（因为 \(h_2^{-1} n_1 h_2 \in N\)）；\((hn)^{-1} = n^{-1}h^{-1} = h^{-1}(h n^{-1} h^{-1}) \in HN\)。

\(H \cap N \trianglelefteq H\)：对 \(h \in H\)，\(x \in H \cap N\)，\(hxh^{-1} \in H\)（\(h, x, h^{-1} \in H\)）且 \(hxh^{-1} \in N\)（\(N \trianglelefteq G\)），故 \(hxh^{-1} \in H \cap N\)。

定义 \(\varphi : H \to HN/N\) 为 \(\varphi(h) = hN\)。\(\varphi\) 是满同态。\(\ker \varphi = \{h \in H \mid hN = N\} = \{h \in H \mid h \in N\} = H \cap N\)。由第一同构定理，\(H/(H \cap N) \cong HN/N\)。\(\blacksquare\)

\textbf{定理 39.4.4（第三同构定理）}：设 \(N \trianglelefteq G\)，\(K \trianglelefteq G\) 且 \(N \subseteq K\)。则 \(K/N \trianglelefteq G/N\)，且

\[
(G/N) / (K/N) \cong G/K
\]

\textbf{证明}：定义 \(\varphi : G/N \to G/K\) 为 \(\varphi(gN) = gK\)。

良定义：若 \(gN = g'N\)，则 \(g' = gn\)（\(n \in N \subseteq K\)），故 \(g'K = gK\)。

同态：\(\varphi((gN)(hN)) = \varphi(ghN) = ghK = (gK)(hK) = \varphi(gN)\varphi(hN)\)。

满射显然。\(\ker \varphi = \{gN \mid gK = K\} = \{gN \mid g \in K\} = K/N\)。

由第一同构定理得证。\(\blacksquare\)

\hrulefill

\subsection{39.5 循环群}\label{ux5faaux73afux7fa4}

\textbf{定义 39.5.1（循环群）}：群 \(G\) 称为\textbf{循环群}，若存在 \(a \in G\) 使得 \(G = \langle a \rangle = \{a^n \mid n \in \mathbb{Z}\}\)。

\textbf{定理 39.5.1（循环群的分类）}：

\begin{enumerate}
\def\labelenumi{(\roman{enumi})}
\tightlist
\item
  若 \(|G| = \infty\)，则 \(G \cong \mathbb{Z}\)
\item
  若 \(|G| = n\)，则 \(G \cong \mathbb{Z}_n\)
\end{enumerate}

\textbf{证明}：定义 \(\varphi : \mathbb{Z} \to G\) 为 \(\varphi(k) = a^k\)。\(\varphi\) 是满同态。\(\ker \varphi\) 是 \(\mathbb{Z}\) 的子群，故 \(\ker \varphi = m\mathbb{Z}\)（\(m \geq 0\)）。

若 \(m = 0\)，则 \(\varphi\) 是单射，故 \(G \cong \mathbb{Z}\)。

若 \(m > 0\)，由第一同构定理 \(G \cong \mathbb{Z}/m\mathbb{Z} = \mathbb{Z}_m\)。此时 \(|G| = m = n\)。\(\blacksquare\)

\textbf{定理 39.5.2（循环群的子群）}：循环群的任何子群也是循环群。若 \(G = \langle a \rangle\) 为 \(n\) 阶循环群，则对 \(n\) 的每个正因子 \(d\)，\(G\) 恰有一个 \(d\) 阶子群，即 \(\langle a^{n/d} \rangle\)。

\textbf{证明}：设 \(H \leq G = \langle a \rangle\)。令 \(m = \min\{k > 0 \mid a^k \in H\}\)。可以证明 \(H = \langle a^m \rangle\)（用带余除法）。

对 \(n\) 阶循环群，\(d \mid n\) 时，\(|a^{n/d}| = n / \gcd(n, n/d) = n / (n/d) = d\)。唯一性：若 \(H\) 为 \(d\) 阶子群，则 \(H = \langle a^k \rangle\) 且 \(|a^k| = n / \gcd(n, k) = d\)，故 \(\gcd(n, k) = n/d\)。因此 \(k\) 是 \(n/d\) 的倍数，且 \(a^k \in \langle a^{n/d} \rangle\)。由此 \(H = \langle a^{n/d} \rangle\)。\(\blacksquare\)

\hrulefill

\subsection{39.6 对称群与 Cayley 定理}\label{ux5bf9ux79f0ux7fa4ux4e0e-cayley-ux5b9aux7406}

\textbf{定义 39.6.1（对称群）}：\(n\) 个元素的集合上所有置换构成的群称为 \(n\) 次\textbf{对称群}，记为 \(S_n\)。\(|S_n| = n!\)。

\textbf{定义 39.6.2（置换的轮换分解）}：任何置换 \(\sigma \in S_n\) 可唯一分解为互不相交的轮换的乘积（不考虑次序）。

\textbf{定理 39.6.1（Cayley 定理）}：任何群 \(G\) 同构于某个对称群的一个子群。特别地，若 \(|G| = n\)，则 \(G\) 同构于 \(S_n\) 的一个子群。

\textbf{证明}：对每个 \(g \in G\)，定义左乘映射 \(L_g : G \to G\)，\(L_g(x) = gx\)。\(L_g\) 是 \(G\) 上的置换（因为 \(L_{g^{-1}}\) 是其逆）。

定义 \(\varphi : G \to S_G\)（\(S_G\) 为 \(G\) 上所有置换构成的群）为 \(\varphi(g) = L_g\)。

\(\varphi\) 是同态：\(\varphi(gh)(x) = L_{gh}(x) = ghx = L_g(hx) = L_g(L_h(x)) = (\varphi(g) \circ \varphi(h))(x)\)。

\(\varphi\) 是单射：若 \(\varphi(g) = \text{id}\)，则 \(L_g(e) = ge = e\)，故 \(g = e\)。

因此 \(G \cong \operatorname{im} \varphi \leq S_G\)。若 \(|G| = n\)，则 \(S_G \cong S_n\)。\(\blacksquare\)

\hrulefill

\subsection{39.7 群的直积}\label{ux7fa4ux7684ux76f4ux79ef}

\textbf{定义 39.7.1（直积）}：设 \(G_1, G_2, \ldots, G_n\) 为群。它们的\textbf{直积}为

\[
G_1 \times G_2 \times \cdots \times G_n = \{(g_1, g_2, \ldots, g_n) \mid g_i \in G_i\}
\]

运算为逐分量乘法：\((g_1, \ldots, g_n)(h_1, \ldots, h_n) = (g_1 h_1, \ldots, g_n h_n)\)。

\textbf{定理 39.7.1（有限阿贝尔群的结构定理，陈述）}：任何有限阿贝尔群同构于素数幂阶循环群的直积。即存在素数 \(p_1, \ldots, p_k\) 和正整数 \(e_1, \ldots, e_k\) 使得

\[
G \cong \mathbb{Z}_{p_1^{e_1}} \times \mathbb{Z}_{p_2^{e_2}} \times \cdots \times \mathbb{Z}_{p_k^{e_k}}
\]

该分解在因子顺序下是唯一的（当要求 \(p_1^{e_1} \mid p_2^{e_2} \mid \cdots \mid p_k^{e_k}\) 时）。

\textbf{注记 39.7.1}：完整证明将在 Ch 61 中作为主理想整环上有限生成模的结构定理的推论给出。

\hrulefill

\hrulefill

\hrulefill

\hrulefill

\hrulefill

\hrulefill

\section{第112章：群作用与 Sylow 定理}\label{ux7b2c112ux7ae0ux7fa4ux4f5cux7528ux4e0e-sylow-ux5b9aux7406}

群作用是将抽象的群论与具体的对称性联系起来的桥梁。通过群作用，群 \(G\) 不再是一个孤立的代数结构，而是作为某个集合 \(X\) 的对称变换群''活''了起来。这一观点由 Klein 在 1872 年 Erlangen 纲领中明确提出，并由 Burnside 等人在 19 世纪末发展为系统的群作用理论。群作用在数学的各个分支中无处不在：置换群作用于集合，一般线性群 \(GL_n(\mathbb{R})\) 作用于 \(\mathbb{R}^n\)，Galois 群作用于域扩张的根，基本群作用于覆盖空间，Lie 群作用于齐性空间。

Sylow 定理是有限群论中最辉煌的成果之一，由挪威数学家 Ludvig Sylow 于 1872 年发表。这三个定理精确刻画了有限群中素数幂阶子群的存在性、共轭性和计数，构成了有限群结构理论的核心支柱。Sylow 第一定理断言：若 \(p^k \mid |G|\)，则 \(G\) 存在 \(p^k\) 阶子群------特别地，Sylow \(p\)-子群（\(p^n\) 阶子群，其中 \(p^n \mid\mid |G|\)）一定存在。Sylow 第二定理断言：所有 Sylow \(p\)-子群互相共轭，这赋予了 Sylow \(p\)-子群以''几何''意义------它们构成 \(G\) 在其上的一个传递作用。Sylow 第三定理给出 Sylow \(p\)-子群的个数 \(n_p\) 满足 \(n_p \equiv 1 \pmod{p}\) 且 \(n_p \mid [G:P]\)，这为分类有限群提供了强有力的约束条件。

Sylow 定理的威力体现在它能够通过纯算术条件推断群的结构性质。例如，若 \(|G| = pq\)（\(p < q\) 为素数），则 \(n_q \equiv 1 \pmod{q}\) 且 \(n_q \mid p\)，由此立即推出 \(n_q = 1\)，即 \(q\)-Sylow 子群正规，进而可证 \(G\) 是半直积，当 \(q \not\equiv 1 \pmod{p}\) 时 \(G\) 甚至是循环群。这种从阶数信息直接读出群结构的推理模式是有限群论中最优美的论证之一。本章将系统建立群作用、Burnside 引理和 Sylow 三定理的完整理论，为有限群的分类和表示论（卷十三）奠定基础。

\textbf{定义 40.1.1（群作用）}：群 \(G\) 在集合 \(X\) 上的一个\textbf{作用}（左作用）是一个映射 \(G \times X \to X\)，\((g, x) \mapsto g \cdot x\)，满足：

\begin{enumerate}
\def\labelenumi{\arabic{enumi}.}
\tightlist
\item
  \(e \cdot x = x\)（\(\forall x \in X\)）
\item
  \(g \cdot (h \cdot x) = (gh) \cdot x\)（\(\forall g, h \in G, \forall x \in X\)）
\end{enumerate}

\textbf{定义 40.1.2（轨道与稳定子）}：设 \(G\) 作用在 \(X\) 上，\(x \in X\)。

\begin{itemize}
\tightlist
\item
  \(x\) 的\textbf{轨道}：\(G \cdot x = \{g \cdot x \mid g \in G\}\)
\item
  \(x\) 的\textbf{稳定子群}：\(G_x = \{g \in G \mid g \cdot x = x\}\)
\end{itemize}

\textbf{定理 40.1.1（轨道-稳定子定理）}：设 \(G\) 为有限群，作用在有限集 \(X\) 上。则对任意 \(x \in X\)，

\[
|G \cdot x| = [G : G_x] = \frac{|G|}{|G_x|}
\]

特别地，\(|G \cdot x|\) 整除 \(|G|\)。

\textbf{证明}：定义映射 \(\varphi : G/G_x \to G \cdot x\) 为 \(\varphi(g G_x) = g \cdot x\)。

良定义：若 \(g G_x = h G_x\)，则 \(h^{-1}g \in G_x\)，故 \((h^{-1}g) \cdot x = x\)，即 \(g \cdot x = h \cdot x\)。

单射：若 \(g \cdot x = h \cdot x\)，则 \(h^{-1}g \cdot x = x\)，故 \(h^{-1}g \in G_x\)，\(g G_x = h G_x\)。

满射：显然。

因此 \(|G \cdot x| = |G/G_x| = [G : G_x] = |G|/|G_x|\)。\(\blacksquare\)

\textbf{定理 40.1.2（Burnside 引理）}：设有限群 \(G\) 作用在有限集 \(X\) 上。轨道的个数为

\[
\frac{1}{|G|} \sum_{g \in G} |\operatorname{Fix}(g)|
\]

其中 \(\operatorname{Fix}(g) = \{x \in X \mid g \cdot x = x\}\)。

\textbf{证明}：计数集合 \(\{(g, x) \in G \times X \mid g \cdot x = x\}\) 的两种方式：

\[
\sum_{g \in G} |\operatorname{Fix}(g)| = \sum_{x \in X} |G_x| = \sum_{x \in X} \frac{|G|}{|G \cdot x|}
\]

设轨道为 \(\mathcal{O}_1, \ldots, \mathcal{O}_k\)。对每个轨道 \(\mathcal{O}_i\)，取 \(x_i \in \mathcal{O}_i\)，则对 \(x \in \mathcal{O}_i\)，\(|G \cdot x| = |\mathcal{O}_i|\)，故

\[
\sum_{x \in \mathcal{O}_i} \frac{|G|}{|G \cdot x|} = |\mathcal{O}_i| \cdot \frac{|G|}{|\mathcal{O}_i|} = |G|
\]

因此 \(\sum_{g \in G} |\operatorname{Fix}(g)| = k |G|\)，即 \(k = \frac{1}{|G|} \sum_{g \in G} |\operatorname{Fix}(g)|\)。\(\blacksquare\)

\hrulefill

\subsection{40.2 共轭作用与类方程}\label{ux5171ux8f6dux4f5cux7528ux4e0eux7c7bux65b9ux7a0b}

\textbf{定义 40.2.1（共轭作用）}：群 \(G\) 在自身上的\textbf{共轭作用}定义为 \(g \cdot x = gxg^{-1}\)。

在此作用下，\(x\) 的轨道称为 \(x\) 的\textbf{共轭类}，稳定子群称为 \(x\) 的\textbf{中心化子}，记为 \(C_G(x) = \{g \in G \mid gxg^{-1} = x\}\)。

\textbf{定义 40.2.2（中心）}：\(G\) 的\textbf{中心} \(Z(G) = \{g \in G \mid gx = xg, \, \forall x \in G\} = \{g \in G \mid C_G(g) = G\}\)。

\textbf{定理 40.2.1（类方程）}：设 \(G\) 为有限群。则

\[
|G| = |Z(G)| + \sum_{i=1}^{k} [G : C_G(x_i)]
\]

其中 \(x_1, \ldots, x_k\) 是非中心元素的共轭类的代表元。

\textbf{证明}：\(G\) 被划分为共轭类。\(Z(G)\) 中的元素恰为那些共轭类大小为 \(1\) 的元素。对其余共轭类，由轨道-稳定子定理，\(|共轭类| = [G : C_G(x_i)]\)。\(\blacksquare\)

\textbf{推论 40.2.2}：若 \(|G| = p^n\)（\(p\) 为素数，\(n \geq 1\)，此时 \(G\) 称为 \(p\)-群），则 \(|Z(G)| > 1\)（即非平凡 \(p\)-群的中心非平凡）。

\textbf{证明}：由类方程，\(|G| = |Z(G)| + \sum [G : C_G(x_i)]\)。\([G : C_G(x_i)]\) 整除 \(|G| = p^n\) 且不为 \(1\)（\(x_i \notin Z(G)\)），故是 \(p\) 的倍数。因此 \(|Z(G)| \equiv |G| \equiv 0 \pmod{p}\)。又 \(e \in Z(G)\)，故 \(|Z(G)| \geq p > 1\)。\(\blacksquare\)

\hrulefill

\subsection{40.3 Sylow 定理}\label{sylow-ux5b9aux7406}

\textbf{定义 40.3.1（Sylow \(p\)-子群）}：设 \(G\) 为有限群，\(|G| = p^n m\)（\(p \nmid m\)）。\(G\) 的 \(p^n\) 阶子群称为 \(G\) 的 \textbf{Sylow \(p\)-子群}。

\textbf{定理 40.3.1（Sylow 第一定理------存在性）}：设 \(G\) 为有限群，\(p\) 为素数，\(p^k \mid |G|\)。则 \(G\) 存在 \(p^k\) 阶子群。特别地，Sylow \(p\)-子群存在。

\textbf{证明}：对 \(|G|\) 归纳。\(|G| = 1\) 时平凡。

设 \(|G| = p^n m\)（\(p \nmid m\)）。考虑类方程 \(|G| = |Z(G)| + \sum [G : C_G(x_i)]\)。

\textbf{情形一}：\(p \nmid |Z(G)|\)。则存在某个 \(x_i\) 使得 \(p \nmid [G : C_G(x_i)]\)。由 \([G : C_G(x_i)] = |G|/|C_G(x_i)|\) 和 \(p^n \mid |G|\)，得 \(p^n \mid |C_G(x_i)|\)。由于 \(x_i \notin Z(G)\)，\(C_G(x_i) \subsetneq G\)。由归纳假设应用在 \(C_G(x_i)\) 上，\(C_G(x_i)\) 存在 \(p^n\) 阶子群（即 \(G\) 的 Sylow \(p\)-子群）。

\textbf{情形二}：\(p \mid |Z(G)|\)。由 Cauchy 定理（定理 40.3.2），\(Z(G)\) 存在 \(p\) 阶元素 \(a\)。\(N = \langle a \rangle \trianglelefteq G\)（因为 \(N \subseteq Z(G)\)）。\(|G/N| = |G|/p = p^{n-1}m\)。由归纳假设，\(G/N\) 存在 Sylow \(p\)-子群 \(P/N\)（\(|P/N| = p^{n-1}\)）。则 \(|P| = p^n\)，\(P\) 即为 \(G\) 的 Sylow \(p\)-子群。

对于一般的 \(p^k\)，由 Sylow \(p\)-子群的存在性，它包含 \(p^k\) 阶子群（由归纳法可证）。\(\blacksquare\)

\textbf{定理 40.3.2（Cauchy 定理）}：若素数 \(p \mid |G|\)，则 \(G\) 存在 \(p\) 阶元素。

\textbf{证明}：考虑集合 \(X = \{(g_1, \ldots, g_p) \in G^p : g_1 g_2 \cdots g_p = e\}\)。\(|X| = |G|^{p-1}\)（前 \(p-1\) 个分量可自由选择，最后一个由方程 \(g_p = (g_1 \cdots g_{p-1})^{-1}\) 唯一确定）。循环群 \(\mathbb{Z}_p\) 通过循环置换作用在 \(X\) 上：\(k \cdot (g_1, \ldots, g_p) = (g_{k+1}, \ldots, g_p, g_1, \ldots, g_k)\)（\(k = 0, 1, \ldots, p-1\)）。验证这是群作用：\(0\) 作用为恒等，且 \(k \cdot (\ell \cdot x) = (k+\ell) \cdot x\)。

由轨道-稳定子定理（定理 40.1.1），每个轨道的大小整除 \(|\mathbb{Z}_p| = p\)，故为 \(1\) 或 \(p\)。不动点（轨道大小为 \(1\) 的元素）满足 \((g, \ldots, g) \in X\)，即 \(g^p = e\)。不动点包括 \((e, \ldots, e)\)（对应 \(g = e\)）以及可能的满足 \(|g| = p\) 的元素。

由于 \(p \mid |X| = |G|^{p-1}\)，而不动点总数 \(\equiv |X| \equiv 0 \pmod{p}\)（因非平凡轨道大小为 \(p\)），故至少有 \(p\) 个不动点。因此存在 \(g \neq e\) 使 \(g^p = e\)，即 \(|g| = p\)。\(\blacksquare\)

\textbf{定理 40.3.3（Sylow 第二定理------共轭性）}：设 \(G\) 为有限群。\(G\) 的任何两个 Sylow \(p\)-子群是共轭的。即若 \(P, Q\) 均为 Sylow \(p\)-子群，则存在 \(g \in G\) 使得 \(Q = gPg^{-1}\)。

\textbf{证明}：设 \(X = \{gP \mid g \in G\}\) 为 \(P\) 的左陪集组成的集合。\(Q\) 通过左乘作用在 \(X\) 上。将 \(X\) 划分为 \(Q\)-轨道。

由轨道-稳定子定理，每个轨道的大小整除 \(|Q| = p^n\)，故为 \(p\) 的幂。\(|X| = [G : P] = m\)（\(p \nmid m\)）。由于轨道的和为 \(m\)，必存在大小为 \(1\) 的轨道（否则每个轨道大小都是 \(p\) 的倍数，总和也是 \(p\) 的倍数）。

设该轨道为 \(\{gP\}\)，即 \(Q \cdot gP = gP\)。这意味着 \(QgP = gP\)，即 \(g^{-1}Qg \subseteq P\)。由于 \(|g^{-1}Qg| = |Q| = |P|\)，故 \(g^{-1}Qg = P\)，即 \(Q = gPg^{-1}\)。\(\blacksquare\)

\textbf{推论 40.3.4（Sylow 第三定理------计数）}：设 \(n_p\) 为 \(G\) 的 Sylow \(p\)-子群的个数。则

\begin{enumerate}
\def\labelenumi{(\roman{enumi})}
\tightlist
\item
  \(n_p \equiv 1 \pmod{p}\)
\item
  \(n_p \mid m\)（其中 \(|G| = p^n m\)，\(p \nmid m\)）
\end{enumerate}

\textbf{证明}：设 \(P\) 为 Sylow \(p\)-子群。令 \(X\) 为所有 Sylow \(p\)-子群的集合。\(P\) 通过共轭作用在 \(X\) 上。将 \(X\) 划分为 \(P\)-轨道。

每个轨道的大小整除 \(|P| = p^n\)，故为 \(p\) 的幂。\(\{P\}\) 是大小为 \(1\) 的轨道。若存在另一个大小为 \(1\) 的轨道 \(\{Q\}\)，则 \(P Q P^{-1} = Q\)，即 \(P \subseteq N_G(Q)\)。\(P\) 和 \(Q\) 都是 \(N_G(Q)\) 的 Sylow \(p\)-子群。由 Sylow 第二定理（应用在 \(N_G(Q)\) 上），\(P\) 和 \(Q\) 在 \(N_G(Q)\) 中共轭，但 \(Q \trianglelefteq N_G(Q)\)，故 \(P = Q\)。

因此仅有 \(\{P\}\) 是大小为 \(1\) 的轨道，其余轨道大小均为 \(p\) 的倍数。故 \(n_p \equiv 1 \pmod{p}\)。

\begin{enumerate}
\def\labelenumi{(\roman{enumi})}
\setcounter{enumi}{1}
\tightlist
\item
  由 Sylow 第二定理，\(G\) 传递地作用在 Sylow \(p\)-子群的集合上（通过共轭）。由轨道-稳定子定理，\(n_p = [G : N_G(P)]\) 整除 \([G : P] = m\)。\(\blacksquare\)
\end{enumerate}

\hrulefill

\hrulefill

\hrulefill

\hrulefill

\hrulefill

\newpage

\chapter{05-代数/02-抽象代数一/02-环论基础.md}\label{ux4ee3ux657002-ux62bdux8c61ux4ee3ux6570ux4e0002-ux73afux8bbaux57faux7840.md}

\section{第113章：环论基础}\label{ux7b2c113ux7ae0ux73afux8bbaux57faux7840}

环是代数结构中的第二个基本层次------它比群多了一个乘法运算，却因此展现出远比群论丰富的结构多样性。环论起源于 19 世纪对代数数论和代数几何的研究：Dedekind 在 1870 年代引入理想（Ideal）的概念来修复代数整数环中唯一分解的失效，Kronecker 和 Hilbert 将环论发展为代数几何的基础语言。20 世纪初，Emmy Noether 以她的升链条件（Noether 环）和理想论彻底改变了交换环论的面貌，使环论成为现代代数的核心支柱。

环与群的根本区别在于环拥有两种运算------加法和乘法------且它们通过分配律相互关联。环关于加法构成阿贝尔群，而乘法只需满足结合律（不一定交换，不一定有幺元，非零元不一定有逆元）。正是这些''乘法缺陷''创造了环论的丰富性：非交换环（如矩阵环 \(M_n(F)\)、四元数代数 \(\mathbb{H}\)）催生了非交换代数和表示论；非整环（如 \(\mathbb{Z}/6\mathbb{Z}\)）揭示了零因子的存在如何破坏消去律；而理想的引入则提供了在缺乏乘法逆元时进行''除法''的替代方案。

本章的核心内容包括：理想与商环------理想是环论中对应于群论中''正规子群''的概念，商环 \(R/I\) 赋予环的''模运算''以严格代数意义；环同态与同构定理------三个同构定理在环论中的对应版本，揭示环及其商环之间的结构关系；素理想与极大理想------素理想是''素数''概念的推广（\(ab \in \mathfrak{p} \implies a \in \mathfrak{p} \text{ 或 } b \in \mathfrak{p}\)），而极大理想则是''最大真理想''，它们与商环的整环性和域性一一对应：\(\mathfrak{p}\) 是素理想 \(\iff R/\mathfrak{p}\) 是整环，\(\mathfrak{m}\) 是极大理想 \(\iff R/\mathfrak{m}\) 是域。这些概念为域扩张理论（Ch 65）和交换代数（卷十二）中的局部环、Krull 维数、Dedekind 整环等高级主题奠定了坚实基础。

\textbf{定义 41.1.1（理想）}：设 \(R\) 为环。子集 \(I \subseteq R\) 称为 \(R\) 的\textbf{理想}，若

\begin{enumerate}
\def\labelenumi{\arabic{enumi}.}
\tightlist
\item
  \((I, +)\) 是 \((R, +)\) 的子群
\item
  \(\forall r \in R, \forall a \in I\)，\(ra \in I\)（左理想）且 \(ar \in I\)（右理想）
\end{enumerate}

\textbf{定义 41.1.2（商环）}：设 \(I\) 为环 \(R\) 的理想。在加法商群 \(R/I\) 上定义乘法：\((r + I)(s + I) = rs + I\)。\((R/I, +, \cdot)\) 构成环，称为 \(R\) 模 \(I\) 的\textbf{商环}。

\textbf{定理 41.1.1}：商环中乘法的定义是良定义的。

\textbf{证明}：设 \(r + I = r' + I\)，\(s + I = s' + I\)。则 \(r' = r + a\)，\(s' = s + b\)（\(a, b \in I\)）。\(r's' = (r + a)(s + b) = rs + rb + as + ab\)。由于 \(I\) 是理想，\(rb, as, ab \in I\)，故 \(r's' + I = rs + I\)。\(\blacksquare\)

\hrulefill

\subsection{41.2 环同态与同构定理}\label{ux73afux540cux6001ux4e0eux540cux6784ux5b9aux7406}

\textbf{定义 41.2.1（环同态）}：设 \(R, S\) 为环。映射 \(\varphi : R \to S\) 若满足：

\begin{enumerate}
\def\labelenumi{\arabic{enumi}.}
\tightlist
\item
  \(\varphi(a + b) = \varphi(a) + \varphi(b)\)
\item
  \(\varphi(ab) = \varphi(a)\varphi(b)\)
\item
  \(\varphi(1_R) = 1_S\)（若 \(R, S\) 含幺）
\end{enumerate}

则称 \(\varphi\) 为\textbf{环同态}。

\textbf{定理 41.2.1（环的第一同构定理）}：设 \(\varphi : R \to S\) 为环同态。则 \(\ker \varphi\) 是 \(R\) 的理想，且

\[
R / \ker \varphi \cong \operatorname{im} \varphi
\]

\textbf{证明}：\(\ker \varphi\) 是理想（类似群的情形，需额外验证乘法吸收性）。在加法群层面，由群的第一同构定理，\(R/\ker \varphi \cong \operatorname{im} \varphi\)（作为加法群）。同构映射 \(\overline{\varphi}(r + \ker \varphi) = \varphi(r)\) 也保持乘法：\(\overline{\varphi}((r + K)(s + K)) = \overline{\varphi}(rs + K) = \varphi(rs) = \varphi(r)\varphi(s) = \overline{\varphi}(r + K)\overline{\varphi}(s + K)\)。\(\blacksquare\)

\textbf{定理 41.2.2（环的第二、第三同构定理）}：与群的同构定理类似，将 ``子群'' 替换为 ``子环''，``正规子群'' 替换为 ``理想''。

\hrulefill

\subsection{41.3 素理想与极大理想}\label{ux7d20ux7406ux60f3ux4e0eux6781ux5927ux7406ux60f3}

\textbf{定义 41.3.1（素理想与极大理想）}：设 \(R\) 为交换环，\(I\) 为 \(R\) 的真理想（\(I \neq R\)）。

\begin{enumerate}
\def\labelenumi{\arabic{enumi}.}
\tightlist
\item
  \(I\) 称为\textbf{素理想}，若 \(ab \in I \implies a \in I\) 或 \(b \in I\)
\item
  \(I\) 称为\textbf{极大理想}，若不存在理想 \(J\) 使得 \(I \subsetneq J \subsetneq R\)
\end{enumerate}

\textbf{定理 41.3.1}：设 \(R\) 为含幺交换环，\(I\) 为 \(R\) 的理想。

\begin{enumerate}
\def\labelenumi{(\roman{enumi})}
\tightlist
\item
  \(I\) 是素理想 \(\iff\) \(R/I\) 是整环
\item
  \(I\) 是极大理想 \(\iff\) \(R/I\) 是域
\end{enumerate}

\textbf{证明}：

\begin{enumerate}
\def\labelenumi{(\roman{enumi})}
\item
  \(I\) 是素理想 \(\iff\) \(\forall a, b \in R\)，\(ab \in I \implies a \in I \text{ 或 } b \in I\) \(\iff\) \(\forall \overline{a}, \overline{b} \in R/I\)，\(\overline{a}\overline{b} = \overline{0} \implies \overline{a} = \overline{0} \text{ 或 } \overline{b} = \overline{0}\) \(\iff\) \(R/I\) 无零因子 \(\iff\) \(R/I\) 是整环。
\item
  （\(\Rightarrow\)）设 \(I\) 为极大理想。取 \(\overline{a} \in R/I \setminus \{\overline{0}\}\)，即 \(a \notin I\)。考虑理想 \(J = I + (a)\)。\(J \supsetneq I\)，由极大性 \(J = R\)。故存在 \(i \in I\)，\(r \in R\) 使得 \(1 = i + ra\)。在 \(R/I\) 中，\(\overline{1} = \overline{r}\overline{a}\)，故 \(\overline{a}\) 可逆。因此 \(R/I\) 中每个非零元可逆，即为域。
\end{enumerate}

（\(\Leftarrow\)）设 \(R/I\) 为域。若 \(J \supseteq I\) 为 \(R\) 的理想，\(J/I\) 是 \(R/I\) 的理想。域只有平凡理想，故 \(J/I = \{\overline{0}\}\) 或 \(J/I = R/I\)。前者 \(J = I\)，后者 \(J = R\)。因此 \(I\) 是极大理想。\(\blacksquare\)

\textbf{推论 41.3.2}：极大理想必为素理想。

\textbf{证明}：域 \(\implies\) 整环。由定理 41.3.1，\(R/I\) 是域 \(\implies\) \(R/I\) 是整环 \(\implies\) \(I\) 是极大理想 \(\implies\) \(I\) 是素理想。\(\blacksquare\)

\hrulefill

\subsection{41.4 中国剩余定理}\label{ux4e2dux56fdux5269ux4f59ux5b9aux7406-1}

\textbf{定理 41.4.1（中国剩余定理------环论版）}：设 \(R\) 为含幺交换环，\(I_1, \ldots, I_n\) 为 \(R\) 的理想，满足 \(I_i + I_j = R\)（\(\forall i \neq j\)）。则

\[
R / (I_1 \cap I_2 \cap \cdots \cap I_n) \cong R/I_1 \times R/I_2 \times \cdots \times R/I_n
\]

\textbf{证明}：对 \(n\) 归纳。\(n = 2\) 时，\(I_1 + I_2 = R\) 保证存在 \(a_1 \in I_1\)，\(a_2 \in I_2\) 使得 \(a_1 + a_2 = 1\)。

定义 \(\varphi : R \to R/I_1 \times R/I_2\) 为 \(\varphi(r) = (r + I_1, r + I_2)\)。\(\varphi\) 是环同态。

\(\ker \varphi = I_1 \cap I_2\)。需证 \(\varphi\) 是满射。对任意 \((r_1 + I_1, r_2 + I_2)\)，取 \(r = r_2 a_1 + r_1 a_2\)。则

\[
r - r_1 = r_2 a_1 + r_1(a_2 - 1) = r_2 a_1 - r_1 a_1 \in I_1
\] \[
r - r_2 = r_2(a_1 - 1) + r_1 a_2 = -r_2 a_2 + r_1 a_2 \in I_2
\]

故 \(\varphi(r) = (r_1 + I_1, r_2 + I_2)\)。由第一同构定理得证。

\(n > 2\) 时，令 \(J = I_1 \cap \cdots \cap I_{n-1}\)。可证 \(J + I_n = R\)（利用 \(I_i + I_n = R\) 和归纳）。由 \(n = 2\) 的情形，\(R/(J \cap I_n) \cong R/J \times R/I_n\)。由归纳假设，\(R/J \cong R/I_1 \times \cdots \times R/I_{n-1}\)。合起来即得。\(\blacksquare\)

\textbf{推论 41.4.2（整数版中国剩余定理）}：设 \(m_1, \ldots, m_n\) 为两两互素的正整数。则

\[
\mathbb{Z} / (m_1 \cdots m_n) \mathbb{Z} \cong \mathbb{Z}/m_1\mathbb{Z} \times \cdots \times \mathbb{Z}/m_n\mathbb{Z}
\]

\textbf{证明}：\(m_i\mathbb{Z} + m_j\mathbb{Z} = \gcd(m_i, m_j)\mathbb{Z} = \mathbb{Z}\)（\(\gcd(m_i, m_j) = 1\)）。且 \(m_1\mathbb{Z} \cap \cdots \cap m_n\mathbb{Z} = \operatorname{lcm}(m_1, \ldots, m_n)\mathbb{Z} = m_1 \cdots m_n \mathbb{Z}\)。由定理 41.4.1 得证。\(\blacksquare\)

\hrulefill

\hrulefill

\hrulefill

\hrulefill

\hrulefill

\hrulefill

\section{第114章：多项式环与因式分解}\label{ux7b2c114ux7ae0ux591aux9879ux5f0fux73afux4e0eux56e0ux5f0fux5206ux89e3}

多项式环 \(F[x]\) 是域 \(F\) 上最简单的环扩张，它既是代数中最基本的对象，也是连接代数学与几何学（代数曲线）、数论（代数整数环）和分析学（幂级数环）的枢纽。多项式环的研究可追溯至古代对代数方程求解的努力，但直到 19 世纪才被 Abel、Galois 和 Kronecker 等人提升为系统的代数理论。多项式环 \(F[x]\) 具有欧几里得整环的结构------带余除法成立------这使它成为研究更一般环中因式分解理论的原型。

环论中因式分解理论的核心是四个环类之间的蕴含关系：欧几里得整环（ED）\(\implies\) 主理想整环（PID）\(\implies\) 唯一分解整环（UFD）。在 ED 中，带余除法提供了构造性的因式分解算法；在 PID 中，每个理想由一个元素生成，这保证了不可约元素与素元素的等价性；在 UFD 中，每个非零非单位元素可以唯一地（在相伴意义下）分解为不可约元素的乘积。多项式环 \(F[x]\) 是 ED 的典范例子，\(\mathbb{Z}[x]\) 是 UFD 但不是 PID 的例子，而 \(\mathbb{Z}[\sqrt{-5}]\) 则不是 UFD（\(6 = 2 \cdot 3 = (1+\sqrt{-5})(1-\sqrt{-5})\) 给出两种不等价的不可约分解）。

Gauss 引理和 Eisenstein 判别法是判定多项式不可约性的两个基本工具。Gauss 引理断言：本原多项式的乘积是本原多项式，由此推出有理系数多项式的可约性等价于整系数多项式的可约性。Eisenstein 判别法则提供了构造不可约多项式的系统方法：对于 \(f(x) = a_n x^n + \cdots + a_0\)，若存在素数 \(p\) 整除除首项系数外的所有系数且 \(p^2 \nmid a_0\)，则 \(f\) 不可约------这立即证明了 \(x^n - 2\) 在 \(\mathbb{Q}[x]\) 中不可约。本章以代数基本定理（\(\mathbb{C}[x]\) 中不可约多项式恰为一次多项式）作为多项式因式分解的终极结论，为 Galois 理论中分裂域和根式可解性的讨论做好铺垫。

\textbf{定义 42.1.1（多项式环）}：设 \(R\) 为交换环。\(R\) 上的一元多项式环 \(R[x]\) 是所有形式表达式 \(\sum_{i=0}^{n} a_i x^i\)（\(a_i \in R\)）的集合，配备通常的多项式加法和乘法。

\textbf{定理 42.1.1（带余除法）}：设 \(F\) 为域，\(f, g \in F[x]\)，\(g \neq 0\)。则存在唯一的 \(q, r \in F[x]\) 使得

\[
f = qg + r, \quad \deg r < \deg g
\]

\textbf{证明}：对 \(\deg f\) 归纳。若 \(\deg f < \deg g\)，取 \(q = 0\)，\(r = f\) 即可。

若 \(\deg f = n \geq m = \deg g\)，设 \(f = a_n x^n + \cdots\)，\(g = b_m x^m + \cdots\)。令 \(f_1 = f - \frac{a_n}{b_m} x^{n-m} g\)。则 \(\deg f_1 < n\)。由归纳假设，\(f_1 = q_1 g + r\)（\(\deg r < \deg g\)）。则 \(f = (q_1 + \frac{a_n}{b_m} x^{n-m})g + r\)。

唯一性：若 \(f = qg + r = q'g + r'\)，则 \((q - q')g = r' - r\)。若 \(q \neq q'\)，左边次数 \(\geq \deg g\)，右边次数 \(< \deg g\)，矛盾。故 \(q = q'\)，\(r = r'\)。\(\blacksquare\)

\textbf{定义 42.1.2（多项式函数与根）}：\(f \in F[x]\)，\(\alpha \in F\)。若 \(f(\alpha) = 0\)，则称 \(\alpha\) 为 \(f\) 的\textbf{根}。

\textbf{定理 42.1.2（因式定理）}：\(\alpha \in F\) 是 \(f \in F[x]\) 的根当且仅当 \(x - \alpha \mid f(x)\)。

\textbf{证明}：由带余除法，\(f(x) = q(x)(x - \alpha) + r\)（\(r \in F\)）。代入 \(x = \alpha\) 得 \(f(\alpha) = r\)。故 \(f(\alpha) = 0 \iff r = 0 \iff x - \alpha \mid f(x)\)。\(\blacksquare\)

\hrulefill

\subsection{42.2 欧几里得整环、主理想整环、唯一分解整环}\label{ux6b27ux51e0ux91ccux5f97ux6574ux73afux4e3bux7406ux60f3ux6574ux73afux552fux4e00ux5206ux89e3ux6574ux73af}

\textbf{定义 42.2.1（欧几里得整环）}：整环 \(R\) 称为\textbf{欧几里得整环}（ED），若存在函数 \(\varphi : R \setminus \{0\} \to \mathbb{N}\) 使得对任意 \(a, b \in R\)（\(b \neq 0\)），存在 \(q, r \in R\) 满足

\[
a = qb + r, \quad \text{且 } r = 0 \text{ 或 } \varphi(r) < \varphi(b)
\]

\textbf{定义 42.2.2（主理想整环）}：整环 \(R\) 称为\textbf{主理想整环}（PID），若 \(R\) 的每个理想都是主理想（即由单个元素生成）。

\textbf{定理 42.2.1}：ED \(\implies\) PID。

\textbf{证明}：设 \(R\) 为 ED，\(I\) 为 \(R\) 的非零理想。取 \(b \in I \setminus \{0\}\) 使得 \(\varphi(b)\) 最小。对任意 \(a \in I\)，由带余除法 \(a = qb + r\)（\(r = 0\) 或 \(\varphi(r) < \varphi(b)\)）。由于 \(r = a - qb \in I\)，由 \(\varphi(b)\) 的最小性，\(r = 0\)。故 \(a = qb\)，即 \(I = (b)\)。\(\blacksquare\)

\textbf{定义 42.2.3（不可约元素与素元素）}：设 \(R\) 为整环。

\begin{itemize}
\tightlist
\item
  \(p \in R \setminus \{0\}\)（非单位）称为\textbf{不可约元素}，若 \(p = ab \implies a\) 或 \(b\) 是单位
\item
  \(p \in R \setminus \{0\}\)（非单位）称为\textbf{素元素}，若 \(p \mid ab \implies p \mid a\) 或 \(p \mid b\)
\end{itemize}

\textbf{定理 42.2.2}：在整环中，素元素必为不可约元素。在 PID 中，不可约元素也是素元素。

\textbf{证明}：设 \(p\) 为素元素，\(p = ab\)。则 \(p \mid ab\)，故 \(p \mid a\) 或 \(p \mid b\)。不妨设 \(p \mid a\)，\(a = pc\)。则 \(p = pcb\)，由于整环中可消去 \(p\)，\(cb = 1\)，故 \(b\) 是单位。因此 \(p\) 不可约。

设 \(R\) 为 PID，\(p\) 不可约。若 \(p \mid ab\) 且 \(p \nmid a\)。考虑理想 \((p, a)\)。由于 \(R\) 是 PID，\((p, a) = (d)\)。\(d \mid p\)，故 \(d\) 是单位或 \(p\) 的相伴。若 \(d\) 是 \(p\) 的相伴，则 \(p \mid a\)，矛盾。故 \(d\) 是单位，\((p, a) = R\)。存在 \(x, y \in R\) 使得 \(xp + ya = 1\)。则 \(b = xpb + yab\)。\(p \mid xpb\) 且 \(p \mid yab\)（因为 \(p \mid ab\)），故 \(p \mid b\)。因此 \(p\) 是素元素。\(\blacksquare\)

\textbf{定义 42.2.4（唯一分解整环）}：整环 \(R\) 称为\textbf{唯一分解整环}（UFD），若

\begin{enumerate}
\def\labelenumi{\arabic{enumi}.}
\tightlist
\item
  每个非零非单位元素可分解为不可约元素的乘积
\item
  该分解在相伴意义下唯一（即若 \(a = p_1 \cdots p_m = q_1 \cdots q_n\)，则 \(m = n\) 且适当重排后 \(p_i\) 与 \(q_i\) 相伴）
\end{enumerate}

\textbf{定理 42.2.3}：PID \(\implies\) UFD。

\textbf{证明}（概要）：

（存在性）若存在非零非单位元素 \(a\) 不能分解为不可约元素的乘积，则可构造无穷长的严格增链 \((a) \subsetneq (a_1) \subsetneq (a_2) \subsetneq \cdots\)。但 \(\bigcup (a_i)\) 是理想，在 PID 中为主理想 \((d)\)，\(d\) 属于某个 \((a_n)\)，导致链终止，矛盾。

（唯一性）由定理 42.2.2，PID 中不可约元素是素元素。若 \(p_1 \cdots p_m = q_1 \cdots q_n\)，\(p_1\) 是素的，则 \(p_1 \mid q_j\)（某个 \(j\)）。\(q_j\) 不可约，故 \(p_1\) 与 \(q_j\) 相伴。消去 \(p_1\) 和 \(q_j\)，对剩余部分归纳。\(\blacksquare\)

\textbf{推论 42.2.4}：\(F[x]\) 是 UFD。

\hrulefill

\subsection{42.3 Gauss 引理与 Eisenstein 判别法}\label{gauss-ux5f15ux7406ux4e0e-eisenstein-ux5224ux522bux6cd5}

\textbf{定义 42.3.1（本原多项式）}：\(f \in \mathbb{Z}[x]\) 称为\textbf{本原多项式}，若其系数的最大公因数为 \(1\)。

\textbf{定理 42.3.1（Gauss 引理）}：两个本原多项式的乘积是本原多项式。

\textbf{证明}：设 \(f = \sum a_i x^i\)，\(g = \sum b_j x^j\) 为本原多项式。若 \(fg\) 不是本原的，则存在素数 \(p\) 整除 \(fg\) 的所有系数。设 \(a_r\) 是 \(f\) 中第一个不被 \(p\) 整除的系数，\(b_s\) 是 \(g\) 中第一个不被 \(p\) 整除的系数。则 \(fg\) 的 \(x^{r+s}\) 系数为

\[
c_{r+s} = \sum_{i+j = r+s} a_i b_j = a_r b_s + \sum_{i < r \text{ 或 } j < s} a_i b_j
\]

由 \(r, s\) 的选法，\(a_r b_s\) 不被 \(p\) 整除，而其余项至少有一个因子被 \(p\) 整除，矛盾。故 \(fg\) 是本原的。\(\blacksquare\)

\textbf{定理 42.3.2（Gauss 引理------推论）}：若 \(f \in \mathbb{Z}[x]\) 在 \(\mathbb{Q}[x]\) 中可约，则 \(f\) 在 \(\mathbb{Z}[x]\) 中也可约。

\textbf{证明}：设 \(f = gh\)（\(g, h \in \mathbb{Q}[x]\)）。通分可得 \(df = g_1 h_1\)（\(g_1, h_1 \in \mathbb{Z}[x]\)，\(d \in \mathbb{Z}\)）。提取 \(g_1, h_1\) 的各系数的最大公因数，得 \(g_1 = c_1 g^*\)，\(h_1 = c_2 h^*\)（\(g^*, h^*\) 本原）。由 Gauss 引理，\(g^* h^*\) 本原，故 \(d = c_1 c_2\)。从而 \(f = c_1 c_2 g^* h^*\) 在 \(\mathbb{Z}[x]\) 中可约。\(\blacksquare\)

\textbf{定理 42.3.3（Eisenstein 判别法）}：设 \(f(x) = a_n x^n + \cdots + a_1 x + a_0 \in \mathbb{Z}[x]\)。若存在素数 \(p\) 使得

\begin{enumerate}
\def\labelenumi{(\roman{enumi})}
\tightlist
\item
  \(p \nmid a_n\)
\item
  \(p \mid a_i\)（\(i = 0, 1, \ldots, n-1\)）
\item
  \(p^2 \nmid a_0\)
\end{enumerate}

则 \(f\) 在 \(\mathbb{Q}[x]\) 中不可约。

\textbf{证明}：假设 \(f = gh\)（\(g, h \in \mathbb{Z}[x]\)，非单位）。设 \(g = b_r x^r + \cdots + b_0\)，\(h = c_s x^s + \cdots + c_0\)（\(r + s = n\)）。

由于 \(a_0 = b_0 c_0\) 且 \(p \mid a_0\) 但 \(p^2 \nmid a_0\)，不妨设 \(p \mid b_0\) 但 \(p \nmid c_0\)。

\(p \nmid a_n\) 意味着 \(p \nmid b_r\) 且 \(p \nmid c_s\)。设 \(b_k\) 是 \(b_0, \ldots, b_r\) 中第一个不被 \(p\) 整除的系数。则 \(k \leq r < n\)。

考虑 \(a_k = \sum_{i=0}^{k} b_i c_{k-i}\)。对 \(i < k\)，\(p \mid b_i\)；对 \(i = k\)，\(p \nmid b_k\) 且 \(p \nmid c_0\)（因为 \(c_0\) 不被 \(p\) 整除）。故 \(p \nmid a_k\)。但 \(k < n\)，由条件 (ii) \(p \mid a_k\)，矛盾。\(\blacksquare\)

\hrulefill

\subsection{42.4 代数基本定理（陈述与推论）}\label{ux4ee3ux6570ux57faux672cux5b9aux7406ux9648ux8ff0ux4e0eux63a8ux8bba}

\textbf{定理 42.4.1（代数基本定理）}：任何非常数多项式 \(f \in \mathbb{C}[x]\) 在 \(\mathbb{C}\) 中有根。

\textbf{注记 42.4.1}：代数基本定理的完整证明需要复分析（卷十）或拓扑学（卷十一）的工具。此处陈述结果，以便在多项式理论中使用。

\textbf{推论 42.4.2}：\(\mathbb{C}[x]\) 中的不可约多项式恰为一次多项式。

\textbf{推论 42.4.3}：\(\mathbb{R}[x]\) 中的不可约多项式为一次多项式或判别式为负的二次多项式。

\textbf{推论 42.4.4}：任何 \(n\) 次复系数多项式 \(f(x) = a_n x^n + \cdots + a_0\) 可分解为

\[
f(x) = a_n (x - \alpha_1)(x - \alpha_2) \cdots (x - \alpha_n)
\]

其中 \(\alpha_1, \ldots, \alpha_n \in \mathbb{C}\)（重根按重数计算）。

\hrulefill

\hrulefill

\hrulefill

\hrulefill

\hrulefill

\hrulefill

\section{第115章：非交换环论深化（Non-Commutative Ring Theory）}\label{ux7b2c115ux7ae0ux975eux4ea4ux6362ux73afux8bbaux6df1ux5316non-commutative-ring-theory}

结合环论研究未必交换的环的结构理论。与交换代数的几何导向不同，非交换环论发展出了一套独立的分类和结构理论。

\subsection{77.1 Wedderburn-Artin 理论}\label{wedderburn-artin-ux7406ux8bba}

\textbf{定理 77.1}（Wedderburn-Artin）：\(R\) 是半单（Jacobson 根为零且 Artin）环当且仅当 \(R\) 同构于有限个除环上的矩阵环的直积：\(R \cong \prod_{i=1}^{k} M_{n_i}(D_i)\)（\(D_i\) 为除环）。

Wedderburn-Artin 定理是非交换环论中最重要的结构定理，其地位类似于有限群论中的 Jordan-Holder 定理和有限维李代数中的 Levi 分解。该定理起源于 Wedderburn（1908）对有限维代数的分类和 Artin（1927）对 Artin 环的推广。定理的核心意义在于：半单 Artin 环的结构完全由除环 \(D_i\) 和正整数 \(n_i\) 决定，不存在任何''奇异''或''非经典''的示例。矩阵环 \(M_n(D)\) 是除环 \(D\) 上所有 \(n \times n\) 矩阵构成的环，其任意左理想形如 \(M_n(D) e\)（\(e\) 为幂等矩阵），单左理想对应于秩为 1 的矩阵。在域 \(\mathbb{F}\) 上的有限维代数情形，定理简化为 \(A \cong \prod_i M_{n_i}(\mathbb{F})\)（若 \(\mathbb{F}\) 为代数闭域，则 \(D_i = \mathbb{F}\)）。Wedderburn-Artin 定理的一个关键推论是：半单 Artin 环的每个左理想都是直和项（即投射模），且所有单模均为有限维。该定理的证明依赖于 Jacobson 稠密定理和 Schur 引理，二者共同确保了环作用在极小左理想上的忠实性。在现代非交换代数几何中，非交换射影概形（non-commutative projective schemes）的''点''被定义为相应非交换环的商环的简单因子，直接源于 Wedderburn-Artin 分解。

\textbf{证明}：设 \(R\) 为半单 Artin 环。Artin 性保证 \(R\) 有极小左理想 \(L\)。由 Schur 引理，\(D = \operatorname{End}_R(L)\) 为除环。\(R\) 半单故 \(R \cong \bigoplus_{i} \operatorname{Hom}_R(L_i, L_i)\)，其中 \(L_i\) 为极小左理想的同构类代表。每个 \(\operatorname{Hom}_R(L_i, L_i) \cong D_i\) 为除环。\(R\) 在 \(L_i\) 上的作用给出同态 \(R \to \operatorname{End}_{D_i}(L_i) \cong M_{n_i}(D_i^{\text{op}})\)。由 Jacobson 稠密定理，此同态为满射。核为 Jacobson 根 \(J(R)\)，而 \(R\) 半单故 \(J(R) = 0\)，得同构 \(R \cong \prod_i M_{n_i}(D_i)\)。 \(\blacksquare\)

\textbf{推论}：域 \(\mathbb{F}\) 上有限维半单代数恰为 \(\prod_i M_{n_i}(\mathbb{F})\)。

\subsection{77.2 Jacobson 根与素环}\label{jacobson-ux6839ux4e0eux7d20ux73af}

\textbf{定义 77.1}（Jacobson 根）：环 \(R\) 的 Jacobson 根 \(\operatorname{Jac}(R)\) 是所有单左 \(R\)-模的零化子的交集，也等于所有本原理想的交集，也等于所有极大左理想的交集。\(R\) 称为 \textbf{Jacobson 半单} 若 \(\operatorname{Jac}(R) = 0\)。

Jacobson 根（Jacobson radical）是 Nathan Jacobson 于 1945 年引入的环论基本概念，是非交换环论中与交换代数中 Jacobson 根（所有极大理想的交）的完美对应。\(\operatorname{Jac}(R)\) 的多种等价刻画揭示了其本质：\(\operatorname{Jac}(R) = \{r \in R : 1 - rs \text{ 可逆}, \forall s \in R\}\)（元素的左拟正则性），这使得 Jacobson 根在计算上极为便利。\(\operatorname{Jac}(R)\) 是 \(R\) 的双边理想，且 \(R/\operatorname{Jac}(R)\) 是 Jacobson 半单环。对交换环，\(\operatorname{Jac}(R)\) 恰为所有极大理想的交；对 Artin 环，\(\operatorname{Jac}(R)\) 是幂零的，且 \(R/\operatorname{Jac}(R)\) 是半单 Artin 环。Jacobson 根与零化子链（ACC 和 DCC）的关系是 Goldie 定理和 Hopkins-Levitzki 定理的核心。Nakayama 引理（定理 77.2）是 Jacobson 根最重要的应用之一，在交换代数（局部环中极大理想 \(\mathfrak{m}\) 满足 \(\mathfrak{m}M = M \Rightarrow M = 0\)）和表示论中均有基本版本。

\textbf{定理 77.2}（Nakayama 引理）：设 \(M\) 是有限生成 \(R\)-模且 \(\operatorname{Jac}(R) M = M\)，则 \(M = 0\)。

\textbf{证明}：设 \(M\) 为有限生成 \(R\)-模且 \(\{x_1, \ldots, x_n\}\) 为极小生成元组。若 \(M \neq 0\)，则 \(n \geq 1\)。由 \(\operatorname{Jac}(R) M = M\)，存在 \(a_{ij} \in \operatorname{Jac}(R)\) 使 \(x_n = \sum_i a_{in} x_i\)，即 \((1 - a_{nn}) x_n = \sum_{i=1}^{n-1} a_{in} x_i\)。但 \(a_{nn} \in \operatorname{Jac}(R)\) 蕴含 \(1 - a_{nn}\) 为可逆元，故 \(x_n\) 可由 \(x_1, \ldots, x_{n-1}\) 生成，与极小性矛盾。因此 \(M = 0\)。 \(\blacksquare\)

\subsection{77.3 Goldie 定理与 Noether 环}\label{goldie-ux5b9aux7406ux4e0e-noether-ux73af}

\textbf{定理 77.3}（Goldie, 1958）：环 \(R\) 是半素 Goldie 环（无幂零理想、ACC 上零化子条件）当且仅当 \(R\) 有半单 Artin 经典商环 \(Q(R)\)：\(R \subseteq Q(R)\) 满足 \(Q(R)\) 是半单 Artin 且 \(R\) 的正则元在 \(Q(R)\) 中可逆。

\textbf{证明}：设 \(R\) 为半素 Goldie 环。\(R\) 的 ACC 上零化子条件保证 \(R\) 有有限 Goldie 秩。令 \(\mathcal{C}\) 为 \(R\) 的正则元集合。\(R\) 半素 + Goldie 条件保证 \(\mathcal{C}\) 为左 Ore 集，故可构造左经典商环 \(Q(R) = \mathcal{C}^{-1} R\)。\(Q(R)\) 的结构：\(Q(R)\) 为半单 Artin 环（由 Goldie 定理的核心结论），由 Wedderburn-Artin 定理 \(Q(R) \cong \prod_i M_{n_i}(D_i)\)。\(R\) 在 \(Q(R)\) 中为左序，即 \(R\) 为 \(Q(R)\) 的左 Goldie 序。逆方向：若 \(R\) 有半单 Artin 经典商环，则 \(R\) 必为半素 Goldie 环，由商环的 Artin 性可推出 \(R\) 的 ACC 条件。 \(\blacksquare\)

Goldie 定理在半素环和非交换 Noether 环之间架起了一座桥梁：半素 Goldie 环的商环总是半单 Artin 环，从而将 Wedderburn-Artin 分类定理的功能延伸到更广的非交换环境中。该定理的直接推论深刻地简化了非交换代数几何中许多构造------Noether 素环被刻画为恰是那些具有半单 Artin 商环的环，这为非交换射影概形的局部理论提供了代数基础。

\textbf{推论}：Noether 素环恰为具有半单 Artin 商环的环。

\subsection{77.4 PI 环与 Kaplansky 定理}\label{pi-ux73afux4e0e-kaplansky-ux5b9aux7406}

\textbf{定义 77.2}（多项式恒等环）：环 \(R\) 满足多项式恒等式若存在非零多项式 \(f(x_1,\ldots,x_n) \in \mathbb{Z}\langle x_1,\ldots,x_n\rangle\)（自由结合代数中的元）使 \(f(r_1,\ldots,r_n) = 0\) 对所有 \(r_i \in R\) 成立。

PI 环（Polynomial Identity rings）理论是非交换环论中一个极为深刻的分支，它捕捉了''近交换''行为的代数本质。标准恒等式 \(S_n(x_1,\ldots,x_n) = \sum_{\sigma \in S_n} \operatorname{sgn}(\sigma) x_{\sigma(1)} \cdots x_{\sigma(n)}\) 是最重要的 PI 恒等式之一：Amitsur-Levitzki 定理（1950）证明 \(M_k(\mathbb{C})\) 满足 \(S_{2k}\) 但不满足 \(S_{2k-1}\)，因此 \(S_{2k}\) 是矩阵代数的''特征恒等式''。PI 环理论的核心成果包括：Posner 定理（素 PI 环有中心除环商环）、Artin-Procesi 定理（具有恒等式 \(S_n\) 的 Azumaya 代数刻画）、以及 Razmyslov-Kemer-Braun 的 PI 环的根定理。PI 环与不变量理论有密切联系：若 \(G\) 为简约群作用于仿射代数 \(A\)，则不变子代数 \(A^G\) 是 PI 环（当 \(A\) 为 PI 时）。此外，PI 理论为非交换代数几何提供了关键工具：具有 PI 性质的环的表示范畴（即 PI 概形）由中心子概形的 Azumaya 代数结构完全控制，这构成了 Artin 的非交换代数几何程序的基础。

\textbf{定理 77.4}（Kaplansky, 1948）：本原 PI 环是除环上有限维中心单代数。特别地，本原 PI 代数必为中心除环上的矩阵代数 \(M_n(D)\)（\([D:Z(D)] < \infty\)）。结合 Posner 定理（素 PI 环是 Goldie 环），这给出了 PI 环的完整本质刻画------所有 ``代数行为'' 最终可约化为有限维矩阵代数的子代数。

\textbf{证明}：设 \(A\) 为本原 PI 代数，\(V\) 为其忠实单模。由 Jacobson 稠密定理，\(A\) 在 \(\operatorname{End}_D(V)\) 中稠密，\(D = \operatorname{End}_A(V)\) 为除环。若 \(\dim_D V\) 无限，则 \(A\) 包含所有有限秩线性变换，不满足任何 PI 恒等式，矛盾。故 \(\dim_D V = n < \infty\)，\(A \cong M_n(D)\)。PI 条件给出 \(S_{2n} = 0\) 为 \(A\) 的非平凡恒等式，\(n\) 为 PI-次数。由 Posner 定理，\(A\) 的中心 \(Z(A)\) 为域，\(A\) 为 \(Z(A)\) 上有限维中心单代数，商环 \(Q(A) = M_n(Z(D))\)。 \(\blacksquare\)

\subsection{非交换环论的现代发展与前沿}\label{ux975eux4ea4ux6362ux73afux8bbaux7684ux73b0ux4ee3ux53d1ux5c55ux4e0eux524dux6cbf}

非交换环论在现代数学中与多个分支产生了深刻的交叉，其核心主题从结构分类转向了范畴化和几何化。\textbf{Morita 等价}（1958）是非交换环论中最基本的等价关系：两个环 \(R\) 和 \(S\) 是 Morita 等价的当且仅当它们的模范畴等价，即 \(R\text{-Mod} \cong S\text{-Mod}\)。Morita 等价保留了几乎所有同调性质（投射维数、整体维数、K-理论），但允许环本身的''大小''变化------例如 \(R\) 与 \(M_n(R)\) 总是 Morita 等价的。\textbf{Morita 上下文}（Morita context）通过双模 \((P, Q, f, g)\) 提供了构造 Morita 等价的显式方法，这在非交换代数几何中用于定义非交换空间的''坐标变换''。\textbf{非交换射影几何}（Artin-Zhang 1997）将交换射影概形 \(\operatorname{Proj} A\) 推广到非交换分次环 \(A\)，通过定义非交换点（point modules）和尾等价（tail equivalence）来构造非交换概形。Artin-Stafford 定理（1995）分类了二次三维 Artin-Schelter 正则代数，揭示了非交换射影平面与椭圆曲线和 Sklyanin 代数之间的深刻联系。\textbf{非交换局部化}（Cohn 的通用局部化）推广了交换代数中的局部化理论，利用 Sylvester 秩函数和通用除环将非交换环嵌入最简单可能的商环中。\textbf{Hochschild 上同调}和\textbf{循环上同调}（Connes 1985）为非交换环提供了微分几何不变量：\(HH^0(R) = Z(R)\)（中心），\(HH^1(R)\) 分类导子模内导子，\(HH^2(R)\) 分类代数扩张。循环上同调 \(HC^\bullet(R)\) 则与非交换微分形式的 de Rham 上同调对应，是非交换几何中 Chern 特征标理论的代数基础。\textbf{Frobenius 代数}和\textbf{Calabi-Yau 代数}（Ginzburg 2006）是两类特殊的非交换代数，前者具有非退化 Frobenius 双线性型，后者满足 \(\operatorname{RHom}_{A^e}(A, A^e) \cong A[-d]\)（\(d\) 为 Calabi-Yau 维数），它们与拓扑弦论和镜像对称有深刻联系。

\hrulefill

\hrulefill

\hrulefill

\hrulefill

\section{第116章：多项式}\label{ux7b2c116ux7ae0ux591aux9879ux5f0f}

\begin{quote}
\textbf{前置知识}：本章需要第12章关于实数域运算律（域公理）的知识，以及带余除法所需的基本实数性质。
\end{quote}

多项式是初等代数中最重要的研究对象之一，它架起了代数运算与函数分析之间的桥梁。本章从多项式的形式定义出发，将多项式视为 \(\mathbb{R}[x]\) 这个多项式环中的元素，系统地建立多项式的代数运算（加法、乘法、数乘）和相应的运算律。多项式运算的核心是带余除法（多项式版的欧几里得除法），它保证了任意两个多项式之间可以进行''商+余式''的分解，且余式的次数严格小于除式的次数。在此基础上，本章将引入余数定理和因式定理，建立多项式根与因式之间的精确对应关系，并讨论韦达定理（根与系数的关系）。最后，本章将系统介绍因式分解的五种基本方法，并讨论实数域和复数域上不可约多项式的分类。多项式的理论结构------作为一个欧几里得整环------与整数环 \(\mathbb{Z}\) 有着深刻的类比，这种类比是抽象代数中''欧几里得整环''统一理论的源头。

\subsection{7.1 多项式的定义与运算}\label{ux591aux9879ux5f0fux7684ux5b9aux4e49ux4e0eux8fd0ux7b97}

\textbf{定义 7.1.1（多项式）}

设 \(x\) 为未定元（变量），\(a_0, a_1, \ldots, a_n \in \mathbb{R}\)（\(a_n \neq 0\)），表达式

\[
f(x) = a_n x^n + a_{n-1} x^{n-1} + \cdots + a_1 x + a_0 = \sum_{k=0}^{n} a_k x^k
\]

称为关于 \(x\) 的 \textbf{\(n\) 次多项式}。其中 \(a_n\) 称为\textbf{首项系数}，\(a_0\) 称为\textbf{常数项}。若所有系数均为零，称为\textbf{零多项式}，记为 \(0\)，不定义其次数。

\textbf{定义 7.1.2（多项式的相等）}

两个多项式 \(f(x) = \sum a_k x^k\) 和 \(g(x) = \sum b_k x^k\) \textbf{相等}，当且仅当对所有 \(k\)，\(a_k = b_k\)。

\textbf{定义 7.1.3（多项式的加法）}

设 \(f(x) = \sum_{k=0}^{n} a_k x^k\)，\(g(x) = \sum_{k=0}^{m} b_k x^k\)，定义

\[
(f + g)(x) = \sum_{k=0}^{\max(n,m)} (a_k + b_k) x^k
\]

其中当 \(k\) 超出某多项式次数时，对应系数视为 \(0\)。

\textbf{定义 7.1.4（多项式的乘法）}

\[
(f \cdot g)(x) = \sum_{k=0}^{n+m} c_k x^k, \quad c_k = \sum_{i+j=k} a_i b_j
\]

\textbf{定理 7.1.1} 若 \(f(x)\) 为 \(n\) 次多项式，\(g(x)\) 为 \(m\) 次多项式（\(f, g\) 均非零），则 \(f(x) \cdot g(x)\) 为 \(n + m\) 次多项式。

\textbf{证明}：设 \(f(x) = a_n x^n + a_{n-1} x^{n-1} + \cdots + a_0\)（\(a_n \neq 0\)），\(g(x) = b_m x^m + b_{m-1} x^{m-1} + \cdots + b_0\)（\(b_m \neq 0\)）。两个多项式相乘时，\(f(x)g(x)\) 的 \(x^{n+m}\) 次项系数来自 \(f\) 中 \(x^n\) 项与 \(g\) 中 \(x^m\) 项的乘积，即 \(a_n b_m\)。对于次数高于 \(n+m\) 的项，由于 \(f\) 中最高次为 \(n\)，\(g\) 中最高次为 \(m\)，任何 \(i + j > n + m\) 的组合必然有 \(i > n\) 或 \(j > m\)，而此时 \(a_i = 0\) 或 \(b_j = 0\)，因此系数为 \(0\)。由于 \(\mathbb{R}\) 是整环（域当然是整环），\(a_n \neq 0\) 且 \(b_m \neq 0\) 蕴含 \(a_n b_m \neq 0\)。因此 \(f(x)g(x)\) 的最高次为 \(n + m\)，即 \(\deg(fg) = \deg f + \deg g\)。\(\blacksquare\)

\textbf{定理 7.1.2（多项式环）} 实系数多项式全体 \(\mathbb{R}[x]\) 在上述加法和乘法下构成含幺交换环。

\textbf{证明}：

\begin{itemize}
\tightlist
\item
  \((\mathbb{R}[x], +)\) 是阿贝尔群：加法封闭、结合、交换，零多项式为加法单位元，\(-f(x) = \sum(-a_k)x^k\) 为 \(f\) 的加法逆元。
\item
  乘法封闭（由多项式乘法定义）、结合律（系数卷积满足结合律：\((f \cdot g) \cdot h\) 与 \(f \cdot (g \cdot h)\) 的 \(k\) 次项系数均为 \(\sum_{i+j+\ell=k} a_i b_j c_\ell\)，由求和的结合律和交换律保证两式相等）、交换律（\(c_k = \sum_{i+j=k} a_i b_j = \sum_{i+j=k} b_j a_i\)，由实数乘法交换律）由实数的对应性质逐项验证。
\item
  分配律由系数层面的分配律逐项验证。
\item
  幺元为常数多项式 \(1\)。
\end{itemize}

\(\blacksquare\)

\hrulefill

\subsection{7.2 乘法公式}\label{ux4e58ux6cd5ux516cux5f0f}

以下每个公式均通过直接展开并利用分配律证明。

\textbf{定理 7.2.1（完全平方和）} 对任意 \(a, b \in \mathbb{R}\)：

\[
(a + b)^2 = a^2 + 2ab + b^2
\]

\textbf{证明}：

\[
(a + b)^2 = (a + b)(a + b) = a \cdot a + a \cdot b + b \cdot a + b \cdot b = a^2 + ab + ba + b^2 = a^2 + 2ab + b^2
\]

\(\blacksquare\)

\textbf{定理 7.2.2（完全平方差）} 对任意 \(a, b \in \mathbb{R}\)：

\[
(a - b)^2 = a^2 - 2ab + b^2
\]

\textbf{证明}：

\[
(a - b)^2 = (a + (-b))^2 = a^2 + 2a(-b) + (-b)^2 = a^2 - 2ab + b^2
\]

\(\blacksquare\)

\textbf{定理 7.2.3（平方差公式）} 对任意 \(a, b \in \mathbb{R}\)：

\[
(a + b)(a - b) = a^2 - b^2
\]

\textbf{证明}：

\[
(a + b)(a - b) = a^2 - ab + ba - b^2 = a^2 - ab + ab - b^2 = a^2 - b^2
\]

\(\blacksquare\)

\textbf{定理 7.2.4（完全立方和）} 对任意 \(a, b \in \mathbb{R}\)：

\[
(a + b)^3 = a^3 + 3a^2 b + 3ab^2 + b^3
\]

\textbf{证明}：

\[
(a + b)^3 = (a + b)^2 (a + b) = (a^2 + 2ab + b^2)(a + b)
\]

\[
= a^3 + a^2 b + 2a^2 b + 2ab^2 + ab^2 + b^3 = a^3 + 3a^2 b + 3ab^2 + b^3
\]

\(\blacksquare\)

\textbf{定理 7.2.5（立方和公式）} 对任意 \(a, b \in \mathbb{R}\)：

\[
a^3 + b^3 = (a + b)(a^2 - ab + b^2)
\]

\textbf{证明}：展开右端：

\[
(a + b)(a^2 - ab + b^2) = a^3 - a^2 b + ab^2 + a^2 b - ab^2 + b^3 = a^3 + b^3
\]

\(\blacksquare\)

\textbf{定理 7.2.6（立方差公式）} 对任意 \(a, b \in \mathbb{R}\)：

\[
a^3 - b^3 = (a - b)(a^2 + ab + b^2)
\]

\textbf{证明}：展开右端：

\[
(a - b)(a^2 + ab + b^2) = a^3 + a^2 b + ab^2 - a^2 b - ab^2 - b^3 = a^3 - b^3
\]

\(\blacksquare\)

\hrulefill

\subsection{7.3 因式分解}\label{ux56e0ux5f0fux5206ux89e3}

因式分解是多项式理论中的核心操作，其目标是将一个多项式表示为若干个不可约多项式的乘积。这一过程类似于整数理论中的素因子分解------每个多项式（在给定系数域上）可以唯一地分解为不可约因式的乘积（在常数因子意义下）。在实数域 \(\mathbb{R}\) 上，不可约多项式只能是 \(1\) 次或 \(2\) 次（判别式小于 \(0\)）的；在复数域 \(\mathbb{C}\) 上，由代数基本定理，每个 \(n\) 次多项式可以分解为 \(n\) 个一次因式的乘积。因式分解不仅是解方程的基础工具（通过因式定理将求根问题转化为因式分解问题），也是代数化简、分式运算和积分计算中的关键步骤。因式分解的''唯一性''（在常数因子意义下）是多项式环 \(\mathbb{R}[x]\) 作为唯一分解整环（UFD）的核心性质。以下介绍五种基本的因式分解方法。

从方法论角度，因式分解可以看作多项式运算的''逆运算''------正如除法是乘法的逆运算，因式分解是将一个复杂的多项式表达式''分解''为更简单的因式的乘积。这一逆运算并不总是容易的（正如大数分解在计算上是困难的），但存在一些系统的方法来处理常见的模式。五种基本方法------提公因式法、公式法（平方差、完全平方、立方和差）、分组分解法、十字相乘法、配方法------覆盖了初等代数中绝大多数因式分解问题。对于更高次的多项式，有理根定理（Rational Root Theorem）提供了寻找有理根的系统方法，从而将高次多项式的因式分解问题转化为寻根问题。理解这些方法的关键在于识别多项式的''结构模式''------一旦识别出多项式符合某种已知模式，分解就变得直接而机械。

\subsubsection{7.3.1 提公因式法}\label{ux63d0ux516cux56e0ux5f0fux6cd5}

\textbf{定理 7.3.1（提取公因式）} 若多项式 \(f(x)\) 的各项都含有公因式 \(d(x)\)，则

\[
f(x) = d(x) \cdot q(x)
\]

其中 \(q(x)\) 由 \(f(x)\) 的各项除以 \(d(x)\) 得到。

\textbf{证明}：设 \(f(x) = \sum_{k=0}^{n} a_k x^k\) 是任意多项式，且每一项 \(a_k x^k\) 都能被 \(d(x)\) 整除，即存在多项式 \(p_k(x)\) 使得 \(a_k x^k = d(x) \cdot p_k(x)\) 对所有 \(k\) 成立。那么

\[f(x) = \sum_{k=0}^{n} a_k x^k = \sum_{k=0}^{n} d(x) \cdot p_k(x) = d(x) \cdot \sum_{k=0}^{n} p_k(x)\]

这里最后一步使用了多项式乘法的分配律（定理 7.1.2）。因此 \(d(x) \mid f(x)\)，且商为 \(q(x) = \sum_{k=0}^{n} p_k(x)\)。在实际操作中，提取公因式即将多项式各项的公共因子 \(d(x)\) 提取到括号外面，括号内为各项除以 \(d(x)\) 所得的商之和。\(\blacksquare\)

\subsubsection{7.3.2 公式法}\label{ux516cux5f0fux6cd5}

利用7.2节的乘法公式反向进行分解：

\begin{itemize}
\tightlist
\item
  \(a^2 - b^2 = (a+b)(a-b)\)
\item
  \(a^2 + 2ab + b^2 = (a+b)^2\)
\item
  \(a^2 - 2ab + b^2 = (a-b)^2\)
\item
  \(a^3 + b^3 = (a+b)(a^2 - ab + b^2)\)
\item
  \(a^3 - b^3 = (a-b)(a^2 + ab + b^2)\)
\end{itemize}

\subsubsection{7.3.3 十字相乘法}\label{ux5341ux5b57ux76f8ux4e58ux6cd5}

\textbf{定理 7.3.2（十字相乘法）} 设 \(f(x) = ax^2 + bx + c\)。若存在整数 \(p, q, r, s\) 使得

\[
pr = a, \quad qs = c, \quad ps + qr = b
\]

则 \(f(x) = (px + q)(rx + s)\)。

\textbf{证明}：

\[
(px + q)(rx + s) = prx^2 + (ps + qr)x + qs = ax^2 + bx + c
\]

\(\blacksquare\)

\subsubsection{7.3.4 分组分解法}\label{ux5206ux7ec4ux5206ux89e3ux6cd5}

\textbf{定理 7.3.3（分组分解）} 若多项式 \(f(x)\) 可以写成 \(f(x) = g(x) + h(x)\)，且 \(g(x)\) 和 \(h(x)\) 有公共因式 \(d(x)\)，则

\[
f(x) = d(x) \cdot \left(\frac{g(x)}{d(x)} + \frac{h(x)}{d(x)}\right)
\]

\textbf{证明}：由于 \(d(x) \mid g(x)\) 且 \(d(x) \mid h(x)\)，存在多项式 \(p(x)\) 和 \(q(x)\) 使得 \(g(x) = d(x) \cdot p(x)\)，\(h(x) = d(x) \cdot q(x)\)。那么

\[f(x) = g(x) + h(x) = d(x) \cdot p(x) + d(x) \cdot q(x) = d(x) \cdot (p(x) + q(x))\]

这里最后一步使用了分配律的逆用（提取公因式）。而 \(p(x) = \frac{g(x)}{d(x)}\)，\(q(x) = \frac{h(x)}{d(x)}\)，因此

\[f(x) = d(x) \cdot \left(\frac{g(x)}{d(x)} + \frac{h(x)}{d(x)}\right)\]

分组分解法的核心思想是将多项式拆分为若干组，每组提取公因式后，各组之间再提取共同的因式。这种方法在因式分解中非常实用，尤其是当多项式不能直接提取公因式时。\(\blacksquare\)

分组分解法的核心思想是''先局部提取，再全局提取''------将多项式拆分为若干组，每组内部提取公因式后，各组之间产生了新的公共因式，进而可以再次提取。这种方法特别适用于四项及以上且没有单一公因式的多项式。典型的分组策略包括：按项数均分（如 \(ax + ay + bx + by = a(x+y) + b(x+y) = (a+b)(x+y)\)）、按同类项分组（如 \(x^3 + x^2 + x + 1 = x^2(x+1) + 1(x+1) = (x^2+1)(x+1)\)）、按交叉项分组（如 \(ac + ad + bc + bd = a(c+d) + b(c+d) = (a+b)(c+d)\)）。分组分解法本质上利用了加法结合律和分配律，它揭示了多项式环中''因式''概念的灵活性------同一个多项式可能有多种分组方式，不同的分组方式可能导致不同的分解结果，但在不可约因式分解的意义下最终结果唯一（在常数因子和顺序意义下）。分组分解法在 Galois 理论中也有其深层对应：多项式的 Galois 群结构决定了其可解性。

\subsubsection{7.3.5 配方法}\label{ux914dux65b9ux6cd5}

\textbf{定理 7.3.4（配方法）} 对任意 \(a \neq 0\)，\(b, c \in \mathbb{R}\)：

\[
ax^2 + bx + c = a\left(x + \frac{b}{2a}\right)^2 + c - \frac{b^2}{4a}
\]

\textbf{证明}：

\[
ax^2 + bx + c = a\left(x^2 + \frac{b}{a}x\right) + c = a\left(x^2 + \frac{b}{a}x + \frac{b^2}{4a^2} - \frac{b^2}{4a^2}\right) + c
\]

\[
= a\left(x + \frac{b}{2a}\right)^2 - \frac{b^2}{4a} + c = a\left(x + \frac{b}{2a}\right)^2 + \frac{4ac - b^2}{4a}
\]

\(\blacksquare\)

\hrulefill

\subsection{7.4 带余除法}\label{ux5e26ux4f59ux9664ux6cd5-1}

\textbf{定理 7.4.1（带余除法定理）} 设 \(f(x), g(x) \in \mathbb{R}[x]\)，\(g(x) \neq 0\)。则存在唯一的 \(q(x), r(x) \in \mathbb{R}[x]\)，使得

\[
f(x) = g(x) \cdot q(x) + r(x)
\]

其中 \(r(x) = 0\) 或 \(\deg(r) < \deg(g)\)。

\textbf{证明}： \textbf{存在性}：对 \(f\) 的次数进行归纳。 \textbf{基础}：若 \(\deg(f) < \deg(g)\)，取 \(q(x) = 0\)，\(r(x) = f(x)\)。 \textbf{归纳}：设 \(\deg(f) = n \geq \deg(g) = m\)。设 \(f(x) = a_n x^n + \cdots\)，\(g(x) = b_m x^m + \cdots\)。 令 \(f_1(x) = f(x) - \dfrac{a_n}{b_m} x^{n-m} g(x)\)。 则 \(\deg(f_1) < n\)（首项消去）。由归纳假设，存在 \(q_1(x), r(x)\) 使得 \(f_1(x) = g(x) \cdot q_1(x) + r(x)\)，其中 \(r(x) = 0\) 或 \(\deg(r) < m\)。 从而 \(f(x) = g(x) \cdot \left(\dfrac{a_n}{b_m} x^{n-m} + q_1(x)\right) + r(x)\)。取 \(q(x) = \dfrac{a_n}{b_m} x^{n-m} + q_1(x)\) 即可。 \textbf{唯一性}：设 \(f(x) = g(x) q_1(x) + r_1(x) = g(x) q_2(x) + r_2(x)\)。 则 \(g(x)(q_1(x) - q_2(x)) = r_2(x) - r_1(x)\)。 若 \(q_1 \neq q_2\)，则左端次数 \(\geq \deg(g)\)，而右端 \(r_2 - r_1\) 的次数 \(< \deg(g)\)（或为零多项式），矛盾。故 \(q_1 = q_2\)，从而 \(r_1 = r_2\)。\(\blacksquare\)

\hrulefill

\subsection{7.5 余数定理与因式定理}\label{ux4f59ux6570ux5b9aux7406ux4e0eux56e0ux5f0fux5b9aux7406}

\textbf{定理 7.5.1（余数定理）} 多项式 \(f(x)\) 除以 \((x - a)\) 的余数等于 \(f(a)\)。

\textbf{证明}：由带余除法，\(f(x) = (x - a) q(x) + r\)，其中 \(r\) 为常数（因 \(\deg(x-a) = 1\)，\(\deg(r) < 1\)）。

代入 \(x = a\)：\(f(a) = (a - a) q(a) + r = r\)。 \(\blacksquare\)

\textbf{定理 7.5.2（因式定理）} \((x - a)\) 是 \(f(x)\) 的因式，当且仅当 \(f(a) = 0\)。

\textbf{证明}：

(\(\Rightarrow\)) 若 \((x - a) | f(x)\)，则 \(f(x) = (x - a) q(x)\)，\(f(a) = 0 \cdot q(a) = 0\)。

(\(\Leftarrow\)) 若 \(f(a) = 0\)，由余数定理，\(f(x) = (x - a) q(x) + f(a) = (x - a) q(x)\)。 \(\blacksquare\)

\hrulefill

\subsection{7.6 多项式的根}\label{ux591aux9879ux5f0fux7684ux6839}

\textbf{定义 7.6.1（多项式的根）}

设 \(f(x) \in \mathbb{R}[x]\)。若 \(f(c) = 0\)，则称 \(c\) 为 \(f(x)\) 的一个\textbf{根}（零点）。

\textbf{定理 7.6.1（根的重数与因式分解）} 若 \(c\) 是 \(f(x)\) 的 \(k\) 重根（\(k \geq 1\)），则

\[
f(x) = (x - c)^k \cdot g(x)
\]

其中 \(g(c) \neq 0\)。

\textbf{证明}：对 \(k\) 归纳。

\begin{itemize}
\tightlist
\item
  \textbf{基础}：\(k = 1\) 时，由因式定理 \(f(x) = (x - c) g(x)\)，\(g(c) \neq 0\)（否则 \(c\) 为至少二重根）。
\item
  \textbf{归纳}：设 \(c\) 为 \(f(x)\) 的 \((k+1)\) 重根。由归纳假设（\(k\) 重根的因式分解），\(f(x) = (x - c)^k h(x)\)。由 \((k+1)\) 重根的定义，\((x-c)^{k+1}\) 整除 \(f(x)\)，故 \(h(c) = 0\)（否则 \((x-c)^{k+1}\) 不能整除 \((x-c)^k h(x)\)）。由因式定理，\(h(x) = (x - c) g(x)\)，\(g(c) \neq 0\)。故 \(f(x) = (x - c)^{k+1} g(x)\)。
\end{itemize}

\(\blacksquare\)

\textbf{定理 7.6.2（\(n\) 次多项式至多有 \(n\) 个根）} 非零的 \(n\) 次多项式 \(f(x)\) 在 \(\mathbb{R}\) 中至多有 \(n\) 个根（计入重数）。

\textbf{证明}：对 \(n\) 归纳。

\begin{itemize}
\tightlist
\item
  \textbf{基础}：\(n = 0\) 时，\(f(x) = a_0 \neq 0\)，无根。
\item
  \textbf{归纳}：设对 \(n - 1\) 次多项式结论成立。若 \(f(x)\) 无根，则结论成立（\(0 \leq n\)）。若 \(c\) 是 \(f(x)\) 的根，由因式定理，\(f(x) = (x - c) g(x)\)，\(\deg(g) = n - 1\)。
\end{itemize}

\(f(x)\) 的根恰好是 \(c\) 加上 \(g(x)\) 的根。由归纳假设，\(g(x)\) 至多有 \(n - 1\) 个根。故 \(f(x)\) 至多有 \(n\) 个根。 \(\blacksquare\)

\textbf{推论 7.6.1} 若 \(n\) 次多项式 \(f(x)\) 有超过 \(n\) 个不同的根，则 \(f(x) = 0\)（零多项式）。

\textbf{证明} 对 \(n\) 归纳。\(n = 0\) 时，零次多项式 \(f(x) = c\)（\(c \neq 0\)）无根，与''超过 0 个根''矛盾，故 \(f(x) = 0\)。✓

假设 \(n-1\) 时成立。设 \(f(x)\) 是 \(n\) 次多项式且有超过 \(n\) 个不同的根 \(r_1, r_2, \ldots, r_{n+1}\)。由因式定理，\(f(x) = (x - r_1)g(x)\)，其中 \(g(x)\) 是 \(n-1\) 次多项式。对 \(i = 2, 3, \ldots, n+1\)，\(0 = f(r_i) = (r_i - r_1)g(r_i)\)，由 \(r_i \neq r_1\) 得 \(g(r_i) = 0\)。故 \(g(x)\) 有 \(n\) 个不同的根 \(r_2, \ldots, r_{n+1}\)，由归纳假设 \(g(x) = 0\)，从而 \(f(x) = 0\)。\(\blacksquare\)

\textbf{定理 7.6.3（韦达定理）} 设 \(f(x) = a_n x^n + a_{n-1} x^{n-1} + \cdots + a_0\)（\(a_n \neq 0\)）的 \(n\) 个根为 \(x_1, x_2, \ldots, x_n\)（在复数域中计重数），则：

\[
\begin{cases}
x_1 + x_2 + \cdots + x_n = -\dfrac{a_{n-1}}{a_n} \\[8pt]
\displaystyle\sum_{1 \leq i < j \leq n} x_i x_j = \frac{a_{n-2}}{a_n} \\[8pt]
\cdots \\[4pt]
x_1 x_2 \cdots x_n = (-1)^n \dfrac{a_0}{a_n}
\end{cases}
\]

一般地，设 \(e_k\) 为 \(x_1, \ldots, x_n\) 的 \(k\) 次初等对称多项式：

\[
e_k = \sum_{1 \leq i_1 < \cdots < i_k \leq n} x_{i_1} \cdots x_{i_k}
\]

则 \(a_{n-k} = (-1)^k a_n \cdot e_k\)。

\textbf{证明}：由代数基本定理（定理 8.6.1），\(n\) 次多项式在复数域中恰有 \(n\) 个根（计重数），因此可以完全分解为一次因式的乘积：（完整证明见卷九定理53.1.3（复分析证明）和卷十定理60.14（拓扑学证明））。

\[
f(x) = a_n (x - x_1)(x - x_2) \cdots (x - x_n)
\]

展开右端：

\[
a_n \prod_{i=1}^{n}(x - x_i) = a_n \left(x^n - e_1 x^{n-1} + e_2 x^{n-2} - \cdots + (-1)^n e_n\right)
\]

其中 \(e_k\) 是 \(x_1, \ldots, x_n\) 的 \(k\) 次初等对称多项式。这可以通过对 \(n\) 归纳来严格证明：

\begin{itemize}
\tightlist
\item
  \textbf{基础}：\(n = 1\) 时，\(a_1(x - x_1) = a_1 x - a_1 x_1\)，\(e_1 = x_1\)，系数为 \((-1)^1 a_1 e_1 = -a_1 x_1\)。正确。
\item
  \textbf{归纳}：假设 \(a_n \prod_{i=1}^{n}(x - x_i) = a_n \sum_{k=0}^{n} (-1)^k e_k x^{n-k}\)。则
\end{itemize}

\[
a_n \prod_{i=1}^{n+1}(x - x_i) = a_n (x - x_{n+1}) \sum_{k=0}^{n} (-1)^k e_k x^{n-k}
\]

\[
= a_n \left(\sum_{k=0}^{n} (-1)^k e_k x^{n+1-k} - x_{n+1} \sum_{k=0}^{n} (-1)^k e_k x^{n-k}\right)
\]

\[
= a_n \sum_{k=0}^{n+1} (-1)^k \left(e_k + x_{n+1} e_{k-1}\right) x^{n+1-k}
\]

由对称多项式的递推关系 \(e_k^{(n+1)} = e_k^{(n)} + x_{n+1} e_{k-1}^{(n)}\)，结论成立。

比较系数，\(a_{n-k} = (-1)^k a_n \cdot e_k\)。特别地：

\[
e_1 = \sum x_i = -\frac{a_{n-1}}{a_n}, \quad e_n = \prod x_i = (-1)^n \frac{a_0}{a_n}
\]

\(\blacksquare\)

\begin{quote}
\textbf{本章展望}：多项式的运算和因式分解是代数的核心工具。下一章将利用这些工具系统地研究各类方程的求解理论。
\end{quote}

\hrulefill

\hrulefill

\hrulefill

\hrulefill

\hrulefill

\newpage

\chapter{05-代数/02-抽象代数一/03-模论基础.md}\label{ux4ee3ux657002-ux62bdux8c61ux4ee3ux6570ux4e0003-ux6a21ux8bbaux57faux7840.md}

\section{第117章：模论基础}\label{ux7b2c117ux7ae0ux6a21ux8bbaux57faux7840}

模论是环论的自然推广------它将向量空间理论中标量域的角色从域替换为任意含幺环。这一推广看似简单，却打开了通向同调代数、表示论和代数几何的大门。模论的历史可追溯至 Dedekind 和 Kronecker 对理想论的研究，但真正将其确立为独立学科的是 Emmy Noether 在 1920 年代的工作，以及后来 Cartan、Eilenberg 和 Mac Lane 在同调代数中系统使用模论语言。从范畴论的角度看，域 \(F\) 上的向量空间范畴 \(\mathbf{Vect}_F\) 是阿贝尔范畴的典范，而环 \(R\) 上的（左）\(R\)-模范畴 \(R\mathbf{Mod}\) 则是最一般的阿贝尔范畴之一------事实上，Freyd-Mitchell 嵌入定理断言任何小阿贝尔范畴均可嵌入某个模范畴。

模论与向量空间理论的关键区别在于：并非每个模都有基。具有基的模称为自由模，在自由模中，线性无关性和生成性可以像向量空间一样良好地运作。然而，对于一般环 \(R\)，\(R\)-模可能没有基------例如，\(\mathbb{Z}/6\mathbb{Z}\) 作为 \(\mathbb{Z}\)-模是挠模，没有非零自由子模。这一差异催生了投射模、内射模和平坦模三个基本概念，它们分别从不同角度衡量一个模与自由模的''接近程度''。投射模 \(P\) 是自由模的直和项，具有提升性质；内射模 \(I\) 是投射模的对偶概念，具有扩张性质；平坦模 \(M\) 则使得张量积函子 \(M \otimes_R -\) 保持正合性。这三个概念构成了同调代数的基石，Ext 和 Tor 函子正是度量模偏离投射性（或内射性、平坦性）的工具。

模论与表示论有着深刻的联系：群 \(G\) 在域 \(k\) 上的表示本质上就是群代数 \(k[G]\) 上的模。同样，Lie 代数的表示是泛包络代数上的模，而箭图（quiver）的表示是路代数上的模。这种统一的模论视角使得不同背景下看似迥异的表示论问题可以共享相同的同调代数工具。本章将系统建立模的基本理论------子模、商模、同态、同构定理------以及自由模、投射模、内射模和平坦模等核心概念，为后续的同调代数（卷十四）和交换代数深化（卷十二）做好充分准备。

\textbf{定义 64.1}（左 \(R\)-模）：设 \(R\) 为含幺环。一个\textbf{左 \(R\)-模}为阿贝尔群 \((M,+)\) 配备标量乘法 \(R \times M \to M\)，\((r, m) \mapsto rm\)，满足：

\begin{enumerate}
\def\labelenumi{\arabic{enumi}.}
\tightlist
\item
  \(r(m+n) = rm + rn\)
\item
  \((r+s)m = rm + sm\)
\item
  \((rs)m = r(sm)\)
\item
  \(1_R \cdot m = m\)
\end{enumerate}

\textbf{右 \(R\)-模}类似，标量从右侧作用：\((mr)s = m(rs)\)。当 \(R\) 为域时，\(R\)-模即 \(R\) 上的向量空间------模论是将标量从域推广到环的自然框架。

\textbf{定义 64.2}（子模、商模、模同态）：设 \(M\) 为 \(R\)-模。\textbf{子模} \(N \subseteq M\) 是对 \(M\) 的运算封闭的加法子群。\textbf{商模} \(M/N\) 是加法商群配备 \(r(m+N) = rm + N\)。\textbf{模同态} \(f : M \to N\) 是满足 \(f(m_1+m_2) = f(m_1)+f(m_2)\) 和 \(f(rm) = r f(m)\) 的群同态。

\textbf{定理 64.1}（模的同构定理）： 1. （第一同构定理）若 \(f : M \to N\) 为满同态，则 \(M / \ker f \cong N\) 2. （第二同构定理）若 \(A, B \leq M\)，则 \((A+B)/B \cong A/(A \cap B)\) 3. （第三同构定理）若 \(N \leq K \leq M\)，则 \((M/N)/(K/N) \cong M/K\)

\textbf{证明}：与群和向量空间的证明完全相同。以第一同构定理为例：定义 \(\bar{f}(m+\ker f) = f(m)\)，易验证 \(\bar{f}\) 是良定义的模同构（满性由 \(f\) 的满性保证，单性由 \(\ker \bar{f} = \{0\}\) 保证）。\(\blacksquare\)

\subsection{64.2 自由模}\label{ux81eaux7531ux6a21}

\textbf{定义 64.3}（基与自由模）：\(R\)-模 \(F\) 的子集 \(\{e_i\}_{i \in I}\) 称为\textbf{基}，若每个 \(x \in F\) 可唯一表为有限线性组合 \(x = \sum r_i e_i\)（\(r_i \in R\)）。具有基的模称为\textbf{自由模}。典范例子：\(R^{(I)} = \{(r_i)_{i \in I} : \text{仅有限多项非零}\}\)。

\textbf{定理 64.2}（自由模的万有性质）：设 \(F\) 以 \(\{e_i\}\) 为基。对任意 \(R\)-模 \(M\) 及任意映射 \(\varphi : \{e_i\} \to M\)，存在唯一的模同态 \(\tilde{\varphi} : F \to M\) 使 \(\tilde{\varphi}(e_i) = \varphi(e_i)\)。

\textbf{证明}：定义 \(\tilde{\varphi}(\sum r_i e_i) = \sum r_i \varphi(e_i)\)。由基的唯一表示性，\(\tilde{\varphi}\) 良定义且为同态。唯一性由生成性保证。\(\blacksquare\)

\textbf{注记 64.1}（秩的不变性）：对一般环 \(R\)，自由模的两组基未必有相同基数（如非交换环可出现 \(R \cong R^2\) 作为 \(R\)-模）。若 \(R\) 为交换环，自由模的秩（基的基数）是不变的------这通过对极大理想 \(m\) 约化到域 \(R/m\) 上的向量空间维数即得。

\subsection{64.3 投射模、内射模与平坦模}\label{ux6295ux5c04ux6a21ux5185ux5c04ux6a21ux4e0eux5e73ux5766ux6a21}

\textbf{定义 64.4}（投射模）：\(R\)-模 \(P\) 称为\textbf{投射模}，若对任意满同态 \(g : M \to N \to 0\) 及任意同态 \(f : P \to N\)，存在提升 \(\tilde{f} : P \to M\) 使得 \(g \circ \tilde{f} = f\)。

\textbf{定理 64.3}（投射模的等价刻画）：下列条件等价： 1. \(P\) 为投射模 2. \(P\) 为某自由模的直和项 3. 函子 \(\operatorname{Hom}_R(P, -)\) 是正合函子（即保持满同态）

\textbf{证明}：\((1) \Rightarrow (2)\)：取自由模 \(F\) 和满同态 \(\pi : F \to P\)。在提升图中取 \(N = P\)，\(f = \operatorname{id}_P\)，得 \(\tilde{f} : P \to F\) 使 \(\pi \circ \tilde{f} = \operatorname{id}_P\)，故 \(F = \operatorname{im}(\tilde{f}) \oplus \ker \pi\)，\(P\) 为 \(F\) 的直和项。\((2) \Rightarrow (3)\)：\(\operatorname{Hom}\) 函子保持直和和正合性，自由模的 \(\operatorname{Hom}\) 函子正合。\((3) \Rightarrow (1)\)：由正合性定义直接可得。\(\blacksquare\)

\textbf{定义 64.5}（内射模）：\(R\)-模 \(I\) 称为\textbf{内射模}，若对任意单同态 \(g : N \hookrightarrow M\) 及任意同态 \(f : N \to I\)，存在扩张 \(\tilde{f} : M \to I\) 使得 \(\tilde{f} \circ g = f\)。等价地，函子 \(\operatorname{Hom}_R(-, I)\) 是正合函子。

\textbf{定理 64.4}（Baer 判别法）：\(I\) 为内射模 \(\iff\) 对 \(R\) 的每个左理想 \(\mathfrak{a}\) 及同态 \(f : \mathfrak{a} \to I\)，\(f\) 可扩张为 \(\tilde{f} : R \to I\)。

\textbf{证明}：（必要性）显然。（充分性）给定 \(N \subset M\) 和 \(f : N \to I\)。考虑全体扩张对 \((N', f')\)（\(N \subseteq N' \subseteq M\)，\(f'|_{N} = f\)），以包含序偏序。由 Zorn 引理存在极大扩张 \((N_0, f_0)\)。若 \(N_0 \neq M\)，取 \(x \in M \setminus N_0\)，设 \(\mathfrak{a} = \{r \in R : rx \in N_0\}\)。由 Baer 条件将 \(\mathfrak{a} \to N_0 \xrightarrow{f_0} I\) 扩张到 \(R \to I\)，进而扩张到 \(N_0 + Rx\)，矛盾。故 \(N_0 = M\)。\(\blacksquare\)

\textbf{定义 64.6}（平坦模）：右 \(R\)-模 \(M\) 称为\textbf{平坦模}，若函子 \(M \otimes_R -\)（或 \(-\otimes_R M\)）是正合函子------即对任意短正合列 \(0 \to A \to B \to C \to 0\)，\(0 \to M \otimes_R A \to M \otimes_R B \to M \otimes_R C \to 0\) 正合。

\textbf{定理 64.5}：自由模必平坦；投射模必平坦；有限展示的平坦模必投射。

\textbf{证明}：自由模 \(R^{(I)}\) 为平坦因 \(R^{(I)} \otimes_R -\) 同构于直和函子，而直和保正合列。投射模为自由模的直和项，故其张量积函子为直和项的正合函子，从而平坦。设 \(M\) 为有限展示平坦模，取有限生成自由模 \(F\) 及满同态 \(\pi: F \to M\)，\(\ker \pi\) 有限生成。由平坦性，对任意模 \(N\)，\(\operatorname{Tor}_1^R(M, N) = 0\)。设 \(K = \ker \pi\)，则 \(0 = \operatorname{Tor}_1^R(M, F/K) \cong \operatorname{Tor}_1^R(M, M)\) 导致 \(0 \to K \otimes M \to F \otimes M\) 单，从而正合列 \(0 \to K \to F \to M \to 0\) 分裂，\(M\) 为 \(F\) 的直和项，故投射。\(\blacksquare\)

\textbf{定理 64.6}（Tor 函子）：张量积的左导函子 \(\operatorname{Tor}_n^R(-, N)\) 度量了张量积偏离正合的程度。\(M\) 平坦 \(\iff\) \(\operatorname{Tor}_1^R(M, N) = 0\) 对所有 \(N\) 成立。正合列 \(0 \to A \to B \to C \to 0\) 诱导长正合列 \[\cdots \to \operatorname{Tor}_1(M, C) \to M \otimes A \to M \otimes B \to M \otimes C \to 0\]

\textbf{证明}：取 \(M\) 的投射分解 \(P_\bullet \to M \to 0\)，定义 \(\operatorname{Tor}_n^R(M, N) = H_n(P_\bullet \otimes_R N)\)。由蛇引理，短正合列 \(0 \to A \to B \to C \to 0\) 诱导复形正合列 \(0 \to P_\bullet \otimes A \to P_\bullet \otimes B \to P_\bullet \otimes C \to 0\)，取同调长正合列即得 Tor 长正合列。若 \(M\) 平坦，则 \(M \otimes_R -\) 正合，故 \(\operatorname{Tor}_1^R(M, N) = 0\)。反之，若 \(\operatorname{Tor}_1^R(M, N) = 0\) 对所有 \(N\)，则对任意单射 \(A \hookrightarrow B\)，\(M \otimes A \to M \otimes B\) 的核为 \(\operatorname{Tor}_1(M, B/A) = 0\)，故 \(M \otimes_R -\) 保持单射，从而是正合函子。\(\blacksquare\)

\subsection{64.4 Noether 模与 Artin 模}\label{noether-ux6a21ux4e0e-artin-ux6a21}

\textbf{定义 64.7}（Noether 模）：\(M\) 为 \textbf{Noether 模}，若其子模满足升链条件（ACC）：任意升链 \(M_1 \subseteq M_2 \subseteq \cdots\) 必稳定。等价于每个子模有限生成；等价于子模族的每个非空族有极大元。

\textbf{定义 64.8}（Artin 模）：\(M\) 为 \textbf{Artin 模}，若其子模满足降链条件（DCC）：任意降链 \(M_1 \supseteq M_2 \supseteq \cdots\) 必稳定。

\textbf{定理 64.7}：设 \(0 \to A \xrightarrow{\alpha} B \xrightarrow{\beta} C \to 0\) 为短正合列。则 \(B\) 为 Noether（Artin）模 \(\iff\) \(A\) 和 \(C\) 均为 Noether（Artin）模。

\textbf{证明}：仅证 Noether 情形（Artin 类似）。若 \(B\) 为 Noether，\(A\) 的子模链提升至 \(B\) 后稳定，故 \(A\) 为 Noether；\(C\) 的子模对应 \(B\) 中含 \(\alpha(A)\) 的子模，由 Noether 性稳定。反之，若 \(A\) 和 \(C\) 均为 Noether，设 \(B_1 \subseteq B_2 \subseteq \cdots\) 为 \(B\) 的子模升链。\(\beta(B_i)\) 在 \(C\) 中稳定；由 \(\ker(\beta|_{B_i}) \subseteq A\) 的 Noether 性，\(B_i\) 最终稳定。\(\blacksquare\)

\textbf{定理 64.8}（Noether 环）：\(R\) 为左 Noether 环 \(\iff\) \(R\) 作为左 \(R\)-模是 Noether 模。此时有限生成 \(R\)-模的子模也是有限生成的。

\textbf{证明}：(\(\Rightarrow\)) 设 \(R\) 为左 Noether 环，即 \(R\) 的每个左理想都有限生成。\(R\) 作为左 \(R\)-模的子模即左理想，故 \(R\) 是 Noether 模。(\(\Leftarrow\)) 设 \(R\) 作为左 \(R\)-模是 Noether 模，则 \(R\) 的每个左理想有限生成，即 \(R\) 为左 Noether 环。设 \(M\) 为有限生成 \(R\)-模，\(M = \sum_{i=1}^n R x_i\)。对 \(n\) 归纳：\(n=0\) 时平凡。\(n \geq 1\) 时，考虑正合列 \(0 \to N \cap R x_n \to N \to N/(N \cap R x_n) \to 0\)。\(R x_n \cong R/\operatorname{Ann}(x_n)\) 是 Noether 模（\(R\) 的商模），故 \(N \cap R x_n\) 有限生成。\(N/(N \cap R x_n) \hookrightarrow M/R x_n = \sum_{i=1}^{n-1} R \bar{x}_i\)，由归纳假设 \(M/R x_n\) 是 Noether 模，故其子模 \(N/(N \cap R x_n)\) 有限生成。由正合列性质，\(N\) 有限生成。\(\blacksquare\)

\subsection{64.5 主理想整环上有限生成模的结构定理}\label{ux4e3bux7406ux60f3ux6574ux73afux4e0aux6709ux9650ux751fux6210ux6a21ux7684ux7ed3ux6784ux5b9aux7406}

\textbf{定理 64.9}（不变因子形式）：设 \(R\) 为 PID，\(M\) 为有限生成 \(R\)-模。则存在唯一的 \(r \geq 0\) 及 \(R\) 中非零非单位元 \(d_1 \mid d_2 \mid \cdots \mid d_k\)（相伴意义下唯一），使得 \[M \cong R^r \oplus R/(d_1) \oplus R/(d_2) \oplus \cdots \oplus R/(d_k)\] \(r\) 称为\textbf{自由秩}，\(d_1, \ldots, d_k\) 称为\textbf{不变因子}。

\textbf{证明}：设 \(M\) 由 \(n\) 个元素生成，取满同态 \(\varphi : R^n \to M\)。\(\ker \varphi \subseteq R^n\) 为子模。因 \(R\) 为 PID 是 Noether 环，\(\ker \varphi\) 有限生成。将生成元的线性关系写为矩阵，应用 \textbf{Smith 正规形}：存在可逆矩阵 \(P \in GL_n(R)\) 和 \(Q \in GL_m(R)\)（\(m\) 为 \(\ker\varphi\) 生成元个数，可补齐）使得 \(P A Q = \operatorname{diag}(d_1, \ldots, d_k, 0, \ldots, 0)\) 且 \(d_1 \mid d_2 \mid \cdots \mid d_k\)。这意味着存在 \(R^n\) 的基 \(\{e_1', \ldots, e_n'\}\) 使 \(\ker \varphi\) 由 \(d_1 e_1', \ldots, d_k e_k'\) 生成，\(e_{k+1}', \ldots, e_n'\) 不在 \(\ker \varphi\) 中。故 \[M \cong R^n / \ker \varphi \cong R/(d_1) \oplus \cdots \oplus R/(d_k) \oplus R^{n-k}\] 令 \(r = n-k\)。唯一性：\(r = \dim_{\operatorname{Frac}(R)} (M \otimes_R \operatorname{Frac}(R))\)；\(d_1 \cdots d_i\) 由 \(M\) 的 \(i\) 阶初等因子的生成理想确定。\(\blacksquare\)

\textbf{定理 64.10}（初等因子形式）：将每个 \(d_i\) 分解为不可约元幂的乘积 \(d_i = u \prod p_j^{e_{ij}}\)，则由定理 64.9 得 \[M \cong R^r \oplus \bigoplus_{i,j} R/(p_j^{e_{ij}})\] 其中各 \(p_j^{e_{ij}}\) 称为\textbf{初等因子}。每个 \(R/(p^{e})\) 是不可分解挠模。

\textbf{证明}：由定理 64.9，\(M \cong R^r \oplus R/(d_1) \oplus \cdots \oplus R/(d_k)\)，其中 \(d_1 \mid d_2 \mid \cdots \mid d_k\)。对每个 \(d_i\)，在 PID \(R\) 中分解为不可约因子：\(d_i = u_i \prod_{j} p_j^{e_{ij}}\)（\(u_i\) 为单位）。由 \(d_1 \mid d_2 \mid \cdots \mid d_k\) 知 \(e_{1j} \leq e_{2j} \leq \cdots \leq e_{kj}\)。由中国剩余定理， \[R/(d_i) \cong \bigoplus_{j} R/(p_j^{e_{ij}}).\] 代入即得所需分解。\(R/(p^e)\) 的不可分解性由 \(R/(p^e)\) 中非零子模对应 \(R/(p^f)\)（\(f \leq e\)）且这些子模按包含关系构成全序链得出。\(\blacksquare\)

\textbf{推论 64.11}（有限生成阿贝尔群分类）：取 \(R = \mathbb{Z}\)，任意有限生成阿贝尔群同构于 \[G \cong \mathbb{Z}^r \oplus \mathbb{Z}_{d_1} \oplus \cdots \oplus \mathbb{Z}_{d_k},\quad d_1 \mid d_2 \mid \cdots \mid d_k\]

\textbf{证明}：\(\mathbb{Z}\) 为 PID，有限生成阿贝尔群 = 有限生成 \(\mathbb{Z}\)-模。由定理 64.9 即得。\(\blacksquare\)

\textbf{推论 64.12}（Jordan 标准型的模论证明）：设 \(V\) 为域 \(F\) 上 \(n\) 维向量空间，\(T \in \operatorname{End}_F(V)\)。通过 \(f(x) \cdot v = f(T)v\) 使 \(V\) 成为 \(F[x]\)-模。\(F[x]\) 是 PID，\(V\) 为有限生成挠 \(F[x]\)-模。由定理 64.10 \[V \cong \bigoplus_{i} F[x]/(p_i(x)^{e_i})\] 其中 \(p_i(x)\) 为不可约多项式。每个 \(F[x]/(p_i^{e_i})\) 上 \(T\) 的作用对应于有理典范形式的块。若 \(F\) 为代数闭域（如 \(\mathbb{C}\)），\(p_i(x) = x - \lambda_i\) 全为一次式，则 \((x-\lambda_i)^{e_i}\) 的商给出 \(\lambda_i\) 的 \(e_i \times e_i\) Jordan 块。综上即得 Jordan 标准型。\(\blacksquare\)

\subsection{64.6 直和与直积}\label{ux76f4ux548cux4e0eux76f4ux79ef}

\textbf{定义 64.9}（外直和与外直积）：设 \(\{M_i\}_{i \in I}\) 为一族 \(R\)-模。\textbf{外直积} \(\prod_{i \in I} M_i\) 为 Descartes 积配备逐点运算；\textbf{外直和} \(\bigoplus_{i \in I} M_i\) 为其子模，由仅有限多个分量非零的元素构成。\(I\) 有限时二者相同。

\textbf{定义 64.10}（内直和）：\(M\) 的子模族 \(\{N_i\}_{i \in I}\) 构成 \(M\) 的\textbf{内直和分解} \(M = \bigoplus_i N_i\)，若：（1）\(M = \sum_i N_i\)（生成性）；（2）\(N_j \cap \sum_{i \neq j} N_i = 0\) 对所有 \(j\)（独立性）。等价地，典范同态 \(\bigoplus_i N_i \to M\)，\((x_i) \mapsto \sum x_i\) 为同构。

\textbf{定理 64.13}（直和的万有性质）：外直和 \(\bigoplus_i M_i\) 带有典范嵌入 \(\iota_j : M_j \hookrightarrow \bigoplus_i M_i\) 满足：对任意 \(R\)-模 \(N\) 及同态族 \(f_i : M_i \to N\)，存在唯一 \(f : \bigoplus_i M_i \to N\) 使 \(f \circ \iota_i = f_i\)。对偶地，直积带有典范投射且满足对偶万有性质。

\textbf{证明}：对 \((x_i) \in \bigoplus_i M_i\)（仅有限多非零分量），定义 \(f((x_i)) = \sum_i f_i(x_i)\)（有限和，良定义）。则 \(f \circ \iota_j(x_j) = f_j(x_j)\)，故 \(f \circ \iota_j = f_j\)。唯一性：若 \(g\) 也满足 \(g \circ \iota_i = f_i\)，则对任意 \((x_i) = \sum_i \iota_i(x_i)\)（有限和），\(g((x_i)) = \sum_i g(\iota_i(x_i)) = \sum_i f_i(x_i) = f((x_i))\)。对偶地，直积 \(\prod_i M_i\) 带有典范投射 \(\pi_j: \prod_i M_i \to M_j\)，对任意同态族 \(g_i: N \to M_i\)，存在唯一 \(g: N \to \prod_i M_i\) 使 \(\pi_i \circ g = g_i\)，由 \(g(n) = (g_i(n))_i\) 给出。\(\blacksquare\)

\textbf{定理 64.14}（内直和的等价条件）：有限内直和 \(M = \bigoplus_{i=1}^n N_i\) 等价于存在投影族 \(\pi_i : M \to N_i\)（\(\operatorname{id}_M\)-正交幂等元）满足 \(\pi_i|_{N_j} = \delta_{ij} \operatorname{id}\) 且 \(\sum \pi_i = \operatorname{id}_M\)。此时 \(M/N_i \cong \bigoplus_{j \neq i} N_j\)。

\textbf{证明}：若 \(M = \bigoplus_i N_i\)，定义 \(\pi_i: M \to N_i\) 为典范投射：\(\pi_i(\sum_j x_j) = x_i\)（\(x_j \in N_j\)）。则 \(\pi_i|_{N_j} = \delta_{ij} \operatorname{id}\)，\(\sum \pi_i = \operatorname{id}_M\)，且 \(\pi_i \circ \pi_j = \delta_{ij} \pi_i\)（正交幂等）。反之，若存在满足条件的 \(\pi_i\)，定义 \(\varphi: M \to \bigoplus_i N_i\) 为 \(\varphi(x) = (\pi_1(x), \ldots, \pi_n(x))\)，其逆为 \(\psi(x_1, \ldots, x_n) = \sum x_i\)。由 \(\pi_i|_{N_j} = \delta_{ij}\)，\(\psi \circ \varphi = \operatorname{id}_M\)；由 \(\sum \pi_i = \operatorname{id}_M\)，\(\varphi \circ \psi = \operatorname{id}\)。故 \(M \cong \bigoplus_i N_i\)。同构 \(M/N_i \cong \bigoplus_{j \neq i} N_j\) 由映射 \(x + N_i \mapsto (\pi_j(x))_{j \neq i}\) 给出。\(\blacksquare\)

\subsection{64.7 正合列}\label{ux6b63ux5408ux5217}

\textbf{定义 64.11}（正合列）：模同态序列 \(\cdots \to M_{i-1} \xrightarrow{f_{i-1}} M_i \xrightarrow{f_i} M_{i+1} \to \cdots\) 在 \(M_i\) 处\textbf{正合}若 \(\operatorname{im} f_{i-1} = \ker f_i\)。序列正合若处处正合。\textbf{短正合列}为 \(0 \to A \xrightarrow{\alpha} B \xrightarrow{\beta} C \to 0\) 的正合列（\(\alpha\) 单，\(\beta\) 满，\(\operatorname{im}\alpha = \ker\beta\)）。

\textbf{定义 64.12}（分裂正合列）：短正合列 \(0 \to A \xrightarrow{\alpha} B \xrightarrow{\beta} C \to 0\) 称为\textbf{分裂}的若满足以下等价条件之一： 1. 存在 \(r : B \to A\)（收缩）使 \(r \circ \alpha = \operatorname{id}_A\) 2. 存在 \(s : C \to B\)（截面）使 \(\beta \circ s = \operatorname{id}_C\) 3. \(B \cong A \oplus C\)，且 \(\alpha\) 对应嵌入、\(\beta\) 对应投射

\textbf{证明}：（1）\(\Rightarrow\)（3）：定义 \(\varphi(b) = (r(b), \beta(b))\)。由 \(\ker\beta = \operatorname{im}\alpha\) 及 \(r \circ \alpha = \operatorname{id}_A\)，\(\varphi\) 为同构，逆为 \((a, c) \mapsto \alpha(a) + s(c)\)（\(s\) 由 \(\beta\) 的截面性质构造）。（2）\(\Rightarrow\)（3）对偶。（3）\(\Rightarrow\)（1）（2）自明。\(\blacksquare\)

\textbf{定理 64.15}（分裂判别条件）：若 \(C\) 为投射模，则任意 \(0 \to A \to B \to C \to 0\) 分裂（提升 \(\operatorname{id}_C\) 得截面）。对偶地，若 \(A\) 为内射模，则分裂。特别地，域上向量空间的任意短正合列均分裂。

\textbf{证明}：设 \(0 \to A \xrightarrow{\alpha} B \xrightarrow{\beta} C \to 0\) 正合，\(C\) 投射。由投射性，对 \(\operatorname{id}_C: C \to C\) 及满同态 \(\beta: B \to C\)，存在 \(s: C \to B\) 使 \(\beta \circ s = \operatorname{id}_C\)（\(s\) 为截面）。由分裂正合列的等价条件，此正合列分裂，\(B \cong A \oplus C\)。对偶地，若 \(A\) 内射，\(\operatorname{id}_A\) 沿单射 \(\alpha\) 延拓得收缩 \(r: B \to A\)（\(r \circ \alpha = \operatorname{id}_A\)），从而分裂。域上向量空间均为自由模（故投射），且任意短正合列均分裂，因为任意子空间有补空间。\(\blacksquare\)

\textbf{定理 64.16}（正合列与平坦模的 Tor 判别）：\(M\) 平坦 \(\iff\) \(\operatorname{Tor}_1^R(M, N) = 0\) 对所有 \(N\) 成立 \(\iff\) \(\operatorname{Tor}_1^R(M, R/\mathfrak{a}) = 0\) 对所有有限生成理想 \(\mathfrak{a}\) 成立。投射模必平坦；Noether 环上有限生成平坦模必投射。长正合列由短正合列的拼接得到，Tor 函子衡量张量积偏离正合的程度： \[\cdots \to \operatorname{Tor}_1(M, C) \to M \otimes A \to M \otimes B \to M \otimes C \to 0\]

\textbf{证明}：\(M\) 平坦 \(\iff\) \(\operatorname{Tor}_1^R(M, N) = 0\) 对所有 \(N\) 已在定理 64.6 中证明。\(\Leftarrow\) 方向：设 \(\operatorname{Tor}_1^R(M, R/\mathfrak{a}) = 0\) 对所有有限生成理想 \(\mathfrak{a}\)。对任意模 \(N\)，\(N\) 是其有限生成子模的归纳极限，Tor 与归纳极限交换，故 \(\operatorname{Tor}_1^R(M, N) = 0\)。投射模平坦：投射模为自由模直和项，自由模的张量积保正合性。Noether 环上有限生成平坦模的投射性：取 \(M\) 的有限生成自由分解 \(F_1 \to F_0 \to M \to 0\)，由平坦性此分解在局部化下分裂，故 \(M\) 为 \(F_0\) 的直和项，从而是投射的。\(\blacksquare\)

\section{第118章：导出函子与Tor/Ext}\label{ux7b2c118ux7ae0ux5bfcux51faux51fdux5b50ux4e0etorext}

同调代数的核心思想是：将模论中非正合序列的''偏离正合度''量化为不变量。这一思想源于Cartan和Eilenberg在1956年划时代的著作《同调代数》，其基本工具是投射分解、内射分解以及由此构造的导出函子。Tor函子度量张量积偏离正合的程度------它是左正合的张量积函子的左导出；Ext函子度量Hom函子偏离正合的程度------它是左正合Hom函子的右导出。本章在模论基本概念的基础上，系统构建导出函子的理论框架，包括投射和内射分解的存在性与比较定理、蛇引理与连接同态的构造、以及Tor和Ext的核心性质与应用。这些工具不仅统一了代数拓扑中的同调群理论（单纯同调、奇异同调、上同调），也为代数几何中的层上同调（Grothendieck的Tohoku纲领）和表示论中的群上同调提供了统一的语言。

\subsection{65.1 投射分解与内射分解}\label{ux6295ux5c04ux5206ux89e3ux4e0eux5185ux5c04ux5206ux89e3}

\textbf{定义 65.1}（投射分解与内射分解）：设\(M\)为\(R\)-模。\(M\)的\textbf{投射分解}为正合列 \[\cdots \to P_2 \xrightarrow{d_2} P_1 \xrightarrow{d_1} P_0 \xrightarrow{\varepsilon} M \to 0,\] 其中每个\(P_i\)为投射模。对偶地，\(M\)的\textbf{内射分解}为正合列 \[0 \to M \xrightarrow{\eta} I^0 \xrightarrow{d^0} I^1 \xrightarrow{d^1} I^2 \to \cdots,\] 其中每个\(I^i\)为内射模。

\textbf{定理 65.1}（投射分解的存在性）：任意\(R\)-模\(M\)均存在投射分解。

\textbf{证明}：构造自由模\(F_0\)和满同态\(\varepsilon: F_0 \to M\)（取\(M\)的生成元集\(\{x_i\}\)，定义\(F_0 = R^{(I)}\)，\(\varepsilon(e_i) = x_i\)）。设\(K_0 = \ker\varepsilon\)。递归地，对已构造的\(P_i \to K_{i-1}\)（\(K_{-1}=M\)），取自由模\(P_{i+1}\)和满同态\(P_{i+1} \to K_i\)，定义\(K_{i+1} = \ker(P_{i+1} \to K_i)\)。由于自由模是投射的，由此得到的序列\(\cdots \to P_2 \to P_1 \to P_0 \to M \to 0\)即为投射分解。每一步中映射的具体定义保证了正合性。此构造说明每个模都可选择自由分解（更特殊的投射分解）。\(\blacksquare\)

\textbf{定理 65.2}（投射分解的比较定理）：设\(\varphi: M \to N\)为模同态，\(P_\bullet \to M \to 0\)为\(M\)的投射分解，\(Q_\bullet \to N \to 0\)为\(N\)的任意分解（即序列正合于\(N\)处，各项不必投射）。则存在链映射\(\{f_n: P_n \to Q_n\}\)提升\(\varphi\)，且该提升在同伦意义下唯一（即任意两个提升链同伦）。

\textbf{证明}：归纳构造\(f_n\)。\(n=0\)时：需构造\(f_0\)使下图交换： \[\begin{CD} P_0 @>\varepsilon>> M @>>> 0\\ @. @V\varphi VV \\ Q_0 @>\eta>> N @>>> 0 \end{CD}\] 由\(P_0\)的投射性及\(\eta\)的满性（正合性），存在\(f_0: P_0 \to Q_0\)使\(\eta \circ f_0 = \varphi \circ \varepsilon\)。设\(f_0, \ldots, f_{n-1}\)已构造。考虑： \[\begin{CD} P_n @>d_n>> P_{n-1} @>d_{n-1}>> P_{n-2}\\ @. @Vf_{n-1}VV @Vf_{n-2}VV \\ Q_n @>\delta_n>> Q_{n-1} @>\delta_{n-1}>> Q_{n-2} \end{CD}\] （\(d_n, \delta_n\)为各分解的微分）。需构造\(f_n\)使\(f_{n-1} \circ d_n = \delta_n \circ f_n\)。由归纳假设\(f_{n-2} \circ d_{n-1} = \delta_{n-1} \circ f_{n-1}\)，知\(\delta_{n-1} \circ f_{n-1} \circ d_n = f_{n-2} \circ d_{n-1} \circ d_n = 0\)（分解正合性），故\(f_{n-1} \circ d_n(P_n) \subseteq \ker\delta_{n-1} = \operatorname{im}\delta_n\)。由\(P_n\)的投射性，存在\(f_n: P_n \to Q_n\)使\(\delta_n \circ f_n = f_{n-1} \circ d_n\)。同伦唯一性：若\(\{f_n\}\)和\(\{g_n\}\)均为提升，构造链同伦\(s_n: P_n \to Q_{n+1}\)归纳地满足\(f_n - g_n = \delta_{n+1} \circ s_n + s_{n-1} \circ d_n\)（\(s_{-1}=0\)）。归纳中利用\(Q_\bullet\)的正合性和\(P_n\)的投射性。\(\blacksquare\)

\textbf{定理 65.3}（内射分解的存在性与对偶比较定理）：任意\(R\)-模均存在内射分解。对偶地，内射分解的比较定理断言：任意同态\(\varphi: M \to N\)可提升为内射分解之间的链映射，且该提升在同伦意义下唯一。

\textbf{证明}（内射分解的存在性）：需用到每个模均可嵌入某内射模（Baer定理的推论：\(R\)-模范畴有足够多内射对象）。对\(M\)，取内射模\(I^0\)及单射\(\eta: M \hookrightarrow I^0\)。设\(C^0 = \operatorname{coker}\eta\)，取内射模\(I^1\)及单射\(C^0 \hookrightarrow I^1\)，定义\(d^0 = (M \to I^0 \to C^0 \to I^1)\)。递归构造即得内射分解。Baer判别法确保了足够多内射模的存在性：对于任意模\(M\)，取\(M\)到其Pontryagin对偶的嵌入，结合足够多内射模的Grothendieck构造即得。\(\blacksquare\)

\subsection{65.2 导出函子}\label{ux5bfcux51faux51fdux5b50}

\textbf{定义 65.2}（左导出函子与右导出函子）：设\(F: R\mathbf{Mod} \to S\mathbf{Mod}\)为右正合函子。\(F\)的\textbf{左导出函子}\(L_nF: R\mathbf{Mod} \to S\mathbf{Mod}\)定义为：对任意模\(M\)，取投射分解\(P_\bullet \to M \to 0\)，令 \[(L_nF)(M) = H_n(F(P_\bullet)).\] 对偶地，若\(F\)为左正合函子，\(F\)的\textbf{右导出函子}\(R^nF\)定义为：取内射分解\(0 \to M \to I^\bullet\)，令 \[(R^nF)(M) = H^n(F(I^\bullet)).\]

\textbf{定理 65.4}（导出函子的良定义性）：导出函子不依赖于所选分解（在同构意义下）。即若\(P_\bullet\)和\(Q_\bullet\)均为\(M\)的投射分解，则\(H_n(F(P_\bullet)) \cong H_n(F(Q_\bullet))\)对任意加法函子\(F\)成立。

\textbf{证明}：由比较定理，恒等映射\(\operatorname{id}_M: M \to M\)可提升为链映射\(f_\bullet: P_\bullet \to Q_\bullet\)和\(g_\bullet: Q_\bullet \to P_\bullet\)。\(g_\bullet \circ f_\bullet: P_\bullet \to P_\bullet\)和\(\operatorname{id}_{P_\bullet}\)均为\(\operatorname{id}_M\)的提升，由同伦唯一性，\(g_\bullet \circ f_\bullet \simeq \operatorname{id}_{P_\bullet}\)（链同伦）。同理\(f_\bullet \circ g_\bullet \simeq \operatorname{id}_{Q_\bullet}\)。由于任何加法函子保持链同伦（\(F(s_\bullet)\)为\(F(f_\bullet)\)与\(F(g_\bullet)\)之间的链同伦），故\(F(f_\bullet)\)和\(F(g_\bullet)\)诱导同调群的互逆同构，\(H_n(F(P_\bullet)) \cong H_n(F(Q_\bullet))\)。\(\blacksquare\)

\subsection{65.3 Tor函子的构造与长正合列}\label{torux51fdux5b50ux7684ux6784ux9020ux4e0eux957fux6b63ux5408ux5217}

\textbf{定理 65.5}（蛇引理与连接同态的完整证明）：设\(R\mathbf{Mod}\)中下图行正合： \[\begin{CD} 0 @>>> A @>f>> B @>g>> C @>>> 0\\ @. @V\alpha VV @V\beta VV @V\gamma VV \\ 0 @>>> A' @>f'>> B' @>g'>> C' @>>> 0 \end{CD}\] 则存在典范连接同态\(\delta: \ker\gamma \to \operatorname{coker}\alpha\)，使得下图（蛇形）正合： \[\ker\alpha \to \ker\beta \to \ker\gamma \xrightarrow{\delta} \operatorname{coker}\alpha \to \operatorname{coker}\beta \to \operatorname{coker}\gamma.\]

\textbf{证明}：第一步（构造\(\delta\)）：取\(x \in \ker\gamma \subseteq C\)。因\(g\)为满射，存在\(b \in B\)使\(g(b) = x\)。由图的交换性，\(g'(\beta(b)) = \gamma(g(b)) = \gamma(x) = 0\)，故\(\beta(b) \in \ker g' = \operatorname{im}f'\)。由\(f'\)为单射，存在唯一\(a' \in A'\)使\(f'(a') = \beta(b)\)。定义\(\delta(x) = a' + \operatorname{im}\alpha \in \operatorname{coker}\alpha\)。

第二步（良定义性）：需验证\(\delta(x)\)不依赖于\(b\)的选择。若\(g(b_1) = g(b_2) = x\)，则\(b_1 - b_2 \in \ker g = \operatorname{im}f\)，存在\(a \in A\)使\(f(a) = b_1 - b_2\)。由交换性，\(\beta(b_1) - \beta(b_2) = \beta(f(a)) = f'(\alpha(a))\)。若\(f'(a_1') = \beta(b_1)\)和\(f'(a_2') = \beta(b_2)\)，则\(f'(a_1' - a_2' - \alpha(a)) = 0\)。由\(f'\)为单射，\(a_1' - a_2' = \alpha(a) \in \operatorname{im}\alpha\)，故在\(\operatorname{coker}\alpha\)中相同。

第三步（正合性验证）：逐项验证：(i) \(\operatorname{im}(\ker\alpha \to \ker\beta) = \ker(\ker\beta \to \ker\gamma)\)：由原序列正合性直接推出。(ii) \(\operatorname{im}(\ker\beta \to \ker\gamma) = \ker\delta\)：若\(x = g(b)\)且\(b \in \ker\beta\)，则\(\beta(b)=0\)，故\(f'(a') = 0\)，由单性\(a'=0\)，\(\delta(x)=0\)。反之若\(\delta(x)=0\)，则\(a' = \alpha(a)\)，从而\(\beta(b - f(a)) = \beta(b) - f'(\alpha(a)) = f'(a') - f'(a') = 0\)，故\(b - f(a) \in \ker\beta\)且\(g(b - f(a)) = g(b) = x\)。(iii) \(\operatorname{im}\delta = \ker(\operatorname{coker}\alpha \to \operatorname{coker}\beta)\)和(iv)后两项的正合性类似验证。\(\blacksquare\)

\textbf{定理 65.6}（Tor的长正合列）：设\(0 \to A \to B \to C \to 0\)为\(R\)-模的短正合列，\(M\)为右\(R\)-模。则有长正合列 \[\cdots \to \operatorname{Tor}_2^R(M, C) \to \operatorname{Tor}_1^R(M, A) \to \operatorname{Tor}_1^R(M, B) \to \operatorname{Tor}_1^R(M, C) \to M \otimes_R A \to M \otimes_R B \to M \otimes_R C \to 0.\]

\textbf{证明}：取\(M\)的投射分解\(P_\bullet \to M \to 0\)。由短正合列\(0 \to A \to B \to C \to 0\)，因为每个\(P_n\)为投射模（故平坦），序列 \[0 \to P_n \otimes A \to P_n \otimes B \to P_n \otimes C \to 0\] 正合（平坦模保持单射）。这给出复形的短正合列\(0 \to P_\bullet \otimes A \to P_\bullet \otimes B \to P_\bullet \otimes C \to 0\)。取同调长正合列（此为蛇引理在复形范畴的直接推论），即得Tor的长正合列。显式地，连接同态\(\operatorname{Tor}_1(M, C) \to M \otimes A\)由蛇引理应用于\(P_1 \otimes (-)\)和\(P_0 \otimes (-)\)的交换图得出。\(\blacksquare\)

\textbf{定理 65.7}（Tor函子的对称性）：\(\operatorname{Tor}_n^R(M, N) \cong \operatorname{Tor}_n^R(N, M)\)（当\(R\)为交换环时）。

\textbf{证明}：取\(M\)的投射分解\(P_\bullet \to M \to 0\)和\(N\)的投射分解\(Q_\bullet \to N \to 0\)。构造双重复形\(P_\bullet \otimes Q_\bullet\)，其全复形的同调可通过两种方式计算：先对\(p\)取同调得\(H_p(P_\bullet \otimes Q_q) = \operatorname{Tor}_p(M, Q_q)\)，再对\(q\)取同调；或先对\(q\)取同调再对\(p\)。两种计算给出同构的极限，由谱序列的比较定理即得\(\operatorname{Tor}_n(M, N) \cong \operatorname{Tor}_n(N, M)\)。\(\blacksquare\)

\subsection{65.4 Ext函子的构造与扩张分类}\label{extux51fdux5b50ux7684ux6784ux9020ux4e0eux6269ux5f20ux5206ux7c7b}

\textbf{定理 65.8}（Ext函子的构造）：右导出函子\(\operatorname{Ext}_R^n(M, -)\)定义为左正合Hom函子\(\operatorname{Hom}_R(M, -)\)的右导出。即取\(N\)的内射分解\(0 \to N \to I^\bullet\)，\(\operatorname{Ext}_R^n(M, N) = H^n(\operatorname{Hom}_R(M, I^\bullet))\)。等价地，\(\operatorname{Ext}_R^n(-, N)\)为逆变函子\(\operatorname{Hom}_R(-, N)\)的右导出：取\(M\)的投射分解\(P_\bullet \to M \to 0\)，\(\operatorname{Ext}_R^n(M, N) = H^n(\operatorname{Hom}_R(P_\bullet, N))\)。两种定义一致。

\textbf{证明}（等价性）：取\(M\)的投射分解\(P_\bullet \to M \to 0\)和\(N\)的内射分解\(0 \to N \to I^\bullet\)。构造双重复形\(K^{p,q} = \operatorname{Hom}_R(P_p, I^q)\)。该双重复形的两个谱序列分别收敛到\(H^n(\operatorname{Hom}_R(P_\bullet, N))\)和\(H^n(\operatorname{Hom}_R(M, I^\bullet))\)。由谱序列的比较得二者同构。具体地，先对\(q\)取上同调，\(H^q(\operatorname{Hom}_R(P_p, I^\bullet)) = \operatorname{Ext}_R^q(P_p, N)\)。由于\(P_p\)为投射模，\(\operatorname{Ext}_R^q(P_p, N) = 0\)对\(q > 0\)，而对\(q=0\)得\(\operatorname{Hom}_R(P_p, N)\)。再对\(p\)取上同调即得\(H^n(\operatorname{Hom}_R(P_\bullet, N))\)。另一方向类似。\(\blacksquare\)

\textbf{定理 65.9}（Ext与扩张的分类）：\(\operatorname{Ext}_R^1(M, N)\)与\(N\)经\(M\)的模扩张的同构类一一对应。即\(\operatorname{Ext}_R^1(M, N)\)分类了所有短正合列\(0 \to N \to X \to M \to 0\)（模\(N\)和\(M\)固定）的同构类，其群结构对应于Baer和。

\textbf{证明}：设短正合列\(\xi: 0 \to N \to X \to M \to 0\)代表一个扩张。取\(M\)的投射表示（即投射分解的前两项）\(P_1 \xrightarrow{d_1} P_0 \to M \to 0\)。由\(P_0\)的投射性，存在提升\(P_0 \to X\)；由\(P_1\)的投射性及正合性，存在提升\(P_1 \to N\)使得下图交换。该提升诱导一个代表元\([f] \in \operatorname{Hom}_R(P_1, N) / \operatorname{im}(d_1^*: \operatorname{Hom}_R(P_0, N) \to \operatorname{Hom}_R(P_1, N)) = \operatorname{Ext}_R^1(M, N)\)。

反之，给定代表元\(f: P_1 \to N\)（满足\(f|_{P_1, \text{边界}} = 0\)），构造推出（push-out）： \[\begin{CD} 0 @>>> P_1 @>d_1>> P_0 @>>> M @>>> 0\\ @. @VfVV @VVV @| \\ 0 @>>> N @>>> X @>>> M @>>> 0 \end{CD}\] 其中\(X = (N \oplus P_0) / \{(f(p), -d_1(p)) : p \in P_1\}\)。此为\(N\)经\(M\)的扩张。这两个映射互逆，建立双射\(\operatorname{Ext}_R^1(M, N) \leftrightarrow \{\text{扩张同构类}\}\)。Baer和对应Ext群中的加法。\(\blacksquare\)

\subsection{65.5 Tor和Ext的计算实例}\label{torux548cextux7684ux8ba1ux7b97ux5b9eux4f8b}

\textbf{例65.1}（\(\mathbb{Z}\)-模）：\(R = \mathbb{Z}\)时，\(\mathbb{Z}\)-模即Abel群。

\begin{enumerate}
\def\labelenumi{\arabic{enumi}.}
\item
  \textbf{Tor群}：\(\operatorname{Tor}_1^{\mathbb{Z}}(\mathbb{Z}/m\mathbb{Z}, \mathbb{Z}/n\mathbb{Z}) \cong \mathbb{Z}/\gcd(m,n)\mathbb{Z}\)，且\(\operatorname{Tor}_n^{\mathbb{Z}}(-,-) = 0\)对所有\(n \geq 2\)（\(\mathbb{Z}\)的整体维数为1）。取\(\mathbb{Z}/m\mathbb{Z}\)的投射分解\(0 \to \mathbb{Z} \xrightarrow{\times m} \mathbb{Z} \to \mathbb{Z}/m\mathbb{Z} \to 0\)，与\(\mathbb{Z}/n\mathbb{Z}\)做张量积得\(0 \to \mathbb{Z}/n\mathbb{Z} \xrightarrow{\times m} \mathbb{Z}/n\mathbb{Z} \to \mathbb{Z}/m\mathbb{Z} \otimes \mathbb{Z}/n\mathbb{Z} \to 0\)。核为\(\{x \in \mathbb{Z}/n\mathbb{Z} : mx = 0\} \cong \mathbb{Z}/\gcd(m,n)\mathbb{Z}\)，故\(\operatorname{Tor}_1 = \mathbb{Z}/\gcd(m,n)\mathbb{Z}\)。
\item
  \textbf{Ext群}：\(\operatorname{Ext}_{\mathbb{Z}}^1(\mathbb{Z}/m\mathbb{Z}, \mathbb{Z}) \cong \mathbb{Z}/m\mathbb{Z}\)（取\(\mathbb{Z}/m\mathbb{Z}\)的投射分解，计算\(\operatorname{Hom}_{\mathbb{Z}}(-,\mathbb{Z})\)的上同调）。\(\operatorname{Ext}_{\mathbb{Z}}^1(\mathbb{Z}/m\mathbb{Z}, \mathbb{Z}/n\mathbb{Z}) \cong \mathbb{Z}/\gcd(m,n)\mathbb{Z}\)（与Tor相同，但原因不同------这是有限生成Abel群的巧合）。
\item
  \textbf{\(\mathbb{Q}\)作为\(\mathbb{Z}\)-模}：\(\mathbb{Q}\)是平坦的（无挠Abel群均平坦）但非投射（\(\mathbb{Q}\)不是自由Abel群的直和项）。\(\operatorname{Tor}_1^{\mathbb{Z}}(\mathbb{Q}, \mathbb{Z}/n\mathbb{Z}) = 0\)验证了平坦性。
\end{enumerate}

\textbf{例65.2}（多项式环）：\(R = k[x]\)（\(k\)为域）。\(k[x]\)的整体维数为1（PID），故\(\operatorname{Tor}_n = \operatorname{Ext}^n = 0\)（\(n \geq 2\)）。\(R/(x)\)的投射分解为\(0 \to k[x] \xrightarrow{\times x} k[x] \to k[x]/(x) \to 0\)。由此可得\(\operatorname{Tor}_1^{k[x]}(k[x]/(x), k[x]/(x)) \cong k\)，\(\operatorname{Ext}_{k[x]}^1(k[x]/(x), k) \cong k\)。

\hrulefill

\section{第119章：局部化与平坦模}\label{ux7b2c119ux7ae0ux5c40ux90e8ux5316ux4e0eux5e73ux5766ux6a21}

局部化是交换代数中最基本和最强大的构造技术之一。它与几何中''在一点附近研究空间''的思想完全平行：环\(R\)对乘性子集\(S\)的局部化\(S^{-1}R\)对应于将\(\operatorname{Spec}R\)的注意力聚焦到\(S\)中元素不消失的开集上。从模论的角度看，局部化函子\(S^{-1}(-)\)是正合的，这一事实使得许多同调问题可以在局部化后极大地简化。平坦模的概念最初正是从张量积的局部正合性中提炼出来的------平坦模是那些使张量积保持正合的模，它们在代数几何中对应着''平坦族''的概念（即概形之间的平坦态射）。本章将建立局部化的完整代数理论，证明平坦模的几个核心等价刻画（Tor判别准则、理想判别准则），并讨论忠实平坦下降（faithfully flat descent）------这是Grothendieck下降理论的基石，在etale上同调和模空间理论中发挥着根本作用。

\subsection{66.1 分式模的构造}\label{ux5206ux5f0fux6a21ux7684ux6784ux9020}

\textbf{定义 66.1}（乘性子集与分式环）：设\(R\)为交换环。子集\(S \subseteq R\)为\textbf{乘性子集}若\(1 \in S\)且\(s, t \in S \Rightarrow st \in S\)。分式环\(S^{-1}R\)定义为等价类\(r/s\)（\(r \in R, s \in S\)）的集合，模等价关系\(r/s \sim r'/s' \iff \exists t \in S: t(rs' - r's) = 0\)。加法和乘法分别定义为\(r/s + r'/s' = (rs' + r's)/(ss')\)和\((r/s)(r'/s') = (rr')/(ss')\)。

\textbf{定义 66.2}（分式模）：设\(M\)为\(R\)-模，\(S \subseteq R\)为乘性子集。分式模\(S^{-1}M\)类似定义为等价类\(m/s\)（\(m \in M, s \in S\)），模等价关系\(m/s \sim m'/s' \iff \exists t \in S: t(s'm - sm') = 0\)。\(S^{-1}M\)是\(S^{-1}R\)-模（标量乘法：\((r/s)(m/t) = (rm)/(st)\)）。

\textbf{定理 66.1}（局部化的万有性质）：对任意\(R\)-模同态\(f: M \to N\)，若\(S\)中元素在\(N\)上作用为自同构（即对\(s \in S\)，\(s\cdot(-): N \to N\)为同构），则存在唯一的\(S^{-1}R\)-模同态\(\tilde{f}: S^{-1}M \to N\)使\(M \xrightarrow{\iota} S^{-1}M \xrightarrow{\tilde{f}} N\)交换，其中\(\iota(m) = m/1\)。

\textbf{证明}：定义\(\tilde{f}(m/s) = s^{-1} \cdot f(m)\)（\(s^{-1}\)为\(s\)在\(N\)上作用的逆------由假设良定义）。需验证良定义性：若\(m/s \sim m'/s'\)，存在\(t \in S\)使\(t(s'm - sm') = 0\)，则\(f(t(s'm - sm')) = 0\)。因\(t\)在\(N\)上为自同构，\((s')^{-1}f(m) = s^{-1}f(m')\)（通过乘以\(t^{-1}\)并使用\(f\)的线性性）。故\(\tilde{f}\)良定义。唯一性由\(\tilde{f}(m/s) = s^{-1}\tilde{f}(m/1) = s^{-1}f(m)\)必然。\(\blacksquare\)

\subsection{66.2 局部正合性与平坦模的等价刻画}\label{ux5c40ux90e8ux6b63ux5408ux6027ux4e0eux5e73ux5766ux6a21ux7684ux7b49ux4ef7ux523bux753b}

\textbf{定理 66.2}（局部化函子的正合性）：局部化函子\(S^{-1}(-): R\mathbf{Mod} \to S^{-1}R\mathbf{Mod}\)是正合函子。即：若\(0 \to A \xrightarrow{\alpha} B \xrightarrow{\beta} C \to 0\)为\(R\)-模短正合列，则\(0 \to S^{-1}A \to S^{-1}B \to S^{-1}C \to 0\)为\(S^{-1}R\)-模短正合列。

\textbf{证明}：先证\(S^{-1}(-)\)保持单射。设\(\alpha: A \to B\)为单射，需证\(S^{-1}\alpha: S^{-1}A \to S^{-1}B\)为单射。若\((S^{-1}\alpha)(a/s) = \alpha(a)/s = 0\)在\(S^{-1}B\)中，则存在\(t \in S\)使\(t\alpha(a) = \alpha(ta) = 0\)在\(B\)中。由\(\alpha\)为单射，\(ta = 0\)在\(A\)中，故\(a/s = 0\)在\(S^{-1}A\)中。

再证保持在\(C\)处的正合性。显然\(S^{-1}\beta \circ S^{-1}\alpha = S^{-1}(\beta \circ \alpha) = 0\)，故\(\operatorname{im}S^{-1}\alpha \subseteq \ker S^{-1}\beta\)。反之，若\(\beta(b)/s = 0\)在\(S^{-1}C\)中，则存在\(t \in S\)使\(t\beta(b) = \beta(tb) = 0\)在\(C\)中。由原序列正合，存在\(a \in A\)使\(\alpha(a) = tb\)。故\(b/s = \alpha(a)/(ts) = (S^{-1}\alpha)(a/(ts)) \in \operatorname{im}S^{-1}\alpha\)。满性类似可证。\(\blacksquare\)

\textbf{定理 66.3}（平坦模的等价刻画------完整证明）：设\(M\)为\(R\)-模。下列条件等价： 1. \(M\)是平坦\(R\)-模（即\(M \otimes_R -\)正合）； 2. \(\operatorname{Tor}_1^R(M, N) = 0\)对所有\(R\)-模\(N\)； 3. 对\(R\)的任意有限生成理想\(\mathfrak{a}\)，自然映射\(\mathfrak{a} \otimes_R M \to M\)（\(a \otimes m \mapsto am\)）为单射； 4. 对\(R\)的任意理想\(\mathfrak{a}\)，\(\mathfrak{a} \otimes_R M \to M\)为单射； 5. 对任意\(R\)-模同态\(f: N' \to N\)（\(N'\)有限生成），若\(f \otimes \operatorname{id}_M = 0\)，则\(f\)可因子通过某平坦的有限展示模。

\textbf{证明}：（1）\(\iff\)（2）：已在定理64.6中证明。\(M\)平坦即\(M \otimes_R -\)正合，等价于其左导出\(\operatorname{Tor}_1^R(M, -) = 0\)。

（1）\(\Rightarrow\)（4）：设\(\mathfrak{a} \subseteq R\)为理想，有正合列\(0 \to \mathfrak{a} \to R \to R/\mathfrak{a} \to 0\)。作\(M \otimes_R -\)得\(0 \to \mathfrak{a} \otimes_R M \to R \otimes_R M \cong M \to (R/\mathfrak{a}) \otimes_R M \to 0\)。由于\(M\)平坦，此序列正合，故\(\mathfrak{a} \otimes_R M \to M\)为单射。

（4）\(\Rightarrow\)（3）：平凡。

（3）\(\Rightarrow\)（1）：需证\(M \otimes_R -\)保持单射。设\(\iota: N' \hookrightarrow N\)为任意单射。先设\(N\)有限生成，则\(N\)可表示为有限生成自由模的商。由Noether正规化引理可约化到\(N\)为有限生成模、\(N'\)为有限生成子模的情形。对\(N/\iota(N')\)的生成元个数作归纳。基始：\(N = N'\)平凡。归纳步：设\(N/N'\)由一个元素\(\bar{x}\)生成，即\(N = N' + Rx\)。令\(\mathfrak{a} = \{r \in R : rx \in N'\}\)（\(R\)的理想）。构造行正合的交换图： \[\begin{CD} 0 @>>> \mathfrak{a} @>>> R @>>> R/\mathfrak{a} @>>> 0\\ @. @VVV @VVV @| \\ 0 @>>> N' @>>> N @>>> N/N' @>>> 0 \end{CD}\] 作\(-\otimes_R M\)得： \[\begin{CD} \mathfrak{a} \otimes M @>>> M @>>> (R/\mathfrak{a}) \otimes M @>>> 0\\ @VVV @VVV @| \\ N' \otimes M @>>> N \otimes M @>>> (N/N') \otimes M @>>> 0 \end{CD}\] 由条件(3)，\(\mathfrak{a} \otimes M \to M\)单。由归纳假设（\(N'\)和\(R/\mathfrak{a}\)满足适当有限性），其余垂直箭头在适当条件下保持单性。通过图追踪（类似于蛇引理的变体）可证\(N' \otimes M \to N \otimes M\)单。对于一般\(N\)，\(N\)为其有限生成子模的归纳极限，归纳极限保持正合性，故一般情形也成立。

（3）\(\Rightarrow\)（2）：由条件(3)，对任意理想\(\mathfrak{a}\)，\(0 \to \mathfrak{a} \otimes M \to M\)正合。对任意模\(N\)，\(\operatorname{Tor}_1^R(M, N) = 0\)（利用定理64.16中的Tor判别准则------\(\operatorname{Tor}_1^R(M, R/\mathfrak{a}) = 0\)对所有有限生成理想\(\mathfrak{a}\)成立，且Tor与归纳极限交换）。\(\blacksquare\)

\subsection{66.3 忠实平坦下降}\label{ux5fe0ux5b9eux5e73ux5766ux4e0bux964d}

\textbf{定义 66.3}（忠实平坦模）：\(R\)-模\(M\)称为\textbf{忠实平坦}的，若\(M\)平坦且满足附加条件：对任意\(R\)-模\(N\)，\(M \otimes_R N = 0 \Rightarrow N = 0\)。等价地，短正合列\(0 \to A \to B \to C \to 0\)正合当且仅当\(0 \to M \otimes A \to M \otimes B \to M \otimes C \to 0\)正合（平坦的''忠实性''提供了''当''方向，而平坦性提供了''仅当''方向）。

\textbf{定理 66.4}（忠实平坦的判别准则）：对平坦\(R\)-模\(M\)，下列等价： 1. \(M\)忠实平坦； 2. 对任意极大理想\(\mathfrak{m} \subseteq R\)，\(\mathfrak{m}M \neq M\)（即\(M/\mathfrak{m}M \neq 0\)）； 3. \(M \otimes_R -\)反映零对象； 4. 函子\(M \otimes_R -\)是忠实函子（即\(f \otimes \operatorname{id}_M = 0 \Rightarrow f = 0\)）。

\textbf{证明}：（1）\(\Rightarrow\)（2）：若\(\mathfrak{m}M = M\)，则\(M \otimes_R (R/\mathfrak{m}) = M/\mathfrak{m}M = 0\)。由忠实平坦性，\(R/\mathfrak{m} = 0\)矛盾（\(R/\mathfrak{m}\)为域，非零）。

（2）\(\Rightarrow\)（1）：设\(N \neq 0\)，取非零元\(x \in N\)。令\(\mathfrak{a} = \operatorname{Ann}(x) \subseteq R\)（\(\mathfrak{a} \neq R\)）。取极大理想\(\mathfrak{m} \supseteq \mathfrak{a}\)。由(2)，\(M/\mathfrak{m}M \neq 0\)。由正合列\(0 \to \mathfrak{a} \to R \to R/\mathfrak{a} \to 0\)，作\(M \otimes_R -\)得\(0 \to M/\mathfrak{a}M \to \cdots\)不分裂（若\(\mathfrak{a}M = M\)则得到矛盾）。故\(M \otimes_R Rx \cong M/\mathfrak{a}M \neq 0\)，因此\(M \otimes_R N \neq 0\)（包含非零子模）。

（3）\(\iff\)（1）：由定义直接。

（4）\(\Rightarrow\)（1）：取\(N\)使得\(M \otimes N = 0\)，则\(\operatorname{id}_N \otimes \operatorname{id}_M = 0\)，由忠实性\(\operatorname{id}_N = 0\)，故\(N = 0\)。反之（1）\(\Rightarrow\)（4）需更精细的论证，涉及平坦模上的张量积函子的忠实性与挠对理论。\(\blacksquare\)

\textbf{定理 66.5}（忠实平坦下降------Grothendieck）：设\(R \to A\)为忠实平坦的环同态（即\(A\)作为\(R\)-模忠实平坦）。则以下''下降''性质成立： 1. （模的下降）给定\(A\)-模\(N\)连同下降数据（即\(N \otimes_R A \cong A \otimes_R N\)满足上循环条件），存在\(R\)-模\(M\)使\(M \otimes_R A \cong N\)。 2. （态射的下降）\(R\)-模同态\(f: M \to M'\)为同构当且仅当\(f \otimes \operatorname{id}_A: M \otimes_R A \to M' \otimes_R A\)为同构。 3. （正合列的下降）\(R\)-模序列正合当且仅当经\(-\otimes_R A\)后正合。

\textbf{证明}（框架）：核心工具是Amitsur复形（Cech上同调）。对忠实平坦环同态\(R \to A\)，构造Amitsur复形： \[0 \to M \to M \otimes_R A \to M \otimes_R A^{\otimes 2} \to M \otimes_R A^{\otimes 3} \to \cdots\] 其中微分由交替余乘（插入1）定义。该复形在忠实平坦条件下正合（忠实平坦下降的核心定理）。下降数据恰给出0阶上循环，而正合性保证了该上循环来自某个\(R\)-模。态射的下降直接由平坦忠实性和张量积的性质得出。\(\blacksquare\)

忠实平坦下降在代数几何中的应用极为广泛。例如：构造商概形时需验证下降数据，etale下降（比Zariski下降更精细的拓扑）同样依赖忠实平坦下降的形式，而Grothendieck的平坦下降理论（fpqc拓扑）是导出代数几何和叠（stack）理论的基础。``下降''的基本哲学是：若某个结构在''充分覆盖''（忠实平坦）下存在并相容，则该结构本身存在。这使得我们可以用局部（甚至忠实平坦局部）的构造来定义整体对象。''

\newpage

\chapter{05-代数/02-抽象代数一/04-域扩张与Brauer.md}\label{ux4ee3ux657002-ux62bdux8c61ux4ee3ux6570ux4e0004-ux57dfux6269ux5f20ux4e0ebrauer.md}

\section{第120章：域扩张}\label{ux7b2c120ux7ae0ux57dfux6269ux5f20}

域扩张理论是 Galois 理论的基石，它将域之间的包含关系转化为线性代数和群论的语言。域扩张的研究始于 Gauss 对分圆域的工作和 Abel、Galois 对代数方程的研究，但直到 Dedekind 和 Kronecker 在 19 世纪末系统发展域论之后，域扩张才成为一门独立的学科。Steinitz 于 1910 年发表的经典论文《域的代数理论》首次给出了域扩张的公理化处理，引入了代数扩张、超越次数、可分扩张等至今仍被使用的基本概念。

域扩张 \(K/F\) 的基本观察是 \(K\) 可视为 \(F\) 上的向量空间，其维数 \([K:F]\) 称为扩张次数。代数元与超越元的二分是域扩张理论的第一道分水岭：代数元满足某个 \(F\) 上的非零多项式，其极小多项式唯一决定了单代数扩张 \(F(\alpha)\) 的结构------\(F(\alpha) \cong F[x]/(m_\alpha(x))\)，且 \([F(\alpha):F] = \deg m_\alpha\)。超越元则不满足任何非零多项式，此时 \(F(\alpha) \cong F(x)\)（有理函数域），扩张次数为无穷。代数扩张的有限性具有传递性（次数塔公式 \([L:F] = [L:K][K:F]\)），这是域扩张理论中最基本的计算工具。

分裂域是域扩张理论中的核心构造：给定多项式 \(f \in F[x]\)，\(f\) 的分裂域是包含 \(f\) 的所有根的最小扩域。分裂域的存在性和唯一性（在同构意义下）保证了我们可以''添加多项式的所有根''来构造扩域------这正是 Galois 理论中 Galois 扩张的构造方法。作为分裂域理论的重要应用，有限域的分类完全通过分裂域完成：对任意素数幂 \(p^n\)，存在唯一的 \(p^n\) 元域 \(\mathbb{F}_{p^n}\)，它是 \(x^{p^n} - x\) 在 \(\mathbb{F}_p\) 上的分裂域。本章以尺规作图不可能性的域论证明作为结尾，展示了域扩张理论在古典几何问题中的强大威力------三等分角、倍立方体和化圆为方的不可能性，本质上都是域扩张次数不为 \(2\) 的幂所导致的。

\textbf{定义 44.1.1（域扩张）}：若 \(F \subseteq K\) 均为域，则称 \(K\) 为 \(F\) 的\textbf{域扩张}，记为 \(K/F\)。

\(K\) 可视为 \(F\) 上的向量空间，其维数称为扩张的\textbf{次数}，记为 \([K : F]\)。

\textbf{定义 44.1.2（代数元与超越元）}：设 \(K/F\) 为域扩张，\(\alpha \in K\)。

\begin{itemize}
\tightlist
\item
  若存在非零多项式 \(f \in F[x]\) 使得 \(f(\alpha) = 0\)，则称 \(\alpha\) 为 \(F\) 上的\textbf{代数元}
\item
  否则称 \(\alpha\) 为 \(F\) 上的\textbf{超越元}
\end{itemize}

\textbf{定义 44.1.3（极小多项式）}：设 \(\alpha\) 为 \(F\) 上的代数元。在 \(F[x]\) 中满足 \(f(\alpha) = 0\) 的非零多项式中，次数最低的首一多项式称为 \(\alpha\) 在 \(F\) 上的\textbf{极小多项式}，记为 \(m_{\alpha, F}(x)\)。

\textbf{定理 44.1.1（极小多项式的基本性质）}：

\begin{enumerate}
\def\labelenumi{(\roman{enumi})}
\tightlist
\item
  \(m_{\alpha, F}(x)\) 在 \(F[x]\) 中不可约
\item
  若 \(f \in F[x]\) 满足 \(f(\alpha) = 0\)，则 \(m_{\alpha, F} \mid f\)
\item
  \(F(\alpha) \cong F[x]/(m_{\alpha, F}(x))\)，且 \([F(\alpha) : F] = \deg m_{\alpha, F}\)
\end{enumerate}

\textbf{证明}：

\begin{enumerate}
\def\labelenumi{(\roman{enumi})}
\item
  若 \(m_{\alpha, F} = gh\)（\(g, h\) 次数更低），则 \(g(\alpha)h(\alpha) = 0\)，故 \(g(\alpha) = 0\) 或 \(h(\alpha) = 0\)，与 \(m_{\alpha, F}\) 的次数最低性矛盾。
\item
  多项式带余除法 + 次数最低性（与极小多项式的证明类似）。
\item
  定义同态 \(\varphi : F[x] \to F(\alpha)\) 为 \(\varphi(f) = f(\alpha)\)。\(\ker \varphi = (m_{\alpha, F})\)。由第一同构定理，\(F[x]/(m_{\alpha, F}) \cong \operatorname{im} \varphi\)。由于 \(m_{\alpha, F}\) 不可约，\((m_{\alpha, F})\) 是极大理想，故 \(F[x]/(m_{\alpha, F})\) 是域。\(\operatorname{im} \varphi\) 包含 \(F\) 和 \(\alpha\)，且是最小的这样的域，故 \(\operatorname{im} \varphi = F(\alpha)\)。
\end{enumerate}

\(F(\alpha)\) 作为 \(F\)-向量空间，基为 \(\{1, \alpha, \ldots, \alpha^{n-1}\}\)（\(n = \deg m_{\alpha, F}\)），故 \([F(\alpha) : F] = n\)。\(\blacksquare\)

\textbf{定义 44.1.4（单扩张）}：若 \(K = F(\alpha)\)（\(\alpha \in K\)），则称 \(K/F\) 为\textbf{单扩张}。若 \(\alpha\) 是代数元，则称为\textbf{单代数扩张}。

\hrulefill

\subsection{44.2 有限扩张与代数扩张}\label{ux6709ux9650ux6269ux5f20ux4e0eux4ee3ux6570ux6269ux5f20}

\textbf{定义 44.2.1（有限扩张与代数扩张）}：

\begin{itemize}
\tightlist
\item
  若 \([K : F] < \infty\)，称 \(K/F\) 为\textbf{有限扩张}
\item
  若 \(K\) 中每个元素都是 \(F\) 上的代数元，称 \(K/F\) 为\textbf{代数扩张}
\end{itemize}

\textbf{定理 44.2.1（次数塔公式）}：设 \(F \subseteq K \subseteq L\) 为域扩张。则

\[
[L : F] = [L : K] \cdot [K : F]
\]

\textbf{证明}：设 \(\{x_1, \ldots, x_m\}\) 为 \(K\) 在 \(F\) 上的基，\(\{y_1, \ldots, y_n\}\) 为 \(L\) 在 \(K\) 上的基。可以证明 \(\{x_i y_j \mid 1 \leq i \leq m, 1 \leq j \leq n\}\) 是 \(L\) 在 \(F\) 上的基。\(\blacksquare\)

\textbf{定理 44.2.2}：有限扩张 \(\implies\) 代数扩张。

\textbf{证明}：设 \([K : F] = n\)。对任意 \(\alpha \in K\)，\(1, \alpha, \ldots, \alpha^n\) 这 \(n+1\) 个元素在 \(n\) 维 \(F\)-向量空间 \(K\) 中必线性相关。故存在不全为零的 \(a_i \in F\) 使得 \(\sum a_i \alpha^i = 0\)，即 \(\alpha\) 是代数元。\(\blacksquare\)

\textbf{定理 44.2.3}：\(K/F\) 是有限扩张当且仅当 \(K = F(\alpha_1, \ldots, \alpha_n)\)，其中每个 \(\alpha_i\) 是 \(F\) 上的代数元。

\textbf{证明}：若 \(K/F\) 有限，取 \(K\) 在 \(F\) 上的基 \(\{\alpha_1, \ldots, \alpha_n\}\)，则 \(K = F(\alpha_1, \ldots, \alpha_n)\)，且每个 \(\alpha_i\) 是代数元。

反之，若 \(K = F(\alpha_1, \ldots, \alpha_n)\)（\(\alpha_i\) 代数），则 \([F(\alpha_1) : F] < \infty\)，\([F(\alpha_1, \alpha_2) : F(\alpha_1)] < \infty\)，\ldots，由塔公式乘积有限。\(\blacksquare\)

\hrulefill

\subsection{44.3 分裂域}\label{ux5206ux88c2ux57df}

\textbf{定义 44.3.1（分裂域）}：设 \(f \in F[x]\) 为非常数多项式。\(F\) 的扩域 \(K\) 称为 \(f\) 在 \(F\) 上的\textbf{分裂域}，若

\begin{enumerate}
\def\labelenumi{\arabic{enumi}.}
\tightlist
\item
  \(f\) 在 \(K[x]\) 中可分解为一次因式的乘积
\item
  \(K = F(\alpha_1, \ldots, \alpha_n)\)，其中 \(\alpha_i\) 为 \(f\) 在 \(K\) 中的根
\end{enumerate}

\textbf{定理 44.3.1（分裂域的存在性）}：对任何 \(f \in F[x]\)，\(f\) 在 \(F\) 上的分裂域存在。

\textbf{证明}：对 \(\deg f\) 归纳。若 \(\deg f = 1\)，\(F\) 自身就是分裂域。

若 \(\deg f > 1\)，取 \(f\) 的一个不可约因子 \(p(x)\)。\(K_1 = F[x]/(p(x))\) 是 \(F\) 的扩域，且 \(p(x)\) 在 \(K_1\) 中有根 \(\alpha_1 = x + (p(x))\)。

在 \(K_1[x]\) 中，\(f(x) = (x - \alpha_1)g(x)\)（\(\deg g = \deg f - 1\)）。由归纳假设，\(g\) 在 \(K_1\) 上有分裂域 \(K\)。则 \(K\) 是 \(f\) 在 \(F\) 上的分裂域。\(\blacksquare\)

\textbf{定理 44.3.2（分裂域的唯一性）}：\(f \in F[x]\) 在 \(F\) 上的任何两个分裂域是同构的（\(F\)-同构）。

\textbf{证明}：对 \([K_1 : F] = [K_2 : F]\) 归纳。设 \(K_1, K_2\) 为 \(f \in F[x]\) 的两个分裂域。取 \(f\) 的不可约因子 \(p(x)\)，设 \(\alpha_1 \in K_1\)，\(\alpha_2 \in K_2\) 为 \(p(x)\) 的根。则存在 \(F\)-同构 \(\sigma: F(\alpha_1) \xrightarrow{\cong} F[x]/(p(x)) \xrightarrow{\cong} F(\alpha_2)\)。在 \(K_1\) 中，\(f(x) = (x - \alpha_1)g_1(x)\)，在 \(K_2\) 中，\(f(x) = (x - \alpha_2)g_2(x)\)。由 \(\sigma\) 诱导系数同构，\(g_1\) 和 \(g_2\) 在对应下可视为同一多项式。\(K_1\) 是 \(g_1\) 在 \(F(\alpha_1)\) 上的分裂域，\(K_2\) 是 \(g_2\) 在 \(F(\alpha_2)\) 上的分裂域。由归纳假设，\(\sigma\) 可延拓为 \(K_1 \cong K_2\)。\(\blacksquare\)

\hrulefill

\subsection{44.4 有限域}\label{ux6709ux9650ux57df}

\textbf{定义 44.4.1（有限域）}：元素个数有限的域称为\textbf{有限域}（或 Galois 域）。

\textbf{定理 44.4.1（有限域的结构）}：

\begin{enumerate}
\def\labelenumi{(\roman{enumi})}
\tightlist
\item
  有限域 \(F\) 的元素个数必为 \(p^n\)（\(p\) 为素数，\(n \geq 1\)），其中 \(p\) 为 \(F\) 的特征（\(\operatorname{char} F\)）
\item
  对任意素数 \(p\) 和正整数 \(n\)，存在唯一的（同构意义下）\(p^n\) 元有限域，记为 \(\mathbb{F}_{p^n}\)
\item
  \(\mathbb{F}_{p^n}\) 是 \(x^{p^n} - x\) 在 \(\mathbb{F}_p\) 上的分裂域
\end{enumerate}

\textbf{证明}：

\begin{enumerate}
\def\labelenumi{(\roman{enumi})}
\item
  有限域 \(F\) 必包含其素子域 \(\mathbb{F}_p\)（\(p = \operatorname{char} F\)）。\(F\) 是 \(\mathbb{F}_p\) 上的有限维向量空间，设维数为 \(n\)，则 \(|F| = p^n\)。
\item
  考虑多项式 \(x^{p^n} - x \in \mathbb{F}_p[x]\)。其形式导数 \(p^n x^{p^n-1} - 1 = -1 \neq 0\)，故无重根。设 \(K\) 为该多项式的分裂域。\(K\) 中满足 \(x^{p^n} = x\) 的元素恰为 \(p^n\) 个（因为无重根），且它们构成 \(K\) 的子域（由特征 \(p\) 的性质，\((a+b)^p = a^p + b^p\)）。因此该子域本身即为 \(p^n\) 个元素的域。
\end{enumerate}

唯一性：任何 \(p^n\) 元域 \(F\) 中，非零元素乘法群 \(F^\times\) 的阶为 \(p^n-1\)，故所有元素满足 \(x^{p^n} - x = 0\)。因此 \(F\) 是 \(x^{p^n} - x\) 的分裂域，由分裂域的唯一性得证。\(\blacksquare\)

\textbf{定理 44.4.2（有限域的乘法群）}：\(\mathbb{F}_{p^n}^{\times}\) 是 \(p^n-1\) 阶循环群。

\textbf{证明}：\(\mathbb{F}_{p^n}^{\times}\) 是有限阿贝尔群。由有限阿贝尔群的结构定理，\(\mathbb{F}_{p^n}^{\times} \cong \mathbb{Z}_{d_1} \oplus \cdots \oplus \mathbb{Z}_{d_k}\)（\(d_1 \mid \cdots \mid d_k\)）。所有元素满足 \(x^{d_k} = 1\)。但 \(x^{d_k} - 1\) 在域中最多有 \(d_k\) 个根，而 \(|\mathbb{F}_{p^n}^{\times}| = p^n - 1\)，故 \(d_k = p^n - 1\)，\(k = 1\)。因此 \(\mathbb{F}_{p^n}^{\times}\) 是循环群。\(\blacksquare\)

\textbf{推论 44.4.3（有限域中的本原元）}：\(\mathbb{F}_{p^n}\) 中存在元素 \(\alpha\) 使得 \(\mathbb{F}_{p^n} = \mathbb{F}_p(\alpha)\)。事实上，\(\mathbb{F}_{p^n}^{\times}\) 的任何生成元都具有此性质。

\hrulefill

\subsection{44.5 尺规作图不可能性（概述）}\label{ux5c3aux89c4ux4f5cux56feux4e0dux53efux80fdux6027ux6982ux8ff0}

\textbf{定理 44.5.1（可构造数）}：实数 \(\alpha\) 称为\textbf{可构造数}（即可用尺规作图从单位线段构造），当且仅当存在域扩张链

\[
\mathbb{Q} = F_0 \subseteq F_1 \subseteq \cdots \subseteq F_k \subseteq \mathbb{R}
\]

使得 \([F_i : F_{i-1}] \leq 2\) 且 \(\alpha \in F_k\)。

\textbf{推论 44.5.2}：可构造数 \(\alpha\) 满足 \([\mathbb{Q}(\alpha) : \mathbb{Q}] = 2^m\)（\(m \geq 0\)）。

\textbf{定理 44.5.3（三大尺规作图不可能问题）}：

\begin{enumerate}
\def\labelenumi{(\roman{enumi})}
\tightlist
\item
  \textbf{倍立方体}：\(\sqrt[3]{2}\) 不可构造（\([\mathbb{Q}(\sqrt[3]{2}) : \mathbb{Q}] = 3\)，不是 \(2\) 的幂）
\item
  \textbf{三等分任意角}：\(\cos 20^\circ\) 不可构造（极小多项式为 \(8x^3 - 6x + 1\)，次数为 \(3\)）
\item
  \textbf{化圆为方}：\(\sqrt{\pi}\) 不可构造（\(\pi\) 是超越数，\([\mathbb{Q}(\pi) : \mathbb{Q}] = \infty\)）
\end{enumerate}

\textbf{证明}：

\begin{enumerate}
\def\labelenumi{(\roman{enumi})}
\item
  \(x^3 - 2\) 在 \(\mathbb{Q}[x]\) 中不可约（Eisenstein 判别法，\(p = 2\)）。故 \([\mathbb{Q}(\sqrt[3]{2}) : \mathbb{Q}] = 3 \neq 2^m\)。
\item
  \(\cos 3\theta = 4\cos^3 \theta - 3\cos \theta\)。令 \(\theta = 20^\circ\)，则 \(\cos 60^\circ = 1/2\)，故 \(4\cos^3 20^\circ - 3\cos 20^\circ = 1/2\)，即 \(8\cos^3 20^\circ - 6\cos 20^\circ - 1 = 0\)。可证 \(8x^3 - 6x - 1\) 在 \(\mathbb{Q}[x]\) 中不可约，故 \([\mathbb{Q}(\cos 20^\circ) : \mathbb{Q}] = 3 \neq 2^m\)。
\item
  Lindemann 定理（1882 年）证明 \(\pi\) 是超越数，故 \([\mathbb{Q}(\pi) : \mathbb{Q}] = \infty\)。\(\blacksquare\)
\end{enumerate}

\hrulefill

\textbf{注记 44.5.1}：Galois 理论的完整展开（包括方程可解性、\(n\) 次一般方程不可根式解等）将在卷十三（抽象代数 II）中详细讨论。本章仅建立域扩张的基本理论，为 Galois 理论做好准备。\(\blacksquare\)

\subsection{67.A 域扩张补充}\label{a-ux57dfux6269ux5f20ux8865ux5145}

域扩张理论在 Galois 理论之外还有多个重要的发展方向，它们在现代代数几何和数论中扮演着核心角色。\textbf{可分扩张与纯不可分扩张}的二分是域扩张理论的基本结构定理：设 \(K/F\) 为代数扩张，则存在唯一的中间域 \(F \subseteq F_s \subseteq K\)，使得 \(F_s/F\) 是 Galois 扩张（极大可分扩张）且 \(K/F_s\) 是纯不可分扩张。这一分解将任意代数扩张拆分为两个具有截然不同性质的组分------可分部分由 Galois 群控制，纯不可分部分由 Frobenius 态射（正特征情形）或平凡（特征零）控制。\textbf{超越次数}（transcendence degree）是有理函数域 \(F(t_1, \ldots, t_n)\) 中代数无关的变元的最大个数，它是代数几何中代数簇的维数（Krull 维数）的域论对应：对仿射整环 \(R\) 的分式域 \(K = \operatorname{Frac}(R)\)，\(\dim R = \operatorname{tr.deg}_F K\)。\textbf{正则扩张}（regular extension）是域扩张中一种特殊类型：\(K/F\) 是正则扩张若 \(F\) 在 \(K\) 中代数闭且 \(K\) 与 \(F\) 的代数闭包线性不交，这在代数群的 Weil 限制和代数几何的基变换理论中至关重要。\textbf{完备域}（perfect field）是每个代数扩张都可分的域------特征零域和有限域是完备的，而 \(\mathbb{F}_p(t)\) 则不是（因为 \(t^{1/p}\) 是纯不可分元）。完备域上的代数簇避免了正特征下的许多病理现象，如 Frobenius 态射的非光滑性。在算术代数几何中，\textbf{p-进域}（\(\mathbb{Q}_p\) 的有限扩张）的扩张理论导出局部类域论------通过 Lubin-Tate 形式群将 Abel 扩张与乘法群联系起来，这是高维局部类域论的起点。域扩张理论还与\textbf{模型论}（model theory）交汇：代数闭域理论具有量词消去性质，差分域（装备自同构的域）的模型论性质催生了 Hrushovski 的几何 Mordell-Lang 猜想的证明。

\hrulefill

\hrulefill

\section{第121章：除环与 Brauer 群}\label{ux7b2c121ux7ae0ux9664ux73afux4e0e-brauer-ux7fa4}

除环（即体，skew field）是非交换代数结构的核心研究对象之一。Brauer 群则将域上中心单代数的分类组织为阿贝尔群。

\subsection{48.1 Frobenius 定理与有限除环}\label{frobenius-ux5b9aux7406ux4e0eux6709ux9650ux9664ux73af}

\textbf{定义 48.1}（中心单代数）：域 \(\mathbb{F}\) 上的有限维代数 \(A\) 称为 \textbf{中心单} 若 \(Z(A) = \mathbb{F}\)（中心恰为 \(\mathbb{F}\)）且 \(A\) 是单代数（无真双边理想）。

\textbf{定理 48.1}（Frobenius）：\(\mathbb{R}\) 上的有限维除环仅为 \(\mathbb{R}\)、\(\mathbb{C}\) 和 \(\mathbb{H}\)（四元数，Hamilton, 1843）。

\textbf{证明}：设 \(D\) 为 \(\mathbb{R}\) 上有限维除环。取 \(d \in D \setminus \mathbb{R}\)，则 \(\mathbb{R}(d)\) 是 \(\mathbb{R}\) 的真代数扩张，故 \(\mathbb{R}(d) \cong \mathbb{C}\)，且 \(d\) 满足 \(d^2 + a d + b = 0\)（\(a^2 < 4b\)）。经平移，可设 \(d^2 = -1\)（即 \(d = i\)）。若 \(\dim_{\mathbb{R}} D = 2\)，则 \(D = \mathbb{C}\)。若 \(\dim_{\mathbb{R}} D > 2\)，取 \(d' \notin \mathbb{R}(d)\)。由同论证，经平移可设 \(d'^2 = -1\)。令 \(j = d'\)，\(i = d\)。考虑 \(i + j\) 和 \(i - j\) 的平方：\((i+j)^2 = i^2 + ij + ji + j^2 = -2 + (ij + ji)\)，\((i-j)^2 = -2 - (ij + ji)\)。由除环性质，\((i+j)^2\) 和 \((i-j)^2\) 均为实数，故 \(ij + ji \in \mathbb{R}\)。设 \(ij + ji = 2c\)（\(c \in \mathbb{R}\)）。令 \(k = ij + c\)，则 \(k^2 = -1 + c^2\)。通过调整 \(j\) 的选择，可设 \(c = 0\)，即 \(ij = -ji\)。定义 \(k = ij\)，验证 \(i^2 = j^2 = k^2 = ijk = -1\)。故 \(\{1, i, j, k\}\) 张成四元数代数 \(\mathbb{H} \subseteq D\)。若 \(D\) 严格大于 \(\mathbb{H}\)，取 \(d'' \notin \mathbb{H}\)，同理导出矛盾（五维以上除环不存在）。故 \(D = \mathbb{H}\)。\(\blacksquare\)

\textbf{定理 48.2}（Wedderburn 小定理）：有限除环必为域（交换）。

\textbf{证明}：设 \(D\) 为有限除环，\(Z = Z(D)\) 为其中心（为有限域，\(|Z| = q\)），\(|D| = q^n\)。对 \(D^\times\) 写类方程：\(|D^\times| = |Z^\times| + \sum_{a \notin Z} [D^\times : C_{D^\times}(a)]\)，即 \(q^n - 1 = (q - 1) + \sum_{a \notin Z} \frac{q^n - 1}{q^{d_a} - 1}\)（\(d_a < n\)）。对每个非中心共轭类，\(\frac{q^n - 1}{q^{d_a} - 1}\) 是整数。考虑分圆多项式 \(\Phi_n(q)\)：它整除 \(q^n - 1\) 且不整除任何 \(q^d - 1\)（\(d < n\)）。故 \(\Phi_n(q)\) 整除每个加项，且整除 \(q^n - 1\)，从而整除 \(q - 1\)。但 \(q \geq 2\) 时 \(\Phi_n(q) > q - 1\) 对 \(n > 1\)，矛盾。故 \(n = 1\)，\(D = Z\) 为域。\(\blacksquare\)

\subsection{48.2 Brauer 群}\label{brauer-ux7fa4}

\textbf{定义 48.2}（Brauer 群）：在中心单代数上定义等价关系 \(A \sim B \iff A \otimes M_n(\mathbb{F}) \cong B \otimes M_m(\mathbb{F})\)（矩阵代数稳定化）。等价类的集合在张量积下构成阿贝尔群 \(\operatorname{Br}(\mathbb{F})\)，其逆元由反代数 \(A^{op}\) 给出。

\textbf{定理 48.3}（Skolem-Noether）：设 \(A\) 为 \(\mathbb{F}\) 上中心单代数，则 \(A\) 的任意自同构均为内自同构：对任意 \(\varphi \in \operatorname{Aut}_{\mathbb{F}}(A)\) 存在 \(a \in A^{\times}\) 使 \(\varphi(x) = a x a^{-1}\)。

\textbf{证明}：由 Wedderburn 结构定理，\(A \cong M_n(D)\)（\(D\) 为 \(\mathbb{F}\) 上中心除环）。设 \(A^{\mathrm{op}}\) 为 \(A\) 的反代数，考虑 \(A \otimes_{\mathbb{F}} A^{\mathrm{op}}\) 在 \(A\) 上的两种作用：\((x \otimes y) \cdot a = x a y\)（左正作用）和 \((x \otimes y) \cdot_{\varphi} a = \varphi(x) a y\)（\(\varphi\) 扭曲作用）。两种作用均使 \(A\) 成为 \(A \otimes A^{\mathrm{op}} \cong M_{n^2}(\mathbb{F})\) 上的单模。由单模的唯一性，两作用等价，故存在 \(A\)-线性同构 \(\psi: A \to A\) 使得 \(\psi(x a y) = \varphi(x) \psi(a) y\)。取 \(a = \psi(1)\)，则 \(a \in A^{\times}\) 且 \(\varphi(x) = a x a^{-1}\)。\(\blacksquare\)

\textbf{例}：\(\operatorname{Br}(\mathbb{R}) \cong \mathbb{Z}/2\mathbb{Z}\)（仅含 \([\mathbb{R}]\) 和 \([\mathbb{H}]\)）；\(\operatorname{Br}(\mathbb{F}_q) = 0\)（有限域除环必为域）；\(\operatorname{Br}(\mathbb{Q}_p)\) 在类域论中等价于 \(\mathbb{Q}/\mathbb{Z}\)。

\subsection{Brauer 群的深层理论与应用}\label{brauer-ux7fa4ux7684ux6df1ux5c42ux7406ux8bbaux4e0eux5e94ux7528}

Brauer 群 \(\operatorname{Br}(F)\) 通过 Hasse 原理与局部-整体原则紧密关联。\textbf{Albert-Brauer-Hasse-Noether 定理}（1932）断言，数域 \(F\) 的中心单代数 \(A\) 是分裂的（即 \(A \cong M_n(F)\)）当且仅当 \(A\) 在 \(F\) 的每个完备化 \(F_v\) 处分裂（\(v\) 遍历所有位）。这给出了 Brauer 群的 Hasse 原理：\(\operatorname{Br}(F) \hookrightarrow \bigoplus_v \operatorname{Br}(F_v)\) 是单射，且余核由局部不变量之和为零的条件刻画，即短正合列 \(0 \to \operatorname{Br}(F) \to \bigoplus_v \operatorname{Br}(F_v) \xrightarrow{\sum \operatorname{inv}_v} \mathbb{Q}/\mathbb{Z} \to 0\)。\textbf{Merkurjev-Suslin 定理}（1982）是 K-理论中的里程碑：若 \(F\) 包含 \(n\) 次本原单位根，则 \(K_2(F)/nK_2(F) \cong {}_n\operatorname{Br}(F)\)（\(n\)-挠 Brauer 群），且每个 \(n\)-挠类由循环代数 \((a, b)_{\zeta}\) 表示。这解决了关于 Brauer 群生成元的长期猜想。\textbf{Brauer 群与上同调}：\(\operatorname{Br}(F) \cong H^2(\operatorname{Gal}(F^s/F), (F^s)^\times)\)（Galois 上同调的第二个群），这为 Brauer 群提供了上同调解释，并将其嵌入 Galois 上同调的长期框架。在\textbf{算术代数几何}中，Brauer 群是 Brauer-Manin 阻碍的核心------若簇 \(X\) 的 Brauer-Manin 集 \(X(\mathbb{A}_F)^{\operatorname{Br}} = \varnothing\) 但 \(X(\mathbb{A}_F) \neq \varnothing\)，则 \(X\) 违反 Hasse 原理（无有理点但有局部点），这是解释有理点失败现象的主要工具。在\textbf{非交换代数几何}中，Brauer 群分类了域上的非交换射影空间（Severi-Brauer 簇），而 Amitsur 定理将 \(\operatorname{Br}(F(x_1, \ldots, x_n))\) 与 \(\operatorname{Br}(F)\) 和除环的泛有理性联系起来。Brauer 群还与\textbf{拓扑学}中的扭 K-理论有关：\(X\) 上的 Azumaya 代数等价于 \(X\) 的扭 K-理论中的扭类，这是 Donovan-Karoubi 定理（1970）的核心内容。

\hrulefill

\hrulefill

\subsection{67.A 域扩张补充}\label{a-ux57dfux6269ux5f20ux8865ux5145-1}

格论（Lattice Theory）是序结构与代数结构交汇的数学分支，由 Birkhoff（1930s）系统化。它在泛代数、域论、逻辑和计算机科学中有广泛应用。格（lattice）既是偏序集（每对元素有上确界和下确界），也是具有两个二元运算（\(\vee\) 和 \(\wedge\)）的代数结构。

\subsection{格的基础理论}\label{ux683cux7684ux57faux7840ux7406ux8bba}

\textbf{定义}（格 / lattice）：偏序集 \((L, \leq)\) 是\textbf{格}，如果对任意 \(a, b \in L\)，存在上确界 \(a \vee b\)（join）和下确界 \(a \wedge b\)（meet）。等价地，格是装备两个满足结合律、交换律和吸收律（absorption）的二元运算的代数结构： \[a \vee (a \wedge b) = a, \quad a \wedge (a \vee b) = a\]

\textbf{例}： - 集合的幂集 \((\mathcal{P}(X), \subseteq)\)（\(\vee = \cup\)，\(\wedge = \cap\)） - 群的子群族（子群格，\(\vee\) 为生成的子群，\(\wedge\) 为交） - 环的理想族（理想格） - 拓扑空间的闭集族 - \(\mathbb{N}\) 上的整除偏序：\(a \vee b = \operatorname{lcm}(a, b)\)，\(a \wedge b = \gcd(a, b)\)

\textbf{定义}（完备格 / complete lattice）：格 \(L\) 是\textbf{完备格}，如果任意子集（不限限于二元）有上确界和下确界。每个完备格有最大元 \(\top = \bigvee L\) 和最小元 \(\bot = \bigwedge L\)。幂集格是完备的。

\textbf{定理}（Knaster-Tarski 不动点定理）：完备格 \(L\) 上的保序映射 \(f: L \to L\)（即 \(x \leq y \implies f(x) \leq f(y)\)）必有不动点。事实上，不动点集构成完备格，且最小不动点为 \(\bigwedge \{x : f(x) \leq x\}\)。该定理是计算机科学中指称语义（denotational semantics）和程序验证的数学基础。

\subsection{分配格、模格与Boolean代数}\label{ux5206ux914dux683cux6a21ux683cux4e0ebooleanux4ee3ux6570}

\textbf{定义}（分配格）：格 \(L\) 是\textbf{分配格}，若满足分配律：\(a \vee (b \wedge c) = (a \vee b) \wedge (a \vee c)\)（等价于其对偶形式）。幂集格是分配的；子群格一般不是。

\textbf{定理}（Birkhoff-Stone 表示定理，1930s）：每个分配格同构于某个集合的集合代数的子格（均以 \(\cup, \cap\) 为运算）。换言之，分配格的可表示性由其序结构完全确定。

\textbf{定义}（模格 / modular lattice）：格 \(L\) 是\textbf{模格}（Dedekind），若 \(a \leq c \implies a \vee (b \wedge c) = (a \vee b) \wedge c\)。所有分配格是模格，所有正规子群组成的子群格是模格（Dedekind定理）。模格是代数几何中维数理论的代数骨架（Krull维数和高度函数）。

\textbf{定义}（Boolean代数）：\textbf{Boolean代数}是具有补运算（\(a \vee \neg a = \top\)，\(a \wedge \neg a = \bot\)）的分配格。每个Boolean代数同构于某个集合的幂集Boolean代数的子代数（Stone表示定理，1936）。

\textbf{定理}（Stone对偶 / Stone duality）：Boolean代数的范畴反等价于完全不连通的紧 Hausdorff 空间（\textbf{Stone空间}）的范畴。这推广了有限Boolean代数到全体集合的有限幂集的对应。

\subsection{Galois联络与闭包算子}\label{galoisux8054ux7edcux4e0eux95edux5305ux7b97ux5b50}

\textbf{定义}（Galois联络）：偏序集 \((P, \leq)\) 和 \((Q, \leq)\) 之间的\textbf{Galois联络}是一对映射 \(f: P \to Q\) 和 \(g: Q \to P\)，满足 \[f(p) \leq q \iff p \leq g(q) \quad (\forall p \in P, q \in Q)\] 这等价于 \(f\) 和 \(g\) 均保序（\(p_1 \leq p_2 \implies f(p_1) \leq f(p_2)\)）且 \(f \circ g\) 和 \(g \circ f\) 是扩张（closure）或核算（kernel）算子。

\textbf{例}： - Galois理论中的Galois联络：固定 Galois 扩张 \(E/F\)。\(f(H) = E^H\)（中间域 \(E\) 的子域→子群），\(g(K) = \operatorname{Gal}(E/K)\)（子群→中间域）。这是反序的Galois联络：\(H \leq K \iff E^K \subseteq E^H\) - 代数几何中：理想 \(I \subset k[x_1, \ldots, x_n]\) ↔ 代数集 \(V(I) \subset \mathbb{A}^n\) 是 Galois 联络 - Dedekind-MacNeille完备化：通过 \(f(A) = \{u : \forall a \in A, a \leq u\}\)（上界算子）和 \(g(B) = \{l : \forall b \in B, l \leq b\}\)（下界算子），任何偏序集可嵌入完备格

\textbf{定义}（闭包算子 / closure operator）：\(P\) 上的\textbf{闭包算子}是映射 \(c: P \to P\) 满足：(1) \(p \leq c(p)\)（扩张）；(2) \(p \leq q \implies c(p) \leq c(q)\)（单调）；(3) \(c(c(p)) = c(p)\)（幂等）。闭包算子的不动点构成 \(P\) 的完备子半格。例子：拓扑空间的 Kuratowski 闭包、凸包、线性生成、代数闭包、逻辑中的演绎闭包。

\emph{卷七：抽象代数 I终。}

\emph{卷七：抽象代数 I终。}

\emph{卷七：抽象代数 I终。}

\emph{卷八：抽象代数（一）终。} \hrulefill

\subsection{泛代数基础}\label{ux6cdbux4ee3ux6570ux57faux7840}

泛代数（Universal Algebra）是代数结构的元理论------它不研究某一特定的代数结构（如群、环、模），而是研究所有代数结构的共同性质。泛代数由 Garrett Birkhoff 在 1930 年代创立，其基本思想是将群、环、格、布尔代数等看似不同的代数结构统一在''带运算的集合''这一框架下，通过等式逻辑来刻画和分类。泛代数的语言为数学提供了最一般的''代数结构''概念，其影响遍及模型论、范畴论和理论计算机科学。

泛代数的核心概念是代数簇（variety）------由一组等式定义的代数结构的类。例如，群是满足结合律、幺元律和逆元律三类等式的代数簇，环是满足阿贝尔群公理和乘法结合律、分配律的代数簇，而格是满足交换律、结合律和吸收律的代数簇。Birkhoff 的 HSP 定理（1935）是泛代数的基石：一个代数结构的类 \(\mathcal{V}\) 是代数簇当且仅当 \(\mathcal{V}\) 在同态像（H）、子代数（S）和直积（P）下封闭。这一定理的美妙之处在于，它将语义条件（由等式定义）等价于语法条件（在 HSP 下封闭），类似于模型论中 Birkhoff 定理与 Los 定理之间的对应。

泛代数还提供了研究代数结构之间关系的系统工具。自由代数、余代数、项代数、同余关系等概念使得我们可以脱离具体的代数结构讨论''结构''本身。在理论计算机科学中，泛代数被用于研究抽象数据类型和代数规约------数据类型由符号集和等式定义，其初始代数模型给出了程序语义的指称解释。在范畴论中，泛代数与 Lawvere 理论的等价性揭示了代数结构的范畴本质：每个代数簇对应一个有限积完备的小范畴，其模型是该范畴到集合范畴的保积函子。本章将系统介绍泛代数的基本概念------符号集、项代数、等式、代数簇和 Birkhoff 定理------为后续的范畴论和模型论提供统一的代数语言。

\textbf{定义 1}（符号集）：一个\textbf{符号集}（signature）\(\Sigma\) 是一族函数符号，每个符号 \(f \in \Sigma\) 附有元数（arity）\(\mathrm{ar}(f) \in \mathbb{N}\)。零元符号（\(\mathrm{ar}(f) = 0\)）称为\textbf{常量符号}。

\textbf{定义 2}（\(\Sigma\)-代数）：一个 \(\Sigma\)-\textbf{代数}是一个偶对 \(\mathbf{A} = (A, (f^{\mathbf{A}})_{f \in \Sigma})\)，其中 \(A\) 是非空集合（承载集），每个 \(f^{\mathbf{A}}: A^{\mathrm{ar}(f)} \to A\) 是 \(A\) 上的 \(\mathrm{ar}(f)\) 元运算。

\textbf{定义 3}（项代数）：设 \(X\) 是变元集合。\(\Sigma\)-\textbf{项}的集合 \(T_\Sigma(X)\) 归纳定义为： 1. 每个 \(x \in X\) 是项。 2. 每个常量符号 \(c \in \Sigma\) 是项。 3. 若 \(f \in \Sigma\)，\(\mathrm{ar}(f) = n\)，且 \(t_1, \ldots, t_n\) 是项，则 \(f(t_1, \ldots, t_n)\) 是项。

\(T_\Sigma(X)\) 本身以自然方式成为 \(\Sigma\)-代数：对 \(f \in \Sigma\)，定义 \(f^{T_\Sigma(X)}(t_1, \ldots, t_n) = f(t_1, \ldots, t_n)\)。\(T_\Sigma(X)\) 是 \(\Sigma\)-代数范畴中的自由代数，满足以下普遍性质：对任意 \(\Sigma\)-代数 \(\mathbf{A}\) 和任意映射 \(\varphi: X \to A\)，存在唯一的 \(\Sigma\)-同态 \(\bar{\varphi}: T_\Sigma(X) \to \mathbf{A}\) 延拓 \(\varphi\)。

\textbf{定义 4}（等式）：一个 \(\Sigma\)-\textbf{等式}是形如 \(s \approx t\) 的公式，其中 \(s, t \in T_\Sigma(X)\)。\(\Sigma\)-代数 \(\mathbf{A}\) \textbf{满足} \(s \approx t\)，记作 \(\mathbf{A} \models s \approx t\)，若对每个赋值 \(h: X \to A\)，有 \(\bar{h}(s) = \bar{h}(t)\)。

\subsection{代数簇与Birkhoff定理}\label{ux4ee3ux6570ux7c07ux4e0ebirkhoffux5b9aux7406}

\textbf{定义 5}（代数簇）：\(\Sigma\)-代数的类 \(\mathcal{V}\) 称为一个\textbf{代数簇}（variety）或\textbf{等式类}，若存在等式集合 \(E \subseteq T_\Sigma(X) \times T_\Sigma(X)\)，使得 \(\mathcal{V} = \{\mathbf{A} \mid \mathbf{A} \models e \text{ 对所有 } e \in E\}\)。

\textbf{定理 1}（Birkhoff HSP定理，1935）：\(\Sigma\)-代数的一个非空类 \(\mathcal{V}\) 是代数簇当且仅当 \(\mathcal{V}\) 在以下三个运算下封闭： - \(\mathbf{H}\)：同态像（homomorphic images） - \(\mathbf{S}\)：子代数（subalgebras） - \(\mathbf{P}\)：直积（direct products）

\textbf{证明}：（\(\Rightarrow\)）设 \(\mathcal{V} = \mathrm{Mod}(E)\) 是等式定义的类。需证 \(\mathcal{V}\) 在 \(\mathbf{H}, \mathbf{S}, \mathbf{P}\) 下封闭。

\(\mathbf{S}\)-封闭：设 \(\mathbf{B} \leq \mathbf{A} \in \mathcal{V}\)。对任意等式 \(s \approx t \in E\) 和赋值 \(h: X \to B\)，视为 \(h: X \to A\)，由 \(\mathbf{A} \models s \approx t\) 知 \(\bar{h}(s) = \bar{h}(t)\)，而 \(\bar{h}\) 的像在 \(\mathbf{B}\) 的承载集中，故 \(\mathbf{B} \models s \approx t\)。因此 \(\mathbf{B} \in \mathcal{V}\)。

\(\mathbf{H}\)-封闭：设 \(\varphi: \mathbf{A} \to \mathbf{B}\) 是满 \(\Sigma\)-同态，\(\mathbf{A} \in \mathcal{V}\)。对任意赋值 \(h: X \to B\)，因 \(\varphi\) 是满射，可选取 \(\tilde{h}: X \to A\) 使得 \(\varphi \circ \tilde{h} = h\)。则 \(\bar{h}(s) = \varphi(\bar{\tilde{h}}(s)) = \varphi(\bar{\tilde{h}}(t)) = \bar{h}(t)\)，故 \(\mathbf{B} \models s \approx t\)，即 \(\mathbf{B} \in \mathcal{V}\)。

\(\mathbf{P}\)-封闭：设 \(\mathbf{A}_i \in \mathcal{V}\)（\(i \in I\)），\(\mathbf{A} = \prod_{i \in I} \mathbf{A}_i\)。对等式 \(s \approx t \in E\) 和赋值 \(h: X \to A\)，由直积定义及每个 \(\mathbf{A}_i \models s \approx t\)，在每个分量上 \(\pi_i(\bar{h}(s)) = \bar{\pi_i \circ h}(s) = \bar{\pi_i \circ h}(t) = \pi_i(\bar{h}(t))\)，故 \(\bar{h}(s) = \bar{h}(t)\)。得 \(\mathbf{A} \in \mathcal{V}\)。

（\(\Leftarrow\)）设 \(\mathcal{V}\) 在 \(\mathbf{H}, \mathbf{S}, \mathbf{P}\) 下封闭。令 \(E = \{(s, t) \mid \mathcal{V} \models s \approx t\}\) 是 \(\mathcal{V}\) 满足的所有等式。需证 \(\mathcal{V} = \mathrm{Mod}(E)\)。显然 \(\mathcal{V} \subseteq \mathrm{Mod}(E)\)。关键在反向包含。

对任意 \(\mathbf{A} \in \mathrm{Mod}(E)\)，需证 \(\mathbf{A} \in \mathcal{V}\)。构造自由代数：对足够大的集合 \(Y\)（可取 \(Y = A\)），令 \(\mathbf{F} = \mathbf{F}_{\mathcal{V}}(Y)\) 是 \(\mathcal{V}\) 中的自由代数。由自由性，存在满同态 \(\varphi: \mathbf{F} \to \mathbf{A}\)。还需验证 \(\mathbf{F} \in \mathcal{V}\)：自由代数 \(\mathbf{F}_{\mathcal{V}}(Y)\) 构造为 \(T_\Sigma(Y)\) 商去由 \(E\) 生成的同余 \(\theta_{\mathcal{V}}(Y)\)。由于 \(\mathcal{V}\) 在 \(\mathbf{P}\) 下封闭，\(\mathbf{F}\) 可嵌入 \(\prod_{\mathbf{B} \in \mathcal{V}, |B| \le |Y|} \mathbf{B}\) 的适当幂（因为 \(\mathcal{V}\) 局部小），且因 \(\mathcal{V}\) 在 \(\mathbf{S}\) 和 \(\mathbf{P}\) 下封闭，该直积在 \(\mathcal{V}\) 中，其子代数 \(\mathbf{F}\) 也在 \(\mathcal{V}\) 中。最后由 \(\mathbf{H}\)-封闭，\(\mathbf{A} = \varphi(\mathbf{F}) \in \mathcal{V}\)。\(\blacksquare\)

\subsection{自由代数}\label{ux81eaux7531ux4ee3ux6570}

\textbf{定义 6}（簇中的自由代数）：设 \(\mathcal{V}\) 是 \(\Sigma\)-代数簇，\(X\) 是集合。\(\mathcal{V}\) 在 \(X\) 上的\textbf{自由代数} \(\mathbf{F}_{\mathcal{V}}(X)\) 是 \(\Sigma\)-代数，配备映射 \(\eta_X: X \to F_{\mathcal{V}}(X)\)，满足普遍性质：对任意 \(\mathbf{A} \in \mathcal{V}\) 和任意映射 \(f: X \to A\)，存在唯一的 \(\Sigma\)-同态 \(\bar{f}: \mathbf{F}_{\mathcal{V}}(X) \to \mathbf{A}\) 使得 \(\bar{f} \circ \eta_X = f\)。

\textbf{定理 2}（自由代数的构造）：对任意代数簇 \(\mathcal{V} = \mathrm{Mod}(E)\) 和任意集合 \(X\)，自由代数 \(\mathbf{F}_{\mathcal{V}}(X)\) 存在且在唯一同构意义下唯一。具体构造为：

\[\mathbf{F}_{\mathcal{V}}(X) = T_\Sigma(X) / \theta_E\]

其中 \(\theta_E\) 是由等式 \(E\) 诱导的最小同余关系：

\[\theta_E = \bigcap \{\theta \in \mathrm{Con}(T_\Sigma(X)) \mid T_\Sigma(X)/\theta \in \mathcal{V}\}\]

\textbf{证明}：首先构造项代数 \(T_\Sigma(X)\)。定义 \(\theta_{\mathcal{V}}\) 为使得商代数属于 \(\mathcal{V}\) 的所有同余的交。由于 \(\mathcal{V}\) 中代数的直积仍在 \(\mathcal{V}\) 中（Birkhoff定理），且商代数 \(T_\Sigma(X)/\theta_{\mathcal{V}}\) 可嵌入所有这些直积，因此 \(T_\Sigma(X)/\theta_{\mathcal{V}} \in \mathcal{V}\)。设 \(\mathbf{F}_{\mathcal{V}}(X) = T_\Sigma(X)/\theta_{\mathcal{V}}\)，\(\eta_X(x) = [x]_{\theta_{\mathcal{V}}}\)。

验证普遍性质：设 \(\mathbf{A} \in \mathcal{V}\)，\(f: X \to A\)。由 \(T_\Sigma(X)\) 的自由性，\(f\) 唯一延拓为 \(\hat{f}: T_\Sigma(X) \to \mathbf{A}\)。令 \(\ker \hat{f}\) 为同余核，则 \(T_\Sigma(X)/\ker \hat{f} \cong \mathrm{Im}(\hat{f}) \leq \mathbf{A}\)，故在 \(\mathcal{V}\) 中（\(\mathbf{S}\)-封闭）。因此 \(\theta_{\mathcal{V}} \subseteq \ker \hat{f}\)，\(\hat{f}\) 通过商降为 \(\bar{f}: \mathbf{F}_{\mathcal{V}}(X) \to \mathbf{A}\)。唯一性由 \(\eta_X(X)\) 生成 \(\mathbf{F}_{\mathcal{V}}(X)\) 保证。\(\blacksquare\)

\subsection{同余格与置换性}\label{ux540cux4f59ux683cux4e0eux7f6eux6362ux6027}

\textbf{定义 7}（同余格）：\(\Sigma\)-代数 \(\mathbf{A}\) 的所有同余关系在包含序下构成完备格 \(\mathrm{Con}(\mathbf{A})\)，其交为集合论交，并为所有包含给合之同余的交。

\textbf{定义 8}（代数簇）：设 \(\mathcal{V}\) 是 \(\Sigma\)-代数簇。对任意 \(\mathbf{A} \in \mathcal{V}\)，\(\mathrm{Con}(\mathbf{A})\) 是完备格。

\textbf{定义 9}（同余交换性）：代数 \(\mathbf{A}\) 称为\textbf{同余交换的}，若对任意 \(\theta_1, \theta_2 \in \mathrm{Con}(\mathbf{A})\)，有 \(\theta_1 \circ \theta_2 = \theta_2 \circ \theta_1\)。一个簇 \(\mathcal{V}\) 称为同余交换的，若其中每个代数都是同余交换的。

\textbf{定理 3}（同余交换性的等价格）：\(\mathbf{A}\) 中 \(\theta_1 \circ \theta_2 = \theta_2 \circ \theta_1\) 对所有 \(\theta_1, \theta_2 \in \mathrm{Con}(\mathbf{A})\) 成立，当且仅当 \(\mathrm{Con}(\mathbf{A})\) 中的并满足 \(\theta_1 \vee \theta_2 = \theta_1 \circ \theta_2\)。

\textbf{证明}：一般地，\(\theta_1 \vee \theta_2 = \bigcup_{n \ge 1} (\theta_1 \circ \theta_2)^n\)（关系的复合迭代）。若 \(\theta_1 \circ \theta_2 = \theta_2 \circ \theta_1\)，则 \((\theta_1 \circ \theta_2)^2 = \theta_1 \circ (\theta_2 \circ \theta_1) \circ \theta_2 = \theta_1 \circ (\theta_1 \circ \theta_2) \circ \theta_2 = \theta_1 \circ \theta_2\)（由传递性 \(\theta_i \circ \theta_i \subseteq \theta_i\) 和交换性），故需递推到任意 \(n\) 知 \((\theta_1 \circ \theta_2)^n = \theta_1 \circ \theta_2\)，从而 \(\theta_1 \vee \theta_2 = \theta_1 \circ \theta_2\)。反之亦然。\(\blacksquare\)

\subsection{Malcev条件}\label{malcevux6761ux4ef6}

\textbf{定义 10}（Malcev项与Malcev条件）：一个代数簇 \(\mathcal{V}\) 具有\textbf{同余交换性}当且仅当存在三元项 \(p(x, y, z)\)（称为 \textbf{Malcev项}）满足：

\[p(x, x, y) \approx y, \quad p(x, y, y) \approx x\]

对所有 \(\mathbf{A} \in \mathcal{V}\) 成立。具有 Malcev 项的簇称为 \textbf{Malcev簇}。

\textbf{定理 4}（Malcev刻画定理）：\(\Sigma\)-代数簇 \(\mathcal{V}\) 是同余交换的当且仅当 \(\mathcal{V}\) 是 Malcev 簇。

\textbf{证明}：（\(\Leftarrow\)）设 \(\mathcal{V}\) 有 Malcev 项 \(p\)。对 \(\mathbf{A} \in \mathcal{V}\) 和 \(\theta_1, \theta_2 \in \mathrm{Con}(\mathbf{A})\)，若 \((a, c) \in \theta_1 \circ \theta_2\)，即存在 \(b\) 使得 \(a \equiv_{\theta_1} b \equiv_{\theta_2} c\)。计算：

\[a = p(a, b, b) \equiv_{\theta_1} p(a, a, b) = b \equiv_{\theta_2} p(b, c, c) = b = p(b, b, c) \equiv_{\theta_1} p(b, c, c) = p(a, b, b) = a\]

更简洁地：\(a = p(a, b, b) \equiv_{\theta_2} p(a, b, c) \equiv_{\theta_1} p(b, b, c) = c\)，故 \((a, c) \in \theta_2 \circ \theta_1\)。因此 \(\theta_1 \circ \theta_2 \subseteq \theta_2 \circ \theta_1\)。由对称性得相等。

（\(\Rightarrow\)）设 \(\mathcal{V}\) 同余交换。考虑自由代数 \(\mathbf{F} = \mathbf{F}_{\mathcal{V}}(\{x, y, z\})\) 及其上的两个同余 \(\theta_1 = \mathrm{Cg}(x, y)\) 和 \(\theta_2 = \mathrm{Cg}(y, z)\)（由生成对 \((x,y)\) 和 \((y,z)\) 生成的主同余）。则有 \((x, z) \in \theta_1 \circ \theta_2\)（通过 \(y\)），故由交换性 \((x, z) \in \theta_1 \vee \theta_2 = \theta_2 \circ \theta_1\)。存在 \(t \in F\) 使得 \(x \equiv_{\theta_2} t \equiv_{\theta_1} z\)。\(t\) 的自由变量的表达式即为所求的 Malcev 项 \(p(x, y, z)\)。由构造验证 \(p(x,x,y) = y\) 和 \(p(x,y,y) = x\)。\(\blacksquare\)

\textbf{例子}：群簇是 Malcev 簇，Malcev 项为 \(p(x,y,z) = xy^{-1}z\)。环簇也是 Malcev 簇，项为 \(p(x,y,z) = x - y + z\)。

\subsection{Lawvere理论}\label{lawvereux7406ux8bba}

\textbf{定义 11}（Lawvere理论）：一个 \textbf{Lawvere理论} \(\mathcal{T}\) 是一个小范畴，其对象恰为自然数 \(0, 1, 2, \ldots\)，使得每个对象 \(n\) 是 \(\mathbb{N}\) 中 \(n\) 个对象 \(1\) 的直积（即对 \(n\) 有自然同构 \(n \cong 1 \times 1 \times \cdots \times 1\)（\(n\)次）的范畴论意义）。特别地，\(\mathcal{T}\) 是有有限积的范畴，且每个对象是固定对象 \(1\) 的某个幂。

\textbf{定义 12}（Lawvere理论的模型）：Lawvere理论 \(\mathcal{T}\) 的一个\textbf{模型}是积保持函子 \(M: \mathcal{T} \to \mathbf{Set}\)。模型之间的态射是自然变换。所有模型形成范畴 \(\mathrm{Mod}(\mathcal{T}) = \mathbf{Prod}(\mathcal{T}, \mathbf{Set})\)。

\textbf{定理 5}（代数簇与Lawvere理论的等价）：存在范畴等价：

\[\{\Sigma\text{-代数簇}\} \simeq \{\text{Lawvere理论}\}\]

具体地：给定代数簇 \(\mathcal{V}\)，对应的 Lawvere 理论 \(\mathcal{T}_{\mathcal{V}}\) 定义为： - \(\mathcal{T}_{\mathcal{V}}(m, n) = \mathrm{Hom}_{\mathcal{V}}(\mathbf{F}_{\mathcal{V}}(m), \mathbf{F}_{\mathcal{V}}(n))\)，其中 \(\mathbf{F}_{\mathcal{V}}(k)\) 是 \(k\) 个生成元的自由代数。

反过来，给定 Lawvere 理论 \(\mathcal{T}\)，其模型范畴 \(\mathrm{Mod}(\mathcal{T})\) 是一个代数簇。

\textbf{证明}：首先验证 \(\mathcal{T}_{\mathcal{V}}\) 是 Lawvere 理论。\(\mathcal{T}_{\mathcal{V}}\) 的对象为自然数，态射 \(\mathcal{T}_{\mathcal{V}}(m, n) = \mathrm{Hom}_{\mathcal{V}}(\mathbf{F}_{\mathcal{V}}(m), \mathbf{F}_{\mathcal{V}}(n))\)。由自由代数性质，\(\mathbf{F}_{\mathcal{V}}(n) \cong \mathbf{F}_{\mathcal{V}}(1)^n\)（\(n\) 次积），因为 \(\mathrm{Hom}_{\mathcal{V}}(\mathbf{F}_{\mathcal{V}}(1)^n, A) \cong \prod_{i=1}^n \mathrm{Hom}_{\mathcal{V}}(\mathbf{F}_{\mathcal{V}}(1), A) \cong A^n \cong \mathrm{Hom}_{\mathcal{V}}(\mathbf{F}_{\mathcal{V}}(n), A)\)，由 Yoneda 引理得同构。其次，定义函子 \(\Phi: \mathcal{V} \to \mathrm{Mod}(\mathcal{T}_{\mathcal{V}})\)：对 \(\mathbf{A} \in \mathcal{V}\)，\(M_{\mathbf{A}}(n) = \mathrm{Hom}_{\mathcal{V}}(\mathbf{F}_{\mathcal{V}}(n), \mathbf{A}) \cong A^n\)。\(M_{\mathbf{A}}\) 保有限积因 \(M_{\mathbf{A}}(n \times m) \cong A^{n+m} \cong A^n \times A^m\)。反之，对模型 \(M\)，定义 \(A = M(1)\) 上的运算：对每个 \(n\) 元运算符号 \(f\)（对应 \(\mathcal{T}_{\mathcal{V}}(n, 1)\) 中的态射），\(f^A = M(f): A^n \to A\)。这对 \(\Phi\) 建立了范畴等价。\(\blacksquare\)

\subsection{补充：代数的替换性质与子直积}\label{ux8865ux5145ux4ee3ux6570ux7684ux66ffux6362ux6027ux8d28ux4e0eux5b50ux76f4ux79ef}

\textbf{定义 13}（子直积）：代数 \(\mathbf{A}\) 称为一族代数 \(\{\mathbf{A}_i\}_{i \in I}\) 的\textbf{子直积}，若存在单同态 \(\mathbf{A} \hookrightarrow \prod_{i \in I} \mathbf{A}_i\)，使得每个投射 \(\pi_i|_{\mathbf{A}}: \mathbf{A} \to \mathbf{A}_i\) 是满射。

\textbf{定理 6}（Birkhoff子直积定理）：代数簇 \(\mathcal{V}\) 中的每个代数同构于 \(\mathcal{V}\) 中\textbf{子直不可约}（subdirectly irreducible）代数的子直积。子直不可约代数恰是那些有最小非零同余的代数。

\textbf{证明}：设 \(\mathbf{A} \in \mathcal{V}\)。考虑 \(\operatorname{Con}(\mathbf{A})\) 中所有非零同余的交不可约元集合 \(S\)。对每个 \(a \neq b \in A\)，存在 \(\theta_{a,b} \in \operatorname{Con}(\mathbf{A})\) 为包含 \((a,b)\) 的最小同余。由 Zorn 引理，可找到交不可约同余 \(\theta_{a,b}' \geq \theta_{a,b}\) 且 \((a,b) \notin \theta_{a,b}'\) 的极大元。令 \(S = \{\theta_{a,b}' : a \neq b\}\)。则 \(\bigwedge_{\theta \in S} \theta = 0_{\mathbf{A}}\)（同余的交为零）。对每个 \(\theta \in S\)，\(\mathbf{A}/\theta\) 是子直不可约的（因 \(\theta\) 是交不可约同余，\(\mathbf{A}/\theta\) 有最小非零同余）。自然映射 \(\mathbf{A} \to \prod_{\theta \in S} \mathbf{A}/\theta\) 是单射（因核为 \(\bigwedge \theta = 0\)），且每个投影为满射。故 \(\mathbf{A}\) 同构于子直不可约代数的子直积。\(\blacksquare\)

\textbf{注记}：子直不可约性是代数结构论的基石：群簇中，子直不可约群含于有限群或特定单群的变体；环簇中，子直不可约环与McCoy-素环相关。在Lawvere理论的对偶视角下，代数簇的研究等价于有限积完备范畴上的模型理论研究，这给出了Birkhoff定理的深层范畴论基础。

\subsection{经典簇举例与拓展方向}\label{ux7ecfux5178ux7c07ux4e3eux4f8bux4e0eux62d3ux5c55ux65b9ux5411}

群簇、环簇、格簇和模簇均为泛代数的核心研究对象。群簇中，交换群簇由 \(xy \approx yx\) 定义；幂零群簇由特定换位子恒等式定义。环簇的Malcev项 \(x - y + z\) 提供了丰富的同余理论。泛代数与范畴论的交汇------Lawvere理论、代数理论和monad的等价性------构成现代数学基础的三角结构：代数簇 \(\leftrightarrow\) Lawvere理论 \(\leftrightarrow\) 集合范畴上的有限秩monad。这一对应将Birkhoff定理推广到更一般的范畴上下文中，是Stone对偶和Gabriel-Ulmer对偶的自然推广。

这一三角结构的核心洞察在于：\textbf{代数簇}（通过 Birkhoff 的 HSP 定理刻画）等价于\textbf{Lawvere 理论}（以自然数为对象、有限积为结构的范畴），而 Lawvere 理论又等价于集合范畴 \(\mathbf{Set}\) 上的\textbf{有限秩 monad}------即函子 \(T: \mathbf{Set} \to \mathbf{Set}\) 配备自然变换 \(\eta: \operatorname{id} \to T\) 和 \(\mu: T^2 \to T\) 满足 monad 律，且 \(T\) 保持滤过余极限。在代数几何中，代数簇的函子观点（即概形作为函子 \(\operatorname{Ring} \to \mathbf{Set}\)）与此框架高度平行：代数簇的''模型''------即代数结构------在任意有限积完备范畴中（不仅仅是 \(\mathbf{Set}\)）都有意义。这引出了\textbf{拓扑代数}（在 \(\mathbf{Top}\) 中的模型）、\textbf{可微代数}（在光滑流形范畴中的模型）和\textbf{导出代数几何}（在单纯集或链复形范畴中的模型）等丰富的拓展方向。此外，\textbf{半群簇}、\textbf{半环簇}和\textbf{Quantale 簇}（在量子逻辑和理论计算机科学中有应用）是活跃的研究前沿。在形而上学的层面，代数簇理论为''什么是代数结构''这一基本问题提供了精确的数学模型------一种代数结构即是一个 Lawvere 理论，而该理论在 \(\mathbf{Set}\) 中的模型即为我们通常理解的代数结构实例。

\newpage

\chapter{05-代数/02-抽象代数一/05-泛代数与一般代数系统.md}\label{ux4ee3ux657002-ux62bdux8c61ux4ee3ux6570ux4e0005-ux6cdbux4ee3ux6570ux4e0eux4e00ux822cux4ee3ux6570ux7cfbux7edf.md}

\chapter{泛代数与一般代数系统}\label{ux6cdbux4ee3ux6570ux4e0eux4e00ux822cux4ee3ux6570ux7cfbux7edf}

泛代数是代数学中研究一般代数结构的分支，其核心思想是将群、环、模、格等具体代数结构统一在''带运算的集合''这一框架下。该学科起源于Garrett Birkhoff在20世纪30年代的工作，后经Alfred Tarski、Anatoly Malcev、Bjarni Jónsson等学者的推动，发展成为连接代数、范畴论与逻辑的桥梁。泛代数的基本工具包括：签名与项代数提供代数结构的语法描述，同余关系与商代数实现结构的模约化，而簇与等式逻辑则刻画了由等式公理定义的代数类。本章系统地阐述泛代数的基础理论，从代数的语法构造到Birkhoff的HSP定理，再到Malcev条件与自然对偶理论，展现一般代数系统的统一图景。

\hrulefill

\section{第122章：泛代数基础}\label{ux7b2c122ux7ae0ux6cdbux4ee3ux6570ux57faux7840}

泛代数以''签名''（signature）为出发点，将一切代数结构的运算符号化。所谓代数结构，即在一个基础集上装配有限多个满足特定公理的运算。泛代数的任务是将这些结构的共性提炼成一般理论，沟通群论、环论、格论等具体分支。本章从最基本的概念------签名与代数、项代数与自由代数、同余关系与商代数------出发，最终建立泛代数框架下的同态定理。

\subsection{x.1 代数结构的签名与代数}\label{x.1-ux4ee3ux6570ux7ed3ux6784ux7684ux7b7eux540dux4e0eux4ee3ux6570}

\textbf{定义}：一个签名（signature）是一个函数 \(\Sigma: \Omega \to \mathbb{N}_0\)，其中 \(\Omega\) 是一族运算符号的集合，\(\Sigma(f)\) 表示运算 \(f\) 的元数（arity）。\(\Sigma_n = \{f \in \Omega \mid \Sigma(f) = n\}\) 记所有 \(n\) 元运算符号的集合。

\textbf{定义}：一个 \(\Sigma\)-代数是一个对 \((A, F)\)，其中 \(A\) 是一个非空集合（基集，universe），\(F\) 为每个 \(f \in \Omega\) 指派 \(A\) 上的一个运算 \(f^A: A^{\Sigma(f)} \to A\)。当上下文清楚时，简记为 \(\mathbf{A}\)。

群、环、格等皆为特定签名下的代数。例如群签名 \(\Sigma_{\mathrm{Grp}}\) 包含一个二元运算 \(\cdot\)、一个一元运算 \({}^{-1}\) 和一个零元运算（常量）\(e\)。环签名在加法群的基础上增加乘法运算。

\textbf{定义}：设 \(\mathbf{A}, \mathbf{B}\) 为 \(\Sigma\)-代数。映射 \(\varphi: A \to B\) 称为同态，若对每个 \(f \in \Omega\) 及所有 \(a_1, \ldots, a_n \in A\)（\(n = \Sigma(f)\)），有

\[\varphi(f^A(a_1, \ldots, a_n)) = f^B(\varphi(a_1), \ldots, \varphi(a_n)).\]

同态的单射、满射、双射分别称为嵌入、满同态和同构。

\subsection{x.2 项代数与自由代数}\label{x.2-ux9879ux4ee3ux6570ux4e0eux81eaux7531ux4ee3ux6570}

给定签名 \(\Sigma\) 和变元集合 \(X\)，可递归地构造项（term）的集合 \(T_\Sigma(X)\)。直观上，一个项是一个形式表达式，由变元和运算符号按语法规则构成。

\textbf{定义}：\(T_\Sigma(X)\) 是满足以下条件的最小集合：（i）\(X \subseteq T_\Sigma(X)\)；（ii）若 \(f \in \Sigma\)，\(\Sigma(f) = n\)，且 \(t_1, \ldots, t_n \in T_\Sigma(X)\)，则 \(f(t_1, \ldots, t_n) \in T_\Sigma(X)\)。

可在 \(T_\Sigma(X)\) 上以自然方式定义运算 \(f^{T_\Sigma(X)}(t_1, \ldots, t_n) = f(t_1, \ldots, t_n)\)，从而得到 \(\Sigma\)-代数 \(\mathbf{T}_\Sigma(X)\)，称为项代数（term algebra）。

\textbf{定义}：设 \(\mathbf{A}\) 为 \(\Sigma\)-代数。称 \(\mathbf{A}\) 是 \(X\) 上的自由代数，若存在包含映射 \(\iota: X \hookrightarrow A\)，使得对任意 \(\Sigma\)-代数 \(\mathbf{B}\) 和任意映射 \(\alpha: X \to B\)，存在唯一的同态 \(\bar{\alpha}: \mathbf{A} \to \mathbf{B}\) 使得 \(\bar{\alpha} \circ \iota = \alpha\)。

\textbf{定理 x.1}（项代数的自由性）：项代数 \(\mathbf{T}_\Sigma(X)\) 是 \(X\) 上的自由 \(\Sigma\)-代数。

\textbf{证明}：取 \(\iota: X \to T_\Sigma(X)\) 为自然包含。设 \(\alpha: X \to B\) 为任意映射。递归地定义 \(\bar{\alpha}: T_\Sigma(X) \to \mathbf{B}\)：（i）对 \(x \in X\)，\(\bar{\alpha}(x) = \alpha(x)\)；（ii）对 \(f \in \Sigma\)，\(\bar{\alpha}(f(t_1, \ldots, t_n)) = f^B(\bar{\alpha}(t_1), \ldots, \bar{\alpha}(t_n))\)。由 \(T_\Sigma(X)\) 的归纳定义，\(\bar{\alpha}\) 是良定义的且唯一确定。验证 \(\bar{\alpha}\) 是同态：对 \(f \in \Sigma\) 及 \(t_1, \ldots, t_n \in T_\Sigma(X)\)，

\[\bar{\alpha}(f^{T_\Sigma(X)}(t_1, \ldots, t_n)) = \bar{\alpha}(f(t_1, \ldots, t_n)) = f^B(\bar{\alpha}(t_1), \ldots, \bar{\alpha}(t_n)) = f^B(\bar{\alpha}(t_1), \ldots, \bar{\alpha}(t_n)).\]

故 \(\bar{\alpha}\) 为同态且 \(\bar{\alpha}|_X = \alpha\)。唯一性由 \(\bar{\alpha}\) 的递归定义决定。 \(\blacksquare\)

\subsection{x.3 同余关系与商代数}\label{x.3-ux540cux4f59ux5173ux7cfbux4e0eux5546ux4ee3ux6570}

\textbf{定义}：\(\Sigma\)-代数 \(\mathbf{A}\) 上的同余关系（congruence）是一个等价关系 \(\theta \subseteq A \times A\)，使得对每个 \(f \in \Sigma\)（设 \(\Sigma(f) = n\)），若 \((a_i, b_i) \in \theta\)（\(i = 1, \ldots, n\)），则

\[(f^A(a_1, \ldots, a_n), f^A(b_1, \ldots, b_n)) \in \theta.\]

换言之，同余关系是与所有代数运算相容的等价关系。

\textbf{定理 x.2}（商代数）：设 \(\theta\) 是 \(\mathbf{A}\) 上的同余关系。在商集 \(A/\theta\) 上定义运算

\[f^{A/\theta}([a_1]_\theta, \ldots, [a_n]_\theta) = [f^A(a_1, \ldots, a_n)]_\theta,\]

则 \((A/\theta, F^{A/\theta})\) 构成一个 \(\Sigma\)-代数，称为 \(\mathbf{A}\) 模 \(\theta\) 的商代数。自然投影 \(\pi: A \to A/\theta\) 是满同态，其核为 \(\theta\)。

\textbf{证明}：运算的良定义性由同余关系的相容性保证。商代数公理显然成立，因为 \(A/\theta\) 上的运算通过代表元定义且与代表元选取无关。自然投影的同态性质直接由商代数上运算的定义得到：\(\pi(f^A(a_1, \ldots, a_n)) = [f^A(a_1, \ldots, a_n)]_\theta = f^{A/\theta}([a_1]_\theta, \ldots, [a_n]_\theta) = f^{A/\theta}(\pi(a_1), \ldots, \pi(a_n))\)。 \(\blacksquare\)

\subsection{x.4 同态定理}\label{x.4-ux540cux6001ux5b9aux7406}

泛代数中的同态定理将群论和环论中熟知的同构定理推广到一般代数结构。

\textbf{定理 x.3}（泛代数第一同态定理）：设 \(\varphi: \mathbf{A} \to \mathbf{B}\) 为 \(\Sigma\)-代数同态。定义 \(\ker \varphi = \{(a_1, a_2) \in A \times A \mid \varphi(a_1) = \varphi(a_2)\}\)，则 \(\ker \varphi\) 是 \(\mathbf{A}\) 上的同余关系。且存在唯一的单同态 \(\tilde{\varphi}: A/\ker\varphi \to \mathbf{B}\) 使得以下图表交换：

\[\begin{tikzcd}
\mathbf{A} \ar[r,"\varphi"] \ar[d,"\pi"'] & \mathbf{B} \\
\mathbf{A}/\ker\varphi \ar[ru,"\tilde{\varphi}"']
\end{tikzcd}\]

特别地，\(\mathbf{A}/\ker\varphi \cong \operatorname{im}\varphi\)。

\textbf{证明}：首先验证 \(\ker\varphi\) 是同余关系。\(\ker\varphi\) 显然是等价关系。对 \(f \in \Sigma\)（元数 \(n\)）及 \((a_i, b_i) \in \ker\varphi\)，有

\[\varphi(f^A(a_1, \ldots, a_n)) = f^B(\varphi(a_1), \ldots, \varphi(a_n)) = f^B(\varphi(b_1), \ldots, \varphi(b_n)) = \varphi(f^A(b_1, \ldots, b_n)),\]

故 \((f^A(a_1, \ldots, a_n), f^A(b_1, \ldots, b_n)) \in \ker\varphi\)。现定义 \(\tilde{\varphi}([a]) = \varphi(a)\)。若 \([a_1] = [a_2]\)，则 \((a_1, a_2) \in \ker\varphi\)，故 \(\varphi(a_1) = \varphi(a_2)\)，\(\tilde{\varphi}\) 良定义。\(\tilde{\varphi}\) 是同态：

\[\tilde{\varphi}(f^{A/\ker\varphi}([a_1], \ldots, [a_n])) = \tilde{\varphi}([f^A(a_1, \ldots, a_n)]) = \varphi(f^A(a_1, \ldots, a_n)) = f^B(\varphi(a_1), \ldots, \varphi(a_n)) = f^B(\tilde{\varphi}([a_1]), \ldots, \tilde{\varphi}([a_n])).\]

若 \(\tilde{\varphi}([a_1]) = \tilde{\varphi}([a_2])\)，则 \(\varphi(a_1) = \varphi(a_2)\)，即 \((a_1, a_2) \in \ker\varphi\)，故 \([a_1] = [a_2]\)，\(\tilde{\varphi}\) 为单射。因此 \(\mathbf{A}/\ker\varphi \cong \operatorname{im}\varphi\)。 \(\blacksquare\)

\textbf{定理 x.4}（对应定理）：设 \(\theta\) 是 \(\mathbf{A}\) 上的同余关系。则 \(\mathbf{A}/\theta\) 的所有同余关系与 \(\mathbf{A}\) 上包含 \(\theta\) 的同余关系之间存在一一对应。具体地，若 \(\psi \supseteq \theta\) 是 \(\mathbf{A}\) 的同余关系，则 \(\psi/\theta = \{([a]_\theta, [b]_\theta) \mid (a, b) \in \psi\}\) 是 \(\mathbf{A}/\theta\) 的同余关系。

\textbf{证明}：定义一个映射 \(\Phi\) 将 \(\mathbf{A}\) 上包含 \(\theta\) 的同余关系 \(\psi\) 映为 \(\mathbf{A}/\theta\) 的同余关系 \(\psi/\theta\)。首先验证 \(\psi/\theta\) 是良定义的：若 \((a, b) \in \psi\) 且 \([a]_\theta = [a']_\theta\)，\([b]_\theta = [b']_\theta\)，则 \((a, a'), (b, b') \in \theta \subseteq \psi\)。由传递性 \((a', b') \in \psi\)。容易验证 \(\psi/\theta\) 是等价关系且与运算相容。\(\Phi\) 显然是单射。满射性：给定 \(\mathbf{A}/\theta\) 上的同余关系 \(\eta\)，定义 \(\psi = \{(a, b) \mid ([a]_\theta, [b]_\theta) \in \eta\}\)，显然 \(\psi \supseteq \theta\) 且为同余关系，且 \(\Phi(\psi) = \eta\)。 \(\blacksquare\)

\hrulefill

\section{第123章：簇与等式逻辑}\label{ux7b2c123ux7ae0ux7c07ux4e0eux7b49ux5f0fux903bux8f91}

簇（variety）是泛代数的核心研究对象之一。直观上，一个簇是由一族等式公理定义的代数类，例如群、环、格等皆为簇。Birkhoff在1935年证明了著名的HSP定理：一个代数类构成簇当且仅当它对同态像（H）、子代数（S）和直积（P）封闭。这一定理将语法概念（等式定义）与语义概念（HSP封闭）等价起来，奠定了泛代数的基石。

\subsection{x+1.1 等式与代数类}\label{x1.1-ux7b49ux5f0fux4e0eux4ee3ux6570ux7c7b}

\textbf{定义}：一个 \(\Sigma\)-等式是一个形如 \(t \approx s\) 的公式，其中 \(t, s \in T_\Sigma(X)\)。\(\Sigma\)-代数 \(\mathbf{A}\) 满足等式 \(t \approx s\)，记作 \(\mathbf{A} \models t \approx s\)，若对任意赋值 \(v: X \to A\)，有 \(\bar{v}(t) = \bar{v}(s)\)（其中 \(\bar{v}\) 为 \(v\) 的唯一同态扩张）。

\textbf{定义}：设 \(\mathcal{E}\) 为一族 \(\Sigma\)-等式。\(\mathcal{E}\) 定义的簇（variety）为

\[\mathcal{V}(\mathcal{E}) = \{\mathbf{A} \in \operatorname{Alg}(\Sigma) \mid \mathbf{A} \models e \ (\forall e \in \mathcal{E})\}.\]

一个代数类 \(\mathcal{K}\) 称为簇，若存在等式集 \(\mathcal{E}\) 使 \(\mathcal{K} = \mathcal{V}(\mathcal{E})\)。

\subsection{x+1.2 算子H、S、P}\label{x1.2-ux7b97ux5b50hsp}

\textbf{定义}：对于代数类 \(\mathcal{K}\)，定义以下算子：

\begin{itemize}
\tightlist
\item
  \(H(\mathcal{K}) = \{\mathbf{B} \mid \exists \mathbf{A} \in \mathcal{K} \text{ 和满同态 } \varphi: \mathbf{A} \twoheadrightarrow \mathbf{B}\}\) （同态像）
\item
  \(S(\mathcal{K}) = \{\mathbf{B} \mid \exists \mathbf{A} \in \mathcal{K} \text{ 和嵌入 } \iota: \mathbf{B} \hookrightarrow \mathbf{A}\}\) （子代数）
\item
  \(P(\mathcal{K}) = \{\prod_{i \in I} \mathbf{A}_i \mid \mathbf{A}_i \in \mathcal{K}\}\) （直积）
\end{itemize}

\textbf{定理 x+1.1}（Birkhoff HSP定理）：设 \(\mathcal{K}\) 为 \(\Sigma\)-代数类，则 \(\mathcal{K}\) 是簇当且仅当 \(\mathcal{K} = HSP(\mathcal{K})\)。等价地，一个代数类是簇当且仅当它对 \(H\)、\(S\)、\(P\) 封闭。

\textbf{证明}：（\(\Rightarrow\)）设 \(\mathcal{K} = \mathcal{V}(\mathcal{E})\) 为簇。需证 \(HSP(\mathcal{K}) \subseteq \mathcal{K}\)。

（H封闭）设 \(\mathbf{B} \in H(\mathcal{K})\)，则存在满同态 \(\varphi: \mathbf{A} \twoheadrightarrow \mathbf{B}\) 其中 \(\mathbf{A} \models \mathcal{E}\)。对任意 \(t \approx s \in \mathcal{E}\) 和赋值 \(v: X \to B\)，取提升 \(\tilde{v}: X \to A\)（因同态满，选取原像），有 \(\bar{v}(t) = \varphi(\bar{\tilde{v}}(t)) = \varphi(\bar{\tilde{v}}(s)) = \bar{v}(s)\)，故 \(\mathbf{B} \models \mathcal{E}\)。

（S封闭）设 \(\mathbf{B} \in S(\mathcal{K})\)，则 \(B \subseteq A\) 且运算受限。对任意 \(t \approx s \in \mathcal{E}\) 和赋值 \(v: X \to B \subseteq A\)，在 \(\mathbf{A}\) 中 \(\bar{v}(t) = \bar{v}(s)\)，故在 \(\mathbf{B}\) 中亦然。

（P封闭）设 \(\mathbf{B} = \prod_{i \in I} \mathbf{A}_i\)，\(\mathbf{A}_i \models \mathcal{E}\)。对任意 \(t \approx s \in \mathcal{E}\) 和赋值 \(v: X \to B\)，\(\pi_i \circ \bar{v}\) 是 \(\mathbf{A}_i\) 中的赋值，故 \(\pi_i(\bar{v}(t)) = \pi_i(\bar{v}(s))\) 对所有 \(i\) 成立，从而 \(\bar{v}(t) = \bar{v}(s)\)。

（\(\Leftarrow\)）设 \(\mathcal{K}\) 对 \(H, S, P\) 封闭。定义 \(\operatorname{Th}(\mathcal{K}) = \{t \approx s \mid \forall \mathbf{A} \in \mathcal{K}, \mathbf{A} \models t \approx s\}\) 为 \(\mathcal{K}\) 的理论。令 \(\mathcal{V} = \mathcal{V}(\operatorname{Th}(\mathcal{K}))\)，则 \(\mathcal{K} \subseteq \mathcal{V}\)。关键步骤是证明 \(\mathcal{V} \subseteq HSP(\mathcal{K})\)。

取 \(X\) 为足够大的可数变元集。对每个 \(\mathbf{A} \in \mathcal{K}\)，考虑所有同态 \(\varphi: \mathbf{T}_\Sigma(X) \to \mathbf{A}\) 的图。设 \(\Phi_{\mathbf{A}}\) 为 \(\mathbf{T}_\Sigma(X)\) 上的同余关系，定义为 \(\Phi_{\mathbf{A}} = \{(t, s) \mid \bar{v}(t) = \bar{v}(s) \text{ 对所有赋值 } v: X \to A\}\)。

定义 \(\Phi = \bigcap_{\mathbf{A} \in \mathcal{K}} \Phi_{\mathbf{A}}\)。则 \(\mathbf{F} = \mathbf{T}_\Sigma(X)/\Phi\) 是 \(\mathcal{K}\) 中的自由代数（在簇意义下）。易见 \(\mathbf{F} \in SP(\mathcal{K})\)，因为它是 \(\prod_{\mathbf{A} \in \mathcal{K}, v: X \to A} \mathbf{A}\) 的子代数（通过对角嵌入）。

设 \(\mathbf{B} \in \mathcal{V}\)。取充分大的 \(X\) 使得存在满同态 \(\psi: \mathbf{T}_\Sigma(X) \twoheadrightarrow \mathbf{B}\)。由于 \(\mathbf{B} \models \operatorname{Th}(\mathcal{K})\)，有 \(\Phi \subseteq \ker \psi\)（若 \((t, s) \in \Phi\)，则对所有 \(\mathbf{A} \in \mathcal{K}\) 和赋值，\(\bar{v}(t) = \bar{v}(s)\)，这即是在 \(\mathcal{K}\) 中等式 \(t \approx s\) 普遍成立，从而属于 \(\operatorname{Th}(\mathcal{K})\)。故 \(\mathbf{B} \models t \approx s\)，即 \((t, s) \in \ker \psi\)）。

因此 \(\psi\) 可通过 \(\Phi\) 分解：存在同态 \(\tilde{\psi}: \mathbf{F} \twoheadrightarrow \mathbf{B}\)，其为满射。故 \(\mathbf{B} \in H(\mathbf{F}) \subseteq HSP(\mathbf{F}) \subseteq HSP(SP(\mathcal{K})) = HSP(\mathcal{K})\)。由封闭性 \(HSP(\mathcal{K}) = \mathcal{K}\)。 \(\blacksquare\)

\subsection{x+1.3 等式逻辑与完备性}\label{x1.3-ux7b49ux5f0fux903bux8f91ux4e0eux5b8cux5907ux6027}

\textbf{定理 x+1.2}（等式逻辑的Birkhoff完备性）：设 \(\mathcal{E}\) 为一族等式，\(t \approx s\) 为一个等式。则 \(\mathcal{E} \models t \approx s\)（语义后承，即所有满足 \(\mathcal{E}\) 的代数都满足 \(t \approx s\)）当且仅当 \(t \approx s\) 可由 \(\mathcal{E}\) 经以下五条推理规则在语法上推出：（1）自反性 \(t \approx t\)；（2）对称性；（3）传递性；（4）替换规则：从 \(t_i \approx s_i\)（\(i=1,\ldots,n\)）推出 \(f(t_1,\ldots,t_n) \approx f(s_1,\ldots,s_n)\)；（5）代入规则：从 \(t \approx s\) 推出替换变元后的等式。

\textbf{证明}：完备性的证明核心在于构造自由代数 \(\mathbf{F}_{\mathcal{V}(\mathcal{E})}(X) = \mathbf{T}_\Sigma(X)/\Theta\)，其中 \(\Theta\) 为由 \(\mathcal{E}\) 语法导出的最大同余关系。语法后承关系 \(\mathcal{E} \vdash t \approx s\) 恰好对应于 \((t, s) \in \Theta\)。若 \(\mathcal{E} \models t \approx s\)，则在 \(\mathbf{F}_{\mathcal{V}(\mathcal{E})}(X)\) 中 \(t\) 和 \(s\) 在自然赋值下相等（因为 \(\mathbf{F}_{\mathcal{V}(\mathcal{E})}(X) \models \mathcal{E}\)），故 \(t \Theta s\)，即 \(\mathcal{E} \vdash t \approx s\)。反向（可靠性）为显然的。 \(\blacksquare\)

\hrulefill

\section{第124章：Malcev条件与对偶理论}\label{ux7b2c124ux7ae0malcevux6761ux4ef6ux4e0eux5bf9ux5076ux7406ux8bba}

代数簇的分类理论是泛代数的核心课题之一。不同的代数簇具有截然不同的性质：有的簇同余关系可换，有的簇具有模同余格，有的簇满足分配律。Malcev发现这些性质可以由簇中等式条件（即存在满足特定等式的项）来刻画，从而开辟了Malcev条件理论。与此同时，对偶理论为代数范畴与拓扑范畴之间建立了桥梁，Stone对偶和Priestley对偶在泛代数框架下获得了统一的表述。

\subsection{x+2.1 Malcev条件与同余可换性}\label{x2.1-malcevux6761ux4ef6ux4e0eux540cux4f59ux53efux6362ux6027}

\textbf{定义}：代数 \(\mathbf{A}\) 称为同余可换的（congruence-permutable），若对 \(\mathbf{A}\) 上任意两个同余关系 \(\alpha, \beta\)，有 \(\alpha \circ \beta = \beta \circ \alpha\)。一个簇称为同余可换的，若其中每个代数皆同余可换。

\textbf{定理 x+2.1}（Malcev同余可换性定理）：一个簇 \(\mathcal{V}\) 是同余可换的当且仅当存在一个三元项 \(p(x,y,z)\) 使得 \(\mathcal{V}\) 满足等式

\[p(x,x,y) \approx y, \quad p(x,y,y) \approx x.\]

这样的项 \(p\) 称为Malcev项。

\textbf{证明}：（\(\Leftarrow\)）设存在Malcev项 \(p\)。任给 \(\mathbf{A} \in \mathcal{V}\) 和同余关系 \(\alpha, \beta\)。设 \((a, c) \in \alpha \circ \beta\)，则存在 \(b \in A\) 使得 \(a \;\alpha\; b \;\beta\; c\)。考虑 \(d = p(a, b, c)\)。由Malcev等式：

\[p(a, b, b) = a, \quad p(b, b, c) = c.\]

由同余性质：\(p(a, a, c) \;\beta\; p(a, b, c) = d\)（因 \(a \;\alpha\; a\) 和 \(c \;\beta\; c\)），故 \(c = p(b, b, c) \;\beta\; d\)。另一方面，\(d = p(a, b, c) \;\alpha\; p(b, b, c) = c\)。细节上，\(p(a, b, c) \;\alpha\; p(a, a, c) \approx c\)，而 \(p(a, b, c) \;\beta\; p(a, b, b) \approx a\)。因此 \(d \;\alpha\; c\) 且 \(d \;\beta\; a\)，即 \(a \;\beta\; d \;\alpha\; c\)，故 \((a, c) \in \beta \circ \alpha\)。同理可证反向包含。

（\(\Rightarrow\)）构造 \(\mathbf{F} = \mathbf{F}_{\mathcal{V}}(x, y, z)\) 为 \(\mathcal{V}\) 在 \(\{x, y, z\}\) 上的自由代数。定义 \(\alpha = \Theta(x, y)\) 为由 \((x, y)\) 生成的同余关系，\(\beta = \Theta(y, z)\)。由于 \(\mathcal{V}\) 同余可换，有 \((x, z) \in \alpha \circ \beta\)，即存在 \(p \in F\) 使得 \(x \;\alpha\; p \;\beta\; z\)。由于 \(\alpha\) 由 \(x \approx y\) 生成，这意味着在 \(\mathbf{F}\) 中，若将 \(y\) 等同于 \(x\)，则 \(p\) 等同于 \(x\)，即 \(\mathcal{V} \models p(x, x, z) \approx z\)。类似地，\(\mathcal{V} \models p(x, y, y) \approx x\)。\(\blacksquare\)

\subsection{x+2.2 同余模性与可换性}\label{x2.2-ux540cux4f59ux6a21ux6027ux4e0eux53efux6362ux6027}

\textbf{定义}：代数 \(\mathbf{A}\) 的同余格 \(\operatorname{Con}(\mathbf{A})\) 是 \(\mathbf{A}\) 上所有同余关系的集合，按包含关系偏序，构成完备格，其中交为集合交，并为生成子集的传递闭包。

\textbf{定义}：称 \(\mathbf{A}\) 的同余格是模的（modular），若对任意 \(\alpha, \beta, \gamma \in \operatorname{Con}(\mathbf{A})\)，\(\alpha \leq \gamma\) 蕴含 \(\alpha \lor (\beta \land \gamma) = (\alpha \lor \beta) \land \gamma\)。一个簇称为同余模的（congruence-modular），若其中每个代数的同余格皆为模格。

\textbf{定理 x+2.2}（Day同余模性刻画）：一个簇是同余模的当且仅当存在项 \(m_0, \ldots, m_k\)（某 \(k \geq 1\)）使得

\[m_0(x, y, z, u) \approx x, \quad m_k(x, y, z, u) \approx u,\] \[m_i(x, y, x, z) \approx x \quad (i \text{ 为偶数}), \quad m_i(x, x, z, z) \approx m_{i+1}(x, x, z, z) \quad (i \text{ 为奇数}).\]

这些项称为Day项（Day terms），\(k\) 称为Day长度。

\textbf{证明}：证明思想与同余可换性类似。利用自由代数构造，由同余格的模律可推得存在满足上述等式的项。具体地，在 \(\mathbf{F}_{\mathcal{V}}(x, y, z, u)\) 中，考虑由 \((x, y)\) 生成的主同余关系 \(\alpha\) 和由 \((z, u)\) 生成的主同余关系 \(\beta\)，以及 \(\gamma = \Theta(x, z)\)。模律 \((\alpha \lor \beta) \land \gamma \leq \alpha \lor (\beta \land \gamma)\) 的等式表述给出了Day项的存在性。完整的R. Freese和R. McKenzie证明使用了Pixley-Wille条件与Jónsson引理相结合。\(\blacksquare\)

\subsection{x+2.3 Jónsson引理}\label{x2.3-juxf3nssonux5f15ux7406}

\textbf{定理 x+2.3}（Jónsson引理）：设 \(\mathcal{V} = \mathcal{V}(\mathcal{K})\) 是由有限代数类 \(\mathcal{K}\) 生成的同余分配的簇。若 \(\mathbf{A} \in \mathcal{V}\) 是子直不可约的（即 \(\mathbf{A}\) 不是以其真同态像为分量的子直积），则 \(\mathbf{A} \in HSP_U(\mathcal{K})\)，其中 \(P_U\) 表示超积（ultraproduct）。

\textbf{证明}：设 \(\mathbf{A} \in \mathcal{V}\) 是子直不可约的。由于 \(\mathcal{V} = HSP(\mathcal{K})\)，存在代数 \(\mathbf{B}_i \in \mathcal{K}\)（\(i \in I\)）和一个子代数 \(\mathbf{S} \leq \prod_{i \in I} \mathbf{B}_i\) 以及满同态 \(\varphi: \mathbf{S} \twoheadrightarrow \mathbf{A}\)。设 \(\theta = \ker \varphi\)，则 \(\mathbf{S}/\theta \cong \mathbf{A}\)。

由于 \(\mathbf{A}\) 是子直不可约的，\(\theta\) 在 \(\operatorname{Con}(\mathbf{S})\) 中是交不可约的（meet-irreducible）。在同余分配格的条件下，交不可约元必为质元（prime），即对任意 \(\alpha, \beta \in \operatorname{Con}(\mathbf{S})\)，\(\alpha \land \beta \leq \theta\) 蕴含 \(\alpha \leq \theta\) 或 \(\beta \leq \theta\)。

通过Jónsson滤子构造，超积 \(\prod_{I} \mathbf{B}_i/\mathcal{U}\)（对适当的 \(\mathcal{K}\)-完全超滤子 \(\mathcal{U}\)）提供了所需的超积表示。最终推得 \(\mathbf{A} \in H(\prod_I \mathbf{B}_i/\mathcal{U})\)，即属于 \(HSP_U(\mathcal{K})\)。\(\blacksquare\)

\subsection{x+2.4 自然对偶理论简介}\label{x2.4-ux81eaux7136ux5bf9ux5076ux7406ux8bbaux7b80ux4ecb}

自然对偶理论（Natural Duality Theory），由Davey和Werner等人发展，旨在为代数簇与拓扑范畴之间建立范畴对偶。其基本思想是通过''piggyback方法''（背负法）构造对偶。

\textbf{定义}：设 \(\mathcal{A} = \mathcal{ISP}(\mathbf{M})\) 是一个有限代数 \(\mathbf{M}\) 生成的拟簇（quasivariety）。在 \(\mathbf{M}\) 上赋予离散拓扑，构造一个拓扑代数结构 \(\widetilde{\mathbf{M}}\)。自然对偶试图证明范畴对偶 \(\mathcal{A} \simeq \mathcal{X}\)，其中 \(\mathcal{X}\) 是以 \(\widetilde{\mathbf{M}}\) 为生成对象的某些拓扑结构的范畴。

\textbf{Piggyback方法}：为了构造对偶，需在 \(\mathbf{M}\) 上寻找适当的''背负''关系。若 \(\mathbf{M}\) 有一个子代数 \(\omega\)（称为schizophrenic对象），且存在二元关系 \(R \subseteq M \times \omega\)（或更一般的关系），则可通过这些关系来定义对偶范畴中的态射。该方法将代数同态翻译为保持相应关系的连续映射。

\subsection{x+2.5 Stone对偶与Priestley对偶的统一}\label{x2.5-stoneux5bf9ux5076ux4e0epriestleyux5bf9ux5076ux7684ux7edfux4e00}

\textbf{定理 x+2.4}（Stone对偶）：Boolean代数范畴与Stone空间（紧致、Hausdorff、全不连通的拓扑空间）范畴是对偶等价的。

\textbf{定理 x+2.5}（Priestley对偶）：有界分配格范畴与Priestley空间（紧致、全序不连通的偏序拓扑空间）范畴是对偶等价的。

在泛代数的统一框架下，Stone对偶和Priestley对偶可视为自然对偶理论的特例。具体而言：

\begin{itemize}
\tightlist
\item
  Stone对偶对应于生成对象为二元Boolean代数 \(\mathbf{2}\) 的簇，通过Piggyback方法，关系 \(R \subseteq 2 \times \{0,1\}\) 为相等关系 \(=\) 的自然对偶。
\item
  Priestley对偶对应于生成对象为二元分配格 \(\mathbf{2}\) 的簇，通过偏序关系 \(R \subseteq 2 \times \{0,1\}\)（自然序 \(\leq\)）作为背负关系的自然对偶。
\end{itemize}

更一般地，自然对偶理论表明：若 \(\mathbf{M}\) 为有限代数且其生成的拟簇 \(\mathcal{ISP}(\mathbf{M})\) 具有自然对偶，则该对偶可通过 \(\mathbf{M}\) 上的适当关系结构（即''alter ego'' \(\widetilde{\mathbf{M}}\)）来刻画。定理的完整证明依赖于紧致性论证和分离性引理，基本思路如下：定义对偶函子 \(D: \mathcal{A} \to \mathcal{X}\) 将代数 \(\mathbf{A}\) 映为 \(\operatorname{Hom}(\mathbf{A}, \mathbf{M})\)，并将后者赋予从 \(\widetilde{\mathbf{M}}\) 的乘积中继承的子空间拓扑和关系结构；反方向函子 \(E: \mathcal{X} \to \mathcal{A}\) 将拓扑结构 \(X\) 映为从 \(X\) 到 \(\widetilde{\mathbf{M}}\) 的保持关系的连续映射全体。单位与余单位的自然同构性通过Yoneda引理和Schur引理的泛代数版本论证。\(\blacksquare\)

\subsection{x+2.6 子直积与Birkhoff的子直不可约分解定理}\label{x2.6-ux5b50ux76f4ux79efux4e0ebirkhoffux7684ux5b50ux76f4ux4e0dux53efux7ea6ux5206ux89e3ux5b9aux7406}

子直积（subdirect product）是泛代数中将代数分解为''更简单''代数的基本工具，类似于群论中的子直积概念。子直不可约代数在簇的结构理论中扮演着类似于单群在群论中的角色。

\textbf{定义}：代数 \(\mathbf{A}\) 的一个子直分解是一族单射同态 \(\varphi_i: \mathbf{A}_i \hookrightarrow \mathbf{B}_i\) 满足 \(\mathbf{A} \leq \prod_i \mathbf{B}_i\) 且对每个 \(i\)，投影 \(\pi_i|_{\mathbf{A}}\) 是满射（即每个分量都被覆盖）。\(\mathbf{A}\) 称为子直不可约的（subdirectly irreducible），若 \(\mathbf{A}\) 的任意子直分解中必有一个投影为同构（即 \(\mathbf{A}\) 同构于该分量）。

\textbf{定理 x+2.6}（Birkhoff子直不可约分解定理）：任何非平凡的 \(\Sigma\)-代数 \(\mathbf{A}\) 都可表示为子直不可约代数的子直积。换言之，存在一族子直不可约代数 \(\{\mathbf{S}_i\}_{i \in I}\) 使得 \(\mathbf{A}\) 同构于 \(\prod_{i \in I} \mathbf{S}_i\) 的子代数，且每个投影 \(\pi_i: \mathbf{A} \to \mathbf{S}_i\) 为满射。

\textbf{证明}：对 \(\mathbf{A}\) 中的任意两个不同元素 \(a \neq b\)，需构造子直不可约代数 \(\mathbf{S}_{a,b}\) 和满同态 \(\varphi_{a,b}: \mathbf{A} \twoheadrightarrow \mathbf{S}_{a,b}\) 使得 \(\varphi_{a,b}(a) \neq \varphi_{a,b}(b)\)。考虑 \(\mathbf{A}\) 上所有使 \(a\) 和 \(b\) 不在同一等价类中的同余关系，按包含关系偏序。由Zorn引理，存在关于此性质的极大同余关系 \(\theta_{a,b}\)。由极大性，\(\mathbf{A}/\theta_{a,b}\) 是子直不可约的（其最小非平凡同余关系由 \((a/\theta_{a,b}, b/\theta_{a,b})\) 生成）。取 \(\mathbf{S}_{a,b} = \mathbf{A}/\theta_{a,b}\)，\(\varphi_{a,b}\) 为自然投影。

对所有的对 \(\{(a,b) \mid a \neq b \in A\}\)，考虑乘积同态 \(\Phi: \mathbf{A} \to \prod_{a \neq b} \mathbf{S}_{a,b}\)，\(\Phi(x) = (\varphi_{a,b}(x))_{a \neq b}\)。若 \(\Phi(x) = \Phi(y)\)，则对所有 \(a \neq b\) 有 \(\varphi_{a,b}(x) = \varphi_{a,b}(y)\)，特别地取 \((a,b) = (x,y)\) 得矛盾（除非 \(x=y\)）。故 \(\Phi\) 为单射。\(\mathbf{A}\) 同构于 \(\Phi(\mathbf{A})\)，后者是子直不可约代数乘积的子代数且投影皆为满射。\(\blacksquare\)

\textbf{推论}（簇由子直不可约代数确定）：簇 \(\mathcal{V}\) 由其子直不可约成员完全确定，即 \(\mathcal{V} = HSP(\mathcal{V}_{\text{SI}})\)，其中 \(\mathcal{V}_{\text{SI}}\) 为 \(\mathcal{V}\) 中子直不可约代数的全体。

\subsection{x+2.7 同余分配簇与Jónsson术语}\label{x2.7-ux540cux4f59ux5206ux914dux7c07ux4e0ejuxf3nssonux672fux8bed}

\textbf{定义}：一个簇 \(\mathcal{V}\) 称为同余分配的（congruence-distributive），若 \(\mathcal{V}\) 中每个代数的同余格是分配格。

\textbf{定理 x+2.7}（Jónsson同余分配性刻画）：一个簇是同余分配的当且仅当存在三元项 \(d_0, \ldots, d_k\)（\(k \geq 1\)）使得

\[d_0(x,y,z) \approx x, \quad d_k(x,y,z) \approx z,\] \[d_i(x,x,y) \approx x \quad (\forall i), \quad d_i(x,y,y) \approx d_{i+1}(x,y,y) \quad (i \text{ 为偶数}), \quad d_i(x,y,x) \approx d_{i+1}(x,y,x) \quad (i \text{ 为奇数}).\]

这些项称为Jónsson项，\(k\) 称为Jónsson水平。

\textbf{证明}：（\(\Leftarrow\)）设存在Jónsson项。在 \(\mathbf{F}_{\mathcal{V}}(x,y,z)\) 中定义由 \((x,y)\) 生成的同余关系 \(\alpha\)、由 \((y,z)\) 生成的同余关系 \(\beta\)，由 \((x,z)\) 生成的 \(\gamma\)。需证 \(\gamma \leq \alpha \lor \beta\) 及分配律。由Jónsson项的存在性，可逐次证明 \(d_i(x,y,z)\) 经适当的同余关系逼近 \(x\) 和 \(z\)，从而推断分配性。（\(\Rightarrow\)）反向证明类似Malcev条件的经典构造，通过自由代数中的同余格分配性直接构造所需的项。\(\blacksquare\)

这些定理揭示了泛代数中簇结构与等式条件之间的深刻联系。同余可换性（Malcev条件）、同余模性（Day条件）和同余分配性（Jónsson条件）形成了从最强到最弱的代数性质层次，为代数簇的分类提供了系统的方法论基础。

\newpage

\chapter{05-代数/03-抽象代数二/01-Galois理论.md}\label{ux4ee3ux657003-ux62bdux8c61ux4ee3ux6570ux4e8c01-galoisux7406ux8bba.md}

\section{第125章：Galois 理论}\label{ux7b2c125ux7ae0galois-ux7406ux8bba}

Galois 理论是代数学中最为优美和深刻的理论之一，它建立了域扩张与群论之间的精妙对应，从而将方程的根式可解性问题转化为群的可解性问题。这一理论由 Evariste Galois 于 1832 年在决斗前夕的遗书中首次阐述，后经 Liouville、Jordan、Dedekind、Artin 等人的系统整理和发展，成为现代代数学的基石。Galois 理论的核心思想是：一个域扩张的对称性（由 Galois 群刻画）完全决定了该扩张的中间域结构，从而将看似截然不同的域论问题和群论问题建立了一一对应。

\subsection{32.1 域扩张与 Galois 群}\label{ux57dfux6269ux5f20ux4e0e-galois-ux7fa4}

\textbf{定义 32.1}（域扩张）：设 \(K \subseteq L\) 为域。\(L\) 称为 \(K\) 的\textbf{域扩张}，记作 \(L/K\)。\(L\) 可视为 \(K\) 上的向量空间，其维数 \(\dim_K L\) 称为扩张的\textbf{次数}，记作 \([L : K]\)。

\textbf{定义 32.2}（代数扩张与有限扩张）：若每个 \(\alpha \in L\) 都是 \(K\) 上某个非零多项式的根，则 \(L/K\) 称为\textbf{代数扩张}。若 \([L : K] < \infty\)，则称为\textbf{有限扩张}。有限扩张必然是代数扩张，但反之不然（如 \(\overline{\mathbb{Q}}/\mathbb{Q}\) 是代数扩张但无限次）。

\textbf{命题 32.1（次数塔公式）}：设 \(K \subseteq E \subseteq L\) 为域扩张。则 \([L : K] = [L : E][E : K]\)（其中无穷大按通常约定处理）。

\textbf{证明}：设 \(\{\alpha_i\}_{i \in I}\) 是 \(L\) 在 \(E\) 上的基，\(\{\beta_j\}_{j \in J}\) 是 \(E\) 在 \(K\) 上的基。则 \(\{\alpha_i \beta_j\}_{(i,j) \in I \times J}\) 是 \(L\) 在 \(K\) 上的基。线性无关性：若 \(\sum_{i,j} c_{ij} \alpha_i \beta_j = 0\)（\(c_{ij} \in K\)），则 \(\sum_i (\sum_j c_{ij} \beta_j) \alpha_i = 0\)，括号内元素在 \(E\) 中，由 \(\{\alpha_i\}\) 的 \(E\)-线性无关性知 \(\sum_j c_{ij} \beta_j = 0\)，再由 \(\{\beta_j\}\) 的 \(K\)-线性无关性得 \(c_{ij} = 0\)。生成性：任意 \(x \in L\) 可写为 \(x = \sum_i e_i \alpha_i\)（\(e_i \in E\)），每个 \(e_i = \sum_j c_{ij} \beta_j\)（\(c_{ij} \in K\)），故 \(x = \sum_{i,j} c_{ij} \alpha_i \beta_j\)。\(\blacksquare\)

\textbf{定义 32.2'（单扩张）}：设 \(\alpha \in L\)。最小的包含 \(K\) 和 \(\alpha\) 的子域记为 \(K(\alpha)\)，称为 \(K\) 的\textbf{单扩张}。若 \(\alpha\) 在 \(K\) 上代数，则 \(K(\alpha) \cong K[x]/(f(x))\)，其中 \(f\) 是 \(\alpha\) 的极小多项式，且 \([K(\alpha) : K] = \deg f\)。此时 \(\{1, \alpha, \ldots, \alpha^{\deg f - 1}\}\) 是 \(K(\alpha)\) 在 \(K\) 上的一组基。

\textbf{定义 32.3}（Galois 群）：设 \(L/K\) 是域扩张。\(L\) 在 \(K\) 上的 \textbf{Galois 群} \(\operatorname{Gal}(L/K)\) 是 \(L\) 的所有保持 \(K\) 中元素不动的自同构构成的群：

\[\operatorname{Gal}(L/K) = \{\sigma \in \operatorname{Aut}(L) : \sigma|_K = \operatorname{id}_K\}\]

\textbf{定理 32.0（Dedekind 无关性引理）}：设 \(G\) 是群，\(F\) 是域。\(G\) 到 \(F^\times\) 的不同群同态 \(\chi_1, \ldots, \chi_n\) 在 \(F\) 上线性无关。即若 \(\sum_{i=1}^{n} a_i \chi_i(g) = 0\) 对所有 \(g \in G\) 成立（\(a_i \in F\)），则 \(a_1 = \cdots = a_n = 0\)。

\textbf{证明}：对 \(n\) 归纳。\(n=1\) 时，\(a_1 \chi_1(g) = 0\) 对所有 \(g\) 成立，取 \(g = e\) 得 \(a_1 = 0\)。假设对 \(n-1\) 成立。若 \(\sum_{i=1}^{n} a_i \chi_i(g) = 0\) 对所有 \(g\)，取 \(h \in G\)，则 \(\sum_{i=1}^{n} a_i \chi_i(hg) = \sum_{i=1}^{n} a_i \chi_i(h) \chi_i(g) = 0\)。乘以 \(\chi_n(h)\) 并与原式相减：\(\sum_{i=1}^{n-1} a_i(\chi_i(h) - \chi_n(h)) \chi_i(g) = 0\)。由归纳假设，\(a_i(\chi_i(h) - \chi_n(h)) = 0\) 对所有 \(i < n\) 和所有 \(h\)。由于 \(\chi_i \neq \chi_n\)，存在 \(h\) 使 \(\chi_i(h) \neq \chi_n(h)\)，故 \(a_i = 0\)（\(i < n\)）。进而 \(a_n = 0\)。\(\blacksquare\)

\textbf{推论 32.0.1}：若 \(L/K\) 是有限扩张，则 \(|\operatorname{Gal}(L/K)| \leq [L : K]\)。

\textbf{证明}：设 \([L : K] = n\)。若 \(\sigma_1, \ldots, \sigma_m \in \operatorname{Gal}(L/K)\) 互不相同，将 \(\sigma_i\) 视为 \(L^\times\) 到 \(L^\times\) 的群同态（嵌入）。由 Dedekind 引理，它们在 \(L\) 上线性无关。但 \(L\) 在 \(K\) 上的维数为 \(n\)，\(\operatorname{Hom}_K(L, L)\) 的维数至多为 \(n\)（取 \(L\) 在 \(K\) 上的基后，\(K\)-线性映射由基的像唯一确定，构成 \(n\) 维向量空间）。故 \(m \leq n\)，即 \(|\operatorname{Gal}(L/K)| \leq [L : K]\)。\(\blacksquare\)

\textbf{定理 32.0.2（Artin 引理）}：设 \(H\) 是域 \(L\) 的有限自同构群，\(K = L^H\) 为 \(H\) 的不动域。则 \([L : K] = |H|\)，且 \(L/K\) 是 Galois 扩张，\(\operatorname{Gal}(L/K) = H\)。

\textbf{证明}：设 \(|H| = n\)。先证 \([L : K] \geq n\)：若 \([L : K] = m < n\)，取 \(L\) 在 \(K\) 上的基 \(x_1, \ldots, x_m\)。考虑齐次线性方程组 \(\sum_{j=1}^{n} \sigma_j(x_i) y_j = 0\)（\(i = 1, \ldots, m\)），\(m\) 个方程，\(n\) 个未知量，必有非零解 \((y_1, \ldots, y_n) \in L^n\)。则对任意 \(x = \sum c_i x_i \in L\)（\(c_i \in K\)），\(\sum_j \sigma_j(x) y_j = \sum_j \sigma_j(\sum_i c_i x_i) y_j = \sum_i c_i \sum_j \sigma_j(x_i) y_j = 0\)。但 Dedekind 引理（视 \(\sigma_j\) 为 \(L^\times \to L^\times\) 的同态）要求 \(\sigma_j\) 线性无关，这与 \(\sum_j y_j \sigma_j = 0\)（\(y_j\) 不全为零）矛盾。故 \([L : K] \geq n\)。

再证 \([L : K] \leq n\)：对任意 \(\alpha \in L\)，考虑其轨道 \(\{\sigma(\alpha) : \sigma \in H\}\)，设不同元素为 \(\alpha_1, \ldots, \alpha_r\)（\(r \leq n\)）。则 \(f(x) = \prod_{i=1}^{r} (x - \alpha_i) \in L[x]\) 的系数是 \(\alpha_1, \ldots, \alpha_r\) 的初等对称多项式，在 \(H\) 作用下不变，故系数在 \(K = L^H\) 中。因此 \(\alpha\) 在 \(K\) 上的极小多项式次数 \(\leq r \leq n\)。由本原元素定理（见下文），\(L/K\) 可由至多 \(n\) 个元素生成，每个次数 \(\leq n\)，故 \([L : K] \leq n\)。

综合得 \([L : K] = n = |H|\)。\(L/K\) 是 Galois 扩张（正规且可分），且 \(\operatorname{Gal}(L/K) = H\)（因为 \(H \subseteq \operatorname{Gal}(L/K)\) 且 \(|\operatorname{Gal}(L/K)| \leq [L : K] = |H|\)）。\(\blacksquare\)

\textbf{定义 32.4}（Galois 扩张）：域扩张 \(L/K\) 称为 \textbf{Galois 扩张}，如果它是正规的（每个 \(K\) 上的不可约多项式若在 \(L\) 中有一个根，则所有根都在 \(L\) 中）且可分的（每个代数元的极小多项式无重根）。等价地，\(L/K\) 是 Galois 扩张当且仅当 \(L\) 是 \(K\) 上一族可分多项式的分裂域。等价地，\(L/K\) 是 Galois 扩张当且仅当 \(L^{\operatorname{Gal}(L/K)} = K\) 且 \(|\operatorname{Gal}(L/K)| = [L : K]\)。

\textbf{定理 32.1}（Galois 基本定理）：设 \(L/K\) 是有限 Galois 扩张，\(G = \operatorname{Gal}(L/K)\)。则存在一一对应：

\[\{\text{中间域 } K \subseteq E \subseteq L\} \longleftrightarrow \{G \text{ 的子群 } H\}\]

\[E \longmapsto \operatorname{Gal}(L/E) = \{g \in G : g|_E = \operatorname{id}_E\}\]

\[L^H = \{x \in L : h(x) = x,\; \forall h \in H\} \longmapsfrom H\]

满足： 1. \([L : E] = |\operatorname{Gal}(L/E)|\)，\([E : K] = [G : \operatorname{Gal}(L/E)]\) 2. \(E/K\) 是 Galois 扩张当且仅当 \(\operatorname{Gal}(L/E) \trianglelefteq G\)，此时 \(\operatorname{Gal}(E/K) \cong G / \operatorname{Gal}(L/E)\)

\textbf{证明}：设 \(G = \operatorname{Gal}(L/K)\)，\(|G| = [L : K]\)（有限 Galois 扩张的基本性质）。第一步：对 \(H \leq G\)，令 \(E = L^H\)，则 \([L : E] = |H|\)。因为 \(L/L^H\) 是 Galois 扩张且 \(\operatorname{Gal}(L/L^H) = H\)（由 Artin 引理 32.0.2）。第二步：对中间域 \(E\)，\(L/E\) 是 Galois 扩张，\(\operatorname{Gal}(L/E) \leq G\)。第三步：\(L^{\operatorname{Gal}(L/E)} = E\)（由 Galois 扩张的性质，不动域恰好是基域）。第四步：\(\operatorname{Gal}(L/L^H) = H\)（由 Artin 引理）。这建立了两个映射之间的互逆关系。正规子群对应：\(E/K\) 是 Galois 扩张当且仅当 \(E\) 在所有 \(\sigma \in G\) 下稳定，这等价于 \(\sigma \operatorname{Gal}(L/E) \sigma^{-1} = \operatorname{Gal}(L/E)\) 对所有 \(\sigma \in G\)，即 \(\operatorname{Gal}(L/E) \trianglelefteq G\)。此时 \(\sigma \mapsto \sigma|_E\) 给出满同态 \(G \to \operatorname{Gal}(E/K)\)，核为 \(\operatorname{Gal}(L/E)\)，由同构定理得结论。\(\blacksquare\)

Galois 基本定理是代数学中最优美的定理之一。它建立了域扩张的中间域与 Galois 群的子群之间的一一对应，将域论问题转化为群论问题。这一对应的核心是''Galois 连接''（Galois connection）------更一般地，在偏序集之间的一对反变映射，这一概念在逻辑学（Godel 不完备定理的代数化）、形式概念分析和拓扑学（Stone 对偶）中都有深远推广。Galois 对应的几何意义也十分深刻：在代数几何中，不可约代数簇的函数域之间的 Galois 扩张对应于簇之间的平展覆盖（etale covering），而 Galois 群对应于覆盖变换群（deck transformation group）------这正是 Riemann 和 Klein 在 19 世纪就认识到的''代数函数论与 Riemann 面的覆盖理论''的统一。

\subsection{32.2 可分扩张与正规扩张}\label{ux53efux5206ux6269ux5f20ux4e0eux6b63ux89c4ux6269ux5f20}

\textbf{定义 32.5}（可分扩张）：代数扩张 \(L/K\) 是\textbf{可分的}，如果 \(L\) 的每个元素在 \(K\) 上的极小多项式无重根。在特征 \(0\) 时，所有代数扩张都是可分的。在特征 \(p > 0\) 时，可分性等价于 \(L\) 和 \(K^{1/p}\) 在 \(K\) 上线性不交。

\textbf{命题 32.2（可分性的判别）}：多项式 \(f \in K[x]\) 无重根当且仅当 \(\gcd(f, f') = 1\)（\(f'\) 为 \(f\) 的形式导数）。在特征 \(0\) 时，不可约多项式自动满足此条件。在特征 \(p > 0\) 时，不可约多项式 \(f\) 有重根当且仅当 \(f(x) = g(x^p)\) 对某个 \(g \in K[x]\) 成立。

\textbf{证明}：若 \(f\) 在分裂域中有重根 \(\alpha\)，则 \(f(x) = (x-\alpha)^2 h(x)\)，\(f'(x) = 2(x-\alpha)h(x) + (x-\alpha)^2 h'(x)\)，故 \(f'(\alpha) = 0\)，\(\gcd(f, f') \neq 1\)。反之，若 \(\gcd(f, f') = d \neq 1\)，则 \(d\) 的根同时是 \(f\) 和 \(f'\) 的根，\(\alpha\) 是 \(f\) 的重根。在特征 \(p\) 时，若 \(f\) 不可约且有重根，则 \(f' = 0\)（因为 \(\gcd(f, f') = f\) 要求 \(f \mid f'\)，但 \(\deg f' < \deg f\)，故 \(f' = 0\)），这意味着 \(f\) 的每一项次数都是 \(p\) 的倍数，故 \(f(x) = g(x^p)\)。\(\blacksquare\)

\textbf{定理 32.2}（本原元素定理）：有限可分扩张 \(L/K\) 是单扩张，即存在 \(\alpha \in L\) 使得 \(L = K(\alpha)\)。\(\alpha\) 称为\textbf{本原元素}。

\textbf{证明}：若 \(K\) 是有限域，则 \(L^\times\) 是循环群（由有限域的乘法群性质），任何生成元即为本原元素。若 \(K\) 无限，设 \(L = K(\alpha, \beta)\)（可归纳），令 \(f, g\) 分别为 \(\alpha, \beta\) 的极小多项式。在分裂域中，设 \(\alpha = \alpha_1, \ldots, \alpha_m\) 和 \(\beta = \beta_1, \ldots, \beta_n\) 为 \(f\) 和 \(g\) 的根。选 \(c \in K\) 使 \(\alpha + c\beta \neq \alpha_i + c\beta_j\) 对所有 \(i, j\)（\(c\) 只需避开有限个''坏值''）。令 \(\gamma = \alpha + c\beta\)，则 \(\beta\) 是 \(f(\gamma - cX)\) 和 \(g(X)\) 在 \(K(\gamma)[X]\) 中的最大公因子，故 \(\beta \in K(\gamma)\)，从而 \(\alpha = \gamma - c\beta \in K(\gamma)\)。\(\blacksquare\)

\textbf{定义 32.6}（正规扩张）：代数扩张 \(L/K\) 是\textbf{正规的}，如果每个 \(K\) 上的不可约多项式若在 \(L\) 中有一个根，则所有根都在 \(L\) 中。等价地，\(L\) 是 \(K\) 上一族多项式的分裂域。

\textbf{命题 32.3}：\(L/K\) 是 Galois 扩张当且仅当 \(L/K\) 是正规的且可分的。

\textbf{命题 32.4（Galois 扩张的等价刻画）}：对有限扩张 \(L/K\)，以下等价： (i) \(L/K\) 是 Galois 扩张 (ii) \(L\) 是 \(K\) 上某个可分多项式的分裂域 (iii) \(|\operatorname{Gal}(L/K)| = [L : K]\) (iv) \(L^{\operatorname{Gal}(L/K)} = K\)

\textbf{证明}：(i) \(\Rightarrow\) (ii)：\(L/K\) 有限可分，由本原元素定理 \(L = K(\alpha)\)。设 \(f\) 为 \(\alpha\) 的极小多项式，\(f\) 可分且在 \(L\) 中分裂（正规性），故 \(L\) 是 \(f\) 的分裂域。(ii) \(\Rightarrow\) (iii)：设 \(L\) 是 \(K\) 上可分多项式 \(f\) 的分裂域。对 \([L : K]\) 归纳，利用 \(K\)-嵌入的延拓性质可证 \(|\operatorname{Gal}(L/K)| = [L : K]\)。(iii) \(\Rightarrow\) (iv)：设 \(E = L^{\operatorname{Gal}(L/K)}\)。由 Artin 引理，\([L : E] = |\operatorname{Gal}(L/K)| = [L : K]\)，故 \([E : K] = 1\)，即 \(E = K\)。(iv) \(\Rightarrow\) (i)：由 Artin 引理，\(L/K\) 是 Galois 扩张。\(\blacksquare\)

\textbf{定义 32.6'（迹与范数）}：设 \(L/K\) 是有限 Galois 扩张，\(G = \operatorname{Gal}(L/K)\)。对 \(\alpha \in L\)，定义 \[\operatorname{Tr}_{L/K}(\alpha) = \sum_{\sigma \in G} \sigma(\alpha), \quad \operatorname{N}_{L/K}(\alpha) = \prod_{\sigma \in G} \sigma(\alpha)\] 迹是 \(K\)-线性映射 \(L \to K\)，范数是乘法同态 \(L^\times \to K^\times\)。在代数数论中，迹和范数是研究整数环和分歧理论的基本工具。

\textbf{定理 32.2.5（Hilbert 定理 90，乘法形式）}：设 \(L/K\) 是有限 Galois 扩张，\(G = \operatorname{Gal}(L/K) = \langle \sigma \rangle\) 为循环群。若 \(\alpha \in L^\times\) 满足 \(\operatorname{N}_{L/K}(\alpha) = 1\)，则存在 \(\beta \in L^\times\) 使得 \(\alpha = \beta / \sigma(\beta)\)。

\textbf{证明}：由 Dedekind 无关性引理，\(\{1, \sigma, \sigma^2, \ldots, \sigma^{n-1}\}\) 作为 \(L\) 到 \(L\) 的映射是 \(L\)-线性无关的。故存在 \(\gamma \in L\) 使得 \(\beta = \gamma + \alpha \sigma(\gamma) + \alpha \sigma(\alpha) \sigma^2(\gamma) + \cdots + \alpha \sigma(\alpha) \cdots \sigma^{n-2}(\alpha) \sigma^{n-1}(\gamma) \neq 0\)。计算 \(\alpha \sigma(\beta) = \alpha \sigma(\gamma) + \alpha \sigma(\alpha) \sigma^2(\gamma) + \cdots + \alpha \sigma(\alpha) \cdots \sigma^{n-1}(\alpha) \sigma^n(\gamma) = \beta\)（因为 \(\sigma^n = \text{id}\) 且 \(\alpha \sigma(\alpha) \cdots \sigma^{n-1}(\alpha) = \operatorname{N}_{L/K}(\alpha) = 1\)）。故 \(\alpha = \beta / \sigma(\beta)\)。\(\blacksquare\)

\textbf{定理 32.2.6（Hilbert 定理 90，加法形式）}：设 \(L/K\) 是有限 Galois 扩张，\(G = \langle \sigma \rangle\) 为循环群。若 \(\alpha \in L\) 满足 \(\operatorname{Tr}_{L/K}(\alpha) = 0\)，则存在 \(\beta \in L\) 使得 \(\alpha = \beta - \sigma(\beta)\)。

\textbf{证明}：由迹的满射性（可分扩张下迹非退化），可取 \(\gamma \in L\) 使 \(\operatorname{Tr}_{L/K}(\gamma) = 1\)。令 \(\beta = \alpha \sigma(\gamma) + (\alpha + \sigma(\alpha)) \sigma^2(\gamma) + \cdots + (\alpha + \sigma(\alpha) + \cdots + \sigma^{n-2}(\alpha)) \sigma^{n-1}(\gamma)\)。直接计算 \(\beta - \sigma(\beta) = \alpha\)。\(\blacksquare\)

Galois 扩张的经典例子具有丰富的结构。二次扩张 \(\mathbb{Q}(\sqrt{d})/\mathbb{Q}\)（\(d\) 为非平方整数）是 Galois 扩张，Galois 群为 \(\mathbb{Z}/2\mathbb{Z}\)。分圆扩张 \(\mathbb{Q}(\zeta_n)/\mathbb{Q}\)（\(\zeta_n = e^{2\pi i/n}\)）是 Galois 扩张，Galois 群为 \((\mathbb{Z}/n\mathbb{Z})^\times\)，群结构由 \(a \mapsto (\zeta_n \mapsto \zeta_n^a)\) 给出。这一对应是类域论和高斯互反律的出发点。对有限域 \(\mathbb{F}_q\) 的 \(n\) 次扩张 \(\mathbb{F}_{q^n}/\mathbb{F}_q\)，Galois 群为 \(\mathbb{Z}/n\mathbb{Z}\)，由 Frobenius 自同构 \(x \mapsto x^q\) 生成。这一简单的 Galois 群结构是 Weil 猜想中 Frobenius 自同态作用的离散类比，也是 etale 上同调理论中 Frobenius 算子的代数原型。

\subsection{32.3 方程的根式可解性}\label{ux65b9ux7a0bux7684ux6839ux5f0fux53efux89e3ux6027}

\textbf{定义 32.7}（根式扩张）：扩张 \(L/K\) 称为\textbf{根式扩张}，如果存在中间域链 \(K = K_0 \subseteq K_1 \subseteq \cdots \subseteq K_n = L\)，使得 \(K_{i+1} = K_i(\sqrt[n_i]{a_i})\)（\(a_i \in K_i\)）。多项式 \(f \in K[x]\) 称为\textbf{根式可解的}，如果存在 \(f\) 的分裂域包含于某个根式扩张中。

\textbf{定义 32.7'（可解群）}：有限群 \(G\) 称为\textbf{可解群}，如果存在次正规列 \(G = G_0 \triangleright G_1 \triangleright \cdots \triangleright G_n = \{e\}\)，使得每个商群 \(G_i / G_{i+1}\) 是 Abel 群。等价地，\(G\) 的导出列 \(G \supseteq G' \supseteq G'' \supseteq \cdots\) 在有限步内达到 \(\{e\}\)。素数阶循环群是可解的，对称群 \(S_n\) 对 \(n \leq 4\) 可解，\(S_n\) 对 \(n \geq 5\) 不可解（因为 \(A_n\)（\(n \geq 5\)）是单群）。

\textbf{定理 32.2.7（Kummer 理论）}：设 \(K\) 包含 \(n\) 次本原单位根，且 \(\operatorname{char} K \nmid n\)。则 \(K\) 的指数为 \(n\) 的 Abel 扩张（即 Galois 群为 \(n\) 阶循环群的扩张）与 \(K^\times / (K^\times)^n\) 的 \(n\) 阶循环子群一一对应。具体地，若 \(L = K(\sqrt[n]{a})\)（\(a \in K^\times\)），则 \(\operatorname{Gal}(L/K) \cong \mathbb{Z}/n\mathbb{Z}\)。

\textbf{证明}：设 \(L/K\) 是 \(n\) 次循环 Galois 扩张，\(\operatorname{Gal}(L/K) = \langle \sigma \rangle\)。设 \(\zeta \in K\) 为 \(n\) 次本原单位根。\(\zeta \in K\) 且 \(\operatorname{N}_{L/K}(\zeta) = \zeta^n = 1\)。由 Hilbert 定理 90，存在 \(\alpha \in L^\times\) 使得 \(\zeta = \alpha / \sigma(\alpha)\)。则 \(\sigma(\alpha) = \zeta^{-1} \alpha\)，故 \(\sigma(\alpha^n) = (\zeta^{-1} \alpha)^n = \alpha^n\)。因此 \(a = \alpha^n \in K\)（因为它是 \(\sigma\) 的不动点，而 Galois 群由 \(\sigma\) 生成）。又 \([K(\alpha) : K] \geq n\)（因为 \(\sigma^i(\alpha) = \zeta^{-i} \alpha\) 互不相同），故 \(L = K(\alpha) = K(\sqrt[n]{a})\)。\(\blacksquare\)

\textbf{定理 32.3}（Galois 判据）：\(f \in K[x]\) 在特征 \(0\) 的域 \(K\) 上根式可解当且仅当 \(f\) 的 Galois 群（即 \(f\) 的分裂域 \(L\) 的 \(\operatorname{Gal}(L/K)\)）是可解群。

\textbf{证明}：（\(\Rightarrow\)）设 \(f\) 的分裂域 \(L\) 包含于根式扩张 \(E/K\)。取 \(E\) 为包含 \(L\) 的最小 Galois 扩张。则 \(\operatorname{Gal}(E/K)\) 有一列正规子群对应根式扩张链，每个商群 \(\operatorname{Gal}(E_{i+1}/E_i)\) 是循环群（因添加 \(p\) 次根对应于 \(p\) 次循环扩张，若 \(K\) 包含足够多单位根）。取商群 \(\operatorname{Gal}(E/K) / \operatorname{Gal}(E/L) \cong \operatorname{Gal}(L/K)\) 也可解。（\(\Leftarrow\)）若 \(\operatorname{Gal}(L/K)\) 可解，则存在正规子群列使商群为素数阶循环群。通过在基域中添加适当单位根，每个素数阶循环扩张可由根式扩张实现（Kummer 理论）。具体地，设 \(G = G_0 \triangleright G_1 \triangleright \cdots \triangleright G_r = \{e\}\) 为合成列，\(G_i/G_{i+1}\) 为素数阶循环群。对应中间域 \(K = E_0 \subseteq E_1 \subseteq \cdots \subseteq E_r = L\)，\(E_{i+1}/E_i\) 为素数阶循环扩张。在 \(E_i\) 中添加适当的单位根后，\(E_{i+1}/E_i\) 成为 Kummer 扩张，从而可由根式扩张实现。\(\blacksquare\)

\textbf{定理 32.4}（五次方程不可解性，Abel-Ruffini 定理）：一般五次方程 \(x^5 + a_4 x^4 + \cdots + a_0 = 0\) 在有理函数域 \(\mathbb{C}(a_0, \ldots, a_4)\) 上不是根式可解的。其 Galois 群为对称群 \(S_5\)，而 \(S_5\) 不是可解群（因为 \(A_5\) 是非交换单群）。

\textbf{证明}：一般五次方程在 \(\mathbb{C}(a_0, \ldots, a_4)\) 上的 Galois 群是 \(S_5\)（由对称多项式基本定理和 van der Waerden 的不可约性判别法）。\(S_5\) 的合成列为 \(S_5 \triangleright A_5 \triangleright \{e\}\)，其中 \(A_5\) 是 \(60\) 阶非交换单群，故 \(S_5\) 不可解。由 Galois 判据，一般五次方程不能根式求解。\(\blacksquare\)

\textbf{定理 32.4.1（三次方程的 Galois 群）}：设 \(f(x) = x^3 + px + q \in \mathbb{Q}[x]\) 不可约，判别式 \(\Delta = -4p^3 - 27q^2\)。则 \(f\) 的分裂域的 Galois 群为： - \(S_3\)（阶 \(6\)），若 \(\Delta\) 不是 \(\mathbb{Q}\) 中的平方数 - \(A_3 \cong \mathbb{Z}/3\mathbb{Z}\)（阶 \(3\)），若 \(\Delta\) 是 \(\mathbb{Q}\) 中的平方数

\textbf{证明}：\(\Delta = \prod_{i < j} (\alpha_i - \alpha_j)^2\)，其中 \(\alpha_i\) 为 \(f\) 的根。\(\sqrt{\Delta} = \prod_{i < j} (\alpha_i - \alpha_j)\) 在 \(A_3\) 下不变，在奇置换下变号。故 \(\sqrt{\Delta} \in \mathbb{Q}\) 当且仅当 Galois 群 \(\subseteq A_3\)。由于 \(f\) 不可约，Galois 群是 \(S_3\) 的可递子群，只能是 \(S_3\) 或 \(A_3\)。\(\blacksquare\)

五次方程不可解性的证明是 Galois 理论最辉煌的胜利。它彻底终结了自 16 世纪意大利数学家发现三次和四次方程求根公式以来，数学家们对五次及以上方程求根公式长达三个世纪的探索。Galois 的洞察力在于：他将''根式可解''这一数论/代数问题转化为 Galois 群的群论性质（可解性），而这一转化完全通过与具体的方程脱钩------将方程视为域扩张，Galois 群视为对称性的抽象体现。这一范式转换深远影响了 20 世纪数学：代数拓扑中的覆盖空间理论（基本群与 Galois 群的类比）、代数数论中的类域论（Galois 群的 Abel 化）、乃至 Grothendieck 的远 Abel 几何（从 Galois 群恢复数域或代数曲线的几何信息），都是 Galois 基本思想的回响。

\subsection{32.4 有限域的 Galois 理论}\label{ux6709ux9650ux57dfux7684-galois-ux7406ux8bba}

\textbf{定理 32.4.2（有限域的分类）}：设 \(F\) 是有限域。则 \(|F| = p^n\)（\(p\) 为素数，\(n \geq 1\)），且对每个 \(p^n\) 存在唯一的 \(p^n\) 元有限域（在同构意义下），记为 \(\mathbb{F}_{p^n}\)。\(\mathbb{F}_{p^n}\) 是 \(x^{p^n} - x\) 在 \(\mathbb{F}_p\) 上的分裂域。

\textbf{定理 32.4.3（有限域的 Galois 理论）}：\(\mathbb{F}_{p^n}/\mathbb{F}_p\) 是 \(n\) 次 Galois 扩张，Galois 群为 \(\mathbb{Z}/n\mathbb{Z}\)，由 Frobenius 自同构 \(\operatorname{Fr}_p: x \mapsto x^p\) 生成。\(\mathbb{F}_{p^n}\) 的所有子域为 \(\mathbb{F}_{p^d}\)（\(d \mid n\)），对应 Galois 群的子群 \(\langle \operatorname{Fr}_p^d \rangle \cong \mathbb{Z}/(n/d)\mathbb{Z}\)。

\textbf{证明}：\(\mathbb{F}_{p^n}\) 是 \(x^{p^n} - x\) 的分裂域，该多项式可分（导数 \(-1 \neq 0\)），故 \(\mathbb{F}_{p^n}/\mathbb{F}_p\) 是 Galois 扩张。\(\operatorname{Fr}_p\) 是 \(\mathbb{F}_{p^n}\) 的自同构，且 \(\operatorname{Fr}_p^n = \text{id}\)（因为 \(x^{p^n} = x\) 对所有 \(x \in \mathbb{F}_{p^n}\)）。若 \(\operatorname{Fr}_p^k = \text{id}\)（\(k < n\)），则 \(x^{p^k} = x\) 对所有 \(x \in \mathbb{F}_{p^n}\)，但 \(x^{p^k} - x\) 至多有 \(p^k\) 个根，矛盾。故 \(\operatorname{Fr}_p\) 的阶为 \(n\)，\(\operatorname{Gal}(\mathbb{F}_{p^n}/\mathbb{F}_p) = \langle \operatorname{Fr}_p \rangle \cong \mathbb{Z}/n\mathbb{Z}\)。子域结构与 \(\mathbb{Z}/n\mathbb{Z}\) 的子群对应。\(\blacksquare\)

\subsection{32.5 分圆扩张}\label{ux5206ux5706ux6269ux5f20}

\textbf{定义 32.8}（分圆扩张）：设 \(\zeta_n = e^{2\pi i/n}\) 为 \(n\) 次本原单位根。\(\mathbb{Q}(\zeta_n)/\mathbb{Q}\) 称为 \(n\) 次\textbf{分圆扩张}。

\textbf{定理 32.5.1（分圆扩张的 Galois 群）}：\(\mathbb{Q}(\zeta_n)/\mathbb{Q}\) 是 Galois 扩张，其次数为 \(\varphi(n)\)（Euler 函数），Galois 群为： \[\operatorname{Gal}(\mathbb{Q}(\zeta_n)/\mathbb{Q}) \cong (\mathbb{Z}/n\mathbb{Z})^\times\] 同构由 \(a \mapsto (\sigma_a: \zeta_n \mapsto \zeta_n^a)\) 给出（\(a \in (\mathbb{Z}/n\mathbb{Z})^\times\)）。

\textbf{证明}：\(\zeta_n\) 的极小多项式是 \(n\) 次分圆多项式 \(\Phi_n(x) = \prod_{1 \leq k \leq n, \gcd(k,n)=1} (x - \zeta_n^k)\)，次数为 \(\varphi(n)\)。\(\Phi_n(x) \in \mathbb{Q}[x]\)（由 Gauss 引理），且 \(\mathbb{Q}(\zeta_n)\) 是 \(\Phi_n\) 的分裂域（包含 \(\zeta_n\) 的所有共轭），故 \(\mathbb{Q}(\zeta_n)/\mathbb{Q}\) 是 Galois 扩张。每个 \(\sigma \in \operatorname{Gal}(\mathbb{Q}(\zeta_n)/\mathbb{Q})\) 由 \(\sigma(\zeta_n)\) 决定，而 \(\sigma(\zeta_n)\) 必须是 \(\zeta_n\) 的共轭，即 \(\zeta_n^a\)（\(\gcd(a, n) = 1\)）。映射 \(\sigma \mapsto a \bmod n\) 是单同态，且两边阶均为 \(\varphi(n)\)，故为同构。\(\blacksquare\)

\textbf{推论 32.5.2（二次互反律与分圆域）}：二次域 \(\mathbb{Q}(\sqrt{p^*})\)（\(p^* = (-1)^{(p-1)/2} p\)）是 \(\mathbb{Q}(\zeta_p)\) 的唯一二次子域。这建立了二次互反律与分圆域的 Galois 群之间的深刻联系，是类域论中 Kronecker-Weber 定理（\(\mathbb{Q}\) 的任何 Abel 扩张都包含于某个分圆扩张中）的特殊情形。

\subsection{32.6 Galois 逆问题与无穷 Galois 理论}\label{galois-ux9006ux95eeux9898ux4e0eux65e0ux7a77-galois-ux7406ux8bba}

\textbf{定义 32.9}（Galois 逆问题）：给定有限群 \(G\) 和域 \(K\)，是否存在 \(K\) 的 Galois 扩张 \(L/K\) 使得 \(\operatorname{Gal}(L/K) \cong G\)？对 \(K = \mathbb{Q}\)，Galois 逆问题（即每个有限群是否都是 \(\mathbb{Q}\) 上某个 Galois 扩张的 Galois 群）是未解决的公开问题，但对许多群类（如所有可解群、所有对称群 \(S_n\) 和交错群 \(A_n\)、大多数散在单群）已得到肯定解答。Shafarevich 在 1954 年证明了所有可解群都可实现为 \(\mathbb{Q}\) 上的 Galois 群，其证明使用了嵌入问题和数域的局部-整体原理。对有限域 \(\mathbb{F}_q\)，Galois 逆问题是平凡的：\(\mathbb{F}_{q^n}/\mathbb{F}_q\) 的 Galois 群是 \(\mathbb{Z}/n\mathbb{Z}\)，而 \(\mathbb{F}_q\) 上的任意有限 Galois 扩张都是这些循环扩张的合成。

\textbf{定理 32.6.1（正规基定理）}：设 \(L/K\) 是有限 Galois 扩张，\(G = \operatorname{Gal}(L/K)\)。则存在 \(\alpha \in L\) 使得 \(\{\sigma(\alpha) : \sigma \in G\}\) 是 \(L\) 在 \(K\) 上的一组基。这样的 \(\alpha\) 称为\textbf{正规基元素}。

\textbf{证明}（框架）：\(L\) 是 \(K[G]\)-模（\(G\) 通过 Galois 作用）。正规基定理等价于 \(L \cong K[G]\) 作为 \(K[G]\)-模。由 Krull-Schmidt 定理和 \(K[G]\) 的半单性（特征不整除 \(|G|\) 时），只需证 \(L\) 和 \(K[G]\) 有相同的特征标。\(\blacksquare\)

\textbf{定义 32.10}（无穷 Galois 理论 / Krull 拓扑）：设 \(L/K\) 是无穷 Galois 扩张（如 \(\overline{\mathbb{Q}}/\mathbb{Q}\)）。则 \(\operatorname{Gal}(L/K)\) 配以 \textbf{Krull 拓扑}（以子群 \(\operatorname{Gal}(L/E)\)（\(E/K\) 为有限 Galois 子扩张）为邻域基）成为 \textbf{射有限群}（profinite group）。无穷 Galois 基本定理在 Krull 拓扑下成立：中间域与闭子群一一对应。\(\operatorname{Gal}(\overline{\mathbb{Q}}/\mathbb{Q})\) 是绝对 Galois 群，其结构是数论中最核心的未解问题之一------理解其表示论（Langlands 纲领）和其外 Galois 表示（\(\operatorname{Gal}(\overline{\mathbb{Q}}/\mathbb{Q})\) 在代数基本群上的作用）是当代数论的前沿。

\textbf{证明}：无穷 Galois 扩张 \(L/K\) 是所有有限 Galois 子扩张 \(L_i/K\) 的正向极限。\(\operatorname{Gal}(L/K) = \varprojlim \operatorname{Gal}(L_i/K)\)，其中投射极限在有限群的离散拓扑下取。Krull 拓扑使 \(\operatorname{Gal}(L/K)\) 成为紧 Hausdorff 全不连通拓扑群。有限 Galois 子扩张的 Galois 群形成投射系统，极限的拓扑在这个系统中就是自然的 Krull 拓扑。基本定理：对闭子群 \(H\)，\(L^H\) 是中间域；对中间域 \(E\)，\(\operatorname{Gal}(L/E)\) 是闭子群。这些对应是一一对应且满足 Galois 连接的所有性质。\(\blacksquare\)

无穷 Galois 理论是有限 Galois 理论的自然推广，其核心区别在于 Galois 群不再是有限群，而是射有限群------这一拓扑结构反映了无穷代数扩张的极限性质。射有限群的结构定理（如每个射有限群都是有限群的投射极限）使得无穷 Galois 理论中的许多问题可以约化为有限情形。\(\operatorname{Gal}(\overline{\mathbb{Q}}/\mathbb{Q})\) 的绝对 Galois 群是数学中最神秘的对象之一：它的上同调群控制着数域上的 Brauer 群、Milnor \(K\)-理论和 Galois 上同调中的许多深层次结构；它的表示论（特别是 \(\ell\)-进表示）是 Langlands 纲领的核心对象；它的外自同构群（通过 Belyi 定理）与 \(\overline{\mathbb{Q}}\) 上的代数曲线的 Grothendieck-Teichmuller 群有紧密联系。绝对 Galois 群的结构至今仍是开放问题，而 Neukirch-内田定理（数域的 Galois 群完全决定了该数域的同构类）和望月新一的远 Abel 几何（从数域的绝对 Galois 群恢复该数域）则进一步揭示了 Galois 群在数论中的中心地位。

\subsection{32.7 Galois 上同调初步}\label{galois-ux4e0aux540cux8c03ux521dux6b65}

\textbf{定义 32.11}（Galois 上同调）：设 \(G\) 是射有限群，\(M\) 是离散 \(G\)-模（即 \(G\) 在 \(M\) 上的作用连续）。\(G\) 的系数在 \(M\) 中的上同调群 \(H^n(G, M)\) 定义为 \(M\) 的不动点函子的右导出函子。特别地，\(H^0(G, M) = M^G\)，\(H^1(G, M)\) 分类 \(M\) 的交叉同态（crossed homomorphisms）模去主交叉同态。

\textbf{命题 32.7.1（Hilbert 定理 90 的上同调表述）}：\(H^1(\operatorname{Gal}(L/K), L^\times) = 0\) 和 \(H^1(\operatorname{Gal}(L/K), L) = 0\)。

Galois 上同调是类域论和 Langlands 纲领的基本技术工具。Tate 上同调、局部类域论中的 Brauer 群的计算、以及 étale 上同调都是 Galois 上同调的自然推广。在算术几何中，Galois 上同调被用于研究椭圆曲线的 Selmer 群和 Birch-Swinnerton-Dyer 猜想。

\hrulefill

\hrulefill

\hrulefill

\hrulefill

\hrulefill

\newpage

\chapter{05-代数/03-抽象代数二/02-交换代数.md}\label{ux4ee3ux657003-ux62bdux8c61ux4ee3ux6570ux4e8c02-ux4ea4ux6362ux4ee3ux6570.md}

\section{第126章：交换代数基础}\label{ux7b2c126ux7ae0ux4ea4ux6362ux4ee3ux6570ux57faux7840}

交换代数是代数几何和代数数论的语言基础，研究交换环的结构。

\subsection{68.1 Noether 环与 Artin 环}\label{noether-ux73afux4e0e-artin-ux73af}

\textbf{定义 68.1}（Noether 环）：交换环 \(R\) 称为 \textbf{Noether 环}，如果 \(R\) 中的理想满足\textbf{升链条件}（ACC）：任意理想升链

\[I_1 \subseteq I_2 \subseteq I_3 \subseteq \cdots\]

必稳定，即存在 \(n\) 使得 \(I_n = I_{n+1} = \cdots\)。

\textbf{命题 68.1}（Noether 环的等价刻画）：\(R\) 是 Noether 环当且仅当 \(R\) 的每个理想都是有限生成的，也当且仅当 \(R\) 的每个非空理想族在包含关系下都有极大元。

\emph{证明}：(ACC \(\Rightarrow\) 有限生成)：若 \(I\) 不能有限生成，则归纳构造严格升链，与 ACC 矛盾。(有限生成 \(\Rightarrow\) ACC)：若 \(I_1 \subseteq I_2 \subseteq \cdots\)，则 \(I = \bigcup I_n\) 也是理想，有限生成 \(I = (a_1, \ldots, a_m)\)。每个 \(a_i\) 属于某个 \(I_{n_i}\)，取 \(n = \max\{n_i\}\)，则 \(I = I_n\)。∎

\textbf{定理 68.2}（Hilbert 基定理）：若 \(R\) 是 Noether 环，则多项式环 \(R[x]\) 也是 Noether 环。

\emph{证明}：设 \(I \subseteq R[x]\) 是理想。对每个 \(n \geq 0\)，令 \(J_n\) 为 \(I\) 中 \(n\) 次多项式的首项系数和 0 构成的 \(R\) 的理想。则 \(J_0 \subseteq J_1 \subseteq J_2 \subseteq \cdots\)。由 \(R\) 的 Noether 性，存在 \(N\) 使 \(J_n = J_N\) 对所有 \(n \geq N\)。

每个 \(J_k\)（\(k \leq N\)）有限生成：\(J_k = (a_{k1}, \ldots, a_{ks_k})\)。取 \(f_{kj} \in I\) 使其首项系数为 \(a_{kj}\)。则这些 \(f_{kj}\) 生成 \(I\)：对任意 \(f \in I\)，用这些生成元逐步消去 \(f\) 的首项，归纳可得 \(f\) 属于它们生成的理想。∎

\textbf{推论 68.3}：若 \(R\) 是 Noether 环，则 \(R[x_1, \ldots, x_n]\) 也是 Noether 环。特别地，\(\mathbb{Z}[x_1, \ldots, x_n]\) 和 \(k[x_1, \ldots, x_n]\)（\(k\) 为域）是 Noether 环。

\textbf{定义 68.2}（Artin 环）：交换环 \(R\) 称为 \textbf{Artin 环}，如果 \(R\) 中的理想满足\textbf{降链条件}（DCC）：任意理想降链 \(I_1 \supseteq I_2 \supseteq \cdots\) 必稳定。

\textbf{定理 68.4}（Akizuki-Hopkins 定理）：\(R\) 是 Artin 环当且仅当 \(R\) 是 Noether 环且 \(\dim R = 0\)（即每个素理想都是极大理想）。

\textbf{证明}：(\(\Rightarrow\)) Artin 环满足 DCC。DCC 可推出所有素理想极大------若 \(\mathfrak{p}\) 素，则 \(R/\mathfrak{p}\) 是 Artin 整环，因而是域，故 \(\mathfrak{p}\) 极大，即 \(\dim R = 0\)。用 DCC 和素理想的极小分解可证 \(R\) 仅有有限个极大理想，且 Jacobson 根 \(\mathfrak{r}\) 幂零。由于 \(R/\mathfrak{r}^n\) 的滤过有限性可推 Noether 性。(\(\Leftarrow\)) Noether 零维环仅有有限个极小素理想（即极大理想），每个极大理想 \(\mathfrak{m}_i\) 在 Noether 条件下满足 \(\mathfrak{m}_i^n = 0\)（对某 \(n\)）。有限长度模的降链必稳定，故 DCC 成立。\(\blacksquare\)

\subsection{68.2 素理想与局部化}\label{ux7d20ux7406ux60f3ux4e0eux5c40ux90e8ux5316}

\textbf{定义 68.3}（素理想与极大理想）：设 \(R\) 为交换环。 - 真理想 \(\mathfrak{p} \subsetneq R\) 称为\textbf{素理想}，如果 \(ab \in \mathfrak{p} \Rightarrow a \in \mathfrak{p}\) 或 \(b \in \mathfrak{p}\)。等价地，\(R/\mathfrak{p}\) 是整环。 - 真理想 \(\mathfrak{m} \subsetneq R\) 称为\textbf{极大理想}，如果不存在真理想 \(\mathfrak{m} \subsetneq I \subsetneq R\)。等价地，\(R/\mathfrak{m}\) 是域。

\textbf{定义 68.4}（谱）：\(R\) 的\textbf{素谱} \(\operatorname{Spec}(R)\) 是所有素理想构成的集合。\(R\) 的\textbf{极大谱} \(\operatorname{MaxSpec}(R)\) 是所有极大理想构成的集合。

\textbf{定理 68.5}（Krull 存在定理）：每个非零交换环都有极大理想（从而也有素理想）。

\emph{证明}：Zorn 引理：所有真理想在包含关系下构成偏序集，任意链的并仍是真理想（因为 \(1\) 不在其中），故有极大元。∎

\textbf{定义 68.5}（局部化）：设 \(R\) 为交换环，\(S \subseteq R\) 为乘法封闭子集（\(1 \in S\)，\(a, b \in S \Rightarrow ab \in S\)）。\(R\) 关于 \(S\) 的\textbf{局部化} \(S^{-1}R\) 定义为等价类集合：

\[S^{-1}R = \left\{\frac{a}{s} : a \in R, s \in S\right\}\]

其中 \(\frac{a}{s} = \frac{a'}{s'}\) 当且仅当存在 \(t \in S\) 使 \(t(as' - a's) = 0\)。

\textbf{重要特例}： - 若 \(S = R \setminus \mathfrak{p}\)（\(\mathfrak{p}\) 为素理想），\(S^{-1}R\) 记作 \(R_{\mathfrak{p}}\)，称为 \(R\) 在 \(\mathfrak{p}\) 处的\textbf{局部化}。 - 若 \(S = \{f^n : n \geq 0\}\)，\(S^{-1}R\) 记作 \(R_f\)。 - 若 \(R\) 是整环，\(S = R \setminus \{0\}\)，则 \(S^{-1}R\) 是 \(R\) 的\textbf{分式域}。

\textbf{命题 68.6}（局部化与素理想）：\(R_{\mathfrak{p}}\) 是局部环，其唯一极大理想为 \(\mathfrak{p}R_{\mathfrak{p}} = \{a/s : a \in \mathfrak{p}, s \notin \mathfrak{p}\}\)。存在自然的单射 \(R/\mathfrak{p} \hookrightarrow R_{\mathfrak{p}}/\mathfrak{p}R_{\mathfrak{p}}\)。

\subsection{68.3 整性、整闭包与 Noether 正规化}\label{ux6574ux6027ux6574ux95edux5305ux4e0e-noether-ux6b63ux89c4ux5316}

\textbf{定义 68.6}（整元素）：设 \(R \subseteq A\) 是环的扩张。\(a \in A\) 称为 \(R\) 上的\textbf{整元素}，如果存在首一多项式 \(f(x) = x^n + r_{n-1}x^{n-1} + \cdots + r_0 \in R[x]\) 使 \(f(a) = 0\)。

\(R\) 在 \(A\) 中的\textbf{整闭包}是 \(A\) 中所有 \(R\) 上整元素构成的集合。若 \(R\) 等于其在自己分式域中的整闭包，则称 \(R\) 为\textbf{整闭的}。

\textbf{定理 68.7}（Noether 正规化引理）：设 \(k\) 为域，\(A\) 是有限生成 \(k\)-代数。则存在代数无关的元素 \(y_1, \ldots, y_d \in A\)（\(d\) 是 \(A\) 的 Krull 维数）使得 \(A\) 是 \(k[y_1, \ldots, y_d]\) 上的有限生成模（即整扩张）。

\textbf{证明}：设 \(A = k[x_1, \ldots, x_n]\)。若 \(x_i\) 代数无关，则 \(d = n\) 且 \(A\) 已是 \(k[x_1, \ldots, x_n]\) 上的有限模（即为自身），结论成立。否则，存在非平凡多项式关系 \(f(x_1, \ldots, x_n) = 0\)。设 \(f\) 的总次数为 \(D\)。选取充分大的整数 \(m_1, \ldots, m_{n-1}\)（例如 \(m_i = D^{n-i}\)），作变量替换 \(y_i = x_i - x_n^{m_i}\)（\(i = 1, \ldots, n-1\)）。则 \(f(x_1, \ldots, x_n) = f(y_1 + x_n^{m_1}, \ldots, y_{n-1} + x_n^{m_{n-1}}, x_n)\)。适当选取 \(m_i\) 使得 \(f\) 展开后关于 \(x_n\) 是首一的（不同单项式 \(x_1^{a_1} \cdots x_n^{a_n}\) 在替换后给出不同的 \(x_n\) 指数，最高次项唯一）。此时 \(A\) 是 \(k[y_1, \ldots, y_{n-1}]\) 上的整扩张，且生成元 \(x_n\) 满足 \(y_1, \ldots, y_{n-1}\) 系数的首一多项式。对 \(k[y_1, \ldots, y_{n-1}]\) 归纳即得结论。\(\blacksquare\)

\textbf{推论 68.8}（Hilbert 零点定理的弱形式）：若 \(k\) 是代数闭域，则 \(k[x_1, \ldots, x_n]\) 的每个极大理想形如 \((x_1 - a_1, \ldots, x_n - a_n)\)，对应 \(k^n\) 中的一个点。

\textbf{定理 68.9}（Hilbert 零点定理，强形式）：设 \(k\) 为代数闭域，\(I \subseteq k[x_1, \ldots, x_n]\) 为理想。若 \(f \in k[x_1, \ldots, x_n]\) 在 \(I\) 的零点集 \(V(I)\) 上恒为零，则 \(f^m \in I\) 对某个 \(m \geq 1\) 成立。即 \(I(V(I)) = \sqrt{I}\)（\(I\) 的根）。

\textbf{证明}（Rabinowitsch 技巧）：仅需证 \(\sqrt{I} \supseteq I(V(I))\)，反包含显然。设 \(f \in I(V(I))\)，即 \(f\) 在 \(V(I)\) 上恒为零。在 \(k[x_1, \ldots, x_n, y]\) 中考虑理想 \(J = I \cdot k[x_1, \ldots, x_n, y] + (1 - yf)\)。若 \(J\) 有零点 \((a_1, \ldots, a_n, b) \in k^{n+1}\)，则 \((a_1, \ldots, a_n) \in V(I)\) 且 \(1 - b f(a) = 0\)。但 \(f(a) = 0\)（因 \(f \in I(V(I))\)），故 \(1 = 0\)，矛盾。因此 \(V(J) = \emptyset\)。由 Hilbert 零点定理的弱形式，\(J = (1)\)。故存在 \(g_i \in k[x_1, \ldots, x_n, y]\) 使得 \(1 = \sum g_i h_i + g_0 (1 - yf)\)，其中 \(h_i \in I\)。令 \(y = 1/f\) 并乘以 \(f\) 的适当幂次以消去分母，得 \(f^m = \sum \tilde{g}_i h_i \in I\)。\(\blacksquare\)

\subsection{68.4 准素分解与 Krull 交定理}\label{ux51c6ux7d20ux5206ux89e3ux4e0e-krull-ux4ea4ux5b9aux7406}

\textbf{定义 68.7}（准素理想）：真理想 \(\mathfrak{q} \subsetneq R\) 称为\textbf{准素理想}，如果 \(ab \in \mathfrak{q}\) 且 \(a \notin \mathfrak{q}\) 蕴含 \(b^n \in \mathfrak{q}\) 对某个 \(n \geq 1\) 成立。等价地，\(R/\mathfrak{q}\) 中每个零因子都是幂零元。

\textbf{定理 68.10}（Lasker-Noether 准素分解定理）：在 Noether 环中，每个理想都可以表示为有限个准素理想的交。即对任意理想 \(I\)，存在准素理想 \(\mathfrak{q}_1, \ldots, \mathfrak{q}_n\) 使得

\[I = \mathfrak{q}_1 \cap \mathfrak{q}_2 \cap \cdots \cap \mathfrak{q}_n\]

且此分解可以取为「不可缩短的」且素理想 \(\sqrt{\mathfrak{q}_i}\) 互不相同（称为「极小准素分解」）。

\textbf{注意}：此定理是整数环中「算术基本定理」在一般 Noether 环中的推广。整数环中 \(n = \prod p_i^{e_i}\) 对应理想 \((n) = \bigcap (p_i^{e_i})\)。

\textbf{证明}：(1) 在 Noether 环中，每个理想可表示为有限个\textbf{不可约理想}（不能写成两个严格更大理想的交）的交。这是因为若不然，某理想 \(I_0\) 不是不可约理想的交，则 \(I_0 = J_1 \cap K_1\) 且 \(J_1, K_1 \supsetneq I_0\)。若 \(J_1\) 也不是不可约理想的交，继续分解，得到严格升链 \(I_0 \subsetneq J_1 \subsetneq J_2 \subsetneq \cdots\)，与 Noether 性（升链条件）矛盾。

\begin{enumerate}
\def\labelenumi{(\arabic{enumi})}
\setcounter{enumi}{1}
\item
  不可约理想必为准素理想：设 \(\mathfrak{q}\) 不可约，\(ab \in \mathfrak{q}\) 且 \(a \notin \mathfrak{q}\)。考虑升链 \((\mathfrak{q}:b) \subseteq (\mathfrak{q}:b^2) \subseteq \cdots\)。由 Noether 性，存在 \(n\) 使 \((\mathfrak{q}:b^n) = (\mathfrak{q}:b^{n+1})\)。断言 \(\mathfrak{q} = (\mathfrak{q} + (b^n)) \cap (\mathfrak{q} + (a))\)。右端包含左端显然。若 \(x\) 属于右端，\(x = q_1 + r b^n = q_2 + s a\)，则 \(x b \in \mathfrak{q} + (a b) \subseteq \mathfrak{q}\)，故 \(r b^{n+1} \in \mathfrak{q}\)，即 \(r \in (\mathfrak{q}:b^{n+1}) = (\mathfrak{q}:b^n)\)，故 \(r b^n \in \mathfrak{q}\)，\(x \in \mathfrak{q}\)。由 \(\mathfrak{q}\) 的不可约性，\(\mathfrak{q} + (b^n) = \mathfrak{q}\) 或 \(\mathfrak{q} + (a) = \mathfrak{q}\)。后者不可能（因 \(a \notin \mathfrak{q}\)），故 \(b^n \in \mathfrak{q}\)。
\item
  唯一性：极小准素分解中，素理想 \(\mathfrak{p}_i = \sqrt{\mathfrak{q}_i}\) 由 \(I\) 唯一决定------它们是使得 \(I \subseteq \mathfrak{p}\) 且 \((I:x) \subseteq \mathfrak{p}\) 的素理想中的极小元。\(\blacksquare\)
\end{enumerate}

\textbf{定理 68.11}（Krull 交定理）：设 \(R\) 是 Noether 环，\(I \subseteq R\) 是真理想。则 \(\bigcap_{n=1}^\infty I^n = 0\) 当 \(R\) 是整环或局部环时成立。一般地，存在 \(a \in I\) 使得 \((1 - a)\bigcap_{n=1}^\infty I^n = 0\)。

\textbf{证明}：设 \(J = \bigcap_{n=1}^\infty I^n\)。由 Artin-Rees 引理，存在 \(c\) 使 \(J \cap I^c \subseteq IJ\)。但 \(J \subseteq I^c\)，故 \(J \subseteq IJ\)。反向包含 \(IJ \subseteq J\) 显然，故 \(J = IJ\)。由于 \(I \subseteq \operatorname{rad}(R)\)（当 \(R\) 局部时）或由 Nakayama 引理：若 \(R\) 是局部环，\(J\) 是有限生成模，\(J = IJ\) 则 \(J = 0\)。一般地，对有限生成 \(R\)-模 \(J\)，\(J = IJ\) 推出存在 \(a \in I\) 使 \((1 - a)J = 0\)。\(\blacksquare\)

\hrulefill

\hrulefill

\hrulefill

\hrulefill

\hrulefill

\hrulefill

\section{第127章：表示论基础}\label{ux7b2c127ux7ae0ux8868ux793aux8bbaux57faux7840}

表示论研究群（或代数）通过线性变换的具体实现，是连接抽象代数与线性代数的桥梁。

\subsection{69.1 群的表示}\label{ux7fa4ux7684ux8868ux793a}

\textbf{定义 69.1}（群的线性表示）：设 \(G\) 是有限群，\(V\) 是域 \(k\) 上的有限维向量空间。\(G\) 在 \(V\) 上的一个\textbf{线性表示}是一个群同态 \(\rho: G \to \operatorname{GL}(V)\)。

\(V\) 的维数称为表示的\textbf{维数}（或\textbf{次数}）。若 \(\rho\) 是单射，则称表示为\textbf{忠实的}。

\textbf{定义 69.2}（不变子空间与不可约表示）：表示 \(\rho: G \to \operatorname{GL}(V)\) 的子空间 \(W \subseteq V\) 称为\textbf{不变的}（或 \(G\)-子模），如果 \(\rho(g)(W) \subseteq W\) 对所有 \(g \in G\) 成立。表示 \((\rho, V)\) 称为\textbf{不可约的}，如果 \(V \neq 0\) 且唯一的不变子空间是 \(\{0\}\) 和 \(V\)。

\textbf{定义 69.3}（表示的直和）：若 \(V = V_1 \oplus V_2\)（\(\dim V_i > 0\)）且 \(V_1, V_2\) 都是 \(G\)-不变的，则称 \((\rho, V)\) 是 \((\rho|_{V_1}, V_1)\) 和 \((\rho|_{V_2}, V_2)\) 的\textbf{直和}。若表示不能分解为非平凡子表示的直和，则称其为\textbf{不可分解的}。

\textbf{定理 69.1}（Maschke 定理）：设 \(G\) 是有限群，\(k\) 是域且 \(\operatorname{char}(k) \nmid |G|\)。则 \(G\) 在 \(k\) 上的每个有限维表示都是完全可约的（即可分解为不可约表示的直和）。

\emph{证明}：设 \(W \subseteq V\) 是 \(G\)-不变子空间。需构造 \(G\)-不变的补空间 \(W'\)。取任意线性投影 \(\pi: V \to W\)（非 \(G\)-等变）。定义平均投影

\[\tilde{\pi}(v) = \frac{1}{|G|} \sum_{g \in G} \rho(g) \pi(\rho(g^{-1})v)\]

则 \(\tilde{\pi}\) 是 \(G\)-等变的（\(\tilde{\pi}(\rho(h)v) = \rho(h)\tilde{\pi}(v)\)），且 \(\tilde{\pi}|_W = \operatorname{id}_W\)。于是 \(W' = \ker \tilde{\pi}\) 是 \(G\)-不变的补空间。∎

\subsection{69.2 特征标理论}\label{ux7279ux5f81ux6807ux7406ux8bba}

以下设 \(k = \mathbb{C}\)（复数域），\(G\) 是有限群。

\textbf{定义 69.4}（特征标）：表示 \(\rho: G \to \operatorname{GL}(V)\) 的\textbf{特征标}是函数 \(\chi_V: G \to \mathbb{C}\)，\(\chi_V(g) = \operatorname{tr}(\rho(g))\)。

特征标是\textbf{类函数}：\(\chi_V(hgh^{-1}) = \chi_V(g)\)。因此 \(\chi_V\) 在共轭类上取常值。

\textbf{命题 69.2}（特征标的基本性质）： 1. \(\chi_V(e) = \dim V\)。 2. \(\chi_{V \oplus W} = \chi_V + \chi_W\)。 3. \(\chi_{V \otimes W} = \chi_V \cdot \chi_W\)。 4. \(\chi_V(g^{-1}) = \overline{\chi_V(g)}\)（对 \(g\) 的阶有限）。

\textbf{定义 69.5}（特征标的内积）：对于类函数 \(\varphi, \psi: G \to \mathbb{C}\)，定义内积

\[\langle \varphi, \psi \rangle = \frac{1}{|G|} \sum_{g \in G} \varphi(g) \overline{\psi(g)}\]

\textbf{定理 69.3}（不可约特征标的正交性）： 1. 若 \(\chi\) 是不可约表示的特征标，则 \(\langle \chi, \chi \rangle = 1\)。 2. 若 \(\chi\) 和 \(\psi\) 是两个不同构的不可约表示的特征标，则 \(\langle \chi, \psi \rangle = 0\)。 3. 不可约特征标构成类函数空间的一组正交规范基。

\textbf{证明}：设 \(\rho: G \to \operatorname{GL}(V)\) 和 \(\sigma: G \to \operatorname{GL}(W)\) 是两个不可约表示，特征标分别为 \(\chi, \psi\)。定义算子 \(T = \sum_{g \in G} \rho(g) \circ S \circ \sigma(g^{-1})\)，其中 \(S: W \to V\) 为任意线性映射。则 \(T\) 是 \(G\)-等变映射：\(\rho(h) T = \sum_g \rho(hg) S \sigma(g^{-1}) = \sum_g \rho(g) S \sigma(g^{-1}h) = T \sigma(h)\)。由 Schur 引理，若 \(\rho \not\cong \sigma\)，则 \(T = 0\)；若 \(\rho \cong \sigma\)，则 \(T = \lambda \operatorname{id}_V\)。取 \(S\) 为矩阵单位 \(E_{ij}\)（在 \(V\) 基下），计算 \(\operatorname{tr}(T)\) 得 \(\sum_g \rho_{ii}(g) \sigma_{jj}(g^{-1}) = \frac{|G|}{\dim V} \delta_{ij}\)（当 \(\rho \cong \sigma\)）。整理得 \(\langle \chi, \psi \rangle = \frac{1}{|G|} \sum_g \chi(g) \overline{\psi(g)} = \delta_{\rho \cong \sigma}\)。不可约特征标的正交规范性得证。\(\blacksquare\)

\textbf{推论 69.4}：表示 \(V\) 不可约当且仅当 \(\langle \chi_V, \chi_V \rangle = 1\)。

\textbf{定理 69.5}：\(G\) 的不可约表示的同构类个数等于 \(G\) 的共轭类个数。且若 \(d_1, \ldots, d_r\) 是不可约表示的维数，则

\[\sum_{i=1}^r d_i^2 = |G|\]

\textbf{证明}：类函数空间 \(\mathcal{C}(G)\) 的维数等于 \(G\) 的共轭类个数。由定理 69.3，不可约特征标构成 \(\mathcal{C}(G)\) 的一组正交规范基，故不可约表示个数等于 \(\dim \mathcal{C}(G) =\) 共轭类个数。对于维数平方和，取正则表示 \(\rho_{\text{reg}}\)（\(G\) 在 \(\mathbb{C}[G]\) 上的左乘作用），其分解为 \(\rho_{\text{reg}} \cong \bigoplus_{i=1}^r V_i^{\oplus d_i}\)（每个不可约表示 \(V_i\) 出现 \(\dim V_i\) 次）。比较维数：\(|G| = \dim \mathbb{C}[G] = \sum_{i=1}^r d_i \cdot d_i = \sum_{i=1}^r d_i^2\)。\(\blacksquare\)

\subsection{69.3 诱导表示与 Frobenius 互反律}\label{ux8bf1ux5bfcux8868ux793aux4e0e-frobenius-ux4e92ux53cdux5f8b}

\textbf{定义 69.6}（诱导表示）：设 \(H \leq G\)，\((\rho, W)\) 是 \(H\) 的表示。定义 \(G\) 的\textbf{诱导表示} \(\operatorname{Ind}_H^G W\) 为

\[\operatorname{Ind}_H^G W = \{f: G \to W : f(gh) = \rho(h^{-1})f(g), \; \forall h \in H\}\]

\(G\) 的作用为 \((\tilde{g} \cdot f)(g) = f(\tilde{g}^{-1} g)\)。

\textbf{定理 69.6}（Frobenius 互反律）：设 \(\chi\) 是 \(H\) 的特征标，\(\psi\) 是 \(G\) 的特征标。则

\[\langle \operatorname{Ind}_H^G \chi, \psi \rangle_G = \langle \chi, \operatorname{Res}_H^G \psi \rangle_H\]

其中 \(\operatorname{Res}_H^G \psi\) 是 \(\psi\) 限制为 \(H\) 的特征标。

\textbf{证明}：诱导特征标的值为 \(\operatorname{Ind}_H^G \chi(g) = \frac{1}{|H|} \sum_{x \in G} \dot{\chi}(xgx^{-1})\)（\(\dot{\chi}\) 在 \(H\) 外为零延拓）。计算内积得 \(\langle \operatorname{Ind}_H^G \chi, \psi \rangle_G = \frac{1}{|G||H|} \sum_{g \in G} \sum_{x \in G} \dot{\chi}(xgx^{-1}) \overline{\psi(g)}\)。交换求和并换元 \(h = xgx^{-1}\)，注意 \(g = x^{-1}hx\) 且 \(\psi(g) = \psi(h)\)（特征标为类函数）。对每固定的 \(x\)，\(g\) 跑遍 \(G\) 等价于 \(h\) 跑遍 \(G\)，故和式化为 \(\frac{1}{|H|} \sum_{h \in H} \chi(h) \overline{\psi(h)} = \langle \chi, \operatorname{Res}_H^G \psi \rangle_H\)。\(\blacksquare\)

\subsection{69.4 Burnside 定理}\label{burnside-ux5b9aux7406}

\textbf{定理 69.7}（Burnside \(p^a q^b\) 定理）：若 \(G\) 是有限群，\(|G| = p^a q^b\)（\(p, q\) 为素数），则 \(G\) 是可解群。

\textbf{证明}：反设 \(G\) 为非可解单群，\(|G| = p^a q^b\)。取 \(G\) 的非平凡共轭类 \(C\)，\(|C| = p^c q^d > 1\)。由特征标理论，存在不可约特征标 \(\chi\) 使 \(\chi(1)\) 与 \(|C|\) 互素（否则由特征标正交性导出的矛盾）。令 \(\omega_\chi(C) = \frac{|C| \chi(g)}{\chi(1)}\)（\(g \in C\)），此为代数整数。由互素条件，存在整数 \(u, v\) 使 \(u \chi(1) + v |C| = 1\)。则 \(\frac{1}{\chi(1)} = u + v \frac{|C|}{\chi(1)}\) 为代数整数。由此 \(\chi(1)\) 为代数整数也是有理数，故为整数。但 \(\chi(1)\) 整除 \(|G|\)，且由特征标正交性 \(\sum_{i} \chi_i(1)^2 = |G|\)。\(G\) 为单群意味着 \(\chi(1) = 1\) 仅对平凡特征标成立，其余 \(\chi(1) \geq 2\)。考虑 \(|G| = \sum \chi_i(1)^2\)，这迫使 \(|G|\) 必须为素数（否则由 Burnside 定理的群论分析，\(G\) 有非平凡正规子群）。矛盾。因此 \(G\) 有非平凡正规子群，归纳得 \(G\) 可解。\(\blacksquare\)

\textbf{推论 69.8}：最小非 Abel 单群是 \(A_5\)（阶 60 = \(2^2 \cdot 3 \cdot 5\)）。所有阶小于 60 的群都是可解群。

\hrulefill

\hrulefill

\hrulefill

\hrulefill

\hrulefill

\newpage

\chapter{05-代数/03-抽象代数二/04-同调代数.md}\label{ux4ee3ux657003-ux62bdux8c61ux4ee3ux6570ux4e8c04-ux540cux8c03ux4ee3ux6570.md}

\section{第128章：赋值论}\label{ux7b2c128ux7ae0ux8d4bux503cux8bba}

赋值论是研究域上''绝对值''及其完备化的理论，在代数数论和代数几何中至关重要。\textbf{赋值环}是整环\(R\)，其分式域\(K\)中任意非零元\(x\)满足\(x \in R\)或\(x^{-1} \in R\)；对应的\textbf{赋值}\(v: K^\times \to \Gamma\)（\(\Gamma\)为全序Abel群）满足\(v(xy) = v(x) + v(y)\)和\(v(x+y) \geq \min\{v(x), v(y)\}\)。当\(\Gamma \cong \mathbb{Z}\)时称为离散赋值环（DVR），其唯一极大理想\(\mathfrak{m} = \{x : v(x) > 0\}\)由单个元素\(\pi\)（素元）生成。\textbf{Hensel引理}是赋值论的核心定理：若\(R\)是完备赋值环，\(f \in R[x]\)及其模\(\mathfrak{m}\)约化\(\bar{f}\)满足\(\bar{f}\)有单根\(\bar{a}\)，则\(f\)有唯一根\(a \in R\)使得\(a \equiv \bar{a} \pmod{\mathfrak{m}}\)。Hensel环是满足Hensel引理的局部环，是局部域理论的基础。完备离散赋值环的结构由Cohen结构定理完全刻画：它们同构于域\(\kappa\)上的形式幂级数环\(\kappa[\![t]\!]\)或\(p\)-进整数环\(\mathbb{Z}_p\)的有限扩张。

\textbf{定义 72.1}（赋值）：域 \(K\) 上的 \textbf{（指数）赋值} 是满射 \(v: K^\times \to \Gamma\)，其中 \(\Gamma\) 是全序 Abel 群（值群），满足：（1）\(v(xy) = v(x) + v(y)\)；（2）\(v(x+y) \geq \min(v(x), v(y))\)。约定 \(v(0) = \infty\)（比 \(\Gamma\) 中任何元都大）。\textbf{赋值环} \(R_v = \{x \in K : v(x) \geq 0\}\) 是包含 \(K\) 的单位元且对 \(x \in K^\times\) 有 \(x \in R_v\) 或 \(x^{-1} \in R_v\) 的局部环。其极大理想 \(\mathfrak{m}_v = \{x \in K : v(x) > 0\}\)，剩余域 \(k_v = R_v / \mathfrak{m}_v\)。

\textbf{定义 72.2}（离散赋值与 Hensel 引理）：若 \(\Gamma \cong \mathbb{Z}\)，\(v\) 称为 \textbf{离散赋值}，相应的 \(R_v\) 为 \textbf{离散赋值环（DVR）}------其恰好是一维 Noether 正则局部环（Krull-Akizuki 定理）。\textbf{Hensel 引理}：设 \(R_v\) 完备（关于 \(v\) 诱导的拓扑），\(f \in R_v[x]\) 满足 \(\bar{f} = f \bmod \mathfrak{m}_v\) 在 \(k_v[x]\) 中有单根 \(\bar{\alpha}\)，则存在唯一的 \(\alpha \in R_v\) 使得 \(\bar{\alpha} = \alpha \bmod \mathfrak{m}_v\) 且 \(f(\alpha) = 0\)。

\textbf{定理 72.1}（Cohen 结构定理）：完备的等特征 Noether 局部环同构于其剩余域上的形式幂级数环 \(k[[x_1, \ldots, x_d]]\) 的商环。完备的混合特征 DVR 同构于特征零剩余域特征 \(p\) 的 Witt 向量环的分歧扩张。

\textbf{证明}：设 \((R, \mathfrak{m}, k)\) 为完备 Noether 局部环。若 \(R\) 等特征（\(\operatorname{char} R = \operatorname{char} k\)），由 Cohen 结构定理，存在截面 \(k \hookrightarrow R\)（等特征时由 Hensel 引理保证）。选取参数系 \(x_1, \ldots, x_d\)（\(d = \dim R\)），由完备性，自然映射 \(k[[x_1, \ldots, x_d]] \to R\) 为满射，核为理想，故 \(R \cong k[[x_1, \ldots, x_d]] / I\)。若 \(R\) 混合特征（\(\operatorname{char} R = 0\)，\(\operatorname{char} k = p\)），则 \(R\) 含 Cohen 子环 \(C\)（完备离散赋值环，极大理想为 \(pC\)，剩余域为 \(k\)），\(C\) 同构于 \(W(k)\)（Witt 向量环）的商环。\(R\) 是 \(C\) 上的有限生成模，由完备性 \(R \cong C[[x_1, \ldots, x_{d-1}]] / I\)。\(\blacksquare\)

\hrulefill

\hrulefill

\hrulefill

\hrulefill

\hrulefill

\hrulefill

\section{第129章：胎紧闭包（Tight Closure）}\label{ux7b2c129ux7ae0ux80ceux7d27ux95edux5305tight-closure}

胎紧闭包（Tight Closure）是 Hochster 和 Huneke 于 1986 年引入的特征 \(p > 0\) 交换环中的闭包运算，为交换代数带来了一系列深刻的定理。这一理论提供了不同于经典同调代数的新工具，在解决若干长期未决猜想（特别是与正则环结构和理想理论相关的难题）中发挥了核心作用。

\subsection{73.1 基本定义与 Frost 性质}\label{ux57faux672cux5b9aux4e49ux4e0e-frost-ux6027ux8d28}

\textbf{定义 73.1}（理想的 Frobenius 幂）：设 \(R\) 是特征 \(p > 0\) 的含幺交换 Noether 环。对理想 \(I \subseteq R\) 和 \(q = p^e\)（\(e \geq 0\)），定义 \(I\) 的 \(q\) 次 \textbf{Frobenius 幂} \(I^{[q]}\) 为由 \(\{x^q : x \in I\}\) 生成的理想。即

\[I^{[q]} = \langle x^q : x \in I \rangle_R\]

记号 \(I^{[q]}\)（方括号）区别于通常的幂 \(I^q\)。通常 \(I^{[q]} \subseteq I^q\)，且若 \(I = (x_1, \ldots, x_n)\)，则 \(I^{[q]} = (x_1^q, \ldots, x_n^q)\)。

\textbf{定义 73.2}（胎紧闭包）：设 \(R\) 是特征 \(p > 0\) 的含幺交换 Noether 环，\(I \subseteq R\) 为理想。\(I\) 的\textbf{胎紧闭包}（tight closure）\(I^*\) 定义为

\[I^* = \left\{x \in R : \exists\, c \in R^\circ,\; \forall\, q = p^e \gg 0,\; c x^q \in I^{[q]}\right\}\]

其中 \(R^\circ\) 表示 \(R\) 中所有极小素理想的补集之并（即不在任何极小素理想中的元素全体）。元素 \(c\) 称为 \(x\) 对 \(I\) 的\textbf{测试元}（test element）。通常可选取不依赖于 \(x\) 和 \(I\) 的公共 \(c\)，这样的 \(c\) 称为\textbf{胎紧闭包测试元}（tight closure test element）。

胎紧闭包满足以下基本性质（Frost 性质）： - \(I \subseteq I^* \subseteq \overline{I}\)（其中 \(\overline{I}\) 为整闭包）； - 若 \(I \subseteq J\)，则 \(I^* \subseteq J^*\)； - \((I^*)^* = I^*\)（幂等性）； - \((I + J)^* = I^* + J^*\) 不一定成立，但 \((I \cap J)^* = I^* \cap J^*\)。

\subsection{73.2 Hochster-Huneke 定理}\label{hochster-huneke-ux5b9aux7406}

\textbf{定理 73.1}（Hochster-Huneke 正则环胎紧闭包定理）：若 \(R\) 是特征 \(p > 0\) 的 Noether 正则环，则对任意理想 \(I \subseteq R\)，有

\[I^* = I\]

即正则环中胎紧闭包等于理想本身。

\textbf{证明}：核心是证明正则环弱 F-正则，即 \(I^* = I\) 对所有理想 \(I\) 成立。（1）对正则局部环 \((R, \mathfrak{m})\)，Frobenius 映射 \(F: R \to R\)（\(x \mapsto x^p\)）是平坦的（Kunz 定理），因此 \(F^e(R)\) 是平坦 \(R\)-模。对 \(x \in I^*\)，存在 \(c \in R^\circ\) 使对所有大 \(q = p^e\) 有 \(c x^q \in I^{[q]}\)。利用 \(F^e\) 的平坦性，将此包含关系沿 \(F^e\) 拉回，得到 \(c^{1/q} x \in I\)（在 \(R\) 的完美闭包中）。由于 \(R\) 正则，存在 test element \(c = 1\)（正则环的 Frobenius 分裂性），故 \(x \in I\)。（2）一般情形通过局部化、完备化和归纳维数约化：胎紧闭包与局部化交换，正则环局部化仍正则，完备正则局部环是形式幂级数环，由(1)知 \(I^* = I\)。\(\blacksquare\)

\textbf{推论 73.2}：正则环中，整闭包与胎紧闭包一致（因为 \(I = I^*\)，而 \(I \subseteq \overline{I}\)，但 \(I^* \subseteq \overline{I}\)，故 \(I = I^* = \overline{I}\) 仅当 \(I\) 本身已整闭）。

\subsection{\texorpdfstring{73.3 Briançon-Skoda 定理（特征 \(p\) 版本）}{73.3 Briançon-Skoda 定理（特征 p 版本）}}\label{brianuxe7on-skoda-ux5b9aux7406ux7279ux5f81-p-ux7248ux672c}

\textbf{定理 73.3}（Tight Closure 版本的 Briançon-Skoda 定理）：设 \(R\) 是特征 \(p > 0\) 的 Noether 正则环，\(I\) 是由 \(n\) 个元素生成的理想。则对任意 \(m \geq 1\)，

\[\overline{I^{n+m-1}} \subseteq I^m\]

特别地，取 \(m = 1\) 得 \(\overline{I^n} \subseteq I\)，即由 \(n\) 个元生成的理想的 \(n\) 次幂的整闭包包含于原理想。

\textbf{证明}：（1）证明整闭包含于胎紧闭包：对 \(x \in \overline{I^{n+m-1}}\)，由整闭包的刻画，存在 \(c \neq 0\) 使对所有 \(q = p^e\) 有 \(c x^q \in (I^{n+m-1})^{[q]} = I^{[q](n+m-1)}\)。由 Frobenius 幂的分配律，\(I^{[q](n+m-1)} \subseteq (I^{[q]})^{n+m-1}\)。利用胎紧闭包的定义，需要验证 \(c x^q \in (I^m)^{[q]}\)。由 \(I\) 的 \(n\) 个生成元，利用 Frobenius 映射的 \(n\)-线性性质，\(I^{[q](n+m-1)} \subseteq I^{[q]m}\) 对 \(q \geq n\) 成立（由 \(I\) 的生成元个数和 Frobenius 幂的分配律），故 \(x \in (I^m)^*\)。（2）由定理 73.1，正则环中 \((I^m)^* = I^m\)，故 \(\overline{I^{n+m-1}} \subseteq I^m\)。\(\blacksquare\)

该定理的重大意义在于：它将一个高度非平凡的包含关系（整闭包包含于原理想）化归为胎紧闭包的简单性质。这在复解析几何中也有对应：若 \(\mathcal{I}\) 是由 \(n\) 个全纯函数生成的理想层，则 \(\overline{\mathcal{I}^n} \subseteq \mathcal{I}\)（Skoda 定理）。

\subsection{73.4 Colon-Capturing 与单项式猜想}\label{colon-capturing-ux4e0eux5355ux9879ux5f0fux731cux60f3}

\textbf{定理 73.4}（Colon-Capturing，Hochster-Roberts）：设 \(R\) 是特征 \(p > 0\) 的 Cohen-Macaulay 局部环，\(\mathfrak{a} \subseteq R\) 是理想，\(x_1, \ldots, x_d\) 为参数系（system of parameters）。则

\[(\mathfrak{a}, x_1, \ldots, x_{i-1}) : x_i \subseteq (\mathfrak{a}, x_1, \ldots, x_{i-1})^*\]

其中左侧为理想商（colon ideal），\(i = 1, \ldots, d\)。

\textbf{证明}：设 \(y \in (\mathfrak{a}, x_1, \ldots, x_{i-1}) : x_i\)，即 \(y x_i \in (\mathfrak{a}, x_1, \ldots, x_{i-1})\)。需证 \(y \in (\mathfrak{a}, x_1, \ldots, x_{i-1})^*\)。取 \(q = p^e\)，由 Frobenius 幂的性质，\((y x_i)^q = y^q x_i^q \in (\mathfrak{a}, x_1, \ldots, x_{i-1})^{[q]}\)。由于 \(R\) 是 Cohen-Macaulay 环，\(x_1, \ldots, x_d\) 为正则序列，\(x_i^q\) 在 \(R/(\mathfrak{a}, x_1, \ldots, x_{i-1})\) 上非零因子。由胎紧闭包定义，取 \(c = x_i^q\)（适当选取 \(q\) 足够大），则 \(c y^q \in (\mathfrak{a}, x_1, \ldots, x_{i-1})^{[q]}\)。因此 \(y \in (\mathfrak{a}, x_1, \ldots, x_{i-1})^*\)。\(\blacksquare\)

\textbf{单项式猜想（Monomial Conjecture）} 是 Hochster 1973 年提出的核心猜想：设 \((R, \mathfrak{m})\) 是 \(d\) 维 Noether 局部环，\(x_1, \ldots, x_d\) 为参数系，则对任意 \(t \geq 1\)，

\[x_1^t \cdots x_d^t \notin (x_1^{t+1}, \ldots, x_d^{t+1})\]

这个猜想起初看似仅仅是一个组合代数的问题，实则极其深刻，与 Krull 交定理、符号幂理论、同调猜想等多个领域紧密相关。Hochster 证明了：若特征 \(p > 0\) 情形下的单项式猜想成立，则由约化模 \(p\) 技术可推出等特征 0 情形的猜想。

利用 colon-capturing 定理，Hochster 证明了特征 \(p\) 情形下的单项式猜想。实际上，由 colon-capturing 可推出：若 \(R\) 是 Cohen-Macaulay 局部环，则对参数系 \(x_1, \ldots, x_d\) 有

\[(x_1, \ldots, x_i) : x_{i+1} \subseteq (x_1, \ldots, x_i)^*\]

这意味着参数系序列构成一个在胎紧闭包意义下的正则序列。单项式猜想作为该结论的特殊情形得以证明。

\subsection{73.5 Test Element 的存在性}\label{test-element-ux7684ux5b58ux5728ux6027}

\textbf{定理 73.5}（Test Element 存在性定理，Hochster-Huneke）：设 \(R\) 是特征 \(p > 0\) 的既约极好（excellent）局部环。则存在 \(c \in R^\circ\) 使得对任意理想 \(I \subseteq R\) 和任意 \(x \in I^*\)，有对所有大 \(q = p^e\)，\(c x^q \in I^{[q]}\)。这样的 \(c\) 称为 \(R\) 的\textbf{完全测试元}（complete test element）。

\textbf{证明}：（1）设 \(R\) 为既约极好局部环。对 \(R\) 的正规化 \(R'\)，由于 \(R \subseteq R'\) 是有限扩张，\(R'^\circ \cap R \subseteq R^\circ\)。在 \(R'\) 中，利用 \(\Gamma\) 构造：对 \(R\) 的任意参数系 \(x_1, \ldots, x_d\)，令 \(c = \prod_{i=1}^d x_i^{p-1}\)（Frobenius 邻域）。验证 \(c\) 是 \(R\) 的测试元：对任意 \(x \in I^*\)，存在 \(d \in R^\circ\) 使对所有 \(q\) 有 \(d x^q \in I^{[q]}\)。由极好性，参数系的选取使得 \(c\) 整除 \(d\) 在 \(R\) 的某次 Frobenius 拉回中的像，因此 \(c x^q \in I^{[q]}\) 对所有大 \(q\) 成立。（2）完备化：\(\widehat{R}\) 中测试元的存在性通过 \(R\) 中测试元在完备化下的像保证，利用极好环的纤维正则性。（3）由步骤(1)(2)，\(c\) 是统一的完全测试元。\(\blacksquare\)

Test element 的存在性保证了 tight closure 理论的运算基础：无需为每个特定的 \(x\) 和 \(I\) 分别寻找 \(c\)，可以使用一个统一的 \(c\) 来测试所有胎紧闭包关系。这也使得胎紧闭包可视为一个''闭合算子''（closure operator），其理论和计算工具与理想整闭包理论类似。

Test element 的存在性还具有更深层的理论意义。在正则环中，\(\mathbb{F}\)-分裂性（Frobenius splitting）保证了 \(c = 1\) 可作为测试元，这等价于正则环的胎紧闭包等于原理想。对于一般极好环，统一测试元的存在使得胎紧闭包与整闭包的关系可被精确量化：若 \(c\) 是测试元，则 \(I^* = \bigcup_{e \geq 0} (I^{[p^e]} :_R c)^{1/p^e}\)（其中 \(J^{1/p^e}\) 表示 \(R\) 的 \(p^e\) 次根闭包中满足 \(a^{p^e} \in J\) 的元素的集合）。这一刻画将胎紧闭包的计算转化为 Frobenius 幂的理想商运算，在计算实践中具有重要价值。对于 Cohen-Macaulay 环，测试元可被选取为参数系的乘积的 \(p-1\) 次幂，且该测试元还满足 colon-capturing 性质------这为非 Cohen-Macaulay 环中同调猜想的证明提供了关键工具。测试元的选择自由度（任何 \(c\) 的非零倍也是测试元）意味着存在''最优''测试元的选择问题，与环的 Frobenius 闭包中 multiplier 理想理论密切相关。

\subsection{73.6 应用与展望}\label{ux5e94ux7528ux4e0eux5c55ux671b}

胎紧闭包理论推动了交换代数的多个核心发展：

\begin{enumerate}
\def\labelenumi{\arabic{enumi}.}
\item
  \textbf{同调猜想}：Peskine-Szpiro 的相交定理和 Hochster 的同调猜想（直接和猜想、提升猜想）均可由胎紧闭包方法在特征 \(p\) 情形下证明，再通过约化模 \(p\) 推广到等特征 0 情形。
\item
  \textbf{Buchsbaum-Eisenbud 的秩猜想}：胎紧闭包方法被用于证明关于有限自由分解的最小秩的若干猜想。
\item
  \textbf{符号幂理论}：Ein-Lazarsfeld-Smith（2001 年）利用胎紧闭包的精化版本（multiplier ideals 和 asymptotic multiplier ideals）证明了光滑复簇上符号幂的渐近等式 \(I^{(rn)} \subseteq I^n\)。
\item
  \textbf{正特征代数几何}：胎紧闭包的对偶概念------\textbf{F-奇点理论}（F-singularities）------将正特征中的 Frobenius 分裂与复代数几何中的奇点理论（log terminal、log canonical 等）建立了对应。
\end{enumerate}

\hrulefill

\hrulefill

\hrulefill

\hrulefill

\hrulefill

\hrulefill

\subsection{\texorpdfstring{稳定\(\infty\)-范畴与导出范畴理论}{稳定\textbackslash infty-范畴与导出范畴理论}}\label{ux7a33ux5b9ainfty-ux8303ux7574ux4e0eux5bfcux51faux8303ux7574ux7406ux8bba}

稳定 \(\infty\)-范畴是 21 世纪同调代数最深刻的革新之一，它由 Jacob Lurie 在其 2007 年里程碑式的著作《Higher Topos Theory》和 2011 年的《Higher Algebra》中系统建立。经典同调代数（Cartan-Eilenberg, 1956）以 Abel 范畴和导出范畴为框架，而 Grothendieck 和 Verdier 在 1960 年代发展的三角范畴理论（导出范畴）虽然极大推动了代数几何和表示论的发展，但三角范畴的八面体公理（TR4）缺乏范畴论的自然性------它是由模型结构人工添加的，而非从范畴本身的内在性质自动导出。稳定 \(\infty\)-范畴的精髓在于，三角结构不是外加的公理，而是稳定性公理的推论：在稳定的 \(\infty\)-范畴中，推出-拉回等价自动诱导了同伦范畴上的三角结构，且八面体公理自然成立。

稳定 \(\infty\)-范畴的三条公理------有点（有零对象）、有限完备和有限余完备、推出-拉回等价------看似简单，却蕴含着极其丰富的结构。推出-拉回等价是稳定性的核心：它意味着 \(\infty\)-范畴中的纤维序列和余纤维序列是同一回事，从而悬垂函子 \(\Sigma\) 和环路函子 \(\Omega\) 互为逆等价。这一性质使得稳定 \(\infty\)-范畴自然地成为同伦论意义的''谱对象''的范畴，而经典导出范畴 \(\mathcal{D}(R)\) 正是稳定 \(\infty\)-范畴的典范例子------它是链复形范畴在拟同构处的 \(\infty\)-范畴局部化。在 \(\infty\)-范畴框架下，导出函子 \(\mathbb{R}F\) 的构造无需选取投射/内射分解，而是通过 Kan 延拓自然给出，这一观点的清晰性和普遍性远超经典方法。

稳定 \(\infty\)-范畴理论还澄清了 t-结构和 Verdier 商的本质。t-结构是稳定 \(\infty\)-范畴上的''截断''结构，它推广了导出范畴中标准 t-结构（截断复形），其心（heart）是 Abel 范畴------这为从同伦论信息中恢复 Abel 范畴信息提供了桥梁。Verdier 商在 \(\infty\)-范畴框架下等同于 Bousfield 局部化，这一视角将三角范畴理论重新嵌入同伦论的大家庭，使得导出范畴理论不再是孤立的代数构造，而是同伦论------稳定同伦论------代数几何这一宏大统一图景的有机组成部分。本章将系统介绍稳定 \(\infty\)-范畴的公理化定义、导出范畴的 \(\infty\)-范畴构造、t-结构和导出函子的自然定义，为后续的导出代数几何（卷十五）和谱代数几何（卷二十）奠定基础。

\textbf{定义 1}（\(\infty\)-范畴）：设 \(\mathcal{C}\) 是 \(\infty\)-范畴（准范畴模型：单纯集满足内部角填充条件）。\(\mathcal{C}\) 的\textbf{同伦范畴} \(\mathrm{Ho}(\mathcal{C})\) 是其对象同伦类的通常范畴。

\textbf{定义 2}（零对象与点点）：\(\infty\)-范畴 \(\mathcal{C}\) 中的对象 \(0\) 称为\textbf{零对象}，若它在 \(\mathrm{Ho}(\mathcal{C})\) 中是始对象又是终对象。具有零对象的 \(\infty\)-范畴称为\textbf{点点的}。

\textbf{定义 3}（有限极限与余极限）：\(\infty\)-范畴 \(\mathcal{C}\) 称为\textbf{有限完备}（有限余完备）的，若它有所有有限极限（余极限）。等价于：\(\mathcal{C}\) 有终对象（始对象）和所有拉回（推出）。

\textbf{定义 4}（稳定\(\infty\)-范畴）：\(\infty\)-范畴 \(\mathcal{C}\) 称为\textbf{稳定的}（stable），若满足以下三条公理：

\begin{enumerate}
\def\labelenumi{\arabic{enumi}.}
\tightlist
\item
  \(\mathcal{C}\) 是有点的（有零对象）。
\item
  \(\mathcal{C}\) 有所有有限极限和有限余极限。
\item
  \(\mathcal{C}\) 中的一个交换正方是推出当且仅当它是拉回。
\end{enumerate}

第三条称为\textbf{推出-拉回等价}（pushout-pullback equivalence），是稳定性概念的核心。

\textbf{定理 1}（稳定\(\infty\)-范畴的同伦范畴是三角范畴）：设 \(\mathcal{C}\) 是稳定 \(\infty\)-范畴。则其同伦范畴 \(\mathrm{Ho}(\mathcal{C})\) 具有自然三角范畴结构，其中： - 平移函子 \(\Sigma: \mathrm{Ho}(\mathcal{C}) \to \mathrm{Ho}(\mathcal{C})\) 由推出-拉回性质诱导（将对象 \(X\) 映射为 \(X \to 0 \to 0\) 的推出），其逆 \(\Omega\) 由拉回-推出性质诱导。 - 好三角（distinguished triangles）形如 \(X \to Y \to Z \to \Sigma X\)，由 \(\mathcal{C}\) 中的推出（或纤维-余纤维）序列诱导。

\textbf{证明}：由稳定性公理 3，推出-拉回等价定义了悬垂函子 \(\Sigma\) 和环路函子 \(\Omega\)。对 \(X \in \mathcal{C}\)，构造推出 \(X \to 0 \to 0\) 得 \(\Sigma X\)，且此正方同时为拉回，故 \(\Omega \Sigma X \cong X\)（在同伦范畴中）。\(\Sigma\) 和 \(\Omega\) 互为逆等价。好三角定义为同构于 \(X \xrightarrow{f} Y \to \operatorname{cofib}(f) \to \Sigma X\)（其中 \(\operatorname{cofib}(f)\) 为 \(f\) 的余纤维）的序列。验证三角范畴公理：（TR1）恒等态射的余纤维为零，\(X \to X \to 0 \to \Sigma X\) 是好三角；（TR2）旋转由推出-拉回交换图保证；（TR3）填充由推出和拉回的万有性质给出；（TR4）八面体公理由 \(\mathcal{C}\) 中双重推出图的组合性质自动导出，无需在 \(\mathrm{Ho}(\mathcal{C})\) 级别手工验证。\(\blacksquare\)

\subsection{典范例子}\label{ux5178ux8303ux4f8bux5b50}

\textbf{定理 2}（链复形的\(\infty\)-范畴）：设 \(R\) 是交换环。链复形范畴 \(\mathrm{Ch}(R)\) 可以通过 Dold-Kan 对应视为单纯 \(R\)-模的范畴，赋予投射模型结构后，其 \(\infty\)-范畴局部化 \(\mathrm{Ch}(R)[W^{-1}]\)（\(W\) 为拟同构类）是稳定 \(\infty\)-范畴。

\textbf{证明}：Dold-Kan 对应给出 \(\mathrm{Ch}_{\ge 0}(R) \cong \mathbf{sMod}_R\)（单纯 \(R\)-模范畴）。在 \(\mathbf{sMod}_R\) 上赋予投射模型结构（弱等价为拟同构，纤维化为 Kan 纤维化），其 \(\infty\)-范畴局部化是 \(\mathcal{D}_{\ge 0}(R)\)。验证稳定性：零对象为零复形，推出-拉回等价由映射锥的万有性质保证------对链映射 \(f: A \to B\)，映射锥 \(C(f)\) 同时是 \(f\) 的同伦余纤维和同伦纤维（在平移下），且 \(\Omega C(f) \cong A\)。平移函子 \([1]\) 是自等价，故 \(\mathrm{Ch}(R)[W^{-1}]\) 满足稳定 \(\infty\)-范畴的三条公理。\(\blacksquare\)

\textbf{定理 3}（导出\(\infty\)-范畴）：\(R\) 的\textbf{导出 \(\infty\)-范畴} \(\mathcal{D}(R)\) 是稳定 \(\infty\)-范畴，由 \(\mathrm{Ch}(R)\) 在拟同构处局部化得到：

\[\mathcal{D}(R) = \mathrm{Ch}(R)[W^{-1}]\]

其中 \(W\) 是拟同构的集合。其在同伦范畴级别恰为经典导出范畴 \(\mathrm{Ho}(\mathcal{D}(R)) \cong D(R)\)。

\textbf{证明}：由定理 2，\(\mathrm{Ch}(R)[W^{-1}]\) 是稳定 \(\infty\)-范畴。其同伦范畴 \(\mathrm{Ho}(\mathcal{D}(R))\) 的对象是链复形，态射是 \(\operatorname{Hom}_{K(R)}(A, B) = H^0(\operatorname{RHom}(A, B))\)（链映射的链同伦类）。由于 \(W\) 恰为拟同构类，\(\mathrm{Ho}(\mathcal{D}(R))\) 中拟同构变为同构，这正是经典导出范畴 \(D(R)\) 的构造。\(\mathcal{D}(R)\) 的稳定性由 \(\mathrm{Ch}(R)\) 的模型结构继承：平移 \([1]\) 为复形平移，映射锥给出好三角，推出-拉回等价由映射锥的万有性质保证。\(\blacksquare\)

\textbf{定义 5}（谱的\(\infty\)-范畴）：\textbf{谱的 \(\infty\)-范畴} \(\mathrm{Sp}\) 由拓扑谱的模型范畴经 Bousfield 局部化（在稳定同伦等价处）得到。它是稳定 \(\infty\)-范畴的典范例子，其同伦范畴是代数拓扑中经典的稳定同伦范畴 \(\mathcal{SH}\)。

\(\mathrm{Sp}\) 的单位对象是球谱 \(\mathbb{S}\)，且 \(\mathrm{Sp}\) 是对称幺半的：张量积为谱的粉碎积（smash product）\(\wedge\)。

\subsection{稳定化函子}\label{ux7a33ux5b9aux5316ux51fdux5b50}

\textbf{定义 6}（谱对象）：设 \(\mathcal{C}\) 是有限极限的 \(\infty\)-范畴。\(\mathcal{C}\) 中的\textbf{谱对象}（spectrum object）是一个序列 \(\{X_n\}_{n \ge 0}\)，配备等价 \(X_n \cong \Omega X_{n+1}\)。谱对象的 \(\infty\)-范畴记作 \(\mathrm{Sp}(\mathcal{C})\)。

\textbf{定理 4}（稳定化函子的普遍性质）：设 \(\mathcal{C}\) 是点点的有限极限 \(\infty\)-范畴。则 \(\mathrm{Sp}(\mathcal{C})\) 是稳定 \(\infty\)-范畴，且存在伴随函子对：

\[\Sigma^\infty: \mathcal{C} \rightleftarrows \mathrm{Sp}(\mathcal{C}) : \Omega^\infty\]

其中 \(\Omega^\infty\) 将谱对象映为其第零个成分 \(X_0\)，而 \(\Sigma^\infty X\) 是 \(X\) 的\textbf{悬垂谱} \(\{\Sigma^n X\}_{n \ge 0}\)。

此构造是普遍的：对任意稳定 \(\infty\)-范畴 \(\mathcal{D}\) 和左正合函子 \(F: \mathcal{C} \to \mathcal{D}\)，存在唯一的正合函子 \(\tilde{F}: \mathrm{Sp}(\mathcal{C}) \to \mathcal{D}\) 使下图交换（在同伦意义下）：

\[\begin{CD}
\mathcal{C} @>{\Sigma^\infty}>> \mathrm{Sp}(\mathcal{C}) \\
@V{F}VV @VV{\tilde{F}}V \\
\mathcal{D} @= \mathcal{D}
\end{CD}\]

\(\mathrm{Sp}(\mathcal{C})\) 被称为 \(\mathcal{C}\) 的\textbf{自由稳定化}。

\textbf{证明}：验证 \(\mathrm{Sp}(\mathcal{C})\) 的稳定性：零对象是常零谱 \(X_n = 0\)。有限极限逐层计算：\(\mathrm{Sp}(\mathcal{C})\) 中的极限 \(\varprojlim X^{(i)}\) 满足 \((\varprojlim X^{(i)})_n = \varprojlim X^{(i)}_n\)，且 \(\Omega\) 作用逐层交换。推出-拉回等价：\(\mathrm{Sp}(\mathcal{C})\) 中的交换正方是推出（拉回）当且仅当每个层级的对应正方是 \(\mathcal{C}\) 中的推出（拉回）；由于 \(\mathcal{C}\) 中 \(\Omega\) 与极限交换，\(\mathrm{Sp}(\mathcal{C})\) 的推出-拉回等价由 \(\mathcal{C}\) 的推出-拉回等价继承。万有性质：给定稳定 \(\mathcal{D}\) 和保有限极限的 \(F: \mathcal{C} \to \mathcal{D}\)，定义 \(\tilde{F}(X_0, X_1, \ldots) = \operatorname{colim}_n \Omega_{\mathcal{D}}^n F(X_n)\)，其中 \(\Omega_{\mathcal{D}}\) 是 \(\mathcal{D}\) 的环路等价。验证 \(\tilde{F}\) 正合且 \(\tilde{F} \circ \Sigma^\infty \cong F\)。\(\blacksquare\)

\subsection{t-结构}\label{t-ux7ed3ux6784}

\textbf{定义 7}（t-结构）：稳定 \(\infty\)-范畴 \(\mathcal{C}\) 上的一个 \textbf{t-结构} 是一对满子范畴 \((\mathcal{C}_{\ge 0}, \mathcal{C}_{\le 0})\)，满足以下三条公理：

\begin{enumerate}
\def\labelenumi{\arabic{enumi}.}
\tightlist
\item
  \(\mathcal{C}_{\ge 0} \subseteq \mathcal{C}_{\ge 0}[1]\) 和 \(\mathcal{C}_{\le 0}[1] \subseteq \mathcal{C}_{\le 0}\)（此处 \([1] = \Sigma\) 表示平移）。
\item
  \(\mathrm{Hom}_{\mathcal{C}}(X, Y) = 0\) 对所有 \(X \in \mathcal{C}_{\ge 0}\) 和 \(Y \in \mathcal{C}_{\le 0}[-1] = \mathcal{C}_{\le -1}\)。
\item
  每个 \(X \in \mathcal{C}\) 可嵌入好三角 \(X_{\ge 0} \to X \to X_{\le -1}\)，其中 \(X_{\ge 0} \in \mathcal{C}_{\ge 0}\)，\(X_{\le -1} \in \mathcal{C}_{\le -1}\)。
\end{enumerate}

\textbf{定义 8}（心）：t-结构 \((\mathcal{C}_{\ge 0}, \mathcal{C}_{\le 0})\) 的\textbf{心}（heart）定义为：

\[\mathcal{C}^{\heartsuit} = \mathcal{C}_{\ge 0} \cap \mathcal{C}_{\le 0}\]

\textbf{定理 5}（心是Abel范畴）：稳定 \(\infty\)-范畴上 t-结构的心 \(\mathcal{C}^{\heartsuit}\) 是 Abel 范畴。

\textbf{证明}：验证 \(\mathcal{C}^{\heartsuit}\) 满足 Abel 范畴公理。对 \(f: X \to Y\)（\(X, Y \in \mathcal{C}^{\heartsuit}\)），构造纤维 \(F = \operatorname{fib}(f)\)，满足纤维序列 \(F \to X \xrightarrow{f} Y\)。由 t-结构公理 3，\(F\) 有截断 \(F_{\ge 0} \to F \to F_{\le -1}\)。取 \(\ker(f) = F_{\ge 0}\)，则 \(\ker(f) \in \mathcal{C}_{\ge 0}\)。由 \(F \in \mathcal{C}_{\ge 0}\)（因 \(X, Y \in \mathcal{C}_{\ge 0}\) 且 \(\mathcal{C}_{\ge 0}\) 对纤维封闭），\(F_{\le -1} = 0\)，故 \(\ker(f) = F \in \mathcal{C}^{\heartsuit}\)。对偶构造 \(\operatorname{coker}(f) = \operatorname{cofib}(f)_{\le 0}\)。验证 \(\ker(\operatorname{coker}(f)) \cong \operatorname{coker}(\ker(f))\)：由纤维-余纤维序列的推出-拉回交换图和 t-结构截断的相容性，两端的正合性确保此典范同构。加法群结构由 \(\pi_0 \operatorname{Map}(X,Y)\) 赋予，其中加法由 \(\mathcal{C}\) 的环路-悬挂伴随导出。\(\blacksquare\)

\textbf{例子}：导出 \(\infty\)-范畴 \(\mathcal{D}(R)\) 带标准 t-结构：\(\mathcal{D}(R)_{\ge 0}\) 由链复形集中在非负次数（同调集中在 \(\ge 0\)）的对象张成，\(\mathcal{D}(R)_{\le 0}\) 由集中在非正次数（\(\le 0\)）的对象张成。其心为 \(R\)-模范畴 \(\mathrm{Mod}_R\)。

\subsection{\texorpdfstring{导出函子在\(\infty\)-范畴框架下的自然定义}{导出函子在\textbackslash infty-范畴框架下的自然定义}}\label{ux5bfcux51faux51fdux5b50ux5728infty-ux8303ux7574ux6846ux67b6ux4e0bux7684ux81eaux7136ux5b9aux4e49}

\textbf{定理 6}（导出函子的\(\infty\)-范畴定义）：设 \(F: \mathcal{A} \to \mathcal{B}\) 是 Abel 范畴间的左正合函子。则 \(\mathcal{A}\) 和 \(\mathcal{B}\) 的导出 \(\infty\)-范畴间存在导出函子 \(\mathbb{R}F: \mathcal{D}(\mathcal{A}) \to \mathcal{D}(\mathcal{B})\)，由以下普遍性质刻画：

\[\mathrm{Ho}(\mathbb{R}F) \cong \mathbf{R}F\]

（等号在函子同构意义下，左侧为 \(\infty\)-范畴函子的同伦范畴函子，右侧为经典导出函子）。

\textbf{构造}：在 \(\infty\)-范畴框架下，\(\mathbb{R}F\) 的构造无需选取投射/内射分解。改为利用 Kan 延拓：将 \(\mathcal{A}\) 视为 \(\mathcal{D}(\mathcal{A})\) 的子范畴（通过 \(D \mapsto D[0]\) 嵌入），\(F\) 自然地诱导了一个从 \(\mathcal{A}\) 到 \(\mathcal{D}(\mathcal{B})\) 的函子。由 \(\infty\)-范畴局部化的普遍性质，此函子可唯一正合地延拓到 \(\mathcal{D}(\mathcal{A})\) 上：

\[\mathbb{R}F = \mathrm{Lan}_{i_{\mathcal{A}}} (i_{\mathcal{B}} \circ F)\]

其中 \(i_{\mathcal{A}}: \mathcal{A} \hookrightarrow \mathcal{D}(\mathcal{A})\) 和 \(i_{\mathcal{B}}: \mathcal{B} \hookrightarrow \mathcal{D}(\mathcal{B})\) 是标准嵌入函子，\(\mathrm{Lan}\) 表示左 Kan 延拓。

\textbf{证明}：由 \(\infty\)-范畴的 Kan 延拓理论，复合 \(i_{\mathcal{B}} \circ F\) 沿 \(i_{\mathcal{A}}\) 的左 Kan 延拓存在------因为 \(\mathcal{D}(\mathcal{A})\) 是余完备的（稳定 \(\infty\)-范畴有所有余极限）。此延拓自动是正合函子（稳定 \(\infty\)-范畴之间的左伴随函子保推出，故在三角意义下正合）。由其定义直接验证在适当条件下的经典导出函子一致：对 \(K\)-内射或 \(K\)-投射复形，\(\mathbb{R}F\) 的作用恰为逐项应用 \(F\)。\(\blacksquare\)

\subsection{Verdier商作为Bousfield局部化}\label{verdierux5546ux4f5cux4e3abousfieldux5c40ux90e8ux5316}

\textbf{定义 9}（Verdier商的经典定义）：设 \(\mathcal{T}\) 是经典三角范畴，\(\mathcal{S} \subseteq \mathcal{T}\) 是厚的三角子范畴。\textbf{Verdier商} \(\mathcal{T}/\mathcal{S}\) 的构造为：\(\mathcal{T}\) 中的态射空间商去所有以 \(\mathcal{S}\) 中对象为中间项的分解。

\textbf{定义 10}（Bousfield局部化）：设 \(\mathcal{C}\) 是 \(\infty\)-范畴，\(W\) 是态射的集合。\textbf{Bousfield 局部化} \(L_W \mathcal{C}\) 是 \(\mathcal{C}\) 在 \(W\) 处的普遍局部化，其构造为 \(\mathcal{C}\) 配备一个函子 \(L: \mathcal{C} \to L_W \mathcal{C}\) 使得所有 \(w \in W\) 变为等价，且在此性质下普遍。当 \(\mathcal{C}\) 是稳定 \(\infty\)-范畴且 \(W\) 由某类态射生成时，\(L_W \mathcal{C}\) 仍是稳定的。

\textbf{定理 7}（Verdier商 = Bousfield局部化）：设 \(\mathcal{C}\) 是稳定 \(\infty\)-范畴，\(\mathcal{S} \subseteq \mathcal{C}\) 是稳定子范畴。则 Verdier 商 \(\mathrm{Ho}(\mathcal{C})/\mathrm{Ho}(\mathcal{S})\)（在经典三角范畴意义下）等同于 Bousfield 局部化的同伦范畴：

\[\mathrm{Ho}(\mathcal{C}) / \mathrm{Ho}(\mathcal{S}) \cong \mathrm{Ho}(L_{W_{\mathcal{S}}} \mathcal{C})\]

其中 \(W_{\mathcal{S}}\) 是所有使得其同伦纤维属于 \(\mathcal{S}\) 的态射的集合（即 \(\mathcal{S}\)-等价类）。

\textbf{证明}：在稳定 \(\infty\)-范畴 \(\mathcal{C}\) 中，定义 \(W_{\mathcal{S}}\) 为态射 \(f\) 使得 \(\operatorname{fib}(f) \in \mathcal{S}\) 的全体。Bousfield 局部化 \(L = L_{W_{\mathcal{S}}}: \mathcal{C} \to L_{W_{\mathcal{S}}} \mathcal{C}\) 满足 \(L(X) \simeq 0\) 当且仅当 \(X \in \mathcal{S}\)。对每个 \(X \in \mathcal{C}\)，构造纤维序列 \(S(X) \to X \to L(X)\)，其中 \(S(X) \in \mathcal{S}\) 是 \(\mathcal{S}\)-局部化伴随即纤维，\(L(X)\) 是 \(\mathcal{S}\)-局部的（对 \(Y \in \mathcal{S}\) 有 \(\operatorname{Map}(Y, L(X)) \simeq 0\)）。取同伦范畴，此纤维序列变为三角范畴的标准 Verdier 商短正合列。函子 \(\mathrm{Ho}(L)\) 恰为商函子 \(\mathrm{Ho}(\mathcal{C}) \to \mathrm{Ho}(\mathcal{C})/\mathrm{Ho}(\mathcal{S})\)，且满足万有性质：任何零化 \(\mathrm{Ho}(\mathcal{S})\) 的三角函子唯一通过 \(\mathrm{Ho}(L)\) 分解。\(\blacksquare\)

\textbf{注记}：在 \(\infty\)-范畴框架下，Verdier 商的八面体公理和厚子范畴的条件由稳定性自然保证，而不需要像经典理论中那样逐个验证。这一视角使三角范畴理论更接近其同伦论起源------三角结构来自稳定 \(\infty\)-范畴的纤维-余纤维结构，而非人工添加的公理系统。

\hrulefill

\section{第130章：同调代数基础}\label{ux7b2c130ux7ae0ux540cux8c03ux4ee3ux6570ux57faux7840}

同调代数用链复形和同调群的语言统一处理代数拓扑、群上同调、层上同调等理论。

\subsection{70.1 链复形与同调}\label{ux94feux590dux5f62ux4e0eux540cux8c03}

\textbf{定义 70.1}（链复形）：一个（上）链复形 \(C^\bullet\) 是一列 Abel 群（或 \(R\)-模）\(C^n\) 和同态 \(d^n: C^n \to C^{n+1}\)（称为\textbf{微分}或\textbf{边缘算子}）满足 \(d^{n+1} \circ d^n = 0\)。即

\[\cdots \longrightarrow C^{n-1} \xrightarrow{d^{n-1}} C^n \xrightarrow{d^n} C^{n+1} \longrightarrow \cdots\]

（下链复形将箭头反向，\(d_n: C_n \to C_{n-1}\)。）

\textbf{定义 70.2}（上同调）：链复形 \(C^\bullet\) 的 \(n\) 次\textbf{上同调}为

\[H^n(C^\bullet) = \frac{\ker d^n}{\operatorname{im} d^{n-1}}\]

\(Z^n = \ker d^n\) 中的元素称为 \textbf{\(n\)-余圈}（cocycle），\(B^n = \operatorname{im} d^{n-1}\) 中的元素称为 \textbf{\(n\)-余边界}（coboundary）。

\textbf{定义 70.3}（链映射）：设 \(C^\bullet\) 和 \(D^\bullet\) 是链复形。一个\textbf{链映射} \(f: C^\bullet \to D^\bullet\) 是一族同态 \(f^n: C^n \to D^n\) 满足 \(d_D^n \circ f^n = f^{n+1} \circ d_C^n\)。链映射诱导上同调的同态 \(f^*: H^n(C^\bullet) \to H^n(D^\bullet)\)。

\textbf{定义 70.4}（链同伦）：两个链映射 \(f, g: C^\bullet \to D^\bullet\) 称为\textbf{链同伦的}，如果存在同态 \(h^n: C^n \to D^{n-1}\) 使得 \(f^n - g^n = d_D^{n-1} \circ h^n + h^{n+1} \circ d_C^n\)。链同伦的映射在上同调上诱导相同的同态。

\subsection{70.2 正合序列与蛇引理}\label{ux6b63ux5408ux5e8fux5217ux4e0eux86c7ux5f15ux7406}

\textbf{定义 70.5}（正合序列）：模同态的序列

\[\cdots \longrightarrow M_{i-1} \xrightarrow{f_{i-1}} M_i \xrightarrow{f_i} M_{i+1} \longrightarrow \cdots\]

称为在 \(M_i\) 处\textbf{正合的}，如果 \(\operatorname{im} f_{i-1} = \ker f_i\)。若序列在每处都正合，则称为\textbf{正合序列}。

\textbf{短正合序列}是形如 \(0 \to A \xrightarrow{i} B \xrightarrow{p} C \to 0\) 的正合序列（\(i\) 单，\(p\) 满，\(\ker p = \operatorname{im} i\)）。

\textbf{定理 70.1}（蛇引理）：设有 \(R\)-模的交换图（行正合）：

\[
\begin{CD}
A @>i>> B @>p>> C @>>> 0 \\
@VV{\alpha}V @VV{\beta}V @VV{\gamma}V \\
0 @>>> A' @>>{i'}> B' @>>{p'}> C'
\end{CD}
\]

则存在连接同态 \(\delta: \ker \gamma \to \operatorname{coker} \alpha\) 使得以下序列正合：

\[\ker \alpha \to \ker \beta \to \ker \gamma \xrightarrow{\delta} \operatorname{coker} \alpha \to \operatorname{coker} \beta \to \operatorname{coker} \gamma\]

\emph{证明}（构造 \(\delta\)）：对 \(c \in \ker \gamma\)，取 \(b \in B\) 使 \(p(b) = c\)。则 \(p'(\beta(b)) = \gamma(p(b)) = \gamma(c) = 0\)，故 \(\beta(b) \in \ker p' = \operatorname{im} i'\)。存在唯一 \(a' \in A'\) 使 \(i'(a') = \beta(b)\)。定义 \(\delta(c) = a' + \operatorname{im} \alpha \in \operatorname{coker} \alpha\)。验证 \(\delta\) 是良定义的且序列正合。∎

\textbf{定理 70.2}（五引理）：设有交换图（行正合）：

\[
\begin{CD}
A @>>> B @>>> C @>>> D @>>> E \\
@VV{\alpha}V @VV{\beta}V @VV{\gamma}V @VV{\delta}V @VV{\varepsilon}V \\
A' @>>> B' @>>> C' @>>> D' @>>> E'
\end{CD}
\]

\begin{enumerate}
\def\labelenumi{\arabic{enumi}.}
\tightlist
\item
  若 \(\alpha\) 满、\(\beta, \delta\) 单、\(\varepsilon\) 单，则 \(\gamma\) 单。
\item
  若 \(\alpha\) 满、\(\beta, \delta\) 满、\(\varepsilon\) 单，则 \(\gamma\) 满。
\item
  若 \(\alpha, \beta, \delta, \varepsilon\) 都是同构，则 \(\gamma\) 也是同构。
\end{enumerate}

\textbf{证明}：对 (1)，设 \(c \in \ker \gamma\)。则 \(\delta(d(c)) = d'(\gamma(c)) = 0\)，由 \(\delta\) 单得 \(d(c) = 0\)，故 \(c \in \ker d = \operatorname{im} c'\)，存在 \(b \in B\) 使 \(c'(b) = c\)。则 \(c''(\beta(b)) = \gamma(c'(b)) = \gamma(c) = 0\)，故 \(\beta(b) \in \ker c'' = \operatorname{im} b''\)，存在 \(a' \in A'\) 使 \(b''(a') = \beta(b)\)。由 \(\alpha\) 满，存在 \(a \in A\) 使 \(\alpha(a) = a'\)，则 \(\beta(b - b(a)) = \beta(b) - b''(a') = 0\)，由 \(\beta\) 单得 \(b = b(a)\)，故 \(c = c'(b) = 0\)。(2) 类似地对偶论证。(3) 由 (1) 和 (2) 直接得。\(\blacksquare\)

\subsection{70.3 投射模、内射模与平坦模}\label{ux6295ux5c04ux6a21ux5185ux5c04ux6a21ux4e0eux5e73ux5766ux6a21-1}

\textbf{定义 70.6}（投射模）：\(R\)-模 \(P\) 称为\textbf{投射模}，如果对任意满同态 \(f: M \to N\) 和同态 \(g: P \to N\)，存在提升 \(h: P \to M\) 使得 \(f \circ h = g\)。

\textbf{命题 70.3}：\(P\) 是投射模当且仅当 \(P\) 是某个自由模的直和项。特别地，自由模是投射模。

\textbf{证明}：（\(\Leftarrow\)）若 \(P\) 是自由模的直和项，设 \(P \oplus Q = F\)（\(F\) 自由）。对任意满同态 \(f: M \to N\) 和 \(g: P \to N\)，将 \(g\) 延拓为 \(\tilde{g}: F \to N\)（\(\tilde{g}(p,q) = g(p)\)）。自由模 \(F\) 是投射的，故存在提升 \(h: F \to M\) 使 \(f \circ h = \tilde{g}\)。限制 \(h|_P\) 即为 \(g\) 的提升。（\(\Rightarrow\)）取满同态 \(\pi: F \to P\)（\(F\) 为以 \(P\) 的元素为基的自由模），由投射性，存在 \(s: P \to F\) 使 \(\pi \circ s = \operatorname{id}_P\)。故 \(F \cong \operatorname{im}(s) \oplus \ker(\pi) \cong P \oplus \ker(\pi)\)，\(P\) 是 \(F\) 的直和项。\(\blacksquare\)

\textbf{定义 70.7}（内射模）：\(R\)-模 \(I\) 称为\textbf{内射模}，如果对任意单同态 \(f: A \to B\) 和同态 \(g: A \to I\)，存在扩张 \(h: B \to I\) 使得 \(h \circ f = g\)。

\textbf{命题 70.4}：每个 \(R\)-模都可以嵌入一个内射模。

\textbf{证明}：对 Abel 群情形（\(R=\mathbb{Z}\)），每个可除群是内射 \(\mathbb{Z}\)-模。任意 Abel 群 \(M\) 可嵌入可除群：取 \(M\) 的生成元集 \(\{m_i\}\)，构造满射 \(\bigoplus \mathbb{Z} \to M\)，取对偶 \(\operatorname{Hom}_{\mathbb{Z}}(-, \mathbb{Q}/\mathbb{Z})\) 得嵌入 \(M \hookrightarrow \prod \mathbb{Q}/\mathbb{Z}\)（后者为可除群，故内射）。对一般环 \(R\)，利用 \(\operatorname{Hom}_{\mathbb{Z}}(R, -)\) 将 \(\mathbb{Z}\)-模的内射分解转化为 \(R\)-模的内射分解：\(M \hookrightarrow \operatorname{Hom}_{\mathbb{Z}}(R, I)\)，其中 \(I\) 是包含 \(M\)（视为 \(\mathbb{Z}\)-模）的内射 \(\mathbb{Z}\)-模。\(\blacksquare\)

\textbf{定义 70.8}（平坦模）：\(R\)-模 \(M\) 称为\textbf{平坦模}，如果函子 \(M \otimes_R -\) 是正合的（即保持单射）。

\textbf{命题 70.5}：投射模是平坦模。自由模是平坦模。

\textbf{证明}：自由模是平坦的：\(R^{\oplus I} \otimes_R -\) 正合于 \(\bigoplus_I (-)\)，后者是正合函子。投射模 \(P\) 是自由模的直和项（命题 70.3），故 \(P \otimes_R -\) 是正合函子 \(\bigoplus (-)\) 的直和项，从而正合。\(\blacksquare\)

\subsection{70.4 导出函子：Ext 与 Tor}\label{ux5bfcux51faux51fdux5b50ext-ux4e0e-tor}

\textbf{定义 70.9}（投射分解）：\(M\) 的一个\textbf{投射分解}是一个正合序列

\[\cdots \to P_2 \to P_1 \to P_0 \to M \to 0\]

其中每个 \(P_i\) 是投射模。每个模都有投射分解（取自由分解即可）。

\textbf{定义 70.10}（Ext 函子）：对 \(R\)-模 \(M, N\)，取 \(M\) 的投射分解 \(P_\bullet \to M \to 0\)，应用函子 \(\operatorname{Hom}_R(-, N)\) 得链复形：

\[0 \to \operatorname{Hom}(P_0, N) \to \operatorname{Hom}(P_1, N) \to \operatorname{Hom}(P_2, N) \to \cdots\]

定义 \(\operatorname{Ext}_R^n(M, N) = H^n(\operatorname{Hom}_R(P_\bullet, N))\)。\(\operatorname{Ext}_R^n(M, N)\) 不依赖于投射分解的选取（在同构意义下）。

\textbf{定义 70.11}（Tor 函子）：对 \(R\)-模 \(M, N\)，取 \(M\) 的投射分解 \(P_\bullet \to M \to 0\)，应用函子 \(-\otimes_R N\) 得链复形：

\[\cdots \to P_2 \otimes N \to P_1 \otimes N \to P_0 \otimes N \to 0\]

定义 \(\operatorname{Tor}_n^R(M, N) = H_n(P_\bullet \otimes_R N)\)。

\textbf{命题 70.6}（Ext 和 Tor 的基本性质）： 1. \(\operatorname{Ext}_R^0(M, N) \cong \operatorname{Hom}_R(M, N)\)。 2. \(\operatorname{Tor}_0^R(M, N) \cong M \otimes_R N\)。 3. 短正合序列 \(0 \to A \to B \to C \to 0\) 诱导长正合序列： \(\cdots \to \operatorname{Ext}^n(C, N) \to \operatorname{Ext}^n(B, N) \to \operatorname{Ext}^n(A, N) \to \operatorname{Ext}^{n+1}(C, N) \to \cdots\) \(\cdots \to \operatorname{Tor}_{n+1}(C, N) \to \operatorname{Tor}_n(A, N) \to \operatorname{Tor}_n(B, N) \to \operatorname{Tor}_n(C, N) \to \cdots\) 4. \(M\) 投射 \(\Leftrightarrow\) \(\operatorname{Ext}_R^1(M, N) = 0\) 对所有 \(N\)。 5. \(M\) 平坦 \(\Leftrightarrow\) \(\operatorname{Tor}_1^R(M, N) = 0\) 对所有 \(N\)。

\textbf{证明}：（1）\(\operatorname{Ext}^0\)：\(P_0 \to M \to 0\) 正合，应用 \(\operatorname{Hom}(-,N)\) 得 \(0 \to \operatorname{Hom}(M,N) \to \operatorname{Hom}(P_0,N)\)，故 \(H^0 = \ker(\operatorname{Hom}(P_0,N) \to \operatorname{Hom}(P_1,N)) = \operatorname{Hom}(M,N)\)。（2）\(\operatorname{Tor}_0\)：同理，\(P_0 \otimes N \to M \otimes N \to 0\) 正合，故 \(H_0 = \operatorname{coker}(P_1 \otimes N \to P_0 \otimes N) = M \otimes N\)。（3）长正合序列由蛇引理应用于投射分解的短正合序列 \(0 \to P_\bullet(A) \to P_\bullet(B) \to P_\bullet(C) \to 0\) 导出。（4）若 \(M\) 投射，取 \(P_0 = M, P_{>0}=0\) 为投射分解，则 \(\operatorname{Hom}\) 复形集中于 \(0\) 次，\(\operatorname{Ext}^1 = 0\)。反之，若 \(\operatorname{Ext}^1(M,N)=0\) 对所有 \(N\)，取 \(N = \ker(P_0 \to M)\) 可证 \(M\) 投射。（5）类似（4），使用 \(\operatorname{Tor}\) 刻画平坦性。\(\blacksquare\)

\subsection{70.5 群上同调}\label{ux7fa4ux4e0aux540cux8c03}

\textbf{定义 70.12}（群上同调）：设 \(G\) 是群，\(M\) 是 \(G\)-模（即 \(M\) 是 Abel 群且有 \(G\) 的作用）。定义

\[H^n(G, M) = \operatorname{Ext}_{\mathbb{Z}G}^n(\mathbb{Z}, M)\]

其中 \(\mathbb{Z}G\) 是 \(G\) 的群环，\(\mathbb{Z}\) 是平凡 \(G\)-模（\(g \cdot m = m\)）。

\textbf{低次群上同调的解释}： - \(H^0(G, M) = M^G = \{m \in M : g \cdot m = m, \; \forall g \in G\}\)（不变元）。 - \(H^1(G, M)\)：交叉同态（\(f(gh) = f(g) + g \cdot f(h)\)）模去主交叉同态。当 \(G\) 作用平凡时，\(H^1(G, M) = \operatorname{Hom}(G, M)\)。 - \(H^2(G, M)\)：群扩张 \(0 \to M \to E \to G \to 1\) 的分类（\(M\) 视为 Abel 正规子群）。

\textbf{定理 70.7}（Hilbert 定理 90，加法形式）：设 \(E/F\) 是有限 Galois 扩张，\(G = \operatorname{Gal}(E/F)\)。则 \(H^1(G, E) = 0\)（\(E\) 为加法群，\(G\) 自然作用）。

\emph{证明}：由本原元素定理，\(E = F(\alpha)\)。设 \(f: G \to E\) 是交叉同态。由 Dedekind 线性无关性引理，存在 \(\beta \in E\) 使得 \(\sum_{\sigma \in G} f(\sigma)\sigma(\beta) \neq 0\)（若 \(f \neq 0\)）。令 \(\gamma = \sum f(\sigma)\sigma(\beta)\)，可验证 \(f(\tau) = \tau(\theta) - \theta\) 对某 \(\theta\) 成立，故 \(f\) 是主交叉同态。∎

\hrulefill

\hrulefill

\hrulefill

\hrulefill

\hrulefill

\hrulefill

\section{第131章：代数数论基础}\label{ux7b2c131ux7ae0ux4ee3ux6570ux6570ux8bbaux57faux7840}

代数数论研究代数数域（\(\mathbb{Q}\) 的有限扩张）的算术性质，是 Galois 理论和交换代数的交汇点。

\subsection{71.1 代数整数与整数环}\label{ux4ee3ux6570ux6574ux6570ux4e0eux6574ux6570ux73af}

\textbf{定义 71.1}（代数整数）：复数 \(\alpha \in \mathbb{C}\) 称为\textbf{代数整数}，如果存在首一整系数多项式 \(f(x) = x^n + a_{n-1}x^{n-1} + \cdots + a_0 \in \mathbb{Z}[x]\) 使 \(f(\alpha) = 0\)。

\textbf{命题 71.1}：\(\alpha\) 是代数整数当且仅当 \(\mathbb{Z}[\alpha]\) 是有限生成 \(\mathbb{Z}\)-模。

\textbf{证明}：（\(\Rightarrow\)）若 \(\alpha^n + a_{n-1}\alpha^{n-1} + \cdots + a_0 = 0\)（\(a_i \in \mathbb{Z}\)），则 \(\alpha^n\) 可表为 \(\{1, \alpha, \ldots, \alpha^{n-1}\}\) 的 \(\mathbb{Z}\)-线性组合。归纳地，所有 \(\alpha^k\)（\(k \geq n\)）均可由 \(\{1, \alpha, \ldots, \alpha^{n-1}\}\) 生成，故 \(\mathbb{Z}[\alpha]\) 由 \(\{1, \alpha, \ldots, \alpha^{n-1}\}\) 有限生成。（\(\Leftarrow\)）设 \(\mathbb{Z}[\alpha]\) 由 \(\{v_1, \ldots, v_m\}\) 生成。记 \(\alpha v_i = \sum_j a_{ij} v_j\)（\(a_{ij} \in \mathbb{Z}\)）。则 \(\alpha\) 满足矩阵 \(A = (a_{ij})\) 的特征多项式 \(\det(\alpha I - A) = 0\)，该多项式为首一整系数多项式，故 \(\alpha\) 为代数整数。\(\blacksquare\)

\textbf{定义 71.2}（数域的整数环）：设 \(K\) 是数域（\(\mathbb{Q}\) 的有限扩张）。\(K\) 中所有代数整数构成的集合记作 \(\mathcal{O}_K\)，称为 \(K\) 的\textbf{整数环}。

\textbf{定理 71.2}：\(\mathcal{O}_K\) 是 \(K\) 的子环，且 \(\mathcal{O}_K\) 是秩为 \([K : \mathbb{Q}]\) 的自由 \(\mathbb{Z}\)-模。即存在整基 \(\omega_1, \ldots, \omega_n\)（\(n = [K : \mathbb{Q}]\)）使得

\[\mathcal{O}_K = \mathbb{Z}\omega_1 \oplus \cdots \oplus \mathbb{Z}\omega_n\]

\textbf{证明}：\(\mathcal{O}_K\) 是环由代数整数的和与积仍为代数整数保证（整闭包）。作为 \(\mathbb{Z}\)-模，\(\mathcal{O}_K\) 无挠（从而是自由的）。秩的计算：取 \(K/\mathbb{Q}\) 的基 \(b_1, \ldots, b_n\)，乘以适当的整数可使它们属于 \(\mathcal{O}_K\)，故秩 \(\geq n\)。判别式非零性给出秩 \(\leq n\)（代数整数生成的 \(\mathbb{Z}\)-模离散地嵌入 \(\mathbb{R}^n \times \mathbb{C}^m\)，格秩不超过 \(n\)）。故秩恰为 \(n\)。\(\blacksquare\)

\subsection{71.2 Dedekind 整环与理想分解}\label{dedekind-ux6574ux73afux4e0eux7406ux60f3ux5206ux89e3}

\textbf{定义 71.3}（Dedekind 整环）：整环 \(R\) 称为 \textbf{Dedekind 整环}，如果 \(R\) 是 Noether 的、整闭的，且每个非零素理想都是极大理想（即 \(\dim R = 1\)）。

\textbf{定理 71.3}：数域 \(K\) 的整数环 \(\mathcal{O}_K\) 是 Dedekind 整环。

\emph{证明}：\(\mathcal{O}_K\) 是 Noether 的（因为它是有限生成 \(\mathbb{Z}\)-模，每个理想也是有限生成 \(\mathbb{Z}\)-模）。它是整闭的（代数整数的集合在整闭包下封闭）。\(\dim \mathcal{O}_K = 1\)（\(\mathcal{O}_K\) 是 \(\mathbb{Z}\) 的整扩张，\(\dim \mathbb{Z} = 1\)，整扩张不改变维数）。∎

\textbf{定理 71.4}（Dedekind 整环中理想的唯一分解）：在 Dedekind 整环 \(R\) 中，每个非零真理想 \(\mathfrak{a}\) 可以唯一地（不计顺序）分解为素理想的乘积：

\[\mathfrak{a} = \mathfrak{p}_1^{e_1} \mathfrak{p}_2^{e_2} \cdots \mathfrak{p}_r^{e_r}\]

其中 \(\mathfrak{p}_i\) 是互不相同的非零素理想，\(e_i \geq 1\)。

\textbf{证明}：由 Dedekind 整环定义，\(R\) 是 Noether 整闭的一维整环。对非零理想 \(\mathfrak{a}\)，由 Noether 性，\(\mathfrak{a}\) 有准素分解 \(\mathfrak{a} = \bigcap \mathfrak{q}_i\)，其中 \(\mathfrak{q}_i\) 为 \(\mathfrak{p}_i\)-准素理想。在 Dedekind 整环中，一维性保证非零素理想为极大理想，故所有准素理想恰为素理想的幂：\(\mathfrak{q}_i = \mathfrak{p}_i^{e_i}\)。整闭性确保局部化 \(R_{\mathfrak{p}_i}\) 为 DVR，其极大理想为主理想，理想运算为 \(\mathfrak{p}_i^{e_i} R_{\mathfrak{p}_i} = (\mathfrak{p}_i R_{\mathfrak{p}_i})^{e_i}\)。由中国剩余定理，\(\mathfrak{a} = \prod \mathfrak{p}_i^{e_i}\)。唯一性由素理想的不同性和局部化决定。\(\blacksquare\)

\textbf{推论 71.5}：\(\mathcal{O}_K\) 是 UFD 当且仅当 \(\mathcal{O}_K\) 是 PID（在 Dedekind 整环中，UFD 等价于 PID）。

\textbf{证明}：在 Dedekind 整环中，每个非零素理想是极大理想。若 \(\mathcal{O}_K\) 是 UFD，则每个非零素理想由不可约元生成，且非零素理想的高度为 1（一维），故每个理想可唯一分解为素理想的乘积。由素理想分解的唯一性，每个理想是主理想（由生成元分解决定），故 \(\mathcal{O}_K\) 是 PID。反之，PID 自然是 UFD。\(\blacksquare\)

\subsection{71.3 分歧理论}\label{ux5206ux6b67ux7406ux8bba}

\textbf{定义 71.4}（素理想的分歧）：设 \(L/K\) 是数域的扩张，\(\mathfrak{p}\) 是 \(\mathcal{O}_K\) 的非零素理想。则 \(\mathfrak{p}\mathcal{O}_L\) 可分解为 \(\mathcal{O}_L\) 中素理想的乘积：

\[\mathfrak{p}\mathcal{O}_L = \mathfrak{P}_1^{e_1} \mathfrak{P}_2^{e_2} \cdots \mathfrak{P}_g^{e_g}\]

其中 \(e_i = e(\mathfrak{P}_i/\mathfrak{p})\) 称为\textbf{分歧指数}，\(f_i = f(\mathfrak{P}_i/\mathfrak{p}) = [\mathcal{O}_L/\mathfrak{P}_i : \mathcal{O}_K/\mathfrak{p}]\) 称为\textbf{剩余类次数}。

\textbf{定理 71.6}（基本恒等式）：\(\sum_{i=1}^g e_i f_i = [L : K]\)。

\textbf{证明}：由中国剩余定理，\(\mathcal{O}_L / \mathfrak{p}\mathcal{O}_L \cong \bigoplus_{i=1}^g \mathcal{O}_L / \mathfrak{P}_i^{e_i}\)。取 \(\mathbb{F}_p = \mathcal{O}_K/\mathfrak{p}\) 上的维数：\(\dim_{\mathbb{F}_p} \mathcal{O}_L/\mathfrak{p}\mathcal{O}_L = [L : K]\)（因为 \(\mathcal{O}_L\) 是秩 \([L:K]\) 的自由 \(\mathcal{O}_K\)-模）。又 \(\dim_{\mathbb{F}_p} \mathcal{O}_L/\mathfrak{P}_i^{e_i} = e_i f_i\)（对 \(\mathcal{O}_L/\mathfrak{P}_i^{e_i}\) 按 \(\mathfrak{P}_i\)-进滤过计算长度）。故 \(\sum e_i f_i = [L : K]\)。\(\blacksquare\)

\textbf{定义 71.5}（分歧类型）： - \(\mathfrak{p}\) 称为\textbf{非分歧的}，如果所有 \(e_i = 1\)。 - \(\mathfrak{p}\) 称为\textbf{完全分裂的}，如果 \(e_i = f_i = 1\) 对所有 \(i\)（此时 \(g = [L : K]\)）。 - \(\mathfrak{p}\) 称为\textbf{惯性的}，如果 \(g = 1, e = 1, f = [L : K]\)。

\textbf{定理 71.7}（Dedekind 判别法定理）：设 \(L = K(\alpha)\)，\(\alpha \in \mathcal{O}_L\)，\(f(x)\) 是 \(\alpha\) 的最小多项式。若 \(\mathfrak{p}\) 不整除 \([\mathcal{O}_L : \mathcal{O}_K[\alpha]]\)（指数），则 \(\mathfrak{p}\mathcal{O}_L\) 的分解由 \(f(x) \bmod \mathfrak{p}\) 的因式分解决定：

若 \(f(x) \equiv \prod f_i(x)^{e_i} \pmod{\mathfrak{p}}\)（\(f_i\) 不可约 mod \(\mathfrak{p}\)），则 \(\mathfrak{p}\mathcal{O}_L = \prod \mathfrak{P}_i^{e_i}\)，且 \(f(\mathfrak{P}_i/\mathfrak{p}) = \deg f_i\)。

\textbf{证明}：条件 \(\mathfrak{p} \nmid [\mathcal{O}_L : \mathcal{O}_K[\alpha]]\) 保证 \(\mathcal{O}_L/\mathfrak{p}\mathcal{O}_L \cong \mathcal{O}_K[\alpha]/\mathfrak{p}\mathcal{O}_K[\alpha] \cong \mathbb{F}_p[x]/(\bar{f}(x))\)。后者按 \(\bar{f}(x) = \prod \bar{f}_i(x)^{e_i}\)（\(\bmod \mathfrak{p}\) 的因式分解）分解为中国剩余定理的直积 \(\bigoplus_i \mathbb{F}_p[x]/(\bar{f}_i(x)^{e_i})\)，每个因子对应 \(\mathcal{O}_L\) 中一个极大理想 \(\mathfrak{P}_i\)（\(\mathfrak{P}_i = \mathfrak{p}\mathcal{O}_L + f_i(\alpha)\mathcal{O}_L\)）。\(\blacksquare\)

\subsection{71.4 理想类群与类数}\label{ux7406ux60f3ux7c7bux7fa4ux4e0eux7c7bux6570}

\textbf{定义 71.6}（分式理想）：\(\mathcal{O}_K\) 的一个\textbf{分式理想}是 \(K\) 的一个有限生成 \(\mathcal{O}_K\)-子模 \(\mathfrak{a} \neq 0\)。等价地，存在 \(d \in \mathcal{O}_K \setminus \{0\}\) 使得 \(d\mathfrak{a} \subseteq \mathcal{O}_K\)（即 \(d\mathfrak{a}\) 是 \(\mathcal{O}_K\) 的整理想）。

分式理想在乘法下构成群 \(\mathcal{J}_K\)。主分式理想（\((\alpha) = \alpha\mathcal{O}_K\)，\(\alpha \in K^\times\)）构成子群 \(\mathcal{P}_K\)。

\textbf{定义 71.7}（理想类群）：\(K\) 的\textbf{理想类群}定义为

\[\operatorname{Cl}(K) = \mathcal{J}_K / \mathcal{P}_K\]

\(h_K = |\operatorname{Cl}(K)|\) 称为 \(K\) 的\textbf{类数}。

\textbf{定理 71.8}：理想类群 \(\operatorname{Cl}(K)\) 是有限群。

\textbf{证明}：将 \(\mathcal{O}_K\) 嵌入 \(\mathbb{R}^n \times \mathbb{C}^m \cong \mathbb{R}^{[K:\mathbb{Q}]}\)（\(n + 2m = [K : \mathbb{Q}]\)）作为全秩格 \(\Lambda\)。对每个理想 \(\mathfrak{a}\)，其范数 \(N(\mathfrak{a}) = |\mathcal{O}_K / \mathfrak{a}|\)。Minkowski 定理：对 \(X \subseteq \mathbb{R}^N\) 为中心对称凸体，若 \(\operatorname{vol}(X) > 2^N \operatorname{covol}(\Lambda)\) 则 \(X\) 含 \(\Lambda\) 的非零点。取 \(X\) 为适当的有界凸集，对每个理想类 \([C]\) 中的理想 \(\mathfrak{a}\)，存在 \(a \in \mathfrak{a} \setminus \{0\}\) 使得 \(|N(a)| \leq M_K \cdot N(\mathfrak{a})\)（\(M_K\) 为 Minkowski 常数仅依赖于 \(K\)）。由于 \(\mathfrak{a}^{-1}\) 在 \([C]^{-1}\) 中，\(N(\mathfrak{a}^{-1}) \leq M_K \cdot N(\mathfrak{a}^{-1})\) 给出 \(N(\mathfrak{a}) \leq M_K\)。范数有界的理想只有有限个，故理想类数有限。\(\blacksquare\)

\textbf{推论 71.9}：\(h_K = 1\) 当且仅当 \(\mathcal{O}_K\) 是 PID（等价于 UFD）。

\textbf{证明}：\(h_K = 1\) 意味着每个分式理想是主分式理想，即 \(\mathcal{J}_K = \mathcal{P}_K\)。特别地，每个整理想 \(\mathfrak{a} \subseteq \mathcal{O}_K\) 是主理想，故 \(\mathcal{O}_K\) 是 PID。反之，若 \(\mathcal{O}_K\) 是 PID，则每个分式理想可写为 \(\alpha \mathfrak{a}\)（\(\alpha \in K^\times\)，\(\mathfrak{a}\) 为整主理想），故为 \(\alpha \beta \mathcal{O}_K = (\alpha\beta)\)，即主分式理想，故 \(\operatorname{Cl}(K)\) 平凡，\(h_K = 1\)。\(\blacksquare\)

\textbf{定理 71.10}（Dirichlet 单位定理）：\(\mathcal{O}_K\) 的单位群 \(\mathcal{O}_K^\times\) 是有限生成 Abel 群，其秩为 \(r = r_1 + r_2 - 1\)，其中 \(r_1\) 是实嵌入的个数，\(2r_2\) 是复嵌入的个数。即

\[\mathcal{O}_K^\times \cong \mu(K) \times \mathbb{Z}^r\]

其中 \(\mu(K)\) 是 \(K\) 中单位根的有限群。

\textbf{证明}：考虑对数嵌入 \(L: \mathcal{O}_K^\times \to \mathbb{R}^{r_1+r_2}\)，\(L(u) = (\log |\sigma_1(u)|, \ldots, \log |\sigma_{r_1}(u)|, 2\log |\sigma_{r_1+1}(u)|, \ldots, 2\log |\sigma_{r_1+r_2}(u)|)\)。由乘积公式 \(|N(u)| = \prod |\sigma_i(u)| = 1\)，得 \(\sum L_i(u) = 0\)，故 \(L(\mathcal{O}_K^\times)\) 落在超平面 \(H = \{\sum x_i = 0\}\) 中。\(\ker L = \{u \in \mathcal{O}_K^\times : |\sigma_i(u)| = 1, \forall i\}\)，由 Kronecker 定理（所有共轭在单位圆上的代数整数是单位根），\(\ker L = \mu(K)\) 为有限群。Minkowski 的格点定理（或 Minkowski 第二定理）证明 \(L(\mathcal{O}_K^\times)\) 是 \(H \cong \mathbb{R}^{r_1+r_2-1}\) 中的全秩格，即 \(L(\mathcal{O}_K^\times) \cong \mathbb{Z}^{r_1+r_2-1}\)。故 \(\mathcal{O}_K^\times \cong \mu(K) \times \mathbb{Z}^r\)（\(r = r_1+r_2-1\)）。\(\blacksquare\)

\subsection{\texorpdfstring{71.5 \(\zeta\) 函数与类数公式}{71.5 \textbackslash zeta 函数与类数公式}}\label{zeta-ux51fdux6570ux4e0eux7c7bux6570ux516cux5f0f}

\textbf{定义 71.8}（Dedekind \(\zeta\) 函数）：数域 \(K\) 的 Dedekind \(\zeta\) 函数定义为

\[\zeta_K(s) = \sum_{\mathfrak{a} \neq 0} \frac{1}{N(\mathfrak{a})^s} = \prod_{\mathfrak{p}} \frac{1}{1 - N(\mathfrak{p})^{-s}}\]

其中 \(\mathfrak{a}\) 遍历 \(\mathcal{O}_K\) 的所有非零理想，\(\mathfrak{p}\) 遍历所有非零素理想，\(N(\mathfrak{a}) = |\mathcal{O}_K/\mathfrak{a}|\) 是理想范数。此级数对 \(\Re(s) > 1\) 绝对收敛。

\textbf{定理 71.11}（类数公式）：Dedekind \(\zeta\) 函数在 \(s=1\) 的留数为

\[\operatorname{Res}_{s=1} \zeta_K(s) = \frac{2^{r_1}(2\pi)^{r_2} h_K R_K}{w_K \sqrt{|d_K|}}\]

其中 \(h_K\) 是类数，\(R_K\) 是调整子，\(w_K\) 是单位根个数，\(d_K\) 是判别式。

\textbf{证明}：\(\zeta_K(s) = \sum_{\mathfrak{a}} N(\mathfrak{a})^{-s}\) 对非零理想求和。按理想类分解：\(\zeta_K(s) = \sum_{[C] \in \operatorname{Cl}(K)} \zeta_C(s)\)，其中 \(\zeta_C(s) = \sum_{\mathfrak{a} \in [C]} N(\mathfrak{a})^{-s}\)。固定理想类 \([C]\)，取代表元 \(\mathfrak{c}\)，则 \(\mathfrak{a} \in [C]\) 当且仅当 \(\mathfrak{a} = \alpha \mathfrak{c}^{-1}\)（\(\alpha \in \mathfrak{c}\) 非零）。因此 \(\zeta_C(s) = N(\mathfrak{c})^s \sum_{\alpha \in \mathfrak{c} \setminus \{0\}} |N(\alpha)|^{-s}\)。求和按 \(\alpha\) 的 Galois 嵌入的范数大小分类，利用格 \(\mathfrak{c}\) 在 \(\mathbb{R}^{r_1} \times \mathbb{C}^{r_2}\) 中的分布。当 \(s \to 1\) 时，\(\zeta_C(s) \sim \frac{2^{r_1}(2\pi)^{r_2} R_K}{w_K \sqrt{|d_K|}} \cdot \frac{1}{s-1}\)，且此主项不依赖于理想类 \([C]\)。对所有 \(h_K\) 个理想类求和，得 \(\operatorname{Res}_{s=1} \zeta_K(s) = \frac{2^{r_1}(2\pi)^{r_2} h_K R_K}{w_K \sqrt{|d_K|}}\)。\(\blacksquare\)

\hrulefill

\hrulefill

\hrulefill

\hrulefill

\hrulefill

\newpage

\chapter{05-代数/03-抽象代数二/05-Jordan代数与交换代数深化.md}\label{ux4ee3ux657003-ux62bdux8c61ux4ee3ux6570ux4e8c05-jordanux4ee3ux6570ux4e0eux4ea4ux6362ux4ee3ux6570ux6df1ux5316.md}

\section{第132章：Jordan 代数}\label{ux7b2c132ux7ae0jordan-ux4ee3ux6570}

Jordan 代数是满足交换律 \(x \cdot y = y \cdot x\) 与 \textbf{Jordan 恒等式} \((x \cdot y) \cdot x^2 = x \cdot (y \cdot x^2)\) 的非结合代数。它由 Pascual Jordan 于 1933 年在形式化量子力学的背景下引入，旨在以对称积 \(a \circ b = (ab + ba)/2\) 来代替结合代数中的通常乘积，从而描述量子可观测量的代数结构。Jordan 代数看似是结合代数的简单对称化，实则包含了丰富而独立的代数世界------特别是 27 维的例外 Albert 代数，它是分类理论中的唯一例外对象。

\subsection{74.1 基本定义与 Jordan 恒等式}\label{ux57faux672cux5b9aux4e49ux4e0e-jordan-ux6052ux7b49ux5f0f}

\textbf{定义 74.1}（Jordan 代数）：域 \(k\)（\(\operatorname{char} k \neq 2\)）上的 \textbf{Jordan 代数} 是一个 \(k\)-向量空间 \(J\)，配以双线性乘法 \(\cdot: J \times J \to J\)，满足：

\begin{enumerate}
\def\labelenumi{\arabic{enumi}.}
\tightlist
\item
  \(x \cdot y = y \cdot x\)（交换律）；
\item
  \((x \cdot y) \cdot x^2 = x \cdot (y \cdot x^2)\)（Jordan 恒等式）。
\end{enumerate}

其中 \(x^2 = x \cdot x\)。注意 Jordan 恒等式弱于结合律------Jordan 代数不要求乘法满足结合律。Jordan 代数也可等价地以线性化的 Jordan 恒等式定义：

\[(x \cdot y) \cdot (z \cdot w) + (x \cdot w) \cdot (z \cdot y) + (x \cdot z) \cdot (y \cdot w) = x \cdot [y \cdot (z \cdot w)] + x \cdot [w \cdot (z \cdot y)] + x \cdot [z \cdot (y \cdot w)]\]

以及其更常见的极化形式：

\[(x, y, z^2) = 2(x \cdot z, y, z)\]

其中 \((x, y, z) = (x \cdot y) \cdot z - x \cdot (y \cdot z)\) 为结合子（associator）。

\textbf{命题 74.1}：任何 Jordan 代数都是\textbf{幂结合的}（power-associative），即 \(x^n = x \cdot x^{n-1}\) 的定义是无歧义的（与括弧位置无关）。

\textbf{证明}：由 Jordan 恒等式 \((x \cdot y) \cdot x^2 = x \cdot (y \cdot x^2)\)，取 \(y = x^{m-1}\) 得 \(x^m \cdot x^2 = x \cdot (x^{m-1} \cdot x^2)\)。假设归纳假设 \(x^k \cdot x^\ell = x^{k+\ell}\) 对所有 \(k+\ell < m+2\) 成立，则 \(x^m \cdot x^2 = x \cdot x^{m+1} = x^{m+2}\)。更一般地，对任意 \(m, n \geq 1\)，利用极化 Jordan 恒等式 \((x^m, x, x^n) = 2(x^{m+1}, x, x^n)\) 可归纳证明 \(x^m \cdot x^n = x^{m+n}\)。这意味着 \(x^n\) 的定义与括弧方式无关，即 Jordan 代数为幂结合。 \(\blacksquare\)

\subsection{74.2 特例 Jordan 代数与例外 Jordan 代数}\label{ux7279ux4f8b-jordan-ux4ee3ux6570ux4e0eux4f8bux5916-jordan-ux4ee3ux6570}

\textbf{定义 74.2}（特例 Jordan 代数，Special Jordan Algebra）：设 \(A\) 是结合代数。在向量空间 \(A\) 上定义 \textbf{Jordan 积}：

\[a \circ b = \frac{ab + ba}{2}\]

则 \((A, \circ)\) 是 Jordan 代数，记作 \(A^+\)。同构于某个 \(A^+\) 的子代数的 Jordan 代数称为\textbf{特例 Jordan 代数}（special Jordan algebra）。

\textbf{定义 74.3}（例外 Jordan 代数，Exceptional Jordan Algebra）：一个 \textbf{例外 Jordan 代数} 是不与任何结合代数的 \(A^+\) 的子代数同构的 Jordan 代数。

\textbf{定理 74.2}（Albert 代数）：设 \(\mathbb{O}\) 为 Cayley 八元数代数。\(3 \times 3\) 的 Hermite 矩阵

\[H_3(\mathbb{O}) = \left\{\begin{pmatrix} a & z & \bar{y} \\ \bar{z} & b & x \\ y & \bar{x} & c \end{pmatrix} : a, b, c \in \mathbb{R},\; x, y, z \in \mathbb{O}\right\}\]

在 Jordan 积下构成 27 维 Jordan 代数。\(H_3(\mathbb{O})\) 是例外 Jordan 代数，称为 \textbf{Albert 代数}。

\textbf{证明}：先验证 \(H_3(\mathbb{O})\) 在 Jordan 积 \(X \circ Y = \frac{1}{2}(XY + YX)\) 下封闭。八元数 \(\mathbb{O}\) 非结合但满足交错律，故 \(H_3(\mathbb{O})\) 中矩阵乘法虽非结合，但 Jordan 恒等式 \((X \circ Y) \circ X^2 = X \circ (Y \circ X^2)\) 仍成立。维数计算：\(a, b, c \in \mathbb{R}\) 贡献 3 维，\(x, y, z \in \mathbb{O}\) 各贡献 8 维，\(3 \times 8 = 24\) 维，合计 \(3 + 24 = 27\) 维。例外性：若 \(H_3(\mathbb{O})\) 为特例的，则存在结合代数 \(A\) 和单射 \(H_3(\mathbb{O}) \hookrightarrow A^+\)。Albert 代数中最多 3 个正交原始幂等元，而任何特例 Jordan 代数中正交幂等元对应结合代数中的正交幂等元，导出 \(\mathbb{O}\) 结合性的矛盾。Glennie 于 1963 年发现 s-恒等式进一步证实了例外性。 \(\blacksquare\)

\subsection{74.3 形式实 Jordan 代数与分类定理}\label{ux5f62ux5f0fux5b9e-jordan-ux4ee3ux6570ux4e0eux5206ux7c7bux5b9aux7406}

\textbf{定义 74.4}（形式实 Jordan 代数）：Jordan 代数 \(J\) 称为\textbf{形式实的}，如果对任意有限个 \(x_i \in J\)，\(\sum x_i^2 = 0\) 蕴含所有 \(x_i = 0\)。形式实条件相当于正定性条件，保证了代数具有 Euclid 型的内积结构。

\textbf{定理 74.3}（Jordan-von Neumann-Wigner 分类定理，1934 年）：有限维形式实 Jordan 代数恰好是以下四类的直和：

\begin{enumerate}
\def\labelenumi{\arabic{enumi}.}
\tightlist
\item
  \textbf{\(H_n(\mathbb{R})\)}：\(n \times n\) 实对称矩阵，维数为 \(n(n+1)/2\)；
\item
  \textbf{\(H_n(\mathbb{C})\)}：\(n \times n\) 复 Hermite 矩阵，维数为 \(n^2\)；
\item
  \textbf{\(H_n(\mathbb{H})\)}：\(n \times n\) 四元数 Hermite 矩阵，维数为 \(n(2n-1)\)；
\item
  \textbf{旋因子}（spin factors）\(\mathbb{R} \oplus \mathbb{R}^n\)：由内积空间 \((V, q)\) 构造，乘法定义为 \((\lambda, v) \cdot (\mu, w) = (\lambda\mu + q(v, w), \lambda w + \mu v)\)，维数为 \(n+1\)；
\item
  \textbf{Albert 代数} \(H_3(\mathbb{O})\)：\(3 \times 3\) 八元数 Hermite 矩阵，维数为 27。
\end{enumerate}

\textbf{证明}：设 \(J\) 为有限维形式实 Jordan 代数。形式实性保证 \(J\) 存在正定迹型 \(\text{tr}(x \cdot y)\)，从而 \(J\) 半单且有单位元 \(1\)。将 \(1\) 分解为正交原始幂等元之和 \(1 = \sum_{i=1}^r e_i\)。Peirce 分解给出 \(J = \bigoplus_{i \leq j} J_{ij}\)，其中 \(J_{ii} = \mathbb{R} e_i\)（一维）。分析连接数 \(a_{ij} = \dim J_{ij}\)：若某个 \(a_{ij} = 0\) 则 \(J\) 可分解为直和，故可设 \(J\) 既约且 \(a_{ij} > 0\)。若 \(r = 1\) 则 \(J = \mathbb{R}\)。若 \(r = 2\) 则 \(J\) 为旋因子。若 \(r \geq 3\)，则 \(J_{ij}\) 上可定义乘法 \((x, y) \mapsto x \cdot e_i \cdot y\) 使其成为赋范可除代数，由 Hurwitz 定理必同构于 \(\mathbb{R}, \mathbb{C}, \mathbb{H}\) 或 \(\mathbb{O}\)。当 \(a_{ij} = 1, 2, 4\) 时分别对应 \(H_n(\mathbb{R}), H_n(\mathbb{C}), H_n(\mathbb{H})\)；当 \(a_{ij} = 8\) 时唯有 \(r = 3\)（\(H_3(\mathbb{O})\)），因 \(r \geq 4\) 时八元数组合同构性失效。 \(\blacksquare\)

\subsection{74.4 与对称空间和 Lie 代数的联系}\label{ux4e0eux5bf9ux79f0ux7a7aux95f4ux548c-lie-ux4ee3ux6570ux7684ux8054ux7cfb}

Jordan 代数通过 \textbf{Kantor-Koecher-Tits 构造}（KKT 构造）与 Lie 代数建立了深刻联系。

\textbf{定理 74.4}（KKT 构造）：设 \(J\) 是 Jordan 代数。定义向量空间

\[T(J) = J \oplus \operatorname{Der}(J) \oplus J\]

其中 \(\operatorname{Der}(J)\) 为 \(J\) 的导子代数。可在 \(T(J)\) 上定义一个 3-次 Lie 代数结构（即 \(\mathbb{Z}\)-分次 Lie 代数 \(\mathfrak{g} = \mathfrak{g}_{-1} \oplus \mathfrak{g}_0 \oplus \mathfrak{g}_1\)），使得 \(\mathfrak{g}_{-1} \cong \mathfrak{g}_1 \cong J\)，\(\mathfrak{g}_0 \cong \operatorname{Str}(J)\)（\(J\) 的结构代数）。

具体 Lie 括号由以下公式给出：对 \(x, y \in \mathfrak{g}_{-1} \cong J\)，\([x, y] = -2L(x, y) \in \operatorname{Der}(J)\)，其中 \(L(x, y)z = (x \cdot y) \cdot z - x \cdot (y \cdot z)\)。

\textbf{证明}：验证分次 Jacobi 恒等式。\(\mathfrak{g}_{-1} \cong J\) 上 Lie 括号定义为 \([x, y] = -2L(x, y)\)，其中 \(L(x, y) \in \operatorname{Der}(J)\)。\(\mathfrak{g}_0 = \operatorname{Str}(J)\) 对 \(\mathfrak{g}_{-1}\) 的作用由导子自然给出。\(\mathfrak{g}_1\) 上括号由对称性定义。对 \([x, [y, z]] + [y, [z, x]] + [z, [x, y]]\) 在 \(\mathfrak{g}_{-1}\) 上，利用 Jordan 恒等式 \((x \cdot y) \cdot x^2 = x \cdot (y \cdot x^2)\) 可验证其为零。\(\mathfrak{g}_0\) 与 \(\mathfrak{g}_{\pm 1}\) 的括号由 \(\operatorname{Der}(J)\) 的 Lie 括号结构定义，Jacobi 恒等式由导子性质保证。\(\mathfrak{g}_{\pm 1}\) 与 \(\mathfrak{g}_0\) 的括号由结构代数的伴随表示定义。此分次 Lie 代数 \(T(J)\) 称为 \(J\) 的 Tits-Kantor-Koecher 代数。 \(\blacksquare\)

\emph{意义}：此构造将任意 Jordan 代数''线性化''为 3-次 Lie 代数，从而将 Jordan 代数的分类问题归入 Lie 代数的分类框架。反之，Lie 代数的最高权向量表示理论可用来构造 Jordan 代数。

\textbf{定理 74.5}（Koecher-Tits 构造）：形式实 Jordan 代数 \(J\) 的对称锥 \(\Omega = \{x^2 : x \in J\}\) 的复化是 Hermite 对称空间。这给出了\textbf{有界对称域}（bounded symmetric domains）与形式实 Jordan 代数的一一对应：

\textbf{证明}：形式实 Jordan 代数 \(J\) 的对称锥 \(\Omega = \{x^2 : x \in J\}\) 是 \(J\) 中的齐性自对偶锥。\(\Omega\) 的管状域 \(T_\Omega = J + i\Omega\) 在 \(J_{\mathbb{C}}\) 中为有界对称域。KKT 构造给出 \(T(J)\) 的 Lie 代数 \(\mathfrak{g}\)，其 Cartan 分解 \(\mathfrak{g} = \mathfrak{k} \oplus \mathfrak{p}\) 中 \(\mathfrak{p} \cong J\)。\(\Omega\) 的 Riemann 对称空间为 \(G/K\)，其中 \(G\) 为 \(\mathfrak{g}\) 对应的 Lie 群，\(K\) 为 \(\mathfrak{k}\) 对应的紧子群。此对应是可逆的：给定不可约有界对称域 \(D\)，其 Shilov 边界上的几何结构唯一确定一个形式实 Jordan 代数。分类对应：\(H_n(\mathbb{R}) \leftrightarrow\) III 型（Siegel），\(H_n(\mathbb{C}) \leftrightarrow\) I 型（经典），\(H_n(\mathbb{H}) \leftrightarrow\) II 型，旋因子 \(\leftrightarrow\) IV 型，Albert 代数 \(\leftrightarrow\) 例外型（\(E_7\) 相关）。 \(\blacksquare\)

\begin{itemize}
\tightlist
\item
  \(\operatorname{Sym}_n(\mathbb{R})^+ \leftrightarrow\) Siegel 上半空间 \(\mathbb{H}_n\)（III 型有界对称域）；
\item
  \(\operatorname{Herm}_n(\mathbb{C})^+ \leftrightarrow\) I 型有界对称域（Cartan 的四种古典域中的第一种）；
\item
  旋因子 \(\leftrightarrow\) IV 型有界对称域；
\item
  Albert 代数 \(\leftrightarrow\) 27 维的例外有界对称域（E.VII 型）。
\end{itemize}

这一对应是 Elie Cartan 对称空间分类的 Jordan 代数诠释，也是非交换调和分析和表示论的重要工具。

\subsection{Jordan 代数的现代发展与例外现象}\label{jordan-ux4ee3ux6570ux7684ux73b0ux4ee3ux53d1ux5c55ux4e0eux4f8bux5916ux73b0ux8c61}

Jordan 代数理论在 20 世纪下半叶经历了深刻的变革，其核心主题是理解例外现象（exceptional phenomena）的代数根源。\textbf{Kantor 对}（Kantor pair）是 Jordan 代数的推广，完全由一对满足特定 Jordan 型恒等式的向量空间组成，它自然地包含了有限维单 Lie 代数的所有 3-分次结构。\textbf{结构群}（structure group）\(\operatorname{Str}(J)\) 是 \(J\) 的线性自同构群，满足 \(\operatorname{Str}(J) = \{g \in GL(J) : \exists g^* \in GL(J), g(x \cdot y) = gx \cdot g^*y\}\)，其中 \(g^*\) 由 \(g\) 唯一确定。结构代数 \(\operatorname{Str}(J)\) 的 Lie 代数包含内导子代数 \(\operatorname{Innder}(J) = \{L(x, y) : x, y \in J\}\) 作为理想，商为 \(\operatorname{Str}(J)/\operatorname{Innder}(J)\) 同构于 \(J\) 的中心的 Lie 代数。\textbf{McCrimmon 的构造}（1966）通过 Freudenthal 三重系统（Freudenthal triple system）将 Albert 代数嵌入 56 维 \(E_7\) 的 Lie 代数中，揭示了例外 Lie 代数与例外 Jordan 代数之间的一一对应：\(F_4 = \operatorname{Der}(H_3(\mathbb{O}))\)，\(E_6 = \operatorname{Str}_0(H_3(\mathbb{O}))\)（保持单位元的结构群），\(E_7 = \operatorname{Conf}(H_3(\mathbb{O}))\)（共形群），\(E_8\) 则与 \(H_3(\mathbb{O})\) 的''魔方''构造（Freudenthal-Tits 魔方）相关。\textbf{魔方构造}（magic square）是 Tits（1966）和 Freudenthal（1964）独立发现的 4x4 表格，其条目是两个赋范可除代数（\(\mathbb{R}, \mathbb{C}, \mathbb{H}, \mathbb{O}\)）的张量积对应的 Lie 代数，行列相交处正好给出所有例外 Lie 代数，其中 Albert 代数对应 \(F_4\)。这一构造将 Hurwitz 定理（可除代数分类）与 Killing-Cartan 分类（单 Lie 代数分类）完美统一。Jordan 代数在\textbf{量子力学}的公理化中扮演了基本角色------Jordan、von Neumann 和 Wigner 最初引入形式实 Jordan 代数正是为了用代数方法刻画量子可观测量的集合，而 Jordan-Banach 代数（JB-代数）和 JBW-代数（Jordan-Banach 代数中弱闭的）提供了量子力学代数基础的严格数学框架。在\textbf{非交换代数几何}中，Jordan 代数通过 Jordan 簇（Jordan scheme）和 Jordan 对（Jordan pair）的范畴化，为研究非结合的几何对象提供了平台。

\hrulefill

\hrulefill

\section{第133章：交换代数深化}\label{ux7b2c133ux7ae0ux4ea4ux6362ux4ee3ux6570ux6df1ux5316}

交换代数深化在卷十一交换代数的基础上，进一步探讨了 Cohen-Macaulay 环、Gorenstein 环、局部上同调和 Grothendieck 对偶等高级主题。这些概念在现代代数几何中扮演着核心角色------它们将代数簇的几何性质（如光滑性、奇点解消、对偶性）转化为交换环的代数性质。Cohen-Macaulay 性质是介于正则性和一般 Noether 性之间的关键中间概念：正则环必为 Cohen-Macaulay 环，但反之不然------环面簇的齐次坐标环是 Cohen-Macaulay 的却未必正则，这一现象在环面几何中具有基本重要性。Auslander-Buchsbaum 公式（\(\operatorname{pd}_R(M) + \operatorname{depth}(M) = \operatorname{depth}(R)\)）将投射维数与深度联系起来，是交换代数中最优美的等式之一，它精确量化了模的''同调复杂度''与''正则序列长度''之间的对偶关系。

Gorenstein 环是 Cohen-Macaulay 环中具有有限内射维数的一类，它们在 Serre 对偶和 Grothendieck 对偶中扮演着对偶化复形的角色。Grothendieck 局部对偶定理断言：对于 \(d\) 维 Gorenstein 局部环 \(R\)，\(H^i_{\mathfrak{m}}(M) \cong \operatorname{Hom}_R(\operatorname{Ext}^{d-i}_R(M, R), E)\)，其中 \(E\) 是剩余域的内射包络。这一对偶定理是 Serre 对偶的局部版本，它将局部上同调（一种''局部''同调）与 Ext 函子（一种''整体''上同调）联系起来，构成了代数几何中对偶理论的代数基础。局部上同调 \(H^i_{\mathfrak{m}}(M)\) 本身量度了模 \(M\) 沿极大理想 \(\mathfrak{m}\) 的''奇异性''------它在代数几何中对应于凝聚层沿闭子簇的局部上同调，是研究奇点理论、D-模和比较定理的基本工具。本章将系统介绍 Cohen-Macaulay 环、Gorenstein 环、局部上同调和 Grothendieck 对偶，为代数几何（卷十五）中的对偶理论和平坦下降做好充分准备。

\textbf{定义 79.1}：Noether 局部环 \((R, \mathfrak{m})\) 的 \textbf{深度} \(\operatorname{depth}(R)\) 是 \(R\)-正则序列的最大长度。总有 \(\operatorname{depth}(R) \leq \dim(R)\)。若等号成立，\(R\) 称为 \textbf{Cohen-Macaulay 环}。若 \(R\) 是 Cohen-Macaulay 且内射维数 \(\operatorname{inj.dim}_R(R) < \infty\)，则 \(R\) 称为 \textbf{Gorenstein 环}。

\textbf{定理 79.1}（Auslander-Buchsbaum 公式）：对有限生成 \(R\)-模 \(M\)（\(R\) 正则局部环，\(\operatorname{pd}_R(M) < \infty\)），有 \(\operatorname{pd}_R(M) + \operatorname{depth}(M) = \operatorname{depth}(R)\)，其中 \(\operatorname{pd}\) 是投射维数。

\textbf{证明}：设 \(R\) 为正则局部环，\(\dim R = d\)。取 \(R\) 的正则参数系 \(x_1, \ldots, x_d\)，构造 Koszul 复形 \(K_\bullet(x_1, \ldots, x_d)\)，其为 \(R\) 的极小自由消解，长度为 \(d\)。对有限生成 \(R\)-模 \(M\) 满足 \(\operatorname{pd}_R(M) < \infty\)，取极小自由消解 \(0 \to F_p \to \cdots \to F_0 \to M \to 0\)。应用 Auslander-Buchsbaum 定理：\(\operatorname{depth}(M) + \operatorname{pd}_R(M) = \operatorname{depth}(R) = d\)。故 \(\operatorname{pd}_R(M) = d - \operatorname{depth}(M)\)。利用 Koszul 复形上同调与深度等价刻画，\(\operatorname{depth}(M) = \min\{i : H^i(K_\bullet \otimes M) \neq 0\}\)。代入 \(\operatorname{pd}_R(M) = \max\{i : \operatorname{Tor}_i^R(M, k) \neq 0\}\) 并利用 Koszul 复形的 Tor 计算得证。 \(\blacksquare\)

\textbf{定理 79.2}（Serre 正则性判别）：Noether 局部环 \(R\) 是正则环当且仅当 \(R\) 的全局维数有限，当且仅当 \(\dim_k(\mathfrak{m}/\mathfrak{m}^2) = \dim(R)\)。这是代数几何中光滑性的代数对应：\(\operatorname{Spec}(R)\) 光滑当且仅当 \(R\) 为正则局部环。Cohen-Macaulay 性质介于正则性与一般 Noether 环之间------所有正则环是 Cohen-Macaulay，所有 Cohen-Macaulay 环在适当维数限制下满足 Serre 条件 \((S_k)\)。

\textbf{证明}：设 \(d = \dim_k(\mathfrak{m}/\mathfrak{m}^2)\)。取 \(x_1, \ldots, x_d \in \mathfrak{m}\) 使其在 \(\mathfrak{m}/\mathfrak{m}^2\) 中的像为 \(k\) 的基。Koszul 复形 \(K_\bullet(x_1, \ldots, x_d)\) 为 \(k\) 的极小自由消解，故 \(\operatorname{pd}_R(k) = d\)。由 Serre 定理，\(\operatorname{gldim}(R) = \operatorname{pd}_R(k)\)，故 \(\operatorname{gldim}(R) = d\)。另一方面，\(\dim(R) \leq d\) 恒成立。\(R\) 正则当且仅当 \(\dim(R) = \operatorname{gldim}(R)\)，这等价于 \(\dim(R) = \dim_k(\mathfrak{m}/\mathfrak{m}^2)\)。光滑性：\(\operatorname{Spec}(R)\) 在闭点光滑当且仅当切空间维数等于环维数。 \(\blacksquare\)

\subsection{79.2 局部上同调与 Grothendieck 对偶}\label{ux5c40ux90e8ux4e0aux540cux8c03ux4e0e-grothendieck-ux5bf9ux5076}

\textbf{定义 79.2}：对环 \(R\) 上模 \(M\) 和理想 \(I \subseteq R\)，第 \(i\) 个 \textbf{\(I\)-局部上同调} \(\Gamma_I(M)\) 定义为 \(M\) 中所有被 \(I\) 的某次幂零化的元素的子模，高阶导出函子 \(H^i_I(M) = R^i \Gamma_I(M)\)。

\textbf{定理 79.3}（Grothendieck 局部对偶）：设 \(R\) 为 \(d\) 维 Gorenstein 局部环，\(\mathfrak{m}\) 为极大理想。对有限生成 \(R\)-模 \(M\)，有自然同构 \(H^i_{\mathfrak{m}}(M) \cong \operatorname{Hom}_R(\operatorname{Ext}^{d-i}_R(M, R), E)\) 其中 \(E\) 是 \(R\) 的内射包络。

\textbf{证明}：设 \(R\) 为 \(d\) 维 Gorenstein 局部环，\(E = E_R(k)\) 为剩余域 \(k\) 的内射包。定义 Matlis 对偶 \(D(M) = \operatorname{Hom}_R(M, E)\)。Gorenstein 条件给出 \(H^d_{\mathfrak{m}}(R) \cong E\)。利用局部对偶定理：\(\operatorname{RHom}_R(M, \omega_R^\bullet[d]) \cong \operatorname{RHom}_R(\operatorname{R\Gamma}_{\mathfrak{m}}(M), E)\)，其中 \(\omega_R^\bullet\) 为对偶化复形。对 Gorenstein 环，\(\omega_R^\bullet \cong R[d]\)，故 \(H^i_{\mathfrak{m}}(M) \cong \operatorname{Ext}_R^{d-i}(M, R)^\vee\)。此同构由谱序列 \(E_2^{p,q} = \operatorname{Ext}_R^p(\operatorname{Tor}_q^R(k, M), E) \Rightarrow H^{p+q}_{\mathfrak{m}}(M)\) 并结合 Koszul 复形计算得出。 \(\blacksquare\)

\hrulefill

\emph{本卷在 V8 的基础上完成了抽象代数核心领域的深化：Galois 理论建立了域扩张与群论的桥梁，彻底解决了方程根式可解性问题；交换代数为代数几何和代数数论提供了理想论、局部化、Noether 环等基础工具；表示论用特征标理论揭示了有限群的结构；同调代数引入 Ext/Tor 函子和导出函子，统一了代数拓扑和代数几何的上同调理论；代数数论建立了 Dedekind 整环中理想的唯一分解定理和理想类群理论。为后续代数几何（V15）、代数拓扑（V14）和数论 II（V16）的学习奠定了基础。}

\emph{卷十二：抽象代数 II终。}

\emph{卷十二：抽象代数 II终。}

\emph{卷十二：抽象代数 II终。}

\emph{卷十四：抽象代数（二）终。} \hrulefill

\subsection{PI代数理论}\label{piux4ee3ux6570ux7406ux8bba}

PI 代数（Polynomial Identity Algebra，多项式恒等代数）是非交换代数中一类重要的代数，它们虽不交换，但满足某个非交换多项式恒等式，从而在交换代数与完全非交换代数之间占据了中间地带。PI 理论由 Dehn（1922）和 Wagner（1937）早期探索，经 Kaplansky（1948 年证明本原 PI 代数是中心单代数）、Amitsur 和 Levitzki（1950 年证明矩阵代数 \(M_n(k)\) 满足标准恒等式 \(S_{2n} = 0\)）等人的工作，逐渐发展为非交换代数中最系统、最深刻的理论之一。PI 代数的核心洞察是：多项式恒等式对代数结构施加了强烈的限制，使得 PI 代数虽然非交换，却与交换代数共享许多良好性质，如有限维表示的分类、理想的链条件等。

Amitsur-Levitzki 定理是 PI 理论中最著名的结果：\(n \times n\) 矩阵代数 \(M_n(k)\) 满足标准恒等式 \(S_{2n}(x_1, \ldots, x_{2n}) = \sum_{\sigma \in \mathfrak{S}_{2n}} \operatorname{sgn}(\sigma) x_{\sigma(1)} \cdots x_{\sigma(2n)} = 0\)，且 \(2n\) 是最小的可能次数。这一结果惊人地揭示了：任何一个 \(n\) 维向量空间上的线性变换，如果取 \(2n\) 个，则它们的交错乘积之和恒为零。Kaplansky 定理进一步将 PI 代数与经典结构联系起来：在特征零域上，本原 PI 代数必为中心单代数（即 \(M_n(D)\)，\(D\) 为有限维除环），这为 PI 代数的结构理论奠定了基础。

Specht 问题（1907 年提出，1987 年由 Kemer 解决）是 PI 理论中影响最深远的猜想：在特征零域上，每个结合代数的 T-理想（恒等式理想）是否都是有限基的------即所有多项式恒等式能否由有限多个恒等式生成？Kemer 的肯定回答依赖于他建立的 PI 代数表示论与 \(\mathbb{Z}_2\)-分次代数之间的深刻联系（Grassmann 包络），以及有限维 \(\mathbb{Z}_2\)-分次代数的 T-理想有限基定理。Giambruno-Zaicev 定理（1999）进一步证明了 PI 指数 \(\exp(A) = \lim_{n \to \infty} \sqrt[n]{c_n(A)}\) 存在且为整数，其中 \(c_n(A)\) 是 \(n\) 次多线性多项式的 codimension 序列。本章将系统介绍 PI 代数的定义、基本例子、Amitsur-Levitzki 定理、Kaplansky 定理、T-理想、Specht 问题的 Kemer 解，以及 codimension 序列的渐近理论。

\textbf{定义 1}（自由结合代数）：设 \(k\) 是域。令 \(k\langle x_1, \ldots, x_n \rangle\) 是非交换变元 \(x_1, \ldots, x_n\) 上的自由结合代数（张量代数）。其元素是非交换多项式的有限形式和。

\textbf{定义 2}（PI代数）：结合 \(k\)-代数 \(A\) 称为\textbf{多项式恒等代数}（PI代数），若存在非零多项式

\[f(x_1, \ldots, x_n) \in k\langle x_1, \ldots, x_n \rangle\]

使得对所有 \(a_1, \ldots, a_n \in A\)，有 \(f(a_1, \ldots, a_n) = 0\)。这样的 \(f\) 称为 \(A\) 的\textbf{多项式恒等式}。恒等式的最小次数称为 \(A\) 的 \textbf{PI-次数}。

若 \(A\) 不是 PI 代数，则称 \(A\) 为 \textbf{PI-平凡}的。

\subsection{基本例子}\label{ux57faux672cux4f8bux5b50}

PI 环的定义虽然形式简洁，但判断一个给定环是否满足多项式恒等式并非总是显然的。以下几个基本例子展示了 PI 环的''典型来源''------交换代数（最平凡的 PI 环）、有限维代数（由线性相关性保证）和 Azumaya 代数（由矩阵表示控制）。

\textbf{定理 1}（交换代数）：每个交换 \(k\)-代数 \(A\) 都是 PI 代数，满足 \([x_1, x_2] = x_1 x_2 - x_2 x_1 = 0\)。此恒等式的次数为 2。

\textbf{证明}：\(A\) 交换意味着对任意 \(a, b \in A\)，\(ab = ba\)。定义多项式 \(f(x_1, x_2) = [x_1, x_2] = x_1 x_2 - x_2 x_1\)。对任意 \(a_1, a_2 \in A\)，\(f(a_1, a_2) = a_1 a_2 - a_2 a_1 = 0\)。故 \(f \in \operatorname{Id}(A)\) 为非平凡 PI 恒等式，次数为 \(\deg f = 2\)。\(f\) 是 \(A\) 的 T-理想的生成元，且 \(A\) 的 PI-次数为 \(1\)。 \(\blacksquare\)

\textbf{定理 2}（有限维代数是PI代数）：设 \(A\) 是域 \(k\) 上的有限维结合代数，\(\dim_k A = d\)。则 \(A\) 是 PI 代数，且所有次数为 \(d+1\) 的单项式的交错和给出一个恒等式。更精确地，\(A\) 满足所有次数为 \(d+1\) 的标准多项式恒等式。

\textbf{证明}：任取 \(d+1\) 个元素 \(a_1, \ldots, a_{d+1} \in A\)。由于 \(\dim_k A = d\)，这 \(d+1\) 个元素在 \(k\) 上线性相关。存在不全为零的 \(\lambda_1, \ldots, \lambda_{d+1} \in k\) 使得 \(\sum \lambda_i a_i = 0\)。将这个线性相关性编码为多项式恒等式需要对所有代换封闭。由 Amitsur-Levitzki 的一般论证（见下），标准恒等式 \(S_{2d}\) 对所有 \(d\) 维代数成立，其证明不依赖矩阵结构的特殊性（仅需线性相关的维数论证和交错性）。\(\blacksquare\)

\textbf{定理 3}（Azumaya代数）：设 \(A\) 是交换环 \(R\) 上的 Azumaya 代数（中心可分离 \(R\)-代数），且 \(A\) 作为 \(R\)-模有限生成忠实。则 \(A\) 是 \(R\) 上的 PI 代数，其 PI-次数等于 \(A\) 的次数（degree）。

\textbf{证明}：\(A\) 为 \(R\) 上 Azumaya 代数，故 \(A \otimes_R A^{\text{op}} \cong \operatorname{End}_R(A)\)。对任意剩余域 \(k = R/\mathfrak{m}\)，\(A \otimes_R k\) 为 \(k\) 上中心单代数，维数为 \(n^2\)，其中 \(n\) 为 \(A\) 的次数。由 Amitsur-Levitzki 定理，\(A \otimes_R k\) 满足 \(S_{2n} = 0\)。由于此恒等式对所有剩余域成立，且 \(A\) 忠实平坦，\(A\) 自身满足 \(S_{2n} = 0\)，故 \(A\) 为 PI 代数，PI-次数为 \(n\)。 \(\blacksquare\)

\subsection{Amitsur-Levitzki定理}\label{amitsur-levitzkiux5b9aux7406}

\textbf{定义 3}（标准多项式）：\(m\) 阶\textbf{标准多项式} \(S_m\) 定义为：

\[S_m(x_1, \ldots, x_m) = \sum_{\sigma \in \mathfrak{S}_m} \operatorname{sgn}(\sigma) \, x_{\sigma(1)} x_{\sigma(2)} \cdots x_{\sigma(m)}\]

\textbf{定理 4}（Amitsur-Levitzki，1950）：矩阵代数 \(M_n(k)\) 满足标准恒等式 \(S_{2n} = 0\)，即对任意 \(n \times n\) 矩阵 \(X_1, \ldots, X_{2n}\)，

\textbf{证明}：因 \(S_{2n}\) 为多线性，只需验证 \(X_i\) 为矩阵单位 \(E_{pq}\) 时成立。将 \(S_{2n}(E_{p_1 q_1}, \ldots, E_{p_{2n} q_{2n}})\) 展开为 \(2n!\) 项的和。关键观察：每一项为 \(E_{p_{\sigma(1)} q_{\sigma(1)}} \cdots E_{p_{\sigma(2n)} q_{\sigma(2n)}}\)，非零仅当 \(q_{\sigma(i)} = p_{\sigma(i+1)}\) 对所有 \(i\) 成立。因仅有 \(n\) 个不同指标，鸽巢原理保证 \(2n\) 个指标中必有重复，但在交错和 \(\sum \operatorname{sgn}(\sigma)\) 中，配对项抵消。严格论证利用 \(S_{2n}\) 可写为 \([X_1, \ldots, X_{2n}]\) 的迹为零的性质，结合 Cayley-Hamilton 定理的线性化版本。 \(\blacksquare\)

\[S_{2n}(X_1, \ldots, X_{2n}) = \sum_{\sigma \in \mathfrak{S}_{2n}} \operatorname{sgn}(\sigma) X_{\sigma(1)} \cdots X_{\sigma(2n)} = 0\]

且 \(2n\) 是最小的可能次数：存在 \(2n-1\) 个矩阵使 \(S_{2n-1} \neq 0\)。

\subsection{Kaplansky定理}\label{kaplanskyux5b9aux7406}

\textbf{定义 4}（本原代数）：环 \(R\) 称为\textbf{右本原的}，若它有忠实单右模。等价地，\(R\) 有极大右理想不包含任何非零理想。

\textbf{定理 5}（Kaplansky，1948）：设 \(A\) 是域 \(k\) 上的本原 PI 代数。则 \(A\) 是有限维中心单代数，具体地：

\textbf{证明}：\(A\) 本原故有忠实不可约模 \(V\)，\(A\) 在 \(\operatorname{End}_k(V)\) 中稠密。若 \(\dim_k V = \infty\)，则 \(A\) 包含任意有限秩线性变换，不满足任何 PI 恒等式。故 \(\dim_k V = n < \infty\)，\(A \cong M_n(D)\)，\(D\) 为除环。PI 条件给出 \(S_{2n} = 0\) 为 \(A\) 的非平凡恒等式。\(D\) 的中心 \(Z(D)\) 为 \(k\) 的有限扩张。由 Posner 定理，\(A\) 的中心 \(Z(A) \cong Z(D)\) 为域，\([A:Z(A)] = n^2 [D:Z(D)] < \infty\)。 \(\blacksquare\)

\[A \cong M_n(D)\]

其中 \(D\) 是 \(k\) 上的有限维可除代数（即中心单代数），且 \(n = \mathrm{PIdeg}(A)\)。

\subsection{T-理想}\label{t-ux7406ux60f3}

\textbf{定义 5}（T-理想）：自由结合代数 \(k\langle X \rangle = k\langle x_1, x_2, \ldots \rangle\) 的一个理想 \(I\) 称为 \textbf{T-理想}，若对每个代数同态 \(\varphi: k\langle X \rangle \to k\langle X \rangle\)，有 \(\varphi(I) \subseteq I\)。等价地，\(I\) 在所有变元代换下封闭。

\textbf{定理 6}（PI恒等式与T-理想的对应）：\(k\langle X \rangle\) 的 T-理想与 PI 代数的恒等式类之间存在一一对应：

\[I \longmapsto \{A \mid I \subseteq \mathrm{Id}(A)\}\]

其中 \(\mathrm{Id}(A) = \{f \in k\langle X \rangle \mid f(a_1, \ldots, a_n) = 0, \forall a_i \in A\}\) 是 \(A\) 的 \textbf{T-理想}（恒等式理想）。

\textbf{证明}：需验证 \(\mathrm{Id}(A)\) 确实是 T-理想。设 \(f(x_1, \ldots, x_n) \in \mathrm{Id}(A)\)，对任意代换 \(\varphi\)（即任意 \(g_1, \ldots, g_n \in k\langle X \rangle\)），\(f(g_1, \ldots, g_n)\) 仍是恒等式：对任意赋值 \(a_i \mapsto g_i(a_1, \ldots)\)，结果为零因 \(f\) 本身为零对任意论证。包含方向和理想性质平凡验证。\(\blacksquare\)

\subsection{Specht问题与Kemer解}\label{spechtux95eeux9898ux4e0ekemerux89e3}

T-理想的语言将 PI 代数的分类问题转化为自由代数中理想生成元的有限性问题。Specht 于 1907 年提出的核心问题是：特征零域上，每个 PI 代数的所有恒等式是否都可以由有限个''基本恒等式''通过代换生成？这一问题悬而未决八十年，直到 Kemer 在 1987 年通过引入有限维 \(\mathbb{Z}_2\)-分次代数的表示论给出了肯定的解答。

\textbf{定义 6}（有限基）：T-理想 \(I\) 称为\textbf{有限基的}（finitely based），若存在有限集合 \(F \subseteq k\langle X \rangle\) 使得 \(I\) 是由 \(F\) 生成的（最小）T-理想。等价地，\(I\) 中所有恒等式可由有限多个恒等式通过代换和线性组合导出。

\textbf{定理 7}（Specht问题，1907年提出）：在特征零域上，每个 T-理想是否都是有限基的？

\textbf{证明}：Specht 问题为开放性猜想，直到 1987 年由 Kemer 解决。证明思路：将 T-理想与 \(\mathbb{Z}_2\)-分次代数关联。Kemer 构造了从 T-理想到有限维 \(\mathbb{Z}_2\)-分次代数（Grassmann 包络）的对应。有限维 \(\mathbb{Z}_2\)-分次代数的 T-理想有有限基，通过此对应推得原 T-理想也有有限基。Kemer 证明的关键步骤：(1) 多项式恒等式的表示理论；(2) PI 代数的结构定理；(3) 有限基性质在 Grassmann 包络下的保持。 \(\blacksquare\)

\textbf{定理 8}（Kemer解，1987）：在特征零域上，每个结合代数的 T-理想是有限基的。即：特征零域上的结合 PI 代数满足 Specht 性质------任意 PI 代数的所有多项式恒等式可由有限多个恒等式生成。

\textbf{证明}：Kemer 的证明分四步。(1) 表示论归约：将自由代数 \(k\langle X \rangle\) 的多线性部分视为 \(\mathfrak{S}_n\)-模。T-理想 \(I\) 对应 \(\mathfrak{S}_n\)-子表示序列 \(I_n \subseteq P_n\)。特征零下 \(\mathfrak{S}_n\) 全可约，故 \(P_n = I_n \oplus Q_n\)，T-理想的有限基等价于 \(Q_n\) 序列满足 Noether 性质。(2) Grassmann 包络：对任意 PI 代数 \(A\)，构造 \(G(A) = A \otimes_k E\)（\(E\) 为 Grassmann 代数），\(G(A)\) 为 \(\mathbb{Z}_2\)-分次代数。Kemer 证明了 \(G(A)\) 的 T-理想与 \(A\) 的 T-理想一一对应。(3) 有限维归约：\(G(A)\) 的 T-理想由有限维 \(\mathbb{Z}_2\)-分次代数 \(B\)（\(G(A)\) 的基代数）的恒等式完全决定。\(B\) 有限维，其 T-理想由幂零指标和半单部分控制。(4) 终结：有限维 \(\mathbb{Z}_2\)-分次代数的 T-理想经有限步后稳定，即存在有限生成元集。通过 Grassmann 包络对应推回 \(A\)，得 \(A\) 的 T-理想有有限基。 \(\blacksquare\)

\subsection{Codimension序列与Regev定理}\label{codimensionux5e8fux5217ux4e0eregevux5b9aux7406}

Kemer 的定理解决了 Specht 问题在定性层面的存在性------恒等式由有限基生成------但它没有给出如何从代数 \(A\) 本身提取定量信息的途径。Codimension 序列 \(c_n(A)\) 正是为此设计的：它度量了 \(n\) 次多线性多项式中模去 \(A\) 的恒等式后所''存活''的空间维数，从而将 PI 代数的复杂度转化为一个数值序列。

\textbf{定义 7}（Codimension序列）：设 \(A\) 是 \(k\) 上的 PI 代数。令 \(P_n = \{f \in k\langle x_1, \ldots, x_n \rangle \mid f \text{ 是多线性的}\}\) 是 \(n\) 次多线性多项式的空间。定义 \(n\) 次 \textbf{codimension}：

\[c_n(A) = \dim_k \frac{P_n}{P_n \cap \mathrm{Id}(A)}\]

\(c_n(A)\) 度量 \(A\) 的恒等式在 \(n\) 次多线性层面的''截断''程度。

\textbf{定理 9}（Regev，1972）：若 \(A\) 是 PI 代数，则 codimension 序列以指数速度增长，即存在常数 \(a > 0\) 使得

\textbf{证明}：\(A\) 为 PI 代数，故存在 \(m\) 满足标准恒等式 \(S_{2m} \in \operatorname{Id}(A)\)。\(S_{2m}\) 的存在性限制 \(\mathfrak{S}_n\)-模 \(P_n/P_n \cap \operatorname{Id}(A)\) 中不可约表示的类型。具体地，仅有 Hook 长度 \(\leq m\) 的 Young 图对应的表示可存活。\(n\) 次多线性多项式中可存活的分割数量为 \(\sum_{\lambda \vdash n, \lambda_1' \leq m} 1\)，其渐近增长为 \(O(m^{2n})\)。故 \(c_n(A) = \dim(P_n/P_n \cap \operatorname{Id}(A)) \leq C \cdot a^n\)，其中 \(a = m^2\)。这给出了 \(c_n(A)\) 的指数上界，而 Regev 定理进一步证明了指数增长是精确的。 \(\blacksquare\)

\[c_n(A) \le a^n, \quad \forall n \ge 1\]

\subsection{PI-exponent定理}\label{pi-exponentux5b9aux7406}

\textbf{定义 8}（PI-exponent）：PI 代数 \(A\) 的 \textbf{PI-指数}定义为：

\[\exp(A) = \lim_{n \to \infty} \sqrt[n]{c_n(A)}\]

Regev 定理保证该极限有界；Giambruno-Zaicev 定理保证该极限存在。

\textbf{定理 10}（Giambruno-Zaicev，1999）：在特征零域上，任意 PI 代数 \(A\) 的 PI-指数存在且是整数：

\textbf{证明}：由 Kemer 的结构理论，\(A\) 的 T-理想对应有限维 \(\mathbb{Z}_2\)-分次代数 \(B\)（Grassmann 包络）。\(c_n(A)\) 的渐近行为由 \(B\) 的半单部分 \(B_{\text{ss}}\) 的表示论决定。具体地，\(\exp(A) = \dim B_{\text{ss}}\)。因 \(B_{\text{ss}}\) 为有限维半单代数，其维数为整数。极限 \(\exp(A) = \lim_{n \to \infty} \sqrt[n]{c_n(A)}\) 的存在性由 \(\mathfrak{S}_n\)-表示渐近理论保证：\(c_n(A) = \sum_{\lambda \vdash n} m_\lambda d_\lambda\)，其中 \(d_\lambda\) 为 \(\lambda\) 的不可约表示的维数，\(m_\lambda\) 为重数。渐近分析给出 \(c_n(A) \sim C \cdot n^t \cdot \exp(A)^n\)，其中 \(t\) 为半整数。 \(\blacksquare\)

\[\exp(A) \in \mathbb{N} \cup \{0\}\]

具体地，\(\exp(A) = \dim_k \overline{A}\)，其中 \(\overline{A}\) 是与 \(A\) 关联的有限维（超）代数的半单部分之维数（在 Kemer 理论意义下）。若 \(A\) 是有限维，则 \(\exp(A) \le \dim_k A\)。

\subsection{张量积的Regev定理}\label{ux5f20ux91cfux79efux7684regevux5b9aux7406}

\textbf{定义 9}（PI代数张量积）：设 \(A, B\) 是域 \(k\) 上的结合代数。称 \(A \otimes_k B\) 为 \(A\) 和 \(B\) 的张量积代数，其乘法为 \((a_1 \otimes b_1)(a_2 \otimes b_2) = a_1 a_2 \otimes b_1 b_2\)。

\textbf{定理 11}（Regev张量积定理，1972）：若 \(A\) 和 \(B\) 都是 \(k\) 上的 PI 代数，则 \(A \otimes_k B\) 也是 PI 代数。

\textbf{证明}：\(A\) 满足某标准恒等式 \(S_{2a} = 0\)，\(B\) 满足某标准恒等式 \(S_{2b} = 0\)。需构造 \(A \otimes B\) 的一个显式 PI。

令 \(m = 2a + 2b\)。考虑标准多项式 \(S_m\)。对任意 \(m\) 个张量积元 \(a_1 \otimes b_1, \ldots, a_m \otimes b_m\)，计算：

\[S_m(a_1 \otimes b_1, \ldots, a_m \otimes b_m) = \sum_{\sigma \in \mathfrak{S}_m} \operatorname{sgn}(\sigma) a_{\sigma(1)}\cdots a_{\sigma(m)} \otimes b_{\sigma(1)}\cdots b_{\sigma(m)}\]

由于 \(m = 2a + 2b\)，对每个排列 \(\sigma\)，因子 \(a_{\sigma(1)}\cdots a_{\sigma(m)}\) 中任何连续 \(2a\) 个因子可以交错化给出 \(S_{2a}\) 的贡献------但因为 \(m\) 足够大，可以利用鸽巢原理：\(m\) 个因子中必然有 \(2a\) 个来自 \(A\) 的 ``交替''块或 \(2b\) 个来自 \(B\) 的交替块。更精确地，对整个和采用 Fourier 变换（将 \(S_m\) 写成关于 Capelli 恒等式的分解），利用 Regev 的原始引理：存在 \(l\)（依赖于 \(a, b\)）使得某交错化多项式恒为零。由此得 \(A \otimes B\) 满足某显式恒等式。\(\blacksquare\)

\textbf{注记}：此定理在结合代数的 PI 理论中具有根本重要性------它表明 PI 性质在张量积下封闭，这意味着 PI 代数构成一个幺半范畴的满子范畴。结合 Amitsur-Levitzki 和 Kaplansky 定理，这进一步构建了''半单 PI 代数''------即有限个矩阵代数张量可除代数的直积------的完整分类框架。

\newpage

\chapter{05-代数/03-抽象代数二/07-表示论深化与Langlands纲领.md}\label{ux4ee3ux657003-ux62bdux8c61ux4ee3ux6570ux4e8c07-ux8868ux793aux8bbaux6df1ux5316ux4e0elanglandsux7eb2ux9886.md}

\section{第134章：有限群表示深化}\label{ux7b2c134ux7ae0ux6709ux9650ux7fa4ux8868ux793aux6df1ux5316}

有限群表示论的经典理论建立在复数域上，其中 Maschke 定理保证了群代数的半单性------每个表示完全可约，不可约表示数目等于共轭类数。然而，当基域的特征 \(p\) 整除群阶 \(|G|\) 时，这一优美的理论框架彻底崩塌：群代数不再是半单的，不可约表示的数目少于 \(p\)-正则共轭类数，出现了块分解（block decomposition）和投射模与不可约模之间的复杂关系。这一现象最早由 Richard Brauer 在 1930-1940 年代系统研究，他引入了 Brauer 特征标、分解矩阵和 Cartan 矩阵等核心工具，建立了模表示论的基本框架。模表示论与有限群论的深层联系体现在 Brauer 的三个主定理中：第一主定理建立了 \(G\) 的 \(p\)-块与 \(p\)-局部子群 \(N_G(D)\) 的块之间的双射对应（Brauer 对应）；第二主定理将特征标的诱导与限制通过 Brauer 对应联系起来；第三主定理则为块论提供了计算工具。Green 的不可分解模理论（1959）和 Green 对应（1964）进一步揭示了模表示论与 \(p\)-局部子群之间的深刻关系。本章将在 V12 第 69 章经典表示论的基础上，系统建立模表示论、块论和 Green 对应的理论框架，并介绍 McKay 猜想、Alperin 权猜想等现代前沿问题。

\subsection{222.1 特征标理论的深化}\label{ux7279ux5f81ux6807ux7406ux8bbaux7684ux6df1ux5316}

\textbf{定理 222.1}（特征标的正交关系，回顾）：\(G\) 的不可约复特征标 \(\operatorname{Irr}(G)\) 满足：

\[\frac{1}{|G|} \sum_{g \in G} \chi_i(g) \overline{\chi_j(g)} = \delta_{ij}\]

\textbf{证明}：对正则表示 \(\mathbb{C}[G] \cong \bigoplus_{i} V_i^{\oplus \dim V_i}\)，其特征标满足 \(\langle \chi_i, \chi_j \rangle = \frac{1}{|G|}\sum_{g} \chi_i(g)\overline{\chi_j(g)}\)。由 Schur 引理，\(\operatorname{Hom}_G(V_i, V_j) = 0\)（\(i \neq j\)）而 \(\operatorname{End}_G(V_i) = \mathbb{C}\)，故 \(\langle \chi_i, \chi_j \rangle = \dim \operatorname{Hom}_G(V_i, V_j) = \delta_{ij}\)。 \(\blacksquare\)

\textbf{定义 222.1}（特征标表）：有限群 \(G\) 的\textbf{特征标表}是 \(k \times k\) 方阵（\(k\) 为共轭类数），行标记不可约特征标 \(\chi_i\)，列标记共轭类 \(C_j\)。

\textbf{定理 222.2}（Frobenius-Schur 指示子）：设 \(\chi\) 是 \(G\) 的不可约复特征标。\textbf{Frobenius-Schur 指示子} \(\nu_2(\chi) = \frac{1}{|G|} \sum_{g \in G} \chi(g^2)\) 取值为 \(1\)（\(\chi\) 可实现为实表示）、\(-1\)（\(\chi\) 为实值但不可实现为实表示）或 \(0\)（\(\chi\) 非实值）。

\textbf{证明}：考虑表示 \(V\) 上的非退化 \(G\)-不变双线性型。\(\nu_2(\chi)\) 是平凡表示在 \(\operatorname{Sym}^2 V\) 中的重数减去在 \(\wedge^2 V\) 中的重数。具体地，\(\nu_2(\chi) = \dim_G \operatorname{Sym}^2 V - \dim_G \wedge^2 V\)。由 Schur 引理，若 \(V \cong V^*\)（即 \(\chi\) 实值），则存在 \(G\)-不变非退化双线性型，其对称或反对称性由 Frobenius-Schur 定理确定：\(\nu_2(\chi) = 1\)（对称型，正交表示），\(\nu_2(\chi) = -1\)（辛表示），\(\nu_2(\chi) = 0\) 当 \(V \not\cong V^*\)。 \(\blacksquare\)

\textbf{定义 222.2}（Artin 定理）：任何特征标可表示为诱导自循环子群的一维特征标的 \(\mathbb{Q}\)-线性组合。\textbf{Brauer 定理}（更强）：任何特征标可表示为诱导自初等子群的一维特征标的 \(\mathbb{Z}\)-线性组合。

\subsection{222.2 模表示论基础}\label{ux6a21ux8868ux793aux8bbaux57faux7840}

\textbf{定义 222.3}（模表示 / Modular Representation）：设 \(p\) 是素数，\(k\) 是特征 \(p\) 的代数闭域。\(G\) 在 \(k\) 上的\textbf{模表示}是群同态 \(G \to \operatorname{GL}(V)\)（\(V\) 是 \(k\) 上的有限维向量空间）。当 \(p \mid |G|\) 时，模表示论与复表示论有本质不同。

\textbf{定义 222.4}（\(p\)-模系统 / \(p\)-Modular System）：\((K, \mathcal{O}, k)\) 是 \textbf{\(p\)-模系统}，其中 \(\mathcal{O}\) 是完备离散赋值环（特征 \(0\)），\(K\) 是 \(\mathcal{O}\) 的分式域（特征 \(0\)），\(k\) 是 \(\mathcal{O}\) 的剩余域（特征 \(p\)）。\(K\) 包含 \(|G|\) 次本原单位根。

\textbf{定理 222.3}（Brauer-Nesbitt 定理）：在 \(p\)-模系统 \((K, \mathcal{O}, k)\) 下，\(K[G]\)-模和 \(k[G]\)-模的不可约表示之间存在由约化映射 \(\mathcal{O}[G] \to k[G]\) 给出的对应关系。

\textbf{证明}：取 \(K[G]\) 的不可约表示 \(V\)，取其 \(\mathcal{O}\)-整结构 \(V_{\mathcal{O}}\)（\(\mathcal{O}\)-自由 \(G\)-不变子模使 \(K \otimes_{\mathcal{O}} V_{\mathcal{O}} = V\)），则 \(\overline{V} = V_{\mathcal{O}} / \mathfrak{m} V_{\mathcal{O}}\) 为 \(k[G]\)-模。反之，\(k[G]\) 的不可约投射覆盖可提升为 \(\mathcal{O}[G]\) 的投射模，后者张量 \(K\) 给出 \(K[G]\) 的不可约表示。此提升不唯一（依赖于整结构选取），但其合成因子与分解矩阵 \(D\) 精确描述了对应。 \(\blacksquare\)

\textbf{定义 222.5}（Brauer 特征标 / Brauer Character）：设 \(\rho: G \to \operatorname{GL}_n(k)\) 是模表示。对 \(G\) 的 \(p\)-正则元 \(g\)（阶不被 \(p\) 整除），\(\rho(g)\) 的特征值在 \(K\) 中有唯一的提升（Teichmüller 提升）。\textbf{Brauer 特征标} \(\varphi(g)\) 是这些提升特征值的和。\(\operatorname{IBr}(G)\) 是所有不可约 Brauer 特征标的集合。

\textbf{定理 222.4}（Brauer 特征标的性质）：\(|\operatorname{IBr}(G)|\) 等于 \(G\) 的 \(p\)-正则共轭类数。Brauer 特征标在 \(p\)-正则共轭类上满足正交关系。

\textbf{证明}：对 \(p\)-正则共轭类数等于 \(|\operatorname{IBr}(G)|\)：注意 \(k[G]\) 的 Jacobson 根恰为所有 \(p\)-奇异元生成的理想，故 \(k[G]/J(k[G])\) 的维数（即不可约 \(k[G]\)-模个数）等于 \(p\)-正则共轭类数。正交关系：定义内积 \(\langle \varphi, \psi \rangle = \frac{1}{|G|} \sum_{g \text{ }p\text{-reg}} \varphi(g)\overline{\psi(g)}\)，由 Brauer 的模特征标理论，\(\varphi_i\) 形成 \(p\)-正则类函数空间的正交基。 \(\blacksquare\)

\textbf{定义 222.6}（分解矩阵与 Cartan 矩阵）：\textbf{分解矩阵} \(D\) 的行标记不可约普通特征标，列标记不可约 Brauer 特征标，\(D_{\chi,\varphi}\) 是 \(\chi\) 约化到特征 \(p\) 时 \(\varphi\) 的重数。\textbf{Cartan 矩阵} \(C = D^T D\)，其元素是 \(p\)-正则类上的标量积。

\subsection{222.3 块理论}\label{ux5757ux7406ux8bba}

\textbf{定义 222.7}（块 / Block）：群代数 \(k[G]\) 可分解为不可分解双理想的直和：\(k[G] = B_1 \oplus B_2 \oplus \cdots \oplus B_r\)。每个 \(B_i\) 称为 \(G\) 的一个 \textbf{\(p\)-块}（或 \textbf{Brauer 块}）。块的个数等于 \(\mathcal{O}[G]\) 分解为不可分解 \(\mathcal{O}\)-代数的个数。

\textbf{定义 222.8}（亏群 / Defect Group）：\(p\)-块 \(B\) 的\textbf{亏群} \(D\)（定义到共轭）是 \(B\) 的「远离块论平凡性」的度量：\(|D| = p^d\)，其中 \(d\) 是 \(B\) 的\textbf{亏数}（defect）。亏数为 \(0\) 的块是单代数（矩阵代数）。

\textbf{定理 222.5}（Brauer 第一主定理）：存在 \(G\) 的 \(p\)-块与 \(N_G(D)\) 的 \(p\)-块（具有相同亏群 \(D\)）之间的一一对应（Brauer 对应）。

\textbf{证明}：Brauer 构造了映射 \(b \mapsto b^G\)（块诱导），将 \(N_G(D)\) 中以 \(D\) 为亏群的块 \(b\) 映为 \(G\) 中以 \(D\) 为亏群的块。关键利用 Green 函数：对 \(G\) 中以 \(D\) 为亏群的块 \(B\)，定义 \(B\) 的 Brauer 对应为 \(N_G(D)\) 中以 \(D\) 为亏群的块，使得 \(B\) 与 \(b^{N_G(D)}\) 的普通特征标在 \(p\)-正则类上的限制一致。经 Brauer 的括号积（block orthogonality relations）验证此对应为双射。 \(\blacksquare\)

\textbf{定理 222.6}（Brauer 第二主定理）：Brauer 对应将特征标诱导和限制联系起来，为块论提供了核心计算工具。

\textbf{证明}：设 \(b\) 是 \(N_G(D)\) 的块，\(B = b^G\) 为其 Brauer 诱导块。对任意 \(p\)-子群 \(Q\) 使 \(C_G(Q) \subseteq N_G(D)\)，广义分解映射 \(d_Q\) 满足 \(d_Q^G \circ \text{Ind} = \text{Ind} \circ d_Q^{N_G(D)}\)，将 \(G\) 中块 \(B\) 的特征标在 \(Q\) 处的分解与 \(N_G(D)\) 中块 \(b\) 的特征标联系起来。该交换图是 Brauer 块论的核心技术工具。 \(\blacksquare\)

\textbf{定义 222.9}（Green 对应 / Green Correspondence，1964）：设 \(D\) 是 \(G\) 的 \(p\)-子群，\(H\) 是包含 \(N_G(D)\) 的子群。Green 对应在 \(k[G]\)-模和 \(k[H]\)-模（具有顶点 \(D\)）之间建立了一一对应，是模表示论中最强大的工具之一。

\textbf{定理 222.7}（Alperin 权猜想，1986）：\(p\)-块 \(B\) 的不可约 Brauer 特征标数等于 \(B\) 的局部子群中 \(p\)-权（Alperin 权重）的数目。这是模表示论中最重要的未完全解决的猜想之一。

\textbf{证明}：Alperin 权重定义为 \(N_G(Q)/Q\) 的亏零块中不可约射影模的数目（\(Q\) 遍历所有 \(p\)-子群）。Knorr-Robinson (1989) 证明了 Alperin 权猜想的简化版：若 \(B\) 是 \(G\) 的块，则 \(l(B) \leq \operatorname{wa}(B)\)（Alperin 权重数）。Cabanes (1988) 证明了 \(l(B) = \operatorname{wa}(B)\) 对可解群成立。一般群的情形由 Navarro-Tiep 于 2011 年简化为单广义 Fitting 子群的情形，至今对一般有限群仍在研究中。 \(\blacksquare\)

\subsection{有限群表示论的现代发展}\label{ux6709ux9650ux7fa4ux8868ux793aux8bbaux7684ux73b0ux4ee3ux53d1ux5c55}

有限群表示论在 Brauer 和 Green 的奠基工作之后，在 20 世纪下半叶继续向纵深发展。\textbf{McKay 猜想}（1972）是有限群表示论中最著名的未解决问题之一：若 \(G\) 是有限群，\(p\) 是素数，则 \(|\operatorname{Irr}_{p'}(G)| = |\operatorname{Irr}_{p'}(N_G(P))|\)，其中 \(\operatorname{Irr}_{p'}(G)\) 表示维数不被 \(p\) 整除的不可约复特征标的集合，\(P\) 是 \(G\) 的 Sylow \(p\)-子群。McKay 猜想在 2024 年由 Malle-Späth 完成证明（对 \(p=2\) 情形由 Isaacs-Malle-Navarro 于 2007 年证明），其证明依赖于约化到有限单群并通过 CFSG 分类验证。\textbf{Alperin-McKay 猜想}（1976）是 McKay 猜想的块论版本：若 \(B\) 是 \(G\) 的 \(p\)-块，\(b\) 是 \(N_G(P)\) 中与其 Brauer 对应的块，则 \(|\operatorname{Irr}_0(B)| = |\operatorname{Irr}_0(b)|\)（\(\operatorname{Irr}_0\) 表示高为零的不可约特征标）。\textbf{Broue 猜想}（1990）进一步断言：若 \(B\) 的亏群 \(D\) 是 Abel 群，则 \(B\) 与 \(N_G(D)\) 中对应的块 \(b\) 之间的导出范畴等价（splendid Rickard 等价），这是一类高度非平凡的范畴等价，已在可解群和对称群等情形下得到验证。\textbf{Gabriel 定理}和\textbf{箭图表示论}（Auslander-Reiten 1970s）为有限维代数的表示论提供了 Auslander-Reiten 箭图和几乎分裂序列（almost split sequences）的组合工具，将有限群表示论中的模结构翻译为箭图的组合数据。\textbf{Dade 猜想}（1990s）则通过投射模的链复形给出了有限群表示论与同伦论之间的深刻联系，Dade 群的构造是有限群论中最重要的不变量之一。

\hrulefill

\hrulefill

\section{第135章：代数群表示}\label{ux7b2c135ux7ae0ux4ee3ux6570ux7fa4ux8868ux793a}

代数群表示论是表示论从有限群到代数几何的自然推广。线性代数群是代数闭域上的仿射代数簇，同时具有群结构，典型的例子包括一般线性群 \(\operatorname{GL}_n\)、特殊线性群 \(\operatorname{SL}_n\)、正交群和辛群等。与有限群表示论不同，代数群的表示论充分利用了代数几何的工具（如层上同调、Borel-Weil-Bott 定理），并与根系和 Weyl 群等组合结构紧密相连。代数群的不可约表示分类由最高权理论完成（Cartan-Weyl）：每个支配权 \(\lambda\) 对应唯一的不可约模 \(L(\lambda)\)，其维数由 Weyl 维数公式给出，特征标由 Weyl 特征标公式刻画。这一理论是 20 世纪数学的巅峰成就之一，E. Cartan、H. Weyl、C. Chevalley 和 A. Borel 等人做出了奠基性贡献。代数群表示论的另一重大进展是特征 \(p>0\) 域上的正特征表示论------Frobenius 态射的引入使得表示论呈现出全新的现象，Steinberg 张量积定理和 Lusztig 猜想是这一方向的核心成果。代数群表示论还与数论（Deligne-Lusztig 理论将有限约化群的表示与 \(\ell\)-进上同调联系起来）、几何表示论（Springer 对应通过幂零轨道实现 Weyl 群的表示）和 Langlands 纲领有着深刻的交叉。

\subsection{223.1 线性代数群基础}\label{ux7ebfux6027ux4ee3ux6570ux7fa4ux57faux7840}

\textbf{定义 223.1}（线性代数群）：\textbf{线性代数群} \(G\) 是代数闭域 \(k\) 上的仿射代数簇，带有群运算（乘法 \(G \times G \to G\) 和取逆 \(G \to G\)），使得运算是代数簇的态射。

\textbf{定义 223.2}（Borel 子群与极大环面）：代数群 \(G\) 的 \textbf{Borel 子群} \(B\) 是 \(G\) 的极大连通可解子群。\(B\) 的\textbf{极大环面} \(T \subset B\) 是 \(B\) 的极大环面子群（同构于 \((\mathbb{G}_m)^n\)）。\(T\) 的\textbf{特征标群} \(X(T) = \operatorname{Hom}(T, \mathbb{G}_m)\) 是自由 Abel 群（秩 \(= \dim T\)）。

\textbf{定理 223.1}（Borel 不动点定理）：若连通可解代数群 \(G\) 作用在完备代数簇 \(X\) 上，则 \(G\) 在 \(X\) 上有不动点。

\textbf{证明}：对 \(\dim G\) 归纳。若 \(\dim G = 0\)，则 \(G \cong \{e\}\)，平凡。设 \(G_1 \subset G\) 为余维数 1 的正规子群（存在性由 \(G\) 可解得），则 \(G/G_1 \cong \mathbb{G}_a\) 或 \(\mathbb{G}_m\)。归纳假设 \(G_1\) 在 \(X\) 上有不动点集 \(X^{G_1} \neq \varnothing\)，且 \(X^{G_1}\) 为完备（闭子簇）。\(G/G_1\) 作用于 \(X^{G_1}\)。若 \(G/G_1 \cong \mathbb{G}_a\)：完备簇上的 \(\mathbb{G}_a\) 轨道必有不动点（否则完备曲线 \(\mathbb{P}^1\) 的仿射像矛盾）。若 \(G/G_1 \cong \mathbb{G}_m\)：完备簇上 \(\mathbb{G}_m\) 作用的极限点给出不动点。 \(\blacksquare\)

\textbf{定理 223.2}（Bruhat 分解）：\(G = \bigsqcup_{w \in W} B w B\)，其中 \(W = N_G(T)/T\) 是 Weyl 群。Bruhat 分解是代数群表示论的核心工具。

\textbf{证明}：由 Borel 子群结构 \(B = T \cdot U\)（\(U\) 幂幺）。利用 \(G\) 在旗簇 \(G/B\) 上的作用及 Schubert 胞腔分解：\(G/B = \bigsqcup_{w \in W} B w B / B\)。每个 \(B w B / B \cong \mathbb{A}^{\ell(w)}\)。取原像即得 \(G = \bigsqcup_w B w B\)。不交性来自 Weyl 群的 Bruhat 序与双陪集的严格分层。 \(\blacksquare\)

\subsection{223.2 最高权理论}\label{ux6700ux9ad8ux6743ux7406ux8bba}

\textbf{定义 223.3}（支配权与最高权模）：设 \(G\) 是连通约化代数群，选定 Borel 子群 \(B \supset T\)。\textbf{权} \(\lambda \in X(T)\) 是\textbf{支配的}（dominant），如果 \(\langle \lambda, \alpha^\vee \rangle \geq 0\)（对所有正根 \(\alpha\)）。对每个支配权 \(\lambda\)，存在唯一的不可约 \(G\)-模 \(L(\lambda)\)，其最高权为 \(\lambda\)。

\textbf{定义 223.4}（诱导模 / Weyl 模与对偶 Weyl 模）：对支配权 \(\lambda\)，\textbf{Weyl 模} \(\Delta(\lambda)\) 定义为 \(\operatorname{ind}_B^G(k_\lambda)\)（\(k_\lambda\) 是 \(B\) 的一维表示）。\(L(\lambda)\) 是 \(\Delta(\lambda)\) 的唯一不可约商。

\textbf{定理 223.3}（Weyl 特征标公式）：不可约表示 \(L(\lambda)\) 的特征标为

\[\chi(\lambda) = \frac{\sum_{w \in W} (-1)^{\ell(w)} e^{w(\lambda + \rho)}}{\sum_{w \in W} (-1)^{\ell(w)} e^{w\rho}}\]

其中 \(\rho\) 是正根的一半之和，\(\ell(w)\) 为 \(w\) 的长度。（这里 \(\rho\) 为所有正根之和的一半，Weyl 群的 dot action 定义为 \(w \cdot \lambda = w(\lambda + \rho) - \rho\)。）这是表示论中最优美的公式之一。

\textbf{证明}：Euler-Poincare 原理用于 Borel-Weil-Bott 定理：\(\chi(\lambda) = \sum_i (-1)^i \operatorname{ch} H^i(G/B, \mathcal{L}(\lambda))\)。由 Bott 定理，仅当 \(i = \ell(w)\) 时 \(H^i\) 非零（\(w \in W\) 满足 \(w \cdot \lambda\) 为支配权的对偶）。重排求和得 \(\chi(\lambda) = \frac{\sum_w (-1)^{\ell(w)} e^{w(\lambda+\rho)-\rho}}{\prod_{\alpha>0}(1-e^{-\alpha})}\)，分子分母同乘 \(e^\rho\)，分母用 Weyl 分母公式 \(\sum_w (-1)^{\ell(w)} e^{w\rho} = e^\rho \prod_{\alpha>0} (1-e^{-\alpha})\) 即得。 \(\blacksquare\)

\textbf{定理 223.4}（Weyl 维数公式）：\(\dim L(\lambda) = \prod_{\alpha > 0} \frac{\langle \lambda + \rho, \alpha^\vee \rangle}{\langle \rho, \alpha^\vee \rangle}\)。

\textbf{证明}：在 Weyl 特征标公式中取 \(e^\mu = 1\)（求值于单位元），应用 L'Hopital 规则得 \(\dim L(\lambda) = \frac{\prod_{\alpha>0} \langle \lambda+\rho, \alpha^\vee \rangle}{\prod_{\alpha>0} \langle \rho, \alpha^\vee \rangle}\)。由 Weyl 分母公式 \(\sum_w (-1)^{\ell(w)} e^{w\rho} = \prod_{\alpha>0} (e^{\alpha/2} - e^{-\alpha/2})\)，在 \(e^\mu \to 1\) 时各因式 \(\langle \rho, \alpha^\vee \rangle\) 自然出现。 \(\blacksquare\)

\subsection{223.3 Kempf 消失定理与上同调}\label{kempf-ux6d88ux5931ux5b9aux7406ux4e0eux4e0aux540cux8c03}

\textbf{定义 223.5}（旗簇与诱导层）：\(G/B\) 是\textbf{旗簇}（flag variety），是光滑射影簇。对权 \(\lambda\)，\(G \times^B k_\lambda\) 是 \(G/B\) 上的可逆层 \(\mathcal{L}(\lambda)\)。

\textbf{定理 223.5}（Borel-Weil-Bott 定理）：不可约 \(G\)-模 \(L(\lambda)\) 可通过旗簇上可逆层的上同调实现： - \(H^0(G/B, \mathcal{L}(\lambda)) \cong L(\lambda)^*\)（若 \(\lambda\) 是支配权） - 对非支配权，上同调消失或给出其他不可约表示（Bott 的精细公式）

\textbf{证明}（Borel-Weil-Bott）：支配权 \(\lambda\) 时，由 Kempf 消失 \(H^{>0}(G/B, \mathcal{L}(\lambda)) = 0\)，而 \(H^0 \cong \nabla(\lambda)\)（对偶 Weyl 模），其对偶为 \(L(\lambda)^*\)。非支配权 \(\mu\) 时，由 Bott 定理，唯一非零上同调在 \(i = \ell(w)\)（\(w \cdot \mu\) 支配），此时 \(H^{\ell(w)}(G/B, \mathcal{L}(\mu)) \cong L(w \cdot \mu)^*\)。 \(\blacksquare\)

\textbf{定理 223.6}（Kempf 消失定理，1976）：若 \(\lambda\) 是支配权，则 \(H^i(G/B, \mathcal{L}(\lambda)) = 0\)（对所有 \(i > 0\)）。这是 Borel-Weil-Bott 定理的特例，是射影簇上丰沛线丛消失定理的推广。

\textbf{证明}：\(\lambda\) 支配时 \(\mathcal{L}(\lambda)\) 为丰沛线丛。特征 \(0\) 由 Kodaira 消失定理直接给出 \(H^{>0}(G/B, \mathcal{L}(\lambda))=0\)。正特征下，Kempf 利用 Frobenius 分裂证明：支配权 \(\lambda\) 对应的 \(\mathcal{L}(\lambda)\) 在 \(\operatorname{Gr}\) 上诱导 Frobenius 分裂结构，从而所有正上同调消失。Andersen 随后推广到任意支配权。 \(\blacksquare\)

\textbf{定义 223.6}（Kazhdan-Lusztig 猜想，1979/证明 1981）：对特征 \(0\) 域上的 Verma 模，不可约模的合成因子重数由 Kazhdan-Lusztig 多项式给出。Brylynski-Kashiwara（1981）和 Beilinson-Bernstein（1981）通过 \(\mathcal{D}\)-模和相交上同调证明了该猜想。

\subsection{223.4 正特征表示论}\label{ux6b63ux7279ux5f81ux8868ux793aux8bba}

\textbf{定义 223.7}（正特征表示论 / 特征 \(p\) 的代数群表示）：在特征 \(p > 0\) 的代数闭域上，代数群表示论与特征 \(0\) 有本质不同。Frobenius 态射 \(F: G \to G\) 引入新的表示（Frobenius 扭）。

\textbf{定理 223.7}（Steinberg 张量积定理）：对特征 \(p\) 域上的单连通半单代数群 \(G\)，不可约 \(G\)-模的结构为

\[L(\lambda) = L(\lambda_0) \otimes L(\lambda_1)^{(1)} \otimes L(\lambda_2)^{(2)} \otimes \cdots\]

其中 \(\lambda = \lambda_0 + p\lambda_1 + p^2\lambda_2 + \cdots\) 是 \(\lambda\) 的 \(p\)-adic 展开，\(L(\mu)^{(r)}\) 是 \(F^r\)-扭表示。

\textbf{证明}：设 \(\lambda = \sum_{i=0}^m p^i\lambda_i\) 为 \(p\)-adic 展开（\(\lambda_i\) 为限制支配权）。Frobenius 态射 \(F: G \to G\) 诱导扭函子 \((-)^{(1)}\)。由 Steinberg 定理，自然映射 \(L(\lambda_0) \otimes L(\lambda_1)^{(1)} \otimes \cdots \otimes L(\lambda_m)^{(m)} \to L(\lambda)\)（由各因子的最高权向量在张量积中的像生成）为同构。证明关键：借助 Frobenius 核 \(G_1 = \ker(F)\) 的表示论，\(L(\lambda_0)\) 恰为 \(L(\lambda)|_{G_1}\) 的唯一不可约和项，其余扭因子由 \(F\)-稳定子群逐渐提升得到。 \(\blacksquare\)

\textbf{定理 223.8}（Lusztig 猜想，1980/部分证明）：对于特征 \(p \gg h\)（\(h\) 是 Coxeter 数），特征 \(p\) 的不可约模的合成因子重数由 Kazhdan-Lusztig 多项式决定。Lusztig 猜想在 \(p \gg 0\) 时成立（Andersen-Jantzen-Soergel 1994），但存在反例证明 \(p\) 需要相当大。

\textbf{证明}：Andersen-Jantzen-Soergel 通过 Soergel 双模和量子群在 \(\ell\)-次本原单位根处的限制对偶建立正特征表示论到特征 \(0\) Kazhdan-Lusztig 理论的传递。核心为 Soergel 的 Struktursatz：射影-入射包络的自同态环同构于 coinvariant 代数 \(C_W\) 的特定商。利用此环的 Koszul 性质及 Kazhdan-Lusztig 多项式的组合解释，在 \(p \gg h\) 时证明特征 \(p\) 下的分解数与特征 \(0\) 下的分解数匹配。Williamson 后来的反例表明对某些类型的李代数 \(p\) 需要远大于 Coxeter 数。 \(\blacksquare\)

\subsection{代数群表示论的现代前沿}\label{ux4ee3ux6570ux7fa4ux8868ux793aux8bbaux7684ux73b0ux4ee3ux524dux6cbf}

代数群表示论在最高权理论的基础上，20 世纪下半叶发展出了多个与代数几何和数论交叉的前沿方向。\textbf{Deligne-Lusztig 理论}（1976）是有限约化群表示论的一场革命：对有限域 \(\mathbb{F}_q\) 上的约化群 \(G(\mathbb{F}_q)\)，Deligne 和 Lusztig 通过 \(\ell\)-进上同调构造了虚拟表示 \(R_T^\theta\)（\(T\) 为极大环面，\(\theta\) 为 \(T(\mathbb{F}_q)\) 的特征标），其不可约分量构成了 \(G(\mathbb{F}_q)\) 的几乎所有不可约表示。这一理论将有限群表示论与代数几何的深层结构联系起来，Lusztig 的分类定理（1984）最终完成了 \(G(\mathbb{F}_q)\) 所有不可约特征标的分类------这被公认为 20 世纪最伟大的数学成就之一。\textbf{Kazhdan-Lusztig 多项式}（1979）不仅解决了 Verma 模的合成因子问题，还通过 Hecke 代数的基的构造（Kazhdan-Lusztig 基）为 Coxeter 群的表示论提供了全新的组合工具。\textbf{Soergel 双模}（1990s）是范畴化（categorification）表示论的关键工具，它将 Hecke 代数范畴化为 \(\mathbb{Z}\)-分次双模范畴，从而将 Kazhdan-Lusztig 多项式提升为分次维数（graded dimension）。\textbf{几何表示论}中的\textbf{仿射 Hecke 代数}和\textbf{分次仿射 Hecke 代数}（Lusztig 1980s-1990s）将代数群表示论与 Springer 对应和仿射 Weyl 群的表示论联系起来，为 Langlands 纲领中的局部 Langlands 对应提供了代数模型。\textbf{正特征表示论}的最新进展（Williamson 2010s）揭示了 Lusztig 猜想在 ``小 \(p\)'' 情形下的反例，修正了我们对模表示论中分解数的理解，并催生了\textbf{双 Sobolev 模}（Tilting 模）和\textbf{p-正则表示论}的新理论。

\hrulefill

\hrulefill

\section{第136章：几何表示论}\label{ux7b2c136ux7ae0ux51e0ux4f55ux8868ux793aux8bba}

几何表示论是 20 世纪末表示论领域的一场范式革命，其核心思想是：通过代数几何的工具------旗簇、幂零锥、仿射 Grassmann 流形和箭图簇等几何对象的拓扑和层论------来构造和研究代数结构（群、李代数、量子群）的表示。这一方向的起源可以追溯到 Springer 在 1970 年代的奠基性工作：他发现 Weyl 群在 Springer 纤维的上同调上的作用（Springer 表示）实现了 Weyl 群的所有不可约表示，从而将代数群的幂零轨道与 Weyl 群的表示论联系起来。Kazhdan-Lusztig 理论（1979）则通过 Schubert 簇的相交上同调彻底解决了 Verma 模的合成因子问题------Kazhdan-Lusztig 多项式恰好编码了奇异空间的局部拓扑信息。几何 Satake 等价（Lusztig、Ginzburg、Mirkovic-Vilonen）是几何表示论中最优美的成果之一：它建立了仿射 Grassmann 流形上的反常层范畴与 Langlands 对偶群 \(G^\vee\) 的表示范畴之间的等价，为数论中的 Langlands 纲领提供了几何模型。Nakajima 的箭图簇理论（1998）则通过辛约化构造了 Kac-Moody 代数的可积表示，将几何表示论推广到无限维李代数的范畴。这些进展使得几何表示论成为连接表示论、代数几何、辛几何和数学物理的核心桥梁。

\subsection{224.1 Springer 理论}\label{springer-ux7406ux8bba}

\textbf{定义 224.1}（Springer 纤维 / Springer 解消）：设 \(G\) 是约化代数群，\(\mathfrak{g} = \operatorname{Lie}(G)\)，\(\mathcal{N} \subset \mathfrak{g}\) 是幂零锥。\textbf{Springer 解消}是态射 \(\pi: \tilde{\mathcal{N}} = \{(x, B) : x \in \operatorname{Lie}(B) \cap \mathcal{N}\} \to \mathcal{N}\)。对 \(x \in \mathcal{N}\)，\(\pi^{-1}(x)\) 是 \textbf{Springer 纤维} \(\mathcal{B}_x\)（包含 \(x\) 的 Borel 子代数的旗簇）。

\textbf{定理 224.1}（Springer 对应，1976-1978）：Weyl 群 \(W\) 的不可约表示与 \(G\) 的幂零轨道的等变可构造层（带局部系统）之间存在一一对应。\(W\) 在 \(H^*(\mathcal{B}_x)\) 上的作用（Springer 表示）实现了 \(W\) 的所有不可约表示。

\textbf{证明}：Springer 解消 \(\pi: \tilde{\mathcal{N}} \to \mathcal{N}\) 的直像层 \(R\pi_*\underline{\mathbb{Q}}_\ell\) 是 \(\mathcal{N}\) 上的反常层。\(G\)-等变结构给出 Weyl 群 \(W\) 在 \(R\pi_*\underline{\mathbb{Q}}_\ell\) 上的作用。对幂零元 \(x \in \mathcal{N}\)，Springer 纤维 \(\mathcal{B}_x = \pi^{-1}(x)\) 的上同调 \(H^*(\mathcal{B}_x)\) 携带 \(W \times C_G(x)\) 的作用。该作用按 \(x\) 的 \(G\)-轨道分解，每个轨道上的局部系统贡献 \(W\) 的不可约表示。特别地，正则幂零轨道对应 \(W\) 的平凡表示，次正则轨道对应反射表示，零轨道给出正则表示。 \(\blacksquare\)

\textbf{定义 224.2}（Green 函数与 Springer 表示）：有限域 \(G(\mathbb{F}_q)\) 的不可约特征标可通过 Springer 对应和 Green 函数（Deligne-Lusztig 理论）参数化。

\subsection{224.2 相交上同调与 Kazhdan-Lusztig 理论}\label{ux76f8ux4ea4ux4e0aux540cux8c03ux4e0e-kazhdan-lusztig-ux7406ux8bba}

\textbf{定义 224.3}（相交上同调 / Intersection Cohomology，Goresky-MacPherson 1980）：对奇异空间 \(X\)，\textbf{相交上同调} \(IH^*(X)\) 是满足 Poincaré 对偶的「修正」上同调理论。由中间扩展（intermediate extension）的可构造层 \(\mathbf{IC}(X)\) 的超上同调给出。

\textbf{定义 224.4}（Schubert 簇与 Kazhdan-Lusztig 多项式）：旗簇 \(G/B\) 的 \textbf{Schubert 簇} \(X_w = \overline{B w B / B}\)（\(w \in W\)）。\textbf{Kazhdan-Lusztig 多项式} \(P_{x,w}(q)\) 编码了 \(X_w\) 沿 \(B x B / B\) 的局部相交上同调 Poincaré 多项式。

\textbf{定理 224.2}（Kazhdan-Lusztig 猜想，1980）：Verma 模 \(M(\lambda)\) 的不可约合成因子重数 \([M(x \cdot 0) : L(y \cdot 0)]\) 等于 Kazhdan-Lusztig 多项式 \(P_{x,y}(1)\)。由 Beilinson-Bernstein 和 Brylynski-Kashiwara（1981）通过 \(\mathcal{D}\)-模和 Riemann-Hilbert 对应证明。

\textbf{证明}：Beilinson-Bernstein 局部化将 \(\mathfrak{g}\)-模范畴等同于旗簇 \(G/B\) 上 \(\mathcal{D}\)-模的范畴，Verma 模对应于 \(\delta\)-函数 \(\mathcal{D}\)-模沿 Schubert 胞腔的标准延拓。Riemann-Hilbert 对应将这些 \(\mathcal{D}\)-模的解层翻译为 Schubert 簇的反常层，其 stalk 由相交上同调给出。而 Schubert 簇沿小胞腔的局部相交上同调的 Poincare 多项式恰为 Kazhdan-Lusztig 多项式 \(P_{x,w}\)，其常数项 \(P_{x,w}(1)\) 即为所求的重数。 \(\blacksquare\)

\textbf{定义 224.5}（\(\mathcal{D}\)-模与表示论）：旗簇 \(G/B\) 上的 \(\mathcal{D}\)-模范畴与 \(\mathfrak{g}\)-模（具有局部有限 \(B\)-作用）的范畴等价（Beilinson-Bernstein 局部化）。这是几何表示论的核心结果。

\subsection{224.3 几何 Satake 对应}\label{ux51e0ux4f55-satake-ux5bf9ux5e94}

\textbf{定义 224.6}（仿射 Grassmann 流形）：对代数闭域 \(k\) 上的约化群 \(G\)，\textbf{仿射 Grassmann 流形} \(\operatorname{Gr}_G = G(k((t))) / G(k[[t]])\)（\(k((t))\) 是 Laurent 级数域，\(k[[t]]\) 是形式幂级数环）。\(\operatorname{Gr}_G\) 是无穷维 ind-射影簇。

\textbf{定义 224.7}（几何 Satake 等价，Lusztig 1983, Ginzburg 1995, Mirković-Vilonen 2007）：存在范畴等价

\[\operatorname{Perv}_{G(k[[t]])}(\operatorname{Gr}_G) \cong \operatorname{Rep}(G^\vee)\]

其中 \(G^\vee\) 是 \(G\) 的 Langlands 对偶群，\(\operatorname{Perv}\) 是反常层（perverse sheaves）范畴。几何 Satake 将 \(G\) 的几何对象与 \(G^\vee\) 的表示范畴联系起来。

\textbf{定理 224.3}（几何 Satake 的推论）：仿射 Grassmann 流形上 \(G(k[[t]])\)-等变反常层的超上同调给出 \(G^\vee\) 的表示的几何构造。卷积积对应于表示的张量积。

\textbf{证明}：由几何 Satake 等价的函子性，反常层 \(\mathcal{F} \in \operatorname{Perv}_{G(k[[t]])}(\operatorname{Gr}_G)\) 的全局截面超上同调 \(H^*(\operatorname{Gr}_G, \mathcal{F})\) 自然具有 \(G^\vee\)-模结构。卷积积 \(\mathcal{F}_1 \star \mathcal{F}_2 = m_*(\mathcal{F}_1 \tilde{\otimes} \mathcal{F}_2)\)（\(m: G((t)) \times_{G[[t]]} \operatorname{Gr}_G \to \operatorname{Gr}_G\) 为卷积映射）在等价的函子下对应于表示张量积。球面函数层的卷积环同构于 \(G^\vee\) 的表示环。 \(\blacksquare\)

\subsection{224.4 箭图簇与表示论}\label{ux7badux56feux7c07ux4e0eux8868ux793aux8bba}

\textbf{定义 224.8}（箭图簇 / Quiver Variety，Nakajima 1998）：\textbf{箭图簇} \(\mathfrak{M}(\mathbf{v}, \mathbf{w})\) 是箭图 \(Q\) 的双重（double）表示模空间的辛约化。箭图簇是超 Kähler 流形。

\textbf{定理 224.4}（Nakajima 定理）：箭图簇的拓扑 Borel-Moore 同调实现了 Kac-Moody 代数的可积最高权表示。具体地，\(\bigoplus_{\mathbf{v}} H_*^{\operatorname{BM}}(\mathfrak{M}(\mathbf{v}, \mathbf{w}))\) 具有 \(\mathfrak{g}(Q)\) 的最高权 \(\mathbf{w}\) 的可积表示结构。

\textbf{证明}：Nakajima 构造了箭图簇 \(\mathfrak{M}(\mathbf{v}, \mathbf{w})\) 之间的 Hecke 对应 \(\mathfrak{M}(\mathbf{v}, \mathbf{v}\prime, \mathbf{w})\)，在其 Borel-Moore 同调上定义 Chevalley 生成元 \(E_i, F_i\) 的作用（通过对应映射的上同调拉回和推前）。\(K_i\) 的作用由维数向量 \(\mathbf{v}\) 的移位给出。验证这些算子满足 Kac-Moody 代数的 Serre 关系，关键利用箭图簇的辛几何结构及 Nakajima 的消没定理。最高权向量对应 \(\mathbf{v} = 0\) 的单点 Borel-Moore 同调，其权为 \(\mathbf{w}\)。 \(\blacksquare\)

\textbf{定义 224.9}（箭图簇与几何表示论的联系）：箭图簇为仿射量子群、杨振宁代数和 \(\mathcal{W}\)-代数的表示提供了统一的几何框架。Nakajima 的箭图簇理论是几何表示论的基础工具。

\subsection{几何表示论的前沿扩展}\label{ux51e0ux4f55ux8868ux793aux8bbaux7684ux524dux6cbfux6269ux5c55}

几何表示论在 21 世纪继续向多个方向扩展，成为连接代数、几何和物理的核心桥梁。\textbf{范畴化}（categorification）是几何表示论中最深刻的思想之一：将数（如维数、重数、特征标值）提升为向量空间或范畴，将代数等式（如 \(U_q(\mathfrak{g})\) 的 Serre 关系）提升为函子之间的自然同构。Khovanov 和 Lauda（2006-2009）通过 \textbf{Khovanov-Lauda-Rouquier 代数}（KLR 代数或箭图 Hecke 代数）范畴化了量子群 \(U_q(\mathfrak{g})\) 及其可积表示，将量子群提升为分次代数的表示范畴。\textbf{仿射 Springer 纤维}（Goresky-Kottwitz-MacPherson 2006）是 Springer 理论在仿射 Weyl 群情形的推广，其 \(\ell\)-进上同调实现了仿射 Hecke 代数的表示，与局部 Langlands 对应中的深度零（depth zero）表示密切相关。\textbf{Bezrukavnikov 的仿射 Springer 对应}（2006）建立了仿射 Springer 纤维的等变 K-理论到仿射 Hecke 代数表示的范畴等价，为几何 Langlands 纲领中的局部部分提供了关键构造。\textbf{Coulomb 分支}（Braverman-Finkelberg-Nakajima 2016）将箭图簇的 Borel-Moore 同调与 Coulomb 分支的相交上同调联系起来，建立了 3d \(\mathcal{N}=4\) 超对称规范理论中 Higgs 分支和 Coulomb 分支的数学对偶，这是 Symplectic 对偶（symplectic duality）------一种新的数学对偶原理------的核心实例。\textbf{BPS 代数}（Davison 2016）通过箭图簇的 Borel-Moore 同调上的 Hall 代数结构构造了 Kac-Moody 代数的量子包络代数的范畴化，将箭图簇的计数不变量（如 Donaldson-Thomas 不变量）与量子群表示论联系起来。这些发展表明，几何表示论已经成为现代数学中一个高度统一的领域，融合了代数几何、辛几何、表示论和数学物理的方法。

\hrulefill

\hrulefill

\section{第137章：Langlands 纲领深化}\label{ux7b2c137ux7ae0langlands-ux7eb2ux9886ux6df1ux5316}

Langlands 纲领是 Robert Langlands 在 1967 年提出的宏大数学猜想体系，被公认为当代数学中最具统一性和深远影响力的纲领之一。其核心思想可以简洁地表述为：数论中的 Galois 表示与调和分析中的自守表示之间存在深刻的双射对应------即''Galois 对称性''与''谱对称性''的统一。局部 Langlands 对应（LLC）是这一纲领在局部域上的精确表述：对于 \(p\)-进域 \(F\) 上的约化群 \(G(F)\)，其不可约光滑表示与 Weil-Deligne 群到 Langlands 对偶群 \(G^\vee\) 的 \(L\)-参数之间存在一一对应。对于 \(G=\operatorname{GL}_n\)，该对应已被 Harris-Taylor（2001）和 Henniart（2000）完全证明。整体 Langlands 纲领则进一步要求这些局部对应在整体自守表示的所有局部成分上兼容，并通过自守 \(L\)-函数与动机 \(L\)-函数的等式建立起数论与分析之间的桥梁。Arthur-Selberg 迹公式是这一纲领的核心计算工具------它将 \(L^2(G(\mathbb{Q})\backslash G(\mathbb{A}))\) 上正则表示的迹分解为谱和（自守表示的特征标）与几何和（轨道积分），其稳定化是 Ngo（2010 年 Fields 奖）证明基本引理的关键技术。几何 Langlands 纲领（Beilinson-Drinfeld）则通过代数曲线上的 \(\mathcal{D}\)-模与 \(G^\vee\)-局部系统的范畴等价，为传统 Langlands 纲领提供了深刻的几何解释。

\subsection{226.1 局部 Langlands 对应}\label{ux5c40ux90e8-langlands-ux5bf9ux5e94}

\textbf{定义 226.1}（局部域与局部 Langlands 对应）：设 \(F\) 是局部域（\(\mathbb{R}\)、\(\mathbb{C}\) 或 \(\mathbb{Q}_p\) 的有限扩张）。\textbf{局部 Langlands 对应}（LLC）断言 \(G(F)\) 的不可约光滑表示与 Weil-Deligne 群 \(W_F'\) 到 Langlands 对偶群 \(G^\vee(\mathbb{C})\) 的 \(L\)-参数之间存在一一对应。

\textbf{定理 226.1}（\(\operatorname{GL}_n\) 的局部 Langlands 对应，Harris-Taylor 2001, Henniart 2000）：\(\operatorname{GL}_n(F)\) 的不可约光滑超尖表示与 \(n\) 维 Frobenius 半单 Weil-Deligne 表示之间存在一一对应，保持 \(L\)-因子和 \(\varepsilon\)-因子。

\textbf{证明}：Harris-Taylor 通过 Shtuka 模空间和 p-adic 完备 unitary 型 Shimura 簇的整体-局部兼容性证明。局部方面，构造函子将 Weil-Deligne 表示 \(\rho\) 与 \(\operatorname{GL}_n(F)\) 的表示 \(\pi(\rho)\) 关联：\(\rho\) 的 Rankin-Selberg \(L\)-函数与 \(\pi(\rho)\) 的 Godement-Jacquet \(L\)-函数一致。Henniart 的数值局部 Langlands 定理通过 Brauer 形式化验证了该对应的双射性，并证明了 \(L\) 和 \(\varepsilon\)-因子在对应下的兼容性。 \(\blacksquare\)

\textbf{定义 226.2}（\(L\)-参数与 \(L\)-包）：一般约化群 \(G\) 的 \(L\)-参数 \(\varphi: W_F' \to {}^L G\) 确定一个 \(L\)-包 \(\Pi_\varphi\)（不可约表示的有限集）。\(L\)-包内的表示由 \(\varphi\) 的中心化子 \(S_\varphi\) 的不可约表示参数化。

\textbf{定理 226.2}（\(\operatorname{GL}_n\) 的局部 Langlands 对应对 \(L\)-函数和 \(\varepsilon\)-因子的保持）：对 \(\operatorname{GL}_n \times \operatorname{GL}_m\) 的 Rankin-Selberg \(L\)-函数，局部 Langlands 对应保持局部 \(L\)-函数和 \(\varepsilon\)-因子。

\textbf{证明}：设 \(\pi_1, \pi_2\) 为 \(\operatorname{GL}_n(F)\) 和 \(\operatorname{GL}_m(F)\) 的不可约光滑表示，\(\rho_1, \rho_2\) 为对应的 Weil-Deligne 表示。由局部 Langlands 对应的构造，Rankin-Selberg \(L\)-函数满足 \(L(s, \pi_1 \times \pi_2) = L(s, \rho_1 \otimes \rho_2)\)。\(\varepsilon\)-因子的匹配通过非分歧情形（Satake 参数与 Frobenius 特征值的匹配）及分歧情形的 Jacquet-Piatetski-Shapiro-Shalika 积分表示完成。 \(\blacksquare\)

\subsection{226.2 整体 Langlands 纲领}\label{ux6574ux4f53-langlands-ux7eb2ux9886}

\textbf{定义 226.3}（自守表示与 \(L\)-函数）：设 \(G\) 是 \(\mathbb{Q}\) 上的约化代数群，\(\mathbb{A}\) 是 \(\mathbb{Q}\) 的 Adele 环。\textbf{自守表示} \(\pi\) 是 \(G(\mathbb{A})\) 在 \(L^2(G(\mathbb{Q}) \backslash G(\mathbb{A}))\) 上的不可约成分。\(\pi\) 的\textbf{自守 \(L\)-函数} \(L(s, \pi, r)\) 是 Euler 乘积（\(r\) 是 \({}^L G\) 的有限维表示）。

\textbf{定理 226.3}（整体 Langlands 函子性猜想）：对 \(L\)-同态 \(\eta: {}^L H \to {}^L G\)，存在自守表示的「提升」映射（函子性转移），将 \(H(\mathbb{A})\) 的自守表示映为 \(G(\mathbb{A})\) 的自守表示，并保持 \(L\)-函数。

\textbf{证明}（已证情形的思路）：对 \(H = \operatorname{GL}_n\)，Arthur 的稳定迹公式将 \(L\)-同态 \(\eta\) 转化为迹公式几何边中轨道积分的比较。Ngo 的基本引理（2010）解决了轨道积分匹配的关键障碍。特定情形中，\(\operatorname{GL}_n \times \operatorname{GL}_m \to \operatorname{GL}_{nm}\) 对应 Rankin-Selberg 积的函子性（已证）；\(\operatorname{SO}_{2n+1} \to \operatorname{GL}_{2n}\) 由 Arthur 2013 分类定理解决；函数域情形由 Lafforgue 2002 完成。 \(\blacksquare\)

\emph{函子性猜想的简要描述}：Langlands 的函子性原理断言，任意两个约化群 \(H\) 和 \(G\) 的 \(L\)-群之间的代数同态 \(\eta: {}^L H \to {}^L G\) 诱导自守表示范畴之间的函子 \(\eta_*: \mathcal{A}(H) \to \mathcal{A}(G)\)，满足 \(L(s, \pi, r \circ \eta) = L(s, \eta_*(\pi), r)\) 对任意 \({}^L G\) 的有限维表示 \(r\)。这意味着 \(L\)-群的同态''提升''了自守形式，将较小的群上的自守表示''转移''到较大的群上。当 \(G = \operatorname{GL}_n\) 时，该猜想已由 Arthur 等人通过稳定迹公式在若干情形下得到证明；对一般的 \(G\)，Langlands 函子性原理仍然是 Langlands 纲领中最核心的未完全解决的问题，Ngô 的基本引理证明（2010）为此提供了关键技术基础。∎

\textbf{定义 226.4}（Langlands 纲领的核心猜想）： 1. \textbf{局部-整体相容性}：整体自守表示在所有局部点的局部分量通过局部 Langlands 对应与整体 Galois 表示兼容 2. \textbf{Ramanujan-Petersson 猜想}：尖点自守表示在非分歧点处的局部 Satake 参数是幺模的 3. \textbf{广义 Riemann 猜想}：自守 \(L\)-函数 \(L(s, \pi)\) 的所有非平凡零点在 \(\Re(s) = 1/2\) 上

\subsection{226.3 Arthur-Selberg 迹公式}\label{arthur-selberg-ux8ff9ux516cux5f0f}

\textbf{定义 226.5}（Arthur-Selberg 迹公式）：\textbf{Arthur-Selberg 迹公式}是 \(L^2(G(\mathbb{Q}) \backslash G(\mathbb{A}))\) 上正则表示的迹的等式，一边是谱和（自守表示的特征标），另一边是几何和（轨道积分）。它是 Poisson 求和公式在非交换调和分析中的推广。

\textbf{定理 226.4}（Arthur 的稳定迹公式，1990-2000s）：James Arthur 发展了稳定迹公式，将迹公式的几何边分解为稳定轨道积分和「内蕴形」（endoscopy）的贡献。这为一般群的 Langlands 函子性提供了基础框架。

\textbf{证明}：Arthur 将 Selberg 迹公式的几何边 \(I_{\text{geom}}(f)\) 按共轭类分解为轨道积分 \(\Phi_M(\gamma, f)\) 的加权和。稳定迹公式将轨道积分分解为稳定部分 \(\Phi_M^{\text{st}}(\gamma, f)\) 与内蕴形传递因子 \(\Delta(\gamma_H, \gamma)\)。核心结构为 \(I_{\text{geom}}(f) = \sum_{M, \gamma} |W^G(M)|^{-1} \Phi_M(\gamma, f) = \sum_{H \text{ endo}} \iota(G, H) I_{\text{geom}}^H(f^H)\)，其中 \(H\) 遍历 \(G\) 的所有椭圆内蕴群，\(\iota(G,H)\) 为 Langlands-Shelstad 传递因子。 \(\blacksquare\)

\textbf{定理 226.5}（Arthur 的分类定理，2013）：使用稳定迹公式，Arthur 分类了经典群（辛群、正交群）的自守表示谱，将自守表示的参数化为 \(L\)-参数。

\textbf{证明}：Arthur 将稳定迹公式应用于 \(\operatorname{SO}_{2n+1}, \operatorname{Sp}_{2n}, \operatorname{SO}_{2n}\) 与 \(\operatorname{GL}_N\)（其中 \(N = 2n\) 或 \(2n+1\)）的比较。利用扭曲迹公式（twisted trace formula），构造从经典群到 \(\operatorname{GL}_N\) 的函子性提升。通过分析谱边的剩余谱和尖点谱，Arthur 证明了每个离散自守表示 \(\pi\) 对应一个参数 \(\psi: \mathcal{L}_F \times \operatorname{SL}_2(\mathbb{C}) \to {}^L G\)，且 \(\pi\) 的局部分量由 \(\psi\) 通过局部 Langlands 对应唯一确定。 \(\blacksquare\)

\textbf{定义 226.6}（迹公式与 Langlands 纲领的联系）：迹公式是 Langlands 函子性猜想的核心计算工具。特别是，通过迹公式的比较（Fundamental Lemma），可以证明自守表示的函子性转移。

\subsection{226.4 几何 Langlands 纲领}\label{ux51e0ux4f55-langlands-ux7eb2ux9886}

\textbf{定义 226.7}（几何 Langlands 纲领，Beilinson-Drinfeld 1990s）：\textbf{几何 Langlands 纲领}是 Langlands 纲领的代数几何类比。设 \(X\) 是 \(\mathbb{C}\) 上的光滑射影代数曲线，\(G\) 是约化代数群。几何 Langlands 猜想断言存在范畴等价：

\[\mathcal{D}\text{-}\operatorname{Mod}(\operatorname{Bun}_G) \cong \operatorname{QCoh}(\operatorname{LocSys}_{G^\vee})\]

其中 \(\operatorname{Bun}_G\) 是 \(X\) 上 \(G\)-主丛的模叠，\(\operatorname{LocSys}_{G^\vee}\) 是 \(X\) 上 \(G^\vee\)-局部系统的模叠。

\textbf{定理 226.6}（几何 Langlands 的进展）： - \(\operatorname{GL}_1\) 的几何 Langlands：由 Laumon（1987）和 Rothstein（1996）证明（等价于 Abel-Jacobi 映射和 Fourier-Mukai 变换） - 一般 \(G\) 的几何 Langlands：Arinkin-Gaitsgory（2015）证明了 de Rham 版本的几何 Langlands 猜想（奇异支集定理） - 几何 Langlands 与 S-对偶性：Kapustin-Witten（2007）将几何 Langlands 解释为 \(\mathcal{N}=4\) 超 Yang-Mills 理论的电磁对偶

\textbf{证明}：\(\operatorname{GL}_1\) 情形：\(\operatorname{Bun}_{\operatorname{GL}_1} \cong \operatorname{Pic}(X)\)，\(\operatorname{LocSys}_{\operatorname{GL}_1} \cong \operatorname{Pic}^\tau(X) \times H^1(X, \mathcal{O}_X)\)，等价的函子为 Abel-Jacobi 映射与 Fourier-Mukai 变换。一般 \(G\) 的 de Rham 版本：Arinkin-Gaitsgory 证明了从 \(\operatorname{QCoh}(\operatorname{LocSys}_{G^\vee})\) 到 \(\mathcal{D}\text{-Mod}(\operatorname{Bun}_G)\) 的函子是满足连续性的忠实函子；反过来，利用奇异支集的条件证明了对 Hecke eigen-\(\mathcal{D}\)-模的满性，从而建立等价。Kapustin-Witten 的物理解释提供了该对应源自 4d \(\mathcal{N}=4\) SYM 紧化的电磁对偶。 \(\blacksquare\)

\textbf{定义 226.8}（基本引理 / Fundamental Lemma）：Ngô Bảo Châu（2010）证明了 Langlands-Shelstad 基本引理（Fundamental Lemma），这是稳定迹公式和 Langlands 函子性猜想的核心技术障碍。Ngô 因此获得 2010 年 Fields 奖。

\textbf{定理 226.7}（Ngô 的基本引理证明，2010）：通过 Hitchin 纤维化（Hitchin fibration）的几何，Ngô 将基本引理转化为 Hitchin 纤维化各纤维上的上同调等式，用相交上同调和 Springer 理论证明。

\textbf{证明}：Ngô 考虑 Hitchin 纤维化 \(h: \mathcal{M}_G \to \mathcal{A}_G\)（\(\mathcal{M}_G\) 为 Higgs 丛模叠，\(\mathcal{A}_G\) 为 Hitchin 基）。基本引理的轨道积分等式等价于 Hitchin 纤维化在椭圆各向异性点附近的纤维上的点计数公式。Ngô 引入弧的概念，将纤维化分解为 \(\delta\)-正则半单部分与幂零部分。通过 Laumon 的 \(\ell\)-进上同调与 Springer 纤维的相交上同调的比较，证明了纤维上的上同调纯度（cleanness），从而将点计数化为上同调的迹公式比较，利用 Grothendieck-Lefschetz 迹公式完成。 \(\blacksquare\)

\hrulefill

\emph{本卷在 V12（抽象代数 II）和 V23（李理论）中表示论基础之上，深化了有限群模表示论（Brauer 特征标、块论、Green 对应）、代数群表示论（最高权理论、Kempf 消失定理、正特征表示论）、几何表示论（Springer 理论、Kazhdan-Lusztig 理论、几何 Satake、箭图簇）、仿射与量子群表示论（仿射 Kac-Moody 代数、顶点算子代数、量子仿射代数）和 Langlands 纲领深化（局部/整体 Langlands 对应、Arthur-Selberg 迹公式、几何 Langlands 纲领）。共 5 章。}

\hrulefill

\emph{本卷建立了李理论的核心框架：李群和李代数的基本关系（指数映射、CBH 公式、Lie 对应）将连续对称性转化为代数结构；半单李代数的分类理论（根系、Dynkin 图、Cartan-Killing 分类）是 19 世纪数学最伟大的成就之一；紧李群的表示论（Peter-Weyl 定理、最高权理论、Weyl 特征公式）为调和分析和量子力学提供了基础；齐性空间、对称空间和主丛上的联络理论将李群应用于微分几何。李理论是现代数学和物理学的核心语言，从基本粒子物理（规范理论）到数论（Langlands 纲领）都有深刻应用。\textbf{补充部分}整合了现代表示论深化，涵盖有限群模表示论（Brauer特征标、块论、Green对应）、代数群表示论（最高权理论、正特征表示论）、几何表示论（Springer理论、几何Satake、箭图簇）、仿射与量子群表示论，以及Langlands纲领深化（局部/整体对应、迹公式、几何Langlands），构建了从经典表示论到现代前沿的完整桥梁。}

\emph{卷二十三：李理论终。}

\emph{卷二十三：李理论终。}

\emph{卷二十三：李理论终。}

\emph{卷二十六：李理论终。} \hrulefill

\subsection{Hopf代数与量子群导引}\label{hopfux4ee3ux6570ux4e0eux91cfux5b50ux7fa4ux5bfcux5f15}

Hopf 代数是同时具有代数结构和余代数结构的代数对象，两者通过双代数公理相容，并通过对极（antipode）映射实现''逆元''的对偶概念。Hopf 代数的概念由 Heinz Hopf 在 1941 年关于流形上同调环的工作中首次出现，随后在 1960-70 年代由 Sweedler、Milnor、Moore 等人发展为系统的代数理论。Hopf 代数的核心思想是自我对偶性：代数结构（乘法 \(\mu\)、单位 \(\eta\)）刻画了''如何将两个元素合并为一个''，而余代数结构（余乘法 \(\Delta\)、余单位 \(\varepsilon\)）刻画了''如何将一个元素分解为两个部分''。当这两个结构在双代数中相容时，Hopf 代数就自然地出现在群代数、包络代数、坐标环等基本例子中。

量子群是 Hopf 代数的非交换、非余交换变体，由 Drinfeld 和 Jimbo 在 1985 年独立发现，最初作为量子可积系统（Yang-Baxter 方程）的代数框架。量子群 \(U_q(\mathfrak{g})\) 是李代数 \(\mathfrak{g}\) 的泛包络代数 \(U(\mathfrak{g})\) 的 \(q\)-形变，当 \(q \to 1\) 时恢复经典包络代数。量子群的关键创新在于它打破了 Hopf 代数的余交换性，引入了拟三角结构（泛 R-矩阵 \(\mathcal{R} \in H \otimes H\)），该结构满足 Yang-Baxter 方程 \(\mathcal{R}_{12}\mathcal{R}_{13}\mathcal{R}_{23} = \mathcal{R}_{23}\mathcal{R}_{13}\mathcal{R}_{12}\)。这一结构使得量子群的表示范畴成为辫子幺半范畴，辫子由 \(\mathcal{R}\) 通过表示作用给出。

Hopf 代数与量子群在数学物理中有着深远的影响。辫子幺半范畴提供了构造纽结和链环不变量的代数框架------通过丝带 Hopf 代数的量子迹和 Markov 迹，每个定向链环 \(L\) 对应一个标量不变量 \(J_{H,V}(L)\)。Jones 多项式（1984）和 HOMFLY-PT 多项式正是从 \(U_q(\mathfrak{sl}_2)\) 和 \(U_q(\mathfrak{sl}_N)\) 的适当表示中得到的特例。此外，量子群还在共形场论（Wess-Zumino-Witten 模型）、可积系统（量子逆散射方法）和三维拓扑量子场论（Reshetikhin-Turaev 不变量）中扮演核心角色。本章将系统介绍 Hopf 代数的公理化定义、Sweedler 记号、基本例子（群代数、包络代数、坐标环）、量子群 \(U_q(\mathfrak{sl}_2)\) 的构造、拟三角 Hopf 代数与泛 R-矩阵、辫子幺半范畴，以及量子迹与纽结不变量的构造。

\textbf{定义 1}（代数）：设 \(k\) 是交换环。一个 \(k\)-\textbf{代数}是一个三元组 \((A, \mu, \eta)\)，其中 \(A\) 是 \(k\)-模，\(\mu: A \otimes_k A \to A\) 是乘法（乘法映射），\(\eta: k \to A\) 是单位映射，满足以下交换图： \[\mu \circ (\mu \otimes \mathrm{id}_A) = \mu \circ (\mathrm{id}_A \otimes \mu) \quad (\text{结合律})\] \[\mu \circ (\eta \otimes \mathrm{id}_A) = \mathrm{id}_A = \mu \circ (\mathrm{id}_A \otimes \eta) \quad (\text{单位律})\]

\textbf{定义 2}（余代数）：一个 \(k\)-\textbf{余代数}是一个三元组 \((C, \Delta, \varepsilon)\)，其中 \(C\) 是 \(k\)-模，\(\Delta: C \to C \otimes_k C\) 是余乘法，\(\varepsilon: C \to k\) 是余单位，满足以下交换图（结合律与单位律的对偶）： \[(\Delta \otimes \mathrm{id}_C) \circ \Delta = (\mathrm{id}_C \otimes \Delta) \circ \Delta \quad (\text{余结合律})\] \[(\varepsilon \otimes \mathrm{id}_C) \circ \Delta = \mathrm{id}_C = (\mathrm{id}_C \otimes \varepsilon) \circ \Delta \quad (\text{余单位律})\]

\textbf{定义 3}（双代数）：一个 \(k\)-\textbf{双代数}是一个五元组 \((B, \mu, \eta, \Delta, \varepsilon)\)，使得 \((B, \mu, \eta)\) 是 \(k\)-代数，\((B, \Delta, \varepsilon)\) 是 \(k\)-余代数，且 \(\Delta\) 和 \(\varepsilon\) 是代数同态。等价地，以下四个交换图成立：

\[\Delta \circ \mu = (\mu \otimes \mu) \circ (\mathrm{id}_B \otimes \tau \otimes \mathrm{id}_B) \circ (\Delta \otimes \Delta) \quad (\text{乘法与余乘法相容})\] \[\Delta \circ \eta = \eta \otimes \eta \quad (\text{单位与余乘法相容})\] \[\varepsilon \circ \mu = \varepsilon \otimes \varepsilon \quad (\text{余单位是代数同态})\] \[\varepsilon \circ \eta = \mathrm{id}_k \quad (\text{余单位与单位相容})\]

其中 \(\tau: B \otimes B \to B \otimes B\) 是交换张量因子的扭转映射 \(\tau(a \otimes b) = b \otimes a\)。

\subsection{Hopf代数的完整公理化定义}\label{hopfux4ee3ux6570ux7684ux5b8cux6574ux516cux7406ux5316ux5b9aux4e49}

\textbf{定义 4}（Hopf代数）：一个 \textbf{Hopf代数} 是域 \(k\) 上的六元组 \((H, \mu, \eta, \Delta, \varepsilon, S)\)，满足： 1. \((H, \mu, \eta, \Delta, \varepsilon)\) 是双代数。 2. \(S: H \to H\) 是 \(k\)-线性映射，称为\textbf{对极}（antipode），满足 Hopf 公理： \[\mu \circ (S \otimes \mathrm{id}_H) \circ \Delta = \eta \circ \varepsilon = \mu \circ (\mathrm{id}_H \otimes S) \circ \Delta\]

在图示意义下，以下交换图交换： \[\begin{CD}
H \otimes H @>{S \otimes \mathrm{id}}>> H \otimes H \\
@A{\Delta}AA @VV{\mu}V \\
H @>{\varepsilon}>> k @>{\eta}>> H \\
@V{\Delta}VV @AA{\mu}A \\
H \otimes H @>{\mathrm{id} \otimes S}>> H \otimes H
\end{CD}\]

对极 \(S\) 是唯一的（若存在），且满足 \(S \circ \mu = \mu \circ \tau \circ (S \otimes S)\) 和 \(S \circ \eta = \eta\) 以及 \(\varepsilon \circ S = \varepsilon\)。若 \(H\) 是交换的或余交换的，则 \(S^2 = \mathrm{id}_H\)。

\subsection{Sweedler记号及其严格性}\label{sweedlerux8bb0ux53f7ux53caux5176ux4e25ux683cux6027}

\textbf{定义 5}（Sweedler记号）：对 \(h \in H\)，记余乘法的像为有限和：

\[\Delta(h) = \sum_{(h)} h_{(1)} \otimes h_{(2)}\]

余结合律在此记号下表现为：

\[\sum_{(h)} \left( \sum_{(h_{(1)})} h_{(1)(1)} \otimes h_{(1)(2)} \right) \otimes h_{(2)} = \sum_{(h)} h_{(1)} \otimes \left( \sum_{(h_{(2)})} h_{(2)(1)} \otimes h_{(2)(2)} \right)\]

这使得我们可以无歧义地记作 \(\sum_{(h)} h_{(1)} \otimes h_{(2)} \otimes h_{(3)}\)。在 Sweedler 记号下，Hopf 公理读作：

\[\sum_{(h)} S(h_{(1)}) h_{(2)} = \varepsilon(h) 1_H = \sum_{(h)} h_{(1)} S(h_{(2)})\]

此记号虽非严格形式化（和式依赖于代表元的选取），但极大地简化了计算。所有使用 Sweedler 记号的恒等式均可还原为标准交换图语言。

\subsection{三个基础例子}\label{ux4e09ux4e2aux57faux7840ux4f8bux5b50}

\textbf{定理 1}（群代数）：设 \(G\) 是群，\(k[G]\) 是群代数。定义： - \(\mu(g \otimes h) = gh\)，\(\eta(1_k) = 1_G\) - \(\Delta(g) = g \otimes g\)，\(\varepsilon(g) = 1_k\) - \(S(g) = g^{-1}\)

则 \((k[G], \mu, \eta, \Delta, \varepsilon, S)\) 是 Hopf 代数。

\textbf{证明}：代数公理由群乘法结合律和单位元保证。余结合律：\((\Delta \otimes \mathrm{id}) \Delta(g) = g \otimes g \otimes g = (\mathrm{id} \otimes \Delta) \Delta(g)\)。余单位律：\((\varepsilon \otimes \mathrm{id}) \Delta(g) = 1 \otimes g = g\)，同理另一侧。双代数条件由 \(\Delta(gh) = gh \otimes gh = (g \otimes g)(h \otimes h) = \Delta(g)\Delta(h)\) 验证。Hopf 公理：\(\mu(S \otimes \mathrm{id})\Delta(g) = g^{-1}g = e = \varepsilon(g)1_H = gg^{-1} = \mu(\mathrm{id} \otimes S)\Delta(g)\)。\(\blacksquare\)

\textbf{定理 2}（包络代数）：设 \(\mathfrak{g}\) 是李代数，\(U(\mathfrak{g})\) 是泛包络代数。定义 \(\Delta(x) = x \otimes 1 + 1 \otimes x\) 对 \(x \in \mathfrak{g}\)，线性延拓为代数同态；\(\varepsilon(x) = 0\) 对 \(x \in \mathfrak{g}\)；\(S(x) = -x\) 对 \(x \in \mathfrak{g}\)，并延拓为反代数同态。则 \((U(\mathfrak{g}), \mu, \eta, \Delta, \varepsilon, S)\) 是 Hopf 代数。

\textbf{证明}：余结合律对生成元验证：对 \(x \in \mathfrak{g}\)，\((\Delta \otimes \mathrm{id})\Delta(x) = x \otimes 1 \otimes 1 + 1 \otimes x \otimes 1 + 1 \otimes 1 \otimes x = (\mathrm{id} \otimes \Delta)\Delta(x)\)。双代数条件因 \(\Delta\) 定义为代数同态而自动满足。Hopf 公理：\(\mu(S \otimes \mathrm{id})\Delta(x) = -x + x = 0 = \varepsilon(x)1 = \mu(\mathrm{id} \otimes S)\Delta(x)\)。\(\blacksquare\)

\textbf{定理 3}（坐标环）：设 \(G\) 是仿射代数群，\(\mathcal{O}(G)\) 是其正则函数环。定义 \(\mu\) 为逐点乘法，\(\eta(1) = 1_G\)（常值函数），\(\Delta(f)(g, h) = f(gh)\)，\(\varepsilon(f) = f(e)\)，\(S(f)(g) = f(g^{-1})\)。则 \((\mathcal{O}(G), \mu, \eta, \Delta, \varepsilon, S)\) 是交换 Hopf 代数。

\textbf{证明}：\(\Delta\) 的余结合律由 \(G\) 中群乘法的结合律保证：\(((\Delta \otimes \mathrm{id})\Delta(f))(g, h, k) = f((gh)k) = f(g(hk)) = ((\mathrm{id} \otimes \Delta)\Delta(f))(g, h, k)\)。双代数性质因 \(\Delta\) 是代数同态（\(\Delta(f_1 f_2) = \Delta(f_1)\Delta(f_2)\) 来自群乘法为代数同态）。Hopf 公理：\((\mu(S \otimes \mathrm{id})\Delta(f))(g) = f(g^{-1}g) = f(e) = \varepsilon(f)1_G\)。\(\blacksquare\)

\subsection{\texorpdfstring{量子群：Drinfeld-Jimbo \(U_q(\mathfrak{sl}_2)\)}{量子群：Drinfeld-Jimbo U\_q(\textbackslash mathfrak\{sl\}\_2)}}\label{ux91cfux5b50ux7fa4drinfeld-jimbo-u_qmathfraksl_2}

\textbf{定义 6}（\(q\)-整数与\(q\)-阶乘）：对 \(q \in k^\times\)，\(q \neq \pm 1\)，定义 \(q\)-整数和 \(q\)-阶乘：

\[[n]_q = \frac{q^n - q^{-n}}{q - q^{-1}}, \quad [n]_q! = [1]_q [2]_q \cdots [n]_q\]

自然地 \([0]_q! = 1\)。注意 \(\lim_{q \to 1} [n]_q = n\)。

\textbf{定义 7}（\(U_q(\mathfrak{sl}_2)\)）：Drinfeld-Jimbo 量子群 \(U_q(\mathfrak{sl}_2)\) 是域 \(k\) 上由生成元 \(E, F, K, K^{-1}\) 生成的结合代数，满足以下六条关系：

\[KK^{-1} = 1 = K^{-1}K\] \[KEK^{-1} = q^2 E\] \[KFK^{-1} = q^{-2} F\] \[[E, F] = \frac{K - K^{-1}}{q - q^{-1}}\]

其中 \(q\)-整数 \([2]_q = q + q^{-1}\) 自然出现在高阶结构中。\(U_q(\mathfrak{sl}_2)\) 上的 Hopf 代数结构定义如下：

\[\Delta(E) = E \otimes 1 + K \otimes E, \quad \Delta(F) = F \otimes K^{-1} + 1 \otimes F\] \[\Delta(K) = K \otimes K, \quad \Delta(K^{-1}) = K^{-1} \otimes K^{-1}\] \[\varepsilon(E) = \varepsilon(F) = 0, \quad \varepsilon(K) = \varepsilon(K^{-1}) = 1\] \[S(E) = -K^{-1}E, \quad S(F) = -FK, \quad S(K) = K^{-1}, \quad S(K^{-1}) = K\]

当 \(q \to 1\) 时，\([E, F] \to H\)（其中 \(H = \lim_{q \to 1} (K - K^{-1})/(q - q^{-1})\)），恢复经典包络代数 \(U(\mathfrak{sl}_2)\)。

\subsection{拟三角Hopf代数与泛R-矩阵}\label{ux62dfux4e09ux89d2hopfux4ee3ux6570ux4e0eux6cdbr-ux77e9ux9635}

\textbf{定义 8}（拟三角Hopf代数）：Hopf 代数 \((H, \mu, \eta, \Delta, \varepsilon, S)\) 称为\textbf{拟三角的}（quasitriangular），如果存在可逆元 \(\mathcal{R} \in H \otimes H\)，称为\textbf{泛R-矩阵}，满足：

\[(\Delta \otimes \mathrm{id})(\mathcal{R}) = \mathcal{R}_{13}\mathcal{R}_{23}\] \[(\mathrm{id} \otimes \Delta)(\mathcal{R}) = \mathcal{R}_{13}\mathcal{R}_{12}\] \[\tau \circ \Delta(h) = \mathcal{R} \Delta(h) \mathcal{R}^{-1}, \quad \forall h \in H\]

其中 \(\mathcal{R}_{12} = \mathcal{R} \otimes 1\)，\(\mathcal{R}_{23} = 1 \otimes \mathcal{R}\)，\(\mathcal{R}_{13} = (\mathrm{id} \otimes \tau)(\mathcal{R} \otimes 1)\)，而 \(\tau\) 如前是张量因子的扭转映射。

\textbf{定理 4}（Yang-Baxter方程）：泛R-矩阵自动满足 Yang-Baxter 方程：

\[\mathcal{R}_{12}\mathcal{R}_{13}\mathcal{R}_{23} = \mathcal{R}_{23}\mathcal{R}_{13}\mathcal{R}_{12}\]

\textbf{证明}：计算 \(\mathcal{R}_{12}\mathcal{R}_{13}\mathcal{R}_{23} = \mathcal{R}_{12} \cdot (\Delta \otimes \mathrm{id})(\mathcal{R}) = (\tau \otimes \mathrm{id})[(\Delta \otimes \mathrm{id})(\mathcal{R})] \cdot \mathcal{R}_{23}\)。利用拟三角条件 \((\Delta \otimes \mathrm{id})(\mathcal{R}) = \mathcal{R}_{13}\mathcal{R}_{23}\) 和 \((\mathrm{id} \otimes \Delta)(\mathcal{R}) = \mathcal{R}_{13}\mathcal{R}_{12}\)，经过标准计算可得两方向相等。详细验证：\(\mathcal{R}_{12}\mathcal{R}_{13}\mathcal{R}_{23} = \mathcal{R}_{12}(\Delta \otimes \mathrm{id})(\mathcal{R}) = (\tau \otimes \mathrm{id}) \circ (\Delta^{\mathrm{op}} \otimes \mathrm{id})(\mathcal{R}) \cdot \mathcal{R}_{23}\)，再由拟三角第三条件的张量形式，即得 \(\mathcal{R}_{23}\mathcal{R}_{13}\mathcal{R}_{12}\)。\(\blacksquare\)

\subsection{辫子幺半范畴}\label{ux8fabux5b50ux5e7aux534aux8303ux7574}

\textbf{定义 9}（辫子幺半范畴）：幺半范畴 \((\mathcal{C}, \otimes, I, \alpha, \lambda, \rho)\) 称为\textbf{辫子的}（braided），若存在自然同构 \(c_{X,Y}: X \otimes Y \to Y \otimes X\) 满足六角公理：

\[c_{X \otimes Y, Z} = (c_{X,Z} \otimes \mathrm{id}_Y) \circ (\mathrm{id}_X \otimes c_{Y,Z})\] \[c_{X, Y \otimes Z} = (\mathrm{id}_Y \otimes c_{X,Z}) \circ (c_{X,Y} \otimes \mathrm{id}_Z)\]

\textbf{定理 5}（从R-矩阵构造辫子）：设 \((H, \mathcal{R})\) 是拟三角Hopf代数，则其表示范畴 \(\mathrm{Rep}(H)\) 是辫子幺半范畴，辫子定义如下：对任意表示 \(V, W\) 和 \(v \in V, w \in W\)，

\[c_{V,W}(v \otimes w) = \tau(\mathcal{R} \cdot (v \otimes w))\]

其中 \(\mathcal{R}\) 通过表示作用在 \(V \otimes W\) 上，\(\tau\) 是通常的张量因子交换。

\textbf{证明}：需验证 \(c_{V,W}\) 是 \(H\)-模同态。对 \(h \in H\)，利用拟三角条件 \(\tau \circ \Delta(h) = \mathcal{R} \Delta(h) \mathcal{R}^{-1}\) 可得 \(c_{V,W} \circ h = h \circ c_{V,W}\)。六角公理由 \((\Delta \otimes \mathrm{id})(\mathcal{R})\) 和 \((\mathrm{id} \otimes \Delta)(\mathcal{R})\) 的条件逐项验证：\(c_{U \otimes V, W} = (c_{U,W} \otimes \mathrm{id}_V)(\mathrm{id}_U \otimes c_{V,W})\) 直接来自 \((\Delta \otimes \mathrm{id})(\mathcal{R}) = \mathcal{R}_{13}\mathcal{R}_{23}\)。同理第二个六角公理由另一拟三角条件保证。自然性来自 \(\mathcal{R}\) 为 \(H \otimes H\) 中元。\(\blacksquare\)

\subsection{量子迹与纽结不变量}\label{ux91cfux5b50ux8ff9ux4e0eux7ebdux7ed3ux4e0dux53d8ux91cf}

\textbf{定义 10}（量子迹）：设 \(H\) 是丝带 Hopf 代数（拟三角 + 中心元 \(\nu\) 满足特定条件），\(V\) 是有限维 \(H\)-模。定义 \(V\) 的\textbf{量子迹}（quantum trace）为：

\[\mathrm{Tr}_q^V(f) = \mathrm{Tr}(f \circ \nu^{-1})\]

其中 \(\mathrm{Tr}\) 是通常的向量空间迹。量子迹满足循环性和部分迹标准性质。

\textbf{定义 11}（Markov迹）：设 \(\{B_n\}_{n \ge 1}\) 是辫子群序列，\(\mathrm{Rep}(H)\) 中对象 \(V\) 诱导了表示 \(\rho_n: k[B_n] \to \mathrm{End}(V^{\otimes n})\)。\textbf{Markov迹} 是一族线性函数 \(\{\mathrm{MT}_n: \mathrm{End}_H(V^{\otimes n}) \to k\}_{n \ge 1}\)，满足： 1. \(\mathrm{MT}_n(fg) = \mathrm{MT}_n(gf)\)（迹性） 2. \(\mathrm{MT}_{n+1}(f \otimes \mathrm{id}_V) = \mathrm{MT}_n(f)\) 和 \(\mathrm{MT}_{n+1}(f \circ c_{V,V}^{\pm 1} \otimes \mathrm{id}_V) = \mathrm{MT}_n(f)\)（Markov性质）

\textbf{定理 6}（纽结不变量构造）：给定丝带 Hopf 代数 \(H\) 和有限维表示 \(V\)，对每个定向链环 \(L\)（表示为辫子的闭包），定义：

\[J_{H,V}(L) = \mathrm{MT}_n(\rho_n(\beta))\]

其中 \(\beta \in B_n\) 满足其闭包同痕于 \(L\)。则 \(J_{H,V}(L)\) 是定向链环的同痕不变量。Jones 多项式和 HOMFLY-PT 多项式可分别从适当参数的 \(U_q(\mathfrak{sl}_2)\) 和 \(U_q(\mathfrak{sl}_N)\) 的表示中得到。

\textbf{证明}：需验证 \(J_{H,V}(L)\) 在 Markov 移动下不变。Markov 移动 I：将 \(\beta \in B_n\) 换为 \(\gamma \beta \gamma^{-1}\)（\(\gamma \in B_n\)）。由 Markov 迹的迹性 \(\mathrm{MT}_n(\rho_n(\gamma \beta \gamma^{-1})) = \mathrm{MT}_n(\rho_n(\gamma^{-1}\gamma \beta)) = \mathrm{MT}_n(\rho_n(\beta))\)，故值不变。Markov 移动 II：将 \(\beta \in B_n\) 换为 \(\beta \sigma_n^{\pm 1} \in B_{n+1}\)。由 Markov 迹的性质 \(\mathrm{MT}_{n+1}(\rho_{n+1}(\beta \sigma_n^{\pm 1})) = \mathrm{MT}_n(\rho_n(\beta))\)，故值不变。由于任意两个表示同一链环的辫子可通过 Markov 移动互相转化，\(J_{H,V}(L)\) 良定且为同痕不变量。对 \(U_q(\mathfrak{sl}_2)\) 取 \(V\) 为二维标准表示，得到 Jones 多项式；对 \(U_q(\mathfrak{sl}_N)\) 取 \(V\) 为 \(N\) 维标准表示，得到 HOMFLY-PT 多项式。 \(\blacksquare\)

\newpage

\chapter{05-代数/04-李理论/01-李群与李代数.md}\label{ux4ee3ux657004-ux674eux7406ux8bba01-ux674eux7fa4ux4e0eux674eux4ee3ux6570.md}

\section{第138章：李群与李代数}\label{ux7b2c138ux7ae0ux674eux7fa4ux4e0eux674eux4ee3ux6570}

李群是微分流形和群的结合，李代数是李群在单位元处的切空间，带有 Lie 括号运算。本章建立李群和李代数的基本理论，以及指数映射。

\subsection{119.1 李群的定义}\label{ux674eux7fa4ux7684ux5b9aux4e49}

\textbf{定义 119.1}（李群）：\textbf{李群} \(G\) 是一个光滑流形（微分流形），同时是一个群，使得群运算（乘法 \(G \times G \to G\) 和取逆 \(G \to G\)）都是光滑映射。

\textbf{定义 119.2}（左不变向量场）：\(G\) 上的向量场 \(X\) 称为\textbf{左不变}的，如果 \((L_g)_* X = X\) 对所有 \(g \in G\)，其中 \(L_g(h) = gh\) 是左平移。

\textbf{命题 119.1}：左不变向量场构成 \(G\) 的切空间 \(T_e G\) 的同构。左不变向量场的 Lie 括号仍为左不变，因此在 \(T_e G\) 上诱导了 Lie 括号。

\textbf{证明}：左不变向量场 \(X\) 由 \(X_e \in T_e G\) 唯一确定：\(X_g = (L_g)_* X_e\)。反之，对任意 \(v \in T_e G\)，定义 \(X_g = (L_g)_* v\)，则 \(X\) 是光滑的（因 \(L_g\) 光滑）且左不变：\((L_h)_* X_g = (L_h)_* (L_g)_* v = (L_{hg})_* v = X_{hg}\)。故 \(X \mapsto X_e\) 给出左不变向量场与 \(T_e G\) 的线性同构。对 Lie 括号，\((L_g)_* [X, Y] = [(L_g)_* X, (L_g)_* Y] = [X, Y]\)，故左不变向量场的 Lie 括号仍为左不变。在 \(T_e G\) 上定义 \([v, w] = [X, Y]_e\)（\(X, Y\) 分别为 \(v, w\) 对应的左不变向量场），即得 \(T_e G\) 上的 Lie 代数结构。\(\blacksquare\)

\subsection{119.2 李代数的定义}\label{ux674eux4ee3ux6570ux7684ux5b9aux4e49}

\textbf{定义 119.3}（李代数）：（实或复）向量空间 \(\mathfrak{g}\) 连同双线性映射 \([\cdot, \cdot]: \mathfrak{g} \times \mathfrak{g} \to \mathfrak{g}\)（\textbf{Lie 括号}）称为\textbf{李代数}，如果满足： 1. 反称性：\([X, Y] = -[Y, X]\) 2. Jacobi 恒等式：\([X, [Y, Z]] + [Y, [Z, X]] + [Z, [X, Y]] = 0\)

\textbf{定义 119.4}（李群对应的李代数）：李群 \(G\) 的\textbf{李代数} \(\mathfrak{g} = \operatorname{Lie}(G) = T_e G\)，带有从左不变向量场的 Lie 括号诱导的括号。

\textbf{定理 119.1}（Lie 第三定理的局部版本）：每个有限维李代数 \(\mathfrak{g}\) 是某个李群 \(G\) 的李代数。\(G\) 在连通且单连通时唯一确定。

\textbf{证明}（含 Ado 定理）：Ado 定理（1947）断言任意有限维 Lie 代数 \(\mathfrak{g}\) 可忠实表示为 \(\mathfrak{gl}(n)\) 的子代数，即存在 injective Lie 代数同态 \(\rho: \mathfrak{g} \hookrightarrow \mathfrak{gl}(n)\)。由于 \(\mathfrak{gl}(n) = \operatorname{Lie}(\operatorname{GL}(n))\)，Lie 代数 \(\mathfrak{gl}(n)\) 对应 Lie 群 \(\operatorname{GL}(n)\)。在 \(\operatorname{GL}(n)\) 的单位元邻域 \(U\) 上，指数映射 \(\exp: \mathfrak{gl}(n) \to \operatorname{GL}(n)\) 是微分同胚。取 \(\rho(\mathfrak{g}) \subseteq \mathfrak{gl}(n)\) 在 \(\exp\) 下的像 \(\exp(\rho(\mathfrak{g}) \cap U')\)（\(U'\) 充分小），由 CBH 公式，该像在乘法下封闭，构成 \(\operatorname{GL}(n)\) 的局部子群。取该局部子群生成的连通李群 \(G\)（或取单连通覆盖），则 \(\operatorname{Lie}(G) \cong \mathfrak{g}\)。单连通性条件下的唯一性由 Lie 对应（定理 119.3）保证。\(\blacksquare\)

\subsection{119.3 指数映射}\label{ux6307ux6570ux6620ux5c04}

\textbf{定义 119.5}（指数映射）：李群 \(G\) 上的\textbf{指数映射} \(\exp: \mathfrak{g} \to G\) 定义为 \(\exp(X) = \gamma_X(1)\)，其中 \(\gamma_X: \mathbb{R} \to G\) 是满足 \(\gamma_X(0) = e\) 和 \(\gamma_X'(t) = X_{\gamma_X(t)}\) 的唯一单参数子群。

\textbf{命题 119.2}（指数映射的性质）： 1. \(\exp(0) = e\) 2. \(\exp((t+s)X) = \exp(tX) \exp(sX)\) 3. \(\frac{d}{dt}\big|_{t=0} \exp(tX) = X\) 4. \(\exp\) 是单位元 \(e \in G\) 附近的微分同胚 5. 对矩阵李群，\(\exp(A) = \sum_{k=0}^\infty \frac{A^k}{k!}\)（矩阵指数）

\textbf{证明}：（1）由定义 \(\exp(0) = \gamma_0(1) = e\)（\(\gamma_0(t) \equiv e\) 是平凡单参数子群）。（2）固定 \(s\)，令 \(\gamma_1(t) = \exp((t+s)X)\) 和 \(\gamma_2(t) = \exp(tX)\exp(sX)\)。两者均为满足 \(\gamma(0) = \exp(sX)\) 和 \(\gamma'(t) = X_{\gamma(t)}\) 的单参数子群，由唯一性 \(\gamma_1 = \gamma_2\)，取 \(t=1\) 即得。（3）由定义 \(\frac{d}{dt}\big|_{t=0} \exp(tX) = \gamma_X'(0) = X_e\)。（4）\(\exp\) 在 \(0 \in \mathfrak{g}\) 处的微分 \(d\exp_0: T_0\mathfrak{g} \cong \mathfrak{g} \to T_e G \cong \mathfrak{g}\) 是恒等映射，由反函数定理，\(\exp\) 在 \(0\) 附近是微分同胚。（5）对矩阵李群，单参数子群 \(\gamma(t) = e^{tA}\) 满足 \(\gamma'(t) = A e^{tA} = A \gamma(t)\)，且 \(\gamma(0) = I\)，故 \(\exp(A) = \gamma(1) = e^A\)。\(\blacksquare\)

\textbf{定理 119.2}（Campbell-Baker-Hausdorff 公式）：对充分小的 \(X, Y \in \mathfrak{g}\)，

\[\exp(X) \exp(Y) = \exp(X + Y + \frac{1}{2}[X, Y] + \frac{1}{12}[X, [X, Y]] - \frac{1}{12}[Y, [X, Y]] + \cdots)\]

CBH 公式表明李群在单位元附近的乘法完全由李代数决定。

\textbf{证明}：设 \(Z(t) = \log(\exp(tX) \exp(tY))\)（对充分小 \(t\)）。计算 \(Z(t)\) 的导数。由指数映射的微分公式： \[\frac{d}{dt}\exp(Z(t)) = \exp(Z(t))\frac{1 - e^{-\operatorname{ad}_{Z(t)}}}{\operatorname{ad}_{Z(t)}}(Z'(t)).\] 同时直接计算 \(\frac{d}{dt}(\exp(tX)\exp(tY)) = \exp(tX)X\exp(tY) + \exp(tX)Y\exp(tY) = \exp(tX)\exp(tY)(\operatorname{Ad}_{\exp(-tY)}X + Y)\)。令两者相等，并利用 \(\operatorname{Ad}_{\exp(-tY)} = e^{-t\operatorname{ad}_Y}\)，得 \[Z'(t) = \frac{\operatorname{ad}_{Z(t)}}{1 - e^{-\operatorname{ad}_{Z(t)}}}(e^{-t\operatorname{ad}_Y}X + Y).\] 将 \(Z(t)\) 展开为 \(Z(t) = tZ_1 + t^2Z_2 + \cdots\)，右端展开为 \(t\) 的幂级数。比较系数：\(Z_1 = X + Y\)，\(Z_2 = \frac{1}{2}[X, Y]\)，\(Z_3 = \frac{1}{12}([X,[X,Y]] + [Y,[Y,X]])\)，依此类推。级数在 \(\|X\|, \|Y\| < \frac{1}{2}\log 2\) 时绝对收敛（由 Dynkin 的范数估计）。取 \(t=1\) 即得 CBH 公式。\(\blacksquare\)

\subsection{119.4 李群同态与李代数同态}\label{ux674eux7fa4ux540cux6001ux4e0eux674eux4ee3ux6570ux540cux6001}

\textbf{定理 119.3}（Lie 对应）：设 \(G, H\) 是李群，\(G\) 连通单连通。李群同态 \(\varphi: G \to H\) 与李代数同态 \(d\varphi: \mathfrak{g} \to \mathfrak{h}\) 一一对应。

\textbf{证明}：设 \(\mathfrak{g} = \operatorname{Lie}(G)\)，\(\mathfrak{h} = \operatorname{Lie}(H)\)。给定李群同态 \(\varphi: G \to H\)，微分 \(d\varphi_e: \mathfrak{g} \to \mathfrak{h}\) 是李代数同态（由 \(\varphi\) 的函子性）。反之，给定李代数同态 \(\psi: \mathfrak{g} \to \mathfrak{h}\)，需构造李群同态 \(\phi: G \to H\)。由 Lie 第三定理，\(G\) 是 \(\mathfrak{g}\) 的单连通李群。考虑积李代数 \(\mathfrak{g} \times \mathfrak{h}\) 的子代数 \(\mathfrak{k} = \{(X, \psi(X)) \mid X \in \mathfrak{g}\}\)（\(\psi\) 的图）。由 Lie 第三定理，存在单连通李群 \(K\) 以 \(\mathfrak{k}\) 为李代数。自然投影 \(\pi_1: \mathfrak{k} \to \mathfrak{g}\) 是李代数同构（因 \(\psi\) 为单射则 \(\mathfrak{k} \cong \mathfrak{g}\)，一般情形由 \(\psi\) 的图是 \(\mathfrak{g} \times \mathfrak{h}\) 的闭子代数）。由 Lie 对应，\(\pi_1\) 积分得李群同态 \(\Pi_1: K \to G\)，其为覆盖映射；因 \(G\) 单连通，\(\Pi_1\) 为同构。定义 \(\phi = \Pi_2 \circ \Pi_1^{-1}: G \to H\)（\(\Pi_2: K \to H\) 为第二投影），则 \(d\phi = \psi\)。此构造给出所需双射。\(\blacksquare\)

\textbf{定义 119.6}（伴随表示）：李群 \(G\) 的\textbf{伴随表示} \(\operatorname{Ad}: G \to \operatorname{GL}(\mathfrak{g})\) 定义为 \(\operatorname{Ad}(g) = d(I_g)_e\)，其中 \(I_g(h) = ghg^{-1}\) 是内自同构。其微分 \(\operatorname{ad}: \mathfrak{g} \to \mathfrak{gl}(\mathfrak{g})\) 为 \(\operatorname{ad}(X)(Y) = [X, Y]\)。

\textbf{定理 119.4}（李群-李代数对应）：连通李群的子群与子代数、理想与正规子群、商群与商代数之间有一一对应。单连通李群范畴等价于有限维李代数范畴。

\textbf{证明}：由 Lie 对应（定理 119.3），连通单连通李群 \(G\) 和李代数 \(\mathfrak{g}\) 一一对应。子李群 \(H \subseteq G\) 通过切映射对应子李代数 \(\mathfrak{h} \subseteq \mathfrak{g}\)（Frobenius 定理）。正规闭子群对应理想：\(H \trianglelefteq G\) 当且仅当 \(\mathfrak{h}\) 是 \(\mathfrak{g}\) 的理想。商群 \(G/H\) 的李代数为 \(\mathfrak{g}/\mathfrak{h}\)（由指数映射的函子性）。这给出了范畴等价 \(\mathbf{Lie}: \mathbf{SLieGrp} \to \mathbf{LieAlg}_{\text{fin}}\) 及其拟逆（Cartan 整体化定理）。\(\blacksquare\)

\hrulefill

\hrulefill

\hrulefill

\hrulefill

\hrulefill

\hrulefill

\section{第139章：半单李代数}\label{ux7b2c139ux7ae0ux534aux5355ux674eux4ee3ux6570}

半单李代数是李理论的核心研究对象，具有完全可约性和丰富的结构理论。本章讨论 Killing 型、Cartan 子代数、根系和 Dynkin 图分类。

\subsection{120.1 Killing 型与半单性}\label{killing-ux578bux4e0eux534aux5355ux6027}

\textbf{定义 120.1}（Killing 型）：李代数 \(\mathfrak{g}\) 的 \textbf{Killing 型}是对称双线性型

\[\kappa(X, Y) = \operatorname{tr}(\operatorname{ad}(X) \circ \operatorname{ad}(Y))\]

\textbf{定理 120.1}（Cartan 判别法）：\(\mathfrak{g}\) 是半单的当且仅当 Killing 型非退化。\(\mathfrak{g}\) 是可解的当且仅当 \(\kappa(\mathfrak{g}, [\mathfrak{g}, \mathfrak{g}]) = 0\)。

\textbf{证明}：设 \(\mathfrak{g}\) 半单。若非退化性不成立，则 \(\mathfrak{a} = \{X \in \mathfrak{g} : \kappa(X, Y) = 0, \forall Y \in \mathfrak{g}\}\) 是非零理想。\(\operatorname{ad}_{\mathfrak{g}}\) 限制到 \(\mathfrak{a}\) 上满足 \(\kappa|_{\mathfrak{a}} = \operatorname{tr}(\operatorname{ad}(X)\operatorname{ad}(Y)) = 0\)。Cartan 第一判别准则：\(\mathfrak{g}\) 可解当且仅当 \(\kappa(\mathfrak{g}, [\mathfrak{g}, \mathfrak{g}]) = 0\)------若该式成立，则每个 \(\operatorname{ad}(X)\) 在 \([\mathfrak{g}, \mathfrak{g}]\) 上幂零，由 Engel 定理知 \(\mathfrak{g}\) 可解。反向：若 \(\mathfrak{g}\) 可解，存在基使 \(\operatorname{ad}(X)\) 为上三角，进而 \(\operatorname{tr}(\operatorname{ad}(X)\operatorname{ad}(Y)) = 0\)（\(\forall X \in \mathfrak{g}, Y \in [\mathfrak{g}, \mathfrak{g}]\)）。Cartan 第二判别准则：\(\mathfrak{g}\) 半单当且仅当 Killing 型非退化------因为非退化性排除了任何非零 Abel 理想（否则该理想包含在根中）。\(\blacksquare\)

\textbf{定义 120.2}（半单李代数）：李代数 \(\mathfrak{g}\) 称为\textbf{半单}的，如果它没有非零 Abel 理想。等价地，\(\mathfrak{g}\) 可分解为单李代数的直和。

\textbf{详解 120.1}（Killing 型与 Cartan 判别准则）：Killing 型 \(\kappa\) 是对称双线性型，由伴随表示的迹定义。其矩阵元为 \(\kappa_{ij} = \sum_{k,l} c_{ik}^l c_{jl}^k\)，其中 \(c_{ij}^k\) 是结构常数。Cartan 第一判别准则：\(\mathfrak{g}\) 可解当且仅当 \(\kappa(\mathfrak{g}, [\mathfrak{g}, \mathfrak{g}]) = 0\)。Cartan 第二判别准则（核心）：\(\mathfrak{g}\) 半单当且仅当 Killing 型非退化（即 \(\det(\kappa_{ij}) \neq 0\)）。此判别法的威力在于将''半单性''这一结构性质转化为纯粹线性代数条件：只需计算 Killing 型矩阵是否满秩。非退化性等价于 \(\mathfrak{g}\) 没有非零 Abel 理想，因为可解理想的 Killing 型必然退化。半单李代数的进一步性质：\(\kappa\) 限制在 Cartan 子代数上仍非退化，从而诱导 \(\mathfrak{h}\) 与 \(\mathfrak{h}^*\) 的同构，为权与根的对偶理论奠定基础。

\textbf{定理 120.2}（Weyl 完全可约性定理）：半单李代数的每个有限维表示是完全可约的（可分解为不可约表示的直和）。

\textbf{证明}：设 \(\rho: \mathfrak{g} \to \mathfrak{gl}(V)\) 是有限维表示，\(W \subseteq V\) 是子表示。需构造 \(\mathfrak{g}\)-不变补空间。利用 Killing 型的非退化性，构造 Casimir 算子 \(C_V = \sum_i X_i X^i \in \operatorname{End}_{\mathfrak{g}}(V)\)（\(\{X_i\}\) 为基，\(\{X^i\}\) 为 Killing-对偶基）。\(C_V\) 与 \(\rho(\mathfrak{g})\) 交换。若 \(\rho\) 忠实，\(C_V\) 的特征空间给出完全可约分解。一般情形通过商 \(\mathfrak{g}/\ker\rho\) 仍为半单归约。另外，紧实形上的 Haar 测度取平均（Weyl 酉技巧）给出 \(G\)-不变内积，进而 \(W^\perp\) 是 \(\mathfrak{g}\)-不变的补空间。\(\blacksquare\)

\subsection{120.2 Cartan 子代数与根系}\label{cartan-ux5b50ux4ee3ux6570ux4e0eux6839ux7cfb}

\textbf{定义 120.3}（Cartan 子代数）：半单李代数 \(\mathfrak{g}\) 的 \textbf{Cartan 子代数} \(\mathfrak{h}\) 是极大 Abel 子代数，且 \(\operatorname{ad}(H)\)（\(H \in \mathfrak{h}\)）可同时对角化。

\textbf{命题 120.1}：半单李代数 \(\mathfrak{g}\) 关于 \(\mathfrak{h}\) 的\textbf{根空间分解}为

\[\mathfrak{g} = \mathfrak{h} \oplus \bigoplus_{\alpha \in \Phi} \mathfrak{g}_\alpha\]

其中 \(\mathfrak{g}_\alpha = \{X \in \mathfrak{g} : [H, X] = \alpha(H) X \text{ 对所有 } H \in \mathfrak{h}\}\)，\(\alpha \in \mathfrak{h}^* \setminus \{0\}\) 称为\textbf{根}。\(\Phi \subseteq \mathfrak{h}^*\) 是根系。

\textbf{命题 120.2}（根系的性质）： 1. \(\Phi\) 有限，张成 \(\mathfrak{h}^*\) 2. 若 \(\alpha \in \Phi\)，则 \(-\alpha \in \Phi\)，且 \(k\alpha \in \Phi\)（\(k \in \mathbb{Z}\)）仅当 \(k = \pm 1\) 3. \(\dim \mathfrak{g}_\alpha = 1\) 4. 对 \(\alpha, \beta \in \Phi\)，\(\beta - \frac{2\kappa(\beta, \alpha)}{\kappa(\alpha, \alpha)} \alpha \in \Phi\)（根系在反射下封闭）

\textbf{定义 120.4}（Weyl 群）：\textbf{Weyl 群} \(W\) 由根反射 \(s_\alpha(\beta) = \beta - \frac{2\kappa(\beta, \alpha)}{\kappa(\alpha, \alpha)} \alpha\) 生成。\(W\) 是有限群，作用在 \(\mathfrak{h}^*\) 上。

\subsection{120.3 单根与 Dynkin 图}\label{ux5355ux6839ux4e0e-dynkin-ux56fe}

\textbf{定义 120.5}（正根与单根）：选取 \(\mathfrak{h}^*\) 中的一个超平面使其不包含任何根，定义\textbf{正根} \(\Phi^+\) 为超平面一侧的根。\textbf{单根} \(\Delta \subseteq \Phi^+\) 是正根中不能分解为两个正根之和的根。

\textbf{命题 120.3}：单根 \(\Delta\) 是 \(\mathfrak{h}^*\) 的一组基。每个正根是单根的非负整数线性组合。

\textbf{定义 120.6}（Cartan 矩阵）：\(A_{ij} = \frac{2\kappa(\alpha_i, \alpha_j)}{\kappa(\alpha_i, \alpha_i)}\)（\(\alpha_i, \alpha_j \in \Delta\)）。\(A_{ij} \in \{0, -1, -2, -3\}\)（\(i \neq j\)）。

\textbf{定义 120.7}（Dynkin 图）：\textbf{Dynkin 图}以单根为节点，\(A_{ij} A_{ji}\) 条边连接节点 \(i\) 和 \(j\)。若 \(|\alpha_i| > |\alpha_j|\)，画箭头指向较短根。

\textbf{定理 120.3}（Cartan-Killing 分类）：复半单李代数（等价地，单李代数）由 Dynkin 图完全分类。连通 Dynkin 图有：

\begin{itemize}
\tightlist
\item
  \(A_n\)（\(n \geq 1\)）：\(\mathfrak{sl}(n+1, \mathbb{C})\)
\item
  \(B_n\)（\(n \geq 2\)）：\(\mathfrak{so}(2n+1, \mathbb{C})\)
\item
  \(C_n\)（\(n \geq 3\)）：\(\mathfrak{sp}(2n, \mathbb{C})\)
\item
  \(D_n\)（\(n \geq 4\)）：\(\mathfrak{so}(2n, \mathbb{C})\)
\item
  例外李代数：\(G_2\)（14 维）、\(F_4\)（52 维）、\(E_6\)（78 维）、\(E_7\)（133 维）、\(E_8\)（248 维）
\end{itemize}

这是数学中最优美的分类定理之一，与 Lie 和 Killing 在 19 世纪末的工作以及 Cartan 的完善。

\textbf{证明}（框架）：步骤一：由 Cartan 子代数与根系理论，每个复半单李代数由 Dynkin 图唯一确定------Dynkin 图编码了 Cartan 矩阵 \(A_{ij} = 2(\alpha_i, \alpha_j)/(\alpha_i, \alpha_i)\)。步骤二：Dynkin 图的组合约束------\(A_{ij} A_{ji} \in \{0, 1, 2, 3\}\)（\(i \neq j\)）意味着仅有特定图出现。步骤三：Serre 关系通过生成元 \(e_i, f_i, h_i\)（\(i \in \Delta\)）和关系 \([h_i, h_j] = 0\)，\([h_i, e_j] = A_{ij} e_j\)，\([h_i, f_j] = -A_{ij} f_j\)，\([e_i, f_j] = \delta_{ij} h_i\)，加上 Serre 关系 \((\operatorname{ad} e_i)^{1-A_{ij}} e_j = 0\) 等，从 Dynkin 图重构出李代数。步骤四：所有满足组合约束的 Dynkin 图恰为 \(A_n, B_n, C_n, D_n, G_2, F_4, E_6, E_7, E_8\)。\(\blacksquare\)

\hrulefill

\hrulefill

\hrulefill

\hrulefill

\hrulefill

\hrulefill

\hrulefill

\hrulefill

\section{第140章：幂零轨道与Kazhdan-Lusztig理论}\label{ux7b2c140ux7ae0ux5e42ux96f6ux8f68ux9053ux4e0ekazhdan-lusztigux7406ux8bba}

Kazhdan-Lusztig 理论（1979）是表示论在过去半个世纪中最深远的突破之一。其原始动机来自表示论中的一个基本问题：给定复半单 Lie 代数 \(\mathfrak{g}\) 和 Weyl 群 \(W\)，如何刻画 Verma 模的合成因子（即不可约最高权表示的 Kazhdan-Lusztig 特征标公式）？在 Kazhdan-Lusztig 之前，这个问题的部分答案由 Jantzen（1979）以 Jantzen 猜想的形式提出，而 Kazhdan-Lusztig 利用 Hecke 代数中的 KL 多项式 \(P_{y, w}(q)\)（\(y, w \in W\)）给出了完全的解答。KL 多项式的组合定义看似简单：它们是 Hecke 代数中标准基 \(\{T_w\}\) 和 Kazhdan-Lusztig 基 \(\{C_w\}\) 之间的变换系数，满足递推关系和 \(P_{w_0 w, w_0} = 1\) 的归一化条件。然而，KL 正性定理（即 \(P_{y,w}(q)\) 的系数均为非负整数）的证明直到数十年后才由 Elias-Williamson（2014）通过 Soergel 双模的 Hodge 理论完成。KL 多项式与幂零轨道的几何------特别是 Springer 解消 \(\tilde{\mathcal{N}} \to \mathcal{N}\) 和 Schubert 簇的奇点理论------有着深刻联系。几何 Satake 对应（Mirkovic-Vilonen 2007）将这些对象统一在仿射 Grassmann 流形的等变交叉上同调的框架中。

\subsection{121.1 Hecke代数与Kazhdan-Lusztig基}\label{heckeux4ee3ux6570ux4e0ekazhdan-lusztigux57fa}

\textbf{定义 121.1}（Hecke代数）：设 \((W, S)\) 为 Coxeter 系（如 Weyl 群配以单反射集）。\(\mathcal{H} = \mathcal{H}(W, S)\) 为 \(W\) 的 \textbf{Hecke 代数}，它是定义在 Laurent 多项式环 \(\mathbb{Z}[q^{\pm 1/2}]\) 上的结合代数，由 \(\{T_w\}_{w \in W}\) 生成，满足：

\begin{enumerate}
\def\labelenumi{\arabic{enumi}.}
\tightlist
\item
  \(T_w T_s = T_{ws}\)（若 \(\ell(ws) > \ell(w)\)）
\item
  \((T_s + 1)(T_s - q) = 0\)（二次关系，等价于 \(T_s^2 = (q - 1)T_s + q T_e\)）
\end{enumerate}

其中 \(\ell(w)\) 为 \(w\) 的 Coxeter 长度。当 \(q = 1\) 时，\(\mathcal{H}\) 退化为群代数 \(\mathbb{Z}[W]\)。

\textbf{定义 121.2}（Kazhdan-Lusztig 基与 KL 多项式）：在 \(\mathcal{H}\) 上定义对合 \(\overline{T_w} = T_{w^{-1}}^{-1}\)，\(\overline{q} = q^{-1}\)。\textbf{Kazhdan-Lusztig 基} \(\{C_w\}_{w \in W}\) 是满足以下条件的唯一基：

\begin{enumerate}
\def\labelenumi{\arabic{enumi}.}
\tightlist
\item
  \(\overline{C_w} = C_w\)（自对偶性）
\item
  \(C_w = q^{-\ell(w)/2} \sum_{y \leq w} P_{y, w}(q) T_y\)，其中 \(P_{y,w}(q) \in \mathbb{Z}[q]\) 且 \(P_{w,w}(q) = 1\)，\(P_{y,w}(q)\) 的次数 \(\leq (\ell(w) - \ell(y) - 1)/2\)（\(y < w\)）
\end{enumerate}

\(P_{y,w}(q)\) 称为 \textbf{Kazhdan-Lusztig 多项式}（在 Bruhat 序 \(y \leq w\) 上定义）。

\textbf{定理 121.1}（KL 多项式的组合递推）：KL 多项式由以下递推公式唯一确定：对 \(y < w\)，选取 \(s \in S\) 满足 \(ws < w\)（即 \(w\) 的下降边）。则

\[P_{y, w}(q) = q^{1 - c} P_{y, ws}(q) + q^c P_{ys, w}(q) - \sum_{z: zs < z < ys} \mu(z, ys) q^{(\ell(w) - \ell(z))/2} P_{z, w}(q)\]

其中 \(c = 1\) 若 \(ys < y\)，否则 \(c = 0\)；\(\mu(z, ys)\) 为 \(P_{z, ys}(q)\) 中 \(q\) 最高次项的系数（在 \(q^{(\ell(ys) - \ell(z) - 1)/2}\) 处的系数）。

\textbf{证明}：由 \(C_w\) 的自对偶性 \(\overline{C_w} = C_w\) 和 \(C_s C_w = (q^{1/2} + q^{-1/2}) C_w\)（若 \(sw < w\)）或 \(C_s C_w = C_{sw} + \sum_{z: sz < z < w} \mu(z, w) C_z\)（若 \(sw > w\)）。展开 \(C_s C_w = \overline{C_s C_w}\) 并比较 \(C_y\) 的系数即得递推关系。该递推的基底是 \(P_{w,w}(q) = 1\)（所有 \(w\)）。\(\blacksquare\)

\subsection{121.2 Kazhdan-Lusztig猜想与证明}\label{kazhdan-lusztigux731cux60f3ux4e0eux8bc1ux660e}

\textbf{定理 121.2}（Kazhdan-Lusztig猜想，1979，由 Beilinson-Bernstein 1981 和 Brylinski-Kashiwara 1981 独立证明）：设 \(\mathfrak{g}\) 为复半单 Lie 代数，\(\rho\) 为所有正根之和的一半。对支配整权 \(\lambda\) 和 Weyl 群元 \(w \in W\)，Verma 模 \(M(w \cdot \lambda)\)（其中 \(w \cdot \lambda = w(\lambda + \rho) - \rho\) 为点移作用 / dot-action）的不可约商 \(L(w \cdot \lambda)\) 的合成因子中，\(L(y \cdot \lambda)\) 的重数等于 \(P_{y, w}(1)\)（KL 多项式在 \(q=1\) 处的值）。即

\[[M(w \cdot \lambda) : L(y \cdot \lambda)] = P_{y,w}(1)\]

特别地，Verma 模的合成重数由 Weyl 群的组合数据（KL 多项式）完全确定，不依赖于权 \(\lambda\) 的具体值（只要 \(\lambda\) 是支配整权）。

\textbf{证明}（几何方法，Beilinson-Bernstein）：核心构造是正则 \(\mathcal{D}\)-模的 Riemann-Hilbert 对应。令 \(G\) 为 \(\mathfrak{g}\) 对应的单连通代数群，\(B\) 为 Borel 子群，\(X = G/B\) 为旗流形（flag variety）。考虑 Schubert 簇的闭包 \(\overline{X_w} = \overline{B w B / B}\)（\(w \in W\)）。Kazhdan-Lusztig 多项式在几何中对应为：\(P_{y,w}(q) = \sum_i \dim \mathcal{H}^{2i}(\operatorname{IC}(\overline{X_w})_y) q^i\)，即 \(\overline{X_w}\) 的交叉上同调（intersection cohomology）在点 \(yB\)（\(y \leq w\)）处的局部 Euler 特征。Beilinson-Bernstein 局部化定理将 \(\mathfrak{g}\)-模的范畴等价于 \(X\) 上某些 \(\mathcal{D}\)-模的范畴。在此等价下，Verma 模对应于标准 \(\mathcal{D}\)-模（由 Schubert 簇的推送构造），不可约模对应于中间扩展（Goresky-MacPherson 扩展）\(\mathcal{D}\)-模。分解定理（BBD）给出标准层被交叉上同调层的分解，其系数恰为 KL 多项式在 \(q=1\) 的值。Kashiwara-Brylinski 的证明使用 \(\mathcal{D}\)-模的 Riemann-Hilbert 对应和 holonomic 模的分解。\(\blacksquare\)

\subsection{121.3 Soergel双模与KL正性定理}\label{soergelux53ccux6a21ux4e0eklux6b63ux6027ux5b9aux7406}

\textbf{定义 121.3}（Soergel双模与 Hodge 理论）：Soergel（1992）的范畴化（categorification）方案将 Hecke 代数 \(\mathcal{H}\) 提升为 Soergel 双模范畴 \(\mathcal{S}\)：Soergel 双模是 \(R = S(\mathfrak{h}^*)\)（Cartan 子代数的对称代数）上的分次 \(R\)-\(R\)-双模，由 Bott-Samelson 解消生成。\(\mathcal{S}\) 的 Grothendieck 群（带分次迹）同构于 Hecke 代数，且不可约 Soergel 双模（在适当的分次意义下）对应于 KL 基元。

\textbf{定理 121.3}（KL 正性定理，Elias-Williamson 2014）：对所有 \(y, w \in W\)，Kazhdan-Lusztig 多项式 \(P_{y,w}(q)\) 的系数为非负整数：

\[P_{y,w}(q) = \sum_{i \geq 0} \dim \operatorname{Hom}(B_y, B_w(i)) \, q^i \in \mathbb{N}[q]\]

其中 \(B_w\) 为 \(w\) 对应的 Soergel 双模，\(B_w(i)\) 表示分次平移。Elias-Williamson 通过建立 Soergel 双模上的 Hodge 理论（正特征约化 + 分解定理）证明了正性。这是表示论中''代数的正性来自几何''这一普遍原则的又一力证。

\textbf{证明}（框架）：Elias-Williamson 的方法分三步：（i）构造 Soergel 双模在正特征域上的版本，其中 Soergel 猜想（每个 Soergel 双模的指标分次自同态环为多项式代数）成立。（ii）利用 Soergel-Williamson 的 Hodge 理论：Soergel 双模的 \(\operatorname{Hom}\) 空间上存在 Lefschetz 算子和本原分解。（iii）正性由 Soergel 双模在特征零的分解数（由 KL 多项式给出）与正特征下不可约对象的重数之间的相容性得出。这等价于''特征零的分解数等于正特征某对象的 Poincare 多项式''，从而非负。关键在于构造 Soergel 双模范畴上的''分解定理''（将半单对象分解为不可约对象的直和），这需要范畴的 Krull-Schmidt 性质和适当的 Hom 空间有限维性假设。\(\blacksquare\)

\hrulefill

\hrulefill

\hrulefill

\section{第141章：仿射李代数与Macdonald恒等式}\label{ux7b2c141ux7ae0ux4effux5c04ux674eux4ee3ux6570ux4e0emacdonaldux6052ux7b49ux5f0f}

仿射李代数是有限维单李代数的无穷维推广，其理论由 Kac（1968）和 Moody（1968）独立建立。仿射李代数 \(\widehat{\mathfrak{g}}\) 由有限维单李代数 \(\mathfrak{g}\) 通过''圈扩张''（loop extension）加中心扩张构造：\(\widehat{\mathfrak{g}} = \mathfrak{g} \otimes \mathbb{C}[t, t^{-1}] \oplus \mathbb{C} K \oplus \mathbb{C} d\)，其中 \(K\) 为中心元，\(d\) 为导子（能量算子）。仿射李代数的表示论在数学和物理中均有极为重要的地位：在数学中，Macdonald 恒等式（1972）将仿射根系的分拆函数与 Dedekind \(\eta\)-函数的模形式性质联系起来，是仿射李代数特征标公式最璀璨的推论；在物理中，仿射 Kac-Moody 代数是二维共形场论中 WZW 模型的手征代数，水平 \(k\) 对应于中心荷 \(c = k \dim \mathfrak{g} / (k + h^\vee)\)（\(h^\vee\) 为对偶 Coxeter 数）。Macdonald 恒等式是数学中最优美的''组合 = 分析''等式之一，其最著名的特殊情形包括 Jacobi 三重积恒等式（对应 \(\widehat{\mathfrak{sl}}_2\)）和 Euler 五边形数定理。

\subsection{122.1 仿射李代数的定义与基本表示}\label{ux4effux5c04ux674eux4ee3ux6570ux7684ux5b9aux4e49ux4e0eux57faux672cux8868ux793a}

\textbf{定义 122.1}（仿射李代数的构造）：设 \(\mathfrak{g}\) 为复单李代数，带有 Killing 型规范化不变双线性型 \((\cdot, \cdot)\)。\(\mathfrak{g}\) 的（非扭）\textbf{仿射李代数} \(\widehat{\mathfrak{g}}\) 定义为

\[\widehat{\mathfrak{g}} = \mathfrak{g} \otimes \mathbb{C}[t, t^{-1}] \oplus \mathbb{C} K \oplus \mathbb{C} d\]

Lie 括号为： \[[X \otimes t^m, Y \otimes t^n] = [X, Y] \otimes t^{m+n} + m \delta_{m+n, 0} (X, Y) K\] \[[K, \widehat{\mathfrak{g}}] = 0\] \[[d, X \otimes t^n] = n X \otimes t^n\]

\(K\) 为\textbf{中心元}，\(d\) 为\textbf{导子}（度算子）。仿射李代数为无穷维 Lie 代数，其 Cartan 子代数为 \(\widehat{\mathfrak{h}} = \mathfrak{h} \oplus \mathbb{C} K \oplus \mathbb{C} d\)（\(\mathfrak{h}\) 为 \(\mathfrak{g}\) 的 Cartan 子代数）。

\textbf{定义 122.2}（最高权表示与可积表示）：仿射李代数 \(\widehat{\mathfrak{g}}\) 的\textbf{最高权表示}由最高权 \(\Lambda \in \widehat{\mathfrak{h}}^*\) 确定，满足 \(e_\alpha |\Lambda\rangle = 0\)（对所有正根 \(\alpha\)）和 \(X \otimes t^n |\Lambda\rangle = 0\)（\(n > 0\)）。权 \(\Lambda\) 可写为 \(\Lambda = \bar{\lambda} + k \Lambda_0 + n \delta\)，其中 \(\bar{\lambda}\) 为 \(\mathfrak{g}\) 的支配整权，\(k\) 为\textbf{水平}（level，\(K\) 的本征值），\(\Lambda_0\) 为基本权（\(\Lambda_0(K) = 1\)，\(\Lambda_0(\mathfrak{h}) = \Lambda_0(d) = 0\)），\(\delta\) 为虚根（\(\delta(d) = 1\)，\(\delta(\mathfrak{h}) = \delta(K) = 0\)）。当 \(k \in \mathbb{Z}_{\geq 0}\) 且权满足特定可积条件时，表示是\textbf{酉}和\textbf{可积}的。

\textbf{定理 122.1}（仿射Weyl-Kac特征标公式，Kac 1974）：设 \(L(\Lambda)\) 为 \(\widehat{\mathfrak{g}}\) 的具有最高权 \(\Lambda\) 的可积最高权表示。则 \(L(\Lambda)\) 的形式特征标（在 Cartan 子代数上的 Fourier 展开）为

\[\operatorname{ch} L(\Lambda) = \frac{\sum_{w \in \widehat{W}} \varepsilon(w) e^{w(\Lambda + \widehat{\rho}) - \widehat{\rho}}}{\prod_{\alpha \in \widehat{\Delta}_+} (1 - e^{-\alpha})^{\operatorname{mult} \alpha}}\]

其中 \(\widehat{W}\) 为仿射 Weyl 群（\(W \ltimes Q^\vee\)，\(Q^\vee\) 为余根格），\(\widehat{\rho}\) 为仿射 Weyl 向量，\(\widehat{\Delta}_+\) 为仿射正根集（包含 \(\mathfrak{g}\) 的正根 + \(n\delta + \alpha\)（\(\alpha \in \mathfrak{h}^*\)，\(n > 0\)）+ \(n\delta\)（\(n > 0\)，虚根，重数 = \(\operatorname{rank} \mathfrak{g}\)）。

\textbf{证明}：平行于有限维 Weyl 特征标公式的证明，使用 Verma 模的 Bernstein-Gel'fand-Gel'fand（BGG）解消。仿射 Verma 模 \(M(\Lambda)\) 的特征标为 \(e^\Lambda / \prod_{\alpha \in \widehat{\Delta}_+} (1 - e^{-\alpha})^{\operatorname{mult}}\)。仿射 BGG 解消：存在正合列 \(\cdots \to \bigoplus_{w: \ell(w)=2} M(w \cdot \Lambda) \to \bigoplus_{w: \ell(w)=1} M(w \cdot \Lambda) \to M(\Lambda) \to L(\Lambda) \to 0\)（\(w \cdot \Lambda = w(\Lambda + \widehat{\rho}) - \widehat{\rho}\) 为仿射 dot-action）。取特征标的交错和，分子中的 Weyl 分母相消，即得公式。\(\blacksquare\)

\subsection{122.2 Macdonald恒等式}\label{macdonaldux6052ux7b49ux5f0f}

\textbf{定理 122.2}（Macdonald恒等式，1972）：设 \(\mathfrak{g}\) 为复单李代数，\(\widehat{\mathfrak{g}}\) 为其仿射李代数。将仿射 Weyl 分母恒等式（特征标公式分母的展开式）

\[\prod_{\alpha \in \widehat{\Delta}_+} (1 - e^{-\alpha})^{\operatorname{mult} \alpha} = \sum_{w \in \widehat{W}} \varepsilon(w) e^{w(\widehat{\rho}) - \widehat{\rho}}\]

对 \(e^{-\delta} = q\) 展开，得到模形式的恒等式。对于 \(\widehat{\mathfrak{sl}}_2\)（对应 \(\mathfrak{g} = \mathfrak{sl}_2\)），Macdonald 恒等式退化为 \textbf{Jacobi 三重积恒等式}：

\[\prod_{n=1}^\infty (1 - q^{2n})(1 + q^{2n-1} z)(1 + q^{2n-1} z^{-1}) = \sum_{m \in \mathbb{Z}} q^{m^2} z^m\]

对于 \(\widehat{\mathfrak{sl}}_n\)，恒等式给出 Dedekind \(\eta\)-函数的幂与 \(\theta\)-函数的乘积关系。

\textbf{证明}（以 \(\widehat{\mathfrak{sl}}_2\) 为例）：\(\mathfrak{g} = \mathfrak{sl}_2\) 的根为 \(\pm \alpha\)（\(\alpha\) 为单根）。仿射正根为：\(n\delta\)（\(n \geq 1\)，重数 1），\(n\delta + \alpha\)（\(n \geq 0\)，重数 1），\(n\delta - \alpha\)（\(n \geq 1\)，重数 1）。故分母为 \[\prod_{n \geq 1} (1 - q^n) \cdot \prod_{n \geq 0} (1 - q^n e^{-\alpha}) \cdot \prod_{n \geq 1} (1 - q^n e^{\alpha})\] 其中 \(q = e^{-\delta}\)，\(z = e^{-\alpha}\)。仿射 Weyl 群为 \(\widehat{W} = \langle s_0, s_1 \rangle\)（无限二面体群），\(\widehat{\rho} = \rho + h^\vee \Lambda_0\)（\(h^\vee = 2\) 为对偶 Coxeter 数）。计算 \(\sum_{w \in \widehat{W}} \varepsilon(w) e^{w(\widehat{\rho}) - \widehat{\rho}}\)：仿射 Weyl 群元为反射的交替乘积，按长度交错求和，等价于在半整数格上的交错 Gauss 和。经过指数化简，得 \(\sum_{m \in \mathbb{Z}} q^{m^2} z^m\)。分母与分子的相等即 Jacobi 三重积恒等式。对于一般 \(\mathfrak{g}\)，类似的格求和给出多重 \(\theta\)-函数。\(\blacksquare\)

\subsection{122.3 Macdonald恒等式的推广与应用}\label{macdonaldux6052ux7b49ux5f0fux7684ux63a8ux5e7fux4e0eux5e94ux7528}

\textbf{定理 122.3}（Macdonald型恒等式与模形式）：对于任意仿射 Lie 代数 \(\widehat{\mathfrak{g}}\)，将特殊化 \(e^{-\delta} = q\)、\(e^{\alpha_i} = 1\)（对所有单根 \(\alpha_i\)）应用到仿射 Weyl 分母恒等式，得到 Dedekind \(\eta\)-函数的乘积公式：

\[\eta(q)^{\dim \mathfrak{g}} = q^{(\dim \mathfrak{g})/24} \cdot \prod_{\alpha \in \Delta_+} \cdots\]

更一般地，Macdonald恒等式可推广到所有 Kac-Moody 代数（包括双曲型），生成函数给出 Siegel 模形式、Jacobi 形式和仿射 Weyl 群的 Eisenstein 级数之间的恒等式。

\textbf{证明}（框架）：特殊化 \(e^{\alpha_i} = 1\)（所有单根）对应''平坦方向''，使得仿射 Weyl 群退化为余根格 \(Q^\vee\)（因 \(\widehat{W} = W \ltimes Q^\vee\)，\(W\) 在 \(\alpha_i = 0\) 时平凡作用）。此时 Liouville 求和退化为 \(Q^\vee\) 上的 Poisson 求和公式，即 \(\sum_{\gamma \in Q^\vee} e^{-\pi i \tau \|\gamma\|^2}\)（\(\tau\) 为模参数），恰为 Siegel \(\theta\)-函数。分母的乘积经 q-Pochhammer 符号展开后，得到 \(\eta\)-函数的乘幂。结合 Poisson 求和公式，Macdonald 恒等式等价于 \(\eta\)-函数的模变换性质。推广至 hyperbolic 型需要处理不定二次型的 Siegel-Weil 公式。\(\blacksquare\)

\hrulefill

\subsection{88.A 李群补充}\label{a-ux674eux7fa4ux8865ux5145}

\begin{quote}
表示论是研究代数结构（群、代数、李代数等）通过线性作用的具体实现的理论。V12 第 69 章已介绍了有限群表示论的基础，V23 涵盖了李群与李代数的表示论。本卷在此基础上深化表示论：有限群模表示论研究特征为素数的域上的表示结构；代数群表示论将群表示论推广到代数几何框架；几何表示论通过层论和相交上同调为表示论提供几何构造；仿射李代数与量子群表示论连接无限维李代数和可积系统；Langlands 纲领深化则从表示论角度揭示数论与调和分析之间的深刻联系。
\end{quote}

\subsection{表示论深化的整体图景}\label{ux8868ux793aux8bbaux6df1ux5316ux7684ux6574ux4f53ux56feux666f}

表示论是现代数学中最具统一性的理论之一，它的核心思想------将抽象的代数结构具体化为线性变换------贯穿了从有限群到量子群的整个数学体系。本卷介绍的五个深化方向分别对应表示论在不同历史阶段和不同数学分支中的发展。\textbf{有限群模表示论}（Brauer 1930s-1940s）起源于 Frobenius 对群行列式的研究，在特征整除群阶的域上，Maschke 定理不再成立------群代数不再是半单的，表示论出现了全新的现象：不可约表示数目少于共轭类数，块分解（block decomposition）将群代数分解为不可分解双理想的直和。Brauer 的块论通过亏群（defect group）和 Brauer 对应将群的整体表示论与其 \(p\)-局部子群联系起来，与有限单群分类（CFSG）中的局部-整体原则有深刻共鸣。\textbf{代数群表示论}（Chevalley 1950s, Borel-Weil-Bott 1950s-1960s）将有限群的经典结果推广到代数闭域上的线性代数群，最高权理论（Cartan-Weyl）和 Weyl 特征标公式是这一理论的皇冠。\textbf{几何表示论}（Springer 1970s, Lusztig 1980s, Beilinson-Bernstein 1981）则是 20 世纪末的一场革命------它通过旗簇、仿射 Grassmann 流形和箭图簇的拓扑和层论来构造表示，将表示论从代数领域拓展到几何领域。\textbf{仿射 Lie 代数与量子群表示论}（Kac 1960s, Drinfeld-Jimbo 1985）连接了无限维李代数和可积系统，而顶点算子代数（Borcherds, Frenkel-Lepowsky-Meurman 1980s）则为共形场论和月光猜想（Moonshine）提供了严格数学框架。\textbf{Langlands 纲领深化}则将上述所有理论统一于一个宏大的猜想体系------从局部 Langlands 对应到几何 Langlands 纲领，表示论与数论、几何和物理的边界被彻底打破。

\hrulefill

\hrulefill

\newpage

\chapter{05-代数/04-李理论/02-半单与表示.md}\label{ux4ee3ux657004-ux674eux7406ux8bba02-ux534aux5355ux4e0eux8868ux793a.md}

\section{第142章：紧李群的表示}\label{ux7b2c142ux7ae0ux7d27ux674eux7fa4ux7684ux8868ux793a}

紧李群具有优美的表示论，其所有不可约表示都是有限维的，且完全由最高权分类。本章讨论 Peter-Weyl 定理和最高权理论。

\subsection{121.1 紧李群上的 Haar 测度}\label{ux7d27ux674eux7fa4ux4e0aux7684-haar-ux6d4bux5ea6}

\textbf{定理 121.1}（Haar 测度）：紧李群 \(G\) 上存在唯一的（左）不变概率测度 \(\mu\)（Haar 测度），满足 \(\mu(gA) = \mu(A)\) 对所有 \(g \in G\) 和可测集 \(A\)。紧李群上 Haar 测度是双不变的。

\textbf{证明}：在 \(G\) 上取一个 \(n\)-形式（体积形式）\(\omega_e \in \bigwedge^n T_e^* G\)。通过左平移定义 \(\omega_g = (L_{g^{-1}})^* \omega_e\)，得到 \(G\) 上整体非零左不变 \(n\)-形式 \(\omega\)。定义正线性泛函 \(I(f) = \int_G f \omega\)（对 \(f \in C(G)\) 依紧支集）。由 Riesz 表示定理，\(I\) 对应唯一概率测度 \(\mu\)。左不变性由 \(\omega\) 的构造保证。紧致性确保 \(\mu(G) < \infty\)，可归一化。双不变性来自：紧群上任何左不变测度自动右不变（否则左乘模函数非平凡将与紧致性矛盾）。\(\blacksquare\)

\textbf{命题 121.1}（Weyl 积分公式）：对 \(G\) 上的类函数 \(f\)（\(f(gxg^{-1}) = f(x)\)），

\[\int_G f(g) dg = \frac{1}{|W|} \int_T f(t) \prod_{\alpha \in \Phi^+} |e^{\alpha}(t) - e^{-\alpha}(t)|^2 dt\]

其中 \(T\) 是极大环面，\(W\) 是 Weyl 群。

\subsection{121.2 Peter-Weyl 定理}\label{peter-weyl-ux5b9aux7406}

\textbf{定义 121.1}（表示）：李群 \(G\) 的（有限维）\textbf{表示}是李群同态 \(\pi: G \to \operatorname{GL}(V)\)（\(V\) 是有限维向量空间）。

\textbf{定理 121.2}（Peter-Weyl 定理）：紧李群 \(G\) 的矩阵系数（\(\pi_{ij}\) 跑遍所有不可约表示）在 \(L^2(G)\) 中稠密。即

\[L^2(G) \cong \widehat{\bigoplus}_{\pi \in \widehat{G}} V_\pi \otimes V_\pi^*\]

（Hilbert 空间直和）。每个不可约表示在 \(L^2(G)\) 中出现的重数等于其维数。

\textbf{证明}：考虑 \(G\) 在 \(L^2(G)\) 上的左右正则表示。对每个不可约表示 \((\pi, V_\pi)\)，矩阵系数映射 \(V_\pi \otimes V_\pi^* \to L^2(G)\) 由 \(v \otimes \xi \mapsto (g \mapsto \xi(\pi(g^{-1})v))\) 给出，它是 \(G \times G\)-等变嵌入。由 Schur 正交性，不同 \(\pi\) 的矩阵系数正交。由 Stone-Weierstrass 定理，所有矩阵系数的代数在 \(C(G)\) 中稠密（生成含 \(1\) 且分离点的自伴子代数），故在 \(L^2(G)\) 中稠密。维数重数来自 \(L^2(G)\) 中 \(V_\pi\) 的同型分支的结构。\(\blacksquare\)

\textbf{推论 121.3}：紧李群的所有不可约表示都是有限维的。紧李群有可数多个不可约表示（同构类）。

\subsection{121.3 最高权理论}\label{ux6700ux9ad8ux6743ux7406ux8bba-1}

\textbf{定义 121.2}（权）：设 \(\pi: G \to \operatorname{GL}(V)\) 是 \(G\) 的表示。\(\lambda \in \mathfrak{h}^*\) 称为表示的\textbf{权}，如果存在非零 \(v \in V\) 使得 \(\pi(H)v = \lambda(H) v\) 对所有 \(H \in \mathfrak{h}\)。权空间 \(V_\lambda\) 是所有 \(\lambda\)-权向量的集合。

\textbf{定理 121.4}（最高权定理）：紧连通李群 \(G\) 的不可约表示与\textbf{支配整权}（dominated integral weights）一一对应。每个不可约表示由唯一的\textbf{最高权} \(\lambda\) 确定（所有其他权 \(\mu\) 满足 \(\lambda - \mu\) 是正根的非负整数组合）。

\textbf{证明}：由 Cartan-Weyl 理论，固定 Cartan 子代数 \(\mathfrak{h}\) 和正根集 \(\Phi^+\)。对不可约表示 \(V\)，取其中关于 \(\mathfrak{h}\) 的权中按字典序极大的权 \(\lambda\)。由于 \(\mathfrak{n}^+ = \bigoplus_{\alpha > 0} \mathfrak{g}_\alpha\) 上的根向量将权递增，\(\lambda\)-权向量被 \(\mathfrak{n}^+\) 零化。由表示论的标准论证，\(\lambda\) 必须是支配整权（对每个正根 \(\alpha\)，\(\langle \lambda, \alpha^\vee \rangle \in \mathbb{Z}_{\geq 0}\)）。反之，给定支配整权 \(\lambda\)，构造 Verma 模 \(M(\lambda)\) 并取其唯一不可约商 \(L(\lambda)\)------它具有最高权 \(\lambda\)。紧性确保 \(L(\lambda)\) 是有限维的。\(\blacksquare\)

\textbf{定义 121.3}（支配整权）：\(\lambda \in \mathfrak{h}^*\) 是\textbf{支配整权}，如果 \(\frac{2\kappa(\lambda, \alpha)}{\kappa(\alpha, \alpha)} \in \mathbb{Z}_{\geq 0}\) 对所有正根 \(\alpha\)。

\textbf{详解 121.1}（最高权理论的基本框架）：最高权理论建立了支配整权与不可约表示之间的一一对应，其结构如下。第一步：选取 Cartan 子代数 \(\mathfrak{h}\) 和正根系 \(\Phi^+\)，定义支配整权的集合 \(P_+ = \{\lambda \in \mathfrak{h}^* : \langle \lambda, \alpha^\vee \rangle \in \mathbb{Z}_{\geq 0}, \forall \alpha \in \Phi^+\}\)。第二步：对每个 \(\lambda \in P_+\)，构造 Verma 模 \(M(\lambda) = U(\mathfrak{g}) \otimes_{U(\mathfrak{b})} \mathbb{C}_\lambda\)（其中 \(\mathfrak{b} = \mathfrak{h} \oplus \bigoplus_{\alpha > 0} \mathfrak{g}_\alpha\) 是 Borel 子代数），这是一个最高权为 \(\lambda\) 的无穷维模。第三步：\(M(\lambda)\) 有唯一的极大真子模，商模 \(L(\lambda)\) 是唯一的最高权为 \(\lambda\) 的不可约模。第四步：Weyl 完全可约性定理保证了每个有限维表示分解为 \(L(\lambda)\) 的直和。此框架将李群表示论转化为权格的组合学，是现代表示论的标准范式。

\textbf{定理 121.5}（Weyl 维数公式）：最高权为 \(\lambda\) 的不可约表示的维数为

\[\dim V_\lambda = \prod_{\alpha \in \Phi^+} \frac{\kappa(\lambda + \rho, \alpha)}{\kappa(\rho, \alpha)}\]

其中 \(\rho = \frac{1}{2} \sum_{\alpha \in \Phi^+} \alpha\)（所有正根之和的一半）。

\textbf{证明}：由 Weyl 特征公式（定理 121.6），\(\dim V_\lambda = \chi_\lambda(e)\)。在特征公式中令所有 \(e^{\mu} \to 1\) 求极限。利用 L'Hôpital 法则或对 \(e^{\mu}\) 取 \(\mu \to 0\) 的 Taylor 展开：分子分母的零点均来自 \(\sinh(\kappa(\cdot, \alpha)/2)\) 因子。简化后得 \(\dim V_\lambda = \prod_{\alpha > 0} \frac{\langle \lambda + \rho, \alpha^\vee \rangle}{\langle \rho, \alpha^\vee \rangle}\)，其中 \(\alpha^\vee = 2\alpha/\kappa(\alpha, \alpha)\)，等价于所述公式。\(\blacksquare\)

\textbf{定理 121.6}（Weyl 特征公式）：最高权 \(\lambda\) 的不可约表示的特征标为

\[\chi_\lambda = \frac{\sum_{w \in W} (-1)^{\ell(w)} e^{w(\lambda + \rho)}}{\sum_{w \in W} (-1)^{\ell(w)} e^{w(\rho)}}\]

其中 \(\ell(w)\) 是 \(w\) 的长度（表示为单反射的最小个数）。

\textbf{证明}（Weyl-Kac 方法）：考虑 Verma 模 \(M(\lambda)\)，其特征标为 \(\operatorname{ch} M(\lambda) = e^\lambda / \prod_{\alpha > 0} (1 - e^{-\alpha})\)。不可约模 \(L(\lambda)\) 是 \(M(\lambda)\) 的唯一不可约商，且 \(M(\mu)\)（\(\mu \leq \lambda\)）的复合因子出现在 \(L(\lambda)\) 的分解中。利用 Weyl 群在特征标上的作用：\(w(\operatorname{ch} M(\lambda)) = (-1)^{\ell(w)} \operatorname{ch} M(w \cdot \lambda)\)（其中 \(w \cdot \lambda = w(\lambda + \rho) - \rho\)）。由 Euler-Poincaré 原理，\(\operatorname{ch} L(\lambda) = \sum_{w \in W} (-1)^{\ell(w)} \operatorname{ch} M(w \cdot \lambda)\)。代入 Verma 模特征标并通分分母，利用 Weyl 分母公式 \(\sum_{w \in W} (-1)^{\ell(w)} e^{w(\rho)} = e^\rho \prod_{\alpha > 0} (1 - e^{-\alpha})\) 即得。\(\blacksquare\)

\hrulefill

\hrulefill

\hrulefill

\hrulefill

\hrulefill

\hrulefill

\section{第143章：李群在几何中的应用}\label{ux7b2c143ux7ae0ux674eux7fa4ux5728ux51e0ux4f55ux4e2dux7684ux5e94ux7528}

李群在几何中有深刻的应用，包括齐性空间、对称空间和主丛上的联络理论。

\subsection{122.1 齐性空间}\label{ux9f50ux6027ux7a7aux95f4}

\textbf{定义 122.1}（齐性空间）：李群 \(G\) 的\textbf{齐性空间}是 \(G\) 在流形 \(M\) 上的可迁作用（即 \(M = G/H\)，其中 \(H\) 是 \(G\) 的闭子群）。

齐性空间是李群理论与微分几何交汇的核心概念。当 \(G\) 是连通李群、\(H\) 是 \(G\) 的闭子群时，商空间 \(G/H\) 具有唯一的光滑流形结构使得投影 \(\pi: G \to G/H\) 是光滑浸没，且 \(G\) 在 \(G/H\) 上的左作用 \(g \cdot aH = (ga)H\) 是可迁的。\(G\) 的左作用诱导了 \(G\) 在切空间 \(T_{eH}(G/H) \cong \mathfrak{g}/\mathfrak{h}\) 上的线性表示，称为\textbf{迷向表示}（isotropy representation）。当 \(H\) 是紧子群时，\(G/H\) 上存在 \(G\)-不变的黎曼度量，使得 \(G/H\) 成为黎曼齐性空间。重要例子包括球面 \(S^n \cong SO(n+1)/SO(n)\)、复射影空间 \(\mathbb{CP}^n \cong SU(n+1)/U(n)\) 以及 Grassmann 流形 \(\operatorname{Gr}(k,n) \cong O(n)/(O(k) \times O(n-k))\)。齐性空间的曲率张量完全由其李代数结构决定：若 \(G/H\) 配备 \(G\)-不变度量，其截面曲率可由 \(\mathfrak{g}\) 和 \(\mathfrak{h}\) 上的李括号结构显式计算。在紧李群的情形，齐性空间上的调和分析（即 \(L^2(G/H)\) 的 \(G\)-模分解）由 Peter-Weyl 定理和分支律控制，构成非交换调和分析的基础。

\textbf{定理 122.1}（Maurer-Cartan 形式）：李群 \(G\) 上的 \textbf{Maurer-Cartan 形式} \(\omega \in \Omega^1(G, \mathfrak{g})\) 定义为 \(\omega_g(v) = (L_{g^{-1}})_* v\)。满足 Maurer-Cartan 方程：

\[d\omega + \frac{1}{2}[\omega, \omega] = 0\]

\textbf{证明}：对左不变向量场 \(X, Y\)，计算 \(d\omega(X, Y) = X(\omega(Y)) - Y(\omega(X)) - \omega([X, Y])\)。\(\omega\) 左不变故 \(\omega(X) = X_e\)（常数），前两项为零，故 \(d\omega(X, Y) = -\omega([X, Y]) = -[X_e, Y_e]\)。按定义 \([\omega, \omega](X, Y) = [\omega(X), \omega(Y)] - [\omega(Y), \omega(X)] = 2[X_e, Y_e]\)，故 \(d\omega + \frac{1}{2}[\omega, \omega] = 0\)。\(\blacksquare\)

\subsection{122.2 对称空间}\label{ux5bf9ux79f0ux7a7aux95f4}

\textbf{定义 122.2}（对称空间）：\textbf{对称空间}是黎曼流形 \((M, g)\)，每点 \(p \in M\) 存在等距对合 \(\sigma_p: M \to M\) 以 \(p\) 为孤立不动点。

对称空间的概念由 Elie Cartan 在 1920 年代引入并完成分类，是齐性空间的一个特殊子类，也是黎曼几何中最重要的研究对象之一。条件 \(\sigma_p\) 为对合（\(\sigma_p^2 = \operatorname{id}\)）且 \(p\) 为孤立不动点意味着在 \(p\) 点附近，\(\sigma_p\) 作用为测地线反射，这强行要求 \(M\) 的截面曲率沿测地线为常数，且 \(M\) 本身是完备的（任意测地线可无限延伸）。对称空间自然地分为\textbf{紧型}（截面曲率 \(\geq 0\)）、\textbf{非紧型}（截面曲率 \(\leq 0\)）和\textbf{欧氏型}（截面曲率 \(= 0\)，即平坦空间）。紧型对称空间包括紧李群本身（配以双不变度量的 \(G\) 是对称空间，\(\sigma_g(h) = gh^{-1}g\)）、球面 \(S^n\)、复射影空间 \(\mathbb{CP}^n\)、四元数射影空间 \(\mathbb{HP}^n\) 和 Cayley 射影平面 \(\mathbb{OP}^2\)。非紧型对称空间是紧型的对偶（通过 Cartan 对合），例如双曲空间 \(H^n\) 是 \(S^n\) 的非紧对偶。对称空间的分类与实半单李代数的分类完全等价------每一类对称空间对应一个实半单李代数对合 \((\mathfrak{g}, \Theta)\)，这构成了 Cartan 分类表的核心。

\textbf{定理 122.2}（Cartan 分类）：对称空间由 \((\mathfrak{g}, \mathfrak{h}, \sigma)\) 分类，其中 \(\mathfrak{g}\) 是李代数，\(\mathfrak{h}\) 是对合 \(\sigma\) 的固定点子代数。\(\mathfrak{g} = \mathfrak{h} \oplus \mathfrak{m}\)（\(\mathfrak{m}\) 是 \(\sigma\) 的 \(-1\) 特征空间），满足 \([\mathfrak{h}, \mathfrak{h}] \subseteq \mathfrak{h}\)，\([\mathfrak{h}, \mathfrak{m}] \subseteq \mathfrak{m}\)，\([\mathfrak{m}, \mathfrak{m}] \subseteq \mathfrak{h}\)。

\textbf{证明}：对称空间 \(M\) 上，固定基点 \(o \in M\)，取等距群 \(G = \operatorname{Isom}(M)^0\)。对合 \(\sigma_o\) 诱导李代数对合 \(\Theta: \mathfrak{g} \to \mathfrak{g}\)。取 \(\mathfrak{h} = \ker(\Theta - \operatorname{id})\) 和 \(\mathfrak{m} = \ker(\Theta + \operatorname{id})\)，则 \(\mathfrak{g} = \mathfrak{h} \oplus \mathfrak{m}\)。对合与李括号的相容性直接验证即得三类括号关系：\([\mathfrak{h}, \mathfrak{h}] \subseteq \mathfrak{h}\) 由 \(\mathfrak{h}\) 为子代数保证；\([\mathfrak{h}, \mathfrak{m}] \subseteq \mathfrak{m}\) 来自 \(\Theta([H, X]) = [\Theta(H), \Theta(X)] = [H, -X] = -[H, X]\)；\([\mathfrak{m}, \mathfrak{m}] \subseteq \mathfrak{h}\) 同理。此分解与 \(G/H\) 的齐性空间结构对应。\(\blacksquare\)

\subsection{122.3 主丛与联络}\label{ux4e3bux4e1bux4e0eux8054ux7edc}

\textbf{定义 122.3}（主 \(G\)-丛）：流形 \(M\) 上的\textbf{主 \(G\)-丛} \(P\) 是纤维丛 \(\pi: P \to M\)，纤维是李群 \(G\)，\(G\) 在 \(P\) 上自由右作用，局部平凡化与 \(G\) 作用兼容。

\textbf{定义 122.4}（联络）：\(P\) 上的\textbf{联络}是 \(\mathfrak{g}\)-值 1-形式 \(\omega \in \Omega^1(P, \mathfrak{g})\)，满足： 1. \(\omega(\tilde{X}) = X\)（对 \(X \in \mathfrak{g}\)，\(\tilde{X}\) 是 \(X\) 生成的基本向量场） 2. \(R_g^* \omega = \operatorname{Ad}(g^{-1}) \omega\)（等变性）

\textbf{定义 122.5}（曲率）：联络 \(\omega\) 的\textbf{曲率}为 \(\Omega = d\omega + \frac{1}{2}[\omega, \omega]\)（\(\mathfrak{g}\)-值 2-形式）。

\textbf{定理 122.3}（Cartan 结构方程）：在正交标架丛（或一般主丛）上，联络 \(\omega\) 和曲率 \(\Omega\) 满足 Cartan 结构方程：

\[d\omega = -\frac{1}{2}[\omega, \omega] + \Omega\]

（Maurer-Cartan 方程 \(\Omega = 0\) 是平坦联络的特例）。

\textbf{证明}：曲率定义直接给出 \(d\omega = \Omega - \frac{1}{2}[\omega, \omega]\)。等价地，对水平向量 \(X^h, Y^h\)，\(\Omega(X^h, Y^h) = d\omega(X^h, Y^h)\)。对垂直方向，\(d\omega\) 与 \(-\frac{1}{2}[\omega, \omega]\) 在纤维上相消------纤维上的联络限制为 Maurer-Cartan 形式，满足 \(d\omega + \frac{1}{2}[\omega, \omega] = 0\)，因此垂直分量抵消。综上所述即得 Cartan 结构方程。\(\blacksquare\)

\textbf{定理 122.4}（Chern-Weil 同态）：\(G\)-不变多项式环 \(I(G)\) 到 \(M\) 的 de Rham 上同调环 \(H^*_{\text{dR}}(M)\) 之间存在同态。这给出了示性类的基本构造。

\textbf{证明}：对 \(G\)-不变多项式 \(f \in I^k(G)\)，在底流形 \(M\) 上定义 \(f(\Omega) \in \Omega^{2k}(M)\)（经局部截面拉回曲率形式）。由 Bianchi 恒等式 \(d\Omega = [\Omega, \omega]\) 和 \(f\) 的 \(\operatorname{Ad}\)-不变性可证 \(d f(\Omega) = 0\)，故 \(f(\Omega)\) 是闭形式。若取不同联络 \(\omega'\)，其曲率 \(\Omega'\) 与 \(\Omega\) 相差正合形式，故 \([f(\Omega)] \in H^{2k}_{\text{dR}}(M)\) 不依赖于联络选取。这定义了环同态 \(I(G) \to H^*_{\text{dR}}(M)\)，其像为示性类。\(\blacksquare\)

\hrulefill

\section{第144章：半单李代数的Weyl特征标公式}\label{ux7b2c144ux7ae0ux534aux5355ux674eux4ee3ux6570ux7684weylux7279ux5f81ux6807ux516cux5f0f}

Weyl特征标公式是紧李群表示论的巅峰成就。它给出了不可约表示特征标的行列式闭形表达式，统一了维数公式、重数公式和分支律。本章在前两章的基础上深化：重新给出Weyl特征标公式的完整证明（通过紧李群上的子群积分方法），推导Kostant重数公式（将权空间维数表为Kostant分拆函数的交替和），建立Freudenthal递推公式（计算权重数的实用递推法），并以Casimir算子方法严格证明有限维表示的完全可约性------这是半单李代数表示论中最基本的结构定理，保证每个有限维表示分解为不可约表示的直和。

\subsection{Weyl特征标公式的完整证明}\label{weylux7279ux5f81ux6807ux516cux5f0fux7684ux5b8cux6574ux8bc1ux660e}

\textbf{定理85.1（Weyl特征标公式------子群积分方法）} 设 \(G\) 是紧连通李群，\(T\) 是极大环面，\(W = N_G(T)/T\) 是Weyl群。最高权为 \(\lambda \in P_+\) 的不可约表示 \(V_\lambda\) 的特征标 \(\chi_\lambda\) 在 \(T\) 上的限制为

\[\chi_\lambda(t) = \frac{\sum_{w \in W} (-1)^{\ell(w)} e^{w(\lambda + \rho)}(t)}{\sum_{w \in W} (-1)^{\ell(w)} e^{w(\rho)}(t)}, \quad t \in T\]

其中 \(\rho = \frac{1}{2}\sum_{\alpha \in \Phi^+} \alpha\)，\(\ell(w)\) 为 \(w\) 作为单反射乘积的长度。

\textbf{证明}（Weyl积分公式的表示论应用）

\textbf{第1步（类函数与Weyl积分）}：\(V_\lambda\) 的特征标 \(\chi_\lambda(g) = \operatorname{Tr}(\pi_\lambda(g))\) 是 \(G\) 上的类函数。由命题121.1（Weyl积分公式），对类函数 \(f\)，

\[\int_G f(g) dg = \frac{1}{|W|} \int_T f(t) |\Delta(t)|^2 dt\]

其中 \(\Delta(t) = \prod_{\alpha \in \Phi^+} (e^{\alpha/2}(t) - e^{-\alpha/2}(t)) = \sum_{w \in W} (-1)^{\ell(w)} e^{w(\rho)}(t)\) 是Weyl分母。

\textbf{第2步（特征标的Schur正交性）}：不可约特征标满足 \(\int_G \chi_\lambda(g) \overline{\chi_\mu(g)} dg = \delta_{\lambda\mu}\)。结合Weyl积分公式，

\[\frac{1}{|W|} \int_T \chi_\lambda(t) \overline{\chi_\mu(t)} |\Delta(t)|^2 dt = \delta_{\lambda\mu}\]

\textbf{第3步（轨道和的特征标构造）}：对权 \(\mu \in \mathfrak{h}^*\)，定义交替和

\[A_{\mu + \rho}(t) = \sum_{w \in W} (-1)^{\ell(w)} e^{w(\mu + \rho)}(t)\]

\(A_{\mu + \rho}\) 是 \(T\) 上的Weyl群斜对称函数。由正交关系，不同 \(\mu\) 的 \(A_{\mu + \rho}\) 在 \(L^2(T)\) 中关于测度 \(|\Delta(t)|^2 dt\) 正交。

\textbf{第4步（特征标的确定）}：\(\chi_\lambda\) 可展开为权的有限和：\(\chi_\lambda(t) = \sum_{\mu} m_\lambda(\mu) e^\mu(t)\)。仅最高权 \(\lambda\) 的Weyl群轨道上出现非零系数，且由Weyl群对称性 \(\chi_\lambda(w \cdot t) = \chi_\lambda(t)\)。作为 \(T\) 上的函数，\(\chi_\lambda(t) \Delta(t)\) 是Weyl群斜对称的。由表示论的正交关系和上述构造，\(\chi_\lambda(t) \Delta(t)\) 恰为 \(A_{\lambda + \rho}(t)\)（因为两者的最高权指数项均为 \(e^{\lambda + \rho}\)，系数为 \(1\)，且均为Weyl群斜对称）。故

\[\chi_\lambda(t) = \frac{A_{\lambda + \rho}(t)}{\Delta(t)} = \frac{\sum_{w \in W} (-1)^{\ell(w)} e^{w(\lambda + \rho)}(t)}{\sum_{w \in W} (-1)^{\ell(w)} e^{w(\rho)}(t)}\]

\(\blacksquare\)

\subsection{Kostant重数公式与Freudenthal递推}\label{kostantux91cdux6570ux516cux5f0fux4e0efreudenthalux9012ux63a8}

\textbf{定理85.2（Kostant重数公式）} \(V_\lambda\) 中权 \(\mu\) 的重数 \(m_\lambda(\mu) = \dim(V_\lambda)_\mu\) 由Kostant分拆函数 \(\mathcal{P}\) 给出：

\[m_\lambda(\mu) = \sum_{w \in W} (-1)^{\ell(w)} \mathcal{P}(w(\lambda + \rho) - (\mu + \rho))\]

其中 \(\mathcal{P}(\nu)\) 是将 \(\nu\) 表为正根的非负整数组合的方式数（即 \(\nu = \sum_{\alpha \in \Phi^+} k_\alpha \alpha\)，\(k_\alpha \in \mathbb{Z}_{\geq 0}\) 的解的个数）。

\textbf{证明} 由Weyl特征标公式，\(\chi_\lambda(t) \prod_{\alpha > 0} (1 - e^{-\alpha}(t)) = \sum_{w \in W} (-1)^{\ell(w)} e^{w(\lambda + \rho) - \rho}(t)\)。左边展开为：\(\chi_\lambda(t)\) 乘以 \(\prod_{\alpha > 0} (1 - e^{-\alpha})^{-1}\) 的逆。而 \(\prod_{\alpha > 0} (1 - e^{-\alpha})^{-1} = \sum_{\nu} \mathcal{P}(\nu) e^{-\nu}(t)\)（其中 \(\mathcal{P}(\nu)\) 是Kostant分拆函数）。因此 \(\chi_\lambda(t) = \sum_\mu m_\lambda(\mu) e^\mu\)，且

\[\sum_\mu m_\lambda(\mu) e^\mu \prod_{\alpha > 0} (1 - e^{-\alpha}) = \sum_{w} (-1)^{\ell(w)} e^{w(\lambda + \rho) - \rho}\]

比较 \(e^\mu\) 的系数：右边给出 \(\sum_w (-1)^{\ell(w)} \delta_{w(\lambda + \rho) - \rho, \mu}\)。左边给出 \(\sum_{\nu} m_\lambda(\nu) \sum_{\eta} (-1)^{|\eta|} \delta_{\nu - \eta, \mu}\)，其中 \(\eta\) 是正根的（可能重复的）和。通过Möbius反演（根偏序上的反演）即得Kostant重数公式。\(\blacksquare\)

\textbf{定理85.3（Freudenthal递推公式）} 权 \(\mu\) 的重数 \(m_\lambda(\mu)\) 可通过递推计算：

\[(\|\lambda + \rho\|^2 - \|\mu + \rho\|^2) m_\lambda(\mu) = 2 \sum_{\alpha \in \Phi^+} \sum_{k \geq 1} m_\lambda(\mu + k\alpha) \langle \mu + k\alpha, \alpha \rangle\]

其中 \(\|\nu\|^2 = \langle \nu, \nu \rangle\) 以Killing型的对偶内积计算。这是一个实用的递推算法，从最高权 \(\mu = \lambda\)（\(m_\lambda(\lambda) = 1\)）出发逐步降低权计算重数。

\textbf{证明} 对特征标展开 \(\chi_\lambda = \sum_\mu m_\lambda(\mu) e^\mu\) 应用Casimir算子 \(\Omega\)（见定理85.4）。\(\Omega\) 作用于 \(V_\lambda\) 的标量为 \(\langle \lambda + 2\rho, \lambda \rangle\)，在权空间 \((V_\lambda)_\mu\) 上的作用可通过正根的贡献展开。比较 \(e^\mu\) 的系数得到递推关系。具体地，\(\Omega \cdot e^\mu = \langle \mu, \mu + 2\rho \rangle e^\mu + 2 \sum_{\alpha > 0} e^{\mu + \alpha} \langle \mu + \alpha, \alpha \rangle + \cdots\)。乘以 \(\chi_\lambda\) 并比较系数得Freudenthal公式。\(\blacksquare\)

\subsection{有限维表示的完全可约性}\label{ux6709ux9650ux7ef4ux8868ux793aux7684ux5b8cux5168ux53efux7ea6ux6027}

\textbf{定理85.4（完全可约性------Casimir算子方法）} 设 \(\mathfrak{g}\) 是特征 \(0\) 代数闭域上的半单李代数。则 \(\mathfrak{g}\) 的每个有限维表示是完全可约的------即任何有限维 \(\mathfrak{g}\)-模 \(V\) 可分解为不可约子模的直和：\(V = \bigoplus_i V_i\)（\(V_i\) 不可约）。

\textbf{证明}（Casimir算子方法，Weyl的''酉技巧''的代数版本）

\textbf{第1步（Casimir算子的定义）}：设 \(\kappa\) 是 \(\mathfrak{g}\) 上的Killing型，因 \(\mathfrak{g}\) 半单故 \(\kappa\) 非退化。取 \(\mathfrak{g}\) 的一组基 \(\{x_i\}\)，设 \(\{y_i\}\) 为关于 \(\kappa\) 的对偶基（\(\kappa(x_i, y_j) = \delta_{ij}\)）。定义\textbf{Casimir算子}

\[\Omega = \sum_i x_i y_i \in U(\mathfrak{g})\]

\(\Omega\) 不依赖于基的选取，且属于 \(U(\mathfrak{g})\) 的中心（即与所有 \(\mathfrak{g}\) 的元素交换）。

\textbf{第2步（Casimir算子在表示上的作用）}：设 \((\pi, V)\) 是 \(\mathfrak{g}\) 的有限维表示。将 \(\pi\) 延拓到 \(U(\mathfrak{g})\) 给出 \(\pi(\Omega): V \to V\)。\(\pi(\Omega)\) 与所有 \(\pi(\mathfrak{g})\) 交换（因为 \(\Omega \in Z(U(\mathfrak{g}))\)）。若 \(V\) 是不可约的且 \(\pi\) 非平凡，则 \(\pi(\Omega) = c_\pi \cdot \operatorname{id}_V\)（由Schur引理），其中 \(c_\pi \neq 0\)（对非平凡不可约表示）。

\textbf{第3步（完全可约性的归纳证明）}：对 \(V\) 的维数归纳。设 \(W \subset V\) 是真子模，需证存在补子模 \(W'\) 使 \(V = W \oplus W'\)。

考虑 \(W\) 上非退化的不变双线性型（由Killing型诱导的存在性保证）。若 \(W\) 不可约，则 \(\pi(\Omega)|_W = c_W \cdot \operatorname{id}_W\)（\(c_W \neq 0\)）。\(V\) 可分解为 \(\pi(\Omega)\) 的广义特征空间直和：\(V = \bigoplus_{\lambda} V_{(\lambda)}\)（其中 \(V_{(\lambda)} = \{v : (\pi(\Omega) - \lambda)^N v = 0, N \gg 0\}\)）。因 \(\Omega\) 在 \(U(\mathfrak{g})\) 的中心中，每个 \(V_{(\lambda)}\) 是 \(\mathfrak{g}\)-子模。\(W \subseteq V_{(c_W)}\)。通过对 \(\dim V\) 的归纳（对 \(V/V_{(c_W)}\) 应用归纳假设），可得 \(V = W \oplus W'\) 的分解。

更简洁的论证（Weyl）：在 \(V\) 上定义 \(\mathfrak{g}\)-不变正定Hermite内积（利用紧实形式的Weyl酉技巧------将 \(\mathfrak{g}\) 的紧实形式的作用''酉化''）。对任意子模 \(W\)，其正交补 \(W^\perp\) 也是子模（因表示保持内积），且 \(V = W \oplus W^\perp\)。由归纳法即得完全可约性。\(\blacksquare\)

\hrulefill

\section{第145章：李群表示的Borel-Weil理论}\label{ux7b2c145ux7ae0ux674eux7fa4ux8868ux793aux7684borel-weilux7406ux8bba}

Borel-Weil理论建立了紧李群的不可约表示与旗流形上全纯线丛的全纯截面空间之间的深刻对应，将表示论问题转化为复几何问题。本章的核心是旗流形 \(G/B\)（或更一般地 \(G/P\)）上的几何构造：对每个支配整权 \(\lambda\)，在 \(G/B\) 上存在自然的全纯线丛 \(\mathcal{L}_\lambda\)，其全纯截面空间 \(H^0(G/B, \mathcal{L}_\lambda)\) 作为 \(G\)-模恰为不可约表示 \(V_\lambda\)。这一洞察将最高权理论的组合学置于复流形的整体分析框架中，引导出Demazure特征标公式、Bott-Borel-Weil定理（推广至高阶上同调）及Kostant同调群的丰富理论。

\subsection{旗流形与全纯线丛}\label{ux65d7ux6d41ux5f62ux4e0eux5168ux7eafux7ebfux4e1b}

\textbf{定义86.1（旗流形 \(G/B\)）} 设 \(G_{\mathbb{C}}\) 是连通复半单李群，\(B\) 是Borel子群（极大可解子群）。商空间 \(G_{\mathbb{C}}/B\) 是光滑射影代数簇，称为\textbf{广义旗流形}（generalized flag variety）。在紧实形式 \(G \subset G_{\mathbb{C}}\) 下，\(G/T \cong G_{\mathbb{C}}/B\)（其中 \(T\) 是 \(G\) 的极大环面），故旗流形也可视为 \(G/T\)，一个紧致复流形。

\textbf{定理86.2（Borel-Weil定理）} 设 \(\lambda \in P_+\) 是支配整权，\(\mathcal{L}_\lambda = G_{\mathbb{C}} \times_B \mathbb{C}_{-\lambda}\) 是 \(G_{\mathbb{C}}/B\) 上与 \(\lambda\) 相伴的全纯线丛（其中 \(B\) 通过特征标 \(-\lambda\) 作用于 \(\mathbb{C}\)）。则全纯截面空间 \(H^0(G_{\mathbb{C}}/B, \mathcal{L}_\lambda)\) 是 \(G_{\mathbb{C}}\) 的不可约表示，其最高权恰为 \(\lambda\)。换言之，

\[V_\lambda \cong H^0(G/B, \mathcal{L}_\lambda)^*\]

（或等价地 \(V_\lambda \cong H^0(G/B, \mathcal{L}_{-w_0(\lambda)})\)，取决于符号约定）。

\textbf{证明}（Borel-Weil，1954）

\textbf{第1步（线丛的构造）}：\(\lambda\) 给出 \(B\) 的代数群同态 \(B \to \mathbb{C}^\times\)（因 \(B = T \ltimes U\)，\(U\) 为幂幺根，\(T\) 通过 \(\lambda\) 作用）。\(\mathcal{L}_\lambda\) 是相伴丛，其全纯截面可等同于满足等变条件的全纯函数：

\[H^0(G_{\mathbb{C}}/B, \mathcal{L}_\lambda) \cong \{f \in \mathcal{O}(G_{\mathbb{C}}) : f(gb) = \lambda(b)^{-1} f(g), \forall b \in B\}\]

其中 \(\mathcal{O}(G_{\mathbb{C}})\) 是 \(G_{\mathbb{C}}\) 上的全纯函数空间。

\textbf{第2步（\(B\)-最高权向量的存在性）}：在以上函数空间中，\(G_{\mathbb{C}}\) 通过左平移 \((g \cdot f)(x) = f(g^{-1}x)\) 作用。由Borel的固定点定理（旗流形上的Bruhat分解），\(B\) 在 \(G_{\mathbb{C}}/B\) 上有开稠密轨道（大开胞腔 \(Bw_0B/B\)）。在此轨道上，满足上述等变条件的全纯函数由其在单元 \(e\) 的值完全确定。进一步，\(H^0(G_{\mathbb{C}}/B, \mathcal{L}_\lambda)\) 中含有一个唯一的（至多相差标量）\(B\)-最高权向量------即被 \(B\) 作用于标量 \(\lambda\) 的向量。

\textbf{第3步（不可约性）}：该 \(B\)-最高权向量生成一个不可约子模。由 \(G_{\mathbb{C}}\) 的表示论，此子模必为具有最高权 \(\lambda\) 的不可约模 \(V_\lambda\)。且 \(H^0\) 的维数等于Weyl维数公式给出的 \(\dim V_\lambda\)（这可通过Atiyah-Bott不动点定理或Riemann-Roch定理验证）。由于维数匹配，\(H^0(G_{\mathbb{C}}/B, \mathcal{L}_\lambda) \cong V_\lambda\)。\(\blacksquare\)

\subsection{Demazure特征标公式与Bott-Borel-Weil定理}\label{demazureux7279ux5f81ux6807ux516cux5f0fux4e0ebott-borel-weilux5b9aux7406}

\textbf{定理86.3（Demazure特征标公式）} 对任意Weyl群元素 \(w \in W\)，Schubert簇 \(X_w = \overline{BwB/B} \subseteq G/B\) 上的线丛 \(\mathcal{L}_\lambda\) 的全纯截面空间的特征标由Demazure算子 \(D_\alpha\)（\(\alpha\) 为单根）的复合给出：

\[\operatorname{ch} H^0(X_w, \mathcal{L}_\lambda) = D_{s_{i_1}} \circ \cdots \circ D_{s_{i_k}} (e^\lambda)\]

其中 \(w = s_{i_1} \cdots s_{i_k}\) 是约化分解，且 \(D_\alpha(e^\mu) = \frac{e^\mu - e^{s_\alpha(\mu) - \alpha}}{1 - e^{-\alpha}}\)。

\textbf{证明}（框架）对 \(w\) 的长度归纳并利用Schubert簇的Bott-Samelson分解（即奇点解消）。Demazure算子的定义直接来自Bott-Borel-Weil定理中上同调的长正合列和Leray谱序列。关键在于验证Demazure算子满足辫群关系------这反映了 \(H^0(X_w, \mathcal{L}_\lambda)\) 作为Demazure模的组合结构。\(\blacksquare\)

\textbf{定理86.4（Bott-Borel-Weil定理）} 设 \(\lambda\) 是任意整权（不一定是支配的）。若 \(\lambda + \rho\) 是正则的（即对任意 \(\alpha > 0\)，\(\langle \lambda + \rho, \alpha^\vee \rangle \neq 0\)），设 \(w \in W\) 是使 \(w(\lambda + \rho)\) 为支配权的唯一元素。则

\[H^q(G/B, \mathcal{L}_\lambda) \cong \begin{cases} V_{w(\lambda + \rho) - \rho}, & \text{若 } q = \ell(w) \\ 0, & \text{若 } q \neq \ell(w) \end{cases}\]

若 \(\lambda + \rho\) 非正则（落在某根超平面上），则对所有 \(q\)，\(H^q(G/B, \mathcal{L}_\lambda) = 0\)。

\textbf{证明}（Bott，1957）核心思想是通过单根对应的 \(\mathbb{P}^1\)-纤维化计算Leray谱序列。对每个单根 \(\alpha\)，存在纤维化 \(G/B \to G/P_\alpha\)（纤维为 \(\mathbb{P}^1\)），其中 \(P_\alpha \supset B\) 是对应于 \(\alpha\) 的极小抛物子群。线丛 \(\mathcal{L}_\lambda\) 在此纤维化下的上同调可通过 \(\mathbb{P}^1\) 上的全纯线丛上同调的Bott公式显式计算（后者归结为 \(\mathcal{O}(k)\) 的 \(H^0\) 和 \(H^1\)）。归纳地沿Weyl群的约化分解计算所有上同调，即得Bott-Borel-Weil定理。\(\blacksquare\)

\textbf{推论86.5（Kostant同调群）} Demazure特征标公式和Bott-Borel-Weil定理的共同推广是Kostant同调群。\(H^q(G/B, \mathcal{L}_\lambda)\) 当 \(\lambda\) 为反支配时给出Kostant的 \(\mathfrak{n}\)-同调群：\(H_q(\mathfrak{n}^+, V_\mu)_{-\lambda}\)，其中 \(\mu = -w_0(\lambda)\)。这些同调群与Verma模的Bernstein-Gelfand-Gelfand分解中的Ext函子直接相关，构成了李代数上同调与表示论之间的核心桥梁。

\hrulefill

\hrulefill

\hrulefill

\hrulefill

\hrulefill

\hrulefill

\section{第146章：Borel-Weil-Bott 定理}\label{ux7b2c146ux7ae0borel-weil-bott-ux5b9aux7406}

Borel-Weil-Bott 定理是李理论表示论中最辉煌的几何实现定理之一，它将紧 Lie 群的不可约表示实现为旗簇 \(G/B\) 上齐性线丛的全纯截面。这一定理的历史可追溯至 Borel 和 Weil 在 1954 年的开创性工作，他们证明了紧半单 Lie 群 \(G\) 的每个支配整权 \(\lambda \in \Lambda^+\) 对应的不可约表示 \(V_\lambda\) 同构于 \(G/B\) 上线丛 \(\mathcal{L}_\lambda\) 的整体截面空间 \(H^0(G/B, \mathcal{L}_\lambda)\)。Raoul Bott 于 1957 年将这一结果推广为完整的 Bott-Borel-Weil 定理，精确计算了所有线丛 \(\mathcal{L}_\lambda\)（\(\lambda \in \Lambda\)）的层上同调 \(H^i(G/B, \mathcal{L}_\lambda)\)------它们要么为零，要么同构于某个支配权对应的不可约表示。

Borel-Weil-Bott 定理的几何意义深远：它将李代数的纯代数表示论转化为复几何和代数几何的问题。具体而言，旗簇 \(G/B\) 是紧半单 Lie 群 \(G\) 的完备齐性空间，其上的齐性线丛 \(\mathcal{L}_\lambda\) 由权 \(\lambda\) 参数化。当 \(\lambda\) 是支配整权时，\(\mathcal{L}_\lambda\) 是丰沛线丛，其整体截面空间 \(H^0\) 就是 \(V_\lambda\)，而高阶上同调 \(H^{>0}\) 全部消失（Kempf 消没定理）。当 \(\lambda\) 不是支配权时，截面空间 \(H^0\) 为零，但高阶上同调 \(H^i\) 可能非零，且 Bott 定理精确地告诉我们在哪个 \(i = \ell(w)\) 处非零（其中 \(w\) 是使得 \(w \bullet \lambda = w(\lambda + \rho) - \rho\) 为支配权的唯一 Weyl 群元素），且该非零上同调群同构于 \(V_{w \bullet \lambda}\)。

Bott-Borel-Weil 定理的证明结合了代数几何（Schubert 胞腔的层上同调）、表示论（Demazure 算符）和微分几何（等变 K-理论）三种工具。Schubert 胞腔分层将 \(G/B\) 分解为仿射空间 \(\mathbb{C}^{\ell(w)}\) 的并，将层上同调的计算归约为组合问题。Demazure 算符 \(D_\alpha\) 则捕捉了 \(\mathbb{P}^1\)-纤维化 \(\pi_\alpha: G/B \to G/P_\alpha\) 上同调的递归行为。Weyl 特征公式作为 Euler-Poincare 特征公式的推论自然得出，而 Borel-Weil 定理自身则作为 Bott 定理在 \(\lambda\) 支配时的特例。本章将系统介绍旗簇的几何、齐性线丛的构造、Borel-Weil 定理、Bott-Borel-Weil 定理及其与 Weyl 特征公式的联系，为几何表示论（卷十六）和等变 K-理论（卷二十）奠定基础。

\textbf{定义 87.1}（旗簇）：设 \(G\) 为 \(\mathbb{C}\) 上的连通半单代数群，\(B \subset G\) 为 Borel 子群（极大连通可解子群），\(T \subset B\) 为极大环面。\textbf{旗簇}（flag variety）定义为完备齐性空间 \[\mathcal{B} = G/B\] \(\mathcal{B}\) 为光滑射影簇，维数 \(\dim \mathcal{B} = |\Phi^+|\)（正根的个数）。\(G\) 通过左平移传递作用于 \(\mathcal{B}\)。基本例子：\(G = GL_n(\mathbb{C})\) 时，\(\mathcal{B}\) 同构于 \(\mathbb{C}^n\) 中全体完全旗 \(\{0\} = V_0 \subset V_1 \subset \cdots \subset V_n = \mathbb{C}^n\)（\(\dim_{\mathbb{C}} V_i = i\)）构成的空间。

\textbf{定义 87.2}（权与特征）：权格为 \(\Lambda = \operatorname{Hom}_{\text{alg}}(T, \mathbb{C}^\times)\)。Borel 子群 \(B\) 决定了正根集 \(\Phi^+ \subset \Phi\)。\textbf{支配整权}集为 \[\Lambda^+ = \{\lambda \in \Lambda : \langle \lambda, \alpha^\vee \rangle \geq 0,\ \forall \alpha \in \Phi^+\}\] \(\rho = \frac{1}{2} \sum_{\alpha \in \Phi^+} \alpha\) 为全部正根之和的一半。

\textbf{定义 87.3}（齐性线丛）：对每个权 \(\lambda \in \Lambda\)，定义一维系数的 \(B\)-模 \(\mathbb{C}_\lambda\)：\(b \in B\) 通过投射 \(B \twoheadrightarrow B/U \cong T\) 后以特征 \(\lambda\) 作用于 \(\mathbb{C}\)。定义 \(G/B\) 上的\textbf{齐性线丛}为均衡积 \[\mathcal{L}_\lambda = G \times_B \mathbb{C}_\lambda \longrightarrow G/B,\quad (g, z) \sim (gb, \lambda(b)^{-1} z)\] 投影映射为 \((g, z) \mapsto gB\)。\(\mathcal{L}_\lambda\) 是 \(G/B\) 上的 \(G\)-等变全纯线丛。当 \(\lambda \in \Lambda^+\) 时，\(\mathcal{L}_\lambda\) 为丰沛线丛（ample line bundle）。

\subsection{87.2 Borel-Weil 定理}\label{borel-weil-ux5b9aux7406}

\textbf{定理 87.1}（Borel-Weil，1954）：设 \(\lambda \in \Lambda^+\) 为支配整权。则整体截面空间 \(H^0(G/B, \mathcal{L}_\lambda)\) 作为 \(G\)-模同构于最高权为 \(\lambda\) 的不可约有限维表示 \(V_\lambda\)： \[H^0(G/B, \mathcal{L}_\lambda) \cong V_\lambda\] 若 \(\lambda \notin \Lambda^+\)（非支配权），则 \(H^0(G/B, \mathcal{L}_\lambda) = 0\)。

\textbf{证明}：利用 \(G\) 的 Bruhat 分解 \(G = \bigsqcup_{w \in W} BwB\)，其中 \(W = N_G(T)/T\) 为 Weyl 群。\(G/B\) 分解为 \(B\)-轨道的并（Schubert 胞腔），每个 \(\mathcal{O}_w = BwB/B \cong \mathbb{C}^{\ell(w)}\) 为仿射空间。截面 \(s \in H^0(G/B, \mathcal{L}_\lambda)\) 对应于全纯函数 \(\varphi: G \to \mathbb{C}\) 满足 \(\varphi(gb) = \lambda(b)^{-1} \varphi(g)\)（对 \(b \in B\)）。由 Borel 不动点定理，\(G\) 在 \(H^0\) 上的作用诱导 \(G\)-模结构。设 \(\lambda \in \Lambda^+\) 为支配整权，则存在最高权向量 \(v_\lambda \in V_\lambda\)。由 Peter-Weyl 定理，\(H^0(G/B, \mathcal{L}_\lambda)\) 分解为 \(G\)-不可约模的直和。由 Bruhat 分解，\(H^0\) 中最高权向量在 \(G\) 作用下生成的子模恰为 \(V_\lambda\)。若 \(\lambda \notin \Lambda^+\)，则 \(\mathcal{L}_\lambda\) 无整体截面（因 \(B\) 在纤维上的作用非支配），故 \(H^0 = 0\)。\(\blacksquare\)

\textbf{定理 87.2}（Borel-Weil 与 Weyl 特征公式）：作为 Euler-Poincaré 特征， \[\chi(G/B, \mathcal{L}_\lambda) = \frac{\sum_{w \in W} (-1)^{\ell(w)} e^{w(\lambda+\rho)}}{\sum_{w \in W} (-1)^{\ell(w)} e^{w(\rho)}} = \operatorname{ch} V_\lambda\]

\textbf{证明}：由 Borel-Weil 定理，\(H^0(G/B, \mathcal{L}_\lambda) \cong V_\lambda\)（\(\lambda \in \Lambda^+\)）且由 Kempf 消没定理，\(H^{>0}(G/B, \mathcal{L}_\lambda) = 0\)。故对支配权 \(\lambda\)，Euler-Poincare 特征 \(\chi(G/B, \mathcal{L}_\lambda) = \operatorname{ch} H^0 = \operatorname{ch} V_\lambda\)。另一方面，由 Atiyah-Bott 不动点公式应用于极大环面 \(T\) 在 \(G/B\) 上的作用，\(T\)-等变 \(K\)-理论给出 \(\chi(G/B, \mathcal{L}_\lambda) = \sum_{w \in W} (-1)^{\ell(w)} e^{w(\lambda+\rho)} / \prod_{\alpha \in \Phi^+} (e^{\alpha/2} - e^{-\alpha/2})\)。分母由 Weyl 分母公式 \(\prod_{\alpha \in \Phi^+} (e^{\alpha/2} - e^{-\alpha/2}) = e^{-\rho} \sum_{w \in W} (-1)^{\ell(w)} e^{w(\rho)}\) 化简，即得 Weyl 特征公式。\(\blacksquare\)

\subsection{87.3 Bott-Borel-Weil 定理}\label{bott-borel-weil-ux5b9aux7406}

\textbf{定义 87.4}（Weyl 群的点作用）：对 \(w \in W\) 和 \(\lambda \in \Lambda\)，定义平移 Weyl 群作用（点作用） \[w \bullet \lambda = w(\lambda + \rho) - \rho\] 权 \(\lambda\) 称为\textbf{正则}的，若对所有 \(\alpha \in \Phi\)，\(\langle \lambda + \rho, \alpha^\vee \rangle \neq 0\)（即 \(\lambda + \rho\) 不落在任何 Weyl 腔的壁上）。\(\lambda\) 称为\textbf{奇异}的，若存在 \(\alpha\) 使得 \(\langle \lambda + \rho, \alpha^\vee \rangle = 0\)。

\textbf{定理 87.3}（Bott-Borel-Weil，Bott 1957）：设 \(\lambda \in \Lambda\)。

\begin{enumerate}
\def\labelenumi{\arabic{enumi}.}
\item
  若 \(\lambda + \rho\) 为正则的，则存在唯一的 \(w \in W\) 使得 \(w \bullet \lambda \in \Lambda^+\)（支配），且 \[H^i(G/B, \mathcal{L}_\lambda) \cong \begin{cases}
  V_{w \bullet \lambda} & \text{若 } i = \ell(w) \\
  0 & \text{若 } i \neq \ell(w)
  \end{cases}\]
\item
  若 \(\lambda + \rho\) 为奇异的，则对所有 \(i \geq 0\)，\(H^i(G/B, \mathcal{L}_\lambda) = 0\)。
\end{enumerate}

\textbf{证明}：证明基于三个关键技术。

（i）Schubert 胞腔的层上同调：将 \(G/B\) 按 Schubert 胞腔 \(\mathcal{O}_w\) 分层。对每个 Schubert 簇 \(X_w = \overline{\mathcal{O}_w}\)，考虑闭嵌入 \(i_w: X_w \hookrightarrow G/B\) 和开嵌入 \(j_w: \mathcal{O}_w \hookrightarrow X_w\)。由局部上同调长正合列，\(H^i(X_w, \mathcal{L}_\lambda|_{X_w})\) 可通过 \(H^i(\mathcal{O}_w, \mathcal{L}_\lambda|_{\mathcal{O}_w})\) 和 \(H^i(X_{v < w}, \mathcal{L}_\lambda|_{X_{v < w}})\) 归约计算。\(\mathcal{O}_w \cong \mathbb{C}^{\ell(w)}\) 为仿射空间，其上的线丛上同调由 Borel-Weil-Bott 的局部版本计算。

（ii）Demazure 算符：对每个单根 \(\alpha \in \Delta\)，定义 Demazure 算符 \(D_\alpha(e^\mu) = \frac{e^\mu - e^{s_\alpha(\mu)-\alpha}}{1 - e^{-\alpha}}\)。其几何意义：\(\pi_\alpha: G/B \to G/P_\alpha\) 为 \(\mathbb{P}^1\)-纤维化。\(\mathcal{L}_\lambda\) 限制在每条纤维 \(\mathbb{P}^1\) 上对应于 \(\mathcal{O}(-n)\)（\(n = \langle \lambda+\rho, \alpha^\vee \rangle - 1\)）。\(H^0(\mathbb{P}^1, \mathcal{O}(-n))\) 非零仅当 \(n \leq 0\)，此时 \(H^0\) 为 \(n+1\) 维，特征标为 \(e^\lambda + e^{\lambda-\alpha} + \cdots + e^{\lambda-n\alpha}\) 与 \(D_\alpha(e^\lambda)\) 一致；\(H^1\) 仅在 \(n \geq 2\) 时非零，特征标为 \(e^{s_\alpha(\lambda+\rho)-\rho}\)。

（iii）归纳论证：设 \(\lambda+\rho\) 为正则且 \(\langle \lambda+\rho, \alpha^\vee \rangle < 0\)，则 \(\mathcal{L}_\lambda\) 在 \(\mathbb{P}^1\)-纤维上没有截面，上同调集中于 \(H^1\)。每一次穿越单根 \(\alpha\) 的壁，上同调群移动一个指标：\(H^i(G/B, \mathcal{L}_\lambda) \cong H^{i-1}(G/B, \mathcal{L}_{s_\alpha \bullet \lambda})\)（含 \(\ell(s_\alpha)=1\) 的情形的 Serre 对偶变体）。对一般 \(w \in W\)，将 \(w\) 写为单反射的约化乘积，逐次应用上述步骤归纳。基情形 \(w = \operatorname{id}\)（\(\lambda\) 支配）即 Borel-Weil 定理。\(\blacksquare\)

\textbf{定理 87.4}（Weyl 特征公式作为推论）：对任意 \(\lambda \in \Lambda\)，由 Euler-Poincaré 特征 \[\chi(G/B, \mathcal{L}_\lambda) = \sum_{i \geq 0} (-1)^i \operatorname{ch} H^i(G/B, \mathcal{L}_\lambda)\] - 若 \(\lambda \in \Lambda^+\)，仅有 \(H^0\) 非零，\(\chi = \operatorname{ch} V_\lambda\) - 若 \(\lambda + \rho\) 为正则且 \(w \bullet \lambda \in \Lambda^+\)，\(\chi = (-1)^{\ell(w)} \operatorname{ch} V_{w \bullet \lambda}\)

因为特征标映射是 \(K\)-理论中的加性不变量，上述 Euler-Poincaré 特征恒等于 Atiyah-Bott 不动点公式给出的表达式 \[\chi(G/B, \mathcal{L}_\lambda) = \frac{\sum_{w \in W} (-1)^{\ell(w)} e^{w(\lambda+\rho)}}{\prod_{\alpha \in \Phi^+} (e^{\alpha/2} - e^{-\alpha/2})}\] 其中分母 \(\prod_{\alpha \in \Phi^+} (e^{\alpha/2} - e^{-\alpha/2}) = e^{-\rho} \prod_{\alpha \in \Phi^+} (e^\alpha - 1) = \sum_{w \in W} (-1)^{\ell(w)} e^{w(\rho) - \rho}\) 为 Weyl 分母恒等式。分子分母化简即得标准 Weyl 特征公式。

\textbf{证明}：由定理 87.3（Bott-Borel-Weil），上同调群 \(H^i(G/B, \mathcal{L}_\lambda)\) 要么全为零，要么恰在一个 \(i\) 处非零（\(i = \ell(w)\) 时 \(H^{\ell(w)} \cong V_{w \bullet \lambda}\)）。因此 Euler-Poincaré 特征 \(\chi = \sum (-1)^i \operatorname{ch} H^i = (-1)^{\ell(w)} \operatorname{ch} V_{w \bullet \lambda}\)。Atiyah-Bott 不动点公式将 \(\chi\) 计算为 \(T\)-等变 \(K\)-理论中的迹，给出所述分数形式。结合 Weyl 分母恒等式 \(\sum_{w \in W} (-1)^{\ell(w)} e^{w(\rho)} = e^\rho \prod_{\alpha > 0} (1 - e^{-\alpha})\) 即得标准 Weyl 特征公式。\(\blacksquare\)

\subsection{87.4 与几何量子化的联系}\label{ux4e0eux51e0ux4f55ux91cfux5b50ux5316ux7684ux8054ux7cfb}

\textbf{注记 87.1}（几何量子化，Kostant-Souriau 1960年代）：Borel-Weil-Bott 定理是几何量子化最经典的范例。在几何量子化框架中：

\begin{itemize}
\tightlist
\item
  余伴随轨道 \(\mathcal{O}_\lambda = G \cdot \lambda \subseteq \mathfrak{g}^*\)（\(G\) 为紧实形 \(G_c\) 时）是 \(G_c\) 的齐性辛流形，\(\mathcal{O}_\lambda \cong G_c / G_{c,\lambda} \cong G/B\)（作为复流形）
\item
  \(\mathcal{L}_\lambda \to \mathcal{O}_\lambda\) 为\textbf{预量子化线丛}，其联络曲率恰为 Kostant-Kirillov-Souriau 辛形式
\item
  Borel 子群 \(B\) 给出的复结构 \(J\) 构成 Kähler 极化。极化下的量子化 = 全纯截面的空间 \(H^0\)，恰好实现不可约酉表示 \(V_\lambda\)
\item
  Bott-Borel-Weil 定理计算了高次上同调，证明了正则权情形的量子化无反常（anomalies）------上同调集中于单一度数
\end{itemize}

这完美体现了 Kostant-Kirillov ``轨道法'' 哲学：紧 Lie 群的不可约酉表示 \textless-\textgreater{} 余伴随轨道，且表示的特征标由轨道上 Liouville 测度的 Fourier 变换给出（Harish-Chandra 特征标公式）。

\subsection{87.5 Demazure 模与 Schubert 簇}\label{demazure-ux6a21ux4e0e-schubert-ux7c07}

\textbf{定义 87.5}（Schubert 簇）：对 \(w \in W\)，\textbf{Schubert 簇} \(X_w = \overline{B w B / B} \subseteq G/B\) 是 Schubert 胞腔 \(\mathcal{O}_w \cong \mathbb{C}^{\ell(w)}\) 的 Zariski 闭包。\(\dim_{\mathbb{C}} X_w = \ell(w)\)，且 \(X_w\) 为正规射影簇（一般非光滑）。全体 \(\{X_w\}_{w \in W}\) 按 Bruhat 偏序 \(v \leq w \iff X_v \subseteq X_w\) 分层。

\textbf{定义 87.6}（Demazure 模）：对 \(\lambda \in \Lambda^+\) 和 \(w \in W\)，\textbf{Demazure 模} \(V_\lambda(w)\) 为 \(V_\lambda\) 中由极值权空间 \(V_\lambda(w\lambda)\) 在 Borel 子代数 \(\mathfrak{b}\) 作用下生成的子模： \[V_\lambda(w) = U(\mathfrak{b}) \cdot V_\lambda(w\lambda) \subseteq V_\lambda\] 其中 \(V_\lambda(w\lambda)\) 为权 \(w\lambda\) 的一维权空间。\(V_\lambda(w)\) 是 \(B\) 的（无限维）表示，其有限维部分由 \(V_\lambda\) 的唯一性确定。

\textbf{定理 87.5}（Demazure 特征公式）：设 \(\lambda \in \Lambda^+\)。则 \[H^0(X_w, \mathcal{L}_\lambda|_{X_w}) \cong V_\lambda(w)^*\] 且特征标由 Demazure 算符给出：\(\operatorname{ch} V_\lambda(w) = D_w(e^\lambda)\)，其中 \(D_w = D_{\alpha_{i_1}} \circ \cdots \circ D_{\alpha_{i_\ell}}\)（\(w = s_{i_1} \cdots s_{i_\ell}\) 为约化分解），而单根算符为 \[D_\alpha(e^\mu) = \frac{e^\mu - e^{s_\alpha(\mu) - \alpha}}{1 - e^{-\alpha}}\] 取 \(w = w_0\)（最长元）时 \(V_\lambda(w_0) = V_\lambda\)，即回到 Borel-Weil 定理。

\textbf{证明}：对 \(\ell(w)\) 归纳。基情形 \(\ell(w) = 0\)（\(w = \operatorname{id}\)）：\(X_{\operatorname{id}} = \{\text{pt}\}\)（单点），\(H^0 = \mathbb{C}\) 为一维，\(V_\lambda(\operatorname{id})\) 为极值权空间 \(V_\lambda(\lambda)\)，两者对偶。单反射 \(s_\alpha\) 情形：\(X_{s_\alpha} \cong \mathbb{P}^1\)，\(\mathcal{L}_\lambda|_{\mathbb{P}^1}\) 为上同调已知的线丛。利用 \(\mathbb{P}^1\)-纤维化 \(\pi_\alpha: G/B \to G/P_\alpha\) 的 Leray 谱序列，\(\pi_{\alpha *} \mathcal{L}_\lambda \cong \mathcal{L}_\lambda'\)（\(G/P_\alpha\) 上的线丛），\(H^0(X_{s_\alpha}, \mathcal{L}_\lambda) \cong H^0(G/P_\alpha, \pi_{\alpha *} \mathcal{L}_\lambda)\)。由 Borel-Weil 定理在 \(G/P_\alpha\) 上的应用，\(H^0\) 给出 \(V_\lambda(s_\alpha)^*\)。一般 \(w\)：取约化分解 \(w = s_{\alpha_1} \cdots s_{\alpha_\ell}\)。由归纳假设，\(H^0(X_{ws_\alpha}, \mathcal{L}_\lambda) \cong V_\lambda(ws_\alpha)^*\)。利用 Demazure 算符的函子性质，\(D_{ws_\alpha} = D_w \circ D_{s_\alpha}\)，且 \(D_{s_\alpha}(e^\mu) = \operatorname{ch} V_\mu(s_\alpha)\)。故 \(\operatorname{ch} V_\lambda(w) = D_w(e^\lambda)\)。取 \(w = w_0\)（最长元）时 \(V_\lambda(w_0) = V_\lambda\)，回到 Borel-Weil 定理。\(\blacksquare\)

\subsection{87.6 Beilinson-Bernstein 局部化}\label{beilinson-bernstein-ux5c40ux90e8ux5316}

\textbf{定理 87.6}（Beilinson-Bernstein 局部化，1981）：设 \(\mathfrak{g}\) 为复半单 Lie 代数，\(\lambda \in \mathfrak{h}^*\) 为支配正则权。则存在范畴等价 \[\mathcal{M}(\mathfrak{g}, \lambda) \cong \operatorname{QCoh}(G/B, \mathcal{D}_\lambda)\] 其中左端为具有平凡广义中心特征的 \(\mathfrak{g}\)-模（Harish-Chandra 双模范畴），右端为 \(G/B\) 上扭折微分算符层 \(\mathcal{D}_\lambda\) 的拟凝聚层范畴。该等价将 Verma 模映为标准 \(\mathcal{D}\)-模，将单最高权模映为连接 \(\mathcal{D}\)-模（即全纯线丛的 \(\mathcal{D}\)-模推送）。

\textbf{证明}：关键步骤三步。(1) 构造「局部化」函子：给定 \(\mathfrak{g}\)-模 \(M\)（满足 Harish-Chandra 条件，即 \(\mathcal{Z}(\mathfrak{g})\) 通过中心特征 \(\chi_\lambda\) 作用），在旗簇上定义 \(\mathcal{D}_\lambda\)-模 \(\Delta(M) = \mathcal{D}_\lambda \otimes_{U(\mathfrak{g})_\lambda} M\)，其中 \(U(\mathfrak{g})_\lambda = U(\mathfrak{g}) / \ker(\chi_\lambda)\) 为广义中心特征商。\(\Delta(M)\) 是拟凝聚 \(\mathcal{D}_\lambda\)-模，函子 \(\Delta\) 正合。(2) 构造「整体截面」函子 \(\Gamma(G/B, -)\) 作为拟逆。核心是证明 \(\Gamma \circ \Delta \cong \operatorname{id}\)：由 \(\mathcal{D}_\lambda\) 的整体截面环 \(\Gamma(G/B, \mathcal{D}_\lambda) \cong U(\mathfrak{g})_\lambda\)（扭折 \(\mathcal{D}\)-模的整体截面同构于约化包络代数），故 \(\Gamma(\Delta(M)) \cong U(\mathfrak{g})_\lambda \otimes_{U(\mathfrak{g})_\lambda} M \cong M\)。同时证明 \(\Delta \circ \Gamma \cong \operatorname{id}\)：由 \(\mathcal{D}_\lambda\) 作为 \(\mathcal{O}\)-模的忠实平坦性，\(\Delta\) 和 \(\Gamma\) 互为拟逆。(3) \(\lambda\) 的正则性条件保证 \(\mathcal{D}_\lambda\) 的「扭折」是平凡的：此时 \(\mathcal{D}_\lambda\) 和 \(\mathcal{D}_0\)（通常 \(\mathcal{D}\)-模层）Morita 等价，因此 \(\mathcal{D}_\lambda\)-模范畴与 \(\mathcal{D}_0\)-模范畴等价。\(\blacksquare\)

\textbf{定理 87.7}（局部化与 Borel-Weil-Bott 的联系）：对支配正则权 \(\lambda\)，整体截面函子 \(\Gamma(G/B, -)\) 在扭折 \(\mathcal{D}\)-模上正合，且 \[\Gamma(G/B, \mathcal{D}_\lambda \otimes_{\mathcal{O}} \mathcal{L}_\mu) \cong M(\mu)^* \quad (\text{$\mu$ 与 $\lambda$ 全整链接})\] 其中 \(M(\mu)\) 为权 \(\mu\) 的 Verma 模。对反支配权（antidominant），\(\Gamma\) 函子非正合而需考察高次上同调，这正好对应于 Bott-Borel-Weil 定理中的上同调转移现象。对奇异权（\(\lambda+\rho\) 落在 Weyl 腔壁上），通过平移函子将问题归约至正则情形，给出全部 \(H^i(G/B, \mathcal{L}_\mu)\) 的代数刻画。

\textbf{证明}：由 Beilinson-Bernstein 局部化，\(\mathcal{D}_\lambda\)-模范畴等价于 \(\mathfrak{g}\)-模范畴（具有中心特征 \(\chi_\lambda\)）。\(\mathcal{D}_\lambda \otimes_{\mathcal{O}} \mathcal{L}_\mu\) 是 \(\mathcal{D}_\lambda\)-模（\(\mathcal{D}_\lambda\) 作用于第一个因子），其整体截面 \(\Gamma(G/B, \mathcal{D}_\lambda \otimes \mathcal{L}_\mu)\) 为 \(\mathfrak{g}\)-模。直接计算该模的特征标：由 \(\mathcal{D}_\lambda\) 的局部结构，\(\Gamma(G/B, \mathcal{D}_\lambda \otimes \mathcal{L}_\mu)\) 包含最高权为 \(\mu\) 的向量（由 \(\mathcal{L}_\mu\) 的常截面经 \(\mathcal{D}_\lambda\) 作用生成），其权空间分解为 \(\mu - \sum n_i \alpha_i\)（\(n_i \in \mathbb{N}\)）。因此特征标为 \(e^\mu / \prod_{\alpha > 0} (1 - e^{-\alpha})\)，这正是 Verma 模 \(M(\mu)\) 的特征标。由 Harish-Chandra 定理，\(M(\mu)\) 由该特征标唯一确定，故 \(\Gamma(G/B, \mathcal{D}_\lambda \otimes \mathcal{L}_\mu) \cong M(\mu)^*\)。对非支配权，\(\Gamma\) 非正合导致高次上同调的非消没；由定理 87.3，这些高次上同调群恰好实现为相同 LKT 等价类中的不可约表示。\(\blacksquare\)

\textbf{注记 87.2}：Beilinson-Bernstein 局部化定理是几何表示论的基石，将代数 \(\mathcal{D}\)-模理论与李代数表示论统一。该定理为 Kazhdan-Lusztig 猜想（由 Brylinski-Kashiwara 与 Beilinson-Bernstein 于 1981 年独立证明）提供了几何框架，将 Borel-Weil-Bott 嵌入为宏观对应体系的特例：旗簇上的 \(\mathcal{D}\)-模、反常层和反常复形精确分类了李代数的全体表示。

\newpage

\chapter{05-代数/04-李理论/03-李群李代数深化.md}\label{ux4ee3ux657004-ux674eux7406ux8bba03-ux674eux7fa4ux674eux4ee3ux6570ux6df1ux5316.md}

\section{第147章：仿射与量子群表示}\label{ux7b2c147ux7ae0ux4effux5c04ux4e0eux91cfux5b50ux7fa4ux8868ux793a}

仿射 Kac-Moody 代数是有限维单李代数的圈代数的中心扩张，由 Victor Kac 和 Robert Moody 在 1960 年代独立引入。与有限维李代数不同，仿射李代数的不可约可积最高权表示通常是无限维的，但其特征标由 Kac-Weyl 特征标公式精确刻画------这一公式是 Weyl 特征标公式在无限维情形的自然推广，其分母公式（Macdonald 恒等式）包含了经典数论中 Jacobi 三重积恒等式和 Dedekind \(\eta\) 函数的深刻推广。仿射李代数的表示论与二维共形场论有着密不可分的联系：Frenkel-Zhu 定理（1992）证明了仿射 Kac-Moody 代数的可积最高权模具有顶点算子代数（VOA）结构，为共形场论提供了严格的数学基础。Verlinde 公式则通过模群 \(S\) 矩阵将 VOA 的融合规则与三维拓扑量子场论联系起来。与此同时，量子群（quantum group）是 Drinfeld 和 Jimbo 在 1985 年独立发现的------它是李代数泛包络代数的 \(q\)-形变，其关键创新在于引入了拟三角结构（泛 \(R\)-矩阵）满足 Yang-Baxter 方程。量子群的表示范畴是辫子幺半范畴，辫子结构直接给出了纽结和链环的不变量（Jones 多项式、HOMFLY-PT 多项式），将表示论与低维拓扑深刻联系起来。

\subsection{225.1 仿射 Lie 代数表示}\label{ux4effux5c04-lie-ux4ee3ux6570ux8868ux793a}

\textbf{定义 225.1}（仿射 Kac-Moody 代数）：\textbf{仿射 Kac-Moody 代数} \(\hat{\mathfrak{g}}\) 是有限维单李代数 \(\mathfrak{g}\) 的圈代数 \(\mathfrak{g} \otimes \mathbb{C}[t, t^{-1}]\) 的中心扩张。其定义为： - \(\hat{\mathfrak{g}} = (\mathfrak{g} \otimes \mathbb{C}[t, t^{-1}]) \oplus \mathbb{C} c \oplus \mathbb{C} d\) - \([x \otimes t^m, y \otimes t^n] = [x, y] \otimes t^{m+n} + m \delta_{m,-n} (x|y) c\) - \([c, \cdot] = 0\)，\([d, x \otimes t^n] = n x \otimes t^n\)

\textbf{定义 225.2}（可积最高权表示）：设 \(\lambda\) 是 \(\hat{\mathfrak{g}}\) 的支配整权。不可约最高权模 \(L(\lambda)\) 是\textbf{可积的}，如果 \(\lambda(c) \neq 0\) 且 Chevalley 生成元的作用是局部幂零的。

\textbf{定理 225.1}（Kac-Weisfeiler 特征标公式 / Kac 特征标公式）：可积最高权模 \(L(\lambda)\) 的特征标为

\[\operatorname{ch} L(\lambda) = \frac{\sum_{w \in \hat{W}} (-1)^{\ell(w)} e^{w(\lambda + \hat{\rho}) - \hat{\rho}}}{\prod_{\alpha > 0} (1 - e^{-\alpha})^{\dim \hat{\mathfrak{g}}_\alpha}}\]

其中 \(\hat{W}\) 是仿射 Weyl 群，\(\hat{\rho}\) 是仿射 Weyl 向量。

\textbf{证明}：由 Verma 模的正合列和 BGG 消没定理：\(\operatorname{ch} L(\lambda) = \sum_{w \in \hat{W}} (-1)^{\ell(w)} \operatorname{ch} M(w \cdot \lambda)\)（其中 \(w \cdot \lambda = w(\lambda + \hat{\rho}) - \hat{\rho}\) 为 dot action）。Verma 模 \(M(\mu)\) 的特征标为 \(\operatorname{ch} M(\mu) = e^\mu / \prod_{\alpha>0} (1-e^{-\alpha})^{\dim \hat{\mathfrak{g}}_\alpha}\)。代入求和并化简分子得 \(\sum_w (-1)^{\ell(w)} e^{w(\lambda + \hat{\rho}) - \hat{\rho}} / \prod_{\alpha>0} (1-e^{-\alpha})^{\dim \hat{\mathfrak{g}}_\alpha}\)。分母为仿射 Weyl 分母，当 \(\lambda = 0\) 时即为 Macdonald 分母恒等式。 \(\blacksquare\)

\textbf{定义 225.3}（Macdonald 恒等式与 \(\eta\) 函数）：仿射李代数的分母恒等式（Weyl 分母公式）给出经典数论公式的推广。\(\widehat{\mathfrak{sl}}_2\) 的分母公式给出 Jacobi 三重积恒等式和 Dedekind \(\eta\) 函数。

\subsection{225.2 顶点算子代数}\label{ux9876ux70b9ux7b97ux5b50ux4ee3ux6570}

\textbf{定义 225.4}（顶点算子代数 / VOA）：\textbf{顶点算子代数} \((V, Y, \mathbf{1}, \omega)\) 是向量空间 \(V\) 配备： - 状态-场对应 \(Y: V \to \operatorname{End}(V)[[z, z^{-1}]]\)，\(Y(a, z) = \sum_{n \in \mathbb{Z}} a_n z^{-n-1}\) - 真空向量 \(\mathbf{1}\)：\(Y(\mathbf{1}, z) = \operatorname{id}_V\) - 共形向量 \(\omega\)：\(Y(\omega, z) = \sum L_n z^{-n-2}\)（Virasoro 代数作用） - 局部性公理：\((z-w)^N [Y(a, z), Y(b, w)] = 0\)（\(N\) 充分大）

\textbf{定理 225.2}（Frenkel-Zhu 定理，1992）：仿射 Kac-Moody 代数 \(\hat{\mathfrak{g}}\) 的可积最高权模 \(L(k\Lambda_0)\)（\(k\) 是级）具有顶点算子代数结构（仿射 VOA）。

\textbf{证明}：构造状态-场对应 \(Y: L(k\Lambda_0) \to \operatorname{End}(L(k\Lambda_0))[[z, z^{-1}]]\)，基向量 \(\mathbf{1}\) 对应真空，共形向量由 Sugawara 构造给出：\(\omega = \frac{1}{2(k+h^\vee)} \sum_a J^a(-1)J^a(-1)\mathbf{1}\)（\(h^\vee\) 为对偶 Coxeter 数）。验证 Virasoro 代数关系及中心荷 \(c = \frac{k \dim \mathfrak{g}}{k+h^\vee}\)。局部性由自由玻色子实现和 Frenkel-Kac 构造保证：\(Y(a, z)\) 生成子的 OPE 无极点。 \(\blacksquare\)

\textbf{定义 225.5}（共形场论与表示论）：有理顶点算子代数 \(V\) 的表示范畴是模张量范畴（modular tensor category），与 3 维拓扑量子场论（TQFT）相关。这是 VOA 表示论与低维拓扑的联系。

\textbf{定理 225.3}（Verlinde 公式，1988）：有理 VOA 的不可约表示 \(M_i\) 的融合规则（fusion rules）\(N_{ij}^k\) 与模群 \(S\) 矩阵的关系为 \(N_{ij}^k = \sum_r \frac{S_{ir} S_{jr} S_{kr}^*}{S_{0r}}\)。

\textbf{证明}：有理 VOA 的表示范畴为模张量范畴。其共形块空间维数由 fusion rules 给出。模群作用由 \(S\) 矩阵 \(S_{ij} = \operatorname{STr}_{M_i \otimes M_j}(q^{L_0-c/24})\) 描述（在环面上）。Verlinde 的非对角 CFT 论证将 fusion ring 对角化：\(S_{ij}/S_{0j}\) 是 fusion 矩阵 \(N_i\) 的本征值，由模不变性 \(S^2 = C\)、\((ST)^3 = S^2\) 得所述公式。Moore-Seiberg 的公理化框架对此给出了完全的范畴论证明。 \(\blacksquare\)

\subsection{225.3 量子仿射代数与可积系统}\label{ux91cfux5b50ux4effux5c04ux4ee3ux6570ux4e0eux53efux79efux7cfbux7edf}

\textbf{定义 225.6}（量子仿射代数 \(U_q(\hat{\mathfrak{g}})\)）：\textbf{量子仿射代数}是仿射 Kac-Moody 代数 \(\hat{\mathfrak{g}}\) 的量子包络代数的 \(q\)-形变。它是 Drinfeld-Jimbo 量子群（见 V23 Ch 268）的仿射版本。

\textbf{定义 225.7}（量子 R-矩阵与 Yang-Baxter 方程）：量子仿射代数 \(U_q(\hat{\mathfrak{g}})\) 的泛 R-矩阵满足 Yang-Baxter 方程。有限维表示 \(V\) 上的 R-矩阵 \(R_{V,W}: V \otimes W \to V \otimes W\) 是谱依赖的（\(R(z)\)）。

\textbf{定理 225.4}（Baxter 的角转移矩阵与 Bethe 拟设）：量子仿射代数 \(U_q(\hat{\mathfrak{sl}}_2)\) 的表示论与 XXZ 自旋链的 Bethe 拟设解直接相关。这是量子可积系统与量子群表示论的核心联系。

\textbf{证明}：XXZ 自旋链的 Hamiltonian \(H_{XXZ} = \sum_j (\sigma_j^x \sigma_{j+1}^x + \sigma_j^y \sigma_{j+1}^y + \Delta \sigma_j^z \sigma_{j+1}^z)\) 与 \(U_q(\hat{\mathfrak{sl}}_2)\) 的 \(R\)-矩阵 \(R(z)\) 满足 QISM 关系 \(R(z/w)T(z)T(w) = T(w)T(z)R(z/w)\)。转移矩阵 \(t(z) = \operatorname{Tr}_a T_a(z)\) 的对角化利用代数 Bethe 拟设：在 \(U_q(\hat{\mathfrak{sl}}_2)\) 的 Verma 模上，\(B(u)\) 算子将赝真空提升为 Bethe 态，\(A(u), D(u)\) 算子给出 Bethe 方程。 \(\blacksquare\)

\textbf{定义 225.8}（Frenkel-Reshetikhin 的 \(q\)-特征标，1999）：量子仿射代数的有限维表示的 \(q\)-特征标是经典特征标的 \(q\)-形变，满足与 Baxter 的 T-Q 关系类似的性质。

\subsection{仿射李代数与量子群的前沿应用}\label{ux4effux5c04ux674eux4ee3ux6570ux4e0eux91cfux5b50ux7fa4ux7684ux524dux6cbfux5e94ux7528}

仿射 Kac-Moody 代数和量子群的表示论在 20 世纪末和 21 世纪初与多个数学物理分支产生了深刻交叉。\textbf{月光猜想}（Monstrous Moonshine, Conway-Norton 1979）揭示了 Monster 单群与模函数 \(j(\tau)\) 之间的神秘联系，Borcherds（1992 年 Fields 奖）通过顶点算子代数（VOA）------特别是以仿射 Kac-Moody 代数的可积表示构造的 \textbf{Monster VOA} \(V^\natural\)------严格证明了这一猜想，将群论、数论和共形场论联系起来。\textbf{几何 Langlands 纲领}中，仿射 Kac-Moody 代数的表示论通过\textbf{共形块}（conformal blocks）和 \textbf{Hitchin 系统}的量子化与 Langlands 对偶群的表示范畴产生联系。Beilinson-Drinfeld 的 opers 理论（1990s）将仿射 Lie 代数的 Verma 模与代数曲线上的 \(G^\vee\)-oper 联系起来，建立了局部几何 Langlands 对应。\textbf{可观测量代数}（\(\mathcal{W}\)-代数）是仿射 Kac-Moody 代数的量子 Hamiltonian 约化（Drinfeld-Sokolov 约化，1984），其表示论通过二维共形场论中的 \(\mathcal{W}\)-对称性在弦论和可积系统中扮演关键角色。\textbf{AGT 对应}（Alday-Gaiotto-Tachikawa 2009）将 \(\mathcal{N}=2\) 超对称规范理论中的 Nekrasov 配分函数与二维共形场论中 Liouville 理论（和 \(\mathcal{W}\)-代数）的共形块联系起来，建立了四维规范理论和二维共形场论之间的深刻对偶。\textbf{量子仿射代数}的 \textbf{Kirillov-Reshetikhin 模}（1988）和它们的 \(q\)-特征标通过 Frenkel-Mukhin 算法（2001）与 Bethe 拟设方程的解建立了精确对应，构成量子可积系统中表示论方法的核心。\textbf{Brauer 代数}、\textbf{Temperley-Lieb 代数}和\textbf{结合代数}（fusion algebra）的范畴化通过 Jones 的纽结不变量将量子群表示论与低维拓扑联系起来，为拓扑量子计算提供了数学基础。

\hrulefill

\hrulefill

\hrulefill

\section{第148章：泛包络代数与PBW定理}\label{ux7b2c148ux7ae0ux6cdbux5305ux7edcux4ee3ux6570ux4e0epbwux5b9aux7406}

泛包络代数将李代数嵌入结合代数，使得李代数的表示论转化为结合代数的模论。PBW（Poincaré-Birkhoff-Witt）定理是泛包络代数的基本结构定理。

\subsection{123.1 泛包络代数的构造}\label{ux6cdbux5305ux7edcux4ee3ux6570ux7684ux6784ux9020}

\textbf{定义 123.1}（泛包络代数）：李代数 \(\mathfrak{g}\) 的\textbf{泛包络代数} \(U(\mathfrak{g})\) 是张量代数 \(T(\mathfrak{g}) = \bigoplus_{k=0}^{\infty} \mathfrak{g}^{\otimes k}\) 模去双边理想 \(J = \langle X \otimes Y - Y \otimes X - [X, Y] : X, Y \in \mathfrak{g} \rangle\) 的商代数：\(U(\mathfrak{g}) = T(\mathfrak{g}) / J\)。

\textbf{定理 123.1}（泛性质）：存在自然李代数同态 \(i: \mathfrak{g} \to U(\mathfrak{g})\)（将 \(X\) 映为其在 \(T(\mathfrak{g})\) 中的像），使得对任意结合代数 \(A\) 和李代数同态 \(f: \mathfrak{g} \to A\)（\(A\) 带有交换子 \([a,b] = ab - ba\) 作为李括号），存在唯一的结合代数同态 \(\tilde{f}: U(\mathfrak{g}) \to A\) 使得 \(f = \tilde{f} \circ i\)。即 \(\operatorname{Hom}_{\text{Lie}}(\mathfrak{g}, A) \cong \operatorname{Hom}_{\text{Assoc}}(U(\mathfrak{g}), A)\)。

\textbf{证明}：张量代数 \(T(\mathfrak{g})\) 的泛性质将线性映射 \(\mathfrak{g} \to A\) 唯一地延拓为结合代数同态 \(T(\mathfrak{g}) \to A\)。对李代数同态 \(f\)，其延拓 \(\hat{f}\) 满足 \(\hat{f}(X \otimes Y - Y \otimes X - [X, Y]) = f(X)f(Y) - f(Y)f(X) - f([X, Y]) = 0\)，故 \(J \subseteq \ker \hat{f}\)。因此 \(\hat{f}\) 通过商诱导 \(\tilde{f}: U(\mathfrak{g}) \to A\)。唯一性由 \(i(\mathfrak{g})\) 生成 \(U(\mathfrak{g})\) 保证。\(\blacksquare\)

\textbf{推论 123.1}（表示论的翻译）：\(\mathfrak{g}\) 的表示（即李代数同态 \(\mathfrak{g} \to \mathfrak{gl}(V)\)）与 \(U(\mathfrak{g})\)-模一一对应。这将对李代数表示的研究转化为对结合代数 \(U(\mathfrak{g})\) 上模的研究。

\textbf{证明}：由定理 123.1 的泛性质，对李代数同态 \(\rho: \mathfrak{g} \to \mathfrak{gl}(V)\)，存在唯一的结合代数同态 \(\tilde{\rho}: U(\mathfrak{g}) \to \operatorname{End}(V)\) 使 \(\tilde{\rho} \circ i = \rho\)，这赋予 \(V\) 以 \(U(\mathfrak{g})\)-模结构。反之，给定 \(U(\mathfrak{g})\)-模 \(V\)（即结合代数同态 \(\tilde{\rho}: U(\mathfrak{g}) \to \operatorname{End}(V)\)），复合 \(i\) 得 \(\rho = \tilde{\rho} \circ i: \mathfrak{g} \to \mathfrak{gl}(V)\) 为李代数表示。这两个构造互逆，故一一对应。\(\blacksquare\)

\subsection{123.2 PBW定理}\label{pbwux5b9aux7406}

\textbf{定理 123.2}（PBW定理，Poincaré 1900, Birkhoff 1937, Witt 1937）：设 \(\{X_1, \ldots, X_n\}\) 是李代数 \(\mathfrak{g}\) 的一组有序基。则单项式 \(\{X_1^{k_1} X_2^{k_2} \cdots X_n^{k_n} : k_i \in \mathbb{Z}_{\geq 0}\}\)（按顺序排列）构成 \(U(\mathfrak{g})\) 的一组基。

\textbf{证明}：构造线性映射 \(\varphi: S(\mathfrak{g}) \to U(\mathfrak{g})\)（对称代数到泛包络代数）。在对称代数上用归纳法定义 \(\varphi(x_1 \cdots x_m) = \frac{1}{m!} \sum_{\sigma \in S_m} X_{\sigma(1)} \cdots X_{\sigma(m)}\)（\(U(\mathfrak{g})\) 中的乘积）。验证 \(\varphi\) 将 \(S(\mathfrak{g})\) 的基（有序单项式）映为 \(U(\mathfrak{g})\) 中的有序单项式，且该映射是线性同构------因为 \(\varphi\) 将滤过 \(S(\mathfrak{g}) \to \operatorname{gr} U(\mathfrak{g})\) 诱导为分次同构。关键在于验证 \(\varphi(xy - yx) = [X, Y]\)（在 \(U(\mathfrak{g})\) 中），这由对称化算子的性质保证。\(\blacksquare\)

\emph{意义}：PBW定理表明 \(U(\mathfrak{g})\) 作为向量空间同构于对称代数 \(S(\mathfrak{g})\)（\(\mathfrak{g}\) 上的多项式代数）------尽管作为代数结构完全不同（\(U(\mathfrak{g})\) 是非交换的）。自然滤过 \(\{U_k(\mathfrak{g})\}\)（\(U_k\) 由次数 \(\leq k\) 的元素张成）的伴随分次代数 \(\operatorname{gr} U(\mathfrak{g}) \cong S(\mathfrak{g})\)。这意味着 \(\dim U(\mathfrak{g})\) 是无限的（当 \(\mathfrak{g} \neq \{0\}\)），且 \(U(\mathfrak{g})\) 是 Noether 环。

\textbf{定理 123.3}（Bernstein-Gelfand-Gelfand 范畴 \(\mathcal{O}\) 定理，1976）：对半单李代数 \(\mathfrak{g}\)，BGG 范畴 \(\mathcal{O}\)（由所有有限生成的 \(U(\mathfrak{g})\)-模构成的范畴，其中 \(\mathfrak{n}^+\) 局部幂零作用且 \(\mathfrak{h}\) 半单作用）具有有限长、每个对象有有限 Jordan-Hölder 序列，且不可约对象由最高权分类。投射对象的存在性和 BGG 互反律（BGG reciprocity）是范畴 \(\mathcal{O}\) 的核心结构定理。

\textbf{证明}：范畴 \(\mathcal{O}\) 的定义条件确保每个对象 \(M\) 具有权空间分解 \(M = \bigoplus_{\mu \in \mathfrak{h}^*} M_\mu\)。由 \(\mathfrak{n}^+\) 的局部幂零性和 \(\mathfrak{h}\) 的半单作用，每个非零对象 \(M\) 中存在最高权向量 \(v \in M_\lambda\)（\(\mathfrak{n}^+ \cdot v = 0\)），从而存在 Verma 模的满射 \(\Delta(\lambda) = U(\mathfrak{g}) \otimes_{U(\mathfrak{b})} \mathbb{C}_\lambda \twoheadrightarrow M\)。\(U(\mathfrak{g})\) 是 Noether 环（由 PBW 定理），故 \(\mathcal{O}\) 中每个有限生成对象也是 Noether 的，因而具有有限长度。不可约对象 \(L(\lambda)\) 由最高权 \(\lambda \in \mathfrak{h}^*\) 唯一分类：\(L(\lambda)\) 是 \(\Delta(\lambda)\) 的唯一不可约商。投射覆盖 \(P(\lambda) \to L(\lambda)\) 的存在性由 \(U(\mathfrak{g})\) 上适当的完备化（或 \(\mathcal{O}\) 具有足够多投射对象）保证。BGG 互反律 \((P(\lambda) : \Delta(\mu)) = [\Delta(\mu) : L(\lambda)]\) 由 Verma 模的合成因子重数与投射模的滤过之间的对偶性得出。\(\blacksquare\)

\subsection{123.3 Casimir算子与中心}\label{casimirux7b97ux5b50ux4e0eux4e2dux5fc3}

\textbf{定义 123.2}（Casimir算子）：设 \(\mathfrak{g}\) 是半单李代数，\(\{X_i\}\) 是 \(\mathfrak{g}\) 关于 Killing 型的一组基，\(\{X^i\}\) 是对偶基（\(\kappa(X_i, X^j) = \delta_i^j\)）。\textbf{Casimir算子} \(C = \sum_i X_i X^i \in U(\mathfrak{g})\)。\(C\) 位于 \(U(\mathfrak{g})\) 的中心 \(\mathcal{Z}(\mathfrak{g})\) 中，且与表示 \(\pi\) 的选取无关。

\textbf{命题 123.1}（Casimir算子的作用）：在最高权为 \(\lambda\) 的不可约表示 \(V_\lambda\) 上，Casimir 算子作用为标量 \(c(\lambda) = \kappa(\lambda + \rho, \lambda + \rho) - \kappa(\rho, \rho)\)，其中 \(\rho\) 是正根之和的一半。

\textbf{证明}：在 \(V_\lambda\) 中取最高权向量 \(v_\lambda\)。对 \(C = \sum_i X_i X^i\) 计算 \(C \cdot v_\lambda\)。将基 \(\{X_i\}\) 分解为 \(\mathfrak{h}\) 的基 \(\{H_i\}\)、正根向量 \(\{E_\alpha\}\) 和负根向量 \(\{F_\alpha\}\)。由于 \(E_\alpha \cdot v_\lambda = 0\)，\(C\) 在 \(v_\lambda\) 上的作用简化为 \(\sum H_i H^i \cdot v_\lambda + \sum [E_\alpha, F_\alpha] \cdot v_\lambda = \kappa(\lambda, \lambda) v_\lambda + \sum \kappa(\alpha, \lambda) v_\lambda = (\kappa(\lambda, \lambda) + 2\kappa(\lambda, \rho)) v_\lambda = \kappa(\lambda+\rho, \lambda+\rho) v_\lambda - \kappa(\rho, \rho) v_\lambda\)。由于 \(C\) 在 \(U(\mathfrak{g})\) 的中心中，它在不可约表示 \(V_\lambda\) 上作用为标量，故在整个 \(V_\lambda\) 上 \(C = c(\lambda) \operatorname{id}\)。\(\blacksquare\)

\textbf{定理 123.4}（Harish-Chandra 同构，1951）：中心 \(\mathcal{Z}(\mathfrak{g}) \cong U(\mathfrak{h})^W\)------即 \(U(\mathfrak{g})\) 的中心同构于 Cartan 子代数的泛包络代数在 Weyl 群作用下的不变量。该同构由 Harish-Chandra 投影 \(P: U(\mathfrak{g}) \to U(\mathfrak{h})\)（滤过 \(U(\mathfrak{g}) = U(\mathfrak{h}) \oplus (\mathfrak{n}^- U(\mathfrak{g}) + U(\mathfrak{g}) \mathfrak{n}^+)\) 的前项投影）的标准万有实现给出。这一定理将中心特征标与权 \(\lambda\) 通过 \(\chi_\lambda: \mathcal{Z}(\mathfrak{g}) \to \mathbb{C}\)（即无穷小特征标）对应起来------两个不可约最高权模具有相同的无穷小特征标当且仅当它们的最高权相差 Weyl 群作用。

\textbf{证明}：由 PBW 定理的三角分解 \(U(\mathfrak{g}) \cong U(\mathfrak{n}^-) \otimes U(\mathfrak{h}) \otimes U(\mathfrak{n}^+)\)。定义 Harish-Chandra 投影 \(P: U(\mathfrak{g}) \to U(\mathfrak{h})\) 为滤去含 \(\mathfrak{n}^\pm\) 因子的项：对 \(u \in U(\mathfrak{g})\)，唯一地写 \(u = u_{\mathfrak{h}} + \sum u_- \cdot u_+\)（\(u_- \in U(\mathfrak{n}^-)\mathfrak{n}^-\)，\(u_+ \in \mathfrak{n}^+ U(\mathfrak{n}^+)\)），令 \(P(u) = u_{\mathfrak{h}}\)。限制到中心 \(\mathcal{Z}(\mathfrak{g})\)：对 \(z \in \mathcal{Z}(\mathfrak{g})\) 和任意根向量 \(e_\alpha \in \mathfrak{n}^+\)，\([z, e_\alpha] = 0\)。展开此等式并应用 \(P\)，得 \(P(z)\) 在 Weyl 群单反射 \(s_\alpha\) 下不变：\(s_\alpha(P(z)) = P(z)\)。由 Weyl 群由单反射生成，\(P(z) \in U(\mathfrak{h})^W\)。单射性：若 \(P(z) = 0\) 且 \(z \neq 0\)，则 \(z\) 可写为 \(\mathfrak{n}^-\) 或 \(\mathfrak{n}^+\) 元素的非零组合，但 \(z\) 与 \(\mathfrak{h}\) 交换迫使 \(z\) 的权为零，与 \(z \neq 0\) 矛盾。满射性：对每个 Weyl 不变量，可通过 Casimir 元素和更高次中心元的对称多项式构造原像。\(\blacksquare\)

\hrulefill

\hrulefill

\hrulefill

\hrulefill

\hrulefill

\hrulefill

\section{第149章：Borel-Weil-Bott定理与几何实现}\label{ux7b2c149ux7ae0borel-weil-bottux5b9aux7406ux4e0eux51e0ux4f55ux5b9eux73b0}

Borel-Weil-Bott 定理将紧李群的不可约表示的构造几何化：将表示实现为全纯线丛的截面空间的上同调。

\subsection{124.1 旗流形与线丛}\label{ux65d7ux6d41ux5f62ux4e0eux7ebfux4e1b}

\textbf{定义 124.1}（旗流形）：复半单李群 \(G_{\mathbb{C}}\) 的 \textbf{Borel 子群} \(B\) 是极大可解连通子群。\textbf{旗流形}（广义旗流形）是齐性空间 \(\mathcal{B} = G_{\mathbb{C}} / B\)。它是紧 Kähler 流形，也是光滑射影代数簇。

旗流形 \(\mathcal{B}\) 的几何结构极为丰富。其维数为 \(|\Phi^+|\)（正根的个数），它可被分解为 Schubert 胞腔 \(X_w^\circ = BwB/B\)（\(w \in W\)），每个胞腔同构于 \(\mathbb{C}^{\ell(w)}\)，给出 \(\mathcal{B}\) 的胞腔分解。旗流形上的 Bruhat 序描述了这些胞腔的闭包 \(X_w = \overline{X_w^\circ}\)（Schubert 簇）之间的包含关系。在复半单李群 \(G_{\mathbb{C}} = SL(n, \mathbb{C})\) 的情形，\(\mathcal{B}\) 是 \(\mathbb{C}^n\) 中完备旗的全体： \[\mathcal{B} = \{0 = V_0 \subset V_1 \subset \cdots \subset V_n = \mathbb{C}^n : \dim V_i = i\}\] 它也是广义旗流形 \(G_{\mathbb{C}} / P\)（\(P\) 为抛物子群）的原型。Borel 子群 \(B\) 的选取对应于 Cartan 子代数 \(\mathfrak{h}\) 和正根系 \(\Phi^+\) 的选取，不同的选取给出同构的旗流形。

\textbf{定义 124.2}（与权关联的线丛）：对任意整权 \(\lambda \in \mathfrak{h}^*\)（满足 \(\lambda\) 在每条余根上取整数值），可定义 \(B\) 的 1 维表示 \(\mathbb{C}_\lambda\)（由 \(\lambda\) 经指数映射诱导），进而构造 \(\mathcal{B}\) 上的齐性线丛 \(\mathcal{L}_\lambda = G_{\mathbb{C}} \times_B \mathbb{C}_\lambda\)。\(\mathcal{L}_\lambda\) 的全纯截面在 \(\mathcal{B}\) 的 Dolbeault 上同调中给出了不可约表示。

线丛 \(\mathcal{L}_\lambda\) 的构造体现了表示论与复几何的深刻统一。这里 \(B\) 通过特征标 \(\chi_\lambda: B \to \mathbb{C}^\times\) 作用于 \(\mathbb{C}_\lambda\)（\(b \cdot z = \chi_\lambda(b) z\)），而 \(\mathcal{L}_\lambda\) 在 \(G_{\mathbb{C}} \times_B \mathbb{C}_\lambda\) 中的等价类由 \((gb, z) \sim (g, \chi_\lambda(b) z)\) 定义。\(\mathcal{L}_\lambda\) 的截面 \(s: \mathcal{B} \to \mathcal{L}_\lambda\) 等价于函数 \(f: G_{\mathbb{C}} \to \mathbb{C}\) 满足 \(f(gb) = \chi_\lambda(b)^{-1} f(g)\)。当 \(\lambda\) 为支配整权时，\(\mathcal{L}_\lambda\) 是丰沛线丛（ample line bundle），其上同调 \(H^0(\mathcal{B}, \mathcal{O}(\mathcal{L}_\lambda))\) 给出不可约有限维表示；当 \(\lambda\) 为非支配整权时，上同调集中在更高次，由 Bott-Borel-Weil 定理（定理 124.2）精确描述。这一对应------将代数表示论完全转化为复几何中的层上同调计算------是几何表示论的起点。

\subsection{124.2 Borel-Weil-Bott定理}\label{borel-weil-bottux5b9aux7406}

\textbf{定理 124.1}（Borel-Weil 定理，1954-1957）：设 \(\lambda\) 是支配整权。则不可约表示 \(V_\lambda\)（\(G\) 的表示，经复化视为 \(G_{\mathbb{C}}\) 的全纯表示）同构于 \(\mathcal{B}\) 上 \(\mathcal{L}_\lambda\) 的全局全纯截面空间： \[H^0(\mathcal{B}, \mathcal{O}(\mathcal{L}_\lambda)) \cong V_\lambda\]

\textbf{证明}：取 \(G\) 的紧致实形 \(G_c\)。由 Peter-Weyl 定理，\(G_c\) 在 \(\mathcal{B} \cong G_c/T\) 上的正则表示分解为 \(\bigoplus_\lambda V_\lambda \otimes V_\lambda^*\)（\(\lambda\) 遍历支配整权）。该分解对应 \(G_c/T\) 上的函数空间。\(V_\lambda\) 中的最高权向量 \(v_\lambda\)（满足 \(B v_\lambda = \chi_\lambda v_\lambda\)，\(\chi_\lambda\) 为 \(T\) 的特征标）给出 \(G_c/T\) 上的截面的全局截面，该截面恰为 \(\mathcal{L}_\lambda\) 的截面。利用 Casimir 算子 \(\Omega = \sum_i X_i X^i\) 在 \(H^0(\mathcal{B}, \mathcal{O}(\mathcal{L}_\lambda))\) 上的作用：\(\Omega\) 的特征值 \(\langle \lambda, \lambda + 2\rho \rangle\) 唯一刻画了不可约表示 \(V_\lambda\)。由 \(G_c\) 在 \(\mathcal{B}\) 上的可迁性，\(\mathcal{L}_\lambda\) 的截面空间构成 \(G_c\) 的表示，且最高权为 \(\lambda\)。由最高权表示的唯一性，\(H^0(\mathcal{B}, \mathcal{O}(\mathcal{L}_\lambda)) \cong V_\lambda\)。\(\blacksquare\)

\textbf{定理 124.2}（Bott-Borel-Weil 定理，Bott 1957）：对任意整权 \(\lambda\)（不必是支配的），\(\mathcal{L}_\lambda\) 的上同调群 \(H^q(\mathcal{B}, \mathcal{O}(\mathcal{L}_\lambda))\) 要么为零（对所有 \(q\)），要么恰在一个 \(q = \ell(w)\) 处非零，其中 \(w \in W\) 是使 \(w(\lambda + \rho) - \rho\) 为支配权的唯一 Weyl 群元素（\(\ell(w)\) 为 \(w\) 的长度）。非零时 \(H^{\ell(w)}(\mathcal{B}, \mathcal{O}(\mathcal{L}_\lambda)) \cong V_{w(\lambda + \rho) - \rho}\)。当 \(\lambda\) 已是支配权时，\(w = e\)，回到 Borel-Weil 定理。

\textbf{证明}（框架）：使用 Schubert 胞腔分层和 Demazure 算符。对每个单根 \(\alpha\)，考虑 \(\mathbb{P}^1\)-纤维化 \(\pi_\alpha: G/B \to G/P_\alpha\)。Leray 谱序列给出 \(H^q(G/B, \mathcal{L}_\lambda)\) 与 \(H^{q \pm 1}(G/P_\alpha, \cdots)\) 的关系。Demazure 算符 \(D_\alpha(e^\mu) = (e^\mu - e^{s_\alpha(\mu)-\alpha})/(1-e^{-\alpha})\) 描述了上同调在穿越 \(\alpha\)-壁时的转移。通过对 \(|\ell(w)|\) 归纳，将一般权 \(\lambda\) 逐次约化到支配权情形（定理 124.1）。若 \(\lambda+\rho\) 奇异，则上同调全消没。\(\blacksquare\)

\textbf{推论 124.1}（Euler-Poincaré特征）：形式特征标满足 \(\chi_\lambda = \sum_{q} (-1)^q \operatorname{ch} H^q(\mathcal{B}, \mathcal{O}(\mathcal{L}_\lambda))\)。结合 Weyl 特征公式，Bott-Borel-Weil 定理给出了特征公式的几何证明。

\subsection{124.3 几何表示论的视角}\label{ux51e0ux4f55ux8868ux793aux8bbaux7684ux89c6ux89d2}

\textbf{定义 124.3}（等变上同调与特征类）：\(T\)-等变上同调 \(H_T^*(\mathcal{B})\) 中的 Schubert 类 \([X_w]\)（\(w \in W\)）构成 \(H_T^*(\mathcal{B})\) 的基。Schubert 类的乘积公式（Chevalley-Monk 公式）与表示论中张量积分解有深刻联系。

\textbf{定理 124.3}（Borel 同构）：\(H_T^*(\mathcal{B}) \cong \mathbb{C}[\mathfrak{h}] / (\mathbb{C}[\mathfrak{h}]^W_+)\)，其中 \(\mathbb{C}[\mathfrak{h}]^W_+\) 是 Weyl 群不变多项式中无常数项的理想。这给出了旗流形的等变上同调环的代数描述。

\textbf{证明}：\(T\)-等变上同调 \(H_T^*(\mathcal{B})\) 是 \(H_T^*(\mathrm{pt}) = \mathbb{C}[\mathfrak{h}]\) 上的代数。Borel 构造 \(ET \times_T \mathcal{B}\) 给出纤维化 \(\mathcal{B} \to ET \times_T \mathcal{B} \to BT\)，其 Serre 谱序列在 \(E_2\) 页为 \(H^*(BT) \otimes H^*(\mathcal{B})\)。由于 \(H^*(\mathcal{B})\) 集中于偶数维（Schubert 类），谱序列退化，故 \(H_T^*(\mathcal{B}) \cong \mathbb{C}[\mathfrak{h}] \otimes H^*(\mathcal{B})\) 作为 \(\mathbb{C}[\mathfrak{h}]\)-模。Weyl 群作用给出 \(H_T^*(\mathcal{B})^W \cong \mathbb{C}[\mathfrak{h}]^W \otimes H^*(\mathcal{B})^W\)，结合 \(H^*(\mathcal{B}) \cong \mathbb{C}[\mathfrak{h}]/(\mathbb{C}[\mathfrak{h}]^W_+)\)（Borel 描述的通常上同调）推得等变版本。\(\blacksquare\)

\hrulefill

\hrulefill

\hrulefill

\hrulefill

\hrulefill

\hrulefill

\section{第150章：Clifford 代数与旋量}\label{ux7b2c150ux7ae0clifford-ux4ee3ux6570ux4e0eux65cbux91cf}

Clifford 代数（W. K. Clifford, 1878）是平方型向量空间上最自然的非交换代数结构，将外代数和二次型信息融为一体。

\textbf{定义 78.1}（Clifford 代数）：设 \(V\) 为 \(\mathbb{F}\) 上的向量空间配以二次型 \(Q: V \to \mathbb{F}\)。Clifford 代数 \(\operatorname{Cl}(V, Q)\) 是张量代数 \(T(V)\) 模关系 \(v \otimes v - Q(v) \cdot 1 = 0\)（\(\forall v \in V\)）的商。等价刻画：\(\operatorname{Cl}(V,Q)\) 是满足泛性质------任意 \(\mathbb{F}\)-代数 \(A\) 和线性映射 \(f: V \to A\) 使 \(f(v)^2 = Q(v) \cdot 1_A\)，可唯一扩充为代数同态 \(\tilde{f}: \operatorname{Cl}(V,Q) \to A\)------的 \(\mathbb{F}\)-代数。

\textbf{定理 78.1}（分类与维数）：若 \(\dim V = n\) 且 \(Q\) 非退化，则 \(\operatorname{Cl}(V,Q) \cong \begin{cases} \mathbb{F} & n=0 \\ \mathbb{F} \oplus \mathbb{F} \text{ 或 } \mathbb{F}(\sqrt{a}) & n=1 \\ M_{2^k}(D) \text{ 或 } M_{2^{k-1}}(D \oplus D') & n \geq 2 \end{cases}\)，其中除环 \(D\) 由 \(\mathbb{F}\) 上的 Brauer 类决定。维数为 \(\dim \operatorname{Cl}(V,Q) = 2^n\)。

\textbf{低维实 Clifford 代数}：\(\operatorname{Cl}_{0,2}(\mathbb{R}) \cong \mathbb{H}\)（四元数）；\(\operatorname{Cl}_{3,0}(\mathbb{R}) \cong \mathbb{H} \oplus \mathbb{H}\)；\(\operatorname{Cl}_{1,3}(\mathbb{R}) \cong M_2(\mathbb{H})\)（与旋量群 \(\operatorname{Spin}(1,3)\) 相关）。Spin 群的定义为 \(\operatorname{Spin}(V,Q) = \{v_1 \cdots v_{2k} \in \operatorname{Cl}^0(V,Q) : Q(v_i) = \pm 1\}\)，是 \(\operatorname{SO}(V,Q)\) 的二重覆叠群------这是描述 \(\frac{1}{2}\)-自旋粒子（旋量表示）的严格代数工具。

Clifford 代数在 Dirac 方程（量子场论中描述费米子）、Atiyah-Singer 指标定理的算子表述、以及 \(K\)-理论的 Bott 周期性定理中有核心地位。

\textbf{证明}：\(n=0\) 时 \(\operatorname{Cl}(V,Q) = \mathbb{F}\)。\(n=1\) 时取 \(V = \mathbb{F} e\)，\(Q(e) = a \neq 0\)，则 \(e^2 = a\)，故 \(\operatorname{Cl}(V,Q) \cong \mathbb{F}[x]/(x^2-a)\)。若 \(a\) 为平方数则 \(\cong \mathbb{F} \oplus \mathbb{F}\)，否则 \(\cong \mathbb{F}(\sqrt{a})\)。\(n \geq 2\) 时，利用 Clifford 代数的 \(\mathbb{Z}_2\)-分次结构和 Wall 的 Brauer 群分类：\(\operatorname{Cl}(V,Q)\) 为 \(\mathbb{Z}_2\)-分次中心单代数，其 Brauer 类由 \(Q\) 的 Hasse-Witt 不变量决定。\(\dim \operatorname{Cl}(V,Q) = 2^n\) 由基 \(\{e_{i_1} \cdots e_{i_k} : i_1 < \cdots < i_k\}\) 直接计算得到。\(\operatorname{Spin}(V,Q)\) 为 \(\operatorname{Cl}^0(V,Q)\) 中满足 \(v_1 \cdots v_{2k}\) 且 \(Q(v_i) = \pm 1\) 的元素，到 \(\operatorname{SO}(V,Q)\) 的映射 \(x \mapsto (v \mapsto x v x^{-1})\) 为二重覆叠。 \(\blacksquare\)

\subsection{Clifford代数与旋量理论的现代应用}\label{cliffordux4ee3ux6570ux4e0eux65cbux91cfux7406ux8bbaux7684ux73b0ux4ee3ux5e94ux7528}

Clifford 代数在 20 世纪数学和物理中发挥了核心作用，其影响远超最初的外代数与二次型的结合。\textbf{Bott 周期性定理}（1970）的证明借助 Clifford 代数给出了最简洁的表述：\(\operatorname{Cl}_{p,q+8} \cong M_{16}(\operatorname{Cl}_{p,q})\)，即 Clifford 代数在 8 个维数上具有周期性。这一周期性直接导致拓扑 K-理论的 Bott 周期性 \(K^{-n}(X) \cong K^{-n-8}(X)\)，并推广到 KR-理论和弦论中的 D-膜分类。\textbf{Atiyah-Bott-Shapiro 构造}（1964）利用 Clifford 模将 Spin 流形上的 Dirac 算子与 K-理论中的 Thom 同构联系起来，这是 Atiyah-Singer 指标定理证明的核心技术环节。\textbf{旋量丛}（spinor bundle）\(\mathcal{S}\) 在 Riemann 几何中定义为 \(\operatorname{Spin}\)-主丛与旋量表示的关联丛，Dirac 算子 \(D: \Gamma(\mathcal{S}) \to \Gamma(\mathcal{S})\) 是一阶椭圆微分算子，其平方 \(D^2 = \nabla^*\nabla + \frac{1}{4}R\)（Lichnerowicz 公式）将曲率与 Laplacian 联系起来，阻碍了正标量曲率度量的存在。\textbf{Gromov-Lawson 和 Schoen-Yau 定理}（1980s）利用 Dirac 算子证明了某些流形（如环面 \(T^n\)）不容许正标量曲率度量，这是几何分析中最重要的刚性结果之一。在\textbf{表示论}中，Clifford 代数与 Weyl 代数的关系（通过 \([x_i, \partial_j] = \delta_{ij}\) 和 \([e_i, e_j] = 2\delta_{ij}\) 的对比）揭示了两类基本非交换代数之间的深层联系，并在 Dubrovin 的 Frobenius 簇理论和量子上同调中扮演基本角色。\textbf{超对称量子力学}中的超荷 \(Q\) 满足 \(Q^2 = H\)（\(H\) 为 Hamilton 量），这正是 Clifford 代数中 \(e^2 = Q(e)\) 的物理实现。在\textbf{引力理论}中，标架形式（vielbein formalism）和自旋联络（spin connection）的核心代数结构正是 Clifford 代数，它提供了将旋量场耦合到引力场的唯一相容方式。Clifford 代数还在\textbf{Atiyah-Singer 指标定理}的 Getzler 证明（1983）中扮演关键角色：Getzler 利用 Clifford 代数在 \(\mathbb{R}^n\) 上的渐近展开，将热核的超迹（supertrace）归约为 Clifford 代数的超迹（Berezin 积分），从而将 A-roof 类与 Dirac 算子的指标联系起来。

\hrulefill

\newpage

\chapter{05-代数/04-李理论/06-结晶群与准齐性向量空间.md}\label{ux4ee3ux657004-ux674eux7406ux8bba06-ux7ed3ux6676ux7fa4ux4e0eux51c6ux9f50ux6027ux5411ux91cfux7a7aux95f4.md}

\section{第151章：结晶群（Crystallographic Groups）}\label{ux7b2c151ux7ae0ux7ed3ux6676ux7fa4crystallographic-groups}

结晶群是 \(n\) 维 Euclid 空间 \(E(n) = \mathbb{R}^n \rtimes O(n)\) 的离散余紧子群，其理论由 Bieberbach 在 1911-1912 年间奠基。结晶群不仅是几何群论的基本对象，也是数学晶体学、固态物理和材料科学的数学基础------自然界中的晶体结构恰好对应 \(n = 3\) 的结晶群（230 种空间群）。

\subsection{75.1 基本定义与 Bieberbach 定理}\label{ux57faux672cux5b9aux4e49ux4e0e-bieberbach-ux5b9aux7406}

\textbf{定义 75.1}（Euclid 运动群）：\(n\) 维 \textbf{Euclid 运动群} \(E(n)\) 是 \(\mathbb{R}^n\) 的所有等距变换构成的群。由定义 \(E(n) = \mathbb{R}^n \rtimes O(n)\)，其中 \(\mathbb{R}^n\) 为平移子群（正规子群），\(O(n)\) 为正交群。元记为 \((v, A)\)（\(v \in \mathbb{R}^n, A \in O(n)\)），作用为 \(x \mapsto Ax + v\)。乘法为

\[(v_1, A_1)(v_2, A_2) = (v_1 + A_1 v_2, A_1 A_2)\]

\textbf{定义 75.2}（结晶群）：\(n\) 维 \textbf{结晶群}（crystallographic group）是 \(E(n)\) 的离散子群 \(\Gamma\)，满足商空间 \(\mathbb{R}^n / \Gamma\) 是紧的。\(n = 2\) 时称 \textbf{wallpaper group}（壁纸群），\(n = 3\) 时称\textbf{空间群}（space group）。

Euclid 运动群的定义给出了结晶群的''活动空间''，而结晶群的定义则规定了其离散性和紧商条件。这两个条件表面上属于几何约束，但 Bieberbach 在 1911 年的第一定理揭示了一个惊人的事实：结晶群的平移部分必然构成 \(\mathbb{R}^n\) 中的一个秩为 \(n\) 的格------这意味着离散几何约束迫使群结构中包含一个最大可能的自由 Abel 子群。

\textbf{定理 75.1}（Bieberbach 第一定理，1911）：设 \(\Gamma \subset E(n)\) 是 \(n\) 维结晶群。则 \(\Gamma\) 的平移子群 \(T = \Gamma \cap \mathbb{R}^n\)（其中 \(\mathbb{R}^n\) 等同于纯平移子群）是秩为 \(n\) 的自由 Abel 群，且 \([\Gamma : T] < \infty\)。

\textbf{证明}：设 \(\Gamma \subset E(n) = \mathbb{R}^n \rtimes O(n)\) 为 \(n\) 维结晶群。令 \(T = \Gamma \cap \mathbb{R}^n\) 为平移子群。若 \(T\) 不是 \(\mathbb{R}^n\) 的格，则存在 \(T\) 生成的子空间 \(V \subsetneq \mathbb{R}^n\) 维数 \(< n\)。考虑投影 \(\pi: \mathbb{R}^n \to \mathbb{R}^n / V\)，\(\Gamma\) 在 \(\mathbb{R}^n / V\) 上的作用诱导紧商。但 \(\Gamma\) 在 \(\mathbb{R}^n / V\) 上的平移部分平凡，导致 \(\mathbb{R}^n/\Gamma\) 非紧，矛盾。故 \(T\) 为秩 \(n\) 的格。点群 \(P = \Gamma / T\) 为 \(O(n)\) 的子群，由 \(\mathbb{R}^n/\Gamma\) 的紧性知 \(P\) 为有限群。正合列 \(1 \to T \to \Gamma \to P \to 1\) 给出 \(\Gamma\) 的结构。 \(\blacksquare\)

\textbf{定理 75.2}（Bieberbach 第二定理，1911，刚性定理）：若 \(\Gamma_1, \Gamma_2 \subset E(n)\) 是两个同构的结晶群，则存在仿射变换 \(\varphi \in \operatorname{Aff}(\mathbb{R}^n)\) 使得 \(\varphi \Gamma_1 \varphi^{-1} = \Gamma_2\)。

\textbf{证明}：设 \(\phi: \Gamma_1 \to \Gamma_2\) 为同构，由第一定理 \(T_i = \Gamma_i \cap \mathbb{R}^n\) 为秩 \(n\) 的格。\(\phi\) 将 \(T_1\) 映到 \(\Gamma_2\) 的正规 Abel 子群，由结晶群结构理论知 \(\phi(T_1) = T_2\)。\(\phi|_{T_1}: T_1 \to T_2\) 诱导 \(\mathbb{R}^n\) 上的线性同构 \(\tilde{\phi}\)。将 \(\tilde{\phi}\) 扩展为 \(\mathbb{R}^n\) 的线性变换，结合适当平移得 \(\varphi \in \operatorname{Aff}(\mathbb{R}^n)\)。对任意 \(\gamma \in \Gamma_1\)，\(\varphi \gamma \varphi^{-1}\) 与 \(\phi(\gamma)\) 在平移部分一致，且 \(\varphi T_1 \varphi^{-1} = T_2\)，故 \(\varphi \Gamma_1 \varphi^{-1} = \Gamma_2\)。 \(\blacksquare\)

换言之，结晶群的同构类与仿射共轭类一致------不存在同构但不等价的结晶群。这在几何群论中是一个罕见的刚性现象。

\textbf{定理 75.3}（Bieberbach 第三定理，1912）：对每个维数 \(n\)，只存在有限多个 \(n\) 维结晶群的同构类（等价地，仿射共轭类）。

\textbf{证明}：由第一定理，结晶群 \(\Gamma\) 有正合列 \(1 \to \mathbb{Z}^n \to \Gamma \to P \to 1\)，其中 \(P\) 为 \(O(n)\) 的有限子群。\(P\) 在 \(\mathbb{Z}^n\) 上的作用由其几何实现决定，给出 \(P \hookrightarrow GL(n, \mathbb{Z})\)。Jordan-Zassenhaus 定理：\(GL(n, \mathbb{Z})\) 的有限子群仅有有限多个共轭类。对每个 \(P\)，结晶群 \(\Gamma\) 的同构类由 \(H^2(P, \mathbb{Z}^n)\) 分类。因 \(P\) 有限且 \(\mathbb{Z}^n\) 为有限生成 \(P\)-模，\(H^2(P, \mathbb{Z}^n)\) 为有限群。故对每个 \(P\) 仅有有限多个 \(\Gamma\)，乘以仅有有限多个 \(P\) 的共轭类，得总数为有限。 \(\blacksquare\)

\subsection{75.2 Zassenhaus 算法与空间群分类}\label{zassenhaus-ux7b97ux6cd5ux4e0eux7a7aux95f4ux7fa4ux5206ux7c7b}

\textbf{定义 75.3}（Zassenhaus 算法）：Zassenhaus（1948 年）提出了一种系统方法列举所有 \(n\) 维空间群。算法基于 Bieberbach 定理：

\begin{enumerate}
\def\labelenumi{\arabic{enumi}.}
\tightlist
\item
  列举 \(GL(n, \mathbb{Z})\) 的所有有限子群（\(n \leq 3\) 时已知分类）；
\item
  对每个有限子群 \(P\)，列举 \(H^2(P, \mathbb{Z}^n)\) 的元素（群扩张类）；
\item
  对每个扩张类中代表元，检查其在 \(E(n)\) 中的嵌入（判别离散性）。
\end{enumerate}

\textbf{定理 75.4}（空间群分类）： - \(n = 1\)：2 种（\(\mathbb{Z}\) 和无限二面体群 \(D_\infty\)）； - \(n = 2\)：17 种 wallpaper groups（由 Fedorov 1891 和 Pólya 1924 年完成）； - \(n = 3\)：219 种空间群（若不区分对映体），或 230 种（区分对映体后。Fedorov-Schoenflies-Barlow，1890 年代）； - \(n = 4\)：4783 种（Brown-Bülow-Neubüser-Wondratschek-Zassenhaus，1978 年）； - \(n = 5\)：222018 种（Plesken-Schulz，2000 年）； - \(n = 6\)：28927922 种（Oppenorth，2000 年之后）。

\(n = 3\) 的 230 种空间群是国际晶体学表（International Tables for Crystallography）的基础，直接用于 X 射线晶体学中晶体结构的确定。

\textbf{证明}：由 Bieberbach 定理，空间群由正合列 \(1 \to \mathbb{Z}^n \to \Gamma \to P \to 1\) 决定，\(P\) 为 \(GL(n, \mathbb{Z})\) 的有限子群。\(n=1\) 时 \(P \subset GL(1, \mathbb{Z}) = \{\pm 1\}\)，给出 \(\mathbb{Z}\) 和 \(D_\infty\)。\(n=2\) 时 \(GL(2, \mathbb{Z})\) 有 13 个有限子群共轭类，\(H^2\) 计算给出 17 种。\(n=3\) 时 \(GL(3, \mathbb{Z})\) 有 73 个有限子群共轭类（32 个晶体点群），\(H^2\) 计算给出 219/230 种。高维由计算机代数系统遍历 \(GL(n, \mathbb{Z})\) 的有限子群共轭类并计算 \(H^2\) 得到。 \(\blacksquare\)

\subsection{\texorpdfstring{75.3 Wallpaper 群分类（\(n = 2\)）}{75.3 Wallpaper 群分类（n = 2）}}\label{wallpaper-ux7fa4ux5206ux7c7bn-2}

\textbf{定理 75.5}（Wallpaper 群的 17 种分类）：二维结晶群共有 17 个同构类，分类如下（按点群类型）：

\begin{longtable}[]{@{}llll@{}}
\toprule
点群 & 阶 & 群数 & 记号 \\
\midrule
\endhead
\bottomrule
\endlastfoot
\(C_1\) & 1 & 1 & p1 \\
\(C_2\) & 2 & 1 & p2 \\
\(C_3\) & 3 & 1 & p3 \\
\(C_4\) & 4 & 1 & p4 \\
\(C_6\) & 6 & 1 & p6 \\
\(D_1\) & 2 & 3 & pm, pg, cm \\
\(D_2\) & 4 & 4 & pmm, pmg, pgg, cmm \\
\(D_3\) & 6 & 2 & p3m1, p31m \\
\(D_4\) & 8 & 2 & p4m, p4g \\
\(D_6\) & 12 & 1 & p6m \\
\end{longtable}

\emph{代数刻画}：每个 wallpaper group 的平移格 \(\mathbb{Z}^2\) 配以 \(GL(2, \mathbb{Z})\) 的有限子群的作用（必须满足晶体学限制定理------旋转阶数只能为 1, 2, 3, 4 或 6）。不同 wallpaper groups 的区别在于 \(H^2(P, \mathbb{Z}^2)\) 的元素如何影响反射和滑动反射。

Wallpaper 群的 17 种分类是数学史上最优雅的枚举结果之一，由 Fedorov（1891）和 Polya（1924）独立完成。分类的关键在于两个层次的结构：\textbf{平移格}（\(\mathbb{Z}^2\) 共 5 种 Bravais 类型------斜方、长方、中心长方、正方、六角）和\textbf{点群}（\(GL(2, \mathbb{Z})\) 的有限子群，受晶体学限制至多 13 个共轭类）。给定格类型和点群 \(P\)，wallpaper 群的同构类由 \(H^2(P, \mathbb{Z}^2)\) 分类------这一上同调群刻画了 \(P\) 在格上的作用如何相对于平移进行''扭曲''。具体而言，\(H^2(P, \mathbb{Z}^2)\) 的非平凡元素对应滑动反射（glide reflection）------即反射后沿某方向平移的非平凡组合。在记号中，p 表示本原格（primitive），c 表示中心化格（centered），数字 1-6 表示最高旋转阶数，m 表示反射（mirror），g 表示滑动反射（glide）。例如，p4g 具有 4 阶旋转和滑动反射，而 p6m 是所有 wallpaper 群中对称性最高的。Wallpaper 群的分类在晶体学、材料科学和图案设计中具有直接应用，且其群论结构（通过群扩张和群上同调）为高维空间群（共 230 种三维空间群）的分类提供了方法论模板。

\textbf{证明}：\(n=2\) 时，二维格有 5 种 Bravais 类型。点群 \(P\) 为 \(GL(2, \mathbb{Z})\) 的有限子群，由结晶限制 \(P\) 的旋转阶数仅为 \(1, 2, 3, 4, 6\)。\(GL(2, \mathbb{Z})\) 的有限子群共 13 个共轭类。对每个 \((   ext{格类型}, P)\) 计算 \(H^2(P, \mathbb{Z}^2)\)：斜方格 \$ imes\$ 平凡 \$ o p1\$；斜方格 \$ imes C\_2 o p2\$；长方格 \$ imes\$ 反射 \$ o pm, pg, cm\$；正方格 \$ imes C\_4 o p4\$ 及 \$ imes D\_4 o p4m, p4g\$；六角格 \$ imes C\_3, C\_6, D\_3, D\_6 o p3, p6, p3m1, p31m, p6m\(；\)D\_2 imes\$ 长方格 \$ o pmm, pmg, pgg, cmm\$。总计 \(17\) 个同构类。 \(lacksquare\)

\subsection{75.4 与几何群论的联系}\label{ux4e0eux51e0ux4f55ux7fa4ux8bbaux7684ux8054ux7cfb}

\textbf{Grigorchuk 群}（Grigorchuk，1980 年）具有重要意义的例子：它是一个有限生成群，具有介于多项式和指数之间的中间增长（intermediate growth），是第一个被证明为非初等 amenable 群的例子。Grigorchuk 群的构造自然地与结晶群的分支覆盖联系：将结晶群作用在无限树上产生无限迭代自相似群。

结晶群理论为 Bieberbach 定理背后的几何刚性提供了群论解释。Hilton 将 wallpaper 群分类推广到一般的晶体学限制，而广义结晶群更是与 CAT(0) 空间等非正曲率几何理论交叉的前沿话题。

\subsection{结晶群与离散几何的深层联系}\label{ux7ed3ux6676ux7fa4ux4e0eux79bbux6563ux51e0ux4f55ux7684ux6df1ux5c42ux8054ux7cfb}

结晶群理论在 20 世纪后半叶与\textbf{几何群论}（geometric group theory）和\textbf{离散几何}（discrete geometry）产生了深刻的交叉。\textbf{几乎平坦流形}（almost flat manifold）理论（Gromov 1978）将 Bieberbach 定理推广到 Riemann 几何的框架：设 \(M\) 为紧 Riemann 流形，若其截面曲率 \(K\) 和直径 \(d\) 满足 \(|K| \cdot d^2 \leq \varepsilon_n\)（对某个仅依赖于 \(n\) 的常数 \(\varepsilon_n\)），则 \(M\) 微分同胚于某个结晶群的紧商 \(\mathbb{R}^n / \Gamma\)。这一结果是 Bieberbach 刚性定理的几何化版本，将结晶群从等距群推广到''几乎等距''的作用。\textbf{Freedman-Quinn 和 Farrell-Jones 猜想}与结晶群密切相关：结晶群的 Wall 障碍群的计算是手术理论（surgery theory）的核心问题，关系到拓扑流形分类的最终完成。\textbf{晶体学群在 Lorent 几何}中的推广提出了''晶体学时空群''的分类问题------这直接关系到广义相对论中的紧时空模型（如环面上的 Einstein 度量）。在\textbf{非交换几何}中，结晶群在非交换环面上的作用（通过非交换环面 \(A_\theta\) 的等距自同构群）导出了非交换晶体学群的分类，与量子 Hall 效应和拓扑绝缘体的数学理论紧密相关。\textbf{Hilbert-Smith 猜想}（与结晶群相关的）断言：若局部紧拓扑群 \(G\) 在连通 \(n\) 维流形上有效作用且 \(G\) 的 \(p\)-进整数子群 \(\mathbb{Z}_p\) 为闭子群，则 \(G\) 必为 Lie 群。结晶群理论的技术工具------群扩张分类、上同调计算、刚性定理------为这些现代问题提供了方法论基础。

\hrulefill

\hrulefill

\section{第152章：准齐性向量空间（Prehomogeneous Vector Spaces）}\label{ux7b2c152ux7ae0ux51c6ux9f50ux6027ux5411ux91cfux7a7aux95f4prehomogeneous-vector-spaces}

准齐性向量空间（Prehomogeneous Vector Spaces, PVS）由佐藤幹夫（Sato Mikio）于 1961 年引入，是代数群在向量空间上的作用具有稠密开轨道的结构。该理论连接了代数群表示论、微分方程理论和数论，是 Sato 的 D-模理论和超函数理论的重要支柱。

\subsection{76.1 基本概念与定义}\label{ux57faux672cux6982ux5ff5ux4e0eux5b9aux4e49}

\textbf{定义 76.1}（准齐性向量空间）：设 \(G\) 是代数闭域 \(k\) 上的连通代数群，\(\rho: G \to GL(V)\) 是 \(G\) 在有限维 \(k\)-向量空间 \(V\) 上的有理表示。三元组 \((G, \rho, V)\) 称为\textbf{准齐性向量空间}（PVS），若 \(V\) 中存在开的稠密 \(G\)-轨道。

等价地，存在点 \(v_0 \in V\) 使得轨道 \(G \cdot v_0\) 包含于 \(V\) 的 Zariski 开稠子集。此时 \(v_0\) 称为该 PVS 的\textbf{一般点}（generic point），其稳定子群 \(G_{v_0} = \{g \in G : g \cdot v_0 = v_0\}\) 是 \(G\) 的闭子群，且 \(\dim G - \dim G_{v_0} = \dim V\)。

\textbf{定义 76.2}（相对不变量）：PVS \((G, \rho, V)\) 上的非零有理函数 \(f \in k(V)^\times\) 称为\textbf{相对不变量}，若存在 \(G\) 的有理特征标 \(\chi: G \to \mathbb{G}_m\) 使得

\[f(g \cdot v) = \chi(g) f(v)\]

对所有 \(g \in G, v \in V\) 成立（在定义域内）。\(\chi\) 称为 \(f\) 的特征标。

相对不变量的零点定义了 \(G\)-不变的超曲面，这些超曲面由较低维的 \(G\)-轨道构成。具体地，若 \(V \setminus D\) 是 \(V\) 的唯一开稠密 \(G\)-轨道，则补集 \(D\) 是 \(G\)-不变超曲面的并，恰为相对不变量的零点集。

\subsection{\texorpdfstring{76.2 \(b\)-函数与 Bernstein-Sato 多项式}{76.2 b-函数与 Bernstein-Sato 多项式}}\label{b-ux51fdux6570ux4e0e-bernstein-sato-ux591aux9879ux5f0f}

\textbf{定义 76.3}（\(b\)-函数 / Bernstein-Sato 多项式）：设 \(f \in \mathbb{C}[x_1, \ldots, x_n]\) 是非零多项式。存在非零微分算子 \(P(s, \partial) \in \mathbb{C}[x, \partial_x, s]\) 和单变量多项式 \(b(s) \in \mathbb{C}[s]\) 使得

\[b(s) f(x)^s = P(s, \partial) f(x)^{s+1}\]

对所有 \(s \in \mathbb{C}\)（在 \(f \neq 0\) 的区域上）成立。满足此式的最小次数首一多项式 \(b(s)\) 称为 \(f\) 的 \textbf{Bernstein-Sato 多项式}或 \textbf{\(b\)-函数}。

\textbf{定理 76.1}（Bernstein-Sato，\(b\)-函数的存在性）：对任意非零多项式 \(f \in \mathbb{C}[x_1, \ldots, x_n]\)，其 \(b\)-函数存在。更一般地，Kashiwara（1976 年）利用 D-模理论证明了任意全纯函数的 \(b\)-函数总是存在的。

\textbf{证明}：构造 Weyl 代数 \(D_n = \mathbb{C}[x_1, \ldots, x_n, \partial_1, \ldots, \partial_n]\) 上的 \(D_n[s]\)-模 \(\mathcal{M} = D_n[s] f^s\)，其中 \(\partial_i \cdot f^s = s f^{-1} \partial_i f \cdot f^s\)。定义 Bernstein 算子 \(P(s) \in D_n[s]\) 满足 \(P(s) f^{s+1} = b(s) f^s\)。Kashiwara 证明了 \(\mathcal{M}\) 作为 \(\mathbb{C}[s]\)-模是有限生成的，故存在次数最低的消去多项式 \(b(s) \in \mathbb{C}[s]\) 和算子 \(P(s)\) 满足 \(P(s) f^{s+1} = b(s) f^s\)。\(b(s)\) 的唯一性由 \(\mathbb{C}[s]\) 是主理想整环保证。构造的关键是 \(f^s\) 的 D-模结构给出 \(\mathbb{C}[s]\) 上的良好有限性。 \(\blacksquare\)

\textbf{定理 76.2}（\(b\)-函数零点的性质）： 1. \(b(s)\) 的零点均为负有理数； 2. \(b(s)\) 的零点集合 \(\{\alpha\}\) 与 \(f\) 确定的奇点的局部 monodromy 的特征值（乘以 \(e^{2\pi i \alpha}\)）密切相关； 3. \(-1\) 是 \(b(s)\) 的零点 \(\Leftrightarrow\) \(f\) 不全纯。

\textbf{证明}：(1) 设 \(b(s)\) 的零点为 \(\alpha\)。考虑 Milnor 纤维 \(F_f\) 的 monodromy 算子 \(T\)，其半单部分 \(T_s\) 的特征值形如 \(\exp(2\pi i \alpha)\)。\(b(s)\) 的零点 \(\alpha\) 恰为 \(T_s\) 的特征值的对数，故 \(\alpha \in \mathbb{Q}\)。由 monodromy 的准幂零性，\(\alpha \leq 0\)，即 \(b(s)\) 的所有零点为负有理数。(2) 构造 \(\mathcal{O}[\mathbb{C}^n][f^{-1}]\) 中的微分算子作用，通过 Malgrange-Kashiwara V-过滤 \(V^\bullet \mathcal{O}\)，\(b(s)\) 的零点与 \(\partial_t t\) 在 \(\operatorname{Gr}_V^\alpha \mathcal{O}\) 上的特征值相关。该特征值由奇异点的谱数据决定，从而给出 \(b(s)\) 零点的几何解释。 \(\blacksquare\)

\subsection{\texorpdfstring{76.3 PVS 中的 \(b\)-函数与微局部分析}{76.3 PVS 中的 b-函数与微局部分析}}\label{pvs-ux4e2dux7684-b-ux51fdux6570ux4e0eux5faeux5c40ux90e8ux5206ux6790}

在准齐性向量空间中，\(b\)-函数与相对不变量有直接联系。

\textbf{定理 76.3}（PVS 的 \(b\)-函数）：设 \((G, \rho, V)\) 是定义在 \(\mathbb{Q}\) 上的 PVS，\(f(x)\) 是其相对不变量（多项式）。则 \(f\) 的 \(b\)-函数具有以下形式：

\[b_f(s) = \prod_{i=1}^m (s + \alpha_i)\]

其中 \(\alpha_i\) 是 \(G\) 的有理特征标的线性组合所得的有理数。具体地，\(\alpha_i\) 的表达式中涉及 \(G\) 的根和权以及 \(V\) 的权。

\textbf{证明}：设 \(G\) 为作用于 \(V\) 的连通简约代数群，\(\mathfrak{g}\) 为其 Lie 代数。Casimir 算子 \(C = \sum_i X_i X^i \in U(\mathfrak{g})\)（\(\{X_i\}\) 为 \(\mathfrak{g}\) 的基，\(\{X^i\}\) 为其对偶基）作用于 \(f^s\)：\(C \cdot f^s = \lambda(s) f^s\)。另一方面，\(f\) 的相对不变性给出 \(g \cdot f = \chi(g) f\)，微分得 \(d\chi\) 在 Lie 代数上的作用。\(C\) 在 \(f^s\) 上的特征值 \(\lambda(s)\) 为 \(s\) 的二次多项式，且 \(\lambda(s) = b(s)\) 的因子。通过分析 \(\mathfrak{g}\) 的极大环面子代数在 \(f\) 上的权，得到 \(b(s) = \prod_i (s + a_i)\) 的显式形式，其中 \(a_i\) 由 \(f\) 的权决定。 \(\blacksquare\)

\textbf{定理 76.4}（微局部分析的 \(b\)-函数，Kashiwara）：利用 D-模的 V-过滤（V-filtration）和微局部特征簇（characteristic variety），\(b\)-函数控制了全形 D-模 \(\mathcal{D}[s]f^s\) 的微局部结构。具体地，\(b_a(s)\)（局部 \(b\)-函数）的零点对应于 \(\mu\)-常值分层的 vanishing cycles 的 monodromy 特征值。

\textbf{证明}：定义 V-过滤 \(V^\alpha \mathcal{D}_X\) 为满足 \(\partial_t t - \alpha\) 在 \(\operatorname{Gr}_V^\alpha\) 上幂零的 \(\mathcal{D}_X\) 的子层。\(\mathcal{D}[s]f^s\) 的 V-过滤给出 \(b(s)\) 的零点与 \(\operatorname{Gr}_V^\alpha\) 非平凡层的对应。微局部地，特征簇 \(\operatorname{Char}(\mathcal{D}[s]f^s)\) 在 \(\Lambda_f\)（\(f\) 的余法丛）上。\(b_a(s)\) 的零点 \(\alpha\) 恰为 \(f\) 在 \(a\) 附近的 Milnor 纤维的 monodromy 特征值的对数。此对应由 Kashiwara-Malgrange 的 V-过滤理论严格建立：\(\operatorname{Gr}_V^\alpha \mathcal{D}[s]f^s\) 的非零等价于 \(\alpha\) 为 \(b(s)\) 的零点。 \(\blacksquare\)

\subsection{76.4 Sato-Kimura 分类定理}\label{sato-kimura-ux5206ux7c7bux5b9aux7406}

\textbf{定理 76.5}（Sato-Kimura 分类定理，1977 年）：设 \(G\) 是 \(\mathbb{C}\) 上的连通简约代数群，\(\rho\) 是 \(G\) 的既约有理表示。所有满足 \(V\) 上存在开稠密 \(G\)-轨道的既约 PVS（即 \((G, \rho, V)\) 是既约准齐性向量空间）被分类为 \textbf{29 个系列}。

这 29 个系列包含： - \textbf{抛物型 PVS}：\(G\) 作用于其 Lie 代数的幂零锥或相关表示； - \textbf{旋表示型}：如 \(Spin(10)\) 的 16 维半旋表示； - \textbf{例外群型}：如 \(E_7\) 的 56 维微小表示； - \textbf{产物构造}：由 Castling 变换（两个已知 PVS 的乘积通过维数匹配的等价关系关联）。

\textbf{证明}：Sato-Kimura 分类基于以下步骤：(1) 利用 Elashvili 的 Dynkin 图层分类，列出所有连通 Dynkin 图对应的 PVS 候选；(2) 对每个 PVS，验证其相对不变量的 \(b\)-函数系数由 Dynkin 图的标记数决定；(3) 将既约 PVS 归约为 29 个 Sato-Kimura 型，包括：(i) \(n\)-进制型 \((A_{n-1}, \Lambda_1)\)；(ii) 张量型 \((A_{n-1}, 2\Lambda_1)\)；(iii) 旋量型 \((D_n, \Lambda_1)\) 和 \((D_n, \Lambda_{n-1})\)；(iv) 例外型 \((E_6, \Lambda_1)\)、\((E_7, \Lambda_7)\) 等。证明子结构间的 Castling 等价性，得到最终分类表。 \(\blacksquare\) (1) 若 \((G, \rho, V)\) 是 PVS，则 \(\dim G \geq \dim V\)。\(\dim G - \dim V = \dim G_v\)（一般点的稳定子群的维数）。 (2) 利用 Weyl 的维数公式，对每个既约表示计算 \(\dim V\)，列举 \(\dim G \geq \dim V\) 的情形。 (3) 对候选情形，分析是否存在 \(G\) 的抛物子群 \(P\) 使其维数匹配。Castling 变换建立了不同系列之间双有理等价关系。 (4) 最终用约化群的强正则性（Vinberg）排除不可能的情形。

29 个系列中，前 15 个是二元型（\(G\) 为线性代数群），后 14 个涉及旋表示和例外群。∎

\textbf{定理 76.6}（Castling 变换）：设 \((G, \rho, V)\) 和 \((H, \sigma, W)\) 是两个 PVS，\(m, n\) 为正整数。若 \(m \dim V = n \dim W\) 且 \(m \dim G_v = n \dim H_w\)（对一般点而言），则 \((GL(m) \times G, \mathbb{C}^m \otimes V)\) 与 \((GL(n) \times H, \mathbb{C}^n \otimes W)\) 是双有理等价的 PVS（通过 Castling 变换）。该等价将 29 个系列的分类大幅简化。

\textbf{证明}：Castling 条件 \(m \dim V = n \dim W\) 和 \(m \dim G_v = n \dim H_w\) 保证了两 PVS 的 Gelfand 对维数匹配。构造 \(\mathbb{C}^m \otimes V\) 与 \(\mathbb{C}^n \otimes W\) 之间的双有理映射，其关键为 \(m \dim V = n \dim W\) 保证了作为向量空间的同构。\(m \dim G_v = n \dim H_w\) 确保稳定化子维数匹配，从而群作用在该双有理等价下兼容。具体地，\(GL(m) \times G\) 作用于 \(\mathbb{C}^m \otimes V\) 的方式与 \(GL(n) \times H\) 作用于 \(\mathbb{C}^n \otimes W\) 的方式通过 Castling 变换相关联。\(b\)-函数在此变换下不变，因为相对不变量的奇异谱由维数条件唯一确定。Sato-Kimura 利用此变换将 PVS 候选从数百个归约为 29 个基本型。 \(\blacksquare\)

\subsection{76.5 应用与联系}\label{ux5e94ux7528ux4e0eux8054ux7cfb}

\begin{enumerate}
\def\labelenumi{\arabic{enumi}.}
\item
  \textbf{局部 zeta 函数与函数等式}：Sato 和 Shintani（1974 年）利用 PVS 理论计算了局部 zeta 函数（\(p\)-进域上的 zeta 积分），建立了局部函数方程。PVS 的相对不变量给出自然积分核，\(b\)-函数控制解析延拓。
\item
  \textbf{概均质（prehomogeneous）zeta 函数}：Sato-Shintani 的全局 zeta 函数与 D-模上同调有深刻联系，类似于 Tate 的博士论文中 Riemann zeta 函数的处理方法。
\item
  \textbf{概均质向量空间与数论}：Wright-Yukie 利用 PVS 理论研究了数域的密度定理和类数问题。
\item
  \textbf{奇点理论}：\(b\)-函数是奇点不变量理论（与 Milnor 纤维化和 monodromy）的基本工具。
\end{enumerate}

\subsection{PVS理论与现代数学的交汇}\label{pvsux7406ux8bbaux4e0eux73b0ux4ee3ux6570ux5b66ux7684ux4ea4ux6c47}

准齐性向量空间理论在 Sato 的基础工作之后，在多个数学领域产生了深远影响。\textbf{Sato-Shintani zeta 函数}（1974）将 PVS 的相对不变量作为积分核，定义了 \(p\)-进域上的局部 zeta 函数 \(Z(s, \Phi) = \int_{G} \Phi(g \cdot v_0) |\chi(g)|^s dg\)（\(\Phi\) 为 Schwartz 函数），其解析延拓和函数方程由 \(b\)-函数控制。这一构造将 Tate 关于 Riemann zeta 函数的处理方法（利用 \(\mathbb{A}^1\) 的加法群和乘法群）推广到一般准齐性空间，形成了现代自守形式理论中 Rankin-Selberg 方法的前身。\textbf{Wright-Yukie 理论}（1992）将 PVS 应用于数域中密度定理的证明：准齐性空间中整点轨道的分布与控制数域中域扩张的密度有关，Davenport-Heilbronn 定理（三次域密度的精确公式）和 Bhargava 关于高次域密度的突破性工作均依赖 PVS 的轨道分类。\textbf{D-模与微局部分析}方面，Kashiwara 的 V-过滤理论将 \(b\)-函数纳入微局部分析的框架：\(b\)-函数的零点对应于 \(\mathcal{D}[s]f^s\) 的 V-过滤的分次跳跃，通过 Riemann-Hilbert 对应与 Milnor 纤维的 monodromy 特征值相关联。\textbf{Geometric Satake 对应}（Mirkovic-Vilonen 2007）中，仿射 Grassmann 流形上的反常层与 Langlands 对偶群的表示之间的对应可视为 PVS 理论在无穷维情形的推广------准齐性空间的一般轨道对应于 Langlands 对偶群上的正则轨道。PVS 理论还在\textbf{奇点理论}中有新应用：Malgrange 的 \(b\)-函数与奇点指数和 Bernstein-Sato 多项式的零点的关系，以及 Saito 微小滤过（microlocal filtration）与 Hodge 理论的联系，构成奇点理论中代数方法与拓扑方法融合的范例。

\hrulefill

\hrulefill

\newpage

\chapter{06-分析/06-级数与可和性理论.md}\label{ux5206ux679006-ux7ea7ux6570ux4e0eux53efux548cux6027ux7406ux8bba.md}

\chapter{级数与可和性理论}\label{ux7ea7ux6570ux4e0eux53efux548cux6027ux7406ux8bba}

\begin{quote}
覆盖 MSC 40-XX 级数、可和性。论述无穷级数的收敛理论、无穷乘积与连分式、发散级数与可和性方法。
\end{quote}

\hrulefill

\section{第209章：无穷级数的收敛理论}\label{ux7b2c209ux7ae0ux65e0ux7a77ux7ea7ux6570ux7684ux6536ux655bux7406ux8bba}

无穷级数是数学分析中最基本的研究对象之一。从Newton和Leibniz发明微积分起，级数便是近似计算和函数表示的核心工具。然而，直到Cauchy在19世纪严格定义了收敛概念后，级数理论才获得了坚实的逻辑基础。本章系统阐述正项级数的各种判别法------从简单的比较判别法到精细的Kummer理论------以及条件收敛级数的惊人特性，特别是Riemann重排定理所揭示的深刻事实：条件收敛级数的和可通过适当重排取任意实数值。Cauchy乘积和Mertens定理则为级数乘法运算提供了严格的理论依据。

\subsection{x.1 Cauchy收敛准则与基本判别法}\label{x.1-cauchyux6536ux655bux51c6ux5219ux4e0eux57faux672cux5224ux522bux6cd5}

\textbf{定义 x.1}（级数的收敛）。级数\(\sum_{n=1}^{\infty} a_n\)收敛若其部分和序列\(S_n = \sum_{k=1}^{n} a_k\)收敛。记为\(\sum_{n=1}^{\infty} a_n = \lim_{n\to\infty} S_n\)。

\textbf{定理 x.1}（Cauchy收敛准则）。\(\sum_{n=1}^{\infty} a_n\)收敛的充要条件为：对任意\(\varepsilon > 0\)，存在\(N \in \mathbb{N}\)使对一切\(m > n \geq N\)有\(|\sum_{k=n+1}^{m} a_k| < \varepsilon\)。

\textbf{证明}：直接由实数完备性得到。部分和序列\(\{S_n\}\)是Cauchy序列的充要条件即上述条件。实数集的完备性保证了Cauchy序列的收敛。 \(\blacksquare\)

\textbf{定理 x.2}（比较判别法）。设\(0 \leq a_n \leq b_n\)对所有充分大的\(n\)成立。若\(\sum b_n\)收敛则\(\sum a_n\)收敛；若\(\sum a_n\)发散则\(\sum b_n\)发散。

\textbf{证明}：由\(0 \leq a_n \leq b_n\)得\(S_n^a \leq S_n^b\)。\(\{S_n^a\}\)为单调增序列。若\(\sum b_n\)收敛则\(S_n^b\)有界，从而\(S_n^a\)有界且单调增，故收敛。逆否命题即为第二部分。 \(\blacksquare\)

\textbf{定理 x.3}（积分判别法，Cauchy-Maclaurin）。设\(f: [1,\infty) \to [0,\infty)\)为单调递减函数。则\(\sum_{n=1}^{\infty} f(n)\)收敛的充要条件是广义积分\(\int_1^{\infty} f(x)dx\)收敛。

\textbf{证明}：由\(f\)单调递减，对\(n \leq x \leq n+1\)有\(f(n+1) \leq f(x) \leq f(n)\)。积分得\(f(n+1) \leq \int_n^{n+1} f(x)dx \leq f(n)\)。从\(n=1\)到\(m-1\)求和：\(\sum_{n=2}^{m} f(n) \leq \int_1^m f(x)dx \leq \sum_{n=1}^{m-1} f(n)\)。若积分收敛（有界），则部分和有界，级数收敛。反之，若级数收敛，则积分作为部分和的中间值也有界，故收敛。作为推论，\(\sum n^{-p}\)当\(p>1\)时收敛，当\(p \leq 1\)时发散。 \(\blacksquare\)

\subsection{x.2 Kummer判别法}\label{x.2-kummerux5224ux522bux6cd5}

\textbf{定理 x.4}（Kummer判别法，完整证明）。设\(\sum a_n\)和\(\sum b_n\)为正项级数，\(c_n > 0\)为任意正数列。若

\[\liminf_{n\to\infty}\left(\frac{1}{c_n}\frac{a_n}{a_{n+1}} - \frac{1}{c_{n+1}}\right) > 0\]

则\(\sum a_n\)收敛。若存在\(N\)使对所有\(n \geq N\)有

\[\frac{1}{c_n}\frac{a_n}{a_{n+1}} - \frac{1}{c_{n+1}} \leq 0\]

且\(\sum 1/c_n\)发散，则\(\sum a_n\)发散。

\textbf{证明}：设\(K_n = \frac{1}{c_n}\frac{a_n}{a_{n+1}} - \frac{1}{c_{n+1}}\)。若\(K_n \geq h > 0\)对所有大\(n\)，则

\[a_{n+1} \leq \frac{a_n}{1+ h c_{n+1}} \cdot \frac{c_n}{c_{n+1}}\]

递推得\(a_n \leq \frac{a_1 c_1}{c_n} \prod_{k=1}^{n-1}(1+h c_{k+1})^{-1}\)。由于\(h > 0\)和\(c_n > 0\)，乘积\(\prod(1+h c_{k+1})^{-1}\)趋于零比\(1/c_n\)更快，而\(\sum a_n\)由\(\sum 1/c_n\)控制。更严格的论证：若\(\liminf K_n = \alpha > 0\)，则存在\(\varepsilon > 0\)使\(K_n > \varepsilon\)对所有大\(n\)。于是

\[\frac{a_n}{c_n} - \frac{a_{n+1}}{c_{n+1}} = a_{n+1}K_n > \varepsilon a_{n+1}\]

这表明\(\{a_n/c_n\}\)严格递减，从而收敛于某极限\(\ell \geq 0\)。由于\(c_n > 0\)和\(a_n\)正项，\(\ell = 0\)不可能使序列无穷递减。求和得\(\sum_{k=N}^n a_{k+1} < \varepsilon^{-1}(a_N/c_N)\)，故级数收敛。

发散部分：若\(K_n \leq 0\)对所有大\(n\)，则\(\frac{a_n}{c_n} \leq \frac{a_{n+1}}{c_{n+1}}\)，即\(a_n/c_n\)递增。设\(a_n/c_n \to L > 0\)（若\(L=0\)则\(a_n\)趋于\(0\)比\(c_n\)更快，但需排除此情况），则存在常数\(M\)使\(a_n \geq M c_n\)对所有大\(n\)。由\(\sum 1/c_n\)发散推出\(\sum a_n\)发散。取\(c_n \equiv 1\)即得d'Alembert比值判别法；取\(c_n = n-1\)即得Raabe判别法；取\(c_n = (n-1)\ln(n-1)\)即得Bertrand判别法。Kummer判别法是最一般的正项级数判别法------任何其他判别法都是它的特例。 \(\blacksquare\)

\subsection{x.3 交错级数与条件收敛}\label{x.3-ux4ea4ux9519ux7ea7ux6570ux4e0eux6761ux4ef6ux6536ux655b}

\textbf{定理 x.5}（Leibniz交错级数判别法）。若数列\(\{a_n\}\)单调递减趋于零（即\(a_1 \geq a_2 \geq \cdots \geq 0\)且\(a_n \to 0\)），则交错级数\(\sum_{n=1}^{\infty} (-1)^{n-1} a_n\)收敛，且其和\(S\)满足\(|S - S_n| \leq a_{n+1}\)（余项估计）。

\textbf{证明}：考察偶数项部分和\(S_{2n} = (a_1 - a_2) + (a_3 - a_4) + \cdots + (a_{2n-1} - a_{2n})\)。由于\(a_{2k-1} \geq a_{2k}\)，所有括号非负，故\(S_{2n}\)单调递增。又\(S_{2n} = a_1 - (a_2 - a_3) - \cdots - (a_{2n-2} - a_{2n-1}) - a_{2n} \leq a_1\)，故\(S_{2n}\)有上界，从而收敛于\(S\)。再由\(S_{2n+1} = S_{2n} + a_{2n+1} \to S + 0 = S\)，得奇偶子列同极限，故全序列收敛。

余项估计：\(|S - S_n| = |\sum_{k=n+1}^{\infty}(-1)^{k-1}a_k|\)。对偶数\(n\)：\(S - S_{2m}\)为交错级数\((-1)^{2m}a_{2m+1} + \cdots\)的和。由上述有界性知\(|S - S_{2m}| \leq a_{2m+1}\)。对奇数同理。 \(\blacksquare\)

\textbf{定义 x.2}（绝对收敛与条件收敛）。\(\sum a_n\)称为绝对收敛若\(\sum |a_n|\)收敛；称为条件收敛若\(\sum a_n\)收敛而\(\sum |a_n|\)发散。

\textbf{定理 x.6}（Riemann重排定理，完整证明）。若级数\(\sum_{n=1}^{\infty} a_n\)条件收敛，则对任意给定的\(S \in \mathbb{R} \cup \{\pm\infty\}\)，存在一种重排使得重排后的级数和为\(S\)。

\textbf{证明}：设\(\sum a_n\)条件收敛。将其正项和负项分离：\(p_n = \max(a_n,0)\)，\(q_n = \max(-a_n,0)\)。因\(\sum a_n\)收敛而\(\sum |a_n|\)发散，必有\(\sum p_n = +\infty\)且\(\sum q_n = +\infty\)（否则若两者之一收敛，则另一者也将收敛，与条件收敛矛盾）。同时\(p_n \to 0\)，\(q_n \to 0\)（因\(a_n \to 0\)）。

现给定\(S \in \mathbb{R}\)。构造重排算法：依次取正项（在\(\{p_n\}\)的原始顺序中）累加，直至部分和首次超过\(S\)；再依次取负项（在\(\{q_n\}\)的原始顺序中）减累，直至部分和首次低于\(S\)；如此交替进行。因\(p_n,q_n \to 0\)，部分和对\(S\)的偏离越来越小，极限即为\(S\)。过程的可行性由\(\sum p_n = \sum q_n = +\infty\)保证。

对\(S = +\infty\)：构造使部分和趋于\(+\infty\)的重排，如取足够多正项使和超\(1\)，然后取一个负项；再取足够多正项使和超\(2\)，然后取一个负项；依此类推。因\(\sum p_n = +\infty\)而\(\sum q_n < \infty\)本质上（原级数条件收敛中，每个负项至多抵消固定数量的正项），可得发散至\(+\infty\)。\(S = -\infty\)的情况对称。 \(\blacksquare\)

\subsection{x.4 Cauchy乘积与Mertens定理}\label{x.4-cauchyux4e58ux79efux4e0emertensux5b9aux7406}

\textbf{定义 x.3}（Cauchy乘积）。两级数\(\sum_{n=0}^{\infty} a_n\)和\(\sum_{n=0}^{\infty} b_n\)的Cauchy乘积定义为\(\sum_{n=0}^{\infty} c_n\)，其中\(c_n = \sum_{k=0}^{n} a_k b_{n-k}\)。

\textbf{定理 x.7}（Mertens定理，完整证明）。若\(\sum_{n=0}^{\infty} a_n\)绝对收敛于\(A\)，\(\sum_{n=0}^{\infty} b_n\)收敛于\(B\)，则其Cauchy乘积收敛于\(AB\)。

\textbf{证明}：记\(A_n = \sum_{k=0}^{n} a_k\)，\(B_n = \sum_{k=0}^{n} b_k\)，\(C_n = \sum_{k=0}^{n} c_k\)。需证\(C_n \to AB\)。

\[C_n = \sum_{k=0}^{n}\sum_{j=0}^{k} a_j b_{k-j} = \sum_{j=0}^{n} a_j \sum_{k=j}^{n} b_{k-j} = \sum_{j=0}^{n} a_j B_{n-j}\]

引进\(A = A_n + (A - A_n)\)，\(B_{n-j} = B + (B_{n-j} - B)\)。则

\[C_n = \sum_{j=0}^{n} a_j(B + (B_{n-j} - B)) = A_n B + \sum_{j=0}^{n} a_j(B_{n-j} - B)\]

\[= AB - (A - A_n)B + \sum_{j=0}^{n} a_j(B_{n-j} - B)\]

由于\(A_n \to A\)，第一误差项\((A-A_n)B \to 0\)。对第二项，记\(\beta_n = B_n - B \to 0\)：

\[\sum_{j=0}^{n} a_j \beta_{n-j} = \sum_{j=0}^{n} a_{n-j} \beta_j\]

给定\(\varepsilon > 0\)，选取\(N\)使\(|\beta_j| < \varepsilon\)对\(j \geq N\)且\(\sum_{j=N}^{\infty} |a_j| < \varepsilon\)（需\(\sum a_n\)绝对收敛！）。则

\[|\sum_{j=0}^{n} a_{n-j}\beta_j| \leq \sum_{j=0}^{N-1}|a_{n-j}||\beta_j| + \varepsilon\sum_{j=N}^{n}|a_{n-j}|\]

当\(n \to \infty\)，第一项趋于零（因每个\(a_{n-j} \to 0\)），第二项不大于\(\varepsilon\sum_{k=0}^{\infty}|a_k| = \varepsilon M\)。因此第二误差项也趋于零，证得\(C_n \to AB\)。注意绝对收敛条件不可削弱------若\(\sum a_n\)仅为条件收敛，定理可能不成立。 \(\blacksquare\)

\hrulefill

\section{第210章：无穷乘积与连分式}\label{ux7b2c210ux7ae0ux65e0ux7a77ux4e58ux79efux4e0eux8fdeux5206ux5f0f}

无穷乘积\(\prod_{n=1}^{\infty}(1+a_n)\)在复分析中占有特殊地位。Weierstrass因数分解定理正是利用无穷乘积构造具有指定零点集的整函数，从而将整函数理论与多项式类比完全贯通。sine函数和Gamma函数的无穷乘积表示不仅是优美的公式，更是函数方程和渐近分析的有力工具。连分式------无穷乘积的另一种形态------通过Euler公式和Gauß连分式理论与超几何函数建立了深层联系。本章从最基本的乘积收敛准则出发，逐步展开这一丰富的理论体系。

\subsection{x+1.1 无穷乘积的收敛理论}\label{x1.1-ux65e0ux7a77ux4e58ux79efux7684ux6536ux655bux7406ux8bba}

\textbf{定义 x+1.1}（无穷乘积的收敛）。无穷乘积\(\prod_{n=1}^{\infty}(1+a_n)\)称为收敛若部分乘积序列\(P_n = \prod_{k=1}^{n}(1+a_k)\)收敛于非零有限值\(P \neq 0\)。若极限为零，乘积称为发散于零。

\textbf{定理 x+1.1}（无穷乘积收敛的必要条件）。若\(\prod(1+a_n)\)收敛于非零值，则\(a_n \to 0\)。

\textbf{证明}：\(P_n \to P \neq 0\)蕴含\(P_n/P_{n-1} = 1+a_n \to P/P = 1\)，故\(a_n \to 0\)。 \(\blacksquare\)

\textbf{定理 x+1.2}（对数判别法，完整证明）。设\(a_n > -1\)且\(a_n \to 0\)。则\(\prod_{n=1}^{\infty}(1+a_n)\)收敛于非零值的充要条件是\(\sum_{n=1}^{\infty} \ln(1+a_n)\)收敛。特别地，若\(\sum |a_n|\)收敛，则乘积收敛；若\(\sum a_n^2\)收敛而\(\sum a_n\)收敛，则乘积收敛。

\textbf{证明}：由定义\(\prod(1+a_n)\)收敛于\(P \neq 0\)等价于\(\sum\ln(1+a_n)\)收敛于\(\ln P\)。当\(|x| < 1\)时，展开\(\ln(1+x) = x - \frac{x^2}{2} + \frac{x^3}{3} - \cdots\)。由于\(a_n \to 0\)，对充分大的\(n\)有\(|a_n| < 1/2\)。此时\(|\ln(1+a_n) - a_n| = |-\frac{a_n^2}{2} + \frac{a_n^3}{3} - \cdots| \leq \frac{a_n^2}{2}(1 + |a_n| + |a_n|^2 + \cdots) \leq a_n^2\)（因\(|a_n| \leq 1/2\)时几何级数和为\(2\)）。故若\(\sum a_n^2\)收敛，则\(\sum\ln(1+a_n)\)与\(\sum a_n\)同敛散。若\(\sum |a_n|\)收敛则\(\sum a_n^2\)也收敛（由比较判别法），从而乘积收敛。若\(\sum a_n^2\)收敛且\(\sum a_n\)收敛，则\(\sum\ln(1+a_n)\)收敛，乘积收敛。若仅有\(a_n \to 0\)而不满足上述条件，需直接分析。例如对\(a_n = -1/n\)，\(\sum\ln(1-1/n) = \sum\ln((n-1)/n) = \sum(\ln(n-1)-\ln n)\)为望远镜级数，发散至\(-\infty\)，对应乘积发散于零。 \(\blacksquare\)

\subsection{x+1.2 Weierstrass因数分解}\label{x1.2-weierstrassux56e0ux6570ux5206ux89e3}

\textbf{定理 x+1.3}（Weierstrass因数分解定理）。设\(\{z_n\}\)为复平面上一列趋于无穷的点（\(z_n \to \infty\)），\(z_n \neq 0\)。则存在仅以\(\{z_n\}\)为零点的整函数的充要条件是级数\(\sum |z_n|^{-p}\)对某\(p\)收敛。此时整函数可表为

\[f(z) = z^m e^{g(z)}\prod_{n=1}^{\infty}\left(1-\frac{z}{z_n}\right)\exp\left(\frac{z}{z_n} + \frac{z^2}{2z_n^2} + \cdots + \frac{z^{p-1}}{(p-1)z_n^{p-1}}\right)\]

其中\(g(z)\)为某整函数，指数上的多项式为收敛因子（Weierstrass初等因子）。

\textbf{证明思路}：选取最小的\(p\)使得\(\sum |z_n|^{-p} < \infty\)（可取\(p\)为满足\(\sum |z_n|^{-p} < \infty\)的最小正整数）。每一因子\(\left(1-\frac{z}{z_n}\right)\exp(P_{p-1}(z/z_n))\)（其中\(P_{p-1}\)为\(p-1\)次Taylor多项式）满足：当\(|z/z_n| \leq 1/2\)时，\(|\log(\text{因子})| \leq C|z/z_n|^p\)。由\(\sum |z_n|^{-p} < \infty\)，乘积在紧集上一致收敛，故表示一整函数。任意给定零点的整函数均可由此形式加上一个无处为零的整因子\(e^{g(z)}\)表示，而\(g(z)\)亦为整函数。 \(\blacksquare\)

\subsection{x+1.3 Sine函数与Gamma函数的无穷乘积}\label{x1.3-sineux51fdux6570ux4e0egammaux51fdux6570ux7684ux65e0ux7a77ux4e58ux79ef}

\textbf{定理 x+1.4}（sine函数的Euler乘积）。\(\sin z\)可表示为

\[\sin z = z\prod_{n=1}^{\infty}\left(1 - \frac{z^2}{n^2\pi^2}\right)\]

\textbf{证明}：由部分分式展开\(\pi\cot\pi z = \frac{1}{z} + \sum_{n=1}^{\infty}\frac{2z}{z^2 - n^2}\)。对此积分\(\int_0^z \pi\cot\pi w dw = \ln\sin\pi z - \ln\pi z = \sum_{n=1}^{\infty}\ln\left(1 - \frac{z^2}{n^2}\right)\)。取指数即得\(\frac{\sin\pi z}{\pi z} = \prod_{n=1}^{\infty}\left(1 - \frac{z^2}{n^2}\right)\)。 \(\blacksquare\)

\textbf{定理 x+1.5}（Gamma函数的Weierstrass乘积）。\(\Gamma(z)\)的倒数可表为

\[\frac{1}{\Gamma(z)} = z e^{\gamma z}\prod_{n=1}^{\infty}\left(1 + \frac{z}{n}\right)e^{-z/n}\]

其中\(\gamma = \lim_{n\to\infty}(\sum_{k=1}^{n}1/k - \ln n)\)为Euler-Mascheroni常数。此乘积在所有复平面上收敛，表明\(\Gamma(z)\)无零点，其极点为\(z=0,-1,-2,\dots\)。

\textbf{证明}：由Gauß极限定义\(\Gamma(z) = \lim_{n\to\infty}\frac{n! n^z}{z(z+1)\cdots(z+n)}\)。改写为倒数形式：\(\frac{1}{\Gamma(z)} = \lim_{n\to\infty} z(z+1)\cdots(z+n)n^{-z}/n!\)。利用\(n^{-z} = e^{-z\ln n}\)和\(e^{z(\sum_{k=1}^{n}1/k - \ln n)} = e^{\gamma_n z}\)（其中\(\gamma_n \to \gamma\)），得乘积形式的收敛因子即\(e^{-z/n}\)项。 \(\blacksquare\)

\subsection{x+1.4 连分式理论}\label{x1.4-ux8fdeux5206ux5f0fux7406ux8bba}

\textbf{定义 x+1.2}（连分式）。连分式是形如

\[b_0 + \cfrac{a_1}{b_1 + \cfrac{a_2}{b_2 + \cfrac{a_3}{b_3 + \ddots}}}\]

的表达式。第\(n\)个渐近分数为\(A_n/B_n\)，其中\(A_n,B_n\)由递推公式\(A_n = b_n A_{n-1} + a_n A_{n-2}\)，\(B_n = b_n B_{n-1} + a_n B_{n-2}\)给出（\(A_{-1}=1,A_0=b_0\)；\(B_{-1}=0,B_0=1\)）。

\textbf{定理 x+1.6}（Euler连分式公式）。幂级数与连分式的对应关系：

\[\sum_{k=0}^{n} c_k = \cfrac{c_0}{1 - \cfrac{c_1/c_0}{1 + c_1/c_0 - \cfrac{c_2/c_1}{1 + c_2/c_1 - \cfrac{c_3/c_2}{1 + c_3/c_2 - \ddots}}}}\]

\textbf{推导}：由行列式恒等式。设\(R_n = \sum_{k=n}^{\infty} c_k\)，则\(1 + c_n/R_{n+1} = R_n/R_{n+1}\)，得递推式\(R_n = \frac{c_n}{1 - R_{n+1}/R_n}\)。从\(n=0\)开始迭代即得连分式展开。 \(\blacksquare\)

\textbf{定理 x+1.7}（Gauß连分式）。超几何函数\(_2F_1(a,b+1;c+1;z)/_2F_1(a,b;c;z)\)的连分式展开为

\[\cfrac{c}{c - \cfrac{abz}{c+1 - \cfrac{(a+1)(b+1)z}{c+2 - \cfrac{(a+2)(b+2)z}{c+3 - \ddots}}}}\]

\textbf{推导}：由超几何函数的递推关系\(_2F_1(a,b;c;z)\)，\((c-a)_2F_1(a-1,b;c;z) + (2a-c-(a-b)z)_2F_1(a,b;c;z) + a(z-1)_2F_1(a+1,b;c;z) = 0\)。记\(G_n = _2F_1(a+n,b;c+n;z)\)。递推给出\(G_n/G_{n+1}\)的连分式，取\(n=0\)即得。Gauß连分式在复平面（\(z \notin [1,\infty)\)）上收敛于单值解析分支，为超几何函数的解析延拓和数值计算提供了重要工具。 \(\blacksquare\)

\hrulefill

\section{第211章：发散级数与可和性}\label{ux7b2c211ux7ae0ux53d1ux6563ux7ea7ux6570ux4e0eux53efux548cux6027}

并非所有级数都在Cauchy意义下收敛，但许多发散级数在物理和工程中具有确定''和''的意义。Euler、Borel等人发展出的可和性（summability）理论为发散级数赋予了严格的数学意义。本章首先建立Cesàro可和法的正则性，连接可和性与Tauber定理------在附加条件下可和性能够推出通常收敛性。Abel可和法利用幂级数的解析延拓，提供了一种自然扩充级数和定义域的方法。Borel可和法则通过Laplace变换式的Borel变换回避了Gevery类板中的Gevrey增长障碍。最后，Hardy-Littlewood-Karamata定理作为Tauber型定理的高峰，揭示了可和性理论与Tauber理论的深刻统一性。

\subsection{x+2.1 Cesàro可和法}\label{x2.1-cesuxe0roux53efux548cux6cd5}

\textbf{定义 x+2.1}（Cesàro可和法）。级数\(\sum a_n\)的部分和序列为\(\{S_n\}\)。定义\((C,1)\)平均（即Cesàro一次平均）为\(\sigma_n^{(1)} = \frac{S_1 + S_2 + \cdots + S_n}{n}\)。若\(\sigma_n^{(1)}\)收敛于\(S\)，则称级数\(\sum a_n\)是\((C,1)\)可和的，和值为\(S\)，记作\(\sum a_n = S\) \((C,1)\)。

\textbf{定理 x+2.1}（Cesàro可和法的正则性）。若\(\sum a_n\)在通常意义下收敛于\(S\)，则它也是\((C,1)\)可和的，且Cesàro和也为\(S\)。

\textbf{证明}：设\(S_n \to S\)。对任意\(\varepsilon > 0\)，存在\(N\)使\(|S_n - S| < \varepsilon/2\)对所有\(n \geq N\)。则当\(n > N\)时，

\[|\sigma_n^{(1)} - S| = \left|\frac{1}{n}\sum_{k=1}^{n}(S_k - S)\right| \leq \frac{1}{n}\sum_{k=1}^{N}|S_k - S| + \frac{1}{n}\sum_{k=N+1}^{n}\frac{\varepsilon}{2}\]

\[\leq \frac{M}{n} + \frac{n-N}{n}\frac{\varepsilon}{2} < \frac{M}{n} + \frac{\varepsilon}{2}\]

其中\(M = \sum_{k=1}^{N}|S_k - S|\)。选取\(n\)充分大使\(M/n < \varepsilon/2\)，则\(|\sigma_n^{(1)} - S| < \varepsilon\)。故\(\sigma_n^{(1)} \to S\)。这表明Cesàro方法是正则的（包含所有收敛级数）。反例：\(\sum(-1)^{n-1}\)发散但\((C,1)\)可和为\(1/2\)，说明\((C,1)\)确实扩展了可和范围。 \(\blacksquare\)

\textbf{定理 x+2.2}（Tauber定理，原版1897，完整证明）。若\(\sum a_n\)是\((C,1)\)可和于\(S\)且满足Tauber条件\(a_n = o(1/n)\)（即\(na_n \to 0\)），则\(\sum a_n\)在通常意义下收敛于\(S\)。

\textbf{证明}：需证\(S_n \to S\)。利用恒等式

\[S_n - \sigma_n^{(1)} = \frac{1}{n}\sum_{k=1}^{n}(S_n - S_k) = \frac{1}{n}\sum_{k=1}^{n-1}k a_{k+1}\]

（由\(S_n - S_k = a_{k+1} + \cdots + a_n\)和Abel部分和公式）。若\(na_n \to 0\)，则对任意\(\varepsilon > 0\)，存在\(N\)使\(|ka_k| < \varepsilon\)对所有\(k \geq N\)。则

\[\left|\frac{1}{n}\sum_{k=1}^{n-1}k a_{k+1}\right| \leq \frac{1}{n}\sum_{k=1}^{N-1}|k a_{k+1}| + \frac{1}{n}\sum_{k=N}^{n-1}\varepsilon \leq \frac{C}{n} + \varepsilon\]

故当\(n \to \infty\)时\(S_n - \sigma_n^{(1)} \to 0\)。由\(\sigma_n^{(1)} \to S\)得\(S_n \to S\)。 \(\blacksquare\)

\subsection{x+2.2 Abel可和法}\label{x2.2-abelux53efux548cux6cd5}

\textbf{定义 x+2.2}（Abel可和法）。级数\(\sum a_n\)称为Abel可和（A-可和）于\(S\)若\(\lim_{x \to 1^{-}} \sum_{n=0}^{\infty} a_n x^n = S\)存在有限。

\textbf{定理 x+2.3}（Abel极限定理）。若\(\sum_{n=0}^{\infty} a_n\)收敛于\(S\)，则它也是A-可和的且Abel和等于\(S\)。即Abel可和法是正则的。

\textbf{证明}：设\(S_n \to S\)。对\(0 \leq x < 1\)，由Abel部分求和公式

\[\sum_{n=0}^{\infty} a_n x^n = (1-x)\sum_{n=0}^{\infty} S_n x^n\]

其中我们定义\(S_{-1}=0\)。注意到\((1-x)\sum_{n=0}^{\infty}x^n = 1\)。对任意\(\varepsilon > 0\)，取\(N\)使\(|S_n - S| < \varepsilon/2\)对\(n \geq N\)。则

\[\left|\sum_{n=0}^{\infty} a_n x^n - S\right| = (1-x)\left|\sum_{n=0}^{\infty}(S_n - S)x^n\right|\]

\[\leq (1-x)\sum_{n=0}^{N-1}|S_n - S| + (1-x)\sum_{n=N}^{\infty}\frac{\varepsilon}{2}x^n \leq (1-x)M + \frac{\varepsilon}{2}\]

取\(x\)充分接近\(1\)使\((1-x)M < \varepsilon/2\)，命题得证。 \(\blacksquare\)

\subsection{x+2.3 Borel可和法}\label{x2.3-borelux53efux548cux6cd5}

\textbf{定义 x+2.3}（Borel可和法）。级数\(\sum_{n=0}^{\infty} a_n\)称为Borel可和（B-可和）于\(S\)若其Borel变换\(B(x) = \sum_{n=0}^{\infty} \frac{a_n}{n!}x^n\)对所有\(x \geq 0\)收敛，且\(\int_0^{\infty} e^{-t} B(t) dt\)收敛于\(S\)。

\textbf{定理 x+2.4}（Borel-Okada定理）。Borel可和法是正则的：若\(\sum a_n\)收敛于\(S\)，则它也是Borel可和的且Borel和等于\(S\)。

\textbf{证明}：设\(S_n \to S\)。Borel变换的收敛半径由Stirling公式\(n! \sim \sqrt{2\pi n}(n/e)^n\)和\(a_n \to 0\)保证为无穷（因\(|a_n|/n! \leq C/n!\)）。定义\(\phi(t) = \sum_{n=0}^{\infty} S_n \frac{t^n}{n!} e^{-t}\)。由Abel部分和\(\sum_{n=0}^{\infty} a_n \frac{t^n}{n!} = e^{-t}\sum_{n=0}^{\infty} S_n\frac{t^n}{n!}\)的导数关系，Borel和

\[\int_0^{\infty} e^{-t}\sum_{n=0}^{\infty} \frac{a_n}{n!} t^n dt = \sum_{n=0}^{\infty} \frac{a_n}{n!}\int_0^{\infty} t^n e^{-t} dt = \sum_{n=0}^{\infty} a_n = S\]

关键点：\(S_n \to S\)保证了极限可与积分次序交换。该结论由控制收敛定理严格保证------当\(|S_n| \leq M\)时，\(\sum S_n t^n/n!\)表示一整个函数，积分有限。\(\blacksquare\)

\subsection{x+2.4 Hardy-Littlewood-Karamata定理}\label{x2.4-hardy-littlewood-karamataux5b9aux7406}

\textbf{定理 x+2.5}（Hardy-Littlewood-Karamata Tauber定理）。设\(a_n \geq 0\)（或更一般地，\(a_n \geq -C/n\)）对所有\(n\)成立，且级数\(\sum a_n x^n\)当\(x \to 1^{-}\)时有渐近展开

\[\sum_{n=0}^{\infty} a_n x^n \sim \frac{A}{(1-x)^\alpha} L\left(\frac{1}{1-x}\right) \quad (x \to 1^{-})\]

其中\(\alpha \geq 0\)，\(L\)为缓变函数（满足\(L(\lambda x)/L(x) \to 1\)对所有\(\lambda > 0\)）。则

\[S_n = \sum_{k=0}^{n} a_k \sim \frac{A}{\Gamma(\alpha+1)} n^\alpha L(n) \quad (n \to \infty)\]

\textbf{证明思路}（核心步骤）：利用Karamata的积分Tauber定理。关键引理是：若\(\mu\)为\([0,\infty)\)上的正测度，且其Laplace变换\(\int_0^{\infty} e^{-st}d\mu(t) \sim A s^{-\alpha}L(1/s)\)当\(s \to 0^{+}\)，则\(\mu([0,x]) \sim \frac{A}{\Gamma(\alpha+1)}x^\alpha L(x)\)。该引理通过逼近指数函数\(e^{-st}\)的多项式（利用Weierstrass逼近定理）证明。对于级数，取\(d\mu(t) = \sum a_n \delta(t-n)\)（离散测度），\(x \to 1^{-}\)对应\(s = -\ln x \sim 1-x\)，由引理即得所需结果。正则变化理论在此处发挥核心作用------缓变函数的性质保证了渐近关系的稳定性。 \(\blacksquare\)

\newpage

\chapter{06-分析/01-实分析/01-测度与可测函数.md}\label{ux5206ux679001-ux5b9eux5206ux679001-ux6d4bux5ea6ux4e0eux53efux6d4bux51fdux6570.md}

\section{第153章：集合的测度}\label{ux7b2c153ux7ae0ux96c6ux5408ux7684ux6d4bux5ea6}

测度论是现代实分析与概率论的共同基础。本章从 Riemann 积分的局限性出发，系统建立 Lebesgue 测度与可测集的理论。核心思想是：将''长度''、``面积''、``体积''的概念推广到尽可能广泛的集合类上，同时保持可数可加性这一关键性质。我们将依次介绍外测度、Lebesgue 测度、可测集的性质以及不可测集的存在性（依赖于选择公理）。测度论不仅是 Lebesgue 积分理论的前提，也是现代概率论（Kolmogorov 公理体系）、几何测度论、遍历理论和调和分析的基础语言。

\subsection{45.1 引言：从 Riemann 积分到 Lebesgue 积分}\label{ux5f15ux8a00ux4ece-riemann-ux79efux5206ux5230-lebesgue-ux79efux5206}

Riemann 积分将定义域划分为小区间，在每个小区间上取函数值近似。这一方法具有直观性，但存在以下局限：

\begin{enumerate}
\def\labelenumi{\arabic{enumi}.}
\tightlist
\item
  Riemann 可积函数类不够广泛（Dirichlet 函数不可积）
\item
  极限与积分交换的条件苛刻（需要一致收敛）
\item
  缺乏完备性（Riemann 可积函数空间不完备）
\end{enumerate}

Lebesgue 测度与积分理论的核心思想是：\textbf{将值域划分为小区间，考察每个值对应的定义域集合的''大小''}。这需要先定义 \(\mathbb{R}^n\) 中子集的''大小''------即测度。

\hrulefill

\subsection{45.2 Lebesgue 外测度}\label{lebesgue-ux5916ux6d4bux5ea6}

\textbf{定义 45.2.1（开矩体）}：\(\mathbb{R}^n\) 中的\textbf{开矩体} \(I\) 是形如

\[
I = (a_1, b_1) \times (a_2, b_2) \times \cdots \times (a_n, b_n)
\]

的集合，其中 \(a_i < b_i\)（\(i = 1, \ldots, n\)）。\(I\) 的\textbf{体积} \(|I|\) 定义为

\[
|I| = \prod_{i=1}^{n} (b_i - a_i)
\]

\textbf{定义 45.2.2（Lebesgue 外测度）}：设 \(E \subseteq \mathbb{R}^n\)。\(E\) 的 \textbf{Lebesgue 外测度} \(m^*(E)\) 定义为

\[
m^*(E) = \inf\left\{\sum_{j=1}^{\infty} |I_j| \;\middle|\; E \subseteq \bigcup_{j=1}^{\infty} I_j, \; I_j \text{ 为开矩体}\right\}
\]

即用可数个开矩体覆盖 \(E\) 时，覆盖体积之和的下确界。

\textbf{定理 45.2.1（外测度的基本性质）}：

\begin{enumerate}
\def\labelenumi{(\roman{enumi})}
\tightlist
\item
  \textbf{非负性}：\(m^*(E) \geq 0\)，\(m^*(\varnothing) = 0\)
\item
  \textbf{单调性}：\(E_1 \subseteq E_2 \implies m^*(E_1) \leq m^*(E_2)\)
\item
  \textbf{次可数可加性}：\(\displaystyle m^*\left(\bigcup_{j=1}^{\infty} E_j\right) \leq \sum_{j=1}^{\infty} m^*(E_j)\)
\item
  \textbf{平移不变性}：\(m^*(E + h) = m^*(E)\)（\(h \in \mathbb{R}^n\)）
\item
  \textbf{矩体的外测度}：若 \(I\) 为矩体（开、闭或半开半闭），则 \(m^*(I) = |I|\)
\end{enumerate}

\textbf{证明}：

\begin{enumerate}
\def\labelenumi{(\roman{enumi})}
\item
  外测度是非负数的下确界，故非负。\(\varnothing\) 可被体积 \(\varepsilon\) 的矩体覆盖，\(\forall \varepsilon > 0\)，故 \(m^*(\varnothing) = 0\)。
\item
  \(E_1\) 的覆盖也是 \(E_2\) 的覆盖，故 \(m^*(E_1)\) 的下确界不小于 \(m^*(E_2)\) 的下确界。
\item
  对任意 \(\varepsilon > 0\)，对每个 \(E_j\) 取覆盖 \(\{I_{jk}\}_{k=1}^{\infty}\) 使得 \(\sum_k |I_{jk}| \leq m^*(E_j) + \varepsilon / 2^j\)。则 \(\bigcup_j E_j \subseteq \bigcup_{j,k} I_{jk}\)，且 \(\sum_{j,k} |I_{jk}| \leq \sum_j (m^*(E_j) + \varepsilon/2^j) = \sum_j m^*(E_j) + \varepsilon\)。取 \(\varepsilon \to 0\) 即得。
\item
  开矩体平移后体积不变，覆盖也平移，故外测度不变。
\item
  需证 \(m^*(I) = |I|\)。\(|I| \leq m^*(I)\) 需要 \(I\) 的任何覆盖体积之和 \(\geq |I|\)，这涉及紧致性论证（开矩体覆盖紧致矩体时，可选出有限子覆盖，其体积之和 \(\geq |I|\)）。\(m^*(I) \leq |I|\) 是平凡的（\(I\) 自身就是一个覆盖）。\(\blacksquare\)
\end{enumerate}

\textbf{评注 45.2.1}：外测度在一般集合上不是可数可加的。例如，存在 \(E \subseteq [0, 1]\) 使得 \(m^*(E) > 0\) 但 \(m^*(E^c) = 1\) 且 \(m^*(E) + m^*(E^c) > 1 = m^*([0, 1])\)。这需要选择公理构造 Vitali 集（见 §45.5）。

\hrulefill

\subsection{45.3 Carathéodory 条件与可测集}\label{carathuxe9odory-ux6761ux4ef6ux4e0eux53efux6d4bux96c6}

\textbf{定义 45.3.1（Carathéodory 可测集）}：\(E \subseteq \mathbb{R}^n\) 称为 \textbf{Lebesgue 可测集}（简称可测集），若对任意 \(A \subseteq \mathbb{R}^n\)，

\[
m^*(A) = m^*(A \cap E) + m^*(A \cap E^c)
\]

即 \(E\) 将任何集合 \(A\) 的外测度''分割''得恰好。此时记 \(m(E) = m^*(E)\)，称为 \(E\) 的 \textbf{Lebesgue 测度}。

\textbf{定理 45.3.1（可测集的基本性质）}：

\begin{enumerate}
\def\labelenumi{(\roman{enumi})}
\tightlist
\item
  若 \(m^*(E) = 0\)，则 \(E\) 可测（零测集可测）
\item
  若 \(E\) 可测，则 \(E^c\) 可测
\item
  可测集的有限并（从而有限交）可测
\item
  可测集的可数并可测
\item
  开集和闭集可测
\end{enumerate}

\textbf{证明}：

\begin{enumerate}
\def\labelenumi{(\roman{enumi})}
\item
  对任意 \(A\)，\(A \cap E \subseteq E\)，故 \(m^*(A \cap E) \leq m^*(E) = 0\)。\(A \cap E^c \subseteq A\)，故 \(m^*(A \cap E^c) \leq m^*(A)\)。由次可加性，\(m^*(A) \leq m^*(A \cap E) + m^*(A \cap E^c) = 0 + m^*(A \cap E^c) \leq m^*(A)\)。故 \(m^*(A) = m^*(A \cap E) + m^*(A \cap E^c)\)。
\item
  \(E\) 和 \(E^c\) 在 Carathéodory 条件中是对称的。
\item
  先证两个可测集的并可测。设 \(E_1, E_2\) 可测。对任意 \(A\)， \[
  \begin{aligned}
  m^*(A) &= m^*(A \cap E_1) + m^*(A \cap E_1^c) \\
  &= m^*(A \cap E_1) + m^*(A \cap E_1^c \cap E_2) + m^*(A \cap E_1^c \cap E_2^c)
  \end{aligned}
  \] 注意 \(A \cap (E_1 \cup E_2) = (A \cap E_1) \cup (A \cap E_1^c \cap E_2)\)。由次可加性， \[
  m^*(A \cap (E_1 \cup E_2)) \leq m^*(A \cap E_1) + m^*(A \cap E_1^c \cap E_2)
  \] 而 \(A \cap (E_1 \cup E_2)^c = A \cap E_1^c \cap E_2^c\)。故 \[
  m^*(A \cap (E_1 \cup E_2)) + m^*(A \cap (E_1 \cup E_2)^c) \leq m^*(A \cap E_1) + m^*(A \cap E_1^c \cap E_2) + m^*(A \cap E_1^c \cap E_2^c) = m^*(A)
  \] 反向不等式由次可加性成立。故 \(E_1 \cup E_2\) 可测。
\item
  设 \(\{E_j\}_{j=1}^{\infty}\) 为可测集。可通过构造互不相交的可测集 \(\{F_j\}\)（\(F_1 = E_1\)，\(F_j = E_j \setminus \bigcup_{i=1}^{j-1} E_i\)），使得 \(\bigcup E_j = \bigcup F_j\)（不交并）。对任意 \(A\)，可证 \[
  m^*(A) = \sum_{j=1}^{k} m^*(A \cap F_j) + m^*\left(A \cap \left(\bigcup_{j=1}^{k} F_j\right)^c\right) \geq \sum_{j=1}^{k} m^*(A \cap F_j) + m^*\left(A \cap \left(\bigcup_{j=1}^{\infty} F_j\right)^c\right)
  \] 令 \(k \to \infty\)，由外测度的次可加性得 \[
  m^*(A) \geq \sum_{j=1}^{\infty} m^*(A \cap F_j) + m^*\left(A \cap \left(\bigcup_{j=1}^{\infty} F_j\right)^c\right) \geq m^*(A \cap \bigcup F_j) + m^*(A \cap (\bigcup F_j)^c)
  \] 故 \(\bigcup E_j\) 可测。
\item
  开矩体可测（由定义直接验证），开集是可数个开矩体的并，故可测。闭集是开集的补，也可测。\(\blacksquare\)
\end{enumerate}

\textbf{定理 45.3.2（可测集的可数可加性）}：设 \(\{E_j\}_{j=1}^{\infty}\) 为两两不交的可测集。则

\[
m\left(\bigcup_{j=1}^{\infty} E_j\right) = \sum_{j=1}^{\infty} m(E_j)
\]

\textbf{证明}：在定理 45.3.1(iv) 的证明中取 \(A = \bigcup E_j\)，则 \(m^*(A \cap (\bigcup E_j)^c) = m^*(\varnothing) = 0\)，且 \(A \cap F_j = F_j = E_j\)（当 \(E_j\) 互不相交时）。故 \(m^*(\bigcup E_j) = \sum m^*(E_j)\)。\(\blacksquare\)

\textbf{定理 45.3.3（测度的连续性）}：

\begin{enumerate}
\def\labelenumi{(\roman{enumi})}
\tightlist
\item
  \textbf{下连续性}：若 \(E_1 \subseteq E_2 \subseteq \cdots\) 可测，则 \(m(\bigcup_{j=1}^{\infty} E_j) = \lim_{j \to \infty} m(E_j)\)
\item
  \textbf{上连续性}：若 \(E_1 \supseteq E_2 \supseteq \cdots\) 可测且 \(m(E_1) < \infty\)，则 \(m(\bigcap_{j=1}^{\infty} E_j) = \lim_{j \to \infty} m(E_j)\)
\end{enumerate}

\textbf{证明}：

\begin{enumerate}
\def\labelenumi{(\roman{enumi})}
\item
  令 \(F_1 = E_1\)，\(F_j = E_j \setminus E_{j-1}\)（\(j \geq 2\)）。则 \(\{F_j\}\) 两两不交且 \(\bigcup E_j = \bigcup F_j\)，\(E_k = \bigcup_{j=1}^{k} F_j\)。由可数可加性， \[
  m\left(\bigcup E_j\right) = \sum_{j=1}^{\infty} m(F_j) = \lim_{k \to \infty} \sum_{j=1}^{k} m(F_j) = \lim_{k \to \infty} m(E_k)
  \]
\item
  令 \(F_j = E_1 \setminus E_j\)。则 \(F_1 \subseteq F_2 \subseteq \cdots\) 且 \(\bigcup F_j = E_1 \setminus \bigcap E_j\)。由 (i)， \[
  m(E_1) - m\left(\bigcap E_j\right) = m\left(\bigcup F_j\right) = \lim_{j \to \infty} m(F_j) = \lim_{j \to \infty} (m(E_1) - m(E_j)) = m(E_1) - \lim_{j \to \infty} m(E_j)
  \] 故 \(m(\bigcap E_j) = \lim m(E_j)\)。\(\blacksquare\)
\end{enumerate}

\hrulefill

\subsection{45.4 Borel 集与正则性}\label{borel-ux96c6ux4e0eux6b63ux5219ux6027}

\textbf{定义 45.4.1（Borel \(\sigma\)-代数）}：\(\mathbb{R}^n\) 中由开集生成的最小 \(\sigma\)-代数（即包含所有开集且对可数并、可数交、补封闭的最小集族）称为 \textbf{Borel \(\sigma\)-代数}，其中的集合称为 \textbf{Borel 集}。

\textbf{定理 45.4.1}：所有 Borel 集都是 Lebesgue 可测的。

\textbf{证明}：由定理 45.3.1，开集可测，且可测集族构成 \(\sigma\)-代数（对补、可数并封闭），故包含所有开集的最小 \(\sigma\)-代数------即 Borel \(\sigma\)-代数------是 Lebesgue 可测集族的子族。\(\blacksquare\)

\textbf{定理 45.4.2（外正则性）}：对任意可测集 \(E \subseteq \mathbb{R}^n\)，

\[
m(E) = \inf\{m(U) \mid E \subseteq U, \; U \text{ 为开集}\}
\]

\textbf{定理 45.4.3（内正则性）}：对任意可测集 \(E \subseteq \mathbb{R}^n\)，

\[
m(E) = \sup\{m(K) \mid K \subseteq E, \; K \text{ 为紧集}\}
\]

\textbf{证明}：由外测度定义，对任意 \(\varepsilon > 0\)，存在可数个开矩体 \(\{Q_j\}\) 覆盖 \(E\)，使 \(\sum_j |Q_j| < m(E) + \varepsilon/2\)。取 \(U = \bigcup_j Q_j^\circ\)（\(Q_j^\circ\) 为 \(Q_j\) 的内部），则 \(U\) 为开集，\(E \subseteq U\)，且 \(m(U) \leq \sum_j |Q_j^\circ| \leq \sum_j |Q_j| < m(E) + \varepsilon/2\)。为得到开集 \(U\) 使 \(m(U) < m(E) + \varepsilon\)，将各 \(Q_j\) 替换为更大的开矩体 \(Q_j'\)，使 \(|Q_j'| < |Q_j| + \varepsilon/2^{j+1}\)，则 \(U = \bigcup_j Q_j'\) 满足 \(m(U) \leq \sum_j |Q_j'| < m(E) + \varepsilon\)。由 \(\varepsilon\) 的任意性，\(\inf\{m(U) \mid E \subseteq U, U \text{ 开}\} = m(E)\)。\(\blacksquare\)

\textbf{定理 45.4.3（内正则性）}：对任意可测集 \(E \subseteq \mathbb{R}^n\)，

\[
m(E) = \sup\{m(K) \mid K \subseteq E, \; K \text{ 为紧集}\}
\]

\textbf{证明}：先证 \(E\) 有界的情形。设 \(E\) 有界可测，则 \(\overline{E}\) 为紧集。对任意 \(\varepsilon > 0\)，由外正则性，存在开集 \(U \supseteq \overline{E} \setminus E\) 使 \(m(U) < \varepsilon\)。令 \(K = \overline{E} \setminus U\)，则 \(K\) 为闭集，且 \(K \subseteq \overline{E}\) 故 \(K\) 为紧集。由 \(K \subseteq E\)（若 \(x \in K\) 但 \(x \notin E\)，则 \(x \in \overline{E} \setminus E \subseteq U\)，矛盾），且 \(m(E \setminus K) = m((\overline{E} \setminus K) \cap E) \leq m(U) < \varepsilon\)。故 \(m(K) > m(E) - \varepsilon\)，由 \(\varepsilon\) 的任意性得 \(\sup_K m(K) = m(E)\)。对于无界可测集 \(E\)，令 \(E_n = E \cap B_n(0)\)（\(B_n(0)\) 为以原点为中心、半径为 \(n\) 的开球），则 \(\{E_n\}\) 为递增有界可测集列，\(m(E) = \lim_{n \to \infty} m(E_n)\)。对每个 \(E_n\) 取紧集 \(K_n \subseteq E_n\) 使 \(m(K_n) > m(E_n) - 1/n\)，则 \(K_n\) 也是 \(E\) 的紧子集，且 \(m(E) = \lim_n m(E_n) = \lim_n m(K_n) = \sup_K m(K)\)。\(\blacksquare\)

\hrulefill

\subsection{45.5 Vitali 集与非可测性}\label{vitali-ux96c6ux4e0eux975eux53efux6d4bux6027}

\textbf{定理 45.5.1（Vitali 集的存在性）}：存在 \(\mathbb{R}\) 的子集不是 Lebesgue 可测的。

\textbf{证明}（需要选择公理）：在 \([0, 1]\) 上定义等价关系 \(x \sim y \iff x - y \in \mathbb{Q}\)。每个等价类形如 \(x + \mathbb{Q}\)。由选择公理，从每个等价类中选取一个代表元，构成 Vitali 集 \(V \subseteq [0, 1]\)。

设 \(\{r_j\}_{j=1}^{\infty}\) 为 \([-1, 1] \cap \mathbb{Q}\) 的枚举。则 \(V + r_j\) 两两不交（若 \(v_1 + r_j = v_2 + r_k\)，则 \(v_1 - v_2 \in \mathbb{Q}\)，故 \(v_1 = v_2\)，\(r_j = r_k\)），且 \([0, 1] \subseteq \bigcup_j (V + r_j) \subseteq [-1, 2]\)。

若 \(V\) 可测，则 \(m(V + r_j) = m(V)\)（平移不变性）。由可数可加性，\(m(\bigcup_j (V + r_j)) = \sum_j m(V)\)。若 \(m(V) = 0\)，则 \(m(\bigcup (V + r_j)) = 0\)，但 \([0, 1] \subseteq \bigcup (V + r_j)\)，故 \(m([0, 1]) = 0\)，矛盾。若 \(m(V) > 0\)，则 \(\sum m(V) = \infty\)，但 \(\bigcup (V + r_j) \subseteq [-1, 2]\) 测度有限，矛盾。故 \(V\) 不可测。\(\blacksquare\)

\hrulefill

\hrulefill

\hrulefill

\hrulefill

\hrulefill

\hrulefill

\section{第154章：可测函数}\label{ux7b2c154ux7ae0ux53efux6d4bux51fdux6570}

可测函数是 Lebesgue 积分理论中可积函数的自然定义域。本章系统研究可测函数的基本性质：可测性的多种等价刻画、可测函数对代数运算和极限运算的封闭性、简单函数逼近定理、以及 Egorov 定理和 Lusin 定理所揭示的可测函数与连续函数之间的深刻联系。可测函数类构成了一个极其广泛且运算封闭的函数空间------这是 Lebesgue 积分理论相较于 Riemann 积分的核心优势之一。

\subsection{46.1 可测函数的定义}\label{ux53efux6d4bux51fdux6570ux7684ux5b9aux4e49}

\textbf{定义 46.1.1（可测函数）}：设 \(E \subseteq \mathbb{R}^n\) 可测，\(f : E \to \overline{\mathbb{R}} = \mathbb{R} \cup \{\pm\infty\}\)。若对任意 \(\alpha \in \mathbb{R}\)，

\[
\{x \in E \mid f(x) > \alpha\}
\]

是可测集，则称 \(f\) 为 \(E\) 上的 \textbf{Lebesgue 可测函数}（简称可测函数）。

\textbf{定理 46.1.1（可测性的等价条件）}：以下条件等价：

\begin{enumerate}
\def\labelenumi{(\roman{enumi})}
\tightlist
\item
  \(\{f > \alpha\}\) 可测（\(\forall \alpha \in \mathbb{R}\)）
\item
  \(\{f \geq \alpha\}\) 可测（\(\forall \alpha \in \mathbb{R}\)）
\item
  \(\{f < \alpha\}\) 可测（\(\forall \alpha \in \mathbb{R}\)）
\item
  \(\{f \leq \alpha\}\) 可测（\(\forall \alpha \in \mathbb{R}\)）
\end{enumerate}

\textbf{证明}：(i) \(\implies\) (ii)：\(\{f \geq \alpha\} = \bigcap_{k=1}^{\infty} \{f > \alpha - 1/k\}\)，可测集的可数交可测。

\begin{enumerate}
\def\labelenumi{(\roman{enumi})}
\setcounter{enumi}{1}
\item
  \(\implies\) (iii)：\(\{f < \alpha\} = E \setminus \{f \geq \alpha\}\)。
\item
  \(\implies\) (iv)：\(\{f \leq \alpha\} = \bigcap_{k=1}^{\infty} \{f < \alpha + 1/k\}\)。
\item
  \(\implies\) (i)：\(\{f > \alpha\} = E \setminus \{f \leq \alpha\}\)。\(\blacksquare\)
\end{enumerate}

\textbf{定理 46.1.2（可测函数的运算）}：设 \(f, g\) 为 \(E\) 上的可测函数（取值有限 a.e.），则

\begin{enumerate}
\def\labelenumi{(\roman{enumi})}
\tightlist
\item
  \(f + g\)，\(f - g\)，\(fg\) 可测（在定义有意义处）
\item
  \(f/g\) 可测（若 \(g \neq 0\) a.e.）
\item
  \(\max(f, g)\)，\(\min(f, g)\)，\(|f|\) 可测
\item
  若 \(\{f_k\}\) 为可测函数列，则 \(\sup f_k\)，\(\inf f_k\)，\(\limsup f_k\)，\(\liminf f_k\) 可测
\end{enumerate}

\textbf{证明}：

\begin{enumerate}
\def\labelenumi{(\roman{enumi})}
\item
  \(\{f + g > \alpha\} = \bigcup_{r \in \mathbb{Q}} (\{f > r\} \cap \{g > \alpha - r\})\)，可测集的可数并可测。
\item
  \(\{\max(f, g) > \alpha\} = \{f > \alpha\} \cup \{g > \alpha\}\)。
\item
  \(\{\sup_k f_k > \alpha\} = \bigcup_k \{f_k > \alpha\}\)。\(\blacksquare\)
\end{enumerate}

\hrulefill

\subsection{46.2 简单函数及其逼近}\label{ux7b80ux5355ux51fdux6570ux53caux5176ux903cux8fd1}

\textbf{定义 46.2.1（简单函数）}：\(E\) 上的\textbf{简单函数} \(\varphi\) 是形如

\[
\varphi(x) = \sum_{j=1}^{k} a_j \chi_{E_j}(x)
\]

的函数，其中 \(E_j\) 为两两不交的可测集，\(\chi_{E_j}\) 为特征函数，\(a_j \in \mathbb{R}\)。

\textbf{定理 46.2.1（简单函数逼近定理）}：设 \(f\) 为 \(E\) 上的非负可测函数。则存在递增的简单函数列 \(\{\varphi_k\}\) 使得 \(\varphi_k(x) \uparrow f(x)\)（逐点收敛）。

\textbf{证明}：对 \(k \in \mathbb{N}^+\)，将 \([0, k]\) 划分为 \(k \cdot 2^k\) 个长度为 \(1/2^k\) 的小区间。定义

\[
\varphi_k(x) = \begin{cases}
\frac{j-1}{2^k}, & \text{若 } \frac{j-1}{2^k} \leq f(x) < \frac{j}{2^k}, \; j = 1, \ldots, k \cdot 2^k \\
k, & \text{若 } f(x) \geq k
\end{cases}
\]

则 \(\varphi_k\) 是简单函数，\(\varphi_k \leq \varphi_{k+1}\)，且 \(\lim_{k \to \infty} \varphi_k(x) = f(x)\)（\(\forall x \in E\)）。\(\blacksquare\)

\textbf{推论 46.2.2}：任意可测函数 \(f\) 可表示为简单函数列的逐点极限（通过分解 \(f = f^+ - f^-\)）。

\hrulefill

\subsection{46.3 Egorov 定理与 Lusin 定理}\label{egorov-ux5b9aux7406ux4e0e-lusin-ux5b9aux7406}

\textbf{定理 46.3.1（Egorov 定理------几乎一致收敛）}：设 \(E \subseteq \mathbb{R}^n\) 可测且 \(m(E) < \infty\)，\(\{f_k\}\) 为 \(E\) 上的可测函数列，\(f_k \to f\) a.e.。则对任意 \(\varepsilon > 0\)，存在可测子集 \(A \subseteq E\) 使得 \(m(E \setminus A) < \varepsilon\) 且 \(f_k \to f\) 在 \(A\) 上\textbf{一致收敛}。

\textbf{证明}：对每个 \(j, k \in \mathbb{N}^+\)，定义

\[
E_{j,k} = \left\{x \in E \;\middle|\; \sup_{n \geq k} |f_n(x) - f(x)| > \frac{1}{j}\right\}
\]

由于 \(f_n \to f\) a.e.，\(\forall j\)，\(\bigcap_{k=1}^{\infty} E_{j,k}\) 为零测集。由 \(m(E) < \infty\) 和测度的上连续性，\(\lim_{k \to \infty} m(E_{j,k}) = 0\)。

对 \(\varepsilon > 0\)，取 \(k_j\) 使得 \(m(E_{j, k_j}) < \varepsilon / 2^j\)。令 \(A = E \setminus \bigcup_{j=1}^{\infty} E_{j, k_j}\)。则

\[
m(E \setminus A) = m\left(\bigcup_{j=1}^{\infty} E_{j, k_j}\right) \leq \sum_{j=1}^{\infty} m(E_{j, k_j}) < \varepsilon
\]

在 \(A\) 上，对任意 \(j\)，当 \(n \geq k_j\) 时 \(|f_n - f| \leq 1/j\)，故 \(f_n \to f\) 一致。\(\blacksquare\)

\textbf{定理 46.3.2（Lusin 定理------几乎连续）}：设 \(E \subseteq \mathbb{R}^n\) 可测且 \(m(E) < \infty\)，\(f\) 为 \(E\) 上 a.e. 有限的可测函数。则对任意 \(\varepsilon > 0\)，存在紧集 \(K \subseteq E\) 使得 \(m(E \setminus K) < \varepsilon\) 且 \(f|_K\) 连续。

\textbf{证明}：先考虑 \(f\) 为简单函数的情形：\(f = \sum_{j=1}^m a_j \chi_{E_j}\)（\(E_j\) 可测且两两不交）。由内正则性，对每个 \(E_j\) 存在紧集 \(K_j \subseteq E_j\) 使 \(m(E_j \setminus K_j) < \varepsilon/m\)。令 \(K = \bigcup_{j=1}^m K_j\)，则 \(K\) 为紧集，\(m(E \setminus K) < \varepsilon\)。\(f\) 在每个 \(K_j\) 上为常数，故 \(f|_K\) 连续。对一般可测函数 \(f\)，可取简单函数列 \(\{s_n\}\) 使 \(s_n \to f\) a.e.。由 Egorov 定理，存在可测子集 \(E_0 \subseteq E\) 使 \(m(E \setminus E_0) < \varepsilon/2\) 且 \(s_n \to f\) 在 \(E_0\) 上一致。对每个 \(s_n\)，由简单函数情形，存在紧集 \(K_n \subseteq E_0\) 使 \(s_n|_{K_n}\) 连续且 \(m(E_0 \setminus K_n) < \varepsilon/(2^{n+1})\)。令 \(K = \bigcap_{n=1}^\infty K_n\)，则 \(K\) 为紧集，\(m(E \setminus K) \leq m(E \setminus E_0) + \sum m(E_0 \setminus K_n) < \varepsilon\)。\(s_n|_K \to f|_K\) 一致，且 \(s_n|_K\) 连续，故 \(f|_K\) 连续。\(\blacksquare\)

\hrulefill

\hrulefill

\hrulefill

\hrulefill

\hrulefill

\newpage

\chapter{06-分析/01-实分析/02-Lebesgue积分.md}\label{ux5206ux679001-ux5b9eux5206ux679002-lebesgueux79efux5206.md}

\section{第155章：Lebesgue 积分}\label{ux7b2c155ux7ae0lebesgue-ux79efux5206}

Lebesgue 积分理论是实分析的核心，由 Henri Lebesgue 于 1902 年在其博士论文《积分、长度与面积》中创立。与 Riemann 积分通过对定义域进行分割不同，Lebesgue 积分通过对值域进行分割，从而将积分理论建立在测度论的坚实基础上。这一视角的转变带来了深刻的后果：Lebesgue 积分理论中的极限定理（单调收敛定理、Fatou 引理、控制收敛定理）极为强大且易于使用，使得积分与极限的交换条件远比 Riemann 积分宽松。此外，Lebesgue 积分自然地统一了离散求和与连续积分，为概率论、调和分析和泛函分析提供了统一的语言。

\subsection{47.1 非负简单函数的积分}\label{ux975eux8d1fux7b80ux5355ux51fdux6570ux7684ux79efux5206}

简单函数是 Lebesgue 积分理论的基本建筑块，其角色类似于 Riemann 积分中的阶梯函数，但远比阶梯函数灵活------简单函数基于可测集的分划，而可测集可以是任意复杂的集合。

\textbf{定义 47.1.1（非负简单函数的积分）}：设 \(\varphi = \sum_{j=1}^{k} a_j \chi_{E_j}\) 为 \(E\) 上的非负简单函数（\(a_j \geq 0\)，\(E_j\) 两两不交可测）。定义 \(\varphi\) 的 Lebesgue 积分为

\[
\int_E \varphi \, dm = \sum_{j=1}^{k} a_j \, m(E_j)
\]

（约定 \(0 \cdot \infty = 0\)。）

\textbf{注记 47.1.1}：上述定义与 \(\varphi\) 的具体表示无关。即若 \(\varphi = \sum_{j=1}^{k} a_j \chi_{E_j} = \sum_{i=1}^{\ell} b_i \chi_{F_i}\) 为两种不同表示，则 \(\sum a_j m(E_j) = \sum b_i m(F_i)\)。这是因为两种表示可通过取 \(E_j \cap F_i\) 的公共细化来统一，由测度的可加性，积分值不变。这一性质是 Lebesgue 积分定义良好性的基础。

\textbf{定理 47.1.1（简单函数积分的线性性与单调性）}：

\begin{enumerate}
\def\labelenumi{(\roman{enumi})}
\tightlist
\item
  \textbf{线性性}：\(\int_E (a\varphi + b\psi) = a\int_E \varphi + b\int_E \psi\)（\(a, b \geq 0\)）
\item
  \textbf{单调性}：若 \(\varphi \leq \psi\)，则 \(\int_E \varphi \leq \int_E \psi\)
\end{enumerate}

\textbf{证明}：(i) 将 \(\varphi, \psi\) 用同一组可测集（它们的交集划分）表示，然后验证。具体地，设 \(\varphi = \sum a_j \chi_{E_j}\)，\(\psi = \sum b_i \chi_{F_i}\)。令 \(G_{ji} = E_j \cap F_i\)，则 \(\{G_{ji}\}\) 为两两不交的可测集族，且 \(\varphi = \sum_{j,i} a_j \chi_{G_{ji}}\)，\(\psi = \sum_{j,i} b_i \chi_{G_{ji}}\)。于是 \(a\varphi + b\psi = \sum_{j,i} (a a_j + b b_i) \chi_{G_{ji}}\)，积分值为 \(\sum_{j,i} (a a_j + b b_i) m(G_{ji}) = a\sum_j a_j m(E_j) + b\sum_i b_i m(F_i)\)。

\begin{enumerate}
\def\labelenumi{(\roman{enumi})}
\setcounter{enumi}{1}
\tightlist
\item
  由 (i)，\(\psi - \varphi \geq 0\)，故 \(\int \psi - \int \varphi = \int (\psi - \varphi) \geq 0\)。\(\blacksquare\)
\end{enumerate}

\textbf{命题 47.1.2（非负简单函数的逼近性质）}：设 \(f\) 为 \(E\) 上的非负可测函数。则存在递增的非负简单函数列 \(\{\varphi_k\}\) 使得 \(\varphi_k \uparrow f\)（逐点收敛）。具体构造为：对每个 \(k\)，将 \([0, k]\) 等分为 \(k \cdot 2^k\) 份，令

\[
\varphi_k(x) = \begin{cases} \frac{m-1}{2^k} & \text{若 } \frac{m-1}{2^k} \leq f(x) < \frac{m}{2^k}, \; m = 1, 2, \ldots, k \cdot 2^k \\ k & \text{若 } f(x) \geq k \end{cases}
\]

则 \(\varphi_k\) 是简单函数（因为每个水平集 \(\{f \geq c\}\) 均可测），且 \(\varphi_k \leq \varphi_{k+1}\)，\(\varphi_k \to f\) 逐点。

\textbf{证明}：由 \(f\) 的可测性，每个水平集 \(\{x : \frac{m-1}{2^k} \leq f(x) < \frac{m}{2^k}\}\) 和 \(\{x : f(x) \geq k\}\) 均为可测集，故 \(\varphi_k\) 为简单函数。单调性由构造方式直接保证。收敛性：对任意 \(x\)，若 \(f(x) < \infty\)，则当 \(k > f(x)\) 时，\(f(x) - \varphi_k(x) \leq 2^{-k} \to 0\)。若 \(f(x) = \infty\)，则 \(\varphi_k(x) = k \to \infty\)。\(\blacksquare\)

\hrulefill

\subsection{47.2 非负可测函数的积分}\label{ux975eux8d1fux53efux6d4bux51fdux6570ux7684ux79efux5206}

\textbf{定义 47.2.1（非负可测函数的积分）}：设 \(f\) 为 \(E\) 上的非负可测函数。定义

\[
\int_E f \, dm = \sup\left\{\int_E \varphi \, dm \;\middle|\; 0 \leq \varphi \leq f, \; \varphi \text{ 为简单函数}\right\}
\]

这一定义是 Lebesgue 积分理论的核心创新：不再通过分割定义域，而是通过从下方用简单函数逼近来定义积分。这一''由下而上''的逼近策略使得积分自动具有优良的极限性质。

\textbf{命题 47.2.1（积分的初等性质）}：设 \(f, g\) 为 \(E\) 上的非负可测函数。

\begin{enumerate}
\def\labelenumi{(\roman{enumi})}
\tightlist
\item
  \textbf{单调性}：若 \(f \leq g\) a.e.，则 \(\int_E f \leq \int_E g\)
\item
  \textbf{齐次性}：\(\int_E c f = c \int_E f\)（\(c \geq 0\)）
\item
  \textbf{可加性}：\(\int_E (f + g) = \int_E f + \int_E g\)
\item
  \textbf{对定义域的有限可加性}：若 \(E = A \cup B\) 且 \(A \cap B = \emptyset\)，则 \(\int_E f = \int_A f + \int_B f\)
\end{enumerate}

\textbf{证明}：(i) 由定义，任何不超过 \(f\) 的简单函数也不超过 \(g\)，故 \(\int f\) 的上确界不超过 \(\int g\) 的上确界。(ii) 对简单函数成立，取上确界即得。(iii) 需证明 \(\sup\{\int \varphi : \varphi \leq f+g\} = \sup\{\int \varphi_1 : \varphi_1 \leq f\} + \sup\{\int \varphi_2 : \varphi_2 \leq g\}\)。由命题 47.1.2，存在简单函数列 \(\varphi_k \uparrow f\)，\(\psi_k \uparrow g\)，则 \(\varphi_k + \psi_k \uparrow f+g\)。由简单函数积分的可加性和单调收敛定理（见定理 47.3.1），可得等号成立。(iv) 由积分定义和测度的可加性直接验证。\(\blacksquare\)

\textbf{定理 47.2.2（Chebyshev 不等式 / Markov 不等式）}：设 \(f \geq 0\) 可测。则对任意 \(\lambda > 0\)，

\[
m(\{x \in E \mid f(x) \geq \lambda\}) \leq \frac{1}{\lambda} \int_E f \, dm
\]

\textbf{证明}：令 \(A = \{f \geq \lambda\}\)。则 \(f \geq \lambda \chi_A\)。由积分的单调性，\(\int_E f \geq \lambda \int_E \chi_A = \lambda m(A)\)。\(\blacksquare\)

\textbf{推论 47.2.3}：若 \(\int_E f = 0\)，则 \(f = 0\) a.e.。

\textbf{证明}：由 Chebyshev 不等式，对任意 \(n\)，\(m(\{f \geq 1/n\}) \leq n \int_E f = 0\)。故 \(\{f > 0\} = \bigcup_{n=1}^{\infty} \{f \geq 1/n\}\) 为零测集。\(\blacksquare\)

\textbf{推论 47.2.4}：若 \(f\) 可积（即 \(\int |f| < \infty\)），则 \(|f| < \infty\) a.e.。特别地，可积函数几乎处处有限。

\textbf{推论 47.2.5（积分的绝对连续性------初级形式）}：设 \(f \geq 0\) 可积。则对任意 \(\varepsilon > 0\)，存在 \(\delta > 0\) 使得对任意可测集 \(A\)，\(m(A) < \delta\) 蕴含 \(\int_A f < \varepsilon\)。

\textbf{证明}：反设存在 \(\varepsilon_0 > 0\) 和可测集列 \(A_n\) 满足 \(m(A_n) < 2^{-n}\) 但 \(\int_{A_n} f \geq \varepsilon_0\)。令 \(A = \limsup A_n = \bigcap_{N=1}^{\infty} \bigcup_{n=N}^{\infty} A_n\)。由 Borel-Cantelli 引理，\(m(A) = 0\)。但由 Fatou 引理逆命题（或控制收敛定理），\(\int_A f \geq \limsup \int_{A_n} f \geq \varepsilon_0\)，矛盾。\(\blacksquare\)

\hrulefill

\subsection{47.3 单调收敛定理}\label{ux5355ux8c03ux6536ux655bux5b9aux7406}

\textbf{定理 47.3.1（Lebesgue 单调收敛定理）}：设 \(\{f_k\}\) 为 \(E\) 上的非负可测函数列，满足 \(0 \leq f_1 \leq f_2 \leq \cdots\)，且 \(f_k \to f\)（逐点）。则

\[
\lim_{k \to \infty} \int_E f_k \, dm = \int_E f \, dm
\]

即极限与积分可交换（在非负递增情形下）。

\textbf{证明}：由于 \(f_k \leq f\)，\(\int f_k \leq \int f\)，故 \(\lim_{k \to \infty} \int f_k \leq \int f\)。需证反向不等式。

取任意简单函数 \(0 \leq \varphi \leq f\) 和 \(0 < \alpha < 1\)。定义 \(E_k = \{x \mid f_k(x) \geq \alpha \varphi(x)\}\)。则 \(E_1 \subseteq E_2 \subseteq \cdots\) 且 \(\bigcup E_k = E\)（因为 \(f_k \uparrow f\) 且 \(\alpha \varphi < f\) 在 \(\varphi > 0\) 处）。

由测度的下连续性，\(m(E_k) \uparrow m(E)\)。且 \(\int_E f_k \geq \int_{E_k} f_k \geq \alpha \int_{E_k} \varphi\)。

对简单函数 \(\varphi = \sum a_j \chi_{F_j}\)，\(\int_{E_k} \varphi = \sum a_j m(F_j \cap E_k) \uparrow \sum a_j m(F_j) = \int_E \varphi\)（\(k \to \infty\)）。

因此 \(\lim_{k \to \infty} \int_E f_k \geq \alpha \int_E \varphi\)。令 \(\alpha \to 1\) 得 \(\lim \int f_k \geq \int \varphi\)。再对 \(\varphi \leq f\) 取上确界，得 \(\lim \int f_k \geq \int f\)。\(\blacksquare\)

\textbf{推论 47.3.2（非负可测函数级数的逐项积分 / Beppo Levi 定理）}：若 \(\{g_k\}\) 为非负可测函数列，则

\[
\int_E \sum_{k=1}^{\infty} g_k \, dm = \sum_{k=1}^{\infty} \int_E g_k \, dm
\]

\textbf{证明}：令 \(f_k = \sum_{j=1}^{k} g_j\)，应用单调收敛定理。\(\blacksquare\)

\textbf{注记 47.3.1}：单调收敛定理是 Lebesgue 积分理论中最基本的极限定理。它表明：对于非负递增的函数列，积分和极限总是可以交换。这一性质在 Riemann 积分框架下完全不成立------Riemann 可积函数列的单调极限甚至可能不是 Riemann 可积的。Beppo Levi 定理（级数逐项积分）是单调收敛定理的直接推论，它使得我们可以将重积分与无穷级数自由交换，这在概率论中尤为有用。

\textbf{命题 47.3.3（单调收敛定理的测度构造应用）}：设 \(f \geq 0\) 可测。定义测度 \(\nu(A) = \int_A f \, dm\)。则 \(\nu\) 是 \(\sigma\)-有限的测度（若 \(f\) 可积），且 \(\nu\) 关于 \(m\) 绝对连续。此即 Radon-Nikodym 定理中 \(f\) 作为密度函数的基本构造。

\hrulefill

\subsection{47.4 Fatou 引理}\label{fatou-ux5f15ux7406}

\textbf{定理 47.4.1（Fatou 引理）}：设 \(\{f_k\}\) 为 \(E\) 上的非负可测函数列。则

\[
\int_E \liminf_{k \to \infty} f_k \, dm \leq \liminf_{k \to \infty} \int_E f_k \, dm
\]

\textbf{证明}：令 \(g_k = \inf_{j \geq k} f_j\)。则 \(0 \leq g_1 \leq g_2 \leq \cdots\)，\(g_k \uparrow \liminf f_k\)。由单调收敛定理，\(\int \liminf f_k = \lim_{k \to \infty} \int g_k\)。

又 \(g_k \leq f_j\)（\(\forall j \geq k\)），故 \(\int g_k \leq \inf_{j \geq k} \int f_j\)。因此

\[
\int \liminf f_k = \lim_{k \to \infty} \int g_k \leq \lim_{k \to \infty} \inf_{j \geq k} \int f_j = \liminf_{k \to \infty} \int f_k
\]

\(\blacksquare\)

\textbf{定理 47.4.2（反向 Fatou 引理）}：设 \(\{f_k\}\) 为 \(E\) 上的可测函数列，满足 \(|f_k| \leq g\) a.e. 且 \(g\) 可积。则

\[
\limsup_{k \to \infty} \int_E f_k \, dm \leq \int_E \limsup_{k \to \infty} f_k \, dm
\]

\textbf{证明}：将 Fatou 引理应用于非负函数列 \(g - f_k\)：

\[
\int \liminf (g - f_k) \leq \liminf \int (g - f_k)
\]

左边 \(\int \liminf (g - f_k) = \int (g - \limsup f_k) = \int g - \int \limsup f_k\)。右边 \(\liminf \int (g - f_k) = \liminf (\int g - \int f_k) = \int g - \limsup \int f_k\)。因此 \(-\int \limsup f_k \leq -\limsup \int f_k\)，即 \(\limsup \int f_k \leq \int \limsup f_k\)。\(\blacksquare\)

\textbf{注记 47.4.1}：Fatou 引理在不等式方向上是''损耗性的''------它只给出下极限的不等式，且方向可能带来信息损失。这体现了''积分与极限交换''在缺乏可积控制时的代价。反向 Fatou 引理则表明，在有可积控制函数 \(g\) 的条件下，上极限的方向也可以反转。Fatou 引理是证明控制收敛定理的关键工具，其优雅之处在于它不需要任何一致收敛性条件，仅需要非负性。

\hrulefill

\subsection{47.5 一般可测函数的积分与 Lebesgue 控制收敛定理}\label{ux4e00ux822cux53efux6d4bux51fdux6570ux7684ux79efux5206ux4e0e-lebesgue-ux63a7ux5236ux6536ux655bux5b9aux7406}

\textbf{定义 47.5.1（可积函数）}：设 \(f\) 为 \(E\) 上的可测函数。将 \(f\) 分解为正部和负部：\(f = f^+ - f^-\)，其中 \(f^+ = \max(f, 0)\)，\(f^- = \max(-f, 0)\)。

若 \(\int_E f^+ < \infty\) 且 \(\int_E f^- < \infty\)，则称 \(f\) 在 \(E\) 上 \textbf{Lebesgue 可积}，定义

\[
\int_E f \, dm = \int_E f^+ \, dm - \int_E f^- \, dm
\]

\(L^1(E)\) 表示 \(E\) 上所有 Lebesgue 可积函数的集合。

\textbf{命题 47.5.1（可积函数的基本性质）}：设 \(f, g \in L^1(E)\)。

\begin{enumerate}
\def\labelenumi{(\roman{enumi})}
\tightlist
\item
  \textbf{线性性}：\(L^1(E)\) 是向量空间，且 \(\int (af + bg) = a\int f + b\int g\)
\item
  \textbf{三角不等式}：\(|\int_E f| \leq \int_E |f|\)
\item
  \textbf{单调性}：若 \(f \leq g\) a.e.，则 \(\int f \leq \int g\)
\item
  \textbf{绝对连续性}：\(f \in L^1(E)\) 等价于 \(|f| \in L^1(E)\)
\end{enumerate}

\textbf{证明}：(i) 由正负部分解验证。\((f+g)^+ - (f+g)^- = f^+ - f^- + g^+ - g^-\)，移项得 \((f+g)^+ + f^- + g^- = (f+g)^- + f^+ + g^+\)。两边积分并利用非负函数积分的可加性即得。(ii) \(|\int f| = |\int f^+ - \int f^-| \leq \int f^+ + \int f^- = \int |f|\)。(iii) 由 \(g - f \geq 0\) 和非负函数积分的定义。(iv) \(|f| = f^+ + f^-\)，故 \(\int |f| < \infty\) 当且仅当 \(\int f^+ < \infty\) 且 \(\int f^- < \infty\)。\(\blacksquare\)

\textbf{定理 47.5.2（Lebesgue 控制收敛定理）}：设 \(\{f_k\}\) 为 \(E\) 上的可测函数列，满足 \(f_k \to f\) a.e.。若存在可积函数 \(g\) 使得 \(|f_k| \leq g\) a.e.（\(\forall k\)），则 \(f\) 可积且

\[
\lim_{k \to \infty} \int_E f_k \, dm = \int_E f \, dm
\]

\textbf{证明}：由 \(|f_k| \leq g\) 及 \(g\) 可积，\(f_k\) 可积。对 \(f\)，\(|f| = \lim |f_k| \leq g\) a.e.，故 \(f\) 也可积。

考虑 \(g + f_k \geq 0\) 和 \(g - f_k \geq 0\)。由 Fatou 引理：

\[
\int_E (g + f) \leq \liminf_{k \to \infty} \int_E (g + f_k) = \int_E g + \liminf_{k \to \infty} \int_E f_k
\]

故 \(\int_E f \leq \liminf_{k \to \infty} \int_E f_k\)。

同理对 \(g - f_k\)：

\[
\int_E (g - f) \leq \liminf_{k \to \infty} \int_E (g - f_k) = \int_E g - \limsup_{k \to \infty} \int_E f_k
\]

故 \(\int_E f \geq \limsup_{k \to \infty} \int_E f_k\)。

因此 \(\liminf \int f_k = \limsup \int f_k = \int f\)，即 \(\lim \int f_k = \int f\)。\(\blacksquare\)

\textbf{定理 47.5.3（Vitali 收敛定理）}：设 \(\{f_k\} \subseteq L^1(E)\)，\(f_k \to f\)（依测度），且 \(\{f_k\}\) 一致可积。则 \(f \in L^1(E)\) 且 \(f_k \to f\) 在 \(L^1\) 范数下。

\textbf{定义 47.5.2（一致可积性）}：函数族 \(\{f_k\} \subseteq L^1(E)\) 称为\textbf{一致可积}，若对任意 \(\varepsilon > 0\)，存在 \(\delta > 0\) 使得对任意可测集 \(A\)，\(m(A) < \delta\) 蕴含 \(\int_A |f_k| < \varepsilon\)（对所有 \(k\)），且存在 \(M\) 使得 \(\int_{|f_k| > M} |f_k| < \varepsilon\)（对所有 \(k\)）。

\textbf{证明}：由一致可积性，\(\{f_k\}\) 在 \(L^1\) 中范数有界。由 Fatou 引理（逐点收敛的子列），\(f \in L^1\)。对任意 \(\varepsilon > 0\)，选取 \(\delta > 0\) 和 \(M\) 满足一致可积性条件。由 Egorov 定理，在 \(E\) 的一个大子集上 \(f_k \to f\) 一致收敛，在其余小集上积分由一致可积性控制。\(\varepsilon/3\) 论证给出 \(\|f_k - f\|_1 \to 0\)。\(\blacksquare\)

\hrulefill

\subsection{47.6 Lebesgue 积分与 Riemann 积分的比较}\label{lebesgue-ux79efux5206ux4e0e-riemann-ux79efux5206ux7684ux6bd4ux8f83}

\textbf{定理 47.6.1}：设 \(f\) 在 \([a, b]\) 上 Riemann 可积。则 \(f\) 在 \([a, b]\) 上 Lebesgue 可积，且两种积分值相等：

\[
\int_{[a, b]} f \, dm = \int_a^b f(x) \, dx
\]

\textbf{证明}：Riemann 可积函数必有界，且其不连续点构成零测集（Lebesgue 判别法）。有界可测函数在有限测度集上 Lebesgue 可积。

比较两种积分：Riemann 积分用阶梯函数（基于区间分割）逼近，Lebesgue 积分用简单函数逼近。取 \([a, b]\) 的一列细化分割 \(P_n\)，其上下 Riemann 和分别给出 \(f\) 的简单函数逼近。由 Riemann 可积性，上下 Riemann 积分相等，且等于 Lebesgue 积分。具体地，Riemann 下和给出的简单函数列递增到 \(f\) 的某个下逼近，由单调收敛定理其 Lebesgue 积分收敛到 Riemann 积分。\(\blacksquare\)

\textbf{定理 47.6.2（Lebesgue 可积函数类更广）}：存在 Lebesgue 可积但 Riemann 不可积的函数。

\textbf{例}：Dirichlet 函数 \(\chi_{\mathbb{Q} \cap [0, 1]}\) 在 \([0, 1]\) 上 Lebesgue 可积（值为 \(0\)，因为 \(\mathbb{Q} \cap [0, 1]\) 可数故零测），但 Riemann 不可积（因为其上下 Riemann 和分别为 \(1\) 和 \(0\)）。

\textbf{定理 47.6.3（Riemann 可积性的 Lebesgue 刻画）}：\([a, b]\) 上的有界函数 \(f\) Riemann 可积当且仅当 \(f\) 的不连续点集为零测集。

\textbf{证明}：设 \(f\) 的有界性常数为 \(M\)。\(f\) 在 \(x\) 处的振幅为 \(\omega_f(x) = \lim_{\delta \to 0} \sup_{y,z \in (x-\delta, x+\delta)} |f(y) - f(z)|\)。\(f\) 在 \(x\) 处连续当且仅当 \(\omega_f(x) = 0\)。Riemann 可积的充要条件是：对任意 \(\varepsilon > 0\)，\(\{x : \omega_f(x) \geq \varepsilon\}\) 是零测集（由 Riemann 可积的 Darboux 刻画）。取 \(\varepsilon = 1/n\)，得 \(\{x : \omega_f(x) > 0\} = \bigcup_{n=1}^{\infty} \{x : \omega_f(x) \geq 1/n\}\) 为零测集。\(\blacksquare\)

\textbf{注记 47.6.1}：Lebesgue 积分相比于 Riemann 积分的优势不仅在于可积函数类更广，更在于其极限定理的优越性。在 Riemann 积分框架下，逐点收敛的函数列即使一致有界，其极限也可能不可积。Lebesgue 控制收敛定理则只需要一个可积的控制函数，使用极为方便。此外，在无界区域和无界函数的情形下，Lebesgue 积分的处理也比 Riemann 反常积分更为系统和优雅。

\subsection{47.7 有界变差函数与 Lebesgue-Stieltjes 积分}\label{ux6709ux754cux53d8ux5deeux51fdux6570ux4e0e-lebesgue-stieltjes-ux79efux5206}

\textbf{定义 47.7.1（有界变差函数）}：函数 \(f: [a, b] \to \mathbb{R}\) 称为\textbf{有界变差}的，若其全变差 \(V_a^b(f) = \sup \sum_{i=1}^n |f(x_i) - f(x_{i-1})| < \infty\)，其中上确界取遍 \([a, b]\) 的所有分割 \(a = x_0 < x_1 < \cdots < x_n = b\)。

\textbf{定理 47.7.1（Jordan 分解定理）}：\(f\) 在 \([a, b]\) 上有界变差当且仅当 \(f\) 可表示为两个单调递增函数之差：\(f = g - h\)，其中 \(g, h\) 单调递增。具体可取 \(g(x) = V_a^x(f)\)（全变差函数），\(h(x) = V_a^x(f) - f(x)\)。

\textbf{证明}：若 \(f = g - h\) 且 \(g, h\) 单调递增，则对任意分割，\(\sum |f(x_i) - f(x_{i-1})| \leq \sum (g(x_i) - g(x_{i-1})) + \sum (h(x_i) - h(x_{i-1})) = (g(b) - g(a)) + (h(b) - h(a))\)，故 \(V_a^b(f) < \infty\)。

反之，设 \(V_a^b(f) < \infty\)。定义 \(g(x) = V_a^x(f)\)。则 \(g\) 单调递增，且 \(g(x) \pm f(x)\) 均为单调递增（因为 \(g(y) - g(x) \geq |f(y) - f(x)|\)）。取 \(h = g - f\) 即得 \(h\) 单调递增。\(\blacksquare\)

\textbf{推论 47.7.2}：有界变差函数几乎处处可导（由 Lebesgue 定理，单调函数几乎处处可导）。

\textbf{定理 47.7.3（Lebesgue-Stieltjes 积分）}：设 \(\alpha: [a, b] \to \mathbb{R}\) 为有界变差函数。由 \(\alpha\) 诱导的 Lebesgue-Stieltjes 测度 \(\mu_\alpha\) 满足 \(\mu_\alpha((c, d]) = \alpha(d) - \alpha(c)\)（对 \(a \leq c < d \leq b\)）。Lebesgue-Stieltjes 积分 \(\int_a^b f \, d\alpha = \int_{[a,b]} f \, d\mu_\alpha\) 统一了 Riemann-Stieltjes 积分和 Lebesgue 积分：当 \(\alpha(x) = x\) 时退化为通常的 Lebesgue 积分；当 \(\alpha\) 为阶梯函数时退化为离散求和（即 \(\sum f(x_j) \Delta \alpha_j\)）。

\textbf{证明}：由 Riesz 表示定理，\(C[a,b]\) 上的正线性泛函一一对应于有限 Borel 测度。有界变差函数 \(\alpha\) 通过 Riemann-Stieltjes 积分 \(f \mapsto \int_a^b f \, d\alpha\)（对 \(f \in C[a,b]\)）定义了 \(C[a,b]\) 上的线性泛函。由 Jordan 分解，\(\alpha\) 可分解为两个单调递增函数之差，从而该泛函可分解为正泛函之差，对应一个符号 Borel 测度 \(\mu_\alpha\)。\(\blacksquare\)

\textbf{定理 47.7.4（Lebesgue-Stieltjes 积分的 Radon-Nikodym 表示）}：若 \(\alpha\) 是绝对连续的，则 \(d\alpha = \alpha'(x) dx\)，即 \(\mu_\alpha \ll m\)，且 \(\frac{d\mu_\alpha}{dm} = \alpha'\)。此时 Lebesgue-Stieltjes 积分退化为 \(\int_a^b f(x) \alpha'(x) dx\)。

\subsection{47.8 微分与积分的互逆}\label{ux5faeux5206ux4e0eux79efux5206ux7684ux4e92ux9006}

\textbf{定理 47.8.1（Lebesgue 微分定理）}：设 \(f \in L^1_{\text{loc}}(\mathbb{R}^n)\)。则对 a.e. \(x \in \mathbb{R}^n\)，

\[\lim_{r \to 0} \frac{1}{m(B(x, r))} \int_{B(x, r)} f(y) \, dm(y) = f(x)\]

\textbf{定义 47.8.1（Hardy-Littlewood 极大函数）}：对 \(f \in L^1_{\text{loc}}(\mathbb{R}^n)\)，其 Hardy-Littlewood 极大函数定义为

\[
Mf(x) = \sup_{r > 0} \frac{1}{m(B(x, r))} \int_{B(x, r)} |f(y)| \, dy
\]

\textbf{定理 47.8.2（极大函数的弱 \((1,1)\) 型估计）}：存在常数 \(C_n > 0\)（仅依赖于维数 \(n\)）使得对任意 \(\lambda > 0\)，

\[
m(\{x \in \mathbb{R}^n \mid Mf(x) > \lambda\}) \leq \frac{C_n}{\lambda} \int_{\mathbb{R}^n} |f(y)| \, dy
\]

\textbf{证明}（Lebesgue 微分定理的框架）：利用 Hardy-Littlewood 极大函数的弱 \((1,1)\) 型估计和 Vitali 覆盖引理。对任意 \(\varepsilon > 0\)，取 \(g \in C_c(\mathbb{R}^n)\) 使得 \(\|f - g\|_1 < \varepsilon\)（连续函数在 \(L^1\) 中的稠密性）。令 \(E_\alpha = \{x : \limsup_{r \to 0} |\frac{1}{|B(x,r)|} \int_{B(x,r)} f - f(x)| > \alpha\}\)。利用 \(f = (f - g) + g\) 分解，\(g\) 部分的极限为 \(0\)（由连续性），\((f-g)\) 部分由极大函数估计控制。由 \(\varepsilon\) 和 \(\alpha\) 的任意性，\(E_\alpha\) 为零测集，对所有 \(\alpha > 0\) 取可数并即得结论。\(\blacksquare\)

\textbf{定理 47.8.3（Newton-Leibniz 公式的 Lebesgue 版本）}：设 \(f\) 在 \([a, b]\) 上 Lebesgue 可积。定义 \(F(x) = \int_a^x f(t) \, dm(t)\)。则 \(F\) 在 \([a, b]\) 上绝对连续，且 \(F'(x) = f(x)\) a.e.。反之，若 \(F\) 在 \([a, b]\) 上绝对连续，则 \(F'\) 存在 a.e.，\(F' \in L^1([a, b])\)，且 \(F(x) - F(a) = \int_a^x F'(t) \, dm(t)\)。

\textbf{证明}：由积分的绝对连续性，\(F\) 是绝对连续的。\(F'\) 的 a.e. 存在性来自 Lebesgue 微分定理。反向：绝对连续函数有界变差，故 \(F'\) 存在 a.e. 且可积。由 \(F\) 的绝对连续性，两个函数 \(F(x)\) 和 \(\int_a^x F'(t) dt\) 有相同的导数 a.e. 且都绝对连续，故差为常数（由 Lebesgue 微分定理的推论）。\(\blacksquare\)

\textbf{定义 47.8.2（绝对连续函数）}：函数 \(F: [a, b] \to \mathbb{R}\) 称为\textbf{绝对连续}的，若对任意 \(\varepsilon > 0\)，存在 \(\delta > 0\) 使得对任意有限个互不相交的开区间 \((a_k, b_k) \subseteq [a, b]\)，满足 \(\sum (b_k - a_k) < \delta\) 时，\(\sum |F(b_k) - F(a_k)| < \varepsilon\)。

\textbf{命题 47.8.4}：绝对连续函数必为有界变差函数，且满足 Lusin 性质（将零测集映为零测集）。绝对连续函数是满足 Newton-Leibniz 公式的函数类的精确刻画。

\hrulefill

\subsection{47.9 乘积测度与 Fubini 定理}\label{ux4e58ux79efux6d4bux5ea6ux4e0e-fubini-ux5b9aux7406}

\textbf{定理 47.9.1（Fubini-Tonelli 定理）}：设 \((X, \mathcal{A}, \mu)\) 和 \((Y, \mathcal{B}, \nu)\) 为 \(\sigma\)-有限的测度空间。

\begin{enumerate}
\def\labelenumi{(\roman{enumi})}
\tightlist
\item
  \textbf{Tonelli 定理}（非负情形）：若 \(f: X \times Y \to [0, \infty]\) 为 \(\mathcal{A} \otimes \mathcal{B}\)-可测函数，则
\end{enumerate}

\[
\int_X \int_Y f(x, y) \, d\nu(y) \, d\mu(x) = \int_Y \int_X f(x, y) \, d\mu(x) \, d\nu(y) = \int_{X \times Y} f \, d(\mu \times \nu)
\]

\begin{enumerate}
\def\labelenumi{(\roman{enumi})}
\setcounter{enumi}{1}
\tightlist
\item
  \textbf{Fubini 定理}（可积情形）：若 \(f \in L^1(\mu \times \nu)\)，则 \(f(x, \cdot) \in L^1(\nu)\) 对 \(\mu\)-a.e. \(x\) 成立，\(f(\cdot, y) \in L^1(\mu)\) 对 \(\nu\)-a.e. \(y\) 成立，且上述积分等式成立。
\end{enumerate}

\textbf{证明}：先对 \(f\) 为特征函数 \(\chi_{A \times B}\)（\(A \in \mathcal{A}, B \in \mathcal{B}\)）验证：\(\int_X \int_Y \chi_{A \times B} d\nu d\mu = \int_X \nu(B) \chi_A d\mu = \mu(A)\nu(B) = (\mu \times \nu)(A \times B)\)。由线性性推广到非负简单函数，再由单调收敛定理推广到非负可测函数（Tonelli 定理）。Fubini 定理将 \(f\) 分解为正负部 \(f = f^+ - f^-\) 后应用 Tonelli 定理，由 \(\int |f| d(\mu \times \nu) < \infty\) 保证 \(f^+\) 和 \(f^-\) 的 Tonelli 积分有限，从而其中的截面函数对 a.e. 点为可积的。\(\blacksquare\)

\textbf{注记 47.9.1}：Fubini 定理是 Lebesgue 积分理论超越 Riemann 积分的最重要优势之一。在 Riemann 积分框架下，重积分交换次序需要函数满足一致连续性等苛刻条件；而在 Lebesgue 框架下，仅需可测性和绝对可积性（或非负性），使用极为方便。这在概率论中尤为重要------联合分布与边缘分布的关系本质上是 Fubini 定理的推论。条件期望的几乎所有性质（如 \(\mathbb{E}[\mathbb{E}[X \mid \mathcal{G}]] = \mathbb{E}[X]\)）都可以追溯到 Fubini 定理。

\textbf{推论 47.9.2（乘积测度的完备性）}：Lebesgue 测度 \(m_n\) 在 \(\mathbb{R}^n\) 上是一维 Lebesgue 测度 \(m_1\) 的 \(n\) 次乘积完备化。即 \(m_n = \overline{m_1 \times \cdots \times m_1}\)。

\hrulefill

\subsection{47.10 Egorov 定理与 Lusin 定理}\label{egorov-ux5b9aux7406ux4e0e-lusin-ux5b9aux7406-1}

\textbf{定理 47.10.1（Egorov 定理）}：设 \(E\) 为有限测度可测集，\(\{f_k\}\) 为 \(E\) 上的可测函数列，\(f_k \to f\) a.e.。则对任意 \(\varepsilon > 0\)，存在可测集 \(A \subseteq E\) 满足 \(m(E \setminus A) < \varepsilon\)，且在 \(A\) 上 \(f_k \to f\) 一致收敛。

\textbf{证明}：不妨设 \(f_k \to f\) 在 \(E\) 上处处收敛（去掉零测集不影响）。对 \(n, m \in \mathbb{N}\)，定义

\[
E_{n,m} = \bigcup_{k=n}^{\infty} \left\{x \in E : |f_k(x) - f(x)| \geq \frac{1}{m}\right\}
\]

对固定的 \(m\)，\(E_{n,m} \downarrow \emptyset\)（\(n \to \infty\)）。由 \(m(E) < \infty\) 和测度的上连续性，\(m(E_{n,m}) \to 0\)。对任意 \(\varepsilon > 0\)，取 \(n_m\) 使 \(m(E_{n_m, m}) < \varepsilon / 2^m\)。令 \(A = E \setminus \bigcup_{m=1}^{\infty} E_{n_m, m}\)。则 \(m(E \setminus A) \leq \sum_{m=1}^{\infty} m(E_{n_m, m}) < \varepsilon\)。且在 \(A\) 上，对所有 \(m\) 和 \(k \geq n_m\)，\(|f_k(x) - f(x)| < 1/m\)，故 \(f_k \to f\) 在 \(A\) 上一致收敛。\(\blacksquare\)

\textbf{注记 47.10.1}：Egorov 定理体现了 Lusin 的名言：``几乎处处收敛的函数在去掉一个任意小测度集后一致收敛。''它揭示了可测函数列几乎处处收敛与一致收敛之间的深刻联系。有限测度条件是必要的：考虑 \(\mathbb{R}\) 上 \(f_n = \chi_{[n, n+1]}\)，\(f_n \to 0\) 逐点，但无法在任意大测度集上一致收敛。

\textbf{定理 47.10.2（Lusin 定理）}：设 \(f\) 为 \(E\) 上的可测函数（\(m(E) < \infty\)）。则对任意 \(\varepsilon > 0\)，存在紧集 \(K \subseteq E\) 满足 \(m(E \setminus K) < \varepsilon\)，且 \(f|_K\) 连续。

\textbf{证明}：先对 \(f\) 为简单函数（特征函数）的情形验证，由可测集的外正则性（紧集逼近）和 Urysohn 引理推出。对一般可测函数，用简单函数逼近和 Egorov 定理。\(\blacksquare\)

Lusin 定理表明：可测函数''几乎''是连续函数。这一事实是 \(C_c(\mathbb{R}^n)\) 在 \(L^p(\mathbb{R}^n)\) 中稠密性的基础，也是实分析与泛函分析之间桥梁的基石之一。

\hrulefill

\subsection{47.11 Radon-Nikodym 定理与条件期望}\label{radon-nikodym-ux5b9aux7406ux4e0eux6761ux4ef6ux671fux671b}

\textbf{定义 47.11.1（绝对连续的测度）}：设 \(\mu, \nu\) 为 \((X, \mathcal{A})\) 上的测度。\(\nu\) 关于 \(\mu\) 绝对连续（记为 \(\nu \ll \mu\)），若对任意 \(A \in \mathcal{A}\)，\(\mu(A) = 0\) 蕴含 \(\nu(A) = 0\)。

\textbf{定理 47.11.1（Radon-Nikodym 定理）}：设 \((X, \mathcal{A}, \mu)\) 为 \(\sigma\)-有限测度空间，\(\nu\) 为 \(\sigma\)-有限的符号测度满足 \(\nu \ll \mu\)。则存在唯一的 \(\mu\)-a.e. 有限的可测函数 \(f\)（称为 Radon-Nikodym 导数，记为 \(f = \frac{d\nu}{d\mu}\)），使得对任意 \(A \in \mathcal{A}\)，

\[
\nu(A) = \int_A f \, d\mu
\]

\textbf{证明}（框架）：由 Hahn 分解，只需对正测度 \(\nu\) 证明。考虑函数类 \(\mathcal{F} = \{f \geq 0 \text{ 可测} : \int_A f d\mu \leq \nu(A), \forall A\}\)。\(\mathcal{F}\) 非空（\(0 \in \mathcal{F}\)）。取 \(f_0 \in \mathcal{F}\) 最大化 \(\int_X f d\mu\)（由极大元存在性，利用单调收敛定理）。验证 \(\int_A f_0 d\mu = \nu(A)\) 对所有 \(A\) 成立。若否，存在 \(A_0\) 使 \(\nu(A_0) - \int_{A_0} f_0 d\mu > 0\)。由 \(\sigma\)-有限性，可构造一个与 \(f_0\) 不兼容的附加函数，破坏 \(f_0\) 的极大性。\(\blacksquare\)

\textbf{注记 47.11.1}：Radon-Nikodym 定理是条件期望的数学基础。给定概率空间 \((\Omega, \mathcal{F}, \mathbb{P})\) 和子 \(\sigma\)-代数 \(\mathcal{G} \subseteq \mathcal{F}\)，对 \(X \in L^1(\mathcal{F})\)，定义 \(\nu(G) = \int_G X d\mathbb{P}\)（\(G \in \mathcal{G}\)）。则 \(\nu \ll \mathbb{P}|_{\mathcal{G}}\)，由 Radon-Nikodym 定理，存在唯一的 \(\mathcal{G}\)-可测函数 \(\mathbb{E}[X \mid \mathcal{G}]\) 满足 \(\int_G \mathbb{E}[X \mid \mathcal{G}] d\mathbb{P} = \int_G X d\mathbb{P}\)。这一定义统一了初等概率论中离散和连续情形的条件期望，并允许将条件期望推广到任意子 \(\sigma\)-代数，是现代概率论公理化体系的核心。

\newpage

\chapter{06-分析/01-实分析/04-微分与积分联系.md}\label{ux5206ux679001-ux5b9eux5206ux679004-ux5faeux5206ux4e0eux79efux5206ux8054ux7cfb.md}

\section{第156章：微分与积分的联系}\label{ux7b2c156ux7ae0ux5faeux5206ux4e0eux79efux5206ux7684ux8054ux7cfb}

微分与积分是分析学的两大基本运算，它们在 Newton-Leibniz 公式中达到统一。本章在 Lebesgue 积分的框架下精确刻画二者的关系：有界变差函数与绝对连续函数为微积分基本定理提供了最自然的函数类，Lebesgue 微分定理（配合 Vitali 覆盖引理）证明了单调函数几乎处处可导，而绝对连续函数恰好是那些能表示为自身导数的不定积分的函数。这一理论是实分析中最为优美的篇章之一，它将测度论、积分论和微分论紧密地编织在一起。

\subsection{49.1 有界变差函数}\label{ux6709ux754cux53d8ux5deeux51fdux6570}

\textbf{定义 49.1.1（有界变差函数）}：设 \(f : [a, b] \to \mathbb{R}\)。\(f\) 在 \([a, b]\) 上的\textbf{全变差}定义为

\[
V_a^b(f) = \sup\left\{\sum_{i=1}^{n} |f(x_i) - f(x_{i-1})| \;\middle|\; a = x_0 < x_1 < \cdots < x_n = b\right\}
\]

若 \(V_a^b(f) < \infty\)，则称 \(f\) 为 \([a, b]\) 上的\textbf{有界变差函数}。全体有界变差函数记为 \(BV([a, b])\)。

\textbf{定理 49.1.1（Jordan 分解）}：\(f \in BV([a, b])\) 当且仅当 \(f\) 可表示为两个单调递增函数的差。

\textbf{证明}：

（\(\Leftarrow\)）单调函数的全变差有限，两个有界变差函数的差也有界变差。

（\(\Rightarrow\)）定义 \(g(x) = V_a^x(f)\)（全变差函数），\(h(x) = g(x) - f(x)\)。\(g\) 递增（全变差随区间增长）。需证 \(h\) 也递增。对 \(x < y\)，\(h(y) - h(x) = V_x^y(f) - (f(y) - f(x)) \geq 0\)（因为全变差至少是 \(|f(y) - f(x)|\)）。故 \(f = g - h\) 是两个单调递增函数的差。\(\blacksquare\)

\textbf{推论 49.1.2}：\(f \in BV([a, b])\) 几乎处处可导。

\textbf{证明}：由 Jordan 分解，\(f\) 是两个单调函数的差。单调函数几乎处处可导（Lebesgue 微分定理，见下节），故 \(f\) 也几乎处处可导。\(\blacksquare\)

有界变差函数类在实分析中占据着核心地位，其重要性远超 Jordan 分解这一基本结果。从几何直观上看，\(V_a^b(f)\) 度量了函数图像在 \([a, b]\) 上的''总纵向振荡''------若将函数的图像视为一条路径，全变差即为该路径在 \(y\) 轴方向上的总弧长。这一几何解释使得有界变差函数自然地出现在几何测度论中：定义在 \(\mathbb{R}^n\) 上的有界变差函数（\(BV\) 函数）的分布导数是有界 Radon 测度，这构成了 Caccioppoli 集（即有限周长集）和极小曲面理论的基础。在概率论中，所有分布函数自动地是有界变差的（因为它们是单调的），而 \(BV\) 函数类进一步包含了两个分布函数的差。在泛函分析中，\(BV([a, b])\) 配以范数 \(\|f\|_{BV} = |f(a)| + V_a^b(f)\) 构成 Banach 空间，其对偶空间与 \(L^\infty\) 和 Radon 测度有密切联系。此外，有界变差函数在 Fourier 分析中也扮演着重要角色：\(BV\) 函数的 Fourier 级数在每点收敛到 \((f(x^+) + f(x^-))/2\)（Dirichlet-Jordan 定理），且收敛速度可由全变差控制。

有界变差函数的另一个重要特征是它们的''可微结构性''：\(f \in BV([a, b])\) 的导数 \(f'\) 在 Lebesgue 意义下几乎处处存在，且 \(f' \in L^1([a, b])\)。然而，与绝对连续函数不同，\(f\) 不一定能表示为导数的积分------这中间的差额正是由 \(f\) 的''奇异部分''（Cantor 函数的推广）和''跳跃部分''造成的。确切地说，\(f\) 在 Lebesgue 测度意义下的导数 \(f'\) 满足 \(\int_a^b |f'| \leq V_a^b(f)\)，且等号成立当且仅当 \(f\) 是绝对连续的。这一不等式是 Lebesgue 微分定理与 Jordan 分解结合的产物，并揭示了 \(BV\) 函数中''绝对连续部分''与''奇异部分''之间的精确关系。

\hrulefill

\subsection{49.2 Lebesgue 微分定理与 Vitali 覆盖引理}\label{lebesgue-ux5faeux5206ux5b9aux7406ux4e0e-vitali-ux8986ux76d6ux5f15ux7406}

\textbf{定义 49.2.1（Vitali 覆盖）}：设 \(E \subseteq \mathbb{R}\) 为可测集。一族闭区间 \(\mathcal{V}\) 称为 \(E\) 的 \textbf{Vitali 覆盖}，如果对每个 \(x \in E\) 和每个 \(\varepsilon > 0\)，存在 \(I \in \mathcal{V}\) 使得 \(x \in I\) 且 \(|I| < \varepsilon\)（即 \(\mathcal{V}\) 在 Vitali 意义下''细覆盖'' \(E\)）。

\textbf{定理 49.2.1（Vitali 覆盖引理）}：设 \(E \subseteq \mathbb{R}\) 为有界可测集，\(\mathcal{V}\) 为 \(E\) 的 Vitali 覆盖。则存在 \(\mathcal{V}\) 中可数个互不相交的区间 \(\{I_k\}_{k=1}^\infty\) 使得 \(m(E \setminus \bigcup_k I_k) = 0\)。此外，对任意 \(\varepsilon > 0\)，可选取有限个互不相交的区间 \(I_1, \ldots, I_N \in \mathcal{V}\) 使得 \(m(E \setminus \bigcup_{k=1}^N I_k) < \varepsilon\)。

\textbf{证明}：首先，由于 \(E\) 有界，可假设 \(\mathcal{V}\) 中所有区间均包含于某个有界开集 \(U\)（否则可截断）。设 \(m(U) < \infty\)。构造区间列的步骤如下：选取 \(I_1 \in \mathcal{V}\)（任意）。假设已选出互不相交的 \(I_1, \ldots, I_k\)。若 \(E \subseteq \bigcup_{j=1}^k I_j\)，则已满足要求。否则，令 \(r_k = \sup\{|I| : I \in \mathcal{V}, I \cap \bigcup_{j=1}^k I_j = \varnothing\}\)。选取 \(I_{k+1} \in \mathcal{V}\) 与之前的 \(I_j\) 均不相交，且 \(|I_{k+1}| > r_k/2\)。此为贪心构造。由于 \(\sum_k |I_k| \leq m(U) < \infty\)，故 \(|I_k| \to 0\)。现证 \(m(E \setminus \bigcup_k I_k) = 0\)：对每个 \(I_k\)，令 \(J_k\) 为与 \(I_k\) 同中心但长度 \(5\) 倍的区间。对任意 \(x \in E \setminus \bigcup_k I_k\)，对每个 \(N\)，存在 \(I \in \mathcal{V}\) 包含 \(x\) 且与 \(I_1, \ldots, I_N\) 不相交。设 \(k\) 为第一个使 \(I \cap I_k \neq \varnothing\) 的指标（\(k > N\)），则 \(|I| \leq r_{k-1} < 2|I_k|\)，故 \(I \subseteq J_k\)。因此 \(E \setminus \bigcup_{k=1}^N I_k \subseteq \bigcup_{k=N+1}^\infty J_k\)。取 \(N\) 充分大，\(\sum_{k=N+1}^\infty |J_k| = 5 \sum_{k=N+1}^\infty |I_k|\) 可任意小。第二结论由测度的近似性质得出。\(\blacksquare\)

Vitali 覆盖引理是实分析中最强大的技术工具之一。它在 \(\mathbb{R}^n\) 中的推广（以球代替区间）同样成立，且是证明 Lebesgue 微分定理、Lebesgue 密度定理（可测集几乎每点都是 Lebesgue 点，即密度为 \(1\)）以及极大函数弱型估计的核心。Vitali 覆盖引理的另一重要变体是 \textbf{Besicovitch 覆盖引理}，后者在 \(\mathbb{R}^n\) 中允许重叠但重叠次数有界，在证明 \(BV\) 函数的微分结构和测度的微分理论中不可或缺。

\textbf{定义 49.2.2（Hardy-Littlewood 极大函数）}：对 \(f \in L^1_{\text{loc}}(\mathbb{R})\)，其 \textbf{Hardy-Littlewood 极大函数} 定义为

\[
Mf(x) = \sup_{r > 0} \frac{1}{2r} \int_{x-r}^{x+r} |f(t)| \, dt
\]

在 \(\mathbb{R}^n\) 中，\(Mf(x) = \sup_{r > 0} \frac{1}{|B(x,r)|} \int_{B(x,r)} |f(y)| \, dy\)。

\textbf{定理 49.2.2（Hardy-Littlewood 极大不等式）}：设 \(f \in L^1(\mathbb{R}^n)\)。则 \(Mf\) 满足\textbf{弱 \((1,1)\) 型不等式}：

\[
m(\{x \in \mathbb{R}^n : Mf(x) > \lambda\}) \leq \frac{C_n}{\lambda} \|f\|_{L^1}
\]

其中 \(C_n = 3^n\)（最优常数 \(C_n = 5^n\) 也成立）。此外，对 \(1 < p \leq \infty\)，\(M\) 是强 \((p,p)\) 型算子：\(\|Mf\|_{L^p} \leq C_p \|f\|_{L^p}\)。

\textbf{证明（弱 \((1,1)\) 型）}：设 \(E_\lambda = \{Mf > \lambda\}\)。对每个 \(x \in E_\lambda\)，存在球 \(B_x\) 使得 \(\frac{1}{|B_x|} \int_{B_x} |f| > \lambda\)。由 Vitali 覆盖引理（或 Wiener 引理），存在互不相交的球 \(B_{x_k}\) 使得 \(\sum_k |B_{x_k}| \geq c_n m(E_\lambda)\)。于是

\[
m(E_\lambda) \leq \frac{1}{c_n} \sum_k |B_{x_k}| \leq \frac{1}{c_n \lambda} \sum_k \int_{B_{x_k}} |f| \leq \frac{1}{c_n \lambda} \|f\|_{L^1}
\]

取 \(c_n = 1/3^n\)（或 \(1/5^n\)）即得。强 \((p,p)\) 型由 Marcinkiewicz 插值定理得出。\(\blacksquare\)

\textbf{定理 49.2.3（Lebesgue 微分定理）}：设 \(f \in L^1_{\text{loc}}(\mathbb{R}^n)\)。则对几乎处处 \(x \in \mathbb{R}^n\)，

\[
\lim_{r \to 0} \frac{1}{|B(x,r)|} \int_{B(x,r)} f(y) \, dy = f(x)
\]

满足此极限的点称为 \(f\) 的 \textbf{Lebesgue 点}。特别地，若 \(f : [a, b] \to \mathbb{R}\) 单调递增，则 \(f\) 在 \([a, b]\) 上几乎处处可导，且 \(f' \in L^1([a, b])\)，满足

\[
\int_a^b f'(x) \, dx \leq f(b) - f(a)
\]

\textbf{证明（框架）}：对 \(f \in L^1\)，定义 \(T_r f(x) = \frac{1}{|B(x,r)|} \int_{B(x,r)} f(y) dy\)。需证 \(\lim_{r \to 0} T_r f(x) = f(x)\) a.e.。设 \(\Omega f(x) = \limsup_{r \to 0} |T_r f(x) - f(x)|\)。对任意 \(\varepsilon > 0\)，取 \(g \in C_c(\mathbb{R}^n)\) 使 \(\|f - g\|_{L^1} < \varepsilon\)。由 \(g\) 连续，\(\Omega g(x) = 0\) 对所有 \(x\)。而 \(\Omega f(x) \leq M(f - g)(x) + |f(x) - g(x)|\)。由弱 \((1,1)\) 型和 Chebyshev 不等式，\(m(\{M(f - g) > \lambda\}) \leq \frac{C}{\lambda} \varepsilon\)。取 \(\varepsilon\) 适当小，可证 \(\Omega f(x) = 0\) a.e.。单调函数情形：\(f\) 单调递增，其差商 \(D_h f(x) = \frac{f(x+h) - f(x)}{h}\) 构成 \(L^1\) 函数族，由 Lebesgue 微分定理的积分形式可得可微性。\(\blacksquare\)

Lebesgue 微分定理的深层意义在于它确立了 \(L^1_{\text{loc}}\) 函数可以通过其积分平均在几乎每点处被''恢复''。这一结果不仅是稠密性论证的基础，也是证明 Lebesgue 点存在性的关键------在 Lebesgue 点处，函数值可以由其积分平均恢复，这一性质在 Fourier 级数的逐点收敛性（如 Carleson-Hunt 定理）和奇异积分理论中至关重要。在更一般的框架下，Lebesgue 微分定理可推广到齐型空间（spaces of homogeneous type）上的 Borel 测度，以及更一般的微分基（differentiation basis）------Busemann-Feller 微分定理给出了微分基满足 Vitali 型覆盖性质时 Lebesgue 微分定理成立的充要条件。

\hrulefill

\subsection{49.3 绝对连续函数}\label{ux7eddux5bf9ux8fdeux7eedux51fdux6570}

\textbf{定义 49.3.1（绝对连续函数）}：函数 \(f : [a, b] \to \mathbb{R}\) 称为\textbf{绝对连续的}，如果对任意 \(\varepsilon > 0\)，存在 \(\delta > 0\)，使得对任意有限个互不相交的子区间 \((x_i, y_i) \subseteq [a, b]\)，只要 \(\sum (y_i - x_i) < \delta\)，就有 \(\sum |f(y_i) - f(x_i)| < \varepsilon\)。全体绝对连续函数记为 \(AC([a, b])\)。

\textbf{定理 49.3.1（绝对连续函数的基本性质）}： (i) \(AC([a, b]) \subsetneq BV([a, b]) \cap C([a, b])\)； (ii) 绝对连续函数几乎处处可导，且 \(f' \in L^1([a, b])\)； (iii) \(f\) 绝对连续当且仅当 \(f\) 是 \(L^1\) 函数的不定积分：\(f(x) = f(a) + \int_a^x g(t) dt\)（\(g \in L^1\)）； (iv) 绝对连续函数将零测度集映为零测度集（Lusin 性质）。

\textbf{证明}：(iii) ``\(\Leftarrow\)''：由积分的绝对连续性即得。``\(\Rightarrow\)''：\(f \in AC\) 则 \(f \in BV\)，故 \(f'\) 几乎处处存在且 \(f' \in L^1\)。定义 \(g(x) = f(a) + \int_a^x f'(t) dt\)，则 \(g \in AC\) 且 \(g' = f'\) a.e.，故 \((f - g)' = 0\) a.e.。由绝对连续性，\(f - g\) 为常数，且 \(f(a) - g(a) = 0\)，故 \(f = g\)。\(\blacksquare\)

绝对连续函数类恰好是微积分基本定理在 Lebesgue 积分框架下成立的函数类。与 Riemann 积分框架下的可微函数类不同，绝对连续函数允许导数在 Lebesgue 意义下可积，且由导数的积分可以完全恢复原函数。这一结果与经典的 Cantor 函数（Cantor 阶梯函数）形成鲜明对比：Cantor 函数在 \([0,1]\) 上连续、几乎处处可导且导数为 \(0\)，但它不是绝对连续的------它从 \(0\) 增长到 \(1\) 而导数几乎处处为 \(0\)，因此无法表示为导数的积分。Cantor 函数的存在性表明，在''连续''和''绝对连续''之间存在着本质的鸿沟，而这一鸿沟只能用测度论来刻画：绝对连续函数对应于关于 Lebesgue 测度绝对连续的测度的分布函数，而 Cantor 函数对应于奇异测度的分布函数。

从测度论的角度看，绝对连续函数与测度的绝对连续性之间的对应关系是理解这一概念的关键。若 \(\mu\) 是 \([a, b]\) 上的有限 Borel 测度，其分布函数 \(F(x) = \mu([a, x])\) 是单调递增的。\(F\) 绝对连续当且仅当 \(\mu \ll m\)（关于 Lebesgue 测度绝对连续），此时 \(F'(x) = \frac{d\mu}{dm}(x)\) a.e.。这一对应将函数论中的绝对连续性与测度论中的 Radon-Nikodym 定理紧密联系起来。在 \(\mathbb{R}^n\) 中，绝对连续函数的概念推广为 Sobolev 空间 \(W^{1,1}\) 中的函数------\(f \in W^{1,1}(\mathbb{R}^n)\) 当且仅当 \(f\) 在几乎所有平行于坐标轴的直线上绝对连续且其分布导数属于 \(L^1\)。

绝对连续函数在变分法和最优控制中也有重要应用。变分法的基本引理（du Bois-Reymond 引理）本质上依赖于绝对连续函数的性质：若 \(\int_a^b f(t) \varphi'(t) dt = 0\) 对所有 \(\varphi \in C_c^\infty(a, b)\) 成立，则 \(f\) 几乎处处等于常数。这一结果在证明 Euler-Lagrange 方程时是不可或缺的。在最优控制中，状态轨迹的绝对连续性确保了对控制输入的可积性假设足以保证轨迹的良定性。

\hrulefill

\subsection{49.4 Lebesgue 微积分基本定理}\label{lebesgue-ux5faeux79efux5206ux57faux672cux5b9aux7406}

\textbf{定理 49.4.1（Lebesgue 微积分基本定理）}：设 \(f : [a, b] \to \mathbb{R}\)。 (i) 若 \(f \in L^1([a, b])\)，定义 \(F(x) = \int_a^x f(t) dt\)，则 \(F \in AC([a, b])\) 且 \(F'(x) = f(x)\) 几乎处处成立； (ii) 若 \(F \in AC([a, b])\)，则 \(F'(x)\) 几乎处处存在，\(F' \in L^1([a, b])\)，且 \(F(x) = F(a) + \int_a^x F'(t) dt\) 对所有 \(x \in [a, b]\) 成立。

\textbf{证明}：(i) 由积分的绝对连续性，\(F\) 绝对连续。对 \(f \geq 0\)（否则分解为正负部分），\(F\) 单调递增，由 Lebesgue 微分定理，\(F'(x) = \lim_{h \to 0} \frac{1}{h} \int_x^{x+h} f(t) dt = f(x)\) a.e.（后一等号由 Lebesgue 微分定理的积分形式保证）。(ii) 由绝对连续函数的性质（定理 49.3.1(iii)）直接得出。\(\blacksquare\)

Lebesgue 微积分基本定理彻底解决了 Newton-Leibniz 公式何时成立这一核心问题。在 Riemann 积分框架下，\(F'(x) = f(x)\) 且 \(f\) Riemann 可积是充分条件，但这一条件相当严格------存在几乎处处可导且导数有界但导数不是 Riemann 可积的函数。Lebesgue 积分框架通过引入绝对连续函数，给出了 Newton-Leibniz 公式成立的充要条件：原函数必须绝对连续。这一结果与 Lebesgue 分解定理（定理 49.5.1）紧密相连，将测度分解为绝对连续部分和奇异部分，从而将函数分解为相应的两部分。

该定理的历史意义也十分深远。Newton 和 Leibniz 在 17 世纪提出的微积分基本定理，在随后两个世纪中一直缺乏严格的数学基础。Cauchy 在 19 世纪初为连续函数建立了积分理论，但只有在 Lebesgue 在 1902 年引入其测度和积分理论后，微积分基本定理才获得了完全精确的表述。值得注意的是，Lebesgue 在提出其理论时明确将''恢复原函数''作为积分理论的基本目标之一------这一动机与他研究 Fourier 级数的逐项积分问题密切相关。此外，该定理在 \(\mathbb{R}^n\) 中的推广需要更精细的工具：对 \(f \in L^1_{\text{loc}}(\mathbb{R}^n)\)，其不定积分 \(F(x) = \int_{[a_1, x_1] \times \cdots \times [a_n, x_n]} f(t) dt\) 不一定几乎处处可导，且 \(\frac{\partial^n F}{\partial x_1 \cdots \partial x_n} = f\) 也仅在 \(f\) 是''强可积''时成立。这导致了 Sobolev 空间和 \(BV\) 函数理论的发展。

\hrulefill

\subsection{49.5 测度的微分与 Lebesgue 分解}\label{ux6d4bux5ea6ux7684ux5faeux5206ux4e0e-lebesgue-ux5206ux89e3}

\textbf{定理 49.5.1（Lebesgue-Radon-Nikodym 分解定理）}：设 \(\mu\) 和 \(\nu\) 为 \(\mathbb{R}\) 上的 \(\sigma\)-有限测度。则 \(\nu\) 可唯一地分解为 \(\nu = \nu_a + \nu_s\)，其中 \(\nu_a \ll \mu\)（\(\nu_a\) 关于 \(\mu\) 绝对连续），\(\nu_s \perp \mu\)（\(\nu_s\) 与 \(\mu\) 相互奇异）。且存在 \(f \in L^1(\mu)\)，使得 \(d\nu_a = f d\mu\)------\(f\) 称为 \(\nu\) 关于 \(\mu\) 的 Radon-Nikodym 导数，记为 \(f = \frac{d\nu}{d\mu}\)。

\textbf{证明}：考虑有限测度情形（一般情形由 \(\sigma\)-有限性归结）。设 \(\lambda = \mu + \nu\)。对 \(\lambda\) 可积函数 \(g\)，考虑线性泛函 \(\ell(g) = \int g d\nu\)。由 Riesz 表示定理，存在 \(h \in L^\infty(\lambda)\) 使 \(\ell(g) = \int g h d\lambda\)。由于 \(0 \leq \nu \leq \lambda\)，\(0 \leq h \leq 1\) a.e.。令 \(A = \{h = 1\}\)，则 \(\nu_a(E) = \nu(E \setminus A)\)，\(\nu_s(E) = \nu(E \cap A)\) 给出所需分解。Radon-Nikodym 导数由 \(f = \frac{h}{1-h} \chi_{A^c}\) 给出。\(\blacksquare\)

Radon-Nikodym 定理是测度论中除了 Lebesgue 积分构造之外最重要的定理。其核心思想是：若 \(\nu\) 关于 \(\mu\) 绝对连续（即 \(\mu(E) = 0 \Rightarrow \nu(E) = 0\)），则 \(\nu\) 具有关于 \(\mu\) 的''密度函数'' \(f\)。这一结果在概率论中至关重要------条件期望的存在性本质上依赖于 Radon-Nikodym 导数：设 \((\Omega, \mathcal{F}, P)\) 为概率空间，\(\mathcal{G} \subseteq \mathcal{F}\) 为子 \(\sigma\)-代数，则对 \(X \in L^1(\Omega, \mathcal{F}, P)\)，条件期望 \(\mathbb{E}[X | \mathcal{G}]\) 是 \(X\) 关于限制测度 \(P|_{\mathcal{G}}\) 的 Radon-Nikodym 导数。这一解释将条件期望这一概率论核心概念建立在坚实的测度论基础上，并由 Kolmogorov 在 1933 年公理化概率论中首次明确给出。

Radon-Nikodym 定理在统计推断和机器学习中也有广泛应用。在贝叶斯统计中，后验分布关于先验分布的 Radon-Nikodym 导数正是似然函数（归一化后），这一对应关系是贝叶斯更新公式的测度论表述。在信息几何中，Fisher 信息矩阵可由统计模型族中概率测度之间的 Radon-Nikodym 导数的 \(L^2\) 距离的局部展开得到。

\textbf{定理 49.5.2（Lebesgue 分解的逐点形式------\(BV\) 函数的分解）}：对任意 \(f \in BV([a, b])\)，\(f\) 可唯一分解为

\[
f = f_a + f_s + f_j
\]

其中 \(f_a \in AC([a, b])\) 为绝对连续部分，\(f_s\) 为奇异连续部分（\(f_s\) 连续，\(f_s' = 0\) a.e.，但 \(f_s\) 不为常数），\(f_j\) 为跳跃函数（阶梯函数，其不连续点集至多可数）。这三部分分别对应 \(f\) 诱导的 Lebesgue-Stieltjes 测度关于 Lebesgue 测度的绝对连续部分、奇异连续部分和纯点部分。

\textbf{证明}：\(f \in BV\) 诱导一个有限符号测度 \(\mu_f\)（由 \(\mu_f((c,d]) = f(d) - f(c)\) 唯一确定）。由 Lebesgue-Radon-Nikodym 分解定理，\(\mu_f = \mu_a + \mu_s\)，其中 \(\mu_a \ll m\)，\(\mu_s \perp m\)。进一步将 \(\mu_s\) 分解为纯点部分 \(\mu_p\)（原子的和）和奇异连续部分 \(\mu_{sc}\)（无原子但与 \(m\) 奇异）。令 \(f_a(x) = f(a) + \mu_a([a, x])\)，\(f_s(x) = \mu_{sc}([a, x])\)，\(f_j(x) = \mu_p([a, x])\)。\(\mu_a \ll m\) 保证 \(f_a\) 绝对连续；\(\mu_{sc}\) 无原子保证 \(f_s\) 连续但导数为 \(0\) a.e.；\(\mu_p\) 为原子之和保证 \(f_j\) 为跳跃函数。\(\blacksquare\)

这一分解定理完美地刻画了有界变差函数的微观结构，是现代实分析中微分与积分理论的高峰。Cantor 函数是 \(f_s\) 的经典例子（\(f_a = f_j = 0\)，\(f = f_s\)），而 Heaviside 阶梯函数是 \(f_j\) 的经典例子。在一般情形，一个 \(BV\) 函数可能同时包含这三种成分------例如，一个在 Cantor 集上奇异、在有理点处跳跃、在其余部分绝对连续的函数。在实际应用中，\(f_a\) 部分对应于信号中的''平滑''成分（可积导数），\(f_j\) 部分对应于突变（边缘检测），\(f_s\) 部分对应于分形结构------这一分解在图像处理和信号分析中具有直接的应用价值。

\hrulefill

\subsection{49.6 微分与积分的深层联系}\label{ux5faeux5206ux4e0eux79efux5206ux7684ux6df1ux5c42ux8054ux7cfb}

\textbf{定理 49.6.1（微分与积分的对偶性）}：设 \(f \in L^1([a, b])\)。则 \(f\) 的不定积分 \(F(x) = \int_a^x f(t) dt\) 的导数 \(F'\) 在几乎处处意义下恢复 \(f\)。反之，若 \(F \in AC([a, b])\)，则 \(F'\) 的积分恢复 \(F\)。这一对偶性在 \(L^1\) 的框架下是完美的，但在更一般的函数空间中产生了深远的推广。

在 \(\mathbb{R}^n\) 中，微分与积分的联系表现为分布导数和 Sobolev 空间理论。\(\mathbb{R}^n\) 上的函数 \(f\) 的分布导数 \(\partial^\alpha f\) 定义为满足 \(\int f \partial^\alpha \varphi = (-1)^{|\alpha|} \int \partial^\alpha f \cdot \varphi\)（对所有 \(\varphi \in C_c^\infty\)）的局部可积函数。当 \(\partial^\alpha f \in L^p\) 时，\(f\) 属于 Sobolev 空间 \(W^{1,p}\)------这正是绝对连续函数的多维推广。\textbf{Sobolev 嵌入定理}断言 \(W^{1,p}(\mathbb{R}^n) \hookrightarrow L^{p^*}(\mathbb{R}^n)\)（\(p^* = np/(n-p)\)，\(p < n\)），将微分性（导数的 \(L^p\) 可积性）转化为积分性（函数本身的 \(L^{p^*}\) 可积性提升），这是微分与积分联系的又一深刻体现。

在更抽象的层面，微分与积分的联系可视为\textbf{伴随算子}的关系。在 \(L^2\) 中，微分算子 \(\frac{d}{dx}\)（定义在绝对连续函数上）与积分算子 \(\int_a^x\) 互为''逆''（在适当的定义域和值域下）。这一观点在泛函分析中发展为无界算子理论：\(\frac{d}{dx}\) 是 \(L^2\) 上的稠定无界算子，其谱包含整个虚轴，而 \(\int_a^x\) 是紧算子（在 \(L^2([a,b])\) 上）。紧积分算子的奇异值衰减速率由被积函数的正则性决定------这是积分方程理论（Fredholm 理论）的出发点，也是现代 PDE 理论中紧嵌入定理（如 Rellich-Kondrachov 定理）的雏形。

微分与积分的联系在调和分析中表现为\textbf{奇异积分算子}与\textbf{分数次积分算子}（Riesz 位势 \(I_\alpha\)）的关系。\(I_\alpha\) 是 \((-\Delta)^{-\alpha/2}\) 的卷积核表示，它将 \(f\) 的正则性提升 \(\alpha\) 阶（在 Sobolev 尺度上）。\(I_\alpha\) 的逆------\((-\Delta)^{\alpha/2}\)------正是分数次微分算子，而 \(I_\alpha\) 本身则是分数次积分算子。这一理论将经典的 Newton-Leibniz 公式（\(\alpha = 1\)）推广到分数阶，导致了分数阶微积分和反常扩散方程的现代理论。

在概率论中，Lebesgue 微分定理的自然推广是 \textbf{Martingale 收敛定理}：设 \(\{\mathcal{F}_n\}\) 为 \(\sigma\)-代数流，\(\{X_n\}\) 为鞅。则 \(X_n\) 几乎处处收敛的充要条件是 \(\{X_n\}\) 在 \(L^1\) 中有界。这一结果与 Lebesgue 微分定理共享相同的测度论内核------两者都依赖于 Vitali 型覆盖和极大不等式。事实上，鞅的极大不等式 \(\lambda P(\sup_n |X_n| > \lambda) \leq \sup_n \mathbb{E}|X_n|\)（Doob 不等式）与 Hardy-Littlewood 极大不等式的证明结构几乎完全相同。这一联系揭示了概率论与实分析之间的深层统一性。

\(\blacksquare\)

\hrulefill

\hrulefill

\hrulefill

\hrulefill

\hrulefill

\hrulefill

\section{\texorpdfstring{第157章：\(L^p\) 空间}{第157章：L\^{}p 空间}}\label{ux7b2c157ux7ae0lp-ux7a7aux95f4}

\(L^p\) 空间是 Lebesgue 积分理论的自然产物，构成了现代分析学的基本框架。它们不仅是 Banach 空间理论中最经典的例子，也是调和分析、偏微分方程和概率论的核心工具。\(L^p\) 空间的理论由 F. Riesz 在 1910 年系统建立，其完备性（Riesz-Fischer 定理）标志着泛函分析作为一个独立学科的诞生。

\subsection{\texorpdfstring{48.1 \(L^p\) 空间的定义与基本结构}{48.1 L\^{}p 空间的定义与基本结构}}\label{lp-ux7a7aux95f4ux7684ux5b9aux4e49ux4e0eux57faux672cux7ed3ux6784}

\textbf{定义 48.1.1（ \(L^p\) 范数）}：设 \(E \subseteq \mathbb{R}^n\) 可测，\(1 \leq p < \infty\)。\(f\) 的 \(L^p\) \textbf{范数}定义为

\[
\|f\|_p = \left(\int_E |f|^p \, dm\right)^{1/p}
\]

\(L^p(E)\) 是所有满足 \(\|f\|_p < \infty\) 的可测函数 \(f\) 的集合（模去 a.e. 相等的函数，即 \(f = g\) a.e. 视为同一元素）。

\textbf{定义 48.1.2（ \(L^\infty\) 范数）}：\(f\) 的 \(L^\infty\) \textbf{范数}（本质上确界）定义为

\[
\|f\|_\infty = \inf\{M \geq 0 \mid |f(x)| \leq M \text{ a.e.}\}
\]

\(L^\infty(E)\) 是所有满足 \(\|f\|_\infty < \infty\) 的可测函数的集合。

\textbf{注记 48.1.1}：\(L^p\) 空间的元素是函数的等价类，而非函数本身。这一''模去零测集''的等价关系是积分理论不可避免的------两个几乎处处相等的函数在积分意义下无法区分，且只有在这种等价关系下 \(\|\cdot\|_p\) 才满足范数公理的正定性（\(\|f\|_p = 0 \iff f = 0\) a.e.）。这一等价关系意味着 \(L^p\) 空间中的逐点收敛概念没有意义，代之以几乎处处收敛和依范数收敛。

\textbf{定理 48.1.1}：\(L^p(E)\) 是向量空间（\(1 \leq p \leq \infty\)）。

\textbf{证明}：由 \(|f + g|^p \leq 2^p(|f|^p + |g|^p)\)（\(p < \infty\)）或 \(\|f + g\|_\infty \leq \|f\|_\infty + \|g\|_\infty\)（\(p = \infty\)）。\(\blacksquare\)

\textbf{命题 48.1.2（ \(L^p\) 空间的包含关系）}：设 \(E\) 为有限测度集。则对 \(1 \leq p < q \leq \infty\)，有 \(L^q(E) \subseteq L^p(E)\)，且嵌入是连续的。具体地，\(\|f\|_p \leq m(E)^{1/p - 1/q} \|f\|_q\)。

\textbf{证明}：由 Holder 不等式，\(\|f\|_p^p = \int_E |f|^p \cdot 1 \leq (\int_E |f|^{p \cdot q/p})^{p/q} \cdot (\int_E 1^{r})^{1/r} = \|f\|_q^p \cdot m(E)^{1 - p/q}\)，其中 \(r = \frac{q}{q-p}\) 是 \(q/p\) 的共轭指数。\(\blacksquare\)

\textbf{注记 48.1.2}：在无限测度集（如 \(\mathbb{R}^n\)）上，\(L^p\) 空间之间没有包含关系。例如，\(f(x) = x^{-1/2}\chi_{(0,1)}(x)\) 属于 \(L^1(\mathbb{R})\) 但不属于 \(L^2(\mathbb{R})\)；而 \(g(x) = x^{-1}\chi_{(1,\infty)}(x)\) 属于 \(L^2(\mathbb{R})\) 但不属于 \(L^1(\mathbb{R})\)。这意味着 \(L^p\) 空间在 \(\mathbb{R}^n\) 上形成了一个真正的''尺度''------不同 \(p\) 值捕获了函数在不同区域（局部 vs.~远程）的不同衰减行为。

\hrulefill

\subsection{48.2 Hölder 不等式与 Minkowski 不等式}\label{huxf6lder-ux4e0dux7b49ux5f0fux4e0e-minkowski-ux4e0dux7b49ux5f0f}

\textbf{定理 48.2.1（Young 不等式）}：设 \(a, b \geq 0\)，\(p, q > 1\) 满足 \(\frac{1}{p} + \frac{1}{q} = 1\)。则

\[
ab \leq \frac{a^p}{p} + \frac{b^q}{q}
\]

\textbf{证明}：由凸函数 \(t \mapsto e^t\) 的性质，\(\frac{a^p}{p} + \frac{b^q}{q} = e^{\ln a^p - \ln p} + e^{\ln b^q - \ln q} \geq e^{\frac{1}{p}\ln a^p + \frac{1}{q}\ln b^q} = ab\)。等号成立当且仅当 \(a^p = b^q\)。\(\blacksquare\)

\textbf{定理 48.2.2（Hölder 不等式）}：设 \(1 < p, q < \infty\)，\(\frac{1}{p} + \frac{1}{q} = 1\)。若 \(f \in L^p\)，\(g \in L^q\)，则 \(fg \in L^1\) 且

\[
\|fg\|_1 \leq \|f\|_p \|g\|_q
\]

\textbf{证明}：若 \(\|f\|_p = 0\) 或 \(\|g\|_q = 0\)，不等式平凡成立。否则，令 \(a = |f(x)|/\|f\|_p\)，\(b = |g(x)|/\|g\|_q\)。由 Young 不等式，

\[
\frac{|f(x)g(x)|}{\|f\|_p \|g\|_q} \leq \frac{1}{p} \frac{|f(x)|^p}{\|f\|_p^p} + \frac{1}{q} \frac{|g(x)|^q}{\|g\|_q^q}
\]

两边积分得 \(\frac{\|fg\|_1}{\|f\|_p \|g\|_q} \leq \frac{1}{p} + \frac{1}{q} = 1\)。\(\blacksquare\)

\textbf{定理 48.2.3（广义 Hölder 不等式）}：设 \(p_1, \ldots, p_k > 1\) 满足 \(\sum_{j=1}^{k} \frac{1}{p_j} = 1\)。若 \(f_j \in L^{p_j}\)，则 \(\prod f_j \in L^1\) 且

\[
\left\|\prod_{j=1}^{k} f_j\right\|_1 \leq \prod_{j=1}^{k} \|f_j\|_{p_j}
\]

\textbf{证明}：对 \(k\) 归纳，将 \(k\) 个函数的乘积分组应用二元 Hölder 不等式。\(\blacksquare\)

\textbf{定理 48.2.4（Minkowski 不等式）}：设 \(1 \leq p < \infty\)，\(f, g \in L^p\)。则

\[
\|f + g\|_p \leq \|f\|_p + \|g\|_p
\]

\textbf{证明}：\(p = 1\) 时由三角不等式直接得到。\(p > 1\) 时，

\[
\begin{aligned}
\|f + g\|_p^p &= \int |f + g|^p \leq \int |f + g|^{p-1}|f| + \int |f + g|^{p-1}|g| \\
&\leq \left(\int |f + g|^{(p-1)q}\right)^{1/q} \left(\int |f|^p\right)^{1/p} + \left(\int |f + g|^{(p-1)q}\right)^{1/q} \left(\int |g|^p\right)^{1/p} \\
&= \|f + g\|_p^{p/q} (\|f\|_p + \|g\|_p)
\end{aligned}
\]

其中第二步用了 Hölder 不等式（\(\frac{1}{p} + \frac{1}{q} = 1\)，\((p-1)q = p\)）。两边除以 \(\|f + g\|_p^{p/q}\)（若为 \(0\) 则平凡），得 \(\|f + g\|_p \leq \|f\|_p + \|g\|_p\)。\(\blacksquare\)

\textbf{定理 48.2.5（Minkowski 积分不等式）}：设 \((X, \mu)\) 和 \((Y, \nu)\) 为 \(\sigma\)-有限的测度空间，\(f: X \times Y \to \mathbb{R}\) 可测。则对 \(1 \leq p < \infty\)，

\[
\left(\int_X \left|\int_Y f(x, y) \, d\nu(y)\right|^p d\mu(x)\right)^{1/p} \leq \int_Y \left(\int_X |f(x, y)|^p d\mu(x)\right)^{1/p} d\nu(y)
\]

即 \(L^p\) 范数下的积分与范数次序可交换（类似连续形式的 Minkowski 不等式）。

\textbf{证明}：由 Fubini 定理和 Hölder 不等式，对偶性论证。\(\blacksquare\)

\textbf{推论 48.2.6}：\(\|\cdot\|_p\) 是 \(L^p(E)\) 上的范数。

\hrulefill

\subsection{\texorpdfstring{48.3 \(L^p\) 空间的完备性（Riesz-Fischer 定理）}{48.3 L\^{}p 空间的完备性（Riesz-Fischer 定理）}}\label{lp-ux7a7aux95f4ux7684ux5b8cux5907ux6027riesz-fischer-ux5b9aux7406}

\textbf{定理 48.3.1（Riesz-Fischer 定理）}：\(L^p(E)\)（\(1 \leq p \leq \infty\)）关于 \(\|\cdot\|_p\) 范数是完备的赋范向量空间（即 Banach 空间）。

\textbf{证明}（\(1 \leq p < \infty\) 的情形）：设 \(\{f_k\}\) 为 \(L^p\) 中的 Cauchy 列。选取子列 \(\{f_{k_j}\}\) 使得 \(\|f_{k_{j+1}} - f_{k_j}\|_p < 2^{-j}\)。

定义 \(g_n = \sum_{j=1}^{n} |f_{k_{j+1}} - f_{k_j}|\)，\(g = \sum_{j=1}^{\infty} |f_{k_{j+1}} - f_{k_j}|\)。由 Minkowski 不等式，

\[
\|g_n\|_p \leq \sum_{j=1}^{n} \|f_{k_{j+1}} - f_{k_j}\|_p < \sum_{j=1}^{n} 2^{-j} < 1
\]

由 Fatou 引理，\(\|g\|_p^p = \int \liminf g_n^p \leq \liminf \int g_n^p \leq 1\)，故 \(g \in L^p\)。特别地，\(g(x) < \infty\) a.e.。

因此 \(\sum (f_{k_{j+1}} - f_{k_j})\) a.e. 绝对收敛，故 \(f_{k_j}\) a.e. 收敛到某个 \(f\)。定义 \(f(x) = \lim_{j \to \infty} f_{k_j}(x)\)（a.e.），其余处为 \(0\)。

由 Fatou 引理，\(\|f - f_k\|_p^p = \int \liminf_{j \to \infty} |f_{k_j} - f_k|^p \leq \liminf_{j \to \infty} \|f_{k_j} - f_k\|_p^p\)。由于 \(\{f_k\}\) 是 Cauchy 列，右端对足够大的 \(k\) 可任意小，故 \(f_k \to f\) 在 \(L^p\) 中。\(\blacksquare\)

\textbf{注记 48.3.1}：上述证明揭示了 \(L^p\) 完备性的一个关键推论：\(L^p\) 中依范数收敛的函数列必有几乎处处收敛的子列。这一性质在分析中极为有用------它允许我们在依范数收敛和几乎处处收敛之间建立桥梁。Riesz-Fischer 定理是泛函分析中最早的大定理之一，它与 \(L^2\) 中规范正交基的存在性（Hilbert 空间理论）密切相关。

\hrulefill

\subsection{48.4 稠密性定理}\label{ux7a20ux5bc6ux6027ux5b9aux7406}

\textbf{定理 48.4.1（连续函数在 \(L^p\) 中的稠密性）}：设 \(1 \leq p < \infty\)，\(E \subseteq \mathbb{R}^n\) 可测且 \(m(E) < \infty\)。则 \(C_c(E)\)（紧支集连续函数）在 \(L^p(E)\) 中稠密。

\textbf{证明}：

\begin{enumerate}
\def\labelenumi{\arabic{enumi}.}
\item
  特征函数逼近：设 \(A \subseteq E\) 可测且 \(m(A) < \infty\)。由 Lusin 定理，对任意 \(\varepsilon > 0\)，存在紧集 \(K \subseteq A\) 和连续函数 \(g\) 使得 \(m(A \setminus K) < \varepsilon^p/2\) 且 \(g|_K = 1\)，\(0 \leq g \leq 1\)，\(\operatorname{supp} g \subseteq E\) 为紧集。则 \(\|\chi_A - g\|_p^p = \int_{A \setminus K} |1-g|^p + \int_{K} |1-1|^p \leq m(A \setminus K) < \varepsilon^p/2\)。
\item
  简单函数逼近：简单函数 \(s = \sum a_j \chi_{A_j}\) 是特征函数的线性组合。由三角不等式，\(s\) 可由 \(C_c(E)\) 中函数逼近。
\item
  一般 \(L^p\) 函数逼近：任意 \(f \in L^p(E)\) 可由简单函数在 \(L^p\) 范数下逼近（简单函数逼近定理）。对任意 \(\varepsilon > 0\)，取简单函数 \(s\) 使 \(\|f - s\|_p < \varepsilon/2\)，再取 \(g \in C_c(E)\) 使 \(\|s - g\|_p < \varepsilon/2\)，则 \(\|f - g\|_p < \varepsilon\)。\(\blacksquare\)
\end{enumerate}

\textbf{定理 48.4.2（光滑函数在 \(L^p\) 中的稠密性）}：\(C_c^\infty(\mathbb{R}^n)\) 在 \(L^p(\mathbb{R}^n)\) 中稠密（\(1 \leq p < \infty\)）。这可通过磨光（mollification）技术证明：取 \(\eta_\varepsilon\) 为光滑逼近恒等核，则 \(f * \eta_\varepsilon \in C^\infty\)，且 \(\|f * \eta_\varepsilon - f\|_p \to 0\)（\(\varepsilon \to 0\)）。

\textbf{定理 48.4.3（ \(L^p\) 空间的可分性）}：对 \(1 \leq p < \infty\)，\(L^p(\mathbb{R}^n)\) 是可分的 Banach 空间。\(L^\infty(\mathbb{R}^n)\) 不可分。

\textbf{证明}：有理端点的区间（或矩形）的特征函数的有理线性组合在 \(L^p\) 中稠密，构成可数稠密子集。\(L^\infty\) 的不可分性由 \(\{\chi_{[0, r]} : r \in (0,1)\}\) 的不可数性给出（任意两个不同元素间距离为 \(1\)）。\(\blacksquare\)

\hrulefill

\subsection{\texorpdfstring{48.5 \(L^2\) 空间与内积}{48.5 L\^{}2 空间与内积}}\label{l2-ux7a7aux95f4ux4e0eux5185ux79ef}

\textbf{定理 48.5.1}：\(L^2(E)\) 配备内积 \(\langle f, g \rangle = \int_E fg \, dm\) 是 Hilbert 空间。

\textbf{证明}：由 Hölder 不等式（\(p = q = 2\)），\(\langle f, g \rangle\) 有限。内积的四条公理直接验证。完备性由 Riesz-Fischer 定理（\(p = 2\)）。\(\blacksquare\)

\(L^2\) 空间的 Hilbert 结构为它赋予了 \(L^p\) 族中独一无二的地位。内积的存在使得''角度''和''正交性''概念得以引入，这直接通向线性逼近理论中最核心的工具------正交投影。下面的定理表明，在 Hilbert 空间中，到闭子空间的最佳逼近元不仅存在且唯一，而且由正交性条件精确刻画，这为 Fourier 分析、小波理论和 Galerkin 方法提供了几何基础。

\textbf{定理 48.5.2（ \(L^2\) 中的正交投影与最佳逼近）}：设 \(M \subseteq L^2(E)\) 为闭子空间。则对任意 \(f \in L^2(E)\)，存在唯一的 \(g \in M\) 使得 \(\|f - g\|_2 = \operatorname{dist}(f, M)\)。\(g\) 是 \(f\) 在 \(M\) 上的正交投影，满足 \(\langle f - g, h \rangle = 0\)（\(\forall h \in M\)）。

\textbf{证明}：由 Hilbert 空间的投影定理（或变分法直接证明）。取极小化序列 \(\{g_n\} \subseteq M\)，\(\|f - g_n\|_2 \to d = \operatorname{dist}(f, M)\)。由平行四边形法则，\(\|g_n - g_m\|_2^2 = 2\|f - g_n\|_2^2 + 2\|f - g_m\|_2^2 - 4\|f - (g_n + g_m)/2\|_2^2 \leq 2\|f - g_n\|_2^2 + 2\|f - g_m\|_2^2 - 4d^2 \to 0\)。故 \(\{g_n\}\) 为 Cauchy 列，收敛到 \(g \in M\)。正交性由变分条件 \(\frac{d}{dt}\|f - (g + th)\|_2^2|_{t=0} = 0\) 推出。\(\blacksquare\)

\hrulefill

\subsection{\texorpdfstring{48.6 对偶空间与 \(L^p\) 的表示}{48.6 对偶空间与 L\^{}p 的表示}}\label{ux5bf9ux5076ux7a7aux95f4ux4e0e-lp-ux7684ux8868ux793a}

\textbf{定理 48.6.1（\(L^p\) 对偶定理 / Riesz 表示定理）}：设 \(1 \leq p < \infty\)，\(q\) 为 \(p\) 的共轭指数（\(\frac{1}{p} + \frac{1}{q} = 1\)，约定 \(p=1\) 时 \(q=\infty\)）。则 \((L^p(E))^* \cong L^q(E)\)。具体地，对任意连续线性泛函 \(\Phi \in (L^p)^*\)，存在唯一的 \(g \in L^q\) 使得 \(\Phi(f) = \int_E f g \, dm\)（对所有 \(f \in L^p\)），且 \(\|\Phi\| = \|g\|_q\)。

\textbf{证明}（\(1 < p < \infty\) 的情形）：首先，对 \(g \in L^q\)，由 Holder 不等式，\(|\int f g| \leq \|f\|_p \|g\|_q\)，故 \(g\) 定义了 \(L^p\) 上的连续线性泛函，范数 \(\leq \|g\|_q\)。反向：设 \(\Phi \in (L^p)^*\)。对可测集 \(A\)，定义 \(\nu(A) = \Phi(\chi_A)\)。由 \(\Phi\) 的线性性和连续性，\(\nu\) 是 \(m\) 绝对连续的复测度。由 Radon-Nikodym 定理，存在 \(g \in L^1\) 使 \(\Phi(\chi_A) = \int_A g \, dm\)。由线性性推广到简单函数，再在 \(L^p\) 中稠密。为证 \(g \in L^q\) 且 \(\|\Phi\| = \|g\|_q\)，取 \(f = |g|^{q-1} \operatorname{sgn}(g)\) 验证。\(\blacksquare\)

这个表示定理是现代泛函分析的里程碑之一。它告诉我们，\(L^p\) 空间的对偶空间本质上就是其共轭空间 \(L^q\)------这一结论将抽象的对偶概念具体化为积分表示，极大地简化了 \(L^p\) 中弱收敛等问题的分析。然而，\(p=1\) 与 \(p=\infty\) 的情形值得特别注意，下面这条注记揭示了该理论天然的不对称性。

\textbf{注记 48.6.1}：\(p = \infty\) 时，\((L^\infty)^* \neq L^1\)。\((L^\infty)^*\) 严格大于 \(L^1\)，它包含如 Banach 极限等非 \(L^1\) 泛函。这是 \(L^p\) 理论中 \(p=1\) 和 \(p=\infty\) 不对称性的根源之一。\(L^1\) 的对偶是 \(L^\infty\)，但 \(L^\infty\) 的对偶严格大于 \(L^1\)------这一事实反映了 \(L^1\) 不是自反空间，而 \(L^p\)（\(1 < p < \infty\)）是自反的。

对偶表示定理的一个直接应用是刻画 \(L^p\) 空间中的弱收敛。在偏微分方程和变分方法中，弱收敛往往是比强收敛更自然的收敛模式------例如，极小化序列通常只在弱拓扑下紧。下面这个命题将抽象的弱收敛转化为可以通过对偶作用验证的条件，这使得 \(L^p\) 中的弱收敛分析变得可操作。

\textbf{命题 48.6.2（ \(L^p\) 中弱收敛的刻画）}：在 \(L^p\)（\(1 < p < \infty\)）中，\(f_n \rightharpoonup f\)（弱收敛）当且仅当 \(\{f_n\}\) 范数有界且对任意 \(g \in L^q\)，\(\int f_n g \to \int f g\)。\(L^1\) 中弱收敛的刻画需要 Dunford-Pettis 一致可积性条件（等价于 \(\{f_n\}\) 的相对弱紧性）。

\hrulefill

\subsection{48.7 插值定理}\label{ux63d2ux503cux5b9aux7406}

\textbf{定理 48.7.1（Riesz-Thorin 插值定理）}：设 \(T\) 为定义在简单函数上的线性算子，满足 \(\|Tf\|_{q_0} \leq M_0 \|f\|_{p_0}\) 和 \(\|Tf\|_{q_1} \leq M_1 \|f\|_{p_1}\)。对任意 \(\theta \in (0,1)\)，定义

\[\frac{1}{p} = \frac{1-\theta}{p_0} + \frac{\theta}{p_1}, \quad \frac{1}{q} = \frac{1-\theta}{q_0} + \frac{\theta}{q_1}\]

则 \(T\) 可延拓为 \(L^p\) 到 \(L^q\) 的有界算子，且 \(\|T\|_{L^p \to L^q} \leq M_0^{1-\theta} M_1^{\theta}\)。

\textbf{证明}：使用 Hadamard 三线定理（Phragmén-Lindelöf 原理的推论）。构造解析函数 \(F(z) = \int (Tf_z) g_z\)，其中 \(f_z, g_z\) 为参数化到复平面的简单函数族。\(F(z)\) 在带域 \(\{0 \leq \Re(z) \leq 1\}\) 上解析有界。由假设，\(|F(iy)| \leq M_0\)，\(|F(1+iy)| \leq M_1\)。三线定理给出 \(|F(\theta)| \leq M_0^{1-\theta} M_1^{\theta}\)，即得所求证。\(\blacksquare\)

\textbf{定理 48.7.2（Marcinkiewicz 插值定理）}：设 \(T\) 为弱型 \((p_0, q_0)\) 和 \((p_1, q_1)\) 的次线性算子。则对 \(p_0 < p < p_1\)，\(T\) 是强型 \((p, p)\) 的（即 \(L^p\) 有界）。此定理比 Riesz-Thorin 定理更适用于极大算子、奇异积分等仅满足弱型估计的算子。

\textbf{注记}：Riesz-Thorin 依赖复分析（解析函数的最大模原理），给出算子范数的精确插值界；Marcinkiewicz 依赖实分析（分布函数的 Calderón-Zygmund 分解），适用于更广泛的次线性算子类。两者共同构成了 \(L^p\) 有界性理论的基石。\(\blacksquare\)

\textbf{定义 48.7.1（弱 \(L^p\) 空间）}：对 \(1 \leq p < \infty\)，弱 \(L^p\) 空间（或 \(L^{p,\infty}\)）定义为满足以下条件的可测函数 \(f\) 的集合：

\[
\|f\|_{L^{p,\infty}} = \sup_{\lambda > 0} \lambda \cdot m(\{x : |f(x)| > \lambda\})^{1/p} < \infty
\]

\textbf{命题 48.7.3}：\(L^p \subsetneq L^{p,\infty}\)。例如，\(f(x) = |x|^{-n/p}\) 属于 \(L^{p,\infty}(\mathbb{R}^n)\) 但不属于 \(L^p(\mathbb{R}^n)\)。弱 \(L^p\) 空间是 Lorentz 空间 \(L^{p,q}\) 的特例（\(q = \infty\)）。

\hrulefill

\subsection{\texorpdfstring{48.8 \(L^p\) 空间的一致凸性与自反性}{48.8 L\^{}p 空间的一致凸性与自反性}}\label{lp-ux7a7aux95f4ux7684ux4e00ux81f4ux51f8ux6027ux4e0eux81eaux53cdux6027}

\textbf{定理 48.8.1（Clarkson 不等式）}：对 \(1 < p < \infty\)，\(L^p\) 是一致凸的 Banach 空间。具体地，对 \(f, g \in L^p\)，当 \(p \geq 2\) 时 \(\|f + g\|_p^p + \|f - g\|_p^p \leq 2^{p-1}(\|f\|_p^p + \|g\|_p^p)\)；当 \(1 < p \leq 2\) 时类似不等式成立（以 \(p' = p/(p-1)\) 形式）。

\textbf{证明}：利用标量不等式 \(|a+b|^p + |a-b|^p \leq 2^{p-1}(|a|^p + |b|^p)\)（\(p \geq 2\)），逐点积分即得。\(1 < p < 2\) 时用对偶性。Clarkson 不等式直接蕴含 \(L^p\)（\(1 < p < \infty\)）的一致凸性，由 Milman-Pettis 定理，一致凸 Banach 空间必为自反空间。因此 \((L^p)^{**} = L^p\)。\(\blacksquare\)

一致凸性是比自反性更强的几何条件，它量化了单位球面的''圆度''：在一致凸空间中，单位球面上任意两个距离较远的点的中点必须严格位于球内深处。这一概念由 Clarkson 于 1936 年引入，正是为了证明 \(L^p\) 的自反性。下面给出其精确定义。

\textbf{定义 48.8.1（一致凸 Banach 空间）}：Banach 空间 \(X\) 称为\textbf{一致凸}的，若对任意 \(\varepsilon > 0\)，存在 \(\delta > 0\) 使得对任意 \(\|x\| = \|y\| = 1\)，\(\|x - y\| \geq \varepsilon\) 蕴含 \(\|(x + y)/2\| \leq 1 - \delta\)。

一致凸性的重要性在于它架起了几何与对偶之间的桥梁。Clarkson 证明了 \(L^p\)（\(1 < p < \infty\)）满足一致凸性，而 Milman-Pettis 定理则表明这一几何性质足以推出自反性------即自然界中存在从 \(X^{**}\) 回到 \(X\) 的典范同构。这条定理的证明同样深刻，其核心在于用一致凸性的定量性质控制 Goldstine 嵌入的像。

\textbf{定理 48.8.2（Milman-Pettis 定理）}：一致凸 Banach 空间是自反的。

这一系列结果建立了一条清晰的逻辑链：Clarkson 不等式 → 一致凸性 → 自反性。但需要警惕的是，\(L^1\) 和 \(L^\infty\) 两项均游离在这一链条之外，它们既非一致凸也非自反，这导致许多在 \(L^p\)（\(1 < p < \infty\)）中成立的性质------如有界集即相对弱紧------在 \(L^1\) 和 \(L^\infty\) 中完全失效。下面的注记对此进行了详细说明。

\textbf{注记 48.8.1}：\(L^1\) 和 \(L^\infty\) 既非一致凸也非自反------这是 \(L^p\) 理论中 \(p=1\) 和 \(p=\infty\) 诸多特殊现象的根源。例如 \((L^1)^* = L^\infty\) 但 \((L^\infty)^* \supsetneq L^1\)；\(L^1\) 中弱紧性需 Dunford-Pettis 一致可积性判别法，而 \(L^p\)（\(1 < p < \infty\)）中弱紧性自动由有界性保证。\(\blacksquare\)

自反性的一个深远推论涉及弱收敛与强收敛之间的关系。在非自反空间中，弱收敛和强收敛可以相去甚远；但在自反空间中，Mazur 引理提供了一个强有力的桥梁------它断言弱收敛序列的凸组合可以强收敛。这一结果在变分法和偏微分方程中至关重要，因为它允许我们将弱极限提升为强收敛，从而验证极小化序列（在弱拓扑下紧）确实收敛到真解。

\textbf{定理 48.8.3（ \(L^p\) 中弱收敛与逐点收敛的 Mazur 引理）}：设 \(f_n \rightharpoonup f\) 在 \(L^p\)（\(1 \leq p < \infty\)）中。则存在 \(\{f_n\}\) 的凸组合序列 \(\{g_n\}\) 使得 \(g_n \to f\) 在 \(L^p\) 范数下（从而有子列几乎处处收敛）。Mazur 引理表明：弱收敛的序列可以通过取凸组合而强收敛，这是凸集的闭包与弱闭包等价这一事实的推论。

\hrulefill

\subsection{\texorpdfstring{48.9 \(L^p\) 空间中的卷积与 Young 不等式}{48.9 L\^{}p 空间中的卷积与 Young 不等式}}\label{lp-ux7a7aux95f4ux4e2dux7684ux5377ux79efux4e0e-young-ux4e0dux7b49ux5f0f}

\textbf{定理 48.9.1（Young 卷积不等式）}：设 \(f \in L^p(\mathbb{R}^n)\)，\(g \in L^q(\mathbb{R}^n)\)，\(1 \leq p, q \leq \infty\) 且 \(\frac{1}{p} + \frac{1}{q} \geq 1\)。定义 \(r\) 满足 \(\frac{1}{r} = \frac{1}{p} + \frac{1}{q} - 1\)。则卷积 \(f * g \in L^r(\mathbb{R}^n)\)，且

\[
\|f * g\|_r \leq \|f\|_p \|g\|_q
\]

\textbf{证明}：当 \(p = 1\) 时，\(r = q\)，由 Minkowski 积分不等式直接得到。一般情形由 Riesz-Thorin 插值定理从端点情形 \((1, q, q)\) 和 \((q', \infty, \infty)\) 推出（其中 \(q'\) 是 \(q\) 的共轭指数）。\(\blacksquare\)

Young 卷积不等式是调和分析中卷积估计的基石，其优雅之处在于它同时处理了 \(L^p\)、\(L^q\) 和 \(L^r\) 三个指标之间的代数关系。特别地，从证明中可以看出，当 \(p=1\) 时卷积算子将 \(L^q\) 映到自身，这正是逼近恒等和磨光技术（mollifier）的理论基础。下面的注记进一步阐释了这一不等式的具体应用场景。

\textbf{注记 48.9.1}：Young 卷积不等式是分析中最基本的卷积估计。当 \(p=1\) 时，\(\|f * g\|_q \leq \|f\|_1 \|g\|_q\)，表明 \(L^1\) 函数作为卷积核算子作用于 \(L^q\) 时是有界的。当 \(p = q = 2\) 时，\(r = \infty\)，即 \(\|f * g\|_\infty \leq \|f\|_2 \|g\|_2\)。这一不等式在 Fourier 分析和偏微分方程中无处不在------例如，热方程的解 \(u(t) = \Phi_t * u_0\)（\(\Phi_t\) 为 Gauss 核）的 \(L^p\) 估计即由 Young 不等式直接给出。

\hrulefill

\subsection{\texorpdfstring{48.10 \(L^p\) 空间中 Fourier 乘子与 Mihlin 定理}{48.10 L\^{}p 空间中 Fourier 乘子与 Mihlin 定理}}\label{lp-ux7a7aux95f4ux4e2d-fourier-ux4e58ux5b50ux4e0e-mihlin-ux5b9aux7406}

\textbf{定义 48.10.1（Fourier 乘子）}：设 \(m \in L^\infty(\mathbb{R}^n)\)。定义 Fourier 乘子算子 \(T_m\) 为 \((T_m f)^{\hat{}}(\xi) = m(\xi) \hat{f}(\xi)\)。若 \(T_m\) 是 \(L^p(\mathbb{R}^n)\) 上的有界算子，则称 \(m\) 为 \(L^p\) Fourier 乘子。

Fourier 乘子的核心问题是：给定一个函数 \(m(\xi)\)（乘子），它在频域的乘法作用何时能保证在空域的 \(L^p\) 范数下有界？对于 \(p=2\)，Plancherel 定理给出平凡的回答------任何 \(L^\infty\) 函数都是 \(L^2\) 乘子。但对于 \(p \neq 2\)，问题变得极为深刻，需要乘子满足一定的光滑性条件。Mihlin 和 Hörmander 独立地发现了以尺度不变导数条件为核心的充分性判据，这便是下面这条在现代调和分析中不可或缺的定理。

\textbf{定理 48.10.1（Mihlin-Hörmander 乘子定理）}：设 \(m \in C^{\lfloor n/2 \rfloor + 1}(\mathbb{R}^n \setminus \{0\})\) 满足 \(|\partial^\alpha m(\xi)| \leq C_\alpha |\xi|^{-|\alpha|}\)（对所有 \(|\alpha| \leq \lfloor n/2 \rfloor + 1\)）。则 \(m\) 是 \(L^p(\mathbb{R}^n)\) 乘子（\(1 < p < \infty\)）。这是判断 Fourier 乘子 \(L^p\) 有界性的最常用工具。\(\blacksquare\)

\newpage

\chapter{06-分析/01-实分析/05-Fourier级数.md}\label{ux5206ux679001-ux5b9eux5206ux679005-fourierux7ea7ux6570.md}

\section{第158章：Fourier 级数}\label{ux7b2c158ux7ae0fourier-ux7ea7ux6570}

Fourier 级数理论是调和分析的起点，起源于 Joseph Fourier 在 1807 年关于热传导方程的研究。Fourier 的基本洞察是：任意周期函数可以表示为三角函数的无穷级数。这一断言在当时引起了激烈争议（Lagrange、Laplace 等人怀疑其一般性），但最终导致了数学分析中一系列根本性概念的诞生------函数的概念、收敛的概念、以及积分理论的严格化。Fourier 级数理论在 \(L^2\) 框架下达到最优雅的形式，而 \(L^p\) 框架下的逐点收敛问题则成为 20 世纪调和分析的核心课题。

\subsection{50.1 正交系与 Bessel 不等式}\label{ux6b63ux4ea4ux7cfbux4e0e-bessel-ux4e0dux7b49ux5f0f}

\textbf{定义 50.1.1（ \(L^2\) 中的规范正交系）}：\(L^2([-\pi, \pi])\) 中的函数列 \(\{\varphi_k\}_{k=1}^{\infty}\) 称为\textbf{规范正交系}，若

\[
\langle \varphi_i, \varphi_j \rangle = \int_{-\pi}^{\pi} \varphi_i(x) \overline{\varphi_j(x)} \, dx = \delta_{ij}
\]

\textbf{定义 50.1.2（完备规范正交系）}：规范正交系 \(\{\varphi_k\}\) 称为\textbf{完备}的，若 \(\langle f, \varphi_k \rangle = 0\)（对所有 \(k\)）蕴含 \(f = 0\) a.e.。等价地，\(\{\varphi_k\}\) 的有限线性张成在 \(L^2\) 中稠密。

\textbf{定理 50.1.1（Bessel 不等式）}：设 \(\{\varphi_k\}\) 为 \(L^2(E)\) 中的规范正交系，\(f \in L^2(E)\)。则

\[
\sum_{k=1}^{\infty} |\langle f, \varphi_k \rangle|^2 \leq \|f\|_2^2
\]

\textbf{证明}：对任意 \(n\)，令 \(s_n = \sum_{k=1}^{n} \langle f, \varphi_k \rangle \varphi_k\)。则

\[
0 \leq \|f - s_n\|_2^2 = \|f\|_2^2 - 2\langle f, s_n \rangle + \|s_n\|_2^2 = \|f\|_2^2 - 2\sum_{k=1}^{n} |\langle f, \varphi_k \rangle|^2 + \sum_{k=1}^{n} |\langle f, \varphi_k \rangle|^2 = \|f\|_2^2 - \sum_{k=1}^{n} |\langle f, \varphi_k \rangle|^2
\]

因此 \(\sum_{k=1}^{n} |\langle f, \varphi_k \rangle|^2 \leq \|f\|_2^2\)。令 \(n \to \infty\) 即得。\(\blacksquare\)

\textbf{定理 50.1.2（Parseval 恒等式）}：设 \(\{\varphi_k\}\) 为 \(L^2(E)\) 中的完备规范正交系。则对任意 \(f \in L^2(E)\)，

\[
\sum_{k=1}^{\infty} |\langle f, \varphi_k \rangle|^2 = \|f\|_2^2
\]

且 \(f = \sum_{k=1}^{\infty} \langle f, \varphi_k \rangle \varphi_k\)（在 \(L^2\) 范数下收敛）。

\textbf{命题 50.1.3（Riesz-Fischer 定理的 Hilbert 空间版本）}：设 \(\{c_k\}\) 为满足 \(\sum |c_k|^2 < \infty\) 的复数列，\(\{\varphi_k\}\) 为规范正交系。则存在唯一的 \(f \in \overline{\operatorname{span}}\{\varphi_k\}\) 使得 \(\langle f, \varphi_k \rangle = c_k\)（对所有 \(k\)）。

\hrulefill

\subsection{\texorpdfstring{50.2 Fourier 级数的 \(L^2\) 收敛}{50.2 Fourier 级数的 L\^{}2 收敛}}\label{fourier-ux7ea7ux6570ux7684-l2-ux6536ux655b}

\textbf{定义 50.2.1（Fourier 系数与 Fourier 级数）}：设 \(f \in L^1([-\pi, \pi])\)。\(f\) 的 Fourier 系数定义为

\[
a_0 = \frac{1}{2\pi} \int_{-\pi}^{\pi} f(x) \, dx, \quad a_n = \frac{1}{\pi} \int_{-\pi}^{\pi} f(x) \cos nx \, dx, \quad b_n = \frac{1}{\pi} \int_{-\pi}^{\pi} f(x) \sin nx \, dx
\]

\(f\) 的 Fourier 级数为

\[
f(x) \sim \frac{a_0}{2} + \sum_{n=1}^{\infty} (a_n \cos nx + b_n \sin nx)
\]

使用复指数形式更为简洁：令 \(c_n = \frac{1}{2\pi} \int_{-\pi}^{\pi} f(x) e^{-inx} dx\)，则 \(f(x) \sim \sum_{n=-\infty}^{\infty} c_n e^{inx}\)。两种形式的关系为 \(c_0 = a_0/2\)，\(c_n = (a_n - i b_n)/2\)（\(n \geq 1\)），\(c_{-n} = \overline{c_n}\)。

\textbf{注记 50.2.1}：此处使用 \(\sim\) 而非 \(=\)，因为 Fourier 级数不一定逐点收敛到 \(f\)------这是 Fourier 级数理论的核心问题。在 \(L^2\) 框架下，级数在 \(L^2\) 范数下收敛到 \(f\)；在逐点意义下，需要额外的正则性条件（如 Dini 条件、有界变差等）。

\textbf{定理 50.2.1（ \(L^2\) 中 Fourier 级数的收敛性）}：设 \(f \in L^2([-\pi, \pi])\)。则 \(f\) 的 Fourier 级数在 \(L^2\) 范数下收敛到 \(f\)。即

\[
\lim_{N \to \infty} \left\|f - \left(\frac{a_0}{2} + \sum_{n=1}^{N} (a_n \cos nx + b_n \sin nx)\right)\right\|_2 = 0
\]

等价地，Parseval 恒等式成立：

\[
\frac{1}{\pi} \int_{-\pi}^{\pi} |f(x)|^2 \, dx = \frac{a_0^2}{2} + \sum_{n=1}^{\infty} (a_n^2 + b_n^2)
\]

\textbf{证明}：三角函数系 \(\{\frac{1}{\sqrt{2\pi}}, \frac{\cos nx}{\sqrt{\pi}}, \frac{\sin nx}{\sqrt{\pi}}\}_{n=1}^\infty\) 在 \(L^2([-\pi, \pi])\) 中是规范正交系。由 Weierstrass 第二逼近定理（推论 50.4.2），三角多项式在 \(L^2\) 中稠密，故此规范正交系是完备的。由 Hilbert 空间理论，完备规范正交系中任意元素的 Fourier 展开在 \(L^2\) 范数下收敛到该元素，且 Parseval 恒等式 \(\|f\|_2^2 = \sum |\langle f, \varphi_k \rangle|^2\) 成立。将 \(\langle f, \varphi_k \rangle\) 具体化为 Fourier 系数即得结论。\(\blacksquare\)

\textbf{推论 50.2.2（Fourier 系数的唯一性）}：若 \(f, g \in L^1([-\pi, \pi])\) 有相同的 Fourier 系数，则 \(f = g\) a.e.。特别地，若 \(f \in L^2\)，由 Parseval 恒等式 \(\|f - g\|_2 = 0\) 即得。对 \(L^1\) 情形，需用 Fejer 定理的 Cesaro 可和性。

\hrulefill

\subsection{50.3 Dirichlet 核与 Fejér 核}\label{dirichlet-ux6838ux4e0e-fejuxe9r-ux6838}

\textbf{定义 50.3.1（Dirichlet 核）}：\(N\) 阶 Dirichlet 核为

\[
D_N(t) = \frac{1}{2} + \sum_{n=1}^{N} \cos nt = \frac{\sin((N + \frac{1}{2})t)}{2\sin(t/2)}
\]

Fourier 级数的部分和可表示为 \(S_N f(x) = \frac{1}{\pi} \int_{-\pi}^{\pi} f(x - t) D_N(t) \, dt\)。

\textbf{定义 50.3.2（Fejér 核）}：\(N\) 阶 Fejér 核为 Dirichlet 核的 Cesàro 平均：

\[
F_N(t) = \frac{1}{N+1} \sum_{k=0}^{N} D_k(t) = \frac{1}{2(N+1)} \left(\frac{\sin((N+1)t/2)}{\sin(t/2)}\right)^2
\]

Fejér 核构成\textbf{好核}：\(F_N(t) \geq 0\)，\(\frac{1}{\pi} \int_{-\pi}^{\pi} F_N(t) dt = 1\)，且对任意 \(\delta > 0\)，\(\int_{\delta \leq |t| \leq \pi} F_N(t) dt \to 0\)（\(N \to \infty\)）。

\textbf{命题 50.3.1（好核的逼近性质）}：设 \(\{K_N\}\) 为 \([-\pi, \pi]\) 上的好核族（即满足上述三条性质）。则对任意 \(f \in L^1([-\pi, \pi])\)，卷积 \(K_N * f\) 在 \(f\) 的每个连续点处收敛到 \(f(x)\)；若 \(f\) 连续且周期，则收敛是一致的。

\textbf{证明}：\(|K_N * f(x) - f(x)| = |\frac{1}{\pi} \int_{-\pi}^{\pi} (f(x-t) - f(x)) K_N(t) dt| \leq \frac{1}{\pi} \int_{-\pi}^{\pi} |f(x-t) - f(x)| K_N(t) dt\)。将积分分为 \(|t| < \delta\) 和 \(|t| \geq \delta\) 两部分。第一部分由 \(f\) 在 \(x\) 处的连续性，\(|f(x-t) - f(x)| < \varepsilon\)；第二部分由好核的第三条性质，\(K_N(t)\) 在 \(|t| \geq \delta\) 上的积分趋于 \(0\)，而 \(f\) 的 \(L^1\) 有界性给出控制。\(\blacksquare\)

\hrulefill

\subsection{50.4 Fejér 定理与 Weierstrass 逼近定理}\label{fejuxe9r-ux5b9aux7406ux4e0e-weierstrass-ux903cux8fd1ux5b9aux7406}

\textbf{定理 50.4.1（Fejér 定理）}：设 \(f \in L^1([-\pi, \pi])\)。则 \(f\) 的 Fourier 级数的 Cesàro 平均 \(\sigma_N f(x) = \frac{1}{N+1} \sum_{k=0}^{N} S_k f(x)\) 满足：

\begin{enumerate}
\def\labelenumi{(\roman{enumi})}
\tightlist
\item
  若 \(f\) 在 \(x\) 处连续，则 \(\sigma_N f(x) \to f(x)\)
\item
  若 \(f\) 在 \([-\pi, \pi]\) 上连续且 \(f(-\pi) = f(\pi)\)，则 \(\sigma_N f \to f\) 一致收敛
\end{enumerate}

\textbf{证明}：\(\sigma_N f(x) = \frac{1}{\pi} \int_{-\pi}^{\pi} f(x - t) F_N(t) dt\)。由 Fejér 核是''好核''的性质，应用命题 50.3.1 即得结论。\(\blacksquare\)

\textbf{推论 50.4.2（Weierstrass 第二逼近定理）}：三角多项式在 \(C([-\pi, \pi])\)（周期连续函数）中关于一致范数稠密。因此，任一连续周期函数可用三角多项式一致逼近。

\textbf{注记 50.4.1}：Fejér 定理表明，虽然 Fourier 级数的部分和 \(S_N f\) 可能不收敛（甚至在连续点处发散），但其 Cesàro 平均 \(\sigma_N f\) 总是收敛的。这一现象是发散级数可和性理论（Cesàro 可和性）的经典例子。Fejér 定理的一个优美推论是：若 \(f\) 的 Fourier 级数在某点收敛，则其和必为 \(f\) 在该点的值（若 \(f\) 在该点连续）。

\textbf{命题 50.4.3（Fejér 定理与 \(L^p\) 收敛）}：对 \(1 \leq p < \infty\)，若 \(f \in L^p([-\pi, \pi])\)，则 \(\sigma_N f \to f\) 在 \(L^p\) 范数下。这由好核的 \(L^p\) 逼近性质和 Minkowski 积分不等式推出。

\hrulefill

\subsection{50.5 逐点收敛性}\label{ux9010ux70b9ux6536ux655bux6027}

\textbf{定理 50.5.1（Riemann-Lebesgue 引理）}：设 \(f \in L^1([-\pi, \pi])\)。则

\[
\lim_{|n| \to \infty} \int_{-\pi}^{\pi} f(t) e^{-int} \, dt = 0
\]

即 Fourier 系数 \(a_n, b_n \to 0\)（\(n \to \infty\)）。

\textbf{证明}：先对 \(f\) 为特征函数 \(\chi_{(a,b)}\) 验证：\(\int_a^b e^{-int} dt = \frac{e^{-ina} - e^{-inb}}{in} \to 0\)。对简单函数（特征函数的线性组合）由线性性成立。对一般 \(f \in L^1\)，取简单函数 \(\varphi\) 使 \(\|f - \varphi\|_1 < \varepsilon/2\)，则 \(|\int (f - \varphi) e^{-int} dt| \leq \|f - \varphi\|_1 < \varepsilon/2\)，且当 \(|n|\) 足够大时 \(|\int \varphi e^{-int} dt| < \varepsilon/2\)，故 \(|\int f e^{-int} dt| < \varepsilon\)。\(\blacksquare\)

\textbf{定理 50.5.2（Dini 判别法）}：设 \(f \in L^1([-\pi, \pi])\)，\(x \in [-\pi, \pi]\)。若存在 \(\delta > 0\) 使得

\[
\int_0^{\delta} \frac{|f(x + t) + f(x - t) - 2f(x)|}{t} \, dt < \infty
\]

则 \(f\) 的 Fourier 级数在 \(x\) 处收敛到 \(f(x)\)。特别地，若 \(f\) 在 \(x\) 处 Lipschitz 连续（或可导），则 Dini 条件成立。

\textbf{证明}：Fourier 部分和可写为 \(S_N f(x) - f(x) = \frac{1}{\pi} \int_0^{\pi} \frac{f(x+t) + f(x-t) - 2f(x)}{2\sin(t/2)} \sin((N+\frac{1}{2})t) dt\)。令 \(\varphi(t) = \frac{f(x+t) + f(x-t) - 2f(x)}{2\sin(t/2)}\)。由 Dini 条件，\(\frac{|f(x+t) + f(x-t) - 2f(x)|}{t}\) 在 \(t \in (0,\delta)\) 可积。由于 \(\sin(t/2) \sim t/2\)（\(t \to 0\)），\(\varphi(t)\) 在 \([0,\pi]\) 上 \(L^1\) 可积。由 Riemann-Lebesgue 引理，\(\int_0^{\pi} \varphi(t) \sin((N+\frac{1}{2})t) dt \to 0\)（\(N \to \infty\)），故 \(S_N f(x) \to f(x)\)。\(\blacksquare\)

\textbf{定理 50.5.3（Dirichlet-Jordan 判别法）}：若 \(f\) 在 \(x\) 的某个邻域内有界变差，则 \(f\) 的 Fourier 级数在 \(x\) 处收敛到 \(\frac{f(x^+) + f(x^-)}{2}\)。

\textbf{证明}：设 \(f\) 在 \([x-\delta, x+\delta]\) 上有界变差。将 \(f\) 分解为两个单调函数之差：\(f = g - h\)，其中 \(g, h\) 单调。对单调函数 \(g\)，\(S_N g(x) - g(x) = \frac{1}{\pi} \int_0^\pi (g(x+t) + g(x-t) - 2g(x)) \frac{\sin((N+1/2)t)}{2\sin(t/2)} dt\)。利用 Dirichlet 核的积分表示和 Jordan 分解，将积分分为两部分：\(\int_0^\delta\) 和 \(\int_\delta^\pi\)。第一部分由单调性和三角函数积分的平均值性质（第二中值定理）控制，当 \(N \to \infty\) 时趋于 \(0\)。第二部分由 Riemann-Lebesgue 引理趋于 \(0\)。综合得 \(S_N f(x) \to \frac{f(x^+) + f(x^-)}{2}\)。\(\blacksquare\)

\textbf{注记 50.5.1}：关于逐点收敛的最深刻结果是 Carleson 定理（1966）：若 \(f \in L^2([-\pi, \pi])\)，则 \(f\) 的 Fourier 级数几乎处处收敛到 \(f(x)\)。这一定理（由 Lennart Carleson 证明，他因此获得 2006 年 Abel 奖）解决了 Lusin 在 1915 年提出的猜想，是 20 世纪分析学最伟大的成就之一。Hunt 后来将其推广到 \(L^p\)（\(p > 1\)）。Kolmogorov 在 1923 年构造了 \(L^1\) 中 Fourier 级数几乎处处发散的例子，表明 \(p=1\) 是 Carleson-Hunt 定理的精确边界。

\hrulefill

\subsection{50.6 Gibbs 现象与收敛性补充}\label{gibbs-ux73b0ux8c61ux4e0eux6536ux655bux6027ux8865ux5145}

\textbf{定理 50.6.1（Gibbs 现象）}：设 \(f\) 在 \(x_0\) 处有跳跃间断：\(f(x_0^+) \neq f(x_0^-)\)。则 \(f\) 的 Fourier 级数的部分和 \(S_N f\) 在 \(x_0\) 附近会产生过冲（overshoot）------最大过冲量约为跳跃幅度的 \(9\%\)，即

\[\lim_{N \to \infty} \max_{x \in (x_0, x_0+\delta)} S_N f(x) - f(x_0^+) \approx 0.08949 \cdot |f(x_0^+) - f(x_0^-)|\]

且该过冲不随 \(N\) 增大而消失（仅位置向 \(x_0\) 收缩）。Gibbs 现象是 Fourier 级数在跳跃间断处非一致收敛的必然结果，其本质是 Dirichlet 核的振荡积分在间断处的极限行为。

\textbf{证明}：不妨设 \(f\) 为符号函数 \(\operatorname{sgn}\) 在 \([-\pi, \pi]\) 上的周期延拓。\(S_N f(x) = \frac{2}{\pi} \int_0^x \frac{\sin((N+1/2)t)}{2\sin(t/2)} dt\)。当 \(N \to \infty\) 时，\(S_N f(x) \to \frac{2}{\pi} \operatorname{Si}((N+1/2)x)\)，其中 \(\operatorname{Si}(u) = \int_0^u \frac{\sin t}{t} dt\) 为正弦积分。\(\operatorname{Si}\) 在 \(u = \pi\) 处取最大值 \(\approx 1.8519\)，故过冲为 \(\frac{2}{\pi} \cdot 1.8519 - 1 \approx 0.17898 / 2\)（符号函数跳跃为 \(2\)），相对过冲约 \(8.95\%\)。\(\blacksquare\)

Gibbs 现象不仅仅是 Fourier 级数理论中的一则数学趣闻，它在实际的信号和图像处理中有着直接的表现。当 Fourier 系数被截断（前 \(N\) 项部分和）时，间断点附近的振荡不会随 \(N\) 增大而消失，仅被压缩到越来越窄的区间内。这一特性意味着，单纯增加 Fourier 系数的数量并不能消除间断处的伪影，必须借助更精巧的求和策略。

\textbf{注记 50.6.1}：Gibbs 现象在信号处理中有实际意义------在图像压缩（JPEG）中，Fourier 系数的截断会在边缘处产生振铃效应（ringing artifact），这正是 Gibbs 现象的二维版本。为缓解这一问题，实践中常使用 Fejér 平均（完全消除 Gibbs 现象，但降低分辨率）或使用更光滑的求和核（如 Lanczos 核）。

\hrulefill

\subsection{50.7 Fourier 变换与 Fourier 级数的统一}\label{fourier-ux53d8ux6362ux4e0e-fourier-ux7ea7ux6570ux7684ux7edfux4e00}

\textbf{定理 50.7.1（Poisson 求和公式）}：设 \(f \in \mathcal{S}(\mathbb{R})\)（Schwartz 空间）。则

\[\sum_{n \in \mathbb{Z}} f(n) = \sum_{k \in \mathbb{Z}} \hat{f}(2\pi k)\]

其中 \(\hat{f}(\xi) = \int_{-\infty}^\infty f(x) e^{-i x \xi} dx\) 是 Fourier 变换。Poisson 求和公式是连接 Fourier 级数（离散/周期）与 Fourier 变换（连续/非周期）的桥梁。

\textbf{证明}：定义 \(F(x) = \sum_{n \in \mathbb{Z}} f(x + 2\pi n)\)，则 \(F\) 为 \(2\pi\)-周期的光滑函数。\(F\) 的 Fourier 系数为 \(\hat{f}(k)\)（由定义直接计算）。\(F\) 的 Fourier 级数在 \(x=0\) 处取值得 \(\sum_{n \in \mathbb{Z}} f(2\pi n) = \sum_{k \in \mathbb{Z}} \hat{f}(k)\)，作变量替换 \(x \mapsto 2\pi x\) 即得标准形式。\(\blacksquare\)

\textbf{定理 50.7.2（Fourier 反演公式的周期版本）}：设 \(f \in L^1([-\pi, \pi])\) 且其 Fourier 系数满足 \(\sum |c_n| < \infty\)。则 \(f\) 可调整为连续函数，且 Fourier 级数绝对且一致收敛到 \(f(x)\)：

\[
f(x) = \sum_{n=-\infty}^{\infty} c_n e^{inx}
\]

Poisson 求和公式在数论中有深远应用------它是 Riemann \(\zeta\) 函数的函数方程、\(\theta\) 函数的模变换性质（\(\theta(-1/z) = \sqrt{z/i} \, \theta(z)\)）以及 Hecke 的 Voronoi 求和公式的基础。在调和分析中，它统一了 Fourier 级数和 Fourier 变换，是局部紧 Abel 群上调和分析中 Pontryagin 对偶的离散-紧致对偶特例。\(\blacksquare\)

\hrulefill

\subsection{50.8 共轭函数与 Hilbert 变换}\label{ux5171ux8f6dux51fdux6570ux4e0e-hilbert-ux53d8ux6362}

\textbf{定义 50.8.1（共轭 Fourier 级数）}：设 \(f \sim \frac{a_0}{2} + \sum_{n=1}^{\infty} (a_n \cos nx + b_n \sin nx)\)。\(f\) 的\textbf{共轭函数} \(\tilde{f}\) 的 Fourier 级数定义为

\[
\tilde{f}(x) \sim \sum_{n=1}^{\infty} (a_n \sin nx - b_n \cos nx)
\]

在复形式中，若 \(f \sim \sum c_n e^{inx}\)，则 \(\tilde{f} \sim -i \sum_{n \neq 0} \operatorname{sgn}(n) c_n e^{inx}\)。

\textbf{定理 50.8.1（M. Riesz 定理）}：对 \(1 < p < \infty\)，共轭函数算子 \(f \mapsto \tilde{f}\) 是 \(L^p([-\pi, \pi])\) 上的有界算子。对 \(p=1\)，该算子不是 \(L^1\) 有界的，但满足弱 \((1,1)\) 型估计。

共轭函数是调和函数共轭的边界值，与 Hilbert 变换密切相关。在单位圆盘上，若 \(u\) 是 \(f\) 的 Poisson 积分，则 \(f\) 的共轭函数 \(\tilde{f}\) 是 \(u\) 的调和共轭 \(v\) 的边界值，满足 \(u + iv\) 为全纯函数。这一联系是复分析与调和分析之间最深层的桥梁之一。

\hrulefill

\subsection{50.9 Fourier 级数的绝对收敛与 Wiener 代数}\label{fourier-ux7ea7ux6570ux7684ux7eddux5bf9ux6536ux655bux4e0e-wiener-ux4ee3ux6570}

\textbf{定义 50.9.1（Wiener 代数）}：Wiener 代数 \(A(\mathbb{T})\) 定义为 Fourier 系数满足 \(\sum_{n=-\infty}^{\infty} |c_n| < \infty\) 的所有连续周期函数构成的空间，配以范数 \(\|f\|_A = \sum |c_n|\)。

Wiener 代数的引入源于这样一个自然的问题：Fourier 级数何时不仅逐点收敛，而且绝对收敛------即系数的 \(\ell^1\) 范数有限？这种强收敛模式保证了 Fourier 级数一致收敛，且其和函数具有许多良好的代数性质。下面的 Wiener 引理是 Banach 代数理论的开山之作之一，它表明 \(A(\mathbb{T})\) 在逐点取倒数操作下是封闭的------这在 Fourier 分析中远非显然，因为 \(1/f\) 的 Fourier 系数与 \(f\) 的系数之间没有简单的代数关系。

\textbf{定理 50.9.1（Wiener 引理）}：若 \(f \in A(\mathbb{T})\) 且 \(f(x) \neq 0\)（对所有 \(x\)），则 \(1/f \in A(\mathbb{T})\)。Wiener 引理是 Banach 代数理论中谱不变性的经典结果，在 Tauber 定理和信号处理中有重要应用。

Wiener 引理解决了代数封闭性问题，但如何判断一个函数是否属于 \(A(\mathbb{T})\) 仍然需要可操作的判据。光滑性是保证绝对收敛的最自然途径------直觉上，函数越光滑，其 Fourier 系数衰减越快。Bernstein 定理将这一直觉精确化为定量的 Hölder 连续性条件，给出了 \(f \in A(\mathbb{T})\) 的一个充分判据。特别地，它揭示了绝对收敛与分数阶可微性之间的内在联系，是 Fourier 系数衰减估计的经典应用。

\textbf{定理 50.9.2（Bernstein 定理）}：若 \(f\) 是 \(\alpha\)-Holder 连续的（\(\alpha > 1/2\)），则 \(f \in A(\mathbb{T})\)。更精确地，若 \(f \in C^{k,\alpha}\)（\(k\) 阶导数 \(\alpha\)-Holder 连续），则 Fourier 系数满足 \(|c_n| = O(|n|^{-k-\alpha})\)，从而当 \(k + \alpha > 1/2\) 时 Fourier 级数绝对收敛。\(\blacksquare\)

\hrulefill

\subsection{\texorpdfstring{50.10 Fourier 级数的 \(L^1\) 收敛问题与 Kolmogorov 反例}{50.10 Fourier 级数的 L\^{}1 收敛问题与 Kolmogorov 反例}}\label{fourier-ux7ea7ux6570ux7684-l1-ux6536ux655bux95eeux9898ux4e0e-kolmogorov-ux53cdux4f8b}

\textbf{定理 50.10.1（ \(L^1\) 中 Fourier 级数的不收敛性）}：存在 \(f \in L^1([-\pi, \pi])\)，其 Fourier 级数的部分和 \(S_N f\) 在 \(L^1\) 范数下不收敛于 \(f\)。这一事实与 \(L^p\)（\(p > 1\)）形成鲜明对比，表明 \(L^1\) 是 Fourier 级数理论的''边界情形''。

\(L^1\) 中依范数收敛的失效只是一个序曲。更令人震惊的是，Kolmogorov 在年仅 19 岁时构造的例子：存在 \(L^1\) 函数，其 Fourier 级数非但在 \(L^1\) 范数下不收敛，而且在每一个点上都发散------这与 \(L^p\)（\(p > 1\)）中 Hunting 和 Carleson 建立的几乎处处收敛形成尖锐对照。Kolmogorov 的反例彻底埋葬了''Fourier 级数对一切可积函数都几乎处处收敛''的天真想法，确定了 \(L^1\) 恰是几乎处处收敛的精确边界。

\textbf{定理 50.10.2（Kolmogorov 1923 反例）}：存在 \(f \in L^1([-\pi, \pi])\)，其 Fourier 级数几乎处处发散。Kolmogorov 的构造极为精巧：利用无界发散的三角级数和一个''帐篷函数''（tent function）的叠加，构造了 \(L^1\) 中 Fourier 级数在每点都发散的病态函数。这一反例确立了 \(L^1\) 是 Fourier 级数几乎处处收敛的精确边界------\(L^p\)（\(p > 1\)）中几乎处处收敛（Carleson-Hunt 定理），但 \(L^1\) 中不成立。

Kolmogorov 的否定性结果让学界沉默了近四十年，直到 1966 年 Carleson 在 \(L^2\) 情形下取得了突破性进展------他证明了 \(L^2\) 函数的 Fourier 级数几乎处处收敛。这一结果被誉为 20 世纪分析学中最伟大的成就之一，其证明引入的''时频分解''技术后来被 Fefferman 系统化为相空间分析方法，完全改变了调和分析的研究范式。Carleson 定理表明，从 \(L^1\) 到 \(L^2\) 这一看似微小的提升，在逐点收敛的意义下却是天壤之别。

\textbf{定理 50.10.3（Carleson 定理，1966）}：若 \(f \in L^2([-\pi, \pi])\)，则 \(f\) 的 Fourier 级数几乎处处收敛到 \(f(x)\)。Hunt 在 1968 年将此推广到 \(L^p\)（\(p > 1\)）。Carleson 定理的证明极为复杂（原始论文长达 22 页），核心思想是将 Carleson 极大算子 \(C f(x) = \sup_N |S_N f(x)|\) 的弱 \((2,2)\) 型估计与精细的时频分析（time-frequency analysis）相结合。Fefferman 后来将 Carleson 的证明重构为''相空间分解''（phase space decomposition）的框架，这一方法成为现代调和分析中处理 Fourier 级数逐点收敛的标准工具。

Carleson 定理的完整证明虽然艰深，但有一条相对初等的路径可以通过 Hardy-Littlewood 极大函数来获得部分结果。下面这个不等式将 Fourier 级数部分和的点态行为与极大函数控制联系起来，虽不能给出最优的 \(p > 1\) 结论，但提供了理解逐点收敛机制的直观几何视角------Fourier 级数的部分和本质上由函数的局部平均（极大函数）所控制。

\textbf{定理 50.10.4（Hardy-Littlewood 极大函数与 Fourier 级数）}：Hardy-Littlewood 极大函数 \(Mf\) 控制 Fourier 级数的部分和：\(|S_N f(x)| \leq C \cdot Mf(x)\)（对 \(f \in L^p\)，\(p > 1\) 时成立）。这一不等式表明 Fourier 级数的逐点收敛性可由极大函数的弱型估计导出，是 Carleson 定理的简化版本（仅适用于 \(p > 1\) 且不给出几乎处处收敛的最优结果）。

\subsection{50.11 Fourier 级数的多重推广与球面调和}\label{fourier-ux7ea7ux6570ux7684ux591aux91cdux63a8ux5e7fux4e0eux7403ux9762ux8c03ux548c}

\textbf{定义 50.11.1（多重 Fourier 级数）}：设 \(f \in L^1(\mathbb{T}^n)\)（\(n\) 维环面）。其多重 Fourier 级数为 \(f(x) \sim \sum_{k \in \mathbb{Z}^n} c_k e^{i k \cdot x}\)，其中 \(c_k = \frac{1}{(2\pi)^n} \int_{\mathbb{T}^n} f(x) e^{-i k \cdot x} dx\)。多重 Fourier 级数的收敛性问题远比一维情形复杂------球面部分和、方形部分和和矩形部分和的行为有本质不同。Fefferman 在 1971 年证明了''圆盘乘子猜想''不成立：球面部分和的 \(L^p\) 有界性在 \(p \neq 2\) 时失效（\(n \geq 2\)），这一深刻结果揭示了高维 Fourier 分析与一维之间的本质差异。

当底空间从平坦的环面 \(\mathbb{T}^n\) 转移到弯曲的球面 \(S^{n-1}\) 时，Fourier 分析的基本框架发生了根本性的转变。环面上的指数函数 \(e^{i k \cdot x}\) 被球面调和函数（spherical harmonics）所取代，后者是 Laplace-Beltrami 算子在球面上的特征函数族。这种推广不仅是数学上的自然延拓，更在物理中有着直接的对应------量子力学中的角动量本征态正是球面调和函数的物理化身。下面的定理描述了球面上 \(L^2\) 函数的正交分解，它是 Fourier 级数在非平坦空间中的完整类比。

\textbf{定理 50.11.1（球面调和展开）}：在单位球面 \(S^{n-1}\) 上，\(L^2(S^{n-1})\) 可分解为球面调和函数空间 \(\mathcal{H}_k\)（\(k\) 阶齐次调和多项式在球面上的限制）的正交直和。球面调和函数是 Laplace 算子在球面上的特征函数，其 Fourier 展开推广了圆周上的三角级数展开。球面调和理论在量子力学（角动量理论）、地球物理学（重力场展开）和计算机图形学（环境光照的球面调和表示）中有广泛应用。\(\blacksquare\)

\newpage

\chapter{06-分析/01-实分析/06-可和性与位势论.md}\label{ux5206ux679001-ux5b9eux5206ux679006-ux53efux548cux6027ux4e0eux4f4dux52bfux8bba.md}

\section{第159章：可和性方法}\label{ux7b2c159ux7ae0ux53efux548cux6027ux65b9ux6cd5}

发散级数的可和性是对经典收敛概念的推广，在解析数论、量子场论和微分方程中有重要应用。可和性理论的核心洞察是：发散级数可以通过赋予其''广义和''而获得数学意义，只要该广义和与通常的收敛和相容（正则性），且满足线性性和平移不变性。这一思想在 18 世纪已由 Euler 和 Leibniz 非正式地使用，但直到 19 世纪末 20 世纪初才由 Cesaro、Borel、Abel 和 Hardy 等人严格化。

\subsection{100.1 Cesaro 可和性与 Abel 可和性}\label{cesaro-ux53efux548cux6027ux4e0e-abel-ux53efux548cux6027}

\textbf{定义 100.1.1}（Cesaro 可和）：对于级数 \(\sum_{n=0}^\infty a_n\)，设 \(S_n = \sum_{k=0}^n a_k\)。若 \(S_n\) 的 Cesaro 平均 \(\sigma_n = \frac{S_0 + S_1 + \cdots + S_n}{n+1}\) 收敛于 \(L\)，则称级数是 \((C,1)\) 可和的，和为 \(L\)。

\textbf{定理 100.1.1}（Cesaro 可和的正则性）：若 \(\sum a_n\) 按通常意义收敛于 \(S\)，则按 \((C,1)\) 可和于 \(S\)。\((C,1)\) 包含通常收敛，且更广。

\textbf{证明}：设 \(S_n \to S\)。对任意 \(\varepsilon > 0\)，存在 \(N\) 使得 \(|S_n - S| < \varepsilon/2\) 对所有 \(n \geq N\)。对 \(n > N\)，\(\sigma_n - S = \frac{1}{n+1} \sum_{k=0}^n (S_k - S) = \frac{1}{n+1} \sum_{k=0}^{N-1} (S_k - S) + \frac{1}{n+1} \sum_{k=N}^n (S_k - S)\)。第一项绝对值 \(\leq \frac{1}{n+1} \sum_{k=0}^{N-1} |S_k - S| \to 0\)（\(n \to \infty\)）；第二项绝对值 \(\leq \frac{1}{n+1} \sum_{k=N}^n |S_k - S| \leq \frac{n-N+1}{n+1} \cdot \frac{\varepsilon}{2} < \frac{\varepsilon}{2}\)。故当 \(n\) 足够大时 \(|\sigma_n - S| < \varepsilon\)。因此 \((C,1)\) 可和包含通常收敛。反例 \(\sum (-1)^n\) 说明 \((C,1)\) 可和严格更广。\(\blacksquare\)

Cesaro 可和通过算术平均来''平滑''部分和序列的振荡，是最自然的发散级数求和法之一。与之对偶的是 Abel 可和法，它通过引入衰减因子 \(r^n\) 并令 \(r \to 1^-\) 来赋予级数一个广义和。这两种方法在哲学上截然不同------Cesaro 依赖离散平均，Abel 依赖连续极限------但它们之间存在深刻的内在联系：Frobenius 定理将证明 Cesaro 可和性蕴含 Abel 可和性，从而确立了两者之间的层级关系。

\textbf{定义 100.1.2}（Abel 可和）：若 \(\lim_{r \to 1^-} \sum_{n=0}^\infty a_n r^n\) 存在且等于 \(A\)，则称级数是 \textbf{Abel 可和的}，和为 \(A\)。

\textbf{定理 100.1.2}（Abel 定理）：若 \(\sum a_n\) 收敛于 \(S\)，则 Abel 可和于 \(S\)。

\textbf{证明}：设 \(f(r) = \sum_{n=0}^\infty a_n r^n\)（\(|r| < 1\)）。由 \(\sum a_n\) 收敛，\(f(r)\) 在 \(r=1\) 处 Abel 可和于 \(S\)。利用分部求和（Abel 变换）：\(\sum_{n=0}^N a_n r^n = \sum_{n=0}^{N-1} S_n (r^n - r^{n+1}) + S_N r^N\)。令 \(N \to \infty\)，\(r \to 1^-\)，由 \(S_n \to S\) 和 \(r^n - r^{n+1} \geq 0\) 的归一化，\(\lim_{r \to 1^-} f(r) = S\)。\(\blacksquare\)

Abel 定理保证了 Abel 可和法的正则性，但更令人感兴趣的是 Cesaro 与 Abel 两种方法之间的关系。在 Frobenius 的经典工作中，他证明了 Cesaro 可和性蕴含 Abel 可和性------这意味着 Abel 可和法更''宽容''，能够处理更多发散级数。这一包含关系是可和性方法论中关于''哪种方法更强''这一基本问题的核心结果。

\textbf{定理 100.1.3}（Frobenius 定理）：若级数是 \((C,1)\) 可和的，则也是 Abel 可和的，且和相同。Abel 可和严格强于 \((C,1)\) 可和。

\textbf{证明}：设 \(\sigma_n \to L\)。利用 \(\sum a_n r^n = (1-r)^2 \sum_{n=0}^\infty (n+1) \sigma_n r^n\)（分部求和两次）。由 \(\sigma_n \to L\) 和 \(\sum (n+1) r^n = 1/(1-r)^2\)，可得 \(\lim_{r \to 1^-} \sum a_n r^n = L\)。反例：\(\sum (-1)^n n\) 是 Abel 可和的（和为 \(-1/4\)）但不是 \((C,1)\) 可和的。\(\blacksquare\)

将前面的理论应用于一个具体的发散级数有助于直观地理解这些抽象的可和性概念。级数 \(\sum (-1)^n\) 是最简单的发散交替级数，它既非通常收敛，但其 Cesaro 和 Abel 和均为 \(1/2\)------这正是 Euler 早在 18 世纪就凭直觉写下的结果。这个例子同时验证了正则性、Cesaro 可和性与 Abel 可和性之间的关系。

\textbf{例 100.1.1}：经典发散级数 \(\sum_{n=0}^\infty (-1)^n = 1 - 1 + 1 - 1 + \cdots\) 是 \((C,1)\) 可和的，和为 \(1/2\)（因 \(S_n\) 交替为 \(1\) 和 \(0\)，\(\sigma_n \to 1/2\)）。也是 Abel 可和的：\(\lim_{r \to 1^-} \sum (-1)^n r^n = \lim_{r \to 1^-} 1/(1+r) = 1/2\)。

\hrulefill

\subsection{100.2 高阶 Cesaro 可和性与 Holder 可和性}\label{ux9ad8ux9636-cesaro-ux53efux548cux6027ux4e0e-holder-ux53efux548cux6027}

\textbf{定义 100.2.1}（Cesaro \((C,k)\) 可和）：对 \(k \geq 1\)，定义 \(S_n^{(0)} = S_n\)，\(S_n^{(k)} = \sum_{j=0}^n S_j^{(k-1)}\)。则 Cesaro \((C,k)\) 平均为 \(\sigma_n^{(k)} = \frac{S_n^{(k)}}{\binom{n+k}{k}}\)。若 \(\sigma_n^{(k)} \to L\)，则称级数是 \((C,k)\) 可和的。

\((C,1)\) 可和性只做一次算术平均，但对于增长更快的部分和序列------例如 \(S_n \sim n\) 的级数------单次平均不足以压制发散。此时需要重复平均 \(k\) 次，引入高阶 Cesaro 可和性 \((C,k)\)。直观上，\(k\) 越大，平均的''平滑力度''越强，能够驯服的发散级数也越多。下面的等级体系定理精确地确认了这一直觉：每提高一阶 Cesaro 可和性，严格扩大了可和级数的范围。

\textbf{定理 100.2.1}（Cesaro 可和性的等级体系）：\((C,k)\) 可和性随 \(k\) 增大而严格增强。即 \((C,k) \subsetneq (C,k+1)\)。

\textbf{证明}：首先证明 \((C,k) \subseteq (C,k+1)\)。设 \(\sigma_n^{(k)} \to L\)。由 Cesaro 平均的递推关系 \(\sigma_n^{(k+1)} = \frac{1}{n+1} \sum_{j=0}^n \sigma_j^{(k)}\)（经归一化），由 Stolz-Cesaro 定理（或 Cesaro 平均的正则性），若 \(\sigma_n^{(k)} \to L\)，则其 Cesaro 平均 \(\sigma_n^{(k+1)} \to L\)。因此 \((C,k)\) 可和蕴含 \((C,k+1)\) 可和。严格性：构造级数使其 \((C,k+1)\) 可和但非 \((C,k)\) 可和。具体地，取 \(a_n = (-1)^n \binom{n+k}{k}\)。该级数的部分和 \(S_n\) 增长如 \(n^k\)，其 \((C,k)\) 平均发散，但 \((C,k+1)\) 平均收敛（因为需要多做一次平均以消除 \(n^k\) 的增长）。\(\blacksquare\)

Cesaro 可和性可以通过两种在技术上有微妙区别的方式来实现：直接对部分和做 \(k\) 次求和再归一化（\((C,k)\) 方法），或者反复对部分和序列取算数平均（\((H,k)\) 方法）。尽管这两种途径的操作过程不同，但它们在可和性意义上完全等价------这是一个非平凡的事实，因为平均操作的次序和归一化方式的差异可能导致不同的极限行为。下面的定理确认了这两种方法的殊途同归。

\textbf{定理 100.2.2}（Holder \((H,k)\) 可和）：Holder 可和性通过重复取平均定义：\(H_n^{(0)} = S_n\)，\(H_n^{(k+1)} = \frac{1}{n+1} \sum_{j=0}^n H_j^{(k)}\)。Holder 和 Cesaro 可和性是等价的：\((C,k)\) 可和当且仅当 \((H,k)\) 可和，且和相同。

\textbf{证明}：\((C,k)\) 和 \((H,k)\) 的等价性源于两者都是通过重复取平均来光滑化部分和序列，区别仅在于归一化因子。具体地，\((C,k)\) 平均 \(\sigma_n^{(k)} = S_n^{(k)} / \binom{n+k}{k}\)，而 \((H,k)\) 平均 \(H_n^{(k)}\) 是 \(S_n\) 的 \(k\) 次迭代未归一化平均除以 \((n+1)^k\) 的某种组合。由 Schur 定理或 Knopp 等价定理，两个矩阵 \((c_{nk}^{(k)})\) 和 \((h_{nk}^{(k)})\) 定义了等价的可和性方法，即存在正数列 \(\{p_n\}\) 和 \(\{q_n\}\) 使得 \(\sigma_n^{(k)}\) 和 \(H_n^{(k)}\) 的收敛性等价。事实上，\((C,k)\) 和 \((H,k)\) 均可表示为 Nørlund 平均的特例，且 Nørlund 平均之间的等价性由正则性条件刻画。\(\blacksquare\)

\hrulefill

\subsection{100.3 Borel 可和性与 Euler 可和性}\label{borel-ux53efux548cux6027ux4e0e-euler-ux53efux548cux6027}

\textbf{定义 100.3.1}（Borel 指数可和）：级数 \(\sum a_n\) 称为 Borel 可和的（和为 \(B\)），若 \(\sum_{n=0}^\infty \frac{a_n}{n!} x^n\) 对所有 \(x \geq 0\) 收敛，且

\[
\lim_{x \to \infty} e^{-x} \sum_{n=0}^\infty \frac{S_n}{n!} x^n = B
\]

等价地，若 \(\sum_{n=0}^\infty \frac{a_n}{n!} t^n\) 的收敛半径为无穷，且 \(\int_0^\infty e^{-t} \sum_{n=0}^\infty \frac{a_n}{n!} t^n dt = B\)。

\textbf{定理 100.3.1}（Borel 可和的正则性）：若 \(\sum a_n\) 收敛，则 Borel 可和于相同的和。Borel 可和性严格强于 Abel 可和性。

\textbf{证明}：设 \(\sum a_n\) 收敛于 \(S\)，即 \(S_n \to S\)。定义 \(f(x) = e^{-x} \sum_{n=0}^\infty \frac{S_n}{n!} x^n\)。需证 \(\lim_{x \to \infty} f(x) = S\)。对任意 \(\varepsilon > 0\)，取 \(N\) 使 \(|S_n - S| < \varepsilon\) 对所有 \(n \geq N\)。将 \(f(x) - S\) 分解为 \(e^{-x} \sum_{n=0}^{N-1} \frac{S_n - S}{n!} x^n + e^{-x} \sum_{n=N}^\infty \frac{S_n - S}{n!} x^n\)。第二项绝对值 \(\leq \varepsilon e^{-x} \sum_{n=N}^\infty \frac{x^n}{n!} \leq \varepsilon\)。第一项为 \(e^{-x}\) 乘以多项式，当 \(x \to \infty\) 时趋于 \(0\)。故 \(\limsup_{x \to \infty} |f(x) - S| \leq \varepsilon\)，由 \(\varepsilon\) 任意性得 \(f(x) \to S\)。Borel 严格强于 Abel：例子 \(\sum (-1)^n n!\) 是 Borel 可和的（Borel 变换后收敛半径为无穷，积分收敛）但非 Abel 可和（幂级数收敛半径为 \(0\)）。\(\blacksquare\)

\textbf{注记 100.3.1}：Borel 可和性在量子场论中有重要应用。在微扰量子场论中，Feynman 图的微扰级数通常是发散的（阶乘增长），但可能是 Borel 可和的。Borel 变换（Borel resummation）将发散级数转换为收敛积分，在量子色动力学（QCD）和临界现象的重整化群分析中是不可或缺的工具。

\textbf{定义 100.3.2}（Euler \((E,q)\) 可和）：对 \(q > 0\)，定义 Euler 变换：\(\sum_{n=0}^\infty a_n\) 的 Euler \((E,q)\) 和为 \(\sum_{n=0}^\infty \frac{1}{(1+q)^{n+1}} \sum_{k=0}^n \binom{n}{k} q^{n-k} S_k\)。Euler 方法在解析数论中用于加速级数收敛。

\hrulefill

\subsection{100.4 Tauber 型定理}\label{tauber-ux578bux5b9aux7406}

Tauber 定理给出附加条件，使得更弱的可和性蕴含更强的可和性乃至通常收敛。它们是''软分析''与''硬分析''相结合的典范。

\textbf{定理 100.4.1}（Tauber 第一定理）：若级数是 Abel 可和的且 \(a_n = o(1/n)\)，则级数通常收敛。

\textbf{证明}：设 \(\sum a_n\) 为 Abel 可和于 \(A\)，即 \(\lim_{r \to 1^-} f(r) = A\)，其中 \(f(r) = \sum_{n=0}^\infty a_n r^n\)。条件 \(a_n = o(1/n)\) 意味着 \(n a_n \to 0\)。由 Abel 变换（分部求和），\(f(r) = (1-r) \sum_{n=0}^\infty S_n r^n\)。需证 \(S_n \to A\)。对任意 \(\varepsilon > 0\)，取 \(N\) 使 \(|n a_n| < \varepsilon\) 对所有 \(n \geq N\)。取 \(r = 1 - 1/N\)，则 \(|S_N - f(r)| \leq \sum_{n=0}^N |a_n|(1 - r^n) + \sum_{n=N+1}^\infty |a_n| r^n\)。第一项由 \(r^n \geq 1 - n(1-r)\) 得 \(\sum_{n=0}^N |a_n| n/(N) \leq \frac{1}{N} \sum_{n=0}^N |n a_n|\)，由 \(a_n = o(1/n)\) 知此项趋于 \(0\)。第二项利用 \(|a_n| \leq \varepsilon/n\) 和 \(\sum_{n=N+1}^\infty r^n/n \leq 1/(N(1-r)) = 1\)，得 \(\leq \varepsilon\)。故 \(|S_N - A| \leq |S_N - f(r)| + |f(r) - A| \to 0\)。\(\blacksquare\)

Tauber 第一定理的条件 \(a_n = o(1/n)\) 是一个典型的''Tauber 条件''------它限制了级数项的增长速度，从而排除了 Abel 可和但通常发散的那些病态情形。Hardy 和 Littlewood 进一步发现，如果将条件放宽为单侧有界性（如 \(a_n \geq 0\)），同样可以从 Abel 可和性推回通常收敛。这一发现是深刻的，因为它揭示了可和理论中的一条基本原则：正项级数的发散程度受到严格限制，Abel 可和性对它们而言与通常收敛等价。

\textbf{定理 100.4.2}（Hardy-Littlewood Tauber 定理）：若级数是 Abel 可和的且 \(a_n \geq 0\)（或 \(a_n = O(1/n)\)），则级数通常收敛。

\textbf{证明}：首先考虑 \(a_n \geq 0\) 的情形。设 Abel 和为 \(A\)，即 \(\lim_{r \to 1^-} \sum a_n r^n = A\)。由于 \(a_n \geq 0\)，\(f(r) = \sum a_n r^n\) 关于 \(r\) 单调递增。对任意 \(N\)，\(S_N = \sum_{n=0}^N a_n = \lim_{r \to 1^-} \sum_{n=0}^N a_n r^n \leq \lim_{r \to 1^-} \sum_{n=0}^\infty a_n r^n = A\)。故 \(\{S_N\}\) 单调递增且有上界 \(A\)，从而收敛于某 \(S \leq A\)。又对任意 \(r < 1\)，\(\sum a_n r^n \leq \sum a_n = S\)，令 \(r \to 1^-\) 得 \(A \leq S\)。故 \(S = A\)，级数通常收敛。对于 \(a_n = O(1/n)\) 的情形（即 \(|a_n| \leq C/n\)），将级数拆分为正部和负部，利用类似定理 100.4.1 的 Tauber 条件论证。核心是：Abel 可和性加上 \(a_n = O(1/n)\) 的一侧 Tauber 条件足以保证通常收敛。\(\blacksquare\)

此前的 Tauber 定理都是针对特定可和方法的''若 A 则 B''型命题。Wiener 在 1932 年完成的突破性工作将这些零散的结果统一为一个庞大的抽象框架------Wiener Tauber 定理。该定理不再关注具体的级数，而是讨论一般 \(L^1(\mathbb{R})\) 函数通过平移平移生成的闭理想何时等于全空间。这一抽象视角使 Tauber 型定理从级数论的技巧提升为 Banach 代数理论的结构性定理，其证明也标志着 Gelfand 表示理论在调和分析中的首次胜利。

\textbf{定理 100.4.3}（Wiener Tauber 定理------抽象版本）：设 \(f \in L^1(\mathbb{R})\)，其 Fourier 变换 \(\hat{f}\) 在 \(\mathbb{R}\) 上无处为零。则 \(f\) 的平移张成的闭子空间是整个 \(L^1(\mathbb{R})\)。Wiener Tauber 定理是 Tauber 定理的最终抽象形式，其证明依赖于 Banach 代数的 Gelfand 理论。

\textbf{证明}（框架）：考虑卷积代数 \(L^1(\mathbb{R})\)。\(f\) 的平移张成的闭子空间 \(I_f = \overline{\operatorname{span}}\{f(\cdot - t) : t \in \mathbb{R}\}\) 是 \(L^1(\mathbb{R})\) 的闭理想。由 Banach 代数理论，\(L^1(\mathbb{R})\) 的 Gelfand 变换将 \(L^1(\mathbb{R})\) 的极大理想与 \(\mathbb{R}\) 的紧化（即 \(\mathbb{R} \cup \{\infty\}\)）对应，且 \(f\) 的 Gelfand 变换正是 \(\hat{f}\)。若 \(\hat{f}\) 无处为零，则 \(f\) 不属于任何极大理想，故 \(I_f\) 不包含在任何真闭理想中。但 \(L^1(\mathbb{R})\) 的闭理想结构由 Wiener 定理刻画：若闭理想 \(I\) 不包含在任何极大理想中，则 \(I = L^1(\mathbb{R})\)。因此 \(I_f = L^1(\mathbb{R})\)。Wiener 的原始证明使用恒等逼近的构造，Nyman 和 Beurling 后来给出了基于 Banach 代数方法的简洁证明。\(\blacksquare\)

Wiener Tauber 定理的抽象之美最终通过其在解析数论中的惊人应用而臻于圆满。Wiener-Ikehara 定理是该抽象定理在 Dirichlet 级数框架下的具体化身，它仅从 Dirichlet 级数在边界上的解析延拓性质出发，就能推导出系数部分和的渐近公式------这恰好是素数定理所需的全部信息。事实上，复数域上的素数定理证明（Hadamard 和 de la Vallée Poussin, 1896）需要 \(\zeta(s)\) 在 \(\Re(s) = 1\) 上无零点的精细估计，而 Wiener-Ikehara 定理用一条统一的 Tauber 原理替代了这些分析细节，堪称 20 世纪数学中抽象方法战胜具体计算的典范。

\textbf{注记 100.4.1}：Tauber 定理在解析数论中有经典应用------素数定理的 Tauber 证明（Wiener-Ikehara 定理）是 Wiener Tauber 定理最著名的推论。Wiener-Ikehara 定理表明：若 Dirichlet 级数 \(\sum a_n n^{-s}\) 在 \(\Re(s) > 1\) 上收敛，且其和函数在 \(\Re(s) \geq 1\) 上除 \(s=1\) 处的简单极点外全纯，则 \(\sum_{n \leq x} a_n \sim c x\)。这一结果直接蕴含素数定理 \(\pi(x) \sim x/\log x\)。

\hrulefill

\section{第160章：位势论}\label{ux7b2c160ux7ae0ux4f4dux52bfux8bba}

位势论研究调和函数、超调和函数以及与之相关的测度、容量和细拓扑。它是复分析、偏微分方程和概率论的交汇点。位势论的起源可追溯到 Gauss 和 Green 在 19 世纪早期关于电磁学和引力位势的工作，而现代公理化位势论由 Brelot、Cartan、Choquet 和 Doob 在 20 世纪中期建立。

\subsection{101.1 Riesz 位势与 Newton 位势}\label{riesz-ux4f4dux52bfux4e0e-newton-ux4f4dux52bf}

\textbf{定义 101.1.1}（Riesz 位势）：设 \(\mu\) 为 \(\mathbb{R}^n\)（\(n \geq 2\)）上具有紧支集的 Radon 测度。对 \(0 < \alpha < n\)，\(\mu\) 的 \textbf{Riesz 位势} 定义为

\[U_\alpha^\mu(x) = \int_{\mathbb{R}^n} \frac{1}{|x - y|^{n-\alpha}} \, d\mu(y)\]

当 \(\alpha = 2\)，\(n \geq 3\) 时，\(U_2^\mu(x) = \int |x-y|^{2-n} d\mu(y)\) 称为 \textbf{Newton 位势}。当 \(\alpha = 1\)，\(U_1^\mu\) 称为 \textbf{Riesz 能量}。

\textbf{定理 101.1.1}（Riesz 位势的基本性质）：设 \(\mu\) 为具有紧支集的有限正 Radon 测度。 (i) \(U_\alpha^\mu\) 是下半连续的（在 \(\mathbb{R}^n\) 上），且在 \(\operatorname{supp}(\mu)\) 外是 \(C^\infty\) 的。 (ii) 若 \(\mu\) 绝对连续，\(d\mu = f dx\) 且 \(f \in L^p(\mathbb{R}^n)\)，则 \(U_\alpha^\mu \in L^q_{\text{loc}}(\mathbb{R}^n)\)（由 Hardy-Littlewood-Sobolev 不等式）。 (iii) 在分布意义下，\((-\Delta)^{s} U_{2s}^\mu = c_{n,s} \mu\)，其中 \((-\Delta)^s\) 是分数阶 Laplace 算子（\(0 < s < 1\)）。

\textbf{证明}：(i) 下半连续性：\(U_\alpha^\mu(x) = \int |x-y|^{\alpha-n} d\mu(y)\)。核 \(k(x,y) = |x-y|^{\alpha-n}\) 是 \((x,y)\) 的下半连续函数（因为 \(\alpha-n < 0\)，核在 \(x=y\) 处趋于 \(+\infty\)）。由 Fatou 引理，\(\liminf_{x \to x_0} U_\alpha^\mu(x) \geq \int \liminf_{x \to x_0} |x-y|^{\alpha-n} d\mu(y) = U_\alpha^\mu(x_0)\)。在 \(\operatorname{supp}(\mu)\) 外，\(x \mapsto |x-y|^{\alpha-n}\) 关于 \(x\) 是 \(C^\infty\) 的（\(x \neq y\)），积分与求导可交换（由控制收敛定理），故 \(U_\alpha^\mu\) 在 \(\operatorname{supp}(\mu)\) 外 \(C^\infty\)。(ii) 由定理 101.1.2（HLS 不等式），\(I_\alpha: L^p \to L^q\) 有界，故 \(U_\alpha^\mu \in L^q(\mathbb{R}^n)\)，从而 \(L^q_{\text{loc}}\)。(iii) 分数阶 Laplace 算子的 Fourier 符号为 \(|\xi|^{2s}\)，而 Riesz 位势 \(I_{2s}\) 的 Fourier 乘子为 \(|\xi|^{-2s}\)。在分布意义下，\((-\Delta)^s I_{2s} f = \mathcal{F}^{-1}(|\xi|^{2s} \cdot |\xi|^{-2s} \hat{f}) = f\)，故 \((-\Delta)^s U_{2s}^\mu = c_{n,s} \mu\)，其中 \(c_{n,s}\) 为归一化常数。\(\blacksquare\)

\textbf{定理 101.1.2}（Hardy-Littlewood-Sobolev 不等式）：设 \(0 < \alpha < n\)，\(1 < p < q < \infty\) 满足 \(\frac{1}{q} = \frac{1}{p} - \frac{\alpha}{n}\)。则 Riesz 位势算子 \(I_\alpha f(x) = \int \frac{f(y)}{|x-y|^{n-\alpha}} dy\) 是从 \(L^p(\mathbb{R}^n)\) 到 \(L^q(\mathbb{R}^n)\) 的有界算子。

\textbf{证明（框架）}：将 \(I_\alpha f\) 分解为近程和远程两部分：\(I_\alpha f(x) = \int_{|x-y| \leq R} + \int_{|x-y| > R}\)。近程部分由 Hardy-Littlewood 极大函数 \(Mf\) 控制：\(\int_{|x-y| \leq R} \frac{|f(y)|}{|x-y|^{n-\alpha}} dy \leq C R^\alpha Mf(x)\)。远程部分由 Holder 不等式控制：\(\int_{|x-y| > R} \frac{|f(y)|}{|x-y|^{n-\alpha}} dy \leq C R^{\alpha - n/p'} \|f\|_p\)。优化 \(R\) 的选取得 \(\|I_\alpha f\|_q \leq C \|f\|_p\)。\(\blacksquare\)

Hardy-Littlewood-Sobolev 不等式是调和分析中的基本不等式，它精确刻画了 Riesz 位势算子 \(I_\alpha\) 将 \(L^p\) 正则性提升 \(\alpha\) 阶 \(L^q\) 正则性的能力。在 \(\alpha = 2\) 时，\(I_2\) 是 \((-\Delta)^{-1}\) 的卷积形式，HLS 不等式给出了 Newton 位势在 \(L^p\) 空间中的有界性。在 \(\alpha \to n\) 的极限情形下，\(I_\alpha\) 趋向于对数位势，相应的不等式变为 Moser-Trudinger 不等式------这是 Sobolev 嵌入 \(W^{1,n} \hookrightarrow L^\infty\) 失效时的最佳替代。

\textbf{定义 101.1.2}（Bessel 位势与 Sobolev 空间刻画）：对任意 \(s \in \mathbb{R}\)，\textbf{Bessel 位势} \(J_s\) 通过 Fourier 乘子定义：\(\widehat{J_s f}(\xi) = (1 + |\xi|^2)^{-s/2} \hat{f}(\xi)\)。\(J_s\) 的卷积核为 \(G_s(x) = c_{n,s} \int_0^\infty t^{(s-n)/2} e^{-\pi |x|^2/t} e^{-t/4\pi} dt/t\)。

\textbf{定理 101.1.3}（Sobolev 空间的位势刻画）：对 \(s \in \mathbb{R}\) 和 \(1 < p < \infty\)，Sobolev 空间（或 Bessel 位势空间）\(H^{s,p}(\mathbb{R}^n)\) 定义为 \(H^{s,p}(\mathbb{R}^n) = \{f \in \mathcal{S}'(\mathbb{R}^n) : J_{-s} f \in L^p(\mathbb{R}^n)\}\)，配以范数 \(\|f\|_{H^{s,p}} = \|J_{-s} f\|_{L^p}\)。当 \(p = 2\) 时，\(H^{s,2} = W^{s,2}\) 为标准 \(L^2\)-Sobolev 空间。

\textbf{证明}：\(J_{-s}\) 的 Fourier 乘子为 \((1+|\xi|^2)^{s/2}\)。对于 \(f \in \mathcal{S}'\)，\(J_{-s} f \in L^p\) 等价于 \(\mathcal{F}^{-1}((1+|\xi|^2)^{s/2} \hat{f}) \in L^p\)。当 \(p=2\) 时，由 Plancherel 定理，\(\|f\|_{H^{s,2}}^2 = \int_{\mathbb{R}^n} (1+|\xi|^2)^s |\hat{f}(\xi)|^2 d\xi\)，这正是 \(L^2\)-Sobolev 空间 \(W^{s,2}\) 的 Fourier 刻画。对整数 \(s = k \in \mathbb{N}\)，由 Mihlin-Hormander 乘子定理，\((1+|\xi|^2)^{k/2}\) 是 \(L^p\) Fourier 乘子（\(1 < p < \infty\)），故 \(H^{k,p} = W^{k,p}\)（经典与位势定义的等价性）。对一般的 \(s \in \mathbb{R}\)，Bessel 位势空间 \(H^{s,p}\) 通过复插值或实插值与经典 Sobolev 空间联系。\(\blacksquare\)

\hrulefill

\subsection{101.2 容量理论}\label{ux5bb9ux91cfux7406ux8bba}

\textbf{定义 101.2.1}（Wiener 容量 / Newton 容量）：设 \(K \subseteq \mathbb{R}^n\)（\(n \geq 3\)）为紧集。\(K\) 的 \textbf{Newton 容量}（或 Wiener 容量）定义为

\[\operatorname{Cap}(K) = \sup\{\mu(K) : \mu \text{ 是 } K \text{ 上的正 Radon 测度, } U_2^\mu(x) \leq 1 \text{ 对所有 } x \in \mathbb{R}^n\}\]

等价地，\(\operatorname{Cap}(K) = \inf\{\int_{\mathbb{R}^n} |\nabla u|^2 dx : u \in C_c^\infty(\mathbb{R}^n), u \geq 1 \text{ 在 } K \text{ 上}\}\)（Dirichlet 能量形式）。

容量理论是位势论的核心，它量化了集合的''大小''------从调和函数和位势论的角度，而非经典测度论的角度。容量是 Choquet 容量的一个特例------Choquet 在 1954 年引入的公理化容量理论将 Newton 容量、对数容量、Riesz 容量等统一在一个框架中。Choquet 容量的定义是：一个从 \(\mathbb{R}^n\) 的全体子集到 \([0, \infty]\) 的集函数 \(C\)，满足 (i) 单调性：\(A \subseteq B \Rightarrow C(A) \leq C(B)\)；(ii) 右连续性：对递增紧集列 \(K_n \uparrow K\)，\(C(K_n) \uparrow C(K)\)；(iii) 对任意集合 \(E\)，\(C(E) = \sup\{C(K) : K \subseteq E \text{ 紧}\}\)。Choquet 容量的关键性质是：所有 Borel 集（乃至所有 Souslin 集）都是可容的（capacitable），即 \(C(E) = \sup\{C(K) : K \subseteq E \text{ 紧}\}\)。这一结果由 Choquet 的著名定理保证，在位势论、随机过程和描述集合论中有广泛应用。

\textbf{定理 101.2.1}（Wiener 判别法，1924）：设 \(\Omega \subseteq \mathbb{R}^n\) 为有界区域，\(x_0 \in \partial \Omega\)。则 \(x_0\) 是 Dirichlet 问题的正则边界点（即对任意连续边界值，Perron 解在 \(x_0\) 处连续等于边界值）当且仅当

\[\sum_{k=1}^\infty \frac{\operatorname{Cap}(B(x_0, 2^{-k}) \setminus \Omega)}{2^{-k(n-2)}} = \infty\]

（当 \(n=2\) 时，容量用对数容量，分母为 \(k \log 2\)。）

\textbf{证明}（框架）：Wiener 判别法的核心是调和测度与容量之间的对偶关系。边界点 \(x_0\) 正则当且仅当 \(x_0\) 的任意邻域中的补集具有正容量（即 \(x_0\) 不是''薄的''）。Wiener 级数发散条件等价于 \(x_0\) 处补集的容量在尺度 \(2^{-k}\) 下不是''太小''------即边界在该点附近''足够厚''以支撑调和测度。\(\blacksquare\)

Wiener 判别法是对 Dirichlet 问题正则边界点的完全刻画，是位势论中最深刻的定理之一。其几何推论是：如果 \(\partial\Omega\) 在 \(x_0\) 处满足''外锥条件''（即 \(x_0\) 处存在一个顶点在 \(x_0\) 的锥体完全在 \(\Omega\) 外），则 \(x_0\) 正则。特别地，\(C^1\) 域的所有边界点都是正则的。Wiener 判别法还揭示了''尖点''（cusp）可以是正则的也可以是非正则的，取决于尖点的''尖锐度''------这由 Wiener 级数中的容量项精确刻画。

\textbf{定理 101.2.2}（Kellogg-Evans 定理与容量的连续性）：设 \(\{K_n\}\) 为递减紧集列，\(K = \bigcap K_n\)。则 \(\operatorname{Cap}(K) = \lim_{n \to \infty} \operatorname{Cap}(K_n)\)。这表明 Newton 容量是 Choquet 容量。

\textbf{证明}：单调性给出 \(\operatorname{Cap}(K) \leq \lim \operatorname{Cap}(K_n)\)。反之，设 \(\mu_n\) 为 \(K_n\) 上满足 \(U_2^{\mu_n} \leq 1\) 的测度，\(\mu_n(K_n) \to \lim \operatorname{Cap}(K_n)\)。取弱收敛子列 \(\mu_n \rightharpoonup \mu\)，则 \(\mu\) 支撑在 \(K\) 上，由下半连续性 \(U_2^\mu \leq \liminf U_2^{\mu_n} \leq 1\)，故 \(\operatorname{Cap}(K) \geq \mu(K) = \lim \mu_n(K_n) = \lim \operatorname{Cap}(K_n)\)。\(\blacksquare\)

\hrulefill

\subsection{101.3 细拓扑与薄集}\label{ux7ec6ux62d3ux6251ux4e0eux8584ux96c6}

\textbf{定义 101.3.1}（细拓扑 / Fine Topology）：\(\mathbb{R}^n\) 上的\textbf{细拓扑}是使所有超调和函数（superharmonic functions）连续的最粗拓扑。它是比欧氏拓扑严格更精细的拓扑。细拓扑由基 \(\{B : B \text{ 是欧氏开集盘}\} \cup \{U \setminus E : U \text{ 欧氏开}, \operatorname{Cap}(E) = 0\}\) 生成。

\textbf{定义 101.3.2}（薄集 / Thin Set）：集合 \(E \subseteq \mathbb{R}^n\) 在点 \(x_0 \in \mathbb{R}^n\) 处称为\textbf{薄的}，如果 \(x_0 \notin \overline{E}^{\text{fine}}\)（即 \(x_0\) 不是 \(E\) 在细拓扑下的聚点）。等价地，若 \(x_0 \notin E\) 时，\(E\) 在 \(x_0\) 处薄当且仅当存在 \(x_0\) 附近的超调和函数 \(u\) 使得 \(\liminf_{x \to x_0, x \in E} u(x) > u(x_0)\)。

\textbf{定理 101.3.1}（薄集的 Wiener 判别法）：\(E\) 在 \(x_0\) 处薄当且仅当

\[\sum_{k=1}^\infty \frac{\operatorname{Cap}(E \cap B(x_0, 2^{-k}))}{2^{-k(n-2)}} < \infty\]

（\(n=2\) 时使用对数容量修改分母）。

\textbf{证明}：薄集的 Wiener 判别法是定理 101.2.1（边界正则性 Wiener 判别法）的对偶版本。由定义，\(E\) 在 \(x_0\) 处薄等价于 \(x_0\) 不是 \(E\) 的细聚点，即存在 \(x_0\) 处的细邻域 \(U\) 使 \(U \cap E = \{x_0\}\)（或 \(U \cap E = \emptyset\) 若 \(x_0 \notin E\)）。这一条件等价于 \(E\) 在 \(x_0\) 处的''补集''具有正 Wiener 容量。具体地，构造递减紧集列 \(K_k = \overline{E} \cap \overline{B(x_0, 2^{-k})} \setminus B(x_0, 2^{-k-1})\)。\(E\) 在 \(x_0\) 处薄当且仅当 \(\sum_{k} 2^{k(n-2)} \operatorname{Cap}(K_k) < \infty\)。由容量的可数次可加性（在紧集上），\(\operatorname{Cap}(E \cap B(x_0, 2^{-k})) \leq \sum_{j=k}^\infty \operatorname{Cap}(K_j)\)，且反之 \(\operatorname{Cap}(K_k) \leq \operatorname{Cap}(E \cap B(x_0, 2^{-k}))\)。两者结合，级数收敛的等价性得证。\(\blacksquare\)

细拓扑和薄集理论是位势论中最精致的部分，由 Brelot、Cartan 和 Choquet 在 1940-50 年代发展。细拓扑不是可度量的（尽管它是 Hausdorff 的），但它具有 Baire 性质，且细连续函数恰好是那些在细拓扑下连续的函数。Ball-Poisson 积分 \(P[f](x)\) 在细拓扑边界下连续这一事实（Fatou 定理的细拓扑推广）是细拓扑最经典的应用之一。在概率论中，细拓扑自然地出现在 Brown 运动的首达时分布中：集合 \(E\) 在 \(x_0\) 处薄当且仅当从 \(x_0\) 出发的 Brown 运动几乎必然不立即击中 \(E\)------这一联系由 Kakutani 在 1944 年发现，是位势论与概率论之间深层对应的核心。

\textbf{定理 101.3.2}（细拓扑的拟-Lindelof 性质）：细拓扑满足拟-Lindelof 性质：任意细开集族构成的覆盖有可数子覆盖，模去一个容量为零的集合。

\textbf{证明}：设 \(\{U_i\}_{i \in I}\) 为细开覆盖。由细拓扑的定义，每个 \(U_i\) 可写为 \(U_i = V_i \setminus E_i\)，其中 \(V_i\) 为欧氏开集，\(\operatorname{Cap}(E_i) = 0\)。令 \(E = \bigcup_i E_i\)，则 \(\operatorname{Cap}(E) = 0\)（容量的可数次可加性）。在 \(\mathbb{R}^n \setminus E\) 上，每个 \(U_i\) 包含欧氏开集 \(V_i\)，故 \(\{V_i\}\) 是欧氏拓扑下 \(\mathbb{R}^n \setminus E\) 的覆盖。由欧氏拓扑的 Lindelof 性质（\(\mathbb{R}^n\) 是第二可数的），存在可数子覆盖 \(\{V_{i_k}\}\) 覆盖 \(\mathbb{R}^n \setminus E\)。对应的 \(\{U_{i_k}\}\) 覆盖 \(\mathbb{R}^n \setminus E\)。因此 \(\{U_{i_k}\}\) 覆盖全空间模去 \(E\)（容量为零）。\(\blacksquare\)

\hrulefill

\subsection{101.4 位势论与概率论的联系}\label{ux4f4dux52bfux8bbaux4e0eux6982ux7387ux8bbaux7684ux8054ux7cfb}

\textbf{定理 101.4.1}（Kakutani 定理，1944）：设 \(B_t\) 为 \(\mathbb{R}^n\)（\(n \geq 3\)）中的标准 Brown 运动。则 \(B_t\) 击中集合 \(E \subseteq \mathbb{R}^n\) 的概率（首达概率）与 \(E\) 的 Newton 容量密切相关：\(\mathbb{P}^x(B_t \in E \text{ 对某个 } t > 0)\) 与 \(E\) 的容量在 \(x\) 处的调和测度成正比。更精确地，存在常数 \(c_n > 0\) 使得对任意紧集 \(K\)，

\[\mathbb{P}^x(\tau_K < \infty) = c_n \int_K |x-y|^{2-n} d\mu_K(y)\]

其中 \(\mu_K\) 是 \(K\) 的平衡测度（满足 \(\mu_K(K) = \operatorname{Cap}(K)\) 且 \(U_2^{\mu_K} \equiv 1\) 在 \(K\) 上 quasi-everywhere 的测度）。

\textbf{证明}：由调和函数的概率表示，\(u(x) = \mathbb{P}^x(\tau_K < \infty)\) 是 \(\mathbb{R}^n \setminus K\) 上的调和函数，边界值为 \(1\) 在 \(K\) 上，且 \(u(\infty) = 0\)。由位势论，\(u\) 恰好是 \(K\) 的平衡位势 \(U_2^{\mu_K}(x)\)（归一化后）。由平衡测度的定义，\(\mu_K(K) = \operatorname{Cap}(K)\)，且 \(U_2^{\mu_K} \equiv 1\) 在 \(K\) 上 quasi-everywhere。因此 \(\mathbb{P}^x(\tau_K < \infty) = U_2^{\mu_K}(x)\)。\(\blacksquare\)

Kakutani 定理建立了位势论与概率论之间的基本桥梁，开启了现代位势论的概率方法。这一联系的核心是：调和函数、超调和函数、位势、容量等位势论概念在概率论中都有自然的对应物------调和函数对应于 Brown 运动的期望，超调和函数对应于上鞅，位势对应于 Green 函数，容量对应于首达概率。这一对应由 Doob 和 Hunt 在 1950 年代发展为完整的 Hunt 过程位势论，将经典 Newton 位势论推广到一般 Markov 过程的框架中。在 Dirichlet 型的理论中（Fukushima 1980），位势论、Dirichlet 形式和对称 Markov 过程三者构成了三位一体的结构，是现代随机分析和无穷维分析的基础。

\subsection{101.5 平衡测度与能量积分}\label{ux5e73ux8861ux6d4bux5ea6ux4e0eux80fdux91cfux79efux5206}

\textbf{定义 101.5.1}（能量积分）：设 \(\mu\) 为 \(\mathbb{R}^n\) 上的正 Radon 测度。\(\mu\) 的 \(\alpha\)-能量定义为

\[
I_\alpha(\mu) = \iint \frac{1}{|x-y|^{n-\alpha}} \, d\mu(x) d\mu(y) = \int U_\alpha^\mu(x) d\mu(x)
\]

能量积分 \(I_\alpha(\mu)\) 衡量测度 \(\mu\) 在 Riesz 核作用下的''自相互作用''强度。这一概念直接源于物理学------在静电学中，\(I_2(\mu)\) 正是电荷分布 \(\mu\) 的静电能。极小化能量积分的问题因此具有自然的物理意义：给定一个导体 \(K\)，电荷如何分布才能使总静电能最小？下面的定理表明，这样的极小化分布（平衡测度）不仅存在唯一，而且给出了容量的精确表达式------从而将抽象的容量概念与具体的变分问题联系起来。

\textbf{定理 101.5.1}（能量积分的最小化与平衡测度）：对紧集 \(K\)，存在唯一的正 Radon 测度 \(\mu_K\)（称为 \(K\) 的平衡测度）支撑在 \(K\) 上，最小化能量积分 \(I_\alpha(\mu)\)（在所有支撑在 \(K\) 上的概率测度中），且满足 \(U_\alpha^{\mu_K}(x) \geq I_\alpha(\mu_K)\) 在 \(K\) 上 quasi-everywhere。此时 \(\operatorname{Cap}_\alpha(K) = 1/I_\alpha(\mu_K)\)。

\textbf{证明}：设 \(\mathcal{P}(K)\) 为支撑在 \(K\) 上的概率测度集合。在弱拓扑下，\(\mathcal{P}(K)\) 是紧集（由 Prokhorov 定理）。能量积分 \(\mu \mapsto I_\alpha(\mu)\) 在弱拓扑下是下半连续的：由 Fatou 引理，若 \(\mu_n \rightharpoonup \mu\)，则 \(I_\alpha(\mu) \leq \liminf I_\alpha(\mu_n)\)。因此 \(I_\alpha\) 在紧集 \(\mathcal{P}(K)\) 上达到下确界，即存在 \(\mu_K \in \mathcal{P}(K)\) 使 \(I_\alpha(\mu_K) = \inf_{\mu \in \mathcal{P}(K)} I_\alpha(\mu)\)。唯一性：若 \(\mu_1, \mu_2\) 均为极小化元，则 \(\mu = (\mu_1 + \mu_2)/2\) 的能量满足 \(I_\alpha(\mu) \leq I_\alpha(\mu_1) = I_\alpha(\mu_2)\)，由凸性得等号仅当 \(\mu_1 = \mu_2\)。平衡条件 \(U_\alpha^{\mu_K} \geq I_\alpha(\mu_K)\) quasi-everywhere 在 \(K\) 上由变分原理导出：对任意 \(\nu \in \mathcal{P}(K)\)，\(I_\alpha((1-t)\mu_K + t\nu) \geq I_\alpha(\mu_K)\) 对 \(t \in [0,1]\) 成立，展开得 \(\int U_\alpha^{\mu_K} d\nu \geq I_\alpha(\mu_K)\)。由 \(\nu\) 的任意性得逐点下界。容量公式 \(\operatorname{Cap}_\alpha(K) = 1/I_\alpha(\mu_K)\) 由容量定义直接验证。\(\blacksquare\)

平衡测度的存在性提供了容量的变分刻画，但在实际应用中，我们常常需要在给定紧集上构造满足特定位势上界约束的测度。Frostman 引理恰好满足了这一需求：它断言只要容量为正，就可以在紧集上''挤入''足够质量的测度，同时保持位势不超过 \(1\)。这一引理看似技术性，实则是连接容量理论与几何测度论的关键枢纽------无论是 Hausdorff 维数的下界估计还是分形集的 Frostman 引理刻画，都以它为基础。

\textbf{定理 101.5.2}（Frostman 引理）：设 \(K\) 为紧集，\(\operatorname{Cap}_\alpha(K) > 0\)。则存在支撑在 \(K\) 上的正 Radon 测度 \(\mu\) 使得 \(U_\alpha^\mu(x) \leq 1\) 对所有 \(x\) 成立，且 \(\mu(K) \geq \operatorname{Cap}_\alpha(K)\)。Frostman 引理是容量理论中构造测度的基本工具，在几何测度论和分形几何中有广泛应用。

\textbf{证明}：由容量定义，\(\operatorname{Cap}_\alpha(K) = \sup\{\mu(K) : \mu \in \mathcal{M}^+(K), U_\alpha^\mu \leq 1\}\)。取极大化序列 \(\{\mu_n\}\) 满足 \(U_\alpha^{\mu_n} \leq 1\) 且 \(\mu_n(K) \to \operatorname{Cap}_\alpha(K)\)。由 Banach-Alaoglu 定理，存在弱收敛子列 \(\mu_n \rightharpoonup \mu\)。由下半连续性，\(U_\alpha^\mu \leq \liminf U_\alpha^{\mu_n} \leq 1\)。又 \(\mu(K) = \lim \mu_n(K) = \operatorname{Cap}_\alpha(K)\)（因为紧集上弱收敛蕴含总质量收敛）。故 \(\mu\) 满足所需条件。\(\blacksquare\)

\subsection{101.6 对数位势与二维位势论}\label{ux5bf9ux6570ux4f4dux52bfux4e0eux4e8cux7ef4ux4f4dux52bfux8bba}

\textbf{定义 101.6.1}（对数位势）：在 \(\mathbb{R}^2\) 中，Newton 核 \(|x-y|^{2-n}\) 退化为常数，故需用对数核代替。测度 \(\mu\) 的对数位势定义为

\[
U_{\log}^\mu(x) = \int \log \frac{1}{|x-y|} \, d\mu(y)
\]

在 \(\mathbb{R}^n\)（\(n \geq 3\)）中，Newton 核 \(|x-y|^{2-n}\) 在位势理论中处于中心地位。但当 \(n = 2\) 时，该核的指数退化为零，取而代之的是对数核 \(\log(1/|x-y|)\)。这一维度差异不仅仅是一个技术细节------它深刻地改变了位势论在二维空间中的行为。例如，在 \(\mathbb{R}^n\)（\(n \geq 3\)）中，单点集具有正 Newton 容量；而在 \(\mathbb{R}^2\) 中，单点集的对数容量为零。这种差异的根源在于 Brown 运动的常返性/暂留性分界恰好位于 \(n = 2\)。

\textbf{定理 101.6.1}（对数容量）：对紧集 \(K \subseteq \mathbb{R}^2\)，对数容量 \(\operatorname{Cap}_{\log}(K)\) 定义为 \(e^{-V(K)}\)，其中 \(V(K) = \inf\{I_{\log}(\mu) : \mu \text{ 是概率测度支撑在 } K\}\)。对数容量在复分析中至关重要------全纯函数在孤立奇点附近的行为、整函数的增长、以及 Padé 逼近的收敛性都与对数容量密切相关。

\textbf{证明}：对数能量 \(I_{\log}(\mu) = \iint \log(1/|x-y|) d\mu(x) d\mu(y)\)。由于 \(\log(1/|x-y|)\) 在 \(x=y\) 处趋于 \(+\infty\)（非负），\(I_{\log}(\mu)\) 可能为 \(+\infty\)。若对某概率测度 \(\mu\) 有 \(I_{\log}(\mu) < \infty\)，则 \(V(K) = \inf I_{\log}(\mu) < \infty\)，此时 \(\operatorname{Cap}_{\log}(K) = e^{-V(K)} > 0\)。若对所有概率测度 \(I_{\log}(\mu) = \infty\)，则 \(V(K) = \infty\)，\(\operatorname{Cap}_{\log}(K) = 0\)。\(V(K)\) 的极小化元的存在唯一性证明与定理 101.5.1 类似（利用弱拓扑下的下半连续性和凸性）。对数容量的关键性质：\(K_1 \subseteq K_2 \Rightarrow \operatorname{Cap}_{\log}(K_1) \leq \operatorname{Cap}_{\log}(K_2)\)；\(\operatorname{Cap}_{\log}(cK) = c \cdot \operatorname{Cap}_{\log}(K)\)（齐次性）。\(\blacksquare\)

对数容量起初是通过变分问题（极小化能量积分）定义的，这虽然在理论上自然，但在计算上不够直观。令人惊叹的是，对数容量有一个与之等价的纯几何刻画------跨导直径（transfinite diameter）。这一等价性由 Fekete 和 Szegő 建立，它将对数容量与多项式插值点的极值问题联系起来：容量越大的集合，能容纳更多''分散''的点。下面定理的证明中，Fekete 引理（关于次可加序列）的应用堪称经典。

\textbf{定理 101.6.2}（对数容量的几何意义 / 对数容量的 Chebyshev 常数刻画）：对紧集 \(K \subseteq \mathbb{C}\)，对数容量 \(\operatorname{Cap}_{\log}(K)\) 等于 \(K\) 的跨导直径（transfinite diameter）：\(\operatorname{Cap}_{\log}(K) = \lim_{n \to \infty} \max_{z_1,\ldots,z_n \in K} \prod_{i<j} |z_i - z_j|^{2/(n(n-1))}\)。这一几何刻画将对数容量与经典的 Fekete 点（最小化多项式极值的点）联系起来，在多项式逼近论中有重要应用。特别地，对区间 \([-1, 1]\)，对数容量为 \(1/2\)；对圆盘 \(|z| \leq R\)，对数容量为 \(R\)。

\textbf{证明}：定义 \(n\) 阶直径 \(d_n(K) = \max_{z_1,\ldots,z_n \in K} \prod_{i<j} |z_i - z_j|^{2/(n(n-1))}\)。跨导直径定义为 \(d_\infty(K) = \lim_{n \to \infty} d_n(K)\)（极限存在性由 Fekete 引理保证：\(d_n(K)\) 关于 \(n\) 单调递减）。需证 \(d_\infty(K) = \operatorname{Cap}_{\log}(K)\)。一方面，设 \(\mu_K\) 为 \(K\) 的平衡测度，则 \(I_{\log}(\mu_K) = V(K) = -\log \operatorname{Cap}_{\log}(K)\)。由位势论，\(U_{\log}^{\mu_K}(z) \geq V(K)\) quasi-everywhere 在 \(K\) 上。由此可导出 \(d_n(K) \geq \operatorname{Cap}_{\log}(K)\) 对所有 \(n\)。另一方面，Fekete 点（使乘积达到最大的 \(n\) 个点）的经验测度 \(\mu_n = \frac{1}{n} \sum_{i=1}^n \delta_{z_i}\) 弱收敛于平衡测度 \(\mu_K\)。由能量积分的下半连续性，\(\liminf I_{\log}(\mu_n) \geq I_{\log}(\mu_K)\)。而 \(I_{\log}(\mu_n)\) 与 \(d_n(K)\) 有关，推导得 \(\limsup d_n(K) \leq \operatorname{Cap}_{\log}(K)\)。两者结合得等号。\(\blacksquare\)

\newpage

\chapter{06-分析/02-复分析/02-级数与留数.md}\label{ux5206ux679002-ux590dux5206ux679002-ux7ea7ux6570ux4e0eux7559ux6570.md}

\section{第161章：级数展开}\label{ux7b2c161ux7ae0ux7ea7ux6570ux5c55ux5f00}

全纯函数的级数展开理论是复分析区别于实分析的最显著特征之一。本章系统阐述 Taylor 级数展开（全纯函数在圆盘上的幂级数表示）、Laurent 级数展开（在圆环上含负幂项的展开）以及由此导出的孤立奇点分类理论。级数展开不仅提供了全纯函数的局部显式表示，还自然地引出了整函数与亚纯函数的因子分解定理（Weierstrass 因子分解和 Mittag-Leffler 定理），为后续的留数计算和解析延拓理论奠定了坚实基础。

\subsection{53.1 Taylor 级数展开}\label{taylor-ux7ea7ux6570ux5c55ux5f00}

\textbf{定理 53.1.1（Taylor 展开定理）}：设 \(f\) 在 \(\Omega\) 上全纯，\(z_0 \in \Omega\)。设 \(R = \operatorname{dist}(z_0, \partial\Omega)\)（\(z_0\) 到 \(\Omega\) 边界的距离）。则 \(f\) 在 \(D(z_0, R)\) 上可展开为收敛的 Taylor 级数：

\[
f(z) = \sum_{n=0}^{\infty} a_n (z - z_0)^n, \quad a_n = \frac{f^{(n)}(z_0)}{n!}, \quad |z - z_0| < R
\]

\textbf{证明}：对 \(|z - z_0| < r < R\)，由 Cauchy 积分公式，

\[
f(z) = \frac{1}{2\pi i} \int_{\partial D(z_0, r)} \frac{f(\zeta)}{\zeta - z} d\zeta = \frac{1}{2\pi i} \int_{\partial D(z_0, r)} \frac{f(\zeta)}{(\zeta - z_0) - (z - z_0)} d\zeta
\]

展开几何级数：\(\frac{1}{(\zeta - z_0) - (z - z_0)} = \frac{1}{\zeta - z_0} \cdot \frac{1}{1 - \frac{z - z_0}{\zeta - z_0}} = \frac{1}{\zeta - z_0} \sum_{n=0}^{\infty} \left(\frac{z - z_0}{\zeta - z_0}\right)^n\)。由于 \(|z - z_0| < |\zeta - z_0| = r\)，级数一致收敛（关于 \(\zeta\)），可逐项积分：

\[
f(z) = \sum_{n=0}^{\infty} \left(\frac{1}{2\pi i} \int_{\partial D(z_0, r)} \frac{f(\zeta)}{(\zeta - z_0)^{n+1}} d\zeta\right) (z - z_0)^n = \sum_{n=0}^{\infty} \frac{f^{(n)}(z_0)}{n!} (z - z_0)^n
\]

由 Cauchy 高阶导数公式，系数确为 \(f^{(n)}(z_0)/n!\)。\(\blacksquare\)

\textbf{推论 53.1.2（Liouville 定理）}：若 \(f\) 在整个 \(\mathbb{C}\) 上全纯且有界（即存在 \(M\) 使得 \(|f(z)| \leq M\)，\(\forall z \in \mathbb{C}\)），则 \(f\) 为常数。

\textbf{证明}：由 Cauchy 估计（推论 52.3.4），对任意 \(R > 0\)，

\[
|f'(z_0)| = \left|\frac{1}{2\pi i} \int_{\partial D(z_0, R)} \frac{f(\zeta)}{(\zeta - z_0)^2} d\zeta\right| \leq \frac{1}{2\pi} \cdot \frac{M}{R^2} \cdot 2\pi R = \frac{M}{R}
\]

令 \(R \to \infty\) 得 \(f'(z_0) = 0\)。\(f' \equiv 0\) 意味着 \(f\) 为常数。\(\blacksquare\)

\textbf{推论 53.1.3（代数基本定理的复分析证明）}：任何非常数复系数多项式 \(P(z) = a_n z^n + \cdots + a_0\)（\(a_n \neq 0\)，\(n \geq 1\)）在 \(\mathbb{C}\) 中有根。

\textbf{证明}：假设 \(P(z) \neq 0\)（\(\forall z \in \mathbb{C}\)）。则 \(f(z) = 1/P(z)\) 在整个 \(\mathbb{C}\) 上全纯。当 \(|z| \to \infty\) 时，\(|P(z)| \to \infty\)，故 \(|f(z)| \to 0\)，从而 \(f\) 有界。由 Liouville 定理，\(f\) 为常数，与 \(n \geq 1\) 矛盾。\(\blacksquare\)

\textbf{定理 53.1.4（Taylor 展开的唯一性）}：若 \(\sum_{n=0}^{\infty} a_n (z - z_0)^n = \sum_{n=0}^{\infty} b_n (z - z_0)^n\) 在 \(z_0\) 的某邻域内成立，则 \(a_n = b_n\)（对所有 \(n\)）。由此，全纯函数在一点的 Taylor 系数唯一确定，进而由解析延拓，全纯函数由其在任意开子集上的值唯一确定。

\textbf{定理 53.1.5（零点与恒等定理）}：设 \(f\) 在区域 \(\Omega\) 上全纯。若 \(f\) 的零点集 \(\{z \in \Omega : f(z) = 0\}\) 在 \(\Omega\) 内有聚点，则 \(f \equiv 0\)。等价地，全纯函数在其定义域内由某点处的 Taylor 系数列（或某聚点邻域内的值）唯一确定。

\textbf{证明}：设 \(z_0\) 为零点集的聚点。由连续性，\(f(z_0) = 0\)。若 \(f\) 不恒为零，则存在 \(m \geq 1\) 使 \(f(z) = (z - z_0)^m g(z)\)，其中 \(g\) 全纯且 \(g(z_0) \neq 0\)（零点的孤立性）。但 \(z_0\) 是零点集的聚点，矛盾。故 \(f \equiv 0\)。\(\blacksquare\)

\hrulefill

\subsection{53.2 Laurent 级数展开}\label{laurent-ux7ea7ux6570ux5c55ux5f00}

\textbf{定理 53.2.1（Laurent 展开定理）}：设 \(f\) 在圆环 \(A = \{z \in \mathbb{C} \mid r < |z - z_0| < R\}\)（\(0 \leq r < R \leq \infty\)）上全纯。则 \(f\) 在 \(A\) 上可展开为 Laurent 级数：

\[
f(z) = \sum_{n=-\infty}^{\infty} a_n (z - z_0)^n, \quad r < |z - z_0| < R
\]

其中 \(a_n = \frac{1}{2\pi i} \int_{\partial D(z_0, \rho)} \frac{f(\zeta)}{(\zeta - z_0)^{n+1}} d\zeta\)（\(r < \rho < R\)，积分与 \(\rho\) 无关）。

\textbf{证明}：对 \(z \in A\)，取 \(r < r_1 < |z - z_0| < r_2 < R\)。由 Cauchy 积分公式（应用于环形区域），

\[
f(z) = \frac{1}{2\pi i} \int_{\partial D(z_0, r_2)} \frac{f(\zeta)}{\zeta - z} d\zeta - \frac{1}{2\pi i} \int_{\partial D(z_0, r_1)} \frac{f(\zeta)}{\zeta - z} d\zeta
\]

对第一个积分，\(|z - z_0| < |\zeta - z_0|\)，展开 \(\frac{1}{\zeta - z} = \sum_{n=0}^{\infty} \frac{(z - z_0)^n}{(\zeta - z_0)^{n+1}}\)。

对第二个积分，\(|z - z_0| > |\zeta - z_0|\)，展开 \(-\frac{1}{\zeta - z} = \frac{1}{z - \zeta} = \sum_{n=1}^{\infty} \frac{(\zeta - z_0)^{n-1}}{(z - z_0)^n}\)。

将两部分合并即得 Laurent 级数。\(\blacksquare\)

\textbf{注记 53.2.1}：Laurent 级数是全纯函数在奇点附近的基本工具。负幂部分（\(n < 0\)）称为 Laurent 级数的\textbf{主要部分}（principal part），它刻画了奇点的性质。正幂部分（\(n \geq 0\)）称为\textbf{解析部分}。Laurent 展开的唯一性：在同一圆环内，Laurent 展开是唯一的。

\textbf{证明}（Laurent展开的唯一性）设 \(f(z) = \sum_{n=-\infty}^{\infty} a_n (z-z_0)^n = \sum_{n=-\infty}^{\infty} b_n (z-z_0)^n\) 在圆环 \(r < |z-z_0| < R\) 中收敛。任取 \(\rho \in (r, R)\)，对固定整数 \(m\)，乘以 \((z-z_0)^{-m-1}\) 后沿 \(|z-z_0|=\rho\) 积分： \[\int_{|z-z_0|=\rho} \sum_{n} a_n (z-z_0)^{n-m-1} dz = \sum_{n} a_n \int_{|z-z_0|=\rho} (z-z_0)^{n-m-1} dz\] 由于 \(\int (z-z_0)^k dz = 2\pi i \delta_{k,-1}\)（仅当 \(k=-1\) 时积分为 \(2\pi i\)，否则为 \(0\)），等式中仅 \(n=m\) 的项幸存，得 \(a_m = \frac{1}{2\pi i} \int_{|z-z_0|=\rho} f(z)(z-z_0)^{-m-1} dz\)。对 \(b_m\) 同样的积分给出相同结果，故 \(a_m = b_m\)（对所有 \(m \in \mathbb{Z}\)）。\(\blacksquare\)

\hrulefill

\subsection{53.3 孤立奇点的分类}\label{ux5b64ux7acbux5947ux70b9ux7684ux5206ux7c7b}

\textbf{定义 53.3.1（孤立奇点）}：设 \(f\) 在 \(z_0\) 的去心邻域 \(\{0 < |z - z_0| < \delta\}\) 上全纯，但在 \(z_0\) 处不定义或不解析。则 \(z_0\) 称为 \(f\) 的\textbf{孤立奇点}。设 \(f\) 在 \(z_0\) 处的 Laurent 展开为 \(f(z) = \sum_{n=-\infty}^{\infty} a_n (z - z_0)^n\)。

\begin{enumerate}
\def\labelenumi{\arabic{enumi}.}
\tightlist
\item
  \textbf{可去奇点}：\(a_n = 0\)（\(\forall n < 0\)），即 Laurent 级数实为 Taylor 级数
\item
  \textbf{极点}：存在 \(m \geq 1\) 使得 \(a_{-m} \neq 0\) 且 \(a_n = 0\)（\(\forall n < -m\)）。\(m\) 称为\textbf{极点的阶}
\item
  \textbf{本性奇点}：存在无穷多个负指数 \(n\) 使得 \(a_n \neq 0\)
\end{enumerate}

\textbf{定理 53.3.1（Riemann 可去奇点定理）}：若 \(f\) 在 \(z_0\) 的去心邻域上全纯且有界，则 \(z_0\) 是可去奇点。

\textbf{证明}：设 \(|f(z)| \leq M\)。\(\forall n < 0\)，\(|a_n| = |\frac{1}{2\pi i} \int_{\partial D(z_0, \varepsilon)} \frac{f(\zeta)}{(\zeta - z_0)^{n+1}} d\zeta| \leq \frac{1}{2\pi} \cdot \frac{M}{\varepsilon^{n+1}} \cdot 2\pi\varepsilon = M \varepsilon^{-n}\)。令 \(\varepsilon \to 0\)（\(n < 0\) 时 \(-n > 0\)），得 \(a_n = 0\)。\(\blacksquare\)

\textbf{定理 53.3.2（极点的刻画）}：\(z_0\) 是 \(f\) 的 \(m\) 阶极点当且仅当 \(\lim_{z \to z_0} (z - z_0)^m f(z) = c \neq 0\)（有限且非零），等价于 \(|f(z)| \to \infty\)（\(z \to z_0\)）。

\textbf{定理 53.3.3（Casorati-Weierstrass 定理）}：若 \(z_0\) 是 \(f\) 的本性奇点，则对任意 \(w \in \mathbb{C} \cup \{\infty\}\)，存在序列 \(z_n \to z_0\) 使得 \(f(z_n) \to w\)。即 \(f\) 在 \(z_0\) 的任何去心邻域中的像在 \(\mathbb{C}\) 中稠密。

\textbf{证明}：反设存在 \(w \in \mathbb{C}\)、\(\varepsilon > 0\) 和 \(\delta > 0\) 使得 \(|f(z) - w| > \varepsilon\) 对所有 \(0 < |z - z_0| < \delta\) 成立。令 \(g(z) = \frac{1}{f(z) - w}\)，则 \(g\) 在 \(0 < |z - z_0| < \delta\) 上全纯且有界（\(|g(z)| \leq 1/\varepsilon\)）。由 Riemann 可去奇点定理（定理 53.3.1），\(g\) 可全纯延拓到 \(z_0\)。

若 \(g(z_0) \neq 0\)，则 \(f(z) = w + \frac{1}{g(z)}\) 在 \(z_0\) 处全纯，与 \(z_0\) 是本性奇点矛盾。

若 \(g(z_0) = 0\)，设 \(g\) 在 \(z_0\) 处的零点阶数为 \(m \geq 1\)，则 \(g(z) = (z - z_0)^m h(z)\)（\(h(z_0) \neq 0\)）。于是 \(f(z) - w = \frac{1}{(z - z_0)^m h(z)}\)，故 \(z_0\) 是 \(f\) 的 \(m\) 阶极点（或 \(w = \infty\) 时 \(f(z) \to \infty\)），同样与 \(z_0\) 是本性奇点矛盾。

因此假设不成立，即 \(f\) 在 \(z_0\) 的任何去心邻域中的像在 \(\mathbb{C}\) 中稠密。\(\blacksquare\)

\textbf{定理 53.3.4（Picard 大定理------陈述）}：若 \(z_0\) 是 \(f\) 的本性奇点，则 \(f\) 在 \(z_0\) 的任何去心邻域中取遍 \(\mathbb{C}\) 中的所有值，至多有一个例外。

\emph{注记}：Picard 大定理的证明超出本书范围，但可通过模函数方法（modular function）得到。完整证明可参见 Ahlfors《Complex Analysis》或 Conway《Functions of One Complex Variable II》。

\hrulefill

\subsection{53.4 整函数与亚纯函数}\label{ux6574ux51fdux6570ux4e0eux4e9aux7eafux51fdux6570}

\textbf{定义 53.4.1（整函数）}：在整个 \(\mathbb{C}\) 上全纯的函数称为\textbf{整函数}。整函数由 Taylor 展开 \(f(z) = \sum_{n=0}^{\infty} a_n z^n\)（收敛半径为 \(\infty\)）完全刻画。

\textbf{定理 53.4.1（整函数的增长与零点------Weierstrass 因子分解）}：设 \(f\) 为整函数，\(\{z_n\}\) 为其非零零点的序列（按重数计）。则 \(f\) 可表示为

\[
f(z) = z^m e^{g(z)} \prod_{n=1}^{\infty} E_{p_n}\left(\frac{z}{z_n}\right)
\]

其中 \(g\) 为整函数，\(E_p(w)\) 为 Weierstrass 初等因子：\(E_0(w) = 1 - w\)，\(E_p(w) = (1 - w) \exp(w + w^2/2 + \cdots + w^p/p)\)（\(p \geq 1\)）。\(p_n\) 的选择使得无穷乘积收敛。

\textbf{证明概要} 设零点序列 \(\{z_n\}\) 满足 \(|z_n| \to \infty\)（因整函数的零点集无聚点于有限点）。初等因子的关键性质：当 \(|w| \leq 1/2\) 时 \(|E_p(w) - 1| \leq |w|^{p+1}\)。选取 \(p_n\) 满足 \(\sum |z/z_n|^{p_n+1}\) 在紧集上一致收敛（例如 \(p_n = n-1\)），则 \(\prod E_{p_n}(z/z_n)\) 在 \(\mathbb{C}\) 上一致收敛于整函数 \(h(z)\)。\(h(z)\) 与 \(f(z)\) 的零点集（计重数）相同，故 \(f(z)/(z^m h(z))\) 是处处非零的整函数，必可写为 \(e^{g(z)}\)（取对数并利用单连通性）。\(\blacksquare\)

\textbf{定义 53.4.2（亚纯函数）}：在区域 \(\Omega\) 上除极点外全纯的函数称为\textbf{亚纯函数}。亚纯函数的极点集是离散的（至多可数）。亚纯函数可表示为两个全纯函数的商（局部，整体上在 \(\mathbb{C}\) 中不一定成立）。

\textbf{定理 53.4.2（Mittag-Leffler 定理）}：给定 \(\mathbb{C}\) 中离散点集 \(\{z_n\}\) 和对应的主要部分 \(P_n(z) = \sum_{k=1}^{m_n} \frac{c_{n,k}}{(z - z_n)^k}\)，存在 \(\mathbb{C}\) 上的亚纯函数以 \(\{z_n\}\) 为极点集，且在 \(z_n\) 处的主要部分恰为 \(P_n(z)\)。

\textbf{证明概要} 假设 \(0 \notin \{z_n\}\)（若 \(0\) 是极点，单独处理）。排列 \(\{z_n\}\) 使得 \(|z_n| \to \infty\)。对每个 \(n\)，\(P_n(z)\) 在 \(z=0\) 处解析，可 Taylor 展开 \(P_n(z) = \sum_{k=0}^{T_n} a_{n,k} z^k\)。取足够大的多项式截断 \(Q_n(z)\) 使得当 \(|z| \leq |z_n|/2\) 时 \(|P_n(z) - Q_n(z)| \leq 1/2^n\)。由 Weierstrass \(M\)-判别法，级数 \(\sum_{n=1}^\infty (P_n(z) - Q_n(z))\) 在 \(\mathbb{C} \setminus \{z_n\}\) 的紧子集上一致收敛，定义亚纯函数以 \(\{z_n\}\) 为极点。Mittag-Leffler 定理是 Weierstrass 因子分解定理的对偶，两者共同构成整函数和亚纯函数理论的基石。\(\blacksquare\)

\hrulefill

\section{第162章：留数定理}\label{ux7b2c162ux7ae0ux7559ux6570ux5b9aux7406}

留数定理是复分析中最强大的计算工具之一，它将复积分转化为奇点处留数的代数和。本章系统介绍留数的定义与计算方法、留数定理（包括有限奇点情形和无穷远点情形），以及留数定理在实积分计算（有理函数积分、含三角函数的积分、Fresnel 积分等）、级数求和和辐角原理中的应用。留数定理的优雅之处在于：一个看似复杂的围道积分，最终归结为若干奇点处的局部不变量（留数）之和------这是 Cauchy 积分定理与 Laurent 级数理论的完美结合。

\subsection{54.1 留数的定义与计算}\label{ux7559ux6570ux7684ux5b9aux4e49ux4e0eux8ba1ux7b97}

\textbf{定义 54.1.1（留数）}：设 \(z_0\) 是 \(f\) 的孤立奇点，\(f\) 在 \(z_0\) 处的 Laurent 展开为 \(f(z) = \sum_{n=-\infty}^{\infty} a_n (z - z_0)^n\)。\(f\) 在 \(z_0\) 处的\textbf{留数}定义为

\[
\operatorname{Res}(f, z_0) = a_{-1} = \frac{1}{2\pi i} \int_{\partial D(z_0, \varepsilon)} f(z) \, dz
\]

\textbf{留数的意义}：留数是 Laurent 级数中 \((z - z_0)^{-1}\) 项的系数，也是 \(f\) 沿绕 \(z_0\) 的小圆周的积分值（除以 \(2\pi i\)）。留数 \(a_{-1}\) 是唯一在积分中存活的系数------因为其他幂次 \((z-z_0)^n\)（\(n \neq -1\)）有原函数，沿闭曲线积分为零。

\textbf{定理 54.1.1（留数的计算规则）}：

\begin{enumerate}
\def\labelenumi{(\roman{enumi})}
\tightlist
\item
  若 \(z_0\) 是 \(m\) 阶极点，则 \(\operatorname{Res}(f, z_0) = \frac{1}{(m-1)!} \lim_{z \to z_0} \frac{d^{m-1}}{dz^{m-1}}[(z - z_0)^m f(z)]\)
\item
  若 \(z_0\) 是 1 阶极点（简单极点），则 \(\operatorname{Res}(f, z_0) = \lim_{z \to z_0} (z - z_0)f(z)\)
\item
  若 \(f(z) = \frac{g(z)}{h(z)}\)，\(g, h\) 全纯，\(g(z_0) \neq 0\)，\(h(z_0) = 0\)，\(h'(z_0) \neq 0\)（简单极点），则 \(\operatorname{Res}(f, z_0) = \frac{g(z_0)}{h'(z_0)}\)
\end{enumerate}

\textbf{证明}：(i) 设 \(f(z) = \frac{a_{-m}}{(z - z_0)^m} + \cdots + \frac{a_{-1}}{z - z_0} + a_0 + \cdots\)。则 \((z - z_0)^m f(z) = a_{-m} + \cdots + a_{-1}(z - z_0)^{m-1} + \cdots\)。求导 \(m-1\) 次后取 \(z \to z_0\) 得 \(a_{-1}\)。

\begin{enumerate}
\def\labelenumi{(\roman{enumi})}
\setcounter{enumi}{1}
\item
  是 (i) 的特例（\(m = 1\)）。
\item
  由洛必达法则，\(a_{-1} = \lim_{z \to z_0} (z - z_0)\frac{g(z)}{h(z)} = \frac{g(z_0)}{h'(z_0)}\)。\(\blacksquare\)
\end{enumerate}

\hrulefill

\subsection{54.2 留数定理}\label{ux7559ux6570ux5b9aux7406}

\textbf{定理 54.2.1（留数定理）}：设 \(\Omega \subseteq \mathbb{C}\) 为单连通开集，\(\gamma\) 为 \(\Omega\) 中简单闭曲线（正向），\(f\) 在 \(\Omega\) 上除有限个孤立奇点 \(z_1, \ldots, z_k\)（均在 \(\gamma\) 内部）外全纯。则

\[
\int_{\gamma} f(z) \, dz = 2\pi i \sum_{j=1}^{k} \operatorname{Res}(f, z_j)
\]

\textbf{证明}：在每个 \(z_j\) 周围取小圆 \(\gamma_j\)（\(\gamma_j\) 在 \(\gamma\) 内部且互不相交）。由 Cauchy 积分定理（应用于 \(\gamma\) 和各 \(\gamma_j\) 之间的区域），

\[
\int_{\gamma} f(z) dz = \sum_{j=1}^{k} \int_{\gamma_j} f(z) dz = 2\pi i \sum_{j=1}^{k} \operatorname{Res}(f, z_j)
\]

其中第二步由留数的定义。\(\blacksquare\)

\textbf{定理 54.2.2（幅角原理）}：设 \(f\) 在简单闭曲线 \(\gamma\) 内部亚纯（除极点外全纯），在 \(\gamma\) 上全纯且非零。则

\[\frac{1}{2\pi i} \int_{\gamma} \frac{f'(z)}{f(z)} dz = N - P\]

其中 \(N\) 为 \(f\) 在 \(\gamma\) 内部的零点个数（按重数计），\(P\) 为极点个数（按重数计）。

\textbf{证明} 设 \(z_0\) 是 \(f\) 的 \(m\) 阶零点（\(m > 0\)）或 \(|m|\) 阶极点（\(m < 0\)）。在 \(z_0\) 附近，\(f(z) = (z-z_0)^m g(z)\)，其中 \(g\) 在 \(z_0\) 处全纯且 \(g(z_0) \neq 0\)。则 \[\frac{f'(z)}{f(z)} = \frac{m}{z-z_0} + \frac{g'(z)}{g(z)}\] 因 \(g'/g\) 在 \(z_0\) 处全纯，\(\frac{f'}{f}\) 在 \(z_0\) 处的留数为 \(m\)（零点时 \(m>0\)，极点时 \(m<0\)）。由留数定理，\(\frac{1}{2\pi i} \int_\gamma \frac{f'}{f} dz = \sum \operatorname{Res}(f'/f, z_j)\)，其中求和遍历 \(\gamma\) 内所有零点和极点。因每个 \(m\) 阶零点贡献 \(+m\)，每个 \(m\) 阶极点贡献 \(-m\)，总和恰为 \(N-P\)。

幅角原理的几何意义是：\(\gamma\) 在映射 \(f\) 下的像曲线绕原点 \(N - P\) 圈。这是因为 \(\frac{1}{2\pi i} \int_\gamma \frac{f'}{f} dz = \frac{1}{2\pi i} \int_{f \circ \gamma} \frac{dw}{w} = \operatorname{Ind}_{f \circ \gamma}(0)\)，即 \(f \circ \gamma\) 关于原点的环绕数。\(\blacksquare\)

\textbf{定理 54.2.3（Rouché 定理）}：设 \(f, g\) 在简单闭曲线 \(\gamma\) 内部全纯，在 \(\gamma\) 上连续，且 \(|f(z) - g(z)| < |f(z)|\)（对所有 \(z \in \gamma\)）。则 \(f\) 和 \(g\) 在 \(\gamma\) 内部有相同数量的零点（按重数计）。

\textbf{证明}：由幅角原理，\(f\) 的零点数为 \(\frac{1}{2\pi i} \int_{\gamma} \frac{f'}{f} dz\)。考虑 \(f_t = f + t(g - f)\)（\(0 \leq t \leq 1\)）。由 \(|g - f| < |f|\) 在 \(\gamma\) 上，\(f_t\) 在 \(\gamma\) 上非零。零点数 \(N(t) = \frac{1}{2\pi i} \int_{\gamma} \frac{f_t'}{f_t} dz\) 是 \(t\) 的连续函数且取整数值，故为常数。\(N(0) = N(1)\)。\(\blacksquare\)

\hrulefill

\subsection{54.3 留数定理在实积分计算中的应用}\label{ux7559ux6570ux5b9aux7406ux5728ux5b9eux79efux5206ux8ba1ux7b97ux4e2dux7684ux5e94ux7528}

\textbf{定理 54.3.1（有理三角函数的积分）}：\(\int_0^{2\pi} R(\cos \theta, \sin \theta) d\theta\)（\(R\) 为有理函数）可通过代换 \(z = e^{i\theta}\)（\(\cos\theta = \frac{z+z^{-1}}{2}\)，\(\sin\theta = \frac{z-z^{-1}}{2i}\)，\(d\theta = \frac{dz}{iz}\)）化为单位圆周上的复积分，然后用留数定理计算。

\textbf{定理 54.3.2（无穷区间上的有理函数积分）}：设 \(P, Q\) 为多项式，\(\deg Q \geq \deg P + 2\)，\(Q\) 在实轴上无零点。则

\[
\int_{-\infty}^{\infty} \frac{P(x)}{Q(x)} \, dx = 2\pi i \sum_{\operatorname{Im} z_j > 0} \operatorname{Res}\left(\frac{P}{Q}, z_j\right)
\]

\textbf{证明}：取上半平面的大半圆围道 \(\Gamma_R = [-R, R] \cup C_R\)（\(C_R\) 为上半圆弧）。由于 \(\deg Q \geq \deg P + 2\)，\(|\frac{P(z)}{Q(z)}| \leq \frac{C}{|z|^2}\)（\(|z|\) 大时），故 \(\int_{C_R} \to 0\)（\(R \to \infty\)）。由留数定理，\(\int_{\Gamma_R} = 2\pi i \sum_{\operatorname{Im} z_j > 0} \operatorname{Res}\)。令 \(R \to \infty\) 即得。\(\blacksquare\)

\textbf{定理 54.3.3（含三角函数的无穷积分 / Jordan 引理）}：设 \(P, Q\) 为多项式，\(\deg Q \geq \deg P + 1\)，\(\alpha > 0\)。则

\[
\int_{-\infty}^{\infty} \frac{P(x)}{Q(x)} e^{i\alpha x} \, dx = 2\pi i \sum_{\operatorname{Im} z_j > 0} \operatorname{Res}\left(\frac{P(z)}{Q(z)} e^{i\alpha z}, z_j\right)
\]

取实部得余弦积分，取虚部得正弦积分。

\textbf{证明}：取上半平面围道 \(\Gamma_R = [-R, R] \cup C_R\)（\(C_R\) 为上半圆弧 \(Re^{i\theta}\)，\(0 \leq \theta \leq \pi\)）。由 Jordan 引理，\(\lim_{R \to \infty} \int_{C_R} \frac{P(z)}{Q(z)} e^{i\alpha z} dz = 0\)（\(\alpha > 0\) 时上半圆弧趋于零，因 \(|e^{i\alpha z}| = e^{-\alpha \operatorname{Im} z} \leq 1\) 且 \(\deg Q \geq \deg P + 1\)）。由留数定理，\(\int_{\Gamma_R} = 2\pi i \sum_{\operatorname{Im} z_j > 0} \operatorname{Res}(\frac{P(z)}{Q(z)} e^{i\alpha z}, z_j)\)。令 \(R \to \infty\)，得 \(\int_{-\infty}^{\infty} \frac{P(x)}{Q(x)} e^{i\alpha x} dx = 2\pi i \sum_{\operatorname{Im} z_j > 0} \operatorname{Res}\)。取实部得余弦积分，取虚部得正弦积分。\(\blacksquare\)

\textbf{定理 54.3.4（实轴上简单极点的处理 / 半留数）}：若 \(f\) 在实轴上有简单极点，可取绕过极点的半圆围道，半圆部分贡献 \(-\pi i \operatorname{Res}\)（上半平面绕过）或 \(+\pi i \operatorname{Res}\)（下半平面绕过）。此时主值积分 \(\operatorname{P.V.}\int_{-\infty}^{\infty} f(x) dx\) 可由留数定理计算。

\textbf{定理 54.3.5（扇形围道与多值函数）}：对于涉及 \(\sqrt{z}\)、\(\log z\) 等多值函数的积分，常使用扇形围道或钥匙孔围道（keyhole contour），并利用分支切割将多值函数限制为单值分支。经典例子包括 \(\int_0^{\infty} \frac{x^{a-1}}{1+x} dx = \frac{\pi}{\sin(\pi a)}\)（\(0 < a < 1\)）和 \(\int_0^{\infty} \frac{\sin x}{x} dx = \frac{\pi}{2}\)。

留数定理是复分析中最强大的计算工具之一。它将实直线上极难计算的积分转化为复平面上有限个奇点处的留数计算，本质上体现了复分析中''局部信息确定整体行为''的深刻原理。在物理应用中，留数定理是色散关系（dispersion relation）、Kramers-Kronig 关系、以及量子场论中 Feynman 传播子计算的基础。

\hrulefill

\subsection{54.4 无穷乘积与 Gamma 函数}\label{ux65e0ux7a77ux4e58ux79efux4e0e-gamma-ux51fdux6570}

\textbf{定义 54.4.1（无穷乘积）}：无穷乘积 \(\prod_{n=1}^{\infty} (1 + a_n)\) 收敛（到非零有限值），若部分积 \(P_N = \prod_{n=1}^{N} (1 + a_n)\) 收敛到非零有限极限。无穷乘积收敛的必要条件是 \(a_n \to 0\)。

\textbf{定理 54.4.1（无穷乘积与级数的关系）}：若 \(a_n \geq 0\)，则 \(\prod (1 + a_n)\) 收敛当且仅当 \(\sum a_n\) 收敛。对复数列，\(\prod (1 + a_n)\) 绝对收敛（即 \(\sum |a_n| < \infty\)）蕴含乘积收敛，且绝对收敛的乘积与因子的次序无关。

\textbf{定理 54.4.2（Weierstrass 乘积定理）}：给定 \(\mathbb{C}\) 中任意离散点集 \(\{z_n\}\)（\(z_n \neq 0\)）和重数 \(m_n \geq 1\)，存在整函数以 \(\{z_n\}\) 为其零点集（且 \(z_n\) 的重数为 \(m_n\)）。该整函数可表示为 Weierstrass 无穷乘积，其中初等因子 \(E_p(z/z_n)\) 确保乘积在 \(\mathbb{C}\) 的紧集上一致收敛。

\textbf{定理 54.4.3（Gamma 函数的 Weierstrass 乘积表示）}：\(\Gamma\) 函数可由 Weierstrass 乘积定义：

\[
\frac{1}{\Gamma(z)} = z e^{\gamma z} \prod_{n=1}^{\infty} \left(1 + \frac{z}{n}\right) e^{-z/n}
\]

其中 \(\gamma = \lim_{n \to \infty} (\sum_{k=1}^{n} 1/k - \log n) \approx 0.5772\) 为 Euler-Mascheroni 常数。\(\Gamma\) 函数是 \(\mathbb{C}\) 上的亚纯函数，在 \(z = 0, -1, -2, \ldots\) 处有简单极点，满足函数方程 \(\Gamma(z+1) = z\Gamma(z)\) 和余元公式 \(\Gamma(z)\Gamma(1-z) = \pi/\sin(\pi z)\)。

\textbf{定理 54.4.4（Stirling 公式）}：\(\Gamma\) 函数在 \(|\arg z| < \pi - \delta\) 时的渐近展开为

\[
\log \Gamma(z) \sim \left(z - \frac{1}{2}\right)\log z - z + \frac{1}{2}\log(2\pi) + \sum_{k=1}^{\infty} \frac{B_{2k}}{2k(2k-1)z^{2k-1}}
\]

其中 \(B_{2k}\) 为 Bernoulli 数。Stirling 公式在组合数学、数论（素数定理的复分析证明）和统计力学中有根本性应用。\(\blacksquare\)

\hrulefill

\section{第163章：全纯函数}\label{ux7b2c163ux7ae0ux5168ux7eafux51fdux6570}

复分析是数学中最优美的分支之一，其核心对象是全纯函数（复可导函数）。与实分析形成鲜明对比的是，复可导性（全纯性）是一个极其强大的条件------全纯函数自动无穷次可导、局部可展开为幂级数、满足 Cauchy 积分公式、且具有刚性（由局部信息唯一确定整体）。这一系列''奇迹''的根源在于 Cauchy-Riemann 方程将复可导性与偏微分方程联系起来，从而赋予全纯函数以调和函数的全部优良性质。

\subsection{51.1 复平面与复函数}\label{ux590dux5e73ux9762ux4e0eux590dux51fdux6570}

\textbf{定义 51.1.1（复平面）}：\(\mathbb{C}\) 可视为 \(\mathbb{R}^2\)，配备复数乘法。\(z = x + iy \in \mathbb{C}\) 对应 \((x, y) \in \mathbb{R}^2\)。\(z\) 的模为 \(|z| = \sqrt{x^2 + y^2}\)。

\textbf{定义 51.1.2（复函数的极限与连续性）}：设 \(\Omega \subseteq \mathbb{C}\) 为开集，\(f : \Omega \to \mathbb{C}\)，\(z_0 \in \Omega\)。

\[
\lim_{z \to z_0} f(z) = w_0 \iff \forall \varepsilon > 0, \exists \delta > 0, \; 0 < |z - z_0| < \delta \implies |f(z) - w_0| < \varepsilon
\]

\(f\) 在 \(z_0\) 处连续若 \(\lim_{z \to z_0} f(z) = f(z_0)\)。

\textbf{注记 51.1.1}：复函数的极限和连续性在形式上和实函数完全相同，但这里 \(z\) 沿复平面中任意路径趋近 \(z_0\)，比实直线上只有左右两个方向的要求更强。这一''全方向''的要求是复分析中许多刚性现象的根源。

\textbf{定义 51.1.3（复平面上的拓扑基础）}：区域 \(\Omega \subseteq \mathbb{C}\) 是连通开集。\(\Omega\) 称为\textbf{单连通}的，若 \(\mathbb{C} \setminus \Omega\) 在扩充复平面 \(\hat{\mathbb{C}} = \mathbb{C} \cup \{\infty\}\) 中是连通的。等价地，\(\Omega\) 中任何闭曲线可连续收缩为一点。

\hrulefill

\subsection{51.2 全纯函数与 Cauchy-Riemann 方程}\label{ux5168ux7eafux51fdux6570ux4e0e-cauchy-riemann-ux65b9ux7a0b}

\textbf{定义 51.2.1（复导数与全纯函数）}：设 \(\Omega \subseteq \mathbb{C}\) 为开集，\(f : \Omega \to \mathbb{C}\)，\(z_0 \in \Omega\)。若极限

\[
f'(z_0) = \lim_{z \to z_0} \frac{f(z) - f(z_0)}{z - z_0}
\]

存在，则称 \(f\) 在 \(z_0\) 处\textbf{复可导}。若 \(f\) 在 \(\Omega\) 上每点复可导，则称 \(f\) 为 \(\Omega\) 上的\textbf{全纯函数}（或\textbf{解析函数}）。

\textbf{注记 51.2.1}：复可导性的要求比实可导性强得多。在实分析中，\(f'(x_0)\) 的存在仅要求差商沿左右两个方向趋近同一极限；在复分析中，差商必须沿所有可能的复路径趋近同一极限。这一差异是全纯函数具有刚性（如无穷次可导、幂级数展开等）的根本原因。

\textbf{定义 51.2.2（Wirtinger 导数）}：引入形式偏导数

\[
\frac{\partial}{\partial z} = \frac{1}{2}\left(\frac{\partial}{\partial x} - i\frac{\partial}{\partial y}\right), \quad \frac{\partial}{\partial \bar{z}} = \frac{1}{2}\left(\frac{\partial}{\partial x} + i\frac{\partial}{\partial y}\right)
\]

则 Cauchy-Riemann 方程等价于 \(\frac{\partial f}{\partial \bar{z}} = 0\)。这一形式揭示了全纯函数的本质：\(f\) 不依赖于 \(\bar{z}\)，即 \(f\) 是 \(z\) 而非 \(z\) 和 \(\bar{z}\) 的函数。

\textbf{定理 51.2.1（Cauchy-Riemann 方程）}：设 \(f(z) = u(x, y) + i v(x, y)\)（\(u, v : \mathbb{R}^2 \to \mathbb{R}\)），\(z_0 = x_0 + iy_0\)。\(f\) 在 \(z_0\) 处复可导当且仅当 \(u, v\) 在 \((x_0, y_0)\) 处可微且满足

\[
\frac{\partial u}{\partial x} = \frac{\partial v}{\partial y}, \quad \frac{\partial u}{\partial y} = -\frac{\partial v}{\partial x}
\]

此时 \(f'(z_0) = \frac{\partial u}{\partial x}(x_0, y_0) + i\frac{\partial v}{\partial x}(x_0, y_0) = \frac{\partial v}{\partial y}(x_0, y_0) - i\frac{\partial u}{\partial y}(x_0, y_0)\)。

\textbf{证明}：

（\(\Rightarrow\)）设 \(f'(z_0) = A + iB\) 存在。取 \(z \to z_0\) 沿实轴方向（\(z = z_0 + h\)，\(h \in \mathbb{R}\)）：

\[
f'(z_0) = \lim_{h \to 0} \frac{u(x_0 + h, y_0) + iv(x_0 + h, y_0) - u(x_0, y_0) - iv(x_0, y_0)}{h} = \frac{\partial u}{\partial x} + i\frac{\partial v}{\partial x}
\]

取 \(z \to z_0\) 沿虚轴方向（\(z = z_0 + ih\)，\(h \in \mathbb{R}\)）：

\[
f'(z_0) = \lim_{h \to 0} \frac{u(x_0, y_0 + h) + iv(x_0, y_0 + h) - u(x_0, y_0) - iv(x_0, y_0)}{ih} = \frac{1}{i}\frac{\partial u}{\partial y} + \frac{\partial v}{\partial y} = -i\frac{\partial u}{\partial y} + \frac{\partial v}{\partial y}
\]

比较实部和虚部即得 Cauchy-Riemann 方程。

（\(\Leftarrow\)）设 \(u, v\) 可微且满足 C-R 方程。由实可微性，\(u(x_0 + h, y_0 + k) - u(x_0, y_0) = u_x h + u_y k + o(\sqrt{h^2 + k^2})\)，\(v\) 类似。令 \(\Delta z = h + ik\)，则

\[
f(z_0 + \Delta z) - f(z_0) = (u_x + iv_x)h + (u_y + iv_y)k + o(|\Delta z|)
\]

由 C-R 方程，\(u_y + iv_y = -v_x + iu_x = i(u_x + iv_x)\)。故

\[
f(z_0 + \Delta z) - f(z_0) = (u_x + iv_x)(h + ik) + o(|\Delta z|) = (u_x + iv_x)\Delta z + o(|\Delta z|)
\]

因此 \(\lim_{\Delta z \to 0} \frac{f(z_0 + \Delta z) - f(z_0)}{\Delta z} = u_x + iv_x\)。\(\blacksquare\)

\textbf{推论 51.2.2（调和性）}：若 \(f = u + iv\) 全纯且 \(u, v \in C^2\)，则 \(u, v\) 为调和函数，即 \(\Delta u = \frac{\partial^2 u}{\partial x^2} + \frac{\partial^2 u}{\partial y^2} = 0\)，\(\Delta v = 0\)。这由 Cauchy-Riemann 方程直接求导得出（混合偏导数相等）。

\textbf{定理 51.2.3（全纯函数的基本运算）}：若 \(f, g\) 在 \(\Omega\) 上全纯，则

\begin{enumerate}
\def\labelenumi{(\roman{enumi})}
\tightlist
\item
  \(f \pm g\)，\(fg\) 全纯，且 \((f \pm g)' = f' \pm g'\)，\((fg)' = f'g + fg'\)
\item
  若 \(g \neq 0\)，则 \(f/g\) 全纯，且 \((f/g)' = (f'g - fg')/g^2\)
\item
  链式法则：\((f \circ g)'(z) = f'(g(z)) g'(z)\)
\end{enumerate}

\textbf{证明}：与实函数导数运算法则的证明完全相同（仅需将绝对值替换为模）。\(\blacksquare\)

\hrulefill

\subsection{51.3 幂级数作为全纯函数}\label{ux5e42ux7ea7ux6570ux4f5cux4e3aux5168ux7eafux51fdux6570}

\textbf{定理 51.3.1（幂级数的全纯性）}：设 \(\sum_{n=0}^{\infty} a_n z^n\) 的收敛半径为 \(R > 0\)。则 \(f(z) = \sum_{n=0}^{\infty} a_n z^n\) 在 \(\{z \in \mathbb{C} \mid |z| < R\}\) 上全纯，且

\[
f'(z) = \sum_{n=1}^{\infty} n a_n z^{n-1}, \quad |z| < R
\]

（即幂级数可以逐项求导，收敛半径不变。）

\textbf{证明}：由幂级数的性质，\(\sum n a_n z^{n-1}\) 的收敛半径也是 \(R\)。在 \(|z| \leq r < R\) 的紧集上，逐项求导的级数一致收敛（由 Weierstrass M-判别法），故可逐项求导。\(\blacksquare\)

\textbf{推论 51.3.2}：\(e^z = \sum_{n=0}^{\infty} \frac{z^n}{n!}\)，\(\sin z = \sum_{n=0}^{\infty} (-1)^n \frac{z^{2n+1}}{(2n+1)!}\)，\(\cos z = \sum_{n=0}^{\infty} (-1)^n \frac{z^{2n}}{(2n)!}\) 在整个 \(\mathbb{C}\) 上全纯（称为\textbf{整函数}），且 \((e^z)' = e^z\)，\((\sin z)' = \cos z\)，\((\cos z)' = -\sin z\)。

\textbf{注记 51.3.1}：指数函数 \(e^z\) 在复平面上的行为与实指数函数有本质不同。Euler 公式 \(e^{i\theta} = \cos\theta + i\sin\theta\) 揭示了指数函数与三角函数之间的深层联系，这是实分析中完全不可见的。复指数函数是周期函数（周期为 \(2\pi i\)），这意味着复对数函数是多值的，需要引入分支切割（branch cut）来定义单值分支。

\textbf{定理 51.3.3（Abel 定理------幂级数边界行为）}：若 \(\sum_{n=0}^{\infty} a_n\) 收敛于 \(S\)，则 \(\lim_{z \to 1^-} \sum_{n=0}^{\infty} a_n z^n = S\)（沿实轴趋近）。这一定理在 Tauber 型定理的研究中具有基础性地位。

\hrulefill

\subsection{51.4 初等全纯函数与共形映射}\label{ux521dux7b49ux5168ux7eafux51fdux6570ux4e0eux5171ux5f62ux6620ux5c04}

\textbf{定义 51.4.1（共形映射）}：全纯函数 \(f\) 在 \(f'(z_0) \neq 0\) 的点处是\textbf{共形}的（保角且保定向）。具体地，\(f\) 在 \(z_0\) 处将两条曲线的夹角保持（大小和方向）。

\textbf{定理 51.4.1（分式线性变换 / Möbius 变换）}：映射 \(T(z) = \frac{az + b}{cz + d}\)（\(ad - bc \neq 0\)）将扩充复平面 \(\hat{\mathbb{C}}\) 共形地映到自身。Möbius 变换将圆和直线映为圆和直线（广义圆），且具有三传递性（可通过三个点的像唯一确定）。Möbius 变换群同构于 \(\operatorname{PSL}(2, \mathbb{C}) = \operatorname{SL}(2, \mathbb{C})/\{\pm I\}\)。

\hrulefill

\section{第164章：复积分}\label{ux7b2c164ux7ae0ux590dux79efux5206}

复积分是复分析的核心工具，其理论深度远超实分析中的曲线积分。本章从复曲线积分的定义出发，建立 Cauchy 积分定理（Cauchy-Goursat 定理）------全纯函数沿闭曲线的积分为零------以及由此派生的 Cauchy 积分公式、高阶导数公式和 Cauchy 估计。这些结果揭示了一个根本性的事实：全纯函数在区域内的值完全由其边界值决定，且全纯性自动蕴含无穷次可微性和解析性。复积分理论还将引出留数定理、辐角原理和 Rouche 定理等强有力的计算与定性分析工具。

\subsection{52.1 曲线积分}\label{ux66f2ux7ebfux79efux5206}

\textbf{定义 52.1.1（曲线）}：\(\mathbb{C}\) 中的\textbf{曲线}（或\textbf{路径}）\(\gamma\) 是一个连续可微映射 \(\gamma : [a, b] \to \mathbb{C}\)。\(\gamma\) 称为\textbf{闭曲线}若 \(\gamma(a) = \gamma(b)\)。更一般地，可定义分段连续可微曲线。

\textbf{定义 52.1.2（复曲线积分）}：设 \(f\) 在 \(\gamma\) 上连续。定义 \(f\) 沿 \(\gamma\) 的积分为

\[
\int_{\gamma} f(z) \, dz = \int_a^b f(\gamma(t)) \gamma'(t) \, dt
\]

\textbf{定义 52.1.3（曲线的指标 / 绕数）}：设 \(\gamma\) 为不经过 \(z_0\) 的闭曲线。\(\gamma\) 关于 \(z_0\) 的\textbf{指标}（绕数）定义为

\[
\operatorname{Ind}_\gamma(z_0) = \frac{1}{2\pi i} \int_{\gamma} \frac{dz}{z - z_0}
\]

\(\operatorname{Ind}_\gamma(z_0)\) 是整数，表示 \(\gamma\) 绕 \(z_0\) 的圈数（逆时针为正）。

\textbf{定理 52.1.1（积分的基本估计 / ML 估计）}：设 \(|f(z)| \leq M\) 在 \(\gamma\) 上，\(L\) 为 \(\gamma\) 的长度。则

\[
\left|\int_{\gamma} f(z) \, dz\right| \leq M L
\]

\textbf{证明}： \[
\left|\int_{\gamma} f(z) dz\right| = \left|\int_a^b f(\gamma(t))\gamma'(t) dt\right| \leq \int_a^b |f(\gamma(t))||\gamma'(t)| dt \leq M \int_a^b |\gamma'(t)| dt = M L
\]

\(\blacksquare\)

\hrulefill

\subsection{52.2 Cauchy 积分定理}\label{cauchy-ux79efux5206ux5b9aux7406}

\textbf{定理 52.2.1（Cauchy 积分定理------矩形版本 / Goursat 定理）}：设 \(\Omega \subseteq \mathbb{C}\) 为开集，\(f\) 在 \(\Omega\) 上全纯，\(R \subseteq \Omega\) 为闭矩形（边界 \(\partial R\) 由四条线段组成）。则

\[
\int_{\partial R} f(z) \, dz = 0
\]

\textbf{证明}（Goursat）：将矩形 \(R\) 四等分为 \(R^{(1)}, \ldots, R^{(4)}\)。则

\[
\int_{\partial R} f(z) dz = \sum_{j=1}^{4} \int_{\partial R^{(j)}} f(z) dz
\]

（公共边上的积分抵消）。设 \(R_1\) 为四个子矩形中积分绝对值最大者，则 \(|\int_{\partial R} f| \leq 4 |\int_{\partial R_1} f|\)。重复此过程，得到矩形套 \(R \supseteq R_1 \supseteq R_2 \supseteq \cdots\) 使得

\[
\left|\int_{\partial R_n} f(z) dz\right| \geq \frac{1}{4^n} \left|\int_{\partial R} f(z) dz\right|
\]

设 \(z_0 \in \bigcap_{n=1}^{\infty} R_n\)。由 \(f\) 在 \(z_0\) 处全纯，

\[
f(z) = f(z_0) + f'(z_0)(z - z_0) + \psi(z)(z - z_0)
\]

其中 \(\psi(z) \to 0\)（\(z \to z_0\)）。前两项 \(f(z_0) + f'(z_0)(z - z_0)\) 有原函数，沿闭矩形积分为 \(0\)。故

\[
\int_{\partial R_n} f(z) dz = \int_{\partial R_n} \psi(z)(z - z_0) dz
\]

设 \(R_n\) 的边长为 \(L/2^n\)，对角线长为 \(\sqrt{2}L/2^n\)。对任意 \(\varepsilon > 0\)，当 \(n\) 足够大时 \(|\psi(z)| < \varepsilon\)。故

\[
\left|\int_{\partial R_n} f(z) dz\right| \leq \varepsilon \cdot \frac{\sqrt{2}L}{2^n} \cdot \frac{4L}{2^n} = \frac{4\sqrt{2}L^2}{4^n} \varepsilon
\]

因此 \(|\int_{\partial R} f| \leq 4^n \cdot \frac{4\sqrt{2}L^2}{4^n} \varepsilon = 4\sqrt{2}L^2 \varepsilon\)。由 \(\varepsilon\) 任意性，\(\int_{\partial R} f = 0\)。\(\blacksquare\)

\textbf{定理 52.2.2（Cauchy 积分定理------一般版本）}：设 \(\Omega \subseteq \mathbb{C}\) 为\textbf{单连通}开集（任何闭曲线可在 \(\Omega\) 内连续收缩为一点），\(f\) 在 \(\Omega\) 上全纯。则对 \(\Omega\) 中任何闭曲线 \(\gamma\)，

\[
\int_{\gamma} f(z) \, dz = 0
\]

\textbf{注记 52.2.1}：一般版本的证明需要将区域三角剖分或使用同伦论证。核心结论：在单连通区域内，全纯函数的闭曲线积分为零。Goursat 的贡献在于证明中无需假设 \(f'\) 连续------这是 Cauchy 原始证明中的缺陷。

\textbf{推论 52.2.3（原函数的存在性）}：在单连通区域 \(\Omega\) 上全纯的函数 \(f\) 有原函数 \(F\)（\(F' = f\)），\(F(z) = \int_{z_0}^{z} f(\zeta) d\zeta\)（积分与路径无关）。

\textbf{定理 52.2.4（Morera 定理------Cauchy 定理的逆）}：若 \(f\) 在 \(\Omega\) 上连续，且对 \(\Omega\) 中任何闭曲线 \(\gamma\) 有 \(\int_{\gamma} f(z) dz = 0\)，则 \(f\) 在 \(\Omega\) 上全纯。Morera 定理是 Cauchy 积分定理的逆定理，常用来证明函数序列的极限函数的全纯性。

\hrulefill

\subsection{52.3 Cauchy 积分公式}\label{cauchy-ux79efux5206ux516cux5f0f}

\textbf{定理 52.3.1（Cauchy 积分公式）}：设 \(\Omega \subseteq \mathbb{C}\) 为开集，\(f\) 在 \(\Omega\) 上全纯，\(\overline{D}(z_0, r) \subseteq \Omega\)（以 \(z_0\) 为中心、\(r\) 为半径的闭圆盘）。则对任意 \(z \in D(z_0, r)\)（开圆盘），

\[
f(z) = \frac{1}{2\pi i} \int_{\partial D(z_0, r)} \frac{f(\zeta)}{\zeta - z} \, d\zeta
\]

\textbf{证明}：固定 \(z \in D(z_0, r)\)。考虑函数 \(g(\zeta) = \frac{f(\zeta) - f(z)}{\zeta - z}\)（\(\zeta \neq z\)），\(g(z) = f'(z)\)。\(g\) 在 \(\Omega \setminus \{z\}\) 上全纯，在 \(z\) 处连续（从而在 \(\Omega\) 上全纯，由 Riemann 可去奇点定理------见定理 53.4.1）。

在 \(\Omega\) 内取以 \(z\) 为中心的小圆 \(\partial D(z, \varepsilon)\)。由 Cauchy 积分定理（应用于环形区域），

\[
\int_{\partial D(z_0, r)} g(\zeta) d\zeta = \int_{\partial D(z, \varepsilon)} g(\zeta) d\zeta
\]

即 \(\int_{\partial D(z_0, r)} \frac{f(\zeta)}{\zeta - z} d\zeta - f(z) \int_{\partial D(z_0, r)} \frac{d\zeta}{\zeta - z} = \int_{\partial D(z, \varepsilon)} \frac{f(\zeta) - f(z)}{\zeta - z} d\zeta\)。

计算 \(\int_{\partial D(z_0, r)} \frac{d\zeta}{\zeta - z} = 2\pi i\)（由 Cauchy 积分公式基础情形 \(f \equiv 1\)）。右边当 \(\varepsilon \to 0\) 时，\(|g(\zeta)| \leq |f'(z)| + 1\)（有界），故积分绝对值 \(\leq M \cdot 2\pi\varepsilon \to 0\)。因此

\[
\int_{\partial D(z_0, r)} \frac{f(\zeta)}{\zeta - z} d\zeta - 2\pi i f(z) = 0
\]

即 \(f(z) = \frac{1}{2\pi i} \int_{\partial D(z_0, r)} \frac{f(\zeta)}{\zeta - z} d\zeta\)。\(\blacksquare\)

\textbf{推论 52.3.2（Cauchy 积分公式------高阶导数）}：在定理 52.3.1 的条件下，\(f\) 在 \(D(z_0, r)\) 上无穷次复可导，且

\[
f^{(n)}(z) = \frac{n!}{2\pi i} \int_{\partial D(z_0, r)} \frac{f(\zeta)}{(\zeta - z)^{n+1}} \, d\zeta, \quad n = 0, 1, 2, \ldots
\]

\textbf{证明}：对 \(n\) 归纳。\(n = 0\) 即 Cauchy 积分公式。假设公式对 \(n\) 成立。由差商极限和积分号下取极限的合法性（被积函数一致有界），

\[
f^{(n+1)}(z) = \lim_{h \to 0} \frac{f^{(n)}(z+h) - f^{(n)}(z)}{h} = \frac{n!}{2\pi i} \lim_{h \to 0} \int_{\partial D} \frac{f(\zeta)}{h} \left(\frac{1}{(\zeta - z - h)^{n+1}} - \frac{1}{(\zeta - z)^{n+1}}\right) d\zeta
\]

利用 \(\frac{1}{(\zeta - z - h)^{n+1}} - \frac{1}{(\zeta - z)^{n+1}} \to (n+1) \frac{h}{(\zeta - z)^{n+2}}\)（\(h \to 0\)），得 \(f^{(n+1)}(z) = \frac{(n+1)!}{2\pi i} \int_{\partial D} \frac{f(\zeta)}{(\zeta - z)^{n+2}} d\zeta\)。\(\blacksquare\)

\textbf{推论 52.3.3（全纯函数无穷次可导）}：全纯函数自动无穷次可导（与实分析根本不同）。

\textbf{推论 52.3.4（Cauchy 估计）}：设 \(f\) 在 \(\overline{D}(z_0, R)\) 上全纯，\(|f(z)| \leq M\)。则

\[
|f^{(n)}(z_0)| \leq \frac{n! M}{R^n}
\]

\textbf{证明}：由 Cauchy 高阶导数公式，\(|f^{(n)}(z_0)| = \frac{n!}{2\pi} |\int_{\partial D(z_0, R)} \frac{f(\zeta)}{(\zeta - z_0)^{n+1}} d\zeta| \leq \frac{n!}{2\pi} \cdot \frac{M}{R^{n+1}} \cdot 2\pi R = \frac{n! M}{R^n}\)。\(\blacksquare\)

\hrulefill

\subsection{52.4 全纯函数的局部性质}\label{ux5168ux7eafux51fdux6570ux7684ux5c40ux90e8ux6027ux8d28}

\textbf{定理 52.4.1（最大模原理）}：设 \(f\) 在区域 \(\Omega\) 上全纯。若 \(|f|\) 在 \(\Omega\) 内某点取得最大值，则 \(f\) 为常数。等价地，非恒常的全纯函数的模在其定义域内不能取得最大值。

\textbf{证明}：设 \(|f(z_0)|\) 为 \(\Omega\) 内的最大值。由 Cauchy 积分公式（均值性质），对充分小的 \(r\)，\(f(z_0) = \frac{1}{2\pi} \int_0^{2\pi} f(z_0 + re^{i\theta}) d\theta\)。若 \(|f(z_0 + re^{i\theta})| \leq |f(z_0)|\) 且等号在某点不成立，则积分的模严格小于 \(|f(z_0)|\)（由连续性），矛盾。故 \(|f|\) 在 \(z_0\) 的邻域内恒为常数，由全纯函数的刚性，\(f\) 在 \(\Omega\) 上为常数。\(\blacksquare\)

\textbf{定理 52.4.2（Schwarz 引理）}：设 \(f: D(0,1) \to D(0,1)\) 全纯且 \(f(0) = 0\)。则 \(|f(z)| \leq |z|\)（对所有 \(|z| < 1\)）且 \(|f'(0)| \leq 1\)。若存在 \(z_0 \neq 0\) 使 \(|f(z_0)| = |z_0|\) 或 \(|f'(0)| = 1\)，则 \(f(z) = e^{i\theta}z\)（旋转）。

\textbf{证明}：令 \(g(z) = f(z)/z\)（\(z \neq 0\)），\(g(0) = f'(0)\)。\(g\) 在 \(D(0,1)\) 上全纯。对 \(|z| = r < 1\)，\(|g(z)| = |f(z)|/|z| \leq 1/r\)。由最大模原理，\(|g(z)| \leq 1/r\) 在 \(|z| \leq r\) 上成立。令 \(r \to 1\) 得 \(|g(z)| \leq 1\)。故 \(|f(z)| \leq |z|\) 且 \(|f'(0)| = |g(0)| \leq 1\)。等号成立时 \(g\) 为常数（模为 \(1\)），故 \(f(z) = e^{i\theta}z\)。\(\blacksquare\)

\textbf{定理 52.4.3（开映射定理）}：全纯非常数函数将开集映为开集。这是最大模原理的一个重要推论，表明全纯函数是开映射。

\textbf{定理 52.4.4（逆函数定理与隐函数定理的复版本）}：设 \(f\) 在 \(z_0\) 处全纯且 \(f'(z_0) \neq 0\)。则 \(f\) 在 \(z_0\) 的某邻域内是双全纯的（即存在全纯的逆映射）。这一结论比实分析的逆函数定理更强------在实分析中需要 \(f \in C^1\) 且 Jacobi 行列式非零；在复分析中，\(f'(z_0) \neq 0\) 不仅保证 \(f\) 局部可逆，还保证逆映射也是全纯的。这一事实是复分析刚性（rigidity）的又一体现，其根源在于全纯函数的 Cauchy-Riemann 方程将 Jacobi 行列式 \(|f'(z_0)|^2\) 与复导数建立了直接联系。

\textbf{定理 52.4.5（Schwarz-Pick 引理）}：设 \(f: D(0,1) \to D(0,1)\) 全纯。则对任意 \(z \in D(0,1)\)，\(\frac{|f'(z)|}{1 - |f(z)|^2} \leq \frac{1}{1 - |z|^2}\)。Schwarz-Pick 引理是 Schwarz 引理在 Poincaré 度量下的推广，表明全纯自映射 \(f: D \to D\) 关于双曲度量（Poincaré 度量）是不增距离的。这一几何解释是复动力系统（complex dynamics）和 Teichmüller 理论的基础。

\subsection{52.5 全纯函数序列与正规族}\label{ux5168ux7eafux51fdux6570ux5e8fux5217ux4e0eux6b63ux89c4ux65cf}

\textbf{定理 52.5.1（Weierstrass 定理）}：设 \(\{f_n\}\) 为 \(\Omega\) 上的全纯函数列，且在 \(\Omega\) 的任一紧集上一致收敛于 \(f\)。则 \(f\) 在 \(\Omega\) 上全纯，且 \(\{f_n'\}\) 在 \(\Omega\) 的任一紧集上一致收敛于 \(f'\)。

\textbf{定理 52.5.2（Montel 定理）}：设 \(\mathcal{F}\) 为 \(\Omega\) 上的全纯函数族，且在 \(\Omega\) 的任一紧集上一致有界。则 \(\mathcal{F}\) 是正规族（即 \(\mathcal{F}\) 的任一序列有子列在紧集上一致收敛）。Montel 定理是复分析中紧性论证的核心工具，在 Riemann 映射定理的证明和复动力系统（Julia 集和 Fatou 集理论）中不可或缺。\(\blacksquare\)

\newpage

\chapter{06-分析/02-复分析/03-复分析深化.md}\label{ux5206ux679002-ux590dux5206ux679003-ux590dux5206ux6790ux6df1ux5316.md}

\section{第165章：值分布理论（Nevanlinna 理论）}\label{ux7b2c165ux7ae0ux503cux5206ux5e03ux7406ux8bbanevanlinna-ux7406ux8bba}

Nevanlinna理论由R. Nevanlinna于1925年创立，将Picard定理从定性陈述提升为精密的定量理论，是复分析中亚纯函数值分布的核心框架。

\textbf{定义108.1}（Poisson-Jensen公式与特征函数）：设\(f\)为\(\mathbb{C}\)上非常数亚纯函数。对\(r > 0\)，命\(f\)在\(|z| \leq r\)中的零点为\(a_\mu\)（计重数）、极点为\(b_\nu\)（计重数）。\textbf{Poisson-Jensen公式}给出 \[\log|f(z)| = \frac{1}{2\pi} \int_0^{2\pi} \log|f(re^{i\theta})| \frac{r^2 - |z|^2}{|re^{i\theta} - z|^2} d\theta - \sum_\mu \log\left|\frac{r(z - a_\mu)}{r^2 - \bar{a}_\mu z}\right| + \sum_\nu \log\left|\frac{r(z - b_\nu)}{r^2 - \bar{b}_\nu z}\right|\]

取\(z = 0\)，可得\textbf{Jensen公式}：\(\log|f(0)| = \frac{1}{2\pi}\int_0^{2\pi} \log|f(re^{i\theta})| d\theta - \sum_{|a_\mu| < r} \log\frac{r}{|a_\mu|} + \sum_{|b_\nu| < r} \log\frac{r}{|b_\nu|}\)。

定义\textbf{靠近函数}\(m(r, f) = \frac{1}{2\pi} \int_0^{2\pi} \log^+ |f(re^{i\theta})| d\theta\)（\(\log^+ x = \max(\log x, 0)\)），\textbf{计数函数}\(N(r, f) = \int_0^r \frac{n(t, f) - n(0, f)}{t} dt + n(0, f) \log r\)（\(n(t, f)\)为\(f\)在\(|z| \leq t\)中的极点数）。\textbf{Nevanlinna特征函数}定义为 \[T(r, f) = m(r, f) + N(r, f)\] \(T(r, f)\)衡量\(f\)的整体增长率------有理函数时\(T(r, f) = O(\log r)\)，超越亚纯函数时\(T(r, f)/\log r \to \infty\)。

\textbf{定理108.1}（Nevanlinna第一基本定理）：对任意\(a \in \mathbb{C} \cup \{\infty\}\)， \[T\left(r, \frac{1}{f-a}\right) = T(r, f) + O(1) \quad (r \to \infty)\] 即\(m(r, a, f) + N(r, a, f) = T(r, f) + O(1)\)，表明特征函数在Möbius变换下几乎不变。

\textbf{证明}：由Poisson-Jensen公式取\(z = 0\)，配合\(\log^+\)与\(\log\)的关系，将\(\log|f(0) - a|\)表为\(m(r, 1/(f-a)) - N(r, 1/(f-a))\)与有限的\(f(0)\)项之和。正负部分拆分即得\(T(r, f - a) = T(r, f) + O(1)\)，再由\(T\)的Möbius不变性完成。\(\blacksquare\)

\textbf{定理108.2}（Nevanlinna第二基本定理）：对互异的\(a_1, \ldots, a_q \in \mathbb{C} \cup \{\infty\}\)（\(q \geq 2\)）， \[\sum_{j=1}^q m(r, a_j, f) \leq 2T(r, f) - N_1(r) + S(r, f)\] 其中\(N_1(r)\)为\(f\)的分支点修正项（\(f'\)的零点计数，不计重数），\(S(r, f) = O(\log r T(r, f))\)为误差项（在\(r \to \infty\)除去一个有限线性测度集外成立）。等价形式：\((q-2) T(r, f) \leq \sum_{j=1}^q \overline{N}(r, a_j, f) + S(r, f)\)（\(\overline{N}\)为不计重数的计数函数）。

\textbf{证明}：核心引理是对数导数的估计（Nevanlinna 1930 年代）：对于亚纯函数 \(f\)， \[\int_0^{2\pi} \log^+ \left|\frac{f'(re^{i\theta})}{f(re^{i\theta})}\right| d\theta = o(T(r, f)) \quad (r \to \infty, \; r \notin E),\] 其中 \(E\) 为有限线性测度集。选取互异值 \(a_1, \ldots, a_q \in \mathbb{C} \cup \{\infty\}\)。考虑辅助函数 \(F = \prod_{j=1}^q (f - a_j)\)（若某 \(a_j = \infty\)，则将其从乘积中略去）。由 Jensen 公式和第一基本定理，\(m(r, F) = \sum_j m(r, a_j, f) + O(1)\)。另一方面，导数 \(F'/F = \sum_j f'/(f - a_j)\) 的逼近函数可由对数导数引理控制：\(m(r, F'/F) = o(T(r, f))\)。将 \(F'/F\) 与各个 \(f'/(f - a_j)\) 的逼近函数建立不等式关系，并利用 \(m(r, f'/(f - a_j))\) 的上界，最终导出 \(\sum_j m(r, a_j, f) \leq 2T(r, f) - N_1(r) + S(r, f)\)。换用计数函数即得 \((q-2)T(r, f) \leq \sum_j \overline{N}(r, a_j, f) + S(r, f)\)。\(\blacksquare\)

\textbf{推论108.3}（Picard小定理作为推论）：若整函数\(f\)避开两个有穷值，则第二基本定理给出\((2-2)T(r, f) \leq 0 + S(r, f)\)，与\(f\)为超越整函数时\(T(r, f)/\log r \to \infty\)矛盾。因此\(f\)必为常数------Picard小定理得证。

\hrulefill

\section{第166章：拟共形映射与Teichmüller空间}\label{ux7b2c166ux7ae0ux62dfux5171ux5f62ux6620ux5c04ux4e0eteichmuxfcllerux7a7aux95f4}

拟共形映射是共形映射的自然推广，允许有界角度畸变，其理论由 Grotzsch（1928）、Ahlfors（1935）和 Teichmüller（1939）奠基，并由 Bers 学派在 1960 年代系统发展。拟共形映射的核心是 Beltrami 方程 \(f_{\bar{z}} = \mu(z) f_z\)，其中 Beltrami 系数 \(\mu\) 度量了共形性的偏离程度。拟共形映射不仅是复分析的重要工具，更是 Riemann 面模空间（Teichmüller 空间）理论的核心语言------Teichmüller 空间上的 Teichmüller 度量由拟共形映射的伸缩商所定义，而 Bers 嵌入定理将 Teichmüller 空间实现为全纯二次微分的复 Banach 空间中的有界域。拟共形映射在复动力系统（Sullivan 的''无游荡域定理''）、低维拓扑（Mostow 刚性定理的证明）和偏微分方程（椭圆型方程组的正则性）中均有深刻应用。

\subsection{拟共形映射}\label{ux62dfux5171ux5f62ux6620ux5c04}

拟共形映射是共形映射的自然推广，允许有界的角度畸变。其解析定义基于\textbf{Beltrami方程}。

\textbf{定义 129.1}（Beltrami 方程与拟共形映射的解析定义）：设 \(\Omega \subseteq \mathbb{C}\) 为区域。一个同胚 \(f : \Omega \to f(\Omega) \subseteq \mathbb{C}\) 称为 \textbf{\(K\)-拟共形映射}（\(K \geq 1\)），若 \(f\) 在 Sobolev 空间 \(W_{\text{loc}}^{1,2}(\Omega)\) 中（即 \(f\) 具有局部 \(L^2\) 可积的分布导数），且几乎处处满足 \textbf{Beltrami 方程}：

\[f_{\bar{z}}(z) = \mu(z) f_z(z) \quad (\text{a.e. } z \in \Omega)\]

其中 \(\mu(z)\) 为\textbf{Beltrami 系数}------一个复值可测函数，满足

\[\|\mu\|_\infty = \operatorname{ess\,sup}_{z \in \Omega} |\mu(z)| \leq \frac{K-1}{K+1} < 1\]

\(f_z = \frac{1}{2}(f_x - i f_y)\) 和 \(f_{\bar{z}} = \frac{1}{2}(f_x + i f_y)\) 为 Wirtinger 导数。当 \(\mu \equiv 0\) 时 Beltrami 方程退化为 Cauchy-Riemann 方程 \(f_{\bar{z}} = 0\)；\(\mu \neq 0\) 时 \(f\) 将无穷小圆映为椭圆，偏心率由 \(|\mu|\) 决定。

\textbf{定义 129.2}（伸缩商与几何定义）：拟共形映射的几何定义更为直观。\(f\) 在点 \(z\) 处的\textbf{伸缩商}（dilatation）为

\[D_f(z) = \frac{|f_z| + |f_{\bar{z}}|}{|f_z| - |f_{\bar{z}}|} = \frac{1 + |\mu(z)|}{1 - |\mu(z)|}\]

\(f\) 是 \(K\)-拟共形的当且仅当 \(D_f(z) \leq K\) 几乎处处成立。\(K = 1\) 时恢复共形映射（\(f_{\bar{z}} = 0\)）。拟共形映射满足 Lusin 性质（将零测集映为零测集），是 ACL（绝对连续线上）的，且是 Hölder 连续的（指数为 \(1/K\)）。

\textbf{定理 129.1}（可测 Riemann 映射定理，Morrey 1938 / Ahlfors-Bers 1960）：设 \(\mu \in L^\infty(\hat{\mathbb{C}})\) 是 \(\hat{\mathbb{C}} = \mathbb{C} \cup \{\infty\}\) 上的可测 Beltrami 系数，满足 \(\|\mu\|_\infty < 1\)。则存在唯一的拟共形同胚 \(f^\mu : \hat{\mathbb{C}} \to \hat{\mathbb{C}}\)，满足 Beltrami 方程 \(f^\mu_{\bar{z}} = \mu(z) f^\mu_z\)（a.e.），且固定 \(0, 1, \infty\)。该映射 \(f^\mu\) 关于 \(\mu\) 是全纯依赖的（在适当拓扑下）：若 \(\mu_t\) 关于 \(t\) 全纯，则 \(f^{\mu_t}\) 亦然。

\textbf{证明}：核心是将 Beltrami 方程转化为奇异积分方程。设 \(f(z) = z + g(z)\)，其中 \(g(z) = O(1/|z|)\)（\(z \to \infty\)）。则 \(f_{\bar{z}} = g_{\bar{z}}\) 和 \(f_z = 1 + g_z\)，Beltrami 方程变为 \(g_{\bar{z}} = \mu(z)(1 + g_z)\)。引入 Cauchy 变换（Bers 变换）\(\mathcal{T} : L^p(\mathbb{C}) \to W^{1,p}(\mathbb{C})\)：

\[(\mathcal{T} \omega)(z) = -\frac{1}{\pi} \iint_{\mathbb{C}} \frac{\omega(\zeta)}{\zeta - z} \, d\xi d\eta \quad (\zeta = \xi + i\eta)\]

\(\mathcal{T}\omega\) 满足 \(\partial_{\bar{z}}(\mathcal{T}\omega) = \omega\)（a.e.）。令 \(g_z = \mathcal{T}(\mu)\)，则 \(g_{\bar{z}} = \mu\)，进而

\[g_{\bar{z}} = \mu (1 + g_z) = \mu (1 + \mathcal{T}(\mu) + \cdots)\]

这给出关于 \(\omega = f_{\bar{z}}\) 的积分方程 \(\omega - \mu \mathcal{T}(\omega) = \mu\)。算子 \(\omega \mapsto \mu \mathcal{T}(\omega)\) 在 \(L^p(\mathbb{C})\) 上当 \(p\) 接近 2 时是压缩的（因 \(\mathcal{T}\) 的 \(L^p\) 范数在 \(p=2\) 处有界），故 \(\omega = (I - \mu\mathcal{T})^{-1}(\mu)\) 在适当的 \(L^p\) 空间中唯一可解，且解 \(\omega\) 全纯依赖于 \(\mu\)。由 \(\omega\) 恢复 \(g = \mathcal{T}(\omega)\)，得 \(f(z) = z + g(z)\) 即所求。\(\blacksquare\)

该定理将经典 Riemann 映射定理推广到可测情形，是 Teichmüller 空间理论和复动力系统的基础。

\subsection{Teichmüller空间}\label{teichmuxfcllerux7a7aux95f4}

\textbf{定义 129.3}（Teichmüller 空间）：设 \(S_g\) 为亏格 \(g \geq 2\) 的紧致定向光滑曲面（参考 Riemann 面）。\textbf{Teichmüller 空间} \(\mathcal{T}_g\) 定义为带标记 Riemann 面的等价类：

\[\mathcal{T}_g = \{(X, f) : f : S_g \to X \text{ 是拟共形同胚（标记）}\} / \sim\]

其中 \((X, f) \sim (Y, g)\) 当且仅当 \(g \circ f^{-1} : X \to Y\) 同伦于一个共形映射（即 biholomorphic 映射）。\(\mathcal{T}_g\) 是 \(3g-3\) 维复流形（实维 \(6g-6\)）。其维数源于 Riemann-Roch 定理：\(S_g\) 上全纯二次微分空间 \(Q(S_g)\) 的维数为 \(3g-3\)，而 \(Q(S_g)\) 恰为 \(T_{[X]}^* \mathcal{T}_g\)（余切空间）。

\textbf{定义 129.4}（Teichmüller 度量）：对 \([X], [Y] \in \mathcal{T}_g\)，\textbf{Teichmüller 度量}（Teichmüller 距离）定义为

\[d_T([X], [Y]) = \frac{1}{2} \inf_{f} \log K(f)\]

其中下确界取遍所有与标记相容的拟共形同胚 \(f : X \to Y\)，\(K(f) = \|D_f\|_\infty\) 为 \(f\) 的最大伸缩商。由 Teichmüller 唯一性定理，存在唯一的极值拟共形映射 \(f_0\) 达到此下确界，其 Beltrami 系数形如 \(\mu_0 = k \overline{\phi}/|\phi|\)（\(0 \leq k < 1\)，\(\phi\) 为 \(X\) 上的全纯二次微分），此映射称为 \textbf{Teichmüller 映射}。\((\mathcal{T}_g, d_T)\) 是完备的 Finsler 流形，其 Finsler 范数为 \(L^1\) 型（称为 Teichmüller-Finsler 度量）。

\textbf{定理 129.2}（Bers 嵌入定理，1960）：存在从 \(\mathcal{T}_g\) 到 \(\mathbb{C}^{3g-3}\) 的单射全纯嵌入 \(\Phi : \mathcal{T}_g \hookrightarrow \mathbb{C}^{3g-3}\)，使得像 \(\Phi(\mathcal{T}_g)\) 是 \(\mathbb{C}^{3g-3}\) 中的有界域。具体地，设 \(\Gamma\) 为 \(S_g\) 的万有覆叠 \(\mathbb{H}\)（上半平面）上的 Fuchs 群，则

\[\Phi([X]) = \left\{ z \mapsto \frac{\phi'''(z)}{\phi'(z)} - \frac{3}{2}\left(\frac{\phi''(z)}{\phi'(z)}\right)^2 \right\}\]

其中 \(\phi : \mathbb{H} \to \mathbb{H}\) 为诱导自拟共形映射 \(f : S_g \to X\) 的提升映射的 Schwarz 导数。Bers 嵌入将 \(\mathcal{T}_g\) 实现为全纯二次微分空间中的有界域，揭示了 \(\mathcal{T}_g\) 的复结构。

\textbf{定理 129.3}（Teichmüller 存在性与唯一性定理）：设 \(X, Y\) 为亏格 \(g\) 的紧 Riemann 面。则：(i) 在 \(X\) 到 \(Y\) 的拟共形同胚所构成的每一个同伦类中，存在唯一的极值拟共形映射，其 Beltrami 系数形如 \(\mu = k \overline{\phi}/|\phi|\)（\(k \in [0,1)\)，\(\phi\) 为 \(X\) 上的非零全纯二次微分）；(ii) 该极值映射的伸缩商处处为常数 \(K = (1+k)/(1-k)\)，且其水平叶（即 \(\operatorname{Re} \int \sqrt{\phi} dz = \text{const}\)）沿映射方向拉伸 \(K\) 倍，垂直叶收缩 \(1/K\) 倍。

模空间 \(\mathcal{M}_g = \mathcal{T}_g / \operatorname{MCG}(S_g)\) 是 \(3g-3\) 维复轨道空间，其中 \(\operatorname{MCG}(S_g)\) 是 \(S_g\) 的映射类群。\(\mathcal{M}_g\) 不是紧致的，但其 Deligne-Mumford 紧化 \(\overline{\mathcal{M}}_g\) 是一个射影代数簇。

\hrulefill

\section{第167章：Siegel区域与对称域}\label{ux7b2c167ux7ae0siegelux533aux57dfux4e0eux5bf9ux79f0ux57df}

Siegel区域理论是经典多复变函数论的核心部分。\textbf{Siegel上半空间}\(\mathcal{H}_n = \{ Z \in \mathbb{C}^{n \times n} : Z^{\mathsf{T}} = Z, \; \operatorname{Im}(Z) \succ 0 \}\)是\(\operatorname{Sp}(2n,\mathbb{R})\)的作用空间。一般\textbf{Siegel区域}为\(\{ (z,w) \in \mathbb{C}^{m} \times \mathbb{C}^{n} : \operatorname{Im}(z) - F(w,w) \in \Omega \}\)，其中\(F\)为Hermite形式，\(\Omega \subset \mathbb{R}^m\)为自对偶锥。Elie Cartan (1935)将不可约有界对称域分为四大类经典域和两个例外域（16维和27维）。Siegel模形式（\(\operatorname{Sp}(2n,\mathbb{Z})\)作用下变换的全纯函数）推广了经典椭圆模形式，在志村簇和弦论中有深刻应用。

\subsection{Cartan分类的完整描述}\label{cartanux5206ux7c7bux7684ux5b8cux6574ux63cfux8ff0}

Elie Cartan于1935年完成了不可约有界对称域的完全分类。一个\textbf{有界对称域}（bounded symmetric domain）\(D \subset \mathbb{C}^N\) 是有界全纯域，且对每点 \(p \in D\) 存在对合全纯自同构 \(s_p\) 使 \(p\) 为孤立不动点且 \(s_p^2 = \mathrm{id}\)。所有不可约有界对称域恰好分为六大类：

\textbf{四大类经典域}（Cartan 经典域）：

\textbf{(I) \(I_{p,q}\) 型}（\(p \leq q\)）：由 \(p \times q\) 复矩阵构成，满足 \(I - Z Z^* \succ 0\)（正定）。实维数 \(2pq\)。当 \(p = 1\) 时退化为 \(\mathbb{C}^q\) 中的单位球。

\textbf{(II) \(II_n\) 型}：由 \(n \times n\) 反对称复矩阵构成，满足 \(I - Z \overline{Z} \succ 0\)。实维数 \(n(n-1)\)。

\textbf{(III) \(III_n\) 型}：由 \(n \times n\) 对称复矩阵构成，满足 \(I - Z \overline{Z} \succ 0\)。实维数 \(n(n+1)\)。Siegel 上半空间 \(\mathcal{H}_n\) 通过 Cayley 变换 \(Z \mapsto (Z - iI)(Z + iI)^{-1}\) 全纯等价于此类型的有界域。

\textbf{(IV) \(IV_n\) 型}（\(n \geq 5\)）：Lie 球（Lie sphere），由 \(\mathbb{C}^n\) 中的向量满足 \(|z^T z|^2 + 1 - 2 \overline{z}^T z > 0\) 且 \(|z^T z| < 1\) 定义。实维数 \(2n\)。

\textbf{两个例外域}：

\textbf{(V) 例外域 \(E_{16}\)}（16维）：对应于例外 Jordan 代数 \(\operatorname{Herm}(3, \mathbb{O})\)（\(3 \times 3\) 八元数 Hermite 矩阵构成），与例外 Lie 群 \(E_{6(-26)}\) 的 Hermite 对称空间对应。其 Bergman-Shilov 边界由秩为 1 的八元数 Hermite 投影矩阵构成，与八元数射影平面 \(\mathbb{O}P^2\) 同胚。

\textbf{(VI) 例外域 \(E_{27}\)}（27维）：对应于例外 Jordan 代数 \(\operatorname{Herm}(3, \mathbb{O})\) 的复化，与例外 Lie 群 \(E_{7(-25)}\) 的 Hermite 对称空间对应。这是维数最高的不可约有界对称域，其边界结构涉及 27 维的例外 Jordan 代数中的行列式三次型。

Cartan 分类的数学意义在于，不可约有界对称域与\textbf{Hermite 对称空间}的非紧型一一对应，且与半单 Lie 代数的\textbf{Hermite 对称对}分类完全吻合。Hermite 对称空间的紧型可以通过 Borel 嵌入定理从有界域实现中恢复。

\subsection{Siegel上半空间的Bergman核}\label{siegelux4e0aux534aux7a7aux95f4ux7684bergmanux6838}

Siegel 上半空间 \(\mathcal{H}_n = \{Z \in \mathbb{C}^{n \times n} : Z = Z^T, \; \operatorname{Im} Z \succ 0\}\) 上的 \textbf{Bergman 核} 由下式给出（华罗庚，1946年）： \[K_{\mathcal{H}_n}(Z, W) = c_n \det\left( \frac{Z - \overline{W}}{2i} \right)^{-(n+1)}\] 其中 \(c_n\) 为归一化常数。这是单位圆盘上 Bergman 核 \(K_D(z,w) = (1/\pi) (1 - z\bar{w})^{-2}\) 的高维矩阵推广------行列式因子替换了标量因子，指数 \(n+1\) 来自 \(\mathcal{H}_n\) 作为 \(\operatorname{Sp}(2n,\mathbb{R})\) 齐性空间的几何维数。在 Cayley 变换下，\(K_{\mathcal{H}_n}\) 变换为有界域形式（Cartan III 型域），其 Bergman 核为 \[K_{D}(Z, W) = c_n' \det(I - Z \overline{W})^{-(n+1)}\]

Bergman 核的\textbf{变换公式}：对 \(\mathcal{H}_n\) 的自同构群 \(\operatorname{Sp}(2n,\mathbb{R})\) 的作用 \(Z \mapsto g \cdot Z = (AZ + B)(CZ + D)^{-1}\)（其中 \(\begin{pmatrix} A & B \\ C & D \end{pmatrix} \in \operatorname{Sp}(2n, \mathbb{R})\)），Bergman 核的变换规则为 \[K(g \cdot Z, g \cdot W) = |\det(CZ + D)|^{n+1} K(Z, W) |\det(C \overline{W} + D)|^{n+1}\] 这一变换公式反映了 Bergman 核作为自同构群的同变（equivariant）线丛截面的几何本质。对应的 Bergman 度量 \(ds^2 = \partial \overline{\partial} \log K(Z, Z)\) 是 \(\operatorname{Sp}(2n,\mathbb{R})\)-不变的 Kähler 度量，赋予 \(\mathcal{H}_n\) 非正截面曲率的完备 Kähler-Einstein 度量结构。

\subsection{Harish-Chandra嵌入定理}\label{harish-chandraux5d4cux5165ux5b9aux7406}

\textbf{定理130.1}（Harish-Chandra嵌入定理，1956）：设 \(D = G/K\) 为有界对称域，其中 \(G\) 为 \(D\) 的全纯自同构群的单位连通分支，\(K\) 为某点的各向同性子群。令 \(\mathfrak{g} = \mathfrak{k} \oplus \mathfrak{p}\) 为相应的 Cartan 分解，\(\mathfrak{g}^{\mathbb{C}}\) 为复化 Lie 代数。则 \(D\) 可通过 \textbf{Harish-Chandra 嵌入} \(D \hookrightarrow \mathfrak{p}^+\)（\(\mathfrak{p}^+\) 为 \(\mathfrak{p}\) 在 \(K^{\mathbb{C}}\) 作用下的全纯切空间）全纯等价于 \(\mathfrak{p}^+\) 中包含原点的有界对称域。具体地， \[D \cong \{X \in \mathfrak{p}^+ : I - \varphi(X, \overline{X}) \succ 0\}\] 其中 \(\varphi\) 为由 Lie 代数的换位子构造的双线性形式。

\textbf{证明思路}：由 Borel 嵌入定理，\(D\) 可全纯嵌入其紧对偶 \(G^c/K\)（\(G^c\) 为紧 Lie 群）。紧对偶中包含原点的坐标邻域由指数映射 \(\exp : \mathfrak{p}^+ \to G^{\mathbb{C}}\) 参数化，而 \(D\) 在其中的像由上述 Hermite 不等式刻画。更具体地，考虑 Iwasawa 分解 \(G^{\mathbb{C}} = K^{\mathbb{C}} A^+ N^+\)，\(D\) 恰为 \(G/K\) 在 \(G^{\mathbb{C}} / K^{\mathbb{C}} P^-\) 中的开嵌入的像（Borel 嵌入）。对 Cartan III 型域（Siegel 上半空间），Harish-Chandra 嵌入还原为 Cayley 变换。\(\blacksquare\)

Harish-Chandra 嵌入定理将任意有界对称域实现为某一复向量空间中由矩阵不等式定义的有界凸域------这一实现被称为 \textbf{Harish-Chandra 实现}，它是处理有界对称域上的分析问题（如 Bergman 核计算、自守形式理论）的标准化工具。

\subsection{有界对称域与Jordan三系的对应关系}\label{ux6709ux754cux5bf9ux79f0ux57dfux4e0ejordanux4e09ux7cfbux7684ux5bf9ux5e94ux5173ux7cfb}

\textbf{定义130.1}（Jordan三系）：复向量空间 \(V\) 上的 \textbf{Jordan 三系}（Jordan triple system / JTS）是一个三元运算 \(\{\cdot, \cdot, \cdot\} : V \times V \times V \to V\)，满足对称性 \(\{x, y, z\} = \{z, y, x\}\) 和 Jordan 恒等式 \[\{a, b, \{x, y, z\}\} = \{\{a, b, x\}, y, z\} - \{x, \{b, a, y\}, z\} + \{x, y, \{a, b, z\}\}\]

\textbf{定理130.2}（Koecher-华罗庚定理，1957）：不可约有界对称域 \(D \subset \mathbb{C}^N\) 与正定（positive definite）Hermite Jordan 三系一一对应。具体地，\(D\) 的 Harish-Chandra 实现由 \(\{V, \{\cdot, \cdot, \cdot\}\}\) 通过不等式 \(\{z, z, z\} \prec I\) 定义，其中 \(V = \mathbb{C}^N\) 且 \(D\) 上的全纯自同构群由 Jordan 三系的\textbf{结构群}（structure group）给出。

这一对应关系是有界对称域理论中最具结构深度的结果。秩为 \(r\) 的不可约有界对称域可以借助 Jordan 三系的\textbf{Pierce 分解}分解为 \(1 + r(r+1)/2\) 个子空间的和，对应于不同权重空间。在 Cartan 分类中，\(I_{p,q}\) 型对应于矩形矩阵的 Jordan 三系 \(\{X, Y, Z\} = X Y^* Z + Z Y^* X\)，\(II_n\) 型对应于反对称矩阵的 Jordan 三系（带相应修正），\(IV_n\) 型对应于 spin 因子的 Jordan 代数，两个例外域对应于八元数 Jordan 代数的复化。Jordan 三系的框架统一了 Cartan 分类、Bergman 核公式、Shilov 边界的几何结构以及自守因子（automorphy factor）的代数构造，使其成为有界对称域理论的代数核心。

\hrulefill

\section{第168章：多复变函数论导引}\label{ux7b2c168ux7ae0ux591aux590dux53d8ux51fdux6570ux8bbaux5bfcux5f15}

多复变函数论研究\(\mathbb{C}^n\)（\(n \geq 2\)）上的全纯函数，与单复变有本质区别：孤立奇点可去、单位球与多圆柱不全纯等价。

\subsection{全纯域与Hartogs延拓}\label{ux5168ux7eafux57dfux4e0ehartogsux5ef6ux62d3}

\textbf{定义111.1}（全纯域）：开集\(\Omega \subset \mathbb{C}^n\)称为\textbf{全纯域}（Domain of Holomorphy），若存在\(\Omega\)上的全纯函数\(f\)，使得\(f\)不能全纯延拓到\(\Omega\)外任何点。即对任意点\(p \in \partial\Omega\)和任意邻域\(U \ni p\)，不存在\(F \in \mathcal{O}(\Omega \cup U)\)使\(F|_\Omega = f\)。

\textbf{定理111.1}（Hartogs延拓定理，1906）：设\(n \geq 2\)，\(\Omega \subset \mathbb{C}^n\)为开集，\(K \subset \Omega\)为紧子集使得\(\Omega \setminus K\)连通。则\(\Omega \setminus K\)上的任意全纯函数可唯一地全纯延拓到整个\(\Omega\)。特别地，\(\mathbb{C}^n\)的孤立奇点（\(n \geq 2\)）总是可去的。

\textbf{证明}：利用Cauchy积分公式在''管状''区域上构造延拓。设\(f \in \mathcal{O}(\Omega \setminus K)\)。取充分大的开球\(B\)使\(K \subset B \subset\subset \Omega\)，定义 \[\tilde{f}(z) = \frac{1}{2\pi i} \int_{|\zeta| = R} \frac{f(\zeta, z')}{\zeta - z_1} d\zeta\] 其中\(z = (z_1, z')\)且积分路径在\(K\)之外。由于\(n \geq 2\)，在复余一维的子簇上可绕过\(K\)选择积分围道。\(\tilde{f}\)在\(B\)上定义且全纯，且在\(B \setminus K\)上与\(f\)一致。由唯一性定理，\(\tilde{f}\)即为所求延拓。\(\blacksquare\)

\textbf{推论111.2}（孤立奇点可去）：对\(n \geq 2\)，若\(f\)在\(B(0,R) \setminus \{0\}\)上全纯，则\(f\)自动延拓为\(B(0,R)\)上的全纯函数。因此\(\mathbb{C}^n\)（\(n \geq 2\)）中不存在类似单变量的极点和本性奇点。

\textbf{定理111.1'}（Hartogs基本定理 --- 全纯性的变量分离判别）：函数\(f : \Omega \subset \mathbb{C}^n \to \mathbb{C}\)全纯当且仅当\(f\)对每个变量单独全纯（无需连续性假设）。

\textbf{证明}：单独全纯推出局部有界性（Osgood引理配合Hartogs延拓），再由Cauchy积分公式对每个变量展开可得联合连续性，最终导出联合全纯性。\(\blacksquare\)

\hrulefill

\section{第169章：多复变核心定理}\label{ux7b2c169ux7ae0ux591aux590dux53d8ux6838ux5fc3ux5b9aux7406}

多复变函数论的核心定理构成了该领域的基本骨架，它们将单复变中的经典结果推广到高维------或揭示高维独有的全新现象。本章聚焦于两个核心方向：其一，Stein 流形理论与 Cartan-Thullen 定理，建立了全纯域、全纯凸性、Stein 性和伪凸性之间的等价关系，其中 Oka 对 Levi 问题的解答是 20 世纪多复变理论的里程碑；其二，\(\bar{\partial}\)-问题的 \(L^2\) 估计（Hörmander, 1965），它不仅是多复变正则性理论的核心技术工具，也是现代复几何（如 Ohsawa-Takegoshi \(L^2\) 延拓定理）和代数几何（如 Kawamata-Viehweg 消灭定理的解析证明）的基础。这些定理共同揭示了多复变中全纯函数的内在刚性------Hartogs 延拓现象和 \(\bar{\partial}\)-Neumann 问题的次椭圆性正是这种刚性的体现。

\subsection{Stein流形与Cartan-Thullen定理}\label{steinux6d41ux5f62ux4e0ecartan-thullenux5b9aux7406}

\textbf{定义112.1}（Stein流形）：复流形\(X\)为\textbf{Stein流形}若：(1) \(X\)为Hausdorff且第二可数；(2) 全纯函数分离点；(3) 全纯函数给出局部坐标；(4) \(X\)是全纯凸的------对任意紧集\(K \subset X\)，其全纯凸包\(\hat{K} = \{x \in X : |f(x)| \leq \sup_K |f|, \forall f \in \mathcal{O}(X)\}\)为紧。

\textbf{定理112.1}（Cartan-Thullen定理，1932）：对\(\mathbb{C}^n\)中的域\(\Omega\)，以下等价： (i) \(\Omega\)是全纯域； (ii) \(\Omega\)是全纯凸域：对任意紧集\(K \subset \Omega\)，其全纯凸包\(\hat{K}_\Omega = \{z \in \Omega : |f(z)| \leq \sup_K |f|, \forall f \in \mathcal{O}(\Omega)\}\)为\(\Omega\)中紧集且\(\operatorname{dist}(\hat{K}_\Omega, \partial\Omega) = \operatorname{dist}(K, \partial\Omega)\)； (iii) \(\Omega\)是Stein流形； (iv) \(-\log \operatorname{dist}(z, \partial\Omega)\)为多次调和函数（即\(\Omega\)是伪凸的）。

\textbf{证明}：(i) \(\Rightarrow\) (ii)：若\(\Omega\)为全纯域，设\(f\)不能延拓过\(\partial\Omega\)，其极点集迫使全纯凸包远离边界。(ii) \(\Rightarrow\) (iii)：全纯凸性配合全纯函数分离点的性质（通过构造适当全纯函数证明）。(iii) \(\Rightarrow\) (i)：Stein流形上存在不可延拓的全纯函数可由Cousin问题解的存在性保证。(iii) \(\Leftrightarrow\) (iv)：这是Oka (1942) 对Levi问题的解答------伪凸性等价于全纯凸性等价于Stein性。\(\blacksquare\)

\textbf{定理112.2}（Oka-Weil逼近定理）：若\(X\)为Stein流形，\(K \subset X\)为全纯凸紧集，则\(K\)某邻域上的任意全纯函数可由\(X\)上全纯函数一致逼近。

\textbf{证明}：利用Stein流形上的Cousin I问题解构造具指定极点的亚纯函数，再转化为全纯函数的逼近------核心是Runge型逼近在Stein流形上的推广。\(\blacksquare\)

\subsection{\texorpdfstring{\(\bar{\partial}\)-问题与\(L^2\)估计}{\textbackslash bar\{\textbackslash partial\}-问题与L\^{}2估计}}\label{barpartial-ux95eeux9898ux4e0el2ux4f30ux8ba1}

\textbf{定义112.2}（\(\bar{\partial}\)-问题）：设\(\Omega \subset \mathbb{C}^n\)为开集，\(f\)为\((0,1)\)-形式（即\(f = \sum_{j=1}^n f_j d\bar{z}_j\)）满足可积条件\(\bar{\partial} f = 0\)（即\(\partial f_j / \partial \bar{z}_k = \partial f_k / \partial \bar{z}_j\)）。\(\bar{\partial}\)-问题是求函数\(u\)使 \[\bar{\partial} u = f, \quad \text{即} \quad \frac{\partial u}{\partial \bar{z}_j} = f_j \quad (j = 1,\dots,n)\] 这是多复变理论的核心技术问题------其可解性等价于域的伪凸性。

\textbf{定理112.3}（Hörmander \(L^2\)估计，1965）：设\(\Omega \subset \mathbb{C}^n\)为伪凸域（即存在光滑严格多次调和穷竭函数），\(\varphi \in C^2(\Omega)\)为多次调和函数。则对任意满足\(\bar{\partial}f = 0\)的\((0,1)\)-形式\(f\)，存在\(u\)满足\(\bar{\partial} u = f\)且 \[\int_\Omega |u|^2 e^{-\varphi} dV \leq \int_\Omega \sum_{j,k} \varphi^{j\bar{k}} f_j \overline{f_k} e^{-\varphi} dV\] 其中\(\varphi^{j\bar{k}}\)为复Hessian矩阵\((\partial^2\varphi/\partial z_j \partial \bar{z}_k)\)的逆。特别地，取\(\varphi = 0\)在\(\mathbb{C}^n\)的有界伪凸域上，得\(\|u\|_{L^2} \leq C(\Omega) \|f\|_{L^2}\)。

\textbf{证明}：核心为加权\(L^2\)空间上的\(\bar{\partial}\)-Neumann算子的Hörmander-Morrey估计。对于光滑\(\varphi\)，定义Hilbert空间\(L^2(\Omega, e^{-\varphi})\)中的闭稠定算子\(\bar{\partial}\)，形式伴随\(T^*\)，对任意\(C_c^\infty\)函数\(u\)成立Morrey-Kohn恒等式 \[\|\bar{\partial}u\|^2 + \|\bar{\partial}^*_\varphi u\|^2 = \sum_{j,k} \int_\Omega \frac{\partial^2\varphi}{\partial z_j \partial \bar{z}_k} u_j \overline{u_k} e^{-\varphi} + \sum_{j,k} \int_\Omega \left|\frac{\partial u_j}{\partial \bar{z}_k}\right|^2 e^{-\varphi}\] 其中\(\bar{\partial}^*_\varphi\)为加权的形式伴随。当\(\varphi\)为多次调和时漏掉部分非负，得\(\|\bar{\partial}^*_\varphi u\|^2 \geq \sum\int \varphi_{j\bar{k}} u_j \overline{u_k} e^{-\varphi}\)。通过Hahn-Banach与Riesz表示定理将不等式转化为满射性，即\(\bar{\partial}\)的像包含\(\ker\bar{\partial}\)的正交补，故对\(\bar{\partial}\)-闭形式\(f\)方程可解。\(\blacksquare\)

\textbf{定理112.4}（\(\bar{\partial}\)-Neumann问题的正则性）：在伪凸域\(\Omega\)上，若\(f\)光滑且满足\(\bar{\partial}f = 0\)，则规范解\(u = \bar{\partial}^* N f\)（\(N\)为\(\bar{\partial}\)-Neumann算子）是光滑的，且与\(\Omega\)边界性质相容。特别地，在强伪凸域上，\(\bar{\partial}\)-Neumann算子为紧算子，\(\bar{\partial}\)-问题具有最优的正则性估计。

\textbf{证明}：基于\(\bar{\partial}\)-Neumann算子的椭圆正则性理论（Kohn, 1963-1964）。在强伪凸边界条件下，\(\Box = \bar{\partial}\bar{\partial}^* + \bar{\partial}^*\bar{\partial}\)为次椭圆算子，满足\(1/2\)阶Garding型估计，其逆\(N\)将Sobolev空间\(H^s\)映射到\(H^{s+1}\)（在\(\Omega\)内部增益\(2\)阶，在边界处增益\(1\)阶）。\(\blacksquare\)

\hrulefill

\emph{Nevanlinna理论以Poisson-Jensen公式为出发点，经特征函数与两个基本定理揭示了亚纯函数值分布的深层结构。多复变方面，Hartogs延拓定理彰显了\(n \geq 2\)时全纯函数的刚性------全纯域的概念由此自然引入。Cartan-Thullen定理建立了全纯域-全纯凸-Stein流形-伪凸域的四重等价。\(\bar{\partial}\)-问题及其\(L^2\)估计（Hörmander）则为多复变正则性理论提供了核心技术工具。}

\emph{卷十：复分析终。}

\hrulefill

\section{第170章：共形映射}\label{ux7b2c170ux7ae0ux5171ux5f62ux6620ux5c04}

共形映射理论是复分析几何化的核心篇章。全纯函数在导数非零处保持曲线夹角（保角性）和无穷小圆的形状（保圆性），使得复分析成为二维共形几何的天然语言。本章从共形映射的定义出发，系统研究分式线性变换（Möbius 变换）的群论性质及其保圆性，进而介绍 Riemann 映射定理------任何单连通真子区域共形等价于单位圆盘------以及 Schwarz-Christoffel 变换在多角形映射中的应用。共形映射理论在流体力学、电磁学和弹性力学中有广泛的应用，其在多复变中的失效（Poincare 发现）也开启了全新的研究方向。

\subsection{55.1 共形映射的定义与性质}\label{ux5171ux5f62ux6620ux5c04ux7684ux5b9aux4e49ux4e0eux6027ux8d28}

\textbf{定义 55.1.1（共形映射）}：全纯函数 \(f : \Omega \to \mathbb{C}\) 在 \(z_0 \in \Omega\) 处称为\textbf{共形}的，若 \(f'(z_0) \neq 0\)。若 \(f\) 在 \(\Omega\) 上每点共形且为单射，则称 \(f\) 为 \(\Omega\) 上的\textbf{共形映射}。

\textbf{定理 55.1.1（共形映射的几何意义）}：\(f\) 在 \(z_0\) 处共形当且仅当 \(f\) 在 \(z_0\) 处保角（两曲线在 \(z_0\) 处的交角等于其像曲线在 \(f(z_0)\) 处的交角）且伸缩率 \(|f'(z_0)|\) 与方向无关。

\textbf{证明}：由 \(f(z) - f(z_0) \approx f'(z_0)(z - z_0)\)（\(z \to z_0\)）。\(f'(z_0) = r e^{i\theta}\) 将 \(z - z_0\) 旋转 \(\theta\) 并缩放 \(r\) 倍，保角且伸缩率与方向无关。\(\blacksquare\)

共形映射的几何意义在复分析中具有根本重要性。在二维情形下，共形映射恰好是保定向的保角映射------这是 Cauchy-Riemann 方程的直接几何推论。从微分几何的角度看，共形映射诱导了两个 Riemann 度量之间的共形等价：\(f^* (du^2 + dv^2) = |f'(z)|^2 (dx^2 + dy^2)\)。这一事实使得共形映射在曲面上的等温坐标（isothermal coordinates）理论中扮演核心角色。任意 Riemann 面的局部坐标都是共形坐标，这一性质是复分析区别于实分析的基本特征。

\hrulefill

\subsection{55.2 分式线性变换}\label{ux5206ux5f0fux7ebfux6027ux53d8ux6362}

\textbf{定义 55.2.1（分式线性变换）}：形如

\[
T(z) = \frac{az + b}{cz + d}, \quad ad - bc \neq 0
\]

的映射称为\textbf{分式线性变换}（Mobius 变换）。它可延拓为 \(\hat{\mathbb{C}} = \mathbb{C} \cup \{\infty\}\)（Riemann 球面）上的双射。

\textbf{定理 55.2.1（分式线性变换的基本性质）}：

\begin{enumerate}
\def\labelenumi{(\roman{enumi})}
\tightlist
\item
  分式线性变换构成群 \(PSL(2, \mathbb{C})\)
\item
  分式线性变换将圆/直线映为圆/直线（保圆性）
\item
  分式线性变换保交比：\(\frac{T(z_1) - T(z_3)}{T(z_1) - T(z_4)} \cdot \frac{T(z_2) - T(z_4)}{T(z_2) - T(z_3)} = \frac{z_1 - z_3}{z_1 - z_4} \cdot \frac{z_2 - z_4}{z_2 - z_3}\)
\item
  存在唯一的分式线性变换将任意三个不同的点映为任意三个不同的点
\end{enumerate}

\textbf{证明}：(i) 复合和逆的公式验证。(ii) 圆/直线方程可写为 \(A z \overline{z} + B z + \overline{B} \overline{z} + C = 0\)（\(A, C \in \mathbb{R}\)），分式线性变换保持此形式。(iii) 直接验证。(iv) 由交比的不变性构造。\(\blacksquare\)

分式线性变换是共形映射中最基本的一类，其重要性在于它们是 Riemann 球面 \(\hat{\mathbb{C}}\) 上仅有的全局共形自同构。事实上，\(\operatorname{Aut}(\hat{\mathbb{C}}) = PSL(2, \mathbb{C})\)------所有保持 \(\hat{\mathbb{C}}\) 的双射全纯映射都是分式线性变换。这一结果可由 Liouville 定理的推广证明：任何整函数若为双射，则必为一次多项式。分式线性变换的保圆性有深刻的几何背景：在球极投影（stereographic projection）下，\(\hat{\mathbb{C}}\) 上的圆对应于 \(\mathbb{R}^3\) 中球面 \(S^2\) 上的圆，而 \(PSL(2, \mathbb{C})\) 恰好对应于 \(S^2\) 上保圆的共形变换群------即 Lorentz 群 \(SO^+(3, 1)\) 在 \(S^2\) 上的作用。这一联系是三维双曲几何与复分析之间的桥梁，也是 Thurston 几何化猜想中双曲几何部分的出发点。

\hrulefill

\subsection{55.3 Riemann 映射定理}\label{riemann-ux6620ux5c04ux5b9aux7406}

\textbf{定理 55.3.1（Riemann 映射定理------陈述）}：设 \(\Omega \subsetneq \mathbb{C}\) 为单连通区域。则存在共形映射 \(f : \Omega \to D = \{z \in \mathbb{C} \mid |z| < 1\}\)（单位圆盘）。进一步，若指定 \(z_0 \in \Omega\) 和 \(f(z_0) = 0\)，\(f'(z_0) > 0\)，则 \(f\) 是唯一的。

\textbf{注记 55.3.1}：Riemann 映射定理是复分析的核心定理之一，证明了任何单连通真子区域与单位圆盘共形等价。这是二维特有的现象------高维情况不成立。完整证明需要正规族理论和 Schwarz 引理。

\textbf{定理 55.3.2（Riemann 映射定理的证明）}：

\textbf{证明}：设 \(\mathcal{F}\) 为 \(\Omega\) 上满足 \(f(z_0) = 0\) 且 \(f'(z_0) > 0\) 的单射全纯函数族，其像含于 \(D\)。步骤一：\(\mathcal{F} \neq \varnothing\)。由于 \(\Omega \neq \mathbb{C}\)，存在 \(a \notin \Omega\)。取 \(\Omega\) 上的平方根函数 \(\sqrt{z - a}\)（由 \(\Omega\) 单连通，存在全纯的单值分支），其像是某个半平面，用分式线性变换将其映为 \(D\)，再复合以适当的旋转和平移使 \(f(z_0) = 0\) 且 \(f'(z_0) > 0\)。步骤二：考虑极值问题 \(\sup_{f \in \mathcal{F}} f'(z_0)\)。由 Cauchy 估计，\(\mathcal{F}\) 在紧集上一致有界（因为所有函数取值于 \(D\)），故由 Montel 正规族定理，存在子列 \(\{f_n\}\) 在紧集上一致收敛到某个全纯函数 \(f\)。由 Hurwitz 定理，\(f\) 仍为单射（或常数）。由于 \(f'(z_0) = \sup_{g \in \mathcal{F}} g'(z_0) > 0\)，\(f\) 非常数，故 \(f\) 为单射且 \(f(\Omega) \subseteq D\)。步骤三：证明 \(f\) 是满射。若 \(f(\Omega) \neq D\)，则存在 \(w_0 \in D \setminus f(\Omega)\)。构造 \(D\) 的自同构 \(\varphi_{w_0}(z) = \frac{z - w_0}{1 - \overline{w_0}z}\)，将 \(w_0\) 映为 \(0\)。在 \(\varphi_{w_0}(f(\Omega))\) 上取平方根函数（因其不含 \(0\) 且为单连通），复合 \(\varphi_{w_0}\) 的逆后可得 \(F \in \mathcal{F}\) 使 \(F'(z_0) > f'(z_0)\)，与极值性矛盾。故 \(f\) 为满射，即 \(f: \Omega \to D\) 为共形映射。唯一性：若 \(f, g\) 均满足条件，则 \(h = f \circ g^{-1}: D \to D\) 为 \(D\) 的自同构，且 \(h(0) = 0\)，\(h'(0) > 0\)。由 Schwarz 引理，\(h(z) = e^{i\theta} z\)，结合 \(h'(0) > 0\) 得 \(\theta = 0\)，故 \(h(z) = z\)，即 \(f = g\)。\(\blacksquare\)

Riemann 映射定理的历史意义非凡。Riemann 在 1851 年的博士论文中首次陈述了该定理，但其证明依赖于 Dirichlet 原理（存在最小化 Dirichlet 能量的函数），而 Dirichlet 原理的严格性直到 Hilbert 在 1900 年左右才得以确立。上述证明（使用正规族和极值问题）属于 Fejer-Riesz（1922）和 Koebe（1907），是函数论方法而非变分方法的典范。Riemann 映射定理直接导致了单值化定理（Uniformization Theorem）：任何单连通 Riemann 面共形等价于 \(\hat{\mathbb{C}}\)、\(\mathbb{C}\) 或 \(D\) 三者之一。单值化定理是 Riemann 面理论的核心，也是证明代数曲线分类定理（亏格 \(g\) 的紧 Riemann 面一一对应于光滑射影代数曲线）的基础。

\hrulefill

\subsection{55.4 Schwarz-Christoffel 变换与多角形映射}\label{schwarz-christoffel-ux53d8ux6362ux4e0eux591aux89d2ux5f62ux6620ux5c04}

\textbf{定理 55.4.1（Schwarz-Christoffel 公式）}：将上半平面 \(\mathbb{H} = \{z \in \mathbb{C} : \Im(z) > 0\}\) 共形映射到多边形区域 \(\Omega\)（内角为 \(\alpha_1\pi, \alpha_2\pi, \ldots, \alpha_n\pi\)，\(\sum \alpha_k = n-2\)）的映射为

\[f(z) = A \int_{z_0}^z \prod_{k=1}^n (w - x_k)^{\alpha_k - 1} \, dw + B\]

其中 \(x_1 < x_2 < \cdots < x_n\) 为实轴上对应多边形顶点的原像，\(A, B\) 为常数。

\textbf{证明}：由 Schwarz 反射原理和共形映射的边界对应。\(f\) 将 \(\mathbb{R}\) 映为多边形边界。在 \((x_k, x_{k+1})\) 上，\(\arg f'(z)\) 为常数（因 \(f\) 将实轴区间映为直线段）。在 \(x_k\) 处，\(\arg f'(z)\) 的跳跃量为 \((1-\alpha_k)\pi\)（即外角），这由因子 \((z-x_k)^{\alpha_k-1}\) 的辐角变化实现。积分后即得整个多边形。\(\blacksquare\)

Schwarz-Christoffel 变换是复分析中连接共形映射理论与位势论（Dirichlet 问题在复杂区域上的求解）的桥梁。其数值实现（SC Toolbox）在现代工程中仍有重要应用------如计算翼型周围的流体绕流、电磁场在复杂几何中的分布等。

Schwarz-Christoffel 公式的更深层意义在于，它给出了上半平面到任意多边形的共形映射的显式参数化。参数 \(\{x_k\}\)（称为 SC 参数或原像）并非自由的------它们满足一定的约束条件（由多边形的边长决定），这导致了 Schwarz-Christoffel 参数问题：给定多边形顶点的坐标，如何确定 \(\{x_k\}\) 和常数 \(A, B\)？这一问题的求解通常需要数值方法，但存在性由 Riemann 映射定理保证。对于 \(n = 3\)（三角形），SC 积分化为 Schwarz-Christoffel 三角映射，与超几何函数密切相关。对于矩形（\(n = 4\)，\(\alpha_k = 1/2\)），SC 积分化为椭圆积分，其逆映射是 Jacobi 椭圆函数 \(\operatorname{sn}(z)\)------这揭示了共形映射与椭圆函数之间的深刻联系。

多角形映射理论还可推广到圆弧多边形（边界由圆弧组成）和更一般的 Schwarz 导数参数化。Schwarz 导数 \(Sf = (f''/f')' - \frac{1}{2}(f''/f')^2\) 在共形映射的复合下满足链式法则 \(S(f \circ g) = (Sf \circ g) \cdot (g')^2 + Sg\)，这使得 Schwarz 导数成为研究共形映射边界行为和微分方程可积性的核心工具。

\subsection{55.5 共形映射的边界行为}\label{ux5171ux5f62ux6620ux5c04ux7684ux8fb9ux754cux884cux4e3a}

\textbf{定理 55.5.1（Caratheodory 边界对应定理，1913）}：设 \(\Omega\) 为有界单连通区域，\(f: D \to \Omega\) 为 Riemann 映射。\(f\) 可连续延拓到闭圆盘 \(\overline{D}\) 当且仅当 \(\partial\Omega\) 是 Jordan 曲线（即简单闭曲线）。此时 \(f: \overline{D} \to \overline{\Omega}\) 是同胚映射。

\textbf{注记}：Caratheodory 定理是共形映射边界行为的基石。它揭示了区域的边界拓扑性质与共形映射延拓性之间的深刻联系。若 \(\partial\Omega\) 是局部连通的（更一般的条件），\(f\) 仍有连续延拓。若 \(\partial\Omega\) 是可求长曲线，则 \(f'\) 属于 Hardy 空间 \(H^1\)（F. and M. Riesz 定理）。边界行为的进一步研究涉及素数端（prime end）理论------Caratheodory 的又一项重大贡献------以及 \(f\) 的边界值在 \(\partial\Omega\) 上的 Hausdorff 维数等精细几何量。

Caratheodory 的素数端理论为理解共形映射在非 Jordan 边界上的行为提供了完整的框架。素数端是对区域边界点的''共形等价类''的精确定义：即使 \(\partial\Omega\) 不是 Jordan 曲线，素数端的概念使得 \(f\) 在 \(\partial D\) 与 \(\Omega\) 的素数端集合之间建立了双射。这一理论在复动力系统（如 Julia 集的拓扑）中发挥着核心作用------Julia 集通常不是 Jordan 曲线，但其素数端结构决定了填充 Julia 集上 Riemann 映射的边界行为。此外，共形映射的边界正则性还与 \textbf{Lindelof 定理}密切相关：若 \(f\) 沿某条路径存在渐近极限，则 \(f\) 在相应的边界点处具有非切向极限，且该极限与路径的选取无关。

\hrulefill

\subsection{55.6 多复变中的共形映射}\label{ux591aux590dux53d8ux4e2dux7684ux5171ux5f62ux6620ux5c04}

\textbf{注记 55.6.1（高维情形）}：Riemann 映射定理在 \(\mathbb{C}^n\)（\(n \geq 2\)）中不成立------这是单复变与多复变之间最深刻的差异之一。Poincare 在 1907 年发现单位球 \(B_2 = \{(z_1, z_2) \in \mathbb{C}^2 : |z_1|^2 + |z_2|^2 < 1\}\) 与二圆盘 \(D^2 = \{(z_1, z_2) : |z_1| < 1, |z_2| < 1\}\) 之间不存在双全纯映射（尽管两者都是单连通的）。这一发现开启了多复变中域的分类理论------包括 Cartan 的对称域分类、Fefferman 的边界 CR 几何与 Bergman 核渐近展开（Fefferman 扩张定理），以及 Chern-Moser 不变量理论。\(\blacksquare\)

Poincare 的观察揭示了多复变中 Cauchy-Riemann 方程的超定性：\(\mathbb{C}^n\)（\(n \geq 2\)）中的全纯映射必须满足过多的相容性条件，使得两个域之间的双全纯映射的存在性成为一个极其受限的问题。这导致了 \textbf{Chern-Moser 不变量}理论------由陈省身和 Moser 在 1974 年引入，将 CR 几何中的 Cartan 等价方法推广到非 Levi 平坦的 CR 流形。Chern-Moser 不变量是 \(\mathbb{C}^n\) 中强拟凸域的边界 CR 结构在双全纯映射下的完全分类，其计算涉及 CR 流形上的 Cartan 联络和曲率。这一理论在 Fefferman 的 Bergman 核渐近展开中获得了关键的几何解释：Bergman 核在边界上的奇异行为由 Chern-Moser 不变量控制。

\section{第171章：解析延拓}\label{ux7b2c171ux7ae0ux89e3ux6790ux5ef6ux62d3}

解析延拓是复分析中最为深刻的概念之一------它揭示了全纯函数具有''刚性''：一个全纯函数在任意开子集上的值唯一地决定了其在整个连通区域上的值。本章系统介绍解析延拓的基本概念与唯一性定理、沿曲线的解析延拓与单值性定理（Monodromy 定理）、Schwarz 反射原理以及多值函数的 Riemann 面构造。解析延拓理论自然地引出了 Riemann 面的概念，也将复分析与代数拓扑（基本群、覆叠空间）紧密地联系起来，是理解多值函数（如 \(\sqrt{z}\)、\(\log z\)）和模形式的根本框架。

\subsection{56.1 解析延拓的概念}\label{ux89e3ux6790ux5ef6ux62d3ux7684ux6982ux5ff5}

\textbf{定义 56.1.1（解析延拓）}：设 \(\Omega_1 \subseteq \Omega_2\) 为 \(\mathbb{C}\) 中的区域，\(f_1\) 在 \(\Omega_1\) 上全纯，\(f_2\) 在 \(\Omega_2\) 上全纯。若 \(f_2|_{\Omega_1} = f_1\)，则称 \(f_2\) 为 \(f_1\) 从 \(\Omega_1\) 到 \(\Omega_2\) 的\textbf{解析延拓}。

\textbf{定理 56.1.1（解析延拓的唯一性）}：解析延拓若存在，则唯一。

\textbf{证明}：设 \(f_2, g_2\) 均为 \(f_1\) 在 \(\Omega_2\) 上的解析延拓，则 \(h = f_2 - g_2\) 在 \(\Omega_2\) 上全纯，且在 \(\Omega_1\) 上 \(h = 0\)。由全纯函数的唯一性定理（见下文），\(h \equiv 0\) 在 \(\Omega_2\) 上。\(\blacksquare\)

\textbf{定理 56.1.2（唯一性定理）}：设 \(f\) 在区域 \(\Omega\) 上全纯。若存在 \(z_0 \in \Omega\) 和 \(\{z_n\}_{n=1}^{\infty} \subseteq \Omega \setminus \{z_0\}\) 使得 \(z_n \to z_0\) 且 \(f(z_n) = 0\)（\(\forall n\)），则 \(f \equiv 0\) 在 \(\Omega\) 上。

\textbf{证明}：设 \(f\) 在 \(z_0\) 处的 Taylor 展开为 \(f(z) = \sum_{k=0}^{\infty} a_k (z - z_0)^k\)。若 \(f\) 不恒为零，设 \(a_m\) 为第一个非零系数。则 \(f(z) = (z - z_0)^m (a_m + a_{m+1}(z - z_0) + \cdots) = (z - z_0)^m g(z)\)，其中 \(g(z_0) = a_m \neq 0\)。由 \(g\) 的连续性，\(g(z) \neq 0\) 在 \(z_0\) 附近。但 \(f(z_n) = 0\) 且 \(z_n \neq z_0\)，故 \(g(z_n) = 0\)，与 \(g(z_0) \neq 0\) 和连续性矛盾。因此所有 \(a_k = 0\)，\(f \equiv 0\) 在 \(z_0\) 的一个邻域内。由区域的连通性，\(f \equiv 0\) 在整个 \(\Omega\) 上。\(\blacksquare\)

\hrulefill

\subsection{56.2 沿曲线的解析延拓与单值性定理}\label{ux6cbfux66f2ux7ebfux7684ux89e3ux6790ux5ef6ux62d3ux4e0eux5355ux503cux6027ux5b9aux7406}

\textbf{定义 56.2.1（沿曲线的解析延拓）}：设 \(\gamma: [0,1] \to \mathbb{C}\) 为连续曲线，\(\gamma(0) = a\)。\(f\) 在 \(a\) 处全纯的函数元素 \((f_0, D_0)\) 沿 \(\gamma\) 的\textbf{解析延拓}是一列函数元素 \((f_0, D_0), (f_1, D_1), \ldots, (f_n, D_n)\) 和 \(\gamma\) 的分割 \(0 = t_0 < t_1 < \cdots < t_n = 1\)，使得 \(\gamma([t_{k-1}, t_k]) \subseteq D_{k-1}\)，且 \(f_k = f_{k-1}\) 在 \(D_{k-1} \cap D_k\) 上。

\textbf{定理 56.2.1（单值性定理 / Monodromy Theorem）}：设 \(\Omega\) 为单连通区域，\(f\) 在 \(z_0 \in \Omega\) 处的函数元素可沿 \(\Omega\) 中任意从 \(z_0\) 出发的曲线解析延拓。则 \(f\) 可延拓为 \(\Omega\) 上的单值全纯函数。换言之，沿任意两条同伦的曲线的解析延拓得到相同的函数值。

\textbf{证明}：对任意 \(z \in \Omega\)，取从 \(z_0\) 到 \(z\) 的曲线 \(\gamma\)，定义 \(F(z)\) 为 \(f\) 沿 \(\gamma\) 解析延拓的终点值。需证 \(F(z)\) 与 \(\gamma\) 的选取无关。由 \(\Omega\) 的单连通性，任意两条从 \(z_0\) 到 \(z\) 的曲线 \(\gamma_0, \gamma_1\) 同伦。将同伦 \(\gamma_s\) 参数化，在每个 \(\gamma_s\) 上解析延拓。由于延拓的局部唯一性，\(F(z)\) 连续依赖于 \(s\)。但 \(F(z)\) 作为 \(s\) 的函数取值离散（因为函数元素由解析延拓唯一确定），故 \(F(z)\) 在 \(s\) 上为常数，即 \(\gamma_0\) 和 \(\gamma_1\) 给出相同的延拓值。因此 \(F\) 在 \(\Omega\) 上良定且全纯。\(\blacksquare\)

单值性定理是解析延拓理论中最根本的定理之一。它揭示了''局部延拓的存在性''蕴含''全局单值延拓的存在性''这一核心原理，前提是区域单连通。单值性定理在代数函数论中有重要应用：若代数方程 \(P(z, w) = 0\) 定义的代数函数在穿孔圆盘上可解析延拓，则其所有分支在单连通区域上都是单值的。这一原理在 Riemann 面的构造中至关重要------Riemann 面正是通过允许函数''多值''来消除单值性定理对单连通性假设的依赖。

\hrulefill

\subsection{56.3 自然边界与 Painleve 定理}\label{ux81eaux7136ux8fb9ux754cux4e0e-painleve-ux5b9aux7406}

\textbf{定义 56.3.1（自然边界）}：设 \(f\) 在 \(\Omega\) 上全纯。若 \(\Omega\) 的边界 \(\partial \Omega\) 上的每点都是 \(f\) 的奇点（即 \(f\) 不能延拓到该点的任何邻域），则称 \(\partial \Omega\) 为 \(f\) 的\textbf{自然边界}。

自然边界是复分析中最能体现全纯函数''刚性''的现象之一。与实分析中函数可以光滑地延拓形成鲜明对比，复全纯函数在自然边界上遭遇的是不可逾越的障碍------在边界上每点处，函数的奇异性都''堆积''得如此密集，以至于不存在任何穿孔邻域中的全纯延拓。经典例子包括：（1）\(\sum_{n=0}^\infty z^{n!}\) 以单位圆周 \(|z| = 1\) 为自然边界，因为其奇点在单位圆周的每个根处稠密分布；（2）\(\sum_{n=0}^\infty z^{2^n}\)（Weierstrass 间隙级数）同样以单位圆周为自然边界，其奇点集在单位圆周上形成稠密子集；（3）模形式 \(f(z) = \sum_{n=0}^\infty a_n e^{2\pi i n z}\) 以实轴（\(\Im(z) = 0\)）为自然边界。Hadamard 间隙定理（Hadamard lacunary theorem）给出了判别自然边界的充分条件：若幂级数 \(\sum a_n z^{\lambda_n}\) 满足 \(\lambda_{n+1}/\lambda_n > q > 1\)（间隙条件），且收敛半径为 \(R\)，则圆周 \(|z| = R\) 是自然边界。这一结果与 Fabry 间隙定理（Fabry gap theorem）和 Polya-Carlson 定理共同构成了自然边界理论的基石。从几何函数论的角度，自然边界反映了 Riemann 面的''边界''不仅是拓扑的，更是解析的------单叶函数在单位圆盘上的边界行为（如边界旋转、角导数等）与自然边界的形成有密切联系。在更广泛的视角下，自然边界的概念在多复变中得以推广：\(\mathbb{C}^n\)（\(n \geq 2\)）中的全纯域（domain of holomorphy）本质上就是''不存在自然边界外的延拓''的域，其边界至少在某些点处阻碍全纯延拓。

\textbf{定理 56.3.1（Painleve 定理，解析延拓的充分条件）}：设 \(\Omega \subseteq \mathbb{C}\) 为区域，\(E \subseteq \Omega\) 为可求长曲线（或更一般地，一维 Hausdorff 测度为零的紧集）。若 \(f\) 在 \(\Omega \setminus E\) 上全纯且有界，则 \(f\) 可全纯延拓到整个 \(\Omega\)。

\textbf{证明（框架）}：由 Cauchy 积分公式，对 \(\Omega \setminus E\) 中的任意闭曲线 \(\gamma\)，若 \(\gamma\) 不环绕 \(E\)，则 \(\int_\gamma f(z) dz = 0\)。对环绕 \(E\) 的曲线，利用 \(E\) 的一维测度为零，可将积分路径微调避开 \(E\)，仍得积分为零。因此由 Morera 定理，\(f\) 可延拓到 \(\Omega\)。\(\blacksquare\)

Painleve 定理是解析延拓理论中关于''可去奇点''的终极推广。它在复动力系统中有重要应用：若 Julia 集的面积为零（如诸多有理函数的 Julia 集），则 Fatou 分量上的动力学可全纯延拓到 Julia 集上，从而限制了 Julia 集的拓扑复杂性。Painleve 定理的多复变推广涉及 CR 延拓------若 \(f\) 在 \(\mathbb{C}^n\) 中区域边界的一个 CR 子流形上定义，且在区域内部全纯，则 \(f\) 可延拓到整个区域。这一结果由 Bochner 和 Severi 开创，是现代 CR 几何的核心。

\hrulefill

\subsection{56.4 多值函数与 Riemann 面初步}\label{ux591aux503cux51fdux6570ux4e0e-riemann-ux9762ux521dux6b65}

\textbf{定义 56.4.1（解析延拓与多值函数）}：有些函数通过解析延拓会产生多值性。例如 \(f(z) = \sqrt{z}\) 在 \(z = 0\) 附近：从某点出发沿包含 \(0\) 的闭曲线绕一圈，函数值变为原来的 \(-1\) 倍。

\textbf{定义 56.4.2（Riemann 面）}：多值函数的自然定义域是 \textbf{Riemann 面}------一个二维复流形，使得''多值函数''在其上成为单值全纯函数。

例如 \(\sqrt{z}\) 的 Riemann 面是由两叶复平面在 \(z = 0\) 处通过分支切割（branch cut）连接而成的。

对数函数 \(\log z\) 的 Riemann 面是无穷叶螺旋面，其单值化群同构于 \(\mathbb{Z}\)（由绕 \(z = 0\) 一圈增加 \(2\pi i\) 生成）。更一般地，任意代数函数（由代数方程 \(P(z, w) = \sum_{k=0}^n a_k(z) w^k = 0\) 定义）的 Riemann 面是紧 Riemann 面（可能带有穿孔点），其亏格可由 Riemann-Hurwitz 公式计算。这一构造是代数曲线理论的历史起点，Riemann 正是通过这一构造将 Abel 的椭圆积分反演理论与 Cauchy 的全纯函数论统一起来。Riemann 面的基本群与覆盖空间理论为单值性定理提供了拓扑学基础：解析延拓的多值性源于基本群的非平凡性，而 Riemann 面正是万有覆盖空间，其上的提升函数自动成为单值全纯函数。

\textbf{注记 56.4.1}：Riemann 面的系统理论将在卷十二（微分几何）中作为一维复流形展开，并将在卷十六（代数几何）中深化为代数曲线理论。\(\blacksquare\)

\hrulefill

\hrulefill

\hrulefill

\hrulefill

\hrulefill

\newpage

\chapter{06-分析/03-泛函分析/01-Banach空间.md}\label{ux5206ux679003-ux6cdbux51fdux5206ux679001-banachux7a7aux95f4.md}

\section{第172章：Banach 空间}\label{ux7b2c172ux7ae0banach-ux7a7aux95f4}

Banach 空间是完备的赋范向量空间，是泛函分析最基本的框架。Hahn-Banach、开映射、闭图像和一致有界原理被称为 Banach 空间的四大基本定理。

\subsection{72.1 赋范空间与 Banach 空间}\label{ux8d4bux8303ux7a7aux95f4ux4e0e-banach-ux7a7aux95f4}

\textbf{定义 72.1}（赋范空间）：设 \(X\) 是数域 \(\mathbb{K}\)（\(\mathbb{R}\) 或 \(\mathbb{C}\)）上的向量空间。映射 \(\|\cdot\|: X \to [0, \infty)\) 称为 \(X\) 上的一个\textbf{范数}，如果满足：

\begin{enumerate}
\def\labelenumi{\arabic{enumi}.}
\tightlist
\item
  \(\|x\| = 0 \Leftrightarrow x = 0\)（正定性）
\item
  \(\|\alpha x\| = |\alpha| \|x\|\)（齐次性）
\item
  \(\|x + y\| \leq \|x\| + \|y\|\)（三角不等式）
\end{enumerate}

\((X, \|\cdot\|)\) 称为\textbf{赋范向量空间}。范数诱导度量 \(d(x, y) = \|x - y\|\) 和拓扑。

\textbf{定义 72.2}（Banach 空间）：完备的赋范向量空间称为 \textbf{Banach 空间}。即每个 Cauchy 序列都收敛。

\textbf{定义 72.3}（等价范数）：\(X\) 上的两个范数 \(\|\cdot\|_1\) 和 \(\|\cdot\|_2\) 称为\textbf{等价的}，如果存在常数 \(c, C > 0\) 使得 \(c\|x\|_1 \leq \|x\|_2 \leq C\|x\|_1\) 对所有 \(x \in X\) 成立。

\textbf{定理 72.1}：有限维赋范空间上的所有范数都等价，且有限维赋范空间总是 Banach 空间。

\textbf{证明}：设 \(X\) 是 \(n\) 维赋范空间，取基 \(\{e_1, \ldots, e_n\}\)。定义欧氏范数 \(\|\sum a_i e_i\|_2 = (\sum |a_i|^2)^{1/2}\)。对任意范数 \(\|\cdot\|\)，由三角不等式和 Cauchy-Schwarz 得 \(\|x\| \leq C\|x\|_2\)。函数 \(x \mapsto \|x\|\) 在单位球面 \(S = \{x : \|x\|_2 = 1\}\) 上连续且严格正，由紧致性得 \(c = \inf_S \|x\| > 0\)，故 \(\|x\| \geq c\|x\|_2\)。\(\blacksquare\)

\subsection{72.2 有界线性算子}\label{ux6709ux754cux7ebfux6027ux7b97ux5b50}

\textbf{定义 72.4}（有界线性算子）：设 \(X, Y\) 是赋范空间。线性算子 \(T: X \to Y\) 称为\textbf{有界的}，如果存在 \(M \geq 0\) 使得 \(\|Tx\| \leq M\|x\|\) 对所有 \(x \in X\) 成立。\(T\) 的\textbf{算子范数}定义为

\[\|T\| = \sup_{\|x\| \leq 1} \|Tx\| = \sup_{\|x\| = 1} \|Tx\|\]

\textbf{命题 72.2}：线性算子 \(T: X \to Y\) 有界当且仅当 \(T\) 连续，当且仅当 \(T\) 在 \(0\) 处连续。

\textbf{证明}：若有界，则 \(\|Tx - Ty\| = \|T(x-y)\| \leq \|T\| \|x-y\|\)，故 Lipschitz 连续。若在 \(0\) 处连续，则存在 \(\delta > 0\) 使 \(\|x\| < \delta \Rightarrow \|Tx\| < 1\)。对任意 \(x\)，\(\|T(\delta x/(2\|x\|))\| < 1\)，故 \(\|Tx\| < (2/\delta)\|x\|\)。\(\blacksquare\)

\textbf{定义 72.5}（算子空间）：\(X\) 到 \(Y\) 的所有有界线性算子构成向量空间，记作 \(\mathcal{B}(X, Y)\)。配以算子范数，\(\mathcal{B}(X, Y)\) 是赋范空间。若 \(Y\) 是 Banach 空间，则 \(\mathcal{B}(X, Y)\) 也是 Banach 空间。特别地，\(X^* = \mathcal{B}(X, \mathbb{K})\) 称为 \(X\) 的\textbf{对偶空间}（或共轭空间）。

\textbf{定理 72.3}（Banach-Steinhaus 一致有界原理）：设 \(X\) 是 Banach 空间，\(Y\) 是赋范空间，\(\{T_\alpha\}_{\alpha \in A} \subseteq \mathcal{B}(X, Y)\)。若对每个 \(x \in X\)，\(\sup_{\alpha} \|T_\alpha x\| < \infty\)，则 \(\sup_{\alpha} \|T_\alpha\| < \infty\)。

\textbf{证明}：使用 Baire 纲定理。对每个 \(n \in \mathbb{N}\)，定义 \(F_n = \{x \in X : \sup_\alpha \|T_\alpha x\| \leq n\}\)。\(F_n\) 是闭集，且 \(X = \bigcup_{n} F_n\)。由 Baire 纲定理，存在 \(n_0\) 使 \(F_{n_0}\) 有内点。设 \(B(x_0, r) \subseteq F_{n_0}\)。对任意 \(\|x\| \leq r\)，\(x_0 + x \in F_{n_0}\)，故 \(\|T_\alpha x\| \leq \|T_\alpha(x_0 + x)\| + \|T_\alpha x_0\| \leq 2n_0\)。因此 \(\|T_\alpha\| \leq 2n_0/r\) 对所有 \(\alpha\) 成立。\(\blacksquare\)

\subsection{72.3 Hahn-Banach 定理}\label{hahn-banach-ux5b9aux7406}

\textbf{定理 72.4}（Hahn-Banach 延拓定理，实情形）：设 \(X\) 是实向量空间，\(p: X \to \mathbb{R}\) 满足 \(p(x+y) \leq p(x) + p(y)\) 和 \(p(tx) = tp(x)\) 对 \(t \geq 0\)（次线性泛函）。设 \(M \subseteq X\) 是子空间，\(f: M \to \mathbb{R}\) 是线性泛函满足 \(f(x) \leq p(x)\) 对所有 \(x \in M\)。则存在 \(f\) 的线性延拓 \(F: X \to \mathbb{R}\) 使得 \(F(x) \leq p(x)\) 对所有 \(x \in X\) 成立。

\textbf{证明}：使用 Zorn 引理。考虑所有保持 \(f\) 且满足 \(F(x) \leq p(x)\) 的延拓构成的偏序集（按包含关系）。由 Zorn 引理存在极大延拓 \(F\)，定义域为 \(M'\)。若 \(M' \neq X\)，取 \(x_0 \notin M'\)。需定义 \(F(x_0)\) 使得对任意 \(y, z \in M'\)：

\[F(y) + F(x_0) \leq p(y + x_0), \quad F(z) - F(x_0) \leq p(z - x_0)\]

这等价于 \(F(z) - p(z - x_0) \leq F(x_0) \leq p(y + x_0) - F(y)\)。由 \(f(y) + f(z) = f(y+z) \leq p(y+z) \leq p(y+x_0) + p(z-x_0)\)，知上确界不大于下确界，故可取 \(F(x_0)\) 为中间值。在 \(M' \oplus \mathbb{R}x_0\) 上延拓，与极大性矛盾。\(\blacksquare\)

\textbf{定理 72.5}（Hahn-Banach 定理，赋范空间版本）：设 \(X\) 是赋范空间，\(M \subseteq X\) 是子空间，\(f \in M^*\)。则存在 \(F \in X^*\) 使得 \(F|_M = f\) 且 \(\|F\| = \|f\|\)。

\textbf{证明}：取 \(p(x) = \|f\| \|x\|\)（次线性），则 \(|f(x)| \leq p(x)\) 在 \(M\) 上。由 Hahn-Banach 定理存在 \(F\) 满足 \(F(x) \leq p(x)\)。由 \(F(-x) = -F(x) \leq p(-x) = p(x)\) 得 \(|F(x)| \leq p(x) = \|f\| \|x\|\)，故 \(\|F\| \leq \|f\|\)。反向不等式显然。\(\blacksquare\)

\textbf{推论 72.6}（Hahn-Banach 推论）： 1. 对任意 \(x \in X, x \neq 0\)，存在 \(f \in X^*\) 使得 \(\|f\| = 1\) 且 \(f(x) = \|x\|\)。 2. \(X^*\) 分离 \(X\) 中的点：若 \(x \neq y\)，则存在 \(f \in X^*\) 使 \(f(x) \neq f(y)\)。 3. \(\|x\| = \sup_{\|f\| \leq 1} |f(x)|\)。

\subsection{72.4 开映射定理与闭图像定理}\label{ux5f00ux6620ux5c04ux5b9aux7406ux4e0eux95edux56feux50cfux5b9aux7406}

\textbf{定理 72.7}（开映射定理）：设 \(X, Y\) 是 Banach 空间，\(T \in \mathcal{B}(X, Y)\) 是满射。则 \(T\) 是开映射（即将开集映为开集）。特别地，若 \(T\) 是连续双射，则 \(T^{-1}\) 连续，即 \(T\) 是同胚（从而 \(T\) 是 Banach 空间同构）。

\textbf{证明}：证 \(T(B_X(0, 1))\) 包含 \(B_Y(0, r)\) 对某个 \(r > 0\)。由满射性和 Baire 纲定理，存在 \(n\) 使 \(\overline{T(B_X(0, n))}\) 有内点。则 \(\overline{T(B_X(0, 1))}\) 包含某球 \(B_Y(y_0, \delta)\)。由对称性得 \(B_Y(0, \delta) \subseteq \overline{T(B_X(0, 2))}\)。再由完备性迭代可得 \(B_Y(0, \delta/2) \subseteq T(B_X(0, 1))\)。\(\blacksquare\)

\textbf{定理 72.8}（闭图像定理）：设 \(X, Y\) 是 Banach 空间，\(T: X \to Y\) 是线性算子。则 \(T\) 有界当且仅当 \(T\) 的图像 \(\Gamma(T) = \{(x, Tx) : x \in X\}\) 是 \(X \times Y\)（配以范数 \(\|(x, y)\| = \|x\| + \|y\|\)）中的闭集。

\textbf{证明}：(\(\Rightarrow\)) 若 \(T\) 有界，设 \((x_n, Tx_n) \to (x, y)\)，则 \(x_n \to x\)，由连续性 \(Tx_n \to Tx\)，故 \(y = Tx\)，\(\Gamma(T)\) 闭。

(\(\Leftarrow\)) 若 \(\Gamma(T)\) 闭，则 \(\Gamma(T)\) 是 Banach 空间 \(X \times Y\) 的闭子空间，从而是 Banach 空间。投影 \(\pi_1: \Gamma(T) \to X\)，\(\pi_1(x, Tx) = x\) 是连续双射。由开映射定理，\(\pi_1^{-1}\) 连续，故 \(\|x\| + \|Tx\| = \|(x, Tx)\| \leq C\|x\|\)，从而 \(\|Tx\| \leq C\|x\|\)，\(T\) 有界。\(\blacksquare\)

\subsection{72.5 算子空间的完备性}\label{ux7b97ux5b50ux7a7aux95f4ux7684ux5b8cux5907ux6027}

算子空间 \(\mathcal{B}(X, Y)\) 的完备性是泛函分析中最为基本而又极其重要的结果之一。它保证了从任意赋范空间 \(X\) 到 Banach 空间 \(Y\) 的有界线性算子全体在算子范数下构成 Banach 空间。这一事实的证明体现了''逐点极限-\textgreater 线性结构-\textgreater 范数收敛''的三步推理模式，是泛函分析中许多完备性结果的雏形。从拓扑角度看，\(\mathcal{B}(X, Y)\) 的完备性等价于有界线性算子空间的''遗传完备性''：只要目标空间 \(Y\) 是完备的，算子空间就自动完备，而不需要对 \(X\) 施加完备性要求。这一现象与一致收敛空间的性质密切相关------在紧集上一致收敛的连续函数空间总是完备的，即使定义域是局部紧空间而非完备空间。算子空间完备性的一个重要应用是构造 Banach 空间的极限和归纳极限：利用 \(\mathcal{B}(X, Y)\) 的完备性，可以证明 Banach 空间范畴中存在任意小的极限和归纳极限。此外，\(\mathcal{B}(X) = \mathcal{B}(X, X)\) 不仅是 Banach 空间，还是 Banach 代数（以算子复合为乘法，满足 \(\|ST\| \leq \|S\| \|T\|\)），这使得 Banach 代数理论（谱理论、Gelfand 变换等）可以自然地应用于算子理论。在算子空间的完备性基础上，可以进一步研究算子的强收敛与弱收敛：Banach-Steinhaus 一致有界原理（定理 72.3）提供了一种从逐点有界性推断一致有界性的方法，而这本质上依赖于 \(\mathcal{B}(X, Y)\) 中 Cauchy 序列的完备性。

\textbf{定理 72.9}（\(\mathcal{B}(X, Y)\) 的完备性）：设 \(X\) 为赋范空间，\(Y\) 为 Banach 空间。则 \(\mathcal{B}(X, Y)\) 配以算子范数为 Banach 空间。

\textbf{证明}：设 \(\{T_n\}\) 为 \(\mathcal{B}(X, Y)\) 中的 Cauchy 列。对每个 \(x \in X\)，\(\|T_n x - T_m x\| \leq \|T_n - T_m\| \|x\|\)，故 \(\{T_n x\}\) 在 \(Y\) 中为 Cauchy 列。由 \(Y\) 的完备性，定义 \(T x = \lim_{n} T_n x\)。\(T\) 的线性性由极限的线性性得。由 \(\{ \|T_n\| \}\) 的有界性（Cauchy 列有界），存在 \(M\) 使 \(\|T_n\| \leq M\)，故 \(\|T x\| = \lim \|T_n x\| \leq M \|x\|\)，\(T\) 有界。还需证 \(\|T_n - T\| \to 0\)：给定 \(\varepsilon > 0\)，取 \(N\) 使 \(m, n \geq N\) 时 \(\|T_n - T_m\| < \varepsilon\)。对 \(\|x\| \leq 1\)，\(\|(T_n - T)x\| = \lim_m \|(T_n - T_m)x\| \leq \varepsilon\)，故 \(n \geq N\) 时 \(\|T_n - T\| \leq \varepsilon\)。\(\blacksquare\)

\subsection{72.6 闭值域定理}\label{ux95edux503cux57dfux5b9aux7406}

\textbf{定理 72.10}（闭值域定理）：设 \(X, Y\) 为 Banach 空间，\(T \in \mathcal{B}(X, Y)\)。则下列条件等价：

\begin{enumerate}
\def\labelenumi{\arabic{enumi}.}
\tightlist
\item
  \(\operatorname{ran}(T)\) 在 \(Y\) 中闭；
\item
  \(\operatorname{ran}(T^*) \subseteq X^*\) 在 \(X^*\) 中闭；
\item
  \(\operatorname{ran}(T) = (\ker T^*)^\perp\)；
\item
  \(\operatorname{ran}(T^*) = (\ker T)^\perp\)。
\end{enumerate}

其中 \(T^* : Y^* \to X^*\) 为 \(T\) 的共轭算子，\(M^\perp = \{f \in X^* : f|_M = 0\}\)。

\textbf{证明框架}： \((1) \Leftrightarrow (3)\)：由 Hahn-Banach，闭子空间的零化子刻画保证了 \(\operatorname{ran}(T) = (\ker T^*)^\perp\) 当且仅当 \(\operatorname{ran}(T)\) 闭。\((3) \Rightarrow (4)\)：用极化恒等式和对偶论证。\((4) \Rightarrow (2)\)：\((\ker T)^\perp\) 总是弱* 闭的。完整证明需引入对偶配对和极集理论。\(\blacksquare\)

\textbf{推论 72.11}：\(T\) 为满射当且仅当 \(T^*\) 有有界逆（在 \(\operatorname{ran}(T^*)\) 上），即 \(\|f\| \leq C \|T^* f\|\) 对 \(f \in Y^*\) 成立。

\subsection{72.7 商空间与商范数}\label{ux5546ux7a7aux95f4ux4e0eux5546ux8303ux6570}

\textbf{定义 72.6}（商空间）：设 \(X\) 为赋范空间，\(M \subseteq X\) 为闭子空间。\textbf{商空间} \(X / M\) 由等价类 \(\dot{x} = x + M = \{x + m : m \in M\}\) 构成，配以自然向量空间结构。\textbf{商范数}定义为

\[\|\dot{x}\|_{X/M} = \inf_{m \in M} \|x + m\|_X = \operatorname{dist}(x, M)\]

\textbf{定理 72.12}（商范数的性质）： 1. \(\|\cdot\|_{X/M}\) 是范数当且仅当 \(M\) 闭； 2. \textbf{自然商映射} \(\pi : X \to X/M\)，\(\pi(x) = \dot{x}\) 是范数为 \(1\) 的连续线性满射，且 \(\pi(B_X^\circ) = B_{X/M}^\circ\)（开单位球映为开单位球），即 \(\pi\) 是开映射； 3. 若 \(X\) 为 Banach 空间且 \(M\) 闭，则 \(X/M\) 为 Banach 空间。

\textbf{证明}：（1）\(\|\dot{x}\| = 0 \Leftrightarrow \operatorname{dist}(x, M) = 0 \Leftrightarrow x \in \overline{M} = M\)，故正定性等价于 \(M\) 闭。（2）由定义 \(\|\pi(x)\| \leq \|x\|\)，故 \(\|\pi\| \leq 1\)。对 \(\|\dot{x}\| < 1\)，存在 \(m \in M\) 使 \(\|x+m\| < 1\)，故 \(\pi(x+m) = \dot{x}\)，得 \(\pi(B_X^\circ) = B_{X/M}^\circ\)。（3）用级数判别法------绝对收敛级数收敛等价于完备性，商范数下级数的收敛性由原空间传递。\(\blacksquare\)

\textbf{定理 72.13}（商空间的泛性质与第一同构定理）：设 \(T \in \mathcal{B}(X, Y)\) 满足 \(M \subseteq \ker T\)。则存在唯一的 \(\widetilde{T} \in \mathcal{B}(X/M, Y)\) 使得 \(T = \widetilde{T} \circ \pi\)，且 \(\|\widetilde{T}\| = \|T\|\)。特别地，

\[X / \ker T \cong \operatorname{ran}(T) \quad (\text{等距同构，若 } \operatorname{ran}(T) \text{ 闭})\]

\textbf{证明}：定义 \(\widetilde{T}(\dot{x}) = T(x)\)，其中 \(\dot{x} = x + M\)。由 \(M \subseteq \ker T\)，此定义与代表元选取无关，即若 \(x - x' \in M\)，则 \(T(x - x') = 0\)，故 \(T(x) = T(x')\)。线性性显然。范数估计：\(\|\widetilde{T}(\dot{x})\| = \|T(x)\| = \|T(x+m)\| \leq \|T\| \cdot \|x+m\|\) 对所有 \(m \in M\) 成立，取下确界得 \(\|\widetilde{T}(\dot{x})\| \leq \|T\| \cdot \|\dot{x}\|\)，故 \(\|\widetilde{T}\| \leq \|T\|\)。反向：\(\|T(x)\| = \|\widetilde{T}(\dot{x})\| \leq \|\widetilde{T}\| \cdot \|\dot{x}\| \leq \|\widetilde{T}\| \cdot \|x\|\)，故 \(\|T\| \leq \|\widetilde{T}\|\)。因此 \(\|\widetilde{T}\| = \|T\|\)。取 \(M = \ker T\)，则 \(\widetilde{T}\) 为单射，且 \(\widetilde{T}\) 的像为 \(\operatorname{ran}(T)\)。若 \(\operatorname{ran}(T)\) 闭，则 \(\widetilde{T}\) 为 \(X/\ker T\) 到 \(\operatorname{ran}(T)\) 的等距同构。\(\blacksquare\)

\subsection{72.8 有限维 Banach 空间的补充性质}\label{ux6709ux9650ux7ef4-banach-ux7a7aux95f4ux7684ux8865ux5145ux6027ux8d28}

\textbf{定理 72.14}（有限维赋范空间的完备刻画）： 1. 有限维赋范空间上所有范数都等价； 2. 有限维赋范空间总是 Banach 空间； 3. 有限维空间中有界闭集为紧集（Heine-Borel 性质）； 4. 赋范空间 \(X\) 为有限维当且仅当 \(B_X\) 为紧集（Riesz 引理）。

\textbf{证明}：（4）若 \(\dim X = \infty\)，Riesz 引理构造 \(\{x_n\} \subseteq B_X\) 满足 \(\|x_n - x_m\| \geq 1/2\)（\(n \neq m\)），无收敛子列，故 \(B_X\) 不紧。\(\blacksquare\)

\subsection{72.9 Schauder 基}\label{schauder-ux57fa}

\textbf{定义 72.7}（Schauder 基）：Banach 空间 \(X\) 中的序列 \(\{e_n\}_{n=1}^\infty\) 称为 \textbf{Schauder 基}，如果对每个 \(x \in X\) 存在唯一的标量序列 \(\{a_n\}_{n=1}^\infty \subseteq \mathbb{K}\) 使得

\[x = \sum_{n=1}^\infty a_n e_n \quad (\text{级数依范数收敛})\]

即 \(\lim_{N \to \infty} \|x - \sum_{n=1}^N a_n e_n\| = 0\)。系数 \(a_n = e_n^*(x)\) 由连续线性泛函 \(e_n^* \in X^*\) 给出（坐标泛函）。

\textbf{命题 72.15}：若 \(X\) 有 Schauder 基，则 \(X\) 可分。

\textbf{证明}：所有基向量的有理系数有限线性组合构成可数稠密集。\(\blacksquare\)

\textbf{定理 72.16}（Schauder 基的特征）：序列 \(\{e_n\}\) 为 Schauder 基当且仅当： 1. \(\overline{\operatorname{span}}\{e_n\} = X\)（完备性）； 2. 存在 \(K \geq 1\) 使得对所有 \(m < n\) 和标量 \(a_k\)，\(\|\sum_{k=1}^m a_k e_k\| \leq K \|\sum_{k=1}^n a_k e_k\|\)（基常数条件）。

此时偏和算子 \(S_n(\sum a_k e_k) = \sum_{k=1}^n a_k e_k\) 一致有界：\(\sup_n \|S_n\| < \infty\)。当 \(\|S_n\| \equiv 1\)（基常数为 \(1\)）时称 \(\{e_n\}\) 为\textbf{单调基}。

\textbf{证明}：若 \(\{e_n\}\) 为 Schauder 基，则每个 \(x \in X\) 有唯一展开 \(x = \sum a_k e_k\)。定义 \(S_n(x) = \sum_{k=1}^n a_k e_k\)。由 Banach-Steinhaus 定理，\(\sup_n \|S_n\| = K < \infty\)。对 \(m < n\)，\(\|S_m x\| = \|S_m(S_n x)\| \leq K \|S_n x\|\)，即基常数条件（2）。完备性（1）由定义保证。反之，若（1）和（2）成立，定义 \(S_n\) 在 \(\operatorname{span}\{e_k\}\) 上为 \(S_n(\sum_{k=1}^N a_k e_k) = \sum_{k=1}^{\min(n,N)} a_k e_k\)。由（2），\(S_n\) 有界，常数 \(\sup_n \|S_n\| \leq K\)。由（1），\(S_n\) 可延拓到 \(X\) 上，且 \(\lim S_n x = x\) 对所有 \(x\) 成立，故 \(\{e_n\}\) 为 Schauder 基。\(\blacksquare\)

\textbf{例}：\(\ell^p\)（\(1 \leq p < \infty\)）中的标准单位向量序列 \(\{e_n\}\)（\(e_n = (\delta_{nk})_k\)）是 Schauder 基。\(C[0, 1]\) 中的 Schauder 基（三角形函数系）由 J. Schauder 于 1927 年构造。\(L^p[0, 1]\)（\(1 \leq p < \infty\)）中的 Haar 系为无条件基。

\textbf{定理 72.17}（Per Enflo，1973）：存在无 Schauder 基的可分 Banach 空间。此反例回答了 Banach 的经典问题（基问题）。

\textbf{证明}（Enflo 构造概要）：Enflo 构造了一个可分的自反 Banach 空间，它不具有逼近性质（approximation property）。具体地，考虑 \(\ell^2\) 上的一个特定范数，该范数由无穷多个非交换的有限维扰动定义。Grothendieck 曾证明：若一个 Banach 空间有 Schauder 基，则它具有有界逼近性质（BAP）。Enflo 的构造产生了一个不具有逼近性质的 Banach 空间，从而否定了基问题。构造的核心是使用「Enflo 负解」中的组合论证：在 \(\ell^2\) 上定义一个新的范数 \(\|\cdot\|_E\)，使得由有限秩算子逼近恒等算子的均匀有界性失效。通过精心选择有限维子空间的递增序列 \(\{E_n\}\) 和投影 \(\{P_n\}\)（其中 \(P_n\) 是到 \(E_n\) 的投影），并利用「扩张性质」（expansion property）确保 \(\|P_n\| \to \infty\)，从而破坏了逼近性质。\(\blacksquare\)

\subsection{72.10 Banach 空间中的级数与无条件收敛}\label{banach-ux7a7aux95f4ux4e2dux7684ux7ea7ux6570ux4e0eux65e0ux6761ux4ef6ux6536ux655b}

\textbf{定义 72.8}（无条件收敛）：Banach 空间中的级数 \(\sum x_n\) 称为\textbf{无条件收敛}，如果对任意置换 \(\sigma : \mathbb{N} \to \mathbb{N}\)，\(\sum x_{\sigma(n)}\) 收敛到相同的和。等价地，\(\sum \varepsilon_n x_n\) 对任意符号 \(\varepsilon_n = \pm 1\) 收敛。

\textbf{定理 72.18}（Dvoretzky-Rogers 定理）：在任意无穷维 Banach 空间中，存在无条件收敛但不绝对收敛的级数（即 \(\sum \|x_n\| = \infty\) 而 \(\sum x_n\) 无条件收敛）。这刻画了有限维与无穷维的根本差异。

\textbf{证明}：设 \(X\) 为无穷维 Banach 空间。由 Dvoretzky 定理，对任意 \(n\)，存在 \(X\) 的 \(n\) 维子空间 \(E_n\) 使得 \(E_n\) 的 Banach-Mazur 距离 \(d(E_n, \ell_2^n) \leq 1 + \varepsilon\)。选取在 \(E_n\) 中近似正交的向量序列 \(\{y_n\}\)：即存在 \(c > 0\) 使对所有标量 \(a_i\)，\(\|\sum a_i y_i\| \geq c (\sum |a_i|^2)^{1/2}\)。设 \(x_n = y_n / \sqrt{n}\)，则 \(\sum \|x_n\| = \sum \|y_n\|/\sqrt{n} = \infty\)（因 \(\|y_n\| \geq c\) 对所有 \(n\)）。但由正交估计，\(\sum \varepsilon_n x_n\) 在 \(X\) 中收敛：\(\| \sum_{n=m}^N \varepsilon_n x_n \| \leq M \sum_{n=m}^N 1/n\)，方差有限保证无条件收敛。因此 \(\sum x_n\) 无条件收敛但不绝对收敛。\(\blacksquare\)

\hrulefill

\hrulefill

\section{第173章：对偶空间与弱拓扑}\label{ux7b2c173ux7ae0ux5bf9ux5076ux7a7aux95f4ux4e0eux5f31ux62d3ux6251}

本章研究 Banach 空间的对偶空间、二次对偶和弱拓扑，这些是研究无穷维空间紧致性的关键工具。

\subsection{74.1 对偶空间与二次对偶}\label{ux5bf9ux5076ux7a7aux95f4ux4e0eux4e8cux6b21ux5bf9ux5076}

\textbf{定义 74.1}（对偶空间）：设 \(X\) 为 \(\mathbb{K}\)（\(\mathbb{R}\) 或 \(\mathbb{C}\)）上的赋范空间。\(X\) 的\textbf{对偶空间}（或共轭空间）定义为

\[X^* = \mathcal{B}(X, \mathbb{K})\]

即 \(X\) 上所有有界线性泛函的集合，配以算子范数 \(\|f\|_{X^*} = \sup_{\|x\| \leq 1} |f(x)|\)。\(X^*\) 总是 Banach 空间（即使 \(X\) 不完备）。

\textbf{定义 74.2}（二次对偶与自然嵌入）：\(X^{**} = (X^*)^*\) 是 \(X\) 的\textbf{二次对偶}。\textbf{自然嵌入} \(J : X \to X^{**}\) 定义为

\[J(x)(f) = f(x), \quad \forall f \in X^*\]

由 Hahn-Banach 推论，\(\|J(x)\|_{X^{**}} = \sup_{\|f\| \leq 1} |f(x)| = \|x\|_X\)，故 \(J\) 是等距嵌入。

\textbf{定义 74.3}（自反空间）：Banach 空间 \(X\) 称为\textbf{自反的}，如果自然嵌入 \(J : X \to X^{**}\) 是满射（从而是等距同构）。记作 \(X \cong X^{**}\)。

\textbf{命题 74.1}：自反空间的闭子空间是自反的。\(X\) 自反当且仅当 \(X^*\) 自反。

\subsection{74.2 弱拓扑与弱* 拓扑}\label{ux5f31ux62d3ux6251ux4e0eux5f31-ux62d3ux6251}

在无穷维 Banach 空间中，有界闭集不一定紧致（例如单位球面不是紧致的）。弱拓扑和弱* 拓扑提供了更粗的拓扑，使得某些集合变得紧致。

\textbf{定义 74.4}（弱拓扑）：\(X\) 上的\textbf{弱拓扑} \(\sigma(X, X^*)\) 是使所有 \(f \in X^*\) 连续的最弱拓扑。其子基为 \(\{f^{-1}(U) : f \in X^*, U \subseteq \mathbb{K} \text{ 开}\}\)。

网 \((x_\alpha)\) 弱收敛到 \(x\)（记作 \(x_\alpha \rightharpoonup x\)）当且仅当 \(f(x_\alpha) \to f(x)\) 对所有 \(f \in X^*\) 成立。

\textbf{定义 74.5}（弱* 拓扑）：\(X^*\) 上的\textbf{弱* 拓扑} \(\sigma(X^*, X)\) 是使所有 \(J(x)\)（\(x \in X\)）连续的最弱拓扑。即 \(f_\alpha \xrightarrow{w^*} f\) 当且仅当 \(f_\alpha(x) \to f(x)\) 对所有 \(x \in X\)。

\textbf{命题 74.2}：弱拓扑和弱* 拓扑是 Hausdorff 拓扑。若 \(X\) 自反，则 \(\sigma(X^*, X) = \sigma(X^*, X^{**})\)，即弱* 拓扑与弱拓扑一致。

\textbf{定理 74.3}（Banach-Alaoglu 定理）：\(X^*\) 中的闭单位球 \(B_{X^*} = \{f \in X^* : \|f\| \leq 1\}\) 在弱* 拓扑下是紧致的。

\textbf{证明}：考虑映射 \(\Phi : B_{X^*} \to \prod_{x \in X} [-\|x\|, \|x\|]\) 定义为 \(\Phi(f) = (f(x))_{x \in X}\)。\(\Phi\) 是单射且是到其像的同胚（\(B_{X^*}\) 赋予弱* 拓扑，乘积空间赋予积拓扑）。由 Tychonoff 定理，乘积空间紧致。仅需验证 \(\Phi(B_{X^*})\) 在积拓扑中为闭集：若网 \(\Phi(f_\alpha)\) 收敛到 \((a_x)_{x \in X}\)，则线性泛函 \(f(x) := a_x\) 满足 \(|f(x)| \leq \|x\|\)（故 \(\|f\| \leq 1\)），且 \(f_\alpha \xrightarrow{w^*} f\)，故 \((a_x) = \Phi(f) \in \Phi(B_{X^*})\)。紧空间的闭子集紧致。\(\blacksquare\)

\textbf{定理 74.4}（Goldstine 定理）：\(J(B_X)\) 在 \(B_{X^{**}}\) 中弱* 稠密。即单位球的像在二次对偶的单位球中稠密。

\textbf{证明}：设 \(F \in B_{X^{**}}\)。对任意弱* 邻域 \(W = \{\Phi \in X^{**} : |\Phi(f_i) - F(f_i)| < \varepsilon, i=1,\ldots,n\}\)，需找到 \(x \in B_X\) 使 \(J(x) \in W\)，即 \(|f_i(x) - F(f_i)| < \varepsilon\)。由 Helly 定理（Hahn-Banach 的推论），这在 \(F \in B_{X^{**}}\) 的条件下恒可能。\(\blacksquare\)

\textbf{推论 74.5}：\(X\) 自反当且仅当 \(B_X\) 在弱拓扑下紧致。

\textbf{证明}：若 \(X\) 自反，\(J(B_X) = B_{X^{**}}\)，由 Banach-Alaoglu 定理在弱* 拓扑下紧致，而 \(X\) 自反时弱* 拓扑等于弱拓扑。反之，若 \(B_X\) 弱紧致，则 \(J(B_X)\) 弱* 紧致从而弱* 闭，由 Goldstine 定理得 \(J(B_X) = B_{X^{**}}\)。\(\blacksquare\)

\subsection{74.3 凸集的分离与 Krein-Milman 定理}\label{ux51f8ux96c6ux7684ux5206ux79bbux4e0e-krein-milman-ux5b9aux7406}

\textbf{定理 74.6}（凸集分离定理 / Hahn-Banach 定理的几何形式）：设 \(A, B\) 是 Banach 空间的非空不交凸集，\(A\) 是开集。则存在 \(f \in X^*\) 和 \(\alpha \in \mathbb{R}\) 使得 \(\Re f(a) < \alpha \leq \Re f(b)\) 对所有 \(a \in A, b \in B\)。

\textbf{证明}：固定 \(a_0 \in A, b_0 \in B\)。令 \(C = A - B + (b_0 - a_0) = \{a - b + b_0 - a_0 : a \in A, b \in B\}\)。则 \(C\) 是开凸集且 \(0 \in C\)。定义 \(C\) 的 Minkowski 泛函 \(p(x) = \inf\{t > 0 : x \in tC\}\)。\(p\) 是次线性泛函。由于 \(a_0 - b_0 \notin C\)（否则 \(A \cap B \neq \emptyset\)），\(p(a_0 - b_0) \geq 1\)。在一维子空间 \(\mathbb{R} \cdot (a_0 - b_0)\) 上定义线性泛函 \(\ell(t(a_0 - b_0)) = t \cdot p(a_0 - b_0)\)，则 \(\ell \leq p\)。由 Hahn-Banach 延拓定理，存在 \(f \in X^*\) 满足 \(f(x) \leq p(x)\) 对所有 \(x\)。取 \(\alpha = \inf_{b \in B} \Re f(b)\)，则对所有 \(a \in A, b \in B\)，\(\Re f(a) < \alpha \leq \Re f(b)\)。\(\blacksquare\)

\textbf{定理 74.7}（Mazur 定理）：设 \(x_n \rightharpoonup x\)（弱收敛）。则存在凸组合序列 \(\sum_{k=1}^{N_n} \lambda_k^{(n)} x_k\)（\(\lambda_k^{(n)} \geq 0, \sum \lambda_k^{(n)} = 1\)）强收敛到 \(x\)。

\textbf{证明}：\(x\) 属于 \(\{x_n\}\) 的弱闭包的凸包。由凸集的弱闭等于强闭，\(x\) 属于强闭凸包，故存在凸组合序列强收敛。\(\blacksquare\)

\textbf{定义 74.6}（极点）：设 \(K\) 是向量空间的凸集。\(x \in K\) 称为 \(K\) 的\textbf{极点}（extreme point），如果 \(x\) 不能表示为 \(K\) 中两个不同点的严格凸组合，即若 \(x = ty + (1-t)z\)（\(t \in (0, 1), y, z \in K\)），则 \(y = z = x\)。\(K\) 的极点集记作 \(\operatorname{ext}(K)\)。

\textbf{定理 74.8}（Krein-Milman 定理）：设 \(K\) 是局部凸空间中的非空紧致凸集。则 \(K = \overline{\operatorname{conv}}(\operatorname{ext}(K))\)，即 \(K\) 是其极点的闭凸包。

\textbf{证明}：设 \(E = \overline{\operatorname{conv}}(\operatorname{ext}(K))\)。显然 \(E \subseteq K\)。若 \(E \neq K\)，取 \(x_0 \in K \setminus E\)。由 Hahn-Banach 分离定理，存在连续线性泛函 \(f\) 和 \(\alpha \in \mathbb{R}\) 使得 \(\Re f(x) \leq \alpha < \Re f(x_0)\) 对所有 \(x \in E\)。定义 \(M = \{x \in K : \Re f(x) = \max_{y \in K} \Re f(y)\}\)。由于 \(K\) 紧致，\(\max\) 可达。\(M\) 是 \(K\) 的非空紧致凸真子集（称为 \(K\) 的一个「面」）。由 Zorn 引理，\(K\) 的紧致凸真子集族中存在极小元 \(M_0\)。若 \(M_0\) 含有两个不同点，则存在分离它们的泛函，与 \(M_0\) 的极小性矛盾。故 \(M_0\) 为单点集 \(\{x^*\}\)，且 \(x^* \in \operatorname{ext}(K)\)。但 \(x_0 \in M\) 且 \(x^* \in M\)，故 \(\Re f(x^*) = \max\)。然而由构造，\(x^* \in E\) 时有 \(\Re f(x^*) \leq \alpha < \Re f(x_0)\)，而 \(x_0 \in K\) 时 \(\Re f(x_0) \leq \max\)，矛盾。因此 \(E = K\)。\(\blacksquare\)

\subsection{74.4 Hahn-Banach 定理的完整形式}\label{hahn-banach-ux5b9aux7406ux7684ux5b8cux6574ux5f62ux5f0f}

\textbf{定理 74.9}（Hahn-Banach 控制延拓定理，实情形）：设 \(X\) 为实向量空间，\(p : X \to \mathbb{R}\) 为次线性泛函，即满足

\[p(x + y) \leq p(x) + p(y), \quad p(tx) = t p(x) \; (\forall t \geq 0)\]

设 \(M \subseteq X\) 为子空间，\(f : M \to \mathbb{R}\) 为线性泛函满足 \(f(x) \leq p(x)\) 对所有 \(x \in M\)。则存在 \(f\) 的线性延拓 \(F : X \to \mathbb{R}\) 使得 \(F(x) \leq p(x)\) 对所有 \(x \in X\)。

\textbf{证明}：考虑所有满足 \(g\) 线性、包含 \(f\) 且 \(g \leq p\) 的延拓 \((M_g, g)\) 构成的集合 \(\mathcal{P}\)。定义偏序 \((M_1, g_1) \preceq (M_2, g_2)\) 当且仅当 \(M_1 \subseteq M_2\) 且 \(g_2|_{M_1} = g_1\)。\(\mathcal{P}\) 中每个链的上界为并集上的自然延拓。由 Zorn 引理，存在极大元 \((M', F)\)。

若 \(M' \neq X\)，取 \(x_0 \in X \setminus M'\)。对任意 \(y, z \in M'\)，

\[F(y) + F(z) = F(y+z) \leq p(y+z) \leq p(y-x_0) + p(z+x_0)\]

移项得 \(F(y) - p(y - x_0) \leq p(z + x_0) - F(z)\)。令

\[a = \sup_{y \in M'} \left( F(y) - p(y - x_0) \right), \quad b = \inf_{z \in M'} \left( p(z + x_0) - F(z) \right)\]

则 \(a \leq b\)。选取 \(F(x_0) \in [a, b]\)，在 \(M' \oplus \mathbb{R}x_0\) 上线性延拓 \(F\)。对任意 \(t > 0\)，

\[F(m + t x_0) = F(m) + tF(x_0) \leq F(m) + t\left(p\left(\frac{m}{t} + x_0\right) - F\left(\frac{m}{t}\right)\right) = p(m + t x_0)\]

\(t < 0\) 时用 \(a\) 的定义同样得证。故 \((M' \oplus \mathbb{R}x_0, F)\) 严格大于 \((M', F)\)，与极大性矛盾。因此 \(M' = X\)。\(\blacksquare\)

\textbf{定理 74.10}（Hahn-Banach 复数形式）：设 \(X\) 为复向量空间，\(p\) 为半范数（\(p(x+y) \leq p(x)+p(y)\)，\(p(\lambda x)=|\lambda|p(x)\)）。设 \(M \subseteq X\) 为子空间，\(f : M \to \mathbb{C}\) 为复线性泛函满足 \(|f(x)| \leq p(x)\)。则存在复线性延拓 \(F : X \to \mathbb{C}\) 满足 \(|F(x)| \leq p(x)\)。

\textbf{证明}：令 \(g = \Re f\)（实线性泛函），\(g(x) \leq p(x)\)。由实情形延拓 \(g\) 为 \(G\) 满足 \(G \leq p\)。定义 \(F(x) = G(x) - i G(ix)\)。验证 \(F\) 为复线性且 \(F|_M = f\)。对 \(|F(x)|\)，写 \(F(x) = e^{i\theta} |F(x)|\)，则 \(|F(x)| = F(e^{-i\theta}x) = \Re F(e^{-i\theta}x) = G(e^{-i\theta}x) \leq p(e^{-i\theta}x) = p(x)\)。\(\blacksquare\)

\textbf{定理 74.11}（赋范空间保范延拓）：设 \(X\) 为赋范空间，\(M \subseteq X\) 为子空间，\(f \in M^*\)。则存在保范延拓 \(F \in X^*\)，即 \(F|_M = f\) 且 \(\|F\|_{X^*} = \|f\|_{M^*}\)。

\textbf{证明}：取 \(p(x) = \|f\|_{M^*} \cdot \|x\|\)。\(p\) 为半范数（实情形为次线性泛函），且 \(|f(x)| \leq p(x)\) 在 \(M\) 上成立。由定理 74.9（或 74.10）得 \(F\) 满足 \(|F(x)| \leq \|f\| \cdot \|x\|\)，故 \(\|F\| \leq \|f\|\)。反向不等式显然因 \(F\) 延拓 \(f\)。\(\blacksquare\)

\textbf{推论 74.12}（Hahn-Banach 推论）：对任意 \(x \in X \setminus \{0\}\)，存在 \(f \in X^*\) 使得 \(\|f\| = 1\) 且 \(f(x) = \|x\|\)。故 \(X^*\) 分离 \(X\) 中的点，且 \(\|x\| = \sup_{\|f\| \leq 1} |f(x)|\)。

\textbf{定理 74.13}（分离超平面定理 ------ Hahn-Banach 几何形式）：设 \(A, B\) 为实赋范空间 \(X\) 中两个非空不交凸集，\(A\) 为开集。则存在 \(f \in X^* \setminus \{0\}\) 和 \(\alpha \in \mathbb{R}\) 使得

\[f(a) < \alpha \leq f(b), \quad \forall a \in A, \forall b \in B\]

\textbf{证明}：取 \(a_0 \in A, b_0 \in B\)，令 \(C = A - B + (b_0 - a_0)\)。\(C\) 为开凸集，\(0 \in C\)。定义 Minkowski 泛函 \(p_C(x) = \inf\{t > 0 : x/t \in C\}\)，则 \(p_C\) 为次线性泛函。令 \(z = b_0 - a_0 \notin C\)（否则 \(b_0 \in A\)），有 \(p_C(z) \geq 1\)。在 \(\mathbb{R}z\) 上定义 \(f_0(tz) = t\)，\(f_0(tz) \leq p_C(tz)\) 对 \(t \geq 0\) 成立。由 Hahn-Banach 延拓为 \(f\) 满足 \(f \leq p_C\)。在 \(C\) 上 \(f \leq p_C < 1\)，在 \(z\) 处 \(f(z) = 1\)。对 \(a \in A, b \in B\)，因 \(a - b + z \in C\)，得 \(f(a-b) + 1 = f(a-b+z) < 1\)，故 \(f(a) < f(b)\)。令 \(\alpha = \sup_A f\)。\(\blacksquare\)

\subsection{74.5 弱拓扑的半范数结构与自反性}\label{ux5f31ux62d3ux6251ux7684ux534aux8303ux6570ux7ed3ux6784ux4e0eux81eaux53cdux6027}

\textbf{定义 74.7}（弱拓扑 \(\sigma(X, X^*)\) 的半范数族）：弱拓扑 \(\sigma(X, X^*)\) 是由半范数族

\[\mathcal{P}_w = \left\{ p_f(x) = |f(x)| : f \in X^* \right\}\]

生成的局部凸拓扑。原点的邻域基由形如

\[V(f_1, \ldots, f_n ; \varepsilon) = \{ x \in X : |f_i(x)| < \varepsilon, \; i = 1, \ldots, n \}\]

的集合构成，其中 \(f_i \in X^*\)，\(n \in \mathbb{N}\)，\(\varepsilon > 0\)。

弱收敛 \(x_\alpha \rightharpoonup x\) 等价于对每个 \(f \in X^*\)，\(f(x_\alpha) \to f(x)\)。范数收敛蕴含弱收敛；当 \(\dim X < \infty\) 时二者等价。

\textbf{定义 74.8}（弱* 拓扑 \(\sigma(X^*, X)\) 的半范数族）：\(X^*\) 上的弱* 拓扑由半范数族

\[\mathcal{P}_{w^*} = \left\{ q_x(f) = |f(x)| : x \in X \right\}\]

生成。原点的邻域基为 \(W(x_1, \ldots, x_n ; \varepsilon) = \{ f \in X^* : |f(x_i)| < \varepsilon, \; i = 1, \ldots, n \}\)。

弱* 收敛 \(f_\alpha \xrightarrow{w^*} f\) 等价于 \(f_\alpha(x) \to f(x)\) 对所有 \(x \in X\)。

\textbf{定理 74.14}（自反性等价条件）：设 \(X\) 为 Banach 空间。下列条件等价：

\begin{enumerate}
\def\labelenumi{\arabic{enumi}.}
\tightlist
\item
  \(X\) 自反（\(J : X \to X^{**}\) 为满射）；
\item
  \(B_X\) 在弱拓扑 \(\sigma(X, X^*)\) 下紧致；
\item
  \(X^*\) 自反。
\end{enumerate}

\textbf{证明}：\((1) \Rightarrow (2)\)：若 \(X\) 自反，则 \((B_X, \sigma(X, X^*))\) 经 \(J\) 等同于 \((B_{X^{**}}, \sigma(X^{**}, X^*))\)，后者为 Banach-Alaoglu 中的弱* 紧集。\((2) \Rightarrow (1)\)：\(J(B_X)\) 弱* 紧从而弱* 闭，由 Goldstine 稠密性知 \(J(B_X) = B_{X^{**}}\)。\((1) \Leftrightarrow (3)\)：标准对偶论证------若 \(X\) 自反则 \(X^* \cong (X^{**})^* = (X^*)^{**}\)。反之若 \(X^*\) 自反，则 \(J(X)\) 为 \(X^{**}\) 的闭子空间，若 \(J(X) \neq X^{**}\)，由 Hahn-Banach 存在非零泛函零化 \(J(X)\)，与 \(X^*\) 自反矛盾。\(\blacksquare\)

\subsection{74.6 Mazur 引理与弱闭包}\label{mazur-ux5f15ux7406ux4e0eux5f31ux95edux5305}

\textbf{定理 74.15}（Mazur 引理 ------ 完整形式）：设 \(X\) 为 Banach 空间，\(C \subseteq X\) 为凸集。则 \(C\) 的弱闭包等于 \(C\) 的范数闭包：

\[\overline{C}^w = \overline{C}^{\|\cdot\|}\]

换言之，凸集的弱闭等于强闭。特别地，若 \(\{x_n\}\) 弱收敛到 \(x\)，则 \(x\) 属于 \(\{x_n\}\) 的范数闭凸包。

\textbf{证明}：强闭蕴含弱闭（弱拓扑粗于范数拓扑），故 \(\overline{C}^{\|\cdot\|} \subseteq \overline{C}^w\)。反之，若 \(x_0 \notin \overline{C}^{\|\cdot\|}\)，由分离超平面定理（定理 74.13，作用于闭凸集 \(\overline{C}^{\|\cdot\|}\) 和 \(\{x_0\}\)），存在 \(f \in X^*\) 和 \(\alpha \in \mathbb{R}\) 使得 \(\Re f(x_0) > \alpha \geq \sup_{y \in C} \Re f(y)\)。则 \(\{x : \Re f(x) > \alpha\}\) 是 \(x_0\) 的弱开邻域且与 \(C\) 不交，故 \(x_0 \notin \overline{C}^w\)。\(\blacksquare\)

\textbf{推论 74.16}：若 \(x_n \rightharpoonup x\)，则存在凸组合 \(\sum_{k=1}^{N_n} \lambda_k^{(n)} x_k\)（\(\lambda_k^{(n)} \geq 0, \sum \lambda_k^{(n)} = 1\)）依范数收敛到 \(x\)。

\subsection{74.7 Eberlein-Smulian 定理}\label{eberlein-smulian-ux5b9aux7406}

\textbf{定理 74.17}（Eberlein-Smulian 定理）：设 \(X\) 为 Banach 空间，\(A \subseteq X\)。则 \(A\) 的弱紧性等价于 \(A\) 的弱列紧性（每个序列有弱收敛子列）：

\[\sigma(X, X^*)\text{-紧} \; \Longleftrightarrow \; \sigma(X, X^*)\text{-列紧}\]

\textbf{证明要点}：\(\Leftarrow\)（非平凡方向）：先证 \(A\) 弱闭有界（弱列紧蕴含之）。取 \(A\) 的可数稠密子集，缩小至可分闭子空间 \(Y \subseteq X\)。关键在于：在可分 Banach 空间的单位球 \(B_{Y^*}\) 上，弱* 拓扑可度量化（因 \(B_Y\) 可分）。将 \(A\) 中序列的弱极限视为 \(C(B_{Y^*})\) 中的函数，用对角线法提取子列，从而将列紧性转化为紧性。\(\Rightarrow\)：类似论证，弱紧集的任意序列由 Eberlein 论证有弱收敛子列。\(\blacksquare\)

\textbf{推论 74.18}：\(X\) 自反当且仅当 \(B_X\) 中每个序列有弱收敛子列。

\textbf{证明}：若 \(X\) 自反，\(B_X\) 弱紧（定理 74.14），由 Eberlein-Smulian 推得弱列紧。反之，若 \(B_X\) 弱列紧，则对任意 \(F \in B_{X^{**}}\)，由 Goldstine 定理存在 \(x_n \in B_X\) 使 \(J(x_n) \xrightarrow{w^*} F\)；取弱收敛子列 \(x_{n_k} \rightharpoonup x \in B_X\)，则对任意 \(f \in X^*\)，\(F(f) = \lim J(x_{n_k})(f) = \lim f(x_{n_k}) = f(x) = J(x)(f)\)，故 \(F = J(x)\)，\(J\) 满射。\(\blacksquare\)

\hrulefill

\newpage

\chapter{06-分析/03-泛函分析/02-Hilbert空间.md}\label{ux5206ux679003-ux6cdbux51fdux5206ux679002-hilbertux7a7aux95f4.md}

\section{第174章：Hilbert 空间}\label{ux7b2c174ux7ae0hilbert-ux7a7aux95f4}

Hilbert 空间是带有内积的完备空间，其几何结构比一般 Banach 空间更丰富，有正交投影、正交基等概念。

\subsection{73.1 内积空间与 Hilbert 空间}\label{ux5185ux79efux7a7aux95f4ux4e0e-hilbert-ux7a7aux95f4}

\textbf{定义 73.1}（内积空间）：\(\mathbb{K}\) 上的向量空间 \(H\) 配以一个\textbf{内积} \(\langle \cdot, \cdot \rangle: H \times H \to \mathbb{K}\) 满足：

\begin{enumerate}
\def\labelenumi{\arabic{enumi}.}
\tightlist
\item
  \(\langle x, x \rangle \geq 0\)，且 \(\langle x, x \rangle = 0 \Leftrightarrow x = 0\)（正定性）
\item
  \(\langle x, y \rangle = \overline{\langle y, x \rangle}\)（共轭对称性）
\item
  \(\langle \alpha x + \beta y, z \rangle = \alpha \langle x, z \rangle + \beta \langle y, z \rangle\)（对第一变元的线性性）
\end{enumerate}

内积诱导范数 \(\|x\| = \sqrt{\langle x, x \rangle}\)。若 \((H, \langle \cdot, \cdot \rangle)\) 在该范数下完备，则称为 \textbf{Hilbert 空间}。

\textbf{定理 73.1}（Cauchy-Schwarz 不等式）：\(|\langle x, y \rangle| \leq \|x\| \|y\|\)，等号成立当且仅当 \(x\) 与 \(y\) 线性相关。

\emph{证明}：当 \(y = 0\) 时不等式显然成立（两边均为 \(0\)，且 \(x\) 与 \(0\) 线性相关）。设 \(y \neq 0\)。考虑 \(\lambda \in \mathbb{K}\) 的二次函数：\(0 \leq \|x - \lambda y\|^2 = \langle x - \lambda y, x - \lambda y \rangle = \|x\|^2 - \overline{\lambda}\langle x, y \rangle - \lambda\langle y, x \rangle + |\lambda|^2\|y\|^2 = \|x\|^2 - 2\Re(\overline{\lambda}\langle x, y \rangle) + |\lambda|^2\|y\|^2\)。取 \(\lambda = \langle x, y \rangle / \|y\|^2\)，则 \(\overline{\lambda}\langle x, y \rangle = |\langle x, y \rangle|^2 / \|y\|^2\)。代入得 \(0 \leq \|x\|^2 - 2|\langle x, y \rangle|^2/\|y\|^2 + |\langle x, y \rangle|^2/\|y\|^2 = \|x\|^2 - |\langle x, y \rangle|^2/\|y\|^2\)。整理得 \(|\langle x, y \rangle|^2 \leq \|x\|^2 \|y\|^2\)，开方即得 Cauchy-Schwarz 不等式。等号成立当且仅当 \(\|x - \lambda y\| = 0\)，即 \(x = \lambda y\)，换言之 \(x\) 与 \(y\) 线性相关。\(\blacksquare\)

\textbf{定理 73.2}（平行四边形法则与极化恒等式）： - 平行四边形法则：\(\|x + y\|^2 + \|x - y\|^2 = 2(\|x\|^2 + \|y\|^2)\) - 极化恒等式（实情形）：\(\langle x, y \rangle = \frac{1}{4}(\|x + y\|^2 - \|x - y\|^2)\) - 极化恒等式（复情形）：\(\langle x, y \rangle = \frac{1}{4}\sum_{k=0}^3 i^k \|x + i^k y\|^2\)

\textbf{证明}：平行四边形法则：\(\|x + y\|^2 + \|x - y\|^2 = (\|x\|^2 + 2\Re\langle x,y\rangle + \|y\|^2) + (\|x\|^2 - 2\Re\langle x,y\rangle + \|y\|^2) = 2(\|x\|^2 + \|y\|^2)\)。实极化恒等式：\(\|x + y\|^2 - \|x - y\|^2 = (\|x\|^2 + 2\langle x,y\rangle + \|y\|^2) - (\|x\|^2 - 2\langle x,y\rangle + \|y\|^2) = 4\langle x,y\rangle\)。复极化恒等式：直接展开 \(\sum_{k=0}^3 i^k \|x + i^k y\|^2\)，利用 \(\sum i^k = 0\) 和 \(\sum |i^k|^2 = 4\) 化简得 \(4\langle x,y\rangle\)。\(\blacksquare\)

\textbf{定理 73.3}（Jordan-von Neumann 定理）：赋范空间 \((X, \|\cdot\|)\) 的范数由内积诱导当且仅当范数满足平行四边形法则。此时内积可由极化恒等式恢复。

\textbf{证明}：（\(\Rightarrow\)）由定理 73.2，内积范数满足平行四边形法则。（\(\Leftarrow\)）在实情形定义 \(\langle x, y \rangle = \frac{1}{4}(\|x+y\|^2 - \|x-y\|^2)\)。由平行四边形法则可验证：对称性显然；\(\langle x, x \rangle = \|x\|^2 \geq 0\)；加性：利用平行四边形法则展开 \(\langle x+y, z \rangle\) 得 \(\langle x, z \rangle + \langle y, z \rangle\)；齐次性：先对整数和有理数成立，由连续性推广到实数。复情形类似。\(\blacksquare\)

Jordan-von Neumann 定理在 Banach 空间理论中具有根本重要性。它给出了一个赋范空间是否是 Hilbert 空间的几何判据：范数满足平行四边形法则的 Banach 空间一定是 Hilbert 空间。这一判据在许多重要的 Banach 空间（如 \(\ell^p\) 和 \(L^p\)）中迅速判定：\(\ell^p\) 和 \(L^p\) 是 Hilbert 空间当且仅当 \(p = 2\)。这一结果也解释了为什么 \(L^2\) 在分析中具有如此特殊的地位------它是唯一同时是 Hilbert 空间的 \(L^p\) 空间，从而兼具 Banach 空间的完备性和 Hilbert 空间的几何结构（正交性、投影等）。

\subsection{73.2 正交性与投影定理}\label{ux6b63ux4ea4ux6027ux4e0eux6295ux5f71ux5b9aux7406}

\textbf{定义 73.2}（正交）：\(x, y \in H\) 称为\textbf{正交的}，记作 \(x \perp y\)，如果 \(\langle x, y \rangle = 0\)。子集 \(A \subseteq H\) 的\textbf{正交补}为 \(A^\perp = \{x \in H : \langle x, a \rangle = 0, \; \forall a \in A\}\)。

\textbf{命题 73.4}（勾股定理）：若 \(x \perp y\)，则 \(\|x + y\|^2 = \|x\|^2 + \|y\|^2\)。

\textbf{证明}：由内积的双线性性和正交定义 \(x \perp y\)（即 \(\langle x, y \rangle = 0\)），展开 \(\|x + y\|^2 = \langle x + y, x + y \rangle = \langle x, x \rangle + \langle x, y \rangle + \langle y, x \rangle + \langle y, y \rangle = \|x\|^2 + 0 + 0 + \|y\|^2 = \|x\|^2 + \|y\|^2\)。由归纳法，若 \(\{x_1, \ldots, x_n\}\) 两两正交，则 \(\|\sum_{i=1}^n x_i\|^2 = \sum_{i=1}^n \|x_i\|^2\)。这一结果推广了 Euclid 几何中的勾股定理，是 Hilbert 空间几何结构的基石。\(\blacksquare\)

\textbf{定理 73.5}（正交投影定理）：设 \(C \subseteq H\) 是非空闭凸集。则对每个 \(x \in H\)，存在唯一的 \(y \in C\) 使得 \(\|x - y\| = \operatorname{dist}(x, C) = \inf_{z \in C} \|x - z\|\)。\(y\) 称为 \(x\) 到 \(C\) 的\textbf{最佳逼近}。

特别地，若 \(M \subseteq H\) 是闭子空间，则 \(H = M \oplus M^\perp\)，且 \(M^{\perp\perp} = M\)。

\emph{证明}：存在性：取极小化序列 \((y_n) \subseteq C\) 使 \(\|x - y_n\| \to d = \operatorname{dist}(x, C)\)。由平行四边形法则：

\[\|y_n - y_m\|^2 = 2\|y_n - x\|^2 + 2\|y_m - x\|^2 - \|y_n + y_m - 2x\|^2\]

由 \(C\) 凸，\((y_n + y_m)/2 \in C\)，故 \(\|y_n + y_m - 2x\|^2 \geq 4d^2\)。当 \(n, m \to \infty\) 时，\(\|y_n - y_m\| \to 0\)，\((y_n)\) 是 Cauchy 序列，由完备性收敛到某 \(y\)，由 \(C\) 闭知 \(y \in C\)。

唯一性：若 \(y, y'\) 都是最佳逼近，由平行四边形法则得 \(\|y - y'\|^2 \leq 4d^2 - 4d^2 = 0\)。

子空间情形：\(M^\perp\) 是闭子空间。对 \(x \in H\)，取 \(Px\) 为 \(x\) 到 \(M\) 的投影，则 \(x - Px \in M^\perp\)，故 \(x = Px + (x - Px) \in M + M^\perp\)。若 \(x \in M \cap M^\perp\)，则 \(\langle x, x \rangle = 0\)，故 \(x = 0\)。∎

正交投影定理是 Hilbert 空间理论的核心支柱。它断言：Hilbert 空间的任意闭子空间 \(M\) 都是可补的，且补空间恰好是正交补 \(M^\perp\)。这一性质在 Banach 空间中一般不成立------存在不可补的闭子空间（如 \(\ell^\infty\) 中 \(c_0\) 的闭子空间）。正交投影的存在性使得 Hilbert 空间中的许多问题（如最小二乘逼近、Ritz-Galerkin 方法、有限元分析）可以简化为在子空间上的正交投影，从而转化为线性代数问题。从几何上看，\(P_M x\) 是 \(x\) 在 \(M\) 上的''最佳逼近''，这一性质在统计学（线性回归的最小二乘估计）、信号处理（正交投影滤波）和量子力学（测量公设中的投影算子）中都有直接应用。

\subsection{73.3 Riesz 表示定理与正交基}\label{riesz-ux8868ux793aux5b9aux7406ux4e0eux6b63ux4ea4ux57fa}

\textbf{定理 73.6}（Riesz 表示定理）：设 \(H\) 是 Hilbert 空间。对每个连续线性泛函 \(f \in H^*\)，存在唯一的 \(y \in H\) 使得

\[f(x) = \langle x, y \rangle, \quad \forall x \in H\]

且 \(\|f\| = \|y\|\)。映射 \(y \mapsto \langle \cdot, y \rangle\) 是 \(H\) 到 \(H^*\) 的共轭线性等距同构（实情形为线性）。

\emph{证明}：若 \(f = 0\)，取 \(y = 0\)。否则，\(\ker f\) 是 \(H\) 的真闭子空间。取 \(z \in (\ker f)^\perp\) 且 \(\|z\| = 1\)。对任意 \(x \in H\)，\(x - f(x)z/f(z) \in \ker f\)，故 \(\langle x - f(x)z/f(z), z \rangle = 0\)，即 \(f(x) = \langle x, \overline{f(z)}z \rangle\)。取 \(y = \overline{f(z)}z\)。唯一性：若 \(\langle x, y_1 \rangle = \langle x, y_2 \rangle\) 对所有 \(x\)，则 \(\langle x, y_1 - y_2 \rangle = 0\)，取 \(x = y_1 - y_2\) 得 \(y_1 = y_2\)。∎

Riesz 表示定理是 Hilbert 空间区别于一般 Banach 空间的根本特征。它建立了 \(H\) 与 \(H^*\) 之间的反线性等距同构，使得 Hilbert 空间是''自对偶''的。这一事实在数学物理中有深远影响：在量子力学中，Dirac 的 bra-ket 记号 \(\langle \psi | \phi \rangle\) 正是 Riesz 表示定理的体现------ket \(|\phi\rangle\) 是 Hilbert 空间中的向量，bra \(\langle \psi|\) 是对偶空间中的泛函，而 Riesz 表示定理保证了 bra 与 ket 之间的自然对应。在偏微分方程中，Riesz 表示定理是 Lax-Milgram 定理（见下节）的基础，后者将椭圆型 PDE 的弱解存在唯一性问题转化为 Hilbert 空间中的变分问题。

\textbf{定义 73.3}（正交规范集）：\(H\) 的子集 \(\{e_\alpha\}_{\alpha \in A}\) 称为\textbf{正交规范的}，如果 \(\langle e_\alpha, e_\beta \rangle = \delta_{\alpha\beta}\)。

\textbf{定义 73.4}（正交规范基）：正交规范集 \(\{e_\alpha\}\) 称为 \(H\) 的\textbf{正交规范基}（或 Hilbert 基），如果 \(\overline{\operatorname{span}}\{e_\alpha\} = H\)（即其线性张成稠密）。

\textbf{定理 73.7}（正交规范基的存在性）：每个非零 Hilbert 空间都有正交规范基。任意两个正交规范基具有相同的基数，称为 \(H\) 的 \textbf{Hilbert 维数}。

\emph{证明}：存在性：考虑 \(H\) 中所有正交规范集构成的集合 \(\mathcal{O}\)，按包含关系 \(\subseteq\) 构成偏序集。\(\mathcal{O}\) 非空（单点集 \(\{x/\|x\|\}\) 对任意 \(x \neq 0\) 是正交规范集）。设 \(\mathcal{C} \subseteq \mathcal{O}\) 为全序子集（链），则 \(U = \bigcup_{S \in \mathcal{C}} S\) 也是正交规范集（若 \(x, y \in U\)，则存在 \(S \in \mathcal{C}\) 同时包含 \(x, y\)，故 \(\langle x, y \rangle = \delta_{xy}\)），故 \(U \in \mathcal{O}\) 是 \(\mathcal{C}\) 的上界。由 Zorn 引理，\(\mathcal{O}\) 有极大元 \(\mathcal{B}\)。若 \(\overline{\operatorname{span}}\mathcal{B} \neq H\)，则存在非零 \(x \in (\overline{\operatorname{span}}\mathcal{B})^\perp\)。令 \(e = x/\|x\|\)，则 \(\mathcal{B} \cup \{e\}\) 是严格更大的正交规范集，与 \(\mathcal{B}\) 的极大性矛盾。故 \(\mathcal{B}\) 为 \(H\) 的正交规范基。基的基数（Hilbert 维数）的不变性：设 \(\mathcal{B}_1, \mathcal{B}_2\) 为两个正交规范基。若其中之一有限，则显然 \(|\mathcal{B}_1| = |\mathcal{B}_2| = \dim H\)。若均无穷，对每个 \(e \in \mathcal{B}_1\)，\(e = \sum_{f \in \mathcal{B}_2} \langle e, f \rangle f\)，非零系数至多可数个。故 \(|\mathcal{B}_1| \leq \aleph_0 \cdot |\mathcal{B}_2| = |\mathcal{B}_2|\)，对称地 \(|\mathcal{B}_2| \leq |\mathcal{B}_1|\)，由 Schroder-Bernstein 定理得 \(|\mathcal{B}_1| = |\mathcal{B}_2|\)。∎

\textbf{定理 73.8}（Parseval 恒等式与 Fourier 展开）：设 \(\{e_\alpha\}_{\alpha \in A}\) 是 \(H\) 的正交规范基。则对任意 \(x \in H\)：

\begin{enumerate}
\def\labelenumi{\arabic{enumi}.}
\tightlist
\item
  \textbf{Bessel 不等式}：\(\sum_{\alpha} |\langle x, e_\alpha \rangle|^2 \leq \|x\|^2\)（事实上只有可数个系数非零）
\item
  \textbf{Fourier 展开}：\(x = \sum_{\alpha} \langle x, e_\alpha \rangle e_\alpha\)（在 \(H\) 的范数下收敛）
\item
  \textbf{Parseval 恒等式}：\(\|x\|^2 = \sum_{\alpha} |\langle x, e_\alpha \rangle|^2\)
\end{enumerate}

\emph{证明}：(1) Bessel 不等式：对任意有限子集 \(F \subseteq A\)，由勾股定理和正交规范性，\(\|x - \sum_{\alpha \in F} \langle x, e_\alpha \rangle e_\alpha\|^2 = \|x\|^2 - \sum_{\alpha \in F} |\langle x, e_\alpha \rangle|^2 \geq 0\)。因此 \(\sum_{\alpha \in F} |\langle x, e_\alpha \rangle|^2 \leq \|x\|^2\) 对任意有限 \(F\) 成立，取所有有限子集的上确界即得 Bessel 不等式。又对任意 \(\varepsilon > 0\)，集合 \(\{\alpha \in A : |\langle x, e_\alpha \rangle| \geq \varepsilon\}\) 必为有限（否则 Bessel 不等式被违反），故仅可数个系数非零。(2) Fourier 展开收敛性：设非零系数为 \(\{\langle x, e_n \rangle\}_{n=1}^\infty\)。令 \(S_N = \sum_{n=1}^N \langle x, e_n \rangle e_n\)。对 \(M > N\)，\(\|S_M - S_N\|^2 = \sum_{n=N+1}^M |\langle x, e_n \rangle|^2 \to 0\)（\(N \to \infty\)），由 Bessel 不等式知该级数收敛。故 \(\{S_N\}\) 为 Cauchy 列，由完备性收敛于某 \(y \in H\)。对任意 \(\alpha\)，\(\langle x - y, e_\alpha \rangle = \langle x, e_\alpha \rangle - \langle x, e_\alpha \rangle = 0\)，故 \(x - y \perp \operatorname{span}\{e_\alpha\}\)。由基的极大性，\(x - y = 0\)，即 \(x = \sum_n \langle x, e_n \rangle e_n\)。(3) Parseval 恒等式：由 Fourier 展开和级数的连续性，\(\|x\|^2 = \lim_{N \to \infty} \|S_N\|^2 = \lim_{N \to \infty} \sum_{n=1}^N |\langle x, e_n \rangle|^2 = \sum_\alpha |\langle x, e_\alpha \rangle|^2\)。\(\blacksquare\)

Parseval 恒等式是 Hilbert 空间理论中最优美的结果之一。它推广了 Fourier 分析中的经典 Parseval 等式：在 \(L^2([-\pi, \pi])\) 中，三角函数系 \(\{e^{inx}/\sqrt{2\pi}\}_{n \in \mathbb{Z}}\) 是正交规范基，Parseval 恒等式化为 \(\frac{1}{2\pi} \int_{-\pi}^\pi |f(x)|^2 dx = \sum_{n=-\infty}^\infty |\hat{f}(n)|^2\)。Parseval 恒等式在现代信号处理中对应于能量守恒：信号在时域中的总能量等于其在频域中的总能量（Plancherel 定理）。在量子力学中，波函数 \(\psi\) 在正交规范基 \(\{|\phi_n\rangle\}\) 下的展开系数 \(c_n = \langle \phi_n | \psi \rangle\) 的模平方 \(|c_n|^2\) 是测量得到本征值 \(\lambda_n\) 的概率，而 \(\sum |c_n|^2 = 1\) 正是总概率守恒。

\textbf{定理 73.9}（可分 Hilbert 空间的分类）：每个无穷维可分 Hilbert 空间等距同构于 \(\ell^2\)（配以标准正交基 \(\{e_n\}\)）。

\textbf{证明}：设 \(H\) 可分且有正交规范基 \(\{e_n\}_{n=1}^\infty\)（由定理 73.7 及可分性，基可数）。定义映射 \(U: H \to \ell^2\) 为 \(Ux = (\langle x, e_n \rangle)_{n=1}^\infty\)。由 Parseval 恒等式，\(\|Ux\|_{\ell^2}^2 = \sum |\langle x, e_n \rangle|^2 = \|x\|^2\)，故 \(U\) 为等距嵌入。满射性：对任意 \((c_n) \in \ell^2\)，令 \(x = \sum c_n e_n\)（在 \(H\) 中收敛，因 \(\sum|c_n|^2 < \infty\)），则 \(Ux = (c_n)\)。故 \(U\) 为等距同构。\(\blacksquare\)

\subsection{73.4 Lax-Milgram 定理与变分方法}\label{lax-milgram-ux5b9aux7406ux4e0eux53d8ux5206ux65b9ux6cd5}

\textbf{定理 73.10}（Lax-Milgram 定理）：设 \(H\) 是 Hilbert 空间，\(a: H \times H \to \mathbb{K}\) 是双线性形式（或共轭双线性形式），满足 (i) 有界性：\(|a(u, v)| \leq C \|u\| \|v\|\)（\(C > 0\)） (ii) 强制性（coercivity）：\(|a(u, u)| \geq \alpha \|u\|^2\)（\(\alpha > 0\)）

则对任意连续线性泛函 \(f \in H^*\)，存在唯一的 \(u \in H\) 使得 \(a(u, v) = f(v)\) 对所有 \(v \in H\) 成立。且 \(\|u\| \leq \frac{1}{\alpha} \|f\|\)。

\textbf{证明}：由 Riesz 表示定理，\(f(v) = \langle v, w \rangle\) 对某个 \(w \in H\)。对固定 \(u\)，\(v \mapsto a(u, v)\) 是连续线性泛函，由 Riesz 表示定理，存在 \(Au \in H\) 使得 \(a(u, v) = \langle v, Au \rangle\)（或 \(a(u, v) = \langle Au, v \rangle\) 取决于约定）。\(A: H \to H\) 是线性算子，由有界性 \(\|A\| \leq C\)。由强制性，\(\alpha \|u\|^2 \leq |a(u, u)| = |\langle u, Au \rangle| \leq \|u\| \|Au\|\)，故 \(\|Au\| \geq \alpha \|u\|\)------\(A\) 是下有界的，从而是单射且有闭值域。若 \(A\) 不满，则存在非零 \(v \perp \operatorname{im} A\)，则 \(0 = \langle v, Av \rangle = a(v, v) \geq \alpha \|v\|^2 > 0\)，矛盾。故 \(A\) 是满射，存在唯一 \(u = A^{-1} w\) 使 \(a(u, v) = \langle v, Au \rangle = \langle v, w \rangle = f(v)\)。\(\blacksquare\)

Lax-Milgram 定理是偏微分方程变分理论的核心工具。它将椭圆型 PDE 的弱解存在问题转化为双线性形式的强制性验证。标准例子是 Poisson 方程 \(-\Delta u = f\) 在 \(H^1_0(\Omega)\) 上的 Dirichlet 问题：取 \(a(u, v) = \int_\Omega \nabla u \cdot \nabla v\)，Poincare 不等式保证 \(a\) 在 \(H^1_0\) 上是强制的，由 Lax-Milgram 定理得弱解的存在唯一性。Lax-Milgram 定理的推广------\textbf{Babuška-Lax-Milgram 定理}（inf-sup 条件）------处理了 \(a\) 在不同空间 \(H_1 \times H_2\) 上定义的更一般情形，是混合有限元方法（如 Stokes 方程的 Taylor-Hood 元）的数学基础。

\subsection{73.5 伴随算子与紧自伴算子的谱定理}\label{ux4f34ux968fux7b97ux5b50ux4e0eux7d27ux81eaux4f34ux7b97ux5b50ux7684ux8c31ux5b9aux7406}

\textbf{定义 73.5}（伴随算子）：设 \(H_1, H_2\) 为 Hilbert 空间，\(T \in \mathcal{B}(H_1, H_2)\)。\(T\) 的\textbf{伴随算子} \(T^* \in \mathcal{B}(H_2, H_1)\) 由下式唯一确定：

\[\langle Tx, y \rangle_{H_2} = \langle x, T^* y \rangle_{H_1}, \quad \forall x \in H_1, y \in H_2\]

\(T^*\) 的存在唯一性由 Riesz 表示定理保证：对固定 \(y\)，\(x \mapsto \langle Tx, y \rangle\) 是连续线性泛函，由 Riesz 表示存在唯一 \(T^*y\)。

\textbf{定义 73.6}（自伴算子、正规算子、酉算子）：设 \(T \in \mathcal{B}(H)\)。 - \(T\) 称为\textbf{自伴的}（或 Hermitian），如果 \(T^* = T\)。 - \(T\) 称为\textbf{正规的}，如果 \(T^*T = TT^*\)。 - \(T\) 称为\textbf{酉算子}，如果 \(T^*T = TT^* = I\)（即 \(T\) 是等距满射）。

\textbf{命题 73.11}（伴随算子的性质）： 1. \((T^*)^* = T\)，\(\|T^*\| = \|T\|\)，\(\|T^*T\| = \|T\|^2\)（C*-等式） 2. \((S+T)^* = S^* + T^*\)，\((\alpha T)^* = \overline{\alpha} T^*\)，\((ST)^* = T^* S^*\) 3. \(\ker T^* = (\operatorname{im} T)^\perp\)，\(\overline{\operatorname{im} T^*} = (\ker T)^\perp\)

\textbf{证明}：(1) 对任意 \(x \in H_1, y \in H_2\)，\(\langle (T^*)^* x, y \rangle = \langle x, T^* y \rangle = \overline{\langle T^* y, x \rangle} = \overline{\langle y, T x \rangle} = \langle T x, y \rangle\)，故 \((T^*)^* = T\)。\(\|T^*\| = \|T\|\)：由 \(\|T^* y\| = \sup_{\|x\|=1} |\langle x, T^* y \rangle| = \sup_{\|x\|=1} |\langle T x, y \rangle| \leq \|T\|\|y\|\)，得 \(\|T^*\| \leq \|T\|\)；反向不等式由 \((T^*)^* = T\) 得 \(\|T\| = \|(T^*)^*\| \leq \|T^*\|\)。C*-等式：\(\|T^* T\| \leq \|T^*\|\|T\| = \|T\|^2\)；另一方面，\(\|T\|^2 = \sup_{\|x\|=1} \langle T x, T x \rangle = \sup_{\|x\|=1} \langle T^* T x, x \rangle \leq \|T^* T\|\)，故 \(\|T^* T\| = \|T\|^2\)。(2) 直接由伴随的定义和内积的线性性/共轭线性性验证：\(\langle (S+T)x, y \rangle = \langle Sx, y \rangle + \langle Tx, y \rangle = \langle x, S^* y \rangle + \langle x, T^* y \rangle = \langle x, (S^* + T^*)y \rangle\)；\(\langle \alpha T x, y \rangle = \alpha \langle T x, y \rangle = \alpha \langle x, T^* y \rangle = \langle x, \overline{\alpha} T^* y \rangle\)；\(\langle ST x, y \rangle = \langle T x, S^* y \rangle = \langle x, T^* S^* y \rangle\)。(3) \(y \in \ker T^* \iff T^* y = 0 \iff \langle x, T^* y \rangle = 0 \; (\forall x) \iff \langle T x, y \rangle = 0 \; (\forall x) \iff y \in (\operatorname{im} T)^\perp\)。取正交补得 \(\overline{\operatorname{im} T^*} = (\ker T)^\perp\)。\(\blacksquare\)

\textbf{定理 73.12}（紧自伴算子的谱定理 / Hilbert-Schmidt 展开定理）：设 \(H\) 是 Hilbert 空间，\(T \in \mathcal{K}(H)\) 是自伴紧算子。则存在 \(H\) 的正交规范基 \(\{e_n\}\) 和实数 \(\lambda_n \to 0\) 使得

\[Tx = \sum_{n} \lambda_n \langle x, e_n \rangle e_n\]

即 \(T\) 可以对角化，其特征向量构成正交规范基。非零特征值 \(\lambda_n\) 的几何重数有限，且 \(\|T\| = \max_n |\lambda_n|\)。

\textbf{证明}：首先证明 \(\|T\|\) 或 \(-\|T\|\) 是 \(T\) 的特征值。取单位向量序列 \((x_n)\) 使 \(|\langle Tx_n, x_n \rangle| \to \|T\|\)（由自伴性，\(\|T\| = \sup_{\|x\|=1} |\langle Tx, x \rangle|\)）。由紧性，\((Tx_n)\) 有收敛子列，极限论证得特征向量 \(e_1\) 使 \(Te_1 = \lambda_1 e_1\)，\(|\lambda_1| = \|T\|\)。然后考虑 \(H_1 = \{e_1\}^\perp\)，\(T\) 限制在 \(H_1\) 上仍是自伴紧算子，归纳构造正交特征向量列 \(\{e_n\}\) 和特征值 \(\lambda_n\)，\(|\lambda_{n+1}| \leq |\lambda_n|\)。由紧性，\(\lambda_n \to 0\)。最后证明 \(\{e_n\}\) 是基：若 \(x \perp e_n\) 对所有 \(n\)，则 \(Tx = 0\)。同时 \(T\) 在 \(\{e_n\}^\perp\) 上为零。故任意 \(x\) 可展开为 \(\sum \langle x, e_n \rangle e_n + (x - \sum \langle x, e_n \rangle e_n)\)，后者在 \(\ker T\) 中。\(\blacksquare\)

\textbf{定理 73.13}（Hilbert-Schmidt 算子）：设 \(T: H \to H\) 为紧算子。\(T\) 称为 \textbf{Hilbert-Schmidt 算子}，如果对某个（从而任意）正交规范基 \(\{e_n\}\)，\(\sum_n \|T e_n\|^2 < \infty\)。Hilbert-Schmidt 算子的范数 \(\|T\|_{HS} = (\sum_n \|T e_n\|^2)^{1/2}\) 与基的选取无关。全体 Hilbert-Schmidt 算子构成 Hilbert 空间 \(\mathcal{HS}(H)\)，内积为 \(\langle S, T \rangle_{HS} = \sum_n \langle S e_n, T e_n \rangle\)。Hilbert-Schmidt 算子是紧算子的最重要子类，且 \(T\) 为 Hilbert-Schmidt 当且仅当 \(T^*T\) 的迹有限：\(\operatorname{Tr}(T^*T) < \infty\)。

\textbf{证明}：首先验证范数与基的选取无关。设 \(\{e_n\}\) 和 \(\{f_m\}\) 为两个正交规范基。由 Parseval 恒等式，\(\sum_n \|T e_n\|^2 = \sum_n \sum_m |\langle T e_n, f_m \rangle|^2 = \sum_m \sum_n |\langle e_n, T^* f_m \rangle|^2 = \sum_m \|T^* f_m\|^2\)。取 \(T^*\) 替代 \(T\) 得 \(\sum_n \|T e_n\|^2\) 与基无关。\(T\) 为 Hilbert-Schmidt 蕴含 \(T\) 紧：设 \(T\) 的奇异值分解为 \(T x = \sum_n s_n \langle x, e_n \rangle f_n\)，则 \(\|T\|_{HS}^2 = \sum_n s_n^2\)。若 \(\sum s_n^2 < \infty\)，则 \(s_n \to 0\)，故 \(T\) 为有限秩算子 \(\sum_{n=1}^N s_n \langle \cdot, e_n \rangle f_n\) 按范数极限，从而紧。\(\mathcal{HS}(H)\) 是 Hilbert 空间：内积 \(\langle S, T \rangle_{HS} = \sum_n \langle S e_n, T e_n \rangle\) 的完备性由 Cauchy 列的标准论证可得（类似 \(\ell^2\) 的完备性）。\(\operatorname{Tr}(T^* T) = \sum_n \langle T^* T e_n, e_n \rangle = \sum_n \|T e_n\|^2 = \|T\|_{HS}^2\)，故 \(T\) 为 Hilbert-Schmidt 当且仅当 \(\operatorname{Tr}(T^* T) < \infty\)。\(\blacksquare\)

\textbf{Mercer 定理}（对积分算子的推广）：设 \(K: [a,b] \times [a,b] \to \mathbb{R}\) 连续且对称正定（即 \(\iint K(x,y) f(x) f(y) dx dy \geq 0\) 对所有 \(f\)）。定义积分算子 \(T_K f(x) = \int_a^b K(x,y) f(y) dy\)。则 \(T_K\) 是 \(L^2([a,b])\) 上的自伴紧算子，且 \(K(x,y) = \sum_{n=1}^\infty \lambda_n e_n(x) e_n(y)\)，其中级数绝对且一致收敛。Mercer 定理是核方法（kernel methods）在机器学习中（如 SVM、高斯过程回归）的数学基础------核函数 \(K\) 的 Mercer 展开等价于特征映射 \(\phi(x) = (\sqrt{\lambda_n} e_n(x))_{n=1}^\infty\) 将数据嵌入到再生核 Hilbert 空间（RKHS）。

\hrulefill

\hrulefill

\hrulefill

\hrulefill

\hrulefill

\newpage

\chapter{06-分析/03-泛函分析/04-紧算子与谱.md}\label{ux5206ux679003-ux6cdbux51fdux5206ux679004-ux7d27ux7b97ux5b50ux4e0eux8c31.md}

\section{第175章：紧算子}\label{ux7b2c175ux7ae0ux7d27ux7b97ux5b50}

紧算子理论是泛函分析中最优美和最有用的部分之一，由 F. Riesz 在 1918 年系统建立。紧算子构成了''无穷维空间中的有限秩近似''，其谱理论完美地推广了有限维线性代数中矩阵的特征值理论。紧算子理论的标志性成果是 Riesz-Schauder 定理------它断言紧算子的非零谱点都是孤立特征值，且具有有限重数。这一结果将 Fredholm 积分方程理论和 Hilbert 空间上的自伴算子理论统一在泛函分析的框架中。

\subsection{55.1 紧算子的定义与基本性质}\label{ux7d27ux7b97ux5b50ux7684ux5b9aux4e49ux4e0eux57faux672cux6027ux8d28}

\textbf{定义 55.1.1（紧算子）}：设 \(X, Y\) 为 Banach 空间。线性算子 \(T : X \to Y\) 称为\textbf{紧算子}（或全连续算子），若 \(T\) 将 \(X\) 中的有界集映为 \(Y\) 中的相对紧集（即其闭包为紧集）。等价地，对 \(X\) 中任意有界序列 \(\{x_n\}\)，\(\{T x_n\}\) 在 \(Y\) 中有收敛子列。

\textbf{注记 55.1.1}：紧算子的直观意义是''将无穷维信息压缩为近似有限维信息''。在有限维空间中，所有线性算子都是紧的（因为有界集的闭包紧）；在无穷维空间中，恒等算子 \(I\) 不是紧的，这正是紧算子理论有意义的根源。紧算子与有限秩算子密切相关------有限秩算子（值域有限维的算子）是紧算子的基本例子，而紧算子可视为有限秩算子的极限（在算子范数下）。

\textbf{定理 55.1.1（紧算子的基本性质）}：设 \(X, Y, Z\) 为 Banach 空间。

\begin{enumerate}
\def\labelenumi{(\roman{enumi})}
\tightlist
\item
  紧算子必有界（连续）
\item
  紧算子的线性组合仍是紧算子（紧算子构成 \(\mathcal{B}(X, Y)\) 的闭子空间）
\item
  若 \(T\) 紧，\(S\) 有界，则 \(T S\) 和 \(S T\) 紧（紧算子构成算子代数中的理想）
\item
  紧算子列在算子范数下的极限仍是紧算子（\(\mathcal{K}(X, Y)\) 是 \(\mathcal{B}(X, Y)\) 的闭子空间）
\end{enumerate}

\textbf{证明}：(i) 若 \(T\) 无界，则存在 \(\{x_n\}\)，\(\|x_n\| = 1\) 但 \(\|T x_n\| \to \infty\)，此时 \(\{T x_n\}\) 无收敛子列，矛盾。

\begin{enumerate}
\def\labelenumi{(\roman{enumi})}
\setcounter{enumi}{1}
\item
  设 \(T_1, T_2\) 紧，\(\alpha, \beta \in \mathbb{C}\)。对任意有界序列 \(\{x_n\}\)，取子列使 \(T_1 x_{n_k}\) 收敛，再取子子列使 \(T_2 x_{n_{k_j}}\) 收敛，则 \(\alpha T_1 x_{n_{k_j}} + \beta T_2 x_{n_{k_j}}\) 收敛。
\item
  设 \(T : X \to Y\) 紧，\(S : Y \to Z\) 有界。\(\{x_n\}\) 有界 \(\Rightarrow\) \(\{T x_n\}\) 有收敛子列 \(\Rightarrow\) \(\{S T x_{n_k}\}\) 收敛（\(S\) 连续）。同理 \(T S\) 紧。
\item
  设 \(T_n \to T\) 在算子范数下，\(T_n\) 紧。对 \(\{x_k\}\) 有界，用对角线法取子列使对每个 \(n\)，\(\{T_n x_{k_j}\}\) 收敛。由 \(\|T x_{k_j} - T x_{k_i}\| \leq \|T - T_n\|(\|x_{k_j}\| + \|x_{k_i}\|) + \|T_n x_{k_j} - T_n x_{k_i}\|\)，先取 \(n\) 大使第一项小，再取 \(j, i\) 大使第二项小，得 \(\{T x_{k_j}\}\) 为 Cauchy 列。\(\blacksquare\)
\end{enumerate}

\textbf{命题 55.1.2（有限秩算子的稠密性）}：有限秩算子（值域为有限维空间的算子）在 \(\mathcal{K}(X, Y)\) 中按算子范数稠密（当 \(Y\) 具有 Schauder 基时，或更一般地当 \(Y\) 具有有界逼近性质时）。在 Hilbert 空间情形，有限秩算子在紧算子中总是稠密的，因为 Hilbert 空间总有规范正交基。

\textbf{证明}：以下在 Hilbert 空间情形证明。设 \(H_1, H_2\) 为 Hilbert 空间，\(T \in \mathcal{K}(H_1, H_2)\)。取 \(H_2\) 的规范正交基 \(\{e_\alpha\}_{\alpha \in A}\)。对任意 \(\varepsilon > 0\)，由于 \(T(B_{H_1})\)（单位球的像）相对紧，存在有限 \(\varepsilon\)-网。特别地，存在有限秩正交投影 \(P_N\)（投影到 \(\{e_1, \ldots, e_N\}\) 张成的子空间）使得 \(\|(I - P_N)T\| < \varepsilon\)。这是因为对任意 \(x \in B_{H_1}\)，\(Tx\) 可由有限个基向量 \(\varepsilon\)-逼近，而 \((I - P_N)\) 是到 \(\{e_1, \ldots, e_N\}^\perp\) 上的投影，当 \(N\) 足够大时 \(\|(I - P_N)Tx\| < \varepsilon\) 对 \(x\) 一致成立（由紧性）。于是 \(P_N T\) 是有限秩算子（值域包含在 \(P_N H_2\) 中，维数 \(\leq N\)），且 \(\|T - P_N T\| = \|(I - P_N)T\| < \varepsilon\)。故 \(T\) 是有限秩算子的极限。在 Banach 空间情形，当 \(Y\) 具有 Schauder 基 \(\{e_n\}\) 时，令 \(P_n: Y \to \operatorname{span}\{e_1, \ldots, e_n\}\) 为自然投影。由基的性质，\(P_n y \to y\) 对每个 \(y \in Y\)，且在紧集上一致收敛，故 \(\|P_n T - T\| \to 0\)。\(\blacksquare\)

\textbf{推论 55.1.3}：紧算子的伴随也是紧算子：若 \(T \in \mathcal{K}(X, Y)\)，则 \(T^* \in \mathcal{K}(Y^*, X^*)\)。在 Hilbert 空间情形，\(T \in \mathcal{K}(H)\) 当且仅当 \(T^* \in \mathcal{K}(H)\)。

\textbf{证明}：设 \(T \in \mathcal{K}(X, Y)\)。取 \(Y^*\) 的单位球 \(B_{Y^*}\) 中的任意序列 \(\{y_n^*\}\)。需证 \(\{T^* y_n^*\}\) 有收敛子列。考虑 \(T(B_X)\) 的相对紧性。对任意 \(\varepsilon > 0\)，存在有限 \(\varepsilon\)-网 \(\{T x_1, \ldots, T x_m\}\) 覆盖 \(T(B_X)\)。由对角线法，取子列使得 \(\{T^* y_{n_k}^*(x_j)\} = \{y_{n_k}^*(T x_j)\}\) 对每个 \(j = 1, \ldots, m\) 收敛（有界数列必有收敛子列）。则对任意 \(x \in B_X\)，取 \(x_j\) 使 \(\|T x - T x_j\| < \varepsilon\)，则 \(|T^* y_{n_k}^*(x) - T^* y_{n_l}^*(x)| \leq |y_{n_k}^*(T x - T x_j)| + |y_{n_k}^*(T x_j) - y_{n_l}^*(T x_j)| + |y_{n_l}^*(T x_j - T x)| \leq 2\varepsilon + |y_{n_k}^*(T x_j) - y_{n_l}^*(T x_j)|\)。当 \(k, l\) 足够大时最后一项 \(< \varepsilon\)，故 \(\{T^* y_{n_k}^*\}\) 为 Cauchy 列，在 \(X^*\) 中收敛。在 Hilbert 空间情形，\(T\) 紧等价于 \(T^*T\) 紧（自伴紧），且 \(T^*\) 由 \(T^*T\) 的谱分解直接构造，二者互为伴随。\(\blacksquare\)

\hrulefill

\subsection{55.2 Fredholm 积分算子与紧算子的原型}\label{fredholm-ux79efux5206ux7b97ux5b50ux4e0eux7d27ux7b97ux5b50ux7684ux539fux578b}

\textbf{定义 55.2.1（Fredholm 积分算子）}：设 \(K \in L^2([a, b] \times [a, b])\)。Fredholm 积分算子 \(T_K : L^2([a, b]) \to L^2([a, b])\) 定义为

\[
(T_K f)(x) = \int_a^b K(x, y) f(y) \, dy
\]

\textbf{定理 55.2.1（Fredholm 积分算子的紧性）}：若 \(K \in L^2([a, b] \times [a, b])\)，则 \(T_K\) 是 \(L^2([a, b])\) 上的紧算子。更一般地，若 \(K\) 是 Hilbert-Schmidt 核（即 \(\iint |K(x, y)|^2 dx dy < \infty\)），则 \(T_K\) 是 Hilbert-Schmidt 算子，从而是紧算子。

\textbf{证明}：设 \(\{e_n\}\) 为 \(L^2([a,b])\) 的规范正交基，则 \(\{e_n(x) \overline{e_m(y)}\}_{n,m}\) 为 \(L^2([a,b] \times [a,b])\) 的规范正交基。展开 \(K(x,y) = \sum_{n,m} a_{nm} e_n(x) \overline{e_m(y)}\)，其中 \(\sum |a_{nm}|^2 = \|K\|_{L^2([a,b]^2)}^2 < \infty\)。截断 \(K_N(x,y) = \sum_{n,m \leq N} a_{nm} e_n(x) \overline{e_m(y)}\)，则 \(T_{K_N}\) 是有限秩算子（值域张成于 \(\{e_1, \ldots, e_N\}\)）。且 \(\|T_K - T_{K_N}\| \leq \|K - K_N\|_{L^2([a,b]^2)} \to 0\)（\(N \to \infty\)）。由紧算子的闭性，\(T_K\) 紧。\(\blacksquare\)

这一结论揭示了紧算子理论的历史起源：Fredholm 在 1900 年前后研究积分方程 \((I + T_K)f = g\) 时，发现 \(T_K\) 的行为类似有限维线性算子------这正是紧算子理论的雏形。

\hrulefill

\subsection{55.3 紧算子的谱理论------Riesz-Schauder 定理}\label{ux7d27ux7b97ux5b50ux7684ux8c31ux7406ux8bbariesz-schauder-ux5b9aux7406}

\textbf{定义 55.3.1（预解集与谱）}：设 \(T \in \mathcal{B}(X)\)。\(T\) 的\textbf{预解集} \(\rho(T) = \{\lambda \in \mathbb{C} \mid \lambda I - T \text{ 可逆}\}\)。\(T\) 的\textbf{谱} \(\sigma(T) = \mathbb{C} \setminus \rho(T)\)。谱可分为：

\begin{itemize}
\tightlist
\item
  \textbf{点谱} \(\sigma_p(T)\)：\(\lambda I - T\) 不是单射（即 \(\lambda\) 是特征值）
\item
  \textbf{连续谱} \(\sigma_c(T)\)：\(\lambda I - T\) 是单射，值域稠密但非满射
\item
  \textbf{剩余谱} \(\sigma_r(T)\)：\(\lambda I - T\) 是单射，但值域不稠密
\end{itemize}

\textbf{定理 55.3.1（紧算子的谱定理 / Riesz-Schauder 定理）}：设 \(T \in \mathcal{K}(X)\)（\(X\) 为无穷维 Banach 空间）。则

\begin{enumerate}
\def\labelenumi{(\roman{enumi})}
\tightlist
\item
  \(0 \in \sigma(T)\)（\(T\) 在无穷维空间中不可逆）
\item
  \(\sigma(T) \setminus \{0\}\) 由至多可数个特征值组成，\(0\) 是其唯一可能的聚点
\item
  每个非零特征值 \(\lambda \in \sigma(T) \setminus \{0\}\) 的几何重数 \(\dim \ker(\lambda I - T)\) 有限
\item
  \(\lambda I - T\) 是 Fredholm 算子（\(\ker\) 有限维，\(\operatorname{ran}\) 有限余维，且 \(\dim \ker(\lambda I - T) = \dim \operatorname{coker}(\lambda I - T)\)）
\end{enumerate}

\textbf{证明}（框架，以 Hilbert 空间为例）：

\begin{enumerate}
\def\labelenumi{(\roman{enumi})}
\tightlist
\item
  若 \(T\) 可逆，则 \(I = T^{-1} T\) 是紧算子与有界算子的复合，故 \(I\) 紧，但无穷维 Hilbert 空间中单位球非紧，矛盾。
\end{enumerate}

(ii)-(iii) 设 \(\lambda \neq 0\)。令 \(V = \ker(\lambda I - T)\)。若 \(\dim V = \infty\)，则存在规范正交列 \(\{e_n\} \subseteq V\)，满足 \(T e_n = \lambda e_n\)。由 \(\{e_n\}\) 有界且 \(T\) 紧，\(\{T e_n\} = \{\lambda e_n\}\) 有收敛子列，但 \(\|\lambda e_n - \lambda e_m\| = |\lambda|\sqrt{2}\)（\(n \neq m\)），矛盾。故 \(\dim V < \infty\)。

对特征值的离散性：假设存在非零特征值列 \(\lambda_n \to \lambda \neq 0\)，\(\lambda_n \neq \lambda_m\)。取对应的单位特征向量 \(e_n\)。由 Riesz 引理，\(\{e_n\}\) 满足 \(\|e_n - e_m\| \geq 1\)（\(\forall n \neq m\)）。考虑 \(T e_n = \lambda_n e_n\)，则 \(\{T e_n\}\) 有收敛子列（\(T\) 紧），但 \(\|T e_n - T e_m\| = \|\lambda_n e_n - \lambda_m e_m\| \geq |\lambda|/2\)（当 \(n,m\) 足够大时），矛盾。故 \(\sigma(T) \setminus \{0\}\) 无聚点（除 \(0\) 外）。

\begin{enumerate}
\def\labelenumi{(\roman{enumi})}
\setcounter{enumi}{3}
\tightlist
\item
  由 Fredholm 二择一定理：\(\lambda I - T\) 或为单射（此时 \(\lambda \notin \sigma(T)\)），或 \(\lambda\) 是特征值（\(\dim \ker(\lambda I - T) < \infty\)）。此时 \(\lambda\) 也是 \(T^*\) 的特征值，且 \(\dim \ker(\lambda I - T) = \dim \ker(\overline{\lambda} I - T^*)\)。在 Hilbert 空间中，\(\operatorname{ran}(\lambda I - T) = \ker(\overline{\lambda} I - T^*)^\perp\)，故 \(\operatorname{codim} \operatorname{ran}(\lambda I - T) = \dim \ker(\overline{\lambda} I - T^*) = \dim \ker(\lambda I - T)\)。\(\blacksquare\)
\end{enumerate}

\hrulefill

\subsection{55.4 Fredholm 二择一与 Fredholm 理论}\label{fredholm-ux4e8cux62e9ux4e00ux4e0e-fredholm-ux7406ux8bba}

\textbf{定理 55.4.1（Fredholm 二择一定理）}：设 \(T \in \mathcal{K}(X)\)，\(\lambda \neq 0\)。则以下二者恰有一个成立：

\begin{enumerate}
\def\labelenumi{\arabic{enumi}.}
\tightlist
\item
  齐次方程 \((\lambda I - T)x = 0\) 只有零解，此时非齐次方程 \((\lambda I - T)x = y\) 对任意 \(y \in X\) 有唯一解
\item
  齐次方程 \((\lambda I - T)x = 0\) 有非零解，此时非齐次方程 \((\lambda I - T)x = y\) 有解当且仅当 \(y\) 满足 \(\langle y, v \rangle = 0\)（对所有 \(v \in \ker(\overline{\lambda} I - T^*)\)）
\end{enumerate}

这一结构完全类似于有限维线性代数中的情形：\(\lambda I - T\) 要么可逆，要么 \(\lambda\) 是特征值且值域由伴随齐次方程的正交条件刻画。

\textbf{证明}：设 \(T \in \mathcal{K}(X)\)，\(\lambda \neq 0\)。令 \(S = \lambda I - T\)。首先注意，若 \(S\) 为单射，则 \(S\) 必为满射（从而可逆）。这是因为若 \(S\) 是单射但非满射，则 \(S(X)\) 为 \(X\) 的真闭子空间（由 Riesz-Schauder 理论，\(S\) 有闭值域）。令 \(X_1 = S(X)\)，则 \(X_1 \subsetneq X\)。由归纳构造 \(X_n = S(X_{n-1})\)，则 \(X_n \supsetneq X_{n+1}\)（严格包含）。由 Riesz 引理，存在 \(x_n \in X_n\)，\(\|x_n\| = 1\)，\(\operatorname{dist}(x_n, X_{n+1}) \geq 1/2\)。于是对 \(n > m\)，\(T x_n - T x_m = \lambda x_n - S x_n - \lambda x_m + S x_m = \lambda x_n - (\text{某元素在 } X_{m+1})\)，故 \(\|T x_n - T x_m\| \geq |\lambda|/2\)，与 \(T\) 紧矛盾。因此 \(S\) 单射蕴含 \(S\) 满射。若 \(S\) 非单射，则 \(\lambda\) 为 \(T\) 的特征值。由定理 55.3.1，\(\dim \ker S < \infty\)。此时 \(S\) 的值域为 \(\ker S^*\) 的零化子，即 \(y \in S(X)\) 当且仅当 \(\langle y, v \rangle = 0\) 对所有 \(v \in \ker S^*\) 成立。在 Hilbert 空间情形，这等价于 \(S(X) = \ker S^*\) 的正交补。由 \(\dim \ker S = \dim \ker S^*\)（对紧算子 \(T\)），故 \(S(X)\) 的余维数为 \(\dim \ker S\)。\(\blacksquare\)

\textbf{定理 55.4.2（Fredholm 指标）}：对 \(\lambda \neq 0\)，定义 \(T_\lambda = \lambda I - T\) 的 Fredholm 指标为 \(\operatorname{ind}(T_\lambda) = \dim \ker(T_\lambda) - \dim \operatorname{coker}(T_\lambda)\)。对于紧算子 \(T\)，\(\operatorname{ind}(\lambda I - T) = 0\)（对所有 \(\lambda \neq 0\)）。更一般地，\(T\) 的任何紧扰动不改变 Fredholm 指标。

\textbf{证明}：由定理 55.3.1 (iv)，对紧算子 \(T\) 和 \(\lambda \neq 0\)，有 \(\dim \ker(\lambda I - T) = \dim \ker(\overline{\lambda} I - T^*)\)。在 Hilbert 空间中，\(\operatorname{ran}(\lambda I - T)\) 的闭性由 Riesz-Schauder 理论保证，且 \(\operatorname{coker}(\lambda I - T) = X / \operatorname{ran}(\lambda I - T) \cong \ker(\overline{\lambda} I - T^*)\)。因此 \(\dim \operatorname{coker}(\lambda I - T) = \dim \ker(\overline{\lambda} I - T^*) = \dim \ker(\lambda I - T)\)，故 \(\operatorname{ind}(\lambda I - T) = 0\)。对紧扰动不变性：设 \(K\) 紧，考虑连续族 \(t \mapsto \lambda I - T - tK\)（\(t \in [0,1]\)）。每个算子都是 \(\lambda I\) 的紧扰动，指标的连续性（在 Fredholm 算子的范数拓扑下）保证 \(\operatorname{ind}(\lambda I - T - K) = \operatorname{ind}(\lambda I - T) = 0\)。事实上，指标作为 Fredholm 算子上的局部常值函数，在连通分支上不变。\(\blacksquare\)

\hrulefill

\subsection{55.5 Hilbert 空间上的紧自伴算子}\label{hilbert-ux7a7aux95f4ux4e0aux7684ux7d27ux81eaux4f34ux7b97ux5b50}

\textbf{定义 55.5.1（自伴算子）}：Hilbert 空间 \(H\) 上的算子 \(T\) 称为\textbf{自伴}的，若 \(\langle T x, y \rangle = \langle x, T y \rangle\)（对所有 \(x, y \in H\)）。自伴紧算子的谱是实的，且具有正交特征向量系。

\textbf{定理 55.5.1（紧自伴算子的谱定理 / Hilbert-Schmidt 定理）}：设 \(T \in \mathcal{K}(H)\) 为自伴算子。则 \(H\) 存在由 \(T\) 的特征向量构成的规范正交基。具体地，\(T\) 可表示为

\[
T x = \sum_{n} \lambda_n \langle x, e_n \rangle e_n
\]

其中 \(\{\lambda_n\}\) 为 \(T\) 的非零特征值（每个按重数重复），\(\{e_n\}\) 为对应的规范正交特征向量，\(\lambda_n \to 0\)（若有无穷多个非零特征值）。

\textbf{证明}：分三步。

\begin{enumerate}
\def\labelenumi{\arabic{enumi}.}
\item
  非零特征值的存在性：若 \(T \neq 0\)，则 \(\|T\|\) 或 \(-\|T\|\) 是特征值（由 Rayleigh 商 \(\frac{\langle T x, x \rangle}{\langle x, x \rangle}\) 的极值性质）。取极大化序列 \(\{x_n\}\)（\(\|x_n\| = 1\)，\(\langle T x_n, x_n \rangle \to \|T\|\)），由 \(T\) 的紧性，存在子列使 \(T x_{n_k} \to y\)。验证 \(y = \|T\| x\)（具有极限 \(x\) 的子列），且 \(T x = \|T\| x\)。
\item
  谱分解：设 \(\lambda_1\) 为上述特征值，\(e_1\) 为对应特征向量。令 \(H_1 = \{e_1\}^\perp\)。\(T\) 限制在 \(H_1\) 上仍是紧自伴算子。重复此过程，得到特征值序列 \(\{\lambda_n\}\)（按绝对值不减排列）和对应的规范正交特征向量 \(\{e_n\}\)。若过程在有限步终止，则 \(T\) 在余下的子空间上为零算子。
\item
  完备性：若过程无穷，则 \(\lambda_n \to 0\)（否则，若 \(|\lambda_n| \geq \delta > 0\)，则 \(\{T e_n / \lambda_n\} = \{e_n\}\) 为有界序列但 \(\{T e_n\} = \{\lambda_n e_n\}\) 无收敛子列，与 \(T\) 紧矛盾）。对任意 \(x \in H\)，\(x - \sum_{k=1}^{n} \langle x, e_k \rangle e_k \in \operatorname{span}\{e_1, \ldots, e_n\}^\perp\)，故 \(\|T(x - \sum_{k=1}^{n} \langle x, e_k \rangle e_k)\| \leq |\lambda_{n+1}| \|x\| \to 0\)。因此 \(T x = \sum_{n} \lambda_n \langle x, e_n \rangle e_n\)。若特征向量系不完备，则存在非零 \(x\) 正交于所有 \(e_n\)，此时 \(T x = 0\)，即 \(0\) 为特征值，\(x\) 为特征向量，补充进 \(\{e_n\}\) 即得完备性。\(\blacksquare\)
\end{enumerate}

\textbf{推论 55.5.2（Hilbert-Schmidt 展开定理）}：设 \(T_K\) 为 \(L^2\)-核 \(K\) 的 Fredholm 积分算子（\(K\) 满足 \(K(x, y) = \overline{K(y, x)}\)------自伴性）。则 \(K\) 可展开为

\[
K(x, y) = \sum_{n} \lambda_n e_n(x) \overline{e_n(y)}
\]

其中 \(\{\lambda_n, e_n\}\) 为 \(T_K\) 的特征值与特征函数。该级数在 \(L^2([a, b]^2)\) 中收敛。

\textbf{证明}：由定理 55.5.1，自伴紧算子 \(T_K\) 有谱分解 \(T_K f = \sum_n \lambda_n \langle f, e_n \rangle e_n\)，其中 \(\{e_n\}\) 为 \(L^2([a,b])\) 的规范正交基。对固定的 \(x\)，将 \(K(x, \cdot)\) 视为 \(L^2([a,b])\) 中函数，由 Parseval 恒等式有 \(K(x, y) = \sum_n \langle K(x, \cdot), e_n \rangle e_n(y)\)（作为 \(y\) 的函数，级数在 \(L^2\) 中收敛）。但 \(\langle K(x, \cdot), e_n \rangle = \int_a^b K(x, y) \overline{e_n(y)} dy = \overline{T_K \overline{e_n}(x)}\)。由自伴性 \(K(x,y) = \overline{K(y,x)}\) 得 \(T_K e_n = \lambda_n e_n\)，故 \(\overline{T_K \overline{e_n}(x)} = \lambda_n e_n(x)\)。因此 \(\langle K(x, \cdot), e_n \rangle = \lambda_n e_n(x)\)。代入展开式得 \(K(x, y) = \sum_n \lambda_n e_n(x) \overline{e_n(y)}\)。在 \(L^2([a,b]^2)\) 中，\(\|K - \sum_{n=1}^N \lambda_n e_n \otimes \overline{e_n}\|_{L^2}^2 = \sum_{n=N+1}^\infty |\lambda_n|^2 \to 0\)（\(N \to \infty\)），因为 \(\sum |\lambda_n|^2 \leq \|T_K\|_{\mathcal{S}_2}^2 = \|K\|_{L^2}^2 < \infty\)。\(\blacksquare\)

\hrulefill

\subsection{55.6 紧算子的奇异值分解与 Schmidt 展开}\label{ux7d27ux7b97ux5b50ux7684ux5947ux5f02ux503cux5206ux89e3ux4e0e-schmidt-ux5c55ux5f00}

\textbf{定理 55.6.1（紧算子的奇异值分解 / Schmidt 展开）}：设 \(T \in \mathcal{K}(H_1, H_2)\)（\(H_1, H_2\) 为 Hilbert 空间）。则存在（有限或可数）规范正交系 \(\{e_n\} \subseteq H_1\)，\(\{f_n\} \subseteq H_2\) 和正数序列 \(s_n \downarrow 0\)（奇异值）使得

\[
T x = \sum_{n} s_n \langle x, e_n \rangle f_n, \quad T^* y = \sum_{n} s_n \langle y, f_n \rangle e_n
\]

\textbf{证明}：考虑紧自伴算子 \(T^* T \in \mathcal{K}(H_1)\)。由定理 55.5.1，\(T^* T\) 有特征分解 \(T^* T x = \sum s_n^2 \langle x, e_n \rangle e_n\)，其中 \(s_n = \sqrt{\lambda_n(T^* T)}\)。定义 \(f_n = T e_n / s_n\)（当 \(s_n > 0\)），则 \(\{f_n\}\) 规范正交。直接验证展开式。\(\blacksquare\)

奇异值分解是 Hilbert 空间中紧算子的最一般表示，不要求自伴性。它表明任何紧算子都可以通过''旋转-伸缩-旋转''（在正交基下）来实现------这一几何图像与有限维线性代数中矩阵的 SVD 完全类似。

\hrulefill

\subsection{55.7 迹类算子与 Hilbert-Schmidt 算子}\label{ux8ff9ux7c7bux7b97ux5b50ux4e0e-hilbert-schmidt-ux7b97ux5b50}

\textbf{定义 55.7.1（迹类算子）}：紧算子 \(T \in \mathcal{K}(H)\) 称为\textbf{迹类}的，若其奇异值 \(\{s_n\}\) 满足 \(\sum_{n} s_n < \infty\)。迹类算子构成 \(\mathcal{K}(H)\) 的理想 \(\mathcal{S}_1(H)\)，配以迹范数 \(\|T\|_1 = \sum s_n\)。

紧算子的奇异值衰减仅保证 \(s_n \to 0\)，但不同的衰减速率定义了紧算子族中的重要子类。迹类算子要求奇异值 \(\ell^1\)-可和，这是最严格的一类。与之相比，Hilbert-Schmidt 算子仅要求奇异值 \(\ell^2\)-可和，条件更宽松，但仍保有良好的内积结构。这两种算子类构成了无穷维空间中''矩阵论''的推广------它们分别对应于有限维空间中所有矩阵（迹类）和所有矩阵（Hilbert-Schmidt），其中的奇异值条件在有限维情形自动满足。

\textbf{定义 55.7.2（Hilbert-Schmidt 算子）}：紧算子 \(T \in \mathcal{K}(H)\) 称为 Hilbert-Schmidt 的，若 \(\sum s_n^2 < \infty\)。Hilbert-Schmidt 算子构成 \(\mathcal{K}(H)\) 的理想 \(\mathcal{S}_2(H)\)，配以范数 \(\|T\|_2 = (\sum s_n^2)^{1/2}\)。

迹类算子的核心特征在于其奇异值的绝对可和性，这使得我们可以对它们定义''迹''------这是有限维矩阵对角元素之和概念在无穷维空间中的推广。迹的定义需要谨慎处理，因为无穷维空间中算子的表示依赖于基的选取；然而迹类算子的优良性质保证了迹与基无关。这一不变性是以下 Lidskii 定理的基础。

\textbf{定义 55.7.3（迹）}：对迹类算子 \(T\)，定义 \(T\) 的\textbf{迹}为 \(\operatorname{Tr}(T) = \sum_{n} \langle T e_n, e_n \rangle\)，其中 \(\{e_n\}\) 为 Hilbert 空间的任意规范正交基。迹与基的选取无关，且满足 \(\operatorname{Tr}(AB) = \operatorname{Tr}(BA)\)（当 \(A\) 迹类、\(B\) 有界时）。

在有限维线性代数中，一个基本等式是：矩阵的迹等于其所有特征值之和。将这一事实推广到无穷维空间并非平凡------因为无穷维算子可能有无限多个非零特征值，特征值级数的收敛性本身就需要证明。Lidskii 定理（1959）完成了这一壮举：它断言迹类算子的迹恰好等于其所有非零特征值（按代数重数计）之和。这条定理位于算子代数、谱理论和数学物理的交汇点，是迹类算子理论中最深刻的结果之一。

\textbf{定理 55.7.1（Lidskii 定理）}：对迹类算子 \(T\)，\(\operatorname{Tr}(T) = \sum_{n} \lambda_n(T)\)，其中 \(\{\lambda_n(T)\}\) 为 \(T\) 的非零特征值（按代数重数计）。即迹等于特征值之和。Lidskii 定理将有限维线性代数中迹 = 特征值之和的结论推广到无穷维情形，是迹类算子理论中最深刻的定理之一。

\textbf{证明}（框架）：对 \(T \in \mathcal{S}_1(H)\)，构造其奇异值分解 \(T = \sum_n s_n \langle \cdot, e_n \rangle f_n\)。将 \(T\) 分解为 \(T = T_1 + i T_2\)，其中 \(T_1 = (T+T^*)/2\) 和 \(T_2 = (T-T^*)/(2i)\) 均为自伴迹类算子。对自伴迹类算子，由谱定理 \(T_1 = \sum \lambda_n^+ \langle \cdot, e_n^+ \rangle e_n^+ - \sum \lambda_n^- \langle \cdot, e_n^- \rangle e_n^-\)，其中 \(\lambda_n^\pm \geq 0\)。迹 \(\operatorname{Tr}(T_1) = \sum \lambda_n^+ - \sum \lambda_n^- = \sum \lambda_n(T_1)\)。对一般迹类算子，利用 Weyl 不等式 \(\sum_{k=1}^n |\lambda_k(T)| \leq \sum_{k=1}^n s_k(T)\)（对所有 \(n\)），且 \(\sum s_k(T) = \|T\|_1 < \infty\)，故特征值级数绝对收敛。进一步，\(\sum_{k=1}^\infty \lambda_k(T) = \lim_{n \to \infty} \sum_{k=1}^n \lambda_k(T)\) 存在。由迹与基无关的性质，可取基使矩阵表示为上三角形式（Schur 形式），此时对角元恰好为特征值，故迹等于特征值之和。\(\blacksquare\)

上述定义和定理揭示了 \(\mathcal{S}_1(H)\)、\(\mathcal{S}_2(H)\)、\(\mathcal{K}(H)\) 和 \(\mathcal{B}(H)\) 之间的包含层级。这一层级可以类比于序列空间中的 \(\ell^1 \subset \ell^2 \subset c_0 \subset \ell^\infty\)：奇异值可和性越强，算子类越小。特别值得注意的是，\(\mathcal{S}_2(H)\) 本身也是一个 Hilbert 空间，它的内积结构使得 Hilbert-Schmidt 算子在量子场论和随机分析中成为自然的算子空间。下面的命题系统总结了这些包含关系和理想性质。

\textbf{命题 55.7.2}：\(\mathcal{S}_1(H) \subseteq \mathcal{S}_2(H) \subseteq \mathcal{K}(H) \subseteq \mathcal{B}(H)\)，且 \(\mathcal{S}_1(H)\) 和 \(\mathcal{S}_2(H)\) 都是 \(\mathcal{B}(H)\) 中的双边理想。\(\mathcal{S}_2(H)\) 本身是 Hilbert 空间，内积为 \(\langle A, B \rangle_{\mathcal{S}_2} = \operatorname{Tr}(B^* A)\)。

\textbf{证明}：包含关系：若 \(T \in \mathcal{S}_1(H)\)，则 \(\sum s_n < \infty\)，从而 \(\sum s_n^2 \leq (\sum s_n)^2 < \infty\)，故 \(T \in \mathcal{S}_2(H)\)。又 \(\sum s_n^2 < \infty\) 蕴含 \(s_n \to 0\)，由奇异值分解知 \(T\) 为有限秩算子按范数极限，故 \(T \in \mathcal{K}(H)\)。紧算子显然有界，故 \(\mathcal{K}(H) \subseteq \mathcal{B}(H)\)。理想性质：设 \(A \in \mathcal{B}(H)\)，\(T \in \mathcal{S}_1(H)\)。则 \(AT\) 的奇异值 \(s_n(AT) \leq \|A\| s_n(T)\)，故 \(\sum s_n(AT) \leq \|A\| \sum s_n(T) < \infty\)，\(AT \in \mathcal{S}_1(H)\)。同理 \(TA \in \mathcal{S}_1(H)\)。对 \(\mathcal{S}_2(H)\) 同理。Hilbert 空间结构：\(\langle A, B \rangle_{\mathcal{S}_2} = \operatorname{Tr}(B^* A) = \sum_n \langle B^* A e_n, e_n \rangle = \sum_n \langle A e_n, B e_n \rangle\)，此内积诱导的范数 \(\|A\|_{\mathcal{S}_2}^2 = \sum_n \|A e_n\|^2 = \sum_n s_n^2\)。完备性：若 \(\{A_k\}\) 是 \(\mathcal{S}_2\) 中的 Cauchy 列，则对每个 \(n\)，\(\{A_k e_n\}\) 在 \(H\) 中 Cauchy，极限定义了算子 \(A\)。由 Fatou 引理，\(\sum \|A e_n\|^2 \leq \liminf \sum \|A_k e_n\|^2 < \infty\)，故 \(A \in \mathcal{S}_2(H)\)，且 \(\|A_k - A\|_{\mathcal{S}_2} \to 0\)。\(\blacksquare\)

\hrulefill

\subsection{55.8 紧算子的应用------Fredholm 积分方程}\label{ux7d27ux7b97ux5b50ux7684ux5e94ux7528fredholm-ux79efux5206ux65b9ux7a0b}

\textbf{定理 55.8.1（第二类 Fredholm 积分方程）}：考虑积分方程

\[
f(x) - \lambda \int_a^b K(x, y) f(y) \, dy = g(x)
\]

其中 \(K \in L^2([a, b]^2)\)。由 Fredholm 二择一定理（定理 55.4.1），对 \(\lambda \neq 0\)，要么齐次方程 \((f = \lambda T_K f)\) 只有零解，此时非齐次方程对任意 \(g \in L^2([a,b])\) 有唯一解；要么 \(1/\lambda\) 是 \(T_K\) 的特征值，此时非齐次方程有解当且仅当 \(g\) 正交于伴随齐次方程的核空间。

\textbf{证明}：将方程写为 \((I - \lambda T_K) f = g\)，其中 \(T_K\) 为 Fredholm 积分算子。由定理 55.2.1，\(T_K\) 是 \(L^2([a,b])\) 上的紧算子。对 \(\lambda \neq 0\)，令 \(\mu = 1/\lambda\)，则方程化为 \((\mu I - T_K) f = \mu g\)。若 \(\mu\) 不是 \(T_K\) 的特征值，则 \(\mu I - T_K\) 可逆（由 Fredholm 二择一定理），故 \(f = \mu (\mu I - T_K)^{-1} g\) 给出唯一解。若 \(\mu\) 是 \(T_K\) 的特征值，则齐次方程 \((\mu I - T_K) f = 0\) 有非零解。由定理 55.4.1，非齐次方程 \((\mu I - T_K) f = \mu g\) 有解当且仅当 \(\mu g \in \operatorname{ran}(\mu I - T_K)\)。在 Hilbert 空间中，\(\operatorname{ran}(\mu I - T_K) = \ker(\overline{\mu} I - T_K^*)^\perp\)。故 \(g\) 必须正交于 \(\ker(\overline{\mu} I - T_K^*)\) 的所有元素。由于 \(K\) 的伴随核为 \(K^*(x, y) = \overline{K(y, x)}\)，这恰好是伴随齐次方程的解空间。\(\blacksquare\)

上述结论将 Fredholm 积分方程的求解完全归结为紧算子的 Fredholm 二择一定理，从而在算子理论的抽象框架中获得了统一的处理。在此基础上，Fredholm 进一步将有限维线性代数中的行列式理论推广到了无穷维情形------他定义了 Fredholm 行列式 \(\det(I - \lambda T)\) 作为整函数，并发现预解式 \((I - \lambda T)^{-1}\) 恰好是两个广义行列式（分子和分母）的商。这种''无穷维行列式''在可积系统（如 KdV 方程的 Fredholm 行列式表示）和随机矩阵理论中有深远应用。

\textbf{定理 55.8.2（Fredholm 行列式与预解式）}：对迹类算子 \(T\)，可定义 Fredholm 行列式 \(\det(I - \lambda T) = \prod_{n} (1 - \lambda \lambda_n(T))\)（整函数）。预解式 \((I - \lambda T)^{-1}\) 可表示为 \(\det(I - \lambda T)\) 的相关经式（Fredholm 第一类级数）之比。Fredholm 行列式理论是有限维线性代数中 \(\det(I - \lambda A)\) 的无穷维推广。

\textbf{证明}：对 \(T \in \mathcal{S}_1(H)\)，特征值 \(\{\lambda_n(T)\}\) 满足 \(\sum |\lambda_n(T)| < \infty\)（由 Lidskii 定理和 Weyl 不等式）。无穷乘积 \(\det(I - \lambda T) = \prod_{n=1}^\infty (1 - \lambda \lambda_n(T))\) 收敛当且仅当 \(\sum |\lambda \lambda_n(T)| < \infty\)，而 \(\sum |\lambda_n(T)| \leq \|T\|_1 < \infty\)，故该乘积对任意 \(\lambda \in \mathbb{C}\) 绝对收敛，定义了 \(\lambda\) 的整函数。该整函数的零点恰为 \(\lambda = 1/\lambda_n(T)\)（\(T\) 的非零特征值的倒数）。预解式 \((I - \lambda T)^{-1}\) 在 \(\lambda \notin \sigma(T)^{-1}\) 处解析。由 Fredholm 第一类级数展开：\((I - \lambda T)^{-1} = I + \sum_{n=1}^\infty \lambda^n \bigwedge^n T / \det(I - \lambda T)\)，其中 \(\bigwedge^n T\) 为 \(T\) 的 \(n\) 次外幂（在 \(\bigwedge^n H\) 上的迹类算子）。这给出了预解式作为两个整函数之比的明确表达式，完全类似于有限维情形中 \((\lambda I - A)^{-1} = \operatorname{adj}(\lambda I - A) / \det(\lambda I - A)\)。\(\blacksquare\)

\hrulefill

\subsection{55.9 紧算子的 Calkin 代数与 Fredholm 算子的 Atiyah-Singer 指数}\label{ux7d27ux7b97ux5b50ux7684-calkin-ux4ee3ux6570ux4e0e-fredholm-ux7b97ux5b50ux7684-atiyah-singer-ux6307ux6570}

\textbf{定义 55.9.1（Calkin 代数）}：设 \(H\) 为 Hilbert 空间。Calkin 代数定义为 \(\mathcal{C}(H) = \mathcal{B}(H) / \mathcal{K}(H)\)（有界算子模去紧算子的商代数）。Calkin 代数是 \(C^*\)-代数，其元素 \([T]\) 可视为 \(T\) 的''本质部分''------即忽略紧算子扰动后的等价类。

Calkin 代数的引入源于这样一个深刻洞察：紧算子在某种意义上是可以被''忽略''的------它们是无穷维空间中最接近有限秩算子的理想，任何紧算子的扰动不应改变算子的本质谱学性质。将全体有界算子模去紧算子所得到的商代数 \(\mathcal{C}(H)\) 恰好捕捉了算子的''本质''部分。Atkinson 定理是 Calkin 代数理论的第一个伟大胜利：它揭示了 Fredholm 算子的本质特征正是在 Calkin 代数中可逆。这一等价性将分析中的 Fredholm 指标问题转化为代数中的可逆性问题，为 Atiyah-Singer 指标定理铺平了道路。

\textbf{定理 55.9.1（Atkinson 定理）}：\(T \in \mathcal{B}(H)\) 是 Fredholm 算子（\(\ker T\) 有限维，\(\operatorname{ran} T\) 闭且 \(\operatorname{coker} T\) 有限维）当且仅当 \([T]\) 在 Calkin 代数 \(\mathcal{C}(H)\) 中可逆。Fredholm 指标 \(\operatorname{ind}(T) = \dim \ker T - \dim \operatorname{coker} T\) 在紧扰动下不变，且是连续函数（在 Fredholm 算子的范数拓扑下）。

\textbf{证明}：（\(\Rightarrow\)）设 \(T\) 为 Fredholm 算子。则 \(\ker T\) 有限维，\(\operatorname{ran} T\) 闭且余维数有限。令 \(P\) 为到 \(\ker T\) 上的正交投影，\(Q\) 为到 \(\operatorname{ran} T\) 上的正交投影。则 \(T: (\ker T)^\perp \to \operatorname{ran} T\) 为双射（单射由构造，满射由值域限制）。由 Banach 逆算子定理，存在有界逆 \(S_0: \operatorname{ran} T \to (\ker T)^\perp\)。定义 \(S = S_0 Q\)，则 \(ST = I - P\) 且 \(TS = Q\)。由于 \(P\) 和 \(I - Q\) 均为有限秩算子（故紧），在 Calkin 代数中 \([S][T] = [I]\) 且 \([T][S] = [I]\)，故 \([T]\) 可逆。（\(\Leftarrow\)）设 \([T]\) 在 Calkin 代数中可逆，即存在 \(S \in \mathcal{B}(H)\) 和紧算子 \(K_1, K_2\) 使得 \(ST = I + K_1\)，\(TS = I + K_2\)。由 \(ST = I + K_1\) 知 \(\ker T \subseteq \ker(I + K_1)\)，而 \(I + K_1\) 为 Fredholm 算子（\(\ker\) 有限维），故 \(\dim \ker T < \infty\)。由 \(TS = I + K_2\) 知 \(\operatorname{ran} T \supseteq \operatorname{ran}(I + K_2)\)，而 \(I + K_2\) 有闭值域且余维数有限，故 \(\operatorname{ran} T\) 也闭且余维数有限。指标不变性：设 \(K\) 紧，\(t \mapsto T + tK\) 为 Fredholm 算子的连续族，指标在此族上为局部常值函数，故 \(\operatorname{ind}(T + K) = \operatorname{ind}(T)\)。\(\blacksquare\)

Atkinson 定理的核心价值在于它架起了一座桥梁，一端是 Hilbert 空间上的 Fredholm 分析理论，另一端是 \(C^*\)-代数的 \(K\)-理论。Fredholm 指标在紧扰动下的不变性暗示了一个深层的拓扑结构：\(\operatorname{ind}\) 实际上是从 Fredholm 算子空间到整数群的一个连续同伦不变量。这一观察最终催生了 Atiyah-Singer 指标定理（1963）------它把椭圆偏微分算子的 Fredholm 指标与底流形的拓扑示性数（如 Euler 示性数、Hirzebruch 符号差等）等同起来，被誉为 20 世纪数学的最高成就之一。

\textbf{注记 55.9.1}：Atkinson 定理将 Fredholm 算子的分析性质转化为 Calkin 代数中的代数可逆性，这是算子代数和 \(K\)-理论方法的经典应用。Fredholm 指标的不变性导致了 Atiyah-Singer 指标定理------将椭圆微分算子的解析指标（Fredholm 指标）与流形的拓扑不变量（Chern 类、Todd 类）联系起来，是 20 世纪数学最伟大的成就之一。

\subsection{55.10 紧算子的谱映射与泛函演算}\label{ux7d27ux7b97ux5b50ux7684ux8c31ux6620ux5c04ux4e0eux6cdbux51fdux6f14ux7b97}

\textbf{定理 55.10.1（紧算子的谱映射定理）}：设 \(T \in \mathcal{K}(X)\)。若 \(p\) 为多项式且 \(p(0) = 0\)，则 \(\sigma(p(T)) = p(\sigma(T))\)。更一般地，对在 \(0\) 处为零的解析函数 \(f\)（在 \(\sigma(T)\) 的邻域上全纯），\(f(T)\)（由 Riesz 泛函演算定义）也是紧算子，且 \(\sigma(f(T)) = f(\sigma(T))\)。

\textbf{证明}：首先，\(p(0) = 0\) 确保 \(p(T)\) 也是紧算子（因为 \(p(T) = T \cdot q(T)\) 为紧算子与有界算子的复合）。对多项式情形，\(\sigma(p(T)) = p(\sigma(T))\) 是 Gelfand 谱映射定理的直接推论（对任意有界算子成立）。对紧性的验证：若 \(p(z) = \sum_{k=1}^n a_k z^k\)，则 \(p(T) = \sum_{k=1}^n a_k T^k\)。每个 \(T^k\)（\(k \geq 1\)）是紧算子（由定理 55.1.1(iii)），故 \(p(T)\) 是紧算子的线性组合，从而紧。对解析函数 \(f\) 满足 \(f(0) = 0\)，由 Riesz 泛函演算，\(f(T) = \frac{1}{2\pi i} \int_\Gamma f(\zeta)(\zeta I - T)^{-1} d\zeta\)，其中 \(\Gamma\) 为包围 \(\sigma(T)\) 的围道。\(f(T)\) 是紧算子的积分极限（用紧算子的逼近序列），故紧。谱映射 \(\sigma(f(T)) = f(\sigma(T))\) 由一般 Gelfand 理论保证。\(\blacksquare\)

谱映射定理刻画了紧算子通过函数演算后的整体谱行为，但它并未告诉我们如何在单个特征值附近提取算子的局部信息。这就引出了 Riesz 投影的概念------这是一种通过复围道积分（类似于复分析中的留数计算）将算子''投影''到特定特征值的广义特征空间上的技术。Riesz 投影在紧算子谱理论中扮演的角色类似于线性代数中广义特征向量的构造，它是理解紧算子局部谱结构的核心工具。

\textbf{定理 55.10.2（紧算子的 Riesz 投影）}：设 \(\lambda \neq 0\) 为 \(T \in \mathcal{K}(X)\) 的特征值。定义 Riesz 投影 \(P_\lambda = \frac{1}{2\pi i} \int_{\partial D(\lambda, \varepsilon)} (\zeta I - T)^{-1} d\zeta\)（\(\varepsilon\) 足够小使 \(\lambda\) 是圆盘内唯一谱点）。则 \(P_\lambda\) 是有限秩投影算子，\(\operatorname{ran} P_\lambda\) 是 \(T\) 的广义特征空间（代数重数有限），且 \(T\) 在 \(\operatorname{ran} P_\lambda\) 上的限制的谱仅为 \(\{\lambda\}\)。

\textbf{证明}：由紧算子的谱性质，\(\lambda \neq 0\) 是 \(\sigma(T)\) 的孤立点。取 \(\varepsilon > 0\) 足够小使 \(D(\lambda, \varepsilon) \cap \sigma(T) = \{\lambda\}\)。预解式 \((\zeta I - T)^{-1}\) 在 \(\partial D(\lambda, \varepsilon)\) 上解析，故围道积分 \(P_\lambda\) 有定义。\(P_\lambda\) 是投影：\(P_\lambda^2 = \frac{1}{(2\pi i)^2} \int_{\Gamma_1} \int_{\Gamma_2} (\zeta I - T)^{-1} (\eta I - T)^{-1} d\eta d\zeta\)。利用预解恒等式 \((\zeta I - T)^{-1} - (\eta I - T)^{-1} = (\eta - \zeta)(\zeta I - T)^{-1}(\eta I - T)^{-1}\)，可证 \(P_\lambda^2 = P_\lambda\)。\(P_\lambda\) 的有限秩性：\(P_\lambda = \frac{1}{2\pi i} \int_{\partial D} ((\zeta - \lambda)I + (\lambda I - T))^{-1} d\zeta\)。展开 \((\zeta I - T)^{-1} = (\zeta - \lambda)^{-1} I + (\zeta - \lambda)^{-2} (\lambda I - T) + \cdots\)，代入得 \(P_\lambda = I + \sum_{n=1}^\infty c_n (\lambda I - T)^n\)。由于 \(\lambda I - T\) 是 Fredholm 算子，\(P_\lambda\) 的有限秩性由 \(\operatorname{ran} P_\lambda\) 是广义特征空间（有限维）得出。\(T\) 在 \(\operatorname{ran} P_\lambda\) 上的限制：对 \(x \in \operatorname{ran} P_\lambda\)，\(T x = T P_\lambda x = P_\lambda T x \in \operatorname{ran} P_\lambda\)，故该子空间为 \(T\)-不变子空间。\(\sigma(T|_{\operatorname{ran} P_\lambda}) = \{\lambda\}\) 由谱的局部化性质得出。\(\blacksquare\)

Riesz 投影将紧算子限制到有限维的不变子空间上，而在这个有限维子空间内，算子无非是一个特征值为 \(\lambda\) 的线性变换。由有限维线性代数的标准结果，这样的算子可以分解为 \(\lambda I\) 加上一个幂零部分 \(N\)------这就是紧算子的广义 Jordan 分解。这个结果在本质上说：尽管紧算子是作用在无穷维空间上的，但其非零谱部分的局部行为完全由有限维线性代数所描述。这一事实大幅度降低了紧算子谱分析的复杂度。

\textbf{定理 55.10.3（紧算子的 Jordan 标准形的推广）}：设 \(\lambda \neq 0\) 为 \(T \in \mathcal{K}(X)\) 的特征值。则 \(T\) 在 \(\operatorname{ran} P_\lambda\) 上的限制可分解为 \(\lambda I + N\)，其中 \(N\) 为幂零算子（\(N^k = 0\) 对某个 \(k\) 成立）。因此紧算子在其广义特征空间上的行为与有限维 Jordan 标准形完全类似，代数量数 \(\dim \operatorname{ran} P_\lambda\) 为有限值。

\textbf{证明}：设 \(V = \operatorname{ran} P_\lambda\)，\(\dim V = m < \infty\)（由定理 55.10.2）。\(T|_V: V \to V\) 是有限维空间上的线性算子，其谱为 \(\{\lambda\}\)。令 \(N = T|_V - \lambda I_V\)。则 \(N\) 的谱为 \(\{0\}\)（因为 \(\sigma(T|_V) = \{\lambda\}\) 蕴含 \(\sigma(N) = \{0\}\)）。在有限维空间中，谱仅含 \(\{0\}\) 的算子必为幂零算子：由 Cayley-Hamilton 定理，\(N\) 的特征多项式为 \(t^m\)，故 \(N^m = 0\)。因此 \(T|_V = \lambda I_V + N\)，\(N\) 幂零。这给出了 \(T\) 在广义特征空间上的 Jordan 分解。代数重数（即 \(\dim \operatorname{ran} P_\lambda\)）的有限性由定理 55.3.1 直接保证。\(\blacksquare\)

\subsection{55.11 紧算子的逼近性质与 Banach 空间中的基}\label{ux7d27ux7b97ux5b50ux7684ux903cux8fd1ux6027ux8d28ux4e0e-banach-ux7a7aux95f4ux4e2dux7684ux57fa}

\textbf{定义 55.11.1（逼近性质）}：Banach 空间 \(X\) 称为具有\textbf{逼近性质}（approximation property），若对任意紧集 \(K \subseteq X\) 和 \(\varepsilon > 0\)，存在有限秩算子 \(F \in \mathcal{B}(X)\) 使得 \(\|F x - x\| < \varepsilon\)（对所有 \(x \in K\)）。等价地，\(X\) 上的恒等算子可由有限秩算子按紧致开拓扑逼近。

逼近性质询问这样一个基本问题：Banach 空间上的恒等算子能否用有限秩算子（一致地在紧子集上）逼近？如果答案是肯定的，那么从某种意义上说，这个无穷维 Banach 空间的行为与有限维空间相差不大。经典 Banach 空间如 \(\ell^p\)、\(L^p\) 和 \(C[0,1]\) 都具有逼近性质，因为它们拥有自然的 Schauder 基。Grothendieck 在 1955 年证明了这个方向的一般结论------所有具有 Schauder 基的空间都有逼近性质------但他提出的''是否所有可分 Banach 空间都具有逼近性质''这一著名问题，直到 1973 年才被 Enflo 以否定形式解决。

\textbf{定理 55.11.1（Grothendieck 定理）}：所有具有 Schauder 基的 Banach 空间都具有逼近性质。1973 年，Enflo 构造了第一个不具有逼近性质的可分 Banach 空间，解决了 Grothendieck 提出的逼近问题（approximation problem），这是 Banach 空间理论中的里程碑式成果。

\textbf{证明}：设 \(X\) 具有 Schauder 基 \(\{e_n\}_{n=1}^\infty\)，即对任意 \(x \in X\)，\(x = \sum_{n=1}^\infty a_n(x) e_n\)（唯一展开）。定义自然投影 \(P_N: X \to X\) 为 \(P_N(x) = \sum_{n=1}^N a_n(x) e_n\)。由 Schauder 基的定义，\(P_N\) 为有界线性算子，且 \(\sup_N \|P_N\| < \infty\)（一致有界原理）。对任意 \(x \in X\)，\(\|P_N x - x\| \to 0\)（\(N \to \infty\)）。现在对任意紧集 \(K \subseteq X\) 和 \(\varepsilon > 0\)：取有限 \(\varepsilon\)-网 \(\{y_1, \ldots, y_m\} \subseteq K\)（由紧性）。对每个 \(y_j\)，存在 \(N_j\) 使 \(\|P_N y_j - y_j\| < \varepsilon\) 对所有 \(N \geq N_j\)。取 \(N = \max_j N_j\)，则对任意 \(x \in K\)，取 \(y_j\) 使 \(\|x - y_j\| < \varepsilon\)，有 \(\|P_N x - x\| \leq \|P_N x - P_N y_j\| + \|P_N y_j - y_j\| + \|y_j - x\| \leq (\|P_N\| + 1)\varepsilon + \varepsilon = (\|P_N\| + 2)\varepsilon\)。由于 \(\|P_N\|\) 一致有界，\(P_N\) 为有限秩算子，故 \(X\) 具有逼近性质。Enflo 的反例构造了无 Schauder 基且无逼近性质的可分 Banach 空间，证明了并非所有 Banach 空间都具有逼近性质。\(\blacksquare\)

逼近性质虽然是 Banach 空间几何中的高度抽象概念，但紧算子和迹类算子的理论在物理学中有着非常具体的应用。在量子统计力学中，量子态由迹类正算子（密度矩阵）表示，其迹天然对应概率守恒。紧自伴算子的谱定理为密度矩阵的对角化提供了依据，而迹的定义和 Lidskii 定理则确保了 von Neumann 熵等量子信息论基本量的良好定义。这是纯数学理论在物理学中获得精确实现的绝佳例证。

\textbf{定理 55.11.2（紧算子的迹与 Fredholm 行列式在量子统计力学中的应用）}：在量子统计力学中，单粒子密度矩阵 \(\rho\) 是迹类算子，满足 \(\operatorname{Tr}(\rho) = 1\)（归一化），\(\rho \geq 0\)（正定性）。由紧自伴算子的谱定理，\(\rho = \sum \lambda_n |\psi_n\rangle\langle\psi_n|\)（\(\lambda_n \geq 0\)，\(\sum \lambda_n = 1\)）。von Neumann 熵 \(S(\rho) = -\operatorname{Tr}(\rho \log \rho) = -\sum \lambda_n \log \lambda_n\) 度量了量子态的混合程度。紧算子的谱理论为量子力学中的密度矩阵形式提供了严格的数学基础。

\textbf{证明}：\(\rho\) 为迹类正算子，故 \(\rho\) 为紧自伴算子。由定理 55.5.1（Hilbert-Schmidt 谱定理），\(\rho\) 可对角化：存在规范正交基 \(\{|\psi_n\rangle\}\) 使得 \(\rho = \sum_n \lambda_n |\psi_n\rangle\langle\psi_n|\)，其中 \(\lambda_n \geq 0\) 为特征值。由迹的定义和 Lidskii 定理，\(\operatorname{Tr}(\rho) = \sum_n \lambda_n = 1\)（归一化条件）。\(\rho \geq 0\) 等价于 \(\lambda_n \geq 0\)（对所有 \(n\)）。\(\log \rho\) 由泛函演算定义为 \(\log \rho = \sum_n (\log \lambda_n) |\psi_n\rangle\langle\psi_n|\)（约定 \(0 \log 0 = 0\)）。于是 \(\rho \log \rho = \sum_n \lambda_n \log \lambda_n |\psi_n\rangle\langle\psi_n|\)，取其迹得 \(S(\rho) = -\operatorname{Tr}(\rho \log \rho) = -\sum_n \lambda_n \log \lambda_n\)。由 \(\lambda_n \in [0,1]\) 和 \(\sum \lambda_n = 1\)，\(S(\rho) \geq 0\)，且 \(S(\rho) = 0\) 当且仅当 \(\rho\) 为纯态（恰好一个 \(\lambda_n = 1\)，其余为 \(0\)）。\(\blacksquare\)

\newpage

\chapter{06-分析/03-泛函分析/05-C代数与深化.md}\label{ux5206ux679003-ux6cdbux51fdux5206ux679005-cux4ee3ux6570ux4e0eux6df1ux5316.md}

\section{第176章：C*-代数初步}\label{ux7b2c176ux7ae0c-ux4ee3ux6570ux521dux6b65}

C*-代数是 Hilbert 空间上有界算子构成的、在伴随和范数拓扑下封闭的代数，是量子力学和非交换几何的数学语言。

\subsection{76.1 Banach 代数与 C*-代数}\label{banach-ux4ee3ux6570ux4e0e-c-ux4ee3ux6570}

\textbf{定义 76.1}（Banach 代数）：\(\mathbb{C}\) 上的一个 \textbf{Banach 代数}是一个结合 \(\mathbb{C}\)-代数 \(\mathcal{A}\) 配以范数 \(\|\cdot\|\)，使得 \((\mathcal{A}, \|\cdot\|)\) 是 Banach 空间，且满足 \(\|ab\| \leq \|a\| \|b\|\) 对所有 \(a, b \in \mathcal{A}\)。

若 \(\mathcal{A}\) 有单位元 \(1\) 且 \(\|1\| = 1\)，则称 \(\mathcal{A}\) 为\textbf{有单位元的 Banach 代数}。

\textbf{定义 76.2}（对合）：\(\mathcal{A}\) 上的一个\textbf{对合}是一个映射 \(*: \mathcal{A} \to \mathcal{A}\)，\(a \mapsto a^*\)，满足： 1. \((a^*)^* = a\) 2. \((a + b)^* = a^* + b^*\) 3. \((\lambda a)^* = \overline{\lambda} a^*\) 4. \((ab)^* = b^* a^*\)

\textbf{定义 76.3}（C*-代数）：一个 \textbf{C\emph{-代数\textbf{是一个带有对合的 Banach 代数 \(\mathcal{A}\)，满足 }C}-等式}：

\[\|a^* a\| = \|a\|^2, \quad \forall a \in \mathcal{A}\]

\subsection{76.2 交换 C*-代数与 Gelfand 变换}\label{ux4ea4ux6362-c-ux4ee3ux6570ux4e0e-gelfand-ux53d8ux6362}

\textbf{定理 76.1}（Gelfand-Mazur 定理）：若 \(\mathcal{A}\) 是有单位元的 Banach 代数，且每个非零元都可逆（即为除环），则 \(\mathcal{A} \cong \mathbb{C}\)。

\emph{证明}：对 \(a \in \mathcal{A}\)，\(\sigma(a) \neq \emptyset\)（谱非空）。取 \(\lambda \in \sigma(a)\)，则 \(a - \lambda 1\) 不可逆，故 \(a - \lambda 1 = 0\)，即 \(a = \lambda 1\)。∎

\textbf{定义 76.4}（特征与 Gelfand 谱）：设 \(\mathcal{A}\) 是交换 Banach 代数。\(\mathcal{A}\) 上的一个\textbf{特征}是一个非零代数同态 \(h: \mathcal{A} \to \mathbb{C}\)。所有特征构成的集合 \(\Omega(\mathcal{A})\) 称为 \(\mathcal{A}\) 的 \textbf{Gelfand 谱}（或极大理想空间）。

\textbf{命题 76.1}：对交换 Banach 代数 \(\mathcal{A}\)，存在自然双射：

\[\Omega(\mathcal{A}) \longleftrightarrow \{\text{极大理想 } \mathfrak{m} \subseteq \mathcal{A}\}\]

\(h \mapsto \ker h\)。每个特征自动连续且 \(\|h\| = 1\)（若 \(\mathcal{A}\) 有单位元）。

\textbf{定义 76.5}（Gelfand 变换）：对 \(a \in \mathcal{A}\)，定义 \(\hat{a}: \Omega(\mathcal{A}) \to \mathbb{C}\) 为 \(\hat{a}(h) = h(a)\)。映射 \(a \mapsto \hat{a}\) 称为 \textbf{Gelfand 变换}，记作 \(\Gamma: \mathcal{A} \to C(\Omega(\mathcal{A}))\)。

\(\Omega(\mathcal{A})\) 赋予弱* 拓扑（作为 \(\mathcal{A}^*\) 的子空间），它是紧致 Hausdorff 空间（若 \(\mathcal{A}\) 有单位元；否则为局部紧致）。Gelfand 变换是代数同态。

\textbf{定理 76.2}（Gelfand-Naimark 定理，交换版本）：设 \(\mathcal{A}\) 是交换 C\emph{-代数（有单位元）。则 Gelfand 变换 \(\Gamma: \mathcal{A} \to C(\Omega(\mathcal{A}))\) 是等距 }-同构。即

\[\|\hat{a}\|_\infty = \|a\|, \quad \widehat{a^*} = \overline{\hat{a}}, \quad \Gamma(\mathcal{A}) = C(\Omega(\mathcal{A}))\]

\textbf{证明}：对 \(a \in \mathcal{A}\)，\(\|\hat{a}\|_\infty = \sup_{h \in \Omega} |h(a)| = r(a)\)（谱半径）。在 C\emph{-代数中，若 \(a\) 是自伴的（\(a = a^*\)），则 \(\|a^2\| = \|a\|^2\)，归纳得 \(\|a^{2^n}\| = \|a\|^{2^n}\)，故 \(r(a) = \lim \|a^n\|^{1/n} = \|a\|\)。对一般 \(a\)，\(\|a\|^2 = \|a^*a\| = r(a^*a) = \|\widehat{a^*a}\|_\infty = \|\overline{\hat{a}}\hat{a}\|_\infty = \|\hat{a}\|_\infty^2\)。满射性：\(\Gamma(\mathcal{A})\) 是 \(C(\Omega)\) 的闭 }-子代数，含常数，分离点（因 \(h_1\neq h_2\) 时存在 \(a\) 使 \(h_1(a)\neq h_2(a)\)），由 Stone-Weierstrass 定理知 \(\Gamma(\mathcal{A}) = C(\Omega(\mathcal{A}))\)。故 Gelfand 变换是等距 *-同构。\(\blacksquare\)

\textbf{推论 76.3}：交换 C*-代数（有单位元）范畴对偶于紧致 Hausdorff 空间范畴。这是 Gelfand-Naimark 对偶。

\subsection{76.3 一般 C*-代数的表示}\label{ux4e00ux822c-c-ux4ee3ux6570ux7684ux8868ux793a}

\textbf{定理 76.4}（Gelfand-Naimark-Segal 构造 / GNS 构造）：设 \(\mathcal{A}\) 是 C*-代数，\(\varphi\) 是 \(\mathcal{A}\) 上的\textbf{正线性泛函}（即 \(\varphi(a^*a) \geq 0\) 对所有 \(a\)）。则存在 Hilbert 空间 \(H_\varphi\)、\(*\)-同态 \(\pi_\varphi: \mathcal{A} \to \mathcal{B}(H_\varphi)\) 和循环向量 \(\xi_\varphi \in H_\varphi\) 使得

\[\varphi(a) = \langle \pi_\varphi(a) \xi_\varphi, \xi_\varphi \rangle, \quad \forall a \in \mathcal{A}\]

\emph{证明构造}：在 \(\mathcal{A}\) 上定义内积 \(\langle a, b \rangle_\varphi = \varphi(b^* a)\)。这是半正定的。令 \(N_\varphi = \{a \in \mathcal{A} : \varphi(a^*a) = 0\}\)，则 \(N_\varphi\) 是左理想。\(H_\varphi\) 是 \(\mathcal{A}/N_\varphi\) 在该内积下的完备化。\(\pi_\varphi(a)(b + N_\varphi) = ab + N_\varphi\)。\(\xi_\varphi = 1 + N_\varphi\)。∎

\textbf{定理 76.5}（Gelfand-Naimark 定理，非交换版本）：每个 C\emph{-代数 \(\mathcal{A}\) 等距 }-同构于某个 Hilbert 空间 \(\mathcal{B}(H)\) 的 C*-子代数。

\textbf{证明}：考虑 \(\mathcal{A}\) 上所有态的集合 \(\mathcal{S}(\mathcal{A})\)。对每个态 \(\varphi \in \mathcal{S}(\mathcal{A})\)，由 GNS 构造得 Hilbert 空间 \(H_\varphi\)、\(*\)-同态 \(\pi_\varphi: \mathcal{A} \to \mathcal{B}(H_\varphi)\) 和循环向量 \(\xi_\varphi\)。令 \(H = \bigoplus_{\varphi \in \mathcal{S}(\mathcal{A})} H_\varphi\) 为所有 \(H_\varphi\) 的 Hilbert 直和，定义 \(\pi = \bigoplus_{\varphi} \pi_\varphi: \mathcal{A} \to \mathcal{B}(H)\)，即 \(\pi(a)((v_\varphi)_\varphi) = (\pi_\varphi(a)v_\varphi)_\varphi\)。\(\pi\) 为 \(*\)-同态。对任意 \(a \neq 0\)，\(a^*a \neq 0\)（因 \(\mathcal{A}\) 为 C\emph{-代数），由正线性泛函的 Hahn-Banach 分离定理，存在态 \(\varphi\) 使 \(\varphi(a^*a) > 0\)。则 \(\|\pi_\varphi(a)\xi_\varphi\|^2 = \varphi(a^*a) > 0\)，故 \(\pi(a) \neq 0\)，\(\pi\) 为单射。由 C}-代数中单射 \(*\)-同态自动等距（\(\|\pi(a)\| = \|a\|\)），\(\pi\) 为等距 \(*\)-同构到其像 \(\pi(\mathcal{A})\)，即 \(\mathcal{A} \cong \pi(\mathcal{A}) \subseteq \mathcal{B}(H)\)。\(\blacksquare\)

\subsection{76.4 正元与序结构}\label{ux6b63ux5143ux4e0eux5e8fux7ed3ux6784}

\textbf{定义 76.6}（正元）：C*-代数 \(\mathcal{A}\) 中的元素 \(a\) 称为\textbf{正的}（记作 \(a \geq 0\)），如果 \(a = a^*\) 且 \(\sigma(a) \subseteq [0, \infty)\)。所有正元构成的集合记作 \(\mathcal{A}^+\)。

\textbf{命题 76.2}：\(a \geq 0\) 当且仅当存在 \(b \in \mathcal{A}\) 使得 \(a = b^*b\)。\(\mathcal{A}^+\) 是 \(\mathcal{A}\) 中的闭凸锥且 \(\mathcal{A}^+ \cap (-\mathcal{A}^+) = \{0\}\)。

\textbf{定义 76.7}（态）：C\emph{-代数 \(\mathcal{A}\) 上的一个\textbf{态}是一个正线性泛函 \(\varphi\) 满足 \(\|\varphi\| = 1\)（若 \(\mathcal{A}\) 有单位元，等价于 \(\varphi(1) = 1\)）。所有态构成的集合 \(\mathcal{S}(\mathcal{A})\) 是 \(\mathcal{A}^*\) 中弱} 紧致的凸集。

\textbf{定理 76.6}（Krein-Milman 定理在态空间中的应用）：\(\mathcal{S}(\mathcal{A})\) 的极点称为\textbf{纯态}。GNS 构造中，纯态对应不可约表示。

\textbf{证明}：由 Krein-Milman 定理，紧凸集 \(\mathcal{S}(\mathcal{A})\)（弱*紧致）是其端点集的闭凸包。若态 \(\varphi\) 是纯态，假设 GNS 表示 \(\pi_\varphi\) 可约：存在非平凡投影 \(P \in \pi_\varphi(\mathcal{A})'\)，则 \(\varphi = \lambda \varphi_1 + (1-\lambda)\varphi_2\)（\(\varphi_i\) 为态），与 \(\varphi\) 为端点矛盾。反之，若 \(\pi_\varphi\) 不可约，由 Schur 引理 \(\pi_\varphi(\mathcal{A})' = \mathbb{C}I\)，若 \(\varphi = t\varphi_1 + (1-t)\varphi_2\)（\(0<t<1\)），则 \(\varphi_i \leq t^{-1}\varphi\)，由 Radon-Nikodym 型定理存在 \(T \in \pi_\varphi(\mathcal{A})'\) 使 \(\varphi_1(a) = \langle \pi_\varphi(a)T\xi_\varphi, \xi_\varphi \rangle\)，而 \(\pi_\varphi(\mathcal{A})' = \mathbb{C}I\) 推出 \(\varphi_1 = \varphi_2 = \varphi\)，故 \(\varphi\) 为端点。\(\blacksquare\)

\subsection{76.5 量子力学诠释}\label{ux91cfux5b50ux529bux5b66ux8be0ux91ca}

C\emph{-代数为量子力学提供了严格的数学框架： - 物理系统的\textbf{可观测量}对应 C}-代数中的自伴元。 - 系统的\textbf{态}对应 C\emph{-代数上的态（正线性泛函）。 - 对态 \(\varphi\) 和可观测量 \(a = a^*\)，\(\varphi(a)\) 是测量 \(a\) 的期望值。 - 交换 C}-代数对应经典力学（可观测量可同时测量），非交换 C*-代数对应量子力学（存在不确定性原理）。

\hrulefill

\hrulefill

\hrulefill

\hrulefill

\hrulefill

\hrulefill

\section{第177章：有序线性空间（Ordered Vector Spaces）}\label{ux7b2c177ux7ae0ux6709ux5e8fux7ebfux6027ux7a7aux95f4ordered-vector-spaces}

有序线性空间是泛函分析中处理不等式、正性和序结构的框架。定义：实向量空间 \(V\) 配以凸锥 \(V_+ \subset V\)（正锥）满足 \(V_+ \cap (-V_+) = \{0\}\)，诱导偏序 \(x \leq y \iff y - x \in V_+\)。\textbf{Riesz 空间}（向量格）：若 \(V\) 配以上述偏序后每对有上确界（和下确界），则 \(V\) 为 Riesz 空间------算子理论的序基础。\textbf{核心定理}：（1）Krein-Rutman 定理------Banach 格上的紧正算子有正特征值及对应的正特征向量（Perron-Frobenius 定理的无穷维推广），在 PDE 和 Markov 过程的基本解估计中至关重要；（2）有序线性空间上的正线性泛函扩张定理（Krein 延拓定理）；（3）Banach 格是等距序同构于某个 \(L^p(\mu)\) 空间的 \(L^p\)-空间理论。在数理经济学和优化理论中，有序线性空间为一般均衡理论和凸对偶理论提供了基础框架。

\textbf{定义 77.1}（Riesz 空间 / 向量格）：实向量空间 \(V\) 配以偏序 \(\leq\) 满足：（1）\(x \leq y \implies x+z \leq y+z\)；（2）\(x \geq 0 \implies \alpha x \geq 0\)（\(\forall \alpha \geq 0\)）；（3）任意两个元 \(x, y\) 有上确界 \(x \vee y\) 和下确界 \(x \wedge y\)。称 \((V, \leq)\) 为 \textbf{Riesz 空间} 或向量格。Banach 空间兼为 Riesz 空间且满足 \(\|x\| \leq \|y\|\) 当 \(|x| \leq |y|\)（其中 \(|x| = x \vee (-x)\)）时称为 \textbf{Banach 格}。

\textbf{定理 77.1}（Krein-Rutman 定理）：设 \(X\) 为 Banach 格，\(T: X \to X\) 为紧正算子且谱半径 \(r(T) > 0\)。则 \(r(T)\) 是 \(T\) 的特征值，对应特征向量可在 \(X_+\)（正锥）中选取。

\textbf{证明}：设 \(r = r(T) > 0\)。由谱半径公式，存在 \(\lambda_n \in \sigma(T)\) 使 \(|\lambda_n| \to r\)。利用 Pringsheim 型论证，可选取子列使 \(\lambda_n \to r\)。由于 \(T\) 是正算子，其预解式 \((zI - T)^{-1}\) 当 \(z > r\) 时为正算子。取 \(z_m \downarrow r\)，构造向量 \(x_m = (z_m I - T)^{-1} y\)（\(y \in X_+\) 固定），由紧性，\(T x_m\) 有收敛子列。极限过程给出 \(T x = r x\) 且 \(x \in X_+ \setminus \{0\}\)。这是 Perron-Frobenius 定理在无穷维情形的推广。\(\blacksquare\)

\hrulefill

\hrulefill

\hrulefill

\hrulefill

\hrulefill

\section{第178章：插值空间（Interpolation Spaces）}\label{ux7b2c178ux7ae0ux63d2ux503cux7a7aux95f4interpolation-spaces}

插值空间理论（Aronszajn-Gagliardo, 1960年代；Peetre, Lions）在 Banach 空间对之间构造中间空间，调和分析中精细的 Sobolev 嵌入和算子有界性证明依赖于此。\textbf{核心方法}：（1）\textbf{实插值（\(K\)-方法）}------设 \((X_0, X_1)\) 为相容 Banach 空间对，对 \(x \in X_0 + X_1\) 定义 \(K(t, x) = \inf\{\|x_0\|_{X_0} + t\|x_1\|_{X_1} : x = x_0 + x_1\}\)，则 \((X_0, X_1)_{\theta,q}\) 由 \(\|x\| = \left(\int_0^\infty (t^{-\theta} K(t,x))^q \frac{dt}{t}\right)^{1/q}\) 定义；（2）\textbf{复插值（Calderón 方法）}------利用解析函数在全纯带状区域边值的范数插值。\textbf{应用}：（1）Riesz-Thorin 定理的抽象化------算子 \(T: X_0 \to Y_0\) 和 \(T: X_1 \to Y_1\) 有界，则在插值空间上也有界；（2）Marcinkiewicz 插值定理------从弱 \(L^p\) 到强 \(L^p\) 的插值；（3）Sobolev 空间的分数阶推广（Besov 空间和 Triebel-Lizorkin 空间）。

\textbf{定义 78.1}（\(K\)-泛函与实插值空间）：设 \((X_0, X_1)\) 为相容 Banach 空间对（即两者连续嵌入同一 Hausdorff 拓扑向量空间）。对 \(x \in X_0 + X_1\) 定义 \textbf{\(K\)-泛函} \(K(t, x) = \inf\{\|x_0\|_{X_0} + t\|x_1\|_{X_1} : x = x_0 + x_1\}\)（\(t > 0\)）。对 \(0 < \theta < 1\) 和 \(1 \leq q \leq \infty\)，\textbf{实插值空间} \((X_0, X_1)_{\theta,q} = \{x \in X_0 + X_1 : \|x\|_{\theta,q} = (\int_0^\infty [t^{-\theta} K(t,x)]^q \frac{dt}{t})^{1/q} < \infty\}\)。

\textbf{定理 78.1}（插值定理）：设 \(T: X_0 + X_1 \to Y_0 + Y_1\) 为线性算子且 \(T|_{X_j}: X_j \to Y_j\) 有界（\(j = 0, 1\)）。则对任意 \(\theta \in (0,1)\) 和 \(q \in [1, \infty]\)，\(T: (X_0, X_1)_{\theta,q} \to (Y_0, Y_1)_{\theta,q}\) 有界。

\textbf{证明}：设 \(M_j = \|T\|_{X_j \to Y_j}\)（\(j=0,1\)）。对任意分解 \(x = x_0 + x_1\)（\(x_j \in X_j\)），\(Tx = Tx_0 + Tx_1\) 且 \(\|Tx_0\|_{Y_0} \leq M_0 \|x_0\|_{X_0}\)，\(\|Tx_1\|_{Y_1} \leq M_1 \|x_1\|_{X_1}\)。由 \(K\)-泛函定义，\(K_Y(t, Tx) \leq \|Tx_0\|_{Y_0} + t\|Tx_1\|_{Y_1} \leq M_0\|x_0\|_{X_0} + t M_1\|x_1\|_{X_1}\)。取 \(\inf\) 于 \(x = x_0 + x_1\) 得 \(K_Y(t, Tx) \leq M_0 K_X(\frac{M_1}{M_0} t, x)\)。代入实插值范数，变量代换 \(s = (M_1/M_0) t\)，得 \(\|Tx\|_{(Y_0,Y_1)_{\theta,q}} \leq M_0^{1-\theta} M_1^\theta \|x\|_{(X_0,X_1)_{\theta,q}}\)。该定理统一了 Riesz-Thorin 定理（复插值）和 Marcinkiewicz 插值定理（实插值）的抽象框架。\(\blacksquare\)

\hrulefill

\hrulefill

\hrulefill

\hrulefill

\hrulefill

\section{第179章：线性算子摄动理论（Operator Perturbation Theory）}\label{ux7b2c179ux7ae0ux7ebfux6027ux7b97ux5b50ux6444ux52a8ux7406ux8bbaoperator-perturbation-theory}

线性算子摄动理论研究当算子受微小扰动后谱、预解式和不变子空间的变化，是量子力学（摄动展开）和数值分析的核心工具。\textbf{核心内容}：（1）\textbf{Kato-Rellich 定理}------若 \(A\) 自伴，\(B\) 对称且关于 \(A\) 相对有界（\(\|Bx\| \leq a\|x\| + b\|Ax\|\)，\(b < 1\)），则 \(A + B\) 自伴且保持 \(A\) 的本质谱；（2）\textbf{预解式级数}------对充分小的 \(\lambda\)，\((A + \lambda B - z)^{-1} = R(z) \sum_{n=0}^{\infty} (-\lambda B R(z))^n\) 给出摄动展开；（3）\textbf{解析摄动理论（Kato-Rellich）}------孤立特征值作为摄动参数的解析函数（分支不分裂条件）；（4）\textbf{奇异摄动理论}------当''微小参数''乘在最高阶导数上时（如 \(\epsilon u'' + u' + u = 0\)），边界层现象要求匹配渐近展开（Prandtl 边界层理论的基础）；（5）\textbf{Fredholm 行列式和限界式}------用于迹类算子的谱摄动精确定量分析（如半经典分析框架）。在量子力学的 Rayleigh-Schrödinger 摄动展开、散射理论和 Birkhoff 标准形理论中有核心应用。

\subsection{Kato-Rellich定理}\label{kato-rellichux5b9aux7406}

\textbf{定义142.1}（相对有界性）：设 \(A, B\) 为 Hilbert 空间 \(H\) 上的线性算子，\(D(A) \subseteq D(B)\)。称 \(B\) 关于 \(A\) \textbf{相对有界}，若存在常数 \(a, b \geq 0\) 使 \[\|Bx\| \leq a\|x\| + b\|Ax\|, \quad \forall x \in D(A)\] \(b\) 的下确界称为 \(B\) 关于 \(A\) 的\textbf{相对界}（relative bound）。当 \(b < 1\) 时称 \(B\) 为 \(A\)-\textbf{有界}的（在 Kato 意义下）。

\textbf{定理142.1}（Kato-Rellich定理，1939-1951）：设 \(A\) 为 \(H\) 上（下）自伴算子，\(B\) 为对称算子且 \(D(A) \subseteq D(B)\)，且 \(B\) 关于 \(A\) 相对有界，相对界 \(b < 1\)。则 \(A + B\)（定义域 \(D(A+B) = D(A)\)）是自伴算子。若 \(A\) 本质自伴于某个核心 \(D\)，且 \(D \subseteq D(B)\)，则 \(A + B\) 本质自伴于 \(D\)。

\textbf{证明}：核心步骤是利用预解式展开证明 \(A+B\) 的预解式的存在性。对充分大的 \(\mu > 0\)，考虑算子恒等式 \[(A+B+\mu i)^{-1} = (A+\mu i)^{-1} [I + B(A+\mu i)^{-1}]^{-1}\] 只需证 \(I + B(A+\mu i)^{-1}\) 可逆。由相对有界性，对 \(x \in D(A)\) 及 \(y \in H\)，取 \(x = (A+\mu i)^{-1}y\)，计算 \(\|B(A+\mu i)^{-1}y\| \leq a\|(A+\mu i)^{-1}y\| + b\|A(A+\mu i)^{-1}y\|\)。由自伴算子的谱表示，\(\|(A+\mu i)^{-1}\| \leq 1/|\mu|\)，\(\|A(A+\mu i)^{-1}\| \leq 1\)。故 \(\|B(A+\mu i)^{-1}\| \leq a/|\mu| + b\)。取 \(\mu\) 充分大使 \(a/|\mu| + b < 1\)，则 \(\|B(A+\mu i)^{-1}\| < 1\)，Neumann 级数 \(\sum_{n=0}^\infty (-B(A+\mu i)^{-1})^n\) 收敛，得 \(I + B(A+\mu i)^{-1}\) 可逆。因此 \((A+B+\mu i)\) 有有界逆，定义域 \(\operatorname{Ran}(A+B+\mu i)^{-1} = H\)，由自伴算子的基本判别准则，\(A+B\) 自伴。\(\blacksquare\)

Kato-Rellich 定理的条件 \(b < 1\) 是精确的：存在反例（Wust, 1973）表明 \(b = 1\) 时 \(A+B\) 不必自伴甚至不必闭。

\subsection{Weyl不等式与本质谱}\label{weylux4e0dux7b49ux5f0fux4e0eux672cux8d28ux8c31}

\textbf{定义142.2}（本质谱）：自伴算子 \(A\) 的\textbf{本质谱} \(\sigma_{\text{ess}}(A)\) 是谱中除去有限重数的孤立特征值后的部分。等价定义：\(\lambda \in \sigma_{\text{ess}}(A)\) 当且仅当存在 Weyl 序列 \(\{x_n\} \subseteq D(A)\) 满足 \(\|x_n\| = 1\)，\(x_n \rightharpoonup 0\)（弱收敛），且 \(\|(A-\lambda)x_n\| \to 0\)。

\textbf{定理142.2}（Weyl不等式，1909/1950）：设 \(A\) 为下半有界自伴算子，其离散特征值（计重数）排为 \(\lambda_1(A) \leq \lambda_2(A) \leq \cdots \leq \inf \sigma_{\text{ess}}(A)\)。若 \(B\) 为对称算子且相对 \(A\) 紧（即 \(B(A+i)^{-1}\) 为紧算子），则 \[\lambda_n(A+B) \leq \lambda_n(A) + \|B\|, \quad \lambda_n(A+B) \geq \lambda_n(A) - \|B\|\] 对 \(A\) 和 \(A+B\) 均在 \(\sigma_{\text{ess}}\) 以内的所有 \(n\) 成立。

\textbf{证明}：由极小-极大原理（Courant-Fischer），\(\lambda_n(A) = \inf_{V_n} \sup_{x \in V_n, \|x\|=1} \langle Ax, x \rangle\)，其中 \(V_n\) 取遍 \(D(A)\) 的 \(n\) 维子空间。由 \(\langle (A+B)x, x \rangle = \langle Ax, x \rangle + \langle Bx, x \rangle \leq \langle Ax, x \rangle + \|B\| \|x\|^2\)，代入极小-极大公式即得上界。下界同理，将 \(A = (A+B) - B\) 代入。\(\blacksquare\)

\textbf{定理142.3}（Weyl本质谱定理）：若 \(B\) 关于 \(A\) 相对紧（即 \(B(A+i)^{-1}\) 为紧算子且 \(A\) 自伴、\(B\) 对称、\(D(A) \subseteq D(B)\)），则 \(A\) 与 \(A+B\) 具有相同的本质谱：\(\sigma_{\text{ess}}(A+B) = \sigma_{\text{ess}}(A)\)。特别地，当 \(A\) 的谱在 \(\sigma_{\text{ess}}\) 以下仅有离散特征值时，\(A+B\) 也具有相同性质。

\textbf{证明}：设 \(\lambda \in \sigma_{\text{ess}}(A)\)，取 Weyl 序列 \(\{x_n\} \subseteq D(A)\) 满足 \(\|x_n\| = 1\)，\(x_n \rightharpoonup 0\)，\(\|(A-\lambda)x_n\| \to 0\)。考虑 \(\|(A+B-\lambda)x_n\| \leq \|(A-\lambda)x_n\| + \|B x_n\|\)。由于 \(B(A+i)^{-1}\) 紧且 \((A+i)^{-1}\) 将弱收敛映为强收敛，可证 \(\|B x_n\| = \|B(A+i)^{-1}(A+i)x_n\| \to 0\)（因为 \((A+i)x_n\) 弱收敛）。故 \(\{x_n\}\) 也是 \(A+B\) 在 \(\lambda\) 处的 Weyl 序列，\(\lambda \in \sigma_{\text{ess}}(A+B)\)。反方向由对称性得证。\(\blacksquare\)

\subsection{解析扰动理论与Friedrichs方法}\label{ux89e3ux6790ux6270ux52a8ux7406ux8bbaux4e0efriedrichsux65b9ux6cd5}

\textbf{定理142.4}（Kato-Rellich解析扰动定理）：设 \(A\) 为自伴算子，\(\lambda_0\) 为 \(A\) 的孤立非退化特征值（重数为 1），\(B\) 为对称算子且关于 \(A\) 相对有界。则对充分小的 \(|\kappa|\)，扰动算子 \(A(\kappa) = A + \kappa B\) 在 \(\lambda_0\) 附近有唯一的孤立特征值 \(\lambda(\kappa)\)，且 \(\lambda(\kappa)\) 和相应的规范特征向量 \(\psi(\kappa)\) 都是 \(\kappa\) 的解析函数： \[\lambda(\kappa) = \lambda_0 + \kappa \lambda_1 + \kappa^2 \lambda_2 + \cdots, \quad \psi(\kappa) = \psi_0 + \kappa \psi_1 + \kappa^2 \psi_2 + \cdots\] 其中 \(\lambda_1 = \langle B\psi_0, \psi_0 \rangle\)（一阶 Hellmann-Feynman 定理），\(\lambda_2 = \sum_{k \neq 0} \frac{|\langle B\psi_0, \psi_k \rangle|^2}{\lambda_0 - \lambda_k}\)（二阶 Rayleigh-Schrödinger 公式，\(\{\psi_k\}\) 为 \(A\) 的完全规范正交特征向量系）。

\textbf{证明}：考虑算子值函数 \(T(\kappa, z) = A + \kappa B - z\)，对固定的 \(z = \lambda_0\)，\(T(0, \lambda_0)\) 有一维零空间。将 \(T(\kappa, \lambda)\) 限制到 \(\ker T(0, \lambda_0)\) 及其正交补上，由隐函数定理（对 \(\kappa\) 解析）在 \(\kappa = 0\) 附近解出 \(\lambda(\kappa)\) 的解析支。在有限维情形（Rayleigh-Schrödinger 摄动）中直接验证级数展开；在无穷维情形中需要控制预解式的范数以证明级数的一致收敛性（Kato 使用预解式的二重复合级数展开实现）。相对有界条件 \(b < 1\) 确保 \(B(A+i)^{-1}\) 有界，从而预解式级数在充分小的 \(|\kappa|\) 下收敛。\(\blacksquare\)

\textbf{Friedrichs方法}（Friedrichs, 1948）提供了一种处理''小分母''问题（即 \(\lambda_0\) 嵌入连续谱的情形）的技术。基本思想是：寻找一个适当的逆变换（dressing transformation）\(U(\kappa)\) 使 \(U(\kappa) (A+\kappa B) U(\kappa)^{-1}\) 具有可分解的结构。形式上，将扰动 \(B\) 分解为 \(B = B_{PP} + B_{PC} + B_{CP} + B_{CC}\)（P 为点谱、C 为连续谱部分），通过求解 Fredholm 型积分方程消除交叉项 \(B_{PC}\) 和 \(B_{CP}\)，使变换后的算子对角化。当重整化后的点谱特征值变为复数时，它从连续谱中消失------这一现象在物理中称为共振态（resonance），是 Friedrichs 模型对量子不稳定性的数学刻画。Friedrichs 变换的具体形式为 \((U(\kappa)f)(\lambda) = f(\lambda) + \kappa \int K(\lambda, \omega) f(\omega) d\omega\)，其中积分核 \(K\) 满足适当的奇性增长条件。解析扰动理论与 Friedrichs 方法共同构成了量子力学中稳态微扰理论和共振散射理论（Fermi 黄金规则）的严格数学基础。

\hrulefill

\hrulefill

\hrulefill

\hrulefill

\hrulefill

\hrulefill

\section{第180章：Banach 代数与调和分析}\label{ux7b2c180ux7ae0banach-ux4ee3ux6570ux4e0eux8c03ux548cux5206ux6790}

Banach 代数框架将泛函分析与调和分析统一起来。

\textbf{定义 80.1}（Banach 代数）：Banach 代数 \(A\) 是 Banach 空间配以结合的乘法运算满足 \(\|ab\| \leq \|a\| \|b\|\)。

\textbf{定理 80.1}（Gelfand 变换）：设 \(A\) 为交换 Banach 代数。Gelfand 谱空间 \(\Phi_A\)（\(A\) 上所有非零乘性线性泛函的集合，配以弱-*拓扑）是局部紧 Hausdorff 空间。\textbf{Gelfand 变换} \(\widehat{a}(\varphi) = \varphi(a)\) 是 \(A \to C_0(\Phi_A)\) 的压缩同态。若 \(A\) 含幺元，\(\Phi_A\) 紧致。

\textbf{证明}：每个 \(\varphi \in \Phi_A\) 满足 \(\|\varphi\| \leq 1\)，故 \(\Phi_A \subseteq \operatorname{Ball}(A^*)\)。由 Banach-Alaoglu 定理，\(\operatorname{Ball}(A^*)\) 弱\emph{紧致。验证 \(\Phi_A \cup \{0\}\) 在弱}拓扑下闭（利用乘性条件的连续性），即得 \(\Phi_A\) 为局部紧 Hausdorff。Gelfand 变换 \(\widehat{a}(\varphi) = \varphi(a)\) 是 \(\Phi_A\) 上的连续函数且 \(\|\widehat{a}\|_\infty = \sup_{\varphi} |\varphi(a)| \leq \|a\|\)，故为压缩同态。若 \(A\) 含幺元，由 \(\{0\}\) 与 \(\Phi_A\) 的分离性知 \(\Phi_A\) 紧致。\(\blacksquare\)

\textbf{定理 80.2}（群代数 \(L^1(G)\) 与乘性线性泛函）：设 \(G\) 为局部紧交换群。则 \(L^1(G)\)（卷积代数）的 Gelfand 谱空间自然同胚于对偶群 \(\widehat{G}\)。Gelfand 变换退化为 Fourier 变换：\(\widehat{f}(\chi) = \int_G f(g) \overline{\chi(g)} dg\)。\textbf{证明}：\(L^1(G)\) 的每个乘性线性泛函 \(\varphi\) 满足 \(\varphi(f*g) = \varphi(f)\varphi(g)\)。由 Riesz 表示定理，存在 \(h \in L^\infty(G)\) 使 \(\varphi(f) = \int_G f(g)h(g)dg\)。乘性条件等价于 \(h(g_1+g_2) = h(g_1)h(g_2)\) 几乎处处，即 \(h\) 是 \(G\) 上的特征标。故 \(\Phi_{L^1(G)} \cong \widehat{G}\)（对偶群）。Gelfand 变换 \(\widehat{f}(\chi) = \int_G f(g)\overline{\chi(g)}dg\) 恰为 Fourier 变换。由此建立了 Banach 代数谱理论与局部紧交换群上调和分析的等价性。\(G=\mathbb{R}\) 时 \(\widehat{G}\cong\mathbb{R}\) 给出经典 Fourier 变换；\(G=\mathbb{T}\) 时 \(\widehat{G}\cong\mathbb{Z}\) 给出 Fourier 级数；\(G=\mathbb{Z}\) 时 \(\widehat{G}\cong\mathbb{T}\) 给出离散 Fourier 分析。\(\blacksquare\)

该观点统一了经典 Fourier 分析（\(G = \mathbb{R}\)）、Fourier 级数（\(G = \mathbb{T}\)）和离散 Fourier 分析（\(G = \mathbb{Z}\)）：三者分别对应于 \(L^1(\mathbb{R})\)、\(\ell^1(\mathbb{Z})\)（卷积）和 \(L^1(\mathbb{T})\) 的 Gelfand 变换。

\hrulefill

\hrulefill

\hrulefill

\hrulefill

\hrulefill

\hrulefill

\subsection{122.A C*代数补充}\label{a-cux4ee3ux6570ux8865ux5145}

\begin{quote}
von Neumann 代数（\(W^*\)-代数）是算子代数的核心理论，由 John von Neumann 在 1929 年引入。本卷在 C*-代数（V14 泛函分析）和 Banach 代数（V31 上半部分）的基础上，建立 von Neumann 代数的基本理论：因子分类（I、II\(_1\)、II\(_\infty\)、III 型）、Tomita-Takesaki 模理论和 III 型因子的 Connes 分类。von Neumann 代数在量子场论和自由概率论中有深刻应用。
\end{quote}

\begin{quote}
\textbf{前置知识}：本子卷核心为算子代数理论，需先掌握本卷第128-134章中的Banach空间、Hilbert空间和C*-代数等核心内容，以及第55章（线性代数）中关于内积空间和谱定理的部分。
\end{quote}

\hrulefill

\hrulefill

\hrulefill

\hrulefill

\hrulefill

\hrulefill

\section{第181章：von Neumann 代数基础}\label{ux7b2c181ux7ae0von-neumann-ux4ee3ux6570ux57faux7840}

von Neumann 代数由 John von Neumann 于 1929 年引入，最初称为''算子环''（rings of operators），是量子力学的数学基础的一部分。本章建立 von Neumann 代数的基础理论，包括定义、交换子定理和 von Neumann 双换位定理。

\subsection{213.1 弱算子拓扑与 von Neumann 代数}\label{ux5f31ux7b97ux5b50ux62d3ux6251ux4e0e-von-neumann-ux4ee3ux6570}

\textbf{定义 213.1}（弱算子拓扑 / Weak Operator Topology, WOT）：\(B(H)\) 上的\textbf{弱算子拓扑}是使所有形如 \(T \mapsto \langle Tx, y \rangle\)（\(x, y \in H\)）的泛函连续的最弱拓扑。具体地，网 \(T_\alpha\) 弱收敛于 \(T\)（记作 \(T_\alpha \xrightarrow{WOT} T\)）当且仅当 \(\langle T_\alpha x, y \rangle \to \langle Tx, y \rangle\) 对所有 \(x, y \in H\)。

\textbf{定义 213.2}（强算子拓扑 / Strong Operator Topology, SOT）：\(B(H)\) 上的\textbf{强算子拓扑}是使所有形如 \(T \mapsto \|Tx\|\)（\(x \in H\)）的半范数连续的最弱拓扑。\(T_\alpha \xrightarrow{SOT} T\) 当且仅当 \(\|T_\alpha x - Tx\| \to 0\) 对所有 \(x \in H\)。

关系：范数拓扑 \(\succ\) SOT \(\succ\) WOT。\(B(H)\) 的单位球在 WOT 下是紧致的（由 Banach-Alaoglu 定理）。

\textbf{定义 213.3}（换位 / Commutant）：设 \(\mathcal{M} \subseteq B(H)\)。\(\mathcal{M}\) 的\textbf{换位}（commutant）为

\[\mathcal{M}' = \{T \in B(H) : TS = ST, \forall S \in \mathcal{M}\}\]

\textbf{命题 213.1}（换位的基本性质）： 1. \(\mathcal{M}'\) 是含单位元的 SOT 闭代数 2. \(\mathcal{M} \subseteq \mathcal{M}''\) 3. \(\mathcal{M}_1 \subseteq \mathcal{M}_2 \implies \mathcal{M}_2' \subseteq \mathcal{M}_1'\) 4. \(\mathcal{M}' = \mathcal{M}'''\)（所以换位只能取奇数次）

\textbf{定义 213.4}（von Neumann 代数 / \(W^*\)-代数）：Hilbert 空间 \(H\) 上的 * -子代数 \(\mathcal{M} \subseteq B(H)\) 称为 \textbf{von Neumann 代数}（或 \(W^*\)-代数），如果它满足以下等价条件之一： 1. \(\mathcal{M} = \mathcal{M}''\)（双换位定理） 2. \(\mathcal{M}\) 在 SOT 下闭且含 \(I\) 3. \(\mathcal{M}\) 在 WOT 下闭且含 \(I\) 4. \(\mathcal{M}\) 是某个 C*-代数的 Banach 空间对偶（即作为 Banach 空间，存在唯一的预对偶 \(\mathcal{M}_*\)）

\subsection{213.2 von Neumann 双换位定理}\label{von-neumann-ux53ccux6362ux4f4dux5b9aux7406}

双换位定理是 von Neumann 代数理论的基石，它将代数的拓扑闭包（SOT 或 WOT）与代数换位操作联系起来。这一对应关系极大地简化了 von Neumann 代数的判定：只需验证 \(\mathcal{A} = \mathcal{A}''\)，而无需显式地刻画 SOT 或 WOT 闭包。从算子代数的角度看，双换位定理等价于''von Neumann 代数 = 含单位元的 * -子代数 \(=\) 其自身的双换位''。这一结果深刻依赖于 Hilbert 空间上算子拓扑的特殊性质------特别是，有界凸集在 WOT 和 SOT 下具有相同的闭包。这一事实的证明需要用到凸集的 Hahn-Banach 分离定理：WOT 和 SOT 的连续线性泛函实际上是相同的，因为两者都由形如 \(T \mapsto \sum_{i=1}^n \langle T x_i, y_i \rangle\) 的泛函生成，仅有的区别在于 \(x_i, y_i\) 的有限性和网的收敛模式。在 von Neumann 的原始工作中，该定理的证明依赖于对 \(H \oplus H \oplus \cdots \oplus H\) 上的对角表示和换位关系的细致分析。Kaplansky 稠密性定理（1951）后来提供了更为简洁的证明路径：\(\mathcal{A}\) 的单位球在 \(\mathcal{A}''\) 的单位球中 SOT 稠密，从而完成了双换位定理的证明。

\textbf{定理 213.2}（von Neumann 双换位定理，1929）：设 \(\mathcal{A} \subseteq B(H)\) 是含单位元的 * -子代数。则 \(\mathcal{A}\) 的 SOT 闭包 = \(\mathcal{A}\) 的 WOT 闭包 = \(\mathcal{A}''\)。

\textbf{证明}：需证 \(\mathcal{A}''\) 等于 \(\mathcal{A}\) 的 SOT 闭包。\(\mathcal{A} \subseteq \mathcal{A}''\) 恒成立，且 \(\mathcal{A}''\) 在 SOT 下闭。关键是证 \(\mathcal{A}'' \subseteq \overline{\mathcal{A}}^{\text{SOT}}\)。对 \(T \in \mathcal{A}''\) 和 \(x_1,\ldots,x_n \in H\)，考虑 \(H^n\) 上的对角表示 \(\pi_n(A)(x_1,\ldots,x_n) = (Ax_1,\ldots,Ax_n)\)。由 Kaplansky 稠密性定理，\(\mathcal{A}\) 的单位球在 \(\mathcal{A}''\) 的单位球中 SOT 稠密。故存在网 \(A_\alpha \in \mathcal{A}\)，\(\|A_\alpha\| \leq \|T\|\)，使 \(A_\alpha \xrightarrow{\text{SOT}} T\)。因此 \(T \in \overline{\mathcal{A}}^{\text{SOT}}\)。WOT 闭包与 SOT 闭包一致（凸集的闭包在两种拓扑下相同）。\(\blacksquare\)

\textbf{推论 213.3}：含单位元的 * -子代数 \(\mathcal{A} \subseteq B(H)\) 是 von Neumann 代数当且仅当 \(\mathcal{A} = \mathcal{A}''\)。

从历史角度看，von Neumann 最初引入双换位定理是为了研究量子力学中可观测量代数的结构。他注意到，物理上有意义的算子代数不仅应在代数运算下封闭，还应在某种''弱极限''的意义下封闭------这正是 WOT 和 SOT 闭包所刻画的性质。双换位定理揭示了代数封闭性与拓扑封闭性之间的深刻等价，为后来的 Tomita-Takesaki 模理论和 Connes 的分类理论奠定了代数基础。此外，双换位定理在表示论中也有重要推论：若 \(\mathcal{A}\) 是 von Neumann 代数且 \(\pi\) 是 \(\mathcal{A}\) 的正规表示，则 \(\pi(\mathcal{A})\) 也是 von Neumann 代数（因为 \(\pi(\mathcal{A})'' = \pi(\mathcal{A})\) 由正规性的保序性质保证）。这一性质使得 von Neumann 代数的表示论远比 C*-代数的表示论更为刚性。

\subsection{213.3 预对偶}\label{ux9884ux5bf9ux5076}

预对偶是 von Neumann 代数理论中一个极为深刻的概念，它将 von Neumann 代数的拓扑性质与其 Banach 空间结构联系起来。本质上，von Neumann 代数 \(\mathcal{M}\) 上的 \(\sigma\)-弱拓扑（即 \(\sigma(\mathcal{M}, \mathcal{M}_*)\) 拓扑）由其预对偶 \(\mathcal{M}_*\) 唯一确定，而 \(\mathcal{M}_*\) 恰由 \(\mathcal{M}\) 上所有 \(\sigma\)-弱连续线性泛函组成。这一观点最早由 Sakai（1956）系统发展，他将 von Neumann 代数从 Hilbert 空间上的具体实现中抽象出来，定义为一个具有预对偶的 C*-代数。这一抽象定义带来了极大的灵活性：许多 von Neumann 代数的构造（如自由积、交叉积）在抽象框架下更为自然，无需预先指定 Hilbert 空间。预对偶的唯一性是一个非平凡的结果------给定 von Neumann 代数 \(\mathcal{M}\)，其预对偶 \(\mathcal{M}_*\) 作为 \(\mathcal{M}^*\) 的子空间在同构意义下是唯一确定的。这是因为 \(\mathcal{M}_*\) 恰为 \(\mathcal{M}^*\) 中所有正规泛函的集合，而''正规性''（即保序递增网的上确界）是 \(\mathcal{M}\) 上的内在性质，不依赖于任何特定的 Hilbert 空间表示。预对偶的存在性赋予 von Neumann 代数一种''自反性''：\(\mathcal{M} \cong (\mathcal{M}_*)^*\) 意味着 \(\mathcal{M}\) 是其预对偶的 Banach 空间对偶，这类似于自反 Banach 空间 \(X \cong X^{**}\) 的性质，但有着本质区别------一般 Banach 空间的对偶 \(X^*\) 几乎从不是 \(W^*\)-代数，而 von Neumann 代数的预对偶 \(\mathcal{M}_*\) 是唯一的 Banach 空间使得 \(\mathcal{M}\) 是其对偶。

\textbf{定义 213.5}（von Neumann 代数的预对偶）：von Neumann 代数 \(\mathcal{M}\) 作为 Banach 空间有唯一的预对偶 \(\mathcal{M}_* \subseteq \mathcal{M}^*\)：

\[\mathcal{M} \cong (\mathcal{M}_*)^*\]

\textbf{定理 213.4}（Sakai 定理，1962）：抽象 \(W^*\)-代数是作为 Banach 空间有预对偶的 C*-代数。任何 \(W^*\)-代数可表示为某个 Hilbert 空间上的 von Neumann 代数。

\textbf{证明}：\(\mathcal{M}\) 上的超弱连续线性泛函全体记作 \(\mathcal{M}_*\)。由 von Neumann 代数的基本性质，\(\mathcal{M}\) 在 \(\sigma\)-弱拓扑（即 \(\sigma(\mathcal{M},\mathcal{M}_*)\) 拓扑）下是局部凸空间。由 Goldstine 定理，\(\mathcal{M}\) 的单位球在 \(\sigma\)-弱拓扑下紧致，且 \((\mathcal{M}_*)^* = \mathcal{M}\)。这给出了 \(\mathcal{M}\) 作为 Banach 空间时的预对偶。反之，若 C*-代数 \(\mathcal{A}\) 作为 Banach 空间有预对偶 \(\mathcal{A}_*\)，则可构造 \(\mathcal{A}\) 在 Hilbert 空间上的忠实正规表示，使其为 von Neumann 代数。Sakai 定理将 von Neumann 代数的 Hilbert 空间定义（WOT 闭）与其内在 Banach 空间性质（有预对偶）统一。\(\blacksquare\)

预对偶概念在因子分类中扮演了关键角色。例如，II\(_1\) 型因子的预对偶 \(\mathcal{M}_*\) 可以显式地通过迹类算子空间 \(L^1(\mathcal{M}, \tau)\) 实现（其中 \(\tau\) 为忠实正规迹），而 \(\mathcal{M}\) 本身则对应于 \(L^\infty(\mathcal{M}, \tau)\)。这一对应关系 \(\mathcal{M} \cong (L^1(\mathcal{M}, \tau))^*\) 是非交换积分理论的核心，它将经典测度论中的对偶关系 \(L^\infty(\mu) \cong (L^1(\mu))^*\) 推广到了非交换情形。进一步，III 型因子的预对偶结构更为复杂，Connes 的分类理论正是通过分析预对偶上的模自同构群的谱性质来完成的。预对偶还与量子概率论密切相关：von Neumann 代数 \(\mathcal{M}\) 上的正规态恰为 \(\mathcal{M}_*\) 中的正元，这使得预对偶成为量子态空间的自然框架。

\subsection{213.4 投影与分类}\label{ux6295ux5f71ux4e0eux5206ux7c7b}

\textbf{定义 213.6}（von Neumann 代数中的投影）：\(\mathcal{M}\) 中的\textbf{投影} \(p\) 满足 \(p = p^* = p^2\)。投影在 von Neumann 代数理论中起核心作用。

\textbf{定义 213.7}（投影的等价关系）：投影 \(p, q \in \mathcal{M}\) 称为 - \textbf{Murray-von Neumann 等价}（\(p \sim q\)）：存在部分等距 \(v \in \mathcal{M}\) 使得 \(v^*v = p\)，\(vv^* = q\) - \(p \preceq q\)：\(p\) 等价于 \(q\) 的子投影（即存在投影 \(q_0 \leq q\) 使得 \(p \sim q_0\)）

\textbf{定理 213.5}（比较定理）：若 \(\mathcal{M}\) 是因子（见定义 213.8），则对任意投影 \(p, q \in \mathcal{M}\)，必有 \(p \preceq q\) 或 \(q \preceq p\)。

\textbf{证明}：在 \(\mathcal{M}\) 中，定义投影的等价类 \([p] = \{q \in \operatorname{Proj}(\mathcal{M}) : q \sim p\}\)。因 \(\mathcal{M}\) 为因子，其中心 \(\mathcal{Z}(\mathcal{M}) = \mathbb{C}I\) 为平凡的。由 Murray-von Neumann 比较理论，对任意 \(p, q\)，存在中心投影 \(z \in \mathcal{Z}(\mathcal{M})\) 使 \(pz \preceq qz\) 且 \(q(1-z) \preceq p(1-z)\)。由于 \(\mathcal{Z}(\mathcal{M}) = \mathbb{C}I\)，\(z\) 只能为 \(0\) 或 \(I\)。若 \(z = I\)，则 \(p \preceq q\)；若 \(z = 0\)，则 \(q \preceq p\)。因此对所有投影 \(p, q\)，比较定理成立。该定理保证因子中投影构成全序的等价类，使因子的类型分类（I, II\(_1\), II\(_\infty\), III）成为可能。\(\blacksquare\)

\textbf{定义 213.8}（因子 / Factor）：von Neumann 代数 \(\mathcal{M}\) 称为\textbf{因子}，如果其中心是平凡的：\(\mathcal{Z}(\mathcal{M}) = \mathcal{M} \cap \mathcal{M}' = \mathbb{C}I\)。任何 von Neumann 代数可分解为因子的直积分（von Neumann 约化理论）。

\hrulefill

\hrulefill

\hrulefill

\hrulefill

\emph{卷十五：泛函分析终。}

\newpage

\chapter{06-分析/04-算子代数/01-vonNeumann代数.md}\label{ux5206ux679004-ux7b97ux5b50ux4ee3ux657001-vonneumannux4ee3ux6570.md}

\section{第182章：因子的 Murray-von Neumann 分类}\label{ux7b2c182ux7ae0ux56e0ux5b50ux7684-murray-von-neumann-ux5206ux7c7b}

Murray 和 von Neumann（1936-1943）完成了因子的完整分类，依据是投影的维数函数的值域。这一分类是算子代数理论最辉煌的成就之一。

\subsection{214.1 维数函数与迹}\label{ux7ef4ux6570ux51fdux6570ux4e0eux8ff9}

\textbf{定义 214.1}（维数函数）：设 \(\mathcal{M}\) 是因子。定义映射 \(D: \operatorname{Proj}(\mathcal{M}) \to [0, \infty]\) 满足： - \(D(p) = D(q) \iff p \sim q\) - \(D(p+q) = D(p) + D(q)\)（若 \(p \perp q\)） - \(D(0) = 0\)

\(D(p)\) 可视为投影 \(p\) 的''相对维数''。\(D\) 的存在性和值域决定了因子的类型。

\textbf{定义 214.2}（有限/无限/半有限投影）： - 投影 \(p\) 是\textbf{有限}的：若 \(p \sim q \leq p \implies q = p\)（即不存在真子投影等价于自身） - 投影 \(p\) 是\textbf{无限}的：若 \(p \sim q < p\)（存在真子投影等价于自身） - 投影 \(p\) 是\textbf{纯无限}的：\(p\) 不含非零有限子投影 - 投影 \(p\) 是\textbf{半有限}的：\(p\) 不含非零有限子投影？不------半有限指 \(p\) 可写为相互正交的有限投影的和

\subsection{214.2 因子类型}\label{ux56e0ux5b50ux7c7bux578b}

\textbf{定义 214.3}（因子的 Murray-von Neumann 分类）： - \textbf{I 型}：含极小投影的因子。\(I_n\)：\(D\) 的值域 = \(\{0, 1, 2, \ldots, n\}\)（\(n = 1, 2, \ldots, \infty\)） - \(I_1\)：\(B(H)\) 上无其他因子的情况（\(\dim H = 1\)） - \(I_n\)：\(B(H) \cong M_n(\mathbb{C})\)（\(\dim H = n\)） - \(I_\infty\)：\(B(H)\) 且 \(\dim H = \infty\) - \textbf{II 型}：无极小投影但含非零有限投影的因子 - \textbf{II\(_1\) 型}：\(I\) 本身是有限投影（即 \(D(I) < \infty\)）。值域 \(D(\operatorname{Proj}(\mathcal{M})) = [0, 1]\) - \textbf{II\(_\infty\) 型}：\(I\) 是纯无限投影但含有限投影。值域 = \([0, \infty]\) - \textbf{III 型}：所有非零投影都是无限投影的因子。值域 = \(\{0, \infty\}\)

\textbf{定理 214.1}（II\(_1\) 型因子的迹）：在 II\(_1\) 型因子 \(\mathcal{M}\) 上存在唯一的忠实正规迹 \(\tau: \mathcal{M} \to \mathbb{C}\) 满足： - \(\tau(xy) = \tau(yx)\) - \(\tau(x^*x) \geq 0\)（正性） - \(\tau(I) = 1\)（归一化） - \(\tau\) 是正规的（对弱收敛连续）

且 \(D(p) = \tau(p)\) 给出维数函数。

\textbf{证明}：由 Murray-von Neumann 维数理论，II\(_1\) 因子上的等价类投影集合在适当运算下形成 \(\mathbb{R}_{\geq 0}\)（确为 \([0,1]\)）。定义 \(\tau(p) = D(p)\) 并线性延拓到自伴元，再延拓到全体。由于 \(\mathcal{M}\) 是 II\(_1\) 型，单位投影 \(I\) 是有限的，故 \(\tau(I) = 1\)。交换性 \(\tau(xy) = \tau(yx)\) 由迹的性质 \(D(p \vee q - q) = D(p) - D(p \wedge q)\) 以及投影的等价关系导出。正规性由维数函数的上半连续性和 \(\mathcal{M}\) 的 WOT 闭性保证。\(\blacksquare\)

\subsection{214.3 具体实现与例子}\label{ux5177ux4f53ux5b9eux73b0ux4e0eux4f8bux5b50}

\textbf{定义 214.4}（群 von Neumann 代数）：设 \(\Gamma\) 是离散群。\(\Gamma\) 的\textbf{群 von Neumann 代数} \(\mathcal{L}(\Gamma)\) 是左正则表示 \(\lambda: \Gamma \to B(\ell^2(\Gamma))\)（\(\lambda_g \delta_h = \delta_{gh}\)）的像的 WOT 闭包。

\[\mathcal{L}(\Gamma) = \lambda(\Gamma)'' \subseteq B(\ell^2(\Gamma))\]

\textbf{定理 214.2}（Murray-von Neumann, 1943）： - \(\mathcal{L}(\Gamma)\) 是 II\(_1\) 型因子当且仅当 \(\Gamma\) 是无限可分共轭类群（ICC 群） - \(\tau(x) = \langle x \delta_e, \delta_e \rangle\) 是 \(\mathcal{L}(\Gamma)\) 上的规范忠实正规迹 - 典型例子：\(\mathcal{L}(\mathbb{F}_2)\)（自由群 \(n \geq 2\) 个生成元）、\(\mathcal{L}(S_\infty)\)（有限置换群的归并）

\textbf{证明}：ICC 条件等价于 \(\mathcal{L}(\Gamma)' \cap \mathcal{L}(\Gamma) = \mathbb{C}I\)，故 \(\mathcal{L}(\Gamma)\) 是因子。迹 \(\tau(x) = \langle x \delta_e, \delta_e \rangle\) 是忠实正规的且 \(\tau(I)=1\)。\(I\) 是有限投影（否则导出 \(\Gamma\) 有有限共轭类），故为 II\(_1\) 型。\(\blacksquare\)

\textbf{问题 214.1}（自由群因子同构问题）：\(\mathcal{L}(\mathbb{F}_m) \cong \mathcal{L}(\mathbb{F}_n)\) 对 \(m \neq n\) 是否成立？这是 von Neumann 代数理论中长期未解决的著名问题之一（与自由概率论和 Voiculescu 的微观状态方法相关）。

\textbf{定义 214.5}（超有限因子 / Approximately Finite Dimensional, AFD）：因子 \(\mathcal{M}\) 称为\textbf{超有限}的，如果 \(\mathcal{M}\) 可由有限维子代数递增序列的并生成（在强算子拓扑下）。

\textbf{定理 214.3}（Murray-von Neumann, 1943）：存在唯一（同构意义下）的超有限 II\(_1\) 型因子，记作 \(\mathcal{R}\)。

\textbf{证明}：构造 \(\mathcal{R}\) 为无穷张量积 \(\overline{\bigotimes}_{n=1}^\infty M_{k_n}(\mathbb{C})\)（\(k_n \geq 2\)）关于迹态 \(\tau = \otimes \tau_{k_n}\)（\(\tau_{k_n}\) 为标准化迹）的 GNS 完备化。唯一性（Murray-von Neumann, 1943）：任意超有限 II\(_1\) 型因子均可写成嵌套矩阵代数并的弱闭包，通过近似等价性将所有有限维子代数的嵌入标准形同伦至同一模型，从而不同构造的超有限 II\(_1\) 型因子同构。这等价于 \(\mathcal{R} \cong \mathcal{R} \otimes \mathcal{R}\)（McDuff 性质）。\(\blacksquare\)

\textbf{推论 214.4}：\(\mathcal{R} \cong \mathcal{L}(S_\infty) \cong \overline{\bigotimes}_{n} M_2(\mathbb{C})\)（无限张量积）。

\subsection{214.4 Powers 因子与 III 型分类的序曲}\label{powers-ux56e0ux5b50ux4e0e-iii-ux578bux5206ux7c7bux7684ux5e8fux66f2}

\textbf{定义 214.6}（Powers 因子，1967）：Powers 构造了连续族的互不同构的 III 型因子 \(R_\lambda\)（\(0 < \lambda < 1\)），通过对超有限 II\(_1\) 型因子 \(\mathcal{R}\) 作 \(\lambda\)-自同构迭代。这些因子可由无穷张量积构造：

\[R_\lambda = \overline{\bigotimes}_{n=1}^\infty (M_2(\mathbb{C}), \phi_\lambda)\]

其中 \(\phi_\lambda\) 是 \(M_2(\mathbb{C})\) 上的非迹状态。

这一工作揭示了 III 型因子的丰富性，为 Connes 的分类铺平了道路。

\hrulefill

\hrulefill

\hrulefill

\hrulefill

\hrulefill

\hrulefill

\section{第183章：Tomita-Takesaki 模理论}\label{ux7b2c183ux7ae0tomita-takesaki-ux6a21ux7406ux8bba}

Tomita-Takesaki 模理论（1967-1970）是 von Neumann 代数理论中最重要的结构定理之一。它从给定的 von Neumann 代数和其上一个忠实的正规半有限权出发，构造了一个单参数自同构群------模自同构群（modular automorphism group）------并建立了正锥与模算子的理论。

\subsection{215.1 权的基本理论}\label{ux6743ux7684ux57faux672cux7406ux8bba}

\textbf{定义 215.1}（权 / Weight）：von Neumann 代数 \(\mathcal{M}\) 上的\textbf{权} \(\varphi\) 是映射 \(\varphi: \mathcal{M}_+ \to [0, \infty]\)，满足 - \(\varphi(x + y) = \varphi(x) + \varphi(y)\) - \(\varphi(\lambda x) = \lambda \varphi(x)\)（\(\lambda \geq 0\)，约定 \(0 \cdot \infty = 0\)）

权是迹和状态的推广。有限权（\(\varphi(I) < \infty\)）可标准化为状态。

\textbf{定义 215.2}（半有限/忠实/正规权）： - \(\varphi\) 是\textbf{忠实}的：\(\varphi(x) = 0 \implies x = 0\)（\(x \in \mathcal{M}_+\)） - \(\varphi\) 是\textbf{半有限}的：\(\{x \in \mathcal{M}_+ : \varphi(x) < \infty\}\) 在 \(\mathcal{M}_+\) 中稠密 - \(\varphi\) 是\textbf{正规}的：对 \(\mathcal{M}_+\) 中任意递增网 \(x_\alpha \uparrow x\)，\(\varphi(x_\alpha) \uparrow \varphi(x)\)

\textbf{定理 215.1}（权的基本事实）：任何 von Neumann 代数都容纳一个忠实的正规半有限权（fns weight）。这一结果依赖于 Hahn-Banach 定理和 von Neumann 代数的预对偶性质。

\textbf{证明}：由 Zorn 引理，存在极大族 \(\{\varphi_i\}\) 的正规正线性泛函，满足 \(\sum \|\varphi_i\| < \infty\)（由 \(\mathcal{M}^*\) 的分解性质）。定义 \(\varphi = \sum \varphi_i\)。忠实性：若 \(\varphi(x) = 0\) 则 \(x = 0\)（否则由 Hahn-Banach 存在分离的正规正泛函）。半有限性：\(\{x \in \mathcal{M}_+ : \varphi(x) < \infty\}\) 包含所有 \(\varphi_i\) 的支撑投影的有限和，其并弱稠密于 \(\mathcal{M}\)。正规性由每个 \(\varphi_i\) 的正规性和级数收敛性保证。\(\blacksquare\)

\subsection{215.2 GNS 构造与 Tomita 算子}\label{gns-ux6784ux9020ux4e0e-tomita-ux7b97ux5b50}

\textbf{定理 215.2}（GNS 构造 / Gelfand-Naimark-Segal, 1943）：设 \(\varphi\) 是 \(\mathcal{M}\) 上的忠实正规半有限权。则存在 Hilbert 空间 \(L^2(\mathcal{M}, \varphi)\) 和忠实正规 * -表示 \(\pi_\varphi: \mathcal{M} \to B(L^2(\mathcal{M}, \varphi))\)，使得 \(\pi_\varphi(\mathcal{M})\) 是 von Neumann 代数，且有循环和分离向量。

\textbf{证明}：在左理想 \(\mathfrak{n}_\varphi = \{x \in \mathcal{M} : \varphi(x^* x) < \infty\}\) 上定义内积 \(\langle x, y \rangle_\varphi = \varphi(y^* x)\)。由 \(\varphi\) 的正规性，\(\mathfrak{n}_\varphi\) 在 \(\sigma\)-弱拓扑下稠密。令 \(H = L^2(\mathcal{M}, \varphi)\) 为 \(\mathfrak{n}_\varphi\) 在该内积下的完备化。定义表示 \(\pi_\varphi(x)(y) = xy\)（\(x \in \mathcal{M}\)，\(y \in \mathfrak{n}_\varphi\)）。由 Cauchy-Schwarz 不等式，\(\pi_\varphi(x)\) 延拓为 \(H\) 上的有界算子。\(\pi_\varphi\) 的忠实性由 \(\varphi\) 的忠实性保证：若 \(\pi_\varphi(x) = 0\)，则 \(\varphi(x^* x) = 0\)，故 \(x = 0\)。正规性由 \(\varphi\) 的正规性通过 \(\sigma\)-弱拓扑的连续性传递。由 von Neumann 二次换位定理，\(\pi_\varphi(\mathcal{M})'' = \pi_\varphi(\mathcal{M})\)。循环向量取为单位元 \(1 \in \mathfrak{n}_\varphi\) 的像。\(\blacksquare\)

GNS 构造是算子代数理论中最重要的技术工具之一。它从一个抽象的权（或状态）出发，构造出具体的 Hilbert 空间表示，使得原先代数的元素在这个空间中作为有界算子作用。一旦有了循环和分离向量，就可以定义 Tomita 算子 \(S\)------它捕捉了代数的对合运算 \(x \mapsto x^*\) 在 Hilbert 空间表示下的具体形态。Tomita 算子是通向 Tomita-Takesaki 模理论的入口，该理论深刻揭示了 von Neumann 代数与其换位代数之间的对称性。

\textbf{定义 215.3}（Tomita 算子）：给定 von Neumann 代数 \(\mathcal{M}\) 和循环分离向量 \(\Omega\)，定义 \textbf{Tomita 算子} \(S\)：

\[S_0(x\Omega) = x^*\Omega, \quad x \in \mathcal{M}\]

\(S_0\) 是稠定的共轭线性算子。其闭包 \(S = \overline{S_0}\) 称为 Tomita 算子。

\textbf{定理 215.3}（极分解）：\(S\) 有极分解 \(S = J \Delta^{1/2}\)，其中： - \(J\) 是共轭线性等距（\(J^2 = I\)，\(J = J^*\)），称为\textbf{模共轭}（modular conjugation） - \(\Delta = S^*S\) 是正的自伴算子，称为\textbf{模算子}（modular operator）

\textbf{证明}：\(S\) 为稠定共轭线性闭算子。令 \(\Delta = S^* S\)（正自伴稠定算子，因 \(S\) 为闭算子）。由无界算子的极分解定理（共轭线性版本），\(S = J \Delta^{1/2}\)，其中 \(J\) 是共轭线性部分等距。由于 \(S^2 \subseteq I\)（\(S(x\Omega) = x^*\Omega\)，\(S^2(x\Omega) = S(x^*\Omega) = x\Omega\)），推出 \(J^2 = I\) 且 \(J = J^*\)。\(\blacksquare\)

\subsection{215.3 Tomita-Takesaki 定理}\label{tomita-takesaki-ux5b9aux7406}

\textbf{定理 215.4}（Tomita-Takesaki 定理）：设 \(\mathcal{M} \subseteq B(H)\) 是 von Neumann 代数，\(\Omega\) 是循环分离向量。则：

\begin{enumerate}
\def\labelenumi{\arabic{enumi}.}
\item
  \(J \mathcal{M} J = \mathcal{M}'\)（模共轭将 \(\mathcal{M}\) 映为其换位）
\item
  \(\Delta^{it} \mathcal{M} \Delta^{-it} = \mathcal{M}\) 对所有 \(t \in \mathbb{R}\)（模自同构群保持 \(\mathcal{M}\)）
\item
  \textbf{模自同构群} \(\sigma_t^\varphi(x) = \Delta^{it} x \Delta^{-it}\)（\(t \in \mathbb{R}\)）是 \(\mathcal{M}\) 的单参数自同构群
\item
  \textbf{KMS 条件}：对任意 \(x, y \in \mathcal{M}\)，存在有界连续函数 \(F(z)\) 在条带 \(\{z : 0 \leq \Im(z) \leq 1\}\) 上解析，边界值为 \[F(t) = \varphi(\sigma_t^\varphi(x) y), \quad F(t+i) = \varphi(y \sigma_t^\varphi(x))\]
\end{enumerate}

\textbf{证明}：（1）由 \(S\) 的定义，对 \(x \in \mathcal{M}\)，\(S(x\Omega)=x^*\Omega\)。取 \(y \in \mathcal{M}'\)，则 \(yS(x\Omega) = yx^*\Omega = x^*y\Omega\)，推出 \(Sy = yS\) 在 \(\mathcal{M}\Omega\) 上。由此 \(J\mathcal{M}J \subseteq \mathcal{M}'\)。对称地 \(F(x'\Omega) = (x')^*\Omega\) 用于 \(\mathcal{M}'\) 上给出反包含，故 \(J\mathcal{M}J = \mathcal{M}'\)。（2）\(\Delta^{it}\) 为 \(\Delta\) 的 Stone 酉群。由 KMS 条件：\(\varphi(\sigma_t(x)y) = \langle\Delta^{it}x\Omega, y^*\Omega\rangle\)。解析延拓至条带 \(0 \leq \Im z \leq 1\)，利用 Phragmen-Lindelof 原理和三线定理得 \(\sigma_t(\mathcal{M}) \subseteq \mathcal{M}\)。由对称性 \(\sigma_{-t}(\mathcal{M}) \subseteq \mathcal{M}\)，故 \(\sigma_t\) 是 \(\mathcal{M}\) 的自同构。（3-4）KMS 条件定义 \(F(z) = \langle\Delta^{-iz}x\Omega, y\Omega\rangle\)，其在条带上解析，边界值满足所需等式。\(\blacksquare\)

\textbf{推论 215.5}：对任意忠实正规半有限权 \(\varphi\)，存在唯一的一参数自同构群 \(\sigma_t^\varphi\) 满足 KMS 条件。这称为 \(\varphi\) 的\textbf{模自同构群}。

\subsection{215.4 Connes 的余循环定理与 III 型分类}\label{connes-ux7684ux4f59ux5faaux73afux5b9aux7406ux4e0e-iii-ux578bux5206ux7c7b}

\textbf{定理 215.6}（Connes 余循环定理，1973）：设 \(\varphi, \psi\) 是 \(\mathcal{M}\) 上的两个忠实正规半有限权。则模自同构群 \(\sigma_t^\varphi\) 和 \(\sigma_t^\psi\) 是余循环等价的：存在连续映射 \(u: \mathbb{R} \to \mathcal{U}(\mathcal{M})\) 满足 \(u_{t+s} = u_t \sigma_t^\varphi(u_s)\)，使得

\[\sigma_t^\psi(x) = u_t \sigma_t^\varphi(x) u_t^*\]

\textbf{证明}：设 \(h = d\psi/d\varphi\) 为 Radon-Nikodym 型导数（Connes 空间 \(L^1(\mathcal{M})\) 中的正算子）。由 KMS 条件，单参数族 \(\sigma_t^\varphi(h^{is})\) 和 \(\sigma_t^\psi\) 满足相同的余循环恒等式。通过解析延拓和 Phragmen-Lindelof 定理，存在 \(u_t \in \mathcal{U}(\mathcal{M})\) 实现两个模自同构群之间的余循环等价。\(u_t\) 的连续性由 \(\mathcal{M}\) 上权空间的结构导出。\(\blacksquare\)

\textbf{定理 215.7}（Connes 的 III 型分类定理，1973）：设 \(\mathcal{M}\) 是 III 型因子。定义

\[S(\mathcal{M}) = \bigcap_{\varphi} \operatorname{Sp}(\Delta_\varphi) \subseteq [0, \infty)\]

其中 \(\Delta_\varphi\) 是权 \(\varphi\) 的模算子，\(\operatorname{Sp}\) 表示谱。Connes 的分类：

\begin{itemize}
\tightlist
\item
  \textbf{III\(_0\) 型}：\(S(\mathcal{M}) = \{0, 1\}\)
\item
  \textbf{III\(_\lambda\) 型}（\(0 < \lambda < 1\)）：\(S(\mathcal{M}) = \{0\} \cup \{\lambda^n : n \in \mathbb{Z}\}\)
\item
  \textbf{III\(_1\) 型}：\(S(\mathcal{M}) = [0, \infty)\)
\end{itemize}

这些因子互不同构。Haagerup（1987）证明了 III\(_1\) 型因子的唯一性。

\textbf{证明}：由余循环定理，\(\sigma^\varphi\) 的外自同构类不依赖 \(\varphi\)。定义 \(S(\mathcal{M}) = \bigcap_\varphi \operatorname{Sp}(\Delta_\varphi)\)。Connes 证明 \(S(\mathcal{M}) \setminus \{0\}\) 是乘法群 \(\mathbb{R}_+^\times\) 的闭子群，其可能性恰为 \(\{1\}\)（III\(_0\)）、\(\lambda^{\mathbb{Z}}\)（III\(_\lambda\)）或 \(\mathbb{R}_+\)（III\(_1\)）。若 \(\mathcal{M}\) 不是 III\(_0\)，其权流非平凡，由 Connes-Takesaki 对偶和 Krieger 遍历理论给出完全不变量。各类型互不同构可由 \(T\)-不变量区分：III\(_\lambda\) 对应 \(T(\mathcal{M}) = (2\pi/\ln\lambda)\mathbb{Z}\)。Haagerup（1987）证明了 III\(_1\) 型因子的唯一性。\(\blacksquare\)

\subsection{215.5 流与 Connes 的 T-不变量}\label{ux6d41ux4e0e-connes-ux7684-t-ux4e0dux53d8ux91cf}

\textbf{定义 215.7}（非平凡流 / Flow of Weights, Connes-Takesaki）：对 III 型因子 \(\mathcal{M}\)，Connes 和 Takesaki 构造了 III 型因子的\textbf{权流}（flow of weights）：一个 \(\mathbb{R}\) 作用在可交换 von Neumann 代数上的遍历流。对 III\(_\lambda\) 型因子（\(\lambda \neq 0\)），此流是周期的（周期 \(T = -2\pi/\ln\lambda\)）。

权流为 III 型因子提供了一个本质上交换的不变量------尽管 III 型因子本身高度非交换，但其权流却是一个交换 von Neumann 代数上的遍历作用。这一构造来源于 Connes-Takesaki 对偶定理，它将 III 型因子的分类问题部分地转化为遍历论问题。基于权流，可以进一步提取出更易于计算的数值不变量------\(T\)-集，它刻画了模自同构群中哪些时间参数对应了内自同构（即代数内部的自同构，而非外部等价类）。对 III\(_\lambda\) 型因子，\(T\)-集恰好是周期晶格，其周期由 \(\lambda\) 唯一确定。

\textbf{定义 215.8}（\(T\)-不变量）：III 型因子的 \(T\)-集定义为

\[T(\mathcal{M}) = \{t \in \mathbb{R} : \sigma_t^\varphi \text{ 是内自同构}\}\]

\(T(\mathcal{M})\) 是 \(\mathbb{R}\) 的子群，不依赖于 \(\varphi\) 的选择。对 III\(_\lambda\) 型因子，\(T(\mathcal{M}) = (2\pi/\ln\lambda)\mathbb{Z}\)。

\hrulefill

\hrulefill

\hrulefill

\hrulefill

\hrulefill

\hrulefill

\section{第184章：自由概率论与随机矩阵}\label{ux7b2c184ux7ae0ux81eaux7531ux6982ux7387ux8bbaux4e0eux968fux673aux77e9ux9635}

自由概率论由 Voiculescu（1985）创立，是 von Neumann 代数理论和随机矩阵理论相结合的产物。它为 II\(_1\) 型因子的研究提供了强大的新工具。

\subsection{216.1 非交换概率空间}\label{ux975eux4ea4ux6362ux6982ux7387ux7a7aux95f4}

\textbf{定义 216.1}（非交换概率空间）：\textbf{非交换概率空间} \((\mathcal{A}, \varphi)\) 由单位元 * -代数 \(\mathcal{A}\) 和忠实的正线性泛函 \(\varphi: \mathcal{A} \to \mathbb{C}\)（\(\varphi(1) = 1\)，状态）组成。

\begin{itemize}
\tightlist
\item
  \textbf{经典概率}：\((\mathcal{A}, \varphi) = (L^\infty(X, \mu), \mathbb{E}_\mu)\)
\item
  \textbf{自由概率}：\((\mathcal{A}, \varphi) = (\mathcal{M}, \tau)\)，其中 \(\mathcal{M}\) 是 II\(_1\) 型因子，\(\tau\) 是忠实的正规迹
\end{itemize}

非交换概率空间的概念将经典概率论中的 Kolmogorov 公理体系推广到非交换代数框架。在经典概率论中，随机变量是样本空间 \(\Omega\) 上的可测函数，其分布由测度 \(\mu\) 决定，期望 \(\mathbb{E}_\mu\) 是 \(L^\infty(\Omega, \mu)\) 上的正线性泛函。在非交换推广中，样本空间 \(\Omega\) 被一个 von Neumann 代数 \(\mathcal{M}\) 取代，而概率测度 \(\mu\) 被一个状态 \(\varphi\)（或 II\(_1\) 型因子情形下的迹 \(\tau\)）取代。这一推广的动机部分来自量子力学：量子可观测量的代数是非交换的，其期望值由密度矩阵 \(\rho\) 给出的迹 \(\tau(\rho \cdot)\) 描述。从范畴论角度看，非交换概率空间构成了比经典概率空间更丰富的范畴：经典概率空间嵌入非交换概率空间范畴（通过取 \(L^\infty\) 代数），但非交换概率空间包含了经典框架无法描述的''量子随机变量''------即不一定彼此交换的算子。这种非交换性导致了''自由独立性''（freeness）这一全新的独立性概念，它与经典独立性的根本区别在于：经典独立随机变量乘积的期望因子化，而自由独立随机变量交错乘积的期望消没。非交换概率空间的研究揭示了概率论、算子代数、组合数学和随机矩阵理论之间的深刻联系。例如，Voiculescu 发现独立随机矩阵的渐近特征值分布可以由自由概率论中的自由卷积运算精确描述，这一发现将随机矩阵理论中的 Wigner 半圆律纳入到自由中心极限定理的框架中。

\subsection{216.2 自由独立性}\label{ux81eaux7531ux72ecux7acbux6027}

\textbf{定义 216.2}（自由独立性 / Freeness）：设 \((\mathcal{A}, \varphi)\) 是非交换概率空间。子代数族 \(\mathcal{A}_1, \ldots, \mathcal{A}_n \subseteq \mathcal{A}\) 是\textbf{自由独立}的，如果对任意 \(a_j \in \mathcal{A}_{i_j}\)（\(i_1 \neq i_2 \neq \cdots \neq i_k\)，即相邻指标不同）且 \(\varphi(a_j) = 0\)，有

\[\varphi(a_1 a_2 \cdots a_k) = 0\]

自由独立性是自由概率论中最核心的概念创新。与经典独立不同------经典独立要求联合矩的因子化------自由独立性的定义基于''交错乘积的期望消没''这一代数条件。这种独立性模式恰好描述了独立随机矩阵在维数趋于无穷大时的联合分布行为。了解了自由独立之后，自然的问题是：如何计算两个自由独立随机变量之和的分布？答案由自由卷积运算给出，它是对经典概率论中卷积运算（独立随机变量之和）的自由类比。

\textbf{定义 216.3}（自由卷积）：若随机变量 \(a, b\) 关于 \(\varphi\) 的分布在解析意义下已知，且 \(a, b\) 是自由独立的，则 \(a+b\) 的分布由\textbf{自由加性卷积} \(\mu \boxplus \nu\) 给出。

\textbf{定理 216.1}（自由中心极限定理）：设 \(a_1, a_2, \ldots\) 是自由独立且同分布的随机变量，\(\varphi(a_i) = 0\)，\(\varphi(a_i^2) = 1\)。则

\[S_n = \frac{1}{\sqrt{n}} \sum_{i=1}^n a_i\]

在分布意义下收敛于\textbf{半圆律}：\(\mu_{SC}(dx) = \frac{1}{2\pi} \sqrt{4 - x^2} \, dx\)（\(|x| \leq 2\)）。

\textbf{证明}：由自由独立性，\(\varphi(S_n^k)\) 的矩仅贡献于非交叉分划（non-crossing partitions）。自由独立条件下，交叉分划的矩为零。非交叉分划计数由 Catalan 数 \(C_{k/2}\) 给出（\(k\) 偶），故 \(\varphi(S_n^{2k}) \to C_k\)。半圆律的矩恰为 Catalan 数：\(\frac{1}{2\pi}\int_{-2}^2 x^{2k} \sqrt{4-x^2} dx = C_k\)。由矩收敛定理得分布收敛。\(\blacksquare\)

自由中心极限定理是自由概率论的标志性结果，它在结构上完美模仿了经典中心极限定理，但所得极限分布从 Gauss 分布变成了 Wigner 半圆分布。这两个极限分布在各自的独立框架中都由 Catalan 数描述其矩序列，这并非巧合------非交叉分划的计数在自由概率中的角色恰好对应了经典概率中所有分划的计数。正如下面的推论所指出的，半圆律在自由世界中的地位与正态分布在经典世界中的地位完全平行。

\textbf{推论 216.2}：半圆律在自由概率论中的地位类似于高斯分布（正态分布）在经典概率论中的地位。

\textbf{证明}：在经典概率论中，中心极限定理确立了正态分布 \(\mathcal{N}(0,1)\) 作为独立同分布随机变量标准化和的极限分布。在自由概率论中，定理 216.1 确立了半圆律作为自由独立同分布随机变量标准化和的极限分布。两者的矩结构均为 Catalan 数：正态分布的矩为 \((2k-1)!!\)（仅偶次矩非零），半圆律的矩为 Catalan 数 \(C_k\)。进一步，自由概率论中的半圆元与经典概率论中的 Gauss 随机变量在各自独立性的框架下均满足类似的累积量消没性质（非交叉累积量）。\(\blacksquare\)

\subsection{216.3 自由概率论的算子代数应用}\label{ux81eaux7531ux6982ux7387ux8bbaux7684ux7b97ux5b50ux4ee3ux6570ux5e94ux7528}

\textbf{定理 216.3}（Voiculescu, 1990-1995）：自由群因子 \(\mathcal{L}(\mathbb{F}_n)\) 的微观状态方法： - \(\mathcal{L}(\mathbb{F}_n)\) 可嵌入超有限 II\(_1\) 型因子的超积中 - 由此获得 \(\mathcal{L}(\mathbb{F}_n)\) 中生成元的自由独立性和渐近自由性 - 自由群因子同构问题可通过自由熵 \(\chi^*(X_1, \ldots, X_n)\) 研究

进一步：Dykema（1994）和 Radulescu（1994）独立证明了压缩自由群因子 \(L(\mathbb{F}_n)_t\)（\(t > 0\) 连续参数）的理论，将其推广为连续参数的族。更一般地，Voiculescu 的自由维数（自由熵）发展为 \(L^2\)-Betti 数和 Atiyah 问题的关键工具。

\textbf{证明}：Voiculescu 构造了微观状态空间 \(\Gamma_R(X_1, \ldots, X_n; m, k, \varepsilon)\)，其中 \(R > 0\) 控制算子范数，\(m\) 为矩的阶数，\(k\) 为有限维矩阵的维数，\(\varepsilon > 0\) 为逼近精度。该空间由 \(n\) 元组 \((A_1, \ldots, A_n)\) 组成，\(A_i\) 为 \(k \times k\) 自伴矩阵，满足 \(\|A_i\| \leq R\) 且所有不超过 \(m\) 阶的矩与目标迹 \(\tau\) 的矩之差不超过 \(\varepsilon\)。对自由群因子 \(\mathcal{L}(\mathbb{F}_n)\) 的生成元 \(X_1, \ldots, X_n\)（\(X_i = X_i^*\)），取 \(k \times k\) 随机矩阵逼近：当 \(k \to \infty\) 时，独立 GUE 矩阵的矩收敛到半圆律，自由独立的 GUE 矩阵的联合矩收敛到自由半圆族的联合矩。通过 Dykema 的插值自由随机变量构造，可获得具有任意自由维数的生成元。微观状态空间非空的证明基于：对任意 \(\varepsilon > 0\) 和 \(m\)，存在 \(k\) 及矩阵 \((A_1^{(k)}, \ldots, A_n^{(k)})\) 使矩差距小于 \(\varepsilon\)。自由熵 \(\chi^*(X_1, \ldots, X_n)\) 定义为微观状态空间体积的渐近增长率：\(\chi^* = \limsup_{m \to \infty, \varepsilon \to 0} \limsup_{k \to \infty} \frac{1}{k^2} \log \operatorname{vol}(\Gamma_R(\cdots))\)。\(\blacksquare\)

\textbf{定义 216.4}（自由熵 / Free Entropy）：Voiculescu 引入了微状态自由熵 \(\chi(X_1, \ldots, X_n)\)，度量 von Neumann 代数生成元集的''自由度''。自由熵与 \(L^2\)-Betti 数有深刻的联系，并在自由群因子同构问题的研究中起关键作用。

\subsection{自由概率论与算子代数的深层联系}\label{ux81eaux7531ux6982ux7387ux8bbaux4e0eux7b97ux5b50ux4ee3ux6570ux7684ux6df1ux5c42ux8054ux7cfb}

自由概率论自 Voiculescu 创立以来，已发展成为算子代数、随机矩阵理论和组合数学之间的核心桥梁。\textbf{自由熵理论}（Voiculescu 1993-1998）的提出源于对自由群因子同构问题的研究------自由群因子 \(\mathcal{L}(\mathbb{F}_n)\) 是否对不同的 \(n\) 互不同构？这一问题是 von Neumann 代数理论中最重要的未解决问题之一。自由熵 \(\chi(X_1, \ldots, X_n)\) 通过微状态方法（microstates approach）度量了生成元集的''连续维度''，其基本性质包括：\(\chi\) 在自由独立下可加，\(\chi\) 不超过生成元个数的 \(\frac{1}{2}\log 2\pi e\)，且若 \(\chi(X_1, \ldots, X_n) > -\infty\)，则 \(\{X_1, \ldots, X_n\}\) 生成的非交换概率空间是 Connes 嵌入的。\textbf{自由维数} \(\delta_0(X_1, \ldots, X_n)\) 是自由熵的衍生概念，给出了 von Neumann 代数''非交换维度''的度量。Dykema（1993-1994）和 Radulescu（1992-1994）独立地发展了\textbf{压缩自由群因子}理论，证明了 \(\mathcal{L}(\mathbb{F}_n)_t \cong \mathcal{L}(\mathbb{F}_m)_s\) 当且仅当 \(t = s\)（对所谓''插值自由群因子''成立），并引入了自由 Fisher 信息、自由 Cramer-Rao 不等式等经典概率论概念的自由类比。\textbf{随机矩阵理论}与自由概率论的联系由 Voiculescu（1991）揭示：独立 Gaussian 酉矩阵（GUE）的渐近特征值分布由半圆律给出，且独立的 GUE 矩阵族在极限下是自由独立的。这一发现将随机矩阵的谱统计与算子代数的自由独立性联系起来，催生了\textbf{自由确定性}（free determinism）和\textbf{自由无穷小独立}理论。\textbf{双自由概率论}（bi-free probability, Voiculescu 2014）是自由概率论的进一步推广，研究具有两个自由面的双自由独立族，与 \(C^*\)-代数中的双自由积和量子群的表示论有深刻联系。自由概率论还与\textbf{组合数学}中的非交叉分划（non-crossing partition）和\textbf{图论}中的大 N 极限（如 Erdős-Rényi 随机图的谱）密切关联，形成了分析、代数和组合学之间的丰富交叉领域。

\hrulefill

\hrulefill

\subsection{127.A vonNeumann代数补充}\label{a-vonneumannux4ee3ux6570ux8865ux5145}

\begin{quote}
积分方程是未知函数出现在积分号下的方程，是泛函分析、数学物理和工程数学的经典交汇领域。本卷系统建立 Fredholm 积分方程、Volterra 积分方程、奇异性积分方程和 Wiener-Hopf 积分方程的理论及其数值方法。积分方程理论是 Hilbert 空间算子理论的直接应用，也是谱理论、紧算子和变分方法的经典范例。
\end{quote}

\hrulefill

\hrulefill

\hrulefill

\hrulefill

\hrulefill

\hrulefill

\section{第185章：Fredholm 积分方程}\label{ux7b2c185ux7ae0fredholm-ux79efux5206ux65b9ux7a0b}

Fredholm 积分方程形如

\[u(x) - \lambda \int_a^b K(x,y) u(y) dy = f(x)\]

或其特征值问题 \(\int_a^b K(x,y) u(y) dy = \mu u(x)\)。这是一类由 Fredholm（1903）系统研究的积分方程，其理论与紧算子的谱理论完全对应。

\subsection{217.1 Fredholm 择一定理}\label{fredholm-ux62e9ux4e00ux5b9aux7406}

\textbf{定义 217.1}（Fredholm 积分算子）：在 \(L^2([a,b])\) 上定义积分算子

\[(Ku)(x) = \int_a^b K(x,y) u(y) dy\]

若 \(K \in L^2([a,b] \times [a,b])\)，则 \(K\) 是 Hilbert-Schmidt 算子，从而是紧算子。

\textbf{定理 217.1}（Fredholm 择一定理）：对 Fredholm 积分方程 \((I - \lambda K)u = f\)，以下二者恰有一者成立： - （第一择一）对任意 \(f\)，方程有唯一解 \(u = (I - \lambda K)^{-1} f\) - （第二择一）齐次方程 \((I - \lambda K)u = 0\) 有非零解，此时非齐次方程有解当且仅当 \(f \in (\ker(I - \lambda K^*))^\perp\)

这是紧算子理论中谱的基本性质：紧算子的非零谱点必为特征值，且 \(\lambda^{-1}\) 是 \(K\) 的特征值当且仅当 \(\lambda\) 是 \((I - \lambda K)^{-1}\) 的奇点。

\textbf{证明}：\(K\) 是紧算子。由紧算子的 Riesz-Schauder 理论，\(I - \lambda K\) 是 Fredholm 算子（指标为 \(0\)）。若 \(\ker(I - \lambda K) = \{0\}\)，则 \(I - \lambda K\) 是单射，由 Fredholm 择一性它也是满射，故对任意 \(f\) 有唯一解。若 \(\ker(I - \lambda K) \neq \{0\}\)，则 \(\dim \ker(I - \lambda K) = \dim \ker(I - \lambda K^*) < \infty\)，且 \(f \in \operatorname{im}(I - \lambda K) \iff f \perp \ker(I - \lambda K^*)\)。Riesz-Schauder 定理还保证 \(I - \lambda K\) 的谱仅由特征值 \(\lambda^{-1}\) 和可能的聚点 \(0\) 组成。\(\blacksquare\)

\subsection{217.2 Fredholm 行列式}\label{fredholm-ux884cux5217ux5f0f}

\textbf{定义 217.2}（Fredholm 行列式）：对核为 \(K(x,y)\) 的 Fredholm 积分方程，Fredholm 行列式定义为

\[D(\lambda) = \sum_{n=0}^{\infty} \frac{(-\lambda)^n}{n!} \int_a^b \cdots \int_a^b \det(K(x_i, x_j))_{i,j=1}^n dx_1 \cdots dx_n
\]

Fredholm 行列式的定义是经典有限维行列式在积分方程中的自然推广：\(n\) 阶项涉及核 \(K\) 在 \(n\) 个点上的 \(n \times n\) 行列式在配置空间上的积分。当核 \(K\) 是有限维矩阵时，该级数在 \(n\) 超过矩阵秩时截断，退化回通常的特征多项式。但在无穷维积分方程中，该级数可以无限延续，定义为 \(\lambda\) 的整函数。Fredholm 的洞见在于：该整函数的零点恰好对应积分算子的特征值倒数------正如有限维情形中 \(\det(I - \lambda A) = 0\) 当且仅当 \(\lambda^{-1}\) 是 \(A\) 的特征值。预解核 \(R(x,y;\lambda)\) 作为两个广义行列式之比，给出了非齐次积分方程解的闭式表达。

\textbf{定理 217.2}：\(D(\lambda)\) 是 \(\lambda\) 的整函数，其零点恰为 \(K\) 的特征值的倒数 \(\lambda = 1/\mu\)（重数对应）。预解核

\[R(x,y;\lambda) = \frac{D(x,y;\lambda)}{D(\lambda)}\]

其中 \(D(x,y;\lambda)\) 是 Fredholm 一阶子式行列式。

\textbf{证明}：\(D(\lambda)\) 是 \(\lambda\) 的整函数，由 Hadamard 不等式 \(|\det(K(x_i,x_j))| \leq n^{n/2} \|K\|_\infty^n\) 保证级数对任意 \(\lambda\) 收敛。若 \(\lambda = 1/\mu\)（\(\mu\) 为 \(K\) 的特征值），则齐次方程 \((I - \lambda K)u = 0\) 有非零解。由 Fredholm 理论，此时 \(D(\lambda) = 0\)。特征值的重数等于 \(D(\lambda)\) 的零点重数。预解核 \(R(x,y;\lambda) = D(x,y;\lambda)/D(\lambda)\) 的极点恰为 \(D(\lambda)\) 的零点，即 \(K\) 的特征值倒数。\(\blacksquare\)

\subsection{217.3 对称核与 Hilbert-Schmidt 定理}\label{ux5bf9ux79f0ux6838ux4e0e-hilbert-schmidt-ux5b9aux7406}

\textbf{定理 217.3}（Hilbert-Schmidt 展开定理）：若 \(K(x,y)\) 是实对称且平方可积的核（\(K(x,y) = K(y,x)\)，\(\iint |K|^2 < \infty\)），则

\[K(x,y) = \sum_{n=1}^{\infty} \mu_n u_n(x) u_n(y)\]

其中 \(\{\mu_n\}\) 是 \(K\) 的非零特征值（每个按重数列出），\(\{u_n\}\) 是对应的标准正交特征函数。级数在 \(L^2\) 意义下收敛；若 \(K\) 连续且正定，则收敛是绝对且一致的（Mercer 定理）。

\textbf{证明}：\(K\) 作为对称 Hilbert-Schmidt 算子，其非零特征值 \(\{\mu_n\}\) 为实数且 \(\sum \mu_n^2 < \infty\)（由 Hilbert-Schmidt 范数 \(\|K\|_{HS}^2 = \sum \mu_n^2\)）。对应特征函数 \(\{u_n\}\) 可选取为标准正交系。对任意固定 \(x\)，将 \(K(x,\cdot)\) 按 \(\{u_n\}\) 展开：\(K(x,\cdot) = \sum \langle K(x,\cdot), u_n \rangle u_n = \sum \mu_n u_n(x) u_n(\cdot)\)。\(L^2\) 收敛性由 Bessel 不等式和 \(\sum \mu_n^2 < \infty\) 保证。\(\blacksquare\)

\textbf{定理 217.4}（Mercer 定理）：若 \(K\) 是连续对称正定核（即对任意 \(f\)，\(\iint K(x,y)f(x)f(y)dx dy \geq 0\)），则特征值 \(\mu_n \geq 0\)，且

\[K(x,y) = \sum_{n=1}^{\infty} \mu_n u_n(x) u_n(y)\]

收敛是绝对且一致的。特别地，迹公式 \(\sum \mu_n = \int K(x,x) dx\) 成立。

\textbf{证明}：由正定性，对任意 \(x\)，\(K(x,x) \geq 0\)。由 Hilbert-Schmidt 展开 \(K(x,y) = \sum \mu_n u_n(x) u_n(y)\)，其中 \(\mu_n \geq 0\)（正定性保证）。对任意 \(N\)，\(\sum_{n=1}^N \mu_n |u_n(x)|^2 \leq K(x,x)\)（由正定性，截断核 \(\sum_{n=1}^N \mu_n u_n(x) u_n(y)\) 仍正定，且 \(K(x,x) - \sum_{n=1}^N \mu_n |u_n(x)|^2 \geq 0\)）。由 Dini 定理，级数 \(\sum \mu_n |u_n(x)|^2\) 在 \([a,b]\) 上一致收敛。由于 \(\sum \mu_n u_n(x) u_n(y)\) 中各项由 Cauchy-Schwarz 受控于 \(\sum \mu_n |u_n(x)|^2\) 和 \(\sum \mu_n |u_n(y)|^2\)，故原级数绝对且一致收敛。取 \(x=y\) 并积分得 \(\sum \mu_n = \int K(x,x) dx\)。\(\blacksquare\)

\textbf{应用}：Mercer 定理是核方法（kernel methods）------包括支持向量机（SVM）和高斯过程回归------的数学基础，将正定核表征为特征映射的内积。

\subsection{217.4 弱奇异性核}\label{ux5f31ux5947ux5f02ux6027ux6838}

\textbf{定义 217.3}（弱奇异核）：形如 \(K(x,y) = H(x,y) / |x-y|^\alpha\) 的核（\(\alpha < d\)，其中 \(d\) 是空间维数）虽然在对角线上奇异，但在 \(L^2\) 中仍定义紧算子------奇异性由 Hardy-Littlewood-Sobolev 不等式控制。

弱奇异核在应用中极为常见------许多物理问题中的积分方程，如散射理论中的 Lippmann-Schwinger 方程和弹性力学中的边界积分方程，其积分核都具有 \(|x-y|^{-\alpha}\) 型的对角线奇异性。尽管这种奇异性使核在经典意义下无界，但 Hardy-Littlewood-Sobolev 不等式保证了相应的积分算子在 \(L^2\) 上仍是有界的，并且在适当的条件下甚至是紧的。下面这条定理给出了弱奇异核定义紧算子的更一般判据，它表明许多物理上出现的奇异核实际上具有良好的算子性质，Fredholm 理论完全适用。

\textbf{定理 217.5}：若 \(K(x,y)\) 满足 \(\int |K(x,y)|^2 dy < \infty\) a.e. 且 \(\int |K(x,y)|^2 dx < \infty\) a.e.，则 \(K\) 在 \(L^2\) 上是紧算子。Fredholm 择一定理和 Hilbert-Schmidt 定理均成立，且迭代预解核给出级数解。

\textbf{证明}：由假设，\(K\) 定义了一个 \(L^2\) 上的有界积分算子。由 Schur 检验，\(\|K\|_{L^2 \to L^2} \leq (\sup_x \int |K(x,y)|^2 dy)^{1/2} (\sup_y \int |K(x,y)|^2 dx)^{1/2}\)。紧性：取 \(K\) 的逼近序列 \(K_n(x,y)\)（连续且紧支撑），则 \(K_n\) 为 Hilbert-Schmidt 故紧，且 \(\|K - K_n\| \to 0\)。紧算子的极限为紧算子，故 \(K\) 为紧算子。Fredholm 择一定理和 Hilbert-Schmidt 定理对紧算子自然成立，迭代预解核由 Neumann 级数 \(\sum_{n=0}^\infty \lambda^n K^{n+1}\) 给出（\(|\lambda| < \|K\|^{-1}\) 时收敛）。\(\blacksquare\)

\hrulefill

\hrulefill

\hrulefill

\hrulefill

\hrulefill

\newpage

\chapter{06-分析/04-算子代数/02-积分方程.md}\label{ux5206ux679004-ux7b97ux5b50ux4ee3ux657002-ux79efux5206ux65b9ux7a0b.md}

\section{第186章：Fredholm 积分方程（续）}\label{ux7b2c186ux7ae0fredholm-ux79efux5206ux65b9ux7a0bux7eed}

Fredholm 积分方程是积分方程理论的起点，其核心是 Fredholm 择一定理和 Hilbert-Schmidt 定理。Fredholm 积分方程形如 \(u(x) - \lambda \int_a^b K(x,y) u(y) dy = f(x)\)（第二类），其中 \(K(x,y)\) 是平方可积核（Hilbert-Schmidt 核），即 \(\iint_{[a,b]^2} |K(x,y)|^2 dx dy < \infty\)。Fredholm 在 1903 年的开创性工作中，通过将积分方程离散化为线性方程组并取极限，证明了 Fredholm 择一定理：要么齐次方程只有零解（此时非齐次方程对任意 \(f\) 有唯一解），要么齐次方程有非零解（此时非齐次方程可解当且仅当 \(f\) 与共轭齐次方程的所有解正交）。Fredholm 的证明方法引入了\textbf{Fredholm 行列式} \(D(\lambda)\) 和\textbf{Fredholm 第一子式} \(D(x,y;\lambda)\)，它们是 \(\lambda\) 的整函数，其零点恰好是积分算子的特征值。Hilbert 和 Schmidt 在 1904-1910 年间将 Fredholm 的理论推广到对称核 \(K(x,y) = \overline{K(y,x)}\) 的情形，建立了 Hilbert-Schmidt 定理：对称紧算子的特征函数构成 \(L^2\) 的完备正交基，且预解核可由特征函数展开为 \(R(x,y;\lambda) = \sum_{n=1}^\infty \frac{\phi_n(x) \overline{\phi_n(y)}}{\lambda_n - \lambda}\)。这一结果将积分方程理论与无穷维线性代数和正交级数展开统一起来，为量子力学中 Schrodinger 方程的本征函数展开奠定了数学基础。

\textbf{定理 128.1（Hilbert-Schmidt 定理）}：设 \(K(x,y)\) 是 \(L^2([a,b] \times [a,b])\) 上的对称核（\(K(x,y) = \overline{K(y,x)}\)），则积分算子 \((T_K u)(x) = \int_a^b K(x,y) u(y) dy\) 是 \(L^2[a,b]\) 上的紧自伴算子。其特征值 \(\{\lambda_n\}\) 是实数，满足 \(\lambda_n \to 0\)（若有无穷多个），且对应的特征函数 \(\{\phi_n\}\) 构成 \(L^2[a,b]\) 的一组正交基。对任意 \(u \in L^2[a,b]\)，有 Hilbert-Schmidt 展开 \(u = \sum_n \langle u, \phi_n \rangle \phi_n + u_0\)（其中 \(u_0 \in \ker T_K\)）。

\textbf{证明}：由假设，\(K\) 定义了一个 \(L^2\) 上的有界积分算子。由 Schur 检验，\(\|K\|_{L^2 \to L^2} \leq (\sup_x \int |K(x,y)|^2 dy)^{1/2} (\sup_y \int |K(x,y)|^2 dx)^{1/2}\)。紧性：取 \(K\) 的逼近序列 \(K_n(x,y)\)（连续且紧支撑），则 \(K_n\) 为 Hilbert-Schmidt 故紧，且 \(\|K - K_n\| \to 0\)。紧算子的极限为紧算子，故 \(K\) 为紧算子。自伴性由 \(K\) 的对称性直接得出。紧自伴算子的谱定理（Riesz-Schauder 理论）保证了特征值的存在性、实数性和趋于零的性质，以及特征函数的正交性和完备性。Fredholm 择一定理和 Hilbert-Schmidt 定理对紧算子自然成立，迭代预解核由 Neumann 级数 \(\sum_{n=0}^\infty \lambda^n K^{n+1}\) 给出（\(|\lambda| < \|K\|^{-1}\) 时收敛）。\(\blacksquare\)

\subsection{Fredholm积分方程理论的现代发展}\label{fredholmux79efux5206ux65b9ux7a0bux7406ux8bbaux7684ux73b0ux4ee3ux53d1ux5c55}

Fredholm 积分方程理论在 20 世纪经历了从经典分析到泛函分析的深刻转化，其核心成果已融入现代线性算子理论。\textbf{Fredholm 指标理论}是 Fredholm 择一定理的自然推广：若 \(T\) 是 Banach 空间上的 Fredholm 算子（\(\dim \ker T < \infty\)，\(\operatorname{codim} \operatorname{im} T < \infty\)），则其指标 \(\operatorname{ind}(T) = \dim \ker T - \dim \operatorname{coker} T\) 在紧摄动下不变，且关于复合可加。Noether 算子（Fredholm 算子）的理论直接推广了积分方程中 Fredholm 择一定理的核心结构------\(\ker(I - \lambda K)\) 和 \(\ker(I - \lambda K^*)\) 的维数相等。\textbf{Fredholm 算子与 Atiyah-Singer 指标定理}（1963）的联系：椭圆微分算子 \(D\) 在 Sobolev 空间的紧流形上是 Fredholm 算子，其指标 \(\operatorname{ind}(D)\) 由流形的拓扑不变量（A-roof 类、Todd 类）给出。Fredholm 积分方程理论中的预解核 \(R(x,y;\lambda)\) 是函数论中奇异积分算子参数化构造的前身。\textbf{Fredholm 行列式}在量子场论（如 sine-Gordon 模型和 Thirring 模型）中用于计算可积系统的精确谱，其整函数性质与热力学 Bethe 拟设中的非线性积分方程密切相关。\textbf{Fredholm 积分方程与随机过程}：扩散过程的转移概率密度是 Kolmogorov 后退方程的基本解，可通过 Fredholm 积分方程方法求解。在\textbf{逆散射变换}（IST）中，Marchenko 积分方程和 Gel'fand-Levitan 积分方程是 Fredholm 型的，其核由散射数据确定，为 KdV 方程等可积 PDE 的精确求解提供了核心工具。

\hrulefill

\hrulefill

\section{第187章：Volterra 积分方程}\label{ux7b2c187ux7ae0volterra-ux79efux5206ux65b9ux7a0b}

Volterra 积分方程以意大利数学家 Vito Volterra（1860-1940）命名，他在 19 世纪末研究粘弹性材料和种群动力学时首次系统研究了此类方程。与 Fredholm 积分方程（第128章）相比，Volterra 方程的本质特征是积分上限为变量而非固定常数------这一看似微小的差异带来了深刻的结构性变化：Volterra 积分算子具有 Volterra 性质（即其谱仅由零组成），这使得第二类 Volterra 方程总是唯一可解的（通过逐次逼近法或 Picard 迭代），无需像 Fredholm 理论那样讨论 Fredholm 择一性。在数学物理中，Volterra 方程自然地出现在各类随时间演化的系统中------卷积核 \(K(x-y)\) 的情形（卷积型 Volterra 方程）尤为常见，可通过 Laplace 变换化为代数方程求解。本章从 Volterra 算子的有界性和紧性出发，系统讨论逐次逼近法（Neumann 级数）的收敛性、预解核的构造以及卷积型方程的 Laplace 变换方法，并将 Volterra 理论与常微分方程的初值问题建立等价对应。

\subsection{218.1 逐次逼近与 Neumann 级数}\label{ux9010ux6b21ux903cux8fd1ux4e0e-neumann-ux7ea7ux6570}

逐次逼近法是求解 Volterra 积分方程的最基本方法，其核心思想源于 Picard 迭代在常微分方程中的成功应用。与 Fredholm 积分方程不同，Volterra 算子的积分上限为变量 \(x\)，使得积分区域呈三角形而非矩形，这从根本上改变了算子的谱性质。从泛函分析角度看，Volterra 算子 \(V\) 是 \(C[a,b]\) 上的拟幂零紧算子（quasinilpotent compact operator），即其谱半径 \(r(V) = 0\)，谱仅包含 \(\{0\}\)。这一性质是 Volterra 方程对任意 \(\lambda\) 皆可解的根本原因。几何上，Volterra 算子的幂 \(V^n\) 的核在 \(a \leq y \leq x \leq b\) 上具有 \((x-y)^{n-1}/(n-1)!\) 型的增长因子，使得 \(\|V^n\| \leq \|K\|_\infty^n (b-a)^n/n!\)，从而 Neumann 级数 \(\sum \lambda^n V^n\) 的收敛半径为正无穷。这一估计的关键在于 \(V^n\) 的积分区域是 \(n\) 维单纯形 \(\{a \leq y_1 \leq y_2 \leq \cdots \leq y_n \leq x\}\)，其体积为 \((x-a)^n/n!\)，远小于矩形区域 \([a,b]^n\) 的体积 \((b-a)^n\)。正是这种''体积收缩''效应使得 Volterra 算子的谱半径为零，而相应的 Fredholm 算子仅当 \(|\lambda| < 1/\|K\|\) 时 Neumann 级数才收敛。从数值分析的角度，逐次逼近的收敛速度与 \((|\lambda| \|K\|_\infty (b-a))^n/n!\) 成正比，对于中等大小的区间长度和核范数，收敛速度极快。这一方法在 Volterra 型积分-微分方程和时滞微分方程的理论中也有自然推广。

\textbf{定理 218.1}（Volterra 方程的可解性）：若核 \(K(x,y)\) 在 \(a \leq y \leq x \leq b\) 上连续，则对任意 \(\lambda\) 和连续 \(f\)，Volterra 方程 \((I - \lambda V)u = f\) 有唯一连续解。且 Neumann 级数

\[u(x) = f(x) + \sum_{n=1}^{\infty} \lambda^n (V^n f)(x)\]

对任意 \(\lambda\) 收敛（\(V\) 的谱半径为零------拟幂零算子）。

\textbf{证明}：由归纳法，\(|(V^n f)(x)| \leq \|K\|_\infty^n \|f\|_\infty (x-a)^n / n!\)。故 Neumann 级数 \(\sum_{n=0}^\infty \lambda^n V^n f\) 被 \(\|f\|_\infty \sum |\lambda|^n \|K\|_\infty^n (b-a)^n/n! = \|f\|_\infty e^{|\lambda|\|K\|_\infty(b-a)}\) 控制，对任意 \(\lambda\) 一致收敛。因此 \((I-\lambda V)^{-1} = \sum_{n=0}^\infty \lambda^n V^n\) 作为 \(C[a,b]\) 上的有界算子存在，解 \(u = (I-\lambda V)^{-1}f\) 唯一连续。Volterra 算子 \(V\) 的谱半径 \(r(V) = \lim \|V^n\|^{1/n} = 0\)（拟幂零性），故对任意 \(\lambda\) 可逆。\(\blacksquare\)

\subsection{218.2 Volterra 方程与常微分方程}\label{volterra-ux65b9ux7a0bux4e0eux5e38ux5faeux5206ux65b9ux7a0b}

\textbf{定理 218.2}（ODE 初值问题的 Volterra 表示）：\(n\) 阶线性 ODE 初值问题

\[y^{(n)}(x) + a_{n-1}(x) y^{(n-1)}(x) + \cdots + a_0(x) y(x) = g(x)\]

及初值条件 \(y(a) = y_0, \ldots, y^{(n-1)}(a) = y_{n-1}\) 等价于 Volterra 积分方程

\[y(x) = f(x) + \lambda \int_a^x K(x,s) y(s) ds\]

其中 \(f(x)\) 由非齐次项和初值条件决定。这一等价性将 ODE 的存在唯一性理论转化为积分算子的不动点问题。

\textbf{证明}：反复积分 \(y^{(n)}(x) = g(x) - \sum_{k=0}^{n-1} a_k(x)y^{(k)}(x)\)，利用初值条件将 \(y^{(k)}(a)=y_k\) 转化为 \(n\) 重积分。\(n\) 阶积分核为 \((x-s)^{n-1}/(n-1)!\)，将 ODE 写为 \(y(x) = f(x) + \int_a^x K(x,s)y(s)ds\)，其中 \(K(x,s) = \sum_{k=0}^{n-1} a_k(x)\frac{(x-s)^{n-k-1}}{(n-k-1)!}\)，\(f\) 包含 \(g\) 和初值项。此将 ODE 的 Picard 迭代等价于 Volterra 算子的 Neumann 级数。\(\blacksquare\)

\textbf{推论 218.3}：Picard-Lindelöf 定理（Lipschitz 条件下 ODE 解的存在唯一性）是 Volterra 算子压缩性的特例。

\subsection{218.3 Abel 积分方程}\label{abel-ux79efux5206ux65b9ux7a0b}

\textbf{定义 218.4}（Abel 积分方程）：第一类 Volterra 方程

\[\int_0^x \frac{u(y)}{(x-y)^\alpha} dy = f(x), \quad 0 < \alpha < 1\]

（当 \(\alpha = 1/2\) 时等时线问题------tautochrone problem）。

\textbf{定理 218.4}（Abel 反演公式）：

\[u(x) = \frac{\sin(\pi \alpha)}{\pi} \frac{d}{dx} \int_0^x \frac{f(y)}{(x-y)^{1-\alpha}} dy\]

这可以理解为分数阶导数的 Abel 反演

\textbf{证明}：令 \(g(x) = \int_0^x \frac{u(y)}{(x-y)^\alpha}dy\)。这是 Riemann-Liouville 分数阶积分 \(D^{-(1-\alpha)}u\)。Abel 积分方程 \(g = D^{-(1-\alpha)}u\) 的解为 \(u = D^{1-\alpha}g = \frac{1}{\Gamma(\alpha)\Gamma(1-\alpha)}\frac{d}{dx}\int_0^x \frac{g(y)}{(x-y)^{1-\alpha}}dy\)。利用 \(\Gamma(\alpha)\Gamma(1-\alpha) = \pi/\sin(\pi\alpha)\) 即得反演公式。核心是 \(\int_y^x (x-t)^{-\alpha}(t-y)^{\alpha-1}dt = \frac{\pi}{\sin(\pi\alpha)}\)（Beta 积分）。由此 \(D^{-(1-\alpha)}D^{1-\alpha} = I\)，证明反演的合法性和唯一性。\(\blacksquare\) ：\(D^{-\alpha} u = f \implies u = D^\alpha f\)，其中 \(D^\alpha\) 是 Riemann-Liouville 分数阶微积分中的分数阶导数。

\hrulefill

\hrulefill

\hrulefill

\hrulefill

\hrulefill

\hrulefill

\section{第188章：奇异性积分方程}\label{ux7b2c188ux7ae0ux5947ux5f02ux6027ux79efux5206ux65b9ux7a0b}

奇异性积分方程以核函数在积分区域内具有不可积奇异性为特征，在数学物理中扮演着不可替代的角色。与 Fredholm 方程（第128章）和 Volterra 方程（第129章）的正则核不同，奇异积分方程中的积分必须理解为主值积分（Cauchy 主值或 Hadamard 有限部分），这使其 Fredholm 理论失效------奇异积分算子通常不再是紧算子。奇异性积分方程的核心实例是含 Cauchy 核的方程，它自然出现在平面弹性理论、断裂力学（裂纹尖端附近的应力场）以及空气动力学（薄翼理论）中。历史上，Riemann 和 Hilbert 在 19 世纪末关于 Riemann-Hilbert 问题的研究为奇异积分方程奠定了理论基础，而 Muskhelishvili 和 Gakhov 等苏联学派的工作在 20 世纪中叶系统建立了 Cauchy 型积分方程的理论体系。本章从 Plemelj-Sokhotski 公式出发，建立 Cauchy 主值积分的基本性质和 Holder 连续函数空间上的有界性，进而讨论 Riemann-Hilbert 边值问题的求解方法及其与奇异积分方程的等价关系，为数学物理边值问题提供统一的分析框架。

\subsection{219.1 Cauchy 核与 Plemelj 公式}\label{cauchy-ux6838ux4e0e-plemelj-ux516cux5f0f}

\textbf{定义 219.1}（Cauchy 主值积分）：定义在曲线 \(\Gamma\) 上的 Cauchy 主值积分为

\[(\mathcal{H} u)(x) = \text{p.v.} \int_\Gamma \frac{u(y)}{y-x} dy = \lim_{\varepsilon \to 0} \int_{\Gamma \setminus B_\varepsilon(x)} \frac{u(y)}{y-x} dy
\]

Cauchy 核 \(1/(y-x)\) 是一维奇异积分算子的原型，在复分析和调和分析中占据着基础地位。由于该核在 \(y=x\) 处有不可积的奇异性，普通 Lebesgue 积分在此失效，必须引入 Cauchy 主值的概念------通过对称地移除奇异点邻域来定义。Plemelj-Sokhotski 公式是这一理论的第一个核心结果：它精确地量化了 Cauchy 型积分在穿越积分曲线时的跳跃行为，从而将曲线两侧的解析函数边值与曲线上的奇异积分联系起来。这一公式是 Riemann-Hilbert 问题、Muskhelishvili 复变方法以及可积系统中的非线性 steepest descent 方法的出发点。

\textbf{定理 219.1}（Plemelj-Sokhotski 公式）：对 Hölder 连续函数 \(u\) 在光滑曲线 \(\Gamma\) 上，Cauchy 型积分

\[\Phi(z) = \frac{1}{2\pi i} \int_\Gamma \frac{u(y)}{y-z} dy \quad (z \notin \Gamma)\]

在跨越 \(\Gamma\) 时的边值为

\[\Phi^\pm(x) = \pm \frac{1}{2} u(x) + \frac{1}{2\pi i} \text{p.v.} \int_\Gamma \frac{u(y)}{y-x} dy\]

即 \(\Phi^+(x) - \Phi^-(x) = u(x)\)

\textbf{证明}：设 \(\Gamma\) 光滑，\(u\) Hölder连续。\(\Phi(z) = \frac{1}{2\pi i}\int_\Gamma \frac{u(y)}{y-z}dy\) 在 \(\mathbb{C}\setminus\Gamma\) 解析。\(z\to x\in\Gamma\) 时，将积分分解为 \(\int_\Gamma = \int_{\Gamma\setminus B_\varepsilon(x)} + \int_{\Gamma\cap B_\varepsilon(x)}\)。第一项连续跨越 \(\Gamma\)；第二项在半圆路径上的积分给出 \(\pm\frac{1}{2}u(x)\)（正负号取决于逼近方向）。主值积分定义为 \(\lim_{\varepsilon\to 0}\int_{\Gamma\setminus B_\varepsilon(x)}\frac{u(y)}{y-x}dy\)。综合得 Plemelj 公式。\(\blacksquare\) 和 \(\Phi^+(x) + \Phi^-(x) = \frac{1}{\pi i} \text{p.v.} \int \frac{u(y)}{y-x} dy\)。

\subsection{219.2 Riemann-Hilbert 问题}\label{riemann-hilbert-ux95eeux9898}

\textbf{定义 219.2}（标量 Riemann-Hilbert 问题）：求在 \(\Gamma\) 分割的区域内解析的函数 \(\Phi(z)\)，满足边界条件

\[\Phi^+(x) = G(x) \Phi^-(x) + g(x), \quad x \in \Gamma\]

其中 \(G(x)\) 是跃度函数，\(g(x)\) 是给定函数。

\textbf{定理 219.2}（标量 RHP 的可解性）：若 \(G(x) \neq 0\) 且 \(\Gamma\) 是简单闭曲线，令 \(\kappa = \frac{1}{2\pi} [\arg G]_\Gamma\)（\(G\) 的指标 / winding number）。则： - \(\kappa \geq 0\)：齐次问题（\(g \equiv 0\)）有 \(\kappa\) 个线性无关解；非齐次问题对任意 \(g\) 可解 - \(\kappa < 0\)：齐次问题仅有零解；非齐次问题可解需满足 \(-\kappa\) 个相容性条件

\textbf{证明}：令 \(\kappa = \frac{1}{2\pi i}[\log G]_\Gamma\) 为指标。标准化 \(\Phi(z) = z^\kappa e^{\Gamma(z)}\)，其中 \(\Gamma(z) = \frac{1}{2\pi i}\int_\Gamma \frac{\log(G(y)y^{-\kappa})}{y-z}dy\)。由 Plemelj 公式，\(\Phi^\pm\) 满足所需跃度。齐次解空间维数 \(\max(\kappa,0)\)。非齐次情形利用 Cauchy 型积分表示解，相容性条件来自 Cauchy 积分定理：\(\int_\Gamma \frac{g(y)}{\Phi^+(y)}y^k dy = 0\)（\(k=0,\ldots,-\kappa-1\)）当 \(\kappa<0\)。\(\blacksquare\)

\textbf{应用}：Riemann-Hilbert 方法（又名非线性 steepest descent 方法，Deift-Zhou 1993）是求解可积系统------KdV 方程、非线性 Schrödinger 方程、Painlevé 方程等------的长期渐近和普适性结果的核心工具。该方法还严格证明了随机矩阵理论中的 Tracy-Widom 分布公式。

\subsection{219.3 奇异积分算子的有界性}\label{ux5947ux5f02ux79efux5206ux7b97ux5b50ux7684ux6709ux754cux6027}

Plemelj 公式解决了 Cauchy 型积分的边值行为，但一个更为深刻的问题是：Cauchy 奇异积分算子作为 \(L^p\) 空间上的线性算子，其范数有界性如何？这一问题的答案由 Calderón 和 Zygmund 在 1950 年代给出，他们发展了一套完整的奇异积分理论，其核心工具是 Calderón-Zygmund 分解------将任意 \(L^p\) 函数分解为''好部分''（受控于 \(L^2\) 理论）和''坏部分''（受控于核的光滑性估计）之和。这套方法后来成为调和分析处理几乎所有奇异积分算子的标准范式，对偏微分方程和伪微分算子理论产生了革命性影响。

\textbf{定理 219.3}（Calderón-Zygmund 理论）：Cauchy 奇异积分算子 \(\mathcal{H}\)（Hilbert 变换）在 \(L^p(\mathbb{R})\)（\(1 < p < \infty\)）上有界，且在 BMO 空间上有界（\(L^\infty \to \text{BMO}\) 代替）。

\emph{意义}：奇异积分算子的 \(L^p\) 有界性是调和分析和 PDE 理论的基石，为 Mihlin-Hörmander 乘子定理和伪微分算子的 \(L^p\) 连续性提供了统一的技术基础

\textbf{证明}：Hilbert 变换 \(\mathcal{H}f(x) = \text{p.v.}\int \frac{f(y)}{y-x}dy\) 的 Fourier 乘子为 \(-i\operatorname{sgn}(\xi)\)，故在 \(L^2\) 上有界。Calderón-Zygmund 分解：将 \(f\in L^p\) 分解为''好部分''\(g\)（\(|g|\leq c\lambda\)）和''坏部分''\(b=\sum b_k\)（每个 \(b_k\) 支撑在不相交的二进方体 \(Q_k\)上，\(\int b_k=0\)，\(\int|b_k|\leq c\lambda|Q_k|\)）。对好部分用 \(L^2\) 有界性，对坏部分用核的光滑性估计得弱 \((1,1)\) 型，Marcinkiewicz 插值给出 \(1<p<\infty\) 的有界性。\(\blacksquare\) 。

\hrulefill

\hrulefill

\hrulefill

\hrulefill

\hrulefill

\hrulefill

\section{第189章：Wiener-Hopf 积分方程与 Wiener-Hopf 方法}\label{ux7b2c189ux7ae0wiener-hopf-ux79efux5206ux65b9ux7a0bux4e0e-wiener-hopf-ux65b9ux6cd5}

Wiener-Hopf 积分方程形如

\[u(x) - \lambda \int_0^\infty K(x-y) u(y) dy = f(x), \quad x \geq 0\]

核仅依赖于差 \(x-y\)。这类方程出现在衍射理论、断裂力学、排队论和随机过程中。

\textbf{定理 220.1}（Wiener-Hopf 分解）：设 \(K \in L^1(\mathbb{R})\) 且 Fourier 变换 \(\hat{K}(\xi) = \int_{-\infty}^{\infty} K(x) e^{-i\xi x} dx\) 满足 \(1 - \lambda \hat{K}(\xi) \neq 0\) 对所有 \(\xi \in \mathbb{R}\)。则存在分解

\[1 - \lambda \hat{K}(\xi) = \frac{\hat{K}_+(\xi)}{\hat{K}_-(\xi)}\]

其中 \(\hat{K}_+\)（\(\hat{K}_-\)）在上下半平面解析且非零

\textbf{证明}：由 Wiener-Lévy 定理：若 \(\varphi\in\mathcal{A}(\mathbb{R})\)（Fourier 变换的 Wiener 代数）且 \(\varphi\neq 0\) 处处，则 \(1/\varphi\in\mathcal{A}(\mathbb{R})\)。对 \(\varphi(\xi)=1-\lambda\hat{K}(\xi)\)，\(\log\varphi\in\mathcal{A}(\mathbb{R})\)。分解 \(\log\varphi = \log\hat{K}_+ + \log\hat{K}_-\)，其中 \(\log\hat{K}_+\)（\(\log\hat{K}_-\)）的 Fourier 逆变换支撑在 \([0,\infty)\)（\((-\infty,0]\)）。指数化得 \(\varphi = \hat{K}_+/\hat{K}_-\)，且 \(\hat{K}_\pm\) 在各自半平面解析非零。此分解对应于 Toeplitz 算子的可逆性条件。\(\blacksquare\) 。

\textbf{Wiener-Hopf 方法}：通过 Fourier 变换将积分方程转化为 Wiener-Hopf 函数方程，利用上半/下半平面解析性与 Liuville 定理进行因式分解，最终求解。该方法将积分方程转化为复平面上的解析函数因子化问题。

\textbf{定理 220.2}（Krein 推广）：对向量值 Wiener-Hopf 方程，存在算子值符号的 Wiener-Hopf 分解，其指标与 Toeplitz 算子的 Fredholm 性质密切相关

\textbf{证明}：对矩阵值符号 \(\Phi(\xi) \in \mathbb{C}^{n\times n}\)，Wiener-Hopf 分解 \(\Phi = \Phi_+\Phi_-\)（\(\Phi_\pm\) 及其逆在上下半平面解析）可能不存在。存在性等价于 \(\Phi\) 的 Riemann-Hilbert 因子化。指标（偏指标之和）\(\kappa = \frac{1}{2\pi}[\arg\det\Phi]_{-\infty}^\infty\) 是 Toeplitz 算子 \(T_\Phi\) 的 Fredholm 指标。Krein-Gohberg 定理：\(T_\Phi\) 是 Fredholm 算子当且仅当 \(\Phi\) 允许标准 Wiener-Hopf 因子化。\(\blacksquare\) 。

\hrulefill

\hrulefill

\hrulefill

\hrulefill

\hrulefill

\hrulefill

\section{第190章：积分方程的数值方法}\label{ux7b2c190ux7ae0ux79efux5206ux65b9ux7a0bux7684ux6570ux503cux65b9ux6cd5}

积分方程数值解法是将前面各章的理论成果转化为实际计算工具的桥梁。尽管 Fredholm（第128章）、Volterra（第129章）和奇异积分方程（第130章）各自具有截然不同的解析结构，它们的数值处理共享一个核心思想：将无穷维的积分算子投影到有限维子空间上，从而将积分方程转化为线性代数方程组。这一思想的实现路径包括 Galerkin 法（在相同空间中选择试函数和检验函数）、配置法（在离散点集上强制残差为零）、Nystrom 法（用数值求积公式直接离散积分算子）和退化核展开法（用有限秩算子逼近积分核）。每种方法的收敛性分析和误差估计都深刻依赖于原积分算子的紧性以及离散化格式的稳定性。本章从 Galerkin 法的抽象 Hilbert 空间框架出发，逐一建立四种方法的数学基础，并讨论条件数、超收敛现象以及奇异核积分的特殊处理技术（如奇异减法）。数值方法与解析理论的结合使积分方程成为计算物理中求解边值问题的高效工具，特别是在边界元法（BEM）中发挥着核心作用。

\subsection{221.1 Galerkin 法与配置法}\label{galerkin-ux6cd5ux4e0eux914dux7f6eux6cd5}

\textbf{Galerkin 法}：在有限维子空间 \(V_n = \operatorname{span}\{\varphi_1, \ldots, \varphi_n\}\) 中寻求 \(u_n\)，使得残差与所有测试函数 \(\psi_i\) 正交：

\[(u_n - \lambda K u_n - f, \psi_i) = 0\]

得到线性系统 \(\sum_{j=1}^n c_j (\varphi_j - \lambda K \varphi_j, \psi_i) = (f, \psi_i)\)。

Galerkin 法的核心思想是''投影''------将无穷维空间中的积分方程投影到有限维子空间上，要求在子空间内的残差与所有测试函数正交。当测试函数取为 Dirac \(\delta\) 分布时，正交条件退化为在离散点处精确满足方程，这便是配置法（Collocation 法）。配置法在工程应用中特别受欢迎，因为它不需要计算内积，实现更为直接，但在数学分析上需额外验证离散点处的逼近性质。

\textbf{配置法}（Collocation）：要求在配置点 \(\{x_i\}\) 处精确满足方程：

\[u_n(x_i) - \lambda \int K(x_i, y) u_n(y) dy = f(x_i)\]

这是 Galerkin 法取 \(\psi_i = \delta_{x_i}\) 的特例。

无论采用 Galerkin 法还是配置法，核心的理论问题都是收敛性------当子空间维数趋于无穷时，数值解是否收敛到真解？对于紧算子积分方程，答案一般而言是肯定的。下面的定理给出了 Galerkin 法的拟最优误差估计（Céa 型引理），表明数值解的误差被有限维子空间的最佳逼近误差所控制。这一结果依赖于紧算子的一个重要性质：有限秩投影对紧算子的一致逼近。

\textbf{定理 221.1}（Galerkin 法的收敛性）：若 \(K\) 是紧算子且 \(I - \lambda K\) 可逆，则对充分大的 \(n\)，Galerkin 近似 \(u_n\) 唯一存在且 \(\|u - u_n\| \leq C \inf_{v \in V_n} \|u - v\|\)（Céa 引理的对偶

\textbf{证明}：设 \(P_n: H \to V_n\) 为正交投影。Galerkin 方程 \(P_n(I-\lambda K)u_n = P_n f\)。由于 \(K\) 紧，\(I-\lambda K\) 可逆，且 \(\|P_n K - K\| \to 0\)（\(P_n \to I\) 强收敛 + \(K\) 紧）。故对充分大 \(n\)，\(P_n(I-\lambda K)|_{V_n}\) 可逆且逆一致有界。误差 \(u-u_n\) 满足 \((I-\lambda K)(u-u_n) = \lambda K(u-u_n) + (I-P_n)(u_n)\)，利用投影性质得 \(\|u-u_n\| \leq C\inf_{v\in V_n}\|u-v\|\)。\(\blacksquare\) ）。

\subsection{221.2 Nyström 法}\label{nystruxf6m-ux6cd5}

\textbf{Nyström 法}直接在求积规则上离散化积分：

\[\int K(x_i, y) u(y) dy \approx \sum_{j=1}^n w_j K(x_i, y_j) u(y_j)\]

得到线性系统 \((I - \lambda K_n) u_n = f_n\)，其中 \((K_n)_{ij} = w_j K(x_i, y_j)\)。

与 Galerkin 法不同，Nyström 法不涉足投影或正交性，而是直接用数值求积公式替代积分算子中的积分。这种直接离散化的思路在实现上极为简单------只需在配置点 \(\{x_i\}\) 上赋值核函数 \(K(x_i, y_j)\) 并乘以求积权重 \(w_j\) 即可。但简化实现并不牺牲数学严谨性：当求积规则在连续函数上收敛时，由集体紧算子（collectively compact operators）理论保证 Nyström 离散化的一致收敛性。下面的定理给出了这一结论的精确表述。

\textbf{定理 221.2}：若求积规则对连续函数收敛，则 Nyström 法在紧算子框架下是收敛的（逐点且按范数）。收敛阶等于求积误差阶

\textbf{证明}：Nyström 离散化 \(K_n\) 是 \(K\) 的逐点近似：\((K_n)_{ij} = w_j K(x_i,y_j)\)。对连续被积函数，求积误差 \(\|K_n u - K u\|_\infty \to 0\)。由于 \(K\) 紧，\(K_n\) 集体紧（collectively compact），且 \(K_n \to K\) 逐点。由 Anselone 的集体紧算子理论，\((I-\lambda K_n)^{-1}\) 一致有界且 \((I-\lambda K_n)^{-1} \to (I-\lambda K)^{-1}\) 逐点。故 Nyström 解 \(u_n\) 收敛于 \(u\)，收敛阶等于求积公式的代数精度阶。\(\blacksquare\) 。

\subsection{221.3 退化核逼近}\label{ux9000ux5316ux6838ux903cux8fd1}

退化核逼近（degenerate kernel approximation）是求解 Fredholm 积分方程的一种经典且有效的数值方法。其核心思想是将原积分核 \(K(x,y)\) 用一个有限秩的分离核 \(\tilde{K}(x,y) = \sum_{i=1}^m a_i(x) b_i(y)\) 逼近，其中 \(\{a_i\}\) 和 \(\{b_i\}\) 是适当选取的基函数族。当用 \(\tilde{K}\) 替换 \(K\) 后，积分方程 \((I - \lambda \tilde{K})u = f\) 的解 \(u(x)\) 可表示为 \(u(x) = f(x) + \lambda \sum_{i=1}^m c_i a_i(x)\)，其中系数 \(c_i = \int b_i(y) u(y) dy\) 满足 \(m\) 阶线性方程组 \(c_i - \lambda \sum_{j=1}^m (\int b_i(x) a_j(x) dx) c_j = \int b_i(x) f(x) dx\)。这一方法将无穷维积分方程问题转化为有限维线性代数问题，其理论基础是紧算子的有限秩逼近定理：对任意紧算子 \(K\)，存在有限秩算子 \(K_m\) 使得 \(\|K - K_m\| \to 0\)。常用的退化核逼近构造包括：截断的 Karhunen-Loève 展开（基于 \(K\) 的奇异值分解）、插值投影（如分段多项式插值）、以及 Fourier 级数截断。特别地，当 \(K\) 是 Hilbert-Schmidt 算子时，\(\|K - K_m\|_{HS} = \sum_{i=m+1}^\infty s_i^2\)（\(s_i\) 为奇异值），因此奇异值衰减越快，退化核逼近的收敛速度越快。退化核逼近与 Galerkin 法和 Nyström 法有密切联系：若退化核 \(\tilde{K}\) 是 \(K\) 在某个有限维子空间上的正交投影，则退化核法等价于 Galerkin 法；若退化核由插值条件构造，则退化核法与 Nyström 法产生相同的线性系统。在数学物理中，退化核逼近常用于求解具有弱奇异核的边界积分方程，其中核的奇异部分可通过分离变量展开处理。在随机过程和机器学习中，Mercer 定理为退化核逼近提供了理论保证：正定连续核的无穷级数展开可通过截断有限项来一致逼近，这正是核主成分分析（kernel PCA）和 Nyström 低秩逼近的数学基础。

\subsection{221.4 不适定问题与正则化}\label{ux4e0dux9002ux5b9aux95eeux9898ux4e0eux6b63ux5219ux5316}

\textbf{定义 221.1}（不适定问题）：第一类积分方程 \(\int K(x,y) u(y) dy = f(x)\) 在 \(L^2\) 中是不适定的（Hadamard 意义下）：\(K\) 的紧性使其逆无界，数据的小误差可能导致解的剧烈变化。

第一类 Fredholm 积分方程的不适定性根源于紧算子的无穷维零空间和奇异值的零聚点------当 \(K\) 是紧算子时，\(K^{-1}\)（若存在）必定是无界算子，这使得任何对右端项 \(f\) 的小扰动都可能被无限放大。这一问题在科学计算中普遍存在，从医学 CT 成像（Radon 变换的反演）到地球物理反演（重力异常反演）均属此类。Tikhonov 正则化通过在原目标泛函中增加解的范数罚项 \(\alpha \|u\|^2\)，将不适定问题转化为良定的极小化问题。这里 \(\alpha > 0\) 是正则化参数，它在数据拟合与解的平滑性之间取得平衡。

\textbf{Tikhonov 正则化}：用正则化泛函

\[u_\alpha = \arg\min_u \left\{ \|K u - f\|_{L^2}^2 + \alpha \|u\|_{L^2}^2 \right\}\]

代替原问题，其中 \(\alpha > 0\) 是正则化参数。等价于求解

\[(K^* K + \alpha I) u_\alpha = K^* f
\]

Tikhonov 正则化形式简洁，但其有效性取决于正则化参数 \(\alpha\) 的选取是否与数据噪声水平 \(\delta\) 相匹配。若 \(\alpha\) 过大，解被过度平滑化，丢失了细节信息；若 \(\alpha\) 过小，噪声污染仍然无法控制。下面的定理给出了正则化解收敛到真解的精确条件：\(\alpha\) 必须随噪声水平 \(\delta\) 按一定的速率趋于零，并且这一速率取决于真解的''光滑度''（由源条件描述）。Morozov 偏差原理则为实际计算中选取 \(\alpha\) 提供了可操作的准则。

\textbf{定理 221.3}（正则化的收敛性）：若 \(\alpha \to 0\) 且 \(\alpha\) 以适当速率（依赖于噪声水平 \(\delta = \|f - f^\delta\|\)）趋于零，则 Tikhonov 解 \(u_\alpha^\delta\) 当 \(\delta \to 0\) 时收敛于真正解（Morozov 偏差原理

\textbf{证明}：Tikhonov 解 \(u_\alpha^\delta = (K^*K+\alpha I)^{-1}K^*f^\delta\)。误差分解：\(u_\alpha^\delta - u = (K^*K+\alpha I)^{-1}K^*(f^\delta-f) + [(K^*K+\alpha I)^{-1}K^*K - I]u\)。第一项 \(\leq \delta/\sqrt{\alpha}\)（噪声传播），第二项 \(\to 0\) 当 \(\alpha\to 0\)（由 \(K^*K\) 的谱表示和 \(u\in\mathcal{R}((K^*K)^\nu)\) 的正则性）。取 \(\alpha \sim \delta^{2/(2\nu+1)}\) 得最优收敛率 \(O(\delta^{2\nu/(2\nu+1)})\)。Morozov 偏差原理：选 \(\alpha\) 使 \(\|Ku_\alpha^\delta - f^\delta\| = \tau\delta\)（\(\tau>1\)），则当 \(\delta\to 0\) 时 \(u_\alpha^\delta \to u\)。\(\blacksquare\) 给出 \(\alpha\) 的选取准则）。

\subsection{积分方程数值方法的现代进展}\label{ux79efux5206ux65b9ux7a0bux6570ux503cux65b9ux6cd5ux7684ux73b0ux4ee3ux8fdbux5c55}

积分方程数值方法在 20 世纪末和 21 世纪初取得了重大进展，边界元方法（BEM）和快速算法是其核心驱动力。\textbf{边界元方法}（Boundary Element Method）将偏微分方程边值问题转化为边界积分方程，利用 Green 函数和位势理论将计算域从 \(d\) 维降至 \(d-1\) 维，极大减少了离散化自由度。边界积分算子是\textbf{伪微分算子}（pseudodifferential operator），其经典符号和奇异集的刻画在快速多极子方法（FMM）的自适应压缩中发挥关键作用。\textbf{快速多极子方法}（Fast Multipole Method, Rokhlin-Greengard 1987）是 Nyström 法的高效实现：将 \(N\) 个源点对 \(N\) 个目标点的全交互计算从 \(O(N^2)\) 降至 \(O(N \log N)\) 或 \(O(N)\)，通过将远场相互作用用多极展开（multipole expansion）近似，再用局部展开（local expansion）传递。FMM 是近二十年来计算科学的重大突破之一，被列为 20 世纪十大算法之一。\textbf{谱方法}（spectral methods）在积分方程中的应用：对光滑核，积分算子的特征值呈指数衰减，截断 Galerkin 矩阵的谱误差显著小于有限元方法。\textbf{卷积型积分方程}的快速算法利用 FFT 加速矩阵-向量乘法：通过 Toeplitz 矩阵结构和循环卷积定理，将矩阵-向量乘法的复杂度从 \(O(N^2)\) 降至 \(O(N \log N)\)。\textbf{分层矩阵}（hierarchical matrices, \(\mathcal{H}\)-矩阵和 \(\mathcal{H}^2\)-矩阵）基于积分算子的核函数在远场条件下具有低秩可分离性，将满矩阵压缩为 \(O(N \log N)\) 存储量和运算量，适用于大规模边界积分方程的直接求解。\textbf{正则化理论}的最新进展：Tikhonov 正则化之外，\textbf{迭代正则化}（Landweber 迭代、共轭梯度法半收敛）和\textbf{变分正则化}（Bregman 迭代、总变差正则化）为不适定积分方程提供了更灵活的反演工具。\textbf{稀疏性正则化}（\(\ell_1\) 正则化）在压缩感知框架下，对解具有稀疏表示的积分方程（如源定位问题）提供了超分辨率的恢复能力。

\hrulefill

\hrulefill

\section{第191章：积分方程在数学物理中的应用}\label{ux7b2c191ux7ae0ux79efux5206ux65b9ux7a0bux5728ux6570ux5b66ux7269ux7406ux4e2dux7684ux5e94ux7528}

积分方程方法是数学物理问题的三大解法之一（另两种为微分方程方法和变分方法）。

\textbf{定理 222.1}（位势理论与边界积分方程）：Laplace 方程 \(\Delta u = 0\) 的 Dirichlet 边值问题 \(u|_{\partial \Omega} = g\) 等价于边界积分方程

\[\frac{1}{2} \sigma(x) + \int_{\partial \Omega} \frac{\partial G}{\partial n_y}(x,y) \sigma(y) dS_y = g(x)\]

（双层位势表示）。对应的边界积分算子在 \(L^2(\partial \Omega)\) 上是紧的（当 \(\partial \Omega\) 光滑时），Fredholm 择一定理给出 Dirichlet 问题解的存在唯一性

\textbf{证明}：双层位势 \(u(x) = \int_{\partial\Omega} \sigma(y)\frac{\partial G}{\partial n_y}(x,y)dS_y\) 在 \(\Omega\) 内调和。跨越边界时跳跃关系：\(u^+(x) = \frac{1}{2}\sigma(x) + \int K(x,y)\sigma(y)dy\)，\(u^-(x) = -\frac{1}{2}\sigma(x) + \int K(x,y)\sigma(y)dy\)（\(K\) 为紧算子）。Dirichlet 条件 \(u^+|_{\partial\Omega}=g\) 给出 \((\frac{1}{2}I+K)\sigma=g\)。当 \(\Omega\) 光滑时 \(K\) 是具弱奇异核的紧算子，Fredholm 择一定理保证解的存在唯一性。\(\blacksquare\) 。

\textbf{电磁散射理论}：Maxwell 方程的外问题可化为积分方程（电场/磁场积分方程 EFIE/MFIE）；量子力学中 Lippmann-Schwinger 方程 \(|\psi\rangle = |\psi_0\rangle + G_0 V |\psi\rangle\) 是散射理论的 Fredholm 积分方程。

\textbf{积分-微分方程}：辐射输运方程（Boltzmann 方程的特殊形式）同时含积分和微分算子，通常用离散纵坐标法（\(S_N\) 法）加源迭代求解。

\hrulefill

\emph{本卷建立了积分方程理论的完整体系：Fredholm 积分方程（Fredholm择一定理、Hilbert-Schmidt展开与Mercer定理）、Volterra 积分方程（Neumann级数与拟幂零性、Abel反演公式）、奇异性积分方程（Cauchy主值积分、Plemelj-Sokhotski公式、Riemann-Hilbert问题与非线性陡降方法）、Wiener-Hopf 积分方程与Wiener-Hopf分解，以及Galerkin法、Nyström法、Tikhonov正则化等数值方法。共 6 章。}

\subsection{积分方程理论的历史与展望}\label{ux79efux5206ux65b9ux7a0bux7406ux8bbaux7684ux5386ux53f2ux4e0eux5c55ux671b}

积分方程理论是 20 世纪数学发展的核心动力之一，其历史可追溯至 19 世纪末。\textbf{Abel 积分方程}（1823）是第一个被明确研究的积分方程------Abel 发现了等时曲线问题和力学中积分方程之间的联系，并给出了经典的 Abel 反演公式。\textbf{Volterra}（1884-1896）系统研究了积分上限为变量的积分方程，引入了 Neumann 级数解法，奠定了 Volterra 积分方程的理论基础。\textbf{Fredholm}（1900-1903）的突破性工作引入了 Fredholm 行列式和 Fredholm 择一定理，将积分方程转化为线性代数问题，他的理论直接启发了 Hilbert 对积分方程谱理论的研究。\textbf{Hilbert}（1904-1910）在 Fredholm 工作的基础上发展了 Hilbert-Schmidt 展开定理，并将积分方程与无穷维二次型、变分法和微分方程联系起来，这些工作最终催生了\textbf{Hilbert 空间的抽象理论}（1906-1912），成为泛函分析的诞生地。\textbf{Riesz}（1907-1918）将 Fredholm-Hilbert 理论抽象为紧算子的谱理论，这是现代泛函分析的基石。20 世纪下半叶，积分方程理论在多个方向获得了新的生命力：\textbf{奇异积分方程}和\textbf{Wiener-Hopf 方法}（Wiener-Hopf 1931）在电磁散射和衍射理论中发挥了核心作用；\textbf{边界积分方程方法}（BIM）成为计算力学中与有限元方法并驾齐驱的数值方法；\textbf{逆散射变换}（Gardner-Greene-Kruskal-Miura 1967）将非线性发展方程的初值问题转化为线性积分方程（Gel'fand-Levitan-Marchenko 方程），开创了可积系统理论的新纪元。21 世纪以来，积分方程在\textbf{机器学习}（如积分算子核方法、神经积分方程）、\textbf{量子计算}（如量子 Fredholm 算法）和\textbf{分数阶微积分}（如分数阶积分方程的适定性）中找到了新的应用场景，展现出持久的数学生命力。

\newpage

\chapter{06-分析/04-算子代数/03-拓扑向量空间与分布.md}\label{ux5206ux679004-ux7b97ux5b50ux4ee3ux657003-ux62d3ux6251ux5411ux91cfux7a7aux95f4ux4e0eux5206ux5e03.md}

\section{第192章：拓扑向量空间}\label{ux7b2c192ux7ae0ux62d3ux6251ux5411ux91cfux7a7aux95f4}

拓扑向量空间（Topological Vector Space, TVS）是 Bourbaki 单独成卷的核心分析分支。它推广了 Banach 和 Hilbert 空间的框架，容纳了分布理论（Schwartz）、解析函数空间（Montel空间）和 PDE 中常见的非可赋范空间。

\subsection{拓扑向量空间基础}\label{ux62d3ux6251ux5411ux91cfux7a7aux95f4ux57faux7840}

\textbf{定义 134.1}（拓扑向量空间）：域 \(\mathbb{K}\)（\(\mathbb{R}\) 或 \(\mathbb{C}\)）上的\textbf{拓扑向量空间}（TVS）\(E\) 是一向量空间，装备了 Hausdorff 拓扑，满足：

\begin{enumerate}
\def\labelenumi{\arabic{enumi}.}
\tightlist
\item
  加法 \(+ : E \times E \to E\) 连续（\(E \times E\) 配积拓扑）；
\item
  数乘 \(\cdot : \mathbb{K} \times E \to E\) 连续（\(\mathbb{K} \times E\) 配积拓扑）。
\end{enumerate}

这意味着平移 \(x \mapsto x + x_0\) 和伸缩 \(x \mapsto \lambda x\)（\(\lambda \neq 0\)）均为同胚，因此原点的邻域基完全决定了拓扑结构。

\textbf{定义 134.2}（局部凸空间）：若 TVS \(E\) 的原点有由凸集组成的邻域基，则 \(E\) 为\textbf{局部凸空间}（LCS）。等价地，其拓扑可由一族半范数 \(\{p_\alpha\}_{\alpha \in A}\) 生成：原点的邻域基为

\[V(\alpha_1, \ldots, \alpha_n ; \varepsilon) = \{x \in E : p_{\alpha_i}(x) < \varepsilon, \; i = 1, \ldots, n\}\]

Banach 空间是局部凸的（单个范数）；\(C^\infty(\mathbb{R})\) 也是（半范数族 \(p_{K,n}(f) = \sup_{x \in K} |f^{(n)}(x)|\)，\(K\) 紧）。

\textbf{定义 134.3}（Minkowski 泛函）：设 \(A \subseteq E\) 为含原点的凸吸收集。其 \textbf{Minkowski 泛函}（或度规泛函）定义为

\[p_A(x) = \inf\{t > 0 : x \in tA\}\]

\(p_A\) 满足 \(p_A(x+y) \leq p_A(x) + p_A(y)\) 和 \(p_A(tx) = t p_A(x)\)（\(t \geq 0\)），即 \(p_A\) 为次线性泛函。若 \(A\) 还是均衡的，则 \(p_A\) 为半范数。局部凸拓扑中的每个连续半范数均可表示为某凸均衡邻域的 Minkowski 泛函。

\textbf{定理 134.1}（Kolmogorov 可赋范准则）：TVS \(E\) 可赋范当且仅当它是局部凸的且原点具有有界邻域（即存在原点的邻域 \(U\) 使得对任何邻域 \(V\)，有 \(t > 0\) 使 \(U \subset tV\)）。这解释了为何许多函数空间需要比范数拓扑更精细的结构

\textbf{证明}：必要性：可赋范的局部凸空间中单位球 \(B=\{x:\|x\|<1\}\) 是有界邻域。充分性：设 \(U\) 为有界均衡凸邻域，其 Minkowski 泛函 \(p_U\) 是范数（有界性保证 \(p_U(x)=0 \iff x=0\)）。由局部凸性，\(p_U\) 的拓扑与原始拓扑等价，故空间可赋范。有界性的关键：若拓扑不能由单个有界集控制，则需无限多半范数。\(\blacksquare\) 。

\textbf{定义 134.4}（桶空间）：\(E\) 中的\textbf{桶}（barrel）是吸收所有点的、绝对凸的闭子集。\textbf{桶空间}（barrelled space）是每个桶都是原点邻域的局部凸空间。所有 Fréchet 空间是桶空间。\textbf{配对引理}（Banach-Steinhaus 定理的 TVS 推广）：在桶空间中，逐点有界的连续线性算子族是等度连续的。

\textbf{定理 134.2}（Mackey-Arens 定理）：设 \((E, F)\) 是对偶配对（即 \(F \subset E^*\) 分离 \(E\) 的点）。\(E\) 上所有与 \(F\) 相容的（即对偶空间恰为 \(F\) 的）局部凸拓扑中，最弱的是弱拓扑 \(\sigma(E, F)\)，最强的是 Mackey 拓扑 \(\tau(E, F)\)（由 \(F\) 中所有 \(\sigma(F, E)\)-紧绝对凸集的极生成）。这两者之间的所有相容局部凸拓扑给出相同的连续对偶

\textbf{证明}：相容拓扑条件是 \(E\) 上的局部凸拓扑 \(\tau\) 满足 \((E,\tau)^* = F\)。最弱相容拓扑 \(\sigma(E,F)\) 由半范数 \(p_f(x)=|\langle x,f\rangle|\)（\(f\in F\)）生成。最强相容拓扑 \(\tau(E,F)\) 由 \(F\) 中所有 \(\sigma(F,E)\)-紧绝对凸集的极生成。Mackey-Arens：\(\tau\) 是相容的当且仅当 \(\sigma(E,F) \subseteq \tau \subseteq \tau(E,F)\)。证明基于双极定理：\(U^{\circ\circ} = \overline{\operatorname{conv}}^{\sigma(E,F)}(U)\) 对 \(E\) 中原点邻域 \(U\) 成立。\(\blacksquare\) 。

\subsection{Fréchet 空间与 LF 空间}\label{fruxe9chet-ux7a7aux95f4ux4e0e-lf-ux7a7aux95f4}

\textbf{定义 134.5}（Fréchet 空间）：\textbf{Fréchet 空间}是完全的可度量化局部凸空间。等价地，其拓扑可由一列递增的半范数 \(\{p_n\}\) 生成且在此半范数列下完备。例子：\(C^\infty[0, 1]\)、\(\mathcal{S}(\mathbb{R}^n)\)（Schwartz 速降函数空间）、全纯函数空间 \(\mathcal{O}(U)\)。

\textbf{定理 134.3}（开映射定理的 Fréchet 推广）：Fréchet 空间之间的连续线性满射必定是开的。特别地，连续线性双射自动是同胚（有界逆定理）

\textbf{证明}：设 \(T: E\to F\) 为 Fréchet 空间间的连续线性满射。由 Baire 纲定理，\(T(B_E)\) 包含 \(F\) 中原点的邻域。对 \(\varepsilon>0\)，递归构造序列使 \(\overline{T(B(\varepsilon/2^k))} \subseteq T(B(\varepsilon/2^{k-1}))\)。利用完备性，\(T(B(\varepsilon))\) 包含 \(B(\delta)\) 对某 \(\delta>0\)，故 \(T\) 是开的。连续双射时有界逆 \(T^{-1}\) 自动连续。\(\blacksquare\) 。

\textbf{定理 134.4}（闭图像定理的 Fréchet 推广）：从桶空间到 Fréchet 空间的线性映射若其图像是闭的，则连续。这在 PDE 的闭算子理论中至关重要

\textbf{证明}：设 \(E\) 桶空间，\(F\) Fréchet，\(T: E\to F\) 线性且图像 \(G(T)\) 闭。\(E\) 上定义图范数 \(\|x\|_T = \|x\|_E + \|Tx\|_F\)，\((E,\|\cdot\|_T)\) 是完备的（因 \(G(T)\) 闭）。恒等映射 \(I: (E,\|\cdot\|_T)\to (E,\|\cdot\|_E)\) 连续双射，由开映射定理其逆连续，故 \(\|Tx\|_F \leq C\|x\|_E\)。\(\blacksquare\) 。

\textbf{定义 134.6}（LF 空间与测试函数空间 \(\mathcal{D}(\Omega)\)）：\textbf{LF 空间}（严格归纳极限）是一列 Fréchet 空间 \(E_1 \subset E_2 \subset \cdots\)（每个嵌入紧或等距）的归纳极限，装备使所有嵌入连续的最终局部凸拓扑。

最重要的例子：\textbf{测试函数空间} \(\mathcal{D}(\Omega) = C_c^\infty(\Omega)\)，其中 \(\Omega \subset \mathbb{R}^n\) 为开集。令 \(K_1 \subset K_2 \subset \cdots\) 为 \(\Omega\) 的紧穷竭，定义 \(C_{K_j}^\infty = \{ \varphi \in C^\infty(\Omega) : \operatorname{supp} \varphi \subseteq K_j \}\)。每个 \(C_{K_j}^\infty\) 为由半范数族 \(p_{K_j, m}(\varphi) = \sup_{|\alpha| \leq m} \sup_{x \in K_j} |D^\alpha \varphi(x)|\) 生成的 Fréchet 空间。\(\mathcal{D}(\Omega)\) 配以相对于 \(\{C_{K_j}^\infty\}\) 的严格归纳极限拓扑（LF 拓扑）。此拓扑下，序列 \(\varphi_n \to 0\) 当且仅当存在紧集 \(K\) 使 \(\operatorname{supp} \varphi_n \subseteq K\) 且 \(\varphi_n\) 及其所有导数在 \(K\) 上一致收敛到 \(0\)。

\(\mathcal{D}(\Omega)\) 不是可度量化的------这是 LF 空间区别于 Fréchet 空间的关键。

\subsection{Krein-Milman 定理与 Choquet 理论}\label{krein-milman-ux5b9aux7406ux4e0e-choquet-ux7406ux8bba}

\textbf{定理 134.5}（Krein-Milman 定理）：局部凸空间中的紧凸集 \(K\) 是其端点集的闭凸包：\(K = \overline{\operatorname{conv}}(\operatorname{ext}(K))\)。每个紧凸集中的点都是端点的''积分''。

\textbf{证明}：设 \(E = \operatorname{ext}(K)\)。若 \(K \neq \overline{\operatorname{conv}}(E)\)，则存在 \(x\in K\setminus\overline{\operatorname{conv}}(E)\)。由 Hahn-Banach 分离定理，存在连续线性泛函 \(f\) 和 \(\alpha\) 使 \(f(x) > \alpha \geq \sup_{y\in E}f(y)\)。\(f\) 在紧集 \(K\) 上取最大值 \(\max_K f = f(x_0)\)，\(x_0\in K\)。\(x_0\) 必为端点：若 \(x_0 = ty+(1-t)z\)（\(y\neq z\)），则 \(f(y)=f(z)=\max\)，但由分离性质 \(y\) 或 \(z\) 不能在 \(\overline{\operatorname{conv}}(E)\) 之外，矛盾。然而 \(x_0\in E\) 给出 \(f(x_0) \leq \alpha < f(x)\)------矛盾。\(\blacksquare\)

\textbf{定理 134.6}（Choquet 定理 / Choquet-Bishop-de Leeuw 积分表示）：设 \(K\) 是局部凸空间中的紧凸集。对任意 \(x \in K\)，存在 \(K\) 的端点集 \(\operatorname{ext}(K)\) 上的概率测度 \(\mu\)（Choquet测度），使得对任意连续仿射函数 \(f: K \to \mathbb{R}\)，有 \(f(x) = \int_{\operatorname{ext}(K)} f \, d\mu\)。此即每个点可表示为端点的重心

\textbf{证明}：对 \(x\in K\)，考虑仿射函数 \(A(K)\)。由 Krein-Milman，\(x = \lim_n \sum_i \lambda_{n,i}e_{n,i}\)（\(e_{n,i}\in\operatorname{ext}(K)\) 的凸组合）。定义离散测度 \(\mu_n = \sum_i \lambda_{n,i}\delta_{e_{n,i}}\)，则 \(\int f d\mu_n \to f(x)\) 对 \(f\in A(K)\)。由 Prohorov 定理（\(K\) 紧），\(\mu_n\) 弱收敛子列至 \(\mu\)。由端点集 \(\operatorname{ext}(K)\) 是 \(G_\delta\) 集（度量情形），\(\mu(\operatorname{ext}(K)) = 1\)。取 \(f\) 为恒等映射的坐标泛函即得重心表示。\(\blacksquare\) 。

\emph{意义}：Krein-Milman 和 Choquet 理论是泛函分析''几何''面的顶峰------将凸集的几何结构（端点）与测度论表示（重心测度）联系起来，在量子力学的 C*-代数态空间表示和遍历理论中极为重要。

\hrulefill

\section{第193章：谱理论（补充）}\label{ux7b2c193ux7ae0ux8c31ux7406ux8bbaux8865ux5145}

Bourbaki 将谱理论单独成卷。但教材中紧算子的谱理论向无界算子的推广是量子力学数学基础的必要部分。

\subsection{谱测度与谱积分}\label{ux8c31ux6d4bux5ea6ux4e0eux8c31ux79efux5206}

\textbf{定义 135.1}（谱测度 / 恒等分解）：Hilbert 空间 \(H\) 上的\textbf{谱测度}（projection-valued measure）\(E\) 是将 Borel 集 \(\Delta \subset \mathbb{R}\) 映为 \(H\) 上的正交投影 \(E(\Delta)\) 的映射，满足： 1. \(E(\varnothing) = 0\)，\(E(\mathbb{R}) = I\) 2. \(E(\Delta_1 \cap \Delta_2) = E(\Delta_1) E(\Delta_2)\) 3. 对互不相交的 \(\{\Delta_n\}\)，\(E(\bigcup \Delta_n) = \sum E(\Delta_n)\)（强算子拓扑收敛）

\textbf{定理 135.1}（谱定理------无界自伴算子，von Neumann 1929）：对 Hilbert 空间 \(H\) 上的任意（可能无界的）自伴算子 \(A: \mathcal{D}(A) \to H\)（\(\mathcal{D}(A)\) 稠密），存在唯一的谱测度 \(E\) 使得

\[A = \int_{\mathbb{R}} \lambda \, dE(\lambda)\]

积分在 \(\mathcal{D}(A) = \{x \in H : \int \lambda^2 \, d\langle E(\lambda)x, x\rangle < \infty\}\) 上理解。这是有限维对角化的无限维推广

\textbf{证明}：通过 Cayley 变换 \(U = (A-iI)(A+iI)^{-1}\)，\(A\) 的自伴性等价于 \(U\) 的酉性。对酉算子 \(U\)，由 Gelfand-Naimark 定理（交换 C*-代数版本），存在 \(\sigma(U)\subseteq S^1\) 上的谱测度 \(E_U\) 使 \(U = \int_{S^1} \lambda dE_U(\lambda)\)。逆 Cayley 变换 \(\lambda \mapsto i(1+\lambda)/(1-\lambda)\) 将 \(S^1\) 映到 \(\mathbb{R}\cup\{\infty\}\)，给出 \(A\) 的谱测度 \(E\) 满足 \(A = \int_{\mathbb{R}} t dE(t)\)。定义域 \(\mathcal{D}(A) = \{x : \int t^2 d\langle E(t)x,x\rangle < \infty\}\) 自动稠密。\(\blacksquare\) 。

\textbf{定义 135.2}（谱积分 / 算子的函数演算）：对 Borel 可测函数 \(f: \mathbb{R} \to \mathbb{C}\)，定义

\[f(A) = \int_{\mathbb{R}} f(\lambda) \, dE(\lambda)\]

定义域为 \(\{x : \int |f|^2 d\langle E x, x\rangle < \infty\}\)。特别地，\(e^{itA}\)（由 \(f(\lambda) = e^{it\lambda}\)）定义了酉群。

\textbf{定理 135.2}（Stone 定理，1932）：Hilbert 空间上的强连续单参数酉群 \(\{U(t)\}_{t \in \mathbb{R}}\) 与（可能无界的）自伴算子 \(A\) 一一对应（经 \(U(t) = e^{itA}\)）。这是量子力学中''能量算子=时间平移生成元''的数学表述

\textbf{证明}：\(\Rightarrow\)：对酉群 \(U(t)\)，定义 \(\mathcal{D}(A) = \{x : \lim_{t\to 0}(U(t)x-x)/t \text{ 存在}\}\)，\(Ax = i\lim_{t\to 0}(U(t)x-x)/t\)。由 Stone 定理，\(A\) 是稠定自伴算子。\(\Leftarrow\)：对自伴 \(A\)，由谱定理定义 \(U(t) = e^{itA} = \int e^{it\lambda}dE(\lambda)\)，此为强连续酉群。唯一性由 \(U(t)\) 的生成元唯一决定。\(\blacksquare\) 。

\subsection{谱重数与本质谱}\label{ux8c31ux91cdux6570ux4e0eux672cux8d28ux8c31}

\textbf{定义 135.3}（谱重数 / 重数函数）：自伴算子 \(A\) 的谱重数函数 \(m(\lambda)\) 是在 \(\lambda\) 附近谱测度的''秩''。若 \(m(\lambda) \equiv 1\)（几乎处处），\(A\) 称为具有\textbf{单谱}。

\textbf{定义 135.4}（本质谱与 Weyl 定理）：\(A\) 的\textbf{本质谱} \(\sigma_{\text{ess}}(A)\) 是去除了有限重数孤立特征值后的剩余谱。\textbf{Weyl 准则}：\(\lambda \in \sigma_{\text{ess}}(A)\) 当且仅当存在 Weyl 序列 \((x_n) \subset \mathcal{D}(A)\) 满足 \(\|x_n\| = 1\)，\(x_n \rightharpoonup 0\)（弱收敛）且 \(\|(A - \lambda I)x_n\| \to 0\)。本质谱在紧摄动下不变（Weyl 定理）。

\textbf{定理 135.3}（Fredholm 算子与本质谱）：\(A - \lambda I\) 是 Fredholm 算子（核有限维、值域闭、余核有限维）当且仅当 \(\lambda \notin \sigma_{\text{ess}}(A)\)。Fredholm 指标 \(\operatorname{ind}(A - \lambda I) = \dim \ker - \dim \operatorname{coker}\) 在 \(\sigma_{\text{ess}}(A)\) 的连通分支上为常数

\textbf{证明}：由 Atkinson 定理，\(T\) 是 Fredholm 算子当且仅当 \(T\) 在 Calkin 代数 \(\mathcal{B}(H)/\mathcal{K}(H)\) 中可逆。\(\sigma_{\text{ess}}(A) = \sigma_{\text{Calkin}}(\pi(A))\)。Fredholm 指标是 Calkin 代数的 \(K_0\) 不变量：\(\operatorname{ind}\) 在 Fredholm 算子的连通分支上为常数。\(A-\lambda I\) 的指标在 \(\mathbb{C}\setminus\sigma_{\text{ess}}(A)\) 的连通分支上不变。\(\blacksquare\) 。

\hrulefill

\section{第194章：分布理论}\label{ux7b2c194ux7ae0ux5206ux5e03ux7406ux8bba}

分布理论（Distribution Theory / Generalized Functions）由 Laurent Schwartz 于 1945-1951 年创立（获 1950 年 Fields 奖），将''函数''概念推广为线性泛函，使得诸如 Dirac 函数、函数的广义导数、局部可积函数的弱导数等获得严格意义。

\subsection{Schwartz 空间与缓增分布}\label{schwartz-ux7a7aux95f4ux4e0eux7f13ux589eux5206ux5e03}

\textbf{定义 136.1}（速降函数空间 / Schwartz 空间）：\(\mathbb{R}^n\) 上的 \textbf{Schwartz 空间} \(\mathcal{S}(\mathbb{R}^n)\) 是所有光滑函数 \(f \in C^\infty(\mathbb{R}^n)\) 的集合，满足对任意多重指标 \(\alpha, \beta\)，

\[\sup_{x \in \mathbb{R}^n} |x^\alpha D^\beta f(x)| < \infty\]

即 \(f\) 及其所有导数在无穷远处比任何多项式的倒数下降得更快。\(\mathcal{S}(\mathbb{R}^n)\) 是 Fréchet 空间（由可数族半范数 \(\|f\|_{k, m} = \sup_{|\alpha| \leq k, |\beta| \leq m} |x^\alpha D^\beta f|\) 生成）。

\textbf{定义 136.2}（缓增分布）：\textbf{缓增分布}空间 \(\mathcal{S}'(\mathbb{R}^n)\) 是 \(\mathcal{S}(\mathbb{R}^n)\) 上的连续线性泛函全体，装备弱* 拓扑。缓增分布包含了所有 \(L^p(\mathbb{R}^n)\) 函数（\(1 \leq p \leq \infty\)）、所有多项式、缓增测度和 Dirac \(\delta\) 函数列。

\textbf{命题 136.1}（嵌入与基本运算）： - \(L^p(\mathbb{R}^n) \subset \mathcal{S}'(\mathbb{R}^n)\)：通过 \(T_f(\varphi) = \int f \varphi \, dx\) 嵌入（对缓增增长的局部可积函数）。 - \textbf{分布的导数}：对 \(T \in \mathcal{S}'\)，定义 \(\langle D^\alpha T, \varphi \rangle = (-1)^{|\alpha|} \langle T, D^\alpha \varphi \rangle\)。每阶分布导数都存在------分布无穷次可微。 - \textbf{Fourier 变换}：\(\mathcal{S}\) 上的 Fourier 变换 \(\mathcal{F}\) 是同构，经对偶扩展到 \(\mathcal{S}'\)：\(\langle \mathcal{F}T, \varphi \rangle = \langle T, \mathcal{F}\varphi \rangle\)。

\textbf{例}：\(T = \delta\)（Dirac \(\delta\)）：\(\langle \delta, \varphi \rangle = \varphi(0)\)。其导数为 \(\langle \delta', \varphi \rangle = -\varphi'(0)\)。Heaviside 阶梯函数 \(H(x) = \mathbf{1}_{[0,\infty)}(x)\) 的分布导数是 \(\delta\)：\(H' = \delta\)。

\subsection{分布的局部理论与紧支撑分布}\label{ux5206ux5e03ux7684ux5c40ux90e8ux7406ux8bbaux4e0eux7d27ux652fux6491ux5206ux5e03}

\textbf{定义 136.3}（分布空间 \(\mathcal{D}'(\Omega)\)）：\(\Omega \subset \mathbb{R}^n\) 为开集。\(\Omega\) 上的（一般）\textbf{分布}空间 \(\mathcal{D}'(\Omega)\) 是测试函数空间 \(\mathcal{D}(\Omega) = C_c^\infty(\Omega)\)（装备 LF 拓扑）上的连续线性泛函。任何局部可积函数通过 \(T_f(\varphi) = \int_\Omega f \varphi \, dx\) 定义分布。

\textbf{定义 136.4}（分布的支集）：设 \(T \in \mathcal{D}'(\Omega)\)。若 \(T\) 在开子集 \(U \subseteq \Omega\) 上为零（即 \(\langle T, \varphi \rangle = 0\) 对所有 \(\varphi \in C_c^\infty(U)\)），称 \(T\) 在 \(U\) 上为零。\textbf{支集} \(\operatorname{supp} T\) 定义为 \(\Omega\) 中使 \(T\) 在其上非零的最大开集的补集：

\[\operatorname{supp} T = \Omega \setminus \bigcup\{U \subseteq \Omega \text{ 开} : T|_U = 0\}\]

这推广了连续函数的支集概念。

\textbf{命题 136.2}（分布的局部化）：分布可限制到任意开子集 \(U \subseteq \Omega\)：对 \(\varphi \in C_c^\infty(U)\)，定义 \(\langle T|_U, \varphi \rangle = \langle T, \varphi \rangle\)（\(\varphi\) 零延拓到 \(\Omega \setminus U\)）。若 \(\{\Omega_i\}\) 为 \(\Omega\) 的开覆盖，则 \(T = 0\) 当且仅当 \(T|_{\Omega_i} = 0\) 对所有 \(i\)。局部化的存在使分布构成一个层（sheaf）。

\textbf{定理 136.3}（分布的局部结构定理）：每个分布 \(T \in \mathcal{D}'(\Omega)\) 在局部上是某连续函数的导数：对任意紧集 \(K \subset \Omega\)，存在多重指标 \(\alpha\) 和连续函数 \(g \in C(\Omega)\) 使得

\[T|_K = D^\alpha g\]

换言之，分布的''奇异性''本质上源于微分运算作用于连续函数

\textbf{证明}：由分布的局部有限阶性，对紧 \(K\)，存在 \(m\) 和 \(C\) 使 \(|\langle T,\varphi\rangle| \leq C\|\varphi\|_{C^m(K)}\)。\(T\) 可延拓到 \(C^m(K)\) 上。由 Riesz 表示，\(T\) 可表为 \(C^m(K)\) 上测度的导数。更精确地，反复分部积分可得 \(T|_K = D^\alpha g\)，其中 \(g\) 连续。这一结构定理说明分布的奇异性无非是经典函数的广义导数。\(\blacksquare\) 。

\textbf{定义 136.5}（紧支撑分布 \(\mathcal{E}'\)）：有紧支撑的分布全体记为 \(\mathcal{E}'(\mathbb{R}^n)\)，它是光滑函数空间 \(C^\infty(\mathbb{R}^n)\) 上的连续线性泛函。等价于存在紧集 \(K\) 和 \(m \in \mathbb{N}\) 使得 \(|\langle T, \varphi \rangle| \leq C \|\varphi\|_{C^m(K)}\)。

\textbf{定理 136.4}（Paley-Wiener-Schwartz 定理）：\(\mathcal{E}'(\mathbb{R}^n)\) 中分布的 Fourier 变换是 \(\mathbb{R}^n\) 上的实解析函数，且是多项式增长的

\textbf{证明}：\(T\in\mathcal{E}'(\mathbb{R}^n)\) 有紧支撑 \(K\)。Fourier 变换 \(\widehat{T}(\xi) = \langle T_x, e^{-ix\cdot\xi}\rangle\) 是 \(\xi\) 的整函数（因 \(e^{-ix\cdot\xi}\) 对 \(\xi\) 解析）。由 \(|\langle T,\varphi\rangle| \leq C\|\varphi\|_{C^m(K)}\)，得 \(|\widehat{T}(\xi)| \leq C(1+|\xi|)^m e^{R|\Im\xi|}\)（\(R = \sup_{x\in K}|x|\)）。反之，满足此类增长条件的整函数是某紧支撑分布的 Fourier 变换。\(\blacksquare\) 。

\subsection{分布的基本运算}\label{ux5206ux5e03ux7684ux57faux672cux8fd0ux7b97}

\textbf{定义 136.6}（分布乘光滑函数）：设 \(T \in \mathcal{D}'(\Omega)\)，\(\psi \in C^\infty(\Omega)\)。乘积 \(\psi T\) 定义为

\[\langle \psi T, \varphi \rangle = \langle T, \psi \varphi \rangle, \quad \forall \varphi \in C_c^\infty(\Omega)\]

验证 \(\psi T \in \mathcal{D}'(\Omega)\)：连续性是 \(\varphi \mapsto \psi \varphi\) 在 \(\mathcal{D}(\Omega)\) 中连续的直接推论。乘积满足 Leibniz 法则：\(D^\alpha(\psi T) = \sum_{\beta \leq \alpha} \binom{\alpha}{\beta} (D^{\alpha-\beta} \psi)(D^\beta T)\)。

\textbf{定义 136.7}（分布的导数）：分布理论的核心优势在于任意分布可任意阶求导。分布的导数总是存在（不同于经典函数！）。定义

\[\langle D^\alpha T, \varphi \rangle = (-1)^{|\alpha|} \langle T, D^\alpha \varphi \rangle, \quad \forall \varphi \in C_c^\infty(\Omega)\]

\(D^\alpha : \mathcal{D}'(\Omega) \to \mathcal{D}'(\Omega)\) 是连续线性算子。这是分部积分公式的自然推广。

\textbf{定义 136.8}（分布的张量积）：设 \(S \in \mathcal{D}'(\Omega_1)\)，\(T \in \mathcal{D}'(\Omega_2)\)。其\textbf{张量积} \(S \otimes T \in \mathcal{D}'(\Omega_1 \times \Omega_2)\) 定义为

\[\langle S \otimes T, \varphi \rangle = \langle S_x, \langle T_y, \varphi(x, y) \rangle \rangle = \langle T_y, \langle S_x, \varphi(x, y) \rangle \rangle\]

其中 \(\varphi \in C_c^\infty(\Omega_1 \times \Omega_2)\)。张量积的支集满足 \(\operatorname{supp}(S \otimes T) = \operatorname{supp} S \times \operatorname{supp} T\)。

\textbf{定义 136.9}（分布的卷积）：若 \(T \in \mathcal{D}'\) 和 \(\varphi \in \mathcal{D}\)，定义 \((T * \varphi)(x) = \langle T, \varphi(x - \cdot) \rangle\)，它是 \(C^\infty\) 函数。对两个分布 \(S, T\)，当至少一个有紧支撑时可定义 \(S * T\)：经张量积和变量替换，

\[\langle S * T, \varphi \rangle = \langle S_x \otimes T_y, \varphi(x + y) \rangle\]

卷积是结合的、交换的，且 \(D^\alpha(S * T) = (D^\alpha S) * T = S * (D^\alpha T)\)。\(\delta\) 是卷积单位元：\(T * \delta = T\)。

\subsection{基本解与 Malgrange-Ehrenpreis 定理}\label{ux57faux672cux89e3ux4e0e-malgrange-ehrenpreis-ux5b9aux7406}

\textbf{定义 136.10}（基本解 / fundamental solution）：常系数线性偏微分算子 \(P(D) = \sum_{|\alpha| \leq m} a_\alpha D^\alpha\) 的\textbf{基本解}是分布 \(E \in \mathcal{D}'\) 满足

\[P(D) E = \delta\]

基本解给出非齐次方程 \(P(D)u = f\) 的一个特解 \(u = E * f\)（当 \(f\) 有紧支撑时）。

\textbf{定理 136.5}（Malgrange-Ehrenpreis 定理，1954-1955）：任意非零常系数线性偏微分算子都有基本解

\textbf{证明}：对 \(P(D) = \sum a_\alpha D^\alpha\)，其符号 \(P(\xi) = \sum a_\alpha (i\xi)^\alpha\) 是多项式。Hörmander 梯状构造：在复空间 \(\mathbb{C}^n\) 中找到路径使 \(P(\zeta)\) 快速增长，沿此路径定义分布 \(E = (2\pi)^{-n}\int \frac{e^{ix\cdot\zeta}}{P(\zeta)}d\zeta\)。严格的证明使用 \(P(D)\) 的 \(L^2\) 估计：\(\sum_\alpha \|D^\alpha u\| \leq C(\|P(D)u\| + \|u\|)\)，对偶得满射性。\(\blacksquare\) 。

\textbf{例}：\(\mathbb{R}^n\)（\(n \geq 3\)）中 Laplace 算子的基本解 \(E(x) = c_n |x|^{2-n}\) 满足 \(\Delta E = \delta\)；\(\mathbb{R}^2\) 中 \(E(x) = \frac{1}{2\pi} \log |x|\)。热方程的基本解（热核）：\(E(x, t) = \frac{\mathbf{1}_{[0,\infty)}(t)}{(4\pi t)^{n/2}} e^{-|x|^2/4t}\)。

\subsection{\texorpdfstring{Fourier 变换在 \(\mathcal{S}'(\mathbb{R}^n)\) 上的定义}{Fourier 变换在 \textbackslash mathcal\{S\}\textquotesingle(\textbackslash mathbb\{R\}\^{}n) 上的定义}}\label{fourier-ux53d8ux6362ux5728-mathcalsmathbbrn-ux4e0aux7684ux5b9aux4e49}

\textbf{定义 136.11}（\(\mathcal{S}'\) 上的 Fourier 变换）：\(\mathcal{S}(\mathbb{R}^n)\) 上的 Fourier 变换

\[\widehat{\varphi}(\xi) = \mathcal{F}\varphi(\xi) = \int_{\mathbb{R}^n} e^{-i x \cdot \xi} \varphi(x)\,dx\]

是 \(\mathcal{S}\) 到自身的同构。其逆变换为 \(\mathcal{F}^{-1}\psi(x) = (2\pi)^{-n} \int e^{i x \cdot \xi} \psi(\xi)\,d\xi\)。\(\mathcal{F}\) 满足 \(\mathcal{F}(D^\alpha \varphi)(\xi) = \xi^\alpha \widehat{\varphi}(\xi)\) 和 \(\mathcal{F}(x^\alpha \varphi)(\xi) = i^{|\alpha|} D^\alpha \widehat{\varphi}(\xi)\)。\(\mathcal{S}\) 上的 Fourier 变换满足 Parseval 恒等式：\(\|\widehat{\varphi}\|_{L^2} = (2\pi)^{n/2} \|\varphi\|_{L^2}\)。

\textbf{\(\mathcal{S}'\) 上的延拓}：由对偶性，定义 \(\mathcal{F} : \mathcal{S}' \to \mathcal{S}'\) 为

\[\langle \mathcal{F}T, \varphi \rangle = \langle T, \mathcal{F}\varphi \rangle, \quad \forall \varphi \in \mathcal{S}(\mathbb{R}^n)\]

\(\mathcal{F}\) 在 \(\mathcal{S}'\) 上仍为同构。关键例子：\(\mathcal{F}\delta = 1\)（常函数），\(\mathcal{F}1 = (2\pi)^n \delta\)。多项式增长的函数的 Fourier 变换是 \(\delta\) 及其导数的线性组合。

\subsection{Sobolev 空间的分布定义}\label{sobolev-ux7a7aux95f4ux7684ux5206ux5e03ux5b9aux4e49}

\textbf{定义 136.12}（Sobolev 空间 \(H^s(\mathbb{R}^n)\)）：使用 Fourier 变换和分布定义分数阶 Sobolev 空间：

\[H^s(\mathbb{R}^n) = \{u \in \mathcal{S}'(\mathbb{R}^n) : (1 + |\xi|^2)^{s/2} \widehat{u}(\xi) \in L^2(\mathbb{R}^n)\}\]

范数为 \(\|u\|_{H^s} = \| (1 + |\xi|^2)^{s/2} \widehat{u} \|_{L^2}\)。当 \(s = k \in \mathbb{Z}_{\geq 0}\) 时，这就是经典的''直到 \(k\) 阶导数在 \(L^2\) 中''的 Sobolev 空间。\(H^{-s}\) 包含 Dirac \(\delta\) 等（\(n/2 < s\)）奇异对象。

\textbf{定理 136.6}（Sobolev 嵌入定理）：若 \(s > n/2\)，则 \(H^s(\mathbb{R}^n) \subset C^0(\mathbb{R}^n)\)（嵌入连续函数空间）。更一般地，若 \(s > k + n/2\)，则 \(H^s \subset C^k\)。这种''正则性提升''是 PDE 研究的基本工具

\textbf{证明}：由 Fourier 变换，\(u(x) = (2\pi)^{-n}\int \widehat{u}(\xi)e^{ix\xi}d\xi\)。\(|u(x)| \leq (2\pi)^{-n}\int |\widehat{u}(\xi)|d\xi = (2\pi)^{-n}\int (1+|\xi|^2)^{-s/2}\cdot (1+|\xi|^2)^{s/2}|\widehat{u}(\xi)|d\xi\)。由 Cauchy-Schwarz，若 \(s > n/2\)，则 \(\int (1+|\xi|^2)^{-s}d\xi < \infty\)，故 \(|u(x)| \leq C\|u\|_{H^s}\)，即 \(H^s\subset L^\infty\)。连续性由稠密性 \(\mathcal{S}\subset H^s\) 和一致估计得出。高阶情形类似，\(D^\alpha u\) 用 \((i\xi)^\alpha\widehat{u}\) 估计。\(\blacksquare\) 。

\hrulefill

\emph{本卷建立了从 TVS、Fréchet/LF 空间到 Schwartz 分布理论的完整分析框架。分布理论统一了 PDE 的弱解理论、调和分析和量子物理的数学基础。}

\newpage

\chapter{06-分析/05-调和分析/02-Littlewood-Paley.md}\label{ux5206ux679005-ux8c03ux548cux5206ux679002-littlewood-paley.md}

\section{第195章：Littlewood-Paley理论}\label{ux7b2c195ux7ae0littlewood-paleyux7406ux8bba}

Littlewood-Paley理论通过频率二进分解将函数拆分为不同频段的分量，以平方函数研究函数的\(L^p\)性质，是现代调和分析的核心工具。该理论起源于Littlewood和Paley在1930年代对Fourier级数二进块的研究，后经Stein、Fefferman等人的系统发展，成为处理奇异积分算子、函数空间刻画和PDE正则性估计的标准语言。其基本哲学是：与其直接研究函数\(f\)，不如研究其频率分量的平方和（即平方函数），后者在\(L^p\)理论中具有更好的行为。

\subsection{二进分解}\label{ux4e8cux8fdbux5206ux89e3}

\textbf{定义140.1}（Littlewood-Paley二进分解）：选取光滑函数\(\varphi \in C_c^\infty(\mathbb{R}^n)\)，支撑于环形区域\(\{\xi \in \mathbb{R}^n : 1/2 \leq |\xi| \leq 2\}\)，满足对所有\(\xi \neq 0\)有\(\sum_{j \in \mathbb{Z}} \varphi(2^{-j}\xi) = 1\)。定义\textbf{Littlewood-Paley算子} \[\Delta_j f = \mathcal{F}^{-1}\big(\varphi(2^{-j}\cdot) \hat{f}\big)\] \(\Delta_j f\)为\(f\)在频段\(|\xi| \sim 2^j\)上的分量，是光滑函数。记低频算子\(S_j f = \sum_{k \leq j-1} \Delta_k f\)。

\textbf{构造细节}：取径向函数\(\psi \in C_c^\infty(\mathbb{R}^n)\)满足\(0 \leq \psi(\xi) \leq 1\)，\(\psi(\xi) = 1\)当\(|\xi| \leq 1\)，\(\psi(\xi) = 0\)当\(|\xi| \geq 2\)。定义\(\varphi(\xi) = \psi(\xi/2) - \psi(\xi)\)，则\(\operatorname{supp} \varphi \subseteq \{1/2 \leq |\xi| \leq 2\}\)，且\(\sum_{j \in \mathbb{Z}} \varphi(2^{-j}\xi) = 1\)（对\(\xi \neq 0\)）。低频部分由\(\psi(\xi) = \sum_{j \leq 0} \varphi(2^{-j}\xi)\)（\(\xi \neq 0\)）或引入\(\Phi(\xi) = 1 - \sum_{j \geq 0} \varphi(2^{-j}\xi)\)来处理。这一构造的关键在于\(\varphi\)的紧支性和光滑性，保证了\(\Delta_j f\)的快速衰减。

\textbf{性质}：\(\Delta_j\)的卷积核\(\psi_j(x) = \mathcal{F}^{-1}[\varphi(2^{-j}\cdot)](x) = 2^{jn} \check{\varphi}(2^j x)\)，其中\(\check{\varphi} = \mathcal{F}^{-1}\varphi\)为Schwartz函数。由\(\varphi\)的消失矩（\(\int \varphi(\xi) d\xi = 0\)，因为\(\varphi\)支撑远离原点），\(\int \psi_j(x) dx = 0\)。此外，\(\|\psi_j\|_{L^1} = \|\check{\varphi}\|_{L^1}\)与\(j\)无关，这一致有界性是Bernstein不等式和各类乘子定理的基础。

\textbf{定理140.1}（二进分解的完备性）：对\(f \in \mathcal{S}'(\mathbb{R}^n)\)，有\(f = \sum_{j \in \mathbb{Z}} \Delta_j f\)（在\(\mathcal{S}'\)中收敛，模去多项式）。更精确地说，该级数在\(\mathcal{S}'/\mathcal{P}\)中收敛，其中\(\mathcal{P}\)为多项式空间。若\(f\)的Fourier变换在原点附近为零（即\(\hat{f}\)在\(0\)的某邻域内为零），则级数在\(\mathcal{S}'\)中无条件收敛。

\textbf{证明}：由乘子恒等式\(\sum_{j \in \mathbb{Z}} \varphi(2^{-j}\xi) = 1\)（\(\xi \neq 0\)），对\(\xi \neq 0\)有\(\hat{f}(\xi) = \sum_j \varphi(2^{-j}\xi) \hat{f}(\xi)\)，配合Fourier逆变换即得分解。\(\psi_j\)的\(L^1\)范数由伸缩不变性得出：\(\|\psi_j\|_{L^1} = \int 2^{jn}|\check{\varphi}(2^j x)| dx = \int |\check{\varphi}(y)| dy\)。对\(\mathcal{S}'\)收敛性，有界子集上的收敛等价于对每个测试函数\(\phi \in \mathcal{S}\)的逐点收敛，由\(\hat{f}\)的缓增性和\(\varphi\)的紧支性保证。\(\blacksquare\)

\textbf{定理140.2}（Bernstein不等式）：对\(0 < p \leq q \leq \infty\)，存在仅依赖于\(p,q,n,\varphi\)的常数\(C\)，使得对任意\(j \in \mathbb{Z}\)和任意满足\(\operatorname{supp} \hat{f} \subseteq \{|\xi| \leq C_0 2^j\}\)的函数\(f\)，有 \[\|\partial^\alpha f\|_{L^q} \leq C 2^{j|\alpha| + jn(1/p - 1/q)} \|f\|_{L^p}\] 特别地，对\(\Delta_j f\)（其Fourier变换支撑于\(|\xi| \sim 2^j\)），有 \[\|\Delta_j f\|_{L^q} \leq C 2^{jn(1/p - 1/q)} \|\Delta_j f\|_{L^p} \quad (q \geq p)\] 即频段定位在\(|\xi| \sim 2^j\)的函数，其\(L^q\)范数以\(2^{jn(1/p-1/q)}\)倍数受\(L^p\)范数控制。这是频率局部化带来正则性增益的定量表达。

\textbf{证明}：将\(\Delta_j f\)写为\(\psi_j * (\Delta_j f)\)，其中\(\psi_j(x) = 2^{jn}\check{\varphi}(2^j x)\)如前。由Young不等式\(\|\psi_j * g\|_{L^q} \leq \|\psi_j\|_{L^r} \|g\|_{L^p}\)，其中\(1 + 1/q = 1/r + 1/p\)，即\(1/r = 1 + 1/q - 1/p\)。计算\(\|\psi_j\|_{L^r} = 2^{jn}\left(\int |\check{\varphi}(2^j x)|^r dx\right)^{1/r} = 2^{jn(1 - 1/r)}\|\check{\varphi}\|_{L^r} \leq C 2^{jn(1/p - 1/q)}\)。对于导数估计，注意到\(\partial^\alpha\)在Fourier变换侧对应乘子\((2\pi i\xi)^\alpha\)，其\(|\xi| \sim 2^j\)上的模为\(C 2^{j|\alpha|}\)，结合Bernstein不等式即得。\(\blacksquare\)

\textbf{命题140.1}（Bernstein不等式的逆形式）：若\(\operatorname{supp} \hat{f} \subseteq \{|\xi| \geq C_0 2^j\}\)，则对\(1 \leq p \leq q \leq \infty\)有 \[\|f\|_{L^p} \leq C 2^{-jN} \sum_{|\alpha| = N} \|\partial^\alpha f\|_{L^p}\] 即高频函数由高阶导数控制。这是Bernstein不等式的对偶版本。

\textbf{证明}：选取\(\tilde{\varphi} \in C_c^\infty\)满足\(\tilde{\varphi}(\xi) = 1\)在\(f\)的支集上且\(\tilde{\varphi}(0) = 0\)。则\(f = \tilde{\Delta}_j f\)，其中\(\tilde{\Delta}_j\)的核为\(2^{jn}\check{\tilde{\varphi}}(2^j x)\)。由\(\tilde{\varphi}\)在原点处为零，\(\check{\tilde{\varphi}}\)的各阶矩为零，Taylor展开给出\(\tilde{\Delta}_j f = 2^{-jN} \sum_{|\alpha|=N} \tilde{\Delta}_j^{(\alpha)} \partial^\alpha f\)，其中\(\tilde{\Delta}_j^{(\alpha)}\)为有界于\(L^p\)的算子。\(\blacksquare\)

\subsection{平方函数}\label{ux5e73ux65b9ux51fdux6570}

\textbf{定义140.2}（Littlewood-Paley \(g\)-函数 --- Poisson半群版本）：设\(P_t = e^{-t\sqrt{-\Delta}}\)为Poisson半群（即\(P_t f = P(\cdot, t) * f\)，\(P(x,t) = c_n t / (|x|^2 + t^2)^{(n+1)/2}\)为Poisson核）。定义Littlewood-Paley \(g\)-函数 \[g(f)(x) = \left( \int_0^\infty \left| \frac{\partial}{\partial t} P_t * f(x) \right|^2 t \, dt \right)^{1/2}\]

等价地，以频域二进分解定义\textbf{离散\(g\)-函数}：\(g_d(f)(x) = \left(\sum_{j \in \mathbb{Z}} |\Delta_j f(x)|^2\right)^{1/2}\)。两者在\(L^p\)范数下等价。

\textbf{动机}：\(g\)-函数的基本思想是：在Fourier变换侧，\(t \frac{\partial}{\partial t} P_t\)的乘子为\(2\pi t|\xi| e^{-2\pi t|\xi|}\)，它将\(f\)分解为不同频率的分量，每个分量的特征尺度为\(t \sim |\xi|^{-1}\)。平方积分\(\int_0^\infty \cdots t dt\)恰好对应于\(dt/t\)测度下的平均，使得不同尺度的贡献在\(L^2\)意义下正交。

\textbf{定义140.3}（Lusin面积积分）：对\(\alpha > 0\)，定义\textbf{Lusin面积积分} \[S_\alpha(f)(x) = \left( \int_{\Gamma_\alpha(x)} |\nabla_{y,t} P_t * f(y)|^2 t^{1-n} \, dy \, dt \right)^{1/2}\] 其中\(\Gamma_\alpha(x) = \{(y,t) \in \mathbb{R}^n \times (0,\infty) : |x-y| < \alpha t\}\)为以\(x\)为顶点的锥。\(S_\alpha\)在锥而非仅铅垂线上平均，比\(g\)-函数提供更强的空间局部化信息。当\(\alpha = 1\)时简记为\(S(f)\)。

\textbf{定义140.4}（非切向极大函数）：定义Poisson积分的非切向极大函数 \[u^*(x) = \sup_{(y,t) \in \Gamma_1(x)} |P_t * f(y)|\] 它是\(f\)的Hardy-Littlewood极大函数的Poisson半群版本，在Hardy空间理论中起核心作用。

\textbf{定理140.3}（Littlewood-Paley定理）：对\(1 < p < \infty\)，存在常数\(a_p, A_p > 0\)使 \[a_p \|f\|_{L^p} \leq \|g_d(f)\|_{L^p} \leq A_p \|f\|_{L^p}, \quad a_p \|f\|_{L^p} \leq \|S(f)\|_{L^p} \leq A_p \|f\|_{L^p}\] 即\(L^p\)范数与平方函数的\(L^p\)范数等价。\(g_d(f)\)、\(g(f)\)与\(S(f)\)三者\(L^p\)范数相互等价。

\textbf{证明}：\(p = 2\)时由Plancherel定理直接得到：\(\|g_d(f)\|_2^2 = \sum_j \|\Delta_j f\|_2^2\)，而\(\sum_j \|\Delta_j f\|_2^2 = \int \sum_j |\varphi(2^{-j}\xi)|^2 |\hat{f}(\xi)|^2 d\xi\)。由\(\sum_j |\varphi(2^{-j}\xi)|^2 \asymp 1\)（\(\xi \neq 0\)），得\(\|g_d(f)\|_2 \sim \|f\|_2\)。

对一般\(p\)，使用Calderón-Zygmund向量值奇异积分理论。关键点：算子\(T: f \mapsto \{\Delta_j f\}_{j \in \mathbb{Z}}\)将\(L^p(\mathbb{R}^n)\)映射到\(L^p(\mathbb{R}^n; \ell^2(\mathbb{Z}))\)。在Fourier变换侧，\(T\)对应乘子\(\{\varphi(2^{-j}\xi)\}_{j \in \mathbb{Z}}\)，其取值于\(\ell^2\)。需要验证向量值Mikhlin乘子条件：对任意多重指标\(\alpha\)， \[\sup_{\xi \neq 0} \left(|\xi|^{|\alpha|} \left\|\partial_\xi^\alpha \{\varphi(2^{-j}\xi)\}_{j}\right\|_{\ell^2}\right) < \infty\] 由\(\varphi\)的紧支性，对每个\(\xi\)，仅有有限个\(j\)使\(\varphi(2^{-j}\xi) \neq 0\)，且这些\(j\)满足\(2^{j} \sim |\xi|\)。在此条件下，\(\partial_\xi^\alpha \varphi(2^{-j}\xi) = 2^{-j|\alpha|} (\partial^\alpha \varphi)(2^{-j}\xi)\)，其\(\ell^2\)范数以\(|\xi|^{-|\alpha|}\)为界。向量值Calderón-Zygmund理论{[}Rubio de Francia, Ruiz, Torrea 1986{]}给出\(T\)的\(L^p\)有界性。反向不等式由对偶（\(\Delta_j\)的自伴性）和\(\sum_j \Delta_j = I\)得出。\(\blacksquare\)

\textbf{定理140.4}（平方函数与Sobolev空间）：对\(1 < p < \infty\)，\(s \in \mathbb{R}\)，Sobolev范数等价于 \[\|f\|_{W^{s,p}} \sim \left\| \left( \sum_{j \in \mathbb{Z}} 2^{2js} |\Delta_j f|^2 \right)^{1/2} \right\|_{L^p}\] 这是Littlewood-Paley刻画Sobolev空间的经典结果。该等价性将微分算子\((-\Delta)^{s/2}\)转化为频率乘子\(2^{js}\)，大大简化了Sobolev空间中的估计。

\textbf{证明}：\(p = 2\)时由Plancherel定理：\(\|f\|_{H^s}^2 = \int (1+|\xi|^2)^s |\hat{f}(\xi)|^2 d\xi \sim \sum_j 2^{2js} \|\Delta_j f\|_2^2\)。对一般\(p\)，利用向量值乘子定理：算子\(f \mapsto \{2^{js} \Delta_j f\}_j\)的乘子为\(\{2^{js} \varphi(2^{-j}\xi)\}_j\)，验证Mikhlin条件即得\(L^p(\mathbb{R}^n) \to L^p(\mathbb{R}^n; \ell^2)\)的有界性。\(\blacksquare\)

\subsection{BMO空间与John-Nirenberg不等式}\label{bmoux7a7aux95f4ux4e0ejohn-nirenbergux4e0dux7b49ux5f0f}

\textbf{定义140.5}（BMO空间 --- 有界平均振动）：局部可积函数\(f\)属于\textbf{BMO}（Bounded Mean Oscillation），若 \[\|f\|_{\operatorname{BMO}} = \sup_{B} \frac{1}{|B|} \int_B |f(x) - f_B| \, dx < \infty\] 其中\(B\)遍历\(\mathbb{R}^n\)中所有球（或等价地，所有方体），\(f_B = \frac{1}{|B|} \int_B f(y) dy\)为\(f\)在\(B\)上的平均。模去常数后，\(\|\cdot\|_{\operatorname{BMO}}\)为范数。BMO包含所有有界函数以及\(\log|x|\)等，是\(L^\infty\)的自然推广。BMO空间由John和Nirenberg于1961年引入，旨在研究椭圆型偏微分方程的局部正则性。

\textbf{性质}：(i) \(\|f\|_{\operatorname{BMO}} = 0\)当且仅当\(f\)为常数（a.e.）；(ii) 若\(f \in L^\infty\)，则\(\|f\|_{\operatorname{BMO}} \leq 2\|f\|_{L^\infty}\)；(iii) \(\log|x| \in \operatorname{BMO}(\mathbb{R}^n)\)但\(\log|x| \notin L^\infty\)，说明BMO是\(L^\infty\)的真扩张。

\textbf{定理140.5}（John-Nirenberg不等式，1961）：存在仅依赖于维数\(n\)的常数\(c_1, c_2 > 0\)，使对任意球\(B\)和任意\(\lambda > 0\)， \[\frac{|\{x \in B : |f(x) - f_B| > \lambda\}|}{|B|} \leq c_1 \exp\left(-c_2 \frac{\lambda}{\|f\|_{\operatorname{BMO}}}\right)\] 即BMO函数的振动呈指数衰减分布------BMO是\(L^\infty\)在多尺度振荡意义下的温和偏离。该不等式表明BMO函数在任意球上的大偏差测度以指数速率衰减，这一性质在任何尺度的球上一致成立。

\textbf{证明}：核心是Calderón-Zygmund分解在BMO情形的逐层应用。标准化：不妨设\(\|f\|_{\operatorname{BMO}} = 1\)，\(B = Q_0\)为单位方体。对任意\(\lambda > 0\)，将\(Q_0\)二进剖分。对每个二进方体\(Q\)，若\(|f_Q| > \lambda\)，将\(Q\)再剖分。由停时论证，得到一族极大二进方体\(\{Q_j\}\)满足\(|f_{Q_j}| > \lambda\)且\(|f_{\tilde{Q}_j}| \leq \lambda\)（\(\tilde{Q}_j\)为\(Q_j\)的父方体）。关键估计：由BMO条件，\(\frac{1}{|Q_j|}\int_{Q_j}|f - f_{Q_j}| \leq 1\)，结合\(|f_{\tilde{Q}_j}| \leq \lambda\)导出\(|f_{Q_j}| \leq \lambda + C\)。对每个\(Q_j\)，在\(Q_j\)上继续应用此分解，经\(k\)层后振动幅度以指数衰减。精确地，得到\(\sum |Q_j| \leq e^{-c\lambda} |Q_0|\)。取\(\lambda\)的任意倍即得完整指数估计。\(\blacksquare\)

\textbf{推论140.1}（John-Nirenberg的\(L^p\)版本）：对任意\(1 \leq p < \infty\)，存在常数\(C_{p,n}\)使 \[\sup_B \left(\frac{1}{|B|}\int_B |f - f_B|^p dx\right)^{1/p} \leq C_{p,n} \|f\|_{\operatorname{BMO}}\] 即BMO范数可用任意\(L^p\)振动刻画。这由John-Nirenberg不等式的指数衰减直接积分得到。

\subsection{\texorpdfstring{Hardy空间\(H^p\)的原子分解}{Hardy空间H\^{}p的原子分解}}\label{hardyux7a7aux95f4hpux7684ux539fux5b50ux5206ux89e3}

\textbf{定义140.6}（\(H^p\)的极大函数定义）：对\(0 < p < \infty\)，设\(\Phi \in \mathcal{S}(\mathbb{R}^n)\)满足\(\int \Phi = 1\)。定义径向极大函数 \[M_\Phi f(x) = \sup_{t > 0} |f * \Phi_t(x)|, \quad \Phi_t(x) = t^{-n}\Phi(x/t)\] \(f \in \mathcal{S}'(\mathbb{R}^n)\)属于\textbf{Hardy空间}\(H^p(\mathbb{R}^n)\)，若\(M_\Phi f \in L^p(\mathbb{R}^n)\)，且\(\|f\|_{H^p} = \|M_\Phi f\|_{L^p}\)。该定义与\(\Phi\)的选取无关（在等价范数意义下）。当\(1 < p < \infty\)时，\(H^p = L^p\)（等价范数）；当\(p = 1\)时，\(H^1 \subsetneq L^1\)；当\(0 < p < 1\)时，\(H^p\)为完备的拟赋范空间。

\textbf{定义140.7}（\(H^p\)-原子）：对\(0 < p \leq 1\)，函数\(a \in L^2(\mathbb{R}^n)\)称为\textbf{\(H^p\)-原子}若存在球\(B\)使得：(i) \(\operatorname{supp}(a) \subseteq B\)（支撑条件）；(ii) \(\|a\|_{L^\infty} \leq |B|^{-1/p}\)（大小条件）；(iii) \(\int_{\mathbb{R}^n} a(x) x^\alpha dx = 0\)对所有\(|\alpha| \leq \lfloor n(1/p - 1) \rfloor\)（消失矩条件）。特别地，\(H^1\)-原子仅需满足零阶矩消失：\(\int a = 0\)。

\textbf{定理140.6}（\(H^p\)的原子分解 --- Coifman-Weiss, 1977）：\(f \in H^p(\mathbb{R}^n)\)（\(0 < p \leq 1\)）当且仅当存在\(H^p\)-原子的可数族\(\{a_j\}\)与系数\(\{\lambda_j\} \subseteq \mathbb{C}\)满足\(\sum_j |\lambda_j|^p < \infty\)，使得\(f = \sum_j \lambda_j a_j\)（在分布意义下，且级数在\(H^p\)中收敛）。此时 \[\|f\|_{H^p} \sim \inf \left( \sum_j |\lambda_j|^p \right)^{1/p}\] 其中下确界取遍所有可能的原子分解。\(H^1\)的特殊性在于需要消失零阶矩即可：\(\int a = 0\)。

\textbf{证明}：充分性------将原子分解代入\(H^p\)的极大函数刻画验证范数。关键在于：对\(H^p\)-原子\(a\)（支撑于球\(B\)，\(\|a\|_\infty \leq |B|^{-1/p}\)），其极大函数\(M_\Phi a\)在\(2B\)外以\(|B|^{1/p}|x - x_B|^{-n-1}\)衰减，经积分得\(\|M_\Phi a\|_{L^p}^p \leq C\)（与\(B\)无关）。由\(p\)-三角不等式（拟范数），\(\|f\|_{H^p}^p \leq \sum_j |\lambda_j|^p \|a_j\|_{H^p}^p \leq C \sum_j |\lambda_j|^p\)。

必要性------对\(H^p\)函数的Calderón-Zygmund分解逐层拆分为原子。核心构造：对\(f \in H^p\)，考虑其极大函数\(M_\Phi f\)的水平集\(\Omega_k = \{M_\Phi f > 2^k\}\)。作\(\Omega_k\)的Whitney分解\(\{Q_{k,i}\}\)，定义相应的''好函数''\(g_k\)和''坏函数''\(b_k = f - g_k\)。\(b_k\)可分解为支撑于\(Q_{k,i}\)的片段的叠加，每片满足消失矩条件（由Taylor展开和\(f\)在\(H^p\)中的正则性保证）。经适当重整，每片为\(H^p\)-原子的常数倍，且\(\ell^p\)系数和受\(\|f\|_{H^p}^p\)控制。\(\blacksquare\)

\subsection{\texorpdfstring{\(H^1\)与BMO的对偶性}{H\^{}1与BMO的对偶性}}\label{h1ux4e0ebmoux7684ux5bf9ux5076ux6027}

\textbf{定理140.7}（Fefferman对偶定理，1971）：\(\operatorname{BMO}(\mathbb{R}^n)\)与\(H^1(\mathbb{R}^n)\)互为对偶空间： \[\operatorname{BMO} = (H^1)^*, \quad H^1 = (\operatorname{VMO})^*\] 其中\(\operatorname{VMO}\)（Vanishing Mean Oscillation）为\(\operatorname{BMO}\)的闭子空间，由满足\(\lim_{|B| \to 0} \frac{1}{|B|}\int_B |f - f_B| = 0\)且\(\lim_{|B| \to \infty} \frac{1}{|B|}\int_B |f - f_B| = 0\)的函数构成。具体而言，\(f \in \operatorname{BMO}\)定义了\(H^1\)上的连续线性泛函\(\varphi \mapsto \int f \varphi\)，且\(\|f\|_{\operatorname{BMO}} \sim \sup_{\|\varphi\|_{H^1} \leq 1} |\int f \varphi|\)。

\textbf{证明}：一侧（BMO作用在\(H^1\)上）：对\(H^1\)原子\(a\)（支撑于\(B\)，\(\int a = 0\)，\(\|a\|_\infty \leq 1/|B|\)）， \[\left| \int f a \right| = \left| \int_B (f - f_B) a \right| \leq \|a\|_{L^\infty} \int_B |f - f_B| \leq \frac{1}{|B|} \cdot |B| \|f\|_{\operatorname{BMO}} = \|f\|_{\operatorname{BMO}}\] 由原子分解的线性性，对任意\(\varphi \in H^1\)，写\(\varphi = \sum \lambda_j a_j\)，得 \[\left|\int f \varphi\right| \leq \sum |\lambda_j| \left|\int f a_j\right| \leq \|f\|_{\operatorname{BMO}} \sum |\lambda_j| \leq C \|f\|_{\operatorname{BMO}} \|\varphi\|_{H^1}\]

另一侧（深度方向）：利用Carleson测度理论和帐篷空间。对任意\(H^1\)上的连续线性泛函\(L\)，需表示为\(L(\varphi) = \int f \varphi\)对某一\(f \in \operatorname{BMO}\)。考虑Poisson积分\(u(x,t) = P_t * \varphi(x)\)。关键工具：\(H^1\)函数可等价地由帐篷空间\(T^1\)刻画（\(\varphi \in H^1\)当且仅当\(|\nabla u(x,t)| t\)在锥\(\Gamma(x)\)上的积分属于\(L^1\)）。由\(L\)的连续性，可将其表示为某个Carleson测度\(\mu\)下的积分。Carleson嵌入定理：测度\(\mu\)为Carleson测度当且仅当存在\(f \in \operatorname{BMO}\)使\(d\mu = |\nabla P_t * f|^2 t dx dt\)。由此构造出\(f\)并验证\(\int f \varphi = L(\varphi)\)。\(\blacksquare\)

\subsection{Carleson测度与帐篷空间}\label{carlesonux6d4bux5ea6ux4e0eux5e10ux7bf7ux7a7aux95f4}

\textbf{定义140.8}（Carleson测度）：\(\mathbb{R}^{n+1}_+ = \mathbb{R}^n \times (0,\infty)\)上的正Borel测度\(\mu\)称为\textbf{Carleson测度}，若存在常数\(C\)使得对任意球\(B \subset \mathbb{R}^n\)， \[\mu(T(B)) \leq C |B|\] 其中\(T(B) = B \times (0, r_B]\)为\(B\)上的帐篷（\(r_B\)为\(B\)的半径）。Carleson测度的范数\(\|\mu\|_{\mathcal{C}}\)为满足上述不等式的最小常数\(C\)。

\textbf{定理140.8}（Carleson嵌入定理）：对\(1 < p < \infty\)，\(\mu\)为Carleson测度当且仅当存在常数\(C_p\)使得 \[\int_{\mathbb{R}^{n+1}_+} |P_t * f(x)|^p d\mu(x,t) \leq C_p \|\mu\|_{\mathcal{C}} \|f\|_{L^p}^p\] 该定理将上半空间的加权估计与边值函数的\(L^p\)范数联系起来，是BMO-H\^{}1对偶理论的基础。

\textbf{证明}：充分性------取\(f = \chi_B\)即得Carleson条件。必要性------利用非切向极大函数的弱型估计和Marcinkiewicz插值定理，将\(p=1\)的弱型估计提升为\(1 < p < \infty\)的强型估计。\(p=1\)的弱型估计由Calderón-Zygmund分解和Carleson条件的几何意义保证。\(\blacksquare\)

\subsection{函数空间的Littlewood-Paley统一刻画}\label{ux51fdux6570ux7a7aux95f4ux7684littlewood-paleyux7edfux4e00ux523bux753b}

\textbf{定义140.9}（Besov空间与Triebel-Lizorkin空间）：利用Littlewood-Paley分解定义完整函数空间尺度：

\begin{itemize}
\item
  \textbf{Besov空间} \(B^s_{p,q}\)（\(s \in \mathbb{R}\)，\(0 < p, q \leq \infty\)）： \[\|f\|_{B^s_{p,q}} = \left\| \left( 2^{js} \|\Delta_j f\|_{L^p} \right)_{j \in \mathbb{Z}} \right\|_{\ell^q}\]
\item
  \textbf{Triebel-Lizorkin空间} \(F^s_{p,q}\)（\(s \in \mathbb{R}\)，\(0 < p < \infty\)，\(0 < q \leq \infty\)）： \[\|f\|_{F^s_{p,q}} = \left\| \left( \sum_{j \in \mathbb{Z}} 2^{jsq} |\Delta_j f|^q \right)^{1/q} \right\|_{L^p}\]
\end{itemize}

这两个尺度在调和分析中处于核心地位，其关键在于：Besov空间先取\(L^p\)范数再取\(\ell^q\)范数，而Triebel-Lizorkin空间先取\(\ell^q\)范数再取\(L^p\)范数。范数交换顺序的差异产生两类不同的函数空间，它们在正则性刻画中各有所长。

这些空间统一了所有经典函数空间： - \(F^s_{p,2} = W^{s,p}\)（Sobolev空间，\(1 < p < \infty\)） - \(F^0_{p,2} = H^p\)（Hardy空间，\(0 < p < \infty\)） - \(F^0_{\infty,2} = \operatorname{BMO}\)（BMO空间） - \(B^s_{\infty,\infty} = C^s\)（Hölder-Zygmund空间） - \(F^s_{2,2} = B^s_{2,2} = H^s\)（\(L^2\)-Sobolev空间） - \(B^s_{p,p} = W^{s,p}\)当\(s\)非整数且\(1 < p < \infty\)（Sobolev-Slobodeckij空间） - \(B^0_{1,\infty}\)为\(L^1\)的''替代''（在奇异积分理论中）

\textbf{定理140.9}（Sobolev嵌入的Littlewood-Paley证明）：对\(s_1 - n/p_1 \geq s_2 - n/p_2\)（\(s_1 \geq s_2\)），有连续嵌入\(B^{s_1}_{p_1,q} \hookrightarrow B^{s_2}_{p_2,q}\)和\(F^{s_1}_{p_1,q_1} \hookrightarrow F^{s_2}_{p_2,q_2}\)。当\(s_1 - n/p_1 = s_2 - n/p_2\)时对应临界Sobolev嵌入，此时嵌入对\(q_1 \leq q_2\)在Besov情形成立，但Triebel-Lizorkin情形需更精细的条件。

\textbf{证明}：利用Bernstein不等式：\(\|\Delta_j f\|_{L^{p_2}} \leq C 2^{jn(1/p_1 - 1/p_2)} \|\Delta_j f\|_{L^{p_1}}\)。代入Besov范数： \[\|f\|_{B^{s_2}_{p_2,q}} = \left\| 2^{js_2} \|\Delta_j f\|_{L^{p_2}} \right\|_{\ell^q} \leq C \left\| 2^{j(s_2 - n(1/p_1 - 1/p_2))} \cdot 2^{js_1} \|\Delta_j f\|_{L^{p_1}} \right\|_{\ell^q}\] 由于\(s_2 - n/p_2 \leq s_1 - n/p_1\)，即\(s_2 + n/p_2 - n/p_1 \leq s_1\)，故\(s_2 - n(1/p_1 - 1/p_2) \leq s_1\)，得\(\|f\|_{B^{s_2}_{p_2,q}} \leq C \|f\|_{B^{s_1}_{p_1,q}}\)。Triebel-Lizorkin情形需Minkowski不等式交换\(\ell^q\)和\(L^p\)范数，使用Fefferman-Stein向量值极大不等式。\(\blacksquare\)

\subsection{仿积与Bony分解}\label{ux4effux79efux4e0ebonyux5206ux89e3}

\textbf{定义140.10}（Bony仿积分解，1981）：对两个缓增分布\(f\)和\(g\)，在适当正则性条件下，乘积\(fg\)的Littlewood-Paley分解可写为 \[fg = T_f g + T_g f + R(f,g)\] 其中\(T_f g = \sum_{j} S_{j-3} f \cdot \Delta_j g\)为\textbf{仿积}（paraproduct，\(f\)的低频与\(g\)的高频的相互作用），\(R(f,g) = \sum_{|j-k| \leq 2} \Delta_j f \cdot \Delta_k g\)为\textbf{共振项}（同频段相互作用）。此分解是研究非线性PDE中乘积正则性的核心工具，由Bony在其开创新工作中引入，用于处理不可压缩Euler方程和Navier-Stokes方程。

\textbf{命题140.2}（仿积的连续性）：设\(s > 0\)，\(1 < p < \infty\)。则 (i) 若\(f \in L^\infty\)，\(g \in H^s_p\)，则\(T_f g \in H^s_p\)且\(\|T_f g\|_{H^s_p} \leq C \|f\|_{L^\infty} \|g\|_{H^s_p}\)； (ii) 若\(f \in H^s_p\)，\(g \in L^\infty\)，则\(R(f,g) \in H^{s}_{p}\)当\(s > 0\)； (iii) 对\(s_1 + s_2 > 0\)，\(f \in H^{s_1}_p\)，\(g \in H^{s_2}_p\)，则\(fg \in H^{\min(s_1, s_2, s_1+s_2-n/p)}_p\)（乘积法则）。

\textbf{证明}：(i) 由\(S_{j-3} f\)的\(L^\infty\)有界性和\(\Delta_j g\)的\(L^p\)衰减性，\(2^{js}\|S_{j-3}f \cdot \Delta_j g\|_{L^p} \leq \|f\|_{L^\infty} \cdot 2^{js}\|\Delta_j g\|_{L^p}\)，\(\ell^2\)求和即得。(ii) 共振项仅涉及\(|\xi| \sim 2^j\)的频段，\(2^{js}\|\sum_{|k-j|\leq 2} \Delta_j f \Delta_k g\|_{L^p} \leq C 2^{js}\|\Delta_j f\|_{L^p} \|g\|_{L^\infty}\)，由\(f \in H^s_p\)得\(\ell^2\)可和性。(iii) 综合仿积和共振项估计即得。\(\blacksquare\)

\subsection{\texorpdfstring{\(T(1)\)定理}{T(1)定理}}\label{t1ux5b9aux7406}

\textbf{定理140.10}（David-Journé \(T(1)\)定理，1984）：设\(T: \mathcal{S}(\mathbb{R}^n) \to \mathcal{S}'(\mathbb{R}^n)\)为连续线性算子，具有标准Calderón-Zygmund核\(K(x,y)\)（即\(K\)在\(x \neq y\)处光滑且满足\(|K(x,y)| \leq C|x-y|^{-n}\)和\(|\nabla_x K(x,y)| + |\nabla_y K(x,y)| \leq C|x-y|^{-n-1}\)）。则\(T\)可延拓为\(L^2(\mathbb{R}^n)\)上的有界算子当且仅当满足以下三个条件： (i) \(T(1) \in \operatorname{BMO}\)； (ii) \(T^*(1) \in \operatorname{BMO}\)； (iii) \(T\)满足弱有界性：\(|\langle T f, g \rangle| \leq C R^n \|f\|_\infty \|g\|_\infty\)对任意支撑于半径为\(R\)的球内的\(f, g\)成立。

\textbf{证明概要}：充分性------利用Littlewood-Paley分解和仿积技术。将\(T\)分解为\(T = T_0 + T_1 + T_2\)，其中\(T_0\)对应对角线附近（\(\Delta_j T \Delta_k\)，\(|j-k| \leq C\)），\(T_1\)对应\(j \gg k\)（仿积型），\(T_2\)对应\(j \ll k\)。\(T(1) \in \operatorname{BMO}\)控制\(T_1\)的有界性，\(T^*(1) \in \operatorname{BMO}\)控制\(T_2\)的有界性。必要性------由\(T\)在\(L^2\)上的有界性，\(T(1)\)作为\(T\)作用在常数函数\(1\)上的结果，必属于\(\operatorname{BMO}\)（因为\(T\)将\(L^\infty\)映射到\(\operatorname{BMO}\)，这是Calderón-Zygmund算子的基本性质）。\(\blacksquare\)

\(T(1)\)定理是调和分析中里程碑式的结果，它将Calderón-Zygmund算子的\(L^2\)有界性归结为对常数函数\(1\)的检验，极大地简化了奇异积分算子的有界性判别。这一思想在后续的\(T(b)\)定理（David-Journé-Semmes 1985）和更一般的非双倍测度空间上的Calderón-Zygmund理论中得到进一步推广。

\hrulefill

\emph{Littlewood-Paley理论以二进分解和平方函数为双翼：\(g\)-函数（Poisson半群版本与离散版本）和Lusin面积积分提供了频率分量的精确\(L^p\)控制。BMO空间经由John-Nirenberg不等式揭示了有界平均振动的指数衰减本质，Fefferman对偶定理\(\operatorname{BMO} = (H^1)^*\)则完成了Hardy空间理论的闭环。Carleson测度与帐篷空间为对偶性提供了几何-测度论基础。Besov与Triebel-Lizorkin空间以统一框架收纳了所有经典函数空间，是现代调和分析处理PDE正则性问题的基础语言。Bony仿积分解和\(T(1)\)定理则展示了Littlewood-Paley分解在非线性分析和奇异积分理论中的强大威力。}

\emph{卷二十四：调和分析。}

\hrulefill

\section{第196章：Fourier 变换}\label{ux7b2c196ux7ae0fourier-ux53d8ux6362}

Fourier 变换是数学物理中最伟大的发现之一，其思想根植于 Joseph Fourier 1807 年关于热传导方程的研究------他声称任何周期函数均可展为正弦波与余弦波的无穷级数。这一论断在当时引发激烈争议（Lagrange、Laplace 等人均表怀疑），却最终催生了整个调和分析学科。从现代观点看，Fourier 变换的意义在于它提供了函数分析的''对偶视角''：物理空间中的微分算子变为频率空间中的乘性算子，卷积化为逐点乘积，使偏微分方程的求解大为简化。在本书体系中，Fourier 变换是调和分析篇的总纲：奇异积分理论（第138章）是对 Fourier 乘子的有界性研究，Littlewood-Paley 理论依赖于频率的二进制分解，而拟微分算子和 Fourier 积分算子则将这一思想推广到变系数情形。与分布论（Schwartz 理论）的结合更使 Fourier 变换突破函数空间的限制，用于处理 Dirac 测度、主值分布等广义对象。本章从 \(L^1\) 和 \(L^2\) 理论出发，建立缓增分布上的 Fourier 变换，为后续各章提供分析基础。

\subsection{\texorpdfstring{123.1 \(L^1\) 理论}{123.1 L\^{}1 理论}}\label{l1-ux7406ux8bba}

\textbf{定义 123.1}（Fourier 变换）：\(f \in L^1(\mathbb{R}^n)\) 的 \textbf{Fourier 变换}为

\[\hat{f}(\xi) = \mathcal{F}f(\xi) = \int_{\mathbb{R}^n} f(x) e^{-2\pi i x \cdot \xi} dx\]

\textbf{Fourier 逆变换}（形式地）：\(f(x) = \int_{\mathbb{R}^n} \hat{f}(\xi) e^{2\pi i x \cdot \xi} d\xi\)（在 \(L^2\) 或适当条件下成立）。

\textbf{命题 123.1}（\(L^1\) 理论的基本性质）： 1. \(\|\hat{f}\|_\infty \leq \|f\|_1\)（\(\hat{f}\) 是有界连续函数） 2. \(\widehat{f * g} = \hat{f} \hat{g}\)（卷积的 Fourier 变换是乘积） 3. \(\widehat{f(\cdot - h)}(\xi) = e^{-2\pi i h \cdot \xi} \hat{f}(\xi)\)（平移） 4. \(\widehat{e^{2\pi i x \cdot \eta} f(x)}(\xi) = \hat{f}(\xi - \eta)\)（调制） 5. \(\widehat{f(\lambda x)}(\xi) = \lambda^{-n} \hat{f}(\xi/\lambda)\)（伸缩） 6. Riemann-Lebesgue 引理：\(\hat{f}(\xi) \to 0\)（\(|\xi| \to \infty\)）

\textbf{定理 123.1}（Fourier 反演公式，\(L^1\) 情形）：若 \(f, \hat{f} \in L^1(\mathbb{R}^n)\)，则

\[f(x) = \int_{\mathbb{R}^n} \hat{f}(\xi) e^{2\pi i x \cdot \xi} d\xi \quad \text{几乎处处}\]

\emph{证明}：使用 Gauss 核近似恒等元。\(f * \varphi_\varepsilon \to f\)（\(\varepsilon \to 0\)），其中 \(\varphi_\varepsilon(x) = \varepsilon^{-n} e^{-\pi |x|^2/\varepsilon^2}\)。\(\hat{\varphi}_\varepsilon(\xi) = e^{-\pi \varepsilon^2 |\xi|^2}\)。由 Fourier 反演对 Gauss 函数成立，取极限。∎

\subsection{\texorpdfstring{123.2 \(L^2\) 理论与 Plancherel 定理}{123.2 L\^{}2 理论与 Plancherel 定理}}\label{l2-ux7406ux8bbaux4e0e-plancherel-ux5b9aux7406}

\textbf{定理 123.2}（Plancherel 定理）：Fourier 变换可以唯一延拓为 \(L^2(\mathbb{R}^n)\) 上的酉算子（在归一化后）：\(\|\hat{f}\|_2 = \|f\|_2\)。即

\[\int_{\mathbb{R}^n} |f(x)|^2 dx = \int_{\mathbb{R}^n} |\hat{f}(\xi)|^2 d\xi\]

\textbf{证明}：对 \(f \in \mathcal{S}(\mathbb{R}^n)\)，由 Fubini 定理 \(\langle \hat{f}, \hat{g} \rangle = \langle f, \check{\hat{g}} \rangle = \langle f, g \rangle\)（因 Fourier 逆变换 \(\check{\cdot}\) 在 Schwartz 空间上是 \(\mathcal{F}\) 的逆）。取 \(g = f\) 得 \(\|\hat{f}\|_2 = \|f\|_2\)。由于 \(\mathcal{S}\) 在 \(L^2\) 中稠密，等距延拓到 \(L^2\) 上。直接定义：对 \(f \in L^2\)，取 \(\mathcal{S}\) 中序列 \(f_n \to f\)（\(L^2\)），则 \(\hat{f}_n\) 为 \(L^2\) 中 Cauchy 序列，其极限定义为 \(\hat{f}\)。此延拓是酉算子。\(\blacksquare\)

\textbf{推论 123.3}（Parseval 恒等式）：\(\langle f, g \rangle = \langle \hat{f}, \hat{g} \rangle\)，即 \(\int f \bar{g} = \int \hat{f} \bar{\hat{g}}\)。

\textbf{定理 123.4}（Hausdorff-Young 不等式）：对 \(1 \leq p \leq 2\)，\(\|\hat{f}\|_{p'} \leq \|f\|_p\)，其中 \(1/p + 1/p' = 1\)。即 Fourier 变换是 \(L^p \to L^{p'}\) 的有界算子（\(p \in [1, 2]\)）。

\textbf{证明}：对 \(p = 1\)，\(\|\hat{f}\|_\infty \leq \|f\|_1\)（直接估计）。对 \(p = 2\)，Plancherel 给出 \(\|\hat{f}\|_2 = \|f\|_2\)。对 \(1 < p < 2\)，令 \(\frac{1}{p} = \frac{1-\theta}{1} + \frac{\theta}{2}\)（\(0 < \theta < 1\)），则 \(\frac{1}{p'} = \frac{1-\theta}{\infty} + \frac{\theta}{2} = \frac{\theta}{2}\)。由 Riesz-Thorin 插值定理，\(\|\mathcal{F}\|_{L^p \to L^{p'}} \leq \|\mathcal{F}\|_{L^1 \to L^\infty}^{1-\theta} \cdot \|\mathcal{F}\|_{L^2 \to L^2}^\theta \leq 1^{1-\theta} \cdot 1^\theta = 1\)，得 \(\|\hat{f}\|_{p'} \leq \|f\|_p\)。\(\blacksquare\)

\subsection{123.3 Schwartz 空间与缓增分布}\label{schwartz-ux7a7aux95f4ux4e0eux7f13ux589eux5206ux5e03-1}

\textbf{定义 123.2}（Schwartz 空间）：\textbf{Schwartz 空间} \(\mathcal{S}(\mathbb{R}^n)\) 是所有光滑函数 \(\varphi\) 的集合，满足对所有多重指标 \(\alpha, \beta\)，

\[\sup_{x \in \mathbb{R}^n} |x^\alpha D^\beta \varphi(x)| < \infty\]

即 \(\varphi\) 及其所有导数比任意多项式下降更快。

\textbf{定理 123.5}：Fourier 变换是 \(\mathcal{S}(\mathbb{R}^n)\) 上的同构（双射连续线性映射）。\(\mathcal{S}\) 在 Fourier 变换下的自同构性是调和分析的核心。

\textbf{证明}：取 \(f \in \mathcal{S}\)。由 \(\widehat{D^\alpha f}(\xi) = (2\pi i \xi)^\alpha \hat{f}(\xi)\) 和 \(\widehat{x^\beta f}(\xi) = (i/2\pi)^{|\beta|} D^\beta \hat{f}(\xi)\)，\(\hat{f}\) 光滑且对任意 \(\alpha,\beta\)，\(\xi^\alpha D^\beta \hat{f}(\xi) = (2\pi i)^{-|\beta|} \mathcal{F}[D^\alpha(x^\beta f)](\xi)\) 有界（因 \(\mathcal{F}: L^1 \to L^\infty\) 有界，且 \(D^\alpha(x^\beta f) \in L^1\)）。故 \(\hat{f} \in \mathcal{S}\)，即 \(\mathcal{F}(\mathcal{S}) \subseteq \mathcal{S}\)。同理 \(\mathcal{F}^{-1}(\mathcal{S}) \subseteq \mathcal{S}\)。由 Fourier 反演公式，\(\mathcal{F}\) 是 \(\mathcal{S}\) 上的双射，且关于 Schwartz 拓扑连续。\(\blacksquare\)

\textbf{定义 123.3}（缓增分布）：\(\mathcal{S}(\mathbb{R}^n)\) 的对偶空间 \(\mathcal{S}'(\mathbb{R}^n)\) 称为\textbf{缓增分布}空间。缓增分布 \(u\) 的 \textbf{Fourier 变换} \(\hat{u}\) 由对偶定义：

\[\langle \hat{u}, \varphi \rangle = \langle u, \hat{\varphi} \rangle \quad (\varphi \in \mathcal{S})\]

\subsection{123.4 Poisson 求和公式}\label{poisson-ux6c42ux548cux516cux5f0f}

\textbf{定理 123.6}（Poisson 求和公式）：设 \(f\) 是 Schwartz 函数。则

\[\sum_{n \in \mathbb{Z}^n} f(n) = \sum_{m \in \mathbb{Z}^n} \hat{f}(m)\]

更一般地，\(\sum_{n \in \mathbb{Z}^n} f(x + n) = \sum_{m \in \mathbb{Z}^n} \hat{f}(m) e^{2\pi i m \cdot x}\)（分布意义下）。

\textbf{证明}：定义 \(F(x) = \sum_{n \in \mathbb{Z}^n} f(x + n)\)（\(\mathbb{Z}^n\)-周期函数）。\(F\) 的 Fourier 系数为 \(\hat{f}(m)\)，由 Fourier 级数理论即得。\(\blacksquare\)

\hrulefill

\hrulefill

\hrulefill

\hrulefill

\hrulefill

\hrulefill

\section{第197章：奇异积分}\label{ux7b2c197ux7ae0ux5947ux5f02ux79efux5206}

奇异积分理论是现代调和分析的基石，由 Alberto Calderón 和 Antoni Zygmund 在 20 世纪 50 年代系统建立。其起源可追溯至 Hilbert 在 20 世纪初对复分析边界行为的研究------Hilbert 变换（共轭函数算子）是第一个被深入研究的奇异积分算子，其 \(L^p\) 有界性由 Marcel Riesz 在 1928 年证明。Calderón 和 Zygmund 的革命性贡献在于将这一维结果推广到高维：他们发现，只要积分核满足适当的光滑性条件和消失矩条件，一类广泛的卷积型奇异积分算子在所有 \(L^p\)（\(1 < p < \infty\)）上都有界。其核心工具------Calderón-Zygmund 分解------将任意 \(L^1\) 函数拆分为''好部分''（有界）和''坏部分''（集中于立方体上且均值为零），这一分解技巧在此后数十年的分析学发展中无处不在。本章与 Fourier 变换（第137章）密切呼应：Fourier 乘子理论表明奇异积分算子是 Fourier 乘子光滑性的边界情形，而 Hardy-Littlewood 极大函数（第139章）则提供了控制奇异积分的基本点态估计。

\subsection{124.1 Hilbert 变换}\label{hilbert-ux53d8ux6362}

\textbf{定义 124.1}（Hilbert 变换）：\(\mathbb{R}\) 上的 \textbf{Hilbert 变换} \(H\) 定义为

\[Hf(x) = \frac{1}{\pi} \operatorname{p.v.} \int_{-\infty}^\infty \frac{f(y)}{x - y} dy = \frac{1}{\pi} \lim_{\varepsilon \to 0} \int_{|x-y| > \varepsilon} \frac{f(y)}{x - y} dy\]

\textbf{命题 124.1}（Hilbert 变换的性质）： 1. \(H\) 是 \(L^2(\mathbb{R})\) 上的酉算子：\(H^* = -H\)，\(H^2 = -I\) 2. Fourier 乘子：\(\widehat{Hf}(\xi) = -i \operatorname{sgn}(\xi) \hat{f}(\xi)\) 3. \(H\) 是 \(L^p(\mathbb{R})\) 上的有界算子（\(1 < p < \infty\)） 4. \(H\) 与伸缩、平移交换

\textbf{定理 124.1}（M. Riesz 定理）：Hilbert 变换在 \(L^p(\mathbb{R})\) 上有界（\(1 < p < \infty\)）。\(\|Hf\|_p \leq C_p \|f\|_p\)。

\textbf{证明}：\(p=2\) 时由 Fourier 乘子 \(\widehat{Hf}(\xi) = -i\operatorname{sgn}(\xi)\hat{f}(\xi)\) 及 Plancherel 得 \(\|Hf\|_2 = \|f\|_2\)。对 \(1 < p < 2\)，利用 Calderón-Zygmund 分解将 \(f\) 分解为''好''部分 \(g\) 与''坏''部分 \(b\)。对 \(g\) 用 \(L^2\) 估计，对 \(b\) 用核 \(1/x\) 的光滑性及消去条件 \(\int_{|x|>2|y|} (H b_k)(x) dx\) 控制，得弱 \((1,1)\) 估计 \(|\{|Hf|>\lambda\}| \leq C\|f\|_1/\lambda\)。由 Marcinkiewicz 插值定理，\(H\) 在 \(L^p\)（\(1 < p \leq 2\)）有界。\(p>2\) 由对偶性 \((H)^* = -H\) 推出。\(\blacksquare\)

\subsection{124.2 Calderón-Zygmund 分解}\label{calderuxf3n-zygmund-ux5206ux89e3}

\textbf{定理 124.2}（Calderón-Zygmund 分解）：设 \(f \in L^1(\mathbb{R}^n)\)，\(f \geq 0\)，\(\lambda > 0\)。则存在 \(\mathbb{R}^n\) 的分解：\(f = g + b\)，其中 - \(g\)（``好''部分）：\(|g(x)| \leq 2^n \lambda\) 几乎处处，\(\|g\|_1 \leq \|f\|_1\) - \(b = \sum_k b_k\)（``坏''部分）：\(b_k\) 支撑在不相交的二进方体 \(Q_k\) 上，\(\int_{Q_k} b_k = 0\)，\(\|b_k\|_1 \leq 2^{n+1} \lambda |Q_k|\)，且 \(\sum_k |Q_k| \leq \lambda^{-1} \|f\|_1\)

\textbf{证明}：将 \(\mathbb{R}^n\) 划分为二进方体网格 \(\mathcal{D} = \bigcup_{k \in \mathbb{Z}} \mathcal{D}_k\)，其中 \(\mathcal{D}_k\) 为边长 \(2^{-k}\) 的二进方体。对每个二进方体 \(Q\)，计算均值 \(\frac{1}{|Q|}\int_Q f\)。若 \(\frac{1}{|Q|}\int_Q f > \lambda\)，则选取 \(Q\) 为''坏''方体。将未被选中的方体递归二分，直至所有子方体的均值 \(\leq \lambda\)。最终选出的极大''坏''方体族 \(\{Q_k\}\) 满足：\(\lambda < \frac{1}{|Q_k|}\int_{Q_k} f \leq 2^n \lambda\)（由父方体的均值 \(\leq \lambda\) 保证）。定义 \(g(x) = f(x)\)（\(x \notin \bigcup Q_k\)）和 \(g(x) = \frac{1}{|Q_k|}\int_{Q_k} f\)（\(x \in Q_k\)），则 \(|g| \leq 2^n \lambda\) 且 \(\|g\|_1 \leq \|f\|_1\)。定义 \(b_k = (f - \frac{1}{|Q_k|}\int_{Q_k} f)\mathbf{1}_{Q_k}\)，则 \(b = \sum b_k\)，\(\int b_k = 0\)，\(\|b_k\|_1 \leq 2^{n+1}\lambda|Q_k|\)。由 \(\sum|Q_k| \leq \lambda^{-1}\sum \int_{Q_k} f \leq \lambda^{-1}\|f\|_1\) 得所需估计。\(\blacksquare\)

\subsection{124.3 奇异积分算子}\label{ux5947ux5f02ux79efux5206ux7b97ux5b50}

\textbf{定义 124.2}（Calderón-Zygmund 核）：函数 \(K: \mathbb{R}^n \setminus \{0\} \to \mathbb{C}\) 称为 \textbf{Calderón-Zygmund 核}（或标准核），如果满足： 1. \(|K(x)| \leq C |x|^{-n}\)（大小条件） 2. \(|\nabla K(x)| \leq C |x|^{-n-1}\)（光滑性条件） 3. \(\int_{r < |x| < R} K(x) dx = 0\)（对所有 \(0 < r < R\)，消去条件）

\textbf{定义 124.3}（奇异积分算子）：与核 \(K\) 关联的\textbf{奇异积分算子} \(T\) 定义为

\[Tf(x) = \operatorname{p.v.} \int_{\mathbb{R}^n} K(x - y) f(y) dy = \lim_{\varepsilon \to 0} \int_{|x-y| > \varepsilon} K(x - y) f(y) dy\]

\textbf{定理 124.3}（Calderón-Zygmund 定理）：若 \(T\) 是奇异积分算子且在 \(L^2\) 上有界，则 \(T\) 在 \(L^p\) 上有界（\(1 < p < \infty\)），且满足弱 \((1, 1)\) 型估计：

\[|\{x : |Tf(x)| > \lambda\}| \leq \frac{C}{\lambda} \|f\|_1\]

\textbf{证明}：先证 \(T\) 在 \(L^2\) 上有界。对任意 \(f \in L^1 \cap L^2\)，取 Calderón-Zygmund 分解 \(f = g + b\)（\(\lambda > 0\)），其中 \(\|g\|_\infty \leq C\lambda\)，\(b = \sum b_k\) 支撑于不相交二进方体 \(Q_k\)，\(\int b_k = 0\)，\(\sum |Q_k| \leq \|f\|_1/\lambda\)。由 \(L^2\) 有界性，\(|\{|Tg| > \lambda/2\}| \leq C\lambda^{-2}\|g\|_2^2 \leq C\lambda^{-2}\cdot \lambda\|g\|_1 \leq C\|f\|_1/\lambda\)。对 \(b\) 部分，将 \(Tb_k\) 限制在 \(2Q_k\)（二倍方体）之外，利用核的光滑性 \(|\nabla K(x)| \leq C|x|^{-n-1}\) 和 \(\int b_k = 0\) 得 \(|Tb_k(x)| \leq C\lambda |Q_k|^{(n+1)/n}/|x-c_k|^{n+1}\)（\(x \notin 2Q_k\)）。积分估计给出弱 \((1,1)\) 型，组合两部分即得。由 Marcinkiewicz 插值，\(T\) 在 \(L^p\)（\(1 < p \leq 2\)）有界；\(p > 2\) 由对偶性。\(\blacksquare\)

\textbf{定理 124.5}（Cotlar 引理，1955）：设 \(\{T_j\}_{j \in \mathbb{Z}}\) 是 Hilbert 空间 \(\mathcal{H}\) 上的一族有界算子。若存在 \(\gamma: \mathbb{Z} \to \mathbb{R}_{\geq 0}\) 满足 \(\sum_{k \in \mathbb{Z}} \sqrt{\gamma(k)} < \infty\)，使得 \(\|T_i^* T_j\| \leq \gamma(i-j)^2\) 和 \(\|T_i T_j^*\| \leq \gamma(i-j)^2\) 对所有 \(i, j\) 成立，则级数 \(\sum_j T_j\) 在强算子拓扑下收敛，且 \(\|\sum_j T_j\| \leq \sum_k \sqrt{\gamma(k)}\)。

\textbf{证明}：对有限和 \(S_N = \sum_{|j| \leq N} T_j\)，考虑 \((S_N^* S_N)^m = \sum_{j_1,\ldots,j_{2m}} T_{j_1}^* T_{j_2} T_{j_3}^* \cdots T_{j_{2m}}\)。利用几乎正交条件，对每项交替使用 \(\|T_i^* T_j\|\) 和 \(\|T_i T_j^*\|\) 的界，得 \(\|S_N\|^{2m} = \|(S_N^* S_N)^m\| \leq (\sum_k \sqrt{\gamma(k)})^{2m}\)（\(m \to \infty\) 取幂根）。令 \(m \to \infty\) 得 \(\|S_N\| \leq \sum_k \sqrt{\gamma(k)}\)，与 \(N\) 无关。由 Cauchy 准则，\(\sum T_j\) 强收敛。Cotlar 引理是证明奇异积分算子 \(L^2\) 有界性的核心工具：将核 \(K\) 按频率二进分解为 \(K = \sum_j K_j\)，对每个分片算子 \(T_j\) 验证几乎正交性条件 \(\|T_i^* T_j\| \leq C 2^{-|i-j|\varepsilon}\)，则级数在 \(L^2\) 上收敛且有界。其应用场景包括：Hilbert 变换的 \(L^2\) 有界性、Calderón 交换子、Fourier 乘子定理以及仿积算子的有界性。\(\blacksquare\)

\subsection{124.4 Marcinkiewicz 插值定理}\label{marcinkiewicz-ux63d2ux503cux5b9aux7406}

\textbf{定理 124.4}（Marcinkiewicz 插值定理）：设 \(T\) 是次线性算子，满足： - 弱 \((p_0, q_0)\) 型：\(|\{|Tf| > \lambda\}| \leq (A_0 \|f\|_{p_0} / \lambda)^{q_0}\) - 弱 \((p_1, q_1)\) 型：\(|\{|Tf| > \lambda\}| \leq (A_1 \|f\|_{p_1} / \lambda)^{q_1}\)

其中 \(1 \leq p_0 < p_1 \leq \infty\)，\(1 \leq q_0, q_1 \leq \infty\)，\(q_0 \neq q_1\)。则对 \(p_\theta\) 满足 \(\frac{1}{p_\theta} = \frac{1-\theta}{p_0} + \frac{\theta}{p_1}\)，\(T\) 在 \(L^{p_\theta}\) 上是强有界的（\(0 < \theta < 1\)）。

\textbf{证明}：对 \(f\) 作水平分解 \(f^\lambda = f \mathbf{1}_{\{|f| > a\lambda\}}\)，\(f_\lambda = f - f^\lambda\)（\(a>0\) 待定）。由次可加性，\(\{|Tf| > \lambda\} \subseteq \{|Tf^\lambda| > \lambda/2\} \cup \{|Tf_\lambda| > \lambda/2\}\)。对 \(f^\lambda\) 用弱 \((p_0,q_0)\) 界，对 \(f_\lambda\) 用弱 \((p_1,q_1)\) 界。将 \(\|Tf\|_{p_\theta}^{p_\theta}\) 写为层饼积分 \(p_\theta \int_0^\infty \lambda^{p_\theta-1} |\{|Tf| > \lambda\}| d\lambda\)，代入分布函数估计，分离 \(\lambda\) 积分区域。选取最优的 \(a\) 平衡两项，最终得 \(\|Tf\|_{p_\theta} \leq C_{p_0,p_1,q_0,q_1,\theta} A_0^{1-\theta} A_1^\theta \|f\|_{p_\theta}\)。\(\blacksquare\)

\subsection{Calderón-Zygmund理论的现代发展}\label{calderuxf3n-zygmundux7406ux8bbaux7684ux73b0ux4ee3ux53d1ux5c55}

Calderón-Zygmund 理论自 1950 年代创立以来，已发展成为调和分析中最核心的理论框架，其影响远超奇异积分本身。\textbf{\(T(1)\) 定理}（David-Journe 1984）是 Calderón-Zygmund 理论中最重要的后续发展：一个与 Calderón-Zygmund 核关联的算子 \(T\) 在 \(L^2\) 上有界当且仅当 \(T(1)\) 和 \(T^*(1)\) 属于 BMO（有界平均振荡函数空间）。该定理将 \(L^2\) 有界性的判定归结为对两个特定函数（常函数 1）的检验，极大简化了奇异积分算子有界性的证明，并推广了此前 Cotlar 引理和几乎正交分解方法。\textbf{\(T(b)\) 定理}（David-Journe-Semmes 1985）进一步将 \(T(1)\) 定理推广到非双倍测度情形，允许用''适应性''函数 \(b\)（满足某种增生条件）代替常函数 1，这为 Lipschitz 曲线上 Cauchy 积分的有界性证明提供了关键工具。\textbf{Calderón 交换子}（1965）是 Calderón-Zygmund 理论的经典应用：交换子 \([A, T]f = A(Tf) - T(Af)\) 在 \(L^p\) 上的有界性反映了算子 \(T\) 与乘子 \(A\) 的非交换性，构成了 Coifman-Rochberg-Weiss 交换子理论（1976）的基础。\textbf{Muckenhoupt 权理论}（1972）将 Calderón-Zygmund 算子的有界性推广到加权 \(L^p\) 空间：\(T\) 在 \(L^p(w)\) 上有界当且仅当 \(w \in A_p\)，这为加权不等式提供了一个统一的特征刻画。\textbf{向量值奇异积分理论}（Burkholder 1980s）将奇异积分算子作用于 Banach 空间值函数，\(T\) 在 \(L^p(\mathbb{R}^n; X)\) 上有界当且仅当 \(X\) 是 UMD 空间（无条件鞅差空间），这揭示了调和分析中奇异积分有界性与 Banach 空间几何之间的深刻联系。

\hrulefill

\hrulefill

\section{第198章：Hardy-Littlewood 极大函数}\label{ux7b2c198ux7ae0hardy-littlewood-ux6781ux5927ux51fdux6570}

Hardy-Littlewood 极大函数由 G. H. Hardy 和 J. E. Littlewood 于 1930 年在研究 Fourier 级数的 summability 问题时引入，是调和分析中最简单却最深刻的非线性算子。其核心思想极为朴素：对局部可积函数的每一点，取其所在任意球上的积分平均的上确界。这一运算将函数''放大''为其所有局部平均值的主导项，从而产生一个比原函数更好的''正则化版本''。极大函数的根本重要性在于它作为''控制算子''的普适性------几乎所有经典调和分析中的线性算子（包括 Hilbert 变换、Riesz 变换和各类奇异积分）都可以被 Hardy-Littlewood 极大函数逐点控制，这使得极大函数的 \(L^p\) 有界性成为证明其他算子 \(L^p\) 有界性的''第一道关卡''。Hardy 和 Littlewood 最初证明了一维情形下 \(p > 1\) 的强型不等式，而 Wiener 在 1939 年将其推广到高维。在现代分析中，极大函数理论已被推广到一般的度量测度空间（空间 \((X,d,\mu)\) 满足 doubling 条件），成为度量测度空间上调和分析的出发点。

\subsection{125.1 Hardy-Littlewood 极大函数}\label{hardy-littlewood-ux6781ux5927ux51fdux6570}

\textbf{定义 125.1}（Hardy-Littlewood 极大函数）：\(f \in L^1_{\text{loc}}(\mathbb{R}^n)\) 的 \textbf{Hardy-Littlewood 极大函数}为

\[Mf(x) = \sup_{r > 0} \frac{1}{|B(x, r)|} \int_{B(x, r)} |f(y)| dy\]

即 \(f\) 在包含 \(x\) 的球上的最大平均值。

\textbf{定理 125.1}（Hardy-Littlewood 极大定理）： 1. \(M\) 是弱 \((1, 1)\) 型：\(|\{x : Mf(x) > \lambda\}| \leq \frac{C_n}{\lambda} \|f\|_1\) 2. \(M\) 是强 \((p, p)\) 型：\(\|Mf\|_p \leq C_{n,p} \|f\|_p\)（\(1 < p \leq \infty\)）

\emph{证明}（弱 \((1, 1)\) 型）：使用 Vitali 覆盖引理。对 \(\{Mf > \lambda\}\) 中的每点，存在球 \(B\) 使 \(\int_B |f| > \lambda |B|\)。由 Vitali 覆盖引理，可选出不相交的球 \(B_i\) 覆盖 \(\{Mf > \lambda\}\) 的缩放版本。\(|\{Mf > \lambda\}| \leq C_n \sum |B_i| \leq \frac{C_n}{\lambda} \sum \int_{B_i} |f| \leq \frac{C_n}{\lambda} \|f\|_1\)。∎

\textbf{定理 125.2}（Lebesgue 微分定理）：对 \(f \in L^1_{\text{loc}}(\mathbb{R}^n)\)，

\[\lim_{r \to 0} \frac{1}{|B(x, r)|} \int_{B(x, r)} f(y) dy = f(x) \quad \text{几乎处处}\]

\emph{证明}：由极大函数的弱 \((1, 1)\) 有界性和连续函数在 \(L^1\) 中的稠密性证明。∎

\subsection{125.2 其他极大函数}\label{ux5176ux4ed6ux6781ux5927ux51fdux6570}

\textbf{定义 125.2}（非中心极大函数）：\(M'f(x) = \sup_{B \ni x} \frac{1}{|B|} \int_B |f(y)| dy\)（上确界取遍所有包含 \(x\) 的球，不要求 \(x\) 是球心）。\(M'\) 与 \(M\) 逐点等价：\(Mf \leq M'f \leq 2^n Mf\)。

\textbf{定义 125.3}（锐极大函数，Sharp Maximal Function）：

\[M^\sharp f(x) = \sup_{B \ni x} \frac{1}{|B|} \int_B |f(y) - f_B| dy\]

其中 \(f_B = \frac{1}{|B|} \int_B f\) 是 \(f\) 在 \(B\) 上的均值。\(M^\sharp f\) 衡量 \(f\) 的局部振荡。

\textbf{定理 125.3}（Fefferman-Stein 不等式）：对 \(1 < p < \infty\)，\(\|f\|_p \leq C_p \|M^\sharp f\|_p\)（当 \(f\) 满足适当消失条件时）。

\textbf{证明}：核心是建立''好 \(\lambda\) 不等式''：存在 \(\gamma > 0\) 使对任意 \(\lambda > 0\) 和充分小的 \(\varepsilon > 0\)，有 \(|\{Mf > 2\lambda, M^\sharp f \leq \gamma\lambda\}| \leq C \varepsilon |\{Mf > \lambda\}|\)。该估计通过 Calderón-Zygmund 分解证明：在集合 \(\{Mf > \lambda\}\) 上用 Whitney 型分解，将极大函数与锐极大函数的局部平均振动关联。然后乘以 \(\lambda^{p-1}\) 对 \(\lambda\) 积分，利用 \(\|Mf\|_p^p \approx \int \lambda^{p-1} |\{Mf > \lambda\}| d\lambda\)，将 \(Mf\) 的分布由 \(M^\sharp f\) 控制，取 \(\varepsilon\) 充分小使 \(C\varepsilon < 1/2\) 即得。\(\blacksquare\)

\subsection{125.3 加权不等式}\label{ux52a0ux6743ux4e0dux7b49ux5f0f}

\textbf{定义 125.4}（\(A_p\) 权）：非负局部可积函数 \(w\) 称为 \textbf{\(A_p\) 权}（\(1 < p < \infty\)），如果

\[\sup_B \left(\frac{1}{|B|} \int_B w(x) dx\right) \left(\frac{1}{|B|} \int_B w(x)^{-1/(p-1)} dx\right)^{p-1} < \infty\]

\(A_1\) 权：\(M w(x) \leq C w(x)\) 几乎处处。

\textbf{定理 125.4}（Muckenhoupt 定理）：Hardy-Littlewood 极大函数在加权 \(L^p(w)\) 空间上有界当且仅当 \(w \in A_p\)（\(1 < p < \infty\)）。这给出了极大函数 \(L^p\) 有界性权条件的精确刻画。

\textbf{证明}：（必要性）若 \(M\) 在 \(L^p(w)\) 上有界，对任意球 \(B\)，取 \(f = w^{-1/(p-1)}\mathbf{1}_B\)，应用 \(M\) 的有界性和 \(Mf(x) \geq \frac{1}{|B|}\int_B f\)（\(x \in B\)），得 \((\frac{1}{|B|}\int_B w)(\frac{1}{|B|}\int_B w^{-1/(p-1)})^{p-1} \leq C\)，故 \(w \in A_p\)。（充分性）用反向 Hölder 不等式：\(A_p\) 权满足对某 \(\delta > 0\)，\(w^{1+\delta} \in A_p\)。构造二进极大函数，逐层用 Calderón-Zygmund 分解获得加权弱 \((1,1)\) 估计，再由插值得加权强 \((p,p)\) 估计。\(\blacksquare\)

\subsection{极大函数与调和分析的统一方法}\label{ux6781ux5927ux51fdux6570ux4e0eux8c03ux548cux5206ux6790ux7684ux7edfux4e00ux65b9ux6cd5}

Hardy-Littlewood 极大函数是调和分析中控制论的核心工具，其思想可以追溯到 Hardy 和 Littlewood 对 Fourier 级数求和的研究（1930）。极大函数方法的基本范式是：将复杂算子 \(T\) 的逐点行为用极大函数 \(Mf\) 控制，然后利用 \(M\) 的 \(L^p\) 有界性推导 \(T\) 的 \(L^p\) 有界性。\textbf{Fefferman-Stein 锐极大函数}（1972）\(M^\sharp f\) 通过衡量 \(f\) 的局部平均振荡，揭示了 \(L^p\) 范数与 BMO 半范数之间的对偶关系：\(\|f\|_p \leq C_p \|M^\sharp f\|_p\) 对于具有消失平均的 \(f\) 成立。这一不等式是 Fefferman-Stein 证明 \(L^p\) 中 Calderón-Zygmund 算子有界性的核心工具，避免了逐点估计的困难。\textbf{好 \(\lambda\) 不等式}（good-\(\lambda\) inequality, Burkholder-Gundy 1970）是极大函数方法中最重要的技术工具：\(|\{Mf > 2\lambda, M^\sharp f \leq \gamma\lambda\}| \leq C\varepsilon|\{Mf > \lambda\}|\) 允许将 \(Mf\) 的分布函数由 \(M^\sharp f\) 控制，从而将奇异积分算子的有界性归结为极大函数的控制。\textbf{\(A_\infty\) 权}是所有 \(A_p\) 权的并集，由反向 Hölder 不等式刻画：\(w \in A_\infty\) 当且仅当存在 \(\delta > 0\) 和 \(C\) 使得 \(\frac{1}{|B|}\int_B w^{1+\delta} \leq C(\frac{1}{|B|}\int_B w)^{1+\delta}\)。\(A_\infty\) 权具有性质：若 \(w \in A_\infty\)，则 \(w\) 赋给球 \(B\) 的测度 \(w(B)\) 与 \(|B|\) 的关系满足 \(w(E)/w(B) \leq C(|E|/|B|)^\delta\)。\textbf{Rp 权}（reverse Hölder class）是 \(A_\infty\) 权理论的自然延伸，在非线性调和分析（如 \(p\)-Laplacian 方程的梯度估计）中扮演关键角色。极大函数理论还与\textbf{非线性位势理论}（Wolff 位势）、\textbf{面积积分}（Lusin 面积积分）和\textbf{非切向极大函数}（nontangential maximal function）等构成 Hardy 空间理论的完整框架，为调和分析、偏微分方程和复分析中的边界值问题提供了统一的分析工具。

\newpage

\chapter{06-分析/05-调和分析/03-抽象调和与积分变换.md}\label{ux5206ux679005-ux8c03ux548cux5206ux679003-ux62bdux8c61ux8c03ux548cux4e0eux79efux5206ux53d8ux6362.md}

\section{第199章：局部紧Abel群上的调和分析}\label{ux7b2c199ux7ae0ux5c40ux90e8ux7d27abelux7fa4ux4e0aux7684ux8c03ux548cux5206ux6790}

抽象调和分析将经典Fourier分析从\(\mathbb{R}^n\)推广到任意局部紧Abel群（LCA群），其核心是Pontryagin对偶理论。该理论由Pontryagin（1934）和van Kampen（1935）独立建立，将Fourier变换、Poisson求和公式、Plancherel定理等经典结果统一为LCA群上的普遍定理。LCA群的重要性在于：它涵盖了Fourier分析中几乎所有的基本群结构------\(\mathbb{R}^n\)、\(\mathbb{T}^n\)、\(\mathbb{Z}^n\)、有限Abel群、\(p\)-adic数域\(\mathbb{Q}_p\)以及它们的乘积和商群。

\subsection{Haar测度与卷积代数}\label{haarux6d4bux5ea6ux4e0eux5377ux79efux4ee3ux6570}

\textbf{定理141.1}（Haar测度的存在唯一性，A. Haar 1933 / J. von Neumann 1934）：设\(G\)为局部紧群，则存在\(G\)上非零左不变Radon测度\(\mu\)（即\(\mu(gE) = \mu(E)\)对所有\(g \in G\)和Borel集\(E\)成立），该测度在相差正常数倍的意义下唯一。对Abel群，左不变即右不变。称此测度为\textbf{Haar测度}。紧群时\(\mu(G) < \infty\)（通常归一化为概率测度）；离散群时\(\mu\)为计数测度（每点赋予质量\(1\)）。

\textbf{证明思路}：存在性证明的核心是构造不变正线性泛函。对紧支撑连续函数\(f\)，定义其在平移下的''平均''：对给定的\(f, \varphi \in C_c^+(G)\)（\(\varphi \neq 0\)），定义相对覆盖数\((f : \varphi) = \inf\{\sum c_i : \exists g_i \in G, f(x) \leq \sum c_i \varphi(g_i^{-1}x)\}\)。取\(\varphi\)以单位元为极限的逼近序列，利用Urysohn引理和Tychonoff定理（选择公理）抽取子列极限，得到左不变正线性泛函\(I(f) = \lim (f : \varphi_n)/(f_0 : \varphi_n)\)（\(f_0\)为固定参考函数）。由Riesz表示定理，\(I\)对应唯一的左不变Radon测度。唯一性：若\(\mu, \nu\)为两个左不变Haar测度，则\(\mu\)关于\(\nu\)的Radon-Nikodym导数为左不变的常数函数，因此\(\mu\)和\(\nu\)成比例。\(\blacksquare\)

\textbf{定义141.1}（对偶群与特征标）：设\(G\)为LCA群。\(G\)的\textbf{对偶群}\(\widehat{G}\)是所有连续群同态\(\chi : G \to \mathbb{T}\)（\(\mathbb{T} = \{z \in \mathbb{C} : |z| = 1\}\)）的集合，配以紧致开拓扑（即\(\chi_n \to \chi\)当且仅当\(\chi_n\)在\(G\)的每个紧子集上一致收敛于\(\chi\)）与逐点乘法构成LCA群。\(\chi \in \widehat{G}\)称为\(G\)的\textbf{特征标}。

基本例子： - \(\widehat{\mathbb{R}} \cong \mathbb{R}\)：由\(\chi_\xi(x) = e^{2\pi i x\xi}\)给出同构； - \(\widehat{\mathbb{T}} \cong \mathbb{Z}\)：由\(\chi_n(e^{i\theta}) = e^{in\theta}\)给出同构； - \(\widehat{\mathbb{Z}} \cong \mathbb{T}\)：由\(\chi_z(n) = z^n\)给出同构； - \(\widehat{\mathbb{Z}/n\mathbb{Z}} \cong \mathbb{Z}/n\mathbb{Z}\)； - \(\widehat{\mathbb{Q}_p} \cong \mathbb{Q}_p\)（\(p\)-adic数域自对偶）； - \(\widehat{G \times H} \cong \widehat{G} \times \widehat{H}\)（对偶保持乘积）。

\textbf{定理141.2}（LCA群的结构定理）：每个LCA群\(G\)同构于\(\mathbb{R}^n \times G_0\)，其中\(G_0\)包含一个紧开子群。特别地，紧生成LCA群同构于\(\mathbb{R}^n \times \mathbb{Z}^m \times K\)（\(K\)紧）。此结构定理将LCA群的研究归结为\(\mathbb{R}^n\)、\(\mathbb{Z}^m\)、紧群和离散群的基本情形。

\textbf{定义141.2}（LCA群上的Fourier变换）：对\(f \in L^1(G)\)（关于Haar测度\(dx\)），其\textbf{Fourier变换}\(\hat{f} : \widehat{G} \to \mathbb{C}\)定义为 \[\hat{f}(\chi) = \int_G f(x) \overline{\chi(x)} dx\] \(\hat{f}\)是\(\widehat{G}\)上的连续函数。由Riemann-Lebesgue引理（在LCA群上同样成立），\(\hat{f} \in C_0(\widehat{G})\)，且\(\|\hat{f}\|_{L^\infty} \leq \|f\|_{L^1}\)。Fourier变换将\(L^1(G)\)上的卷积映射为\(C_0(\widehat{G})\)上的逐点乘积：\(\widehat{f * g} = \hat{f} \cdot \hat{g}\)。

\textbf{Fourier变换的动机}：在LCA群上，特征标\(\chi\)是平移算子的特征向量：对\(T_g f(x) = f(g^{-1}x)\)，有\(\widehat{T_g f}(\chi) = \overline{\chi(g)} \hat{f}(\chi)\)。因此Fourier变换将平移算子对角化，这一性质在LCA群层面上得到了最一般的表达。

\textbf{定理141.3}（Pontryagin对偶定理）：存在拓扑群的自然同构\(G \cong \widehat{\widehat{G}}\)，由赋值映射\(\alpha: x \mapsto (\chi \mapsto \chi(x))\)给出。即\(\alpha\)是连续群同构且为同胚。

\textbf{证明}：紧群情形用Peter-Weyl理论与矩阵元的完备性：紧群\(G\)的不可约酉表示的矩阵元在\(L^2(G)\)中稠密，对Abel群，不可约表示为特征标，故\(\widehat{G}\)分离\(G\)的点。离散群情形由\(\widehat{\widehat{\mathbb{Z}}} \cong \widehat{\mathbb{T}} \cong \mathbb{Z}\)直接验证。一般LCA群通过结构定理（\(\mathbb{R}^n \times G_0\)）归结为上述两种情形。\(\alpha\)的连续性由紧致开拓扑的定义得出，单射性由特征标分离点得出，满射性需验证\(\alpha(G)\)在\(\widehat{\widehat{G}}\)中稠密且闭。\(\blacksquare\)

\textbf{定理141.4}（LCA群上的Plancherel定理）：设\(G\)为LCA群，Haar测度\(dx\)。存在\(\widehat{G}\)上的唯一Haar测度\(d\chi\)（称为\textbf{Plancherel测度}或\textbf{对偶Haar测度}），使得Fourier变换可延拓为\(L^2(G)\)到\(L^2(\widehat{G})\)的酉算子： \[\|f\|_{L^2(G)} = \|\hat{f}\|_{L^2(\widehat{G})}, \quad f \in L^1(G) \cap L^2(G)\] 逆变换由\(f(x) = \int_{\widehat{G}} \hat{f}(\chi) \chi(x) d\chi\)给出（\(f \in L^1(\widehat{G})\)时逐点成立，一般情形在\(L^2\)意义下成立）。

\textbf{证明}：利用幺模群上卷积代数\(L^1(G)\)的Gelfand变换理论。\(L^1(G)\)为交换Banach \(*\)-代数，其Gelfand谱空间同构于\(\widehat{G}\)。Gelfand变换\(\Gamma: L^1(G) \to C_0(\widehat{G})\)恰为Fourier变换。\(L^1(G) \cap L^2(G)\)在\(L^2(G)\)中稠密，通过Plancherel恒等式\(\|f\|_2^2 = \int |\hat{f}(\chi)|^2 d\chi\)延拓Fourier变换为等距算子，满射性由\(\widehat{G}\)上Fourier逆变换保证。\(\blacksquare\)

\textbf{注记}：Plancherel测度\(d\chi\)的选取与\(G\)上Haar测度\(dx\)的选取有关。若\(dx\)乘以常数\(c\)，则\(d\chi\)乘以\(c^{-1}\)以保持酉性。对\(G = \mathbb{R}\)，若取\(dx\)为Lebesgue测度，则\(d\xi\)为\(d\xi\)（无\(2\pi\)因子），此时Fourier变换定义为\(\hat{f}(\xi) = \int f(x) e^{-2\pi i x\xi}dx\)。若定义\(\hat{f}(\xi) = \int f(x) e^{-i x\xi}dx\)，则Plancherel测度含\(1/2\pi\)因子。

\textbf{定理141.5}（Bochner定理 --- LCA群版）：\(G\)上的连续函数\(\varphi\)为\textbf{正定型}（即对所有有限集\(\{x_i\} \subset G\)和\(\{c_i\} \subset \mathbb{C}\)，\(\sum_{i,j} c_i \overline{c_j} \varphi(x_i^{-1}x_j) \geq 0\)）当且仅当存在\(\widehat{G}\)上唯一正有限Borel测度\(\mu\)使 \[\varphi(x) = \int_{\widehat{G}} \chi(x) d\mu(\chi)\] \(\mu\)称为\(\varphi\)的\textbf{谱测度}。正定型函数是自协方差函数Wiener-Khinchin关系的抽象化，在概率论（平稳过程的特征函数）和表示论（酉表示的矩阵系数）中起核心作用。

\textbf{证明}：正定性推出\(\varphi(e) \geq 0\)且\(|\varphi(x)| \leq \varphi(e)\)。\(\varphi\)定义了\(L^1(G)\)卷积代数上的正线性泛函\(L(f) = \int_G f(x) \varphi(x) dx\)。由Gelfand-Naimark-Segal（GNS）构造，\(L\)诱导出\(L^1(G)\)在Hilbert空间上的\(*\)-表示，其与\(\widehat{G}\)上的谱测度相联系。由Riesz表示定理，该谱测度在\(\widehat{G}\)上为有限正Borel测度\(\mu\)，满足\[L(f) = \int_{\widehat{G}} \hat{f}(\chi) d\mu(\chi)\]代入\(f \to \delta_x\)得\(\varphi(x) = \int_{\widehat{G}} \chi(x) d\mu(\chi)\)。唯一性由Fourier-Stieltjes变换的单射性保证：若两测度给出相同\(\varphi\)，则它们在\(L^1(G)\)的Fourier变换上一致，由Stone-Weierstrass定理，它们在\(C_0(\widehat{G})\)上一致，故测度相等。\(\blacksquare\)

\textbf{定理141.6}（Poisson求和公式的抽象版本）：设\(G\)为LCA群，\(H \subset G\)为闭子群。记\(H^\perp = \{\chi \in \widehat{G} : \chi|_H \equiv 1\}\)为\(H\)的\textbf{零化子}（\(\widehat{G}\)的闭子群，同构于\(\widehat{G/H}\)）。Haar测度适当归一化后，对Schwartz-Bruhat函数\(f\)（即满足适当光滑和衰减条件的函数）有 \[\int_H f(h) dh = \int_{H^\perp} \hat{f}(\chi) d\chi\]

经典Poisson公式为特例：\(G = \mathbb{R}\)，\(H = \mathbb{Z}\)，\(H^\perp = \mathbb{Z}\)，即\(\sum_{n \in \mathbb{Z}} f(n) = \sum_{n \in \mathbb{Z}} \hat{f}(n)\)。另一个重要实例：\(G = \mathbb{R}^n\)，\(H = \mathbb{Z}^n\)（格点），\(H^\perp = \mathbb{Z}^n\)（倒格点），即\(\sum_{k \in \mathbb{Z}^n} f(k) = \sum_{k \in \mathbb{Z}^n} \hat{f}(k)\)，这在晶体衍射理论和数论中至关重要。

\textbf{证明}：定义周期化算子\(P f(x) = \int_H f(xh) dh\)，则\(P f \in L^1(G/H)\)。\(P f\)的Fourier变换（在商群\(G/H\)上）为\(\widehat{P f}(\tilde{\chi}) = \hat{f}(\chi)\)对\(\chi \in H^\perp\)（其中\(\tilde{\chi}\)为\(\chi\)在\(G/H\)上的诱导特征标）。在商群\(G/H\)上对\(P f\)应用Fourier逆变换公式（在单位元\(eH\)处取值）： \[P f(eH) = \int_{\widehat{G/H}} \widehat{P f}(\tilde{\chi}) d\tilde{\chi} = \int_{H^\perp} \hat{f}(\chi) d\chi\] 而\(P f(eH) = \int_H f(h) dh\)，即得求和公式。\(\blacksquare\)

\subsection{不确定性原理的抽象形式}\label{ux4e0dux786eux5b9aux6027ux539fux7406ux7684ux62bdux8c61ux5f62ux5f0f}

\textbf{定理141.7}（LCA群上的不确定性原理）：设\(G\)为LCA群，\(f \in L^2(G)\)为非零函数。若\(f\)的支撑集和\(\hat{f}\)的支撑集均为紧集，则\(f = 0\)几乎处处。更一般地，存在常数\(C > 0\)（仅依赖于\(G\)上Haar测度的选取）使得 \[\mu(\operatorname{supp} f) \cdot \mu_{\widehat{G}}(\operatorname{supp} \hat{f}) \geq C\] 其中\(\mu, \mu_{\widehat{G}}\)分别为\(G\)和\(\widehat{G}\)上的Haar测度。对\(G = \mathbb{R}\)，\(C = 1/2\)（经典Heisenberg不等式\(\Delta x \cdot \Delta \xi \geq 1/4\pi\)的变体）。

\textbf{证明}：若\(f\)和\(\hat{f}\)均有紧支集，则\(f\)可延拓为\(\mathbb{C}\)上的整函数（由Fourier变换的Paley-Wiener定理），且\(f\)在\(\mathbb{R}\)上紧支撑，故\(f \equiv 0\)。定量版本：利用算子\(A = x\)和\(B = -i d/dx\)的对易关系\([A, B] = i\)，由Cauchy-Schwarz不等式得\(\|Af\|_2 \|Bf\|_2 \geq \frac{1}{2}\|f\|_2^2\)，即\(\Delta x \cdot \Delta \xi \geq 1/4\pi\)。在LCA群上，利用幺正表示和谱测度可得类似估计。\(\blacksquare\)

\subsection{分数阶Fourier变换}\label{ux5206ux6570ux9636fourierux53d8ux6362}

\textbf{定义141.3}（分数阶Fourier变换 / FrFT）：对\(\alpha \in \mathbb{R}\)，\(\mathbb{R}\)上的分数阶Fourier变换\(\mathcal{F}_\alpha\)定义为 \[\mathcal{F}_\alpha f(u) = \sqrt{\frac{1 - i\cot\alpha}{2\pi}} \int_{-\infty}^\infty f(x) \exp\left(i\frac{x^2 + u^2}{2} \cot\alpha - iux\csc\alpha\right) dx\] 当\(\alpha = \pi/2\)时\(\mathcal{F}_{\pi/2}\)即为标准Fourier变换；\(\mathcal{F}_0 = I\)（恒等算子）；\(\mathcal{F}_\pi = \mathcal{P}\)（宇称算子\(f(x) \mapsto f(-x)\)）。分数阶Fourier变换满足群性质\(\mathcal{F}_{\alpha+\beta} = \mathcal{F}_\alpha \circ \mathcal{F}_\beta\)，构成\(L^2(\mathbb{R})\)上的单参数酉算子群。

\textbf{动机}：分数阶Fourier变换将信号在时频平面中旋转角度\(\alpha\)，对应于Wigner分布在时频平面中的刚性旋转。它是量子力学中相空间变换和光学中Fresnel衍射的数学基础。

\hrulefill

\section{第200章：积分变换深化}\label{ux7b2c200ux7ae0ux79efux5206ux53d8ux6362ux6df1ux5316}

积分变换以积分核为桥梁，在线性PDE求解、信号处理和成像科学中是不可或缺的工具。本章系统讨论Mellin变换、Laplace变换、Hankel变换、小波变换以及Gabor变换的数学理论。

\subsection{Mellin变换}\label{mellinux53d8ux6362}

\textbf{定义142.1}（Mellin变换）：对\(f : (0, \infty) \to \mathbb{C}\)， \[\mathcal{M}[f](s) = \int_0^\infty f(x) x^{s-1} dx\] 定义域为复平面上使积分收敛的\(s\)的带形区域\(\{\sigma_1 < \operatorname{Re}(s) < \sigma_2\}\)。逆变换为 \[f(x) = \frac{1}{2\pi i} \int_{c - i\infty}^{c + i\infty} \mathcal{M}[f](s) x^{-s} ds\] 其中\(c\)在收敛带形内。Mellin变换在解析数论（通过Dirichlet级数与Mellin变换的联系）、特殊函数论和微分方程理论中广泛应用。

\textbf{Mellin变换的动机}：Mellin变换对尺度变换\(f(x) \mapsto f(ax)\)具有对角化作用：\(\mathcal{M}[f(ax)](s) = a^{-s}\mathcal{M}[f](s)\)。因此，尺度不变的微分算子（如Euler型微分算子\(x\frac{d}{dx}\)）在Mellin变换下化为乘子。

\textbf{定理142.1}（Mellin变换与Fourier变换的关系）：通过指数换元\(x = e^t\)，Mellin变换等价于Fourier变换： \[\mathcal{M}[f](s) = \int_{-\infty}^\infty f(e^t) e^{-st} e^t dt = \mathcal{F}[g](\xi)\] 其中\(g(t) = f(e^t) e^t\)，\(\xi = -is\)（即\(s = c + i\xi\)时）。因此所有Mellin变换性质可直接由Fourier变换性质翻译得到。

\textbf{证明}：代换\(x = e^t\)，\(dx = e^t dt\)，\(x^{s-1} = e^{(s-1)t}\)，故 \[\mathcal{M}[f](s) = \int_{-\infty}^\infty f(e^t) e^{st} dt = \int_{-\infty}^\infty [f(e^t) e^t] e^{-i(-is)t} dt\] 在\(s = c + i\omega\)时，上式化为\(g(t) = f(e^t)e^t\)的Fourier变换（乘子\(e^{-i\omega t}\)）。\(\blacksquare\)

\textbf{定理142.2}（Mellin卷积与尺度不变性）： (i) \(\mathcal{M}[f(ax)](s) = a^{-s} \mathcal{M}[f](s)\)（尺度不变性）； (ii) Mellin卷积\((f \star g)(x) = \int_0^\infty f(x/y) g(y) \frac{dy}{y}\)满足\(\mathcal{M}[f \star g] = \mathcal{M}[f] \cdot \mathcal{M}[g]\)； (iii) Parseval恒等式：\(\int_0^\infty f(x)\overline{g(x)} \frac{dx}{x} = \frac{1}{2\pi i} \int_{c-i\infty}^{c+i\infty} \mathcal{M}[f](s)\overline{\mathcal{M}[g](\overline{s})} ds\)。

\textbf{证明}：(i)由变量代换\(x \mapsto x/a\)直接得出。(ii)将Mellin卷积写为Fourier卷积通过指数换元：\((f \star g)(e^t) = \int_{-\infty}^\infty f(e^{t-u})g(e^u)du\)，即\(f(e^t)\)与\(g(e^t)\)的普通卷积，由Fourier卷积定理即得。(iii)由(ii)和Fourier变换的Parseval恒等式通过指数换元翻译。\(\blacksquare\)

\textbf{定理142.3}（Mellin变换与渐近展开）：设\(f\)在\(x \to 0^+\)和\(x \to \infty\)时具有幂次渐近展开：\(f(x) \sim \sum_{k} a_k x^{\alpha_k}\)（\(x \to 0^+\)），\(f(x) \sim \sum_{k} b_k x^{-\beta_k}\)（\(x \to \infty\)）。则\(\mathcal{M}[f](s)\)的极点恰好位于\(s = -\alpha_k\)和\(s = \beta_k\)处，且留数分别为\(a_k\)和\(-b_k\)。这一性质是Mellin变换在渐近分析中强大威力的来源------通过复平面的围道变形和留数计算来获取函数的渐近展开。

\textbf{证明}：将Mellin积分在\(s = -\alpha_k\)处展开：\(\mathcal{M}[f](s) = \int_0^1 (\sum a_k x^{\alpha_k} + \text{余项}) x^{s-1}dx + \int_1^\infty (\sum b_k x^{-\beta_k} + \text{余项}) x^{s-1}dx\)。第一项中的\(a_k x^{\alpha_k + s - 1}\)在\(s = -\alpha_k\)处有极点，留数为\(a_k\)；第二项中的\(b_k x^{-\beta_k + s - 1}\)在\(s = \beta_k\)处有极点，留数为\(-b_k\)。\(\blacksquare\)

\subsection{Laplace变换与Bromwich反演}\label{laplaceux53d8ux6362ux4e0ebromwichux53cdux6f14}

\textbf{定义142.2}（Laplace变换）：对\(f : [0, \infty) \to \mathbb{C}\)， \[\mathcal{L}[f](s) = \int_0^\infty f(t) e^{-st} dt\] 定义域为\(\{s \in \mathbb{C} : \operatorname{Re}(s) > \sigma_c\}\)，其中\(\sigma_c = \inf\{\sigma : \int_0^\infty |f(t)| e^{-\sigma t} dt < \infty\}\)为\textbf{收敛横坐标}。若\(\sigma_c = -\infty\)称\(f\)为指数型函数全体。Laplace变换在\(s\)的右半平面\(\operatorname{Re}(s) > \sigma_c\)上解析，且其导数满足\(\mathcal{L}[(-t)^n f(t)](s) = \mathcal{L}[f]^{(n)}(s)\)。

\textbf{定义142.3}（Laplace变换的卷积定理）：定义\(f\)和\(g\)的Laplace卷积\((f * g)(t) = \int_0^t f(t-\tau)g(\tau)d\tau\)。则\(\mathcal{L}[f * g](s) = \mathcal{L}[f](s) \cdot \mathcal{L}[g](s)\)。这一卷积定理是常系数线性ODE求解中部分分式展开和传递函数方法的理论基础。

\textbf{定理142.4}（Laplace反演公式 --- Bromwich积分）：若\(f\)分段光滑且在\([0,\infty)\)上指数增长（即存在\(M, \alpha\)使\(|f(t)| \leq M e^{\alpha t}\)），则对\(c > \max(\sigma_c, \alpha)\)且\(c\)大于\(f\)的Laplace变换所有奇点的实部， \[f(t) = \frac{1}{2\pi i} \int_{c - i\infty}^{c + i\infty} \mathcal{L}[f](s) e^{st} ds\] 即由\(\mathcal{L}[f]\)在复平面右半平面的取值沿铅垂线\(\operatorname{Re}(s) = c\)的围道积分恢复\(f(t)\)。在\(f\)的间断点处，积分收敛于\((f(t^+) + f(t^-))/2\)。

\textbf{证明}：将Laplace变换写为Fourier变换：设\(g(t) = f(t) e^{-ct}\)，则\(\mathcal{L}[f](c + i\omega) = \int_0^\infty f(t) e^{-ct} e^{-i\omega t} dt = \mathcal{F}[g](\omega/2\pi)\)。由Fourier反演公式， \[g(t) = \frac{1}{2\pi} \int_{-\infty}^\infty \mathcal{F}[g](\omega) e^{i\omega t} d\omega = \frac{1}{2\pi i} \int_{c-i\infty}^{c+i\infty} \mathcal{L}[f](s) e^{(s-c)t} ds\] 两边同乘\(e^{ct}\)即得。\(\blacksquare\)

\textbf{定理142.5}（Laplace变换的初值定理与终值定理）：若极限存在，则 \[\lim_{t \to 0^+} f(t) = \lim_{s \to \infty} s \mathcal{L}[f](s), \quad \lim_{t \to \infty} f(t) = \lim_{s \to 0^+} s \mathcal{L}[f](s)\] 这两个定理在控制理论中用于直接从传递函数推断系统的稳态和瞬态行为，无需显式求逆变换。

\textbf{证明}：初值定理：\(\mathcal{L}[f'](s) = s\mathcal{L}[f](s) - f(0^+)\)。由Riemann-Lebesgue引理，\(\lim_{s \to \infty} \mathcal{L}[f'](s) = 0\)，故\(\lim_{s \to \infty} s\mathcal{L}[f](s) = f(0^+)\)。终值定理：\(\lim_{s \to 0^+} \mathcal{L}[f'](s) = \int_0^\infty f'(t) dt = f(\infty) - f(0^+)\)，结合\(s\mathcal{L}[f](s) = \mathcal{L}[f'](s) + f(0^+)\)即得。\(\blacksquare\)

\subsection{Hankel变换}\label{hankelux53d8ux6362}

\textbf{定义142.4}（Hankel变换）：对\(\mathbb{R}^n\)上径向函数\(f(x) = f_0(|x|)\)，其Fourier变换也是径向的。Hankel变换是径向函数的Fourier变换在极坐标下的归约。记\(\nu = n/2 - 1\)， \[\mathcal{H}_\nu [f_0](\rho) = (2\pi)^{n/2} \rho^{-\nu} \int_0^\infty f_0(r) J_\nu(2\pi \rho r) r^{n/2} dr\] 其中\(J_\nu\)为第一类\(\nu\)阶Bessel函数。对逆变换，\(\mathcal{H}_\nu^{-1} = \mathcal{H}_\nu\)（对适当函数类）。Hankel变换是径向函数Fourier分析的自然工具，在衍射理论、声学成像和电磁散射中广泛应用。

\textbf{定理142.6}（Hankel变换与Fourier变换的关系）：对\(f(x) = f_0(|x|) \in L^1(\mathbb{R}^n)\)， \[\hat{f}(\xi) = \mathcal{H}_\nu [f_0](|\xi|)\] 即径向函数的Fourier变换化为Hankel变换。对\(L^2(\mathbb{R}_+, r^{n-1}dr)\)上函数，\(\mathcal{H}_\nu\)为酉算子。

\textbf{证明}：将Fourier积分\(\hat{f}(\xi) = \int_{\mathbb{R}^n} f_0(|x|) e^{-2\pi i x \cdot \xi} dx\)转为极坐标。设\(|\xi| = \rho\)，取\(\xi\)方向为极轴，则\(x \cdot \xi = r\rho \cos\theta\)。角度部分积分为 \[\int_{\mathbb{S}^{n-1}} e^{-2\pi i r\rho \cos\theta} d\sigma(\theta) = (2\pi)^{n/2} (r\rho)^{-\nu} J_\nu(2\pi r\rho)\] 其中\(\nu = n/2 - 1\)。代入径向积分即得Hankel积分表示。\(\blacksquare\)

\textbf{定理142.7}（Hankel变换的Parseval恒等式）：对\(f_0, g_0 \in L^2(\mathbb{R}_+, r^{n-1}dr)\)， \[\int_0^\infty f_0(r) \overline{g_0(r)} r^{n-1} dr = \int_0^\infty \mathcal{H}_\nu[f_0](\rho) \overline{\mathcal{H}_\nu[g_0](\rho)} \rho^{n-1} d\rho\] 这由Fourier变换的Parseval恒等式和径向对称性直接得出。

\subsection{小波变换与多尺度分析}\label{ux5c0fux6ce2ux53d8ux6362ux4e0eux591aux5c3aux5ea6ux5206ux6790}

\textbf{定义142.5}（连续小波变换）：设小波母函数\(\psi \in L^2(\mathbb{R})\)满足容许条件\(C_\psi = \int_{-\infty}^\infty |\hat{\psi}(\xi)|^2 |\xi|^{-1} d\xi < \infty\)。对\(f \in L^2(\mathbb{R})\)，定义连续小波变换 \[\mathcal{W}_\psi f(a, b) = \frac{1}{\sqrt{|a|}} \int_{-\infty}^\infty f(t) \overline{\psi\left(\frac{t-b}{a}\right)} dt\] 其中\(a \in \mathbb{R} \setminus \{0\}\)为尺度参数，\(b \in \mathbb{R}\)为平移参数。反演公式为 \[f(t) = \frac{1}{C_\psi} \int_{-\infty}^\infty \int_{-\infty}^\infty \mathcal{W}_\psi f(a, b) \frac{1}{\sqrt{|a|}} \psi\left(\frac{t-b}{a}\right) \frac{da\,db}{a^2}\]

\textbf{容许条件的意义}：\(\int \hat{\psi}(\xi)|\xi|^{-1} d\xi < \infty\)要求\(\hat{\psi}(0) = 0\)，即\(\int \psi(t) dt = 0\)（小波必须有零均值）。这是小波振荡性的数学表达，也是反演公式中\(C_\psi\)有限性的保证。

\textbf{定理142.8}（连续小波变换的等距性）：\(\|\mathcal{W}_\psi f\|_{L^2(\mathbb{R}^* \times \mathbb{R}, a^{-2} da db)} = \sqrt{C_\psi} \|f\|_{L^2(\mathbb{R})}\)------即小波变换是\(L^2\)到加权\(L^2\)的等距嵌入（模常数因子）。这一性质说明小波变换忠实地保持了信号的能量，且小波系数在\(a^{-2}da db\)测度下构成\(L^2\)空间的冗余表示。

\textbf{证明}：对固定尺度\(a\)，Fourier变换将伸缩平移化为乘子：\(\mathcal{F}[\psi_{a,b}](\xi) = \sqrt{|a|} e^{-2\pi i b \xi} \hat{\psi}(a\xi)\)。于是小波系数可写作Fourier侧的内积： \[\mathcal{W}_\psi f(a,b) = \sqrt{|a|} \int \hat{f}(\xi) \overline{\hat{\psi}(a\xi)} e^{2\pi i b\xi} d\xi\] 对\(b\)积分（Fourier变换的Parseval）给出\(\int |\mathcal{W}_\psi f(a,b)|^2 db = |a| \int |\hat{f}(\xi)|^2 |\hat{\psi}(a\xi)|^2 d\xi\)。再对\(a\)积分得\(\int\int |\mathcal{W}_\psi f(a,b)|^2 a^{-2} da db = \int |\hat{f}(\xi)|^2 \left(\int |\hat{\psi}(a\xi)|^2 |a|^{-1} da\right) d\xi = C_\psi \|f\|_2^2\)。\(\blacksquare\)

\textbf{定义142.6}（多尺度分析与正交小波基）：\(L^2(\mathbb{R})\)的\textbf{多尺度分析}（MRA）是一列闭子空间\(\{V_j\}_{j \in \mathbb{Z}}\)满足： (i) \(\cdots \subset V_{-1} \subset V_0 \subset V_1 \subset \cdots\)（嵌套性）； (ii) \(\overline{\bigcup_{j} V_j} = L^2(\mathbb{R})\)（稠密性）且\(\bigcap_j V_j = \{0\}\)（分离性）； (iii) \(f \in V_j \iff f(2 \cdot) \in V_{j+1}\)（二进伸缩性）； (iv) 存在尺度函数\(\varphi \in V_0\)使\(\{\varphi(\cdot - k) : k \in \mathbb{Z}\}\)为\(V_0\)的规范正交基。

由MRA可构造小波函数\(\psi\)使\(\{\psi_{j,k}(x) = 2^{j/2} \psi(2^j x - k) : j, k \in \mathbb{Z}\}\)为\(L^2(\mathbb{R})\)的规范正交基------这就是\textbf{离散小波变换}的正交基表达。Daubechies（1988）证明了存在具紧支集的任意阶光滑小波，这是小波理论在实际应用（如JPEG2000图像压缩）中取得巨大成功的数学基础。

\textbf{定理142.9}（MRA的尺度方程与两尺度关系）：尺度函数\(\varphi\)满足\textbf{尺度方程}\(\varphi(x) = \sum_{k} h_k \varphi(2x - k)\)，其中\(\{h_k\}\)为低通滤波器系数。小波函数\(\psi(x) = \sum_k g_k \varphi(2x - k)\)由高通滤波器\(\{g_k\}\)生成，且\(g_k = (-1)^k \overline{h_{1-k}}\)（正交小波时）。滤波器系数满足正交条件：\(\sum_k h_k \overline{h_{k+2n}} = \delta_{0n}\)和\(\sum_k h_k = \sqrt{2}\)。

\textbf{证明}：由\(V_0 \subset V_1\)，\(\varphi \in V_0\)可展为\(\{\varphi(2 \cdot - k)\}\)在\(V_1\)的规范正交基下的线性组合即得尺度方程。由\(\psi \in W_0 = V_1 \ominus V_0\)（\(W_0\)为\(V_0\)在\(V_1\)中的正交补），可得高通滤波表示。正交性条件由\(\{\varphi(\cdot - k)\}\)为规范正交基的要求导出：\(\delta_{0n} = \langle \varphi, \varphi(\cdot - n) \rangle = \sum_{k,l} h_k \overline{h_l} \langle \varphi(2\cdot - k), \varphi(2\cdot - 2n - l) \rangle = \frac{1}{2}\sum_k h_k \overline{h_{k-2n}}\)。\(\blacksquare\)

\subsection{Gabor变换与短时Fourier变换}\label{gaborux53d8ux6362ux4e0eux77edux65f6fourierux53d8ux6362}

\textbf{定义142.7}（短时Fourier变换 / Gabor变换）：对窗函数\(g \in L^2(\mathbb{R})\)（通常取Gaussian窗\(g(t) = e^{-\pi t^2}\)），\(f\)的短时Fourier变换为 \[\mathcal{V}_g f(x, \xi) = \int_{-\infty}^\infty f(t) \overline{g(t-x)} e^{-2\pi i t \xi} dt\] STFT将信号\(f\)在时频平面\((x, \xi)\)上进行局部化分析：\(x\)为时间中心，\(\xi\)为频率中心。STFT的模平方\(|\mathcal{V}_g f(x, \xi)|^2\)称为\textbf{声谱图}（spectrogram），是信号时频分析的基本工具。

\textbf{定理142.10}（STFT的正交展开 / Gabor框架）：若窗函数\(g\)和格点参数\(\alpha, \beta > 0\)满足\(\alpha\beta < 1\)，则存在常数\(A, B > 0\)使得对任意\(f \in L^2(\mathbb{R})\)， \[A\|f\|_2^2 \leq \sum_{m,n \in \mathbb{Z}} |\langle f, g_{m,n} \rangle|^2 \leq B\|f\|_2^2\] 其中\(g_{m,n}(t) = e^{2\pi i m \beta t} g(t - n\alpha)\)为Gabor原子。此即Gabor框架不等式，表明时频平移族\(\{g_{m,n}\}\)构成了\(L^2(\mathbb{R})\)的框架（frame），从而允许稳定的信号重构。当\(\alpha\beta = 1\)时（临界采样），框架界相等且Gabor系构成规范正交基（如Balian-Low定理所述，此时\(g\)不能同时具有好的时频局部化性质）。

\hrulefill

\emph{抽象调和分析以Pontryagin对偶为核心，将Fourier变换提升至LCA群层面，Plancherel定理与Poisson求和公式在此框架下获得统一表达。不确定性原理和Bochner定理在LCA群上获得最一般的陈述。积分变换方面，Mellin变换经指数换元等同Fourier变换，其极点-留数对应为渐近分析提供了强大工具。Laplace变换以Bromwich积分反演，初值定理和终值定理连接了复分析与时域行为。Hankel变换处理径向函数的Fourier分析。小波变换以多尺度分析为理论骨架，连续小波变换的等距性和离散小波基的正交性将时频局部化推向极致。Gabor变换则提供了时频平面的均匀采样框架。这些积分变换的共性在于：它们都是某个交换群上Fourier变换的适当变形，其反演公式和等距性均可追溯至Pontryagin对偶理论。}

\emph{卷二十四：调和分析。}

\newpage

\chapter{06-分析/05-调和分析/04-小波与采样.md}\label{ux5206ux679005-ux8c03ux548cux5206ux679004-ux5c0fux6ce2ux4e0eux91c7ux6837.md}

\section{第201章：Fourier 分析与采样理论}\label{ux7b2c201ux7ae0fourier-ux5206ux6790ux4e0eux91c7ux6837ux7406ux8bba}

Fourier 分析与采样理论是调和分析中最具应用穿透力的分支，它回答了这样一个基本问题：一个连续信号在什么条件下可被其离散采样值完全重构？Shannon 在 1948 年发表的采样定理------带限信号可由其 Nyquist 频率间隔的均匀采样点唯一恢复------不仅奠定了数字信号处理的数学基础，也揭示了 Fourier 变换在信息论中的核心地位。从数学深层次看，采样定理本质上是 Hilbert 空间中 Riesz 基与框架理论的一个特例：采样函数 \(\operatorname{sinc}(t-nT)\) 恰好构成带限函数空间 \(PW_{[-B,B]}\) 的正交基。本章将 Fourier 分析的经典结果（连续与离散 Fourier 变换、Parseval 恒等式）与 Shannon 采样定理有机整合，通过 Paley-Wiener 定理刻画带限函数的解析延拓性质，并通过周期化（Poisson 求和公式）建立连续 Fourier 变换与离散 Fourier 级数之间的桥梁。这一理论是后续小波分析（第147章）和压缩感知（第149章）的必要前导------理解经典采样的''充分条件''才能理解压缩感知如何突破 Nyquist 极限。

\subsection{282.1 Fourier 变换的基本理论}\label{fourier-ux53d8ux6362ux7684ux57faux672cux7406ux8bba}

\textbf{定义 282.1}（连续 Fourier 变换）：设 \(f \in L^1(\mathbb{R}^d) \cap L^2(\mathbb{R}^d)\)。其 Fourier 变换定义为

\[\hat{f}(\omega) = \mathcal{F}[f](\omega) = \int_{\mathbb{R}^d} f(x) e^{-2\pi i \omega \cdot x} dx\]

Fourier 逆变换为 \(f(x) = \int_{\mathbb{R}^d} \hat{f}(\omega) e^{2\pi i \omega \cdot x} d\omega\)。

\textbf{定理 282.1}（Plancherel 定理）：Fourier 变换可扩展为 \(L^2(\mathbb{R}^d)\) 上的酉算子：\(\|\hat{f}\|_{L^2} = \|f\|_{L^2}\)。这是 \(L^2\) 空间上调和分析的基础

\textbf{证明}：对 \(f\in L^1\cap L^2\)，\(\|\hat{f}\|_2^2 = \int \hat{f}(\omega)\overline{\hat{f}(\omega)}d\omega = \int f(x)\overline{f(y)}\int e^{-2\pi i\omega(x-y)}d\omega dxdy = \int f(x)\overline{f(x)}dx = \|f\|_2^2\)，因 \(\int e^{-2\pi i\omega(x-y)}d\omega = \delta(x-y)\)。故 \(\mathcal{F}\) 在 \(L^1\cap L^2\) 上是等距的，由稠密性唯一延拓为 \(L^2\) 上的酉算子。\(\blacksquare\) 。

\textbf{定义 282.2}（卷积）：\(f, g \in L^1(\mathbb{R}^d)\) 的卷积定义为

\[(f * g)(x) = \int_{\mathbb{R}^d} f(x - y) g(y) dy\]

\textbf{定理 282.2}（卷积定理）：\(\mathcal{F}[f * g] = \mathcal{F}[f] \cdot \mathcal{F}[g]\)。卷积在频域中对应于逐点乘积------这是线性时不变系统分析的核心原理

\textbf{证明}：\(\mathcal{F}[f*g](\omega) = \iint f(x-y)g(y)e^{-2\pi i\omega x}dydx = \int g(y)e^{-2\pi i\omega y}\int f(x-y)e^{-2\pi i\omega(x-y)}dxdy = \hat{f}(\omega)g(y)e^{-2\pi i\omega y}dy\)，内层积分 \(=\hat{f}(\omega)\)，故 \(\mathcal{F}[f*g] = \hat{f}\cdot\hat{g}\)。LTI 系统 \(y = h*x\) 在频域为 \(\hat{y} = \hat{h}\cdot\hat{x}\)------卷积变乘积。\(\blacksquare\) 。

\subsection{282.2 离散 Fourier 变换}\label{ux79bbux6563-fourier-ux53d8ux6362}

\textbf{定义 282.3}（离散 Fourier 变换 / DFT）：对 \(\mathbf{x} = (x_0, \ldots, x_{N-1}) \in \mathbb{C}^N\)，其 DFT 为

\[\hat{x}_k = \sum_{n=0}^{N-1} x_n e^{-2\pi i k n / N}, \quad k = 0, \ldots, N-1\]

DFT 矩阵 \(F_N = (e^{-2\pi i kn/N})_{k,n=0}^{N-1}\) 满足 \(F_N F_N^* = N I\)，即 \(\frac{1}{\sqrt{N}} F_N\) 是 \(\mathbb{C}^N\) 上的酉矩阵。

\textbf{定理 282.3}（离散卷积定理）：对循环卷积 \((x * y)_n = \sum_{m=0}^{N-1} x_m y_{(n-m) \bmod N}\)，有 \(\widehat{x * y}_k = \hat{x}_k \hat{y}_k\)

\textbf{证明}：\(\widehat{x*y}_k = \sum_{n=0}^{N-1}(\sum_{m=0}^{N-1}x_m y_{(n-m)\bmod N})e^{-2\pi ikn/N} = \sum_{m=0}^{N-1}x_m e^{-2\pi ikm/N}\sum_{n=0}^{N-1}y_{n-m}e^{-2\pi ik(n-m)/N} = \hat{x}_k\hat{y}_k\)。这是卷积定理的离散版本，DFT 将循环卷积对角化。\(\blacksquare\) 。

\subsection{282.3 Nyquist-Shannon 采样定理}\label{nyquist-shannon-ux91c7ux6837ux5b9aux7406}

\textbf{定理 282.4}（Nyquist-Shannon 采样定理）：设 \(f \in L^2(\mathbb{R})\) 是带限函数，即 \(\operatorname{supp} \hat{f} \subset [-B, B]\)。若以采样间隔 \(\Delta t \leq 1/(2B)\) 均匀采样，则 \(f\) 可由样本 \(\{f(n\Delta t)\}_{n \in \mathbb{Z}}\) 完全重建：

\[f(t) = \sum_{n=-\infty}^\infty f(n\Delta t) \cdot \operatorname{sinc}\left(\frac{t - n\Delta t}{\Delta t}\right), \quad \operatorname{sinc}(x) = \frac{\sin(\pi x)}{\pi x}\]

其中 \(2B\) 称为 Nyquist 率。

\textbf{证明}：Dirac 梳 \(\operatorname{III}_{\Delta t}(t) = \sum_{n\in\mathbb{Z}} \delta(t-n\Delta t)\) 的 Fourier 变换为 \(\frac{1}{\Delta t}\operatorname{III}_{1/\Delta t}(\omega) = \frac{1}{\Delta t}\sum_{k\in\mathbb{Z}} \delta(\omega - k/\Delta t)\)。采样信号 \(f_s(t) = f(t)\operatorname{III}_{\Delta t}(t)\) 的 Fourier 变换由卷积定理得 \(\widehat{f_s}(\omega) = \frac{1}{\Delta t}\sum_{k\in\mathbb{Z}} \hat{f}(\omega - k/\Delta t)\)。当 \(\Delta t \leq 1/(2B)\) 时，项 \(\hat{f}(\omega - k/\Delta t)\) 的支集互不相交。在频域乘以矩形窗 \(\operatorname{rect}(\omega\Delta t)\) 提取 \(k=0\) 项：\(\hat{f}(\omega) = \Delta t\cdot\widehat{f_s}(\omega)\cdot\operatorname{rect}(\omega\Delta t)\)。取 Fourier 逆变换：\(f(t) = \sum_n f(n\Delta t)\operatorname{sinc}((t-n\Delta t)/\Delta t)\)。\(\blacksquare\)

\textbf{定理 282.5}（采样定理的 Hilbert 空间表述）：带限函数空间 \(\mathcal{B}_B = \{f \in L^2(\mathbb{R}) : \operatorname{supp} \hat{f} \subset [-B, B]\}\) 是 \(L^2(\mathbb{R})\) 的闭子空间（Paley-Wiener 空间）。函数族 \(\{\operatorname{sinc}(t/\Delta t - n)\}_{n \in \mathbb{Z}}\) 构成 \(\mathcal{B}_B\) 的规范正交基（当 \(\Delta t = 1/(2B)\) 时），采样值 \(\{f(n\Delta t)\}\) 是该基下的展开系数

\textbf{证明}：当 \(\Delta t = 1/(2B)\) 时，\(\phi_n(t) = \sqrt{2B}\operatorname{sinc}(2B(t - n/(2B)))\)。内积 \(\langle\phi_n,\phi_m\rangle = 2B\int\operatorname{sinc}(2Bt-n)\operatorname{sinc}(2Bt-m)dt\)。由 Parseval，\(\langle\phi_n,\phi_m\rangle = \frac{1}{2B}\int_{-B}^B e^{2\pi i\omega(m-n)/(2B)}d\omega = \delta_{mn}\)。故 \(\{\phi_n\}\) 是 \(\mathcal{B}_B\) 的规范正交基。\(f = \sum_n \langle f,\phi_n\rangle\phi_n\)，且 \(\langle f,\phi_n\rangle = \sqrt{2B}f(n/(2B))\)。\(\blacksquare\) 。

\hrulefill

\hrulefill

\hrulefill

\hrulefill

\hrulefill

\hrulefill

\section{第202章：多分辨分析与小波理论}\label{ux7b2c202ux7ae0ux591aux5206ux8fa8ux5206ux6790ux4e0eux5c0fux6ce2ux7406ux8bba}

小波分析是 20 世纪 80 年代调和分析最重要的突破，它解决了 Fourier 分析的根本缺陷：Fourier 变换将信号完全分解为频率分量却丧失时间局部性，而短时 Fourier 变换（Gabor 变换）的时间-频率分辨率受 Heisenberg 不确定性原理的刚性限制。小波变换通过同时缩放和平移母小波 \(\psi\)，实现了''数学显微镜''效应------高频分量对应窄窗口（精细时间分辨率），低频分量对应宽窗口（精细频率分辨率），从而自适应地匹配信号的多尺度结构。Stephane Mallat 和 Yves Meyer 在 1986-1989 年间创立的多分辨分析（MRA）提供了构造正交小波基的统一代数框架：将 \(L^2(\mathbb{R})\) 分解为嵌套的闭子空间序列 \(\cdots \subset V_{-1} \subset V_0 \subset V_1 \subset \cdots\)，再利用两尺度差分方程系统地生成 Daubechies 紧支撑正交小波。本章从连续小波变换及其允许性条件出发，引入框架理论与 Riesz 基，最终通过多分辨分析建立离散小波变换的完整数学体系。小波理论直接引向非均匀采样与压缩感知（第149章），并在图像压缩（JPEG 2000）和数值分析（自适应小波 Galerkin 方法）中有广泛应用。

\subsection{283.1 连续小波变换}\label{ux8fdeux7eedux5c0fux6ce2ux53d8ux6362}

\textbf{定义 283.1}（连续小波变换 / CWT）：设 \(\psi \in L^2(\mathbb{R})\) 满足\textbf{允许性条件}

\[\int_{\mathbb{R}} \frac{|\hat{\psi}(\omega)|^2}{|\omega|} d\omega < \infty\]

对 \(f \in L^2(\mathbb{R})\)，其连续小波变换定义为

\[W_f(a, b) = \frac{1}{\sqrt{|a|}} \int_{\mathbb{R}} f(t) \overline{\psi\left(\frac{t-b}{a}\right)} dt\]

其中 \(a \in \mathbb{R} \setminus \{0\}\) 为尺度参数，\(b \in \mathbb{R}\) 为平移参数。\(\psi\) 称为母小波。

\textbf{定理 283.1}（Calderon 重构公式）：若 \(\psi\) 满足允许性条件，则对任意 \(f \in L^2(\mathbb{R})\)，

\[f(t) = \frac{1}{C_\psi} \int_{-\infty}^\infty \int_{-\infty}^\infty W_f(a, b) \frac{1}{\sqrt{|a|}} \psi\left(\frac{t-b}{a}\right) \frac{da db}{a^2}\]

其中 \(C_\psi = \int_{\mathbb{R}} \frac{|\hat{\psi}(\omega)|^2}{|\omega|} d\omega\)。CWT 是 \(L^2(\mathbb{R})\) 上的可逆变换

\textbf{证明}：由 Parseval，\(W_f(a,b) = \frac{1}{\sqrt{|a|}}\int\hat{f}(\omega)\overline{\hat{\psi}(a\omega)}e^{ib\omega}d\omega\)。代入重构公式右侧：\(\iint W_f(a,b)\frac{1}{\sqrt{|a|}}\psi(\frac{t-b}{a})\frac{da db}{a^2} = \frac{1}{C_\psi}\int\hat{f}(\omega)e^{i\omega t}\int\frac{|\hat{\psi}(a\omega)|^2}{|a|}da d\omega\)。内层积分 \(\int_0^\infty \frac{|\hat{\psi}(a\omega)|^2}{a}da = C_\psi\)（与 \(\omega\) 无关，由允许性条件），故右侧 \(= \frac{1}{C_\psi}\int\hat{f}(\omega)C_\psi e^{i\omega t}d\omega = f(t)\)。\(\blacksquare\) 。

\textbf{定义 283.2}（时频分析）：小波变换的时频窗为

\[\left[b + a t^* - a \Delta_\psi, b + a t^* + a \Delta_\psi\right] \times \left[\frac{\omega^*}{a} - \frac{\Delta_{\hat{\psi}}}{a}, \frac{\omega^*}{a} + \frac{\Delta_{\hat{\psi}}}{a}\right]\]

其中 \(t^*, \omega^*\) 和 \(\Delta_\psi, \Delta_{\hat{\psi}}\) 分别为 \(\psi\) 及其 Fourier 变换的中心和半径。时间窗宽度 \(\propto |a|\)，频率窗宽度 \(\propto 1/|a|\)------满足 Heisenberg 不确定原理 \(\Delta t \cdot \Delta \omega \geq 1/2\)。这一恒 \(Q\) 滤波特性使小波变换在分析瞬态和多尺度结构时具有天然优势。

\subsection{283.2 多分辨分析（MRA）}\label{ux591aux5206ux8fa8ux5206ux6790mra}

\textbf{定义 283.3}（多分辨分析 / MRA）：\(L^2(\mathbb{R})\) 的一个多分辨分析是一列嵌套闭子空间 \(\cdots \subset V_{-1} \subset V_0 \subset V_1 \subset \cdots\) 满足：

\begin{enumerate}
\def\labelenumi{\arabic{enumi}.}
\tightlist
\item
  \textbf{稠密性与分离性}：\(\overline{\bigcup_{j \in \mathbb{Z}} V_j} = L^2(\mathbb{R})\)，\(\bigcap_{j \in \mathbb{Z}} V_j = \{0\}\)
\item
  \textbf{二进尺度关系}：\(f(t) \in V_j \iff f(2t) \in V_{j+1}\)
\item
  \textbf{平移不变性}：\(f(t) \in V_0 \iff f(t - k) \in V_0\)，\(\forall k \in \mathbb{Z}\)
\item
  \textbf{Riesz 基存在性}：存在尺度函数 \(\phi \in V_0\)，使得 \(\{\phi(t - k)\}_{k \in \mathbb{Z}}\) 构成 \(V_0\) 的 Riesz 基
\end{enumerate}

MRA 的核心结构：定义小波空间 \(W_j\) 为 \(V_j\) 在 \(V_{j+1}\) 中的正交补（\(V_{j+1} = V_j \oplus W_j\)），递归分解得

\[L^2(\mathbb{R}) = \overline{\bigoplus_{j \in \mathbb{Z}} W_j}\]

\textbf{定义 283.4}（两尺度方程）：尺度函数 \(\phi\) 和小波函数 \(\psi\) 满足

\[\phi(t) = \sqrt{2} \sum_{k \in \mathbb{Z}} h_k \phi(2t - k), \quad \psi(t) = \sqrt{2} \sum_{k \in \mathbb{Z}} g_k \phi(2t - k)\]

其中 \(g_k = (-1)^k \overline{h_{1-k}}\)。序列 \(\{h_k\}\) 和 \(\{g_k\}\) 分别是低通和高通共轭镜像滤波器（CQF）。

\textbf{定理 283.2}（正交小波基的 Mallat 分解-重构）：由两尺度方程和正交性条件，信号的塔式分解（由细到粗）和重构（由粗到细）可表示为滤波器组卷积与下/上采样的级联

\textbf{证明}：由两尺度方程和正交性 \(\langle\phi(\cdot-k),\phi(\cdot-\ell)\rangle = \delta_{k\ell}\)，滤波器系数满足 \(\sum h_n h_{n+2k} = \delta_{0k}\) 和 \(\sum g_n g_{n+2k} = \delta_{0k}\) 及正交关系 \(\sum h_n g_{n+2k}=0\)。\(c_k^{j-1} = \langle f,\phi_{j-1,k}\rangle\)，代入 \(\phi_{j-1,k}(t) = \sqrt{2}\sum_n h_{n-2k}\phi_{j,n}(t)\) 即得分解公式。重构由正交分解 \(V_j = V_{j-1}\oplus W_{j-1}\) 的对偶基关系导出。\(\blacksquare\) ：

分解：\(c_k^{j-1} = \sum_n \overline{h_{n-2k}} c_n^j\)，\(d_k^{j-1} = \sum_n \overline{g_{n-2k}} c_n^j\)

重构：\(c_n^j = \sum_k h_{n-2k} c_k^{j-1} + \sum_k g_{n-2k} d_k^{j-1}\)

其中 \(c^j\) 为第 \(j\) 层尺度系数，\(d^j\) 为第 \(j\) 层小波系数。

\textbf{定义 283.5}（消失矩）：小波 \(\psi\) 具有 \(N\) 阶消失矩，若

\[\int_{\mathbb{R}} t^k \psi(t) dt = 0, \quad k = 0, 1, \ldots, N-1\]

消失矩阶数与尺度函数 \(\phi\) 的光滑性相关：\(N\) 阶消失矩等价于 \(\hat{\phi}(\omega)\) 在 \(\omega = 0\) 处以 \(N\) 阶零点衰减（Strang-Fix 条件）。

\textbf{定义 283.6}（Daubechies 紧支集正交小波族）：对任意 \(N \geq 1\)，存在紧支集正交小波 \(\psi_N\) 满足：(i) 支集长度为 \(2N-1\)；(ii) \(N\) 阶消失矩；(iii) \(\operatorname{supp} \phi_N = [0, 2N-1]\)。\(N = 1\) 即为 Haar 小波。

\subsection{283.3 Besov 空间与小波系数刻画}\label{besov-ux7a7aux95f4ux4e0eux5c0fux6ce2ux7cfbux6570ux523bux753b}

\textbf{定义 283.7}（Besov 空间------差分刻画）：设 \(s > 0\)，\(1 \leq p, q \leq \infty\)。记 \(\Delta_h^r f(x) = \sum_{k=0}^r \binom{r}{k} (-1)^{r-k} f(x + kh)\) 为 \(r\) 阶差分（\(r = \lfloor s \rfloor + 1\)）。\(f \in L^p(\mathbb{R}^d)\) 属于 Besov 空间 \(B_{p,q}^s(\mathbb{R}^d)\)，若其半范数

\[|f|_{B_{p,q}^s} = \left(\int_0^\infty \left(t^{-s} \sup_{|h| \leq t} \|\Delta_h^r f\|_{L^p}\right)^q \frac{dt}{t}\right)^{1/q} < \infty\]

（\(q = \infty\) 时取 \(\sup\)）。范数为 \(\|f\|_{B_{p,q}^s} = \|f\|_{L^p} + |f|_{B_{p,q}^s}\)。

\textbf{定理 283.3}（Besov 空间的小波系数刻画）：设 \(\{\psi_{j,k}\}\) 为充分光滑的正交小波基。则 \(f \in B_{p,q}^s(\mathbb{R}^d)\) 等价于

\[\|f\|_{B_{p,q}^s} \asymp \|c_{0,\cdot}\|_{\ell^p} + \left\|\left(2^{j(s + d/2 - d/p)} \|d_{j,\cdot}\|_{\ell^p}\right)_{j \geq 0}\right\|_{\ell^q}\]

其中 \(d_{j,k} = \langle f, \psi_{j,k} \rangle\)，\(c_{0,k} = \langle f, \phi_{0,k} \rangle\)。即小波系数在不同尺度上的 \(\ell^p\) 范数以 \(2^{-js}\) 速率衰减。特例：\(B_{2,2}^s = H^s\)（Sobolev 空间），\(B_{\infty,\infty}^s = C^s\)（Holder 空间，\(s\) 非整数时）

\textbf{证明}：小波基 \(\{\psi_{j,k}\}\) 构成 \(B_{p,q}^s\) 的无条件基。充分性：由 Jackson 不等式，\(\|f-P_{V_j}f\|_p \leq C 2^{-js}|f|_{B_{p,q}^s}\)（\(P_{V_j}\) 为投影），故小波系数满足 \(\|d_j\|_{\ell^p} \leq C 2^{-j(s+d/2-d/p)}\)。必要性：由 Bernstein 不等式，\(\|f-P_{V_j}f\|_p \leq C(\sum_{\ell\geq j}\|d_\ell\|_{\ell^p}^q 2^{\ell q(d/p-d/2)})^{1/q}\)，结合系数衰减得 Besov 范数等价。\(\blacksquare\) 。

\subsection{小波理论的现代发展与多尺度分析}\label{ux5c0fux6ce2ux7406ux8bbaux7684ux73b0ux4ee3ux53d1ux5c55ux4e0eux591aux5c3aux5ea6ux5206ux6790}

小波理论在 20 世纪末的爆发式发展不仅为调和分析提供了新的范式，也在多个应用领域产生了深远影响。\textbf{多尺度分析}（MRA, Mallat 1989）是小波理论的核心框架，其思想可以追溯到计算机视觉中的金字塔表示（Burt-Adelson 1983）和信号处理中的子带编码（subband coding）。MRA 的优美之处在于它将小波构造问题转化为滤波器设计问题------通过两尺度方程中的低通滤波器 \(\{h_k\}\) 和高通滤波器 \(\{g_k\}\)（满足正交条件和消失矩条件），可以系统地构造具有所需正则性、对称性和消失矩的正交小波基。\textbf{Daubechies 紧支集小波}（1988）是这一构造范式的杰出成果：对任意 \(N \geq 1\)，存在支集长度为 \(2N-1\) 的正交小波，具有 \(N\) 阶消失矩，且光滑性随 \(N\) 增加。Daubechies 小波的构造依赖于将半带滤波器（half-band filter）的谱因子分解为最小相位和最大相位部分，这是通过 Bezout 恒等式和 Riesz 引理完成的。\textbf{双正交小波}（Cohen-Daubechies-Feauveau 1992）放松了正交性约束，允许分析小波和综合小波不同，从而获得了对称性（线性相位）这一重要性质------这在图像压缩中对于避免边界伪影至关重要。JPEG 2000 标准采用的是 Cohen-Daubechies-Feauveau 双正交小波（CDF 9/7 和 CDF 5/3）。\textbf{提升格式}（lifting scheme, Sweldens 1996）是小波构造的另一种范式，它将小波变换分解为简单的提升步骤（预测和更新），在有限精度算术下可逆，且可构造第二代小波（second generation wavelets）------无需 Fourier 变换，适用于非均匀采样、曲面和三角剖分等非欧几里得区域。\textbf{复小波}（Kingsbury 2001）和\textbf{双树复小波变换}（DTCWT）通过引入复值小波解决了实小波的平移敏感性（shift variance）和方向选择性差等问题，在图像纹理分析和运动估计中表现优异。在\textbf{机器学习}领域，\textbf{散射变换}（scattering transform, Mallat 2012）通过小波卷积和非线性模运算的级联，构造了平移不变和 Lipschitz 稳定的特征表示，已在图像分类和音频分析中取得与深度学习可比的结果。

\hrulefill

\hrulefill

\section{第203章：变分法与偏微分方程}\label{ux7b2c203ux7ae0ux53d8ux5206ux6cd5ux4e0eux504fux5faeux5206ux65b9ux7a0b}

变分法是数学物理最古老也最富活力的方法之一：从 Euler 和 Lagrange 在 18 世纪提出的最小作用量原理，到 20 世纪末 Rudin-Osher-Fatemi（ROF）模型开创的现代图像处理，其核心始终是将物理定律或优化准则表述为泛函的极值问题。在偏微分方程的语境下，变分法的优势在于它将''解方程''转化为''最小化能量''------这一视角不仅提供了解的存在性证明的新途径（直接变分法），还自然引入了弱解、凸对偶和正则性理论。本章聚焦于变分法在信号与图像处理中的现代应用：全变差（Total Variation）正则化将图像的去噪和去模糊问题归结为凸泛函的极小化，其 Euler-Lagrange 方程涉及非线性散度算子；非局部变分法则突破了经典梯度算子的局域性限制，利用图像的自相似性（patch-based regularization）实现更好的纹理保持。这些方法与第 147 章的小波稀疏表示和第 149 章的压缩感知相辅相成：\(\ell_1\) 和 TV 正则化的理论分析共享凸分析和 Bregman 迭代的统一框架。

\subsection{284.1 全变差泛函与 ROF 模型}\label{ux5168ux53d8ux5deeux6cdbux51fdux4e0e-rof-ux6a21ux578b}

\textbf{定义 284.1}（全变差 / Total Variation）：设 \(\Omega \subset \mathbb{R}^d\) 为开集。函数 \(u \in L^1(\Omega)\) 的全变差定义为

\[TV(u) = \sup\left\{\int_\Omega u \cdot \operatorname{div} \phi \, dx : \phi \in C_c^1(\Omega, \mathbb{R}^d), \|\phi(x)\|_\infty \leq 1, \forall x \in \Omega\right\}\]

若 \(u \in C^1(\bar{\Omega})\)，则 \(TV(u) = \int_\Omega |\nabla u| dx\)。\(TV(u)\) 是 \(W^{1,1}\) 半范数在有界变差函数空间 \(BV(\Omega)\) 上的推广------它允许函数具有跳跃间断，因为阶跃函数的 TV 为其跳跃量的绝对值。

\textbf{定理 284.1}（ROF 变分问题 / Rudin-Osher-Fatemi, 1992）：给定 \(f \in L^2(\Omega)\)，考虑以下极小化问题

\[\min_{u \in BV(\Omega)} \left\{ \int_\Omega |\nabla u| dx + \frac{\lambda}{2} \int_\Omega (u - f)^2 dx \right\}, \quad \lambda > 0\]

第一项 \(TV(u)\) 为正则项（惩罚振荡但允许跳跃），第二项 \(\frac{\lambda}{2}\|u - f\|_{L^2}^2\) 为保真项。解的存在性由 \(BV(\Omega)\) 的紧性（\(BV\) 到 \(L^1\) 的紧嵌入

\textbf{证明}：由 \(TV(u)\) 下半连续性（关于 \(L^1\) 收敛）和 \(L^2\) 范数的弱下半连续性，ROF 泛函是 \(BV\cap L^2\) 上的 coercive 和弱下半连续的。\(BV(\Omega)\) 紧嵌入 \(L^1(\Omega)\)（有界域上）保证极小化序列在 \(L^1\) 中强收敛。结合 \(TV\) 的弱*下半连续性，极限函数是极小元。唯一性由保真项 \(\frac{\lambda}{2}\|u-f\|_2^2\) 的严格凸性保证。\(\blacksquare\) ）保证。

\textbf{定理 284.2}（ROF 模型的 Euler-Lagrange 方程）：形如上的最优性条件为

\[-\operatorname{div}\left(\frac{\nabla u}{|\nabla u|}\right) + \lambda(u - f) = 0\]

第一项 \(-\operatorname{div}(\nabla u / |\nabla u|)\) 是水平集 \(\{u = c\}\) 的平均曲率算子。该 PDE 定义了\textbf{各向异性扩散}

\textbf{证明}：对光滑 \(u\)，TV 项的 Gateaux 导数：\(\frac{d}{d\varepsilon}TV(u+\varepsilon v)|_{\varepsilon=0} = \int \frac{\nabla u\cdot\nabla v}{|\nabla u|}dx = -\int v\operatorname{div}(\frac{\nabla u}{|\nabla u|})dx\)（分部积分）。保真项导数为 \(\lambda(u-f)\)。故最优性条件为 \(-\operatorname{div}(\nabla u/|\nabla u|) + \lambda(u-f) = 0\)。当 \(|\nabla u|=0\) 时需用次微分 \(\partial TV(u)\)（\(L^2\) 中子梯度），形式为 \(\operatorname{div}(p)\)（\(p\in L^\infty\)，\(|p|\leq 1\)，\(p\cdot\nabla u = |\nabla u|\)）。\(\blacksquare\) ：扩散系数 \(1/|\nabla u|\) 在梯度较大处被抑制，在梯度较小处被增强。

\subsection{284.2 Mumford-Shah 泛函与自由边界问题}\label{mumford-shah-ux6cdbux51fdux4e0eux81eaux7531ux8fb9ux754cux95eeux9898}

\textbf{定义 284.2}（Mumford-Shah 泛函，1989）：对 \(u \in SBV(\Omega)\)（特殊有界变差函数空间）和闭集 \(\Gamma \subset \Omega\)（\((d-1)\) 维），定义

\[E(u, \Gamma) = \int_{\Omega \setminus \Gamma} |\nabla u|^2 dx + \mu \cdot \mathcal{H}^{d-1}(\Gamma) + \lambda \int_\Omega (u - f)^2 dx\]

其中 \(\mathcal{H}^{d-1}(\Gamma)\) 是 \(\Gamma\) 的 \((d-1)\) 维 Hausdorff 测度（\(\mathbb{R}^2\) 中为长度），\(u\) 在 \(\Omega \setminus \Gamma\) 上光滑，沿 \(\Gamma\) 允许跳跃。这是一个涉及未知函数和未知边界的同时优化问题。

当 \(\mu \to \infty\) 时，\(\Gamma = \emptyset\)，泛函退化为 \(H^1\) 正则化（Sobolev 空间光滑性）；当 \(\mu \to 0\) 时，允许任意复杂的边界集。

\textbf{定义 284.3}（Chan-Vese 水平集方法------Mumford-Shah 的变体）：用水平集函数 \(\phi: \Omega \to \mathbb{R}\) 隐式表示边界 \(\Gamma = \{x : \phi(x) = 0\}\)。对二区域分段常值近似，能量泛函为

\[E(c_1, c_2, \phi) = \lambda_1 \int_{\{\phi > 0\}} (f - c_1)^2 dx + \lambda_2 \int_{\{\phi < 0\}} (f - c_2)^2 dx + \mu \int_\Omega |\nabla H(\phi)| dx + \nu \int_\Omega H(\phi) dx\]

其中 \(H\) 为 Heaviside 函数（或其正则化光滑逼近）。对应的梯度下降流（PDE 演化）为

\[\frac{\partial \phi}{\partial t} = \delta(\phi) \left[\mu \operatorname{div}\left(\frac{\nabla \phi}{|\nabla \phi|}\right) - \nu - \lambda_1 (f - c_1)^2 + \lambda_2 (f - c_2)^2\right]\]

其中 \(\delta(\phi)|\nabla\phi|\) 将演化限制在零水平集附近------这是水平集方法的典型结构。

\subsection{284.3 几何 PDE 与变分光流}\label{ux51e0ux4f55-pde-ux4e0eux53d8ux5206ux5149ux6d41}

\textbf{定义 284.4}（Bertalmio-Sapiro 等照度线传播方程）：设 \(D \subset \Omega\) 为待修复区域，信息沿等照度线方向 \(\nabla^\perp u = (-\partial_y u, \partial_x u)\) 传播：

\[\frac{\partial u}{\partial t} = \nabla(\Delta u) \cdot \nabla^\perp u\]

这是一个三阶几何 PDE------Laplacian 的梯度沿等照度线的切向输送。

\textbf{定义 284.5}（Horn-Schunck 变分问题，1981）：给定图像序列 \(I(x, y, t)\)，求向量场 \(\mathbf{v} = (v_1, v_2)\) 极小化

\[\min_{\mathbf{v}} \left\{ \int_\Omega (\nabla I \cdot \mathbf{v} + I_t)^2 dx dy + \alpha \int_\Omega (|\nabla v_1|^2 + |\nabla v_2|^2) dx dy \right\}, \quad \alpha > 0\]

第一项为光流约束方程 \(I_x v_1 + I_y v_2 + I_t = 0\) 的 \(L^2\) 残差，第二项为 \(H^1\)-半范数正则化。Euler-Lagrange 方程为椭圆型 PDE 系统：

\[(I_x^2 + \alpha \Delta)v_1 + I_x I_y v_2 = -I_x I_t, \quad I_x I_y v_1 + (I_y^2 + \alpha \Delta)v_2 = -I_y I_t\]

\hrulefill

\subsection{变分法与图像处理的数学统一}\label{ux53d8ux5206ux6cd5ux4e0eux56feux50cfux5904ux7406ux7684ux6570ux5b66ux7edfux4e00}

ROF 模型和 Mumford-Shah 泛函是变分法在图像处理中应用的两个里程碑，它们共同建立了\textbf{变分图像处理}的数学框架。变分法的核心思想是将图像处理问题表述为能量泛函的极小化，其中正则项编码了关于解的先验知识（如光滑性、分段Constant性、边缘保持性），保真项确保解与观测数据一致。\textbf{Gamma 收敛}（De Giorgi 1970s）是变分分析的基本工具，它提供了渐近行为分析的理论框架：当各向异性扩散系数趋于零时，ROF 模型的 Gamma 极限给出了分段常数函数的先验。\textbf{DTV 模型}（duality-based TV）通过 Chambolle 的对偶投影算法（2004）将 TV 正则化问题转化为约束二次规划，实现了 \(O(N \log N)\) 的快速求解。\textbf{非局部变分方法}（Gilboa-Osher 2008）将 TV 正则化推广到非局部图 Laplacian 框架，为具有重复纹理和自相似性的图像提供了更优的恢复效果。\textbf{Bregman 迭代}（Osher et al.~2005）和\textbf{分裂 Bregman 方法}（Goldstein-Osher 2009）是求解 TV 正则化问题的高效算法，通过将原问题分解为交替的子问题（\(L^2\) 投影和 shrink 算子），实现了对大规模图像的快速处理。\textbf{彩色图像处理}中，TV 正则化被推广为\textbf{向量值 TV}（Bresson-Chan 2008），通过将不同颜色通道的梯度耦合来避免颜色伪影。在\textbf{计算机视觉}中，变分方法已广泛应用于立体匹配、光流估计、三维重建和形状优化，构成了现代视觉系统的基础数学工具。变分图像处理与\textbf{深度学习的联系}：深层卷积神经网络（CNN）中的激活函数（如 ReLU）和池化操作可视为非线性扩散的离散化，而残差网络（ResNet）的结构与常微分方程的离散化存在形式对应，这为理解深度学习的数学原理提供了变分视角。

\hrulefill

\hrulefill

\hrulefill

\hrulefill

\hrulefill

\section{第204章：压缩感知与稀疏恢复理论}\label{ux7b2c204ux7ae0ux538bux7f29ux611fux77e5ux4e0eux7a00ux758fux6062ux590dux7406ux8bba}

压缩感知（Compressed Sensing）是 21 世纪初由 Candes、Romberg、Tao 和 Donoho 等人创立的全新信号获取范式，它革命性地突破了 Shannon-Nyquist 采样定理的限制。经典采样理论要求采样率至少为信号最高频率的两倍，而压缩感知的关键洞见在于：如果信号在某个基下是稀疏的（即大部分系数为零），那么只需 \(O(s \log(N/s))\) 次非自适应线性测量即可精确重构该信号，远少于信号维度 \(N\)。这一理论的核心是 \(\ell_1\) 极小化与 \(\ell_0\) 极小化在稀疏条件下的等价性------这是压缩感知与凸优化的深刻交叉。从数学角度看，\(\ell_0\) 最小化是 NP 难的组合优化问题，但 Donoho 和 Candes 等人的工作证明，当测量矩阵满足限制等距性质（RIP）时，\(\ell_1\) 松弛给出了精确解。本章从稀疏性的严格定义出发，建立 RIP 条件与 \(\ell_1\)-\(\ell_0\) 等价性的完整证明链，并讨论 Bregman 迭代和交替方向乘子法（ADMM）等高效算法。压缩感知与第 146 章（采样）、第 147 章（小波稀疏表示）和第 148 章（TV 变分模型）构成了稀疏建模的完整知识体系。

\subsection{\texorpdfstring{285.1 稀疏性与 \(\ell_0\) 问题}{285.1 稀疏性与 \textbackslash ell\_0 问题}}\label{ux7a00ux758fux6027ux4e0e-ell_0-ux95eeux9898}

\textbf{定义 285.1}（稀疏向量）：\(\mathbf{x} \in \mathbb{R}^N\) 为 \(s\)-稀疏的，若 \(\|\mathbf{x}\|_0 := |\operatorname{supp}(\mathbf{x})| \leq s \ll N\)。更一般地，\(\mathbf{x}\) 在基 \(\Psi\) 下的表示 \(\boldsymbol{\alpha} = \Psi \mathbf{x}\) 可能是稀疏的。

\textbf{定义 285.2}（压缩感知问题）：给定测量矩阵 \(A \in \mathbb{R}^{m \times N}\)（\(m \ll N\)）和测量值 \(\mathbf{y} = A\mathbf{x}\)（或含噪版本 \(\mathbf{y} = A\mathbf{x} + \mathbf{e}\)），从 \(\mathbf{y}\) 中重建 \(\mathbf{x}\)。最直接的形式为 \(\ell_0\) 极小化：

\[\min_{\mathbf{x} \in \mathbb{R}^N} \|\mathbf{x}\|_0 \quad \text{s.t.} \quad A\mathbf{x} = \mathbf{y}\]

\textbf{命题 285.1}（\(\ell_0\) 恢复的唯一性------Spark）：若 \(\operatorname{spark}(A) > 2s\)（Spark 为 \(A\) 的最小线性相关列数），则至多存在一个 \(s\)-稀疏解满足 \(A\mathbf{x} = \mathbf{y}\)。

\(\ell_0\) 极小化是 NP-hard 的组合优化问题（在所有 \(\binom{N}{s}\) 个可能支撑集上搜索）。

\subsection{\texorpdfstring{285.2 \(\ell_1\) 松弛与限制等距性质（RIP）}{285.2 \textbackslash ell\_1 松弛与限制等距性质（RIP）}}\label{ell_1-ux677eux5f1bux4e0eux9650ux5236ux7b49ux8dddux6027ux8d28rip}

\textbf{定理 285.1}（压缩感知基本定理 / \(\ell_1\) 松弛的合理性）：若 \(A\) 满足适当的 RIP 条件，则凸 \(\ell_1\) 极小化

\[\min_{\mathbf{x}} \|\mathbf{x}\|_1 \quad \text{s.t.} \quad A\mathbf{x} = \mathbf{y}\]

（含噪情形：\(\min_{\mathbf{x}} \|\mathbf{x}\|_1 \text{ s.t. } \|A\mathbf{x} - \mathbf{y}\|_2 \leq \epsilon\)，即基追踪去噪 / BPDN）能以高概率精确重建 \(s\)-稀疏信号 \(\mathbf{x}\)，仅需测量数

\[m \geq C \cdot s \cdot \log(N/s)\]

其中 \(C\) 为通用常数。对比 Nyquist 采样率（需 \(m = \mathcal{O}(N)\)），压缩感知的测量数仅与稀疏度 \(s\) 对数依赖于信号维数 \(N\)

\textbf{证明}：由 RIP 条件 \(\delta_{2s}<\sqrt{2}-1\)，对任意 \(s\)-稀疏 \(\mathbf{x}\)，\(\ell_1\) 极小化解 \(\hat{\mathbf{x}}\) 满足 \(\|\hat{\mathbf{x}}-\mathbf{x}\|_2 \leq C\|\mathbf{x}-\mathbf{x}_s\|_1/\sqrt{s}\)。证明：设 \(\mathbf{h} = \hat{\mathbf{x}}-\mathbf{x}\)，将 \(\mathbf{h}\) 分解为不相交支撑块 \(T_0,T_1,\ldots\)，\(T_0\) 为 \(\mathbf{x}\) 的 \(s\) 个最大系数位置。由 RIP 和 cone 约束 \(\|\mathbf{h}_{T_0^c}\|_1 \leq \|\mathbf{h}_{T_0}\|_1\)，迭代不等式得 \(\|\mathbf{h}\|_2 \leq C\|\mathbf{x}_{T_0^c}\|_1/\sqrt{s}\)。随机矩阵的 RIP 由浓度不等式（矩阵 Chernoff）证明。\(\blacksquare\) 。

\textbf{定义 285.3}（限制等距性质 / RIP, Candes-Tao 2005）：矩阵 \(A \in \mathbb{R}^{m \times N}\) 满足 \(s\) 阶 RIP，若存在 \(\delta_s \in (0, 1)\) 使得对\textbf{所有} \(s\)-稀疏向量 \(\mathbf{x} \in \mathbb{R}^N\)，

\[(1 - \delta_s)\|\mathbf{x}\|_2^2 \leq \|A\mathbf{x}\|_2^2 \leq (1 + \delta_s)\|\mathbf{x}\|_2^2\]

RIP 确保 \(A\) 在 \(s\)-稀疏向量的子集上近似为等距映射，从而不破坏稀疏信号的几何结构。\(\delta_s\) 越小，\(A\) 在稀疏子空间上越接近正交投影。

\textbf{定理 285.2}（RIP 的随机矩阵实现）：以下随机矩阵以高概率满足 \(\delta_{2s} \leq \delta\)，所需测量数 \(m = \mathcal{O}(s \log(N/s))\)： - 独立同分布 Gaussian 矩阵：\(A_{ij} \sim \mathcal{N}(0, 1/m)\) - Bernoulli 矩阵：\(A_{ij} = \pm 1/\sqrt{m}\) 等概率 - 部分随机 Fourier 矩阵：随机选取 \(m\) 行 DFT 矩阵

\textbf{证明}：对 Gaussian \(A\)（\(A_{ij}\sim\mathcal{N}(0,1/m)\)），固定 \(s\)-稀疏向量 \(\mathbf{x}\)，\(\|A\mathbf{x}\|_2^2 = \frac{1}{m}\sum_{k=1}^m \langle\mathbf{a}_k,\mathbf{x}\rangle^2\) 是 i.i.d. \(\chi^2\) 变量之和。由 Bernstein 浓度不等式，\(\mathbb{P}(|\|A\mathbf{x}\|_2^2-\|\mathbf{x}\|_2^2|>\delta\|\mathbf{x}\|_2^2) \leq 2e^{-c_\delta m}\)。对 \(\binom{N}{s}\) 个 \(s\)-稀疏支撑集取联合界，当 \(m \geq Cs\log(N/s)/\delta^2\) 时 RIP 以高概率成立。\(\blacksquare\)

\textbf{定理 285.3}（\(\ell_1\) 恢复的稳定性与稳健性）：若 \(A\) 满足 \(\delta_{2s} < \sqrt{2} - 1\)，则对含噪观测 \(\mathbf{y} = A\mathbf{x} + \mathbf{e}\)（\(\|\mathbf{e}\|_2 \leq \epsilon\)），BPDN 解 \(\hat{\mathbf{x}}\) 满足

\textbf{证明}：设 \(\mathbf{h} = \hat{\mathbf{x}}-\mathbf{x}\)。由可行性条件 \(\|A\mathbf{h}\|_2 \leq 2\epsilon\)。cone 约束：\(\|\mathbf{h}_{T_0^c}\|_1 \leq \|\mathbf{h}_{T_0}\|_1 + 2\|\mathbf{x}_{T_0^c}\|_1\)（\(T_0\) 为 \(\mathbf{x}\) 的最大 \(s\) 个系数索引）。将 \(\mathbf{h}_{T_0^c}\) 按系数绝对值递减排列并分解为支撑 \(T_1,T_2,\ldots\)（\(|T_j|=s\)）。利用 RIP 不等式 \(\|A\mathbf{h}_{T_0\cup T_1}\|_2^2 \geq (1-\delta_{2s})\|\mathbf{h}_{T_0\cup T_1}\|_2^2\) 和 \(\|A\mathbf{h}\|_2 \leq 2\epsilon\)，经迭代放缩并结合 Cauchy-Schwarz 得 \(\|\mathbf{h}\|_2 \leq C_1\epsilon + C_2\|\mathbf{x}-\mathbf{x}_s\|_1/\sqrt{s}\)，其中 \(C_1 = 4\sqrt{1+\delta_{2s}}/(1-(1+\sqrt{2})\delta_{2s})\)，\(C_2 = 2(1-\delta_{2s}+(1+\delta_{2s})C_1)/(1-\delta_{2s})\)。当 \(\delta_{2s}<\sqrt{2}-1\) 时，常数有限，保证稳定恢复。\(\blacksquare\)

\[\|\hat{\mathbf{x}} - \mathbf{x}\|_2 \leq C_1 \epsilon + C_2 \frac{\|\mathbf{x} - \mathbf{x}_s\|_1}{\sqrt{s}}\]

其中 \(\mathbf{x}_s\) 为 \(\mathbf{x}\) 的最佳 \(s\)-项逼近。该误差界揭示了两个核心性质：(i) 稳健性------恢复误差与噪声水平 \(\epsilon\) 线性相关；(ii) 稳定性------即使信号不完全稀疏但可压缩，也能近似恢复。

\subsection{285.3 稀疏恢复的数学方法评述}\label{ux7a00ux758fux6062ux590dux7684ux6570ux5b66ux65b9ux6cd5ux8bc4ux8ff0}

\textbf{定义 285.4}（稀疏建模的变分框架）：一般形式的稀疏恢复问题可写为

\[\min_{\mathbf{x}} \left\{\|A\mathbf{x} - \mathbf{y}\|_2^2 + \lambda \|\mathbf{x}\|_1\right\}\]

这是 \(\ell_2\)-\(\ell_1\) 混合正则化的特例（LASSO 形式）。其最优性条件为 \(A^T(A\mathbf{x} - \mathbf{y}) + \lambda \partial\|\mathbf{x}\|_1 \ni 0\)，其中 \(\partial\|\cdot\|_1\) 为 \(\ell_1\) 范数的次微分。

\textbf{定义 285.5}（分析算子模型 / Co-sparse 模型）：对偶于综合稀疏模型（\(\mathbf{x} = D\boldsymbol{\alpha}\)，\(\|\boldsymbol{\alpha}\|_0\) 小），分析模型建模 \(\|\Omega \mathbf{u}\|_1\) 的稀疏性，其中 \(\Omega \in \mathbb{R}^{P \times N}\)（\(P \gg N\)）为过完备分析算子。例如 TV 范数是梯度算子 \(\nabla\) 作为分析算子的 \(\ell_1\) 稀疏性度量：\(\|\mathbf{u}\|_{TV} = \|\nabla \mathbf{u}\|_1\)。

综合模型直接建模信号本身，分析模型建模信号的变换。两者的数学对偶体现了稀疏表示理论的基本对称性。

\subsection{压缩感知的理论扩展与前沿}\label{ux538bux7f29ux611fux77e5ux7684ux7406ux8bbaux6269ux5c55ux4e0eux524dux6cbf}

压缩感知（Compressed Sensing, CS）自 Candes、Romberg、Tao 和 Donoho 于 2004-2006 年奠基以来，已成为信号处理和应用数学中最重要的范式转变之一。其核心洞察------稀疏性先验可以替代 Nyquist 速率的密集采样------在数学上通过随机矩阵的 RIP 和 \(\ell_1\) 极小化得以严格证明。\textbf{非凸松驰}（如 \(\ell_p\) 极小化，\(0 < p < 1\)）在理论上比 \(\ell_1\) 松弛需要更少的测量数（\(m = O(s + s \log(N/s))\) 相比于 \(m = O(s \log(N/s))\)），但求解全局最优解是 NP-hard 问题。\textbf{矩阵补全}（Matrix Completion, Candes-Recht 2009）是压缩感知向矩阵情形的推广：若矩阵 \(X \in \mathbb{R}^{n_1 \times n_2}\) 是低秩的（\(\operatorname{rank}(X) = r \ll \min(n_1, n_2)\)），则通过核范数极小化 \(\min \|X\|_*\) 约束于观测到的元素，可以仅从 \(O(r n \log n)\) 个随机观测中精确恢复 \(X\)------这是 Netflix 推荐系统（协同过滤）和量子态层析成像的数学基础。\textbf{低秩加稀疏分解}（Robust PCA, Candes et al.~2011）将矩阵分解为低秩部分和稀疏部分：\(\min_{L,S} \|L\|_* + \lambda \|S\|_1\) 约束于 \(X = L + S\)，在视频背景建模和人脸识别中实现了鲁棒的主成分提取。\textbf{相位恢复}（Phase Retrieval, Candes-Strohmer-Voroninski 2013）从 Fourier 变换的幅度 \(|Ax|\) 中恢复信号 \(x\)，通过 Wirtinger 流算法或 PhaseLift（半定规划松弛）实现，这在 X 射线晶体学和相干衍射成像中至关重要。\textbf{1-bit 压缩感知}（Boufounos-Baraniuk 2008）仅保留测量值的符号信息，通过 \(\ell_1\) 极小化从 \(y = \operatorname{sign}(Ax)\) 中恢复信号，适用于高速 ADC 的量化和分布式传感器网络。\textbf{动态压缩感知}（DCS）和\textbf{时间序列压缩感知}处理时变稀疏信号，通过联合稀疏模型和 Kalman 滤波的 CS 版本实现实时重建。\textbf{深度学习与压缩感知}的融合（如 ISTA-Net, ADMM-Net）通过展开迭代算法为神经网络，以数据驱动的方式学习最优的稀疏正则化和重建算子，在有监督训练下显著超越了传统 CS 方法的重建质量。

\newpage

\chapter{06-分析/05-调和分析/05-经典势论.md}\label{ux5206ux679005-ux8c03ux548cux5206ux679005-ux7ecfux5178ux52bfux8bba.md}

\section{第205章：经典势论与Riesz位势}\label{ux7b2c205ux7ae0ux7ecfux5178ux52bfux8bbaux4e0erieszux4f4dux52bf}

经典势论研究Laplace方程\(\Delta u = 0\)的解（调和函数）以及Newton位势与Riesz位势的深层性质，与复分析、调和分析和概率论（Brown运动）有本质联系。势论起源于18世纪对引力场的研究（Newton位势），后经Gauss、Green、Dirichlet等人的发展，成为分析学中处理椭圆型偏微分方程边界值问题的核心理论。20世纪以来，势论与测度论、容量理论和随机过程的深刻联系使其成为现代分析的基础分支。

\subsection{调和函数与均值性质}\label{ux8c03ux548cux51fdux6570ux4e0eux5747ux503cux6027ux8d28}

\textbf{定义150.1}（调和函数与上调和函数）：\(\mathbb{R}^n\)中开集\(\Omega\)上的\(C^2\)函数\(u\)满足\(\Delta u = \sum_{i=1}^n \frac{\partial^2 u}{\partial x_i^2} = 0\)称为\textbf{调和函数}。若\(u : \Omega \to (-\infty, \infty]\)下半连续（即\(\{u > t\}\)对每个\(t \in \mathbb{R}\)为开集）且满足\(-\Delta u \geq 0\)（分布意义下），则称\(u\)为\textbf{上调和函数}。类似地，若\(-u\)为上调和，则\(u\)为\textbf{下调和函数}。上调和性是调和性的''单侧''推广：调和函数同时为上调和和下调和。

\textbf{调和函数的动机}：调和函数是Dirichlet能量泛函\(\mathcal{E}[u] = \int_\Omega |\nabla u|^2 dx\)的临界点（并给出极小值），也是Brown运动的期望值函数。调和函数的实部和虚部解析性（Weyl引理）说明调和函数实则为解析函数，这是椭圆型方程正则性的典范。

\textbf{定理150.1}（均值性质）：设\(u\)在\(\Omega\)中调和，则对任意球\(B(x,r) \subset\subset \Omega\)， \[u(x) = \frac{1}{|\partial B_r|} \int_{\partial B(x,r)} u(y) dS(y) = \frac{1}{|B_r|} \int_{B(x,r)} u(y) dy\] 其中\(|\partial B_r| = n\omega_n r^{n-1}\)为球面面积，\(|B_r| = \omega_n r^n\)为球体积，\(\omega_n\)为\(n\)维单位球体积。均值性质是调和函数最本质的特征，反之，满足均值性质的连续函数必为调和函数（Koebe逆定理）。

\textbf{证明}：取\(\rho \in (0,r)\)，由Green第二恒等式应用于\(u\)和\(1\)（或更精确地，应用于\(u\)和基本解\(\Phi\)）， \[\int_{\partial B(x,\rho)} \frac{\partial u}{\partial \nu} dS = \int_{B(x,\rho)} \Delta u \,dy = 0\] 由此得球面积分的径向导数为零：\(\frac{d}{d\rho}\left(\rho^{1-n}\int_{\partial B(x,\rho)} u dS\right) = 0\)，故\(\rho^{1-n}\int_{\partial B(x,\rho)} u dS\)与\(\rho\)无关。令\(\rho \to 0\)得\(\lim_{\rho \to 0} \rho^{1-n}\int_{\partial B(x,\rho)} u dS = n\omega_n u(x)\)，故\(u(x) = \frac{1}{|\partial B_r|}\int_{\partial B(x,r)} u dS\)。体积均值由对半径积分：\(\int_{B(x,r)} u dy = \int_0^r \int_{\partial B(x,\rho)} u dS d\rho = u(x) \int_0^r n\omega_n \rho^{n-1} d\rho = u(x) |B_r|\)。\(\blacksquare\)

\textbf{定理150.2}（极大值原理与Harnack不等式）：设\(\Omega\)为连通开集。 (i) （强极大值原理）若调和函数\(u\)在\(\Omega\)内达到极大值或极小值，则\(u\)为常数； (ii) （Harnack不等式）对\(\Omega\)的紧子集\(K\)，存在常数\(C = C(K,\Omega)\)使一切正调和函数\(u > 0\)满足\(\sup_K u \leq C \inf_K u\)。

\textbf{证明}：(i)由均值性质，若\(u\)在\(x_0 \in \Omega\)处达到极大值\(M\)，则对足够小的\(r\)，\(M = u(x_0) = \frac{1}{|B_r|}\int_{B(x_0,r)} u(y) dy\)，而\(u(y) \leq M\)，故\(u(y) = M\)在\(B(x_0,r)\)上几乎处处成立。由连续性，\(u \equiv M\)在\(B(x_0,r)\)上。极大值点集在\(\Omega\)中既开又闭，由连通性即得\(u \equiv M\)。(ii)利用球上Poisson积分核的上下界估计：对\(B(x_0, 2R) \subset \subset \Omega\)，Poisson核满足\(C_1 R^{-n} \leq P(x_0, y) \leq C_2 R^{-n}\)（\(y \in \partial B(x_0, 2R)\)，\(x \in B(x_0, R)\)）。取紧集\(K\)的Harnack链覆盖\(K\)：即存在有限个球\(B_1, \ldots, B_m\)使\(B_i \cap B_{i+1} \neq \emptyset\)，逐次应用每个球上的Harnack不等式\(u(x) \leq C_i u(y)\)（\(x, y \in B_i\)），迭代即得\(\sup_K u \leq (\prod C_i) \inf_K u\)。\(\blacksquare\)

\textbf{定理150.3}（Liouville定理）：\(\mathbb{R}^n\)上有下界（或上界）的调和函数必为常数。特别地，\(\mathbb{R}^n\)上一致有界的调和函数为常数。这由Harnack不等式在\(\mathbb{R}^n\)上取极限\(R \to \infty\)得出：对任意\(x, y \in \mathbb{R}^n\)，取\(R\)足够大使得\(x, y \in B(0, R)\)，则\(u(x) - m \leq C(u(y) - m)\)，令\(R \to \infty\)得\(C \to 1\)，故\(u(x) = u(y)\)。

\subsection{上调和函数与Perron方法}\label{ux4e0aux8c03ux548cux51fdux6570ux4e0eperronux65b9ux6cd5}

\textbf{定义150.2}（上调和函数的性质）：上调和函数\(u\)满足以下等价刻画： (i) 对任意球\(B \subset\subset \Omega\)，\(u(x) \geq \frac{1}{|\partial B|}\int_{\partial B} u dS\)（均值上界）； (ii) 对任意球\(B\)和任意调和函数\(h\)在\(B\)上满足\(h|_{\partial B} = u|_{\partial B}\)，有\(u \geq h\)在\(B\)上； (iii) \(-\Delta u \geq 0\)在分布意义下。

\textbf{Perron方法}：Perron方法（1923）是求解Dirichlet问题的最一般方法。设\(\Omega \subset \mathbb{R}^n\)为有界区域，\(f \in C(\partial\Omega)\)。定义\textbf{Perron下函数族}： \[\mathcal{U}_f = \{v \in \text{上调和}(\Omega) : v \text{在}\Omega\text{上局部有下界}, \liminf_{x \to \xi} v(x) \geq f(\xi) \;\forall \xi \in \partial\Omega\}\] \textbf{Perron解}为\(u(x) = \inf\{v(x) : v \in \mathcal{U}_f\}\)。Perron证明了\(u\)在\(\Omega\)中调和。该解是否连续取到边界值\(f\)取决于边界点\(\xi\)的正则性，这正是Wiener判别法（定理150.7）所刻画的内容。

\subsection{Newton位势与Riesz位势}\label{newtonux4f4dux52bfux4e0erieszux4f4dux52bf}

\textbf{定义150.3}（Newton位势）：当\(n \geq 3\)时，对\(\mathbb{R}^n\)上有限Radon测度\(\mu\)，定义其\textbf{Newton位势}为 \[U^\mu(x) = \int_{\mathbb{R}^n} \frac{1}{|x-y|^{n-2}} d\mu(y)\] 当\(n = 2\)时， \[U^\mu(x) = -\frac{1}{2\pi} \int_{\mathbb{R}^2} \log|x-y| d\mu(y)\] \(U^\mu\)在\(\operatorname{supp}\mu\)之外调和，且\(\Delta U^\mu = -c_n \mu\)（分布意义下），其中\(c_n = (n-2)|\partial B_1| = (n-2)n\omega_n\)（\(n \geq 3\)）或\(c_2 = 2\pi\)。Newton位势是分布意义下Poisson方程\(\Delta u = -\mu\)的基本解卷积。

\textbf{定义150.4}（Riesz位势\(I_\alpha\)）：对\(0 < \alpha < n\)，\textbf{Riesz位势}\(I_\alpha\)形式上定义为\(I_\alpha = (-\Delta)^{-\alpha/2}\)，其Fourier乘子表示为 \[\widehat{I_\alpha f}(\xi) = |\xi|^{-\alpha} \hat{f}(\xi)\] 物理空间中\(I_\alpha f(x) = \gamma_\alpha \int_{\mathbb{R}^n} |x-y|^{\alpha-n} f(y) dy\)，其中\(\gamma_\alpha = \pi^{n/2} 2^\alpha \Gamma(\alpha/2) / \Gamma((n-\alpha)/2)\)。当\(\alpha = 2\)时\(I_2\)即为Newton位势算子\((-\Delta)^{-1}\)；Newton位势是Riesz位势的特例。当\(\alpha \to 0\)时，\(I_\alpha\)趋近于恒等算子；当\(\alpha \to n\)时，\(I_\alpha\)趋近于局部算子。

\textbf{定理150.4}（Riesz位势的半群性质）：对\(0 < \alpha, \beta\)且\(\alpha + \beta < n\)， \[I_\alpha \circ I_\beta = I_{\alpha+\beta}\] 这一半群性质是Riesz位势理论的基石，表明分数阶积分算子\(I_\alpha\)关于\(\alpha\)构成半群。

\textbf{证明}：在Fourier变换侧， \[\widehat{I_\alpha I_\beta f}(\xi) = |\xi|^{-\alpha} \widehat{I_\beta f}(\xi) = |\xi|^{-\alpha} \cdot |\xi|^{-\beta} \hat{f}(\xi) = |\xi|^{-(\alpha+\beta)} \hat{f}(\xi) = \widehat{I_{\alpha+\beta} f}(\xi)\] Fourier变换在\(\mathcal{S}'(\mathbb{R}^n)\)上为同构，故\(I_\alpha \circ I_\beta = I_{\alpha+\beta}\)。物理空间可通过Riesz核的卷积恒等式直接验证，其核心是Gamma函数恒等式：\(\gamma_\alpha \gamma_\beta \int_{\mathbb{R}^n} |x-z|^{\alpha-n} |z-y|^{\beta-n} dz = \gamma_{\alpha+\beta} |x-y|^{\alpha+\beta-n}\)。\(\blacksquare\)

\textbf{定理150.5}（Hardy-Littlewood-Sobolev不等式）：对\(0 < \alpha < n\)，\(1 < p < q < \infty\)满足\(1/q = 1/p - \alpha/n\)，存在常数\(C = C(p,q,n,\alpha)\)使 \[\|I_\alpha f\|_{L^q(\mathbb{R}^n)} \leq C \|f\|_{L^p(\mathbb{R}^n)}\] 该不等式是Sobolev嵌入定理的位势形式，等价于\(\|(-\Delta)^{-\alpha/2} f\|_{L^q} \leq C \|f\|_{L^p}\)。

\textbf{证明}：将Riesz核分解为\(|x|^{\alpha-n} = K_0 + K_\infty\)，其中\(K_0(x) = |x|^{\alpha-n} \chi_{\{|x| \leq R\}}\)（局部奇异部分），\(K_\infty(x) = |x|^{\alpha-n} \chi_{\{|x| > R\}}\)（长程衰减部分）。对\(K_0\)，由Young不等式得\(\|K_0 * f\|_{L^q} \leq \|K_0\|_{L^r} \|f\|_{L^p}\)，其中\(1/q = 1/r + 1/p - 1\)。计算\(\|K_0\|_{L^r} = (\int_{|x|\leq R} |x|^{r(\alpha-n)}dx)^{1/r} = C R^{\alpha - n/r}\)。对\(K_\infty\)，\(\|K_\infty * f\|_{L^\infty} \leq \|K_\infty\|_{L^\infty} \|f\|_{L^1} = R^{\alpha-n} \|f\|_{L^1}\)。将两项分解代入，利用Marcinkiewicz插值定理将两个弱型估计组合为强\((p,q)\)型估计。优化参数\(R\)（取\(R\)使两项平衡）即得最优常数。\(\blacksquare\)

\textbf{定理150.6}（Riesz位势的端点估计）：(i) 当\(p = 1\)时，\(I_\alpha\)将\(L^1\)映射到弱\(L^{n/(n-\alpha)}\)：\(\|I_\alpha f\|_{L^{n/(n-\alpha),\infty}} \leq C \|f\|_{L^1}\)；(ii) 当\(q = \infty\)时，\(I_\alpha\)仅当\(\alpha = 0\)时为\(L^p \to L^\infty\)有界。端点情形\(p = n/\alpha\)时\(I_\alpha: L^{n/\alpha} \to \operatorname{BMO}\)（Muckenhoupt-Wheeden, 1974）。

\subsection{Green函数与Dirichlet问题}\label{greenux51fdux6570ux4e0edirichletux95eeux9898}

\textbf{定义150.5}（Green函数）：区域\(\Omega \subset \mathbb{R}^n\)的\textbf{Green函数}\(G_\Omega(x,y) = \Phi(x-y) - H_y(x)\)，其中\(\Phi\)为Laplace方程的基本解（\(n \geq 3\)时\(\Phi(x) = c_n|x|^{2-n}\)，\(n=2\)时\(\Phi(x) = -\frac{1}{2\pi}\log|x|\)），\(H_y(x)\)为调和函数、满足边界条件\(H_y(x)|_{\partial\Omega} = \Phi(x-y)|_{\partial\Omega}\)（在正则边界点的极限意义下）。\(G_\Omega\)对称：\(G_\Omega(x,y) = G_\Omega(y,x)\)；正性：\(G_\Omega(x,y) > 0\)（对连通\(\Omega\)）。

Green 函数是求解 Dirichlet 问题的核心工具。其构造思想来源于物理中的''镜像法''：基本解 \(\Phi(x-y)\) 描述了 \(y\) 处的点电荷在 \(x\) 处产生的位势，而调和修正项 \(H_y(x)\) 可视为由边界 \(\partial\Omega\) 上感应电荷产生的位势，其作用恰好抵消基本解在边界上的值，从而满足 Dirichlet 边界条件。Green 函数的对称性 \(G_\Omega(x,y) = G_\Omega(y,x)\) 是自伴椭圆算子 Green 函数的普遍性质，源于基本解 \(\Phi\) 的对称性和 \(H_y(x)\) 的构造。在实际计算中，Green 函数仅在少数对称区域上可显式写出：对于上半空间 \(\mathbb{R}^n_+\)，\(G(x,y) = \Phi(x-y) - \Phi(x-\tilde{y})\)（镜像法），其中 \(\tilde{y}\) 是 \(y\) 关于 \(\partial\mathbb{R}^n_+\) 的反射；对于球 \(B(0,R)\)，\(G(x,y) = \Phi(x-y) - \Phi(\frac{|y|}{R}(x - \frac{R^2}{|y|^2}y))\)（Kelvin 变换）。对于一般区域，Green 函数的存在性由 Perron 方法（或变分方法）保证，但显式表达式通常不可得。

\textbf{定理150.7}（Green表示公式）：对\(u \in C^2(\overline{\Omega})\)， \[u(x) = \int_\Omega G_\Omega(x,y) \Delta u(y) dy + \int_{\partial\Omega} \frac{\partial G_\Omega}{\partial n_y}(x,y) u(y) dS(y)\] 其中\(\frac{\partial G_\Omega}{\partial n_y}\)为\(G_\Omega\)关于\(y\)的外法向导数（在边界上的极限）。当\(u\)为调和函数时，体积项消失，退化为Poisson积分公式。

\textbf{证明}：在\(\Omega_\varepsilon = \Omega \setminus B(x,\varepsilon)\)上应用Green第二恒等式于\(u(\cdot)\)与\(G_\Omega(x,\cdot)\)： \[\int_{\Omega_\varepsilon} (u\Delta_y G_\Omega - G_\Omega \Delta_y u) dy = \int_{\partial\Omega_\varepsilon} \left(u\frac{\partial G_\Omega}{\partial n_y} - G_\Omega\frac{\partial u}{\partial n_y}\right) dS\] 在\(\Omega_\varepsilon\)上\(\Delta_y G_\Omega(x,\cdot) = 0\)（因为\(G_\Omega\)关于\(y\)调和）。令\(\varepsilon \to 0\)，在\(\partial B(x,\varepsilon)\)上\(G_\Omega(x,y) \sim \Phi(x-y)\)，\(\frac{\partial G_\Omega}{\partial n_y} \sim \frac{\partial \Phi}{\partial n_y}\)。基本解奇性贡献：\(\lim_{\varepsilon \to 0} \int_{\partial B(x,\varepsilon)} u(y) \frac{\partial \Phi}{\partial n_y}(x-y) dS(y) = u(x)\)（由\(\Delta\Phi = \delta_0\)）。\(\partial\Omega\)上的边界项保持不变，整理即得Green表示公式。\(\blacksquare\)

Green 表示公式揭示了解析结构与边界数据的深刻联系：第一项为 Poisson 方程 \(\Delta u = f\) 的特解（体积位势），第二项为 Dirichlet 问题的解（边界位势）。特别地，当 \(u\) 为调和函数时（\(\Delta u = 0\)），Green 表示公式退化为 Poisson 积分公式，其中 Poisson 核 \(P_\Omega(x,y) = \frac{\partial G_\Omega}{\partial n_y}(x,y)\) 将边界数据延拓为区域内部的调和函数。Poisson 核的正性（对充分光滑的边界）是调和函数极大值原理和 Harnack 不等式的几何基础。

\textbf{定理150.8}（调和测度）：对Borel集\(E \subset \partial\Omega\)，定义\(\omega^x(E) = \int_E P_\Omega(x,y) dS(y)\)为\(E\)在\(x\)处的\textbf{调和测度}。\(\omega^x\)是\(\partial\Omega\)上的概率测度，且对固定\(E\)，\(x \mapsto \omega^x(E)\)为\(\Omega\)上的调和函数。调和测度在概率论中对应Brown运动从\(x\)出发首次击中\(\partial\Omega\)时落于\(E\)内的概率，建立了势论与随机过程的深刻联系。

\subsection{容量理论}\label{ux5bb9ux91cfux7406ux8bba-1}

\textbf{定义150.6}（Newton容量 --- 能量泛函定义）：紧集\(K \subset \mathbb{R}^n\)（\(n \geq 3\)）的\textbf{Newton容量}为 \[\operatorname{cap}(K) = \sup\left\{ \mu(K) : \mu\text{为支撑于}K\text{的正Radon测度}, \; I(\mu) \leq 1 \right\}\] 其中\(I(\mu) = \int_K \int_K |x-y|^{2-n} d\mu(x)d\mu(y)\)为\(\mu\)的能量。等价的能量极小化定义为\(\operatorname{cap}(K) = \sup\{\mu(K)^2 / I(\mu) : I(\mu) > 0\}\)。容量是测度论中''质量''概念的推广，量度了集合容纳电荷的能力。

\textbf{定义150.7}（Wiener容量 --- Dirichlet能量极小化）：对紧集\(K\)， \[\operatorname{cap}_W(K) = \inf\left\{ \int_{\mathbb{R}^n} |\nabla u|^2 dx : u \in C_c^\infty(\mathbb{R}^n), \; u \geq 1 \text{在}K\text{上} \right\}\] 对一般Borel集\(E\)，\(\operatorname{cap}_W(E) = \inf\{\operatorname{cap}_W(U) : U\text{开}, E \subset U\}\)。Wiener容量与Newton容量等价：\(\operatorname{cap}_W \sim \operatorname{cap}\)（相差仅依赖于\(n\)的常数因子）。

\textbf{定义150.8}（平衡测度与平衡位势）：对紧集\(K\)（\(\operatorname{cap}(K) > 0\)），存在唯一概率测度\(\mu_K\)（称为\textbf{平衡测度}）使\(I(\mu_K) = \min\{I(\mu) : \mu\text{为}K\text{上概率测度}\}\)。平衡位势\(U^{\mu_K}\)在\(K\)上几乎处处等于常数\(\operatorname{cap}(K)^{-1}\)（在\(\operatorname{cap}(K)\)的精细连续性意义下）。对球\(K = B(0,R)\)，平衡测度为边界上的均匀分布\(\frac{1}{|\partial B_R|} dS\)。

\textbf{定理150.9}（电容量的基本性质）：容量满足： (i) \textbf{单调性}：\(E_1 \subset E_2 \Rightarrow \operatorname{cap}(E_1) \leq \operatorname{cap}(E_2)\)； (ii) \textbf{可数次可加性}：\(\operatorname{cap}(\bigcup_{i=1}^\infty E_i) \leq \sum_{i=1}^\infty \operatorname{cap}(E_i)\)； (iii) \textbf{外正则性}：\(\operatorname{cap}(E) = \inf\{\operatorname{cap}(U) : U\text{开}, E \subset U\}\)； (iv) \textbf{单点零容}：\(\operatorname{cap}(\{x\}) = 0\)（\(n \geq 2\)）； (v) \textbf{球的容量}：\(\operatorname{cap}(B(0,R)) = c_n R^{n-2}\)（\(n \geq 3\)），\(\operatorname{cap}(B(0,R)) = c_2\)（\(n=2\)，对数容量）。 (vi) \textbf{伸缩性}：\(\operatorname{cap}(tE) = t^{n-2} \operatorname{cap}(E)\)（\(n \geq 3\)）。

\textbf{证明}：(i)测度类的包含关系直接导致。(ii)对每个\(E_i\)取近似紧集、利用能量泛函的次可加性。(iii)由Wiener容量的外正则定义导出。(iv)单点处能量发散（\(|x-y|^{2-n}\)在\(n \geq 2\)时在原点处不可积）。(v)通过对称重排和等周不等式，球的平衡测度为其边界上均匀分布，计算\(I(\mu) = \frac{1}{|\partial B_R|^2}\int_{\partial B_R}\int_{\partial B_R}|x-y|^{2-n}dS(x)dS(y) = R^{2-n}\)（尺度分析），故\(\operatorname{cap}(B(0,R)) = R^{n-2}\)。(vi)由伸缩变换\(x \mapsto t x\)和能量泛函的齐次性得出。\(\blacksquare\)

\subsection{极集与细拓扑}\label{ux6781ux96c6ux4e0eux7ec6ux62d3ux6251}

\textbf{定义150.9}（极集）：\(E \subset \mathbb{R}^n\)是\textbf{极集}若\(\operatorname{cap}(E) = 0\)。等价于存在上调和的局部可积函数\(v\)使\(E \subset \{v = +\infty\}\)。极集是调和函数可去奇异集的精确刻画：若\(u\)在\(\Omega \setminus E\)上调和且有界（或更一般地，在\(E\)附近有界），且\(E\)为极集，则\(u\)可延拓为\(\Omega\)上的调和函数。极集在\(\mathbb{R}^n\)中为Lebesgue零测集，但反之不真------Cantor型集可具有正容量但零Lebesgue测度。

\textbf{定义150.10}（细拓扑）：\(\mathbb{R}^n\)上的\textbf{细拓扑}（Fine Topology）是使所有上调和函数连续的最弱拓扑。确切地说，\(V \subset \mathbb{R}^n\)是细开集，若对每点\(x \in V\)，存在上调和函数\(u\)和\(\varepsilon > 0\)使\(\{u < u(x) + \varepsilon\} \subset V\)。细拓扑严格细于Euclid拓扑，\(x_n \to x\)（细拓扑）当且仅当去掉一个极集后\(x_n\)在Euclid拓扑下收敛于\(x\)。细拓扑是势论中''自然''的拓扑，上调和函数在细拓扑下连续，且容量具有细拓扑下的紧性。

\textbf{定理150.10}（Kellogg定理 --- Poisson积分的弱边界极限）：设\(\Omega \subset \mathbb{R}^n\)为有界区域，\(f \in C(\partial\Omega)\)，\(u = H_\Omega f\)为Poisson积分解（或Perron解）。则对\(\partial\Omega\)的每个Wiener正则点\(\xi\)， \[\lim_{\substack{x \to \xi \\ x \in \Omega}} u(x) = f(\xi)\] 边界点\(\xi\)为正则点当且仅当\textbf{Wiener判别法}成立： \[\sum_{k=1}^\infty 2^{k(n-2)} \operatorname{cap}\big(B(\xi, 2^{-k}) \setminus \Omega\big) = \infty\] （\(n=2\)时以对数容量\(\operatorname{cap}_{\log}\)替代Newton容量）。

\textbf{证明}：正则性等价于\(\xi\)处存在\textbf{闸函数}（barrier function）------在\(\xi\)趋于零且在\(\Omega^c\)上有正下界的上调和函数。Wiener判别法通过逐层容量估计给出闸函数存在的充要条件。具体地，定义环带\(A_k = B(\xi, 2^{-k}) \setminus B(\xi, 2^{-k-1})\)，构造上调和函数\(w_k\)在\(A_k \cap \Omega^c\)上\(\geq 1\)且能量受限于\(\operatorname{cap}(A_k \setminus \Omega)\)。若Wiener级数发散，则\(\sum w_k\)在\(\xi\)附近收敛于零，构造出闸函数；若级数收敛，则不存在闸函数，\(\xi\)为不规则点。\(\blacksquare\)

\textbf{定理150.11}（Choquet容量定理）：Newton容量在解析集（从而在所有Borel集）上满足\textbf{Choquet容量公理}： (i) \textbf{升序连续}：\(E_1 \subset E_2 \subset \cdots\)递增时，\(\operatorname{cap}(\bigcup_{i=1}^\infty E_i) = \lim_{i \to \infty} \operatorname{cap}(E_i)\)； (ii) \textbf{内正则}：对任意\(E\)，\(\operatorname{cap}(E) = \sup\{\operatorname{cap}(K) : K \subset E, K\text{紧}\}\)。

\textbf{证明}：(i)令\(E = \bigcup E_i\)。对任意容许测度\(\mu\)支撑于\(E\)，截断得\(\mu|_{E_i}\)支撑于\(E_i\)、能量不超过\(I(\mu)\)，取上确界即得\(\lim \operatorname{cap}(E_i) \geq \operatorname{cap}(E)\)。反向不等式由单调性给出。(ii)对Borel集\(E\)，由外正则性配合取\(K_n \subset E\)紧使\(\operatorname{cap}(E \setminus K_n) < 1/n\)，则\(\operatorname{cap}(K_n) \to \operatorname{cap}(E)\)。解析集情形通过Souslin运算保持此性质。Choquet定理的重要意义在于：它表明容量是具有良好连续性公理的''测度型''集函数，这为容量的概率论解释（与Brown运动的首中分布联系）提供了理论基础。\(\blacksquare\)

\subsection{势论中的能量方法}\label{ux52bfux8bbaux4e2dux7684ux80fdux91cfux65b9ux6cd5}

\textbf{定义150.11}（能量空间）：对\(n \geq 3\)，定义\textbf{能量空间}\(\mathcal{E}\)为所有有限能量Radon测度的集合： \[\mathcal{E} = \{\mu : I(\mu) = \int_{\mathbb{R}^n}\int_{\mathbb{R}^n} |x-y|^{2-n} d\mu(x)d\mu(y) < \infty\}\] \(\mathcal{E}\)赋予内积\(\langle \mu, \nu \rangle_E = \int\int |x-y|^{2-n} d\mu(x)d\nu(y)\)构成一个准Hilbert空间（在商去零能量测度后为Hilbert空间）。能量空间是势论的泛函分析基础，其中Gauss变分原理描述平衡测度。

\textbf{定理150.12}（Gauss变分原理）：设\(K\)为紧集，\(\operatorname{cap}(K) > 0\)。则存在唯一测度\(\mu_K \in \mathcal{E}\)（支撑于\(K\)）使\(I(\mu_K) = \min\{I(\mu) : \mu\text{支撑于}K, \mu(K)=1\}\)。\(\mu_K\)为平衡测度，且\(U^{\mu_K} = \operatorname{cap}(K)^{-1}\)在\(K\)上拟处处成立（quasi-everywhere，即除去一个极集后成立）。此外，\(\mu_K\)满足变分不等式：\(\int U^{\mu_K} d\nu \geq \operatorname{cap}(K)^{-1}\)对所有\(K\)上概率测度\(\nu\)成立。

\textbf{证明}：能量泛函\(I\)是\(\mathcal{E}\)上严格凸的弱下半连续泛函。在凸集\(\{\mu \text{支撑于}K, \mu(K)=1\}\)上，\(I\)达到唯一极小值\(\mu_K\)。由Lagrange乘子法，极小值\(\mu_K\)满足\(\int U^{\mu_K} d\nu \geq \lambda\)对任意\(K\)上概率测度\(\nu\)，且等号对\(\nu = \mu_K\)成立。由此得\(U^{\mu_K} \geq \lambda\)在\(K\)上拟处处成立，且由\(\int U^{\mu_K} d\mu_K = \lambda\)得\(U^{\mu_K} = \lambda\)在\(\operatorname{supp}\mu_K\)上拟处处成立。\(\lambda = I(\mu_K) = \operatorname{cap}(K)^{-1}\)。\(\blacksquare\)

\subsection{Martin边界与正调和函数}\label{martinux8fb9ux754cux4e0eux6b63ux8c03ux548cux51fdux6570}

\textbf{定义150.12}（Martin边界）：设\(\Omega \subset \mathbb{R}^n\)为区域。固定\(x_0 \in \Omega\)，定义\textbf{Martin核}\(K(x,y) = G_\Omega(x,y) / G_\Omega(x_0,y)\)（\(y \neq x_0\)）。\(\Omega\)的\textbf{Martin紧化}\(\widehat{\Omega}\)是\(\Omega\)在\(K(x,\cdot)\)诱导的拓扑下的紧化。\textbf{Martin边界}\(\partial_M \Omega = \widehat{\Omega} \setminus \Omega\)。Martin边界上的每个点对应一个极小正调和函数（通过极限\(\lim_{y \to \xi} K(x,y)\)），且每个正调和函数可唯一表示为Martin边界上这些极小调和函数的积分（Martin表示定理）。Martin边界理论将正调和函数的表示与区域的边界几何联系起来，是势论中最深刻的结果之一。

\textbf{定理150.13}（Martin表示定理，1941）：\(\Omega\)上的每个正调和函数\(u\)可唯一表示为 \[u(x) = \int_{\partial_M \Omega} K_\xi(x) d\mu(\xi)\] 其中\(K_\xi(x) = \lim_{y \to \xi} K(x,y)\)为\(\xi \in \partial_M \Omega\)对应的极小调和函数，\(\mu\)为\(\partial_M \Omega\)上唯一有限Borel测度。

\hrulefill

\emph{经典势论从调和函数的均值性质出发，经Newton位势和Riesz位势的半群结构与HLS不等式，进入容量理论。Perron方法、Wiener容量与能量泛函的等价性给出了Dirichlet问题边界正则的完整刻画（Wiener判别法）。细拓扑的引入使极集理论系统化，Choquet容量定理完成容量公理体系的建立。Gauss变分原理和能量空间提供了势论与泛函分析的联系。Martin边界理论则将正调和函数的表示与区域的理想边界几何化，为现代位势论的测度论基础画上句号。}

\emph{卷二十七：调和分析。}

\newpage

\chapter{06-分析/05-调和分析/06-积分变换.md}\label{ux5206ux679005-ux8c03ux548cux5206ux679006-ux79efux5206ux53d8ux6362.md}

\chapter{积分变换}\label{ux79efux5206ux53d8ux6362}

\hrulefill

\section{第206章：Fourier变换与Laplace变换}\label{ux7b2c206ux7ae0fourierux53d8ux6362ux4e0elaplaceux53d8ux6362}

积分变换（Integral Transforms）是将函数映射到另一类函数空间的线性算子，其核心思想是利用积分核将微分方程转化为代数方程，或将卷积转化为乘积。Fourier变换和Laplace变换是最基本也是最重要的两类积分变换，构成了调和分析、偏微分方程和信号处理的理论基石。

\subsection{153.1 Fourier变换在L1和L2上的理论}\label{fourierux53d8ux6362ux5728l1ux548cl2ux4e0aux7684ux7406ux8bba}

\textbf{定义 153.1.1（Fourier变换------L1理论）} 对 \(f \in L^1(\mathbb{R}^n)\)，其Fourier变换 \(\widehat{f} = \mathcal{F}f\) 定义为：

\[\widehat{f}(\xi) = \int_{\mathbb{R}^n} f(x) e^{-2\pi i \langle x, \xi \rangle} dx, \quad \xi \in \mathbb{R}^n.\]

\textbf{命题 153.1.2（Fourier变换的基本性质）} 对 \(f, g \in L^1(\mathbb{R}^n)\)：

\begin{enumerate}
\def\labelenumi{\arabic{enumi}.}
\tightlist
\item
  （线性）\(\widehat{\alpha f + \beta g} = \alpha \widehat{f} + \beta \widehat{g}\)；
\item
  （平移）\(\widehat{f(\cdot - h)}(\xi) = e^{-2\pi i \langle h, \xi \rangle} \widehat{f}(\xi)\)；
\item
  （调制）\(\widehat{e^{2\pi i \langle x, \eta \rangle} f(x)}(\xi) = \widehat{f}(\xi - \eta)\)；
\item
  （伸缩）若 \(f_\lambda(x) = f(x/\lambda)\)，则 \(\widehat{f_\lambda}(\xi) = |\lambda|^n \widehat{f}(\lambda \xi)\)；
\item
  \(\widehat{f}\) 是一致连续的，且 \(\|\widehat{f}\|_{L^\infty} \leq \|f\|_{L^1}\)。
\end{enumerate}

\textbf{定理 153.1.3（Riemann-Lebesgue引理）} 对 \(f \in L^1(\mathbb{R}^n)\)，

\[\lim_{|\xi| \to \infty} \widehat{f}(\xi) = 0.\]

即 \(\widehat{f} \in C_0(\mathbb{R}^n)\)（在无穷远处趋于零的连续函数空间）。

\textbf{证明}：首先对特征函数 \(\chi_{[a, b]}\)（在 \(\mathbb{R}\) 上）直接计算得其Fourier变换为 \(e^{-2\pi i a\xi} - e^{-2\pi i b\xi} / (2\pi i \xi)\)，在无穷远处趋于零。阶梯函数的有限线性组合也具有此性质。然后对一般的 \(f \in L^1\)，用阶梯函数序列逼近（\(L^1\) 范数意义下），利用 \(\|\widehat{f}\|_{L^\infty} \leq \|f\|_{L^1}\) 的一致连续性得结论。多维情形由Fubini定理化为一维。\(\blacksquare\)

\textbf{定义 153.1.4（Fourier变换------L2理论）} 对 \(f \in L^1(\mathbb{R}^n) \cap L^2(\mathbb{R}^n)\)，其Fourier变换如上述定义。由Plancherel定理可延拓为 \(L^2(\mathbb{R}^n)\) 上的酉算子。

\textbf{定理 153.1.5（Plancherel定理）} Fourier变换 \(\mathcal{F}\) 在 \(L^1 \cap L^2\) 上的限制可唯一延拓为 \(L^2(\mathbb{R}^n)\) 上的酉算子（不计常数因子 \(2\pi\) 的约定）。对任意 \(f \in L^2(\mathbb{R}^n)\)，

\[\|\widehat{f}\|_{L^2} = \|f\|_{L^2}.\]

\textbf{证明概要}：首先证明对 \(f, g \in \mathcal{S}(\mathbb{R}^n)\)（Schwartz速降函数空间），\(\langle \widehat{f}, \widehat{g} \rangle_{L^2} = \langle f, g \rangle_{L^2}\)。\(\mathcal{S}\) 在 \(L^2\) 中稠密，故 \(\mathcal{F}\) 按等距延拓为 \(L^2\) 上的等距同构。满射性由反演公式保证。\(\blacksquare\)

\textbf{定义 153.1.6（Schwartz空间与缓增分布）} Schwartz空间 \(\mathcal{S}(\mathbb{R}^n)\) 由所有满足 \(\sup_{x \in \mathbb{R}^n} |x^\alpha \partial^\beta f(x)| < \infty\) 对任意多重指标 \(\alpha, \beta\) 的光滑函数组成。\(\mathcal{F}\) 将 \(\mathcal{S}\) 同构地映射到自身。缓增分布 \(\mathcal{S}'(\mathbb{R}^n)\) 上的Fourier变换由对偶定义：\(\langle \widehat{u}, \varphi \rangle = \langle u, \widehat{\varphi} \rangle\)。

\subsection{153.2 Fourier反演公式与Parseval-Plancherel定理}\label{fourierux53cdux6f14ux516cux5f0fux4e0eparseval-plancherelux5b9aux7406}

\textbf{定理 153.2.1（Fourier反演公式）} 若 \(f \in L^1(\mathbb{R}^n)\) 且 \(\widehat{f} \in L^1(\mathbb{R}^n)\)，则对几乎处处的 \(x\)，

\[f(x) = \int_{\mathbb{R}^n} \widehat{f}(\xi) e^{2\pi i \langle x, \xi \rangle} d\xi.\]

更一般地，对 \(f \in L^1(\mathbb{R}^n)\)，引入Gauss-Weierstrass核（或Fejér核）的可和性方法可得几乎处处收敛的Fourier反演。

\textbf{完整证明}：引入Gauss函数 \(g_\varepsilon(x) = e^{-\pi \varepsilon |x|^2}\)，其Fourier变换为 \(\widehat{g_\varepsilon}(\xi) = \varepsilon^{-n/2} e^{-\pi |\xi|^2 / \varepsilon}\)。注意到 \(\int_{\mathbb{R}^n} \widehat{g_\varepsilon}(\xi) d\xi = 1\)（Gauss积分）。则

\[\int_{\mathbb{R}^n} \widehat{f}(\xi) \widehat{g_\varepsilon}(\xi) e^{2\pi i \langle x, \xi \rangle} d\xi = \int_{\mathbb{R}^n} f(y) \left(\int_{\mathbb{R}^n} \widehat{g_\varepsilon}(\xi) e^{2\pi i \langle x - y, \xi \rangle} d\xi\right) dy = (f * g_\varepsilon)(x).\]

左边在 \(\varepsilon \to 0^+\) 时趋于 \(\int \widehat{f}(\xi) e^{2\pi i \langle x, \xi \rangle} d\xi\)（因 \(\widehat{g_\varepsilon} \to 1\) 逐点且由控制收敛定理）。右边 \(f * g_\varepsilon\) 是 \(f\) 的恒等逼近，在 \(f\) 的Lebesgue点处趋于 \(f(x)\)，故几乎处处收敛。\(\blacksquare\)

\textbf{定理 153.2.2（Parseval-Plancherel定理的完整形式）} 对 \(f, g \in L^2(\mathbb{R}^n)\)，

\[\int_{\mathbb{R}^n} f(x) \overline{g(x)} dx = \int_{\mathbb{R}^n} \widehat{f}(\xi) \overline{\widehat{g}(\xi)} d\xi.\]

特别地，\(\|f\|_{L^2} = \|\widehat{f}\|_{L^2}\)。Fourier变换 \(\mathcal{F}: L^2(\mathbb{R}^n) \to L^2(\mathbb{R}^n)\) 是酉算子（在适当归一化下）。

\textbf{证明}：首先在Schwartz空间 \(\mathcal{S}(\mathbb{R}^n)\) 上建立恒等式。对 \(f, g \in \mathcal{S}(\mathbb{R}^n)\)，由Fubini定理， \[\int_{\mathbb{R}^n} \widehat{f}(\xi) \overline{\widehat{g}(\xi)} d\xi = \int_{\mathbb{R}^n} \widehat{f}(\xi) \left(\int_{\mathbb{R}^n} \overline{g(x)} e^{2\pi i \langle x, \xi \rangle} dx\right) d\xi\] \[= \int_{\mathbb{R}^n} \overline{g(x)} \left(\int_{\mathbb{R}^n} \widehat{f}(\xi) e^{2\pi i \langle x, \xi \rangle} d\xi\right) dx = \int_{\mathbb{R}^n} \overline{g(x)} f(x) dx,\] 最后一步使用了Fourier反演公式（定理153.2.1），由于 \(f \in \mathcal{S}\) 有 \(\widehat{f} \in \mathcal{S} \subseteq L^1\)，反演成立。取 \(g = f\) 即得 \(\|f\|_{L^2} = \|\widehat{f}\|_{L^2}\)。

由于 \(\mathcal{S}\) 在 \(L^2\) 中稠密，\(\mathcal{F}\) 可唯一地等距延拓为 \(L^2(\mathbb{R}^n)\) 到自身的等距同构。满射性由 \(\mathcal{F}^4 = \text{id}\)（在 \(\mathcal{S}\) 上）和稠密性保证。在适当的归一化下，\(\mathcal{F}\) 是酉算子。\(\blacksquare\)

\textbf{定理 153.2.3（卷积定理）} 对 \(f, g \in L^1(\mathbb{R}^n)\)，

\[\widehat{f * g}(\xi) = \widehat{f}(\xi) \cdot \widehat{g}(\xi).\]

对Fourier逆变换亦有类似性质。卷积定理是Fourier变换将分析问题化为代数问题的核心机制：微分方程中的卷积运算在Fourier域中化为逐点乘法。

\textbf{证明}：对 \(f, g \in L^1(\mathbb{R}^n)\)，由Fubini定理和变量替换， \[\widehat{f * g}(\xi) = \int_{\mathbb{R}^n} \left(\int_{\mathbb{R}^n} f(x - y) g(y) dy\right) e^{-2\pi i \langle x, \xi \rangle} dx\] \[= \int_{\mathbb{R}^n} g(y) e^{-2\pi i \langle y, \xi \rangle} \left(\int_{\mathbb{R}^n} f(x - y) e^{-2\pi i \langle x - y, \xi \rangle} dx\right) dy\] \[= \int_{\mathbb{R}^n} g(y) e^{-2\pi i \langle y, \xi \rangle} \widehat{f}(\xi) dy = \widehat{f}(\xi) \cdot \widehat{g}(\xi).\] 其中第二个等号用了Fubini定理（\(f, g \in L^1\) 保证 \((x, y) \mapsto f(x-y)g(y)\) 在 \(\mathbb{R}^n \times \mathbb{R}^n\) 上可积），第三个等号对内层积分作变量替换 \(z = x - y\)。\(\blacksquare\)

\subsection{153.3 Laplace变换的定义与收敛域}\label{laplaceux53d8ux6362ux7684ux5b9aux4e49ux4e0eux6536ux655bux57df}

\textbf{定义 153.3.1（Laplace变换）} 函数 \(f: [0, \infty) \to \mathbb{C}\) 的Laplace变换定义为：

\[\mathcal{L}\{f\}(s) = F(s) = \int_0^\infty f(t) e^{-st} dt,\]

其中 \(s = \sigma + i\omega \in \mathbb{C}\)。绝对收敛的横坐标 \(\sigma_a = \inf \{\sigma \in \mathbb{R} : \int_0^\infty |f(t)| e^{-\sigma t} dt < \infty\}\) 定义了半平面 \(\operatorname{Re}(s) > \sigma_a\) 中的收敛域。

\textbf{命题 153.3.2（解析性）} \(F(s) = \mathcal{L}\{f\}(s)\) 在收敛半平面 \(\operatorname{Re}(s) > \sigma_a\) 内是解析函数（全纯函数），且其导数可通过积分号下求导得到：

\[F^{(k)}(s) = (-1)^k \int_0^\infty t^k f(t) e^{-st} dt.\]

\textbf{定义 153.3.3（Laplace反演公式------Bromwich积分）} 若 \(f\) 在 \(t \geq 0\) 上分段光滑且其Laplace变换 \(F(s)\) 在 \(\operatorname{Re}(s) > \sigma_a\) 收敛，则对 \(\gamma > \sigma_a\)，

\[f(t) = \frac{1}{2\pi i} \lim_{R \to \infty} \int_{\gamma - iR}^{\gamma + iR} F(s) e^{st} ds, \quad t > 0.\]

\textbf{完整证明}：将Bromwich积分写为

\[\frac{1}{2\pi i} \int_{\gamma - iR}^{\gamma + iR} F(s) e^{st} ds = \frac{1}{2\pi i} \int_{\gamma - iR}^{\gamma + iR} \left(\int_0^\infty f(u) e^{-su} du\right) e^{st} ds.\]

交换积分次序（由Fubini定理及指数衰减保证）：\(= \int_0^\infty f(u) \left(\frac{1}{2\pi i} \int_{\gamma - iR}^{\gamma + iR} e^{s(t-u)} ds\right) du = \int_0^\infty f(u) \frac{\sin(R(t-u))}{\pi (t-u)} e^{\gamma(t-u)} du\)。

此即 \(f(t) e^{\gamma t}\) 的Dirichlet核卷积，在 \(f\) 的连续点处，当 \(R \to \infty\) 时趋于 \(f(t)\)（因 \(\frac{\sin(Rx)}{\pi x} \to \delta(x)\)）。\(\blacksquare\)

\textbf{定理 153.3.4（Laplace变换的基本性质）}

\begin{enumerate}
\def\labelenumi{\arabic{enumi}.}
\tightlist
\item
  （导数）\(\mathcal{L}\{f'\}(s) = s F(s) - f(0^+)\)；一般地 \(\mathcal{L}\{f^{(n)}\}(s) = s^n F(s) - \sum_{k=0}^{n-1} s^{n-1-k} f^{(k)}(0^+)\)；
\item
  （积分）\(\mathcal{L}\{\int_0^t f(\tau) d\tau\}(s) = \frac{F(s)}{s}\)；
\item
  （卷积）\(\mathcal{L}\{f * g\}(s) = F(s) G(s)\)，其中 \((f * g)(t) = \int_0^t f(\tau) g(t - \tau) d\tau\) 为时域卷积。
\end{enumerate}

\textbf{证明}：（1）由分部积分，\(\mathcal{L}\{f'\}(s) = \int_0^\infty f'(t) e^{-st} dt = [f(t) e^{-st}]_0^\infty + s \int_0^\infty f(t) e^{-st} dt = sF(s) - f(0^+)\)，前提是 \(\lim_{t \to \infty} f(t)e^{-st} = 0\)（在收敛域内成立）。对高阶导数反复应用即得一般公式。

（2）记 \(g(t) = \int_0^t f(\tau) d\tau\)，则 \(g'(t) = f(t)\) 且 \(g(0) = 0\)。由导数性质，\(F(s) = \mathcal{L}\{g'\}(s) = s \mathcal{L}\{g\}(s) - g(0) = s \mathcal{L}\{g\}(s)\)，故 \(\mathcal{L}\{g\}(s) = F(s)/s\)。

（3）由定义和积分次序交换， \[\mathcal{L}\{f * g\}(s) = \int_0^\infty \left(\int_0^t f(\tau) g(t - \tau) d\tau\right) e^{-st} dt\] \[= \int_0^\infty \int_\tau^\infty f(\tau) g(t - \tau) e^{-st} dt d\tau = \int_0^\infty f(\tau) e^{-s\tau} \left(\int_0^\infty g(u) e^{-su} du\right) d\tau = F(s) G(s),\] 其中第二个等号交换了积分次序（Fubini定理），第三个等号做了变量替换 \(u = t - \tau\)。\(\blacksquare\)

\subsection{153.4 卷积定理与Abel-Askey定理}\label{ux5377ux79efux5b9aux7406ux4e0eabel-askeyux5b9aux7406}

\textbf{定理 153.4.1（卷积定理在微分方程求解中的应用）} 对常系数线性微分方程

\[a_n y^{(n)} + \cdots + a_1 y' + a_0 y = f(t),\]

取Laplace变换得 \(P(s) Y(s) - Q(s) = F(s)\)，其中 \(P(s) = \sum_{k=0}^n a_k s^k\) 为特征多项式，\(Q(s)\) 为初值项。故 \(Y(s) = (F(s) + Q(s)) / P(s)\)，取反演即得解。此方法将微分方程化为代数方程，对于非齐次项 \(f\) 的Laplace变换为有理函数或易计算函数时尤为高效。

\textbf{证明}：对微分方程两边取Laplace变换。由线性性质和导数公式（定理153.3.4）， \[\mathcal{L}\left\{\sum_{k=0}^n a_k y^{(k)}\right\}(s) = \sum_{k=0}^n a_k \mathcal{L}\{y^{(k)}\}(s) = \sum_{k=0}^n a_k \left(s^k Y(s) - \sum_{j=0}^{k-1} s^{k-1-j} y^{(j)}(0^+)\right).\] 将所有包含初值的项合并为 \(Q(s)\)，即得 \(P(s) Y(s) - Q(s) = F(s)\)。由此解得 \(Y(s) = (F(s) + Q(s)) / P(s)\)。对其取Laplace逆变换（Bromwich积分或部分分式展开）即得原方程的解 \(y(t)\)。特别地，若 \(F(s), Q(s), P(s)\) 均为有理函数，可利用部分分式分解和Laplace变换表直接写出显式解。\(\blacksquare\)

\textbf{定理 153.4.2（Laplace变换的大t行为------Abel-Askey定理）} 若 \(\lim_{t \to \infty} f(t) = A\)（有限极限），则

\[\lim_{s \to 0^+} s F(s) = A.\]

更一般地，若 \(f\) 在无穷远处的Cesàro平均为 \(A\)，则 \(\lim_{s \to 0^+} s F(s) = A\)。

\textbf{证明}：\(\int_0^\infty f(t) s e^{-st} dt\) 中，作变量替换 \(u = st\)，得 \(\int_0^\infty f(u/s) e^{-u} du\)。当 \(s \to 0^+\)，\(f(u/s) \to A\) 对每个 \(u > 0\)，由控制收敛定理（\(f\) 有界时）得极限为 \(A \int_0^\infty e^{-u} du = A\)。\(\blacksquare\)

\textbf{定理 153.4.3（Tauber定理）} 若 \(f(t) \geq 0\)（或更一般地 \(f(t) \geq -C\)）且 \(\lim_{s \to 0^+} s F(s) = A\)，则 \(\lim_{t \to \infty} \frac{1}{t} \int_0^t f(\tau) d\tau = A\)（Cesàro意义下的Tauber收敛）。Wiener的一般Tauber定理将此推广到了更广泛的卷积代数框架中。

\textbf{证明}：设 \(f(t) \geq 0\) 且 \(\lim_{s \to 0^+} sF(s) = A\)。引入变量替换，对任意固定的 \(\lambda > 0\)， \[s F(s) = \int_0^\infty f(t) s e^{-st} dt = \int_0^\infty f\left(\frac{u}{s}\right) e^{-u} du.\] 记 \(g(u) = \int_0^u f(\tau) d\tau\)。分部积分可得 \[sF(s) = s^2 \int_0^\infty g(t) e^{-st} dt = \int_0^\infty \frac{g(u/s)}{u^2/s^2} \cdot \frac{u e^{-u}}{s} du.\] 更直接地，利用 \(sF(s) = \int_0^\infty f(t/s) e^{-t} dt\)（变量替换 \(u = st\)）。由Karamata的Tauber定理（或对单调函数 \(g\) 使用正则变化理论）：条件 \(\lim_{s \to 0^+} sF(s) = A\) 蕴涵 \(\int_0^\infty f(u/s) e^{-u} du \to A\)，由 \(f \geq 0\) 和Wiener-Ikehara型论证（或直接使用Hardy-Littlewood-Karamata Tauber定理），得 \(g(t)/t \to A\) 当 \(t \to \infty\)。对下界条件 \(f(t) \geq -C\) 的情形，用 \(f(t) + C\) 替代并由线性性归结为非负情形。\(\blacksquare\)

\subsection{153.5 Fourier变换与Laplace变换的关系}\label{fourierux53d8ux6362ux4e0elaplaceux53d8ux6362ux7684ux5173ux7cfb}

\textbf{定理 153.5.1（Fourier-Laplace对偶）} 对 \(f\) 支撑在 \([0, \infty)\)，

\[\mathcal{L}\{f\}(\sigma + i\omega) = \mathcal{F}\{t \mapsto f(t) e^{-\sigma t} \mathbf{1}_{[0, \infty)}(t)\}\left(\frac{\omega}{2\pi}\right).\]

即单边Laplace变换本质上是加窗Fourier变换的解析延拓。

\textbf{证明}：由定义， \[\mathcal{L}\{f\}(\sigma + i\omega) = \int_0^\infty f(t) e^{-(\sigma + i\omega)t} dt = \int_{-\infty}^\infty f(t) e^{-\sigma t} \mathbf{1}_{[0,\infty)}(t) \cdot e^{-i\omega t} dt.\] 将Fourier变换定义为 \(\mathcal{F}\{g\}(\xi) = \int_{-\infty}^\infty g(t) e^{-2\pi i t \xi} dt\)（注意此处约定），则有 \[\mathcal{L}\{f\}(\sigma + i\omega) = \mathcal{F}\{t \mapsto f(t) e^{-\sigma t} \mathbf{1}_{[0,\infty)}(t)\}\left(\frac{\omega}{2\pi}\right).\] 此式说明Laplace变换在直线 \(\operatorname{Re}(s) = \sigma\) 上的值恰为函数 \(f(t)e^{-\sigma t}\)（延拓到整个实轴为零）的Fourier变换（适当缩放频率变量）。\(\blacksquare\)

\textbf{定理 153.5.2（Paley-Wiener定理）} \(L^2(0, \infty)\) 中函数的Laplace变换在右半平面 \(\operatorname{Re}(s) > 0\) 上属于Hardy空间 \(H^2\)，且边界值（\(\operatorname{Re}(s) = 0\)）恰好是Fourier-Plancherel变换。更一般地，Paley-Wiener定理刻画了紧支撑分布和速降函数的Fourier-Laplace变换的解析性质。

\textbf{证明}：设 \(f \in L^2(0, \infty)\)。定义 \(F(s) = \int_0^\infty f(t) e^{-st} dt\)（\(\operatorname{Re}(s) > 0\)）。对固定的 \(\sigma > 0\)，令 \(g_\sigma(t) = f(t) e^{-\sigma t}\)。由Cauchy-Schwarz不等式，\(\|g_\sigma\|_{L^1} \leq \|f\|_{L^2} \cdot \|e^{-\sigma t}\|_{L^2(0,\infty)} = \|f\|_{L^2} / \sqrt{2\sigma} < \infty\)，故 \(F(\sigma + i\omega) = \mathcal{F}\{g_\sigma\}(\omega/2\pi)\)（适当归一化下）。

由Plancherel定理，\(\int_{-\infty}^\infty |F(\sigma + i\omega)|^2 d\omega = 2\pi \int_0^\infty |f(t)|^2 e^{-2\sigma t} dt \leq 2\pi \|f\|_{L^2}^2\)。这表明 \(F\) 属于Hardy空间 \(H^2\)（右半平面上的解析函数且满足 \(\sup_{\sigma > 0} \int |F(\sigma + i\omega)|^2 d\omega < \infty\)）。当 \(\sigma \to 0^+\) 时，\(F(\sigma + i\omega)\) 在 \(L^2\) 意义下收敛到 \(\mathcal{F}\{f\}(\omega/2\pi)\)。紧支集版本的Paley-Wiener定理刻画了：\(f \in L^2\) 支撑在 \([-R, R]\) 当且仅当 \(\widehat{f}\) 可延拓为指数型 \(\leq 2\pi R\) 的整函数。\(\blacksquare\)

\hrulefill

\section{第207章：其他重要积分变换}\label{ux7b2c207ux7ae0ux5176ux4ed6ux91cdux8981ux79efux5206ux53d8ux6362}

\subsection{154.1 Mellin变换}\label{mellinux53d8ux6362-1}

\textbf{定义 154.1.1（Mellin变换）} 对 \(f: (0, \infty) \to \mathbb{C}\)，其Mellin变换定义为：

\[\mathcal{M}\{f\}(s) = \int_0^\infty f(x) x^{s-1} dx,\]

其中 \(s = \sigma + it \in \mathbb{C}\)。Mellin变换可视为Fourier变换在乘法群 \(\mathbb{R}^+\) 上的版本：作指数变量替换 \(x = e^t\)，则 \(\mathcal{M}\{f\}(s) = \mathcal{F}\{t \mapsto f(e^t) e^{\sigma t}\}(t/2\pi)\)。

\textbf{定理 154.1.2（Mellin反演公式）} 在适当的解析性条件下，

\[f(x) = \frac{1}{2\pi i} \int_{c - i\infty}^{c + i\infty} \mathcal{M}\{f\}(s) x^{-s} ds,\]

其中积分路径 \(\operatorname{Re}(s) = c\) 位于Mellin变换的收敛带内。

\textbf{证明}：Mellin变换通过指数变量替换 \(x = e^t\) 与Fourier变换联系。令 \(t = \log x\)，\(g(t) = f(e^t) e^{c t}\)。则 \[\mathcal{M}\{f\}(c + i\omega) = \int_0^\infty f(x) x^{c + i\omega - 1} dx = \int_{-\infty}^\infty f(e^t) e^{(c + i\omega)t} dt = \int_{-\infty}^\infty g(t) e^{i\omega t} dt.\] 这是 \(g(t)\) 的Fourier变换（适当归一化）。由Fourier反演公式， \[g(t) = \frac{1}{2\pi} \int_{-\infty}^\infty \mathcal{M}\{f\}(c + i\omega) e^{-i\omega t} d\omega.\] 代回 \(f(x) = x^{-c} g(\log x)\)，令 \(s = c + i\omega\)，\(ds = i d\omega\)，则 \[f(x) = \frac{1}{2\pi i} \int_{c - i\infty}^{c + i\infty} \mathcal{M}\{f\}(s) x^{-s} ds.\] 这正是Mellin反演公式。\(\blacksquare\)

\textbf{定理 154.1.3（Mellin卷积）} Mellin卷积定义为 \((f \star g)(x) = \int_0^\infty f(y) g(x/y) \frac{dy}{y}\)，满足 \(\mathcal{M}\{f \star g\}(s) = \mathcal{M}\{f\}(s) \cdot \mathcal{M}\{g\}(s)\)。

\textbf{证明}：由定义和积分次序交换， \[\mathcal{M}\{f \star g\}(s) = \int_0^\infty \left(\int_0^\infty f(y) g(x/y) \frac{dy}{y}\right) x^{s-1} dx\] \[= \int_0^\infty f(y) \left(\int_0^\infty g(x/y) x^{s-1} dx\right) \frac{dy}{y}.\] 对内层积分作变量替换 \(u = x/y\)，\(dx = y du\)，得 \(x^{s-1} dx = y^{s-1} u^{s-1} \cdot y du = y^s u^{s-1} du\)。故 \[\int_0^\infty g(x/y) x^{s-1} dx = y^s \int_0^\infty g(u) u^{s-1} du = y^s \mathcal{M}\{g\}(s).\] 代回外层积分， \[\mathcal{M}\{f \star g\}(s) = \int_0^\infty f(y) \cdot y^s \mathcal{M}\{g\}(s) \cdot \frac{dy}{y} = \mathcal{M}\{g\}(s) \int_0^\infty f(y) y^{s-1} dy = \mathcal{M}\{f\}(s) \cdot \mathcal{M}\{g\}(s).\] 这完全类似于Fourier变换的卷积定理，只是乘法群上的卷积取不同的Haar测度 \(dy/y\)。\(\blacksquare\)

\subsection{154.2 Mellin变换在数论中的应用}\label{mellinux53d8ux6362ux5728ux6570ux8bbaux4e2dux7684ux5e94ux7528}

\textbf{定理 154.2.1（Mellin变换与Dirichlet级数）} Dirichlet级数 \(D(s) = \sum_{n=1}^\infty a_n n^{-s}\) 可视为测度 \(\sum a_n \delta_{\log n}\) 的Mellin-Stieltjes变换。特别地，Riemann zeta函数满足：

\[\pi^{-s/2} \Gamma\left(\frac{s}{2}\right) \zeta(s) = \int_0^\infty \theta(it) t^{s/2 - 1} \frac{dt}{2},\]

其中 \(\theta(z) = \sum_{n=-\infty}^\infty e^{\pi i n^2 z}\) 为Jacobi theta函数。此Mellin变换表示是 \(\zeta(s)\) 函数方程证明的核心工具。

\textbf{证明}：将 \(\zeta(s)\) 的定义 \(\zeta(s) = \sum_{n=1}^\infty n^{-s}\)（\(\operatorname{Re}(s) > 1\)）改写为Mellin变换的形式。利用Gamma函数的积分表示 \(\Gamma(s/2) = \int_0^\infty t^{s/2-1} e^{-t} dt\)，有 \[\pi^{-s/2} \Gamma(s/2) n^{-s} = \int_0^\infty t^{s/2-1} e^{-\pi n^2 t} dt.\] 对所有 \(n\) 求和并交换积分与求和次序（由绝对收敛性保证）， \[\pi^{-s/2} \Gamma(s/2) \zeta(s) = \int_0^\infty t^{s/2-1} \sum_{n=1}^\infty e^{-\pi n^2 t} dt.\] 定义 \(\psi(t) = \sum_{n=1}^\infty e^{-\pi n^2 t}\)，则 \(\theta(it) = 1 + 2\psi(t)\)（其中 \(\theta(z) = \sum_{n=-\infty}^\infty e^{\pi i n^2 z}\)）。由此即得所需的Mellin积分表示。\(\blacksquare\)

\textbf{定理 154.2.2（解析延拓与函数方程）} 利用上述Mellin积分表示和theta函数的模变换性质，可得 \(\zeta(s)\) 的解析延拓至全复平面（除 \(s = 1\) 处的单极点外）并满足函数方程 \(\xi(s) = \xi(1 - s)\)。一般自守 \(L\)-函数的解析性质也通过类似Mellin变换技术建立。

\textbf{证明}：将前述Mellin积分分解为 \(\int_0^1 + \int_1^\infty\)。在 \(\int_0^1\) 部分应用Jacobi theta函数的模变换性质 \(\theta(-1/z) = \sqrt{z/i} \theta(z)\)，可得 \(\psi(1/t) = t^{1/2} \psi(t) + \frac{1}{2}(t^{1/2} - 1)\)。代入积分并作变量替换 \(t \mapsto 1/t\) 得全复平面上的解析延拓表示。具体地，定义完备的 \(\xi\) 函数 \(\xi(s) = \frac{1}{2} s(s-1) \pi^{-s/2} \Gamma(s/2) \zeta(s)\)，则 \(\xi(s)\) 为整函数且满足 \(\xi(s) = \xi(1-s)\)。此函数方程的对称性来源于 \(s \leftrightarrow 1-s\) 时Mellin积分中 \(\int_0^1\) 与 \(\int_1^\infty\) 部分的互换及theta函数模变换的精确抵消。\(\blacksquare\)

\textbf{定理 154.2.3（Hecke的逆定理）} 一个Dirichlet级数若满足一定函数方程和解析性条件，则其系数恰好对应某个模形式的Fourier系数。此Hecke逆定理将Mellin变换的解析性质与模形式的代数-几何结构建立了等价关系，是Langlands纲领中自守表示理论的技术基础。

\textbf{证明}：设 \(D(s) = \sum_{n=1}^\infty a_n n^{-s}\) 为Dirichlet级数，满足：（1）在某个半平面内绝对收敛；（2）可解析延拓至全复平面且满足有限阶增长条件；（3）存在正整数 \(N\)、权 \(k\) 和符号 \(\varepsilon = \pm 1\) 使得函数方程 \[\left(\frac{2\pi}{\sqrt{N}}\right)^{-s} \Gamma(s) D(s) = \varepsilon \left(\frac{2\pi}{\sqrt{N}}\right)^{-(k-s)} \Gamma(k-s) D(k-s)\] 成立。由Mellin反演公式，系数 \(a_n\) 可表为 \(D(s)\) 沿适当垂直线的围道积分。将积分路径向左平移并利用函数方程的对称性，可得 \(a_n\) 满足模变换关系。Hecke证明：定义 \(f(z) = \sum_{n=1}^\infty a_n e^{2\pi i n z}\)（\(z\) 在上半平面），则由函数方程可导出 \(f(-1/(Nz)) = \varepsilon N^{k/2} z^k f(z)\)，即 \(f\) 是权 \(k\)、级 \(N\) 的模形式。Hecke逆定理的本质在于：Mellin变换将模形式的自守性（\(z \mapsto -1/(Nz)\) 下的变换）等价地转化为Dirichlet级数的函数方程，这一双向对应是Langlands纲领中 \(L\)-函数与自守表示对应的雏形。\(\blacksquare\)

\subsection{154.3 Hankel变换与径向函数的Fourier变换}\label{hankelux53d8ux6362ux4e0eux5f84ux5411ux51fdux6570ux7684fourierux53d8ux6362}

\textbf{定义 154.3.1（Hankel变换）} 对径向函数 \(f(x) = f_0(|x|)\)（在 \(\mathbb{R}^n\) 上），其Fourier变换也是径向的，且可由一维Hankel变换给出：

\[\widehat{f}(\xi) = (2\pi)^{n/2} |\xi|^{-(n-2)/2} \int_0^\infty f_0(r) J_{(n-2)/2}(2\pi |\xi| r) r^{n/2} dr,\]

其中 \(J_\nu\) 为第一类Bessel函数。\(\nu\)-阶Hankel变换（或Fourier-Bessel变换）定义为：

\[\mathcal{H}_\nu\{g\}(k) = \int_0^\infty g(r) J_\nu(kr) r dr.\]

\textbf{定理 154.3.2（Hankel变换的反演）} \(\mathcal{H}_\nu\) 是 \(L^2((0, \infty), r dr)\) 上的对合（即 \(\mathcal{H}_\nu^2 = \text{id}\)），反演公式为对称形式。Hankel变换在具有圆对称性的物理问题中（如圆形波导、轴对称热传导）广泛应用。

\textbf{证明}：Hankel变换 \(\mathcal{H}_\nu\) 的对合性源于Bessel函数的正交性：\(\int_0^\infty J_\nu(kr) J_\nu(k'r) r dr = \delta(k - k') / k\)。对 \(g \in L^2((0,\infty), r dr)\)，令 \(G(k) = \mathcal{H}_\nu\{g\}(k) = \int_0^\infty g(r) J_\nu(kr) r dr\)。再次应用Hankel变换并交换积分次序： \[\mathcal{H}_\nu\{G\}(r) = \int_0^\infty G(k) J_\nu(kr) k dk = \int_0^\infty \left(\int_0^\infty g(s) J_\nu(ks) s ds\right) J_\nu(kr) k dk\] \[= \int_0^\infty g(s) \left(\int_0^\infty J_\nu(ks) J_\nu(kr) k dk\right) s ds = \int_0^\infty g(s) \frac{\delta(s-r)}{s} s ds = g(r).\] 因此 \(g(r) = \int_0^\infty G(k) J_\nu(kr) k dk\) 即为反演公式。等距性由Parseval恒等式 \(\int_0^\infty |g(r)|^2 r dr = \int_0^\infty |G(k)|^2 k dk\) 保证。\(\blacksquare\)

\subsection{154.4 Radon变换与CT扫描的数学原理}\label{radonux53d8ux6362ux4e0ectux626bux63cfux7684ux6570ux5b66ux539fux7406}

\textbf{定义 154.4.1（Radon变换）} 对 \(\mathbb{R}^n\) 上的函数 \(f\)，其Radon变换为沿超平面的积分：

\[\mathcal{R}\{f\}(p, \theta) = \int_{\{x : \langle x, \theta \rangle = p\}} f(x) d\sigma(x),\]

其中 \(\theta \in \mathbb{S}^{n-1}\)，\(p \in \mathbb{R}\)。在二维情形（\(n = 2\)），这是沿直线的积分，\(\theta\) 可用角度 \(\phi\) 参数化：\(\theta = (\cos \phi, \sin \phi)\)。

\textbf{定理 154.4.2（投影切片定理------Fourier切片定理）} Radon变换与Fourier变换通过以下恒等式联系：

\[\widehat{\mathcal{R}\{f\}(\cdot, \theta)}(r) = \widehat{f}(r \theta),\]

其中左边为关于 \(p\) 的一维Fourier变换（频率 \(r\)），右边为 \(n\) 维Fourier变换在射线 \(r\theta\) 上的值。

\textbf{证明}：由定义，

\[\widehat{\mathcal{R}\{f\}(\cdot, \theta)}(r) = \int_{\mathbb{R}} \left(\int_{\langle x, \theta \rangle = p} f(x) d\sigma(x)\right) e^{-2\pi i r p} dp = \int_{\mathbb{R}^n} f(x) e^{-2\pi i r \langle x, \theta \rangle} dx = \widehat{f}(r \theta).\]

第二步利用了将 \(\mathbb{R}^n\) 分解为与 \(\theta\) 正交的超平面族的Fubini型积分。\(\blacksquare\)

\textbf{定理 154.4.3（Radon反演公式）} 对 \(n\) 为奇数，

\[f(x) = \frac{1}{2} (2\pi)^{1-n} (-1)^{(n-1)/2} \int_{\mathbb{S}^{n-1}} \frac{\partial^{n-1}}{\partial p^{n-1}} \mathcal{R}\{f\}(p, \theta) \bigg|_{p = \langle x, \theta \rangle} d\theta.\]

对 \(n\) 为偶数，反演公式涉及Hilbert变换。二维情形下：

\[f(x, y) = \frac{1}{2\pi} \int_0^\pi \int_{-\infty}^\infty \frac{\partial_p \mathcal{R}\{f\}(p, \phi)}{(x \cos \phi + y \sin \phi) - p} dp d\phi.\]

\textbf{证明}：由投影切片定理和Fourier反演，\(f(x) = \int_{\mathbb{R}^n} \widehat{f}(\xi) e^{2\pi i \langle x, \xi \rangle} d\xi = \int_0^\infty \int_{\mathbb{S}^{n-1}} \widehat{f}(r\theta) e^{2\pi i r \langle x, \theta \rangle} r^{n-1} d\theta dr\)。利用对称性延拓到负频率：\(f(x) = \frac{1}{2} \int_{-\infty}^\infty \int_{\mathbb{S}^{n-1}} \widehat{\mathcal{R}f(\cdot,\theta)}(r) |r|^{n-1} e^{2\pi i r \langle x,\theta \rangle} d\theta dr\)。当 \(n\) 为奇数时，\(|r|^{n-1} = r^{n-1}\)，因子 \(r^{n-1}\) 在Fourier逆变换中对应 \(\frac{1}{(2\pi i)^{n-1}} \partial_p^{n-1}\)，即得所述公式。当 \(n\) 为偶数时，\(|r|^{n-1} = r^{n-1} \operatorname{sgn}(r)\)，\(\operatorname{sgn}(r)\) 对应Hilbert变换，由此导出二维反演中的Hilbert核。\(\blacksquare\)

\textbf{定理 154.4.4（CT扫描的数学原理）} 计算机断层扫描（CT）的数学基础恰为二维Radon反演：X射线在不同角度下的衰减测量恰为物体密度函数 \(f(x, y)\) 在各条直线上的Radon变换。由投影切片定理，可通过对所有角度的测量进行Fourier变换、径向插值、二维反Fourier变换重建原始密度图（滤波反投影算法）。Cormack和Hounsfield因发展CT技术（1979年诺贝尔生理学或医学奖）深刻地证明了应用数学对医学的革命性影响。

\textbf{证明}：X射线穿过物体时，强度按Beer-Lambert定律衰减：\(I = I_0 \exp(-\int_L \mu(x,y) ds)\)，其中 \(\mu\) 为线性衰减系数。测量值 \(\ln(I_0/I) = \int_L \mu ds = \mathcal{R}\{\mu\}(p,\phi)\)。由投影切片定理，\(\widehat{\mathcal{R}\mu(\cdot,\phi)}(r) = \widehat{\mu}(r \cos\phi, r \sin\phi)\)。采集所有 \(\phi \in [0,\pi)\) 的投影并取一维Fourier变换，即得 \(\widehat{\mu}\) 在整个二维频域的径向采样。滤波反投影算法等价于：\(\mu(x,y) = \int_0^\pi (\mathcal{R}\mu(\cdot,\phi) * h)(x \cos\phi + y \sin\phi) d\phi\)，其中 \(h\) 的Fourier变换为 \(\widehat{h}(r) = |r|\)（斜坡滤波器，源自极坐标Fourier反演的Jacobian因子）。\(\blacksquare\)

\subsection{154.5 小波变换}\label{ux5c0fux6ce2ux53d8ux6362}

\textbf{定义 154.5.1（连续小波变换）} 设 \(\psi \in L^2(\mathbb{R})\) 满足容许性条件：

\[C_\psi = \int_{-\infty}^\infty \frac{|\widehat{\psi}(\xi)|^2}{|\xi|} d\xi < \infty.\]

对 \(f \in L^2(\mathbb{R})\)，其连续小波变换定义为：

\[\mathcal{W}_\psi f(a, b) = \frac{1}{\sqrt{|a|}} \int_{-\infty}^\infty f(t) \overline{\psi\left(\frac{t - b}{a}\right)} dt, \quad a \neq 0.\]

参数 \(a\) 为尺度（scale），\(b\) 为平移（translation）。

\textbf{定理 154.5.2（小波反演公式）} 对满足容许性条件的母小波 \(\psi\)，

\[f(t) = \frac{1}{C_\psi} \int_{-\infty}^\infty \int_{-\infty}^\infty \mathcal{W}_\psi f(a, b) \frac{1}{\sqrt{|a|}} \psi\left(\frac{t - b}{a}\right) \frac{da db}{a^2}.\]

\textbf{证明概要}：小波变换可表示为卷积和伸缩的组合。在Fourier域中，容许性条件保证了重构核的恒等性。本质上是Calderón再生公式：\(\int_{-\infty}^\infty \widehat{\psi}(a\xi) \overline{\widehat{\psi}(a\xi)} da/|a| = C_\psi\) 对 \(\xi \neq 0\)。\(\blacksquare\)

\textbf{定义 154.5.3（多分辨率分析与离散小波）} 多分辨率分析（MRA）是一个嵌套子空间序列 \(\cdots \subset V_{-1} \subset V_0 \subset V_1 \subset \cdots\) 满足 \(\overline{\bigcup_j V_j} = L^2(\mathbb{R})\)，\(\bigcap_j V_j = \{0\}\)，且 \(f(x) \in V_j \iff f(2x) \in V_{j+1}\)。存在尺度函数 \(\phi\) 使得 \(\{\phi(x - k)\}_{k \in \mathbb{Z}}\) 构成 \(V_0\) 的规范正交基。由Mallat算法，可快速计算正交小波系数。Daubechies构造了紧支撑光滑正交小波族，这构成了JPEG2000图像压缩和信号去噪的数学基础。

\textbf{定理 154.5.4（小波在奇异性检测中的应用）} 对Holder指数为 \(\alpha\) 的函数，其小波系数在细尺度下衰减为 \(|\mathcal{W}_\psi f(a, b)| = O(|a|^{\alpha + 1/2})\)。此性质使得小波变换成为奇异性检测和分形分析的有力工具，优于仅捕获全局正则性的Fourier方法。

\textbf{证明}：设 \(f\) 在点 \(t_0\) 处具有点态Holder指数 \(\alpha\)，即存在多项式 \(P\) 使得 \(|f(t) - P(t-t_0)| \leq C|t-t_0|^\alpha\)。选取具有 \(N > \alpha\) 阶消失矩的小波 \(\psi\)（\(\int t^k \psi(t) dt = 0\)，\(k=0,\ldots,N-1\)）。由于消失矩性质，多项式 \(P\) 在小波变换中贡献为零，故 \[|\mathcal{W}_\psi f(a, b)| \leq \frac{1}{\sqrt{|a|}} \int |f(t) - P(t-t_0)| \cdot |\psi((t-b)/a)| dt \leq \frac{C}{\sqrt{|a|}} \int |t-t_0|^\alpha |\psi((t-b)/a)| dt.\] 作变量替换 \(u = (t-b)/a\)，\(dt = |a| du\)，得 \(|t-t_0| = |a u + (b-t_0)| \leq |a|(|u| + |b-t_0|/|a|)\)。当 \(|b-t_0| = O(|a|)\) 时，积分的主要贡献为 \(O(|a|^\alpha \cdot |a|^{1/2}) = O(|a|^{\alpha+1/2})\)。这意味着小波系数的局部衰减速率直接给出了逐点Holder指数的估计，而Fourier变换仅能反映全局 \(L^2\) 正则性（\(\widehat{f}(\xi) = O(|\xi|^{-\alpha-1})\) 仅当 \(f\) 在所有点具有一致的Holder正则性）。\(\blacksquare\)

\hrulefill

\section{第208章：积分变换在偏微分方程中的应用}\label{ux7b2c208ux7ae0ux79efux5206ux53d8ux6362ux5728ux504fux5faeux5206ux65b9ux7a0bux4e2dux7684ux5e94ux7528}

积分变换是偏微分方程（PDE）研究的核心工具之一，将微分方程转化为代数方程或更简单的微分方程，从而求得精确解或分析解的性质。

\textbf{定理 155.1（热方程的基本解）} 热方程 \(\partial_t u = \Delta u\) 的基本解（热核）为

\[\Phi(x, t) = \frac{1}{(4\pi t)^{n/2}} e^{-|x|^2 / (4t)}, \quad t > 0.\]

对空间变量取Fourier变换，\(\widehat{\Phi}(\xi, t) = e^{-4\pi^2 |\xi|^2 t}\)，对时间变量取Laplace变换（关于 \(t\)），\(\widehat{\Phi}(\xi, s) = 1 / (s + 4\pi^2 |\xi|^2)\)。初始值问题 \(u(x, 0) = f(x)\) 的解为 \(u = \Phi * f\)。 \textbf{证明}：对热方程 \(\partial_t u = \Delta u\) 关于 \(x\) 取Fourier变换，得 \(\partial_t \widehat{u}(\xi, t) = -4\pi^2 |\xi|^2 \widehat{u}(\xi, t)\)。对固定 \(\xi\)，这是关于 \(t\) 的ODE，解为 \(\widehat{u}(\xi, t) = \widehat{u}(\xi, 0) e^{-4\pi^2 |\xi|^2 t}\)。取 \(u(x,0) = \delta(x)\)，\(\widehat{u}(\xi,0)=1\)，得 \(\widehat{\Phi}(\xi,t) = e^{-4\pi^2 |\xi|^2 t}\)。Fourier逆变换：\(\Phi(x,t) = \int_{\mathbb{R}^n} e^{-4\pi^2 |\xi|^2 t} e^{2\pi i \langle x,\xi\rangle} d\xi = \prod_{j=1}^n \int_{-\infty}^\infty e^{-4\pi^2 \xi_j^2 t + 2\pi i x_j \xi_j} d\xi_j = (4\pi t)^{-n/2} e^{-|x|^2/(4t)}\)。对初值问题 \(u(x,0)=f(x)\)，由叠加原理 \(u=\Phi * f\)。\(\blacksquare\)

\textbf{定理 155.2（波动方程与Fourier变换）} 波动方程 \(\partial_{tt} u = c^2 \Delta u\) 在Fourier域中化为 \(\partial_{tt} \widehat{u} = -4\pi^2 c^2 |\xi|^2 \widehat{u}\)，即谐振子方程，解为 \(\widehat{u}(\xi, t) = \widehat{f}(\xi) \cos(2\pi c |\xi| t) + \widehat{g}(\xi) \sin(2\pi c |\xi| t) / (2\pi c |\xi|)\)。在三维空间中，Kirchhoff公式和Huygens原理均可由Fourier反演导出。 \textbf{证明}：对波动方程关于 \(x\) 取Fourier变换，\(\partial_{tt}\widehat{u} = -4\pi^2 c^2 |\xi|^2 \widehat{u}\)。这是频率为 \(2\pi c|\xi|\) 的谐振子方程，通解为 \(\widehat{u}(\xi,t) = A(\xi)\cos(2\pi c|\xi|t) + B(\xi)\sin(2\pi c|\xi|t)\)。由初值 \(\widehat{u}(\xi,0)=\widehat{f}(\xi)\)，\(\partial_t\widehat{u}(\xi,0)=\widehat{g}(\xi)\) 得 \(A=\widehat{f}\)，\(B=\widehat{g}/(2\pi c|\xi|)\)。Fourier逆变换即得d'Alembert型解。三维空间中的Kirchhoff公式由球面平均和Fourier反演导出。\(\blacksquare\)

\textbf{定理 155.3（Laplace方程与Mellin变换）} 在扇形区域或角域中的Laplace方程（或更一般的椭圆方程）可利用Mellin变换分离变量求解。Mellin变换将径向坐标的Euler型微分算子 \(\partial_x (x \partial_x)\) 对角化，特别适用于处理在原点处有奇异性的问题（如裂缝尖端的应力场分析）。 \textbf{证明}：在极坐标下Laplace算子为 \(\Delta = \partial_r^2 + r^{-1}\partial_r + r^{-2}\partial_\theta^2\)。Mellin变换 \(\mathcal{M}\{u(r)\}(s) = \int_0^\infty u(r) r^{s-1} dr\) 将Euler型算子对角化：分部积分得 \(\mathcal{M}\{r\partial_r u\}(s) = -s \mathcal{M}\{u\}(s)\)，\(\mathcal{M}\{\partial_r(r\partial_r u)\}(s) = s^2 \mathcal{M}\{u\}(s)\)。将Mellin变换应用于径向方程，偏微分方程化为关于 \(s\) 的代数方程，从而分离变量求解角域中的Laplace方程。\(\blacksquare\)

\textbf{定理 155.4（分数阶微分方程与Laplace变换）} 分数阶导数 \(D_t^\alpha\)（Caputo或Riemann-Liouville定义）的Laplace变换为 \(\mathcal{L}\{D_t^\alpha f\}(s) = s^\alpha F(s) - \text{初值项}\)。这一性质使得分数阶扩散方程 \(\partial_t^\alpha u = \Delta u\)（反常扩散的数学模型）在Laplace域中的分析成为可能，其基本解由Mittag-Leffler函数表达。 \textbf{证明}：Caputo分数阶导数 \(D_t^\alpha f(t) = \frac{1}{\Gamma(n-\alpha)} \int_0^t \frac{f^{(n)}(\tau)}{(t-\tau)^{\alpha-n+1}} d\tau\)（\(n=\lceil\alpha\rceil\)）的Laplace变换由卷积定理给出：\(\mathcal{L}\{D_t^\alpha f\}(s) = \frac{1}{s^{n-\alpha}} \cdot (s^n F(s) - \sum_{k=0}^{n-1} s^{n-1-k} f^{(k)}(0^+)) = s^\alpha F(s) - \sum_{k=0}^{n-1} s^{\alpha-1-k} f^{(k)}(0^+)\)。对分数阶扩散方程 \(\partial_t^\alpha u = \Delta u\) 取此替换得 \(s^\alpha \widehat{U} - s^{\alpha-1}\widehat{u}_0 = -4\pi^2|\xi|^2 \widehat{U}\)，解出 \(\widehat{U}\) 后取Laplace逆变换，涉及Mittag-Leffler函数 \(E_\alpha\)。\(\blacksquare\)

\textbf{定理 155.5（Hankel变换在轴对称PDE中的应用）} 对于柱坐标或球坐标下具有轴对称性的PDE，Hankel变换将径向Bessel型算子 \(r^{-1} \partial_r (r \partial_r) - \nu^2 / r^2\) 对角化。在处理圆形膜振动、圆形波导中电磁波传播和轴对称弹性力学问题时，Hankel变换是必不可少的工具。 \textbf{证明}：Bessel型算子 \(\mathcal{B}_\nu = \partial_r^2 + r^{-1}\partial_r - \nu^2/r^2\) 满足 \(\mathcal{B}_\nu J_\nu(kr) = -k^2 J_\nu(kr)\)（Bessel微分方程）。因此Hankel变换将其对角化：\(\mathcal{H}_\nu\{\mathcal{B}_\nu u\}(k) = -k^2 \mathcal{H}_\nu\{u\}(k)\)。对轴对称Helmholtz方程 \((\Delta + k_0^2)u = 0\) 的径向部分应用Hankel变换，即得代数方程，简化了圆形波导和圆膜振动问题的求解。\(\blacksquare\)

\textbf{定理 155.6（Radon变换与偏微分方程的反问题）} 在许多成像反问题（CT、MRI、地震勘探）中，测量数据是未知物理量沿直线或曲面的积分（Radon变换或其广义形式）。求解PDE约束下的Radon反演问题是数学反问题理论的核心。Fourier积分算子（FIO）和微局部分析（wavefront set）理论为这些反问题的适定性提供了严格的数学框架。 \textbf{证明}：CT反问题即给定 \(g(p,\theta) = \mathcal{R}f(p,\theta)\) 求 \(f\)。由投影切片定理 \(\widehat{g}(r,\theta) = \widehat{f}(r\theta)\)，需从径向采样重构 \(\widehat{f}\)。Fourier积分算子理论将Radon反演算子视为 \(I^{-(n-1)/2}\) 类FIO，其典范关系由投影几何决定。微局部分析表明Radon变换的波前集沿Hamilton流传播，为有限角度和含噪重建提供了奇异性检测的严格框架。\(\blacksquare\)

\textbf{定理 155.7（热方程的非齐次问题与Duhamel原理）} 非齐次热方程 \(\partial_t u - \Delta u = f(x, t)\) 在齐次初值条件 \(u(x, 0) = 0\) 下的解可由Duhamel原理表示为：

\[u(x, t) = \int_0^t \int_{\mathbb{R}^n} \Phi(x - y, t - s) f(y, s) dy ds,\]

其中 \(\Phi\) 为热核。对空间变量取Fourier变换得到关于时间的常微分方程 \(\partial_t \widehat{u} + 4\pi^2 |\xi|^2 \widehat{u} = \widehat{f}(\xi, t)\)，由常数变易法求解即得上述表达式。此方法体现了积分变换将偏微分方程转化为可解的常微分方程的核心威力。 \textbf{证明}：对非齐次热方程 \(\partial_t u - \Delta u = f\) 取空间Fourier变换，\(\partial_t \widehat{u} + 4\pi^2 |\xi|^2 \widehat{u} = \widehat{f}\)。常数变易法解得 \(\widehat{u}(\xi,t) = \int_0^t e^{-4\pi^2 |\xi|^2 (t-s)} \widehat{f}(\xi,s) ds\)（\(\widehat{u}(\xi,0)=0\)）。Fourier逆变换并利用卷积定理：\(u(x,t) = \int_0^t (\Phi(\cdot,t-s) * f(\cdot,s))(x) ds\)，其中 \(\Phi\) 为热核。此即Duhamel原理。\(\blacksquare\)

\textbf{定理 155.8（Schrödinger方程与Fourier积分算子）} 自由Schrödinger方程 \(i \partial_t \psi = -\frac{1}{2} \Delta \psi\) 在Fourier域中化为 \(i \partial_t \widehat{\psi} = 2\pi^2 |\xi|^2 \widehat{\psi}\)，解为 \(\widehat{\psi}(\xi, t) = e^{-2\pi^2 i |\xi|^2 t} \widehat{\psi}_0(\xi)\)。Fourier反演给出：

\[\psi(x, t) = \frac{1}{(2\pi i t)^{n/2}} \int_{\mathbb{R}^n} e^{i|x-y|^2 / (2t)} \psi_0(y) dy.\]

更一般地，变系数的偏微分方程可通过Fourier积分算子（FIO）理论处理。FIO是形如 \(\int e^{i \varphi(x, y, \xi)} a(x, y, \xi) \widehat{f}(\xi) d\xi\) 的振荡积分算子，由Hörmander于1970年代系统建立。FIO在奇性传播定理（Hörmander定理）和波动方程解的渐近分析中起着不可替代的作用。 \textbf{证明}：自由Schrodinger方程 \(i\partial_t\psi = -\frac{1}{2}\Delta\psi\) 取Fourier变换，\(i\partial_t\widehat{\psi} = 2\pi^2|\xi|^2\widehat{\psi}\)。解得 \(\widehat{\psi}(\xi,t) = e^{-2\pi^2 i |\xi|^2 t} \widehat{\psi}_0(\xi)\)。Fourier逆变换：\(\psi(x,t) = \int e^{-2\pi^2 i |\xi|^2 t} (\int \psi_0(y) e^{-2\pi i \langle y,\xi\rangle} dy) e^{2\pi i \langle x,\xi\rangle} d\xi\)。交换积分次序并计算Gauss型积分（配方完成平方）得 \(\psi(x,t) = (2\pi i t)^{-n/2} \int e^{i|x-y|^2/(2t)} \psi_0(y) dy\)。\(\blacksquare\)

\textbf{定理 155.9（散射理论与逆散射变换）} 反散射变换（IST）是求解KdV方程等可积非线性偏微分方程的Fourier变换的''非线性类比''。对KdV方程 \(u_t + u_{xxx} + 6u u_x = 0\)，定义与初值 \(u(x, 0)\) 相关的Schrödinger算子的散射数据，该数据在时间演化下按简单的线性规律变化。通过Gel'fand-Levitan-Marchenko积分方程反演即得任意时刻的解。IST方法将非线性PDE的求解转化为线性积分方程------这是积分变换思想的最高境界，也是现代可积系统理论的起点。 \textbf{证明}：对KdV方程 \(u_t + u_{xxx} + 6uu_x = 0\)，关联Schrodinger算子 \(L = -\partial_x^2 - u(x,t)\) 在KdV流下保持谱不变（等谱形变）。散射数据为反射系数 \(R(k,t)\) 和束缚态 \(\{\kappa_j, c_j(t)\}\)。关键发现（GGKM, 1967）：散射数据按线性规律演化------\(R(k,t)=R(k,0)e^{8ik^3t}\)，\(c_j(t)=c_j(0)e^{8\kappa_j^3t}\)。反散射步骤由Gel'fand-Levitan-Marchenko线性积分方程从散射数据重构势函数 \(u(x,t)\)，从而将非线性PDE化为线性问题。\(\blacksquare\)

\textbf{定理 155.10（Fourier变换在数值分析中的应用------拟谱方法）} 拟谱方法（pseudospectral method）利用快速Fourier变换（FFT）在Fourier空间中离散化空间导数，将偏微分方程约化为关于时间的常微分方程组。对周期边界条件的方程，谱精度（指数收敛）远优于有限差分和有限元方法。在湍流直接数值模拟、天气预报和量子动力学等高性能计算领域，拟谱方法是最优选择。 \textbf{证明}：对周期函数 \(u(x)\)（周期 \(L\)），Fourier级数 \(u(x) = \sum_k \widehat{u}_k e^{2\pi i k x/L}\)。空间导数在Fourier域化为乘法：\(\partial_x^m u\) 的系数为 \((2\pi i k/L)^m \widehat{u}_k\)。拟谱方法在Fourier空间精确计算导数，通过FFT在物理与Fourier空间之间切换以计算非线性项。对Burgers方程 \(u_t + uu_x = \nu u_{xx}\)，在Fourier空间中离散化空间导数（指数精度），在物理空间中计算乘积 \(u \cdot \partial_x u\)，收敛速度为谱精度（指数级）。\(\blacksquare\)

\textbf{定理 155.11（Hilbert变换与色散关系）} Hilbert变换定义为 \(\mathcal{H}\{f\}(t) = \frac{1}{\pi} \text{p.v.} \int_{-\infty}^\infty \frac{f(\tau)}{t - \tau} d\tau\)。在Fourier域中 \(\widehat{\mathcal{H}f}(\omega) = -i \operatorname{sgn}(\omega) \widehat{f}(\omega)\)。因果信号的实部和虚部由Kramers-Kronig关系（色散关系）联系------此关系是Hilbert变换的直接推论。在光学、电子学和信号处理中，色散关系表明了系统的吸收谱和折射率之间的内在联系，是因果律的频域表示。 \textbf{证明}：Hilbert变换 \(\mathcal{H}f(t) = \frac{1}{\pi} \text{p.v.} \int \frac{f(\tau)}{t-\tau} d\tau\) 是卷积 \(f * (1/(\pi t))\)。\(1/(\pi t)\) 的Fourier变换为 \(-i \operatorname{sgn}(\omega)\)，故 \(\widehat{\mathcal{H}f}(\omega) = -i \operatorname{sgn}(\omega) \widehat{f}(\omega)\)。对因果信号 \(f(t)=0\)（\(t<0\)），\(\widehat{f}(\omega)\) 在上半平面解析。由留数定理得Kramers-Kronig关系：\(\operatorname{Im}\widehat{f} = -\mathcal{H}\{\operatorname{Re}\widehat{f}\}\)，\(\operatorname{Re}\widehat{f} = \mathcal{H}\{\operatorname{Im}\widehat{f}\}\)。此即吸收谱与折射率的色散关系。\(\blacksquare\)

\textbf{定理 155.12（Z变换与离散时间信号处理）} Z变换是Laplace变换的离散时间版本，将序列 \(\{x_n\}_{n=0}^\infty\) 映射为复变函数 \(X(z) = \sum_{n=0}^\infty x_n z^{-n}\)。Z变换的收敛域为某个环形区域，其极点位置决定了序列的增长和振荡性质。Z变换的逆变换由围道积分 \(x_n = \frac{1}{2\pi i} \oint_C X(z) z^{n-1} dz\) 给出，或通过部分分式展开和查表法计算。在数字信号处理中，Z变换将差分方程（线性时不变系统的递推描述）转化为复频域的代数方程，其系统函数 \(H(z)\) 的极零点分布直接决定了系统的稳定性和频率响应特性。 \textbf{证明}：Z变换 \(X(z) = \sum_{n=0}^\infty x_n z^{-n}\) 在 \(|z|>R\) 内解析（\(R\) 为收敛半径）。逆Z变换由Cauchy积分给出：\(x_n = \frac{1}{2\pi i} \oint_C X(z) z^{n-1} dz\)（\(C\) 包围原点）。因为 \(\oint_C z^{-k-1} dz = 2\pi i \delta_{k,0}\)，故被积式中幂级数展开后仅 \(n\) 项幸存。对线性时不变系统 \(\sum a_k y_{n-k} = \sum b_j x_{n-j}\)，取Z变换得 \(H(z) = Y(z)/X(z) = (\sum b_j z^{-j})/(\sum a_k z^{-k})\)。系统BIBO稳定当且仅当 \(H(z)\) 的所有极点满足 \(|z|<1\)。\(\blacksquare\)

\textbf{定理 155.13（Gabor变换与短时Fourier变换）} Gabor变换（或短时Fourier变换，STFT）在Fourier变换中引入滑动窗口以获取信号的局部频谱信息：

\[S_f(\tau, \omega) = \int_{-\infty}^\infty f(t) \overline{g(t - \tau)} e^{-i\omega t} dt,\]

其中 \(g\) 为窗口函数（通常取Gauss函数，此时得到最优的时频分辨率）。STFT的平方模量 \(|S_f(\tau, \omega)|^2\) 称为谱图（spectrogram）。Heisenberg测不准原理约束了STFT的时频分辨率乘积：\(\Delta_t \cdot \Delta_\omega \geq 1/2\)。Gabor展开将信号表示为时频平移的基函数 \(g_{m,n}(t) = g(t - nT) e^{i m\Omega t}\) 的线性组合，其中 \(T\Omega \leq 2\pi\) 的必要条件体现了信息论和时频分析的深层联系。 \textbf{证明}：STFT定义为 \(S_f(\tau,\omega) = \int f(t) \overline{g(t-\tau)} e^{-i\omega t} dt\)。固定 \(\tau\)，这是加窗函数的Fourier变换。重构：若 \(\|g\|=1\)，则 \(f(t) = \frac{1}{2\pi} \iint S_f(\tau,\omega) g(t-\tau) e^{i\omega t} d\omega d\tau\)。时频分辨率约束为Heisenberg不等式 \(\Delta_t \cdot \Delta_\omega \geq 1/2\)，等号仅当 \(g\) 为Gauss函数。Gabor展开 \(f = \sum_{m,n} c_{m,n} g(t-nT)e^{im\Omega t}\) 的稳定性要求 \(T\Omega \leq 2\pi\)（框架条件）。\(\blacksquare\)

\textbf{定理 155.14（Wigner-Ville分布与二次时频表示）} Wigner-Ville分布（WVD）是信号的二次时频表示：

\[W_f(t, \omega) = \int_{-\infty}^\infty f\left(t + \frac{\tau}{2}\right) \overline{f\left(t - \frac{\tau}{2}\right)} e^{-i\omega\tau} d\tau.\]

WVD具有理想的瞬时频率和群延迟性质，且满足边缘分布：\(\int W_f(t,\omega) d\omega = |f(t)|^2\)，\(\int W_f(t,\omega) dt = |\widehat{f}(\omega)|^2\)。然而WVD存在交叉项（cross-terms）的问题------两个信号之和的WVD会产生振荡干扰项。Cohen类通过平滑核方法 \(\Phi(\theta, \tau)\) 统一了所有二次时频表示，其中Choi-Williams分布和锥形核分布等通过特定平滑策略在分辨率与交叉项抑制之间取得平衡。 \textbf{证明}：WVD定义为 \(W_f(t,\omega) = \int f(t+\tau/2) \overline{f(t-\tau/2)} e^{-i\omega\tau} d\tau\)。边缘分布：\(\int W_f(t,\omega) d\omega = 2\pi |f(t)|^2\)（因 \(\int e^{-i\omega\tau} d\omega = 2\pi\delta(\tau)\)）；频域边缘类似。Cohen类通过平滑核 \(\Phi(\theta,\tau)\) 统一表示：\(C_f(t,\omega) = \iiint f(u+\tau/2)\overline{f(u-\tau/2)} \Phi(\theta,\tau) e^{-i(\theta t + \omega\tau - \theta u)} du d\tau d\theta\)。WVD取 \(\Phi\equiv 1\)（最优分辨率但含交叉项），Choi-Williams取 \(\Phi = e^{-\theta^2\tau^2/\sigma}\)（抑制交叉项）。\(\blacksquare\)

\textbf{定理 155.15（框架理论------小波基的推广）} Hilbert空间 \(\mathcal{H}\) 中的一族向量 \(\{\varphi_i\}_{i \in I}\) 称为一个框架（frame），若存在常数 \(0 < A \leq B < \infty\) 使得对所有 \(f \in \mathcal{H}\)：

\[A \|f\|^2 \leq \sum_{i \in I} |\langle f, \varphi_i \rangle|^2 \leq B \|f\|^2.\]

框架是规范正交基的推广------它允许冗余表示，从而在信号处理中提供对传输误差和有限精度计算的鲁棒性。紧框架（\(A = B\)）给出类似于Parseval恒等式的完美重构。Daubechies、Grossmann和Meyer在1980年代建立的框架理论为Gabor分析和小波分析提供了统一的数学语言，将时频分析的工程直觉与泛函分析的严格性融为一体。

\newpage

\chapter{07-拓扑与几何/06-整体分析与指标理论.md}\label{ux62d3ux6251ux4e0eux51e0ux4f5506-ux6574ux4f53ux5206ux6790ux4e0eux6307ux6807ux7406ux8bba.md}

\chapter{整体分析与指标理论}\label{ux6574ux4f53ux5206ux6790ux4e0eux6307ux6807ux7406ux8bba}

\begin{quote}
覆盖 MSC 58-XX 整体分析、流形上的分析。论述Atiyah-Singer指标定理、Morse理论与变分法、Chern-Weil理论与示性类。
\end{quote}

\hrulefill

\section{第277章：Atiyah-Singer指标定理}\label{ux7b2c277ux7ae0atiyah-singerux6307ux6807ux5b9aux7406}

Atiyah-Singer指标定理（1963）是20世纪数学最伟大的成就之一，它将椭圆微分算子的解析指标（定义为核与余核维数之差）与底流形的拓扑不变量（由示性类的积分给出）联系起来。这一定理的深远意义在于它实现了分析、拓扑和几何的统一。在物理上，指标定理为量子反常（如手征反常和引力反常）提供了数学基础。本章从椭圆算子的Fredholm性质出发，建立指标的拓扑不变性，然后分别介绍K理论证明框架和热核证明方法，最后说明Gauß-Bonnet-Chern定理如何作为特例被包含在指标定理之中。

\subsection{x.1 椭圆算子与Fredholm性质}\label{x.1-ux692dux5706ux7b97ux5b50ux4e0efredholmux6027ux8d28}

\textbf{定义 x.1}（椭圆微分算子）。设\(M\)为紧致无边\(n\)维光滑流形，\(E,F \to M\)为复向量丛。\(m\)阶微分算子\(P: C^{\infty}(M,E) \to C^{\infty}(M,F)\)称为椭圆的，若其主象征\(\sigma_P(x,\xi): E_x \to F_x\)对所有\((x,\xi) \in T^*M\setminus 0\)都是线性同构。

\textbf{定理 x.1}（椭圆算子的Fredholm性质）。紧致流形上的椭圆算子\(P: H^s(M,E) \to H^{s-m}(M,F)\)（作为Sobolev空间之间的映射）是Fredholm算子：\(\ker P\)维数有限，\(\operatorname{range} P\)是闭子空间，\(\operatorname{coker} P = H^{s-m}/\operatorname{range}P\)维数有限。

\textbf{证明框架}：核心为构造拟逆（parametrix）\(Q\)使得\(QP - I\)和\(PQ - I\)为光滑化算子（从而是紧算子）。用象征演算------在局部坐标下，将\(P\)的象征\(\sigma_P(x,\xi)\)展开为齐次分量\(p_m + p_{m-1} + \cdots\)。由椭圆性，\(p_m(x,\xi)\)对\(\xi \neq 0\)可逆。构造\(Q\)的象征\(q_{-m} + q_{-m-1} + \cdots\)满足\(\sum_{|\alpha|\geq 0} \frac{1}{\alpha!}\partial_\xi^\alpha \sigma_P \cdot \partial_x^\alpha \sigma_Q = I\)和\(\sum_{|\alpha|\geq 0} \frac{1}{\alpha!}\partial_\xi^\alpha \sigma_Q \cdot \partial_x^\alpha \sigma_P = I\)（象征的星积）。逐次求解\(q_{-m} = p_m^{-1}\)，\(q_{-m-1} = -p_m^{-1}(p_{m-1}q_{-m} + \sum_{|\alpha|=1}\partial_\xi^\alpha p_m \cdot \partial_x^\alpha q_{-m})\)等。用Borel引理将形式级数提升为拟微分算子的象征。拟逆的存在保证了\(P\)具有有限维核和余核（Fredholm性质），因为\(Q\)提供了有限秩扰动下的双面逆。 \(\blacksquare\)

\textbf{定义 x.2}（解析指标）。椭圆算子\(P\)的指标定义为

\[\operatorname{ind}(P) = \dim\ker P - \dim\operatorname{coker} P\]

\subsection{x.2 指标的定义与同伦不变性}\label{x.2-ux6307ux6807ux7684ux5b9aux4e49ux4e0eux540cux4f26ux4e0dux53d8ux6027}

\textbf{定理 x.2}（指标的变形不变性与同伦不变性，完整证明）。指标\(\operatorname{ind}(P)\)在椭圆算子的连续形变下保持不变。特别地，若\(P_t\)是一族连续依赖于\(t \in [0,1]\)的椭圆算子，则\(\operatorname{ind}(P_0) = \operatorname{ind}(P_1)\)。

\textbf{证明}：记\(K(H_1,H_2)\)为从Hilbert空间\(H_1\)到\(H_2\)的紧算子集合（在算子范数下为闭子空间）。Fredholm算子\(F(H_1,H_2)\)为Calkin代数\(\mathcal{B}(H_1,H_2)/K(H_1,H_2)\)中的可逆元，指标函数\(\operatorname{ind}: F(H_1,H_2) \to \mathbb{Z}\)局部常值（因小扰动下维数\(\dim\ker\)和\(\dim\operatorname{coker}\)不变）。而映射\(t \mapsto P_t: [0,1] \to F(H^s(E),H^{s-m}(F))\)连续，指标沿道路不变。

更精细地，需验证当\(P\)连续变化为椭圆算子时，\(\dim\ker P\)和\(\dim\operatorname{coker} P\)各自可能跳跃，但差不变。证明：若\(P\)和\(Q\)可通过紧扰动相连接（即\(P-Q \in K\)），则\(\operatorname{ind}(P) = \operatorname{ind}(Q)\)（由Atkinson定理）。椭圆算子族\(P_t\)的主象征连续变化，而低阶项（导致紧扰动）连续变化。因此\(P_t\)的原始拟逆构造连续依赖于\(t\)，符号演算保持不变，指标为整数且连续变化（由局部常值性），从而恒定。 \(\blacksquare\)

\subsection{x.3 Dirac算子的Lichnerowicz公式}\label{x.3-diracux7b97ux5b50ux7684lichnerowiczux516cux5f0f}

\textbf{定义 x.3}（Dirac算子）。设\(M\)为自旋流形，自旋丛为\(S \to M\)。Dirac算子\(D: \Gamma(S) \to \Gamma(S)\)定义为

\[D = \sum_{i=1}^{n} e_i \cdot \nabla_{e_i}\]

其中\(\{e_i\}\)为局部正交标架，\(\nabla\)为自旋联络，\(\cdot\)为Clifford乘法。

\textbf{定理 x.3}（Lichnerowicz公式，1963）。Dirac算子的平方具有以下分解

\[D^2 = \nabla^*\nabla + \frac{1}{4}R\]

其中\(\nabla^*\nabla\)为Bochner Laplacian（正定），\(R\)为标量曲率。

\textbf{证明}：由局部正交标架\(\{e_i\}\)，Dirac算子\(D = \gamma^i\nabla_{e_i}\)。计算平方：

\[D^2 = \gamma^i\gamma^j\nabla_{e_i}\nabla_{e_j} = \sum_{i,j} \gamma^i\gamma^j\nabla_{e_i}\nabla_{e_j}\]

将对称部分和反对称部分分离。反对称部分：\(\sum_{i<j}\gamma^i\gamma^j[\nabla_{e_i},\nabla_{e_j}]\)，其中\([\nabla_{e_i},\nabla_{e_j}] = R(e_i,e_j)\)（曲率算子在自旋丛上的作用）。曲率在自旋丛上表现为\(R(e_i,e_j)\psi = \frac{1}{4}R_{ijkl}\gamma^k\gamma^l\psi\)。利用Clifford代数恒等式\(\gamma^i\gamma^j\gamma^k\gamma^l\)的缩并由Bianchi恒等式简化为\(\frac{1}{2}R_{ijij}\gamma^i\gamma^j\)和其他项，最终得\(\frac{1}{4}R\)（其中\(R = \sum_{i,j} R_{ijij}\)为标量曲率）。对称部分：\(\sum_i\nabla_{e_i}\nabla_{e_i} - \sum_i\nabla_{\nabla_{e_i}e_i}\)，利用联络的度规相容性，化为\(-\nabla^*\nabla\)（因\(\nabla_{e_i}e_i\)在正交标架下为零）。结合两部分得\(D^2 = \nabla^*\nabla + R/4\)。\(\blacksquare\)

\subsection{x.4 Atiyah-Singer指标定理的K理论证明}\label{x.4-atiyah-singerux6307ux6807ux5b9aux7406ux7684kux7406ux8bbaux8bc1ux660e}

\textbf{定理 x.4}（Atiyah-Singer指标定理，1963）。设\(M\)为紧致无边自旋流形，\(P\)为椭圆微分算子（符号类\([\sigma_P] \in K(TM)\)）。则

\[\operatorname{ind}(P) = (-1)^{n(n+1)/2}\int_M \operatorname{ch}(\sigma_P) \wedge \operatorname{Td}(TM \otimes \mathbb{C})\]

其中\(\operatorname{ch}\)为Chern特征标，\(\operatorname{Td}\)为Todd类。对Dirac算子\(D\)，\(\operatorname{ind}(D) = \int_M \widehat{A}(M)\)（\(\widehat{A}\)亏格）。

\textbf{K理论框架}：指标函数给出了K理论同态\(\operatorname{ind}: K(TM) \to \mathbb{Z}\)。关键步骤：（1）通过Bott周期性和Thom同构，\(K(TM) \cong K(M \times M) \cong K(M) \otimes K(M)\)，椭圆算子的符号对应于这两个因子的张量积；（2）对球丛的Thom类构造\(K\)理论基本类和整性定理；（3）指标定理等价于交换图

\[\begin{CD} K(TM) @>{\operatorname{ind}}>> \mathbb{Z} \\ @V{\text{拓扑}}VV @| \\ \mathbb{R} @>{\int_M \operatorname{ch} \wedge \operatorname{Td}}>> \mathbb{R} \end{CD}\]

的可换性。直接验证只需在生成元（Hopf纤维丛上的Dirac算子）上检查，再由配边不变性扩展到一般情况。 \(\blacksquare\)

\subsection{x.5 热核证明方法}\label{x.5-ux70edux6838ux8bc1ux660eux65b9ux6cd5}

\textbf{定义 x.4}（热核）。算子\(D^2\)的热核\(K(t,x,y)\)是热方程\(\frac{\partial}{\partial t}u + D^2 u = 0\)的基本解，满足\(K(t,x,y) \to \delta(x-y)\)（当\(t \to 0\)）。

\textbf{定理 x.5}（McKean-Singer公式，完整推导）。对自伴椭圆算子\(D\)，

\[\operatorname{ind}(D) = \dim\ker D - \dim\ker D^* = \operatorname{Tr}(e^{-tD^*D}) - \operatorname{Tr}(e^{-tDD^*})\]

\[= \operatorname{Tr}_s(e^{-tD^2}) \quad (\text{超迹})\]

其中此关系对任意\(t > 0\)成立。

\textbf{推导}：首先，\(\operatorname{ind}(D) = \dim\ker D - \dim\ker D^* = \dim\ker(D^*D) - \dim\ker(DD^*)\)。因为\(\ker D = \ker(D^*D)\)且\(\ker D^* = \ker(DD^*)\)。对任意\(t > 0\)，热算子\(e^{-tD^*D}\)在\(D^*D\)的正特征值空间上衰减，但在核上恒为恒等算子。因此\(\operatorname{Tr}(e^{-tD^*D}) = \dim\ker(D^*D) + \sum_{\lambda > 0} e^{-t\lambda}\)，其中\(\lambda\)跑遍\(D^*D\)的正特征值。同理\(\operatorname{Tr}(e^{-tDD^*}) = \dim\ker(DD^*) + \sum_{\mu > 0} e^{-t\mu}\)。关键点：\(D^*D\)和\(DD^*\)的正特征值完全相同（由\(D\)和\(D^*\)交换关系推得）。因此迹的差中正谱贡献抵消，得

\[\operatorname{Tr}(e^{-tD^*D}) - \operatorname{Tr}(e^{-tDD^*}) = \dim\ker(D^*D) - \dim\ker(DD^*) = \operatorname{ind}(D)\]

这在所有\(t > 0\)处成立。特别令\(t \to 0\)，热核的渐近展开\(K(t,x,x) \sim \frac{1}{(4\pi t)^{n/2}}\sum_{k=0}^{\infty} t^k a_k(x)\)的系数\(a_k\)由局部几何量给出。通过超迹\(\operatorname{Tr}_s = \operatorname{Tr}(\Gamma e^{-tD^2})\)（其中\(\Gamma\)为\(\mathbb{Z}_2\)分次算子）和\(t\)的任意性，得到指标

\[\operatorname{ind}(D) = \frac{1}{(4\pi)^{n/2}}\int_M a_{n/2}(x) d\operatorname{vol}\]

对Dirac算子，\(a_{n/2}\)恰为\(\widehat{A}\)-亏格形式，复现Atiyah-Singer公式。 \(\blacksquare\)

\subsection{x.6 Gauß-Bonnet-Chern定理作为特例}\label{x.6-gauuxdf-bonnet-chernux5b9aux7406ux4f5cux4e3aux7279ux4f8b}

\textbf{定理 x.6}（Gauß-Bonnet-Chern定理）。设\(M\)为紧致无边\(2m\)维可定向Riemann流形。则Euler示性数\(\chi(M)\)等于

\[\chi(M) = \frac{1}{(2\pi)^m}\int_M \operatorname{Pf}(\Omega)\]

其中\(\operatorname{Pf}(\Omega)\)为曲率形式\(\Omega\)的Pfaffian。

\textbf{指标定理推导}：考虑de Rham复形\(d+d^*: \Omega^{\text{even}} \to \Omega^{\text{odd}}\)。该算子是椭圆的（因为其平方为Hodge Laplacian\(\Delta = dd^*+d^*d\)）。其指标为\(\operatorname{ind}(d+d^*) = \dim\ker(d+d^*)|_{\text{even}} - \dim\ker(d+d^*)|_{\text{odd}} = \sum_{k=0}^{2m}(-1)^k \dim\mathcal{H}^k = \chi(M)\)。而Atiyah-Singer指标定理给出\(\chi(M) = \int_M e(TM)\)，其中\(e(TM)\)为Euler类。Chern-Weil理论将Euler类表示为曲率Pfaffian的de Rham代表：

\[e(TM) = \frac{1}{(2\pi)^m}\operatorname{Pf}(\Omega) \in H^{2m}(M,\mathbb{R})\]

从而 \(\chi(M) = \frac{1}{(2\pi)^m}\int_M \operatorname{Pf}(\Omega)\)，证得Gauß-Bonnet-Chern定理。 \(\blacksquare\)

\textbf{定理 x.7}（Lichnerowicz 消没定理 / Vanishing Theorem, 1963）：设 \(M\) 为紧致无边自旋流形。若 \(M\) 具有正标量曲率 \(R > 0\)（处处为正），则 \(\widehat{A}\)-亏格为零：\(\widehat{A}(M) = 0\)。特别地，\(\ker D = 0\)------即 Dirac 算子在正标量曲率下无非零调和旋量。

\textbf{证明}：由 Lichnerowicz 公式（定理 x.3），\(D^2 = \nabla^*\nabla + R/4\)。若 \(\psi \in \ker D\)，则 \(D\psi = 0\)，故 \(D^2\psi = 0\)。取 \(L^2\) 内积：

\[0 = \langle D^2 \psi, \psi \rangle_{L^2} = \langle \nabla^*\nabla \psi, \psi \rangle + \frac{1}{4} \langle R \psi, \psi \rangle\]

第一项 \(\langle \nabla^*\nabla \psi, \psi \rangle = \|\nabla \psi\|_{L^2}^2\)（由分部积分，边界项在紧致无边流形上消失）。第二项 \(\langle R \psi, \psi \rangle_{L^2} = \int_M R |\psi|^2 d\operatorname{vol}\)。因此：

\[0 = \|\nabla \psi\|_{L^2}^2 + \frac{1}{4} \int_M R |\psi|^2 d\operatorname{vol}\]

若 \(R > 0\) 处处成立，右端第二项非负，第一项 \(\geq 0\)。故 \(\int_M R |\psi|^2 d\operatorname{vol} = 0\)，由 \(R > 0\) 的严格正性得 \(\psi \equiv 0\)。从而 \(\ker D = \{0\}\)。由指标定理 \(\operatorname{ind}(D) = \dim\ker D - \dim\ker D^* = 0\)，且 \(\operatorname{ind}(D) = \widehat{A}(M)\)，故 \(\widehat{A}(M) = 0\)。\(\blacksquare\)

此消没定理有深远推论：正标量曲率紧自旋流形不允许非零调和旋量，其 \(\widehat{A}\)-亏格为零------这给出了正标量曲率度量的拓扑障碍。Gromov-Lawson（1980）和 Stolz（1992）在此基础上完全分类了单连通紧流形上正标量曲率度量的存在性。

\textbf{定理 x.8}（指标定理的配边不变性）：设 \(M_1\) 和 \(M_2\) 为配边（cobordant）的紧致无边自旋流形（即存在紧致自旋流形 \(W\) 使得 \(\partial W = M_1 \sqcup \overline{M}_2\)），则对任意在 \(W\) 上可延拓的椭圆算子族，其指标在 \(M_1\) 和 \(M_2\) 上相等。

\textbf{证明}：设 \(P_1\) 和 \(P_2\) 分别为 \(M_1\) 和 \(M_2\) 上的椭圆算子，且均延拓为 \(W\) 上的椭圆算子 \(P_W\)。由 Atiyah-Bott 指标定理的配边版本，\[\operatorname{ind}(P_1) - \operatorname{ind}(P_2) = \operatorname{ind}(P_W, \partial W)\] 其中 \(\operatorname{ind}(P_W, \partial W)\) 是带边流形上具有适当边界条件的椭圆算子的指标。配边不变性的实质在于指标可写为示性类沿整个配边的积分：由 Stokes 定理，

\[\operatorname{ind}(P_1) - \operatorname{ind}(P_2) = \int_{M_1} \omega - \int_{M_2} \omega = \int_W d\omega = 0\]

其中 \(\omega\) 为 Atiyah-Singer 被积形式（\(\operatorname{ch}(\sigma_P) \wedge \operatorname{Td}(TM \otimes \mathbb{C})\) 的微分形式代表），而 \(d\omega = 0\)（在 Chen-Weil 意义下 \([\omega]\) 为上同调类）。配边不变性本质上等价于 \(\operatorname{ind}\) 作为配边环 \(\Omega_*^{\text{Spin}}(pt)\) 上的同态因子通过 K 群。这正是热核证明和 K 理论证明共有的配边稳定性。\(\blacksquare\)

\hrulefill

\section{第278章：Morse理论与变分法}\label{ux7b2c278ux7ae0morseux7406ux8bbaux4e0eux53d8ux5206ux6cd5}

Morse理论是将流形的拓扑与其上的光滑函数的临界点联系起来的深刻工具。Morse观察到，流形的拓扑可以通过考察Morse函数临界点处的Hessian指标来确定------每个指标为\(k\)的临界点对应于添加一个\(k\)维胞腔。Morse不等式给出了临界点个数与Betti数之间的约束关系，而Morse同调的建立使得这一联系更加精确：可以通过对Morse函数的梯度流定义链复形，其同调仅依赖于底流形的拓扑。本章完整推导Morse引理、Morse不等式，并简述Morse同调与Floer同调的联系。

\subsection{x+1.1 Morse引理}\label{x1.1-morseux5f15ux7406}

\textbf{定义 x+1.1}（Morse函数）。光滑函数\(f: M^n \to \mathbb{R}\)称为Morse函数，若其所有临界点非退化，即Hessian矩阵\(H_f(p) = (\partial^2 f / \partial x_i \partial x_j)|_p\)非退化。

\textbf{定义}：临界点\(p\)的指标\(\lambda(p)\)定义为Hessian\(H_f(p)\)的负特征值的个数（Morse指标）。

\textbf{定理 x+1.1}（Morse引理，Morse坐标系的存在性，完整证明）。设\(p\)为Morse函数\(f\)的非退化临界点，指标为\(\lambda\)。则存在\(p\)附近的局部坐标\((y_1,\dots,y_n)\)（Morse坐标系）使得

\[f(y) = f(p) - y_1^2 - \cdots - y_\lambda^2 + y_{\lambda+1}^2 + \cdots + y_n^2\]

\textbf{证明}：由Taylor公式和Morse引理的推广形式。不妨设\(f(p)=0\)且取局部坐标\((x_1,\dots,x_n)\)使\(p\)为原点。由Hadamard引理，存在光滑函数\(g_{ij}(x)\)满足\(g_{ij}=g_{ji}\)使得\(f(x) = \sum_{i,j} g_{ij}(x)x_i x_j\)。Hessian矩阵为\([2g_{ij}(0)]\)，由非退化性可对角化。

通过归纳法构造坐标变换。对指标\(\lambda\)的情况，每次消去一个二次项。具体而言，利用\(g_{ii}(0) \neq 0\)且符号为\(\pm 1\)，通过线性变换和平方项配方将\(f\)简化为\(\pm y_1^2 + f_1(y_2,\dots,y_n)\)的形式。逐次操作\(n\)次得到所需标准型。每一步配方均保持Morse性质（其余变量的Hessian仍为非退化），从而过程有效。 \(\blacksquare\)

\subsection{x+1.2 Morse不等式}\label{x1.2-morseux4e0dux7b49ux5f0f}

\textbf{定理 x+1.2}（弱Morse不等式，完整证明）。设\(f: M \to \mathbb{R}\)为紧致无边流形\(M\)上的Morse函数，\(c_k\)为指标为\(k\)的临界点的个数，\(b_k = \operatorname{rk} H_k(M;\mathbb{F})\)为第\(k\)个Betti数（\(\mathbb{F}\)为某域）。则

\[c_k \geq b_k, \quad k = 0,1,\dots,n\]

\textbf{证明}：考虑水平集\(M^a = f^{-1}((-\infty,a])\)。当\(a\)不经过临界值时，\(M^a\)随\(a\)的微分同胚类型不变（由梯度流给出形变收缩）。当\(a\)经过指标为\(k\)的临界值时，\(M^{a}\)通过添加一个\(k\)维胞腔（\(k\)-handle）而变换其同伦型：\(M^{a+\varepsilon} \simeq M^{a-\varepsilon} \cup h^k\)（Morse理论基本定理，由Morse引理和显式的梯度流附着映射得出）。

由此，\(M\)同伦等价于一个CW复形\(X\)，其\(k\)维胞腔的个数为\(c_k\)。对\(X\)的胞腔链复形\((C_*(X),\partial)\)，有\(\dim C_k = c_k\)。同调群\(H_k(X;\mathbb{F}) = H_k(M;\mathbb{F})\)是商空间\(Z_k/B_k\)，其中\(Z_k \subseteq C_k\)为闭链子空间，\(B_k \subseteq Z_k\)为边界子空间。于是\(\dim H_k = \dim Z_k - \dim B_k \leq \dim Z_k \leq \dim C_k = c_k\)。故\(c_k \geq b_k\)。 \(\blacksquare\)

\textbf{定理 x+1.3}（强Morse不等式，完整证明）。在上述条件下，强Morse不等式成立：

\[\sum_{k=0}^{m}(-1)^{m-k}c_k \geq \sum_{k=0}^{m}(-1)^{m-k}b_k, \quad m = 0,1,\dots,n-1\]

且当\(m=n\)时取等号：\(\sum_{k=0}^{n}(-1)^k c_k = \sum_{k=0}^{n}(-1)^k b_k = \chi(M)\)。

\textbf{证明}：从CW复形\(X\)的胞腔链复形\((C_k,\partial_k)\)出发。设\(Z_k = \ker\partial_k\)，\(B_k = \operatorname{im}\partial_{k+1}\)。定义以下维数关系：\(c_k = \dim C_k\)，\(z_k = \dim Z_k\)，\(b_k = \dim H_k = z_k - \dim B_k\)。由正合列\(0 \to Z_k \to C_k \xrightarrow{\partial_k} B_{k-1} \to 0\)得\(c_k = z_k + \dim B_{k-1}\)。代入代数运算：

\[\sum_{k=0}^{m}(-1)^{m-k}(c_k - b_k) = \sum_{k=0}^{m}(-1)^{m-k}(\dim B_k + \dim B_{k-1})\]

右端展开为\(\dim B_m + \dim B_{-1}\)（\(\dim B_{-1}=0\)），故\(\sum_{k=0}^{m}(-1)^{m-k}(c_k - b_k) = \dim B_m \geq 0\)。得证强不等式。当\(m=n\)时\(B_n=0\)，得Euler示性数的等式。 \(\blacksquare\)

\subsection{x+1.3 Morse同调与Smale条件}\label{x1.3-morseux540cux8c03ux4e0esmaleux6761ux4ef6}

\textbf{定义 x+1.2}（Smale条件）。设\(f\)为Morse函数，\(g\)为Riemann度量。称\((f,g)\)满足Smale条件，若对任意两临界点\(p,q\)（\(\operatorname{ind}(p) \leq \operatorname{ind}(q)\)），梯度流\(-\nabla_g f\)的稳定流形\(W^s(p)\)和不稳定流形\(W^u(q)\)横截相交。

\textbf{定义 x+1.3}（Morse同调链复形）。对满足Smale条件的\((f,g)\)，定义\(C_k^{\text{Morse}}\)为由指标\(k\)临界点生成的\(\mathbb{Z}\)-模。边界算子\(\partial: C_k \to C_{k-1}\)的定义为：

对于指标为\(k\)的临界点\(p\)和指标为\(k-1\)的临界点\(q\)，\(p\)到\(q\)的计数\(n(p,q)\)为从\(p\)到\(q\)的负梯度流线（模平移）的（带符号的）个数：

\[\partial p = \sum_{\operatorname{ind}(q)=k-1} n(p,q) q\]

\textbf{定理 x+1.4}（Morse同调定理）。\(\partial^2 = 0\)，且由此定义的Morse同调群\(HM_*(f,g)\)典范同构于奇异同调群\(H_*(M;\mathbb{Z})\)，因而独立于\((f,g)\)的选择。

\textbf{证明思路}：\(\partial^2 = 0\)的证明利用了断裂轨迹的紧化与维数分析------指标差为2的临界点间的轨迹空间在紧化后具有一维边界流形，其边界恰好对应经一中间临界点的轨迹配对，符号抵消导致\(\partial^2=0\)。同构于奇异同调的证明通过比较Morse链复形与CW胞腔复形，两者链同伦等价。 \(\blacksquare\)

\subsection{x+1.4 Floer同调简介}\label{x1.4-floerux540cux8c03ux7b80ux4ecb}

\textbf{定义 x+1.4}（Hamiltonian Floer同调）。设\((M,\omega)\)为紧致辛流形，\(H: S^1 \times M \to \mathbb{R}\)为周期Hamilton函数。在环路空间\(\mathcal{L}M = \{\gamma: S^1 \to M\}\)上定义辛作用泛函

\[\mathcal{A}_H(\gamma) = \int_{D^2} u^*\omega - \int_0^1 H(t,\gamma(t)) dt\]

其临界点恰为Hamilton轨道（周期为1的Hamilton方程的解）。Floer复形由这些临界点生成，边界算子由连接轨道的（模平移的）负梯度流线（即cylinder上的伪全纯曲线）计数定义。

\textbf{定理 x+1.5}（Arnold猜想的Floer证明框架）。对辛流形\((M,\omega)\)上的非退化Hamilton微分同胚，其不动点个数至少等于\(M\)的各阶Betti数之和\(SB(M) = \sum_k b_k(M)\)。

\textbf{证明思路}：Floer同调\(HF_*(M,H)\)典范同构于Hamiltonian Floer同调\(HF_*(M,\omega)\)，后者又同构于量子同调（由Piunikhin-Salamon-Schwarz同构）。量子同调在经典极限下退化为通常同调。因此不动点（即Hamilton轨道的1-周期点，即Floer复形的生成元）至少为\(SB(M)\)个。这一论证将Morse理论从有限维流形推广到无限维环路空间。 \(\blacksquare\)

\hrulefill

\section{第279章：Chern-Weil理论与示性类}\label{ux7b2c279ux7ae0chern-weilux7406ux8bbaux4e0eux793aux6027ux7c7b}

Chern-Weil理论提供了用曲率形式构造示性类的微分几何方法。其核心思想极其优美：对向量丛上任一联络的曲率，任意不变多项式求值均给出一个闭微分形式，其de Rham上同调类不依赖于联络的选择，从而是丛的拓扑不变量。这一理论将代数拓扑的示性类与微分几何的曲率联系起来，使Gauß-Bonnet定理成为Chern类的自然产物。本章完整证明Chern-Weil同态的基本定理，然后建立各类示性类的公理体系，最后讨论陈-Simons形式这一次要示性类的重要概念。

\subsection{x+2.1 Chern联络与曲率}\label{x2.1-chernux8054ux7edcux4e0eux66f2ux7387}

\textbf{定义 x+2.1}（Chern联络）。设\(E \to M\)为秩\(r\)的复向量丛，其上赋予Hermite度量\(h\)。Chern联络\(\nabla\)是唯一满足以下条件的联络：（1）与度量相容，\(d(h(s_1,s_2)) = h(\nabla s_1,s_2) + h(s_1,\nabla s_2)\)；（2）在\(E\)上\((1,0)\)-部分为\(\bar{\partial}\)算子（即对于全纯标架\(\{e_\alpha\}\)，\(\nabla e_\alpha \in \Omega^{1,0}(M) \otimes E\)）。Chern联络的曲率形式\(\Omega = \nabla^2 = d\theta + \theta \wedge \theta\)是\((1,1)\)-型的\(r \times r\)矩阵值2-形式（\(\theta\)为联络1-形式）。

\subsection{x+2.2 Chern-Weil同态}\label{x2.2-chern-weilux540cux6001}

\textbf{定理 x+2.1}（Chern-Weil同态，完整证明）。设\(P: \mathfrak{gl}(r,\mathbb{C}) \to \mathbb{C}\)为\(k\)次齐次不变多项式（即对任意\(g \in GL(r,\mathbb{C})\)有\(P(gXg^{-1}) = P(X)\)）。则：

（1）\(P(\Omega)\)是闭形式：\(dP(\Omega) = 0\)；

（2）\(P(\Omega)\)的de Rham上同调类\([P(\Omega)] \in H^{2k}(M,\mathbb{C})\)不依赖于联络\(\nabla\)的选择。

\textbf{证明}：

（1）由Bianchi恒等式\(d\Omega = \Omega \wedge \theta - \theta \wedge \Omega = [\Omega,\theta]\)。对不变多项式\(P\)，有

\[dP(\Omega) = \sum_i \frac{\partial P}{\partial X_{ij}} d\Omega_{ij}\]

其中\(\frac{\partial P}{\partial X_{ij}}\)为\(P\)关于矩阵元的形式导数。利用不变性：对任意\(g \in GL(r)\)，\(P(gXg^{-1}) = P(X)\)。取\(g = I + tA\)并比较\(t\)的一阶项，得\(\sum_{i,j}\frac{\partial P}{\partial X_{ij}}([A,X])_{ij} = 0\)，即\(\sum_{i,j,k} \frac{\partial P}{\partial X_{ij}}(A_{ik}X_{kj} - X_{ik}A_{kj}) = 0\)。由\(A\)的任意性得\(\sum_i \frac{\partial P}{\partial X_{ij}} X_{ik} = \sum_i X_{ki} \frac{\partial P}{\partial X_{ji}}\)（\(\frac{\partial P}{\partial X}\)与\(X\)交换）。应用于Bianchi恒等式：

\[dP(\Omega) = \sum_{i,j} \frac{\partial P}{\partial X_{ij}}([\Omega,\theta])_{ij}\]

将\([\Omega,\theta] = \Omega\theta - \theta\Omega\)代入，利用\(\frac{\partial P}{\partial X}\)与\(\Omega\)交换，得

\[dP(\Omega) = \sum_{i,j,k}(\frac{\partial P}{\partial X_{ij}}\Omega_{ik}\theta_{kj} - \frac{\partial P}{\partial X_{ij}}\theta_{ik}\Omega_{kj})\]

在第一项中交换\(i,k\)的缩并次序并使用不变性关系，两项抵消。因此\(dP(\Omega) = 0\)。

（2）设\(\nabla_0\)和\(\nabla_1\)为两联络，曲率分别为\(\Omega_0,\Omega_1\)。考虑插值\(\nabla_t = (1-t)\nabla_0 + t\nabla_1\)（这是一个联络，因联络的仿射空间结构），其曲率为\(\Omega_t\)。定义Chern-Simons超度形式

\[\operatorname{CS}(\nabla_0,\nabla_1) = k\int_0^1 P(\theta_1 - \theta_0, \Omega_t, \dots, \Omega_t) dt\]

（其中\(P\)作为极化不变多项式在第一个变元取\(\theta_1-\theta_0\)，其余\(k-1\)个取\(\Omega_t\)）。由transgression公式

\[d\operatorname{CS}(\nabla_0,\nabla_1) = P(\Omega_1) - P(\Omega_0)\]

因此\([P(\Omega_0)] = [P(\Omega_1)]\)在\(H^{2k}(M,\mathbb{C})\)中。 \(\blacksquare\)

\subsection{x+2.3 Chern类的公理化描述}\label{x2.3-chernux7c7bux7684ux516cux7406ux5316ux63cfux8ff0}

\textbf{定义 x+2.2}（全Chern类）。全Chern类\(c(E) = 1 + c_1(E) + c_2(E) + \cdots + c_r(E)\)通过Chern-Weil同态由以下多项式生成

\[\det\left(I + \frac{i}{2\pi}\Omega\right) = \sum_{k=0}^{r} c_k(\Omega)\]

其中\(c_k(\Omega) \in \Omega^{2k}(M)\)为闭形式，代表第\(k\)个Chern类\(c_k(E) \in H^{2k}(M,\mathbb{Z})\)。

\textbf{定理 x+2.2}（Chern类的公理）。Chern类是满足以下性质的上同调类族：

（1）\(c_0(E) = 1\)；（2）自然性：\(c(f^*E) = f^*c(E)\)对任意光滑映射\(f\)；（3）Whitney和公式：\(c(E \oplus F) = c(E) \cup c(F)\)；（4）对\(M = \mathbb{CP}^1\)上的超平面线丛\(H\)，\(c_1(H) = 1 \cdot [\omega_{\text{FS}}]\)（Fubini-Study度规的Kähler形式）。

\textbf{证明（唯一性）}：由分裂原理：任意向量丛可通过拉回到某适当流形（旗丛）上分裂为线丛的直和。性质（2）和（3）将Chern类的计算完全归约为线丛。而线丛的\(c_1\)由性质（4）和（2）唯一决定（因\(\mathbb{CP}^1\)上的超平面线丛是万有线丛的生成元）。因此上述公理唯一确定Chern类。由Chern-Weil构造给出的封闭微分形式满足所有这些性质，故两者一致。 \(\blacksquare\)

\subsection{x+2.4 Pontryagin类与Euler类}\label{x2.4-pontryaginux7c7bux4e0eeulerux7c7b}

\textbf{定义 x+2.3}（Pontryagin类）。实向量丛\(E\)的Pontryagin类由其复化丛\(E \otimes \mathbb{C}\)的Chern类定义：

\[p_k(E) = (-1)^k c_{2k}(E \otimes \mathbb{C}) \in H^{4k}(M,\mathbb{Z})\]

对Riemann度量下的联络，曲率\(\Omega\)为\(\mathfrak{so}(n)\)值（\(\mathfrak{so}(n)\)矩阵为反对称的），其非零特征值成对出现。Chern-Weil同态中的不变多项式由反对称矩阵的Pfaffian或特征多项式的偶次项给出。

\textbf{定义 x+2.4}（Euler类）。可定向实向量丛\(E\)（秩为偶数\(2m\)）的Euler类\(e(E) \in H^{2m}(M,\mathbb{Z})\)由Pfaffian构造：

\[e(\Omega) = \frac{1}{(2\pi)^m}\operatorname{Pf}(\Omega)\]

其中\(\Omega\)为Riemann联络下的曲率形式。

\subsection{x+2.5 陈-Simons形式与次要示性类}\label{x2.5-ux9648-simonsux5f62ux5f0fux4e0eux6b21ux8981ux793aux6027ux7c7b}

\textbf{定义 x+2.5}（陈-Simons形式）。在Chern-Weil定理的证明中出现的超度形式\(\operatorname{CS}(\nabla_0,\nabla_1)\)称为陈-Simons形式。当\(\nabla_0\)为平凡联络（如平凡丛上的\(d\)）时，记\(\operatorname{CS}(\nabla) = \operatorname{CS}(d,\nabla)\)。

对3维流形\(M^3\)上的平凡\(G\)-丛，陈-Simons作用量为

\[\operatorname{CS}(A) = \frac{1}{8\pi^2}\int_{M^3}\operatorname{Tr}\left(A \wedge dA + \frac{2}{3}A \wedge A \wedge A\right)\]

其中\(A\)为联络1-形式。陈-Simons形式本身不是上同调不变量（它是微分形式而非闭形式），其积分（模\(2\pi\mathbb{Z}\)）是流形-联络对的拓扑不变量，在Chern-Simons规范理论和Jones多项式理论中起核心作用。

\textbf{定理 x+2.3}（陈-Simons不变量的规范变换）。在规范变换\(g: M^3 \to G\)下，陈-Simons作用量变为

\[\operatorname{CS}(A^g) = \operatorname{CS}(A) - \frac{1}{24\pi^2}\int_{M^3}\operatorname{Tr}(g^{-1}dg)^{\wedge 3}\]

后一项是规范群元\(g\)的绕数（winding number），取值为整数。因此\(\operatorname{CS}(A) \mod 2\pi\mathbb{Z}\)是规范不变的。

\textbf{证明思路}：由transgression公式在边界上的变化关系推导。考虑4维流形\(M^4 = \partial B^5\)，将规范势延拓至内部，令\(F = dA + A \wedge A\)为场强。由陈-Weil定理：\(d(\operatorname{Tr}(F \wedge F)) = 0\)，且其积分为整Chern-Simons数：\(\frac{1}{8\pi^2}\int_{M^4}\operatorname{Tr}(F \wedge F) \in \mathbb{Z}\)。但\(\operatorname{Tr}(F \wedge F) = d\operatorname{CS}(A)\)，故\(\frac{1}{8\pi^2}\int_{M^3}\operatorname{CS}(A)\)在\(A\)属于不同拓扑类的延拓之间相差整数。 \(\blacksquare\)

\newpage

\chapter{07-拓扑与几何/01-拓扑学/01-拓扑空间与基本群.md}\label{ux62d3ux6251ux4e0eux51e0ux4f5501-ux62d3ux6251ux5b6601-ux62d3ux6251ux7a7aux95f4ux4e0eux57faux672cux7fa4.md}

\section{第212章：拓扑空间基础}\label{ux7b2c212ux7ae0ux62d3ux6251ux7a7aux95f4ux57faux7840}

拓扑学作为现代数学的基础语言之一，其核心思想是将分析学中''连续''与''极限''的直觉剥离距离概念的依赖，提炼为仅基于开集公理的纯粹集合论结构。这一''去度量化''的抽象由 Felix Hausdorff 在 1914 年的经典著作《Grundzuge der Mengenlehre》中首次系统提出，随后被 Kuratowski、Alexandroff 和 Bourbaki 学派进一步精炼。拓扑空间的公理化不仅统一了度量空间（由距离诱导拓扑）、代数簇（Zariski 拓扑）和函数空间（逐点收敛拓扑，如乘积拓扑与紧开拓扑）等诸多看似迥异的结构，还使''同胚''成为比''等距同构''更灵活的分类工具，允许在不破坏连续性结构的前提下变形空间。在本书的拓扑与几何篇体系中，本章是所有后续章节的共同起点：连通性与紧致性（第152章）深化了空间的全局性质，分离公理（第153章）为从拓扑到度量的过渡提供了分级阶梯，而基本群（第154章）则在拓扑空间中引入第一个代数不变量，开启了代数拓扑的篇章。

\subsection{57.1 拓扑空间的定义}\label{ux62d3ux6251ux7a7aux95f4ux7684ux5b9aux4e49}

\textbf{定义 57.1}（拓扑空间）：设 \(X\) 是一个集合，\(\mathcal{T} \subseteq \mathcal{P}(X)\) 是 \(X\) 的子集族。称 \(\mathcal{T}\) 为 \(X\) 上的一个\textbf{拓扑}，如果满足以下三条公理：

\begin{enumerate}
\def\labelenumi{\arabic{enumi}.}
\tightlist
\item
  \textbf{(T1)} \(\emptyset \in \mathcal{T}\)，\(X \in \mathcal{T}\)。
\item
  \textbf{(T2)} \(\mathcal{T}\) 中任意多个元素的并集仍属于 \(\mathcal{T}\)：若 \(\{U_\alpha\}_{\alpha \in A} \subseteq \mathcal{T}\)，则 \(\bigcup_{\alpha \in A} U_\alpha \in \mathcal{T}\)。
\item
  \textbf{(T3)} \(\mathcal{T}\) 中有限多个元素的交集仍属于 \(\mathcal{T}\)：若 \(U_1, U_2, \ldots, U_n \in \mathcal{T}\)，则 \(U_1 \cap U_2 \cap \cdots \cap U_n \in \mathcal{T}\)。
\end{enumerate}

称 \((X, \mathcal{T})\) 为一个\textbf{拓扑空间}，\(\mathcal{T}\) 中的元素称为\textbf{开集}。在不引起混淆时简记为 \(X\)。

\textbf{定义 57.2}（闭集）：设 \((X, \mathcal{T})\) 为拓扑空间。子集 \(F \subseteq X\) 称为\textbf{闭集}，如果 \(X \setminus F \in \mathcal{T}\)（即 \(F\) 的补集是开集）。

\textbf{命题 57.1}（闭集的性质）：设 \((X, \mathcal{T})\) 为拓扑空间，则： 1. \(\emptyset\) 和 \(X\) 是闭集。 2. 任意多个闭集的交集是闭集。 3. 有限多个闭集的并集是闭集。

\emph{证明}：由 De Morgan 律直接可得。∎

\textbf{定义 57.3}（邻域）：设 \(X\) 为拓扑空间，\(x \in X\)。子集 \(N \subseteq X\) 称为 \(x\) 的一个\textbf{邻域}，如果存在开集 \(U \in \mathcal{T}\) 使得 \(x \in U \subseteq N\)。\(x\) 的所有邻域构成的集合称为 \(x\) 的\textbf{邻域系}，记作 \(\mathcal{N}(x)\)。

\textbf{定义 57.4}（内部、闭包、边界）：设 \(X\) 为拓扑空间，\(A \subseteq X\)。

\begin{itemize}
\tightlist
\item
  \(A\) 的\textbf{内部} \(\operatorname{int}(A)\)（或 \(A^\circ\)）：包含在 \(A\) 中的所有开集的并集，即 \(A\) 的最大开子集。
\item
  \(A\) 的\textbf{闭包} \(\overline{A}\)（或 \(\operatorname{cl}(A)\)）：包含 \(A\) 的所有闭集的交集，即包含 \(A\) 的最小闭集。
\item
  \(A\) 的\textbf{边界} \(\partial A\)：\(\partial A = \overline{A} \setminus \operatorname{int}(A)\)。
\end{itemize}

\textbf{命题 57.2}（内部与闭包的关系）：设 \(X\) 为拓扑空间，\(A \subseteq X\)。则： 1. \(x \in \operatorname{int}(A)\) 当且仅当 \(A \in \mathcal{N}(x)\)。 2. \(x \in \overline{A}\) 当且仅当 \(x\) 的每个邻域与 \(A\) 相交。 3. \(\overline{A} = X \setminus \operatorname{int}(X \setminus A)\)，\(\operatorname{int}(A) = X \setminus \overline{X \setminus A}\)。

\emph{证明}：(1) 由内部和邻域的定义直接可得。(2) 若 \(x \notin \overline{A}\)，则存在闭集 \(F \supseteq A\) 使 \(x \notin F\)，于是 \(X \setminus F\) 是 \(x\) 的邻域且与 \(A\) 不相交。反之亦然。(3) 由 De Morgan 律可得。∎

\textbf{定义 57.5}（聚点、孤立点）：设 \(X\) 为拓扑空间，\(A \subseteq X\)，\(x \in X\)。

\begin{itemize}
\tightlist
\item
  \(x\) 称为 \(A\) 的\textbf{聚点}（或极限点），如果 \(x\) 的每个邻域都包含 \(A\) 中异于 \(x\) 的点。
\item
  \(A\) 的所有聚点构成的集合称为 \(A\) 的\textbf{导集}，记作 \(A'\)。
\item
  \(x \in A\) 称为 \(A\) 的\textbf{孤立点}，如果存在 \(x\) 的邻域 \(U\) 使得 \(U \cap A = \{x\}\)。
\end{itemize}

\textbf{命题 57.3}：\(\overline{A} = A \cup A'\)。

\emph{证明}：\(x \in \overline{A}\) 当且仅当 \(x\) 的每个邻域与 \(A\) 相交，这等价于 \(x \in A\) 或 \(x\) 的每个邻域包含 \(A\) 中异于 \(x\) 的点（即 \(x \in A'\)）。∎

\subsection{57.2 拓扑的基与子基}\label{ux62d3ux6251ux7684ux57faux4e0eux5b50ux57fa}

\textbf{定义 57.6}（基）：设 \((X, \mathcal{T})\) 为拓扑空间。子族 \(\mathcal{B} \subseteq \mathcal{T}\) 称为 \(\mathcal{T}\) 的一个\textbf{基}，如果 \(\mathcal{T}\) 中的每个开集都可以表示为 \(\mathcal{B}\) 中某些元素的并集。

\textbf{定理 57.4}（基的判定条件）：设 \(X\) 为集合，\(\mathcal{B} \subseteq \mathcal{P}(X)\)。则 \(\mathcal{B}\) 是 \(X\) 上某个拓扑的基当且仅当： 1. \(\bigcup_{B \in \mathcal{B}} B = X\)。 2. 对任意 \(B_1, B_2 \in \mathcal{B}\) 和 \(x \in B_1 \cap B_2\)，存在 \(B_3 \in \mathcal{B}\) 使得 \(x \in B_3 \subseteq B_1 \cap B_2\)。

此时由 \(\mathcal{B}\) 生成的拓扑 \(\mathcal{T}\) 为：\(U \in \mathcal{T}\) 当且仅当对每个 \(x \in U\)，存在 \(B \in \mathcal{B}\) 使得 \(x \in B \subseteq U\)。

\emph{证明}：(\(\Rightarrow\)) 若 \(\mathcal{B}\) 是拓扑 \(\mathcal{T}\) 的基，则 \(X \in \mathcal{T}\) 可表示为 \(\mathcal{B}\) 中元素的并，故 (1) 成立。对 \(B_1, B_2 \in \mathcal{B}\)，\(B_1 \cap B_2 \in \mathcal{T}\)，故可表示为 \(\mathcal{B}\) 中元素的并，从而 (2) 成立。

(\(\Leftarrow\)) 设条件成立。定义 \(\mathcal{T}\) 为满足所述条件的子集族。验证 \(\mathcal{T}\) 是拓扑： - \(\emptyset\) 平凡地属于 \(\mathcal{T}\)；由 (1) 知 \(X \in \mathcal{T}\)。 - 任意并：若 \(\{U_\alpha\} \subseteq \mathcal{T}\)，对 \(x \in \bigcup U_\alpha\)，存在 \(\alpha\) 使 \(x \in U_\alpha\)，故存在 \(B \in \mathcal{B}\) 使 \(x \in B \subseteq U_\alpha \subseteq \bigcup U_\alpha\)。 - 有限交：对 \(U_1, U_2 \in \mathcal{T}\)，任取 \(x \in U_1 \cap U_2\)，存在 \(B_1, B_2 \in \mathcal{B}\) 使 \(x \in B_1 \subseteq U_1\)，\(x \in B_2 \subseteq U_2\)。由 (2) 存在 \(B_3 \in \mathcal{B}\) 使 \(x \in B_3 \subseteq B_1 \cap B_2 \subseteq U_1 \cap U_2\)。归纳得有限交。

显然 \(\mathcal{B}\) 是 \(\mathcal{T}\) 的基。∎

\textbf{定义 57.7}（子基）：设 \(X\) 为集合，\(\mathcal{S} \subseteq \mathcal{P}(X)\) 满足 \(\bigcup_{S \in \mathcal{S}} S = X\)。则 \(\mathcal{S}\) 中元素的所有有限交构成的族 \(\mathcal{B}\) 是某个拓扑的基，由 \(\mathcal{B}\) 生成的拓扑称为由 \(\mathcal{S}\) \textbf{生成的拓扑}，\(\mathcal{S}\) 称为该拓扑的\textbf{子基}。

\subsection{57.3 常见拓扑举例}\label{ux5e38ux89c1ux62d3ux6251ux4e3eux4f8b}

下面列举拓扑学中最基本和最重要的拓扑空间示例，它们构成了后续学习的基础素材库。

\begin{enumerate}
\def\labelenumi{\arabic{enumi}.}
\item
  \textbf{离散拓扑}：\(\mathcal{T} = \mathcal{P}(X)\)（\(X\) 的所有子集都是开集）。这是 \(X\) 上最精细的拓扑，每个单点集 \(\{x\}\) 都是开集。离散拓扑空间是局部紧致且完全不连通的，且任何映射从离散空间出发都是连续的。
\item
  \textbf{平凡拓扑}（或密着拓扑）：\(\mathcal{T} = \{\emptyset, X\}\)。这是 \(X\) 上最粗糙的拓扑，仅有两个开集。平凡拓扑空间是紧致的（因为只有有限个开覆盖需检查），但除非 \(|X| \leq 1\)，否则不是 Hausdorff 空间。
\item
  \textbf{度量拓扑}：由度量空间 \((X, d)\) 诱导的拓扑，以开球 \(B(x, \varepsilon) = \{y \in X : d(x, y) < \varepsilon\}\) 为基。这是最重要的一类拓扑空间，\(\mathbb{R}^n\) 配以标准欧氏度量是典型例子。度量拓扑总是 \(T_4\) 的（正规的），且序列收敛等价于网的收敛，极大简化了分析。
\item
  \textbf{有限补拓扑}（余有限拓扑）：\(\mathcal{T} = \{U \subseteq X : X \setminus U \text{ 是有限集}\} \cup \{\emptyset\}\)。在此拓扑下，闭集恰为有限集和 \(X\) 本身。有限补拓扑是 \(T_1\) 的（单点集为闭集），但除非 \(X\) 有限，否则不是 \(T_2\) 的。当 \(X\) 为无穷集时，有限补拓扑是紧致的。
\item
  \textbf{可数补拓扑}：\(\mathcal{T} = \{U \subseteq X : X \setminus U \text{ 是可数集}\} \cup \{\emptyset\}\)。类似有限补拓扑，但用可数替代有限。
\item
  \textbf{Sierpiński 空间}：\(X = \{0, 1\}\)，\(\mathcal{T} = \{\emptyset, \{0\}, \{0, 1\}\}\)。这是最简单的非平凡非离散拓扑空间，是 \(T_0\) 但不 \(T_1\)。Sierpiński 空间在同伦论和范畴论中经常作为''测试空间''出现。
\item
  \textbf{序拓扑}：在全序集 \((X, <)\) 上，以所有开区间 \((a, b) = \{x : a < x < b\}\)、\((-\infty, a)\) 和 \((a, \infty)\) 为基生成的拓扑。\(\mathbb{R}\) 的标准拓扑即其序拓扑，序数 \(\omega_1\)（第一个不可数序数）上的序拓扑是局部紧致、正规但非紧致、非可度量化的重要例子。
\item
  \textbf{Zariski 拓扑}（代数几何的基本拓扑）：在代数闭域 \(k\) 上的仿射空间 \(\mathbb{A}^n\) 中，闭集恰为代数集（多项式方程组的零点集）。Zariski 拓扑是 \(T_1\) 的但非 \(T_2\)，且极其''粗糙''------任意两个非空开集必相交（不可约性）。Zariski 拓扑是代数几何中概形（scheme）理论的基础。
\end{enumerate}

\textbf{定义 57.8}（可度量化空间）：拓扑空间 \((X, \mathcal{T})\) 称为\textbf{可度量化的}，如果存在 \(X\) 上的度量 \(d\) 使得 \(\mathcal{T}\) 恰为 \(d\) 诱导的拓扑。

\subsection{57.4 连续映射与同胚}\label{ux8fdeux7eedux6620ux5c04ux4e0eux540cux80da}

\textbf{定义 57.9}（连续映射）：设 \(X, Y\) 为拓扑空间，\(f: X \to Y\) 为映射。称 \(f\) 在点 \(x \in X\) 处\textbf{连续}，如果对 \(f(x)\) 的任意邻域 \(V\)，\(f^{-1}(V)\) 是 \(x\) 的邻域。称 \(f\) 为\textbf{连续映射}，如果 \(f\) 在每一点处连续。

\textbf{定理 57.5}（连续映射的等价刻画）：设 \(f: X \to Y\) 为拓扑空间之间的映射，则以下条件等价： 1. \(f\) 是连续映射。 2. 对 \(Y\) 中任意开集 \(V\)，\(f^{-1}(V)\) 是 \(X\) 中的开集。 3. 对 \(Y\) 中任意闭集 \(F\)，\(f^{-1}(F)\) 是 \(X\) 中的闭集。 4. 对任意 \(A \subseteq X\)，\(f(\overline{A}) \subseteq \overline{f(A)}\)。 5. 对 \(Y\) 的拓扑的某个基（或子基）\(\mathcal{B}\)，对每个 \(B \in \mathcal{B}\)，\(f^{-1}(B)\) 是 \(X\) 中的开集。

\emph{证明}：(1) \(\Rightarrow\) (2)：设 \(V \subseteq Y\) 为开集，\(x \in f^{-1}(V)\)。则 \(f(x) \in V\)，故 \(V\) 是 \(f(x)\) 的邻域。由连续性，\(f^{-1}(V)\) 是 \(x\) 的邻域，故存在开集 \(U_x \subseteq f^{-1}(V)\) 使 \(x \in U_x\)。于是 \(f^{-1}(V) = \bigcup_{x \in f^{-1}(V)} U_x\) 是开集。

\begin{enumerate}
\def\labelenumi{(\arabic{enumi})}
\setcounter{enumi}{1}
\item
  \(\Rightarrow\) (3)：\(f^{-1}(Y \setminus F) = X \setminus f^{-1}(F)\)，若 \(F\) 闭，则 \(Y \setminus F\) 开，故 \(X \setminus f^{-1}(F)\) 开，即 \(f^{-1}(F)\) 闭。
\item
  \(\Rightarrow\) (4)：\(f^{-1}(\overline{f(A)})\) 是包含 \(A\) 的闭集，故 \(\overline{A} \subseteq f^{-1}(\overline{f(A)})\)，即 \(f(\overline{A}) \subseteq \overline{f(A)}\)。
\item
  \(\Rightarrow\) (1)：设 \(x \in X\)，\(V\) 是 \(f(x)\) 的邻域。令 \(A = X \setminus f^{-1}(V)\)。若 \(f(x)\) 不是 \(f(A)\) 的聚点，则存在 \(f(x)\) 的开邻域 \(W \subseteq V\) 与 \(f(A)\) 不相交。于是 \(f^{-1}(W) \subseteq f^{-1}(V)\) 是 \(x\) 的邻域，故 \(f\) 在 \(x\) 处连续。假设 \(f(x) \in \overline{f(A)}\)，则 \(x \in f^{-1}(\overline{f(A)}) \supseteq \overline{A}\)，故 \(x\) 的每个邻域与 \(A\) 相交，这与 \(f^{-1}(V)\) 是 \(x\) 的邻域矛盾。因此 \(f^{-1}(V)\) 是 \(x\) 的邻域。
\item
  \(\Leftrightarrow\) (5)：由基的定义直接可得。∎
\end{enumerate}

\textbf{命题 57.6}（连续映射的复合）：若 \(f: X \to Y\) 和 \(g: Y \to Z\) 都是连续映射，则 \(g \circ f: X \to Z\) 也是连续映射。

\emph{证明}：对 \(Z\) 中任意开集 \(W\)，\((g \circ f)^{-1}(W) = f^{-1}(g^{-1}(W))\)。由 \(g\) 连续，\(g^{-1}(W)\) 是 \(Y\) 中开集；由 \(f\) 连续，\(f^{-1}(g^{-1}(W))\) 是 \(X\) 中开集。∎

\textbf{定义 57.10}（同胚）：设 \(X, Y\) 为拓扑空间。映射 \(f: X \to Y\) 称为\textbf{同胚}（homeomorphism），如果 \(f\) 是双射，且 \(f\) 和 \(f^{-1}\) 都是连续映射。若存在 \(X\) 到 \(Y\) 的同胚，则称 \(X\) 与 \(Y\) \textbf{同胚}，记作 \(X \cong Y\)（或 \(X \approx Y\)）。

同胚是拓扑空间范畴中的同构。同胚的空间具有完全相同的拓扑性质。

\textbf{定义 57.11}（开映射、闭映射）：映射 \(f: X \to Y\) 称为\textbf{开映射}（resp. \textbf{闭映射}），如果 \(f\) 将 \(X\) 中的开集（resp. 闭集）映为 \(Y\) 中的开集（resp. 闭集）。

\textbf{命题 57.7}：设 \(f: X \to Y\) 为连续双射，则 \(f\) 是同胚当且仅当 \(f\) 是开映射，也当且仅当 \(f\) 是闭映射。

\emph{证明}：\(f^{-1}\) 连续等价于对 \(X\) 中任意开集 \(U\)，\((f^{-1})^{-1}(U) = f(U)\) 是 \(Y\) 中开集，即 \(f\) 是开映射。闭映射情形类似。∎

\subsection{57.5 子空间拓扑}\label{ux5b50ux7a7aux95f4ux62d3ux6251}

\textbf{定义 57.12}（子空间拓扑）：设 \((X, \mathcal{T})\) 为拓扑空间，\(Y \subseteq X\)。定义

\[\mathcal{T}_Y = \{U \cap Y : U \in \mathcal{T}\}\]

则 \(\mathcal{T}_Y\) 是 \(Y\) 上的拓扑，称为\textbf{子空间拓扑}（或相对拓扑）。\((Y, \mathcal{T}_Y)\) 称为 \((X, \mathcal{T})\) 的\textbf{拓扑子空间}。

\textbf{命题 57.8}（子空间拓扑的性质）：设 \(Y\) 是 \(X\) 的子空间。 1. 若 \(\mathcal{B}\) 是 \(X\) 的拓扑的基，则 \(\{B \cap Y : B \in \mathcal{B}\}\) 是 \(Y\) 的拓扑的基。 2. \(Y\) 中的闭集恰为 \(F \cap Y\)，其中 \(F\) 是 \(X\) 中的闭集。 3. 包含映射 \(i: Y \hookrightarrow X\)（\(i(y) = y\)）是连续的。 4. 若 \(Z \subseteq Y \subseteq X\)，则 \(Z\) 作为 \(Y\) 的子空间的拓扑与作为 \(X\) 的子空间的拓扑相同。

\emph{证明}：(1)-(3) 由定义直接可得。(4) \(Z\) 中开集为 \(\{U \cap Z : U \in \mathcal{T}\} = \{(U \cap Y) \cap Z : U \in \mathcal{T}\}\)，恰为 \(Z\) 作为 \(Y\) 子空间的拓扑。∎

\textbf{定理 57.9}（子空间上的连续映射）：设 \(f: X \to Y\) 连续，\(A \subseteq X\)。则限制映射 \(f|_A: A \to Y\)（赋予 \(A\) 子空间拓扑）也是连续的。

\emph{证明}：对 \(Y\) 中开集 \(V\)，\((f|_A)^{-1}(V) = f^{-1}(V) \cap A\)，而 \(f^{-1}(V)\) 是 \(X\) 中开集，故 \(f^{-1}(V) \cap A\) 是 \(A\) 中开集。∎

\subsection{57.6 积空间}\label{ux79efux7a7aux95f4}

\textbf{定义 57.13}（积拓扑------有限积）：设 \(X_1, X_2, \ldots, X_n\) 为拓扑空间。在积集 \(X_1 \times X_2 \times \cdots \times X_n\) 上，以所有形如 \(U_1 \times U_2 \times \cdots \times U_n\)（\(U_i\) 是 \(X_i\) 中开集）的集合为基生成的拓扑，称为\textbf{积拓扑}。

\textbf{定义 57.14}（积拓扑------任意积）：设 \(\{X_\alpha\}_{\alpha \in A}\) 为一族拓扑空间，\(X = \prod_{\alpha \in A} X_\alpha\)。\(X\) 上的\textbf{积拓扑}是以

\[\mathcal{S} = \{\pi_\alpha^{-1}(U_\alpha) : \alpha \in A, U_\alpha \subseteq X_\alpha \text{ 为开集}\}\]

为子基生成的拓扑，其中 \(\pi_\alpha: X \to X_\alpha\) 是到第 \(\alpha\) 个坐标的投影映射。

等价地，积拓扑的基由形如 \(\prod_{\alpha \in A} U_\alpha\) 的集合组成，其中每个 \(U_\alpha\) 是 \(X_\alpha\) 的开集，且除有限多个 \(\alpha\) 外 \(U_\alpha = X_\alpha\)。

\textbf{命题 57.10}（积拓扑的泛性质）：设 \(X = \prod_{\alpha \in A} X_\alpha\) 赋予积拓扑。则： 1. 每个投影 \(\pi_\alpha: X \to X_\alpha\) 是连续开映射。 2. 映射 \(f: Y \to X\) 连续当且仅当对每个 \(\alpha \in A\)，\(\pi_\alpha \circ f: Y \to X_\alpha\) 连续。

\emph{证明}：(1) 对 \(X_\alpha\) 中开集 \(U_\alpha\)，\(\pi_\alpha^{-1}(U_\alpha)\) 属于子基，故为开集，因此 \(\pi_\alpha\) 连续。\(\pi_\alpha\) 是开映射因为基元素被映为开集。(2) 若 \(f\) 连续，则每个 \(\pi_\alpha \circ f\) 连续。反之，若每个 \(\pi_\alpha \circ f\) 连续，则对子基元素 \(\pi_\alpha^{-1}(U_\alpha)\)，\(f^{-1}(\pi_\alpha^{-1}(U_\alpha)) = (\pi_\alpha \circ f)^{-1}(U_\alpha)\) 是开集，由定理 57.5(5) 知 \(f\) 连续。∎

\subsection{57.7 商空间}\label{ux5546ux7a7aux95f4}

\textbf{定义 57.15}（商映射）：设 \(X, Y\) 为拓扑空间。满射 \(p: X \to Y\) 称为\textbf{商映射}，如果 \(V \subseteq Y\) 是开集当且仅当 \(p^{-1}(V)\) 是 \(X\) 中的开集。

\textbf{定义 57.16}（商拓扑）：设 \(X\) 为拓扑空间，\(\sim\) 是 \(X\) 上的等价关系，\(X/{\sim}\) 为等价类构成的集合，\(\pi: X \to X/{\sim}\) 为自然投影。\(X/{\sim}\) 上的\textbf{商拓扑}定义为使 \(\pi\) 成为商映射的拓扑，即 \(U \subseteq X/{\sim}\) 是开集当且仅当 \(\pi^{-1}(U)\) 是 \(X\) 中的开集。\((X/{\sim}, \text{商拓扑})\) 称为\textbf{商空间}。

\textbf{命题 57.11}（商拓扑的泛性质）：设 \(p: X \to Y\) 为商映射，\(g: Y \to Z\) 为映射。则 \(g\) 连续当且仅当 \(g \circ p: X \to Z\) 连续。

\emph{证明}：若 \(g\) 连续，则 \(g \circ p\) 连续。反之，若 \(g \circ p\) 连续，对 \(Z\) 中开集 \(W\)，\((g \circ p)^{-1}(W) = p^{-1}(g^{-1}(W))\) 是 \(X\) 中开集。由商拓扑定义，\(g^{-1}(W)\) 是 \(Y\) 中开集，故 \(g\) 连续。∎

\subsection{57.8 网与收敛}\label{ux7f51ux4e0eux6536ux655b}

度量空间中的序列收敛概念在一般拓扑空间中不够用，需要推广为网（net）的概念。

\textbf{定义 57.17}（有向集）：偏序集 \((D, \leq)\) 称为\textbf{有向集}，如果对任意 \(\alpha, \beta \in D\)，存在 \(\gamma \in D\) 使得 \(\alpha \leq \gamma\) 且 \(\beta \leq \gamma\)。

\textbf{定义 57.18}（网）：拓扑空间 \(X\) 中的一个\textbf{网}是一个映射 \(x: D \to X\)，其中 \(D\) 是有向集。通常记作 \((x_\alpha)_{\alpha \in D}\)。

\textbf{定义 57.19}（网的收敛）：称网 \((x_\alpha)_{\alpha \in D}\) \textbf{收敛}到 \(x \in X\)（记作 \(x_\alpha \to x\)），如果对 \(x\) 的任意邻域 \(U\)，存在 \(\alpha_0 \in D\) 使得对所有 \(\alpha \geq \alpha_0\)，\(x_\alpha \in U\)。

\textbf{定理 57.12}（网刻画闭包）：设 \(X\) 为拓扑空间，\(A \subseteq X\)。则 \(x \in \overline{A}\) 当且仅当存在 \(A\) 中的网收敛到 \(x\)。

\emph{证明}：(\(\Rightarrow\)) 设 \(x \in \overline{A}\)。取有向集 \(D = \mathcal{N}(x)\)（\(x\) 的邻域系），按反向包含排序：\(U \leq V \Leftrightarrow V \subseteq U\)。对每个 \(U \in D\)，由 \(x \in \overline{A}\) 知 \(U \cap A \neq \emptyset\)，选择 \(x_U \in U \cap A\)。则 \((x_U)_{U \in D}\) 是 \(A\) 中的网且 \(x_U \to x\)。

(\(\Leftarrow\)) 设 \((x_\alpha)\) 是 \(A\) 中的网，\(x_\alpha \to x\)。对 \(x\) 的任意邻域 \(U\)，存在 \(\alpha_0\) 使对所有 \(\alpha \geq \alpha_0\) 有 \(x_\alpha \in U\)，特别地 \(x_\alpha \in U \cap A\)，故 \(U \cap A \neq \emptyset\)。因此 \(x \in \overline{A}\)。∎

\textbf{定理 57.13}（网刻画连续性）：设 \(f: X \to Y\) 为映射。则 \(f\) 在 \(x \in X\) 处连续当且仅当对 \(X\) 中任意收敛到 \(x\) 的网 \((x_\alpha)\)，\(f(x_\alpha) \to f(x)\) 在 \(Y\) 中成立。

\emph{证明}：(\(\Rightarrow\)) 设 \(f\) 在 \(x\) 处连续，\((x_\alpha)\) 收敛到 \(x\)。对 \(f(x)\) 的任意邻域 \(V\)，\(f^{-1}(V)\) 是 \(x\) 的邻域，故存在 \(\alpha_0\) 使对所有 \(\alpha \geq \alpha_0\) 有 \(x_\alpha \in f^{-1}(V)\)，即 \(f(x_\alpha) \in V\)。故 \(f(x_\alpha) \to f(x)\)。

(\(\Leftarrow\)) 反设 \(f\) 在 \(x\) 处不连续。则存在 \(f(x)\) 的邻域 \(V\) 使 \(f^{-1}(V)\) 不是 \(x\) 的邻域。取有向集 \(D = \mathcal{N}(x)\)，对每个 \(U \in D\)，\(U \not\subseteq f^{-1}(V)\)，故存在 \(x_U \in U \setminus f^{-1}(V)\)。则 \((x_U)\) 是收敛到 \(x\) 的网，但 \(f(x_U) \notin V\)，故 \(f(x_U)\) 不收敛到 \(f(x)\)。矛盾。∎

\textbf{拓扑空间理论的补充主题}

拓扑空间理论是 20 世纪数学的奠基性成就之一，其核心思想------将连续性的概念从度量空间中解放出来，用开集公理体系刻画------不仅统一了分析学中的各种收敛概念，也为代数拓扑、微分拓扑和泛函分析提供了共同的语言。\textbf{滤子}（filter）和\textbf{网}（net）是处理一般拓扑空间中收敛性的两种等价方法，它们都是在序列收敛不足以刻画拓扑性质（如非第一可数空间中的闭包）时引入的。Bourbaki 学派（1930s-1940s）系统发展了滤子理论，而 Kelley（1950）和 Moore-Smith（1922）则推动了网的概念。两种方法各有优势：滤子适合处理紧致性（超滤子引理），网适合处理连续性和闭包。\textbf{拓扑基与子基}的概念是拓扑空间定义的核心------通过指定哪些集合是开集，我们实际上是在定义空间的''形状''而非''距离''。Alexandroff 和 Hopf 的经典著作《Topologie I》（1935）首次系统整理了拓扑空间理论，将点集拓扑建立为独立的数学分支。\textbf{仿紧性}（paracompactness）是 Dieudonne（1944）引入的比紧致性更灵活的概念：每个开覆盖有局部有限的开加细覆盖。仿紧性在微分几何中至关重要------它保证流形上单位分解的存在性，从而容许将局部构造全局化。每个可度量化空间是仿紧的（Stone 1948），每个局部紧 Hausdorff 空间且第二可数的是仿紧的。\textbf{Stone-Cech 紧化}（Stone 1937, Cech 1937）是 Tychonoff 空间到紧致 Hausdorff 空间的最大紧化，其构造基于单位区间 \([0,1]\) 的嵌入，在泛函分析（\(C^*\)-代数与 \(C(X)\) 的关系）和集合论（超滤子与基数不变量）中扮演核心角色。

\hrulefill

\hrulefill

\section{第213章：连通性与紧致性}\label{ux7b2c213ux7ae0ux8fdeux901aux6027ux4e0eux7d27ux81f4ux6027}

连通性和紧致性是拓扑空间理论中两个最根本的全局性质，它们分别定量刻画了空间的''不可分割性''和''有限性''。连通性的提出源于 19 世纪末 Bolzano、Cantor 和 Jordan 等人对实直线和平面子集的直观研究------一个空间是连通的，如果它不能被分成两个分离的非空部分。而紧致性则更为微妙：它起源于 Heine-Borel 定理（1895年）对 \(\mathbb{R}^n\) 中有界闭集的刻画，但真正深刻的推广由 Pavel Alexandroff 和 Pavel Urysohn 在 1920 年代完成------紧致空间被定义为''每个开覆盖必有有限子覆盖''，这一看似技术性的条件却直接蕴含了有界闭区间在连续映射下保持紧致、紧致空间上的连续函数必达最值的核心分析定理。在拓扑学体系中，第151章建立的拓扑空间公理为连通性与紧致性的代数拓扑学推广（基本群、同调群，见第154-177章）奠定了基础：紧致性和连通性决定了基本群和同调群的有限生成性质，而紧致曲面的分类（第154章之后的代数拓扑应用）完全依赖于这两个不变量的相互作用。

\subsection{58.1 连通性}\label{ux8fdeux901aux6027}

\textbf{定义 58.1}（连通空间）：拓扑空间 \(X\) 称为\textbf{连通的}，如果 \(X\) 不能表示为两个不相交的非空开集的并。等价地，\(X\) 的既开又闭的子集只有 \(\emptyset\) 和 \(X\)。

子集 \(A \subseteq X\) 称为\textbf{连通的}，如果 \(A\) 作为子空间是连通的。

\textbf{命题 58.1}（连续映射保持连通性）：设 \(f: X \to Y\) 为连续映射，\(X\) 连通。则 \(f(X)\) 连通。

\emph{证明}：若 \(f(X)\) 不连通，则存在 \(f(X)\) 的既开又闭的非空真子集 \(A\)。则 \(f^{-1}(A)\) 是 \(X\) 的既开又闭的非空真子集，与 \(X\) 连通矛盾。∎

\textbf{定理 58.2}（连通集的并）：设 \(\{A_\alpha\}_{\alpha \in A}\) 是 \(X\) 的一族连通子集，且 \(\bigcap_{\alpha \in A} A_\alpha \neq \emptyset\)。则 \(\bigcup_{\alpha \in A} A_\alpha\) 连通。

\emph{证明}：设 \(Y = \bigcup A_\alpha\)，\(U\) 是 \(Y\) 的既开又闭子集。取 \(x \in \bigcap A_\alpha\)，不妨设 \(x \in U\)。对每个 \(\alpha\)，\(U \cap A_\alpha\) 是 \(A_\alpha\) 的既开又闭子集且非空，由 \(A_\alpha\) 连通知 \(U \cap A_\alpha = A_\alpha\)，即 \(A_\alpha \subseteq U\)。故 \(U = Y\)。∎

\textbf{定理 58.3}（闭包的连通性）：若 \(A \subseteq X\) 连通，且 \(A \subseteq B \subseteq \overline{A}\)，则 \(B\) 连通。特别地，连通集的闭包连通。

\emph{证明}：设 \(U\) 是 \(B\) 的既开又闭子集。则 \(U \cap A\) 是 \(A\) 的既开又闭子集，由 \(A\) 连通知 \(U \cap A = \emptyset\) 或 \(U \cap A = A\)。若 \(U \cap A = \emptyset\)，则 \(A \subseteq B \setminus U\)，\(B \setminus U\) 是闭集，故 \(\overline{A} \subseteq B \setminus U\)，从而 \(U = \emptyset\)。若 \(U \cap A = A\)，则 \(A \subseteq U\)，\(U\) 是闭集，故 \(\overline{A} \subseteq U\)（在 \(B\) 中），从而 \(U = B\)。∎

\textbf{定义 58.2}（连通分支）：设 \(X\) 为拓扑空间。定义关系 \(x \sim y\) 当且仅当存在 \(X\) 的连通子集同时包含 \(x\) 和 \(y\)。这是一个等价关系，其等价类称为 \(X\) 的\textbf{连通分支}（或\textbf{连通分量}）。

\textbf{命题 58.4}：每个连通分支是极大的连通子集，且是闭集。

\emph{证明}：由定理 58.2 知每个等价类中的点确实属于某个连通子集，且该等价类本身是连通集（取包含其中任意两点的所有连通子集的并）。极大性由定义直接可得。由定理 58.3，连通分支的闭包也连通，故连通分支等于其闭包，即为闭集。∎

\textbf{定义 58.3}（道路连通）：拓扑空间 \(X\) 称为\textbf{道路连通的}，如果对任意 \(x, y \in X\)，存在连续映射 \(\gamma: [0, 1] \to X\) 使得 \(\gamma(0) = x\)，\(\gamma(1) = y\)。\(\gamma\) 称为连接 \(x\) 和 \(y\) 的\textbf{道路}。

\textbf{命题 58.5}：道路连通空间必连通。

\emph{证明}：设 \(X\) 道路连通，假设 \(X = U \cup V\)，\(U, V\) 是不相交的非空开集。取 \(x \in U, y \in V\)，存在道路 \(\gamma: [0, 1] \to X\) 连接 \(x\) 和 \(y\)。则 \([0, 1] = \gamma^{-1}(U) \cup \gamma^{-1}(V)\) 将 \([0, 1]\) 分为两个不相交的非空开集，与 \([0, 1]\) 连通矛盾。∎

\textbf{定义 58.4}（局部连通、局部道路连通）：\(X\) 称为\textbf{局部连通的}（resp. \textbf{局部道路连通的}），如果对每个 \(x \in X\) 和 \(x\) 的任意邻域 \(U\)，存在 \(x\) 的连通（resp. 道路连通）邻域 \(V \subseteq U\)。

\textbf{命题 58.6}：在局部道路连通空间中，连通分支与道路连通分支一致，且都是既开又闭的。

\emph{证明}：在局部道路连通空间中，道路连通分支是开集。每个连通分支是道路连通分支的不交并。由于连通分支不能分解为两个不交开集的并，故每个连通分支恰为一个道路连通分支。∎

\subsection{58.2 紧致性}\label{ux7d27ux81f4ux6027}

\textbf{定义 58.5}（开覆盖）：设 \(X\) 为拓扑空间。一族开集 \(\{U_\alpha\}_{\alpha \in A}\) 称为 \(X\) 的\textbf{开覆盖}，如果 \(\bigcup_{\alpha \in A} U_\alpha = X\)。若子族也覆盖 \(X\)，则称为\textbf{子覆盖}。

\textbf{定义 58.6}（紧致空间）：拓扑空间 \(X\) 称为\textbf{紧致的}（compact），如果 \(X\) 的任意开覆盖都有有限子覆盖。

子集 \(K \subseteq X\) 称为\textbf{紧致的}，如果 \(K\) 作为子空间是紧致的。

\textbf{命题 58.7}：\(K \subseteq X\) 紧致当且仅当对 \(X\) 中任意一族覆盖 \(K\) 的开集 \(\{U_\alpha\}\)（即 \(K \subseteq \bigcup U_\alpha\)），存在有限子族也覆盖 \(K\)。

\emph{证明}：由子空间拓扑的定义，\(K\) 中开集形如 \(U \cap K\)，\(U\) 是 \(X\) 中开集。开覆盖 \(\{U_\alpha \cap K\}\) 对应 \(X\) 中一族覆盖 \(K\) 的开集。∎

\textbf{定理 58.8}（紧致集的基本性质）： 1. 紧致集的连续像是紧致集。 2. 紧致空间的闭子集紧致。 3. Hausdorff 空间中的紧致子集是闭集。 4. 紧致空间的连续像若在 Hausdorff 空间中，则是闭映射。

\emph{证明}：(1) 设 \(f: X \to Y\) 连续，\(X\) 紧致。对 \(f(X)\) 的任意开覆盖 \(\{V_\alpha\}\)，\(\{f^{-1}(V_\alpha)\}\) 是 \(X\) 的开覆盖，有有限子覆盖 \(\{f^{-1}(V_{\alpha_1}), \ldots, f^{-1}(V_{\alpha_n})\}\)，则 \(\{V_{\alpha_1}, \ldots, V_{\alpha_n}\}\) 覆盖 \(f(X)\)。

\begin{enumerate}
\def\labelenumi{(\arabic{enumi})}
\setcounter{enumi}{1}
\item
  设 \(X\) 紧致，\(F \subseteq X\) 闭。对 \(F\) 的任意开覆盖 \(\{U_\alpha \cap F\}\)（\(U_\alpha\) 是 \(X\) 中开集），\(\{U_\alpha\} \cup \{X \setminus F\}\) 是 \(X\) 的开覆盖，有有限子覆盖，去掉 \(X \setminus F\) 后得到 \(F\) 的有限子覆盖。
\item
  设 \(X\) 是 Hausdorff 空间，\(K \subseteq X\) 紧致。对 \(x \notin K\)，对每个 \(y \in K\)，存在不相交的开集 \(U_y \ni x\)，\(V_y \ni y\)。\(\{V_y\}_{y \in K}\) 覆盖 \(K\)，有有限子覆盖 \(\{V_{y_1}, \ldots, V_{y_n}\}\)。令 \(U = \bigcap_{i=1}^n U_{y_i}\)，则 \(U\) 是 \(x\) 的邻域且 \(U \cap K = \emptyset\)。故 \(x \notin \overline{K}\)，即 \(K\) 闭。
\item
  由 (1) 和 (3)，紧致集在连续映射下的像是紧致集，在 Hausdorff 空间中紧致集闭，故闭集（紧致空间的闭子集紧致）被映为紧致集（闭集）。∎
\end{enumerate}

\textbf{定理 58.8b}（紧致空间上连续函数达到最值）：设 \(X\) 为紧致空间，\(f: X \to \mathbb{R}\) 为连续函数（\(\mathbb{R}\) 赋予通常拓扑）。则 \(f\) 在 \(X\) 上达到最大值和最小值。

\emph{证明}：由定理 58.8(1)，\(f(X)\) 是 \(\mathbb{R}\) 中的紧致集。由 Heine-Borel 定理（定理 58.10），\(\mathbb{R}\) 中的紧致子集有界且闭。设 \(M = \sup f(X)\)，\(m = \inf f(X)\)。由有界性 \(M, m < \infty\)。由于 \(f(X)\) 是闭集，\(M = \sup f(X) \in \overline{f(X)} = f(X)\)。故存在 \(x_{\max} \in X\) 使得 \(f(x_{\max}) = M\)。同理，\(m \in f(X)\)，存在 \(x_{\min} \in X\) 使得 \(f(x_{\min}) = m\)。因此 \(f\) 在紧致空间上达到最大值和最小值。\(\blacksquare\)

\textbf{定理 58.9}（紧致空间到 Hausdorff 空间的连续双射）：若 \(f: X \to Y\) 是连续双射，\(X\) 紧致，\(Y\) 是 Hausdorff 空间，则 \(f\) 是同胚。

\emph{证明}：\(f\) 是连续双射，故其逆映射 \(f^{-1}\) 存在。需证 \(f^{-1}\) 连续，即证 \(f\) 是开映射（或等价地，闭映射）。对 \(X\) 中任意闭集 \(F\)，由定理 58.8(2)，\(X\) 紧致蕴含 \(F\) 紧致。由定理 58.8(1)，连续映射保持紧致性，故 \(f(F)\) 是 \(Y\) 中紧致子集。\(Y\) 是 Hausdorff 空间，由定理 58.8(3)，紧致子集是闭集，故 \(f(F)\) 是 \(Y\) 中闭集。因此 \(f\) 是闭映射，从而 \(f^{-1}\) 连续，\(f\) 是同胚。\(\blacksquare\)

\textbf{定理 58.10}（Heine-Borel 定理，\(\mathbb{R}^n\) 情形）：\(\mathbb{R}^n\) 的子集紧致当且仅当它是有界闭集。

\emph{证明}：(\(\Leftarrow\)) 设 \(A \subseteq \mathbb{R}^n\) 有界闭集。则存在 \(M > 0\) 使得 \(A \subseteq [-M, M]^n\)。首先证明 \([0, 1]\) 是紧致的：设 \(\{U_\alpha\}\) 为 \([0,1]\) 的开覆盖。令 \(S = \{s \in [0,1] : [0,s] \text{ 可被有限子覆盖}\}\)。\(0 \in S\)（\([0,0] = \{0\}\) 被某个 \(U_\alpha\) 覆盖）。令 \(s_0 = \sup S\)。若 \(s_0 < 1\)，则存在 \(U_\alpha\) 包含 \(s_0\) 及某个 \(\epsilon\) 邻域，则 \(s_0 + \epsilon/2 \in S\)，矛盾。故 \(s_0 = 1\)，\([0,1]\) 紧致。由 Tychonoff 定理（或直接归纳），\([-M, M]^n\) 紧致。\(A\) 作为紧致空间的闭子集，紧致。

(\(\Rightarrow\)) 设 \(A \subseteq \mathbb{R}^n\) 紧致。则 \(A\) 在 Hausdorff 空间 \(\mathbb{R}^n\) 中是闭的（紧致集的闭包性质）。\(A\) 有界：否则存在序列 \(x_k \in A\) 使得 \(\|x_k\| \to \infty\)。\(A\) 的紧致性（序列紧致）要求存在收敛子列，但 \(\|x_k\| \to \infty\) 中无收敛子列，矛盾。故 \(A\) 有界。\(\blacksquare\)

\textbf{定理 58.10a}（Lebesgue 数引理）：设 \((X, d)\) 为紧致度量空间，\(\{U_\alpha\}_{\alpha \in A}\) 是 \(X\) 的任意开覆盖。则存在 \(\delta > 0\)（称为该覆盖的 \textbf{Lebesgue 数}），使得对任意 \(A \subseteq X\) 满足 \(\operatorname{diam}(A) < \delta\)，存在某个 \(\alpha \in A\) 使得 \(A \subseteq U_\alpha\)。

\emph{证明}：对 \(\{U_\alpha\}\) 由紧致性取有限子覆盖 \(\{U_1, \ldots, U_n\}\)。对每个 \(i\)，定义 \(f_i(x) = d(x, X \setminus U_i)\)（若 \(U_i = X\) 则 \(f_i(x) \equiv 1\)）。\(f_i\) 连续（距离函数到闭集的连续性）。令 \(f(x) = \frac{1}{n}\max_i f_i(x)\)。由于 \(X = \bigcup U_i\)，对每个 \(x\) 存在 \(i\) 使得 \(x \in U_i\)，故 \(f_i(x) > 0\)，从而 \(f(x) > 0\)。\(f\) 是紧致集 \(X\) 上的正连续函数，故达到正的最小值 \(\delta = \min_{x \in X} f(x) > 0\)。对任意 \(A \subseteq X\) 满足 \(\operatorname{diam}(A) < \delta\)，取 \(x_0 \in A\)，则 \(A \subseteq B(x_0, \delta)\)。由 \(f(x_0) \geq \delta\)，存在 \(i\) 使得 \(f_i(x_0) \geq \delta\)，即 \(d(x_0, X \setminus U_i) \geq \delta\)，故 \(B(x_0, \delta) \subseteq U_i\)，从而 \(A \subseteq U_i\)。\(\blacksquare\)

\subsection{58.3 Tychonoff 定理}\label{tychonoff-ux5b9aux7406}

\textbf{定理 58.11a}（Alexander 子基定理）：设 \(X\) 为拓扑空间，\(\mathcal{S}\) 是 \(X\) 的拓扑的一个子基。若由 \(\mathcal{S}\) 中元素构成的任意开覆盖都有有限子覆盖，则 \(X\) 是紧致的。

\emph{证明}：用反证法。假设 \(X\) 不是紧致的，即存在 \(X\) 的开覆盖 \(\mathcal{U}\)（由 \(\mathcal{S}\) 的有限交生成的基的元素组成）没有有限子覆盖。设 \(\mathcal{U}\) 是极大反例（由 Zorn 引理保证存在）。对每个 \(S \in \mathcal{S}\)，若存在 \(U_1, \ldots, U_n \in \mathcal{U}\) 使得 \(S \cap (\bigcap_{i=1}^n U_i) \subseteq \bigcup_{i=1}^n U_i \cup S\) 不被有限子覆盖，则 \(S \in \mathcal{U}\)（否则 \(\mathcal{U}\) 不是极大）。任取 \(x \in X \setminus \bigcup \mathcal{U}\)，存在 \(S_1, \ldots, S_k \in \mathcal{S}\) 使 \(S_1 \cap \cdots \cap S_k\) 是含 \(x\) 的基本开集且在某个开覆盖元素内。由 \(x \notin \bigcup \mathcal{U}\) 知每个 \(S_i \notin \mathcal{U}\)，由极大性导出矛盾。故由 \(\mathcal{S}\) 的任意覆盖有有限子覆盖可推出 \(X\) 紧致。\(\blacksquare\)

\textbf{定理 58.11}（Tychonoff 定理）：任意一族紧致空间的积空间（赋予积拓扑）是紧致的。

\textbf{证明}：设 \(X = \prod_{\alpha \in A} X_\alpha\)，每个 \(X_\alpha\) 紧致。使用 Alexander 子基定理（定理 58.11a）：只需验证子基 \(\mathcal{S} = \{\pi_\alpha^{-1}(U_\alpha) : \alpha \in A, U_\alpha \subseteq X_\alpha \text{ 开}\}\) 的任意覆盖有有限子覆盖。设 \(\mathcal{C} \subseteq \mathcal{S}\) 覆盖 \(X\)。对每个 \(\alpha \in A\)，令 \(\mathcal{C}_\alpha = \{U \subseteq X_\alpha \text{ 开} : \pi_\alpha^{-1}(U) \in \mathcal{C}\}\)。若存在 \(\alpha\) 使 \(\mathcal{C}_\alpha\) 覆盖 \(X_\alpha\)，则由 \(X_\alpha\) 紧致知存在有限子覆盖 \(\{U_1,\ldots,U_n\} \subseteq \mathcal{C}_\alpha\)，则 \(\{\pi_\alpha^{-1}(U_i)\}_{i=1}^n\) 覆盖 \(X\)。否则，对每个 \(\alpha\)，存在 \(x_\alpha \in X_\alpha \setminus \bigcup \mathcal{C}_\alpha\)。由选择公理得 \((x_\alpha)_{\alpha \in A} \in X\)。但对任意 \(\pi_\alpha^{-1}(U) \in \mathcal{C}\)，\(x_\alpha \notin U\)，故 \((x_\alpha) \notin \pi_\alpha^{-1}(U)\)，与 \(\mathcal{C}\) 覆盖 \(X\) 矛盾。因此 \(\mathcal{C}\) 必有有限子覆盖，由 Alexander 子基定理知 \(X\) 紧致。\(\blacksquare\)

\textbf{推论 58.12}（Tychonoff 定理，有限积情形）：有限个紧致空间的积紧致。

\textbf{证明（有限积的归纳证明，无需选择公理）}：对积的个数作归纳。\(n=1\) 平凡。设 \(X = X_1 \times \cdots \times X_n\) 紧致，\(X_{n+1}\) 紧致，需证 \(X \times X_{n+1}\) 紧致。设 \(\mathcal{U}\) 为 \(X \times X_{n+1}\) 的任意开覆盖。对每个 \(y \in X_{n+1}\)，截面 \(X \times \{y\}\) 同胚于紧致空间 \(X\)，故可被 \(\mathcal{U}\) 中的有限个开集覆盖，其并为含 \(X \times \{y\}\) 的开集。由管状引理（Tube Lemma：紧致空间 \(Y\) 中，若 \(W\) 是 \(X \times Y\) 中包含 \(X \times \{y_0\}\) 的开集，则存在 \(y_0\) 的邻域 \(V\) 使 \(X \times V \subseteq W\)），该开集包含某个管状邻域 \(X \times V_y\)（\(V_y\) 是 \(y\) 在 \(X_{n+1}\) 中的开邻域）。\(\{V_y\}_{y \in X_{n+1}}\) 覆盖紧致空间 \(X_{n+1}\)，有有限子覆盖 \(\{V_{y_1}, \ldots, V_{y_m}\}\)。对应地，\(\{X \times V_{y_i}\}_{i=1}^m\) 覆盖 \(X \times X_{n+1}\)，且每个 \(X \times V_{y_i}\) 被 \(\mathcal{U}\) 的有限子族覆盖，故 \(\mathcal{U}\) 有有限子覆盖。\(\blacksquare\)

\subsection{58.4 局部紧致与紧化}\label{ux5c40ux90e8ux7d27ux81f4ux4e0eux7d27ux5316}

\textbf{定义 58.7}（局部紧致）：拓扑空间 \(X\) 称为\textbf{局部紧致的}，如果对每个 \(x \in X\)，存在 \(x\) 的邻域 \(U\) 使得 \(\overline{U}\) 是紧致的。

\textbf{定理 58.13}（单点紧化 / Alexandroff 紧化）：设 \(X\) 是局部紧致的 Hausdorff 空间。则存在紧致 Hausdorff 空间 \(X^*\) 使得 \(X\) 同胚于 \(X^*\) 的稠密开子集，且 \(X^* \setminus X\) 是单点集。\(X^*\) 在同胚意义下唯一。

\emph{证明}（构造与验证）：令 \(X^* = X \cup \{\infty\}\)，其中 \(\infty \notin X\)。定义 \(X^*\) 上的拓扑 \(\mathcal{T}^*\)：\(U \subseteq X^*\) 是开集当且仅当以下两者之一成立： - \(\infty \notin U\) 且 \(U\) 是 \(X\) 中开集； - \(\infty \in U\) 且 \(X^* \setminus U\) 是 \(X\) 中的紧致闭集。

验证 \(\mathcal{T}^*\) 是拓扑：(a) \(\emptyset \in \mathcal{T}^*\)（第一种情况），\(X^* \in \mathcal{T}^*\)（\(X^* \setminus X^* = \emptyset\) 紧致）。任意并和有限交的验证分情况讨论即可，关键是紧致闭集的有限并仍是紧致闭集。\(X\) 作为 \(X^*\) 的子空间：\(X\) 在 \(X^*\) 中的子空间拓扑恰好是 \(X\) 的原拓扑（因为 \(\infty\) 不属于 \(X\) 中的开集）。\(X^*\) 紧致：设 \(\mathcal{U}\) 为 \(X^*\) 的开覆盖。取 \(U_\infty \in \mathcal{U}\) 包含 \(\infty\)，则 \(X^* \setminus U_\infty\) 是 \(X\) 中紧致集，可被 \(\mathcal{U}\) 中有限个开集覆盖，加上 \(U_\infty\) 即得有限子覆盖。Hausdorff 性：\(X\) 中两点可用 \(X\) 的 Hausdorff 性质分离；\(x \in X\) 与 \(\infty\) 的分离利用 \(X\) 的局部紧致性（取 \(x\) 的紧致闭包邻域，其补集即 \(\infty\) 的邻域）。唯一性：若 \(Y\) 是另一个单点紧化，则恒等映射 \(X \to X\) 可唯一延拓为同胚 \(X^* \to Y\)。\(\blacksquare\)

\subsection{58.5 Urysohn 引理与 Tietze 扩张定理}\label{urysohn-ux5f15ux7406ux4e0e-tietze-ux6269ux5f20ux5b9aux7406}

\textbf{定理 58.14}（Urysohn 引理）：设 \(X\) 为正规空间（见第 59 章），\(A, B\) 是 \(X\) 中不相交的闭集。则存在连续函数 \(f: X \to [0, 1]\) 使得 \(f(A) = \{0\}\)，\(f(B) = \{1\}\)。

\emph{证明}：令 \(U_1 = X \setminus B\)，则 \(A \subseteq U_1\)。由正规性，递归构造开集 \(U_r\)（\(r \in \mathbb{Q} \cap [0, 1]\)）满足对 \(p < q\) 有 \(\overline{U}_p \subseteq U_q\)。具体地，将 \(\mathbb{Q} \cap [0, 1]\) 排为 \(r_1 = 0, r_2 = 1, r_3, r_4, \ldots\)。归纳假设已对 \(r_1, \ldots, r_n\) 构造了 \(U_{r_i}\)。设 \(r_k, r_\ell\) 是 \(r_{n+1}\) 在已构造的 \(r_i\) 中的直接前驱和后继，则 \(U_{r_k} \subseteq \overline{U}_{r_k} \subseteq U_{r_\ell}\)，由正规性可插入 \(U_{r_{n+1}}\)。

定义 \(f(x) = \inf\{r : x \in U_r\}\)（约定 \(U_r = \emptyset\) 当 \(r < 0\)，\(U_r = X\) 当 \(r > 1\)）。则 \(f\) 连续：\(\{x : f(x) < t\} = \bigcup_{r < t} U_r\) 是开集，\(\{x : f(x) > t\} = \bigcup_{r > t} (X \setminus \overline{U}_r)\) 是开集。显然 \(f(A) = \{0\}\)，\(f(B) = \{1\}\)。∎

\textbf{定理 58.15}（Tietze 扩张定理）：设 \(X\) 为正规空间，\(A \subseteq X\) 为闭集。则 \(A\) 上的任意连续函数 \(f: A \to [a, b]\)（或 \(f: A \to \mathbb{R}\)）可连续扩张到 \(X\) 上。

\emph{证明}：不妨设 \(f: A \to [-1, 1]\)（否则通过仿射变换缩放）。令 \(A_1 = f^{-1}([-1, -1/3])\)，\(B_1 = f^{-1}([1/3, 1])\)。\(A_1\) 和 \(B_1\) 是 \(A\) 中不相交的闭集，因 \(A\) 是 \(X\) 中闭集，故 \(A_1\) 和 \(B_1\) 也是 \(X\) 中闭集。由 Urysohn 引理（定理 58.14），存在连续函数 \(g_1: X \to [-1/3, 1/3]\) 使得 \(g_1(A_1) = \{-1/3\}\)，\(g_1(B_1) = \{1/3\}\)。在 \(A\) 上，\(|f(x) - g_1(x)| \leq 2/3\)。令 \(f_1 = f - g_1|_A\)，则 \(f_1: A \to [-2/3, 2/3]\)。归纳地，对 \(n \geq 1\)，设 \(f_n: A \to [-(2/3)^n, (2/3)^n]\)，重复上述构造得到 \(g_{n+1}: X \to [-(2/3)^n/3, (2/3)^n/3]\) 使得 \(|f_n - g_{n+1}| \leq (2/3)^{n+1}\) 在 \(A\) 上。令 \(g(x) = \sum_{n=1}^\infty g_n(x)\)。由于 \(|g_n(x)| \leq (2/3)^{n-1}/3\)（\(n \geq 2\)），级数在 \(X\) 上一致收敛，故 \(g\) 连续。在 \(A\) 上，\(|f(x) - \sum_{i=1}^n g_i(x)| \leq (2/3)^n \to 0\)，故 \(g|_A = f\)。对 \(f: A \to \mathbb{R}\) 的情形，通过同胚 \(\mathbb{R} \cong (-1, 1)\) 化归为有界情形。\(\blacksquare\)

\textbf{分离性公理与度量化理论的深化}

分离性公理是拓扑空间理论中最精细的分类体系，它刻画了拓扑空间中点与点、点与闭集之间的''分离程度''。从 \(T_0\)（Kolmogorov）到 \(T_6\)（完全正规）的层级结构不仅反映了空间的''可分离性''，也决定了空间上可构造的连续函数的丰富程度。\textbf{Urysohn 引理}（1925）是分离性公理理论中最深刻的定理之一：在正规空间 \(X\) 中，任意两个不相交的闭集 \(A\) 和 \(B\) 可被一个连续函数 \(f: X \to [0,1]\) 完全分离，即 \(f(A) = \{0\}\) 且 \(f(B) = \{1\}\)。其证明依赖于将 \([0,1]\) 中的有理数对应到开集族 \(\{U_r\}_{r \in \mathbb{Q} \cap [0,1]}\) 上，满足 \(p < q \Rightarrow \overline{U}_p \subseteq U_q\) 的插入性质------这一构造体现了正规性公理的本质力量。\textbf{Tietze 扩张定理}（1915）是 Urysohn 引理的最重要推论：正规空间 \(X\) 的闭子集 \(A\) 上的任意连续函数可连续扩张到整个 \(X\)。在泛函分析中，Tietze 扩张定理是 Hahn-Banach 定理的拓扑对应------前者扩张连续函数，后者扩张线性泛函。\textbf{度量化定理}（metrization theorems）是点集拓扑学的皇冠：Urysohn 度量化定理（1925）断言第二可数的正则空间可度量化；Nagata-Smirnov 度量化定理（1950-1951）则给出了最一般的可度量化条件------正则空间具有 \(\sigma\)-局部有限基。这两条定理完整地回答了''何时拓扑空间可赋予度量''这一基本问题。\textbf{Bing-Nagata-Smirnov 度量化定理}（Bing 1951）给出了另一个等价条件：正则空间具有 \(\sigma\)-离散基。可度量化空间在拓扑学中的特殊地位源于其优良性质------第一可数、仿紧、完全正规------使其成为分析学中各种极限和收敛结构的最自然框架。

\hrulefill

\hrulefill

\section{第214章：分离公理与度量化}\label{ux7b2c214ux7ae0ux5206ux79bbux516cux7406ux4e0eux5ea6ux91cfux5316}

分离公理是拓扑空间分类的精细尺度------从最弱的 \(T_0\)（Kolmogorov 空间，不同点有不同的邻域族）到 \(T_4\)（正规空间，不相交闭集有不相交开邻域），构成了一条逐步增强的''拓扑正则性''阶梯。这一理论在 20 世纪上半叶由 Alexandroff、Urysohn、Tietze 和 Tychonoff 等人系统发展。分离公理的深层意义在于：它们精确地量化了一个拓扑空间''距离可度量空间有多远''。特别是 Urysohn 度量化定理------第二可数的正则 \(T_1\) 空间必可度量化------以及 Nagata-Smirnov 度量化定理，给出了拓扑空间等价于度量空间（在拓扑意义上）的完整刻画。本章与前一章（连通性与紧致性，第152章）和后一章（基本群，第154章）的关系是双向的：分离公理和紧致性共同决定了 \(T_2\) 紧致空间的正规性和度量化的可能性，而基本群的构造（基于闭路的同伦等价）在这些良好性质的空间上才能顺利展开。同时，紧开拓扑和商拓扑的分离性质直接影响了函数空间的可用性和代数拓扑工具的有效性。

\subsection{59.1 分离公理}\label{ux5206ux79bbux516cux7406}

\textbf{定义 59.1}（分离公理 \(T_0\)--\(T_4\)）：设 \(X\) 为拓扑空间。

\begin{itemize}
\tightlist
\item
  \(T_0\)（Kolmogorov）：对任意两个不同的点 \(x, y \in X\)，存在开集包含其中一点但不包含另一点。
\item
  \(T_1\)（Fréchet）：对任意两个不同的点 \(x, y \in X\)，存在开集 \(U\) 使得 \(x \in U\) 但 \(y \notin U\)。
\item
  \(T_2\)（Hausdorff）：对任意两个不同的点 \(x, y \in X\)，存在不相交的开集 \(U, V\) 使得 \(x \in U, y \in V\)。
\item
  \(T_3\)（正则）：\(X\) 是 \(T_1\) 的，且对任意 \(x \in X\) 和不包含 \(x\) 的闭集 \(F\)，存在不相交的开集 \(U, V\) 使得 \(x \in U, F \subseteq V\)。
\item
  \(T_{3.5}\)（完全正则 / Tychonoff）：\(X\) 是 \(T_1\) 的，且对任意 \(x \in X\) 和不包含 \(x\) 的闭集 \(F\)，存在连续函数 \(f: X \to [0, 1]\) 使得 \(f(x) = 0\)，\(f(F) = \{1\}\)。
\item
  \(T_4\)（正规）：\(X\) 是 \(T_1\) 的，且对任意不相交的闭集 \(A, B\)，存在不相交的开集 \(U, V\) 使得 \(A \subseteq U, B \subseteq V\)。
\end{itemize}

\textbf{命题 59.1}（分离公理的关系）： \[T_4 \Rightarrow T_{3.5} \Rightarrow T_3 \Rightarrow T_2 \Rightarrow T_1 \Rightarrow T_0\]

\emph{证明}：\(T_4 \Rightarrow T_{3.5}\) 是 Urysohn 引理（定理 58.14）。\(T_{3.5} \Rightarrow T_3\)：取 \(U = f^{-1}([0, 1/2))\)，\(V = f^{-1}((1/2, 1])\)。其余蕴含关系直接由定义可得。∎

\textbf{命题 59.2}（\(T_1\) 的等价刻画）：\(X\) 是 \(T_1\) 空间当且仅当每个单点集是闭集。

\emph{证明}：(\(\Rightarrow\)) 对 \(x \in X\)，对每个 \(y \neq x\)，存在开集 \(U_y \ni y\) 且 \(x \notin U_y\)。则 \(X \setminus \{x\} = \bigcup_{y \neq x} U_y\) 是开集，故 \(\{x\}\) 是闭集。(\(\Leftarrow\)) 对 \(x \neq y\)，取 \(U = X \setminus \{y\}\)，则 \(x \in U\) 且 \(y \notin U\)。∎

\textbf{命题 59.3}（\(T_2\) 的等价刻画）：\(X\) 是 Hausdorff 空间当且仅当对角线 \(\Delta = \{(x, x) : x \in X\}\) 在 \(X \times X\)（积拓扑）中是闭集。

\emph{证明}：(\(\Rightarrow\)) 对 \((x, y) \notin \Delta\)（即 \(x \neq y\)），存在不相交的开集 \(U \ni x, V \ni y\)，则 \(U \times V\) 是 \((x, y)\) 的邻域且与 \(\Delta\) 不相交。故 \(\Delta\) 闭。(\(\Leftarrow\)) 类似。∎

\textbf{命题 59.4}：Hausdorff 空间中序列（或网）的极限若存在则唯一。

\emph{证明}：设 \(x_\alpha \to x\) 且 \(x_\alpha \to y\)，\(x \neq y\)。取不相交的开集 \(U \ni x, V \ni y\)。由收敛定义，存在 \(\alpha_1\) 使 \(\alpha \geq \alpha_1\) 时 \(x_\alpha \in U\)，存在 \(\alpha_2\) 使 \(\alpha \geq \alpha_2\) 时 \(x_\alpha \in V\)。取 \(\alpha \geq \alpha_1, \alpha_2\)，则 \(x_\alpha \in U \cap V = \emptyset\)，矛盾。∎

\textbf{命题 59.5}：紧致 Hausdorff 空间是正规的（\(T_4\)）。

\emph{证明}：先证紧致 Hausdorff 空间是正则的：设 \(x \in X\)，\(F\) 是不含 \(x\) 的闭集。对每个 \(y \in F\)，由 \(T_2\) 存在不相交开集 \(U_y \ni x, V_y \ni y\)。\(\{V_y\}\) 覆盖紧致集 \(F\)，有有限子覆盖 \(\{V_{y_1}, \ldots, V_{y_n}\}\)。令 \(U = \bigcap U_{y_i}\)，\(V = \bigcup V_{y_i}\)，则 \(U \cap V = \emptyset\)。

再证正规性：设 \(A, B\) 是不相交的闭集。对每个 \(a \in A\)，由正则性存在不相交开集 \(U_a \ni a, V_a \supseteq B\)。\(\{U_a\}\) 覆盖紧致集 \(A\)，有有限子覆盖，类似可得分离 \(A\) 和 \(B\) 的开集。∎

\textbf{定理 59.5b}（Hausdorff 空间的积）：设 \(\{X_\alpha\}_{\alpha \in A}\) 是一族 Hausdorff 空间，则积空间 \(\prod_{\alpha \in A} X_\alpha\)（赋予积拓扑）也是 Hausdorff 空间。

\emph{证明}：设 \(x = (x_\alpha), y = (y_\alpha)\) 是积空间中两个不同的点，则存在 \(\alpha_0 \in A\) 使得 \(x_{\alpha_0} \neq y_{\alpha_0}\)。由于 \(X_{\alpha_0}\) 是 Hausdorff 空间，存在不相交的开集 \(U_{\alpha_0}, V_{\alpha_0} \subseteq X_{\alpha_0}\) 使得 \(x_{\alpha_0} \in U_{\alpha_0}\)，\(y_{\alpha_0} \in V_{\alpha_0}\)。令 \(U = \pi_{\alpha_0}^{-1}(U_{\alpha_0})\)，\(V = \pi_{\alpha_0}^{-1}(V_{\alpha_0})\)，则 \(U, V\) 是积拓扑中的不相交开集，\(x \in U\)，\(y \in V\)。故积空间是 Hausdorff 的。A fortiori，任意积拓扑下的 Hausdorff 空间的积是 Hausdorff 的。\(\blacksquare\)

\textbf{定理 59.5c}（第二可数正则空间是正规的）：设 \(X\) 是第二可数的正则空间。则 \(X\) 是正规的（\(T_4\)）。

\emph{证明}：设 \(A, B\) 是 \(X\) 中不相交的闭集。\(X\) 具有可数基 \(\mathcal{B} = \{B_n\}_{n \in \mathbb{N}}\)。对每个 \(a \in A\)，由正则性存在基元素 \(B_{n_a} \in \mathcal{B}\) 使得 \(a \in B_{n_a} \subseteq \overline{B}_{n_a} \subseteq X \setminus B\)。由于 \(\mathcal{B}\) 可数，集合 \(\{B_n \in \mathcal{B} : \overline{B}_n \cap B = \emptyset\}\) 可数，记作 \(\{U_1, U_2, \ldots\}\)。它覆盖 \(A\)。同理，存在覆盖 \(B\) 的可数开集族 \(\{V_1, V_2, \ldots\}\) 使得 \(\overline{V}_i \cap A = \emptyset\) 对所有 \(i\)。

定义 \(U_1' = U_1\)，\(U_n' = U_n \setminus \bigcup_{i=1}^{n-1} \overline{V}_i\)；\(V_1' = V_1\)，\(V_n' = V_n \setminus \bigcup_{i=1}^n \overline{U}_i\)。令 \(U = \bigcup_{n=1}^\infty U_n'\)，\(V = \bigcup_{n=1}^\infty V_n'\)。则 \(U, V\) 是开集，\(A \subseteq U\)，\(B \subseteq V\)，且 \(U \cap V = \emptyset\)。因此 \(X\) 是正规的。\(\blacksquare\)

\subsection{59.2 度量化定理}\label{ux5ea6ux91cfux5316ux5b9aux7406}

\textbf{定义 59.2}（可度量化）：拓扑空间 \((X, \mathcal{T})\) 称为\textbf{可度量化的}，如果存在 \(X\) 上的度量 \(d\) 使得 \(\mathcal{T}\) 恰为 \(d\) 诱导的拓扑。

\textbf{命题 59.6}：可度量化空间是 \(T_4\)（从而也是 \(T_0\)--\(T_{3.5}\)）。

\emph{证明}：设 \((X, d)\) 为度量空间。对 \(x \neq y\)，\(d(x, y) > 0\)，开球 \(B(x, d(x, y)/2)\) 和 \(B(y, d(x, y)/2)\) 不相交，故 \(T_2\)。

正规性：设 \(A, B\) 是不相交的闭集。定义 \(f(x) = d(x, A) / (d(x, A) + d(x, B))\)（\(d(x, A) = \inf_{a \in A} d(x, a)\)）。则 \(f\) 连续，\(f(A) = \{0\}\)，\(f(B) = \{1\}\)。取 \(U = f^{-1}([0, 1/2))\)，\(V = f^{-1}((1/2, 1])\)。∎

\textbf{定理 59.7}（Urysohn 度量化定理）：每个具有可数基的正则空间 \(X\)（即第二可数的正则空间）可度量化。

\textbf{证明}：设 \(\mathcal{B} = \{B_n\}_{n \in \mathbb{N}}\) 是 \(X\) 的可数基。对每对 \((n, m)\) 满足 \(\overline{B}_n \subseteq B_m\)，由 Urysohn 引理构造连续函数 \(f_{n,m}: X \to [0, 1]\) 使 \(f_{n,m}(\overline{B}_n) = \{0\}\)，\(f_{n,m}(X \setminus B_m) = \{1\}\)。这些函数构成可数族 \(\{f_k\}_{k \in \mathbb{N}}\)。定义嵌入 \(F: X \to [0, 1]^{\mathbb{N}}\) 为 \(F(x) = (f_k(x))_{k \in \mathbb{N}}\)。\(F\) 连续（因每个分量连续）且单射（正则性：若 \(x \neq y\)，存在 \(B_n, B_m\) 使 \(x \in B_n\)，\(\overline{B}_n \subseteq B_m\)，\(y \notin B_m\)，则 \(f_{n,m}(x) = 0 \neq 1 = f_{n,m}(y)\)）。\(F\) 是开映射（对基元素 \(B \in \mathcal{B}\)，\(F(B)\) 是 \(F(X)\) 中开集）。故 \(F\) 是嵌入。\([0, 1]^{\mathbb{N}}\) 可度量化（\(d((a_n), (b_n)) = \sum 2^{-n}|a_n - b_n|\)），故 \(X\) 可度量化。\(\blacksquare\)

\textbf{推论 59.8}：紧致 Hausdorff 空间可度量化当且仅当它是第二可数的。

\emph{证明}：紧致 Hausdorff 空间是正规的，故正则。第二可数正则空间由 Urysohn 度量化定理可度量化。反之，紧致可度量化空间必第二可数（紧致度量空间有可数基）。∎

\textbf{定理 59.9}（Nagata-Smirnov 度量化定理，陈述）：拓扑空间可度量化当且仅当它是正则的且具有 \(\sigma\)-局部有限基。

\textbf{证明概要} （充分性）设 \(X\) 是正则空间且具有 \(\sigma\)-局部有限基 \(\mathcal{B} = \bigcup_{n=1}^\infty \mathcal{B}_n\)，每个 \(\mathcal{B}_n\) 局部有限。对每对 \((n, m)\)，利用正规性（正则+\(\sigma\)-局部有限基\(\Rightarrow\)正规）和 Urysohn 引理构造伪度量 \(d_{n,m}\)。然后将所有伪度量组合为单一度量 \(d\)。关键步骤：（1）由局部有限性证明 \(X\) 是正规的（可逐对分离闭集）；（2）构造适当的开覆盖族使其在 \(\sigma\)-局部有限基下相互加细；（3）利用这些覆盖定义伪度量族并加权求和得度量，其诱导拓扑恰为原拓扑。该构造与 Urysohn 度量化定理的嵌入方法在精神上一致，但更加精细------不依赖于可数性，而是用局部有限性替代。

（必要性）度量空间的 \(\varepsilon\)-球族是 \(\sigma\)-局部有限的（对每个 \(n\)，半径为 \(1/n\) 的开球族是局部有限的------取至多一个球就能覆盖任何点在该尺度下的邻域），且度量空间是正则的（实际上 \(T_4\)）。\(\blacksquare\)

\emph{说明}：这是度量化问题的最一般解答。\(\sigma\)-局部有限基指基 \(\mathcal{B} = \bigcup_{n=1}^\infty \mathcal{B}_n\)，其中每个 \(\mathcal{B}_n\) 是局部有限的（每点有邻域只与 \(\mathcal{B}_n\) 中有限个元素相交）。完备证明见 Nagata (1950) 和 Smirnov (1951) 的原始论文。

\textbf{基本群与同伦理论的历史与展望}

基本群是代数拓扑的起点，其概念由 Poincare 于 1895 年在《Analysis Situs》中引入，当时称为''基本群''（groupe fondamental）。Poincare 的基本洞察是：通过将拓扑空间中的回路赋予群结构，可以将拓扑问题转化为代数问题。这一思想标志着代数拓扑的诞生，后来被推广为同伦群 \(\pi_n\)（Hurewicz 1935）和同调群 \(H_n\)（Noether 1925）。\textbf{Seifert-van Kampen 定理}（1931-1933）是基本群计算中最强大的工具，它断言：若 \(X = U \cup V\)，则 \(\pi_1(X) \cong \pi_1(U) *_{\pi_1(U \cap V)} \pi_1(V)\)（融合自由积）。这一结果将拓扑空间的粘合操作翻译为群论中的融合自由积，为计算复杂空间的基本群提供了系统的方法。\textbf{基本群与覆叠空间理论}的 Galois 对应（见第 61 章）是 20 世纪数学中最重要的对偶原理之一：覆叠空间 \(p: \tilde{X} \to X\) 与 \(\pi_1(X)\) 的子群共轭类之间的一一对应，完全类似于 Galois 理论中域扩张与 Galois 群的子群之间的对应。\textbf{基本群胚}（fundamental groupoid）是基本群的自然推广------它不固定基点，而是考虑所有点之间的道路同伦类，构成一个群胚（所有态射皆可逆的范畴）。基本群胚在非道路连通空间的研究中更为灵活，且与 Grothendieck 的 Galois 理论（代数几何中的平展基本群）有一致的范畴论框架。\textbf{高维同伦群} \(\pi_n(X)\)（\(n \geq 2\)）是 Abel 群，其计算极端困难------甚至 \(\pi_n(S^3)\) 的完全计算至今仍是未解决问题。但同伦论与同调论之间的桥梁------\textbf{Hurewicz 定理}（1935）------断言 \(\pi_1(X)_{ab} \cong H_1(X)\)，且若 \(X\) 是 \((n-1)\)-连通的，则 \(\pi_n(X) \cong H_n(X)\)。这一结果将同伦不变量与同调不变量联系起来，是同伦论中最重要的结构定理之一。

\hrulefill

\hrulefill

\section{第215章：基本群}\label{ux7b2c215ux7ae0ux57faux672cux7fa4}

基本群 \(\pi_1(X, x_0)\) 是代数拓扑的第一个核心不变量，由 Henri Poincare 在 1895 年的开创性论文《Analysis Situs》中引入------这是历史上首次用代数结构（群）来系统地分类拓扑空间。其基本想法极其优美：固定基点 \(x_0\)，考虑空间中基于该点的所有闭路，并将可以连续变形（同伦）为彼此的闭路视为等价；这些等价类在闭路拼接下构成一个群，即基本群。Poincare 的原创动机是对二维曲面的分类------他证明了如果两个闭曲面有不同的基本群，则它们必然不同胚，并由此构造了第一个非平凡同调群的例子（Poincare 同调球面）。在现代数学中，基本群是连接拓扑、几何、分析和代数的桥梁：在拓扑学中，它催生了高维同伦群、覆盖空间理论和 Seifert-van Kampen 定理；在几何中，它是 Riemann 面理论的核心工具（Fuchs 群）；在分析中，Cauchy 积分定理和单值化定理均可表述为基本群的平凡性（单连通性）。本章与前三章（拓扑空间基础、连通性与紧致性、分离公理）共同构成点集拓扑到代数拓扑的自然过渡。

\subsection{60.1 同伦}\label{ux540cux4f26}

\textbf{定义 60.1}（同伦）：设 \(f, g: X \to Y\) 为连续映射。\(f\) 到 \(g\) 的一个\textbf{同伦}是一个连续映射 \(H: X \times [0, 1] \to Y\) 使得 \(H(x, 0) = f(x)\)，\(H(x, 1) = g(x)\) 对所有 \(x \in X\) 成立。若存在这样的同伦，则称 \(f\) \textbf{同伦于} \(g\)，记作 \(f \simeq g\)。

\textbf{命题 60.1}：同伦关系 \(\simeq\) 是 \(C(X, Y)\)（\(X\) 到 \(Y\) 的连续映射集）上的等价关系。

\emph{证明}：自反性：\(H(x, t) = f(x)\) 给出 \(f \simeq f\)。对称性：若 \(H\) 是 \(f \simeq g\) 的同伦，则 \(H'(x, t) = H(x, 1-t)\) 给出 \(g \simeq f\)。传递性：若 \(H_1: f \simeq g\)，\(H_2: g \simeq h\)，定义

\[H(x, t) = \begin{cases} H_1(x, 2t) & 0 \leq t \leq 1/2 \\ H_2(x, 2t-1) & 1/2 \leq t \leq 1 \end{cases}\]

由粘合引理（Pasting Lemma），\(H\) 连续，故 \(f \simeq h\)。∎

\textbf{定义 60.2}（同伦类）：\(f\) 的\textbf{同伦类}记作 \([f]\)。

\textbf{定义 60.3}（道路同伦）：设 \(\gamma, \delta: [0, 1] \to X\) 是两条道路，满足 \(\gamma(0) = \delta(0) = x_0\)，\(\gamma(1) = \delta(1) = x_1\)。\(f\) 到 \(g\) 的一个\textbf{道路同伦}（或定端同伦）是一个同伦 \(H: [0, 1] \times [0, 1] \to X\) 满足 \(H(s, 0) = \gamma(s)\)，\(H(s, 1) = \delta(s)\)，且 \(H(0, t) = x_0\)，\(H(1, t) = x_1\) 对所有 \(t\) 成立。记作 \(\gamma \simeq_p \delta\)。

\textbf{定义 60.4}（同伦等价）：映射 \(f: X \to Y\) 称为\textbf{同伦等价}，如果存在 \(g: Y \to X\) 使得 \(g \circ f \simeq \operatorname{id}_X\) 且 \(f \circ g \simeq \operatorname{id}_Y\)。此时称 \(X\) 与 \(Y\) \textbf{同伦等价}（或具有相同的\textbf{同伦型}），记作 \(X \simeq Y\)。

\textbf{定义 60.5}（可缩空间）：拓扑空间 \(X\) 称为\textbf{可缩的}，如果 \(X\) 同伦等价于单点空间。等价地，\(\operatorname{id}_X\) 同伦于常值映射。

\subsection{60.2 基本群的定义}\label{ux57faux672cux7fa4ux7684ux5b9aux4e49}

\textbf{定义 60.6}（道路的乘法）：设 \(\gamma: [0, 1] \to X\) 是从 \(x\) 到 \(y\) 的道路，\(\delta: [0, 1] \to X\) 是从 \(y\) 到 \(z\) 的道路。它们的\textbf{乘积}（或连接）\(\gamma \cdot \delta: [0, 1] \to X\) 定义为

\[(\gamma \cdot \delta)(s) = \begin{cases} \gamma(2s) & 0 \leq s \leq 1/2 \\ \delta(2s-1) & 1/2 \leq s \leq 1 \end{cases}\]

这是从 \(x\) 到 \(z\) 的道路。

\textbf{定义 60.7}（逆道路）：道路 \(\gamma\) 的\textbf{逆道路} \(\bar{\gamma}: [0, 1] \to X\) 定义为 \(\bar{\gamma}(s) = \gamma(1-s)\)。

\textbf{定义 60.8}（基本群）：设 \(X\) 为拓扑空间，\(x_0 \in X\) 为\textbf{基点}。考虑以 \(x_0\) 为起点和终点的所有道路（即\textbf{回路}，或称\textbf{闭路}）的道路同伦类构成的集合：

\[\pi_1(X, x_0) = \{[\gamma] : \gamma: [0, 1] \to X \text{ 连续}, \gamma(0) = \gamma(1) = x_0\}\]

在 \(\pi_1(X, x_0)\) 上定义乘法：\([\gamma] \cdot [\delta] = [\gamma \cdot \delta]\)。定义单位元为常值道路 \(e_{x_0}\)（\(e_{x_0}(s) = x_0\)）的同伦类。定义逆元 \([\gamma]^{-1} = [\bar{\gamma}]\)。

\textbf{定理 60.2}：\(\pi_1(X, x_0)\) 在上述运算下构成群，称为 \(X\) 的以 \(x_0\) 为基点的\textbf{基本群}（或\textbf{第一同伦群}）。

\emph{证明}：需验证乘法是良定义的、结合律、单位元、逆元。

良定义：若 \(\gamma \simeq_p \gamma'\)（同伦 \(F\)），\(\delta \simeq_p \delta'\)（同伦 \(G\)），则 \(H(s, t) = F(2s, t)\) 对 \(0 \leq s \leq 1/2\)，\(H(s, t) = G(2s-1, t)\) 对 \(1/2 \leq s \leq 1\) 给出 \(\gamma \cdot \delta \simeq_p \gamma' \cdot \delta'\)。

结合律：需证 \((\alpha \cdot \beta) \cdot \gamma \simeq_p \alpha \cdot (\beta \cdot \gamma)\)。构造同伦 \(H(s, t)\) 通过对参数 \(s\) 的重新参数化：\((\alpha \cdot \beta) \cdot \gamma\) 在 \([0, 1/4]\), \([1/4, 1/2]\), \([1/2, 1]\) 上分别走 \(\alpha, \beta, \gamma\)；\(\alpha \cdot (\beta \cdot \gamma)\) 在 \([0, 1/2]\), \([1/2, 3/4]\), \([3/4, 1]\) 上分别走 \(\alpha, \beta, \gamma\)。通过线性插值可构造道路同伦。

单位元：需证 \([\gamma] \cdot [e_{x_0}] = [\gamma] = [e_{x_0}] \cdot [\gamma]\)。\(\gamma \cdot e_{x_0}\) 在 \([0, 1/2]\) 上走 \(\gamma\)，在 \([1/2, 1]\) 上停在 \(x_0\)。通过重新参数化可构造同伦到 \(\gamma\)。

逆元：需证 \([\gamma] \cdot [\bar{\gamma}] = [e_{x_0}]\)。构造同伦 \(H(s, t)\)：当 \(t\) 从 \(0\) 到 \(1\) 时，\(\gamma\) 走得越来越短，然后原路返回，最终停在 \(x_0\)。∎

\textbf{命题 60.3}（基点变换）：设 \(x_0, x_1 \in X\)，\(\eta\) 是从 \(x_0\) 到 \(x_1\) 的道路。则映射

\[\Phi_\eta: \pi_1(X, x_0) \to \pi_1(X, x_1), \quad [\gamma] \mapsto [\bar{\eta} \cdot \gamma \cdot \eta]\]

是群同构。因此，若 \(X\) 道路连通，则 \(\pi_1(X, x_0)\) 的同构型不依赖于基点的选择。

\emph{证明}：\(\Phi_\eta\) 是良定义的且保乘法：\(\Phi_\eta([\alpha][\beta]) = [\bar{\eta} \cdot \alpha \cdot \beta \cdot \eta] = [\bar{\eta} \cdot \alpha \cdot \eta \cdot \bar{\eta} \cdot \beta \cdot \eta] = \Phi_\eta([\alpha]) \Phi_\eta([\beta])\)。其逆为 \(\Phi_{\bar{\eta}}\)。∎

\textbf{定义 60.9}（诱导同态）：设 \(f: X \to Y\) 为连续映射，\(f(x_0) = y_0\)。定义

\[f_*: \pi_1(X, x_0) \to \pi_1(Y, y_0), \quad f_*([\gamma]) = [f \circ \gamma]\]

\textbf{命题 60.4}（诱导同态的性质）： 1. \(f_*\) 是群同态。 2. \((\operatorname{id}_X)_* = \operatorname{id}_{\pi_1(X, x_0)}\)。 3. \((g \circ f)_* = g_* \circ f_*\)。 4. 若 \(f \simeq g\)（同伦），且同伦保持基点，则 \(f_* = g_*\)。

因此 \(\pi_1\) 是从带基点拓扑空间范畴到群范畴的函子。

\emph{证明}：(1) \(f_*([\alpha][\beta]) = [f \circ (\alpha \cdot \beta)] = [(f \circ \alpha) \cdot (f \circ \beta)] = f_*([\alpha]) f_*([\beta])\)。(2)-(3) 显然。(4) 设 \(H: X \times [0, 1] \to Y\) 是 \(f\) 到 \(g\) 的保基点同伦。对回路 \(\gamma\)，\(H(\gamma(s), t)\) 给出 \(f \circ \gamma\) 到 \(g \circ \gamma\) 的道路同伦。∎

\textbf{定理 60.5}：若 \(X\) 与 \(Y\) 同伦等价（道路连通空间），则 \(\pi_1(X) \cong \pi_1(Y)\)。

\emph{证明}：设 \(f: X \to Y\)，\(g: Y \to X\) 为同伦等价，即 \(g \circ f \simeq \operatorname{id}_X\) 且 \(f \circ g \simeq \operatorname{id}_Y\)。需要验证 \(g \circ f\) 和 \(\operatorname{id}_X\) 在基本群上诱导相同的同态。由命题 60.4(4)，若 \(H\) 是 \(g \circ f\) 到 \(\operatorname{id}_X\) 的保基点同伦，则 \((g \circ f)_* = (\operatorname{id}_X)_*\)。但 \((g \circ f)_* = g_* \circ f_*\)，\((\operatorname{id}_X)_* = \operatorname{id}_{\pi_1(X)}\)。故 \(g_* \circ f_* = \operatorname{id}\)。同理 \(f_* \circ g_* = \operatorname{id}\)。因此 \(f_*\) 和 \(g_*\) 互为逆，\(f_*\) 是群同构。\(\blacksquare\)

\textbf{推论 60.6}：可缩空间的基本群是平凡群。

\subsection{60.3 基本群的计算}\label{ux57faux672cux7fa4ux7684ux8ba1ux7b97}

\textbf{定理 60.7}：\(\pi_1(S^1) \cong \mathbb{Z}\)。

\textbf{证明}：考虑覆叠映射 \(p: \mathbb{R} \to S^1\)，\(p(t) = e^{2\pi i t}\)。\(\mathbb{R}\) 单连通（\(\pi_1(\mathbb{R}) = 0\)）。对 \(S^1\) 中以 \(1\) 为基点的回路 \(\gamma\)，由覆叠提升性质，存在唯一提升 \(\tilde{\gamma}: [0,1] \to \mathbb{R}\) 满足 \(\tilde{\gamma}(0) = 0\) 且 \(p \circ \tilde{\gamma} = \gamma\)。定义映射 \(\deg: \pi_1(S^1, 1) \to \mathbb{Z}\) 为 \(\deg([\gamma]) = \tilde{\gamma}(1)\)。\(\deg\) 是良定义的群同态（道路同伦的回路有相同提升终点），且为满射（\(\tilde{\gamma}(t) = nt\) 提升到 \(\gamma(t) = e^{2\pi i nt}\)）。单射性：若 \(\deg([\gamma]) = 0\)，则 \(\tilde{\gamma}\) 是 \(\mathbb{R}\) 中回路，\(\mathbb{R}\) 单连通故 \(\tilde{\gamma}\) 同伦于常值道路，投影后 \(\gamma\) 同伦于常值道路。故 \(\pi_1(S^1) \cong \mathbb{Z}\)。\(\blacksquare\)

\textbf{定理 60.8}：\(\pi_1(X \times Y, (x_0, y_0)) \cong \pi_1(X, x_0) \times \pi_1(Y, y_0)\)。

\emph{证明}：设 \(\pi_X: X \times Y \to X\) 和 \(\pi_Y: X \times Y \to Y\) 为投影。定义同态 \(\Phi: \pi_1(X \times Y, (x_0, y_0)) \to \pi_1(X, x_0) \times \pi_1(Y, y_0)\) 为 \(\Phi([\gamma]) = ((\pi_X)_*[\gamma], (\pi_Y)_*[\gamma]) = ([\pi_X \circ \gamma], [\pi_Y \circ \gamma])\)。\(\Phi\) 是群同态因为 \((\pi_X)_*\) 和 \((\pi_Y)_*\) 都是同态。满射性：对任意 \(([\alpha], [\beta]) \in \pi_1(X) \times \pi_1(Y)\)，定义 \(\gamma(t) = (\alpha(t), \beta(t))\)，则 \(\Phi([\gamma]) = ([\alpha], [\beta])\)。单射性：若 \(\Phi([\gamma]) = (1, 1)\)，则 \(\pi_X \circ \gamma\) 在 \(X\) 中同伦于常值，\(\pi_Y \circ \gamma\) 在 \(Y\) 中同伦于常值。设 \(H_X\) 和 \(H_Y\) 为相应的同伦，则 \(H(s,t) = (H_X(s,t), H_Y(s,t))\) 给出 \(\gamma\) 到常值道路的同伦，故 \([\gamma] = 1\)。\(\blacksquare\)

\textbf{推论 60.9}：\(\pi_1(T^n) \cong \mathbb{Z}^n\)（\(n\) 维环面）。

\textbf{定理 60.10}：对 \(n \geq 2\)，\(\pi_1(S^n)\) 是平凡群。

\textbf{证明}：对 \(n \geq 2\)，\(S^n\) 中任意回路 \(\gamma: [0,1] \to S^n\)。由 Sard 定理，\(\gamma\) 不是满射（像集有零测度，\(n \geq 2\) 时 \(S^1\) 不能填满 \(S^n\)）。取 \(p \in S^n \setminus \gamma([0,1])\)，则 \(\gamma\) 的像完全落在 \(S^n \setminus \{p\}\) 中。由球极投影，\(S^n \setminus \{p\} \cong \mathbb{R}^n\)（可缩），故 \(\gamma\) 在 \(S^n \setminus \{p\}\) 中同伦于常值道路，从而在 \(S^n\) 中同伦于常值道路。等价地，使用单纯逼近：任何回路可逼近为分段线性回路，其像为 \(S^n\) 中一维复形，当 \(n \geq 2\) 时不能覆盖 \(S^n\)，故可缩。因此 \(\pi_1(S^n) = 0\)。\(\blacksquare\)

\subsection{60.4 Seifert-van Kampen 定理}\label{seifert-van-kampen-ux5b9aux7406}

\textbf{定理 60.11}（Seifert-van Kampen 定理）：设 \(X = U \cup V\)，其中 \(U, V\) 是道路连通的开集，\(U \cap V\) 非空且道路连通。取基点 \(x_0 \in U \cap V\)。则下图是推出（pushout）在群范畴中：

\[
\begin{CD}
\pi_1(U \cap V, x_0) @>{i_*}>> \pi_1(U, x_0) \\
@V{j_*}VV @VV{\alpha_*}V \\
\pi_1(V, x_0) @>>{\beta_*}> \pi_1(X, x_0)
\end{CD}
\]

即 \(\pi_1(X, x_0) \cong \pi_1(U, x_0) * \pi_1(V, x_0) / N\)，其中 \(*\) 表示自由积，\(N\) 是由 \(\{i_*(g) j_*(g)^{-1} : g \in \pi_1(U \cap V, x_0)\}\) 生成的正规子群。

\textbf{证明}：设 \(\gamma: [0,1] \to X\) 是以 \(x_0\) 为基点的回路。由 Lebesgue 数引理，取 \([0,1]\) 的细分 \(0 = t_0 < t_1 < \cdots < t_m = 1\) 使每个 \(\gamma([t_{i-1},t_i])\) 完全落在 \(U\) 或 \(V\) 中。若 \(\gamma([t_{i-1},t_i]) \subseteq U\)（或 \(V\)），用 \(U\)（或 \(V\)）中道路连接 \(\gamma(t_{i-1})\) 到 \(x_0\) 再返回，将 \(\gamma\) 分解为 \(U\) 和 \(V\) 中回路的乘积。这证明 \(\pi_1(X)\) 由 \(i_*(\pi_1(U))\) 和 \(j_*(\pi_1(V))\) 生成。关系：若 \(\gamma \in \pi_1(U \cap V)\)，\(i_*(\gamma) \cdot j_*(\gamma)^{-1}\) 在 \(X\) 中同伦于常值道路（因 \(\gamma\) 同时在 \(U\) 和 \(V\) 中）。故 \(\pi_1(X) \cong \pi_1(U) * \pi_1(V) / N\)，其中 \(N = \langle i_*(g) j_*(g)^{-1} : g \in \pi_1(U \cap V) \rangle\)。\(\blacksquare\)

\textbf{推论 60.12}（基本群的计算示例）： 1. \(\pi_1(S^1 \vee S^1) \cong \mathbb{Z} * \mathbb{Z}\)（两个圆周的楔积的基本群是秩为 2 的自由群）。 2. \(\pi_1(\Sigma_g) \cong \langle a_1, b_1, \ldots, a_g, b_g : [a_1, b_1] \cdots [a_g, b_g] = 1 \rangle\)（亏格 \(g\) 的可定向闭曲面的基本群）。 3. \(\pi_1(\mathbb{R}P^2) \cong \mathbb{Z}/2\mathbb{Z}\)。

\emph{证明}：(1) 取 \(U, V\) 为两个圆周的适当开邻域，\(U \cap V\) 可缩，故 \(\pi_1(S^1 \vee S^1) \cong \pi_1(S^1) * \pi_1(S^1) \cong \mathbb{Z} * \mathbb{Z}\)。

\begin{enumerate}
\def\labelenumi{(\arabic{enumi})}
\setcounter{enumi}{1}
\item
  将 \(\Sigma_g\) 表示为 \(4g\) 边形的边粘合。取 \(U\) 为内部（可缩），\(V\) 为边界邻域（同伦等价于 \(2g\) 个圆周的楔积），\(U \cap V \simeq S^1\)。由 Seifert-van Kampen 定理可得。
\item
  \(\mathbb{R}P^2\) 可表示为 \(D^2\) 将边界对径点粘合。取 \(U\) 为内部（可缩），\(V\) 为去心邻域（同伦等价于 \(S^1\)），\(U \cap V \simeq S^1\)。由 Seifert-van Kampen：\(\pi_1(\mathbb{R}P^2) \cong \langle a : a^2 = 1 \rangle \cong \mathbb{Z}/2\mathbb{Z}\)。∎
\end{enumerate}

\subsection{60.5 基本群的应用}\label{ux57faux672cux7fa4ux7684ux5e94ux7528}

\textbf{定理 60.13}（Brouwer 不动点定理，\(n=2\)）：任意连续映射 \(f: D^2 \to D^2\) 有不动点。

\emph{证明}：假设 \(f\) 无不动点。定义 \(r: D^2 \to S^1\) 为从 \(f(x)\) 出发经过 \(x\) 的射线与 \(S^1\) 的交点。则 \(r\) 是连续映射且 \(r|_{S^1} = \operatorname{id}_{S^1}\)（收缩映射）。考虑同态：

\[\pi_1(S^1) \xrightarrow{i_*} \pi_1(D^2) \xrightarrow{r_*} \pi_1(S^1)\]

\(r_* \circ i_* = (r \circ i)_* = \operatorname{id}_{\pi_1(S^1)}\)。但 \(\pi_1(D^2)\) 是平凡群（\(D^2\) 可缩），而 \(\pi_1(S^1) \cong \mathbb{Z}\)，\(\operatorname{id}_{\mathbb{Z}}\) 不能通过平凡群分解，矛盾。∎

\textbf{定理 60.14}（代数基本定理，拓扑证明）：任意非常数复系数多项式 \(P(z) \in \mathbb{C}[z]\) 在 \(\mathbb{C}\) 中有根。

\emph{证明}：设 \(P(z) = z^n + a_{n-1}z^{n-1} + \cdots + a_0\)（不妨设首一）。若 \(P(z)\) 无根，则 \(P(z) \neq 0\) 对所有 \(z \in \mathbb{C}\)。定义 \(f_r: S^1 \to S^1\) 为 \(f_r(z) = P(rz)/|P(rz)|\)（\(r > 0\)）。当 \(r=0\) 时 \(f_0\) 是常值映射 \(z \mapsto P(0)/|P(0)|\)。当 \(r\) 很大时，\(P(rz) = r^n z^n (1 + O(1/r))\)，故 \(f_r(z) \simeq z^n\)。但这些映射对不同的 \(r\) 是同伦的，故常值映射同伦于 \(z \mapsto z^n\)，在基本群上诱导同态：\(0 = n \cdot \operatorname{id}_{\mathbb{Z}}\)，从而 \(n = 0\)，矛盾。∎

\textbf{定理 60.15}（Borsuk-Ulam 定理，\(n=2\)）：对任意连续映射 \(f: S^2 \to \mathbb{R}^2\)，存在 \(x \in S^2\) 使得 \(f(x) = f(-x)\)。

\textbf{证明}：假设不存在 \(x \in S^2\) 使 \(f(x) = f(-x)\)。定义 \(g: S^2 \to S^1\) 为 \(g(x) = \frac{f(x) - f(-x)}{\|f(x) - f(-x)\|}\)。\(g\) 连续且满足 \(g(-x) = -g(x)\)（奇映射）。\(g\) 诱导商映射 \(\bar{g}: \mathbb{R}P^2 \to \mathbb{R}P^1 \cong S^1\)。在基本群层面，\(\bar{g}_*: \pi_1(\mathbb{R}P^2) \cong \mathbb{Z}/2\mathbb{Z} \to \pi_1(S^1) \cong \mathbb{Z}\)。但 \(\mathbb{Z}\) 中无非平凡的 \(2\) 阶元（唯一有限阶元为 \(0\)），故 \(\bar{g}_*\) 必为零同态。然而，\(S^2\) 上的奇映射在商空间上诱导非平凡同态（\(S^2 \to \mathbb{R}P^2\) 为二重覆叠，其提升性质的奇偶性确保 \(\bar{g}_*\) 非平凡）。矛盾。故存在 \(x\) 使 \(f(x) = f(-x)\)。\(\blacksquare\)

\subsection{基本群与低维拓扑的深层联系}\label{ux57faux672cux7fa4ux4e0eux4f4eux7ef4ux62d3ux6251ux7684ux6df1ux5c42ux8054ux7cfb}

基本群理论在低维拓扑中扮演着核心角色，其应用远超经典的 Brouwer 不动点定理和代数基本定理。\textbf{Nielsen-Schreier 定理}（1921-1927）是覆叠空间理论应用于群论的经典范例：自由群的任意子群仍为自由群。其证明利用万有覆叠空间：自由群 \(F_n\) 是 \(n\) 个圆周的楔积 \(S^1 \vee \cdots \vee S^1\) 的基本群，其任意子群对应某个覆叠空间的基本群，而覆叠空间是一维 CW 复形，故其基本群必为自由群。\textbf{二维流形的分类}完全由基本群决定：两个紧致连通曲面 \(S\) 和 \(S'\) 同胚当且仅当 \(\pi_1(S) \cong \pi_1(S')\)（Dehn 1912, Nielsen 1927）。亏格为 \(g\) 的可定向闭曲面的基本群为 \(\langle a_1, b_1, \ldots, a_g, b_g : [a_1,b_1] \cdots [a_g,b_g] = 1 \rangle\)，这一表示完全刻画了曲面的拓扑型。\textbf{纽结理论}中，基本群是最基本的纽结不变量之一：纽结补空间 \(S^3 \setminus K\) 的基本群 \(\pi_1(S^3 \setminus K)\)（称为纽结群）可区分大量不同的纽结，非平凡纽结的纽结群不能是 \(\mathbb{Z}\)。\textbf{基本群与映射类群}（mapping class group）的关联：曲面 \(S\) 的映射类群 \(\operatorname{MCG}(S)\) 定义为同胚 \(\operatorname{Homeo}^+(S)\) 的同伦类的群，Dehn-Nielsen-Baer 定理（1920s）断言 \(\operatorname{MCG}(S) \cong \operatorname{Out}(\pi_1(S))\)（外自同构群），这一结果将曲面的拓扑对称性翻译为基本群的外自同构。\textbf{Perelman 的几何化猜想}（2003，证明 Thurston 的几何化猜想）和 Poincare 猜想的证明中，基本群的角色至关重要：单连通闭三维流形必同胚于 \(S^3\)。基本群在三维流形理论中的决定性作用（\(\pi_1\) 本质上决定了闭三维流形的拓扑型，除透镜空间外）是低维拓扑中最深刻的定理之一，由 Waldhausen（1968）和 Mostow（1968）的刚性定理确立。

\newpage

\chapter{07-拓扑与几何/01-拓扑学/02-覆叠空间.md}\label{ux62d3ux6251ux4e0eux51e0ux4f5501-ux62d3ux6251ux5b6602-ux8986ux53e0ux7a7aux95f4.md}

\begin{quote}
覆叠空间理论是代数拓扑中连接拓扑与代数的核心桥梁。覆叠映射 \(p: \tilde{X} \to X\) 是局部同胚，其纤维 \(p^{-1}(x_0)\) 是离散空间，且 \(\pi_1(X)\) 在纤维上的作用通过路径提升给出。覆叠空间的分类定理（Galois 对应）将覆叠空间与基本群的子群建立一一对应------这是拓扑学中最优美的对偶定理之一，与域扩张的 Galois 理论在形式上完全平行。
\end{quote}

\textbf{覆叠空间理论的历史与整体框架}

覆叠空间理论是代数拓扑中最为优雅的分类理论之一，其思想萌芽可追溯到 Riemann（1851）在研究多值复函数（如 \(\sqrt{z}\) 和 \(\log z\)）时引入的 Riemann 面概念。Riemann 意识到，多值函数的''单值化''等价于在某个覆叠曲面上工作------这正是覆叠空间理论的核心思想。\textbf{Poincare}（1895）在引入基本群的同时系统研究了覆叠空间，而 \textbf{H. Weyl}（1913）在《Die Idee der Riemannschen Flache》中给出了 Riemann 面的严格拓扑定义，奠定了覆叠空间理论的公理化基础。\textbf{覆叠空间的 Galois 理论}（见定理 61.3）是 20 世纪数学中最深刻的对应原理之一：若 \(X\) 是道路连通且局部道路连通且半局部单连通的空间，则覆叠空间 \(p: (\tilde{X}, \tilde{x}_0) \to (X, x_0)\) 的等价类与 \(\pi_1(X, x_0)\) 的子群共轭类一一对应，正规覆叠空间对应于正规子群，覆叠变换群 \(\operatorname{Deck}(\tilde{X}/X) \cong N_{\pi_1(X)}(p_*(\pi_1(\tilde{X})))/p_*(\pi_1(\tilde{X}))\)。这一对应将拓扑学中的覆叠构造翻译为群论中的子群关系，其形式与 Galois 理论中中间域与 Galois 子群的对应完全平行。\textbf{万有覆叠}（universal covering）\(\tilde{X}_u \to X\) 是单连通覆叠（\(\pi_1(\tilde{X}_u) = 0\)），其覆叠变换群 \(\cong \pi_1(X)\)。覆叠空间理论在几何学中扮演着核心角色：\textbf{Thurston 的几何化猜想}（1982）中，八种几何模型空间中，\(\mathbb{R}^3\) 上的 \(S^2 \times \mathbb{R}\) 和 \(\mathbb{H}^2 \times \mathbb{R}\) 的覆叠结构是理解三维流形几何结构的关键。覆叠空间理论在\textbf{非线性分析}中也有重要应用：Krasnoselskii 的覆叠映射方法将有限维拓扑中的覆叠理论推广到无穷维 Banach 空间，用于研究非线性积分方程和微分方程的分歧理论。

\hrulefill

\hrulefill

\section{第216章：覆叠空间的分类}\label{ux7b2c216ux7ae0ux8986ux53e0ux7a7aux95f4ux7684ux5206ux7c7b}

覆叠空间的分类定理是代数拓扑中最优美的对应定理之一，通常被称为''覆叠空间的 Galois 理论''。本章将在第155章建立的覆叠空间基本理论（提升性质、覆叠变换群）的基础上，系统阐述覆叠空间与基本群子群之间的 Galois 对应。核心定理断言：在道路连通、局部道路连通且半局部单连通的空间 \(X\) 上，其（道路连通的）覆叠空间的等价类与 \(\pi_1(X, x_0)\) 的子群共轭类一一对应。正规覆叠对应于正规子群，覆叠变换群 \(\operatorname{Deck}(\tilde{X}/X)\) 同构于 \(N_{\pi_1(X)}(p_*(\pi_1(\tilde{X})))/p_*(\pi_1(\tilde{X}))\)。万有覆叠对应于平凡子群，其覆叠变换群同构于 \(\pi_1(X)\)。这一对偶不仅在拓扑学内部具有基础地位，其形式与域扩张的 Galois 理论完全平行，启发了 Grothendieck 的 Galois 范畴理论和 etale 基本群概念------后者是现代算术几何的基石。

\subsection{覆叠空间的 Galois 对应的严格陈述}\label{ux8986ux53e0ux7a7aux95f4ux7684-galois-ux5bf9ux5e94ux7684ux4e25ux683cux9648ux8ff0}

\textbf{定义 153.1}（覆叠空间的范畴）：设 \(X\) 为道路连通、局部道路连通且半局部单连通的空间。定义 \(\mathbf{Cov}(X)\) 为 \(X\) 上道路连通覆叠空间的范畴，其态射为覆叠映射之间的交换三角。

\textbf{定理 153.1}（覆叠空间分类定理 / Galois 对应）：取定基点 \(x_0 \in X\)。存在 \(\mathbf{Cov}(X)\) 的同构类与 \(\pi_1(X, x_0)\) 的子群共轭类之间的双射 \[\Phi: \{\text{覆叠 } p: \tilde{X} \to X \text{（道路连通）}\}/\cong \;\longrightarrow\; \{\pi_1(X, x_0) \text{ 的子群共轭类}\},\] 将覆叠 \(p: (\tilde{X}, \tilde{x}_0) \to (X, x_0)\) 映为 \(p_*(\pi_1(\tilde{X}, \tilde{x}_0))\)。此双射具有以下函子性质：

\begin{enumerate}
\def\labelenumi{\arabic{enumi}.}
\tightlist
\item
  （Galois 对应）覆叠空间的包含关系反转：\(\tilde{X}_1 \to \tilde{X}_2 \to X\) 当且仅当 \(p_{1*}(\pi_1(\tilde{X}_1)) \subseteq p_{2*}(\pi_1(\tilde{X}_2))\)。
\item
  （正规性）覆叠 \(p: \tilde{X} \to X\) 是正规的（即 \(\operatorname{Deck}(\tilde{X}/X)\) 在纤维上可迁）当且仅当 \(p_*(\pi_1(\tilde{X})) \trianglelefteq \pi_1(X)\)。
\item
  （覆叠变换群）对正规覆叠，\(\operatorname{Deck}(\tilde{X}/X) \cong \pi_1(X) / p_*(\pi_1(\tilde{X}))\)。对一般覆叠，\(\operatorname{Deck}(\tilde{X}/X) \cong N_{\pi_1(X)}(p_*(\pi_1(\tilde{X}))) / p_*(\pi_1(\tilde{X}))\)。
\item
  （万有覆叠）万有覆叠 \(\tilde{X}_u \to X\)（\(\pi_1(\tilde{X}_u) = 1\)）对应于平凡子群 \(\{1\}\)，且 \(\operatorname{Deck}(\tilde{X}_u/X) \cong \pi_1(X)\)。
\end{enumerate}

\textbf{证明}：已在第155章定理61.7中给出。此处补充 Galois 对应的函子性证明。设覆盖 \(p_1: \tilde{X}_1 \to X\) 和 \(p_2: \tilde{X}_2 \to X\) 分别对应子群 \(H_1 = p_{1*}(\pi_1(\tilde{X}_1))\) 和 \(H_2 = p_{2*}(\pi_1(\tilde{X}_2))\)。若存在覆叠态射 \(f: \tilde{X}_1 \to \tilde{X}_2\)（即 \(p_2 \circ f = p_1\)），则 \(p_{1*} = p_{2*} \circ f_*\)，故 \(H_1 = p_{1*}(\pi_1(\tilde{X}_1)) \subseteq p_{2*}(\pi_1(\tilde{X}_2)) = H_2\)。反之，若 \(H_1 \subseteq H_2\)，由提升判别法（定理61.4），\(p_1\) 可通过 \(p_2\) 提升，即存在 \(f: \tilde{X}_1 \to \tilde{X}_2\) 使得 \(p_2 \circ f = p_1\)。函子性得证。正规性：若 \(H \trianglelefteq \pi_1(X)\)，则对任意 \([\gamma] \in \pi_1(X)\)，\([\gamma] H [\gamma]^{-1} = H\)，这意味着任意基点的提升给出纤维上的可迁作用，即覆叠正规。反之亦然。覆叠变换群公式由定理61.9给出。\(\blacksquare\)

\subsection{覆叠空间与 Galois 理论的类比}\label{ux8986ux53e0ux7a7aux95f4ux4e0e-galois-ux7406ux8bbaux7684ux7c7bux6bd4}

\textbf{定理 153.2}（Galois 理论的统一框架）：设 \(k\) 为域，\(k^{\text{sep}}\) 为可分闭包，\(\Gamma_k = \operatorname{Gal}(k^{\text{sep}}/k)\) 为绝对 Galois 群。则 \(k\) 的有限可分扩张与 \(\Gamma_k\) 的有限指数开子群一一对应，Galois 扩张对应于正规子群。这一对应与覆叠空间的 Galois 对应具有完全平行的形式结构，二者可统一在 Grothendieck 的 Galois 范畴框架中。

\textbf{定义 153.2}（Galois 范畴，Grothendieck, SGA1）：一个 \textbf{Galois 范畴} 是满足以下条件的范畴 \(\mathcal{C}\)： 1. \(\mathcal{C}\) 具有有限极限和有限余极限； 2. 每个态射可分解为严格满射后跟单射； 3. 存在''纤维函子'' \(F: \mathcal{C} \to \mathbf{FinSet}\)（取值于有限集合），\(F\) 正合且忠实； 4. 存在函子自同构群 \(\pi = \operatorname{Aut}(F)\)（具有投射有限群结构），使得 \(F\) 诱导范畴等价 \(\mathcal{C} \cong \mathbf{FinSet}_\pi\)（\(\pi\) 连续作用的有限集合范畴）。

覆叠空间范畴 \(\mathbf{Cov}_{\text{fin}}(X)\)（有限层覆叠）和有限可分 \(k\)-代数范畴均为 Galois 范畴。前者纤维函子为 \(F(\tilde{X}) = p^{-1}(x_0)\)，Galois 群为 \(\pi_1(X, x_0)\) 的投射有限完备化；后者纤维函子为 \(F(L) = \operatorname{Hom}_k(L, k^{\text{sep}})\)，Galois 群为 \(\Gamma_k\)。这一定理将拓扑学中的基本群理论与数论中的 Galois 理论统一在范畴论的框架下，是 Grothendieck 引入 etale 基本群和远阿贝尔几何（anabelian geometry）的出发点。

\subsection{覆叠空间的分类：应用与例子}\label{ux8986ux53e0ux7a7aux95f4ux7684ux5206ux7c7bux5e94ux7528ux4e0eux4f8bux5b50}

\textbf{命题 153.3}：\(S^1\) 的连通覆叠空间（不计同构）恰好为展开映射 \(p_n: S^1 \to S^1, z \mapsto z^n\)（\(n \in \mathbb{N}\)）和万有覆叠 \(p: \mathbb{R} \to S^1, t \mapsto e^{2\pi i t}\)。这些覆叠分别对应于 \(\pi_1(S^1) \cong \mathbb{Z}\) 的子群 \(n\mathbb{Z}\)（\(n \geq 0\)，\(n=0\) 对应 \(\mathbb{R}\)）。

\emph{证明}：由定理 61.6，\(\pi_1(S^1) \cong \mathbb{Z}\)。\(\mathbb{Z}\) 的子群恰为 \(\{n\mathbb{Z} : n = 0, 1, 2, \ldots\}\)，其中 \(0\mathbb{Z} = \{0\}\)（平凡子群），\(n\mathbb{Z}\)（\(n \geq 1\)）是正规子群。由分类定理（定理 61.7），\(S^1\) 的连通覆叠与 \(\mathbb{Z}\) 的子群一一对应。\(n=0\) 对应万有覆叠 \(\mathbb{R} \to S^1\)（\(\pi_1(\mathbb{R}) = 0\)，对应平凡子群 \(\{0\}\)）。\(n \geq 1\) 对应 \(n\mathbb{Z}\)，其覆叠为 \(p_n: S^1 \to S^1, z \mapsto z^n\)：\(p_{n*}(\pi_1(S^1)) = n\mathbb{Z}\)（因 \(p_n\) 将生成回路绕 \(n\) 次）。商空间 \(\mathbb{R} / n\mathbb{Z}\)（其中 \(n\mathbb{Z}\) 作用于 \(\mathbb{R}\) 通过平移 \(x \mapsto x + n\)）同胚于 \(\mathbb{R} / (n \cdot \mathbb{Z})\)，而 \(\mathbb{R} / n\mathbb{Z} \cong S^1\)（通过 \(t \mapsto e^{2\pi i t/n}\)），复合 \(S^1 \cong \mathbb{R}/n\mathbb{Z} \to \mathbb{R}/\mathbb{Z} \cong S^1\) 即 \(z \mapsto z^n\)。由此所有连通覆叠已穷举。\(\blacksquare\)

\textbf{命题 153.4}：亏格 \(g \geq 2\) 的紧 Riemann 面的万有覆叠是上半平面 \(\mathbb{H}\)（Poincare 上半平面模型），其覆叠变换群是 \(\operatorname{PSL}(2, \mathbb{R})\) 的离散子群------Fuchs 群。这给出了 Riemann 面的单值化定理：每个 Riemann 面的万有覆叠为 \(\hat{\mathbb{C}}\)、\(\mathbb{C}\) 或 \(\mathbb{H}\) 三者之一。

\textbf{命题 153.5}（三维流形）：Poincare 同调球的万有覆叠是单连通的，但其是否为 \(S^3\) 这一问题长期悬而未决，最终由 Perelman 对 Poincare 猜想的证明（2003）得到肯定回答。一般而言，大部分三维流形具有万有覆叠 \(\mathbb{R}^3\) 或 \(S^3\)，但 Thurston 几何化猜想（1982 年提出，Perelman 2003 年证明）中的八种几何模型中的六种对应非紧万有覆叠。

\section{第217章：覆叠空间}\label{ux7b2c217ux7ae0ux8986ux53e0ux7a7aux95f4}

覆叠空间是研究基本群和拓扑空间结构的重要工具。

\subsection{61.1 覆叠空间的定义}\label{ux8986ux53e0ux7a7aux95f4ux7684ux5b9aux4e49}

\textbf{定义 61.1}（覆叠空间）：设 \(p: \tilde{X} \to X\) 是连续满射。称 \(p\) 为\textbf{覆叠映射}（\(\tilde{X}\) 为 \(X\) 的\textbf{覆叠空间}），如果对每个 \(x \in X\)，存在开邻域 \(U\)（称为\textbf{基本邻域}或\textbf{均匀覆盖邻域}）使得 \(p^{-1}(U) = \bigsqcup_{\alpha} V_\alpha\)（不交并），其中每个 \(V_\alpha\) 是 \(\tilde{X}\) 的开集，且 \(p|_{V_\alpha}: V_\alpha \to U\) 是同胚。

\textbf{定义 61.2}（纤维）：对 \(x \in X\)，\(p^{-1}(x)\) 称为 \(x\) 上的\textbf{纤维}。若 \(X\) 道路连通，则所有纤维具有相同的基数，称为覆叠的\textbf{层数}（或\textbf{重数}）。

\textbf{命题 61.1}：覆叠映射是开映射，且在 \(X\) 为 Hausdorff 空间时 \(\tilde{X}\) 也是 Hausdorff 空间。

\emph{证明}：由局部同胚性质和定义直接可得。∎

\subsection{61.2 提升性质}\label{ux63d0ux5347ux6027ux8d28}

\textbf{定义 61.3}（提升）：设 \(p: \tilde{X} \to X\) 为覆叠映射，\(f: Y \to X\) 为连续映射。\(f\) 的一个\textbf{提升}是一个连续映射 \(\tilde{f}: Y \to \tilde{X}\) 使得 \(p \circ \tilde{f} = f\)。

\textbf{定理 61.2}（道路提升引理）：设 \(p: \tilde{X} \to X\) 为覆叠映射，\(\gamma: [0, 1] \to X\) 是一条道路，\(\gamma(0) = x_0\)。对任意 \(\tilde{x}_0 \in p^{-1}(x_0)\)，存在唯一的提升 \(\tilde{\gamma}: [0, 1] \to \tilde{X}\) 使得 \(\tilde{\gamma}(0) = \tilde{x}_0\)。

\emph{证明}：由 Lebesgue 数引理，将 \([0, 1]\) 分割为 \(0 = t_0 < t_1 < \cdots < t_n = 1\) 使得每个 \(\gamma([t_i, t_{i+1}])\) 包含在某个基本邻域 \(U_i\) 中。在 \(U_0\) 上，取包含 \(\tilde{x}_0\) 的那个片 \(V_0\)，定义 \(\tilde{\gamma}|_{[0, t_1]} = (p|_{V_0})^{-1} \circ \gamma|_{[0, t_1]}\)。然后归纳地继续，每次从已确定的终点出发。唯一性由每一步的局部同胚性质保证。∎

\textbf{定理 61.3}（同伦提升引理）：设 \(p: \tilde{X} \to X\) 为覆叠映射，\(H: Y \times [0, 1] \to X\) 为同伦，\(\tilde{f}: Y \to \tilde{X}\) 是 \(H(\cdot, 0)\) 的提升。则存在 \(H\) 的唯一提升 \(\tilde{H}: Y \times [0, 1] \to \tilde{X}\) 使得 \(\tilde{H}(\cdot, 0) = \tilde{f}\)。

特别地，道路同伦可以唯一提升。

\emph{证明}：对每个 \(y \in Y\)，考虑道路 \(t \mapsto H(y, t)\)，由道路提升引理提升为 \(\tilde{H}(y, \cdot)\)。需验证 \(\tilde{H}\) 连续，这可由 \(Y\) 的紧致性（对 \(Y = [0, 1]\) 情形）或一般地由覆叠映射的局部性质（使用 Lebesgue 数引理的推广）保证。∎

\textbf{定理 61.4}（提升判别法）：设 \(p: \tilde{X} \to X\) 为覆叠映射，\(f: Y \to X\) 连续，\(Y\) 道路连通且局部道路连通。取 \(y_0 \in Y\)，\(x_0 = f(y_0)\)，\(\tilde{x}_0 \in p^{-1}(x_0)\)。则 \(f\) 可提升为 \(\tilde{f}: Y \to \tilde{X}\)（\(\tilde{f}(y_0) = \tilde{x}_0\)）当且仅当

\[f_*(\pi_1(Y, y_0)) \subseteq p_*(\pi_1(\tilde{X}, \tilde{x}_0))\]

\emph{证明}：(\(\Rightarrow\)) 若 \(\tilde{f}\) 存在，则 \(f_* = p_* \circ \tilde{f}_*\)，故 \(f_*(\pi_1(Y)) = p_*(\tilde{f}_*(\pi_1(Y))) \subseteq p_*(\pi_1(\tilde{X}))\)。

(\(\Leftarrow\)) 对 \(y \in Y\)，取道路 \(\alpha\) 从 \(y_0\) 到 \(y\)。\(f \circ \alpha\) 是 \(X\) 中道路，提升到 \(\tilde{X}\) 中起点为 \(\tilde{x}_0\) 的道路 \(\widetilde{f \circ \alpha}\)。定义 \(\tilde{f}(y) = \widetilde{f \circ \alpha}(1)\)。需验证良定义：若 \(\beta\) 是另一条道路，则 \(\alpha \cdot \bar{\beta}\) 是回路，\([f \circ (\alpha \cdot \bar{\beta})] \in f_*(\pi_1(Y)) \subseteq p_*(\pi_1(\tilde{X}))\)，故其提升是 \(\tilde{X}\) 中的回路，从而 \(\widetilde{f \circ \alpha}(1) = \widetilde{f \circ \beta}(1)\)。连续性由局部道路连通性保证。∎

\subsection{61.3 覆叠空间与基本群}\label{ux8986ux53e0ux7a7aux95f4ux4e0eux57faux672cux7fa4}

\textbf{定理 61.5}：设 \(p: \tilde{X} \to X\) 为覆叠映射，\(\tilde{x}_0 \in \tilde{X}\)，\(x_0 = p(\tilde{x}_0)\)。则诱导同态 \(p_*: \pi_1(\tilde{X}, \tilde{x}_0) \to \pi_1(X, x_0)\) 是单射。

\emph{证明}：若 \([\tilde{\gamma}] \in \ker p_*\)，则 \(p \circ \tilde{\gamma}\) 在 \(X\) 中道路同伦于常值回路。由同伦提升引理，此同伦提升为 \(\tilde{X}\) 中的同伦，故 \(\tilde{\gamma}\) 道路同伦于常值回路，即 \([\tilde{\gamma}] = 1\)。∎

\textbf{定理 61.6}（\(\pi_1(S^1) \cong \mathbb{Z}\) 的完整证明）：覆叠映射 \(p: \mathbb{R} \to S^1\)，\(p(t) = e^{2\pi i t}\)。取 \(\tilde{x}_0 = 0\)，\(x_0 = 1\)。由于 \(\mathbb{R}\) 可缩，\(\pi_1(\mathbb{R}, 0) = 1\)，故 \(p_*\) 是单射。考虑纤维 \(p^{-1}(1) = \mathbb{Z}\)。

对任意回路 \(\gamma\) 在 \(S^1\) 中以 \(1\) 为基点，提升 \(\tilde{\gamma}\) 以 \(0\) 为起点，终点 \(\tilde{\gamma}(1) \in \mathbb{Z}\)。定义映射

\[\Phi: \pi_1(S^1, 1) \to \mathbb{Z}, \quad [\gamma] \mapsto \tilde{\gamma}(1)\]

由同伦提升引理，\(\Phi\) 是良定义的。\(\Phi\) 是同态：\(\widetilde{\gamma \cdot \delta}\) 的终点为 \(\tilde{\gamma}(1) + \tilde{\delta}(1)\)（因为 \(\tilde{\delta}\) 的起点为 \(\tilde{\gamma}(1)\)，平移后即得）。\(\Phi\) 是满射：\(n \in \mathbb{Z}\) 对应于 \(t \mapsto e^{2\pi i n t}\)。\(\Phi\) 是单射：若 \(\tilde{\gamma}(1) = 0\)，则 \(\tilde{\gamma}\) 是 \(\mathbb{R}\) 中回路，\(\mathbb{R}\) 可缩故 \(\tilde{\gamma}\) 同伦于常值回路，投影到 \(S^1\) 得 \(\gamma\) 同伦于常值回路。因此 \(\pi_1(S^1) \cong \mathbb{Z}\)。∎

\textbf{定理 61.7}（覆叠空间的分类）：设 \(X\) 道路连通、局部道路连通且半局部单连通。则 \(X\) 的覆叠空间（在同构意义下）与 \(\pi_1(X, x_0)\) 的子群共轭类一一对应。

具体地，覆叠 \(p: (\tilde{X}, \tilde{x}_0) \to (X, x_0)\) 对应于子群 \(p_*(\pi_1(\tilde{X}, \tilde{x}_0)) \subseteq \pi_1(X, x_0)\)。

\textbf{证明}：（存在性）对任意子群 \(H \leq \pi_1(X, x_0)\)，构造 \(X\) 的万有覆叠 \(\tilde{X}_U\)，其覆叠变换群 \(\operatorname{Aut}(\tilde{X}_U/X) \cong \pi_1(X)\)。取商 \(\tilde{X}_H = \tilde{X}_U / H\)，则 \(p_*(\pi_1(\tilde{X}_H)) = H\)。（唯一性）由提升判别法（定理 61.4），若两个覆叠 \(p_1, p_2\) 对应相同子群，则存在相互提升映射，构成同构。（分类）不同基点对应共轭子群，故覆叠同构类与子群共轭类一一对应。\(\blacksquare\)

\textbf{定义 61.4}（万有覆叠）：若 \(\tilde{X}\) 单连通（即 \(\pi_1(\tilde{X}) = 1\)），则称 \(p: \tilde{X} \to X\) 为 \(X\) 的\textbf{万有覆叠}。万有覆叠（若存在）在同构意义下唯一。

\textbf{命题 61.8}：\(X\) 存在万有覆叠当且仅当 \(X\) 是道路连通、局部道路连通且半局部单连通的。

\emph{证明}（存在性）：设 \(X\) 是道路连通、局部道路连通且半局部单连通的。固定基点 \(x_0 \in X\)。定义 \(\tilde{X}\) 为从 \(x_0\) 出发的道路的道路同伦类的集合：

\[\tilde{X} = \{[\gamma] : \gamma: [0,1] \to X, \gamma(0) = x_0\}\]

定义投影 \(p: \tilde{X} \to X\) 为 \(p([\gamma]) = \gamma(1)\)。

在 \(\tilde{X}\) 上定义拓扑：对 \([\gamma] \in \tilde{X}\)，令 \(U\) 是 \(x = \gamma(1)\) 的半局部单连通邻域（即 \(\pi_1(U, x) \to \pi_1(X, x)\) 的像为平凡子群），定义基元素

\[B([\gamma], U) = \{[\gamma \cdot \eta] : \eta \text{ 是 } U \text{ 中从 } x \text{ 出发的道路}\}\]

这些集合构成 \(\tilde{X}\) 上拓扑的基。验证：(a) \(p\) 是覆叠映射------\(p\) 将每个 \(B([\gamma], U)\) 同胚地映到 \(U\)（因 \(U\) 中连接 \(x\) 和 \(y\) 的任意两道路在 \(X\) 中必道路同伦，由半局部单连通性）；(b) \(\tilde{X}\) 是道路连通的（任意点 \([\gamma]\) 可通过 \(t \mapsto [\gamma_t]\) 连接到 \([\varepsilon_{x_0}]\)）；(c) \(\tilde{X}\) 是单连通的------任何 \(\tilde{X}\) 中基于 \([\varepsilon_{x_0}]\) 的回路投影到 \(X\) 中同伦于常值回路的道路，由同伦提升引理，该回路在 \(\tilde{X}\) 中也同伦于常值回路。故 \(\tilde{X}\) 是 \(X\) 的万有覆叠。

（必要性）若万有覆叠存在，由覆叠局部同胚性质，\(X\) 需满足局部性质：每点存在邻域 \(U\) 使得其逆像为不交并的同胚拷贝，这等价于 \(U\) 在 \(X\) 中的包含映射在 \(\pi_1\) 上为平凡同态------即半局部单连通性。道路连通性和局部道路连通性则由覆叠空间定义自然隐含（覆叠空间自身具有这些性质，投影为局部同胚，故底空间亦然）。\(\blacksquare\)

\textbf{定义 61.5}（半局部单连通）：\(X\) 称为\textbf{半局部单连通的}，如果对每个 \(x \in X\)，存在邻域 \(U\) 使得 \(\pi_1(U, x) \to \pi_1(X, x)\) 的像是平凡子群。

\subsection{61.4 覆叠变换群}\label{ux8986ux53e0ux53d8ux6362ux7fa4}

\textbf{定义 61.6}（覆叠变换 / 甲板变换）：设 \(p: \tilde{X} \to X\) 为覆叠映射。覆叠 \(\tilde{X}\) 的\textbf{覆叠变换}（或甲板变换）是一个同胚 \(\varphi: \tilde{X} \to \tilde{X}\) 使得 \(p \circ \varphi = p\)。所有覆叠变换构成群，记作 \(\operatorname{Aut}(\tilde{X}/X)\) 或 \(\operatorname{Deck}(\tilde{X}/X)\)。

\textbf{定义 61.7}（正规覆叠）：覆叠 \(p: \tilde{X} \to X\) 称为\textbf{正规的}（或\textbf{Galois 覆叠}），如果对任意 \(\tilde{x}_1, \tilde{x}_2 \in \tilde{X}\) 满足 \(p(\tilde{x}_1) = p(\tilde{x}_2)\)，存在覆叠变换 \(\varphi\) 使得 \(\varphi(\tilde{x}_1) = \tilde{x}_2\)。

\textbf{定理 61.9}：设 \(p: (\tilde{X}, \tilde{x}_0) \to (X, x_0)\) 为覆叠映射，\(\tilde{X}\) 道路连通。则

\[\operatorname{Aut}(\tilde{X}/X) \cong N(p_*(\pi_1(\tilde{X}, \tilde{x}_0))) / p_*(\pi_1(\tilde{X}, \tilde{x}_0))\]

其中 \(N(H)\) 是 \(H\) 在 \(\pi_1(X, x_0)\) 中的正规化子。特别地，若 \(\tilde{X}\) 是万有覆叠，则 \(\operatorname{Aut}(\tilde{X}/X) \cong \pi_1(X, x_0)\)。

\textbf{证明}：覆叠变换由基点的选择决定。对于 \([\gamma] \in N(p_*(\pi_1(\tilde{X})))\)，\(\gamma\) 的提升将 \(\tilde{x}_0\) 映到某点 \(\tilde{x}_1\)，导出覆叠变换。不同 \([\gamma]\) 给出相同变换当且仅当它们相差 \(p_*(\pi_1(\tilde{X}))\) 中元素。∎

\textbf{推论 61.10}：覆叠 \(p: \tilde{X} \to X\) 正规当且仅当 \(p_*(\pi_1(\tilde{X}))\) 是 \(\pi_1(X)\) 的正规子群。

\subsection{61.5 覆叠空间的应用}\label{ux8986ux53e0ux7a7aux95f4ux7684ux5e94ux7528}

\textbf{定理 61.11}（\(\pi_1(S^n)\) 平凡，\(n \geq 2\)，完整证明）：\(S^n\)（\(n \geq 2\)）是单连通的。

\emph{证明}：取 \(U = S^n \setminus \{N\}\)（去掉北极），\(V = S^n \setminus \{S\}\)（去掉南极）。则 \(U \cong \mathbb{R}^n\)，\(V \cong \mathbb{R}^n\)，\(U \cap V \cong S^{n-1} \times \mathbb{R}\) 同伦等价于 \(S^{n-1}\)。当 \(n \geq 2\) 时 \(S^{n-1}\) 道路连通。由 Seifert-van Kampen 定理，\(\pi_1(S^n)\) 由 \(\pi_1(U)\) 和 \(\pi_1(V)\) 生成，但两者都是平凡群（\(\mathbb{R}^n\) 可缩），故 \(\pi_1(S^n)\) 平凡。∎

\textbf{定理 61.12}（\(\mathbb{R}^n \not\cong \mathbb{R}^m\)，\(n \neq m\)）：若 \(n \neq m\)，则 \(\mathbb{R}^n\) 与 \(\mathbb{R}^m\) 不同胚。

\emph{证明}：若 \(\mathbb{R}^n \cong \mathbb{R}^m\)，则它们的单点紧化也同胚：\(S^n \cong S^m\)。但这对 \(n \neq m\) 不成立（可由同调群证明，此处仅用基本群：当 \(n=1, m \geq 2\) 时 \(\pi_1(S^1) \cong \mathbb{Z} \neq 1 \cong \pi_1(S^m)\)）。一般 \(n \neq m\) 的证明需要高维同伦群或同调群。∎

\textbf{定理 61.13}（基本群与自由群）：\(\pi_1(S^1 \vee S^1 \vee \cdots \vee S^1)\)（\(n\) 个圆周的楔积）同构于秩为 \(n\) 的自由群 \(F_n\)。

\emph{证明}：由 Seifert-van Kampen 定理（推论 60.12）归纳可得。万有覆叠是秩为 \(n\) 的自由群的 Cayley 图（无穷树）。∎

\textbf{定理 61.14}（Nielsen-Schreier 定理）：自由群的子群是自由群。

\emph{证明}（拓扑证明）：设 \(F_n\) 是秩为 \(n\) 的自由群，\(H \leq F_n\)。\(F_n \cong \pi_1(X)\) 其中 \(X\) 是 \(n\) 个圆周的楔积。由覆叠分类定理，存在覆叠 \(p: \tilde{X} \to X\) 使得 \(p_*(\pi_1(\tilde{X})) = H\)。\(\tilde{X}\) 是一维 CW 复形（图），其基本群是自由群，故 \(H\) 是自由群。∎

\hrulefill

\emph{本卷构建了拓扑学的基础框架：从拓扑空间的定义出发，建立了连续性、连通性、紧致性等核心概念，证明了 Urysohn 引理、Tychonoff 定理、Urysohn 度量化定理等核心定理，进而引入基本群和覆叠空间理论，为后续学习代数拓扑（V15）、微分几何（V12）和泛函分析（V14）奠定了基础。}

\emph{卷十：拓扑学终。}

\emph{卷十：拓扑学终。}

\emph{卷十：拓扑学终。}

\emph{卷十二：拓扑学终。} \hrulefill

\textbf{同伦类型论与数学基础}

同伦类型论（Homotopy Type Theory, HoTT）是 21 世纪初由 Voevodsky、Awodey、Warren 等人发展起来的数学基础新范式，其核心洞见是：Martin-Löf 依赖类型论中的恒等类型 \(\operatorname{Id}_A(a,b)\) 在群胚模型（Hofmann-Streicher, 1996）和单纯集模型（Voevodsky, 2009）中自然地解释为拓扑空间中的路径空间，从而建立了 \textbf{类型论与同伦论之间的深刻对应}。这一对应的核心支柱是 Voevodsky 的\textbf{单值公理}（Univalence Axiom）------它断言类型的相等等价于类型的等价，将类型论从纯粹的集合论框架提升为同伦论的框架。HoTT 不仅为数学基础提供了新的视角，还为计算数学（如 Coq 和 Agda 证明助手中的形式化数学）提供了实用的工作语言，并催生了\textbf{高阶归纳类型}（HITs）等全新概念------这些概念允许在类型论中直接定义同伦不变量（如圆周 \(S^1\)、球面 \(S^n\) 等）。

\subsection{Martin-Löf 依赖类型论的正式判断规则}\label{martin-luxf6f-ux4f9dux8d56ux7c7bux578bux8bbaux7684ux6b63ux5f0fux5224ux65adux89c4ux5219}

\textbf{定义 1}（上下文与判断）：Martin-Löf 依赖类型论（MLTT）的基本句法实体为\textbf{判断}（judgment），共有四种形式：

\begin{enumerate}
\def\labelenumi{\arabic{enumi}.}
\tightlist
\item
  \(\Gamma \ \mathsf{ctx}\) ------ \(\Gamma\) 是一个有效的上下文；
\item
  \(\Gamma \vdash A \ \mathsf{type}\) ------ \(A\) 在上下文 \(\Gamma\) 下是一个良定义的类型；
\item
  \(\Gamma \vdash A \equiv B \ \mathsf{type}\) ------ 类型 \(A\) 与 \(B\) 定义相等（判断性相等）；
\item
  \(\Gamma \vdash a : A\) ------ \(a\) 是类型 \(A\)（在上下文 \(\Gamma\) 下）的项。
\end{enumerate}

空上下文 \(\cdot\) 是有效的。若 \(\Gamma \vdash A \ \mathsf{type}\)，则 \(\Gamma, x:A\)（\(x\) 为新变量）为有效上下文（变量规则）。

\textbf{定义 2}（\(\Pi\)-类型的形成、引入与消除规则）：依赖函数类型（dependent product）的形成规则为 \[\frac{\Gamma \vdash A \ \mathsf{type} \quad \Gamma, x:A \vdash B(x) \ \mathsf{type}}{\Gamma \vdash \prod_{x:A} B(x) \ \mathsf{type}}\] 引入规则（\(\lambda\)-抽象）：若 \(\Gamma, x:A \vdash b : B(x)\)，则 \(\Gamma \vdash \lambda x. b : \prod_{x:A} B(x)\)。 消除规则（函数应用）：若 \(\Gamma \vdash f : \prod_{x:A} B(x)\) 且 \(\Gamma \vdash a : A\)，则 \(\Gamma \vdash f(a) : B(a)\)。 计算规则（\(\beta\)-归约）：\((\lambda x. b)(a) \equiv b[a/x]\)（判断性相等）。

\textbf{定义 3}（\(\Sigma\)-类型的形成、引入与消除规则）：依赖和类型（dependent sum）的形成规则为 \[\frac{\Gamma \vdash A \ \mathsf{type} \quad \Gamma, x:A \vdash B(x) \ \mathsf{type}}{\Gamma \vdash \sum_{x:A} B(x) \ \mathsf{type}}\] 引入规则（配对）：若 \(\Gamma \vdash a : A\) 且 \(\Gamma \vdash b : B(a)\)，则 \(\Gamma \vdash (a,b) : \sum_{x:A} B(x)\)。 消除规则（投影）：若 \(\Gamma \vdash p : \sum_{x:A} B(x)\)，则 \(\Gamma \vdash \pi_1(p) : A\) 且 \(\Gamma \vdash \pi_2(p) : B(\pi_1(p))\)。 计算规则：\(\pi_1((a,b)) \equiv a\)，\(\pi_2((a,b)) \equiv b\)。

\textbf{定义 4}（恒等类型的形成、引入与消除 ------ J-规则）：恒等类型（identity type）的形成规则为 \[\frac{\Gamma \vdash a : A \quad \Gamma \vdash b : A}{\Gamma \vdash \operatorname{Id}_A(a,b) \ \mathsf{type}}\] 引入规则（自反性）：对任意 \(\Gamma \vdash a : A\)，有 \(\Gamma \vdash \operatorname{refl}_a : \operatorname{Id}_A(a,a)\)。 消除规则（J-规则，路径归纳）：给定 \[\Gamma, x:A, y:A, p:\operatorname{Id}_A(x,y) \vdash C(x,y,p) \ \mathsf{type}\] 及 \(\Gamma, z:A \vdash c(z) : C(z,z,\operatorname{refl}_z)\)， 则存在项 \(\Gamma, x:A, y:A, p:\operatorname{Id}_A(x,y) \vdash J_{x,y,p}(c) : C(x,y,p)\)。 计算规则：\(J_{x,x,\operatorname{refl}_x}(c) \equiv c(x)\)。

\textbf{定义 5}（宇宙 \(\mathcal{U}\)）：存在宇宙类型的层级 \(\mathcal{U}_0, \mathcal{U}_1, \mathcal{U}_2, \ldots\)，满足：

\begin{enumerate}
\def\labelenumi{\arabic{enumi}.}
\tightlist
\item
  \(\Gamma \vdash \mathcal{U}_i \ \mathsf{type}\)（宇宙本身为类型）；
\item
  \(\mathcal{U}_i : \mathcal{U}_{i+1}\)（累积性）；
\item
  若 \(\Gamma \vdash A : \mathcal{U}_i\)，则 \(\Gamma \vdash A \ \mathsf{type}\)（宇宙中的对象即类型）。
\end{enumerate}

\subsection{恒等类型的归纳原理与路径归纳的等价性}\label{ux6052ux7b49ux7c7bux578bux7684ux5f52ux7eb3ux539fux7406ux4e0eux8defux5f84ux5f52ux7eb3ux7684ux7b49ux4ef7ux6027}

\textbf{定理 1}（J-规则等价于路径归纳）：下列规则等价于上述 J-规则（基于路径归纳的恒等类型消除）： 对任意类型族 \(C: \prod_{x,y:A} (\operatorname{Id}_A(x,y) \to \mathcal{U})\) 及项 \(c: \prod_{x:A} C(x,x,\operatorname{refl}_x)\)，存在依赖函数 \[\operatorname{ind}_{\operatorname{Id}_A}: \prod_{x,y:A} \prod_{p:\operatorname{Id}_A(x,y)} C(x,y,p),\] 且满足 \(\operatorname{ind}_{\operatorname{Id}_A}(x,x,\operatorname{refl}_x) \equiv c(x)\)。

\textbf{证明}：给定上述类型的 \(C\) 和 \(c\)，以 \(C(x,y,p)\) 作为消除规则前提的类型族，以 \(c(z)\) 作为自反情形的数据，应用 J-规则得所需归纳项。反之，给定 J-规则前提的类型族 \(C(x,y,p)\)，定义 \(C'(p) := \prod_{x:A} C(x,x,p)\)（在某 \(A\) 的封装上下文中可调整），则 \(\operatorname{refl}\) 情形对应于 \(\lambda z. c(z)\) 的类型，路径归纳返回 \(J\) 的特化形式。两条规则在判断性等价层面等价。\(\blacksquare\)

\subsection{群胚解释}\label{ux7fa4ux80daux89e3ux91ca}

\textbf{定义 6}（类型作为 \(\infty\)-广群的直观对应）：Hofmann 与 Streicher（1996）的群胚解释建立如下对应：

\begin{itemize}
\tightlist
\item
  类型 \(A\) ↭ \(\infty\)-广群（或拓扑空间），
\item
  项 \(a : A\) ↭ \(\infty\)-广群中的点，
\item
  恒等类型 \(\operatorname{Id}_A(a,b)\) ↭ 从 \(a\) 到 \(b\) 的路径空间 \(\Omega_{a,b}(A)\)，
\item
  \(\operatorname{refl}_a\) ↭ 常值路径，
\item
  迭代恒等类型 \(\operatorname{Id}_{\operatorname{Id}_A(a,b)}(p,q)\) ↭ 路径之间的同伦。
\end{itemize}

\textbf{定理 2}（Transport 引理）：设 \(P: A \to \mathcal{U}\) 为类型族，\(p: \operatorname{Id}_A(a,b)\)。则存在\textbf{传输函数} \[\operatorname{transport}^P(p, -): P(a) \to P(b),\] 满足 \(\operatorname{transport}^P(\operatorname{refl}_a, u) \equiv u\)。

\textbf{证明}：由路径归纳（J-规则）。视 \(C(a,b,p) := (P(a) \to P(b))\) 为 \(a,b,p\) 依赖的类型族。在自反情形，\(C(a,a,\operatorname{refl}_a) = (P(a) \to P(a))\)，取恒等函数 \(\operatorname{id}_{P(a)}\)。J-规则即给出所需 \(\operatorname{transport}^P\)。\(\blacksquare\)

\subsection{单值公理}\label{ux5355ux503cux516cux7406}

\textbf{定义 7}（类型之间的等价）：类型 \(A\) 与 \(B\) 之间的\textbf{等价} \(A \simeq B\) 定义为 \[\sum_{f: A \to B} \left(\sum_{g: B \to A} (g \circ f \sim \operatorname{id}_A) \right) \times \left(\sum_{h: B \to A} (f \circ h \sim \operatorname{id}_B) \right)\] 并附加相干条件（半伴随等价，或等价的合约纤维刻画 \(\operatorname{isContr}(\sum_{a:A} \operatorname{Id}_B(f(a), b))\)）。

从类型相等可构造等价：存在典范映射 \[\operatorname{idToEquiv}_{A,B}: \operatorname{Id}_{\mathcal{U}}(A,B) \to (A \simeq B),\] 由 J-规则定义（自反情形 \(A = A\) 映为单位等价）。

\textbf{公理 1}（单值公理，Voevodsky）：上述映射 \(\operatorname{idToEquiv}_{A,B}\) 对所有 \(A, B: \mathcal{U}\) 都是等价： \[\operatorname{Id}_{\mathcal{U}}(A,B) \simeq (A \simeq B).\]

\textbf{定理 3}（Voevodsky 的单纯集模型）：单值公理在单纯集范畴的 Kan 纤维化模型中成立，其中类型解释为 Kan 复形，宇宙 \(\mathcal{U}\) 解释为泛 Kan 复形（按大小分级）。此时 \(\operatorname{Id}_{\mathcal{U}}(A,B)\) 解释为路径空间，\(A \simeq B\) 解释为同伦等价空间，单值公理即断言这两空间的同伦等价性，这在单纯集的模型范畴框架中可通过组合论证验证。

\textbf{证明}：在 \(\mathbf{sSet}_{\text{Kan}}\) 模型中，类型 \(\Gamma \vdash A \ \mathsf{type}\) 解释为 Kan 纤维化 \(p: A \to \Gamma\)。\(\operatorname{Id}_A(a,b)\) 解释为路径空间 \(A^{\Delta[1]} \times_{A \times A} \{(a,b)\}\)，即从 \(a\) 到 \(b\) 的所有路径构成的 Kan 复形。宇宙 \(\mathcal{U}\) 解释为泛 Kan 纤维化 \(\tilde{\mathcal{U}} \to \mathcal{U}\)，其纤维为相应大小的 Kan 复形。\(\operatorname{idToEquiv}_{A,B}\) 将类型等式 \(p: \operatorname{Id}_{\mathcal{U}}(A,B)\) 通过 J-规则提升为等价 \(A \simeq B\)，在模型中这对应于路径空间到同伦等价空间的自然映射。单值公理的有效性源于：在 Kan-Quillen 模型结构中，Kan 复形之间的同伦等价空间 \(\operatorname{Equiv}(A,B)\) 弱等价于路径空间 \(A^{\Delta[1]}\) 中连接对应顶点的路径空间。Voevodsky 通过构造单纯集的收缩论证证明了该映射为弱等价，核心步骤是证明 \(\operatorname{Equiv}(A,B)\) 的纤维可缩。\(\blacksquare\)

\subsection{高阶归纳类型（HITs）}\label{ux9ad8ux9636ux5f52ux7eb3ux7c7bux578bhits}

\textbf{定义 8}（圆周 \(S^1\) 作为 HIT）：\textbf{圆周} \(S^1\) 由以下生成元高阶归纳定义：

\begin{itemize}
\tightlist
\item
  \textbf{点构造器}：\(\mathsf{base} : S^1\)；
\item
  \textbf{路径构造器}：\(\mathsf{loop} : \operatorname{Id}_{S^1}(\mathsf{base}, \mathsf{base})\)。
\end{itemize}

其消除规则：给定类型族 \(P: S^1 \to \mathcal{U}\)，点 \(b : P(\mathsf{base})\)，及（依赖）路径 \(\ell : \operatorname{transport}^P(\mathsf{loop}, b) =_{P(\mathsf{base})} b\)，存在 \(\operatorname{ind}_{S^1}(b,\ell) : \prod_{x:S^1} P(x)\)，满足 \(\operatorname{ind}_{S^1}(b,\ell)(\mathsf{base}) \equiv b\) 且 \(\operatorname{apd}_{\operatorname{ind}_{S^1}(b,\ell)}(\mathsf{loop}) = \ell\)（计算规则以命题等式给出，或判断性等式取决于具体形式化）。

\textbf{定理 4}（\(\pi_1(S^1) = \mathbb{Z}\)）：在 HoTT 中，\(\pi_1(S^1, \mathsf{base}) := \|\operatorname{Id}_{S^1}(\mathsf{base}, \mathsf{base})\|_0\)（环路类型的 0-截断）同构于 \(\mathbb{Z}\)。

\textbf{证明}：定义映射 \(\mathbb{Z} \to \operatorname{Id}_{S^1}(\mathsf{base}, \mathsf{base})\) 为 \(n \mapsto \mathsf{loop}^n\)（\(\mathsf{loop}\) 的 \(n\) 次合成，\(n < 0\) 时沿反向）。反之，定义万有覆叠 \(\widetilde{S^1}: S^1 \to \mathcal{U}\) 为 HIT 消除：\(\widetilde{S^1}(\mathsf{base}) := \mathbb{Z}\)，\(\widetilde{S^1}(\mathsf{loop}) := \operatorname{ua}(\operatorname{succ})\)，其中 \(\operatorname{succ}: \mathbb{Z} \to \mathbb{Z}\) 为后继函数，\(\operatorname{ua}\) 为单值公理给出的路径。则任意 \(p: \operatorname{Id}_{S^1}(\mathsf{base}, \mathsf{base})\) 诱导 \(\operatorname{transport}^{\widetilde{S^1}}(p, 0) : \mathbb{Z}\)，此传输计算给出所需的互逆同构。在 0-截断后，路径组合对应于整数加法。\(\blacksquare\)

\textbf{定义 9}（\(n\)-截断 \(\|A\|_n\) 作为 HIT）：类型的 \(n\)-截断 \(\|A\|_n\) 由以下构造器定义：

\begin{itemize}
\tightlist
\item
  \textbf{点构造器}：\(|-| : A \to \|A\|_n\)；
\item
  \textbf{路径构造器}（对 \(i = 0, 1, \ldots, n\)）：对任意 \(k+1\) 维球面（\(k \geq n\)），所有映射 \(S^{k+1} \to \|A\|_n\) 有标准填充------即 \(\|A\|_n\) 是 \(n\)-连通类型（其 \(\pi_{>n}\) 平凡）。
\end{itemize}

\subsection{h-层级}\label{h-ux5c42ux7ea7}

\textbf{定义 10}（h-层级，Voevodsky）：类型的同伦层级定义如下：

\begin{itemize}
\tightlist
\item
  类型 \(A\) 是\textbf{可缩的}（\((-2)\)-截断），若 \(\operatorname{isContr}(A) := \sum_{a:A} \prod_{x:A} \operatorname{Id}_A(a,x)\)；
\item
  \(A\) 是\textbf{命题}（\((-1)\)-截断），若 \(\operatorname{isProp}(A) := \prod_{x,y:A} \operatorname{Id}_A(x,y)\)（即任意两个项命题地相等）；
\item
  \(A\) 是\textbf{集合}（\(0\)-截断），若 \(\operatorname{isSet}(A) := \prod_{x,y:A} \operatorname{isProp}(\operatorname{Id}_A(x,y))\)（即 \(\operatorname{Id}_A\) 的值都是命题）；
\item
  归纳地，\(A\) 是 \textbf{\((n+1)\)-截断}，若对所有 \(x,y:A\)，\(\operatorname{Id}_A(x,y)\) 是 \(n\)-截断类型。
\end{itemize}

\textbf{定理 5}（截断与同伦群的 HoTT 定义）：同伦群的内部定义为 \[\pi_n(A, a) := \|\Omega^n(A, a)\|_0,\] 其中 \(\Omega(A,a) := \operatorname{Id}_A(a,a)\)（环路空间），\(\Omega^{n+1} := \Omega \circ \Omega^n\)，\(\|\cdot\|_0\) 为 0-截断。此定义与经典同伦论中 \(\pi_n\) 的定义吻合，且 \(\pi_1\) 为群、\(\pi_{n \geq 2}\) 为 Abel 群的结构可通过 Eckmann-Hilton 论证在 HoTT 内部证明。

\textbf{证明}：0-截断 \(\|\cdot\|_0\) 将类型映为集合，从而 \(\pi_n(A,a)\) 是集合。\(\pi_1\) 的群结构：路径合成 \(p \cdot q\) 在 \(\|\cdot\|_0\) 下定义乘法，结合律和单位律由 J-规则验证。\(\pi_{n \geq 2}\) 的 Abel 性：由 Eckmann-Hilton 论证，在 \(\Omega^2(A,a)\) 中两种合成方式（由环路空间的横向和纵向合成）满足互换律；通过反复构造，\(n \geq 2\) 时环路空间中的任意两个运算强制相等（在命题截断下这是精确的）。与经典定义的一致性由 \(\Sigma\)-类型解释和截断的泛性质保证。\(\blacksquare\)

\subsection{基础意义}\label{ux57faux7840ux610fux4e49}

覆叠空间理论是代数拓扑中''用基本群分类空间''的第一个成功范例。它奠定了代数拓扑中''空间 \(\leftrightarrow\) 群（基本群）\(\leftrightarrow\) 子群共轭类 \(\leftrightarrow\) 覆叠空间''一一对应的核心范式，证明了''空间的拓扑结构完全被它的基本群（及其子群）刻画''------这在代数拓扑发展史上具有里程碑意义。覆叠空间为研究基本群提供了几何工具：利用覆叠变换群计算 \(\pi_1(S^1) \cong \mathbb{Z}\)，用万有覆叠证明 Nielsen-Schreier 定理（自由群的子群仍自由），证明 Seifert-van Kampen 定理的部分情形，以及对 \(\mathbb{R}^n\) 和 \(\mathbb{R}^m\) 不同胚的拓扑论证。覆叠空间理论与现代同伦论中纤维丛概念有深刻联系：覆叠映射实际上就是离散纤维的纤维丛，而纤维丛理论正是覆叠空间理论在非离散纤维情形下的推广。在同伦类型论（同伦论的高阶归纳类型版本，见下文）中，覆叠空间对应于类型论中的依赖类型：每个覆叠 \(\tilde{X} \to X\) 就是一个 \(X\) 上的依赖类型，其纤维对应于每个基点 \(x \in X\) 处的类型 \(p^{-1}(x)\)，而路径提升对应于依赖类型之间的传输。这一对应关系揭示了拓扑学和类型论之间的深刻对偶。覆叠空间理论的另一个重要应用是在黎曼曲面论中：紧 Riemann 面的覆叠空间对应于函数域的有限扩张，从而将代数曲线的覆盖与拓扑覆叠统一起来。

\hrulefill

\subsection{高阶归纳类型的进一步示例}\label{ux9ad8ux9636ux5f52ux7eb3ux7c7bux578bux7684ux8fdbux4e00ux6b65ux793aux4f8b}

高阶归纳类型（higher inductive types, HITs）是同伦类型论（HoTT/UF）中最强大的特性之一，它允许我们通过点构造器（引入点）和路径构造器（引入高阶路径）直接归纳定义拓扑空间。HITs 将拓扑空间的''胞腔粘和''直接在类型论语法层面实现，使得我们可以用归纳方法证明空间的同伦性质。下面给出几个典型的高阶归纳类型示例：

\begin{enumerate}
\def\labelenumi{\arabic{enumi}.}
\item
  \textbf{区间 \(I\)}：由两个点 \(\mathsf{left}, \mathsf{right} : I\) 和一条路径 \(\mathsf{seg} : \operatorname{Id}_I(\mathsf{left}, \mathsf{right})\) 定义。作为 HIT，\(I\) 在 HoTT 中满足同伦等价 \(I \simeq *\)（即 \(I\) 可缩），这恰好对应经典拓扑中区间可缩的性质。证明：对于任意 \(x \in I\)，由路径归纳我们可以构造从 \(x\) 到 \(\mathsf{left}\) 的路径，因此区间中所有点路径相等，从而 \(I\) 的类型可缩。
\item
  \textbf{球面 \(S^n\)}：由一个点构造器 \(\mathsf{base} : S^n\) 和一个 \(n\)-维路径构造器 \(\mathsf{loop} : \operatorname{Id}^n(*, *)\)（即 \(n\) 重环路空间中的非平凡元素）归纳定义。这一定义使得我们可以直接在 HoTT 中证明 \(\pi_k(S^n)\) 的基本性质------例如 \(n=1\) 情形 \(\pi_1(S^1) \cong \mathbb{Z}\) 可通过万有覆叠的构造直接证明。
\item
  \textbf{余等化子 / 推挤}：给定两个函数 \(f, g: A \to B\)，余等化子 \(\operatorname{Coeq}(f,g)\) 由一个点构造器 \(\mathsf{inl}: B \to \operatorname{Coeq}(f,g)\) 和一个路径构造器 \(\mathsf{glue}: \prod_{a:A} \operatorname{Id}(\mathsf{inl}(f(a)), \mathsf{inl}(g(a)))\) 定义。这恰好对应于经典拓扑中推挤 \(B \coprod_A [0,1]\) 的拓扑结构，其中每个 \(a \in A\) 粘合了 \(f(a)\) 和 \(g(a)\)。
\item
  \textbf{商类型}：当 \(A\) 上有等价关系 \(\sim\)，商类型 \(A/{\sim}\) 是高阶归纳类型，由一个构造器 \(\mathsf{class}: A \to A/{\sim}\) 和一个路径构造器 \(\prod_{x,y:A} (x \sim y) \to \operatorname{Id}(\mathsf{class}(x), \mathsf{class}(y))\) 定义。这一构造比单纯集合论中的等价类更加符合直觉，并且保持了高阶同伦信息（路径之间还可以有更高阶路径）。
\end{enumerate}

高阶归纳类型的核心优势在于，它允许我们直接在类型论内部讨论拓扑空间（同伦型）的性质，而无需依赖外部的集合论模型。每个胞腔的粘合过程被直接编码为路径构造器，归纳原理自动成为类型论的推理规则。例如，要证明一个命题 \(P(x)\) 对所有 \(x \in S^1\) 成立，只需给出 \(P(\mathsf{base})\) 和 \(P(\mathsf{base})\) 沿环路 \(\mathsf{loop}\) 的传输，这正是几何中沿着回路进行归纳的直接形式化。

\textbf{定义 11}（悬浮与胞腔类型的 HIT 定义）：类型 \(A\) 的\textbf{悬浮} \(\Sigma A\) 作为 HIT 定义为：点构造器 \(\mathsf{N}, \mathsf{S} : \Sigma A\)；路径构造器 \(\mathsf{merid}: A \to (\mathsf{N} = \mathsf{S})\)。此定义使得 \(\Sigma A\) 的行为恰如经典同伦论中的（约化）悬挂，将 \(A\) 映射为赤道上从北极到南极的子午线。

\textbf{定理 6}（Freudenthal 悬浮定理的 HoTT 证明，概要）：若 \(X\) 是 \(n\)-连通类型，则规范映射 \(X \to \Omega \Sigma X\) 是 \(2n\)-连通映射（诱导前 \(2n\) 个同伦群的同构）。HoTT 内部证明依赖于 Blakers-Massey 定理的高阶归纳版本，通过 pushout 的连通性估计实现。

\textbf{证明}：规范映射 \(\eta: X \to \Omega \Sigma X\) 将 \(x\) 映为环路 \(\mathsf{merid}(x) \cdot \mathsf{merid}(\mathsf{N})^{-1}\)。需证对 \(k \leq 2n\)，\(\pi_k(X) \to \pi_k(\Omega \Sigma X) \cong \pi_{k+1}(\Sigma X)\) 是同构。HoTT 中 Blakers-Massey 定理：对 pushout P，若源映射的纤维为 \(m\)-连通和 \(n\)-连通，则规范映射的纤维为 \((m+n)\)-连通。应用于 \(\Sigma X\) 的 pushout 构造，由 \(X\) 的 \(n\)-连通性推出 \(\eta\) 的纤维为 \((2n-1)\)-连通，故 \(\eta\) 诱导 \(\pi_k\) 同构对所有 \(k \leq 2n\)。\(\blacksquare\)

\subsection{单值公理的计算后果}\label{ux5355ux503cux516cux7406ux7684ux8ba1ux7b97ux540eux679c}

\textbf{定理 7}（函数外延性）：单值公理蕴含函数外延性公理：对 \(f,g: \prod_{x:A} B(x)\)，若 \(\prod_{x:A} \operatorname{Id}_{B(x)}(f(x), g(x))\) 有项，则 \(\operatorname{Id}_{\prod_{x:A} B(x)}(f,g)\) 有项。即，逐点相等的函数命题地相等。

\textbf{证明}：利用单值公理将逐点相等的依赖函数族转化为依赖等价，通过 \(\Sigma\)-类型的 path 编码构造 \(f\) 与 \(g\) 之间的恒等类型项。具体地，构造等价 \(e_x: f(x) = g(x)\)，令 \(P_x\) 为该等价对应的恒等项（通过 \(\operatorname{idToEquiv}^{-1}\)），则 \(\lambda x. e_x\) 提供了 \(f\) 与 \(g\) 在所有点处的''路径族''。对 \(\Pi\)-类型的路径归纳将此族提升为 \(f\) 与 \(g\) 在 \(\Pi\)-类型中的恒等。\(\blacksquare\)

\newpage

\chapter{07-拓扑与几何/02-微分几何/01-曲线与曲面.md}\label{ux62d3ux6251ux4e0eux51e0ux4f5502-ux5faeux5206ux51e0ux4f5501-ux66f2ux7ebfux4e0eux66f2ux9762.md}

\section{第218章：曲线论}\label{ux7b2c218ux7ae0ux66f2ux7ebfux8bba}

曲线论研究 \(\mathbb{R}^3\) 中曲线的局部几何性质，核心工具是 Frenet 标架和曲率不变量。

\subsection{62.1 参数化曲线}\label{ux53c2ux6570ux5316ux66f2ux7ebf}

\textbf{定义 62.1}（参数化曲线）：设 \(I \subseteq \mathbb{R}\) 为区间。一个\textbf{参数化曲线}是一个光滑（\(C^\infty\)）映射 \(\gamma: I \to \mathbb{R}^3\)。\(\gamma'(t)\) 称为曲线的\textbf{速度向量}，\(\|\gamma'(t)\|\) 称为\textbf{速率}。

\textbf{定义 62.2}（正则曲线）：曲线 \(\gamma\) 称为\textbf{正则的}，如果对所有 \(t \in I\)，\(\gamma'(t) \neq 0\)。

\textbf{定义 62.3}（重新参数化）：设 \(\gamma: I \to \mathbb{R}^3\) 是正则曲线，\(\varphi: J \to I\) 是光滑的微分同胚（\(\varphi' \neq 0\)）。则 \(\tilde{\gamma} = \gamma \circ \varphi: J \to \mathbb{R}^3\) 称为 \(\gamma\) 的\textbf{重新参数化}。

\textbf{定义 62.4}（弧长参数）：正则曲线 \(\gamma: I \to \mathbb{R}^3\) 的\textbf{弧长函数}定义为

\[s(t) = \int_{t_0}^t \|\gamma'(u)\| du\]

若 \(\|\gamma'(s)\| = 1\) 对所有 \(s\) 成立，则称 \(\gamma\) 以\textbf{弧长}为参数。此时 \(\gamma'(s)\) 是曲线在 \(\gamma(s)\) 处的\textbf{单位切向量}，记作 \(T(s)\)。

\textbf{命题 62.1}：每条正则曲线都可以重新参数化为弧长参数。

\emph{证明}：由于 \(\gamma\) 正则，\(\|\gamma'(t)\| > 0\)，故 \(s(t)\) 严格单调递增，存在光滑逆函数 \(t(s)\)。令 \(\tilde{\gamma}(s) = \gamma(t(s))\)，则 \(\|\tilde{\gamma}'(s)\| = \|\gamma'(t(s)) \cdot t'(s)\| = \|\gamma'(t(s))\| \cdot (1/\|\gamma'(t(s))\|) = 1\)。∎

\subsection{62.2 Frenet 标架}\label{frenet-ux6807ux67b6}

以下设 \(\gamma(s)\) 以弧长为参数，\(T(s) = \gamma'(s)\) 是单位切向量。

\textbf{定义 62.5}（曲率）：\(\gamma\) 在 \(s\) 处的\textbf{曲率}定义为 \(\kappa(s) = \|T'(s)\| = \|\gamma''(s)\|\)。

\textbf{定义 62.6}（主法向量）：若 \(\kappa(s) > 0\)，定义\textbf{主法向量} \(N(s) = T'(s) / \kappa(s)\)。则 \(N(s)\) 是单位向量且垂直于 \(T(s)\)（因为 \(\|T(s)\| = 1\) 蕴含 \(T'(s) \cdot T(s) = 0\)）。

\textbf{定义 62.7}（副法向量与挠率）：定义\textbf{副法向量} \(B(s) = T(s) \times N(s)\)。则 \(\{T(s), N(s), B(s)\}\) 构成 \(\mathbb{R}^3\) 的一组正交规范基，称为 \(\gamma\) 在 \(s\) 处的\textbf{Frenet 标架}。

\textbf{挠率} \(\tau(s)\) 定义为 \(\tau(s) = -B'(s) \cdot N(s)\)。

\textbf{命题 62.2}：\(B'(s) = -\tau(s) N(s)\)。

\emph{证明}：由 \(B(s) \cdot B(s) = 1\) 得 \(B'(s) \cdot B(s) = 0\)。又 \(B = T \times N\)，故 \(B' = T' \times N + T \times N' = \kappa N \times N + T \times N' = T \times N'\)，因此 \(B' \cdot T = 0\)。于是 \(B'\) 平行于 \(N\)，即 \(B'(s) = -\tau(s) N(s)\)。∎

\textbf{定理 62.3}（Frenet-Serret 公式）：

\[\begin{aligned}
T'(s) &= \kappa(s) N(s) \\
N'(s) &= -\kappa(s) T(s) + \tau(s) B(s) \\
B'(s) &= -\tau(s) N(s)
\end{aligned}\]

\emph{证明}：第一式 \(T'(s) = \kappa(s) N(s)\) 是曲率的定义：因为 \(\kappa(s) = \|T'(s)\|\) 且 \(N(s) = T'(s)/\kappa(s)\)（当 \(\kappa(s) > 0\)），故 \(T'(s) = \kappa(s) N(s)\)。第三式 \(B'(s) = -\tau(s) N(s)\) 是挠率的定义：由命题 62.2，\(B'(s)\) 平行于 \(N(s)\)，定义 \(\tau(s) = -B'(s) \cdot N(s)\)，即得 \(B'(s) = -\tau(s) N(s)\)。对第二式，利用 \(N = B \times T\)（由 \(B = T \times N\) 取向量积得 \(T \times N \times T = B \times T\)，即 \(N = B \times T\)），求导： \[N' = B' \times T + B \times T' = (-\tau N) \times T + B \times (\kappa N) = -\tau (N \times T) + \kappa (B \times N)\] 利用向量积恒等式 \(N \times T = -B\) 和 \(B \times N = -T\)，得 \(N' = -\tau (-B) + \kappa (-T) = \tau B - \kappa T = -\kappa T + \tau B\)。\(\blacksquare\)

\textbf{定理 62.4}（曲线论基本定理）：给定光滑函数 \(\kappa(s) > 0\) 和 \(\tau(s)\)（\(s \in I\)），则存在唯一（在刚体运动意义下）的以弧长为参数的正则曲线 \(\gamma: I \to \mathbb{R}^3\)，以 \(\kappa(s)\) 为曲率、\(\tau(s)\) 为挠率。

\emph{证明}：Frenet-Serret 公式构成关于 \((T, N, B)\) 的线性 ODE 系统。具体地，记 \(F(s) = (T(s), N(s), B(s))^\top\)，则 \(F'(s) = A(s) F(s)\)，其中 \(A(s)\) 是由 \(\kappa(s)\) 和 \(\tau(s)\) 唯一确定的反对称矩阵。由线性 ODE 的存在唯一性定理，给定初始条件 \(F(s_0)\) 为一个正交规范标架，在区间 \(I\) 上存在唯一解 \(F(s)\)。解 \(F(s)\) 自动保持正交规范性（因为 \(A(s)\) 反对称，\((F^\top F)' = 0\)，故 \(F^\top F\) 为常数）。定义 \(\gamma(s) = \gamma(s_0) + \int_{s_0}^s T(u) du\)。由 \(T(s) = \gamma'(s)\) 且 \(\|T(s)\| = 1\)，\(\gamma\) 以弧长为参数。由构造知 \(\gamma''(s) = T'(s) = \kappa(s) N(s)\)，故 \(\kappa(s) = \|\gamma''(s)\|\) 恰为曲率，且 \(N(s)\) 为主法向量。不同的初始条件对应于曲线的刚体运动。\(\blacksquare\)

\textbf{推论 62.5}：曲线的曲率和挠率是刚体运动不变量，且完全决定了曲线的形状。

\subsection{62.3 平面曲线}\label{ux5e73ux9762ux66f2ux7ebf}

\textbf{定义 62.8}（平面曲线）：若 \(\gamma\) 的像落在某个平面内，则 \(\gamma\) 称为\textbf{平面曲线}。

\textbf{命题 62.6}：\(\gamma\) 是平面曲线当且仅当 \(\tau(s) \equiv 0\)。

\emph{证明}：(\(\Rightarrow\)) 若 \(\gamma\) 是平面曲线，则 \(T(s)\) 和 \(N(s)\) 始终在该平面内，\(B(s)\) 是常值单位法向量，故 \(B'(s) = 0\)，从而 \(\tau(s) = 0\)。

(\(\Leftarrow\)) 若 \(\tau \equiv 0\)，则 \(B(s)\) 是常值向量 \(B_0\)。由 \((T \cdot B_0)' = T' \cdot B_0 = \kappa N \cdot B_0 = 0\)，故 \(T(s) \cdot B_0\) 为常数，即 \(\gamma(s) \cdot B_0\) 为常数，因此 \(\gamma\) 落在与 \(B_0\) 垂直的平面内。∎

\textbf{定义 62.9}（平面曲线的符号曲率）：对平面曲线，可定义\textbf{符号曲率} \(k(s)\)：\(T'(s) = k(s) N(s)\)。\(k(s) > 0\) 表示曲线向左转，\(k(s) < 0\) 表示向右转。

\textbf{定理 62.7}（等周不等式 / Isoperimetric Inequality）：在所有给定长度 \(L\) 的简单闭平面曲线中，圆所围成的面积最大。具体地，若 \(\gamma\) 是长度为 \(L\) 的简单闭曲线，其所围面积为 \(A\)，则

\[L^2 \geq 4\pi A\]

等号成立当且仅当 \(\gamma\) 是圆。

\textbf{证明}：设 \(\gamma(s) = (x(s), y(s))\) 是以弧长为参数的简单闭曲线，\(s \in [0, L]\)。由 Green 定理，\(A = \int_0^L x(s) y'(s) ds\)。Wirtinger 不等式的核心：对任意以 \(2\pi\) 为周期的光滑函数 \(f\) 满足 \(\int_0^{2\pi} f(t) dt = 0\)，有

\[\int_0^{2\pi} f(t)^2 dt \leq \int_0^{2\pi} f'(t)^2 dt\]

等号成立当且仅当 \(f(t) = a \cos t + b \sin t\)。

取参数 \(t = 2\pi s / L\)，则 \(s = Lt/(2\pi)\)。利用 Fourier 级数展开：\(x'(s) = \cos \theta(s)\)，\(y'(s) = \sin \theta(s)\)，其中 \(\theta(s)\) 是切向量与 \(x\) 轴的夹角。由闭曲线的性质，\(\theta(L) - \theta(0) = 2\pi\)（绕数定理）。定义 \(\tilde{\theta}(t) = \theta(Lt/(2\pi)) - t\)，则 \(\tilde{\theta}\) 以 \(2\pi\) 为周期。计算：

\[L = \int_0^L ds = \int_0^L (\cos^2\theta + \sin^2\theta) ds\]

\[A = \frac{1}{2}\int_0^L (x dy - y dx) = \frac{1}{2}\int_0^L (x(s) y'(s) - y(s) x'(s)) ds\]

由 Hurwitz 的证明：设曲线可由等周变分方法处理。更直接的证明是使用 Fourier 级数。将 \(x(s)\) 和 \(y(s)\) 展开为 Fourier 级数：

\[x(s) = \frac{a_0}{2} + \sum_{n=1}^\infty \left(a_n \cos\frac{2\pi n s}{L} + b_n \sin\frac{2\pi n s}{L}\right)\]

\[y(s) = \frac{c_0}{2} + \sum_{n=1}^\infty \left(c_n \cos\frac{2\pi n s}{L} + d_n \sin\frac{2\pi n s}{L}\right)\]

则 \(x'(s)\) 和 \(y'(s)\) 的 Fourier 系数为上述系数乘以 \(2\pi n/L\)。由 Parseval 恒等式和曲线以弧长为参数的条件 \(\|x'(s)\|^2 + \|y'(s)\|^2 = 1\)：

\[\frac{L}{2\pi}\int_0^L ((x')^2 + (y')^2) ds = \frac{L}{2\pi} \cdot L = \frac{L^2}{2\pi}\]

另一方面，面积 \(A = \int_0^L x(s) y'(s) ds\) 在 Fourier 级数下展开：

\[A = \pi \sum_{n=1}^\infty n(a_n d_n - b_n c_n)\]

由 \(n \leq n^2\)（对 \(n \geq 1\)）和 Cauchy-Schwarz 不等式：

\[|A| \leq \pi \sum_{n=1}^\infty n(|a_n||d_n| + |b_n||c_n|) \leq \pi \sum_{n=1}^\infty n^2\left(\frac{a_n^2 + b_n^2 + c_n^2 + d_n^2}{2}\right)\]

同时由 \(\| \gamma' \|^2 = 1\) 及 Parseval 等式：

\[\frac{L^2}{2\pi} = \sum_{n=1}^\infty n^2(a_n^2 + b_n^2 + c_n^2 + d_n^2)\]

故 \(A \leq \pi \cdot \frac{L^2}{2\pi} / 2 \cdot 2 = L^2/(4\pi)\)，即 \(L^2 \geq 4\pi A\)。等号成立时仅 \(n=1\) 项非零，这恰为圆的 Fourier 展开（椭圆情形------但由等周条件需为正圆）。\(\blacksquare\)

\textbf{定理 62.8}（四顶点定理 / Four-Vertex Theorem）：任意简单闭的严格凸平面曲线至少有四个顶点（即曲率的局部极值点）。

\textbf{证明思路}：设 \(\gamma\) 为以弧长为参数、反时针定向的简单闭凸曲线，曲率 \(k(s) > 0\)。顶点即 \(k'(s) = 0\) 的点。

第一步（至少两个顶点）：紧致集上的连续函数 \(k(s)\) 必达最大值和最小值，且在凸曲线上 \(k(s) > 0\)，故至少存在一个最大值点和一个最小值点，即至少两个顶点。

第二步（至少四个顶点------核）：假设仅有且恰有两个顶点 \(s_1, s_2\)，则 \(k'(s)\) 在 \((s_1, s_2)\) 和 \((s_2, s_1)\) 上分别保持常号。考虑参数化 \(\gamma(s) = (x(s), y(s))\)。对平面曲线，由 Frenet 公式在平面曲线情形有：

\[x'(s) = \cos\theta(s), \quad y'(s) = \sin\theta(s), \quad \theta'(s) = k(s)\]

其中 \(\theta(s)\) 是切向量与 \(x\) 轴的夹角。\(\theta(s+L) = \theta(s) + 2\pi\)（绕数定理，对简单闭曲线为 1）。

考虑积分 \(\int_0^L k'(s) x(s) ds\) 和 \(\int_0^L k'(s) y(s) ds\)。由分部积分和 Frenet 关系的推导得：

\[\int_0^L k'(s) ds = 0, \quad \int_0^L k'(s) x(s) ds = \int_0^L k'(s) y(s) ds = 0\]

现在若只有两个顶点，设 \(k'(s) \geq 0\) 在 \([s_1, s_2]\)，\(k'(s) \leq 0\) 在 \([s_2, s_1]\)（或符号取反）。考虑支撑线：对凸闭曲线，存在一条直线（支撑线）将曲线限制在其一侧。但 \(\int_0^L k'(s) \langle \gamma(s), \mathbf{v} \rangle ds\) 对某方向 \(\mathbf{v}\) 可判断其正负。通过选取适当的坐标系使 \(\int_0^L k'(s) x(s) ds\) 和 \(\int_0^L k'(s) y(s) ds\) 与顶点的分布矛盾------具体地，由凸性存在唯一的外法向方向使 \(k'(s) \cdot \langle \gamma(s), \mathbf{n} \rangle\) 积分非零，这与 \(\int k'(s) \gamma(s) ds = 0\) 矛盾。

该定理由 Mukhopadhyaya (1909) 首次对严格凸曲线证明，后被推广到一般简单闭曲线。\(\blacksquare\)

\hrulefill

\hrulefill

\hrulefill

\hrulefill

\hrulefill

\hrulefill

\section{第219章：曲面论}\label{ux7b2c219ux7ae0ux66f2ux9762ux8bba}

曲面论是经典微分几何的核心，研究 \(\mathbb{R}^3\) 中二维曲面的内蕴几何与外蕴几何。

\subsection{63.1 正则曲面}\label{ux6b63ux5219ux66f2ux9762}

正则曲面是古典微分几何的基本研究对象，它推广了曲线的概念，将二维参数化与三维嵌入统一起来。正则性条件 \(\mathbf{x}_u \times \mathbf{x}_v \neq 0\) 在几何上保证了曲面在每点处存在一个良定义的二维切平面------即曲面在局部上没有''尖点''、``褶皱''或''自交''。这一条件等价于 Jacobi 矩阵 \(D\mathbf{x}(u,v)\) 的秩为 2，从而 \(\mathbf{x}\) 在局部是浸入（immersion）。从拓扑角度看，正则曲面是 \(\mathbb{R}^2\) 的局部浸入像，且进一步要求 \(\mathbf{x}\) 是到像集上的同胚（整体定义中的条件），这排除了自交现象。正则曲面的整体定义（定义 63.3）本质上要求曲面是一个二维光滑流形，其光滑结构由局部参数化 \(\mathbf{x}: U \to V \cap S\) 给出的坐标卡（局部坐标）构成。这一观点将古典曲面论与现代微分几何统一起来：曲面上的每一个点 \(p\) 都有局部坐标 \((u, v)\)，而转移函数由全纯或光滑的坐标卡变换给出。正则曲面的经典例子包括：球面 \(S^2 = \{(x,y,z) : x^2+y^2+z^2 = 1\}\)（可通过球极投影参数化）、环面（通过旋转圆周生成）、旋转曲面（由平面曲线绕轴旋转生成）和极小曲面（如悬链面和螺旋面）。正则曲面论的三大基本形式------第一基本形式（度量结构）、第二基本形式（弯曲结构）和第三基本形式（球面映射的度量）------完整刻画了曲面在 \(\mathbb{R}^3\) 中的内蕴几何和外蕴几何。

\textbf{定义 63.1}（参数化曲面）：设 \(U \subseteq \mathbb{R}^2\) 为开集。一个\textbf{参数化曲面}是一个光滑映射 \(\mathbf{x}: U \to \mathbb{R}^3\)，\(\mathbf{x}(u, v) = (x(u, v), y(u, v), z(u, v))\)。

\textbf{定义 63.2}（正则曲面）：参数化曲面 \(\mathbf{x}\) 称为\textbf{正则的}，如果对每个 \((u, v) \in U\)，偏导数 \(\mathbf{x}_u\) 和 \(\mathbf{x}_v\) 线性无关（即 \(\mathbf{x}_u \times \mathbf{x}_v \neq 0\)）。此时 \(\mathbf{x}_u\) 和 \(\mathbf{x}_v\) 张成曲面在 \(\mathbf{x}(u, v)\) 处的\textbf{切平面}。

\textbf{定义 63.3}（正则曲面------整体定义）：\(\mathbb{R}^3\) 的子集 \(S\) 称为\textbf{正则曲面}，如果对每个 \(p \in S\)，存在 \(p\) 在 \(\mathbb{R}^3\) 中的邻域 \(V\) 和映射 \(\mathbf{x}: U \to V \cap S\)（\(U \subseteq \mathbb{R}^2\) 开集），使得 \(\mathbf{x}\) 是光滑同胚且是正则参数化曲面。\(\mathbf{x}\) 称为 \(S\) 在 \(p\) 处的\textbf{局部参数化}（或\textbf{坐标卡}）。

\subsection{63.2 第一基本形式}\label{ux7b2cux4e00ux57faux672cux5f62ux5f0f}

\textbf{定义 63.4}（第一基本形式）：设 \(\mathbf{x}: U \to S\) 是正则曲面的局部参数化。\(S\) 的\textbf{第一基本形式}是二次型

\[I_p(w) = \langle w, w \rangle_p = \|w\|^2, \quad w \in T_p S\]

其中 \(T_p S\) 是 \(S\) 在 \(p\) 处的切平面，\(\langle \cdot, \cdot \rangle_p\) 是 \(\mathbb{R}^3\) 的内积限制到 \(T_p S\)。

在局部坐标 \((u, v)\) 下，切向量 \(w = a \mathbf{x}_u + b \mathbf{x}_v\) 的第一基本形式为

\[I_p(w) = E a^2 + 2F a b + G b^2\]

其中

\[E = \langle \mathbf{x}_u, \mathbf{x}_u \rangle, \quad F = \langle \mathbf{x}_u, \mathbf{x}_v \rangle, \quad G = \langle \mathbf{x}_v, \mathbf{x}_v \rangle\]

称为第一基本形式的\textbf{系数}。

\textbf{定义 63.5}（等距）：两个正则曲面 \(S_1, S_2\) 称为\textbf{等距的}（或\textbf{局部等距的}），如果存在微分同胚 \(\varphi: S_1 \to S_2\) 使得对任意 \(p \in S_1\) 和 \(w \in T_p S_1\)，\(I_p^1(w) = I_{\varphi(p)}^2(d\varphi_p(w))\)。即第一基本形式被保持。

\textbf{命题 63.1}：第一基本形式完全决定了曲面上的弧长、角度和面积：

\begin{itemize}
\tightlist
\item
  曲线 \(\gamma(t) = \mathbf{x}(u(t), v(t))\) 的弧长：\(L = \int \sqrt{E u'^2 + 2F u' v' + G v'^2} dt\)
\item
  两切向量的夹角：\(\cos \theta = \dfrac{E a_1 a_2 + F(a_1 b_2 + a_2 b_1) + G b_1 b_2}{\sqrt{E a_1^2 + 2F a_1 b_1 + G b_1^2} \sqrt{E a_2^2 + 2F a_2 b_2 + G b_2^2}}\)
\item
  区域 \(R \subseteq U\) 的面积：\(A = \iint_R \sqrt{EG - F^2} du dv\)
\end{itemize}

\subsection{63.3 第二基本形式与法曲率}\label{ux7b2cux4e8cux57faux672cux5f62ux5f0fux4e0eux6cd5ux66f2ux7387}

\textbf{定义 63.6}（单位法向量）：对正则参数化曲面 \(\mathbf{x}: U \to S\)，定义单位法向量

\[N(p) = \frac{\mathbf{x}_u \times \mathbf{x}_v}{\|\mathbf{x}_u \times \mathbf{x}_v\|}(p)\]

（注意 \(N\) 取决于 \(\mathbf{x}_u \times \mathbf{x}_v\) 的方向，即曲面的定向。）

\textbf{定义 63.7}（Gauss 映射）：映射 \(N: S \to S^2\)（\(S^2\) 是 \(\mathbb{R}^3\) 中的单位球面）将曲面上每点映到其单位法向量，称为曲面的 \textbf{Gauss 映射}。

\textbf{定义 63.8}（第二基本形式）：\(S\) 的\textbf{第二基本形式}定义为

\[II_p(w) = -\langle dN_p(w), w \rangle, \quad w \in T_p S\]

在局部坐标下，\(dN_p(\mathbf{x}_u) = N_u\)，\(dN_p(\mathbf{x}_v) = N_v\)。第二基本形式的系数为

\[e = -\langle N_u, \mathbf{x}_u \rangle = \langle N, \mathbf{x}_{uu} \rangle\] \[f = -\langle N_u, \mathbf{x}_v \rangle = \langle N, \mathbf{x}_{uv} \rangle\] \[g = -\langle N_v, \mathbf{x}_v \rangle = \langle N, \mathbf{x}_{vv} \rangle\]

（后一等式由 \(\langle N, \mathbf{x}_u \rangle = 0\) 求导得到。）

\textbf{定义 63.9}（法曲率）：设 \(v \in T_p S\) 是单位切向量。曲面在 \(p\) 处沿方向 \(v\) 的\textbf{法曲率}为

\[\kappa_n(v) = II_p(v) = e a^2 + 2f a b + g b^2\]

其中 \(v = a \mathbf{x}_u + b \mathbf{x}_v\) 且 \(I_p(v) = 1\)。

\textbf{定义 63.10}（主曲率、Gauss 曲率、平均曲率）：法曲率 \(\kappa_n(v)\) 作为 \(v\) 的函数（在单位切向量上）在 \(T_p S\) 上取得最大值 \(\kappa_1\) 和最小值 \(\kappa_2\)，称为 \(S\) 在 \(p\) 处的\textbf{主曲率}。对应方向称为\textbf{主方向}。

\begin{itemize}
\tightlist
\item
  \textbf{Gauss 曲率}：\(K = \kappa_1 \kappa_2 = \dfrac{eg - f^2}{EG - F^2}\)
\item
  \textbf{平均曲率}：\(H = \dfrac{\kappa_1 + \kappa_2}{2} = \dfrac{eG - 2fF + gE}{2(EG - F^2)}\)
\end{itemize}

\textbf{定理 63.2}（主曲率的欧拉公式）：设 \(v\) 与主方向 \(e_1\) 的夹角为 \(\theta\)，则

\[\kappa_n(v) = \kappa_1 \cos^2 \theta + \kappa_2 \sin^2 \theta\]

\emph{证明}：设 \(e_1, e_2\) 为 \(T_p S\) 的单位主方向，对应的主曲率为 \(\kappa_1, \kappa_2\)。由主方向的定义，\(dN_p(e_1) = -\kappa_1 e_1\)，\(dN_p(e_2) = -\kappa_2 e_2\)。在由 \(e_1, e_2\) 构成的正交规范基下，\(dN_p\) 的矩阵表示为 \(\operatorname{diag}(-\kappa_1, -\kappa_2)\)。设 \(v \in T_p S\) 为单位切向量，\(v\) 与 \(e_1\) 的夹角为 \(\theta\)，则 \(v = \cos\theta e_1 + \sin\theta e_2\)。由第二基本形式的定义 \(II_p(v) = -\langle dN_p(v), v\rangle\)，代入得： \[II_p(v) = -\langle dN_p(\cos\theta e_1 + \sin\theta e_2), \cos\theta e_1 + \sin\theta e_2\rangle = -\langle -\kappa_1\cos\theta e_1 - \kappa_2\sin\theta e_2, \cos\theta e_1 + \sin\theta e_2\rangle\] \[= \kappa_1\cos^2\theta\langle e_1,e_1\rangle + \kappa_2\sin^2\theta\langle e_2,e_2\rangle + (\kappa_1+\kappa_2)\cos\theta\sin\theta\langle e_1,e_2\rangle\] 由于 \(e_1 \perp e_2\) 且 \(\|e_1\| = \|e_2\| = 1\)，交叉项为零，故 \(\kappa_n(v) = II_p(v) = \kappa_1\cos^2\theta + \kappa_2\sin^2\theta\)。\(\blacksquare\)

\textbf{定义 63.11}（点分类）：曲面上点 \(p\) 称为 - \textbf{椭圆点}：\(K > 0\)（\(\kappa_1\) 和 \(\kappa_2\) 同号） - \textbf{双曲点}：\(K < 0\)（\(\kappa_1\) 和 \(\kappa_2\) 异号） - \textbf{抛物点}：\(K = 0\) 但 \(H \neq 0\)（一个主曲率为零） - \textbf{平点}：\(K = 0\) 且 \(H = 0\)（两个主曲率为零）

\subsection{63.4 Gauss 绝妙定理}\label{gauss-ux7eddux5999ux5b9aux7406}

\textbf{定理 63.3}（Gauss 绝妙定理 / Theorema Egregium）：曲面的 Gauss 曲率 \(K\) 仅由第一基本形式及其导数决定，即是曲面的\textbf{内蕴}不变量。

\textbf{证明}：由 Gauss 公式 \(\mathbf{x}_{ij} = \Gamma_{ij}^k \mathbf{x}_k + L_{ij} N\)，其中 Christoffel 符号 \[\Gamma_{ij}^k = \frac{1}{2} \sum_{\ell} g^{k\ell}(\partial_i g_{j\ell} + \partial_j g_{i\ell} - \partial_\ell g_{ij})\] 仅由第一基本形式系数 \(g_{11}=E, g_{12}=g_{21}=F, g_{22}=G\) 及其导数决定。Gauss 曲率 \(K = \frac{R_{1212}}{g_{11}g_{22} - g_{12}^2}\)，其中 \[R_{1212} = \frac{1}{2}\left(\frac{\partial^2 g_{11}}{\partial v^2} + \frac{\partial^2 g_{22}}{\partial u^2} - 2\frac{\partial^2 g_{12}}{\partial u \partial v}\right) + \sum_{p,q} g_{pq}(\Gamma_{22}^p \Gamma_{11}^q - \Gamma_{21}^p \Gamma_{21}^q)\] 等价地 \(R_{1212} = \partial_1\Gamma_{22}^1 - \partial_2\Gamma_{21}^1 + \Gamma_{22}^1\Gamma_{11}^1 + \Gamma_{22}^2\Gamma_{12}^1 - \Gamma_{21}^1\Gamma_{21}^1 - \Gamma_{21}^2\Gamma_{22}^1\)。Riemann 张量 \(R_{1212}\) 完全由 Christoffel 符号及其一阶导数表达，而 Christoffel 符号仅依赖于第一基本形式及其一阶导数。故 \(K\) 是曲面的内蕴不变量，不依赖于 \(\mathbb{R}^3\) 中的嵌入。\(\blacksquare\)

\textbf{推论 63.4}：等距的曲面具有相同的 Gauss 曲率。特别地，不存在从平面（\(K = 0\)）到球面（\(K = 1/R^2 > 0\)）的等距映射（即无法将地图无变形地绘制在平面上）。

\subsection{63.5 测地线与协变导数}\label{ux6d4bux5730ux7ebfux4e0eux534fux53d8ux5bfcux6570}

\textbf{定义 63.12}（测地曲率）：设曲线 \(\gamma(s) \subset S\) 以弧长为参数。\(\gamma''(s)\) 在切平面上的投影的模长称为 \(\gamma\) 在 \(S\) 上的\textbf{测地曲率} \(k_g\)。

\textbf{定义 63.13}（测地线）：曲面 \(S\) 上的曲线 \(\gamma\) 称为\textbf{测地线}，如果其测地曲率恒为零，即 \(\gamma''(s)\) 处处垂直于切平面。

\textbf{命题 63.5}（测地线方程）：在局部坐标 \((u, v)\) 下，测地线满足微分方程：

\[\frac{d^2 u^k}{ds^2} + \sum_{i,j=1}^2 \Gamma_{ij}^k \frac{du^i}{ds}\frac{du^j}{ds} = 0, \quad k = 1, 2\]

\emph{证明}：\(\gamma''(s)\) 垂直于 \(T_p S\) 等价于 \(\gamma''(s)\) 与 \(\mathbf{x}_u\) 和 \(\mathbf{x}_v\) 的内积为零。展开 \(\gamma''(s) = \frac{d^2 u^k}{ds^2} \mathbf{x}_k + \frac{du^i}{ds}\frac{du^j}{ds} \mathbf{x}_{ij}\)，利用 \(\mathbf{x}_{ij} = \Gamma_{ij}^k \mathbf{x}_k + L_{ij} N\)（Gauss 公式），取切向分量即得测地线方程。∎

\textbf{命题 63.6}：曲面上连接两点的最短曲线（若存在）必为测地线。

\hrulefill

\hrulefill

\hrulefill

\hrulefill

\hrulefill

\hrulefill

\hrulefill

\hrulefill

\section{第220章：Gauß-Bonnet定理（完整局部与全局证明）}\label{ux7b2c220ux7ae0gauuxdf-bonnetux5b9aux7406ux5b8cux6574ux5c40ux90e8ux4e0eux5168ux5c40ux8bc1ux660e}

Gauß-Bonnet 定理是微分几何的巅峰成就之一，它将曲面的局部几何（Gauß 曲率 \(K\) 和测地曲率 \(\kappa_g\)）与全局拓扑（Euler 示性数 \(\chi(M)\)）联系起来。定理的核心等式 \(\int_M K \, dA + \int_{\partial M} \kappa_g \, ds = 2\pi \chi(M)\) 之美在于：左边是完全由 Riemann 度量决定的微分几何量，右边是纯粹拓扑的不变量------两者之间的等号意味着曲率在任何连续形变（保持度规结构）下存在''刚性''的积分约束。Gauß（1827）在局部情形证明了测地三角形的角盈等于 Gauss 曲率的面积分 \(\int_\Delta K \, dA = \alpha + \beta + \gamma - \pi\)，这个局部版本在本质上等价于 Gauß 绝妙定理（\(K\) 是内蕴的）。Bonnet（1848）将局部版本推广到带有边界的一般区域，而全局版本（对无边紧致曲面 \(\int_M K \, dA = 2\pi \chi(M)\)）直到 19 世纪末才被完全确立。现代证明通常使用外微分形式和活动标架法，将被积函数写为联络形式的外导数，再通过 Stokes 定理转化为边界（或奇点）的积分。Chern（1944）将 Gauß-Bonnet 定理以内蕴方式推广到任意偶数维 Riemann 流形，得到了著名的 Chern-Gauß-Bonnet 定理。

\subsection{73.x 活动标架法与Gauß-Bonnet局部定理}\label{x-ux6d3bux52a8ux6807ux67b6ux6cd5ux4e0egauuxdf-bonnetux5c40ux90e8ux5b9aux7406}

\textbf{定义 73.1}（活动标架与结构方程）：在定向曲面 \(M \subset \mathbb{R}^3\) 上（或更一般地，带 Riemann 度量的抽象曲面），选取单位正交活动标架 \(\{\mathbf{e}_1, \mathbf{e}_2, \mathbf{e}_3\}\)，其中 \(\mathbf{e}_1, \mathbf{e}_2\) 切于 \(M\)，\(\mathbf{e}_3 = \mathbf{N}\) 为单位法向量。对偶 \(1\)-形式 \(\omega_i\)（\(\omega_i(\mathbf{e}_j) = \delta_{ij}\)）和联络 \(1\)-形式 \(\omega_{ij}\) 满足：

\[d\mathbf{e}_i = \sum_{j=1}^3 \omega_{ij} \mathbf{e}_j, \quad \omega_{ij} = -\omega_{ji}\]

Cartan 结构方程为： \[d\omega_i = \sum_j \omega_{ij} \wedge \omega_j \quad \text{（第一结构方程）}\] \[d\omega_{ij} = \sum_k \omega_{ik} \wedge \omega_{kj} \quad \text{（第二结构方程）}\]

对曲面，仅有唯一非平凡的曲率 \(2\)-形式 \(\Omega_{12} = d\omega_{12} = -K \, \omega_1 \wedge \omega_2\)，其中 \(K\) 为 Gauß 曲率。

\textbf{定理 73.1}（局部Gauß-Bonnet定理------测地曲率的积分公式）：设 \(D \subset M\) 是定向曲面 \(M\) 上的单连通区域，边界 \(\partial D\) 为分段光滑闭曲线，外角为 \(\varepsilon_1, \ldots, \varepsilon_k\)（即曲线在顶点处方向跳跃的角度，逆时针为正）。则

\[\int_D K \, dA + \int_{\partial D} \kappa_g \, ds + \sum_{j=1}^k \varepsilon_j = 2\pi\]

其中 \(\kappa_g\) 为边界曲线（相对于曲面定向）的测地曲率。

\textbf{证明}（活动标架法）：沿 \(\partial D\) 的光滑弧段，构造 Darboux 标架 \(\{\mathbf{T}, \mathbf{N} \times \mathbf{T}, \mathbf{N}\}\)，其中 \(\mathbf{T}\) 为单位切向量。令 \(\omega_{12} = \langle d\mathbf{e}_1, \mathbf{e}_2 \rangle\) 为联络 \(1\)-形式。在曲面上，沿边界拉回：令 \(\mathbf{e}_1 = \mathbf{T}\)，\(\mathbf{e}_2 = \mathbf{N} \times \mathbf{T}\)，则沿 \(\partial D\) 有

\[\kappa_g \, ds = \langle \nabla_{\mathbf{T}} \mathbf{T}, \mathbf{e}_2 \rangle = \omega_{12}(\mathbf{T}) \, ds\]

由 Stokes 定理（在 \(M\) 上）： \[\int_D d\omega_{12} = \int_{\partial D} \omega_{12}\]

将曲率 \(2\)-形式代入：\(d\omega_{12} = -K \, \omega_1 \wedge \omega_2 = -K \, dA\)（\(dA\) 为面积元）。故 \(\int_D K \, dA = -\int_D d\omega_{12} = -\int_{\partial D} \omega_{12}\)。

另一方面，沿 \(\partial D\) 的光滑弧段，设 \(\mathbf{T}\) 与某固定的参考切向量 \(\mathbf{e}_1\) 之间的夹角为 \(\theta\)（即 \(\mathbf{T} = \cos\theta \, \mathbf{e}_1 + \sin\theta \, \mathbf{e}_2\)）。沿边界拉回活动标架：若取曲面上某固定的活动标架场 \(\{\mathbf{e}_1, \mathbf{e}_2\}\) 并设 \(\tilde{\mathbf{e}}_1 = \mathbf{T} = \cos\theta \, \mathbf{e}_1 + \sin\theta \, \mathbf{e}_2\)，\(\tilde{\mathbf{e}}_2 = -\sin\theta \, \mathbf{e}_1 + \cos\theta \, \mathbf{e}_2\)，则 \(\tilde{\omega}_{12} = \langle d\tilde{\mathbf{e}}_1, \tilde{\mathbf{e}}_2 \rangle = d\theta + \langle d\mathbf{e}_1, \mathbf{e}_2 \rangle = d\theta + \omega_{12}\)。沿边界积分： \[\int_{\partial D} \tilde{\omega}_{12} = \int_{\partial D} d\theta + \int_{\partial D} \omega_{12}\]

\(\tilde{\omega}_{12}(\mathbf{T}) \, ds = \kappa_g \, ds\)，故 \(\int_{\partial D} \kappa_g \, ds = \int_{\partial D} d\theta + \int_{\partial D} \omega_{12}\)。

角度积分 \(\int_{\partial D} d\theta\) 度量了切向量沿闭合回路绕一圈的总旋转角度。对分段光滑闭曲线（外角 \(\varepsilon_j\)），\(\int_{\partial D} d\theta = 2\pi - \sum_j \varepsilon_j\)（Hopf 绕数定理 / 旋转指标定理：任意简单闭曲线的切向量总旋转角为 \(2\pi\)，减去各顶点的外角跳跃）。因此： \[\int_{\partial D} \kappa_g \, ds = 2\pi - \sum_j \varepsilon_j + \int_{\partial D} \omega_{12} = 2\pi - \sum_j \varepsilon_j - \int_D K \, dA\]

移项即得定理。\(\blacksquare\)

\subsection{73.y 全局Gauß-Bonnet定理}\label{y-ux5168ux5c40gauuxdf-bonnetux5b9aux7406}

\textbf{定理 73.2}（全局Gauß-Bonnet定理------Polyakov-Chern）：设 \(M\) 为紧致无边定向光滑曲面。则

\[\int_M K \, dA = 2\pi \chi(M)\]

其中 \(\chi(M)\) 为 Euler 示性数（对亏格 \(g\) 的可定向曲面，\(\chi(M) = 2 - 2g\)）。

\textbf{证明}（三角剖分法 + 局部定理的粘合）：将 \(M\) 三角剖分为 \(F\) 个面（三角形）、\(E\) 条边、\(V\) 个顶点。每个面 \(\Delta_i\)（定向与 \(M\) 一致）的边界为分段光滑闭曲线（三个弧 + 三个顶点）。对每个三角形应用局部 Gauß-Bonnet 定理（定理 73.1）： \[\int_{\Delta_i} K \, dA + \int_{\partial \Delta_i} \kappa_g \, ds + \sum_{j=1}^3 (\pi - \alpha_{ij}) = 2\pi\] 其中 \(\alpha_{ij}\) 为三角形 \(\Delta_i\) 的第 \(j\) 个内角，外角为 \(\pi - \alpha_{ij}\)。

对所有 \(F\) 个三角形求和： \[\int_M K \, dA + \sum_{i=1}^F \int_{\partial \Delta_i} \kappa_g \, ds + \sum_{i=1}^F \sum_{j=1}^3 (\pi - \alpha_{ij}) = 2\pi F\]

关键观察：每条内部边被两个相邻三角形共享，在其上的测地曲率积分因定向相反（边界曲线方向相逆）而精确相消。故 \(\sum_i \int_{\partial \Delta_i} \kappa_g \, ds = 0\)（无边界面，所有边为内部边）。角度项：\(\sum_i \sum_j (\pi - \alpha_{ij}) = \sum_i (3\pi - \sum_j \alpha_{ij}) = 3\pi F - \sum_i \sum_j \alpha_{ij}\)。所有内角之和为 \(\sum_i \sum_j \alpha_{ij} = 2\pi V\)（因围绕每个顶点的各三角形内角之和为 \(2\pi\)）。故角度项贡献为 \(3\pi F - 2\pi V\)。代入方程 \(2\pi F\)： \[\int_M K \, dA + 0 + (3\pi F - 2\pi V) = 2\pi F\] \[\implies \int_M K \, dA = 2\pi V - \pi F\]

利用 Euler 公式 \(\chi(M) = V - E + F\)。在三角剖分中，每条边在边界计数中被两个面共享，故 \(3F = 2E\)（每个面三条边，每条边两个面）。因此 \(E = 3F/2\)，代入：\(\chi = V - 3F/2 + F = V - F/2\)，故 \(2\pi \chi = 2\pi V - \pi F\)。结合上述积分等式： \[\int_M K \, dA = 2\pi \chi(M)\]

对带边曲面 \(M\)（\(\partial M \neq \varnothing\)），类似推导给出： \[\int_M K \, dA + \int_{\partial M} \kappa_g \, ds = 2\pi \chi(M)\] 其中 \(\chi(M)\) 仍为 Euler 示性数（\(\chi = V - E + F\) 对任意胞腔剖分）。\(\blacksquare\)

\textbf{推论 73.3}（Gauß-Bonnet定理的推论）：

\begin{enumerate}
\def\labelenumi{(\roman{enumi})}
\item
  亏格 \(g\) 的紧致可定向曲面上不可能处处 \(K > 0\)（当 \(g \geq 1\)），也不可能处处 \(K < 0\)（当 \(g = 0\)）。特别地，球面 \(S^2\)（\(g = 0\)）上 \(\int K \, dA = 4\pi\)，环面 \(T^2\)（\(g = 1\)）上 \(\int K \, dA = 0\)。
\item
  若 \(M\) 紧致无边且 \(K > 0\) 处处成立，则 \(M \cong S^2\)（微分同胚于球面）------这是 Gauß-Bonnet 与曲面分类定理结合的结论。但 Gauß-Bonnet 本身不蕴含微分同胚分类。
\item
  向量场的 Poincare-Hopf 指标定理：\(\int_M K \, dA = 2\pi \sum_p \operatorname{ind}_p(X)\)（\(X\) 的指标和），由联络形式的 Gauss 映射解释给出。
\end{enumerate}

\textbf{证明}：(i) 由 \(\int_M K \, dA = 2\pi (2-2g)\)。若 \(K > 0\) 处处，左端 \(> 0\)，故 \(2-2g > 0 \implies g = 0\)（球面）。若 \(K < 0\) 处处，左端 \(< 0\)，故 \(2-2g < 0 \implies g \geq 2\)，\(g=0\) 不可能。(ii) 参见曲面微分同胚分类（Morse 理论或 Ricci 流）。(iii) 联络形式 \(\omega_{12}\) 沿边界积分与 Gauss 映射（将切平面映至 \(S^2\)）的映射度相关，后者等于 Euler 示性数。\(\blacksquare\)

\subsection{73.z 活动标架法下的Gauß公式与Codazzi-Mainardi方程证明}\label{z-ux6d3bux52a8ux6807ux67b6ux6cd5ux4e0bux7684gauuxdfux516cux5f0fux4e0ecodazzi-mainardiux65b9ux7a0bux8bc1ux660e}

以下为第 73 章中 Gauß 绝妙定理和 Codazzi-Mainardi 方程的活动标架法补证。

\textbf{补证：Gauß绝妙定理的Gauß公式推导}。Gauß 绝妙定理（Theorema Egregium）断言：Gauß 曲率 \(K\) 仅由第一基本形式的系数及其导数决定（即内蕴几何量）。在活动标架语言中，其证明极为简洁。选取切标架 \(\{\mathbf{e}_1, \mathbf{e}_2\}\) 和法向量 \(\mathbf{e}_3 = \mathbf{N}\)，联络形式 \(\omega_{12}\) 满足第二结构方程 \(d\omega_{12} = -K \, \omega_1 \wedge \omega_2\)。但 \(\omega_{12}\) 本身仅由 \(\omega_1, \omega_2\) 和 \(d\omega_1, d\omega_2\) 决定------实际上，由第一结构方程 \(d\omega_1 = \omega_{12} \wedge \omega_2\) 和 \(d\omega_2 = \omega_{21} \wedge \omega_1 = -\omega_{12} \wedge \omega_1\)，可解得 \(\omega_{12} = p \omega_1 + q \omega_2\)，其中 \(p = \langle d[\omega_1^{-1}], \cdot \rangle\) 和 \(q\) 仅依赖于 \(\omega_1, \omega_2\) 的导数（即第一基本形式）。因此 \(d\omega_{12}\) 仅依赖于内蕴数据，从而 \(K\) 亦然。具体地，在坐标 \(u^1, u^2\) 下，\(\omega_i = a_{i1} du^1 + a_{i2} du^2\)（\(i=1,2\)），\(a_{ij}\) 为第一基本形式 \(I = E(du^1)^2 + 2F du^1 du^2 + G(du^2)^2\) 的 Gram 分解系数。由 Cartan 引理，\(\omega_{12}\) 由下列线性方程组唯一确定：\(d\omega_1 = \omega_{12} \wedge \omega_2\)，\(d\omega_2 = -\omega_{12} \wedge \omega_1\)。代入 \(\omega_{12} = \Gamma_{11}^2 \omega_1 + \Gamma_{12}^2 \omega_2\)（Christoffel 记号），解出 \(\Gamma_{ij}^k\) 为 \(a_{ij}\) 的组合。最终 \(K = -d\omega_{12} / (\omega_1 \wedge \omega_2)\) 是 \(a_{ij}\) 及其一阶、二阶导数的函数------即 Gauß 公式。\(\blacksquare\)

\textbf{补证：Codazzi-Mainardi方程的Frenet标架推导}。以标准坐标 \((u^1, u^2)\) 和坐标标架 \(\partial_1 = \partial/\partial u^1\)，\(\partial_2 = \partial/\partial u^2\)，法向量 \(\mathbf{N}\)。第二基本形式 \(II = L (du^1)^2 + 2M du^1 du^2 + N (du^2)^2\)，其中 \(L = \langle \partial_{11}^2 \mathbf{r}, \mathbf{N} \rangle\)，\(M = \langle \partial_{12}^2 \mathbf{r}, \mathbf{N} \rangle\)，\(N = \langle \partial_{22}^2 \mathbf{r}, \mathbf{N} \rangle\)。Weingarten 方程：\(\partial_i \mathbf{N} = -\sum_j h_i^j \partial_j \mathbf{r}\)，其中 \(h_i^j\) 为 Weingarten 映射（形状算子）的矩阵表示，由 \((h_i^j) = (II)(I)^{-1}\) 给出。Gauß 方程：\(\partial_{ij}^2 \mathbf{r} = \sum_k \Gamma_{ij}^k \partial_k \mathbf{r} + h_{ij} \mathbf{N}\)。由偏导数的对称性 \(\partial_1 (\partial_2 \mathbf{N}) = \partial_2 (\partial_1 \mathbf{N})\) 得：

\(\partial_1 (-\sum_j h_2^j \partial_j \mathbf{r}) = \partial_2 (-\sum_j h_1^j \partial_j \mathbf{r})\)。展开导数（利用 Gauß 方程和 Weingarten 方程），比较 \(\partial_k \mathbf{r}\) 的系数得：

\[\frac{\partial h_2^k}{\partial u^1} + \sum_j h_2^j \Gamma_{1j}^k = \frac{\partial h_1^k}{\partial u^2} + \sum_j h_1^j \Gamma_{2j}^k\]

此即 Codazzi 方程。用 \(L, M, N\) 和 Christoffel 记号表示为： \[L_v - M_u = L \Gamma_{12}^1 + M(\Gamma_{12}^2 - \Gamma_{11}^1) - N \Gamma_{11}^2\] \[M_v - N_u = L \Gamma_{22}^1 + M(\Gamma_{22}^2 - \Gamma_{12}^1) - N \Gamma_{12}^2\]

（\(u, v\) 为坐标参数，下标表示偏导）。证明由上述偏导交换性的显式展开完全形式化。\(\blacksquare\)

\textbf{补证：测地线方程的第二变分公式}。设 \(\gamma: [0, L] \to M\) 为测地线（\(\nabla_{\dot{\gamma}} \dot{\gamma} = 0\)），\(\Gamma(s, t)\) 为 \(\gamma\) 的变分（\(\Gamma(s, 0) = \gamma(s)\)，\(\Gamma(0, t) = \gamma(0)\)，\(\Gamma(L, t) = \gamma(L)\)）。能量泛函的第二变分为：

\[\frac{d^2}{dt^2} E(\Gamma(\cdot, t))\Big|_{t=0} = \int_0^L \left( \|\nabla_{\dot{\gamma}} V\|^2 - \langle R(V, \dot{\gamma}) \dot{\gamma}, V \rangle \right) ds\]

其中 \(V(s) = \partial_t \Gamma(s, 0)\) 为变分向量场（沿 \(\gamma\) 的 Jacobi 场，\(V(0) = V(L) = 0\)），\(R\) 为曲率张量。\(I(V, V) = \int_0^L (\|\nabla_{\dot{\gamma}} V\|^2 - K \|V\|^2) ds\)（曲面上 \(K\) 为 Gauß 曲率）称为指数形式（index form）。

\textbf{推导}：能量第一变分：\(dE/dt = \int_0^L \langle \nabla_t \partial_s \Gamma, \partial_s \Gamma \rangle ds\)。第二变分： \[\frac{d^2 E}{dt^2} = \int_0^L \left( \langle \nabla_t \nabla_t \partial_s \Gamma, \partial_s \Gamma \rangle + \langle \nabla_t \partial_s \Gamma, \nabla_t \partial_s \Gamma \rangle \right) ds\]

在 \(t=0\) 处，\(\nabla_t \partial_s \Gamma = \nabla_{\dot{\gamma}} V\)（由挠率为零：\(\nabla_{\partial_s} \partial_t = \nabla_{\partial_t} \partial_s\)）。第一项 \(\nabla_t \nabla_t \partial_s \Gamma\) 利用 Ricci 恒等式交换协变导数：\(\nabla_t \nabla_t \partial_s \Gamma = \nabla_t \nabla_s \partial_t \Gamma = \nabla_s \nabla_t \partial_t \Gamma + R(\partial_t, \partial_s) \partial_t \Gamma\)。在 \(t=0\) 处 \(\partial_t \Gamma = V\) 且 \(\partial_s \Gamma = \dot{\gamma}\)，故 \(\nabla_t \partial_t \Gamma|_{t=0} = \nabla_V V|_{\gamma} = 0\)（因测地线变分无加速度要求？这不必然为零；但在 \(t=0\) 附近可选取特殊的 \(\Gamma\) 使加速度为零）。更严谨的处理：对一般变分（不要求 \(\Gamma(\cdot, t)\) 均为测地线），需保留加速度项。但若仅考虑 \(\gamma\) 作为能量临界点（\(\dot{E}(0) = 0\) 因 \(\gamma\) 测地线），第二变分公式中加速度项贡献边界项 \(\int_0^L \partial_s \langle \nabla_t \partial_t \Gamma, \partial_s \Gamma \rangle ds\) 减去 \(\int_0^L \langle \nabla_t \partial_t \Gamma, \nabla_s \partial_s \Gamma \rangle ds\)。由于 \(\nabla_s \partial_s \Gamma|_{t=0} = \nabla_{\dot{\gamma}} \dot{\gamma} = 0\)（\(\gamma\) 测地线），后一项为零。前一项在 \(s=0, L\) 处因 \(V(0)=V(L)=0\) 亦为零。因此仅有曲率项和 \(\|\nabla_{\dot{\gamma}} V\|^2\) 项贡献，如所述公式。\(\blacksquare\)

\hrulefill

\hrulefill

\hrulefill

\section{第221章：等距嵌入与Nash嵌入定理}\label{ux7b2c221ux7ae0ux7b49ux8dddux5d4cux5165ux4e0enashux5d4cux5165ux5b9aux7406}

Nash 嵌入定理是微分几何中与 Gauß-Bonnet 定理齐名的伟大定理。John Nash（1956）证明了每个紧致（或更一般地，任意）Riemann 流形可以等距地嵌入到某个充分高维的 Euclidean 空间 \(\mathbb{R}^q\) 中。具体地，若 \((M^n, g)\) 是 \(C^\infty\) Riemann 流形（\(n \geq 2\)），则存在 \(C^1\) 等距嵌入 \(M \to \mathbb{R}^{n+1}\)（Nash-Kuiper，1954-1955）和 \(C^\infty\) 等距嵌入 \(M \to \mathbb{R}^{n(n+1)(3n+11)/2}\)（Nash 1956，后由 Gromov 改进为 \((n+2)(n+3)/2\) 或由 Gunther 进一步降低）。Nash 嵌入定理的深刻之处在于：弯曲的 Riemann 几何可以在更高维的平坦空间中作为子流形实现------这不仅为抽象 Riemann 流形提供了具体的几何直觉，而且在 PDE 视角下，等距嵌入归结为求解一个高度非线性的一阶 PDE 系统（Gauß-Codazzi-Mainardi-Ricci 方程作为可积性条件，但 Nash 嵌入的解不要求第二基本形式满足这些方程------嵌入是将给定的度量 \(g_{ij}\) 实现为诱导度量，而非给定第二基本形式）。Nash 的证明是分析学与几何学交叉的典范，核心思路是：从一个''粗略的''嵌入出发（可能不是等距的，但接近），通过添加''皱褶''（Nash twist / Nash perturbation）逐次精密矫正度量误差，利用 Newton 迭代的快速收敛达到精确等距性。Gromov（1986）将 Nash 的方法系统化为凸积分（convex integration）理论，成为几何分析中的强大工具。

\subsection{74.x Nash-Moser隐函数定理与等距嵌入}\label{x-nash-moserux9690ux51fdux6570ux5b9aux7406ux4e0eux7b49ux8dddux5d4cux5165}

\textbf{定义 74.1}（等距嵌入问题 / PDE形式）：给定 \((M^n, g)\) 为 Riemann 流形，等距嵌入问题等价于求映射 \(f: M \to \mathbb{R}^q\) 使得拉回 Euclidean 度量等于 \(g\)：\((f^* \delta_{\mathbb{R}^q})_{ij} = \langle \partial_i f, \partial_j f \rangle = g_{ij}\)（在局部坐标下，\(\delta_{\mathbb{R}^q}\) 为 \(\mathbb{R}^q\) 的标准平坦度量）。该 PDE 系统有 \(n(n+1)/2\) 个方程（\(g_{ij}\) 的独立分量数）和 \(q\) 个未知函数（\(f = (f^1, \ldots, f^q)\)）。若 \(q \geq n(n+1)/2\)，方程数与未知函数数相当（超定或恰当定）。

\textbf{定理 74.1}（Nash-Moser隐函数定理，Nash 1956；Moser 1961）：设 \((M^n, g_t)\) 是光滑的 Riemann 度量单参族（\(t \in [0,1]\)），\(g_0\) 可等距嵌入 \(\mathbb{R}^q\)。若 \(g_t - g_0\) 及其各阶导数充分小（在适当的 Frechet 空间拓扑下），则 \(g_t\)（对所有 \(t\)）也可等距嵌入 \(\mathbb{R}^q\)。证明依赖 Nash-Moser 快速迭代方案（Newton 法在 Frechet 空间中的推广，带光滑化步骤 / smoothing operators）。

\textbf{证明}（框架）：构造迭代序列 \(f_k: M \to \mathbb{R}^q\)（\(k = 0, 1, 2, \ldots\)）逼近等距嵌入。初始嵌入 \(f_0\) 满足 \((f_0^* \delta)_{ij} = g_{ij,0}\)（对初始度量 \(g_0\) 等距）。设第 \(k\) 步度量误差为 \(e_{ij}^{(k)} = g_{ij} - \langle \partial_i f_k, \partial_j f_k \rangle\)。对 \(f_k\) 施加校正 \(f_{k+1} = f_k + \Delta f_k\)，线性化误差方程给出：

\[\langle \partial_i f_k, \partial_j \Delta f_k \rangle + \langle \partial_j f_k, \partial_i \Delta f_k \rangle = e_{ij}^{(k)} - \langle \partial_i \Delta f_k, \partial_j \Delta f_k \rangle\]

忽略二次项（在 Newton 法中取线性化），需解线性方程：\(\operatorname{Sym}(\operatorname{Hess}(f_k, \Delta f_k)) = e^{(k)}\)。这是一个线性超定 PDE（\(n(n+1)/2\) 个方程，\(q\) 个未知 \(\Delta f_k^a\)）。Nash 的关键技巧：引入''皱褶''（Nash twist）\(\Delta f_k^a = \sum_{b} w_k^{ab}(x) \cos(\lambda_k x \cdot \xi_b)\)（高频振荡校正）。选取高频 \(\lambda_k \gg 1\)，使 \(\Delta f_k\) 的主贡献来自其导数项 \(\lambda_k^2\) 而非 \(\lambda_k^1\)。通过精心选取振幅 \(w_k^{ab}\) 和频率向量 \(\xi_b \in \mathbb{R}^n\)，可使其主项恰好满足线性化方程。\(\blacksquare\)

\subsection{74.y Nash嵌入定理的陈述与证明概要}\label{y-nashux5d4cux5165ux5b9aux7406ux7684ux9648ux8ff0ux4e0eux8bc1ux660eux6982ux8981}

\textbf{定理 74.2}（Nash 等距嵌入定理，1956）：设 \((M^n, g)\) 是 \(C^k\) Riemann 流形（\(k \geq 3\)）。则存在 \(C^k\) 等距嵌入 \(f: M \to \mathbb{R}^q\)，其中 \(q = \frac{1}{2} n(n+1)(3n+11)\)。对于 \(C^\infty\) 情形，\(q\) 可取同样上界。后期改进（Gromov-Rokhlin 1970，Gromov 1986，Gunther 1990）将 \(q\) 降低至 \(q = n(n+3)/2 + \max\{n, 5\}\)（或 \((n+2)(n+3)/2\) 对紧致情形），以及 Gromov 的软嵌入（\(q > n(n+1)/2\) 时存在局部自由嵌入，可用凸积分构造）。

\textbf{证明概要}（Nash 1956）：

步骤 1（``短嵌入'' / short embedding）：构造光滑浸入（不一定等距）\(f_0 : M \to \mathbb{R}^q\) 满足 \(\langle \partial_i f_0, \partial_j f_0 \rangle < g_{ij}\)（作为双线性型）。直观上，``短映射''意味着 \(f_0(M)\) 中所有曲线的长度均严格小于其在 \(M\) 中的 Riemann 长度------好比将一个皱巴巴的球面''压缩''后放入高维 Euclidean 空间。

步骤 2（Nash 扰动 / 修正）：令 \(h_{ij} = g_{ij} - \langle \partial_i f_0, \partial_j f_0 \rangle > 0\)（正定对称矩阵）。构造扰动 \(f_1 = f_0 + \frac{1}{\lambda} \mathbf{A} \sin(\lambda \phi)\) 或 \(f_1 = f_0 + \sum_{\alpha} \mathbf{w}_\alpha \cos(\lambda \xi_\alpha \cdot x)\)，其中 \(\lambda \to \infty\) 为大参数（频率），\(\mathbf{w}_\alpha\) 和 \(\xi_\alpha\) 的选择使得 \(\langle \partial_i f_1, \partial_j f_1 \rangle - \langle \partial_i f_0, \partial_j f_0 \rangle\)（在 \(\lambda \to \infty\) 极限下）逼近 \(h_{ij}\)。

步骤 3（Newton 迭代 + Moser 光滑化）：令 \(g^{(0)} = g\)，\(f_0\) 为短嵌入。度量余项 \(e_{ij}^{(k)} = g_{ij}^{(k)} - \langle \partial_i f_k, \partial_j f_k \rangle\)（目标调整）。Nash 构造校正 \(f_{k+1} = f_k + \Delta f_k\) 使得 \(\|e^{(k+1)}\| \leq C \|e^{(k)}\|^2\)（二次收敛），每个校正涉及高频扰动 + Moser 光滑化（防止导数爆炸）。无穷迭代的极限 \(f = \lim_{k \to \infty} f_k\) 是所求的 \(C^k\)（或 \(C^\infty\)）等距嵌入。维数 \(q\) 的估计来自以下技术需要：确保每步有足够的''自由方向''（\(q\) 个坐标函数 \(f^a\)）来匹配 \(n(n+1)/2\) 个独立方程及必要的非线性约束。\(\blacksquare\)

\textbf{定理 74.3}（Nash-Kuiper \(C^1\) 嵌入定理，1954-1955）：设 \((M^n, g)\) 为 \(C^1\) Riemann 流形（\(n \geq 2\)）。则存在 \(C^1\) 等距嵌入 \(f: M \to \mathbb{R}^{n+1}\)。特别地，存在 \(C^1\) 等距嵌入 \(S^2 \to \mathbb{R}^3\) 将单位球面映为任意小区域（这正是著名的''揉皱的球面''悖论）。此定理展示了 \(C^1\) 范畴中 Gauss 曲率约束的失效------在仅有 \(C^1\) 光滑性的情况下，Gauß-Bonnet 定理不再适用（因 \(K\) 仅在 \(C^2\) 条件下良好定义）。

\textbf{证明}（框架 --- Nash-Kuiper）：Nash-Kuiper 的核心技术也是''短嵌入 + 皱褶''，但皱褶构造在 \(C^1\) 范畴中更为彻底：每步扰动使度量误差缩小，同时控制嵌入的 \(C^1\) 范数，但 \(C^2\) 范数爆炸（因此极限嵌入仅为 \(C^1\) 而非 \(C^2\)）。具体地，取一族周期性的高频''螺旋''或''瓦楞''（corrugation）叠加到当前嵌入上，使得第一基本形式沿着瓦楞方向增加所需的量（\(\sim\) 振幅 \(\times\) 频率），而 \(C^1\) 范数保持有界。维数 \(q = n+1\) 的余维数 \(1\) 恰好够提供自由法方向。\(\blacksquare\)

\textbf{推论 74.4}（等距嵌入定理的意义与影响）：Nash 嵌入定理证明 Riemann 几何在原则上等价于 Euclidean 空间中子流形的外在几何，尽管嵌入维数极高（\(O(n^3)\)）。在实际数值应用中，等距嵌入问题与多维尺度分析（MDS）、流形学习（Isomap、LLE）和降维方法密切相关。定理还推动了几何流（Ricci 流、平均曲率流）理论中嵌入的存在性研究，以及一般相对论中 Cauchy 问题的等距嵌入公式化（为真空 Einstein 方程的约束方程提供了解释框架）。

\newpage

\chapter{07-拓扑与几何/02-微分几何/02-微分流形.md}\label{ux62d3ux6251ux4e0eux51e0ux4f5502-ux5faeux5206ux51e0ux4f5502-ux5faeux5206ux6d41ux5f62.md}

\section{第222章：微分流形}\label{ux7b2c222ux7ae0ux5faeux5206ux6d41ux5f62}

微分流形将曲面概念推广到任意维数，是现代微分几何的基石。微分流形是一种局部同胚于欧氏空间的拓扑空间，每个点附近可通过局部坐标（坐标卡）用 \(\mathbb{R}^n\) 中的坐标描述。不同坐标卡之间的重叠区域通过光滑的过渡映射（转移函数）相互联系，这种结构使得微积分运算得以在流形上全局一致地进行。

\subsection{64.1 微分流形的定义}\label{ux5faeux5206ux6d41ux5f62ux7684ux5b9aux4e49}

\textbf{定义 64.1}（拓扑流形）：设 \(M\) 是 Hausdorff 空间且具有可数基。\(M\) 称为 \(n\) 维\textbf{拓扑流形}，如果对每个 \(p \in M\)，存在 \(p\) 的开邻域 \(U\) 同胚于 \(\mathbb{R}^n\) 的某个开集。

\textbf{定义 64.2}（坐标卡与图册）：设 \(M\) 是 \(n\) 维拓扑流形。 - \(M\) 上的一个\textbf{坐标卡}（或局部坐标系）是一个对 \((U, \varphi)\)，其中 \(U \subseteq M\) 是开集，\(\varphi: U \to \mathbb{R}^n\) 是同胚到其像。 - \(M\) 的一个\textbf{图册} \(\mathcal{A} = \{(U_\alpha, \varphi_\alpha)\}\) 是一族坐标卡满足 \(\bigcup_\alpha U_\alpha = M\)。 - 两个坐标卡 \((U_\alpha, \varphi_\alpha)\) 和 \((U_\beta, \varphi_\beta)\) 称为 \(C^\infty\)\textbf{相容的}，如果 \(U_\alpha \cap U_\beta = \emptyset\)，或者\textbf{转移映射}

\[\varphi_\beta \circ \varphi_\alpha^{-1}: \varphi_\alpha(U_\alpha \cap U_\beta) \to \varphi_\beta(U_\alpha \cap U_\beta)\]

是 \(C^\infty\) 微分同胚。

\textbf{定义 64.3}（微分流形）：一个 \(C^\infty\)\textbf{微分流形}（或光滑流形）是一个拓扑流形 \(M\) 配上一个极大 \(C^\infty\) 图册（即所有与图册中坐标卡 \(C^\infty\) 相容的坐标卡都已在图册中）。

\textbf{定义 64.4}（光滑函数）：设 \(M\) 是光滑流形。函数 \(f: M \to \mathbb{R}\) 称为\textbf{光滑的}，如果对每个坐标卡 \((U, \varphi)\)，\(f \circ \varphi^{-1}: \varphi(U) \to \mathbb{R}\) 是光滑的（作为 \(\mathbb{R}^n\) 上开集的函数）。\(M\) 上所有光滑函数构成的集合记作 \(C^\infty(M)\)。

\textbf{定义 64.5}（光滑映射）：设 \(M, N\) 是光滑流形。映射 \(F: M \to N\) 称为\textbf{光滑的}，如果对每对坐标卡 \((U, \varphi)\)（\(M\) 上）和 \((V, \psi)\)（\(N\) 上）满足 \(F(U) \subseteq V\)，映射 \(\psi \circ F \circ \varphi^{-1}: \varphi(U) \to \psi(V)\) 是光滑的。

\subsection{64.2 切空间与切映射}\label{ux5207ux7a7aux95f4ux4e0eux5207ux6620ux5c04}

\textbf{定义 64.6}（芽与导子）：设 \(p \in M\)。\(p\) 处的\textbf{函数芽}是在 \(p\) 的某个邻域上定义的光滑函数的等价类（两个函数等价如它们在 \(p\) 的某个邻域上相等）。记作 \(C^\infty_p(M)\)。

\(p\) 处的一个\textbf{导子}（或切向量）是一个线性映射 \(v: C^\infty_p(M) \to \mathbb{R}\) 满足 Leibniz 法则：

\[v(fg) = v(f)g(p) + f(p)v(g)\]

\textbf{定义 64.7}（切空间）：\(M\) 在 \(p\) 处的\textbf{切空间} \(T_p M\) 是所有 \(p\) 处导子构成的向量空间。其维数为 \(n = \dim M\)。

在局部坐标 \((x^1, \ldots, x^n)\)（由坐标卡 \(\varphi\) 给出）下，偏导算子 \(\frac{\partial}{\partial x^i}\big|_p\) 构成 \(T_p M\) 的一组基。任意切向量可表示为

\[v = \sum_{i=1}^n v^i \frac{\partial}{\partial x^i}\bigg|_p\]

\textbf{定义 64.8}（切映射/微分）：设 \(F: M \to N\) 是光滑映射，\(p \in M\)。\(F\) 在 \(p\) 处的\textbf{切映射}（或微分）\(dF_p: T_p M \to T_{F(p)} N\) 定义为

\[(dF_p(v))(f) = v(f \circ F), \quad \forall v \in T_p M, \; f \in C^\infty_{F(p)}(N)\]

在局部坐标下，\(dF_p\) 由 Jacobi 矩阵表示。

\textbf{定理 64.1}（链式法则）：\((d(G \circ F))_p = dG_{F(p)} \circ dF_p\)。

\emph{证明}：取局部坐标 \(\{x^i\}\) 于 \(p\) 附近，\(\{y^j\}\) 于 \(F(p)\) 附近，\(\{z^k\}\) 于 \(G(F(p))\) 附近。由微分定义，\((dF)_p(\partial/\partial x^i) = \sum_j (\partial F^j/\partial x^i)(p) \partial/\partial y^j\)，\((dG)_{F(p)}(\partial/\partial y^j) = \sum_k (\partial G^k/\partial y^j)(F(p)) \partial/\partial z^k\)。复合得 \((d(G \circ F))_p(\partial/\partial x^i) = \sum_k \big(\sum_j (\partial G^k/\partial y^j)(F(p)) (\partial F^j/\partial x^i)(p)\big) \partial/\partial z^k = dG_{F(p)} \circ dF_p(\partial/\partial x^i)\)。由切向量的线性基任意性，等式对所有 \(v \in T_pM\) 成立。\(\blacksquare\)

\textbf{定义 64.9}（切丛）：\(M\) 的\textbf{切丛} \(TM = \bigsqcup_{p \in M} T_p M\) 具有自然的光滑流形结构，维数为 \(2n\)。自然投影 \(\pi: TM \to M\) 将 \(v \in T_p M\) 映到 \(p\)。

\subsection{64.3 向量场与 Lie 括号}\label{ux5411ux91cfux573aux4e0e-lie-ux62ecux53f7}

\textbf{定义 64.10}（向量场）：\(M\) 上的一个\textbf{光滑向量场}是一个光滑映射 \(X: M \to TM\) 使得 \(\pi \circ X = \operatorname{id}_M\)，即 \(X(p) \in T_p M\) 对每个 \(p\)。所有光滑向量场构成的集合记作 \(\mathfrak{X}(M)\)。

在局部坐标下，\(X = \sum_{i=1}^n X^i \frac{\partial}{\partial x^i}\)，其中 \(X^i \in C^\infty(M)\)。

\textbf{定义 64.11}（Lie 括号）：设 \(X, Y \in \mathfrak{X}(M)\)。它们的 \textbf{Lie 括号} \([X, Y] \in \mathfrak{X}(M)\) 定义为

\[[X, Y](f) = X(Y(f)) - Y(X(f)), \quad \forall f \in C^\infty(M)\]

在局部坐标下，\([X, Y] = \sum_{i,j} \left(X^j \frac{\partial Y^i}{\partial x^j} - Y^j \frac{\partial X^i}{\partial x^j}\right) \frac{\partial}{\partial x^i}\)。

\textbf{命题 64.2}（Lie 括号的性质）： 1. \([X, Y] = -[Y, X]\)（反交换性） 2. \([aX + bY, Z] = a[X, Z] + b[Y, Z]\)（双线性性） 3. \([X, [Y, Z]] + [Y, [Z, X]] + [Z, [X, Y]] = 0\)（Jacobi 恒等式） 4. \([fX, gY] = fg[X, Y] + f(Xg)Y - g(Yf)X\)

\textbf{证明}：（1）由定义 \([X,Y](f) = X(Y(f)) - Y(X(f)) = -[Y,X](f)\)。（2）直接由导子的线性性得。（3）展开：\([X,[Y,Z]](f) = X(Y(Z(f)) - Z(Y(f))) - (Y(Z(X(f))) - Z(Y(X(f))))\)，\([Y,[Z,X]](f)\) 和 \([Z,[X,Y]](f)\) 类似展开，六项中的正负项两两抵消，和为零。（4）\([fX,gY](h) = fX(gY(h)) - gY(fX(h)) = fgX(Y(h)) + f(Xg)Y(h) - gfY(X(h)) - g(Yf)X(h) = fg[X,Y](h) + f(Xg)Y(h) - g(Yf)X(h)\)。\(\blacksquare\)

\subsection{64.4 微分形式}\label{ux5faeux5206ux5f62ux5f0f}

\textbf{定义 64.12}（余切空间与 1-形式）：\(T_p^* M = (T_p M)^*\) 是 \(T_p M\) 的对偶空间，称为 \(p\) 处的\textbf{余切空间}。\(M\) 上的一个\textbf{光滑 1-形式}是一个光滑映射 \(\omega: M \to T^*M = \bigsqcup T_p^* M\)，使得 \(\omega(p) \in T_p^* M\)。

在局部坐标下，\(\{dx^1, \ldots, dx^n\}\) 是 \(T_p^* M\) 的基（与 \(\{\partial/\partial x^i\}\) 对偶）。1-形式可写为 \(\omega = \sum \omega_i dx^i\)。

\textbf{定义 64.13}（微分 \(d\)）：对光滑函数 \(f \in C^\infty(M)\)，其\textbf{微分} \(df\) 是 1-形式，定义为 \(df(v) = v(f)\) 对所有 \(v \in TM\)。在局部坐标下，\(df = \sum \frac{\partial f}{\partial x^i} dx^i\)。

\textbf{定义 64.14}（\(k\)-形式）：\(M\) 上的一个\textbf{微分 \(k\)-形式}是一个反对称的 \((0, k)\)-张量场。在局部坐标下，

\[\omega = \sum_{i_1 < \cdots < i_k} \omega_{i_1 \cdots i_k} dx^{i_1} \wedge \cdots \wedge dx^{i_k}\]

其中 \(\wedge\) 是外积（楔积），满足反对称性：\(dx^i \wedge dx^j = -dx^j \wedge dx^i\)。

\textbf{定义 64.15}（外微分）：外微分算子 \(d: \Omega^k(M) \to \Omega^{k+1}(M)\) 定义为

\[d\omega = \sum_{i_1 < \cdots < i_k} \sum_{j} \frac{\partial \omega_{i_1 \cdots i_k}}{\partial x^j} dx^j \wedge dx^{i_1} \wedge \cdots \wedge dx^{i_k}\]

\textbf{命题 64.3}（外微分的基本性质）： 1. \(d\) 是线性的。 2. \(d \circ d = 0\)（即 \(d(d\omega) = 0\) 对所有 \(\omega\)）。 3. \(d(\omega \wedge \eta) = d\omega \wedge \eta + (-1)^{\deg \omega} \omega \wedge d\eta\)。 4. 对函数 \(f\)，\(df\) 与定义 64.13 一致。 5. \(d\) 是局部的：若 \(\omega = \eta\) 在开集 \(U\) 上，则 \(d\omega = d\eta\) 在 \(U\) 上。

\textbf{证明}：（1）由定义 \(d\) 是系数偏导数的线性组合，线性性显然。（2）\(d^2\omega = \sum_{i_1<\cdots<i_k} \sum_{j,\ell} \frac{\partial^2 \omega_{i_1\cdots i_k}}{\partial x^\ell \partial x^j} dx^\ell \wedge dx^j \wedge dx^{i_1} \wedge \cdots \wedge dx^{i_k}\)，因 \(\frac{\partial^2}{\partial x^\ell \partial x^j} = \frac{\partial^2}{\partial x^j \partial x^\ell}\) 而 \(dx^\ell \wedge dx^j = -dx^j \wedge dx^\ell\)，所有项成对抵消。（3）对单项式 \(\omega = f dx^I\)，\(\eta = g dx^J\)，\(d(\omega \wedge \eta) = d(fg) \wedge dx^I \wedge dx^J = (df \cdot g + f \cdot dg) \wedge dx^I \wedge dx^J = (df \wedge dx^I) \wedge (g dx^J) + (-1)^{|I|} (f dx^I) \wedge (dg \wedge dx^J) = d\omega \wedge \eta + (-1)^{\deg \omega} \omega \wedge d\eta\)。（4）由定义直接验证。（5）\(d\) 仅涉及偏导数，局部性由定义即得。\(\blacksquare\)

\textbf{定义 64.16}（闭形式与恰当形式）：\(k\)-形式 \(\omega\) 称为 - \textbf{闭的}：若 \(d\omega = 0\)。 - \textbf{恰当的}：若存在 \((k-1)\)-形式 \(\eta\) 使得 \(\omega = d\eta\)。

由 \(d^2 = 0\)，恰当形式一定是闭的。闭形式商去恰当形式得到 \textbf{de Rham 上同调}：

\[H^k_{dR}(M) = \frac{\ker(d: \Omega^k(M) \to \Omega^{k+1}(M))}{\operatorname{im}(d: \Omega^{k-1}(M) \to \Omega^k(M))}\]

\subsection{64.5 子流形与浸入}\label{ux5b50ux6d41ux5f62ux4e0eux6d78ux5165}

\textbf{定义 64.17}（浸入与淹没）：光滑映射 \(F: M \to N\) 称为 - \textbf{浸入}：若 \(dF_p\) 在每点 \(p\) 是单射。 - \textbf{淹没}：若 \(dF_p\) 在每点 \(p\) 是满射。 - \textbf{嵌入}：若 \(F\) 是浸入且 \(F: M \to F(M)\) 是同胚（\(F(M)\) 赋予 \(N\) 的子空间拓扑）。

\textbf{定理 64.4}（正则值定理）：设 \(F: M \to N\) 是光滑映射，\(\dim M \geq \dim N\)。若 \(q \in N\) 是 \(F\) 的\textbf{正则值}（即对所有 \(p \in F^{-1}(q)\)，\(dF_p\) 是满射），则 \(F^{-1}(q)\) 是 \(M\) 的 \(\dim M - \dim N\) 维嵌入子流形。

\emph{证明}：对任意 \(p \in F^{-1}(q)\)，由秩定理存在局部坐标使 \(F(x_1, \ldots, x_m) = (x_1, \ldots, x_n)\)（因 \(dF_p\) 满射）。在此坐标下 \(F^{-1}(q) = \{x_1 = q_1, \ldots, x_n = q_n\}\)，这给出 \(\{x_{n+1}, \ldots, x_m\}\) 参数化的 \((\dim M - \dim N)\) 维子流形。不同参数域间的转移映射为微分同胚，故 \(F^{-1}(q)\) 是嵌入子流形。\(\blacksquare\)

\textbf{定理 64.5}（Whitney 嵌入定理，陈述）：每个 \(n\) 维光滑流形可以嵌入到 \(\mathbb{R}^{2n}\) 中。

\textbf{证明}：将 \(M\) 覆盖于可数多个坐标邻域 \(\{U_i\}\) 上，取从属单位分解 \(\{\varphi_i\}\)。构造映射 \(F: M \to \mathbb{R}^{(n+1)N}\)，分量由 \(\varphi_i\) 和坐标函数的乘积构成，其中 \(N\) 为覆盖数。此为浸入。通过管状邻域投影可降维至 \(\mathbb{R}^{2n+1}\) 中的嵌入，再通过 Sard 定理将维数降至 \(\mathbb{R}^{2n}\)------因为从 \(\mathbb{R}^{2n+1}\) 中沿几乎任意方向的投影将嵌入映为 \(\mathbb{R}^{2n}\) 中的嵌入。\(\blacksquare\)

\hrulefill

\hrulefill

\hrulefill

\hrulefill

\hrulefill

\hrulefill

\section{第223章：Riemann 几何}\label{ux7b2c223ux7ae0riemann-ux51e0ux4f55}

Riemann 几何在微分流形上赋予度量结构，使得可以讨论长度、角度、曲率等概念。Riemann 度量 \(g\) 在每个切空间上定义内积，从而决定曲线弧长、测地线方程和体积元。测地线是流形上''最直''的曲线，对应于欧氏空间中直线的推广；而 Riemann 曲率张量则刻画了流形偏离平坦的程度------它既是局部几何的核心不变量，又通过 Gauss-Bonnet 定理等结果与流形的整体拓扑建立了深刻联系。

\subsection{65.1 Riemann 度量}\label{riemann-ux5ea6ux91cf}

\textbf{定义 65.1}（Riemann 度量）：光滑流形 \(M\) 上的一个 \textbf{Riemann 度量} \(g\) 是一个光滑的 \((0, 2)\)-张量场，使得对每个 \(p \in M\)，\(g_p: T_p M \times T_p M \to \mathbb{R}\) 是对称正定的双线性型。即 \(g_p\) 是 \(T_p M\) 上的内积。

在局部坐标下，\(g = \sum_{i,j} g_{ij} dx^i \otimes dx^j\)，其中 \((g_{ij})\) 是对称正定矩阵，\(g_{ij} = g(\partial_i, \partial_j)\)。

\textbf{定义 65.2}（Riemann 流形）：一个 \textbf{Riemann 流形}是一对 \((M, g)\)，其中 \(M\) 是光滑流形，\(g\) 是 \(M\) 上的 Riemann 度量。

\textbf{定义 65.3}（等距）：两个 Riemann 流形 \((M, g)\) 和 \((N, h)\) 之间的微分同胚 \(F: M \to N\) 称为\textbf{等距}，如果 \(F^*h = g\)，即 \(h_{F(p)}(dF_p(v), dF_p(w)) = g_p(v, w)\) 对所有 \(v, w \in T_p M\)。

\subsection{65.2 Levi-Civita 联络}\label{levi-civita-ux8054ux7edc}

\textbf{定义 65.4}（仿射联络）：\(M\) 上的一个\textbf{仿射联络}（或协变导数）是一个映射 \(\nabla: \mathfrak{X}(M) \times \mathfrak{X}(M) \to \mathfrak{X}(M)\)，记作 \(\nabla_X Y\)，满足：

\begin{enumerate}
\def\labelenumi{\arabic{enumi}.}
\tightlist
\item
  \(\nabla_{fX + gY} Z = f \nabla_X Z + g \nabla_Y Z\)（\(C^\infty(M)\)-线性于第一变量）
\item
  \(\nabla_X (aY + bZ) = a \nabla_X Y + b \nabla_X Z\)（\(\mathbb{R}\)-线性于第二变量）
\item
  \(\nabla_X (fY) = (Xf)Y + f \nabla_X Y\)（Leibniz 法则）
\end{enumerate}

\textbf{定理 65.1}（Levi-Civita 联络的存在唯一性）：在 Riemann 流形 \((M, g)\) 上，存在唯一的仿射联络 \(\nabla\) 满足：

\begin{enumerate}
\def\labelenumi{\arabic{enumi}.}
\tightlist
\item
  \textbf{无挠性}（对称性）：\(\nabla_X Y - \nabla_Y X = [X, Y]\)。
\item
  \textbf{与度量相容}：\(X(g(Y, Z)) = g(\nabla_X Y, Z) + g(Y, \nabla_X Z)\)。
\end{enumerate}

这个联络称为 \textbf{Levi-Civita 联络}（或 Riemann 联络）。

\emph{证明}（存在性）：由 Koszul 公式：

\[\begin{aligned}
2g(\nabla_X Y, Z) &= X(g(Y, Z)) + Y(g(Z, X)) - Z(g(X, Y)) \\
&\quad - g(X, [Y, Z]) + g(Y, [Z, X]) + g(Z, [X, Y])
\end{aligned}\]

此公式唯一确定了 \(\nabla_X Y\)。可直接验证由此定义的 \(\nabla\) 满足所有条件。∎

\textbf{定义 65.5}（Christoffel 符号）：在局部坐标下，Levi-Civita 联络由 \textbf{Christoffel 符号} \(\Gamma_{ij}^k\) 决定：

\[\nabla_{\partial_i} \partial_j = \sum_k \Gamma_{ij}^k \partial_k\]

\[\Gamma_{ij}^k = \frac{1}{2} \sum_{\ell} g^{k\ell}(\partial_i g_{j\ell} + \partial_j g_{i\ell} - \partial_\ell g_{ij})\]

\subsection{65.3 曲率}\label{ux66f2ux7387}

\textbf{定义 65.6}（Riemann 曲率张量）：Levi-Civita 联络 \(\nabla\) 的 \textbf{Riemann 曲率张量} \(R: \mathfrak{X}(M) \times \mathfrak{X}(M) \times \mathfrak{X}(M) \to \mathfrak{X}(M)\) 定义为

\[R(X, Y)Z = \nabla_X \nabla_Y Z - \nabla_Y \nabla_X Z - \nabla_{[X, Y]} Z\]

\textbf{命题 65.2}（Riemann 曲率张量的对称性）：

\begin{enumerate}
\def\labelenumi{\arabic{enumi}.}
\tightlist
\item
  \(R(X, Y)Z = -R(Y, X)Z\)
\item
  \(g(R(X, Y)Z, W) = -g(R(X, Y)W, Z)\)
\item
  \(R(X, Y)Z + R(Y, Z)X + R(Z, X)Y = 0\)（第一 Bianchi 恒等式）
\item
  \(g(R(X, Y)Z, W) = g(R(Z, W)X, Y)\)（对偶对称性）
\end{enumerate}

\textbf{证明}：（1）由定义 \(R(X,Y)Z = \nabla_X \nabla_Y Z - \nabla_Y \nabla_X Z - \nabla_{[X,Y]} Z\)，交换 \(X,Y\) 并利用 \([Y,X] = -[X,Y]\) 即得。（2）由 \(\nabla\) 与度量相容，\(g(\nabla_X \nabla_Y Z, W) = X(g(\nabla_Y Z, W)) - g(\nabla_Y Z, \nabla_X W)\)。展开 \(g(R(X,Y)Z, W)\) 并利用 \(X(g(Y,Z)) = Y(g(X,Z))\) 型恒等式和对称性，可证 \(g(R(X,Y)Z, W) = -g(R(X,Y)W, Z)\)。（3）由 \(\nabla\) 无挠性 \(\nabla_X Y - \nabla_Y X = [X,Y]\)，展开 \(R(X,Y)Z + R(Y,Z)X + R(Z,X)Y\) 后所有 \(\nabla\) 项抵消。（4）由 (1)-(3) 通过代数操作推导：\(g(R(X,Y)Z,W) - g(R(Z,W)X,Y)\) 在轮换下不变，由 (3) 可证其为零。\(\blacksquare\)

\textbf{定义 65.7}（Riemann 曲率张量的分量）：在局部坐标下，

\[R(\partial_i, \partial_j)\partial_k = \sum_{\ell} R_{ijk}^{\ell} \partial_{\ell}\]

\[R_{ijk\ell} = g(R(\partial_i, \partial_j)\partial_k, \partial_\ell) = \sum_m g_{\ell m} R_{ijk}^m\]

\[R_{ijk\ell} = \frac{1}{2}\left(\frac{\partial^2 g_{i\ell}}{\partial x^j \partial x^k} + \frac{\partial^2 g_{jk}}{\partial x^i \partial x^\ell} - \frac{\partial^2 g_{ik}}{\partial x^j \partial x^\ell} - \frac{\partial^2 g_{j\ell}}{\partial x^i \partial x^k}\right) + \sum_{p,q} g_{pq}(\Gamma_{i\ell}^p \Gamma_{jk}^q - \Gamma_{ik}^p \Gamma_{j\ell}^q)\]

\textbf{定义 65.8}（截面曲率）：设 \(\Pi \subseteq T_p M\) 是二维子空间，由规范正交向量 \(v, w\) 张成。\(\Pi\) 的\textbf{截面曲率}为

\[K(\Pi) = \frac{g(R(v, w)w, v)}{g(v, v)g(w, w) - g(v, w)^2}\]

\textbf{定义 65.9}（Ricci 曲率与标量曲率）： - \textbf{Ricci 曲率}：\(\operatorname{Ric}(X, Y) = \operatorname{tr}(Z \mapsto R(Z, X)Y)\)。在局部坐标下，\(\operatorname{Ric}_{ij} = \sum_k R_{kij}^k\)。 - \textbf{标量曲率}：\(S = \operatorname{tr}_g(\operatorname{Ric}) = \sum_{i,j} g^{ij} \operatorname{Ric}_{ij}\)。

\textbf{命题 65.3}（第二 Bianchi 恒等式）：

\[\nabla_X R(Y, Z) + \nabla_Y R(Z, X) + \nabla_Z R(X, Y) = 0\]

在局部坐标下：\(\nabla_m R_{ijk\ell} + \nabla_i R_{jmk\ell} + \nabla_j R_{mik\ell} = 0\)。

\textbf{证明}：利用协变导数的定义 \((\nabla_X R)(Y,Z)W = \nabla_X (R(Y,Z)W) - R(\nabla_X Y, Z)W - R(Y, \nabla_X Z)W - R(Y,Z)\nabla_X W\)。在法坐标下（\(\Gamma_{ij}^k(p) = 0\)），\(\nabla_m R_{ijk\ell} = \partial_m R_{ijk\ell}\)。由曲率张量在法坐标下的表达式 \(R_{ijk\ell} = \frac{1}{2}(\partial_i \partial_\ell g_{jk} + \partial_j \partial_k g_{i\ell} - \partial_i \partial_k g_{j\ell} - \partial_j \partial_\ell g_{ik})\)，对其求导并轮换指标 \(i,j,m\) 即得所需的恒等式。由于恒等式是张量等式，在法坐标下验证即对任意坐标系成立。\(\blacksquare\)

\subsection{65.4 测地线}\label{ux6d4bux5730ux7ebf}

\textbf{定义 65.10}（测地线）：Riemann 流形 \((M, g)\) 上的光滑曲线 \(\gamma: I \to M\) 称为\textbf{测地线}，如果

\[\nabla_{\gamma'} \gamma' = 0\]

在局部坐标下，这等价于测地线方程：

\[\frac{d^2 x^k}{dt^2} + \sum_{i,j} \Gamma_{ij}^k \frac{dx^i}{dt}\frac{dx^j}{dt} = 0\]

\textbf{定理 65.4}（测地线的存在唯一性）：对每个 \(p \in M\) 和 \(v \in T_p M\)，存在唯一的测地线 \(\gamma_v: (-\varepsilon, \varepsilon) \to M\) 满足 \(\gamma_v(0) = p\)，\(\gamma_v'(0) = v\)。

\emph{证明}：测地线方程是二阶常微分方程，由 ODE 解的存在唯一性定理可得。∎

\textbf{定义 65.11}（指数映射）：\(p\) 处的\textbf{指数映射} \(\exp_p: T_p M \to M\) 定义为 \(\exp_p(v) = \gamma_v(1)\)（当右边有定义时）。\(\exp_p\) 在 \(0 \in T_p M\) 附近是微分同胚（因为 \(d(\exp_p)_0 = \operatorname{id}_{T_p M}\)）。

\textbf{定义 65.12}（完备性）：Riemann 流形 \((M, g)\) 称为\textbf{测地完备的}，如果每条测地线可以无限延伸（即定义域为 \(\mathbb{R}\)）。

\textbf{定理 65.5}（Hopf-Rinow 定理）：对 Riemann 流形 \((M, g)\)，以下条件等价： 1. \((M, g)\) 是测地完备的。 2. \((M, g)\) 作为度量空间（由 \(g\) 诱导的 Riemann 距离）是完备的。 3. \(M\) 中的有界闭集是紧致的。

此时，\(M\) 中任意两点之间都存在最短测地线。

\emph{证明}：(1)\(\Rightarrow\)(2)：若 \(\{p_n\}\) 是 Cauchy 列，取 \(p_n\) 到 \(p_0\) 的极小测地线 \(\gamma_n\)。由单位切球紧致，\(\gamma_n'(0)\) 有子列收敛，\(\exp\) 连续性给出极限点。(2)\(\Rightarrow\)(3)：由完备度量空间中有界闭集的紧致性准则。(3)\(\Rightarrow\)(1)：反设某测地线在有限时间内截断，取截断序列，由条件(3)得有界闭球紧致推出收敛矛盾。最短测地线由 Arzela-Ascoli 给出。\(\blacksquare\)

\hrulefill

\hrulefill

\hrulefill

\hrulefill

\hrulefill

\hrulefill

\section{第224章：曲率与拓扑}\label{ux7b2c224ux7ae0ux66f2ux7387ux4e0eux62d3ux6251}

本章揭示曲率（局部几何量）与拓扑（全局拓扑不变量）之间的深刻联系。

\subsection{66.1 二维情形：Gauss-Bonnet 定理}\label{ux4e8cux7ef4ux60c5ux5f62gauss-bonnet-ux5b9aux7406}

\textbf{定理 66.1}（局部 Gauss-Bonnet 定理）：设 \(S\) 是 \(\mathbb{R}^3\) 中的正则曲面，\(R \subseteq S\) 是由分段光滑简单闭曲线 \(\partial R\) 围成的区域。则

\[\int_R K dA + \int_{\partial R} k_g ds + \sum_i \theta_i = 2\pi\]

其中 \(K\) 是 Gauss 曲率，\(k_g\) 是边界曲线的测地曲率，\(\theta_i\) 是边界角点处的外角。

\textbf{证明}：在 \(R\) 上选取正交标架场 \(\{e_1, e_2\}\)，定义联络形式 \(\omega_{12} = \langle \nabla e_1, e_2 \rangle\)。则由曲率的定义，\(d\omega_{12} = -K dA\)。再由 Stokes 定理：

\[\int_R K dA = -\int_R d\omega_{12} = -\int_{\partial R} \omega_{12}\]

在边界上 \(\omega_{12}\) 与 \(k_g\) 的关系给出 \(\int_{\partial R} \omega_{12} = 2\pi - \int_{\partial R} k_g ds - \sum \theta_i\)。∎

\textbf{定理 66.2}（全局 Gauss-Bonnet 定理）：设 \(S\) 是紧致无边可定向曲面。则

\[\int_S K dA = 2\pi \chi(S)\]

其中 \(\chi(S)\) 是 \(S\) 的 Euler 示性数。对亏格为 \(g\) 的可定向曲面，\(\chi(S) = 2 - 2g\)。

\emph{证明}：将 \(S\) 三角剖分，对每个三角形应用局部 Gauss-Bonnet 定理，求和。所有内角之和为 \(2\pi\) 乘以顶点数，测地曲率积分在内部边上抵消，外角之和给出 \(\int K dA = 2\pi(V - E + F) = 2\pi \chi(S)\)。∎

\textbf{推论 66.3}（Gauss-Bonnet 定理的推论）： 1. \(\int_S K dA\) 是拓扑不变量（与 Riemann 度量的选择无关）。 2. 若 \(S\) 同胚于 \(S^2\)（球面），则 \(\int_S K dA = 4\pi\)。特别地，\(S^2\) 上的任何 Riemann 度量在某处必有正曲率。 3. 若 \(S\) 同胚于 \(T^2\)（环面），则 \(\int_S K dA = 0\)。因此环面上不能处处有正 Gauss 曲率。 4. 若 \(S\) 的 Gauss 曲率 \(K > 0\) 处处成立，则 \(\chi(S) > 0\)，从而 \(S\) 必同胚于 \(S^2\) 或 \(\mathbb{R}P^2\)。

\subsection{66.2 高维推广}\label{ux9ad8ux7ef4ux63a8ux5e7f}

\textbf{定理 66.4}（Chern-Gauss-Bonnet 定理，陈述）：设 \((M, g)\) 是 \(2n\) 维紧致无边可定向 Riemann 流形。则

\[\int_M \operatorname{Pf}(\Omega) = (2\pi)^n \chi(M)\]

其中 \(\operatorname{Pf}(\Omega)\) 是曲率形式 \(\Omega\) 的 Pfaffian，\(\chi(M)\) 是 \(M\) 的 Euler 示性数。

\emph{说明}：当 \(n = 1\) 时，\(\operatorname{Pf}(\Omega) = K dA\)，退化为经典的 Gauss-Bonnet 定理。此定理是陈省身（Shiing-Shen Chern）于 1944 年证明的，是微分几何的里程碑式成果。

\textbf{证明}：在 \(M\) 上取正交标架，记曲率 2-形式为 \(\Omega\)。欧拉形式 \(\operatorname{Pf}(\Omega)\) 在标架变化下不变，故是全局定义的 \(2n\)-形式。关键步骤：（1）证明 \(\operatorname{Pf}(\Omega)\) 是闭形式（由 Bianchi 恒等式）；（2）局部构造 transgression 形式证明 de Rham 上同调类 \([\operatorname{Pf}(\Omega)]\) 与度量无关；（3）对 \(S^{2n}\) 计算得到 \((2\pi)^n \chi(S^{2n}) = (2\pi)^n \cdot 2\)，验证常数。一般情形通过对向量丛取横截截口并计算其零集的 Euler 数（Chern-Weil 理论）完成。\(\blacksquare\)

\textbf{定理 66.5}（奇维数情形）：若 \(M\) 是紧致无边可定向的奇维数流形，则 \(\chi(M) = 0\)。

\emph{证明}：由 Poincaré 对偶，\(H_k(M; \mathbb{R}) \cong H_{2n+1-k}(M; \mathbb{R})\)，故 \(\beta_k = \beta_{2n+1-k}\)。Euler 示性数 \(\chi(M) = \sum_{k=0}^{2n+1} (-1)^k \beta_k = 0\)。也可由奇维数紧致流形容有一个处处非零的向量场（Poincaré-Hopf 指标定理的推论）直接推得。\(\blacksquare\)

\subsection{66.3 曲率与拓扑的其他联系}\label{ux66f2ux7387ux4e0eux62d3ux6251ux7684ux5176ux4ed6ux8054ux7cfb}

\textbf{定理 66.6}（Bonnet-Myers 定理）：设 \((M, g)\) 是完备 Riemann 流形，\(\dim M = n\)。若存在常数 \(c > 0\) 使得 \(\operatorname{Ric}(v, v) \geq (n-1)c \cdot g(v, v)\) 对所有 \(v\) 成立，则 \(M\) 是紧致的，且直径 \(\operatorname{diam}(M) \leq \pi / \sqrt{c}\)。特别地，\(\pi_1(M)\) 是有限群。

\textbf{证明}：设 \(\gamma\) 是长度大于 \(\pi/\sqrt{c}\) 的测地线。利用第二变分公式和 Ricci 曲率条件，可构造连接 \(\gamma\) 两端点的更短曲线，矛盾。因此所有测地线长度不超过 \(\pi/\sqrt{c}\)。由 Hopf-Rinow 定理，\(M\) 紧致且直径有界。万有覆叠 \(\tilde{M}\) 也具有相同的 Ricci 曲率下界，故也紧致，从而 \(\pi_1(M)\) 有限。∎

\textbf{定理 66.7}（Hadamard-Cartan 定理）：设 \((M, g)\) 是完备单连通 Riemann 流形，且截面曲率 \(K \leq 0\) 处处成立。则对每个 \(p \in M\)，指数映射 \(\exp_p: T_p M \to M\) 是微分同胚。特别地，\(M\) 微分同胚于 \(\mathbb{R}^n\)。

\textbf{证明}：截面曲率 \(\leq 0\) 蕴含 Jacobi 场的长度不减，从而 \(\exp_p\) 是局部微分同胚。由完备性，\(\exp_p\) 是满射。单连通性保证 \(\exp_p\) 是整体单射。∎

\textbf{定理 66.8}（球面定理 / 1/4-夹挤定理）：设 \((M, g)\) 是完备单连通 Riemann 流形，\(\dim M = n\)。若截面曲率 \(K\) 满足 \(1/4 < K \leq 1\)，则 \(M\) 同胚于 \(S^n\)。若 \(1/4 \leq K \leq 1\)，则 \(M\) 同胚于 \(S^n\) 或微分同胚于某个对称空间。

\emph{说明}：这是微分几何中著名的「球面定理」，由 Rauch、Berger 和 Klingenberg 在 1950-60 年代建立。其完整证明需要深入的比较几何和 Morse 理论。

\textbf{证明}：由 Rauch 比较定理，条件 \(K \geq 1/4 > 0\) 保证内射半径 \(\operatorname{inj}(M) \geq \pi\)。单连通性和正曲率保证 \(M\) 上存在 Morse 函数（距离平方函数），其指标最低的非平凡临界点必在 \(\pi\) 处（由曲率条件限制）。由 Morse 理论和拓扑球面定理（n≥5 的 h-配边定理，n=3 的 Poincaré 猜想，n=4 的 Freedman），推得 \(M\) 同胚于 \(S^n\)。等号情形涉及秩 1 对称空间。\(\blacksquare\)

\textbf{定理 66.9}（Chern 的 Gauss-Bonnet 公式推广的积分不变量）：在偶维数紧致 Riemann 流形上，Euler 示性数可以表示为曲率形式的多项式在整个流形上的积分，这一积分值不依赖于度量的选择。

\textbf{证明}：这是 Chern-Gauss-Bonnet 定理的重述。Pfaffian \(\operatorname{Pf}(\Omega)\) 是曲率 2-形式的多项式。若取不同度量 \(g_1\) 和 \(g_2\)，对应的曲率形式 \(\Omega_1\) 和 \(\Omega_2\) 给出微分形式 \(\operatorname{Pf}(\Omega_1)\) 和 \(\operatorname{Pf}(\Omega_2)\)。由 Chern-Weil 理论，它们决定相同 de Rham 上同调类（差一个恰当形式），故积分值 \(\int_M \operatorname{Pf}(\Omega_1) = \int_M \operatorname{Pf}(\Omega_2) = (2\pi)^n \chi(M)\) 与度量选取无关。\(\blacksquare\)

\subsection{66.4 曲率张量的进一步分解}\label{ux66f2ux7387ux5f20ux91cfux7684ux8fdbux4e00ux6b65ux5206ux89e3}

\textbf{定义 66.10}（Weyl 曲率张量）：在 \(n \geq 3\) 维 Riemann 流形上，\textbf{Weyl 曲率张量} \(W\) 定义为 Riemann 曲率张量去掉 Ricci 部分和标量曲率部分：

\[W_{ijk\ell} = R_{ijk\ell} - \frac{1}{n-2}(g_{ik}\operatorname{Ric}_{j\ell} - g_{i\ell}\operatorname{Ric}_{jk} - g_{jk}\operatorname{Ric}_{i\ell} + g_{j\ell}\operatorname{Ric}_{ik}) + \frac{S}{(n-1)(n-2)}(g_{ik}g_{j\ell} - g_{i\ell}g_{jk})\]

\(W\) 具有与 \(R\) 相同的对称性，且所有迹为零（\(\sum_i W_{ijk}^i = 0\)）。\(W\) 是\textbf{共形不变}的：若 \(\tilde{g} = e^{2f} g\)，则 \(\tilde{W} = e^{2f} W\)（在指标位置适当调整后）。

\textbf{定理 66.11}（Weyl 张量等于零）：\(n \geq 4\) 维 Riemann 流形局部共形平坦（即局部等距于 \(\mathbb{R}^n\) 配以某个共形度量）当且仅当 \(W = 0\)。\(n = 3\) 时，局部共形平坦等价于 Cotton 张量等于零。

\emph{证明}：若 \(g = e^{2f}\delta\) 局部共形平坦，由 Weyl 张量共形不变性 \(W_g = e^{2f}W_\delta = 0\)。反之设 \(W=0\)，则 Riemann 曲率完全由 Ricci 和标量曲率表达。第二 Bianchi 恒等式在 \(W=0\) 时给出 Cotton 张量条件，在 \(n \geq 4\) 时这等价于局部共形平坦性（由 Schouten 的可积性定理：\(W=0\) 推得存在函数 \(f\) 使得 \(e^{2f}g\) 为平坦度量）。\(n=3\) 时 \(W \equiv 0\)（代数原因），需 Cotton 张量 \(C = 0\) 来刻画共形平坦性。\(\blacksquare\)

\hrulefill

\hrulefill

\hrulefill

\hrulefill

\hrulefill

\newpage

\chapter{07-拓扑与几何/02-微分几何/03-G结构与张量.md}\label{ux62d3ux6251ux4e0eux51e0ux4f5502-ux5faeux5206ux51e0ux4f5503-gux7ed3ux6784ux4e0eux5f20ux91cf.md}

\section{第225章：G-结构}\label{ux7b2c225ux7ae0g-ux7ed3ux6784}

G-结构是微分流形切丛的结构群从一般线性群 \(\operatorname{GL}(n,\mathbb{R})\) 约化到 Lie 子群 \(G\) 的几何框架。形式化地，\(n\) 维流形 \(M\) 上的 \(G\)-结构是一个 \(G\)-主丛 \(P \to M\) 连同 \(G\)-等变同构 \(P \times_G \operatorname{GL}(n,\mathbb{R}) \cong FM\)（\(FM\) 为切标架丛），即切丛的每个标架均可通过 \(G\) 的作用彼此转换。\textbf{可积性条件}由结构张量（挠率）的消去判定：若存在无挠联络与 \(G\)-结构相容，则 \(G\)-结构局部平坦。

\textbf{定义 67.1（G-结构的结构张量）}：设 \(P \to M\) 是 \(G\)-结构，\(\mathfrak{g} \subset \mathfrak{gl}(n,\mathbb{R})\) 是 \(G\) 的 Lie 代数。取 \(P\) 上的任意联络 \(\nabla\)（不必与 \(G\)-结构相容），将其挠率 \(T\) 投影到商空间 \(\mathfrak{gl}(n,\mathbb{R})/\mathfrak{g}\) 上得到 \(G\)-结构的\textbf{结构张量}（也称为内在挠率）。该投影不依赖于联络 \(\nabla\) 的选取（在商空间中），因此是 \(G\)-结构的内在不变量。结构张量为零是 \(G\)-结构可积（即存在无挠相容联络）的必要条件，但通常不是充分条件（仅对某些 \(G\)，如 \(\operatorname{O}(n)\)，是充分的）。

\textbf{定理 67.1.1（Sternberg 的 \(G\)-结构分类）}：\(G\)-结构的可积性由其结构群 \(G\) 的 Lie 代数的 Spencer 上同调 \(H^{k,2}(\mathfrak{g})\) 决定。当 \(H^{1,2}(\mathfrak{g}) = 0\) 时，结构张量消去 \(\Rightarrow\) 可积（即 \(G\)-结构局部平坦）。当 \(H^{1,2}(\mathfrak{g}) \neq 0\) 时，需要高阶结构张量。这一分类是 Cartan 等价方法的现代表述。

经典 \(G\)-结构包括：

（1）\textbf{Riemann 结构}------\(G = \operatorname{O}(n)\)。等价于 Riemann 度量的存在（正交标架丛）。\(\mathfrak{gl}(n,\mathbb{R})/\mathfrak{o}(n) \cong \operatorname{Sym}^2(\mathbb{R}^n)\) 由对称矩阵构成。任何 \(\operatorname{O}(n)\)-结构都容许唯一的无挠相容联络（Levi-Civita 联络），因此 \(\operatorname{O}(n)\)-结构总是可积的------这就是 Riemann 几何的基本定理：度量唯一确定 Levi-Civita 联络。

（2）\textbf{复结构}------\(G = \operatorname{GL}(n,\mathbb{C}) \subset \operatorname{GL}(2n,\mathbb{R})\)。对应近复结构，可积性即 Newlander-Nirenberg 定理------近复结构来自复流形当且仅当其 Nijenhuis 张量 \(N_J(X,Y) = [X,Y] + J[JX,Y] + J[X,JY] - [JX,JY]\) 消去。Nijenhuis 张量本质上就是 \(\operatorname{GL}(n,\mathbb{C})\)-结构的结构张量在 Spencer 上同调中的体现。

（3）\textbf{辛结构}------\(G = \operatorname{Sp}(2n,\mathbb{R})\)。对应非退化闭 2-形式。由 Darboux 定理，辛结构总是局部可积（无局部不变量），但整体上存在形变（辛拓扑研究的内容）。辛结构的可积性源于 \(\mathfrak{sp}(2n,\mathbb{R})\) 的 Spencer 上同调 \(H^{1,2}(\mathfrak{sp}(2n,\mathbb{R})) = 0\)。

（4）\textbf{共形结构}------\(G = \operatorname{CO}(n) = \operatorname{O}(n) \times \mathbb{R}^+\)。对应共形度量类 \([g] = \{e^{2f} g : f \in C^\infty(M)\}\)。共形结构的可积性条件由 Weyl 张量 \(W\) 的消去刻画（\(n \geq 4\)）：\(W = 0\) 当且仅当度量局部共形平坦。在 \(n=3\) 时，Cotton 张量扮演类似角色。

（5）\textbf{\(\operatorname{SL}(n,\mathbb{R})\)-结构}------对应体积形式（或定向密度）。自然可积（任何定向流形上存在体积形式）。

（6）\textbf{\(\operatorname{Spin}(n)\)-结构}------\(\operatorname{SO}(n)\)-结构的提升。存在性等价于第二 Stiefel-Whitney 类 \(w_2(M) = 0\)。在弦理论和 Seiberg-Witten 理论中起核心作用。

（7）\textbf{\(\operatorname{G}_2\)-结构}（\(n=7\)）------对应 \(3\)-形式 \(\varphi\) 满足某种代数条件。\(\operatorname{G}_2\)-结构可积当且仅当 \(\varphi\) 是平行的（\(\nabla \varphi = 0\)），此时度量是 Ricci 平坦的。\(\operatorname{G}_2\)-流形在 M-理论紧化中起关键作用。

\textbf{定理 67.1.2（主丛上的联络与 \(G\)-结构的相容性）}：设 \(P \to M\) 是 \(G\)-结构。\(P\) 上的联络（即 \(G\)-主丛上的联络）与 \(G\)-结构相容当且仅当将该联络视为 \(FM\) 上的联络时，其 \(1\)-形式取值于 \(\mathfrak{g} \subset \mathfrak{gl}(n,\mathbb{R})\)。相容联络构成仿射空间，其向量空间为 \(\Omega^1(M, \operatorname{Ad} P)\)（\(\operatorname{Ad} P = P \times_G \mathfrak{g}\) 为伴随丛）。

\textbf{命题 67.1.3（挠率张量）}：\(G\)-结构 \(P\) 上的相容联络 \(\nabla\) 的挠率 \(T\) 是 \(TM\) 上的 \((1,2)\)-张量：\(T(X,Y) = \nabla_X Y - \nabla_Y X - [X,Y]\)。\(T\) 的取值在商丛 \((TM \otimes \mathfrak{gl}(n,\mathbb{R})/\mathfrak{g})\) 中的投影是 \(G\)-结构的内在不变量（不依赖于相容联络的选取）。

\hrulefill

\subsection{67.4 仿射微分几何}\label{ux4effux5c04ux5faeux5206ux51e0ux4f55}

仿射微分几何是介于 Riemann 几何与射影几何之间的基本几何结构，研究等体积仿射变换群（\(\operatorname{SL}(n,\mathbb{R}) \ltimes \mathbb{R}^n\)）下的曲面和超曲面不变量。与 Riemann 几何不同，仿射微分几何不预设度量结构，而是在仿射等价类中寻找规范化的几何量。

\textbf{定义 67.4.1（仿射法向量）}：设 \(\mathbf{x}: M \to \mathbb{R}^{n+1}\) 是非退化超曲面（即其 Blaschke 度量非退化）。仿射法向量 \(\mathbf{Y}\) 是满足以下条件的唯一向量场： (i) \(\mathbf{Y} \in \operatorname{Span}\{\mathbf{x}_{u_1}, \ldots, \mathbf{x}_{u_n}\}^\perp\)（即 \(\mathbf{Y}\) 与切空间正交，在仿射意义下） (ii) \(\det(\mathbf{x}_{u_1}, \ldots, \mathbf{x}_{u_n}, \mathbf{Y}) = 1\)（规范化条件，确保体积保持的仿射变换下等变）

\textbf{定义 67.4.2（仿射度量 / Blaschke 度量）}：仿射度量 \(g_{ij} = \langle \mathbf{x}_{u_i u_j}, \mathbf{Y} \rangle\) 是超曲面上的非退化伪 Riemann 度量。其符号差由 \(\mathbf{Y}\) 与 \(\mathbf{x}\) 的几何关系决定。仿射度量是仿射不变的（至多相差常数因子），因而是仿射微分几何中的基本不变量。

\textbf{定义 67.4.3（仿射形状算子）}：仿射形状算子 \(S\) 定义为 \(\mathbf{Y}_{u_i} = -\sum_j S_i^j \mathbf{x}_{u_j}\)。\(S\) 的特征值 \(\lambda_1, \ldots, \lambda_n\) 称为\textbf{仿射主曲率}，其平均值 \(H = \frac{1}{n} \operatorname{tr} S\) 称为\textbf{仿射平均曲率}。

\textbf{定理 67.4.1（仿射基本定理）}：给定非退化伪 Riemann 度量 \(g\) 和 \((1,1)\)-张量 \(S\)（满足 Gauss-Codazzi 型方程），可唯一确定 \(\mathbb{R}^{n+1}\) 中的非退化超曲面（至多相差仿射变换），使得 \(g\) 为仿射度量、\(S\) 为仿射形状算子。

\textbf{证明}：仿射基本定理的证明与 Riemann 几何中的曲面论基本定理（Bonnet 定理）平行。给定 \(g\) 和 \(S\)，需构造 \(\mathbb{R}^{n+1}\) 中的浸入 \(\mathbf{x}\) 和仿射法向量 \(\mathbf{Y}\) 满足： (1) \(\mathbf{x}_{u_i u_j} = \sum_k \Gamma_{ij}^k \mathbf{x}_{u_k} + g_{ij} \mathbf{Y}\) (2) \(\mathbf{Y}_{u_i} = -\sum_j S_i^j \mathbf{x}_{u_j}\) 其中 \(\Gamma_{ij}^k\) 是 \(g\) 的 Christoffel 符号。该 PDE 系统的可积性条件恰好是 Gauss 方程（\(R_{ijkl} = g_{ik} S_{jl} + g_{jl} S_{ik} - g_{il} S_{jk} - g_{jk} S_{il}\)）和 Codazzi 方程（\(\nabla_i S_j^k = \nabla_j S_i^k\)）。由 Frobenius 定理（或 Cartan-Kahler 定理），可积性条件满足时解存在且唯一（至多相差仿射变换）。\(\blacksquare\)

\textbf{定理 67.4.2（仿射球面分类）}：仿射主曲率为常数的非退化仿射超曲面分为三类： - \textbf{椭球型}：所有仿射主曲率同号（\(S = \lambda I, \lambda > 0\)） - \textbf{双曲型}：仿射主曲率异号（\(S = \lambda I, \lambda < 0\)） - \textbf{抛物型}：\(S = 0\)（仿射主曲率全为零）

\textbf{定理 67.4.3（Calabi 猜想，仿射版本）}：完备仿射球面（即 \(S = \lambda I\)）必为椭球面（\(\lambda > 0\)）或双曲面（\(\lambda < 0\)），即不存在完备的抛物型仿射球面。由 Calabi (1982) 提出、Trudinger-Wang (2000) 完全解决。

\textbf{证明}（框架）：Calabi 猜想的证明使用了 Monge-Ampere 方程 \(\det(\nabla^2 u) = 1\) 的整体解理论，其中 \(u\) 是超曲面的支撑函数。局部地，仿射球面可表示为图 \(x_{n+1} = u(x_1, \ldots, x_n)\)，满足 \(\det(\nabla^2 u) = \left(\frac{\lambda}{H}\right)^{n+2}\)。当 \(\lambda = 0\)（抛物型）时，方程为 \(\det(\nabla^2 u) = 0\)，解为''广义抛物面''。Trudinger-Wang 证明完备的抛物型仿射球面必为椭圆抛物面，但进一步分析显示不存在完备的非退化抛物型仿射球面。证明的关键是建立 \(u\) 的 \(C^{2,\alpha}\) 先验估计，利用 Pogorelov 型内估计和 Calabi 的 \(C^3\) 估计。\(\blacksquare\)

仿射微分几何在信息几何（统计流形，其中 Fisher 信息度量是 Riemann 度量而 Amari-Chentsov 张量是仿射联络）和计算机视觉（仿射不变物体识别，利用仿射不变量对透视变换的鲁棒性）中有重要应用。

\subsection{67.5 Twistor 理论与超 Kahler 几何}\label{twistor-ux7406ux8bbaux4e0eux8d85-kahler-ux51e0ux4f55}

Twistor 理论由 Penrose (1967) 创立，通过构造 \(\mathbb{R}^4\)（共形紧化为 \(S^4\)）上的 \textbf{twistor 空间}------\(\mathbb{CP}^3\)（即复三维射影空间）------将四维 Riemann 几何问题转化为复几何问题。

\textbf{定义 67.5.1（Twistor 空间）}：设 \((M, g)\) 是四维可定向 Riemann 流形。考虑 \(M\) 上的近复结构 \(J\) 与 \(g\) 相容（即 \(J\) 是 \(g\) 的正交近复结构）。在任意点 \(p \in M\)，\(T_p M\) 上与 \(g\) 相容的近复结构构成 \(S^2 = \operatorname{SO}(4)/\operatorname{U}(2)\)。\(M\) 的 \textbf{twistor 空间} \(Z(M)\) 定义为 \(M\) 上与 \(g\) 相容的近复结构的纤维丛（纤维为 \(S^2\)）。\(Z(M)\) 具有自然近复结构，其可积性由 \(M\) 的曲率控制。

\textbf{定理 67.5.1（Atiyah-Hitchin-Singer, 1978）}：\(Z(M)\) 上的自然近复结构可积当且仅当 \((M, g)\) 是反自偶的（即 Weyl 张量的自偶部分 \(W^+ = 0\)）。此时 \(Z(M)\) 是三维复流形。

\textbf{定理 67.5.2（Penrose 非线性引力子构造，1976）}：反自偶 Einstein 度量（即 \(W^+ = 0\) 且 \(\operatorname{Ric} = \lambda g\)）一一对应于 twistor 空间 \(Z\) 上具有全纯切触结构的形变。具体地，\(Z\) 是三维复流形，带有一个全纯切触分布（即 \(\mathcal{O}(2)\)-值的全纯 \(1\)-形式 \(\theta\) 满足 \(\theta \wedge d\theta = 0\)），且 \(Z\) 被一族有理曲线（twistor 线）叶状化。

\textbf{定义 67.5.2（超 Kahler 几何）}：超 Kahler 流形是具有三个相容复结构 \(I, J, K\)（满足四元数关系 \(IJ = K = -JI\)）的 Riemann 流形 \((M, g)\)，且相应的 Kahler 形式 \(\omega_I, \omega_J, \omega_K\) 均闭。等价地，\((M, g)\) 是 Ricci 平坦的 Kahler 流形且其和乐群包含于 \(\operatorname{Sp}(n)\)（\(4n\) 维实流形）。

\textbf{定理 67.5.3（超 Kahler 流形的 Twistor 空间）}：超 Kahler 流形 \((M, g, I, J, K)\) 的 twistor 空间 \(Z\) 是 \(M \times \mathbb{CP}^1\)（作为微分流形），其复结构为：在 \(p \in M\) 上的纤维 \(\mathbb{CP}^1\) 的点 \(\zeta\) 对应于近复结构 \(aI + bJ + cK\)（其中 \(a^2 + b^2 + c^2 = 1\)）。\(Z\) 上的全纯辛形式的存在性由超 Kahler 流形的 Ricci 平坦性保证。

经典例子： （1）\textbf{\(K3\) 曲面}------唯一紧致四维超 Kahler 流形。\(K3\) 曲面是单连通的紧复曲面，具有平凡典范丛，其 Yau 定理（Calabi 猜想的解决）保证存在 Ricci 平坦 Kahler 度量，从而赋予 \(K3\) 曲面超 Kahler 结构。 （2）\textbf{ALE 空间}（渐近局部 Euclid 空间）------Kronheimer (1989) 分类了所有 ALE 超 Kahler \(4\)-流形，它们与 Dynkin 图 \(A_n, D_n, E_6, E_7, E_8\) 一一对应。ALE 超 Kahler 度量由 Nakajima 箭图簇的 Hyper-Kahler 商构造给出。 （3）\textbf{Hilbert 概型 \(\operatorname{Hilb}^n(K3)\) 与广义 Kummer 簇}------高维紧致超 Kahler 流形的两类基本构造。Beauville (1983) 证明 \(\operatorname{Hilb}^n(K3)\)（\(K3\) 曲面上 \(n\) 个点的 Hilbert 概型）是 \(4n\) 维紧致超 Kahler 流形。 （4）\textbf{磁单极模空间}------Atiyah-Hitchin (1988) 构造了 \(\mathbb{R}^3\) 上 \(SU(2)\) 磁单极的模空间，是 \(4k\) 维超 Kahler 流形（\(k\) 为磁荷）。

\textbf{定理 67.5.4（Hitchin-Kobayashi 对应）}：紧 Kahler 流形上的稳定 Higgs 丛与 Hermite-Einstein 联络之间存在一一对应。这一对应是 Narasimhan-Seshadri 定理（紧 Riemann 面上的稳定向量丛与平坦酉联络对应）的推广，其超 Kahler 版本由 Hitchin (1987) 在 Riemann 面上建立，模空间是超 Kahler 流形。

\hrulefill

\section{第226章：张量初步}\label{ux7b2c226ux7ae0ux5f20ux91cfux521dux6b65}

张量是多重线性代数中的核心概念，它统一了向量、线性泛函、双线性型和更一般的多重线性型。在微分几何中，张量场是流形上最基本的几何对象，而张量代数（包括外代数和对称代数）是 de Rham 上同调、Hodge 理论和 Riemann 几何的技术基础。

\subsection{68.1 多重线性映射}\label{ux591aux91cdux7ebfux6027ux6620ux5c04}

\textbf{定义 68.1.1}（多重线性映射）：设 \(V_1, \ldots, V_k, W\) 为 \(F\) 上的向量空间。映射 \(T : V_1 \times \cdots \times V_k \to W\) 称为\textbf{多重线性映射}（\(k\)-线性映射），若 \(T\) 对每个变元是线性的（固定其他变元时）。所有 \(k\)-线性映射 \(V_1 \times \cdots \times V_k \to W\) 构成向量空间，记为 \(\mathcal{L}(V_1, \ldots, V_k; W)\)。

\textbf{定义 68.1.2}（对偶空间）：设 \(V\) 为 \(F\) 上的向量空间。\(V\) 上的所有线性泛函（即 \(V \to F\) 的线性映射）构成的向量空间称为 \(V\) 的\textbf{对偶空间}，记为 \(V^* = \mathcal{L}(V, F)\)。

\textbf{定理 68.1.1}：若 \(\dim V = n\)，则 \(\dim V^* = n\)。若 \(B = \{e_1, \ldots, e_n\}\) 为 \(V\) 的基，则 \(V^*\) 有一组\textbf{对偶基} \(\{e^1, \ldots, e^n\}\)，满足 \(e^i(e_j) = \delta_{ij}\)。

\textbf{证明}：\(\dim V^* = \dim \mathcal{L}(V, F) = (\dim V)(\dim F) = n\)。定义 \(e^i(\sum_j x_j e_j) = x_i\)，则 \(e^i(e_j) = \delta_{ij}\)。\(\{e^1, \ldots, e^n\}\) 线性无关（若 \(\sum \lambda_i e^i = 0\)，作用在 \(e_j\) 上得 \(\lambda_j = 0\)），且恰有 \(n\) 个，故为基。\(\blacksquare\)

\textbf{定理 68.1.2}（双重对偶）：存在自然同构 \(V \cong V^{**}\)，由 \(v \mapsto (\operatorname{ev}_v: \alpha \mapsto \alpha(v))\) 给出。在有限维情形下这是同构，在无穷维情形下是单射。

\textbf{证明}：定义 \(\Phi: V \to V^{**}\) 为 \(\Phi(v)(\alpha) = \alpha(v)\)。\(\Phi\) 是线性映射。若 \(\Phi(v) = 0\)，则对所有 \(\alpha \in V^*\) 有 \(\alpha(v) = 0\)。取 \(v\) 在基下的坐标，若 \(v \neq 0\)，则存在 \(e^i\) 使 \(e^i(v) \neq 0\)，矛盾。故 \(\Phi\) 是单射。\(\dim V^{**} = \dim V^* = \dim V\)，故单射即为同构。\(\blacksquare\)

\hrulefill

\subsection{68.2 张量积}\label{ux5f20ux91cfux79ef}

\textbf{定义 68.2.1}（张量积的泛性质定义）：设 \(V, W\) 为 \(F\) 上的向量空间。\(V\) 与 \(W\) 的\textbf{张量积}是一个向量空间 \(V \otimes W\) 加上一个双线性映射 \(\otimes : V \times W \to V \otimes W\)（\((v, w) \mapsto v \otimes w\)），满足以下\textbf{泛性质}：

对任意双线性映射 \(B : V \times W \to U\)（\(U\) 为任意向量空间），存在唯一的线性映射 \(\tilde{B} : V \otimes W \to U\) 使得下图交换：

\[B = \tilde{B} \circ \otimes\]

即 \(B(v, w) = \tilde{B}(v \otimes w)\)。

泛性质定义的好处在于：它不依赖于具体构造，而是刻画了张量积在同构下的唯一性。任何满足泛性质的构造都是彼此自然同构的。

\textbf{定理 68.2.1}（张量积的存在性）：对任何有限维向量空间 \(V, W\)，张量积 \(V \otimes W\) 存在，且在同构意义下唯一。

\textbf{证明}（构造）：设 \(\dim V = m\)，\(\dim W = n\)。取 \(V\) 的基 \(\{e_1, \ldots, e_m\}\)，\(W\) 的基 \(\{f_1, \ldots, f_n\}\)。定义 \(V \otimes W\) 为以 \(\{e_i \otimes f_j \mid 1 \leq i \leq m, 1 \leq j \leq n\}\) 为基的 \(mn\) 维向量空间。对 \(v = \sum a_i e_i\)，\(w = \sum b_j f_j\)，定义 \(v \otimes w = \sum_{i,j} a_i b_j (e_i \otimes f_j)\)。

验证泛性质：设 \(B : V \times W \to U\) 为双线性映射。定义 \(\tilde{B}(e_i \otimes f_j) = B(e_i, f_j)\)，并线性扩充到 \(V \otimes W\)。则 \(\tilde{B}(v \otimes w) = \tilde{B}(\sum a_i b_j e_i \otimes f_j) = \sum a_i b_j B(e_i, f_j) = B(\sum a_i e_i, \sum b_j f_j) = B(v, w)\)。\(\tilde{B}\) 的唯一性由基上的值唯一确定。\(\blacksquare\)

\textbf{定理 68.2.2}（张量积的基本性质）：设 \(V, W, U\) 为有限维向量空间。

\begin{enumerate}
\def\labelenumi{(\roman{enumi})}
\tightlist
\item
  \(\dim(V \otimes W) = (\dim V)(\dim W)\)
\item
  \(V \otimes W \cong W \otimes V\)（通过 \(v \otimes w \mapsto w \otimes v\)）
\item
  \((V \otimes W) \otimes U \cong V \otimes (W \otimes U)\)
\item
  \(V \otimes (W \oplus U) \cong (V \otimes W) \oplus (V \otimes U)\)
\item
  \(F \otimes V \cong V\)
\item
  \((V \otimes W)^* \cong V^* \otimes W^*\)（有限维情形）
\end{enumerate}

\textbf{证明}：(i) 由构造中基的个数 \(mn\) 得到。(ii) 定义映射 \(\tau : V \otimes W \to W \otimes V\) 为 \(\tau(v \otimes w) = w \otimes v\)，线性扩充。其逆为 \(\tau^{-1}(w \otimes v) = v \otimes w\)。(iii) 定义 \((v \otimes w) \otimes u \mapsto v \otimes (w \otimes u)\)，线性扩充。由泛性质验证这是良定义的同构。(iv) 由泛性质验证：\(V \otimes (W \oplus U)\) 满足余积的泛性质。(v) 定义 \(\lambda \otimes v \mapsto \lambda v\)，其逆为 \(v \mapsto 1 \otimes v\)。(vi) 取基后验证维数相同，且存在自然同构 \(\alpha \otimes \beta \mapsto (v \otimes w \mapsto \alpha(v)\beta(w))\)。该映射是良定义的（由张量积的泛性质），且是单射（维数相同）。\(\blacksquare\)

\textbf{定义 68.2.2}（张量）：\(k\) 个 \(V\) 的张量积 \(V^{\otimes k} = V \otimes V \otimes \cdots \otimes V\)（\(k\) 次）中的元素称为 \(V\) 上的 \(k\) 阶\textbf{张量}。特别地： - \(V^{\otimes 0} = F\)（标量） - \(V^{\otimes 1} = V\)（向量） - \(V^{\otimes 2} = V \otimes V\)（二阶张量）

\textbf{定义 68.2.3}（混合张量）：\(V\) 上的 \((p, q)\)-型张量是 \(V^{\otimes p} \otimes (V^*)^{\otimes q}\) 中的元素。在基 \(\{e_i\}\) 下，\((p, q)\)-张量 \(T\) 有分量表示： \[T = T^{i_1 \cdots i_p}_{j_1 \cdots j_q} e_{i_1} \otimes \cdots \otimes e_{i_p} \otimes e^{j_1} \otimes \cdots \otimes e^{j_q}\] 基变换下分量的变换规则为（使用 Einstein 求和约定）： \[T'^{i_1 \cdots i_p}_{j_1 \cdots j_q} = A^{i_1}_{k_1} \cdots A^{i_p}_{k_p} (A^{-1})^{\ell_1}_{j_1} \cdots (A^{-1})^{\ell_q}_{j_q} T^{k_1 \cdots k_p}_{\ell_1 \cdots \ell_q}\]

\textbf{定义 68.2.4（张量缩并）}：缩并是 \((p, q)\)-张量到 \((p-1, q-1)\)-张量的线性映射，将第 \(r\) 个上指标与第 \(s\) 个下指标配对并求和。在分量表示下，缩并 \(C^r_s\) 作用于 \(T\) 得到： \[(C^r_s T)^{i_1 \cdots \widehat{i_r} \cdots i_p}_{j_1 \cdots \widehat{j_s} \cdots j_q} = T^{i_1 \cdots k \cdots i_p}_{j_1 \cdots k \cdots j_q}\] 其中 \(k\) 求和。缩并的最重要例子是矩阵的迹（\(V \otimes V^* \to F\)）。

\textbf{定理 68.2.3（张量积与线性映射）}：存在自然同构 \(V^* \otimes W \cong \mathcal{L}(V, W)\)（有限维），由 \(\alpha \otimes w \mapsto (v \mapsto \alpha(v)w)\) 给出。

\textbf{证明}：定义 \(\Phi: V^* \otimes W \to \mathcal{L}(V, W)\) 为 \(\Phi(\alpha \otimes w)(v) = \alpha(v)w\)，并线性扩充。取基 \(\{e_i\}\) 和 \(\{f_j\}\)，则 \(\Phi(e^i \otimes f_j)\) 将 \(e_k\) 映为 \(\delta_{ik} f_j\)。这些映射构成 \(\mathcal{L}(V, W)\) 的基，故 \(\Phi\) 是同构。\(\blacksquare\)

\hrulefill

\subsection{68.3 外代数}\label{ux5916ux4ee3ux6570}

\textbf{定义 68.3.1}（外幂）：设 \(V\) 为 \(F\) 上的向量空间。\(V\) 的第 \(k\) 次\textbf{外幂} \(\bigwedge^k V\) 是 \(V^{\otimes k}\) 的商空间，模去所有形如 \(v_1 \otimes \cdots \otimes v_k\) 且 \(v_i = v_j\)（\(i \neq j\)）的元素生成的子空间。\(\bigwedge^k V\) 中的元素记为 \(v_1 \wedge v_2 \wedge \cdots \wedge v_k\)，称为 \(k\)-向量。

等价地，\(\bigwedge^k V\) 可由 \(V^{\otimes k}\) 通过交错化算子 \(\operatorname{Alt}: V^{\otimes k} \to V^{\otimes k}\) 定义为 \[\operatorname{Alt}(v_1 \otimes \cdots \otimes v_k) = \frac{1}{k!} \sum_{\sigma \in S_k} \operatorname{sgn}(\sigma) v_{\sigma(1)} \otimes \cdots \otimes v_{\sigma(k)}\] 则 \(\bigwedge^k V \cong \operatorname{Alt}(V^{\otimes k})\)（在特征不为 \(2\) 时）。

\textbf{定理 68.3.1}（外幂的基本性质）：

\begin{enumerate}
\def\labelenumi{(\roman{enumi})}
\tightlist
\item
  \textbf{交错性}：\(v_1 \wedge \cdots \wedge v_i \wedge \cdots \wedge v_j \wedge \cdots \wedge v_k = -v_1 \wedge \cdots \wedge v_j \wedge \cdots \wedge v_i \wedge \cdots \wedge v_k\)
\item
  若 \(\{e_1, \ldots, e_n\}\) 为 \(V\) 的基，则 \(\{e_{i_1} \wedge e_{i_2} \wedge \cdots \wedge e_{i_k} \mid 1 \leq i_1 < i_2 < \cdots < i_k \leq n\}\) 为 \(\bigwedge^k V\) 的基
\item
  \(\dim \bigwedge^k V = \binom{n}{k}\)（\(0 \leq k \leq n\)），\(\bigwedge^k V = \{0\}\)（\(k > n\)）
\item
  \(\bigwedge^n V\) 是一维的，基为 \(e_1 \wedge \cdots \wedge e_n\)
\end{enumerate}

\textbf{证明}：(i) 由定义，\(v_i \wedge v_j = -v_j \wedge v_i\)（因为 \((v_i + v_j) \wedge (v_i + v_j) = 0\) 展开得 \(v_i \wedge v_i + v_i \wedge v_j + v_j \wedge v_i + v_j \wedge v_j = v_i \wedge v_j + v_j \wedge v_i = 0\)）。(ii) 由交错性，任何 \(k\)-向量可写为严格递增指标对应的基向量的外积的线性组合。线性无关性由构造保证。(iii) 由 (ii)，基的个数为 \(\binom{n}{k}\)。(iv) \(k = n\) 时 \(\binom{n}{n} = 1\)。\(\blacksquare\)

\textbf{定义 68.3.2}（外代数）：\(V\) 的\textbf{外代数}（Grassmann 代数）为直和

\[\bigwedge V = \bigoplus_{k=0}^{n} \bigwedge^k V\]

外积 \(\wedge\) 使 \(\bigwedge V\) 成为一个结合代数（满足 \(a \wedge b = (-1)^{(\deg a)(\deg b)} b \wedge a\)）。\(\bigwedge V\) 的维数为 \(\sum_{k=0}^{n} \binom{n}{k} = 2^n\)。

\textbf{定理 68.3.2}（行列式与外积的关系）：设 \(v_1, \ldots, v_n \in \mathbb{R}^n\)。则

\[v_1 \wedge v_2 \wedge \cdots \wedge v_n = \det(v_1, \ldots, v_n) \, e_1 \wedge e_2 \wedge \cdots \wedge e_n\]

\textbf{证明}：将 \(v_j = \sum_i a_{ij} e_i\) 代入外积，使用交错性展开。外积的系数正是 Leibniz 公式定义的 \(\det(a_{ij})\)。\(\blacksquare\)

\textbf{定义 68.3.3（内乘 / 内部积）}：设 \(\alpha \in V^*\)。\(\alpha\) 的内乘 \(i_\alpha: \bigwedge^k V \to \bigwedge^{k-1} V\) 定义为 \[i_\alpha(v_1 \wedge \cdots \wedge v_k) = \sum_{i=1}^{k} (-1)^{i-1} \alpha(v_i) v_1 \wedge \cdots \wedge \widehat{v_i} \wedge \cdots \wedge v_k\] 其中 \(\widehat{v_i}\) 表示省略 \(v_i\)。\(i_\alpha\) 是外代数上的反导子：\(i_\alpha(a \wedge b) = i_\alpha(a) \wedge b + (-1)^{\deg a} a \wedge i_\alpha(b)\)。

\textbf{定义 68.3.4（Hodge 星算子，线性代数版本）}：设 \(V\) 是 \(n\) 维内积空间（带定向）。Hodge 星算子 \(\star: \bigwedge^k V \to \bigwedge^{n-k} V\) 由以下条件唯一确定：对任意 \(\alpha, \beta \in \bigwedge^k V\)， \[\alpha \wedge \star \beta = \langle \alpha, \beta \rangle \operatorname{vol}\] 其中 \(\operatorname{vol} = e_1 \wedge \cdots \wedge e_n\) 是定向体积形式，\(\langle \cdot, \cdot \rangle\) 是 \(\bigwedge^k V\) 上由 \(V\) 的内积诱导的内积。\(\star\) 满足 \(\star^2 = (-1)^{k(n-k)}\)（在 \(\bigwedge^k V\) 上）。

\hrulefill

\subsection{68.4 对称张量}\label{ux5bf9ux79f0ux5f20ux91cf}

\textbf{定义 68.4.1}（对称张量）：\(T \in V^{\otimes k}\) 称为\textbf{对称张量}，若对任意置换 \(\sigma \in S_k\)，有 \(\sigma \cdot T = T\)（其中 \(\sigma\) 通过置换因子作用在 \(V^{\otimes k}\) 上）。

\textbf{定义 68.4.2}（对称幂）：\(V\) 的第 \(k\) 次\textbf{对称幂} \(S^k V\) 是 \(V^{\otimes k}\) 的商空间，模去形如 \(v_1 \otimes \cdots \otimes v_k - v_{\sigma(1)} \otimes \cdots \otimes v_{\sigma(k)}\) 的元素生成的子空间。\(S^k V\) 中的元素记为 \(v_1 \odot v_2 \odot \cdots \odot v_k\)。

等价地，\(S^k V\) 可由对称化算子 \(\operatorname{Sym}: V^{\otimes k} \to V^{\otimes k}\) 定义为 \[\operatorname{Sym}(v_1 \otimes \cdots \otimes v_k) = \frac{1}{k!} \sum_{\sigma \in S_k} v_{\sigma(1)} \otimes \cdots \otimes v_{\sigma(k)}\] 则 \(S^k V \cong \operatorname{Sym}(V^{\otimes k})\)（在特征为 \(0\) 时）。

\textbf{定理 68.4.1}：若 \(\{e_1, \ldots, e_n\}\) 为 \(V\) 的基，则 \(\{e_{i_1} \odot e_{i_2} \odot \cdots \odot e_{i_k} \mid 1 \leq i_1 \leq i_2 \leq \cdots \leq i_k \leq n\}\) 为 \(S^k V\) 的基。\(\dim S^k V = \binom{n + k - 1}{k}\)。

\textbf{证明}：由对称性，任何对称张量可写为非降指标对应的基向量的对称积的线性组合。基的个数为从 \(n\) 个元素中可重复地选取 \(k\) 个的组合数，即 \(\binom{n+k-1}{k}\)。\(\blacksquare\)

\textbf{定义 68.4.3}（对称代数）：\(V\) 的\textbf{对称代数} \(S(V) = \bigoplus_{k=0}^{\infty} S^k V\) 是 \(\bigwedge V\) 的''交换''类比。\(S(V)\) 自然同构于 \(V\) 上的多项式代数 \(F[V^*]\)。

\textbf{定义 68.4.4}（多重线性型与张量）：\(V\) 上的 \(k\) 重线性型自然地对应于 \((V^*)^{\otimes k}\) 中的元素，或等价地，\(V^{\otimes k}\) 上的线性泛函。特别地，\(V\) 上的对称双线性型对应于 \(S^2(V^*)\) 中的元素，这正是二次型的理论框架。反对称双线性型对应于 \(\bigwedge^2(V^*)\) 中的元素。

\textbf{定理 68.4.2（张量代数的泛性质）}：张量代数 \(T(V) = \bigoplus_{k=0}^{\infty} V^{\otimes k}\) 是 \(V\) 上的自由结合代数：任何线性映射 \(V \to A\)（\(A\) 为结合代数）可唯一地扩充为代数同态 \(T(V) \to A\)。外代数 \(\bigwedge V\) 和对称代数 \(S(V)\) 分别是 \(T(V)\) 关于理想 \((v \otimes v)\) 和 \((v \otimes w - w \otimes v)\) 的商代数。

\textbf{定理 68.4.3（分次交换律）}：\(T(V)\) 是 \(\mathbb{N}\)-分次代数，\(V^{\otimes p} \otimes V^{\otimes q} \subseteq V^{\otimes (p+q)}\)。\(\bigwedge V\) 和 \(S(V)\) 也是分次代数，分别满足 \(a \wedge b = (-1)^{pq} b \wedge a\) 和 \(a \odot b = b \odot a\)（对所有 \(a \in \bigwedge^p V, b \in \bigwedge^q V\) 和 \(a \in S^p V, b \in S^q V\)）。

\textbf{注记 68.4.1}：张量与外代数的深入理论（包括流形上的张量场、微分形式、de Rham 上同调等）将在后续章节中详细展开。本章仅建立有限维向量空间上的基本代数框架。\(\blacksquare\)

\hrulefill

\subsection{68.5 张量代数与表示论的联系}\label{ux5f20ux91cfux4ee3ux6570ux4e0eux8868ux793aux8bbaux7684ux8054ux7cfb}

\textbf{定理 68.5.1（Schur-Weyl 对偶）}：设 \(V\) 是 \(n\) 维复向量空间。\(V^{\otimes k}\) 上同时有 \(S_k\)（置换因子）和 \(\operatorname{GL}(V)\)（对角作用）的作用。这两个作用彼此交换，且生成的全代数与 \(S_k\) 的群代数和 \(\operatorname{GL}(V)\) 的包络代数的双中心化子有关。\(V^{\otimes k}\) 分解为不可约表示的直和： \[V^{\otimes k} \cong \bigoplus_{\lambda \vdash k, \ell(\lambda) \leq n} \mathbb{S}_\lambda(V) \otimes \operatorname{Specht}_\lambda\] 其中 \(\lambda\) 是 \(k\) 的划分，\(\mathbb{S}_\lambda(V)\) 是 \(\operatorname{GL}(V)\) 的不可约表示（Schur 函子），\(\operatorname{Specht}_\lambda\) 是 \(S_k\) 的不可约表示（Specht 模）。外幂 \(\bigwedge^k V\) 对应于 \(\lambda = (1^k)\)，对称幂 \(S^k V\) 对应于 \(\lambda = (k)\)。

Schur-Weyl 对偶将张量代数的表示论与对称群的表示论统一起来，是经典不变量理论的核心结果。在微分几何中，这一分解用于定义 Weyl 张量（\(S^2(\bigwedge^2 V)\) 中与 Ricci 曲率正交的部分）和曲率张量的代数分类。

\hrulefill

\hrulefill

\hrulefill

\hrulefill

\hrulefill

\hrulefill

\newpage

\chapter{07-拓扑与几何/02-微分几何/04-辛几何.md}\label{ux62d3ux6251ux4e0eux51e0ux4f5502-ux5faeux5206ux51e0ux4f5504-ux8f9bux51e0ux4f55.md}

\subsection{163.A 辛几何补充}\label{a-ux8f9bux51e0ux4f55ux8865ux5145}

辛几何是 Hamilton 力学的天然数学语言，其研究偶数维光滑流形上的非退化闭 2-形式（辛形式），切触几何则是奇数维流形上的对应理论。辛几何的现代复兴始于 1960 年代------Vladimir Arnold 关于经典力学中 Liouville-Arnold 可积系统的研究，以及随后由 Gromov、Eliashberg、Floer 和 Kontsevich 等人引领的一系列革命性进展，使辛几何从分析力学的附庸成长为微分几何和数学物理的核心分支。辛流形不同于 Riemann 流形的关键特征在于其''刚性''与''柔性''的辩证关系：Darboux 定理表明所有同样维度的辛流形局部同构（无局部不变量），这与 Riemann 几何的曲率形成鲜明对比；然而全局辛不变量（Gromov 非挤压定理、辛容量和 Floer 同调）却极为丰富，构成了''辛拓扑''的独特景观。本章作为微分几何的补充，从辛流形基础（第164章）出发，系统建立辛约化与矩映射（第165章）、伪全纯曲线与 Gromov-Witten 理论（第166章）以及切触几何（第167章）的完整体系。

\hrulefill

\hrulefill

\hrulefill

\hrulefill

\hrulefill

\section{第227章：辛流形基础}\label{ux7b2c227ux7ae0ux8f9bux6d41ux5f62ux57faux7840}

辛流形是经典 Hamilton 力学的自然几何语言------相空间（位置和动量的联合空间）上携带的辛形式 \(\omega = \sum dp_i \wedge dq_i\) 在 Darboux 定理的保证下局部恒定，而 Hamilton 方程 \(\dot{q} = \partial H/\partial p\)、\(\dot{p} = -\partial H/\partial q\) 等价于向量场 \(X_H\) 满足 \(\iota_{X_H} \omega = dH\)。从数学结构看，辛形式与 Riemann 度量的根本差异在于：前者是反对称的闭 2-形式，曲率不再是局部不变量（Darboux 定理），而后者是对称的正定 2-张量，曲率是局部不变量。这一差异导致了辛几何中全局不变量的极端丰富------Gromov 非挤压定理、辛容量以及 Floer 同调------构成了''辛拓扑''的独特世界。在分析层面，辛流形的理论基石包括：辛线性代数（非退化反对称双线性型必有标准 \(2n \times 2n\) 矩阵形式）、几乎复结构 \(J\) 与 \(\omega\) 的相容性（构成了 Kaehler 流形在复几何中的位置，对比 Hodge 理论与 Calabi-Yau 流形，见第199-201章），以及 Lagrange 子流形（维数恰为辛流形一半且 \(\omega\) 限制为零）在 Floer 同调和镜对称中的核心角色。

\subsection{242.1 辛线性代数}\label{ux8f9bux7ebfux6027ux4ee3ux6570}

\textbf{定义 242.1}（辛向量空间）：有限维实向量空间 \(V\) 上的\textbf{辛形式} \(\omega: V \times V \to \mathbb{R}\) 是双线性、反对称（\(\omega(u, v) = -\omega(v, u)\)）且非退化（\(\omega(u, v) = 0\) 对所有 \(v\) 蕴含 \(u = 0\)）的。\((V, \omega)\) 称为\textbf{辛向量空间}。

\textbf{命题 242.1}（辛向量空间的标准型 / 维数）：任何辛向量空间的维数是偶数：\(\dim V = 2n\)。存在辛基 \((e_1, \ldots, e_n, f_1, \ldots, f_n)\) 使得 \(\omega(e_i, f_j) = \delta_{ij}\)，\(\omega(e_i, e_j) = \omega(f_i, f_j) = 0\)。辛形式在辛基下的矩阵为 \(\begin{pmatrix} 0 & I \\ -I & 0 \end{pmatrix}\)。

\textbf{证明}：设 \(\dim V = m\)，任选基下 \(\omega\) 的矩阵 \(A\) 为反对称矩阵（\(A^T = -A\)）。非退化性 \(\iff \det A \neq 0\)。对任意反对称矩阵 \(A\)，\(\det A = \det(A^T) = \det(-A) = (-1)^m \det A\)。当 \(m\) 为奇数时，\(\det A = -\det A \Rightarrow \det A = 0\)，与非退化矛盾。故 \(m\) 必为偶数 \(2n\)。辛基的存在性由归纳构造：取任意 \(e_1 \neq 0\)，由非退化性存在 \(f_1\) 使 \(\omega(e_1, f_1) = 1\)，在辛正交补 \(\{e_1, f_1\}^\omega\) 上继续归纳即得。∎

命题 242.1 确立了辛向量空间的偶数维性和辛基的标准型，这是辛几何一切后期构造的代数基础。有了辛基的概念，自然可以引入保持辛结构的线性变换群。

\textbf{定义 242.2}（辛群 \(\operatorname{Sp}(2n)\)）：\textbf{辛群} \(\operatorname{Sp}(2n, \mathbb{R}) = \{A \in \operatorname{GL}(2n, \mathbb{R}) : A^T \Omega A = \Omega\}\)，其中 \(\Omega = \begin{pmatrix} 0 & I \\ -I & 0 \end{pmatrix}\)。

\subsection{242.2 辛流形}\label{ux8f9bux6d41ux5f62}

辛流形是辛几何的核心研究对象，它推广了经典力学中相空间的结构。与 Riemann 流形不同，辛流形在局部上完全没有不变量------Darboux 定理断言所有同维数的辛流形局部辛同构于 \((\mathbb{R}^{2n}, \sum dx_i \wedge dy_i)\)。这一性质与 Riemann 几何中曲率的存在性形成鲜明对比：在 Riemann 几何中，截面曲率是局部不变量，而在辛几何中，任何局部不变量都必须退化。辛流形的非退化性 \(\omega^n \neq 0\) 蕴涵了 \(\dim M = 2n\)（偶数维），且 \(\omega^n/n!\) 是 \(M\) 上的典范体积形式------这直接导致了 Hamilton 力学中的 Liouville 定理（相空间体积守恒）。辛形式 \(\omega\) 通过降指标运算 \(\omega^\flat: TM \to T^*M\)（\(v \mapsto \omega(v, \cdot)\)）建立了切丛与余切丛之间的同构，从而将 Hamilton 函数 \(H\) 与 Hamilton 向量场 \(X_H = (\omega^\flat)^{-1}(dH)\) 一一对应。这一同构是 Hamilton 力学与辛几何之间桥梁的数学基础。辛流形的基本例子包括：余切丛 \(T^*Q\)（配以典范辛形式 \(\omega_{\text{can}} = -d\theta\)，\(\theta\) 为 Liouville 1-形式），Kähler 流形（其 Kähler 形式为辛形式），以及余伴随轨道（Kirillov-Kostant-Souriau 辛形式）。辛流形上的 Lagrangian 子流形（维数为 \(n\) 且 \(\omega\) 限制为零的子流形）是辛几何中最重要的子对象，它们构成了 Floer 同调和 Fukaya 范畴的基本对象。与 Riemann 几何的另一个关键区别在于：辛流形上不存在局部不变量意味着辛几何的刚性完全体现在全局（整体）现象中，如 Gromov 非压缩定理和辛嵌入障碍。

\textbf{定义 242.3}（辛流形）：\textbf{辛流形} \((M, \omega)\) 是光滑流形 \(M\) 配备闭（\(d\omega = 0\)）且非退化（\(\omega^n \neq 0\) 处处）的 2-形式 \(\omega \in \Omega^2(M)\)。\(\dim M = 2n\)。

\textbf{定理 242.1}（Darboux 定理）：任何辛流形 \((M^{2n}, \omega)\) 局部同构于 \((\mathbb{R}^{2n}, \omega_0)\)。即辛几何无局部不变量------所有 2n 维辛流形局部看起来都一样。

\textbf{证明}（Moser 道路法）：设 \(p \in M\)，取 \(p\) 附近的局部坐标使 \(\omega|_p = \omega_0|_p\)（先通过线性代数找到辛基，再用指数映射拉平到一阶）。构造辛形式的线性同伦 \(\omega_t = (1-t)\omega_0 + t\omega\)，在 \(p\) 的某邻域上 \(\omega_t\) 均为闭的且 \(\omega_t|_p = \omega_0|_p\)。由 Poincare 引理，\(\omega - \omega_0 = d\alpha\)（\(\alpha\) 为 1-形式，可选取使得 \(\alpha|_p = 0\)）。目标是寻找一族微分同胚 \(\phi_t\)（\(\phi_0 = \operatorname{id}\)）满足 \(\phi_t^* \omega_t = \omega_0\)。对 \(t\) 求导得 \(\phi_t^*(\mathcal{L}_{X_t} \omega_t + \dot{\omega}_t) = 0\)，其中 \(X_t\) 是 \(\phi_t\) 的生成向量场。需解 \(\mathcal{L}_{X_t} \omega_t + \dot{\omega}_t = 0\)。利用 Cartan 公式 \(\mathcal{L}_{X_t} \omega_t = d(\iota_{X_t} \omega_t)\) 和 \(\dot{\omega}_t = d\alpha\)，方程化为 \(d(\iota_{X_t} \omega_t + \alpha) = 0\)。由于 \(\omega_t\) 非退化（在 \(p\) 的充分小邻域上），可从 \(\iota_{X_t} \omega_t = -\alpha\) 唯一解出 \(X_t\)。因 \(\alpha|_p = 0\)，\(X_t|_p = 0\)，故 \(X_t\) 的流 \(\phi_t\) 在小邻域上定义至 \(t=1\)。于是 \(\phi_1^* \omega = \phi_1^* \omega_1 = \omega_0\)，即 Darboux 坐标存在。∎

Darboux 定理揭示了辛几何与 Riemann 几何之间最本质的差异：辛流形没有局部曲率不变量，所有 \(2n\) 维辛流形在局部上完全相同。这一事实迫使辛几何将注意力转向全局现象------而其中最基本的全局对象便是 Hamilton 力学系统，它与辛结构通过 Hamilton 向量场的定义紧密相连。

\subsection{242.3 Hamilton 力学与辛几何}\label{hamilton-ux529bux5b66ux4e0eux8f9bux51e0ux4f55}

\textbf{定义 242.4}（Hamilton 向量场）：对函数 \(H \in C^\infty(M)\)，其 \textbf{Hamilton 向量场} \(X_H\) 由 \(\iota_{X_H} \omega = dH\) 唯一确定。

Hamilton 向量场将光滑函数与流形上的动力学建立了一一对应，而这一对应自然地引出了光滑函数空间上的二元运算------Poisson 括号------它刻画了不同物理量在 Hamilton 演化下的相互影响。

\textbf{定义 242.5}（Poisson 括号）：\(C^\infty(M)\) 上的 \textbf{Poisson 括号}为 \(\{F, G\} = \omega(X_F, X_G) = dF(X_G)\)。Poisson 括号满足 Jacobi 恒等式，使 \((C^\infty(M), \{\cdot, \cdot\})\) 成为 Poisson 代数。

Poisson 括号的反对称性直接蕴含了一个根本的几何事实：Hamilton 流保持辛体积形式不变。这正是经典力学中 Liouville 定理------相空间体积在演化中守恒------的精确数学表述。

\textbf{推论 242.2}（Liouville 定理 / 相空间体积守恒）：Hamilton 流 \(\phi_H^t\) 保持辛体积形式 \(\omega^n/n!\)：\((\phi_H^t)^* \omega^n = \omega^n\)。特别是，Hamilton 相流不可压缩。

\textbf{定理 242.3}（Noether 定理的辛表述）：若 \(G\) 通过辛同胚作用在 \((M, \omega)\) 上，且作用有矩映射 \(\mu: M \to \mathfrak{g}^*\)（\(\langle \mu, \xi \rangle\) 的 Hamilton 向量场生成 \(G\) 作用），则 \(\mu\) 沿 Hamilton 流守恒：\(\mu \circ \phi_H^t = \mu\)（若 \(\{H, \mu\} = 0\)）。

\textbf{证明}：设 \(\mu_\xi = \langle \mu, \xi \rangle\) 对 \(\xi \in \mathfrak{g}\)。由矩映射定义，\(d\mu_\xi = \iota_{X_\xi}\omega\)。沿 Hamilton 流 \(\phi_H^t\) 求导：\(\frac{d}{dt} \mu_\xi(\phi_H^t) = X_H(\mu_\xi) = d\mu_\xi(X_H) = \omega(X_\xi, X_H) = \{H, \mu_\xi\}\)。若 \(\{H, \mu_\xi\} = 0\) 对所有 \(\xi\) 成立，则 \(\mu_\xi\) 守恒，从而 \(\mu\) 守恒。若 \(G\) 作用与 Hamilton 流交换（即 \(\phi_H^t \circ g = g \circ \phi_H^t\)），由矩映射的等变性，\(\{H, \mu_\xi\} = 0\) 自动成立。此为 Noether 定理的辛几何版本。\(\blacksquare\)

定理 242.3 的证明表明，矩映射的守恒性与 Poisson 可交换性是同一几何事实的两种表述。Liouville 定理和 Noether 定理共同构成了 Hamilton 力学的两块基石，它们之间的关联------辛体积守恒与对称性守恒的统一------正是 Marsden-Weinstein 辛约化理论的出发点，这将在第 165 章中详细展开。

\hrulefill

\hrulefill

\hrulefill

\hrulefill

\hrulefill

\hrulefill

\section{第228章：辛约化与矩映射}\label{ux7b2c228ux7ae0ux8f9bux7ea6ux5316ux4e0eux77e9ux6620ux5c04}

辛约化（Marsden-Weinstein 约化，1974）是构造新辛流形的最基本方法，其数学本质是经典力学中消除对称性的抽象化------如果 Hamilton 系统具有连续对称群 \(G\)（由矩映射 \(\mu: M \to \mathfrak{g}^*\) 的保存体现为 Noether 定理的守恒量），那么将相空间沿群轨道''约化''后得到的新空间 \(\mu^{-1}(0)/G\) 仍是一个辛流形，但自由度减少了两倍群维数。矩映射的几何意义远不限于力学：Atiyah-Guillemin-Sternberg 凸性定理（1982）证明紧连通 Abel 群作用的矩映射的像是 \(\mathfrak{t}^*\) 中的凸多面体，而 Kirwan 的后续工作将此推广到非 Abel 群并与 Mumford 的几何不变量理论（GIT）建立了惊人的对应------GIT 商对应于 \(0\) 处的辛约化，而一般水平 \(a\) 处的商对应于稳定点的概念。这一联系将代数几何（模空间理论）与辛几何深度融为一体，并通过 Kempf-Ness 定理在复约化和 Kaehler 商之间架设桥梁。本章在后继的第166章（伪全纯曲线）中扮演关键角色：Gromov-Witten 不变量的构造需要约化辛流形上的模空间紧化，而镜对称的许多预测都是通过约化辛流形的量子化表述的。

\subsection{243.1 矩映射与 Hamilton 群作用}\label{ux77e9ux6620ux5c04ux4e0e-hamilton-ux7fa4ux4f5cux7528}

矩映射（moment map）是辛几何中对称性与守恒量之间联系的精确数学表述，是 Noether 定理在辛几何框架下的几何化。给定李群 \(G\) 在辛流形 \((M, \omega)\) 上的辛作用，若每个无穷小生成元 \(X_\xi\)（\(\xi \in \mathfrak{g}\)）都是 Hamilton 向量场，则存在矩映射 \(\mu: M \to \mathfrak{g}^*\) 满足 \(d\langle \mu, \xi \rangle = \iota_{X_\xi} \omega\)。这一条件的几何意义是：\(\langle \mu, \xi \rangle\) 是 \(X_\xi\) 的 Hamilton 函数，而 \(\mu\) 将这些 Hamilton 函数''打包''为一个 \(\mathfrak{g}^*\) 值函数。矩映射的等变性（equivariance）\(\mu(g \cdot p) = \operatorname{Ad}_g^* \mu(p)\) 是 Hamilton 作用可积性条件的体现------它保证了 Poisson 括号 \(\{\langle \mu, \xi \rangle, \langle \mu, \eta \rangle\} = \langle \mu, [\xi, \eta] \rangle\)，即 \(\mu\) 是 Lie 代数同态 \(C^\infty(M) \leftarrow \mathfrak{g}\) 的''提升''。在物理学中，矩映射直接对应了守恒量：若 \(H\) 在 \(G\) 作用下不变，则 \(\mu\) 沿 Hamilton 流守恒------这正是经典 Noether 定理（对称性 \(\Rightarrow\) 守恒律）的辛几何表述。矩映射的典型例子包括：角动量映射（\(SO(3)\) 作用于 \(\mathbb{R}^6\) 的相空间）、线性动量映射（\(\mathbb{R}^n\) 的平移对称性）、以及环面作用的矩映射（其像为凸多胞体，见 Atiyah-Guillemin-Sternberg 凸性定理）。矩映射的零纤维 \(\mu^{-1}(0)\) 是辛约化（Marsden-Weinstein 约化）的起点，而一般纤维 \(\mu^{-1}(a)\) 的商空间 \(M /\!/ G(a)\) 构成了构造新辛流形的基本方法。

\textbf{定义 243.1}（Hamilton 群作用 / 辛作用）：李群 \(G\) 在辛流形 \((M, \omega)\) 上的辛作用称为 \textbf{Hamilton 作用}，如果每个 \(\xi \in \mathfrak{g}\) 对应的无穷小生成元 \(X_\xi\) 是某个函数 \(H_\xi\) 的 Hamilton 向量场。

\textbf{定义 243.2}（矩映射 / Moment Map）：\(G\)-等变映射 \(\mu: M \to \mathfrak{g}^*\) 称为 \textbf{矩映射}（或动量映射），如果对所有 \(\xi \in \mathfrak{g}\)，\(\langle \mu, \xi \rangle\) 是 \(X_\xi\) 的 Hamilton 函数：

\[d\langle \mu, \xi \rangle = \iota_{X_\xi} \omega\]

\subsection{243.2 Marsden-Weinstein 约化}\label{marsden-weinstein-ux7ea6ux5316}

\textbf{定理 243.1}（Marsden-Weinstein-Meyer 辛约化定理，1974）：设 \(G\) 以 Hamilton 作用在 \((M, \omega)\) 上，\(\mu: M \to \mathfrak{g}^*\) 是矩映射。若 \(a \in \mathfrak{g}^*\) 是 \(\mu\) 的正则值且 \(G_a\)（\(a\) 的迷向子群）自由作用于 \(\mu^{-1}(a)\)，则商空间

\[M /\!/ G(a) = \mu^{-1}(a) / G_a\]

具有唯一的辛结构 \(\omega_{\text{red}}\)，满足 \(\pi^* \omega_{\text{red}} = \iota^* \omega\)（\(\iota: \mu^{-1}(a) \hookrightarrow M\)，\(\pi: \mu^{-1}(a) \to M /\!/ G(a)\) 是投影）。

\textbf{证明}：由 \(\mu\) 的正则性和 \(G_a\) 的自由作用，\(\mu^{-1}(a)\) 为 \(M\) 的子流形。对 \(p \in \mu^{-1}(a)\)，切空间 \(T_p \mu^{-1}(a) = \ker d\mu_p\)。核满足 \(T_p(G_a\cdot p) = (T_p \mu^{-1}(a))^\omega\)（辛正交补）。\(\iota^*\omega\) 在纤维方向退化，核恰为 \(G_a\) 轨道，故 descends 到商空间 \(M/\!/G(a)\) 为非退化 2-形式。闭性由 \(d\omega=0\) 传承。\(\blacksquare\)

\textbf{定理 243.2}（约化的维数）：\(\dim(M /\!/ G(a)) = \dim M - 2 \dim G_a\)。特别地，若 \(G\) 以有效自由作用在 \(\mu^{-1}(a)\) 上，则 \(\dim(M /\!/ G(a)) = \dim M - 2 \dim G\)。

\textbf{证明}：由 Marsden-Weinstein 约化定义，\(M /\!/ G(a) = \mu^{-1}(a)/G_a\)，其中 \(G_a = \{g \in G : \operatorname{Ad}_g^* a = a\}\) 是 \(a\) 的迷向子群。由矩映射的等变性，\(G_a\) 保持 \(\mu^{-1}(a)\) 不变。\(\mu^{-1}(a)\) 是 \(M\) 的子流形，其维数为 \(\dim M - \dim G\)（因为 \(a\) 是 \(\mu\) 的正则值，而 \(\mu\) 的值域为 \(\mathfrak{g}^*\)，维数为 \(\dim G\)）。\(G_a\) 在 \(\mu^{-1}(a)\) 上的作用是自由的（当 \(a\) 为正则值时），故商空间的维数为 \(\dim \mu^{-1}(a) - \dim G_a = \dim M - \dim G - \dim G_a\)。由于 \(G/G_a \cong \mathcal{O}_a\)（余伴随轨道），\(\dim G = \dim G_a + \dim \mathcal{O}_a\)。对于正则值 \(a\)，\(\dim \mathcal{O}_a = \dim G_a\)（\(G_a\) 为极大环面时成立），故 \(\dim(M /\!/ G(a)) = \dim M - 2 \dim G_a\)。特别地，若 \(G\) 自由作用，则 \(G_a = G\) 且 \(\dim = \dim M - 2\dim G\)。\(\blacksquare\)

\subsection{243.3 Delzant 定理与环面簇}\label{delzant-ux5b9aux7406ux4e0eux73afux9762ux7c07}

\textbf{定义 243.3}（环面辛流形 / Toric Symplectic Manifold）：紧致连通辛流形 \((M^{2n}, \omega)\) 是\textbf{环面的}，如果它容许 \(\mathbb{T}^n\) 的有效 Hamilton 作用（最大可能的环面维数）。

\textbf{定理 243.3}（Delzant 定理，1988）：紧致环面辛流形与 Delzant 多面体（有理多面体，满足在每个顶点处边方向形成 \(\mathbb{Z}\)-基）之间存在一一对应。矩映射的像 \(\mu(M)\) 恰好是 Delzant 多面体。

\textbf{证明}：由 \(T^n\) 的 Hamilton 作用，矩映射 \(\mu: M \to \mathbb{R}^n \cong \mathfrak{t}^*\) 的像为凸多胞体（Atiyah 凸性定理），其顶点对应 \(T^n\) 的 fixed points。Delzant 条件确保每个顶点邻域辛同构于标准环面作用的 \(\mathbb{C}^n\)。反过来，从 Delzant 多面体 P 通过辛商构造可得环面辛流形：取 \(\mathbb{C}^d\) 上一个由多面体不等式定义的 T\^{}d 作用的辛约化即得。\(\blacksquare\)

\hrulefill

\hrulefill

\hrulefill

\hrulefill

\hrulefill

\hrulefill

\section{第229章：伪全纯曲线与 Gromov-Witten 理论}\label{ux7b2c229ux7ae0ux4f2aux5168ux7eafux66f2ux7ebfux4e0e-gromov-witten-ux7406ux8bba}

伪全纯曲线是由 Mikhail Gromov 在 1985 年的里程碑论文《Pseudo-holomorphic curves in symplectic manifolds》中引入的，其革命性在于用 Cauchy-Riemann 方程的椭圆正则性来探测辛流形的全局结构------这在 Darboux 定理所暗示的''局部无趣''背景下尤其引人注目。核心技巧是：在辛流形上选取与 \(\omega\) 相容的近复结构 \(J\)，然后研究 \(J\)-全纯映射 \(u: (\Sigma, j) \to (M, J)\)（即满足 \(\bar{\partial}_J u = 0\) 的映射），其模空间 \(\mathcal{M}_{g,k}(M, A)\) 在适当的紧化下成为一个良好的（实）orbifold，从而通过积分拉回其上同调类定义 Gromov-Witten 不变量。这一理论与第 165 章的辛约化及矩映射通过 Kirwan 映射和量子上同调产生深刻联系，并在 Kontsevich 的镜对称猜想和 Givental 等人在 toric 流形上 Gromov-Witten 不变量的计算中取得了巨大成功------后者直接证明了 Calabi-Yau 五维三次超曲面的镜对称预测（Candelas 等人的公式），成为 20 世纪数学物理最辉煌的成果之一。

\subsection{244.1 近复结构与伪全纯曲线}\label{ux8fd1ux590dux7ed3ux6784ux4e0eux4f2aux5168ux7eafux66f2ux7ebf}

\textbf{定义 244.1}（与辛形式相容的近复结构）：流形 \(M\) 上的\textbf{近复结构} \(J\) 满足 \(J^2 = -\operatorname{id}\)。\(J\) 与辛形式 \(\omega\) \textbf{相容}（compatible），如果 \(g_J(v, w) = \omega(v, Jw)\) 是 Riemann 度量。\footnote{注：\textbf{tame} 是更弱的条件------仅要求 \(\omega(v, Jv) > 0\) 对所有 \(v \neq 0\) 成立，而不要求 \(g_J\) 是对称的。compatible 等价于 tame + \(J\) 是 \(\omega\)-不变的（即 \(\omega(Jv, Jw) = \omega(v, w)\)）。}

\textbf{定理 244.1}（相容近复结构的存在性）：辛流形 \((M, \omega)\) 上相容近复结构的空间 \(\mathcal{J}(M, \omega)\) 是非空且可缩的。

\textbf{证明}：选取 \(M\) 上的 Riemann 度量 \(g_0\)。定义 \((1,1)\)-张量 \(A\) 满足 \(g_0(Au, v) = \omega(u, v)\)。\(A\) 为反对称。令 \(P = \sqrt{-A^2}\)（正定对称），定义 \(J = P^{-1}A\)。验证 \(J^2 = -\operatorname{id}\) 且 \(\omega(u, Jv) = g_0(Pu, v)\) 为正定对称双线性型，故 \(J\) 为相容近复结构。可缩性：\(\mathcal{J}(M,\omega)\) 是 \(\operatorname{Sp}(2n,\mathbb{R})/\operatorname{U}(n)\) 的截面空间，后者的纤维可缩，故空间可缩。\(\blacksquare\)

\textbf{定义 244.2}（\(J\)-全纯曲线）：映射 \(u: (\Sigma, j) \to (M, J)\)（\(\Sigma\) 是 Riemann 面）是 \textbf{\(J\)-全纯的}，如果

\[\bar{\partial}_J u = \frac{1}{2}(du + J \circ du \circ j) = 0\]

即 \((u, J)\)-全纯曲线的 Cauchy-Riemann 方程。

\textbf{定理 244.2}（能量恒等式）：\(J\)-全纯曲线 \(u\) 的能量 \(E(u) = \frac{1}{2} \int_\Sigma |du|^2 d\operatorname{vol}_\Sigma = \int_\Sigma u^* \omega\)。因此 \(J\)-全纯曲线的面积等于辛面积（拓扑量）。

\textbf{证明}：在 \(\Sigma\) 上取局部共形坐标 \(z = s+it\)，\(j(\partial_s) = \partial_t\)。\(J\)-全纯条件 \(\bar{\partial}_J u = 0\) 即 \(\partial_s u + J \partial_t u = 0\)。则 \(|du|^2 = |\partial_s u|^2 + |\partial_t u|^2\)。利用相容性 \(g_J(\cdot,\cdot) = \omega(\cdot, J\cdot)\)，计算 \(|du|^2 = |\partial_s u|^2 + |J\partial_s u|^2 = 2|\partial_s u|^2\)。同时 \(u^*\omega(\partial_s, \partial_t) = \omega(\partial_s u, \partial_t u) = \omega(\partial_s u, -J\partial_s u) = g_J(\partial_s u, \partial_s u) = |\partial_s u|^2\)。故 \(\frac{1}{2}|du|^2 d\operatorname{vol} = |\partial_s u|^2 ds \wedge dt = u^*\omega\)。积分即得 \(E(u) = \int_\Sigma u^*\omega\)。\(\blacksquare\)

\subsection{244.2 Gromov 紧致性与非压缩定理}\label{gromov-ux7d27ux81f4ux6027ux4e0eux975eux538bux7f29ux5b9aux7406}

\textbf{定理 244.3}（Gromov 紧致性定理，1985）：具有一致有界能量的 \(J\)-全纯曲线的序列可以「气泡化」收敛到节点 \(J\)-全纯曲线（nodal \(J\)-holomorphic curve），除有限个「气泡」点外在 \(\Sigma\) 上光滑收敛。

\textbf{证明}：设 \(\{u_n\}\) 为 \(J\)-全纯曲线序列，能量 \(E(u_n) \leq C\)。由能量恒等式，\(u_n\) 的面积有界。在 \(\Sigma\) 上，若 \(u_n\) 的梯度在 \(p\) 点附近无界，由 \(\varepsilon\)-正则性，能量在 \(p\) 处集中（气泡）。重新缩放 \(p\) 附近得非平凡 \(J\)-全纯球面 \(S^2 \to M\)（气泡）。移除气泡后余下区域光滑收敛到极限 \(u_\infty\)。气泡形成在有限个孤立点（每次气泡消耗最少能量 \(\hbar = \inf\{E(v) : v \text{ 非平凡 } J\text{-全纯球面}\}\)）。极限为节点曲线，节点对应气泡附着点。\(\blacksquare\)

Gromov 紧致性定理是伪全纯曲线理论的基石：它保证了在能量有界条件下，模空间在经过适当的紧化后是紧致的，从而可以通过积分给出良定义的数值不变量。这一紧致性结果直接引向了辛几何中最著名的刚性现象------Gromov 非压缩定理，它以一种令人震惊的方式揭示了辛嵌入与体积嵌入之间的本质区别。

\textbf{定理 244.4}（Gromov 非压缩定理 / Non-Squeezing Theorem）：若存在辛嵌入 \(B^{2n}(R) \hookrightarrow B^2(r) \times \mathbb{R}^{2n-2}\)（柱面），则 \(R \leq r\)。即一个大的辛球不能被「挤入」一个细的辛柱面中------这是辛刚性的里程碑。

\textbf{证明}：使用伪全纯曲线。在 \(\mathbb{C}P^1\) 的辛嵌入中，\(J\)-全纯球的不变量限制了 \(R\) 和 \(r\) 的关系。能量 \(\geq \pi R^2\) 但面积受限于 \(\pi r^2\)。∎

\subsection{244.3 Gromov-Witten 不变量}\label{gromov-witten-ux4e0dux53d8ux91cf}

\textbf{定义 244.3}（Gromov-Witten 不变量）：辛流形 \((M, \omega)\) 的 \textbf{Gromov-Witten 不变量}是映射

\[\operatorname{GW}_{g, k, A}: H^*(M)^{\otimes k} \to \mathbb{Q}\]

计数（在适当意义下）亏格 \(g\)、\(k\) 个标记点且同调类为 \(A \in H_2(M; \mathbb{Z})\) 的 \(J\)-全纯曲线的数目。形式上，

\[\operatorname{GW}_{g, k, A}(\alpha_1, \ldots, \alpha_k) = \int_{[\overline{\mathcal{M}}_{g,k}(M, A)]^{\text{vir}}} \prod_{i=1}^k \operatorname{ev}_i^* \alpha_i\]

其中 \(\overline{\mathcal{M}}_{g,k}(M, A)\) 是稳定映射的模空间（紧化），\([\cdot]^{\text{vir}}\) 是虚基本类。

\textbf{定理 244.5}（GW 不变量与计数几何）：对合适的 \(\alpha_i\)（如 Poincaré 对偶于点类的上同调类），GW 不变量计算特定几何约束下的 \(J\)-全纯曲线数。这些数字是辛不变量（不依赖于 \(J\) 的选择）。

\textbf{证明}：Gromov-Witten 不变量 \(\langle\tau_{a_1}(\alpha_1),\ldots,\tau_{a_k}(\alpha_k)\rangle_{g,\beta}^X\) 定义为稳定映射模空间 \(\overline{\mathcal{M}}_{g,k}(X,\beta)\) 上评估类与 \(\psi\)-类（余切线丛的 Chern 类）的乘积在虚基本类上的积分。当 \(\alpha_i\) 取为点类的 Poincaré 对偶且 \(a_i=0\) 时，不变量精确计数通过 \(k\) 个固定点的亏格 \(g\)、同调类 \(\beta\) 的 \(J\)-全纯曲线数目。这些数字不依赖于近复结构 \(J\) 的选取（在辛同痕下不变），是辛流形的辛不变量。该性质由模空间的紧致性（Gromov 紧致性）和 Gromov-Witten 类在变形下的不变性保证。\(\blacksquare\)

\subsection{244.4 量子上同调与 WDVV 方程}\label{ux91cfux5b50ux4e0aux540cux8c03ux4e0e-wdvv-ux65b9ux7a0b}

\textbf{定义 244.4}（量子上同调环 / Quantum Cohomology）：\(M\) 的\textbf{量子上同调} \(QH^*(M)\) 是 \(H^*(M; \mathbb{C})\) 作为向量空间，但乘法量子变形：

\[\alpha * \beta = \sum_{A \in H_2(M)} (\alpha * \beta)_A \cdot q^A\]

其中 \((\alpha * \beta)_A\) 由三个点的 GW 不变量给出：\(\langle \alpha, \beta, \gamma \rangle_A = \operatorname{GW}_{0,3,A}(\alpha, \beta, \gamma)\)。

\textbf{定理 244.6}（WDVV 方程 / Witten-Dijkgraaf-Verlinde-Verlinde 方程）：量子上同调的关联函数满足 WDVV 方程，这是 GW 势的完全可积系统的条件。WDVV 方程等价于量子上同调的结合性。

\textbf{证明}：GW 势 \(\Phi(t) = \sum_{k,A} \frac{1}{k!} \langle t, \ldots, t \rangle_{0,k,A}\) 的三阶导数 \(c_{ijk} = \partial_i \partial_j \partial_k \Phi\) 定义量子乘法。结合律 \((t_i * t_j) * t_k = t_i * (t_j * t_k)\) 在 \(\partial_i \partial_j \partial_k \partial_\ell \Phi\) 的层面给出：\(\sum_{a,b} \partial_i \partial_j \partial_a \Phi \cdot g^{ab} \cdot \partial_b \partial_k \partial_\ell \Phi = \sum_{a,b} \partial_i \partial_\ell \partial_a \Phi \cdot g^{ab} \cdot \partial_b \partial_j \partial_k \Phi\)（对所有 \(i,j,k,\ell\)）。其中 \(g_{ab} = \partial_a \partial_b \partial_0 \Phi\) 为拓扑度量。此即 WDVV 方程，表明 GW 势是 Frobenius 流形结构的势函数。\(\blacksquare\)

\hrulefill

\hrulefill

\hrulefill

\hrulefill

\hrulefill

\hrulefill

\section{第230章：切触几何}\label{ux7b2c230ux7ae0ux5207ux89e6ux51e0ux4f55}

切触几何是奇数维流形上的几何结构，与辛几何通过''辛化''过程构成自然的配对关系。其物理根源可追溯到 19 世纪：Carnot 在热力学和 Huygens 在几何光学中的工作早已暗示了切触结构的必要性，但直到 Sophus Lie 创立变换群理论时，才将切触变换作为一阶偏微分方程解法几何化的核心工具。从现代观点看，\(2n+1\) 维切触流形 \((M, \xi)\) 是整个辛几何体系不可分割的组成部分------其辛化流形 \((M \times \mathbb{R}, d(e^t \alpha))\) 携带了自然的 \(\mathbb{R}\)-等变辛结构，而切触结构的 Legendre 子流形（最大各向同性子流形，维数恰为 \(n\)）则在辛化后的 Lagrange 子流形层面上与 Floer 理论相呼应。在低维拓扑中，三维流形上的切触结构（由 Bennequin、Eliashberg 和 Giroux 系统分类）是理解 tight 与 overtwisted 结构、Weinstein 猜想以及 Heegaard Floer 同调的重要途径；在物理中，切触几何是几何热力学、非完整约束力学系统和超弦理论中 \(S^1\)-纤维化描述的自然语言。

\subsection{245.1 切触结构}\label{ux5207ux89e6ux7ed3ux6784}

\textbf{定义 245.1}（切触流形）：奇数维流形 \(M^{2n+1}\) 上的\textbf{切触结构} \(\xi\) 是余维 1 的极大非可积超平面场。局部地，\(\xi = \ker \alpha\)，其中 \textbf{切触形式} \(\alpha \in \Omega^1(M)\) 满足 \(\alpha \wedge (d\alpha)^n \neq 0\) 处处。

\textbf{定义 245.2}（Legendre 子流形）：切触流形 \((M^{2n+1}, \xi)\) 的最大维积分子流形（维数为 \(n\)）称为 \textbf{Legendre 子流形}。它们是辛几何中 Lagrangian 子流形的切触类比。

\subsection{245.2 切触同调}\label{ux5207ux89e6ux540cux8c03}

\textbf{定义 245.3}（Reeb 向量场）：切触形式 \(\alpha\) 的 \textbf{Reeb 向量场} \(R_\alpha\) 由 \(\iota_{R_\alpha} d\alpha = 0\) 和 \(\alpha(R_\alpha) = 1\) 唯一确定。Reeb 向量场是 Hamilton 向量场的切触类比。

Reeb 向量场在切触几何中扮演着与 Hamilton 向量场在辛几何中类似的角色------它定义了切触流形上的''动力学''。一旦掌握了 Reeb 向量场的概念，便可以类比辛几何中 Hamiltonian Floer 同调的构造，在切触流形上对 Reeb 轨道的模空间进行计数，从而定义切触同调这一核心不变量。

\textbf{定义 245.4}（切触同调 / Contact Homology，Eliashberg-Givental-Hofer 2000）：切触流形 \((M, \xi)\) 的\textbf{切触同调} \(HC_*(M, \xi)\) 由 Reeb 轨道的链复形生成，边界算子计数伪全纯曲线。切触同调是切触结构的不变量。

切触同调的定义将切触几何纳入了 Floer 型同调理论的广阔框架。在这个框架下，切触结构的刚性与柔性不再仅仅是直觉概念，而是可以被同调群精确量化的几何不变量。这一背景下，Eliashberg 的分类定理揭示了切触几何中刚性（tight）与柔性（overtwisted）之间的根本对偶。

\textbf{定理 245.1}（切触几何的刚性 / Eliashberg 1989）：在维数 \(> 3\) 的闭流形上，所有过度扭转（overtwisted）切触结构是同痕的，但紧切触（tight）结构形成更丰富的分类。即切触几何将流形的切触结构分为两类：柔性的（过度扭转）和刚性的（紧的）。

\textbf{证明}：对过度扭转切触结构，Eliashberg 证明 \(h\)-原理成立：任意过度扭转切触结构之间存在切触同痕。核心构造使用切触结构的柔性 \(h\)-原理：过度扭转圆盘的存在性允许将切触平面沿任意同伦类变形。紧切接触结构（tight）则相反------它们在 \(h\)-原理下表现出刚性，其分类由基于全纯曲线计数的不变量（如切触同调、嵌入切触同调 ECH）决定。例如 \(S^3\) 上仅有一个紧切接触结构（标准结构），而 \(T^3\) 上有可数无穷多个互不同痕的紧切接触结构。\(\blacksquare\)

\subsection{245.3 三维切接触拓扑}\label{ux4e09ux7ef4ux5207ux63a5ux89e6ux62d3ux6251}

\textbf{定义 245.5}（三维切触拓扑的特殊性）：维数 3 中切触结构 \(\xi\) 的平面场可以用特征叶理（characteristic foliation）和开卷分解（open book decomposition）研究。Giroux 定理建立了开卷分解与切触结构之间的对应。

\textbf{定理 245.2}（Giroux 定理，2000）：三维闭流形上的切触结构（模同痕）与开卷分解（模正稳定化）一一对应。这是理解和分类三维切触结构的基本工具。

\textbf{证明}：三维闭流形上切触结构 \(\xi\) 的开卷分解：取 \(M\) 中亏格 \(g\) 曲面 \(\Sigma\) 沿其边界曲线 Dehn 扭转，构造流形使得 \(\xi\) 由 \(\Sigma\) 上的一族曲线（特征叶状结构）沿横向 \(S^1\) 方向编码。Giroux 证明该对应是双射（模正稳定化）：每个切触结构唯一对应开卷分解，且正稳定化（增加亏格但不改变切触同痕类）等价于切触结构级别的平凡化操作。该对应实现了切触结构的组合分类，是 Giroux 程序（高维切触结构的开卷分解理论）的基础。\(\blacksquare\)

\subsection{切触几何与低维拓扑的深层联系}\label{ux5207ux89e6ux51e0ux4f55ux4e0eux4f4eux7ef4ux62d3ux6251ux7684ux6df1ux5c42ux8054ux7cfb}

切触几何在 20 世纪末和 21 世纪初成为辛几何和低维拓扑的核心交叉领域。\textbf{切触结构}（contact structure）的拓扑分类是三维流形理论中最活跃的研究方向之一。Eliashberg（1989）证明了过度扭转切触结构满足 \(h\)-原理（Gromov 1986）------这意味着它们在给定同伦类中是唯一的，而紧切触结构的分类则极其丰富。\textbf{Giroux 对应}（2000）建立了三维流形上切触结构与开卷分解之间的一一对应，这是切触拓扑中最深刻的结构定理。\textbf{Ozsvath-Szabo 的 Heegaard Floer 同调}（2001）和 \textbf{Kronheimer-Mrowka 的 Seiberg-Witten Floer 同调}（2007）为切触结构提供了强大的同调不变量，可用于区分紧切触结构和过度扭转切触结构。\textbf{嵌入切触同调}（Embedded Contact Homology, ECH）由 Hutchings（2002）引入，是三维切触流形的 Reeb 轨道的 Floer 同调，与 Seiberg-Witten Floer 同调同构（Taubes 2010）------这一深刻结果将切触动力系统与规范理论联系起来。\textbf{Weinstein 猜想}（Weinstein 1978）断言闭切触流形上的 Reeb 向量场总有闭轨道，在三维情形由 Taubes（2007）通过 ECH = SWF 等价性证明。\textbf{Legendrian 纽结理论}是切触几何在纽结理论中的推广：Legendrian 纽结是切触三维流形 \(S^3\)（带标准切触结构）中处处切于切触平面场的一维子流形，其 Legendrian 不变量（如 Thurston-Bennequin 数、旋转数、Chekanov-Eliashberg 微分分次代数）为经典纽结理论提供了更精细的分类工具。

\hrulefill

\hrulefill

\section{第231章：Fukaya 范畴与镜对称}\label{ux7b2c231ux7ae0fukaya-ux8303ux7574ux4e0eux955cux5bf9ux79f0}

Fukaya 范畴是辛流形中 Lagrangian 子流形的 \(A_\infty\)-范畴，是 Kontsevich 的同调镜对称猜想中的核心对象。

\subsection{246.1 Lagrangian 子流形与 Floer 同调}\label{lagrangian-ux5b50ux6d41ux5f62ux4e0e-floer-ux540cux8c03}

\textbf{定义 246.1}（Lagrangian 子流形）：辛流形 \((M^{2n}, \omega)\) 的 \(n\) 维子流形 \(L\) 是 \textbf{Lagrangian} 的，如果 \(\omega|_L = 0\)。Lagrangian 子流形是辛几何中最基本的对象。

Lagrangian 子流形之所以占据辛几何的核心地位，是因为它们既是辛结构''退化''的极限情形（辛形式限制为零），又恰好具有辛流形一半的维数------这一维数特征使得它们在镜对称中扮演着连接辛几何与复几何的桥梁角色。Lagrangian 子流形之间有交点时，可以仿照 Hamiltonian Floer 同调的构造，通过计数连接交点的伪全纯带定义 Lagrangian Floer 同调。

\textbf{定义 246.2}（Lagrangian Floer 同调）：对两个横截相交的 Lagrangian 子流形 \(L_0, L_1\)，\textbf{Lagrangian Floer 同调} \(HF(L_0, L_1)\) 由交点生成，边界算子计数连接交点的伪全纯带（strip）。\(HF(L_0, L_1)\) 在 Hamiltonian 同痕下不变。

\textbf{定理 246.1}（Floer 的 Lagrangian 交猜想 / Arnold 猜想）：在适当条件下，\(\operatorname{rank} HF(L_0, L_1) \leq \sum_k \dim H^k(L_0)\)，且若 \(L_0 \pitchfork L_1\)，则 \(|L_0 \cap L_1| \geq \sum_k \dim H^k(L_0)\)。

\textbf{证明}：Lagrangian Floer 同调 \(HF(L_0,L_1)\) 的生成元为 \(L_0\cap L_1\) 的交点，微分计数连接交点的伪全纯带（strip）（边缘分别落在 \(L_0\) 和 \(L_1\) 上）。Arnold 猜想：\(|L_0\cap L_1|\ge\sum\dim H^k(L_0;\mathbb{Z}/2)\)（\(L_0\) 模 Hamiltonian 同痕后与 \(L_1\) 横截）。Floer 证明 \(\dim HF(L_0,L_1)\le\sum\dim H^k(L_0)\)，且存在谱序列从 \(HF(L_0,L_1)\) 收敛至 \(H^*(L_0)\)。Arnold 猜想的证明（由 Fukaya-Ono、Liu-Tian 等完成）需要处理模空间的紧化和横截性（通过 Kuranishi 结构或虚拟基本链技术）。\(\blacksquare\)

\subsection{\texorpdfstring{246.2 Fukaya \(A_\infty\)-范畴}{246.2 Fukaya A\_\textbackslash infty-范畴}}\label{fukaya-a_infty-ux8303ux7574}

\textbf{定义 246.3}（\(A_\infty\)-范畴）：一个 \textbf{\(A_\infty\)-范畴} \(\mathcal{C}\) 由以下数据组成： - 一组对象 \(\operatorname{Ob}(\mathcal{C})\) - 对每对对象 \(X, Y\)，一个分次向量空间 \(\operatorname{Hom}_{\mathcal{C}}(X, Y)\) - 对每个 \(k \geq 1\)，\(k\) 阶乘法运算 \(\mu^k: \operatorname{Hom}(X_{k-1}, X_k) \otimes \cdots \otimes \operatorname{Hom}(X_0, X_1) \to \operatorname{Hom}(X_0, X_k)\)，次数为 \(\deg \mu^k = 2 - k\)

这些运算满足 \(A_\infty\) 关系：对所有 \(k \geq 1\)， \[\sum_{r+s+t = k} (-1)^{r+st} \mu^{r+1+t}(\mathrm{id}^{\otimes r} \otimes \mu^s \otimes \mathrm{id}^{\otimes t}) = 0\]

其中符号规则中 \((-1)^{r+st}\) 的指数由 Koszul 符号约定确定。特别地： - \(\mu^1 \circ \mu^1 = 0\)（\(\mu^1\) 是微分） - \(\mu^1(\mu^2(a, b)) = \mu^2(\mu^1(a), b) + (-1)^{|a|} \mu^2(a, \mu^1(b))\)（Leibniz 法则） - \(\mu^2\) 的结合律在同伦意义下成立：\(\mu^2(\mu^2(a, b), c) - \mu^2(a, \mu^2(b, c)) = \pm \mu^1(\mu^3(a, b, c)) \pm \mu^3(\mu^1(a), b, c) \pm \mu^3(a, \mu^1(b), c) \pm \mu^3(a, b, \mu^1(c))\)

\textbf{定义 246.4}（Fukaya 范畴 / Fukaya 1993）：辛流形 \((M, \omega)\) 的 \textbf{Fukaya 范畴} \(\mathcal{F}uk(M)\) 是一个 \(A_\infty\)-范畴，其中： - 对象：合适的 Lagrangian 子流形（带额外的结构：自旋结构、分次、局部系统） - 态射空间：\(\operatorname{Hom}(L_0, L_1) = CF^*(L_0, L_1)\)（Floer 链复形），其生成元为 \(L_0\) 与 \(L_1\) 的横截交点，分次由 Maslov 指标给出 - \(A_\infty\) 运算 \(\mu^k\)：\(\mu^1\) 是 Floer 微分（计数全纯带），\(\mu^k\)（\(k \geq 2\)）由 \(k+1\) 个边界标记点的全纯 \(k\)-边形的计数给出，即全纯映射 \(u: D^2 \setminus \{z_0, \ldots, z_k\} \to M\) 的模空间在 \(\mathbb{R}\) 平移下的商

\textbf{定理 246.2}（\(A_\infty\)-关系的几何来源）：Floer 微分的平方 \(\mu^1 \circ \mu^1 = 0\) 由模空间的边界给出（气泡化和断裂）。\(\mu^2\) 结合律的偏差 \(\mu^2(\mu^2(a, b), c) - \mu^2(a, \mu^2(b, c))\) 由 \(\mu^3\) 对边界算子的同伦给出。所有 \(A_\infty\) 关系来自模空间紧化边界的几何。

\textbf{证明}：Fukaya 范畴 \(\mathcal{F}uk(X)\) 的对象为 Lagrangian 子流形，态射空间为 Floer 复形 \(CF(L_0,L_1)\)。\(A_\infty\)-结构映射 \(\mu^k: CF(L_{k-1},L_k)\otimes\cdots\otimes CF(L_0,L_1)\to CF(L_0,L_k)\) 由带边界标记点的伪全纯 \((k+1)\)-边形计数定义（模空间 \(\mathcal{M}_{k+1}\)）。\(A_\infty\)-关系 \(\sum_{r,s,t}(-1)^{*}\mu^{r+1+t}(a_k,\ldots,\hat{a}_{s},\ldots,\hat{a}_{s+r},\ldots,a_0)=0\) 来自 \(\mathcal{M}_{k+1}\) 的边界分解：模空间的 Gromov 边界由气泡化和断裂的圆盘构成，恰好对应 \(\mu^i\) 的复合------即所有可能分割 \((r,s,t)\) 的和在代数和意义下为零。这反映了模空间边界的装配代数结构。\(\blacksquare\)

\subsection{246.3 Kontsevich 同调镜对称}\label{kontsevich-ux540cux8c03ux955cux5bf9ux79f0}

\textbf{定义 246.5}（同调镜对称猜想 / Homological Mirror Symmetry，Kontsevich 1994）：对镜面对 \((X, \check{X})\)（\(X\) 是 Calabi-Yau 辛流形，\(\check{X}\) 是其镜面复流形），存在范畴等价：

\[D^b \mathcal{F}uk(X) \cong D^b \operatorname{Coh}(\check{X})\]

即 \(X\) 的 Fukaya 范畴的导出范畴（辛几何）等价于镜面 \(\check{X}\) 的相干层范畴的导出范畴（复代数几何）。

\textbf{定理 246.3}（镜对称的部分证明）： - 椭圆曲线：Polishchuk-Zaslow（1998）证明了 \(T^2\) 的同调镜对称 - 四次曲面（quartic surface）：Seidel（2015）和 Sheridan（2015）证明 - 环面簇的衍生范畴等价：由 Abouzaid、Fang、Liu 等人证明 - 完全交和超曲面的镜对称：多个作者推进

\textbf{证明}：镜对称断言 Calabi-Yau 流形 \(X\) 的 Gromov-Witten 理论（A-模型）等价于镜流形 \(\check{X}\) 的周期积分和变分 Hodge 结构（B-模型）。数学证明进展：(1) Candelas 等（1991）用镜对称预测五次三重形 \(X_5\subset\mathbb{P}^4\) 的有理曲线计数，由 Givental（1996）和 Lian-Liu-Yau（1997）通过局部化和退化技术严格证明；(2) Kontsevich 的同调镜对称猜想（1994）：\(D^b\mathcal{F}uk(X)\cong D^b\operatorname{Coh}(\check{X})\)，由 Seidel 在若干低维情形验证；(3) SYZ 猜想（1996）将镜对称归约为特殊 Lagrangian \(T^n\)-纤维化上的 T-对偶。完整的一般性证明仍为开放问题。\(\blacksquare\)

\subsection{246.4 SYZ 猜想与 Strominger-Yau-Zaslow}\label{syz-ux731cux60f3ux4e0e-strominger-yau-zaslow}

\textbf{定义 246.6}（SYZ 猜想，Strominger-Yau-Zaslow 1996）：接近大体积极限的 Calabi-Yau 流形 \(X\) 是特殊 Lagrangian 环面 \(T^n\) 的纤维化（\(T^n\)-对偶 / T-duality）。镜面 \(\check{X}\) 通过对偶纤维构造。

\textbf{定理 246.4}（SYZ 的 T-对偶图像）：在 SYZ 图像下，\(X\) 上的 Lagrangian 子流形对应 \(\check{X}\) 上的全纯向量丛（通过纤维 \(T^n\) 上的 Fourier-Mukai 变换），\(X\) 上的伪全纯曲线对应 \(\check{X}\) 上相干层的修正复形。

\textbf{证明}：SYZ 假设 \(X\) 存在特殊 Lagrangian \(T^n\)-纤维化 \(f:X\to B\)。镜流形 \(\check{X}\) 由对偶环面纤维化 \(\check{f}:\check{X}\to B\) 构造（每个纤维为对偶环面 \(\check{T}^n=\operatorname{Hom}(\pi_1(T^n),\operatorname{U}(1))\)）。在此图像下，\(X\) 上的 Lagrangian 子流形（截面和多截面对应的多值截面）通过沿纤维的 Fourier-Mukai 变换对应 \(\check{X}\) 上的全纯向量丛；\(X\) 上的伪全纯圆盘（边界落在 Lagrangian 上）对应 \(\check{X}\) 上相干层的 Ext 类。Kontsevich-Soibelman 和 Gross-Siebert 的代数化程序将 SYZ 几何翻译为 tropical 几何和散射图（scattering diagram）的语言。\(\blacksquare\)

\hrulefill

\emph{本卷在微分几何（V12）和代数拓扑（V15）的基础上，系统建立了辛几何（Darboux 定理、Hamilton 力学、Noether 定理）、辛约化（Marsden-Weinstein 约化、Delzant 定理）、Gromov-Witten 理论（伪全纯曲线、量子上同调、WDVV）、切触几何（Reeb 动力学、切触同调、Giroux 定理）和 Fukaya 范畴与镜对称（Lagrangian Floer 同调、Kontsevich 猜想、SYZ 猜想）的理论体系。共 5 章。}

\subsection{辛几何的数学全景与未来方向}\label{ux8f9bux51e0ux4f55ux7684ux6570ux5b66ux5168ux666fux4e0eux672aux6765ux65b9ux5411}

辛几何自 20 世纪 80 年代以来经历了革命性的发展，从经典力学中的 Hamilton 形式化演变为现代数学中最活跃的核心领域之一。\textbf{Gromov-Witten 理论}（1990s）将代数几何中曲线计数的经典问题推广到辛流形上，通过伪全纯曲线的模空间定义了 Gromov-Witten 不变量，其生成函数满足 \textbf{WDVV 方程}（Witten-Dijkgraaf-Verlinde-Verlinde），等价于量子上同调环的关联性------这一结构是二维拓扑引力理论和弦论中配分函数的数学基础。\textbf{Floer 同调}（1988）是辛几何中最重要的工具之一，它将 Morse 理论的无穷维推广应用于辛流形上的作用泛函，建立了 Hamiltonian Floer 同调 \(HF_*(H)\) 和 Lagrangian Floer 同调 \(HF(L_0, L_1)\)。Arnold 猜想的证明（Fukaya-Ono 1999, Liu-Tian 1998）是 Floer 同调的巅峰成就。\textbf{镜面对称}（mirror symmetry）是辛几何中最深刻的开放问题之一：Kontsevich 的同调镜对称猜想（1994）断言 \(D^b \mathcal{F}uk(X) \cong D^b \operatorname{Coh}(\check{X})\)，将辛几何的 Fukaya 范畴与复代数几何的凝聚层范畴等价起来。SYZ 猜想（Strominger-Yau-Zaslow 1996）通过特殊 Lagrangian 环面纤维化上的 T-对偶提供了镜面对称的几何解释。\textbf{辛场论}（Symplectic Field Theory, SFT, Eliashberg-Givental-Hofer 2000）是辛几何中最大的统一框架，它通过 Reeb 轨道的全纯曲线计数将 Gromov-Witten 理论、Floer 同调和切触同调统一为单一的代数结构。辛几何与\textbf{表示论}（通过莽形辛约化和几何量子化）、\textbf{可积系统}（通过 Liouville 可积性和 Arnold-Liouville 定理）、\textbf{代数几何}（通过深谷范畴和凝聚层范畴的镜对称）以及\textbf{低维拓扑}（通过切触三维流形和 Legendrian 纽结）的交叉，使辛几何成为 21 世纪数学中最具综合性的领域之一。

\newpage

\chapter{07-拓扑与几何/02-微分几何/05-非交换几何.md}\label{ux62d3ux6251ux4e0eux51e0ux4f5502-ux5faeux5206ux51e0ux4f5505-ux975eux4ea4ux6362ux51e0ux4f55.md}

\subsection{169.A 非交换几何补充}\label{a-ux975eux4ea4ux6362ux51e0ux4f55ux8865ux5145}

\begin{quote}
非交换几何由 Alain Connes 在 1980 年代创立，将经典的微分几何推广到非交换代数框架。其核心思想是：空间与函数代数是等价的（Gelfand-Naimark 定理），非交换代数对应''非交换空间''。本卷在 C*-代数（V14）和 von Neumann 代数（V31）的基础上，建立谱三元组、循环上同调、非交换微分形式和标准模型几何。Connes 因算子代数理论（特别是 III 型因子的分类与结构定理）以及对叶状结构和微分几何的应用获 1982 年 Fields 奖。
\end{quote}

\hrulefill

\hrulefill

\hrulefill

\hrulefill

\hrulefill

\hrulefill

\section{第232章：谱三元组}\label{ux7b2c232ux7ae0ux8c31ux4e09ux5143ux7ec4}

谱三元组（spectral triple）\((\mathcal{A}, \mathcal{H}, D)\) 是 Alain Connes 在 1980 年代末引入的非交换几何基本对象，其设计灵感直接来源于经典 Riemann 流形上 Dirac 算子的谱几何。在经典情形中，紧自旋流形 \(M\) 自然生成一个谱三元组：\(\mathcal{A} = C^\infty(M)\)（光滑函数代数），\(\mathcal{H} = L^2(M, S)\)（自旋子空间的平方可积截影），\(D\) 为 Dirac 算子。谱三元组的深刻之处在于------通过 Connes 关于紧流形的重构定理------(1) 算子 \(D\) 的谱数据与 \(\mathcal{A}\) 在其上的作用完全编码了流形 \(M\) 的度量结构、自旋结构和微分拓扑；(2) 这一编码在 \(\mathcal{A}\) 为非交换代数时仍然有效，从而为''非交换空间''赋予了 Riemann 几何学意义。在物理应用中，Connes-Lott 模型和 Chamseddine-Connes 的谱作用原理将粒子物理标准模型的 Higgs 机制和 Yang-Mills 理论统一为一个纯几何框架：Higgs 场被解释为离散空间中 Dirac 算子的联络分量。本章与后继的循环上同调（第171章）和非交换微分形式（第172章）构成三位一体：谱三元组提供了几何数据，循环上同调提供了积分与示性类的代数工具，而非交换微分形式则在两者之间建立了微分结构的桥梁。

\subsection{192.1 从 Gelfand-Naimark 到非交换空间}\label{ux4ece-gelfand-naimark-ux5230ux975eux4ea4ux6362ux7a7aux95f4}

\textbf{定理 192.1}（Gelfand-Naimark 定理回顾）：紧 Hausdorff 空间 \(X\) 上的连续函数代数 \(C(X)\) 是交换 C\emph{-代数。反之，任何交换 C}-代数 \(\cong C(X)\)（\(X\) 是其谱空间）。这建立了空间与交换代数之间的范畴等价。

\textbf{证明}：设 \(\mathcal{A}\) 为交换 C\emph{-代数。定义其谱空间 \(\widehat{\mathcal{A}} = \{\varphi: \mathcal{A} \to \mathbb{C} \mid \varphi \text{ 是非零乘性线性泛函（特征）}\}\)。每个 \(\varphi \in \widehat{\mathcal{A}}\) 满足 \(\|\varphi\| = 1\) 和 \(\varphi(a^*) = \overline{\varphi(a)}\)。\(\widehat{\mathcal{A}}\) 赋以弱-} 拓扑成为紧 Hausdorff 空间。Gelfand 变换 \(\Gamma: \mathcal{A} \to C(\widehat{\mathcal{A}})\)，\((\Gamma a)(\varphi) = \varphi(a)\)，是等距 \emph{-同构。这是因为 \(\operatorname{sp}(a) = \{\varphi(a) : \varphi \in \widehat{\mathcal{A}}\}\)，且 \(\|a\| = \sup_{\varphi} |\varphi(a)| = \|\Gamma a\|_\infty\)，故 \(\Gamma\) 保范且单射；满射性由 Stone-Weierstrass 定理保证。反之，对紧 Hausdorff 空间 \(X\)，\(C(X)\) 的谱空间同胚于 \(X\)。这给出了紧 Hausdorff 空间范畴与交换 C}-代数范畴的范畴反等价。此定理是 Connes 非交换几何的出发点：非交换 C*-代数对应''非交换空间''。\(\blacksquare\)

\textbf{Connes 的核心洞察}：将''空间''推广为''非交换 C*-代数''，从而''非交换空间''上的几何由非交换代数上的结构定义。

\textbf{定义 192.1}（谱三元组）：\textbf{谱三元组} \((\mathcal{A}, \mathcal{H}, \mathcal{D})\) 由以下组成： 1. \(\mathcal{A}\)：*-代数，在 Hilbert 空间 \(\mathcal{H}\) 上有界表示 \(\pi: \mathcal{A} \to \mathcal{B}(\mathcal{H})\) 2. \(\mathcal{D}\)：\(\mathcal{H}\) 上的自伴算子（未必有界），称为 \textbf{Dirac 算子} 3. \([\mathcal{D}, \pi(a)]\) 对所有 \(a \in \mathcal{A}\) 是有界算子 4. \(\mathcal{D}\) 具有紧预解式：\((\mathcal{D} - \lambda)^{-1}\) 对所有 \(\lambda \notin \mathbb{R}\) 是紧算子

谱三元组是 Connes 非交换几何的基本数据，它由三个对象组成：非交换代数 \(\mathcal{A}\)（取代函数空间）、Hilbert 空间 \(\mathcal{H}\)（表示空间）和 Dirac 算子 \(\mathcal{D}\)（编码度量几何）。这三个要素分别对应了经典几何中''空间上的函数''\,``旋量场空间''和''Dirac 算子''三重角色。条件 (3) 保证了 \(\mathcal{A}\) 与 \(\mathcal{D}\) 的交换子有界------即代数元素在度量意义下是 Lipschitz 连续的；条件 (4) 保证了 \(\mathcal{D}\) 的逆（在某种迹类理想中）具有紧性------对应几何空间具有有限的''测度''。谱三元组的魅力在于，它仅用泛函分析和算子代数语言就完全刻画了 Riemann 自旋流形的全部几何数据：度量、旋量结构、示性类和指标理论皆可由 \((\mathcal{A}, \mathcal{H}, \mathcal{D})\) 恢复。

\textbf{定义 192.2}（谱三元组的维数）：谱三元组有\textbf{维数} \(n\)（或更一般地，有\textbf{维数谱}），如果 \(|\mathcal{D}|^{-1}\) 属于 Dixmier 理想 \(\mathcal{L}^{(n,\infty)}\)，即 \(\mu_k(|\mathcal{D}|^{-1}) = O(k^{-1/n})\)（\(\mu_k\) 是奇异值）。

\subsection{192.2 Dirac 算子与非交换度量}\label{dirac-ux7b97ux5b50ux4e0eux975eux4ea4ux6362ux5ea6ux91cf}

\textbf{定理 192.2}（Connes 的距离公式，1989）：在经典情形的谱三元组 \((C^\infty(M), L^2(S), \mathcal{D}_M)\) 上，Riemann 距离被 Dirac 算子完全恢复：

\[d(x, y) = \sup\{|f(x) - f(y)| : f \in C^\infty(M), \|[D, f]\| \leq 1\}\]

这给出了非交换几何中''度量''的纯代数定义！

\textbf{证明}：对任意 \(f \in C^\infty(M)\)，\([D, f] = df\)（Clifford 乘法），其算子范数 \(\|[D, f]\| = \|df\|_\infty\)（Lipschitz 范数）。由中值定理，\(|f(x) - f(y)| \leq \|df\|_\infty \cdot d(x,y)\)。若 \(\|[D, f]\| \leq 1\)，则 \(|f(x) - f(y)| \leq d(x,y)\)。取 \(f(z) = d(x,z)\)（距离函数，在 \(x\) 附近光滑），则 \(\|[D, f]\| = 1\)（几乎处处），且 \(|f(x) - f(y)| = d(x,y)\)，故上确界达到 \(d(x,y)\)。因此距离公式给出 Riemann 距离的精确恢复。\(\blacksquare\)

Connes 距离公式的标志性成就在于：它仅用交换子 \([D, f]\) 的算子范数这一纯代数量就完全恢复了 Riemann 度量，从而给出了''非交换度量''的代数化定义。这意味着即使在代数 \(\mathcal{A}\) 为非交换的情形，只要给定了一个谱三元组，就可以通过同一上确界公式在 \(\mathcal{A}\) 的态空间上赋予一个距离。

\textbf{定义 192.3}（非交换度量）：对谱三元组 \((\mathcal{A}, \mathcal{H}, \mathcal{D})\)，\(\mathcal{A}\) 的态空间（纯态）上的距离定义为

\[d(\varphi, \psi) = \sup\{|\varphi(a) - \psi(a)| : a \in \mathcal{A}, \|[D, a]\| \leq 1\}\]

\subsection{192.3 正则性与有限性}\label{ux6b63ux5219ux6027ux4e0eux6709ux9650ux6027}

\textbf{定义 192.4}（正则谱三元组）：谱三元组是\textbf{正则}的，如果 \(\mathcal{A}\) 和 \([\mathcal{D}, \mathcal{A}]\) 属于 \(\bigcap_{k \geq 1} \operatorname{Dom}(\delta^k)\)，其中 \(\delta(T) = [|\mathcal{D}|, T]\)。

\textbf{定义 192.5}（有限谱三元组 / Fredholm 模）：谱三元组是\textbf{偶}的（even），如果存在 \(\mathbb{Z}_2\)-分次算子 \(\gamma\)（\(\gamma = \gamma^*\)，\(\gamma^2 = 1\)）使得 \(\gamma \mathcal{D} = -\mathcal{D}\gamma\) 和 \(\gamma a = a\gamma\) 对所有 \(a \in \mathcal{A}\)。否则是\textbf{奇}的（odd）。

\textbf{定义 192.6}（实结构）：谱三元组的\textbf{实结构}是反线性等距 \(J: \mathcal{H} \to \mathcal{H}\)，满足 \(J^2 = \pm 1\)，\(J\mathcal{D} = \pm \mathcal{D}J\)，\(J\gamma = \pm \gamma J\)（偶情形），以及 \([a, Jb^*J^{-1}] = 0\) 对所有 \(a, b \in \mathcal{A}\)（交换性条件）。

8 种可能的符号组合（\(J^2, J\mathcal{D}, J\gamma\) 的符号）对应 KO-维数模 8 的 8 种情形（Bott 周期性）。

\textbf{定理 192.3}（Connes 的谱三元组公理）：实谱三元组满足 7 条公理，包括正则性、有限性、一阶条件（\([[\mathcal{D}, a], Jb^*J^{-1}] = 0\)）、定向性、Poincaré 对偶等。这些公理精确刻画了''自旋流形上的 Dirac 算子''的代数量子化。

\emph{各公理的几何动机}：\textbf{(1) 正则性}：\(\mathcal{A}\) 及 \([\mathcal{D},\mathcal{A}]\) 对 \(\delta(\cdot)=[|\mathcal{D}|,\cdot]\) 的光滑性，对应流形上 \(C^\infty\) 函数在Dirac算子下的光滑作用。\textbf{(2) 有限性}：\(\bigcap_k \operatorname{Dom}(\mathcal{D}^k)\) 是有限投射 \(\mathcal{A}\)-模，推广了自旋丛截面的有限生成性。\textbf{(3) 一阶条件}：\([[\mathcal{D},a], Jb^*J^{-1}]=0\) 确保Dirac算子是一阶微分算子------\([\mathcal{D},a]\) 与右乘作用对易，这是流形上Dirac算子的本质特征。\textbf{(4) 定向性}：存在Hochschild \(n\)-圈 \(c\) 满足 \(\pi(c)=\gamma\)，编码了体积形式的存在性。\textbf{(5) Poincaré对偶}：Fredholm指标与 \(K\)-理论的配对非退化，对应经典流形上的Poincaré对偶。\textbf{(6) 实结构符号}：\(J^2, J\mathcal{D}, J\gamma\) 的八种符号组合对应KO-维数（Bott周期性模8），反映旋量表示的实结构分类。\textbf{(7) 谱维数}：\(|\mathcal{D}|^{-1} \in \mathcal{L}^{(n,\infty)}\) 指定了非交换空间的维数。∎

\hrulefill

\hrulefill

\hrulefill

\hrulefill

\hrulefill

\hrulefill

\section{第233章：循环上同调}\label{ux7b2c233ux7ae0ux5faaux73afux4e0aux540cux8c03}

循环上同调（cyclic cohomology）是 Alain Connes 在 1980 年代为非交换空间量身打造的上同调理论，它推广了经典 de Rham 上同调，并在非交换世界中扮演着示性类和积分理论的双重角色。经典的 de Rham 上同调 \(H^*_{dR}(M)\) 基于流形 \(M\) 上微分形式的代数结构，而 Connes 的洞察在于：在代数层面上，de Rham 上同调等价于光滑函数代数 \(C^\infty(M)\) 的循环上同调 \(HC^*(C^\infty(M))\)------因此，循环上同调是 de Rham 理论的''代数化''，可直接推广到任意（非交换）代数。循环上同调的技术核心是 Connes 的 \((b, B)\) 双复形：Hochschild 上边缘算子 \(b\)（编码代数的结合律）与 Connes 边界算子 \(B\)（编码循环置换的不变性）相互作用，形成一个混合复形，其全上同调即为循环上同调。在第170章谱三元组的语境下，循环上同调的具体应用包括：(1) Connes-Chern 特征标------将 K 理论映射到循环上同调的函子，相当于经典微分几何中的 Chern-Weil 同态；(2) 局部指标公式------将 Dirac 算子的 Fredholm 指标表示为循环上同调中元素的配对，这是 Atiyah-Singer 指标定理的非交换推广。

\subsection{193.1 Hochschild 与循环上同调}\label{hochschild-ux4e0eux5faaux73afux4e0aux540cux8c03}

\textbf{定义 193.1}（Hochschild 上链）：代数 \(\mathcal{A}\) 上的 \(n\)-\textbf{Hochschild 上链}是 \((n+1)\)-线性泛函 \(\varphi: \mathcal{A}^{n+1} \to \mathbb{C}\)。Hochschild 上边缘算子 \(b\) 为

\[(b\varphi)(a_0, \ldots, a_{n+1}) = \sum_{i=0}^n (-1)^i \varphi(a_0, \ldots, a_i a_{i+1}, \ldots, a_{n+1}) + (-1)^{n+1} \varphi(a_{n+1}a_0, a_1, \ldots, a_n)\]

\(b^2 = 0\)，定义 \textbf{Hochschild 上同调} \(HH^n(\mathcal{A})\)。

\textbf{定义 193.2}（循环上链）：\textbf{循环 \(n\)-上链}是满足循环条件的 Hochschild 上链：

\[\varphi(a_n, a_0, a_1, \ldots, a_{n-1}) = (-1)^n \varphi(a_0, a_1, \ldots, a_n)\]

\textbf{循环上边缘算子} \(B\)（Connes 算子）为

\[(B\varphi)(a_0, \ldots, a_{n-1}) = \sum_{i=0}^{n-1} (-1)^{(n-1)i} \varphi(1, a_i, a_{i+1}, \ldots, a_{i-1})\]

\(B^2 = 0\)，\(bB + Bb = 0\)。\textbf{循环上同调} \(HC^n(\mathcal{A})\) 是 \((b, B)\)-双复形的上同调。

Hochschild 上同调度量了代数的''结合性障碍''------它本质上是代数扩张和形变理论的计算工具；而循环上同调通过施加循环置换条件使理论适应于非交换微积分中迹类运算的需要。两者之间的桥梁是由 Connes 建立的 SBI 长正合列，它精确地量化了从 Hochschild 到循环上同调的信息升降关系。

\textbf{定理 193.1}（Connes 的 SBI 长正合列）：Hochschild 上同调和循环上同调之间存在长正合列：

\[\cdots \to HH^n(\mathcal{A}) \xrightarrow{I} HC^n(\mathcal{A}) \xrightarrow{S} HC^{n-2}(\mathcal{A}) \xrightarrow{B} HH^{n+1}(\mathcal{A}) \to \cdots\]

其中 \(S\) 是周期映射，\(I\) 是包含映射，\(B\) 是 Connes 算子。

\textbf{证明}：由 \((b, B)\)-双复形构造，循环上同调 \(HC^n\) 是双复形 \(\operatorname{Tot}(C^{\bullet,\bullet})\) 的上同调。双复形的谱序列 \(E_2^{p,q} = H_b^p(H_B^q)\) 收敛到 \(HC^{p+q}\)。因 \(B\) 是非循环的（acyclic），\(H_B^q \cong H_b^q\) 当 \(q=0\)，否则为零。故谱序列退化，给出长正合列。\(I\) 将 Hochschild 上循环视为特殊循环上循环，\(B\) 是 Connes 的 \(B\)-算子，\(S\) 是周期映射 \(HC^n \to HC^{n-2}\)。\(\blacksquare\)

\subsection{193.2 Chern 特征标}\label{chern-ux7279ux5f81ux6807}

循环上同调的关键价值在于它与 \(K\)-理论的对偶配对，而实现这一配对的工具正是 Chern 特征标。在经典微分几何中，Chern-Weil 理论将向量丛的拓扑不变量（Chern 类）表示为曲率形式的微分形式积分；在非交换几何中，Chern 特征标将 Fredholm 模的 \(K\)-理论类映射到循环上同调类，从而建立了非交换指标公式的代数基础。

\textbf{定义 193.3}（Fredholm 模的 Chern 特征标）：设 \((\mathcal{A}, \mathcal{H}, F)\) 是奇 Fredholm 模（\(F = \operatorname{sign}(\mathcal{D})\)）。其 \textbf{Chern 特征标}是循环上循环 \(\tau_n\)（\(n\) 为奇数）：

\[\tau_n(a_0, \ldots, a_n) = \operatorname{Tr}(F[F, a_0][F, a_1] \cdots [F, a_n])\]

其中 \(\operatorname{Tr}\) 是 Dixmier 迹（或其推广）。

\textbf{定理 193.2}（Chern 特征标与 \(K\)-理论）：Chern 特征标诱导从 \(K\)-理论到循环上同调的配对（Chern-Connes 配对）：

\[\langle [e], [\varphi] \rangle = \frac{1}{n!} \varphi_{2n}(e, e, \ldots, e)\]

其中 \([e] \in K_0(\mathcal{A})\)，\([\varphi] \in HC^{2n}(\mathcal{A})\)。这是非交换几何中指标定理的核心。

\textbf{证明}：对投影 \(e \in M_N(\mathcal{A})\)，Chern 特征标定义为 \(\operatorname{Ch}_n(e) = \frac{1}{n!} \operatorname{Tr}((e - \frac{1}{2}) de \cdots de)\)（\(2n\) 个 \(de\)）。配对 \(\langle [e], [\varphi] \rangle = \varphi_{2n}(e, \ldots, e)\) 诱导良定义的 \(K_0(\mathcal{A}) \times HC^{2n}(\mathcal{A}) \to \mathbb{C}\) 双线性型。由于 \(\varphi\) 是循环上循环，\(\varphi_{2n}\) 满足循环条件，该配对在 \(K_0\) 的稳定等价关系下不变，故诱导了 \(K_0(\mathcal{A}) \times HC^{2n}(\mathcal{A})\) 的配对。\(\blacksquare\)

\textbf{定理 193.3}（局部指标公式，Connes-Moscovici 1995）：对正则谱三元组，Chern 特征标在循环上同调中的类可由局部公式（涉及 \(\mathcal{D}\) 的预解式和 \(\mathcal{A}\) 的泛函演算）表示。这推广了 Atiyah-Singer 指标定理。

\textbf{证明}：设 \((\mathcal{A}, \mathcal{H}, \mathcal{D})\) 为 \(n\) 维正则谱三元组。Chern 特征标 \(\tau_n(a_0, \dots, a_n) = \operatorname{Tr}_\omega(\gamma a_0 [\mathcal{D}, a_1] \cdots [\mathcal{D}, a_n] |\mathcal{D}|^{-n})\) 在循环上同调中定义了一个类 \([\tau_n] \in HC^n(\mathcal{A})\)。局部指标公式将此循环上循环表示为 \(\tau_n(a_0, \dots, a_n) = \sum c_\alpha \int a_0 \nabla^{\alpha_1}(a_1) \cdots \nabla^{\alpha_n}(a_n) \nu\) 的形式，其中 \(\nabla\) 是某种非交换联络，\(\nu\) 是非交换体积形式，积分由非交换留数给出。关键步骤：利用 Getzler 渐近微积分，将 \(\operatorname{Tr}_\omega(\gamma a_0 e^{-t\mathcal{D}^2})\) 的热核展开与 \(\mathcal{D}\) 的预解式联系起来，提取 \(t \to 0\) 时的 \(t^{-n/2}\) 阶系数。正则性条件保证 \(\mathcal{A}\) 和 \([\mathcal{D}, \mathcal{A}]\) 对 \(|\mathcal{D}|\) 的交换子迭代足够光滑，使得该展开收敛。对经典 Dirac 算子的情形，局部公式恢复 Atiyah-Singer 被积式（\(\hat{A}\)-亏格与陈特征形式之积）。\(\blacksquare\)

\subsection{193.3 周期循环上同调}\label{ux5468ux671fux5faaux73afux4e0aux540cux8c03}

\textbf{定义 193.4}（周期循环上同调）：\textbf{周期循环上同调} \(HP^*(\mathcal{A})\)（\(* = 0, 1\)）定义为

\[HP^0(\mathcal{A}) = \varinjlim HC^{2n}(\mathcal{A}), \quad HP^1(\mathcal{A}) = \varinjlim HC^{2n+1}(\mathcal{A})\]

其中正向极限由 \(S\) 映射给出。周期循环上同调是 \(K\)-理论的对偶理论。

\textbf{定理 193.4}（Connes 的周期精确性）：设 \(0 \to \mathcal{J} \to \mathcal{A} \to \mathcal{A}/\mathcal{J} \to 0\) 是代数的扩张。则在 \(K\)-理论和周期循环上同调之间存在自然的六项正合列，推广了拓扑 \(K\)-理论的 Bott 周期性。

\textbf{证明}：代数扩张 \(0 \to \mathcal{J} \to \mathcal{A} \to \mathcal{A}/\mathcal{J} \to 0\) 诱导 \(K\)-理论的长正合列： \[K_0(\mathcal{J}) \to K_0(\mathcal{A}) \to K_0(\mathcal{A}/\mathcal{J}) \xrightarrow{\partial} K_1(\mathcal{J}) \to K_1(\mathcal{A}) \to K_1(\mathcal{A}/\mathcal{J})\] 同时，周期循环上同调函子 \(HP^*\) 对扩张具有 excision 性质：\(HP^*(\mathcal{A}/\mathcal{J}) \to HP^*(\mathcal{A}) \to HP^*(\mathcal{J})\) 是正合的。Chern-Connes 特征标 \(\operatorname{ch}: K_*(\mathcal{A}) \to HP^*(\mathcal{A})\) 提供了从 \(K\)-理论正合列到 \(HP^*\) 正合列的自然变换。由于 \(HP^*\) 具有周期 2（\(HP^* \cong HP^{*+2}\)），可以将六项裁剪为循环六项正合列。Connes 证明了此特征标在 \(HP^*\) 中诱导同构 \(K_*(\mathcal{A}) \otimes \mathbb{C} \cong HP^*(\mathcal{A})\)，这推广了拓扑 \(K\)-理论中的 Chern 特征标同构和 Bott 周期性。\(\blacksquare\)

\hrulefill

\hrulefill

\hrulefill

\hrulefill

\hrulefill

\hrulefill

\section{第234章：非交换微分形式}\label{ux7b2c234ux7ae0ux975eux4ea4ux6362ux5faeux5206ux5f62ux5f0f}

非交换微分形式是将经典微分几何中外代数 \(\Omega^\bullet(M)\) 的概念推广为非交换代数的系统方法，其意义在于为非交换空间赋予了可操作的''微分结构''------这是直接应用微分算子、定义联络和曲率、计算示性类的前提。经典情形中，流形 \(M\) 上的 de Rham 复形 \((\Omega^\bullet(M), d)\) 由函数代数通过外微分 \(d\) 逐次生成：\(d(fg) = df \cdot g + f \cdot dg\)（Leibniz 律）。在非交换框架下，Karoubi 引入了泛微分分次代数 \(\Omega^\bullet(\mathcal{A})\) 及其泛微分包络的思想：对任意代数 \(\mathcal{A}\)，存在一个''最一般的''分次微分代数，使得所有 \(\mathcal{A}\) 上的微分同态（到任意分次微分代数）唯一地通过 \(\Omega^\bullet(\mathcal{A})\) 分解。然而，泛微分的缺点在于过高阶（包含过多信息），实际几何应用通常需要通过商去适当理想来获得''几何相关的''微分形式。本章的主要构造包括：基于导子的微分形式（最接近经典 intuition）、非交换 de Rham 定理与 Hochschild 同调的联系，以及 Connes 定义的规范非交换微分形式------它在谱三元组的框架下与 \(|D|^{-1}\)（Dirac 算子的逆）的求和条件相协调。微分形式的理论为第170章和第171章之间的''几何-分析-代数''三角关系提供了封闭性。

\subsection{194.1 泛微分分次代数}\label{ux6cdbux5faeux5206ux5206ux6b21ux4ee3ux6570}

\textbf{定义 194.1}（泛微分分次代数）：代数 \(\mathcal{A}\) 的\textbf{泛微分分次代数}（universal differential graded algebra）\(\Omega^\bullet(\mathcal{A})\) 是 \(\mathcal{A}\) 上的自由分次代数，由 \(\mathcal{A}\)（次数 0）和 \(da\)（\(a \in \mathcal{A}\)，次数 1）生成，满足 \(d(ab) = da \cdot b + a \cdot db\) 和 \(d(1) = 0\)。

\(\Omega^k(\mathcal{A})\) 的元素是 \(a_0 da_1 da_2 \cdots da_k\) 的线性组合。

泛微分分次代数 \(\Omega^\bullet(\mathcal{A})\) 满足一个万有性质：从 \(\mathcal{A}\) 到任意分次微分代数的任何代数同态都可以唯一地通过 \(\Omega^\bullet(\mathcal{A})\) 提升为分次微分代数的同态。这使得 \(\Omega^\bullet(\mathcal{A})\) 成为研究非交换微分结构的''母对象''------所有具体的微分形式构造都可以通过商去适当理想从它出发得到。它与 Hochschild 同调的联系揭示了微分形式与代数循环的内在对应。

\textbf{命题 194.1}（泛微分形式与 Hochschild 同调）：\(\Omega^k(\mathcal{A})\) 与 Hochschild 同调 \(HH_k(\mathcal{A})\) 之间有自然联系，由映射

\[a_0 da_1 \cdots da_k \mapsto a_0 \otimes a_1 \otimes \cdots \otimes a_k\]

诱导。

\subsection{194.2 非交换积分}\label{ux975eux4ea4ux6362ux79efux5206}

微分形式给出了非交换空间上的''被积函数''，而完成积分理论则需要一个替代经典测度 \(\int_M f \, dV\) 的非交换迹。Connes 发现，Dixmier 迹------一种定义在紧算子理想上的广义迹------恰好提供了经典积分的非交换推广。其核心机制在于 Weyl 渐近公式：Dirac 算子的 \(|D|^{-n}\) 次幂的奇异值衰减速率 \(\mu_k \sim c \cdot k^{-1/n}\) 精确编码了流形的体积，而 Dixmier 迹从中提取出积分值。

\textbf{定义 194.2}（Dixmier 迹）：设 \(T\) 是紧算子，\(\mu_n(T)\) 为其奇异值。若 \(\frac{1}{\log N} \sum_{n=1}^N \mu_n(T)\) 有极限 \(\lim_{N \to \infty}\)，则此极限为 \(T\) 的 \textbf{Dixmier 迹} \(\operatorname{Tr}_\omega(T)\)（依赖于广义极限 \(\omega\)）。若该极限存在且与 \(\omega\) 无关，\(T\) 称为\textbf{可测}的。

\textbf{定义 194.3}（非交换积分）：设 \((\mathcal{A}, \mathcal{H}, \mathcal{D})\) 是 \(n\) 维谱三元组。\textbf{非交换积分}（noncommutative integral）定义为

\[\int a \, ds^n = \operatorname{Tr}_\omega(a |\mathcal{D}|^{-n})\]

当 \(a|\mathcal{D}|^{-n}\) 可测时。这推广了经典 Riemann 流形上的积分 \(\int_M f \, dV\)。

\textbf{定理 194.1}（Connes 的迹定理，1988）：对紧 \(n\) 维自旋流形 \(M\)，Dirac 算子 \(\mathcal{D}_M\) 满足

\[\operatorname{Tr}_\omega(f |\mathcal{D}_M|^{-n}) = c_n \int_M f \, dV\]

其中 \(c_n = 2^{[n/2]} \pi^{n/2} / (n \Gamma(n/2))\)。Dixmier 迹恢复了经典的积分！

\textbf{证明}：由 Weyl 渐近公式，\(|\mathcal{D}_M|^{-n}\) 的奇异值 \(\mu_k \sim (c_n \operatorname{Vol}(M)/k)^{-1/n}\)。对 \(f \in C^\infty(M)\)，\(f|\mathcal{D}_M|^{-n}\) 的奇异值满足 \(\mu_k(f|\mathcal{D}_M|^{-n}) \sim c_n^p \int_M |f|^p dV \cdot k^{-p/n}\)。取 \(p=1\)，Dixmier 迹 \(\operatorname{Tr}_\omega(f|\mathcal{D}_M|^{-n}) = \lim_{N \to \infty} \frac{1}{\log N} \sum_{k=1}^N \mu_k(f|\mathcal{D}_M|^{-n})\) 由渐近展开的主项给出 \(c_n \int_M f dV\)。常数 \(c_n\) 由 \(n\) 维球面 Dirac 算子的热核展开确定。\(\blacksquare\)

\textbf{定理 194.2}（非交换留数，Wodzicki 1984）：伪微分算子 \(P\) 的\textbf{非交换留数}（Wodzicki residue）定义为

\[\operatorname{Res}(P) = \operatorname{Tr}_\omega(P |\mathcal{D}|^{-n})\]

\(\operatorname{Res}\) 是迹（trace）------即 \(\operatorname{Res}([A, B]) = 0\)。在经典流形上，\(\operatorname{Res}(P)\) 等于 \(P\) 的符号的非交换留数（Guillemin 留数）。

\textbf{证明}：\(\operatorname{Res}([A,B]) = \operatorname{Tr}_\omega([A,B]|\mathcal{D}|^{-n}) = \operatorname{Tr}_\omega(A B |\mathcal{D}|^{-n} - B A |\mathcal{D}|^{-n})\)。由 Dixmier 迹的迹性质，\(\operatorname{Tr}_\omega(AB|\mathcal{D}|^{-n}) = \operatorname{Tr}_\omega(BA|\mathcal{D}|^{-n})\)（两算子乘积的奇异值渐近展开的可交换性），故 \(\operatorname{Res}([A,B]) = 0\)。在经典流形上，\(|\mathcal{D}|^{-n}\) 是 \(-n\) 阶伪微分算子，其符号与 \(P\) 的符号的乘积的 \(-n\) 阶部分为 \(P\) 的留数密度，积分即得 Guillemin 留数。\(\blacksquare\)

\subsection{194.3 非交换示性类}\label{ux975eux4ea4ux6362ux793aux6027ux7c7b}

\textbf{定义 194.4}（Chern 特征标，非交换形式）：在非交换几何中，投影 \(e \in M_N(\mathcal{A})\) 的 Chern 特征标由循环上循环

\[\operatorname{Ch}_n(e) = \operatorname{Tr}(e - \frac{1}{2}) de \cdots de \quad (2n \text{ 个 } de)\]

给出。这推广了微分几何中的 Chern-Weil 理论。

\textbf{定理 194.3}（非交换指标定理，Connes 1979）：对 \(p\)-可求和 Fredholm 模 \((\mathcal{A}, \mathcal{H}, F)\)，

\[\operatorname{Index}(eFe) = \langle \operatorname{Ch}_*(e), \tau_* \rangle\]

其中 \(\tau_*\) 是 Fredholm 模的 Chern 特征标。这是 Atiyah-Singer 指标定理的非交换推广。

\textbf{证明}：对 \(p\)-可求和 Fredholm 模，\(\tau_n\) 是良定义的循环上循环（因 \([F, a]\) 属于 Schatten 理想 \(\mathcal{L}^p\)）。\(eFe\) 的指标由 \(\operatorname{Index}(eFe) = \frac{1}{2} \operatorname{Tr}(F[F, e][F, e])\)（奇情形）或类似公式给出。将 \(\operatorname{Ch}_*(e)\) 展开为 \(\operatorname{Tr}((e - \frac{1}{2}) de \cdots de)\)，与 \(\tau_*\) 配对即得 \(\operatorname{Tr}(F[F, e] \cdots [F, e])\)。由迹的循环性，该配对不依赖于 \(e\) 在 \(K_0\) 中的代表元选择。对经典 Dirac 算子，\(\tau_*\) 局部化为 \(\hat{A}\)-类，恢复 Atiyah-Singer 公式。\(\blacksquare\)

\hrulefill

\hrulefill

\hrulefill

\hrulefill

\hrulefill

\newpage

\chapter{07-拓扑与几何/02-微分几何/06-专题.md}\label{ux62d3ux6251ux4e0eux51e0ux4f5502-ux5faeux5206ux51e0ux4f5506-ux4e13ux9898.md}

\section{第235章：标准模型的非交换几何}\label{ux7b2c235ux7ae0ux6807ux51c6ux6a21ux578bux7684ux975eux4ea4ux6362ux51e0ux4f55}

Connes 与 Lott（1991）和 Connes（1996）发现，粒子物理的标准模型（带 Higgs 场）可以自然地由非交换几何导出。这是非交换几何最壮观的应用。

\subsection{195.1 有限几何与内部空间}\label{ux6709ux9650ux51e0ux4f55ux4e0eux5185ux90e8ux7a7aux95f4}

\textbf{定义 195.1}（有限谱三元组）：\textbf{有限谱三元组} \((\mathcal{A}_F, \mathcal{H}_F, \mathcal{D}_F)\) 的 \(\mathcal{A}_F\) 是有限维代数，\(\mathcal{H}_F\) 是有限维 Hilbert 空间。有限谱三元组描述了''内部空间''（非交换空间）的几何。

\textbf{定义 195.2}（标准模型的有限代数）：标准模型的有限代数为

\[\mathcal{A}_F = \mathbb{C} \oplus \mathbb{H} \oplus M_3(\mathbb{C})\]

（\(\mathbb{H}\) 是四元数代数）。这对应标准模型的规范群 \(U(1) \times SU(2) \times SU(3)\)（模去离散子群）。

\textbf{定义 195.3}（有限 Hilbert 空间）：\(\mathcal{H}_F\) 是一代费米子的 Hilbert 空间（\(N\) 代，\(N=3\)），包含所有费米子场（轻子、夸克，分为左手和右手分量）。

\textbf{定义 195.4}（有限 Dirac 算子）：\(\mathcal{D}_F\) 是 \(\mathcal{H}_F\) 上的 \(N \times N\) 矩阵，其非零项是 Yukawa 耦合矩阵（给出费米子质量矩阵）和 Majorana 质量项（中微子）。

\subsection{195.2 谱作用量}\label{ux8c31ux4f5cux7528ux91cf}

\textbf{定义 195.5}（乘积谱三元组）：标准模型的整体谱三元组是经典几何与有限几何的乘积：

\[\mathcal{A} = C^\infty(M) \otimes \mathcal{A}_F, \quad \mathcal{H} = L^2(S) \otimes \mathcal{H}_F, \quad \mathcal{D} = \mathcal{D}_M \otimes 1 + \gamma_5 \otimes \mathcal{D}_F\]

\textbf{定理 195.1}（谱作用量原理，Connes-Chamseddine 1996-1997）：\textbf{谱作用量}（spectral action）定义为

\[S = \operatorname{Tr}_\omega\left(f\left(\frac{\mathcal{D}^2}{\Lambda^2}\right)\right) + \langle \psi, \mathcal{D} \psi \rangle\]

其中 \(f\) 是截断函数（cutoff function），\(\Lambda\) 是能量截断。展开谱作用量得到：

\[S = \int_M \left(\frac{1}{2g^2} F_{\mu\nu}F^{\mu\nu} + \cdots\right) d^4x + \text{Higgs 项} + \text{Yukawa 耦合}\]

即 Einstein-Hilbert 作用量 + Yang-Mills 作用量 + Higgs 作用量 + 费米子作用量------\textbf{整个标准模型}（包括 Higgs 场）从纯几何原理自动导出！

\textbf{证明}：由热核展开 \(f(\mathcal{D}^2/\Lambda^2) \sim \sum_{k=0}^\infty f_k \Lambda^{4-2k} a_k(\mathcal{D}^2)\)，其中 \(a_k\) 为 Seeley-DeWitt 系数。\(a_0\) 给出宇宙学常数项，\(a_2\) 给出 Einstein-Hilbert 作用量 \(\int R \sqrt{g} d^4x\)，\(a_4\) 给出 Yang-Mills 作用量、Higgs 动能项和势能项。\(\mathcal{D}\) 的构造中 \(\mathcal{D}_F\)（有限 Dirac 算子）的矩阵元编码了 Yukawa 耦合，而 \([\mathcal{D}, a]\) 的平方给出规范玻色子的动能项。费米子部分 \(\langle \psi, \mathcal{D} \psi \rangle\) 直接给出 Dirac 作用量。\(\blacksquare\)

\textbf{定理 195.2}（Higgs 场的几何起源）：在非交换几何中，Higgs 场自然地出现为内部空间的''联络''的分量------它是规范场在非交换方向上的分量。这解释了 Higgs 场的几何本质。

\textbf{证明}：在乘积谱三元组 \(\mathcal{D} = \mathcal{D}_M \otimes 1 + \gamma_5 \otimes \mathcal{D}_F\) 中，一般联络形式为 \(\nabla = d + A + \gamma_5 \otimes \Phi\)，其中 \(A\) 是 \(M\) 上的规范联络（取值于 \(\mathcal{A}_F\) 的 Lie 代数），\(\Phi \in \mathcal{A}_F\) 是标量场。\(\mathcal{D}\) 的涨落 \(\mathcal{D}_A = \mathcal{D} + A + JAJ^{-1}\) 展开后，\(\Phi\) 出现在非对角位置。谱作用量展开中 \(\Phi\) 的动能项和四次势能项精确给出 Higgs 作用量，且 \(\Phi\) 的真空期望值由 \(\mathcal{D}_F\) 的矩阵元确定。因此 Higgs 场是 \(\mathcal{D}_F\) 的''联络''在内部空间方向的分量。\(\blacksquare\)

\subsection{195.3 质量关系与预言}\label{ux8d28ux91cfux5173ux7cfbux4e0eux9884ux8a00}

\textbf{定理 195.3}（标准模型的约束）：非交换几何方法给出了标准模型参数之间的约束，包括： - Higgs 质量 \(m_H \approx 125\) GeV（实验值，LHC 2012） - 预测了 Higgs 自耦合常数 \(\lambda\) 与规范耦合常数之间的关系 - 轻微偏离标准模型：中微子必须是 Majorana 粒子，且存在质量

\textbf{证明}：谱作用量 \(S = \operatorname{Tr}_\omega(f(\mathcal{D}^2/\Lambda^2)) + \langle \psi, \mathcal{D} \psi \rangle\) 的 Seeley-DeWitt 展开至 \(a_4\) 项给出 Higgs 势能项 \(V(\Phi) = \lambda |\Phi|^4 - \mu^2 |\Phi|^2\)。其中 \(\lambda\) 和 \(\mu^2\) 由 \(\mathcal{D}_F\)（有限 Dirac 算子）的矩阵元通过热核系数 \(a_4\) 显式表达。在截断能标 \(\Lambda\) 处，跑动耦合常数满足 \(\lambda(\Lambda) = \frac{g_3^2(\Lambda)}{4}\)（由 \(\mathcal{D}_F\) 的代数结构导出），其中 \(g_3\) 是 \(SU(3)\) 规范耦合常数。由重整化群方程演化至电弱能标，此关系给出 \(m_H \approx 125\) GeV 的预测。此外，\(\mathcal{D}_F\) 的非对角块包含 Majorana 质量项，要求中微子为 Majorana 粒子。\(\mathcal{A}_F = \mathbb{C} \oplus \mathbb{H} \oplus M_3(\mathbb{C})\) 的表示结构决定了规范群 \(U(1) \times SU(2) \times SU(3)\) 及其耦合常数之间的关系。\(\blacksquare\)

\hrulefill

\hrulefill

\hrulefill

\hrulefill

\hrulefill

\hrulefill

\section{第236章：非交换 Riemann 几何与量子群}\label{ux7b2c236ux7ae0ux975eux4ea4ux6362-riemann-ux51e0ux4f55ux4e0eux91cfux5b50ux7fa4}

非交换 Riemann 几何将经典 Riemann 几何推广到非交换代数，量子群是其中的核心对象。

\subsection{196.1 非交换环面}\label{ux975eux4ea4ux6362ux73afux9762}

\textbf{定义 196.1}（非交换环面 \(\mathbb{T}_\theta^n\)）：\textbf{非交换 \(n\)-环面} \(\mathbb{T}_\theta^n\) 由旋转代数 \(A_\theta^n\) 描述，其生成元 \(U_1, \ldots, U_n\) 满足

\[U_k U_j = e^{2\pi i \theta_{jk}} U_j U_k\]

其中 \(\theta_{jk} = -\theta_{kj} \in \mathbb{R}/\mathbb{Z}\) 是反对称矩阵。\(A_\theta^n\) 是 \(C(\mathbb{T}^n)\) 的形变量子化。

\textbf{定理 196.1}（非交换环面的几何）：当 \(\theta\) 是''无理''矩阵（即 \(\theta_{jk}\) 是无理数）时，\(A_\theta^n\) 是单代数（无真双边理想）。非交换环面的 \(K\)-理论：\(K_0(A_\theta^n) \cong K_1(A_\theta^n) \cong \mathbb{Z}^{2^{n-1}}\)。

\textbf{证明}：以 \(n=2\) 为例，\(A_\theta^2\) 的 \(K_0\) 群由投射的等价类生成。利用 Pimsner-Voiculescu 六项正合列（将非交换环面嵌入为 \(C(\mathbb{T}) \rtimes_\theta \mathbb{Z}\) 的交叉积），得 \(K_0(A_\theta^2) \cong \mathbb{Z} \oplus \mathbb{Z}\)。当 \(\theta\) 为无理数时，迹映射 \(\tau: K_0 \to \mathbb{R}\) 的像为 \(\mathbb{Z} + \theta\mathbb{Z}\)，这是稠密子群。投射的分类由迹和 Chern 特征标完全确定。一般 \(n\) 维情形，\(K_*(A_\theta^n) \cong \bigwedge^* \mathbb{Z}^n \cong \mathbb{Z}^{2^{n-1}}\)，由生成元 \(U_k\) 的 Künneth 分解得到。非交换环面的 \(K\)-理论是 Connes 非交换微分几何的经典范例。\(\blacksquare\)

\textbf{定理 196.2}（非交换环面的 Gauss-Bonnet 定理，Connes-Tretkoff 2010）：非交换 2-环面 \(\mathbb{T}_\theta^2\) 上存在 Gauss-Bonnet 定理的推广，使用了非交换共形几何和局部指标公式。

\textbf{证明}：设 \((\mathcal{A}_\theta, \mathcal{H}, \mathcal{D})\) 为 \(\mathbb{T}_\theta^2\) 上的谱三元组，其中 \(\mathcal{A}_\theta = C^\infty(\mathbb{T}_\theta^2)\) 为非交换环面的光滑子代数，\(\mathcal{D}\) 为 Dirac 算子。Connes-Tretkoff 构造了非交换共形结构：对正元 \(k \in \mathcal{A}_\theta\)（共形因子），定义扭曲的 Dirac 算子 \(\mathcal{D}_k = k^{1/2} \mathcal{D} k^{1/2}\)。由局部指标公式，\(\operatorname{Res}_{z=0} \operatorname{Tr}(a \mathcal{D}_k^{-z})\) 给出 Einstein-Hilbert 作用量的非交换推广。计算 Dixmier 迹 \(\operatorname{Tr}_\omega(a \mathcal{D}_k^{-2})\) 得曲率泛函 \(\int_{\mathbb{T}_\theta^2} R \sqrt{g} d^2x\) 的非交换类比。对常值曲率情形，该迹等于 \(2\pi \chi(\mathbb{T}_\theta^2)\)（其中 \(\chi = 0\) 为 Euler 示性数），与经典 Gauss-Bonnet 定理一致。证明的关键是 Connes 的伪微分算子演算在非交换设定下的推广，以及局部指标公式中第二项的消失。\(\blacksquare\)

\subsection{196.2 量子群}\label{ux91cfux5b50ux7fa4}

\textbf{定义 196.2}（量子群 / Hopf 代数）：\textbf{量子群}（quantum group）是 Hopf 代数，带有余乘法 \(\Delta: \mathcal{A} \to \mathcal{A} \otimes \mathcal{A}\)、余单位 \(\varepsilon: \mathcal{A} \to \mathbb{C}\) 和对极 \(S: \mathcal{A} \to \mathcal{A}\)，满足相容性条件。

\textbf{定理 196.2.1}（Hopf代数对极的唯一性）：若 \(S, S'\) 均为 Hopf 代数 \(\mathcal{A}\) 的对极映射，则 \(S = S'\)。即对极映射由代数结构和余代数结构唯一确定。

\textbf{证明}：设 \(S, S'\) 均为对极，对任意 \(a \in \mathcal{A}\)，利用余乘的 Sweedler 记号 \(\Delta(a) = \sum a_{(1)} \otimes a_{(2)}\)。由对极公理 \(m \circ (S \otimes \operatorname{id}) \circ \Delta = u \circ \varepsilon\) 和 \(m \circ (\operatorname{id} \otimes S') \circ \Delta = u \circ \varepsilon\)，计算： \[S(a) = \sum S(a_{(1)}) \varepsilon(a_{(2)}) = \sum S(a_{(1)}) a_{(2)} S'(a_{(3)}) = \sum \varepsilon(a_{(1)}) S'(a_{(2)}) = S'(a)\] 其中利用了余结合律和余单位性质。故 \(S = S'\)。这一定理表明 Hopf 代数的''取逆''操作（对极）由代数-余代数结构强制确定，如同群中逆元由其乘积和单位唯一确定。∎

\textbf{详解 196.1}（Hopf 代数的完整公理体系）：Hopf 代数 \((\mathcal{A}, m, u, \Delta, \varepsilon, S)\) 同时是代数和余代数，结构映射满足六条公理：(1) \textbf{余乘法公理}：\(\Delta\) 是代数同态，即 \(\Delta(ab) = \Delta(a)\Delta(b)\)；(2) \textbf{余单位公理}：\(\varepsilon\) 是代数同态，且 \((\varepsilon \otimes \operatorname{id}) \circ \Delta = (\operatorname{id} \otimes \varepsilon) \circ \Delta = \operatorname{id}\)（左右余单位律）；(3) \textbf{余结合律}：\((\Delta \otimes \operatorname{id}) \circ \Delta = (\operatorname{id} \otimes \Delta) \circ \Delta\)，即两次余乘的顺序无关；(4) \textbf{对极公理}：\(S\) 是反代数同态，满足 \(m \circ (S \otimes \operatorname{id}) \circ \Delta = m \circ (\operatorname{id} \otimes S) \circ \Delta = u \circ \varepsilon\)，其中 \(m\) 是乘法，\(u\) 是单位映射；(5) \textbf{相容性条件}：\(\Delta(1) = 1 \otimes 1\)，\(\varepsilon(1) = 1\)，\(S(1) = 1\)；(6) 若 \(S\) 可逆（常用要求），则对极为对合：\(S \circ S = \operatorname{id}\)。此公理体系将群的三个运算（乘法、单位、取逆）双对偶化为余乘（描述''分裂''）、余单位（描述''删除''）和对极（描述''反转''），经典群的函数代数 \(C(G)\) 是交换 Hopf 代数，而非交换 Hopf 代数即为量子群。当 \(\mathcal{A}\) 是交换 Hopf 代数时，对应经典代数群；当 \(\mathcal{A}\) 是余交换 Hopf 代数时，对应李双代数。

\textbf{定义 196.3}（量子 \(SL(2)\) / \(SL_q(2)\)）：量子 \(SL(2)\)（Drinfeld-Jimbo 1985）由生成元 \(a, b, c, d\) 和关系

\[ab = qba, \quad ac = qca, \quad bd = qdb, \quad cd = qdc\]

\[bc = cb, \quad ad - da = (q - q^{-1})bc, \quad ad - q^{-1}bc = 1\]

定义。当 \(q = 1\) 时恢复经典 \(SL(2)\)。

\textbf{定理 196.3}（量子群与非交换几何）：量子群 \(SL_q(2)\) 上的非交换几何由谱三元组描述，其 Dirac 算子由 Dabrowski 等（2003）构造。这给出了量子群空间上的非交换 Riemann 几何。

\textbf{证明}：\(SL_q(2)\) 的 Hopf *-代数结构由余乘 \(\Delta\)、余单位 \(\varepsilon\) 和对极 \(S\) 定义。Dabrowski 等构造了 \(SL_q(2)\) 上的等变谱三元组 \((\mathcal{A}(SL_q(2)), \mathcal{H}, \mathcal{D})\)，其中 \(\mathcal{H}\) 是 \(L^2(SL_q(2))\) 的适当完备化，\(\mathcal{D}\) 借用经典 \(SU(2)\) 的 Dirac 算子经 \(q\)-变形得到。具体地，\(SU(2)\) 的 Dirac 算子 \(\mathcal{D}_{SU(2)}\) 作用在旋量丛截面上，其谱为 \(\{ \pm (j + 1/2) : j \in \frac{1}{2}\mathbb{N}_0 \}\)。\(q\)-变形后，\(\mathcal{D}_q\) 的谱变为 \(\{ \pm q^{j+1/2} [j+1/2]_q \}\)（\([n]_q = (q^n - q^{-n})/(q - q^{-1})\)）。\(\mathcal{D}_q\) 满足 \([\mathcal{D}_q, a]\) 有界对所有 \(a \in \mathcal{A}(SL_q(2))\) 成立，且 \(\mathcal{D}_q^{-1}\) 属于 Dixmier 理想 \(\mathcal{L}^{(3,\infty)}\)（维数 \(n=3\)）。该谱三元组具有 \(SU_q(2)\)-等变性，即 \(SU_q(2)\) 的 Hopf 代数作用与 Dirac 算子对易。这一构造证明了量子群空间上存在完整的非交换 Riemann 几何结构。\(\blacksquare\)

\subsection{196.3 Podleś 量子球面}\label{podleux15b-ux91cfux5b50ux7403ux9762}

Podleś 量子球面 \(S_q^2\)（Podleś, 1987）是量子齐性空间理论中最重要的示例之一，它展示了如何将经典对称空间的概念推广到量子群框架。经典球面 \(S^2 = SU(2)/U(1)\) 是 \(SU(2)\) 的齐性空间，而量子球面 \(S_q^2\) 是量子群 \(SU_q(2)\) 的量子齐性空间------即 \(C(S_q^2)\) 是 \(C(SU_q(2))\) 的某个子代数，在 \(SU_q(2)\) 的 Hopf 代数作用下不变。Podleś 的构造极为巧妙：他考虑 \(SU_q(2)\) 的生成元 \(a, c\) 满足 \(a^*a + c^*c = 1\)（\(q\)-变形关系），并定义 \(S_q^2\) 的坐标代数由 \(a^*a\) 和 \(a c^*\) 等特定的二次组合生成，这些组合在 \(U(1)\) 的量子变形下不变。这一构造表明量子球面的代数结构并非平凡的 \(q\)-变形------它有非交换的坐标环，但在 \(q \to 1\) 时还原为经典球面 \(S^2\) 上的多项式函数环。量子球面 \(S_q^2\) 的另一个重要性质是其 K 理论：\(K_0(C(S_q^2)) \cong \mathbb{Z} \oplus \mathbb{Z}\)，\(K_1(C(S_q^2)) \cong 0\)，与经典 \(S^2\) 的拓扑 K 理论完全一致------这体现了 \(q\)-变形在 K 理论层面的稳定性。此外，Podleś 量子球面是量子主丛 \(SU_q(2) \to S_q^2\) 的底空间，其纤维为量子 \(U(1)\)（即经典圆周 \(S^1\) 的 \(q\)-变形），这一主丛结构是非交换几何中''量子 Hopf 纤维化''的核心例子。量子球面在非交换微分几何中扮演了类似于经典球面在微分几何中的角色------它是研究非交换曲率、非交换 Dirac 算子和非交换指标定理的''试验场''。

\textbf{定义 196.4}（Podleś 量子球面，1987）：\textbf{Podleś 量子球面} \(S_q^2\) 是 \(SU_q(2)\) 的齐性空间（量子商空间）。\(C(S_q^2)\) 是 \(C(SU_q(2))\) 的子代数，由 \(SU_q(2)\) 的特定元素生成。

\textbf{定理 196.4}（量子球面的谱三元组，Dabrowski-Sitarz 2003）：\(S_q^2\) 上存在 \(SU_q(2)\)-等变谱三元组，其 Dirac 算子的谱具有明确的 \(q\)-变形形式。这是非交换 Riemann 几何在量子齐性空间上的重要实现。

\textbf{证明}：\(S_q^2\) 的坐标代数由 \(SU_q(2)\) 的生成元 \(a, c\) 满足 \(a^*a + c^*c = 1\)（\(q\)-变形关系）生成。Dabrowski-Sitarz 构造了 \(S_q^2\) 上的 Dirac 算子 \(\mathcal{D}_q\)，其作用在 \(L^2(S_q^2) \otimes \mathbb{C}^2\) 上。\(\mathcal{D}_q\) 的谱为 \(\{ \pm q^{-1/2} [n+1/2]_q : n \in \frac{1}{2}\mathbb{N}_0 \}\)（\([n]_q = (q^n - q^{-n})/(q - q^{-1})\)）。验证 \(SU_q(2)\)-等变性：\(\operatorname{Ad}_{U_q} \mathcal{D}_q = \mathcal{D}_q \otimes 1\)（\(U_q\) 为 \(SU_q(2)\) 的万有 \(R\)-矩阵的表示）。局部指标公式给出 \(\operatorname{Res}_{z=0} \operatorname{Tr}(a \mathcal{D}_q^{-z})\) 的非交换计算，恢复 \(q\)-变形的 Chern 特征标。当 \(q \to 1\) 时，该谱三元组还原为经典 \(S^2\) 上的 Dirac 算子。\(\blacksquare\)

\subsection{196.4 非交换几何的未来方向}\label{ux975eux4ea4ux6362ux51e0ux4f55ux7684ux672aux6765ux65b9ux5411}

非交换几何的前沿方向正在多个学科交叉处蓬勃发展。\textbf{导出非交换几何}将非交换几何与导出代数几何（Toën-Vezzosi 和 Lurie 的导出代数几何框架）结合，研究非交换空间的导出范畴和导出模空间。这一方向的核心思想是：非交换空间应被理解为''非交换 dg-代数''的谱（在导出意义下），而非交换代数几何的基本对象就是这些导出范畴。在导出非交换几何中，Kontsevich 的非交换代数几何纲领（基于形式 dg-范畴和非交换 Hodge 理论）与 Connes 的谱三元组方法通过 Chern 特征标的循环上同调解释相互连接，形成了统一的框架。\textbf{非交换几何与拓扑量子场论（TQFT）}的深层联系是另一个活跃方向：谱三元组中的 Dirac 算子通过 Chern-Simons 理论和辫子群表示自然地与三维拓扑量子场论相关联。具体而言，Connes-Kreimer 的量子场论重整化方案将 Feynman 图的 Hopf 代数结构与 Riemann-Hilbert 对应联系起来，而 Riemann-Hilbert 对应在非交换几何中由谱三元组的循环上同调编码。\textbf{非交换几何与数论}的联系也日益深化：Bost-Connes 系统（1995）将类域论的显式公式通过非交换动力系统实现，而 Connes-Marcolli 的 \(\mathbb{Q}\)-格点理论将 GL(2) 系统的 Langlands 对应与模曲线上的非交换空间联系起来。这些方向表明，非交换几何不仅是数学内部的统一框架，也是连接数学物理、数论和量子场论的核心桥梁。

\textbf{定理 196.5}（Baum-Connes猜想 {[}含等变KK-理论{]}）：对任意局部紧群 \(G\) 和 \(G\)-\(C^*\)-代数 \(A\)，\textbf{组装映射} \(\mu: K^G_*( \mathcal{E}G; A) \to K_*(C^*_r(G, A))\) 是群同构。其中 \(\mathcal{E}G\) 是 \(G\) 的万有真作用空间。

\textbf{证明}：等变 \(KK\)-理论 \(KK^G_*(A, B)\) 是 Kasparov 为研究 \(C^*\)-代数上群作用的指标问题引入的双函子，推广了拓扑 \(K\)-理论。Baum-Connes 猜想的左端是等变 \(K\)-同调（以 \(\mathcal{E}G\) 为空间参数的 Bredon 同调），右端是约化群 \(C^*\)-代数的 \(K\)-理论。组装映射将几何/拓扑数据（\(G\) 在 \(\mathcal{E}G\) 上的作用）映射到解析数据（约化交叉积 \(C^*_r(G,A)\) 的 \(K\)-理论）。该猜想蕴含了 Novikov 猜想、Gromov-Lawson-Rosenberg 正标量曲率猜想等。Higson-Kasparov（2001）对顺从群证明了该猜想；Lafforgue（2002）对双曲群证明了它。此猜想是非交换几何与几何拓扑之间最深刻的桥梁之一。\(\blacksquare\)

\hrulefill

\emph{本卷建立了非交换几何的核心理论框架：谱三元组（\(\mathcal{A}, \mathcal{H}, \mathcal{D}\)）将 Dirac 算子推广到非交换代数，Connes 距离公式恢复了度量；循环上同调（Hochschild/循环上同调、Chern 特征标、局部指标公式）是非交换空间的上同调理论；非交换微分形式（泛微分分次代数、Dixmier 迹、非交换积分）为非交换几何提供了积分和示性类；标准模型的非交换几何是 Connes 理论最壮观的应用，从纯几何原理导出了整个标准模型（包括 Higgs 场）；非交换环面、量子群和量子球面展示了非交换空间的具体例子。非交换几何是当代数学最深刻和最具野心的统一理论之一。}

\hrulefill

\emph{本卷从经典曲线论和曲面论出发，建立了 Frenet 标架、第一/第二基本形式、Gauss 曲率等核心概念，证明了 Gauss 绝妙定理和 Gauss-Bonnet 定理。随后引入微分流形和 Riemann 几何的现代框架，包括切空间、微分形式、Levi-Civita 联络、曲率张量等。讨论了曲率与拓扑的深层联系（Bonnet-Myers 定理、Hadamard-Cartan 定理、Chern-Gauss-Bonnet 定理）。补充内容涵盖辛几何与切触几何（Darboux 定理、Marsden-Weinstein 约化、Gromov-Witten 理论、Fukaya 范畴与镜对称）、非交换几何（Connes 理论、谱三元组、循环上同调、非交换环面、标准模型几何重构），为后续学习代数拓扑（V14）和李理论（V23）奠定了基础。}

\hrulefill

\hrulefill

\hrulefill

\hrulefill

\hrulefill

\hrulefill

\begin{quote}
微分几何在 Riemann 几何和辛几何的经典理论之外，还包含多个重要的专题方向，包括 Finsler 几何（Riemann 度量的各向异性推广）、谱几何（Laplace 算子的谱与几何的关系）、积分几何（几何量在群作用下的积分平均）、子流形几何（嵌入流形的外在曲率）和信息几何（统计流形上的 Riemann 度量）。这些专题在不同领域都有重要应用，从生物形态学（Finsler 几何）到量子力学和医学成像（谱几何和积分几何），构成了现代微分几何的丰富图景。
\end{quote}

\textbf{微分几何专题的现代发展}

微分几何的专题方向在 20 世纪下半叶经历了深刻的交叉融合，其共同特征是将 Riemann 几何的框架推广到更一般的几何结构上。\textbf{Finsler 几何}（Finsler 1918, Busemann 1950s, Chern 1948）是 Riemann 几何最自然的推广：将切空间的二次型度量替换为更一般的 Finsler 度量 \(F(x, y)\)（满足正齐性和强凸性）。Finsler 几何中的弧长和测地线不再由 Riemann 度量唯一确定，而是依赖于方向------这使其在优化控制理论（如 Zermelo 导航问题）和生物形态学（如植物生长模式）中具有自然的应用。\textbf{谱几何}（Kac 1966）研究 Laplace-Beltrami 算子 \(\Delta_g\) 的谱 \(\{\lambda_k\}\) 与 Riemann 流形 \((M, g)\) 的几何性质之间的关系。Weyl 渐近公式 \(N(\lambda) \sim \frac{\omega_n}{(2\pi)^n} \text{Vol}(M) \lambda^{n/2}\) 表明谱决定了体积和维数，但 Kac 问题的否定回答（Gordon-Webb-Wolpert 1992）表明谱不能唯一决定等距类。\textbf{积分几何}（Blaschke 1930s, Santalo 1976）的核心公式是 Crofton 公式：平面曲线的长度可通过随机直线与曲线的交点数的期望来表示。Radon 变换（1917）将函数映射为其在超平面上的积分，其反演公式是现代 CT 扫描的数学基础。\textbf{子流形几何}研究等距浸入 \(f: M \to \bar{M}\) 的外在曲率，其基本方程是 Gauss 方程、Codazzi 方程和 Ricci 方程，将子流形的内蕴曲率（Riemann 曲率张量）与外在曲率（第二基本形式）联系起来。\textbf{信息几何}（Amari 1985）将统计流形（概率分布族的参数空间）赋予 Fisher 信息度量 \(g_{ij}(\theta) = \mathbb{E}_\theta[\partial_i \log p(x|\theta) \partial_j \log p(x|\theta)]\)，使统计推断中的 Cramer-Rao 不等式和最大似然估计具有了自然的几何解释。

\hrulefill

\hrulefill

\section{第237章：其他几何专题}\label{ux7b2c237ux7ae0ux5176ux4ed6ux51e0ux4f55ux4e13ux9898}

本章汇集了微分几何中若干具有独立研究方向的重要专题，它们从不同角度拓展和深化了 Riemann 几何的框架。Finsler 几何将切空间的二次型度量替换为更一般的 Minkowski 范数，允许各向异性的度量结构；谱几何研究 Laplace-Beltrami 算子的特征值谱与底流形几何性质之间的深刻联系------``能否听出鼓的形状''（Kac 问题）是其中最著名的问题；积分几何通过群不变测度下的积分平均将几何量（长度、面积、曲率）转化为概率期望，Radon 变换的反演公式是现代 CT 扫描的数学基础。这些专题的共同特征是：它们将 Riemann 几何的核心思想推广到更广泛的数学结构之中，同时与物理学（量子力学、广义相对论）、工程学（医学成像）和统计学（信息几何）产生了深刻的交叉。

\subsection{67.1 Finsler 几何简介}\label{finsler-ux51e0ux4f55ux7b80ux4ecb}

Finsler 几何是 Riemann 几何的重要推广，由 Paul Finsler 在其 1918 年的博士论文中系统建立。

\textbf{定义 200.1}（Finsler 度量）：设 \(M\) 为 \(n\) 维光滑流形，\(TM\) 为其切丛。\(M\) 上的 \textbf{Finsler 度量}是切丛上的连续函数 \(F : TM \to [0, +\infty)\)，在 \(TM \setminus \{0\}\)（去零截面）上光滑，满足：

\begin{enumerate}
\def\labelenumi{(\arabic{enumi})}
\item
  \textbf{正齐性}：\(F(x, \lambda y) = \lambda F(x, y)\) 对所有 \(\lambda > 0\) 成立；
\item
  \textbf{强凸性}：对每个 \((x, y) \in TM \setminus \{0\}\)，基本张量（fundamental tensor）
\end{enumerate}

\[g_{ij}(x, y) = \frac{1}{2}\frac{\partial^2 F^2}{\partial y^i \partial y^j}(x, y)\]

是一个正定矩阵。\((M, F)\) 称为 \textbf{Finsler 流形}。

Finsler 流形上的\textbf{弧长}定义为 \(\ell(\gamma) = \int_a^b F(\gamma(t), \dot{\gamma}(t)) dt\)，其测地线由 Euler-Lagrange 方程确定。

\textbf{定义 200.2}（旗曲率 / Flag Curvature）：Finsler 几何中最重要的曲率不变量是\textbf{旗曲率}（flag curvature）。对切向量 \(y \in T_xM \setminus \{0\}\) 和与 \(y\) 线性无关的单位向量 \(v \in T_xM\)（``旗杆'' \(y\) + ``旗面'' \(v\)），旗曲率定义为

\[K(y, v) = \frac{\langle R_y(v), v \rangle}{\langle v, v \rangle_y \langle y, y \rangle_y - \langle v, y \rangle_y^2}\]

其中 \(R_y\) 为 Riemann 曲率（由 Chern 联络的 \(hh\)-曲率张量导出），\(\langle\cdot,\cdot\rangle_y\) 为 \(g_{ij}(x, y)\) 诱导的内积。当 \(F\) 为 Riemann 度量时，\(K(y, v)\) 退化到经典的截面曲率。

\subsection{67.2 谱几何（Spectral Geometry）}\label{ux8c31ux51e0ux4f55spectral-geometry}

谱几何研究紧致 Riemann 流形 \((M, g)\) 上 \textbf{Laplace-Beltrami 算子} \(\Delta_g\) 的谱与几何性质之间的对应关系。

\textbf{定义 200.3}（Laplace-Beltrami 算子的谱）：在紧致无边 Riemann 流形 \((M, g)\) 上（可带边界，此时配以 Dirichlet 或 Neumann 边界条件），Laplace-Beltrami 算子 \(\Delta_g = -\operatorname{div}_g \circ \nabla_g\) 是 \(\mathcal{C}^\infty(M)\) 上的二阶椭圆算子。其谱由可数个离散特征值组成：

\[0 = \lambda_0 < \lambda_1 \leq \lambda_2 \leq \cdots \nearrow \infty\]

（Dirichlet 条件下 \(\lambda_0 > 0\)）。对应的特征函数 \(\{\phi_k\}\) 构成 \(L^2(M, dV_g)\) 的完备标准正交基：\(\Delta_g \phi_k = \lambda_k \phi_k\)。

\textbf{定理 200.1}（Weyl 渐近律，1911）：设 \(N(\lambda) = \#\{k : \lambda_k \leq \lambda\}\) 为特征值计数函数。则

\[N(\lambda) \sim \frac{\omega_n}{(2\pi)^n} \operatorname{Vol}(M) \, \lambda^{n/2}, \quad \lambda \to \infty\]

其中 \(\omega_n = \pi^{n/2} / \Gamma(n/2 + 1)\) 为 \(\mathbb{R}^n\) 中单位球的体积。该渐近式表明谱的前导项决定了流形的维数 \(n\) 和体积 \(\operatorname{Vol}(M)\)。

\textbf{定理 200.2}（热核渐近展开 / Minakshisundaram-Pleijel, 1949）：热核 \(K(t, x, y) = \sum_{k=0}^\infty e^{-\lambda_k t} \phi_k(x)\phi_k(y)\) 在对角线 \(x = y\) 上具有渐近展开

\[K(t, x, x) \sim \frac{1}{(4\pi t)^{n/2}} \sum_{j=0}^\infty a_j(x) \, t^j, \quad t \to 0^+\]

其中 \(a_0(x) = 1\)，\(a_1(x) = \frac{1}{6}R(x)\)（\(R\) 为标量曲率），\(a_2(x)\) 涉及曲率张量的二次组合。积分给出热迹 \(\operatorname{Tr}(e^{-t\Delta_g}) \sim (4\pi t)^{-n/2} \sum_{j=0}^\infty a_j t^j\)（\(a_j = \int_M a_j(x) dV_g\)），其各阶系数从谱中可确定------因此谱决定了体积（\(a_0\)）、总标量曲率（\(a_1\)）等全局几何不变量。

\subsection{67.3 积分几何（Integral Geometry）}\label{ux79efux5206ux51e0ux4f55integral-geometry}

积分几何研究几何量在群作用下的积分平均，由 Blaschke 学派于 20 世纪初创立。

\textbf{定义 200.4}（运动不变测度与 Crofton 公式）：在 \(\mathbb{R}^2\) 中，直线的运动不变测度由 \(d\ell = dp \, d\theta\) 给出（Hesse 法式表示）。核心公式------\textbf{Crofton 公式}（1868）将平面可求长曲线 \(\Gamma\) 的长度 \(L\) 表示为

\[L = \frac{1}{2} \iint_{\ell} n_\ell(\Gamma) \, dp \, d\theta\]

其中 \(n_\ell(\Gamma)\) 为直线 \(\ell\) 与 \(\Gamma\) 的交点个数，积分对平面上的全体直线进行。等价地，曲线长度等于其与随机直线相交的期望数乘以 \(\pi/2\)。

\textbf{定义 200.5}（Radon 变换与反演公式）：\textbf{Radon 变换}将 \(\mathbb{R}^n\) 上的函数 \(f\) 映射为其在超平面上的积分值：

\[(\mathcal{R}f)(H) = \int_H f \, d\sigma\]

其中 \(H\) 为 \(\mathbb{R}^n\) 中的超平面，\(d\sigma\) 为 \(H\) 上的 Lebesgue 测度。Johann Radon 于 1917 年建立了反演公式：在适当的衰减条件下，

\[f(x) = c_n \iint_{\{H : x \in H\}} (\mathcal{R}f)(H) \, d\mu(H)\]

该公式是 X 射线 CT 扫描的理论基础------\(\mathcal{R}f\) 对应每个方向上的线积分投影，反演公式通过逆向投影重建原始密度函数 \(f\)。

\textbf{定理 200.3}（Blaschke-Santalo 不等式）：设 \(K \subset \mathbb{R}^n\) 为凸体，\(K^\circ = \{y : \langle x, y \rangle \leq 1, \forall x \in K\}\) 为其极体。以 Santalo 点为中心时，

\[\operatorname{Vol}(K) \cdot \operatorname{Vol}(K^\circ) \leq (\operatorname{Vol}(B^n))^2\]

等号成立当且仅当 \(K\) 是椭球。对上述乘积的下界，Mahler 猜想（1948）断言 \(\operatorname{Vol}(K) \cdot \operatorname{Vol}(K^\circ) \geq \frac{(n+1)^{n+1}}{(n!)^2}\)，最小乘积极值由单纯形达到。

积分几何将度量几何性质（长度、面积、曲率）转化为群不变测度下的积分平均值，深刻体现了几何与分析、概率论的对偶关系，在凸几何、随机几何和医学成像中均有重要应用。

\emph{卷十一：微分几何终。}

\emph{卷十一：微分几何终。}

\emph{卷十一：微分几何终。}

\emph{卷十三：微分几何终。}

\newpage

\chapter{07-拓扑与几何/03-代数拓扑/01-同调与上同调.md}\label{ux62d3ux6251ux4e0eux51e0ux4f5503-ux4ee3ux6570ux62d3ux625101-ux540cux8c03ux4e0eux4e0aux540cux8c03.md}

\section{第238章：奇异同调}\label{ux7b2c238ux7ae0ux5947ux5f02ux540cux8c03}

同调论是代数拓扑的核心------它将拓扑空间的几何直觉（``洞''和''边界''）转换为代数不变量（Abel 群）。奇异同调（singular homology）由 Samuel Eilenberg 在 1944 年引入，其绝妙之处在于：不依赖空间的任何三角剖分或胞腔结构，而是通过从标准单形到空间的所有连续映射（奇异单形）来编码空间的拓扑信息。这一''万能的''构造使奇异同调成为定义在任何拓扑空间上的函子，且满足 Eilenberg-Steenrod 公理体系的全部条件：同伦不变性、切除公理和维数公理。从计算角度看，奇异链复形 \(C_*(X)\) 通常是无穷维的，但通过胞腔分解（CW 复形）和 Mayer-Vietoris 序列可以高效计算具体空间的同调群。在代数拓扑的整体框架中，本章与后续的上同调与杯积（第177章）密不可分：同调群的系数序列 \(\cdots \to H_n(X) \to H_{n-1}(X) \to \cdots\) 提供了''洞''的计数，而上同调 \(H^n(X)\) 通过万有系数定理作为其对偶而出现，额外的杯积乘法结构则将上同调提升为一个分次交换环，为空间分类提供了更强的判别工具。

\subsection{77.1 奇异单形与链复形}\label{ux5947ux5f02ux5355ux5f62ux4e0eux94feux590dux5f62}

\textbf{定义 77.1}（标准单形）：\(n\) 维\textbf{标准单形} \(\Delta^n \subseteq \mathbb{R}^{n+1}\) 定义为

\[\Delta^n = \left\{(t_0, \ldots, t_n) \in \mathbb{R}^{n+1} : t_i \geq 0, \sum_{i=0}^n t_i = 1\right\}\]

\(\Delta^n\) 的顶点为 \(e_0 = (1, 0, \ldots, 0), \ldots, e_n = (0, \ldots, 0, 1)\)。

\textbf{定义 77.2}（面映射）：第 \(i\) 个\textbf{面映射} \(F_i^n: \Delta^{n-1} \to \Delta^n\) 定义为

\[F_i^n(t_0, \ldots, t_{n-1}) = (t_0, \ldots, t_{i-1}, 0, t_i, \ldots, t_{n-1})\]

即将 \(\Delta^{n-1}\) 嵌入为 \(\Delta^n\) 中与顶点 \(e_i\) 相对的面。

\textbf{定义 77.3}（奇异单形）：拓扑空间 \(X\) 中的一个\textbf{奇异 \(n\)-单形}是一个连续映射 \(\sigma: \Delta^n \to X\)。所有奇异 \(n\)-单形生成的自由 Abel 群记作 \(C_n(X)\)（系数在 \(\mathbb{Z}\) 中，若系数在 \(R\) 中则记作 \(C_n(X; R)\)）。

\textbf{定义 77.4}（边缘算子）：边缘算子 \(\partial_n: C_n(X) \to C_{n-1}(X)\) 在生成元上定义为

\[\partial_n(\sigma) = \sum_{i=0}^n (-1)^i \sigma \circ F_i^n\]

其中 \(\sigma \circ F_i^n\) 是 \(\sigma\) 的第 \(i\) 个面。

\textbf{命题 77.1}：\(\partial_{n-1} \circ \partial_n = 0\)。即 \((C_\bullet(X), \partial)\) 是链复形。

\emph{证明}：由面映射的恒等式 \(F_i^{n+1} \circ F_j^n = F_{j+1}^{n+1} \circ F_i^n\)（\(i \leq j\)），展开 \(\partial_{n-1} \circ \partial_n\) 后各项成对抵消。∎

\textbf{定义 77.5}（奇异同调群）：\(X\) 的 \(n\) 次\textbf{奇异同调群}为

\[H_n(X) = \frac{\ker \partial_n}{\operatorname{im} \partial_{n+1}}\]

\(Z_n(X) = \ker \partial_n\) 中的元素称为\textbf{闭链}（cycle），\(B_n(X) = \operatorname{im} \partial_{n+1}\) 中的元素称为\textbf{边缘链}（boundary）。

\subsection{77.2 同调的基本性质}\label{ux540cux8c03ux7684ux57faux672cux6027ux8d28}

\textbf{命题 77.2}（\(H_0\) 的计算）：若 \(X\) 非空且道路连通，则 \(H_0(X) \cong \mathbb{Z}\)。一般地，\(H_0(X)\) 是自由 Abel 群，秩等于 \(X\) 的道路连通分支数。

\emph{证明}：\(\partial_1(\sigma) = \sigma(1) - \sigma(0)\)。对道路连通空间，所有 \(0\)-单形（即点）互相等价，生成一个 \(\mathbb{Z}\)。∎

\textbf{定义 77.6}（诱导同态）：连续映射 \(f: X \to Y\) 诱导链映射 \(f_\#: C_n(X) \to C_n(Y)\)，\(f_\#(\sigma) = f \circ \sigma\)。\(f_\#\) 与边缘算子交换，故诱导同调群同态 \(f_*: H_n(X) \to H_n(Y)\)。

\textbf{命题 77.3}（同伦不变性）：若 \(f \simeq g: X \to Y\)（同伦），则 \(f_* = g_*: H_n(X) \to H_n(Y)\)。特别地，同伦等价的空间具有同构的同调群。

\textbf{证明}：设 \(H: X \times [0,1] \to Y\) 为 \(f\) 到 \(g\) 的同伦。构造棱柱算子 \(P: C_n(X) \to C_{n+1}(Y)\)。将 \(\Delta^n \times [0,1]\) 剖分为 \(n+1\) 个 \((n+1)\)-单形：顶点为 \((e_i,0)\) 和 \((e_i,1)\)，剖分为 \([v_0 v_1 \cdots v_i v_i' v_{i+1}' \cdots v_n']\)（\(i=0,\ldots,n\)）。定义 \(P(\sigma) = \sum_{i=0}^n (-1)^i H \circ (\sigma \times \operatorname{id})|_{[v_0\cdots v_i v_i'\cdots v_n']}\)。验证 \(\partial P(\sigma) + P(\partial\sigma) = H \circ (\sigma \times \operatorname{id})|_{[v_0'\cdots v_n']} - H \circ (\sigma \times \operatorname{id})|_{[v_0\cdots v_n]} = g_\#(\sigma) - f_\#(\sigma)\)。故 \(f_\#\) 和 \(g_\#\) 链同伦，在链复形上诱导相同的同调映射。\(\blacksquare\)

\textbf{推论 77.4}：可缩空间的同调群：\(H_0(X) \cong \mathbb{Z}\)，\(H_n(X) = 0\)（\(n \geq 1\)）。

\subsection{77.3 相对同调与切除定理}\label{ux76f8ux5bf9ux540cux8c03ux4e0eux5207ux9664ux5b9aux7406}

\textbf{定义 77.7}（相对同调群）：设 \(A \subseteq X\) 是子空间。定义相对链群 \(C_n(X, A) = C_n(X) / C_n(A)\)。边缘算子诱导 \(\partial: C_n(X, A) \to C_{n-1}(X, A)\)。\textbf{相对同调群} \(H_n(X, A)\) 是此链复形的同调。

\textbf{命题 77.5}（长正合序列）：对空间对 \((X, A)\)，存在自然的同调长正合序列：

\[\cdots \to H_n(A) \to H_n(X) \to H_n(X, A) \xrightarrow{\partial} H_{n-1}(A) \to \cdots \to H_0(X, A) \to 0\]

\emph{证明}：由短正合序列 \(0 \to C_\bullet(A) \to C_\bullet(X) \to C_\bullet(X, A) \to 0\) 应用蛇引理（V12 定理 70.1）。∎

\textbf{定理 77.6}（切除定理）：设 \(U \subseteq A \subseteq X\) 满足 \(\overline{U} \subseteq \operatorname{int}(A)\)。则包含映射 \((X \setminus U, A \setminus U) \hookrightarrow (X, A)\) 诱导同调同构：

\[H_n(X \setminus U, A \setminus U) \cong H_n(X, A)\]

\textbf{证明}：使用重心重分算子 \(S: C_n(X) \to C_n(X)\)。对奇异单形 \(\sigma: \Delta^n \to X\)，\(S(\sigma)\) 为 \(\sigma\) 的重心重分（将 \(\Delta^n\) 细分为 \((n+1)!\) 个小单形）。\(S\) 是链映射且链同伦于恒等：存在 \(T: C_n(X) \to C_{n+1}(X)\) 使 \(\partial T + T\partial = \operatorname{id} - S\)。迭代 \(S^m\) 得足够细的重分，使每个小单形要么落在 \(A\) 中要么落在 \(X \setminus U\) 中（因 \(\overline{U} \subseteq \operatorname{int}(A)\) 且重分直径趋于零）。\(S^m\) 与恒等链同伦，故 \(H_n(X \setminus U, A \setminus U) \to H_n(X, A)\) 诱导同构。\(\blacksquare\)

\textbf{推论 77.7}（Mayer-Vietoris 序列）：设 \(X = U \cup V\)，其中 \(U, V\) 是开集（或 \(\operatorname{int}(U) \cup \operatorname{int}(V) = X\)）。则存在长正合序列：

\[\cdots \to H_n(U \cap V) \to H_n(U) \oplus H_n(V) \to H_n(X) \xrightarrow{\partial} H_{n-1}(U \cap V) \to \cdots\]

\emph{证明}：考虑短正合序列 \(0 \to C_\bullet(U \cap V) \to C_\bullet(U) \oplus C_\bullet(V) \to C_\bullet(U) + C_\bullet(V) \to 0\)。由切除定理，\(C_\bullet(U) + C_\bullet(V) \hookrightarrow C_\bullet(X)\) 是链同伦等价。应用蛇引理。∎

\subsection{77.4 球面的同调}\label{ux7403ux9762ux7684ux540cux8c03}

\textbf{定理 77.8}（球面的同调群）：

\[H_n(S^k) \cong \begin{cases} \mathbb{Z} & n = 0 \text{ 或 } n = k \\ 0 & \text{其他} \end{cases}\]

\emph{证明}：对 \(k\) 归纳。\(S^0\) 是两个点，\(H_0(S^0) \cong \mathbb{Z}^2\)，高阶同调为零。归纳步骤：将 \(S^k\) 分解为两个半球 \(D^k_+\) 和 \(D^k_-\)，交为 \(S^{k-1}\)。应用 Mayer-Vietoris 序列：

\[\cdots \to H_n(S^{k-1}) \to H_n(D^k) \oplus H_n(D^k) \to H_n(S^k) \to H_{n-1}(S^{k-1}) \to \cdots\]

\(D^k\) 可缩，故 \(H_n(D^k) = 0\)（\(n \geq 1\)）。因此 \(H_n(S^k) \cong H_{n-1}(S^{k-1})\) 对 \(n \geq 2\) 成立。由归纳假设得结论。∎

\textbf{推论 77.9}（Brouwer 不动点定理，一般维数）：任意连续映射 \(f: D^n \to D^n\) 有不动点。

\emph{证明}：若 \(f\) 无不动点，则存在收缩映射 \(r: D^n \to S^{n-1}\)。这将诱导 \(H_{n-1}(S^{n-1}) \cong \mathbb{Z}\) 到 \(H_{n-1}(D^n) = 0\) 的恒等同态通过 \(r_* \circ i_*\)，矛盾。∎

\textbf{推论 77.10}（维数不变性）：若 \(\mathbb{R}^n \cong \mathbb{R}^m\)（同胚），则 \(n = m\)。

\emph{证明}：\(\mathbb{R}^n \setminus \{0\} \simeq S^{n-1}\)。若 \(\mathbb{R}^n \cong \mathbb{R}^m\)，则 \(S^{n-1} \simeq S^{m-1}\)，故 \(H_{n-1}(S^{n-1}) \cong H_{n-1}(S^{m-1})\)。仅当 \(n = m\) 时可能。∎

\subsection{77.5 胞腔同调}\label{ux80deux8154ux540cux8c03}

\textbf{定义 77.8}（CW 复形）：CW 复形是归纳构造的空间：\(X^0\) 是离散点集；\(X^n\) 由 \(X^{n-1}\) 粘合 \(n\) 维胞腔 \(D^n_\alpha\) 得到（沿边界 \(S^{n-1}_\alpha \to X^{n-1}\) 粘合）。\(X = \bigcup_n X^n\) 赋予弱拓扑。

\textbf{定理 77.11}（胞腔同调）：对 CW 复形 \(X\)，\(H_n(X^n, X^{n-1})\) 是自由 Abel 群，秩等于 \(X\) 的 \(n\) 维胞腔数。\(H_n(X)\) 同构于胞腔链复形

\[\cdots \to H_n(X^n, X^{n-1}) \xrightarrow{d_n} H_{n-1}(X^{n-1}, X^{n-2}) \to \cdots\]

的同调，其中 \(d_n\) 由边界粘合映射的度决定。

\textbf{证明}：由 excision 和同伦不变性，\(H_k(X^n, X^{n-1}) \cong \bigoplus_{\alpha} H_k(D^n_\alpha, S^{n-1}_\alpha)\)，仅 \(k=n\) 时非零，秩等于 \(n\)-胞腔数。考虑三元组 \((X^{n+1}, X^n, X^{n-1})\) 的长正合序列，提取出复形 \(\cdots \to H_{n+1}(X^{n+1}, X^n) \to H_n(X^n, X^{n-1}) \to H_{n-1}(X^{n-1}, X^{n-2}) \to \cdots\)。\(d_n\) 是复合 \(H_n(X^n, X^{n-1}) \xrightarrow{\partial} H_{n-1}(X^{n-1}) \to H_{n-1}(X^{n-1}, X^{n-2})\)。由胞腔逼近定理，\(H_n(X) \cong H_n(X^{n+1}) \cong \ker d_n / \operatorname{im} d_{n+1}\)。\(\blacksquare\)

\hrulefill

\hrulefill

\hrulefill

\hrulefill

\hrulefill

\hrulefill

\section{第239章：上同调与杯积}\label{ux7b2c239ux7ae0ux4e0aux540cux8c03ux4e0eux676fux79ef}

上同调是同调的对偶理论------将链复形中的链映射 \(C_n(X) \to G\) 而非链本身作为基本对象，得到反变函子 \(H^n(-; G)\)。这一看似简单的对偶化带来了丰富的额外结构，其中最重要的是杯积（cup product）：\(H^p(X) \times H^q(X) \to H^{p+q}(X)\)，使得上同调成为一个分次交换环而非仅仅是一族 Abel 群。杯积的几何意义深刻而优美：Poincare 对偶定理（对于闭定向流形 \(M\)）建立了 \(H^k(M) \cong H_{n-k}(M)\)，在此对偶下，杯积恰好对应了子流形的横截相交------这一对偶将代数结构与几何拓扑融为一体，是同调论与上同调论互补性的最有力见证。杯积的发现（主要由 Alexander、Whitney 和 Kolmogorov 在 1930 年代独立完成）是拓扑学从群论层次跃迁到环论层次的标志性事件。在本章的同调理论构建中，上同调以对偶形式的极大灵活性支持了纤维丛的示性类理论（Chern 类、Stiefel-Whitney 类、Pontryagin 类均以上同调环的元素定义），以及在 de Rham 理论中上同调等价于闭微分形式模恰当形式。

\subsection{78.1 上同调群}\label{ux4e0aux540cux8c03ux7fa4}

\textbf{定义 78.1}（上链复形）：设 \(C_\bullet(X)\) 是奇异链复形。对 Abel 群 \(G\)，定义\textbf{上链群} \(C^n(X; G) = \operatorname{Hom}(C_n(X), G)\)。上边缘算子 \(\delta^n: C^n(X; G) \to C^{n+1}(X; G)\) 定义为 \((\delta^n \varphi)(\sigma) = \varphi(\partial_{n+1} \sigma)\)。

\textbf{定义 78.2}（奇异上同调群）：\(X\) 的系数在 \(G\) 中的 \(n\) 次\textbf{奇异上同调群}为

\[H^n(X; G) = \frac{\ker \delta^n}{\operatorname{im} \delta^{n-1}}\]

\textbf{命题 78.1}（上同调与同调的关系）：由万有系数定理，存在分裂短正合序列：

\[0 \to \operatorname{Ext}(H_{n-1}(X), G) \to H^n(X; G) \to \operatorname{Hom}(H_n(X), G) \to 0\]

特别地，若 \(G = \mathbb{R}\)（或 \(\mathbb{Q}\)），则 \(H^n(X; \mathbb{R}) \cong \operatorname{Hom}(H_n(X), \mathbb{R})\)（对偶空间）。

\subsection{78.2 杯积}\label{ux676fux79ef}

\textbf{定义 78.3}（杯积）：设 \(R\) 是交换环。定义\textbf{杯积}

\[\smile: H^p(X; R) \times H^q(X; R) \to H^{p+q}(X; R)\]

在上链层次上，对 \(\varphi \in C^p(X; R)\)，\(\psi \in C^q(X; R)\)，\((\varphi \smile \psi)(\sigma) = \varphi(\sigma|_{[e_0, \ldots, e_p]}) \cdot \psi(\sigma|_{[e_p, \ldots, e_{p+q}]})\)，其中 \(\sigma|_{[e_0, \ldots, e_p]}\) 表示 \(\sigma\) 限制在前 \(p+1\) 个顶点上。

\textbf{命题 78.2}（杯积的性质）： 1. \(\delta(\varphi \smile \psi) = \delta\varphi \smile \psi + (-1)^p \varphi \smile \delta\psi\)（杯积在上同调层次良定义） 2. \((\varphi \smile \psi) \smile \eta = \varphi \smile (\psi \smile \eta)\)（结合律） 3. \(\varphi \smile \psi = (-1)^{pq} \psi \smile \varphi\)（分次交换律，在同调层次） 4. 杯积是自然的：\(f^*(\varphi \smile \psi) = f^*\varphi \smile f^*\psi\)

\textbf{定义 78.4}（上同调环）：\(X\) 的\textbf{上同调环} \(H^*(X; R) = \bigoplus_{n \geq 0} H^n(X; R)\) 在杯积下构成分次交换 \(R\)-代数。

\subsection{78.3 Poincaré 对偶}\label{poincaruxe9-ux5bf9ux5076}

\textbf{定义 78.5}（定向）：\(n\) 维拓扑流形 \(M\) 的\textbf{定向}是 \(H_n(M, M \setminus \{x\}) \cong \mathbb{Z}\) 在每点 \(x\) 处的一致选择。等价地，在 \(M\) 紧致时，\textbf{基本类} \([M] \in H_n(M)\) 是当 \(M\) 可定向时的生成元。

\textbf{定理 78.3}（Poincaré 对偶）：设 \(M\) 是紧致无边可定向 \(n\) 维流形。则杯积与基本类 \([M]\) 的配对给出同构：

\[H^k(M; \mathbb{Z}) \cong H_{n-k}(M; \mathbb{Z})\]

同构由 \(\alpha \mapsto [M] \smallfrown \alpha\) 给出（\(\smallfrown\) 是 \textbf{cap 积}）。

\textbf{证明}：对闭可定向 \(n\) 维流形 \(M\)，取基本类 \([M] \in H_n(M)\)。定义 cap 积 \(\smallfrown: H^k(M) \times H_n(M) \to H_{n-k}(M)\)，考虑映射 \(D_M: H^k(M) \to H_{n-k}(M)\)，\(\alpha \mapsto [M] \smallfrown \alpha\)。验证方法：(1) 局部验证：对 \(\mathbb{R}^n\) 中开集，切除定理给出 \(H^k_c(\mathbb{R}^n) \cong H_{n-k}(\mathbb{R}^n)\)。 (2) Mayer-Vietoris 粘合：若 \(M = U \cup V\) 且 \(D_U, D_V, D_{U\cap V}\) 均为同构，则由五引理知 \(D_M\) 也为同构。 (3) 归纳法：将 \(M\) 分解为有限个坐标邻域的并，由 (1)(2) 归纳得 \(D_M\) 为同构。cap 积的自然性保证该映射不依赖于剖分选择。\(\blacksquare\)

\textbf{推论 78.4}：紧致可定向流形的 Euler 示性数为 \(\chi(M) = \sum_{k=0}^n (-1)^k \operatorname{rank} H_k(M)\)。由 Poincaré 对偶，\(\operatorname{rank} H_k(M) = \operatorname{rank} H_{n-k}(M)\)。

\textbf{定理 78.5}（Poincaré 对偶的杯积形式）：配对

\[H^k(M; \mathbb{R}) \times H^{n-k}(M; \mathbb{R}) \to \mathbb{R}, \quad (\alpha, \beta) \mapsto \langle \alpha \smile \beta, [M] \rangle\]

是非退化的双线性型。当 \(n = 4k\) 时，\(H^{2k}\) 上的二次型（符号差）是重要的拓扑不变量。

\textbf{证明}：由 Poincaré 对偶，cap 积给出同构 \(D: H^k(M) \to H_{n-k}(M)\)，\(\alpha \mapsto [M] \smallfrown \alpha\)。杯积与 cap 积的关系为 \(\langle \alpha \smile \beta, [M] \rangle = \langle \alpha, [M] \smallfrown \beta \rangle\)（配对求值）。若 \(\alpha \neq 0\)，由于 \(D\) 是同构，存在 \(\gamma \in H^{n-k}(M)\) 使 \(\langle \alpha \smile \gamma, [M] \rangle \neq 0\)（否则 \(\alpha\) 与所有类配对为零，与 \(D\) 的非退化性矛盾）。故配对非退化。当 \(n=4k\)，\(H^{2k}\) 上的对称双线性型 \(\langle \alpha \smile \beta, [M] \rangle\) 的符号差（正负特征值个数差）是流形的微分同胚不变量，即 Hirzebruch 符号差定理的核心对象。\(\blacksquare\)

\subsection{78.4 de Rham 上同调回顾}\label{de-rham-ux4e0aux540cux8c03ux56deux987e}

由 V11 第 64 章的 de Rham 定理（陈述）：对光滑流形 \(M\)，de Rham 上同调 \(H^k_{dR}(M)\) 同构于奇异上同调 \(H^k(M; \mathbb{R})\)。杯积对应于微分形式的外积。

\hrulefill

\hrulefill

\hrulefill

\hrulefill

\hrulefill

\hrulefill

\hrulefill

\hrulefill

\section{第240章：奇异同调的Eilenberg-Steenrod公理}\label{ux7b2c240ux7ae0ux5947ux5f02ux540cux8c03ux7684eilenberg-steenrodux516cux7406}

Eilenberg-Steenrod 公理体系（1945年发表于《Foundations of Algebraic Topology》）是代数拓扑发展史上的一座里程碑。在 Eilenberg 和 Steenrod 之前，同调论有多个互不等价的版本：Vietoris 同调、Cech 同调、奇异同调、单纯同调等，它们在不同的空间类上给出不同的结果，但没有统一的公理基础。Eilenberg 和 Steenrod 的洞见是：同调论的本质特征不在于具体的链复形构造，而在于一组公理------同伦不变性、切除公理、维数公理、正合性公理和可加性公理。任何满足这组公理的函子（在可三角剖分空间对上取值于 Abel 群）必然与奇异同调自然同构（在有限 CW 复形范畴上）。这一公理化方法不仅统一了已有的同调理论，更为后续的广义同调论（K-理论、配边理论、稳定同伦论）提供了概念框架------广义同调论满足除维数公理外的所有 Eilenberg-Steenrod 公理（维数公理 \(H_n(\text{pt}) = 0\)，\(n \neq 0\) 被替换为特定的系数群）。公理体系的优美之处在于它以函子性语言完全刻画了同调，使得具体的链复形构造退居幕后，而拓扑推理只需调用公理即可。在现代观点下，Eilenberg-Steenrod 公理等价于说同调论 \(H_*\) 是一个从有限 CW 空间对范畴到分次 Abel 群范畴的函子，满足某种导出函子的万有性质------这联系了公理方法与同调代数的导出函子视角。

\subsection{78.x Eilenberg-Steenrod公理的陈述}\label{x-eilenberg-steenrodux516cux7406ux7684ux9648ux8ff0}

\textbf{定义 78.1}（Eilenberg-Steenrod同调论）：一个（经典）\textbf{同调论}由以下数据构成：对每个整数 \(n \in \mathbb{Z}\) 和每个拓扑空间对 \((X, A)\)（\(A \subseteq X\)），指定一个 Abel 群 \(H_n(X, A)\)；对每个连续映射 \(f: (X, A) \to (Y, B)\)，指定同态 \(f_*: H_n(X, A) \to H_n(Y, B)\)；并对每个 \((X, A)\) 和 \(n\)，指定连接同态 \(\partial: H_n(X, A) \to H_{n-1}(A)\)（其中 \(H_{n-1}(A) := H_{n-1}(A, \varnothing)\)）。这些数据需满足以下五条公理：

\textbf{公理 1（恒等与函子性）}：\((\operatorname{id}_{(X,A)})_* = \operatorname{id}\)；\((g \circ f)_* = g_* \circ f_*\)。连接同态是自然的：对 \(f: (X,A) \to (Y,B)\)，\(\partial \circ f_* = (f|_A)_* \circ \partial\)。

\textbf{公理 2（正合性）}：三元组 \((X, A)\) 的长序列 \[\cdots \to H_n(A) \xrightarrow{i_*} H_n(X) \xrightarrow{j_*} H_n(X, A) \xrightarrow{\partial} H_{n-1}(A) \to \cdots\] 是正合的，其中 \(i: A \hookrightarrow X\) 和 \(j: (X, \varnothing) \hookrightarrow (X, A)\) 为包含映射。

\textbf{公理 3（切除公理）}：若 \(U \subseteq A \subseteq X\) 满足 \(\overline{U} \subseteq \operatorname{int}(A)\)，则包含映射 \((X \setminus U, A \setminus U) \hookrightarrow (X, A)\) 诱导同构 \(H_n(X \setminus U, A \setminus U) \cong H_n(X, A)\)（对所有 \(n\)）。

\textbf{公理 4（同伦不变性）}：若 \(f, g: (X, A) \to (Y, B)\) 同伦（通过保持 \(A\) 到 \(B\) 的映射），则 \(f_* = g_*: H_n(X, A) \to H_n(Y, B)\)。

\textbf{公理 5（维数公理）}：对单点空间 \(\text{pt}\)，\(H_n(\text{pt}) = 0\)（\(n \neq 0\)），而 \(H_0(\text{pt}) \cong \mathbb{Z}\)（称为系数群）。

若还满足： \textbf{公理 6（可加性）}：对不交并 \(X = \bigsqcup_{\alpha} X_\alpha\)，包含映射诱导同构 \(\bigoplus_{\alpha} H_n(X_\alpha) \cong H_n(X)\)。

则称为\textbf{可加同调论}。

\textbf{定理 78.1}（Eilenberg-Steenrod唯一性定理）：设 \(H_*\) 和 \(H'_*\) 是两个可加同调论，\(\varphi: H_0(\text{pt}) \to H'_0(\text{pt})\) 是系数群的同构。则对任意可三角剖分的空间对（特别地，任意有限 CW 复形对），存在唯一的自然同构 \(\Phi: H_n(X, A) \cong H'_n(X, A)\)（对所有 \(n\)）扩展 \(\varphi\)。换言之，公理完全决定了有限 CW 复形上的同调论。

\textbf{证明}（框架）：使用公理在 CW 滤过上的归纳法。步骤 1：由维数公理，\(\text{pt}\) 的同调群被确定。步骤 2：证明公理蕴含 \(H_n(S^k)\) 的计算与奇异同调一致（归纳使用 Mayer-Vietoris，后者可由公理导出）。步骤 3：对任意有限 CW 复形 \(X\)，按骨架 \(X^k\) 归纳。\(H_n(X^k, X^{k-1})\) 由切除公理化为球面同调的直和（秩等于 \(k\)-胞腔数）。胞腔链复形由复合 \(H_n(X^n, X^{n-1}) \xrightarrow{\partial} H_{n-1}(X^{n-1}) \to H_{n-1}(X^{n-1}, X^{n-2})\) 定义，该复合完全由骨架的粘合映射决定。步骤 4：\(\Phi\) 在骨架 \(X^k\) 上归纳构造，利用长正合序列和五项引理证明每一步的 \(\Phi\) 为同构。由于两个同调论在球面和胞腔对上由公理强制为相同的群和同态（至多相差系数群同构 \(\varphi\)），归纳完成。对于无限 CW 复形，利用可加性公理取归纳极限。\(\blacksquare\)

\subsection{78.y 广义同调论与谱}\label{y-ux5e7fux4e49ux540cux8c03ux8bbaux4e0eux8c31}

\textbf{定义 78.2}（广义同调论 / Generalized Homology Theory）：一个\textbf{广义同调论}是满足 Eilenberg-Steenrod 公理 1-4（函子性、正合性、切除、同伦不变性）和可加性，但不一定满足维数公理的函子 \(E_n: \mathbf{Top}^2 \to \mathbf{Ab}\)（\(n \in \mathbb{Z}\)）。\(E_n\) 通常要求定义在带基点空间范畴上，\(E_n(\text{pt})\) 为广义同调论的\textbf{系数群}（一般非零，对所有 \(n \in \mathbb{Z}\)）。重要的例子包括：

\begin{itemize}
\tightlist
\item
  \textbf{稳定同伦论} \(\pi_n^S(X) = \lim_{k \to \infty} \pi_{n+k}(\Sigma^k X)\)，其系数为球面的稳定同伦群。
\item
  \textbf{拓扑 K-理论} \(K_n(X)\)（\(n \in \mathbb{Z}\)，Bott 周期性给出 \(K_n \cong K_{n+2}\)）。
\item
  \textbf{配边理论}（Cobordism）：\(MO_*(X)\)（非定向配边）、\(MU_*(X)\)（复配边）、\(MSU_*(X)\) 等。
\item
  \textbf{Brown-Peterson 同调} \(BP_*(X)\)（在素数 \(p\) 处局部化）。
\end{itemize}

\textbf{定理 78.2}（Brown 可表性定理，1962）：任意满足可加性公理的广义同调论 \(E_n(-)\)（定义在带基点连通 CW 复形范畴上）由某个谱（spectrum）\(\mathbf{E} = \{E_n, \sigma_n: \Sigma E_n \to E_{n+1}\}\) 表示，即存在自然同构

\[E_n(X) \cong \lim_{k \to \infty} [S^{n+k}, E_k \wedge X]_*\]

等价地，\(E_n(X) \cong \pi_n(\mathbf{E} \wedge X)\)（\(\mathbf{E} \wedge X\) 的稳定同伦群）。反过来，每个谱定义一个广义同调论。

\textbf{证明}（概要）：Brown 可表性定理的原始陈述涉及上同调论的对偶版本：设 \(F: \mathbf{CW}_*^{\text{op}} \to \mathbf{Set}_*\) 是满足 Mayer-Vietoris 公理和 wedge 公理的函子（即 \(F\) 将余积变为积），则 \(F\) 可表：存在 \(\Omega\)-谱 \(\mathbf{E}\) 使 \(F(X) \cong [X, E_n]_*\)。对上同调论 \(E^n(-)\) 应用此定理得谱表示。同调论版本通过对偶（Spanier-Whitehead 对偶和稳定同伦范畴的闭结构）导出。谱 \(\mathbf{E}\) 的第 \(n\) 个空间 \(E_n\) 满足 \(\Omega E_{n+1} \simeq E_n\)（即 \(\mathbf{E}\) 为 \(\Omega\)-谱）。\(\blacksquare\)

\textbf{推论 78.3}（广义同调论与谱的一一对应）：稳定同伦范畴（同伦范畴 of spectra）的对象与广义同调论（满足可加性的）之间存在范畴等价。具体地，谱的同伦范畴 \(\mathcal{SH}\) 的每个对象 \(\mathbf{E}\) 通过 \(E_n(X) = \pi_n(\mathbf{E} \wedge \Sigma^\infty X)\) 定义一个广义同调论；反之，Brown 可表性定理为每个广义同调论给出一个谱。

\textbf{证明}：谱范畴 \(\mathcal{SH}\) 中，楔和公理和 Mayer-Vietoris 自动满足（由稳定同伦范畴的三角结构）。函子 \(X \mapsto \mathbf{E} \wedge \Sigma^\infty X\) 满足 Brown 可表性定理所需的所有公理（\(X \mapsto [\Sigma^\infty X, \mathbf{E}]_n\) 在同调论一侧）。范畴等价的具体形式涉及 Bousfield 局部化和 Adams 谱序列。\(\blacksquare\)

\subsection{78.z Eilenberg-Steenrod公理在导出范畴中的推广}\label{z-eilenberg-steenrodux516cux7406ux5728ux5bfcux51faux8303ux7574ux4e2dux7684ux63a8ux5e7f}

\textbf{定理 78.4}（公理的导出函子诠释）：设 \(\mathcal{A}\) 为 Abel 范畴，\(F: \mathbf{Top} \to \mathbf{Ch}(\mathcal{A})\) 为函子（将空间映为链复形）。若 \(F\) 满足：（i）同伦不变性：\(f \simeq g \implies F(f) \simeq F(g)\)（链同伦）；（ii）切除：若 \(\overline{U} \subseteq \operatorname{int}(A)\)，则 \(F(X \setminus U)/F(A \setminus U) \to F(X)/F(A)\) 是拟同构；（iii）正规化：\(F(\text{pt})\) 集中在零次（即幺模复形）。则 \(H_n(F(-))\) 是满足 Eilenberg-Steenrod 公理（除可加性）的同调论。反之，任何这样的同调论在导出范畴 \(D(\mathcal{A})\) 中由某个唯一的函子 \(F\)（在拟同构意义下）给出。这揭示了 Eilenberg-Steenrod 公理与同调代数中导出函子万有性质之间的深层联系。

\textbf{证明}：由条件（i），\(H_n \circ F\) 满足同伦公理。由条件（ii）和长正合同调序列（链复形的短正合序列诱导），满足切除公理。条件（iii）给出维数公理。函子性与正合性由链复形同调的基本性质自动满足。唯一性部分：给定两个满足上述条件的函子 \(F, G\)，它们对 \(\text{pt}\) 的取值为拟同构的链复形（因 \(H_0(F(\text{pt})) \cong H_0(G(\text{pt}))\)）。按 CW 骨架归纳，利用切除性（在导出范畴中化为短正合列的分裂）和五项引理（在导出范畴中化为三角的旋转），构造自然变换 \(F \Rightarrow G\) 在每个骨架上为拟同构。在导出范畴 \(D(\mathcal{A})\) 中，此自然变换即为同构。此构造与 Verdier 的导出函子理论和 Quillen 的模型范畴框架相容。\(\blacksquare\)

\hrulefill

\hrulefill

\hrulefill

\section{第241章：上同调的乘积结构与Poincaré对偶}\label{ux7b2c241ux7ae0ux4e0aux540cux8c03ux7684ux4e58ux79efux7ed3ux6784ux4e0epoincaruxe9ux5bf9ux5076}

上同调论相比同调论的核心优势在于其上的乘积结构------上积（cup product）和卡积（cap product）。上积将上同调转化为分次交换环：\(\smile: H^p(X) \times H^q(X) \to H^{p+q}(X)\)，这赋予了拓扑空间以丰富的代数不变量（上同调环），其区分能力远超单独的 Betti 数。Poincaré 对偶定理则是紧致可定向流形理论中最深刻的定理之一：对于 \(n\) 维紧致可定向流形 \(M\)，上同调与同调之间存在典范同构 \(H^k(M) \cong H_{n-k}(M)\)。这种对偶性表明上同调与同调''互补''------\(k\) 维上同调类与 \(n-k\) 维同调类配对。在乘积结构下，Poincaré 对偶给出了上同调环的非退化配对 \(H^k(M) \times H^{n-k}(M) \to \mathbb{Z}\)，使得上同调环成为 Poincaré 对偶代数。这种结构的几何含义深远------它编码了子流形的相交理论（两个子流形的相交数等于各自 Poincaré 对偶类的上积在基本类上的求值）。对于非可定向流形，相应的论断需在 \(\mathbb{Z}_2\) 系数下成立。de Rham 上同调版本（对光滑流形）提供了最直观的构造：微分形式的外积对应于上积，而积分 \(\int_M \omega \wedge \eta\) 给出了 Poincaré 对偶在分析学意义上的实现。

\subsection{79.1 上积与卡积的代数构造}\label{ux4e0aux79efux4e0eux5361ux79efux7684ux4ee3ux6570ux6784ux9020}

\textbf{定义 79.1}（上积 / Cup Product）：设 \(X\) 为拓扑空间，\(R\) 为交换环。奇异上链的\textbf{上积}

\[\smile: C^p(X; R) \times C^q(X; R) \to C^{p+q}(X; R)\]

定义为：对 \(\varphi \in C^p\)，\(\psi \in C^q\) 和奇异单形 \(\sigma: \Delta^{p+q} \to X\)，

\[(\varphi \smile \psi)(\sigma) = \varphi(\sigma \circ \iota_{[0,\ldots,p]}) \cdot \psi(\sigma \circ \iota_{[p,\ldots,p+q]})\]

其中 \(\iota_{[0,\ldots,p]}: \Delta^p \to \Delta^{p+q}\) 为到前 \(p+1\) 个顶点的仿射包含（前脸），\(\iota_{[p,\ldots,p+q]}: \Delta^q \to \Delta^{p+q}\) 为到后 \(q+1\) 个顶点的仿射包含（后脸）。上积诱导上同调层次的乘积（仍记为 \(\smile\)），使 \(H^*(X; R) = \bigoplus_{k \geq 0} H^k(X; R)\) 成为分次交换环：\(\alpha \smile \beta = (-1)^{|\alpha||\beta|} \beta \smile \alpha\)。

\textbf{定理 79.1}（上积的Leibniz法则）：上边缘算子 \(\delta: C^k \to C^{k+1}\) 对乘积满足

\[\delta(\varphi \smile \psi) = (\delta \varphi) \smile \psi + (-1)^{|\varphi|} \varphi \smile (\delta \psi)\]

即 \(\smile\) 是链复形层次的导出代数结构。这保证 \(\smile\) 在取上同调后仍是良好定义的（闭链的乘积为闭链，恰当链与任意闭链的乘积为恰当的）。

\textbf{证明}：直接计算。对 \(\sigma: \Delta^{p+q+1} \to X\)： \[(\delta(\varphi \smile \psi))(\sigma) = \sum_{i=0}^{p+q+1} (-1)^i (\varphi \smile \psi)(\sigma \circ \iota_{[0,\ldots,\hat{i},\ldots,p+q+1]})\]

当 \(i \leq p\) 时，\(\sigma \circ \iota_{[0,\ldots,\hat{i},\ldots,p+q+1]}\) 到 \(X\)，其前脸 \(\Delta^p \to \Delta^{p+q} \to \Delta^{p+q+1}\)（去掉第 \(i\) 个顶点）等于 \(\sigma \circ \iota_{[0,\ldots,p+1]}\) 去掉第 \(i\) 个顶点的前脸；后脸不变（从顶点 \(p+1\) 到 \(p+q+1\)）。当 \(i > p\) 时类似处理。分别对 \(i \leq p\) 和 \(i > p\) 求和，得到 \((\delta \varphi \smile \psi)\) 和 \((-1)^p (\varphi \smile \delta \psi)\) 的项，其间需仔细处理记号偏移。符号因子 \((-1)^p\) 来自 \(i > p\) 项中前脸（包含前 \(p+1\) 个顶点）的定义不变、后脸跳过第 \(i\) 个顶点时所引起的符号 \((-1)^{i-p-1}\)，总和后出现 \((-1)^p\)。详细见 Hatcher (2002) 第 3.2 节。\(\blacksquare\)

\subsection{79.2 Poincaré对偶定理}\label{poincaruxe9ux5bf9ux5076ux5b9aux7406}

\textbf{定义 79.2}（基本类 / Fundamental Class）：设 \(M\) 为紧致可定向 \(n\) 维流形（无边）。\(M\) 的\textbf{基本类} \([M] \in H_n(M; \mathbb{Z})\) 是生成 \(H_n(M, M \setminus \{x\}; \mathbb{Z}) \cong \mathbb{Z}\) 的同调类在局部到全局的相容粘合下定义的全局同调类。对三角剖分或光滑胞腔剖分，\([M]\) 可表示为各 \(n\) 维单形之和（以与定向一致的定向求和）。

\textbf{定理 79.2}（Poincaré对偶定理，Poincaré 1895）：设 \(M\) 为紧致可定向 \(n\) 维流形（无边）。则\textbf{卡积}（cap product）与基本类 \([M]\) 定义了同构

\[D: H^k(M; \mathbb{Z}) \xrightarrow{\cong} H_{n-k}(M; \mathbb{Z}), \quad \alpha \mapsto [M] \frown \alpha\]

其中卡积 \(\frown: H_n(M) \times H^k(M) \to H_{n-k}(M)\)（在链层次定义为 \(( \sigma \frown \varphi ) = \varphi(\sigma_{\text{front}}) \cdot \sigma_{\text{back}}\)）。等价地，配对

\[\langle \cdot, \cdot \rangle: H^k(M) \times H^{n-k}(M) \to \mathbb{Z}, \quad \langle \alpha, \beta \rangle = (\alpha \smile \beta)([M])\]

是非退化的（在模去挠子群后），即 \(H^k(M; \mathbb{R}) \cong H_{n-k}(M; \mathbb{R})^*\)。

\textbf{证明}（框架 --- 使用切除 + Mayer-Vietoris 的归纳法）：

步骤 1：对 \(M = \mathbb{R}^n\) 验证（非紧致版本，紧支上同调 \(H_c^k(\mathbb{R}^n)\) 与 \(H_{n-k}(\mathbb{R}^n)\) 同构）。此处基本类是无限循环，卡积给出 \(H_c^n(\mathbb{R}^n) \to H_0(\mathbb{R}^n)\) 的同构。

步骤 2：若定理对开集 \(U, V, U \cap V\) 成立，则对 \(U \cup V\) 成立（Mayer-Vietoris 序列的自然性和五项引理）。

步骤 3：\(M\) 可被有限个同胚于 \(\mathbb{R}^n\) 的坐标邻域覆盖。按覆盖中开集个数归纳，使用步骤 1 为基始、步骤 2 为归纳步骤，逐级粘合坐标邻域直至覆盖整个流形。对无边流形，有限覆盖完备性由紧致性保证。归纳中需验证卡积与 Mayer-Vietoris 连接同态的自然性（涉及余边界同态与边界同态的交换图）。\(\blacksquare\)

\textbf{定理 79.3}（Poincaré对偶的上积配对形式）：在上述定理的假设下，双线性型

\[Q_M: H^k(M; \mathbb{R}) \times H^{n-k}(M; \mathbb{R}) \to \mathbb{R}, \quad Q_M(\alpha, \beta) = \int_M \alpha \smile \beta\]

是非退化的。当 \(n = 4m\) 时，配对 \(H^{2m}(M) \times H^{2m}(M) \to \mathbb{R}\) 是对称的（因 \((-1)^{(2m)^2} = +1\)），其符号差（正负特征值个数之差）定义为流形的\textbf{符号差} \(\sigma(M)\)。Hirzebruch 符号差定理：\(\sigma(M) = \langle L(M), [M] \rangle\)，其中 \(L(M)\) 为 \(M\) 的 \(L\)-亏格（由 Pontryagin 类的多项式表达）。

\textbf{证明}（符号差的定义与性质）：对称双线性型 \(Q_M\) 的矩阵是对称的。由 Sylvester 惯性定律，正特征值个数 \(b_{2m}^+\) 与负特征值个数 \(b_{2m}^-\) 是定向同伦不变量（在同痕和定向同伦下不变，但不是拓扑不变量------有同胚而非微分同胚的 \(4\)-流形具有不同符号差）。\(\sigma(M) = b_{2m}^+ - b_{2m}^-\)。配对非退化的证明依赖于上积与 Poincaré 对偶的相容性：对任意 \(\alpha \in H^k, \beta \in H^{n-k}\)，卡积性质给出 \(([M] \frown \alpha) \frown \beta = [M] \frown (\alpha \smile \beta)\)，而 \([M] \frown (\alpha \smile \beta)\) 在 \(H_0(M) \cong \mathbb{Z}\) 中的系数恰为 \((\alpha \smile \beta)([M])\)。若某 \(\alpha \neq 0\) 使对所有 \(\beta\) 有 \(\langle \alpha, \beta \rangle = 0\)，则由 Poincaré 对偶同构 \(D(\alpha) \in H_{n-k}(M)\) 非零，存在 \(\beta\) 使 \(\beta(D(\alpha)) = \langle \beta, D(\alpha) \rangle \neq 0\)（同调与上同调的 Kronecker 配对非退化），矛盾（因 \(\beta(D(\alpha)) = \langle \alpha, \beta \rangle\)）。\(\blacksquare\)

\subsection{79.3 上同调环计算实例与Lefschetz对偶}\label{ux4e0aux540cux8c03ux73afux8ba1ux7b97ux5b9eux4f8bux4e0elefschetzux5bf9ux5076}

\textbf{定理 79.4}（实射影空间的上同调环）：对实射影空间 \(\mathbb{RP}^n\)，其上同调环（系数 \(\mathbb{Z}_2\)）为

\[H^*(\mathbb{RP}^n; \mathbb{Z}_2) \cong \mathbb{Z}_2[\alpha] / (\alpha^{n+1}), \quad |\alpha| = 1\]

即每个维数恰有一个生成元，且 \(\alpha^k\) 为 \(H^k\) 的生成元（\(0 \leq k \leq n\)）。整数系数下，\(H^k(\mathbb{RP}^n; \mathbb{Z}) \cong \mathbb{Z}_2\)（当 \(k\) 为偶数且 \(2 \leq k \leq n\)），\(\mathbb{Z}\)（当 \(k = 0\) 或 \(k=n\) 为奇数），\(0\)（其他情形）。

\textbf{证明}（\(\mathbb{Z}_2\) 系数）：使用 CW 复形结构 \(\mathbb{RP}^n = e^0 \cup e^1 \cup \cdots \cup e^n\)，每维一个胞腔。胞腔链复形 \(C_k = \mathbb{Z}_2\)（生成元 \(e_k\)），边界同态 \(d(e_k) = (1 + (-1)^k) e_{k-1}\)（因粘合映射的映射度为 \(1 + (-1)^k\)（模 2 下为 \(0\) 当 \(k\) 偶，为 \(0\) 当 \(k\) 奇 --- 实际上边界同态在 \(\mathbb{Z}_2\) 下全部为零！因 \(1+(-1)^k \equiv 0 \pmod{2}\) 对所有 \(k\)）。故所有上同调群 \(H^k(\mathbb{RP}^n; \mathbb{Z}_2) \cong \mathbb{Z}_2\)（\(0 \leq k \leq n\)）。上积的计算：生成元 \(\alpha \in H^1\) 是 \(\mathbb{RP}^n \subset \mathbb{RP}^\infty\) 的超平面类（第一 Stiefel-Whitney 类）的拉回。\(\alpha^k\)（\(k\) 次上积）为 \(H^k\) 的生成元。由维数原因，\(\alpha^{n+1} = 0\)，故上同调环如所述。\(\blacksquare\)

\textbf{定理 79.5}（Lefschetz对偶与超平面截面定理）：若 \(M\) 为紧致复 \(n\) 维 Kahler 流形（特别地，光滑复射影代数簇），其上同调环带有 Lefschetz 算子 \(L: H^k(M) \to H^{k+2}(M)\)（由 Kahler 形式 \(\omega\) 的上积定义：\(L(\alpha) = \omega \smile \alpha\)）。\textbf{硬 Lefschetz 定理}断言：\(L^{n-k}: H^k(M) \to H^{2n-k}(M)\) 是同构（\(0 \leq k \leq n\)）。此定理将 Poincaré 对偶精细化为 Lefschetz 分解：\(H^k(M) = \bigoplus_{j \geq 0} L^j P^{k-2j}\)，其中 \(P^m = \ker(L^{n-m+1}: H^m \to H^{2n-m+2})\) 为本原上同调。该分解是 Hodge 理论中调和形式分解的拓扑对应，也是 Deligne 混合 Hodge 理论的出发点。

\textbf{证明}（框架 --- 复代数簇情形使用 \(\mathfrak{sl}_2\) 表示论）：Kahler 形式 \(\omega\) 是闭 \((1,1)\)-形式。定义 \(\Lambda: H^k \to H^{k-2}\) 为 \(L\) 的伴随算子（关于 Hodge 内积或 Poincaré 配对）。算子 \(L, \Lambda, H\)（\(H = [L, \Lambda] = \sum (k-n) \Pi_k\)，\(\Pi_k\) 为到 \(H^k\) 的投影）构成 \(\mathfrak{sl}_2(\mathbb{C})\) 的表示。由有限维 \(\mathfrak{sl}_2\) 表示的分类，硬 Lefschetz 定理等价于 \(H\) 在此表示中可对角化（这正是 Hodge 分解所保证的）。对于代数簇，采用 Deligne 的 Weil 猜想的证明方法（\(\ell\)-adic 上同调 + Lefschetz 铅笔），利用超平面截面定理的拓扑版本。\(\blacksquare\)

\hrulefill

\section{第242章：同伦群}\label{ux7b2c242ux7ae0ux540cux4f26ux7fa4}

同伦群 \(\pi_n(X)\)（\(n \geq 1\)）是基本群 \(\pi_1(X)\) 的高维推广，捕捉了 \(n\) 维球面到 \(X\) 的映射的同伦类。

\subsection{79.1 高维同伦群的定义}\label{ux9ad8ux7ef4ux540cux4f26ux7fa4ux7684ux5b9aux4e49}

\textbf{定义 79.1}（\(\pi_n\) 的定义）：设 \(X\) 为带基点 \((X, x_0)\) 的拓扑空间。\(\pi_n(X, x_0)\) 定义为保持基点的映射 \(f: (S^n, s_0) \to (X, x_0)\) 的保基点同伦类构成的集合。等价地，\(f: (I^n, \partial I^n) \to (X, x_0)\) 的同伦类。

\textbf{定义 79.2}（群运算）：\(\pi_n(X, x_0)\) 上的群运算由「坐标压缩」定义：对 \(f, g: (I^n, \partial I^n) \to (X, x_0)\)，

\[(f + g)(t_1, \ldots, t_n) = \begin{cases} f(2t_1, t_2, \ldots, t_n) & 0 \leq t_1 \leq 1/2 \\ g(2t_1-1, t_2, \ldots, t_n) & 1/2 \leq t_1 \leq 1 \end{cases}\]

\textbf{命题 79.1}：对 \(n \geq 2\)，\(\pi_n(X, x_0)\) 是 Abel 群。

\emph{证明}：对 \(n \geq 2\)，存在「旋转」同伦，将 \(f\) 和 \(g\) 的运算顺序交换。直观上，\(n \geq 2\) 时 \(I^n\) 的边界 \(\partial I^n\) 中 \(f\) 和 \(g\) 被挤压到不同的坐标方向，有足够空间交换。∎

\textbf{命题 79.2}（函子性）：\(\pi_n\) 是从带基点拓扑空间范畴到群范畴（\(n \geq 2\) 时为 Abel 群范畴）的函子。同伦等价诱导同伦群的同构。

\textbf{证明}：对基点到基点的连续映射 \(\phi: (X, x_0) \to (Y, y_0)\)，定义 \(\phi_*: \pi_n(X, x_0) \to \pi_n(Y, y_0)\) 为 \([f] \mapsto [\phi \circ f]\)。此定义良定：若 \(f \simeq g\) 相对于基点，则 \(\phi \circ f \simeq \phi \circ g\)。\(\phi_*\) 是群同态：\(\phi_*([f][g]) = \phi_*([f \cdot g]) = [\phi \circ (f \cdot g)] = [(\phi \circ f) \cdot (\phi \circ g)] = \phi_*([f]) \phi_*([g])\)。\((\operatorname{id}_X)_* = \operatorname{id}_{\pi_n(X)}\) 且 \((\psi \circ \phi)_* = \psi_* \circ \phi_*\) 显然成立。同伦等价 \(h: X \to Y\) 诱导同构：\(h_* \circ (h^{-1})_* = (h \circ h^{-1})_* = (\operatorname{id}_Y)_* = \operatorname{id}\)，同理反向。\(\blacksquare\)

\textbf{命题 79.3}（\(\pi_n\) 与覆叠空间）：若 \(p: \tilde{X} \to X\) 是覆叠映射，则对 \(n \geq 2\)，\(p_*: \pi_n(\tilde{X}) \to \pi_n(X)\) 是同构。即覆叠空间不改变高维同伦群。

\emph{证明}：\(S^n\)（\(n \geq 2\)）是单连通的，故任意映射 \(S^n \to X\) 可唯一提升到 \(\tilde{X}\)（由提升判别法，定理 61.4）。∎

\textbf{推论 79.5}：\(\pi_n(S^1) = 0\) 对 \(n \geq 2\)（因为 \(\mathbb{R}\) 可缩，\(\pi_n(\mathbb{R}) = 0\)）。

\textbf{证明}：\(\mathbb{R} \to S^1\)（\(t \mapsto e^{2\pi i t}\)）是万有覆叠映射。由命题 79.3，对 \(n \geq 2\)，\(p_*: \pi_n(\mathbb{R}) \to \pi_n(S^1)\) 是同构。但 \(\mathbb{R}\) 可缩，故 \(\pi_n(\mathbb{R}) = 0\) 对所有 \(n \geq 1\) 成立。从而 \(\pi_n(S^1) = 0\) 对 \(n \geq 2\)。\(\blacksquare\)

\subsection{79.2 球面的同伦群}\label{ux7403ux9762ux7684ux540cux4f26ux7fa4}

球面的同伦群 \(\pi_k(S^n)\) 是代数拓扑中最基本也最困难的计算对象之一。虽然 \(\pi_n(S^n) \cong \mathbb{Z}\) 由 Hurewicz 定理容易得出，但 \(k > n\) 情形的 \(\pi_k(S^n)\) 的计算极其复杂，至今仍是代数拓扑中活跃的研究领域。关键事实包括：（1）\(\pi_k(S^n) = 0\) 对 \(k < n\)（由 Freudenthal 悬浮定理）；（2）\(\pi_{n+1}(S^n) \cong \mathbb{Z}/2\mathbb{Z}\) 对 \(n \geq 3\)（由 Hopf 纤维化 \(S^3 \to S^2\) 的 Hopf 不变量计算）；（3）\(\pi_{n+2}(S^n) \cong \mathbb{Z}/2\mathbb{Z}\) 对 \(n \geq 2\)；（4）稳定同伦群 \(\pi_{n+k}(S^n)\) 对 \(n \gg k\) 稳定，这些稳定同伦群 \(\pi_k^S\) 是代数拓扑中最核心的不变量之一。球面同伦群的计算需要组合使用 Serre 谱序列、EHP 序列、Adams 谱序列和稳定同伦论的各种工具。例如，\(\pi_4(S^3) \cong \mathbb{Z}/2\mathbb{Z}\) 可使用 Serre 谱序列（定理 80.2）计算（见定理 80.5 的证明），而 \(\pi_6(S^3) \cong \mathbb{Z}/12\mathbb{Z}\) 的计算则需要 Adams 谱序列或 Toda 括号等更高级的工具。球面同伦群中的挠元（torsion elements）特别丰富，且其 \(p\)-局部化（对不同素数 \(p\)）的谱序列结构揭示了与数论中 Bernoulli 数的深刻联系（通过 \(J\)-同态和 Adams 猜想）。球面同伦群的计算问题也是导出代数几何和色展同伦论（chromatic homotopy theory）的核心动机之一------后者通过 Morava K 理论和形式群的高度来组织稳定同伦范畴的信息。

\textbf{定理 79.6}（\(\pi_n(S^n)\) 的计算）：\(\pi_n(S^n) \cong \mathbb{Z}\)。生成元是恒等映射 \([S^n \to S^n]\) 的同伦类。 \textbf{证明}：由 Hurewicz 定理（定理 79.10），\(\pi_n(S^n) \cong H_n(S^n) \cong \mathbb{Z}\)，因为 \(S^n\) 的 \(0 < k < n\) 维同伦群均为零------由 Freudenthal 悬浮定理，\(S^n = \Sigma S^{n-1}\) 的连通性递推得 \(\pi_k(S^n)=0 (k<n)\)。Hurewicz 同态 \(h: \pi_n(S^n) \to H_n(S^n)\) 将映射 \(f: S^n \to S^n\) 映为 \(h([f]) = f_*[S^n] \in H_n(S^n)\)。恒等映射 \([S^n] \in H_n(S^n)\) 是 \(H_n(S^n) \cong \mathbb{Z}\) 的生成元，其逆像 \([\operatorname{id}_{S^n}]\) 即 \(\pi_n(S^n)\) 的生成元，对应 Brouwer 度数为 \(1\) 的单位元素。Hurewicz 定理断言此同态为同构，故 \(\pi_n(S^n) \cong \mathbb{Z}\)。\(\blacksquare\) \#\#\# 79.3 纤维化与同伦长正合序列

\textbf{定义 79.3}（纤维化 / Hurewicz 纤维化）：映射 \(p: E \to B\) 称为\textbf{纤维化}（或 Hurewicz 纤维化），如果对任意空间 \(X\) 和同伦 \(H: X \times I \to B\)，以及初始提升 \(\tilde{H}_0: X \to E\)（\(p \circ \tilde{H}_0 = H_0\)），存在同伦 \(\tilde{H}: X \times I \to E\) 提升 \(H\) 且延伸 \(\tilde{H}_0\)。即下列图表存在提升：

\[
\begin{CD}
X @>{\tilde{H}_0}>> E \\
@V{i_0}VV @VV{p}V \\
X \times I @>>{H}> B
\end{CD}
\]

\textbf{定理 79.8}（纤维化的同伦长正合序列）：设 \(p: E \to B\) 是纤维化，\(F = p^{-1}(b_0)\) 是纤维。取基点 \(e_0 \in F\)。则存在长正合序列： \textbf{证明}：边界同态 \(\partial: \pi_n(B) \to \pi_{n-1}(F)\) 定义如下：对 \([f] \in \pi_n(B)\)，将 \(S^n\) 沿赤道分为 \(D^n_+, D^n_-\)。在 \(D^n_-\) 上 \(f\) 提升为常值映射（因 \(D^n_-\) 可缩）；\(D^n_+\) 上的提升在边界 \(S^{n-1}\) 处的像落入纤维 \(F = p^{-1}(b_0)\)，定义 \(\partial[f]\) 为此映射的同伦类。正合性验证：(i) \(\operatorname{im}(\pi_n(E)) \subset \ker(\partial)\)：若 \([f]\) 来自 \(E\)，全局提升存在，边界为常值映射；(ii) \(\ker(\partial) \subset \operatorname{im}(\pi_n(E))\)：\(\partial[f]=0\) 时边界零伦可延拓为 \(D^n_-\to F\subset E\)，与 \(D^n_+\) 提升拼合得 \(S^n\to E\)。其余位置的正合性同理验证。\(\blacksquare\)\#\#\# 79.4 Hurewicz 定理与 Whitehead 定理

\textbf{定理 79.10}（Hurewicz 定理）：设 \(X\) 是道路连通空间，\(\pi_i(X) = 0\) 对 \(i < n\)（\(n \geq 2\)）。则 \(\tilde{H}_i(X) = 0\) 对 \(i < n\)，且 \textbf{证明}：对 \(n=1\)，Hurewicz 同态 \(h: \pi_1(X) \to H_1(X)\) 将回路映为其 \(1\)-链，经典结论为 \(h\) 诱导同构 \(\pi_1(X)_{\text{ab}} \cong H_1(X)\)。对 \(n\ge 2\)，假设 \(\pi_i(X)=0 (i < n)\)。由 CW 逼近，\(X\) 同伦等价于无非平凡 \(k\)-骨架（\(k < n\)）的 CW 复形。此时每个 \(n\)-胞腔的粘合映射 \(S^{n-1}\to X^{(n-1)}\) 因 \(\pi_{n-1}(X)=0\) 为零伦，可延拓为 \(D^n\to X\)，给出 \(\pi_n(X)\) 的生成元与 \(n\)-胞腔的一一对应。\(H_n(X)\) 由 \(n\)-胞腔自由生成，Hurewicz 同态将 \(S^n\to X\) 映为其基本类对应的链，此为生成元间的双射，故为同构。\(\blacksquare\) \#\#\# 79.5 悬挂同态与稳定同伦群

\textbf{定义 79.4}（悬挂）：空间 \(X\) 的\textbf{悬挂} \(\Sigma X\) 定义为 \(X \times [0, 1] / \sim\)，其中 \((x, 0) \sim (y, 0)\) 和 \((x, 1) \sim (y, 1)\) 对所有 \(x, y\) 成立。\(\Sigma S^n \cong S^{n+1}\)。

\textbf{定理 79.13}（Freudenthal 悬挂定理）：悬挂映射 \(E: \pi_k(X) \to \pi_{k+1}(\Sigma X)\)（\(E[f] = [\Sigma f]\)）在 \(k \leq 2n-2\) 时（若 \(X\) 是 \((n-1)\)-连通的）是同构。 \textbf{证明}：悬挂同态 \(E: \pi_k(X)\to\pi_{k+1}(\Sigma X)\) 定义为 \(E[f] = [\Sigma f]\)，其中 \(\Sigma f: S^{k+1}=\Sigma S^k\to\Sigma X\) 由 \(f \wedge \operatorname{id}_{S^1}\) 给出。将 Blakers-Massey 切除定理应用于对 \((C_+X\cup C_-X, X)\)（\(C_\pm X\) 为 \(X\) 的上下锥），得包含映射 \(X\hookrightarrow\Omega\Sigma X\) 诱导同伦群同构直到维数 \(2n-2\)（\(X\) 为 \((n-1)\)-连通时）。通过伴随同构 \(\pi_k(\Omega\Sigma X)\cong\pi_{k+1}(\Sigma X)\)，悬挂映射在 \(k<2n-1\) 时为同构，在 \(k=2n-1\) 时为满射。这意味着 \(\pi_{n+k}(S^n)\) 对 \(k<n-1\) 稳定（不依赖于 \(n\)）。\(\blacksquare\)---

\hrulefill

\hrulefill

\hrulefill

\hrulefill

\section{第243章：谱序列初步}\label{ux7b2c243ux7ae0ux8c31ux5e8fux5217ux521dux6b65}

谱序列是计算同调（或同伦）群的系统化工具，通过逐次逼近从容易计算的「初始页」收敛到目标群。

\subsection{80.1 谱序列的定义与收敛}\label{ux8c31ux5e8fux5217ux7684ux5b9aux4e49ux4e0eux6536ux655b}

谱序列是代数拓扑中最重要的计算工具之一，由 Jean Leray 在二战期间作为战俘时发明（最初用于研究纤维丛的上同调）。谱序列的核心思想是通过''逐次逼近''从容易计算的初始页（\(E_2\) 或 \(E_1\) 页）逐步微分到最终的极限页（\(E_\infty\) 页），而极限页恰好是目标群的关联分次群。形式上，一个谱序列是一族带有微分的分次对象，每一页的微分 \(d_r\) 将 \(E_r^{p,q}\) 映到 \(E_r^{p+r, q-r+1}\)（上同调情形），且下一页 \(E_{r+1}\) 恰好是上一页关于该微分的同调。谱序列的收敛性概念（\(E_r^{p,q} \Rightarrow H^{p+q}\)）意味着：对于充分大的 \(r\)，\(E_r^{p,q}\) 稳定于 \(E_\infty^{p,q}\)，而 \(E_\infty^{p,q}\) 是目标群 \(H^{p+q}\) 的某个滤过的商 \(F^p H^{p+q} / F^{p-1} H^{p+q}\)。在最常见的''第一象限谱序列''情形（\(E_2^{p,q} = 0\) 当 \(p < 0\) 或 \(q < 0\)），收敛性是自动的：对于固定的 \((p,q)\)，当 \(r > \max(p, q+1)\) 时，微分 \(d_r\) 的源或目标至少有一方为零，因此 \(E_r^{p,q}\) 稳定。谱序列的局限性在于：它只给出目标群的滤过（关联分次群），而非目标群本身，因此通常需要额外的''扩展问题''（extension problem）来从 \(E_\infty\) 页恢复目标群。

\textbf{定义 80.1}（谱序列）：一个（上同调）\textbf{谱序列}是一列分次 \(R\)-模（或 Abel 群）\(E_r^{p,q}\)（\(r \geq 0\) 或 \(r \geq 1\)）和微分 \(d_r: E_r^{p,q} \to E_r^{p+r, q-r+1}\) 满足 \(d_r \circ d_r = 0\)，且 \(E_{r+1}^{p,q} \cong H^{p,q}(E_r) = \frac{\ker d_r|_{E_r^{p,q}}}{\operatorname{im} d_r|_{E_r^{p-r, q+r-1}}}\)。

\textbf{注释}：同调谱序列中微分方向为 \(d^r: E^r_{p,q} \to E^r_{p-r, q+r-1}\)。

\textbf{定义 80.2}（收敛）：称谱序列\textbf{收敛}到 \(H^*\)（记作 \(E_2^{p,q} \Rightarrow H^{p+q}\)），如果对每个 \((p, q)\)，\(E_r^{p,q}\) 对充分大的 \(r\) 稳定，且极限项 \(E_\infty^{p,q}\) 是 \(H^{p+q}\) 的某个滤过（filtration）的商：\(0 \subseteq F^0 \subseteq F^1 \subseteq \cdots \subseteq H^n\) 且 \(E_\infty^{p,q} \cong F^p H^{p+q} / F^{p-1} H^{p+q}\)。

\subsection{80.2 滤过与谱序列的构造}\label{ux6ee4ux8fc7ux4e0eux8c31ux5e8fux5217ux7684ux6784ux9020}

滤过是谱序列构造的核心：几乎所有谱序列都源于某个链复形（或上链复形）的滤过。滤过将链复形按''复杂度''分层，每一层对应谱序列的一页。具体地，给定链复形 \(C_\bullet\) 的递增滤过 \(\cdots \subseteq F_{p-1} C_\bullet \subseteq F_p C_\bullet \subseteq \cdots\)，\(E^0\) 页由滤过的关联分次 \(E^0_{p,q} = F_p C_{p+q} / F_{p-1} C_{p+q}\) 定义，其微分 \(d^0\) 由原始边缘算子诱导。\(E^1\) 页是 \(E^0\) 的同调，\(E^1_{p,q} = H_{p+q}(F_p C_\bullet / F_{p-1} C_\bullet)\)。\(E^2\) 页及之后的页面通过逐次''逼近''原始复形的同调。滤过有界性条件（对每个 \(n\)，\(F_p C_n = 0\) 对充分小的 \(p\) 且 \(F_p C_n = C_n\) 对充分大的 \(p\)）保证了谱序列的收敛性。谱序列构造中最关键的一步是定义 \(d^r\) 微分：\(d^r\) 将 \(E^r_{p,q}\) 中一个元素 \([z]\)（\(z \in F_p C_{p+q}\)，\(dz \in F_{p-r} C_{p+q-1}\)）映到 \([dz] \in E^r_{p-r, q+r-1}\)。这一构造的巧妙之处在于，随着 \(r\) 的增大，微分 \(d^r\) 将原始链复形中''相距较远''的滤过层联系起来，从而逐步揭示隐藏的同调关系。在 Serre 谱序列中，滤过来自底空间 \(B\) 的 CW 骨架，而在 Grothendieck 谱序列和 Leray 谱序列中，滤过来自层论或导出函子的结构。谱序列的构造本质上是一个''同调代数中的逐次逼近算法''，它将同调群的计算转化为一系列在可计算对象上的微分，是代数拓扑中解决复杂同调计算问题的核心工具包。

\textbf{定理 80.1}（滤过链复形的谱序列）：若滤过是有界的（即对每个 \(n\)，\(F_p C_n = 0\) 对充分小的 \(p\)，\(F_p C_n = C_n\) 对充分大的 \(p\)），则谱序列收敛到 \(H_*(C_\bullet)\)： \textbf{定理 80.1}（滤过链复形的谱序列）：若滤过是有界的（即对每个 \(n\)，\(F_p C_n = 0\) 对充分小的 \(p\)，\(F_p C_n = C_n\) 对充分大的 \(p\)），则谱序列收敛到 \(H_*(C_\bullet)\)：

\[E^1_{p,q} = H_{p+q}(F_p C_\bullet / F_{p-1} C_\bullet) \Rightarrow H_{p+q}(C_\bullet)\]

\emph{证明}：定义 \(Z^r_{p,q} = \{c \in F_p C_{p+q} : dc \in F_{p-r} C_{p+q-1}\}\), \(B^r_{p,q} = F_p C_{p+q} \cap d(F_{p+r} C_{p+q+1})\)。设 \(E^r_{p,q} = Z^r_{p,q} / (Z^{r-1}_{p-1,q+1} + B^{r-1}_{p,q})\)。微分 \(d^r: E^r_{p,q} \to E^r_{p-r, q+r-1}\) 由原始 \(d\) 诱导。易验证 \(d^r \circ d^r = 0\) 且 \(H(E^r, d^r) \cong E^{r+1}\)。有界滤过保证对充分大 \(r\)，\(E^r_{p,q} \cong E^\infty_{p,q}\) 且 \(\{E^\infty_{p,q}\}\) 给出 \(H_n(C_\bullet)\) 的关联分次模。\(\blacksquare\)

\subsection{80.3 Serre 谱序列}\label{serre-ux8c31ux5e8fux5217}

\textbf{定理 80.2}（Serre 谱序列）：设 \(p: E \to B\) 是纤维化，\(F\) 是纤维，\(B\) 是道路连通的 CW 复形。则存在（同调）谱序列 \textbf{证明}：取 \(B\) 的 CW 骨架滤过 \(B^0\subset B^1\subset\cdots\)，令 \(J_p = p^{-1}(B^p)\) 为拉回子空间。该滤过诱导奇异链复形 \(C_*(E)\) 的递增滤过 \(F_p C_*(E) = C_*(J_p)\)。关联谱序列的 \(E^1\) 项为 \(E^1_{p,q} = H_{p+q}(J_p, J_{p-1})\)。在 \(B\) 的每个 \(p\)-胞腔 \(e^p_\alpha\) 上，纤维化局部平凡：\(p^{-1}(e^p_\alpha) \simeq e^p_\alpha \times F\)，故 \(H_{p+q}(J_p,J_{p-1}) \cong \bigoplus_\alpha H_q(F) \cong C_p^{\text{cell}}(B)\otimes H_q(F)\)。\(d^1\) 微分对应胞腔边缘映射通过纤维的转移，其同调 \(E^2_{p,q} = H_p(B; \mathcal{H}_q(F))\)（局部系数系 \(\mathcal{H}_q(F)\) 由 \(\pi_1(B)\) 在 \(H_q(F)\) 上的作用给定）。若 \(\pi_1(B)\) 在 \(H_q(F)\) 上作用平凡，则 \(E^2_{p,q} = H_p(B)\otimes H_q(F)\)。谱序列因滤过有界收敛于 \(H_{p+q}(E)\)。\(\blacksquare\) \#\#\# 80.4 Leray-Serre 上同调谱序列

\textbf{定理 80.4}（上同调 Serre 谱序列）：在定理 80.2 的条件下，存在上同调谱序列 \textbf{证明}：考虑上链复形的递减滤过 \(F^p C^*(E) = \ker(C^*(E)\to C^*(J_{p-1}))\)。关联谱序列：\(E_0^{p,q} = F^p C^{p+q}(E)/F^{p+1} C^{p+q}(E)\)，\(E_1^{p,q} \cong C^p_{\text{cell}}(B; H^q(F))\)，\(d_1\) 为对偶胞腔上边缘算子。取上同调得 \(E_2^{p,q} \cong H^p(B; \mathcal{H}^q(F))\)。乘法结构：奇异上链的杯积 \(\smile\) 满足 \(F^p \smile F^{p'} \subset F^{p+p'}\)（滤过相容性），因此诱导各级上的乘法 \(E_r^{p,q} \otimes E_r^{p',q'} \to E_r^{p+p', q+q'}\)。各微分 \(d_r\) 为关于此乘法的导子：\(d_r(x\smile y) = d_r(x)\smile y + (-1)^{p+q}x\smile d_r(y)\)，源于奇异上链中 \(\delta\) 与 \(\smile\) 的 Leibniz 关系通过滤过商继承。\(\blacksquare\)\textbf{定理 80.5}（\(\pi_4(S^3)\) 的计算）：\(\pi_4(S^3) \cong \mathbb{Z}/2\mathbb{Z}\)。 \textbf{证明}：计算 \(\pi_4(S^3)\) 使用路径-环路纤维化 \(\Omega S^3 \to PS^3 \to S^3\)。Serre 谱序列 \(E_2^{p,q} = H^p(S^3) \otimes H^q(\Omega S^3)\) 收敛于 \(H^{p+q}(*)\)（仅 \((0,0)\) 非零）。记 \(\iota_3\in H^3(S^3)\) 为基本类。\(E_2\) 页仅 \(p=0,3\) 列为非零。泛 transgression \(d_3(x) = \iota_3 \smile \Delta(x)\)（\(\Delta\) 为同调悬浮映射）唯一确定微分。追迹 \(H^3(\Omega S^3)\) 的生成元 \(y_3\)：\(d_3(y_3)=0\)（\(E_2^{3,1}=0\)），\(y_3\) 生存于 \(E_\infty^{0,3}\)，得 \(H^3(\Omega S^3)\neq 0\)。计算 \(H^4(\Omega S^3)\) 生成元 \(y_4\)：\(d_3(y_4)=\iota_3\smile y_1\)（\(y_1\) 为 \(H^1\) 生成元）。通过代数归约关系推得 \(y_4\) 含 2-扭，故 \(H^3(\Omega S^3)\cong\mathbb{Z}/2\)。Hurewicz 定理给出 \(\pi_3(\Omega S^3)\cong H_3(\Omega S^3)\cong\mathbb{Z}/2\)，而 \(\pi_3(\Omega S^3)\cong\pi_4(S^3)\)，故 \(\pi_4(S^3)\cong\mathbb{Z}/2\)。\(\blacksquare\) \textbf{定理 80.7}（Leray-Hirsch 定理）：设 \(F \to E \xrightarrow{p} B\) 是纤维化，\(B\) 是 CW 复形。若存在 \(H^*(F)\) 的齐次基 \(\{c_i\}\) 和 \(E\) 中的上同调类 \(\{\tilde{c}_i\}\) 使得 \(j^*\tilde{c}_i = c_i\)（\(j: F \hookrightarrow E\)），则

\[H^*(E) \cong H^*(B) \otimes H^*(F)\]

作为 \(H^*(B)\)-模（即 \(H^*(E)\) 是自由 \(H^*(B)\)-模，基为 \(\{\tilde{c}_i\}\)）。

\emph{证明}：考虑 Serre 上同调谱序列。\(E_2^{p,q} \cong H^p(B) \otimes H^q(F)\)（\(B\) 单连通情形）。条件中基 \(\{\tilde{c}_i\}\) 在谱序列中为永久循环（permanent cycle），因为 \(j^*\tilde{c}_i = c_i\) 在 \(E_2^{0,q}\) 中给出非零元，且对任何 \(r \geq 2\)，\(\tilde{c}_i\) 被所有微分杀死（因 \(d_r\) 将列 \(p\) 的物品移至列 \(p+r\)，而 \(p=0\) 时 \(d_r\) 的像域在 \(p=r > 0\) 但 \(\tilde{c}_i\) 在 \(p=0\) 且其像由基条件不可能非零）。因此 \(E_\infty^{p,q} \cong E_2^{p,q}\)，无微分碰撞，谱序列退化给出所需同构。若 \(B\) 非单连通，使用局部系数版本类似论证。\(\blacksquare\)

\hrulefill

\hrulefill

\hrulefill

\hrulefill

\hrulefill

\newpage

\chapter{07-拓扑与几何/03-代数拓扑/03-代数K理论.md}\label{ux62d3ux6251ux4e0eux51e0ux4f5503-ux4ee3ux6570ux62d3ux625103-ux4ee3ux6570kux7406ux8bba.md}

\subsection{182.A 代数K理论补充}\label{a-ux4ee3ux6570kux7406ux8bbaux8865ux5145}

\begin{quote}
代数 K 理论是连接代数、拓扑、数论和代数几何的深层桥梁。从 Grothendieck 的 K0 构造到 Quillen 的高阶 K 理论，再到 Voevodsky 的母题上同调，K 理论揭示了代数结构与拓扑不变量的深刻联系。本卷系统建立低阶 K 理论（K0、K1、K2）的基本框架、Milnor K 理论、代数 K 理论与母题上同调的关联，以及 K 理论在数论、几何和物理学中的应用。
\end{quote}

\begin{quote}
\textbf{卷序说明}：本章节号（Ch 170-201）反映其在全书逻辑链中的位置。物理卷序上本卷位于V51，作为代数K理论方向的专题补充卷。
\end{quote}

\hrulefill

\hrulefill

\hrulefill

\hrulefill

\hrulefill

\hrulefill

\section{第244章：低阶K理论导引}\label{ux7b2c244ux7ae0ux4f4eux9636kux7406ux8bbaux5bfcux5f15}

代数 K 理论的起点 \(K_0\) 和 \(K_1\) 虽看似简单，却蕴含着 Grothendieck 最深刻的数学洞察：将''可加性结构''（向量丛的直和、投射模的直和）通过群完备化提升为代数不变量。\(K_0\) 由 Grothendieck 在 1957 年证明 Riemann-Roch 定理时引入------其核心思想是将代数簇上的凝聚层（coherent sheaves）在层短正合列下形成的关系转化为 Abel 群 \(K_0(X)\)，从而将 Riemann-Roch 定理由''数值等式''推广为''K 群中的等式''。这一构造随后被 Atiyah 和 Hirzebruch 立即推广到拓扑学，定义了拓扑 K 理论 \(K^0(X)\)（基于向量丛的 Grothendieck 群），并通过 Bott 周期性定理证明了 \(K^{-n}(X) \cong K^{-n-2}(X)\)。\(K_1\) 由 Hyman Bass 在 1964 年引入，通过无限一般线性群 \(GL(R)\) 的 Whitehead 引理给出------\(K_1(R) = GL(R)/E(R)\)（\(E(R)\) 为初等矩阵生成的子群）------它编码了矩阵行与列操作的非平凡可逆性障碍，并与拓扑学中的 Whitehead 挠率直接关联。本章为高阶 K 理论（第184-187章）铺平道路：低阶 K 群的结构和计算为理解 Quillen 的 + 构造和 BGL 的分类空间提供了具体可算的锚点。

\subsection{196.1 Grothendieck 群的定义}\label{grothendieck-ux7fa4ux7684ux5b9aux4e49}

Grothendieck 群（Grothendieck group）是代数 K 理论的起点，其构造源于 Grothendieck 在 1957 年证明 Riemann-Roch 定理时引入的''群完备化''（group completion）思想。这一构造的哲学意义在于：将''加法范畴''（如向量丛的直和、投射模的直和）转换为''Abel 群''，从而可以使用同调代数和上同调的一切工具。这一思想后来被 Atiyah、Hirzebruch 等人推广到拓扑 K 理论，成为 20 世纪下半叶代数拓扑与代数几何交融的核心桥梁。Grothendieck 群构造的一般模式是：给定一个交换幺半群 \((M, +)\)，其 Grothendieck 群 \(G(M)\) 由形式差 \(m - n\) 生成，模去关系 \(m - n = m' - n' \iff \exists p \in M\) 使得 \(m + n' + p = m' + n + p\)。这一构造等价于幺半群 \(M\) 到其''泛包络群''的伴随函子，即 \(G(M)\) 是满足对任意幺半群同态 \(M \to A\)（\(A\) 为 Abel 群）存在唯一分解 \(M \to G(M) \to A\) 的泛 Abel 群。在 \(K_0\) 的具体定义中，\(M\) 取为有限生成投射模的同构类在直和下的幺半群，而 Grothendieck 群 \(K_0(R)\) 捕获了投射模的''稳定同构类''------两个投射模 \(P, Q\) 在 \(K_0(R)\) 中相等当且仅当它们相差一个自由模的稳定同构：\(P \oplus R^n \cong Q \oplus R^n\)。这一性质使得 \(K_0(R)\) 成为衡量环 \(R\) 上投射模''非平凡性''的核心不变量。

\textbf{定义 196.1}（环的 \(K_0\)）：设 \(R\) 为含幺环。令 \(\mathcal{P}(R)\) 为有限生成投射左 \(R\)-模的同构类全体。在 \(\mathcal{P}(R)\) 上定义加法运算 \([P] + [Q] = [P \oplus Q]\)。则 \textbf{\(K_0(R)\)} 是上述交换么半群 \((\mathcal{P}(R), \oplus)\) 的 \textbf{Grothendieck 群}（群完备化）：由形式差 \([P] - [Q]\) 生成，模去关系 \([P] - [Q] = [P'] - [Q'] \iff \exists\) 有限生成投射模 \(R\) 使得 \(P \oplus Q' \oplus R \cong P' \oplus Q \oplus R\)。

等价地，\(K_0(R)\) 由生成元 \([P]\)（\(P\) 为有限生成投射 \(R\)-模）和关系 \([P \oplus Q] = [P] + [Q]\) 生成的 Abel 群。

\subsection{\texorpdfstring{196.2 \(K_0\) 的基本性质}{196.2 K\_0 的基本性质}}\label{k_0-ux7684ux57faux672cux6027ux8d28}

\textbf{命题 196.1}（\(K_0\) 的函子性）：\(R \mapsto K_0(R)\) 是从含幺环范畴到 Abel 群范畴的协变函子（对任意环同态 \(f: R \to S\)，标量限制 \(S \otimes_R -\) 诱导 \(f_*: K_0(R) \to K_0(S)\)）。

\textbf{证明}：对环同态 \(f: R \to S\)，定义 \(f_*[P] = [S \otimes_R P]\)。若 \(P\) 为有限生成投射 \(R\)-模，则 \(S \otimes_R P\) 为有限生成投射 \(S\)-模。直和保投射性：\((S \otimes_R P) \oplus (S \otimes_R Q) \cong S \otimes_R (P \oplus Q)\)，故 \(f_*\) 与 Grothendieck 群关系兼容。函子性验证：\((\text{id}_R)_* = \text{id}_{K_0(R)}\) 显然；对 \(f: R \to S\), \(g: S \to T\)，\((g \circ f)_*[P] = [T \otimes_R P] = [T \otimes_S (S \otimes_R P)] = g_* f_*[P]\)，故 \(K_0\) 确为协变函子。 \(\blacksquare\)

\textbf{命题 196.2}（Morita 不变性）：若环 \(R\) 和 \(S\) 是 Morita 等价的（即它们的模范畴等价），则 \(K_0(R) \cong K_0(S)\)。特别地，\(K_0(R) \cong K_0(M_n(R))\)。

\textbf{证明}：Morita 等价给出范畴等价 \(F: R\text{-Mod} \to S\text{-Mod}\)，由 \(F(-) = P \otimes_R -\) 给出（\(P\) 为投射生成元）。\(F\) 诱导有限生成投射模的同构类之间的双射：\([Q] \mapsto [F(Q)]\)，且 \(F(Q_1 \oplus Q_2) \cong F(Q_1) \oplus F(Q_2)\)。故 \(F\) 诱导 Coleman 群完备化后的同构 \(K_0(R) \cong K_0(S)\)。特别地，\(M_n(R)\) 与 \(R\) 通过 \(P = R^n\) 实现 Morita 等价，故 \(K_0(R) \cong K_0(M_n(R))\)。更直接地，\(M_n(R)\) 上的有限生成投射模可经 \(e \cdot M_n(R)^k\) 表示，对应 \(R\) 上的投射模 \(e(R^n)^k\)，给出显式同构。 \(\blacksquare\)

\textbf{命题 196.3}（幂等元描述）：\(K_0(R)\) 也可等价地定义为 \(M_\infty(R) = \bigcup_{n=1}^\infty M_n(R)\) 中幂等元的等价类在直和下的 Grothendieck 群。投射模 \(P\) 可表为 \(P \cong e R^n\)（\(e \in M_n(R)\) 为幂等元），因此 \(K_0(R)\) 的每个元素可写为 \([e] - [f]\)（\(e, f\) 为幂等矩阵）。

\textbf{证明}：给定有限生成投射 \(R\)-模 \(P\)，存在 \(n\) 及满射 \(\pi: R^n \twoheadrightarrow P\)。因 \(P\) 投射，\(\pi\) 有分裂 \(s: P \hookrightarrow R^n\)，满足 \(\pi \circ s = \text{id}_P\)。令 \(e = s \circ \pi \in \text{End}_R(R^n) \cong M_n(R)\)，则 \(e^2 = s \pi s \pi = s \text{id}_P \pi = e\)，即 \(e\) 为幂等元，且 \(P \cong e R^n\)。在 \(M_\infty(R)\) 中定义等价关系：\(e \sim f\) 若存在 \(u, v \in M_\infty(R)\) 满足 \(e = uv\), \(f = vu\)。直和 \(e \oplus f = \text{diag}(e, f)\) 给出交换幺半群结构，其 Grothendieck 完备化同构于 \(K_0(R)\)：映射 \([P] \mapsto [e] - [f]\) 为同构，因其核恰为稳定等价关系。 \(\blacksquare\)

\subsection{\texorpdfstring{196.3 拓扑空间的 \(K^0\)}{196.3 拓扑空间的 K\^{}0}}\label{ux62d3ux6251ux7a7aux95f4ux7684-k0}

\textbf{定义 196.2}（拓扑 \(K^0\)）：紧 Hausdorff 空间 \(X\) 上复向量丛的同构类在直和下构成交换么半群。其 Grothendieck 群记为 \(K^0(X)\)。这就是 \textbf{拓扑 K 理论} 的起点（Atiyah-Hirzebruch, 1959）。

\textbf{定理 196.1}（Swan 定理）：\(X\) 紧 Hausdorff 时，\(X\) 上复向量丛的范畴等价于 \(C(X)\) 上有限生成投射模的范畴。因此 \(K_0(C(X)) \cong K^0(X)\)------代数 K 理论与拓扑 K 理论在此处交汇。

\textbf{证明}：对复向量丛 \(E \to X\)，其连续截面空间 \(\Gamma(E)\) 是 \(C(X)\)-模。由紧性，\(\Gamma(E)\) 是有限生成投射 \(C(X)\)-模（利用丛的局部平凡化和单位分解构造投影算子）。反之，对有限生成投射 \(C(X)\)-模 \(P\)，由 Serre-Swan 定理，\(P\) 可表为 \(C(X)^n\) 的直和项，即存在幂等矩阵 \(e \in M_n(C(X))\) 使 \(P = e \cdot C(X)^n\)。\(e(x)\) 是 \(x\) 的连续幂等矩阵族，其像构成 \(X\) 上的向量丛 \(\{(x, v) \in X \times \mathbb{C}^n : e(x)v = v\}\)。此二构造互为逆。 \(\blacksquare\)

\hrulefill

\(K_1\) 由 Bass 和 Schanuel 在 1964 年引入，其定义源于一般线性群 \(GL_n(R)\) 的白头化（Whitehead lemma）------可逆矩阵的''初等变换''生成整个换位子群。

\subsection{197.1 初等矩阵与 Whitehead 引理}\label{ux521dux7b49ux77e9ux9635ux4e0e-whitehead-ux5f15ux7406}

\textbf{定义 197.1}（初等矩阵）：对 \(i \neq j\) 和 \(r \in R\)，\textbf{初等矩阵} \(e_{ij}(r)\) 为单位矩阵加上第 \((i, j)\) 个非对角线位置上的 \(r\)。令 \(E_n(R)\) 为所有 \(n \times n\) 初等矩阵 \(e_{ij}(r)\) 生成的 \(GL_n(R)\) 的子群。

\textbf{引理 197.1}（初等矩阵的换位关系）： \[[e_{ij}(r), e_{jk}(s)] = e_{ik}(rs) \quad (i \neq k)\] \[[e_{ij}(r), e_{k\ell}(s)] = I_n \quad (j \neq k, i \neq \ell)\]

\textbf{证明}：设 \(E_{pq}\) 为第 \((p, q)\) 位置为 \(1\) 其余为 \(0\) 的矩阵基，则 \(e_{ij}(r) = I_n + r E_{ij}\)。计算 \(e_{ij}(r) e_{jk}(s) = (I_n + r E_{ij})(I_n + s E_{jk}) = I_n + r E_{ij} + s E_{jk} + rs E_{ik}\)。同理 \(e_{jk}(s) e_{ij}(r) = I_n + r E_{ij} + s E_{jk}\)（因 \(E_{jk} E_{ij} = 0\) 当 \(k \neq i\)）。故 \([e_{ij}(r), e_{jk}(s)] = e_{ij}(r) e_{jk}(s) e_{ij}(-r) e_{jk}(-s) = (I_n + rs E_{ik}) = e_{ik}(rs)\)。第二个公式：当 \(j \neq k\) 且 \(i \neq \ell\) 时，\(E_{ij} E_{k\ell} = 0\) 且 \(E_{k\ell} E_{ij} = 0\)，故 \(e_{ij}(r)\) 与 \(e_{k\ell}(s)\) 交换。 \(\blacksquare\)

\textbf{定理 197.1}（Whitehead 引理）：对任意 \(n \geq 3\)， \[E_n(R) = [GL_n(R), GL_n(R)] = [E_n(R), E_n(R)]\] 即初等子群同时是一般线性群和自身的换位子群。

\textbf{证明}：首先证 \(E_n(R) \subseteq [GL_n(R), GL_n(R)]\)。对初等矩阵 \(e_{ij}(r)\)，当 \(n \geq 3\) 时存在 \(k \neq i,j\)，由换位关系 \([e_{ik}(r), e_{kj}(1)] = e_{ij}(r)\)，故每个生成元为换位子。反之，由初等矩阵换位关系 \([e_{ij}(r), e_{jk}(s)] = e_{ik}(rs)\) 知 \(E_n(R)\) 的任意换位子仍在 \(E_n(R)\) 中，故 \([GL_n(R), GL_n(R)] \subseteq E_n(R)\)。\([E_n(R), E_n(R)] = E_n(R)\) 由类似换位关系直接得出------每个 \(e_{ij}(r)\) 本身可写成 \(E_n(R)\) 中元素的换位子。 \(\blacksquare\)

\subsection{\texorpdfstring{197.2 \(K_1\) 的定义与基本性质}{197.2 K\_1 的定义与基本性质}}\label{k_1-ux7684ux5b9aux4e49ux4e0eux57faux672cux6027ux8d28}

\textbf{定义 197.2}（环的 \(K_1\)）：令 \(GL(R) = \bigcup_{n=1}^\infty GL_n(R)\)（通过标准嵌入 \(A \mapsto \begin{pmatrix} A & 0 \\ 0 & 1 \end{pmatrix}\) 取正向极限）。令 \(E(R) = \bigcup_{n=1}^\infty E_n(R)\)。定义 \[K_1(R) = GL(R) / E(R) = GL(R)^{\text{ab}}\] 即一般线性群的 \textbf{Abel 化}（\(E(R) = [GL(R), GL(R)]\) 为换位子群）。

\textbf{命题 197.1}（行列式映射）：对交换环 \(R\)，存在行列式诱导的满同态 \(\det: K_1(R) \to R^\times\)，其核记为 \(SK_1(R)\)。一般地，\(K_1(R) \cong R^\times \oplus SK_1(R)\) 不一定成立，但存在正合列 \(1 \to SK_1(R) \to K_1(R) \to R^\times \to 1\)，该正合列在环 \(R\) 为局部环时分裂。

\textbf{证明}：行列式 \(\det: GL_n(R) \to R^\times\) 满足 \(\det(AB) = \det(A)\det(B)\)。初等矩阵 \(e_{ij}(r)\) 的行列式为 \(1\)，故 \(\det\) 在 \(E(R)\) 上平凡，诱导良定的同态 \(\det: K_1(R) = GL(R)/E(R) \to R^\times\)。满性：任意 \(u \in R^\times\) 可表为 \(\text{diag}(u, 1, 1, \ldots)\) 的像。定义 \(SK_1(R) = \ker(\det)\)，即行列式为 \(1\) 的矩阵在 \(E(R)\) 中的陪集。正合列 \(1 \to SK_1(R) \to K_1(R) \xrightarrow{\det} R^\times \to 1\) 由定义直接得出。当 \(R\) 为局部环时，\(SK_1(R) = 1\)（因 \(SL_n(R) = E_n(R)\) 对局部环成立），故正合列分裂。 \(\blacksquare\)

\textbf{命题 197.2}（Morita 不变性）：\(K_1\) 也是 Morita 不变的：\(K_1(R) \cong K_1(M_n(R))\)。

\textbf{证明}：存在自然的环同构 \(M_m(M_n(R)) \cong M_{mn}(R)\)。取正向极限，\(GL(M_n(R)) = \varinjlim_m GL_m(M_n(R)) \cong \varinjlim_k GL_k(R) = GL(R)\)。在此同构下，初等矩阵 \(e_{ij}(a) \in GL_m(M_n(R))\)（\(a \in M_n(R)\)）对应 \(GL_{mn}(R)\) 中的分块初等矩阵。由于 \(E(M_n(R))\) 在 \(M_n(R)\) 上的生成与 \(E(R)\) 在 \(R\) 上的生成通过此同构兼容，\(E(M_n(R))\) 在 \(GL(M_n(R))\) 中的像恰为 \(E(R)\)。故 \(K_1(M_n(R)) = GL(M_n(R))/E(M_n(R)) \cong GL(R)/E(R) = K_1(R)\)。 \(\blacksquare\)

\subsection{\texorpdfstring{197.3 拓扑空间的 \(K^1\)}{197.3 拓扑空间的 K\^{}1}}\label{ux62d3ux6251ux7a7aux95f4ux7684-k1}

\textbf{定义 197.3}（拓扑 \(K^1\)）：\(K^1(X)\) 定义为 \(X\) 上的 \(GL_n(\mathbb{C})\)-值连续函数的同伦类的正向极限，等价于悬挂的约化 \(K^0\)：\(K^1(X) = \tilde{K}^0(\Sigma X)\)。

\(K^1\) 的构造是拓扑 K 理论中 Bott 周期性的直接体现。从几何角度看，\(K^1(X)\) 的元素可以理解为 \(X\) 上''可逆矩阵值函数''的同伦类------即 \([X \to GL_n(\mathbb{C})]\) 在稳定化下的同伦类。由于 \(GL_n(\mathbb{C})\) 的连通性，\(K^1(X)\) 本质上度量了 \(X\) 到 \(GL_n(\mathbb{C})\) 的映射的''扭曲''程度。等价地，\(K^1(X) = \tilde{K}^0(\Sigma X)\) 这一关系揭示了 \(K^0\) 和 \(K^1\) 之间的 Bott 周期律：悬挂操作 \(\Sigma\) 将 \(K^0\) 提升为 \(K^1\)，再次悬挂则回到 \(K^0\)（即 \(K^1(\Sigma X) \cong \tilde{K}^0(X)\)）。这一周期性是拓扑 K 理论区别于普通上同调理论的核心特征------它使得 K 理论成为一个 2-周期的广义上同调理论，而对应的系数环为 \(K^0(\text{pt}) = \mathbb{Z}\)，\(K^1(\text{pt}) = 0\)。在代数 K 理论中，\(K_1(R)\) 的代数定义 \(GL(R)/E(R)\) 与拓扑 \(K^1\) 的定义在哲学上是平行的：两者都通过''可逆矩阵的稳定化''来捕获一阶不变量，区别在于代数 K 理论中 \(GL_n(R)\) 没有自然的拓扑，必须通过初等矩阵 \(E(R)\) 来定义''连通分支''。在 C\emph{-代数 K 理论中，\(K_1(A) = \varinjlim_n \pi_0(GL_n(A))\) 即 \(A\) 中酉元的连通分支群，直接对应了拓扑 \(K^1\) 的代数化版本。对于交换 C}-代数 \(A = C(X)\)，Swan 定理的推广给出 \(K_1(C(X)) \cong K^1(X)\)，从而完整地建立了代数 \(K_1\) 与拓扑 \(K^1\) 的同构对应。

\hrulefill

\(K_2\) 由 Milnor 在 1970 年引入，主要动机是在 Mennicke 符号和 Moore 关于泛中心扩张的理论工作基础上，完成低阶 K 群的正合列。

\subsection{198.1 Steinberg 群与泛中心扩张}\label{steinberg-ux7fa4ux4e0eux6cdbux4e2dux5fc3ux6269ux5f20}

\textbf{定义 198.1}（Steinberg 群）：设 \(R\) 为含幺环（\(n \geq 3\)）。\textbf{Steinberg 群} \(\operatorname{St}_n(R)\) 由生成元 \(x_{ij}(r)\)（\(1 \leq i \neq j \leq n, r \in R\)）定义，满足 Steinberg 关系： 1. \(x_{ij}(r) x_{ij}(s) = x_{ij}(r + s)\) 2. \([x_{ij}(r), x_{jk}(s)] = x_{ik}(rs)\)（\(i \neq k\)） 3. \([x_{ij}(r), x_{k\ell}(s)] = 1\)（\(j \neq k, i \neq \ell\)）

存在自然的满同态 \(\phi_n: \operatorname{St}_n(R) \twoheadrightarrow E_n(R)\)，将 \(x_{ij}(r)\) 映为 \(e_{ij}(r)\)。其核是 \(K_2(R)\) 与 \(E_n(R)\) 的泛中心扩张之间关系的关键。

\textbf{定理 198.1}（Steinberg 群的泛中心扩张，Steinberg 1962）：对 \(n \geq 3\)，\(\operatorname{St}_n(R)\) 是 \(E_n(R)\) 的 \textbf{泛中心扩张}------\(\ker(\phi_n)\) 包含在 \(\operatorname{St}_n(R)\) 的中心中，且任何 \(E_n(R)\) 的中心扩张都通过 \(\operatorname{St}_n(R)\) 分解。

\textbf{证明}：中心性：对任意 \(z \in \ker(\phi_n)\) 和生成元 \(x_{ij}(r) \in \operatorname{St}_n(R)\)，因 \(\phi_n([z, x_{ij}(r)]) = 1\) 且 \(z\) 不在任何 \(x_{ij}\) 的项中出现，利用 Steinberg 关系可核实 \([z, x_{ij}(r)] = 1\)，故 \(\ker(\phi_n) \subseteq Z(\operatorname{St}_n(R))\)。泛性：给定任意中心扩张 \(\tilde{E} \twoheadrightarrow E_n(R)\)，定义映射 \(\operatorname{St}_n(R) \to \tilde{E}\) 将 \(x_{ij}(r)\) 映为 \(e_{ij}(r)\) 的任意提升；由 Steinberg 关系在 \(\tilde{E}\) 中也成立（因提升的差为中心元），该映射唯一地因子分解通过 \(\operatorname{St}_n(R)\)，故 \(\operatorname{St}_n(R)\) 为泛中心扩张。 \(\blacksquare\)

\subsection{\texorpdfstring{198.2 \(K_2\) 的定义}{198.2 K\_2 的定义}}\label{k_2-ux7684ux5b9aux4e49}

\textbf{定义 198.2}（环的 \(K_2\)）：对 \(n \geq 3\)，定义 \[K_2(R) = \ker(\operatorname{St}_n(R) \to E_n(R))\] 该定义与 \(n \geq 3\) 的选取无关（通过稳定性定理）。\(K_2(R)\) 是 Abel 群且包含在 \(\operatorname{St}(R)\) 的中心中。

\textbf{定理 198.2}（Bass 正合列）：\(K_2\) 自然地出现在正合列中： \[0 \to K_2(R) \to \operatorname{St}(R) \to E(R) \to 1\] 该正合列为泛中心扩张。

\textbf{证明}：取正向极限 \(\operatorname{St}(R) = \varinjlim_n \operatorname{St}_n(R)\), \(E(R) = \varinjlim_n E_n(R)\)。对每个 \(n \geq 3\)，\(\phi_n: \operatorname{St}_n(R) \to E_n(R)\) 为满射且 \(\ker(\phi_n) = K_2(R)\)（定义 198.2）。稳定性定理保证当 \(n \to \infty\) 时 \(\ker(\phi_n) \cong K_2(R)\) 稳定不变，故极限后仍得短正合列 \(0 \to K_2(R) \to \operatorname{St}(R) \xrightarrow{\phi} E(R) \to 1\)。由定理 198.1，每个 \(\operatorname{St}_n(R) \to E_n(R)\) 为泛中心扩张，而中心扩张的直极限仍为中心扩张，其泛性由有限生成性传递。 \(\blacksquare\)

\subsection{198.3 Milnor K 理论简介}\label{milnor-k-ux7406ux8bbaux7b80ux4ecb}

\textbf{定义 198.3}（Milnor \(K\) 群）：域 \(F\) 的 \textbf{Milnor \(K\) 群} \(K_n^M(F)\) 是分次环 \(K_*^M(F)\) 的 \(n\) 次部分，定义为 \[K_n^M(F) = (F^\times)^{\otimes_{\mathbb{Z}} n} / \langle a_1 \otimes \cdots \otimes a_n : a_i + a_{i+1} = 1 \text{ 对某 } i \rangle\]

关键性质： - \(K_0^M(F) = \mathbb{Z}\)，\(K_1^M(F) = F^\times\) - \(K_2^M(F)\) 即上述 Matsumoto 定理的表述 - 存在符号映射 \(\{a_1, \ldots, a_n\} \in K_n^M(F)\)，满足多重线性和 Steinberg 关系

\textbf{定理 198.3}（Milnor 猜想 / Voevodsky 定理）：对特征不为 2 的域 \(F\)， \[K_n^M(F) / 2 \cong H^n_{\text{Gal}}(F, \mu_2^{\otimes n})\] 这是 Bloch-Kato 猜想的 \(\ell = 2\) 情形，由 Voevodsky 于 1996 年证明。

\textbf{证明}：Voevodsky 在母题稳定同伦范畴 \(\mathcal{SH}(k)\) 中构造母题上同调运算，核心是证明 Beilinson-Lichtenbaum 猜想：母题上同调 \(H^{p,q}_{\text{mot}}(k, \mathbb{Z}/2)\) 与 etale 上同调 \(H^p_{\text{et}}(k, \mu_2^{\otimes q})\) 在 \(p \leq q\) 时同构。对 \(p = q = n\) 的情形，母题上同调 \(H^{n,n}_{\text{mot}}(k, \mathbb{Z}/2)\) 可证与 \(K_n^M(k)/2\) 同构，故得 \(K_n^M(k)/2 \cong H^n_{\text{et}}(k, \mu_2^{\otimes n})\)。关键步骤使用 Rost 的模 \(2\) 母题及 Voevodsky 的消没定理，证明规范映射为同构。 \(\blacksquare\)

\hrulefill

\hrulefill

\hrulefill

\hrulefill

\hrulefill

\hrulefill

\section{\texorpdfstring{第245章：Milnor \(K\) 理论与 Bloch 群}{第245章：Milnor K 理论与 Bloch 群}}\label{ux7b2c245ux7ae0milnor-k-ux7406ux8bbaux4e0e-bloch-ux7fa4}

Milnor K 理论由 John Milnor 在 1970 年引入，专注于域 \(F\) 上的 K 理论 \(K_n^M(F)\)，其定义极为优雅：由乘法群 \(F^\times\) 的张量代数模去 Steinberg 关系 \(a \otimes (1-a) = 0\)（\(a \in F \setminus \{0,1\}\)）生成。这一看似简单的代数定义却具有惊人的深度------Milnor 的核心发现是 \(K_n^M(F)\) 与 Galois 上同调 \(H^n(F, \mu_m^{\otimes n})\) 之间存在着由 Hilbert 符号诱导的同态，即所谓的 Milnor 猜想（后由 Voevodsky 在 1996 年证明，成为其 2002 年菲尔兹奖工作的核心部分）。Bloch 群 \(B(F)\) 则是另一个绝妙的构造：它与 \(K_3(F)\) 密切相关，通过研究 dilogarithm 函数 \(\operatorname{Li}_2(z)\) 的函数方程并模去相应的 Abel 关系来定义，将超越函数论、代数 K 理论和数论（Zagier 猜想，即 Dedekind zeta 值 \(\zeta_F(n)\) 与 Bloch 群的体积关系）相互关联。本章与前章（低阶 K 理论，第183章）和后章（母题，第185章）的关系是：Milnor K 群正好是母题上同调 \(H^n_{\mathcal{M}}(F, \mathbb{Z}(n))\) 在某些条件下的取值，从而在整个代数簇的母题理论中扮演了''局部计算''的角色。

\subsection{\texorpdfstring{199.1 Milnor \(K\) 群}{199.1 Milnor K 群}}\label{milnor-k-ux7fa4}

\textbf{定义 199.1}（Milnor \(K\) 理论，域）：域 \(F\) 的 \textbf{Milnor \(K\) 群} \(K_n^M(F)\) 是分次环 \(K_*^M(F)\) 的 \(n\) 次部分，由

\[K_*^M(F) = T(F^\times) / \langle a \otimes (1-a) : a \in F \setminus \{0, 1\} \rangle\]

定义，其中 \(T(F^\times)\) 是乘法群 \(F^\times\) 的张量代数。\(\{a_1, \ldots, a_n\}\) 记 \(a_1 \otimes \cdots \otimes a_n\) 在 \(K_n^M(F)\) 中的像（Milnor 符号）。

\textbf{命题 199.1}（Milnor \(K\) 群的低阶）： - \(K_0^M(F) = \mathbb{Z}\) - \(K_1^M(F) = F^\times\) - \(K_2^M(F)\) 与 Quillen \(K_2(F)\) 一致（Matsumoto 定理）

\textbf{证明}：\(K_0^M(F) = \mathbb{Z}\)：由定义 \(K_0^M(F)\) 是张量代数 \(T(F^\times)\) 的 \(0\) 次部分，张量积 \(0\) 次为基域 \(k\) 在 \(T^0\) 中的拷贝，无任何 Steinberg 关系涉及 \(0\) 次元，故 \(K_0^M(F) = \mathbb{Z}\)。\(K_1^M(F) = F^\times\)：\(T^1(F^\times) = F^\times\)，Steinberg 关系 \(a \otimes (1-a)\) 在 \(1\) 次为零，故 \(K_1^M(F) = F^\times\)。\(K_2^M(F) \cong K_2(F)\)：Matsumoto 定理给出 \(K_2(F) \cong (F^\times \otimes_{\mathbb{Z}} F^\times)/\langle a \otimes (1-a) \rangle\)，恰为 Milnor \(K_2^M(F)\) 的定义。证明核心为 Steinberg 群 \(\text{St}(F)\) 到 \(K_2^M(F)\) 的符号映射，验证其核恰为 Steinberg 关系的 Milnor 版本。 \(\blacksquare\)

\textbf{定理 199.1}（Milnor \(K\) 理论与 Galois 上同调，Bloch-Kato 猜想 / Voevodsky 定理 2011）：对域 \(F\) 和素数 \(\ell \neq \operatorname{char}(F)\)，存在自然同构

\[K_n^M(F) / \ell \cong H^n_{\text{Gal}}(F, \mu_\ell^{\otimes n})\]

（Galois 上同调与 Milnor \(K\) 理论的联系）。这是 Milnor 猜想（\(\ell = 2\)）和 Bloch-Kato 猜想（一般 \(\ell\)）的融合，由 Voevodsky 证明（2011 年 Fields 奖）。

\textbf{证明}：Voevodsky 在母题稳定同伦范畴中将 Milnor \(K\) 群实现为母题上同调 \(K_n^M(F) \cong H^{n,n}_{\text{mot}}(F, \mathbb{Z})\)。关键为证明母题上同调到 etale 上同调的规范映射 \(H^{n,n}_{\text{mot}}(F, \mathbb{Z}/\ell) \to H^n_{\text{et}}(F, \mu_\ell^{\otimes n})\) 为同构。此即 Beilinson-Lichtenbaum 猜想，Voevodsky 通过引入 Rost 的范数剩余母题并证明其消没定理（消没猜想的证明）而解决。对 \(\ell = 2\)，Voevodsky (1996) 完成证明；一般 \(\ell\) 由 Voevodsky 和 Rost 在 2011 年前完成。 \(\blacksquare\)

\subsection{199.2 Bloch 群与稀释子}\label{bloch-ux7fa4ux4e0eux7a00ux91caux5b50}

\textbf{定义 199.2}（Bloch 群）：域 \(F\) 的 \textbf{Bloch 群} \(\mathcal{B}(F)\) 是由 \([x]\)（\(x \in F \setminus \{0, 1\}\)）生成的 Abel 群，模去关系

\[[x] - [y] + \left[\frac{y}{x}\right] - \left[\frac{1 - x^{-1}}{1 - y^{-1}}\right] + \left[\frac{1 - x}{1 - y}\right] = 0\]

（五项关系，five-term relation）。

\textbf{定理 199.2}（Bloch-Suslin 定理，1986）：对无限域 \(F\)，\(K_3(F)_{\text{ind}} \cong \mathcal{B}(F)\)（\(K_3\) 的不分解部分）。Bloch 群是 \(K_3\) 的精确代数描述。

\textbf{证明}：\(K_3(F)_{\text{ind}}\) 定义为 \(K_3(F)\) 模去分解部分（由乘积映射 \(K_1(F) \otimes K_2(F) \to K_3(F)\) 的像）。Bloch 定义的映射 \(\beta: \mathcal{B}(F) \to K_3(F)_{\text{ind}}\) 将 \([x]\) 映为符号 \(\{x, 1-x\}\) 的某种提升。五项关系恰对应 \(K\) 理论中的 Steinberg 关系。Suslin 证明了 \(\beta\) 的满射性：每个 \(K_3(F)_{\text{ind}}\) 中元素可表为符号的组合；Bloch 定义倒映射 \(K_3(F)_{\text{ind}} \to \mathcal{B}(F)\) 通过稀释子实现，利用 Suslin 的刚性定理验证其为同构。 \(\blacksquare\)

\textbf{定义 199.3}（稀释子，Dilogarithm）：\textbf{Bloch-Wigner 稀释子} \(D: \mathbb{C} \setminus \{0, 1\} \to \mathbb{R}\) 定义为

\[D(z) = \operatorname{Im}(\operatorname{Li}_2(z)) + \log|z| \cdot \arg(1 - z)\]

\(D(z)\) 是单值的且在 \(\mathbb{C}P^1\) 上连续延拓。\(D(z)\) 满足五项关系，因此诱导了 \(\mathcal{B}(\mathbb{C}) \to \mathbb{R}\) 的同态。

\textbf{定理 199.3}（Bloch 群与双对数恒等式）：Bloch 群的结构由稀释子（dilogarithm）和经典的双对数恒等式（Euler, Abel, Kummer）控制。\(\mathcal{B}(\mathbb{C})\) 的秩为不可数无限。

\textbf{证明}：Bloch-Wigner 稀释子 \(D(z)\) 满足五项关系，故诱导满同态 \(\mathcal{B}(\mathbb{C}) \twoheadrightarrow \mathbb{R}\)。同时，将 \(\mathbb{C}\) 视为 \(\mathbb{Q}\) 上的无穷维向量空间，可构造无穷多 \(\mathbb{Q}\)-线性无关的元素 \([x_\alpha] \in \mathcal{B}(\mathbb{C})\)（\(\alpha\) 遍历超越基底）。\(D\) 在这些元上的取值由 \(\operatorname{Li}_2\) 的泛函方程（Euler-Abel-Kummer 恒等式）给出有理关系，但不影响实线性无关性。由 Zagier 猜想（现已证），\(\mathcal{B}(\mathbb{C})\) 作为 \(\mathbb{Q}\)-向量空间的秩为不可数无限。 \(\blacksquare\)

\subsection{199.3 调节子映射}\label{ux8c03ux8282ux5b50ux6620ux5c04}

\textbf{定义 199.4}（Borel 调节子）：Borel 定义了从 \(K_{2n-1}(\mathbb{C})\) 到 \(\mathbb{R}\) 的调节子映射，其像与 \(\zeta(n)\) 相关。更一般地，对数域 \(F\)，

\[K_{2n-1}(\mathcal{O}_F) \otimes \mathbb{Q} \cong \mathbb{Q}^{d_n}\]

其中 \(d_n = \dim_{\mathbb{Q}} H^1_{\text{mot}}(F, \mathbb{Q}(n))\)。

\textbf{定理 199.4}（Beilinson 猜想，1984）：特殊 \(L\) 值 \(L^*(M, n)\) 与母题 \(M\) 的调节子映射的像（在 \(K\) 理论中）有联系。这是代数 \(K\) 理论与数论 \(L\) 函数之间最深刻的猜想。

\textbf{证明}：Beilinson 构造了调节子映射 \(r_{\mathcal{D}}: K_{2n-1}(X)_{\mathbb{Q}} \to H_{\mathcal{D}}^{2n-1}(X_{\mathbb{R}}, \mathbb{R}(n))\)（Deligne-Beilinson 上同调），断言 \(r_{\mathcal{D}}\) 的像是 \(H_{\mathcal{D}}^{2n-1}\) 的有理结构的一部分，且其秩等于 \(L(h^{2n-1}(X), s)\) 在 \(s = n\) 处的零点阶数。Beilinson 还猜测 \(r_{\mathcal{D}}\) 的像与特殊 \(L\) 值 \(L^*(h^{2n-1}(X), n)\) 通过周期关系精确联系。对 \(X = \operatorname{Spec}(\mathcal{O}_F)\)（数域），Borel 证明了 \(K_{2n-1}(\mathcal{O}_F)\) 的秩由 \(\zeta_F(n)\) 的零点阶数给出，验证了 Beilinson 猜想的最简情形。 \(\blacksquare\)

\hrulefill

\hrulefill

\hrulefill

\hrulefill

\hrulefill

\hrulefill

\section{\texorpdfstring{第246章：代数 \(K\) 理论与母题}{第246章：代数 K 理论与母题}}\label{ux7b2c246ux7ae0ux4ee3ux6570-k-ux7406ux8bbaux4e0eux6bcdux9898}

母题（motives）是 Grothendieck 在 1960 年代构思的终极统一理论------其宏伟愿景是为代数簇的所有已知上同调理论（Betti、de Rham、etale、crystalline 等）构建一个''泛上同调''，使得每个具体的上同调理论都通过''实现函子''从这个普遍对象中派生。这一构想的核心类比是：正如有限群的所有线性表示都由其不可约特征标完全决定，代数簇的所有上同调信息都应被某种''母题分解''所编码。Grothendieck 的标准猜想（Lefschetz 型标准猜想和 Hodge 型标准猜想）若成立，将赋予母题范畴一个半单的 Tannakian 结构。Voevodsky 在 1990 年代的革命性工作（导致 2002 年菲尔兹奖）在更为广义的框架下（\(\mathbb{A}^1\)-同伦理论）建立了母题上同调 \(H^*_{\mathcal{M}}(-, \mathbb{Z}(n))\) 与代数 K 理论的谱序列联系：对于光滑代数簇 \(X\)，存在一条从母题上同调到代数 K 理论的收敛谱序列 \(E_2^{p,q} = H_{\mathcal{M}}^{p-q}(X, \mathbb{Z}(-q)) \Rightarrow K_{-p-q}(X)\)。这一联系表明：母题上同调是代数 K 理论的分次版本，从而将 Grothendieck 的 Riemann-Roch 哲学（第183章 \(K_0\)）和 Milnor 的 K 理论（第184章）统一在同一个母题框架之下。

\subsection{200.1 母题范畴}\label{ux6bcdux9898ux8303ux7574}

\textbf{定义 200.1}（母题，Grothendieck）：\textbf{母题}（motive）是代数簇的''普遍上同调''对象。对光滑射影代数簇 \(X\)，其母题 \(h(X)\) 是 Abel 范畴 \(\mathcal{M}(k)\)（母题范畴）中的对象，满足对所有 Weil 上同调理论 \(H^*\) 存在实现函子 \(r_H: \mathcal{M}(k) \to \text{Vec}\) 使得 \(r_H(h(X)) = H^*(X)\)。

\textbf{定义 200.2}（Chow 母题）：\textbf{Chow 母题}范畴 \(\mathcal{M}_{\text{rat}}(k)\) 由光滑射影代数簇的对应（correspondences）构造。对象是 \((X, p, n)\)（\(X\) 是簇，\(p\) 是幂等对应，\(n\) 是 Tate 扭转）。其中： - \textbf{幂等对应} \(p \in \operatorname{CH}^{\dim X}(X \times X)_{\mathbb{Q}}\) 满足 \(p \circ p = p\)（对应合成的意义下），作用相当于在 \(X\) 的上同调中取出某个直和分量------形式上 \(p\) 将母题 \(h(X)\) 投影到其''子母题''。 - \textbf{Tate 扭转} \(n \in \mathbb{Z}\)：对象 \((X, p, n)\) 的''母题上同调''相对于 \((X, p, 0)\) 有一个维数位移。Tate 母题 \(\mathbb{Q}(n)\) 是 \(\operatorname{Spec}(k)\) 的母题经过 \(n\) 次 Tate 扭转，其 Betti 实现给出 \((2\pi i)^n \mathbb{Q}\)，反映了母题范畴中的''维数权重''结构。Tate 扭转是定义母题范畴中 Hom 群的分次结构以及联系母题上同调与代数 \(K\) 理论谱序列的关键成分。

\textbf{定理 200.1}（Jannsen 定理，1992）：Chow 母题范畴是半单 Abel 范畴当且仅当数值等价与有理等价一致（这是 Grothendieck 标准猜想之一）。

\textbf{证明}：Chow 母题范畴 \(\mathcal{M}_{\text{rat}}(k)_{\mathbb{Q}}\) 的对象间 Hom 群通过有理等价模去对应定义。数值等价是比有理等价的更粗的等价关系。若两者一致，则 Jannsen 证明了 \(\mathcal{M}_{\text{rat}}(k)_{\mathbb{Q}}\) 中所有态射的核与余核均分裂（因为数值等价为 0 蕴含有理等价为 0 意味着零态射）。由范畴论的一般原理，任何对象 \(M\) 的自同态代数 \(\operatorname{End}(M)_{\mathbb{Q}}\) 变为有限维半单代数，从而 \(\mathcal{M}_{\text{rat}}(k)_{\mathbb{Q}}\) 为半单 Abel 范畴。反之，若范畴半单，其上的正合函子（上同调实现）必须识别两种等价关系，故它们一致。 \(\blacksquare\)

\subsection{200.2 Voevodsky 的母题上同调}\label{voevodsky-ux7684ux6bcdux9898ux4e0aux540cux8c03}

\textbf{定义 200.3}（母题上同调，Voevodsky 2000）：光滑代数簇 \(X\) 的\textbf{母题上同调} \(H^{p,q}_{\text{mot}}(X, \mathbb{Z})\) 定义为 \(X\) 在母题稳定同伦范畴 \(\mathcal{SH}(k)\) 中的母题 Eilenberg-MacLane 谱 \(H\mathbb{Z}\) 的上同调。

\textbf{定理 200.2}（母题上同调与代数 \(K\) 理论，Voevodsky 2003）：对光滑代数簇 \(X\)，

\[H^{p,q}_{\text{mot}}(X, \mathbb{Z}) \cong \operatorname{CH}^q(X, 2q-p)\]

（Bloch 高阶 Chow 群）。母题上同调给出了 Bloch 高阶 Chow 群的上同调解释。

\textbf{证明}：Voevodsky 将 \(H\mathbb{Z}\) 母题谱定义为母题稳定同伦范畴 \(\mathcal{SH}(k)\) 中的 Eilenberg-MacLane 谱，其同伦层由有限对应（finite correspondences）的层复形给出。对光滑簇 \(X\)，\(H^{p,q}_{\text{mot}}(X, \mathbb{Z}) = \operatorname{Hom}_{\mathcal{SH}(k)}(\Sigma^\infty X_+, \Sigma^{p,q} H\mathbb{Z})\)。另一方面，Bloch 高阶 Chow 群 \(\operatorname{CH}^q(X, n)\) 由 \(X\) 上余维数 \(q\) 的代数圈在 \(\Delta^n\) 上的族模去边界定义。Friedlander-Suslin-Voevodsky 证明了母题复形 \(\mathbb{Z}(q)\)（由 \(\mathbb{A}^q\)-同伦不变量层的链复形给出）的同调同伦群恰为 Bloch 高阶 Chow 群，从而建立了同构。 \(\blacksquare\)

\textbf{定理 200.3}（母题谱序列，Friedlander-Suslin-Voevodsky）：存在从母题上同调到代数 \(K\) 理论的谱序列：

\[E_2^{p,q} = H^{p-q,-q}_{\text{mot}}(X, \mathbb{Z}) \Rightarrow K_{-p-q}(X)\]

这统一了代数 \(K\) 理论和母题上同调。

\textbf{证明}：代数 \(K\) 理论谱 \(KGL\) 在母题稳定同伦范畴 \(\mathcal{SH}(k)\) 中具有切片过滤（slice filtration）\(F^\bullet KGL\)。切片 \(s_q(KGL) \cong \Sigma^{2q,q} H\mathbb{Z}\)（Voevodsky 的关键定理）。此切片塔诱导的谱序列在 \(E_2\) 页的项为母题上同调 \(H^{p-q,-q}_{\text{mot}}(X, \mathbb{Z})\)，收敛于 \(KGL^{p+q}(X) = K_{-p-q}(X)\)。这完整统一了母题上同调与代数 \(K\) 理论。 \(\blacksquare\)

\subsection{200.3 Voevodsky 的 Fields 奖工作}\label{voevodsky-ux7684-fields-ux5956ux5de5ux4f5c}

Voevodsky 的 Fields 奖（2002）主要表彰他在母题上同调理论的开创性工作以及 Milnor 猜想的证明。他的核心贡献在于建立了一个全新的母题稳定同伦范畴 \(\mathcal{SH}(k)\)，在其中代数簇的母题可以被视为''谱''并研究其上同调运算。这一框架将代数几何中的经典问题------如 Milnor K 理论、Galois 上同调、Bloch-Kato 猜想------转化为母题稳定同伦论中的计算问题。Voevodsky 的母题上同调运算理论（特别是母题 Steenrod 代数的构造）使他能够证明 Milnor 猜想（\(K_n^M(k)/2 \cong H_{\text{ét}}^n(k, \mathbb{Z}/2)\)），这是母题方法在经典问题上的首次重大胜利。随后，他与 Rost 合作将这一方法推广到模 \(\ell\) 情形，于 2011 年完成 Bloch-Kato 猜想的证明------这一结果被公认为 21 世纪初代数几何和数论领域最重大的突破之一。Voevodsky 的工作深刻改变了代数 K 理论和数论的面貌：他证明了母题上同调不仅是代数 K 理论的''影子''，而且通过母题谱序列（定理 200.3）与 K 理论直接等价，从而将代数 K 理论的计算问题转化为母题上同调的计算问题。

\textbf{定理 200.4}（Milnor 猜想 / Voevodsky 定理，1996-2003）：对域 \(k\)（\(\operatorname{char}(k) \neq 2\)），Milnor \(K\) 理论模 2 与 étale 上同调之间存在自然同构：

\[K_n^M(k) / 2 \cong H^n_{\text{ét}}(k, \mathbb{Z}/2\mathbb{Z})\]

\textbf{证明}：这等价于定理 198.3 的 etale 上同调版本（因 \(H^n_{\text{ét}}(k, \mathbb{Z}/2\mathbb{Z}) \cong H^n_{\text{Gal}}(k, \mu_2^{\otimes n})\)，对特征不为 2 的域）。Voevodsky 在母题稳定同伦论中引入母题 Steenrod 代数，证明母题上同调模 2 与 Milnor \(K\) 理论模 2 同构。同时利用范数剩余母题的消没定理证明母题上同调到 etale 上同调的映射为同构，结合两者即得。 \(\blacksquare\)

\textbf{定理 200.5}（Bloch-Kato 猜想，已证明，Voevodsky-Rost 2011）：该猜想等价于定理 199.1（Bloch-Kato 猜想的 étale 上同调版本），两者表述不同但实质相同。详见第 199 章。

\textbf{证明}：同定理 199.1 的证明。Bloch-Kato 猜想的一般版本断言 \(K_n^M(k)/\ell \cong H^n_{\text{ét}}(k, \mu_\ell^{\otimes n})\)。Voevodsky-Rost 通过证明范数剩余母题的消没定理完成了一般 \(\ell\) 的证明，核心为 Rost 的 chain lemma 和 Voevodsky 的母题上同调运算理论。 \(\blacksquare\)

\subsection{\texorpdfstring{200.4 代数 \(K\) 理论与 \(L\) 函数}{200.4 代数 K 理论与 L 函数}}\label{ux4ee3ux6570-k-ux7406ux8bbaux4e0e-l-ux51fdux6570}

代数 \(K\) 理论与 \(L\) 函数的联系是现代数论中最深层的猜想之一，由 Lichtenbaum（1973）和 Beilinson（1984）分别独立提出。这些猜想断言：代数簇的 \(K\) 群的秩和特殊 \(L\) 值由调节子映射（regulator map）精确控制，从而将代数 \(K\) 理论、母题上同调和数论 \(L\) 函数统一在一个宏大的框架中。Lichtenbaum 猜想的独特之处在于，它预测了 \(K\) 群与 Hasse-Weil \(L\) 函数在整数点处零点阶数之间的精确关系------这一关系已经被 Borel 在数域情形下完全证明。Borel 的证明使用对称空间的局部化理论，将 \(K_{2n-1}(\mathcal{O}_F)\) 与 Dedekind zeta 函数 \(\zeta_F(s)\) 在 \(s=n\) 处的零点阶数联系起来。对于一般的光滑射影簇，Beilinson 猜想进一步断言调节子映射 \(r_{\mathcal{D}}: K_{2n-1}(X)_{\mathbb{Q}}^{(n)} \to H_{\mathcal{D}}^{2n-1}(X_{\mathbb{R}}, \mathbb{R}(n))\) 的像与特殊 \(L\) 值 \(L^*(h^{2n-1}(X), n)\) 之间通过周期矩阵精确关联。这一猜想已在椭圆曲线、模曲线和 Shimura 簇等特殊情形下得到部分验证，但一般情形仍远未解决。它代表了代数 \(K\) 理论、代数几何和数论三者交汇的最前沿，是 Langlands 纲领和 Grothendieck 母题哲学在 \(K\) 理论中的具体体现。

\textbf{定理 200.6}（Lichtenbaum 猜想，Beilinson 猜想）：对光滑射影代数簇 \(X\) 在整数环 \(\mathcal{O}_F\) 上的模型，\(K\) 群的秩由 Hasse-Weil \(L\) 函数在整点处的零点阶数给出：

\[\dim_{\mathbb{Q}} K_{2n-1}(X) \otimes \mathbb{Q} = \operatorname{ord}_{s=n} L(h^{2n-1}(X), s)\]

这推广了 Borel 关于数域 \(K\) 群的结果。

\textbf{证明}：Borel 对 \(X = \operatorname{Spec}(\mathcal{O}_F)\) 证明了该结果的秩部分：他计算了 \(K_{2n-1}(\mathcal{O}_F) \otimes \mathbb{Q}\) 的维数，并证明其等于 Dedekind zeta 函数 \(\zeta_F(s)\) 在 \(s = n\) 处的零点阶数（当 \(n \geq 2\)）。一般情形的正合陈述属于 Beilinson 猜想：调节子映射 \(r_{\mathcal{D}}: K_{2n-1}(X)_{\mathbb{Q}}^{(n)} \to H_{\mathcal{D}}^{2n-1}(X_{\mathbb{R}}, \mathbb{R}(n))\) 的像是极大秩的，且 \(\dim K_{2n-1}(X)_{\mathbb{Q}}^{(n)} = \operatorname{ord}_{s=n} L(h^{2n-1}(X), s)\)。该猜想在部分情形（如数域、椭圆曲线）已得验证。 \(\blacksquare\)

\hrulefill

\hrulefill

\hrulefill

\hrulefill

\hrulefill

\hrulefill

\section{\texorpdfstring{第247章：\(K\) 理论的其他方向与应用}{第247章：K 理论的其他方向与应用}}\label{ux7b2c247ux7ae0k-ux7406ux8bbaux7684ux5176ux4ed6ux65b9ux5411ux4e0eux5e94ux7528}

代数 K 理论的影响力远超其代数根源------它已成为沟通拓扑、几何、算子代数和数学物理的普适语言。在拓扑 K 理论的方向上，Atiyah 和 Singer 在 1963 年将 \(K^0(X)\) 与椭圆算子的 Fredholm 指标联系起来，证明了著名的 Atiyah-Singer 指标定理：椭圆微分算子的解析指标（其核与余核维数之差）等于其符号在拓扑 K 理论中定义的拓扑指标。这一结果不仅是 20 世纪数学的顶峰之一，还直接催生了非交换几何（Connes，见第170-172章）和整个几何分析学派的发展。在算子代数方向，Cuntz 和 Quillen 建立了算子代数 K 理论（\(KK\)-理论），成为 Kasparov 分类程序的基础。在数论方向，Lichtenbaum 和 Quillen 的猜想将代数 K 群与 Dedekind zeta 函数在整数点的值联系起来，而 Bloch-Kato 猜想（由 Voevodsky 和 Rost 证明）则是 Milnor K 理论的完整 Galois 上同调刻画。在数学物理中，Karoubi 的乘性 K 理论和 Freed-Hopkins 对可逆拓扑场论的分类充分展示了 K 理论在量子场论中的基础性角色。本章是对全书 K 理论三章（第183-186章）的总结与前瞻。代数 \(K\) 理论在拓扑、几何、算子代数和物理学中有广泛的应用。

\subsection{\texorpdfstring{201.1 拓扑 \(K\) 理论与指标定理}{201.1 拓扑 K 理论与指标定理}}\label{ux62d3ux6251-k-ux7406ux8bbaux4e0eux6307ux6807ux5b9aux7406}

\textbf{定义 201.1}（拓扑 \(K\) 理论，Atiyah-Hirzebruch 1959）：紧 Hausdorff 空间 \(X\) 的拓扑 \(K\) 理论 \(K^0(X)\) 是 \(X\) 上复向量丛的 Grothendieck 群。\(K^{-n}(X) = \tilde{K}^0(S^n X)\)（Bott 周期性：\(K^{-n-2}(X) \cong K^{-n}(X)\)）。

\textbf{定理 201.1}（Atiyah-Singer 指标定理，1963）：椭圆算子 \(D\) 的解析指标等于拓扑指标：

\[\operatorname{Index}(D) = \int_M \operatorname{Ch}(\sigma(D)) \wedge \operatorname{Td}(M)\]

其中 \(\operatorname{Ch}\) 是 Chern 特征标（\(K\) 理论到上同调的映射），\(\operatorname{Td}\) 是 Todd 类。

\textbf{证明}：同定理 81.11，这是 Atiyah-Singer 指标定理以 Todd 类表述的版本（对应 Dolbeault 算子情形）。\(\operatorname{Index}(D) = \langle \operatorname{Ch}(\sigma(D)) \wedge \operatorname{Td}(TM), [M] \rangle\)。关键步骤为将 \(K(T^*M)\) 中 \(\sigma(D)\) 的拓扑指标约化为上同调积分。利用嵌入 \(M \hookrightarrow \mathbb{R}^N\) 及其法丛的 Thom 同构，拓扑指标映射 \(K(T^*M) \to \mathbb{Z}\) 等同于上同调中的 \(\langle \operatorname{Ch}(-) \wedge \operatorname{Td}(TM \otimes \mathbb{C}), [TM] \rangle\)。 \(\blacksquare\)

\textbf{定理 201.2}（Swan 定理，1962）：\(C(X)\) 上有限生成投射模的范畴等价于 \(X\) 上复向量丛的范畴。因此 \(K_0(C(X)) \cong K^0(X)\)（代数 \(K\) 理论与拓扑 \(K\) 理论的联系）。

\textbf{证明}：同定理 196.1 的拓扑版本。向量丛 \(E \to X\) 的截面空间 \(\Gamma(E)\) 为有限生成投射 \(C(X)\)-模（紧 Hausdorff \(X\) 上由 Swan 原始论证）。反向构造：投射 \(C(X)\)-模 \(P \subset C(X)^n\) 在每点 \(x \in X\) 处的纤维 \(P_x = \{s(x) : s \in P\} \subset \mathbb{C}^n\) 为向量空间，构成 \(X\) 上的向量丛。两函子互为拟逆，给出范畴等价，继而诱导 \(K\) 群的同构。 \(\blacksquare\)

\subsection{\texorpdfstring{201.2 C*-代数 \(K\) 理论}{201.2 C*-代数 K 理论}}\label{c-ux4ee3ux6570-k-ux7406ux8bba}

\textbf{定义 201.2}（C\emph{-代数 \(K\) 理论）：C}-代数 \(A\) 的 \(K_0(A)\) 定义为 \(A\) 中投影的 Murray-von Neumann 等价类的 Grothendieck 群。\(K_1(A)\) 定义为 \(A\) 中酉元的连通分支群。

\textbf{定理 201.3}（六项正合列，\(K\) 理论的基本正合列）：对 C*-代数的短正合列 \(0 \to J \to A \to A/J \to 0\)，存在 \(K\) 理论的六项正合列：

\[\begin{CD}
K_0(J) @>>> K_0(A) @>>> K_0(A/J) \\
@AAA @. @VVV \\
K_1(A/J) @<<< K_1(A) @<<< K_1(J)
\end{CD}\]

\textbf{证明}：短正合列 \(0 \to J \to A \to A/J \to 0\) 诱导映射锥的正合列，在 \(K\) 理论中产生连接同态 \(\partial: K_1(A/J) \to K_0(J)\) 和 \(\partial: K_0(A/J) \to K_1(J)\)。连接同态通过悬挂和 Bott 周期性构造：将 C*-代数的短正合列与 \(C_0(\mathbb{R})\) 的张量积结合起来，利用 Toeplitz 代数的标准扩张。周期性 \(K_0(SA) \cong K_1(A)\) 和 \(K_1(SA) \cong K_0(A)\) 将六项正合列卷成循环。具体连接映射为指数映射：\(\partial([u]) = [e] - [e_0]\)（\(u \in A/J\) 的酉提升分解为投影差）。 \(\blacksquare\)

\textbf{定理 201.4}（Bott 周期性，C\emph{-代数情形）：\(K_0(SA) \cong K_1(A)\)，\(K_1(SA) \cong K_0(A)\)（其中 \(SA = C_0(\mathbb{R}) \otimes A\) 是悬挂）。这给出了 C}-代数 \(K\) 理论的 2-周期性。

\textbf{证明}：\(K_0(SA)\) 由 \(SA\) 中投影分类，等价于 \(M_n(A)\) 中满足 \(f(0) = f(\infty) = 0\) 的投影值连续函数 \(f: S^1 \to M_n(A)^\sim\)，这恰给出 \(M_n(A)^\sim\) 中酉元的同伦类，即 \(K_1(A)\) 的元素。\(K_1(SA)\) 同构于 \(K_0(A)\) 由 Bott 映射 \(S^1 \wedge A \to A \otimes \mathcal{K}\) 诱导，利用 Bott 投影的显式构造：\(b(z) = \begin{pmatrix} |z|^2 & z\sqrt{1-|z|^2} \\ \bar{z}\sqrt{1-|z|^2} & 1-|z|^2 \end{pmatrix}\) 在单位圆盘边界给出零投影。 \(\blacksquare\)

\subsection{\texorpdfstring{201.3 \(K\) 理论在物理学中的应用}{201.3 K 理论在物理学中的应用}}\label{k-ux7406ux8bbaux5728ux7269ux7406ux5b66ux4e2dux7684ux5e94ux7528}

K 理论在物理学中的应用是 20 世纪末和 21 世纪初最令人瞩目的跨学科进展之一。在凝聚态物理中，拓扑绝缘体（topological insulators）的拓扑相由 K 理论分类------这一发现源于对自由费米子系统的对称性分析。具体而言，动量空间中的绝缘体 Hamilton 量定义了从 Brillouin 区到对称类 Grassmann 流形的连续映射，其同伦分类恰由复 K 理论 \(K^{-d}(\mathbb{T}^d)\) 或实 K 理论 \(KO^{-d}(\mathbb{T}^d)\) 给出。Kitaev 的''十重分类''（2009）将 Altland-Zirnbauer 的十种对称类与 K 理论的 Bott 周期表精确对应，这一分类在拓扑量子计算中对 Majorana 零模和任意子的理解至关重要。在弦论中，D-膜的 Ramond-Ramond 荷由 K 理论分类（而非整数上同调）------这一由 Witten（1998）提出的观点揭示了 D-膜在 K 理论中的稳定性：两个 D-膜可以通过 tachyon 凝聚湮灭当且仅当它们的 K 理论类相同。磁单极子、瞬子（instanton）和反常消除（anomaly cancellation）等物理现象也自然地在 K 理论框架中得到解释。例如，杨-Mills 理论的瞬子数由第二 Chern 类 \(c_2\) 的积分给出，而更一般的拓扑荷则由 K 理论中的 Chern 特征标与上同调配对描述。K 理论还为量子场论中的''指标定理''提供了统一的数学语言，使得 Dirac 算子、Dolbeault 算子和符号差算子的指标可以统一地表示为 K 理论中的拓扑不变量。

\textbf{定义 201.3}（拓扑绝缘体与 \(K\) 理论）：拓扑绝缘体（topological insulator）的拓扑相由 \(K\) 理论分类。时间反演对称的拓扑绝缘体由 \(KR\) 理论（实 \(K\) 理论带对合）分类。

\textbf{定理 201.5}（Kitaev 的周期表，2009）：自由费米子系统的拓扑相由 \(K\) 理论的 Bott 周期间接分类（``十重分类''）。\(K\) 理论是理解拓扑序和拓扑量子计算的数学基础。

\textbf{证明}：Kitaev 考虑了自由费米子 Hamilton 量在时间反演 \(T\)（满足 \(T^2 = \pm 1\)）和粒子-空穴对称 \(C\)（满足 \(C^2 = \pm 1\)）下的十重分类。在动量空间中，\(d\) 维拓扑绝缘体对应于从 Brillouin 区（\(\mathbb{T}^d\)）到某个对称类 Grassmann 流形的映射。该映射的同伦分类由实/复 \(K\) 理论给出：复类（A, AIII）由 \(K^{-d}(\mathbb{T}^d)\) 分类，实类（AI, BDI, D, DIII, AII, CII, C, CI）由 \(KR^{-d}(\mathbb{T}^d)\) 或 \(KO^{-d}(\mathbb{T}^d)\) 分类。Bott 周期表中的 2（复）、8（实）周期性直接来自 \(K\) 理论的 Bott 周期性。 \(\blacksquare\)

\textbf{定义 201.4}（\(K\) 理论与 D-膜）：弦论中 D-膜的荷由 \(K\) 理论分类（而非整数上同调）。\(K^0(X)\) 分类 IIB 理论中 D-膜的荷，\(K^1(X)\) 分类 IIA 理论中 D-膜的荷。

\subsection{\texorpdfstring{201.4 Waldhausen \(K\) 理论与代数 \(K\) 理论空间的构造}{201.4 Waldhausen K 理论与代数 K 理论空间的构造}}\label{waldhausen-k-ux7406ux8bbaux4e0eux4ee3ux6570-k-ux7406ux8bbaux7a7aux95f4ux7684ux6784ux9020}

\textbf{定义 201.5}（Waldhausen \(K\) 理论，1985）：Waldhausen 将 Quillen 的 Q-构造推广到带有余纤维化（cofibration）和弱等价（weak equivalence）的范畴（Waldhausen 范畴）。Waldhausen \(K\) 理论是代数 \(K\) 理论、代数几何 \(K\) 理论和等变稳定同伦论中最重要的统一框架。

\textbf{定理 201.6}（Waldhausen 的加性定理和纤维化定理）：Waldhausen \(K\) 理论满足对 Waldhausen 范畴的分裂正合列和纤维化序列的加性行为，推广了 Quillen 的局部化定理。

\textbf{证明}：加性定理：若 Waldhausen 范畴 \(\mathcal{C}\) 是 \(\mathcal{A}\) 和 \(\mathcal{B}\) 的直和（等价于分裂正合列 \(0 \to \mathcal{A} \to \mathcal{C} \to \mathcal{B} \to 0\)），则诱导同伦等价 \(K(\mathcal{C}) \simeq K(\mathcal{A}) \times K(\mathcal{B})\)。证明通过 Waldhausen 的 \(S_\bullet\) 构造将其转化为单纯范畴的纤维化序列。纤维化定理：对 Waldhausen 范畴的恰当序列 \(\mathcal{A} \to \mathcal{C} \to \mathcal{B}\)（\(\mathcal{A}\) 中弱等价为 \(\mathcal{C}\) 中落在 \(\mathcal{B}\) 的弱等价原像的范畴），有同伦纤维列 \(K(\mathcal{A}) \to K(\mathcal{C}) \to K(\mathcal{B})\)。此推广了 Quillen 局部化定理（\(\mathcal{A}\) 为 \(\mathcal{B}\)-零调对象）。 \(\blacksquare\)

\hrulefill

\hrulefill

\hrulefill

\hrulefill

\hrulefill

\hrulefill

\subsection{187.A Quillen的+构造与高阶K-理论}\label{a-quillenux7684ux6784ux9020ux4e0eux9ad8ux9636k-ux7406ux8bba}

D. Quillen (1973) 的 + 构造将 K-理论从 \(K_0\) 和 \(K_1\) 推广到所有 \(n \geq 0\)，是代数 K-理论成为独立学科的标志（Quillen 因此获 1978 年菲尔兹奖）。

\textbf{定义 201.1}（Quillen + 构造）：设 \(X\) 为连通 CW 复形，\(N \trianglelefteq \pi_1(X)\) 为完全子群。则存在典范映射 \(f: X \to X_N^+\) 满足：（1）\(\pi_1(f)\) 诱导 \(\pi_1(X) / N\) 的同构；（2）对所有 \(\mathbb{Z}[\pi_1(X)/N]\)-模 \(M\) 及所有 \(n \geq 1\)，\(f\) 诱导同调同构 \(H_n(X; f^* M) \cong H_n(X_N^+; M)\)。对环 \(R\)，定义 \(K_n(R) := \pi_n(B \operatorname{GL}(R)^+)\)（\(n \geq 1\)）其中 \(B \operatorname{GL}(R)\) 是无穷一般线性群 \(\operatorname{GL}(R) = \varinjlim_n \operatorname{GL}_n(R)\) 的分类空间。

\textbf{定理 201.1}（Bass-Heller-Swan, Quillen）：\(K_0(R)\)（Grothendieck 群）、\(K_1(R)\)（\(\operatorname{GL}(R) / [\operatorname{GL}(R), \operatorname{GL}(R)]\)）与上述定义一致。\(K_2(R)\)（Milnor K-理论）与 \(H_2(E(R); \mathbb{Z})\)（初等矩阵群 \(E(R)\) 的整系数二阶同调）一致。

\textbf{证明}：\(K_0: \pi_0(B\operatorname{GL}(R)^+) \cong \mathbb{Z} \times B\operatorname{GL}(R)^+\) 的 \(\pi_0\) 同构于 Grothendieck 群（通过 \([P] - [Q]\) 与 \(B\operatorname{GL}(R)\) 的路连通分支对应）。\(K_1: \pi_1(B\operatorname{GL}(R)^+) \cong \operatorname{GL}(R) / E(R)\) 因为 \(+\) 构造将完全子群 \(E(R) = [\operatorname{GL}(R), \operatorname{GL}(R)]\) 切除。\(K_2: \pi_2(B\operatorname{GL}(R)^+) \cong H_2(B\operatorname{GL}(R)^+) \cong H_2(E(R))\) 由 Hurewicz 同构（因 \(B\operatorname{GL}(R)^+\) 为单连通）和同调切除。Matsumoto 定理给出 \(H_2(E(R)) \cong K_2(R)\)。 \(\blacksquare\)

\textbf{定理 201.2}（Quillen 的基本定理与 Adams 操作）： （1）局部化定理：对环 \(R\) 和乘性子集 \(S\)，有长正合序列 \(K_{n+1}(R_S) \to K_n(\mathcal{H}) \to K_n(R) \to K_n(R_S) \to \cdots\)，其中 \(\mathcal{H}\) 是有限生成 \(S\)-挠 \(R\)-模范畴的完全子范畴。（2）\textbf{Adams 操作} \(\psi^k: K_n(R) \to K_n(R)\) 由外幂构造诱导，满足 \(\psi^k \circ \psi^\ell = \psi^{k\ell}\) 且 \(\psi^p \equiv \text{Frob}^* \pmod{p}\)（在特征 \(p\) 上与 Frobenius 一致）。这些操作是 Lichtenbaum-Quillen 猜想（后发展为 Voevodsky-Rost 的 Bloch-Kato 猜想的证明）的核心代数工具。

\textbf{证明}：（1）局部化定理：仿紧空间范畴的纤维化 \(B\operatorname{GL}(\mathcal{H}) \to B\operatorname{GL}(R) \to B\operatorname{GL}(R_S)\)（Quillen 局部化定理）诱导同伦纤维列。取 \(+\) 构造后得到 \(B\operatorname{GL}(\mathcal{H})^+ \to B\operatorname{GL}(R)^+ \to B\operatorname{GL}(R_S)^+\) 为同伦纤维列，其同伦群长正合列即为所述 \(K\) 群长正合列。（2）Adams 操作：外幂 \(\lambda^k: \operatorname{GL}_n(R) \to \operatorname{GL}_{\binom{n}{k}}(R)\) 诱导 \(B\operatorname{GL}(R)^+\) 的映射，在同伦群上给出 \(\psi^k: K_n(R) \to K_n(R)\)。外幂的乘法性质给出 \(\psi^k \circ \psi^\ell = \psi^{k\ell}\)。在特征 \(p\) 上与 Frobenius 的一致性来自 Grothendieck 的 \(\lambda\)-环公理中特征 \(p\) 的 Adams-Riemann-Roch 定理。 \(\blacksquare\)

\hrulefill

\emph{本卷建立了代数 \(K\) 理论的完整框架：\(K_0\)（投射模的 Grothendieck 群）建立了将正合列转化为群关系的代数框架；\(K_1\)（一般线性群的 Abel 化）通过 Whitehead 引理揭示了初等矩阵与换位子的深层联系；\(K_2\)（Steinberg 群与泛中心扩张）通过 Matsumoto 定理将域上的 \(K_2\) 与 Milnor 符号连接；Milnor \(K\) 理论与 Bloch 群（符号、稀释子、调节子）将 \(K\) 理论与数论 \(L\) 函数连接；母题上同调（Voevodsky 的母题稳定同伦论、Bloch-Kato 猜想）是 21 世纪数学的里程碑；拓扑 \(K\) 理论、C}-代数 \(K\) 理论和物理学应用展示了 \(K\) 理论的跨领域影响力。代数 \(K\) 理论是当代数学中连接代数、拓扑、几何、数论和物理的统一语言。*

\hrulefill

\hrulefill

\hrulefill

\hrulefill

\hrulefill

\hrulefill

\subsection{稳定同伦论与色展理论}\label{ux7a33ux5b9aux540cux4f26ux8bbaux4e0eux8272ux5c55ux7406ux8bba}

稳定同伦论是代数拓扑的核心分支，研究在悬挂操作下稳定化的同伦现象。其基本对象是\textbf{谱}（spectrum）------它推广了拓扑空间的概念，允许出现负维数的同伦群。Freudenthal 悬挂定理（1937）揭示了同伦群在高维悬挂下的稳定性，而谱的引入使得稳定同伦范畴 \(\mathcal{SH}\) 成为具有良好性质的三角范畴，其中存在球面谱 \(\mathbb{S}\) 作为基本对象，稳定同伦群 \(\pi_n^S = \pi_n(\mathbb{S})\) 是代数拓扑中最核心的不变量之一。色展理论（chromatic homotopy theory）由 Quillen、Morava、Ravenel 和 Hopkins 等人在 1970-90 年代发展，将稳定同伦范畴按''色层''（chromatic level）进行分层------通过 Morava \(K\)-理论和复配边 \(MU\) 的 Quillen 分解，将稳定同伦论与形式群理论（Honda 形式群律）紧密联系起来。Ravenel 猜想（后由 Devinatz-Hopkins-Smith 证明）的色展滤过定理和厚子范畴定理是这一理论的顶峰，它们给出了稳定同伦范畴的有限谱的完整分类。

\subsection{谱的正式定义}\label{ux8c31ux7684ux6b63ux5f0fux5b9aux4e49}

\textbf{定义 1}（序列谱）：一个\textbf{序列谱}（sequential spectrum）\(E\) 是一族带基点拓扑空间 \(\{E_n\}_{n \geq 0}\) 连同结构映射 \[\sigma_n: \Sigma E_n \longrightarrow E_{n+1},\] 其中 \(\Sigma\) 为约化悬挂。

若每个结构映射的伴随 \(\tilde{\sigma}_n: E_n \to \Omega E_{n+1}\) 是弱同伦等价，则称 \(E\) 为 \textbf{\(\Omega\)-谱}。

序列谱 \(E\) 的第 \(n\) 个稳定同伦群定义为 \[\pi_n(E) = \operatorname{colim}_{k \to \infty} \pi_{n+k}(E_k),\] 其中极限沿映射 \(\pi_{n+k}(E_k) \to \pi_{n+k+1}(\Sigma E_k) \xrightarrow{\sigma_k} \pi_{n+k+1}(E_{k+1})\) 取。

\textbf{定理 1}（\(\Omega\)-谱的上同调表示）：若 \(E\) 为 \(\Omega\)-谱，则对任意 CW 复形 \(X\)， \[E^n(X) := [X, E_n] \cong \operatorname{colim}_k [\Sigma^k X, E_{n+k}]\] 定义了一个约化上同调理论。反之，Brown 可表性定理保证每个上同调理论由此种 \(\Omega\)-谱表示。

\textbf{证明}：验证 Eilenberg-Steenrod 公理：同伦不变性由 \([-, E_n]\) 的函子性直接得出。切除公理：对 CW 对 \((X,A)\)，\(X/A\) 的同伦型决定映射 \(X \to E_n\) 与 \(A \to E_n\) 的差。悬挂同构 \(\tilde{E}^n(\Sigma X) \cong \tilde{E}^{n-1}(X)\) 由伴随 \(\Sigma \dashv \Omega\) 及 \(\Omega\)-谱条件导出：\(E_{n-1} \simeq \Omega E_n\)。\(\blacksquare\)

\subsection{稳定同伦范畴的正式定义}\label{ux7a33ux5b9aux540cux4f26ux8303ux7574ux7684ux6b63ux5f0fux5b9aux4e49}

\textbf{定义 2}（稳定同伦范畴 \(\mathcal{SH}\)）：以 \(\mathbf{Sp}\) 记序列谱范畴（态射为谱映射），在其上赋予稳定模型结构（stable model structure）。\textbf{稳定同伦范畴} \(\mathcal{SH}\) 定义为 \(\mathbf{Sp}\) 的同伦范畴： \[\mathcal{SH} = \operatorname{Ho}(\mathbf{Sp}).\]

等价地，\(\mathcal{SH}\) 可视为带基点拓扑空间范畴 \(\mathbf{Top}_*\) 在 \(\Sigma\) 上的 Bousfield 局部化------即将 \(\Sigma\) 变为范畴等价后得到的同伦范畴。具体而言，考虑序列 \[\mathbf{Top}_* \xrightarrow{\Sigma} \mathbf{Top}_* \xrightarrow{\Sigma} \cdots\] 的同伦余极限，其对象为谱。\(\mathcal{SH}\) 中，悬挂函子 \(\Sigma\) 是可逆的（其逆为环路函子 \(\Omega\)），故 \(\mathcal{SH}\) 为三角范畴，其平移函子 \([1]\) 恰为 \(\Sigma\)。

\subsection{经典谱的正式定义}\label{ux7ecfux5178ux8c31ux7684ux6b63ux5f0fux5b9aux4e49}

\textbf{定义 3}（球谱）：\textbf{球谱} \(\mathbb{S}\) 定义为 \[\mathbb{S}_n = S^n, \quad \sigma_n: \Sigma S^n \xrightarrow{\cong} S^{n+1}.\] \(\mathbb{S}\) 是 \(\mathcal{SH}\) 的幺半单位：\(E \wedge \mathbb{S} \cong E \cong \mathbb{S} \wedge E\)（在适当 smash 积模型下）。

\textbf{定义 4}（Eilenberg-MacLane 谱）：对 Abel 群 \(A\)，\textbf{Eilenberg-MacLane 谱} \(HA\) 由 \[(HA)_n = K(A, n)\] 定义，其中 \(K(A,n)\) 为 Eilenberg-MacLane 空间（\(\pi_n = A\)，其他同伦群消失）。结构映射由标准等价 \(\Sigma K(A,n) \to K(A,n+1)\) 给出。\(HA\) 是 \(\Omega\)-谱，对应经典的上同调理论 \(H^*(-; A)\)。

\textbf{定义 5}（拓扑 K-理论谱）：\textbf{复 K-理论谱} \(KU\) 定义为 \[KU_n = \begin{cases} \mathbb{Z} \times BU, & n \text{ 为偶数} \\ U, & n \text{ 为奇数} \end{cases}\] 配备 Bott 周期性同伦等价 \(\Omega^2(\mathbb{Z} \times BU) \simeq \mathbb{Z} \times BU\) 和 \(\Omega U \simeq \mathbb{Z} \times BU\)。类似地，\textbf{实 K-理论谱} \(KO\) 以 \(BO\) 为基础，具有周期性 8：\(\Omega^8 KO \simeq KO\)。

\textbf{定义 6}（配边谱）：\textbf{复配边谱} \(MU\) 定义为 \[MU_n = MU(n) \quad \text{（万有 Thom 复形）},\] 其中 \(MU(n)\) 是万有 \(n\) 维复向量丛的 Thom 空间。结构映射 \(\Sigma MU(n) \to MU(n+1)\) 由万有丛的 Whitney 和诱导。

\subsection{\texorpdfstring{\(\mathcal{SH}\) 的对称幺半结构}{\textbackslash mathcal\{SH\} 的对称幺半结构}}\label{mathcalsh-ux7684ux5bf9ux79f0ux5e7aux534aux7ed3ux6784}

\textbf{定理 2}（Smash 积的构造困难与解决）：序列谱范畴中，朴素定义的 smash 积不满足严格结合律。主流解决方案如下：

\begin{enumerate}
\def\labelenumi{\arabic{enumi}.}
\tightlist
\item
  \textbf{对称谱}（Hovey-Shipley-Smith）：在对称群 \(\Sigma_n \times \Sigma_m \to \Sigma_{n+m}\) 作用下的谱，smash 积自然地对称幺半。
\item
  \textbf{正交谱}（Mandell-May-Schwede-Shipley）：以正交群 \(O(n)\)-空间为 \(n\) 次分量，smash 积由 \(O(n) \times O(m) \hookrightarrow O(n+m)\) 诱导。
\item
  \textbf{\(\Gamma\)-空间}（Segal）：谱的特殊 \(\Gamma\)-空间模型，smash 积通过 \(E_\infty\)-结构构造。
\item
  \textbf{Lydakis 的 smash 积模型}：在序列谱上直接构造结合且对称的 smash 积（通过变形收缩技术）。
\end{enumerate}

所有上述模型产生等价于 \(\mathcal{SH}\) 的对称幺半范畴。记 smash 积为 \(\wedge\)，则 \((\mathcal{SH}, \wedge, \mathbb{S})\) 为闭对称幺半范畴，内 Hom 为函子谱 \(F(E,F)\)，满足伴随 \(\operatorname{Map}(E \wedge F, G) \cong \operatorname{Map}(E, F(F,G))\)。

\textbf{证明}：对称谱模型中，\((E \wedge F)_n = \bigvee_{p+q=n} \Sigma_n \times_{\Sigma_p \times \Sigma_q} (E_p \wedge F_q)\)，结构映射由 smash 积的悬挂交换诱导。结合律由对称群的作用协调保证，对称性由块置换实现。正交谱模型以 \(O(n)\) 替代 \(\Sigma_n\)，利用 \(O(p) \times O(q) \hookrightarrow O(p+q)\) 给出同样的幺半结构。所有模型通过 Quillen 等价互相关联，故 \(\mathcal{SH}\) 中的幺半结构是典范的。 \(\blacksquare\)

\subsection{色展理论核心}\label{ux8272ux5c55ux7406ux8bbaux6838ux5fc3}

色展同伦论（Chromatic Homotopy Theory）将 \(\mathcal{SH}\) 按素数 \(p\) 及形式群定律的高度分解为互不干扰的''色谱层''。

\textbf{定义 7}（MU 的同调与形式群）：复配边谱 \(MU\) 的系数环为 \[MU_* = \mathbb{Z}[x_1, x_2, \ldots], \quad |x_i| = 2i.\] \(MU\) 上存在一个典范的形式群定律 \(\mathbb{F}\)，其对数级数为 \(\sum_{n \geq 0} [\mathbb{C}P^n] t^{n+1}/(n+1)\)。

\textbf{定理 3}（Quillen）：\(MU\) 上的形式群定律是万有的：对任意交换环 \(R\) 上的形式群定律 \(F\)，存在唯一环同态 \(\phi: MU_* \to R\) 使得 \(F = \phi_* \mathbb{F}\)。等价地，Lazard 环 \(L\)（形式群定律的分类环）典范同构于 \(MU_*\)。

\textbf{证明}：Lazard 定理表明 \(L \cong \mathbb{Z}[t_1, t_2, \ldots]\)，\(|t_i| = 2i\)。Quillen 构造了一个映射 \(MU_* \to L\)，通过对陈类与形式群系数建立对应，验证其在每个次数上为同构。关键步骤：利用 Milnor 的 \(MU_*(MU)\) 计算与形式群定律严格同构环的关系，证明 Thom 谱的同伦论完美编码了形式群定律的代数。\(\blacksquare\)

\textbf{定义 8}（Morava K-理论）：固定素数 \(p\) 及整数 \(n \geq 1\)。第 \(n\) 个 \textbf{Morava K-理论谱} \(K(n)\) 的系数环为 \[K(n)_* = \mathbb{F}_p[v_n^{\pm 1}], \quad |v_n| = 2(p^n - 1).\] 其中 \(v_n\) 对应 \(p\)-典型形式群定律在高度 \(n\) 处的参数。\(K(n)\) 满足 Kunneth 同构：\(K(n)_*(X \times Y) \cong K(n)_*(X) \otimes_{K(n)_*} K(n)_*(Y)\)。

\textbf{定义 9}（色展过滤与 Bousfield 局部化）：给定素数 \(p\)，定义 \(p\)-局部谱范畴 \(\mathcal{SH}_{(p)}\) 上的 Bousfield 局部化函子 \(L_n\)： \[L_0 = L_{H\mathbb{Q}} \quad \text{（有理化）},\] \[L_n = L_{E(n)} \quad \text{（对 Johnson-Wilson 理论 $E(n)$ 的局部化，$n \geq 1$）}.\] 这里 \(E(n)_* = \mathbb{Z}_{(p)}[v_1,\ldots,v_{n-1},v_n^{\pm 1}]\)，\(|v_i| = 2(p^i-1)\)。Bousfield 类之间满足 \[L_n > L_{n-1} > \cdots > L_0,\] 并且存在典范变换 \(L_n \to L_{n-1}\)，其同伦纤维由 Morava K-理论 \(K(n)\)-局部化 \(L_{K(n)}\) 刻画。Bousfield 局部化的层结构称为\textbf{色展过滤}。

\textbf{定义 10}（Height）：\(p\) 处形式群定律的\textbf{高度}（height）\(n\) 是指其 \(p\)-级数 \([p]_F(x)\) 的首项阶为 \(p^n\)。高度严格为正，且是 Morava \(K(n)\)-局部化分类同伦层的基本不变量。

\subsection{Ravenel 猜想}\label{ravenel-ux731cux60f3}

\textbf{猜想 1}（幂零猜想，已证 ------ Devinatz-Hopkins-Smith）：设 \(f: \Sigma^k X \to X\) 为有限谱的自映射。若 \(MU_*(f) = 0\)（即 \(f\) 在 MU-同调中诱导零映射），则 \(f\) 是幂零的：存在 \(m\) 使 \(f^{\circ m} \simeq 0\)。

\textbf{定理 4}（Devinatz-Hopkins-Smith，幂零猜想的证明概要）：核心思想是通过 \(MU\) 的 \(v_n\)-周期性地层分解。对每个 \(n\)，构造 Morava K-理论探测函子：若 \(K(n)_*(f) = 0\) 对所有 \(n\) 成立，则 \(f\) 幂零。证明通过引入 \(p\)-典型 \(v_n\)-映射到环谱的同伦论中，利用类 nilpotence 定理的稳定性论证，最终归约为 \(MU_*\) 上形式群的分歧性质。

\textbf{证明}：设 \(f\) 非幂零，则存在素理想 \(\mathfrak{p} \subset MU_*\) 使 \(MU_*(f)_{\mathfrak{p}} \neq 0\)。由 Quillen 定理，\(\mathfrak{p}\) 对应某个高度 \(n\) 的形式群定律。利用 Morava 变理论（change-of-rings），\(MU_*(f)\) 之局部化信息由 \(E(n)_*(f)\) 捕获。进一步，\(E(n)\) 局部化在 \(K(n)\) 局部化中分裂：\(L_{E(n)}X \simeq \operatorname{holim} L_{K(n)} \cdots L_{K(0)}X\)。若 \(K(n)_*(f) = 0\) 对所有 \(n\)，则在塔的每一步 \(f\) 为零，由收敛性（定理 12）知 \(f\) 幂零。Devatz-Hopkins-Smith 通过构造 \(v_n\)-自映射严格实现了此论证。 \(\blacksquare\)

\textbf{猜想 2}（望远镜猜想，部分公开）：\(v_n\)-望远镜 \(T(n)\)（定义为 \(S^0/p, v_1, \ldots, v_{n-1}\) 的映射望远镜）的 Bousfield 类等于 \(K(n)\) 的 Bousfield 类。此猜想对所有已知质数的高度 \(n=0,1\) 成立（\(n=0\) 为 Miller 定理，\(n=1\) 为 Mahowald 定理），但一般情形仍为现代同伦论最核心的未决问题之一。

\subsection{色展红移现象}\label{ux8272ux5c55ux7ea2ux79fbux73b0ux8c61}

\textbf{定理 5}（Ausoni-Rognes 红移）：代数 K-理论函子 \(K\) 在色展层级上产生红移： \[K(KU_{(p)}) \simeq KU_{(p)} \vee \Sigma^{-3} KU_{(p)} \quad \text{（高度 1）},\] 更一般地，猜想对高度 \(n\) 的 \(E_\infty\)-环谱 \(R\)，\(K(R)\) 的色展信息集中在高度 \(n+1\)。此即\textbf{红移猜想}：\(K\) 将高度提升 1。

\textbf{证明}：Dundas-Goodwillie-McCarthy 定理给出同伦纤维列 \(K(KU) \to \operatorname{TC}(KU) \to \Sigma^{-1} \operatorname{THH}(KU)^{hS^1}\)，其中 \(\operatorname{TC}\) 为拓扑循环同调。计算 \(\operatorname{THH}(KU) \simeq KU \otimes S^1_+ \simeq KU \vee \Sigma KU\)（Bokstedt 定理）。\(\Sigma KU\) 在 \(S^1\) 循环算子作用下经 Frobenius 和限制映射给出 \(\operatorname{TC}(KU)\) 的完整描述。在有理同伦层面，迹映射 \(K(KU)_{\mathbb{Q}} \to \operatorname{TC}(KU)_{\mathbb{Q}}\) 为同构，两者均被识别为 \(KU_{\mathbb{Q}}\) 的特定 wedge。计算得 \(K(KU_{(p)})\) 在 \(p\)-完备下为 \(\Sigma\) 周期性 \(KU\) 的 wedge，验证了高度 1 到高度 2 的红移。 \(\blacksquare\)

\subsection{色展塔与分层细化}\label{ux8272ux5c55ux5854ux4e0eux5206ux5c42ux7ec6ux5316}

\textbf{定理 6}（色展塔，Hopkins-Miller-Lurie）：在每个质数 \(p\) 和高度 \(n \geq 1\) 处，存在 Morava E-理论谱 \(E_n\)（Lubin-Tate 理论），它是 \(E_\infty\)-环谱，其系数环为 \[(E_n)_* = W(\overline{\mathbb{F}}_p)[[u_1, \ldots, u_{n-1}]][u^{\pm 1}], \quad |u| = 2, \ |u_i| = 0,\] 且 \(\pi_0(E_n)\) 参数化 \(\overline{\mathbb{F}}_p\) 上高度 \(n\) 的形式群定律的万有形变空间。\(E_n\) 带有 Morava 稳定化群 \(\mathbb{G}_n = \operatorname{Aut}(\mathbb{F}_{p^n}) \rtimes \operatorname{Gal}(\mathbb{F}_{p^n}/\mathbb{F}_p)\) 的典范作用，使得固定点谱 \(E_n^{h\mathbb{G}_n}\) 精确地恢复第 \(n\) 个 Lubin-Tate 谱，并与球面的同伦群在高高度处有深层的计算联系。

\textbf{证明}：由 Lubin-Tate 形变理论，\(\overline{\mathbb{F}}_p\) 上高度 \(n\) 的形式群定律存在万有形变环 \(LT_n = W(\overline{\mathbb{F}}_p)[[u_1, \ldots, u_{n-1}]]\)。Goerss-Hopkins-Miller 证明存在唯一的 \(E_\infty\)-环谱 \(E_n\) 以 \(LT_n\) 为系数环，其形式群为万有形变。\(\mathbb{G}_n = \operatorname{Aut}(\mathbb{F}_{p^n}) \rtimes \operatorname{Gal}(\mathbb{F}_{p^n}/\mathbb{F}_p)\) 作用于 \(E_n\)，其同伦不动点谱 \(E_n^{h\mathbb{G}_n}\) 恰为 Lubin-Tate 谱在第 \(n\) 层的固定点。色展塔 \(\{L_n \mathbb{S}_{(p)}\}\) 的分层由这些固定点谱通过 descent 谱序列组织。 \(\blacksquare\)

\textbf{定理 7}（Ravenel 的消没线定理）：对任意有限谱 \(X\)，存在某个 \(n\) 使得 \(\operatorname{Ext}\) 图（Adams-Novikov 谱序列的 \(E_2\) 页）在斜率大于某固定常数时消没。此定理对同伦群计算的收敛性提供了精确控制，并完成了色展过滤在有限谱上的全面覆盖。

\textbf{证明}：Adams-Novikov \(E_2\) 项为 \(\operatorname{Ext}_{MU_*(MU)}(MU_*, MU_*(X))\)。由 Morava 变理论，\(MU_*(MU)\) 在 \(p\) 处的结构由形式群定律的自同构控制。给定有限谱 \(X\)，其 \(MU_*\)-同调为有限生成 \(MU_*\)-模，故存在 \(N\) 使 \(v_n^{-1}MU_*(X) = 0\) 对所有 \(n > N\)。在色展谱序列中，这对应''消没线''：\(\operatorname{Ext}\) 项在斜率 \(> 1/(p^N-1)\) 处消没。形式上，利用 Miller-Ravenel 的色展谱序列 \(E_1^{n,*} \Rightarrow \operatorname{Ext}_{MU_*(MU)}(MU_*, MU_*(X))\)，\(v_n\)-周期性层的贡献在 \(E_2\) 页具有特定斜率范围，有限谱的有限性保证了高斜率处无贡献。 \(\blacksquare\)

\subsection{色展理论与母题同伦论}\label{ux8272ux5c55ux7406ux8bbaux4e0eux6bcdux9898ux540cux4f26ux8bba}

\textbf{定理 8}（色展红移与母题稳定同伦论的关系）：在母题稳定同伦论中，代数 K-理论谱 \(KGL\) 的色展红移行为由 Voevodsky 母题范畴的切片过滤（slice filtration）控制。\(KGL\) 的切片在高度方面的行为与拓扑 K-理论谱 \(KU\) 在经典色展过滤中的行为平行，这为代数几何与拓扑的深层统一提供了数论变体的色展视角。

\textbf{证明}：Voevodsky 的切片过滤 \(F^\bullet KGL\) 满足 \(s_q(KGL) \cong \Sigma^{2q,q} H\mathbb{Z}\)（母题 Eilenberg-MacLane 谱）。在母题稳定同伦范畴 \(\mathcal{SH}(k)\) 中，\(KGL\) 的 \(p\)-局部色展过滤由母题 Morava K-理论 \(K(n)^{\text{mot}}\) 组织。切片塔诱导的谱序列显示 \(KGL\) 的 \(v_n\)-周期性切片对应于拓扑 \(KU\) 在经典色展过滤中高度 \(n+1\) 的表现，给出了代数-拓扑红移平行的母题解释。此对应由 \(\mathbb{C}\) 上的 Betti 实现函子 \(\mathcal{SH}(\mathbb{C}) \to \mathcal{SH}\) 精确建立。 \(\blacksquare\)

\subsection{稳定同伦论的前沿方向}\label{ux7a33ux5b9aux540cux4f26ux8bbaux7684ux524dux6cbfux65b9ux5411}

\textbf{定理 9}（Goerss-Hopkins-Miller 定理）：在 Morava E-理论 \(E_n\) 上，存在唯一的 \(E_\infty\)-结构，且由该 \(E_\infty\)-结构构成的模空间是可缩的。等价地，\(E_n\) 作为 Lubin-Tate 理论的自同构群 \(\mathbb{G}_n\) 连续作用于此 \(E_\infty\)-环谱上，其同伦不动点 \(E_n^{h\mathbb{G}_n}\) 的分类问题与球面同伦群在色展塔高层的计算密切相关。

\textbf{证明}：Goerss-Hopkins 建立了 \(E_\infty\)-环谱的障碍理论。给定形式模的 Lubin-Tate 塔，其在谱层面的严格实现对应于 \(E_\infty\) 操作在 \(E_n\) 的上同调上的作用。障碍群位于幂运算（power operations）的导出函子中，对 \(E_n\) 情形，因 \(E_n^*E_n\) 的代数结构足够丰富，所有障碍消没。唯一性源于同伦范畴中的刚性：任何两个 \(E_\infty\)-实现之间的同伦论等价可提升到结构化等价。可缩模空间的结论来自障碍理论的参数化版本------Lubin-Tate 形变空间的刚性使每个切方向无阻碍。 \(\blacksquare\)

\textbf{定理 10}（色展层与 \(v_n\)-周期性）：对素数 \(p\) 和高度 \(n\)，存在小范畴 \(\mathcal{F}_n\)（类型 \(n\) 的有限谱），使得其经 Bousfield 局部化 \(L_n^f\) 后的 \(K(n)\)-局部同伦范畴可由对偶对和 Morava 稳定化群的群上同调完全分类。此为色展理论的''有限性定理''------高度顶层的计算本质上是 Morava 稳定化群在特定余模上的群上同调，同构于球面稳定同伦群在色展过滤最高层的贡献。

\textbf{证明}：Bousfield 局部化 \(L_n^f\) 在有限谱上的作用为：\(L_n^f X\) 的 \(K(n)\)-局部同伦型由 \(X\) 的 \(K(n)_*E_n\)-余模结构完全确定。由 Hopkins 的类型 \(n\) 有限谱的厚子范畴定理，\(\mathcal{F}_n\) 恰为 Kohler 对偶的整对象范畴。色展开塔同伦极限的 \(n\) 阶差分 \(L_n X / L_{n-1}X\) 的同伦群恰好由 \(\mathbb{G}_n\) 在高 Morava 稳定同伦群 \(\pi_*(L_{K(n)}\mathbb{S})\) 上的群上同调作用计算。由 Devinatz-Hopkins 的连续群关于连续链复形的上同调理论，该群上同调与有限谱在同伦范畴中的态射形成严格的对偶性。 \(\blacksquare\)

\subsection{球面同伦群与色展计算}\label{ux7403ux9762ux540cux4f26ux7fa4ux4e0eux8272ux5c55ux8ba1ux7b97}

\textbf{定理 11}（Adams-Novikov 谱序列，ANSS）：基于 MU 的 Adams 谱序列是计算球面稳定同伦群的主要工具： \[E_2^{s,t} = \operatorname{Ext}_{MU_*(MU)}^{s,t}(MU_*, MU_*) \Longrightarrow \pi_{t-s}(\mathbb{S}_{(p)}).\] 其中 \(E_2\) 页在代数上为 \(\operatorname{Ext}\) 群，其色展过滤由 Miller-Ravenel 的色展谱序列进一步组织为按高度 \(n\) 分层的计算方案。

\textbf{证明}：取 \(\mathbb{S}_{(p)}\) 的 \(MU\)-Adams 分解 \(\mathbb{S}_{(p)} \to MU \to \overline{MU} \to \Sigma MU \to \cdots\)，其中 \(\overline{MU}\) 为 \(MU\) 的同伦余纤维。将此塔 smash 自身并取同伦极限得到 \(MU\)-Adams 塔，其极大特 (maximal) 的 \(E_1\) 页由 \(MU_*(MU^{\wedge s})\) 的同伦给出。因 \(MU_*(MU)\) 为平坦 Hopf 代数胚，\(E_2\) 页即其上的 \(\operatorname{Ext}\) 群。进一步，\(\operatorname{Ext}\) 群的色展过滤由 \(v_n\)-局部化塔 \(MU \to L_1 MU \to L_2 MU \to \cdots\) 的代数镜像给出，Miller-Ravenel 在此基础上证明收敛性。 \(\blacksquare\)

\textbf{定理 12}（色展收敛定理，Hopkins-Ravenel）：色展塔在有限谱上是收敛的：对任意有限谱 \(X\) 及任意质数 \(p\)，自然映射 \[X_{(p)} \longrightarrow \operatorname{holim}_n L_n X_{(p)}\] 是 \(p\)-完备化下的等价。色展层的逆极限完整地重建了 \(p\)-局部的稳定同伦论。此定理确保了色展分层方案的完备性------在极限意义上，所有有限谱的信息均被色展过滤完全捕获。

\textbf{证明}：对有限谱 \(X\) 及固定的 \(n\)，构造比较映射 \(X_{(p)} \to L_n X_{(p)}\)。由 Bousfield 局部化的性质，余纤维 \(C_n(X)\) 满足 \(C_n(X) \wedge T = *\) 对所有 \(L_n\)-零调谱 \(T\)（即类型 \(\geq n+1\) 的有限谱）。当 \(n \to \infty\) 时，余纤维的 \(p\)-局部化逐渐变为零：因有限谱在色展塔的高高度处贡献有限，其 \(K(n)\)-局部同伦群对所有足够大 \(n\) 为零。精确地，若 \(X\) 为类型 \(m\) 的有限谱，则 \(K(n)_*(X) = 0\) 当 \(n > m\)，从而 \(L_n X \simeq X\) 对所有 \(n \geq m\)，极限自然同构。 \(\blacksquare\)

色展理论的深远影响超出传统同伦论：其技术（形式群、Morava 理论、Bousfield 局部化）为导出代数几何（通过高度分类 \(E_\infty\)-环谱）、算术几何（通过 Lubin-Tate 空间与志村簇的关系）乃至母题同伦论（通过 \(KGL\) 的色展行为）提供了统一的语言与计算工具。

\newpage

\chapter{07-拓扑与几何/03-代数拓扑/04-微分拓扑.md}\label{ux62d3ux6251ux4e0eux51e0ux4f5503-ux4ee3ux6570ux62d3ux625104-ux5faeux5206ux62d3ux6251.md}

\textbf{第188章（补充）}

\begin{quote}
低维拓扑研究 3 维和 4 维流形的拓扑结构，是当代拓扑学中最活跃的前沿领域之一。纽结理论是 3 维拓扑的核心组成部分，通过代数和组合不变量（Jones 多项式、Khovanov 同调）来区分和分类纽结。Floer 同调为 3 维和 4 维流形提供了强大的不变量，而 Donaldson 理论和 Seiberg-Witten 理论揭示了 4 维流形独特的微分结构现象。本卷在拓扑学（V10）、代数拓扑（V14）和微分几何（V11）的基础上，系统介绍低维拓扑的核心理论。
\end{quote}

\hrulefill

\hrulefill

\hrulefill

\hrulefill

\hrulefill

\hrulefill

\textbf{等变同伦论与等变上同调}

等变同伦论研究带有紧 Lie 群 \(G\) 连续作用的拓扑空间的同伦不变量，是现代代数拓扑与几何拓扑的交汇点。其核心思想是：\(G\)-空间 \(X\) 的''正确''不变量不仅来自 \(X\) 本身的拓扑，还应编码 \(G\) 作用的轨道结构信息。Elmendorf 定理（1983）是等变同伦论的基石------它断言 \(G\)-空间的 \(G\)-同伦范畴等价于从轨道范畴 \(\mathcal{O}_G^{\mathrm{op}}\) 到拓扑空间的函子范畴的同伦范畴，从而将等变同伦论约化为对 \(G\) 的所有闭子群 \(H\) 的固定点空间 \(X^H\) 的单独研究。等变上同调理论（Bredon 上同调、等变 \(K\)-理论、等变配边理论）将经典上同调推广到 \(G\)-空间，其系数系统不再是单个 Abel 群，而是定义在轨道范畴上的函子。等变同伦论在指标定理（Atiyah-Segal 等变 \(K\)-理论指标定理）、几何拓扑（Mostow 刚性）和数学物理（Chern-Simons 理论、等变量子场论）中有重要应用。

\subsection{G-空间与 G-CW 复形}\label{g-ux7a7aux95f4ux4e0e-g-cw-ux590dux5f62}

\textbf{定义 1}（G-空间）：设 \(G\) 为紧 Lie 群。一个 \textbf{G-空间} 是拓扑空间 \(X\) 连同连续左作用 \(G \times X \to X\)，满足 \(e \cdot x = x\) 和 \(g \cdot (h \cdot x) = (gh) \cdot x\)。G-等变映射 \(f: X \to Y\) 满足 \(f(g \cdot x) = g \cdot f(x)\)。G-空间范畴记作 \(G\mathbf{Top}\)。

\textbf{定义 2}（G-CW 复形）：\(G\) 为紧 Lie 群。\textbf{G-CW 复形} 是 G-空间 \(X\)，配备递增滤过 \(\emptyset = X_{-1} \subset X_0 \subset X_1 \subset \cdots\)，使 \(X = \bigcup X_n\) 且每个 \(X_n\) 由 \(X_{n-1}\) 通过粘合 \(n\) 维 G-胞腔 \(G/H \times D^n\) 得到。精确地说，存在推出图表 \[\begin{CD}
\bigsqcup_{i} G/H_i \times S^{n-1} @>>> X_{n-1} \\
@VVV @VVV \\
\bigsqcup_{i} G/H_i \times D^n @>>> X_n
\end{CD}\] 其中 \(H_i \subset G\) 为闭子群。

G-CW 复形继承了 CW 复形的胞腔逼近定理与 Whitehead 定理：若 \(f: X \to Y\) 为 G-CW 复形间的 G-等变映射，且对每个闭子群 \(H \subset G\)，\(f^H: X^H \to Y^H\) 为弱同伦等价，则 \(f\) 为 G-同伦等价。

\subsection{Elmendorf 定理}\label{elmendorf-ux5b9aux7406}

\textbf{定义 3}（轨道范畴）：\(G\) 的\textbf{轨道范畴} \(\mathcal{O}_G\) 是以 \(G/H\)（\(H \subset G\) 为闭子群）为对象，以 G-等变映射为态射的范畴。等价地，\(\mathcal{O}_G\) 可视为传递 G-空间构成的满子范畴。

\textbf{定理 1}（Elmendorf）：存在函子 \[\Phi: G\mathbf{Top} \longrightarrow \operatorname{Fun}(\mathcal{O}_G^{\mathrm{op}}, \mathbf{Top}), \quad X \longmapsto (G/H \mapsto X^H),\] 其中 \(X^H = \{x \in X \mid h \cdot x = x \ \forall h \in H\}\) 为 \(H\)-固定点子空间。\(\Phi\) 在同伦范畴层面是范畴等价：G-空间的 G-同伦范畴与从 \(\mathcal{O}_G^{\mathrm{op}}\) 到拓扑空间的函子的同伦范畴等价。后者赋予投影模型结构（弱等价和纤维化为对象层面的）。

\textbf{证明}：(1) \(\Phi\) 保持同伦等价：若 \(f:X\to Y\) 为 \(G\)-同伦等价，则存在 \(g:Y\to X\) 使得 \(g\circ f\simeq_G\operatorname{id}\)。对每个闭子群 \(H\) 取 \(H\)-固定点得 \(g^H\circ f^H\simeq\operatorname{id}\)，故 \(\Phi(f)\) 在函子范畴中逐对象同伦等价。(2) 构造拟逆 \(\Psi\)：对函子 \(F:\mathcal{O}_G^{\text{op}}\to\mathbf{Top}\)，取 Bar 复形 \(B_n(F)=\bigsqcup_{G/H_0\to\cdots\to G/H_n}F(G/H_n)\)，几何实现 \(|B_\bullet(F)|\) 赋 \(G\)-作用 \(g\cdot[\alpha]=[g\cdot\alpha]\)。(3) \(\Phi\circ\Psi\simeq\operatorname{id}\)：\((\Psi(F))^H\simeq F(G/H)\)，自然同构由恒等态射的简并给出。(4) \(\Psi\circ\Phi\simeq\operatorname{id}\)：等变映射 \(\varepsilon:|B_\bullet(\Phi(X))|\to X\) 将态射链和 \(F\) 值映为作用于 \(x\) 的结果，在固定点上为弱等价。故 \(\Phi\) 为同伦范畴的等价。\(\blacksquare\)

\subsection{Borel 等变上同调}\label{borel-ux7b49ux53d8ux4e0aux540cux8c03}

\textbf{定义 4}（万有 G-丛的 Milnor 构造）：万有 G-主丛 \(EG \to BG\) 由 Milnor 联合构造定义：\(EG = G * G * G * \cdots\) 的有限联合归纳极限（赋予弱拓扑）。\(EG\) 是可缩 G-空间，且 G 自由作用于其上。分类空间 \(BG = EG/G\)。

\textbf{定义 5}（Borel 等变上同调）：对任意系数环 \(R\)，G-空间 \(X\) 的 \textbf{Borel 等变上同调} 定义为 \[H^*_G(X; R) := H^*(EG \times_G X; R),\] 其中 \(EG \times_G X = (EG \times X)/G\)（对角线作用）。Borel 构造是同伦商（homotopy quotient）\(X_{hG}\) 的模型：首先''以同伦自由作用替换''\(X\)，再取商。

\textbf{定理 2}（Borel 上同调的基本性质）：\(H^*_G(-; R)\) 满足同伦不变性、切除公理（对 G-CW 对）和 Mayer-Vietoris 序列。若 \(G\) 在 \(X\) 上自由作用，则 \[H^*_G(X; R) \cong H^*(X/G; R),\] 因为 \(EG \times_G X \simeq X/G\)。

\textbf{证明}：当 \(G\) 自由作用时，投射 \(EG \times_G X \to X/G\) 以 \(EG\)（可缩）为纤维，从而是同伦等价。所有公理由经典上同调的公理继承。\(\blacksquare\)

\subsection{Bredon 等变上同调}\label{bredon-ux7b49ux53d8ux4e0aux540cux8c03}

\textbf{定义 6}（系数系统）：一个 \textbf{Bredon 系数系统} 是函子 \[M: \mathcal{O}_G^{\mathrm{op}} \longrightarrow \mathbf{Ab},\] 即对每个 \(G/H\) 分配一个 Abel 群 \(M(G/H)\)，对每个 G-等变映射 \(G/H \to G/K\) 分配一个群同态 \(M(G/K) \to M(G/H)\)。

\textbf{定义 7}（Bredon 上链复形）：设 \(X\) 为 G-CW 复形，记 \(\underline{I}_n\) 为 \(X\) 的 \(n\) 维 G-胞腔的指标集。定义 \(n\) 维 Bredon 上链群为 \[C^n_G(X; M) = \prod_{G/H \times D^n \in X^{(n)}} M(G/H).\] 边缘同态 \(\delta: C^n_G(X; M) \to C^{n+1}_G(X; M)\) 由胞腔粘合映射通过 \(M\) 的函子性诱导。\textbf{Bredon 等变上同调} 是此上链复形的上同调： \[H^*_G(X; M) = H^*(C^\bullet_G(X; M)).\]

\textbf{定理 3}（Borel 与 Bredon 的比较）：若 \(G\) 在 \(X\) 上自由作用，则存在同构 \[H^*_G(X; R) \cong H^*_G(X; R)_{\mathrm{Bredon}},\] 其中右侧的 Bredon 系数系统为常数函子 \(M(G/H) = R\)（带有态射限制为恒等映射）。

一般情形（非自由作用）下，Bredon 上同调捕捉了 Borel 上同调无法区分的等变信息：例如对 \(G = \mathbb{Z}/2\)，非平凡 \(M\) 可区分两种对合（含固定点集与不含固定点集），而 Borel 上同调对任意同伦等价的 Borel 构造给出相同结果。

\textbf{证明}：当 \(G\) 自由作用时，\(EG \times_G X \simeq X/G\)，Borel 上同调即 \(H^*(X/G; R)\)。自由 G-CW 复形的胞腔链复形与 \(X/G\) 的普通 CW 链复形一致。常数 Mackey 函子 \(R\) 的 Bredon 上链复形在自由作用下退化为普通奇异上链复形。由同伦不变性，\(H^*(EG \times_G X; R) \cong H^*(X/G; R) \cong H^*_G(X; R)_{\text{Bredon}}\)。\(\blacksquare\)

\subsection{真正等变稳定同伦论（Genuine G-Spectra）}\label{ux771fux6b63ux7b49ux53d8ux7a33ux5b9aux540cux4f26ux8bbagenuine-g-spectra}

\textbf{定义 8}（G-表示球面）：设 \(V\) 为 \(G\) 的有限维实正交表示。\textbf{表示球面} \(S^V\) 是 \(V\) 的单点紧化，赋予诱导的 G-作用。当 \(V = \mathbb{R}^n\)（平凡作用）时，\(S^V \cong S^n\)（带平凡 G-作用）。

\textbf{定义 9}（真正 G-谱）：一个\textbf{真正 G-谱}（genuine G-spectrum）\(E\) 由一族 \(G\)-空间 \(\{E_V\}\)（以有限维 \(G\)-表示 \(V\) 为指标）和各含入 \(V \subset W\)（当 \(W-V\) 为 \(G\)-表示时）的结构映射 \[\Sigma^{W-V} E_V \longrightarrow E_W\] 组成，满足传递性条件，且需对所有表示 \(V\) 有规范等价 \(E_V \xrightarrow{\cong} \Omega^{W-V} E_W\) 的伴随版本。

关键区别：在''朴素''（naive）G-谱中，指标仅为平凡表示 \(\mathbb{R}^n\)，等价于普通谱带 G-作用。真正 G-谱的指标包含所有 \(G\)-表示 \(V\)，从而可构造以 \(S^V\) 为迹球的转移映射 \(\operatorname{tr}_H^G: E_V \to E_V\)（Mackey 结构），捕捉等变稳定同伦论的真正内涵。

\subsection{Mackey 函子与等变稳定同伦范畴}\label{mackey-ux51fdux5b50ux4e0eux7b49ux53d8ux7a33ux5b9aux540cux4f26ux8303ux7574}

\textbf{定义 10}（Mackey 函子）：\(G\) 上的 \textbf{Mackey 函子} \(M\) 是一对函子 \(M_*, M^*: \mathcal{O}_G \to \mathbf{Ab}\)（分别对共变和逆变），使得对象层面 \(M_*(G/H) = M^*(G/H) =: M(G/H)\)，且满足：

\begin{enumerate}
\def\labelenumi{\arabic{enumi}.}
\tightlist
\item
  （Mackey 双余积公式）对图表 \[\begin{CD}
  G/H @<<< G/K \\
  @VVV @VVV \\
  G/H' @<<< G/K'
  \end{CD}\] 有 \(M^*(\text{左下}) \circ M_*(\text{右上}) = M_*(\text{右下}) \circ M^*(\text{左上})\)（态射对应回溯-前推的一致性）。
\end{enumerate}

\textbf{定理 4}（Mackey 函子与等变稳定同伦）：真正 G-谱的 \(\pi_0\)（即稳定同伦群的等变推广）构成 Mackey 函子。更一般地，\(\pi_n^G(E)\) 对所有 \(n \in \mathbb{Z}\) 和组织子群构成 Mackey 函子族。等变稳定同伦范畴 \(\mathcal{SH}^G\) 的 Burnside 范畴诠释即为：\(\pi_0^G(\mathbb{S}_G)\) 同构于 Burnside 环 \(A(G)\)，它作为 Mackey 函子编码所有有限 G-集的同构类。

\textbf{证明}：对真正 G-谱 \(E\)，定义 \(\underline{\pi}_n(E)(G/H) = \pi_n^H(E)\)（\(H\)-固定点子谱的同伦群）。共变部分（转移）来自等变稳定同伦论的 Wirthmüller 同构：对闭子群 \(H \subset K\)，转移映射 \(\operatorname{tr}_H^K\) 由 \(G/K_+ \wedge E \to G/H_+ \wedge E\)（诱导自稳定同伦范畴中的对偶和迹映射）实现。Burnside 环 \(A(G)\) 的情形由 Segal 猜想（后由 Carlsson 证明）确认：\(\pi_0^G(\mathbb{S}_G) \otimes \hat{\mathbb{Z}} \cong A(G) \otimes \hat{\mathbb{Z}}\)。\(\blacksquare\)

\subsection{Atiyah-Segal 完备化定理}\label{atiyah-segal-ux5b8cux5907ux5316ux5b9aux7406}

Atiyah-Segal 完备化定理是等变 K 理论中最核心的结构定理之一，由 Atiyah 和 Segal 在 1969 年首次发表。该定理揭示了等变 K 理论 \(K_G^*(X)\) 与 Borel 构造的经典 K 理论 \(K^*(EG \times_G X)\) 之间的深刻关系：等变 K 理论在增广理想 \(I = \ker(K_G^*(pt) \to K^*(pt))\) 的 \(I\)-进完备化下同构于 Borel 构造的 K 理论。这一结果的哲学意义在于：等变 K 理论包含了比 Borel 构造的 K 理论更丰富的信息，而完备化过程恰好''遗忘''了这些额外的信息，仅仅保留''稳定''的等变数据。从几何角度看，\(K_G^*(X)\) 的参数取自 \(G\) 的有限维表示，而 \(K^*(EG \times_G X)\) 的参数取自 \(G\) 的所有表示（包括无限维表示），完备化过程从有限维表示过渡到''所有表示''的极限。Atiyah-Segal 定理的证明关键依赖于 \(EG\) 的有限维逼近 \(E_nG\)（如 Stiefel 流形或 Grassmann 流形），以及 Atiyah-Hirzebruch 谱序列的比较。完备化定理的一个重要推论是：当 \(G\) 在 \(X\) 上自由作用时，\(K_G^*(X) \cong K^*(X/G)\)（无需完备化），因为 \(EG \times_G X \simeq X/G\) 且增广理想在自由作用时为 \(0\)。另一个重要应用出现在 Atiyah-Segal 不动点定理中：完备化定理将等变 K 理论 \(K_G^*(X)\) 与不动点集 \(X^g\) 上的经典 K 理论联系起来，给出等变 K 理论的 Lefschetz 不动点公式。在更广泛的框架中，Atiyah-Segal 完备化定理被推广到等变椭圆上同调（Lurie 的导出代数几何框架）和等变 Morava K 理论，构成了现代等变稳定同伦论中''完备化''思想的源头。

\textbf{定理 5}（Atiyah-Segal 完备化定理）：设 \(G\) 为紧 Lie 群，\(X\) 为有限 G-CW 复形。记 \(I = \ker(K_G^*(pt) \to K(\text{pt}) = \mathbb{Z})\) 为增广理想。则在 \(I\)-进完备化下存在同构 \[K_G^*(X)_I^\wedge \cong K^*(EG \times_G X),\] 其中左侧为等变 K-理论的 \(I\)-进完备化，\(K^*(EG \times_G X)\) 为 Borel 构造的经典 K-理论。

\textbf{证明}：考虑 \(EG\) 的有限维逼近 \(E_nG\)（如 Stiefel 流形 \(V_n(\mathbb{R}^N)\)）。完备化 \(K_G^*(X)^\wedge_I=\varprojlim_n K_G^*(X)/I^nK_G^*(X)\)。映射 \(\varphi\) 由拉回 \(EG\times_GX\to EG\times_G*\simeq BG\) 在完备化上诱导。右侧 \(K^*(EG\times_GX)=\varinjlim_n K^*(E_nG\times_GX)\)。Atiyah-Hirzebruch 谱序列比较完备化极限与归纳极限，在有限型 \(G\)-空间上两者同构------因 \(I\)-进拓扑下 Cauchy 列在有限维截断收敛且 \(K\)-理论满足 Mittag-Leffler 条件。\(\blacksquare\)

\subsection{Green 函子与等变稳定同伦范畴的代数结构}\label{green-ux51fdux5b50ux4e0eux7b49ux53d8ux7a33ux5b9aux540cux4f26ux8303ux7574ux7684ux4ee3ux6570ux7ed3ux6784}

\textbf{定义 11}（Green 函子）：G-等变稳定同伦范畴 \(\mathcal{SH}^G\) 上的 \textbf{Green 函子} 是一个 Mackey 函子 \(M\)，附带额外的乘法结构：对每个 \(G/H\)，\(M(G/H)\) 是交换环，且转移映射 \(\operatorname{tr}_H^K\) 与限制映射 \(\operatorname{res}_H^K\) 及乘法兼容（满足 Frobenius 互反律）。典范的例子如等变 Burnside 环函子 \(\underline{A}_G\) 及等变稳定同伦群 \(\underline{\pi}_n^{\text{st}}\)。

\textbf{定理 6}（等变 Freudenthal 悬挂定理的等变推广）：对于真正 G-谱 \(E\)，若 \(G\) 是有限群，则对足够大的 \(n\)（依赖于 \(E\) 的连通性），映射 \[\pi_V^G(E) \to \pi_{V \oplus W}^G(\Sigma^W E)\] （其中 \(V, W\) 为 G-表示，\(W\) 包含平凡直和项）为同构。这推广了经典 Freudenthal 定理：等变情形的''足够连通''条件需要表示 \(W\) 包含所有不可约表示，而非仅依赖维数。

\textbf{证明}：设 \(V,W\) 为 \(G\) 的正交表示。对每个闭子群 \(H\)，经典 Freudenthal 定理应用于 \(S^{V^H}\)，稳定范围 \(i\le 2|V^H|-2\) 内 \(\pi_i(S^{V^H})\to\pi_{i+1}(S^{V^H+W^H})\) 为同构。取所有 \(H\) 范围下确界得等变稳定范围。等变 Blakers-Massey 切除定理确保悬挂 \(\Sigma^W:\pi_V^G(E)\to\pi_{V\oplus W}^G(\Sigma^WE)\) 在所示范围为同构。当 \(W\) 含所有不可约表示时等变 Freudenthal 范围与普通情形一致。\(\blacksquare\)

\subsection{等变上同调与全局反变层面}\label{ux7b49ux53d8ux4e0aux540cux8c03ux4e0eux5168ux5c40ux53cdux53d8ux5c42ux9762}

\textbf{定理 7}（tom Dieck 分裂定理）：对紧 Lie 群 \(G\) 和 G-空间 \(X\)，Borel 等变上同调的加法分解为 \[H^*_G(X; R) \cong \prod_{(H)} H^*(EW_G(H) \times_{W_G(H)} X^H; R),\] 其中积取遍 \(G\) 的共轭类闭子群 \((H)\)，\(W_G(H) = N_G(H)/H\) 为 Weyl 群。此分裂公式将 Borel 等变上同调完全化归为非等变数据，但代价是丢失了乘法结构的等变相干性。

\textbf{证明}：将 \(EG \times_G X\) 按轨道类型分层。对每个共轭类 \((H)\)，\(X^H\) 上的 \(W_G(H)\) 作用给出分片 \(EW_G(H) \times_{W_G(H)} X^H\)。Borel 构造的胞腔分解按此分片分裂，相应的上同调分裂为乘积（需验证在滤过的 \(E_\infty\) 页上无扩张问题）。严格论证使用等变上同调的谱序列和 tom Dieck 的滤过分裂引理。\(\blacksquare\)

\textbf{定理 8}（全局等变同伦论与超交换 Green 函子）：Schwede 的全局等变同伦论的等价刻画为：全局函子 \(\mathcal{SH}^{\text{global}}\) 的 Burnside 范畴核心等价于超交换 Green 函子构成的心范畴，且 \(\mathcal{SH}^{\text{global}}\) 为稳定、可表的对称幺半 \(\infty\)-范畴，其所有紧对象的子范畴生成整个范畴。这为统一处理所有紧 Lie 群的等变稳定同伦信息提供了系统框架。

\textbf{证明}：Schwede 构造全局等变稳定同伦范畴 \(\mathcal{SH}^{\text{global}}\)，其中对象用（有限群的）谱的正合范畴上的 Day 卷积实现。超交换 Green 函子构成的心范畴等价性来自 Mackey 函子的 Day 卷积和 Green 函子的乘法结构的标准分解。紧生成性由球谱 \(\mathbb{S}\) 及其（解）悬挂的所有有限 G-集的操作的稳定闭包给出。\(\blacksquare\)

\subsection{等变对偶与等变层论}\label{ux7b49ux53d8ux5bf9ux5076ux4e0eux7b49ux53d8ux5c42ux8bba}

\textbf{定理 9}（等变 Spanier-Whitehead 对偶）：对有限 G-CW 复形 \(X\)，存在 G-对偶谱 \(D(X)\)（等变 Spanier-Whitehead 对偶）使得对任意真正 G-谱 \(E\) 有自然同构 \[\tilde{E}^G_n(X) \cong \tilde{E}_G^{-n}(D(X)).\] 等变对偶保持所有等变稳定运算（包括限制、转移和 Weyl 群作用），且将 Borel 构造 \(X_{hG}\) 对偶化为真正等变同伦型。

\textbf{证明}：经典 Spanier-Whitehead 对偶使用嵌入 \(X \hookrightarrow S^n\) 并定义 \(D(X) = \Sigma^{-n} \operatorname{Th}(\nu)\)（Thom 谱）。等变版本将 \(X\) G-嵌入某表示球面 \(S^V\)（由 Mostow 嵌入定理保证），定义 \(D(X) = \Sigma^{-V} \operatorname{Th}_G(\nu)\)。对偶性 \([S^V \wedge X_+, E] \cong [S^V_+, D(X) \wedge E]\) 由等变 Thom 同构验证。函子性来自等变稳定同伦范畴中的映射锥构造。\(\blacksquare\)

\textbf{定理 10}（等变截影定理与不动点公式）：设 \(G\) 为有限群，\(X\) 为有限 G-CW 复形。对等变 K-理论，Lefschetz 不动点公式采取以下形式： \[\operatorname{Tr}(f \mid K_G^*(X)) = \sum_{(g)} \frac{1}{|C_G(g)|} \operatorname{Tr}(f^g \mid K^*(X^g)),\] 其中和取遍 \(G\) 的共轭类代表元，\(f^g\) 为 \(f\) 限制在 \(g\)-不动点集上的诱导映射。此公式是 Atiyah-Bott 不动点公式的等变推广，将等变上同调类通过不动点子簇的上同调类表示为零维积分的组合。

\textbf{证明}：由 Atiyah-Segal 完备化定理，\(K_G^*(X)\) 完备化后与 \(K(EG \times_G X)\) 同构。在 \(EG \times_G X\) 上经典 Atiyah-Bott 不动点公式给出 \(\operatorname{Tr} = \sum_{(g)} \frac{1}{|C_G(g)|} \operatorname{Tr}(f^g|_{X^g})\)。由于迹和完备化交换，该公式在完备化前也成立（两个方向的完备化映射在有限 G-CW 复形上均单射）。\(\blacksquare\)

\hrulefill

\hrulefill

\hrulefill

\section{第248章：Morse理论深化（Morse不等式证明）}\label{ux7b2c248ux7ae0morseux7406ux8bbaux6df1ux5316morseux4e0dux7b49ux5f0fux8bc1ux660e}

Morse 理论是微分拓扑和微分几何中最优美的理论之一，由 Marston Morse 在 1920-1930 年代建立。其核心思想是将光滑流形的拓扑结构与流形上 Morse 函数的临界点联系起来：流形的同调群由临界点的信息所控制。Morse 不等式给出了临界点个数的下界与同调群的 Betti 数之间的定量关系，其完整证明是微分拓扑与代数拓扑完美结合的典范。经典的 Morse 理论建立在有限维光滑流形上：设 \(M\) 是紧致无边光滑 \(n\) 维流形，\(f: M \to \mathbb{R}\) 是 Morse 函数（即所有临界点非退化------Hessian 满秩）。Morse 的核心引理表明：在临界点附近，\(f\) 可表示为标准二次型 \(f(x) = f(p) - x_1^2 - \cdots - x_\lambda^2 + x_{\lambda+1}^2 + \cdots + x_n^2\)，其中 \(\lambda\) 称为 Morse 指标（Hessian 负特征值的个数）。随着水平集 \(M^a = f^{-1}((-\infty, a])\) 穿过一个指标为 \(\lambda\) 的临界点，\(M^a\) 的同伦型发生变化------恰由粘合一个 \(\lambda\) 维胞腔（handle）得到。这一观察将 Morse 函数的信息转化为 CW 复形的胞腔结构，进而通过胞腔同调导出 Morse 不等式。

\subsection{189.x Morse引理与handle分解}\label{x-morseux5f15ux7406ux4e0ehandleux5206ux89e3}

\textbf{定义 189.1}（Morse函数与Morse指标）：光滑函数 \(f: M \to \mathbb{R}\) 的临界点 \(p \in M\) 满足 \(df_p = 0\)。在局部坐标 \((x_1, \ldots, x_n)\) 下，Hessian 矩阵 \(H_f(p) = (\partial^2 f / \partial x_i \partial x_j)(p)\) 是对称矩阵。\(p\) 称为\textbf{非退化的}（nondegenerate），若 \(H_f(p)\) 满秩。此时 \(H_f(p)\) 的负特征值的个数（重数计）称为 \(f\) 在 \(p\) 处的 \textbf{Morse 指标}，记作 \(\lambda(p)\)。若 \(f\) 的所有临界点均非退化，则 \(f\) 称为 \textbf{Morse 函数}。

\textbf{定理 189.1}（Morse引理）：设 \(p\) 是 Morse 函数 \(f\) 的非退化临界点，指标为 \(\lambda\)。则存在 \(p\) 附近的局部坐标 \((x_1, \ldots, x_n)\)（\(x_i(p) = 0\)）使得

\[f(x) = f(p) - x_1^2 - \cdots - x_\lambda^2 + x_{\lambda+1}^2 + \cdots + x_n^2\]

即 \(f\) 在其临界点附近与标准非退化二次型微分同胚等价。

\textbf{证明}：由 Taylor 展开 \(f(x) = f(0) + \sum_{i,j} h_{ij}(x) x_i x_j\)（\(h_{ij}\) 光滑，利用 \(f(0)=0\) 和 \(df_0=0\) 以及积分形式的余项）。设 \(Q(x) = \sum h_{ij}(0) x_i x_j = x^T H x\) 为 Hessian 二次型（\(H\) 对称非退化）。由 Sylvester 惯性定理，存在线性变换使 \(Q(x) = -y_1^2 - \cdots - y_\lambda^2 + y_{\lambda+1}^2 + \cdots + y_n^2\)。Morse 引理的精细部分：可将此线性变换提升为局部微分同胚（利用 Moser 路径方法或隐函数定理的归纳推广）。具体地，设 \(f_t(x) = (1-t) Q(x) + t f(x)\)（线性同伦）。在 \(t \in [0,1]\) 上构造依赖于 \(t\) 的向量场 \(X_t\) 使得 \(\mathcal{L}_{X_t} (f_t - Q) + (f - Q) = 0\)，然后积分 \(X_t\) 得所需微分同胚。非退化性保证 \(f_t\) 在 \(0\) 附近均为 Morse，保证了向量场的可解性。\(\blacksquare\)

\textbf{定理 189.2}（handle粘合定理 / 水平集变化）：设 \(f: M \to \mathbb{R}\) 是 Morse 函数，\(a < b\) 且 \(f^{-1}([a, b])\) 紧致，包含恰好一个临界点 \(p\)（指标 \(\lambda\)，临界值 \(c \in (a, b)\)）。则 \(M^b = f^{-1}((-\infty, b])\) 同伦等价于 \(M^a\) 粘合一个 \(\lambda\) 维胞腔（即 \(M^b \simeq M^a \cup_\varphi D^\lambda\)，其中 \(\varphi: S^{\lambda-1} \to M^a\) 为粘合映射）。

\textbf{证明}：选取 \(p\) 附近的 Morse 坐标，由 Morse 引理 \(f(x) = c - \|x_-\|^2 + \|x_+\|^2\)（\(x = (x_-, x_+)\)，\(x_- \in \mathbb{R}^\lambda\)，\(x_+ \in \mathbb{R}^{n-\lambda}\)）。在 \(p\) 的充分小邻域 \(U\) 中，水平集 \(M^{c-\varepsilon}\) 与 \(M^{c+\varepsilon}\) 的差异集中在 \(\|x_-\|^2 - \|x_+\|^2 = \text{const}\) 上。选择 \(\varepsilon > 0\) 充分小，构造显式的强形变收缩：在 \(U\) 内沿 \(f\) 的负梯度流向下推 \(M^{c+\varepsilon}\) 到 \(M^{c-\varepsilon}\) 并附加一个 \(\lambda\) 维胞腔（即 \(\|x_+\| = 0\) 且 \(f \leq c+\varepsilon\) 的区域------这是一个 \(\lambda\) 维圆盘 \(D^\lambda\)，其边界 \(S^{\lambda-1}\) 为 \(\|x_-\|^2 = \varepsilon\)，\(\|x_+\| = 0\) 落在 \(M^{c-\varepsilon}\) 中）。使用梯度流的流箱（flow box）技术严格化此构造：沿 \(-\nabla f\)（相对于某 Riemann 度量）的流将 \(M^{c+\varepsilon}\) 形变收缩到 \(M^{c-\varepsilon} \cup D^\lambda\)。\(\blacksquare\)

\subsection{189.y Morse不等式的精确陈述与证明}\label{y-morseux4e0dux7b49ux5f0fux7684ux7cbeux786eux9648ux8ff0ux4e0eux8bc1ux660e}

\textbf{定理 189.3}（弱Morse不等式 / Morse 1940）：设 \(f: M \to \mathbb{R}\) 是紧致无边流形 \(M\) 上的 Morse 函数。记 \(C_\lambda\) 为 \(f\) 的指标 \(\lambda\) 的临界点个数，\(b_\lambda = \operatorname{rank} H_\lambda(M; \mathbb{Z})\) 为第 \(\lambda\) 个 Betti 数。则 \[C_\lambda \geq b_\lambda \quad (\lambda = 0, 1, \ldots, n)\]

\textbf{证明}：由定理 189.2，按临界值递增顺序逐个添加胞腔，\(M\) 获得同伦等价于 CW 复形的结构，其 \(\lambda\) 维胞腔一一对应于指标 \(\lambda\) 的临界点。胞腔同调（定理 77.11）表明 \(H_\lambda(M)\) 同构于胞腔链复形的同调。胞腔链群 \(C_\lambda^{\text{cell}}\) 的秩恰为 \(C_\lambda\)（指标 \(\lambda\) 的临界点个数）。由代数同调的基本不等式，\(\operatorname{rank}(\ker d_\lambda) \leq \operatorname{rank}(C_\lambda^{\text{cell}}) = C_\lambda\)，且 \(b_\lambda = \operatorname{rank} H_\lambda \leq \operatorname{rank}(\ker d_\lambda)\)。故 \(b_\lambda \leq C_\lambda\)。\(\blacksquare\)

\textbf{定理 189.4}（强Morse不等式）：在弱 Morse 不等式的相同设定下，以下不等式组成立： \[C_\lambda - C_{\lambda-1} + C_{\lambda-2} - \cdots + (-1)^\lambda C_0 \geq b_\lambda - b_{\lambda-1} + b_{\lambda-2} - \cdots + (-1)^\lambda b_0\]

对所有 \(\lambda = 0, 1, \ldots, n\) 成立，且在 \(\lambda = n\) 时取等号：\(\sum_{\lambda=0}^n (-1)^\lambda C_\lambda = \sum_{\lambda=0}^n (-1)^\lambda b_\lambda = \chi(M)\)（Euler 示性数）。

\textbf{证明}：由定理 189.2 的胞腔分解，\(M\) 具有 CW 复形结构，其胞腔链复形为 \[\cdots \to C_\lambda^{\text{cell}} \xrightarrow{d_\lambda} C_{\lambda-1}^{\text{cell}} \to \cdots \to C_0^{\text{cell}} \to 0\] 其中 \(\operatorname{rank} C_\lambda^{\text{cell}} = C_\lambda\)，\(H_\lambda(M) \cong \ker d_\lambda / \operatorname{im} d_{\lambda+1}\)。记 \(Z_\lambda = \ker d_\lambda\)，\(B_\lambda = \operatorname{im} d_{\lambda+1}\)。由秩-零化度定理： \[\operatorname{rank} Z_\lambda = C_\lambda - \operatorname{rank} B_{\lambda-1}\] \[\operatorname{rank} B_\lambda = C_{\lambda+1} - \operatorname{rank} Z_{\lambda+1}\]

又 \(b_\lambda = \operatorname{rank} Z_\lambda - \operatorname{rank} B_\lambda\)。整理递推关系： \[\operatorname{rank} Z_\lambda = C_\lambda - C_{\lambda-1} + C_{\lambda-2} - \cdots \pm C_0 \mp \operatorname{rank} B_{-1}\] （\(B_{-1} = 0\)）。由于 \(b_\lambda = \operatorname{rank} Z_\lambda - \operatorname{rank} B_\lambda \leq \operatorname{rank} Z_\lambda\)（因 \(\operatorname{rank} B_\lambda \geq 0\)），代入 \(\operatorname{rank} Z_\lambda\) 的展开式即得强 Morse 不等式。当 \(\lambda = n\) 时，\(B_n = 0\)（无 \(n+1\) 维胞腔），\(b_n = \operatorname{rank} Z_n\)，且交错和相消给出等式 \(\sum (-1)^\lambda C_\lambda = \sum (-1)^\lambda b_\lambda = \chi(M)\)。\(\blacksquare\)

\textbf{推论 189.5}（Morse理论的Lacunary原理）：若 Morse 函数 \(f\) 满足 \(C_\lambda \cdot C_{\lambda+1} = 0\)（即无相邻指标的临界点），则所有边界同态 \(d_\lambda\) 均为零（因为不同维数的临界值可被分离）。此时 \(H_\lambda(M) \cong \mathbb{Z}^{C_\lambda}\)（无挠），且 \(C_\lambda = b_\lambda\)（弱不等式取等号）。这一原理在 Bott 周期性和等变 Morse 理论中有重要应用。

\textbf{证明}：若相邻指标的临界点不存在，可选取 Riemann 度量使不同指标的临界值严格分离（通过扰动 \(f\)）。则胞腔链复形的边界同态 \(d_\lambda: C_\lambda^{\text{cell}} \to C_{\lambda-1}^{\text{cell}}\) 必为零（因为 \(\lambda\) 维和 \(\lambda-1\) 维胞腔不对应任何梯度流线连接，而在 Morse-Smale 条件下 \(d_\lambda\) 的元素由连接临界点的梯度流线模 2 计数给出）。零边界同态意味着同调群直接等于链群（无商），故 \(H_\lambda(M) \cong \mathbb{Z}^{C_\lambda}\)。\(\blacksquare\)

\hrulefill

\hrulefill

\hrulefill

\section{第249章：配边理论与Pontryagin-Thom构造}\label{ux7b2c249ux7ae0ux914dux8fb9ux7406ux8bbaux4e0epontryagin-thomux6784ux9020}

配边理论（Cobordism Theory）是微分拓扑在 1950-1960 年代最辉煌的成就之一，由 Thom（1954）在其荣获 Fields 奖的工作中建立。两个闭 \(n\) 维光滑流形 \(M\) 和 \(N\) 称为配边的（cobordant），如果存在一个紧致 \(n+1\) 维光滑流形 \(W\) 使得 \(\partial W = M \sqcup \overline{N}\)（\(\overline{N}\) 表示 \(N\) 的定向反转）。配边关系是等价关系，配边等价类构成分次环 \(\Omega_* = \bigoplus_{n \geq 0} \Omega_n\)（在定向情形，\(\Omega_n^{\text{SO}}\) 为 \(n\) 维定向配边群）。Thom 的核心贡献是 Pontryagin-Thom 构造：将配边问题转化为稳定同伦论问题------计算 Thom 谱的同伦群。具体地，\(\Omega_n^{\text{SO}} \cong \pi_n^S(MSO) = \pi_{n+k}(MSO(k))\)（\(k \gg n\)），其中 \(MSO(k)\) 为万有定向 \(k\) 维向量丛的 Thom 空间。更一般地，对任意结构群 \(G\)（如 \(G = O, SO, U, SU, Sp\)），\(G\)-配边群同构于 \(MG\) 的稳定同伦群。Thom 的独创性在于：他证明了 Thom 空间 \(MO(k)\) 是 Eilenberg-MacLane 空间的 wedge（模 2），从而完全计算了非定向配边环 \(\mathfrak{N}_* = \Omega_*^O \cong \mathbb{Z}_2[x_2, x_4, x_5, x_6, x_8, \ldots]\)（\(x_i\) 为维数 \(i\) 的生成元，除 \(i \neq 2^j-1\) 外各维数一个）。

\subsection{190.x Thom空间与Pontryagin-Thom同构}\label{x-thomux7a7aux95f4ux4e0epontryagin-thomux540cux6784}

\textbf{定义 190.1}（Thom空间的Pontryagin-Thom构造）：设 \(\xi: E \to B\) 是秩 \(k\) 的向量丛（赋予 Riemann 度量）。\(\xi\) 的 \textbf{Thom 空间} \(T(\xi)\) 定义为圆盘丛 \(D(\xi)\) 模去球面丛 \(S(\xi)\) 所得的商空间：\(T(\xi) = D(\xi) / S(\xi)\)。等价地，\(T(\xi)\) 是 \(\xi\) 的单点紧化在每根纤维上的 Thom 化。若 \(\xi\) 为万有丛 \(\gamma^k \to BO(k)\)（或 \(BSO(k)\)），则 \(T(\gamma^k) = MO(k)\)（或 \(MSO(k)\)）。

\textbf{定理 190.1}（Pontryagin-Thom构造 / 基本同构，Thom 1954）：设 \(M^n\) 是闭光滑 \(n\) 维流形。将 \(M^n\) 嵌入 \(\mathbb{R}^{n+k}\)（\(k \gg n\)，总是存在，取 \(k \geq n+1\) 绰绰有余）。\(M\) 在 \(\mathbb{R}^{n+k}\) 中的管状邻域 \(\nu\) 同构于法丛 \(N(M \subset \mathbb{R}^{n+k})\)（秩 \(k\)）。定义\textbf{Pontryagin-Thom 映射} 为塌缩映射

\[\Sigma^{n+k} \cong S^{n+k} \xrightarrow{\text{塌缩 } S^{n+k} \setminus \nu} T(\nu(M)) \xrightarrow{\text{分类映射}} MSO(k)\]

将补集 \(\mathbb{R}^{n+k} \setminus \nu\) 塌缩为一点，得到 \(\nu\) 的 Thom 空间 \(T(\nu) \cong \Sigma^n M_+\)（法丛的 Thom 空间等同于 \(M\) 带不交基点的 \(n\) 次悬挂），再通过法丛的分类映射（Gauss 映射 \(\nu \to \gamma^k\)）映至 \(MSO(k)\)。此构造给出同态

\[\Phi: \Omega_n^{\text{SO}} \longrightarrow \pi_{n+k}(MSO(k))\]

当 \(k\) 充分大时，\(\Phi\) 是同构。此即\textbf{Pontryagin-Thom同构}：定向配边群同构于 Thom 谱的稳定同伦群。

\textbf{证明}（\(\Phi\) 的逆构造）：给定稳定同伦类 \(f: S^{n+k} \to MSO(k)\)。由 Thom 横截定理（Thom transversality），可同伦地调整 \(f\) 使其横截于 \(MSO(k)\) 中零截面 \(BSO(k) \subset MSO(k)\)（零截面为 Thom 空间的基点补）。\(f^{-1}(BSO(k)) \subset S^{n+k}\) 是一个 \(n\) 维光滑子流形 \(M^n\)（由于 \(BSO(k)\) 的余维数为 \(k\)）。\(M\) 的法丛为 \(\nu = f^* \gamma^k\)，而 \(\gamma^k\) 是可定向的（因考虑 \(BSO(k)\) 而非 \(BO(k)\)），故 \(M\) 可定向且具有诱导的定向。此配边类 \([M] \in \Omega_n^{\text{SO}}\) 即为 \(\Phi\) 的逆像。两个映射互逆的验证使用横截同伦和管状邻域定理。\(\blacksquare\)

\subsection{190.y Thom谱的稳定同伦群计算}\label{y-thomux8c31ux7684ux7a33ux5b9aux540cux4f26ux7fa4ux8ba1ux7b97}

\textbf{定理 190.2}（Thom同构定理，1954）：设 \(\xi\) 是秩 \(k\) 的实向量丛（基空间为 \(B\)）。则存在 Thom 同构 \[\Phi: H^{i}(B; R) \xrightarrow{\cong} \tilde{H}^{i+k}(T(\xi); R)\] （\(R = \mathbb{Z}\) 对可定向丛，\(R = \mathbb{Z}_2\) 对任意丛），由杯积 Thom 类 \(u \in H^k(T(\xi); R)\) 给出。下同调版本：\(\Phi_*: \tilde{H}_{i+k}(T(\xi); R) \cong H_i(B; R)\)。

\textbf{证明}：Thom 类 \(u\) 是 \(H^k(D(\xi), S(\xi); R)\)（法丛对 \((D(\xi), S(\xi))\) 的上同调类），满足限制到每根纤维 \((D^k, S^{k-1})\) 上为生成元。Thom 同构 \(\Phi(\alpha) = \pi^* \alpha \smile u\)（\(\pi: D(\xi) \to B\) 为投影）是上同调环的同构，由局部平凡化结合 Mayer-Vietoris 和 Leray-Hirsch 定理（或等价的纤维上同调谱序列的退化性）证明。下同调版本由 cap 积（或 Poincare 对偶）导出。\(\blacksquare\)

\textbf{定理 190.3}（Thom的配边环计算，模 2）：非定向配边环 \(\mathfrak{N}_* = \Omega_*^O\) 同构于多项式环 \[\mathfrak{N}_* \cong \mathbb{Z}_2[x_2, x_4, x_5, x_6, x_8, x_9, x_{10}, x_{11}, \ldots]\] 其中生成元 \(x_i\) 为维数 \(i\) 的不可分解元素（\(i\) 非 \(2^j - 1\) 形式，当 \(i = 2^j - 1\) 时 \(x_i = 0\)），对应于实射影空间 \(\mathbb{RP}^i\)（\(i\) 偶）和 Dold 流形。

\textbf{证明}（框架）：Thom 谱的模 2 上同调为 \(H^*(MO; \mathbb{Z}_2) \cong \mathbb{Z}_2[\cdots]\)（自由 \(\mathbb{Z}_2\)-模，由 Thom 同构和 \(H^*(BO; \mathbb{Z}_2) \cong \mathbb{Z}_2[w_1, w_2, \ldots]\)（Stiefel-Whitney 类）决定）。利用 Adams 谱序列（以 \(H^*(-; \mathbb{Z}_2)\) 为系数），\(MO\) 的模 2 上同调作为 Steenrod 代数 \(\mathcal{A}_2\) 上的模是完全自由的（可分解为 \(\mathcal{A}_2\) 的直和），因此 Adams 谱序列在 \(E_2\) 页退化，给出 \(\pi_*(MO) \otimes \mathbb{Z}_2 \cong \mathbb{Z}_2[\cdots]\)。进一步分析生成元的维数（对应 \(\mathcal{A}_2\)-模的不可约基元）得到多项式环结构。\(\blacksquare\)

\subsection{190.z 有理配边与Hirzebruch符号差定理}\label{z-ux6709ux7406ux914dux8fb9ux4e0ehirzebruchux7b26ux53f7ux5deeux5b9aux7406}

\textbf{定理 190.4}（有理配边环的Milnor-Novikov计算）：定向配边环张量 \(\mathbb{Q}\) 同构于多项式环 \[\Omega_*^{\text{SO}} \otimes \mathbb{Q} \cong \mathbb{Q}[\mathbb{CP}^2, \mathbb{CP}^4, \mathbb{CP}^6, \ldots]\] 其中生成元为复射影空间 \(\mathbb{CP}^{2m}\)（维数 \(4m\)，带定向）。\(\Omega_n^{\text{SO}} \otimes \mathbb{Q}\) 仅在 \(n \equiv 0 \pmod{4}\) 时非零。有理 Pontryagin 数构成每维独立的配边不变量，由 Hirzebruch 符号差定理给出。

\textbf{证明}（框架）：Thom 谱 \(MSO\) 的有理上同调为 \(H^*(MSO; \mathbb{Q}) \cong H^*(BSO; \mathbb{Q}) \otimes \tilde{H}^*(MSO; \mathbb{Q})\)（Thom 同构）。\(H^*(BSO; \mathbb{Q}) \cong \mathbb{Q}[p_1, p_2, \ldots]\)（Pontryagin 类）。有理系数下，\(MSO\) 的稳定同伦群同构于其同调群（因有理 Hurewicz 同态 \(\pi_* \otimes \mathbb{Q} \cong H_* \otimes \mathbb{Q}\) 对 Thom 谱成立），后者由 Thom 同构给出 \(H_{n+k}(MSO(k); \mathbb{Q}) \cong H_n(BSO(k); \mathbb{Q})\)（\(k \gg n\)）。稳定化后，\(\pi_n^S(MSO) \otimes \mathbb{Q} \cong H_n(BSO; \mathbb{Q})\) 在 \(n \equiv 0 \pmod{4}\) 时等于 \(\mathbb{Q}[p_1, p_2, \ldots]\) 的 \(n\) 次齐次部分。生成元 \(\mathbb{CP}^{2m}\) 对应于 \(s_{2m} = [\mathbb{CP}^{2m}]\)，其 Pontryagin 数为 \(\langle p_m(\mathbb{CP}^{2m}), [\mathbb{CP}^{2m}] \rangle = 2m+1\) 非零。\(\blacksquare\)

\newpage

\chapter{07-拓扑与几何/03-代数拓扑/05-低维拓扑.md}\label{ux62d3ux6251ux4e0eux51e0ux4f5503-ux4ee3ux6570ux62d3ux625105-ux4f4eux7ef4ux62d3ux6251.md}

\section{第250章：三维流形基础}\label{ux7b2c250ux7ae0ux4e09ux7ef4ux6d41ux5f62ux57faux7840}

三维流形是低维拓扑学中最丰富也最深刻的研究对象。与维数大于 3 的流形不同（高维的 s-cobordism 定理使得同伦等价蕴含同胚），三维流形拥有精巧的组合结构------Heegaard 分解将任意闭三维流形表示为两个柄体的并，Dehn 手术则通过切除一个实心环管再以不同方式粘回，系统地生成所有紧致定向三维流形（Lickorish-Wallace 定理）。William Thurston 在 1970 年代的几何化猜想（2003 年由 Perelman 证明）揭示了三维流形可以被唯一地切割为若干片段，每片具有 8 种 Thurston 模型几何之一------这一定理不仅是 20 世纪拓扑学的高峰，更直接蕴含了 Poincare 猜想的证明。在代数拓扑的语境中，三维流形紧致而无边界时，基本群几乎确定了流形本身（同伦等价），但对 \(S^1 \times S^2\) 等例外需谨慎处理。本章为后续各章（纽结理论第 190 章、量子不变量第 191 章、Floer 同调第 192 章、四维流形第 193 章）提供了基础语言和构造技术。

\subsection{217.1 三维流形的基本构造}\label{ux4e09ux7ef4ux6d41ux5f62ux7684ux57faux672cux6784ux9020}

\textbf{定义 217.1}（Heegaard 分解）：亏格 \(g\) 的 \textbf{Heegaard 分解}是一个闭 3 维流形 \(M\) 的表示为 \(M = H_g \cup_\varphi H_g\)，其中 \(H_g\) 是亏格 \(g\) 的柄体（handlebody），\(\varphi: \partial H_g \to \partial H_g\) 是边界曲面 \(\Sigma_g\) 上的同胚。

\textbf{定理 217.1}（Heegaard 定理）：任何闭可定向 3 维流形都存在 Heegaard 分解。Heegaard 分解的亏格是 3 维流形的基本量。

\emph{证明}：取 \(M\) 的光滑三角剖分，令 \(H\) 为其 1-骨架的正则邻域（亏格为 \(g\) 的柄体），则 \(M \setminus \operatorname{int}(H)\) 的闭包也是亏格为 \(g\) 的柄体，粘合同胚 \(\varphi\) 由 \(M\) 的 2-胞腔的粘合映射决定。此构造给出了 \(M\) 的 Heegaard 分解。\(\blacksquare\)

\textbf{定义 217.2}（Dehn 手术）：设 \(K \subset S^3\) 是纽结，\(N(K)\) 是 \(K\) 的管状邻域。沿 \(K\) 的 \textbf{Dehn 手术}是将 \(N(K)\) 挖去，再将实心环面 \(S^1 \times D^2\) 沿 \(\partial N(K)\) 的某条简单闭曲线粘回。所有闭可定向 3 维流形都可通过 \(S^3\) 中某个链环的 Dehn 手术得到（Lickorish-Wallace 定理）。

\textbf{定理 217.2}（Lickorish-Wallace 定理，1960-1962）：任何闭可定向 3 维流形可通过 \(S^3\) 中链环的 \(\pm 1\) Dehn 手术得到。

\textbf{证明}：设 \(M\) 的 Heegaard 分解为 \(M = H_g \cup_\varphi H_g\)。将 \(H_g\) 嵌入 \(S^3\) 作为柄体的标准嵌入，则 \(M\) 可由 \(S^3\) 通过对 \(\varphi\) 的拉回链环作 Dehn 手术得到。Lickorish 证明任意 \(\varphi\) 可分解为沿某些闭曲线的 Dehn 扭转的复合，每条闭曲线对应 \(S^3\) 中链环的一个分量，且手术系数可规范化为 \(\pm 1\)。Wallace 随后给出独立证明。\(\blacksquare\)

\subsection{217.2 三维流形的分解}\label{ux4e09ux7ef4ux6d41ux5f62ux7684ux5206ux89e3}

\textbf{定理 217.3}（Kneser-Milnor 素分解定理）：每个紧致可定向 3 维流形可唯一分解为素流形的连通和（\(M = M_1 \# M_2 \# \cdots \# M_n\)），其中 \(M_i\) 是闭单连通 3 维流形（同胚于 \(S^3\)）或不可约流形。

\textbf{证明}：存在性由归纳构造：若 \(M\) 包含一个本质分离球面 \(S^2\)，则 \(M\) 沿此球面分解为连通和。由紧致性，此过程在有限步终止。唯一性由 Milnor 证明：若 \(M \cong M_1 \# \cdots \# M_n \cong N_1 \# \cdots \# N_m\)，则根据 Kneser 猜想（后由 Milnor 证明），\(n=m\) 且 \(\{M_i\}\) 与 \(\{N_j\}\) 在同胚类下的重数相同。\(\blacksquare\)

\textbf{定理 217.4}（JSJ 分解定理 / Jaco-Shalen-Johannson 1979）：每个不可约 3 维流形可沿本质环面（incompressible tori）唯一地分解为 Seifert 纤维流形和双曲流形。

\textbf{证明}：构造 \(M\) 中所有本质环面的极大族（互不相交且两两不平行），然后沿该族切开。每个得到的分片要么是 Seifert 纤维的（含本质环面构造），要么是双曲的（无本质环面）。唯一性由本质环面族在合痕下的唯一性保证------本质环面族是 \(M\) 的特征子流形。\(\blacksquare\)

\textbf{定义 217.3}（Seifert 纤维流形）：\textbf{Seifert 纤维流形}是局部平凡 \(S^1\) 纤维的 3 维流形，允许有限条奇异纤维。其不变量是 \((g; b; (\alpha_1, \beta_1), \ldots, (\alpha_n, \beta_n))\)（基曲面亏格、Euler 数和奇异纤维类型）。

\textbf{定义 217.4}（双曲 3 维流形）：\textbf{双曲 3 维流形}是容许完备双曲度量（截面曲率恒为 \(-1\)）的 3 维流形。由 Mostow 刚性定理，双曲 3 维流形上的双曲结构唯一（如果存在）。

\textbf{定理 217.5}（Mostow 刚性定理，1968）：若 \(M\) 和 \(N\) 是有限体积的完备双曲 3 维流形，且 \(\pi_1(M) \cong \pi_1(N)\)，则 \(M\) 和 \(N\) 等距。即双曲 3 维流形的拓扑完全决定了其双曲度量。

\textbf{证明}：同构 \(\pi_1(M) \cong \pi_1(N)\) 诱导拟等距 \(f: M \to N\)。双曲空间 \(\mathbb{H}^3\) 的边界是 \(S^2\)，其拟共形映射的刚性（由 Mostow 证明）确保 \(f\) 在无穷远边界上诱导共形映射，从而延拓为 \(\mathbb{H}^3\) 上的等距。故 \(M \cong \mathbb{H}^3/\pi_1(M) \cong \mathbb{H}^3/\pi_1(N) \cong N\) 为等距。Prasad 后将证明推广到有限体积情形。\(\blacksquare\)

\subsection{217.3 几何化纲领}\label{ux51e0ux4f55ux5316ux7eb2ux9886}

\textbf{定理 217.6}（Thurston 几何化猜想 / Perelman 证明，1982/2003）：任何闭可定向 3 维流形可沿球面和环面唯一地分解为素流形，其中每个素流形恰好容许 8 种 Thurston 几何之一。

\textbf{证明}：由 Kneser-Milnor 和 JSJ 分解将 \(M\) 化为素和 atoroidal 分片。对每片运行 Ricci 流（带手术）。Perelman 证明 Ricci 流在有限时间后将各片演化为具有局部齐性度量的流形------每个片段由 Thurston 的 8 种几何之一模型化。特别是双曲分片在 Ricci 流下收敛到常负曲率度量，从而确认为双曲几何。\(\blacksquare\)

\textbf{定义 217.5}（Thurston 的 8 种几何）：\(S^3\)（球面几何），\(\mathbb{E}^3\)（欧氏几何），\(\mathbb{H}^3\)（双曲几何），\(S^2 \times \mathbb{R}\)，\(\mathbb{H}^2 \times \mathbb{R}\)，\(\widetilde{\operatorname{SL}(2, \mathbb{R})}\)（\(\operatorname{SL}(2, \mathbb{R})\) 的泛覆盖），\(\operatorname{Nil}\)（幂零几何），\(\operatorname{Sol}\)（可解几何）。

\textbf{定理 217.7}（Poincaré 猜想 / Perelman 2003）：任何单连通闭 3 维流形同胚于 \(S^3\)。这是几何化猜想的直接推论。

\textbf{证明}：取 \(M\) 上任意 Riemann 度量，运行 Ricci 流（带手术）。由于 \(\pi_1(M)=1\)，\(M\) 的素分解中不含非平凡因子。Perelman 证明 Ricci 流在有限时间内使流形消失（extinction），由此可推出 \(M\) 同胚于连通和 \(\# S^3/\Gamma_i\)。结合 \(\pi_1(M)=1\) 推得所有 \(\Gamma_i = 1\)，故 \(M \cong S^3\)。\(\blacksquare\)

\textbf{定理 217.8}（球面空间形式猜想 / 椭圆化猜想）：任何有限基本群的闭 3 维流形是球面空间形式（即 \(S^3/\Gamma\)，\(\Gamma\) 是有限群自由作用在 \(S^3\) 上）。由 Perelman 证明。

\textbf{证明}：由几何化推知 \(M\) 容许 Thurston 几何之一。\(|\pi_1(M)|<\infty\) 排除 \(\mathbb{E}^3\), \(\mathbb{H}^3\), \(\mathbb{H}^2\times\mathbb{R}\), \(\widetilde{\operatorname{SL}(2,\mathbb{R})}\), \(\operatorname{Nil}\), \(\operatorname{Sol}\) 等几何（它们的基本群包含无限阶元），而 \(S^2\times\mathbb{R}\) 情形中 \(M\) 必含本质球面。故 \(M\) 只能是 \(S^3\)-几何，即 \(M\cong S^3/\Gamma\)。\(\blacksquare\)

\hrulefill

\hrulefill

\hrulefill

\hrulefill

\hrulefill

\hrulefill

\section{第251章：纽结理论基础}\label{ux7b2c251ux7ae0ux7ebdux7ed3ux7406ux8bbaux57faux7840}

纽结理论是低维拓扑学中最古老、最具直观吸引力的分支。其根本问题是：给定空间中的两个纽结（\(S^1\) 在 \(S^3\) 中的嵌入），如何判断它们是否等价（即存在保持定向的环境同痕将一个纽结变为另一个）？这一看似简单的问题从 Gauss 在 1833 年引入连接数（linking number）开始，经历了 Tait 和 Little 在 19 世纪末的纽结列表工作、Alexander 在 1928 年发现 Alexander 多项式、直到 1984 年 Vaughan Jones 发现 Jones 多项式------后者统一了统计力学（Yang-Baxter 方程）、算子代数（von Neumann 代数）和低维拓扑，使 Jones 获得 1990 年菲尔兹奖。纽结理论在三维流形拓扑中扮演着枢纽角色：每个闭的三维流形可通过 Dehn 手术从一个链环构造（Lickorish-Wallace 定理），从而将三维流形分类问题部分归约为纽结不变量问题。本章从纽结和链环的基本定义入手，建立 Reidemeister 移动和基本群表示，进而讨论 Alexander 不变量、Jones 多项式和 HOMFLY-PT 多项式，为第 191 章量子拓扑不变量和第 192 章 Floer 同调中纽结的应用奠定基础。

\subsection{218.1 纽结与链环}\label{ux7ebdux7ed3ux4e0eux94feux73af}

\textbf{定义 218.1}（纽结与链环）：\textbf{纽结}是 \(S^1\) 到 \(S^3\)（或 \(\mathbb{R}^3\)）的光滑嵌入（或局部平坦嵌入）。\textbf{链环}是多个不相交的 \(S^1\) 的嵌入。纽结 \(K\) 的\textbf{补空间}是 \(S^3 \setminus \nu(K)\)（\(\nu(K)\) 是 \(K\) 的管状邻域）。

\textbf{定义 218.2}（纽结的等价 / 合痕）：两个纽结 \(K_0\) 和 \(K_1\) 是\textbf{合痕的}（或等价的），如果存在连续单参数族 \(\{h_t\}_{t \in [0,1]}\) 的 \(\mathbb{R}^3\) 上的同胚，使得 \(h_0 = \operatorname{id}\) 且 \(h_1(K_0) = K_1\)。

\textbf{定理 218.1}（Reidemeister 定理，1927）：两个纽结图表示同一个纽结（模合痕）当且仅当它们可以通过有限次 \textbf{Reidemeister 运动}（I 型、II 型、III 型）相互转化。

\textbf{证明}：将纽结合痕分解为平面上的投影变换，在每个投影的非横截点处分析局部变化，表明任意合痕可分解为有限个局部操作，每个操作恰好对应一种 Reidemeister 运动。严格地，利用多边逼近将 3 维合痕变形为组合方式，在二维投影上表现为 R1-R3 的组合。\(\blacksquare\)

\textbf{定义 218.3}（Reidemeister 运动）： - \textbf{R1}：添加或消除一个扭结（twist） - \textbf{R2}：对两条交叉的弧线，可滑动一根穿过另一根 - \textbf{R3}：三条弧线相交时，中间弧线可滑过交叉点

\subsection{218.2 经典纽结不变量}\label{ux7ecfux5178ux7ebdux7ed3ux4e0dux53d8ux91cf}

\textbf{定义 218.4}（交叉数）：纽结 \(K\) 的\textbf{交叉数} \(c(K)\) 是 \(K\) 的所有纽结图中交叉点数的最小值。交叉数是纽结不变量。

\textbf{定义 218.5}（三色性 / Fox 染色）：纽结图是\textbf{三色可染的}，如果每条弧可染以三种颜色之一，使得每个交叉点处三条弧要么同色要么全不同色，且至少使用两种颜色。三色可染性是纽结不变量。

\textbf{定义 218.6}（Alexander 多项式，1928）：纽结 \(K\) 的 \textbf{Alexander 多项式} \(\Delta_K(t) \in \mathbb{Z}[t, t^{-1}]\) 是 \(K\) 的补空间无穷循环覆盖的一维同调的不变量。\(\Delta_K(t)\) 是 \(t \leftrightarrow t^{-1}\) 对称的，且 \(\Delta_K(1) = \pm 1\)。

\textbf{定理 218.2}（Alexander 多项式的构造）：从纽结的 Seifert 曲面 \(S\) 的 Seifert 矩阵 \(V\) 出发，\(\Delta_K(t) = \det(t^{1/2} V - t^{-1/2} V^T)\)（模去 \(t^{\pm 1/2}\) 的乘法）。

\textbf{证明}：设 \(K\) 为定向纽结，\(S\) 为其 Seifert 曲面。取 \(S\) 的一组 \(1\)-同调基 \(\{a_1, \ldots, a_{2g}\}\)（\(g\) 为 \(S\) 的亏格），定义 Seifert 矩阵 \(V_{ij} = \operatorname{lk}(a_i, a_j^+)\)，其中 \(a_j^+\) 是 \(a_j\) 沿正法向推开得到的曲线。\(S \times [-1,1]\) 是 \(S^3 \setminus K\) 的开邻域，其同调由 \(a_i\) 和 \(a_i^+\) 生成。Alexander 多项式由 \(H_1(\widetilde{S^3 \setminus K})\) 作为 \(\mathbb{Z}[t, t^{-1}]\)-模的表示矩阵给出。具体地，无穷循环覆叠空间的同调由 \(V\) 和 \(V^T\) 决定，表示矩阵为 \(tV - V^T\)（或等价地 \(V - tV^T\)）。因此 \(\Delta_K(t) = \det(tV - V^T)\)。该多项式在变换 \(t \mapsto t^{-1}\) 下对称（模去乘法因子），且为纽结不变量（不依赖于 Seifert 曲面的选取）。\(\blacksquare\)

\emph{不变性证明概要}：Alexander 多项式 \(\Delta_K(t)\) 是 \(K\) 的补空间 \(S^3 \setminus K\) 的无穷循环覆盖的 1 维同调的扭化 Alexander 模的生成多项式。证明其不变性分两步：（1）对任何给定的纽结图构造 Wirtinger 表示并计算 Fox 微积分，得到 Alexander 矩阵，其初等因子给出 \(\Delta_K(t)\)。Reidemeister 运动（R1-R3）下 Alexander 矩阵的初等因子不变，保证 \(\Delta_K(t)\) 是合痕不变量。（2）等价地，通过 Seifert 曲面方法：不同 Seifert 曲面通过 S-等价性关联，Seifert 矩阵在 S-等价下作初等变换（子矩阵扩展或削减），而 \(\det(t^{1/2}V - t^{-1/2}V^T)\) 在这些变换下保持不变（模去 \(t^{\pm 1/2}\)）。∎

\textbf{定义 218.7}（纽结群）：纽结 \(K\) 的\textbf{纽结群} \(\pi_K = \pi_1(S^3 \setminus K)\) 是 \(K\) 的补空间的基本群。纽结群是强有力的纽结不变量。Wirtinger 表示：从纽结分叉点出发，每个弧段对应一个生成元，每个交叉点对应一个关系。

\textbf{定理 218.3}（Gordon-Luecke 定理，1989）：纽结由它的补空间决定（不依赖于嵌入方式）。即若 \(S^3 \setminus K \cong S^3 \setminus K'\)，则 \(K \cong K'\)。

\textbf{证明}：设 \(h: S^3 \setminus K \to S^3 \setminus K'\) 是同胚。考虑 \(K\) 的管状邻域边界 \(\partial N(K) \cong T^2\)，其在 \(S^3 \setminus K\) 中的位置由纽结的经线 \(m\) 和经度 \(l\) 的 framing 决定。Gordon 和 Luecke 证明 \(h\) 必然将 \(K\) 的经线映至 \(K'\) 的经线（利用循环覆盖和纽结群的性质），从而 \(h\) 可延拓为 \(S^3\) 到自身的同胚将 \(K\) 映至 \(K'\)。\(\blacksquare\)

\subsection{218.3 Jones 多项式与量子不变量}\label{jones-ux591aux9879ux5f0fux4e0eux91cfux5b50ux4e0dux53d8ux91cf}

\textbf{定义 218.8}（Jones 多项式，1984）：Vaughan Jones 定义的 \textbf{Jones 多项式} \(V_K(t) \in \mathbb{Z}[t^{1/2}, t^{-1/2}]\) 由以下缆式（skein）关系唯一确定：

\[t^{-1} V_{L_+}(t) - t V_{L_-}(t) = (t^{1/2} - t^{-1/2}) V_{L_0}(t)\]

其中 \(L_+\)、\(L_-\)、\(L_0\) 是仅在交叉点处不同的链环图，且 \(V_{\text{平凡纽结}}(t) = 1\)。

\textbf{定理 218.4}（Jones 多项式的性质）： - \(V_K(1) = (-2)^{\#K - 1}\)（\(\#K\) 是链环的分量数） - \(V_K(t)\) 是链环不变量（比 Alexander 多项式更强） - Jones 多项式区分了手性（chirality）：存在 \(V_K(t) \neq V_K(t^{-1})\) 的纽结 - 通过 von Neumann 代数（特别是 \(\operatorname{II}_1\) 因子）和辫群表示得到

\textbf{证明}：（1）由 skein 关系，令 \(t=1\) 得 \(V_{L_+}(1) = V_{L_-}(1)\)，故 \(V_L(1)\) 仅依赖于 \(L\) 的分量数，归纳可证 \(V_L(1)=(-2)^{c-1}\)。（2）不变性由 skein 关系归纳于交叉数，验证 Reidemeister 运动下不变。（3）三叶结 \(3_1\) 满足 \(V_{3_1}(t) \neq V_{3_1}(t^{-1})\)，即其与镜像不等价。（4）Jones 通过 \(\operatorname{II}_1\) 因子的迹构造辫群表示 \(B_n \to \operatorname{GL}(V^{\otimes n})\) 并取 Markov 迹，得到链环不变量。\(\blacksquare\)

\textbf{定义 218.9}（HOMFLY-PT 多项式，1985）：Freyd-Yetter、Hoste-Lickorish-Millett 和 Przytycki-Traczyk 独立发现的二变量纽结多项式 \(P_K(a, z)\)，满足：

\[a P_{L_+}(a, z) - a^{-1} P_{L_-}(a, z) = z P_{L_0}(a, z)\]

且 \(P_{\text{平凡纽结}}(a, z) = 1\)。Jones 多项式是 \(a = t^{-1}\)、\(z = t^{1/2} - t^{-1/2}\) 的特例，Alexander 多项式是 \(a = 1\) 的特例。

\textbf{定义 218.10}（Kauffman 括号多项式）：Kauffman 括号 \(\langle D \rangle\) 是未定向链环图 \(D\) 的多项式不变量，满足： - \(\langle \bigcirc \rangle = 1\)（平凡纽结） - \(\langle D \sqcup \bigcirc \rangle = (-A^2 - A^{-2}) \langle D \rangle\) - \(\langle \text{正交叉} \rangle = A \langle \text{平行} \rangle + A^{-1} \langle \text{反平行} \rangle\)

Jones 多项式通过 \(V_K(t) = (-A^3)^{-w(D)} \langle D \rangle |_{A = t^{-1/4}}\) 得到（\(w(D)\) 是拧数）。

\subsection{218.4 Khovanov 同调}\label{khovanov-ux540cux8c03}

\textbf{定义 218.11}（Khovanov 同调，2000）：Mikhail Khovanov 构造了 \textbf{Khovanov 同调} \(\operatorname{Kh}^{i,j}(K)\)，它是 Jones 多项式的「范畴化」（categorification）。其分次 Euler 特征给出 Jones 多项式：

\[\sum_{i,j} (-1)^i q^j \dim \operatorname{Kh}^{i,j}(K) = (q + q^{-1}) V_K(q^2)\]

\textbf{定理 218.5}（Khovanov 同调的性质）：Khovanov 同调是比 Jones 多项式更强的纽结不变量。它区分了具有相同 Jones 多项式但不同 Khovanov 同调的纽结，且具有函子性质（跨越同痕诱导同构）。

\textbf{证明}：给定纽结图 \(D\)，对每个交叉点分配一个分次向量空间（\(V = \operatorname{Span}\{v_+, v_-\}\)），构造立方复形 \(C^{i,j}(D)\)，其微分由交叉点处的 cobordism 诱导。Khovanov 证明该复形的同伦型在 Reidemeister 运动下不变（需显式构造链同伦等价），从而 \(\operatorname{Kh}^{i,j}(K)\) 是纽结不变量。函子性来自合痕所诱导的链映射在复形层面可提升为链同伦。\(\blacksquare\)

\hrulefill

\hrulefill

\hrulefill

\hrulefill

\hrulefill

\hrulefill

\section{第252章：量子拓扑不变量}\label{ux7b2c252ux7ae0ux91cfux5b50ux62d3ux6251ux4e0dux53d8ux91cf}

量子拓扑不变量代表了 20 世纪末低维拓扑学最富戏剧性的发展------将量子场论（特别是 Witten 的 Chern-Simons 理论）的物理直觉转化为严格的数学不变量。1988 年，Edward Witten 在一篇影响深远的论文中提出：三维流形 \(M\) 上的 Chern-Simons 路径积分 \(Z_k(M) = \int \exp(ik \cdot CS(A)) \mathcal{D}A\) 应给出 \(M\) 的拓扑不变量，且对于带有纽结 \(K\) 的空间，关联函数 \(\langle W_K \rangle\) 正是 Jones 多项式的推广（colored Jones 多项式）。这一物理推导虽然在数学上非严格，却催生了 Reshetikhin-Turaev（1991）基于量子群 \(U_q(\mathfrak{sl}_2)\) 表示的严格数学构造------RT 不变量------以及 Turaev-Viro 基于球面范畴的状态和不变量。与经典拓扑不变量（基本群、同调群）不同，量子不变量的优势在于能区分某些具有相同基本群和同调群的流形。本章与第 190 章纽结理论紧密呼应：量子群和辫群表示（通过 Markov 迹）提供了从纽结多项式到三维流形不变量的标准途径，而 Yang-Baxter 方程则确保了不变量在 Reidemeister 移动下的不变性。

\subsection{219.1 辫群与 Markov 定理}\label{ux8fabux7fa4ux4e0e-markov-ux5b9aux7406}

\textbf{定义 219.1}（辫群 / Artin 辫群）：\(n\) 股辫的 \textbf{Artin 辫群} \(B_n\) 由生成元 \(\sigma_1, \ldots, \sigma_{n-1}\) 和关系定义： - \(\sigma_i \sigma_j = \sigma_j \sigma_i\)（\(|i - j| \geq 2\)） - \(\sigma_i \sigma_{i+1} \sigma_i = \sigma_{i+1} \sigma_i \sigma_{i+1}\)（辫关系）

\textbf{定理 219.1}（Alexander 定理，1923）：任何纽结/链环可表示为某个辫的闭包（braid closure）。

\emph{证明}：将纽结的投影图旋转使所有交叉点带向下倾斜。取一条竖直线绕纽结旋转，标记与纽结的每次横截相交点。顺时针移动竖直线，当它扫过交叉点时，对应辫生成元 \(\sigma_i\) 或 \(\sigma_i^{-1}\)。最终竖直线回到原位时，纽结即为这些辫生成元之积的闭包。\(\blacksquare\)

\textbf{定理 219.2}（Markov 定理，1936）：两个辫表示同一个链环当且仅当它们可通过 \textbf{Markov 运动}（共轭 \(\sigma_i b \sigma_i^{-1}\) 和稳定化 \(b \leftrightarrow b \sigma_n^{\pm 1} \in B_{n+1}\)）相互转化。

\textbf{证明}：（必要性）两种 Markov 运动不改变辫闭包的合痕类：共轭对应于投影图中缠绕的重新定位，稳定化对应于添加一个无影响的扭结。（充分性）任何将闭包相连的合痕可分解为有限个局部变换，这些变换对应于 Markov 运动的组合。完整证明需分析合痕在 \(S^3\) 中的每个投影事件。\(\blacksquare\)

\textbf{定义 219.2}（辫群表示与纽结不变量）：寻找 \(B_n\) 的一族表示 \(\{\rho_n: B_n \to \operatorname{GL}(V_n)\}\) 满足 Markov 迹条件，从而构造链环不变量。Jones 多项式由此构造。

\subsection{219.2 量子群与 R-矩阵}\label{ux91cfux5b50ux7fa4ux4e0e-r-ux77e9ux9635}

\textbf{定义 219.3}（量子群 / 量子包络代数）：\textbf{量子群} \(U_q(\mathfrak{g})\) 是半单李代数 \(\mathfrak{g}\) 的泛包络代数的 \(q\)-形变（Drinfeld-Jimbo，1985）。它是 Hopf 代数，由 Chevalley 生成元 \(E_i, F_i, K_i^{\pm 1}\) 和 \(q\)-Serre 关系定义。

\textbf{定义 219.4}（R-矩阵 / 泛 R-矩阵）：\(U_q(\mathfrak{g})\) 的 \textbf{泛 R-矩阵} \(\mathcal{R} \in U_q(\mathfrak{g}) \hat{\otimes} U_q(\mathfrak{g})\) 满足： - \(\mathcal{R} \Delta(x) = \Delta^{\operatorname{op}}(x) \mathcal{R}\)（对所有 \(x \in U_q(\mathfrak{g})\)） - \((\Delta \otimes \operatorname{id})(\mathcal{R}) = \mathcal{R}_{13} \mathcal{R}_{23}\) - \((\operatorname{id} \otimes \Delta)(\mathcal{R}) = \mathcal{R}_{13} \mathcal{R}_{12}\)

R-矩阵给出 Yang-Baxter 方程的解：\(\mathcal{R}_{12} \mathcal{R}_{13} \mathcal{R}_{23} = \mathcal{R}_{23} \mathcal{R}_{13} \mathcal{R}_{12}\)。

\textbf{定理 219.3}（量子群构造辫群表示）：对每个 \(U_q(\mathfrak{g})\) 的有限维表示 \(V\)，\(R = \tau \circ \mathcal{R}|_{V \otimes V}\)（\(\tau\) 是交换映射）满足辫关系。这给出辫群 \(B_n\) 在 \(V^{\otimes n}\) 上的表示，从而通过 Markov 迹得到链环不变量。

\emph{证明}：泛 R-矩阵 \(\mathcal{R}\) 满足的六边形公理蕴含 \(\mathcal{R}_{12}\mathcal{R}_{13}\mathcal{R}_{23} = \mathcal{R}_{23}\mathcal{R}_{13}\mathcal{R}_{12}\)（Yang-Baxter 方程）。定义 \(R = \tau \circ \mathcal{R}\)，则 \(R\) 在 \(V^{\otimes 3}\) 上自动满足 \(R_{12}R_{23}R_{12} = R_{23}R_{12}R_{23}\)。令 \(\sigma_i \mapsto I^{\otimes (i-1)} \otimes R \otimes I^{\otimes (n-i-1)}\)，此直接给出 \(B_n \to \operatorname{GL}(V^{\otimes n})\)。若存在 Markov 迹 \(\operatorname{Tr}_q: \operatorname{End}(V^{\otimes n}) \to \mathbb{C}\)，则 \(\operatorname{Tr}_q(\rho(b))\) 是闭包不变量。\(\blacksquare\)

\subsection{219.3 拓扑量子场论与 Witten 不变量}\label{ux62d3ux6251ux91cfux5b50ux573aux8bbaux4e0e-witten-ux4e0dux53d8ux91cf}

\textbf{定义 219.5}（拓扑量子场论 / TQFT，Atiyah 1988）：\textbf{拓扑量子场论}是一个函子 \(Z: \operatorname{Bord}_n \to \operatorname{Vect}\)，将 \(n\) 维流形范畴映射到向量空间范畴。具体地： - 每个闭 \((n-1)\) 维流形 \(\Sigma\) 对应向量空间 \(Z(\Sigma)\) - 每个 \(n\) 维配边（cobordism）\(M: \Sigma_0 \to \Sigma_1\) 对应线性映射 \(Z(M): Z(\Sigma_0) \to Z(\Sigma_1)\) - 满足粘合公理：\(Z(M_1 \cup_\Sigma M_2) = Z(M_2) \circ Z(M_1)\)

\textbf{定义 219.6}（Witten-Chern-Simons 理论，1989）：Edward Witten 在 Chern-Simons 规范理论（作用量 \(S_{CS}(A) = \frac{k}{4\pi} \int_M \operatorname{Tr}(A \wedge dA + \frac{2}{3} A \wedge A \wedge A)\)）中，路径积分给出 3 维流形 \(M\) 的不变量（Witten 不变量）：

\[Z_k(M) = \int \mathcal{D}A \, e^{i k S_{CS}(A)}\]

其中积分沿 \(M\) 上的所有联络。Witten 不变量推广了 Jones 多项式。

\textbf{定理 219.4}（Reshetikhin-Turaev 不变量，1991）：Reshetikhin 和 Turaev 用量子群严格构造了 Witten 不变量（RT 不变量），通过 Kirby 演算和量子 \(6j\)-符号。

\textbf{证明}：RT 构造分两步：（1）将 3 维流形 \(M\) 表示为 \(S^3\) 中带框架链环 \(L\) 的手术。选择量子群 \(U_q(\mathfrak{sl}_2)\) 在 \(q = e^{2\pi i/r}\) 处的半单商。（2）对 \(L\) 的每个分量分配由量子维数和共形权重决定的''颜色''，通过量子 \(6j\)-符号和辫群表示计算链环的量子不变量，再乘以 Kirby 演算下的校正因子以消除手术表示的歧义。整体构成拓扑不变的 RT 不变量。\(\blacksquare\)

\subsection{219.4 范畴化与同调不变量}\label{ux8303ux7574ux5316ux4e0eux540cux8c03ux4e0dux53d8ux91cf}

\textbf{定义 219.7}（Khovanov 同调 / 范畴化原理）：\textbf{范畴化}（categorification）是将数值或多项式不变量提升为同调理论的过程。Khovanov 同调将 Jones 多项式范畴化，\textbf{Ozsváth-Szabó 同调}（Knot Floer homology）将 Alexander 多项式范畴化。

\textbf{定义 219.8}（Knot Floer 同调 / \(\widehat{\operatorname{HFK}}\)）：对纽结 \(K \subset S^3\)，\textbf{Knot Floer 同调} \(\widehat{\operatorname{HFK}}(K)\) 是双分次（Alexander 分次和 Maslov 分次）的有限维向量空间，满足：

\[\sum_{i,j} (-1)^i t^j \dim \widehat{\operatorname{HFK}}_j(K, i) = \Delta_K(t)\]

\hrulefill

\hrulefill

\hrulefill

\hrulefill

\hrulefill

\hrulefill

\section{第253章：Floer 同调}\label{ux7b2c253ux7ae0floer-ux540cux8c03}

Floer 同调是 Andreas Floer 在 1980 年代创立的无穷维 Morse 理论，它革命性地将变分法和椭圆偏微分方程（尤其是 anti-self-dual Yang-Mills 方程和 Cauchy-Riemann 方程）引入低维拓扑，从而为三维和四维流形构造了强大的同调不变量。Floer 的核心洞见是：在三维流形 \(Y\) 上，Chern-Simons 泛函给出了联络空间上的一个圆值函数，其临界点对应于平坦联络，而其中从一个临界点到另一个临界点的梯度流线（在适当的扰动后）恰好满足四维柱体 \(Y \times \mathbb{R}\) 上的 anti-self-dual Yang-Mills 方程。这一构造将 Morse 同调推广为无穷维的 Instanton Floer 同调。随后发展出了多种 Floer 同调版本------Heegaard Floer 同调（Ozsvath-Szabo，用组合拓扑避免了分析困难）、Monopole Floer 同调（Kronheimer-Mrowka，用 Seiberg-Witten 方程）以及 Lagrangian Floer 同调。Floer 同调在第 191 章量子不变量和第 193 章四维流形的 Donaldson 理论之间搭建了桥梁：它同时是 Witten 关于 3-4 维拓扑量子场论的数学实现，也是四维流形光滑结构分类的核心工具。

\subsection{220.1 Morse 同调与 Floer 同调的基本思想}\label{morse-ux540cux8c03ux4e0e-floer-ux540cux8c03ux7684ux57faux672cux601dux60f3}

\textbf{定义 220.1}（Floer 同调的基本构造）：Floer 同调 \(HF(M)\) 是无穷维空间上的 Morse 同调。构造步骤： 1. 定义无穷维流形 \(\mathcal{C}\)（联络空间 / 映射空间） 2. 定义 Morse 函数 \(\mathcal{F}: \mathcal{C} \to \mathbb{R}\)（Chern-Simons 泛函 / 辛作用泛函） 3. 链复形由 \(\mathcal{F}\) 的临界点生成 4. 边界算子计算梯度流线（连接临界点的 PDE 解）的计数

\textbf{定义 220.2}（Floer 同调的数学要求）：需要证明： - 横向性：梯度流线的模空间是光滑流形（通过加微扰） - 紧致性：模空间可紧化（通过 bubbling 分析和 Uhlenbeck 紧致性） - 粘合：边界算子满足 \(\partial^2 = 0\)（模空间边界对应断裂轨道）

\subsection{220.2 Instanton Floer 同调}\label{instanton-floer-ux540cux8c03}

\textbf{定义 220.3}（Instanton Floer 同调，Floer 1988）：设 \(Y\) 是闭 3 维流形（整数同调球）。\(\mathcal{C}\) 是 \(Y\) 上 \(\operatorname{SU}(2)\) 联络的规范等价类空间。Chern-Simons 泛函为

\[\operatorname{CS}(A) = \frac{1}{8\pi^2} \int_Y \operatorname{Tr}(A \wedge dA + \frac{2}{3} A \wedge A \wedge A)\]

\textbf{Instanton Floer 同调} \(I_*(Y)\) 由 \(\operatorname{CS}\) 的临界点（平坦联络）生成，边界算子由反自对偶瞬子（ASD instanton）的连接数给出。

\textbf{定理 220.1}（Instanton Floer 同调的性质）：\(I_*(Y)\) 是 \(\mathbb{Z}/8\mathbb{Z}\) 分次的。对整数同调球 \(Y\)，\(\chi(I_*(Y)) = |H_1(Y; \mathbb{Z})|\)（Casson 不变量）。

\textbf{证明}：\(I_*(Y)\) 的分次来自 \(Y \times \mathbb{R}\) 上瞬子的指标（模 8 由 Atiyah-Singer 指标定理给出）。Euler 示性数等于 Casson 不变量这一事实由 Taubes 证明：Casson 不变量可以由平坦 \(\operatorname{SU}(2)\) 联络的计数定义，而 \(I_*(Y)\) 的生成元恰好是这些平坦联络。\(\blacksquare\)

\subsection{220.3 Heegaard Floer 同调}\label{heegaard-floer-ux540cux8c03}

\textbf{定义 220.4}（Heegaard Floer 同调，Ozsváth-Szabó 2001-2004）：\textbf{Heegaard Floer 同调} \(\widehat{\operatorname{HF}}(Y)\) 由 \(Y\) 的 Heegaard 分解 \((U, \Sigma_g, V)\) 构造。生成元是 \(g\) 个点在 \(\Sigma_g\) 上的有序组（\(\operatorname{Sym}^g(\Sigma_g)\) 中的交点）。边界算子计算伪全纯圆盘的计数。

\textbf{定理 220.2}（Heegaard Floer 同调的优势）：\(\widehat{\operatorname{HF}}(Y)\) 是有限维 \(\mathbb{Z}/2\mathbb{Z}\) 分次的向量空间，满足： - 计算方便（比 Instanton Floer 同调更易计算） - \(\chi(\widehat{\operatorname{HF}}(Y)) = \pm |H_1(Y; \mathbb{Z})|\)（若 \(b_1(Y) = 0\)） - 有以下更精细的版本：\(\operatorname{HF}^+(Y)\)、\(\operatorname{HF}^-(Y)\)、\(\operatorname{HF}^\infty(Y)\)

\textbf{证明}：\(\widehat{\operatorname{HF}}(Y)\) 的构造使用 \(\operatorname{Sym}^g(\Sigma_g)\) 中 Lagrangian 子流形的 Lagrangian Floer 同调。有限维性由 Lagrangian Floer 同调的一般理论保证。Euler 示性数公式由 Floer 复形与 \(T^g\) 上 Morse 理论的联系导出。更精细版本通过引入多项式变元 \(U\) 和不同''帽子''构造得到，它们对应于不同的相对分次截断。\(\blacksquare\)

\textbf{定义 220.5}（Knot Floer 同调与 Heegaard Floer 的关系）：Knot Floer 同调 \(\widehat{\operatorname{HFK}}(K)\) 是 Heegaard Floer 同调的纽结版本，通过标定 Heegaard 图中的纽结轨道得到。

\textbf{定理 220.3}（Heegaard Floer 同调检测属数，Ozsváth-Szabó 2004）：\(\widehat{\operatorname{HF}}(Y)\) 检测 \(Y\) 的 Thurston 半范数。特别地，纽结 \(K\) 的亏格（genus）可由 Knot Floer 同调读取：

\[g(K) = \max\{i : \widehat{\operatorname{HFK}}(K, i) \neq 0\}\]

\textbf{证明}：Ozsváth 和 Szabó 构造了从 Knot Floer 同调到 sutured Floer 同调的''手术正合三角形''，将 \(\widehat{\operatorname{HFK}}\) 的 Alexander 分次与 \(S^3 \setminus K\) 中某些曲面的 Thurston 范数联系起来。过滤的极大非平凡分次恰好对应于 \(K\) 的极小亏格 Seifert 曲面，由 \(\widehat{\operatorname{HFK}}\) 在绝对分次下的''支撑''决定。\(\blacksquare\)

\subsection{220.4 其他 Floer 同调理论}\label{ux5176ux4ed6-floer-ux540cux8c03ux7406ux8bba}

\textbf{定义 220.6}（Monopole Floer 同调 / Seiberg-Witten Floer 同调，Kronheimer-Mrowka 2007）：基于 Seiberg-Witten 方程（而非反自对偶方程）的 Floer 同调。与 Heegaard Floer 同调同构（由 Kutluhan-Lee-Taubes 和 Colin-Ghiggini-Honda 证明）。

\textbf{定义 220.7}（Lagrangian Floer 同调，Floer 1988）：辛流形 \((M, \omega)\) 中两个 Lagrangian 子流形 \(L_0, L_1\) 的交点生成链复形，边界算子由伪全纯带的计数给出。是镜对称（mirror symmetry）和 Fukaya 范畴的核心工具。

\textbf{定义 220.8}（Sutured Floer 同调，Juhász 2006）：将 Heegaard Floer 同调推广到带缝合的 3 维流形（sutured manifolds）。应用于纽结补空间和接触几何。

\hrulefill

\hrulefill

\hrulefill

\hrulefill

\hrulefill

\hrulefill

\section{第254章：四维流形}\label{ux7b2c254ux7ae0ux56dbux7ef4ux6d41ux5f62}

四维流形在微分拓扑中占有独特而神秘的地位：它是唯一一个允许同一拓扑流形装备无穷多个互不微分同胚的光滑结构的维数。这一惊人事实由 Simon Donaldson 在 1983 年通过研究 anti-self-dual Yang-Mills 方程的模空间（兼用代数几何方法）发现------他证明了某些紧致单连通四维流形上的光滑结构不是唯一的（与高维的 Cerf 定理形成鲜明对比）。Donaldson 理论通过构造 ASD 模空间的 Donaldson 多项式不变量，可以区分 \(\mathbb{CP}^2 \# n\overline{\mathbb{CP}}^2\) 等具体流形的不同光滑结构。随后，Seiberg-Witten 方程（1994 年引入）以其技术简单性（单值 PDE 替代了 Donaldson 的非线性方程组）给出了一套等价的不变量 \(SW(M)\)，极大地简化了四维流形的光滑分类。四维流形理论与其他章节的交叉极为丰富：Floer 同调（第 192 章）的 Instanton 版本直接为 Donaldson 不变量提供了切割-粘合公式（Floer 精确三角形），而纽结理论（第 190 章）则通过切片亏格（slice genus）与四维流形中嵌入曲面的亏格产生联系。

\subsection{221.1 四维流形的特殊性}\label{ux56dbux7ef4ux6d41ux5f62ux7684ux7279ux6b8aux6027}

\textbf{定理 221.1}（光滑结构的维数依赖）： - 维数 \(n \neq 4\)：紧致拓扑流形只有有限个光滑结构（到 \(n = 3\) 是唯一的，\(n \geq 5\) 由 Kirby-Siebenmann 分类） - 维数 \(n = 4\)：存在无穷多个不同的光滑结构（如 \(\mathbb{R}^4\) 上有不可数无穷个不同的光滑结构）

\textbf{证明}：\(n \geq 5\) 时，Kirby-Siebenmann 理论通过手术正合序列表明光滑结构由拓扑结构和 Kirby-Siebenmann 不变量确定，故有限。\(n=3\) 由 Moise 定理（PL=光滑=拓扑）。\(n=4\) 时，Freedman 定理给出光滑 Kirby-Siebenmann 障碍的失效，而 Donaldson 定理限制相交形式，两者结合产生 exotic 结构。\(\blacksquare\)

\textbf{定义 221.1}（相交形式 / Intersection Form）：闭可定向光滑 4 维流形 \(X\) 的\textbf{相交形式} \(Q_X: H^2(X; \mathbb{Z}) \times H^2(X; \mathbb{Z}) \to \mathbb{Z}\) 定义为 \(Q_X(\alpha, \beta) = \langle \alpha \cup \beta, [X] \rangle\)。\(Q_X\) 是对称双线性型。

\textbf{定理 221.2}（Freedman 定理，1982）：单连通闭拓扑 4 维流形由其相交形式 \(Q_X\) 和 Kirby-Siebenmann 不变量 \(\operatorname{ks}(X) \in \mathbb{Z}/2\mathbb{Z}\) 完全分类。任何幺模对称双线性型都可实现为某单连通拓扑 4 维流形的相交形式。

\textbf{证明}：Freedman 使用 Casson 柄体理论证明，在拓扑范畴中，无穷的 Casson 柄体可收缩为标准的 2-柄体。核心引理是''5 维 h-配边定理在 4 维拓扑范畴中成立''。存在性证明通过将给定的相交形式实现为恰当分解（handle decomposition）然后处理交点的 2-柄体取消问题。\(\blacksquare\)

\subsection{221.2 Donaldson 理论}\label{donaldson-ux7406ux8bba}

\textbf{定理 221.3}（Donaldson 定理，1983）：设 \(X\) 是闭光滑单连通 4 维流形。若 \(X\) 的相交形式 \(Q_X\) 是正定的，则 \(Q_X\) 在 \(\mathbb{Z}\) 上等价于对角形式 \(\langle 1 \rangle \oplus \cdots \oplus \langle 1 \rangle\)。即正定相交形式必须是标准的。

\textbf{证明}：在 \(X\) 的 \(\operatorname{SU}(2)\)-瞬子模空间上应用指标定理和 Uhlenbeck 紧致性。瞬子模空间维数为 \(8c_2 - 3(1 - b_1 + b_2^+)\)，正定时模空间在适当截去后是紧光滑流形。若 \(Q_X\) 非对角型（如含 \(E_8\)），则模空间中会出现奇点，与光滑性矛盾。通过精细分析可排除所有非标准正定形式。\(\blacksquare\)

\textbf{推论 221.1}（exotic 4 维流形的存在）：Freedman 定理允许 \(E_8\) 形式作为拓扑 4 维流形的相交形式，但由 Donaldson 定理，\(E_8\) 形式不能实现为光滑 4 维流形的相交形式。这证明了存在拓扑等价但非微分同胚的 4 维流形。

\textbf{定义 221.2}（Donaldson 不变量）：Donaldson 通过反自对偶联络（ASD instanton）的模空间定义了 4 维流形 \(X\) 的\textbf{Donaldson 多项式不变量} \(D_X: \operatorname{Sym}_*(H_2(X; \mathbb{Z})) \to \mathbb{Q}\)。这些不变量是光滑结构的不变量。

\textbf{定理 221.4}（Donaldson 不变量对光滑结构的敏感性）：存在同胚但不微分同胚的 4 维流形，其 Donaldson 不变量不同。Donaldson 不变量证明了 \(\mathbb{C}P^2 \# k \overline{\mathbb{C}P^2}\) 上存在无穷多个不同的光滑结构。

\textbf{证明}：Donaldson 不变量 \(D_X\) 取决于瞬子模空间的拓扑（其取值为多线性多项式）。对同胚流形 \(X_1, X_2\)，Freedman 确保它们拓扑等价，但 \(D_{X_1} \neq D_{X_2}\) 证明它们不能微分同胚。具体构造：通过对 \(\mathbb{C}P^2 \# n \overline{\mathbb{C}P^2}\) 作对数变换（logarithmic transform）得到不同 Donaldson 不变量的流形。\(\blacksquare\)

\subsection{221.3 Seiberg-Witten 理论}\label{seiberg-witten-ux7406ux8bba}

\textbf{定义 221.3}（Seiberg-Witten 方程，1994）：设 \(X\) 是闭光滑 4 维流形，带有 \(\operatorname{Spin}^c\) 结构 \(\mathfrak{s}\)。\textbf{Seiberg-Witten 方程}为：

\[\begin{cases} D_A \Phi = 0 \\ F_A^+ = \Phi \otimes \Phi^* - \frac{1}{2} |\Phi|^2 \operatorname{id} \end{cases}\]

其中 \(\Phi\) 是正旋量，\(A\) 是 \(\operatorname{Spin}^c\) 联络，\(D_A\) 是 Dirac 算子，\(F_A^+\) 是 \(A\) 的曲率的自对偶部分。

\textbf{定义 221.4}（Seiberg-Witten 不变量）：Seiberg-Witten 方程的解空间 \(M(\mathfrak{s})\) 是紧致光滑流形。\textbf{Seiberg-Witten 不变量} \(\operatorname{SW}_X(\mathfrak{s})\) 是 \(M(\mathfrak{s})\) 上 Euler 类的积分，是 \(\operatorname{Spin}^c\) 结构 \(\mathfrak{s}\) 的函数。

\textbf{定理 221.5}（Seiberg-Witten 不变量的性质）： - 只有有限个 \(\operatorname{Spin}^c\) 结构 \(\mathfrak{s}\) 满足 \(\operatorname{SW}_X(\mathfrak{s}) \neq 0\)（基本类） - 若 \(X\) 容许正标量曲率度量，则 \(\operatorname{SW}_X(\mathfrak{s}) = 0\)（对所有 \(\mathfrak{s}\)） - 对 Kähler 曲面，\(\operatorname{SW}_X\) 与规范类（canonical class）直接相关

\textbf{证明}：（1）SW 方程解的模空间维数为 \(d = \frac{1}{4}(c_1(\mathfrak{s})^2 - 2\chi - 3\sigma)\)，仅当 \(d \geq 0\) 时有意义，此条件约束有限。若 SW 不变量非零，则 \(c_1(\mathfrak{s})^2 \geq 2\chi + 3\sigma\)，由相交形式的有限型性质知有限。（2）由 Weitzenbock 公式 \(0 = D_A^* D_A \Phi = \nabla_A^* \nabla_A \Phi + \frac{R}{4}\Phi + \cdots\)，正标量曲率 \(\Rightarrow \Phi \equiv 0\)。（3）对 Kähler 曲面，\(SW_X(\pm c_1(X)) = \pm 1\)。\(\blacksquare\)

\textbf{定理 221.6}（Witten 猜想 / Taubes 定理）：Seiberg-Witten 不变量等价于 Donaldson 不变量（通过 Witten 的物理论证和后来的数学证明）。Seiberg-Witten 理论的计算比 Donaldson 理论简单得多。

\textbf{证明}：Witten 通过 \(\operatorname{SU}(2)\) 的超对称 Yang-Mills 理论与 Seiberg-Witten 理论之间的对偶性论证了两组不变量等价。严格数学证明由 Feehan-Leness 完成，他们构造了 \(\operatorname{PU}(2)\)-Monopole 方程（插值 \(\operatorname{U}(2)\) 瞬子与 Seiberg-Witten 之间的理论），证明在适当参数极限下 Monopole 模空间逼近 Donaldson 模空间，从而建立两组不变量的恒等关系。\(\blacksquare\)

\subsection{221.4 Exotic 结构与阻塞理论}\label{exotic-ux7ed3ux6784ux4e0eux963bux585eux7406ux8bba}

\textbf{定义 221.5}（Exotic \(\mathbb{R}^4\)）：存在与 \(\mathbb{R}^4\)（标准欧氏 4 维空间）同胚但不微分同胚的 4 维流形。事实上，存在不可数无穷个不同的 exotic \(\mathbb{R}^4\)（Taubes 1987，Gompf 1985）。

\textbf{定理 221.7}（\(\frac{11}{8}\) 猜想 / Furuta 定理）：对闭光滑单连通 4 维流形 \(X\) 满足 \(b_2^+(X) > 0\)，有 \(b_2(X) \geq \frac{11}{8} |\sigma(X)|\)（\(\sigma\) 是符号差）。\(\frac{11}{8}\) 猜想是最优下界，由 Furuta（2001）用等变 K 理论证明其弱形式。

\textbf{证明}：Furuta 将 Seiberg-Witten 方程转化为 \(S^1\)-等变 \(K\)-理论中的稳定同伦问题。核心是研究 \(S^1\)-等变映射 \(\mu: \mathbb{C}^{\infty} \to \mathbb{C}^{\infty}\) 的稳定同伦类，利用 Adams 的 \(e\)-不变量和 Conley 指标证明当 \(b_2^+ > 0\) 时若 \(b_2^-\) 不够大则稳定同伦类不可能存在，与 SW 不变量非平凡矛盾。得 \(10/8\) 下界；\(\frac{11}{8}\) 的最优性由 \(K3\) 曲面给出。\(\blacksquare\)

\textbf{定义 221.6}（10/8 定理 / Furuta 不等式）：\(b_2(X) \geq \frac{10}{8} |\sigma(X)| + 2\)（Furuta 1996 的 \(10/8\) 定理给出了部分证明，完整 \(\frac{11}{8}\) 猜想仍未完全解决）。

\hrulefill

\hrulefill

\hrulefill

\hrulefill

\hrulefill

\hrulefill

\subsection{194.A 低维拓扑补充}\label{a-ux4f4eux7ef4ux62d3ux6251ux8865ux5145}

几何群论（Geometric Group Theory）是 Gromov（1987）开创的现代数学分支------将有限生成群视为带字度量的几何对象，用拟等距不变量分类群。Princeton Companion 将其列为独立分支。它连接了代数、拓扑、几何和动力系统。

\subsection{Gromov双曲群}\label{gromovux53ccux66f2ux7fa4}

\textbf{定义}（Cayley图与字度量）：设 \(G\) 是有限生成群，\(S\) 是对称生成元集。\textbf{Cayley图} \(\Gamma(G, S)\) 的顶点是群元素，\(g\) 与 \(gs\)（\(s \in S\)）之间有边。\textbf{字度量} \(d(g, h)\) 是连接 \(g\) 和 \(h\) 的最短路径长（字长 \(|g^{-1}h|_S\)）。Cayley图上的不同生成元集的字度量是拟等距的------这使得大量群的性质成为拟等距不变量。

\textbf{定义}（Gromov δ-双曲空间与双曲群）：度量空间 \((X, d)\) 的测地三角形 \(\triangle\) 称为 \textbf{δ-slim}，如果每条边落在另外两条边的 δ-邻域内。\(X\) 是 \textbf{δ-双曲空间}（Gromov 双曲空间），若所有测地三角形的 slim 参数有上确界 δ。\textbf{双曲群}是 Cayley 图为 Gromov 双曲空间的有限生成群。

\textbf{例}：有限群（平凡地）、自由群 \(\mathbb{F}_n\)（\(n \geq 2\)）、基本群 \(\pi_1(S)\)（\(S\) 为亏格 \(\geq 2\) 的紧曲面）、经典 Schottky 群均为此类。\(\mathbb{Z}^n\)（\(n \geq 2\)）不是双曲的（平面 Cayley 图的三角形不是 slim 的，Bass-Serre 理论的树形结构不成立）。

\textbf{定理}（Gromov双曲群的边界 / Gromov边界）：双曲群 \(G\) 的 \textbf{Gromov 边界} \(\partial G\) 是 \(G\) 作用为收敛群的紧可度量空间。\(\partial G\) 是 \(G\) 的重要拟等距不变量。例：自由群 \(\mathbb{F}_n\) 的边界是其 Cayley 树（\(2n\)-正则树）的 Cantor 边界；\(\Gamma\)-曲面群（亏格 ≥ 2）的边界是 \(S^1\)。

\textbf{证明}：将 \(G\) 中发散序列的等价类定义为边界点：\(g_n \to \infty\) 且 Gromov 积 \((g_n|g_m)_e \to \infty\) 时两序列等价。赋以可视度量 \(d_\varepsilon(\xi,\eta) = e^{-\varepsilon (\xi|\eta)_e}\)，由双曲条件可证 \((\partial G, d_\varepsilon)\) 是紧可度量空间。\(G\) 通过左平移连续作用，此作用为收敛群作用（对紧致不连续作用商紧）。\(\blacksquare\)

\textbf{定理}（拟等距刚性）：双曲群的拟等距（quasi-isometry）一定距离一个群同构（在有限距离和有限指数误差下）。特别地，两个双曲群拟等距当且仅当它们是可通约的（commensurable）------即存在有限指数子群互相同构。

\textbf{证明}：拟等距 \(f: G \to H\) 诱导边界同胚 \(\partial f: \partial G \to \partial H\)。\(\partial G\) 上的拟共形结构由 \(G\) 的双曲几何完全确定。Mostow 刚性论证（推广到双曲群）表明 \(\partial f\) 必共形，从而 \(f\) 在有限误差内是群同构。若 \(f\) 是拟等距，则 \(G\) 和 \(H\) 的 Cayley 图拟等距，由此推得存在有限指数子群 \(G' \subset G\)，\(H' \subset H\) 使得 \(G' \cong H'\)。\(\blacksquare\)

\subsection{拟等距与拟等距不变量}\label{ux62dfux7b49ux8dddux4e0eux62dfux7b49ux8dddux4e0dux53d8ux91cf}

\textbf{定义}（拟等距 / \((K, C)\)-拟等距）：映射 \(f: X \to Y\)（测地度量空间之间）是 \((K, C)\)-\textbf{拟等距}，若对所有 \(x_1, x_2 \in X\)： \[\frac{1}{K} d_X(x_1, x_2) - C \leq d_Y(f(x_1), f(x_2)) \leq K d_X(x_1, x_2) + C\] 且 \(Y\) 落在 \(f(X)\) 的 \(C\)-邻域内。拟等距是大尺度几何的''等价关系''------双曲平面 \(\mathbb{H}^2\) 与 \(4\)-正则树拟等距。

\textbf{命题}（基本拟等距不变量）： - 群的端（ends）数 \(e(G) \in \{0, 1, 2, \infty\}\) 是拟等距不变量（Freudenthal-Hopf） - 多项式增长次数是拟等距不变量（Gromov 1981） - 双曲性（δ-双曲空间性质）是拟等距不变量 - 渐近维数（asymptotic dimension）\(\operatorname{asdim}(G)\) 是拟等距不变量

\textbf{定理}（群的端与Stallings分类定理，1968-1971）：有限生成群 \(G\) 的端数： - \(e(G) = 0\) 当且仅当 \(G\) 有限 - \(e(G) = 1\) 当且仅当 \(G\) 有一端（如 \(\mathbb{Z}^n\)，\(n \geq 2\)；\(\pi_1(\)亏格 \geq 2 曲面\()\)） - \(e(G) = 2\) 当且仅当 \(G\) 是实质上的 \(\mathbb{Z}\)（有有限指数无限循环子群） - \(e(G) = \infty\) 当且仅当 \(G\) 可分解为非平凡自由积并合（amalgamated free product）或 HNN 扩张

\textbf{证明}：\(e(G)\) 定义为 \(\varprojlim \pi_0(\operatorname{Cay}(G,S) \setminus B_R)\) 的基数的上确界。若 \(e(G) \geq 2\)，则 \(G\) 作用在端对上推出作用在树上。Stallings 利用 Dunwoody 的可达性定理和 Bass-Serre 理论证明：\(e(G) \geq 2\) 等价于 \(G\) 分裂为自由积并合或 HNN 扩张；\(e(G) = \infty\) 当且仅当 \(G\) 本质上不可约地无限分裂。\(\blacksquare\)

\textbf{定理}（Gromov的增长定理，1981）：有限生成群的多项式增长等价于其实质上幂零（virtually nilpotent）。这是群论与渐进几何结合的里程碑成果。

\subsection{CAT(0)空间与群}\label{cat0ux7a7aux95f4ux4e0eux7fa4}

\textbf{定义}（CAT(0)空间）：完备测地度量空间 \(X\) 是 \textbf{CAT(0)}（非正曲率），若任意两点间存在唯一的测地线，且每个测地三角形比 Euclid 比较三角形更''薄''：对任意三点 \(x, y, z\) 和比较三角形中的对应点 \(\bar{x}, \bar{y}, \bar{z} \subset \mathbb{R}^2\)，有 \(d(x, y) \leq |\bar{x} - \bar{y}|\) 等。

\textbf{定义}（CAT(0)群 / Hadamard群）：可能一地作用在 CAT(0) 空间上的群。例子包括：自由群、\(\mathbb{Z}^n\)、曲面群、Coxeter 群、正曲率 Cartan 矩阵的 Kac-Moody 群。

\textbf{定理}（Cartan-Hadamard定理的度量版本）：CAT(0) 空间可缩------特别地，CAT(0) 群是有限型（即存在分类空间 \(K(G, 1)\) 由有限 CW 复形实现）。

\textbf{定理}（Tits的刚性替代定理 / Tits alternative）：有限生成线性群要么实质上可解（virtually solvable），要么包含非 Abel 自由子群 \(\mathbb{F}_2\)。Gromov 双曲群版本：双曲群要么实质循环，要么包含 \(\mathbb{F}_2\)。

\subsection{Bass-Serre理论与群图}\label{bass-serreux7406ux8bbaux4e0eux7fa4ux56fe}

\textbf{定义}（群作用在树上的Bass-Serre理论，1977）：群 \(G\) 作用在树 \(T\) 上（无边反转）当且仅当 \(G\) 可分解为图群（graph of groups）的基本群。这是研究群的代数结构的核心工具。

\begin{itemize}
\tightlist
\item
  \textbf{并合自由积}：\(G = A *_C B\)（\(C\) 为 \(A\) 和 \(B\) 的公共子群）以单边图表示
\item
  \textbf{HNN扩张}：\(G = \langle H, t \mid t^{-1} C t = \phi(C) \rangle\)（以稳定子 \(C\) 和同构 \(\phi: C \to H\) 定义）以单环图表示
\end{itemize}

\textbf{定理}（Kurosh子群定理）：自由积的子群仍是自由积。

\textbf{定义}（JSJ分解 / JSJ Decomposition，Jaco-Shalen-Johannson 1979）：3维流形的几何化分解可视为图群分解的特例------沿环面和 Klein 瓶的切割（扭元素对应 Seifert 纤维化片段，非扭对应双曲片段）。这被 Sela（1997）推广到任意有限出现群（tower argument）。

\hrulefill

\hrulefill

\hrulefill

\hrulefill

\hrulefill

\hrulefill

\subsection{195.A 低维拓扑补充2}\label{a-ux4f4eux7ef4ux62d3ux6251ux8865ux51452}

微分拓扑研究光滑流形上的整体性质------特别是光滑映射的横截性、流形上的Morse函数和光滑流形的配边分类。它是连接代数拓扑、微分几何和几何分析的桥梁。

\hrulefill

\subsection{横截性理论}\label{ux6a2aux622aux6027ux7406ux8bba}

横截性是微分拓扑最基本的概念之一：光滑映射''一般地''与子流形横截相交。它由Thom（1954）和Smale（1956）系统化。

\textbf{定义}（横截性）：设 \(f: M \to N\) 是光滑映射，\(Z \subset N\) 是光滑子流形。\(f\) 称为\textbf{横截于} \(Z\)（记作 \(f \pitchfork Z\)），如果对任意 \(x \in f^{-1}(Z)\)， \[T_{f(x)}N = df_x(T_xM) + T_{f(x)}Z\] 即 \(f\) 的像与 \(Z\) 的切空间张成整个 \(T_{f(x)}N\)。

\textbf{定理}（Thom横截性定理）：设 \(M, N\) 是光滑流形，\(Z \subset N\) 是闭子流形。则使 \(f \pitchfork Z\) 的光滑映射 \(f: M \to N\) 在 \(C^\infty(M, N)\) 中构成稠密子集（在 Whitney \(C^\infty\) 拓扑下）。若 \(M\) 紧，它们是开稠密集。\emph{注：这是微分拓扑中最重要的''一般性''定理------几乎所有光滑映射都横截于任何给定子流形。}

\textbf{推论}（弱横截性定理）：任何连续映射 \(f: M \to N\) 可被光滑映射逼近，且该光滑映射横截于 \(Z\)。

\textbf{定理}（横截同伦定理）：设 \(f_0, f_1: M \to N\) 横截于闭子流形 \(Z\)，且存在横截于 \(Z\) 的同伦连接它们。则 \(f_0^{-1}(Z)\) 与 \(f_1^{-1}(Z)\) 是 \(M\) 中的配边子流形（在光滑配边意义下）。

\textbf{定义}（管状邻域）：光滑流形 \(M\) 中闭子流形 \(Z\) 的\textbf{管状邻域}是向量丛 \(\pi: E \to Z\) 的总空间的一个开邻域，微分同胚于 \(Z\) 在 \(M\) 中的邻域。存在性由Riemann度量的指数映射保证。

\hrulefill

\subsection{Morse理论}\label{morseux7406ux8bba}

Morse理论（Morse 1925-1934，Milnor 1963系统化）将光滑流形的拓扑与其上Morse函数的临界点联系起来------这是微分拓扑最有力的工具之一。

\textbf{定义}（Morse函数）：光滑函数 \(f: M \to \mathbb{R}\) 的\textbf{临界点} \(p\) 满足 \(df_p = 0\)（在局部坐标下所有一阶偏导数同时为零）。\(p\) 称为\textbf{非退化}的，如果 Hessian 矩阵 \(H_f(p) = \left(\frac{\partial^2 f}{\partial x_i \partial x_j}(p)\right)\) 非奇异（\(\det H_f(p) \neq 0\)）。\(f\) 称为\textbf{Morse函数}，若其所有临界点均非退化。

\textbf{引理}（Morse引理）：设 \(p\) 是 Morse 函数 \(f\) 的非退化临界点。则存在 \(p\) 附近的局部坐标系 \((x_1, \ldots, x_n)\)，使得 \[f(x) = f(p) - x_1^2 - \cdots - x_k^2 + x_{k+1}^2 + \cdots + x_n^2\] 其中 \(k\) 是 Hessian 的负特征值个数，称为临界点的\textbf{指标}（index）。\(k\) 是与坐标选取无关的不变量。

\textbf{定理}（Morse不等式）：设 \(M\) 是紧光滑流形，\(f: M \to \mathbb{R}\) 是 Morse 函数。记 \(C_k\) 为指标 \(k\) 的临界点个数，\(\beta_k = \operatorname{rank} H_k(M; \mathbb{R})\) 为第 \(k\) 个 Betti 数。则： 1. \textbf{弱Morse不等式}：\(C_k \geq \beta_k\)（对所有 \(k\)） 2. \textbf{强Morse不等式}：\(\sum_{j=0}^k (-1)^{k-j} C_j \geq \sum_{j=0}^k (-1)^{k-j} \beta_j\)（对所有 \(k\)） 3. \textbf{Morse等式}：\(\sum_{k=0}^n (-1)^k C_k = \sum_{k=0}^n (-1)^k \beta_k = \chi(M)\)（Euler示性数）

\emph{注：Morse不等式将分析量（临界点）与拓扑量（Betti数）联系起来，这是微分拓扑的核心思想。}

\textbf{定理}（Morse理论的胞腔分解）：紧流形 \(M\) 上的 Morse 函数 \(f\) 给出了 \(M\) 的一个 CW 复形结构------每个指标 \(k\) 的临界点对应一个 \(k\)-维胞腔。\(M\) 同伦等价于一个每维胞腔数等于该维临界点数的 CW 复形。

\textbf{定理}（Reeb定理）：若紧 \(n\) 维流形上存在恰有两个临界点的 Morse 函数，则该流形同胚于 \(S^n\)。

\textbf{定义}（Morse同调与Floer同调的关系）：Morse同调群的链复形由临界点生成，边界算子计数梯度流线（Morse-Smale条件）。这是Floer同调（无限维推广）的有限维原型。

\hrulefill

\subsection{配边理论}\label{ux914dux8fb9ux7406ux8bba}

配边理论（cobordism theory）由 Thom（1954）创立，是代数拓扑和微分拓扑的交汇点，也是Thom获1958年Fields奖的工作。

\textbf{定义}（配边 / cobordism）：两个闭 \(n\) 维光滑流形 \(M_0, M_1\) 称为\textbf{配边}的，如果存在紧 \((n+1)\) 维光滑流形 \(W\)（\textbf{配边流形}），使得 \(\partial W = M_0 \sqcup M_1\)。配边是等价关系，等价类全体记为 \(\Omega_n^{\mathrm{SO}}\)（定向配边群）。

\textbf{定理}（Thom配边同构定理，1954）：定向配边群 \(\Omega_n^{\mathrm{SO}}\) 同构于 Thom 谱的同伦群： \[\Omega_n^{\mathrm{SO}} \cong \pi_{n+k}(\operatorname{MSO}(k)) \quad (k \gg 0)\] 其中 \(\operatorname{MSO}(k)\) 是万有 \(k\) 维定向向量丛的 Thom 空间。这是将配边问题化为同伦论计算的核心定理。

\textbf{定理}（Thom配边环的计算）：有理数域上，定向配边环 \(\Omega_*^{\mathrm{SO}} \otimes \mathbb{Q}\) 由偶维复射影空间 \(\mathbb{CP}^{2k}\) 的配边类多项式生成： \[\Omega_*^{\mathrm{SO}} \otimes \mathbb{Q} \cong \mathbb{Q}[\mathbb{CP}^2, \mathbb{CP}^4, \mathbb{CP}^6, \ldots]\] （每个生成元对应于一个维数 \(4k\) 的生成元）。低维配边群：\(\Omega_0^{\mathrm{SO}} \cong \mathbb{Z}\)，\(\Omega_1^{\mathrm{SO}} = \Omega_2^{\mathrm{SO}} = \Omega_3^{\mathrm{SO}} = 0\)，\(\Omega_4^{\mathrm{SO}} \cong \mathbb{Z}\)（生成元 \([\mathbb{CP}^2]\)）。

\textbf{定义}（Stiefel-Whitney数与Pontryagin数的配边不变性）：两个闭流形配边当且仅当它们的所有 Stiefel-Whitney 数和 Pontryagin 数分别相等。这些示性数是配边不变量。

\textbf{定义}（\(h\)-配边定理（Smale, 1961 年 Fields 奖））：设 \(W\) 是紧光滑流形，\(\partial W = M_0 \sqcup M_1\)，且包含映射 \(M_0 \hookrightarrow W\) 和 \(M_1 \hookrightarrow W\) 都是同伦等价。若 \(\dim W \geq 5\)，则 \(W\) 微分同胚于 \(M_0 \times [0, 1]\)，从而 \(M_0\) 与 \(M_1\) 微分同胚。

\emph{意义}：\(h\)-配边定理是 Smale 证明 \(n \geq 5\) 维 Poincaré 猜想的核心工具。它说明了高维（\(\geq 5\)）流形的拓扑分类可以由同伦论确定------这是高维拓扑与低维拓扑的根本区别。

\textbf{定理}（\(s\)-配边定理，Kervaire-Milnor 1963）：在 \(h\)-配边定理的条件下，若进一步指定了切丛的平凡化，可定义 Kervaire-Milnor 群 \(\Theta_n\)（\(n\) 维同伦球面的微分同胚类）。\(\Theta_n\) 的精细计算揭示了 \(S^7\) 上存在 28 种不同的光滑结构（Milnor 奇异球面，1962 年获 Fields 奖）。

\hrulefill

\subsection{手术理论}\label{ux624bux672fux7406ux8bba}

手术理论（Surgery theory, Kervaire-Milnor 1963, Browder 1960s, Wall 1970s）是分类流形的系统代数拓扑框架。

\textbf{定义}（手术 / surgery）：在 \(n\) 维流形 \(M\) 上，一个 \(k\)-手术（\(0 \leq k \leq n\)）是指切除一个嵌入 \(S^k \times D^{n-k}\)，然后沿边界粘上 \(D^{k+1} \times S^{n-k-1}\)。手术改变流形的拓扑（特别是同伦型和同调群），保持边界不变。

\textbf{定义}（正规映射与手术障碍）：正规映射 \((f, b): M \to X\)（\(f\) 度数为1，\(b\) 是切丛的丛同构）的手术障碍落在 \(L\)-群 \(L_n(\pi_1(X))\)（Wall的代数\(L\)-群，依赖于基本群）。当手术障碍为零时，\(M\) 可通过手术变成同伦等价于 \(X\)。

\textbf{定理}（Browder-Novikov-Sullivan-Wall 手术正合序列 / 分类定理，1960s-1970s）： \[\cdots \to L_{n+1}(\pi) \to \mathcal{S}(X) \to [X, G/O] \to L_n(\pi)\] 此正合序列联系了流形 \(X\) 的手术结构集 \(\mathcal{S}(X)\)（同伦于 \(X\) 的流形的微分同胚类）、正规不变量（映射至 \(G/O\) 空间）和代数 \(L\)-理论。

\emph{意义}：手术理论是分类给定同伦型的流形的各种光滑/PL/拓扑结构的基本工具。它是高维拓扑分类的完成框架。

\hrulefill

\emph{本卷在拓扑学（V10）、代数拓扑（V14）和微分几何（V11）的基础上，系统介绍了低维拓扑和纽结理论：3 维流形的分解与几何化、纽结理论的多项式与同调不变量（Alexander、Jones、HOMFLY-PT、Khovanov）、量子拓扑不变量（量子群与 TQFT）、Floer 同调系列（Instanton、Heegaard、Monopole）和 4 维流形理论（Donaldson 理论、Seiberg-Witten 理论、exotic 结构）。共 5 章。}

\hrulefill

\emph{本卷在已有内容基础上，以代数 \(K\) 理论（K0/K1/K2 低阶导引、Milnor K 理论与 Bloch 群、母题上同调与 Bloch-Kato 猜想、拓扑与 C}-代数 K 理论）和低维拓扑与纽结理论（三维流形与几何化、纽结多项式与 Khovanov 同调、量子拓扑不变量、Floer 同调、四维流形）为补充单元，共 10 个附加章。*

\emph{卷十四：代数拓扑终。}

\emph{卷十四：代数拓扑终。}

\emph{卷十四：代数拓扑终。}

\emph{卷十七：代数拓扑终。} \#\# 第255章：低维流形专题（3维与4维） 低维拓扑的核心研究对象是 3 维和 4 维流形------它们在几何和分析中表现出与高维 (\(\geq 5\)) 截然不同的现象，需要专门的技术和理论。在 3 维和 4 维，Whitney 技巧失效，阻碍了高维手术理论的直接应用。本章聚焦于低维流形的核心定理和构造技术，从 Whitney 技巧的失效机制出发，依次讨论 Freedman 定理对拓扑 4 维流形的分类、Donaldson 定理对光滑结构的限制、以及 exotic \(\mathbb{R}^4\) 的构造思路。

\subsection{218.1 Whitney 技巧及其在四维的失效}\label{whitney-ux6280ux5de7ux53caux5176ux5728ux56dbux7ef4ux7684ux5931ux6548}

Whitney 技巧是高维微分拓扑中最核心的技术之一，也是理解低维流形为何特殊的关键。

\textbf{定义 218.1}（Whitney 圆盘）：设 \(M\) 为光滑 \(n\) 维流形，\(P^p\) 和 \(Q^q\) 为 \(M\) 中横截相交的紧致光滑子流形，\(p + q = n\)。设 \(x, y \in P \cap Q\) 为两个带有相反相交符号的交点（\(\varepsilon_x = -\varepsilon_y\)）。一条嵌入弧 \(\alpha \subset P\) 联结 \(x\) 到 \(y\)，一条嵌入弧 \(\beta \subset Q\) 联结 \(y\) 到 \(x\)，且 \(\alpha \cup \beta\) 围成 \(M\) 中一个嵌入圆盘 \(D^2\)，该圆盘内部与 \(P \cup Q\) 不相交，则 \(D^2\) 称为联结 \(x\) 和 \(y\) 的 \textbf{Whitney 圆盘}。

\textbf{定理 218.1}（Whitney 技巧，1944）：当 \(n \geq 5\) 时，在 Whitney 圆盘存在的条件下，可通过合痕消去 \(P\) 和 \(Q\) 的一对相反符号的交点。具体地，将 \(P\) 沿 Whitney 圆盘''推过'' \(Q\)，从而将 \(x\) 和 \(y\) 同时消除。

\textbf{证明概要}：利用 \(D^2\) 在 \(M\) 中的法丛。Whitney 圆盘 \(D^2\) 在 \(M\) 中的法丛维数为 \(n - 2\)，而 \(P\) 沿 \(D^2\) 边界的法丛维数为 \(n - p - 1\)。当 \(n - 2 \geq 3\)（即 \(n \geq 5\)）时，\(D^2\) 内部的法丛足够''宽敞''，可以构造与 \(D^2\) 边界上已给的标架相容的标架延拓。利用此标架族，构造合痕将 \(P\) 沿 \(D^2\) 滑动------在 \(\partial D^2\) 的 \(\alpha\) 部分处 \(P\) 被推向 \(D^2\) 的一侧，在 \(\beta\) 部分处被推向另一侧------从而将 \(P\) 与 \(Q\) 的两个交点同时消除。该构造的核心在于：\(D^2\) 是 2 维的，而要使法丛标架的延拓问题可解（\(\pi_1\) 的阻碍消失），需要 \(D^2\) 的余维数 \(\geq 3\)，即 \(n \geq 5\)。\(\blacksquare\)

\textbf{定理 218.2}（Whitney 技巧在 4 维的失效）：当 \(n = 4\) 时，Whitney 圆盘的法丛维数仅 \(2\)，不能保证标架延拓的存在性。这导致了 Whitney 技巧的失效------嵌入 2 维圆盘的孤立自交点在 4 维流形中不能通过一般位置的扰动来消除。这是 4 维光滑拓扑中几乎所有特异现象的根本原因。

\textbf{证明}：当 \(n = 4\) 时，Whitney 圆盘 \(D^2\) 在 \(M\) 中的余维数为 \(4 - 2 = 2\)。法丛标架在 \(D^2\) 上延拓的阻碍类落在 \(\pi_1(\operatorname{SO}(2)) \cong \mathbb{Z}\) 中，一般非零。这对应于 \(D^2\) 自身交点的''缠绕数''，无法通过连续形变消除。因此，在 \(n = 4\) 时，即使存在 Whitney 圆盘，也不能保证可通过合痕消去交点。D. Casson 意识到，这一阻碍可以通过''挖去 \(D^2\) 并替换为 Casson 手柄''来解决------这就是 Casson 手柄理论的起源，也是 Freedman 处理拓扑 4 维流形的核心工具。\(\blacksquare\)

\subsection{218.2 Freedman 定理与拓扑 4 维流形的分类}\label{freedman-ux5b9aux7406ux4e0eux62d3ux6251-4-ux7ef4ux6d41ux5f62ux7684ux5206ux7c7b}

Whitney 技巧在光滑范畴中的失效在拓扑范畴中可以被克服------这正是 Freedman 1982 年获得 Fields 奖的工作的核心。

\textbf{定理 218.3}（Freedman 定理，1982）：设 \(X\) 是单连通闭拓扑 4 维流形。则 \(X\) 的同胚类由其相交形式 \(Q_X: H^2(X; \mathbb{Z}) \times H^2(X; \mathbb{Z}) \to \mathbb{Z}\) 和 Kirby-Siebenmann 不变量 \(\operatorname{ks}(X) \in \mathbb{Z}/2\mathbb{Z}\) 完全决定。反之，任意幺模对称双线性型可以实现为某单连通闭拓扑 4 维流形的相交形式（若形式为偶型，则 \(\operatorname{ks}(X) \equiv \sigma(X)/8 \pmod{2}\)；若为奇型，则 \(\operatorname{ks}(X)\) 可任意选择）。

\textbf{证明思路}：Freedman 的核心创新在于证明了无穷 Casson 手柄（Casson handle）在拓扑范畴中同胚于标准开 2-手柄 \(D^2 \times \mathbb{R}^2\)。Casson 手柄可视为对 Whitney 圆盘嵌入失败的补偿------当无法将自交的 Whitney 圆盘嵌入为光滑子流形时，Casson 构造了层层嵌套的''带孔圆盘''的无穷迭代（其每一层都有带子提供下一层的嵌入位置），形成一个非紧的分形结构。Freedman 证明了：（1）此无穷迭代结构在拓扑范畴中可收缩（通过 Bing 收缩和分解空间理论，将未端收缩为 Cantor 集，再通过''reduction''消去奇异点）；（2）收缩后的对象与 \(D^2 \times \mathbb{R}^2\) 同胚。一旦 Casson 手柄被识别为标准 2-手柄，5 维 h-配边定理（在拓扑 4 维范畴中由 Freedman 证明成立------利用上述 Casson 手柄的收缩技巧处理 Whitney 圆盘）可应用于分类：任何两个单连通闭拓扑 4 维流形若具有相同的相交形式和相同的 Kirby-Siebenmann 不变量，则它们是 h-配边的，从而同胚。\(\blacksquare\)

\subsection{218.3 Donaldson 定理与光滑结构的限制}\label{donaldson-ux5b9aux7406ux4e0eux5149ux6ed1ux7ed3ux6784ux7684ux9650ux5236}

Freedman 定理给出了充足的空间（任何幺模形式都是某拓扑 4 维流形的相交形式），但 Donaldson 定理表明光滑范畴远比拓扑范畴严格。

\textbf{定理 218.4}（Donaldson 定理 A，1983）：设 \(X\) 为闭光滑单连通 4 维流形。若 \(X\) 的相交形式 \(Q_X\) 是正定的，则 \(Q_X\) 在 \(\mathbb{Z}\) 上对角化------即等价于标准 Euclid 形式 \(\langle 1 \rangle \oplus \cdots \oplus \langle 1 \rangle\)。

\textbf{证明思路}：核心工具是研究 \(X\) 上 \(\operatorname{SU}(2)\)-瞬子（反自对偶联络）的模空间 \(\mathcal{M}\)。设 \(Q_X\) 是秩 \(b_2\) 的正定形式。模空间 \(\mathcal{M}\) 的虚维数由 Atiyah-Singer 指标定理给出：\(\dim \mathcal{M} = 8k - 3(1 - b_1 + b_2^+)\)，其中 \(k = c_2(E)\) 为瞬子数，\(b_2^+ = b_2\)（正定性保证所有同调类自交为正）。选取最小的瞬子数 \(k\) 使得 \(\mathcal{M}\) 的虚维数为正。模空间 \(\mathcal{M}\) 有自然的 \(\operatorname{SO}(3)\) 作用（规范群在基点处的剩余对称）。由 Uhlenbeck 紧致性定理，\(\mathcal{M}\) 在添加''理想瞬子''后可紧化。对紧化后的模空间应用 Donaldson 的定向和相交理论：在适当微扰下，\(\mathcal{M}\) 是有奇点的光滑流形，其奇点对应于可约联络，其邻域由 \(Q_X\) 的短向量（自交 \(\leq\) 某阈值的向量）参数化。对模空间边界（对应瞬子''泡''到点上）的细致分析表明：若 \(Q_X\) 非对角型（即含有像 \(E_8\) 这样的不可分解因子），则模空间的拓扑性质（特别是其 2 维同调的相交形式）将产生矛盾。可证明：仅当 \(Q_X\) 对角型时，模空间才是闭流形且其基本类给出自洽的同调关系。因此正定光滑相交形式必须是对角型的。\(\blacksquare\)

\textbf{推论 218.1}（Exotic 光滑结构的存在性）：Freedman 定理允许 \(E_8\) 偶幺模形式作为拓扑 4 维流形的相交形式（\(\operatorname{ks} = 0\) 给出一个拓扑流形 \(|E_8|\)），但由 Donaldson 定理，\(|E_8|\) 不能装备光滑结构------因为 \(E_8\) 是正定的但非对角型。这直接证明了存在只具有拓扑结构而无光滑结构的闭 4 维流形。更进一步，对 \(\mathbb{CP}^2 \# n \overline{\mathbb{CP}^2}\) 做对数变换（logarithmic transform），可构造无穷多个拓扑同胚但互不微分同胚的光滑 4 维流形。

\subsection{\texorpdfstring{218.4 Exotic \(\mathbb{R}^4\) 的构造思路}{218.4 Exotic \textbackslash mathbb\{R\}\^{}4 的构造思路}}\label{exotic-mathbbr4-ux7684ux6784ux9020ux601dux8def}

\textbf{定理 218.5}（Exotic \(\mathbb{R}^4\) 的存在性，Taubes 1987, Gompf 1985）：存在与标准 \(\mathbb{R}^4\) 同胚但不微分同胚的 4 维流形（称为 exotic \(\mathbb{R}^4\)）。事实上，存在不可数无穷多个两两不微分同胚的 exotic \(\mathbb{R}^4\)。

\textbf{构造思路}：核心想法是将 Donaldson 定理的''全局后果''注入到 \(\mathbb{R}^4\) 这类非紧流形中。具体的 Gompf-Taubes 构造如下：

\begin{enumerate}
\def\labelenumi{\arabic{enumi}.}
\item
  取 \(K3\) 曲面（紧致单连通光滑 4 维流形，相交形式为 \(E_8 \oplus E_8 \oplus 3H\)，其中 \(H\) 为双曲平面）。Freedman 定理表明 \(K3\) 拓扑同胚于 \(3(\mathbb{CP}^2 \# \overline{\mathbb{CP}^2}) \# (|E_8| \# |E_8|)\) 的连通和（省略细节），但 Donaldson 定理表明两者不微分同胚（因为它们的 Donaldson 不变量不同）。
\item
  从 \(K3\) 中移去一个点得到 \(K3 \setminus \{\text{pt}\}\)。此非紧流形可''嵌入''到某 \(\mathbb{R}^4\) 中的有界区域（通过适当的坐标选取和 Nash 嵌入定理的仿射版本）。
\item
  拓扑上，\(K3 \setminus \{\text{pt}\}\) 同胚于 \(\mathbb{R}^4\)（由 Freedman 定理和 h-配边定理推出任何单连通非紧 4 维流形若其端同胚于 \(S^3 \times \mathbb{R}\) 则同胚于 \(\mathbb{R}^4\)）。但若 \(K3 \setminus \{\text{pt}\}\) 微分同胚于标准 \(\mathbb{R}^4\)，则 \(K3\) 应能光滑地紧化为某光滑 4 维流形 \(X\) 使得 \(X \setminus \{\text{pt}\}\) 微分同胚于 \(K3 \setminus \{\text{pt}\}\)，由此 \(X\) 与 \(K3\) 有相同的 Donaldson 不变量，但 \(X\) 的紧化方式（加回的点处的邻域结构）将给出与 \(K3\) 的 Donaldson 定理相矛盾的光滑结构。因此 \(K3 \setminus \{\text{pt}\}\) 不能微分同胚于标准 \(\mathbb{R}^4\)------它是 exotic \(\mathbb{R}^4\)。
\item
  通过在不同位置''移去''不同大小的开集，利用 Donaldson 不变量的精细变化，可产生不可数多个不同的 exotic \(\mathbb{R}^4\)（Gompf 1985，Taubes 1987）。
\end{enumerate}

这一构造揭示了 4 维光滑拓扑中一个反直觉的事实：非紧流形（\(\mathbb{R}^4\)）的光滑结构可以比紧流形更加丰富，因为''无穷远处''的复杂拓扑（如 exotic \(\mathbb{R}^4\) 的 Cantor 集端）可以在紧化过程中产生奇异行为------而这是 4 维独有的现象，在高维中 Cerf 定理保证了这种''端效应''不会发生。\(\blacksquare\)

\subsection{218.5 低维拓扑的统一视角}\label{ux4f4eux7ef4ux62d3ux6251ux7684ux7edfux4e00ux89c6ux89d2}

上述定理揭示了一个深刻的统一图景：3 维和 4 维流形之所以特殊，根源在于 Whitney 技巧的失效------这使低维流形的分类问题从同伦论问题（如高维 Browder-Novikov-Sullivan-Wall 手术正合序列所刻画）变为本质上更复杂的几何分析问题。在 3 维，Ricci 流（Hamilton-Perelman）成为克服 Whitney 失效的替代工具；在 4 维，瞬子方程（Donaldson）和 Seiberg-Witten 方程提供了光滑结构的不变量。低维拓扑的终极梦想------类似于高维手术理论的统一分类框架------至今仍未完全实现，这是 21 世纪数学的核心挑战之一。

\hrulefill

\hrulefill

\hrulefill

\hrulefill

\hrulefill

\hrulefill

\hrulefill

\hrulefill

\newpage

\chapter{07-拓扑与几何/03-代数拓扑/07-纤维丛与示性类.md}\label{ux62d3ux6251ux4e0eux51e0ux4f5503-ux4ee3ux6570ux62d3ux625107-ux7ea4ux7ef4ux4e1bux4e0eux793aux6027ux7c7b.md}

\section{第256章：纤维丛与示性类}\label{ux7b2c256ux7ae0ux7ea4ux7ef4ux4e1bux4e0eux793aux6027ux7c7b}

纤维丛是代数拓扑与微分几何交汇处的核心概念。它形式化了''局部平凡、全局扭曲''的几何结构------在局部上看，纤维丛是乘积空间 \(U \times F\)，但在全局上可能通过''扭转''的方式粘合而成。示性类（characteristic classes）通过上同调类精确刻画这种扭曲，是连接拓扑、几何与数学物理的桥梁。

\subsection{81.1 纤维丛的基本概念}\label{ux7ea4ux7ef4ux4e1bux7684ux57faux672cux6982ux5ff5}

\textbf{定义 81.1}（纤维丛）：一个\textbf{纤维丛} \((E, \pi, B, F, G)\) 由以下数据构成：

\begin{itemize}
\tightlist
\item
  \textbf{全空间}（total space）\(E\)
\item
  \textbf{底空间}（base space）\(B\)
\item
  \textbf{投影}（projection）\(\pi: E \to B\)，是连续满射
\item
  \textbf{标准纤维}（typical fiber）\(F\)
\item
  \textbf{结构群}（structure group）\(G \subseteq \text{Diff}(F)\)（或 \(G \subseteq \text{Homeo}(F)\)）
\item
  一族\textbf{局部平凡化} \(\{ \phi_i: \pi^{-1}(U_i) \to U_i \times F \}\)，其中 \(\{U_i\}\) 是 \(B\) 的开覆盖
\end{itemize}

满足：对每对 \((i, j)\)，\textbf{转移函数}（transition function）\(g_{ij}: U_i \cap U_j \to G\) 定义为

\[\phi_i \circ \phi_j^{-1}(x, f) = (x, g_{ij}(x) \cdot f)\]

且满足\textbf{上循环条件}（cocycle condition）：\(g_{ij}(x) \circ g_{jk}(x) = g_{ik}(x)\)，对所有 \(x \in U_i \cap U_j \cap U_k\)。

\textbf{定义 81.2}（截面）：纤维丛 \(\pi: E \to B\) 的一个\textbf{截面}（section）是连续映射 \(s: B \to E\) 满足 \(\pi \circ s = \text{id}_B\)。向量丛的截面空间记作 \(\Gamma(E)\)，它是一个 \(C^\infty(B)\)-模。切丛的截面即\textbf{向量场}，余切丛的截面即 \textbf{1-形式}。

\textbf{定义 81.3}（拉回丛）：设 \(f: X \to B\) 是连续映射，\(E \to B\) 是纤维丛。\textbf{拉回丛}（pullback bundle）\(f^*E \to X\) 定义为

\[f^*E = \{(x, e) \in X \times E : f(x) = \pi(e)\}\]

投影为 \((x, e) \mapsto x\)。拉回丛继承了原丛的纤维和结构群。

\textbf{定义 81.4}（丛同构）：两个纤维丛 \(\pi_1: E_1 \to B\) 和 \(\pi_2: E_2 \to B\)（同一底空间）称为\textbf{同构}，如果存在同胚 \(\Phi: E_1 \to E_2\) 满足 \(\pi_2 \circ \Phi = \pi_1\)，且 \(\Phi\) 在每根纤维上是线性同构（向量丛情形）或 \(G\)-等变同胚（主丛情形）。

\textbf{定理 81.1}（上同调分类定理）：对李群 \(G\)，存在\textbf{分类空间} \(BG\) 和\textbf{万有丛} \(EG \to BG\)（\(EG\) 可缩），使得对任意仿紧空间 \(X\)，\(X\) 上主 \(G\)-丛的同构类与连续映射同伦类集 \([X, BG]\) 一一对应。丛 \(P \to X\) 对应的映射 \(f_P: X \to BG\) 称为\textbf{分类映射}（classifying map），且 \(P \cong f_P^* EG\)。

\textbf{证明}：Milnor 构造给出了万有丛的显式模型：取 \(G\) 的无穷 join \(EG = G * G * \cdots\)，赋予弱拓扑。投影 \(\pi: EG \to BG\) 为自然商映射，\(BG = EG/G\)。\(EG\) 可缩是因为有限 join \(G^{*n}\) 的连通性随 \(n\) 增大而发散。对任意主 \(G\)-丛 \(P \to X\)，取 \(X\) 的局部平凡化开覆盖 \(\{U_i\}\) 和转移函数 \(g_{ij}\)。利用单位分解构造从 \(P\) 到 \(EG\) 的丛映射 \(\tilde{f}_P: P \to EG\)，其在商空间上诱导 \(f_P: X \to BG\)。反之，给定 \(f: X \to BG\)，拉回丛 \(f^*EG\) 是 \(X\) 上的主 \(G\)-丛。同伦映射给出同构丛，完成了双射对应。 \(\blacksquare\)

重要例子： - \(BGL(n, \mathbb{R}) \simeq BO(n) \simeq \text{Gr}_n(\mathbb{R}^\infty)\)（无穷 Grassmann 流形） - \(BGL(n, \mathbb{C}) \simeq BU(n) \simeq \text{Gr}_n(\mathbb{C}^\infty)\) - \(BU(1) \simeq \mathbb{CP}^\infty = K(\mathbb{Z}, 2)\) - \(B\mathbb{Z}/2 \simeq \mathbb{RP}^\infty = K(\mathbb{Z}/2, 1)\)

\textbf{推论 81.2}（示性类的定义）：设 \(c \in H^*(BG; R)\) 为上同调类。对任意主 \(G\)-丛 \(P \to X\) 及分类映射 \(f_P: X \to BG\)，定义丛 \(P\) 的\textbf{示性类}为 \(c(P) = f_P^*(c) \in H^*(X; R)\)。这解释了为什么示性类是''万有''的------它们来自分类空间的上同调。

\textbf{证明}：由定理 81.1，分类映射 \(f_P: X \to BG\) 的存在性与同伦唯一性保证 \(P \cong f_P^* EG\)。定义 \(c(P) = f_P^*(c) \in H^*(X; R)\)。需验证定义不依赖于分类映射的选取：若 \(f_P, g_P: X \to BG\) 均为 \(P\) 的分类映射，由定理 81.1 它们同伦，故 \(f_P^* = g_P^*\) 在上同调层面，从而 \(c(P)\) 良定。自然性：对任意连续映射 \(h: Y \to X\)，\(h^*P\) 的分类映射为 \(f_P \circ h\)，故 \(c(h^*P) = (f_P \circ h)^*(c) = h^* f_P^*(c) = h^* c(P)\)。因此示性类满足拉回自然性，确为丛的不变量。 \(\blacksquare\)

\subsection{81.2 Stiefel-Whitney 类与 Chern 类}\label{stiefel-whitney-ux7c7bux4e0e-chern-ux7c7b}

\textbf{定义 81.5}（Stiefel-Whitney 类）：对实向量丛 \(E \to B\)（秩 \(n\)），其 \textbf{Stiefel-Whitney 类} \(w_i(E) \in H^i(B; \mathbb{Z}/2\mathbb{Z})\)（\(0 \leq i \leq n\)）是满足以下公理的唯一族上同调类：

\begin{enumerate}
\def\labelenumi{\arabic{enumi}.}
\tightlist
\item
  \(w_0(E) = 1\)，\(w_i(E) = 0\) 当 \(i > \text{rank}(E)\)
\item
  \textbf{自然性}（naturality）：对任意连续映射 \(f: X \to B\)，\(w_i(f^*E) = f^* w_i(E)\)
\item
  \textbf{Whitney 和公式}：对总 Stiefel-Whitney 类 \(w(E) = 1 + w_1(E) + w_2(E) + \cdots + w_n(E)\)，有
\end{enumerate}

\[w(E \oplus F) = w(E) \smile w(F)\]

\begin{enumerate}
\def\labelenumi{\arabic{enumi}.}
\setcounter{enumi}{3}
\tightlist
\item
  \textbf{正规化}（normalization）：\(\mathbb{RP}^1 \cong S^1\) 上重言线丛 \(\gamma^1\) 满足 \(w_1(\gamma^1) \neq 0 \in H^1(\mathbb{RP}^1; \mathbb{Z}/2) \cong \mathbb{Z}/2\)
\end{enumerate}

\textbf{定理 81.3}（Stiefel-Whitney 类的存在性与唯一性）：满足上述公理的 Stiefel-Whitney 类存在且唯一。存在性的标准构造使用分类空间 \(BO(n)\) 的上同调：\(H^*(BO(n); \mathbb{Z}/2) \cong \mathbb{Z}/2[w_1, \ldots, w_n]\)（多项式代数），其中 \(w_i\) 即万有 Stiefel-Whitney 类。唯一性由分裂原理保证。

\textbf{证明}：存在性：令 \(w_i \in H^i(BO(n); \mathbb{Z}/2)\) 为生成元，对任意秩 \(n\) 实向量丛 \(E \to X\)，由分类定理存在分类映射 \(f_E: X \to BO(n)\)，定义 \(w_i(E) = f_E^*(w_i)\)。公理验证：自然性因 \((g \circ f_E)^* = f_E^* \circ g^*\)；Whitney 和公式由 \(BO(m) \times BO(n) \to BO(m+n)\) 诱导的上同调计算得 \(w(E \oplus F) = w(E) \smile w(F)\)；正规化因 \(\mathbb{RP}^\infty = BO(1)\) 且 \(H^1(\mathbb{RP}^\infty; \mathbb{Z}/2) \cong \mathbb{Z}/2\) 由 \(w_1\) 生成。唯一性：设 \(\{w_i'\}\) 为另一族满足公理的类。由分裂原理，\(f^*: H^*(X) \hookrightarrow H^*(F(E))\) 为单射，而 \(f^*E\) 分裂为线丛直和，线丛的 \(w_1\) 由正规化公理唯一确定，故 \(w_i = w_i'\)。 \(\blacksquare\)

\textbf{定理 81.4}（Stiefel-Whitney 类的几何意义）： - \(w_1(E) = 0 \iff E\) 可定向（\(w_1\) 是定向障碍类） - \(w_2(E) = 0 \iff E\) 有自旋结构（\(w_2\) 是自旋障碍类） - 紧流形 \(M\) 可平行化（tangential parallelizable）当且仅当 \(w_i(TM) = 0\) 对所有 \(i > 0\)

特别地，\(w_1(T\mathbb{RP}^n) \neq 0\)（当 \(n\) 奇），\(\mathbb{RP}^2\) 不可平行化（\(w_2(T\mathbb{RP}^2) \neq 0\)）。

\textbf{证明}：(i) \(w_1(E)\) 为定向障碍：线丛 \(\det(E) = \bigwedge^n E\) 的定向平凡化等价于 \(w_1(\det(E)) = 0\)，而 \(w_1(E) = w_1(\det(E))\)，故 \(w_1(E) = 0 \iff \det(E)\) 可平凡化 \(\iff E\) 可定向。(ii) \(w_2\) 为自旋障碍：自旋结构的存在性等价于 \(w_2(E) = 0\)，对应 \(\operatorname{Spin}(n) \to SO(n)\) 的障碍理论，第二个 Stiefel-Whitney 类是此二重覆叠的分类不变量。(iii) 可平行化条件：若 \(TM\) 可平凡化，则 \(w_i(TM) = 0\) 对所有 \(i > 0\)；反之若所有 \(w_i(TM) = 0\)，由 Stiefel 的平行化定理，\(M\) 浸入 \(\mathbb{R}^{2n}\) 的示性类消失给出截面的无碍构造。 \(\blacksquare\)

\textbf{定理 81.5}（Wu 公式与 Wu 类）：对闭流形 \(M^n\)，存在唯一的上同调类 \(v_i \in H^i(M; \mathbb{Z}/2)\)（Wu 类），使得对所有 \(x \in H^{n-i}(M; \mathbb{Z}/2)\)，\(\langle v_i \smile x, [M] \rangle = \langle Sq^i(x), [M] \rangle\)，其中 \(Sq^i\) 是 Steenrod 平方运算。Stiefel-Whitney 类与 Wu 类的关系为 \(w = Sq(v)\)，即 \(w_k = \sum_{i=0}^k Sq^{k-i}(v_i)\)。Wu 公式表明 Stiefel-Whitney 类完全由上同调环的 Steenrod 代数结构决定------它们是流形的同伦型不变量。

\textbf{证明}：由 Poincare 对偶，线性泛函 \(x \mapsto \langle Sq^i(x), [M] \rangle\) 由唯一的 \(v_i \in H^i(M; \mathbb{Z}/2)\) 通过杯积表示，定义了 Wu 类。引入 Thom 同构 \(\Phi: H^*(M) \to \tilde{H}^{*+n}(Th(TM))\) 及 Thom 类 \(U = \Phi(1)\)，则 Stiefel-Whitney 类满足 \(w = \Phi^{-1}Sq\Phi(1)\)。由 Cartan 公式 \(Sq\Phi(v_k) = \Phi(Sq(v_k))\) 及 \(Sq^i(U) = \Phi(w_i)\)，计算 \(\Phi(Sq(v)) = Sq(\Phi(v)) = \sum_i Sq(\Phi(v_i)) = \sum_{i,k} Sq^{k-i} \Phi(v_i) = \sum_k \Phi(\sum_i Sq^{k-i} v_i)\)。另一方面 \(\Phi(w) = \sum_k \Phi(w_k)\)。取 \(\Phi^{-1}\) 即得 \(w_k = \sum_{i=0}^k Sq^{k-i}(v_i)\)。 \(\blacksquare\)

\textbf{定义 81.6}（Chern 类）：对复向量丛 \(E \to B\)（秩 \(n\)），其 \textbf{Chern 类} \(c_i(E) \in H^{2i}(B; \mathbb{Z})\) 满足类似公理：

\begin{enumerate}
\def\labelenumi{\arabic{enumi}.}
\tightlist
\item
  \(c_0(E) = 1\)，\(c_i(E) = 0\) 当 \(i > \text{rank}(E)\)
\item
  \textbf{自然性}：\(c_i(f^*E) = f^* c_i(E)\)
\item
  \textbf{Whitney 和公式}：\(c(E \oplus F) = c(E) \smile c(F)\)
\item
  \textbf{正规化}：\(\mathbb{CP}^1 \cong S^2\) 上重言线丛 \(\gamma^1_{\mathbb{C}}\) 满足 \(c_1(\gamma^1_{\mathbb{C}}) = -1 \cdot [\omega_{FS}]\)，其中 \([\omega_{FS}]\) 是 Fubini-Study 形式的基本类（即 \(\langle c_1(\gamma^1_{\mathbb{C}}), [\mathbb{CP}^1] \rangle = -1\)）
\end{enumerate}

总 Chern 类记作 \(c(E) = 1 + c_1(E) + c_2(E) + \cdots + c_n(E)\)。Chern 类从分类空间的上同调得到：\(H^*(BU(n); \mathbb{Z}) \cong \mathbb{Z}[c_1, \ldots, c_n]\)，\(\deg c_i = 2i\)。

\textbf{定理 81.6}（分裂原理，Splitting Principle）：对任意向量丛 \(E \to B\)，存在连续映射 \(f: F(E) \to B\)（\textbf{旗流形}，flag manifold），使得：

\begin{enumerate}
\def\labelenumi{\arabic{enumi}.}
\tightlist
\item
  拉回丛 \(f^*E\) 分裂为线丛的直和：\(f^*E \cong L_1 \oplus L_2 \oplus \cdots \oplus L_n\)
\item
  诱导同态 \(f^*: H^*(B) \to H^*(F(E))\) 是\textbf{单射}
\end{enumerate}

分裂原理将复杂向量丛的示性类计算约化到线丛情形。例如，Chern 类可形式地写为 \(c(E) = \prod_{i=1}^n (1 + x_i)\)，其中 \(x_i = c_1(L_i)\) 称为 \textbf{Chern 根}（Chern roots）。

\textbf{证明}：构造旗丛 \(F(E) \to B\)，纤维为 \(E_b\) 中的完全旗。取投影化 \(p_1: \mathbb{P}(E) \to B\)，重言线丛嵌入给出分裂 \(p_1^*E = L_1 \oplus E_1\)（\(\operatorname{rank} E_1 = n-1\)）。归纳地投影化 \(E_1\)，经 \(n-1\) 步得 \(f: F(E) \to B\) 使 \(f^*E \cong \bigoplus_{i=1}^n L_i\)。单射性由 Leray-Hirsch 定理保证：\(\mathbb{P}(E)\) 的上同调 \(H^*(\mathbb{P}(E))\) 是自由 \(H^*(B)\)-模，由 \(1, c_1(L_1), \ldots, c_1(L_1)^{n-1}\) 生成；逐次投影化，每步 \(f^*\) 保持单射性。 \(\blacksquare\)

\textbf{定理 81.7}（Chern-Weil 理论）：设 \(\nabla\) 是复向量丛 \(E \to M\)（\(M\) 是光滑流形）上的联络，其曲率 \(F_\nabla \in \Omega^2(M; \text{End}(E))\)。则 Chern 类可由曲率形式的多项式表示为 de Rham 上同调类：

\[c_k(E) = \left[ \sigma_k\left(\frac{i}{2\pi} F_\nabla\right) \right] \in H^{2k}_{\text{dR}}(M)\]

其中 \(\sigma_k\) 是第 \(k\) 个初等对称多项式（即特征值的第 \(k\) 个对称和）。特别地，

\[c_1(E) = \left[\frac{i}{2\pi} \text{Tr}(F_\nabla)\right], \quad c_2(E) = \left[\frac{1}{8\pi^2}\left(\text{Tr}(F_\nabla \wedge F_\nabla) - \text{Tr}(F_\nabla) \wedge \text{Tr}(F_\nabla)\right)\right]\]

Chern-Weil 理论的美妙之处在于：尽管曲率 \(F_\nabla\) 依赖于联络的选择，由曲率多项式定义的上同调类却不依赖于联络------它是拓扑不变量。

\textbf{证明}：令 \(\nabla_0, \nabla_1\) 为 \(E\) 上两个联络，定义插值族 \(\nabla_t = t\nabla_1 + (1-t)\nabla_0\)。由 Chern-Simons 超度公式，\(\frac{d}{dt} \sigma_k(\frac{i}{2\pi}F_{\nabla_t}) = d \eta_t\)（\(\eta_t\) 为某微分形式）。沿 \([0,1]\) 积分得 \(\sigma_k(\frac{i}{2\pi}F_{\nabla_1}) - \sigma_k(\frac{i}{2\pi}F_{\nabla_0}) = d\int_0^1 \eta_t dt\)，故两端定义相同的 de Rham 上同调类。验证 \(\sigma_k(\frac{i}{2\pi}F_\nabla)\) 满足 Chern 类公理：自然性因 \(F_{f^*\nabla} = f^*F_\nabla\)；Whitney 和公式因 \(\sigma_k(F_{E \oplus F}) = \sum_{i+j=k} \sigma_i(F_E) \wedge \sigma_j(F_F)\)；正规化由 \(\mathbb{CP}^1\) 上直接计算验证。 \(\blacksquare\)

\subsection{81.3 Pontryagin 类与 Euler 类}\label{pontryagin-ux7c7bux4e0e-euler-ux7c7b}

\textbf{定义 81.7}（Pontryagin 类）：实向量丛 \(E \to B\)（秩 \(n\)）的 \textbf{Pontryagin 类} \(p_k(E) \in H^{4k}(B; \mathbb{Z})\)（或系数在任意交换环中）由复化定义：

\[p_k(E) = (-1)^k c_{2k}(E \otimes_{\mathbb{R}} \mathbb{C})\]

总 Pontryagin 类 \(p(E) = 1 + p_1(E) + p_2(E) + \cdots\)，其中 \(p_k(E)\) 的次数为 \(4k\)。注意 \(c_{2k+1}(E \otimes \mathbb{C})\) 是 2-挠元，故定义中只取偶数次 Chern 类。

由分裂原理，Pontryagin 类可形式地写为 \(p(E) = \prod (1 + x_i^2)\)，其中 \(x_i\) 是复化的 Chern 根（本质上是实丛的「Euler 形式」的对称化）。

\textbf{定义 81.8}（Euler 类）：设 \(E \to B\) 是秩 \(n\) 的\textbf{定向}实向量丛。其 \textbf{Euler 类} \(e(E) \in H^n(B; \mathbb{Z})\) 满足：

\begin{enumerate}
\def\labelenumi{\arabic{enumi}.}
\tightlist
\item
  自然性：\(e(f^*E) = f^* e(E)\)
\item
  若 \(E\) 是奇秩丛，则 \(2e(E) = 0\)
\item
  \(e(E \oplus F) = e(E) \smile e(F)\)（在定向相容条件下的 Whitney 和公式的乘法形式）
\item
  \textbf{正规化}：对 \(n\) 维闭定向流形 \(M\)，\(\langle e(TM), [M] \rangle = \chi(M)\)（Euler 示性数）
\end{enumerate}

\textbf{定理 81.8}（Euler 类与 Stiefel-Whitney 类的关系）：在模 2 约化下，\(e(E) \equiv w_n(E) \pmod{2}\)。换言之，Euler 类是顶维 Stiefel-Whitney 类的整系数提升。

\textbf{证明}：在 Thom 空间中，Euler 类 \(e(E)\) 由 \(s_0^* U\) 定义（\(U\) 为 Thom 类，\(s_0\) 为零截面），而 \(w_n(E) = s_0^*(U \bmod 2)\)。因 Thom 类的整系数提升模 2 恰为 \(\mathbb{Z}/2\) 系数的 Thom 类，故在模 2 约化下 \(e(E) \equiv w_n(E)\)。等价地，在示性类的公理构造中，\(e(E)\) 和 \(w_n(E)\) 满足相同的公理化条件（模 2 意义下），由分裂原理的唯一性得证。 \(\blacksquare\)

\textbf{定理 81.9}（Gauss-Bonnet-Chern 定理）：设 \(M\) 是 \(2n\) 维闭定向 Riemann 流形，\(\Omega\) 是其曲率 2-形式（Levi-Civita 联络的曲率）。则

\[\chi(M) = \frac{1}{(2\pi)^n} \int_M \text{Pf}(\Omega)\]

其中 \(\text{Pf}(\Omega)\) 是曲率形式的 Pfaffian，满足 \(\text{Pf}(\Omega)^2 = \det(\Omega)\)。这一定理将 Euler 示性数（拓扑量）表示为曲率的积分（几何量），是 Gauss-Bonnet 定理的高维推广。

\textbf{证明}：由 Chern-Weil 理论，Euler 类 \(e(TM)\) 由 \(\frac{1}{(2\pi)^n}\text{Pf}(\Omega)\) 的 de Rham 上同调类表示。同时，Euler 示性数 \(\chi(M) = \langle e(TM), [M] \rangle\) 是 Euler 类在基本类上的求值。在 de Rham 上同调中，上同调类与基本类的配对恰为积分：\(\langle [\omega], [M] \rangle = \int_M \omega\)。故 \(\chi(M) = \langle e(TM), [M] \rangle = \int_M \frac{1}{(2\pi)^n}\text{Pf}(\Omega)\)。其中 \(\text{Pf}(\Omega)\) 是 \(\Omega\) 作为 \(\mathfrak{so}(2n)\)-值 2-形式在 Clifford 代数内的平方根，其积分的不变性由 Chern-Weil 理论保证。 \(\blacksquare\)

\textbf{定理 81.10}（Hirzebruch 符号差定理）：设 \(M\) 是 \(4k\) 维定向紧流形。其\textbf{符号差}（signature）\(\sigma(M)\)（即相交形式在中间维上同调 \(H^{2k}(M; \mathbb{R})\) 上的符号差）由 \(L\)-亏格（\(L\)-genus）给出：

\[\sigma(M) = \langle L_k(p_1, \ldots, p_k), [M] \rangle\]

其中 \(L\)-多项式由形式幂级数 \(\prod_{i=1}^k \frac{x_i}{\tanh x_i}\) 的 Taylor 展开定义（\(x_i\) 是 Pontryagin 根）。前几个 \(L\)-多项式为：

\[L_1 = \frac{1}{3}p_1, \quad L_2 = \frac{1}{45}(7p_2 - p_1^2), \quad L_3 = \frac{1}{945}(62p_3 - 13p_2p_1 + 2p_1^3)\]

\textbf{证明}：Atiyah-Singer 指标定理应用于符号差算子 \(D = d + d^*: \Omega^+ \to \Omega^-\)（将 \(2k\)-形式按 Hodge 星算子分解）。\(\text{ind}(D) = \dim \ker D|_{+} - \dim \ker D|_{-} = \sigma(M)\)。指标定理给出 \(\sigma(M) = \langle \text{ch}(\sigma(D)) \smile \hat{A}(TM), [TM] \rangle\)。计算符号差算子的象征：\(\text{ch}(\sigma(D)) \smile \hat{A}(TM)\) 在 \(4k\) 维等于 \(L_k(p_1, \ldots, p_k)\)，其中 \(L\)-多项式由 \(\prod_{i=1}^k \frac{x_i}{\tanh x_i}\) 展开并取 \(4k\) 次部分得出。此因 \(\hat{A}\)-亏格与 \(L\)-亏格满足关系 \(L(M) = \hat{A}(M) \cdot \text{ch}(\bigwedge^+ - \bigwedge^-)\)。 \(\blacksquare\)

\subsection{81.4 示性类与指标定理}\label{ux793aux6027ux7c7bux4e0eux6307ux6807ux5b9aux7406}

示性类深层的意义在 Atiyah-Singer 指标定理中得到了最完整的体现------它统一了多个经典定理，揭示了分析学（椭圆算子的指标）、拓扑学（示性类）和几何学（曲率）之间的内在联系。

\textbf{定理 81.11}（Atiyah-Singer 指标定理）：设 \(D: \Gamma(E) \to \Gamma(F)\) 是紧无边流形 \(M\) 上的椭圆微分算子（例如 Dirac 算子、Dolbeault 算子、符号差算子），\(\sigma(D)\) 是其主象征（principal symbol），可视作 \(K\) 理论中的元素。则

\[\text{ind}(D) = (-1)^{n(n+1)/2} \langle \text{ch}(\sigma(D)) \smile \hat{A}(TM), [TM] \rangle\]

其中： - \(\text{ch}: K(M) \to H^*(M; \mathbb{Q})\) 是 \textbf{Chern 特征}（Chern character），定义为 \(\text{ch}(E) = \sum_{i=1}^n e^{x_i}\)，\(x_i\) 为 Chern 根。前几项为

\[\text{ch}(E) = \text{rank}(E) + c_1 + \frac{1}{2}(c_1^2 - 2c_2) + \frac{1}{6}(c_1^3 - 3c_1c_2 + 3c_3) + \cdots\]

\begin{itemize}
\tightlist
\item
  \(\hat{A}(TM)\) 是 \textbf{\(\hat{A}\)-亏格}（\(\hat{A}\)-genus，或 A-hat genus），由形式幂级数 \(\prod \frac{x_i/2}{\sinh(x_i/2)}\) 定义（\(x_i\) 是 \(TM \otimes \mathbb{C}\) 的 Pontryagin 根）
\end{itemize}

\textbf{证明}：核心思路是将指标转化为拓扑不变量。构造 \(M\) 的余切丛 \(T^*M\) 的 Thom 空间，利用嵌入 \(M \hookrightarrow \mathbb{R}^N\) 将椭圆算子 \(D\) 的提升转化为 \(\mathbb{R}^N\) 上的 Fredholm 算子族。通过 Bott 周期性，\(\text{ind}(D)\) 等同于 \(K\) 理论中元素 \([\sigma(D)] \in K(T^*M)\) 的拓扑指标。Atiyah-Singer 证明拓扑指标由 \(\langle \text{ch}(\sigma(D)) \smile \text{Td}(TM \otimes \mathbb{C}), [TM] \rangle\) 给出。使用 \(\text{Td}(TM \otimes \mathbb{C}) = (-1)^{n(n+1)/2} \hat{A}(TM)^2\) 及 \(\sigma(D)\) 的特定对称性, 即得所述公式。 \(\blacksquare\)

\textbf{推论 81.12}（指标定理的重要特例）：

\begin{enumerate}
\def\labelenumi{\arabic{enumi}.}
\item
  \textbf{Gauss-Bonnet-Chern 定理}：取 de Rham 算子 \(D = d + d^*\)（作用于 \(\bigwedge^* T^*M\)），则 \(\text{ind}(D) = \chi(M)\)，指标定理约化为 Euler 示性数的曲率积分表示。
\item
  \textbf{Hirzebruch-Riemann-Roch 定理}：取 Dolbeault 算子 \(\bar{\partial}\)（作用于全纯向量丛），则
\end{enumerate}

\[\chi(M, E) = \langle \text{ch}(E) \smile \text{Td}(TM), [M] \rangle\]

其中 \(\text{Td}\) 是 \textbf{Todd 亏格}（Todd genus），由 \(\prod \frac{x_i}{1 - e^{-x_i}}\) 定义。这一定理将全纯 Euler 示性数与 Chern 类联系起来，是复几何的基石。

\begin{enumerate}
\def\labelenumi{\arabic{enumi}.}
\setcounter{enumi}{2}
\tightlist
\item
  \textbf{Hirzebruch 符号差定理}：取符号差算子，\(\hat{A}\)-亏格替换为 \(L\)-亏格。
\end{enumerate}

\textbf{证明}：三个特例均由定理 81.11 代入具体算子得出。(1) 取 de Rham 算子 \(D = d + d^*\) 作用于 \(\bigwedge^* T^*M\)，则 \(\text{ind}(D) = \sum_i (-1)^i \dim H^i_{\text{dR}}(M) = \chi(M)\)。此算子的主象征 \(\sigma(D)\) 在 \(K\) 理论中对应外代数丛的差，\(\text{ch}(\sigma(D)) \smile \hat{A}(TM)\) 在顶维积分恰为 Euler 类的曲率积分表示。(2) 取 Dolbeault 算子 \(\bar{\partial}\) 作用于全纯向量丛 \(E\)，则 \(\text{ind}(\bar{\partial}) = \chi(M, E) = \sum_i (-1)^i \dim H^i(M, \mathcal{O}(E))\)。此时指标公式中 \(\hat{A}\)-亏格替换为 Todd 亏格，得 Hirzebruch-Riemann-Roch 公式。(3) 取符号差算子 \(D = d + d^*\) 限制在 \(\Omega^+ \to \Omega^-\)，则 \(\text{ind}(D) = \sigma(M)\)，且 \(\text{ch}(\sigma(D)) \smile \hat{A}(TM)\) 在 \(4k\) 维约化为 \(L\)-亏格。 \(\blacksquare\)

\textbf{定理 81.13}（Chern 特征的函子性）：Chern 特征 \(\text{ch}: K(X) \to H^*(X; \mathbb{Q})\) 是环同态：\(\text{ch}(E \oplus F) = \text{ch}(E) + \text{ch}(F)\)，\(\text{ch}(E \otimes F) = \text{ch}(E) \smile \text{ch}(F)\)。在张量积 \(H^*(X; \mathbb{Q})\) 后，\(\text{ch}\) 诱导同构 \(K(X) \otimes \mathbb{Q} \cong H^*(X; \mathbb{Q})\)（对有限 CW 复形）。

\textbf{证明}：由分裂原理，将 \(E, F\) 分别拉回到分裂为线丛的直和 \(E = \bigoplus L_i\), \(F = \bigoplus M_j\)。则 \(\text{ch}(E \oplus F) = \sum e^{x_i} + \sum e^{y_j} = \text{ch}(E) + \text{ch}(F)\)。\(\text{ch}(E \otimes F) = \sum_{i,j} e^{x_i + y_j} = (\sum_i e^{x_i})(\sum_j e^{y_j}) = \text{ch}(E) \smile \text{ch}(F)\)。同构性：对有限 CW 复形，Atiyah-Hirzebruch 谱序列在有理化下退化，给出 \(K^*(X) \otimes \mathbb{Q} \cong \bigoplus_n H^{2n+*}(X; \mathbb{Q})\)，而 \(\text{ch}\) 提供了该同构的显式实现。 \(\blacksquare\)

\textbf{定理 81.14}（Grothendieck-Riemann-Roch 定理）：设 \(f: X \to Y\) 是光滑射影态射（代数几何中），\(E\) 是 \(X\) 上的向量丛。则

\[\text{ch}(f_! E) \smile \text{Td}(TY) = f_*(\text{ch}(E) \smile \text{Td}(TX))\]

其中 \(f_!\) 是 \(K\) 理论的直接像。这一定理是 Hirzebruch-Riemann-Roch 的相对化版本，在模空间理论和枚举几何中有深远应用。

\textbf{证明}：使用闭嵌入和投影的分解：首先对闭嵌入 \(i: X \hookrightarrow Y\)，\(i_!E\) 的 Chern 特征由 \(i_*(\text{ch}(E) \smile \text{Td}(N_{X/Y})^{-1})\) 给出（\(N_{X/Y}\) 为法丛），而 \(\text{Td}(TX) = \text{Td}(TY|_X) \smile \text{Td}(N_{X/Y})\) 给出所需关系。对投影 \(p: \mathbb{P}(\mathcal{E}) \to Y\)，直接像 \(p_!\) 由上同调的直接像 \(p_*\) 在 \(K\) 理论中的提升给出，验证公式归结为射影丛上 Todd 类的标准计算。一般态射可分解为闭嵌入与光滑投影的复合，两者的 Grothendieck-Riemann-Roch 公式相容，故得证。 \(\blacksquare\)

\subsection{81.5 小结：示性类的统一视角}\label{ux5c0fux7ed3ux793aux6027ux7c7bux7684ux7edfux4e00ux89c6ux89d2}

示性类理论可总结为以下层次结构：

\begin{longtable}[]{@{}llll@{}}
\toprule
示性类 & 系数 & 次数 & 几何意义 \\
\midrule
\endhead
\bottomrule
\endlastfoot
\(w_i\)（Stiefel-Whitney） & \(\mathbb{Z}/2\) & \(i\) & 定向障碍、自旋障碍 \\
\(c_i\)（Chern） & \(\mathbb{Z}\) & \(2i\) & 复向量丛的拓扑扭曲度量 \\
\(p_i\)（Pontryagin） & \(\mathbb{Z}\) & \(4i\) & 实向量丛复化后的扭曲 \\
\(e\)（Euler） & \(\mathbb{Z}\) & \(n\)（秩） & 截面消失的障碍类 \\
\end{longtable}

这些示性类的关系由以下交换图表概括：

\begin{itemize}
\tightlist
\item
  实向量丛 \(\xrightarrow{\text{复化}}\) 复向量丛：\(p_k(E) = (-1)^k c_{2k}(E \otimes \mathbb{C})\)
\item
  复向量丛 \(\xrightarrow{\text{忘却复结构}}\) 实向量丛：\(c_i(E) \mapsto w_{2i}(E_{\mathbb{R}})\)（模 2）
\item
  Euler 类与顶维 Stiefel-Whitney 类：\(e(E) \equiv w_n(E) \pmod{2}\)
\item
  所有示性类 \(\xrightarrow{\text{Chern-Weil}}\) 曲率形式的多项式（de Rham 上同调）
\end{itemize}

示性类不仅是拓扑学的核心工具，更是现代理论物理中不可或缺的语言------从杨-Mills 理论的瞬子数（由第二 Chern 类 \(c_2\) 的积分给出）到弦论中的反常消除（要求 \(p_1/2 = 0\)），从拓扑绝缘体的 \(\mathbb{Z}_2\) 不变量到量子霍尔效应中的 Chern 数，示性类的思想渗透到物理学的多个前沿领域。

\hrulefill

\hrulefill

\hrulefill

\hrulefill

\hrulefill

\hrulefill

\newpage

\chapter{07-拓扑与几何/04-复几何/02-Hodge与Calabi-Yau.md}\label{ux62d3ux6251ux4e0eux51e0ux4f5504-ux590dux51e0ux4f5502-hodgeux4e0ecalabi-yau.md}

\section{第257章：Hodge 理论}\label{ux7b2c257ux7ae0hodge-ux7406ux8bba}

Hodge 理论是紧 Kahler 流形上同调理论的核心，由 W. V. D. Hodge 在 1930-1940 年代系统创立。其最基本结果是 Hodge 分解定理：在紧 Kahler 流形 \(M\) 上，每个 de Rham 上同调类都有唯一的调和形式代表元。更精确地，复值上同调群 \(H^k(M, \mathbb{C})\) 可以按 \((p,q)\) 类型分解为调和 \((p,q)\)-形式空间的直和 \(H^k(M, \mathbb{C}) = \bigoplus_{p+q=k} \mathcal{H}^{p,q}(M)\)，且满足 Hodge 对偶 \(\mathcal{H}^{p,q} \cong \overline{\mathcal{H}^{q,p}}\)。这一分解的几何意义极为深远：Hodge 数 \(h^{p,q} = \dim \mathcal{H}^{p,q}\) 不仅给出了 Betti 数的精细分解（\(\sum_{p+q=k} h^{p,q} = b_k\)），其 Hodge 菱形对称性还同时编码了 Serre 对偶、复共轭对称以及 Lefschetz 分解（通过 Kahler 类的楔积给出 \(H^{p,q} \to H^{p+1,q+1}\) 的同构）。在代数几何中，Hodge 理论是 Hodge 猜想（每个有理 \((p,p)\) 类必是代数闭链的有理线性组合）的基础，后者是 Clay 数学研究所的七大千禧问题之一。本章与形变理论（第200章）的联系在于：Hodge 结构随复结构的形变而变化，产生了 Griffiths 周期映射，这直接引向 Calabi-Yau 流形（第201章）的镜对称研究。

\subsection{179.1 调和形式与 Hodge 定理}\label{ux8c03ux548cux5f62ux5f0fux4e0e-hodge-ux5b9aux7406}

\textbf{定义 179.1}（调和形式）：\((p,q)\)-形式 \(\alpha\) 称为\textbf{调和}的，如果 \(\Delta_{\bar{\partial}} \alpha = 0\)。记 \(\mathcal{H}^{p,q}(M)\) 为调和 \((p,q)\)-形式空间。

\textbf{定理 179.1}（Hodge 定理，调和表示）：在紧 Kähler 流形上，每个 Dolbeault 上同调类有唯一的调和代表元：

\[H_{\bar{\partial}}^{p,q}(M) \cong \mathcal{H}^{p,q}(M)\]

且 \(\mathcal{H}^{p,q}(M)\) 是有限维的。

\emph{证明}：\(\Delta_{\bar{\partial}}\) 是椭圆算子（在紧流形上），其核是有限维的。\(\bar{\partial}\)-上同调类与 \(\Delta_{\bar{\partial}}\)-调和形式通过正交分解 \(H_{\bar{\partial}}^{p,q} \cong \ker \Delta_{\bar{\partial}}\) 建立同构。∎

\subsection{\texorpdfstring{179.2 Lefschetz 分解与 \(\mathfrak{sl}_2\)-表示}{179.2 Lefschetz 分解与 \textbackslash mathfrak\{sl\}\_2-表示}}\label{lefschetz-ux5206ux89e3ux4e0e-mathfraksl_2-ux8868ux793a}

Hodge 定理在每一上同调层次给出了调和形式的存在性，但它并未揭示不同次数的 Hodge 分量之间的代数关系。Lefschetz 分解正是描述这些关系的基本结构------它表明 Hodge 结构的各个分量实际上构成了 \(\mathfrak{sl}_2\) 表示论中的有限维不可约表示，从而将复几何与李代数的表示理论建立了深刻的联系。

\textbf{定义 179.2}（Lefschetz 算子与对偶 Lefschetz 算子）：定义 \(L: \Omega^{p,q} \to \Omega^{p+1,q+1}\) 为 \(L(\alpha) = \omega \wedge \alpha\)（Kähler 形式的外乘）。\((p,q)\)-形式 \(\alpha\) 称为\textbf{原始}的（primitive），如果 \(L^{n-p-q+1} \alpha = 0\)（\(n = \dim_{\mathbb{C}} M\)）。

\textbf{定理 179.2}（Lefschetz 分解，硬 Lefschetz 定理）：对紧 Kähler 流形，\(L^k: H^{n-k}(M, \mathbb{C}) \to H^{n+k}(M, \mathbb{C})\) 是同构（\(0 \leq k \leq n\)）。每个 \(k\)-形式可唯一分解为

\[\alpha = \sum_{j \geq \max(0, k-n)} L^j \alpha_j\]

其中 \(\alpha_j\) 是原始的 \((k-2j)\)-形式。

\emph{注}：硬 Lefschetz 定理是 Kähler 几何区别于一般复流形的最深刻特征之一。

\textbf{证明}：设 \(\Lambda = L^*\) 为 \(L\) 的伴随算子。\(\mathfrak{sl}_2\)-三元组 \(\{L, \Lambda, H = [L, \Lambda]\}\) 在 \(\Omega^*(M)\) 上有表示，\(H\alpha = (k-n)\alpha\) 对 \(k\)-形式。由 \(\mathfrak{sl}_2\)-表示论，\(\alpha\) 有唯一分解 \(\alpha = \sum L^j \alpha_j\)（\(\alpha_j\) 原始）。硬 Lefschetz 同构来自 \(L^k\) 在调和形式上诱导的同构（由 Kähler 恒等式 \(\Delta_d = 2\Delta_\partial\) 和 Lefschetz \(L\) 与 Laplacian 交换）。\(\blacksquare\)

\textbf{定理 179.3}（Lefschetz 超平面定理）：设 \(X \subset \mathbb{CP}^N\) 是光滑射影簇，\(H\) 是超平面截面。则包含映射 \(H \cap X \hookrightarrow X\) 诱导的同调群同态在维数 \(< \dim_{\mathbb{C}} X - 1\) 上是同构。

\textbf{证明}：由 Morse 理论，\(X \setminus H\) 是仿射簇，具有同伦型至多为 \(\dim_{\mathbb{C}} X\) 的 CW 复形。Lefschetz 铅笔论证表明 \(X \setminus H\) 上的相对同调在维数 \(< n-1\) 时为零。由 Alexander-Lefschetz 对偶，包含映射 \(i: H \cap X \hookrightarrow X\) 诱导同构 \(i_*: H_k(H \cap X) \to H_k(X)\)（\(k < n-1\)）。\(\blacksquare\)

\subsection{179.3 Hodge 菱形与 Hodge 猜想}\label{hodge-ux83f1ux5f62ux4e0e-hodge-ux731cux60f3}

\textbf{定义 179.3}（Hodge 菱形）：Hodge 菱形是 \(h^{p,q}\) 的阵列，行对应 \(p+q\)，列对应 \(p-q\)。Hodge 对称性 \(h^{p,q} = h^{q,p} = h^{n-p,n-q}\) 反映在菱形的对称性上。

\textbf{命题 179.1}（Hodge 指标定理）：紧 Kähler 曲面 \(M\) 的相交形式在 \(H^{1,1}(M) \cap H^2(M, \mathbb{R})\) 上的符号差为 \((1, h^{1,1} - 1)\)。

\textbf{证明}：由 Hodge 分解，\(H^2(M, \mathbb{R}) = H^{1,1}(M) \cap H^2(M, \mathbb{R}) \oplus (H^{2,0} \oplus H^{0,2}) \cap H^2(M, \mathbb{R})\)。Kähler 形式 \(\omega\) 给出 \(H^{1,1}\) 中的一个正类（\(\int_M \omega^2 > 0\)），故相交形式有一个正特征值。由 Lefschetz 分解，\(H^{1,1}\) 的其余部分 \(H^{1,1}_{\text{prim}}\) 中，\(\omega\) 的作用给出负定形式：对 \(\alpha \in H^{1,1}_{\text{prim}}\)，\(\int_M \alpha \wedge \alpha \wedge \omega^{n-2} < 0\)。对于 \(n=2\)（曲面），这给出 \(\int_M \alpha^2 < 0\)。故 \(H^{1,1}\) 的相交形式符号差为 \((1, h^{1,1} - 1)\)。\(\blacksquare\)

\textbf{定义 179.4}（Hodge 猜想）：设 \(X\) 是光滑复射影代数簇。对每个 \(k\)，有理 \((k,k)\)-类（即 \(H^{2k}(X, \mathbb{Q}) \cap H^{k,k}(X)\)）是否由代数子簇的类生成？这是 Clay 千禧年问题之一，至今未解决（\(k=1\) 时由 Lefschetz 的 \((1,1)\)-定理解决）。

\textbf{定理 179.4}（Lefschetz \((1,1)\)-定理）：紧 Kähler 流形 \(M\) 上的每个整 \((1,1)\)-类是线丛的 Chern 类。即 \(H^{1,1}(M) \cap H^2(M, \mathbb{Z}) \cong \operatorname{Pic}(M)\)（Néron-Severi 群）。

\emph{证明}：指数正合列 \(0 \to \mathbb{Z} \to \mathcal{O}_M \xrightarrow{\exp(2\pi i\cdot)} \mathcal{O}_M^* \to 0\) 的长正合序列给出 \(\operatorname{Pic}(M)=H^1(M,\mathcal{O}_M^*) \to H^2(M,\mathbb{Z}) \to H^2(M,\mathcal{O}_M)\)。像落在 \(H^2(M,\mathbb{Z})\) 在 \(H^2(M,\mathbb{C})=H^{2,0}\oplus H^{1,1}\oplus H^{0,2}\) 中的 \((1,1)\)-部分。由 Dolbeault，\(H^2(M,\mathcal{O}_M)=H^{0,2}\)，故核即整 \((1,1)\)-类。\(\blacksquare\)

\subsection{179.4 Hodge 结构的变分}\label{hodge-ux7ed3ux6784ux7684ux53d8ux5206}

Hodge 结构的本质不仅仅在于单个复流形上静态的 \((p,q)\)-分解，更在于当复流形在一族中连续变化时，这些分解如何随之演化。Griffiths 在 1960-70 年代建立的理论表明，Hodge 分解的变分受到严格的微分约束（横截性），而这些约束反过来通过周期映射精确控制了复结构的模空间几何。

\textbf{定义 179.5}（Hodge 结构）：权 \(k\) 的\textbf{有理 Hodge 结构}由有限维有理向量空间 \(V_{\mathbb{Q}}\) 和 \(V_{\mathbb{C}} = V_{\mathbb{Q}} \otimes \mathbb{C}\) 上的分解 \(V_{\mathbb{C}} = \bigoplus_{p+q=k} V^{p,q}\) 组成，满足 \(\overline{V^{p,q}} = V^{q,p}\)。等价地，Hodge 结构可用两个滤过刻画：\textbf{Hodge 滤过} \(F^p = \bigoplus_{i \geq p} V^{i,k-i}\) 是 \(V_{\mathbb{C}}\) 的下降滤过，满足 \(F^p \oplus \overline{F^{k-p+1}} \cong V_{\mathbb{C}}\)；\textbf{权滤过} \(W_\bullet\) 是定义在 \(V_{\mathbb{Q}}\) 上的递增滤过，满足 \(W_m = 0\)（\(m < 0\)）且 \(W_k = V_{\mathbb{Q}}\)。称 \((V_{\mathbb{Q}}, F^\bullet, W_\bullet)\) 为\textbf{混合 Hodge 结构}。当 \(W_{k-1} = 0\) 时即为\textbf{纯 Hodge 结构}。纯 Hodge 结构的基本例子包括：光滑复射影代数簇 \(X\) 的 \(k\) 阶上同调 \(H^k(X, \mathbb{Q})\) 带由其 Kähler 结构诱导的 Hodge 分解 \(H^k(X, \mathbb{C}) \cong \bigoplus_{p+q=k} H^{p,q}(X)\)；以及 Abel 簇（复环面）的 \(H^1\) 上的极化的纯 Hodge 结构。混合 Hodge 结构最典型的例子来自代数簇的补空间或奇异代数簇的上同调------Deligne 证明任何复代数簇的上同调都携带典范的混合 Hodge 结构，其权滤过反映了奇异性的贡献。

\textbf{定理 179.5}（Griffiths 横截性）：在 Hodge 簇的形变中，Hodge 分解满足 \(V^{p,q}\) 的微分位于 \(V^{p-1,q+1} \oplus V^{p,q} \oplus V^{p+1,q-1}\) 中。这是周期映射的基本性质。

\textbf{证明}：设 \(F_t^p = \bigoplus_{i\ge p} H^{i,n-i}_t\) 为 Hodge 滤过。形变方向 \(\frac{\partial}{\partial t}\) 作用下，由 Kodaira-Spencer 理论，速率为 \(\rho(\partial/\partial t) \cup \alpha\)（杯积作用）。因 \(\rho \in H^1(\Theta)\) 的杯积降低全纯次数一位，故 \(\frac{\partial}{\partial t} F^p \subseteq F^{p-1}\)。\(\blacksquare\)

\hrulefill

\hrulefill

\hrulefill

\hrulefill

\hrulefill

\hrulefill

\section{第258章：形变理论}\label{ux7b2c258ux7ae0ux5f62ux53d8ux7406ux8bba}

形变理论（deformation theory）研究复流形在保持复解析结构前提下的连续变动，由 Kunihiko Kodaira 和 Donald Spencer 在 1950-1960 年代通过发展复结构的偏微分方程方法系统建立。核心问题是：给定一个紧复流形 \(M\)，其''近邻''有多少个不同的复结构？Kodaira-Spencer 映射将这一问题转化为上同调的计算：一阶无穷小形变由 \(H^1(M, \Theta_M)\)（全纯切丛的上同调）参数化，而形变的障碍（即无穷小形变是否能积分为真正的形变族）位于 \(H^2(M, \Theta_M)\) 中。当 \(H^2(M, \Theta_M) = 0\) 时，Kodaira 定理保证所有一阶形变都能积分------形变空间在该点是光滑的，且其局部维数等于 \(h^1(M, \Theta_M)\)。形变理论在代数几何和镜对称中的关键应用包括：Calabi-Yau 流形（第201章）的复结构模空间的维数计算（在天野-田刚-Yau 定理中 \(\dim H^1(M, \Theta_M)\) 恰好是复结构模空间的维数），以及 Griffiths 周期映射通过 Hodge 结构（第199章）的形变来刻画模空间的几何。

\subsection{180.1 无穷小形变}\label{ux65e0ux7a77ux5c0fux5f62ux53d8}

形变理论的基本思想是研究复流形在保持其复解析结构的前提下能如何连续变化。这一问题由 Kodaira 和 Spencer 在 1950-60 年代系统地转化为上同调计算------这是一个典型的将几何问题翻译为代数问题（层上同调）的范例。

\textbf{定义 180.1}（形变）：紧复流形 \(M\) 的\textbf{形变}是光滑族 \(\pi: \mathcal{M} \to B\)（\(B\) 是连通复流形），使得 \(\pi^{-1}(0) = M\)（\(0 \in B\)）。\(M_t = \pi^{-1}(t)\) 是形变后的复流形。

\textbf{定义 180.2}（无穷小形变）：\(M\) 的\textbf{无穷小形变}空间是 \(H^1(M, \Theta_M)\)，其中 \(\Theta_M = \mathcal{O}(TM)\) 是全纯切丛的层。

\textbf{定理 180.1}（Kodaira-Spencer 映射）：设 \(\pi: \mathcal{M} \to B\) 是 \(M\) 的形变。则存在 Kodaira-Spencer 映射 \(\rho: T_0 B \to H^1(M, \Theta_M)\)，将形变方向映射到无穷小形变类。

\textbf{证明}：取切向量 \(\partial/\partial t \in T_0B\)。选择局部全纯坐标覆盖，转移函数 \(f_{\alpha\beta}(z,t)\) 的导数 \(\partial f_{\alpha\beta}/\partial t\) 定义 Cech 1-上闭链，值在 \(\Theta_M\) 中。其上同调类 \(\rho(\partial/\partial t) \in H^1(M,\Theta_M)\) 即 Kodaira-Spencer 类。不同坐标选取对应上边缘的改动。\(\blacksquare\)

\subsection{180.2 阻碍与 Kuranishi 族}\label{ux963bux788dux4e0e-kuranishi-ux65cf}

\textbf{定理 180.2}（Kodaira 阻碍定理）：形变存在阻碍，阻碍空间为 \(H^2(M, \Theta_M)\)。若 \(H^2(M, \Theta_M) = 0\)，则每个无穷小形变可积为真正的形变（\(M\) 无阻碍）。

\textbf{证明}：给定 \(\theta \in H^1(M,\Theta_M)\)，逐阶构造形变的幂级数。一阶可积性由 \(\bar{\partial}\theta=0\) 保证。二阶可积性要求 Maurer-Cartan 方程 \([\theta,\theta]=0 \in H^2(M,\Theta_M)\)。一般地，第 \(k\) 阶阻碍均落在 \(H^2(M,\Theta_M)\) 中。若该空间为零，则所有阻碍消失，形变全局可积。\(\blacksquare\)

Kodaira 阻碍定理回答了无穷小形变何时可积的问题，但它仍然留下了一个更精细的问题：当 \(H^2(M, \Theta_M) \neq 0\) 时，形变空间的局部结构是什么？Kuranishi 在 1965 年通过 Banach 尺度上的不动点定理给出了完整的刻画------形变空间在一个形式幂级数的零点附近被精确描述。

\textbf{定理 180.3}（Kuranishi 定理，1965）：紧复流形 \(M\) 存在\textbf{半通用形变}（Kuranishi 族），其参数空间是 \(H^1(M, \Theta_M)\) 在 \(0\) 处的一个邻域中由 \(H^2(M, \Theta_M)\) 定义的解析子集。即形变由

\[\mu: H^1(M, \Theta_M) \to H^2(M, \Theta_M), \quad \mu(\varphi) = [\varphi, \varphi]\]

的零点定义（Kuranishi 映射 \(\mu\) 是二次的）。

\textbf{证明}：选取 Hermitian 度量，定义 \(\bar{\partial}\) 的形式伴随 \(\bar{\partial}^*\) 和 Green 算子 \(G\)。构造映射 \(\Phi(\varphi) = \varphi + \frac{1}{2}\bar{\partial}^* G[\varphi,\varphi]\)，在 Sobolev 空间 \(H^s\) 上用不动点定理求解 \(\Phi(\varphi) = \varphi\)。解集 \(\mu^{-1}(0)\) 为 \(\varphi\) 满足 \(\bar{\partial}\varphi + \frac{1}{2}[\varphi,\varphi] = 0\) 的正则化饱和条件。由椭圆正则性，解是光滑的，给出有限维半通用族。\(\blacksquare\)

\subsection{180.3 形变量子化}\label{ux5f62ux53d8ux91cfux5b50ux5316}

\textbf{定义 180.3}（形变量子化）：形变量子化将光滑函数代数（Poisson 代数）形变为非交换代数。对辛流形 \((M, \omega)\)，形变量子化是形式幂级数环 \(\mathbb{C}[[\hbar]]\) 上的结合代数结构：

\[f \star g = fg + \frac{i\hbar}{2}\{f, g\} + \sum_{k \geq 2} \hbar^k B_k(f, g)\]

其中 \(B_k\) 是双微分算子，\(\star\) 满足结合律，且 \(f \star g - g \star f = i\hbar\{f, g\} + O(\hbar^2)\)。

\textbf{定理 180.4}（Kontsevich 形变量子化定理，1997）：任意 Poisson 流形上存在形变量子化。

\textbf{证明}：Kontsevich 构造星积 \(f \star g = \sum_{n\ge 0} \frac{(i\hbar)^n}{n!} \sum_{\Gamma \in G_n} w_\Gamma B_\Gamma(f,g)\)，其中 \(G_n\) 是含 \(n\) 个顶点的 Kontsevich 图，\(w_\Gamma\) 是上半平面配置空间上的角度积分（Maurer-Cartan 形式），\(B_\Gamma\) 是由 \(\pi\)（Poisson 双向量场）的多阶导数构成的多微分算子。结合性由 Stokes 定理和图的边收缩对应的边界项消去证明。这等价于形式性定理：多向量场的 \(L_\infty\)-代数形式等价于多微分算子的 \(L_\infty\)-代数。\(\blacksquare\)

\hrulefill

\hrulefill

\hrulefill

\hrulefill

\hrulefill

\hrulefill

\section{第259章：Calabi-Yau 流形与镜对称}\label{ux7b2c259ux7ae0calabi-yau-ux6d41ux5f62ux4e0eux955cux5bf9ux79f0}

Calabi-Yau 流形是第一 Chern 类为零（\(c_1(M) = 0\)）的紧 Kahler 流形，其重要性在现代数学物理中无可替代。Calabi 在 1954 年的著名猜想------具有给定 Kahler 类和消失的第一 Chern 类的紧 Kahler 流形上，存在唯一的 Ricci 平坦 Kahler 度量------由丘成桐（Shing-Tung Yau）在 1976 年通过求解复 Monge-Ampere 方程严格证明，这是几何分析最辉煌的胜利之一，丘成桐也因此获得 1982 年菲尔兹奖。Yau 定理的直接推论是：所有 \(c_1 = 0\) 的紧 Kahler 流形必定是 Ricci 平坦的，其 holonomy 群包含在 \(SU(n)\) 中------这正是超弦理论中''紧化''所需的几何条件。在理论物理中，Calabi-Yau 三维流形（复三维，实六维）是超弦理论中隐藏的''额外维''的候选模型，而镜对称（mirror symmetry）则是 20 世纪 90 年代由 Candelas、Greene、Vafa、Kontsevich 等人发展的惊人发现：每对镜对称的 Calabi-Yau 流形 \((M, W)\) 满足 \(h^{p,q}(M) = h^{n-p,q}(W)\)，且 \(M\) 上的 A-模型（Gromov-Witten 不变量）与 \(W\) 上的 B-模型（变分复结构）完全等价。本章与前章（Hodge 理论第 199 章和形变理论第 200 章）密切呼应------镜对称以 Hodge 菱形和周期映射为基本工具展开。

\subsection{181.1 Calabi 猜想与 Yau 定理}\label{calabi-ux731cux60f3ux4e0e-yau-ux5b9aux7406}

Calabi-Yau 流形在现代几何中占据核心地位，因为它们同时满足了代数几何、微分几何和数学物理三重视角下的极端条件：代数几何中 \(c_1=0\) 意味着典范丛是平凡的，微分几何中这等价于存在 Ricci 平坦度量，而物理中这提供了超弦理论紧化所需的六维内部空间。Calabi 在 1954 年提出了一个大胆的猜想：在这些代数约束下，每个 Kähler 类中恰有一个 Ricci 平坦度量。Yau 在 1976 年通过求解复 Monge-Ampere 方程证明了这一猜想，标志着几何分析的顶峰成就之一。

\textbf{定义 181.1}（Calabi-Yau 流形）：紧 Kähler 流形 \(M\) 称为 \textbf{Calabi-Yau 流形}，如果 \(c_1(M) = 0\)（在实上同调中）。等价地，\(M\) 的典范丛 \(K_M = \bigwedge^n T^*M\) 是平凡的。

\textbf{定理 181.1}（Calabi 猜想 / Yau 定理，1976-1978）：设 \(M\) 是紧 Kähler 流形，\(c_1(M) = 0\)。则 \(M\) 上存在唯一的 Ricci 平坦 Kähler 度量，其 Kähler 类等于任意给定的 Kähler 类。

\emph{意义}：Yau 定理证明了 Calabi 猜想，是 20 世纪几何学最伟大的成就之一。它保证了 Calabi-Yau 流形上 Ricci 平坦 Kähler 度量的存在性，为弦论中的 Calabi-Yau 紧化提供了数学基础。Yau 因此获 Fields 奖（1982）。

\textbf{注 181.1}（Yau 定理的证明概要）：Yau 的证明将问题转化为求解复 Monge-Ampere 方程：给定 Kähler 形式 \(\omega_0\) 和体积形式 \(\Omega\)（满足 \(\int_M \Omega = \int_M \omega_0^n\)），求 Kähler 势 \(\varphi\) 使得 \((\omega_0 + i\partial\bar{\partial}\varphi)^n = \Omega\) 且 \(\omega_0 + i\partial\bar{\partial}\varphi > 0\)。这是一个完全非线性的椭圆 PDE。Yau 使用\textbf{连续性方法}：构造从 \(\omega_0^n\) 到 \(\Omega\) 的一族体积形式，证明解的集合既开又闭。开性由椭圆算子的可逆性保证；闭性需要建立解的 \(C^0\)、\(C^2\) 和 \(C^{2,\alpha}\) 先验估计。其中 \(C^0\) 估计（即 \(\varphi\) 的有界性）是最困难的一步，Yau 通过精巧的 Moser 迭代和极大值原理完成。该定理的直接推论是：当 \(c_1(M) = 0\) 时，取 \(\Omega\) 为零曲率体积形式，得到的 Kähler 度量满足 \(\operatorname{Ric}(g) = 0\)，即 Ricci 平坦。

\textbf{证明}：Calabi 猜想等价于求解复 Monge-Ampère 方程 \(\det(g_{i\bar{\jmath}} + \partial_i\partial_{\bar{\jmath}}\varphi) = e^F \det(g_{i\bar{\jmath}})\)，其中 \(\varphi\) 是 Kähler 势、\(F\) 由 Ricci 形式约束给定。Yau 通过连续性方法证明该方程的可解性：(1) 构造一族方程 \((\omega + i\partial\bar{\partial}\varphi_t)^n = e^{tF + c_t}\omega^n\)（\(t\in[0,1]\)）；(2) \(t=0\) 时 \(\varphi_0=0\)（平凡解）；(3) 用 \(C^0\) 估计（Moser 迭代）、\(C^2\) 估计（二阶导数 Laplacian 比较）和 Calabi 的第三估计（\(C^3\) 估计）获得闭性；(4) 开性由线性化椭圆算子的可逆性得出（\(\Delta + \text{低阶项}\) 的指标为零，核平凡）。连续性方法在 \(t=1\) 给出 Ricci 平坦 Kähler 度量，其 Kähler 类等于原始类。唯一性由最大值原理保证。\(\blacksquare\)

\subsection{181.2 Calabi-Yau 3 维流形}\label{calabi-yau-3-ux7ef4ux6d41ux5f62}

\textbf{定义 181.2}（Hodge 菱形，CY 3 维流形）：Calabi-Yau 3 维流形的 Hodge 数为：

\[
\begin{array}{ccccccc}
& & & 1 & & & \\
& & 0 & & 0 & & \\
& 0 & & h^{1,1} & & 0 & \\
1 & & h^{2,1} & & h^{2,1} & & 1 \\
& 0 & & h^{1,1} & & 0 & \\
& & 0 & & 0 & & \\
& & & 1 & & &
\end{array}
\]

其中 \(h^{1,1}\) 和 \(h^{2,1}\) 是仅有的独立 Hodge 数。\(h^{1,1}\) 是 Kähler 模空间维数，\(h^{2,1}\) 是复结构模空间维数。

\textbf{定义 181.3}（五次 3 维流形）：\(\mathbb{CP}^4\) 中一般五次超曲面的 Hodge 数为 \(h^{1,1} = 1\)，\(h^{2,1} = 101\)。Euler 示性数 \(\chi = -200\)。这是研究最多的 Calabi-Yau 3 维流形。

\subsection{181.3 镜对称}\label{ux955cux5bf9ux79f0}

\textbf{定义 181.4}（镜对称猜想）：对每个 Calabi-Yau 3 维流形 \(M\)，存在\textbf{镜流形} \(M^\vee\)（也是 Calabi-Yau 3 维流形），满足：

\[h^{1,1}(M) = h^{2,1}(M^\vee), \quad h^{2,1}(M) = h^{1,1}(M^\vee)\]

即 Hodge 菱形沿 \(45^\circ\) 线翻转。更深层的镜对称断言 \(M\) 上的 \textbf{A-模型}（Gromov-Witten 不变量，与 Kähler 模空间相关）等价于 \(M^\vee\) 上的 \textbf{B-模型}（周期积分，与复结构模空间相关）。

\textbf{定理 181.2}（Candelas-de la Ossa-Green-Parkes, 1991）：五次 3 维流形上的有理曲线计数（通过镜对称计算）与代数几何直接计算一致。这是镜对称的第一个数值验证，开启了枚举几何的新纪元。

\textbf{证明}：五次 3 维流形 \(X\) 的镜 \(X^\vee\) 的周期积分满足 Picard-Fuchs 微分方程。其解在极大单值点（大半径极限）附近的展开系数预言了有理曲线数 \(n_d\)。Candelas 等计算得 \(n_1=2875\), \(n_2=609250\), \(n_3=317206375\) 等，与代数几何直接计算（使用椭圆曲线纤维化和模空间方法）吻合。\(\blacksquare\)

\textbf{定义 181.5}（Gromov-Witten 不变量）：\(M\) 的 \textbf{Gromov-Witten 不变量}计算 \(M\) 中给定亏格和同调类的全纯曲线的''虚拟''数目。它们是辛几何的不变量，在镜对称中对应 B-模型的周期积分。

\textbf{定理 181.3}（Kontsevich 镜对称定理，1994）：对环面簇，Gromov-Witten 不变量满足的微分方程等价于镜流形上周期积分满足的 Picard-Fuchs 方程。这给出了镜对称在环面簇情形下的精确表述。

\textbf{证明}：对光滑环面簇 \(X\)，Gromov-Witten 理论的基本对象是 \(J\)-函数（由稳定映射的模空间的拉回给出）。Kontsevich 证明 \(J\)-函数可表示为超几何级数，且满足的量子微分方程（\(D\)-模结构）与镜流形 \(X^\vee\) 上的 GKZ 超几何系统一致。等价性通过比较 Monodromy 群和初始值验证。\(\blacksquare\)

\subsection{181.4 SYZ 猜想}\label{syz-ux731cux60f3}

\textbf{定理 181.4}（Strominger-Yau-Zaslow 猜想，1996）：镜对称可以通过 Calabi-Yau 流形的 \textbf{T-对偶}实现：\(M\) 和 \(M^\vee\) 都是特殊 Lagrangian 3-环面 \(T^3\) 的纤维丛（在奇异纤维的极限下），且彼此通过交换纤维与基实现镜对称。

\textbf{证明}：在大半径极限下，CY 3-流形 \(M\) 退化为以 \(B\) 为基、\(T^3\) 为纤维的丛。沿纤维的 T-对偶（Fourier-Mukai 变换）生成镜 \(M^\vee\)：\(M\) 上的 Lagrangian 子流形对应于 \(M^\vee\) 上的全纯向量丛。Kontsevich-Soibelman 和 Gross-Siebert 通过热带几何和 log 结构的光滑化使该构造严格化。\(\blacksquare\)

SYZ 猜想将镜对称从代数几何的等价性转化为几何构造，是理解镜对称本质的关键框架。

\hrulefill

\emph{本卷建立了复几何的核心理论框架：复流形（近复结构、可积性、Dolbeault 上同调）将实微分几何推广到全纯范畴；Kähler 几何（Kähler 度量、Kähler 恒等式、Hodge 分解）是连接复结构、Riemann 度量和辛结构的典范框架；Hodge 理论（调和形式、Lefschetz 分解、Hodge 猜想）揭示了紧 Kähler 流形上同调的精细结构；形变理论（无穷小形变、Kuranishi 族）刻画了复结构的连续变化；Calabi-Yau 流形与镜对称（Yau 定理、镜对称）是复几何中最深刻的结果，连接了几何、代数、物理三大领域。复几何是当代数学几何化的核心语言。}

\hrulefill

\emph{本卷建立了复分析的完整理论框架：全纯函数与 Cauchy-Riemann 方程（Ch 73）是复分析的起点；复积分与 Cauchy 积分定理/公式（Ch 74）揭示了全纯函数的深层性质；级数展开与奇点分类（Ch 75）提供了分析工具；留数定理（Ch 76）实现了实积分的复计算；共形映射与 Riemann 映射定理（Ch 77）建立了区域间的解析对应；解析延拓与 Riemann 面（Ch 78）处理了多值性问题；Nevanlinna 值分布理论（Ch 90）量化了亚纯函数的取值行为；拟共形映射与 Teichmüller 空间（Ch 91）将共形几何推广至有界畸变；Siegel 区域与对称域（Ch 92）及多复变函数论（Ch 93-61）将理论推广至高维。补充内容涵盖了复几何（复流形、Kähler 几何、Hodge 理论、形变理论、Calabi-Yau 流形与镜对称），构成了复分析与微分几何的桥梁。}

\emph{卷九：复分析终。}

\emph{卷九：复分析终。}

\emph{卷九：复分析终。}

\emph{卷十一：复几何终。}

\hrulefill

\subsection{196.A 复流形补充}\label{a-ux590dux6d41ux5f62ux8865ux5145}

\begin{quote}
复几何是复分析和微分几何的交汇，研究复流形的几何结构。本部分作为复分析的几何补充，建立复流形、Kahler 几何、Hodge 理论、形变理论及 Calabi-Yau 流形与镜对称的系统理论。
\end{quote}

复几何的核心特征在于复流形的''刚性''------与光滑流形不同，复流形上的全纯函数由极大值原理在紧流形上必然为常数，这使得复流形范畴中的对象比实流形范畴中的对象更加''刚性''。但同时，Kahler 条件又为复流形注入了足够的''柔性''------Kahler 度量与 Ricci 曲率之间的关系（通过 Ricci 形式）产生了丰富的 Hodge 理论结构。复几何与代数几何的关系由 Serre 的 GAGA 原理和 Kodaira 嵌入定理架起桥梁：紧复流形是射影代数的当且仅当它容许一个 Hodge 度量（即 Kahler 形式是整上同调类）。

\hrulefill

\section{第260章：复流形}\label{ux7b2c260ux7ae0ux590dux6d41ux5f62}

复流形是局部同胚于 \(\mathbb{C}^n\) 且转移函数为全纯的拓扑空间。复流形的研究是复几何的起点。

\subsection{177.1 复流形的定义}\label{ux590dux6d41ux5f62ux7684ux5b9aux4e49}

复流形是微分几何与复分析交汇的核心概念，它将 \(\mathbb{C}^n\) 中全纯函数的结构推广到具有全纯转移函数的拓扑空间。复流形最本质的特征在于：它不仅是一个光滑流形（实 \(2n\) 维），而且其坐标卡之间的转移函数满足 Cauchy-Riemann 方程，即它们是全纯的。这一性质赋予了复流形极大的刚性------与光滑流形不同，紧复流形上的全纯函数由极大值原理必然为常数，这意味着复流形上的''函数''范畴比光滑流形上的''函数''范畴小得多。复流形的维数通常以复维数 \(n\) 计，其对应的实维数为 \(2n\)。复流形的基本例子包括：\(\mathbb{C}^n\) 本身、复射影空间 \(\mathbb{CP}^n\)（由 \(\mathbb{C}^{n+1} \setminus \{0\}\) 在标量乘法下的商空间，全纯坐标卡由齐次坐标的仿射片给出）、复环面 \(\mathbb{C}^n / \Lambda\)（\(\Lambda\) 为满秩格）、以及代数簇的光滑点集（Serre 的 GAGA 原理建立了复流形与代数簇之间的深刻联系）。复流形上的全纯映射 \(f: M \to N\) 是满足局部坐标表示为全纯函数的映射，且全纯映射的复合仍为全纯。复流形范畴中，双全纯映射（全纯且逆也为全纯）对应于复流形的同构，而双全纯等价类是复几何中的基本分类问题。紧复流形的基本不变量包括 Hodge 数 \(h^{p,q}\)、Chern 数和 Kodaira 维数，后者是 Enriques-Kodaira 分类理论的核心。

\textbf{定义 177.1}（复流形）：\(n\) 维\textbf{复流形} \(M\) 是一个拓扑空间，带有复坐标卡 \(\{(U_\alpha, \varphi_\alpha)\}\)，其中 \(\varphi_\alpha: U_\alpha \to \mathbb{C}^n\) 是同胚到 \(\mathbb{C}^n\) 中开集，且转移函数 \(\varphi_\alpha \circ \varphi_\beta^{-1}\) 在 \(\varphi_\beta(U_\alpha \cap U_\beta)\) 上是全纯的。

\textbf{定义 177.2}（全纯函数与全纯映射）：\(f: M \to \mathbb{C}\) 是\textbf{全纯}的，如果 \(f \circ \varphi_\alpha^{-1}\) 在每个坐标卡上是全纯的。\(f: M \to N\)（\(N\) 是复流形）是全纯的，如果局部坐标表示是全纯的。

\textbf{命题 177.1}（紧复流形上的全纯函数）：紧连通复流形上的全纯函数必为常数（极大值原理的推广）。这使复几何与实几何有本质区别。

\textbf{证明}：设 \(f: M \to \mathbb{C}\) 全纯，\(M\) 紧连通。则 \(|f|: M \to \mathbb{R}\) 是连续的，在紧集上达到最大值。设 \(|f(p)|\) 最大。取 \(p\) 附近的坐标卡，在该局部坐标下 \(f\) 是全纯函数。由极大值原理，\(f\) 在 \(p\) 的邻域中为常数。由连通性，\(f\) 在整个 \(M\) 上为常数。\(\blacksquare\)

\subsection{177.2 近复结构与可积性}\label{ux8fd1ux590dux7ed3ux6784ux4e0eux53efux79efux6027}

\textbf{定义 177.3}（近复结构）：实流形 \(M\)（\(2n\) 维）上的\textbf{近复结构}是切丛 \(TM\) 的自同态 \(J: TM \to TM\)，满足 \(J^2 = -\operatorname{id}\)。\(J\) 赋予 \(TM\) 复向量空间结构：\((a+ib) \cdot v = av + bJv\)。

复化切丛 \(TM \otimes \mathbb{C}\) 在 \(J\) 作用下分解为 \(\pm i\)-特征空间： \[TM \otimes \mathbb{C} = T^{1,0}M \oplus T^{0,1}M\] 其中 \(T^{1,0}M = \{X - iJX : X \in TM\}\)（全纯切空间），\(T^{0,1}M = \{X + iJX : X \in TM\}\)（反全纯切空间）。在局部坐标 \(\{z_j = x_j + iy_j\}\) 下，\(T^{1,0}M\) 由 \(\partial/\partial z_j\) 张成，\(T^{0,1}M\) 由 \(\partial/\partial \bar{z}_j\) 张成。

\textbf{定义 177.4}（Nijenhuis 张量）：近复结构 \(J\) 的 \textbf{Nijenhuis 张量}为

\[N_J(X, Y) = [X, Y] + J[JX, Y] + J[X, JY] - [JX, JY]\]

\(N_J\) 衡量 \(J\) 与 Lie 括号的不兼容性。\(N_J\) 是 \((1,2)\)-张量，且 \(N_J(JX, Y) = N_J(X, JY) = -J N_J(X, Y)\)。\(N_J = 0\) 等价于 \(T^{0,1}M\) 对 Lie 括号封闭（即 \([T^{0,1}, T^{0,1}] \subseteq T^{0,1}\)），这是 Frobenius 可积性定理的条件。

\textbf{定理 177.1}（Newlander-Nirenberg 定理，1957）：近复结构 \(J\) 来自复流形结构的充要条件是 \(N_J \equiv 0\)。此时 \(J\) 称为\textbf{可积的}。

\emph{意义}：该定理将复流形的存在性转化为近复结构的可积性条件，是复几何的核心定理。证明使用全纯坐标的构造（通过 \(\bar{\partial}\)-方程的可解性）。

\textbf{证明}：\(N_J=0\) 等价于 \(T^{1,0}M\) 对 Lie 括号封闭（Frobenius 定理的条件）。在局部坐标 \(\{z_i = x_i + i y_i\}\) 中，需构造全纯坐标函数 \(w_j\)。Newlander 与 Nirenberg 证明可解性归结为 Cauchy-Riemann 方程 \(\bar{\partial} w = 0\) 对非齐次项的估计，利用 \(N_J=0\) 时 \(\bar{\partial}^2=0\) 的形式可积性和 L\^{}2-估计（或 Newton 迭代）得到收敛的全纯坐标。具体地，利用 \(\bar{\partial}\)-Neumann 问题和 Hormander 的 \(L^2\) 估计，构造 \(\bar{\partial}\)-闭的 \((0,1)\)-形式的解，从而得到全纯坐标函数。\(\blacksquare\)

\textbf{命题 177.2}（实二维情形）：任何 2 维可定向 Riemann 流形上的近复结构（由旋转 \(90^\circ\) 定义）总是可积的。因此任何可定向 Riemann 曲面是复流形（1 维）。

\textbf{证明}：在实二维情形，\(N_J\) 自动为零（因为 \(N_J\) 的对称性导致维数约束）。任何可定向 Riemann 曲面可以通过等温坐标局部地赋予复坐标。\(\blacksquare\)

\subsection{177.3 Dolbeault 上同调}\label{dolbeault-ux4e0aux540cux8c03}

\textbf{定义 177.5}（\((p, q)\)-形式）：复流形 \(M\) 上的\textbf{\((p, q)\)-形式}是 \(p\) 个全纯微分和 \(q\) 个反全纯微分的外积。外微分 \(d\) 分解为 \(d = \partial + \bar{\partial}\)，其中 \(\partial\) 增加全纯次数，\(\bar{\partial}\) 增加反全纯次数。具体地， \[\partial = \sum_j dz_j \wedge \frac{\partial}{\partial z_j}, \quad \bar{\partial} = \sum_j d\bar{z}_j \wedge \frac{\partial}{\partial \bar{z}_j}\]

\(\partial^2 = 0\)，\(\bar{\partial}^2 = 0\)，\(\partial \bar{\partial} + \bar{\partial} \partial = 0\)（由 \(d^2 = 0\) 和类型的分解唯一性）。

\textbf{定义 177.6}（Dolbeault 上同调）：\(\bar{\partial}\)-算子满足 \(\bar{\partial}^2 = 0\)，定义\textbf{Dolbeault 上同调群}：

\[H_{\bar{\partial}}^{p,q}(M) = \frac{\ker(\bar{\partial}: \Omega^{p,q} \to \Omega^{p,q+1})}{\operatorname{im}(\bar{\partial}: \Omega^{p,q-1} \to \Omega^{p,q})}\]

\textbf{定理 177.2}（Dolbeault 定理）：\(H_{\bar{\partial}}^{p,q}(M) \cong H^q(M, \Omega^p)\)，其中 \(\Omega^p\) 是全纯 \(p\)-形式的层。Dolbeault 上同调是层上同调的具体实现。

\textbf{证明}：层 \(\Omega^p\) 的 Dolbeault 分解 \(0 \to \Omega^p \to \mathcal{A}^{p,0} \xrightarrow{\bar{\partial}} \mathcal{A}^{p,1} \xrightarrow{\bar{\partial}} \cdots\) 是 \(\Omega^p\) 的精细消解（fine resolution），因为 \(\mathcal{A}^{p,q}\) 是 \(\bar{\partial}\)-光滑形式层，即细层。由 Dolbeault-Grothendieck 引理（\(\bar{\partial}\)-Poincare 引理），此复形局部正合。层上同调的定义 \(H^q(M, \Omega^p) = R^q\Gamma(M, \Omega^p)\) 即该消解的截面复形的上同调，即 \(H_{\bar{\partial}}^{p,q}(M)\)。\(\blacksquare\)

\textbf{定义 177.7}（Hodge 数）：\(h^{p,q}(M) = \dim_{\mathbb{C}} H_{\bar{\partial}}^{p,q}(M)\)。Hodge 数是复流形的基本数值不变量。

\textbf{定理 177.3（Serre 对偶）}：设 \(M\) 是紧 \(n\) 维复流形。则 \(H^q(M, \Omega^p) \cong H^{n-q}(M, \Omega^{n-p})^*\)（对偶空间）。特别地，\(h^{p,q} = h^{n-p, n-q}\)。

\textbf{证明}：Serre 对偶是凝聚层上同调的 Serre 对偶定理的特殊情形。对偶层 \(\Omega^p\) 的 Serre 对偶为 \(\Omega^{n-p}\)（因为 \(\Omega^n\) 是典范丛，\(\Omega^p\) 与 \(\Omega^{n-p}\) 的对偶由 \(\wedge\) 配对给出）。利用层上同调的 Serre 对偶定理，得 \(H^q(M, \Omega^p)^* \cong \operatorname{Ext}^{n-q}(\Omega^p, \Omega^n) \cong H^{n-q}(M, \Omega^{n-p})\)。\(\blacksquare\)

\hrulefill

\section{第261章：Kahler 几何}\label{ux7b2c261ux7ae0kahler-ux51e0ux4f55}

Kahler 流形是同时具有 Riemann 度量、复结构和辛结构且三者兼容的流形。Kahler 几何是复几何的绝对核心。

\subsection{178.1 Kahler 度量的定义}\label{kahler-ux5ea6ux91cfux7684ux5b9aux4e49}

\textbf{定义 178.1}（Hermite 度量）：复流形 \(M\) 上的 \textbf{Hermite 度量} \(h\) 是每个切空间 \(T_pM\)（视为复向量空间）上的 Hermite 内积，光滑依赖于 \(p\)。分解为 \(h = g - i\omega\)，其中 \(g = \operatorname{Re} h\) 是 Riemann 度量，\(\omega = -\operatorname{Im} h\) 是实 \((1,1)\)-形式。\(g\) 和 \(\omega\) 满足 \(g(JX, JY) = g(X, Y)\) 和 \(\omega(X, Y) = g(JX, Y)\)。

\textbf{定义 178.2}（Kahler 度量）：Hermite 度量 \(h\) 称为 \textbf{Kahler 度量}，如果相应的 \((1,1)\)-形式 \(\omega\) 是闭的：\(d\omega = 0\)。此时 \((M, h)\) 称为 \textbf{Kahler 流形}。

Kahler 条件等价于：复结构 \(J\) 关于 \(g\) 的 Levi-Civita 联络是平行的：\(\nabla J = 0\)。这一等价性揭示了 Kahler 条件的本质：复结构和 Riemann 度量在 Levi-Civita 平行移动下完全兼容。

\textbf{证明}（\(\nabla J = 0 \iff d\omega = 0\)）：\(\omega(X, Y) = g(JX, Y)\)。计算 \(d\omega(X, Y, Z) = (\nabla_X \omega)(Y, Z) + (\nabla_Y \omega)(Z, X) + (\nabla_Z \omega)(X, Y)\)（利用 \(\nabla\) 无挠且度量相容）。\((\nabla_X \omega)(Y, Z) = g((\nabla_X J)Y, Z)\)。故 \(d\omega = 0\) 当且仅当 \(\nabla_X J\) 满足某种循环对称性，这等价于 \(\nabla J = 0\)。\(\blacksquare\)

\textbf{命题 178.1}（Kahler 势）：在局部坐标下，Kahler 形式 \(\omega\) 可写为

\[\omega = \frac{i}{2} \sum_{j,k} h_{j\bar{k}} dz_j \wedge d\bar{z}_k\]

其中 \(h_{j\bar{k}} = \frac{\partial^2 \varphi}{\partial z_j \partial \bar{z}_k}\)，\(\varphi\) 是局部 \textbf{Kahler 势}。\(d\omega = 0\) 自动满足。

Kahler 势的存在性体现了 Kahler 度量的一个关键特征：度量由单个实值函数 \(\varphi\) 的二阶导数给出。在紧 Kahler 流形上，全局 Kahler 势的存在性等价于 \(\omega\) 是恰当形式（即 \([\omega] = 0 \in H^2_{\text{dR}}(M)\)），这通常不成立。但 \(\partial \bar{\partial}\)-引理保证了局部 Kahler 势的存在性，且 \(\omega\) 的形变可由 \(\varphi\) 的形变描述。

\textbf{经典 Kahler 流形的例子}： - \(\mathbb{C}^n\) 配以标准 Hermite 度量 \(\omega = \frac{i}{2} \sum dz_j \wedge d\bar{z}_j\)，Kahler 势 \(\varphi = \frac{1}{2}\|z\|^2\) - \(\mathbb{CP}^n\) 配以 \textbf{Fubini-Study 度量}：在齐次坐标 \([Z_0: \cdots : Z_n]\) 下，\(\omega_{\text{FS}} = \frac{i}{2\pi} \partial \bar{\partial} \log \|Z\|^2\)，Kahler 势 \(\varphi = \log(1 + \|z\|^2)\)（在仿射坐标片 \(Z_0 = 1\) 上） - 复环面 \(\mathbb{C}^n/\Lambda\) 配以 \(\mathbb{C}^n\) 上平移不变的 Kahler 度量 - 光滑射影代数簇（由 \(\mathbb{CP}^n\) 的 Kahler 度量限制得到）

\subsection{178.2 Kahler 恒等式}\label{kahler-ux6052ux7b49ux5f0f}

\textbf{定理 178.1}（Kahler 恒等式）：在 Kahler 流形上，Hodge 星算子 \(\star\)、\(\partial\)、\(\bar{\partial}\) 及其伴随算子满足以下恒等式：

\begin{enumerate}
\def\labelenumi{\arabic{enumi}.}
\tightlist
\item
  \(\Delta_d = 2\Delta_\partial = 2\Delta_{\bar{\partial}}\)（Laplacian 的一致性）
\item
  \([\Lambda, \bar{\partial}] = -i \partial^*\)，\([\Lambda, \partial] = i \bar{\partial}^*\)（其中 \(\Lambda = L^*\) 是 \(\omega \wedge\) 的伴随）
\item
  \([L, \bar{\partial}^*] = i \partial\)，\([L, \partial^*] = -i \bar{\partial}\)
\end{enumerate}

其中 \(L = \omega \wedge \cdot\) 是 Lefschetz 算子，\(\Lambda\) 是其伴随。

这些恒等式是 Hodge 理论在 Kahler 流形上成立的核心技术原因。

\textbf{证明}：由 \(\omega\) 的 Kahler 条件 \(d\omega = 0\) 推出 \(\partial\) 和 \(\bar{\partial}\) 与 \(L\) 的交换子关系。利用局部坐标中 \(\omega = i\sum h_{ij} dz^i \wedge d\bar{z}^j\) 和 \(h_{ij}\) 的 Kahler 势表达式，计算 \(\partial^* = - \star \bar{\partial} \star\) 和 Lefschetz 算子 \(L = \omega \wedge \cdot\) 及其伴随 \(\Lambda = \star^{-1} L \star\) 的交换关系。初等的线性代数计算给出 \([\Lambda, \bar{\partial}] = -i\partial^*\) 和 \([L, \partial^*] = -i\bar{\partial}\)。由这些恒等式和 \(\Delta_d = d d^* + d^* d\) 的展开得 \(\Delta_d = \Delta_\partial + \Delta_{\bar{\partial}} + \text{cross terms}\)，交叉项由恒等式消去，故 \(\Delta_d = 2\Delta_\partial = 2\Delta_{\bar{\partial}}\)。\(\blacksquare\)

\textbf{定理 178.1.5（Lefschetz 分解定理）}：在紧 Kahler 流形 \(M\)（复 \(n\) 维）上，Lefschetz 算子 \(L^k: \bigwedge^{n-k} T^*M \to \bigwedge^{n+k} T^*M\) 诱导上同调层面的同构： \[L^k: H^{n-k}_{\text{dR}}(M, \mathbb{C}) \xrightarrow{\cong} H^{n+k}_{\text{dR}}(M, \mathbb{C})\] 且 de Rham 上同调群有 \textbf{Lefschetz 分解}： \[H^k_{\text{dR}}(M, \mathbb{C}) = \bigoplus_{j \geq \max(0, k-n)} L^j H^{k-2j}_{\text{prim}}(M, \mathbb{C})\] 其中 \(H^r_{\text{prim}}(M, \mathbb{C}) = \ker(L^{n-r+1}: H^r \to H^{2n-r+2})\) 是\textbf{本原上同调}（primitive cohomology）。

\textbf{证明}：Lefschetz 分解的关键在于 \(L\)、\(\Lambda\) 和 \(H = [L, \Lambda]\)（其中 \(H\) 在 \((p,q)\)-形式上作用为 \((p+q-n)\) 倍）满足 \(\mathfrak{sl}(2,\mathbb{C})\) 的交换关系：\([H, L] = 2L\)，\([H, \Lambda] = -2\Lambda\)，\([L, \Lambda] = H\)。de Rham 上同调群成为 \(\mathfrak{sl}(2,\mathbb{C})\) 的有限维表示，由 \(\mathfrak{sl}(2,\mathbb{C})\)-表示论的完全可约性得到 Lefschetz 分解。\(\blacksquare\)

\subsection{178.3 Hodge 分解定理}\label{hodge-ux5206ux89e3ux5b9aux7406}

\textbf{定理 178.2}（Hodge 分解定理，Kahler 情形）：紧 Kahler 流形 \(M\) 上的 de Rham 上同调（复系数）分解为

\[H^k_{\text{dR}}(M, \mathbb{C}) \cong \bigoplus_{p+q=k} H_{\bar{\partial}}^{p,q}(M)\]

且满足 \(H_{\bar{\partial}}^{p,q}(M) \cong \overline{H_{\bar{\partial}}^{q,p}(M)}\)（复共轭对称）。

\textbf{证明}：由 Kahler 恒等式（定理 178.1），\(\Delta_d = 2\Delta_\partial = 2\Delta_{\bar{\partial}}\)，故 \(\Delta_d\)-调和形式与 \(\Delta_{\bar{\partial}}\)-调和形式一致。对 \(\Delta_d\) 应用标准 Hodge 定理：紧 Riemann 流形上每个 de Rham 上同调类有唯一 \(\Delta_d\)-调和代表元。任取 \([\omega] \in H^k_{\text{dR}}(M, \mathbb{C})\)，其调和代表元 \(\alpha \in \ker \Delta_d\) 可按 \((p,q)\) 类型分解为 \(\alpha = \sum_{p+q=k} \alpha^{p,q}\)。因 \(\Delta_d = 2\Delta_{\bar{\partial}}\)，\(\Delta_d \alpha = 0\) 推出 \(\Delta_{\bar{\partial}} \alpha^{p,q} = 0\)（各 \((p,q)\) 分量独立），故每个 \(\alpha^{p,q}\) 是 \(\bar{\partial}\)-调和形式，对应 Dolbeault 上同调类 \([\alpha^{p,q}] \in H_{\bar{\partial}}^{p,q}(M)\)。该分配是良定义的：若另一调和代表元 \(\beta\) 满足 \(\alpha - \beta = d\gamma\)，由 Kahler 恒等式得 \(\alpha^{p,q} - \beta^{p,q} = \bar{\partial} \eta^{p,q-1}\)，故 \([\alpha^{p,q}] = [\beta^{p,q}]\)。逆映射将 \(\bar{\partial}\)-调和类求和即得。这给出同构 \(H^k_{\text{dR}}(M, \mathbb{C}) \cong \bigoplus_{p+q=k} H_{\bar{\partial}}^{p,q}(M)\)。复共轭对称性 \(H_{\bar{\partial}}^{p,q} \cong \overline{H_{\bar{\partial}}^{q,p}}\) 由 \(\Delta_{\bar{\partial}}\) 的实性推出。\(\blacksquare\)

\textbf{推论 178.3}（Hodge 数的对称性）：\(h^{p,q} = h^{q,p} = h^{n-p,n-q} = h^{n-q,n-p}\)。

\textbf{推论 178.4}（Betti 数的奇偶性约束）：紧 Kahler 流形的奇数阶 Betti 数 \(b_{2k+1}\) 为偶数。

\textbf{证明}：\(b_{2k+1} = \sum_{p+q=2k+1} h^{p,q}\)。由 \(h^{p,q} = h^{q,p}\)，\(h^{p,q} + h^{q,p} = 2h^{p,q}\) 是偶数，而 \(p = q\) 时 \(p+q = 2p\) 是偶数，不可能等于 \(2k+1\)。故 \(b_{2k+1}\) 是偶数。\(\blacksquare\)

\subsection{178.4 Calabi 猜想与 Yau 定理}\label{calabi-ux731cux60f3ux4e0e-yau-ux5b9aux7406-1}

\textbf{定理 178.5（Calabi-Yau 定理，Yau, 1978）}：设 \((M, \omega)\) 是紧 Kahler 流形，\(R\) 是任意表示 \(c_1(M)\) 的实 \((1,1)\)-形式。则存在唯一的 Kahler 度量 \(\tilde{\omega} \in [\omega]\)（同一 Kahler 类中），其 Ricci 形式 \(\rho(\tilde{\omega}) = R\)。

特别地，若 \(c_1(M) = 0\)，则 \(M\) 容许唯一的 Ricci 平坦 Kahler 度量（在每个 Kahler 类中）。这样的流形称为 \textbf{Calabi-Yau 流形}。

\textbf{证明}（框架）：Calabi 猜想等价于求解复 Monge-Ampere 方程： \[\frac{(\omega + i\partial \bar{\partial} \varphi)^n}{\omega^n} = e^F, \quad \omega + i\partial \bar{\partial} \varphi > 0\] 其中 \(F\) 由 \(R - \rho(\omega) = i\partial \bar{\partial} F\) 确定。Yau 使用连续性方法：构造一族方程 \(\frac{(\omega + i\partial \bar{\partial} \varphi_t)^n}{\omega^n} = e^{tF + c_t}\)（\(t \in [0,1]\)），并证明 \(t=1\) 时可解。关键步骤是建立 \(\varphi\) 的 \(C^0\) 估计（通过 Moser 迭代和 Sobolev 不等式）、\(C^2\) 估计（通过 Calabi 的第三阶估计技巧）和 \(C^{2,\alpha}\) 估计（通过 Evans-Krylov 理论或 Calabi 的 \(C^3\) 估计）。\(\blacksquare\)

Yau 定理是复几何和 Kahler 几何中最深刻的定理之一，其直接推论包括：\(K3\) 曲面存在 Ricci 平坦 Kahler 度量（从而 \(K3\) 曲面是超 Kahler 流形），以及任意次数的 Calabi-Yau 三维流形（如 \(\mathbb{CP}^4\) 中的五次超曲面）的 Ricci 平坦 Kahler 度量的存在性。这些结果在弦理论（超弦紧化需要 Calabi-Yau 三维流形）和镜对称（mirror symmetry）中具有根本重要性。

\subsection{178.5 Chern 类与 Chern-Weil 理论}\label{chern-ux7c7bux4e0e-chern-weil-ux7406ux8bba}

\textbf{定义 178.6}（Chern 类，复几何视角）：复向量丛 \(E \to M\) 的\textbf{第 \(k\) 个 Chern 类} \(c_k(E) \in H^{2k}(M, \mathbb{Z})\) 由曲率形式 \(\Omega\) 通过

\[\det\left(I + \frac{i}{2\pi} \Omega\right) = 1 + c_1(E) + c_2(E) + \cdots\]

定义。对复流形 \(M\)，\(c_k(M) = c_k(TM)\)。

\textbf{定理 178.7}（Chern-Weil 理论）：Chern 类在 de Rham 上同调中不依赖于联络的选择，且是复流形的微分同胚不变量。

\textbf{证明}：设 \(\nabla_0\) 和 \(\nabla_1\) 是同一向量丛上的两个联络，定义插值联络 \(\nabla_t = (1-t)\nabla_0 + t\nabla_1\)。其曲率形式 \(\Omega_t\) 满足 \(\frac{d}{dt}c_k(\Omega_t) = d(\text{transgression form})\)（通过 Chern-Simons 形式）。对 \(t\) 从 \(0\) 到 \(1\) 积分，得 \(c_k(\Omega_1) - c_k(\Omega_0) = d(\int_0^1 \Phi_t dt)\)，即差一个恰当形式，故 de Rham 类相同。微分同胚不变性来自光滑函子性。\(\blacksquare\)

\textbf{定义 178.8}（Chern 数）：紧复 \(n\) 维流形 \(M\) 的 \textbf{Chern 数}是 \(\int_M c_{i_1}(M) \wedge \cdots \wedge c_{i_k}(M)\)，其中 \(\sum i_j = n\)。

\textbf{定理 178.9（Kodaira 嵌入定理）}：紧复流形 \(M\) 是射影代数的（即可嵌入复射影空间）当且仅当 \(M\) 容许一个 Hodge 度量------即 Kahler 形式 \(\omega\) 是整上同调类（\([\omega] \in H^2(M, \mathbb{Z})\)）。

\textbf{证明}（框架）：若 \(M\) 是射影代数簇，则 \(\mathbb{CP}^N\) 的 Fubini-Study 度量的限制给出 Hodge 度量。反之，若 \(M\) 有 Hodge 度量 \(\omega\)，则 \(\omega\) 对应全纯线丛 \(L \to M\)（\(c_1(L) = [\omega]\)）。\(L\) 是正线丛（因为 \(\omega\) 是 Kahler 形式），由 Kodaira 消没定理，\(L^{\otimes k}\) 的整体截面在 \(k\) 充分大时分离点和切方向，从而给出嵌入 \(M \hookrightarrow \mathbb{P}(H^0(M, L^{\otimes k})^*)\)。\(\blacksquare\)

\textbf{定理 178.10（Kodaira 消没定理）}：设 \(M\) 是紧 Kahler 流形，\(L \to M\) 是正全纯线丛。则 \(H^q(M, \Omega^n(L)) = 0\) 对所有 \(q > 0\)。等价地，\(H^q(M, K_M \otimes L) = 0\)（\(q > 0\)），其中 \(K_M = \Omega^n\) 是典范丛。

\textbf{证明}：Kodaira 消没定理的证明使用 Bochner 技巧。正线丛 \(L\) 的曲率形式 \(\Theta(L)\) 是 Kahler 形式（至多差常数），利用 Weitzenbock 公式和 Hodge 定理，\(\bar{\partial}\)-Laplacian 与曲率项的关系给出 \(\Delta_{\bar{\partial}} \alpha = 0\) 推出 \(\alpha = 0\)（对 \(\alpha \in \Omega^{n,q}(L)\) 且 \(q > 0\)）。\(\blacksquare\)

Kodaira 消没定理是复几何中最重要的消没定理之一，它与 Hodge 分解定理、Kodaira 嵌入定理和 Kodaira-Spencer 形变理论一起构成了复几何的四大支柱。在代数几何中，Kodaira 消没定理的代数版本（Kawamata-Viehweg 消没定理）通过正特征约化方法得到推广，成为极小模型纲领（MMP）的核心技术工具。

\newpage

\chapter{07-拓扑与几何/05-几何分析/01-极小曲面与调和映射.md}\label{ux62d3ux6251ux4e0eux51e0ux4f5505-ux51e0ux4f55ux5206ux679001-ux6781ux5c0fux66f2ux9762ux4e0eux8c03ux548cux6620ux5c04.md}

\section{第262章：极小曲面}\label{ux7b2c262ux7ae0ux6781ux5c0fux66f2ux9762}

极小曲面是局部面积极小的曲面，是几何变分问题最经典的例子。Plateau 问题（1840年代提出）是极小曲面理论的起源。极小曲面理论横跨微分几何、偏微分方程、复分析和几何测度论，是几何分析中最富有成果的研究领域之一。

\subsection{182.1 极小曲面方程}\label{ux6781ux5c0fux66f2ux9762ux65b9ux7a0b}

\textbf{定义 182.1}（面积泛函）：曲面 \(\Sigma\) 的参数化 \(\mathbf{r}: U \to \mathbb{R}^3\) 的\textbf{面积}为

\[A(\mathbf{r}) = \int_U \sqrt{EG - F^2} \, du \, dv\]

其中 \(E, F, G\) 是第一基本形式的系数。\textbf{极小曲面}是面积泛函的临界点。

\textbf{定义 182.2}（平均曲率）：曲面 \(\Sigma\) 的\textbf{平均曲率} \(H = \frac{1}{2}(\kappa_1 + \kappa_2)\)（主曲率的平均值）。极小曲面满足 \(H \equiv 0\)。

\textbf{定理 182.1}（极小曲面方程）：非参数化曲面 \(z = u(x, y)\) 是极小曲面的充要条件是

\[\operatorname{div}\left(\frac{\nabla u}{\sqrt{1 + |\nabla u|^2}}\right) = 0\]

即

\[(1 + u_y^2)u_{xx} - 2u_x u_y u_{xy} + (1 + u_x^2)u_{yy} = 0\]

这是二阶拟线性椭圆型偏微分方程。

\textbf{证明}：计算面积泛函的第一变分：\(\delta A = \int_U \big( \frac{\partial}{\partial u}(\frac{F r_u - G r_v}{\sqrt{EG-F^2}}) + \frac{\partial}{\partial v}(\frac{F r_v - E r_u}{\sqrt{EG-F^2}}) \big) \cdot \delta r \, du dv\)。对图参数化 \(r(x,y) = (x, y, u(x,y))\)，第一基本形式系数代入得 \(E=1+u_x^2, F=u_x u_y, G=1+u_y^2\)。令第一变分为零，经计算得散度方程 \(\operatorname{div}\big(\frac{\nabla u}{\sqrt{1+|\nabla u|^2}}\big)=0\)。展开即得二阶方程。极小曲面方程也可写为 \(\Delta u = \frac{u_{xx}u_x^2 + 2u_{xy}u_x u_y + u_{yy}u_y^2}{1+u_x^2+u_y^2}\) 的形式，右端是 \(u\) 的二阶导数的二次型，体现了方程的非线性特征。\(\blacksquare\)

\textbf{定理 182.2}（Bernstein 定理，1915-1927）：\(\mathbb{R}^2\) 上整体定义的极小曲面图（\(z = u(x, y)\)，\(u\) 定义在整个 \(\mathbb{R}^2\) 上）必为平面。即 \(\mathbb{R}^3\) 中整体定义的极小曲面图只能是平面。

\emph{注}：Bernstein 定理在 \(\mathbb{R}^n\)（\(n \leq 7\)）中成立，在 \(\mathbb{R}^8\) 中不成立（Bombieri-De Giorgi-Giusti, 1969）。这是极小曲面理论与几何测度论的深刻联系。\(n \leq 7\) 的情形由 de Giorgi (1965)、Almgren (1966) 和 Simons (1968) 逐步解决，Simons 发现了 \(n=7\) 时存在稳定极小锥面（锥面 \(C = \{(x_1,\ldots,x_8): x_1^2+x_2^2+x_3^2+x_4^2 = x_5^2+x_6^2+x_7^2+x_8^2\}\)），Bombieri-De Giorgi-Giusti 证明该锥面是面积极小的，从而构造了 \(\mathbb{R}^8\) 中非平面的整体极小曲面图。

\textbf{证明}：由极小曲面方程 \((1+u_y^2)u_{xx}-2u_x u_y u_{xy}+(1+u_x^2)u_{yy}=0\)，定义 \(\theta = \arctan(u_y/u_x)\)（梯度方向角）。可证 \(\operatorname{div}(\nabla \theta/\sqrt{1+|\nabla u|^2}) = 0\)，即 \(\theta\) 满足某一致椭圆方程。由 Liouville 定理（整体有界调和解必为常数），\(\nabla u \equiv C\)（常向量），故 \(u\) 是线性函数（平面）。\(n \leq 7\) 的推广依赖 de Giorgi 的 Bernstein 定理版本和 Simons 的锥面论证。\(\blacksquare\)

\subsection{182.2 Weierstrass 表示公式}\label{weierstrass-ux8868ux793aux516cux5f0f}

\textbf{定理 182.2.5（Weierstrass-Enneper 表示）}：设 \(f, g\) 分别是 Riemann 面 \(M\) 上的全纯函数和亚纯函数，满足 \(g\) 的极点恰好是 \(f\) 的零点（且阶数匹配）。则

\[\mathbf{r}(z) = \operatorname{Re} \int_{z_0}^z \left(\frac{1}{2}f(1-g^2), \frac{i}{2}f(1+g^2), fg\right) d\zeta\]

给出 \(\mathbb{R}^3\) 中的（广义）极小曲面。其中 \(g\) 是此极小曲面的 Gauss 映射的球极投影（即 Gauss 映射与全纯函数的对应），\(f\) 是 Weierstrass 数据。

\textbf{证明}：直接验证：\(\mathbf{r}\) 满足等温条件 \(\mathbf{r}_z \cdot \mathbf{r}_z = 0\)（这等价于极小曲面条件 \(\mathbf{r}_{z\bar{z}} = 0\)）。计算 \(\mathbf{r}_z = (\frac{1}{2}f(1-g^2), \frac{i}{2}f(1+g^2), fg)\)，则 \(\mathbf{r}_z \cdot \mathbf{r}_z = \frac{1}{4}f^2(1-g^2)^2 - \frac{1}{4}f^2(1+g^2)^2 + f^2g^2 = 0\)。\(\blacksquare\)

Weierstrass 表示公式是将极小曲面构造转化为复分析问题的核心工具。通过选择适当的全纯数据 \((f, g)\)，可以构造出任意亏格、任意拓扑类型的极小曲面。Costa 曲面（亏格 1，三个端的嵌入极小曲面）的构造正是基于超椭圆 Riemann 面上的 Weierstrass 数据。

\subsection{182.3 Plateau 问题}\label{plateau-ux95eeux9898}

\textbf{定义 182.3}（Plateau 问题）：给定 \(\mathbb{R}^3\) 中的闭曲线 \(\Gamma\)，求以 \(\Gamma\) 为边界的面积最小的曲面。

\textbf{定理 182.3}（Douglas-Rado 解，1930-1931）：对任意可求长 Jordan 曲线 \(\Gamma\)，存在以 \(\Gamma\) 为边界的极小曲面（可能带有分支点）。Douglas 因此获首届 Fields 奖（1936）。

\emph{方法}：Douglas 将问题转化为 Dirichlet 积分（能量泛函）的极小化：

\[E(\mathbf{r}) = \frac{1}{2} \int_U (|\mathbf{r}_u|^2 + |\mathbf{r}_v|^2) \, du \, dv\]

在共形参数化下，\(E(\mathbf{r}) = A(\mathbf{r})\)，极小曲面等价于调和映射。

\textbf{证明}：利用变分法在 Sobolev 空间 \(H^1(D, \mathbb{R}^3)\) 中极小化 Dirichlet 能量。关键步骤：（1）构造极小化序列 \(\{X_n\}\) 使得 \(E(X_n) \to \inf E\)；（2）利用 Courant-Lebesgue 引理证明 \(X_n\) 在 \(\partial D\) 上等度连续（避免边界退化）；（3）由 Banach-Alaoglu 和紧嵌入得弱收敛子序列 \(X_n \rightharpoonup X\)，且下半连续性保证 \(E(X) = \inf E\)；（4）\(X\) 满足调和方程且共形（\(|X_u|=|X_v|\), \(X_u \cdot X_v=0\)），故为极小曲面。\(\blacksquare\)

\textbf{定理 182.4}（共形表示定理）：任何极小曲面 \(\Sigma \subset \mathbb{R}^3\) 存在等温坐标（共形参数化），使得 \(\mathbf{r}: U \to \Sigma\) 满足 \(\mathbf{r}_{uu} + \mathbf{r}_{vv} = 0\)（即分量是调和函数）且 \(|\mathbf{r}_u| = |\mathbf{r}_v|\)，\(\mathbf{r}_u \cdot \mathbf{r}_v = 0\)。

\textbf{证明}：取 \(\Sigma\) 的任意局部参数化，引入复坐标 \(z = u+iv\)。等温条件 \(|\mathbf{r}_u|^2 = |\mathbf{r}_v|^2\), \(\mathbf{r}_u \cdot \mathbf{r}_v = 0\) 等价于 \(\mathbf{r}_z \cdot \mathbf{r}_z = 0\)。由极小曲面方程 \(H=0\) 等价于 \(\mathbf{r}_{z\bar{z}} = 0\)，故 \(\mathbf{r}_z\) 是全纯函数。若 \(\mathbf{r}_z \neq 0\) 则存在局部全纯坐标变换使 \(\mathbf{r}_z \cdot \mathbf{r}_z = 0\) 成立，从而得到共形参数化。调和性由 \(\mathbf{r}_{uu}+\mathbf{r}_{vv}=4\mathbf{r}_{z\bar{z}}=0\) 得。\(\blacksquare\)

\subsection{182.4 极小曲面的例子}\label{ux6781ux5c0fux66f2ux9762ux7684ux4f8bux5b50}

极小曲面构成了 \(\mathbb{R}^3\) 中一类极为丰富且优美的曲面族。\textbf{平面}（\(\mathbb{R}^2\)）是最简单的极小曲面，也是 Bernstein 定理的平凡情形。\textbf{悬链面}（catenoid）：由悬链线 \(y = \cosh x\) 绕 \(x\) 轴旋转生成，参数化为 \(\mathbf{r}(u, v) = (\cosh u \cos v, \cosh u \sin v, u)\)，是仅有的极小旋转曲面（除平面外）。\textbf{螺旋面}（helicoid）：参数化为 \(\mathbf{r}(u, v) = (u \cos v, u \sin v, v)\)，是仅有的直纹极小曲面（由直线运动生成）。悬链面与螺旋面通过 Bonnet 变换（共形同胚的连续族）局部等距------这一关系揭示了极小曲面的 Weierstrass 表示理论中 Gauss 映射与全纯函数的深层对应。\textbf{Scherk 曲面}：由 \(z = \ln(\cos y / \cos x)\) 给出，是周期极小曲面的基本例子，具有鞍形结构。\textbf{Enneper 曲面}：参数化为 \(\mathbf{r}(u, v) = (u - u^3/3 + uv^2, v - v^3/3 + vu^2, u^2 - v^2)\)，是自交的代数极小曲面。\textbf{Costa 曲面}（Costa, 1982）：亏格 1 的完备嵌入极小曲面，具有三个端（end），拓扑上为环面去掉三个点。Hoffman 和 Meeks 进一步证明 Costa 曲面是嵌入的，并构造了任意高亏格的嵌入极小曲面。

\textbf{定理 182.5}（Costa-Hoffman-Meeks 定理）：存在任意高亏格的完备嵌入极小曲面。

\textbf{证明}：Costa（1982）构造了亏格 1 的 Costa 曲面，Hoffman-Meeks 证明它是嵌入的并通过对称柄的迭代构造推广到任意高亏格。构造基于 Weierstrass 表示公式，用超椭圆 Riemann 面上的全纯数据生成。\(\blacksquare\)

\subsection{182.5 几何测度论方法}\label{ux51e0ux4f55ux6d4bux5ea6ux8bbaux65b9ux6cd5}

\textbf{定义 182.4}（整变分流）：\textbf{整变分流}（integral varifold）是几何测度论中广义曲面的概念，允许奇点存在。变分流是面积泛函的临界点。

\textbf{定理 182.6}（Allard 正则性定理）：变分流在几乎处处意义下是光滑的极小曲面。奇点集是闭的且 Hausdorff 维数不超过 \(n-7\)（\(n\) 为外围空间的维数）。

\textbf{证明}：Allard 证明若变分流在某点处的密度比为 1（等价于该点是正则点），则该点附近变分流是光滑极小曲面。奇点处的密度比 \textgreater{} 1，通过单调性公式和维数约化（dimension reduction）论证奇点集的 Hausdorff 维数估计。\(n=7\) 时出现锥面奇点，\(n \leq 6\) 时无奇点。\(\blacksquare\)

\textbf{定理 182.7}（Almgren 大正则性定理，1983）：面积最小化整变分流在开稠密集上是光滑的。

\textbf{证明}：Almgren 证明奇异集（非正则点集合）的 Hausdorff 维数 \(\leq n-2\)（对 \(n\) 维外围空间中的质量极小化整变分流），远小于总维数，故光滑点集开且稠密。方法是用 blow-up 分析（在奇点附近缩放）结合 Federer 的维数约化技术，将奇异行为约化为平稳锥分类问题。\(\blacksquare\)

\subsection{182.6 极小曲面的稳定性}\label{ux6781ux5c0fux66f2ux9762ux7684ux7a33ux5b9aux6027}

\textbf{定义 182.5（稳定极小曲面）}：极小曲面称为\textbf{稳定的}，如果面积泛函的第二变分 \(\delta^2 A \geq 0\) 对所有紧支集法向变分成立。第二变分公式为：

\[\delta^2 A = \int_\Sigma (|\nabla^\perp \varphi|^2 - (|A|^2 + \operatorname{Ric}(N, N)) \varphi^2) dA\]

其中 \(A\) 是第二基本形式，\(\operatorname{Ric}\) 是外围空间的 Ricci 曲率，\(N\) 是单位法向量。在 \(\mathbb{R}^3\) 中，\(\operatorname{Ric}=0\)，稳定性条件简化为 \(\int_\Sigma |\nabla^\perp \varphi|^2 dA \geq \int_\Sigma |A|^2 \varphi^2 dA\)。

\textbf{定理 182.8（do Carmo-Peng, 1979; Fischer-Colbrie-Schoen, 1980）}：\(\mathbb{R}^3\) 中完备稳定极小曲面必为平面。这一结果建立了稳定极小曲面与 Bernstein 定理之间的深刻联系------在二维情形，稳定极小曲面只能是平面，而高维情形则允许非平面的稳定极小曲面。

\textbf{证明}：取 \(\varphi = |A|^{1+\alpha}\) 型测试函数代入第二变分不等式，利用 Simons 不等式 \(\Delta |A|^2 - 2|\nabla A|^2 + 2|A|^4 \geq 0\)（极小曲面上的 Simons 恒等式），通过积分估计得 \(|A| \equiv 0\)，故曲面为全测地（平面）。\(\blacksquare\)

\hrulefill

\section{第263章：调和映射}\label{ux7b2c263ux7ae0ux8c03ux548cux6620ux5c04}

调和映射是 Riemann 流形之间的映射，极小化能量泛函。它是极小曲面（调和函数）在流形间的推广。

\subsection{183.1 能量泛函与调和映射方程}\label{ux80fdux91cfux6cdbux51fdux4e0eux8c03ux548cux6620ux5c04ux65b9ux7a0b}

\textbf{定义 183.1}（能量泛函）：Riemann 流形 \((M, g)\) 到 \((N, h)\) 的光滑映射 \(f: M \to N\) 的\textbf{能量}（Dirichlet 能量）为

\[E(f) = \frac{1}{2} \int_M \|df\|^2 \, dV_g\]

其中 \(\|df\|^2 = \operatorname{tr}_g(f^*h)\) 是映射的 Hilbert-Schmidt 范数。在局部坐标下，\(\|df\|^2 = g^{ij} h_{\alpha\beta} \partial_i f^\alpha \partial_j f^\beta\)。

\textbf{定义 183.2}（调和映射）：\textbf{调和映射}是能量泛函 \(E(f)\) 的临界点。在局部坐标下，调和映射方程（Eells-Sampson 方程）为

\[\tau(f) = \Delta_g f^\alpha + \Gamma_{\beta\gamma}^\alpha(f) g^{ij} \frac{\partial f^\beta}{\partial x^i} \frac{\partial f^\gamma}{\partial x^j} = 0\]

其中 \(\tau(f)\) 是\textbf{张力场}（tension field），\(\Gamma_{\beta\gamma}^\alpha\) 是 \(N\) 的 Christoffel 符号。调和映射方程是半线性椭圆型方程组，其非线性项来自目标流形的曲率。

\subsection{183.2 Eells-Sampson 定理}\label{eells-sampson-ux5b9aux7406}

\textbf{定理 183.1}（Eells-Sampson 定理，1964）：设 \(M\) 是紧 Riemann 流形，\(N\) 是紧 Riemann 流形且截面曲率非正（\(K_N \leq 0\)）。则任何连续映射 \(f: M \to N\) 同伦于一个调和映射，且该调和映射在同伦类中极小化能量。

\textbf{证明}：使用热流方法。考虑演化方程 \(\frac{\partial f_t}{\partial t} = \tau(f_t)\)。当 \(K_N \leq 0\) 时，计算能量沿热流的演化： \[\frac{d}{dt}E(f_t) = -\int_M \|\tau(f_t)\|^2 dV_g \leq 0\] 故能量单调递减。由 Bochner 公式（见定理 183.2），\(\|df\|^2\) 沿热流满足抛物型不等式，这给出 \(\|df\|^2\) 的一致有界性。利用抛物 Schauder 估计，\(f_t\) 在 \(C^\infty\) 中有界，从而存在子序列 \(t_n \to \infty\) 使得 \(f_{t_n}\) 收敛到 \(f_\infty\)，且 \(f_\infty\) 满足 \(\tau(f_\infty) = 0\)（即调和映射）。\(\blacksquare\)

\textbf{定理 183.2}（Bochner 公式，调和映射情形）：对调和映射 \(f: M \to N\)，

\[\frac{1}{2} \Delta \|df\|^2 = \|\nabla df\|^2 + \sum_i \langle df(\operatorname{Ric}^M(e_i)), df(e_i) \rangle - \sum_{i,j} \langle R^N(df(e_i), df(e_j))df(e_j), df(e_i) \rangle\]

当 \(M\) 的 Ricci 曲率非负且 \(N\) 的截面曲率非正时，\(\|df\|^2\) 是下调和的，从而调和映射的能量极小。

\textbf{证明}：由 Weitzenbock 公式 \(\frac{1}{2}\Delta\|df\|^2 = \|\nabla df\|^2 + \langle df(\operatorname{Ric}^M(e_i)), df(e_i)\rangle - \langle R^N(df(e_i),df(e_j))df(e_j), df(e_i)\rangle\)。当 \(\operatorname{Ric}^M \geq 0\) 和 \(K_N \leq 0\) 时后两项 \(\geq 0\)，故 \(\Delta\|df\|^2 \geq 0\)（下调和）。由强极值原理，若 \(\|df\|^2\) 在 \(M\) 上某点达最大值则必为常数。对能量极小的调和映射，此常数必为零。\(\blacksquare\)

\subsection{183.3 调和映射的存在性与障碍}\label{ux8c03ux548cux6620ux5c04ux7684ux5b58ux5728ux6027ux4e0eux969cux788d}

\textbf{定理 183.3}（Sacks-Uhlenbeck 定理，1981）：设 \(M\) 是紧 Riemann 面，\(\pi_2(N) = 0\)。则任何连续映射 \(f: M \to N\) 同伦于一个能量极小调和映射。

\emph{方法}：引入 \(\alpha\)-能量 \(E_\alpha(f) = \int_M (1 + \|df\|^2)^\alpha \, dV\)（\(\alpha > 1\)），通过 \(\alpha \to 1\) 的极限获得调和映射。这种方法克服了 Sobolev 嵌入 \(W^{1,2} \not\subset L^\infty\) 的困难。

\textbf{证明}：对 \(\alpha > 1\)，\(E_\alpha\) 的 Euler-Lagrange 方程有更好的正则性（\(f \in W^{1,2\alpha}\) 由 Morrey 嵌入得 \(f\) 连续），存在极小化元 \(f_\alpha\)。由 \(\pi_2(N)=0\) 排除''气泡''（bubbling）现象------即 \(f_\alpha\) 在 \(\alpha \to 1\) 时不会产生非平凡调和球。由紧性，\(\{f_\alpha\}\) 子序列收敛到能量极小调和映射 \(f\)。若 \(\pi_2(N) \neq 0\)，气泡可能产生，导致极小子序列在有限点处能量集中。气泡分析（bubble tree construction）是 Sacks-Uhlenbeck 理论的核心创新。\(\blacksquare\)

\textbf{定理 183.4}（Siu 刚性定理，1980）：紧 Kahler 流形之间的调和映射在适当的曲率条件下必为全纯或反全纯的。这是强刚性定理的几何分析证明。

\textbf{证明}：设 \(f: M \to N\) 是调和映射，\(M, N\) 为 Kahler 流形。Bochner 公式给出 \(\int_M \|\bar{\partial} f\|^2 dV = \int_M \langle \partial f \wedge \bar{\partial} f, \partial f \wedge \bar{\partial} f \rangle\) 型积分。当 \(M\) 的曲率足够正且 \(N\) 的曲率足够负时，右边由曲率项控制且必为零。故 \(\bar{\partial} f = 0\) 或 \(\partial f = 0\)，即 \(f\) 为全纯或反全纯映射。关键是用 Kahler 恒等式将能量分为全纯和反全纯部分。\(\blacksquare\)

\textbf{定理 183.5（Schoen-Yau 调和映射与正质量定理）}：调和映射在广义相对论中有重要应用。Schoen 和 Yau (1979) 使用稳定极小曲面和调和映射技术证明了正质量定理：渐近平坦的初始数据集的总质量 \(m \geq 0\)，且 \(m=0\) 当且仅当初始数据集是 Minkowski 时空。证明的关键是构造从初始数据集到 \(\mathbb{R}\) 的调和映射，利用其能量极小性质导出质量非负。

\subsection{183.4 调和映射的奇点理论}\label{ux8c03ux548cux6620ux5c04ux7684ux5947ux70b9ux7406ux8bba}

\textbf{定理 183.6（Hélein 正则性，1990）}：从 Riemann 面到紧 Riemann 流形的弱调和映射（\(W^{1,2}\) 解）是光滑的。Hélein 证明的关键是发现调和映射方程具有''补偿紧性''结构------非线性项可写为 Jacobi 行列式形式，从而利用 Wente 不等式和 Hardy 空间估计得到连续性。

\textbf{定理 183.7（Hardt-Lin 奇点紧致性）}：调和映射序列的奇点集（气泡集）是 Hausdorff 维数至多为 \(m-2\) 的闭集（\(m = \dim M\)）。在奇点处，切映射是 \(\mathbb{R}^{m-2}\) 上的调和映射，且是''同变''的（关于 \(S^1\) 作用）。

\hrulefill

\section{第264章：Plateau 问题（变分视角）}\label{ux7b2c264ux7ae0plateau-ux95eeux9898ux53d8ux5206ux89c6ux89d2}

Plateau问题源于十九世纪比利时物理学家Plateau的肥皂膜实验：给定空间中的闭曲线\(\Gamma\)，求以\(\Gamma\)为边界的极小曲面。数学表述为在\(C^1\)参数化\(X: D \to \mathbb{R}^3\)（\(D\)为单位圆盘）的约束\(X(\partial D) = \Gamma\)下，最小化面积泛函\(\mathcal{A}[X] = \iint_D |X_u \times X_v| du dv\)。Douglas（1931）和Rado（1930）独立地使用变分方法证明了任何可求长Jordan曲线\(\Gamma\)都存在以它为边界的极小曲面，并获得了1936年首届Fields奖。其核心步骤为：(i)将面积泛函转化为Dirichlet能量\(\mathcal{E}[X] = \frac{1}{2}\iint_D (|X_u|^2 + |X_v|^2) du dv\)（极值下两者等价）；(ii)在Sobolev空间\(H^1(D, \mathbb{R}^3)\)中极小化，利用下半连续性和紧性得到弱解；(iii)证明弱解的光滑性。几何测度论将Plateau问题推广为：求整流传（integer multiplicity current）\(T\)使得\(\partial T = \Gamma\)且质量\(\mathbf{M}(T)\)极小，Federer-Fleming的闭包定理和紧性定理为此提供了严格的框架。

\textbf{定理 204.1（高维 Plateau 问题）}：设 \(M^{k-1} \subset \mathbb{R}^n\) 是紧嵌入光滑子流形。则存在以 \(M\) 为边界的面积最小化 \(k\) 维整流传（area-minimizing integral current）。解是光滑的，除了可能一个 Hausdorff 维数至多为 \(k-7\) 的闭奇点集。

\textbf{定理 204.2（Hardt-Simon 边界正则性）}：如果边界 \(\Gamma\) 是光滑的且满足某种''有界平均曲率''条件，则面积最小化解在边界附近是光滑的（直至边界）。这一结果对于理解肥皂膜的物理行为至关重要------肥皂膜的边界是固定的金属丝，而 Hardt-Simon 定理保证了膜在边界附近不会产生奇点。

极小曲面和调和映射理论的现代发展包括：平均曲率流（mean curvature flow）------以极小曲面为静止解的几何流；自由边界极小曲面------在 Riemann 流形内部以某子流形为自由边界的极小曲面（在物理中对应毛细管现象）；调和映射热流与 Yang-Mills 联络的对应（通过 Atiyah-Bott 的二维 Yang-Mills 理论）。这些方向将变分法、几何分析和数学物理紧密结合，构成了当代几何分析的活跃前沿。

\subsection{182.7 极小曲面方程与共形表示}\label{ux6781ux5c0fux66f2ux9762ux65b9ux7a0bux4e0eux5171ux5f62ux8868ux793a}

极小曲面与复分析之间的深刻联系通过共形表示理论得以建立。任何极小曲面 \(\Sigma \subset \mathbb{R}^3\) 可赋予 Riemann 面结构，使得 Gauss 映射 \(N: \Sigma \to S^2\) 是反全纯的（关于 \(S^2 \cong \mathbb{CP}^1\) 的复结构）。等价地，Gauss 映射的球极投影 \(g: \Sigma \to \mathbb{CP}^1\) 是全纯函数。这一全纯性使得 Weierstrass 表示理论成为可能，并且将极小曲面理论嵌入到 Riemann 面的全纯几何中。更为深刻的是，极小曲面的共形结构与其拓扑类型（亏格、端点数）之间存在约束关系------例如，完备嵌入极小曲面的总曲率满足 \(\int_\Sigma K dA \leq 2\pi(\chi(\Sigma) - r)\)，其中 \(r\) 是端点数（Cohn-Vossen 不等式在极小曲面上的推广）。这一不等式表明，极小曲面的总曲率（绝对值）必须足够大以补偿其拓扑复杂性，从而限制了完备嵌入极小曲面的拓扑类型。

\textbf{定理 182.9（Osserman 定理）}：完备极小曲面 \(\Sigma \subset \mathbb{R}^3\) 的全曲率 \(\int_\Sigma |K| dA\) 必为 \(4\pi\) 的整数倍，且 \(\int_\Sigma K dA \leq 2\pi(\chi(\Sigma) - r)\)，其中 \(r\) 为端点数。当等号成立时，\(\Sigma\) 是代数极小曲面（即由 Weierstrass 数据在紧 Riemann 面上定义）。

\textbf{证明}：由 Gauss-Bonnet 定理的推广（针对非紧完备曲面），总曲率 \(\int_\Sigma K dA = 2\pi(\chi(\Sigma) - \sum (\text{端的阶}))\)。每个端（end）的阶是 Gauss 映射在该端处的渐近行为的整数不变量。由 Weierstrass 表示，完备极小曲面的 Gauss 映射是亚纯函数，其奇点对应于端的渐近行为。\(\blacksquare\)

\textbf{定理 182.10（Meeks-Rosenberg 唯一性定理）}：\(\mathbb{R}^3\) 中由闭曲线生成的完备嵌入极小曲面（即具有紧边界且边界为凸曲线的极小曲面）是唯一的（在等距意义下）。这一结果将极小曲面理论与凸几何和曲率流联系起来。

\subsection{183.5 调和映射与 Teichmuller 理论}\label{ux8c03ux548cux6620ux5c04ux4e0e-teichmuller-ux7406ux8bba}

调和映射在 Teichmuller 理论中扮演核心角色。设 \(\Sigma\) 是亏格 \(g \geq 2\) 的闭 Riemann 面，\(\rho: \pi_1(\Sigma) \to \operatorname{PSL}(2,\mathbb{R})\) 是 Fuchsian 表示。则 \(\Sigma\) 上的双曲度量 \(g_\rho\) 是 \(\rho\) 的等变调和映射 \(\tilde{\Sigma} \to \mathbb{H}^2\) 的诱导度量。Eells-Sampson 的热流方法给出了双曲度量的唯一存在性证明，而 Wolf (1989) 证明了热流以指数速度收敛到唯一的调和映射。调和映射的参数化给出了 Teichmuller 空间 \(\mathcal{T}_g\) 与 \(\operatorname{Hom}(\pi_1(\Sigma), \operatorname{PSL}(2,\mathbb{R})) / \operatorname{PSL}(2,\mathbb{R})\) 的连通分支（Hitchin 分量）之间的微分同胚，这一对应是 Higgs 丛模空间的 Hitchin 系统的核心。Hitchin (1987) 证明：在紧 Riemann 面上，Higgs 丛模空间 \(\mathcal{M}_{\text{Higgs}}\) 是超 Kahler 流形，其 Hitchin 纤维化 \(\mathcal{M}_{\text{Higgs}} \to \bigoplus_{i=2}^n H^0(\Sigma, K_\Sigma^i)\) 是代数完全可积系统，其中的纤维是 Abel 簇的 Prym 簇。调和映射理论在这里提供了 Higgs 丛中的 Hermite-Einstein 联络与平坦联络之间的等价性（通过 Donaldson 和 Corlette 的非交换 Hodge 理论），是几何 Langlands 纲领的微分几何基础。

\newpage

\chapter{07-拓扑与几何/05-几何分析/02-Ricci流与共形几何.md}\label{ux62d3ux6251ux4e0eux51e0ux4f5505-ux51e0ux4f55ux5206ux679002-ricciux6d41ux4e0eux5171ux5f62ux51e0ux4f55.md}

\section{第265章：Ricci 流}\label{ux7b2c265ux7ae0ricci-ux6d41}

Ricci 流方程 \(\partial_t g = -2\operatorname{Ric}(g)\) 由 Richard Hamilton 在 1982 年引入，是 20 世纪几何分析最重大的创见之一------其基本想法是将 Ricci 曲率视为度量张量的''热扩散''驱动力，使得曲率在演化过程中趋于均匀，从而实现微分几何与抛物方程理论在最高层次上的统一。Hamilton 本人证明了正 Ricci 曲率的三维流形在 Ricci 流下必然收缩为球面（从而证明此类流形必微分同胚于 \(S^3/\Gamma\)），并在 1980 年代和 1990 年代发展了 Ricci 流的奇点分析、Harnack 不等式和紧致性定理。Perelman 在 2002-2003 年的三篇著名预印本（后由多位数学家证实）引入了熵泛函 \(\mathcal{W}\)、约化体积和 \(\kappa\)-解的关键概念，成功克服了奇点处''外科手术''的技术困难，从而完成了 Thurston 几何化猜想的完整证明------这不仅蕴含了 Poincare 猜想，还给出了任意闭三维流形的标准几何分解。在本书体系中，Ricci 流与 Yamabe 问题（第206章）共享共形几何的理论渊源：Yamabe 问题的变分结构事实上是 Ricci 流在共形类内的''静态''版本。

\subsection{184.1 Ricci 流方程}\label{ricci-ux6d41ux65b9ux7a0b}

\textbf{定义 184.1}（Ricci 流）：Riemann 度量 \(g(t)\) 的 \textbf{Ricci 流}方程为

\[\frac{\partial g_{ij}}{\partial t} = -2 \operatorname{Ric}_{ij}\]

这是度量的非线性扩散方程，类似热方程：曲率在正曲率区域减小，在负曲率区域增大。

\textbf{无量纲化 Ricci 流}（保持体积）：\(\frac{\partial g_{ij}}{\partial t} = -2 \operatorname{Ric}_{ij} + \frac{2}{n} \bar{R} g_{ij}\)，其中 \(\bar{R}\) 是平均标量曲率。

\textbf{命题 184.1}（曲率演化方程）：在 Ricci 流下，曲率张量满足

\[\frac{\partial R_{ijkl}}{\partial t} = \Delta R_{ijkl} + 2(B_{ijkl} - B_{ijlk} + B_{ikjl} - B_{iljk})\]

其中 \(B_{ijkl} = g^{pq} g^{rs} R_{pijr} R_{qkls}\)。标量曲率满足 \(\frac{\partial R}{\partial t} = \Delta R + 2|\operatorname{Ric}|^2\)。

\subsection{184.2 Hamilton 的定理}\label{hamilton-ux7684ux5b9aux7406}

\textbf{定理 184.1}（Hamilton 短时间存在性，1982）：对任意紧 Riemann 流形，Ricci 流在短时间 \([0, T)\) 上存在唯一解。

\textbf{证明}：将 Ricci 流方程 \(\partial_t g = -2\operatorname{Ric}(g)\) 视为非线性抛物型 PDE 系统。作 DeTurck 技巧：引入背景联络 \(\tilde{\nabla}\) 并考虑修正流 \(\partial_t g = -2\operatorname{Ric}(g) + \mathcal{L}_W g\)，其中 \(W^k = g^{ij}(\Gamma_{ij}^k - \tilde{\Gamma}_{ij}^k)\)。修正后方程变为严格抛物型 \(\partial_t g_{ij} = g^{kl}\tilde{\nabla}_k\tilde{\nabla}_l g_{ij} + \text{低阶项}\)。由标准抛物型 PDE 的短时间存在性理论（线性化加非线性迭代），存在 \(T>0\) 和唯一光滑解 \(g(t)\)。DeTurck 向量场 \(W\) 对应时间依赖的微分同胚，原 Ricci 流等价地短时间存在。唯一性由 Hamilton 的能量估计和最大值原理保证。\(\blacksquare\)

\textbf{定理 184.2}（Hamilton 的 3 维正 Ricci 曲率定理，1982）：设 \(M^3\) 是紧 3 维流形，初始度量具有正 Ricci 曲率。则 Ricci 流（经归一化）收敛到常正截面曲率度量。因此 \(M^3\) 微分同胚于球面空间形式 \(S^3 / \Gamma\)。

\textbf{证明}：设初始度量 \(\operatorname{Ric} > 0\)。在归一化 Ricci 流 \(\partial_t g = -2\operatorname{Ric} + \frac{2}{3}\bar{R}g\) 下，体积保持。由曲率张量演化方程和张量最大值原理，\(\operatorname{Ric} > 0\) 在演化中保持。Hamilton 证明曲率夹逼（pinching）估计：\(\frac{|\operatorname{Ric} - \frac{1}{3}Rg|^2}{R^2} \leq C R^{-\delta}\) 沿流单调递减至零。取缩放极限，曲率趋于常正截面曲率。由 Bishop-Gromov 体积比较和 Myers 定理，缩放极限的万有覆盖为 \(S^3\)。故 \(M^3\) 微分同胚于 \(S^3/\Gamma\)。\(\blacksquare\)

\emph{意义}：这是 Hamilton 开创 Ricci 流的动因------证明 3 维 Poincaré 猜想在正 Ricci 曲率下的特例。Hamilton 获 2011 年 Shaw 奖。

Hamilton 的正 Ricci 曲率定理（1982）首次展示了 Ricci 流作为几何化工具的潜力。然而，若要处理一般情形的奇异性，仅仅依赖曲率张量的点态估计是不够的------需要沿 Ricci 流演变的不变量来控制曲率的时空行为。这便是 Hamilton 在 1993 年引入微分 Harnack 不等式的动机。

\textbf{定理 184.3}（Hamilton 的 Harnack 不等式，1993）：对 Ricci 流的正曲率算子解，存在微分 Harnack 不等式，控制曲率的时间演化。

\textbf{证明}：设曲率算子 \(\operatorname{Rm} \geq 0\)。定义 Harnack 量 \(Z_{ijkl} = \Delta R_{ijkl} + 2(B_{ijkl} - B_{ijlk} + B_{ikjl} - B_{iljk}) + \frac{1}{t}R_{ijkl}\)。Hamilton 证明对任意向量场 \(U, V\)，二次型 \(Z(U, V) \geq 0\)。引入 Lie 导数沿梯度方向的时间导数，构造 Hamilton 矩阵型 Harnack 表达式 \(P_{ijk} = \nabla_i R_{jk} - \nabla_j R_{ik}\)，并证明 \(P\) 满足抛物型不等式。由最大值原理得时空梯度估计 \(\frac{|\nabla R|}{R} \leq \frac{C}{\sqrt{t}}\)，继而控制曲率的时间演化，等价于 \(t \cdot \partial_t R + R \geq 0\) 沿所有方向。\(\blacksquare\)

\textbf{定理 184.4}（Hamilton 的紧致性定理）：Ricci 流解序列在一定的曲率、直径和内射半径有界条件下，具有光滑收敛的子序列。

\textbf{证明}：设 \((M_k, g_k(t))\) 为完备 Ricci 流解序列，\(t \in [0, T]\)，满足 \(|\operatorname{Rm}| \leq C\)，\(\operatorname{diam} \leq D\)，\(\operatorname{inj} \geq \iota > 0\)。由 Cheeger-Gromov 紧致性定理，对每个 \(t\)，\((M_k, g_k(t))\) 在 Cheeger-Gromov 意义下子列收敛到光滑极限流形。Shi 的局部导数估计 \(\|\nabla^m \operatorname{Rm}\| \leq C_m t^{-m/2}\) 给出 \(C^\infty\) 有界性。由 Arzela-Ascoli 定理（在各时间片的 \(C^\infty\) 拓扑下），对角论证得子列在 \(C^\infty\) 中时空收敛。极限满足 Ricci 流方程（连续性）。\(\blacksquare\)

Hamilton 的 Harnack 不等式和紧致性定理为分析 Ricci 流的奇点结构提供了关键技术基础，但它们本身不足以处理奇点形成时流形的''塌缩''问题。Perelman 在 2002 年的突破正是从这一点切入的------他引入了一系列沿 Ricci 流单调的熵泛函，成功排除了塌缩现象。

\subsection{184.3 Perelman 的工作与 Poincaré 猜想（概述，详见 V13）}\label{perelman-ux7684ux5de5ux4f5cux4e0e-poincaruxe9-ux731cux60f3ux6982ux8ff0ux8be6ux89c1-v13}

下面给出 Perelman 在 2002-2003 年三篇里程碑式预印本中的核心定义和定理陈述。Grigory Perelman 的贡献是使用 Ricci 流（带手术）证明了 Thurston 几何化猜想（从而 Poincaré 猜想作为推论）。

Perelman 首次引入 \(W\)-熵泛函（2002 年的预印本中 §3），旨在利用 \(W\)-熵的单调性来控制沿 Ricci 流的有界熵推导出\textbf{紧性}（compactness），从而对任何极限中的奇点模型进行分类。\(W\)-熵的单调性是 Perelman 工作体系中最根本的分析工具：它确保了在 Ricci 流奇点处取缩放极限时，极限流形不会发生''塌缩''------即内射半径下有正下界控制。

\textbf{定理 184.5}（Perelman 的 \(\mathcal{F}\)-泛函与 \(\mathcal{W}\)-泛函，2002）：Perelman 引入两个关键泛函：

\[\mathcal{F}(g, f) = \int_M (R + |\nabla f|^2) e^{-f} \, dV\]

\[\mathcal{W}(g, f, \tau) = \int_M [\tau(R + |\nabla f|^2) + f - n] (4\pi\tau)^{-n/2} e^{-f} \, dV\]

\(\mathcal{F}\) 沿 Ricci 流（与共轭热方程耦合）单调不减。\(\mathcal{W}\) 沿 Ricci 流（与 \(\tau\) 的适当演化耦合）单调不减。

\textbf{证明}：\(\mathcal{F}\) 沿 Ricci 流 \(\partial_t g = -2\operatorname{Ric}\) 和共轭热方程 \(\partial_t f = -\Delta f + |\nabla f|^2 - R\) 的演化：\(\frac{d}{dt}\mathcal{F}(g,f) = 2\int_M |\operatorname{Ric} + \nabla^2 f|^2 e^{-f} dV \geq 0\)。\(\mathcal{W}\) 沿 Ricci 流和 \(\partial_t\tau = -1\) 以及 \(\partial_t f = -\Delta f + |\nabla f|^2 - R + \frac{n}{2\tau}\) 的演化：\(\frac{d}{dt}\mathcal{W} = 2\tau\int_M |\operatorname{Ric} + \nabla^2 f - \frac{1}{2\tau}g|^2 (4\pi\tau)^{-n/2}e^{-f}dV \geq 0\)。单调性分别由 \(\mu(g,\tau) = \inf\{\mathcal{W}(g,f,\tau) : \int (4\pi\tau)^{-n/2}e^{-f}dV = 1\}\) 和 \(\lambda(g) = \inf\{\mathcal{F}(g,f) : \int e^{-f}dV = 1\}\) 的定义得到。\(\blacksquare\)

\textbf{定理 184.6}（非塌缩定理）：Perelman 证明在 Ricci 流下，曲率有界的区域体积非塌缩（内射半径有下界）。这排除了 Cheeger-Gromov 意义下的''塌缩''现象。

\textbf{证明}：设 \((M, g(t))\) 为 Ricci 流解，\(p \in M\)，\(r > 0\) 满足 \(|\operatorname{Rm}| \leq r^{-2}\) 在 \(B_t(p, r)\) 上。Perelman 利用 \(\mathcal{W}\)-泛函的单调性：选取 \(f\) 近似于 \(\mathcal{W}\)-极小化函数，定义约化体积 \(\tilde{V}(t) = \int_M (4\pi\tau)^{-n/2} e^{-f} dV\)。\(\mathcal{W}\) 单调性蕴含 \(\tilde{V}\) 在倒向时间下单调递减，导出 \(t=0\) 处 \(B_0(p, r)\) 的体积有下界 \(V \geq c r^n\)（\(c>0\) 为万有常数）。由 Bishop-Gromov 体积比较，内射半径 \(\operatorname{inj}(p) \geq c' r\)。\(\blacksquare\)

非塌缩定理排除了 Ricci 流奇点分析中最棘手的障碍：它确保了在奇点处缩放后所得的规范模型（canonical models）具有正的内射半径，从而可以被有效地分类。在建立了奇点模型的分类型论之后，下一步自然是在 Ricci 流遇到奇点时执行几何手术------切除高曲率区域并粘贴标准帽，使得流可以继续光滑演化。

\textbf{定理 184.7}（Ricci 流的手术理论，2003）：当曲率在局部区域趋向无穷时，Perelman 的手术（surgery）程序切除高曲率区域，粘贴标准帽（standard cap），然后继续 Ricci 流的演化。手术次数有限。

\textbf{证明}：在 Ricci 流演化中，当曲率在点 \(p\) 处满足 \(|\operatorname{Rm}|(p) \geq \Theta\)（手术阈值）时，检查 \(p\) 的邻域是否为 \(\varepsilon\)-颈（\(\varepsilon\)-neck）。若为 \(\varepsilon\)-颈，在其狭窄处切断，粘贴标准解（standard solution，即 \(S^2\times\mathbb{R}\) 上的旋转对称 Ricci 流解）。手术参数 \(\varepsilon, \Theta\) 经精细选取确保：(1) 手术后的度量保持 Ricci 流初始条件；(2) 每次手术消耗一定量的 \(\mathcal{W}\)-熵，但 \(\mathcal{W}\) 有下界（由拓扑决定），故手术仅发生有限次；(3) 手术后的 Ricci 流可继续演化至下一次手术或消失。\(\blacksquare\)

\textbf{定理 184.8}（Poincaré 猜想，Perelman 2002-2003）：任何单连通闭 3 维流形同胚于 \(S^3\)。

\textbf{证明}：任取 \(M^3\) 上的初始度量，运行 Ricci 流（带手术）。Perelman 证明手术后的 Ricci 流在有限时间内消失（extinction），这意味着 \(M^3\) 是连通和 \(S^3 / \Gamma_i\) 的连通和。由基本群条件，\(M^3 \cong S^3\)。∎

\textbf{定理 184.9}（Thurston 几何化猜想，Perelman 2003）：任何闭 3 维流形可沿球面和环面分解为具有 8 种 Thurston 几何之一的片段。

\emph{意义}：Perelman 因证明 Poincaré 猜想获 2006 年 Fields 奖（但拒绝领奖），是 21 世纪数学第一个重大突破。

\textbf{证明}：运行带手术的 Ricci 流。Perelman 证明：（1）有限时间奇性仅在颈（neck）结构处发生，手术可移除颈并恢复光滑流；（2）手术次数有限（由 \(\int R^2\) 的单调下降保证）；（3）长时间极限流易分解为双曲片段（Ricci 流缩放至常负曲率）和 Seifert 片段或球面片段。结合 JSJ 分解的拓扑刚性和 Mostow 刚性完成。\(\blacksquare\)

\subsection{184.4 Ricci 流在复几何中的应用}\label{ricci-ux6d41ux5728ux590dux51e0ux4f55ux4e2dux7684ux5e94ux7528}

\textbf{定理 184.10}（Cao 定理，1985）：紧 Kähler 流形上的 Kähler-Ricci 流 \(\frac{\partial g_{i\bar{j}}}{\partial t} = -R_{i\bar{j}}\) 在 \(c_1(M) \leq 0\) 时收敛到 Kähler-Einstein 度量。

\textbf{证明}：Kähler-Ricci 流在 Kähler 类内演化为 \(\frac{\partial}{\partial t} \omega_t = -\operatorname{Ric}(\omega_t)\)。若 \(c_1(M)=0\)，此为 Yau 的定理（Ricci 平坦度量）。若 \(c_1(M)<0\)，演化为 \(\frac{\partial}{\partial t} \omega_t = -\operatorname{Ric}(\omega_t) - \omega_t\)，通过 Aubin-Yau 连续性方法的抛物化得到指数收敛到唯一的 Kähler-Einstein 度量（\(\operatorname{Ric} = -\omega\)）。\(\blacksquare\)

Cao 定理将 Ricci 流的应用从实几何扩展到了复几何框架。当 \(c_1(M)>0\)（即 Fano 流形情形）时，问题变得更加微妙------Kähler-Ricci 流的收敛性等价于流形的 \(K\)-稳定性，这是一个代数几何概念。这一深刻的联系由 Chen-Donaldson-Sun 在 2014 年完整证明，完成了 Yau-Tian-Donaldson 猜想的最后一环。

\textbf{定理 184.11}（Chen-Donaldson-Sun 定理，2014）：Fano 流形上 Kähler-Einstein 度量存在的充要条件是 \(K\)-稳定性。这是 Yau-Tian-Donaldson 猜想的证明，使用 Kähler-Ricci 流方法。

\textbf{证明}：必要性由 Tian 和 Donaldson 证明（若 KE 度量存在则 \(K\)-稳定）。充分性由 Chen-Donaldson-Sun 通过 Kähler-Ricci 流的解析紧化证明：若 \(K\)-稳定，则沿 Kähler-Ricci 流的奇点有限且可由代数几何的破坏边模（destabilizing test configuration）探测；这给出矛盾故 Ricci 流收敛到 KE 度量。\(\blacksquare\)

\subsection{184.5 Cheeger-Gromoll 分裂定理}\label{cheeger-gromoll-ux5206ux88c2ux5b9aux7406}

Cheeger-Gromoll 分裂定理是关于非负 Ricci 曲率完备 Riemann 流形结构的深刻结果，是 Ricci 流理论的几何背景之一。

\textbf{定理 184.12}（Cheeger-Gromoll 分裂定理，1971）：设 \((M, g)\) 是完备的 \(n\) 维 Riemann 流形，满足 \(\operatorname{Ric} \geq 0\)（Ricci 曲率处处非负）。若 \(M\) 包含一条直线（即等距嵌入 \(\gamma: \mathbb{R} \to M\) 满足 \(d(\gamma(s), \gamma(t)) = |s-t|\) 对所有 \(s,t\)），则 \(M\) 等距分裂为乘积 \(M \cong \mathbb{R} \times N\)，其中 \(N\) 是完备的 \((n-1)\) 维 Riemann 流形且 \(\operatorname{Ric}_N \geq 0\)。该定理使用 Busemann 函数和极大值原理证明：直线 \(\gamma\) 定义了互为反方向的 Busemann 函数 \(b^\pm\)，其和 \(b^+ + b^-\) 在 \(M\) 上恒为零且梯度为平行向量场，从而给出分裂的测地叶状结构。推论包括：\(\mathbb{R}^n\) 是唯一含直线的正截面曲率完备流形；任何 \(\operatorname{Ric} \geq 0\) 的紧流形的万有覆盖等距分裂为 \(\mathbb{R}^k \times N\)，其中 \(N\) 是紧的。

\textbf{证明}：设 \(\gamma:\mathbb{R}\to M\) 为直线（单位速度测地线，其任意子段均极小）。对每个 \(t\in\mathbb{R}\)，定义 Busemann 函数 \(b^+(x)=\lim_{t\to+\infty}(d(x,\gamma(t))-t)\) 和 \(b^-(x)=\lim_{t\to-\infty}(d(x,\gamma(t))+t)\)。由 Ric\(\ge 0\) 和 Laplacian 比较定理，\(b^+\) 和 \(b^-\) 均为调和函数（\(\Delta b^\pm = 0\)），且 \(b^++b^-=0\)。其梯度 \(\nabla b^+\) 为平行向量场，故 \(M\) 局部裂为 \(\mathbb{R}\times N\)。完备性确保全局分裂：\(M\cong\mathbb{R}\times N\) 为等距，其中 \(N\) 是满足 Ric\(\ge 0\) 的完备流形。\(\blacksquare\)

\hrulefill

\hrulefill

\hrulefill

\hrulefill

\hrulefill

\hrulefill

\section{第266章：Yamabe 问题与共形几何}\label{ux7b2c266ux7ae0yamabe-ux95eeux9898ux4e0eux5171ux5f62ux51e0ux4f55}

Yamabe 问题是最具历史深度的几何分析问题之一：给定紧 Riemann 流形 \((M^n, g)\)（\(n \geq 3\)），在度量 \(g\) 的共形类 \([g]\) 中是否存在常标量曲率度量？这一问题由 Yamabe Hidehiko 在 1960 年提出，给出的证明包含一个深刻错误（被 Neil Trudinger 在 1968 年发现），最终由 Thierry Aubin（1976 年，对于非共形平坦且 \(n \geq 6\) 的情形）和 Richard Schoen（1984 年，通过正质量定理完成剩余情形）彻底解决。从分析角度看，Yamabe 问题等价于求解半线性椭圆方程 \(-\Delta_g u + \frac{n-2}{4(n-1)} R_g u = \lambda u^{\frac{n+2}{n-2}}\)（共形 Laplace/Yamabe 方程），其非线性指数 \(\frac{n+2}{n-2}\) 恰是 Sobolev 嵌入的临界指数------这使得变分问题失去了紧性，需要 Aubin 和 Schoen 发展的集中紧性原理和爆破分析技巧。在几何学中，Yamabe 不变的符号是度量共形类的核心分类数据：正 Yamabe 不变量的流形容许正标量曲率度量（与正质量定理和广义相对论中的正能量条件密切相关），而零和负的 Yamabe 不变量则通向非紧致流形的共形几何。本章与 Ricci 流（第205章）之间的关系是：Ricci 流方程在共形类内取定标量曲率常值的''静态解''，恰好对应于 Yamabe 问题的极值度量。

\subsection{185.1 Yamabe 问题}\label{yamabe-ux95eeux9898}

\textbf{定义 185.1}（Yamabe 问题）：设 \((M^n, g)\) 是紧 Riemann 流形（\(n \geq 3\)）。在 \(g\) 的共形类 \([g] = \{u^{4/(n-2)} g : u > 0, u \in C^\infty(M)\}\) 中，是否存在常标量曲率的度量？

\textbf{共形变换下的标量曲率}：若 \(\tilde{g} = u^{4/(n-2)} g\)，则

\[\tilde{R} = u^{-(n+2)/(n-2)} \left(Ru - \frac{4(n-1)}{n-2} \Delta_g u\right)\]

令 \(\tilde{R} = \text{const}\)，得到共形 Yamabe 方程：

\[-\frac{4(n-1)}{n-2} \Delta_g u + Ru = \lambda u^{(n+2)/(n-2)}\]

这是临界 Sobolev 指数的半线性椭圆 PDE。

\textbf{定义 185.2}（Yamabe 不变量）：度量 \(g\) 的共形类的 \textbf{Yamabe 不变量}（Yamabe 商）为

\[Y(g) = \inf_{u \in C^\infty, u > 0} \frac{\int_M (\frac{4(n-1)}{n-2}|\nabla u|^2 + Ru^2) dV_g}{(\int_M u^{2n/(n-2)} dV_g)^{(n-2)/n}}\]

\subsection{185.2 Yamabe 问题的解决}\label{yamabe-ux95eeux9898ux7684ux89e3ux51b3}

\textbf{定理 185.1}（Yamabe 定理 / Trudinger-Aubin-Schoen 定理）：对任何紧 Riemann 流形 \((M^n, g)\)（\(n \geq 3\)），共形类 \([g]\) 中存在常标量曲率度量。

\emph{历史}：Yamabe（1960）声称证明但有错误------他在证明中未正确处理临界 Sobolev 嵌入的紧性缺失。Trudinger（1968）通过对 Yamabe 泛函的精细分析，证明了当 \(Y(g) \leq 0\) 时极小化序列的紧性成立，从而解决了此情形。Aubin（1976）引入局部化方法，证明当 \(Y(g) < Y(S^n)\)（\(Y(S^n)\) 为标准球面的 Yamabe 不变量）时，极小化泛函的紧性恢复，并完成了 \(n \geq 6\) 且 \((M,g)\) 非局部共形平坦情形的证明。剩余的临界情形------特别是 \(Y(g)=Y(S^n)\) 以及低维的情形------则需要排除''气泡''的形成。Schoen（1984）使用\textbf{正质量定理}（Positive Mass Theorem）证明了在剩余所有情形中均有 \(Y(g) < Y(S^n)\)（除非 \((M,g)\) 共形等价于 \(S^n\)），从而彻底解决了 Yamabe 问题。该问题的完整解决历时 24 年，涉及几何分析、临界 Sobolev 理论和广义相对论中正质量猜想的深刻交汇。

\textbf{定理 185.1\(^*\)}（正质量定理，Schoen-Yau 1979, Witten 1981）：设 \((M^n, g)\) 是渐近平坦 Riemann 流形（\(3 \leq n \leq 7\)），满足主导能量条件（\(R_g \geq 0\) 在主方向上）。则 ADM 质量 \(m_{\operatorname{ADM}} \geq 0\)，且 \(m_{\operatorname{ADM}} = 0\) 当且仅当 \((M, g)\) 等距于 \(\mathbb{R}^n\)。

\textbf{证明}（Schoen-Yau方法）：假设 \(m_{\text{ADM}} < 0\)。构造极小曲面 \(\Sigma \subset M\) 为面积泛函的极小化序列的极限。由第二变分公式，\(\Sigma\) 的稳定性给出 \(\int_\Sigma |A|^2 + \operatorname{Ric}_M(\nu,\nu) \leq 0\)。利用 Gauss 方程将 \(\operatorname{Ric}_M(\nu,\nu)\) 表为 \(\Sigma\) 的内蕴曲率，得 \(\int_\Sigma R_\Sigma - |A|^2 \geq 0\)。由 Gauss-Bonnet 和 Gauss 曲率积分估计，导出矛盾（当 \(n=3\) 时 \(\Sigma\) 为极小曲面，面积泛函的渐近行为与 ADM 质量符号冲突）。故 \(m_{\text{ADM}} \geq 0\)。若 \(m_{\text{ADM}} = 0\)，则 \(M\) 为 Ricci 平坦且渐近平坦，由刚性定理必为 \(\mathbb{R}^n\)。\(\blacksquare\)

\textbf{证明}：当 Yamabe 不变量 \(Y(g) < Y(S^n)\) 时，极小化序列的紧性成立（Sobolev 嵌入的临界指数困难通过 Aubin 的局部化方法克服）。Schoen 使用正质量定理证明 \(Y(g) \leq Y(S^n)\) 且等号成立当且仅当 \((M, g)\) 共形等价于 \(S^n\)。∎

\subsection{185.3 预定标量曲率}\label{ux9884ux5b9aux6807ux91cfux66f2ux7387}

\textbf{定义 185.3}（预定标量曲率问题 / Nirenberg 问题）：给定 \(S^2\) 上的函数 \(K(x)\)，是否存在 \(S^2\) 上共形于标准度量的度量，其 Gauss 曲率为 \(K(x)\)？

Yamabe 问题的解决回答了''是否存在常标量曲率度量''这一存在性问题；而 Nirenberg 问题更进一步------它要求度量上的标量曲率恰好取为给定的函数 \(K(x)\)。这是一个本质上更为困难的反问题，其可解性不仅取决于 \(K\) 的正则性，还受到深刻的拓扑障碍制约。

\textbf{定理 185.2}（Kazdan-Warner 障碍，1974）：\(S^2\) 上存在 Gauss 曲率 \(K\) 的必要条件是：对任何一阶球谐函数 \(\psi\)，\(\int_{S^2} \nabla K \cdot \nabla \psi \, dV = 0\)。这给出了 \(S^2\) 上预定 Gauss 曲率问题的拓扑障碍。

\textbf{证明}：\(S^2\) 上的共形 Gauss 曲率方程 \(\Delta u + K e^{2u} = 1\)（\(K\) 为预定曲率）。对任意一阶球谐函数 \(\psi\)（即 \(\Delta_{S^2}\psi = -2\psi\)），将方程两端乘 \(\nabla\psi\) 并在 \(S^2\) 上积分：\(\int \Delta u \nabla\psi + \int K e^{2u} \nabla\psi = \int \nabla\psi\)。分部积分并利用 \(\nabla\psi\) 与共形因子的相容性，得 \(\int \nabla K \cdot \nabla\psi \, e^{2u} dV = 0\)。因 \(e^{2u}>0\)，这等价于 \(\int \nabla K\cdot\nabla\psi \, d\sigma = 0\)（标准球面测度下），即 \(K\) 必须与所有一阶球谐函数正交。此为 \(S^2\) 上 Nirenberg 问题的著名障碍条件。\(\blacksquare\)

Kazdan-Warner 障碍条件表明，并非每个光滑正函数都可以实现为 \(S^2\) 上某个度量的 Gauss 曲率------这是一个纯分析的 PDE 问题中存在全局拓扑障碍的典型案例。在排除了这些障碍之后，能否给出充分条件保证解的存在性？Chang 和 Yang 的工作正是沿着这一方向给出了一组较为一般的充分条件。

\textbf{定理 185.3}（Chang-Yang 解，1987）：\(S^2\) 上的 Nirenberg 问题在 \(K(x)\) 为正函数且满足某些非退化条件时解决。

\textbf{证明}：Chang-Yang 方法基于 Kazdan-Warner 条件的充分性证明。当 \(K>0\) 且满足非退化条件（\(\nabla^2 K\) 在临界点处非退化），用变分方法在 H¹ 空间中构造能量泛函 \(J(u)=\frac{1}{2}\int|\nabla u|^2 + \int u - \frac{1}{2}\log\int K e^{2u}\) 的临界点。Moser-Trudinger 不等式保证 \(J\) 在约束条件下满足 Palais-Smale 条件（排除边界发散和集中紧性缺失）。通过 min-max 和 Morse 理论，临界点数目由 \(K\) 的拓扑信息（临界点 Morse 指标）决定，保证解的存在性。\(\blacksquare\)

\hrulefill

\hrulefill

\hrulefill

\hrulefill

\hrulefill

\hrulefill

\section{第267章：几何 PDE 选讲}\label{ux7b2c267ux7ae0ux51e0ux4f55-pde-ux9009ux8bb2}

本章补充介绍几何分析中与 Ricci 流和 Yamabe 问题并列的其他核心几何偏微分方程。平均曲率流（mean curvature flow，MCF）是超曲面在外法向方向以平均曲率速度演化的几何热流------其方程 \(\partial_t X = -H \mathbf{n}\) 是 Ricci 流在余维数大于 1 情形的自然类比，在曲面论的 Grayson 定理、Huisken 的凸性理论和 Brakke 的 varifold 弱解理论中产生了深刻结果。Yang-Mills 方程则源于规范场论：联络 \(A\) 的 Yang-Mills 泛函 \(\|F_A\|_{L^2}^2\) 的 Euler-Lagrange 方程 \(D_A^* F_A = 0\) 是四维流形 Donaldson 理论（第193章）的出发点，其 anti-self-dual 解（瞬子）在低维拓扑的 Floer 同调（第192章）和四维光滑结构分类中发挥了核心作用。调和映射热流则将映射间的能量泛函极小化问题转化为梯度流，在 Eells-Sampson 关于非正截面曲率存在性定理和 Sacks-Uhlenbeck 的泡泡分析中具有里程碑意义。这三类 PDE 共享的数学结构------能量泛函极小化、单调性公式和奇点形成-分析的正则化------构成了几何分析中''变分法 + 热流''范式的教科书示例。

\subsection{186.1 平均曲率流}\label{ux5e73ux5747ux66f2ux7387ux6d41}

\textbf{定义 186.1}（平均曲率流）：\(n\) 维超曲面 \(M_t \subset \mathbb{R}^{n+1}\) 的\textbf{平均曲率流}（MCF）为

\[\frac{\partial \mathbf{x}}{\partial t} = -H \mathbf{n}\]

其中 \(\mathbf{x}\) 是位置向量，\(H\) 是平均曲率，\(\mathbf{n}\) 是单位法向量。MCF 是面积泛函的梯度流：\(\frac{d}{dt} \operatorname{Area}(M_t) = -\int_{M_t} H^2 dV \leq 0\)。

\textbf{定理 186.1}（Huisken 定理，1984）：\(\mathbb{R}^{n+1}\) 中的紧凸超曲面在平均曲率流下保持凸性，并在有限时间内收缩为一点（``round point''）。经归一化后收敛到标准球面。

\textbf{证明}：设 \(M_t\subset\mathbb{R}^{n+1}\) 为平均曲率流下的紧凸超曲面。凸性保持：由第二基本形式 \(h_{ij}\) 的演化方程 \(\partial_t h_{ij} = \Delta h_{ij} + |A|^2 h_{ij} - 2H h_{ik}h_{kj}\) 和 Hamilton 的张量最大值原理，\(h_{ij}\ge 0\) 在演化中保持。当曲率趋于无穷时缩放极限为自相似收缩解。Huisken 通过单调性公式分析高斯密度，证明收缩极限为球面------因球面是唯一紧致自相似收缩凸超曲面（由 Huisken 的分类定理）。归一化后沿 \(S^n\) 球面形状收敛。\(\blacksquare\)

\textbf{定理 186.2}（Grayson 定理，1987）：\(\mathbb{R}^2\) 中任何嵌入闭曲线在曲线缩短流（平均曲率流的 1 维情形）下变为凸曲线，然后收缩为一点。

\textbf{证明}：平面曲线缩短流（CSF）：曲线 \(\gamma(s,t)\) 满足 \(\partial_t\gamma = \kappa N\)（\(\kappa\) 为曲率）。Gage-Hamilton 证明凸曲线在 CSF 下保持凸性并收缩为圆。Grayson 证明非凸嵌入闭曲线的曲率在有限时间后处处变号仅有限次，最终变为处处正------利用 Sturm 零点比较定理分析曲率零点的演化，每次零点的消失对应曲线的自交消除。一旦变为凸曲线，由 Gage-Hamilton 结果知曲线收缩为一点。\(\blacksquare\)

\textbf{定理 186.3}（Huisken-Sinestrari 手术理论，1999）：对 2-凸平均曲率流（\(\lambda_1 + \lambda_2 > 0\)），可以通过手术理论处理奇点形成。

\textbf{证明}：2-凸（\(\lambda_1+\lambda_2>0\)）平均曲率流的手术：当第二基本形式在某区域超过某阈值 \(H_{\max}\) 时，识别高曲率区域（\(\varepsilon\)-颈状或帽状）。手术步骤：切除高曲率颈部，粘贴标准解（如旋转对称收缩孤子）。关键估计：(1) 凸性估计确保 \(\lambda_1+\lambda_2>0\) 保持；(2) 梯度估计控制奇点附近的几何形状；(3) \(\varepsilon\)-正则性定理确保低于阈值区域光滑。手术仅发生有限次，每次手术后的流在拓扑上简化------最终拓扑为若干球面的连通和。\(\blacksquare\)

\subsection{186.2 Yang-Mills 方程}\label{yang-mills-ux65b9ux7a0b}

\textbf{定义 186.2}（Yang-Mills 泛函）：设 \(P \to M\) 是主 \(G\)-丛，\(A\) 是联络。\textbf{Yang-Mills 泛函}为

\[\mathcal{YM}(A) = \int_M \|F_A\|^2 \, dV\]

其中 \(F_A = dA + A \wedge A\) 是曲率 2-形式。\textbf{Yang-Mills 方程}是

\[D_A^* F_A = 0\]

（\(D_A\) 是协变外微分）。Yang-Mills 联络是 Yang-Mills 泛函的临界点。

\textbf{定理 186.4}（Uhlenbeck 的可去奇点定理，1982）：设 \(A\) 是 \(B^4 \setminus \{0\}\) 上的 Yang-Mills 联络，具有有限 Yang-Mills 作用量。则 \(A\) 可光滑延拓到整个 \(B^4\)。

\textbf{证明}：设 \(A\) 是 \(B^4\setminus\{0\}\) 上的 Yang-Mills 联络，能量 \(\int_{B^4}|F_A|^2<\infty\)。选取规范使 \(A\) 满足 Coulomb 规范条件 \(d^*A=0\)。在此规范下，Yang-Mills 方程变为椭圆型：\(\Delta A = \text{低阶项}\)。由能量有界性和椭圆正则性（Morrey 估计），\(A\) 在原点上具有可去奇点------可光滑延拓过 \(0\)。关键技术：Uhlenbeck 的 Coulomb 规范构造在小球上存在且唯一（模去紧规范群的常值作用），提供了获得正则性的合适坐标系。\(\blacksquare\)

\textbf{定理 186.5}（Donaldson 定理，1983）：单连通光滑 4 维流形上 Yang-Mills 瞬子（自对偶 Yang-Mills 联络）的模空间给出了微分同胚不变量（Donaldson 不变量），揭示了 4 维拓扑的奇异性质。

\textbf{证明}：在单连通光滑 4 维流形 \(X\) 上研究 \(\operatorname{SU}(2)\)-瞬子（反自对偶 Yang-Mills 联络）的模空间 \(\mathcal{M}_k\)。Atiyah-Hitchin-Singer 指标定理给出模空间维数 \(\dim\mathcal{M}_k=8k-3(1-b_1+b_2^+)\)。Uhlenbeck 紧致性定理给出模空间紧化，其边界由理想瞬子构成。在适当维数下，\(\mathcal{M}_k\) 的光滑部分为有限维流形，定义其基本类的杯积结构得 Donaldson 多项式不变量。这些不变量在微分同胚下不变，但可区分同胚但非微分同胚的 4 维流形，揭示 4 维微分拓扑的非凡复杂性。\(\blacksquare\)

\textbf{定理 186.6}（Uhlenbeck 紧致性定理，1982）：Yang-Mills 联络序列在 Yang-Mills 作用量有界条件下，除有限个奇点外，弱收敛到 Yang-Mills 联络（模去规范变换）。

\textbf{证明}：设 \(\{A_n\}\) 为 \(B^4\) 上 Yang-Mills 联络序列，满足 \(\int_{B^4}|F_{A_n}|^2\le C<\infty\)。选取 Coulomb 规范，由弱紧致性（\(W^{1,2}\) 中有界）得子列弱收敛到弱 Yang-Mills 联络 \(A_\infty\)。利用 \(\varepsilon\)-正则性：若某点处曲率密度小于某阈值 \(\varepsilon_0\)，则该点邻域内 \(A_\infty\) 光滑。曲率可集中在有限个孤立点------这些点处曲率密度超过阈值，其总曲率有正下界，由能量界知奇点有限。在光滑区域，\(A_n\) 强收敛到光滑联络 \(A_\infty\)。\(\blacksquare\)

\subsection{186.3 调和映射热流}\label{ux8c03ux548cux6620ux5c04ux70edux6d41}

调和映射热流是 Eells-Sampson 关于调和映射存在性理论的自然抛物化。从变分角度看，调和映射是 Dirichlet 能量的临界点，而热流则是沿此能量的最速下降方向寻找临界点的动力学过程。这一思想与 Ricci 流（面积/体积泛函的梯度流）和平均曲率流（面积泛函的梯度流）一脉相承，共同构成了几何热流理论的三大支柱。

\textbf{定义 186.3}（调和映射热流）：调和映射的热流方程为

\[\frac{\partial f_t}{\partial t} = \tau(f_t)\]

其中 \(\tau(f_t)\) 是 \(f_t\) 的张力场。这是能量泛函的梯度流。

\textbf{定理 186.7}（Struwe 定理，1985）：从紧 Riemann 面到紧 Riemann 流形的调和映射热流，在有限个奇点处形成''气泡''（bubble）后，继续光滑演化。全局弱解存在。

\textbf{证明}：调和映射热流 \(\partial_t u = \tau(u)\)（\(\tau\) 为张力场）的全局弱解存在。Struwe 构造时间离散的变分逼近：在每个时间步 \(t_{k+1}=t_k+h\) 解最小化问题 \(\min \int|\nabla u|^2 + \frac{1}{2h}\int d(u,u_k)^2\)。当气泡（能量集中）形成时，在奇点处移除气泡（对应该处调和球面 \(S^2\to N\) 的生成），在气泡移除后的区域流继续光滑演化。气泡生成仅有限次（因每次消耗最少能量 \(\varepsilon_0\)）。全局弱解在气泡处连续但不一定光滑。\(\blacksquare\)

\textbf{定理 186.8}（气泡分析，Sacks-Uhlenbeck 1981）：调和映射序列的紧性缺失由气泡现象（bubbling）描述：能量在孤立点处集中，形成非平凡的调和球面 \(S^2 \to N\)。

\textbf{证明}：设 \(\{u_n\}\) 为从 Riemann 面 \(\Sigma\) 到紧 Riemann 流形 \(N\) 的调和映射序列，能量有界。气泡分析的结论：(1) \(u_n\) 光滑收敛到光滑调和映射 \(u_\infty\)，除有限个气泡点外；(2) 在每个气泡点 \(p\) 处，重新缩放（放大）\(u_n\) 在 \(p\) 附近的限制，得到非平凡调和球面 \(S^2\to N\)（即气泡）；(3) 能量等式：\(\lim E(u_n)=E(u_\infty)+\sum \text{（各气泡的能量）}\)；(4) 气泡间形成气泡树（bubble tree）结构------连接气泡的 neck 区域能量在极限下消失。\(\blacksquare\)

\hrulefill

\hrulefill

\hrulefill

\hrulefill

\hrulefill

\newpage

\chapter{07-拓扑与几何/05-几何分析/03-谱与等周不等式.md}\label{ux62d3ux6251ux4e0eux51e0ux4f5505-ux51e0ux4f55ux5206ux679003-ux8c31ux4e0eux7b49ux5468ux4e0dux7b49ux5f0f.md}

\section{第268章：谱几何}\label{ux7b2c268ux7ae0ux8c31ux51e0ux4f55}

谱几何研究 Riemann 流形上 Laplace 算子的特征值谱与流形几何之间的深刻关系。``能否听到鼓的形状？''（Kac, 1966）是谱几何的核心问题：流形的 Laplace 谱在多大程度上决定了流形的几何（和拓扑）？

\subsection{188.1 Laplace算子与热核}\label{laplaceux7b97ux5b50ux4e0eux70edux6838}

\textbf{定义 188.1}（Laplace-Beltrami 算子）：紧 Riemann 流形 \((M, g)\) 上的 \textbf{Laplace-Beltrami 算子} \(\Delta_g = -\operatorname{div} \circ \nabla = -g^{ij} \nabla_i \nabla_j\) 是 \(L^2(M)\) 上的正定自伴椭圆算子。其谱 \(\{0 = \lambda_0 < \lambda_1 \leq \lambda_2 \leq \cdots\}\) 是离散的，每个特征值的重数有限，且特征函数构成 \(L^2(M)\) 的标准正交基。

\textbf{定义 188.2}（热核）：热方程 \(\partial_t u + \Delta_g u = 0\) 的基本解 \(K(t, x, y)\)（热核）的迹给出了谱函数： \[\operatorname{Tr}(e^{-t\Delta_g}) = \sum_{j=0}^{\infty} e^{-t\lambda_j} = \int_M K(t, x, x) dV_g\]

\textbf{定理 188.1}（热核渐近展开，Minakshisundaram-Pleijel 1949）：当 \(t \to 0^+\)，热核有以下渐近展开： \[K(t, x, x) \sim \frac{1}{(4\pi t)^{n/2}} \sum_{k=0}^{\infty} a_k(x) t^k\] 其中 \(a_k(x)\) 是局部 Riemann 不变量（曲率及其协变导数的多项式）。前几项为：\(a_0(x) = 1\)，\(a_1(x) = \frac{1}{6} R(x)\)（标量曲率），\(a_2(x) = \frac{1}{360}(12\Delta R + 5R^2 - 2|\operatorname{Ric}|^2 + 2|Rm|^2)\)。

\textbf{推论 188.1}（谱不变量）：迹 \(\operatorname{Tr}(e^{-t\Delta_g})\) 的渐近展开系数是谱不变量。积分 \(\int_M a_k(x) dV_g\) 由谱完全确定。特别地： - \(\operatorname{Vol}(M)\) 由谱确定（Weyl 渐近律） - \(\int_M R dV_g\)（总标量曲率）由谱确定（对曲面，即 Gauss-Bonnet 定理的 Euler 示性数）

\textbf{证明}：热核 \(H(t,x,y)\) 为热方程 \((\partial_t+\Delta)H=0\) 的基本解。在法坐标系下，将 Laplace 算子按测地极坐标展开：\(\Delta = \partial_r^2 + \frac{n-1}{r}\partial_r + \frac{1}{3}\operatorname{Ric}(\partial_r,\partial_r) + O(r^2)\)。热核的渐近展开通过 WKB 方法（抛物型伪微分算子演算）逐项构造：假设 \(H\sim(4\pi t)^{-n/2}e^{-d^2/4t}\sum_{k=0}^\infty t^k u_k(x,y)\)，代入热方程得关于 \(u_k\) 的递推 ODE 系统。初始项 \(u_0(x,x)=1\)，高阶项 \(u_k(x,x)\) 由曲率张量及其协变导数的局部不变量表示（非交换解析几何中的局部指标公式给出）。\(\blacksquare\)

\textbf{定理 188.2}（Weyl 渐近律，1911）：特征值的计数函数 \(N(\lambda) = \#\{j : \lambda_j \leq \lambda\}\) 满足 \[N(\lambda) \sim \frac{\omega_n \operatorname{Vol}(M)}{(2\pi)^n} \lambda^{n/2} \quad (\lambda \to \infty)\] 其中 \(\omega_n\) 是 \(\mathbb{R}^n\) 中单位球的体积。Weyl 律表明流形的体积是谱不变量，且特征值的渐近增长率仅依赖于维数和体积。

\textbf{证明}：设 \(N(\lambda)=\#\{j:\lambda_j\le\lambda\}\)。Weyl 的证明使用极小极大原理：\(\lambda_k = \min_{V_k}\max_{u\in V_k\setminus\{0\}} \frac{\|\nabla u\|^2}{\|u\|^2}\)（\(V_k\) 为 \(k\) 维子空间）。构造试探子空间：将 \(M\) 分割为直径 \(\sim\lambda^{-1/2}\) 的小区域，在每个小区域上取 Fourier 模（局部平坦近似）。由 Dirichlet-Neumann 括号定理，\(\sum_{j} e^{-t\lambda_j} = \int_M H(t,x,x) d\operatorname{vol}\)，由热核渐近展开 \(H(t,0,0)\sim(4\pi t)^{-n/2}\operatorname{vol}(M)\)，应用 Tauber 定理（Karamata）从 Laplace 变换反演得 \(N(\lambda)\sim\frac{\omega_n}{(2\pi)^n}\operatorname{vol}(M)\lambda^{n/2}\)。\(\blacksquare\)

\subsection{188.2 等谱问题与反谱问题}\label{ux7b49ux8c31ux95eeux9898ux4e0eux53cdux8c31ux95eeux9898}

\textbf{定理 188.3}（等谱非等距的流形，Gordon-Webb-Wolpert 1992）：存在一对 2 维平坦环面（16 维环面的商），它们是等谱的（所有特征值相同）但不等距------从而对 Kac 的''能否听到鼓的形状''给出了否定答案。构造使用了数论中的 Sunada 方法（群论等谱判据）。

\textbf{证明}：构造两个等谱 Riemann 流形使用 Sunada 方法：取有限群 \(G\) 及其子群 \(H_1,H_2\) 满足弱共轭条件（\(H_1\) 和 \(H_2\) 在 \(G\) 的每个共轭类中交点数目相同）。在覆盖流形 \(\tilde{M}\) 上赋予 \(G\)-不变的 Riemann 度量，商流形 \(M_i=\tilde{M}/H_i\) 具有相同的 Laplace 谱（因表示的等变性）。Gordon-Webb-Wolpert 在 16 维平坦环面上取适当格子和有限群 \(G\) 的子群对，构造了 \(M_1\not\cong M_2\) 但 \(\operatorname{Spec}(M_1)=\operatorname{Spec}(M_2)\) 的例子。此外不变量如直径（由热核小时间展开的系数部分决定）可区分二者，证明不等距。\(\blacksquare\)

Gordon-Webb-Wolpert 的构造给出了 Kac 问题的否定回答------鼓声不能唯一确定鼓的形状。但谱数据毕竟包含了大量几何信息：特征值本身与流形的曲率、体积和拓扑有着紧密的联系。Tanno 定理从相反的方向切入，探讨了特征值的极值性质能否反过来唯一确定几何。它断言，当 \(\lambda_1\) 达到 Lichnerowicz 定理所允许的下界 \(n\) 时（在 Ricci 曲率 \(\geq (n-1)\) 的条件下），流形必须是标准球面。这是一个典型的''刚性''结果------极值情形唯一确定了几何。

\textbf{定理 188.4}（Tanno 定理，1974）：球面 \(S^n\) 是唯一满足 \(\lambda_1 = n\)（即 \(\lambda_1\) 达到 Lichnerowicz-Obata 上界）的紧 Riemann 流形。特征值的极值性质可用于刻画标准几何。

Tanno 定理和 Lichnerowicz-Obata 定理共同指向一个更深层的现象：在适当的曲率条件下，谱数据（特别是 \(\lambda_1\)）与流形的等距类之间存在刚性对应。这一观察启发 Bérard、Besson 和 Gallot 提出了谱刚性猜想：热核不变量在极值情形应当唯一确定流形的等距类------换言之，当谱不变量达到极值时，流形的几何被''锁定''在标准模型上。这一猜想至今仍是谱几何的核心未解决问题。

\textbf{猜想 188.1}（Bérard-Besson-Gallot 刚性猜想）：某些谱数据（如热核迹的前几个系数和 \(\lambda_1\) 的极值性质）在适当曲率条件下可以唯一确定流形的等距类。谱刚性（spectral rigidity）是谱几何的核心研究方向。

\textbf{证明}：由 Lichnerowicz 定理，\(\operatorname{Ric}\ge (n-1)K>0\) 时 \(\lambda_1\ge nK\)。\(S^n\) 上 \(\operatorname{Ric}\equiv (n-1)\)（标准度量），实际 \(\lambda_1(S^n)=n\)。若某流形满足 \(\lambda_1=n\)，则由 Obata 定理，在 \(\operatorname{Ric}\ge(n-1)\) 条件下 \(\lambda_1=n\) 推出 \(M\) 等距于 \(S^n\)。Tanno 将此推广到更一般的曲率条件下，证明 \(\lambda_1\ge n\) 是球面的特征性质：当 \(\operatorname{Ric}\ge(n-1)\) 且 \(\lambda_1=n\) 时，梯度第一特征函数为共形向量场，由此积分恒等式推出度量必为标准球面度量。\(\blacksquare\)

\subsection{188.3 特征值估计}\label{ux7279ux5f81ux503cux4f30ux8ba1}

\textbf{定理 188.5}（Lichnerowicz 定理，1958）：设 \((M^n, g)\) 是紧 Riemann 流形，满足 \(\operatorname{Ric} \geq (n-1)K > 0\)（Ricci 曲率有正下界）。则第一个非零特征值满足 \[\lambda_1 \geq nK\] \(S^n\) 上 \(\lambda_1 = n\)（当 \(K=1\)）达到等号（Obata 1962 证明仅有球面达到）。

\textbf{证明}：对第一特征函数 \(f\)（\(\Delta f=\lambda_1 f\)），使用 Bochner 公式：\(\frac{1}{2}\Delta|\nabla f|^2 = |\operatorname{Hess} f|^2 + \langle\nabla\Delta f,\nabla f\rangle + \operatorname{Ric}(\nabla f,\nabla f)\)。代入 \(\Delta f=-\lambda_1 f\) 得 \(\frac{1}{2}\Delta|\nabla f|^2 = |\operatorname{Hess} f|^2 - \lambda_1|\nabla f|^2 + \operatorname{Ric}(\nabla f,\nabla f)\)。由 \(\operatorname{Ric}\ge(n-1)K\) 和 \(|\operatorname{Hess}f|^2\ge\frac{1}{n}(\Delta f)^2=\frac{1}{n}\lambda_1^2 f^2\)，在全流形上积分：\(0\ge\int(\frac{1}{n}\lambda_1^2 f^2 - \lambda_1|\nabla f|^2 + (n-1)K|\nabla f|^2)\)。利用 \(\int|\nabla f|^2=\lambda_1\int f^2\)，化简得 \(\lambda_1\ge nK\)。\(\blacksquare\)

Lichnerowicz 定理将 Ricci 曲率下界转化为 \(\lambda_1\) 的下界，体现了曲率对谱的控制。但曲率是一个局部几何量------它在空间每一点都有定义且可能变化。是否存在更为全局的几何量也能控制 \(\lambda_1\)？Cheeger 在 1970 年给出了肯定回答，他引入了等周常数 \(h(M)\)------它衡量的是流形中能将体积一分为二的极小曲面的面积与体积之比------并证明了 \(\lambda_1 \geq h(M)^2/4\)。这条不等式将谱理论与等周几何深刻联系在一起，是几何分析中最优美也最有用的不等式之一。

\textbf{定理 188.6}（Cheeger 不等式，1970）：设 \(h(M) = \inf_{A} \frac{\operatorname{Area}(\partial A)}{\min(\operatorname{Vol}(A), \operatorname{Vol}(M \setminus A))}\) 是 Cheeger 等周常数。则 \[\lambda_1 \geq \frac{h(M)^2}{4}\] 这个不等式将几何量（等周常数）与分析量（第一特征值）联系起来，是谱几何中最基本的不等式之一。

\textbf{证明}：Cheeger 常数 \(h(M)=\inf_A \operatorname{Area}(\partial A)/\min(\operatorname{Vol}(A),\operatorname{Vol}(M\setminus A))\)。由余面积公式，对任意光滑 \(f\)，\(\int_M|\nabla f| dV = \int_{-\infty}^\infty \operatorname{Area}(f^{-1}(t)) dt \ge h(M)\int_{-\infty}^\infty \min(\operatorname{Vol}(\{f>t\}),\operatorname{Vol}(\{f<t\})) dt\)。设 \(f_1\) 为第一 Dirichlet 特征函数（经调整使中位数为 \(0\)），结合等周不等式和 \(L^1\)-\(L^2\) 估计，推得 \(\lambda_1 \ge \frac{h(M)^2}{4}\)。该不等式将谱理论与等周几何联系起来，是高维等周不等式的核心工具。\(\blacksquare\)

Lichnerowicz 和 Cheeger 都是从下方估计 \(\lambda_1\)。反过来，从中方（或上方）估计 \(\lambda_1\) 同样重要------它告诉我们曲率和直径是否对 \(\lambda_1\) 施加了上界约束。Zhong 和 Yang（1984）证明了在正 Ricci 曲率条件下，\(\lambda_1\) 的上界完全由 Ricci 曲率下界和流形直径决定，且极值情形再次唯一刻画了标准球面。这一结果与 Lichnerowicz 下界一起，将 \(\lambda_1\) 牢牢限定在由曲率和直径所确定的紧致区间内，为谱几何的比较理论奠定了完整框架。

\textbf{定理 188.7}（Zhong-Yang 和 Hang-Wang 的特征值比较，1984）：对 \(\operatorname{Ric} \geq (n-1)K\) 的紧 \(n\) 维流形，\(\lambda_1 \leq \bar{\lambda}_1(K, d)\)，其中 \(\bar{\lambda}_1\) 是相同 Ricci 下界和直径上界的模型空间（空间形式）的 \(\lambda_1\)。这说明 \(\lambda_1\) 的上界由曲率和直径控制。

\textbf{证明}：当 \(\operatorname{Ric}\ge(n-1)K>0\) 时，Zhong-Yang 第一特征值比较：\(\lambda_1(M)\le\lambda_1(S^n_K)\)，其中 \(S^n_K\) 是常曲率 \(K\) 的球面。证明使用 Liouville 型变换将流形上的特征函数映为模型空间上的球谐函数，结合 Bishop-Gromov 体积比较定理和梯度估计。通过构造 test 函数 \(\varphi=\Phi\circ r\)（\(r\) 为距某点的距离），极小极大原理给出上界。等号成立当且仅当 \(M\) 为球面空间形式。\(\blacksquare\)

到目前为止的讨论都是围绕单个特征值（主要是 \(\lambda_1\)）展开的。但 Laplace 算子的整个谱 \(\{\lambda_j\}_{j=1}^\infty\) 包含了比任何单个特征值丰富得多的信息。如何将这些信息组织为可分析的数学对象？谱 \(\zeta\) 函数 \(\zeta_\Delta(s) = \sum \lambda_j^{-s}\) 提供了一种自然的途径------它将离散的谱序列编码为一个复变量的亚纯函数。这一构造的思想源自 Riemann \(\zeta\) 函数在解析数论中的成功，由 Minakshisundaram 和 Pleijel 在 1949 年引入谱几何。谱 \(\zeta\) 函数不仅是分析谱分布的有力工具，还通过 \(\zeta\)-正则化定义了 Laplace 算子的行列式 \(\det\Delta\)，后者在量子场论（单圈有效作用量）、弦理论和 Ray-Singer 解析挠率中发挥着核心作用。

\textbf{定理 188.8}（谱 \(\zeta\) 函数与行列式）：\textbf{谱 \(\zeta\) 函数} \(\zeta_\Delta(s) = \sum_{j=1}^\infty \lambda_j^{-s}\) 在 \(\Re(s) > n/2\) 收敛，可亚纯延拓到整个复平面（在 \(s=0\) 处正则）。\textbf{行列式} \(\det \Delta\) 由 \(\zeta_\Delta'(0)\) 定义（经 \(\zeta\) 正则化）。行列式是重要的谱不变量，与 Ray-Singer 解析挠率（analytic torsion）等价（Cheeger-Müller 定理）。

\textbf{证明}：由 Weyl 渐近律 \(\lambda_j\sim c j^{2/n}\)，级数 \(\zeta_\Delta(s)=\sum\lambda_j^{-s}\) 在 \(\Re(s)>n/2\) 绝对收敛。亚纯延拓通过热核展开获得：\(\zeta_\Delta(s)=\frac{1}{\Gamma(s)}\int_0^\infty t^{s-1}\operatorname{Tr}(e^{-t\Delta}) dt\)。将热核渐近展开代入，\(\operatorname{Tr}(e^{-t\Delta})\sim\sum_{k=0}^\infty a_k t^{k-n/2}\)，逐项积分得 \(\zeta_\Delta(s)=\frac{1}{\Gamma(s)}\sum_{k=0}^\infty\frac{a_k}{s+k-n/2}+\text{整函数}\)。\(s=0\) 处正则（\(\Gamma(s)\sim 1/s\) 可抵消极点）。泛函行列式定义为 \(\det \Delta = \exp(-\zeta_\Delta'(0))\)，物理中用于量子场论的单圈有效作用量计算。\(\blacksquare\)

\hrulefill

\hrulefill

\hrulefill

\hrulefill

\hrulefill

\hrulefill

\section{第269章：等周不等式与几何不等式}\label{ux7b2c269ux7ae0ux7b49ux5468ux4e0dux7b49ux5f0fux4e0eux51e0ux4f55ux4e0dux7b49ux5f0f}

等周不等式是几何学中最古老的不等式之一------给定固定体积（面积），最小表面积（周长）由球面（圆）达到。本章从经典等周不等式出发，推广到 Riemann 流形上的 Lévy-Gromov 等不等式和 Brunn-Minkowski 理论。

\subsection{189.1 经典等周不等式}\label{ux7ecfux5178ux7b49ux5468ux4e0dux7b49ux5f0f}

\textbf{定理 189.1}（\(\mathbb{R}^n\) 中的等周不等式）：设 \(\Omega \subset \mathbb{R}^n\) 是有界域，具有光滑边界（或有限周长）。则 \[\frac{\operatorname{Vol}(\partial \Omega)}{\operatorname{Vol}(\Omega)^{(n-1)/n}} \geq \frac{\operatorname{Vol}(S^{n-1})}{\operatorname{Vol}(B^n)^{(n-1)/n}} = n \omega_n^{1/n}\] 其中 \(\omega_n\) 是 \(\mathbb{R}^n\) 中单位球的体积，等号成立当且仅当 \(\Omega\) 是球。

\textbf{证明}（Brunn-Minkowski 不等式途径）：设 \(A = \Omega\)，\(B = B_r(0)\) 为半径 \(r\) 的球。由 Brunn-Minkowski 不等式，\(\operatorname{Vol}(A+B)^{1/n} \geq \operatorname{Vol}(A)^{1/n} + \operatorname{Vol}(B)^{1/n}\)。Minkowski 和 \(A+B = \{x + y : x \in A, y \in B\}\) 的体积有展开：\(\operatorname{Vol}(A+B) = \operatorname{Vol}(A) + r \operatorname{Vol}(\partial A) + o(r)\) 当 \(r \to 0\)。代入得 \(\operatorname{Vol}(\partial \Omega) \geq n \omega_n^{1/n} \operatorname{Vol}(\Omega)^{(n-1)/n}\)。等号成立当且仅当 \(\Omega\) 与球在相差平移和缩放下相同（由 Brunn-Minkowski 不等式的等号条件）。\(\blacksquare\)

等周不等式不仅是几何学中最古老的不等式，也与分析学中最重要的函数空间嵌入------Sobolev 不等式------有着深刻的等价关系。实际上，等周不等式本质上就是 \(BV(\mathbb{R}^n)\)（有界变差函数空间）嵌入 \(L^{n/(n-1)}(\mathbb{R}^n)\) 的 Sobolev 型不等式，其中的表面积对应 \(L^1\) 梯度范数。这一等价性通过余面积公式（Coarea Formula）得以精确建立，它将函数水平集的几何与函数的分析范数联系起来，是将几何不等式与分析不等式双向转化的核心工具。

\textbf{推论 189.1}（等周不等式的最优常数与 Sobolev 不等式）：等周不等式等价于嵌入 \(BV(\mathbb{R}^n) \hookrightarrow L^{n/(n-1)}(\mathbb{R}^n)\) 的 Sobolev 不等式（等周型 Sobolev 不等式）。事实上，面积与体积的极小化问题通过余面积公式与函数的 \(L^1\) 梯度范数关联。

\subsection{189.2 曲率与等周不等式}\label{ux66f2ux7387ux4e0eux7b49ux5468ux4e0dux7b49ux5f0f}

曲率与等周不等式之间的关系是 Riemann 几何中最深刻的结构定理之一，它将局部几何量（Ricci 曲率）与全局几何量（等周剖面）联系起来。Gromov 的 Lévy-Gromov 等周不等式（1980）是这一方向的核心结果：在 \(\operatorname{Ric} \geq (n-1)K > 0\) 的紧流形上，等周剖面被具有相同 Ricci 下界的球面 \(S^n\) 的等周剖面从上方控制。这一不等式的物理直觉是：正 Ricci 曲率意味着流形中的测地球比平坦空间中的球增长得更慢（Bishop-Gromov 体积比较定理），而等周不等式本质上是在比较区域边界面积与区域体积的比率------正 Ricci 曲率使得边界面积相对于体积更大，从而等周常数更大。Gromov 的证明使用了 Almgren 的变分技巧和等周测度比较：定义等周剖面 \(I_M(v) = \inf\{\operatorname{Area}(\partial \Omega) : \operatorname{Vol}(\Omega) = v\}\)，然后对 \(I_M(v)\) 满足的微分不等式进行分析，利用 Bishop-Gromov 不等式将体积比和面积比与模型空间（球面空间形式）进行逐点比较。Lévy-Gromov 不等式的一个直接推论是：若 \(\operatorname{Ric} \geq (n-1)K > 0\)，则 Cheeger 等周常数 \(h(M) \geq h(S^n_K)\)（\(S^n_K\) 为常曲率 \(K\) 的球面），从而由 Cheeger 不等式 \(\lambda_1 \geq h(M)^2/4\) 得到 Lichnerowicz 定理的另一种证明------\(\lambda_1 \geq nK\)。这一推理链（Ricci 下界 \(\Rightarrow\) 体积比较 \(\Rightarrow\) 等周不等式 \(\Rightarrow\) Cheeger 不等式 \(\Rightarrow\) \(\lambda_1\) 下界）展示了谱几何中曲率、等周和谱三者之间的深层联系。在更一般的框架中，Colding-Minicozzi 的等周比较定理和 Gallot 的等周谱比较定理将 Lévy-Gromov 不等式推广到 Ricci 曲率有任意下界（不限于正曲率）的情形，建立了等周剖面与 Ricci 曲率下界之间的精确对应关系。

\textbf{定理 189.2}（Lévy-Gromov 等周不等式，Gromov 1980（推广 Lévy 1951 的结果））：设 \((M^n, g)\) 是紧 Riemann 流形，\(\operatorname{Ric} \geq (n-1)K > 0\)。则 \(M\) 满足与 Ricci 曲率为 \((n-1)K\) 的球面 \(S^n\) 比较的等周不等式： \[\frac{\operatorname{Area}(\partial \Omega)}{\operatorname{Vol}(M)} \geq \frac{\operatorname{Area}(\partial B)}{\operatorname{Vol}(S^n)}\] 其中 \(B \subset S^n\) 是具有相同归一化体积的球冠，\(\Omega \subset M\) 是充分光滑的区域。即正 Ricci 曲率下界使得流形的等周剖面被球面控制。

\textbf{证明}：核心工具是 Gromov 的等周测度比较。对任意 Borel 集 \(A \subset M\)，Minkowski 边界面积 \(\operatorname{Area}^+(\partial A) = \lim_{\epsilon \to 0} \frac{\operatorname{Vol}(A_\epsilon) - \operatorname{Vol}(A)}{\epsilon}\) 满足关于 Bishop-Gromov 体积比较的微分不等式。将体积比转化为模型空间的对应量即得等周不等式的比较形式。

\subsection{189.3 Brunn-Minkowski 不等式}\label{brunn-minkowski-ux4e0dux7b49ux5f0f}

\textbf{定理 189.3}（Brunn-Minkowski 不等式）：设 \(A, B \subset \mathbb{R}^n\) 是紧集。则 \[\operatorname{Vol}(A + B)^{1/n} \geq \operatorname{Vol}(A)^{1/n} + \operatorname{Vol}(B)^{1/n}\] 其中 \(A + B = \{a + b : a \in A, b \in B\}\) 是 Minkowski 和。等号成立当且仅当 \(A\) 和 \(B\) 是相似凸集（或 \(A, B\) 之一是点）。

\textbf{证明}：当 \(A, B\) 是坐标方体（boxes）时平凡成立。对一般情形，使用压缩映射 \(\Phi_t(x) = (1-t)x_0 + t x\) 将 \(A\) 收缩并结合几何测度论中 Brunn 的凹性原理。等号情形由等周不等式的最优性分析得到。\(\blacksquare\)

Brunn-Minkowski 不等式的威力在于，经典等周不等式可以作为它的极限情形推导出来。取 \(B\) 为半径 \(\varepsilon\) 的小球，则 Minkowski 和 \(A + \varepsilon B\) 的体积展开式中的一阶项恰好涉及 \(\partial A\) 的表面积，令 \(\varepsilon \to 0\) 即得等周不等式。这表明 Brunn-Minkowski 不等式在某种意义上更为基本------它是等周不平式的''无穷小''前身，同时比等周不等式更具普适性，因为它对任意两个紧集（无需光滑边界）都成立。

\textbf{推论 189.2}（等周不等式作为 Brunn-Minkowski 的极限）：取 \(B = \epsilon B^n\)（半径 \(\epsilon\) 的小球），则 \(\operatorname{Vol}(A + \epsilon B^n)^{1/n} \gtrsim \operatorname{Vol}(A)^{1/n} + \epsilon \omega_n^{1/n}\)，令 \(\epsilon \to 0\) 即得等周不等式（Minkowski 边界面积公式）。

Brunn-Minkowski 不等式处理的是集合之间的几何加法（Minkowski 和），但通过''层垒表示''（layer cake representation），可以将集合的不等式转化为函数的不等式。Prékopa 和 Leindler 在 1970 年代独立发现了这一函数形式的推广：给定三个函数 \(f, g, h\)，若它们在凸组合点上的值满足对数凹性条件 \(h((1-t)x+ty) \geq f(x)^{1-t} g(y)^t\)，则它们的积分也满足相应的不等式。这一推广不仅蕴含了 Brunn-Minkowski 不等式（取 \(f, g, h\) 为集合的示性函数），还在泛函分析中引发了一系列深刻结果：Borell-Brascamp-Lieb 浓度不等式和对数 Sobolev 不等式均源于此。

\textbf{定理 189.4}（Prékopa-Leindler 不等式，函数形式的 Brunn-Minkowski）：设 \(f, g, h: \mathbb{R}^n \to [0, \infty)\) 可测，满足 \(h((1-t)x + ty) \geq f(x)^{1-t} g(y)^t\) 对所有 \(x, y\)。则 \[\int_{\mathbb{R}^n} h \geq \left(\int_{\mathbb{R}^n} f\right)^{1-t} \left(\int_{\mathbb{R}^n} g\right)^t\] 此函数形式蕴含 Brunn-Minkowski 不等式，且是证明 Borell-Brascamp-Lieb 浓度不等式的基本工具。

\textbf{证明}：Prékopa-Leindler 不等式的证明基于 Brunn-Minkowski 不等式的张量积分表示。核心步骤：由假设 \(h((1-t)x+ty)\ge f(x)^{1-t}g(y)^t\)，令 \(F(s)=\operatorname{Vol}(\{f> s\})\) 和 \(G(s)=\operatorname{Vol}(\{g> s\})\)。对 \(h\) 的水平集 \(\{h>s\}\) 应用几何 Brunn-Minkowski 不等式：\(\{h>s\}\supseteq (1-t)\{f>s_1\} + t\{g>s_2\}\)，选择 \(s_1,s_2\) 使积分优化后得 \(\int h \ge (\int f)^{1-t}(\int g)^t\)（当函数的「和」满足一致下界时）。这一不等式是广义等周不等式的基础，也是 Ledoux、Bobkov 等人证明对数 Sobolev 不等式和传输不等式的核心。\(\blacksquare\)

\subsection{189.4 等周不等式在几何分析中的应用}\label{ux7b49ux5468ux4e0dux7b49ux5f0fux5728ux51e0ux4f55ux5206ux6790ux4e2dux7684ux5e94ux7528}

\textbf{定理 189.5}（Bérard-Meyer 特征值比较，1982）：结合 Cheeger 不等式和 Lévy-Gromov 等周不等式，得到带 Ricci 曲率下界的流形 \(\lambda_1\) 的显式下界估计。这推广了 Lichnerowicz 估计------允许通过等周剖面获得更精细的特征值信息。

\textbf{证明}：在 \(\operatorname{Ric}\ge(n-1)K>0\) 条件下，Bérard-Meyer 发展了等周剖面的热核估计。等周剖面 \(I_M(v)=\inf\{\operatorname{Area}(\partial\Omega):\operatorname{Vol}(\Omega)=v\}\) 满足 \(I_M(v)\ge I_{S^n_K}(v)\)（Lévy-Gromov）。Cheeger 不等式提供 \(\lambda_1\ge \frac{1}{4}\inf_v(I_M(v)/v)^2\)。结合模型空间 \(S^n_K\) 的显式等周剖面和 Cheeger 常数的精细化，得 \(\lambda_1\ge f(n,K,\operatorname{diam}(M))\)，其中 \(f\) 为仅依赖于维数、曲率下界和直径的显式函数。这恢复了 Lichnerowicz 的 \(\lambda_1\ge nK\) 对短直径流形的最优估计，并定量地表现了长直径流形的退化行为。\(\blacksquare\)

将 Cheeger 不等式与 Lévy-Gromov 等周不等式结合，可以给出 \(\lambda_1\) 的显式下界估计。但等周不等式与谱理论之间的联系远不止于 Cheeger 原理。在非紧完备流形上（\(\operatorname{Ric} \geq 0\)），等周不等式与 \(L^1\)-Poincaré 不等式构成了完全的等价关系------这意味着控制''面积/体积''比率的几何条件与控制''函数振荡/梯度''比率的分析条件在本质上是一回事。这一等价性是谱几何与等周几何之间最深刻的结构性定理之一，将前几节中零散出现的联系系统化为统一的理论框架。

\textbf{定理 189.6}（等周不等式与 Poincaré 不等式）：在 \(\operatorname{Ric} \geq 0\) 的完备非紧流形上，等周不等式等价于 \(L^1\)-Poincaré 不等式 \(\int |f - \bar{f}_B| \leq C r \int |\nabla f|\)。这建立了等周剖面与流形的 Sobolev 不等式的等价关系。

\textbf{证明}：在 \(\operatorname{Ric}\ge 0\) 的完备非紧流形上，\(L^1\)-等周不等式 \(\operatorname{Area}(\partial\Omega)\ge C\operatorname{Vol}(\Omega)^{(n-1)/n}\) 等价于 \(L^1\)-Poincaré 不等式 \(\int_{B_R}|f-\bar{f}_{B_R}| \le C'R\int_{B_R}|\nabla f|\)。由 Cheeger 的原理，等周常数和 Poincaré 常数之间满足 \(C_{\text{Poinc}}\le 2/C_{\text{Iso}}\)。反向不等式使用等周区域的构造：对特征函数的水平集应用等周不等式，由 coarea 公式给出积分估计。等价性的核心是 Cheeger 的等周-Poincaré 二重性原理在非紧完备流形上的推广。\(\blacksquare\)

等周不等式与 Poincaré 不等式的等价性属于几何分析中的''软''结果------它将不同量之间的定性关系建立了起来。但等周不等式还能产生更令人惊讶的''硬''结果：集中的测度现象（concentration of measure）。Milman 在 1988 年发现，Lévy-Gromov 等周不等式直接蕴含了正 Ricci 曲率流形上 Lipschitz 函数的 Gaussian 型浓度不等式------这与高维球面上著名的 Lévy 族集中现象完全一致。这一结果的意义在于，它将等周不等式从几何不等式的范畴扩展到了概率论和泛函分析的领域，成为连接几何、测度和函数空间不等式的又一座桥梁。

\textbf{定理 189.7}（Milman 的推论，1988）：Lévy-Gromov 等周不等式隐含紧 Riemann 流形（\(\operatorname{Ric} \geq (n-1)K > 0\)）的 \(1\)-Lipschitz 函数满足 Gaussian 型浓度不等式（类似于球面）。这给出了球面上 Lévy 族的测度集中现象的 Riemann 推广。

\textbf{证明}：Lévy-Gromov 等周不等式：\(\operatorname{Ric}\ge(n-1)K>0\) 时 \(M\) 的等周剖面 \(I_M(v)\ge I_{S^n_K}(v)\)（\(S^n_K\) 为半径 \(1/\sqrt{K}\) 的球面）。Milman 将此与浓度现象联系：对任意 1-Lipschitz 函数 \(f\)（\(|f(x)-f(y)|\le d(x,y)\)），设 \(M_f\) 为中位数。取 \(A=\{f\le M_f\}\)，则 \(\operatorname{Vol}(A)\ge\frac{1}{2}\operatorname{Vol}(M)\)。等周不等式给出 \(\operatorname{Vol}(\{f>M_f+t\})\le \frac{1}{2}e^{-c_n K t^2}\operatorname{Vol}(M)\)。这就是 Gaussian 型浓度不等式，证明了对高维球面上 Lipschitz 函数的极值集中现象。\(\blacksquare\)

\hrulefill

\hrulefill

\hrulefill

\hrulefill

\hrulefill

\hrulefill

\newpage

\chapter{07-拓扑与几何/05-几何分析/05-分形与Plateau.md}\label{ux62d3ux6251ux4e0eux51e0ux4f5505-ux51e0ux4f55ux5206ux679005-ux5206ux5f62ux4e0eplateau.md}

\section{第270章：Plateau 问题与极小曲面}\label{ux7b2c270ux7ae0plateau-ux95eeux9898ux4e0eux6781ux5c0fux66f2ux9762}

Plateau 问题（求以给定闭曲线为边界的极小面积曲面）是变分法的古典问题。几何测度论用电流理论给出了第一个完全严格的解。

\subsection{250.1 经典 Plateau 问题}\label{ux7ecfux5178-plateau-ux95eeux9898}

\textbf{定义 250.1}（Plateau 问题 / 几何测度论形式）：给定 \(\mathbb{R}^3\) 中的（可整流）闭曲线 \(\Gamma\)（一维整系数整流链满足 \(\partial \Gamma = 0\)）。求二维整系数整流链 \(T\) 满足： - \(\partial T = \Gamma\)（\(T\) 以 \(\Gamma\) 为边界） - \(\mathbf{M}(T) \leq \mathbf{M}(S)\)（对所有满足 \(\partial S = \Gamma\) 的 \(S\)）

\textbf{定理 250.1}（Douglas-Radó 解法，1931/1930）：任何 \(\mathbb{R}^3\) 中的可整流 Jordan 曲线都张成某个极小面积曲面（参数化圆盘映射）。Douglas 因此获得首届 Fields 奖（1936）。

\textbf{证明}（Douglas 的能量极小化方法）：设 \(\Gamma\) 为 \(\mathbb{R}^3\) 中的 Jordan 曲线。考虑所有从单位圆盘 \(D\) 到 \(\mathbb{R}^3\) 的 Sobolev 映射 \(u \in W^{1,2}(D, \mathbb{R}^3)\)，满足 \(u|_{\partial D}\) 为 \(\partial D\) 到 \(\Gamma\) 的单调参数化。定义 Dirichlet 能量 \(E(u) = \frac{1}{2} \int_D |\nabla u|^2\)。由面积-能量不等式，\(A(u) \leq E(u)\)，等号当且仅当 \(u\) 为共形映射。取能量极小化序列 \(\{u_n\}\)，由 Courant-Lebesgue 引理，\(u_n\) 在 \(\partial D\) 上等度连续。由 Banach-Alaoglu 定理和 Rellich-Kondrachov 紧嵌入，存在子列弱收敛到 \(u \in W^{1,2}(D, \mathbb{R}^3)\)，且 \(u|_{\partial D}\) 为 \(\Gamma\) 的连续参数化。由 Dirichlet 能量的弱下半连续性，\(E(u) \leq \liminf E(u_n)\)，故 \(u\) 为能量极小元。由 Courant 的边界条件稳定性，\(u\) 的参数化是单调的。最后，由共形不变性，可调整 \(u\) 使其为共形映射，此时 \(A(u) = E(u)\) 且 \(u\) 为极小曲面。\(\blacksquare\)

\textbf{定理 250.2}（Federer-Fleming 解法 / 几何测度论中的 Plateau 问题解，1960）：对任何一维整系数整流链 \(\Gamma\) 满足 \(\partial \Gamma = 0\) 且紧支撑，存在二维整系数整流链 \(T\) 实现极小质量。

\textbf{证明}：取质量最小化序列 \(\{T_i\}\)。用等周不等式控制边界，用紧致性定理（Federer-Fleming 紧致性定理）提取子列收敛到整系数平坦链。由下半连续性（质量在平坦范数下是下半连续的），极限链是最小的。用正则性理论证明极限是整系数整流链。∎

\subsection{250.2 极小曲面的正则性}\label{ux6781ux5c0fux66f2ux9762ux7684ux6b63ux5219ux6027}

Douglas 和 Radó 的存在性定理保证了极小曲面作为某种变分问题的解总是存在的，但它们留下了一个关键问题：这些解在几何上到底有多''光滑''？De Giorgi 在 1961 年给出的正则性定理是几何测度论早期最辉煌的成就之一------它证明了质量最小化的整系数整流链在低余维情形中几乎处处光滑，奇点至多是孤立点。

\textbf{定理 250.3}（De Giorgi 正则性定理，1961）：质量最小化的二维整系数整流链 \(T\) 在 \(\mathbb{R}^n\) 中除了孤立奇点外是光滑解析曲面。在三维中，\(T\) 是嵌入曲面（无奇点）。

\textbf{证明}：设 \(T\) 为 \(\mathbb{R}^n\) 中质量最小的 2 维整系数整流链。由 Morrey 的 Dirichlet 能量增长估计和 epiperimetric 不等式，\(T\) 的支撑在 \(\mathcal{H}^2\)-几乎处处具有 \(C^{1,\alpha}\) 正则性。De Giorgi 证明在 3 维中奇点完全消失：因 2 维最小曲面在 3 维中仅容许可去孤立奇点（由 Bernstein 型定理保证），而最小性迫使这些点正则。高维中奇集为孤立离散点，由 Schoen-Simon-Yau 的维数约化论证给出。\(\blacksquare\)

De Giorgi 的正则性理论聚焦于二维整流链（余维数为 1），而 Allard 在 1972 年将正则性理论推广到任意维数和余维数的情形，建立了以单调解（monotonicity）公式和 \(\varepsilon\)-正则性为核心的普适框架。这一推广使得极小曲面的正则性理论能够处理高维和高余维的几何变分问题。

\textbf{定理 250.4}（Allard 正则性定理，1972）：质量最小的整系数 \(k\) 维整流链在 \(\mathbb{R}^n\) 中具有解析的 Hausdorff 维数 \(k\) 的支撑，除一个 \(\mathcal{H}^{k-2}\) 维的闭奇异集外。

\textbf{定义 250.2}（奇异集的维数）：\(k\) 维面积最小化整系数整流链的奇异集 \(\operatorname{Sing}(T)\) 满足 \(\dim_H \operatorname{Sing}(T) \leq k - 7\)（对超曲面 \(k = n-1\)）。特别是，在 \(\mathbb{R}^8\) 以下维数中无奇异点。

\textbf{证明}：设 \(T\) 为 \(k\) 维质量最小整系数整流链。Allard 的单调解（monotonicity）公式：质量比 \(r^{-k}\mathbf{M}(T\llcorner B_r(x))\) 关于 \(r\) 单调非减。该公式将切锥的存在性（Blow-up 极限）与 \(T\) 的局部行为关联。Allard 的 \(\varepsilon\)-正则性定理：若某点处质量密度充分接近 \(1\)（或某平面密度），则该点正则。由此奇集为闭集且 Hausdorff 维数至多 \(k-2\)（如 \(k=2\) 时奇集为 0 维）。解析性在正则点由标准的椭圆局部理论得出。\(\blacksquare\)

\subsection{250.3 自由边界问题与毛细曲面}\label{ux81eaux7531ux8fb9ux754cux95eeux9898ux4e0eux6bdbux7ec6ux66f2ux9762}

经典 Plateau 问题的边界是事先固定的闭曲线。然而在物理应用中（如液体与容器壁的接触），极小曲面的边界可以沿一个给定的支撑曲面''自由滑动''。这种自由边界条件在数学上引入了额外的复杂性------极小曲面与支撑曲面必须以垂直角度相交------但同时也催生了丰富的几何正则性理论。

\textbf{定义 250.3}（自由边界 Plateau 问题）：极小曲面不仅张成给定曲线，还要与某固定曲面 \(W\)（如容器壁）以自由边界条件相交。极小曲面与 \(W\) 沿边界以垂直角度相交。

\textbf{定理 250.5}（自由边界问题的存在性与正则性）：在适当的条件下，自由边界 Plateau 问题有质量最小化解，其在内部光滑，在自由边界处（除特定维数的奇异集）光滑。

\textbf{定义 250.4}（毛细曲面与 Gauss 自由能）：毛细曲面满足均值曲率与高度成比例的方程（Young-Laplace 方程）。Gauss 自由能泛函结合了面积和重力势能。

\textbf{证明}：自由边界 Plateau 问题：在给定边界 \(\Gamma\) 和自由曲面 \(\Sigma\)（沿某约束流形 \(N\) 自由滑动）条件下寻找面积最小曲面。存在性由紧致性（整流链在质量有界条件下的闭性）和下半连续性（质量泛函的 \(\Gamma\)-收敛）保证。内部正则性由 Allard 定理给出。自由边界处的正则性需额外的反射原理：沿 \(N\) 反射将自由边界曲面偶延拓为无约束最小曲面，应用内部正则性得边界光滑性。奇集分析类似固定边界情形但包含边界奇点（维数至多 \(k-2\)）。\(\blacksquare\)

\hrulefill

\hrulefill

\hrulefill

\hrulefill

\hrulefill

\hrulefill

\section{第271章：变分法的几何测度论}\label{ux7b2c271ux7ae0ux53d8ux5206ux6cd5ux7684ux51e0ux4f55ux6d4bux5ea6ux8bba}

几何测度论为变分法中的许多经典问题提供了严格的数学基础，包括集合的边界变分（Caccioppoli 集）和具有奇异性的最小化子。

\subsection{251.1 Caccioppoli 集与有限周长集}\label{caccioppoli-ux96c6ux4e0eux6709ux9650ux5468ux957fux96c6}

经典变分法中极小曲面由光滑参数化定义，但在处理奇点和一般的存在性问题时，光滑范畴的局限性变得明显。De Giorgi 引入的 Caccioppoli 集（即有限周长集）是几何测度论对变分法最根本的贡献之一：它将集合的''边界面积''概念推广到仅需分布导数有意义的一般 Borel 集，而不要求边界的任何光滑性。

\textbf{定义 251.1}（有限周长集 / Caccioppoli 集）：Borel 集 \(E \subset \mathbb{R}^n\) 是\textbf{有限周长集}（Caccioppoli 集），如果其特征函数 \(\chi_E\) 是有界变差函数（\(BV\)），即分布导数 \(D\chi_E\) 是有限 Radon 测度。周长定义为 \(\operatorname{Per}(E) = \int |D\chi_E|\)。

\textbf{定理 251.1}（De Giorgi 结构定理）：有限周长集 \(E\) 的约化边界 \(\partial^* E\) 是 \(\mathcal{H}^{n-1}\)-可整流的，且 \(D\chi_E = \nu_E \mathcal{H}^{n-1} \llcorner \partial^* E\)（\(\nu_E\) 是广义外法向量）。即周长等于约化边界的 Hausdorff 面积。

\textbf{证明}：对有限周长集 \(E\)（\(\chi_E\in BV(\mathbb{R}^n)\)），分布导数 \(D\chi_E\) 为 Radon 测度。约化边界 \(\partial^*E\) 定义为具有外法向量的点集（即存在单位向量 \(\nu_E(x)\) 使测度密度满足 \(\lim_{r\to 0} D\chi_E(B_r(x))/|D\chi_E|(B_r(x))=\nu_E(x)\)）。De Giorgi 证明 \(\partial^*E\) 的 \(\mathcal{H}^{n-1}\)-可整流性由 blow-up 分析给出：在每个 \(x\in\partial^*E\) 处 blow-up 极限为半空间（超平面），这等价于 \(\partial^*E\) 具有 \(\mathcal{H}^{n-1}\)-a.e. 的 \(n-1\) 维切空间。Gauss-Green 公式 \(\int_E \operatorname{div}\varphi = \int_{\partial^*E} \varphi\cdot\nu_E d\mathcal{H}^{n-1}\) 成立。\(\blacksquare\)

\textbf{定理 251.2}（等周不等式的几何测度论证明）：\(\operatorname{Per}(E) \geq n \omega_n^{1/n} (\mathcal{L}^n(E))^{1 - 1/n}\)，等号成立当且仅当 \(E\) 是球（模去零测集）。

\textbf{证明}：对有限周长集 \(E\)，周长 \(\operatorname{Per}(E)=\mathcal{H}^{n-1}(\partial^*E)\)。余面积公式用于距离函数：设 \(u(x)=d(x,E^c)\)（或光滑化 \(u_\varepsilon\approx\chi_E\)）。对 \(BV\) 函数的 Sobolev 不等式（Gagliardo-Nirenberg）：\(\|u\|_{L^{n/(n-1)}}\le c_n\|Du\|(\mathbb{R}^n)\)。取 \(u=\chi_E\)，\(\|\chi_E\|_{n/(n-1)}=|E|^{(n-1)/n}\)，\(\|D\chi_E\|=\operatorname{Per}(E)\)，得 \(|E|^{(n-1)/n}\le c_n\operatorname{Per}(E)\)。\(c_n=n^{-1}\omega_n^{-1/n}\) 由容斥球论证和 Brunn-Minkowski 不等式推导。等号仅当 \(E\) 为球（由等周极值刚性定理）。\(\blacksquare\)

\subsection{\texorpdfstring{251.2 \(BV\) 函数与总变差}{251.2 BV 函数与总变差}}\label{bv-ux51fdux6570ux4e0eux603bux53d8ux5dee}

Caccioppoli 集的理论建立在特征函数上，而 \(BV\) 函数空间将这一框架从集合推广到一般函数。\(BV\)（有界变差）空间是 Sobolev 空间 \(W^{1,1}\) 的最自然的推广------它允许函数具有跳跃间断，同时保留了足以进行变分分析的正则性结构（如紧嵌入和下半连续性）。

\textbf{定义 251.2}（有界变差函数 / \(BV\) 函数）：\(u \in L^1(\Omega)\) 属于 \(BV(\Omega)\)，如果其分布导数 \(Du\) 是有限 Radon 测度。总变差为 \(|Du|(\Omega)\)。\(BV\) 空间是 Sobolev 空间 \(W^{1,1}\) 的自然推广（允许跳跃间断）。

\textbf{定理 251.3}（\(BV\) 函数的紧致性 / \(BV\) 紧嵌入定理）：若有界开集 \(\Omega\) 具有 Lipschitz 边界且 \(\{u_k\} \subset BV(\Omega)\) 满足 \(\sup_k \|u_k\|_{BV} < \infty\)，则存在子列在 \(L^1(\Omega)\) 中强收敛到 \(u \in BV(\Omega)\)。

\textbf{证明}：\(BV(\Omega)\) 紧嵌入 \(L^1(\Omega)\)（\(\Omega\) 为 Lipschitz 有界域）。证明：(1) \(BV\) 有界序列 \(\|u_k\|_{BV}\le C\) 蕴含 \(\|u_k\|_{L^1}\) 有界和 \(\|Du_k\|(\Omega)\) 有界。(2) 光滑逼近：用磨光核 \(\rho_\varepsilon\) 卷积得光滑函数 \(u_k^\varepsilon\)，其在 \(W^{1,1}\) 中有界。(3) Rellich-Kondrachov：\(W^{1,1}\hookrightarrow L^1\) 紧，得子列在 \(L^1\) 中收敛。(4) 对角论证：取 \(\varepsilon\to 0\) 得原序列子列在 \(L^1\) 中收敛。极限函数属于 \(BV\) 由下半连续性 \(\|D u_\infty\|\le\liminf\|D u_{k_j}\|\) 保证。\(\blacksquare\)

\textbf{定理 251.4}（\(BV\) 函数的精细结构 / Federer-Vol'pert 定理）：\(u \in BV(\Omega)\) 关于 \(\mathcal{H}^{n-1}\)-可整流的跳跃集 \(J_u\) 具有近似极限和跳跃。\(Du\) 可分解为绝对连续部分、跳跃部分和 Cantor 部分。

\textbf{证明}：\(u \in BV(\Omega)\) 的分布导数 \(Du\) 是有限 Radon 测度。由 Radon-Nikodym 定理，\(Du\) 可唯一分解为：\(Du = D^a u + D^j u + D^c u\)。绝对连续部分 \(D^a u = \nabla u \, d\mathcal{L}^n\)，其中 \(\nabla u\) 是 \(u\) 的近似梯度（由 \(L^1\) 密度的 Lebesgue 微分定理给出）。跳跃部分 \(D^j u = (u^+ - u^-) \nu_u \mathcal{H}^{n-1} \llcorner J_u\)，其中 \(J_u\) 是跳跃集，\(u^\pm(x)\) 是沿法向 \(\nu_u(x)\) 的近似极限。Federer 证明了 \(J_u\) 是 \(\mathcal{H}^{n-1}\)-可整流的，这通过 blow-up 分析：在 \(\mathcal{H}^{n-1}\)-几乎每个 \(x \in J_u\) 处，缩放极限 \(u_{x,r}(y) = u(x+ry)\) 当 \(r \to 0\) 时收敛到 Heaviside 阶梯函数。Cantor 部分 \(D^c u\) 与 \(\mathcal{L}^n\) 和 \(\mathcal{H}^{n-1}\) 均奇异，无原子部分，如 Cantor-Vitali 函数的导数。关键工具包括 Besicovitch 覆盖引理和 Blaschke 选择定理。\(\blacksquare\)

\subsection{251.3 图像处理中的变分方法}\label{ux56feux50cfux5904ux7406ux4e2dux7684ux53d8ux5206ux65b9ux6cd5}

\textbf{定义 251.3}（Mumford-Shah 泛函 / 图像分割，1989）：

\[E(u, K) = \int_{\Omega \setminus K} |\nabla u|^2 dx + \alpha \mathcal{H}^{n-1}(K) + \beta \int_{\Omega} (u - g)^2 dx\]

其中 \(u\) 是分段光滑的分割，\(K\) 是边界跳跃集，\(g\) 是观测到的噪声图像。

\textbf{定义 251.4}（Rudin-Osher-Fatemi 模型 / 全变差去噪，1992）：

\[\min_u \left\{ \int_{\Omega} |Du| + \frac{\lambda}{2} \int_{\Omega} (u - g)^2 dx \right\}\]

全变差正则化保持图像边缘（不模糊跳跃），是现代图像处理的基石。

\textbf{定理 251.5}（ROF 模型的解的性质）：ROF 模型的解 \(u\) 是 \(g\) 的「全变差流」在尺度 \(\lambda^{-1}\) 下的结果。大尺度结构被保留，小尺度振荡被滤除。

\textbf{证明}：ROF 模型 \(\min_u \|u-g\|_{L^2}^2 + \lambda\|Du\|(\Omega)\) 的解 \(u\) 是 \(g\) 的全变差流在尺度 \(\lambda^{-1}\) 下的演化结果。由最优性条件（Euler-Lagrange 包含关系）：\(u-g\in\lambda\partial\|Du\|\)（次微分）。当 \(\lambda\) 增大时，\(u\) 趋向于 \(g\) 的局部中位数（或几何中值）。在区域内部 \(u\) 为分段常数（staircasing 效应）------这源于全变差正则化项惩罚振荡但不惩罚常数平移------大尺度特征（如边缘和纹理区域）被保留，而尺度小于 \(\lambda\) 的振荡（如噪声、纹理）被滤除。该性质使 ROF 模型成为图像去噪的基准方法。\(\blacksquare\)

\subsection{251.4 集中紧致性原理与 Gamma 收敛}\label{ux96c6ux4e2dux7d27ux81f4ux6027ux539fux7406ux4e0e-gamma-ux6536ux655b}

\textbf{定义 251.5}（Gamma 收敛 / \(\Gamma\)-收敛，De Giorgi 1970s）：泛函序列 \(F_n\) \textbf{\(\Gamma\)-收敛}到 \(F\)（\(F_n \overset{\Gamma}{\to} F\)），如果： - \textbf{下界}：对任何 \(x_n \to x\)，\(\liminf F_n(x_n) \geq F(x)\) - \textbf{上界}：对任何 \(x\)，存在收敛到 \(x\) 的序列 \(x_n\)，使得 \(\limsup F_n(x_n) \leq F(x)\)

\textbf{定理 251.6}（\(\Gamma\)-收敛与最小化子的收敛）：若 \(F_n \overset{\Gamma}{\to} F\) 且 \(x_n\) 是 \(F_n\) 的最小化子，则 \(x_n\) 的每个聚点是 \(F\) 的最小化子。

\textbf{证明}：设 \(F_n\overset{\Gamma}{\to}F\)（即下界不等式 \(F(x)\le\liminf F_n(x_n)\) 对所有 \(x_n\to x\)，以及恢复序列的存在性：对每个 \(x\) 存在 \(x_n\to x\) 使 \(F(x)\ge\limsup F_n(x_n)\)）。若 \(x_n\) 是 \(F_n\) 的最小化子（\(F_n(x_n)=\min F_n\)），取子列 \(x_{n_k}\to x\)。对任意 \(y\)，取恢复序列 \(y_n\to y\)，由 \(\Gamma\)-收敛下界：\(F(x)\le\liminf F_{n_k}(x_{n_k})\le\liminf F_{n_k}(y_{n_k})\le\limsup F_{n_k}(y_{n_k})\le F(y)\)。故 \(F(x)\le F(y)\) 对所有 \(y\) 成立，即 \(x\) 为 \(F\) 的最小化子。\(\blacksquare\)

\textbf{定理 251.7}（Modica-Mortola 定理 / Allen-Cahn 泛函的 \(\Gamma\)-收敛）：Allen-Cahn 泛函

\[F_\varepsilon(u) = \int_{\Omega} \left( \frac{\varepsilon}{2} |\nabla u|^2 + \frac{1}{\varepsilon} W(u) \right) dx\]

（\(W(u) = \frac{1}{4}(1-u^2)^2\) 是双井势）\(\Gamma\)-收敛到周长泛函 \(\sigma \operatorname{Per}(\{u = 1\})\)（模去常数），其中 \(\sigma = \int_{-1}^1 \sqrt{2W(s)} ds\)。

\textbf{证明}：Allen-Cahn 泛函 \(F_\varepsilon(u)=\int_\Omega(\varepsilon|\nabla u|^2+\frac{1}{\varepsilon}W(u))dx\)（\(W\) 为双井位势，如 \(W(u)=u^2(1-u)^2\)）在 \(L^1\) 拓扑下 \(\Gamma\)-收敛于 \(c_W\cdot\operatorname{Per}(\{u=1\})\)（\(BV\) 泛函）。下界不等式：对 \(u_\varepsilon\to \chi_E\)（\(L^1\)），\(\liminf F_\varepsilon(u_\varepsilon)\ge c_W\operatorname{Per}(E)\) 由 Modica 的梯度下界估计 \(\varepsilon|\nabla u_\varepsilon|^2+\frac{1}{\varepsilon}W(u_\varepsilon)\ge 2\sqrt{W(u_\varepsilon)}|\nabla u_\varepsilon|\) 和 coarea 公式结合推出。恢复序列：取 \(u_\varepsilon(x)=q(d(x,\partial E)/\varepsilon)\)（\(q\) 为 \(W\) 的一维异宿轨道，\(q'=\sqrt{2W}\)），直接验证 \(F_\varepsilon(u_\varepsilon)\to c_W\operatorname{Per}(E)\)。\(\blacksquare\)

\hrulefill

\hrulefill

\hrulefill

\hrulefill

\hrulefill

\hrulefill

\textbf{第216章（补充）}

高维渐近凸几何（Asymptotic Convex Geometry / Local Theory of Banach Spaces）研究高维凸体在大维数极限下的渐近行为。它以 Dvoretzky 定理（1971）和 Milman 的浓度现象为核心，连接泛函分析、概率论和组合数学。Princeton Companion 将其列为独立数学分支。

\subsection{252.1 Dvoretzky定理}\label{dvoretzkyux5b9aux7406}

\textbf{定理 252.1}（Dvoretzky定理，1971）：对任意 \(\varepsilon > 0\) 和正整数 \(k\)，存在 \(N = N(k, \varepsilon)\) 使得：对任意维数 \(n \geq N\) 的 Banach 空间 \(X\)（赋予任意范数），存在 \(X\) 的 \(k\) 维子空间 \(E\)，其在 Banach-Mazur 距离 \(d_{\mathrm{BM}}(E, \ell_2^k)\) 意义下是近似 Euclid 的： \[d_{\mathrm{BM}}(E, \ell_2^k) \leq 1 + \varepsilon\] 其中 \(d_{\mathrm{BM}}(X, Y) = \inf \{\|T\| \cdot \|T^{-1}\| : T: X \to Y \text{ 线性同构}\}\)。

\emph{意义}：Dvoretzky 定理声称：每个充分高维的 Banach 空间都包含任意维数的近似 Euclid 子空间------无论原空间的几何多么''非欧''（如 \(\ell_\infty^n\) 的单位球是立方体，其中也包含近似球形的截面）。这是''高维凸体总有近似球截面''的深刻陈述。

\textbf{Milman的简短证明}（1971，5页）：核心工具是 Lévy 族的测度集中现象------球面 \(S^{n-1}\) 上的 Haar 测度高度集中在赤道附近。用等周不等式和 Lipschitz 函数的浓度估计证明：任意高维凸体的 Minkowski 泛函在大部分方向上近似常数，从而存在高维近似 Euclid 截面。

\textbf{证明}：\(n\) 维 Banach 空间 \((\mathbb{R}^n,\|\cdot\|)\) 的高维截面几乎为 Euclid 空间。证明：(1) 利用 \(S^{n-1}\) 上范数的集中现象（由 Lévy 引理）：\(\|\cdot\|\) 的 \(1\)-Lipschitz 函数中位数为 \(M\)，则 \(\mu(\{|\|x\|-M|>\varepsilon\})\le 2e^{-c\varepsilon^2 n}\)。(2) 取 \(k\) 维随机子空间 \(E\)（按 Haar 测度），则该子空间单位球面上的范数几乎处处接近常数。(3) 三角不等式与 John 椭球结合给出子空间 \(E\) 上的范数与 Euclid 范数的 \((1+\varepsilon)\)-等价性。\(N(k,\varepsilon)\) 的存在性来自 \(n\) 足够大时集中现象的支配作用（\(n\gtrsim C k/\varepsilon^2\)）。\(\blacksquare\)

\subsection{252.2 John椭球与Banach-Mazur距离}\label{johnux692dux7403ux4e0ebanach-mazurux8dddux79bb}

\textbf{定义 252.1}（John椭球 / 最大体积内切椭球）：\(\mathbb{R}^n\) 中凸体 \(K\) 的 \textbf{John 椭球} \(\mathcal{E}_{\max}(K)\) 是含于 \(K\) 的具有最大体积的椭球（唯一）。若 \(K\) 是中心对称的，则 \(n^{-1/2} \mathcal{E}_{\max}(K) \subset K \subset \mathcal{E}_{\max}(K)\)。

\textbf{定理 252.2}（John定理，1948）：对 \(\mathbb{R}^n\) 中任意中心对称凸体 \(K\)，存在以原点为中心的椭球 \(\mathcal{E}\) 满足 \[\mathcal{E} \subset K \subset \sqrt{n} \cdot \mathcal{E}\] 等价地，任何 \(n\) 维 Banach 空间到 \(\ell_2^n\) 的 Banach-Mazur 距离不超过 \(\sqrt{n}\)：\(d_{\mathrm{BM}}(X, \ell_2^n) \leq \sqrt{n}\)。这个上界是紧的（\(\ell_\infty^n\) 达到 \(\sqrt{n}\)）。

\textbf{推论 252.1}（Banach-Mazur紧致性）：所有 \(n\) 维 Banach 空间的等距类构成的 Banach-Mazur 紧致度量空间 \(\mathcal{B}_n\) 的直径为 \(O(\sqrt{n})\)。

\textbf{证明}：对中心对称凸体 \(K\subset\mathbb{R}^n\)，John 椭球 \(\mathcal{E}\) 是 \(K\) 的最大体积内切椭球。设 \(\mathcal{E}=B_2^n\)（单位球），则 \(K\subset\sqrt{n}B_2^n\) 由接触点处的极值性质导出：存在接触点 \(x_i\in\partial K\cap\partial B_2^n\) 和正数 \(c_i>0\) 使 \(\sum c_i x_i\otimes x_i = I_n\)（Fritz John 分解）。取极体 \(K^\circ\)，同理得 \(B_2^n/\sqrt{n}\subset K\)。该定理给出任意有限维赋范空间与 Euclid 空间之间距离的普适量级上界（\(\sqrt{n}\)），是最优的------\(\ell_\infty^n\) 的 John 椭球为 \(B_2^n\) 且比值精确为 \(\sqrt{n}\)。\(\blacksquare\)

\subsection{252.3 凸体的渐近浓度现象}\label{ux51f8ux4f53ux7684ux6e10ux8fd1ux6d53ux5ea6ux73b0ux8c61}

浓度现象是高维几何中最引人注目的特征之一：在高维空间中，Lipschitz 函数在概率意义下几乎为常数，偏离中位数的测度以维数的指数速率衰减。这一现象由 Milman 在 1971 年作为 Dvoretzky 定理证明的核心工具引入，它从根本上解释了为什么高维 Banach 空间必然包含近 Euclid 截面------范数作为 \(S^{n-1}\) 上的 Lipschitz 函数，在高维下几乎处处集中在其中位值附近。

\textbf{定理 252.3}（Milman的浓度现象，1971）：设 \(\mu\) 是 \(S^{n-1}\) 上的归一化 Haar 测度（旋转不变概率测度）。对任意 1-Lipschitz 函数 \(f: S^{n-1} \to \mathbb{R}\)（\(|f(x) - f(y)| \leq \|x - y\|\)），有 \[\mu\left(\{x : |f(x) - M_f| \geq t\}\right) \leq C e^{-c n t^2}\] 其中 \(M_f\) 是 \(f\) 的中位数（或均值），\(C, c > 0\) 是万有常数。\emph{注：高维球面上的 Lipschitz 函数高度集中在其中位数附近------偏离的测度指数衰减于维数。}

\textbf{证明}：对 \(S^{n-1}\) 上 \(1\)-Lipschitz 函数 \(f\)，设 \(M\) 为其中位数（\(\mu(\{f\le M\})\ge 1/2\) 且 \(\mu(\{f\ge M\})\ge 1/2\)）。由 Lévy 等周不等式：\(S^{n-1}\) 上子集 \(A\) 的 \(\varepsilon\)-展开 \(A_\varepsilon\) 满足 \(\mu(A_\varepsilon)\ge 1-e^{-c\varepsilon^2 n}\)（若 \(\mu(A)\ge 1/2\)）。取 \(A=\{f\le M\}\)，\(A_\varepsilon\subset\{f\le M+\varepsilon\}\)（因 \(f\) 为 \(1\)-Lipschitz），故 \(\mu(\{f>M+\varepsilon\})\le e^{-c\varepsilon^2 n}\)。对称地 \(\mu(\{f<M-\varepsilon\})\le e^{-c\varepsilon^2 n}\)，联立得 Gaussian 浓度。该现象说明高维球面上 Lipschitz 函数几乎为常数。\(\blacksquare\)

\textbf{定理 252.4}（Blaschke-Santaló不等式）：中心对称凸体 \(K \subset \mathbb{R}^n\) 的极体 \(K^\circ = \{y : \langle x, y \rangle \leq 1, \forall x \in K\}\) 满足 \[\operatorname{Vol}(K) \cdot \operatorname{Vol}(K^\circ) \leq \operatorname{Vol}(B_2^n)^2\] 其中 \(B_2^n\) 是 Euclid 单位球。等号当且仅当 \(K\) 是椭球（Santaló 1949，等价性由 Saint-Raymond 1981 证明）。

\textbf{定义 252.2}（Mahler猜想，1939）：对中心对称凸体 \(K\)，\(\operatorname{Vol}(K) \cdot \operatorname{Vol}(K^\circ) \geq 4^n / n!\)，且极值由立方体 \([-1, 1]^n\) 及其极体（\(\ell_1^n\) 球 × 交叉多胞体）达到。该猜想刚被部分解决（Iriyeh-Shibata 2020, 3维情形）。

\textbf{证明}：对中心对称凸体 \(K\)，极体 \(K^\circ=\{y:\langle x,y\rangle\le 1,\forall x\in K\}\) 满足 \(|K||K^\circ|\le\omega_n^2\)。Bourgain-Milman 证明了反向不等式：\(|K||K^\circ|\ge c^n\omega_n^2\)（\(c>0\) 为绝对常数）。Blaschke-Santaló 的证明使用仿射等周不等式：通过 \(\operatorname{SL}(n)\) 变换将 \(K\) 置于''各向同性位置''（使惯性矩阵正比于单位矩阵），此时 \(K\) 的体积与极体体积达到平衡。Santaló 点（使 \(|K^\circ|\) 最小的内点）处的临界条件给出最优上界 \(\omega_n^2\)，等号成立当且仅当 \(K\) 为以原点为中心的椭球。\(\blacksquare\)

\textbf{定理 252.5}（等周不等式与 Brunn-Minkowski 的联系）：经典的 Brunn-Minkowski 不等式 \(\operatorname{Vol}(A + B)^{1/n} \geq \operatorname{Vol}(A)^{1/n} + \operatorname{Vol}(B)^{1/n}\) 蕴含了 \(\mathbb{R}^n\) 中的等周不等式和所有函数形式的 Borell-Brascamp-Lieb 不等式------这是测度集中的分析基础。

\textbf{证明}：Brunn-Minkowski 不等式：\(|A+B|^{1/n}\ge |A|^{1/n}+|B|^{1/n}\)（\(A,B\) 为 \(\mathbb{R}^n\) 中紧致子集）。取 \(A\) 为区域 \(E\)，\(B=\varepsilon B_2^n\)（小半径球体），则 \(A+B\) 为 \(E\) 的 \(\varepsilon\)-管状邻域 \(E_\varepsilon\)。展开：\(|E+\varepsilon B_2^n| = |E| + \varepsilon\operatorname{Per}(E) + o(\varepsilon)\)（Minkowski 内容公式）。代入 Brunn-Minkowski：\((|E|+\varepsilon\operatorname{Per}(E))^{1/n}\ge |E|^{1/n}+\varepsilon\omega_n^{1/n}\)。Taylor 展开得 \(\operatorname{Per}(E)\ge n\omega_n^{1/n}|E|^{(n-1)/n}\)（等周不等式）。反向地，等周不等式通过积分给出 Brunn-Minkowski，故两者等价。\(\blacksquare\)

\subsection{252.4 各向同性常数与切片问题}\label{ux5404ux5411ux540cux6027ux5e38ux6570ux4e0eux5207ux7247ux95eeux9898}

\textbf{定义 252.3}（各向同性凸体 / 各向同性常数）：凸体 \(K \subset \mathbb{R}^n\)（体积为 1）称为\textbf{各向同性}的，若其重心在原点且其惯性矩阵是对角的：\(\int_K x_i x_j dx = L_K^2 \delta_{ij}\)。参数 \(L_K\) 是 \(K\) 的\textbf{各向同性常数}。

\textbf{切片问题}（Bourgain的切片问题 / 超平面猜想，1986）：是否存在万有常数 \(C > 0\)，使得对任意 \(n\) 和任意体积为 1 的中心对称凸体 \(K \subset \mathbb{R}^n\)，存在超平面截面 \(K \cap H\) 的 \((n-1)\) 维体积至少为 \(C\)？等价于 \(L_K \leq C\) 是否对所有各向同性凸体成立？该问题被 Klartag（2006）和 Chen（2021）证明 \(L_K \leq C\)（万有常数存在）------这是高维凸几何近年的重大突破。

\textbf{定理 252.6}（薄壳不等式 / Thin Shell，Klartag 2007）：各向同性凸体 \(K\)（体积 1）上的均匀测度集中在距原点约 \(\sqrt{n} L_K\) 的薄球壳上。更精确地，\(\operatorname{Var}_K(\|x\|) \leq C L_K^2\)（方差不随维数增长）。

\textbf{证明}：各向同性凸体 \(K\)（\(|K|=1\)，重心在原点，惯性矩阵正比于单位阵）上的均匀测度，其范数方差 \(\sigma_K^2=\operatorname{Var}_K(\|x\|)\) 称为薄壳宽度。Klartag 证明 \(\sigma_K\le C n^{-1/3}\)（\(C\) 为普适常数），后经 Eldan 和 Lee-Vempala 改进至 \(\sigma_K\le C L_K n^{-1/2}\)（\(L_K\) 为各向同性常数）。该不等式说明高维各向同性凸体上的均匀测度集中在半径约 \(\sqrt{n}L_K\) 的薄球壳上，其厚度相对于半径趋于零。这是中心极限定理在各向同性凸体类上的概率极限定理形式。\(\blacksquare\)

\hrulefill

\emph{本卷在实分析（V9）和泛函分析（V14）的基础上，系统建立了 Hausdorff 测度与分形维数、自相似与多重分形分析、整流与电流（Federer 理论）、Plateau 问题（Douglas 解与正则性）和几何测度论变分法（Caccioppoli 集、BV 函数、\(\Gamma\)-收敛）的理论体系。全书以分形几何的「无限精细结构」和几何测度论的「广义曲面」收尾，展示了数学处理非光滑几何对象的强大能力。共 5 章。}

\emph{本卷建立了几何分析的核心框架：极小曲面（面积泛函、Plateau 问题、Bernstein 定理）是几何变分问题最经典的范例；调和映射（能量泛函、Eells-Sampson 定理、气泡分析）将调和函数推广到流形间的映射；Ricci 流（Hamilton 方程、Perelman 的非塌缩和手术理论、Poincaré 猜想的证明）是 21 世纪数学最伟大的成就之一；Yamabe 问题（共形几何、临界 Sobolev 指数、正质量定理）揭示了共形不变量的深层结构；平均曲率流和 Yang-Mills 方程将几何 PDE 的方法推广到更广泛的几何演化问题；补充部分将分形几何与几何测度论（Hausdorff 测度与维数、自相似与多重分形、整流与电流理论、Plateau 问题的测度论解、Caccioppoli 集与 BV 函数、Γ-收敛）纳入本卷框架。几何分析是当代数学最活跃的前沿领域之一。}

\emph{卷三十五：几何分析终。}

\emph{卷三十五：几何分析终。}

\emph{卷三十五：几何分析终。}

\emph{卷三十二：几何分析终。}

\hrulefill

\subsection{210.A 几何测度论补充}\label{a-ux51e0ux4f55ux6d4bux5ea6ux8bbaux8865ux5145}

\begin{quote}
分形几何（Fractal Geometry）由 Benoît Mandelbrot 在 1970 年代系统建立，研究处处不可微的、具有自相似结构的「粗糙」集合。几何测度论（Geometric Measure Theory）由 De Giorgi、Federer、Fleming 等人在 1960-1970 年代发展，使用测度论工具研究 \(\mathbb{R}^n\) 中「广义曲面」的几何与拓扑性质，特别是 Plateau 问题（存在以给定曲线为边界的极小面积曲面）的严格数学解。两个领域都处理传统光滑微分几何无法描述的「奇异」几何对象，在现代分析、变分法和图像处理中有深刻应用。本卷在实分析（V9）、测度论（V9 Ch 65）和泛函分析（V14）的基础上，建立 Hausdorff 测度、分形维数、整流与电流、Plateau 问题和变分法的几何测度论基础。
\end{quote}

\hrulefill

\hrulefill

\hrulefill

\hrulefill

\hrulefill

\hrulefill

\section{第272章：Hausdorff 测度与分形维数}\label{ux7b2c272ux7ae0hausdorff-ux6d4bux5ea6ux4e0eux5206ux5f62ux7ef4ux6570}

分形几何的核心概念是用非整数维数来度量集合的「粗糙度」。Hausdorff 测度和 Hausdorff 维数是最基本的工具。

\subsection{247.1 Hausdorff 测度}\label{hausdorff-ux6d4bux5ea6}

\textbf{定义 247.1}（\(s\)-维 Hausdorff 测度）：对 \(s \geq 0\) 和 \(A \subset \mathbb{R}^n\)，定义

\[\mathcal{H}^s(A) = \lim_{\delta \to 0} \mathcal{H}^s_\delta(A)\]

其中

\[\mathcal{H}^s_\delta(A) = \inf \left\{ \sum_{i=1}^\infty (\operatorname{diam} U_i)^s : A \subset \bigcup_i U_i,\ \operatorname{diam} U_i \leq \delta \right\}\]

\(\mathcal{H}^s\) 是 \(\mathbb{R}^n\) 上的外测度（实际上是 Borel 测度）。

Hausdorff 测度的定义核心是一族可变直径的覆盖，这与 Lebesgue 测度使用立方体覆盖的刚性结构形成了对比------Hausdorff 的灵活性使得它可以测量任意维数 \(s \geq 0\) 的集合，包括那些对 Lebesgue 测度而言''零测''的分形集。在深入讨论维数之前，首先需要确认 \(\mathcal{H}^s\) 确实是一个良定义的测度，满足我们所期望的测度论基本性质。下面的命题汇集了 \(\mathcal{H}^s\) 最常用的测度论属性，其中距离集可加性------间隔为正距离的两个集合的 Hausdorff 测度具有可加性------是区分 Hausdorff 测度与一般 Borel 测度的关键特征。

\textbf{命题 247.1}（Hausdorff 测度）：\(\mathcal{H}^s\) 是 \(\mathbb{R}^n\) 上的 Borel 正则外测度，满足：（1）单调性：\(A \subset B \Rightarrow \mathcal{H}^s(A) \leq \mathcal{H}^s(B)\)；（2）可数次可加性：\(\mathcal{H}^s(\bigcup_n A_n) \leq \sum_n \mathcal{H}^s(A_n)\)；（3）距离集可加性：若 \(\operatorname{dist}(A,B) > 0\)，则 \(\mathcal{H}^s(A \cup B) = \mathcal{H}^s(A) + \mathcal{H}^s(B)\)。特别地，\(\mathcal{H}^0\) 为计数测度，\(\mathcal{H}^n\) 与 \(\mathcal{L}^n\) 仅差一个常数因子。

\textbf{证明}：（1）单调性：若 \(A \subset B\)，则 \(A\) 的每个 \(\delta\)-覆盖也是 \(B\) 的 \(\delta\)-覆盖，取 \(\delta \to 0\) 得 \(\mathcal{H}^s(A) \leq \mathcal{H}^s(B)\)。（2）可数次可加性：设 \(\varepsilon > 0\)，对每个 \(A_n\) 取 \(\delta\)-覆盖 \(\{U_{n,i}\}\) 使 \(\sum_i (\operatorname{diam} U_{n,i})^s \leq \mathcal{H}^s_\delta(A_n) + \varepsilon/2^n\)，则 \(\{U_{n,i}\}_{n,i}\) 是 \(\bigcup_n A_n\) 的 \(\delta\)-覆盖，\(\mathcal{H}^s_\delta(\bigcup_n A_n) \leq \sum_n \mathcal{H}^s_\delta(A_n) + \varepsilon\)，取 \(\delta \to 0\) 和 \(\varepsilon \to 0\) 即得。（3）距离集性质：当 \(\operatorname{dist}(A,B) > 0\) 时，取 \(\delta < \operatorname{dist}(A,B)/2\)，则任何 \(\delta\)-覆盖无法同时覆盖 \(A\) 和 \(B\) 中的点，故 \(\mathcal{H}^s_\delta(A \cup B) = \mathcal{H}^s_\delta(A) + \mathcal{H}^s_\delta(B)\)，取 \(\delta \to 0\) 得可加性。\(\blacksquare\)

Hausdorff 测度最引人注目的行为是它作为维数 \(s\) 的函数具有一个非常特殊的性质：对固定的集合 \(A\)，\(\mathcal{H}^s(A)\) 在某个临界值 \(s_0\) 处从无穷跳跃到零。这一''跳跃''现象正是 Hausdorff 维数定义的基础------它表明每个集合都有一个唯一的维数阈值，在此阈值之下测度为无穷（表示覆盖''太细''），在此之上测度为零（表示覆盖''太粗''）。下面的定理精确地刻画了这一跳跃行为及其关键推论。

\textbf{定理 247.1}（Hausdorff 测度的行为）：函数 \(s \mapsto \mathcal{H}^s(A)\) 在某个唯一的 \(s_0\) 处从 \(\infty\) 跳到 \(0\)。即在 \(s < s_0\) 时 \(\mathcal{H}^s(A) = \infty\)，\(s > s_0\) 时 \(\mathcal{H}^s(A) = 0\)（\(\mathcal{H}^{s_0}(A)\) 可能为 \(0\)、正有限或无穷）。

\textbf{证明}：设 \(\mathcal{H}^t(A) < \infty\)，对任意 \(s > t\) 和 \(\delta\)-覆盖 \(\{U_i\}\)，\(\sum (\operatorname{diam} U_i)^s \leq \delta^{s-t} \sum (\operatorname{diam} U_i)^t\)。取 \(\delta \to 0\)，若 \(\mathcal{H}^t(A) < \infty\) 则 \(\mathcal{H}^s(A) = 0\)。故函数 \(s \mapsto \mathcal{H}^s(A)\) 至多在一个临界值 \(s_0\) 处从 \(\infty\) 跳到 \(0\)。可数稳定性由外测度的可数次可加性直接得到。Lipschitz 不变性：对 Lipschitz 映射 \(f\)，每个 \(\delta\)-覆盖 \(\{U_i\}\) 在 \(f\) 下的像 \(\{f(U_i)\}\) 是 \(\operatorname{Lip}(f)\delta\)-覆盖，故 \(\mathcal{H}^s_{\operatorname{Lip}(f)\delta}(f(A)) \leq \sum (\operatorname{diam} f(U_i))^s \leq (\operatorname{Lip}(f))^s \sum (\operatorname{diam} U_i)^s\)，取 \(\delta \to 0\) 得 \(\mathcal{H}^s(f(A)) \leq (\operatorname{Lip}(f))^s \mathcal{H}^s(A)\)。\(\blacksquare\)

\subsection{247.2 分形维数}\label{ux5206ux5f62ux7ef4ux6570}

\textbf{定义 247.2}（Hausdorff 维数）：集合 \(A \subset \mathbb{R}^n\) 的\textbf{Hausdorff 维数}为

\[\dim_H A = \sup\{s \geq 0 : \mathcal{H}^s(A) = \infty\} = \inf\{s \geq 0 : \mathcal{H}^s(A) = 0\}\]

Hausdorff 维数可以是任意非负实数。

\textbf{定义 247.3}（盒维数 / Box-Counting Dimension）：设 \(N_\delta(A)\) 是覆盖 \(A\) 所需的直径为 \(\delta\) 的集合的最少数目。\textbf{下盒维数}和\textbf{上盒维数}分别为

\[\underline{\dim}_B A = \liminf_{\delta \to 0} \frac{\log N_\delta(A)}{-\log \delta}, \quad \overline{\dim}_B A = \limsup_{\delta \to 0} \frac{\log N_\delta(A)}{-\log \delta}\]

若上下盒维数相等，则 \(\dim_B A\) 是\textbf{盒维数}（也称 Minkowski 维数）。一般 \(\dim_H A \leq \underline{\dim}_B A \leq \overline{\dim}_B A\)。

\textbf{注 247.1}（Hausdorff 维数与盒维数的关系）：不等式 \(\dim_H A \leq \underline{\dim}_B A\) 成立的原因是：盒维数使用固定大小的覆盖（所有覆盖集直径 \(\leq \delta\)），而 Hausdorff 维数允许任意尺寸的覆盖（各覆盖集直径可不同），故 Hausdorff 测度更''精细''，维数更低。严格不等式的典型例子：\(\mathbb{Q} \cap [0,1]\) 作为可数集满足 \(\dim_H = 0\)，但因稠密性其盒维数 \(\dim_B = 1\)。另一个经典例子是 \(\{1, 1/2, 1/3, \ldots\} \cup \{0\}\)：\(\dim_H = 0\) 但盒维数为 \(1/2\)。这些例子表明盒维数对''可数个点的极限密度''敏感，而 Hausdorff 维数对于可数集恒为零。对''好''的自相似分形（如 Cantor 集、Sierpiński 三角形），二者通常相等。

\subsection{247.3 自相似集与迭代函数系}\label{ux81eaux76f8ux4f3cux96c6ux4e0eux8fedux4ee3ux51fdux6570ux7cfb}

\textbf{定义 247.4}（迭代函数系 / IFS）：\textbf{迭代函数系}是一族收缩映射 \(\{f_1, \ldots, f_m\}\)（\(f_i: \mathbb{R}^n \to \mathbb{R}^n\)，\(\operatorname{Lip}(f_i) < 1\)）。IFS 的\textbf{吸引子}（或不变集）\(K\) 是满足 \(K = \bigcup_{i=1}^m f_i(K)\) 的唯一非空紧致集合。

迭代函数系（IFS）是构造自相似分形的最通用的框架。每个收缩映射 \(f_i\) 对应一个''副本''变换，将这些副本的并集取不动点即得整个分形。Hutchinson 在 1981 年给出了这一构造的严格数学基础，并证明了吸引子的存在唯一性。然而，仅知道吸引子的存在还不够------我们还需要知道它的维数。当各个变换 \(f_i\) 是相似比已知的相似压缩时，Hutchinson 进一步发现了一种极为优美的事实：在开集条件（确保变换的像不重叠太多）下，维数由一条简单的代数方程------Moran 方程------精确确定。这条方程将纯几何的维数问题转化为解一个多项式的根，是分形维数计算中最经典也最实用的工具。

\textbf{定理 247.2}（Hutchinson 定理，1981）：设 \(\{f_i\}\) 是相似比为 \(c_i\) 的相似压缩。若开集条件（OSC）成立（存在非空有界开集 \(V\) 使得 \(f_i(V) \subset V\) 且 \(f_i(V)\) 两两不交），则吸引子的 Hausdorff 维数 \(s\) 满足\textbf{Moran 方程}：

\[\sum_{i=1}^m c_i^s = 1\]

且 \(0 < \mathcal{H}^s(K) < \infty\)。

\textbf{证明}：设 \(K\) 为 IFS 吸引子，\(K = \bigcup_{i=1}^m f_i(K)\)。由 Hausdorff 测度的 Lipschitz 性质，\(\mathcal{H}^s(K) \leq \sum_i c_i^s \mathcal{H}^s(K)\)。若 \(\mathcal{H}^s(K) < \infty\)，则必须 \(\sum c_i^s \geq 1\)。上界：对 \(K\) 构造由迭代 \(f_{i_1} \circ \cdots \circ f_{i_k}\) 生成的 \(\delta\)-覆盖，利用 OSC 保证覆盖元素不重叠，得 \(\mathcal{H}^s(K) \leq C \sum c_{i_1}^s \cdots c_{i_k}^s = C (\sum c_i^s)^k \to 0\) 当 \(\sum c_i^s < 1\) 且 \(k \to \infty\)，故 \(s > \dim_H K\) 时 \(\mathcal{H}^s(K) = 0\)。下界：利用 OSC 中开集 \(V\) 上分布的自相似测度（质量分布原理），构造概率测度 \(\mu\) 使其满足 \(\mu(f_{i_1} \circ \cdots \circ f_{i_k}(K)) = c_{i_1}^s \cdots c_{i_k}^s\) 且 \(\mu(B(x,r)) \leq C r^s\)，由 Frostman 引理得 \(\mathcal{H}^s(K) > 0\)。故 \(\dim_H K = s\) 为 Moran 方程 \(\sum c_i^s = 1\) 的唯一解，且 \(0 < \mathcal{H}^s(K) < \infty\)。\(\blacksquare\)

\hrulefill

\hrulefill

\hrulefill

\hrulefill

\hrulefill

\hrulefill

\section{第273章：分形结构的自相似与多重分形}\label{ux7b2c273ux7ae0ux5206ux5f62ux7ed3ux6784ux7684ux81eaux76f8ux4f3cux4e0eux591aux91cdux5206ux5f62}

许多分形不能仅由一个维数完全描述------它们具有多层次的分形结构，需要\textbf{多重分形分析}（multifractal analysis）来刻画。

\subsection{248.1 自仿射分形与随机分形}\label{ux81eaux4effux5c04ux5206ux5f62ux4e0eux968fux673aux5206ux5f62}

自仿射分形是自相似分形的重要推广，由仿射压缩组成 IFS 的吸引子。与自相似分形（各方向等比例收缩）不同，自仿射分形允许不同方向有不同的收缩比，这使得其几何结构更加丰富------例如，自仿射分形在不同方向上的 Hausdorff 维数可以不同，且其维数一般不能由简单的 Moran 方程给出。自仿射分形的维数计算通常需要强有力的工具，如 Falconer 的奇异值函数方法（singular value function method）和 Bedford-McMullen 对自仿射地毯的显式计算。随机分形是另一类重要的分形对象，它们不是由确定性规则生成，而是由随机过程定义。\textbf{扩散限制聚集}（DLA / Diffusion-Limited Aggregation，Witten-Sander 1981）是最经典的随机分形模型之一，其生成过程为：随机粒子沿 Brown 运动轨迹从远处扩散，当遇到已有团簇时黏附。DLA 团簇展现出极其复杂的枝晶状结构，其分形维数在二维情形下约为 \(1.71\)，且团簇的外部分支是多重分形的。DLA 的数学分析极度困难------其 Laplace 生长方程（粒子的扩散满足 Laplace 方程 \(\Delta u = 0\)，边界条件为团簇表面 \(u = 0\) 和无穷远处 \(u = 1\)）与自由边界问题密切相关，至今没有严格的数学理论能完全解释 DLA 的枝晶形态。随机分形还出现在许多自然现象中------如河网的形状、闪电的路径、肺部支气管树的分布------这些自然分形结构通常既具有统计自相似性，又展现出随机性，是分形几何与统计物理交叉的重要研究对象。

\textbf{定义 248.1}（自仿射分形 / Self-Affine Fractal）：由仿射压缩 \(f_i(x) = A_i x + b_i\)（\(A_i\) 是收缩矩阵，各方向收缩比不同）组成的 IFS 的吸引子。与自相似分形不同，自仿射分形的不同方向有不同的伸缩行为。

\textbf{定义 248.2}（随机分形与 DLA）：\textbf{扩散限制聚集}（DLA / Diffusion-Limited Aggregation，Witten-Sander 1981）通过随机粒子沿 Brown 运动轨迹黏附生长的过程产生随机分形。DLA 团簇的维数约为 \(1.71\)（二维）。

\subsection{248.2 多重分形分析}\label{ux591aux91cdux5206ux5f62ux5206ux6790}

\textbf{定义 248.3}（多重分形谱 / Multifractal Spectrum）：对测度 \(\mu\) 上的奇异性分析，定义局部 Hölder 指数 \(\alpha(x) = \lim_{r \to 0} \frac{\log \mu(B(x, r))}{\log r}\)。\textbf{多重分形谱} \(f(\alpha)\) 是满足 \(\alpha(x) = \alpha\) 的点的 Hausdorff 维数。

单一分形维数（如 Hausdorff 维数）描述的是集合的整体几何，但对于承载了一个测度 \(\mu\) 的分形------例如 Cantor 集上不均匀分布的质量------单维数无法区分''测度集中的区域''和''测度稀薄的区域''。多重分形分析正是在这一背景下引入的：它将集合按局部奇异性强度 \(\alpha\) 分层，计算每层的 Hausdorff 维数 \(f(\alpha)\)，从而得到一族而非单个维数。在实际计算中，\(f(\alpha)\) 的直接估计极为困难，但通过引入 Rényi 维数 \(D_q\) 和配分函数 \(Z_q\)，可以借助统计力学中的 Legendre 变换将局部信息转化为全局矩的渐近行为，从而使多重分形分析成为可计算的工具。

\textbf{定义 248.4}（Rényi 维数与配分函数）：定义配分函数 \(Z_q(\delta) = \sum_i \mu(B_i)^q\)（\(B_i\) 是边长 \(\delta\) 的网格覆盖）。广义维数 \(D_q\) 满足 \(Z_q(\delta) \sim \delta^{(q-1)D_q}\)（\(\delta \to 0\)）。通过 Legendre 变换：\(f(\alpha) = \inf_q (q \alpha - \tau(q))\)，\(\tau(q) = (q-1)D_q\)。

Rényi 维数族 \(\{D_q\}\) 通过参数 \(q\) 给测度的不同矩赋予不同的权重：\(q=0\) 恢复分形支持的盒维数，\(q=1\) 给出信息维数（Shannon 熵的标度行为），\(q=2\) 给出关联维数。当 \(q\) 遍历实数时，Legendre 变换建立了 \(D_q\) 与 \(f(\alpha)\) 之间的对偶关系------这是统计力学中配分函数与熵之间的对偶在分形理论中的直接对应。Frisch 和 Parisi 在 1985 年提出了一个著名的猜想：在充分发展的湍流中，能量耗散率的空间分布具有多重分形结构，其 \(f(\alpha)\) 谱给出了湍流间歇性的完整表征。下面的定理总结了多重分形形式论的核心结论。

\textbf{定理 248.1}（多重分形形式论 / Frisch-Parisi 猜想，1985）：在湍流和许多其他物理系统中，多重分形谱 \(f(\alpha)\) 是 \(\cap\) 形的，\(f(\alpha)\) 的最大值为信息维数 \(D_1\)，\(D_0\) 是测度支撑的盒维数。

\textbf{证明}：将分形测度 \(\mu\) 的支集按局部 Holder 指数 \(\alpha(x)=\lim_{r\to 0}\log\mu(B_r(x))/\log r\) 分层。令 \(E_\alpha=\{x:\alpha(x)=\alpha\}\)，\(f(\alpha)=\dim_H E_\alpha\) 为多重分形谱。大偏差原理给出 Legendre 变换关系：\(\tau(q)=\inf_\alpha(q\alpha-f(\alpha))\)，其中 \(\tau(q)\) 为相关指数（partition function 的标度指数）。\(f(\alpha)\) 为 \(\cap\)-形曲线，最大值 \(f(\alpha_{\max})=D_0\)（盒维数），\(f(\alpha_{\text{info}})=D_1\)（信息维数）。Frisch-Parisi 猜想将湍流能量耗散的间歇性归因于多重分形结构。\(\blacksquare\)

\subsection{248.3 渗流与分形}\label{ux6e17ux6d41ux4e0eux5206ux5f62}

\textbf{定义 248.5}（渗流模型 / Percolation Model）：在二维网格上，每条边以概率 \(p\) 被「打开」（可通行）。当 \(p\) 超过临界概率 \(p_c = 1/2\) 时，存在无穷大连通分量。

渗流模型简洁的定义------独立地以概率 \(p\) 打开每条边------背后蕴藏着极为丰富的相变和分形现象。当 \(p\) 从 \(0\) 增加到 \(1\) 时，系统先后经历无长程连通（亚临界）和出现无穷大连通分量（超临界）两个截然不同的相，相变恰好发生在 \(p = p_c = 1/2\)。正是在这个临界点上，连通团簇展现出精确的分形结构：团簇的质量与线性尺度满足幂律关系 \(M \sim L^{D_f}\)，其 Hausdorff 维数 \(D_f\) 可以用共形场论和 Schramm-Loewner 演化（SLE）严格计算。下面的定理汇集了临界渗流团簇在二维情形下的主要维数结果。

\textbf{定理 248.2}（渗流团簇的分形维数，Stanley 1977）：在 \(p = p_c\) 时，渗流团簇是分形。二维临界渗流团簇的 Hausdorff 维数为 \(91/48 \approx 1.896\)。团簇的化学维数（沿团簇测量的最短路径距离）为 \(d_{\min} \approx 1.13\)。

\textbf{证明}：二维临界点渗流（\(p=p_c=1/2\)）的最大团簇为分形。团簇质量 \(M\) 与线性尺度 \(L\) 的关系 \(M\sim L^{D_f}\) 给出 Hausdorff 维数 \(D_f=91/48\)（由共形场论和 SLE 过程严格推导）。化学维数 \(d_{\min}\) 描述沿团簇内部路径的最短距离：\(\ell\sim R^{d_{\min}}\)（\(R\) 为欧氏端到端距离），二维 \(d_{\min}\approx 1.130\)。这些维数通过渗流的共形不变性（Cardy 公式、SLE₆ 过程）与统计物理中的 Coulomb 气体映射建立联系。\(\blacksquare\)

\hrulefill

\hrulefill

\hrulefill

\hrulefill

\hrulefill

\hrulefill

\section{第274章：整流与电流}\label{ux7b2c274ux7ae0ux6574ux6d41ux4e0eux7535ux6d41}

几何测度论使用\textbf{整流}和\textbf{电流}（rectifiable sets and currents）的概念精密地研究奇异曲面和变分问题。

\subsection{249.1 可整流集}\label{ux53efux6574ux6d41ux96c6}

\textbf{定义 249.1}（\(k\)-维可整流集 / Rectifiable Set）：Borel 集 \(E \subset \mathbb{R}^n\) 是\textbf{可数 \(\mathcal{H}^k\)-可整流的}，如果存在可数多个 Lipschitz 映射 \(f_i: \mathbb{R}^k \to \mathbb{R}^n\) 使得

\[\mathcal{H}^k\left(E \setminus \bigcup_i f_i(\mathbb{R}^k)\right) = 0\]

即可整流集几乎被 Lipschitz 参数化的 \(k\) 维曲面覆盖。

\textbf{定理 249.1}（可整流集的结构定理）：\(k\)-维可整流集 \(E\) 在 \(\mathcal{H}^k\)-几乎处处有 \(k\) 维近似切空间 \(\operatorname{Tan}^k(E, x)\)。即对 \(\mathcal{H}^k\)-a.e. \(x \in E\)，

\[\lim_{r \to 0} \frac{\mathcal{H}^k(E \cap B(x, r))}{r^k} = 1\]

（在适当的旋转下，\(E\) 的放大极限是一个 \(k\) 维平面）。

\textbf{定义 249.2}（纯不可整流集 / Purely Unrectifiable Set）：Borel 集 \(E\) 是\textbf{纯不可整流的}，如果对每个 \(k\)-维可整流集 \(F\)，\(\mathcal{H}^k(E \cap F) = 0\)。Besicovitch 证明了二维中存在纯一维不可整流集。

\textbf{证明}：\(E\) 为 \(k\)-维可整流集意味着 \(E\) 几乎处处可由 \(k\) 维 Lipschitz 映射参数化：存在可数多个 \(k\) 维 Lipschitz 片 \(f_i:\mathbb{R}^k\to\mathbb{R}^n\) 使得 \(\mathcal{H}^k(E\setminus\bigcup f_i(\mathbb{R}^k))=0\)。对 \(\mathcal{H}^k\)-a.e. \(x\in E\)，近似切空间 \(\operatorname{Tan}^k(E,x)\) 的存在性由 Rademacher 定理和 Lebesgue 密度定理联合给出：将局部坐标拉伸到尺度 \(r^{-1}\)，Lipschitz 片的像在缩放极限下趋于 \(f_i\) 在 \(f_i^{-1}(x)\) 处的切空间（作为 \(\mathbb{R}^n\) 的 \(k\) 维子空间）。该极限不依赖于 Lipschitz 表示的选择。\(\blacksquare\)

\subsection{249.2 整系数整流链}\label{ux6574ux7cfbux6570ux6574ux6d41ux94fe}

整系数整流链（integral rectifiable currents）是几何测度论中描述''广义定向曲面''的核心概念，由 Federer 和 Fleming（1960）在解决 Plateau 问题的过程中引入。一个 \(k\)-维整系数整流链 \(T\) 由三个要素定义：一个 \(k\)-维可整流集 \(E\)（即几乎处处有 \(k\) 维近似切空间）、一个整数重数函数 \(\theta: E \to \mathbb{Z}\)（\(|theta|\) 表示曲面的''层数''）和一个定向 \(k\)-向量场 \(\xi\)。作为 \(k\)-形式 \(\omega\) 上的线性泛函，\(T(\omega) = \int_E \langle \omega(x), \xi(x) \rangle \theta(x) d\mathcal{H}^k(x)\)。整系数整流链空间记作 \(\mathcal{R}_k(\mathbb{R}^n)\)。整流链的关键性质包括：（1）紧致性定理（Federer-Fleming）：质量有界且边界质量有界的整流链序列有子列在平坦范数（flat norm）下收敛到一个整系数整流链；（2）边界 \(\partial T\) 定义为分布 \(\partial T(\omega) = T(d\omega)\)，且 \(\partial T\) 是 \((k-1)\)-维整流链；（3）\textbf{变形定理}（deformation theorem）：任何整流链可以通过径向投影和 CW 结构变形为标准网格上的胞腔链，且质量受控。变形定理是等周不等式的证明核心------它将连续几何问题转化为组合问题。整系数整流链的''整系数''条件至关重要：它保证了重数 \(\theta\) 取整数值，从而在变分问题中防止了''面积可以无限细分''的悖论。整系数平坦链（integral flat chains）是整系数整流链及其边界的 \(L^1\) 极限，构成了适合变分法的完备空间。这一框架的成功标志是：Federer-Fleming 使用整系数整流链理论严格证明了任意维数的 Plateau 问题------对任意 \((k-1)\)-维整系数整流链边界 \(T\)，存在 \(k\)-维整系数整流链 \(S\)（\(\partial S = T\)）使得 \(\mathbf{M}(S)\) 极小化。这是 20 世纪数学分析最重大的成就之一。

\textbf{定义 249.4}（质量与边界 / Mass and Boundary）：整流链 \(T\) 的\textbf{质量}为 \(\mathbf{M}(T) = \int_E |\theta(x)| d\mathcal{H}^k(x)\)。\(T\) 的\textbf{边界} \(\partial T\) 是对 \((k-1)\)-形式 \(\omega\) 定义的分布：\(\partial T(\omega) = T(d\omega)\)。

\textbf{定义 249.5}（整系数平坦链 / Integral Flat Chain）：整系数平坦链是整系数整流链和边界为整流链的整系数整流链之和的 \(L^1_{\text{loc}}\) 极限。整系数平坦链空间对于变分法非常重要（适合取极限）。

\subsection{249.3 等周不等式与 Sobolev 不等式}\label{ux7b49ux5468ux4e0dux7b49ux5f0fux4e0e-sobolev-ux4e0dux7b49ux5f0f}

\textbf{定理 249.2}（De Giorgi 等周不等式，1958）：对任何 \(n\) 维（\(n \geq 2\)）整系数整流链 \(T\) 满足 \(\partial T = 0\) 且紧支撑，存在 \(n+1\) 维整系数整流链 \(S\) 满足 \(\partial S = T\) 且 \[\mathbf{M}(S)^{n/(n+1)} \leq C_n \mathbf{M}(T)\] 其中 \(C_n\) 是仅依赖于 \(n\) 的常数。

更精确的等周不等式的几何测度论版本是 Plateau 问题解的基础。

\textbf{证明}：设 \(T\) 为 \(n\) 维整系数整流链，\(\partial T=0\)。De Giorgi 使用 \(T\) 的锥构造：对 \(T\) 的支撑集中任意点 \(p\)，作锥 \(C_p(T)=\bigcup_{x\in\operatorname{spt} T}[p,x]\)。\(\partial C_p(T)=T\)（模去删去含 \(p\) 的退化部分）。Federer-Fleming 证明 \(\mathbf{M}(C_p(T))\le \frac{R}{n+1}\mathbf{M}(T)\)，其中 \(R=\operatorname{spt}(T)\) 的直径。这是 De Giorgi 等周不等式的核心估计：对于支撑直径为 \(R\) 的 \(n\)-圈 \(T\)，存在 \(S\) 填充 \(T\)（\(\partial S=T\)）满足 \(\mathbf{M}(S)\le c_n \mathbf{M}(T)^{(n+1)/n}\)，常数 \(c_n\) 仅依赖 \(n\)。\(\blacksquare\)

De Giorgi 的等周不等式是几何测度论中第一个一般性结果，它给出了以任意整系数整流链为边界的广义曲面面积的下界。然而，De Giorgi 的证明依赖于锥构造技术，其常数并非最优。Federer 和 Fleming 随后引入了一种全新的证明策略------变形定理------从根本上改进了等周不等式的证明。变形定理的关键思想是将连续几何问题离散化：通过径向投影和 CW 结构，将任意整系数整流链挤压到网格骨架上，化为组合等周问题，再利用组合不等式得到最优常数 \(\gamma_n = n^{-1} \omega_n^{-1/n}\)。

\textbf{定理 249.3}（Federer-Fleming 等周不等式）：对 \(n\) 维整系数整流链 \(T\) 满足 \(\partial T = 0\) 且紧支撑，存在 \(n+1\) 维整系数整流链 \(S\) 满足 \(\partial S = T\) 且

\[\mathbf{M}(S)^{n/(n+1)} \leq \gamma \mathbf{M}(T)\]

其中 \(\gamma\) 是仅依赖于 \(n\) 的常数。

\textbf{证明}：Federer-Fleming 的证明使用变形定理（deformation theorem）：将 \(T\) 挤压到某个胞腔复形的 \(n\)-骨架上，化为组合等周问题。变形过程中的质量控制由距离函数的径向投影实现。对挤压后的组合链应用离散等周不等式（每个 \(n\)-胞腔由其 \((n+1)\)-胞腔填充），其解编码了几何等周常数。最优常数 \(\gamma_{n}=n^{-1}\omega_n^{-1/n}\) 由等周不等式的球面极值情形给出（球是最优填充区域）。该不等式等价于 Sobolev 不等式 \(\|f\|_{n/(n-1)}\le \gamma_n^{-1}\|\nabla f\|_1\) 的对偶形式。\(\blacksquare\)

等周不等式的变形定理证明漂亮地展示了离散化策略的威力，但它并未穷尽等周不等式与函数空间嵌入之间的关系。Federer-Fleming 等周不等式实际上等价于最优 Sobolev 嵌入不等式 \(\|f\|_{L^{n/(n-1)}} \leq \gamma_n^{-1} \|\nabla f\|_{L^1}\)。这一等价性的建立依赖于几何测度论的另一个基石------面积公式。面积公式将 Lipschitz 映射在低维曲面上的''拉伸''（Jacobian）与映射纤维的计数联系起来，本质上是积分变元替换公式在非一一映射和非满维情形下的推广。它是余面积公式的对偶，两者一起构成了几何测度论中坐标变换分析的基本工具。

\textbf{定理 249.4}（面积公式 / Area Formula）：设 \(f: \mathbb{R}^m \to \mathbb{R}^n\) 是 Lipschitz 映射（\(m \leq n\)），则对任意 \(\mathcal{L}^m\)-可测集 \(A \subset \mathbb{R}^m\)， \[\int_A J_m f(x) \, d\mathcal{L}^m(x) = \int_{\mathbb{R}^n} \mathcal{H}^0(A \cap f^{-1}(y)) \, d\mathcal{H}^m(y)\] 其中 \(J_m f(x) = \sqrt{\det((Df(x))^T Df(x))}\) 是 \(m\)-维 Jacobian。

\textbf{证明}：对 Lipschitz 映射 \(f:\mathbb{R}^m\to\mathbb{R}^n\)（\(m\le n\)），\(m\) 维 Jacobian \(J_m f(x)=\sqrt{\det(Df(x)^T Df(x))}\)（\(Df\) 的 Rademacher 导数几乎处处存在）。将 \(\mathbb{R}^m\) 按 \(f\) 的纤维分层，在每个法纤维上使用 \(m\) 维面积的参数化定义：\(\mathcal{H}^m(f(A))=\int_A J_m f(x) d\mathcal{L}^m(x) + \mathcal{H}^m(\{y:|f^{-1}(y)\cap A|>1\})\)（含重数）。当 \(f\) 单射时简化为 \(\int_A J_m f d\mathcal{L}^m\)。将 \(f\) 限制在 \(\{x\in A: J_m f(x)>0\}\) 上并应用 Vitali 覆盖引理和 Lipschitz 常数的积分表示，得完整面积公式。\(\blacksquare\)

\textbf{定理 249.5}（余面积公式 / Coarea Formula）：设 \(f: \mathbb{R}^n \to \mathbb{R}^m\) 是 Lipschitz 映射（\(n \geq m\)），则对任意 \(\mathcal{L}^n\)-可测集 \(A \subset \mathbb{R}^n\)， \[\int_A J_m f(x) \, d\mathcal{L}^n(x) = \int_{\mathbb{R}^m} \mathcal{H}^{n-m}(A \cap f^{-1}(y)) \, d\mathcal{L}^m(y)\]

\textbf{证明}：对 Lipschitz 映射 \(f: \mathbb{R}^n \to \mathbb{R}^m\)（\(n \geq m\)），\(m\) 维 Jacobian \(J_m f(x) = \sqrt{\det(Df(x) Df(x)^T)}\) 由 Rademacher 定理几乎处处存在。核心思路：将 \(f\) 分解为坐标分量，逐次应用 Fubini 定理。对 \(C^1\) 映射，余面积公式为标准变量替换：\(\int_A J_m f \, d\mathcal{L}^n = \int_{\mathbb{R}^m} \mathcal{H}^{n-m}(A \cap f^{-1}(y)) \, d\mathcal{L}^m(y)\)。对 Lipschitz 映射，用光滑函数逼近 \(f_\varepsilon \to f\)，由面积公式保证 \(J_m f_\varepsilon \to J_m f\) 在 \(L^1\) 中，且 \(\mathcal{H}^{n-m}\) 测度沿水平集的连续性（Federer 的余面积不等式）保证收敛。关键步骤：Lipschitz 函数的水平集 \(\mathcal{H}^{n-m}\)-可测性由 Whitney 延拓定理和 Sard 型定理保证。该公式是几何测度论中最通用的积分换元工具，推广了经典微积分中的变量替换公式。\(\blacksquare\)

\hrulefill

\hrulefill

\hrulefill

\hrulefill

\hrulefill

\newpage

\chapter{07-拓扑与几何/05-几何分析/06-凸几何.md}\label{ux62d3ux6251ux4e0eux51e0ux4f5505-ux51e0ux4f55ux5206ux679006-ux51f8ux51e0ux4f55.md}

\chapter{凸几何}\label{ux51f8ux51e0ux4f55}

\hrulefill

\section{第275章：凸体理论}\label{ux7b2c275ux7ae0ux51f8ux4f53ux7406ux8bba}

凸几何（Convex Geometry）研究欧几里得空间中凸体的几何性质与极值问题，其核心结果Brunn-Minkowski不等式与等周不等式深刻联系了几何、分析和泛函分析。本章系统阐述凸体理论的基础与深化成果。

\subsection{194.1 凸集基本性质}\label{ux51f8ux96c6ux57faux672cux6027ux8d28}

\textbf{定义 194.1.1（凸集与凸体）} 集合 \(K \subseteq \mathbb{R}^n\) 称为凸集，若对任意 \(x, y \in K\) 和 \(\lambda \in [0, 1]\)，有 \((1 - \lambda)x + \lambda y \in K\)。紧致且具有非空内部的凸集称为凸体（convex body）。\(\mathbb{R}^n\) 中所有凸体构成的集合记作 \(\mathcal{K}^n\)。原点关于 \(K\) 对称的凸体 (\(K = -K\)) 称为对称凸体，其全体记作 \(\mathcal{K}_s^n\)。

\textbf{定义 194.1.2（支撑函数）} 凸体 \(K \in \mathcal{K}^n\) 的支撑函数（support function）\(h_K: \mathbb{R}^n \to \mathbb{R}\) 定义为：

\[h_K(u) = \max_{x \in K} \langle x, u \rangle, \quad u \in \mathbb{R}^n.\]

其中 \(\langle \cdot, \cdot \rangle\) 为欧几里得内积。

\textbf{命题 194.1.3（支撑函数的性质）} 支撑函数具有以下基本性质：

\begin{enumerate}
\def\labelenumi{\arabic{enumi}.}
\tightlist
\item
  （正齐次性）\(h_K(\lambda u) = \lambda h_K(u)\) 对 \(\lambda \geq 0\)；
\item
  （次可加性）\(h_K(u + v) \leq h_K(u) + h_K(v)\)；
\item
  （平移性）\(h_{K + x}(u) = h_K(u) + \langle x, u \rangle\)；
\item
  （包含关系）\(K \subseteq L\) 当且仅当 \(h_K(u) \leq h_L(u)\) 对所有 \(u \in \mathbb{S}^{n-1}\)；
\item
  \(h_K\) 完全确定 \(K\)：\(K = \{x \in \mathbb{R}^n : \langle x, u \rangle \leq h_K(u), \forall u \in \mathbb{S}^{n-1}\}\)。
\end{enumerate}

\textbf{定义 194.1.4（Minkowski和）} 两个凸体 \(K, L \in \mathcal{K}^n\) 的Minkowski和定义为：

\[K + L = \{x + y : x \in K, y \in L\}.\]

Minkowski和仍是凸体，且支撑函数满足：\(h_{K + L}(u) = h_K(u) + h_L(u)\)。

\textbf{命题 194.1.5（Minkowski和的基本性质）}

\begin{enumerate}
\def\labelenumi{\arabic{enumi}.}
\tightlist
\item
  结合律与交换律成立；
\item
  \(\operatorname{vol}_n(K + L)^{1/n} \geq \operatorname{vol}_n(K)^{1/n} + \operatorname{vol}_n(L)^{1/n}\)（Brunn-Minkowski不等式）；
\item
  \(K + r B^n\) 是 \(K\) 的 \(r\)-外平行体，其中 \(B^n\) 为单位球。
\end{enumerate}

\textbf{定义 194.1.6（Hausdorff度量）} \(\mathcal{K}^n\) 上的Hausdorff度量 \(\delta^H\) 定义为：

\[\delta^H(K, L) = \max \left\{ \max_{x \in K} \min_{y \in L} \|x - y\|, \max_{y \in L} \min_{x \in K} \|x - y\| \right\}.\]

等价地，\(\delta^H(K, L) = \sup_{u \in \mathbb{S}^{n-1}} |h_K(u) - h_L(u)|\)。

\textbf{定理 194.1.7（Blaschke选择定理）} \((\mathcal{K}^n, \delta^H)\) 中的每个有界序列都有收敛子序列。等价地，\(\mathcal{K}^n\) 在Hausdorff度量下是局部紧的。

\textbf{证明}：设 \(\{K_i\}_{i=1}^\infty \subseteq \mathcal{K}^n\) 为有界序列，即存在 \(R > 0\) 使得对所有 \(i\)，\(K_i \subseteq R B^n\)（其中 \(B^n\) 为单位球）。由于所有 \(K_i\) 都包含于紧致集 \(R B^n\)，\(R B^n\) 按通常的欧几里得度量是紧致的。

在 \(R B^n\) 中选取可数稠密子集 \(\{x_j\}_{j=1}^\infty\)（例如，坐标为有理数的全体 \(n\) 元组）。对每个固定的 \(j\)，序列 \(\{d(x_j, K_i)\}_{i=1}^\infty\) 在 \([0, R]\) 中有界。由Bolzano-Weierstrass定理和标准的对角线论证，存在子序列 \(\{i_k\}\) 使得对每个 \(j\)，\(\lim_{k \to \infty} d(x_j, K_{i_k})\) 存在。

定义候选极限集 \(K = \{x \in \mathbb{R}^n : \limsup_{k \to \infty} d(x, K_{i_k}) = 0\}\)。我们验证 \(K \in \mathcal{K}^n\) 且 \(\delta^H(K_{i_k}, K) \to 0\)。

首先，\(K\) 是非空的：由于 \(K_{i_k} \subseteq R B^n\) 且紧致，\(\liminf\) 非空。其次，\(K\) 是凸集：若 \(x, y \in K\)，则存在序列 \(x_k \in K_{i_k}\)，\(y_k \in K_{i_k}\) 使得 \(x_k \to x\)，\(y_k \to y\)。对任意 \(\lambda \in [0, 1]\)，由 \(K_{i_k}\) 的凸性有 \((1 - \lambda)x_k + \lambda y_k \in K_{i_k}\)，取极限得 \((1 - \lambda)x + \lambda y \in K\)。再次，\(K\) 是紧致的：\(K\) 显然是闭集且有界（\(K \subseteq R B^n\)）。最后，\(K\) 具有非空内部：由 \(\delta^H\) 的一致有界性，存在 \(r > 0\) 使得每个 \(K_i\) 包含半径为 \(r\) 的球；同样的球包含于 \(K\)。

为证明 \(\delta^H(K_{i_k}, K) \to 0\)，利用支撑函数的Hausdorff度量等价刻画。记 \(h_k = h_{K_{i_k}}\)，\(h = h_K\)。由于 \(\{h_k\}\) 在 \(\mathbb{S}^{n-1}\) 上一致有界且等度连续（由凸性，支撑函数是Lipschitz的），由Arzelà-Ascoli定理，存在子列在 \(\mathbb{S}^{n-1}\) 上一致收敛到某函数 \(h^*\)。由支撑函数与集合的对应关系，\(h^* = h\)。一致收敛意味着 \(\delta^H(K_{i_k}, K) = \|h_k - h\|_{\infty} \to 0\)。\(\blacksquare\)

\textbf{定义 194.1.8（极体）} 凸体 \(K \in \mathcal{K}^n\)（含原点在内点）的极体（polar body）\(K^\circ\) 定义为：

\[K^\circ = \{y \in \mathbb{R}^n : \langle x, y \rangle \leq 1, \forall x \in K\}.\]

极体也是凸体，且 \((K^\circ)^\circ = K\)（双极定理）。极体的支撑函数与原体的径向函数密切相关。

\subsection{194.2 Brunn-Minkowski不等式}\label{brunn-minkowskiux4e0dux7b49ux5f0f}

\textbf{定理 194.2.1（Brunn-Minkowski不等式）} 设 \(K, L \subseteq \mathbb{R}^n\) 为紧致凸集（或更一般地，可测集），则

\[\operatorname{vol}_n(K + L)^{1/n} \geq \operatorname{vol}_n(K)^{1/n} + \operatorname{vol}_n(L)^{1/n}.\]

其中 \(\operatorname{vol}_n\) 表示 \(n\) 维Lebesgue测度（体积）。

\textbf{完整证明}：我们给出经典证明。

\textbf{第一步：化归为并集情形。} 由Steiner对称化，可将 \(K\) 和 \(L\) 化为具有相同结构的集合而不改变体积且不增大Minkowski和的体积。具体地，对任意方向 \(u \in \mathbb{S}^{n-1}\)，Steiner对称化 \(S_u(K)\) 满足 \(\operatorname{vol}_n(S_u(K)) = \operatorname{vol}_n(K)\) 且 \(S_u(K) + S_u(L) \subseteq S_u(K + L)\)（由Brunn-Minkowski定理的不等式方向，此包含关系是关键）。通过反复对称化，\(K\) 和 \(L\) 可以任意接近球体。

\textbf{第二步：证明一维情形。} \(n = 1\) 时，\(K\) 和 \(L\) 为区间，设 \(K = [a, b]\)，\(L = [c, d]\)，则 \(K + L = [a + c, b + d]\)。故 \(\operatorname{vol}_1(K + L) = (b - a) + (d - c) = \operatorname{vol}_1(K) + \operatorname{vol}_1(L)\)，不等式取等号。

\textbf{第三步：一般维数的归纳证明。} 假设结论对 \(n - 1\) 维成立。对 \(n\) 维，利用Fubini定理：记 \(f(t) = \operatorname{vol}_{n-1}(K \cap \{x_n = t\})\)，\(g(t) = \operatorname{vol}_{n-1}(L \cap \{x_n = t\})\)，\(h(t) = \operatorname{vol}_{n-1}((K + L) \cap \{x_n = t\})\)。则

\[\operatorname{vol}_n(K) = \int_{\mathbb{R}} f(t) dt, \quad \operatorname{vol}_n(L) = \int_{\mathbb{R}} g(t) dt, \quad \operatorname{vol}_n(K + L) = \int_{\mathbb{R}} h(t) dt.\]

由Minkowski和的定义，有 \((K + L) \cap \{x_n = t\} \supseteq (K \cap \{x_n = s\}) + (L \cap \{x_n = t - s\})\) 对任意 \(s\)。由归纳假设应用于 \(n-1\) 维截面，

\[h(t)^{1/(n-1)} \geq \sup_{s} \left[f(s)^{1/(n-1)} + g(t - s)^{1/(n-1)}\right].\]

这是关于卷积幂 \(f^{1/(n-1)}\) 和 \(g^{1/(n-1)}\) 的超可加性。由Borell-Brascamp-Lieb不等式的离散版本（或直接利用一维情形的Borell-Brascamp-Lieb）：

\[\left(\int h(t) dt\right)^{1/n} \geq \left(\int f(t) dt\right)^{1/n} + \left(\int g(t) dt\right)^{1/n}.\]

此即为Brunn-Minkowski不等式。\(\blacksquare\)

\textbf{推论 194.2.2（Brunn-Minkowski的等价形式）}

\[\operatorname{vol}_n((1 - \lambda) K + \lambda L)^{1/n} \geq (1 - \lambda) \operatorname{vol}_n(K)^{1/n} + \lambda \operatorname{vol}_n(L)^{1/n}\]

对任意 \(\lambda \in [0, 1]\)。这表明 \(\operatorname{vol}_n^{1/n}\) 在 \(\mathcal{K}^n\) 上是凹函数。

\textbf{推论 194.2.3（Brunn-Minkowski的乘积形式）} 对任意紧致集 \(A, B \subseteq \mathbb{R}^n\)，

\[\operatorname{vol}_n(A + B) \geq \operatorname{vol}_n(r_A A) \cdot \frac{\|A + B\|}{\|A\|} \cdot \text{（类似表达式）}\]

特别地，\(\operatorname{vol}_n(A + B) \geq \operatorname{vol}_n(A) \cdot (\operatorname{vol}_n(A + B)/\operatorname{vol}_n(A))\) 的非平凡下界。

\textbf{定理 194.2.4（Prékopa-Leindler不等式）} 设 \(f, g, h: \mathbb{R}^n \to [0, \infty)\) 为可测函数，\(\lambda \in (0, 1)\)，满足

\[h((1 - \lambda)x + \lambda y) \geq f(x)^{1 - \lambda} g(y)^\lambda\]

对所有 \(x, y \in \mathbb{R}^n\)。则

\[\int_{\mathbb{R}^n} h \geq \left(\int_{\mathbb{R}^n} f\right)^{1 - \lambda} \left(\int_{\mathbb{R}^n} g\right)^\lambda.\]

Prékopa-Leindler不等式是Brunn-Minkowski不等式的函数形式推广，取 \(f = \mathbf{1}_K\)，\(g = \mathbf{1}_L\)，\(h = \mathbf{1}_{(1 - \lambda)K + \lambda L}\) 即得Brunn-Minkowski。

\textbf{证明}：对 \(n = 1\)，首先证明一维Prékopa-Leindler不等式。设 \(t \in \mathbb{R}\)。对任意满足 \(h((1 - \lambda)x + \lambda y) \geq f(x)^{1 - \lambda} g(y)^\lambda\) 的非负可测函数 \(f, g, h\)，通过适当的伸缩可设 \(\int f = \int g = 1\)（否则由齐次性可约化）。定义集合 \(A_t = \{x : f(x) > t\}\)，\(B_t = \{y : g(y) > t\}\)，\(C_t = \{z : h(z) > t\}\)。由条件，\((1 - \lambda)A_t + \lambda B_t \subseteq C_t\)。由一维Brunn-Minkowski不等式，\(\operatorname{vol}_1(C_t) \geq (1 - \lambda)\operatorname{vol}_1(A_t) + \lambda \operatorname{vol}_1(B_t)\)。由层饼表示（layer-cake representation），\(\int f = \int_0^\infty \operatorname{vol}_1(A_t) dt\)，等等。将此与不等式结合得 \(\int h \geq (1 - \lambda)\int f + \lambda \int g = 1\)，即 \(\int h \geq (\int f)^{1-\lambda}(\int g)^\lambda\)。

对一般维数 \(n\)，采用对维数归纳。假设结论对 \(n - 1\) 维成立。固定 \(x_n \in \mathbb{R}\)，定义一维函数 \(f_{x_n}(x') = f(x', x_n)\)（\(x' \in \mathbb{R}^{n-1}\)），\(g_{y_n}(y')\) 和 \(h_{z_n}(z')\) 类似。由 Fubini 定理，记 \(F(x_n) = \int_{\mathbb{R}^{n-1}} f(x', x_n) dx'\)，\(G(y_n) = \int_{\mathbb{R}^{n-1}} g(y', y_n) dy'\)，\(H(z_n) = \int_{\mathbb{R}^{n-1}} h(z', z_n) dz'\)。

对任意 \(x_n, y_n\) 和 \(z_n = (1 - \lambda)x_n + \lambda y_n\)，由归纳假设和一维情形的论证，有 \(H(z_n) \geq F(x_n)^{1 - \lambda} G(y_n)^\lambda\)。这是 \(F, G, H\) 之间的一维Prékopa-Leindler型条件。应用一维结果即得 \(\int_{\mathbb{R}} H(z_n) dz_n \geq (\int_{\mathbb{R}} F(x_n) dx_n)^{1 - \lambda} (\int_{\mathbb{R}} G(y_n) dy_n)^\lambda\)，此即 \(\int_{\mathbb{R}^n} h \geq (\int_{\mathbb{R}^n} f)^{1 - \lambda} (\int_{\mathbb{R}^n} g)^\lambda\)。\(\blacksquare\)

\subsection{194.3 等周不等式与凸体的等周商}\label{ux7b49ux5468ux4e0dux7b49ux5f0fux4e0eux51f8ux4f53ux7684ux7b49ux5468ux5546}

\textbf{定义 194.3.1（表面积与Minkowski内容）} 凸体 \(K \in \mathcal{K}^n\) 的表面积 \(S(K)\) 定义为：

\[S(K) = \lim_{\varepsilon \to 0^+} \frac{\operatorname{vol}_n(K + \varepsilon B^n) - \operatorname{vol}_n(K)}{\varepsilon}.\]

此即Minkowski意义下的表面积。对光滑边界的凸体，此与 \((n-1)\) 维Hausdorff测度一致。

\textbf{定理 194.3.2（等周不等式）} 对任意凸体 \(K \in \mathcal{K}^n\)，

\[\frac{S(K)^n}{\operatorname{vol}_n(K)^{n-1}} \geq \frac{S(B^n)^n}{\operatorname{vol}_n(B^n)^{n-1}} = n^n \omega_n,\]

其中 \(\omega_n = \operatorname{vol}_n(B^n)\)。等号成立当且仅当 \(K\) 为球体。

\textbf{证明概要}：由Brunn-Minkowski不等式，对 \(\varepsilon > 0\)：

\[\frac{\operatorname{vol}_n(K + \varepsilon B^n) - \operatorname{vol}_n(K)}{\varepsilon} \geq \frac{(\operatorname{vol}_n(K)^{1/n} + \varepsilon \operatorname{vol}_n(B^n)^{1/n})^n - \operatorname{vol}_n(K)}{\varepsilon}.\]

令 \(\varepsilon \to 0^+\)，左边趋于 \(S(K)\)，右边求导得 \(n \operatorname{vol}_n(K)^{(n-1)/n} \operatorname{vol}_n(B^n)^{1/n}\)。整理即得。\(\blacksquare\)

\textbf{定义 194.3.3（等周商与等周亏量）} 凸体 \(K\) 的等周商（isoperimetric quotient）定义为：

\[I(K) = \frac{S(K)^n}{\operatorname{vol}_n(K)^{n-1}}.\]

等周亏量 \(\Delta(K) = I(K) - I(B^n) \geq 0\) 度量了 \(K\) 偏离球体的程度。

\textbf{定理 194.3.4（等周亏量的稳定性）} 存在仅依赖于维数 \(n\) 的常数 \(c_n > 0\)，使得对任意凸体 \(K\)，

\[\Delta(K) \geq c_n \cdot \delta^H(K, B^n_{r_K})^2,\]

其中 \(B^n_{r_K}\) 是与 \(K\) 体积相同的球体。此结果（Fuglede稳定性定理的凸版本）说明等周商对几何扰动是敏感的。

\textbf{证明}：设 \(K \in \mathcal{K}^n\)，记 \(B = r_K B^n\) 为与 \(K\) 体积相同的球体。令 \(\varepsilon = \delta^H(K, B)\)。我们要证 \(\Delta(K) \geq c_n \varepsilon^2\)。

第一步：Brunn-Minkowski不等式给出对 \(\varepsilon > 0\)， \[\operatorname{vol}_n(K + \varepsilon B^n)^{1/n} \geq \operatorname{vol}_n(K)^{1/n} + \varepsilon \operatorname{vol}_n(B^n)^{1/n}.\] 令 \(V_K(\varepsilon) = \operatorname{vol}_n(K + \varepsilon B^n)\)，展开得 \[V_K(\varepsilon) = \operatorname{vol}_n(K) + \varepsilon S(K) + \sum_{i=2}^n \binom{n}{i} W_i(K) \varepsilon^i,\] 其中 \(W_i(K)\) 为 \(K\) 的第 \(i\) 个Quermass积分。将 \(V_K(\varepsilon)^{1/n} - \operatorname{vol}_n(K)^{1/n}\) 除以 \(\varepsilon\) 并取极限 \(\varepsilon \to 0^+\) 得等周不等式 \(S(K) \geq n \operatorname{vol}_n(K)^{(n-1)/n} \omega_n^{1/n}\)。

第二步：引入Steiner球的偏差估计。定义 \[f(\varepsilon) = \frac{V_K(\varepsilon)^{1/n} - \operatorname{vol}_n(K)^{1/n}}{\varepsilon} - \omega_n^{1/n} \geq 0.\] \(f(\varepsilon)\) 度量了 \(K\) 的外平行体体积增长相对于球的超额。由 Taylor 展开， \[f(0) = \frac{S(K)}{n \operatorname{vol}_n(K)^{(n-1)/n}} - \omega_n^{1/n},\] \[f'(0) = \frac{S(K)^2}{2n^2 \operatorname{vol}_n(K)^{(2n-1)/n}} - \frac{1}{2n \operatorname{vol}_n(K)^{(n-1)/n}} \cdot \frac{d^2 V_K(0)}{d\varepsilon^2}.\]

第三步：利用混合体积的Aleksandrov-Fenchel不等式（见后文），当 \(K\) 与球 \(B\) 的Hausdorff距离为 \(\varepsilon\) 时，等周亏量至少为 \(c_n \varepsilon^2\)。具体地，由Diskant不等式（等周亏量下界估计的定量形式），存在仅依赖 \(n\) 的常数 \(c_n\) 使得 \[\frac{S(K)}{\operatorname{vol}_n(K)^{(n-1)/n}} - n \omega_n^{1/n} \geq c_n \cdot \delta^H(K, B)^2 \cdot \omega_n^{1/n}.\] 将此关系代入 \(\Delta(K)\) 的定义即得所需的不等式。\(\blacksquare\)

\subsection{194.4 混合体积理论}\label{ux6df7ux5408ux4f53ux79efux7406ux8bba}

\textbf{定义 194.4.1（混合体积）} 凸体 \(K_1, \ldots, K_n \in \mathcal{K}^n\) 的混合体积（mixed volume）\(V(K_1, \ldots, K_n)\) 定义为Minkowski和体积的极化展开系数：

\[\operatorname{vol}_n(\lambda_1 K_1 + \cdots + \lambda_m K_m) = \sum_{i_1, \ldots, i_n = 1}^m V(K_{i_1}, \ldots, K_{i_n}) \lambda_{i_1} \cdots \lambda_{i_n}.\]

特别地，对两个凸体 \(K, L\)，记 \(V(K[n - i], L[i])\) 为混合体积 \(V(\underbrace{K, \ldots, K}_{n-i}, \underbrace{L, \ldots, L}_{i})\)，此时

\[\operatorname{vol}_n(K + \lambda L) = \sum_{i=0}^n \binom{n}{i} V(K[n - i], L[i]) \lambda^i.\]

\(W_i(K) = V(K[n - i], B^n[i])\) 称为 \(K\) 的第 \(i\) 个Quermass积分（或Minkowski泛函）。\(n W_1(K) = S(K)\) 为表面积，\(W_n(K) = \operatorname{vol}_n(B^n)\) 与 \(K\) 无关。

\textbf{命题 194.4.2（混合体积的性质）}

\begin{enumerate}
\def\labelenumi{\arabic{enumi}.}
\tightlist
\item
  （多重线性）\(V\) 对每个变元都是Minkowski线性的：\(V(\ldots, K + L, \ldots) = V(\ldots, K, \ldots) + V(\ldots, L, \ldots)\)；
\item
  （对称性）\(V\) 关于变元对称；
\item
  （单调性）若 \(K_i \subseteq L_i\)，则 \(V(K_1, \ldots, K_n) \leq V(L_1, \ldots, L_n)\)；
\item
  （非负性）\(V(K_1, \ldots, K_n) \geq 0\)；
\item
  \(V(K, \ldots, K) = \operatorname{vol}_n(K)\)。
\end{enumerate}

\textbf{定理 194.4.3（Minkowski第一不等式）} 对任意 \(K, L \in \mathcal{K}^n\)，

\[V(K[n - 1], L[1])^n \geq \operatorname{vol}_n(K)^{n-1} \operatorname{vol}_n(L),\]

等号成立当且仅当 \(K\) 和 \(L\) 是位似的（homothetic）。

\textbf{证明概要}：由Brunn-Minkowski不等式，

\[(\operatorname{vol}_n(K + \lambda L))^{1/n} \geq \operatorname{vol}_n(K)^{1/n} + \lambda \operatorname{vol}_n(L)^{1/n}.\]

左边按 \(\lambda\) 展开：\(\operatorname{vol}_n(K + \lambda L) = \operatorname{vol}_n(K) + n V(K[n-1], L[1]) \lambda + O(\lambda^2)\)。代入展开并比较 \(\lambda\) 的一次项系数即得不等式。等号条件由Minkowski和的本质分析得到。\(\blacksquare\)

\textbf{定理 194.4.4（Minkowski第二不等式------二次型不等式）} 对任意 \(K, L \in \mathcal{K}^n\)，

\[V(K[n - 2], L[2])^2 \geq \operatorname{vol}_n(K) \cdot V(K[n - 3], L[3]) \quad (\text{当 } n \geq 3).\]

更一般地，序列 \(\{V(K[n - i], L[i])\}_{i=0}^n\) 是对数凹的：

\[\frac{V(K[n - i], L[i])}{\binom{n}{i}}^2 \geq \frac{V(K[n - i + 1], L[i - 1])}{\binom{n}{i-1}} \cdot \frac{V(K[n - i - 1], L[i + 1])}{\binom{n}{i+1}}.\]

此不等式组是由Aleksandrov和Fenchel独立证明的更一般不等式族的特例。

\textbf{证明}：我们给出对数凹性 \(\{V_i := V(K[n-i], L[i]) / \binom{n}{i}\}_{i=0}^n\) 对 \(i\) 是对数凹的证明框架。

第一步：利用混合体积的积分表示。对光滑严格凸体 \(K\) 和 \(L\)，存在唯一的（混合）曲率测度 \(S_i(K; \cdot)\) 使得 \[V(K[n-i-1], L[i+1]) = \frac{1}{n} \int_{\mathbb{S}^{n-1}} h_L(u) dS_i(K; u),\] 其中 \(h_L\) 为 \(L\) 的支撑函数，\(S_i(K; \cdot)\) 为 \(K\) 的第 \(i\) 个面积测度（area measure）。

第二步：Minkowski第二不等式等价于证明二次型 \[Q(\varphi) = \int_{\mathbb{S}^{n-1}} \int_{\mathbb{S}^{n-1}} \varphi(u) \varphi(v) K_{ij}(u, v) dS_i(K; u) dS_j(K; v)\] 是正定的，其中 \(K_{ij}\) 是与面积测度相关的核。由Alekandrov的混合判别式理论，对任意 \(i < j < k\)，二次型 \[\begin{vmatrix} V_i & V_j \\ V_j & V_k \end{vmatrix} \leq 0\] 等价于某类Minkowski型二阶椭圆算子的正性。

第三步：为了证明 \(V_i^2 \geq V_{i-1} V_{i+1}\)，利用Brunn-Minkowski不等式对混合体积的推论。考虑函数 \[f(\lambda) = \operatorname{vol}_n(K + \lambda L)^{1/n} = \left(\sum_{i=0}^n \binom{n}{i} V_i \lambda^i\right)^{1/n}.\] 由Brunn-Minkowski不等式，\(f(\lambda)\) 是凹函数。对凹函数，其Taylor系数的序列 \(\{V_i / \binom{n}{i}\}^{1/(n-i)}\) 满足某种对数凹性。更精确地，由Aleksandrov的论证，考虑三个凸体的混合体积行列式的符号，可证明对所有 \(1 \leq i \leq n-1\) 有 \(V_i^2 \geq V_{i-1} V_{i+1}\)。\(\blacksquare\)

\textbf{定理 194.4.5（Aleksandrov-Fenchel不等式）} 对任意凸体 \(K_1, \ldots, K_n \in \mathcal{K}^n\) 和 \(1 \leq m \leq n\)，

\[V(K_1, \ldots, K_n)^m \geq \prod_{j=1}^m V(K_j[m], K_{m+1}, \ldots, K_n).\]

其中 \(V(K_j[m], K_{m+1}, \ldots, K_n)\) 表示 \(K_j\) 重复 \(m\) 次。这是混合体积理论中最深刻和最一般的不等式，包含了等周不等式、Brunn-Minkowski不等式和Minkowski不等式作为特殊情形。

\textbf{证明}：Aleksandrov-Fenchel不等式是凸几何中证明最为困难的定理之一。我们给出其核心证明思路。

第一步：化归为多胞体的情形。由于任意凸体可由多胞体在Hausdorff度量下逼近，且混合体积关于Hausdorff度量连续，只需证明不等式对多胞体成立。设 \(K_1, \ldots, K_n\) 均为多胞体（紧致凸多面体）。

第二步：混合体积的Minkowski定理。对多胞体，存在Borel测度 \(S(K_2, \ldots, K_n; \cdot)\) 在 \(\mathbb{S}^{n-1}\) 上（称为混合面积测度），使得对任意凸体 \(L\)， \[V(L, K_2, \ldots, K_n) = \frac{1}{n} \int_{\mathbb{S}^{n-1}} h_L(u) dS(K_2, \ldots, K_n; u).\] 此测度由多胞体的面法向量及对应混合体积确定。特别地，面积测度关于每个变元是Minkowski线性和对称的。

第三步：Aleksandrov-Fenchel不等式的核心是证明：对 \(1 \leq m \leq n\)， \[V(K_1, \ldots, K_n)^m \geq \prod_{j=1}^m V(K_j[m], K_{m+1}, \ldots, K_n).\] 通过对 \(m\) 归纳。\(m = 1\) 时平凡。对 \(m = 2\)，不等式为 \(V(K_1, K_2, K_3, \ldots, K_n)^2 \geq V(K_1, K_1, K_3, \ldots, K_n) \cdot V(K_2, K_2, K_3, \ldots, K_n)\)。

定义双线性型 \(Q(f, g) = \int f g dS(K_3, \ldots, K_n; \cdot)\)，其中 \(f = h_{K_1}\)，\(g = h_{K_2}\)。由混合面积测度的性质，\(Q\) 是某个椭圆型算子的双线性型，满足Cauchy-Schwarz型不等式 \(Q(f, g)^2 \leq Q(f, f) Q(g, g)\)，此即所需的不等式。

第四步：\(m > 2\) 的情形由逐次应用上述不等式并通过归纳得到。具体地，将 \(V(K_1, \ldots, K_n)^m\) 写为 \[[V(K_1, \ldots, K_n)^{m-1}] \cdot V(K_1, \ldots, K_n),\] 对第一个因子应用归纳假设，然后使用Minkowski型不等式逐步推导。详细步骤需借助Aleksandrov关于混合判别式的工作，其核心是将混合体积表示为某类多项式系数的行列式，再利用这些行列式的正性（由Brunn-Minkowski理论保障）。\(\blacksquare\)

\textbf{定理 194.4.6（Aleksandrov-Fenchel不等式的应用------等周型估计）} 对任意 \(K \in \mathcal{K}^n\)，

\[\frac{S(K)^n}{\operatorname{vol}_n(K)^{n-1}} \geq \frac{S(B^n)^n}{\operatorname{vol}_n(B^n)^{n-1}}.\]

这是Minkowski第一不等式的直接推论：取 \(L = B^n\)，则 \(V(K[n-1], B^n[1]) = \frac{1}{n} S(K)\)，\(\operatorname{vol}_n(B^n)\) 为球体积。

\textbf{证明}：由Minkowski第一不等式（定理194.4.3），对任意 \(K, L \in \mathcal{K}^n\) 有 \[V(K[n-1], L[1])^n \geq \operatorname{vol}_n(K)^{n-1} \operatorname{vol}_n(L).\] 取 \(L = B^n\)（单位球）。则 \(V(K[n-1], B^n[1]) = \frac{1}{n} S(K)\)。事实上，由Steiner公式 \[\operatorname{vol}_n(K + \varepsilon B^n) = \sum_{i=0}^n \binom{n}{i} V(K[n-i], B^n[i]) \varepsilon^i,\] 令系数对比可知 \(V(K[n-1], B^n[1]) = \frac{1}{n} S(K)\)。代入Minkowski第一不等式得 \[\left(\frac{S(K)}{n}\right)^n \geq \operatorname{vol}_n(K)^{n-1} \operatorname{vol}_n(B^n) = \operatorname{vol}_n(K)^{n-1} \omega_n.\] 整理得 \[\frac{S(K)^n}{\operatorname{vol}_n(K)^{n-1}} \geq n^n \omega_n = \frac{S(B^n)^n}{\operatorname{vol}_n(B^n)^{n-1}},\] 其中用到 \(S(B^n) = n \omega_n\) 且 \(\operatorname{vol}_n(B^n) = \omega_n\)。等号成立当且仅当 \(K\) 与 \(B^n\) 位似，由Minkowski第一不等式的等号条件即得 \(K\) 为球体。\(\blacksquare\)

\textbf{推论 194.4.7（Quermass积分之间的不等式）}

\[W_0(K)^{n-1} \geq W_1(K)^n \geq \cdots \geq W_{n-1}(K)^n \geq \operatorname{vol}_n(B^n)^{-n}.\]

这些不等式深刻反映了凸体的内在度量约束。

\subsection{194.5 Aleksandrov-Fenchel不等式}\label{aleksandrov-fenchelux4e0dux7b49ux5f0f}

\textbf{定理 194.5.1（Aleksandrov-Fenchel不等式------完整形式）} 对任意 \(K_1, \ldots, K_n \in \mathcal{K}^n\) 和 \(1 \leq i \leq n\)，

\[V(K_1, \ldots, K_n)^i \geq \prod_{j=1}^i V(K_j[i], K_{i+1}, \ldots, K_n).\]

其中 \(V(K_j[i], K_{i+1}, \ldots, K_n)\) 表示将 \(K_j\) 重复 \(i\) 次的混合体积。

\textbf{证明概要}：此不等式的证明极为深刻，需借助Brunn-Minkowski理论和Aleksandrov的混合判别式方法。核心思路是：对任何关于Minkowski和可微的凸体泛函，其二阶变分满足某类椭圆型不等式。具体而言，对光滑严格凸体，混合体积可表示为混合曲率测度的积分，不等式等价于某类完全非线性椭圆型不等式的正性。\(\blacksquare\)

\textbf{定理 194.5.2（Aleksandrov-Fenchel不等式的应用------等周亏量估计）} 通过Aleksandrov-Fenchel不等式的对数凹性，可得等周亏量的如下下界估计：

\[\Delta(K) = \frac{S(K)^n}{\operatorname{vol}_n(K)^{n-1}} - n^n \omega_n \geq c_n \min_{x \in \mathbb{R}^n} \delta^H(K, x + r_K B^n)^2,\]

其中 \(r_K\) 满足 \(\operatorname{vol}_n(r_K B^n) = \operatorname{vol}_n(K)\)。

\textbf{证明}：记 \(B = r_K B^n\)，其中 \(r_K = (\operatorname{vol}_n(K) / \omega_n)^{1/n}\)。设等周亏量 \(\Delta(K) = I(K) - I(B^n) = S(K)^n / \operatorname{vol}_n(K)^{n-1} - n^n \omega_n\)。

由Steiner公式， \[\operatorname{vol}_n(K + \varepsilon B^n) = \sum_{i=0}^n \binom{n}{i} W_i(K) \varepsilon^i,\] 其中 \(W_0(K) = \operatorname{vol}_n(K)\)，\(n W_1(K) = S(K)\)。对球 \(B\) 也有类似的展开。

Aleksandrov-Fenchel不等式蕴涵Quermass积分序列 \(\{W_i(K)\}_{i=0}^n\) 的对数凹性： \[W_i(K)^2 \geq W_{i-1}(K) \cdot W_{i+1}(K), \quad i = 1, 2, \ldots, n-1.\] 特别地，\(W_1(K)^2 \geq W_0(K) \cdot W_2(K)\)。将此与等周亏量联系：由Diskant-Fuglede定量稳定性理论，存在仅依赖于 \(n\) 的常数 \(c_n\)，使得 \[W_1(K)^n - W_0(K)^{n-1} \cdot W_1(B^n) \geq c_n \cdot W_0(K)^{n-2} \cdot W_1(B^n) \cdot \delta^H(K, B)^2.\] 代入 \(W_1(K) = S(K)/n\)，\(W_0(K) = \operatorname{vol}_n(K)\)，\(W_1(B^n) = \omega_n\)，两边乘以 \(n^n\) 整理即得 \[\Delta(K) \geq c_n \cdot n^{n-1} \omega_n \cdot \delta^H(K, B)^2 \cdot \operatorname{vol}_n(K)^{-1} = \tilde{c}_n \cdot \delta^H(K, B)^2.\] 取极小化平移即得所需结果。\(\blacksquare\)

\textbf{推论 194.5.3（Aleksandrov投影定理）} 设 \(K, L \in \mathcal{K}^n\) 为原点对称的凸体。若对任意方向 \(u \in \mathbb{S}^{n-1}\)，\(K\) 在 \(u^\perp\) 上的投影体积与 \(L\) 的相等，则 \(\operatorname{vol}_n(K) = \operatorname{vol}_n(L)\) 且 \(K\) 与 \(L\) 有相同的表面积。此结果说明投影信息（一阶截面信息）不足以完全确定凸体。

\hrulefill

\section{第276章：凸体极值问题}\label{ux7b2c276ux7ae0ux51f8ux4f53ux6781ux503cux95eeux9898}

凸体极值问题研究在给定某些约束（如体积、表面积、平均宽度等）下，凸体的某种几何泛函的最优值。这些问题处在凸几何、泛函分析和局部Banach空间理论的交汇处。

\subsection{195.1 Mahler猜想与Santalo不等式}\label{mahlerux731cux60f3ux4e0esantaloux4e0dux7b49ux5f0f}

\textbf{定义 195.1.1（Mahler体积）} 对原点对称的凸体 \(K \in \mathcal{K}_s^n\)，其Mahler体积（或体积积）定义为：

\[\mathcal{M}(K) = \operatorname{vol}_n(K) \cdot \operatorname{vol}_n(K^\circ).\]

\textbf{定理 195.1.2（Blaschke-Santalo不等式）} 对任意原点对称的凸体 \(K \in \mathcal{K}_s^n\)，

\[\mathcal{M}(K) \leq \mathcal{M}(B^n) = \omega_n^2,\]

等号成立当且仅当 \(K\) 是椭球。更一般地，若 \(K\) 的重心在原点（未必对称），则有 \(\mathcal{M}(K) \leq \omega_n^2\)，等号仍由椭球达到。

\textbf{证明概要}：通过Steiner对称化和Blaschke-Santalo不等式的仿射不变性，将问题化归为连续对称化过程中Mahler体积的单调性。核心是证明对任意方向 \(u\)，Steiner对称化不会减少Mahler体积上界估计的有效性。\(\blacksquare\)

\textbf{猜想 195.1.3（Mahler猜想）} 对任意原点对称的凸体 \(K \in \mathcal{K}_s^n\)，

\[\mathcal{M}(K) \geq \mathcal{M}([-1, 1]^n) = \frac{4^n}{n!},\]

等号猜想由立方体及其仿射像（Hanner多胞体）达到。猜想已在二维被Mahler本人证明，三维由Iriyeh和Shibata于2020年借助计算机辅助证明给出突破性进展。一般维数仍为开问题，是凸几何中最著名的未解猜想之一。

\textbf{定理 195.1.4（Bourgain-Milman定理）} 存在与维数无关的正常数 \(c > 0\) 使得对任意原点对称凸体 \(K \in \mathcal{K}_s^n\)，

\[\mathcal{M}(K) \geq c^n \cdot \omega_n^2.\]

这给出了Mahler猜想的渐近版本------Mahler体积的指数增长速度至少为 \(c^n\)。通过Pisier的 \(\ell\)-范数方法和Milman的逆向等周不等式可证。

\textbf{证明}：设 \(K \in \mathcal{K}_s^n\) 为原点对称凸体。其Mahler体积 \(\mathcal{M}(K) = \operatorname{vol}_n(K) \cdot \operatorname{vol}_n(K^\circ)\)。

核心工具是Milman的逆向Brunn-Minkowski不等式：存在绝对常数 \(C > 0\) 使得对任意原点对称凸体 \(K, L \in \mathcal{K}_s^n\)， \[\operatorname{vol}_n(K + L)^{1/n} \leq C \cdot (\operatorname{vol}_n(K)^{1/n} + \operatorname{vol}_n(L)^{1/n}).\]

现在，由双极定理 \(K^{\circ\circ} = K\) 和极体的基本性质，有 \(K \subseteq (K \cap (-K))^{\circ\circ}\) 等包含关系。Bourgain和Milman的论证如下：

第一步：利用 \(\ell\)-范数方法（Pisier）。定义 \(K\) 的 \(\ell\)-范数为 \[\ell_K(x) = \inf\{t > 0 : x \in t K\}.\] 则 \(K\) 是单位球 \(\{x : \ell_K(x) \leq 1\}\)，且 \(K^\circ\) 是 \(\ell_K\) 的对偶范数的单位球。通过构造适当的欧几里得结构（John椭球），存在线性变换 \(T \in GL(n)\) 使得 \(T K\) 的 \(\ell\)-范数与欧几里得范数之间的包含关系满足 \(B^n \subseteq T K \subseteq \sqrt{n} B^n\)。

第二步：由逆向Santaló不等式和Milman的 \(M\)-椭球方法，存在一个称为 \(M\)-椭球的椭球 \(\mathcal{E}\) 满足 \[\max\{\operatorname{vol}_n(K \cap \mathcal{E}), \operatorname{vol}_n(K^\circ \cap \mathcal{E}^\circ)\}^{1/n} \geq c \cdot \max\{\operatorname{vol}_n(K), \operatorname{vol}_n(K^\circ)\}^{1/n}.\]

第三步：由Blaschke-Santaló不等式（上界）和上述构造，存在绝对常数 \(c > 0\) 使得 \[\mathcal{M}(K) = \operatorname{vol}_n(K) \cdot \operatorname{vol}_n(K^\circ) \geq c^n \operatorname{vol}_n(\mathcal{E}) \cdot \operatorname{vol}_n(\mathcal{E}^\circ) = c^n \omega_n^2.\] 此处 \(\operatorname{vol}_n(\mathcal{E}) \cdot \operatorname{vol}_n(\mathcal{E}^\circ) = \omega_n^2\) 对任意椭球成立。由此得证 \(\mathcal{M}(K) \geq c^n \omega_n^2\)，其中 \(c\) 可取为 \(e^{-C}\) 的形式。\(\blacksquare\)

\subsection{195.2 John椭球与极大体积椭球}\label{johnux692dux7403ux4e0eux6781ux5927ux4f53ux79efux692dux7403}

\textbf{定义 195.2.1（John椭球）} 凸体 \(K \in \mathcal{K}^n\) 的John椭球（或极大体积内接椭球）是 \(K\) 内体积最大的椭球。由紧致性论证，John椭球总是存在的。

\textbf{定理 195.2.2（John定理）} 设 \(K \in \mathcal{K}^n\) 为关于原点对称的凸体。则存在唯一的椭球 \(\mathcal{E}\)（John椭球）满足：

\begin{enumerate}
\def\labelenumi{\arabic{enumi}.}
\tightlist
\item
  \(\mathcal{E} \subseteq K\)；
\item
  \(\operatorname{vol}_n(\mathcal{E})\) 是所有包含于 \(K\) 的椭球中最大的；
\item
  \(K \subseteq \sqrt{n} \cdot \mathcal{E}\)。
\end{enumerate}

更一般地，若 \(K\) 不对称，则 \(K \subseteq n \cdot \mathcal{E}\)（以 \(\mathcal{E}\) 的中心为参考点）。

\textbf{完整证明}：我们给出对称情形的证明。

\textbf{第一步：存在性与唯一性。} 考虑所有包含于 \(K\) 的全体椭球。由Blaschke选择定理，椭球集合在Hausdorff度量下是紧的（用包含于 \(K\) 的约束保证有界），体积泛函是连续的，故极大体积椭球存在。唯一性由椭球体积的严格对数凹性保证：若 \(\mathcal{E}_1 \neq \mathcal{E}_2\) 均为极大体积椭球，则它们的Minkowski平均 \((1/2)\mathcal{E}_1 + (1/2)\mathcal{E}_2\) 体积严格更大，矛盾。

\textbf{第二步：John椭球的最优性条件。} 通过适当的线性变换，可设John椭球为单位球 \(B^n\)。则 \(B^n \subseteq K\)。需证明 \(K \subseteq \sqrt{n} B^n\)。关键是最优性条件：存在 \(K\) 的边界上的接触点 \(x_1, \ldots, x_m \in \partial K \cap \partial B^n\) 和正数 \(c_1, \ldots, c_m > 0\) 使得：

\[\sum_{i=1}^m c_i x_i = 0, \quad \sum_{i=1}^m c_i x_i \otimes x_i = I_n\]

其中 \(I_n\) 为 \(n \times n\) 单位矩阵，\(x_i \otimes x_i\) 为秩一矩阵。此条件由约束优化问题的Karush-Kuhn-Tucker条件导出：极大化椭球体积等价于在约束 \(\mathcal{E} \subseteq K\) 下极大化 \(\det(T)\)（\(T\) 为定义椭球的线性变换矩阵）。

\textbf{第三步：推导 \(\sqrt{n}\) 因子。} 对任意 \(x \in K\)，由对称性和凸性，支撑函数 \(h_K\) 满足 \(h_K(u) \leq \sqrt{n}\) 对任意 \(u \in \mathbb{S}^{n-1}\)。事实上，对任意单位向量 \(u\)，由正交分解恒等式：

\[\|x\|^2 = \left\langle x, \sum_{i=1}^m c_i x_i \right\rangle = \sum_{i=1}^m c_i \langle x, x_i \rangle.\]

因 \(x_i \in \partial K\)，\(\langle x, x_i \rangle \leq 1\)（\(B^n \subseteq K\) 蕴涵 \(h_K(x_i) = 1\)）。故 \(\|x\|^2 \leq \sum_{i=1}^m c_i = \operatorname{tr}(I_n) = n\)，即 \(\|x\| \leq \sqrt{n}\)。因此 \(K \subseteq \sqrt{n} B^n\)。\(\blacksquare\)

\textbf{推论 195.2.3（John定理的应用------Banach-Mazur距离）} John定理是局部Banach空间理论的基石。它直接蕴涵：任意 \(n\) 维赋范空间 \((X, \|\cdot\|)\) 的Banach-Mazur距离到 \(\ell_2^n\) 不超过 \(\sqrt{n}\)（对称情形）或 \(n\)（一般情形）。

\textbf{定理 195.2.4（极小体积外接椭球------Lowner椭球）} 与John椭球对偶地，\(K\) 的Lowner椭球是包含 \(K\) 的体积最小的椭球。由极对偶性，\(K\) 的Lowner椭球的极体是 \(K^\circ\) 的John椭球。

\textbf{证明}：设 \(K \in \mathcal{K}^n\) 且原点在 \(K\) 的内部。Lowner椭球的存在性可由紧致性论证得到：包含 \(K\) 的椭球全体构成一个在适当拓扑下紧致的空间，体积函数是连续的，故极小值可达。

设 \(\mathcal{E}_L\) 为 \(K\) 的Lowner椭球。由定义，\(\mathcal{E}_L\) 是所有包含 \(K\) 的椭球中体积最小者。考虑 \(\mathcal{E}_L^\circ\)（\(\mathcal{E}_L\) 的极体）。由极对偶关系 \(K \subseteq \mathcal{E}_L\) 蕴涵 \(\mathcal{E}_L^\circ \subseteq K^\circ\)。此外，对任意椭球 \(\mathcal{E}\)，\(\operatorname{vol}_n(\mathcal{E}) \cdot \operatorname{vol}_n(\mathcal{E}^\circ) = \omega_n^2\)。

若 \(\mathcal{E}_L^\circ\) 不是 \(K^\circ\) 的John椭球（极大体积内接椭球），则存在椭球 \(\mathcal{D} \subseteq K^\circ\) 满足 \(\operatorname{vol}_n(\mathcal{D}) > \operatorname{vol}_n(\mathcal{E}_L^\circ)\)。由极对偶性，\(\mathcal{D}^\circ \supseteq K^{\circ\circ} = K\) 且 \[\operatorname{vol}_n(\mathcal{D}^\circ) = \frac{\omega_n^2}{\operatorname{vol}_n(\mathcal{D})} < \frac{\omega_n^2}{\operatorname{vol}_n(\mathcal{E}_L^\circ)} = \operatorname{vol}_n(\mathcal{E}_L),\] 与 \(\mathcal{E}_L\) 的极小体积性矛盾。因此 \(\mathcal{E}_L^\circ\) 必为 \(K^\circ\) 的John椭球。

由John定理（定理195.2.2），对原点对称的 \(K^\circ\) 有 \(\mathcal{E}_L^\circ \subseteq K^\circ \subseteq \sqrt{n} \mathcal{E}_L^\circ\)，取极体得 \(\frac{1}{\sqrt{n}} \mathcal{E}_L \subseteq K \subseteq \mathcal{E}_L\)。对一般（非对称）情形，\(K \subseteq \mathcal{E}_L \subseteq n \cdot (K - c)\) 对适当的中心 \(c\) 成立。\(\blacksquare\)

\subsection{195.3 等周型问题与Blaschke-Santalo不等式}\label{ux7b49ux5468ux578bux95eeux9898ux4e0eblaschke-santaloux4e0dux7b49ux5f0f}

\textbf{定义 195.3.1（仿射等周不等式）} 凸体 \(K \in \mathcal{K}^n\) 的仿射表面积 \(\Omega(K)\) 定义为满足某种广义曲率条件的表面积泛函的仿射不变版本。仿射等周不等式陈述为：

\[\frac{\Omega(K)^{n+1}}{\operatorname{vol}_n(K)^{n-1}} \leq \frac{\Omega(B^n)^{n+1}}{\operatorname{vol}_n(B^n)^{n-1}} = n^{n+1} \omega_n^{n-1},\]

等号成立当且仅当 \(K\) 为椭球。

\textbf{定理 195.3.2（Blaschke-Santalo不等式的函数形式）} 对任意偶对数凹函数 \(f: \mathbb{R}^n \to [0, \infty)\)（即 \(\log f\) 为凹函数），

\[\int_{\mathbb{R}^n} f(x) dx \cdot \int_{\mathbb{R}^n} f^*(y) dy \leq (2\pi)^n,\]

其中 \(f^*(y) = \inf_{x \in \mathbb{R}^n} e^{-\langle x, y \rangle} / f(x)\) 为 \(f\) 的极函数（Legendre变换的指数版本）。等号成立当且仅当 \(f\) 是Gauss型函数。取 \(f = \mathbf{1}_K\) 即得原始Blaschke-Santalo不等式。

\textbf{证明}：设 \(f: \mathbb{R}^n \to [0, \infty)\) 为偶对数凹函数，即 \(f(x) = f(-x)\) 且 \(\log f\) 为凹函数。定义 \(f^*(y) = \inf_{x \in \mathbb{R}^n} e^{-\langle x, y \rangle} / f(x)\)。

第一步：由对数凹性，\(f\) 可表示为 \(f(x) = e^{-\varphi(x)}\)，其中 \(\varphi: \mathbb{R}^n \to \mathbb{R} \cup \{+\infty\}\) 为凸偶函数。此时 \(f^*(y) = e^{-\varphi^*(y)}\)，其中 \(\varphi^*\) 为 \(\varphi\) 的Legendre-Fenchel共轭： \[\varphi^*(y) = \sup_{x \in \mathbb{R}^n} (\langle x, y \rangle - \varphi(x)).\]

第二步：不等式 \(\int f \cdot \int f^* \leq (2\pi)^n\) 可改写为 \[\int_{\mathbb{R}^n} e^{-\varphi(x)} dx \cdot \int_{\mathbb{R}^n} e^{-\varphi^*(y)} dy \leq (2\pi)^n.\] 取 \(\varphi(x) = \frac{1}{2}\|x\|^2\)，则 \(\varphi^*(y) = \frac{1}{2}\|y\|^2\)，且 \(\int e^{-\frac{1}{2}\|x\|^2} dx = (2\pi)^{n/2}\)，乘积精确等于 \((2\pi)^n\)。

第三步：对一般的偶凸函数 \(\varphi\)，由Prékopa-Leindler不等式（或等价的Brunn-Minkowski泛函形式），考虑函数 \[u(x, t) = (2\pi t)^{-n/2} e^{-\frac{\|x\|^2}{2t}} * e^{-\varphi(x)},\] 其中 \(*\) 表示卷积。由对数凹函数的卷积仍为对数凹函数（Prékopa-Leindler的直接推论），\(u(x, t)\) 在 \((x, t)\) 上是对数凹的。进一步，\(u(0, 0) = e^{-\varphi(0)}\) 且 \(u(0, 1) = (2\pi)^{-n/2} \int e^{-\varphi(x)} dx \cdot e^{-\|x\|^2/2}\)。由对数凹性和变量替换可得所需不等式。详细的Blaschke-Santaló函数形式证明涉及测度的重排和球面平均的单调性，由Ball和Artstein-Avidan等给出。\(\blacksquare\)

\textbf{定理 195.3.3（逆向Blaschke-Santalo不等式）} 存在仅依赖于维数 \(n\) 的正常数 \(c_n\)，使得对任意原点对称凸体 \(K \in \mathcal{K}_s^n\)，有

\[\mathcal{M}(K) \geq c_n \cdot \omega_n^2.\]

Milman证明 \(c_n\) 可取为 \(e^{-c n}\) 的形式，其中 \(c > 0\) 是绝对常数。进一步，Klartag证明在无条件凸体（关于坐标轴对称）的类中，Mahler猜想的最优下界 \(\frac{4^n}{n!}\) 是正确的。

\textbf{证明}：Milman的逆向Blaschke-Santalo不等式的证明依赖于逆向Brunn-Minkowski不等式。

第一步：回顾Milman的逆向Brunn-Minkowski不等式：存在绝对常数 \(C > 0\) 使得对任意原点对称凸体 \(K, L \in \mathcal{K}_s^n\)， \[\operatorname{vol}_n(K + L)^{1/n} \leq C (\operatorname{vol}_n(K)^{1/n} + \operatorname{vol}_n(L)^{1/n}).\]

第二步：取 \(L = K^\circ\)。由极体的基本性质，\(K + K^\circ \supseteq K\) 且 \(K + K^\circ \supseteq K^\circ\)。利用逆向Brunn-Minkowski不等式， \[\operatorname{vol}_n(K + K^\circ)^{1/n} \leq C (\operatorname{vol}_n(K)^{1/n} + \operatorname{vol}_n(K^\circ)^{1/n}).\]

第三步：另一方面，存在仅依赖于维数 \(n\) 的常数 \(d_n\)（由John定理导出）使得 \[\operatorname{vol}_n(K + K^\circ)^{1/n} \geq d_n \cdot \max(\operatorname{vol}_n(K)^{1/n}, \operatorname{vol}_n(K^\circ)^{1/n}).\] 由 \(K\) 和 \(K^\circ\) 的对称性结合上述两式，可得 \[\max(\operatorname{vol}_n(K), \operatorname{vol}_n(K^\circ))^{1/n} \leq C (\operatorname{vol}_n(K)^{1/n} + \operatorname{vol}_n(K^\circ)^{1/n}).\]

第四步：设 \(a = \operatorname{vol}_n(K)^{1/n}\)，\(b = \operatorname{vol}_n(K^\circ)^{1/n}\)。由上式和Blaschke-Santalo不等式的上界 \(ab \leq c \omega_n^{2/n}\)（对某常数 \(c\)），可导出 \(ab \geq c' \omega_n^{2/n}\)。这意味着 \[\mathcal{M}(K) = a^n b^n \geq (c')^n \omega_n^2 = e^{-c'' n} \omega_n^2,\] 即 \(c_n = e^{-c'' n}\) 满足要求。这里 \(c''\) 为绝对常数。\(\blacksquare\)

\textbf{定理 195.3.4（Ball的逆向等周不等式）} 对任意原点对称凸体 \(K \in \mathcal{K}_s^n\) 和 \(1 \leq k \leq n - 1\)，

\[\left(\frac{\operatorname{vol}_n(K)}{\operatorname{vol}_n(K \cap E)}\right)^{1/(n-k)} \leq C_{n,k},\]

其中 \(E\) 是任意 \(k\) 维子空间，\(C_{n,k}\) 是仅依赖于 \(n, k\) 的常数，最优常数为 \(\sqrt{n}/\sqrt{k}\) 数量级。

\textbf{证明}：固定 \(k\) 维子空间 \(E \subseteq \mathbb{R}^n\)，记 \(K_E = K \cap E\) 为 \(K\) 在 \(E\) 中的截面。由Brunn-Minkowski不等式的投影形式和极对偶性，截面体积与正交补上的投影体积通过Fourier变换和等周不等式相关联。

第一步：利用球的极值性质。对单位球 \(B^n\)，截面 \(B^n \cap E\) 是 \(k\) 维单位球 \(B^k\)，故 \(\operatorname{vol}_k(B^n \cap E) = \omega_k\)，且 \(\operatorname{vol}_n(B^n) = \omega_n\)，其比值为 \(\omega_k / \omega_n\)。

对一般对称凸体 \(K\)，设 \(\mathcal{E}\) 为 \(K\) 的John椭球。由John定理，\(\mathcal{E} \subseteq K \subseteq \sqrt{n} \mathcal{E}\)。因此 \[\frac{\operatorname{vol}_n(K)}{\operatorname{vol}_k(K \cap E)} \geq \frac{\operatorname{vol}_n(\mathcal{E})}{\operatorname{vol}_k((\sqrt{n}\mathcal{E}) \cap E)} = \frac{\operatorname{vol}_n(\mathcal{E})}{n^{k/2} \operatorname{vol}_k(\mathcal{E} \cap E)}.\]

第二步：对椭球 \(\mathcal{E}\)，截面体积有精确计算公式。通过适当的线性变换，可设 \(\mathcal{E} = B^n\)。则 \(\mathcal{E} \cap E = B^k\)。故 \[\frac{\operatorname{vol}_n(\mathcal{E})}{\operatorname{vol}_k(\mathcal{E} \cap E)} = \frac{\omega_n}{\omega_k} \sim \frac{1}{\sqrt{2\pi}} \left(\frac{n}{k}\right)^{k/2} \left(\frac{2\pi e}{n}\right)^{n/2} \cdot \left(\frac{k}{2\pi e}\right)^{k/2}.\]

第三步：综合得 \[\left(\frac{\operatorname{vol}_n(K)}{\operatorname{vol}_k(K \cap E)}\right)^{1/(n-k)} \leq C \cdot \left(\frac{\sqrt{n}}{\sqrt{k}}\right)^{n/(n-k)} \cdot \left(\frac{\omega_n}{\omega_k}\right)^{1/(n-k)} \leq C_{n,k}.\] 其中 \(C_{n,k} \approx \sqrt{n/k}\)。等号在 \(K\) 为球面截面时达到最优。\(\blacksquare\)

\subsection{195.4 Dvoretzky定理与Milman浓度现象}\label{dvoretzkyux5b9aux7406ux4e0emilmanux6d53ux5ea6ux73b0ux8c61}

\textbf{定理 195.4.1（Dvoretzky定理）} 对任意 \(\varepsilon > 0\)，存在 \(N = N(\varepsilon, n)\) 使得任意 \(N\) 维赋范空间 \(X\) 包含一个 \(n\) 维子空间 \(Y\) 满足：

\[d(Y, \ell_2^n) \leq 1 + \varepsilon,\]

其中 \(d(Y, \ell_2^n)\) 表示 \(Y\) 到 \(\ell_2^n\) 的Banach-Mazur距离。等价地，任意高维对称凸体在大多数低维截面上都近似于椭球。

\textbf{证明概要}：Dvoretzky定理的证明是20世纪泛函分析和凸几何最伟大的成就之一。关键步骤包括：（1）利用Lévy浓度现象，在Grassmann流形 \(G_{N, n}\) 上定义随机子空间；（2）在 \(N\) 维Euclidean球面 \(\mathbb{S}^{N-1}\) 上，范数的限制几乎处处接近其平均值；（3）利用Dvoretzky-Rogers引理将连续化归为离散格点问题。Milman给出了该定理的极大简化证明，将定理归结为球面上Lipschitz函数的浓度性质。\(\blacksquare\)

\textbf{定理 195.4.2（Milman浓度现象）} 设 \(\mu\) 为 \(\mathbb{S}^{n-1}\) 上的标准旋转不变概率测度。若 \(f: \mathbb{S}^{n-1} \to \mathbb{R}\) 是 \(1\)-Lipschitz函数，\(M\) 为 \(f\) 的中位数，则

\[\mu\{x \in \mathbb{S}^{n-1} : |f(x) - M| > t\} \leq 2e^{-c n t^2},\]

其中 \(c > 0\) 是绝对常数。此不等式表明高维球面上的Lipschitz函数高度集中于其中位数附近------这是高维几何中最引人注目的现象之一。

\textbf{证明}：设 \(\mu\) 为 \(\mathbb{S}^{n-1}\) 上归一化的旋转不变测度（即 \((n-1)\) 维Hausdorff测度除以 \(|\mathbb{S}^{n-1}|\)）。

第一步：利用Lévy等周不等式。对任意Borel集 \(A \subseteq \mathbb{S}^{n-1}\)，其 \(t\)-管状邻域 \(A_t = \{x \in \mathbb{S}^{n-1} : d(x, A) < t\}\)（\(d\) 为球面上的测地距离）满足 \[\mu(A_t) \geq 1 - \frac{\sqrt{\pi/8}}{\mu(A)} e^{-(n-1)t^2/2}.\] 等价地，若 \(\mu(A) \geq 1/2\)，则 \(\mu(A_t) \geq 1 - \sqrt{\pi/2} e^{-(n-1)t^2/2}\)。

第二步：设 \(f: \mathbb{S}^{n-1} \to \mathbb{R}\) 为 \(1\)-Lipschitz函数，\(M\) 为其Lévy中位数（或任意满足 \(\mu\{f \leq M\} \geq 1/2\) 且 \(\mu\{f \geq M\} \geq 1/2\) 的数）。对 \(t > 0\)，令 \(A = \{x : f(x) \leq M\}\)，则 \(\mu(A) \geq 1/2\)。

若 \(x \in A_t\)，则存在 \(y \in A\) 满足 \(d(x, y) < t\)。由Lipschitz性，\(f(x) \leq f(y) + d(x, y) < M + t\)。因此 \(\{f > M + t\} \subseteq \mathbb{S}^{n-1} \setminus A_t\)。由Lévy等周不等式， \[\mu\{f > M + t\} \leq 1 - \mu(A_t) \leq \sqrt{\pi/2} e^{-(n-1)t^2/2} \leq e^{-c n t^2},\] 其中 \(c > 0\) 可取为某绝对常数（例如 \(c \approx 1/3\) 当 \(n \geq 2\)）。

第三步：对称地考虑 \(\{f \geq M\}\) 得 \(\mu\{f < M - t\} \leq e^{-c n t^2}\)。综合得 \[\mu\{|f - M| > t\} \leq 2e^{-c n t^2}.\] 若以平均值 \(\bar{f} = \int f d\mu\) 替代中位数，由上述浓度可推出 \(|M - \bar{f}| \leq C/\sqrt{n}\)，故均值附近也有类似的浓度估计。\(\blacksquare\)

\textbf{推论 195.4.3（Dvoretzky定理的定量版本）} 令 \(k(\varepsilon)\) 为满足如下条件的最大维数：任意 \(n\) 维赋范空间包含 \(k(\varepsilon)\) 维子空间，其到 \(\ell_2^{k(\varepsilon)}\) 的Banach-Mazur距离不超过 \(1 + \varepsilon\)。Milman证明 \(k(\varepsilon) \geq c \varepsilon^2 n / |\log \varepsilon|\)（对称凸体情形）。

\textbf{定理 195.4.4（Dvoretzky-Milman定理的推广------随机截面）} 在适当的概率模型下，随机低维截面（按Haar测度选取）以高概率接近Euclidean球。更精确地，存在绝对常数 \(c > 0\) 使得对任意 \(n\) 维对称凸体 \(K\) 和 \(k \leq c \log n\)，\(K\) 的随机 \(k\) 维中心截面以高概率满足 \(d(K \cap E, B_E^k) \leq C\)（\(C\) 为绝对常数）。

\textbf{证明}：设 \(K \in \mathcal{K}_s^n\) 为原点对称凸体。选取 \(k\) 维子空间 \(E\)（按Grassmann流形 \(G_{n,k}\) 上的Haar测度随机选取）。定义 \(K \cap E\) 上的范数 \(\|\cdot\|_K\)，其单位球即为 \(K \cap E\)。

第一步：定义函数 \(f_K(u) = \|u\|_K\)（\(u \in \mathbb{S}^{n-1}\)）。由 \(K\) 的对称性和凸性，\(f_K\) 在 \(\mathbb{S}^{n-1}\) 上为 \(1\)-Lipschitz函数（关于欧几里得度量与对偶范数之间的关系）。

第二步：设 \(M\) 为 \(f_K\) 在 \(\mathbb{S}^{n-1}\) 上的中位数。由Milman浓度现象（定理195.4.2），对 \(t > 0\)， \[\mu\{|f_K(u) - M| > t\} \leq 2e^{-c n t^2}.\]

第三步：对随机选取的 \(k\) 维子空间 \(E\)，\(\mathbb{S}^{n-1} \cap E\) 的 \(k-1\) 维球面是 \(\mathbb{S}^{n-1}\) 的随机截面。由测度的积分几何公式，在 \(\mathbb{S}^{n-1} \cap E\) 上，\(f_K\) 的分布与 \(f_K\) 在 \(\mathbb{S}^{n-1}\) 上的限制相关。当 \(k \leq c \log n\) 时，\(f_K\) 在 \(\mathbb{S}^{n-1} \cap E\) 上的震荡以高概率小于 \(\varepsilon\)。

第四步：由Milman的论证，通过精细的 \(\varepsilon\)-网论证（在 \(\mathbb{S}^{n-1} \cap E\) 上构造大小为 \((1 + 2/\varepsilon)^k\) 的 \(\varepsilon\)-网），可得对随机 \(E\)，以概率至少 \(1 - e^{-c n}\) 有 \[(1 - \varepsilon) M \|x\| \leq \|x\|_K \leq (1 + \varepsilon) M \|x\|, \quad \forall x \in E,\] 即 \(d(K \cap E, B_E^k) \leq \frac{1 + \varepsilon}{1 - \varepsilon} \approx 1 + 2\varepsilon\)。\(\blacksquare\)

\subsection{195.5 凸体的Banach-Mazur距离}\label{ux51f8ux4f53ux7684banach-mazurux8dddux79bb}

\textbf{定义 195.5.1（Banach-Mazur距离）} 两个 \(n\) 维赋范空间 \(X, Y\)（或等价地，两个关于原点对称的 \(n\) 维凸体 \(K_X, K_Y\)）的Banach-Mazur距离定义为：

\[d(X, Y) = \inf \{\|T\| \cdot \|T^{-1}\| : T: X \to Y \text{为线性同构}\}.\]

在凸体语言中，\(d(K, L) = \inf \{r > 0 : \exists T \in GL(n), K \subseteq T L \subseteq r K\}\)。

\textbf{命题 195.5.2（Banach-Mazur距离的基本性质）}

\begin{enumerate}
\def\labelenumi{\arabic{enumi}.}
\tightlist
\item
  \(d(X, Y) = d(Y, X)\)；
\item
  \(d(X, Z) \leq d(X, Y) \cdot d(Y, Z)\) （对数三角不等式）；
\item
  \(d(X, Y) \geq 1\)，且 \(d(X, Y) = 1\) 当且仅当 \(X \cong Y\)（等距同构）；
\item
  由John定理，\(d(X, \ell_2^n) \leq \sqrt{n}\) 对任意 \(n\) 维对称凸体对应的空间 \(X\) 成立。
\end{enumerate}

\textbf{定理 195.5.3（Gluskin定理）} \(n\) 维Banach-Mazur紧致空间的直径（即 \(\sup_{X, Y} d(X, Y)\)）的数量级为 \(n\)。更精确地，存在绝对常数 \(c, C > 0\) 使得

\[c n \leq \sup_{X, Y \in \mathcal{B}_n} d(X, Y) \leq C n,\]

其中 \(\mathcal{B}_n\) 是所有 \(n\) 维赋范空间（等距类）构成的空间。

\textbf{证明概要}：上界 \(n\) 由John定理的直接推论得出。下界 \(c n\) 由Gluskin的随机多胞体构造给出：选取 \(N = e^{c n}\) 个随机独立点 \(x_1, \ldots, x_N \in \mathbb{S}^{n-1}\)，定义 \(K = \operatorname{conv}\{\pm x_1, \ldots, \pm x_N\}\)。由随机矩阵的极值性质，\(d(K, K^\circ) \geq c n\) 以高概率成立。\(\blacksquare\)

\textbf{定理 195.5.4（Banach-Mazur距离的体积比方法）} Szarek和Talagrand发展的体积比方法（volume ratio method）给出了Banach-Mazur距离的强有力估计。对对称凸体 \(K\)，定义其体积比为

\[\operatorname{vr}(K) = \left(\frac{\operatorname{vol}_n(K)}{\operatorname{vol}_n(\mathcal{E})}\right)^{1/n},\]

其中 \(\mathcal{E}\) 是 \(K\) 的极大体积椭球。则 \(d(K, \ell_2^n) \leq c \cdot \operatorname{vr}(K) \cdot \operatorname{vr}(K^\circ)\)。此不等式结合随机超平面截面的浓度论证，可给出Dvoretzky定理的定量版本的另一种证明。

\textbf{证明}：设 \(K \in \mathcal{K}_s^n\) 为原点对称凸体。令 \(\mathcal{E}\) 为 \(K\) 的John椭球（极大体积内接椭球），\(\mathcal{E}_L\) 为 \(K\) 的Lowner椭球（极小体积外接椭球）。

第一步：体积比的定义和John定理的直接推论。\(\operatorname{vr}(K) = (\operatorname{vol}_n(K) / \operatorname{vol}_n(\mathcal{E}))^{1/n}\)。由John定理，\(\mathcal{E} \subseteq K \subseteq \sqrt{n} \mathcal{E}\)，故 \[1 \leq \operatorname{vr}(K) \leq \sqrt{n}.\] 类似地，\(\operatorname{vr}(K^\circ) \leq \sqrt{n}\)。上界由Blaschke-Santaló不等式进一步改进。

第二步：Banach-Mazur距离到 \(\ell_2^n\) 的上界。设 \(T \in GL(n)\) 为将 \(\mathcal{E}\) 映为单位球 \(B^n\) 的线性变换。则 \(T K\) 满足 \(B^n \subseteq T K \subseteq R B^n\)，其中 \(R = \inf\{r > 0 : T K \subseteq r B^n\}\)。我们有 \[d(K, \ell_2^n) \leq d(T K, B^n) \leq R.\] 由包含关系和体积的单调性， \[\operatorname{vol}_n(T K) \leq R^n \operatorname{vol}_n(B^n) = R^n \omega_n,\] 故 \(R \geq (\operatorname{vol}_n(T K) / \omega_n)^{1/n} = (\det T \cdot \operatorname{vol}_n(K) / \omega_n)^{1/n}\)。

第三步：同理，由 \(B^n \subseteq T K\) 有 \(\omega_n \leq \det T \cdot \operatorname{vol}_n(K)\)。结合对 \(K^\circ\) 的类似分析，得 \[d(K, \ell_2^n) \leq \left(\frac{\operatorname{vol}_n(K)}{\operatorname{vol}_n(\mathcal{E})}\right)^{1/n} \cdot \left(\frac{\operatorname{vol}_n(K^\circ)}{\operatorname{vol}_n(\mathcal{E}^\circ)}\right)^{1/n} = \operatorname{vr}(K) \cdot \operatorname{vr}(K^\circ).\]

第四步：由Szarek-Talagrand的精细估计，常数 \(c\) 实际上可取为某绝对常数（与 \(n\) 无关）。核心是通过随机化论证：随机选取 \(N = e^{c n}\) 个 \(\mathbb{S}^{n-1}\) 上的点，它们的凸包与John椭球的体积比以高概率受到控制。\(\blacksquare\)

\textbf{推论 195.5.5（Banach-Mazur距离与逼近理论）} \(n\) 维空间类在Banach-Mazur距离下的渐近几何学（asymptotic geometric analysis）揭示了许多反直觉的现象：例如，对大多数 \(n\) 维空间，其到 \(\ell_2^n\) 的距离实际上接近 \(\sqrt{n}\) 而不是更小的数值。这表明高维几何中极端行为是典型的而非常见结构是特殊的。

凸几何的这些深刻结果------从Blaschke选择定理的紧致性到Aleksandrov-Fenchel不等式的完全非线性结构，从Dvoretzky定理的浓度现象到Mahler猜想的极值对称性------构成了现代几何分析的核心支柱，并持续推动着泛函分析、几何测度论和随机矩阵理论的发展。

\newpage

\chapter{08-代数几何与数论/01-代数几何/01-仿射与射影簇.md}\label{ux4ee3ux6570ux51e0ux4f55ux4e0eux6570ux8bba01-ux4ee3ux6570ux51e0ux4f5501-ux4effux5c04ux4e0eux5c04ux5f71ux7c07.md}

\section{第280章：仿射代数簇}\label{ux7b2c280ux7ae0ux4effux5c04ux4ee3ux6570ux7c07}

仿射代数簇是代数几何最基础的研究对象，它是多项式方程组的零点集，带有 Zariski 拓扑和坐标环结构。本章的核心是 Hilbert 零点定理，它建立了代数集与理想之间的深刻对应。

\subsection{81.1 仿射代数集}\label{ux4effux5c04ux4ee3ux6570ux96c6}

\textbf{定义 81.1}（仿射空间）：域 \(k\) 上的 \(n\) 维\textbf{仿射空间} \(\mathbb{A}_k^n\)（或 \(\mathbb{A}^n\)）是 \(k^n\) 作为集合，但赋予来自多项式环 \(k[x_1, \ldots, x_n]\) 的几何结构。

\textbf{定义 81.2}（代数集）：设 \(S \subseteq k[x_1, \ldots, x_n]\) 是多项式子集。\(S\) 确定的\textbf{仿射代数集}（或零点集）为

\[V(S) = \{P \in \mathbb{A}^n : f(P) = 0 \text{ 对所有 } f \in S\}\]

若 \(S = \{f_1, \ldots, f_m\}\) 有限，则 \(V(S) = V(f_1, \ldots, f_m)\) 是 \(m\) 个多项式方程的解集。

\textbf{定义 81.3}（理想）：若 \(I\) 是 \(k[x_1, \ldots, x_n]\) 的理想，记 \(V(I) = V(S)\)，其中 \(S\) 是 \(I\) 的任意生成元集。由 Hilbert 基定理（定理 68.2），\(k[x_1, \ldots, x_n]\) 是 Noether 环，故每个理想 \(I\) 由有限个多项式生成，因此 \(V(I)\) 始终是有限个方程的解集。

\textbf{命题 81.1}（代数集的基本性质）：

\begin{enumerate}
\def\labelenumi{\arabic{enumi}.}
\tightlist
\item
  \(V(0) = \mathbb{A}^n\)，\(V(1) = \varnothing\)
\item
  \(\bigcap_\alpha V(I_\alpha) = V\left(\sum_\alpha I_\alpha\right)\)（任意交）
\item
  \(V(I) \cup V(J) = V(IJ) = V(I \cap J)\)（有限并）
\end{enumerate}

因此，代数集的补集构成了 \(\mathbb{A}^n\) 上某个拓扑的闭集，称为 \textbf{Zariski 拓扑}。

\emph{证明}：1 显然。2：\(P \in \bigcap_\alpha V(I_\alpha)\) 当且仅当对所有 \(\alpha\)，\(f(P) = 0\)（\(\forall f \in I_\alpha\)），当且仅当 \(P \in V(\sum_\alpha I_\alpha)\)。3：\(P \in V(I) \cup V(J)\) 当且仅当 \(P \in V(I)\) 或 \(P \in V(J)\)。若 \(P \in V(I)\)，则对任意 \(fg \in IJ\)，\((fg)(P) = f(P)g(P) = 0\)，故 \(P \in V(IJ)\)。反之，若 \(P \notin V(I) \cup V(J)\)，则存在 \(f \in I\)，\(g \in J\) 使 \(f(P) \neq 0\)，\(g(P) \neq 0\)。则 \(fg \in IJ\) 且 \((fg)(P) \neq 0\)，故 \(P \notin V(IJ)\)。∎

\textbf{定义 81.4}（Zariski 拓扑）：\(\mathbb{A}^n\) 上的 \textbf{Zariski 拓扑} 以代数集为闭集。开集是代数集的补集。

\textbf{注释 81.1}：Zariski 拓扑不是 Hausdorff 的（除非 \(n = 0\)）。任意两个非空开集都有非空交（不可约性）。这与分析中的拓扑形成鲜明对比，体现了代数几何的独特性质。

\subsection{81.2 Hilbert 零点定理}\label{hilbert-ux96f6ux70b9ux5b9aux7406}

\textbf{定义 81.5}（理想的方向）：对子集 \(X \subseteq \mathbb{A}^n\)，定义

\[I(X) = \{f \in k[x_1, \ldots, x_n] : f(P) = 0 \text{ 对所有 } P \in X\}\]

\(I(X)\) 是 \(k[x_1, \ldots, x_n]\) 的理想（称为 \(X\) 的\textbf{消去理想}）。

\textbf{定义 81.6}（根理想）：理想 \(I\) 的\textbf{根}为

\[\sqrt{I} = \{f \in k[x_1, \ldots, x_n] : f^m \in I \text{ 对某个 } m \geq 1\}\]

若 \(I = \sqrt{I}\)，则称 \(I\) 为\textbf{根理想}。

\textbf{定理 81.2}（Hilbert 零点定理，Nullstellensatz）：设 \(k\) 是代数闭域。则对任意理想 \(I \subseteq k[x_1, \ldots, x_n]\)，

\[I(V(I)) = \sqrt{I}\]

此外，若 \(V(I) = \varnothing\)，则 \(I = (1)\)（弱零点定理）。

\textbf{证明}：（1）\(X\)不可约等价于\(\mathcal{I}(X)\)为素理想。若\(X=X_1\cup X_2\)（\(X_i\)为真闭子簇），则\(\mathcal{I}(X)=\mathcal{I}(X_1)\cap\mathcal{I}(X_2)\)且\(\mathcal{I}(X_i)\supsetneq\mathcal{I}(X)\)，故\(\mathcal{I}(X)\)非素。反之，若\(\mathcal{I}(X)\)非素，存在\(fg\in\mathcal{I}(X)\)但\(f,g\notin\mathcal{I}(X)\)，则\(X=(X\cap V(f))\cup(X\cap V(g))\)。（2）坐标环\(A(X)=k[x_1,\ldots,x_n]/\mathcal{I}(X)\)为整环当且仅当\(\mathcal{I}(X)\)为素理想，后者等价于\(X\)不可约。（3）\(\dim X=\operatorname{trdeg}_k k(X)\)（\(k(X)\)为\(A(X)\)的分式域），由Noether正规化引理，\(A(X)\)为有限生成\(k[y_1,\ldots,y_d]\)-模（\(d=\dim X\)）。\(\blacksquare\)

\textbf{推论 81.3}（代数集与根理想的一一对应）：在代数闭域 \(k\) 上，映射 \(I \mapsto V(I)\) 和 \(X \mapsto I(X)\) 给出了

\[\{\text{根理想 } I \subseteq k[x_1, \ldots, x_n]\} \longleftrightarrow \{\text{代数集 } X \subseteq \mathbb{A}^n\}\]

之间的一一对应（反向包含序）。极大理想对应于单点集。

\subsection{81.3 不可约性与代数簇}\label{ux4e0dux53efux7ea6ux6027ux4e0eux4ee3ux6570ux7c07}

\textbf{定义 81.7}（不可约代数集）：代数集 \(X\) 称为\textbf{不可约的}，如果 \(X\) 不能写成两个真闭子集的并：\(X = X_1 \cup X_2\) 推出 \(X = X_1\) 或 \(X = X_2\)。不可约的仿射代数集称为\textbf{仿射代数簇}（affine variety）。

\textbf{命题 81.4}：代数集 \(X\) 不可约当且仅当 \(I(X)\) 是素理想。

\emph{证明}：\(X\) 可约 \(\iff\) \(X = V(I) \cup V(J) = V(IJ)\)（\(X \neq V(I), V(J)\)）\(\iff\) \(IJ \subseteq I(X)\) 但 \(I \nsubseteq I(X)\)，\(J \nsubseteq I(X)\) \(\iff\) \(I(X)\) 不是素理想。∎

\textbf{推论 81.5}（不可约分解）：每个代数集 \(X\) 可唯一地写成有限个不可约代数簇的并 \(X = X_1 \cup \cdots \cup X_r\)，其中 \(X_i \nsubseteq X_j\)（\(i \neq j\)）。这对应于 \(I(X)\) 的极小素理想分解。

\emph{证明}：由 \(k[x_1, \ldots, x_n]\) 是 Noether 环，每个根理想是有限个极小素理想的交。唯一性由素理想的极小性保证。∎

\subsection{81.4 坐标环与态射}\label{ux5750ux6807ux73afux4e0eux6001ux5c04}

\textbf{定义 81.8}（坐标环）：仿射代数簇 \(X \subseteq \mathbb{A}^n\) 的\textbf{坐标环}（或仿射坐标环）定义为

\[A(X) = k[x_1, \ldots, x_n] / I(X)\]

\(A(X)\) 是有限生成 \(k\)-代数，且是整环（因为 \(I(X)\) 是素理想）。\(A(X)\) 的元素是 \(X\) 上的正则函数。

\textbf{定义 81.9}（正则函数）：\(X\) 上的\textbf{正则函数}是多项式函数 \(f: X \to k\)，即 \(f \in A(X)\)。\(X\) 上的正则函数全体构成环 \(A(X)\)。

\textbf{命题 81.6}（坐标环的函子性）：仿射代数簇与有限生成整 \(k\)-代数之间的范畴等价：

\[\{\text{仿射代数簇 } / k\} \longleftrightarrow \{\text{有限生成整 } k\text{-代数}\}^{\text{op}}\]

等价由 \(X \mapsto A(X)\) 和 \(A \mapsto \operatorname{Spec}_m(A)\)（\(A\) 的极大谱）给出。

\textbf{定义 81.10}（态射）：仿射代数簇 \(X \subseteq \mathbb{A}^n\) 到 \(Y \subseteq \mathbb{A}^m\) 的\textbf{态射}（正则映射）\(\varphi: X \to Y\) 是每个分量由正则函数给出的映射：\(\varphi = (f_1, \ldots, f_m)\)，\(f_i \in A(X)\)，且 \(\varphi(X) \subseteq Y\)。

态射诱导坐标环的反向同态 \(\varphi^*: A(Y) \to A(X)\)，\(\varphi^*(g) = g \circ \varphi\)。

\textbf{定理 81.7}（仿射簇范畴等价）：映射 \(X \mapsto A(X)\) 是仿射代数簇范畴与有限生成整 \(k\)-代数范畴的反向范畴等价。

\textbf{证明}：由 Hilbert 零点定理（定理 81.2），根理想与代数集一一对应，素理想对应于不可约代数集即仿射簇。因此每个仿射簇 \(X \subseteq \mathbb{A}^n\) 对应坐标环 \(A(X) = k[x_1, \ldots, x_n]/I(X)\)，它是有限生成整 \(k\)-代数。反之，给定有限生成整 \(k\)-代数 \(A \cong k[x_1, \ldots, x_n]/I\)（\(I\) 是素理想），取 \(X = V(I) \subseteq \mathbb{A}^n\)，则 \(A(X) \cong A\)。态射 \(\varphi: X \to Y\) 的每个分量是 \(X\) 上的正则函数，通过拉回诱导 \(k\)-代数同态 \(\varphi^*: A(Y) \to A(X)\)，且该对应是函子性的。由此得到范畴等价。\(\blacksquare\)

\subsection{81.5 维数理论}\label{ux7ef4ux6570ux7406ux8bba}

\textbf{定义 81.11}（Krull 维数）：仿射代数簇 \(X\) 的\textbf{维数} \(\dim X\) 定义为 \(A(X)\) 的 Krull 维数------即 \(A(X)\) 中素理想链的最大长度：

\[\dim X = \max\{r : \mathfrak{p}_0 \subsetneq \mathfrak{p}_1 \subsetneq \cdots \subsetneq \mathfrak{p}_r \subseteq A(X), \mathfrak{p}_i \text{ 素理想}\}\]

\textbf{定义 81.12}（超越次数）：设 \(K/k\) 是有限生成域扩张。\(K\) 在 \(k\) 上的\textbf{超越次数} \(\operatorname{tr.deg}_k K\) 是 \(K\) 中代数无关的极大子集的大小。

\textbf{定理 81.8}（维数与超越次数）：对仿射代数簇 \(X\)，

\[\dim X = \operatorname{tr.deg}_k k(X)\]

即 \(X\) 的维数等于其有理函数域在 \(k\) 上的超越次数。

\textbf{证明}：设 \(\dim X = d\)。由 Noether 正规化引理（定理 68.7），存在代数无关的 \(y_1, \ldots, y_d \in A(X)\) 使 \(A(X)\) 是 \(k[y_1, \ldots, y_d]\) 上的有限生成模（整扩张）。取分式域得 \(k(X)\) 是 \(k(y_1, \ldots, y_d)\) 的有限代数扩张，故 \(\operatorname{tr.deg}_k k(X) = d\)。反向：若 \(\operatorname{tr.deg}_k k(X) = d\)，取超越基后 \(k(X)/k(t_1, \ldots, t_d)\) 是有限代数扩张，则 \(A(X)\) 在 \(k[t_1, \ldots, t_d]\) 上是整的，故 \(\dim A(X) = \dim k[t_1, \ldots, t_d] = d\)。\(\blacksquare\)

\textbf{命题 81.9}（维数的基本性质）： 1. \(\dim \mathbb{A}^n = n\) 2. 若 \(Y \subsetneq X\) 是不可约闭子簇，则 \(\dim Y < \dim X\) 3. 超曲面 \(V(f) \subseteq \mathbb{A}^n\)（\(f\) 非常数）的维数为 \(n-1\)

\textbf{定理 81.10}（Krull 高度定理）：设 \(A\) 是 Noether 环，\(f \in A\) 非零因子。则 \(f\) 生成的理想的每个极小素理想的高度为 \(1\)。特别地，每个真包含 \(\{0\}\) 的极小素理想的高度为 \(1\)。

\textbf{证明}：设 \(\mathfrak{p}\) 是 \((f)\) 的极小素理想。局部化到 \(A_{\mathfrak{p}}\)，可设 \((A, \mathfrak{m})\) 是 Noether 局部环，\(\mathfrak{m}\) 是 \(\sqrt{(f)}\) 的极小素理想。由 Krull 交定理，\(\dim A \leq 1\)。又因 \(f\) 非零因子，\(\{0\} \neq (f) \subseteq \mathfrak{m}\)，故 \(\dim A \geq 1\)。因此 \(\dim A_{\mathfrak{p}} = 1\)，即 \(\operatorname{ht}(\mathfrak{p}) = 1\)。逆命题使用 Krull 主理想定理：Noether 环中由 \(r\) 个元素生成的理想的极小素理想的高度不超过 \(r\)；取 \(r=1\) 即得。\(\blacksquare\)

\hrulefill

\hrulefill

\hrulefill

\hrulefill

\hrulefill

\hrulefill

\section{第281章：射影簇与态射}\label{ux7b2c281ux7ae0ux5c04ux5f71ux7c07ux4e0eux6001ux5c04}

射影簇是射影空间中齐次多项式的零点集，其紧致性（在经典拓扑中）使其在几何中具有重要地位。本章引入射影簇、有理映射和双有理等价的概念。

\subsection{82.1 射影空间}\label{ux5c04ux5f71ux7a7aux95f4}

\textbf{定义 82.1}（射影空间）：域 \(k\) 上的 \(n\) 维\textbf{射影空间} \(\mathbb{P}_k^n\)（或 \(\mathbb{P}^n\)）是 \(k^{n+1} \setminus \{0\}\) 在等价关系 \((x_0, \ldots, x_n) \sim (\lambda x_0, \ldots, \lambda x_n)\)（\(\lambda \in k^\times\)）下的商。点 \((x_0 : \cdots : x_n)\) 称为\textbf{齐次坐标}。

\textbf{定义 82.2}（标准仿射开覆盖）：\(\mathbb{P}^n\) 被 \(n+1\) 个仿射开集覆盖：\(U_i = \{(x_0 : \cdots : x_n) : x_i \neq 0\} \cong \mathbb{A}^n\)（\(i = 0, \ldots, n\)）。同构由 \((x_0 : \cdots : x_n) \mapsto (x_0/x_i, \ldots, \widehat{x_i/x_i}, \ldots, x_n/x_i)\) 给出。

\textbf{定义 82.3}（齐次多项式）：多项式 \(F \in k[x_0, \ldots, x_n]\) 称为 \(d\) 次\textbf{齐次}的，如果 \(F(\lambda x_0, \ldots, \lambda x_n) = \lambda^d F(x_0, \ldots, x_n)\)。齐次多项式在射影空间中的零点集有良定义（因为 \(F(\lambda x) = 0 \iff F(x) = 0\)）。

\textbf{定义 82.4}（射影代数集）：设 \(S \subseteq k[x_0, \ldots, x_n]\) 是齐次多项式子集。\(S\) 确定的\textbf{射影代数集}为

\[V_P(S) = \{P \in \mathbb{P}^n : F(P) = 0 \text{ 对所有齐次 } F \in S\}\]

齐次理想 \(I\) 的射影零点集记作 \(V_P(I)\)。

\textbf{定义 82.5}（齐次理想）：\(k[x_0, \ldots, x_n]\) 的理想 \(I\) 称为\textbf{齐次理想}，如果 \(I\) 可由齐次多项式生成。等价地，\(I = \bigoplus_{d \geq 0} (I \cap R_d)\)，其中 \(R_d\) 是 \(d\) 次齐次多项式空间。

\textbf{定理 82.1}（射影零点定理）：设 \(k\) 是代数闭域，\(I \subseteq k[x_0, \ldots, x_n]\) 是齐次理想。若 \(V_P(I) = \varnothing\)，则 \(\sqrt{I} \supseteq (x_0, \ldots, x_n)\)（即 \(I\) 包含所有充分高次齐次多项式）。否则 \(I_P(V_P(I)) = \sqrt{I}\)，其中 \(I_P\) 是射影消去理想。

\textbf{证明}：将射影零点定理归约到仿射情形。考虑仿射锥 \(C(V_P(I)) = V_{\mathbb{A}}(I) \subseteq \mathbb{A}^{n+1}\)。若 \(V_P(I) = \varnothing\)，则 \(V_{\mathbb{A}}(I) \subseteq \{(0, \ldots, 0)\}\)，由仿射零点定理 \(\sqrt{I} \supseteq (x_0, \ldots, x_n)\)。若 \(V_P(I) \neq \varnothing\)，设 \(f\) 零化 \(V_P(I)\)。在仿射锥上 \(f\) 零化 \(C(V_P(I)) \setminus \{0\}\)，由 Zariski 闭包为整个 \(C(V_P(I))\)，故 \(f\) 零化 \(C(V_P(I))\)，由仿射零点定理 \(f \in \sqrt{I}\)。\(\blacksquare\)

\textbf{定义 82.6}（射影代数簇）：不可约的射影代数集称为\textbf{射影代数簇}（projective variety）。

\textbf{定义 82.7}（齐次坐标环）：射影簇 \(X \subseteq \mathbb{P}^n\) 的\textbf{齐次坐标环}为

\[S(X) = k[x_0, \ldots, x_n] / I(X)\]

\(S(X)\) 是分次 \(k\)-代数：\(S(X) = \bigoplus_{d \geq 0} S_d\)。与仿射情形不同，\(S(X)\) 的分次元素不是 \(X\) 上的函数（因为齐次坐标不是良定义的），但 \(S_d\) 中两个元素的比值是 \(X\) 上的有理函数。

\subsection{82.2 有理映射与双有理等价}\label{ux6709ux7406ux6620ux5c04ux4e0eux53ccux6709ux7406ux7b49ux4ef7}

\textbf{定义 82.8}（有理映射）：设 \(X\) 是代数簇。\(X\) 到 \(Y\) 的\textbf{有理映射} \(\varphi: X \dashrightarrow Y\) 是定义在 \(X\) 的某个非空开子集 \(U\) 上的态射。两个有理映射视为相等，如果它们在交上一致。

\textbf{定义 82.9}（双有理等价）：代数簇 \(X\) 和 \(Y\) 称为\textbf{双有理等价}的（记作 \(X \sim_{\text{bir}} Y\)），如果存在有理映射 \(\varphi: X \dashrightarrow Y\) 和 \(\psi: Y \dashrightarrow X\)，使得 \(\psi \circ \varphi = \operatorname{id}_X\) 和 \(\varphi \circ \psi = \operatorname{id}_Y\)（作为有理映射）。

\textbf{命题 82.2}：\(X\) 和 \(Y\) 双有理等价当且仅当它们的有理函数域同构：\(k(X) \cong k(Y)\)（作为 \(k\) 的扩域）。

\textbf{定义 82.10}（双有理几何）：双有理几何研究代数簇在双有理等价下的不变量。维数是双有理不变量。光滑性不是双有理不变量（奇点解消！）。

\textbf{定理 82.3}（Hironaka 奇点解消，陈述）：特征 \(0\) 的域上的每个代数簇都双有理等价于一个光滑射影代数簇。

\textbf{证明}（框架）：核心步骤分为两阶段。第一阶段（奇点解消）：利用容许中心（admissible center）的概念，在奇点集上进行爆破（blow-up）。关键不变量是 Hilbert-Samuel 函数，它在每次爆破下严格下降------这是 Noether 归纳的基础。第二阶段（紧化）：利用 Nagata 紧化定理将解消后的光滑簇嵌入射影空间。Hironaka 证明的关键创新是：在每一归纳步骤中，选取极大接触（maximal contact）超曲面，将问题降维。归纳同时在局部环的余维数上实施，确保算法在有限步内终止于光滑簇。\(\blacksquare\)

\subsection{82.3 Veronese 嵌入与 Segre 嵌入}\label{veronese-ux5d4cux5165ux4e0e-segre-ux5d4cux5165}

\textbf{定义 82.11}（Veronese 嵌入）：\(d\) 次 \textbf{Veronese 嵌入} \(\nu_d: \mathbb{P}^n \to \mathbb{P}^N\)（\(N = \binom{n+d}{d} - 1\)）由所有 \(d\) 次单项式给出：

\[\nu_d(x_0 : \cdots : x_n) = (\cdots : x^I : \cdots)\]

其中 \(x^I\) 遍历所有 \(d\) 次单项式。\(\nu_d\) 是闭嵌入，其像由二次关系定义。

\textbf{定义 82.12}（Segre 嵌入）：\textbf{Segre 嵌入} \(\sigma: \mathbb{P}^n \times \mathbb{P}^m \to \mathbb{P}^{nm + n + m}\) 由

\[\sigma((x_0 : \cdots : x_n), (y_0 : \cdots : y_m)) = (\cdots : x_i y_j : \cdots)_{0 \leq i \leq n, 0 \leq j \leq m}\]

给出。\(\sigma\) 是闭嵌入，将积 \(\mathbb{P}^n \times \mathbb{P}^m\) 实现为 \(\mathbb{P}^{nm+n+m}\) 中的射影簇。

\textbf{命题 82.4}：射影簇的积是射影簇（通过 Segre 嵌入）。仿射簇的积是仿射簇（\(\mathbb{A}^n \times \mathbb{A}^m \cong \mathbb{A}^{n+m}\)）。

\hrulefill

\section{第282章：Gröbner基与计算代数几何}\label{ux7b2c282ux7ae0gruxf6bnerux57faux4e0eux8ba1ux7b97ux4ee3ux6570ux51e0ux4f55}

Gröbner基理论是计算代数几何的核心工具，由Buchberger于1965年在其博士论文中建立（导师为Gröbner）。它为多项式理想的计算提供了算法基础------理想成员判定、消元理论、解多项式方程组等均可通过Gröbner基实现。本章从Hilbert基定理与项序理论出发（证明Dickson引理------每个单项式理想都是有限生成的，这保证了Gröbner基的存在性），建立Buchberger算法与S-多项式的判别标准，讨论消元理论与解多项式方程组的应用，最后给出Hilbert零点定理的计算版本------通过Gröbner基可判定一个多项式是否属于给定理想的根。Gröbner基将代数几何从存在性理论转变为可计算理论，是符号计算领域最重大的成就之一。

\subsection{Hilbert基定理与项序理论}\label{hilbertux57faux5b9aux7406ux4e0eux9879ux5e8fux7406ux8bba}

\textbf{定义219.1（项序）} 在多项式环 \(k[x_1, \ldots, x_n]\) 的单項式集 \(\{x^\alpha = x_1^{\alpha_1} \cdots x_n^{\alpha_n} : \alpha \in \mathbb{N}^n\}\) 上的全序 \(\prec\) 称为\textbf{项序}（term order），如果：（i）\(\prec\) 是良序（每个非空子集有最小元素）；（ii）\(x^\alpha \prec x^\beta \Rightarrow x^{\alpha + \gamma} \prec x^{\beta + \gamma}\)（与乘法相容）。三项最重要的项序为：字典序（lex）、分次字典序（grlex）和分次逆字典序（grevlex）。

\textbf{定理219.2（Dickson引理）} 每个单项式理想 \(I = \langle x^{\alpha} : \alpha \in A \rangle \subseteq k[x_1, \ldots, x_n]\) 是有限生成的。即存在有限子集 \(A_0 \subseteq A\) 使得 \(I = \langle x^{\alpha} : \alpha \in A_0 \rangle\)。

\textbf{证明} 对变元数 \(n\) 归纳。

\textbf{基础} \(n = 1\)：\(k[x]\) 是主理想整环。单項式理想形如 \(\langle x^d \rangle\)，其中 \(d = \min\{\alpha : x^\alpha \in I\}\)，由一个生成元生成。

\textbf{归纳}：设 \(n > 1\) 且对 \(n-1\) 个变元成立。将单项式按 \(x_n\) 的幂次分组：\(I = \bigcup_{k=0}^\infty \langle x^\alpha x_n^k : \alpha \in A_k \rangle\)（其中 \(A_k \subseteq \mathbb{N}^{n-1}\)）。令 \(J_k = \langle x^\alpha : \alpha \in A_k \rangle \subseteq k[x_1, \ldots, x_{n-1}]\) 是由从 \(x_n^k\) 的系数提取的单项式生成的理想。由 \(J_1 \subseteq J_2 \subseteq \cdots\) 的升链条件（\(k[x_1, \ldots, x_{n-1}]\) 是Noether环，由Hilbert基定理），该升链稳定：存在 \(m\) 使 \(J_k = J_m\) 对所有 \(k \geq m\)。由归纳假设，每个 \(J_0, \ldots, J_m\) 由有限多个单项式生成。这些有限生成元（乘以适当的 \(x_n\) 幂）给出了 \(I\) 的有限生成元集。\(\blacksquare\)

\textbf{推论219.3（Gröbner基的存在性）} 固定项序 \(\prec\)。对任意理想 \(I \subseteq k[x_1, \ldots, x_n]\)，定义其\textbf{首项理想} \(\operatorname{LT}(I) = \langle \operatorname{LT}(f) : f \in I \rangle\)。由Dickson引理，\(\operatorname{LT}(I)\) 由有限多个首项生成：\(\operatorname{LT}(I) = \langle \operatorname{LT}(g_1), \ldots, \operatorname{LT}(g_s) \rangle\)。则 \(\{g_1, \ldots, g_s\}\) 是 \(I\) 的\textbf{Gröbner基}------任意 \(f \in I\) 均可通过 \(g_i\) 的线性组合（带多项式系数）和多項式除法（按 \(\prec\) 序归约）归约为 \(0\)。

\subsection{Buchberger算法与S-多项式标准}\label{buchbergerux7b97ux6cd5ux4e0es-ux591aux9879ux5f0fux6807ux51c6}

\textbf{定义219.4（S-多项式）} 设 \(f, g \in k[x_1, \ldots, x_n]\) 为非零多项式。它们的\textbf{S-多项式}定义为

\[S(f, g) = \frac{x^\gamma}{\operatorname{LT}(f)} \cdot f - \frac{x^\gamma}{\operatorname{LT}(g)} \cdot g\]

其中 \(x^\gamma = \operatorname{lcm}(\operatorname{LM}(f), \operatorname{LM}(g))\) 是首项单项式的最小公倍。S-多项式消去了 \(f\) 和 \(g\) 的首项，旨在捕获由 \(f\) 和 \(g\) 生成的理想的可能的新的首项约化关系。

\textbf{定理219.5（Buchberger判别法）} 有限集 \(G = \{g_1, \ldots, g_s\} \subseteq I\) 是理想 \(I\) 的Gröbner基当且仅当对所有 \(i \neq j\)，\(S(g_i, g_j)\) 经 \(G\) 归约后的余式为 \(0\)（即 \(S(g_i, g_j) \xrightarrow{G}_+ 0\)）。

\textbf{证明}（框架）必要性显然：若 \(G\) 是Gröbner基，则任意 \(S(g_i, g_j) \in I\) 经 \(G\) 归约后余式为 \(0\)。充分性的证明基于：若某多项式的归约出现''歧义''（两个不同的归约路径产生不同的余式），则该歧义可归结为某S-多项式的归约问题------这是Buchberger的''带''引理（Buchberger's chain criterion）的实质。通过对多项式次数使用归纳法，证明首项理想的生成关系完全由S-多项式的归约捕捉。\(\blacksquare\)

\textbf{算法219.6（Buchberger算法）}

输入：\(F = \{f_1, \ldots, f_r\} \subseteq k[x_1, \ldots, x_n]\)（理想 \(I = \langle F \rangle\) 的生成元集）。

输出：\(I\) 的Gröbner基 \(G\)。

（1）\(G \leftarrow F\)。 （2）对每对 \(p \neq q \in G\)，计算 \(S(p, q)\) 经当前 \(G\) 归约后的余式 \(r\)。 （3）若 \(r \neq 0\)，则 \(G \leftarrow G \cup \{r\}\)，返回步骤（2）。 （4）当所有 \(S(p, q)\) 的余式均为 \(0\) 时停止。

该算法的终止性由Dickson引理保证：每添加一个非零余式 \(r\)，理想 \(\langle \operatorname{LT}(G) \rangle\) 严格扩大。由Noether性质，此扩大过程只能有有限步。

\subsection{消元理论与解多项式方程组}\label{ux6d88ux5143ux7406ux8bbaux4e0eux89e3ux591aux9879ux5f0fux65b9ux7a0bux7ec4}

\textbf{定理219.7（消元定理）} 设 \(I \subseteq k[x_1, \ldots, x_n]\) 是理想，\(G\) 是 \(I\) 关于字典序（\(x_1 \succ x_2 \succ \cdots \succ x_n\)）的Gröbner基。则对每个 \(k = 1, \ldots, n-1\)，\(G \cap k[x_{k+1}, \ldots, x_n]\) 是第 \(k\) 个消元理想 \(I_k = I \cap k[x_{k+1}, \ldots, x_n]\) 的Gröbner基。

\textbf{证明} 字典序的性质：若 \(f \in k[x_{k+1}, \ldots, x_n]\) 且 \(f \in I\)，则 \(f\) 经 \(G\) 的字典序归约必为 \(0\)。由于 \(f\) 不含 \(x_1, \ldots, x_k\)，归约过程中涉及的 \(g \in G\) 的首项也必须不含 \(x_1, \ldots, x_k\)，即 \(g \in G \cap k[x_{k+1}, \ldots, x_n]\)。故 \(G \cap k[x_{k+1}, \ldots, x_n]\) 是 \(I_k\) 的Gröbner基。\(\blacksquare\)

\textbf{推论219.8（三角化求解法）} 对 \(k = n-1, n-2, \ldots, 1\) 依次应用消元定理，方程组 \(f_1 = \cdots = f_r = 0\) 转化为三角形式：先解 \(I_{n-1} \subseteq k[x_n]\)（单变元理想），代入解到 \(I_{n-2}\) 中求解 \(x_{n-1}\)，依次回代。此即Gauss消元法的多项式推广。

\subsection{Hilbert零点定理的计算版本}\label{hilbertux96f6ux70b9ux5b9aux7406ux7684ux8ba1ux7b97ux7248ux672c}

\textbf{定理219.9（理想成员判定与根成员判定）} 设 \(I = \langle f_1, \ldots, f_r \rangle \subseteq k[x_1, \ldots, x_n]\)（\(k\) 代数闭域），\(G\) 是 \(I\) 的Gröbner基。

（i）\(f \in I\) 当且仅当 \(f\) 经 \(G\) 归约后的余式为 \(0\)。

（ii）\(f \in \sqrt{I}\) 当且仅当 \(1 \in \langle f_1, \ldots, f_r, 1 - y f \rangle \subseteq k[x_1, \ldots, x_n, y]\)（Rabinowitsch技巧，其中 \(y\) 是新变元）。后者可通过计算Gröbner基来判定。

\textbf{证明}（i）Gröbner基的定义保证：归约余式唯一且为 \(0\) 当且仅当 \(f \in I\)。（ii）由Hilbert零点定理，\(f \in \sqrt{I}\) 等价于 \(V(I) \subseteq V(f)\)，即方程组 \(f_1 = \cdots = f_r = 0\) 且 \(f \neq 0\) 无解。引入新变元 \(y\) 将 \(f \neq 0\) 转化为 \(1 - y f = 0\)，则 \(V(I) \subseteq V(f)\) 等价于扩充系统 \(f_1 = \cdots = f_r = 1 - y f = 0\) 无解，由弱零点定理等价于 \(1 \in \langle f_1, \ldots, f_r, 1 - y f \rangle\)。\(\blacksquare\)

\hrulefill

\section{第283章：代数簇的奇点与解消}\label{ux7b2c283ux7ae0ux4ee3ux6570ux7c07ux7684ux5947ux70b9ux4e0eux89e3ux6d88}

代数簇在大多数点处是''光滑''的（局部同构于仿射空间），但在部分点处可能存在奇点------切空间维数异常、局部环非正则之处。奇点理论是代数几何中最深刻的分支之一，其核心问题是：任何代数簇能否通过某种操作（如爆破）转化为光滑簇？Hironaka于1964年证明了特征零域上此问题的正面回答（荣获1970年菲尔兹奖）。本章从Zariski切空间与正则局部环出发，证明奇点集的结构定理（奇点集为真Zariski闭子集），引入正规化与奇点解消的基本概念，最后陈述Hironaka定理并阐释其在双有理几何中的核心意义。

\subsection{Zariski切空间与正则局部环}\label{zariskiux5207ux7a7aux95f4ux4e0eux6b63ux5219ux5c40ux90e8ux73af}

\textbf{定义220.1（Zariski切空间）} 设 \(X\) 是代数簇，\(p \in X\)。令 \(\mathfrak{m}_p \subset \mathcal{O}_{X, p}\) 为局部环的极大理想（即 \(p\) 处消去的正则函数芽）。\(p\) 处的\textbf{Zariski切空间} \(T_p X\) 定义为 \(\mathfrak{m}_p / \mathfrak{m}_p^2\) 的对偶空间：

\[T_p X = (\mathfrak{m}_p / \mathfrak{m}_p^2)^*\]

其维数 \(\dim T_p X\) 等于 \(X\) 在 \(p\) 处的嵌入维数（即 \(X\) 局部嵌入 \(\mathbb{A}^N\) 所需的最小外围空间维数）。

\textbf{定理220.2（维数不等式）} 对代数簇 \(X\) 上的任意点 \(p\)，

\[\dim X \leq \dim T_p X\]

且等号成立当且仅当局部环 \(\mathcal{O}_{X, p}\) 是正则局部环（即 \(\dim \mathcal{O}_{X, p} = \dim \mathfrak{m}_p / \mathfrak{m}_p^2\)）。

\textbf{证明} \(\dim X = \dim \mathcal{O}_{X, p}\)（由定义）。\(\dim \mathcal{O}_{X, p} \leq \dim \mathfrak{m}_p / \mathfrak{m}_p^2\) 是Noether局部环的一般性质（嵌入维数 \(\geq\) Krull维数，由Krull交定理和主理想定理）。正则局部环的定义恰好是等号成立。正规簇上的正则点等价于光滑点。\(\blacksquare\)

\textbf{定义220.3（光滑点与奇点）} 点 \(p \in X\) 称为\textbf{光滑点}（或非奇异点），若 \(\dim T_p X = \dim X\)；否则称为\textbf{奇点}。\(X\) 的\textbf{奇点集}记为 \(\operatorname{Sing}(X) = \{p \in X : p \text{ 是奇点}\}\)。

\subsection{奇点集的结构}\label{ux5947ux70b9ux96c6ux7684ux7ed3ux6784}

\textbf{定理220.4（奇点集为真Zariski闭子集）} 设 \(X\) 是不可约代数簇。则 \(\operatorname{Sing}(X)\) 是 \(X\) 的真Zariski闭子集。特别地，光滑点构成 \(X\) 的一个非空开稠密子集（\(X_{\text{sm}} = X \setminus \operatorname{Sing}(X)\)）。

\textbf{证明}（基于Jacobi判别法）

\textbf{第1步（Jacobi矩阵的秩条件）}：设 \(X \subseteq \mathbb{A}^n\) 是仿射簇，\(I(X) = \langle f_1, \ldots, f_r \rangle\)。在点 \(p \in X\)，Zariski切空间 \(T_p X\) 可等同于Jacobi矩阵 \(J(p) = (\partial f_i / \partial x_j)(p)\) 的零空间。\(p\) 光滑等价于 \(\operatorname{rank} J(p) = n - \dim X\)（Jacobi矩阵的秩达到最大可能值------余维数）。

\textbf{第2步（秩条件的代数刻画）}：考虑所有 \((n - \dim X + 1) \times (n - \dim X + 1)\) 子式的理想。这些子式定义了 \(X\) 上秩严格小于 \(n - \dim X\) 的点的闭集。此闭集恰为 \(\operatorname{Sing}(X)\)。因此 \(\operatorname{Sing}(X)\) 是 \(X\) 的Zariski闭子集。

\textbf{第3步（真子集性）}：需证 \(\operatorname{Sing}(X) \subsetneq X\)。考虑 \(X\) 的有理函数域 \(k(X)\)。由域的分离超越基的性质，\(k(X)\) 作为 \(k\) 的有限生成扩域，存在分离超越基。在分离超越基的代数闭包对应的一般点处，\(X\) 是光滑的。因此光滑点集非空。

\textbf{第4步（一般光滑性）}：令 \(d = \dim X\)。存在一个 \(d \times n\) 的Jacobi子矩阵，其行列式在 \(k(X)\) 中非零。该行列式在 \(X\) 的某个非空开子集上非零，故光滑点集包含一个非空开集，\(\operatorname{Sing}(X)\) 为真闭子集。\(\blacksquare\)

\subsection{正规化与奇点解消}\label{ux6b63ux89c4ux5316ux4e0eux5947ux70b9ux89e3ux6d88}

\textbf{定义220.5（正规化）} 代数簇 \(X\) 的\textbf{正规化}是有限双有理态射 \(\nu: \tilde{X} \to X\)，其中 \(\tilde{X}\) 是正规簇（每个局部环是整闭整环），且 \(\nu\) 满足泛性质：任何从正规簇到 \(X\) 的支配态射唯一通过 \(\nu\) 分解。对仿射簇 \(X = \operatorname{Spec} A\)，\(\tilde{X} = \operatorname{Spec} \tilde{A}\)，其中 \(\tilde{A}\) 是 \(A\) 在其分式域中的整闭包。

\textbf{定理220.6（曲线的奇点解消）} 设 \(C\) 是代数闭域上的不可约代数曲线。则存在光滑射影曲线 \(\tilde{C}\) 和双有理态射 \(\pi: \tilde{C} \to C\)。\(\tilde{C}\) 在同构意义下唯一（即为 \(C\) 的非奇异紧化），称为 \(C\) 的\textbf{奇点解消}（或非奇异模型）。

\textbf{证明}（框架，对曲线的爆破归纳）曲线的奇点是有限集。在每个奇点处进行爆破（blow-up）：爆破后的例外除子是 \(\mathbb{P}^1\) 的并（最多），炸开后奇点被替换为较低重数的点。由于曲线是 \(1\) 维的，每次爆破严格降低奇点的重数或嵌入维数。有限的爆破次数后得到处处光滑曲线。紧化由将曲线嵌入某 \(\mathbb{P}^n\) 并取Zariski闭包实现。\(\blacksquare\)

\textbf{定理220.7（Hironaka奇点解消定理，1964）} 设 \(X\) 是特征 \(0\) 的代数闭域上的代数簇（不一定射影）。则存在光滑代数簇 \(\tilde{X}\) 和真双有理态射 \(\pi: \tilde{X} \to X\)（奇点解消），满足：（i）\(\pi\) 在 \(X_{\text{sm}}\)（光滑点集）上是同构；（ii）\(\pi^{-1}(\operatorname{Sing}(X))\) 是 \(\tilde{X}\) 中的简单正规交叉除子（即例外除子由光滑超曲面组成，任意交为横截）；（iii）\(\pi\) 通过一系列在光滑中心的爆破构造。

\textbf{证明}（结构框架，非完整证明）Hironaka证明的核心是''容许中心''（admissible center）的选取和Hilbert-Samuel函数的严格下降。证明分为两个嵌套的归纳：（a）对外围空间（或簇的嵌入维数）的维数进行归纳；（b）在固定维数内对极大接触（maximal contact）超曲面的阶数进行归纳。通过局部坐标的Weierstrass准备定理将奇点降维，每次爆破使不变量（Hilbert-Samuel函数在字典序下）严格下降，由良基归纳保证有限步终止。\(\blacksquare\)

Hironaka定理是代数几何中里程碑式的成就，其特征零的假设是本质的------正特征域中的奇点解消至今仍是开放问题（除维数 \(\leq 3\) 被Abhyankar和Cossart-Piltant解决外）。奇点解消在双有理几何、极小模型纲领、Hodge理论和D模理论中发挥着基础作用，其算法版本（如Villamayor、Bierstone-Milman、Wlodarczyk等人的工作）已发展为有效的计算工具。

\hrulefill

\hrulefill

\hrulefill

\hrulefill

\hrulefill

\newpage

\chapter{08-代数几何与数论/01-代数几何/02-概形理论.md}\label{ux4ee3ux6570ux51e0ux4f55ux4e0eux6570ux8bba01-ux4ee3ux6570ux51e0ux4f5502-ux6982ux5f62ux7406ux8bba.md}

\section{第284章：概形理论}\label{ux7b2c284ux7ae0ux6982ux5f62ux7406ux8bba}

概形（Scheme）是 Grothendieck 在 1960 年代引入的代数几何的基本对象，它推广了代数簇的概念，统一了代数几何、数论和交换代数。概形理论的核心思想是：将交换环视为某个空间上的''函数环''，而该空间就是环的素理想谱。这一视角的转变是 20 世纪数学最深刻的革命之一------在此之前，代数几何研究的是多项式方程组的零点集（经典代数簇），而 Grothendieck 的概形理论将研究对象从''零点''推广到''素理想''，从而使得超越基域限制的算术几何成为可能。

\subsection{83.1 仿射概形}\label{ux4effux5c04ux6982ux5f62}

\textbf{定义 83.1}（交换环的谱）：设 \(A\) 是交换环（含幺）。\(A\) 的\textbf{谱} \(\operatorname{Spec} A\) 定义为 \(A\) 的所有素理想的集合：

\[\operatorname{Spec} A = \{\mathfrak{p} \subseteq A : \mathfrak{p} \text{ 是素理想}\}\]

\(\operatorname{Spec} A\) 配以 \textbf{Zariski 拓扑}：闭集为 \(V(I) = \{\mathfrak{p} \in \operatorname{Spec} A : I \subseteq \mathfrak{p}\}\)（\(I\) 为 \(A\) 的理想）。注意 \(V(A) = \varnothing\)，\(V(0) = \operatorname{Spec} A\)，\(V(I) \cup V(J) = V(IJ)\)，\(\bigcap V(I_\alpha) = V(\sum I_\alpha)\)。

\textbf{定义 83.2}（结构层）：\(\operatorname{Spec} A\) 上的\textbf{结构层} \(\mathcal{O}_{\operatorname{Spec} A}\) 由以下方式定义：对每个开集 \(U \subseteq \operatorname{Spec} A\)，\(\mathcal{O}(U)\) 由那些在 \(U\) 的每点局部地可表示为 \(A\) 中元素之商（分母不在该点对应的素理想中）的函数构成。具体地，对 \(f \in A\)，定义基本开集 \(D(f) = \{\mathfrak{p} \in \operatorname{Spec} A : f \notin \mathfrak{p}\}\)，则 \(\mathcal{O}(D(f)) = A_f\)（\(A\) 在 \(f\) 处的局部化）。\(\{D(f)\}_{f \in A}\) 构成 Zariski 拓扑的基。

\textbf{定义 83.3}（仿射概形）：环层空间 \((\operatorname{Spec} A, \mathcal{O}_{\operatorname{Spec} A})\) 称为 \textbf{仿射概形}。环层空间 \((X, \mathcal{O}_X)\) 若局部同构于仿射概形，则称为\textbf{概形}。

\textbf{定理 83.1}（仿射概形的基本性质）： 1. \(\operatorname{Spec} A\) 是拟紧的（Zariski 拓扑下任意开覆盖有有限子覆盖） 2. \(\mathcal{O}_{\operatorname{Spec} A}\) 在茎处的局部环为 \(\mathcal{O}_{\mathfrak{p}} = A_{\mathfrak{p}}\)（\(A\) 在素理想 \(\mathfrak{p}\) 处的局部化） 3. \(\Gamma(\operatorname{Spec} A, \mathcal{O}) = A\)（整体截面函子恢复原环） 4. 仿射概形之间的态射 \(\operatorname{Spec} A \to \operatorname{Spec} B\) 与环同态 \(B \to A\) 一一对应（反变等价）

\textbf{证明}：(1) 对开覆盖 \(\operatorname{Spec} A = \bigcup D(f_i)\)，\(1 \in \sum (f_i)\)，故 \(1 = \sum g_i f_i\)（有限和），则 \(\{D(f_i)\}\) 的有限子族即覆盖。(2) \(\mathcal{O}_{\mathfrak{p}} = \varinjlim_{f \notin \mathfrak{p}} A_f = A_{\mathfrak{p}}\)。(3) 由层的定义和 \(A \to \prod A_f\)（\(f\) 遍历所有元素）的正合性。(4) 整体截面函子给出环同态 \(B \to \Gamma(\operatorname{Spec} A, \mathcal{O}) = A\)，反之环同态引发谱的连续映射和结构层的层同态。\(\blacksquare\)

仿射概形是概形理论中最基本的构造块，其核心思想是：将交换环 \(A\) 的所有素理想视为''点''，赋予 Zariski 拓扑和结构层，使得环 \(A\) 可以''几何化''为空间 \(\operatorname{Spec} A\)。这一构造将代数几何从研究多项式方程组的零点集（经典代数簇）推广到研究任意交换环的素理想谱------包括 \(\mathbb{Z}\)、\(\mathbb{Z}_p\)（\(p\)-进整数）、\(\mathbb{F}_p[t]\) 等，从而将数论和代数几何统一在同一个框架下。在 \(\operatorname{Spec} \mathbb{Z}\) 中，素理想 \((p)\) 对应于''曲线''上的''点''（有限域 \(\mathbb{F}_p\)），而泛点 \((0)\) 对应于''一般点''（有理数域 \(\mathbb{Q}\)）------这一类比是 Arakelov 几何和算术代数几何的出发点。仿射概形 \(\operatorname{Spec} A\) 的 Zariski 拓扑虽然粗糙（非 Hausdorff），但其''拟紧性''和''不可约分支''的概念恰好反映了 \(A\) 的代数性质：\(\operatorname{Spec} A\) 不可约当且仅当 \(A\) 的幂零根 \(\sqrt{(0)}\) 为素理想，即 \(A/\sqrt{(0)}\) 为整环。

\textbf{定义 83.3'（概形的开子概形与闭子概形）}：设 \((X, \mathcal{O}_X)\) 是概形。\textbf{开子概形}是形如 \((U, \mathcal{O}_X|_U)\) 的概形（\(U \subseteq X\) 为开集）。\textbf{闭子概形}由拟凝聚理想层 \(\mathcal{I} \subseteq \mathcal{O}_X\) 定义：其底空间为 \(\operatorname{Supp}(\mathcal{O}_X/\mathcal{I})\)，结构层为 \(\mathcal{O}_X/\mathcal{I}\)。闭子概形对应于概形的''子簇''类比。

\textbf{定理 83.1.1（幂零根与约化概形）}：设 \(X = \operatorname{Spec} A\)。\(X\) 是\textbf{约化}的（即 \(\mathcal{O}_X\) 无幂零元）当且仅当 \(A\) 是约化环（\(\sqrt{(0)} = (0)\)）。\(X_{\text{red}} = \operatorname{Spec}(A/\sqrt{(0)})\) 是 \(X\) 的\textbf{约化闭子概形}，且 \(X_{\text{red}} \hookrightarrow X\) 是闭浸入，底空间是同胚。约化概形在代数几何中对应于''没有多重结构的经典代数簇''。在模空间理论中，约化结构用于去除''非约化''的 Artin 环结构（如 \(\operatorname{Spec} k[\varepsilon]/(\varepsilon^2)\) 中的''厚点''）。

\textbf{证明}：首先，\(X\) 约化意味着对每个开集 \(U\)，\(\mathcal{O}_X(U)\) 无幂零元。特别地，取整体截面 \(A = \Gamma(X, \mathcal{O}_X)\)，若 \(X\) 约化则 \(A\) 无幂零元，即 \(\sqrt{(0)} = (0)\)。反之，若 \(A\) 约化，则对每个 \(f \in A\)，局部化 \(A_f\) 也约化（因为 \(A_f\) 中的幂零元必来自 \(A\) 中幂零元除以 \(f\) 的幂），故 \(\mathcal{O}_X(D(f)) = A_f\) 无幂零元。由于 \(\{D(f)\}\) 构成拓扑基，茎 \(\mathcal{O}_{X, \mathfrak{p}} = A_{\mathfrak{p}}\) 也约化，故 \(X\) 约化。其次，商映射 \(A \to A/\sqrt{(0)}\) 诱导闭浸入 \(\operatorname{Spec}(A/\sqrt{(0)}) \hookrightarrow \operatorname{Spec} A\)。由于任何素理想都包含 \(\sqrt{(0)}\)（幂零根是所有素理想的交），映射 \(\mathfrak{p} \mapsto \mathfrak{p}/\sqrt{(0)}\) 给出 \(\operatorname{Spec} A\) 与 \(\operatorname{Spec}(A/\sqrt{(0)})\) 之间的双射，且闭集 \(V(I)\) 在此双射下对应于 \(V((I+\sqrt{(0)})/\sqrt{(0)})\)，故为同胚。结构层 \(\mathcal{O}_{X_{\text{red}}}\) 是 \(\mathcal{O}_X/\mathcal{N}\)（其中 \(\mathcal{N}\) 是幂零根层），显然是约化的。\(\blacksquare\)

\subsection{83.2 层与环层空间}\label{ux5c42ux4e0eux73afux5c42ux7a7aux95f4}

\textbf{定义 83.4}（环层空间）：\textbf{环层空间} \((X, \mathcal{O}_X)\) 是拓扑空间 \(X\) 配以交换环层 \(\mathcal{O}_X\)（结构层）。\((X, \mathcal{O}_X)\) 上点 \(x\) 的\textbf{茎} \(\mathcal{O}_{X,x}\) 是局部环（局部环层空间）。概形是局部环层空间，局部同构于仿射概形。

\textbf{定义 83.5}（\(\mathcal{O}_X\)-模层）：\(\mathcal{O}_X\)-模层 \(\mathcal{F}\) 是 Abel 群层，使得每个 \(\mathcal{F}(U)\) 是 \(\mathcal{O}_X(U)\)-模且限制映射与模结构兼容。\(\mathcal{O}_X\)-模层构成 Abel 范畴 \(\mathcal{O}_X\text{-}\mathbf{Mod}\)。

\textbf{定义 83.6}（拟凝聚层与凝聚层）：\(\mathcal{O}_X\)-模层 \(\mathcal{F}\) 称为\textbf{拟凝聚的}，如果局部地存在正合列 \(\mathcal{O}_X^{\oplus J}|_U \to \mathcal{O}_X^{\oplus I}|_U \to \mathcal{F}|_U \to 0\)。若 \(I, J\) 可取有限，则 \(\mathcal{F}\) 称为\textbf{凝聚层}。

\textbf{定理 83.2}（仿射概形上层与模的等价）：设 \(X = \operatorname{Spec} A\)。则函子 \(\widetilde{M \mapsto M}\)（将 \(A\)-模 \(M\) 对应为 \(\mathcal{O}_X\)-模层 \(\widetilde{M}\)）给出 \(A\text{-}\mathbf{Mod}\) 与 \(\operatorname{QCoh}(X)\)（拟凝聚层范畴）之间的等价。特别地，\(\Gamma(X, \widetilde{M}) = M\)。

\textbf{证明}：对 \(A\)-模 \(M\)，定义 \(\widetilde{M}(D(f)) = M_f\)（\(M\) 在 \(f\) 处的局部化），然后延拓为 \(X\) 上的层。验证 \(\widetilde{M}\) 是拟凝聚层。函子 \(\Gamma(X, -)\) 给出逆：\(\Gamma(X, \mathcal{F})^{\sim} \cong \mathcal{F}\) 当 \(\mathcal{F}\) 拟凝聚时。核心在于：对仿射概形，层上同调由 Serre 定理消没，且整体截面函子忠实地反映了层的结构。\(\blacksquare\)

\textbf{定理 83.2.1（Serre 仿射性判别）}：概形 \(X\) 是仿射的当且仅当 \(X\) 是拟紧的，且存在 \(X\) 上的凝聚层 \(\mathcal{F}\) 使得 \(X\) 同构于 \(\operatorname{Spec}(\Gamma(X, \mathcal{O}_X))\)，且 \(H^i(X, \mathcal{F}) = 0\) 对所有 \(i > 0\) 和任意拟凝聚层 \(\mathcal{F}\) 成立。这是 Serre 的著名判别法，将仿射性归结为上同调的消没性质。

\subsection{83.3 射影概形}\label{ux5c04ux5f71ux6982ux5f62}

\textbf{定义 83.7}（射影概形 \(\operatorname{Proj}\)）：设 \(S = \bigoplus_{d \geq 0} S_d\) 是分次环（\(S_0 = A\) 为环，\(S\) 由 \(S_1\) 作为 \(A\)-代数有限生成）。\(\operatorname{Proj} S\) 是由 \(S\) 中不包含无关理想 \(S_+ = \bigoplus_{d > 0} S_d\) 的齐次素理想构成的集合。\(\operatorname{Proj} S\) 是概形，其结构层由齐次局部化定义：对 \(f \in S_d\)，\(D_+(f) = \{\mathfrak{p} \in \operatorname{Proj} S : f \notin \mathfrak{p}\} \cong \operatorname{Spec} S_{(f)}\)（其中 \(S_{(f)}\) 是 \(S_f\) 的 \(0\) 次部分）。

\textbf{定理 83.3}（\(\operatorname{Proj}\) 的基本性质）： 1. \(\operatorname{Proj} \mathbb{Z}[x_0, \ldots, x_n] = \mathbb{P}^n_{\mathbb{Z}}\) 2. \(\operatorname{Proj} k[x_0, \ldots, x_n] = \mathbb{P}^n_k\) 3. \(\operatorname{Proj} S\) 是真（proper）概形当 \(S\) 是有限生成 \(A\)-代数 4. Serre 捻层 \(\mathcal{O}(d)\) 对应于分次 \(S\)-模 \(S(d)\)（次数平移）

\textbf{证明}：(1) \(\operatorname{Proj} \mathbb{Z}[x_0,\ldots,x_n]\) 的齐次素理想不包含 \((x_0,\ldots,x_n)\) 的条件恰好对应于 \(\mathbb{P}^n_{\mathbb{Z}}\) 的经典定义，而 \(D_+(x_i) \cong \operatorname{Spec} \mathbb{Z}[x_0/x_i,\ldots,x_n/x_i]\) 给出标准仿射覆盖。(2) 将基变换 \(\operatorname{Spec} k \to \operatorname{Spec} \mathbb{Z}\) 应用于 (1)。(3) \(\operatorname{Proj} S\) 的分离性来自对角态射的闭性验证：利用仿射开覆盖上对角态射对应环的满射 \(S_{(f)} \otimes_A S_{(f)} \to S_{(f)}\)。泛闭性由射影态射的赋值判别法（valuative criterion）验证：对任何赋值环 \(R\) 和态射 \(\operatorname{Spec} K \to \operatorname{Proj} S\)（\(K = \operatorname{Frac}(R)\)），存在唯一的提升 \(\operatorname{Spec} R \to \operatorname{Proj} S\)，这由齐次坐标的齐次化给出。(4) Serre 捻层 \(\mathcal{O}(d)\) 定义为 \(S(d)^{\sim}\)，其中 \(S(d)_n = S_{n+d}\)。在 \(D_+(f)\) 上，\(\mathcal{O}(d)(D_+(f)) = S_{(f)}(d) = \{g/f^k : \deg g = kd + \deg f\}\)。\(\blacksquare\)

射影概形 \(\operatorname{Proj} S\) 是概形理论中仿射概形之外最重要的构造，它统一了经典射影代数簇和现代算术几何中的射影空间。\(\operatorname{Proj}\) 与 \(\operatorname{Spec}\) 的根本区别在于：\(\operatorname{Proj} S\) 排除了''无关理想'' \(S_+\)，这对应于射影空间中''无穷远点''在仿射坐标下不可见的事实。在 \(\operatorname{Proj}\) 构造中，分次环 \(S\) 的齐次素理想 \(\mathfrak{p}\) 构成了 \(\operatorname{Proj} S\) 的点，而 \(\mathfrak{p}\) 不包含 \(S_+\) 的条件确保了 \(\operatorname{Proj} S\) 可以被仿射开集 \(\{D_+(f)\}_{f \in S_1}\) 覆盖------这些 \(D_+(f)\) 恰好是 \(\operatorname{Spec} S_{(f)}\)，即射影空间的标准仿射坐标卡。\(\operatorname{Proj}\) 构造在代数几何中用于定义射影丛、Grassmann 概形、Hilbert 概形等基本对象，并通过 Serre 捻层 \(\mathcal{O}(d)\) 为凝聚层上同调的计算提供了标准框架。

\textbf{定理 83.3.1（爆破 / Blow-up）}：设 \(X\) 是概形，\(\mathcal{I} \subseteq \mathcal{O}_X\) 是拟凝聚理想层。\(X\) 沿 \(\mathcal{I}\) 的\textbf{爆破} \(\operatorname{Bl}_{\mathcal{I}}(X) = \operatorname{Proj}(\bigoplus_{d \geq 0} \mathcal{I}^d)\) 是 \(X\) 上的概形，其例外除子（exceptional divisor）对应于 \(\mathcal{I}\) 的零点集。爆破是双有理几何中最基本的操作------它用余维数 1 的子簇（例外除子）替换高余维数的子簇，从而消除奇点或实现极小模型纲领中的''翻转''（flip）操作。

\textbf{证明}：首先验证 \(\mathcal{S} = \bigoplus_{d \geq 0} \mathcal{I}^d\) 是分次 \(\mathcal{O}_X\)-代数层（约定 \(\mathcal{I}^0 = \mathcal{O}_X\)），且 \(\mathcal{I}\) 拟凝聚保证每个 \(\mathcal{I}^d\) 拟凝聚，从而 \(\mathcal{S}\) 是拟凝聚分次代数层。在仿射开集 \(U = \operatorname{Spec} A\) 上，\(\mathcal{I}|_U\) 对应理想 \(I \subseteq A\)，则 \(\operatorname{Bl}_{\mathcal{I}}(X)|_U \cong \operatorname{Proj}(\bigoplus_{d \geq 0} I^d)\)。商映射 \(\mathcal{O}_X \to \mathcal{O}_X/\mathcal{I}\) 诱导闭浸入 \(Z \hookrightarrow X\)（\(Z\) 为 \(\mathcal{I}\) 定义的闭子概形）。在 \(X \setminus Z\) 上，\(\mathcal{I}|_{X \setminus Z} = \mathcal{O}_X|_{X \setminus Z}\)，故 \(\operatorname{Bl}_{\mathcal{I}}(X)|_{X \setminus Z} \cong \operatorname{Proj}(\mathcal{O}_{X \setminus Z}[t]) \cong X \setminus Z\)，即爆破在 \(Z\) 之外是同构。例外除子 \(E = \pi^{-1}(Z)\)（\(\pi: \operatorname{Bl}_{\mathcal{I}}(X) \to X\) 为自然投影）是 \(\operatorname{Proj}(\bigoplus \mathcal{I}^d/\mathcal{I}^{d+1})\)，即射影法锥。\(\blacksquare\)

\subsection{83.4 概形的态射}\label{ux6982ux5f62ux7684ux6001ux5c04}

\textbf{定义 83.8}（概形的态射）：概形之间的态射 \(f: (X, \mathcal{O}_X) \to (Y, \mathcal{O}_Y)\) 由连续映射 \(f: X \to Y\) 和层同态 \(f^\#: \mathcal{O}_Y \to f_* \mathcal{O}_X\) 组成，使得在每点 \(x \in X\)，茎上的诱导映射 \(f^\#_x: \mathcal{O}_{Y, f(x)} \to \mathcal{O}_{X, x}\) 是局部环同态（即极大理想的逆像是极大理想）。

\textbf{定义 83.9}（分离态射）：概形态射 \(f: X \to Y\) 称为\textbf{分离的}，如果对角态射 \(\Delta_{X/Y}: X \to X \times_Y X\) 是闭浸入。\(X\) 称为分离概形，如果 \(X \to \operatorname{Spec} \mathbb{Z}\) 是分离的。分离性在代数几何中对应于拓扑空间中的 Hausdorff 性，但 Zariski 拓扑本身不是 Hausdorff 的------分离性用对角态射的闭性来刻画，这是 Grothendieck 的深刻创新。

\textbf{定义 83.10}（真态射）：\(f: X \to Y\) 称为\textbf{真态射}（proper），如果 \(f\) 是分离的、有限型的，且泛闭的（即对任意基变换 \(Y' \to Y\)，\(X \times_Y Y' \to Y'\) 是闭映射）。真态射是代数几何中''紧性''的类比------射影态射 \(\mathbb{P}^n_Y \to Y\) 是真态射的典型例子，而 \(\mathbb{A}^1 \to \operatorname{Spec} k\) 不是真的（因为仿射直线不是''紧''的）。真态射的另一个关键性质是：真态射下的凝聚层直像是凝聚的（Grothendieck 凝聚性定理），这保证了上同调群的有限性。

\textbf{定义 83.11}（纤维积 / Fibered Product）：概形范畴中存在纤维积。对态射 \(X \to S\) 和 \(Y \to S\)，纤维积 \(X \times_S Y\) 是满足泛性质的概形。在仿射情形，\(\operatorname{Spec} A \times_{\operatorname{Spec} R} \operatorname{Spec} B = \operatorname{Spec}(A \otimes_R B)\)。纤维积是概形范畴中最重要的构造之一，它是基变换（base change）的形式化，也是定义概形族（如曲线族、曲面族）的纤维、平坦态射的纤维的维数相等性等概念的基础。

\textbf{定义 83.12}（光滑态射与 etale 态射）：\(f: X \to Y\) 是\textbf{光滑态射}，如果 \(f\) 局部有限表示，平坦，且几何纤维是光滑的（即正则的）。\(f\) 是 \textbf{etale 态射}，如果 \(f\) 光滑且相对维数为 \(0\)（即几何纤维是 \(0\) 维光滑的，从而是一族离散点）。etale 态射是代数几何中''局部同构''的类比（在微分几何中对应局部微分同胚），但 etale 态射不需要是 Zariski 局部同构------它允许''非分歧覆盖''（如 \(\operatorname{Spec} \mathbb{C} \to \operatorname{Spec} \mathbb{R}\) 或 \(\operatorname{Spec} \mathbb{Z}[\sqrt{2}] \to \operatorname{Spec} \mathbb{Z}\)）。etale 态射是 etale 拓扑和 etale 上同调的基本对象，也是 Weil 猜想证明和 Fermat 大定理证明中不可或缺的工具。

\textbf{定义 83.13（平坦态射）}：态射 \(f: X \to Y\) 是\textbf{平坦的}，如果对每个 \(x \in X\)，\(\mathcal{O}_{X,x}\) 作为 \(\mathcal{O}_{Y,f(x)}\)-模是平坦的。平坦性是代数几何和交换代数中最重要的概念之一------它保证了纤维的''连续性''（例如，平坦族中纤维的 Hilbert 多项式是常数），是模空间理论（如 Hilbert 概形、Picard 概形）和形变理论的基础。``平坦性''一词形象地描述了在平坦态射下，概形的''纤维''不发生''跳跃''------即几何性质（如维数、次数）沿纤维连续变化。

\subsection{83.5 概形上的拟凝聚层}\label{ux6982ux5f62ux4e0aux7684ux62dfux51ddux805aux5c42}

\textbf{定理 83.4}（拟凝聚层的基本性质）：设 \(X\) 是概形。 1. 拟凝聚层范畴 \(\operatorname{QCoh}(X)\) 是 Abel 范畴，且关于核、余核、扩张封闭 2. 若 \(X\) 是 Noether 概形，凝聚层范畴 \(\operatorname{Coh}(X)\) 是 \(\operatorname{QCoh}(X)\) 的 Abel 子范畴 3. 直像函子 \(f_*: \operatorname{QCoh}(X) \to \operatorname{QCoh}(Y)\) 对仿射态射 \(f\) 是正合的 4. 若 \(f: X \to Y\) 是拟紧分离态射，\(R^i f_* \mathcal{F}\) 是拟凝聚的（\(\mathcal{F} \in \operatorname{QCoh}(X)\)）

\textbf{证明}：(1) \(\operatorname{QCoh}(X)\) 是 \(\mathcal{O}_X\text{-}\mathbf{Mod}\) 的全子范畴。对拟凝聚层之间的态射 \(\phi: \mathcal{F} \to \mathcal{G}\)，核 \(\ker \phi\) 在仿射开集 \(\operatorname{Spec} A\) 上对应 \(\widetilde{\ker(\phi_A)}\)（其中 \(\phi_A: \mathcal{F}(\operatorname{Spec} A) \to \mathcal{G}(\operatorname{Spec} A)\)），故拟凝聚。余核类似。扩张封闭性：若 \(0 \to \mathcal{F}' \to \mathcal{F} \to \mathcal{F}'' \to 0\) 中 \(\mathcal{F}', \mathcal{F}''\) 拟凝聚，则局部地 \(\mathcal{F}|_U\) 是 \(\mathcal{F}'|_U\) 和 \(\mathcal{F}''|_U\) 的扩张，由五引理知 \(\mathcal{F}\) 拟凝聚。(2) Noether 概形上凝聚层是局部有限生成的，核和余核的有限生成性由 Noether 性质保证。(3) 仿射态射 \(f\) 可局部化为 \(\operatorname{Spec} B \to \operatorname{Spec} A\)，此时 \(f_*\) 对应 \(A\)-模到 \(B\)-模的标量限制函子，它是正合的（因为忘却函子正合）。(4) 由 Cech 上同调计算：取仿射开覆盖，\(R^i f_* \mathcal{F}\) 的 Cech 复形各项为拟凝聚层，故上同调层也拟凝聚。\(\blacksquare\)

\textbf{定理 83.5}（Serre 定理，概形版本）：设 \(X\) 是 Noether 环 \(A\) 上的射影概形，\(\mathcal{F}\) 是 \(X\) 上的凝聚层。则 1. \(H^i(X, \mathcal{F})\) 是有限 \(A\)-模 2. 存在 \(d_0\) 使得对 \(d \geq d_0\) 和 \(i > 0\)，\(H^i(X, \mathcal{F}(d)) = 0\) 3. \(\mathcal{F}(d)\) 由整体截面生成（\(d \gg 0\)）

\textbf{证明}：将 \(X\) 嵌入射影空间 \(\mathbb{P}^n_A\) 作为闭子概形，令 \(i: X \hookrightarrow \mathbb{P}^n_A\) 为闭浸入。则 \(H^i(X, \mathcal{F}) \cong H^i(\mathbb{P}^n_A, i_* \mathcal{F})\)（因为 \(i_*\) 是正合的且保持凝聚性）。对 \(\mathbb{P}^n_A\)，取标准仿射开覆盖 \(\{D_+(x_j)\}_{j=0}^n\)，计算 Cech 上同调。Cech 复形 \(\check{C}^\bullet\) 的每个项是分次 \(A[x_0,\ldots,x_n]\)-模的局部化，其齐次部分为有限生成 \(A\)-模。取上同调后，\(H^i\) 为这些有限生成模的子商，故为有限 \(A\)-模（证 (1)）。对 \(\mathcal{F}(d)\)，Cech 复形各项在次数 \(d\) 足够大时变为零（因为局部化 \(M_f\) 中次数足够高的元素来自 \(M\) 乘以 \(f\) 的高次幂），故 \(H^i(\mathbb{P}^n_A, i_*\mathcal{F}(d)) = 0\) 对 \(i > 0\) 和 \(d \gg 0\) 成立（证 (2)）。\(d\) 足够大时，\(\mathcal{F}(d)\) 的整体截面在每点处生成茎，因为对每个 \(x \in X\)，存在 \(s \in H^0(X, \mathcal{F}(d))\) 使得 \(s(x) \notin \mathfrak{m}_x \mathcal{F}(d)_x\)（这由 Serre 的对应定理保证：射影概形上凝聚层的整体截面生成性与该层是''丰富''的相关联）（证 (3)）。\(\blacksquare\)

\subsection{83.6 概形理论中的维数论}\label{ux6982ux5f62ux7406ux8bbaux4e2dux7684ux7ef4ux6570ux8bba}

\textbf{定义 83.14}（概形的维数）：概形 \(X\) 的\textbf{维数} \(\dim X\) 定义为 \(X\) 作为拓扑空间的 Krull 维数：\(\dim X = \sup\{\ell : \text{存在不可约闭子集链 } Z_0 \subsetneq Z_1 \subsetneq \cdots \subsetneq Z_\ell \subseteq X\}\)。对于仿射概形 \(\operatorname{Spec} A\)，\(\dim(\operatorname{Spec} A) = \dim A\)（\(A\) 的 Krull 维数）。对于 Noether 概形上的有限型态射 \(f: X \to Y\)，纤维维数满足 \(\dim X_y \geq \dim X - \dim Y\)，且等号在平坦态射下成立（``纤维维数定理''）。

\textbf{定理 83.6（正则局部环与光滑性）}：Noether 局部环 \((A, \mathfrak{m})\) 是\textbf{正则的}，如果 \(\dim A = \dim_{A/\mathfrak{m}} \mathfrak{m}/\mathfrak{m}^2\)。概形 \(X\) 在点 \(x\) 处是\textbf{正则的}（或光滑的），如果 \(\mathcal{O}_{X,x}\) 是正则局部环。在代数闭域上的有限型概形的情形，正则性等价于该点处的 Zariski 切空间的维数 \(\dim_k \mathfrak{m}/\mathfrak{m}^2\) 等于概形的局部维数------这就是经典代数几何中''光滑点''的 Jacobi 判别法。

\textbf{证明}：由 Krull 维数理论，Noether 局部环 \((A, \mathfrak{m})\) 的维数满足 \(\dim A \leq \dim_{A/\mathfrak{m}} \mathfrak{m}/\mathfrak{m}^2\)（嵌入维数不小于 Krull 维数）。等号成立时称 \(A\) 正则。正则局部环是整环（Auslander-Buchsbaum 定理），且是 Cohen-Macaulay 环。在代数闭域 \(k\) 上的有限型概形中，\(\mathcal{O}_{X,x} \cong k[x_1,\ldots,x_n]_{\mathfrak{m}}/I\)，其中 \(I = (f_1,\ldots,f_r)\) 为定义 \(X\) 的局部方程。此时 \(\mathfrak{m}/\mathfrak{m}^2\) 是 \(k\)-向量空间，其维数等于 \(n\) 减去 Jacobi 矩阵 \((\partial f_i/\partial x_j)\) 在 \(x\) 处的秩，而 \(\dim \mathcal{O}_{X,x} = n - \operatorname{ht}(I)\)。正则性等价于 \(\operatorname{ht}(I) = \operatorname{rank} J\)，即 Jacobi 判别法。\(\blacksquare\)

\textbf{定理 83.7（平坦性与维数）}：设 \(f: X \to Y\) 是 Noether 概形之间的平坦态射。则对任意 \(y \in Y\)，\(\dim_x X_y = \dim_x X - \dim_y Y\)（在 \(x\) 处的局部纤维维数）。这一结果将微分几何中''淹没''（submersion）的纤维维数定理推广到代数几何。

\textbf{证明}：设 \(B = \mathcal{O}_{Y,y}\)，\(A = \mathcal{O}_{X,x}\)，则 \(B \to A\) 是平坦局部同态。由 Chevalley 的维数公式，对 Noether 局部环的平坦扩张，\(\dim A = \dim B + \dim(A/\mathfrak{m}_B A)\)。这里 \(\dim(A/\mathfrak{m}_B A) = \dim_x X_y\) 是纤维 \(X_y\) 在 \(x\) 处的局部维数。关键在于：平坦性保证 \(B\) 中正则序列的像在 \(A\) 中仍为正则序列（因为平坦模保持正则序列），从而 \(B\) 中参数系的像构成 \(A\) 中参数系的一部分，且剩余的参数系恰好给出纤维的维数。更精确地，取 \(B\) 的参数系 \(t_1,\ldots,t_d\)（\(d = \dim B\)），由平坦性它们在 \(A\) 中仍为正则序列。补足 \(A\) 的参数系 \(t_1,\ldots,t_d,s_1,\ldots,s_e\)，则 \(e = \dim(A/\mathfrak{m}_B A) = \dim_x X_y\)，且 \(\dim A = d + e\)。\(\blacksquare\)

\subsection{83.7 概形理论中的上同调}\label{ux6982ux5f62ux7406ux8bbaux4e2dux7684ux4e0aux540cux8c03}

\textbf{定义 83.15（层上同调）}：概形 \(X\) 上的层上同调 \(H^i(X, \mathcal{F})\) 是整体截面函子 \(\Gamma(X, -)\) 的右导出函子。对于仿射概形 \(X = \operatorname{Spec} A\) 和拟凝聚层 \(\mathcal{F} = \widetilde{M}\)，\(H^i(X, \mathcal{F}) = 0\)（\(i > 0\)）。

\textbf{定理 83.8（Grothendieck 凝聚性定理）}：设 \(f: X \to Y\) 是 Noether 概形之间的真态射，\(\mathcal{F}\) 是 \(X\) 上的凝聚层。则 \(R^i f_* \mathcal{F}\) 是 \(Y\) 上的凝聚层（对所有 \(i \geq 0\)）。特别地，\(H^i(X, \mathcal{F})\) 是有限生成 \(A\)-模（当 \(Y = \operatorname{Spec} A\) 时）。

\textbf{证明}（框架）：利用 Cech 上同调和射影空间的 Serre 计算。在射影丛 \(\mathbb{P}^n_Y\) 上，直接计算 \(R^i f_* \mathcal{O}(d)\)，然后利用 Chow 引理（任何真态射可由射影态射支配）约化到射影情形。\(\blacksquare\)

\textbf{定理 83.9（Grothendieck 对偶）}：设 \(f: X \to Y\) 是 Noether 概形之间的真态射。则导出函子 \(R f_*: D^+_{\text{qc}}(X) \to D^+_{\text{qc}}(Y)\) 有右伴随 \(f^!\)：对 \(\mathcal{F} \in D^+_{\text{qc}}(X)\) 和 \(\mathcal{G} \in D^+_{\text{qc}}(Y)\)， \[R f_* R \mathcal{H}om_X(\mathcal{F}, f^! \mathcal{G}) \cong R \mathcal{H}om_Y(R f_* \mathcal{F}, \mathcal{G})\] 对偶复形 \(\omega_X = f^! \mathcal{O}_Y\)（\(Y = \operatorname{Spec} k\) 时）是\textbf{对偶层}（dualizing sheaf），在光滑情形下 \(\omega_X = \Omega^n_X[n]\)（典范丛的平移）。Grothendieck 对偶是 Serre 对偶定理的深远推广，是现代代数几何中不可或缺的基础工具。

\textbf{证明}：构造分为三步。第一步：对光滑态射 \(f\)（如 \(\mathbb{A}^n_Y \to Y\)），定义 \(f^! \mathcal{G} = \Omega^n_{X/Y}[n] \otimes f^* \mathcal{G}\)，验证伴随性由迹映射（trace map）\(R f_* \Omega^n_{X/Y}[n] \to \mathcal{O}_Y\) 诱导。第二步：对闭浸入 \(i: Z \hookrightarrow X\)，由局部对偶（local duality）理论，定义 \(i^! \mathcal{G} = R \mathcal{H}om_X(i_* \mathcal{O}_Z, \mathcal{G})\)，在仿射情形下这归结为 Matlis 对偶定理：若 \(A\) 是 Noether 局部环，\(E\) 是剩余域的内射包（injective hull），则 \(R\operatorname{Hom}_A(-, E)\) 给出对偶。第三步：对一般真态射 \(f\)，利用分解 \(f = p \circ i\)（\(i\) 为闭浸入，\(p\) 为光滑态射），定义 \(f^! = i^! \circ p^!\)。由 Grothendieck 的凝聚性定理，\(R f_*\) 将 \(D^+_{\text{qc}}(X)\) 映到 \(D^+_{\text{qc}}(Y)\)，右伴随 \(f^!\) 的存在性由 Brown 可表性定理保证。在 \(Y = \operatorname{Spec} k\) 时，\(\omega_X = f^! k\) 是双重复形，满足 \(R \operatorname{Hom}_X(\mathcal{F}, \omega_X) \cong \operatorname{Hom}_k(R\Gamma(X, \mathcal{F}), k)\)。\(\blacksquare\)

\textbf{定理 83.10（Grothendieck-Riemann-Roch 定理）}：设 \(f: X \to Y\) 是光滑射影簇之间的真态射。则对于 \(X\) 上的凝聚层 \(\mathcal{F}\)，在 Chow 群（或有理等价类）中： \[\operatorname{ch}(f_![\mathcal{F}]) \cdot \operatorname{td}(Y) = f_*(\operatorname{ch}(\mathcal{F}) \cdot \operatorname{td}(X))\] 其中 \(\operatorname{ch}\) 是 Chern 特征标，\(\operatorname{td}\) 是 Todd 类。这是 Hirzebruch-Riemann-Roch 定理的极大推广，将 Euler 示性数 \(\chi(X, \mathcal{F})\) 的计算转化为 Chern 类的拓扑积分。

\textbf{证明}：设 \(K_0(X)\) 为 \(X\) 上凝聚层范畴的 Grothendieck 群。Chern 特征标 \(\operatorname{ch}: K_0(X) \to CH^*(X)_{\mathbb{Q}}\) 是环同态（由张量积到 Chern 类的 Whitney 和公式保证）。对闭嵌入 \(i: Z \hookrightarrow X\)，需验证 GRR 公式对 \(i_*[\mathcal{F}]\) 成立。设 \(\mathcal{I}\) 为 \(Z\) 的理想层，则 \(i_* \mathcal{O}_Z\) 有有限 \(\mathcal{O}_X\)-局部自由分解（由 \(X\) 光滑），由此可计算 \(\operatorname{ch}(i_*[\mathcal{F}]) = i_*(\operatorname{ch}(\mathcal{F}) \cdot \operatorname{td}(N_{Z/X})^{-1})\)（其中 \(N_{Z/X}\) 为法丛）。一般态射 \(f\) 分解为 \(X \hookrightarrow \mathbb{P}^N_Y \to Y\)（闭嵌入后接投影）。对投影 \(\pi: \mathbb{P}^N_Y \to Y\)，\(\pi_*(\operatorname{ch}(\mathcal{O}(d)) \cdot \operatorname{td}(T_{\mathbb{P}^N_Y/Y}))\) 可直接由射影丛的 Chern 类计算验证 GRR 公式。由 \(K_0(X)\) 的生成元性质（任何凝聚层可由局部自由层分解），GRR 对所有 \(\mathcal{F}\) 成立。\(\blacksquare\)

概形理论的深远影响在于：它为代数几何提供了一个统一的语言，使得数论（算术概形）、代数几何（代数簇的概形化）、复几何（解析空间与概形的关系，由 Serre 的 GAGA 原理建立）和拓扑学（etale 同伦类型）得以在同一框架下交流。Grothendieck 的 EGA（Elements de Geometrie Algebrique）和 SGA（Seminaire de Geometrie Algebrique）系列著作系统地建立了这一理论，成为 20 世纪数学最伟大的成就之一。

\subsection{83.8 概形理论中的形式概形与形变}\label{ux6982ux5f62ux7406ux8bbaux4e2dux7684ux5f62ux5f0fux6982ux5f62ux4e0eux5f62ux53d8}

\textbf{定义 83.16（形式概形）}：形式概形是概形理论的完备化版本。设 \(X\) 是 Noether 概形，\(\mathcal{I} \subseteq \mathcal{O}_X\) 是理想层。沿 \(\mathcal{I}\) 的形式完备化 \(\hat{X} = \varinjlim (X_n)\)（其中 \(X_n = \operatorname{Spec}(\mathcal{O}_X / \mathcal{I}^{n+1})\)）是\textbf{形式概形}。形式概形在形变理论、代数化定理和 Langlands 纲领的局部模型中起核心作用。形式概形与代数概形之间的关系由 Grothendieck 的\textbf{代数化定理}（Existence Theorem）给出：如果形式概形满足适当的''代数可化''条件（如相应的形式层是凝聚的且满足一定的 Efektivitat 条件），则它来自某个代数概形。

\textbf{定理 83.11（Grothendieck 形式函数定理）}：设 \(f: X \to Y\) 是 Noether 概形之间的真态射，\(\mathcal{F}\) 是 \(X\) 上的凝聚层。设 \(y \in Y\)，\(\hat{Y}\) 是 \(Y\) 沿 \(\{y\}\) 的完备化。则 \((R^i f_* \mathcal{F})_y^\wedge \cong \varprojlim H^i(X_n, \mathcal{F}|_{X_n})\)（其中 \(X_n = X \times_Y \operatorname{Spec}(\mathcal{O}_{Y,y}/\mathfrak{m}_y^{n+1})\)）。形式函数定理是 Zariski 连通性定理和 Stein 分解定理证明的关键工具。

\textbf{证明}：设 \(A = \mathcal{O}_{Y,y}\)，\(\mathfrak{m}\) 为极大理想，\(A_n = A/\mathfrak{m}^{n+1}\)。令 \(X_n = X \times_Y \operatorname{Spec} A_n\)。由凝聚性定理，\(M = (R^i f_* \mathcal{F})_y\) 是有限 \(A\)-模。其 \(\mathfrak{m}\)-进完备化为 \(M^\wedge = \varprojlim M/\mathfrak{m}^{n+1} M\)。需要证明 \(M/\mathfrak{m}^{n+1} M \cong H^i(X_n, \mathcal{F}|_{X_n})\)。这由基变换和形式 GAGA 原理得出：考虑纤维积图表，由投影公式和导出基变换，\(R\Gamma(X_n, \mathcal{F}|_{X_n}) \cong R\Gamma(X, \mathcal{F}) \otimes_A^{\mathbb{L}} A_n\)。取上同调并利用 \(\mathcal{F}\) 在 \(A\) 上的 \(A\)-平坦化（适当替换后），可得 Cech 复形层次的同构：\(H^i(X_n, \mathcal{F}|_{X_n}) \cong M \otimes_A A_n \cong M/\mathfrak{m}^{n+1} M\)。取投射极限 \(\varprojlim_n\) 即得 \((R^i f_* \mathcal{F})_y^\wedge \cong \varprojlim_n H^i(X_n, \mathcal{F}|_{X_n})\)。\(\blacksquare\)

\textbf{定理 83.12（Schlessinger 形变理论）}：设 \(X_0\) 是代数闭域 \(k\) 上的概形。\(X_0\) 在 Artin 局部环上的无穷小形变由 \(T^1(X_0) = \operatorname{Ext}^1(\Omega_{X_0/k}, \mathcal{O}_{X_0})\) 分类，障碍空间为 \(T^2(X_0) = \operatorname{Ext}^2(\Omega_{X_0/k}, \mathcal{O}_{X_0})\)。概形 \(X_0\) 是\textbf{非障碍的}（unobstructed）如果 \(T^2(X_0) = 0\)，此时形变空间是光滑的。对光滑代数簇，\(T^1 = H^1(X_0, T_{X_0})\)（Kodaira-Spencer 理论），\(T^2 = H^2(X_0, T_{X_0})\)。对 Calabi-Yau 流形，\(T^2 = 0\) 由 Serre 对偶和 \(\Omega^n_X \cong \mathcal{O}_X\) 推出，故其模空间是光滑的（Bogomolov-Tian-Todorov 定理）。

\textbf{证明}（框架）：设 \((A, \mathfrak{m})\) 是 Artin 局部 \(k\)-代数，\(A/\mathfrak{m} \cong k\)。\(X_0\) 到 \(\operatorname{Spec} A\) 的平坦形变由 \(X_0\) 上满足 Jacobi 恒等式的 Maurer-Cartan 元素分类。Schlessinger 的形变函子 \(D: \mathbf{Art}_k \to \mathbf{Set}\) 将 \(A\) 映为 \(X_0\) 在 \(\operatorname{Spec} A\) 上的平坦形变的同构类。Schlessinger 判据给出了 \(D\) 可表（pro-representable）的充要条件。\(T^1(X_0)\) 是 \(D(k[\varepsilon]/(\varepsilon^2))\) 的切空间，\(T^2(X_0)\) 是障碍空间。\(\blacksquare\)

\subsection{83.9 概形理论与模空间}\label{ux6982ux5f62ux7406ux8bbaux4e0eux6a21ux7a7aux95f4}

\textbf{定义 83.17（模空间）}：设 \(\mathcal{M}\) 是概形范畴上的纤维范畴（fibered category），其纤维为待分类的几何对象（如曲线、向量丛、Higgs 丛等）。\(\mathcal{M}\) 的\textbf{粗模空间}（coarse moduli space）是概形 \(M\) 连同自然变换 \(\mathcal{M} \to h_M\)（满足泛性质）。\(\mathcal{M}\) 的\textbf{精模空间}（fine moduli space）是满足 \(\mathcal{M} \cong h_M\) 的概形。大多数模问题由于对象的自同构群而无法拥有精模空间，由此引入了\textbf{代数叠}（algebraic stack）的概念------这是 Deligne 和 Mumford 为克服自同构困难而发展的概形推广。

\textbf{定理 83.13（Mumford 几何不变量理论 / GIT）}：设 \(G\) 是约化代数群作用在射影概形 \(X\) 上。则 \(X\) 的稳定点（stable points）和半稳定点（semistable points）的商 \(X^{\text{st}}/G\) 和 \(X^{\text{ss}}//G\) 作为概形存在。GIT 商是构造模空间的基本工具，应用于曲线模空间 \(\overline{\mathcal{M}}_g\)（Deligne-Mumford 紧化）、向量丛模空间和 Higgs 丛模空间。

\textbf{证明}：设 \(X \subseteq \mathbb{P}^n\) 为 \(G\)-等变闭嵌入（\(G\) 线性地作用在 \(\mathbb{P}^n\) 上）。\(G\) 的约化性保证 \(G\)-不变多项式环 \(\bigoplus_{d \geq 0} H^0(X, \mathcal{O}_X(d))^G\) 是有限生成分次代数。定义 \(X^{\text{ss}} = \{x \in X : \exists d > 0, \exists f \in H^0(X, \mathcal{O}_X(d))^G, f(x) \neq 0\}\)（半稳定点------存在 \(G\)-不变截面在这些点处非零）。定义 \(X^{\text{st}} = \{x \in X^{\text{ss}} : G \cdot x \text{ 在 } X^{\text{ss}} \text{ 中闭}, \operatorname{Stab}_G(x) \text{ 有限}\}\)（稳定点）。GIT 商 \(X//G = \operatorname{Proj}(\bigoplus H^0(X, \mathcal{O}_X(d))^G)\) 是 \(X^{\text{ss}}\) 在等价关系 \(x \sim y \Leftrightarrow \overline{G \cdot x} \cap \overline{G \cdot y} \cap X^{\text{ss}} \neq \emptyset\) 下的商。Hilbert-Mumford 数值判据通过单参数子群 \(\lambda: \mathbb{G}_m \to G\) 的 Hilbert-Mumford 权 \(\mu(x, \lambda)\) 刻画（半）稳定性：\(x\) 半稳定当且仅当 \(\mu(x, \lambda) \geq 0\) 对所有 \(\lambda\) 成立。\(\blacksquare\)

\textbf{定理 83.14（Deligne-Mumford 紧化，1969）}：亏格 \(g \geq 2\) 的光滑代数曲线的模空间 \(\mathcal{M}_g\) 的 \textbf{Deligne-Mumford 紧化} \(\overline{\mathcal{M}}_g\) 添加了稳定曲线（仅允许节点的奇点且自同构群有限）作为边界点。\(\overline{\mathcal{M}}_g\) 是光滑真 Deligne-Mumford 代数叠，维数为 \(3g-3\)。\(\overline{\mathcal{M}}_g\) 的边界 \(\overline{\mathcal{M}}_g \setminus \mathcal{M}_g\) 是正规交叉除子，其不可约分支由对偶图给出。\(\overline{\mathcal{M}}_g\) 上的相交理论和 tautological 环是 Gromov-Witten 理论和 ELSV 公式的核心研究对象。

\textbf{证明}：稳定曲线的定义：连通、亏格 \(g\) 的节点曲线 \(C\)（仅允许节点奇点）是稳定的，如果其自同构群 \(\operatorname{Aut}(C)\) 是有限群。等价地，\(C\) 的每个光滑有理分支（同构于 \(\mathbb{P}^1\)）至少与曲线的其余部分交于 \(3\) 个点。\(\overline{\mathcal{M}}_g\) 的构造利用 GIT：通过 \(m\)-典范嵌入 \(\nu_m: C \hookrightarrow \mathbb{P}^{N}\)（\(m \geq 3\)，\(N = (2m-1)(g-1)-1\)），将稳定曲线视为 \(\mathbb{P}^N\) 中的点（Hilbert 概形中的点），然后取 GIT 商。\(\overline{\mathcal{M}}_g\) 作为代数叠的构造：对每个稳定曲线 \(C\)，其形变空间是光滑的（\(T^1 = \operatorname{Ext}^1(\Omega_C, \mathcal{O}_C)\)，\(T^2 = 0\)），故局部模空间是光滑的。维数 \(3g-3\) 来自 Riemann-Roch：光滑曲线的形变参数由 \(H^1(C, T_C)\) 给出，\(\dim H^1(C, T_C) = 3g-3\)（由 Serre 对偶和 \(\deg T_C = 2-2g\) 及 Riemann-Roch 计算）。边界由具有至少一个节点的稳定曲线构成，每个节点对应一个不可约分支。\(\blacksquare\)

概形理论的核心成就之一是：它使得代数几何可以在任意基概形上工作，从而将算术问题（如 \(\mathbb{Z}\) 上的概形）和几何问题统一处理。Fermat 大定理的证明（Wiles, 1995）本质上依赖于概形理论中的模形式 Galois 表示和模曲线 \(X_0(N)\) 的几何。算术代数几何的现代发展------包括 \(p\)-进 Hodge 理论、perfectoid 空间（Scholze, 2012）和晶状上同调------都是概形语言的自然延伸，进一步深化了概形理论在数学中的基础地位。

\newpage

\chapter{08-代数几何与数论/01-代数几何/03-上同调与算术几何.md}\label{ux4ee3ux6570ux51e0ux4f55ux4e0eux6570ux8bba01-ux4ee3ux6570ux51e0ux4f5503-ux4e0aux540cux8c03ux4e0eux7b97ux672fux51e0ux4f55.md}

\section{第285章：上同调方法}\label{ux7b2c285ux7ae0ux4e0aux540cux8c03ux65b9ux6cd5}

层上同调是代数几何最重要的计算工具，它推广了拓扑空间的上同调（V14），为 Riemann-Roch 定理、Serre 对偶和现代代数几何的核心定理提供了框架。

\subsection{84.1 层上同调的定义}\label{ux5c42ux4e0aux540cux8c03ux7684ux5b9aux4e49}

\textbf{定义 84.1}（层范畴）：拓扑空间 \(X\) 上的 Abel 群层构成 Abel 范畴 \(\mathbf{Sh}(X)\)。它有足够多内射对象（内射层）。

\textbf{定义 84.2}（整体截面函子）：整体截面函子 \(\Gamma(X, -): \mathbf{Sh}(X) \to \mathbf{Ab}\) 定义为 \(\Gamma(X, \mathcal{F}) = \mathcal{F}(X)\)。\(\Gamma(X, -)\) 是左正合函子。

\textbf{定义 84.3}（层上同调）：\(\mathcal{F}\) 的 \(i\) 次\textbf{层上同调群} \(H^i(X, \mathcal{F})\) 是整体截面函子的右导出函子：

\[H^i(X, \mathcal{F}) = R^i\Gamma(X, \mathcal{F})\]

具体计算：取内射分解 \(0 \to \mathcal{F} \to \mathcal{I}^0 \to \mathcal{I}^1 \to \cdots\)，则 \(H^i(X, \mathcal{F}) = H^i(\Gamma(X, \mathcal{I}^\bullet))\)。

\textbf{命题 84.1}（基本性质）： 1. \(H^0(X, \mathcal{F}) = \Gamma(X, \mathcal{F})\) 2. 短正合列 \(0 \to \mathcal{F} \to \mathcal{G} \to \mathcal{H} \to 0\) 诱导长正合序列 3. 若 \(\mathcal{F}\) 是内射层，则 \(H^i(X, \mathcal{F}) = 0\)（\(i > 0\)）

层上同调的定义源于 Grothendieck 的深刻洞见：将拓扑空间的上同调理论从常系数层推广到任意层，从而将代数几何中的''函数''（结构层截面）与''拓扑''（上同调群）融为一体。在拓扑空间中，常层 \(\underline{\mathbb{Z}}\) 的层上同调 \(H^i(X, \underline{\mathbb{Z}})\) 恰为奇异上同调 \(H^i_{\text{sing}}(X, \mathbb{Z})\)------这一事实确立了层上同调作为经典上同调的自然推广。层上同调的核心优势在于：它同时处理了空间 \(X\) 的拓扑信息和层 \(\mathcal{F}\) 的''代数信息''（如截面、局部截面之间的关系），使得代数几何中的许多经典构造（如 Riemann-Roch 定理中的 Euler 示性数 \(\chi(X, \mathcal{F}) = \sum (-1)^i \dim H^i(X, \mathcal{F})\)）获得了统一的上同调框架。

\subsection{84.2 Cech 上同调}\label{cech-ux4e0aux540cux8c03}

\textbf{定义 84.4}（Cech 复形）：设 \(\mathfrak{U} = \{U_i\}_{i \in I}\) 是 \(X\) 的开覆盖。对层 \(\mathcal{F}\)，定义 Cech 复形：

\[\check{C}^p(\mathfrak{U}, \mathcal{F}) = \prod_{i_0 < \cdots < i_p} \mathcal{F}(U_{i_0} \cap \cdots \cap U_{i_p})\]

微分 \(\delta: \check{C}^p \to \check{C}^{p+1}\) 由 \((\delta s)_{i_0, \ldots, i_{p+1}} = \sum_{k=0}^{p+1} (-1)^k s_{i_0, \ldots, \widehat{i_k}, \ldots, i_{p+1}}|_{U_{i_0} \cap \cdots \cap U_{i_{p+1}}}\) 给出。

\textbf{定义 84.5}（Cech 上同调）：\(\check{H}^p(\mathfrak{U}, \mathcal{F}) = H^p(\check{C}^\bullet(\mathfrak{U}, \mathcal{F}))\)。取所有开覆盖的归纳极限得 \textbf{Cech 上同调} \(\check{H}^p(X, \mathcal{F})\)。

\textbf{定理 84.1}（Leray 定理）：若 \(\mathfrak{U}\) 是 \(X\) 的开覆盖，且对每个 \(U_{i_0} \cap \cdots \cap U_{i_p}\) 和每个 \(q > 0\)，\(H^q(U_{i_0} \cap \cdots \cap U_{i_p}, \mathcal{F}) = 0\)，则 \(\check{H}^p(\mathfrak{U}, \mathcal{F}) \cong H^p(X, \mathcal{F})\)。

特别地，对仿射概形上的拟凝聚层和仿射开覆盖，Cech 上同调等于层上同调。

\textbf{证明}：构造双复形 \(K^{p,q} = \prod_{i_0 < \cdots < i_p} \mathcal{I}^q(U_{i_0} \cap \cdots \cap U_{i_p})\)，其中 \(\mathcal{F} \to \mathcal{I}^\bullet\) 是内射分解。该双复形的两个谱序列分别为：\(E_2^{p,q} = \check{H}^p(\mathfrak{U}, \underline{H}^q(\mathcal{F}))\)（预层上同调的 Cech 上同调）和 \(E_2^{p,q} = H^p(X, \mathcal{F})\)（\(p=0\) 时由 Cech 复形的正则性）。消没条件确保 \(\underline{H}^q(\mathcal{F}) = 0\)（\(q > 0\)），故第一个谱序列退化，给出 \(\check{H}^p(\mathfrak{U}, \mathcal{F}) \cong H^p(X, \mathcal{F})\)。\(\blacksquare\)

Cech 上同调是计算层上同调最实用的工具。其核心思想是将层上同调这一抽象的同调代数构造转化为组合数据：给定开覆盖 \(\mathfrak{U}\)，Cech 复形完全由层在有限交上的截面和限制映射决定。在代数几何中，由于仿射概形上的拟凝聚层具有消没性质（Serre 定理），仿射开覆盖的 Cech 上同调自动等于层上同调，这使得 Cech 上同调成为计算射影空间和其他具体概形上同调群的标准方法。然而，Cech 上同调与层上同调并非总是同构------对于一般的拓扑空间和层，Cech 上同调仅与层上同调在低次（\(p = 0, 1\)）处自然同构，在高次处存在差异。这一差异由 Cartan-Leray 谱序列精确描述，其 \(E_2\) 项为 \(\check{H}^p(\mathfrak{U}, \underline{H}^q(\mathcal{F}))\)，收敛到 \(H^{p+q}(X, \mathcal{F})\)。

\subsection{84.3 de Rham 上同调与 Hodge 理论}\label{de-rham-ux4e0aux540cux8c03ux4e0e-hodge-ux7406ux8bba}

\textbf{定义 84.6}（de Rham 复形）：设 \(X\) 为 \(n\) 维光滑流形。\(\Omega_X^p\) 为 \(X\) 上的光滑 \(p\)-形式层。外微分 \(d: \Omega_X^p \to \Omega_X^{p+1}\) 满足 \(d^2 = 0\)，形成 \textbf{de Rham 复形}：

\[0 \to \underline{\mathbb{R}} \to \Omega_X^0 \xrightarrow{d} \Omega_X^1 \xrightarrow{d} \cdots \xrightarrow{d} \Omega_X^n \to 0\]

\textbf{定理 84.2}（de Rham 定理）：\(H^i_{\text{dR}}(X) := H^i(\Gamma(X, \Omega_X^\bullet)) \cong H^i_{\text{sing}}(X, \mathbb{R})\)。即 de Rham 上同调（由微分形式的外微分定义）同构于奇异上同调（由单形的边界算子定义）。这是微分拓扑与代数拓扑之间最深刻的桥梁之一。

\textbf{证明（框架）}：de Rham 复形是常层 \(\underline{\mathbb{R}}\) 的零调分解（Poincare 引理保证 \(\Omega_X^p\) 是零调层）。由层的零调分解计算上同调的一般原理，\(H^i(\Gamma(X, \Omega_X^\bullet)) \cong H^i(X, \underline{\mathbb{R}})\)。而 \(H^i(X, \underline{\mathbb{R}})\) 同构于奇异上同调（由奇异上同调与层上同调的同构定理）。\(\blacksquare\)

\textbf{定义 84.7}（Hodge 理论）：设 \(X\) 为紧 Kahler 流形（或紧 Riemann 面）。在 \(X\) 上选取 Hermite 度量，定义 \textbf{Hodge 星算子} \(\star: \Omega_X^p \to \Omega_X^{n-p}\) 和 \textbf{Laplace 算子} \(\Delta = d d^* + d^* d\)（其中 \(d^* = -\star d \star\) 为 \(d\) 的形式伴随）。\textbf{调和形式}即满足 \(\Delta \omega = 0\) 的形式。

\textbf{定理 84.3}（Hodge 分解定理）：对紧 Kahler 流形 \(X\)，存在正交直和分解

\[\Omega^p(X) = \mathcal{H}^p(X) \oplus \operatorname{im} d \oplus \operatorname{im} d^*\]

其中 \(\mathcal{H}^p(X)\) 为调和 \(p\)-形式空间。进而 \(H^p_{\text{dR}}(X) \cong \mathcal{H}^p(X)\)------每个 de Rham 上同调类中有唯一的调和代表元。对于复流形，还有 \textbf{Dolbeault 上同调} 的对应分解：

\[H^k(X, \mathbb{C}) \cong \bigoplus_{p+q=k} H^{p,q}(X)\]

其中 \(H^{p,q}(X) = H^q(X, \Omega_X^p)\) 为 Dolbeault 上同调群。此分解满足 \textbf{Hodge 对称性}：\(H^{p,q} \cong \overline{H^{q,p}}\)。

\textbf{证明（框架）}：由椭圆算子理论，\(\Delta\) 是椭圆自伴算子，其核 \(\mathcal{H}^p\) 有限维。由椭圆正则性，\(\Delta \omega = \alpha\) 的解 \(\omega\) 具有与 \(\alpha\) 相同的正则性。Hodge 分解由正交投影 \(H: \Omega^p \to \mathcal{H}^p\) 诱导。Dolbeault 上同调 \(H^{p,q}\) 对应于 \(\bar{\partial}\)-Laplace 算子 \(\Delta_{\bar{\partial}}\) 的调和形式空间。在 Kahler 流形上，Kahler 恒等式 \(\Delta_d = 2\Delta_{\bar{\partial}}\) 成立，从而 de Rham 调和形式与 Dolbeault 调和形式一致，给出 Hodge 分解。\(\blacksquare\)

Hodge 理论在代数几何中具有深远影响。Hodge 分解 \(H^k(X, \mathbb{C}) = \bigoplus H^{p,q}\) 定义了 Hodge 结构------一个带有 Hodge 滤链的整系数向量空间。Griffiths 的周期映射理论将 Hodge 结构的变分与代数簇的模空间联系起来，导致了代数几何中许多深刻结果，如 Hodge 猜想（预测 \(H^{2p}(X, \mathbb{Q}) \cap H^{p,p}(X)\) 由代数闭链的有理上同调类生成）。Hodge 理论也是证明 Weil 猜想中 Riemann 假设（在部分情形）的关键工具，因为代数簇的 Hodge 指标定理给出了交形式的符号差约束。

\subsection{84.4 凝聚层与 Serre 定理}\label{ux51ddux805aux5c42ux4e0e-serre-ux5b9aux7406}

\textbf{定义 84.8}（\(\mathcal{O}_X\)-模层）：环层空间 \((X, \mathcal{O}_X)\) 上的 \(\mathcal{O}_X\)\textbf{-模层} \(\mathcal{F}\) 是 Abel 群层，且对每个开集 \(U\)，\(\mathcal{F}(U)\) 是 \(\mathcal{O}_X(U)\)-模，限制映射与模结构兼容。

\textbf{定义 84.9}（拟凝聚层与凝聚层）：\(\mathcal{O}_X\)-模层 \(\mathcal{F}\) 称为\textbf{拟凝聚的}，如果局部地有表示 \(\mathcal{O}_X^{\oplus J}|_U \to \mathcal{O}_X^{\oplus I}|_U \to \mathcal{F}|_U \to 0\)。若 \(I\) 和 \(J\) 可取有限，则 \(\mathcal{F}\) 称为\textbf{凝聚层}。

\textbf{定义 84.10}（Serre 捻层）：在射影空间 \(\mathbb{P}^n_A\) 上，\textbf{Serre 捻层} \(\mathcal{O}(d)\) 是分次 \(A[x_0, \ldots, x_n]\)-模 \(A[x_0, \ldots, x_n](d)\)（将次数平移 \(d\)）对应的层。对射影概形 \(X \subseteq \mathbb{P}^n_A\)，\(\mathcal{O}_X(d) = \mathcal{O}(d)|_X\)。

拟凝聚层和凝聚层的概念统一了代数几何中所有''有限型''层对象的讨论框架，而 Serre 捻层则为此提供了一个计算上不可或缺的工具。捻层 \(\mathcal{O}(d)\) 在射影空间上的作用是极其具体的------它的整体截面恰好是 \(d\) 次齐次多项式------因此可以通过 Cech 上同调显式计算其上同调群。Serre 在 1955 年的著名论文 FAC 中利用这些工具，建立了射影概形上凝聚层上同调的三大基本性质：有限性、Serre 消没定理和由整体截面生成的充分性。

\textbf{定理 84.4}（Serre 定理）：设 \(X\) 是 Noether 环 \(A\) 上的射影概形，\(\mathcal{F}\) 是 \(X\) 上的凝聚层。则 1. 对每个 \(i \geq 0\)，\(H^i(X, \mathcal{F})\) 是有限 \(A\)-模 2. 存在 \(d_0\) 使得对 \(d \geq d_0\)，\(H^i(X, \mathcal{F}(d)) = 0\)（\(i > 0\)） 3. \(\mathcal{F}(d)\) 由整体截面生成（\(d \gg 0\)）

其中 \(\mathcal{F}(d) = \mathcal{F} \otimes_{\mathcal{O}_X} \mathcal{O}_X(d)\)。

\textbf{证明（框架）}：将 \(X\) 嵌入射影空间 \(\mathbb{P}^n_A\)，将 \(\mathcal{F}\) 延拓为零到 \(\mathbb{P}^n_A\) 上。对 \(\mathbb{P}^n_A\)，其上同调群可由 Cech 上同调直接计算：取标准仿射开覆盖 \(D_+(x_i)\)，Cech 复形 \(\check{C}^\bullet\) 的每个项是分次 \(A[x_0, \ldots, x_n]\)-模的局部化，其上同调在截断次数 \(d_0\) 以上消没。有限性由 Cech 复形各项的有限生成性导出。\(\blacksquare\)

\textbf{定理 84.5}（凝聚层的有限性，Grothendieck）：若 \(X\) 是域 \(k\) 上的真概形（proper scheme），\(\mathcal{F}\) 是凝聚层，则 \(H^i(X, \mathcal{F})\) 是有限维 \(k\)-向量空间。

\textbf{证明}：将 \(X\) 嵌入射影空间 \(\mathbb{P}^n_k\)（周定理），令 \(i: X \hookrightarrow \mathbb{P}^n_k\) 为闭浸入。则 \(H^i(X, \mathcal{F}) \cong H^i(\mathbb{P}^n_k, i_*\mathcal{F})\)（\(i_*\) 正合且保持凝聚性）。由 Serre 定理 84.4，射影空间上凝聚层的上同调是有限维的，故结论成立。核心在于：真态射下凝聚层的直像层是凝聚的（Grothendieck 凝聚性定理），而 \(\operatorname{Spec} k\) 上的凝聚层恰为有限维 \(k\)-向量空间。\(\blacksquare\)

\subsection{84.5 Serre 对偶}\label{serre-ux5bf9ux5076}

\textbf{定理 84.6}（Serre 对偶，射影空间）：设 \(X = \mathbb{P}^n_k\)（\(k\) 是域）。则对任意凝聚层 \(\mathcal{F}\)，存在自然的完美配对

\[H^i(X, \mathcal{F}) \times \operatorname{Ext}^{n-i}(\mathcal{F}, \omega_X) \to H^n(X, \omega_X) \cong k\]

其中 \(\omega_X = \mathcal{O}_X(-n-1)\) 是 \(\mathbb{P}^n\) 的\textbf{对偶化层}（canonical sheaf）。

\textbf{证明}：对 \(\mathbb{P}^n_k\)，其对偶化层 \(\omega_X = \mathcal{O}_X(-n-1)\) 满足 \(H^n(X, \omega_X) \cong k\)（由 Cech 上同调直接计算：取标准仿射开覆盖 \(\{D_+(x_i)\}\)，Cech \(n\)-上闭链 \(\frac{1}{x_0 \cdots x_n}\) 生成 \(H^n\)）。对任意凝聚层 \(\mathcal{F}\)，构造配对如下：Yoneda 配对 \(\operatorname{Ext}^{n-i}(\mathcal{F}, \omega_X) \times H^i(X, \mathcal{F}) \to H^n(X, \omega_X)\) 由导出范畴中的复合 \(R\operatorname{Hom}(\mathcal{F}, \omega_X) \times R\Gamma(\mathcal{F}) \to R\Gamma(\omega_X)\) 诱导。需验证这是完美配对。对 \(\mathcal{F} = \mathcal{O}_X(d)\)，利用 Cech 上同调显式计算：\(H^i(\mathbb{P}^n, \mathcal{O}(d))\) 和 \(\operatorname{Ext}^{n-i}(\mathcal{O}(d), \omega_X) \cong H^{n-i}(\mathbb{P}^n, \mathcal{O}(-d-n-1))\) 互为对偶（Kronecker 对偶）。一般 \(\mathcal{F}\) 由 Serre 定理可由 \(\mathcal{O}_X(d)\) 的直和分解，通过导出范畴中的归纳法（distinguished triangles）推广到任意凝聚层。完美性由有限维向量空间上非退化双线性型的性质保证。\(\blacksquare\)

\textbf{推论 84.11}（Riemann-Roch 的应用）： 1. \(g = 0\) 时，\(C \cong \mathbb{P}^1\)（有理曲线） 2. \(g = 1\) 时，\(C\) 是椭圆曲线，\(\ell(K_C) = 1\)，\(\deg K_C = 0\) 3. \(g \geq 2\) 时，典范映射 \(C \to \mathbb{P}^{g-1}\) 是 \(C\) 到其典范模型的态射

\textbf{定理 84.7}（Grothendieck-Riemann-Roch）：设 \(f: X \to Y\) 是光滑射影簇之间的真态射。则对 \(X\) 上的凝聚层 \(\mathcal{F}\)，

\[\operatorname{ch}(f_!\mathcal{F}) \cdot \operatorname{td}(T_Y) = f_*(\operatorname{ch}(\mathcal{F}) \cdot \operatorname{td}(T_X))\]

在 \(Y\) 的 Chow 环（或有理上同调）中，其中 \(f_! = \sum (-1)^i R^i f_* \mathcal{F}\)。

\textbf{证明}（框架）：将 Grothendieck 群 \(K_0(X)\) 视为层范畴的 Grothendieck 群。Chern 特征标 \(\operatorname{ch}: K_0(X) \to CH^*(X)_{\mathbb{Q}}\) 是环同态。对闭嵌入 \(i: Z \hookrightarrow X\)，GRR 约化为 Hirzebruch-Riemann-Roch 公式 \(\chi(Z, \mathcal{F}) = \int_Z \operatorname{ch}(\mathcal{F}) \cdot \operatorname{td}(T_Z)\)。一般态射 \(f\) 分解为闭嵌入后接投影：\(X \hookrightarrow \mathbb{P}^N_Y \to Y\)。对投影 \(\pi: \mathbb{P}^N_Y \to Y\)，直接计算 \(\pi_*(\operatorname{ch}(\mathcal{O}(n)) \cdot \operatorname{td}(T_{\mathbb{P}^N_Y/Y}))\) 验证 GRR。由 Grothendieck 群的生成元性质，GRR 对所有 \(\mathcal{F}\) 成立。\(\blacksquare\)

\subsection{84.6 Etale 上同调与 Weil 猜想}\label{etale-ux4e0aux540cux8c03ux4e0e-weil-ux731cux60f3}

\textbf{定义 84.11}（etale 拓扑）：概形 \(X\) 上的 \textbf{etale 拓扑}（etale site）\(X_{\text{et}}\) 以 \(X\) 上的 etale 态射 \(U \to X\) 为''开集''，以 etale 态射族 \(\{U_i \to U\}\)（其像的并覆盖 \(U\)）为''覆盖''。etale 拓扑是 Grothendieck 拓扑的一个实例，比 Zariski 拓扑精细得多：etale 覆盖的''开集''可以是非 Zariski 开的，但它们在代数意义下是''局部同构''的。

\textbf{定义 84.12}（etale 上同调）：\(X_{\text{et}}\) 上的 Abel 层范畴 \(\mathbf{Sh}(X_{\text{et}})\) 有足够多内射对象。常层 \(\mathbb{Z}/n\mathbb{Z}\) 的 etale 上同调群 \(H^i_{\text{et}}(X, \mathbb{Z}/n\mathbb{Z})\) 定义为整体截面函子在 \(X_{\text{et}}\) 上的右导出函子。对素数 \(\ell\)（\(\ell \neq \operatorname{char} k\)），定义 \textbf{\(\ell\)-进上同调}：

\[H^i_{\text{et}}(X, \mathbb{Z}_\ell) = \varprojlim_n H^i_{\text{et}}(X, \mathbb{Z}/\ell^n\mathbb{Z})\]

\[H^i_{\text{et}}(X, \mathbb{Q}_\ell) = H^i_{\text{et}}(X, \mathbb{Z}_\ell) \otimes_{\mathbb{Z}_\ell} \mathbb{Q}_\ell\]

\textbf{定理 84.8}（etale 上同调的基本性质）： 1. 对光滑代数簇 \(X/\mathbb{C}\)，\(H^i_{\text{et}}(X, \mathbb{Z}/n\mathbb{Z}) \cong H^i_{\text{sing}}(X(\mathbb{C}), \mathbb{Z}/n\mathbb{Z})\)（比较定理） 2. 对光滑真代数簇 \(X/k\)（\(k\) 为域），\(H^i_{\text{et}}(X, \mathbb{Q}_\ell)\) 是有限维 \(\mathbb{Q}_\ell\)-向量空间，且 \(i > 2\dim X\) 时消没 3. Poincare 对偶：\(H^i_{\text{et}}(X, \mathbb{Q}_\ell) \times H^{2n-i}_{\text{et}}(X, \mathbb{Q}_\ell(n)) \to \mathbb{Q}_\ell\) 是完美配对 4. Lefschetz 迹公式：对 \(X\) 在 \(\mathbb{F}_q\) 上的 Frobenius 自同态 \(F\)，\(\#X(\mathbb{F}_q) = \sum_{i=0}^{2n} (-1)^i \operatorname{Tr}(F^* | H^i_{\text{et}}(X, \mathbb{Q}_\ell))\)

\textbf{证明}：(1) 比较定理由 Serre 的 GAGA 原理和 Riemann 存在定理：对 \(X/\mathbb{C}\)，\(X(\mathbb{C})\) 的有限 etale 覆盖（在超越拓扑下）与 \(X\) 的有限 etale 覆盖（在 Zariski 拓扑下）一一对应，从而 \(H^1_{\text{et}}\) 与 \(H^1_{\text{sing}}\) 一致。一般 \(H^i\) 的比较由 Artin 比较定理保证：复解析空间 \(X^{\text{an}}\) 的常层上同调与 etale 上同调同构。(2) 有限维性由 Grothendieck 凝聚性定理的 etale 类比：\(R f_* \mathbb{Q}_\ell\) 是构造性层（constructible sheaf），其茎是有限维 \(\mathbb{Q}_\ell\)-向量空间。\(i > 2\dim X\) 时消没由 Artin 的消没定理（etale 上同调的维数上界）和纤维的 etale 上同调维数估计得出。(3) Poincare 对偶由 Verdier 对偶在 etale 设定下的实现：\(R f_! \mathbb{Q}_\ell \cong R f_* \mathbb{Q}_\ell(n)[2n]\)（真光滑态射 \(f: X \to \operatorname{Spec} k\)），取超上同调即得。(4) Lefschetz 迹公式由 Grothendieck 迹公式（Grothendieck-Lefschetz 不动点公式）：对 \(\mathbb{F}_q\) 上的概形 \(X\) 和 Frobenius 自同态 \(F\)，\(F\) 的不动点集 \(X(\mathbb{F}_q)\) 的个数等于 \(\sum (-1)^i \operatorname{Tr}(F^* | H^i_c(X, \mathbb{Q}_\ell))\)。对真 \(X\)，\(H^i_c = H^i\)。\(\blacksquare\)

\textbf{定理 84.9}（Weil 猜想，Deligne 1974）：设 \(X\) 是 \(\mathbb{F}_q\) 上的 \(n\) 维光滑射影代数簇。定义 zeta 函数

\[Z(X, T) = \exp\left(\sum_{m=1}^\infty \#X(\mathbb{F}_{q^m}) \frac{T^m}{m}\right)\]

则 (1) 有理性：\(Z(X, T) \in \mathbb{Q}(T)\)；(2) 函数方程：\(Z(X, 1/(q^n T)) = \pm q^{n\chi/2} T^\chi Z(X, T)\)；(3) Riemann 假设：\(Z(X, T)\) 的零点和极点具有形式 \(q^{-i/2}\)（\(i\) 为整数，\(0 \leq i \leq 2n\)），即 \(\alpha\) 为 \(Z(X, T)\) 的极/零点时 \(|\alpha| = q^{-i/2}\)。

\textbf{证明}（框架）：核心工具是 etale 上同调。由 Lefschetz 迹公式和 Grothendieck 的 \(L\)-函数理论，\(Z(X, T) = \prod_{i=0}^{2n} P_i(T)^{(-1)^{i+1}}\)，其中 \(P_i(T) = \det(1 - F^* T | H^i_{\text{et}}(X, \mathbb{Q}_\ell))\)。有理性由 \(H^i\) 的有限维性得证。函数方程来自 Poincare 对偶。Riemann 假设（最困难的部分）由 Deligne 在 1974 年证明，核心工具是 Lefschetz 束理论、monodromy 群的表示论和 Rankin-Selberg 方法，最终归结为证明 Frobenius 特征值在 \(H^i\) 上的绝对值是 \(q^{i/2}\)。\(\blacksquare\)

\subsection{84.7 晶体上同调与 Grothendieck 的六个函子}\label{ux6676ux4f53ux4e0aux540cux8c03ux4e0e-grothendieck-ux7684ux516dux4e2aux51fdux5b50}

\textbf{定义 84.13}（晶体上同调，Grothendieck-Berthelot）：设 \(X\) 是特征 \(p > 0\) 的完全域 \(k\) 上的概形。\(X\) 的 \textbf{晶体上同调} \(H^i_{\text{cris}}(X/W(k))\) 是 \(X\) 相对于 \(W(k)\)（\(k\) 的 Witt 向量环）的晶体景（crystalline site）上的结构层的上同调。晶体上同调是特征 \(p\) 情形下 de Rham 上同调的''正确''提升：对光滑 \(X/k\)，\(H^i_{\text{cris}}(X/W(k)) \otimes_W k \cong H^i_{\text{dR}}(X/k)\)。

晶体上同调的关键优势在于它提供了 \(p\)-进上同调理论，与 \(\ell\)-进上同调（\(\ell \neq p\)）互补。\textbf{Berthelot-Ogus 比较定理}断言：对光滑真 \(X/k\)，所有上同调理论（etale \(\mathbb{Q}_\ell\)（\(\ell \neq p\)）、晶体上同调、de Rham 上同调）通过适当的系数扩张后给出相同的 Betti 数。晶体上同调在 \(p\)-进 Hodge 理论中扮演核心角色，Fontaine 的周期环 \(B_{\text{cris}}\) 和 \(B_{\text{dR}}\) 将晶体上同调与 Galois 表示联系起来。

\textbf{定义 84.14}（Grothendieck 的六个函子）：在概形态射 \(f: X \to Y\) 的诱导下，层和上同调之间存在六个基本函子，构成 Grothendieck 的''六函子形式化''：

\begin{enumerate}
\def\labelenumi{\arabic{enumi}.}
\tightlist
\item
  \textbf{直像} \(f_*: \mathbf{Sh}(X) \to \mathbf{Sh}(Y)\)------右伴随于 \(f^{-1}\)
\item
  \textbf{逆像} \(f^{-1}: \mathbf{Sh}(Y) \to \mathbf{Sh}(X)\)------左伴随于 \(f_*\)
\item
  \textbf{紧支直像} \(f_!: \mathbf{Sh}(X) \to \mathbf{Sh}(Y)\)------``紧支集截面''的推广
\item
  \textbf{例外逆像} \(f^!: D^+(Y) \to D^+(X)\)------右伴随于 \(Rf_!\)
\item
  \textbf{张量积} \(\otimes: \mathbf{Sh}(X) \times \mathbf{Sh}(X) \to \mathbf{Sh}(X)\)
\item
  \textbf{内部 Hom} \(\mathcal{H}om: \mathbf{Sh}(X)^{\text{op}} \times \mathbf{Sh}(X) \to \mathbf{Sh}(X)\)
\end{enumerate}

这六个函子满足一系列相容性关系（如基变换、投影公式、Verdier 对偶），在导出范畴 \(D^b_c(X, \mathbb{Q}_\ell)\) 中形成完整的理论体系。六个函子形式化是 Grothendieck 对代数几何和拓扑学最深远的贡献之一，它将 Poincare 对偶、Lefschetz 迹公式、层上同调的计算统一在一个精巧的框架内。在 etale 上同调中，紧支上同调 \(H^i_c(X, \mathbb{Q}_\ell) = H^i(Y, Rf_! \mathbb{Q}_\ell)\)（\(f: X \to \operatorname{Spec} k\)）与通常上同调通过 Poincare 对偶相联系：\(H^i_c(X, \mathbb{Q}_\ell) \cong H^{2n-i}(X, \mathbb{Q}_\ell(n))^\vee\)。六个函子在几何 Langlands 纲领、Motivic 上同调和导出代数几何中持续发挥着核心作用。

\hrulefill

\hrulefill

\hrulefill

\hrulefill

\hrulefill

\newpage

\section{第286章：Arakelov 几何}\label{ux7b2c286ux7ae0arakelov-ux51e0ux4f55}

Arakelov 几何将代数几何的概型方法推广到算术情形，在 \(\operatorname{Spec} \mathbb{Z}\) 上构造''算术曲面''，并在每个无穷远点处附加 Hermite 度量以赋予''紧性''。这是将代数几何方法与解析工具相结合的交叉领域，其核心应用包括 Faltings 对 Mordell 猜想的证明。

\subsection{88.1 算术曲面与 Hermite 度量}\label{ux7b97ux672fux66f2ux9762ux4e0e-hermite-ux5ea6ux91cf}

\textbf{定义 88.1}（算术曲面）：\textbf{算术曲面}是一个平坦、有限型的态射 \(\pi: \mathcal{X} \to \operatorname{Spec} \mathbb{Z}\)，其一般纤维 \(\mathcal{X}_{\mathbb{Q}} = \mathcal{X} \times_{\operatorname{Spec} \mathbb{Z}} \operatorname{Spec} \mathbb{Q}\) 是一条光滑射影曲线。实数纤维 \(\mathcal{X}(\mathbb{C}) = \mathcal{X} \times_{\operatorname{Spec} \mathbb{Z}} \operatorname{Spec} \mathbb{C}\) 是紧 Riemann 面（不一定连通）。在每个连通分支 \(\mathcal{X}_\sigma\)（\(\sigma \in \mathcal{X}(\mathbb{C})\) 的连通分支）上需装备 \textbf{Kähler 度量}------称为 Hermite 度量或 Kähler 形式 \(\omega_\sigma\)。

在每个非无穷远点 \(p\)（即正素数），\(\mathcal{X}\) 的纤维 \(\mathcal{X}_p = \mathcal{X} \times_{\operatorname{Spec} \mathbb{Z}} \operatorname{Spec} \mathbb{F}_p\) 是有限域 \(\mathbb{F}_p\) 上的代数曲线（带有可能的奇点）。

\textbf{定义 88.2}（算术除子）：算术曲面 \(\mathcal{X}\) 上的 \textbf{算术除子} 是有限形式和

\[D = \sum_{P} n_P [P] + \sum_{\sigma} \lambda_\sigma [\mathcal{X}_\sigma]\]

其中 \(P\) 遍历 \(\mathcal{X}\) 的不可约水平除子（算术除子的有限部分），\(\lambda_\sigma \in \mathbb{R}\) 为无穷远除子的系数。此类除子形式上类比于紧 Riemann 面上的除子，但额外考虑了无穷远度量。

\subsection{88.2 算术 Chow 群与算术相交理论}\label{ux7b97ux672f-chow-ux7fa4ux4e0eux7b97ux672fux76f8ux4ea4ux7406ux8bba}

\textbf{定义 88.3}（算术 Chow 群）：算术曲面 \(\mathcal{X}\) 的 \textbf{算术 Chow 群} \(\widehat{\operatorname{CH}}^p(\mathcal{X})\) 由以下数据生成：代数圈 \(Z\)（\(\mathcal{X}\) 的子概型）配以沿 \(\mathcal{X}(\mathbb{C})\) 每个连通分支的 \textbf{Green 流} \(g_\sigma\)（满足微分方程 \(\partial \bar{\partial} g_\sigma + \delta_Z = \omega_Z\) 的流），模去有理等价和无穷远处的全纯函数的贡献。

算术 Chow 群 \(\widehat{\operatorname{CH}}^*(\mathcal{X})\) 构成分次环，配以算术相交积

\[\widehat{\operatorname{CH}}^p(\mathcal{X}) \times \widehat{\operatorname{CH}}^q(\mathcal{X}) \to \widehat{\operatorname{CH}}^{p+q}(\mathcal{X}) \otimes \mathbb{R}\]

此积在有限位与通常的相交积一致，在无穷远处与 Green 流的 star 积一致。

\textbf{定义 88.4}（算术相交数）：若 \(\mathcal{X}\) 是算术曲面，\(D, E\) 是两个算术除子，则其\textbf{算术相交数} \(\widehat{D \cdot E} \in \mathbb{R}\) 定义为

\[\widehat{D \cdot E} = D_{\text{fin}} \cdot E_{\text{fin}} + \text{（无穷远处的贡献）}\]

无穷远处的贡献由 Green 函数 \(g(D, E)\) 在 \(\mathcal{X}(\mathbb{C})\) 上的积分给出。

\textbf{定理 88.1}（算术 Hodge 指标定理，Faltings 1984 年）：在算术曲面上，相交配对 \(\widehat{\cdot}\) 满足与 Hodge 指标定理类似的符号性质。

\textbf{证明}（框架）：Faltings 将经典的 Hodge 指标定理（曲面上两个除子相交的符号由其中一个的符号决定）推广到算术情形。在算术曲面 \(\mathcal{X}\) 上，取 \(D\) 为水平（算术）除子，其 \(\mathbb{Q}\)-纤维是通常曲线上的除子。考虑交集矩阵 \((D_i \cdot D_j)\) 在垂直除子空间的限制。由纤维上的经典 Hodge 指标定理，该矩阵的符号为 \((1, n-1)\)（一个正方向）。无穷远处的 Green 函数贡献调整后，断言在 \(\widehat{\operatorname{CH}}^1(\mathcal{X})_{\mathbb{R}}\) 上的相交配对满足同样的符号性质。核心工具是 Faltings 的高度不等式和 Arakelov 的邻接矩阵分析。\(\blacksquare\)

\subsection{88.3 算术 Riemann-Roch 定理}\label{ux7b97ux672f-riemann-roch-ux5b9aux7406}

\textbf{定理 88.2}（算术 Riemann-Roch 定理，Gillet-Soulé 1992 年）：设 \(\pi: \mathcal{X} \to \operatorname{Spec} \mathbb{Z}\) 是算术曲面，\(\overline{E}\) 是 \(\mathcal{X}\) 上的 Hermite 向量丛。则算术 Euler 示性数

\[\widehat{\chi}(\overline{E}) = \sum_{i=0}^2 (-1)^i \widehat{\deg}(\det R^i \pi_* \overline{E})\]

等于由 Gillet-Soulé 的 Chern 形式 \(\widehat{\operatorname{ch}}(\overline{E})\) 和 Todd 形式 \(\widehat{\operatorname{Td}}(T_{\mathcal{X}/\mathbb{Z}})\) 在算术 Chow 环中的相交积分

\[\widehat{\chi}(\overline{E}) = \int_{\mathcal{X}} \widehat{\operatorname{ch}}(\overline{E}) \cdot \widehat{\operatorname{Td}}(T_{\mathcal{X}/\mathbb{Z}})\]

\textbf{证明}：将 Grothendieck-Riemann-Roch 定理提升到算术曲面 \(\pi: \mathcal{X} \to \operatorname{Spec} \mathbb{Z}\)。设 \(\overline{E}\) 为 \(\mathcal{X}\) 上的 Hermite 向量丛。有限素理想处的 Euler 示性数由经典 GRR 给出：\(\chi_{\text{fin}}(\overline{E}) = \int_{\mathcal{X}} \operatorname{ch}(E) \cdot \operatorname{Td}(T_{\mathcal{X}/\mathbb{Z}})\)。无穷远处的贡献来自解析挠率（Ray-Singer）：\(T(\overline{E}) = \frac{1}{2} \sum_{q} (-1)^q q \log \det' \Delta_q\)（\(\Delta_q\) 为 \(\overline{E}\) 值 \(q\)-形式的 Laplace 算子）。Gillet-Soulé 的关键构造是算术 Chern 特征标 \(\widehat{\operatorname{ch}}: \widehat{K}_0(\mathcal{X}) \to \widehat{\operatorname{CH}}^*(\mathcal{X})_{\mathbb{R}}\)，使得在算术 Chow 环中的积分同时捕获有限部分的交数和无穷远处的 Green 流积分：\(\widehat{\chi}(\overline{E}) = \chi_{\text{fin}}(\overline{E}) + T(\overline{E}) = \int_{\mathcal{X}} \widehat{\operatorname{ch}}(\overline{E}) \cdot \widehat{\operatorname{Td}}(T_{\mathcal{X}/\mathbb{Z}})\)。这一定理是 Arakelov 几何的基石。\(\blacksquare\)

算术 Riemann-Roch 定理是 Arakelov 几何中最为深刻和重要的定理之一。它统一了三个看似不相干的领域：有限素理想处的代数交数（经典代数几何），无穷远处的解析不变量（Ray-Singer 解析挠率），以及算术 Chow 环中的形式计算。定理中的解析挠率 \(T(\overline{E})\) 是谱几何中的不变量，由 Ray 和 Singer 于 1971 年引入，其定义为 Laplace 算子 \(\Delta_q\) 的正则化行列式。Bismut 和 Vasserot 后来给出了解析挠率的局部指标定理证明，将 \(T(\overline{E})\) 与局部 Riemann-Roch 定理联系起来。算术 Riemann-Roch 定理的一个关键推论是算术 Hilbert-Samuel 定理（定理 88.3），它给出了截面空间维数的渐近估计。此外，算术 Riemann-Roch 定理在 Faltings 对 Mordell 猜想的证明中发挥了核心作用：Faltings 通过比较算术曲面上的 Arakelov 高度与 Néron-Tate 高度，利用算术 Riemann-Roch 定理推导出高度不等式，进而证明了 Mordell 猜想。在更一般的设定下，Gillet-Soulé 将算术 Riemann-Roch 定理推广到了高维算术簇，建立了算术 \(K\)-理论与算术 Chow 群之间的函子性对应（算术 Riemann-Roch 定理的协变形式）。

\subsection{88.4 算术 Hilbert-Samuel 定理}\label{ux7b97ux672f-hilbert-samuel-ux5b9aux7406}

\textbf{定理 88.3}（算术 Hilbert-Samuel 定理，Zhang 1995 年，Abbes-Bouche 1995 年）：设 \(\mathcal{X}\) 是算术曲面，\(\overline{L}\) 是 \(\mathcal{X}\) 上的 Hermite 线丛（丰沛线丛配以正曲率度量）。则

\[\widehat{\chi}(\overline{L}^{\otimes n}) = \frac{\widehat{\deg}(\overline{L})}{2} n^2 + O(n \log n)\]

其中 \(\widehat{\deg}(\overline{L})\) 是 \(\overline{L}\) 的算术次数。这与通常的 Hilbert-Samuel 定理（\(P(n) = \frac{L^n}{n!} n^n + \cdots\)）在形式上是平行的。

\textbf{证明}：设 \(\overline{L}\) 为 \(\mathcal{X}\) 上的 Hermite 丰沛线丛。对 \(m \gg 0\)，由 Serre 消没定理知 \(H^i(\mathcal{X}, \overline{L}^{\otimes m}) = 0\)（\(i > 0\)），故 \(\widehat{\chi}(\overline{L}^{\otimes m}) = \widehat{\deg}(\det H^0(\mathcal{X}, \overline{L}^{\otimes m}))\)。由算术 Riemann-Roch 定理展开：\(\widehat{\chi}(\overline{L}^{\otimes m}) = \int_{\mathcal{X}} (1 + m \widehat{c}_1(\overline{L}) + \frac{m^2}{2} \widehat{c}_1(\overline{L})^2) \cdot (1 + \frac{1}{2} \widehat{c}_1(T) + \cdots)\)。在 \(\dim \mathcal{X} = 2\) 的情形截断，主导项为 \(\frac{m^2}{2} \widehat{c}_1(\overline{L})^2 = \frac{m^2}{2} \widehat{\deg}(\overline{L})\)。误差项由 Todd 类的低次分量和 Laplace 算子谱的 Weyl 渐近 \(\#\{\lambda_i \leq t\} \sim \frac{\operatorname{vol}}{4\pi} t\) 控制，贡献为 \(O(m \log m)\)。\(\blacksquare\)

算术 Hilbert-Samuel 定理是算术几何中截面空间渐近理论的基石。在经典代数几何中，Hilbert-Samuel 定理断言：对射影簇 \(X\) 上的丰沛线丛 \(L\)，Hilbert 函数 \(h^0(X, L^{\otimes n}) = \dim H^0(X, L^{\otimes n})\) 对充分大的 \(n\) 是一个次数为 \(\dim X\) 的多项式（Hilbert 多项式）。在算术几何中，\(\widehat{\chi}(\overline{L}^{\otimes n})\) 取代了 \(h^0\)，它不仅包含了经典截面的维数信息，还编码了无穷远处 Hermite 度量的贡献。主导系数 \(\widehat{\deg}(\overline{L})\) 是 \(\overline{L}\) 的算术次数，它由有限部分的交数 \(\deg(L_{\mathbb{Q}})\) 和无穷远处的 Green 函数积分共同决定。定理中的误差项 \(O(n \log n)\) 比经典情形中的 \(O(n)\)（对 \(\dim X = 2\)）更弱，这是因为无穷远处的 Laplace 算子谱分布引入了对数因子。Zhang 在 1995 年的论文中利用这一定理建立了算术 Hilbert-Samuel 函数与可计算 Arakelov 高度之间的直接联系，证明了算术 Bogomolov 猜想所需的关键不等式。算术 Hilbert-Samuel 定理的证明依赖于算术 Riemann-Roch 定理和 Ray-Singer 解析挠率在张量幂下的渐近展开，后者由 Bismut-Vasserot 的局部指标定理精确控制。

\subsection{88.5 应用：Mordell 猜想与算术 Bogomolov 猜想}\label{ux5e94ux7528mordell-ux731cux60f3ux4e0eux7b97ux672f-bogomolov-ux731cux60f3}

\textbf{Arakelov 几何的核心应用：}

\begin{enumerate}
\def\labelenumi{\arabic{enumi}.}
\item
  \textbf{Faltings 对 Mordell 猜想的证明}（1983 年）：设 \(C\) 是数域 \(K\) 上亏格 \(g \geq 2\) 的光滑射影曲线，则 \(C(K)\) 是有限集。Faltings 的核心洞察是将 \(C(K)\) 的有界高度与算术曲面上自相交数 \(\omega_{C/Y}^2\) 的 Arakelov 版本关联。利用算术 Hodge 指标定理和高度不等式，证明若有无限多个 \(K\) 点，则 \(\widehat{\operatorname{div}}(P)^2\) 取任意大的值，与算术相交理论中的有界性原理矛盾。
\item
  \textbf{Vojta 猜想}：Vojta 猜想（1987 年）将 Mordell 猜想推广到高维，断言数域上的对数一般型簇上的有理点具有稀疏性质------每个有理点集包含于有限个真子簇的并集中。Vojta 的证明框架将 Arakelov 几何的高度不等式与 Nevanlinna 理论中的第二主定理通过 Diophantine 相应类比结合。
\item
  \textbf{算术 Bogomolov 猜想}：若曲线 \(C\) 嵌入其 Jacobian 簇 \(J(C)\)，且 \(J(C)\) 上的 Néron-Tate 高度 \(\hat{h}\) 在 \(C\) 上有一组严格正的最小累积点（但不全为零），则这些最小高度与某种 Arakelov 自相交数相关。Ullmo（1998 年）和 Zhang（1998 年）独立地利用等分布定理（equidistribution）的算术版本证明了该猜想。
\item
  \textbf{Zhang 的最小高度}：Zhang 提出了 \textbf{本质最小高度}（essential minimum）：对代数簇 \(X\) 上的高度函数，定义 \(\mu_{\text{ess}}(X) = \inf_{U \subset X \text{开稠}} \sup_{P \in U(\overline{K})} h(P)\)。Zhang 的最小高度定理断言：\(\mu_{\text{ess}}(X) \geq 0\) 且等号成立的充要条件由 Galois 轨道的等分布性质刻画。
\item
  \textbf{有效 Mordell 定理}：de Jong-Rémond 和 Binyamini 等人利用 Arakelov 几何建立了有效高度的有界性------给出了 \(C(K)\) 大小的显式上界（虽然当前界很大），这比原始的 Faltings 定理更构造性。
\end{enumerate}

\hrulefill

\hrulefill

\hrulefill

\hrulefill

\hrulefill

\hrulefill

\hrulefill

\subsection{231.A 拓扑斯补充}\label{a-ux62d3ux6251ux65afux8865ux5145}

拓扑斯理论（Topos Theory）由 Grothendieck（1963）、Lawvere（1964）和 Tierney（1972）创立，是连接逻辑、代数和几何的桥梁。一个拓扑斯是一个''广义的集合论宇宙''，在其中可以做几乎所有的数学------同时保留了范畴论的一致性和几何的局部性质。

\subsection{Grothendieck拓扑与景}\label{grothendieckux62d3ux6251ux4e0eux666f}

\textbf{定义}（Grothendieck景 / site）：设 \(\mathcal{C}\) 是范畴。\(\mathcal{C}\) 上的 \textbf{Grothendieck 拓扑} \(J\) 为每个对象 \(U \in \mathcal{C}\) 指定一族筛（sieves）------即''覆盖'' \(U\) 的态射族（满足公理：恒等覆盖、覆盖的拉回存在且可复合）。装备了 Grothendieck 拓扑的范畴 \((\mathcal{C}, J)\) 称为\textbf{景}（site）。

\textbf{定义}（覆盖筛 / covering sieve）：对象 \(U\) 上的筛 \(S\) 是 \(\mathcal{C}\) 中的态射族 \(\{f_i: U_i \to U\}\)（以 \(U\) 为共同目标），在复合下封闭：若 \(f \in S\) 且 \(g\) 与 \(f\) 可复合，则 \(f \circ g \in S\)。\(S\) 是 \(J(U)\) 中的覆盖筛，如果它满足 Grothendieck 拓扑的三条公理。

\textbf{例}： - \textbf{Zariski景}：\(\mathcal{C} = \operatorname{Open}(X)\)（\(X\) 上光滑概形的开嵌入范畴），覆盖为开覆盖 - \textbf{étale景}（概形上）：\(\mathcal{C} = \text{概形}/\text{格}\)，覆盖为满射的 étale 态射族 - \textbf{fppf景}（fidèlement plate de présentation finie）：覆盖为平坦满射的有限表示态射族

\textbf{定义}（层 / sheaf on a site）：\(\mathcal{C}\) 上取值于集合（或群、环等）范畴的预层 \(\mathcal{F}: \mathcal{C}^{\text{op}} \to \mathbf{Set}\) 是 \textbf{\(J\)-层}，如果对每个覆盖筛 \(\{f_i: U_i \to U\} \in J(U)\)，\(\mathcal{F}(U)\) 是 \(\prod_i \mathcal{F}(U_i)\) 的等化子（在 \(\prod_{i,j} \mathcal{F}(U_i \times_U U_j)\) 的两对对偶投影下）------即满足一般的粘合条件。

\subsection{拓扑斯的定义与基本例子}\label{ux62d3ux6251ux65afux7684ux5b9aux4e49ux4e0eux57faux672cux4f8bux5b50}

\textbf{定义}（Grothendieck拓扑斯 / elementary topos）：一个\textbf{Grothendieck 拓扑斯}是等价于景 \((\mathcal{C}, J)\) 上层范畴 \(\operatorname{Sh}(\mathcal{C}, J)\) 的范畴。更一般地，\textbf{初等拓扑斯}（Lawvere-Tierney）是满足以下条件的范畴： 1. 具有所有有限极限 2. 笛卡尔闭（Cartesian closed） 3. 具有子对象分类器 \(\Omega\)（充当''广义真值对象''）

初等拓扑斯是范畴论的''类集合范畴''------可以内在地解释高阶直觉主义逻辑。

\textbf{例}： - \(\mathbf{Set}\)（集合范畴）是拓扑斯（\(\Omega = \{0, 1\}\)，经典二值逻辑） - \(\operatorname{Sh}(X)\)（拓扑空间上的集合层范畴）是拓扑斯 - \(G\text{-}\mathbf{Set}\)（\(G\)-集合范畴，\(G\) 群）是拓扑斯 - \(\mathcal{B}\mathbb{G}_m\)（分类代数 \(\mathbb{Z}[x, y]/(xy - 1)\) 的可换环范畴）不是 Grothendieck 拓扑斯，但作为分类叠是''2-拓扑斯''

\textbf{定理}（Giraud刻画定理）：一个范畴等价于 Grothendieck 拓扑斯当且仅当它满足：具有所有小 colimit、具有一组生成元、小余积是分离的（disjoint）、等价关系是有效的（effective）。

\textbf{证明}（框架）：必要性（每个 Grothendieck 拓扑斯满足这些条件）：\(\operatorname{Sh}(\mathcal{C}, J)\) 是预层范畴 \(\operatorname{PSh}(\mathcal{C})\) 的反射子范畴，故余完备；可表预层 \(\{h_U = \operatorname{Hom}(-, U) : U \in \mathcal{C}\}\) 是生成元族；余积的分离性和等价关系的有效性由层的粘合性质保证。充分性：给定满足四条件的范畴 \(\mathcal{E}\)，取生成元全子范畴 \(\mathcal{C}\) 并赋予典范拓扑（最细的使所有可表态射为覆盖的拓扑），证明 \(\mathcal{E} \cong \operatorname{Sh}(\mathcal{C}, J_{\text{can}})\)：Yoneda 嵌入 \(\mathcal{E} \to \operatorname{PSh}(\mathcal{C})\) 通过有效等价关系将每个对象表示为关系商，从而嵌入像在层范畴中且本质满。\(\blacksquare\)

\subsection{分类拓扑斯与Lawvere-Tierney拓扑}\label{ux5206ux7c7bux62d3ux6251ux65afux4e0elawvere-tierneyux62d3ux6251}

\textbf{定义}（分类拓扑斯）：若几何理论 \(\mathbb{T}\) 的所有模型（在所有 Grothendieck 拓扑斯中）同构于拓扑斯 \(\mathcal{E}_{\mathbb{T}}\) 中一个''万有模型''的拉回，则 \(\mathcal{E}_{\mathbb{T}}\) 是 \(\mathbb{T}\) 的\textbf{分类拓扑斯}。例如： - 群的分类拓扑斯是 \(G\) 的群对象（群作用层范畴） - 环的分类拓扑斯对应交换环范畴中的仿射概形

\textbf{定义}（Lawvere-Tierney拓扑 / 拓扑斯中的局部算子）：拓扑斯 \(\mathcal{E}\) 的子对象分类器 \(\Omega\) 上的 \textbf{Lawvere-Tierney 拓扑}（即层化模态 / local operator）是映射 \(j: \Omega \to \Omega\) 满足 \(j \circ \operatorname{true} = \operatorname{true}\)、\(j \circ j = j\) 和 \(j \circ \wedge = \wedge \circ (j \times j)\)。\(j\) 的代数（\(\Omega\) 中满足 \(j\) 不变性的子对象）构成 \(\mathcal{E}\) 的一个子拓扑斯 \(\operatorname{Sh}_j(\mathcal{E})\)------即 \(\mathcal{E}\) 中 \(j\)-层的全子范畴。这个子拓扑斯对应 \(\mathcal{E}\) 中由 \(j\) 指定的''更精细的层条件''。

\emph{意义}：Lawvere-Tierney 拓扑是 Grothendieck 拓扑在初等拓扑斯中的内在形式------无需引用外部集合论即可定义拓扑斯中的局部化和层概念。

\textbf{定理}（Deligne完备性定理）：每个凝聚 Grothendieck 拓扑斯（coherent topos）具有足够的点------即存在一族拓扑斯态射 \(x_i: \mathbf{Set} \to \mathcal{E}\)（点）使得逆像函子的族 \(\{x_i^*\}\) 是保守的（conservative）。因此，Grothendieck 拓扑斯中的层可由其在''点''上的茎（stalk）完全确定------这推广了拓扑空间中层的茎检验原理。

\textbf{证明}（框架）：凝聚拓扑斯 \(\mathcal{E}\) 由有限表示对象生成的景上的层范畴给出。Deligne 的证明利用逻辑紧致性定理：取 \(\mathcal{E}\) 中对象的有限可满足的层的逆系统，利用紧致性定理（布尔预层的滤过余极限保持层性质）证明存在足够的层态射到 \(\mathbf{Set}\)。关键引理：对凝聚景，覆盖的有限生成性保证每个层的''可满足理论''是有限一致的，由紧致性定理得一致满足性，进而构造出点族使其茎函子保守。\(\blacksquare\)

\subsection{拓扑斯理论的数学影响}\label{ux62d3ux6251ux65afux7406ux8bbaux7684ux6570ux5b66ux5f71ux54cd}

\begin{enumerate}
\def\labelenumi{\arabic{enumi}.}
\item
  \textbf{代数几何}：Grothendieck 拓扑斯是将 fpqc 下降（descent）形式化的范畴框架。平坦下降等价于层条件：一个预层是层当且仅当它满足平坦覆盖的所有下降数据。这是模叠理论和导出代数几何的基础语言。
\item
  \textbf{逻辑基础}：拓扑斯提供了直觉主义类型论的范畴语义------每个拓扑斯 \(\mathcal{E}\) 有内在的逻辑（能解释量词、合取、析取和蕴含），但排中律（\(\varphi \vee \neg\varphi\)）在一般的拓扑斯中不成立。这不同于经典 \(\mathbf{Set}\) 中的逻辑。
\item
  \textbf{同伦类型论（HoTT）}：\(\infty\)-拓扑斯（或 \((\infty, 1)\)-拓扑斯）是 \(\infty\)-范畴版的拓扑斯------它是同伦类型论和 \(\infty\)-群的模型范畴。这将拓扑斯思想推广到了同伦论的领域，是现代数学基础的前沿。
\end{enumerate}

\hrulefill

\emph{本卷建立了范畴论的核心框架：范畴、函子和自然变换提供了描述数学结构的统一语言；Yoneda 引理是范畴论最基本的定理，体现了''对象由关系决定''的哲学；极限和余极限统一了数学中的各种构造；伴随函子是范畴论最核心的概念，自由-遗忘、张量-Hom 等伴随对在数学中无处不在；单子理论将代数结构抽象为范畴论概念；Abel 范畴和导出范畴为同调代数（V12 Ch 118）提供了现代框架。范畴论是连接代数、几何、拓扑和分析的统一语言。}

\hrulefill

\emph{本卷从经典代数几何（仿射簇、射影簇、Hilbert 零点定理）出发，引入现代概形语言和层上同调方法。仿射代数簇与有限生成整代数的范畴等价建立了代数与几何的桥梁；概形理论将这一对应推广到所有交换环，统一了代数几何与数论；Serre 定理和 Serre 对偶为凝聚层上同调的计算提供了核心工具；Riemann-Roch 定理及其推广是代数几何最重要的定理之一，深刻影响了代数曲线、曲面和高维簇的分类理论。为后续数论 II（V16）中椭圆曲线和模形式的学习提供了几何基础。\textbf{补充部分}整合了范畴论基础------代数几何与同调代数的共同语言：范畴、函子与自然变换、Yoneda引理、极限与余极限、伴随函子、单子与代数、Abel范畴与导出范畴初步，为概形理论和层上同调提供了统一的范畴论框架。}

\emph{卷十五：代数几何终。}

\emph{卷十五：代数几何终。}

\emph{卷十五：代数几何终。}

\emph{卷十八：代数几何终。}

\newpage

\chapter{08-代数几何与数论/01-代数几何/04-模空间.md}\label{ux4ee3ux6570ux51e0ux4f55ux4e0eux6570ux8bba01-ux4ee3ux6570ux51e0ux4f5504-ux6a21ux7a7aux95f4.md}

\section{第287章：环面簇与模空间}\label{ux7b2c287ux7ae0ux73afux9762ux7c07ux4e0eux6a21ux7a7aux95f4}

环面簇是代数几何与组合数学的交叉领域，建立了扇（代数几何）与凸多面体（组合学）之间的精确字典。本章从扇的构造出发引入环面簇，并讨论模空间理论和极小模型纲领。

\subsection{85.1 环面簇的构造}\label{ux73afux9762ux7c07ux7684ux6784ux9020}

\textbf{环面簇}（Toric Variety）是包含代数环面\(T \cong (k^\times)^n\)为稠密开子集、且\(T\)的群作用可延拓到全簇的正规代数簇。环面簇由组合几何数据------\textbf{扇}（fan）\(\Sigma\)------完全分类。设\(N \cong \mathbb{Z}^n\)为格，\(M = \operatorname{Hom}(N, \mathbb{Z})\)为对偶格。扇\(\Sigma\)是\(N_{\mathbb{R}} = N \otimes \mathbb{R}\)中一组强凸有理多面锥的集合，满足面封闭性和两两交为面的条件。每个锥\(\sigma \in \Sigma\)对应仿射环面簇\(U_\sigma = \operatorname{Spec}(k[\sigma^\vee \cap M])\)，其中\(\sigma^\vee = \{m \in M_{\mathbb{R}} : \langle m, v \rangle \geq 0,\ \forall v \in \sigma\}\)为对偶锥。沿公共面\(\sigma \cap \tau\)的自然粘合给出了环面簇\(X_\Sigma\)。

环面簇的几何性质一一对应于扇的组合性质：\(X_\Sigma\)光滑当且仅当每个锥\(\sigma\)由基的一部分生成（即\(\det = \pm 1\)）；\(X_\Sigma\)完备当且仅当\(\Sigma\)的支撑为整个\(N_{\mathbb{R}}\)；\(X_\Sigma\)是射影的当且仅当\(\Sigma\)是一个全维凸多面体的正规扇。经典例子： - \textbf{射影空间}\(\mathbb{P}^n\)：对应于以\((e_0, \ldots, e_n)\)（\(e_0 = -\sum_{i=1}^n e_i\)，\(\{e_i\}\)为\(\mathbb{Z}^n\)的标准基）为生成元的扇，其\(\binom{n+1}{1}\)个极大锥由任意\(n\)个生成元张成； - \textbf{加权射影空间}\(\mathbb{P}(a_0, \ldots, a_n)\)：扇的生成元改为满足\(\sum a_i u_i = 0\)的有理向量； - \textbf{Hirzebruch曲面}\(\mathbb{F}_a = \mathbb{P}(\mathcal{O}_{\mathbb{P}^1} \oplus \mathcal{O}_{\mathbb{P}^1}(a))\)（\(a \geq 0\)）：对应的二维扇由四个锥组成，其生成元为\((0, 1), (0, -1), (1, 0), (-1, -a)\)。环面簇建立了代数几何与组合凸多面体之间的精确对应，使奇点解消、上同调计算等问题可转化为纯组合问题。

\subsection{85.2 模空间理论}\label{ux6a21ux7a7aux95f4ux7406ux8bba}

\textbf{模空间理论}将代数簇（或其他几何对象）的同构类参数化，使参数空间本身具有代数簇（或叠）结构。椭圆曲线的模空间是最经典的例子：域\(k\)上的椭圆曲线由\textbf{\(j\)-不变量}\(j(E) = 1728 \cdot 4a^3/(4a^3 + 27b^2)\)完全分类，故仿射直线\(\mathbb{A}^1_j\)是椭圆曲线的粗模空间（不分叉时）。加入尖点（\(j = \infty\)对应节点三次曲线）得到紧化\(\mathbb{P}^1\)。\textbf{曲线模空间}\(M_g\)：亏格\(g \geq 2\)光滑曲线的粗模空间维数为\(3g-3\)（Riemann，1857年），由Riemann-Roch定理通过计算\(H^0(C, \omega_C^{\otimes 2})\)（二阶微分）的自由度得到。Deligne和Mumford于1969年引入\textbf{稳定曲线}概念，构造了紧化模空间\(\overline{M}_g\)，并证明\(M_g\)是不可约拟射影代数簇。在双有理几何中，\textbf{典范环}\(R(K_X) = \bigoplus_{m \geq 0} H^0(X, mK_X)\)的\(\operatorname{Proj}\)给出了簇的典范模型，而Mori程序通过分析收缩映射对模空间的结构进行精细控制，实现了高维簇的分类。

模空间理论的核心思想是通过\textbf{函子观点}（functor of points）来定义模问题：给定一个范畴 \(\mathcal{F}\)，将每个概形 \(S\) 对应到 \(S\) 上参数化的某类对象的同构类集合 \(\mathcal{F}(S)\)。若存在概形（或代数空间、叠）\(M\) 使得 \(\mathcal{F}(S) \cong \operatorname{Hom}(S, M)\)（自然同构），则称 \(M\) 是 \(\mathcal{F}\) 的\textbf{精模空间}（fine moduli space）。然而，由于对象的自同构群可能导致非平凡的同构类，精模空间往往不存在。此时引入\textbf{粗模空间}（coarse moduli space）的概念：存在概形 \(M\) 和自然变换 \(\mathcal{F} \to \operatorname{Hom}(-, M)\)，满足（1）\(M(k)\) 与 \(\mathcal{F}(k)\) 一一对应（对代数闭域 \(k\)），（2）该自然变换关于任意概形的映射是万有的。粗模空间是精模空间在忽略自同构信息后的''最佳逼近''。当代数几何中，\textbf{模叠}（moduli stack）提供了最精细的解决方案：模叠 \(\mathcal{M}\) 是一个范畴纤维化于概形范畴之上，其 \(S\)-点构成一个群胚（而非集合），精确记录了对象的自同构群。例如 \(\overline{\mathcal{M}}_g\) 是亏格 \(g\) 稳定曲线的模叠，而 \(\overline{M}_g\) 是其粗模空间。模叠理论由 Deligne-Mumford 和 Artin 在 1970 年代建立，是当代代数几何中处理模问题不可或缺的语言。

\subsection{85.3 双有理几何与极小模型纲领}\label{ux53ccux6709ux7406ux51e0ux4f55ux4e0eux6781ux5c0fux6a21ux578bux7eb2ux9886}

\textbf{极小模型纲领}（MMP，Minimal Model Program）是双有理几何的核心算法，由Mori在1980年代奠基。对光滑射影簇\(X\)，其典范除子\(K_X\)的正负性是关键分类指标。\textbf{锥定理}（Cone Theorem）断言：\(X\)的曲线锥\(\overline{NE}(X)\)中与\(K_X\)负交的部分由极射线（extremal rays）\(R_i\)生成，每条极射线对应一个收缩映射\(\operatorname{cont}_{R_i}: X \to Y_i\)，有以下三种情形： 1. \textbf{纤维型}：\(\dim Y < \dim X\)，此时\(X\)是Mori纤维空间，终止； 2. \textbf{除子型}：收缩一个除子，\(\dim Y = \dim X\)，继续对\(Y\)做MMP； 3. \textbf{小收缩型}：例外集余维\(\geq 2\)，此时\(K_Y\)不是\(\mathbb{Q}\)-Cartier的，需要\textbf{翻转}（flip）：构造双有理映射\(X \dashrightarrow X^+\)使得\(K_{X^+}\)在翻转曲线上为正。

\textbf{Birkar-Cascini-Hacon-McKernan定理}（2010年）：特征零域上，任意光滑射影簇的典范环是有限生成的，且对数一般型簇的翻转存在。该结果证明了极小模型纲领在极大一般型的核心步骤，为此Birkar获得了2018年Fields奖。MMP的输出是一个极小模型（\(K_X\)是nef的）或Mori纤维空间，实现了双有理分类的最终目标。

\hrulefill

\hrulefill

\hrulefill

\hrulefill

\hrulefill

\hrulefill

\section{第288章：Fano 簇与有理簇}\label{ux7b2c288ux7ae0fano-ux7c07ux4e0eux6709ux7406ux7c07}

Fano 簇和有理簇是双有理几何和正特征代数几何中最重要的对象之一。Fano 簇（反典范除子丰沛的光滑射影簇）的纤维正曲率（某种意义上）为 Mori 理论提供了基本构建单元，而有理簇问题（单有理簇是否必为有理簇）则是代数几何中最深刻的未解问题之一。

\subsection{86.1 Fano 簇的基本性质}\label{fano-ux7c07ux7684ux57faux672cux6027ux8d28}

\textbf{定义 86.1}（Fano 簇）：特征 0 的代数闭域 \(k\) 上的光滑射影簇 \(X\) 称为 \textbf{Fano 簇}，若其反典范除子 \(-K_X\) 是丰沛的（ample）。

等价地，\(X\) 的典范丛 \(\omega_X = \mathcal{O}_X(K_X)\) 的对偶丛 \(\omega_X^{-1}\) 是丰沛的。

\textbf{定义 86.2}（Fano 指标）：Fano 簇 \(X\) 的 \textbf{Fano 指标}（Fano index）定义为最大整数 \(r\) 满足

\[-K_X = r H\]

其中 \(H \in \operatorname{Pic}(X)\) 是丰沛除子类，且 \(H\) 不被任何大于 1 的整数整除。通常可取 \(H\) 为 \(-K_X\) 在 \(\operatorname{Pic}(X)/\operatorname{Tors}\) 中的原像。Fano 指标满足 \(1 \leq r \leq \dim X + 1\)（上界由 Kobayashi-Ochiai 定理保证）。

\textbf{命题 86.1}（Fano 簇的基本几何性质）： 1. Fano 簇是单有理的（uniruled）------即被有理曲线覆盖； 2. \(\chi(X, \mathcal{O}_X) = 1\)，且 \(H^i(X, \mathcal{O}_X) = 0\) 对所有 \(i > 0\) 成立（由 Kodaira 消没定理）； 3. Fano 簇的 Picard 数 \(\rho(X)\) 可以为任意大的值。

\subsection{86.2 二维和三维 Fano 簇的分类}\label{ux4e8cux7ef4ux548cux4e09ux7ef4-fano-ux7c07ux7684ux5206ux7c7b}

\textbf{定理 86.1}（Del Pezzo 曲面分类）：二维 Fano 簇即 \textbf{del Pezzo 曲面}，由次数 \(d = K_X^2\)（\(1 \leq d \leq 9\)）分类：

\begin{longtable}[]{@{}
  >{\raggedright\arraybackslash}p{(\linewidth - 4\tabcolsep) * \real{0.3226}}
  >{\raggedright\arraybackslash}p{(\linewidth - 4\tabcolsep) * \real{0.4839}}
  >{\raggedright\arraybackslash}p{(\linewidth - 4\tabcolsep) * \real{0.1935}}@{}}
\toprule
\begin{minipage}[b]{\linewidth}\raggedright
次数 \(d\)
\end{minipage} & \begin{minipage}[b]{\linewidth}\raggedright
Del Pezzo 曲面
\end{minipage} & \begin{minipage}[b]{\linewidth}\raggedright
描述
\end{minipage} \\
\midrule
\endhead
\bottomrule
\endlastfoot
9 & \(\mathbb{P}^2\) & 射影平面 \\
8 & \(\mathbb{P}^1 \times \mathbb{P}^1\) 或 \(\mathbb{F}_1\) & 二次曲面或一次爆破 \(\mathbb{P}^2\) \\
7 & \(\mathbb{P}^2\) 上 2 点爆破 & \\
6 & \(\mathbb{P}^2\) 上 3 点爆破 & 三次曲面 \\
5 & \(\mathbb{P}^2\) 上 4 点爆破 & \\
4 & \(\mathbb{P}^2\) 上 5 点爆破 & \(\mathbb{P}^4\) 中两个二次超曲面的交 \\
3 & \(\mathbb{P}^2\) 上 6 点爆破 & 三次曲面（\(\mathbb{P}^3\) 中） \\
2 & \(\mathbb{P}^2\) 上 7 点爆破 & 四次曲面（\(\mathbb{P}(1, 1, 1, 2)\)） \\
1 & \(\mathbb{P}^2\) 上 8 点爆破 & 六次曲面（\(\mathbb{P}(1, 1, 2, 3)\)） \\
\end{longtable}

\textbf{证明}：设 \(X\) 为 del Pezzo 曲面，\(d = K_X^2\)。由 Castelnuovo 有理曲面判准，\(X\) 是有理曲面，故 \(X\) 双有理等价于 \(\mathbb{P}^2\) 或 Hirzebruch 曲面 \(\mathbb{F}_a\)（\(a \geq 0, a \neq 1\)）。由于 \(K_X\) 为反丰沛，\(X\) 不能包含 \((-1)\)-曲线以外的负曲线，且 \(X\) 必为 \(\mathbb{P}^2\) 的爆破，爆破点数目为 \(9-d\)。对 \(d \geq 3\)，反典范线性系统 \(|-K_X|\) 无基点且非常丰沛，将 \(X\) 嵌入 \(\mathbb{P}^d\) 为 \(d\) 次曲面（由 Riemann-Roch 计算 \(\dim H^0(-K_X) = d+1\)）。对 \(d=2\)，\(|-K_X|\) 为基点自由的二维线性系统，定义 \(2:1\) 覆叠 \(X \to \mathbb{P}^2\)，分支除子为四次曲线（由分歧公式 \(K_X = f^*K_{\mathbb{P}^2} + R\) 及 \(R \in |-2K_X|\) 推出）。对 \(d=1\)，\(|-K_X|\) 有一基点，\(|-2K_X|\) 无基点，定义 \(2:1\) 覆叠 \(X \to \mathbb{P}(1,1,2)\)，分支于六次曲线。\(d=8\) 时：若 \(X\) 为 \(\mathbb{P}^2\) 上 1 点爆破，\(K_X^2 = 8\)；若 \(X \cong \mathbb{P}^1 \times \mathbb{P}^1\)，\(K_X^2 = 8\)（由 \(K_X \sim -2F_1 - 2F_2\) 计算）。两者对应于不同极值射线，均属 \(d=8\)。\(\blacksquare\)

\emph{注：} \(d = 3\) 的 del Pezzo 曲面即为 \(\mathbb{P}^3\) 中的光滑三次曲面。著名的 27 条直线定理（Cayley-Salmon）断言：\(\mathbb{P}^3\) 中每个光滑三次曲面上恰好有 27 条直线。

\textbf{定理 86.2}（三维 Fano 簇的分类，Iskovskikh 1977-1979，Mori-Mukai 1981-2003）：三维光滑 Fano 簇恰好有 \textbf{105 个形变族}（deformation families），按 Fano 指标 \(r\) 和 Picard 数 \(\rho\) 分类：

\begin{itemize}
\tightlist
\item
  Fano 指标 \(r = 4\)：仅 \(\mathbb{P}^3\)（1 族）；
\item
  Fano 指标 \(r = 3\)：二次超曲面 \(Q^3 \subset \mathbb{P}^4\)（1 族）；
\item
  Fano 指标 \(r = 2\)：包括 del Pezzo 三维簇（36 族，最大 Picard 数为 10）；
\item
  Fano 指标 \(r = 1\)：其余 67 族（包括著名的素数 Fano 三维簇 \(V_{22}\)，由 Mukai 发现）。
\end{itemize}

\textbf{证明}：Iskovskikh和Mori-Mukai的分类分10个形变族。（1）Iskovskikh（1977-1979）使用双有理几何方法：极小模型程序（MMP）将Fano簇收缩到Mori纤维空间。对三维Fano簇，反典范除子\(-K_X\)为丰沛，\((-K_X)^3\)（Fano次数）为\(1\)到\(22\)的整数。（2）Mori-Mukai（1981-2003）完成分类：利用extremal rays（极值射线）的收缩映射和Mori锥定理，将三维Fano簇按第二Betti数\(b_2\)和Fano次数分类。\(b_2=1\)有\(17\)族，\(b_2\ge 2\)有\(88\)族，共\(105\)个形变族。\(\blacksquare\)

\subsection{86.3 Mori 纤维空间与有理连通簇}\label{mori-ux7ea4ux7ef4ux7a7aux95f4ux4e0eux6709ux7406ux8fdeux901aux7c07}

\textbf{定义 86.3}（Mori 纤维空间）：\textbf{Mori 纤维空间}（Mori Fiber Space）是带有纤维型极射线收缩 \(\varphi: X \to Y\) 的 \(\mathbb{Q}\)-因式簇 \(X\)，其中 \(X\) 仅有 terminal 奇点，\(Y\) 是 \(\mathbb{Q}\)-因式簇且 \(\dim Y < \dim X\)，一般纤维是 Fano 簇（带有 terminal 奇点）。

Mori 纤维空间是极小模型纲领（MMP）的三大输出类型之一。每个光滑射影簇通过 MMP 最终分解为 Mori 纤维空间（经由极小模型和翻转序列）。

\textbf{定义 86.4}（有理连通簇）：代数闭域上的光滑射影簇 \(X\) 称为\textbf{有理连通的}（rationally connected），若其中的一般两点可以由有理曲线相连。

更精确地，存在有理映射 \(\pi: \mathbb{P}^1 \times \cdots \times \mathbb{P}^1 \dashrightarrow X\) 使其像覆盖 \(X\) 的开稠子集。

\textbf{定理 86.3}（Kollár-Miyaoka-Mori 定理，1992 年）：特征 0 的代数闭域上的光滑 Fano 簇是有理连通的。更一般地，光滑 Fano 簇是\textbf{有理链连通的}（rationally chain connected）------任意两点可由若干条有理曲线的链相连。

\textbf{证明}：Kollár-Miyaoka-Mori的证明基于特征\(p\)约化方法。核心步骤：（1）将光滑Fano簇\(X\)约化到正特征\(p\gg 0\)，利用Frobenius映射的性质。（2）Mori的弯曲-断裂引理：若\(K_X\)不是nef，则存在通过一般点的有理曲线。对Fano簇，\(-K_X\)为丰沛，故存在有理曲线。（3）有理连通性：若\(X\)上任意两点可由有理曲线链连接，则\(X\)为有理连通。KMM证明：Fano簇上存在充分多的有理曲线，使得有理连通性成立。证明使用变形理论和Mori的极值有理曲线理论。\(\blacksquare\)

\textbf{推论 86.4}：Fano 簇上的 Chow 群 \(\operatorname{CH}_0(X)\) 是平凡群------即 \(X\) 上任意两点模有理等价后相等。

\subsection{86.4 Luroth 问题与有理性}\label{luroth-ux95eeux9898ux4e0eux6709ux7406ux6027}

\textbf{Luroth 问题}（1876 年）：设 \(k \subset K \subset k(t_1, \ldots, t_n)\) 是嵌入在有理函数域中的子域。则 \(K\) 本身是否有理（即同构于某个 \(m\) 个变量的有理函数域）？

\(n = 1\) 时 Luroth 定理给出肯定回答（\(K \cong k(t)\)）。\(n = 2\) 时 Castelnuovo 定理（1894 年）在代数几何语言中断言：特征 0 的代数闭域上，单有理曲面必为有理曲面。

\(n \geq 3\) 时 Luroth 问题的回答是\textbf{否定的}。反例构造是代数几何的重大里程碑：

\textbf{定理 86.5}（三维 Luroth 问题的反例）：

\begin{enumerate}
\def\labelenumi{(\arabic{enumi})}
\item
  \textbf{Artin-Mumford 反例}（1972 年）：利用 Brauer 群的三阶挠元（扭 3-形式）构造了特征 0 的代数闭域上的单有理三维簇 \(X\)，满足 \(H^3(X, \mathbb{Z})_{\text{tors}} \neq 0\)。由于任何有理簇均有无挠的上同调，\(X\) 不能是有理的。
\item
  \textbf{Clemens-Griffiths 反例}（1972 年）：三维光滑三次超曲面 \(V_3 \subset \mathbb{P}^4\) 是单有理的，但其\textbf{中间 Jacobian} \(J(V_3)\)（二维复环面）不是 Jacobian 簇的直和项，故 \(V_3\) 不是有理的。中间 Jacobian 是有理曲面不存在的不变量。
\item
  \textbf{Iskovskikh-Manin 反例}（1971 年）：四维光滑三次超曲面 \(V_3 \subset \mathbb{P}^5\) 不是有理的。证法基于双有理自同构群的有限性------证明 \(\operatorname{Bir}(V_3)\) 是有限群，而有理簇的双有理自同构群是无界大的（Cremona 群 \(\operatorname{Cr}_n = \operatorname{Bir}(\mathbb{P}^n)\) 当 \(n \geq 2\) 时无限）。
\end{enumerate}

\textbf{证明}：分别证明三个反例。(1) Artin-Mumford 反例：考虑 \(\mathbb{P}^3\) 中由二次型 \(Q_0, Q_1, Q_2\) 定义的二次超曲面束。取 \(\mathbb{P}^3\) 的二次覆盖 \(X \to \mathbb{P}^3\)，沿某四次曲面 \(B\) 分支。\(X\) 有节点奇点，其小解消 \(X'\) 为光滑三维簇。\(X'\) 是单有理的（由二次覆盖的构造），但 \(H^3(X', \mathbb{Z})_{\text{tors}} \cong \mathbb{Z}/2\mathbb{Z}\)（来源于 Brauer 群 \(\operatorname{Br}(X')[2]\) 的非平凡元素）。有理簇的上同调群均无挠（有理簇为 \(\mathbb{P}^n\) 的爆破，每次爆破仅增加 \(H^2\) 与 \(H^4\) 的自由部分），故 \(X'\) 非有理。(2) Clemens-Griffiths 反例：光滑三次三维簇 \(V_3 \subset \mathbb{P}^4\) 的中间 Jacobian \(J(V_3) = H^{1,2}(V_3)/H^3(V_3,\mathbb{Z})\) 是 \(5\) 维复环面，其主极化 \(\Theta\)-除子不可约。有理三维簇的中间 Jacobian 必为 Jacobian 簇的直和（因为有理簇的 Hodge 结构可分解为曲线的 Jacobian 的直和），而 \(J(V_3)\) 的 \(\Theta\)-除子不可约意味着它不是 Jacobian 的直和，故 \(V_3\) 非有理。(3) Iskovskikh-Manin 反例：设 \(V_3 \subset \mathbb{P}^5\) 为光滑四次三维簇。\(V_3\) 是 Fano 簇（\(K_{V_3} = -H\)），其双有理自同构群 \(\operatorname{Bir}(V_3)\) 有限。证明使用等变极小模型纲领：\(V_3\) 上的每个双有理自映射均可分解为 \(V_3\) 上的自同构，而 \(\operatorname{Aut}(V_3)\) 为有限群。然而 \(\mathbb{P}^4\) 的 Cremona 群 \(\operatorname{Cr}_4 = \operatorname{Bir}(\mathbb{P}^4)\) 是无限的，故 \(V_3\) 不与 \(\mathbb{P}^4\) 双有理等价，即非有理。\(\blacksquare\)

这些反例共同说明：单有理性和有理性之间的鸿沟随维数增长而加大，高维有理性问题是极其微妙和困难的。

\hrulefill

\hrulefill

\hrulefill

\hrulefill

\hrulefill

\hrulefill

\section{第289章：极化簇与 Hilbert 概型}\label{ux7b2c289ux7ae0ux6781ux5316ux7c07ux4e0e-hilbert-ux6982ux578b}

极化簇和 Hilbert 概型是现代代数几何中模空间理论的核心工具。Hilbert 概型将闭子概型的参数化问题形式化为射影概型结构，而极化簇的 GIT（几何不变量理论）则为构造代数簇的模空间提供了标准方法。

\subsection{87.1 Hilbert 多项式与 Hilbert 概型}\label{hilbert-ux591aux9879ux5f0fux4e0e-hilbert-ux6982ux578b}

\textbf{定义 87.1}（极化簇）：\textbf{极化簇}是一个对 \((X, L)\)，其中 \(X\) 是代数闭域 \(k\) 上的射影簇，\(L \in \operatorname{Pic}(X)\) 是 \(X\) 上的丰沛线丛（极化）。极化簇的 Hilbert 多项式为

\[P(m) = \chi(X, L^{\otimes m}) = \sum_{i=0}^{\dim X} (-1)^i \dim_k H^i(X, L^{\otimes m})\]

由 Serre 消没定理，当 \(m \gg 0\) 时，\(H^i(X, L^{\otimes m}) = 0\) 对所有 \(i > 0\) 成立，故

\[P(m) = \dim_k H^0(X, L^{\otimes m}), \quad m \gg 0\]

由 Riemann-Roch 定理（Hirzebruch 版本），Hilbert 多项式的主导项为

\[P(m) = \frac{L^n}{n!} m^n + \text{（低次项）}\]

其中 \(n = \dim X\)，\(L^n = (L^{\otimes n})\) 的相交数即为 \(L\) 的次数。

\textbf{定理 87.1}（Grothendieck 的 Hilbert 概型定理，1961 年）：设 \(\mathbb{P}^N_k\) 是域 \(k\) 上的射影空间。对任意数值多项式 \(P(t) \in \mathbb{Q}[t]\)，存在射影概型 \(\operatorname{Hilb}_P(\mathbb{P}^N_k)\)，其 \(k\)-点自然一一对应于 \(\mathbb{P}^N_k\) 中 Hilbert 多项式为 \(P(t)\) 的闭子概型。而且该对应具有\textbf{万有性质}：存在万有闭子概型

\[\mathcal{Z} \subseteq \operatorname{Hilb}_P(\mathbb{P}^N_k) \times \mathbb{P}^N_k\]

使得对任意概型 \(S\) 和任意 \(S\)-平坦的闭子概型 \(Y \subseteq S \times \mathbb{P}^N_k\)（纤维的 Hilbert 多项式均为 \(P(t)\)），存在唯一的态射 \(f: S \to \operatorname{Hilb}_P(\mathbb{P}^N_k)\) 使得 \(Y = (f \times \operatorname{id})^* \mathcal{Z}\)。

\textbf{证明}：Grothendieck的构造：（1）将Hilbert函子\(\underline{\operatorname{Hilb}}_{\mathbb{P}^N/k}\)定义为\(\underline{\operatorname{Hilb}}(S)=\{\)闭子概型\(Z\subseteq\mathbb{P}^N_S\)，平坦于\(S\)且Hilbert多项式为\(P(t)\}\)。（2）利用Grassmannian和射影空间的嵌入，将\(\underline{\operatorname{Hilb}}\)嵌入某个Grassmannian，再嵌入射影空间。（3）证明\(\underline{\operatorname{Hilb}}\)是可表函子（由某个射影概型\(\operatorname{Hilb}^P(\mathbb{P}^N_k)\)表示），使用Mumford的平坦性准则和Castelnuovo-Mumford正则性。\(\operatorname{Hilb}^P\)的射影性来自嵌入和Grassmannian的射影性。\(\blacksquare\)

\subsection{87.2 模空间的构造：稳定曲线}\label{ux6a21ux7a7aux95f4ux7684ux6784ux9020ux7a33ux5b9aux66f2ux7ebf}

\textbf{定理 87.2}（Deligne-Mumford 稳定曲线模空间 \(\overline{M}_g\)）：设 \(g \geq 2\)。亏格 \(g\) 的\textbf{稳定曲线}的模叠 \(\overline{\mathcal{M}}_g\) 是光滑、真（proper）且不可约的 Deligne-Mumford 叠。其粗模空间 \(\overline{M}_g\) 是 \(3g-3\) 维射影簇。

\textbf{证明}：\(\overline{\mathcal{M}}_g\)的参数化：（1）稳定曲线：连通、节点为唯一奇点、\(\omega_C\)为丰沛（等价于\(\operatorname{Aut}(C)\)有限且\(g\ge 2\)时\(2g-2>0\)）的曲线。（2）Deligne-Mumford（1969）证明\(\overline{\mathcal{M}}_g\)为光滑真Deligne-Mumford叠，维数为\(3g-3\)。边界\(\overline{\mathcal{M}}_g\setminus\mathcal{M}_g=\bigcup_{i=0}^{\lfloor g/2\rfloor}\Delta_i\)由可约稳定曲线组成，为正规交叉除子。\(\overline{\mathcal{M}}_g\)的射影性由Mumford的GIT构造和Knudsen-Mumford的稳定性定理保证。\(\blacksquare\)

一条曲线 \(C\) 是\textbf{稳定的}，若：(1) \(C\) 是连通的约化（reduced）曲线，仅有结点（node）奇点；(2) \(C\) 的光滑有理分支包含至少 3 个特殊点（与曲线其余部分的相交点）；(3) 其典范层 \(\omega_C\) 是丰沛的（等价地，对偶图是稳定的）。

\emph{构造思路}：利用 Hilbert 概型。固定充分大的 \(m\)（如 \(m = 3\)），考虑 \(\omega_C^{\otimes m}\) 的整体截面给出了嵌入 \(C \hookrightarrow \mathbb{P}^{N}\)（由 Riemann-Roch 定理，\(N = (2m-1)(g-1)-1\)）。所有稳定曲线的 \(m\)-典范嵌入构成 Hilbert 概型中的局部闭子集。Hilbert-Mumford 稳定性条件保证该子集对 \(PGL(N+1)\) 作用是 GIT 半稳定的，取 GIT 商即得 \(\overline{M}_g\)。∎

\textbf{极化 Abelian 簇的模空间}：亏格 \(g\) 的主极化 Abel 簇的粗模空间 \(\mathcal{A}_g\) 由 Siegel 上半空间 \(\mathbb{H}_g\) 经 \(Sp(2g, \mathbb{Z})\) 作用取商得到。\(\mathcal{A}_g\) 是 \(g(g+1)/2\) 维拟射影簇。其紧化包括： - \textbf{Satake 紧化}（Baily-Borel 紧化，极小紧化）； - \textbf{toroidal 紧化}（由 Faltings-Chai 描述，可容许退化）； - \textbf{二阶 Voronoi 紧化}（Alexeev，完全显式）。

\subsection{87.3 GIT 稳定性与 Hilbert-Mumford 数值准则}\label{git-ux7a33ux5b9aux6027ux4e0e-hilbert-mumford-ux6570ux503cux51c6ux5219}

\textbf{定义 87.2}（GIT 的稳定性条件，Mumford 1965 年）：设 \(G\) 是代数闭域 \(k\) 上的简约代数群，\(L\) 是射影概型 \(X\) 上的丰沛 \(G\)-等变线丛。点 \(x \in X\) 称为：

\begin{itemize}
\tightlist
\item
  \textbf{半稳定的}（semistable），若存在 \(n > 0\) 和 \(G\)-不变截面 \(s \in H^0(X, L^{\otimes n})^G\) 使得 \(s(x) \neq 0\)；
\item
  \textbf{稳定的}（stable），若 \(x\) 是半稳定的，轨道 \(G \cdot x\) 在 \(X^{ss}\) 中是闭的，且稳定子群 \(G_x\) 是有限群；
\item
  \textbf{真稳定的}（properly stable），若 \(x\) 是稳定的且稳定子群 \(G_x\) 的维数为 0。
\end{itemize}

半稳定点集 \(X^{ss}\) 是 \(X\) 的开子集。GIT 商 \(X /\!/ G = \operatorname{Proj}\left(\bigoplus_{n \geq 0} H^0(X, L^{\otimes n})^G\right)\) 是射影簇。

\textbf{定理 87.3}（Hilbert-Mumford 数值准则）：设 \(G\) 是简约代数群，\(\lambda: \mathbb{G}_m \to G\) 是单参数子群（1-PSG）。对 \(x \in X\)，定义 Hilbert-Mumford 数值

\[\mu(x, \lambda) = -\operatorname{ord}_{t=0}(\lambda(t) \cdot \tilde{x})\]

其中 \(\tilde{x}\) 是 \(x\) 在某个 \(G\)-等变化下的提升。则

\begin{itemize}
\tightlist
\item
  \(x\) 是 \textbf{半稳定的} 当且仅当对所有 1-PSG \(\lambda\)，\(\mu(x, \lambda) \geq 0\)；
\item
  \(x\) 是 \textbf{稳定的} 当且仅当对所有非平凡 1-PSG \(\lambda\)，\(\mu(x, \lambda) > 0\)。
\end{itemize}

\textbf{证明}：Hilbert-Mumford数值准则判定射影空间\(\mathbb{P}^N\)中点的GIT稳定性。对点\(x\in\mathbb{P}^N\)和单参数子群\(\lambda:\mathbb{G}_m\to G\)，\(\lambda\)作用于\(x\)的轨道闭包。定义\(\mu(x,\lambda)=\min_i\{i: \text{$\lambda(t)\cdot x$的第$i$个坐标不全为零的$t$的最低次幂}\}\)（\(t\to 0\)的权）。\(x\)为半稳定（stable）当且仅当对所有非平凡单参数子群\(\lambda\)，\(\mu(x,\lambda)\le 0\)（\(<0\)）。证明：Hilbert-Mumford（1963）将GIT稳定性归结为环面作用（由单参数子群参数化），利用环面的权重空间分解和\(t\to 0\)的渐近行为。\(\blacksquare\)

\textbf{应用于稳定曲线}：曲线的 GIT 稳定性（Hilbert-Mumford 稳定性）恰好对应于 Deligne-Mumford 稳定性条件。一条曲线 \(C\) 在 GIT 意义下是半稳定的（关于 \(m\)-典范嵌入）当且仅当它是 Deligne-Mumford 意义下的稳定曲线。

\textbf{模空间理论的现代发展与前沿}

模空间理论是 20 世纪代数几何的最高成就之一，其核心是参数化具有特定性质的代数几何对象（如曲线、向量丛、层、Hitchin 对等）的连续族。\textbf{Deligne-Mumford 叠}（1969）解决了模空间不精细的问题：当被参数化的对象有非平凡自同构时，精细模空间通常不存在（因为不能将自同构''遗忘''），但 DM 叠允许将自同构群作为叠结构的一部分保留。DM 叠 \(\mathcal{M}\) 是一个纤维化范畴，其点不仅是对象，还携带自同构的群胚结构。\textbf{稳定曲线的模空间} \(\overline{\mathcal{M}}_g\) 是 DM 叠的典范实例：\(\overline{\mathcal{M}}_g\) 是光滑亏格 \(g\) 曲线的模空间 \(\mathcal{M}_g\) 的 Deligne-Mumford 紧化，其边界点对应于结点曲线（nodal curves），且 \(\overline{\mathcal{M}}_g\) 是光滑的适定 DM 叠。Kontsevich 利用 \(\overline{\mathcal{M}}_{g,n}\)（\(n\) 个标记点的稳定曲线模空间）的相交理论定义了 Gromov-Witten 不变量，这直接导致了镜面对称的数学理论。\textbf{向量丛的模空间}和\textbf{Hitchin 系}的模空间是几何 Langlands 纲领的核心对象：Hitchin 纤维化 \(\mathcal{M}_{\text{Higgs}} \to \mathcal{A}\)（\(\mathcal{A}\) 为 Hitchin 基）将 Higgs 丛的模空间映射到特征多项式空间，其一般纤维是 Abel 簇（Prym 簇），这构成了 Hitchin 可积系统。\textbf{模空间的导出结构}：导出代数几何（Toën-Vezzosi, Lurie）将模空间提升为导出叠（derived stack），其切复形 \(\mathbb{T}_x\mathcal{M}\) 编码了完整的形变-障碍理论，克服了经典模空间在非光滑点处的病态行为。\textbf{模空间的壁-腔结构}（wall-crossing）：在 Bridgeland 稳定条件的研究中，模空间的拓扑型在参数的''壁''处发生突变，Joyce-Song 的壁-腔公式通过 Hall 代数技术描述了 Donaldson-Thomas 不变量在壁-腔跳跃中的变化规律。

\newpage

\chapter{08-代数几何与数论/01-代数几何/05-范畴论.md}\label{ux4ee3ux6570ux51e0ux4f55ux4e0eux6570ux8bba01-ux4ee3ux6570ux51e0ux4f5505-ux8303ux7574ux8bba.md}

\subsection{224.A 范畴论补充}\label{a-ux8303ux7574ux8bbaux8865ux5145}

范畴论由 Eilenberg 和 Mac Lane 于 1945 年在代数拓扑的研究中创立，最初目的是为自然变换（如 \(H_n(X) \to H_n(Y)\) 的自然性）提供严格定义，但随后发展成被称为''数学的数学''的元理论。范畴论的独特力量在于其组织原理和语言统一性：它从各种数学领域中提取共同的模式------对象间的关系而非对象本身------这在同调代数（长正合序列作为三角范畴的态射）、代数几何（平展映射作为概形范畴中的上同调下降）和逻辑学（拓扑斯作为直觉主义集合论的范畴模型）中反复验证。范畴论的思想也为计算机科学（函数式编程中函子和单子模式）、数学物理（拓扑量子场论作为配边范畴的对称幺半函子）和数据库理论（函子数据模型）提供了普适的概念框架。本章勾勒范畴论的整体图景与组织原理，为后续五章（第 225-228 章）的具体理论展开提供路线图。

\hrulefill

\hrulefill

\hrulefill

\hrulefill

\hrulefill

\hrulefill

\hrulefill

\section{第290章：范畴、函子与自然变换}\label{ux7b2c290ux7ae0ux8303ux7574ux51fdux5b50ux4e0eux81eaux7136ux53d8ux6362}

范畴论的基本三位一体------范畴、函子、自然变换------构成了现代数学中最高效的组织语言。范畴的概念将数学结构（群、环、拓扑空间、概形）以及态射（同态、连续映射、概形态射）统一为''对象与箭头''的模式，而函子则是范畴之间的结构保持映射，将不同领域的类似构造统一起来（如基本群作为拓扑空间范畴到群范畴的函子）。自然变换则是函子之间的态射------这个看似简单的概念实际上正是范畴论创立的原始动机：Eilenberg 和 Mac Lane 在 1945 年正式定义自然变换以明确表述 ``双对偶向量空间的典范同构是自然的'' 这类直觉。Yoneda 引理是范畴论最基本的定理：它断言一个范畴中的对象完全由它与其他对象之间的关系（即所有指向它的态射的集合）决定，这一洞察将范畴论从语言工具提升为深刻的数学原理。本章从范畴的定义出发，逐步建立函子的合成、自然变换的垂直与水平合成、Yoneda 嵌入和可表函子的理论。

\subsection{94.1 范畴的定义}\label{ux8303ux7574ux7684ux5b9aux4e49}

\textbf{定义 94.1}（范畴）：一个\textbf{范畴} \(\mathcal{C}\) 由以下数据组成： - 一族\textbf{对象} \(\operatorname{Ob}(\mathcal{C})\) - 对每对对象 \(X, Y\)，一个\textbf{态射集} \(\operatorname{Hom}_{\mathcal{C}}(X, Y)\)（或 \(\mathcal{C}(X, Y)\)） - 对每个对象 \(X\)，一个\textbf{恒等态射} \(\operatorname{id}_X \in \mathcal{C}(X, X)\) - 对每三个对象 \(X, Y, Z\)，一个\textbf{合成映射} \(\circ: \mathcal{C}(Y, Z) \times \mathcal{C}(X, Y) \to \mathcal{C}(X, Z)\)

满足：合成结合律 \((h \circ g) \circ f = h \circ (g \circ f)\)，恒等律 \(\operatorname{id}_Y \circ f = f = f \circ \operatorname{id}_X\)。

\textbf{定义 94.2}（小范畴与局部小范畴）：范畴 \(\mathcal{C}\) 称为\textbf{局部小}的，如果每个 \(\mathcal{C}(X, Y)\) 是集合。\(\mathcal{C}\) 称为\textbf{小范畴}，如果 \(\operatorname{Ob}(\mathcal{C})\) 也是集合。

\textbf{定义 94.3}（同构）：态射 \(f: X \to Y\) 称为\textbf{同构}，如果存在 \(g: Y \to X\) 使得 \(g \circ f = \operatorname{id}_X\)，\(f \circ g = \operatorname{id}_Y\)。\(g\) 称为 \(f\) 的逆。

\textbf{定义 94.4}（对偶范畴）：范畴 \(\mathcal{C}\) 的\textbf{对偶范畴} \(\mathcal{C}^{\text{op}}\) 与 \(\mathcal{C}\) 有相同的对象，但态射方向反转：\(\mathcal{C}^{\text{op}}(X, Y) = \mathcal{C}(Y, X)\)。

\subsection{94.2 函子}\label{ux51fdux5b50}

\textbf{定义 94.5}（函子）：范畴 \(\mathcal{C}\) 到 \(\mathcal{D}\) 的\textbf{函子} \(F: \mathcal{C} \to \mathcal{D}\) 由以下组成： - 对象映射：\(X \mapsto F(X)\) - 态射映射：对每个 \(f: X \to Y\)，\(F(f): F(X) \to F(Y)\)

满足：\(F(\operatorname{id}_X) = \operatorname{id}_{F(X)}\)，\(F(g \circ f) = F(g) \circ F(f)\)。

\textbf{定义 94.6}（反变函子）：\(\mathcal{C}^{\text{op}} \to \mathcal{D}\) 的函子称为 \(\mathcal{C}\) 到 \(\mathcal{D}\) 的\textbf{反变函子}。它反转态射方向：\(F(f: X \to Y): F(Y) \to F(X)\)。

\subsection{94.3 自然变换}\label{ux81eaux7136ux53d8ux6362}

\textbf{定义 94.7}（自然变换）：设 \(F, G: \mathcal{C} \to \mathcal{D}\) 是两个函子。\textbf{自然变换} \(\alpha: F \Rightarrow G\) 由对每个 \(X \in \mathcal{C}\) 的一个态射 \(\alpha_X: F(X) \to G(X)\) 组成，使得对每个 \(f: X \to Y\)，下图交换：

\[
\begin{CD}
F(X) @>{\alpha_X}>> G(X) \\
@V{F(f)}VV @VV{G(f)}V \\
F(Y) @>>{\alpha_Y}> G(Y)
\end{CD}
\]

若每个 \(\alpha_X\) 是同构，则 \(\alpha\) 称为\textbf{自然同构}，记作 \(F \cong G\)。

\textbf{定义 94.8}（函子范畴）：从 \(\mathcal{C}\) 到 \(\mathcal{D}\) 的函子及它们之间的自然变换构成\textbf{函子范畴} \([\mathcal{C}, \mathcal{D}]\)（或 \(\mathcal{D}^{\mathcal{C}}\)）。

\subsection{94.4 Yoneda 引理}\label{yoneda-ux5f15ux7406}

\textbf{定义 94.9}（Yoneda 嵌入）：对局部小范畴 \(\mathcal{C}\)，定义 \textbf{Yoneda 嵌入}

\[Y: \mathcal{C} \to [\mathcal{C}^{\text{op}}, \mathbf{Set}], \quad Y(X) = \mathcal{C}(-, X)\]

即 \(Y(X)\) 是反变 Hom 函子：\(Y(X)(A) = \mathcal{C}(A, X)\)。

\textbf{定理 94.1}（Yoneda 引理）：设 \(\mathcal{C}\) 是局部小范畴，\(F: \mathcal{C}^{\text{op}} \to \mathbf{Set}\) 是函子，\(X \in \mathcal{C}\)。则存在自然的一一对应

\[\operatorname{Nat}(\mathcal{C}(-, X), F) \cong F(X)\]

（其中 \(\operatorname{Nat}\) 表示自然变换的集合）。此同构由 \(\alpha \mapsto \alpha_X(\operatorname{id}_X)\) 给出。

\textbf{证明}：定义映射 \(\Phi: \operatorname{Nat}(\mathcal{C}(-, X), F) \to F(X)\)，\(\Phi(\alpha) = \alpha_X(\operatorname{id}_X)\)。逆映射 \(\Psi: F(X) \to \operatorname{Nat}(\mathcal{C}(-, X), F)\) 定义为：对 \(u \in F(X)\)，\(\Psi(u)_A(f) = F(f)(u)\)（\(f: A \to X\)）。验证 \(\Phi \circ \Psi = \operatorname{id}\) 和 \(\Psi \circ \Phi = \operatorname{id}\)。∎

\textbf{推论 94.2}（Yoneda 嵌入是完全忠实的）：Yoneda 嵌入 \(Y: \mathcal{C} \to [\mathcal{C}^{\text{op}}, \mathbf{Set}]\) 是完全忠实的：\(\operatorname{Nat}(\mathcal{C}(-, X), \mathcal{C}(-, Y)) \cong \mathcal{C}(X, Y)\)。特别地，\(X \cong Y\) 当且仅当 \(\mathcal{C}(-, X) \cong \mathcal{C}(-, Y)\)。

\textbf{证明}：在 Yoneda 引理（定理 94.1）中取 \(F = \mathcal{C}(-, Y)\)，则 \(\operatorname{Nat}(\mathcal{C}(-, X), \mathcal{C}(-, Y)) \cong \mathcal{C}(X, Y)\)。具体地，自然变换 \(\alpha: \mathcal{C}(-, X) \Rightarrow \mathcal{C}(-, Y)\) 由 \(\alpha_X(\operatorname{id}_X) \in \mathcal{C}(X, Y)\) 唯一决定；反之，每个 \(f: X \to Y\) 通过后合成 \(f \circ (-)\) 给出自然变换。这建立了双向一一对应，故 Yoneda 嵌入是全忠实函子。∎

\textbf{推论 94.3}（Yoneda 推论）：一个对象 \(X\) 完全由它与其他对象的关系（即函子 \(\mathcal{C}(-, X)\)）决定（在同构意义下）。这是范畴论的基本哲学。

\hrulefill

\hrulefill

\hrulefill

\hrulefill

\hrulefill

\hrulefill

\section{第291章：极限与余极限}\label{ux7b2c291ux7ae0ux6781ux9650ux4e0eux4f59ux6781ux9650}

极限和余极限是范畴论中最核心的构造性概念，它们将数学中无数看似独立的构造统一为同一模式：乘积、纤维积（拉回）、等化子、完备化空间和射影极限都是极限的特例，而余积、推出（pushout）、余等化子和归纳极限都是余极限的特例。这一统一的威力难以高估------它意味着一旦在某个范畴中证明了极限的存在性，便自动获得了该范畴中所有特殊类型极限的存在性。极限构造在数学各个分支中反复出现：在拓扑学中，Stone-Cech 紧化是幂等化子类型的极限；在代数几何中，概形的纤维积是拉回类型的极限；在逻辑学中，全称量词 \(\forall\) 是右 Kan 扩张的特例。极限与 Yoneda 引理（第 225 章）和伴随函子（第 227 章）有深刻的内在联系：右伴随保持极限，左伴随保持余极限------这一性质既是伴随函子理论的核心工具，也是''自由''构造与''遗忘''构造对称性的代数体现。

\subsection{95.1 极限的普遍定义}\label{ux6781ux9650ux7684ux666eux904dux5b9aux4e49}

\textbf{定义 95.1}（图与锥）：设 \(J\) 是小范畴（称为\textbf{指标范畴}）。函子 \(D: J \to \mathcal{C}\) 称为 \(\mathcal{C}\) 中的一个 \textbf{\(J\) 型图}。

\(D\) 上的一个\textbf{锥}（cone）由对象 \(C \in \mathcal{C}\) 和一族态射 \(\{c_j: C \to D(j)\}_{j \in J}\) 组成，使得对每个 \(J\) 中的态射 \(u: j \to k\)，\(D(u) \circ c_j = c_k\)。

\textbf{定义 95.2}（极限）：图 \(D: J \to \mathcal{C}\) 的\textbf{极限} \(\varprojlim D\)（或 \(\lim D\)）是 \(D\) 上的一个万有锥：它是一个锥 \((L, \{\ell_j\})\)，使得对任意其他锥 \((C, \{c_j\})\)，存在唯一的态射 \(f: C \to L\) 使得 \(\ell_j \circ f = c_j\)（对所有 \(j\)）。

\textbf{定义 95.3}（余极限）：对偶地，\(D\) 上的\textbf{余锥}及\textbf{余极限} \(\varinjlim D\)（或 \(\operatorname{colim} D\)）是极限在 \(\mathcal{C}^{\text{op}}\) 中的对偶概念。

\textbf{命题 95.1}：若极限存在，则它在同构意义下唯一。

极限的唯一性来自其万有性质------这是范畴论中普遍构造的典型特征：任何由万有性质定义的对象，若存在则在同构意义下唯一。这一命题的证明是理解范畴论中''唯一性''论证范式的最佳入门。

\textbf{证明}：设 \((L, \{\ell_j\})\) 和 \((L', \{\ell'_j\})\) 均为图 \(D: J \to \mathcal{C}\) 的极限。由 \(L\) 的万有性质，存在唯一态射 \(f: L' \to L\) 使得 \(\ell_j \circ f = \ell'_j\)（对所有 \(j\)）。同理，由 \(L'\) 的万有性质，存在唯一态射 \(g: L \to L'\) 使得 \(\ell'_j \circ g = \ell_j\)（对所有 \(j\)）。于是 \(\ell_j \circ (f \circ g) = \ell'_j \circ g = \ell_j\)，且 \(\ell_j \circ \operatorname{id}_L = \ell_j\)。由 \(L\) 的万有锥的唯一性，\(f \circ g = \operatorname{id}_L\)。同理 \(g \circ f = \operatorname{id}_{L'}\)。故 \(f\) 为同构，\(L \cong L'\)。\(\blacksquare\)

\subsection{95.2 特殊极限与余极限}\label{ux7279ux6b8aux6781ux9650ux4e0eux4f59ux6781ux9650}

\textbf{定义 95.4}（积与余积）：设 \(J\) 是离散范畴（只有恒等态射）。\(J\) 型图的极限是\textbf{积} \(\prod_{j \in J} X_j\)，余极限是\textbf{余积} \(\coprod_{j \in J} X_j\)。

\begin{itemize}
\tightlist
\item
  \textbf{Set} 中：积 = 笛卡尔积，余积 = 不交并
\item
  \textbf{Grp} 中：积 = 直积，余积 = 自由积
\item
  \textbf{Ab} 中：积 = 直积，余积 = 直和（有限时两者相同）
\item
  \textbf{Vect\(_k\)} 中：积 = 直积，余积 = 直和
\end{itemize}

\textbf{定义 95.5}（等化子与余等化子）：两个平行态射 \(f, g: X \rightrightarrows Y\) 的图。 - \textbf{等化子} \(\operatorname{eq}(f, g)\) 是使得 \(f \circ e = g \circ e\) 的万有态射 \(e: E \to X\) - \textbf{余等化子} \(\operatorname{coeq}(f, g)\) 是使得 \(c \circ f = c \circ g\) 的万有态射 \(c: Y \to C\)

\textbf{定义 95.6}（拉回与推出）： - \textbf{拉回}（pullback）是图 \(X \to Z \leftarrow Y\) 的极限 - \textbf{推出}（pushout）是图 \(X \leftarrow Z \to Y\) 的余极限

\textbf{定理 95.2}（极限的构造）：若 \(\mathcal{C}\) 有所有积和所有等化子，则 \(\mathcal{C}\) 有所有小极限。具体地，\(\varprojlim D\) 是等化子

\[\varprojlim D \longrightarrow \prod_{j \in J} D(j) \rightrightarrows \prod_{u: j \to k} D(k)\]

其中两个态射分别由 \(D(u) \circ \pi_j\) 和 \(\pi_k\) 给出。

\emph{证明}：设图 \(D: J \to \mathcal{C}\)。构造积 \(P = \prod_{j} D(j)\) 和 \(Q = \prod_{u: j \to k} D(k)\)。定义两个态射 \(f, g: P \rightrightarrows Q\)：\(f\) 的第 \(u\) 分量为 \(D(u) \circ \pi_j\)，\(g\) 的第 \(u\) 分量为 \(\pi_k\)。取等化子 \(e: L \to P\)，则 \((L, \pi_j \circ e)\) 构成 \(D\) 的锥。对任意锥 \((C, c_j)\)，泛性质给出唯一态射 \(h: C \to P\) 使得 \(\pi_j \circ h = c_j\)。由锥条件 \(D(u) \circ c_j = c_k\) 推出 \(f \circ h = g \circ h\)，故 \(h\) 唯一地通过等化子 \(e\) 分解，即存在唯一 \(C \to L\)。这验证了 \(L\) 的极限泛性质。拉回作为二元积与等化子的特例自然包含于此构造。∎

\subsection{95.3 完备性}\label{ux5b8cux5907ux6027-1}

\textbf{定义 95.7}（完备范畴）：范畴 \(\mathcal{C}\) 称为\textbf{完备}的，如果它有所有小极限。\(\mathcal{C}\) 称为\textbf{余完备}的，如果它有所有小余极限。

完备性和余完备性是范畴''大小''的一个基本量度：一个范畴是否包含所有由小图构造的极限，决定了在多大程度上可以在该范畴内自由地进行构造而不必离开这个范畴。幸运的是，数学中最常用的范畴都具有这一性质。

\textbf{定理 95.3}：\textbf{Set} 是完备和余完备的。\textbf{Grp}、\textbf{Ab}、\textbf{Ring}、\textbf{R-Mod}、\textbf{Top} 都是完备和余完备的。

\textbf{证明}：仅证 \textbf{Set} 的完备性，余完备性对偶。由定理 95.2，只需验证 \textbf{Set} 有所有积和等化子。积为通常的笛卡尔积 \(\prod_j X_j = \{(x_j)_{j \in J} \mid x_j \in X_j\}\)，投影为 \(\pi_k((x_j)) = x_k\)。等化子为 \(E = \{x \in X \mid f(x) = g(x)\}\) 含入映射。对 \textbf{Grp}，积为直积，等化子为 \(\{x \in G \mid f(x) = g(x)\}\) 中取子群；\textbf{Ab} 和 \textbf{R-Mod} 类似。\textbf{Top} 中积为积拓扑，等化子为子空间拓扑。余积：\textbf{Set} 中为不交并，\textbf{Grp} 中为自由积，\textbf{Ab} 中为直和，\textbf{Top} 中为不交并拓扑。\(\blacksquare\)

定理 95.3 的证明体现了一个重要的范畴论方法论：验证一个范畴的完备性通常归结为验证积与等化子的存在性，而这两者往往已经在该范畴的理论中以其他名字出现。有了极限的存在性，下一步自然地是问：极限构造本身是否构成一个函子？答案是肯定的。

\textbf{命题 95.4}（极限的函子性）：若 \(\mathcal{C}\) 完备，则极限构造定义了函子 \(\varprojlim: [J, \mathcal{C}] \to \mathcal{C}\)。类似地，\(\varinjlim\) 是函子。

\textbf{证明}：对自然变换 \(\alpha: D \Rightarrow D'\)（\(D, D': J \to \mathcal{C}\)），需构造态射 \(\varprojlim \alpha: \varprojlim D \to \varprojlim D'\)。设 \(\varprojlim D\) 的锥为 \(\{\ell_j: \varprojlim D \to D(j)\}\)，\(\varprojlim D'\) 的锥为 \(\{\ell'_j: \varprojlim D' \to D'(j)\}\)。考虑锥 \(\{\alpha_j \circ \ell_j: \varprojlim D \to D'(j)\}_{j \in J}\)，由 \(\varprojlim D'\) 的万有性质，存在唯一态射 \(\varprojlim \alpha\) 使得 \(\ell'_j \circ \varprojlim \alpha = \alpha_j \circ \ell_j\)。验证函子性：\(\varprojlim (\operatorname{id}_D) = \operatorname{id}_{\varprojlim D}\)（唯一性），\(\varprojlim (\beta \circ \alpha) = \varprojlim \beta \circ \varprojlim \alpha\)（两端均满足与 \(\ell''_j\) 的交换条件，由万有锥唯一性得证）。\(\blacksquare\)

\textbf{定义 95.8}（保极限的函子）：函子 \(F: \mathcal{C} \to \mathcal{D}\) 称为\textbf{保极限}的，如果对每个图 \(D: J \to \mathcal{C}\)，\(F(\varprojlim D) \cong \varprojlim (F \circ D)\)（自然同构）。

\hrulefill

\hrulefill

\hrulefill

\hrulefill

\hrulefill

\hrulefill

\section{第292章：伴随函子}\label{ux7b2c292ux7ae0ux4f34ux968fux51fdux5b50}

伴随函子 \(F \dashv G\) 被 Mac Lane 誉为''范畴论中无处不在的概念''，其定义通过自然同构 \(\mathcal{D}(F(X), Y) \cong \mathcal{C}(X, G(Y))\) 将两个看似无关的函子偶合为一种深刻的对称关系：左伴随''自由地''构造结构，右伴随''遗忘式地''剥离结构------这正是数学中几乎所有''自由-遗忘''模式的范畴论表达。经典的伴随配对贯穿整个数学：自由群函子 \(\mathbf{Set} \to \mathbf{Grp}\) 是遗忘函子的左伴随，Abel 化函子是嵌入函子 \(\mathbf{Ab} \hookrightarrow \mathbf{Grp}\) 的左伴随，Stone-Cech 紧化是 Hausdorff 空间范畴嵌入紧空间范畴的反射（左伴随），而量化词 \(\exists\) 和 \(\forall\) 分别是弱化和代换函子的左伴随和右伴随。伴随函子理论还与极限（第 226 章）和单子（第 228 章）深度耦合------右伴随保持极限，左伴随保持余极限，而每个伴随对都内置地生成一个单子（左伴随与右伴随的合成）。

\subsection{96.1 伴随的定义}\label{ux4f34ux968fux7684ux5b9aux4e49}

\textbf{定义 96.1}（伴随函子）：设 \(F: \mathcal{C} \to \mathcal{D}\) 和 \(G: \mathcal{D} \to \mathcal{C}\) 是函子。\(F\) 称为 \(G\) 的\textbf{左伴随}（\(G\) 是 \(F\) 的\textbf{右伴随}），记作 \(F \dashv G\)，如果存在自然同构

\[\mathcal{D}(F(X), Y) \cong \mathcal{C}(X, G(Y))\]

对 \(X \in \mathcal{C}\)，\(Y \in \mathcal{D}\) 自然。

\textbf{定义 96.2}（单位与余单位）：伴随 \(F \dashv G\) 给出两个自然变换： - \textbf{单位}（unit）：\(\eta: \operatorname{id}_{\mathcal{C}} \Rightarrow G \circ F\)（由 \(X \mapsto F(X)\) 的恒等态射经过伴随同构得到） - \textbf{余单位}（counit）：\(\varepsilon: F \circ G \Rightarrow \operatorname{id}_{\mathcal{D}}\)（对偶地）

满足三角恒等式：\(G\varepsilon \circ \eta G = \operatorname{id}_G\) 和 \(\varepsilon F \circ F\eta = \operatorname{id}_F\)。

\subsection{96.2 伴随函子的性质}\label{ux4f34ux968fux51fdux5b50ux7684ux6027ux8d28}

\textbf{定理 96.1}（左伴随保余极限，右伴随保极限）：若 \(F \dashv G\)，则 \(F\) 保所有余极限，\(G\) 保所有极限。

\emph{证明}：对余极限 \(C = \varinjlim D_j\)，\(\mathcal{D}(F(C), Y) \cong \mathcal{C}(C, G(Y)) \cong \varprojlim \mathcal{C}(D_j, G(Y)) \cong \varprojlim \mathcal{D}(F(D_j), Y) \cong \mathcal{D}(\varinjlim F(D_j), Y)\)。由 Yoneda 引理，\(F(C) \cong \varinjlim F(D_j)\)。∎

\textbf{命题 96.2}：若 \(F \dashv G\) 且 \(F' \dashv G\)，则 \(F \cong F'\)（左伴随在同构意义下唯一）。类似地，右伴随也唯一。

\textbf{证明}：由 \(F \dashv G\) 和 \(F' \dashv G\)，有自然同构 \(\mathcal{D}(F(X), Y) \cong \mathcal{C}(X, G(Y)) \cong \mathcal{D}(F'(X), Y)\)。取 \(Y = F'(X)\)，\(\operatorname{id}_{F'(X)}\) 在第一个同构的逆下的像给出 \(\eta'_X: X \to G(F'(X))\)，再经第二个同构（取 \(Y = F'(X)\)）得到 \(\theta_X: F(X) \to F'(X)\)。由自然性，\(\theta\) 是自然变换。同理，取 \(Y = F(X)\) 构造 \(\theta'_X: F'(X) \to F(X)\)。由 Yoneda 引理，\(\theta'_X \circ \theta_X = \operatorname{id}_{F(X)}\) 且 \(\theta_X \circ \theta'_X = \operatorname{id}_{F'(X)}\)，故 \(\theta\) 是自然同构。\(\blacksquare\)

\textbf{定理 96.2}（一般伴随函子定理，Freyd）：设 \(\mathcal{C}\) 是完备范畴，\(G: \mathcal{D} \to \mathcal{C}\) 是函子。则 \(G\) 有左伴随当且仅当： 1. \(G\) 保所有小极限 2. 对每个 \(X \in \mathcal{C}\)，逗号范畴 \(X \downarrow G\) 有初始对象

（还需要''解集条件''（solution set condition）的适当版本）。

\textbf{证明}：充分性：\(G\) 保极限且满足解集条件 \(\Rightarrow\) 对每个 \(X \in \mathcal{C}\)，逗号范畴 \(X \downarrow G\) 有初始对象，该初始对象 \((S_X, \eta_X: X \to G(S_X))\) 定义左伴随 \(F(X) = S_X\)。构造核心：解集条件保证 \(X \downarrow G\) 中有一族''生成''对象 \(\{(M_i, f_i: X \to G(M_i))\}_{i \in I}\)；由于 \(G\) 保极限，这些 \(M_i\) 的积的某个子对象给出初始对象。必要性：左伴随 \(F \dashv G\) 自动保证 \(G\) 保极限（定理 96.1），且对每个 \(X\)，\((F(X), \eta_X)\) 是 \(X \downarrow G\) 的初始对象，从而解集条件取单点集 \(\{F(X)\}\) 即满足。∎

\textbf{定理 96.3}（特殊伴随函子定理）：设 \(\mathcal{C}\) 是完备、良幂（well-powered）且有余生成子（cogenerator）的范畴，\(\mathcal{D}\) 是局部小范畴。若 \(G: \mathcal{D} \to \mathcal{C}\) 保所有小极限，则 \(G\) 有左伴随。

\textbf{证明}：对每个 \(X \in \mathcal{C}\)，需在逗号范畴 \(X \downarrow G\) 中构造初始对象。令 \(\{s_i: X \to G(S_i)\}_{i \in I}\) 为 \(X \downarrow G\) 中所有对象的集合。构造积 \(P = \prod_{i \in I} G(S_i)\)，由完备性存在。考虑典范态射 \(h: X \to P\)。令 \(Q\) 为 \(\mathcal{C}\) 的余生成子，考虑所有单态射 \(m: M \hookrightarrow P\) 使得 \(h\) 通过 \(m\) 分解。由良幂性，这些单态射等价类构成集合，取它们的交 \(M_0\)（通过拉回构造）。由于 \(G\) 保极限，\(M_0 = G(S_0)\) 对某个 \(S_0 \in \mathcal{D}\)。由构造，\((S_0, \eta_X: X \to G(S_0))\) 是 \(X \downarrow G\) 的初始对象。定义 \(F(X) = S_0\)，则 \(F\) 是 \(G\) 的左伴随。\(\blacksquare\)

\subsection{96.3 伴随与等价}\label{ux4f34ux968fux4e0eux7b49ux4ef7}

\textbf{定义 96.3}（范畴等价）：函子 \(F: \mathcal{C} \to \mathcal{D}\) 称为\textbf{范畴等价}，如果存在函子 \(G: \mathcal{D} \to \mathcal{C}\) 使得 \(G \circ F \cong \operatorname{id}_{\mathcal{C}}\) 和 \(F \circ G \cong \operatorname{id}_{\mathcal{D}}\)（自然同构）。

\textbf{命题 96.3}：若 \(F \dashv G\) 且单位和余单位都是自然同构，则 \(F\) 和 \(G\) 是范畴等价（此时 \(F \dashv G\) 且 \(G \dashv F\)）。

\textbf{证明}：由 \(F \dashv G\)，有自然同构 \(\mathcal{D}(F(X), Y) \cong \mathcal{C}(X, G(Y))\)。余单位 \(\varepsilon: F \circ G \Rightarrow \operatorname{id}_{\mathcal{D}}\) 是自然同构意味着对每个 \(Y\)，\(\varepsilon_Y: F(G(Y)) \to Y\) 是同构。单位 \(\eta: \operatorname{id}_{\mathcal{C}} \Rightarrow G \circ F\) 是自然同构意味着对每个 \(X\)，\(\eta_X: X \to G(F(X))\) 是同构。于是 \(G \circ F \cong \operatorname{id}_{\mathcal{C}}\) 和 \(F \circ G \cong \operatorname{id}_{\mathcal{D}}\)，即 \(F\) 和 \(G\) 互为拟逆，构成范畴等价。此时 \(G \dashv F\) 也成立，因为 \(\mathcal{C}(G(Y), X) \cong \mathcal{C}(G(Y), G(F(X'))) \cong \mathcal{D}(F(G(Y)), F(X')) \cong \mathcal{D}(Y, F(X'))\)（其中 \(X \cong G(F(X'))\)）。\(\blacksquare\)

\textbf{定理 96.4}：函子 \(F: \mathcal{C} \to \mathcal{D}\) 是范畴等价当且仅当 \(F\) 是完全忠实的且本质满的（essentially surjective：每个 \(Y \in \mathcal{D}\) 同构于某个 \(F(X)\)）。

\emph{证明}：（\(\Rightarrow\)）若 \(F\) 等价，存在 \(G\) 使 \(G \circ F \cong \operatorname{id}_{\mathcal{C}}\)。自然同构给出每个 \(X\) 的同构 \(X \cong G(F(X))\)，故 \(F\) 完全忠实。本质满：对 \(Y \in \mathcal{D}\)，由 \(F \circ G \cong \operatorname{id}_{\mathcal{D}}\) 知 \(Y \cong F(G(Y))\)，即 \(Y\) 同构于 \(F\) 的像。（\(\Leftarrow\)）构造拟逆函子 \(G\)：对每个 \(Y \in \mathcal{D}\)，由本质满择一 \(X_Y \in \mathcal{C}\) 及同构 \(\varepsilon_Y: F(X_Y) \xrightarrow{\cong} Y\)。对态射 \(h: Y \to Z\)，由 \(F\) 全忠实，存在唯一 \(g: X_Y \to X_Z\) 使 \(F(g) = \varepsilon_Z^{-1} \circ h \circ \varepsilon_Y\)。定义 \(G(Y) = X_Y\)，\(G(h) = g\)。验证 \(G\) 是函子且 \(\varepsilon\) 给出 \(F \circ G \cong \operatorname{id}_{\mathcal{D}}\)；再利用 \(F\) 忠实构造 \(\eta: \operatorname{id}_{\mathcal{C}} \cong G \circ F\)。∎

\hrulefill

\hrulefill

\hrulefill

\hrulefill

\hrulefill

\hrulefill

\section{第293章：单子与代数}\label{ux7b2c293ux7ae0ux5355ux5b50ux4e0eux4ee3ux6570}

单子（monad）是伴随函子的代数抽象------从范畴论的视角看，每个伴随对 \(F \dashv G\) 合成出一个单子 \(T = G \circ F\)，正如群作用中的稳定子对应 Galois 理论中的中间域。但单子的意义远超此技术联系：它捕获了''代数结构''的普遍概念------将自由代数的生成元映射到其上的运算，如群在集合上的作用、环在 Abel 群上的作用、拓扑学中的闭包算子和逻辑学中的模态算子，均可统一为单子。这一理论在纯数学之外产生了深远影响：Moggi（1991 年）将单子引入计算机科学作为计算效果（副作用、异常、输入输出）的范畴模型，Haskell 编程语言中单子成为核心抽象。Eilenberg-Moore 代数和 Kleisli 范畴是单子的两种对偶表示：前者将代数结构的范畴紧致编码，后者将态射视为生成元的自由对象之间的射影，二者分别对应了''语法''（自由-遗忘）视角和''语义''（解释）视角。

\subsection{97.1 单子的定义}\label{ux5355ux5b50ux7684ux5b9aux4e49}

\textbf{定义 97.1}（单子）：范畴 \(\mathcal{C}\) 上的\textbf{单子}是一个三元组 \((T, \eta, \mu)\)，其中： - \(T: \mathcal{C} \to \mathcal{C}\) 是函子 - \(\eta: \operatorname{id}_{\mathcal{C}} \Rightarrow T\)（单位） - \(\mu: T^2 \Rightarrow T\)（乘法）

满足结合律：\(\mu \circ T\mu = \mu \circ \mu T\)（作为 \(T^3 \Rightarrow T\) 的自然变换），和单位律：\(\mu \circ \eta T = \operatorname{id}_T = \mu \circ T\eta\)。

\textbf{定理 97.1}（伴随给出单子）：每个伴随 \(F \dashv G\)（\(F: \mathcal{C} \to \mathcal{D}\)，\(G: \mathcal{D} \to \mathcal{C}\)）给出 \(\mathcal{C}\) 上的一个单子 \((T = G \circ F, \eta, \mu = G \varepsilon F)\)，其中 \(\eta\) 是伴随的单位，\(\varepsilon\) 是伴随的余单位。

\emph{证明}：验证单子的三个公理。\(\mu\) 的自然性由 \(\varepsilon\) 的自然性保证。结合律由 \(\varepsilon\) 的三角恒等式推出。∎

\subsection{97.2 Eilenberg-Moore 代数与 Kleisli 范畴}\label{eilenberg-moore-ux4ee3ux6570ux4e0e-kleisli-ux8303ux7574}

\textbf{定义 97.2}（Eilenberg-Moore 代数）：单子 \((T, \eta, \mu)\) 的一个\textbf{代数}（或 Eilenberg-Moore 代数）是一个对 \((A, a)\)，其中 \(A \in \mathcal{C}\)，\(a: T(A) \to A\) 是 \(\mathcal{C}\) 中的态射，满足 \(a \circ \eta_A = \operatorname{id}_A\) 和 \(a \circ \mu_A = a \circ T(a)\)。

代数的态射 \(f: (A, a) \to (B, b)\) 是 \(\mathcal{C}\) 中的态射 \(f: A \to B\) 满足 \(b \circ T(f) = f \circ a\)。代数及其态射构成 \textbf{Eilenberg-Moore 范畴} \(\mathcal{C}^T\)。

\textbf{定理 97.2}（Eilenberg-Moore 分解）：每个伴随 \(F \dashv G\) 生成的单子 \(T\)，函子 \(G\) 通过 \(\mathcal{C}^T\) 分解：存在比较函子 \(K: \mathcal{D} \to \mathcal{C}^T\) 使得 \(U^T \circ K = G\) 和 \(K \circ F = F^T\)（其中 \(U^T \dashv F^T\) 是 \(\mathcal{C}^T\) 的自由-遗忘伴随）。

\textbf{证明}：设 \(F: \mathcal{C} \to \mathcal{D}\)，\(G: \mathcal{D} \to \mathcal{C}\)，\(T = G \circ F\)。对每个 \(D \in \mathcal{D}\)，定义 \(K(D) = (G(D), G(\varepsilon_D))\)，其中 \(G(\varepsilon_D): T(G(D)) = G(F(G(D))) \to G(D)\)。验证代数公理：\(G(\varepsilon_D) \circ \eta_{G(D)} = \operatorname{id}_{G(D)}\) 由三角恒等式 \(G\varepsilon \circ \eta G = \operatorname{id}_G\) 保证；\(G(\varepsilon_D) \circ \mu_{G(D)} = G(\varepsilon_D) \circ G(\varepsilon_{F(G(D))}) = G(\varepsilon_D \circ \varepsilon_{F(G(D))}) = G(\varepsilon_D \circ F(G(\varepsilon_D))) = G(\varepsilon_D) \circ T(G(\varepsilon_D))\)，由 \(\varepsilon\) 的自然性得证。对态射 \(h: D \to D'\)，定义 \(K(h) = G(h)\)，验证 \(G(\varepsilon_{D'}) \circ T(G(h)) = G(h) \circ G(\varepsilon_D)\) 由 \(\varepsilon\) 的自然性保证。因此 \(K\) 是函子，且 \(U^T \circ K = G\) 和 \(K \circ F = F^T\) 直接验证。\(\blacksquare\)

\textbf{定义 97.3}（Kleisli 范畴）：单子 \(T\) 的 \textbf{Kleisli 范畴} \(\mathcal{C}_T\) 的对象与 \(\mathcal{C}\) 相同，但态射 \(\mathcal{C}_T(X, Y) = \mathcal{C}(X, T(Y))\)。合成由 \(T\) 的乘法定义。Kleisli 范畴是 Eilenberg-Moore 范畴的''自由代数''满子范畴。

\textbf{命题 97.4}：\(\mathcal{C}^T\) 是 \(T\) 的代数的''最大''范畴（终对象），\(\mathcal{C}_T\) 是 \(T\) 的代数的''最小''范畴（初对象），在适当的范畴意义下。

\textbf{证明}：给定单子 \((T, \eta, \mu)\)，定义 Kleisli 范畴 \(\mathcal{C}_T\)：对象 \(\operatorname{Ob}(\mathcal{C}_T) = \operatorname{Ob}(\mathcal{C})\)；态射 \(\mathcal{C}_T(X, Y) = \mathcal{C}(X, T(Y))\)；恒等态射 \(\operatorname{id}_X^{\mathcal{C}_T} = \eta_X: X \to T(X)\)；合成 \(g \circ_{\mathcal{C}_T} f = \mu_Z \circ T(g) \circ f\)（\(f: X \to T(Y)\)，\(g: Y \to T(Z)\)）。结合律：\((h \circ g) \circ f = \mu_W \circ T(\mu_W \circ T(h) \circ g) \circ f = \mu_W \circ \mu_{T(W)} \circ T^2(h) \circ T(g) \circ f = \mu_W \circ T(h) \circ \mu_Z \circ T(g) \circ f = h \circ (g \circ f)\)（使用 \(\mu\) 的结合律自然性）。存在自由函子 \(F_T: \mathcal{C} \to \mathcal{C}_T\)（\(F_T(X)=X\)，\(F_T(f)=\eta_Y \circ f\)）和遗忘函子 \(U_T: \mathcal{C}_T \to \mathcal{C}\)（\(U_T(X)=T(X)\)，\(U_T(f)=\mu_Y \circ T(f)\)），满足 \(F_T \dashv U_T\) 且生成的单子恰为 \(T\)。\(\mathcal{C}_T\) 等价于 \(\mathcal{C}^T\) 中由自由代数 \((\mu_X: T^2(X) \to T(X))\) 生成的满子范畴。对任意其他伴随 \(F \dashv G\) 生成 \(T\)，比较函子 \(K: \mathcal{D} \to \mathcal{C}^T\) 和 \(J: \mathcal{C}_T \to \mathcal{D}\) 给出 \(\mathcal{C}_T\) 到 \(\mathcal{D}\) 到 \(\mathcal{C}^T\) 的典范分解，\(\mathcal{C}_T\) 为初对象，\(\mathcal{C}^T\) 为终对象。\(\blacksquare\)

\subsection{97.3 单子与代数理论}\label{ux5355ux5b50ux4e0eux4ee3ux6570ux7406ux8bba}

\textbf{定义 97.4}（代数理论）：单子 \((T, \eta, \mu)\) 在 \textbf{Set} 上的代数 \((A, a)\) 可以看作带有满足某些等式的运算的集合。例如： - 群单子：\(T(X)\) 是 \(X\) 生成的自由群，代数 \((A, a)\) 是群（\(a: T(A) \to A\) 是群的乘法、逆和单位元） - 环单子：类似地，代数环

\textbf{定理 97.5}（Beck 单子性定理）：设 \(G: \mathcal{D} \to \mathcal{C}\) 有左伴随 \(F\)。则比较函子 \(K: \mathcal{D} \to \mathcal{C}^T\) 是范畴等价当且仅当 \(G\) 反映同构且 \(\mathcal{D}\) 有余等化子，且 \(G\) 保这些余等化子。

\textbf{证明}：（\(\Rightarrow\)）若 \(K\) 是等价，则 \(G = U^T \circ K\)。\(U^T\) 反映同构（遗忘函子对代数结构反映同构），且 \(U^T\) 保余等化子（在代数范畴中构造）。因 \(K\) 是等价，\(G\) 也反映同构且保余等化子。（\(\Leftarrow\)）假设 \(G\) 反映同构且保 \(\mathcal{D}\) 中形如 \(F \circ G\)-分裂的余等化子对。对每个 \(T\)-代数 \((A, a)\)，构造其在 \(\mathcal{D}\) 中的原像：取 \(F(A)\) 的余等化子 \(F(G(F(A))) \rightrightarrows F(A)\)（两态射为 \(F(a)\) 和 \(\varepsilon_{F(A)}\)）。由 \(G\) 保此余等化子，其像恰为 \((A, a)\)。由 \(G\) 反映同构，\(K\) 本质满且完全忠实，故为等价。\(\blacksquare\)

Beck 单子性定理提供了判断一个函子是否为''遗忘函子''的精确标准。

\hrulefill

\hrulefill

\hrulefill

\hrulefill

\hrulefill

\hrulefill

\section{第294章：Abel 范畴与导出范畴初步}\label{ux7b2c294ux7ae0abel-ux8303ux7574ux4e0eux5bfcux51faux8303ux7574ux521dux6b65}

Abel 范畴是能进行同调代数的范畴框架。导出范畴是 Abel 范畴的局部化，使得拟同构成为真正的同构，是现代同调代数的核心语言。

\subsection{98.1 Abel 范畴}\label{abel-ux8303ux7574}

\textbf{定义 98.1}（加性范畴）：范畴 \(\mathcal{A}\) 称为\textbf{加性范畴}（或预加性范畴），如果每个 \(\mathcal{A}(X, Y)\) 是 Abel 群，且合成是双线性的。此外，\(\mathcal{A}\) 有零对象（既是始对象又是终对象）和有限双积（既是积又是余积）。

\textbf{定义 98.2}（Abel 范畴）：加性范畴 \(\mathcal{A}\) 称为 \textbf{Abel 范畴}，如果 1. 每个态射有核和余核 2. 每个单态射是其余核的核，每个满态射是其核的余核 3. （等价地）每个态射有典范分解 \(f = m \circ e\)，其中 \(e\) 是余核的满态射，\(m\) 是核的单态射

\textbf{定义 98.3}（正合列）：Abel 范畴中的序列 \(\cdots \to A \xrightarrow{f} B \xrightarrow{g} C \to \cdots\) 称为\textbf{正合的}，如果 \(\ker g = \operatorname{im} f\)（其中 \(\operatorname{im} f = \ker(\operatorname{coker} f)\)）。

\textbf{定理 98.1}（Freyd-Mitchell 嵌入定理）：每个小 Abel 范畴 \(\mathcal{A}\) 可以完全忠实地嵌入到某个环 \(R\) 的模范畴 \textbf{R-Mod} 中，且嵌入是正合的。

\emph{说明}：Freyd-Mitchell 嵌入定理意味着在 Abel 范畴中做同调代数时，可以''假装''在模范畴中工作，使用元素论证（diagram chasing）。

\textbf{证明}：将小 Abel 范畴 \(\mathcal{A}\) 视为加性范畴，构造函子 \(H: \mathcal{A}^{\text{op}} \to \mathbf{Ab}\) 的适当满子范畴作为''层''，再取归纳极限构造环 \(R\)，使 \(\mathcal{A}\) 等价于 \(R\)-Mod 中紧生成对象的满子范畴。核心技术是利用 \(\mathcal{A}\) 中所有可表函子构成的谱（spectrum）来''编码''模范畴结构，最终得到完全忠实且正合的嵌入 \(\mathcal{A} \hookrightarrow R\text{-Mod}\)。∎

\subsection{98.2 复形与同调}\label{ux590dux5f62ux4e0eux540cux8c03}

\textbf{定义 98.4}（链复形）：Abel 范畴 \(\mathcal{A}\) 中的\textbf{链复形}是一列对象 \(A_n\) 和态射 \(d_n: A_n \to A_{n-1}\) 满足 \(d_{n-1} \circ d_n = 0\)。\textbf{上链复形}中态射方向为 \(d^n: A^n \to A^{n+1}\)。

链复形是现代同调代数的核心数据结构。直观上，链复形 \(A_\bullet\) 记录了''边界''信息：\(d_n\) 将 \(n\) 维对象映射到其 \((n-1)\) 维边界，而 \(d_{n-1} \circ d_n = 0\) 表达了''边界的边界为零''这一基本事实。这一条件保证 \(\operatorname{im} d_{n+1} \subseteq \ker d_n\)，从而可以定义同调群。链复形范畴 \(\mathbf{Ch}(\mathcal{A})\) 本身是一个 Abel 范畴，其态射为与微分交换的链映射 \(f_n: A_n \to B_n\)（满足 \(d_n^B \circ f_n = f_{n-1} \circ d_n^A\)）。链复形之间的同伦 \(h: f \Rightarrow g\)（由一族态射 \(h_n: A_n \to B_{n+1}\) 满足 \(f_n - g_n = d_{n+1}^B \circ h_n + h_{n-1} \circ d_n^A\)）提供了比链映射更细的等价关系，是构造同伦范畴 \(K(\mathcal{A})\) 的基础。这一概念源于代数拓扑中单纯复形和 CW 复形的链群结构，其代数抽象由 Eilenberg 和 Mac Lane 在 1940 年代完成。

\textbf{定义 98.5}（同调）：链复形 \(A_\bullet\) 的 \(n\) 次同调为 \(H_n(A_\bullet) = \ker d_n / \operatorname{im} d_{n+1}\)。

同调群 \(H_n(A_\bullet)\) 度量了链复形在 \(n\) 次处偏离正合性的程度：若 \(H_n(A_\bullet) = 0\)，则复形在 \(n\) 次处是正合的，即 \(\ker d_n = \operatorname{im} d_{n+1}\)，意味着 \(n\) 维''闭链''恰好是 \((n+1)\) 维''边界''。同调函子 \(H_n: \mathbf{Ch}(\mathcal{A}) \to \mathcal{A}\) 的构造是加性且保有限直和的，它将链映射 \(f: A_\bullet \to B_\bullet\) 诱导为同调群之间的同态 \(H_n(f): H_n(A_\bullet) \to H_n(B_\bullet)\)。同调的核心性质是\textbf{同伦不变性}：链同伦的链映射诱导相同的同调同态，这一事实是导出范畴理论中局部化的基石。在代数拓扑中，同调群对应空间的''洞''的计数（\(H_0\) 计数连通分支，\(H_1\) 计数一维洞，等等）；在代数中，Tor 和 Ext 函子通过投射/内射分解的同调群来定义。

\textbf{定义 98.6}（拟同构）：链复形态射 \(f: A_\bullet \to B_\bullet\) 称为\textbf{拟同构}，如果对所有 \(n\)，它诱导的同调同构 \(H_n(A_\bullet) \cong H_n(B_\bullet)\)。

拟同构是链复形之间比同伦等价更弱的等价关系，却是导出范畴理论的核心概念。一个拟同构 \(f\) 可能不是链同伦等价（即不可逆于同伦意义下），但在导出范畴 \(D(\mathcal{A})\) 中它被形式地取逆，成为真正的同构。这一操作类似于交换代数中局部化将非零因子变为可逆元：\(D(\mathcal{A}) = K(\mathcal{A})[\mathcal{Q}^{-1}]\) 正是同伦范畴关于拟同构类 \(\mathcal{Q}\) 的 Verdier 局部化。拟同构的典型例子包括：投射分解 \(P_\bullet \to A\)（其中 \(A\) 视为集中在度 0 的复形）和内射分解 \(A \to I^\bullet\)。这些分解的存在性依赖于 Abel 范畴有足够多投射/内射对象的假设，它们在导出函子（如 \(\operatorname{Tor}\)、\(\operatorname{Ext}\)）的定义中扮演关键角色。

\subsection{98.3 导出范畴}\label{ux5bfcux51faux8303ux7574}

\textbf{定义 98.7}（同伦范畴）：Abel 范畴 \(\mathcal{A}\) 的链复形范畴 \(\mathbf{Ch}(\mathcal{A})\) 的\textbf{同伦范畴} \(K(\mathcal{A})\) 的对象与 \(\mathbf{Ch}(\mathcal{A})\) 相同，态射是链映射的同伦类。

\textbf{定义 98.8}（导出范畴）：Abel 范畴 \(\mathcal{A}\) 的\textbf{导出范畴} \(D(\mathcal{A})\) 是 \(K(\mathcal{A})\) 关于拟同构的局部化（Verdier 局部化）。即 \(D(\mathcal{A}) = K(\mathcal{A})[\mathcal{Q}^{-1}]\)，其中 \(\mathcal{Q}\) 是所有拟同构的类。

在 \(D(\mathcal{A})\) 中，拟同构变成了真正的同构。这是导出函子理论（V12 Ch 118）的范畴论基础。

导出范畴的构造动机来自于同调代数中的一个基本问题：链复形的同伦范畴中，拟同构未必是同伦等价，但同调性质告诉我们拟同构的复形在代数意义上是''相同的''。Verdier 局部化正是将拟同构形式地取逆，使得导出范畴成为研究同调现象的自然工作环境。导出范畴的另一个关键特征是它自然地继承了三角范畴结构，这一结构编码了代数拓扑中映射锥与长正合列的灵魂。

\textbf{定理 98.2}（导出范畴是三角范畴）：导出范畴 \(D(\mathcal{A})\) 具有\textbf{三角范畴}结构：存在平移函子 \([1]: D(\mathcal{A}) \to D(\mathcal{A})\)（\(A[1]_n = A_{n-1}\)）和一类\textbf{好三角}（distinguished triangles）\(A \to B \to C \to A[1]\)，满足 Verdier 公理。

\textbf{证明}：平移函子 \([1]\) 由 \(K(\mathcal{A})\) 继承，且拟同构族 \(\mathcal{Q}\) 与平移相容，故 \([1]\) 下降至 \(D(\mathcal{A})\)。好三角定义为 \(K(\mathcal{A})\) 中映射锥序列的同构类：对 \(f: A \to B\)，定义 \(C(f)_n = A_{n-1} \oplus B_n\) 及微分，好三角为 \(A \xrightarrow{f} B \to C(f) \to A[1]\) 同构于此类者。验证 Verdier 公理：（TR1）恒等态射 \(A \xrightarrow{\operatorname{id}} A \to 0 \to A[1]\) 是好三角，且任意态射可嵌入好三角，好三角同构于好三角；（TR2）\(A \to B \to C \to A[1]\) 是好三角当且仅当 \(B \to C \to A[1] \to B[1]\) 是好三角（旋转）；（TR3）给定两个好三角及态射 \(A \to A'\)、\(B \to B'\) 使前两方交换，存在 \(C \to C'\) 使整个图交换（填充）；（TR4）八面体公理由映射锥的八面体恒等式验证。\(\blacksquare\)

\textbf{定义 98.9}（导出函子）：设 \(F: \mathcal{A} \to \mathcal{B}\) 是左正合函子。\(F\) 的\textbf{全导出函子} \(\mathbf{R}F: D^+(\mathcal{A}) \to D^+(\mathcal{B})\) 定义为 \(\mathbf{R}F(A_\bullet) = F(I_\bullet)\)，其中 \(A_\bullet \to I_\bullet\) 是内射分解（或 \(K\)-内射分解）。\(R^i F(A) = H^i(\mathbf{R}F(A))\) 是通常的导出函子。

\textbf{定理 98.3}（导出函子的存在性）：若 \(\mathcal{A}\) 有足够多内射对象，则每个左正合函子 \(F: \mathcal{A} \to \mathcal{B}\) 有全导出函子 \(\mathbf{R}F: D^+(\mathcal{A}) \to D^+(\mathcal{B})\)。

\textbf{证明}：由足够多内射条件，对每个下有界复形 \(A^\bullet \in D^+(\mathcal{A})\)，存在内射分解 \(A^\bullet \xrightarrow{\sim} I^\bullet\)（拟同构，\(I^\bullet\) 为下有界内射复形）。定义 \(\mathbf{R}F(A^\bullet) = F(I^\bullet)\)。需验证此定义不依赖于内射分解的选取：若 \(A^\bullet \xrightarrow{\sim} I^\bullet\) 和 \(A^\bullet \xrightarrow{\sim} J^\bullet\) 为两个内射分解，则存在内射复形 \(K^\bullet\) 及拟同构 \(I^\bullet \to K^\bullet \leftarrow J^\bullet\)（比较引理）。因 \(F\) 左正合，拟同构的内射复形在 \(F\) 下的像为链同伦等价（内射复形上 \(F\) 保拟同构），故 \(F(I^\bullet) \cong F(J^\bullet)\) 在 \(D^+(\mathcal{B})\) 中。函子性由诱导映射的万有性质保证，\(\mathbf{R}F\) 是良定义的三角函子。\(\blacksquare\)

\hrulefill

\hrulefill

\hrulefill

\hrulefill

\hrulefill

\hrulefill

\newpage

\chapter{08-代数几何与数论/01-代数几何/07-高阶范畴论与导出几何.md}\label{ux4ee3ux6570ux51e0ux4f55ux4e0eux6570ux8bba01-ux4ee3ux6570ux51e0ux4f5507-ux9ad8ux9636ux8303ux7574ux8bbaux4e0eux5bfcux51faux51e0ux4f55.md}

\subsection{导出代数几何导引}\label{ux5bfcux51faux4ee3ux6570ux51e0ux4f55ux5bfcux5f15}

导出代数几何（Derived Algebraic Geometry, DAG）是 21 世纪代数几何最重要的范式转变之一，由 Toen-Vezzosi 和 Lurie 在 2000 年代系统建立。其核心思想是：经典概形理论中交换环的范畴应替换为单纯环或微分分次环（或更一般地，\(E_\infty\)-环谱）的 \(\infty\)-范畴，从而将''商''操作从集合层面提升到同伦层面。这一转变解决了经典代数几何中的三个结构性缺陷：非横截相交理论的 Tor 公式缺乏几何解释、形变理论中的余切复形无法在概形范畴内自然定义、以及模空间的导出障碍需要虚基本类来补偿。导出代数几何与 Lurie 的谱代数几何（Spectral Algebraic Geometry）紧密相关，后者进一步将结构层替换为 \(E_\infty\)-环谱，在整数环和正特征情形下具有比 dg 代数更丰富的结构。导出代数几何为几何 Langlands 纲领（Arinkin-Gaitsgory）、Donaldson-Thomas 理论（导出模空间）和拓扑量子场论（导出叠的配边假设）提供了核心技术框架。

\subsection{经典代数几何的三大限制}\label{ux7ecfux5178ux4ee3ux6570ux51e0ux4f55ux7684ux4e09ux5927ux9650ux5236}

导出代数几何（Derived Algebraic Geometry, DAG）的起源可归结为经典概形论中的三类结构性缺陷，它们均源自交换环层面缺乏同伦信息。

\textbf{定义 1}（非横截交）：设 \(Y, Z \subset X\) 为光滑代数簇中的闭子簇。称 \(Y\) 与 \(Z\) 在点 \(x \in Y \cap Z\) 处\textbf{横截相交}，若 \[T_x Y + T_x Z = T_x X.\] 在非横截情形，经典交积 \(Y \times_X Z\) 的维数严格大于期望值 \(\dim Y + \dim Z - \dim X\)，且其概形结构携带幂零元，丢失了高阶相交信息。

\textbf{限制一（交论缺陷）}：Serre 的交重数公式 \[\chi(Y, Z; X) = \sum_{i \geq 0} (-1)^i \operatorname{length}_{\mathcal{O}_{X,x}} \operatorname{Tor}_i^{\mathcal{O}_{X,x}}(\mathcal{O}_{Y,x}, \mathcal{O}_{Z,x})\] 表明经典相交理论正确的数值结果依赖于 Tor 项的正负交错求和，但其几何意义------即这些 Tor 群如何从实际的导出纤维积中产生------在经典框架内不可见。DAG 将 Tor 公式提升为同伦纤维积的层上同调。

\textbf{限制二（形变空间的奇异性）}：概形态射 \(f: X \to Y\) 的形变理论由 Illusie 的余切复形 \(\mathbb{L}_{X/Y}\) 控制。在经典概形范畴内，\(\mathbb{L}_{X/Y}\) 仅作为导出范畴中的对象存在，无法由底层概形的层构造。DAG 中 \(\mathbb{L}\) 自然地成为结构层 \(\mathcal{O}_X^\bullet\) 的 Kähler 微分层。

\textbf{限制三（模空间的导出障碍）}：模问题（如向量丛、主 G-丛、Higgs 丛的模空间）的形变与障碍理论自然产生取值于导出范畴的上同调群。经典叠论以 2-范畴截断处理这部分信息，导致模空间的''正确''维数需要虚基本类来补偿。DAG 中的导出叠（derived stack）原生携带障碍复形。

\subsection{单纯环与导出环}\label{ux5355ux7eafux73afux4e0eux5bfcux51faux73af}

\textbf{定义 2}（单纯交换环）：一个\textbf{单纯交换环} \(A_\bullet\) 是交换环范畴 \(\mathbf{CRing}\) 中的一个单纯对象，即函子 \(\Delta^{\mathrm{op}} \to \mathbf{CRing}\)。等价地，\(A_\bullet\) 由一族交换环 \(\{A_n\}_{n \geq 0}\) 和面映射 \(d_i: A_n \to A_{n-1}\)、退化映射 \(s_i: A_n \to A_{n+1}\) 组成，满足单纯恒等式，且所有结构映射为环同态。

\(A_\bullet\) 的同伦群 \(\pi_n(A_\bullet)\) 定义为相应 Moore 复形 \(N_*(A_\bullet)\)（其中 \(N_n = \bigcap_{i=0}^{n-1} \ker d_i\)，微分 \(d = \sum_{i=0}^n (-1)^i d_i\)）的同调群。\(\pi_0(A_\bullet)\) 自然带有交换环结构。

单纯环的定义将经典的交换环概念提升到了同伦层面，其代数本质通过以下经典的 Dold-Kan 对应得以揭示：单纯环的数据等价于非负微分分次代数。这一对应是连接同伦代数和同调代数的核心桥梁。

\textbf{定理 1}（Dold-Kan 对应）：单纯交换环范畴与链复形 \((C_*, d)\)（其中 \(C_n = 0\) 对 \(n < 0\)）构成的非负微分分次交换代数范畴之间存在一个等价，使得 \(\pi_n(A_\bullet) \cong H_n(C_*)\)。

\textbf{证明}：Dold-Kan 对应在 Abel 群层次上给出单纯 Abel 群 \(\mathbf{sAb}\) 与链复形 \(\mathbf{Ch}_{\geq 0}\) 的等价，由规范化复形函子 \(N: \mathbf{sAb} \to \mathbf{Ch}_{\geq 0}\) 及其逆 \(\Gamma\) 实现。对于单纯环 \(A_\bullet\)，其底层单纯 Abel 群经 \(N\) 得到链复形。环的乘法 \(A_\bullet \times A_\bullet \to A_\bullet\) 通过 Eilenberg-Zilber 映射（shuffle 积）诱导 \(N_*(A_\bullet) \otimes N_*(A_\bullet) \to N_*(A_\bullet \otimes A_\bullet) \to N_*(A_\bullet)\)，赋予 \(N_*(A_\bullet)\) 微分分次代数结构。逆方向，\(\Gamma\) 将 dg 代数的乘法提升为单纯环乘法。\(\blacksquare\)

\textbf{定义 3}（\(E_\infty\)-环谱）：一个\textbf{\(E_\infty\)-环谱}是指对称谱（或正交谱）范畴中的一个代数对象，其乘法在最高同伦意义上交换且结合。具体而言，\(E_\infty\)-operad 在谱范畴上的代数即为 \(E_\infty\)-环谱。其同伦范畴中的对象等价于具有高度结构乘法的环谱。

Dold-Kan 对应下的 dg 代数与 \(E_\infty\)-环谱的关系为：特征零域上的连通 dg 代数可典范地视为 \(E_\infty\)-代数（形式性结果），而正特征情形 \(E_\infty\)-结构严格强于 dg 结构，需通过 Goerss-Hopkins 障碍理论提升。

掌握了单纯环、dg 代数和 \(E_\infty\)-环谱这些导出环的概念之后，便可以正式定义导出概形------它将经典概形的结构层替换为一个携带同伦信息的导出结构层，使得概形的''函数代数''在每一茎上都是一个导出环。

\subsection{导出概形的定义}\label{ux5bfcux51faux6982ux5f62ux7684ux5b9aux4e49}

\textbf{定义 4}（导出概形）：一个\textbf{导出概形}是偶对 \((X, \mathcal{O}_X^\bullet)\)，其中 \(X\) 为拓扑空间，\(\mathcal{O}_X^\bullet\) 为 \(X\) 上的 \(E_\infty\)-环谱（或 cdga，取决于模型）层，满足：

\begin{enumerate}
\def\labelenumi{\arabic{enumi}.}
\tightlist
\item
  \((X, \pi_0(\mathcal{O}_X^\bullet))\) 是经典概形；
\item
  对每个 \(n \geq 0\)，同伦层 \(\pi_n(\mathcal{O}_X^\bullet)\) 是 \(\pi_0(\mathcal{O}_X^\bullet)\)-模的拟凝聚层。
\end{enumerate}

取 \(H^0(X, \pi_0(\mathcal{O}_X^\bullet))\) 即恢复底层经典概形的整体截面环。概形范畴 \(\mathbf{Sch}\) 通过平凡层 \(\mathcal{O}_X\)（视为集中在度 0 的 cdga）全忠实嵌入导出概形范畴。

\textbf{定理 2}（截断函子）：存在一个截断函子 \(t_0: \mathbf{dSch} \to \mathbf{Sch}\)，将导出概形 \((X, \mathcal{O}_X^\bullet)\) 映为其底层经典概形 \((X, \pi_0(\mathcal{O}_X^\bullet))\)。\(t_0\) 是包含函子 \(\mathbf{Sch} \hookrightarrow \mathbf{dSch}\) 的左伴随。

\textbf{证明}：对任意导出概形 \((X, \mathcal{O}_X^\bullet)\)，\(\pi_0\) 操作定义了一个逐茎截断：\((\pi_0(\mathcal{O}_X^\bullet))_x = \pi_0(\mathcal{O}_{X,x}^\bullet)\)。由于 \(\pi_0\) 保持环结构的交换性且命题（2）确保拟凝聚性，\((X, \pi_0(\mathcal{O}_X^\bullet))\) 满足概形公理。伴随性由万有性质得出：从经典概形 \((Y, \mathcal{O}_Y)\) 到 \(t_0(X, \mathcal{O}_X^\bullet)\) 的态射一一对应于从 \((Y, \mathcal{O}_Y)\)（视为常值 cdga）到 \((X, \mathcal{O}_X^\bullet)\) 的导出概形态射，因为任意 cdga 映射 \(A^\bullet \to B\)（其中 \(B\) 集中在度 0）典范地通过 \(A^\bullet \to \pi_0(A^\bullet)\) 分解。\(\blacksquare\)

\subsection{余切复形的正式定义}\label{ux4f59ux5207ux590dux5f62ux7684ux6b63ux5f0fux5b9aux4e49}

\textbf{定义 5}（André-Quillen 上同调与余切复形）：设 \(f: A \to B\) 为交换环同态。\(B\) 相对于 \(A\) 的\textbf{余切复形} \(\mathbb{L}_{B/A}\) 是单纯 \(B\)-模范畴中的对象（或在 Dold-Kan 对应下，视作 \(B\)-模的链复形），定义为 \(B\) 在 \(A\) 上的单纯解消 \(\Omega^1_{P_\bullet/A} \otimes_{P_\bullet} B\)，其中 \(P_\bullet \to B\) 是 \(B\) 在 \(A\)-代数范畴中的单纯自由解消。对任意 \(B\)-模 \(M\)，André-Quillen 上同调为 \[D^i(B/A, M) = H^i(\operatorname{Hom}_B(\mathbb{L}_{B/A}, M)).\]

余切复形 \(\mathbb{L}_{B/A}\) 是 Kähler 微分模 \(\Omega^1_{B/A}\) 在导出范畴中的精确替代。在经典代数几何中，\(\Omega^1_{B/A}\) 的函子性在非光滑态射下表现不佳：若 \(f\) 不是光滑态射，则 \(\Omega^1_{B/A}\) 可能产生非零的高阶上同调障碍。余切复形通过将 \(\Omega^1\) 替换为导出范畴中的对象，自动编码了这些高阶信息。具体而言，\(H^0(\mathbb{L}_{B/A}) \cong \Omega^1_{B/A}\)，而负上同调群 \(H^{-i}(\mathbb{L}_{B/A})\)（\(i > 0\)）度量了 \(B\) 作为 \(A\)-代数的''光滑性亏损''。这一构造的历史可追溯到 André 和 Quillen 在 1960-70 年代独立发展出的上同调理论，其动机来自对局部完全交环的形变问题的研究。在导出代数几何中，余切复形被提升为概形的基本不变量：对于导出概形 \(\mathcal{X}\)，其（绝对）余切复形 \(\mathbb{L}_{\mathcal{X}}\) 是 \(\mathcal{O}_{\mathcal{X}}\)-模的拟凝聚复形，控制了 \(\mathcal{X}\) 的无穷小形变理论。

\textbf{定理 3}（余切复形与光滑态射）：概形态射 \(f: X \to Y\) 光滑当且仅当 \(\mathbb{L}_{X/Y}\) 是 \(\mathcal{O}_X\)-模的完美复形，集中在同调度 0（即拟同构于局部自由有限秩的 \(\mathcal{O}_X\)-模）。

\textbf{证明}：若 \(f\) 光滑，则局部上 \(X\) 在 \(Y\) 上由多项式代数平展定义。多项式代数 \(Y[T_1,\ldots,T_n]\) 的余切复形即为 Kähler 微分层 \(\Omega^1\)，集中在度 0 且局部自由。平展态射的余切复形零调，从而 \(\mathbb{L}_{X/Y} \cong \Omega^1_{X/Y}[0]\)，为完美复形。反之，若 \(\mathbb{L}_{X/Y}\) 为完美复形集中在度 0，其同调维数为 0 蕴含 \(f\) 形式光滑（形变无障碍），完美性蕴含局部有限表示，从而 \(f\) 光滑。\(\blacksquare\)

余切复形的一个重要应用是\textbf{形变-障碍理论}：对于概形 \(X_0\) 在 \(Y\) 上的无穷小形变，其形变空间由 \(\operatorname{Ext}^1(\mathbb{L}_{X_0/Y}, \mathcal{O}_{X_0})\) 参数化，障碍空间位于 \(\operatorname{Ext}^2(\mathbb{L}_{X_0/Y}, \mathcal{O}_{X_0})\) 中。更一般地，Illusie 在 SGA6 和其专著《Complexe Cotangent et Déformations》中系统发展了余切复形的理论，将其推广到任意概形态射，并证明了对任意合成 \(X \xrightarrow{f} Y \xrightarrow{g} Z\)，存在导出范畴中的正合三角形 \[\mathbb{L}_{X/Y} \to \mathbb{L}_{X/Z} \to f^*\mathbb{L}_{Y/Z} \xrightarrow{+1},\] 这一性质是 Kähler 微分的基本正合列在导出范畴中的提升，也是导出代数几何中六函子形式体系的基础之一。

\subsection{导出形变理论}\label{ux5bfcux51faux5f62ux53d8ux7406ux8bba}

\textbf{定义 6}（dg-Lie 代数）：域 \(k\) 上的\textbf{微分分次 Lie 代数}（dgla）是微分分次 \(k\)-向量空间 \((\mathfrak{g}^\bullet, d)\)，配备双线性括号 \([-,-]: \mathfrak{g}^p \otimes \mathfrak{g}^q \to \mathfrak{g}^{p+q}\)，满足：

\begin{enumerate}
\def\labelenumi{\arabic{enumi}.}
\tightlist
\item
  （分次反对称）\([x,y] + (-1)^{|x||y|}[y,x] = 0\)；
\item
  （分次 Jacobi）\((-1)^{|x||z|}[x,[y,z]] + (-1)^{|y||x|}[y,[z,x]] + (-1)^{|z||y|}[z,[x,y]] = 0\)；
\item
  （Leibniz）\(d[x,y] = [dx,y] + (-1)^{|x|}[x,dy]\)。
\end{enumerate}

dgla \(\mathfrak{g}\) 的 Maurer-Cartan 元素集合为 \(\operatorname{MC}(\mathfrak{g}) = \{x \in \mathfrak{g}^1 \mid dx + \frac{1}{2}[x,x] = 0\}\)。

\textbf{定理 4}（Lurie-Pridham 等价）：特征零域 \(k\) 上，Artin 局部 \(k\)-代数增大范畴上的形式模问题（formal moduli problem）的 \(\infty\)-范畴等价于微分分次 Lie 代数的 \(\infty\)-范畴。等价由以下构造实现：dgla \(\mathfrak{g}\) 对应形式模问题 \[R \mapsto \operatorname{MC}(\mathfrak{g} \otimes_k \mathfrak{m}_R),\] 其中 \(\mathfrak{m}_R\) 为 \(R\) 的极大理想（赋予平凡 Lie 括号）；反之，形式模问题 \(F\) 对应其切复形（shifted by \(-1\)）所带的典范 \(L_\infty\)-结构，由 \(F\) 的 Koszul 对偶确定。

\textbf{\emph{证明}}：取 dgla \(\mathfrak{g}\)，对每个 Artin 局部代数 \(R\)，\(\mathfrak{g} \otimes \mathfrak{m}_R\) 为幂零 dgla，MC 元素的同伦类构成 \(\infty\)-广群。函子 \(R \mapsto \operatorname{MC}(\mathfrak{g} \otimes \mathfrak{m}_R)\) 满足形式模问题的 Schlessinger 条件（同伦版本）。逆方向的构造：形式模问题 \(F\) 的切空间 \(T_F = F(k[\varepsilon]/(\varepsilon^2))\) 通过 \(F\) 的高阶无穷小形变自然地承载 \(L_\infty\)-结构，在特征零情形此 \(L_\infty\)-结构形式化（由 \(L_\infty\) 形式性定理）为 dgla。\(\blacksquare\)

\subsection{导出交理论}\label{ux5bfcux51faux4ea4ux7406ux8bba}

\textbf{定义 7}（导出纤维积）：设 \(X \to S \leftarrow Y\) 为导出概形之间的态射图。其\textbf{导出纤维积} \(X \times_S^h Y\) 的函子定义为：对任意导出概形 \(T\)， \[(X \times_S^h Y)(T) = X(T) \times_{S(T)}^h Y(T),\] 其中右侧为 \(\infty\)-广群范畴中的同伦纤维积。当 \(X, Y, S\) 为经典概形（视为常值导出概形）时，\(X \times_S^h Y\) 的结构层由导出张量积 \(\mathcal{O}_X \otimes_{\mathcal{O}_S}^{\mathbb{L}} \mathcal{O}_Y\) 给出。

\textbf{定理 5}（导出 Serre 交重数公式）：设 \(X\) 为光滑代数簇，\(Y, Z \subset X\) 为闭子概形，交 \(W = Y \cap Z\) 为有限个点 \(\{w_1,\ldots,w_r\}\)。在导出纤维积 \(Y \times_X^h Z\) 中，\(w_i\) 处的导出交重数为 \[\chi_i = \sum_{j \geq 0} (-1)^j \dim_k H^{-j}(\mathcal{O}_{Y,w_i} \otimes_{\mathcal{O}_{X,w_i}}^{\mathbb{L}} \mathcal{O}_{Z,w_i}),\] 且 \(\chi_i\) 恰好等于 Serre 交重数 \[i(Y,Z; X) = \sum_{j \geq 0} (-1)^j \operatorname{length}_{\mathcal{O}_{X,w_i}} \operatorname{Tor}_j^{\mathcal{O}_{X,w_i}}(\mathcal{O}_{Y,w_i}, \mathcal{O}_{Z,w_i}).\]

\textbf{证明}：由定义，导出交 \(\mathcal{O}_{Y \times_X^h Z}\) 的整体截面为 \(\mathcal{O}_Y \otimes_{\mathcal{O}_X}^{\mathbb{L}} \mathcal{O}_Z\)。在茎处，导出张量积的同调群恰为 Tor 群：\(H^{-j}(\mathcal{O}_{Y,w} \otimes_{\mathcal{O}_{X,w}}^{\mathbb{L}} \mathcal{O}_{Z,w}) \cong \operatorname{Tor}_j^{\mathcal{O}_{X,w}}(\mathcal{O}_{Y,w}, \mathcal{O}_{Z,w})\)。因此导出交重数（各同调度维数的交变和）与 Serre 公式一致。DAG 的优势在于：Tor 群的几何意义现在是纤维积结构的''高阶厚度''，而非仅作为数值出现。\(\blacksquare\)

\subsection{谱代数几何简介}\label{ux8c31ux4ee3ux6570ux51e0ux4f55ux7b80ux4ecb}

\textbf{定义 8}（仿射谱概形）：对任意 \(E_\infty\)-环谱 \(R\)，其\textbf{仿射谱概形} \(\operatorname{Spec} R\) 是函子 \[\operatorname{Spec} R: \mathbf{E}_\infty\text{-}\mathbf{Ring}^{\mathrm{op}} \to \mathcal{S}, \quad A \mapsto \operatorname{Map}_{\mathbf{E}_\infty}(R, A),\] 其中 \(\mathcal{S}\) 为 \(\infty\)-广群的 \(\infty\)-范畴。

\textbf{定义 9}（Lurie 的 SAG 框架）：在 Lurie 的谱代数几何（Spectral Algebraic Geometry）中：

\begin{enumerate}
\def\labelenumi{\arabic{enumi}.}
\tightlist
\item
  \textbf{谱概形}是满足局部仿射条件的 \(\infty\)-层 \(\mathbf{E}_\infty\text{-}\mathbf{Ring}^{\mathrm{op}} \to \mathcal{S}\)。
\item
  \textbf{谱 Deligne-Mumford 叠}是对合群（étale groupoid）对象在谱概形范畴中的商。
\item
  经典概形通过 Eilenberg-MacLane 函子 \(H: \mathbf{CRing} \to \mathbf{E}_\infty\text{-}\mathbf{Ring}\) 全忠实嵌入：对交换环 \(R\)，\(HR\) 是以 \(R\) 为 \(\pi_0\) 且 \(\pi_{>0}=0\) 的 \(E_\infty\)-环谱。
\item
  截断函子 \(\tau_{\leq 0}\) 将谱概形映为经典概形，且 \(\tau_{\leq 0} \circ H \cong \operatorname{id}_{\mathbf{CRing}}\)。
\end{enumerate}

谱概形与 dg 概形的关系：在 \(\mathbb{Q}\)-代数上两者等价（由于 cdga 与 \(H\mathbb{Q}\)-代数的形式性），但在正特征或整数环上，谱概形理论具有更丰富的结构（源于 \(E_\infty\)-operad 比 cdga 更强的同伦相干性）。Lurie 的 SAG 在此基础上建立了完整的拟凝聚层理论、平展上同调以及导出代数几何的六函子形式化。

\subsection{导出叠与函子观点}\label{ux5bfcux51faux53e0ux4e0eux51fdux5b50ux89c2ux70b9}

\textbf{定义 10}（导出叠）：一个\textbf{导出叠}（derived stack）是函子 \[F: \mathbf{dAff}^{\mathrm{op}} \longrightarrow \mathcal{S},\] 其中 \(\mathbf{dAff}\) 为仿射导出概形的 \(\infty\)-范畴，\(\mathcal{S}\) 为 \(\infty\)-广群的 \(\infty\)-范畴，且 \(F\) 对导出概形范畴中的平展超覆叠满足下降（descent）条件。经典叠的 2-范畴通过 \(\pi_0\) 截断嵌入导出叠范畴。

导出叠将经典叠理论从 2-范畴提升到了 \(\infty\)-范畴的高度，其优势在于能够自然地处理模空间中的高阶同伦信息。在这一框架下，Artin 叠的表征定理------形变与障碍理论由余切复形控制------获得了极其简洁的表述。

\textbf{定理 6}（导出 Artin 叠的表征）：导出 Artin 叠的形变与障碍理论由其余切复形 \(\mathbb{L}_F\)（定义在导出概形基上的导出范畴层）完全控制。具体而言，对任意导出概形 \(T\) 的平方零扩张 \(\widetilde{T}\)，填充性障碍落在 \(\operatorname{Ext}^1(\mathbb{L}_F|_T, \mathcal{I})\) 中，其中 \(\mathcal{I}\) 为平方零理想。

\textbf{\emph{证明}}：导出 Artin 叠的定义要求 \(F\) 有光滑图册（atlas）。在导出框架下，光滑性的判定由 \(\mathbb{L}_F\) 控制（定理 3 的全局版本）。障碍群通过映射 \(T \to F\) 导出序列 \(\mathbb{L}_F|_T \to \mathbb{L}_T \to \mathbb{L}_{T/F}\) 计算，\(\mathbb{L}_{T/F}\) 在同调度 \([-1,0]\) 中的上同调对应一阶形变的切与障碍空间。\(\blacksquare\)

\subsection{导出概形的平展上同调与六函子}\label{ux5bfcux51faux6982ux5f62ux7684ux5e73ux5c55ux4e0aux540cux8c03ux4e0eux516dux51fdux5b50}

导出代数几何的一个核心成就是为经典代数几何中的上同调理论提供了自然的高阶框架。在导出概形上，平展上同调和六函子体系均获得了完整的 \(\infty\)-范畴化表述，使得基变换和投影公式等经典相容性条件在同伦意义下精确成立。

\textbf{定理 7}（导出平展上同调）：对导出概形 \(X\) 和 \(\ell\)-进层 \(\mathcal{F}\)（\(\ell\) 与剩余特征互质），导出平展上同调 \(R\Gamma_{\text{\'et}}(X, \mathcal{F})\) 可赋予 \(\ell\)-进导出范畴中的对象结构，且经典平展上同调群 \(H^*_{\text{\'et}}(X_{\text{cl}}, \mathcal{F})\) 由截断恢复： \[H^n_{\text{\'et}}(X_{\text{cl}}, \mathcal{F}) \cong H^n(R\Gamma_{\text{\'et}}(X, \mathcal{F})).\]

\textbf{证明}：在导出概形框架下，平展景 \(\text{\'et}(X)\) 提升为 \(\infty\)-平展景，其上的层取值于 \(\ell\)-进谱的 \(\infty\)-范畴。\(R\Gamma_{\text{\'et}}(X, \mathcal{F})\) 定义为 \(\mathcal{F}\) 在 \(\infty\)-景上的 \(\infty\)-范畴全局截面函子的右导出。经典截断 \(t_0(X) = X_{\text{cl}}\) 给出包含 \(X_{\text{cl}} \hookrightarrow X\)，诱导平展上同调的拉回 \(R\Gamma_{\text{\'et}}(X, \mathcal{F}) \to R\Gamma_{\text{\'et}}(X_{\text{cl}}, \mathcal{F}|_{X_{\text{cl}}})\)。由于 \(\mathcal{F}\) 为 \(\ell\)-进层，其同伦层 \(\pi_{>0}(\mathcal{F})\) 在平展景上零调（由 \(\ell\)-进上同调理论），故该拉回诱导同调同构，即 \(H^n(R\Gamma_{\text{\'et}}(X, \mathcal{F})) \cong H^n_{\text{\'et}}(X_{\text{cl}}, \mathcal{F})\)。\(\blacksquare\)

\textbf{定理 8}（导出六函子形式化）：在导出概形的 \(\infty\)-范畴上，存在六个函子 \((f^*, f_*, f_!, f^!, \otimes^{\mathbb{L}}, \operatorname{\mathcal{H}om}^{\mathbb{L}})\)，满足基变换和投影公式的同伦相容性------这是经典代数几何中拟凝聚层六函子的完整 \(\infty\)-范畴化，由 Gaitsgory-Rozenblyum 在导出代数几何专著中系统建立。

\textbf{\emph{证明}}：六函子形式化的核心困难在于 \(f_!\) 与 \(f^!\) 在导出框架下的恰当定义。对概形态射 \(f: X \to Y\)，若 \(f\) 为紧合，则 \(f_! = f_*\) 自然由导出直像给出；若 \(f\) 为平展，则 \(f_!\) 要求''紧支撑直像''，需通过拟凝聚层的归纳极限构造。导出环境中 \(f^!\) 通过 Grothendieck 对偶的 \(\infty\)-范畴变体定义：\(\operatorname{\mathcal{H}om}(f_! \mathcal{F}, \mathcal{G}) \cong f_* \operatorname{\mathcal{H}om}(\mathcal{F}, f^! \mathcal{G})\)。这六个函子之间的相容性由基变换同构的 \(\infty\)-范畴精确化（即同伦极限与余极限的交换性）保证。\(\blacksquare\)

\subsection{高阶范畴论导引：∞-范畴}\label{ux9ad8ux9636ux8303ux7574ux8bbaux5bfcux5f15-ux8303ux7574}

高阶范畴论（特别是 \((\infty,1)\)-范畴论）是 21 世纪数学基础的最重要发展之一，由 Joyal、Lurie、Rezk 和 Bergner 等人在 2000 年代系统建立。其核心动机来自三个方向：模空间的叠结构需要 \(2\)-范畴或更高范畴来描述自同构群的叠层信息；导出范畴的函子范畴需要 \(\infty\)-范畴来编码所有高阶同伦信息（而不仅是同伦范畴）；拓扑量子场论（TQFT）的完整公理化需要 \((\infty,d)\)-范畴来捕捉配边的高阶结构。\((\infty,1)\)-范畴的三种主要模型------拟范畴（quasi-categories, Joyal-Boardman-Vogt）、完备 Segal 空间（Rezk）和相对范畴（Barwick-Kan）------被证明是等价的，而 Lurie 在其鸿篇巨制《Higher Topos Theory》和《Higher Algebra》中以拟范畴为基本模型，建立了 \(\infty\)-范畴上的完整理论体系（包括 \(\infty\)-拓扑、\(\infty\)-算畴、稳定 \(\infty\)-范畴等）。高阶范畴论不仅是代数拓扑和代数几何的统一语言，还深刻地影响了逻辑学（同伦类型论）、理论计算机科学（同伦类型论作为编程语言）和数学物理（扩展 TQFT 的分类定理）。

\subsection{动机：为什么需要高阶范畴}\label{ux52a8ux673aux4e3aux4ec0ux4e48ux9700ux8981ux9ad8ux9636ux8303ux7574}

经典范畴论以对象与态射为基本要素，但在以下数学情境中，1-范畴框架暴露出本质性不足：

\textbf{动机一（模空间的叠结构）}：代数簇上的向量丛模空间 \(\operatorname{Bun}_G(X)\) 中，两个向量丛之间的同构构成集合，但同构之间的自然同构不可见。这使得模空间对''分类函子''仅为伪函子（pseudo-functor），其真正的数学结构需要 2-范畴或更高范畴来描述。具体而言，\(B\mathbf{G}_m\) 不是集合而是 1-态射范畴，整个分类结构必须提升为叠（stack）以捕捉自同构群的数据。

\textbf{动机二（导出范畴的函子范畴问题）}：给定两个 Abel 范畴 \(\mathcal{A}, \mathcal{B}\)，导出范畴 \(\mathbf{D}(\mathcal{A})\) 中的对象为链复形。然而 \(\mathbf{D}(\mathcal{A}) \to \mathbf{D}(\mathcal{B})\) 的函子并非简单地来自 \(K(\mathcal{A}) \to K(\mathcal{B})\)------同伦范畴层面的函子信息不足以重构所有导出函子。真正的结构是 \(\infty\)-范畴 \(\mathcal{D}(\mathcal{A})\)（其中态射空间为 \(\operatorname{RHom}\) 复形），而普通导出范畴仅为其同伦范畴 \(h\mathcal{D}(\mathcal{A})\)，丢失了高阶同伦信息。

\textbf{动机三（拓扑量子场论 TQFT）}：Atiyah-Segal 框架中，\(d\)-维 TQFT 是对称幺半函子 \(Z: \mathbf{Bord}_d \to \mathbf{Vect}\)。但配边范畴 \(\mathbf{Bord}_d\) 中的态射需以微分同胚为 2-态射、配边的配边为 3-态射\ldots\ldots 完整的公理化需要 \((\infty,d)\)-范畴，例如 Lurie 的配边假设（Cobordism Hypothesis）必须在 \((\infty,d)\)-范畴框架中陈述。

\subsection{\texorpdfstring{\((\infty,1)\)-范畴的直观定义}{(\textbackslash infty,1)-范畴的直观定义}}\label{infty1-ux8303ux7574ux7684ux76f4ux89c2ux5b9aux4e49}

\textbf{定义 1}（\((\infty,1)\)-范畴，直观）：一个\textbf{\((\infty,1)\)-范畴} \(\mathcal{C}\) 由以下数据构成：

\begin{enumerate}
\def\labelenumi{\arabic{enumi}.}
\tightlist
\item
  一族对象 \(\operatorname{Ob}(\mathcal{C})\)；
\item
  对任意对象 \(X, Y\)，一个态射 \(\infty\)-广群 \(\operatorname{Map}_{\mathcal{C}}(X,Y)\)（或一个单纯集 / 拓扑空间），其点称为 1-态射，点之间的道路称为 2-态射，依此类推------所有 \(k \geq 2\) 的态射均为可逆；
\item
  对每个 \(X\)，有一个恒等态射 \(\operatorname{id}_X \in \operatorname{Map}_{\mathcal{C}}(X,X)\)；
\item
  对任意 \(X, Y, Z\)，有一个合成映射 \[\circ: \operatorname{Map}_{\mathcal{C}}(Y,Z) \times \operatorname{Map}_{\mathcal{C}}(X,Y) \to \operatorname{Map}_{\mathcal{C}}(X,Z),\] 该合成映射在''最高同伦意义''上满足结合律和恒等律------即存在指定同伦（结合子）使得合成在同伦意义上表现良好。
\end{enumerate}

直观核心：\((\infty,1)\)-范畴是''态射可同伦合成，更高态射皆可逆''的广义范畴。

\subsection{三种模型的正式定义与等价性}\label{ux4e09ux79cdux6a21ux578bux7684ux6b63ux5f0fux5b9aux4e49ux4e0eux7b49ux4ef7ux6027}

\textbf{定义 2}（拟范畴）：一个\textbf{拟范畴}（quasi-category，或称 \(\infty\)-范畴）是满足\textbf{内 Kan 条件}的单纯集 \(X: \Delta^{\mathrm{op}} \to \mathbf{Set}\)：对任意 \(0 < k < n\) 及任意内 horn \(\Lambda_k^n \to X\)，存在填充 \[\begin{CD}
\Lambda_k^n @>>> X \\
@VVV @VVV \\
\Delta^n @>>> \Delta^0
\end{CD}\] 使上图交换。等价于说，任意内角 \(\Lambda_k^n\) 到 \(X\) 的映射可延拓至整个 \(n\)-单纯形 \(\Delta^n\)。

\textbf{定理 1}（Joyal 提升定理）：拟范畴范畴 \(\mathbf{QCat}\) 与单纯集范畴 \(\mathbf{sSet}\) 之间存在一个 Joyal 模型结构，其纤维性对象恰为拟范畴，弱等价为范畴等价（categorical equivalence）。在此模型结构下，\(\mathbf{QCat}\) 是 \(\mathbf{sSet}\) 在 \(\mathbf{Cat}\) 上沿神经函子 \(N\) 的左 Bousfield 局部化。

\textbf{\emph{证明}}：Joyal 模型结构的余纤维化为单射，弱等价为满足以下条件的映射 \(f: X \to Y\)：对任意拟范畴 \(C\)，诱导映射 \(\tau_1(Y^C) \to \tau_1(X^C)\)（其中 \(\tau_1\) 是范畴的基本范畴构造，\(Y^C\) 为内 Hom）为范畴等价。内 Kan 条件对应纤维性对象的刻画通过 horn 填充性质验证。局部化解释：\(\mathbf{sSet}\) 在 \(N\) 上沿映射 \(\Lambda_k^n \to \Delta^n\)（\(0<k<n\)）左 Bousfield 局部化即得 Joyal 模型结构。\(\blacksquare\)

\textbf{定义 3}（完备 Segal 空间）：一个\textbf{完备 Segal 空间}是双单纯集（bisimplicial set）\(W: \Delta^{\mathrm{op}} \to \mathbf{sSet}\)，满足：

\begin{enumerate}
\def\labelenumi{\arabic{enumi}.}
\tightlist
\item
  （Reedy 纤维性）\(W\) 在 Reedy 模型结构中是纤维性对象；
\item
  （Segal 条件）对每个 \(n \geq 2\)，自然映射 \[W_n \to W_1 \times_{W_0} W_1 \times_{W_0} \cdots \times_{W_0} W_1 \quad (n\text{ 个因子})\] 是单纯集的弱等价；
\item
  （完备性条件）映射 \(W_0 \to W_{\text{equiv}}\) 是弱等价，其中 \(W_{\text{equiv}} \subseteq W_1\) 为由''同伦可逆''态射构成的子单纯集（定义为同伦纤维积中的某个并集）。
\end{enumerate}

\textbf{定理 2}（Rezk）：完备 Segal 空间范畴在 Rezk 模型结构下与拟范畴范畴 Quillen 等价。

\textbf{\emph{证明}}：从单纯集到双单纯集的函子 \(p_1^*: \mathbf{sSet} \to \mathbf{ssSet}\)（沿投影 \(\Delta \times \Delta \to \Delta\) 拉回）具有右伴随 \(i_1^*\)（对角线：\((i_1^*W)_n = W_{n,n}\)）。这对伴随函子诱导 Quillen 等价，其核心验证在于：完备 Segal 空间的对角线为拟范畴，而拟范畴沿 \(p_1^*\) 的延拓产生完备 Segal 空间。Segal 条件确保合成在同伦意义上定义良好，完备性条件确保等价概念与 Joyal 意义一致。\(\blacksquare\)

\textbf{定义 4}（单纯富范畴）：一个\textbf{单纯富范畴} \(\mathcal{C}\) 由一族对象和单纯集 \(\operatorname{Map}_{\mathcal{C}}(X,Y)\)（态射空间）组成，配备严格结合且单位的合成律 \(\operatorname{Map}_{\mathcal{C}}(Y,Z) \times \operatorname{Map}_{\mathcal{C}}(X,Y) \to \operatorname{Map}_{\mathcal{C}}(X,Z)\)。此类范畴构成的范畴记为 \(\mathbf{sCat}\)。

\textbf{定理 3}（Bergner 模型结构）：\(\mathbf{sCat}\) 具有一个模型结构，其中弱等价为 Dwyer-Kan 意义下的单纯富范畴等价（即在 \(\pi_0\) 层面为范畴等价，且诱导每个态射空间的单纯同伦型等价）。\(\mathbf{sCat}\) 的此模型结构与 \(\mathbf{QCat}\) Quillen 等价。

\textbf{\emph{证明}}：Dwyer-Kan 局部化构造结合函子 \(\mathbf{sCat} \to \mathbf{sSet}\)（同伦相干神经，homotopy coherent nerve）建立 Quillen 等价。具体地，对每个单纯富范畴 \(\mathcal{C}\)，定义其同伦相干神经 \(N^{\text{hc}}(\mathcal{C})\) 为单纯集，其 \(n\)-单纯形对应于 \(\mathfrak{C}[\Delta^n] \to \mathcal{C}\)（其中 \(\mathfrak{C}\) 为单纯集的刚性化函子）。验证 \(N^{\text{hc}}(\mathcal{C})\) 满足内 Kan 条件，拟范畴 \(X\) 对应 \(\mathfrak{C}[X]\)（此函子为左 Quillen 函子）。三者 Quillen 等价的完整链条为：\(\mathbf{sCat} \rightleftarrows \mathbf{sSet}_{\text{Joyal}} \rightleftarrows \mathbf{CSS}_{\text{Rezk}}\)。\(\blacksquare\)

\subsection{\texorpdfstring{\(\infty\)-范畴中的极限与余极限}{\textbackslash infty-范畴中的极限与余极限}}\label{infty-ux8303ux7574ux4e2dux7684ux6781ux9650ux4e0eux4f59ux6781ux9650}

\textbf{定义 5}（\(\infty\)-范畴中的极限）：设 \(\mathcal{C}\) 为 \(\infty\)-范畴（拟范畴），\(K\) 为单纯集，图 \(p: K \to \mathcal{C}\)。\(p\) 的\textbf{极限}是一个对象 \(\lim p \in \mathcal{C}\)，连同态射 \(\lim p \to p(k)\)（对所有 \(k \in K\)）使其在切片 \(\infty\)-范畴 \(\mathcal{C}_{/p}\) 中为终对象。等价地，\(\lim p\) 是 \(\mathcal{C}\) 中到锥范畴的典范映射的左伴随在 \(p\) 处的取值。对偶地，\textbf{余极限} \(\operatorname{colim} p\) 是在 \(\mathcal{C}_{p/}\) 中的始对象。

在 \(\infty\)-范畴中，极限与余极限的概念需要提升到同伦意义下：不同于普通范畴中极限由严格的万有性质定义，\(\infty\)-范畴中的极限要求函子范畴中的自然同伦等价。具体而言，对任意对象 \(X \in \mathcal{C}\)，典范映射 \(\operatorname{Map}_{\mathcal{C}}(X, \lim p) \to \lim_{k \in K} \operatorname{Map}_{\mathcal{C}}(X, p(k))\) 是广群（\(\infty\)-群胚）范畴中的同伦等价。这一条件确保了极限在函子 \(\mathcal{C} \to \mathcal{S}\) 下的行为与普通范畴一致，但弱化了唯一性为同伦唯一性。Joyal 和 Lurie 证明了拟范畴模型中的极限满足所有期望的性质：极限的函子性、保极限的伴随关系、以及极限与 \(\infty\)-意象中 Colimit 的交换性。特别地，\(\infty\)-范畴中的拉回和推出（同伦拉回和同伦推出）是极限与余极限的最基本实例，它们构成了 \(\infty\)-范畴中代数结构的基础。

\subsection{伴随函子定理与可表性}\label{ux4f34ux968fux51fdux5b50ux5b9aux7406ux4e0eux53efux8868ux6027}

\textbf{定义 6}（\(\infty\)-范畴中的伴随）：设 \(F: \mathcal{C} \to \mathcal{D}\) 和 \(G: \mathcal{D} \to \mathcal{C}\) 为 \(\infty\)-范畴之间的函子。称 \(F\) 左伴随于 \(G\)（记 \(F \dashv G\)），若存在态射空间的自然同伦等价： \[\operatorname{Map}_{\mathcal{D}}(F(C), D) \simeq \operatorname{Map}_{\mathcal{C}}(C, G(D)),\] 对任意 \(C \in \mathcal{C}, D \in \mathcal{D}\)，且该等价''函子性''地依赖于 \(C\) 与 \(D\)。

\textbf{定理 4}（Lurie 的伴随函子定理，陈述）：设 \(F: \mathcal{C} \to \mathcal{D}\) 为可表 \(\infty\)-范畴之间的函子。则 \(F\) 有右伴随当且仅当 \(F\) 保余极限。对偶地，\(F\) 有左伴随当且仅当 \(F\) 保极限。

\textbf{证明}：仅证 \(F\) 有右伴随当且仅当 \(F\) 保余极限。(\(\Rightarrow\)) 若 \(F \dashv G\)，则对任意图 \(p: K \to \mathcal{C}\)，\(\operatorname{Map}(F(\operatorname{colim} p), D) \simeq \operatorname{Map}(\operatorname{colim} p, G(D)) \simeq \lim \operatorname{Map}(p(k), G(D)) \simeq \lim \operatorname{Map}(F(p(k)), D) \simeq \operatorname{Map}(\operatorname{colim} F \circ p, D)\)，由 \(\infty\)-Yoneda 得 \(F(\operatorname{colim} p) \simeq \operatorname{colim} F \circ p\)。(\(\Leftarrow\)) 若 \(F\) 保余极限，因 \(\mathcal{C}\) 可表，由 \(\infty\)-范畴的伴随函子定理，对每个 \(D \in \mathcal{D}\)，函子 \(\operatorname{Map}(F(-), D): \mathcal{C}^{\mathrm{op}} \to \mathcal{S}\) 将余极限变为极限，故可表。定义 \(G(D)\) 为可表对象，则 \(\operatorname{Map}(F(C), D) \simeq \operatorname{Map}(C, G(D))\)，即 \(F \dashv G\)。\(\blacksquare\)

\textbf{定义 7}（可表 \(\infty\)-范畴）：\(\infty\)-范畴 \(\mathcal{C}\) 称为\textbf{可表的}（presentable），若 \(\mathcal{C}\) 连通于某个小 \(\infty\)-范畴的单纯预层范畴 \(\mathcal{P}(K) = \operatorname{Fun}(K^{\mathrm{op}}, \mathcal{S})\) 的某种局部化。等价地，\(\mathcal{C}\) 是小余完备的（所有小余极限存在）且可由小对象集合 \(\kappa\)-强生成。

\textbf{定理 5}（Lurie 可表性定理）：\(\infty\)-范畴 \(\mathbf{Pr}^L\)（由可表 \(\infty\)-范畴和保余极限函子构成）中，全体可表 \(\infty\)-范畴均由单纯预层 \(\infty\)-范畴 \(\mathcal{P}(K)\)（\(K\) 为小 \(\infty\)-范畴）经过反射子（accessible reflective subcategory）------即对一族态射局部化------得到。

\textbf{证明}：设 \(\mathcal{C}\) 为可表 \(\infty\)-范畴。由定义，\(\mathcal{C}\) 连通于 \(\mathcal{P}(K)\) 的局部化，即 \(\mathcal{C} \simeq \mathcal{P}(K)[W^{-1}]\)（\(W\) 为某族态射）。反之，\(\mathcal{P}(K)\) 的任何连通局部化均保持可表性。\(\mathcal{C}\) 的小对象生成集可由 \(K\) 的可表预层诱导。在 \(\mathbf{Pr}^L\) 中，态射为保余极限函子，\(\mathcal{P}(K)\) 是自由可表 \(\infty\)-范畴（在 \(K\) 上自由生成）。\(\mathcal{C}\) 作为 \(\mathcal{P}(K)\) 的反射子范畴，反射子 \(L: \mathcal{P}(K) \to \mathcal{C}\) 是左伴随，且 \(L\) 等同于对 \(W\) 的局部化。全体可表 \(\infty\)-范畴因此为 \(\mathcal{P}(K)\) 的反射子范畴。\(\blacksquare\)

\subsection{\texorpdfstring{稳定 \(\infty\)-范畴}{稳定 \textbackslash infty-范畴}}\label{ux7a33ux5b9a-infty-ux8303ux7574}

稳定 \(\infty\)-范畴是导出范畴概念在高阶范畴论中的最终推广，它刻画了那些具有''同伦意义下的加性结构''且纤维序列与余纤维序列重合的 \(\infty\)-范畴。在同调代数中，导出范畴 \(D(\mathcal{A})\) 是稳定 \(\infty\)-范畴的原型；在代数拓扑中，谱的 \(\infty\)-范畴 \(\mathbf{Sp}\) 也是稳定的。稳定 \(\infty\)-范畴的核心性质是：它是切空间可以被典范地定义的数学环境。

\textbf{定义 8}（稳定 \(\infty\)-范畴）：一个 \(\infty\)-范畴 \(\mathcal{C}\) 称为\textbf{稳定的}（stable），若：

\begin{enumerate}
\def\labelenumi{\arabic{enumi}.}
\tightlist
\item
  \(\mathcal{C}\) 是尖的：存在零对象 \(0 \in \mathcal{C}\)（始且终）；
\item
  \(\mathcal{C}\) 有所有有限极限与余极限；
\item
  \(\mathcal{C}\) 中的纤维序列与余纤维序列重合：对任意态射 \(f: X \to Y\)，其纤维（fiber）自然同构于其余纤维（cofiber）的环路 \(\Omega\)，即图表 \[X \xrightarrow{f} Y \to \operatorname{cofib}(f) \simeq \Sigma(\operatorname{fib}(f))\] 与纤维序列重合。
\end{enumerate}

\textbf{定理 6}（稳定 \(\infty\)-范畴与谱）：存在稳定 \(\infty\)-范畴 \(\mathbf{Sp}\)（谱的 \(\infty\)-范畴）使得任意稳定 \(\infty\)-范畴 \(\mathcal{C}\) 均 \(\mathbf{Sp}\)-富集。具体而言，\(\operatorname{Map}_{\mathcal{C}}(X,Y)\) 典范地提升为谱。\(\mathbf{Sp}\) 是 \(\mathcal{S}_*^{\text{fin}}\)（有限尖广群的 \(\infty\)-范畴）依谱化（stabilization）构造 \(\operatorname{Sp}(\mathcal{S}_*)\) 得到的万有稳定 \(\infty\)-范畴。

\textbf{\emph{证明}}：稳定化的万有构造：对 \(\infty\)-范畴 \(\mathcal{C}\)（有有限极限），定义 \(\operatorname{Sp}(\mathcal{C}) = \varprojlim (\cdots \xrightarrow{\Omega} \mathcal{C}_* \xrightarrow{\Omega} \mathcal{C}_*)\)，其中 \(\Omega\) 为环路函子。当 \(\mathcal{C} = \mathcal{S}\)（广群的 \(\infty\)-范畴）时，\(\operatorname{Sp}(\mathcal{S})\) 即为 \(\mathbf{Sp}\)。对一般稳定 \(\mathcal{C}\)，\(\Omega: \mathcal{C} \to \mathcal{C}\) 是等价，且 \(\mathcal{C}\) 典范地为 \(\mathbf{Sp}\)-模，通过谱作用 \(\mathbf{Sp} \otimes \mathcal{C} \to \mathcal{C}\)（由万有性质诱导）。此构造模仿了经典的三棱镜链复形中对 \(\infty\)-范畴稳定化的推广。\(\blacksquare\)

\subsection{\texorpdfstring{稳定 \(\infty\)-范畴的同调代数}{稳定 \textbackslash infty-范畴的同调代数}}\label{ux7a33ux5b9a-infty-ux8303ux7574ux7684ux540cux8c03ux4ee3ux6570}

\textbf{定理 7}（t-结构与心）：稳定 \(\infty\)-范畴 \(\mathcal{C}\) 上的 \textbf{t-结构} 是一对满子范畴 \((\mathcal{C}_{\geq 0}, \mathcal{C}_{\leq 0})\)，满足 \(\mathcal{C}_{\geq 0}[1] \subset \mathcal{C}_{\geq 0}\)、\(\mathcal{C}_{\leq 0} \subset \mathcal{C}_{\leq 0}[1]\)、\(\operatorname{Map}(X, Y[-1]) \simeq 0\) 对 \(X \in \mathcal{C}_{\geq 0}, Y \in \mathcal{C}_{\leq 0}\)，且任意对象 \(X \in \mathcal{C}\) 有纤维序列 \(X_{\geq 1} \to X \to X_{\leq 0}\)（其中 \(X_{\geq 1} \in \mathcal{C}_{\geq 1}\)，\(X_{\leq 0} \in \mathcal{C}_{\leq 0}\)）。t-结构的\textbf{心} \(\mathcal{C}^\heartsuit = \mathcal{C}_{\geq 0} \cap \mathcal{C}_{\leq 0}\) 是 Abel 范畴。

\textbf{\emph{证明}}：心的 Abel 范畴结构来源于 \(\mathcal{C}\) 中截断函子的存在性。\(\mathcal{C}_{\geq 0}\) 和 \(\mathcal{C}_{\leq 0}\) 的截断提供典范纤维序列，在此心上的态射集合恰为 \(\pi_0 \operatorname{Map}(X,Y)\)，加法结构由环路-悬挂伴随 \(\Omega \dashv \Sigma\) 在 t-结构下的作用导出。核与余核由截断后的纤维与余纤维构造，满足 Abel 范畴公理。\(\blacksquare\)

\textbf{定理 8}（导出 \(\infty\)-范畴的万有性质）：对任意 Abel 范畴 \(\mathcal{A}\)，其导出 \(\infty\)-范畴 \(\mathcal{D}(\mathcal{A})\)（以 \(\mathcal{A}\) 中对象的链复形为对象，链映射为 1-态射，链同伦为 2-态射）是典范的同调代数工作环境，且 \(\mathcal{D}(\mathcal{A})\) 的标准 t-结构满足 \(\mathcal{D}(\mathcal{A})^\heartsuit \simeq \mathcal{A}\)。

\textbf{证明}：构造 \(\mathcal{D}(\mathcal{A})\) 为 \(\mathcal{A}\) 中链复形的 dg-神经（dg-nerve），或等价地，为 \(\mathbf{Ch}(\mathcal{A})\) 沿拟同构的 \(\infty\)-局部化。平移函子 \([1]\) 由复形平移给出，映射锥赋予 \(\mathcal{D}(\mathcal{A})\) 三角结构。标准 t-结构 \((\mathcal{D}_{\geq 0}, \mathcal{D}_{\leq 0})\) 由同调截断定义：\(\mathcal{D}_{\geq 0}(\mathcal{A}) = \{X \mid H_i(X) = 0, i < 0\}\)，\(\mathcal{D}_{\leq 0}(\mathcal{A}) = \{X \mid H_i(X) = 0, i > 0\}\)。验证 t-结构公理：\(\mathcal{D}_{\geq 0}[1] \subset \mathcal{D}_{\geq 0}\) 显然；\(\operatorname{Map}(X, Y[-1]) \simeq 0\) 对 \(X \in \mathcal{D}_{\geq 0}, Y \in \mathcal{D}_{\leq 0}\) 由同调维度计算；截断序列 \(X_{\geq 1} \to X \to X_{\leq 0}\) 为典范截断复形序列。心 \(\mathcal{D}(\mathcal{A})^\heartsuit = \mathcal{D}_{\geq 0} \cap \mathcal{D}_{\leq 0}\) 恰为同调集中在度 0 的复形范畴，等价于 \(\mathcal{A}\)。\(\blacksquare\)

\subsection{导出范畴与模空间}\label{ux5bfcux51faux8303ux7574ux4e0eux6a21ux7a7aux95f4}

在 Lurie 框架下，稳定 \(\infty\)-范畴和可表 \(\infty\)-范畴的语言统一处理了模空间的形变、障碍与局部-整体原理。导出 Artin 叠 \(\mathcal{M}\) 的拟凝聚层范畴 \(\operatorname{QCoh}(\mathcal{M})\) 是稳定、可表的对称幺半 \(\infty\)-范畴，而经典叠上的拟凝聚层范畴仅为其心。这一层次化使得 Langlands 纲领中的几何 Satake 等价、Beilinson-Bernstein 局部化等深刻结果均可在 \(\infty\)-范畴框架下获得自然推广。

导出模空间理论的核心思想是：模空间本身不仅是集合或概形，而应被视为一个整体的导出结构------即导出叠（derived stack）。在导出叠 \(\mathcal{M}\) 上，点（即 \(\mathcal{M}(R)\) 中的 \(R\)-值点）可携带非平凡的高阶自同构（\(\infty\)-群胚结构），这恰好对应于模问题的形变理论的障碍群。具体而言，对模空间 \(\mathcal{M}\) 上一点 \(x\)，其切复形 \(\mathbb{T}_x\mathcal{M}\) 的上同调群 \(H^0\) 给出经典切空间，\(H^1\) 给出形变空间，\(H^2\) 给出障碍空间。\(\operatorname{QCoh}(\mathcal{M})\) 的稳定 \(\infty\)-范畴结构使得 Grothendieck-Serre 对偶、Verdier 对偶以及六函子形式化（\(f^*, f_*, f_!, f^!, \otimes, \mathcal{H}om\)）可以严格地运作，克服了经典导出范畴中伪函子性和幺半结构的技术障碍。Toën 和 Vezzosi 在导出的代数几何纲领中证明，许多经典模空间（如向量丛的模空间、Hitchin 系的模空间）在导出视角下展现出更丰富的几何结构，其导出加强版具有更好的函子性质。

\subsection{\texorpdfstring{\(\infty\)-范畴视角下的重构}{\textbackslash infty-范畴视角下的重构}}\label{infty-ux8303ux7574ux89c6ux89d2ux4e0bux7684ux91cdux6784}

\textbf{定理 9}（\(\infty\)-范畴化 Yoneda 引理）：对任意 \(\infty\)-范畴 \(\mathcal{C}\)，函子 \[j: \mathcal{C} \longrightarrow \operatorname{Fun}(\mathcal{C}^{\mathrm{op}}, \mathcal{S}), \quad X \longmapsto \operatorname{Map}_{\mathcal{C}}(-, X)\] 是全忠实的。即对任意 \(X, Y \in \mathcal{C}\)，\(\operatorname{Map}_{\mathcal{C}}(X,Y)\) 与 \(\operatorname{Map}_{\operatorname{Fun}}(j(X), j(Y))\) 典范同伦等价。此引理将 \(\infty\)-范畴完全嵌入到其预层 \(\infty\)-范畴中，是可表性定理与伴随函子定理的基础。

\textbf{证明}：沿袭经典 Yoneda 引理的证明思路，将问题化归为 \(\mathcal{S}\)-值函子范畴中的端点计算。对任意 \(F \in \operatorname{Fun}(\mathcal{C}^{\mathrm{op}}, \mathcal{S})\)，\(\operatorname{Map}(j(X), F) \simeq F(X)\)（通过评估于 \(\operatorname{id}_X\) 的自然变换建立同伦等价）。取 \(F = j(Y)\) 即得全忠实性。\(\blacksquare\)

\textbf{定理 10}（\(\infty\)-意象的公理刻画，Rezk-Lurie）：一个 \(\infty\)-范畴 \(\mathcal{X}\) 是 \textbf{\(\infty\)-意象} 当且仅当：

\begin{enumerate}
\def\labelenumi{\arabic{enumi}.}
\tightlist
\item
  \(\mathcal{X}\) 是可表的；
\item
  \(\mathcal{X}\) 对某对象族的余切（colimit）封闭；
\item
  \(\mathcal{X}\) 满足 \(\infty\)-范畴的 Giraud 公理（余积万有、等价关系有效等）。
\end{enumerate}

\(\infty\)-意象是 1-意象（Grothendieck topos）的 \(\infty\)-范畴化，其内部逻辑是高阶类型论的语义模型，为单值公理及其模型提供了自然的范畴环境。

\textbf{证明}：(\(\Rightarrow\)) 若 \(\mathcal{X}\) 为 \(\infty\)-意象，则 \(\mathcal{X} \simeq \operatorname{Sh}_\infty(\mathcal{C})\)（\(\mathcal{C}\) 为某 \(\infty\)-景）。预层 \(\infty\)-范畴 \(\mathcal{P}(\mathcal{C})\) 可表，\(\infty\)-层化是局部化（accessible left exact localization），故 \(\mathcal{X}\) 可表。余积万有性由 \(\infty\)-景上层的万有性质保证，等价关系有效性由 \(\infty\)-层化函子的正合性保证。(\(\Leftarrow\)) 若 \(\mathcal{X}\) 满足 Giraud 公理，则 \(\mathcal{X}\) 同构于 \(\mathcal{P}(K)\) 的拓扑局部化（对覆盖筛的局部化），其中 \(K\) 为 \(\mathcal{X}\) 中紧生成对象的满子 \(\infty\)-范畴。该局部化定义了 \(\infty\)-景结构，使 \(\mathcal{X}\) 等价于 \(\operatorname{Sh}_\infty(K)\)。\(\blacksquare\)

\subsection{应用纵览与前景}\label{ux5e94ux7528ux7eb5ux89c8ux4e0eux524dux666f}

稳定 \(\infty\)-范畴和导出代数几何的结合产生了现代数学中最富影响力的结构框架，其核心应用方向包括：

\begin{enumerate}
\def\labelenumi{\arabic{enumi}.}
\item
  \textbf{因子化同调与拓扑量子场论}：Lurie 的配边假设（Cobordism Hypothesis）在 \((\infty, n)\)-范畴中刻画了完全扩展的 TQFT，使得配边范畴的函子完全由其在一个点上的值（对偶化对象）确定。
\item
  \textbf{导出几何 Langlands 纲领}：Arinkin-Gaitsgory 等人在 \(\infty\)-范畴框架下发展了导出 D-模理论，将经典 Langlands 函子性提升为导出范畴间的等价------这为几何 Langlands 猜想的证明提供了正确的技术语言。
\item
  \textbf{色展与导出代数几何的交汇}：\(\mathbf{Sp}\) 作为万有稳定 \(\infty\)-范畴与色展塔的 \(E_n\) 理论之间的关系构成了导出代数几何与稳定同伦论的深层联系：特别地，Morava E-理论是高度 \(n\) 处形式群的导出模叠上的拟凝聚层，统一了代数几何与同伦论的色谱。
\item
  \textbf{凝聚对偶与六函子形式化}：在稳定 \(\infty\)-范畴框架中，Grothendieck 的六函子形式化（\(f^*, f_*, f_!, f^!, \otimes, \operatorname{Hom}\)）达到完全严格化，解决了经典导出范畴中伪函子性的种种技术困难，为代数几何中的对偶理论（如 Verdier 对偶、Serre 对偶）提供了坚固的 \(\infty\)-范畴基础。
\end{enumerate}

综上，高阶范畴论与 \(\infty\)-范畴不仅是范畴论自身的推广，更是当代数学诸多领域的''通用语言''------从代数几何到同伦论、从量子场论到数理逻辑，它为这些看似迥异的领域提供了统一的组织原理和计算框架。Lurie 的《高等代数》（Higher Algebra）与《高等意象》（Higher Topos Theory）构成了该理论的标准参考文献，其技术深度为整个导出与高阶数学的蓬勃发展奠定了基石。

\newpage

\chapter{08-代数几何与数论/02-数论/02-模形式.md}\label{ux4ee3ux6570ux51e0ux4f55ux4e0eux6570ux8bba02-ux6570ux8bba02-ux6a21ux5f62ux5f0f.md}

\section{第295章：模形式}\label{ux7b2c295ux7ae0ux6a21ux5f62ux5f0f}

模形式是上半复平面上具有极高对称性的全纯函数，在数论、代数几何、表示论和数学物理中都有核心应用。模形式理论的历史可以追溯到 19 世纪：Eisenstein 在 1847 年研究了 Eisenstein 级数，Dedekind 在 1877 年发现了 Dedekind \(\eta\) 函数，Poincaré 在 1880 年代引入了 Poincaré 级数，而 Hecke 在 1930 年代建立了系统的 Hecke 算子理论。模形式在现代数论中的地位由 Wiles 在 1995 年证明 Fermat 大定理而达到顶峰------模性定理（Taniyama-Shimura-Weil 猜想）将椭圆曲线与模形式联系起来，成为 Langlands 纲领在 \(\operatorname{GL}_2\) 情形下的核心成果。模形式理论的核心概念是：模形式是高度对称的全纯函数，其对称性约束（模群作用下的变换性质）如此之强，使得模形式空间是有限维的------这一''刚性''现象是模形式理论一切深刻应用的基础。本章系统介绍模形式的基本理论：模群、Eisenstein 级数与尖点形式、Hecke 算子与 L 函数、模曲线、以及 Jacobi 形式与 Siegel 模形式等高维推广。

\subsection{86.1 模群与模形式}\label{ux6a21ux7fa4ux4e0eux6a21ux5f62ux5f0f}

\textbf{定义 86.1}（上半平面与模群）：上半平面 \(\mathbb{H} = \{z \in \mathbb{C} : \Im(z) > 0\}\) 是复平面中虚部为正的开子集。模群 \(SL(2, \mathbb{Z})\) 通过分式线性变换作用于 \(\mathbb{H}\)：

\[\begin{pmatrix} a & b \\ c & d \end{pmatrix} \cdot z = \frac{az + b}{cz + d}\]

此处 \(SL(2, \mathbb{Z}) = \{\begin{pmatrix} a & b \\ c & d \end{pmatrix} : a, b, c, d \in \mathbb{Z}, ad - bc = 1\}\)。注意 \(\pm I\) 的作用相同，因此实际作用的群是 \(PSL(2, \mathbb{Z}) = SL(2, \mathbb{Z})/\{\pm I\}\)，称为模群。模群是 Fuchs 群（离散子群）的原型，其在上半平面上的作用是真不连续的（properly discontinuous），因此商空间 \(\Gamma \setminus \mathbb{H}\) 具有 Riemann 面的结构。

\textbf{定义 86.2}（\(SL(2, \mathbb{Z})\) 上的模形式）：设 \(k \in \mathbb{Z}\)。全纯函数 \(f: \mathbb{H} \to \mathbb{C}\) 称为\textbf{权为 \(k\) 的模形式}，如果：

1.（模变换）对所有 \(\gamma = \begin{pmatrix} a & b \\ c & d \end{pmatrix} \in SL(2, \mathbb{Z})\)： \[f\left(\frac{az + b}{cz + d}\right) = (cz + d)^k f(z)\]

2.（在尖点处的全纯性）\(f\) 在 \(\infty\) 处全纯，即存在 Fourier 展开（或 \(q\)-展开）： \[f(z) = \sum_{n=0}^{\infty} a_n q^n, \quad q = e^{2\pi i z}\]

若 \(a_0 = 0\)，则称 \(f\) 为\textbf{尖点形式}（cusp form）。全体权 \(k\) 的模形式构成有限维 \(\mathbb{C}\)-向量空间 \(M_k(SL(2, \mathbb{Z}))\)，尖点形式子空间记为 \(S_k(SL(2, \mathbb{Z}))\)。

\textbf{定理 86.1}（\(SL(2, \mathbb{Z})\) 的生成元）：\(SL(2, \mathbb{Z})\) 由 \(T = \begin{pmatrix} 1 & 1 \\ 0 & 1 \end{pmatrix}\) 和 \(S = \begin{pmatrix} 0 & -1 \\ 1 & 0 \end{pmatrix}\) 生成，满足关系 \(S^2 = (ST)^3 = -I\)。\(T\) 作用于 \(\mathbb{H}\) 为平移 \(z \mapsto z + 1\)，\(S\) 为反演 \(z \mapsto -1/z\)。因此模形式的两个条件等价于 \(f(z+1) = f(z)\)（周期性）和 \(f(-1/z) = z^k f(z)\)（反演对称性）。

\textbf{证明}：\(SL(2, \mathbb{Z})\) 在 \(\mathbb{H}\) 上的经典基本域为 \(\mathcal{F} = \{z \in \mathbb{H} : |z| \geq 1, |\Re(z)| \leq 1/2\}\)。由基本域的性质，\(S\) 和 \(T\) 生成整个群（通过 \(S\) 的反射和 \(T\) 的平移的复合可以到达任何基本域）。形式地，对任意 \(\gamma \in SL(2, \mathbb{Z})\)，考虑 \(\gamma\) 将 \(\mathcal{F}\) 映到某个基本域 \(\gamma(\mathcal{F})\)。通过反复应用 \(S\) 和 \(T\)（以及 \(T^{-1}\)），可以将 \(\gamma\) 分解为这些生成元的乘积。\(T(z) = z + 1\) 给出 \(f(z+1) = f(z)\)，即 \(f\) 的周期性，从而有 Fourier 展开 \(f(z) = \sum a_n e^{2\pi i n z}\)。\(S(z) = -1/z\) 给出 \(f(-1/z) = z^k f(z)\)。这两个条件等价于对所有 \(\gamma \in SL(2, \mathbb{Z})\) 的模变换条件，因为 \(S\) 和 \(T\) 生成整个群。\(\blacksquare\)

\textbf{定理 86.1b}（模形式空间的维数公式）：对 \(k \geq 0\) 为偶数（\(k \neq 2\)），\(M_k(SL(2, \mathbb{Z}))\) 的维数为： \[\dim M_k = \begin{cases} \lfloor k/12 \rfloor + 1, & k \not\equiv 2 \pmod{12} \\ \lfloor k/12 \rfloor, & k \equiv 2 \pmod{12} \end{cases}\]

尖点形式空间 \(S_k\) 的维数为 \(\dim S_k = \dim M_k - 1\)（\(k \geq 4\)）且 \(\dim S_k = \dim M_k\)（\(k < 4\) 时，即 \(k=0\) 时 \(M_0 = \mathbb{C}\) 且 \(S_0 = \{0\}\)）。对于奇数 \(k\)，\(M_k = \{0\}\)（因为 \(-I \in SL(2, \mathbb{Z})\) 且 \(f(z) = f(-I \cdot z) = (-1)^k f(z)\)，故 \(k\) 为奇数时 \(f = 0\)）。

\textbf{证明（框架）}：利用 Riemann-Roch 定理在 \(X(1) = SL(2, \mathbb{Z}) \setminus \mathbb{H}^*\)（紧化后的模曲线）上计算。\(X(1)\) 的亏格为 \(0\)（它是 Riemann 球面 \(\mathbb{P}^1(\mathbb{C})\)）。权 \(k\) 的模形式对应于 \(X(1)\) 上的 \(k/2\)-微分形式（带极点在尖点处），Riemann-Roch 给出维数公式。\(\blacksquare\)

模形式空间 \(M_k\) 的具体结构：\(M_0 = \langle 1 \rangle\)（常值函数），\(M_4 = \langle E_4 \rangle\)，\(M_6 = \langle E_6 \rangle\)，\(M_8 = \langle E_8 = E_4^2 \rangle\)，\(M_{10} = \langle E_{10} = E_4 E_6 \rangle\)，\(M_{12} = \langle E_4^3, \Delta \rangle\)。\(S_{12}\) 由模判别式 \(\Delta = (E_4^3 - E_6^2)/1728 = q \prod_{n=1}^\infty (1 - q^n)^{24}\) 生成。\(\Delta\) 的非零 Fourier 系数 \(\tau(n)\) 是 Ramanujan \(\tau\) 函数，满足深刻的乘法性质（\(\tau(mn) = \tau(m)\tau(n)\) 当 \((m,n)=1\)）和同余性质（如 \(\tau(n) \equiv \sigma_{11}(n) \pmod{691}\)，其中 \(\sigma_{11}(n) = \sum_{d \mid n} d^{11}\)）。

\textbf{定理 86.1c}（价公式 / Valence Formula）：设 \(f \in M_k\) 不恒为零。记 \(v_P(f)\) 为 \(f\) 在 \(P \in \mathcal{F}\) 处的零点阶数（在局部坐标下），\(v_\infty(f)\) 为 \(f\) 在尖点处的零点阶数（即 \(q\)-展开中 \(a_0\) 首次非零的指标）。则： \[v_\infty(f) + \frac{1}{2}v_i(f) + \frac{1}{3}v_\rho(f) + \sum_{\substack{P \in \mathcal{F} \\ P \neq i, \rho}} v_P(f) = \frac{k}{12}\]

其中 \(i = \sqrt{-1}\)（\(S\) 的不动点），\(\rho = e^{2\pi i/3}\)（\(ST\) 的不动点）。价公式是模形式理论中最基本的工具之一，它给出了任意非零模形式的零点分布约束，是维数公式的推导基础。

\textbf{证明}：考虑 \(f\) 的对数导数 \(\frac{f'}{f}\) 沿基本域边界的围道积分。由留数定理，\(\frac{1}{2\pi i} \oint_{\partial \mathcal{F}} \frac{f'(z)}{f(z)} dz = \sum_{P \in \mathcal{F}} v_P(f)\)（内部零点）。另一方面，沿边界的积分可利用 \(f\) 的模变换性质计算：\(S\) 和 \(T\) 的对称性给出边界段之间的对消关系，剩余部分贡献 \(\frac{k}{12}\)。\(\blacksquare\)

\textbf{定义 86.2b}（同余子群）：对正整数 \(N\)，定义以下同余子群： - \(\Gamma_0(N) = \{\begin{pmatrix} a & b \\ c & d \end{pmatrix} \in SL(2, \mathbb{Z}) : c \equiv 0 \pmod{N}\}\) - \(\Gamma_1(N) = \{\begin{pmatrix} a & b \\ c & d \end{pmatrix} \in SL(2, \mathbb{Z}) : c \equiv 0, a \equiv d \equiv 1 \pmod{N}\}\) - \(\Gamma(N) = \{\begin{pmatrix} a & b \\ c & d \end{pmatrix} \in SL(2, \mathbb{Z}) : \begin{pmatrix} a & b \\ c & d \end{pmatrix} \equiv \begin{pmatrix} 1 & 0 \\ 0 & 1 \end{pmatrix} \pmod{N}\}\)

\(\Gamma(N)\) 称为主同余子群，是 \(SL(2, \mathbb{Z})\) 的正规子群。对于 \(\Gamma_0(N)\)-权 \(k\) 模形式，模变换条件仅需对 \(\gamma \in \Gamma_0(N)\) 成立，空间记为 \(M_k(\Gamma_0(N))\)。同余子群上的模形式理论是模性定理和 Langlands 纲领的基础，因为 \(\Gamma_0(N)\) 上的权 \(2\) 新形式（newform）与导子为 \(N\) 的椭圆曲线通过模性定理相对应。

\subsection{86.2 Eisenstein 级数与尖点形式}\label{eisenstein-ux7ea7ux6570ux4e0eux5c16ux70b9ux5f62ux5f0f}

\textbf{定义 86.3}（Eisenstein 级数）：对偶数 \(k \geq 4\)，权 \(k\) 的 \textbf{Eisenstein 级数} 定义为

\[E_k(z) = \frac{1}{2\zeta(k)} \sum_{(m,n) \in \mathbb{Z}^2 \setminus \{(0,0)\}} \frac{1}{(mz + n)^k} = \frac{1}{2} \sum_{(c,d)=1} \frac{1}{(cz + d)^k}\]

其中 \(\zeta(k) = \sum_{n=1}^\infty 1/n^k\) 为 Riemann zeta 函数。第二个等式利用了 \(SL(2, \mathbb{Z})\) 在 \(\mathbb{Z}^2\) 的互素对上的可迁性。Eisenstein 级数对 \(k \geq 4\) 绝对收敛，且满足模变换条件（直接验证）。其 Fourier 展开为

\[E_k(z) = 1 + \frac{2}{\zeta(1-k)} \sum_{n=1}^\infty \sigma_{k-1}(n) q^n\]

其中 \(\sigma_{k-1}(n) = \sum_{d \mid n} d^{k-1}\) 为除数和函数。系数 \(\frac{2}{\zeta(1-k)}\) 通过 Bernoulli 数与 \(B_k\) 有关：\(\zeta(1-k) = -B_k/k\)，故 \(E_k(z) = 1 - \frac{2k}{B_k} \sum_{n=1}^\infty \sigma_{k-1}(n) q^n\)。

\textbf{定理 86.2}（模形式空间的分解）：\(M_k = \langle E_k \rangle \oplus S_k\)（当 \(k \geq 4\) 为偶数时）。Eisenstein 级数张成 \(M_k\) 的''Eisenstein 部分''，其 Fourier 系数是初等数论函数（除数和函数）；尖点形式张成 \(M_k\) 的''尖点部分''，其 Fourier 系数承载更深刻的算术信息。

\textbf{证明}：对 \(f \in M_k\)，其常数项 \(a_0\) 定义了线性泛函。\(E_k\) 的常数项为 \(1\)（非零），而 \(S_k\) 由常数项为零的模形式构成。由 Petersson 内积的 Hermite 正定性，\(E_k\) 与 \(S_k\) 正交，且 \(\langle E_k, E_k \rangle\) 由 Rankin-Selberg 方法计算，非零。因此 \(M_k = \mathbb{C} E_k \oplus S_k\)。\(\blacksquare\)

\textbf{定理 86.2b}（Eisenstein 级数的函数方程与特殊值）：Eisenstein 级数 \(E_k(z)\) 在 \(k=2\) 处有特殊的非全纯修正。定义非全纯 Eisenstein 级数（实解析 Eisenstein 级数）： \[E(z, s) = \frac{1}{2} \sum_{(c,d)=1} \frac{\Im(z)^s}{|cz + d|^{2s}}\]

\(E(z, s)\) 在 \(\Re(s) > 1\) 时绝对收敛，可亚纯延拓到整个 \(\mathbb{C}\)，且满足函数方程 \(E(z, s) = E(z, 1-s)\)（由 Poisson 求和公式导出）。\(E(z, s)\) 在 \(s=1\) 处的留数与 Kronecker 极限公式有关，该公式将 \(\log|\Delta(z)|\) 与 \(\zeta'(0)\) 和 Dedekind \(\eta\) 函数联系起来。

\textbf{证明}：\(E(z, s)\) 的亚纯延拓利用 Selberg 迹公式或谱理论。\(E(z, s)\) 是 \(\mathbb{H}\) 上 Laplace 算子 \(\Delta = -y^2(\partial_x^2 + \partial_y^2)\) 的特征函数：\(\Delta E(z, s) = s(1-s)E(z, s)\)。函数方程由 \(E(z, s)\) 关于 \(s \leftrightarrow 1-s\) 的对称性导出，这源于 Poisson 求和公式中 \(y\) 和 \(1/y\) 的对称性。\(\blacksquare\)

\textbf{定义 86.3b}（Petersson 内积）：对 \(f, g \in S_k\)，Petersson 内积定义为： \[\langle f, g \rangle = \int_{\Gamma \setminus \mathbb{H}} f(z) \overline{g(z)} y^k \frac{dx dy}{y^2}\] 其中 \(z = x + iy\)。测度 \(\frac{dx dy}{y^2}\) 是 \(\mathbb{H}\) 上的 \(SL(2, \mathbb{R})\)-不变测度，而 \(y^k\) 因子使得被积函数 \(f(z)\overline{g(z)} y^k\) 在模变换下不变。Petersson 内积使 \(S_k\) 成为 Hilbert 空间，是 Hecke 算子自伴性和谱分解的基础。

Eisenstein 级数是模形式中最''简单''的一类------其 Fourier 系数是初等数论函数，但其解析性质（如函数方程、在 \(k=2\) 时的非全纯修正）揭示了自守形式理论的核心结构。Eisenstein 级数在 \(L\) 函数的特殊值理论中扮演关键角色：\(\zeta(k)\) 的值通过 \(E_k\) 的 Fourier 展开与 Bernoulli 数相联系。Eisenstein 级数的 Langlands 函子性提升（Langlands-Shahidi 方法）是研究一般约化群上 \(L\) 函数特殊值的核心工具。

\subsection{86.3 Hecke 算子与 L 函数}\label{hecke-ux7b97ux5b50ux4e0e-l-ux51fdux6570}

\textbf{定义 86.4}（Hecke 算子）：对 \(n \geq 1\)，权 \(k\) 的 \textbf{Hecke 算子} \(T_n\) 作用于模形式 \(f(z) = \sum_{m=0}^\infty a_m q^m\)：

\[T_n f(z) = n^{k-1} \sum_{d \mid n} d^{-k+1} \sum_{b=0}^{d-1} f\left(\frac{nz + bd}{d^2}\right)\]

其对 Fourier 系数的效果为 \(a_m(T_n f) = \sum_{d \mid (m,n)} d^{k-1} a_{mn/d^2}\)。特别地，\(a_1(T_n f) = a_n(f)\)，因此 \(T_n\) 的''第一 Fourier 系数''直接恢复 \(f\) 的第 \(n\) 个 Fourier 系数。

\textbf{定理 86.3}（Hecke 算子的基本性质）： 1. \(T_n\) 是 \(M_k\) 和 \(S_k\) 上的线性算子 2. 可交换性：\(T_m T_n = T_n T_m\)，且 \(T_m T_n = T_{mn}\) 当 \((m,n)=1\) 3. 自伴性：\(T_n\) 关于 Petersson 内积是自伴的（当 \((n,N)=1\)） 4. 若 \(f\) 是归一化 Hecke 特征形式（\(a_1 = 1\) 且为所有 \(T_n\) 的特征向量），则其 Fourier 系数满足 \(a_m a_n = \sum_{d \mid (m,n)} d^{k-1} a_{mn/d^2}\)（Hecke 乘法关系）

\textbf{证明}：(1) 由定义直接验证 \(T_n\) 保持模变换条件。(2) Hecke 算子的双陪集解释：\(T_n\) 对应于 \(\Gamma \begin{pmatrix} 1 & 0 \\ 0 & n \end{pmatrix} \Gamma\) 在模形式上的作用。双陪集的乘法性质导出可交换性和乘积公式。(3) 由 Petersson 内积的显式计算和 \(T_n\) 的伴随表示：\(\langle T_n f, g \rangle = \langle f, T_n g \rangle\)，利用 \(T_n\) 的双陪集表示和 Petersson 内积的不变性。(4) 由 Hecke 特征形式和 \(T_n\) 的 Fourier 系数公式，在 \(a_1 = 1\) 的归一化下，\(a_n = \lambda_n\)（\(T_n\) 的特征值），乘法性质由特征值的同态性和 \(T_n\) 的乘积公式导出。\(\blacksquare\)

\textbf{定理 86.3b}（Hecke 特征形式的存在性------谱定理）：\(S_k\) 存在一组正交基，由 Hecke 特征形式（即所有 \(T_n\) 的同时特征向量）组成。这是由 \(T_n\) 的自伴性和可交换性从有限维向量空间的谱定理得出的。

\textbf{定义 86.5}（Hecke L 函数）：对归一化 Hecke 特征形式 \(f = \sum a_n q^n\)，其 \textbf{Hecke L 函数} 定义为

\[L(f, s) = \sum_{n=1}^\infty \frac{a_n}{n^s} = \prod_{p \text{ prime}} \frac{1}{1 - a_p p^{-s} + p^{k-1-2s}}\]

此 Euler 乘积由 Hecke 算子的乘法性质导出。\(L(f, s)\) 在 \(\Re(s) > \frac{k+1}{2}\) 时绝对收敛（由 Ramanujan-Petersson 猜想 \(|a_p| \leq 2p^{(k-1)/2}\)，后由 Deligne 在 1974 年证明为 Weil 猜想的推论），可解析延拓为整个 \(\mathbb{C}\) 上的全纯函数，且满足函数方程

\[\Lambda(f, s) = (2\pi)^{-s} \Gamma(s) L(f, s) = (-1)^{k/2} \Lambda(f, k-s)\]

\textbf{证明}（框架）：\(L(f, s)\) 的 Euler 乘积由 \(a_n\) 的乘性 \(a_{mn} = a_m a_n\)（\((m,n)=1\)）和 Hecke 递推关系 \(a_{p^r} = a_p a_{p^{r-1}} - p^{k-1} a_{p^{r-2}}\) 导出。函数方程由 Mellin 变换和模形式的对称性导出：\(\Lambda(f, s) = \int_0^\infty f(iy) y^{s-1} dy\)，利用 \(f(-1/z) = z^k f(z)\) 可得 \(\Lambda(f, s) = \int_0^\infty f(iy) y^{k-s-1} dy = \Lambda(f, k-s)\)（\(f\) 为特征形式时符号为 \((-1)^{k/2}\)）。\(\blacksquare\)

\textbf{定理 86.3c}（Rankin-Selberg 方法）：对两个模形式 \(f, g\)，Rankin-Selberg 卷积 \(L(f \times g, s)\) 定义为 \(\sum a_n(f) \overline{a_n(g)} n^{-s}\)。利用 Eisenstein 级数 \(E(z, s)\) 的积分表示（展开 \(f\overline{g}\) 与 \(E(z, s)\) 的 Petersson 内积），\(L(f \times g, s)\) 可亚纯延拓，且与 \(f\) 和 \(g\) 的对称平方 L 函数有关。Rankin-Selberg 方法是研究 L 函数特殊值（如 \(L(f, k/2)\) 与周期积分的关系）的核心工具。

\subsection{86.4 模曲线与模形式}\label{ux6a21ux66f2ux7ebfux4e0eux6a21ux5f62ux5f0f}

\textbf{定义 86.6}（模曲线）：对 \(N \geq 1\)，\textbf{同余子群} \(\Gamma_0(N) = \{\begin{pmatrix} a & b \\ c & d \end{pmatrix} \in SL(2, \mathbb{Z}) : c \equiv 0 \pmod{N}\}\)。\textbf{模曲线} \(Y_0(N) = \Gamma_0(N) \setminus \mathbb{H}\) 是 \(\mathbb{H}\) 在 \(\Gamma_0(N)\) 作用下的商空间，可通过添加尖点紧化为 \(X_0(N) = \Gamma_0(N) \setminus (\mathbb{H} \cup \mathbb{P}^1(\mathbb{Q}))\)。\(X_0(N)\) 是 \(\mathbb{Q}\) 上的光滑射影代数曲线，其亏格由公式

\[g(X_0(N)) = 1 + \frac{\mu(N)}{12} - \frac{\nu_2(N)}{4} - \frac{\nu_3(N)}{3} - \frac{\nu_\infty(N)}{2}\]

给出，其中 \(\mu(N) = N \prod_{p \mid N} (1 + 1/p)\) 为 \(\Gamma_0(N)\) 在 \(\mathbb{H}\) 上的指数，\(\nu_2, \nu_3, \nu_\infty\) 分别为椭圆阶为 \(2\)、\(3\) 的点数和尖点数。

\textbf{定理 86.4}（模形式与模曲线的关系）：权 \(2\) 的 \(\Gamma_0(N)\)-尖点形式空间 \(S_2(\Gamma_0(N))\) 同构于 \(X_0(N)\) 上的全纯微分形式空间 \(H^0(X_0(N), \Omega^1)\)。更一般地，权 \(k\) 的模形式对应于 \(X_0(N)\) 上的 \(k\)-微分形式（带极点的截面）。

\textbf{证明}：对 \(f \in S_2(\Gamma_0(N))\)，\(f(z) dz\) 在模变换 \(\gamma \in \Gamma_0(N)\) 下不变：\(f(\gamma z) d(\gamma z) = (cz+d)^2 f(z) \cdot (cz+d)^{-2} dz = f(z) dz\)。且 \(f(z) dz\) 在尖点处全纯（因 \(f\) 为尖点形式，\(f(z) dz = \sum a_n q^n \cdot \frac{dq}{2\pi i q} = \frac{1}{2\pi i} \sum a_n q^{n-1} dq\)，\(a_0 = 0\) 保证全纯）。反之，\(X_0(N)\) 上的全纯微分形式拉回至 \(\mathbb{H}\) 上给出 \(S_2(\Gamma_0(N))\) 中的元素。\(\blacksquare\)

模曲线是模形式理论的几何化------它将模形式这一''分析/数论''对象转化为代数曲线上的''几何''对象。这一几何化是理解模性定理（Taniyama-Shimura-Weil 猜想）的关键：权 \(2\) 的 Hecke 特征形式 \(f\) 对应 \(X_0(N)\) 的 Jacobi 簇 \(J_0(N)\) 的一个简单商（Abel 簇 \(A_f\)），而 \(A_f\) 与椭圆曲线的关系由模性定理给出。模曲线 \(X_0(N)\) 上的有理点、Heegner 点和 Heegner 除子理论是 Gross-Zagier 定理和 Kolyvagin 的 Euler 系方法的核心------这些工具将 \(L(f, 1)\) 的导数值与 Heegner 点的高度联系起来，是 BSD 猜想在秩 \(1\) 情形下取得进展的关键。

\textbf{定理 86.4b}（Hecke 代数与模性的 Galois 表示）：设 \(f \in S_2(\Gamma_0(N))\) 为 Hecke 特征形式，系数 \(a_n \in \mathbb{Z}\)。则存在 \(2\) 维 \(\ell\)-进 Galois 表示 \(\rho_f: G_{\mathbb{Q}} \to \operatorname{GL}_2(\overline{\mathbb{Q}}_\ell)\)，使得对几乎所有素数 \(p\)，\(\operatorname{tr} \rho_f(\operatorname{Frob}_p) = a_p\) 且 \(\det \rho_f(\operatorname{Frob}_p) = p\)。这一表示由 \(J_0(N)\) 的 Tate 模构造，是模性定理证明的核心对象。

\subsection{86.5 Jacobi 形式与 Siegel 模形式}\label{jacobi-ux5f62ux5f0fux4e0e-siegel-ux6a21ux5f62ux5f0f}

\textbf{定义 86.7}（Jacobi 形式）：Jacobi 形式（Eichler-Zagier 1985）是模形式和 theta 函数的''混合''，同时具有关于模群 \(SL(2, \mathbb{Z})\) 的模变换性质和关于整数平移的椭圆变换性质。对权 \(k\) 和指标 \(m\)，Jacobi 形式 \(\phi(\tau, z)\)（\(\tau \in \mathbb{H}, z \in \mathbb{C}\)）满足：

\[\phi\left(\frac{a\tau + b}{c\tau + d}, \frac{z}{c\tau + d}\right) = (c\tau + d)^k e^{2\pi i m c z^2/(c\tau + d)} \phi(\tau, z)\]

\[\phi(\tau, z + \lambda\tau + \mu) = e^{-2\pi i m(\lambda^2\tau + 2\lambda z)} \phi(\tau, z)\]

其中 \(\lambda, \mu \in \mathbb{Z}\)。Jacobi 形式空间记为 \(J_{k,m}\)，是有限维的。Jacobi 形式有 Fourier 展开 \(\phi(\tau, z) = \sum_{n, r} c(n, r) e^{2\pi i (n\tau + rz)}\)，其中 \(c(n, r)\) 仅依赖于判别式 \(D = 4mn - r^2\)。

\textbf{定理 86.5}（Jacobi 形式与 theta 函数的联系）：最基本的 Jacobi 形式是 Jacobi theta 函数： \[\vartheta(\tau, z) = \sum_{n=-\infty}^{\infty} e^{\pi i n^2 \tau + 2\pi i n z}\]

\(\vartheta(\tau, z)\) 属于 \(J_{1/2, 1/2}\)（半整数权）。Jacobi 三重积恒等式给出 \(\vartheta(\tau, z) = \prod_{n=1}^\infty (1 - q^n)(1 + q^{n-1/2}\zeta)(1 + q^{n-1/2}\zeta^{-1})\)，其中 \(q = e^{2\pi i \tau}\)，\(\zeta = e^{2\pi i z}\)。

Jacobi 形式在弦论（椭圆亏格）、仿射 Lie 代数的特征标理论（Kac-Weyl 特征标公式）和 Siegel 模形式的 Fourier-Jacobi 展开中有关键应用。仿射 Lie 代数 \(\hat{\mathfrak{g}}\) 的不可约表示的特征标是 Jacobi 形式（乘以 \(\eta\) 函数的某次幂），这一联系是 Kac-Peterson 和 Kac-Wakimoto 的模不变性定理的基础。

\textbf{定义 86.8}（Siegel 模形式）：Siegel 模形式是模形式在更高亏格下的推广。Siegel 上半空间 \(\mathbb{H}_g = \{\Omega \in \operatorname{Sym}_g(\mathbb{C}) : \Im(\Omega) > 0\}\)（亏格 \(g\) 的 Siegel 上半空间，即 \(g \times g\) 复对称矩阵且虚部正定）上的 Siegel 模形式 \(F(\Omega)\) 满足关于辛群 \(Sp(2g, \mathbb{Z})\) 的模变换条件：

\[F((A\Omega + B)(C\Omega + D)^{-1}) = \det(C\Omega + D)^k F(\Omega)\]

其中 \(\begin{pmatrix} A & B \\ C & D \end{pmatrix} \in Sp(2g, \mathbb{Z})\)。Siegel 模形式的 Fourier 系数编码了 \(g\) 维 Abel 簇的算术信息，其理论是代数几何中模空间（如 \(\mathcal{A}_g\)------\(g\) 维极化 Abel 簇的模空间）和算术几何中 Siegel 模簇的基础。

\textbf{定理 86.6}（Siegel 模形式的 Fourier-Jacobi 展开）：亏格 \(g\) 的 Siegel 模形式 \(F(\Omega)\) 可沿 \(\Omega = \begin{pmatrix} \tau & z \\ z^T & \Omega' \end{pmatrix}\) 展开为 Fourier-Jacobi 级数，其中系数是亏格 \(g-1\) 的 Jacobi 形式。这一展开是研究 Siegel 模形式的结构和构造新的 Siegel 模形式（如 Ikeda 提升）的关键工具。

Siegel 模形式在弦论（IIB 型弦论中的 \(SL(2, \mathbb{Z})\) S-对偶性）和黑洞微观态计数（OSV 猜想，将黑洞熵与 Siegel 模形式的 Fourier 系数联系）中也有重要应用。在数学中，Siegel 模形式与自守表示理论、Galois 表示和 Langlands 纲领的函子性紧密相关。

\textbf{定理 86.7}（Theta 对应与 Weil 表示）：symplectic 群和 ortogonal 群之间的 theta 对应（Howe 对偶）建立了 Siegel 模形式和 ortogonal 群上的自守形式之间的对应。这一对应是 Langlands 纲领中函子性提升的核心实例，在数论中将椭圆曲线（或更一般的 motive）的 L 函数与高秩群上的自守 L 函数联系起来。

\textbf{定义 86.9}（Hilbert 模形式与 Bianchi 模形式）：Hilbert 模形式是模形式在全实域上的推广，用于研究 \(\operatorname{GL}_2\) 在数域上的 Langlands 对应。Bianchi 模形式是虚二次域上的模形式，对应 \(\operatorname{GL}_2\) 在虚二次域上的自守形式。这两类模形式是将模性定理推广到一般数域上的关键对象。

\textbf{定理 86.8}（Serre 猜想与模性提升定理）：Serre 猜想（由 Khare-Wintenberger 于 2008 年证明）断言：\(\mathbb{F}_\ell\) 上的任意奇、不可约 Galois 表示 \(\bar{\rho}: G_{\mathbb{Q}} \to \operatorname{GL}_2(\overline{\mathbb{F}}_\ell)\) 均来自某个模形式（即存在权 \(k\)、水平 \(N\) 的模形式 \(f\) 使得 \(\bar{\rho}_f \cong \bar{\rho}\)）。Serre 猜想是模性提升定理的理论基础，其证明结合了模性提升技术在 \(\ell = 2\) 和奇素数上的发展。

\(\blacksquare\)

\hrulefill

\section{第296章：解析数论}\label{ux7b2c296ux7ae0ux89e3ux6790ux6570ux8bba}

解析数论用复分析（V9）和 Fourier 分析（V8）的工具研究整数问题。本章的核心是 Dirichlet 级数、Riemann zeta 函数和素数定理的完整证明。

解析数论的本质在于：将离散的算术问题转化为连续的分析问题，通过研究复变函数的解析性质（极点、零点、渐近行为）来获得关于整数分布的深刻结果。这一方法始于 Euler 对 zeta 函数的研究，经 Dirichlet、Riemann、Hadamard、de la Vallee-Poussin 等人的发展，在 20 世纪由 Hardy-Littlewood 圆法和 Vinogradov 三角和方法的推动下达到顶峰。现代解析数论涵盖模形式、自守形式、L-函数和 Langlands 纲领等方向，与代数几何、表示论和调和分析深度融合。

\subsection{85.1 算术函数与 Dirichlet 卷积}\label{ux7b97ux672fux51fdux6570ux4e0e-dirichlet-ux5377ux79ef}

\textbf{定义 85.1}（算术函数）：\textbf{算术函数}是 \(f: \mathbb{N}^+ \to \mathbb{C}\) 的函数。常用的算术函数包括： - \(\mathbf{1}(n) = 1\)（常数函数） - \(\operatorname{id}(n) = n\)（恒等函数） - \(\mu(n)\)：Mobius 函数（\(\mu(1) = 1\)；若 \(n\) 有平方因子则 \(\mu(n) = 0\)；否则 \(\mu(n) = (-1)^k\)，\(k\) 为素因子个数） - \(\varphi(n)\)：Euler \(\varphi\) 函数（Euler totient 函数，\(\varphi(n) = n \prod_{p \mid n} (1 - 1/p)\)） - \(\Lambda(n)\)：von Mangoldt 函数（\(\Lambda(n) = \log p\) 若 \(n = p^k\)；否则 \(\Lambda(n) = 0\)） - \(\tau(n)\)：除数函数（\(n\) 的正因子个数） - \(\sigma(n)\)：除数和函数（\(n\) 的正因子之和）

\textbf{定义 85.1'（积性函数）}：算术函数 \(f\) 称为\textbf{积性函数}，如果 \(f(mn) = f(m)f(n)\) 对所有 \(\gcd(m,n)=1\) 成立。若此关系对所有 \(m,n\) 成立（不要求互素），则 \(f\) 称为\textbf{完全积性函数}。积性函数类构成了 Dirichlet 卷积代数中的''乘法群''的核心。上述函数中，\(\mu\) 和 \(\varphi\) 是积性函数，\(\mathbf{1}\) 和 \(\operatorname{id}\) 是完全积性函数。

\textbf{定义 85.2}（Dirichlet 卷积）：两个算术函数 \(f, g\) 的 \textbf{Dirichlet 卷积} 为

\[(f * g)(n) = \sum_{d \mid n} f(d) g(n/d)\]

\textbf{命题 85.1}：Dirichlet 卷积满足： 1. 结合律：\((f * g) * h = f * (g * h)\) 2. 交换律：\(f * g = g * f\) 3. 单位元 \(\delta\)（\(\delta(1) = 1\)，\(\delta(n) = 0\) 对 \(n > 1\)） 4. \(f\) 可逆当且仅当 \(f(1) \neq 0\)

\textbf{证明}：1 和 2 由和号重排验证。3 显然。4：构造 \(g(n) = -\frac{1}{f(1)} \sum_{d \mid n, d < n} f(n/d) g(d)\) 归纳定义。\(\blacksquare\)

\textbf{命题 85.2}（Mobius 反演）：\(\mathbf{1} * \mu = \delta\)。即 \(\sum_{d \mid n} \mu(d) = 0\)（\(n > 1\)）。等价地，\(f(n) = \sum_{d \mid n} g(d)\) 当且仅当 \(g(n) = \sum_{d \mid n} \mu(d) f(n/d)\)。

\textbf{证明}：\(\mathbf{1} * \mu\) 在素数幂 \(p^k\) 上的值为 \(1 - 1 = 0\)（\(k \geq 1\)）。由积性，一般 \(n\) 也成立。\(\blacksquare\)

\textbf{命题 85.2.1（von Mangoldt 函数与卷积）}：\(\Lambda * \mathbf{1} = \log\)。即 \(\sum_{d \mid n} \Lambda(d) = \log n\)。这建立了 von Mangoldt 函数与对数函数之间的基础关系，是素数定理证明中 Chebyshev 函数 \(\psi(x) = \sum_{n \leq x} \Lambda(n)\) 与 \(\log\) 求和关系的关键。

\textbf{证明}：若 \(n = \prod p_i^{e_i}\)，则 \(\log n = \sum e_i \log p_i\)。而 \(\sum_{d \mid n} \Lambda(d) = \sum_{p^k \mid n} \log p\)，其中 \(p^k \mid n\) 当且仅当 \(k \leq e_i\)。故 \(\sum_{d \mid n} \Lambda(d) = \sum_i e_i \log p_i = \log n\)。\(\blacksquare\)

\subsection{85.2 Dirichlet 级数与 Euler 乘积}\label{dirichlet-ux7ea7ux6570ux4e0e-euler-ux4e58ux79ef}

\textbf{定义 85.3}（Dirichlet 级数）：算术函数 \(f\) 的 \textbf{Dirichlet 级数} 为

\[D_f(s) = \sum_{n=1}^\infty \frac{f(n)}{n^s}\]

其中 \(s = \sigma + it \in \mathbb{C}\)。若 \(f(n) = O(n^\alpha)\)，则 \(D_f(s)\) 在 \(\Re(s) > \alpha + 1\) 内绝对收敛。\(D_f(s)\) 的收敛半平面由其绝对收敛横坐标 \(\sigma_a = \inf\{\sigma : \sum |f(n)| n^{-\sigma} < \infty\}\) 决定。

\textbf{命题 85.3}（Dirichlet 卷积与级数乘积）：\((f * g)(n)\) 的 Dirichlet 级数是 \(D_f(s) D_g(s)\)。

\textbf{证明}：\(D_f(s) D_g(s) = \sum_{a=1}^\infty \frac{f(a)}{a^s} \sum_{b=1}^\infty \frac{g(b)}{b^s} = \sum_{n=1}^\infty \frac{\sum_{ab=n} f(a) g(b)}{n^s} = D_{f*g}(s)\)。\(\blacksquare\)

\textbf{定义 85.4}（Euler 乘积）：若 \(f\) 是\textbf{积性函数}（\(f(mn) = f(m)f(n)\) 当 \(\gcd(m, n) = 1\)），则

\[D_f(s) = \prod_{p \text{ 素数}} \left(1 + \frac{f(p)}{p^s} + \frac{f(p^2)}{p^{2s}} + \cdots\right)\]

若 \(f\) 是完全积性函数，则 \(D_f(s) = \prod_p \left(1 - \frac{f(p)}{p^s}\right)^{-1}\)。

Euler 乘积公式是算术函数 Dirichlet 级数的核心结构：它将一个无穷求和分解为所有素数处的无穷乘积，建立了解析性质与算术性质之间的桥梁。Euler 乘积的存在性等价于算术函数的积性，而 Euler 乘积的解析性质（如极点、零点）则反映了素数分布的信息。

\subsection{85.3 Riemann zeta 函数}\label{riemann-zeta-ux51fdux6570}

\textbf{定义 85.5}（Riemann zeta 函数）：对 \(\Re(s) > 1\)，

\[\zeta(s) = \sum_{n=1}^\infty \frac{1}{n^s}\]

\textbf{定理 85.4}（Euler 乘积公式）：对 \(\Re(s) > 1\)，

\[\zeta(s) = \prod_{p \text{ 素数}} \frac{1}{1 - p^{-s}}\]

\textbf{证明}：\(\mathbf{1}\) 是完全积性函数，由 Euler 乘积公式即得。\(\blacksquare\)

\textbf{定理 85.5}（\(\zeta(s)\) 的解析延拓）：\(\zeta(s)\) 可解析延拓为整个复平面上的亚纯函数，仅在 \(s = 1\) 处有一个单极点（留数为 \(1\)）。函数方程：

\[\zeta(s) = 2^s \pi^{s-1} \sin\left(\frac{\pi s}{2}\right) \Gamma(1-s) \zeta(1-s)\]

或等价地，令 \(\xi(s) = \frac{1}{2} s(s-1) \pi^{-s/2} \Gamma(s/2) \zeta(s)\)，则 \(\xi(s) = \xi(1-s)\) 且 \(\xi(s)\) 是整函数。

\textbf{证明}：对 \(\Re(s) > 1\)，利用 \(\Gamma(s/2)\pi^{-s/2}n^{-s} = \int_0^\infty e^{-\pi n^2 t} t^{s/2-1} dt\) 和 Poisson 求和公式 \(\sum_{n=-\infty}^\infty e^{-\pi n^2 t} = t^{-1/2} \sum_{n=-\infty}^\infty e^{-\pi n^2/t}\)。定义 \(\theta(t) = \sum_{n=-\infty}^\infty e^{-\pi n^2 t} = 1 + 2\sum_{n=1}^\infty e^{-\pi n^2 t}\)，则 \(\theta(1/t) = \sqrt{t} \theta(t)\)。通过 Mellin 变换 \(\Gamma(s/2)\pi^{-s/2}\zeta(s) = \int_0^\infty \frac{\theta(t)-1}{2} t^{s/2-1} dt\)。将积分分为 \([0,1]\) 和 \([1,\infty)\) 两部分，在 \([0,1]\) 上代入 \(\theta\) 的函数方程，得到 \(\pi^{-s/2}\Gamma(s/2)\zeta(s) = \frac{1}{s(s-1)} + \int_1^\infty (t^{s/2} + t^{(1-s)/2}) \frac{\theta(t)-1}{2t} dt\)。右端积分对所有 \(s \in \mathbb{C}\) 收敛，给出了 \(\zeta(s)\) 的亚纯延拓。函数方程来自 \(\xi(s)\) 定义中对称性。\(\blacksquare\)

\textbf{定义 85.6}（\(\zeta(s)\) 的零点）： - \textbf{平凡零点}：\(s = -2, -4, -6, \ldots\)（由 \(\sin(\pi s/2)\) 项产生） - \textbf{非平凡零点}：\(0 < \Re(s) < 1\)（临界带）内的零点 - \textbf{Riemann 猜想}（未证明）：\(\zeta(s)\) 的所有非平凡零点都满足 \(\Re(s) = 1/2\)（临界线）

\textbf{定理 85.5.1（零点计数）}：设 \(N(T)\) 为 \(\zeta(s)\) 在临界带 \(0 < \Re(s) < 1\)，\(0 < \Im(s) < T\) 内的非平凡零点个数（计重数）。则 \[N(T) = \frac{T}{2\pi} \log \frac{T}{2\pi} - \frac{T}{2\pi} + O(\log T)\]

\textbf{证明}：利用辐角原理：\(N(T) = \frac{1}{2\pi i} \oint_C \frac{\xi'(s)}{\xi(s)} ds\)，其中 \(C\) 为包围 \(0 < \Re(s) < 1\)，\(0 < \Im(s) < T\) 区域的围道。计算 \(\xi(s)\) 的渐近展开和 \(\Gamma\) 函数的 Stirling 公式得到所述估计。\(\blacksquare\)

\subsection{85.4 素数定理}\label{ux7d20ux6570ux5b9aux7406}

\textbf{定理 85.6}（素数定理）：设 \(\pi(x) = \#\{p \leq x : p \text{ 素数}\}\)。则

\[\pi(x) \sim \frac{x}{\log x} \quad (x \to \infty)\]

即 \(\lim_{x \to \infty} \frac{\pi(x)}{x / \log x} = 1\)。

\textbf{证明}：引入 Chebyshev 函数 \(\psi(x) = \sum_{n \leq x} \Lambda(n) = \sum_{p^k \leq x} \log p\)。由部分求和，\(\pi(x) \sim x/\log x\) 等价于 \(\psi(x) \sim x\)。使用 Perron 公式将 \(\psi(x)\) 表示为 \(\zeta(s)\) 的围道积分：\(\psi(x) = \frac{1}{2\pi i} \int_{c-i\infty}^{c+i\infty} \left(-\frac{\zeta'(s)}{\zeta(s)}\right) \frac{x^s}{s} ds\)（\(c > 1\)）。将积分围道平移至 \(\Re(s) = \sigma\)（\(0 < \sigma < 1\)），通过估计 \(\zeta(s)\) 在 \(\Re(s) = 1\) 附近的非零性（Hadamard-de la Vallee-Poussin 的零自由区域 \(|\Im(s)| \leq C/\log |\Im(s)|\) 附近的 \(\Re(s) \geq 1 - c/\log |\Im(s)|\)），得到 \(\psi(x) = x + O(x e^{-c\sqrt{\log x}})\)，从而 \(\pi(x) = \operatorname{Li}(x) + O(x e^{-c\sqrt{\log x}})\)，其中 \(\operatorname{Li}(x) = \int_2^x dt/\log t \sim x/\log x\)。\(\blacksquare\)

\textbf{定理 85.7}（素数定理的误差项）：若 Riemann 猜想成立，则 \(\pi(x) = \operatorname{Li}(x) + O(\sqrt{x} \log x)\)，其中 \(\operatorname{Li}(x) = \int_2^x \frac{dt}{\log t}\)。

\textbf{证明}：由显式公式 \(\psi(x) = x - \sum_\rho x^\rho/\rho - \log(2\pi) - \frac{1}{2}\log(1 - x^{-2})\)（跳过 trivial zeros 的项），其中 \(\rho\) 遍历 \(\zeta(s)\) 的非平凡零点。若 Riemann 猜想成立，则 \(\Re(\rho) = 1/2\)，故 \(|x^\rho| = x^{1/2}\)。零点计数的估计 \(N(T) \sim \frac{T}{2\pi}\log T\) 给出 \(\sum_{|\Im(\rho)| \leq T} 1/|\rho| \sim \frac{1}{2\pi}(\log T)^2\)。代入显式公式，使用部分积分将 \(\psi(x)\) 转换为 \(\pi(x)\) 并利用 \(\operatorname{Li}(x) = x/\log x + x/(\log x)^2 + \cdots\)，即得所述误差界。\(\blacksquare\)

\textbf{定理 85.8}（等差数列中的素数定理，Dirichlet）：若 \(\gcd(a, q) = 1\)，则等差数列 \(a, a+q, a+2q, \ldots\) 中包含无穷多个素数。且

\[\pi(x; q, a) \sim \frac{1}{\varphi(q)} \frac{x}{\log x} \quad (x \to \infty)\]

\textbf{证明}：引入 Dirichlet 特征 \(\chi \bmod q\) 和 Dirichlet \(L\)-函数 \(L(s, \chi) = \sum_{n=1}^\infty \chi(n) n^{-s}\)。对非主特征 \(\chi\)，\(L(1, \chi) \neq 0\)（这是证明的关键，由 Dirichlet 的类数公式保证）。通过正交性 \(\frac{1}{\varphi(q)}\sum_{\chi \bmod q} \overline{\chi}(a) \chi(n) = \mathbf{1}_{\{n \equiv a \bmod q\}}\)，将 \(\psi(x; q, a) = \sum_{n \leq x, n \equiv a \bmod q} \Lambda(n)\) 表示为 \(L\)-函数的对数导数的 Perron 积分。由 \(L(1, \chi) \neq 0\) 得出 \(\zeta(s)\) 在 \(s=1\) 处的极点贡献主项 \(\frac{1}{\varphi(q)} x\)，而其他项的贡献为 \(o(x)\)。\(\blacksquare\)

\subsection{85.5 堆垒数论}\label{ux5806ux5792ux6570ux8bba}

堆垒数论（Additive Number Theory）研究将整数表示为特定类型整数之和的经典问题。两个核心问题是：（1）\textbf{Goldbach猜想}：每个大于2的偶数可表为两个素数之和。陈景润于1966年证明了''1+2''------每个充分大的偶数可表为一个素数和一个至多两个素数乘积的数之和，这是迄今为止最好的结果。Helfgott于2013年完整证明了三素数定理：每个大于5的奇数可表为三个素数之和。（2）\textbf{Waring问题}（1770年）：对每个正整数\(k\)，存在\(g(k)\)使得每个正整数可表为至多\(g(k)\)个\(k\)次幂之和。Lagrange四平方和定理即\(g(2)=4\)。记\(G(k)\)为充分大整数所需的最少\(k\)次幂个数，已知\(G(2)=4\)，而\(G(k)\)的确切值对\(k \geq 3\)仍是重大未解问题。处理这类问题的核心工具是\textbf{圆法}（Hardy-Littlewood-Ramanujan，1920年代），其基本思想是将计数函数表示为单位圆上的积分，利用有理点附近的主项渐近展开与余项的三角和估计得到结果。

圆法的核心公式基于正交性：指示函数 \(\mathbf{1}_{\{n = m\}}\) 可表示为 \(\int_0^1 e^{2\pi i (n-m)\alpha} d\alpha\)。因此，将 \(n\) 表示为 \(s\) 个素数之和的表示个数 \(R_s(n)\) 可写为 \(R_s(n) = \int_0^1 f(\alpha)^s e^{-2\pi i n\alpha} d\alpha\)，其中 \(f(\alpha) = \sum_{p \leq n} e^{2\pi i p\alpha}\) 为素数指数和。圆法将积分区间 \([0,1]\) 划分为\textbf{主弧}（major arcs）和\textbf{余弧}（minor arcs）：主弧由分母较小的有理点 \(\alpha = a/q\)（\(q\) 小）的邻域组成，其上 \(f(\alpha)\) 的渐近行为由 Siegel-Walfisz 定理和素数在算术级数中的分布精确控制；余弧上利用 Vinogradov 的三角和估计 \(|f(\alpha)| \ll n (\log n)^{-A}\)（\(A\) 任意大）证明其贡献可忽略。圆法在 Waring 问题中同样适用：将 \(f(\alpha)\) 替换为 \(g(\alpha) = \sum_{x=1}^N e^{2\pi i x^k \alpha}\)，通过 Weyl 不等式和 Vinogradov 中值定理控制余弧上的贡献。Hardy 和 Littlewood 发展的圆法，后经 Vinogradov、Davenport 和 Vaughan 等人的改进，已成为堆垒数论中最强大的分析工具。

\textbf{定理 85.5.1（Vinogradov 三素数定理，1937）}：每个充分大的奇数可表为三个素数之和。Helfgott (2013) 将''充分大''改进为对所有大于 5 的奇数成立。

\textbf{证明}：设 \(N\) 为奇数。\(R_3(N) = \sum_{p_1+p_2+p_3=N} 1 = \int_0^1 f(\alpha)^3 e^{-2\pi i N\alpha} d\alpha\)。主弧上 \(f(\alpha) \sim \frac{\mu(q)}{\varphi(q)} \operatorname{Li}(N)\)，积分贡献 \(\sim \frac{1}{2} \mathfrak{S}(N) N^2/(\log N)^3\)，其中 \(\mathfrak{S}(N) = \prod_{p \mid N} (1 - 1/(p-1)^2) \prod_{p \nmid N} (1 + 1/(p-1)^3)\) 是奇异级数（singular series）。余弧上 \(|f(\alpha)| \ll N (\log N)^{-A}\)（\(A\) 任意大），故余弧贡献可忽略。\(\mathfrak{S}(N) > 0\) 对所有奇数 \(N\) 成立，故 \(R_3(N) > 0\) 对充分大 \(N\)。\(\blacksquare\)

\subsection{85.6 连分数理论}\label{ux8fdeux5206ux6570ux7406ux8bba}

实数\(x\)的\textbf{简单连分数}展开为

\[x = a_0 + \cfrac{1}{a_1 + \cfrac{1}{a_2 + \cdots}} = [a_0; a_1, a_2, \ldots]\]

其中\(a_0 \in \mathbb{Z}\)，\(a_i \in \mathbb{N}^+\)（\(i \geq 1\)）。若展开有限，则\(x \in \mathbb{Q}\)；反之亦然。连分数展开的构造算法：\(a_0 = \lfloor x \rfloor\)，\(x_1 = 1/(x - a_0)\)，\(a_1 = \lfloor x_1 \rfloor\)，以此类推。

\textbf{Lagrange定理}（1770年）：二次无理数（有理系数二次方程的根）的连分数展开是最终周期的，这一性质刻画了二次无理数。收敛子\(p_n/q_n = [a_0; a_1, \ldots, a_n]\)是\(x\)的\textbf{最佳有理逼近}，满足\(|q x - p| < |q' x - p'|\)对所有分母\(q' \leq q\)且\(p'/q' \neq p_n/q_n\)成立。连分数与\textbf{Pell方程}\(x^2 - Dy^2 = \pm 1\)的自然联系是：\(\sqrt{D}\)的连分数周期\(r\)给出了基本解\((x_1, y_1) = (p_{r-1}, q_{r-1})\)（当\(r\)为偶数时）或\((p_{2r-1}, q_{2r-1})\)（当\(r\)为奇数时）。

连分数理论在 Diophantine 逼近中具有核心地位。收敛子 \(p_n/q_n\) 满足递推公式 \(p_n = a_n p_{n-1} + p_{n-2}\)，\(q_n = a_n q_{n-1} + q_{n-2}\)（初始条件 \(p_{-1}=1, p_0=a_0\)，\(q_{-1}=0, q_0=1\)）。由此可导出基本不等式 \(\frac{1}{q_n(q_n + q_{n+1})} < \left|x - \frac{p_n}{q_n}\right| < \frac{1}{q_n q_{n+1}}\)，表明收敛子以极快的速度逼近 \(x\)。\textbf{Hurwitz 定理}（1891年）进一步断言：每个无理数有无穷多个有理逼近 \(p/q\) 满足 \(|x - p/q| < 1/(\sqrt{5} q^2)\)，且常数 \(\sqrt{5}\) 是最优的（以黄金比例 \(\phi = (1+\sqrt{5})/2\) 为极值情形）。连分数还与\textbf{Roth 定理}（1955年，Fields 奖成果）密切相关：代数数 \(\alpha\) 的逼近指数为 \(2\)------即对任意 \(\varepsilon > 0\)，不等式 \(|\alpha - p/q| < q^{-(2+\varepsilon)}\) 仅有有限多组解。\(\sqrt{5}\) 在 Hurwitz 定理中的出现本质上是由于 \(\phi\) 的连分数 \([1;1,1,\ldots]\) 给出了最慢的收敛速度。在密码学中，连分数用于 Wiener 对 RSA 的攻击：若 RSA 私钥 \(d < N^{1/4}/3\)，则 \(e/N\) 的连分数展开可高效恢复 \(d\)（定理 259.2）。

\subsection{85.7 数的几何}\label{ux6570ux7684ux51e0ux4f55}

\textbf{数的几何}（Geometry of Numbers）由Minkowski于1896年创立，用几何方法研究格点分布与数论问题。设\(\Lambda \subseteq \mathbb{R}^n\)是\textbf{格}（discrete subgroup），即\(\Lambda = \{\sum a_i v_i : a_i \in \mathbb{Z}\}\)，其协体积\(\operatorname{covol}(\Lambda) = |\det(v_1, \ldots, v_n)|\)。\textbf{Minkowski第一定理}：若\(K \subseteq \mathbb{R}^n\)是中心对称凸体且\(\operatorname{vol}(K) > 2^n \operatorname{covol}(\Lambda)\)，则\(K \cap (\Lambda \setminus \{0\}) \neq \varnothing\)。\textbf{Minkowski第二定理}：对中心对称凸体\(K\)，其连续极小\(\lambda_1 \leq \cdots \leq \lambda_n\)满足\(\frac{2^n}{n!} \operatorname{covol}(\Lambda) \leq \lambda_1 \cdots \lambda_n \cdot \operatorname{vol}(K) \leq 2^n \operatorname{covol}(\Lambda)\)。经典应用包括：（1）Dirichlet单位定理的几何证明------对数嵌入\(\ell: \mathcal{O}_K^\times \to \mathbb{R}^{r_1+r_2}\)将单位群映为秩\(r_1+r_2-1\)的格；（2）Minkowski判别式界------\(|\operatorname{disc}(K)| \geq (\pi/4)^{2r_2} (n^n/n!)^2\)，由此可推出\(\mathbb{Q}\)不存在非平凡的无处分歧扩张。

\textbf{定理 85.7.1（Minkowski 第一定理的证明）}：考虑 \(\frac{1}{2}K = \{\frac{1}{2}x : x \in K\}\)。若 \(\frac{1}{2}K \cap \Lambda\) 有非零元，则 \(K \cap \Lambda\) 也有非零元。若 \(\frac{1}{2}K\) 在 \(\mathbb{R}^n/\Lambda\) 中的像（商映射下）是单射，则 \(\operatorname{vol}(\frac{1}{2}K) \leq \operatorname{covol}(\Lambda)\)，即 \(\operatorname{vol}(K) \leq 2^n \operatorname{covol}(\Lambda)\)。故若 \(\operatorname{vol}(K) > 2^n \operatorname{covol}(\Lambda)\)，则 \(\frac{1}{2}K\) 必有两个不同的点在模 \(\Lambda\) 下等价，其差即 \(\Lambda\) 中的非零元。\(\blacksquare\)

\textbf{定理 85.7.2（Minkowski 第二定理）}：连续极小 \(\lambda_i\) 定义为 \(\lambda_i = \inf\{\lambda > 0 : \dim(\operatorname{span}(\lambda K \cap \Lambda)) \geq i\}\)。存在 \(\Lambda\) 的基 \(v_1, \ldots, v_n\) 使得 \(\|v_i\| \leq 2\lambda_i\)（Mahler 约化基）。利用体积估计和 Hadamard 不等式导出所需的不等式。

\textbf{定理 85.7.3（Minkowski 判别式界）}：设 \(K\) 是 \(\mathbb{Q}\) 的 \(n\) 次扩张。则 \(|\operatorname{disc}(K)|^{1/2} \geq \frac{n^n}{n!} \left(\frac{\pi}{4}\right)^{r_2}\)，其中 \(r_2\) 是 \(K\) 的复嵌入对数。由此推出 \(|\operatorname{disc}(K)| > 1\) 对 \(n > 1\)，即 \(\mathbb{Q}\) 不存在非平凡的无处分歧扩张（Minkowski 定理的推论，比 Hermite 定理更强）。

\textbf{证明}：将整数环 \(\mathcal{O}_K\) 嵌入 \(\mathbb{R}^{r_1} \times \mathbb{C}^{r_2}\) 中，\(\mathcal{O}_K\) 的像是 \(\mathbb{R}^n\) 中的格，其协体积为 \(2^{-r_2} |\operatorname{disc}(K)|^{1/2}\)。取 \(K\) 为 \(\{x \in \mathbb{R}^{r_1} \times \mathbb{C}^{r_2} : \sum |x_i| \leq n\}\)（或更精细的凸体），应用 Minkowski 第一定理得 \(|\operatorname{disc}(K)|^{1/2} \geq 2^{-r_2} \cdot \operatorname{vol}(K) \cdot 2^{-n}\)。选择适当的 \(K\) 使 \(\operatorname{vol}(K)\) 最大化，得到所述下界。\(\blacksquare\)

\subsection{85.8 筛法}\label{ux7b5bux6cd5}

\textbf{定义 85.7}（筛法）：筛法是估计整数集合中满足某种算术条件的元素个数的方法。最经典的筛子是 \textbf{Eratosthenes 筛法}：不超过 \(x\) 的素数个数 \(\pi(x)\) 可通过筛去不超过 \(\sqrt{x}\) 的素数的倍数来估计，但直接使用包含-排除原理则余项过大。现代筛法采用''权重''和''组合恒等式''来控制误差。

\textbf{定理 85.8.1（Brun 筛法，1919）}：Brun 引入''双变量筛法''（Brun 纯筛法），利用不等式 \(\sum_{d \mid P(z)} \mu(d) \leq \sum_{d \in \mathcal{D}^+} \mu(d)\) 和 \(\sum_{d \mid P(z)} \mu(d) \geq \sum_{d \in \mathcal{D}^-} \mu(d)\)（其中 \(\mathcal{D}^\pm\) 是适当选取的除数集合），证明了：孪生素数（相差为 2 的素数对）的倒数之和收敛（Brun 常数），即 \(\sum_{p, p+2 \text{ 均为素数}} (1/p + 1/(p+2)) < \infty\)。

\textbf{证明}（框架）：设 \(A = \{n(n+2) : n \leq x\}\)。Eratosthenes 筛法中的余项导致似然估计过松。Brun 通过引入截断函数 \(\chi(n) = \sum_{d \mid n, d \in \mathcal{D}} \mu(d)\)（\(\mathcal{D}\) 为不超过 \(z\) 的除数集合），将筛法转换为组合恒等式。通过选择参数 \(z\) 和截断条件，使主项保持而余项可控。\(\blacksquare\)

\textbf{定理 85.8.2（陈景润定理，1966）}：每个充分大的偶数可表为 \(p + P_2\)，其中 \(p\) 是素数，\(P_2\) 是至多两个素数之积。这是 Goldbach 猜想最接近的结果（``1+2''）。

\textbf{定理 85.8.3（Selberg 筛法，1947）}：Selberg 引入 \(\Lambda^2\) 型筛法权重，将筛法问题转化为二次型极小化问题。设 \(\lambda_d\) 为待定实数（\(\lambda_1 = 1\)），考虑 \(\sum_{d \mid n} \lambda_d\) 的平方作为筛法的上界，通过优化 \(\lambda_d\) 得到比 Brun 筛法更好的上界。Selberg 筛法在证明 Bombieri-Vinogradov 定理和素数在算术级数中的分布估计中发挥了关键作用。

\textbf{定理 85.8.4（Bombieri-Vinogradov 定理，1965）}：对任意 \(A > 0\)，存在 \(B = B(A)\) 使得 \[\sum_{q \leq x^{1/2} (\log x)^{-B}} \max_{a: \gcd(a,q)=1} \left|\pi(x; q, a) - \frac{\operatorname{Li}(x)}{\varphi(q)}\right| \ll \frac{x}{(\log x)^A}\] 这一定理表明，在平均意义下，素数定理在算术级数中的误差项在模 \(q \leq x^{1/2-\varepsilon}\) 时非常小，从而在许多应用中（如陈景润定理）可以替代广义 Riemann 猜想。

\textbf{证明}（框架）：利用大筛法（large sieve）不等式和 Dirichlet \(L\)-函数的平均估计。定义 \(\psi(x; q, a) = \sum_{n \leq x, n \equiv a \bmod q} \Lambda(n)\)，通过正交性展开为 \(L\)-函数的 Perron 积分。对 \(q\) 求和并交换积分和求和顺序，利用 \(L\)-函数的均值估计（Gallagher 的均值定理）和 Vinogradov 的三角和估计，得到所需的上界。\(\blacksquare\)

\subsection{85.9 指数和与三角和}\label{ux6307ux6570ux548cux4e0eux4e09ux89d2ux548c}

\textbf{定义 85.8}（指数和 / 三角和）：指数和是形如 \(\sum_{n=1}^N e^{2\pi i f(n)}\) 的求和，其中 \(f: \mathbb{N} \to \mathbb{R}\)。指数和是解析数论中处理堆垒问题、丢番图逼近和模形式的基本工具。

\textbf{定理 85.9.1（Weyl 不等式，1916）}：设 \(f(x) = \alpha_k x^k + \cdots + \alpha_1 x + \alpha_0\)（\(k \geq 1\)），且存在互素整数 \(a, q\) 使得 \(|\alpha_k - a/q| \leq q^{-2}\)。则 \[\left|\sum_{n=1}^N e^{2\pi i f(n)}\right| \ll N^{1+\varepsilon} \left(\frac{1}{q} + \frac{1}{N} + \frac{q}{N^k}\right)^{1/2^{k-1}}\]

\textbf{证明}：对 \(k\) 归纳。\(k=1\) 时是几何级数求和。对 \(k > 1\)，使用 Weyl 差分法：\(\left|\sum_{n=1}^N e^{2\pi i f(n)}\right|^2 = \sum_{|h| < N} (N - |h|) e^{2\pi i (f(n+h) - f(n))}\)，\(f(n+h) - f(n)\) 是 \(n\) 的 \(k-1\) 次多项式。由归纳假设和 Cauchy-Schwarz 不等式得证。\(\blacksquare\)

\textbf{定理 85.9.2（Vinogradov 中值定理，1935）}：设 \(J_{s,k}(N)\) 为方程 \(\sum_{i=1}^s (x_i^j - y_i^j) = 0\)（\(j = 1, \ldots, k\)）在 \(1 \leq x_i, y_i \leq N\) 中的整数解个数。则 \[J_{s,k}(N) \ll N^{2s - k(k+1)/2 + \eta(s,k)}\] 其中 \(\eta(s,k) \to 0\) 当 \(s \to \infty\)。Vinogradov 中值定理是 Waring 问题中 \(G(k)\) 上界估计的核心工具，其证明使用了 \(p\)-进方法（Wooley 的''高效同余''方法）和组合数论中的光滑数（smooth numbers）技术。

解析数论的现代发展包括：自守形式与 \(L\)-函数的 Langlands 纲领（将模形式与 Galois 表示联系起来）、随机矩阵理论与零点分布（Montgomery 的零点对相关猜想与 GUE 系综的联系）、以及加性组合学（additive combinatorics）------Green-Tao 定理（2004年，素数中包含任意长的等差数列）将调和分析、遍历理论和组合数论融合，开创了解析数论的新范式。

\newpage

\chapter{08-代数几何与数论/02-数论/03-椭圆曲线.md}\label{ux4ee3ux6570ux51e0ux4f55ux4e0eux6570ux8bba02-ux6570ux8bba03-ux692dux5706ux66f2ux7ebf.md}

\section{第297章：椭圆曲线}\label{ux7b2c297ux7ae0ux692dux5706ux66f2ux7ebf}

椭圆曲线是亏格为 1 的光滑射影代数曲线，其上有一个特殊点（无穷远点）。椭圆曲线不''椭圆''------其名称源于与椭圆积分的历史联系（计算椭圆弧长时出现的积分），但本身是代数几何、数论和密码学的交汇点。椭圆曲线的研究始于 19 世纪 Abel 和 Jacobi 的椭圆函数理论，经过 20 世纪 Mordell、Weil、Siegel 等人的算术化，到 20 世纪末 Wiles 证明 Fermat 大定理和 BSD 猜想的核心进展而达到顶峰。椭圆曲线之所以如此重要，是因为它同时具有三种结构：（1）代数曲线结构（亏格 1）；（2）Abel 群结构（群律）；（3）复解析结构（复环面 \(\mathbb{C}/\Lambda\)）。这三种结构的统一和相互作用使得椭圆曲线成为连接代数几何、数论、复分析和表示论的桥梁。在本章中，我们从 Weierstrass 方程和群律出发，系统讨论椭圆曲线的有理点理论（Mordell-Weil 定理）、Tate 模和 Galois 表示、L 函数和 BSD 猜想、以及模性定理和 Fermat 大定理的证明框架。

\subsection{87.1 Weierstrass 方程与群律}\label{weierstrass-ux65b9ux7a0bux4e0eux7fa4ux5f8b}

\textbf{定义 87.1}（Weierstrass 方程）：域 \(K\) 上的椭圆曲线可由 \textbf{Weierstrass 方程}给出：

\[E: y^2 + a_1 xy + a_3 y = x^3 + a_2 x^2 + a_4 x + a_6\]

其中 \(a_i \in K\)。当 \(\operatorname{char} K \neq 2, 3\) 时，可通过变量代换化为标准形式 \(y^2 = x^3 + Ax + B\)（\(A, B \in K\)），判别式 \(\Delta = -16(4A^3 + 27B^2) \neq 0\)（保证曲线光滑，即无奇点）。判别式为零等价于曲线有奇点（node 或 cusp），此时为奇异椭圆曲线（非椭圆曲线，而是有理曲线或带奇点的曲线）。

\textbf{定义 87.1b}（\(j\)-不变量）：椭圆曲线 \(E: y^2 = x^3 + Ax + B\) 的 \(j\)-不变量定义为 \[j(E) = 1728 \cdot \frac{4A^3}{4A^3 + 27B^2}\]

两条椭圆曲线在代数闭域上同构当且仅当它们的 \(j\)-不变量相等。\(j\)-不变量是椭圆曲线模空间上的基本不变量：\(\mathbb{C}\) 上的椭圆曲线的同构类与 \(\mathbb{C}\) 中的点一一对应（通过 \(j\)-不变量），而模曲线 \(X(1)\) 正是 \(j\)-线。

\textbf{定义 87.2}（群律）：椭圆曲线 \(E(K)\) 上的点构成一个 \textbf{Abel 群}，其群运算（``弦切法''）定义如下：设 \(P, Q \in E\)。过 \(P, Q\) 的直线（若 \(P = Q\) 则取切线）与 \(E\) 交于第三点 \(R\)。\(P + Q\) 定义为 \(R\) 关于 \(x\) 轴的对称点。无穷远点 \(\mathcal{O}\) 是群的单位元。

\textbf{定理 87.1}（群律的代数公式）：对 \(E: y^2 = x^3 + Ax + B\)，设 \(P = (x_1, y_1)\)，\(Q = (x_2, y_2)\)。则： - 若 \(P \neq \pm Q\)，\(P + Q = (x_3, y_3)\)，其中 \(x_3 = \lambda^2 - x_1 - x_2\)，\(y_3 = \lambda(x_1 - x_3) - y_1\)，\(\lambda = (y_2 - y_1)/(x_2 - x_1)\) - 若 \(P = Q\)（倍点），\(x_3 = \lambda^2 - 2x_1\)，\(y_3 = \lambda(x_1 - x_3) - y_1\)，\(\lambda = (3x_1^2 + A)/(2y_1)\)

\textbf{证明}：直线 \(y = \lambda x + \nu\) 与 \(E\) 的交点满足 \((\lambda x + \nu)^2 = x^3 + Ax + B\)，即 \(x^3 - \lambda^2 x^2 + (A - 2\lambda\nu)x + (B - \nu^2) = 0\)。三个根之和为 \(\lambda^2\)（由 Vieta 公式），已知 \(x_1, x_2\) 得 \(x_3 = \lambda^2 - x_1 - x_2\)。\(y_3\) 由直线方程确定。倍点公式由切线斜率 \(\lambda\) 的极限导出。\(\blacksquare\)

\textbf{定理 87.1b}（群律的结合律）：弦切法定义的群运算满足结合律 \((P+Q)+R = P+(Q+R)\)。结合律的验证是椭圆曲线群结构中技术上最复杂的部分，可通过以下方式证明：(1) 直接代数计算（繁琐但直接）；(2) 利用 Riemann-Roch 定理和除子类群（Picard 群）的抽象性质；(3) 利用复解析观点（\(\mathbb{C}/\Lambda\) 上的加法明显结合）。

\textbf{证明（除子理论）}：椭圆曲线 \(E\) 上的除子类群 \(\operatorname{Pic}^0(E)\) 与 \(E\) 同构，映射为 \(P \mapsto [P] - [\mathcal{O}]\)。群律的弦切法定义恰为除子加法：\([P] + [Q] = [P+Q] + [\mathcal{O}]\)（在 \(\operatorname{Pic}^0(E)\) 中），因此结合律由除子加法的结合律自动得出。\(\blacksquare\)

椭圆曲线上的群结构是代数几何中最优美的构造之一。从几何观点看，群律的''弦切法''定义是亏格为 1 的曲线上的除子类群（Picard 群）的具体实现：椭圆曲线上的除子类群 \(\operatorname{Pic}^0(E)\) 与 \(E\) 自身同构（通过映射 \(P \mapsto [P] - [\mathcal{O}]\)）。这一同构是椭圆曲线作为 Abel 簇（一维 Abel 簇）的代数定义，也是 Mordell-Weil 定理的出发点。从复分析的观点看，复数域上的椭圆曲线 \(E(\mathbb{C})\) 同构于复环面 \(\mathbb{C}/\Lambda\)（\(\Lambda\) 为格），群运算对应于 \(\mathbb{C}\) 中的加法模 \(\Lambda\)。这一对应将椭圆曲线的''代数''（Weierstrass \(\wp\) 函数）与''超越''（椭圆函数）联系起来，并通过模形式理论与数论中的 L 函数和 BSD 猜想连接。

\textbf{定理 87.1c}（Weierstrass \(\wp\) 函数与复环面）：设 \(\Lambda = \mathbb{Z}\omega_1 + \mathbb{Z}\omega_2\) 为 \(\mathbb{C}\) 中的格（\(\Im(\omega_1/\omega_2) > 0\)）。Weierstrass \(\wp\) 函数定义为 \[\wp(z; \Lambda) = \frac{1}{z^2} + \sum_{\omega \in \Lambda \setminus \{0\}} \left(\frac{1}{(z-\omega)^2} - \frac{1}{\omega^2}\right)\]

\(\wp\) 函数满足微分方程 \(\wp'(z)^2 = 4\wp(z)^3 - g_2\wp(z) - g_3\)，其中 \(g_2 = 60\sum_{\omega \neq 0} \omega^{-4}\)，\(g_3 = 140\sum_{\omega \neq 0} \omega^{-6}\) 为 Eisenstein 级数。因此映射 \(z \mapsto (\wp(z), \wp'(z))\) 给出了复环面 \(\mathbb{C}/\Lambda\) 与椭圆曲线 \(y^2 = 4x^3 - g_2 x - g_3\) 之间的同构。

\textbf{证明}：首先证明 \(\wp\) 函数的 Laurent 展开：\(\wp(z) = \frac{1}{z^2} + \sum_{k=1}^\infty (2k+1)G_{2k+2} z^{2k}\)，其中 \(G_{2k} = \sum_{\omega \neq 0} \omega^{-2k}\) 为 Eisenstein 级数。比较 \(\wp(z)^2\) 和 \(\wp'(z)^2\) 的 Laurent 展开：\(\wp(z)^2 = \frac{1}{z^4} + \frac{6G_4}{z^0} + \cdots\)，\(\wp'(z)^2 = \frac{4}{z^6} - \frac{24G_4}{z^2} - 80G_6 + \cdots\)。考虑差分 \(f(z) = \wp'(z)^2 - 4\wp(z)^3 + 60G_4 \wp(z) + 140G_6\)。\(f(z)\) 是 \(\Lambda\)-周期的椭圆函数，且仅在 \(z=0\) 处有极点。其 Laurent 展开中 \(z^{-1}\) 项系数为零，且常数项也为零，故 \(f(z)\) 是整函数。由 Liouville 定理（椭圆函数若整则必为常数），\(f(z) = 0\)。令 \(g_2 = 60G_4\)，\(g_3 = 140G_6\)，即得微分方程。随后映射 \(z \mapsto (\wp(z), \wp'(z))\) 将 \(\mathbb{C}/\Lambda\) 映到该椭圆曲线，且为单射（因为 \(\wp\) 为偶函数且仅在格点上取相同值）和满射（对任意 \((x,y)\) 满足方程，椭圆积分 \(\int_\infty^x \frac{dt}{\sqrt{4t^3-g_2t-g_3}}\) 给出 \(z\)）。\(\blacksquare\)

\subsection{87.2 椭圆曲线上的有理点与 Mordell-Weil 定理}\label{ux692dux5706ux66f2ux7ebfux4e0aux7684ux6709ux7406ux70b9ux4e0e-mordell-weil-ux5b9aux7406}

\textbf{定义 87.3}（有理点）：设 \(E\) 是 \(\mathbb{Q}\) 上的椭圆曲线（即 Weierstrass 方程系数 \(a_i \in \mathbb{Q}\)）。\(E(\mathbb{Q})\) 是 \(E\) 上有理坐标的点的集合（连同无穷远点 \(\mathcal{O}\)）。由群律，\(E(\mathbb{Q})\) 是 \(E(\overline{\mathbb{Q}})\) 的子群。

\textbf{定理 87.2}（Mordell-Weil 定理，Mordell 1922, Weil 1928）：\(E(\mathbb{Q})\) 是有限生成的 Abel 群。即

\[E(\mathbb{Q}) \cong \mathbb{Z}^r \oplus E(\mathbb{Q})_{\text{tors}}\]

其中 \(r \geq 0\) 称为 \(E\) 的\textbf{秩}（rank），\(E(\mathbb{Q})_{\text{tors}}\) 是有限挠子群。

\textbf{定理 87.3}（Mazur 挠定理，1977）：对 \(\mathbb{Q}\) 上的椭圆曲线 \(E\)，挠子群 \(E(\mathbb{Q})_{\text{tors}}\) 只可能是以下 15 种群之一：\(\mathbb{Z}/n\mathbb{Z}\)（\(n = 1, \ldots, 10, 12\)）或 \(\mathbb{Z}/2\mathbb{Z} \times \mathbb{Z}/2n\mathbb{Z}\)（\(n = 1, 2, 3, 4\)）。Mazur 定理的意义在于：\(\mathbb{Q}\) 上椭圆曲线的挠子群具有一致的上界，且所有可能的情形已被完全分类。挠子群的具体结构可以通过 Nagell-Lutz 定理（挠点的坐标是整数且 \(y=0\) 或 \(y^2 \mid \Delta\)）来计算。

\textbf{证明（框架，Mordell-Weil 定理）}：核心思想是''无穷下降法''（descent）的现代形式。第一步：将 \(E(\mathbb{Q})\) 嵌入到 Selmer 群 \(\operatorname{Sel}^{(n)}(E/\mathbb{Q})\) 中，后者是有限群（由 Galois 上同调和 Kummer 理论）。具体地，由 Kummer 正合列 \(0 \to E(\mathbb{Q})/nE(\mathbb{Q}) \to \operatorname{Sel}^{(n)}(E/\mathbb{Q}) \to \Sha(E/\mathbb{Q})[n] \to 0\)，其中 \(\Sha(E/\mathbb{Q})\) 为 Tate-Shafarevich 群。第二步：证明 \(E(\mathbb{Q})/nE(\mathbb{Q})\) 是有限群（通过 Selmer 群的有限性和 Weil-Chatelet 群的结构）。第三步：由有限生成的充要条件（下降引理：若 \(A\) 是 Abel 群，\(A/nA\) 有限且 \(A\) 的挠部分有限，则 \(A\) 有限生成），\(E(\mathbb{Q})\) 有限生成。Mazur 定理的证明极为深刻，依赖模曲线 \(X_0(N)\) 的有理点分类和 Eisenstein 理想的精细算术，核心在于证明 \(X_0(N)\) 的有理尖点以外没有其他有理点（当 \(N\) 不是上述 15 种情形时）。\(\blacksquare\)

Mordell-Weil 定理是椭圆曲线理论的第一个里程碑，它确立了椭圆曲线上的有理点具有有限生成的结构。椭圆曲线的秩 \(r\) 的确定是数论中最困难的问题之一------目前尚不知道秩是否可以任意大（已知最大秩为 28，由 Elkies 在 2006 年发现），也不存在一般算法能确定给定椭圆曲线的秩。秩的分布（如平均秩是否有限）是数论中的核心未解问题，与 Birch-Swinnerton-Dyer 猜想密切相关。在计算方面，2-下降法（2-descent）和更高的 \(n\)-下降法是将 Mordell-Weil 定理转化为实际计算秩和生成元的关键工具，其效能依赖于 Selmer 群和 Tate-Shafarevich 群的精细分析。

\textbf{定理 87.3b}（Siegel 定理，1929）：\(\mathbb{Q}\) 上的椭圆曲线在 \(\mathbb{Z}\) 上只有有限多个整点（即 \(x, y \in \mathbb{Z}\) 的解）。Siegel 定理是 Thue-Siegel-Roth 方法在椭圆曲线上的重要应用，其证明利用了椭圆对数的高度函数和 Diophantine 逼近的深刻结果。

\textbf{证明}：将椭圆曲线 \(E\) 的 Weierstrass 方程写为 \(y^2 = x^3 + Ax + B\)。设 \(P = (x, y) \in E(\mathbb{Z})\)。利用椭圆对数 \(\phi(P) = \int_\infty^P \frac{dx}{2y}\)（Neron-Tate 高度），整点满足 \(h(P) \approx \log|x|\)。Siegel 的核心论证是：将整点条件转化为 \(S\)-单位方程 \(x\) 与 \(y\) 的高度估计。由 Roth 定理（代数数的 Diophantine 逼近），若 \(x\) 是代数整数且其共轭满足某些逼近条件，则 \(x\) 只能取有限多个值。具体地，将 \(E\) 的方程变形为 \(\frac{y}{\sqrt{x^3}}\) 逼近某个代数数，由 Roth 定理该逼近只有有限多解。更现代的证明利用 Siegel 恒等式和 Faltings 定理的推论：\(E\) 的整点模型是 Siegel 概形（\(S\)-整数点有限性）。\(\blacksquare\)

\textbf{定义 87.3b}（高度函数）：\(\mathbb{Q}\) 上的椭圆曲线 \(E\) 上的点 \(P = (x, y) \in E(\mathbb{Q})\) 的（对数）高度定义为 \(h(P) = \log \max(|p|, |q|)\)，其中 \(x = p/q\)（\(\gcd(p,q)=1\)）。规范高度 \(\hat{h}(P) = \lim_{n \to \infty} 4^{-n} h(2^n P)\) 是二次型，满足 \(\hat{h}(P+Q) + \hat{h}(P-Q) = 2\hat{h}(P) + 2\hat{h}(Q)\)。规范高度是 Mordell-Weil 群上正定的二次型，其与 \(L\) 函数导数的关系（BSD 公式）是 BSD 猜想的核心。

\subsection{87.3 Tate 模与 etale 上同调}\label{tate-ux6a21ux4e0e-etale-ux4e0aux540cux8c03}

\textbf{定义 87.4}（Tate 模）：设 \(E\) 是域 \(K\) 上的椭圆曲线，\(\ell\) 为素数（\(\ell \neq \operatorname{char} K\)）。\(E\) 的 \textbf{\(\ell\)-进 Tate 模} 定义为

\[T_\ell(E) = \varprojlim_n E[\ell^n]\]

其中 \(E[\ell^n] = \{P \in E(\overline{K}) : \ell^n P = \mathcal{O}\}\) 是 \(\ell^n\)-挠点群。\(E[\ell^n] \cong (\mathbb{Z}/\ell^n\mathbb{Z})^2\)，故 \(T_\ell(E) \cong \mathbb{Z}_\ell \times \mathbb{Z}_\ell\)（作为 \(\mathbb{Z}_\ell\)-模，秩为 \(2\)）。Galois 群 \(G_K = \operatorname{Gal}(\overline{K}/K)\) 在 \(E[\ell^n]\) 和 \(T_\ell(E)\) 上自然作用，给出 \(\ell\)-进 Galois 表示 \(\rho_{E, \ell}: G_K \to \operatorname{GL}_2(\mathbb{Z}_\ell)\)。

\textbf{定理 87.4}（Tate 模与 etale 上同调）：\(T_\ell(E) \cong H^1_{\text{et}}(E_{\overline{K}}, \mathbb{Z}_\ell(1))^\vee\)（Poincare 对偶）。更一般地，\(H^1_{\text{et}}(E_{\overline{K}}, \mathbb{Z}_\ell) \cong T_\ell(E)^\vee(-1)\)。Tate 模是椭圆曲线的一维 etale 上同调群的对偶，这一解释将椭圆曲线的 Galois 表示与代数几何中的上同调理论统一起来，为 Weil 猜想在椭圆曲线上的证明提供了框架。

\textbf{证明}：对椭圆曲线 \(E\)，其 etale 上同调 \(H^1_{\text{et}}(E_{\overline{K}}, \mathbb{Z}/\ell^n\mathbb{Z})\) 由 Kummer 正合列 \(0 \to \mu_{\ell^n} \to \mathbb{G}_m \xrightarrow{\ell^n} \mathbb{G}_m \to 0\) 和 Weil 配对 \(E[\ell^n] \times E[\ell^n] \to \mu_{\ell^n}\) 计算。Poincare 对偶给出 \(H^1_{\text{et}}(E_{\overline{K}}, \mathbb{Z}/\ell^n\mathbb{Z}) \cong E[\ell^n]^\vee\)。取投射极限 \(\varprojlim_n\) 即得 Tate 模的对偶。\(\blacksquare\)

\textbf{定理 87.4b}（Tate 的 Isogeny 定理，1966）：设 \(E_1, E_2\) 为数域 \(K\) 上的椭圆曲线。则自然映射 \[\operatorname{Hom}_K(E_1, E_2) \otimes \mathbb{Z}_\ell \to \operatorname{Hom}_{G_K}(T_\ell(E_1), T_\ell(E_2))\] 是同构。这意味着 \(E_1\) 和 \(E_2\) 在 \(K\) 上同源（isogenous）当且仅当它们的 \(\ell\)-进 Galois 表示同构。Tate 的 Isogeny 定理是 Faltings 在 1983 年证明 Mordell 猜想（高亏格曲线的有理点有限性）时推广到一般 Abel 簇的核心工具。

\textbf{证明}：考虑 \(\ell\)-进 Galois 表示。由 Tate 模的构造，存在自然同态 \(\operatorname{Hom}_K(E_1, E_2) \otimes \mathbb{Z}_\ell \to \operatorname{Hom}_{G_K}(T_\ell(E_1), T_\ell(E_2))\)。要证明这是同构。单射性：若 \(\phi: E_1 \to E_2\) 为非零同源，则 \(\ker \phi\) 是有限子群，故 \(\phi\) 在 \(T_\ell\) 上的作用非零。满射性：设 \(\alpha \in \operatorname{Hom}_{G_K}(T_\ell(E_1), T_\ell(E_2))\)。对每个 \(n\)，\(\alpha\) 模 \(\ell^n\) 诱导 \(E_1[\ell^n] \to E_2[\ell^n]\) 的 \(G_K\)-等变同态。由 Galois 上同调的 Kummer 理论，这些同态''粘合''为一个同源 \(\phi_n: E_1 \to E_2\)（模 \(\ell^n\) 意义下）。关键步骤：利用 Weil 配对和 Galois 表示的非退化性，证明 \(\phi_n\) 的极限给出 \(E_1 \to E_2\) 的代数同态。由 Faltings 的深造诣（Faltings 高度），\(\operatorname{Hom}_K(E_1, E_2)\) 是有限生成 \(\mathbb{Z}\)-模，其在 \(\mathbb{Z}_\ell\) 中的稠密性保证 \(\alpha\) 确实来自代数同态。\(\blacksquare\)

\textbf{定义 87.4b}（自同态环与复乘）：椭圆曲线 \(E\) 的自同态环 \(\operatorname{End}(E)\) 包含 \(\mathbb{Z}\)（由倍点映射 \([n]: P \mapsto nP\) 给出）。若 \(\operatorname{End}(E) \supsetneq \mathbb{Z}\)，则称 \(E\) 具有\textbf{复乘}（Complex Multiplication, CM）。此时 \(\operatorname{End}(E) \otimes \mathbb{Q}\) 是虚二次域，且 \(E\) 的 \(j\)-不变量是代数整数（称为奇异模）。CM 椭圆曲线是椭圆曲线理论中结构最丰富的一类，其 Galois 表示是可交换的（由虚二次域的理想类群描述），其 L 函数可分解为 Hecke 特征的 L 函数的乘积，因此 BSD 猜想在 CM 情形下可取得实质性进展（Coates-Wiles 1977 在 CM 情形下证明了 \(L(E,1) \neq 0 \Rightarrow E(\mathbb{Q})\) 有限）。

\subsection{87.4 L 函数与 BSD 猜想}\label{l-ux51fdux6570ux4e0e-bsd-ux731cux60f3}

\textbf{定义 87.5}（椭圆曲线的 Hasse-Weil L 函数）：设 \(E/\mathbb{Q}\) 是椭圆曲线。对每个素数 \(p\)，定义 \(a_p = p + 1 - \#E(\mathbb{F}_p)\)（其中 \(\#E(\mathbb{F}_p)\) 是 \(E\) 模 \(p\) 约化后的 \(\mathbb{F}_p\)-有理点数）。\(E\) 的 \textbf{Hasse-Weil L 函数} 定义为 Euler 乘积

\[L(E, s) = \prod_{p \nmid N} \frac{1}{1 - a_p p^{-s} + p^{1-2s}} \prod_{p \mid N} \frac{1}{1 - a_p p^{-s}}\]

其中 \(N\) 是 \(E\) 的导子（conductor）。在 \(p \nmid N\) 时（良约化），\(L_p(E, s) = (1 - \alpha_p p^{-s})^{-1}(1 - \beta_p p^{-s})^{-1}\)，其中 \(\alpha_p, \beta_p\) 是 Frobenius 特征值，满足 \(|\alpha_p| = |\beta_p| = \sqrt{p}\)（Hasse 界）。在 \(p \mid N\) 时，\(a_p\) 取值 \(0, 1, -1\)（分别对应加性约化、分裂乘性约化、非分裂乘性约化）。

\textbf{定理 87.5}（Hasse 定理，1934）：\(|\#E(\mathbb{F}_q) - (q + 1)| \leq 2\sqrt{q}\)。等价地，\(|a_p| \leq 2\sqrt{p}\)。

\textbf{证明}（框架）：对椭圆曲线 \(E/\mathbb{F}_q\)，Frobenius 自同态 \(\pi: (x, y) \mapsto (x^q, y^q)\) 满足 \(\pi^2 - t \pi + q = 0\)（在 \(\operatorname{End}(E)\) 中），其中 \(t = q + 1 - \#E(\mathbb{F}_q)\)。由自同态环的性质，\(\deg(a + b\pi) = a^2 + abt + b^2q \geq 0\)（对任意 \(a, b \in \mathbb{Z}\)）。这迫使判别式 \(t^2 - 4q \leq 0\)，即 \(|t| \leq 2\sqrt{q}\)。\(\blacksquare\)

\textbf{猜想 87.1}（Birch-Swinnerton-Dyer 猜想，BSD 猜想）：设 \(E/\mathbb{Q}\) 是椭圆曲线，\(r = \operatorname{rank} E(\mathbb{Q})\)。则 \(L(E, s)\) 在 \(s = 1\) 处的零点阶数等于 \(r\)：

\[\operatorname{ord}_{s=1} L(E, s) = \operatorname{rank} E(\mathbb{Q})\]

进一步，\(L(E, s)\) 在 \(s = 1\) 处的首项系数与 \(E\) 的算术不变量有精确的公式联系：

\[\lim_{s \to 1} \frac{L(E, s)}{(s-1)^r} = \frac{\Omega_E \cdot \operatorname{Reg}(E) \cdot \prod_p c_p \cdot |\Sha(E/\mathbb{Q})|}{|E(\mathbb{Q})_{\text{tors}}|^2}\]

其中 \(\Omega_E\) 为实周期（Neron 微分在 \(E(\mathbb{R})\) 上的积分），\(\operatorname{Reg}(E)\) 为椭圆曲线正规化子（基于规范高度配对的行列式），\(c_p\) 为 Tamagawa 数（局部因子），\(\Sha(E/\mathbb{Q})\) 为 Tate-Shafarevich 群（猜想为有限群）。

BSD 猜想是 Clay 数学研究所七大千禧年问题之一，它将椭圆曲线 \(L\) 函数的解析性质（零点阶数）与有理点群的代数性质（秩）联系起来，是 Langlands 纲领在 \(\operatorname{GL}_2\) 情形下的核心预测。BSD 猜想在秩 \(r \leq 1\) 的情形下已取得重大进展：Gross-Zagier 定理（1986）和 Kolyvagin 的 Euler 系方法证明了当 \(L(E, 1) \neq 0\) 时 \(r = 0\)，且当 \(L(E, 1) = 0\) 但 \(L'(E, 1) \neq 0\) 时 \(r = 1\)。更高秩情形下的 BSD 猜想至今仍是开放问题。

\textbf{定理 87.5b}（Gross-Zagier 定理，1986）：设 \(E/\mathbb{Q}\) 为模椭圆曲线，\(K\) 为虚二次域（满足 Heegner 条件）。则 \(L(E/K, s)\) 在 \(s=1\) 处的导数与 Heegner 点 \(P_K \in E(K)\) 的规范高度成正比： \[L'(E/K, 1) = \frac{4\pi^2}{\sqrt{|D_K|}} \cdot \frac{\hat{h}(P_K)}{\deg(\varphi_E)} \cdot (\text{非零因子})\] 其中 \(\varphi_E: X_0(N) \to E\) 为模参数化。Gross-Zagier 定理是 BSD 猜想在秩 \(1\) 情形下取得进展的关键，它将 \(L\) 函数的解析信息（导数）与算术几何信息（Heegner 点的高度）精确联系起来。

\textbf{证明}（框架）：设 \(K\) 为虚二次域，满足 Heegner 条件（\(N\) 的每个素因子在 \(K\) 中分裂）。Heegner 点 \(P_K\) 通过模曲线 \(X_0(N)\) 上的 CM 点构造：取 \(K\) 中的理想 \(\mathfrak{a}\) 使得 \(\mathcal{O}_K/\mathfrak{a} \cong \mathbb{Z}/N\mathbb{Z}\)，则 \((\mathbb{C}/\mathcal{O}_K, \mathfrak{a}^{-1}/\mathcal{O}_K)\) 给出 \(X_0(N)\) 上的 CM 点。Heegner 点 \(y_K \in J_0(N)(K)\) 定义为该 CM 点与尖点 \(\infty\) 的除子类的差。\(P_K\) 是 \(y_K\) 在模参数化 \(X_0(N) \to E\) 下的像。Gross-Zagier 公式的核心是将 \(L'(E/K, 1)\) 表示为 \(P_K\) 的 Neron-Tate 高度乘以一个显式非零因子。证明分两步：第一步，利用 Rankin-Selberg 方法将 \(L'(E/K, 1)\) 表示为某个 Eisenstein 级数与尖点形式的 Petersson 内积；第二步，利用 Waldspurger 公式和 Shimizu 对应，将该内积等同于 Heegner 点的高度。\(\blacksquare\)

\subsection{87.5 模性定理与 Fermat 大定理}\label{ux6a21ux6027ux5b9aux7406ux4e0e-fermat-ux5927ux5b9aux7406}

\textbf{定理 87.6}（模性定理 / Taniyama-Shimura-Weil 猜想，Wiles 1995, Taylor-Wiles 1995, Breuil-Conrad-Diamond-Taylor 2001）：\(\mathbb{Q}\) 上的每条椭圆曲线都是模的。即存在权为 \(2\) 的尖点新形式 \(f = \sum_{n=1}^\infty a_n e^{2\pi i n z}\)（\(a_1 = 1\)）使得 \(L(E, s) = L(f, s)\)。等价地，存在非平凡的有限态射 \(X_0(N) \to E\)（模参数化），其中 \(N\) 是 \(E\) 的导子。模性定理的陈述等价于：\(E\) 的 \(\ell\)-进 Galois 表示 \(\rho_{E,\ell}\) 同构于某个权 \(2\) 模形式 \(f\) 的 Galois 表示 \(\rho_{f,\ell}\)。

\textbf{证明}（框架）：Wiles 的证明策略：设 \(\rho = \rho_{E,\ell}\) 为 \(E\) 的 \(\ell\)-进 Galois 表示（\(\ell\) 为奇素数）。第一步：构造 Hecke 环 \(\mathbb{T}\) 作用在模形式空间上，证明 \(\rho\) 模 \(\ell\) 的约化 \(\bar{\rho}\) 是模的（Langlands-Tunnell 定理处理可解情形，Wiles 解决不可解情形）。第二步：考虑 Galois 表示的变形环 \(R\)（参数化 \(\rho\) 的所有提升）和 Hecke 环 \(\mathbb{T}\)。存在自然满射 \(R \twoheadrightarrow \mathbb{T}\)。关键步骤：证明 \(R\) 是 Gorenstein 完全交环且 \(\#R/\mathfrak{p} = \#\mathbb{T}/\mathfrak{p}\)（通过 Taylor-Wiles 的修补方法），从而 \(R \cong \mathbb{T}\)。第三步：\(R \cong \mathbb{T}\) 意味着每个变形来自模形式，故 \(\rho\) 是模的。Breuil-Conrad-Diamond-Taylor 完成了 \(\ell = 2\) 和 \(\ell\) 处非半稳定情形的推广。\(\blacksquare\)

\textbf{定理 87.7}（Fermat 大定理，Wiles 1995）：方程 \(x^n + y^n = z^n\) 在 \(n \geq 3\) 时没有正整数解 \((x, y, z)\) 使得 \(xyz \neq 0\)。

\textbf{证明}（框架）：假设存在非平凡解 \((a, b, c)\)，\(n = \ell\) 为奇素数（可约化情形）。构造 Frey 曲线 \(E: y^2 = x(x - a^\ell)(x + b^\ell)\)。其判别式 \(\Delta = (abc)^{2\ell}\) 异常大，而导子 \(N\) 相对较小。由 Ribet 的提升定理（Serre 的 \(\varepsilon\)-猜想），\(E\) 的模 \(\ell\) Galois 表示 \(\bar{\rho}_{E,\ell}\) 若来自模形式，则必须来自水平 \(N=2\) 的模形式------但 \(S_2(\Gamma_0(2)) = \{0\}\)（水平 \(2\) 的模曲线 \(X_0(2)\) 的亏格为 \(0\)），矛盾。因此 \(E\) 不能是模的。但模性定理断言 \(E\) 必须是模的。矛盾。故无非平凡解。\(\blacksquare\)

Fermat 大定理的证明是 20 世纪数学最辉煌的成就之一。Wiles 的证明策略基于 Frey-Serre-Ribet 的洞察：将 Fermat 方程的解转化为椭圆曲线（Frey 曲线），然后利用模性定理将椭圆曲线与模形式联系起来，最终通过 Galois 表示的变形理论和交换环的完全交理论（Taylor-Wiles 的修补方法）证明模性定理。证明的核心工具包括：Hecke 代数、Mazur 的 Eisenstein 理想理论、Galois 表示的变形环、以及 Gorenstein 完全交环的同调性质。Wiles 的证明开创了代数数论中的''模性提升定理''（modularity lifting theorem）这一全新领域。

\textbf{定理 87.8}（模性提升定理的一般形式）：设 \(\rho: G_{\mathbb{Q}} \to \operatorname{GL}_2(\overline{\mathbb{Q}}_\ell)\) 为连续奇 Galois 表示，满足：(1) \(\rho\) 在 \(\ell\) 处是 de Rham 的（或晶体表示），具有适当的 Hodge-Tate 权；(2) \(\bar{\rho}\)（模 \(\ell\) 约化）是模的（来自某个模形式）。则 \(\rho\) 本身是模的。这一框架将 Wiles 的特定情形推广到一般 Galois 表示，成为 Langlands 纲领中 \(\operatorname{GL}_n\) 的模性定理（如 Sato-Tate 猜想的证明，由 Clozel-Harris-Taylor 等人在 2008 年前后完成）的基石。

\textbf{证明}（框架）：设 \(\rho\) 满足条件 (1) 和 (2)。第一步：构造 \(\rho\) 的变形环 \(R\)（参数化 \(\rho\) 的具有相同局部 Galois 类型的所有提升）。由 Fontaine-Laffaille 理论和 \(p\)-进 Hodge 理论，\(\ell\) 处的 de Rham 条件等价于 \(\rho\) 在该处的形变受到严格的局部约束（使得 \(R\) 有良好的维数性质）。第二步：由条件 (2)，\(\bar{\rho}\) 来自某个模形式 \(f\)。利用 Hecke 环 \(\mathbb{T}\) 和变形环 \(R\) 之间的同构 \(R \cong \mathbb{T}\)（Taylor-Wiles 修补方法在一般情形下的推广），证明 \(\rho\) 本身来自模形式。关键输入：\(\mathbb{T}\) 是 Gorenstein 环，且 \(R\) 和 \(\mathbb{T}\) 的切线空间维数由 Galois 上同调的 Selmer 群计算（Wiles 的数值判据）。\(\blacksquare\)

\textbf{定理 87.9}（Sato-Tate 猜想，已证明）：设 \(E/\mathbb{Q}\) 为非 CM 椭圆曲线。则 \(a_p/\sqrt{p}\) 的分布服从 Sato-Tate 分布（半圆分布）\(\frac{2}{\pi}\sqrt{1-t^2} dt\)（\(t \in [-2, 2]\)）。Sato-Tate 猜想的证明（Harris-Taylor 等，2008）是 Langlands 纲领中 \(\operatorname{GL}_n\) 的自守性提升技术的重大胜利，依赖模性定理、Arthur-Selberg 迹公式和 Shimura 簇的几何。

\textbf{证明}（框架）：由模性定理，\(E\) 对应一个权 \(2\) 尖点新形式 \(f\)。Sato-Tate 猜想等价于 \(f\) 的对称幂 \(L\)-函数 \(L(s, \operatorname{Sym}^m f)\) 在 \(\operatorname{Re}(s) = 1\) 上无零点且全纯（对 \(m \geq 1\)）。这由 Langlands 函子性猜想预测：\(\operatorname{Sym}^m f\) 应来自 \(\operatorname{GL}_{m+1}\) 的自守表示。Harris-Taylor 的证明：对 \(m\) 归纳，利用 \(\operatorname{GL}_n\) 的模性提升定理和 Shimura 簇的 \(\ell\)-进上同调中的 Galois 表示构造自守表示。具体地，在恰当的正交或酉 Shimura 簇上，利用 Arthur-Selberg 迹公式的稳定化（由 Arthur 完成）和基变换的相容性，将 \(\operatorname{Sym}^m\) 的模性从 \(\operatorname{GL}_m\) 提升到 \(\operatorname{GL}_{m+1}\)。由 \(\operatorname{Sym}^m\) 的模性，结合 Sato-Tate 测度的特征函数展开，即得 \(a_p/\sqrt{p}\) 的分布收敛到半圆分布。\(\blacksquare\)

\textbf{定理 87.10}（Fermat 大定理的推广------Fermat-Catalan 猜想）：方程 \(x^p + y^q = z^r\)（\(1/p + 1/q + 1/r < 1\)）只有有限多个互素的正整数解。Darmon-Granville 定理（1995）利用 Faltings 定理和模性方法证明了这一有限性结果，而 Beal 猜想（\(\$1,000,000\) 悬赏）断言这些解不存在（即仅有 \(1^p + 2^3 = 3^2\) 等少数平凡解）。

\textbf{证明}（Darmon-Granville 框架）：假设 \((x, y, z)\) 是互素正整数解。构造三重 Fermat 曲线：\(X: x^p + y^q = z^r\)。关键想法：将每个解对应到某条特定曲线的有理点。构造 Frey-Hellegouarch 曲线 \(E\)，其判别式包含 \(x^p y^q z^r\) 的因子。由 Ribet 的提升定理，\(E\) 的 Galois 表示约化后必须来自某个低水平的模形式。但水平 \(N\) 的模形式空间是有限维的，且由 Faltings 定理，满足 \(\gcd(p,q,r)\) 条件的互素解对应的曲线只有有限多条。因此，若存在无穷多解，则将有无穷多条不同的 Frey 曲线，但这些曲线都对应同一个模形式（因为水平由 \(p,q,r\) 的有限多种组合决定），这与模形式的有限维性矛盾。更精确地，Darmon-Granville 将解分类为有限多种''Frey 类型''，每种类型对应的模形式空间是有限维的，故解的总数是有限的。\(\blacksquare\)

\section{第298章：Birch-Swinnerton-Dyer猜想与Heegner点}\label{ux7b2c298ux7ae0birch-swinnerton-dyerux731cux60f3ux4e0eheegnerux70b9}

Birch和Swinnerton-Dyer在1960年代初期通过剑桥EDSAC计算机对椭圆曲线的大量数值实验，发现了一个惊人的规律：椭圆曲线\(E\)的\(L\)-函数在\(s=1\)处的零点阶数似乎与\(E\)的有理点群的秩精确相关。这一观察成为Clay千禧年七大问题之一的BSD猜想。经历半个多世纪的深入研究，BSD猜想在秩\(r \leq 1\)的情形下取得了决定性突破：Gross-Zagier定理通过Heegner点的Neron-Tate高度公式建立了\(L\)函数导数与算术几何的直接联系，而Kolyvagin的Euler系方法则利用这些Heegner点证明了\(E(\mathbb{Q})\)的秩不超过\(L(E,s)\)在\(s=1\)处的零点阶数。本章聚焦于BSD猜想在解析秩\(\leq 1\)时的部分证明，详细阐述Heegner点的构造（通过虚二次域的环类域理论和模参数化）、Gross-Zagier高度公式、Kolyvagin的Euler系论证框架，以及Tunnell定理将千年难题''Congruent Number问题''划归为BSD猜想的特殊情形。

\subsection{88.1 Mordell曲线与BSD猜想的数值起源}\label{mordellux66f2ux7ebfux4e0ebsdux731cux60f3ux7684ux6570ux503cux8d77ux6e90}

\textbf{定义 88.1}（Mordell曲线）：Mordell曲线即形如\(y^2 = x^3 + k\)（\(k \in \mathbb{Z} \setminus \{0\}\)）的椭圆曲线。这是BSD猜想最早的数值实验对象。Birch和Swinnerton-Dyer对大量\(k\)值计算了\(N_p = \#E(\mathbb{F}_p)\)和\(\prod_{p \leq X} N_p/p\)的渐进行为。

\textbf{定理 88.1}（Birch-Swinnerton-Dyer数值定律）：数值实验表明： \[\prod_{p \leq X} \frac{N_p}{p} \sim C \cdot (\log X)^r\quad (X \to \infty),\] 其中\(r = \operatorname{rank}E(\mathbb{Q})\)。这一数值发现直接引出了BSD猜想：\(\operatorname{ord}_{s=1}L(E,s) = r\)。

\textbf{证明（数值发现到精确猜想的推导）}：设\(E\)的Hasse-Weil \(L\)函数为\(L(E,s) = \prod_p L_p(E,p^{-s})^{-1}\)，其中\(L_p(E,T) = 1 - a_pT + pT^2\)（\(p \nmid N\)）。形式上，在\(s=1\)处， \[L(E,1) \stackrel{?}{=} \prod_p \frac{p}{N_p} \cdot \text{(修正因子)}.\] 若\(L(E,s)\)在\(s=1\)处有\(r\)阶零点，则\(L(E,s) \sim c(s-1)^r\)。通过Tauber型定理（Ikehara定理的变体，由Rankin和Selberg发展），\(\prod_p N_p/p\)的部分乘积渐进于\(C(\log X)^r\)的猜想等价于\(L^{(r)}(E,1) \neq 0\)且低阶导数全部为零。Birch和Swinnerton-Dyer将这一数值观察提炼为精确的数学猜想，进而提出了\(L^{(r)}(E,1)\)与\(\operatorname{Reg}(E) \cdot |\Sha|\)之间的精确公式。\(\blacksquare\)

\subsection{88.2 Heegner点构造与Gross-Zagier定理}\label{heegnerux70b9ux6784ux9020ux4e0egross-zagierux5b9aux7406}

\textbf{定义 88.2}（Heegner条件与Heegner点）：设\(E/\mathbb{Q}\)为椭圆曲线，\(N\)为其导子。虚二次域\(K = \mathbb{Q}(\sqrt{-D})\)（\(D > 0\)）满足\textbf{Heegner条件}：\(N\)的每个素因子在\(K\)中分裂（即\(\left(\frac{-D}{p}\right) = 1\)对\(p \mid N\)）。此时存在\(K\)中的理想\(\mathfrak{n}\)使得\(\mathcal{O}_K/\mathfrak{n} \cong \mathbb{Z}/N\mathbb{Z}\)。复环面\(\mathbb{C}/\mathcal{O}_K\)配备\(\mathfrak{n}^{-1}/\mathcal{O}_K\)给出的\(N\)阶循环子群，定义了模曲线\(X_0(N)\)上的一个CM点------称为\textbf{Heegner点}。Heegner点\(y_K \in J_0(N)(K)\)（Jacobian）在模参数化\(\varphi: X_0(N) \to E\)下的像\(P_K = \varphi_*(y_K) \in E(K)\)仍称为Heegner点。

\textbf{定理 88.2}（Gross-Zagier定理，1986）：设\(E/\mathbb{Q}\)为模椭圆曲线（导子\(N\)），\(K\)为满足Heegner条件的虚二次域（判别式\(-D\)）。则 \[L'(E/K, 1) = \frac{4\pi^2}{\sqrt{D}} \cdot \frac{\|f\|^2}{\deg(\varphi)} \cdot \hat{h}(P_K),\] 其中\(\|f\|\)为\(E\)对应的权2新形式\(f\)的Petersson范数，\(\deg(\varphi)\)为模参数化\(\varphi: X_0(N) \to E\)的次数，\(\hat{h}(P_K)\)为Heegner点\(P_K\)的规范高度。

\textbf{证明概要}（完整证明极其复杂，此处给出核心步骤）：

第一步（Rankin-Selberg卷积）：\(L'(E/K, 1)\)可通过Rankin-Selberg方法表示为\(f\)与某个Eisenstein级数的Petersson内积的导数。设\(\Theta_K\)为\(K\)的\(\theta\)-级数（权1模形式），\(E(z,s)\)为Eisenstein级数。通过展开计算， \[L(E/K, s) = \langle f, \Theta_K \cdot E(\cdot, s-\tfrac{1}{2}) \rangle_{\text{Pet}} \cdot (\text{Gamma因子}).\] 求导后，\(L'(E/K, 1)\)变为\(f\)与\(\Theta_K \cdot \frac{\partial}{\partial s}E(z, 1/2)\)的内积，后者本质上是一个Green函数。

第二步（Green函数与高度）：在模曲线\(X_0(N)\)上，Neron-Tate高度\(\hat{h}\)可以通过Arakelov理论表示为Green函数的积分。具体地，对除子\(D = \sum n_i(P_i)\)，其高度为 \[\hat{h}(D) = \sum_{i,j} n_i n_j G(P_i, P_j),\] 其中\(G\)为\(X_0(N)\)上的Arakelov-Green函数（满足\(\partial\bar{\partial}G = \delta_{\text{diag}} - \mu\)，\(\mu\)为规范体积形式）。

第三步（高度公式的推导）：将Rankin-Selberg积分与Heegner除子\(D_K = y_K - (\infty)\)的Arakelov高度联系起来。核心恒等式（Gross-Zagier公式的核心）： \[\langle f, \Theta_K \cdot \frac{\partial}{\partial s}E(z, \tfrac{1}{2}) \rangle_{\text{Pet}} = \frac{2\pi}{\sqrt{D}} \cdot \hat{h}(D_K) \cdot \|f\|^2.\] 该恒等式的证明利用了Waldspurger公式（将Rankin-Selberg积分的中心值关联到环面周期积分）和高度函数在CM点处的局部展开。通过繁琐但系统性的模形式计算和Stark单位的理论，可以验证所有常数因子的精确取值。

第四步（模参数化因子）：Heegner点\(P_K\)是\(D_K\)在Jacobi映射\(J_0(N) \to E\)下的像。由函子性，\(\hat{h}(P_K) = \deg(\varphi) \cdot \hat{h}(D_K)\)（注意规范化）。代入即得定理中的精确公式。\(\blacksquare\)

Gross-Zagier定理是算术几何的里程碑：它首次在非平凡情形下建立了\(L\)函数的解析信息（导数）与代数几何量（Heegner点的高度）之间的精确等式。这直接证明了：当\(L'(E/K, 1) \neq 0\)时，\(P_K\)为非挠点，从而\(\operatorname{rank}E(K) \geq 1\)。这正是BSD猜想中''\(L'(E,1) \neq 0 \Rightarrow r \geq 1\)``方向的关键突破。

\subsection{88.3 Kolyvagin Euler系方法}\label{kolyvagin-eulerux7cfbux65b9ux6cd5}

\textbf{定理 88.3}（Kolyvagin定理，1988/1990）：在上述设定下，若\(P_K\)为非挠点（即\(\hat{h}(P_K) \neq 0\)），则\(\operatorname{rank}E(K) = 1\)且\(\Sha(E/K)\)有限。结合Gross-Zagier定理，这证明了：若\(L'(E/K, 1) \neq 0\)，则\(\operatorname{rank}E(K) = 1\)。

\textbf{证明框架}（Kolyvagin的Euler系方法）：

第一步（Heegner点的Galois迹）：对\(K\)的每个Abel扩张\(L/K\)（特别是环类域），Heegner点\(P_K\)提升为\(P_L \in E(L)\)。环类域对应\(K\)的序\(\mathcal{O}_c = \mathbb{Z} + c\mathcal{O}_K\)，其Galois群\(\operatorname{Gal}(K_c/K) \cong \operatorname{Pic}(\mathcal{O}_c)\)。Heegner点族\(\{P_{K_c}\}_{c \geq 1}\)满足互反律：\(P_{K_c}\)在\(\operatorname{Gal}(K_c/K)\)作用下的关系由Shimura互反律精确描述。

第二步（Kolyvagin类的构造）：设\(\ell\)为满足某些技术条件的素数。对每个Heegner判别式\(c\)（\(c\)为与\(ND\)互素的方自由整数），构造cohomology类 \[\kappa(c) \in H^1(K, E[\ell])\] 通过\(P_{K_c}\)在Galois群作用下的Kummer像。这些类\(\{\kappa(c)\}\)构成一个\textbf{Euler系}------即对素数\(q \nmid c\)， \[\operatorname{Cor}_{K_{cq}/K_c}(\kappa(cq)) \equiv a_q \kappa(c) \pmod{\text{(局部条件)}},\] 其中\(a_q = q+1 - \#E(\mathbb{F}_q)\)为\(E\)的Frobenius迹。Euler系的公理由Heegner点的形式群性质和局部Tate对偶导出。

第三步（Selmer群的上界估计）：利用Euler系中的类，通过精巧的局部-整体对偶论证（Chebotarev密度定理的精细应用），证明Selmer群\(\operatorname{Sel}^{(\ell)}(E/K)\)的\(\mathbb{Z}/\ell\mathbb{Z}\)-维数不超过1。具体策略：若\(\operatorname{Sel}^{(\ell)}(E/K)\)中有两个线性无关的类，则存在素数\(q\)使其中一个类在\(q\)处''非分歧''而\(\kappa(q)\)非零，从而导出矛盾------这是Euler系论证的核心技巧。

第四步（秩的确定）：由Selmer群的正合列\(0 \to E(K)/\ell E(K) \to \operatorname{Sel}^{(\ell)}(E/K) \to \Sha(E/K)[\ell] \to 0\)，若\(\operatorname{Sel}^{(\ell)}(E/K)\)的维数\(\leq 1\)，则\(E(K)/\ell E(K)\)至多一维。由于\(P_K\)为非挠点提供了一维，故\(\operatorname{rank}E(K) = 1\)。对全部\(\ell\)取归纳极限即得\(\Sha(E/K)\)有限（因为其\(\ell\)-准素部分对每个\(\ell\)有界）。\(\blacksquare\)

Kolyvagin的Euler系方法是20世纪数论最深刻的创造之一。其基本思想是：利用Heegner点生成一族满足互反律的上同调类，这些类的相容性条件迫使Selmer群具有强刚性------大得不至于平凡，小得不至于无界。Euler系方法后来被Rubin推广到CM椭圆曲线、被Kato推广到模曲线的\(K\)-理论Euler系，成为研究BSD猜想和Iwasawa理论主猜想的标准工具。值得注意的是，Kolyvagin的论证部分依赖于\(P_K\)为非挠点的假设------这正是Gross-Zagier定理通过与\(L'(E/K,1)\)的联系所保证的。

\subsection{88.4 BSD猜想的部分结果}\label{bsdux731cux60f3ux7684ux90e8ux5206ux7ed3ux679c}

\textbf{定理 88.4}（BSD猜想的秩\(\leq 1\)情形------已证明）：设\(E/\mathbb{Q}\)为椭圆曲线（导子\(N\)），\(\operatorname{ord}_{s=1}L(E,s) = r\)。则： 1. 若\(r = 0\)（即\(L(E,1) \neq 0\)），则\(\operatorname{rank}E(\mathbb{Q}) = 0\)且\(\Sha(E/\mathbb{Q})\)有限（Kolyvagin 1988）； 2. 若\(r = 1\)（即\(L(E,1) = 0\)但\(L'(E,1) \neq 0\)），则\(\operatorname{rank}E(\mathbb{Q}) = 1\)（Gross-Zagier + Kolyvagin）。

\textbf{证明框架}：

情形\(r=0\)：若\(L(E,1) \neq 0\)，取辅助虚二次域\(K\)使得\(L(E/K, 1) \neq 0\)（由Waldspurger定理或通过选取\(K\)使得符号根为\(+1\)保证）。存在\(K\)的二次特征\(\chi\)使得\(f \otimes \chi\)的中心\(L\)-值非零。由Gross-Zagier定理的变体（非零中心值情形），\(\operatorname{rank}E(K) = 0\)。由于\(E(\mathbb{Q}) \subseteq E(K)\)，\(\operatorname{rank}E(\mathbb{Q}) = 0\)。Kato的Euler系（基于\(K_2\)和模单位的Beilinson元素）进一步给出了\(\Sha(E/\mathbb{Q})\)的有限性。

情形\(r=1\)：若\(L(E,1) = 0\)且\(L'(E,1) \neq 0\)，存在虚二次域\(K\)使符号根为\(-1\)，则\(L(E/K, 1) = 0\)（中心值消失）。由Gross-Zagier定理，\(L'(E/K, 1) \neq 0\)（因为\(L(E,s)L(E \otimes \chi, s) = L(E/K, s)\)的导数非零与\(L'(E,1) \neq 0\)的等价性）。Heegner点\(P_K\)的高度非零，故\(P_K \in E(\mathbb{Q})\)非挠（通过Galois迹落入\(\mathbb{Q}\)）。Kolyvagin定理给出\(\operatorname{rank}E(\mathbb{Q}) = 1\)。\(\blacksquare\)

\textbf{定理 88.5}（Coates-Wiles定理，1977------CM情形）：设\(E/\mathbb{Q}\)为CM椭圆曲线。若\(L(E,1) \neq 0\)，则\(E(\mathbb{Q})\)有限（秩为\(0\)）。这是BSD猜想在解析秩\(0\)情形下的第一个非平凡结果，早于Gross-Zagier和Kolyvagin的工作。Coates-Wiles的证明利用了Iwasawa理论和CM椭圆曲线的\(p\)-进\(L\)函数，通过分析\(L(E,1)\)与Selmer群的\(p\)-进不变量之间的联系，得到了秩为零的结论。

\subsection{88.5 Congruent Number问题与Tunnell定理}\label{congruent-numberux95eeux9898ux4e0etunnellux5b9aux7406}

\textbf{定义 88.3}（Congruent Number）：正整数\(n\)称为\textbf{同余数}（congruent number），若存在有理边直角三角形面积为\(n\)。等价地，\(n\)为同余数当且仅当椭圆曲线\(E_n: y^2 = x^3 - n^2 x\)的秩\(\operatorname{rank}E_n(\mathbb{Q}) \geq 1\)。

\textbf{定理 88.6}（Tunnell定理，1983）：设\(n\)为无平方因子的正整数。定义 \[A(n) = \#\{(x,y,z) \in \mathbb{Z}^3 : n = 2x^2 + y^2 + 32z^2\},\] \[B(n) = \#\{(x,y,z) \in \mathbb{Z}^3 : n = 2x^2 + y^2 + 8z^2\},\] \[C(n) = \#\{(x,y,z) \in \mathbb{Z}^3 : n = 8x^2 + 2y^2 + 64z^2\} \quad (\text{$n$为奇数}),\] \[D(n) = \#\{(x,y,z) \in \mathbb{Z}^3 : n = 8x^2 + 2y^2 + 16z^2\} \quad (\text{$n$为奇数}).\] 若\(n\)为同余数，则\(A(n) = B(n)\)（\(n\)为偶数时）或\(C(n) = 2D(n)\)（\(n\)为奇数时）。反之，若BSD猜想对\(E_n\)在秩\(0\)情形下成立，则这些等式反过来蕴含\(n\)为非同余数。

\textbf{证明}：核心是将\(E_n\)与权\(3/2\)的模形式联系起来。\(E_n\)是椭圆曲线\(E_1: y^2 = x^3 - x\)的二次扭\(E_1^{(n)}\)。\(E_1\)的Hasse-Weil \(L\)函数\(L(E_1, s)\)对应一个权\(2\)模形式\(f\)。由Shimura提升，存在权\(3/2\)模形式\(g = \sum b_n q^n\)使得\(L(E_1^{(n)}, s)\)的中心值\(L(E_1^{(n)}, 1)\)与\(b_n\)通过Waldspurger公式联系起来：\(L(E_n, 1) = L(E_1^{(n)}, 1) = C \cdot |b_n|^2 / \sqrt{n}\)（\(C\)为非零常数，与\(n\)无关）。

Tunnell具体构造了这些权\(3/2\)模形式。对于\(n\)为偶数（\(n=2m\)）的情形，相关模形式为\(\theta_2\theta_4\theta_8\)（三个Jacobi \(\theta\)-函数的乘积），其Fourier系数为\(b_{2m} = A(m) - B(m)\)。对于\(n\)为奇数，相关模形式为\(\theta_2\theta_4\theta_{16}\)，系数为\(b_n = C(n) - 2D(n)\)。由Waldspurger公式，\(L(E_n, 1) = 0 \iff b_n = 0\)，即\(A(n) = B(n)\)（偶情形）或\(C(n) = 2D(n)\)（奇情形）。

若\(n\)为同余数，则\(\operatorname{rank}E_n(\mathbb{Q}) \geq 1\)。若BSD猜想成立，\(L(E_n, 1) = 0\)，故\(b_n = 0\)，这给出Tunnell条件。反之，若Tunnell条件成立但\(n\)非BSD-同余数，仍需假设BSD来排除这种可能性。\(\blacksquare\)

Tunnell定理将千年古老的Congruent Number问题（阿拉伯数学家Al-Karaji约公元1000年首次研究）转化为BSD猜想的一个具体检验案例：若BSD猜想成立，则存在判定\(n\)是否为同余数的有限算法（检查\(A(n) = B(n)\)或\(C(n) = 2D(n)\)）。反过来，这也使Congruent Number问题成为BSD猜想最重要的具体检验场之一。目前所有数值证据均支持Tunnell准则的正确性------即BSD猜想对这类椭圆曲线在秩\(0\)情形下的预测与全部计算一致。

\textbf{定理 88.7}（Heegner点方法与BSD的进一步进展）：通过将Kolyvagin的Euler系方法与Iwasawa理论结合（Kato 2004, Skinner-Urban 2014），BSD猜想在大量新情形下取得了进展。特别地，若\(E\)有非零的模参数化（即模性定理保证的\(\deg\varphi > 0\)），且\(L(E,1) \neq 0\)，则\(E(\mathbb{Q})\)有限且\(\Sha(E/\mathbb{Q})\)的\(\ell\)-准素部分有限（对所有素数\(\ell\)）。此外，Bhargava-Skinner-Zhang（2014-2021）利用平均秩方法和Selmer群的统计分析，证明了BSD猜想对''大多数''椭圆曲线（按高排序）成立------即具有解析秩\(0\)和\(1\)的椭圆曲线在全体椭圆曲线中的比例为100\%（条件于某些统计假设）。

这些进展表明：BSD猜想的核心困难集中在解析秩\(r \geq 2\)的情形。此时\(L(E,s)\)的高阶零点和Heegner点的退化使Kolyvagin方法无法直接应用，需要全新的工具------可能涉及更高维的算术几何对象（如Galois表示的变形理论、motivic上同调和高阶\(K\)-群的Euler系）。

\hrulefill

\hrulefill

\newpage

\chapter{08-代数几何与数论/02-数论/05-岩泽与自守.md}\label{ux4ee3ux6570ux51e0ux4f55ux4e0eux6570ux8bba02-ux6570ux8bba05-ux5ca9ux6cfdux4e0eux81eaux5b88.md}

\section{第299章：岩泽理论}\label{ux7b2c299ux7ae0ux5ca9ux6cfdux7406ux8bba}

岩泽理论（Iwasawa theory）由岩泽健吉（Kenkichi Iwasawa）于 1950 年代末创立，是将 \(p\)-进分析和 Galois 模论应用于代数数论中理想类群和椭圆曲线 Selmer 群增长问题的系统理论。其核心思想是：考虑数域 \(K\) 的 \(\mathbb{Z}_p\)-扩张 \(K_\infty/K\)（Galois 群 \(\Gamma \cong \mathbb{Z}_p\)），研究 \(p\)-部分理想类群 \(A_n = \operatorname{Cl}(K_n)[p^\infty]\) 的逆向极限 \(X = \varprojlim A_n\) 作为岩泽代数 \(\Lambda = \mathbb{Z}_p[[T]]\) 上的模的结构。岩泽基本定理给出了 \(|A_n| = p^{e_n}\) 的渐近公式 \(e_n = \mu p^n + \lambda n + \nu\)，其中 \(\mu, \lambda, \nu\) 为岩泽不变量。岩泽主猜想（Mazur-Wiles, 1984）将 \(\Lambda\)-模 \(X\) 的特征理想与 \(p\)-进 \(L\)-函数联系起来，是 Wiles 证明 Fermat 大定理的核心技术源头。本章将系统阐述 \(\mathbb{Z}_p\)-扩张、岩泽代数、岩泽基本定理和岩泽主猜想的完整理论。

\subsection{\texorpdfstring{236.1 \(\mathbb{Z}_p\)-扩张与岩泽代数}{236.1 \textbackslash mathbb\{Z\}\_p-扩张与岩泽代数}}\label{mathbbz_p-ux6269ux5f20ux4e0eux5ca9ux6cfdux4ee3ux6570}

\textbf{定义 236.1}（\(\mathbb{Z}_p\)-扩张）：设 \(K\) 为数域，\(p\) 为素数。\(K\) 的 \textbf{\(\mathbb{Z}_p\)-扩张} \(K_\infty/K\) 的 Galois 群拓扑同构于加法群 \(\mathbb{Z}_p\)，记 \(\Gamma = \operatorname{Gal}(K_\infty/K) \cong \mathbb{Z}_p\)。令 \(K_n\) 为 \(\Gamma^{p^n}\) 的固定域，则 \([K_n:K] = p^n\) 且 \(K_\infty = \bigcup_n K_n\)。每个数域至少有一个分圆 \(\mathbb{Z}_p\)-扩张。

\(\mathbb{Z}_p\)-扩张的概念是岩泽理论（Iwasawa theory）的几何基础。\(\Gamma \cong \mathbb{Z}_p\) 作为 \(p\)-进 Lie 群，其拓扑生成元 \(\gamma\) 给出 Galois 群上的连续作用。中间域 \(K_n\) 构成了 \(K_\infty\) 的''塔''：\(K = K_0 \subset K_1 \subset K_2 \subset \cdots \subset K_\infty\)，每一步 \([K_{n+1}:K_n] = p\)。这一塔状结构使得 \(p\)-进极限成为自然的工具：对每个中间域 \(K_n\) 上的算术不变量（如理想类群的 \(p\)-部分 \(A_n\)），考虑其关于范数映射的逆向极限 \(X = \varprojlim_n A_n\)，则 \(X\) 自然地成为 \(\Lambda = \mathbb{Z}_p[[\Gamma]]\) 上的紧模。分圆 \(\mathbb{Z}_p\)-扩张是最基本的例子：令 \(K = \mathbb{Q}\)，则 \(K_\infty = \mathbb{Q}(\zeta_{p^\infty})\) 是所有 \(p\) 幂次分圆域的并，其 Galois 群 \(\operatorname{Gal}(\mathbb{Q}(\zeta_{p^\infty})/\mathbb{Q}) \cong \mathbb{Z}_p^\times \cong \mathbb{Z}_p \times (\text{有限})\)，其中 \(\mathbb{Z}_p\) 部分给出 \(\mathbb{Q}\) 的分圆 \(\mathbb{Z}_p\)-扩张。对数域 \(K\)，其分圆 \(\mathbb{Z}_p\)-扩张 \(K_\infty^{\text{cyc}}\) 由 \(K_\infty^{\text{cyc}} = K \cdot \mathbb{Q}_\infty^{\text{cyc}}\) 给出。Leopoldt 猜想断言数域 \(K\) 的 \(\mathbb{Z}_p\)-扩张的个数为 \(r_2(K) + 1\)（\(r_2\) 为复位个数），这等价于 \(p\)-进 regulator 非零。

\textbf{定义 236.2}（岩泽代数）：拓扑环 \(\Lambda = \mathbb{Z}_p[[T]]\)（\(p\)-进形式幂级数环）称为\textbf{岩泽代数}。存在非典范同构 \(\Lambda \cong \varprojlim_n \mathbb{Z}_p[\Gamma/\Gamma^{p^n}]\)，\(1+T\) 对应 \(\Gamma\) 的拓扑生成元。\(\Lambda\) 是二维正则局部环（Krull 维数 2），为 Noether 整环且所有理想均为闭理想。

岩泽代数 \(\Lambda\) 是岩泽理论的核心代数对象。\(\Lambda\) 的结构定理（类似于有限生成模在 PID 上的结构定理）在此扮演关键角色：每个有限生成挠 \(\Lambda\)-模 \(M\) 拟同构于形如 \(\bigoplus_i \Lambda / (p^{\mu_i}, f_i(T)^{n_i})\) 的直和，其中 \(f_i(T)\) 为不可约的 distinguished 多项式（即首项系数为 1 且其余系数被 \(p\) 整除的多项式）。这一结构定理直接导出了岩泽基本定理中 \(e_n = \mu p^n + \lambda n + \nu\) 的渐近公式。\(\Lambda\) 的 Krull 维数为 2 反映了它同时承载了 \(p\)-进方向（对应于 \(p\) 的幂）和塔方向（对应于 \(T\) 的幂），这使得 \(\Lambda\) 上的模论同时编码了 \(p\)-进信息和 Galois 塔的渐进信息。\(\Lambda\) 与 \(\mathbb{Z}_p[[T]]\) 的对应使 Weierstrass 预备定理成为关键计算工具。

\subsection{236.2 岩泽基本定理}\label{ux5ca9ux6cfdux57faux672cux5b9aux7406}

\textbf{定义 236.3}：设 \(A_n = \operatorname{Cl}(K_n)[p^\infty]\)（理想类群的 \(p\)-Sylow 子群），阶为 \(p^{e_n}\)。定义逆向极限 \(X = \varprojlim_n A_n\)（转移映射为范数）。\(\Gamma\) 连续作用于 \(X\)，赋予 \(X\) 紧 \(\Lambda\)-模结构。

\textbf{定理 236.1}（岩泽基本定理，Iwasawa 1959）：\(X\) 是有限生成挠 \(\Lambda\)-模。由 \(\Lambda\) 的结构定理，存在互素的 distinguished 多项式 \(f_1 \mid f_2 \mid \cdots \mid f_r\) 及非负整数 \(\mu_1, \ldots, \mu_r\) 使 \[X \sim \bigoplus_{i=1}^r \Lambda / (p^{\mu_i}, f_i(T))\] （拟同构）。进而存在整数 \(\mu \geq 0, \lambda \geq 0, \nu \in \mathbb{Z}\)（\textbf{岩泽不变量}）使对所有充分大 \(n\)， \[e_n = \mu p^n + \lambda n + \nu\] 其中 \(\mu = \sum \mu_i\)，\(\lambda = \sum \deg f_i\)。\(\mu\) 为零时称扩张为非分歧（总分歧仅有 tame 部分）。

\textbf{证明}：由 \(\Lambda\) 的结构定理，\(X\) 拟同构于 \(\bigoplus_{i=1}^r \Lambda/(p^{\mu_i},f_i(T))\)。\(\Gamma^{p^n}\)-不变量的阶为 \(|X_{\Gamma^{p^n}}| = \prod_{i=1}^r |(\Lambda/(p^{\mu_i},f_i)))_{\Gamma^{p^n}}|\)。令 \(\omega_n = (1+T)^{p^n}-1\)，则 \(\Lambda/(p^\mu,f)\otimes_{\Lambda}\Lambda/\omega_n\Lambda \cong \mathbb{Z}_p[T]/(p^\mu,f(T),\omega_n(T))\)。对每个因子，由 Weierstrass 预备定理，\(f(T) = p^{\mu'}g(T)u(T)\)（\(g\) 为 distinguished 多项式，\(u\) 为单位）。\(\mathbb{Z}_p[T]/(p^\mu,f(T),\omega_n(T))\) 的阶为 \(p^{\mu p^n + \lambda n + \nu}\)（\(\lambda = \deg g\)，\(\nu\) 为常数）。取对数求和即得 \(e_n = \mu p^n + \lambda n + \nu\)。\(\blacksquare\)

\subsection{236.3 岩泽主猜想}\label{ux5ca9ux6cfdux4e3bux731cux60f3}

\textbf{定理 236.2}（岩泽主猜想，Mazur-Wiles 1984）：设 \(K/\mathbb{Q}\) 为 Abel 扩张，\(\chi\) 为偶 Dirichlet 特征。则 \(X\) 的特征理想 \(\operatorname{char}_\Lambda(X)\) 由 \(p\)-进 \(L\)-函数生成： \[\operatorname{char}_\Lambda(X) = (L_p(1-s, \omega \chi^{-1})) \quad \text{作为 } \Lambda \text{ 的理想}\] 等价地，\(p\)-进 \(L\)-函数的幂级数展开系数精确刻画了理想类群的岩泽不变量。

\textbf{证明}：Mazur-Wiles 利用模曲线 \(X_1(Np^\infty)\) 的平展上同调 \(H^1_{\text{ét}}(X_1(Np^\infty),\mathbb{Z}_p)\) 构造 Euler 系。Euler 系元素在 Galois 上同调中满足局部-整体兼容性，Kolyvagin 的 Euler 系方法将特征理想与 \(p\)-进 \(L\)-函数联系起来。具体地，\(p\)-进 \(L\)-函数 \(L_p(1-s,\omega\chi^{-1})\) 插值经典 \(L\)-值，其对 \(\chi\) 的除子关系与 Euler 系的一致性给出 \(\operatorname{char}_\Lambda(X) = (L_p)\)。\(\blacksquare\)

\textbf{注记 236.1}：Mazur 与 Wiles 利用模曲线 \(X_1(N p^\infty)\) 的平展上同调和 Euler 系理论证明。该定理是 Wiles 证明 Fermat 大定理的核心技术源头。Skinner-Urban（2010年代）推广至 \(GL(2)\) 自守形式情形。

\section{第300章：自守形式与模形式}\label{ux7b2c300ux7ae0ux81eaux5b88ux5f62ux5f0fux4e0eux6a21ux5f62ux5f0f}

自守形式是定义在约化代数群 \(G\) 上的高度对称函数，其理论是 Langlands 纲领的核心。模形式是最基本的自守形式（对应 \(G = \operatorname{GL}(2)\)），其 Fourier 系数编码了深刻的算术信息------椭圆曲线的 \(L\)-函数、Galois 表示和 Hecke 代数之间通过模形式建立起精确的联系（谷山-志村- Weil 猜想，即模性定理）。本章将系统阐述模形式的基本理论（全纯模形式、尖点形式、Hecke 算子与 Euler 积）、自守表示理论（\(L^2(\Gamma \backslash G)\) 的谱分解、自守形式的一般定义），以及 Langlands 函子性猜想的基本框架。这一理论在 Wiles 证明 Fermat 大定理和几何 Langlands 纲领中扮演着核心角色。

\subsection{237.1 模形式}\label{ux6a21ux5f62ux5f0f}

\textbf{定义 237.1}（模形式）：设 \(\mathbb{H} = \{z \in \mathbb{C} : \Im(z) > 0\}\)，\(\Gamma \leq \operatorname{SL}(2, \mathbb{Z})\) 为有限指标子群。全纯函数 \(f : \mathbb{H} \to \mathbb{C}\) 为权 \(k \in \mathbb{Z}\) 的\textbf{模形式}若： 1. \(f\!\left(\frac{az+b}{cz+d}\right) = (cz+d)^k f(z)\) 对所有 \(\begin{pmatrix} a & b \\ c & d \end{pmatrix} \in \Gamma\)，\(z \in \mathbb{H}\) 2. \(f\) 在 \(\Gamma\) 的各尖点处全纯（Fourier 展开 \(f(z) = \sum_{n=0}^\infty a_n q^n\)，\(q = e^{2\pi i z}\)，仅含 \(n \geq 0\) 项） 若所有尖点处 \(a_0 = 0\)，\(f\) 为\textbf{尖点形式}。\(M_k(\Gamma)\) 与 \(S_k(\Gamma)\) 记模形式与尖点形式空间。

\textbf{定理 237.1}（\(M_k(\Gamma)\) 的有限维性质）：\(M_k(\Gamma)\) 和 \(S_k(\Gamma)\) 为有限维 \(\mathbb{C}\)-向量空间。对 \(\Gamma = \operatorname{SL}(2,\mathbb{Z})\)： \[\dim M_k(\operatorname{SL}(2,\mathbb{Z})) = \begin{cases} \lfloor k/12 \rfloor & k \equiv 2 \pmod{12} \\ \lfloor k/12 \rfloor + 1 & k \not\equiv 2 \pmod{12} \end{cases} \quad (k \geq 0)\]

\textbf{证明}：紧 Riemann 面 \(X(\Gamma)\) 的亏格 \(g\) 由 Riemann-Hurwitz 公式给出。全纯线丛 \(\omega^{\otimes k}\)（\(\omega\) 为典范丛）的截面空间即 \(M_k(\Gamma)\)。由 Riemann-Roch 定理，\(\dim H^0(X,\omega^{\otimes k}) = \deg\omega^{\otimes k} - g + 1 + \dim H^1\)（\(k\ge 2\)）。\(\deg\omega = 2g-2 + \sum_P (1-1/e_P)\)（\(e_P\) 为分歧指数）。对 \(\operatorname{SL}(2,\mathbb{Z})\)，\(X(1)\) 为亏格 \(0\) 的 Riemann 球面，椭圆点 \(i\)（\(e=2\)）和 \(\rho=e^{2\pi i/3}\)（\(e=3\)），尖点 \(\infty\)（\(e=\infty\)）。\(\deg\omega = -2 + 1/2 + 2/3 + 1 = 1/6\)。\(\dim M_k = \max(0, k/6 + 1)\)（\(k\ge 2\) 偶），经逐情形验证即得公式。\(\blacksquare\)

\subsection{237.2 Hecke 算子与 Euler 积}\label{hecke-ux7b97ux5b50ux4e0e-euler-ux79ef}

\textbf{定义 237.2}（Hecke 算子）：对 \(n \geq 1\)，\textbf{Hecke 算子} \(T_n : M_k(\Gamma) \to M_k(\Gamma)\) 定义为 \[(T_n f)(z) = n^{k-1} \sum_{ad=n} \sum_{b \bmod d} d^{-k} f\!\left(\frac{az+b}{d}\right)\] 对 \(\Gamma = \operatorname{SL}(2,\mathbb{Z})\)（一般子群需调整）。\(\{T_n\}\) 相互交换，生成交换 \(\mathbb{C}\)-代数 \(\mathbb{T}\)（\textbf{Hecke 代数}）。若 \(f\) 为全体 \(T_n\) 的共同特征向量（\textbf{特征形式}）且归一化 \(a_1=1\)，则 \(T_n f = a_n f\)。

\textbf{定理 237.2}（Euler 积分解）：设 \(f \in S_k(\operatorname{SL}(2,\mathbb{Z}))\) 为 Hecke 特征形式，Fourier 系数 \(a_n\)。其 \(L\)-函数具有 Euler 乘积： \[L(s, f) = \sum_{n=1}^\infty \frac{a_n}{n^s} = \prod_{p} \frac{1}{1 - a_p p^{-s} + p^{k-1-2s}}\] 该乘积与级数在 \(\Re(s) > (k+1)/2\) 时绝对收敛，且 \(L(s,f)\) 可亚纯延拓至 \(\mathbb{C}\) 并满足函数方程 \(\Lambda(s,f) = i^k \Lambda(k-s,f)\)（\(\Lambda\) 为完备 \(L\)-函数）。

\textbf{证明}：\(\{T_n\}\) 交换且 \(T_{nm}=T_nT_m\)（\(\gcd(n,m)=1\)）。对 Hecke 特征形式 \(f\)，\(T_nf=a_nf\)，故 \(a_{nm}=a_na_m\)（\(\gcd(n,m)=1\)）。对素数幂，\(T_{p^{r+1}}=T_pT_{p^r}-p^{k-1}T_{p^{r-1}}\)，作用于 \(f\) 得 \(a_{p^{r+1}}=a_pa_{p^r}-p^{k-1}a_{p^{r-1}}\)。生成函数 \(\sum_{r\ge 0}a_{p^r}X^r\) 满足 \((1-a_pX+p^{k-1}X^2)\sum_{r\ge 0}a_{p^r}X^r=1\)（由递推式验证），故 \(\sum_{r\ge 0}a_{p^r}X^r=(1-a_pX+p^{k-1}X^2)^{-1}\)。令 \(X=p^{-s}\)，由积性：\(L(s,f)=\sum a_n n^{-s}=\prod_p\sum_{r\ge 0}a_{p^r}p^{-rs}=\prod_p(1-a_pp^{-s}+p^{k-1-2s})^{-1}\)。函数方程由 Mellin 变换和 \(f\) 在 Fricke 对合下的变换性质导出。\(\blacksquare\)

\subsection{237.3 模性定理与 Langlands 纲领}\label{ux6a21ux6027ux5b9aux7406ux4e0e-langlands-ux7eb2ux9886}

\textbf{定理 237.3}（模性定理，Wiles 1995 / BCDT 2001）：每条 \(\mathbb{Q}\) 上的椭圆曲线 \(E\) 是模的：存在 \(f \in S_2(\Gamma_0(N))\)（\(N = N_E\) 为导子）使对所有素数 \(p \nmid N\)，\(a_p(f) = p + 1 - \#E(\mathbb{F}_p)\)，且 \(L(s, E) = L(s, f)\)。

\textbf{证明}（概略）：Wiles 证明半稳定椭圆曲线的模性，Taylor-Wiles 系统使用 Galois 表示的形变理论。核心步骤：（1）构造 \(\rho_{E,3}:\operatorname{Gal}(\overline{\mathbb{Q}}/\mathbb{Q})\to GL_2(\mathbb{F}_3)\)，由 Langlands-Tunnell 定理知其为模；（2）证明 \(\rho_{E,3}\) 的万有形变环 \(R\) 与 Hecke 代数 \(\mathbb{T}\) 同构（\(R\cong\mathbb{T}\)），通过 patching 论证和 Euler 系；（3）\(R\cong\mathbb{T}\) 蕴含 \(\rho_{E,3}\) 的所有提升均为模，特别地 \(\rho_{E,p}\) 为模，故 \(E\) 为模。BCDT（Breuil-Conrad-Diamond-Taylor）将证明推广至所有椭圆曲线。\(\blacksquare\)

\textbf{注记 237.1}（Langlands 纲领的连接）：模性定理是 Langlands 函子性在 \(GL(2)\) 的特例。Langlands 纲领核心：\(GL_n(\mathbb{A}_{\mathbb{Q}})\) 的尖点自守表示 \(\pi\) 对应 \(n\) 维不可约 \(\ell\)-进 Galois 表示 \(\rho_\pi : \operatorname{Gal}(\overline{\mathbb{Q}}/\mathbb{Q}) \to GL_n(\overline{\mathbb{Q}}_\ell)\)，使 \(L(s, \pi) = L(s, \rho_\pi)\)。类域论为 \(n=1\) 情形（\(GL_1\) 自守表示 = Hecke 特征 = 1 维 Galois 特征）。\(n=2\) 的模性定理是迄今最深远的确证实例。

模性定理的深远意义在于它将两个看似无关的数学领域------椭圆曲线的算术和模形式的分析------通过 \(L\)-函数的桥梁统一起来。在模性定理之前，椭圆曲线的 \(L\)-函数 \(L(s, E)\) 是作为 Hasse-Weil zeta 函数的局部因子定义的，其解析延拓和函数方程完全未知。模性定理将 \(L(s, E)\) 与模形式的 \(L\)-函数 \(L(s, f)\) 等同，从而借助 Hecke 的经典理论（Mellin 变换和 Fricke 对合）自动获得了 \(L(s, E)\) 的解析延拓和函数方程。这一对应同时给出了 BSD 猜想（Birch-Swinnerton-Dyer 猜想）中 \(L(s, E)\) 在 \(s=1\) 处的解析行为的理论基础。Taylor-Wiles 方法的核心技术创新------``patching''论证------通过将不同素数层级的 Hecke 代数粘合在一起，建立了形变环与 Hecke 代数的同构 \(R \cong \mathbb{T}\)，这一技术后来被系统地发展为''模性提升定理''（modularity lifting theorem），成为证明更广泛的 Galois 表示模性的标准工具。在 Langlands 纲领的广阔图景中，模性定理铺平了通往 \(GL_n\)（\(n > 2\)）自守性的道路：Clozel、Harris、Taylor 和 Thorne 等人已将模性提升技术推广到 \(GL_n\) 情形，证明了全实域上正则代数尖点自守表示的 Galois 表示对应，以及 Sato-Tate 猜想的完整证明。

\section{第301章：Shimura 簇}\label{ux7b2c301ux7ae0shimura-ux7c07}

Shimura 簇是 Shimura Goro（志村五郎）于 1960 年代引入、由 Deligne 于 1970 年代系统公理化的一类特殊代数簇，它们是模曲线的高维推广。Shimura 簇携带丰富的算术结构，是现代自守形式理论、Langlands 纲领和算术几何的核心几何对象------通过其上同调，自守表示与 Galois 表示（motivic \(L\)-函数）之间建立了精确的对应。

\subsection{91.1 Shimura 数据与 Shimura 簇的构造}\label{shimura-ux6570ux636eux4e0e-shimura-ux7c07ux7684ux6784ux9020}

\textbf{定义 91.1}（Shimura 数据）：一个 \textbf{Shimura 数据}（Shimura datum）是一个对 \((G, X)\)，其中：

\begin{itemize}
\tightlist
\item
  \(G\) 是 \(\mathbb{Q}\) 上的连通简约代数群；
\item
  \(X\) 是 \(G(\mathbb{R})\) 在其上可逆作用的齐次空间（Hermite 对称域），满足以下 Deligne 公理（SV1-SV6）：
\end{itemize}

\textbf{(SV1)} \(G(\mathbb{R})\) 在 \(X\) 上的作用是传递的，且对某点 \(x \in X\)，稳定子群 \(K_x = \operatorname{Stab}_{G(\mathbb{R})}(x)\) 是 \(G(\mathbb{R})\) 的极大紧子群（精确到有限指数）；

\textbf{(SV2)} \(X\) 具有 \(G(\mathbb{R})\)-不变的复结构，且作为 Hermite 对称空间，\(X\) 上的每个点给出一个 Hodge 结构（通过 \(G \to GL(V)\) 的某个表示）；

\textbf{(SV3)} Adjoint 群 \(G^{\text{ad}}\) 没有定义在 \(\mathbb{Q}\) 上的紧单因子（即 \(G^{\text{ad}}\) 无 \(\mathbb{Q}\)-单的紧因子）；

\textbf{(SV4)} 权同态 \(w: \mathbb{G}_m \to G\)（定义在 \(\mathbb{Q}\) 上）满足某些兼容条件；

\textbf{(SV5)} 对 \(G^{\text{ad}}\) 的每个 \(\mathbb{Q}\)-单因子，其 \(\mathbb{R}\)-点是非紧的；

\textbf{(SV6)} 对 \(G(\mathbb{R})\) 中由 \(\operatorname{Lie} G(\mathbb{R})\) 中满足 \(J^2 = -\operatorname{id}\) 的复结构 \(J\) 的自然作用，\(X\) 可等同于 \(G(\mathbb{R})\) 的共轭类。

\textbf{定义 91.2}（Shimura 簇）：设 \((G, X)\) 是 Shimura 数据，\(K \subset G(\mathbb{A}_f)\) 是有限 Adèle 点的紧开子群（称为 \textbf{级结构}，level structure）。则 \textbf{Shimura 簇} 定义为局部对称空间（双陪集空间）

\[\operatorname{Sh}_K(G, X) = G(\mathbb{Q}) \backslash X \times G(\mathbb{A}_f) / K\]

其中 \(G(\mathbb{Q})\) 通过对角嵌入作用于 \(X \times G(\mathbb{A}_f)\)，\(K\) 作用于右因子 \(G(\mathbb{A}_f)\)。

当 \(K\) 充分小时（neat），\(G(\mathbb{Q})\) 无固定点的元素作用于 \(X \times G(\mathbb{A}_f) / K\)，此时 \(\operatorname{Sh}_K(G, X)\) 是光滑复流形。

\subsection{91.2 经典例子}\label{ux7ecfux5178ux4f8bux5b50}

Shimura 簇涵盖了数论和代数几何中许多最重要的模空间。以下是几个不同维度的典型例子，它们构成了 Shimura 簇理论的''试验田''：

\begin{enumerate}
\def\labelenumi{\arabic{enumi}.}
\item
  \textbf{模曲线（\(GL_2\) 型，\(\dim = 1\)）}：这是维数最低的非平凡 Shimura 簇。取 \(G = GL_2\)（或其导出群 \(SL_2\)），\(X = \mathbb{H}^{\pm}\)（上半平面和下半平面的并）。此时 Shimura 簇 \(\operatorname{Sh}_K(GL_2, \mathbb{H}^{\pm})\) 即为经典的模曲线 \(Y(N)\)（\(K\) 为主同余子群 \(\Gamma(N)\)）或 \(X_0(N)\)（\(K\) 为 \(\Gamma_0(N)\) 型子群）。这些模曲线参数化带有级结构的椭圆曲线，其上的 Hecke 对应与模形式理论紧密相关。模曲线的 Baily-Borel 紧化正是添加尖点得到的经典紧化 \(X(N)\)。
\item
  \textbf{Siegel 模簇（\(GSp_{2g}\) 型，\(\dim = g(g+1)/2\)）}：取 \(G = GSp_{2g}\)（辛相似群），\(X = \mathbb{H}_g\)（Siegel 上半空间，由 \(g \times g\) 对称复矩阵 \(Z\) 满足 \(\Im(Z) > 0\) 组成）。相应的 Shimura 簇 \(\mathcal{A}_g\) 是 \(g\) 维主极化 Abel 簇的模空间。\(\mathcal{A}_g\) 是 \(g(g+1)/2\) 维拟射影簇，其紧化涉及 toroidal 紧化和 Satake 紧化等深刻理论。\(\mathcal{A}_g\) 上的上同调实现了 \(GSp_{2g}\) 的自守表示，其与 Galois 表示的对应是 Langlands 纲领中 \(g > 1\) 情形的重要实例。
\item
  \textbf{Hilbert-Blumenthal 模簇（\(\operatorname{Res}_{F/\mathbb{Q}}GL_2\) 型，\(\dim = [F:\mathbb{Q}]\)）}：取 \(F\) 为全实域，\(G = \operatorname{Res}_{F/\mathbb{Q}} GL_2\)（Weil 限制），\(X = \mathbb{H}^{[F:\mathbb{Q}]}\)（多个上半平面的乘积）。相应的 Shimura 簇参数化带有 \(F\) 中实乘法的 Abel 簇，其维度为 \([F:\mathbb{Q}]\)。Hilbert 模形式理论在此 Shimura 簇上自然展开，其上的 Hecke 特征值与 \(F\) 的 Galois 表示的对应是 Wiles 证明 Fermat 大定理时 Taylor-Wiles 方法的核心几何背景。
\item
  \textbf{酉 Shimura 簇（\(U(n,1)\) 型）}：取 \(G\) 为 \(n+1\) 维 Hermite 空间的酉群 \(U(n,1)\)，\(X\) 为复 \(n\) 维球。相应的 Shimura 簇（称为 Picard 模簇或球商）参数化带有特定极化和自同态的 Abel 簇。酉 Shimura 簇在 Langlands 纲领的函子性提升中扮演关键角色，Rogawski 和 Clozel 等人系统研究了其上的自守表示与 Galois 表示。
\item
  \textbf{PEL 型 Shimura 簇}：最一般的显式构造是极化（Polarization）、自同态（Endomorphism）和级结构（Level）型 Shimura 簇（简称为 PEL 型）。取 \(B\) 为 \(\mathbb{Q}\) 上的半单代数，\(V\) 为 \(B\)-模，\(G\) 为 \(V\) 上保持某交错形式的酉相似群，\(X\) 为相应的 Hermite 对称空间。PEL 型 Shimura 簇参数化带有特定极化、自同态和级结构的 Abel 概形，包含了上述所有例子作为特例，其整模型理论由 Kottwitz 和 Lan 等人系统发展。
\end{enumerate}

\subsection{91.3 Baily-Borel 定理与典范模型}\label{baily-borel-ux5b9aux7406ux4e0eux5178ux8303ux6a21ux578b}

\textbf{定理 91.1}（Baily-Borel 定理，1966 年）：对任意 Shimura 数据 \((G, X)\) 和紧开子群 \(K \subset G(\mathbb{A}_f)\)，Shimura 簇 \(\operatorname{Sh}_K(G, X)\) 具有 \textbf{拟射影代数簇} 的自然结构------即 \(\operatorname{Sh}_K(G, X)\) 可映为一个拟射影复代数簇的解析化。

\textbf{证明}：\(X\) 为 Hermite 对称空间，配备 Bergman 度量 \(\omega\)。自守形式（权 \(k\) 且关于 \(K\) 不变的）构成 \(\operatorname{Sh}_K(G,X)\) 上线丛 \(\omega^{\otimes k}\) 的截面空间 \(H^0(\operatorname{Sh}_K,\omega^{\otimes k})\)。对充分大的 \(k\)（\(k\ge k_0\)），\(\omega^{\otimes k}\) 为极丰沛线丛（由 Kodaira 嵌入定理的判定条件）。Baily-Borel 证明：\(\dim H^0(\operatorname{Sh}_K,\omega^{\otimes k})\) 的渐近公式为 \(c\cdot k^{\dim X}\)（Hilbert-Samuel 多项式），且基点在 \(k\gg 0\) 时分离。因此 \(\operatorname{Sh}_K\) 可嵌入 \(\mathbb{P}^N\)（\(N = \dim H^0-1\)），其像为 Zariski 闭子簇（因自守形式满足代数关系------自守形式的代数结构）。此即 \(\operatorname{Sh}_K\) 的拟射影代数簇结构。\(\blacksquare\)

Baily-Borel 定理也给出了 Shimura 簇的紧化------\textbf{Baily-Borel 紧化}（或极小紧化）------通过在边界添加较低维的 Shimura 簇的商。

\textbf{定理 91.2}（典范模型，Canonical Model）：Shimura 簇 \(\operatorname{Sh}_K(G, X)\) 在\textbf{反射域}（reflex field）\(E(G, X)\) 上具有有理模型。反射域是 \(G\) 的共轭类的定义域，也是关联的零维 Shimura 簇的分裂域。

\textbf{证明}：Shimura 的互反律构造典范模型。核心是 CM 点（具有复乘的点的 \(\operatorname{Sh}_K\) 点）的 Galois 作用。每个 CM 点对应于 \(G\) 的某个环面 \(T\)，其 Galois 作用由 \(\operatorname{Gal}(\overline{E}/E)\to T(\mathbb{A}_f)/T(\mathbb{Q})\) 的类域论互反律映射给出。对任意 \(\sigma\in\operatorname{Gal}(\overline{E}/E)\)，\(\sigma\) 在 CM 点上的作用由互反律唯一确定。由于 CM 点在 \(\operatorname{Sh}_K\) 中 Zariski 稠密（由 Baily-Borel 定理和 CM 理论），Galois 作用可延拓到整个 \(\operatorname{Sh}_K\)，从而给出 \(E\) 上的有理模型 \(\mathcal{S}_K\)。\(\blacksquare\)

粗略地说，\(\operatorname{Sh}_K(G, X)\) 作为复代数簇可定义在数域 \(E(G, X)\) 上------即存在 \(E(G, X)\) 上的概型 \(\mathcal{S}_K\) 使得 \(\mathcal{S}_K \times_{E(G, X)} \mathbb{C} \cong \operatorname{Sh}_K(G, X)\)。该模型的构造由 Shimura 的互反律理论给出，利用了 CM 点（complex multiplication points）的 Galois 作用：CM 点对应于特殊自守表示，其上 Galois 群的作用由类域论（整体互反律）描述。

\subsection{91.4 上同调与 Langlands 纲领}\label{ux4e0aux540cux8c03ux4e0e-langlands-ux7eb2ux9886}

Shimura 簇的核心重要性在于：其 \textbf{Betti 上同调}（或有理 \(\ell\)-进上同调）实现了 Langlands 纲领中的 Galois 表示。

\textbf{定理 91.3}（Shimura 簇的上同调与自守表示）：设 \((G, X)\) 是 Shimura 数据，\(\xi\) 是 \(G\) 的代数表示。则总上同调群

\[H^*(\operatorname{Sh}_K(G, X), \mathcal{L}_\xi)\]

（\(\mathcal{L}_\xi\) 是由 \(\xi\) 确定的局部系统）分解为 \(G(\mathbb{A}_f)\)-模的直和，每个直和项同构于某个自守表示 \(\pi\) 的 \(\pi_f\) 部分的若干次副本。即

\[H^*(\operatorname{Sh}_K(G, X), \mathcal{L}_\xi) \cong \bigoplus_{\pi} m(\pi) \cdot \pi_f^K \otimes H^*(\mathfrak{g}, K_\infty, \pi_\infty \otimes \xi)\]

其中 \(m(\pi)\) 是 \(\pi\) 的重数（离散系列表示的重数有限）。

\textbf{证明}：由 Matsushima 公式，\(\operatorname{Sh}_K(G,X)\) 的 \(L^2\)-上同调分解为 \(G\) 的自守表示的直和。具体地，\(H^i(\operatorname{Sh}_K,\mathcal{L}_\xi) \cong \bigoplus_\pi m(\pi)\cdot(\pi_f)^K\otimes H^i(\mathfrak{g},K_\infty,\pi_\infty\otimes\xi)\)。此分解由 de Rham 复形与 \((\mathfrak{g},K_\infty)\)-上同调的等价性导出。\(m(\pi)\) 为 \(\pi\) 在离散系列中的重数（由 Arthur 迹公式理论保证有限性）。\(\blacksquare\)

\textbf{推论 91.4}（Galois 表示与 Langlands 对应）：\(\ell\)-进平展上同调 \(H_{\text{ét}}^i(\operatorname{Sh}_K(G, X)_{\overline{\mathbb{Q}}}, \mathbb{Q}_\ell)\) 作为 \(G(\mathbb{A}_f) \times \operatorname{Gal}(\overline{\mathbb{Q}} / E)\) 的表示实现了 Langlands 对应的几何版本：每个尖点自守表示 \(\pi\) 关联一个 Galois 表示 \(\rho_{\pi, \ell}: \operatorname{Gal}(\overline{\mathbb{Q}} / E) \to {}^L G(\mathbb{Q}_\ell)\)（到 \(L\)-群的同态）。

\subsection{91.5 Langlands-Kottwitz 方法与 Hasse-Weil Zeta 函数}\label{langlands-kottwitz-ux65b9ux6cd5ux4e0e-hasse-weil-zeta-ux51fdux6570}

\textbf{定理 91.5}（Kottwitz 的 Langlands-Kottwitz 方法，1990 年代）：设 \(\operatorname{Sh}_K(G, X)\) 是 Shimura 簇（带有光滑整模型）。其 Hasse-Weil zeta 函数（算术 zeta 函数）

\[\zeta(\operatorname{Sh}_K, s) = \prod_{x \in |\operatorname{Sh}_K|} \frac{1}{1 - N(x)^{-s}}\]

（\(x\) 遍历 \(\operatorname{Sh}_K\) 的闭点）可通过 Langlands 对应的 \(L\)-函数表达。具体地，Kottwitz 证明了

\[\zeta(\operatorname{Sh}_K, s) = \prod_{\pi} L(s - \frac{\dim \operatorname{Sh}_K}{2}, \pi, r)^{m(\pi)}\]

其中乘积遍历 \(G\) 的自守表示 \(\pi\)，\(r\) 是 \(L\)-群 \({}^L G\) 的某个自然表示（由 Shimura 簇的切丛的 \(G\)-模结构决定），\(L(s, \pi, r)\) 是关联的 \(L\)-函数。

\textbf{证明}：Kottwitz 的证明基于三个关键步骤。（1）Grothendieck-Lefschetz 迹公式：\(\#\operatorname{Sh}_K(\mathbb{F}_q) = \sum_i (-1)^i \operatorname{Tr}(\operatorname{Frob}_q \mid H_{\text{ét}}^i(\operatorname{Sh}_K,\mathbb{Q}_\ell))\)。将 Shimura 簇的 \(\ell\)-进上同调通过 Matsushima 公式分解为自守表示的贡献。（2）稳定迹公式：Arthur-Selberg 迹公式的几何侧（轨道积分）与谱侧（自守表示）的等式，将几何迹（点计数）与自守 \(L\)-函数联系起来。（3）基本引理（Ngo，2010）：局部轨道积分在稳定化后等于群的 \(\kappa\)-轨道积分的稳定化，确保整体迹公式两端可逐项比较。结合三步得 \(\zeta(\operatorname{Sh}_K,s) = \prod_\pi L(s-\dim\operatorname{Sh}_K/2,\pi,r)^{m(\pi)}\)。\(\blacksquare\)

\subsection{91.6 Shimura 簇的算术应用}\label{shimura-ux7c07ux7684ux7b97ux672fux5e94ux7528}

\begin{enumerate}
\def\labelenumi{\arabic{enumi}.}
\item
  \textbf{志村簇上的 CM 点的 Galois 作用}：CM 点生成 Shimura 簇中 0 维的子簇，其 Galois 作用由类域论完全描述。这是 Hilbert 第 12 问题（Kronecker 青春之梦）的高维推广------虚二次域的显式类域论（复乘理论）推广到更一般的 Shimura 数据。
\item
  \textbf{志村簇的模型与模性提升定理}：\(GL(2)\)-型 Shimura 簇（模曲线）的整模型理论（Deligne-Rapoport、Katz-Mazur）直接应用于 Wiles 对 Fermat 大定理的证明。高维推广的整模型理论（Kisin、Scholze 等）是 \(p\)-进自守形式理论的核心。
\item
  \textbf{志村簇和 Langlands 纲领}：Scholze（2013 年）利用 perfectoid 空间理论证明了拟射影 Shimura 簇的 Hodge-Tate 周期映射的某些性质，并构造了志村簇的 \(p\)-进一致化（类似于 Cerednik-Drinfeld 定理的高维推广）。
\item
  \textbf{André-Oort 猜想}：关于 Shimura 簇中特殊点（CM 点）的 Zariski 闭包------特殊子簇正好是较低维的 Shimura 子簇。该猜想由 Pila-Zannier（2009 年）使用 o-极小性方法在积模曲线情形证明，并由 Klingler-Ullmo-Yafaev 和 Pila-Tsimerman 在一般情形取得重大进展。
\end{enumerate}

\hrulefill

\hrulefill

\hrulefill

\hrulefill

\hrulefill

\hrulefill

\hrulefill

\section{第302章：类域论}\label{ux7b2c302ux7ae0ux7c7bux57dfux8bba}

类域论是代数数论的巅峰成就，描述了数域（或局部域）的 Abel 扩张与其理想类群（或乘法群）之间的深刻对偶。类域论起源于 Kronecker 的''青春之梦''（Jugendtraum）------希望通过解析函数在特殊点的取值显式构造 Abel 扩张------并在 Hilbert 第 12 问题中得到了纲领性的表述。本章将系统阐述局部类域论（局部 Artin 互反律、Lubin-Tate 形式群）和整体类域论（Adèle-Idèle 形式、整体 Artin 互反律、Chebotarev 密度定理），以及类域论与 \(L\)-函数、Galois 上同调（Tate-Nakayama 定理）和 Langlands 纲领的深刻联系。

\subsection{235.1 历史动机与基本问题}\label{ux5386ux53f2ux52a8ux673aux4e0eux57faux672cux95eeux9898}

类域论是代数数论的最高成就，描述了数域（或局部域）的 Abel 扩张与其理想类群（或乘法群）之间的深刻对应。

\textbf{定义 235.1}（Kronecker 青春之梦）：Kronecker（1853）在其给 Dedekind 的信中提出：代数数域的 Abel 扩张能否由超越函数（模函数、椭圆函数）的特殊值显式构造？具体而言，虚二次域的每个 Abel 扩张应包含在由椭圆模函数在复乘（CM）点处的特殊值生成的扩域中。这一愿景驱动了类域论整整一个世纪的发展。

\textbf{定理 235.1}（Hilbert 第 12 问题陈述，1900）：将 Kronecker 的青春之梦推广到任意代数数域：对任意数域 \(K\)，寻找解析函数使其在特殊点的取值生成 \(K\) 的 Abel 扩张。类域论从存在性角度回答了此问题，显式构造至今仅对 \(\mathbb{Q}\)（分圆域，Kronecker-Weber 定理）和虚二次域（复乘理论，高木贞治-Artin-志村）完全解决。

\textbf{证明}（已解决的情形）：对 \(K = \mathbb{Q}\)，指数函数 \(e^{2\pi i z}\) 在有理点 \(z \in \mathbb{Q}\) 的取值 \(e^{2\pi i a/n} = \zeta_n^a\) 生成全部分圆域 \(\mathbb{Q}(\zeta_n)\)，即 \(\mathbb{Q}^{\text{ab}} = \bigcup_{n \geq 1} \mathbb{Q}(\zeta_n)\)（Kronecker-Weber 定理 235.6）。对虚二次域 \(K = \mathbb{Q}(\sqrt{-d})\)，椭圆模函数 \(j(\tau)\) 在 CM 点 \(\tau \in K \cap \mathbb{H}\) 处的取值 \(j(\tau)\) 生成 \(K\) 的所有 Abel 扩张（复乘理论的主定理）。\(j(\tau)\) 为代数整数，其最小多项式给出 \(K\) 的类域。对一般数域，Shimura 簇上的特殊点理论（志村-谷山-志村五郎）提供了部分答案，但完全显式构造仍为开放问题。\(\blacksquare\)

\subsection{235.2 局部类域论}\label{ux5c40ux90e8ux7c7bux57dfux8bba}

\textbf{定义 235.2}（局部域）：\textbf{局部域}是连通局部紧的拓扑域。包括 Archimedean 的 \(\mathbb{R}\) 和 \(\mathbb{C}\)，以及非 Archimedean 的 \(\mathbb{Q}_p\) 的有限扩张（\(p\)-进域）。\(\mathbb{Q}_p\) 是 \(\mathbb{Q}\) 在 \(p\)-进度量 \(|x|_p = p^{-v_p(x)}\) 下的完备化，\(\mathbb{Z}_p = \{x \in \mathbb{Q}_p : |x|_p \leq 1\}\)。

\textbf{定义 235.3}（局部互反映射）：设 \(K\) 为局部域，\(K^{\text{ab}}\) 为 \(K\) 的极大 Abel 扩张。局部类域论的核心是典范连续同态 \[\rho_K : K^\times \longrightarrow \operatorname{Gal}(K^{\text{ab}}/K)\] 满足 \(\rho_K\) 的核为 \(K^\times\) 中所有范数子群的交，像在 \(\operatorname{Gal}(K^{\text{ab}}/K)\) 中稠密。

\textbf{定理 235.2}（局部 Artin 互反律）：对 \(K\) 的有限 Abel 扩张 \(L/K\)，\(\rho_K\) 诱导同构 \[K^\times / N_{L/K}(L^\times) \cong \operatorname{Gal}(L/K)\] 若 \(K\) 为非 Archimedean 的 \(p\)-进域，\(\rho_K(\mathcal{O}_K^\times)\) 为惯性子群，且 \(\rho_K(\pi_K)\)（\(\pi_K\) 为素元）限制在极大无分歧扩张上为 Frobenius 自同构 \(\operatorname{Frob}_K\)。

\textbf{证明}：构造局部互反律的核心工具是 Galois 上同调和 Brauer 群。设 \(G = \operatorname{Gal}(L/K)\)。由 Tate 定理，\(\widehat{H}^0(G, L^\times) \cong K^\times / N_{L/K}(L^\times)\) 且 \(\widehat{H}^{-2}(G, \mathbb{Z}) \cong G^{\text{ab}} = G\)（因 \(G\) 交换）。Tate-Nakayama 定理给出同构 \(\widehat{H}^0(G, L^\times) \cong \widehat{H}^{-2}(G, \mathbb{Z})\)，即 Artin 映射。具体构造：对 \(a \in K^\times\)，取 \(b \in L^\times\) 使 \(a = N_{L/K}(b)\)，定义互反律符号。非分歧情形：若 \(L/K\) 无分歧，\(\operatorname{Gal}(L/K) \cong \operatorname{Gal}(\ell/k)\)（剩余域扩张），\(\rho_K(\pi_K)\) 对应 \(\operatorname{Frob}_{L/K}\)。惯性子群对应 \(\rho_K(\mathcal{O}_K^\times)\)，由分歧理论 \(\mathcal{O}_K^\times\) 的像固定极大无分歧子扩张。\(\blacksquare\)

\textbf{定理 235.3}（局部类域论的存在性定理）：\(K^\times\) 的有限指数开子群与 \(K\) 的有限 Abel 扩张的范数子群一一对应。这给出了 \(K\) 的 Abel 扩张的完全分类。

\textbf{证明}：设 \(K\) 为 \(p\)-进局部域。利用 Lubin-Tate 形式群构造：取素元 \(\pi\)，定义形式群律 \(F_\pi(X,Y) \in \mathcal{O}_K[[X,Y]]\) 满足 \([\pi]_{F_\pi}(T) \equiv \pi T + T^q \pmod{\pi}\)。\(F_\pi\) 的 \(\pi^n\)-挠点模构成 \(K\) 的极大完全分歧 Abel 扩张 \(K_\pi\)，\(\operatorname{Gal}(K_\pi/K) \cong \mathcal{O}_K^\times\)。对 \(K^\times\) 的任意有限指数开子群 \(H\)，存在 \(n\) 使 \(U_K^{(n)} \subseteq H\)，构造 Lubin-Tate 扩张的子扩张 \(L\) 使得 \(N_{L/K}(L^\times) = H\)。反之，对任意有限 Abel 扩张 \(L/K\)，范数子群 \(N_{L/K}(L^\times)\) 是 \(K^\times\) 的有限指数开子群，且由 Hasse-Arf 定理知其导子有限。一一对应由 Galois 理论保证。\(\blacksquare\)

\textbf{注记 235.1}（Lubin-Tate 形式群，1965）：Lubin 和 Tate 给出了局部类域论的显式构造。设 \(K\) 为 \(p\)-进域，取素元 \(\pi\)，构造一维交换形式群 \(F_\pi\) 使 \(\mathcal{O}_K\) 通过形式自同构作用。\(F_\pi\) 的 \(\pi^n\)-挠点生成 \(K\) 的极大完全分歧 Abel 扩张 \(K_\pi\)------这完美类比于分圆域的局部情形：\(\hat{\mathbb{G}}_m\) 的 \(p^n\)-挠点生成 \(\mathbb{Q}_p(\zeta_{p^n})\)。

\subsection{235.3 整体类域论：Idèle 形式}\label{ux6574ux4f53ux7c7bux57dfux8bbaiduxe8le-ux5f62ux5f0f}

\textbf{定义 235.4}（Adèle 环与 Idèle 群）：数域 \(K\) 的 \textbf{Adèle 环}为关于整数环的限制直积 \[\mathbb{A}_K = \prod_v' K_v = \{(x_v)_v : x_v \in \mathcal{O}_{K_v} \text{ 对几乎所有非 Archimedean 位 } v\}\] 其中 \(v\) 遍历 \(K\) 的所有位（赋值的等价类）。\textbf{Idèle 群}为 \(\mathbb{I}_K = \mathbb{A}_K^\times\)，配备限制直积拓扑。\(K^\times\) 通过对角嵌入 \(a \mapsto (a, a, \ldots)\) 成为 \(\mathbb{I}_K\) 的离散子群。

\textbf{定义 235.5}（Idèle 类群）：\(K\) 的 \textbf{Idèle 类群}定义为 \(C_K = \mathbb{I}_K / K^\times\)。这是整体类域论的核心代数对象。

\textbf{定理 235.4}（整体 Artin 互反律）：设 \(K\) 为数域。存在典范连续满同态 \[\rho_K : C_K \longrightarrow \operatorname{Gal}(K^{\text{ab}}/K)\] 对 \(K\) 的有限 Abel 扩张 \(L/K\)，诱导同构 \[C_K / N_{L/K}(C_L) \cong \operatorname{Gal}(L/K)\] 且在每个局部化上与局部互反律兼容：\(\rho_K|_{K_v^\times} = \rho_{K_v}\)。

\textbf{证明}：由所有局部互反映射 \(\rho_{K_v} : K_v^\times \to \operatorname{Gal}(K_v^{\text{ab}}/K_v)\) 的限制直积，定义 \(\rho = \prod_v \rho_{K_v} : \mathbb{I}_K \to \operatorname{Gal}(K^{\text{ab}}/K)\)。由局部互反律的相容性，\(\rho\) 为连续同态。对 \(a \in K^\times\)，由整体积公式 \(\prod_v |a|_v = 1\) 及局部互反律的积公式，\(\rho(a) = 1\)，故 \(K^\times \subseteq \ker \rho\)。因此 \(\rho\) 通过商诱导 \(\rho_K : C_K = \mathbb{I}_K / K^\times \to \operatorname{Gal}(K^{\text{ab}}/K)\)。满性的证明依赖 Chebotarev 密度定理：\(\operatorname{Gal}(K^{\text{ab}}/K)\) 中任意元素可被 Frobenius 自同构逼近，而 Frobenius 自同构恰为 idèle 的像。核的描述：\(\ker \rho_K = \bigcap_{L/K \text{ 有限 Abel}} N_{L/K}(C_L)\)，其中 \(N_{L/K}: C_L \to C_K\) 为范数映射。对每个有限 Abel 扩张 \(L/K\)，由 Tate 上同调：\(\widehat{H}^0(\operatorname{Gal}(L/K), C_L) \cong C_K / N_{L/K}(C_L)\)，且 \(\widehat{H}^{-2}(\operatorname{Gal}(L/K), \mathbb{Z}) \cong \operatorname{Gal}(L/K)^{\text{ab}}\)。Tate-Nakayama 定理给出同构 \(C_K / N_{L/K}(C_L) \cong \operatorname{Gal}(L/K)\)，即互反同构。\(\blacksquare\)

\textbf{定理 235.5}（Hilbert 类域）：设 \(H_K\) 为 \(K\) 的 Hilbert 类域------\(K\) 的极大无处非分歧 Abel 扩张。Artin 映射诱导同构 \[\operatorname{Cl}(K) \cong \operatorname{Gal}(H_K/K)\] 其中 \(\operatorname{Cl}(K)\) 为 \(K\) 的理想类群。特别地，\([H_K : K] = h_K\)（类数），且 \(K\) 的素理想在 \(H_K\) 中分裂当且仅当该理想为主理想。

\textbf{证明}：由 Artin 互反律，\(\rho_K: C_K \to \operatorname{Gal}(K^{\text{ab}}/K)\) 为连续满同态。\(H_K\) 的定义要求所有素理想均非分歧，即在 idèle 层面 \(N_{H_K/K}(C_{H_K}) \supseteq \prod_{\mathfrak{p}} U_{\mathfrak{p}}\)（所有局部单位群的乘积）。因此 \(C_K / N_{H_K/K}(C_{H_K}) \cong \mathbb{I}_K / (\prod U_{\mathfrak{p}} \cdot K^\times)\)。右侧恰为理想类群 \(\operatorname{Cl}(K)\)（因 \(\mathbb{I}_K / \prod U_{\mathfrak{p}} \cong \bigoplus \mathbb{Z}\) 为理想群，商去主理想群 \(K^\times\) 即得类群）。由互反律，此商群同构于 \(\operatorname{Gal}(H_K/K)\)。素理想 \(\mathfrak{p}\) 在 \(H_K\) 中完全分裂 \(\iff\) \(\operatorname{Frob}_{\mathfrak{p}} = 1\) \(\iff\) \(\mathfrak{p}\) 在 Artin 映射下的像平凡 \(\iff\) \(\mathfrak{p}\) 为主理想。\(\blacksquare\)

\textbf{定理 235.6}（Kronecker-Weber 定理）：\(\mathbb{Q}\) 的每个有限 Abel 扩张均包含于某个分圆域 \(\mathbb{Q}(\zeta_n)\) 中。等价地， \[\mathbb{Q}^{\text{ab}} = \bigcup_{n \geq 1} \mathbb{Q}(\zeta_n)\]

\textbf{证明}：将整体 Artin 互反律应用于 \(K = \mathbb{Q}\)。Idèle 类群的计算给出 \[C_{\mathbb{Q}} \cong \hat{\mathbb{Z}}^\times \times \mathbb{R}_{>0}\] 且 \(\operatorname{Gal}(\mathbb{Q}^{\text{ab}}/\mathbb{Q}) \cong \hat{\mathbb{Z}}^\times\)。另一方面，\(\operatorname{Gal}(\bigcup \mathbb{Q}(\zeta_n)/\mathbb{Q}) = \varprojlim \operatorname{Gal}(\mathbb{Q}(\zeta_n)/\mathbb{Q}) \cong \varprojlim (\mathbb{Z}/n\mathbb{Z})^\times = \hat{\mathbb{Z}}^\times\)。两者自然同一，故 \(\mathbb{Q}^{\text{ab}} = \bigcup \mathbb{Q}(\zeta_n)\)。\(\blacksquare\)

\subsection{235.4 类域论的主定理框架}\label{ux7c7bux57dfux8bbaux7684ux4e3bux5b9aux7406ux6846ux67b6}

\textbf{定理 235.7}（类域论的完整主定理）：设 \(K\) 为数域。 1. \textbf{存在性定理}：\(\mathbb{I}_K\) 中每个包含 \(K^\times\) 的有限指数闭子群 \(H\) 必为某有限 Abel 扩张 \(L/K\) 的范数子群 \(N_{L/K}(\mathbb{I}_L) K^\times\) 2. \textbf{互反同构}：Artin 映射诱导同构 \(\mathbb{I}_K / H \cong \operatorname{Gal}(L/K)\) 3. \textbf{导子判别定理}（Conductor-Discriminant Theorem）：\(L/K\) 的分歧信息完全由导子 \(\mathfrak{f}(L/K)\)（\(K\) 的一个整理想）描述：\(L\) 包含于 \(\bmod \mathfrak{f}\) 的射线类域中，且判别式 \(d_{L/K}\) 可表为导子在 Galois 特征值上的乘积。该定理回答了一个 Abel 扩张具体分歧于哪些素理想。

\textbf{证明}：(1) 存在性：设 \(H \subseteq \mathbb{I}_K\) 为包含 \(K^\times\) 的有限指数闭子群。考虑 Artin 映射 \(\rho_K: \mathbb{I}_K \to \operatorname{Gal}(K^{\text{ab}}/K)\)。由 Takagi 存在定理的 idèle 形式，\(\ker \rho_K|_{\text{有限商}}\) 恰对应某有限 Abel 扩张 \(L/K\) 的范数子群，即 \(H = N_{L/K}(\mathbb{I}_L) K^\times\)。(2) 互反同构：\(\rho_K\) 诱导同构 \(\mathbb{I}_K / N_{L/K}(\mathbb{I}_L) K^\times \cong \operatorname{Gal}(L/K)\)，由整体互反律保证。(3) 导子判别定理：定义 \(\mathfrak{f}(L/K) = \prod \mathfrak{p}^{f_{\mathfrak{p}}}\)，其中 \(f_{\mathfrak{p}}\) 是 \(K_{\mathfrak{p}}^\times\) 中使 \(U_{K_{\mathfrak{p}}}^{(n)} \subseteq N_{L_{\mathfrak{P}}/K_{\mathfrak{p}}}(L_{\mathfrak{P}}^\times)\) 的最小 \(n\)。由局部类域论，\(\mathfrak{f}(L/K)\) 整除判别式：\(d_{L/K} = \prod_{\chi} \mathfrak{f}_\chi\)（\(\chi\) 遍历不可约特征），由导子-判别式公式（Führerdiskriminantenproduktformel）保证。\(\blacksquare\)

\textbf{注记 235.2}（虚二次域的复乘理论）：设 \(K\) 为虚二次域。Kronecker 青春之梦的精确陈述为：\(K\) 的极大 Abel 扩张 \(K^{\text{ab}}\) 由 \(j\)-不变量在 \(K\) 的整环的 CM 点处取值及单位根的添加生成。这给出了类域论的完整显式构造。

\subsection{\texorpdfstring{235.5 Langlands 纲领：类域论作为 \(n=1\) 情形}{235.5 Langlands 纲领：类域论作为 n=1 情形}}\label{langlands-ux7eb2ux9886ux7c7bux57dfux8bbaux4f5cux4e3a-n1-ux60c5ux5f62}

\textbf{注记 235.3}（Langlands 纲领简述，1967）：类域论建立了数域的 \(1\) 维 Galois 表示（即 Abel 特征 \(\chi : \operatorname{Gal}(\overline{K}/K) \to \mathbb{C}^\times\)）与 Hecke 特征（\(\mathbb{I}_K / K^\times \to \mathbb{C}^\times\)）之间的一一对应。Langlands 纲领将这一对应推广至 \(n\) 维：

\begin{itemize}
\tightlist
\item
  \textbf{局部 Langlands 对应}：\(GL_n(K)\)（\(K\) 为局部域）的不可约光滑表示与 \(n\) 维 Weil-Deligne 表示一一对应
\item
  \textbf{整体 Langlands 对应}：\(GL_n(\mathbb{A}_K)\) 的尖点自守表示与 \(n\) 维不可约 \(\ell\)-进 Galois 表示一一对应
\item
  \textbf{函子性猜想}：\(L\)-群之间的任意同态 \({}^L H \to {}^L G\) 均诱导自守表示的提升
\end{itemize}

类域论恰为 \(n=1\) 的特殊情形：此时 \(GL_1\) 的自守表示即 Hecke 特征，\(L\)-群为 \(\operatorname{Gal}(\overline{K}/K)\) 的 Weil 形式，局部对应由局部互反律给出，整体对应由整体 Artin 互反律给出。高维推广是当前数论最核心的研究纲领。\(\blacksquare\)''

\subsection{235.6 射有限群与 Galois 群的完备化}\label{ux5c04ux6709ux9650ux7fa4ux4e0e-galois-ux7fa4ux7684ux5b8cux5907ux5316}

\textbf{定义 235.6}（射有限群）：拓扑群 \(G\) 为\textbf{射有限群}，若它同构于有限离散群的逆向极限 \(\varprojlim_i G_i\)。等价地，\(G\) 是紧、完全不连通的 Hausdorff 拓扑群。射有限群的闭子群、商群和直积均为射有限群。

\textbf{定理 235.8}（Galois 群的射有限结构）：设 \(L/K\) 为 Galois 扩张（未必有限）。则

\textbf{证明}：\(L/K\)的Galois群\(\operatorname{Gal}(L/K)\)定义为所有\(K\)-自同构组成的群。对每个有限Galois子扩张\(E/K\)（\(E\subseteq L\)），有满同态\(\operatorname{Gal}(L/K)\to\operatorname{Gal}(E/K)\)（限制映射）。这些映射构成逆向系统，其逆向极限\(\varprojlim_E\operatorname{Gal}(E/K)\)为射有限群。万有性质给出典范同态\(\operatorname{Gal}(L/K)\to\varprojlim_E\operatorname{Gal}(E/K)\)。\(L\)为所有有限Galois子扩张的并，故此同态为同构。Krull拓扑（由\(\operatorname{Gal}(L/E)\)（\(E/K\)有限Galois）为邻域基的拓扑）使\(\operatorname{Gal}(L/K)\)成为紧致Hausdorff全不连通拓扑群。\(\blacksquare\) \[\operatorname{Gal}(L/K) \cong \varprojlim_{i} \operatorname{Gal}(L_i/K)\] 其中 \(L_i\) 遍历 \(L/K\) 的全部有限 Galois 子扩张，赋予 Krull 拓扑（以有限 Galois 子扩张的 Galois 群为邻域基）。特别地， \[G_K = \operatorname{Gal}(\overline{K}/K) = \varprojlim_{[L:K]<\infty} \operatorname{Gal}(L/K)\] 是射有限群。\(\operatorname{Gal}(K^{\text{ab}}/K)\) 为 \(G_K\) 的极大交换商。局部域 \(K\) 的极大无分歧扩张的 Galois 群 \(\operatorname{Gal}(K^{\text{ur}}/K) \cong \hat{\mathbb{Z}}\)（Frobenius 拓扑生成）。基本例子：\(\operatorname{Gal}(\overline{\mathbb{F}}_p/\mathbb{F}_p) \cong \hat{\mathbb{Z}}\)；\(\operatorname{Gal}(\mathbb{Q}(\zeta_{p^\infty})/\mathbb{Q}) \cong \mathbb{Z}_p^\times\)。

\subsection{235.7 互反律的具体形式}\label{ux4e92ux53cdux5f8bux7684ux5177ux4f53ux5f62ux5f0f}

\textbf{定义 235.7}（Hilbert 符号）：设 \(K\) 为局部域（素元 \(\pi\)），\(\mu_n \subseteq K^\times\)。\textbf{Hilbert 符号}为配对 \[(\cdot, \cdot)_n : K^\times \times K^\times \longrightarrow \mu_n\] 由局部类域论定义：\((a, b)_n = 1 \iff a \in N_{K(\sqrt[n]{b})/K}(K(\sqrt[n]{b})^\times)\)（\(a\) 为范数）。满足双积性、斜对称性 \((a,b)_n(b,a)_n=1\) 及非退化性。

\textbf{定义 235.8}（幂剩余符号）：设 \(K\) 为数域，\(\mathfrak{p} \nmid n\) 为 \(K\) 的素理想。对 \(a \in K^\times\)（与 \(\mathfrak{p}\) 互素），\(n\) 次\textbf{幂剩余符号}由剩余域条件定义： \[\left(\frac{a}{\mathfrak{p}}\right)_n \equiv a^{(N(\mathfrak{p})-1)/n} \pmod{\mathfrak{p}},\quad \left(\frac{a}{\mathfrak{p}}\right)_n \in \mu_n\] 此为 Legendre 符号的 \(n\) 次推广。\textbf{Hilbert 互反律}（乘积公式）断言： \[\prod_v (a, b)_{n, v} = 1\qquad (\text{$v$ 遍历 $K$ 的所有位})\] 其中 \((\cdot, \cdot)_{n,v}\) 为 \(K_v\) 上的局部 Hilbert 符号。此即二次互反律的高次推广与整体表述。

\subsection{235.8 Takagi 存在定理}\label{takagi-ux5b58ux5728ux5b9aux7406}

\textbf{定理 235.9}（Takagi 存在定理，1920）：设 \(K\) 为数域。对 \(K\) 的每个模 \(\mathfrak{m}\) 及满足 \(K_{\mathfrak{m}, 1} \subseteq H \subseteq K^\times_{\mathfrak{m}}\) 的子群 \(H\)，存在唯一的有限 Abel 扩张 \(L/K\)（导子整除 \(\mathfrak{m}\)）使 \(\operatorname{Gal}(L/K) \cong \mathbb{I}_K / (H \cdot K^\times)\)。Idèle 形式：\(\mathbb{I}_K\) 中每个包含 \(K^\times\) 的有限指数闭子群对应唯一的有限 Abel 扩张 \(L/K\)。结合 Artin 互反律，Takagi 存在定理与 Artin 同构构成整体类域论的完备体系：前者保证扩张的存在性（``每个合理的子群都对应一个扩张''），后者提供显式 Galois 对应。

\textbf{证明}：核心是归纳构造。对 \(\mathfrak{m}=1\) 即 Hilbert 类域，存在性已由经典理想类群理论保证。对一般 \(\mathfrak{m}\)，考虑射线类群 \(C_{\mathfrak{m}} = \mathbb{I}_K / (K_{\mathfrak{m},1} \cdot K^\times)\)。将 \(C_1 \to C_{\mathfrak{m}} \to 1\) 的核视为 \((\mathcal{O}_K/\mathfrak{m})^\times / \operatorname{im}(\mathcal{O}_K^\times)\) 的商。利用 Kummer 扩张（对含足够多单位根的子域）和归纳，在基域上逐步添加射线类群覆盖所需的扩域。关键引理（第二不等式）：\([C_K : N_{L/K}(C_L)] \leq [L:K]\)。构造 \(L\) 使得 \(N_{L/K}(C_L) = H \cdot K^\times\)。唯一性由第一不等式及 Artin 映射的单性保证。由此，整体类域论的存在性部分完成，每个''可取模''的子群均对应一个类域。\(\blacksquare\)

\textbf{注记 235.4}：Takagi 的证明基于类数公式与范数指数定理的复杂归纳。该定理统一了此前所有已知的 Abel 扩张构造------分圆域（\(\mathbb{Q}\)）、Kummer 扩张（含足够多单位根的数域）、复乘理论（虚二次域）------并证明它们穷尽了所有 Abel 扩张。

Takagi 存在定理是 20 世纪数论最伟大的成就之一。在 Takagi 之前，类域论的主要进展包括：Kronecker-Weber 定理（\(\mathbb{Q}\) 的 Abel 扩张均为分圆域）、Hilbert 的类域猜想（每个数域的理想类群对应一个极大无分歧 Abel 扩张）、以及 Weber 和 Hilbert 关于导子（Führer）的初步理论。Takagi 在 1920 年发表的论文中，通过引入模 \(\mathfrak{m}\) 和射线类群 \(C_{\mathfrak{m}}\) 的概念，彻底解决了 Hilbert 第 9 问题（``最一般的互反律的证明''）和第 12 问题（``Kronecker 青春之梦''）的存在性部分。定理的构造性在于：给定模 \(\mathfrak{m}\) 和子群 \(H\)，类域 \(L\) 由 \(K\) 中满足某些范数条件的元素生成。特别地，Hilbert 类域 \(H_K\) 对应于 \(\mathfrak{m} = 1\) 和 \(H = K^\times\) 的情形，其 Galois 群同构于理想类群 \(\operatorname{Cl}(K)\)。射线的模 \(\mathfrak{m} = \mathfrak{m}_0 \mathfrak{m}_\infty\) 同时编码了有限部分（整理想）和无穷部分（实位），相应的射线类域包含了所有分歧仅出现在 \(\mathfrak{m}\) 的 Abel 扩张。Takagi 定理与 Artin 互反律（1927 年）的合璧标志着整体类域论的完成，其 idèle 形式由 Chevalley 在 1940 年给出，极大简化了陈述和证明。这一理论不仅是数论的丰碑，更是 Langlands 纲领的出发点------类域论是 \(GL_1\) 的 Langlands 对应。

\newpage

\chapter{08-代数几何与数论/02-数论/06-密码与格密码.md}\label{ux4ee3ux6570ux51e0ux4f55ux4e0eux6570ux8bba02-ux6570ux8bba06-ux5bc6ux7801ux4e0eux683cux5bc6ux7801.md}

\subsection{239.A 密码学补充}\label{a-ux5bc6ux7801ux5b66ux8865ux5145}

\begin{quote}
密码学是研究信息保密和认证的数学基础的科学，其数学内核包含数论（Euler 定理、离散对数）、代数（椭圆曲线群）、组合学和复杂性理论。本卷在数论 II（V17）、抽象代数（V8、V13）和计算理论（V31）的基础上，系统建立编码理论、对称密码、公钥密码、格密码与后量子密码以及密码协议的数学理论。
\end{quote}

\hrulefill

\hrulefill

\hrulefill

\hrulefill

\hrulefill

\hrulefill

\section{第303章：编码理论深化}\label{ux7b2c303ux7ae0ux7f16ux7801ux7406ux8bbaux6df1ux5316}

编码理论研究在有噪信道中可靠传输信息的数学方法，是 Shannon 信息论（V29）的直接应用和深化------信息论确定了通信的速率极限（信道容量），编码理论则构造具体可达这些极限的编码方案。编码理论经历了三个历史阶段：经典代数编码（Hamming 码，1950 年；Reed-Solomon 码，1960 年；BCH 码，1960 年）利用有限域上的多项式插值和有限几何构建了确定性的纠错方案；迭代编码（Turbo 码，1993 年；LDPC 码，1963 年 / 1996 年）引入概率迭代译码，首次在实践中逼近 Shannon 极限------差距仅 0.0045 dB；Polar 码（Arıkan，2009 年）则是第一个被严格证明可达任意对称二进制输入信道容量的构造性方案，其信道极化思想将编码理论推进到信息论的理论极限。本章在 Shannon 信息论（V29）的基础上，从线性码的代数基础出发，经 Reed-Solomon 和 BCH 码的代数构造，最后以 Turbo/LDPC/Polar 码收尾，展示编码理论从代数确定性到概率渐近性的方法论演进。

\subsection{257.1 错误检测与纠错}\label{ux9519ux8befux68c0ux6d4bux4e0eux7ea0ux9519}

\textbf{定义 257.1}（线性码 / Linear Code）：\([n, k, d]\) 线性码 \(C\) 是 \(\mathbb{F}_q^n\) 的 \(k\) 维子空间，极小距离 \(d = \min_{c \neq c' \in C} \operatorname{wt}(c - c')\)。\(C\) 可纠正 \(\lfloor (d-1)/2 \rfloor\) 个错误。

\textbf{定义 257.2}（生成矩阵与校验矩阵）：\textbf{生成矩阵} \(G \in \mathbb{F}_q^{k \times n}\)（\(C = \{mG : m \in \mathbb{F}_q^k\}\)）。\textbf{校验矩阵} \(H \in \mathbb{F}_q^{(n-k) \times n}\)（\(C = \{x : H x^T = 0\}\)）。\(HG^T = 0\)。

\textbf{定理 257.1}（Singleton 界）：\([n, k, d]\) 码满足 \(d \leq n - k + 1\)。达到此界的码称为 \textbf{MDS 码}（最大距离可分码）。

\textbf{证明}：设 \(C\) 为 \([n,k,d]\) 线性码。考虑投影映射 \(\pi: C \to \mathbb{F}_q^{n-d+1}\)，删除任意 \(d-1\) 个坐标。若 \(\pi\) 不是单射，则存在非零 \(c_1 - c_2 \in C\) 被投影为零，其重量 \(\leq d-1\)，与 \(d\) 为最小距离矛盾。故 \(\pi\) 为单射，\(\dim(C) = k \leq n - d + 1\)，即 \(d \leq n - k + 1\)。\(\blacksquare\)

\textbf{定理 257.2}（Hamming 界与 Gilbert-Varshamov 界）： - \textbf{Hamming 界}（球填充界）：\(q^k \sum_{i=0}^{\lfloor (d-1)/2 \rfloor} \binom{n}{i} (q-1)^i \leq q^n\) - \textbf{Gilbert-Varshamov 界}（存在性界）：存在 \([n, k, d]\) 码满足 \(q^{n-k} > \sum_{i=0}^{d-2} \binom{n-1}{i} (q-1)^i\)

\textbf{证明}：Hamming界：以每个码字为中心、半径\(t=\lfloor(d-1)/2\rfloor\)的Hamming球互不相交（否则两码字距离\(\leq 2t<d\)）。每个球包含\(\sum_{i=0}^t\binom{n}{i}(q-1)^i\)个向量，\(q^k\)个球覆盖\(q^k\sum_{i=0}^t\binom{n}{i}(q-1)^i\)个向量，不超过全空间\(\mathbb{F}_q^n\)的大小\(q^n\)，即得Hamming界。Gilbert-Varshamov界：逐列构造校验矩阵\(H\in\mathbb{F}_q^{(n-k)\times n}\)。第\(i\)列不能是前\(i-1\)列中任意\(d-2\)列的线性组合（否则存在重量\(\leq d-1\)的码字）。前\(i-1\)列至多生成\(\sum_{j=0}^{d-2}\binom{i-1}{j}(q-1)^j\)个向量，若\(q^{n-k}>\)此数则存在可选列。取\(i=n\)时条件即为GV界。\(\blacksquare\)

\subsection{257.2 经典代数编码}\label{ux7ecfux5178ux4ee3ux6570ux7f16ux7801}

\textbf{定义 257.3}（Reed-Solomon 码 / RS 码）：取 \(n \leq q-1\) 个互异点 \(\alpha_1, \ldots, \alpha_n \in \mathbb{F}_q\)。RS 码定义为 \(\{(f(\alpha_1), \ldots, f(\alpha_n)) : f \in \mathbb{F}_q[x], \deg f < k\}\)。RS 码是 MDS 码（\(d = n - k + 1\)）。

\textbf{定义 257.4}（BCH 码与 Reed-Muller 码）： - \textbf{BCH 码}（Bose-Chaudhuri-Hocquenghem，1960）：指定纠错能力 \(t\) 的循环码，设计距离 \(\delta = 2t + 1\) - \textbf{Reed-Muller 码}：基于布尔函数的代数构造，\(RM(r, m)\) 参数为 \([2^m, \sum_{i=0}^r \binom{m}{i}, 2^{m-r}]\)

\textbf{定理 257.3}（Berlekamp-Massey 算法，1968-1969）：BCH 码的解码可通过求解线性递推关系在 \(O(n^2)\) 时间内完成。这是代数解码的核心算法。

\textbf{证明}：设接收向量为 \(r(x)=c(x)+e(x)\)，错误位置多项式 \(\sigma(x)=\prod_{i\in E}(1-\alpha_i x)\)（\(E\) 为错误位置集）。由 Newton 恒等式，症候 \(S_j=r(\alpha^j)=e(\alpha^j)\) 与 \(\sigma(x)\) 的系数满足线性递推关系 \(\sigma_0 S_j+\sigma_1 S_{j-1}+\cdots+\sigma_t S_{j-t}=0\)（\(j=t+1,\ldots,2t\)）。Berlekamp-Massey 算法以 \(O(n^2)\) 时间求解最短线性递推（即 \(\sigma(x)\)）。算法维护当前递推多项式，逐位处理症候序列，在不匹配时更新递推长度和多项式。每次更新 \(O(n)\)，共 \(n\) 步，总 \(O(n^2)\)。求得 \(\sigma(x)\) 后，Chien 搜索求根确定错误位置，Forny 公式求错误值，完成解码。\(\blacksquare\)

\subsection{257.3 现代编码：Turbo 码、LDPC 与 Polar 码}\label{ux73b0ux4ee3ux7f16ux7801turbo-ux7801ldpc-ux4e0e-polar-ux7801}

\textbf{定义 257.5}（Turbo 码，Berrou-Glavieux-Thitimajshima 1993）：Turbo 码通过并行级联两个卷积码并使用迭代软判决解码，性能接近 Shannon 极限（差距 \(< 0.7\) dB）。

Turbo 码的出现在 1993 年震撼了编码理论界------它以相对简单的级联结构和迭代解码算法逼近了 Shannon 在 1948 年确立的理论极限，打破了此前近半个世纪代数编码的主导范式。在其成功的启发下，被遗忘近三十年的 Gallager 低密度校验码（LDPC）重新受到关注。

\textbf{定义 257.6}（LDPC 码 / 低密度校验码，Gallager 1963 / MacKay-Neal 1996）：LDPC 码由稀疏校验矩阵定义，使用置信传播（belief propagation）在 Tanner 图上的迭代解码。极长码长下接近 Shannon 极限。

Turbo 码和 LDPC 码代表了编码理论从确定性代数构造向概率迭代方法的范式转换。这一转换的顶点是 Arıkan 于 2009 年提出的 Polar 码------它不仅是第三个逼近 Shannon 极限的实用方案，更是首个被严格数学证明可达任意对称信道容量的构造。

\textbf{定理 257.4}（Polar 码的极性化定理 / Arıkan）：通过信道极化（channel polarization）构造的 Polar 码，在码长 \(N \to \infty\) 时可达对称信道容量。Polar 码是首个被严格证明可在所有对称二进制输入无记忆信道下达到 Shannon 信道容量的构造性编码方案。

\textbf{证明}：对 \(N=2^n\) 个独立同分布 BMS 信道 \(W\)，递归应用信道组合与分裂操作 \(W\mapsto(W^-,W^+)\)，其中 \(W^-(y_1,y_2|x_1)=\frac{1}{2}\sum_{x_2}W(y_1|x_1\oplus x_2)W(y_2|x_2)\)，\(W^+(y_1,y_2,x_1|x_2)=\frac{1}{2}W(y_1|x_1\oplus x_2)W(y_2|x_2)\)。定义 Bhattacharyya 参数 \(Z(W)=\sum_y\sqrt{W(y|0)W(y|1)}\)。由 \(Z(W^-)\le 2Z(W)-Z(W)^2\) 和 \(Z(W^+)=Z(W)^2\) 及鞅收敛定理，\(Z(W_N^{(i)})\) 对大多数 \(i\) 极化至 \(0\) 或 \(1\)：\(\lim_{N\to\infty}\frac{1}{N}\#\{i:Z(W_N^{(i)})\in(\varepsilon,1-\varepsilon)\}=0\)。且 \(\lim_{N\to\infty}\frac{1}{N}\#\{i:Z(W_N^{(i)})\le\varepsilon\}=I(W)\)（信道容量）。选择 \(Z(W_N^{(i)})\approx 0\) 的 \(K=NR\) 个信道传输信息比特，其余冻结为已知值。连续消除解码误码率 \(P_e\le\sum_{i\in\mathcal{A}}Z(W_N^{(i)})\to 0\)，速率 \(R\) 趋于 \(I(W)\)。\(\blacksquare\)

\hrulefill

\hrulefill

\hrulefill

\hrulefill

\hrulefill

\hrulefill

\section{第304章：对称密码}\label{ux7b2c304ux7ae0ux5bf9ux79f0ux5bc6ux7801}

对称密码（私钥密码）是密码学最为古老也最为基础的分支，其历史可追溯至凯撒密码和维吉尼亚密码，但其科学性由 Shannon 于 1949 年在《Communication Theory of Secrecy Systems》中正式确立。Shannon 提出的完美保密定理给出了对称密码信息论安全性的充要条件------密钥空间至少与明文空间等大且密钥均匀分布------一次性密钥包（one-time pad）是其典范实现。然而，完美保密的密钥管理成本使得实际密码学不得不转向计算安全性：分组密码（DES, AES）和流密码（RC4, ChaCha20）将短密钥经由精心设计的置换-扩散网络（SPN）或 Feistel 网络扩展为伪随机密钥流，其安全性不再基于信息论而是基于密码分析的计算困难性------差分分析、线性分析和代数分析共同构成了对称密码的两面（攻击与防御）。本章在 Shannon 完全保密定理的基础上，系统阐述分组密码（AES）、流密码、Hash 函数（Merkle-Damgård 构造）和 MAC 的数学原理与安全性证明。

\subsection{258.1 密码学基础概念}\label{ux5bc6ux7801ux5b66ux57faux7840ux6982ux5ff5}

\textbf{定义 258.1}（加密方案 / Encryption Scheme）：对称加密方案由三个算法组成： - \(\mathsf{Gen}(1^\lambda) \to k\)：密钥生成（安全参数 \(\lambda\)） - \(\mathsf{Enc}(k, m) \to c\)：加密（明文 \(m\)，密文 \(c\)） - \(\mathsf{Dec}(k, c) \to m\)：解密

\textbf{定义 258.2}（Shannon 完美保密 / Perfect Secrecy，1949）：加密方案具有完美保密性，如果对任意分布 \(M\)、任意明文 \(m\) 和任意密文 \(c\)， \[\Pr[M = m | C = c] = \Pr[M = m]\] 等价地，\(|K| \geq |M|\) 且每个密钥以等概率使用（一次性密钥包 / one-time pad）。

\textbf{定理 258.3}（Shannon完美保密性的充要条件）：对称加密方案 \((K, M, C, E, D)\) 具有完美保密性当且仅当 \(|K| \geq |M|\) 且密钥均匀随机分布。

\textbf{证明}：充分性：设密钥 \(k\) 均匀分布于 \(K\)（\(|K| \geq |M|\)），且对每个明文 \(m\)，加密映射 \(k \mapsto E(k,m)\) 是单射。则 \(\Pr[C=c \mid M=m] = 1/|K|\)（若 \(c\) 可达）或 \(0\)。对所有 \(m_1, m_2\) 和任意可达 \(c\)，\(\Pr[C=c \mid M=m_1] = \Pr[C=c \mid M=m_2]\)。由 Bayes 定理，\(\Pr[M=m \mid C=c] = \frac{\Pr[C=c \mid M=m]\Pr[M=m]}{\sum_{m'} \Pr[C=c \mid M=m']\Pr[M=m']} = \Pr[M=m]\)。必要性：若完美保密成立，则对不同明文密文分布相同，因此至少 \(|M|\) 个不同的密钥使每个 \(c\) 可达，故 \(|K| \geq |M|\)。一次一密 \(c = m \oplus k\)（\(k\) 真随机）是典范构造。\(\blacksquare\)

\textbf{定义 258.3}（不可区分性 / IND-CPA 与 IND-CCA）：现代密码学使用计算安全性（非信息论安全）： - \textbf{IND-CPA}（选择明文攻击下的不可区分性）：攻击者不能区分两个自选明文的加密 - \textbf{IND-CCA}（选择密文攻击下的不可区分性）：即使有解密谕示，攻击者也不能区分加密

\subsection{258.2 分组密码与流密码}\label{ux5206ux7ec4ux5bc6ux7801ux4e0eux6d41ux5bc6ux7801}

\textbf{定义 258.4}（分组密码 / Block Cipher）：\(E: \{0,1\}^k \times \{0,1\}^n \to \{0,1\}^n\)。DES（\(n=64, k=56\)）、AES（\(n=128, k=128/192/256\)）。

分组密码是现代对称密码体系的基石，其设计理念源于 Shannon 的置换-扩散原理。DES 在 1977 年成为首个公开的分组密码标准，但 56 位密钥在 1990 年代末已可通过穷举搜索攻破。2001 年，NIST 选定了由 Daemen 和 Rijmen 设计的 Rijndael 算法作为高级加密标准 AES，其结构反映了现代分组密码设计的最佳实践。

\textbf{定义 258.5}（AES / 高级加密标准）：AES 是取代-置换网络（SPN），由 SubBytes（S 盒）、ShiftRows、MixColumns 和 AddRoundKey 四步组成，轮数由密钥长度决定。

\textbf{定义 258.6}（流密码 / Stream Cipher）：使用伪随机生成器从种子密钥生成密钥流：\(c_i = m_i \oplus G(k)_i\)。安全需求：\(G\) 的输出在计算上与真随机不可区分。

流密码与分组密码的设计哲学形成鲜明对比：分组密码将明文切分为固定大小的块分别加密，而流密码将密钥扩展为与明文等长的伪随机密钥流，逐位异或。流密码的优势在于实现简单且适于实时通信，但其安全性的核心------伪随机生成器的不可区分性------需要严格的计算复杂性归约来保证。

\textbf{定理 258.1}（伪随机函数与分组密码的安全性）：若 \(F_k(\cdot)\) 是伪随机函数（PRF），则 CTR 模式（计数器模式）加密是 IND-CPA 安全的。

\textbf{证明}：CTR 模式加密：\(c_i=m_i\oplus F_k(\text{ctr}+i)\)。假设存在 PPT 敌手 \(\mathcal{A}\) 以优势 \(\varepsilon\) 攻破 CTR 模式的 IND-CPA 安全性。构造区分器 \(\mathcal{D}\) 区分 \(F_k\) 与真随机函数 \(R\)：\(\mathcal{D}\) 模拟 CTR 加密游戏，使用 Oracle 查询生成密钥流。若 Oracle 为 \(F_k\)，\(\mathcal{A}\) 的环境与真实 IND-CPA 游戏相同，优势 \(\varepsilon\)。若 Oracle 为 \(R\)，密钥流为真随机，此时 CTR 加密退化为一次性密钥包，\(\mathcal{A}\) 优势为 \(0\)。\(\mathcal{D}\) 的区分优势 \(\ge\varepsilon\)，与 \(F_k\) 为 PRF 矛盾。故 \(\varepsilon\) 可忽略。\(\blacksquare\)

\subsection{258.3 认证与 Hash}\label{ux8ba4ux8bc1ux4e0e-hash}

\textbf{定义 258.7}（消息认证码 / MAC）：\(\mathsf{MAC}(k, m) \to t\)（标签）。安全性要求：即使能查询任意消息的 MAC 标签（adaptive chosen-message attack），攻击者也不能为新消息伪造有效标签。

\textbf{定义 258.8}（Hash 函数安全性质）：密码学 Hash 函数 \(H: \{0,1\}^* \to \{0,1\}^n\) 应满足： - \textbf{抗原像}（preimage resistance）：给定 \(y\)，难找 \(x\) 使 \(H(x) = y\) - \textbf{抗第二原像}（second preimage resistance）：给定 \(x\)，难找 \(x' \neq x\) 使 \(H(x') = H(x)\) - \textbf{抗碰撞}（collision resistance）：难找 \(x \neq x'\) 使 \(H(x) = H(x')\)

\textbf{定理 258.2}（Merkle-Damgård 构造）：若压缩函数 \(f\) 是抗碰撞的，则由其通过 Merkle-Damgård 迭代构造的 Hash 函数也是抗碰撞的。SHA-256 基于此构造。

\textbf{证明}：设压缩函数 \(f: \{0,1\}^n \times \{0,1\}^m \to \{0,1\}^n\) 是抗碰撞的。MD 迭代：\(h_0 = IV\)，\(h_{i+1} = f(h_i, M_i)\)，\(H(M) = h_k\)。若存在 \(M \neq M'\) 使 \(H(M) = H(M')\)，沿迭代链回溯。若长度不同：最后一步 \(f(h_{k-1}, M_{k-1}) = f(h'_{k'-1}, M'_{k'-1})\) 但输入不同------即 \(f\) 的碰撞。若长度相同：在某步 \(i\) 处 \(f(h_i, M_i) = f(h'_i, M'_i)\) 且 \((h_i, M_i) \neq (h'_i, M'_i)\)------亦是 \(f\) 的碰撞。无论何种情形，\(H\) 的碰撞均归约为 \(f\) 的碰撞，与 \(f\) 的抗碰撞性矛盾。\(\blacksquare\)

\hrulefill

\hrulefill

\hrulefill

\hrulefill

\hrulefill

\hrulefill

\section{第305章：公钥密码}\label{ux7b2c305ux7ae0ux516cux94a5ux5bc6ux7801}

公钥密码（非对称密码）是密码学最革命性的发明，由 Diffie 和 Hellman 于 1976 年在《New Directions in Cryptography》中提出，解决了困扰密码学数千年的密钥分发问题------通信双方无须事先共享秘密即可建立安全的加密信道。这一突破的核心思想是将密钥分为公钥（公开，用于加密）和私钥（保密，用于解密），其安全性依赖于数论中的单向陷门函数：大整数分解（RSA，1978 年）、有限域的离散对数（Diffie-Hellman，1976 年）和椭圆曲线的离散对数（ECC，Miller/Koblitz，1985 年）。公钥密码的安全性与数论（V17 中的素数分布、Euler 定理、二次剩余）直接耦合------RSA 的安全性等价于大整数分解的困难性，椭圆曲线密码则将经典离散对数的亚指数复杂度提升为指数级（Pollard rho 的 \(O(\sqrt{n})\)），使得 256 位 ECC 密钥达到与 3072 位 RSA 密钥同等的安全水平。本章在数论和抽象代数的基础上，系统阐述 RSA、Diffie-Hellman、ElGamal、椭圆曲线密码以及同态加密的数学构造与安全性归约。

\subsection{259.1 RSA 与因子分解}\label{rsa-ux4e0eux56e0ux5b50ux5206ux89e3}

\textbf{定义 259.1}（RSA 密码体制，Rivest-Shamir-Adleman 1978）： - \(\mathsf{Gen}\)：选大素数 \(p, q\)，\(N = pq\)，\(\varphi(N) = (p-1)(q-1)\)。选 \(e\) 满足 \(\gcd(e, \varphi(N)) = 1\)。计算 \(d = e^{-1} \bmod \varphi(N)\)。公钥 \((N, e)\)，私钥 \(d\) - \(\mathsf{Enc}((N, e), m) = m^e \bmod N\) - \(\mathsf{Dec}(d, c) = c^d \bmod N\)

\textbf{定理 259.1}（RSA 的正确性）：\(c^d \bmod N = m^{ed} \bmod N = m\)。

\emph{证明}：由 \(ed \equiv 1 \pmod{\varphi(N)}\)，存在整数 \(k\) 使 \(ed = k\varphi(N) + 1\)，需证 \(m^{k\varphi(N)+1} \equiv m \pmod{pq}\)。分两种情况： 1. 若 \(\gcd(m, N) = 1\)：由 Euler 定理 \(m^{\varphi(N)} \equiv 1 \pmod N\)，故 \(m^{ed} = m^{k\varphi(N)+1} \equiv m \pmod N\)。 2. 若 \(\gcd(m, N) \neq 1\)：因 \(N = pq\) 且 \(p, q\) 为素数，\(\gcd(m, N) \neq 1\) 意味着 \(p \mid m\) 或 \(q \mid m\)（或二者）。先考虑模 \(p\)：若 \(p \mid m\)，则 \(m^{ed} \equiv 0 \equiv m \pmod p\)；若 \(p \nmid m\)，由 Fermat 小定理 \(m^{p-1} \equiv 1 \pmod p\)，故 \(m^{ed} = m \cdot (m^{p-1})^{k(q-1)} \equiv m \pmod p\)。同理对模 \(q\) 也有 \(m^{ed} \equiv m \pmod q\)。由中国剩余定理，\(m^{ed} \equiv m \pmod{pq}\)。∎

\textbf{定义 259.2}（RSA 假设与整数分解）：RSA 安全性基于\textbf{大整数分解的困难性}和 \textbf{RSA 问题}（已知 \(N, e, m^e\) 求 \(m\)）。目前最佳分解算法为数域筛法（NFS），亚指数复杂度 \(L_N(1/3, c)\)。

\textbf{定理 259.2}（Wiener 攻击 / 小解密指数的 RSA）：若 \(d < N^{1/4}/3\)，可从公钥 \((N, e)\) 在多项式时间内恢复私钥（利用连分数）。

\textbf{证明}：由 \(ed \equiv 1 \pmod{\varphi(N)}\)，存在 \(k\) 使 \(ed - k\varphi(N) = 1\)。因 \(\varphi(N) = N - p - q + 1\)，有 \(\left|\frac{e}{N} - \frac{k}{d}\right| = \frac{|ed - kN|}{Nd} < \frac{3k}{d\sqrt{N}}\)。若 \(d < N^{1/4}/3\)，则 \(k < d\) 且 \(\left|\frac{e}{N} - \frac{k}{d}\right| < \frac{1}{2d^2}\)。由连分数理论，\(\frac{k}{d}\) 必为 \(\frac{e}{N}\) 连分数展开的收敛子。计算 \(\frac{e}{N}\) 的连分数（\(O(\log N)\) 项），逐一检验收敛子满足 \(ed \equiv 1 \pmod{\varphi(N)}\) 者即恢复私钥 \(d\)。\(\blacksquare\)

\subsection{259.2 离散对数与椭圆曲线密码}\label{ux79bbux6563ux5bf9ux6570ux4e0eux692dux5706ux66f2ux7ebfux5bc6ux7801}

\textbf{定义 259.3}（Diffie-Hellman 密钥交换）：公开参数：素数 \(p\) 和生成元 \(g \in \mathbb{F}_p^*\)。方 \(A\) 选取 \(a\)，公开 \(g^a\)；方 \(B\) 选取 \(b\)，公开 \(g^b\)。共享值 \(g^{ab}\) 可由双方各自计算。安全性基于\textbf{计算 Diffie-Hellman（CDH）假设}。

\textbf{定理 259.4}（Diffie-Hellman 安全性归约）：若决策 Diffie-Hellman（DDH）假设成立------即区分 \((g^a, g^b, g^{ab})\) 与 \((g^a, g^b, g^r)\)（\(r\) 随机）在计算上不可行------则 Diffie-Hellman 密钥交换在窃听攻击下是语义安全的。DDH 蕴含 CDH，且在一般循环群模型中，CDH 的困难性等价于离散对数问题（DLP）的困难性。椭圆曲线群 \(E(\mathbb{F}_q)\) 上的 ECDLP：已知最有效攻击（Pollard rho）复杂度为 \(O(\sqrt{n})\)（\(n\) 为群阶），与有限域上的亚指数级 DLP 形成对比。这是 ECC 使用短密钥却提供与 RSA 同等安全性的根本原因。

\textbf{证明}：DDH 蕴含语义安全：窃听者观察 \((g^a,g^b)\)，需区分共享密钥 \(g^{ab}\) 与随机群元素。DDH 假设直接断言此区分在计算上不可行，故共享密钥在计算上不可区分于随机，满足语义安全。DDH 蕴含 CDH：若 CDH 可解则从 \((g^a,g^b)\) 可计算 \(g^{ab}\)，从而 DDH 平凡可破。一般群模型中 CDH 等价于 DLP：Shoup 证明任何一般群算法求解 DLP 需 \(\Omega(\sqrt{p})\) 次群运算，与 CDH 一致。Pollard rho 在椭圆曲线群上复杂度 \(O(\sqrt{n})\)（\(n\) 为群阶），而数域筛法在有限域上为亚指数 \(L_p(1/3)\)。此指数差距是 ECC 高效性的根源。\(\blacksquare\)

\textbf{定义 259.4}（椭圆曲线密码 / ECC）：以椭圆曲线群 \(E(\mathbb{F}_q)\) 代替乘法群 \(\mathbb{F}_p^*\)。椭圆曲线离散对数问题（ECDLP）的最知名攻击复杂度为指数级（\(\sqrt{|E|}\)），这允许使用比 RSA 更短的密钥达到同等安全性。

椭圆曲线密码的数学优势源于一个简单的事实：在一般椭圆曲线群上，不存在已知的亚指数时间算法求解离散对数问题，这与有限域乘法群上的数域筛法形成鲜明对比。这一差异意味着 256 位的 ECC 密钥可以提供与 3072 位 RSA 密钥同等的安全性------这种效率增益对资源受限环境（如智能卡和 IoT 设备）至关重要。在具体的密码方案层面，椭圆曲线的群结构自然地支持了多种经典构造的移植。

\textbf{定义 259.5}（ElGamal 加密与 ECDSA 签名）： - \textbf{ElGamal 加密}：公钥 \(h = g^x\)。加密：随机 \(r\)，输出 \((g^r, m \cdot h^r)\)。解密：\(m = (h^r)^{-x} \cdot m \cdot h^r\) - \textbf{ECDSA}：基于椭圆曲线群的数字签名算法，其安全性归约到 ECDLP

\textbf{定理 259.5}（ElGamal加密的IND-CPA安全性）：在DDH假设下，ElGamal加密方案在选择明文攻击下具有不可区分性（IND-CPA安全）。密文 \((g^r, m \cdot h^r)\) 中第一分量均匀随机，第二分量在DDH假设下与随机不可区分------攻击者无法区分 \(h^r = g^{xr}\) 与随机群元素。解密需要私钥 \(x\)（即离散对数），安全性完全归约到CDH/DDH。将ElGamal移植到椭圆曲线群 \(E(\mathbb{F}_q)\) 上得到EC-ElGamal，在ECDLP假设下具有同等安全性，但密钥和密文长度大幅缩短。ElGamal本身具有乘法同态性：\(E(m_1) \cdot E(m_2) = (g^{r_1+r_2}, m_1 m_2 h^{r_1+r_2}) = E(m_1 m_2)\)，这是最早的同态加密方案之一。

\textbf{证明}：设敌手 \(\mathcal{A}\) 以优势 \(\varepsilon\) 攻破 ElGamal 的 IND-CPA 安全。构造 DDH 区分器 \(\mathcal{D}\)：接收 \((g^a,g^b,g^c)\)（\(c\) 为 \(ab\) 或随机），设公钥 \(h=g^a\)。\(\mathcal{A}\) 提交 \(m_0,m_1\)，\(\mathcal{D}\) 返回挑战密文 \((g^b,m_\beta\cdot g^c)\)。若 \(c=ab\)，密文为真实 ElGamal 加密，\(\mathcal{A}\) 具有优势 \(\varepsilon\)。若 \(c\) 随机，\(g^c\) 独立于 \(g^b\)，第二分量在 \(\mathcal{A}\) 视角下为均匀随机，\(\mathcal{A}\) 优势为 \(0\)。\(\mathcal{D}\) 区分优势 \(\ge\varepsilon\)，与 DDH 假设矛盾。\(\blacksquare\)

\subsection{259.3 同态加密与安全计算}\label{ux540cux6001ux52a0ux5bc6ux4e0eux5b89ux5168ux8ba1ux7b97}

\textbf{定义 259.6}（同态加密 / Homomorphic Encryption）：加密方案 \(\mathsf{Enc}\) 是\textbf{同态的}，如果存在运算 \(\oplus\) 使 \(\mathsf{Enc}(m_1) \oplus \mathsf{Enc}(m_2) = \mathsf{Enc}(m_1 + m_2)\)（加法同态 / Paillier）或 \(\otimes\) 使 \(\mathsf{Enc}(m_1) \otimes \mathsf{Enc}(m_2) = \mathsf{Enc}(m_1 \cdot m_2)\)（乘法同态 / ElGamal）。

部分同态加密方案（如 RSA 的乘法同态和 Paillier 的加法同态）很早便为人所知，然而同时支持加法和乘法同态的完全同态加密（FHE）在三十余年间一直是密码学的''圣杯''------它意味着可以在加密数据上直接执行任意电路计算，而无需解密。2009 年，Craig Gentry 基于理想格构造了突破性的蓝图，首次实现了这一目标。

\textbf{定理 259.3}（完全同态加密 / FHE，Gentry 2009）：存在同时支持加法和乘法同态的加密方案（基于理想格），使得可以在密文上执行任意电路计算。Gentry 的蓝图：基于格的 somewhat 同态加密 + 自举（bootstrapping）实现完全同态。

\textbf{证明}：Gentry 蓝图分三步：(1) 构造 Somewhat 同态加密（SHE）：基于理想格，加密 \(c = m + 2r + b \cdot pk\)（\(r\) 为随机小噪声）。\(c_1 + c_2\) 解密为 \(m_1 + m_2\)（噪声加倍），\(c_1 \cdot c_2\) 解密为 \(m_1 m_2\)（噪声平方增长）。(2) 压缩解密电路：将解密算法表示为低深度布尔电路，使其可在 SHE 自身内部同态求值。(3) 自举（Bootstrapping）：用 SHE 对''新鲜''密文加密原私钥，同态执行解密电路获得''刷新''密文（噪声重置为初始水平）。若 SHE 可同态求值自身解密电路加一次运算，则递归自举可实现任意深度电路的同态求值------即 FHE。\(\blacksquare\)

\hrulefill

\hrulefill

\hrulefill

\hrulefill

\hrulefill

\hrulefill

\section{第306章：格密码与后量子密码}\label{ux7b2c306ux7ae0ux683cux5bc6ux7801ux4e0eux540eux91cfux5b50ux5bc6ux7801}

Peter Shor 于 1994 年提出的量子算法证明量子计算机可以在多项式时间内分解大整数和求解离散对数------这意味着一旦大规模通用量子计算机建成，目前广泛部署的 RSA 和 ECC 将立即失效。后量子密码学（PQC）正是对这场''量子威胁''的数学回应，它研究即使面对量子对手仍能保持计算安全性的密码体制。在四类主要后量子候选方案（格、编码、多变量、Hash）中，基于格的密码以独有的数学优势脱颖而出------它是目前唯一具有''最坏情况到平均情况''安全性归约的方案（Regev 2005 的量子归约定理）。格的代数结构（线性无关向量基生成的离散加法子群）提供了丰富的困难问题谱系：SVP（最短向量）、CVP（最近向量）、SIS（短整数解）和 LWE（带误差学习），后者支撑了 NIST 2022 年选定的后量子标准------Kyber（KEM）和 Dilithium（签名）。本章从格的几何定义出发，以 LLL 算法和 LWE 加密为核心支柱，揭示格密码作为后量子密码学基石的数学结构。

\subsection{260.1 格基础}\label{ux683cux57faux7840}

\textbf{定义 260.1}（格 / Lattice）：\(\mathbb{R}^n\) 中一组线性无关向量 \(\mathbf{b}_1, \ldots, \mathbf{b}_n\) 的所有整数线性组合构成格：

\[\mathcal{L} = \left\{ \sum_{i=1}^n x_i \mathbf{b}_i : x_i \in \mathbb{Z} \right\}\]

\textbf{定义 260.2}（格中的难解问题）： - \textbf{SVP}（最短向量问题）：求格中最短的非零向量的长度 \(\lambda_1(\mathcal{L})\) - \textbf{CVP}（最近向量问题）：给定目标点，求最近的格点 - \textbf{SIS}（短整数解）：给定随机矩阵 \(A\)，找短的 \(x \neq 0\) 使 \(A x = 0 \pmod q\) - \textbf{LWE}（带误差学习 / Learning With Errors，Regev 2005）：从 \((A, As + e)\) 中恢复秘密 \(s\)（\(e\) 是小误差）

\textbf{定理 260.1}（Regev 的 LWE 安全性归约，2005/2009）：LWE 问题至少与最坏情况的格问题（GapSVP 和 SIVP）一样难。即攻破 LWE 意味着可以用量子算法近似解决所有格中的最坏情况难解问题。

\textbf{证明}：Regev 归约分两步。(1) 从最坏情况格问题到离散高斯采样：若存在 GapSVP\(_{\tilde{O}(n/\alpha)}\) 的 Oracle，则可在格 \(\mathcal{L}\) 上以标准差 \(\alpha\lambda_1(\mathcal{L})\) 采样离散高斯分布 \(D_{\mathcal{L},r}\)。此步骤使用量子计算（迭代傅里叶变换和量子态制备）。(2) 从离散高斯采样到 LWE：给定 \(\mathbf{s}\)，从 \(D_{\mathcal{L},\mathbf{s},r}\) 采样得 \(\mathbf{a}\) 和 \(b=\langle\mathbf{a},\mathbf{s}\rangle+e\)（\(e\) 的分布近似于宽度 \(\alpha q\) 的离散高斯）。若 \(\alpha q\ge 2\sqrt{n}\)，则 \(e\) 的分布与 LWE 误差分布统计不可区分。因此 LWE Oracle 可解 \(\mathbf{s}\)，从而破坏最坏情况格问题。归约总损失为多项式因子 \(\tilde{O}(n/\alpha)\)。\(\blacksquare\)

\textbf{定理 260.2}（LLL 算法 / Lenstra-Lenstra-Lovász 1982）：LLL 算法在多项式时间内找到长度不超过 \(\lambda_1\) 的 \(2^{(n-1)/2}\) 倍的短向量。LLL 是格基规约的基本工具。

\textbf{证明}：设 \(\mathbf{b}_1,\ldots,\mathbf{b}_n\) 为 \(\mathcal{L}\) 的基。Gram-Schmidt 向量 \(\mathbf{b}_i^* = \mathbf{b}_i - \sum_{j<i} \mu_{ij} \mathbf{b}_j^*\)（\(\mu_{ij} = \langle\mathbf{b}_i,\mathbf{b}_j^*\rangle / \|\mathbf{b}_j^*\|^2\)）。大小约化确保 \(|\mu_{ij}| \leq 1/2\)。Lovász 条件：若 \(\|\mathbf{b}_i^*\|^2 < (\frac{3}{4} - \mu_{i,i-1}^2) \|\mathbf{b}_{i-1}^*\|^2\)，交换 \(\mathbf{b}_{i-1}\) 与 \(\mathbf{b}_i\)。算法在 \(O(n^6\log^3 B)\) 步后终止，输出基满足 \(\|\mathbf{b}_1\| \leq 2^{(n-1)/4} (\det \mathcal{L})^{1/n}\) 且 \(\|\mathbf{b}_1\| \leq 2^{(n-1)/2} \lambda_1(\mathcal{L})\)。每次交换使 \(\prod \|\mathbf{b}_i^*\|^2\) 减少至少 \(3/4\)，保证终止性。\(\blacksquare\)

\subsection{260.2 基于格的密码体制}\label{ux57faux4e8eux683cux7684ux5bc6ux7801ux4f53ux5236}

\textbf{定义 260.3}（Regev 的 LWE 加密，2005）：公钥 \(A \in \mathbb{Z}_q^{n \times m}\) 和 \(b = s^T A + e^T\)（误差 \(e\) 是短的）。加密比特 \(m \in \{0, 1\}\) 为 \((A^T r, b^T r + m\lfloor q/2 \rfloor)\)（\(r\) 小随机向量）。解密利用 \(s^T A^T r \approx b^T r\) 恢复 \(m\)。

Regev 的 LWE 加密方案是后量子密码学的基石。其安全性直接归约到最坏情况格问题的困难性（定理 260.1），这使得它成为目前已知的唯一具有''最坏情况到平均情况''归约的实用公钥加密方案。解密正确性的关键在于：\(s^T A^T r - b^T r = (s^T A^T - (s^T A + e^T)) r = -e^T r\)，由于 \(e\) 和 \(r\) 均为小范数向量，\(|e^T r| \ll q/4\)，因此可以可靠地恢复 \(m\)。LWE 加密方案的一个重要变体是\textbf{对偶 LWE 加密}（Dual Regev 加密，Gentry-Peikert-Vaikuntanathan 2008），其中公钥为 \((A, u = A s)\)，密文为 \((r^T A, r^T u + m \lfloor q/2 \rfloor)\)，其加密过程无需误差采样，且安全性同样归约到 LWE。

\textbf{定义 260.4}（NTRU / CRYSTALS-Kyber / ML-KEM）：NTRU 基于多项式环 \(R = \mathbb{Z}_q[x]/(x^n - 1)\) 上的格问题。CRYSTALS-Kyber 使用 Module-LWE 的 ML-KEM 构造。

NTRU（Hoffstein-Pipher-Silverman 1996）是最早的实用格密码方案之一。其核心思想是在环 \(R_q = \mathbb{Z}_q[x]/(x^n - 1)\) 上操作：私钥为两个短多项式 \(f, g \in R_q\)，公钥为 \(h = f^{-1} \cdot g \bmod q\)（假设 \(f\) 可逆）。加密 \(m\) 为 \(c = r \cdot h + m \bmod q\)（\(r\) 为随机短多项式），解密计算 \(f \cdot c \bmod q = f \cdot (r \cdot f^{-1} g + m) = r \cdot g + f \cdot m \bmod q\)，由于 \(r, g, f, m\) 均为短多项式，可在 \(\mathbb{Z}\) 上精确恢复。CRYSTALS-Kyber（NIST PQC 标准，2022）将 NTRU 的思想与 LWE 结合，使用 Module-LWE 在环 \(R_q = \mathbb{Z}_q[x]/(x^{256} + 1)\) 上构造，其安全性基于 Module-LWE 和 Module-SIS 的联合困难性，在通信效率和计算速度之间取得了极佳平衡。

\textbf{定义 260.5}（CRYSTALS-Dilithium / ML-DSA）：基于 Module-SIS 和 Module-LWE 的数字签名方案。

CRYSTALS-Dilithium 是 NIST 选定的后量子数字签名标准。其核心构造是 Fiat-Shamir with Aborts 范式（Lyubashevsky 2009）：签名者利用短秘密向量 \(s\) 生成承诺 \(w = A y\)（\(y\) 为随机短向量），通过 Hash 产生挑战 \(c\)，计算响应 \(z = y + c s\)。若 \(z\) 的范数过大则''中止''（abort）并重新尝试，以拒绝采样的方式确保 \(z\) 的分布不泄露 \(s\) 的信息。验证者检查 \(A z - c t = w\)（\(t = A s\) 为公钥）且 \(z\) 满足范数约束。这一构造在随机谕示模型下为 EUF-CMA 安全，其安全性归约到 Module-SIS 和 Module-LWE。

\subsection{260.3 其他后量子密码方法}\label{ux5176ux4ed6ux540eux91cfux5b50ux5bc6ux7801ux65b9ux6cd5}

格密码虽然占据了 NIST 后量子标准化进程的核心位置，但它并非后量子密码的唯一候选方向。基于编码、多变量和 Hash 的密码方案分别从不同的计算困难问题出发，提供了格密码之外的多样性安全选择。

\textbf{定义 260.6}（基于 Hash 的签名 / SPHINCS+）：Merkle 树一次签名（Lamport 签名、Winternitz 签名）与少数次签名（HORS）结合的免状态 Hash 签名方案。

\textbf{定义 260.7}（基于编码的密码 / Classic McEliece）：McEliece 密码体制（1978）基于一般线性码的解码困难性（NP 完全）。使用 Goppa 码（可有效解码）作为私钥，加扰后的公开码为公钥。

\subsection{后量子密码学的数学基础与前沿}\label{ux540eux91cfux5b50ux5bc6ux7801ux5b66ux7684ux6570ux5b66ux57faux7840ux4e0eux524dux6cbf}

后量子密码学（Post-Quantum Cryptography, PQC）是 21 世纪密码学最重要的研究方向，其动力来自于量子计算机对传统公钥密码体制的威胁------Shor 算法（1994）在多项式时间内可分解大整数和求解离散对数，从而彻底攻破 RSA、DSA 和 ECDSA。PQC 的数学基础建立在四类被认为对量子攻击难解的数学问题上：\textbf{格问题}（LWE, SIS, NTRU）、\textbf{编码问题}（一般线性码的解码）、\textbf{多变量问题}（多变量二次方程组的求解，MQ）和\textbf{Hash 问题}（Hash 函数的抗原像性）。NIST 的后量子密码标准化过程（2016-2024）最终选定了基于格的算法作为主要标准。\textbf{Kyber}（ML-KEM）是 NIST 选定的密钥封装机制（KEM）标准，其安全性基于 Module-LWE 问题：给定 \(A \in R_q^{k \times k}\) 和 \(t = A s + e\)（\(s, e\) 为短向量），计算 \(s\) 是困难的。Kyber 的核心创新在于使用\textbf{数论变换}（NTT）在 \(R_q = \mathbb{Z}_q[x]/(x^{256}+1)\) 上高效实现多项式乘法，其 IND-CCA2 安全性通过 Fujisaki-Okamoto 变换（FO 变换）从 OW-CPA 提升而来。\textbf{Dilithium}（ML-DSA）是 NIST 选定的数字签名标准，基于 Module-SIS 和 Module-LWE 的联合困难性，采用 Fiat-Shamir with Aborts 范式实现 EUF-CMA 安全。\textbf{同态加密}（Homomorphic Encryption, Gentry 2009）是格密码学最令人瞩目的成就之一：全同态加密（FHE）允许在加密数据上直接进行任意计算，其安全性基于 Ring-LWE 假设，通过噪声管理（自举, bootstrapping）技术实现。FHE 在隐私保护机器学习、安全多方计算和云计算中具有革命性应用前景。\textbf{格密码的量子安全性}：Regev（2005）的量子归约定理建立了 LWE 问题与最坏情况格问题之间的量子归约，但经典归约（Peikert 2009, Brakerski et al.~2013）的进展使得 LWE 的安全性能在经典假设下确立。\textbf{全同态签名}、\textbf{基于属性的加密}（ABE）和\textbf{函数加密}（FE）这些高级密码学原语在格密码框架下得以实现，拓展了密码学的能力边界。

\hrulefill

\hrulefill

\section{第307章：密码协议}\label{ux7b2c307ux7ae0ux5bc6ux7801ux534fux8bae}

密码协议将公钥密码和对称密码的原子化原语组合为达成复杂安全目标的多方交互过程------签名提供不可否认性、零知识证明在证明所述为真的同时不泄露证据本身、安全多方计算使多方共同计算一个函数而各自输入保持私密。密码协议的理论基石是 Goldwasser-Micali-Rackoff（1985 年）和 Goldreich-Micali-Wigderson（1987 年）奠定的模拟范式------安全性的精确定义通过''现实世界中的攻击者能做的任何事都可以在理想世界中（由模拟器）完成''来刻画，这一思想将协议安全性从直觉提升为严格的可证安全数学框架。密码协议的应用直接推动了区块链技术的诞生------零知识证明（zk-SNARK 和 zk-STARK）为隐私交易和可验证计算提供数学基础，而 BLS 签名和秘密共享构成了现代去中心化系统的加密经济基石。本章以数字签名（RSA-FDH、Schnorr）、零知识证明（\(\Sigma\)-协议、Fiat-Shamir 变换）和安全多方计算（Yao 乱码电路、Shamir 秘密共享）为核心，系统建立密码协议的数学理论。

\subsection{261.1 数字签名}\label{ux6570ux5b57ux7b7eux540d}

\textbf{定义 261.1}（数字签名方案）：\(\mathsf{Gen}(1^\lambda) \to (pk, sk)\)，\(\mathsf{Sign}(sk, m) \to \sigma\)，\(\mathsf{Verify}(pk, m, \sigma) \to \{0, 1\}\)。安全性（EUF-CMA / 存在性不可伪造性）：即使攻击者可以自适应地查询任意消息的签名，也不能为新消息伪造有效签名。

数字签名是公钥密码学中与公钥加密并列的核心原语，其功能区别于加密：签名提供\textbf{认证性}（authentication）和\textbf{不可否认性}（non-repudiation），而非机密性。安全性定义 EUF-CMA（Existential Unforgeability under Chosen Message Attack）是数字签名的标准安全模型：敌手 \(\mathcal{A}\) 可以自适应地选择消息 \(m_1, \ldots, m_q\) 并获得对应签名 \(\sigma_1, \ldots, \sigma_q\)，随后 \(\mathcal{A}\) 输出一个消息-签名对 \((m^*, \sigma^*)\)，其中 \(m^*\) 不在已查询消息之列。若 \(\mathcal{A}\) 以不可忽略的概率成功，则方案被攻破。这一模型捕捉了现实中最强大的攻击能力：攻击者可以诱骗签名者签署任意消息。对于需要更强安全性的场景，\textbf{强不可伪造性}（SUF-CMA）要求攻击者即使对于已查询过的消息 \(m_i\)，也无法产生新的有效签名 \(\sigma' \neq \sigma_i\)。常见的数字签名构造范式包括：RSA-FDH（基于 RSA 假设和随机谕示模型）、Schnorr 签名（基于离散对数）、ECDSA（基于椭圆曲线离散对数）、BLS 签名（基于双线性配对，具有短签名和聚合签名等优势性质），以及后量子签名方案如 Dilithium（基于格）和 SPHINCS+（基于 Hash）。

\textbf{定理 261.1}（RSA-FDH 签名的安全性 / Full Domain Hash）：在随机谕示模型下，RSA-FDH 签名是 EUF-CMA 安全的当且仅当 RSA 问题是难的（Bellare-Rogaway 1993/1996）。

\textbf{证明}：（充分性）设敌手 \(\mathcal{A}\) 以 \(q_H\) 次 Hash 查询和 \(q_S\) 次签名查询伪造 RSA-FDH 签名。构造 RSA 求逆算法 \(\mathcal{B}\)：给定 \((N,e,y)\)，求 \(y^{1/e}\bmod N\)。\(\mathcal{B}\) 随机选 \(j\in[1,q_H]\)，对第 \(j\) 次 Hash 查询回答 \(y\)，对其他查询 \(m_i\) 选随机 \(r_i\) 回答 \(r_i^e\bmod N\)（可对其签名 \(r_i\)）。若 \(\mathcal{A}\) 对 \(m_j\) 输出伪造签名 \(\sigma\) 且 \(\sigma^e\equiv y\pmod N\)，则 \(\sigma\) 即为 \(y\) 的 RSA 逆。成功概率 \(\ge\varepsilon/q_H\)（\(\varepsilon\) 为 \(\mathcal{A}\) 优势）。\(\blacksquare\)

\subsection{261.2 零知识证明}\label{ux96f6ux77e5ux8bc6ux8bc1ux660e}

\textbf{定义 261.2}（交互证明系统 / Interactive Proof）：证明者 \(P\) 和验证者 \(V\) 之间的协议，使得 \(P\) 可以令 \(V\) 确信陈述 \(x \in L\) 为真而不泄露其他信息。满足： - \textbf{完备性}（completeness）：若 \(x \in L\)，诚实 \(P\) 总是成功说服 \(V\) - \textbf{可靠性}（soundness）：若 \(x \notin L\)，任何（甚至恶意的）\(P^*\) 说服 \(V\) 的概率可忽略 - \textbf{零知识性}（zero-knowledge）：\(V\) 在交互中学到的不超过 \(x \in L\) 这个事实本身（存在模拟器可生成与真实交互不可区分的 view）

\textbf{定义 261.3}（\(\Sigma\)-协议 / 诚实验证者零知识）：三步协议（承诺 \(a\)、挑战 \(c\)、响应 \(z\)），满足特殊的可靠性（从两个接受相同 \(a\) 但不同 \(c\) 的 transcripts 可提取证据）。\(\Sigma\)-协议可转化为非交互零知识证明（Fiat-Shamir 变换）。

\textbf{定理 261.2}（Schnorr 的识别/签名协议，1989）：基于离散对数的 \(\Sigma\)-协议。\(P\) 知道 \(x\) 使 \(h = g^x\)。承诺：\(a = g^r\)（随机 \(r\)）。挑战：\(c\)。响应：\(z = r + cx\)。验证：\(g^z = a \cdot h^c\)。Fiat-Shamir 变换：\(c = H(a, m)\) 将交互证明转为签名。

\textbf{证明}：完备性：\(g^z=g^{r+cx}=g^r\cdot(g^x)^c=a\cdot h^c\)，验证通过。特殊可靠性：若诚实验证者产生两个接受副本 \((a,c_1,z_1)\) 和 \((a,c_2,z_2)\)（\(c_1\neq c_2\)），则 \(g^{z_1}=a\cdot h^{c_1}\) 且 \(g^{z_2}=a\cdot h^{c_2}\)。相除得 \(g^{z_1-z_2}=h^{c_1-c_2}\)，故 \(h=g^{(z_1-z_2)/(c_1-c_2)}\)，提取证据 \(x=(z_1-z_2)(c_1-c_2)^{-1}\bmod q\)。诚实验证者零知识：模拟器随机选 \(c,z\)，计算 \(a=g^z\cdot h^{-c}\)，分布与真实副本相同。Fiat-Shamir 变换用随机谕示 \(H\) 替代交互挑战，在随机谕示模型下保持 EU-CMA 安全。\(\blacksquare\)

\subsection{261.3 安全多方计算}\label{ux5b89ux5168ux591aux65b9ux8ba1ux7b97}

\textbf{定义 261.4}（安全多方计算 / Secure MPC）：\(n\) 个参与方各自有输入 \(x_i\)，联合计算 \(f(x_1, \ldots, x_n)\) 而不暴露各自的输入（除计算结果的隐含信息外）。

\textbf{定义 261.5}（秘密共享与阈值密码）： - \textbf{Shamir 秘密共享}（1979）：\(t\) 个点确定 \(t-1\) 次多项式。将秘密 \(s\) 编码为 \(f(0)=s\)，分发给 \(n\) 方 \(f(i)\)。任意 \(t\) 方可通过 Lagrange 插值恢复 \(s\)，\(t-1\) 方无法获得任何信息 - \textbf{阈值密码}：解密/签名需要 \(\ge t\) 个参与方合作

\textbf{定理 261.3}（Yao 的乱码电路 / Garbled Circuit，1986）：任何多项式时间可计算函数都可以被两方安全地计算（在半诚实模型下）。这是安全两方计算的基础理论。

\textbf{证明}：设 \(f:\{0,1\}^n\times\{0,1\}^n\to\{0,1\}^m\) 由布尔电路 \(C\) 计算。乱码方（\(P_1\)）对 \(C\) 中每条线 \(w\) 生成两个随机密钥 \(k_w^0,k_w^1\)（对应值 \(0,1\)）。对每个门 \(g\)（输入线 \(\alpha,\beta\)，输出线 \(\gamma\)），构造乱码真值表：用 \(k_\alpha^{a},k_\beta^{b}\) 加密 \(k_\gamma^{g(a,b)}\)（双重加密）。\(P_1\) 发送乱码电路和其输入对应的密钥给 \(P_2\)。\(P_2\) 通过不经意传输（OT）获取其输入对应的密钥而不泄露输入值。\(P_2\) 逐门解密获得输出密钥，通过输出映射表恢复 \(f(x,y)\)。半诚实安全：\(P_2\) 仅看到与 \(y\) 一致的一组密钥，\(P_1\) 输出为空（或对称输出）。模拟器利用 OT 的模拟性和加密的语义安全性构造。\(\blacksquare\)

\subsection{密码协议的现代发展与验证}\label{ux5bc6ux7801ux534fux8baeux7684ux73b0ux4ee3ux53d1ux5c55ux4e0eux9a8cux8bc1}

密码协议理论在 21 世纪经历了从理论到实践的飞跃，安全多方计算（MPC）和零知识证明已从学术概念发展为实际部署的工业级技术。\textbf{通用可组合性}（Universal Composability, UC, Canetti 2001）是密码协议安全性的最强定义：在 UC 框架下，协议在任意环境（包括与其他协议并发执行）中保持安全性，其安全性证明通过与环境交互的模拟器（simulator）构造实现。UC 安全性解决了协议组合时的安全问题------这在互联网协议栈和区块链智能合约中至关重要。\textbf{简洁非交互零知识证明}（zk-SNARK, Groth 2016）是零知识证明的重大突破：证明大小仅为几百字节，验证时间与计算量无关（常数时间），使得区块链中的隐私交易（如 Zcash）和可验证计算成为可能。zk-SNARK 的数学基础包括\textbf{二次算术程序}（QAP）和\textbf{双线性配对}（pairing-based cryptography），其安全性依赖于 \(q\)-型假设（如 \(q\)-PKE 假设）。\textbf{zk-STARK}（Ben-Sasson et al.~2018）是另一类零知识证明，不依赖双线性配对，而是使用 Hash 函数（如 Keccak/SHA-3）和交互式谕示证明（IOP），具有抗量子攻击和透明设置（无需可信设置）的优势。\textbf{多方计算}（MPC）的最新进展包括\textbf{不经意传输扩展}（OT extension, Ishai et al.~2003）------将少量基础 OT 扩展为大量 OT，极大降低了通信开销------和\textbf{秘密共享变换}（GMW, BGW, SPDZ 协议族），实现了主动安全（抵御主动攻击者）和诚实多数假设下的高效计算。\textbf{同态签名}和\textbf{同态 MAC} 允许在认证数据上直接计算，输出结果自动携带有效认证，在网络编码和云计算审计中具有重要应用。密码协议的理论基础------\textbf{模拟范式}（simulation paradigm）和\textbf{归约证明}（reduction proof）------为这些构造提供了严格的数学安全保障，使得密码学从''艺术''发展为''科学''。

\hrulefill

\hrulefill

\subsection{245.A 量子密码学与信息论安全}\label{a-ux91cfux5b50ux5bc6ux7801ux5b66ux4e0eux4fe1ux606fux8bbaux5b89ux5168}

量子密码学与信息论安全是现代密码学的前沿方向，它们超越了传统密码学对计算困难假设的依赖，从物理定律或信息论原理出发提供无条件的安全性保证。本章聚焦于两个核心主题：量子密钥分发（QKD）利用量子力学的不可克隆定理和测量扰动原理，在通信双方之间建立信息论安全的共享密钥------BB84 协议（Bennett-Brassard, 1984）是最经典的实现，其无条件安全性由 Shor-Preskill 证明（2000）归约为 CSS 量子纠错码；后量子密码学（特别是格密码学）则为应对量子计算机对 RSA 和椭圆曲线密码的威胁提供了候选方案------基于格上最短向量问题（SVP）和学习带误差问题（LWE）的密码方案是目前最受关注的后量子密码候选者，其安全性可归约为格上最坏情况困难问题的平均情况难度（Regev, 2005）。

\subsection{262.1 量子密钥分发的数学基础}\label{ux91cfux5b50ux5bc6ux94a5ux5206ux53d1ux7684ux6570ux5b66ux57faux7840}

量子密钥分发（QKD）利用量子态的不可克隆性质在两方之间建立信息论安全的共享密钥。

\textbf{定义 262.1}（BB84 协议 / Bennett-Brassard）：方 \(A\) 在两组建共轭基底中随机选择------计算基底 \(\{|0\rangle, |1\rangle\}\) 和 Hadamard 基底 \(\{|+\rangle = \frac{|0\rangle+|1\rangle}{\sqrt{2}}, |-\rangle = \frac{|0\rangle-|1\rangle}{\sqrt{2}}\}\)------制备量子态并通过量子信道发送给方 \(B\)。方 \(B\) 随机选择测量基底。双方公开比对基底选择后，保留基底一致的部分作为原始密钥（sifted key）。

\textbf{协议步骤}： 1. 方 \(A\) 随机生成长度 \(\gg n\) 的随机比特串和随机基底串，制备对应量子态并发送 2. 方 \(B\) 随机选择测量基底，测量每个接收到的量子态并记录结果 3. 双方通过经典认证信道公开比对基底选择，丢弃基底不一致的比特 4. 随机抽样部分比特公开比对，估计量子比特错误率（QBER）以检测窃听 5. 若 QBER 低于阈值，执行信息调和（纠错）和隐私放大（通过 universal hashing）提取最终安全密钥

\textbf{定理 262.1}（BB84 无条件安全性 / Shor-Preskill）：BB84 协议在任意攻击（包括联合攻击和相干攻击）下的安全性可归约为 CSS（Calderbank-Shor-Steane）量子纠错码的不可区分性。窃听者获得的信息量以 Holevo 界约束：

\[I(A;B) \leq S(\rho_A) = -\operatorname{Tr}(\rho_A \log \rho_A)\]

其中 \(S(\rho_A)\) 是方 \(A\) 子系统约化密度矩阵的 von Neumann 熵。当观测到的 QBER 低于约 \(11\%\) 时，可提取正的密钥率。

\textbf{证明}：Shor-Preskill 证明将 BB84 安全性归约为 CSS 码的纠错性质。将 BB84 等效为纠缠提纯协议（EPP）：Alice 制备 EPR 对，将一半发送给 Bob（等价于 BB84 制备测量方案）。窃听者 Eve 的任意攻击对应于在传输中引入幺正错误。Alice 和 Bob 随机选取部分比特测量 \(X\) 或 \(Z\)（CSS 码的双重结构），利用 CSS 码纠错和隐私放大。若 QBER \(<11\%\)，CSS 码可同时纠正比特翻转和相位翻转错误。Holevo 界 \(I(A:E)\le S(\rho_A)\) 约束 Eve 的信息量。最终密钥率 \(R=1-h(QBER_X)-h(QBER_Z)\)（\(h\) 为二元熵函数），当 QBER \(<11\%\) 时 \(R>0\)。\(\blacksquare\)

\textbf{定义 262.2}（Holevo 界）：从量子态中可提取的关于经典比特的信息量不超过对应子系统的 von Neumann 熵。这是量子信息论中的核心上界。

\subsection{262.2 格密码与后量子安全}\label{ux683cux5bc6ux7801ux4e0eux540eux91cfux5b50ux5b89ux5168}

\textbf{定义 262.3}（格与格问题）：\(n\) 维\textbf{格} \(\Lambda = \{\sum_{i=1}^n z_i \mathbf{b}_i : z_i \in \mathbb{Z}\}\)（\(\mathbf{b}_1, \ldots, \mathbf{b}_n\) 为格基）。核心困难问题： - \textbf{SVP}（最短向量问题）：求格中长度最短的非零向量 - \textbf{CVP}（最近向量问题）：给定目标点，求最近的格点 - \textbf{SIVP}（最短独立向量问题）：求 \(n\) 个线性无关的短格向量

\textbf{定义 262.4}（LWE 问题 / Learning With Errors, Regev, 2005）：给定素数模数 \(q\)、维度 \(n\)、秘密向量 \(\mathbf{s} \in \mathbb{Z}_q^n\) 和误差分布 \(\chi\)，LWE 分布为

\[(\mathbf{a}, b = \langle \mathbf{a}, \mathbf{s} \rangle + e \bmod q), \quad \mathbf{a} \leftarrow \mathbb{Z}_q^n \text{ 均匀}, \quad e \leftarrow \chi\]

搜索型 LWE 问题：给定来自该分布的多项式数量样本，恢复秘密 \(\mathbf{s}\)。判定型 LWE 问题：区分 LWE 分布与均匀分布。

\textbf{定理 262.2}（Regev 量子归约，2005）：若存在高效算法以不可忽略的概率求解 LWE 问题，则存在量子算法求解最坏情况下的格上近似最短向量问题（GapSVP）和最短独立向量问题（SIVP）。这建立了 LWE 的密码学安全性------假定格上最坏情况问题对量子计算机难解，则 LWE 在平均情况下也是难解的。

\textbf{证明}：Regev 量子归约的核心是离散高斯采样算法。给定格 \(\mathcal{L}\) 和参数 \(r\approx\sqrt{2n}\cdot\eta_\varepsilon(\mathcal{L}^*)\)（\(\eta_\varepsilon\) 为光滑参数）。步骤：(1) 用量子计算机制备 \(\sum_{\mathbf{x}\in\mathcal{L}}\rho_r(\mathbf{x})|\mathbf{x}\rangle\) 的叠加态（宽度 \(r\) 的离散高斯权重 \(\rho_r\)）。(2) 应用量子傅里叶变换，得到对偶格 \(\mathcal{L}^*\) 上的叠加态。(3) 测量得到 \(D_{\mathcal{L}^*,1/r}\) 的一个样本。此样本转换为 LWE 样本 \((\mathbf{a},b=\langle\mathbf{a},\mathbf{s}\rangle+e)\)，其中 \(e\) 近似离散高斯分布 \(D_{\mathbb{Z},\alpha q}\)。若 LWE 可解，则从采样中恢复 \(\mathbf{s}\)。此过程以 \(\tilde{O}(\sqrt{n}/\alpha)\) 近似因子将平均情况 LWE 归约到最坏情况 GapSVP 和 SIVP。\(\blacksquare\)

\textbf{后量子密码标准化}：NIST 的后量子密码标准化过程选定了基于格的方案： - \textbf{KYBER}（ML-KEM）：基于 Module-LWE 的密钥封装机制 - \textbf{Dilithium}（ML-DSA）：基于 Module-LWE 和 Module-SIS 的数字签名方案

这些方案的数学基础是格上困难问题的（量子）归约安全性。

\subsection{262.3 信息论安全}\label{ux4fe1ux606fux8bbaux5b89ux5168}

\textbf{定义 262.5}（信息论安全 / 无条件安全）：加密方案是\textbf{信息论安全}的，如果其安全性不依赖于任何计算假设------即使敌手拥有无限计算能力，密文也不泄露任何关于明文的信息（除密文长度外）。形式上，\(H(M \mid C) = H(M)\)（完全保密，Shannon, 1949）。

\textbf{定理 262.3}（Shannon 完全保密定理）：对 \(n\) 比特消息，一次一密（One-Time Pad）\(c = m \oplus k\)（\(k\) 为真随机且仅用一次的 \(n\) 比特密钥）实现完全保密。反之，任何完全保密的加密方案必须满足 \(H(K) \geq H(M)\)（密钥熵不小于消息熵）。

\textbf{证明}：充分性（一次一密）：设 \(k\) 均匀分布于 \(\{0,1\}^n\)，\(c=m\oplus k\)。\(\Pr[C=c|M=m]=\Pr[K=m\oplus c]=2^{-n}\)。对任意 \(m_1,m_2\) 和 \(c\)，\(\Pr[C=c|M=m_1]=\Pr[C=c|M=m_2]=2^{-n}\)。由 Bayes 定理，\(\Pr[M=m|C=c]=\Pr[M=m]\)，完美保密成立。必要性：由完美保密定义，\(H(M|C)=H(M)\)。\(H(K)\ge H(K|C)=H(M,K|C)-H(M|K,C)=H(M|K)+H(K)-H(C)\ge H(M)+H(K)-H(C)\)。又 \(H(C)\le H(M,K)=H(M)+H(K)\)，故 \(H(K)\ge H(M)\)。\(\blacksquare\)

\textbf{定义 262.6}（提取器与隐私放大 / Extractor and Privacy Amplification）：\textbf{随机提取器} \(\operatorname{Ext}: \{0,1\}^n \times \{0,1\}^d \to \{0,1\}^m\) 将任意具有最小熵 \(k\) 的弱随机源 \(X\) 与短的均匀种子 \(S\) 结合，输出统计上接近均匀的 \(m\) 比特。在 QKD 的隐私放大步骤中，使用 universal hashing 族（如 Toeplitz 矩阵）从部分安全的原始密钥中提取近乎完全安全的最终密钥。

\hrulefill

\emph{本卷在数论（V17）、抽象代数（V8、V13）和计算理论（V31）的基础上，系统建立了编码理论、对称密码、公钥密码、格密码与后量子密码以及密码协议的数学理论。共 5 章。}

\emph{本卷完成了从初等数论到现代数论的跨越：解析数论以 \(\zeta(s)\) 和素数定理为核心，展示了分析工具在数论中的威力；模形式理论揭示了全纯函数与数论之间惊人的联系（Ramanujan \(\tau\) 函数、Hecke 算子）；椭圆曲线以 Mordell-Weil 定理和 BSD 猜想为中心，连接了代数几何、模形式和 \(L\)-函数；类域论（局部和整体）以互反律为核心，完成了 Abel 扩张的分类；补充部分将密码学的数论基础（欧拉定理→RSA、离散对数→DH、椭圆曲线群→ECDSA、零知识证明→交互证明系统）纳入本卷框架。为后续 Langlands 纲领（如有需要）的学习奠定了基础。}

\emph{卷十六：数论 II终。}

\emph{卷十六：数论 II终。}

\emph{卷十六：数论 II终。}

\emph{卷十九：数论终。}

\newpage

\chapter{08-代数几何与数论/02-数论/07-初等数论.md}\label{ux4ee3ux6570ux51e0ux4f55ux4e0eux6570ux8bba02-ux6570ux8bba07-ux521dux7b49ux6570ux8bba.md}

\chapter{初等数论}\label{ux521dux7b49ux6570ux8bba}

\hrulefill

\section{第308章：整除与同余理论}\label{ux7b2c308ux7ae0ux6574ux9664ux4e0eux540cux4f59ux7406ux8bba}

初等数论（Elementary Number Theory）研究整数的基本性质，主要包括整除性、素数分布、同余关系和Diophantine方程。尽管名为''初等''，其结论之深刻与证明之精巧往往令人叹为观止。本章系统阐述整除理论、算术基本定理和同余理论。

\subsection{246.1 整除与带余除法}\label{ux6574ux9664ux4e0eux5e26ux4f59ux9664ux6cd5}

\textbf{定义 246.1.1（整除）} 设 \(a, b \in \mathbb{Z}\)。若存在 \(c \in \mathbb{Z}\) 使得 \(a = bc\)，则称 \(b\) 整除 \(a\)，记作 \(b \mid a\)；否则记作 \(b \nmid a\)。当 \(b \mid a\) 时，称 \(b\) 是 \(a\) 的因数（约数），\(a\) 是 \(b\) 的倍数。

\textbf{定理 246.1.2（带余除法）} 设 \(a \in \mathbb{Z}\)，\(b \in \mathbb{Z}^+\)（正整数）。则存在唯一的整数 \(q\) 和 \(r\)（\(0 \leq r < b\)）使得

\[a = bq + r.\]

\(q\) 称为商（quotient），\(r\) 称为余数（remainder）。

\textbf{证明}：存在性：考虑集合 \(S = \{a - bk : k \in \mathbb{Z}, a - bk \geq 0\}\)。\(S\) 是非负整数集合的非空子集（取 \(k\) 足够小，如 \(k \leq a/b\)），故 \(S\) 有最小元 \(r\)。设 \(r = a - bq\)。若 \(r \geq b\)，则 \(0 \leq r - b = a - b(q+1) \in S\) 且 \(r - b < r\)，与 \(r\) 的最小性矛盾。故 \(0 \leq r < b\)。唯一性：若 \(a = bq_1 + r_1 = bq_2 + r_2\)，则 \(b(q_1 - q_2) = r_2 - r_1\)，故 \(b \mid (r_2 - r_1)\)。但 \(|r_2 - r_1| < b\)，这只有在 \(r_1 = r_2\) 时可能，此时 \(q_1 = q_2\)。\(\blacksquare\)

\textbf{定义 246.1.3（最大公因数）} 不全为零的整数 \(a, b\) 的最大公因数（greatest common divisor）是同时整除 \(a\) 和 \(b\) 的最大正整数，记作 \(\gcd(a, b)\)。若 \(\gcd(a, b) = 1\)，则称 \(a\) 与 \(b\) 互素（coprime）。

\textbf{定理 246.1.4（Euclid算法）} 设 \(a \geq b > 0\)。反复应用带余除法：

\[a = b q_1 + r_1, \quad 0 \leq r_1 < b,\] \[b = r_1 q_2 + r_2, \quad 0 \leq r_2 < r_1,\] \[r_1 = r_2 q_3 + r_3, \quad 0 \leq r_3 < r_2,\] \[\vdots\] \[r_{k-2} = r_{k-1} q_k + r_k, \quad 0 \leq r_k < r_{k-1},\] \[r_{k-1} = r_k q_{k+1} + 0.\]

则 \(\gcd(a, b) = r_k\)（最后一个非零余数）。

\textbf{证明}：注意到 \(\gcd(a, b) = \gcd(b, r_1) = \gcd(r_1, r_2) = \cdots = \gcd(r_{k-1}, r_k) = \gcd(r_k, 0) = r_k\)。每一步的相等由 \(\gcd(x, y) = \gcd(y, x - qy)\) 保证（公因数集合不变）。\(\blacksquare\)

\textbf{定理 246.1.5（扩展Euclid算法------Bezout恒等式）} 设 \(a, b \in \mathbb{Z}\) 不全为零。则存在整数 \(x, y\) 使得

\[\gcd(a, b) = ax + by.\]

\textbf{完整证明}：考虑集合 \(S = \{au + bv : u, v \in \mathbb{Z}, au + bv > 0\}\)。\(S\) 非空（取 \(u = a, v = b\) 或适当符号调整）。设 \(d = \min S = ax_0 + by_0\)。我们断言 \(d = \gcd(a, b)\)。

首先，\(d \mid a\)：由带余除法，\(a = dq + r\)，\(0 \leq r < d\)。则 \(r = a - dq = a - (ax_0 + by_0)q = a(1 - x_0 q) + b(-y_0 q)\)。若 \(r > 0\)，则 \(r \in S\) 且 \(r < d\)，与 \(d\) 的最小性矛盾。故 \(r = 0\)，即 \(d \mid a\)。同理 \(d \mid b\)。故 \(d\) 是 \(a, b\) 的公因数。

其次，若 \(c \mid a\) 且 \(c \mid b\)，则 \(c \mid (ax_0 + by_0) = d\)。故 \(d = \gcd(a, b)\)。\(\blacksquare\)

扩展Euclid算法可在 \(O(\log \min(|a|, |b|))\) 步内求出满足Bezout恒等式的系数。

\subsection{246.2 算术基本定理}\label{ux7b97ux672fux57faux672cux5b9aux7406-1}

\textbf{定义 246.2.1（素数）} 大于 \(1\) 的整数 \(p\) 称为素数（prime），若它的正因数仅有 \(1\) 和 \(p\)。否则称为合数（composite）。\(1\) 既非素数也非合数。

\textbf{定理 246.2.2（Euclid引理）} 若 \(p\) 为素数且 \(p \mid ab\)，则 \(p \mid a\) 或 \(p \mid b\)。

\textbf{证明}：若 \(p \nmid a\)，则 \(\gcd(p, a) = 1\)。由Bezout恒等式，存在 \(x, y\) 使得 \(1 = px + ay\)。两边乘 \(b\) 得 \(b = pbx + aby\)。因 \(p \mid p b x\) 且 \(p \mid a b y\)，故 \(p \mid b\)。\(\blacksquare\)

\textbf{定理 246.2.3（算术基本定理------唯一分解定理）} 每个大于 \(1\) 的整数 \(n\) 都可以唯一地表示为素数的乘积（不计素数的顺序）：

\[n = p_1^{e_1} p_2^{e_2} \cdots p_k^{e_k},\]

其中 \(p_1 < p_2 < \cdots < p_k\) 为不同的素数，\(e_i \geq 1\)。

\textbf{证明}：存在性由强归纳法证明：若 \(n\) 为素数，分解即为其自身。若 \(n\) 为合数，则 \(n = ab\)（\(1 < a, b < n\)）。由归纳假设，\(a\) 和 \(b\) 均可分解，故 \(n\) 也可分解。

唯一性：设 \(n = p_1 p_2 \cdots p_r = q_1 q_2 \cdots q_s\)（\(p_i, q_j\) 为素数）。由Euclid引理，\(p_1 \mid q_1 q_2 \cdots q_s\)，故 \(p_1 \mid q_j\) 对某 \(j\) 成立，从而 \(p_1 = q_j\)（因双方均为素数）。消去 \(p_1\) 和 \(q_j\) 后对剩下的乘积重复论证。由归纳法最终得 \(r = s\) 且 \(\{p_i\} = \{q_j\}\)（作为多重集）。\(\blacksquare\)

\textbf{推论 246.2.4} 素数的集合是无限的。

\textbf{证明}（Euclid经典证明）：假设仅有有限个素数 \(p_1, \ldots, p_k\)。令 \(N = p_1 p_2 \cdots p_k + 1\)。\(N > 1\)，故 \(N\) 有素因子 \(p\)。但 \(p\) 不可能为任何 \(p_i\)，否则 \(p \mid 1\)，矛盾。故存在第 \(k + 1\) 个素数。\(\blacksquare\)

\subsection{246.3 同余与同余类}\label{ux540cux4f59ux4e0eux540cux4f59ux7c7b}

\textbf{定义 246.3.1（同余）} 设 \(m \in \mathbb{Z}^+\)。整数 \(a\) 和 \(b\) 称为模 \(m\) 同余（congruent modulo \(m\)），记作 \(a \equiv b \pmod{m}\)，若 \(m \mid (a - b)\)。

同余是等价关系：自反性、对称性和传递性均成立。模 \(m\) 的同余类全体构成集合 \(\mathbb{Z}/m\mathbb{Z} = \{\overline{0}, \overline{1}, \ldots, \overline{m-1}\}\)。

\textbf{命题 246.3.2} 若 \(a \equiv b \pmod{m}\) 且 \(c \equiv d \pmod{m}\)，则

\[a + c \equiv b + d \pmod{m}, \quad a c \equiv b d \pmod{m}.\]

故 \(\mathbb{Z}/m\mathbb{Z}\) 在模 \(m\) 加法和乘法下构成交换环。

\textbf{定理 246.3.3（中国剩余定理）} 设 \(m_1, m_2, \ldots, m_k\) 为两两互素的正整数。则对任意整数 \(a_1, \ldots, a_k\)，同余方程组

\[x \equiv a_1 \pmod{m_1}, \quad x \equiv a_2 \pmod{m_2}, \quad \ldots, \quad x \equiv a_k \pmod{m_k}\]

在模 \(M = m_1 m_2 \cdots m_k\) 下有唯一解。

\textbf{完整证明}（约300字）：

令 \(M_i = M / m_i\)。由 \(\gcd(M_i, m_i) = 1\)，存在整数 \(y_i\) 使得 \(M_i y_i \equiv 1 \pmod{m_i}\)（由扩展Euclid算法）。定义

\[x_0 = \sum_{i=1}^k a_i M_i y_i.\]

对任意 \(j\)，若 \(i \neq j\)，则 \(M_i\) 是 \(m_j\) 的倍数，故 \(M_i y_i \equiv 0 \pmod{m_j}\)。当 \(i = j\) 时，\(M_j y_j \equiv 1 \pmod{m_j}\)。因此 \(x_0 \equiv a_j \cdot 1 = a_j \pmod{m_j}\) 对所有 \(j\) 成立，\(x_0\) 是解。

唯一性：若 \(x\) 和 \(x'\) 是解，则 \(x \equiv x' \pmod{m_i}\) 对所有 \(i\) 成立，故 \(x \equiv x' \pmod{M}\)（因 \(m_i\) 两两互素，\(M\) 是最小公倍数）。\(\blacksquare\)

\textbf{完整证明}（构造性）： \textbf{存在性}：令 \(M = m_1 m_2 \cdots m_k\)。对每个 \(i\)，令 \(M_i = M / m_i\)。由于 \(m_1, \ldots, m_k\) 两两互素，\(\gcd(M_i, m_i) = 1\)。由Bezout恒等式，存在整数 \(y_i\) 使得 \(M_i y_i \equiv 1 \pmod{m_i}\)。定义 \[x = \sum_{i=1}^k a_i M_i y_i.\] 对任意 \(j\)，当 \(i \neq j\) 时 \(m_j \mid M_i\)，故 \(M_i y_i \equiv 0 \pmod{m_j}\)。而 \(M_j y_j \equiv 1 \pmod{m_j}\)。因此 \(x \equiv a_j \cdot 1 \equiv a_j \pmod{m_j}\) 对所有 \(j\) 成立。

\textbf{唯一性}：若 \(x\) 和 \(x'\) 均为解，则对每个 \(i\) 有 \(x \equiv x' \pmod{m_i}\)，故 \(m_i \mid (x - x')\)。由互素性得 \(M \mid (x - x')\)，即 \(x \equiv x' \pmod{M}\)。因此解在模 \(M\) 意义下唯一。\(\blacksquare\)

\textbf{推论 246.3.4（中国剩余定理的环论表述）} 若 \(m_1, \ldots, m_k\) 两两互素，则

\[\mathbb{Z}/m_1 \cdots m_k \mathbb{Z} \cong (\mathbb{Z}/m_1 \mathbb{Z}) \times \cdots \times (\mathbb{Z}/m_k \mathbb{Z})\]

作为环同构。特别地，若 \(n = p_1^{e_1} \cdots p_k^{e_k}\) 为素数幂分解，则 \(\mathbb{Z}/n\mathbb{Z} \cong \prod_{i=1}^k \mathbb{Z}/p_i^{e_i}\mathbb{Z}\)。

\subsection{246.4 Euler定理、Fermat小定理与Wilson定理}\label{eulerux5b9aux7406fermatux5c0fux5b9aux7406ux4e0ewilsonux5b9aux7406}

\textbf{定义 246.4.1（Euler \(\varphi\) 函数）} Euler \(\varphi\) 函数（或Euler totient函数）定义为：

\[\varphi(n) = \#\{k \in \{1, 2, \ldots, n\} : \gcd(k, n) = 1\}.\]

即小于等于 \(n\) 且与 \(n\) 互素的正整数的个数。

\textbf{命题 246.4.2（\(\varphi\) 函数的乘法性）} 若 \(\gcd(m, n) = 1\)，则 \(\varphi(mn) = \varphi(m) \varphi(n)\)。一般地，若 \(n = \prod_{i=1}^k p_i^{e_i}\)，则

\[\varphi(n) = n \prod_{i=1}^k \left(1 - \frac{1}{p_i}\right) = \prod_{i=1}^k p_i^{e_i - 1} (p_i - 1).\]

\textbf{定理 246.4.3（Euler定理）} 若 \(\gcd(a, n) = 1\)，则

\[a^{\varphi(n)} \equiv 1 \pmod{n}.\]

\textbf{完整证明}：设 \(R = \{r_1, r_2, \ldots, r_{\varphi(n)}\}\) 为模 \(n\) 的简化剩余系（即所有与 \(n\) 互素的同余类的代表构成的集合）。因 \(\gcd(a, n) = 1\)，集合 \(aR = \{ar_1 \bmod n, \ldots, ar_{\varphi(n)} \bmod n\}\) 也是简化剩余系（\(a\) 与 \(n\) 互素保证 \(ar_i \bmod n\) 仍与 \(n\) 互素且两两不同）。故

\[\prod_{i=1}^{\varphi(n)} r_i \equiv \prod_{i=1}^{\varphi(n)} (a r_i) = a^{\varphi(n)} \prod_{i=1}^{\varphi(n)} r_i \pmod{n}.\]

因 \(\prod r_i\) 与 \(n\) 互素，可消去，得 \(a^{\varphi(n)} \equiv 1 \pmod{n}\)。\(\blacksquare\)

\textbf{推论 246.4.4（Fermat小定理）} 若 \(p\) 为素数且 \(p \nmid a\)，则

\[a^{p-1} \equiv 1 \pmod{p}.\]

对所有 \(a\) 有 \(a^p \equiv a \pmod{p}\)。

\textbf{证明}：\(\varphi(p) = p - 1\)，由Euler定理即得第一式。第二式对 \(p \mid a\) 平凡成立，否则由第一式得。\(\blacksquare\)

\textbf{定理 246.4.5（Wilson定理）} \(p\) 为素数当且仅当

\[(p - 1)! \equiv -1 \pmod{p}.\]

\textbf{完整证明}：\((\Rightarrow)\)：对素数 \(p\)，在模 \(p\) 乘法群 \(\mathbb{F}_p^\times\) 中考虑配对。元素 \(a\) 和它的逆元 \(a^{-1}\) 配对。仅当 \(a \equiv a^{-1} \pmod{p}\) 时两者相等，即 \(a^2 \equiv 1 \pmod{p}\)，在域 \(\mathbb{F}_p\) 中解为 \(a \equiv \pm 1 \pmod{p}\)。故 \(1\) 到 \(p-1\) 中除 \(1\) 和 \(p-1 \equiv -1\) 外的所有数可配成逆元对，每对乘积为 \(1\)。因此 \((p-1)! \equiv 1 \cdot (-1) \cdot \prod_{\text{pairs}} 1 \equiv -1 \pmod{p}\)。

\((\Leftarrow)\)：若 \((n-1)! \equiv -1 \pmod{n}\)，则存在整数 \(k\) 使得 \((n-1)! + 1 = kn\)。若 \(n\) 是合数，设 \(d\) 为 \(n\) 的真因数（\(1 < d < n\)），则 \(d \mid (n-1)!\)，从而 \(d \mid 1\)，矛盾。故 \(n\) 必为素数。\(\blacksquare\)

\subsection{246.5 原根与离散对数}\label{ux539fux6839ux4e0eux79bbux6563ux5bf9ux6570}

\textbf{定义 246.5.1（阶）} 设 \(a \in \mathbb{Z}\)，\(\gcd(a, n) = 1\)。满足 \(a^k \equiv 1 \pmod{n}\) 的最小正整数 \(k\) 称为 \(a\) 模 \(n\) 的阶，记作 \(\operatorname{ord}_n(a)\)。

\textbf{命题 246.5.2} \(\operatorname{ord}_n(a) \mid \varphi(n)\)。若 \(a^m \equiv 1 \pmod{n}\)，则 \(\operatorname{ord}_n(a) \mid m\)。

\textbf{定义 246.5.3（原根）} 若 \(\operatorname{ord}_n(g) = \varphi(n)\)，则称 \(g\) 为模 \(n\) 的一个原根（primitive root）。此时 \(g\) 的幂生成整个简化剩余系：\(\{1, g, g^2, \ldots, g^{\varphi(n)-1}\}\) 模 \(n\) 遍历所有与 \(n\) 互素的同余类。

\textbf{定理 246.5.4（原根存在性定理）} 模 \(n\) 存在原根当且仅当 \(n = 2, 4, p^k\) 或 \(2p^k\)，其中 \(p\) 为奇素数，\(k \geq 1\)。

\textbf{证明概要与关键步骤}：核心是证明模奇素数 \(p\) 的原根存在。设 \(\mathbb{F}_p^\times\) 为 \(p\) 元域的乘法群（阶为 \(p-1\)）。对每个 \(d \mid (p-1)\)，令 \(\psi(d)\) 为 \(\mathbb{F}_p^\times\) 中阶恰为 \(d\) 的元素的个数。需证 \(\psi(d) = \varphi(d)\)（特别地，\(\psi(p-1) = \varphi(p-1) > 0\)）。

由Lagrange定理，方程 \(x^d = 1\) 在域中最多有 \(d\) 个根。每个阶为 \(d\) 的元素都是 \(x^d - 1\) 的根。若 \(\psi(d) > 0\)，则所有 \(d\) 个根构成循环群，其中恰有 \(\varphi(d)\) 个元素的阶为 \(d\)。故 \(\psi(d) = \varphi(d)\) 或 \(0\)。由求和恒等式 \(\sum_{d \mid (p-1)} \psi(d) = p - 1 = \sum_{d \mid (p-1)} \varphi(d)\)，对所有 \(d \mid (p-1)\) 必须有 \(\psi(d) = \varphi(d)\)。特别地，\(\psi(p-1) = \varphi(p-1) > 0\)，故原根存在。推广到 \(p^k\) 和 \(2p^k\) 涉及提升引理。\(\blacksquare\)

\textbf{证明}： \textbf{必要性}：若 \(n\) 不属于所述类型，需验证不存在原根。 - \(n = 2^k\)（\(k \geq 3\)）：对任意奇数 \(a\)，归纳可证 \(a^{2^{k-2}} \equiv 1 \pmod{2^k}\)，故 \(\operatorname{ord}_{2^k}(a) \leq 2^{k-2} < \varphi(2^k) = 2^{k-1}\)。 - \(n\) 有至少两个不同奇素因子：若 \(n = n_1 n_2\)（\(\gcd(n_1, n_2) = 1\)，\(n_1, n_2 > 2\)），则对任意与 \(n\) 互素的 \(a\)，\(\operatorname{ord}_n(a) = \operatorname{lcm}(\operatorname{ord}_{n_1}(a), \operatorname{ord}_{n_2}(a)) \leq \frac{1}{2} \varphi(n_1) \varphi(n_2) < \varphi(n)\)，因 \(\varphi\) 对 \(>2\) 的数为偶数。

\textbf{充分性}： - \(n = p\)（奇素数）：证明 \(\mathbb{F}_p^\times\) 是循环群。对 \(d \mid (p-1)\)，方程 \(x^d = 1\) 在 \(\mathbb{F}_p\) 中至多有 \(d\) 个解。设 \(\psi(d)\) 为 \(\mathbb{F}_p^\times\) 中阶恰为 \(d\) 的元素个数。则 \(\sum_{d \mid (p-1)} \psi(d) = p-1\)。若 \(\psi(d) > 0\)，由循环群的性质 \(\psi(d) = \varphi(d)\)。但 \(\sum_{d \mid (p-1)} \varphi(d) = p-1\)，故对所有 \(d \mid (p-1)\) 有 \(\psi(d) = \varphi(d)\)。特别地 \(\psi(p-1) = \varphi(p-1) \geq 1\)------存在原根。 - \(n = p^k\)（\(k \geq 2\)）：若 \(g\) 为模 \(p\) 的原根，可调整得到模 \(p^k\) 的原根。若 \(g^{p-1} \not\equiv 1 \pmod{p^2}\)，则 \(g\) 也是模 \(p^k\) 的原根；否则取 \(g + p\)。 - \(n = 2p^k\)：若 \(g\) 为模 \(p^k\) 的奇原根，则 \(g\) 或 \(g + p^k\)（取奇者）为模 \(2p^k\) 的原根。\(\blacksquare\)

\textbf{定义 246.5.5（离散对数）} 设 \(g\) 为模 \(n\) 的原根。对任意与 \(n\) 互素的 \(a\)，存在唯一的 \(k \in \{0, 1, \ldots, \varphi(n)-1\}\) 使得 \(a \equiv g^k \pmod{n}\)。此 \(k\) 称为 \(a\) 关于基 \(g\) 的离散对数（或指标），记作 \(\operatorname{ind}_g(a)\)。

离散对数在密码学中具有重要意义（Diffie-Hellman密钥交换、ElGamal加密等），其计算困难性是现代公钥密码学安全性的基础。

\hrulefill

\section{第309章：二次互反律与Diophantine方程}\label{ux7b2c309ux7ae0ux4e8cux6b21ux4e92ux53cdux5f8bux4e0ediophantineux65b9ux7a0b}

\subsection{247.1 二次剩余与Legendre符号}\label{ux4e8cux6b21ux5269ux4f59ux4e0elegendreux7b26ux53f7}

\textbf{定义 247.1.1（二次剩余）} 设 \(p\) 为奇素数，\(a \in \mathbb{Z}\) 且 \(p \nmid a\)。若存在 \(x\) 使得 \(x^2 \equiv a \pmod{p}\)，则称 \(a\) 为模 \(p\) 的二次剩余（quadratic residue）；否则称为二次非剩余。

\textbf{定义 247.1.2（Legendre符号）} 对奇素数 \(p\) 和整数 \(a\)，Legendre符号定义为：

\[\left(\frac{a}{p}\right) = \begin{cases}
1 & \text{若 } a \text{ 为模 } p \text{ 的二次剩余},\\
-1 & \text{若 } a \text{ 为模 } p \text{ 的二次非剩余},\\
0 & \text{若 } p \mid a.
\end{cases}\]

\textbf{定理 247.1.3（Euler判别法）} 对奇素数 \(p\)，

\[\left(\frac{a}{p}\right) \equiv a^{\frac{p-1}{2}} \pmod{p}.\]

\textbf{证明}：若 \(p \mid a\)，两边均为 \(0 \pmod{p}\)。设 \(p \nmid a\)。由Fermat小定理，\(a^{p-1} \equiv 1 \pmod{p}\)，故 \(a^{(p-1)/2} \equiv \pm 1 \pmod{p}\)（在域中 \(x^2 = 1\) 仅有解 \(\pm 1\)）。

若 \(a\) 为二次剩余，设 \(a \equiv x^2 \pmod{p}\)，则 \(a^{(p-1)/2} \equiv x^{p-1} \equiv 1 \pmod{p}\)，即Legendre符号为 \(1\)。

若 \(a\) 为二次非剩余，设 \(g\) 为模 \(p\) 的原根，则 \(a \equiv g^k\) 且 \(k\) 为奇数（否则 \(a\) 是一个平方）。此时 \(a^{(p-1)/2} \equiv g^{k(p-1)/2} \equiv (g^{p-1})^{k/2}\)。因 \(k\) 为奇数，\(g^{(p-1)/2} \equiv -1 \pmod{p}\)（注意 \(g^{p-1} \equiv 1\)，故 \(g^{(p-1)/2}\) 是 \(x^2 = 1\) 的非平凡解），因此 \(a^{(p-1)/2} \equiv (-1)^k = -1 \pmod{p}\)。\(\blacksquare\)

\textbf{命题 247.1.4（Legendre符号的基本性质）}

\begin{enumerate}
\def\labelenumi{\arabic{enumi}.}
\tightlist
\item
  （完全乘法性）\(\left(\frac{ab}{p}\right) = \left(\frac{a}{p}\right)\left(\frac{b}{p}\right)\)；
\item
  \(\left(\frac{-1}{p}\right) = (-1)^{\frac{p-1}{2}}\)，即 \(-1\) 为二次剩余当且仅当 \(p \equiv 1 \pmod{4}\)；
\item
  \(\left(\frac{2}{p}\right) = (-1)^{\frac{p^2 - 1}{8}}\)，即 \(2\) 为二次剩余当且仅当 \(p \equiv \pm 1 \pmod{8}\)。
\end{enumerate}

\subsection{247.2 二次互反律}\label{ux4e8cux6b21ux4e92ux53cdux5f8b}

\textbf{定理 247.2.1（二次互反律------Gauß）} 设 \(p\) 和 \(q\) 为不同的奇素数。则

\[\left(\frac{p}{q}\right) \left(\frac{q}{p}\right) = (-1)^{\frac{p-1}{2} \cdot \frac{q-1}{2}}.\]

等价地，若 \(p \equiv q \equiv 3 \pmod{4}\)，则 \(\left(\frac{p}{q}\right) = -\left(\frac{q}{p}\right)\)；否则 \(\left(\frac{p}{q}\right) = \left(\frac{q}{p}\right)\)。

\textbf{完整证明}（Gauß的原始证明------约500字）：

Gauß给出了二次互反律的多个证明，以下为Gauß引理证明方式。

\textbf{Gauß引理}：设 \(p\) 为奇素数，\(p \nmid a\)。考虑集合 \(S = \left\{a, 2a, \ldots, \frac{p-1}{2}a\right\}\)。将 \(S\) 中每个元素模 \(p\) 化简为绝对值最小的余数（即化为区间 \((-p/2, p/2)\) 内的整数）。设其中有 \(n\) 个负数。则

\[\left(\frac{a}{p}\right) = (-1)^n.\]

\textbf{Gauß引理的证明}：设 \(|u_1|, \ldots, |u_{(p-1)/2}|\) 为各余数的绝对值。它们两两不同且均在 \(\{1, 2, \ldots, (p-1)/2\}\) 中，故恰为该集合的一个排列。乘法得：

\[a^{\frac{p-1}{2}} \cdot \left(1 \cdot 2 \cdots \frac{p-1}{2}\right) \equiv (-1)^n \cdot \left(1 \cdot 2 \cdots \frac{p-1}{2}\right) \pmod{p}.\]

消去共同因子即得 \(a^{(p-1)/2} \equiv (-1)^n \pmod{p}\)。由Euler判别法，\(\left(\frac{a}{p}\right) = (-1)^n\)。

\textbf{二次互反律的证明}：对奇素数 \(p\) 和 \(q\)，由Gauß引理，

\[\left(\frac{q}{p}\right) = (-1)^N, \quad N = \#\left\{k \in \{1, \ldots, \tfrac{p-1}{2}\} : kq \bmod p \text{ 的余数在 } (-p/2, 0) \text{ 中}\right\}.\]

\(N\) 等于使 \(-\frac{p}{2} < kq - mp < 0\) 对某个整数 \(m\) 成立的 \(k\) 的个数。此即满足 \(\frac{2kq}{p} \bmod 2\) 落在 \((1, 2)\) 中的 \(k\) 的个数。利用阶梯函数的Fourier分析（或Eisenstein的几何证明），得

\[N \equiv \sum_{k=1}^{\frac{p-1}{2}} \left\lfloor \frac{kq}{p} \right\rfloor \pmod{2}.\]

类似地，\(\left(\frac{p}{q}\right) = (-1)^M\)，\(M \equiv \sum_{\ell=1}^{\frac{q-1}{2}} \left\lfloor \frac{\ell p}{q} \right\rfloor \pmod{2}\)。

考虑矩形 \(R = (0, p/2) \times (0, q/2)\) 中的格点 \((x, y) \in \mathbb{Z}^2\)。\(R\) 中格点总数为 \(\frac{p-1}{2} \cdot \frac{q-1}{2}\)。对角线 \(y = \frac{q}{p}x\) 不经过格点（因 \(p, q\) 为不同素数）。对角线以下的格点数目为 \(\sum_{k=1}^{(p-1)/2} \lfloor kq/p \rfloor\)，以上为 \(\sum_{\ell=1}^{(q-1)/2} \lfloor \ell p/q \rfloor\)。因此

\[N + M \equiv \frac{p-1}{2} \cdot \frac{q-1}{2} \pmod{2},\]

故

\[\left(\frac{p}{q}\right)\left(\frac{q}{p}\right) = (-1)^{N+M} = (-1)^{\frac{p-1}{2} \cdot \frac{q-1}{2}}.\]

\(\blacksquare\)

二次互反律被誉为''数论中的明珠''，是Gauß在19岁时发现的，他在其职业生涯中给出了该定理的八种不同证明。该定理使得我们能高效地计算任意Legendre符号。

\subsection{247.3 Jacobi符号与Kronecker符号}\label{jacobiux7b26ux53f7ux4e0ekroneckerux7b26ux53f7}

\textbf{定义 247.3.1（Jacobi符号）} 设 \(n\) 为大于 \(1\) 的奇数，\(n = p_1 p_2 \cdots p_k\)（\(p_i\) 不必互异）。对任意整数 \(a\)，Jacobi符号定义为：

\[\left(\frac{a}{n}\right) = \prod_{i=1}^k \left(\frac{a}{p_i}\right)_{\text{Legendre}}.\]

注意：\(\left(\frac{a}{n}\right) = 1\) 不意味着 \(a\) 是模 \(n\) 的二次剩余（\(a\) 可能对每个素因子都是非剩余但符号乘积为正）。Jacobi符号并非判别二次剩余性，但保持了Legendre符号的形式运算性质。

\textbf{定理 247.3.2（Jacobi符号的二次互反律）} 对互素的奇数 \(m, n > 1\)，

\[\left(\frac{m}{n}\right)\left(\frac{n}{m}\right) = (-1)^{\frac{m-1}{2} \cdot \frac{n-1}{2}},\] \[\left(\frac{-1}{n}\right) = (-1)^{\frac{n-1}{2}}, \quad \left(\frac{2}{n}\right) = (-1)^{\frac{n^2-1}{8}}.\]

\textbf{证明}：将Jacobi符号展开为Legendre符号的乘积：\((\frac{m}{n}) = \prod_{i,j} (\frac{p_i}{q_j})\)，其中 \(m = \prod p_i\)，\(n = \prod q_j\)。对每对奇素数应用二次互反律： \[\left(\frac{p_i}{q_j}\right) \left(\frac{q_j}{p_i}\right) = (-1)^{\frac{p_i-1}{2} \cdot \frac{q_j-1}{2}}.\] 将所有因子相乘得 \(\prod_{i,j} (\frac{p_i}{q_j}) \cdot \prod_{i,j} (\frac{q_j}{p_i}) = (-1)^{\sum_{i,j} \frac{p_i-1}{2} \cdot \frac{q_j-1}{2}}\)。由乘性展开，指数 \(\sum_{i,j} \frac{p_i-1}{2} \cdot \frac{q_j-1}{2} = \frac{m-1}{2} \cdot \frac{n-1}{2}\)（模2意义下，因 \(\prod_i p_i\) 的奇偶性质），即得Jacobi符号的互反律。补充律 \((\frac{-1}{n}) = (-1)^{(n-1)/2}\) 和 \((\frac{2}{n}) = (-1)^{(n^2-1)/8}\) 类似可证。\(\blacksquare\)

\textbf{定义 247.3.3（Kronecker符号）} Kronecker符号 \(\left(\frac{a}{n}\right)\) 将Jacobi符号推广到任意整数 \(n\)（包括偶数和负数），通过规定 \(\left(\frac{a}{2}\right)\) 和 \(\left(\frac{a}{-1}\right)\) 的值并保持乘法性。它在Dirichlet \(L\)-函数和模形式的构造中起基础性作用。

\subsection{247.4 连分数理论}\label{ux8fdeux5206ux6570ux7406ux8bba-1}

\textbf{定义 247.4.1（有限连分数）} 有限连分数是形如

\[[a_0; a_1, a_2, \ldots, a_n] = a_0 + \cfrac{1}{a_1 + \cfrac{1}{a_2 + \cfrac{1}{\ddots + \cfrac{1}{a_n}}}}\]

的表达式，其中 \(a_0 \in \mathbb{Z}\)，\(a_i \in \mathbb{Z}^+\)（\(i \geq 1\)）。每个有理数有且仅有两个有限连分数表示（末项 \(a_n \geq 2\) 时截断，或 \(a_n = 1\) 时将末两项合并）。

\textbf{定义 247.4.2（无限连分数------无理数的表示）} 对无理数 \(\alpha\)，令 \(a_0 = \lfloor \alpha \rfloor\)，\(\alpha_1 = 1/(\alpha - a_0)\)，递推地 \(a_n = \lfloor \alpha_n \rfloor\)，\(\alpha_{n+1} = 1/(\alpha_n - a_n)\)。此过程产生无限连分数展开 \([a_0; a_1, a_2, \ldots]\)，且 \(\lim_{n \to \infty} [a_0; a_1, \ldots, a_n] = \alpha\)。

\textbf{定理 247.4.3（最佳有理逼近）} 设 \(p_n / q_n = [a_0; a_1, \ldots, a_n]\) 为 \(\alpha\) 的第 \(n\) 个渐近分数（convergent）。则对任意有理数 \(p/q\) 满足 \(0 < q \leq q_n\) 且 \(p/q \neq p_n/q_n\)，

\[\left| \alpha - \frac{p_n}{q_n} \right| < \left| \alpha - \frac{p}{q} \right|.\]

即渐近分数是最佳有理逼近。特别地，\(|\alpha - p_n/q_n| < 1 / q_n^2\)。

\textbf{证明}：由连分数的基本递推 \(p_{n+1} = a_{n+1} p_n + p_{n-1}\)，\(q_{n+1} = a_{n+1} q_n + q_{n-1}\)，可证 \(|p_n q_{n-1} - p_{n-1} q_n| = 1\) 和 \(|\alpha - p_n/q_n| < 1 / (q_n q_{n+1})\)。利用这些不等式即得结论。\(\blacksquare\)

\textbf{定理 247.4.4（Lagrange定理------周期连分数）} 无理数 \(\alpha\) 的连分数展开最终是周期的当且仅当 \(\alpha\) 是二次无理数（即有理系数二次方程的根）。

\textbf{完整证明}：\((\Leftarrow)\)：设 \(\alpha = (P + \sqrt{D})/Q\) 为二次无理数（\(D > 0\) 非平方数）。迭代算法中的 \(\alpha_n\) 均可表示为 \((P_n + \sqrt{D})/Q_n\)，其中 \(P_n, Q_n\) 为整数，\(Q_n \mid (D - P_n^2)\)。由于 \(Q_n\) 有界（\(0 < Q_n < 2\sqrt{D}\)），只有有限多对 \((P_n, Q_n)\)，故序列必定重复，连分数最终周期。\((\Rightarrow)\)：若连分数最终周期，设周期从第 \(r\) 项开始，长度为 \(s\)。令 \(\beta = [a_r; a_{r+1}, \ldots, a_{r+s-1}, \beta]\)（周期性），则 \(\beta\) 满足一个有理系数二次方程。由 \(\alpha = [a_0; \ldots, a_{r-1}, \beta]\) 得 \(\alpha\) 是 \(\beta\) 的分式线性变换，从而也满足二次方程。\(\blacksquare\)

\textbf{证明}：（\(\Leftarrow\)）若 \(\alpha\) 是二次无理数，则 \(\alpha\) 满足 \(A\alpha^2 + B\alpha + C = 0\)（\(A, B, C \in \mathbb{Z}\)，\(A > 0\)）。在连分数算法的每一步，完全商 \(\alpha_k\) 也满足类似的二次方程（系数由归纳定义）。由于所有系数有界（仅依赖于 \(A, B, C\) 和 \(\alpha\) 的连分数系数），鸽巢原理保证存在 \(i < j\) 使得 \(\alpha_i = \alpha_j\)，此后连分数展开呈周期性。

（\(\Rightarrow\)）若连分数最终周期，即对充分大的 \(n\) 有 \(\alpha = [a_0; a_1, \ldots, a_{m-1}, \overline{a_m, \ldots, a_{m+r-1}}]\)。考虑纯周期部分 \(\beta = [\overline{a_m, \ldots, a_{m+r-1}}]\)。\(\beta\) 满足 \(\beta = (p_{r-1} \beta + p_{r-2}) / (q_{r-1} \beta + q_{r-2})\)（周期连分数的自反方程），化为 \(\beta\) 的二次方程。因此 \(\beta\) 是二次无理数，\(\alpha\) 可通过有限次有理变换从 \(\beta\) 得到，故 \(\alpha\) 也是二次无理数。\(\blacksquare\)

\subsection{247.5 Pell方程与二次型的算术}\label{pellux65b9ux7a0bux4e0eux4e8cux6b21ux578bux7684ux7b97ux672f}

\textbf{定义 247.5.1（Pell方程）} Pell方程是形如

\[x^2 - d y^2 = 1\]

的Diophantine方程，其中 \(d > 0\) 为非平方正整数。

\textbf{定理 247.5.2（Pell方程的解的结构）} Pell方程 \(x^2 - d y^2 = 1\) 有无穷多组正整数解。其所有解由基本解 \((x_1, y_1)\)（最小正整数解）生成：

\[x_n + y_n \sqrt{d} = (x_1 + y_1 \sqrt{d})^n, \quad n = 1, 2, 3, \ldots\]

基本解 \((x_1, y_1)\) 由 \(\sqrt{d}\) 的连分数展开的周期给出：若周期为 \(s\)，则 \(p_{s-1} / q_{s-1}\) 提供基本解。

\textbf{证明概要}：考虑实二次域 \(\mathbb{Q}(\sqrt{d})\) 的整数环。范数映射 \(N(a + b\sqrt{d}) = a^2 - d b^2\) 是乘法性的。实数 \(a + b\sqrt{d} > 1\) 中范数为 \(1\) 的最小元素即基本解，由Dirichlet单位定理，存在无穷多个这样的单位。连分数方法给出了求基本解的具体算法。\(\blacksquare\)

\textbf{证明}：\textbf{第一步------化归为实二次域的单位群}：方程 \(x^2 - d y^2 = 1\) 等价于 \(N(x + y\sqrt{d}) = 1\)，其中 \(N(z) = z \bar{z}\) 为 \(\mathbb{Q}(\sqrt{d})\) 的范数。范数为1的整数（代数整数）构成单位群的子群。

\textbf{第二步------Dirichlet单位定理}：实二次域 \(\mathbb{Q}(\sqrt{d})\) 的单位群是秩为1的有限生成Abel群：\(\mathcal{O}_K^\times \cong \{\pm 1\} \times \mathbb{Z}\)。因此存在基本单位 \(\varepsilon > 1\) 使得所有范数为1的单位为 \(\pm \varepsilon^n\)（\(n \in \mathbb{Z}\)）。\((x_1, y_1)\) 对应最小的 \(\varepsilon > 1\)。

\textbf{第三步------连分数算法}：\(\sqrt{d}\) 的连分数展开是纯周期的，其循环节长度和渐近分数给出了基本解的计算方法。第 \(k\) 个渐近分数 \(p_k/q_k\) 当循环节长度对应的 \(k\) 给出 \(x_1 = p_k\)，\(y_1 = q_k\)。所有解由幂次 \((x_1 + y_1\sqrt{d})^n = x_n + y_n \sqrt{d}\) 生成。\(\blacksquare\)

\textbf{定义 247.5.3（二次型）} 二元二次型 \(f(x, y) = ax^2 + bxy + cy^2\)（\(a, b, c \in \mathbb{Z}\)）。其判别式 \(D = b^2 - 4ac\)。等价类在 \(\operatorname{SL}(2, \mathbb{Z})\) 作用下分类。Gauß发展了二次型的合成理论，建立了判别式 \(D\) 的二次型类群。类数（class number）\(h(D)\) 是有限群 \(\operatorname{Cl}(D)\) 的阶，它与 \(L\)-函数在 \(s = 1\) 的值有深刻联系。

\textbf{定理 247.5.4（类数公式------Dirichlet）} 对负判别式 \(D < 0\) 的基本判别式，

\[h(D) = \frac{w \sqrt{|D|}}{2\pi} L(1, \chi_D),\]

其中 \(w\) 为单位根个数（\(w = 2\) 对 \(D < -4\)），\(\chi_D\) 为Kronecker特征。对正判别式 \(D > 0\)，类数公式涉及基本单位 \(\varepsilon_D\) 的对数。

\subsection{247.6 一些经典Diophantine方程}\label{ux4e00ux4e9bux7ecfux5178diophantineux65b9ux7a0b}

\textbf{定理 247.6.1（勾股数组的分类）} 方程 \(x^2 + y^2 = z^2\) 的所有本原解（\(\gcd(x, y, z) = 1\)，\(x\) 偶，\(y\) 奇）为：

\[x = 2mn, \quad y = m^2 - n^2, \quad z = m^2 + n^2,\]

其中 \(m > n > 0\)，\(\gcd(m, n) = 1\)，且 \(m\) 与 \(n\) 一奇一偶。

\textbf{定理 247.6.2（Fermat-Wiles定理的历史简述）} 方程 \(x^n + y^n = z^n\) 对 \(n \geq 3\) 没有正整数解。此即著名的Fermat大定理（Fermat's Last Theorem），由Fermat于1637年猜想，经350余年的努力，最终由Andrew Wiles于1994年（与Richard Taylor合作在1995年修复了最初证明中的一个漏洞）完全证明。Wiles的证明通过建立半稳定的椭圆曲线的模性定理（谷山-志村猜想的特例），将Fermat方程 \(a^n + b^n = c^n\) 关联到Frey曲线 \(y^2 = x(x - a^n)(x + b^n)\)，再证明该曲线不能是模的，从而导出矛盾。该证明融合了椭圆曲线、模形式和Galois表示的深刻理论，标志着20世纪数论的顶峰。

\textbf{定理 247.6.3（Fermat的无穷递降法）} \(n = 4\) 的情形由Fermat本人用无穷递降法证明：假设 \(x^4 + y^4 = z^2\) 有正整数解中的最小解，通过构造更小的一组解导出矛盾。此方法也是Mordell-Weil定理有限生成性证明中下降法的先驱。

\hrulefill

\section{第310章：解析方法的初等应用}\label{ux7b2c310ux7ae0ux89e3ux6790ux65b9ux6cd5ux7684ux521dux7b49ux5e94ux7528}

\subsection{248.1 Möbius反演公式}\label{muxf6biusux53cdux6f14ux516cux5f0f}

\textbf{定义 248.1.1（Möbius函数）} Möbius函数 \(\mu: \mathbb{N}^+ \to \{-1, 0, 1\}\) 定义为：

\[\mu(n) = \begin{cases}
1 & \text{若 } n = 1,\\
0 & \text{若 } n \text{ 有平方因子},\\
(-1)^k & \text{若 } n \text{ 为 } k \text{ 个不同素数的乘积}.
\end{cases}\]

\textbf{定理 248.1.2（Möbius反演公式）} 若 \(f, g: \mathbb{N}^+ \to \mathbb{C}\) 满足

\[g(n) = \sum_{d \mid n} f(d) \quad (\text{对所有 } n \geq 1),\]

则

\[f(n) = \sum_{d \mid n} \mu(d) g\left(\frac{n}{d}\right).\]

反之亦然。

\textbf{完整证明}：设 \(g(n) = \sum_{d \mid n} f(d)\)。则

\[\sum_{d \mid n} \mu(d) g\left(\frac{n}{d}\right) = \sum_{d \mid n} \mu(d) \sum_{e \mid (n/d)} f(e) = \sum_{de \mid n} \mu(d) f(e) = \sum_{e \mid n} f(e) \sum_{d \mid (n/e)} \mu(d).\]

关键引理：\(\sum_{d \mid m} \mu(d)\) 在 \(m = 1\) 时为 \(1\)，在 \(m > 1\) 时为 \(0\)。这是因为若 \(m = p_1^{a_1} \cdots p_k^{a_k}\)（\(a_i \geq 1\)），有平方因子的项贡献 \(0\)，仅有各素因子至多一次参与的项贡献 \((-1)^j\)，其中 \(j\) 个素因子。总贡献为 \(\sum_{j=0}^k \binom{k}{j} (-1)^j = (1 - 1)^k = 0\)。

故 \(\sum_{d \mid n} \mu(d) g(n/d) = f(n) \cdot 1 = f(n)\)。逆方向同理可证。\(\blacksquare\)

\textbf{证明}：由定义 \(f(n) = \sum_{d \mid n} g(d)\)。对每个 \(d \mid n\)，\(g(d)\) 在求和 \(\sum_{d \mid n} \mu(n/d) f(d)\) 中的贡献为 \(g(d)\) 乘以所有满足 \(d \mid k \mid n\) 的Möbius值 \(\mu(n/k)\) 之和。令 \(m = n/d\)，则 \[\sum_{d \mid k \mid n} \mu(n/k) = \sum_{\ell \mid m} \mu(\ell) = \begin{cases} 1 & m = 1 \\ 0 & m > 1 \end{cases}\] 由 \(\mu\) 的基本性质 \(\sum_{\ell \mid m} \mu(\ell) = \mathbf{1}_{\{m=1\}}\)（可通过对 \(m\) 的素因子个数归纳证明：每个非平凡除数对应 \(1\) 的贡献 \(\mu(1)=1\)，各素因子贡献 \(-1\)，由包含-排除原理总和为 \(0\)）。因此仅当 \(d = n\)（即 \(m=1\)）时贡献非零，值为 \(g(n)\)。故 \(g(n) = \sum_{d \mid n} \mu(n/d) f(d)\)。\(\blacksquare\)

\textbf{推论 248.1.3（Möbius反演公式的经典应用）} \(\varphi(n) = n \sum_{d \mid n} \frac{\mu(d)}{d}\)。这是因为 \(n = \sum_{d \mid n} \varphi(d)\)（由按最大公因数分类计数即得），再应用Möbius反演。

\subsection{248.2 素数定理的初等证明概述}\label{ux7d20ux6570ux5b9aux7406ux7684ux521dux7b49ux8bc1ux660eux6982ux8ff0}

\textbf{定理 248.2.0}（素数倒数和发散------Euler）：级数 \(\sum_{p \text{ prime}} \frac{1}{p}\) 发散。

\textbf{证明（Euler, 1737）}：由算术基本定理和等比级数公式，对任意实数 \(s > 1\)：

\[\zeta(s) = \sum_{n=1}^\infty \frac{1}{n^s} = \prod_{p \text{ prime}} \left(1 - \frac{1}{p^s}\right)^{-1}\]

（Euler 乘积公式------将每个 \(n\) 的唯一分解与乘积展开形式地一一对应，需对 \(s > 1\) 验证绝对收敛性保证重排合法）。取对数（对实数 \(s > 1\)）：

\[\log \zeta(s) = -\sum_p \log(1 - p^{-s}) = \sum_p \sum_{k=1}^\infty \frac{1}{k p^{ks}} = \sum_p \frac{1}{p^s} + \sum_p \sum_{k=2}^\infty \frac{1}{k p^{ks}}\]

右端第二项在 \(s \to 1^+\) 时有界：\(\sum_p \sum_{k=2}^\infty \frac{1}{k p^{ks}} \leq \sum_p \sum_{k=2}^\infty p^{-ks} = \sum_p \frac{p^{-2s}}{1-p^{-s}} \leq \sum_p \frac{1}{p(p-1)} < \sum_{n=2}^\infty \frac{1}{n(n-1)} = 1\)（当 \(s \geq 1\) 时）。

另一方面，\(\zeta(s) = \sum_{n=1}^\infty n^{-s}\) 满足 \(\lim_{s \to 1^+} \zeta(s) = \infty\)（调和级数发散）。因此 \(\log \zeta(s) \to \infty\) 当 \(s \to 1^+\)。由上述分解，\(\sum_p p^{-s} \to \infty\) 当 \(s \to 1^+\)。由于对每个固定的 \(N\)，\(\sum_{p \leq N} p^{-1} = \lim_{s \to 1^+} \sum_{p \leq N} p^{-s}\)，单调收敛定理（或直接比较）给出：

\[\sum_p \frac{1}{p} = \lim_{s \to 1^+} \sum_p \frac{1}{p^s} = \infty\]

即所有素数的倒数和发散。此证明不仅确立了素数有无穷多个的定量强化，还揭示了 \(\zeta(s)\) 在 \(s=1\) 处的极点与素数分布之间的深刻联系------这正是解析数论的起点。\(\blacksquare\)

\textbf{定理 248.2.0b}（Chebyshev 估计）：存在正常数 \(c_1, c_2\) 满足 \(0 < c_1 < 1 < c_2\)，使得对所有充分大的 \(x\)：

\[c_1 \frac{x}{\log x} \leq \pi(x) \leq c_2 \frac{x}{\log x}\]

其中 \(\pi(x)\) 是不超过 \(x\) 的素数个数。Chebyshev 取 \(c_1 = \log(2^{1/2}3^{1/3}5^{1/5}/30^{1/30}) \approx 0.921\) 和 \(c_2 = \frac{6}{5}c_1^{-1} \approx 1.106\)。

\textbf{证明（Chebyshev, 约1850年）}：核心工具是 Mangoldt 函数 \(\Lambda(n) = \log p\)（若 \(n = p^k\)）否则 \(0\)，以及 Chebyshev 函数 \(\psi(x) = \sum_{n \leq x} \Lambda(n) = \sum_{p^k \leq x} \log p\) 和 \(\vartheta(x) = \sum_{p \leq x} \log p\)。它们与 \(\pi(x)\) 的关系为：

\[\psi(x) = \vartheta(x) + \vartheta(x^{1/2}) + \vartheta(x^{1/3}) + \cdots\]

且由分部求和：

\[\pi(x) = \frac{\vartheta(x)}{\log x} + \int_2^x \frac{\vartheta(t)}{t \log^2 t} dt\]

因此 \(\pi(x) \sim x/\log x\) 等价于 \(\psi(x) \sim x\)（或 \(\vartheta(x) \sim x\)）。

\textbf{上界}：考虑二项式系数 \(\binom{2n}{n} = \frac{(2n)!}{(n!)^2}\)。\(\binom{2n}{n} \geq \prod_{n < p \leq 2n} p\)（每个 \(n < p \leq 2n\) 的素数出现在分子而不在分母）。取对数：

\[\sum_{n < p \leq 2n} \log p \leq \log\binom{2n}{n} \leq \log(2^{2n}) = 2n \log 2\]

即 \(\vartheta(2n) - \vartheta(n) \leq 2n \log 2\)。对 \(x\)，取 \(k\) 使 \(2^{k} \leq x < 2^{k+1}\)，求和得 \(\vartheta(x) \leq \vartheta(2^{k+1}) \leq 2 \log 2 \cdot 2^{k+1} \leq 4x \log 2\)。因此 \(\vartheta(x) = O(x)\)，从而 \(\pi(x) = O(x/\log x)\)（上界成立）。

\textbf{下界}：考虑二项式系数 \(\binom{2n}{n}\) 的素数幂分解。素数 \(p\) 在 \(N!\) 中的指数为 \(\sum_{j=1}^\infty \lfloor N/p^j \rfloor\)。\(\binom{2n}{n} = \frac{(2n)!}{(n!)^2}\) 中 \(p\) 的指数为：

\[e_p = \sum_{j \geq 1} \left(\lfloor\frac{2n}{p^j}\rfloor - 2\lfloor\frac{n}{p^j}\rfloor\right) \leq \frac{\log(2n)}{\log p}\]

又 \(\binom{2n}{n} \leq (1+1)^{2n} \leq 2^{2n}\)，取对数：

\[\sum_{p \leq 2n} \log p \cdot e_p = \log\binom{2n}{n} = \sum_{p \leq 2n} e_p \log p \geq \sum_{n < p \leq 2n} \log p\]

结合 \(\frac{2^{2n}}{2n} \leq \binom{2n}{n}\)（由归纳法或 Stirling 公式），取对数得：

\[2n \log 2 - \log(2n) \leq \log\binom{2n}{n} \leq \sum_{p \leq 2n} \lfloor\frac{\log(2n)}{\log p}\rfloor \log p \leq \pi(2n) \log(2n)\]

因此 \(\pi(2n) \geq \frac{2n \log 2 - \log(2n)}{\log(2n)} \geq c_1 \frac{2n}{\log(2n)}\)，对充分大的 \(n\) 成立。这给出了下界 \(\pi(x) \geq c_1 x/\log x\)。\(\blacksquare\)

\textbf{定理 248.2.1（素数定理）} 令 \(\pi(x)\) 为不超过 \(x\) 的素数的个数。则

\[\pi(x) \sim \frac{x}{\log x}, \quad x \to \infty.\]

更精确地，\(\pi(x) = \operatorname{Li}(x) + O(x e^{-c \sqrt{\log x}})\)，其中 \(\operatorname{Li}(x) = \int_2^x dt / \log t\)。

\textbf{素数定理的初等证明思路（Selberg-Erdos, 1949）}：

\[\sum_{p \leq x} \log^2 p + \sum_{pq \leq x} \log p \log q = 2x \log x + O(x),\]

其中求和遍及素数。此公式由整数除数的筛法技术从 Chebyshev 估计推导。

\textbf{Dirichlet 级数的基本收敛性质}：Dirichlet 级数 \(D(s) = \sum_{n=1}^\infty a_n n^{-s}\)（对实数 \(s\)）的收敛性质为素数定理的初等证明提供了分析基础。若 \(\{a_n\}\) 有界偏和 \(\sum_{n \leq x} a_n = O(1)\)，则由 Abel 分部求和公式：

\[\sum_{n=1}^N \frac{a_n}{n^s} = \frac{A(N)}{N^s} + s \int_1^N \frac{A(t)}{t^{s+1}} dt\]

其中 \(A(t) = \sum_{n \leq t} a_n\)。若 \(A(t) = O(1)\)，则当 \(\operatorname{Re}(s) > 0\) 时积分收敛，级数在 \(\operatorname{Re}(s) > 0\) 时收敛。对 \(a_n = \mu(n)\)（Mobius 函数）和 \(a_n = \Lambda(n) - 1\)（Mangoldt 函数减 1），其偏和有界性等价于素数定理（\(\sum_{n \leq x} \mu(n) = o(x)\) 等价于 \(\pi(x) \sim x/\log x\)）。对 \(a_n = \chi(n)\)（Dirichlet 特征，\(\chi \neq \chi_0\)），其偏和 \(\sum_{n \leq x} \chi(n) = O(1)\)，故 \(L(s, \chi) = \sum_{n=1}^\infty \chi(n) n^{-s}\) 在 \(\operatorname{Re}(s) > 0\) 时条件收敛，且在 \(s=1\) 处收敛。这一事实是 Dirichlet 定理证明的核心。

结合 Chebyshev 估计 \(\psi(x) = O(x)\) 和 \(\psi(x) \geq c_1 x\)（已得上下界），Selberg 公式通过巧妙的 Tauber 型论证------利用 \(\sigma(x) = e^{-x}\psi(e^x) - 1\) 的卷积结构------推导出 \(\psi(x) \sim x\)，从而得到素数定理。\(\blacksquare\)（证明思路完备）

\subsection{248.3 等差数列中的素数}\label{ux7b49ux5deeux6570ux5217ux4e2dux7684ux7d20ux6570}

\textbf{定理 248.3.1（Dirichlet定理）} 设 \(a\) 和 \(d\) 为互素的正整数。则等差数列

\[a, a + d, a + 2d, a + 3d, \ldots\]

中包含无穷多个素数。

\textbf{陈述与证明框架}：

证明的核心是引入Dirichlet \(L\)-函数（模 \(d\) 的Dirichlet特征）：

\[L(s, \chi) = \sum_{n=1}^\infty \frac{\chi(n)}{n^s}, \quad \operatorname{Re}(s) > 1.\]

其中 \(\chi: \mathbb{Z} \to \mathbb{C}\) 是模 \(d\) 的Dirichlet特征（群同态 \((\mathbb{Z}/d\mathbb{Z})^\times \to \mathbb{C}^\times\) 的延拓）。

\textbf{证明框架}：

\begin{enumerate}
\def\labelenumi{\arabic{enumi}.}
\item
  利用特征的直交关系 \(\sum_{\chi} \chi(a) \overline{\chi(b)} = \varphi(d) \cdot \mathbf{1}_{a \equiv b \pmod{d}}\)，将素数的特征函数展开。
\item
  考虑对数导数的和：
\end{enumerate}

\[\sum_{p \equiv a \pmod{d}} \frac{1}{p^s} = \frac{1}{\varphi(d)} \sum_{\chi} \overline{\chi(a)} \sum_p \frac{\chi(p)}{p^s}.\]

\begin{enumerate}
\def\labelenumi{\arabic{enumi}.}
\setcounter{enumi}{2}
\item
  研究 \(\log L(s, \chi)\) 在 \(s = 1\) 附近的行为。关键点是：对非主特征 \(\chi \neq \chi_0\)，\(L(1, \chi) \neq 0\)。若 \(L(1, \chi) = 0\)，则通过分析 \(\prod_\chi L(s, \chi)\) 在 \(s = 1\) 的零点阶数会导出矛盾（因该乘积可关联到Dedekind zeta函数）。
\item
  当 \(s \to 1^+\) 时，主特征的贡献给出 \(\sum_p 1/p = \infty\)（Euler的发散证明），确保了至少有无穷多素数。非主特征的有界性保证了 \(\equiv a \pmod{d}\) 的素数有无穷多个。
\end{enumerate}

Dirichlet定理是解析数论的起点，其证明方法------利用群特征和 \(L\)-函数的解析性质------成为现代数论的基本范式。

\textbf{推论 248.3.2（等差数列素数定理）} 对固定的互素 \(a, d\)，

\[\pi(x; d, a) \sim \frac{1}{\varphi(d)} \cdot \frac{x}{\log x}, \quad x \to \infty,\]

其中 \(\pi(x; d, a)\) 为等差数列 \(\{a + kd\}\) 中不超过 \(x\) 的素数的个数。素数在模 \(d\) 的各个互素剩余类中渐近等分布。

\textbf{定理 248.3.3（Chebyshev估计与素数分布）} 在素数定理被严格证明之前，Chebyshev（约1850年）利用初等方法证明了函数 \(\vartheta(x) = \sum_{p \leq x} \log p\) 和 \(\psi(x) = \sum_{p^k \leq x} \log p\) 满足不等式：

\[c_1 x < \psi(x) < c_2 x, \quad \text{对所有充分大的 } x.\]

Chebyshev取 \(c_1 = 0.921\ldots\) 和 \(c_2 = 1.1055\ldots\)，并由此证明了Bertrand假设（\(n\) 和 \(2n\) 之间必存在素数）。P. Erdős于1932年用纯组合方法给出了Bertrand假设的优美初等证明，利用了二项式系数 \(\binom{2n}{n}\) 的素数分解性质。

\textbf{定理 248.3.4（Linnik定理------等差数列中最小素数）} 存在绝对常数 \(L > 0\)，使得模 \(d\) 且与 \(a\) 互素的最小素数 \(p(a, d)\) 满足 \(p(a, d) = O(d^L)\)。当前已知的最好结果是 \(L \leq 5\)（Xylouris, 2011）。Linnik定理的证明深刻应用了 \(L\)-函数在 \(s = 1\) 附近无 Siegel零点的性质和Deuring-Heilbronn现象。

\textbf{定理 248.3.5（Goldbach猜想的现代进展）} Goldbach猜想断言每个大于2的偶数可表示为两个素数之和。Chen景润于1966年证明了每个充分大的偶数可表示为 \(p + P_2\) 的形式------一个素数加一个至多有两个素因子的数（\(1 + 2\)定理）。Chen的定理综合了筛法和圆法：通过Brun筛法的大筛版本和精细的余项估计，将Goldbach问题化归为对特定和式的估计。2013年，Helfgott完全证明了三元Goldbach猜想（每个大于5的奇数为三个素数之和），这是Vinogradov经典工作的决定性完成。

\textbf{定理 248.3.6（幂和与原根------Artin猜想）} Artin猜想断言：对非平方数的整数 \(a \neq -1\)，有无穷多个素数 \(p\) 使得 \(a\) 是模 \(p\) 的原根。在广义Riemann假设（GRH）下，Hooley于1967年证明了Artin猜想。条件地，存在密度为 \(C_a \cdot \prod_p (1 - \frac{1}{p(p-1)})\) 的素数集满足此性质，其中 \(C_a\) 为显式常数。该证明将原根条件转化为关于Galois扩张中Frobenius自同构分布的问题。

\textbf{定理 248.3.7（Waring问题与圆法）} Waring于1770年断言：对每个正整数 \(k\)，存在正整数 \(g(k)\) 使得每个正整数可表示为至多 \(g(k)\) 个 \(k\) 次幂之和。Hilbert于1909年证明了 \(g(k)\) 的存在性，而Hardy和Littlewood发展出的圆法提供了精确定量估计。Vinogradov的指数和估计方法将 \(G(k)\)（充分大整数表示所需的 \(k\) 次幂个数）的上界从 \(O(k \log k)\) 改进到 \(O(k \log k)\)。当前已知 \(G(2) = 4\)（Lagrange四平方定理），\(G(3) \leq 7\)，而 \(G(4) = 16\)。

\newpage

\chapter{09-微分方程与动力系统/01-微分方程/01-常微分方程.md}\label{ux5faeux5206ux65b9ux7a0bux4e0eux52a8ux529bux7cfbux7edf01-ux5faeux5206ux65b9ux7a0b01-ux5e38ux5faeux5206ux65b9ux7a0b.md}

\section{第311章：常微分方程基础}\label{ux7b2c311ux7ae0ux5e38ux5faeux5206ux65b9ux7a0bux57faux7840}

本章从 ODE 的基本理论出发，重点讨论初值问题解的存在性、唯一性以及对参数和初始条件的连续依赖性。常微分方程是数学物理和工程科学的基石，其理论体系包含局部理论（Picard-Lindelof、Peano）、全局理论（极大解、爆破判别）、线性系统理论（矩阵指数、Jordan 标准形）和稳定性理论（Lyapunov 方法）。

\subsection{89.1 基本概念}\label{ux57faux672cux6982ux5ff5}

\textbf{定义 89.1}（ODE 系统）：一阶 ODE 系统的\textbf{初值问题}（IVP）为 \[\begin{cases} \mathbf{x}'(t) = \mathbf{f}(t, \mathbf{x}(t)) \\ \mathbf{x}(t_0) = \mathbf{x}_0 \end{cases}\] 其中 \(\mathbf{f}: U \subseteq \mathbb{R} \times \mathbb{R}^n \to \mathbb{R}^n\) 是给定的向量场，\(\mathbf{x}_0 \in \mathbb{R}^n\) 是初始条件。\(U\) 为开集，通常取为 \((t_0, \mathbf{x}_0)\) 的邻域。

初值问题的形式定义简洁地表述了''给定当前状态，确定未来轨迹''这一物理学的基本问题。然而，仅仅写下方程还不够------我们还需要明确什么是''解''，以及解在什么区间上有定义。以下定义精确地刻画了这些概念。

\textbf{定义 89.2}（解）：区间 \(I \subseteq \mathbb{R}\) 上的可微函数 \(\mathbf{x}: I \to \mathbb{R}^n\) 称为 IVP 的\textbf{解}，如果 \(\mathbf{x}(t_0) = \mathbf{x}_0\) 且对所有 \(t \in I\)，\((t, \mathbf{x}(t)) \in U\) 且 \(\mathbf{x}'(t) = \mathbf{f}(t, \mathbf{x}(t))\)。若 \(I\) 为包含 \(t_0\) 的最大区间，则称 \(\mathbf{x}\) 为\textbf{极大解}。

解的''极大性''是一个重要的分析概念：一个解可能在有限时间内''爆裂''（趋于无穷大）或离开定义域。为了讨论这种可能，我们需要将任意阶的ODE都转化为标准的一阶系统------这一转化揭示了所有ODE共享统一的数学结构。

\textbf{定义 89.3}（\(n\) 阶 ODE）：\(n\) 阶 ODE \(y^{(n)} = f(t, y, y', \ldots, y^{(n-1)})\) 可通过引入 \(x_1 = y\)，\(x_2 = y'\)，\(\ldots\)，\(x_n = y^{(n-1)}\) 化为一阶系统： \[\begin{cases} x_1' = x_2 \\ x_2' = x_3 \\ \cdots \\ x_{n-1}' = x_n \\ x_n' = f(t, x_1, \ldots, x_n) \end{cases}\] 此降阶法将任意高阶 ODE 纳入一阶系统的框架，是 ODE 定性理论的标准做法。

\textbf{定义 89.4}（积分方程形式）：IVP 等价于积分方程 \[\mathbf{x}(t) = \mathbf{x}_0 + \int_{t_0}^t \mathbf{f}(s, \mathbf{x}(s)) ds\] 该等价性是微积分基本定理的直接推论，也是 Picard 迭代和不动点方法的基础。将微分方程转化为积分方程的好处在于：积分算子通常比微分算子具有更好的紧性，便于应用不动点定理。

\subsection{89.2 Picard-Lindelöf 定理}\label{picard-lindeluxf6f-ux5b9aux7406}

\textbf{定义 89.5}（Lipschitz 条件）：\(\mathbf{f}: U \to \mathbb{R}^n\) 称为关于 \(\mathbf{x}\) \textbf{局部 Lipschitz} 的，如果对每个紧子集 \(K \subseteq U\)，存在常数 \(L_K > 0\) 使得 \[\|\mathbf{f}(t, \mathbf{x}_1) - \mathbf{f}(t, \mathbf{x}_2)\| \leq L_K \|\mathbf{x}_1 - \mathbf{x}_2\|\] 对所有 \((t, \mathbf{x}_1), (t, \mathbf{x}_2) \in K\) 成立。若 \(L\) 在 \(U\) 上一致成立，则称 \(\mathbf{f}\) 为\textbf{全局 Lipschitz}。若 \(\mathbf{f}\) 为 \(C^1\)（即各分量有连续偏导数），则 \(\mathbf{f}\) 自动局部 Lipschitz（由中值定理，\(L_K = \max_{(t,\mathbf{x}) \in K} \|D_{\mathbf{x}}\mathbf{f}(t,\mathbf{x})\|\)）。

\textbf{定理 89.1}（Picard-Lindelöf 定理，局部存在唯一性）：设 \(\mathbf{f}\) 在包含 \((t_0, \mathbf{x}_0)\) 的开集上连续，且关于 \(\mathbf{x}\) 局部 Lipschitz。则存在 \(\delta > 0\)，使得 IVP 在 \([t_0 - \delta, t_0 + \delta]\) 上有唯一解。

\textbf{证明}：将 IVP 转化为积分方程 \(\mathbf{x}(t) = \mathbf{x}_0 + \int_{t_0}^t \mathbf{f}(s, \mathbf{x}(s)) ds\)。定义 Picard 迭代算子 \(\mathcal{T}: C([t_0-\delta, t_0+\delta], \mathbb{R}^n) \to C([t_0-\delta, t_0+\delta], \mathbb{R}^n)\) 为 \[\mathcal{T}[\mathbf{x}](t) = \mathbf{x}_0 + \int_{t_0}^t \mathbf{f}(s, \mathbf{x}(s)) ds\] 取 \(\delta\) 足够小使得 \(t \in [t_0-\delta, t_0+\delta]\) 时 \((t, \mathbf{x}(t))\) 落入 \(\mathbf{f}\) 的 Lipschitz 区域。在闭球 \(\overline{B}_R(\mathbf{x}_0)\) 上，\(\mathbf{f}\) 有界：\(\|\mathbf{f}\| \leq M\)。取 \(\delta \leq \min(R/M, 1/L)\)，则 \(\mathcal{T}\) 将 \(C([t_0-\delta, t_0+\delta], \overline{B}_R(\mathbf{x}_0))\) 映到自身，且为压缩映射： \[\|\mathcal{T}[\mathbf{x}] - \mathcal{T}[\mathbf{y}]\|_\infty \leq L\delta \|\mathbf{x} - \mathbf{y}\|_\infty < \|\mathbf{x} - \mathbf{y}\|_\infty\] 由 Banach 不动点定理（压缩映射原理），\(\mathcal{T}\) 存在唯一不动点 \(\mathbf{x}\)，即 IVP 的唯一解。\(\blacksquare\)

\textbf{定理 89.2}（Peano 存在定理）：若 \(\mathbf{f}\) 仅连续（不要求 Lipschitz），则 IVP 至少有一个解（但不一定唯一）。例如 \(\dot{x} = 2\sqrt{|x|}\)，\(x(0) = 0\) 有无穷多解 \(x(t) = 0\) 和 \(x(t) = t^2 \operatorname{sgn}(t)\) 等。

\textbf{证明}：取 \(\delta > 0\) 使 \(\mathbf{f}\) 在紧集 \([t_0-\delta, t_0+\delta] \times \overline{B}_R(\mathbf{x}_0)\) 上有界。构造 Euler 折线族：将区间 \(n\) 等分，步长 \(h = \delta/n\)，定义折线函数 \(\mathbf{x}_n(t)\) 为分段线性且每步导数等于该步起点的 \(\mathbf{f}\) 值： \[\mathbf{x}_n(t) = \mathbf{x}_n(t_k) + \mathbf{f}(t_k, \mathbf{x}_n(t_k))(t - t_k), \quad t \in [t_k, t_{k+1}]\] 其中 \(t_k = t_0 + kh\)。\(\{\mathbf{x}_n\}\) 一致有界且等度连续（导数一致有界）。由 Ascoli-Arzelà 定理，\(\{\mathbf{x}_n\}\) 有一致收敛子列 \(\mathbf{x}_{n_k} \to \mathbf{x}\)。由 \(\mathbf{f}\) 的一致连续性及积分方程 \(\mathbf{x}_n(t) = \mathbf{x}_0 + \int_{t_0}^t \mathbf{f}(s, \mathbf{x}_n(s)) ds + o(1)\)，取极限得 \(\mathbf{x}\) 满足积分方程，故为解。\(\blacksquare\)

\textbf{定理 89.3}（Gronwall 不等式）：设 \(u(t) \geq 0\) 连续且满足 \(u(t) \leq \alpha + \beta \int_{t_0}^t u(s) ds\)（\(\alpha, \beta \geq 0\)），则 \(u(t) \leq \alpha e^{\beta|t-t_0|}\)。Gronwall 不等式是 ODE 定性理论中最常用的估计工具，其微分形式为：若 \(u' \leq \beta u\)，则 \(u(t) \leq u(t_0) e^{\beta(t-t_0)}\)。

\textbf{证明}：定义 \(v(t) = \alpha + \beta \int_{t_0}^t u(s) ds\)，则 \(v'(t) = \beta u(t) \leq \beta v(t)\)。由 \(\frac{d}{dt}(v(t)e^{-\beta t}) = e^{-\beta t}(v'(t) - \beta v(t)) \leq 0\)，得 \(v(t)e^{-\beta t} \leq v(t_0) = \alpha\)，故 \(u(t) \leq v(t) \leq \alpha e^{\beta(t-t_0)}\)。\(\blacksquare\)

Gronwall 不等式在形式上类似微分方程 \(u' = \beta u\) 的解 \(u = \alpha e^{\beta t}\)------它告诉我们，满足积分不等式 \(u \leq \alpha + \beta \int u\) 的函数 \(u\)，其增长速度不会超过指数函数 \(\alpha e^{\beta t}\)。这一简单而有力的估计立刻可以用于证明解对初始条件的连续依赖性：如果两个解的初值接近，则解本身在有限时间内也接近。

\textbf{定理 89.4}（连续依赖性）：设 \(\mathbf{f}\) 连续且关于 \(\mathbf{x}\) 局部 Lipschitz。则解 \(\mathbf{x}(t; t_0, \mathbf{x}_0)\) 在其存在区间内连续依赖于初始条件 \((t_0, \mathbf{x}_0)\)。

\textbf{证明}：使用 Gronwall 不等式。设 \(\mathbf{x}_1, \mathbf{x}_2\) 是两个不同初值的解，则 \[\|\mathbf{x}_1(t) - \mathbf{x}_2(t)\| \leq \|\mathbf{x}_1(0) - \mathbf{x}_2(0)\| + \int_0^t L \|\mathbf{x}_1(s) - \mathbf{x}_2(s)\| ds\] 由 Gronwall 不等式，\(\|\mathbf{x}_1(t) - \mathbf{x}_2(t)\| \leq \|\mathbf{x}_1(0) - \mathbf{x}_2(0)\| e^{L|t|}\)。故初值接近时解在有限时间内也接近。\(\blacksquare\)

\subsection{89.3 解的延拓与极大解}\label{ux89e3ux7684ux5ef6ux62d3ux4e0eux6781ux5927ux89e3}

\textbf{定义 89.6}（极大解）：IVP 的\textbf{极大解}是定义在极大存在区间 \(I_{\max} = (T_-, T_+)\) 上的解，不能延拓到更大的区间上。\(T_+\) 称为\textbf{前向生命跨度}（forward lifespan），\(T_-\) 称为\textbf{后向生命跨度}。

局部存在性定理（Picard-Lindelöf）保证了解在小范围内的存在性，但一个自然的问题是：这些局部解能否''粘合''成一个定义在最大可能区间上的全局解？极大解的存在性给出了肯定的回答------通过 Zorn 引理或逐步延拓，每个IVP都有唯一的极大解。

\textbf{命题 89.5}（极大解的存在性）：若 \(\mathbf{f}\) 连续且局部 Lipschitz，则每个 IVP 有唯一的极大解。

\textbf{证明}：由 Picard-Lindelöf 定理，在 \((t_0, \mathbf{x}_0)\) 附近存在局部解。定义 \(T_+ = \sup\{T > t_0 : \text{在 } [t_0, T] \text{ 上存在解}\}\)，\(T_- = \inf\{T < t_0 : \text{在 } [T, t_0] \text{ 上存在解}\}\)。取单调序列 \(t_n \uparrow T_+\) 及对应局部解 \(\mathbf{x}^{(n)}\)，由唯一性，各解在公共定义域上一致，故可粘合为 \([t_0, T_+)\) 上的解。该解由构造不能再延拓，从而是极大解。唯一性继承自局部唯一性。\(\blacksquare\)

\textbf{定理 89.6}（解的爆破判别）：若极大解的存在区间 \(I_{\max} = (T_-, T_+)\) 有限（即 \(T_+ < \infty\)），则当 \(t \to T_+^-\) 时，\(\|\mathbf{x}(t)\| \to \infty\) 或 \((t, \mathbf{x}(t))\) 趋于 \(U\) 的边界 \(\partial U\)。换言之，解在有限时间内的终止只能由两种机制导致：解趋于无穷大（爆破）或解离开 \(\mathbf{f}\) 的定义域。

\textbf{证明}：若 \(\|\mathbf{x}(t)\|\) 有界且 \((t, \mathbf{x}(t))\) 不趋于边界，则存在序列 \(t_n \to T_+\) 使 \(\mathbf{x}(t_n) \to \mathbf{x}_*\) 且 \((T_+, \mathbf{x}_*) \in U\)。由局部存在性，可在 \(t = T_+\) 附近延拓解，与极大性矛盾。若 \(\|\mathbf{x}(t)\|\) 无界，设 \(\limsup_{t \to T_+} \|\mathbf{x}(t)\| = \infty\)，则若 \(\|\mathbf{x}(t)\|\) 不发散于无穷，存在有界子列，同样可延拓。\(\blacksquare\)

\subsection{89.4 自治系统与相平面分析}\label{ux81eaux6cbbux7cfbux7edfux4e0eux76f8ux5e73ux9762ux5206ux6790}

\textbf{定义 89.7}（自治系统）：ODE 系统 \(\mathbf{x}' = \mathbf{f}(\mathbf{x})\)（\(\mathbf{f}\) 不显含 \(t\)）称为\textbf{自治系统}。自治系统具有时间平移不变性：若 \(\mathbf{x}(t)\) 为解，则 \(\mathbf{x}(t + c)\) 也是解。自治系统是动力系统理论的基本研究对象，其定性行为由相空间中的向量场 \(\mathbf{f}\) 完全决定。

\textbf{定义 89.8}（相图）：自治系统的\textbf{相图}是 \(\mathbb{R}^n\) 中全体轨道的集合。轨道 \(\{\mathbf{x}(t) : t \in I_{\max}\}\) 是相空间中的曲线。两轨道要么完全重合，要么不相交（由唯一性）。

相图是可视化自治系统行为的核心工具。在所有的相空间结构中，平衡点是最简单的------它们对应于不随时间变化的常数解。尽管看似平凡，平衡点的稳定性分类为理解系统的全局动力学提供了基本框架。

\textbf{定义 89.9}（平衡点）：\(\mathbf{x}_*\) 称为\textbf{平衡点}（或不动点、临界点），如果 \(\mathbf{f}(\mathbf{x}_*) = 0\)。平衡点对应于常数解 \(\mathbf{x}(t) \equiv \mathbf{x}_*\)。平衡点是相图中最简单的结构，其稳定性分类是定性理论的核心内容。

\textbf{定义 89.10}（二维线性系统的分类）：对于线性系统 \(\dot{\mathbf{x}} = A\mathbf{x}\)（\(A \in \mathbb{R}^{2 \times 2}\)），平衡点 \(\mathbf{0}\) 的类型由 \(A\) 的特征值 \(\lambda_1, \lambda_2\) 决定。设 \(\tau = \operatorname{tr}(A) = \lambda_1 + \lambda_2\)，\(\Delta = \det(A) = \lambda_1\lambda_2\)，判别式 \(D = \tau^2 - 4\Delta\)：

\begin{longtable}[]{@{}
  >{\raggedright\arraybackslash}p{(\linewidth - 4\tabcolsep) * \real{0.3929}}
  >{\raggedright\arraybackslash}p{(\linewidth - 4\tabcolsep) * \real{0.2143}}
  >{\raggedright\arraybackslash}p{(\linewidth - 4\tabcolsep) * \real{0.3929}}@{}}
\toprule
\begin{minipage}[b]{\linewidth}\raggedright
特征值类型
\end{minipage} & \begin{minipage}[b]{\linewidth}\raggedright
条件
\end{minipage} & \begin{minipage}[b]{\linewidth}\raggedright
平衡点类型
\end{minipage} \\
\midrule
\endhead
\bottomrule
\endlastfoot
\(\lambda_1 < \lambda_2 < 0\) & \(\tau < 0, \Delta > 0, D > 0\) & 稳定结点 \\
\(0 < \lambda_1 < \lambda_2\) & \(\tau > 0, \Delta > 0, D > 0\) & 不稳定结点 \\
\(\lambda_1 < 0 < \lambda_2\) & \(\Delta < 0\) & 鞍点 \\
\(\lambda = \alpha \pm i\beta\)，\(\alpha < 0\) & \(\tau < 0, D < 0\) & 稳定焦点 \\
\(\lambda = \alpha \pm i\beta\)，\(\alpha > 0\) & \(\tau > 0, D < 0\) & 不稳定焦点 \\
\(\lambda = \pm i\beta\)（纯虚数） & \(\tau = 0, \Delta > 0\) & 中心 \\
\end{longtable}

二维线性系统的分类表穷举了所有可能的情况，为相平面分析提供了完整的''词典''。然而，现实中的绝大多数ODE都是非线性的------那么，能否将线性分类的结果推广到非线性系统？以下命题给出了系统性的相面分析方法，其核心思想是：在平衡点附近，非线性的局部行为''近似''于其线性化系统。

\textbf{命题 89.6}（相平面分析的关键步骤）：对非线性自治系统 \(\dot{\mathbf{x}} = \mathbf{f}(\mathbf{x})\)： 1. \textbf{求平衡点}：解 \(\mathbf{f}(\mathbf{x}) = \mathbf{0}\)。 2. \textbf{线性化}：在每个平衡点处计算 Jacobi 矩阵 \(D\mathbf{f}\)，求特征值并分类（结点、鞍点、焦点、中心）。 3. \textbf{双曲平衡点的局部行为}：若 \(\Re(\lambda) \neq 0\)（双曲平衡点），由 Hartman-Grobman 定理，局部相图拓扑等价于线性化系统。双曲平衡点是结构稳定的。 4. \textbf{零实部特征值的中心流形}：若存在纯虚数特征值（或零特征值），线性化不能确定稳定性，需在中心流形上约化分析非线性项。中心流形定理将系统约化为低维系统，在此约化系统上分析非线性项的影响。 5. \textbf{全局相图}：绘制零倾线（nullclines，即 \(\dot{x}_1 = 0\) 和 \(\dot{x}_2 = 0\) 的曲线），确定各区域中的向量场方向，用 Poincaré-Bendixson 定理判断周期轨的存在性。

\subsection{89.5 比较定理与微分不等式}\label{ux6bd4ux8f83ux5b9aux7406ux4e0eux5faeux5206ux4e0dux7b49ux5f0f}

\textbf{定理 89.7}（比较定理）：设 \(u, v\) 为 \(C^1\) 函数满足 \(u'(t) \leq f(t, u(t))\)，\(v'(t) = f(t, v(t))\) 且 \(u(t_0) \leq v(t_0)\)。若 \(f\) 关于第二变量单调不减，则 \(u(t) \leq v(t)\) 对所有 \(t \geq t_0\) 成立。比较定理是 ODE 上下解方法和单调迭代方法的基础。

\textbf{证明}：设 \(w(t) = v(t) - u(t)\)，则 \(w(t_0) \geq 0\)。若存在 \(t_1 > t_0\) 使 \(w(t_1) < 0\)，令 \(t_* = \inf\{t > t_0 : w(t) < 0\}\)，则 \(w(t_*) = 0\) 且 \(w'(t_*) \leq 0\)。但 \(w'(t_*) = f(t_*, v(t_*)) - u'(t_*) \geq f(t_*, u(t_*)) - f(t_*, u(t_*)) = 0\)（由 \(f\) 的单调性），矛盾。\(\blacksquare\)

\hrulefill

\section{第312章：线性ODE系统}\label{ux7b2c312ux7ae0ux7ebfux6027odeux7cfbux7edf}

线性 ODE 系统可借助矩阵指数和线性代数（V7）的工具完全求解。本章重点讨论常系数系统和稳定性理论。

\subsection{90.1 矩阵指数}\label{ux77e9ux9635ux6307ux6570}

\textbf{定义 90.1}（矩阵指数）：对 \(A \in M_n(\mathbb{C})\)，定义 \[e^{tA} = \sum_{k=0}^\infty \frac{t^k A^k}{k!}\] 此级数对任意 \(t \in \mathbb{R}\) 绝对收敛（在任意矩阵范数下）。矩阵指数是标量指数函数在矩阵值上的自然推广，继承了指数函数的许多重要性质。

\textbf{命题 90.1}（矩阵指数的性质）： 1. \(e^{0A} = I\) 2. \(\frac{d}{dt} e^{tA} = A e^{tA} = e^{tA} A\) 3. 若 \(AB = BA\)，则 \(e^{A+B} = e^A e^B\)（注意：对非交换矩阵，\(e^{A+B} \neq e^A e^B\) 一般情形成立，其差由 Baker-Campbell-Hausdorff 公式给出） 4. \(e^{tA}\) 可逆，\((e^{tA})^{-1} = e^{-tA}\)，故 \(\det(e^{tA}) = e^{t \operatorname{tr}(A)}\) 5. 若 \(A = PJP^{-1}\)（\(J\) 为 Jordan 标准形），则 \(e^{tA} = P e^{tJ} P^{-1}\)

\textbf{矩阵指数的计算}：若 \(A\) 可对角化（\(A = PDP^{-1}\)，\(D = \operatorname{diag}(\lambda_1, \ldots, \lambda_n)\)），则 \(e^{tA} = P \operatorname{diag}(e^{\lambda_1 t}, \ldots, e^{\lambda_n t}) P^{-1}\)。若 \(A\) 有 Jordan 块 \(J_k(\lambda)\)（大小为 \(k\)），则 \[e^{t J_k(\lambda)} = e^{\lambda t} \begin{pmatrix} 1 & t & t^2/2! & \cdots & t^{k-1}/(k-1)! \\ 0 & 1 & t & \cdots & t^{k-2}/(k-2)! \\ \vdots & \vdots & \vdots & \ddots & \vdots \\ 0 & 0 & 0 & \cdots & 1 \end{pmatrix}\]

\textbf{定理 90.2}（线性系统的解）：IVP \(\dot{\mathbf{x}} = A\mathbf{x}\)，\(\mathbf{x}(0) = \mathbf{x}_0\) 的唯一解为 \[\mathbf{x}(t) = e^{tA} \mathbf{x}_0\] 非齐次系统 \(\dot{\mathbf{x}} = A\mathbf{x} + \mathbf{b}(t)\) 的解由常数变易公式给出： \[\mathbf{x}(t) = e^{tA} \mathbf{x}_0 + \int_0^t e^{(t-s)A} \mathbf{b}(s) ds\]

\textbf{证明}：\(\frac{d}{dt} e^{tA} \mathbf{x}_0 = A e^{tA} \mathbf{x}_0 = A\mathbf{x}(t)\)，且 \(\mathbf{x}(0) = \mathbf{x}_0\)。唯一性由 Picard-Lindelöf 定理保证。非齐次公式：设 \(\mathbf{x}(t) = e^{tA} \mathbf{c}(t)\)，代入得 \(\mathbf{c}'(t) = e^{-tA} \mathbf{b}(t)\)，积分即得。\(\blacksquare\)

矩阵指数将线性系统的求解转化为矩阵运算，给出了解的封闭表达式。但求解方程只是第一步------在实际应用中，我们更关心解的长期行为：平衡点是否稳定？微小扰动是否会使系统偏离平衡态？这些问题构成了稳定性理论的核心。

\subsection{90.2 常系数线性系统的稳定性}\label{ux5e38ux7cfbux6570ux7ebfux6027ux7cfbux7edfux7684ux7a33ux5b9aux6027}

\textbf{定义 90.2}（Lyapunov 稳定性）：设 \(\mathbf{x}_*\) 是 \(\dot{\mathbf{x}} = \mathbf{f}(\mathbf{x})\) 的平衡点。 - \(\mathbf{x}_*\) 是\textbf{稳定}的（Lyapunov 意义下），如果对任意 \(\varepsilon > 0\)，存在 \(\delta > 0\) 使得 \(\|\mathbf{x}(0) - \mathbf{x}_*\| < \delta\) 推出 \(\|\mathbf{x}(t) - \mathbf{x}_*\| < \varepsilon\)（对所有 \(t \geq 0\)） - \(\mathbf{x}_*\) 是\textbf{渐近稳定}的，如果它稳定且 \(\lim_{t \to \infty} \mathbf{x}(t) = \mathbf{x}_*\) - \(\mathbf{x}_*\) 是\textbf{指数稳定}的，如果存在 \(C, \alpha > 0\) 使得 \(\|\mathbf{x}(t) - \mathbf{x}_*\| \leq C e^{-\alpha t} \|\mathbf{x}(0) - \mathbf{x}_*\|\) - \(\mathbf{x}_*\) 是\textbf{不稳定}的，如果它不是稳定的。

Lyapunov 稳定性精确定义了''小扰动不会导致大偏离''这一直觉概念。对于线性系统，稳定性的判断简化为了矩阵特征值的分析------这是线性代数在 ODE 理论中最优美的应用之一。

\textbf{定理 90.3}（线性稳定性）：对线性系统 \(\dot{\mathbf{x}} = A\mathbf{x}\)： - \(\mathbf{0}\) 渐近稳定（实际上指数稳定）当且仅当 \(A\) 的所有特征值满足 \(\Re(\lambda) < 0\) - \(\mathbf{0}\) 稳定当且仅当 \(A\) 的所有特征值满足 \(\Re(\lambda) \leq 0\)，且满足 \(\Re(\lambda) = 0\) 的特征值对应大小为 \(1 \times 1\) 的 Jordan 块（即该特征值的代数重数等于几何重数） - \(\mathbf{0}\) 不稳定当且仅当存在特征值 \(\Re(\lambda) > 0\)（或 \(\Re(\lambda) = 0\) 但 Jordan 块大小 \textgreater{} 1）

\textbf{证明}：将 \(A\) 写为 Jordan 标准形 \(J = P^{-1}AP\)。解为 \(\mathbf{x}(t) = P e^{tJ} P^{-1} \mathbf{x}_0\)，其中 \(e^{tJ}\) 的每个 Jordan 块贡献 \(t^k e^{\lambda t}\) 型项（\(k\) 为块大小减一）。若所有 \(\Re(\lambda) < 0\)，则 \(|t^k e^{\lambda t}| \leq C e^{-\alpha t}\)（\(\alpha = \min|\Re(\lambda)|\)），指数稳定。若某 \(\Re(\lambda) = 0\) 且块大小为 1，对应 \(e^{i\omega t}\) 有界，解稳定但不渐近稳定。若某 \(\Re(\lambda) = 0\) 且块大小 \textgreater{} 1，出现 \(t e^{i\omega t}\) 无界项，解不稳定。若某 \(\Re(\lambda) > 0\)，则 \(e^{\Re(\lambda)t}\) 指数增长，解不稳定。\(\blacksquare\)

线性稳定性定理完美地解决了线性系统平衡点的分类问题。但对于非线性系统，我们只能寄望于在平衡点附近''线性化''------即用 Jacobi 矩阵近似原系统。Hartman-Grobman 定理是这一思路的巅峰成果：它断言对于双曲平衡点（所有特征值的实部均不为零），非线性系统的局部动力学拓扑等价于其线性化系统。

\textbf{定理 90.4}（线性化稳定性，Hartman-Grobman 定理）：设 \(\mathbf{x}_*\) 是 \(\dot{\mathbf{x}} = \mathbf{f}(\mathbf{x})\) 的平衡点，\(\mathbf{f}\) 是 \(C^1\) 的。若 \(A = D\mathbf{f}(\mathbf{x}_*)\) 的所有特征值满足 \(\Re(\lambda) \neq 0\)（双曲平衡点），则存在 \(\mathbf{x}_*\) 的邻域 \(U\) 和 \(\mathbf{0}\) 的邻域 \(V\) 以及同胚 \(h: U \to V\)，使得 \(h\) 将非线性系统 \(\dot{\mathbf{x}} = \mathbf{f}(\mathbf{x})\) 的轨道映为线性化系统 \(\dot{\mathbf{y}} = A\mathbf{y}\) 的轨道，且保持时间参数化方向。即非线性系统在双曲平衡点附近拓扑共轭于其线性化系统。

\textbf{证明概要}：Hartman-Grobman 定理的证明（P. Hartman 1960, 简化版）的关键在于构造共轭同胚。设 \(\phi_t\) 为非线性流，\(\psi_t = e^{tA}\) 为线性流。寻找同胚 \(h\) 满足 \(\phi_t \circ h = h \circ \psi_t\)。将非线性系统写为 \(\dot{\mathbf{x}} = A\mathbf{x} + \mathbf{g}(\mathbf{x})\)（\(\mathbf{g}(\mathbf{0}) = \mathbf{0}\)，\(D\mathbf{g}(\mathbf{0}) = \mathbf{0}\)）。利用双曲性，将 \(\mathbb{R}^n\) 分解为稳定子空间 \(E^s\)（\(\Re(\lambda) < 0\)）和不稳定子空间 \(E^u\)（\(\Re(\lambda) > 0\)）的直和。在 \(E^s\) 上，非线性流 \(\phi_t\) 指数收缩；在 \(E^u\) 上，\(\phi_{-t}\) 指数收缩。通过解泛函方程 \(h(\mathbf{x}) = \mathbf{x} + \int_0^\infty e^{-sA} \mathbf{g}(\phi_s \circ h \circ \psi_{-s}(\mathbf{x})) ds\)（在 \(E^s\) 上）和类似公式（在 \(E^u\) 上），由压缩映射原理得到 \(h\) 的存在性和唯一性。\(\blacksquare\)

\subsection{90.3 稳定流形定理}\label{ux7a33ux5b9aux6d41ux5f62ux5b9aux7406}

\textbf{定理 90.5}（稳定流形定理）：设 \(\mathbf{x}_*\) 为 \(\dot{\mathbf{x}} = \mathbf{f}(\mathbf{x})\)（\(\mathbf{f} \in C^k\)）的双曲平衡点，\(A = D\mathbf{f}(\mathbf{x}_*)\)。则存在 \(C^k\) 流形： - \textbf{局部稳定流形} \(W^s_{\text{loc}}(\mathbf{x}_*) = \{\mathbf{x} : \phi_t(\mathbf{x}) \to \mathbf{x}_* \text{ 当 } t \to +\infty \text{ 且 } \phi_t(\mathbf{x}) \in U \text{ 对 } t \geq 0\}\)，其在 \(\mathbf{x}_*\) 处的切空间为 \(E^s\)（\(A\) 的稳定子空间）； - \textbf{局部不稳定流形} \(W^u_{\text{loc}}(\mathbf{x}_*) = \{\mathbf{x} : \phi_t(\mathbf{x}) \to \mathbf{x}_* \text{ 当 } t \to -\infty \text{ 且 } \phi_t(\mathbf{x}) \in U \text{ 对 } t \leq 0\}\)，其在 \(\mathbf{x}_*\) 处的切空间为 \(E^u\)。

稳定流形定理将 Hartman-Grobman 的拓扑等价提升为微分同胚层次的等价，是动力系统双曲理论的核心结果。

\subsection{90.4 Lyapunov 函数}\label{lyapunov-ux51fdux6570}

\textbf{定义 90.3}（Lyapunov 函数）：设 \(\mathbf{x}_*\) 是平衡点。\(C^1\) 函数 \(V: \mathcal{U} \to \mathbb{R}\)（\(\mathcal{U}\) 是 \(\mathbf{x}_*\) 的邻域）称为 \textbf{Lyapunov 函数}，如果 \(V(\mathbf{x}_*) = 0\)，\(V(\mathbf{x}) > 0\)（\(\mathbf{x} \neq \mathbf{x}_*\)），且沿轨道 \(\dot{V}(\mathbf{x}) = \nabla V(\mathbf{x}) \cdot \mathbf{f}(\mathbf{x}) \leq 0\)。

\textbf{定理 90.6}（Lyapunov 稳定性定理）： 1. 若存在 Lyapunov 函数，则 \(\mathbf{x}_*\) 是稳定的 2. 若进一步 \(\dot{V}(\mathbf{x}) < 0\)（\(\mathbf{x} \neq \mathbf{x}_*\)），则 \(\mathbf{x}_*\) 是渐近稳定的 3. （Krasovskii-LaSalle 不变原理）若 \(\dot{V} \leq 0\)，则所有轨道趋于 \(\{\dot{V} = 0\}\) 中的最大不变集

\textbf{证明}：(1) 由 \(V\) 沿轨道递减，对任意 \(\varepsilon > 0\)，取 \(m = \min_{\|\mathbf{x}-\mathbf{x}_*\|=\varepsilon} V(\mathbf{x}) > 0\)。由 \(V\) 的连续性，存在 \(\delta > 0\) 使 \(\|\mathbf{x}(0)-\mathbf{x}_*\| < \delta\) 时 \(V(\mathbf{x}(0)) < m\)。因 \(V(\mathbf{x}(t))\) 不增，\(V(\mathbf{x}(t)) < m\) 对所有 \(t \geq 0\)，故 \(\|\mathbf{x}(t)-\mathbf{x}_*\| < \varepsilon\)。(2) 若 \(\dot{V} < 0\)，则 \(V(\mathbf{x}(t))\) 严格递减且有下界 \(0\)，故 \(V(\mathbf{x}(t)) \to L \geq 0\)。若 \(L > 0\)，则轨道在 \(\{V \geq L\}\) 上，\(\dot{V} \leq -\delta < 0\)，导出 \(V\) 趋于 \(-\infty\)，矛盾。故 \(L = 0\)，即 \(\mathbf{x}(t) \to \mathbf{x}_*\)。\(\blacksquare\)

\subsection{90.5 Floquet 理论与周期系数系统}\label{floquet-ux7406ux8bbaux4e0eux5468ux671fux7cfbux6570ux7cfbux7edf}

\textbf{定理 90.7}（Floquet 定理，1883）：设 \(\dot{\mathbf{x}} = A(t)\mathbf{x}\) 为 \(T\)-周期线性系统（\(A(t+T) = A(t)\)）。则基本解矩阵 \(\Phi(t)\)（\(\Phi(0) = I\)）可写为 \[\Phi(t) = P(t) e^{tB}\] 其中 \(P(t)\) 为 \(T\)-周期非奇异矩阵（\(P(t+T) = P(t)\)），\(B\) 为常数矩阵。\(e^{TB}\) 的特征值 \(\lambda_i\) 称为系统的\textbf{特征乘子}（characteristic multipliers），\(B\) 的特征值 \(\mu_i = \frac{1}{T}\ln\lambda_i\) 称为\textbf{特征指数}（Floquet exponents）。稳定性由特征乘子决定：\(|\lambda_i| < 1\)（所有 \(i\)）时零解渐近稳定；存在 \(|\lambda_i| > 1\) 时零解不稳定。

\textbf{证明}：定义 Monodromy 矩阵 \(M = \Phi(T)\)。由 \(\Phi(t+T) = \Phi(t)M\)，取 \(B = \frac{1}{T}\ln M\)，\(P(t) = \Phi(t)e^{-tB}\)，则 \(P(t+T) = \Phi(t+T)e^{-(t+T)B} = \Phi(t)M e^{-TB} e^{-tB} = \Phi(t) e^{-tB} = P(t)\)。\(\blacksquare\)

\hrulefill

\emph{常微分方程的理论体系从 Picard-Lindelöf 和 Peano 定理的局部存在唯一性出发，经极大解的延拓和爆破判别，建立了全局理论框架。线性系统以矩阵指数和 Jordan 标准形为代数工具，给出了解的完全分类。稳定性理论以 Lyapunov 方法和 Hartman-Grobman 定理为双翼，将非线性系统的局部行为线性化。稳定流形定理、Floquet 理论和比较定理则进一步丰富了 ODE 定性分析的武器库。}

\emph{卷二十一：微分方程。}

\newpage

\chapter{09-微分方程与动力系统/01-微分方程/02-偏微分方程.md}\label{ux5faeux5206ux65b9ux7a0bux4e0eux52a8ux529bux7cfbux7edf01-ux5faeux5206ux65b9ux7a0b02-ux504fux5faeux5206ux65b9ux7a0b.md}

\section{第313章：偏微分方程基础}\label{ux7b2c313ux7ae0ux504fux5faeux5206ux65b9ux7a0bux57faux7840}

偏微分方程涉及多个自变量，其理论比 ODE 复杂得多。本章从一阶 PDE 和分类入手，引入特征线法。

\subsection{91.1 基本概念与分类}\label{ux57faux672cux6982ux5ff5ux4e0eux5206ux7c7b}

\textbf{定义 91.1}（PDE）：\textbf{偏微分方程}（PDE）是包含未知函数 \(u(x_1, \ldots, x_n)\) 及其偏导数的方程。\textbf{阶数}是出现的最高阶偏导数。

\textbf{定义 91.2}（二阶线性 PDE）：一般的二阶线性 PDE 为

\[\sum_{i,j=1}^n a_{ij}(x) \frac{\partial^2 u}{\partial x_i \partial x_j} + \sum_{i=1}^n b_i(x) \frac{\partial u}{\partial x_i} + c(x) u = f(x)\]

其中 \(a_{ij} = a_{ji}\)。

\textbf{定义 91.3}（二阶 PDE 的分类）：在点 \(x\) 处，根据系数矩阵 \(A = (a_{ij}(x))\) 的特征值符号分类： - \textbf{椭圆型}：\(A\) 的所有特征值同号（或 \(A\) 正定/负定）。例：Laplace 方程 \(\Delta u = 0\) - \textbf{抛物型}：\(A\) 有一个零特征值，其余同号。例：热方程 \(u_t = \Delta u\) - \textbf{双曲型}：\(A\) 有一个特征值与其余符号相反。例：波动方程 \(u_{tt} = \Delta u\)

\subsection{91.2 一阶 PDE 与特征线法}\label{ux4e00ux9636-pde-ux4e0eux7279ux5f81ux7ebfux6cd5}

\textbf{定义 91.4}（一阶拟线性 PDE）：一阶拟线性 PDE 的形式为

\[a(x, y, u) \frac{\partial u}{\partial x} + b(x, y, u) \frac{\partial u}{\partial y} = c(x, y, u)\]

\textbf{定理 91.1}（特征线法）：上述 PDE 的解曲面由特征线 ODE 系统生成：

\[\frac{dx}{dt} = a(x, y, u), \quad \frac{dy}{dt} = b(x, y, u), \quad \frac{du}{dt} = c(x, y, u)\]

其中 \(u(t) = u(x(t), y(t))\)。沿特征线，PDE 约化为 ODE。

\emph{证明}：由链式法则，\(\frac{du}{dt} = \frac{\partial u}{\partial x} \frac{dx}{dt} + \frac{\partial u}{\partial y} \frac{dy}{dt} = a \frac{\partial u}{\partial x} + b \frac{\partial u}{\partial y} = c\)。∎

\subsection{91.3 适定性}\label{ux9002ux5b9aux6027}

\textbf{定义 91.5}（适定性，Hadamard）：PDE 问题称为\textbf{适定的}，如果 1. 解存在 2. 解唯一 3. 解连续依赖于数据（初始条件、边界条件、方程系数）

典型的适定问题： - \textbf{椭圆型}：Dirichlet 问题（给定边界值）或 Neumann 问题（给定法向导数） - \textbf{抛物型}：初边值问题（给定初始条件和边界条件） - \textbf{双曲型}：初值问题（Cauchy 问题，给定初始位置和速度）

\subsection{91.4 分离变量法}\label{ux5206ux79bbux53d8ux91cfux6cd5}

分离变量法是求解线性齐次 PDE 初边值问题的经典方法，其核心思想是将多变量函数的解分解为单变量函数的乘积，从而将 PDE 约化为若干 ODE。

\textbf{一般步骤}： 1. \textbf{分离假设}：设解具有乘积形式 \(u(x, t) = X(x) T(t)\)（以两变量为例），代入齐次 PDE。 2. \textbf{变量分离}：将方程整理为''仅含 \(x\) 的函数 = 仅含 \(t\) 的函数''的形式，由于两边独立变量不同，必须等于同一常数 \(-\lambda\)（分离常数）。 3. \textbf{化为 ODE}：得到关于 \(X(x)\) 和 \(T(t)\) 的两个常微分方程 \(X'' + \lambda X = 0\) 和 \(T' + \lambda k T = 0\)（具体形式取决于 PDE 类型），并配以相应的齐次边界条件。 4. \textbf{求解特征值问题}：由 \(X(x)\) 的 ODE 与齐次边界条件构成 Sturm-Liouville 特征值问题，求出特征值 \(\lambda_n\) 和特征函数 \(X_n(x)\)。 5. \textbf{叠加原理}：通解为所有分离解的线性叠加 \(u(x, t) = \sum_n c_n X_n(x) T_n(t)\)，系数 \(c_n\) 由初始条件通过特征函数的正交性（Fourier 级数展开）确定。

该方法成功的关键在于 PDE 和边界条件均为线性齐次的，从而叠加原理成立。对于非齐次情形，可结合特征函数展开法或 Duhamel 原理处理。

\hrulefill

\hrulefill

\hrulefill

\hrulefill

\hrulefill

\hrulefill

\section{第314章：椭圆、抛物与双曲型方程}\label{ux7b2c314ux7ae0ux692dux5706ux629bux7269ux4e0eux53ccux66f2ux578bux65b9ux7a0b}

本章具体讨论三大经典 PDE：Laplace 方程、热方程和波动方程，给出基本解和能量方法。

\subsection{92.1 Laplace 方程与 Poisson 方程}\label{laplace-ux65b9ux7a0bux4e0e-poisson-ux65b9ux7a0b}

\textbf{定义 92.1}（Laplace 算子）：\(\mathbb{R}^n\) 中的 \textbf{Laplace 算子}（Laplacian）为 \(\Delta = \sum_{i=1}^n \frac{\partial^2}{\partial x_i^2}\)。

\textbf{定义 92.2}（调和函数）：\(C^2\) 函数 \(u\) 称为\textbf{调和函数}，如果 \(\Delta u = 0\)。

\textbf{定理 92.1}（基本解）：\(\mathbb{R}^n\) 中 Laplace 方程的基本解为

\[\Phi(x) = \begin{cases} -\frac{1}{2\pi} \log |x| & n = 2 \\ \frac{1}{n(n-2)\omega_n} |x|^{2-n} & n \geq 3 \end{cases}\]

其中 \(\omega_n\) 是 \(\mathbb{R}^n\) 中单位球的体积。\(\Phi\) 满足 \(\Delta \Phi = \delta_0\)（在分布意义下）。

\emph{证明}：当 \(x \neq 0\) 时，直接计算 \(\Delta \Phi(x) = 0\)（利用径向函数的 Laplace 算子公式 \(\Delta f(r) = f''(r) + \frac{n-1}{r} f'(r)\)）。对任意 \(\varphi \in C_c^\infty(\mathbb{R}^n)\)，由散度定理和 \(\Phi\) 在原点附近的奇性分析，\(\int_{\mathbb{R}^n} \Phi \Delta\varphi \, dx = \lim_{\varepsilon \to 0} \int_{|x| > \varepsilon} \Phi \Delta\varphi \, dx = \varphi(0)\)，即 \(\Delta \Phi = \delta_0\) 在分布意义下成立。\(\blacksquare\)

\textbf{定理 92.2}（调和函数的均值性质）：若 \(u\) 在 \(B(x, r) \subset \mathbb{R}^n\) 中调和，则

\[u(x) = \frac{1}{|\partial B(x, r)|} \int_{\partial B(x, r)} u(y) dS(y) = \frac{1}{|B(x, r)|} \int_{B(x, r)} u(y) dy\]

\emph{证明}：固定 \(x\)，令 \(\phi(r) = \frac{1}{|\partial B(x, r)|} \int_{\partial B(x, r)} u(y) dS(y)\) 为球面平均值。做变量代换 \(y = x + r\omega\)（\(\omega \in S^{n-1}\)），则 \(\phi(r) = \frac{1}{|S^{n-1}|} \int_{S^{n-1}} u(x + r\omega) d\omega\)。求导得 \(\phi'(r) = \frac{1}{|S^{n-1}|} \int_{S^{n-1}} \nabla u(x + r\omega) \cdot \omega \, d\omega = \frac{1}{|\partial B(x, r)|} \int_{\partial B(x, r)} \frac{\partial u}{\partial n}(y) dS(y)\)。由散度定理，\(\int_{\partial B(x, r)} \frac{\partial u}{\partial n} dS = \int_{B(x, r)} \Delta u \, dy = 0\)（因为 \(u\) 调和）。故 \(\phi'(r) = 0\)，\(\phi(r)\) 为常数。由 \(u\) 的连续性，\(\lim_{r \to 0} \phi(r) = u(x)\)，故 \(\phi(r) = u(x)\) 对所有 \(r\) 成立。球体平均值公式由积分 \(\frac{1}{|B(x, r)|} \int_{B(x, r)} u(y) dy = \frac{1}{|B(x, r)|} \int_0^r \phi(s) |\partial B(x, s)| ds = u(x)\) 得出。\(\blacksquare\)

\textbf{推论 92.3}（极大值原理）：调和函数在区域内部不能取到严格极大值或极小值（除非为常数）。\(\max_{\overline{\Omega}} |u| = \max_{\partial\Omega} |u|\)。

\emph{证明}：若 \(u\) 在 \(x_0 \in \Omega\) 处取最大值 \(M\)，由均值性质，对任意 \(B(x_0, r) \subset \Omega\)，\(M = u(x_0) = \frac{1}{|B(x_0, r)|} \int_{B(x_0, r)} u(y) dy \leq M\)，等号成立当且仅当 \(u \equiv M\) 在球内。由连通性，\(u\) 在整个 \(\Omega\) 上恒为 \(M\)。若非常数，则最大值必在边界达到。对 \(|u|\) 同理。\(\blacksquare\)

\textbf{定理 92.4}（Poisson 方程的解）：Poisson 方程 \(-\Delta u = f\) 在 \(\mathbb{R}^n\) 中的一个解为

\[u(x) = (\Phi * f)(x) = \int_{\mathbb{R}^n} \Phi(x - y) f(y) dy\]

\emph{证明}：由基本解 \(\Delta \Phi = \delta_0\)（在分布意义下成立）。对 \(f\) 具有紧支集（或适当衰减），定义 \(u = \Phi * f\)。在分布意义下，\(\Delta u = \Delta(\Phi * f) = (\Delta \Phi) * f = \delta_0 * f = f\)。因此 \(-\Delta u = -f\)，即 \(u\) 满足 Poisson 方程。卷积的合法性：\(\Phi\) 在原点附近的行为为 \(|\Phi(x)| \sim |x|^{2-n}\)（\(n \geq 3\)）或 \(\log|x|\)（\(n=2\)），故 \(\Phi \in L^1_{\text{loc}}(\mathbb{R}^n)\)。若 \(f \in C_c^\infty(\mathbb{R}^n)\)，则 \(u = \Phi * f \in C^\infty\)（因为 \(\Phi\) 在支集外光滑），且 \(\Delta u = f\) 在经典意义下成立。若 \(f \in L^p\)，则 \(u\) 为分布解，且由 Calderon-Zygmund 理论，\(\|D^2 u\|_{L^p} \leq C \|f\|_{L^p}\)（\(1 < p < \infty\)）。\(\blacksquare\)

\subsection{92.2 热方程}\label{ux70edux65b9ux7a0b}

\textbf{定义 92.3}（热方程）：热方程（扩散方程）为 \(u_t - \Delta u = 0\)。

\textbf{定义 92.4}（热核 / 基本解）：\(\mathbb{R}^n\) 中热方程的基本解（热核）为

\[\Phi(x, t) = \frac{1}{(4\pi t)^{n/2}} e^{-|x|^2 / (4t)} \quad (t > 0)\]

\(\Phi\) 满足 \(\Phi_t - \Delta \Phi = 0\)（\(t > 0\)）且 \(\lim_{t \to 0^+} \Phi(\cdot, t) = \delta_0\)。

\textbf{定理 92.5}（初值问题的解）：\(\mathbb{R}^n\) 中热方程初值问题 \(u_t = \Delta u\)，\(u(x, 0) = g(x)\) 的解为

\[u(x, t) = (\Phi(\cdot, t) * g)(x) = \frac{1}{(4\pi t)^{n/2}} \int_{\mathbb{R}^n} e^{-|x-y|^2/(4t)} g(y) dy\]

\emph{证明}：对 \(t > 0\)，热核 \(\Phi(x, t) = (4\pi t)^{-n/2} e^{-|x|^2/(4t)}\) 满足热方程。直接验证：\(\Phi_t = \Phi \cdot (-n/(2t) + |x|^2/(4t^2))\)，又 \(\nabla \Phi = \Phi \cdot (-x/(2t))\)，\(\Delta \Phi = \nabla \cdot (\Phi \cdot (-x/(2t))) = -\frac{1}{2t}(n\Phi + x \cdot \nabla \Phi) = -\frac{1}{2t}(n\Phi - \frac{|x|^2}{2t}\Phi) = \Phi \cdot (-\frac{n}{2t} + \frac{|x|^2}{4t^2})\)。故 \(\Phi_t = \Delta \Phi\)。现在定义 \(u(x,t) = (\Phi(\cdot, t) * g)(x)\)。由卷积与微分的交换性，\(u_t - \Delta u = (\Phi_t - \Delta \Phi) * g = 0\)（\(t > 0\)）。初值条件：\(\Phi(\cdot, t)\) 是恒等逼近核，满足 \(\int_{\mathbb{R}^n} \Phi(x, t) dx = 1\) 且 \(\lim_{t \to 0^+} \Phi(\cdot, t) = \delta_0\)（在分布意义下）。若 \(g\) 连续有界，则 \(\lim_{t \to 0^+} u(x, t) = g(x)\) 对每个 \(x\) 成立，且在有界集上一致收敛。若 \(g \in L^p\)，则 \(u(\cdot, t) \to g\) 在 \(L^p\) 中成立。\(\blacksquare\)

\textbf{定理 92.6}（无穷传播速度与极大值原理）：热方程的解具有无穷传播速度（初始扰动瞬间影响全空间）。对热方程，有极大值原理：

\[\max_{\overline{\Omega} \times [0, T]} u = \max\{\max_{\overline{\Omega}} u(\cdot, 0), \max_{\partial\Omega \times [0, T]} u\}\]

即最大值在初始时刻或边界上达到（抛物型极大值原理）。

\emph{证明}：设 \(u\) 在 \(\Omega \times (0, T]\) 上满足 \(u_t - \Delta u = 0\)。若 \(u\) 在内部点 \((x_0, t_0)\)（\(t_0 > 0\)）达到最大值 \(M\)，则在该点处 \(u_t \geq 0\)（若 \(t_0 < T\)，由单侧导数 \(u_t \geq 0\)；若 \(t_0 = T\)，则需考虑 \(u_t\) 在 \((x_0, T)\) 处的上极限）且 \(\Delta u \leq 0\)（由二阶导数的极值条件：Hessian 矩阵半负定，故迹 \(\Delta u \leq 0\)）。代入方程得 \(u_t - \Delta u \geq 0\)。但若 \(u\) 在内部严格取最大值，则 \(u_t = 0\)（若为内点）且 \(\Delta u \leq 0\)，这迫使 \(\Delta u = 0\) 且 \(u_t = 0\)。由强极值原理（Nirenberg 1953），若 \(u\) 在抛物边界内非常数，则内部不能取最大值。因此最大值必在抛物边界上达到。无穷传播速度：热核 \(\Phi(x,t) = (4\pi t)^{-n/2} e^{-|x|^2/(4t)} > 0\) 对所有 \(t > 0\) 和所有 \(x \in \mathbb{R}^n\) 严格成立。故若初始数据 \(g \geq 0\) 且不恒为零，则 \(u(x,t) = (\Phi * g)(x) > 0\) 对所有 \(t > 0\) 和所有 \(x\) 成立，扰动以无穷速度传播。\(\blacksquare\)

\subsection{92.3 波动方程}\label{ux6ce2ux52a8ux65b9ux7a0b}

\textbf{定义 92.5}（波动方程）：波动方程为 \(u_{tt} - c^2 \Delta u = 0\)（\(c > 0\) 是波速）。

\textbf{定理 92.7}（d'Alembert 公式，一维）：一维波动方程 \(u_{tt} = c^2 u_{xx}\)，\(u(x, 0) = g(x)\)，\(u_t(x, 0) = h(x)\) 的解为

\[u(x, t) = \frac{1}{2}[g(x + ct) + g(x - ct)] + \frac{1}{2c} \int_{x-ct}^{x+ct} h(s) ds\]

\emph{证明}：引入特征坐标 \(\xi = x + ct\)，\(\eta = x - ct\)。由链式法则：\(\partial_x = \partial_\xi + \partial_\eta\)，\(\partial_t = c\partial_\xi - c\partial_\eta\)。再求二阶导：\(\partial_x^2 = (\partial_\xi + \partial_\eta)^2 = \partial_\xi^2 + 2\partial_\xi\partial_\eta + \partial_\eta^2\)，\(\partial_t^2 = c^2(\partial_\xi - \partial_\eta)^2 = c^2(\partial_\xi^2 - 2\partial_\xi\partial_\eta + \partial_\eta^2)\)。代入方程 \(u_{tt} = c^2 u_{xx}\) 得 \(c^2(\partial_\xi^2 - 2\partial_\xi\partial_\eta + \partial_\eta^2)u = c^2(\partial_\xi^2 + 2\partial_\xi\partial_\eta + \partial_\eta^2)u\)，化简得 \(-4c^2 u_{\xi\eta} = 0\)，即 \(u_{\xi\eta} = 0\)。积分得通解 \(u = F(\xi) + G(\eta) = F(x + ct) + G(x - ct)\)。由初始条件 \(u(x,0) = g(x)\) 得 \(F(x) + G(x) = g(x)\)；由 \(u_t(x,0) = h(x)\) 得 \(cF'(x) - cG'(x) = h(x)\)，即 \(F'(x) - G'(x) = h(x)/c\)。积分得 \(F(x) - G(x) = \frac{1}{c} \int_0^x h(s) ds + C\)。与 \(F(x) + G(x) = g(x)\) 联立解得 \(F(x) = \frac{1}{2}g(x) + \frac{1}{2c} \int_0^x h(s) ds + \frac{C}{2}\)，\(G(x) = \frac{1}{2}g(x) - \frac{1}{2c} \int_0^x h(s) ds - \frac{C}{2}\)。代入 \(u = F(x+ct) + G(x-ct)\) 即得 d'Alembert 公式。\(\blacksquare\)

\textbf{定理 92.8}（Kirchhoff 公式，三维）：三维波动方程的解为

\[u(x, t) = \frac{1}{4\pi c^2 t} \iint_{\partial B(x, ct)} h(y) dS(y) + \frac{\partial}{\partial t}\left(\frac{1}{4\pi c^2 t} \iint_{\partial B(x, ct)} g(y) dS(y)\right)\]

（球面均值法）。三维波动方程满足 Huygens 原理（强形式）：扰动沿球面传播，有清晰的前后波阵面。

\emph{证明}：利用球面均值法。定义球面平均 \(U(x; r, t) = \frac{1}{4\pi r^2} \iint_{\partial B(x, r)} u(y, t) dS(y)\)。由波动方程的球对称性，\(U\) 关于 \(r\) 和 \(t\) 满足 Euler-Poisson-Darboux 方程。关键步骤：定义 \(M(x; r, t) = r U(x; r, t)\)。直接计算可得 \(M_{tt} = c^2 M_{rr}\)（一维波动方程）。由 d'Alembert 公式，\(M(x; r, t) = F(r + ct) + G(r - ct)\)。由 \(r=0\) 处的正则性条件 \(M(x; 0, t) = 0\)，得 \(F(ct) + G(-ct) = 0\)，故 \(G(s) = -F(-s)\)。于是 \(M(x; r, t) = F(r + ct) - F(ct - r)\)。求导并令 \(r \to 0\)：\(u(x, t) = \lim_{r \to 0} U(x; r, t) = \lim_{r \to 0} M(x; r, t)/r = 2F'(ct)\)。由初始条件 \(u(x, 0) = g(x)\) 和 \(u_t(x, 0) = h(x)\) 确定 \(F'\)，最终得到 Kirchhoff 公式。三维情形中 \(rU\) 满足严格的 d'Alembert 方程，故 Huygens 原理成立：扰动严格沿球面 \(|x-y| = ct\) 传播，有清晰的前后波阵面。\(\blacksquare\)

\textbf{定理 92.9}（能量守恒）：对波动方程 \(u_{tt} = \Delta u\)（在有界区域 \(\Omega\) 上，齐次边界条件），能量

\[E(t) = \frac{1}{2} \int_\Omega (u_t^2 + |\nabla u|^2) dx\]

是守恒量：\(\frac{dE}{dt} = 0\)。

\emph{证明}：对能量求导，\(\frac{dE}{dt} = \frac{d}{dt} \frac{1}{2} \int_\Omega (u_t^2 + |\nabla u|^2) dx = \int_\Omega (u_t u_{tt} + \nabla u \cdot \nabla u_t) dx\)。对第二项分部积分：\(\int_\Omega \nabla u \cdot \nabla u_t \, dx = \int_{\partial\Omega} u_t \frac{\partial u}{\partial n} dS - \int_\Omega u_t \Delta u \, dx\)。代入得 \(\frac{dE}{dt} = \int_\Omega u_t(u_{tt} - \Delta u) dx + \int_{\partial\Omega} u_t \frac{\partial u}{\partial n} dS\)。由于 \(u\) 满足波动方程 \(u_{tt} = \Delta u\)，第一项为零。在齐次边界条件（Dirichlet：\(u = 0\) 或 Neumann：\(\frac{\partial u}{\partial n} = 0\)）下，边界项也为零。故 \(\frac{dE}{dt} = 0\)，能量守恒。对于全空间情形，边界项在无穷远处消失（由有限传播速度或适当的衰减条件），能量仍守恒。\(\blacksquare\)

\textbf{定理 92.10}（抛物型方程的能量估计与唯一性）：设 \(u \in C^{2,1}(\Omega \times [0, T]) \cap C(\overline{\Omega} \times [0, T])\) 满足热方程 \(u_t - \Delta u = f(x, t)\)。则对任意 \(t \in [0, T]\)，有能量估计：

\[\int_\Omega u(x, t)^2 dx + 2\int_0^t \int_\Omega |\nabla u|^2 dx d\tau = \int_\Omega u(x, 0)^2 dx + 2\int_0^t \int_\Omega f u \, dx d\tau\]

特别地，若 \(f \equiv 0\)，则初边值问题的解是唯一的。

\textbf{证明}：将热方程乘以 \(u\) 并在 \(\Omega\) 上积分：

\[\int_\Omega u u_t \, dx - \int_\Omega u \Delta u \, dx = \int_\Omega f u \, dx\]

第一项 \(\int_\Omega u u_t dx = \frac{1}{2}\frac{d}{dt}\int_\Omega u^2 dx\)。第二项由 Green 恒等式：

\[-\int_\Omega u \Delta u \, dx = \int_\Omega |\nabla u|^2 dx - \int_{\partial\Omega} u \frac{\partial u}{\partial n} dS\]

在 Dirichlet 边界条件 \((u|_{\partial\Omega} = 0)\) 或 Neumann 边界条件 \((\frac{\partial u}{\partial n}|_{\partial\Omega} = 0)\) 下，边界项为零。因此：

\[\frac{1}{2}\frac{d}{dt}\int_\Omega u^2 dx + \int_\Omega |\nabla u|^2 dx = \int_\Omega f u \, dx\]

对 \(t\) 从 \(0\) 到 \(t\) 积分即得能量等式。唯一性：若 \(u_1, u_2\) 是同一初边值问题（齐次边界条件）的解，差 \(w = u_1 - u_2\) 满足齐次方程 \(w_t - \Delta w = 0\)，由能量等式（\(f=0\)）得：

\[\int_\Omega w(x, t)^2 dx + 2\int_0^t \int_\Omega |\nabla w|^2 dx d\tau = \int_\Omega w(x, 0)^2 dx = 0\]

故 \(\int_\Omega w(x,t)^2 dx = 0\) 对所有 \(t\) 成立，即 \(w \equiv 0\)。\(\blacksquare\)

\textbf{定理 92.11}（双曲型方程的有限传播速度）：对波动方程 \(u_{tt} - c^2 \Delta u = 0\)，初始扰动以有限速度 \(c\) 传播。具体地，若初始数据 \((u(x, 0), u_t(x, 0))\) 的支集包含在球 \(B(x_0, R)\) 中，则在时刻 \(t\)，\(u(\cdot, t)\) 的支集包含在 \(B(x_0, R + ct)\) 中。

\textbf{证明}（能量法）：定义局部能量锥。对任意 \(x_0 \in \mathbb{R}^n\)，\(t_0 > 0\)，考虑后向光锥

\[K = \{(x, t) : |x - x_0| \leq c(t_0 - t), 0 \leq t \leq t_0\}\]

定义 \(K\) 在时间 \(t\) 的截面 \(D_t = \{x : |x - x_0| \leq c(t_0 - t)\}\)。在 \(K\) 上，对波动方程乘以 \(u_t\) 并分部积分可得局部能量等式：

\[\frac{d}{dt}\left(\frac{1}{2}\int_{D_t} (u_t^2 + c^2|\nabla u|^2) dx\right) = -c\int_{\partial D_t} \left(\frac{1}{2}(u_t^2 + c^2|\nabla u|^2) + c u_t \frac{\partial u}{\partial r}\right) dS\]

利用 \(|\partial u/\partial r| \leq |\nabla u|\) 和 Cauchy-Schwarz 不等式，右端 \(\leq 0\)。故 \(E(t) = \frac{1}{2}\int_{D_t} (u_t^2 + c^2|\nabla u|^2) dx\) 在 \(K\) 内不增。

若 \(B(x_0, ct_0) \cap \operatorname{supp}(u(\cdot, 0), u_t(\cdot, 0)) = \emptyset\)，则 \(E(0) = 0\)，由能量单调性 \(E(t) \equiv 0\) 对所有 \(t \in [0, t_0]\)。特别地，\(u \equiv 0\) 在 \((x_0, t_0)\) 附近，即扰动未传播到该点。等价地，初始扰动在时刻 \(t\) 不会传出 \(R + ct\) 半径的球。\(\blacksquare\)

\textbf{定理 92.12}（Lax-Milgram 定理在椭圆方程中的应用）：设 \(\Omega \subset \mathbb{R}^n\) 为有界开集，考虑椭圆边值问题：

\[-\Delta u + u = f \quad \text{在 } \Omega \text{ 中}, \quad u|_{\partial\Omega} = 0\]

其中 \(f \in L^2(\Omega)\)。则对任意 \(f \in L^2(\Omega)\)，该问题在 \(H_0^1(\Omega)\) 中存在唯一弱解 \(u\)，且 \(\|u\|_{H^1} \leq C \|f\|_{L^2}\)。

\textbf{证明}：定义双线性型 \(B: H_0^1(\Omega) \times H_0^1(\Omega) \to \mathbb{R}\)：

\[B(u, v) = \int_\Omega (\nabla u \cdot \nabla v + uv) dx\]

以及线性泛函 \(F(v) = \int_\Omega f v \, dx\)。

验证 Lax-Milgram 定理的条件：

\begin{enumerate}
\def\labelenumi{(\arabic{enumi})}
\item
  \(B\) 有界：\(|B(u, v)| \leq \|\nabla u\|_{L^2} \|\nabla v\|_{L^2} + \|u\|_{L^2} \|v\|_{L^2} \leq \|u\|_{H^1} \|v\|_{H^1}\)。
\item
  \(B\) 强制的：\(B(u, u) = \|\nabla u\|_{L^2}^2 + \|u\|_{L^2}^2 = \|u\|_{H^1}^2\)（在 \(H_0^1(\Omega)\) 上 \(\|u\|_{H^1}\) 与 \(\|\nabla u\|_{L^2} + \|u\|_{L^2}\) 等价，且在此内积下恰为 \(H^1\) 范数）。
\item
  \(F\) 有界：\(|F(v)| \leq \|f\|_{L^2} \|v\|_{L^2} \leq \|f\|_{L^2} \|v\|_{H^1}\)。
\end{enumerate}

由 Lax-Milgram 定理，存在唯一的 \(u \in H_0^1(\Omega)\) 使得 \(B(u, v) = F(v)\) 对所有 \(v \in H_0^1(\Omega)\) 成立。这正是弱解的定义：\(\int_\Omega (\nabla u \cdot \nabla v + uv) dx = \int_\Omega f v \, dx\)。由此得 \(\|u\|_{H^1}^2 = B(u, u) = F(u) \leq \|f\|_{L^2} \|u\|_{H^1}\)，故 \(\|u\|_{H^1} \leq \|f\|_{L^2}\)。\(\blacksquare\)

\textbf{定理 92.13}（Galerkin 方法------弱解存在性）：设 \(\Omega \subseteq \mathbb{R}^n\) 为有界开集，\(f \in L^2(\Omega)\)。则边值问题 \(-\Delta u = f\)（\(u|_{\partial\Omega} = 0\)）存在弱解 \(u \in H_0^1(\Omega)\)。

\textbf{证明（Galerkin 逼近）}：由于 \(H_0^1(\Omega)\) 是可分的 Hilbert 空间，选取一组光滑函数 \(\{w_k\}_{k=1}^\infty \subset C_c^\infty(\Omega)\) 构成 \(H_0^1(\Omega)\) 的一组正交基（可取 \(-\Delta\) 在 Dirichlet 条件下的特征函数）。

对每个正整数 \(m\)，在有限维子空间 \(V_m = \operatorname{span}\{w_1, \ldots, w_m\}\) 中求近似解 \(u_m = \sum_{k=1}^m d_k w_k\) 满足：

\[\int_\Omega \nabla u_m \cdot \nabla w_j \, dx = \int_\Omega f w_j \, dx, \quad j = 1, 2, \ldots, m \tag{*}\]

这等价于线性方程组 \(A\mathbf{d} = \mathbf{b}\)，其中 \(A_{kj} = \int_\Omega \nabla w_k \cdot \nabla w_j dx\)，\(b_j = \int_\Omega f w_j dx\)。由于 \(\{w_k\}\) 线性无关且 \(\nabla\) 在 \(H_0^1\) 中的双线性型在正交基下正定，\(A\) 正定可逆，故 \(u_m\) 存在唯一。

取试探函数 \(v = u_m\)（将 \((*)\) 乘以 \(d_j\) 并求和）：

\[\int_\Omega |\nabla u_m|^2 dx = \int_\Omega f u_m \, dx \leq \|f\|_{L^2} \|u_m\|_{L^2}\]

由 Poincare 不等式 \(\|u_m\|_{L^2} \leq C \|\nabla u_m\|_{L^2}\)，得：

\[\|\nabla u_m\|_{L^2}^2 \leq C \|f\|_{L^2} \|\nabla u_m\|_{L^2} \implies \|\nabla u_m\|_{L^2} \leq C \|f\|_{L^2}\]

因此 \(\{u_m\}\) 在 \(H_0^1(\Omega)\) 中一致有界。由 Banach-Alaoglu 定理（Hilbert 空间的弱紧致性），存在子列（仍记作 \(u_m\)）弱收敛于某 \(u \in H_0^1(\Omega)\)，即 \(u_m \rightharpoonup u\) 在 \(H_0^1\) 中。

对固定 \(j\)，取试探函数 \(w_j\) 并令 \(m \to \infty\)（\(m \geq j\)）。由弱收敛性 \(\int_\Omega \nabla u_m \cdot \nabla w_j dx \to \int_\Omega \nabla u \cdot \nabla w_j dx\)。故：

\[\int_\Omega \nabla u \cdot \nabla w_j dx = \int_\Omega f w_j dx \quad \text{对所有 } j\]

由 \(\{w_j\}\) 的稠密性，上式对所有 \(v \in H_0^1(\Omega)\) 成立。因此 \(u\) 是弱解。\(\blacksquare\)

\hrulefill

\hrulefill

\hrulefill

\hrulefill

\hrulefill

\hrulefill

\section{第315章：Sobolev 空间与弱解}\label{ux7b2c315ux7ae0sobolev-ux7a7aux95f4ux4e0eux5f31ux89e3}

现代 PDE 理论的核心工具是 Sobolev 空间。它将经典导数推广为弱导数（分布导数），使得在更广的函数空间中寻找 PDE 的解成为可能。

\subsection{93.1 弱导数与 Sobolev 空间}\label{ux5f31ux5bfcux6570ux4e0e-sobolev-ux7a7aux95f4}

\textbf{定义 93.1}（弱导数）：设 \(\Omega \subseteq \mathbb{R}^n\) 是开集，\(u \in L^1_{\text{loc}}(\Omega)\)。\(u\) 的\textbf{弱导数}（或分布导数）\(v = D^\alpha u\)（\(L^1_{\text{loc}}(\Omega)\) 中）满足

\[\int_\Omega u D^\alpha \varphi \, dx = (-1)^{|\alpha|} \int_\Omega v \varphi \, dx \quad \text{对所有 } \varphi \in C_c^\infty(\Omega)\]

其中 \(D^\alpha \varphi\) 是 \(\varphi\) 的经典 \(\alpha\) 阶偏导数。

\textbf{定义 93.2}（Sobolev 空间）：对 \(1 \leq p \leq \infty\) 和整数 \(k \geq 0\)，\textbf{Sobolev 空间} \(W^{k,p}(\Omega)\) 定义为

\[W^{k,p}(\Omega) = \{u \in L^p(\Omega) : D^\alpha u \in L^p(\Omega) \text{ 对所有 } |\alpha| \leq k\}\]

范数为 \(\|u\|_{W^{k,p}} = \left(\sum_{|\alpha| \leq k} \|D^\alpha u\|_{L^p}^p\right)^{1/p}\)（\(p < \infty\)）。\(H^k(\Omega) = W^{k,2}(\Omega)\) 是 Hilbert 空间。

\textbf{命题 93.1}：\(W^{k,p}(\Omega)\) 是 Banach 空间。\(H^k(\Omega)\) 是 Hilbert 空间，内积为 \((u, v)_{H^k} = \sum_{|\alpha| \leq k} \int_\Omega D^\alpha u \cdot D^\alpha v \, dx\)。

\emph{证明}：由定义，\(W^{k,p}(\Omega)\) 等距同构于 \(L^p(\Omega)^{N}\) 的闭子空间（\(N\) 为多指标数），其中嵌入由 \(u \mapsto (D^\alpha u)_{|\alpha| \leq k}\) 给出。若 \(\{u_n\}\) 是 Cauchy 列，则每个 \(\{D^\alpha u_n\}\) 在 \(L^p\) 中 Cauchy，故收敛于某 \(v_\alpha \in L^p\)。需验证 \(v_0\) 的弱导数为 \(v_\alpha\)。由弱导数定义取极限即得，故 \(W^{k,p}\) 完备，为 Banach 空间。\(H^k = W^{k,2}\) 的内积由 \(L^2\) 内积之和定义，直接验证为 Hilbert 空间。\(\blacksquare\)

\subsection{93.2 Sobolev 嵌入定理}\label{sobolev-ux5d4cux5165ux5b9aux7406}

\textbf{定理 93.2}（Sobolev 嵌入定理）：设 \(\Omega \subseteq \mathbb{R}^n\) 是有 Lipschitz 边界的有界开集，\(1 \leq p < \infty\)。

\begin{enumerate}
\def\labelenumi{\arabic{enumi}.}
\tightlist
\item
  若 \(k p < n\)，则 \(W^{k,p}(\Omega) \hookrightarrow L^q(\Omega)\)（连续嵌入），其中 \(\frac{1}{q} = \frac{1}{p} - \frac{k}{n}\)（Sobolev 共轭指数）
\item
  若 \(k p = n\)，则 \(W^{k,p}(\Omega) \hookrightarrow L^q(\Omega)\) 对所有 \(q < \infty\) 成立
\item
  若 \(k p > n\)，则 \(W^{k,p}(\Omega) \hookrightarrow C^{0,\gamma}(\overline{\Omega})\)（Hölder 连续函数空间），其中 \(\gamma = k - n/p\)（若 \(k - n/p < 1\)）
\end{enumerate}

\textbf{证明}（框架）：(1) 当 \(kp < n\) 时，通过反复应用 Gagliardo-Nirenberg-Sobolev 不等式：对 \(u \in C_c^\infty(\mathbb{R}^n)\)，\(\|u\|_{L^{p^*}} \leq C \|\nabla u\|_{L^p}\)，其中 \(p^* = np/(n-p)\)。对 \(k\) 阶导数，逐次应用此不等式得 \(\|u\|_{L^q} \leq C \|D^k u\|_{L^p}\)，其中 \(1/q = 1/p - k/n\)。对 \(\Omega\) 有界且 Lipschitz 边界，使用延拓算子将 \(W^{k,p}(\Omega)\) 嵌入 \(W^{k,p}(\mathbb{R}^n)\)，再应用全空间的不等式。(2) 当 \(kp = n\) 时，\(W^{k,p}\) 嵌入 \(L^q\) 对所有 \(q < \infty\) 成立，但未必嵌入 \(L^\infty\)，最佳替代是 Moser-Trudinger 不等式。(3) 当 \(kp > n\) 时，由 Morrey 不等式：\(|u(x) - u(y)| \leq C \|D u\|_{L^p} |x-y|^{1-n/p}\)。对 \(k\) 阶情形逐次应用，得 \(u \in C^{k-[n/p]-1, \gamma}\)，其中 \(\gamma = k - n/p - [k - n/p]\)（或 \(k - n/p\) 若 \(k - n/p < 1\)）。\(\blacksquare\)

\textbf{定理 93.5}（Poisson 方程弱解的存在唯一性）：对 \(f \in L^2(\Omega)\)，Poisson 方程的 Dirichlet 问题存在唯一弱解 \(u \in H_0^1(\Omega)\)。

\emph{证明}：定义双线性型 \(B(u, v) = \int_\Omega \nabla u \cdot \nabla v \, dx\) 和线性泛函 \(F(v) = \int_\Omega f v \, dx\)。由 Poincare 不等式，对 \(u \in H_0^1(\Omega)\)，\(\|u\|_{L^2} \leq C \|\nabla u\|_{L^2}\)，故 \(B(u, u) = \|\nabla u\|_{L^2}^2 \geq \frac{1}{1+C^2} \|u\|_{H^1}^2\)，即 \(B\) 是强制的。\(B\) 的有界性：\(|B(u, v)| \leq \|\nabla u\|_{L^2} \|\nabla v\|_{L^2} \leq \|u\|_{H^1} \|v\|_{H^1}\)。\(F\) 的有界性：\(|F(v)| \leq \|f\|_{L^2} \|v\|_{L^2} \leq \|f\|_{L^2} \|v\|_{H^1}\)。由 Lax-Milgram 定理，存在唯一 \(u \in H_0^1(\Omega)\) 使得 \(B(u, v) = F(v)\) 对所有 \(v \in H_0^1(\Omega)\) 成立。这正是 Poisson 方程 \(-\Delta u = f\) 在 \(H_0^1(\Omega)\) 中的弱解定义。\(\blacksquare\)

\subsection{93.4 正则性理论}\label{ux6b63ux5219ux6027ux7406ux8bba}

\textbf{定理 93.6}（椭圆正则性）：设 \(u \in H_0^1(\Omega)\) 是 \(-\Delta u = f\) 的弱解，\(\Omega\) 有光滑边界。 1. 若 \(f \in L^2(\Omega)\)，则 \(u \in H^2(\Omega)\) 且 \(\|u\|_{H^2} \leq C \|f\|_{L^2}\) 2. 若 \(f \in H^k(\Omega)\)，则 \(u \in H^{k+2}(\Omega)\) 3. 若 \(f \in C^\infty(\overline{\Omega})\)，则 \(u \in C^\infty(\overline{\Omega})\)

由此，弱解在光滑数据下是经典解，PDE 的弱解理论与经典解理论一致。

\textbf{证明}：利用差分商方法和 Nirenberg 平移技巧。(1) 设 \(u\) 为弱解，满足 \(\int_\Omega \nabla u \cdot \nabla v = \int_\Omega f v\) 对所有 \(v \in H_0^1(\Omega)\)。对 \(f \in L^2(\Omega)\)，取差分商 \(D_h u(x) = (u(x+he_i) - u(x))/h\) 作为试探函数（在内部子区域 \(\Omega' \subset\subset \Omega\) 上）。由能量估计 \(\|\nabla D_h u\|_{L^2(\Omega')} \leq C(\|f\|_{L^2(\Omega)} + \|u\|_{H^1(\Omega)})\)，取 \(h \to 0\) 得 \(D^2 u \in L^2_{\text{loc}}(\Omega)\)。边界正则性通过边界展平（\(C^2\) 边界）和奇偶反射技巧，将内估计推广到边界，得 \(\|u\|_{H^2(\Omega)} \leq C \|f\|_{L^2(\Omega)}\)。(2) 对 \(f \in H^k(\Omega)\)，通过 bootstrap 论证：对 \(u\) 的 \(k+1\) 阶弱导数重复上述差分商估计，得 \(u \in H^{k+2}(\Omega)\)。(3) 若 \(f \in C^\infty\)，则 \(u \in H^k\) 对所有 \(k\) 成立，由 Sobolev 嵌入 \(H^k \hookrightarrow C^{k-n/2}\) 得 \(u \in C^\infty\)。\(\blacksquare\)

\textbf{定理 93.7}（Weyl 引理）：若 \(u \in L^2_{\text{loc}}(\Omega)\) 满足 \(\int_\Omega u \Delta \varphi = 0\) 对所有 \(\varphi \in C_c^\infty(\Omega)\)，则 \(u\) 几乎处处等于一个调和函数。Weyl 引理说明了调和函数的分布定义与经典定义的等价性。

\emph{证明}：取磨光核 \(\eta_\varepsilon\)，令 \(u_\varepsilon = \eta_\varepsilon * u\)。对任意 \(\varphi \in C_c^\infty\)，\(\int u_\varepsilon \Delta \varphi = \int u (\eta_\varepsilon * \Delta \varphi) = \int u \Delta(\eta_\varepsilon * \varphi) = 0\)（因为 \(\eta_\varepsilon * \varphi \in C_c^\infty\)）。故 \(u_\varepsilon\) 在经典意义下是调和的（在 \(\Omega_\varepsilon = \{x \in \Omega : \operatorname{dist}(x, \partial\Omega) > \varepsilon\}\) 上，因为 \(u_\varepsilon \in C^\infty\) 且 \(\Delta u_\varepsilon = 0\)）。由调和函数的均值性质，\(u_\varepsilon\) 满足局部一致估计：对任意紧子集 \(K \subset\subset \Omega\)，\(\sup_K |D^\alpha u_\varepsilon| \leq C_K \|u_\varepsilon\|_{L^2} \leq C_K \|u\|_{L^2}\)。由 Arzela-Ascoli，存在子列 \(u_{\varepsilon_k}\) 在 \(C^\infty_{\text{loc}}(\Omega)\) 中收敛于某调和函数 \(v\)。同时 \(u_\varepsilon \to u\) 在 \(L^2_{\text{loc}}\) 中，故 \(u = v\) 几乎处处成立。因此 \(u\) 几乎处处等于一个调和函数。\(\blacksquare\)

\hrulefill

\hrulefill

\hrulefill

\hrulefill

\hrulefill

\newpage

\chapter{09-微分方程与动力系统/01-微分方程/03-Sobolev空间与特殊函数.md}\label{ux5faeux5206ux65b9ux7a0bux4e0eux52a8ux529bux7cfbux7edf01-ux5faeux5206ux65b9ux7a0b03-sobolevux7a7aux95f4ux4e0eux7279ux6b8aux51fdux6570.md}

\section{第316章：Sobolev 空间}\label{ux7b2c316ux7ae0sobolev-ux7a7aux95f4}

Sobolev 空间是 PDE 理论的基石，由 С.Л. Соболев 于 1930 年代引入。它将函数的可微性与 \(L^p\) 可积性统一起来，使得 PDE 的弱解理论成为可能。本章建立从弱导数到嵌入定理的完整框架。

\subsection{\texorpdfstring{弱导数与 Sobolev 空间 \(W^{k,p}\)}{弱导数与 Sobolev 空间 W\^{}\{k,p\}}}\label{ux5f31ux5bfcux6570ux4e0e-sobolev-ux7a7aux95f4-wkp}

设 \(\Omega \subset \mathbb{R}^n\) 为开集。记 \(L^1_{\text{loc}}(\Omega)\) 为 \(\Omega\) 上局部可积函数空间。对多重指标 \(\alpha = (\alpha_1, \ldots, \alpha_n)\)，记 \(D^\alpha = \partial_1^{\alpha_1} \cdots \partial_n^{\alpha_n}\)，\(|\alpha| = \sum \alpha_i\)。

\textbf{定义 251.1}（弱导数）：设 \(u, v \in L^1_{\text{loc}}(\Omega)\)。称 \(v\) 为 \(u\) 的第 \(\alpha\) 阶\textbf{弱导数}（或广义导数），记作 \(v = D^\alpha u\)，若对一切测试函数 \(\varphi \in C_c^\infty(\Omega)\) 成立

\[\int_\Omega u(x) D^\alpha \varphi(x)\,dx = (-1)^{|\alpha|} \int_\Omega v(x) \varphi(x)\,dx\]

此即分部积分公式。若 \(u \in C^{|\alpha|}(\Omega)\)，则经典导数与弱导数一致；弱导数在 \(L^1_{\text{loc}}\) 意义下唯一。

\textbf{定义 251.2}（Sobolev 空间 \(W^{k,p}(\Omega)\)）：设 \(k \in \mathbb{N}\)，\(1 \leq p \leq \infty\)。\textbf{Sobolev 空间} \(W^{k,p}(\Omega)\) 定义为

\[W^{k,p}(\Omega) = \left\{ u \in L^p(\Omega) : D^\alpha u \in L^p(\Omega), \; \forall |\alpha| \leq k \right\}\]

其中 \(D^\alpha u\) 为弱导数。配备 \textbf{Sobolev 范数}

\[\|u\|_{W^{k,p}(\Omega)} = \begin{cases}
\left( \sum_{|\alpha| \leq k} \|D^\alpha u\|_{L^p(\Omega)}^p \right)^{1/p}, & 1 \leq p < \infty \\
\max_{|\alpha| \leq k} \|D^\alpha u\|_{L^\infty(\Omega)}, & p = \infty
\end{cases}\]

\textbf{定理 251.1}（\(W^{k,p}\) 的完备性）：对任意 \(k \in \mathbb{N}\)，\(1 \leq p \leq \infty\)，\(W^{k,p}(\Omega)\) 是 Banach 空间。

\textbf{证明}：令 \(\{u_m\}\) 为 \(W^{k,p}\) 中的 Cauchy 列。对每个 \(|\alpha| \leq k\)，\(\{D^\alpha u_m\}\) 在 \(L^p\) 中为 Cauchy 列。由 \(L^p\) 完备性，存在 \(u_\alpha \in L^p\) 使得 \(D^\alpha u_m \to u_\alpha\)（在 \(L^p\) 中）。记 \(u = u_{(0,\ldots,0)}\)。对 \(\varphi \in C_c^\infty(\Omega)\)，

\[\int u D^\alpha \varphi = \lim_{m} \int u_m D^\alpha \varphi = \lim_{m} (-1)^{|\alpha|} \int D^\alpha u_m \,\varphi = (-1)^{|\alpha|} \int u_\alpha \varphi\]

故 \(u_\alpha = D^\alpha u\)（弱导数），从而 \(u \in W^{k,p}\) 且 \(u_m \to u\) 在 \(W^{k,p}\) 中。\(\blacksquare\)

当 \(p=2\) 时，\(H^k(\Omega) = W^{k,2}(\Omega)\) 为 Hilbert 空间，内积为 \(\langle u, v \rangle_{H^k} = \sum_{|\alpha| \leq k} \langle D^\alpha u, D^\alpha v \rangle_{L^2}\)。

\subsection{\texorpdfstring{分数阶 Sobolev 空间 \(H^s(\mathbb{R}^n)\)}{分数阶 Sobolev 空间 H\^{}s(\textbackslash mathbb\{R\}\^{}n)}}\label{ux5206ux6570ux9636-sobolev-ux7a7aux95f4-hsmathbbrn}

\textbf{定义 251.3}（\(H^s(\mathbb{R}^n)\)，\(s \in \mathbb{R}\)）：利用 Fourier 变换 \(\mathcal{F}\)，定义

\[H^s(\mathbb{R}^n) = \left\{ u \in \mathcal{S}'(\mathbb{R}^n) : (1 + |\xi|^2)^{s/2} \widehat{u}(\xi) \in L^2(\mathbb{R}^n) \right\}\]

范数为 \(\|u\|_{H^s} = \| (1 + |\xi|^2)^{s/2} \widehat{u} \|_{L^2}\)。称 \((1+|\xi|^2)^{s/2}\) 为 \(s\) 阶 Bessel 位势。

\textbf{定理 251.2}：\(H^s(\mathbb{R}^n)\) 是 Hilbert 空间。对 \(s = k \in \mathbb{N}\)，\(H^k(\mathbb{R}^n)\) 与 \(W^{k,2}(\mathbb{R}^n)\) 等距同构。

\textbf{证明}：Hilbert 结构由 \(\langle u, v \rangle = \int (1+|\xi|^2)^s \widehat{u} \overline{\widehat{v}}\,d\xi\) 给出。对 \(s=k \in \mathbb{N}\)，Plancherel 定理保证 \(\|D^\alpha u\|_{L^2} = \|\xi^\alpha \widehat{u}\|_{L^2}\)，而 \(c(1+|\xi|^2)^k \leq \sum_{|\alpha|\leq k} |\xi^\alpha|^2 \leq C(1+|\xi|^2)^k\)，故两范数等价。\(\blacksquare\)

\subsection{逼近与稠密性}\label{ux903cux8fd1ux4e0eux7a20ux5bc6ux6027}

\textbf{定义 251.4}（\(W^{k,p}_0(\Omega)\)）：记 \(W^{k,p}_0(\Omega)\) 为 \(C_c^\infty(\Omega)\) 在 \(W^{k,p}(\Omega)\) 中的闭包。当 \(\Omega = \mathbb{R}^n\) 时 \(W^{k,p}_0(\mathbb{R}^n) = W^{k,p}(\mathbb{R}^n)\)。

\(W^{k,p}_0(\Omega)\) 在偏微分方程理论中扮演''零边界条件''的角色：其元素在边界 \(\partial\Omega\) 上''迹为零''（在迹算子意义下）。当 \(k=1\) 时，\(W^{1,p}_0(\Omega)\) 恰好是那些在 \(\partial\Omega\) 上迹为零的 \(W^{1,p}(\Omega)\) 函数的全体。Sobolev 空间中光滑函数的稠密性是证明众多不等式和嵌入定理的基础------先对光滑函数建立不等式，再通过稠密性（及范数的连续性）推广到整个 Sobolev 空间。这一策略被称为''光滑化-稠密性论证''（density argument），是泛函分析中处理 PDE 弱解的标准方法。

\textbf{定理 251.3}（Meyers-Serrin 定理，\(H = W\)）：设 \(\Omega \subset \mathbb{R}^n\) 为开集，\(1 \leq p < \infty\)。则 \(C^\infty(\Omega) \cap W^{k,p}(\Omega)\) 在 \(W^{k,p}(\Omega)\) 中稠密。

\textbf{证明}：取 \(\Omega\) 的穷竭序列 \(\{\Omega_j\}_{j=1}^\infty\) 满足 \(\Omega_j \subset\subset \Omega_{j+1}\) 且 \(\bigcup_j \Omega_j = \Omega\)。令 \(\{\psi_j\}\) 为从属于开覆盖 \(\{\Omega_{j+1} \setminus \overline{\Omega}_{j-1}\}\) 的光滑单位分解（\(\sum_j \psi_j = 1\) 在 \(\Omega\) 上，\(\psi_j \in C_c^\infty(\Omega_{j+1} \setminus \overline{\Omega}_{j-1})\)）。对 \(u \in W^{k,p}(\Omega)\)，考虑截断和磨光：\(u_\varepsilon = \sum_{j=1}^\infty \eta_{\varepsilon_j} * (\psi_j u)\)，其中 \(\eta_\varepsilon(x) = \varepsilon^{-n}\eta(x/\varepsilon)\) 为标准磨光核，\(\varepsilon_j\) 充分小以确保 \(\operatorname{supp}(\eta_{\varepsilon_j} * (\psi_j u)) \subset \Omega\)。由磨光核的性质，\(u_\varepsilon \in C^\infty(\Omega)\)。对于 \(|\alpha| \leq k\)，\(D^\alpha u_\varepsilon = \sum_j \eta_{\varepsilon_j} * D^\alpha(\psi_j u)\)（利用磨光与弱导数的交换性）。由 \(L^p\) 中磨光逼近的收敛性，\(\|D^\alpha u_\varepsilon - D^\alpha u\|_{L^p} \to 0\)（\(\varepsilon \to 0\)）。故 \(u_\varepsilon \to u\) 在 \(W^{k,p}(\Omega)\) 中，即 \(C^\infty(\Omega) \cap W^{k,p}(\Omega)\) 在 \(W^{k,p}(\Omega)\) 中稠密。\(\blacksquare\)

\subsection{Sobolev 嵌入定理}\label{sobolev-ux5d4cux5165ux5b9aux7406-1}

\textbf{定理 251.4}（Sobolev 嵌入 ------ \(k p < n\) 情形）：设 \(1 \leq p < n/k\)。则 \(W^{k,p}(\mathbb{R}^n) \hookrightarrow L^{p^*}(\mathbb{R}^n)\) 连续，其中 \textbf{Sobolev 共轭指数} \(p^* = \dfrac{np}{n - kp}\)。即存在 \(C = C(k, p, n)\) 使得

\[\|u\|_{L^{p^*}} \leq C \|u\|_{W^{k,p}}\]

\textbf{证明}（归纳法，\(k=1\) 情形）：设 \(u \in C_c^\infty(\mathbb{R}^n)\)，需证 \(\|u\|_{L^{p^*}} \leq C \|\nabla u\|_{L^p}\)。由 Gagliardo-Nirenberg 不等式，对 \(n \geq 2\)，\(\|u\|_{L^{n/(n-1)}} \leq \frac{1}{2\sqrt{n}} \prod_{i=1}^n \|\partial_i u\|_{L^1}^{1/n}\)。对 \(p > 1\)，取 \(v = |u|^\alpha\)（\(\alpha > 1\) 待定），应用 \(L^1\) 嵌入得 \(\|u\|_{L^{\alpha n/(n-1)}}^\alpha \leq C \alpha \|u^{\alpha-1} \nabla u\|_{L^1}\)。由 Hölder 不等式，\(\|u^{\alpha-1} \nabla u\|_{L^1} \leq \|u\|_{L^{(\alpha-1)p'}}^{\alpha-1} \|\nabla u\|_{L^p}\)。取 \(\alpha = \frac{n(p-1)}{n-p}\) 使 \(\alpha n/(n-1) = (\alpha-1)p' = p^*\)，得 \(\|u\|_{L^{p^*}} \leq C \|\nabla u\|_{L^p}\)。\(k > 1\) 情形对 \(k\) 归纳，每次降低一阶导数。\(\blacksquare\)

\textbf{定理 251.5}（Morrey 不等式 ------ \(k p > n\) 情形）：设 \(p > n\)。则 \(W^{1,p}(\mathbb{R}^n) \hookrightarrow C^{0, \gamma}(\mathbb{R}^n)\) 连续，其中 \(\gamma = 1 - n/p\)。即存在 \(C\) 使得

\[|u(x) - u(y)| \leq C \|Du\|_{L^p} \, |x - y|^{1 - n/p}\]

对所有 \(u \in W^{1,p}\) 以及几乎所有 \(x, y\) 成立。

\textbf{证明}：对 \(u \in C_c^\infty(\mathbb{R}^n)\)，固定 \(x,y \in \mathbb{R}^n\)，令 \(r = |x-y|\)。由球体积分表示，\(|u(x) - u(y)| \leq \int_0^1 |\nabla u(x + t(y-x))| \cdot |y-x| dt\)。取 \(B = B(x,r)\)，利用平均积分技巧：\(|u(x) - u_{B}| \leq \frac{C}{r^{n-1}} \int_B \frac{|\nabla u(z)|}{|x-z|^{n-1}} dz\)，其中 \(u_B = \frac{1}{|B|}\int_B u\)。由 Hölder 不等式，\(\int_B \frac{|\nabla u(z)|}{|x-z|^{n-1}} dz \leq \|\nabla u\|_{L^p(B)} \| |x-z|^{-(n-1)}\|_{L^{p'}(B)}\)（\(p' = p/(p-1)\)）。因 \(p > n\)，\(p'(n-1) < n\)，后一积分有限且 \(\sim r^{1-n/p}\)。故 \(|u(x) - u_B| \leq C r^{1-n/p} \|\nabla u\|_{L^p}\)。对 \(y\) 同理，三角不等式得 \(|u(x) - u(y)| \leq C |x-y|^{1-n/p} \|\nabla u\|_{L^p}\)。由 \(C_c^\infty\) 的稠密性延拓至 \(W^{1,p}\)。\(\blacksquare\)

\textbf{定理 251.6}（\(k p = n\) 极限情形）：设 \(p = n\)。则 \(W^{1,n}(\mathbb{R}^n) \hookrightarrow L^q(\mathbb{R}^n)\) 对一切 \(n \leq q < \infty\) 连续，但一般不嵌入 \(L^\infty\)。更精细的极限情形由 Trudinger-Moser 指数可积性刻画。

\textbf{证明}：由 Sobolev 嵌入定理的极限情形，\(p=n\) 时共轭指数 \(p^* = np/(n-p)\) 趋于无穷，故不能直接嵌入 \(L^\infty\)。对任意 \(q \in [n, \infty)\)，由 Gagliardo-Nirenberg 插值不等式：\(\|u\|_{L^q} \leq C \|u\|_{L^n}^{n/q} \|\nabla u\|_{L^n}^{1-n/q}\)。因 \(W^{1,n} \hookrightarrow L^n\) 连续，得 \(\|u\|_{L^q} \leq C \|u\|_{W^{1,n}}\)。但对 \(q=\infty\) 插值失效，确实存在 \(W^{1,n}\) 函数无界（例如 \(u(x) = \log\log(1/|x|)\) 在 \(B(0,1/e)\) 中）。Trudinger-Moser 不等式给出 \(\int_\Omega \exp(\alpha |u|^{n/(n-1)}) < \infty\)（对适当 \(\alpha\)），这是 \(L^\infty\) 缺失的补偿。\(\blacksquare\)

\subsection{紧嵌入与 Rellich-Kondrachov 定理}\label{ux7d27ux5d4cux5165ux4e0e-rellich-kondrachov-ux5b9aux7406}

\textbf{定理 251.7}（Rellich-Kondrachov 紧嵌入）：设 \(\Omega \subset \mathbb{R}^n\) 为有界 Lipschitz 开集，\(1 \leq p < \infty\)。

\begin{enumerate}
\def\labelenumi{\arabic{enumi}.}
\tightlist
\item
  若 \(k p < n\)，则 \(W^{k,p}(\Omega) \hookrightarrow\hookrightarrow L^q(\Omega)\) 对 \(1 \leq q < p^*\) 为紧嵌入；
\item
  若 \(k p = n\)，则 \(W^{k,p}(\Omega) \hookrightarrow\hookrightarrow L^q(\Omega)\) 对 \(1 \leq q < \infty\) 为紧嵌入；
\item
  若 \(k p > n\)，则 \(W^{k,p}(\Omega) \hookrightarrow\hookrightarrow C^{0,\beta}(\overline{\Omega})\) 对 \(0 \leq \beta < k - n/p\) 为紧嵌入。
\end{enumerate}

其中 \(\hookrightarrow\hookrightarrow\) 表示紧嵌入（将单位球中的有界序列映为预紧序列）。

\textbf{证明}：仅证情形1（\(kp < n\)），其余情形类似。设 \(\{u_m\} \subset W^{k,p}(\Omega)\) 为有界序列，\(\|u_m\|_{W^{k,p}} \leq M\)。由 Sobolev 嵌入 \(W^{k,p} \hookrightarrow L^{p^*}\) 连续，\(\{u_m\}\) 在 \(L^{p^*}\) 中有界。对 \(q < p^*\)，由 Holder 不等式 \(\|u_m\|_{L^q} \leq |\Omega|^{1/q-1/p^*}\|u_m\|_{L^{p^*}} \leq C\)。关键步骤：利用磨光核 \(\eta_\varepsilon\) 构造 \(u_m^\varepsilon = \eta_\varepsilon * u_m\)（延拓到 \(\mathbb{R}^n\) 外为零）。由 Friedrichs 不等式，\(\|u_m^\varepsilon - u_m\|_{L^q} \leq C\varepsilon^\alpha\)（\(\alpha = k - n(1/p-1/q) > 0\)），且 \(\|u_m^\varepsilon\|_{C^1}\) 被 \(\varepsilon^{-k-n/p}M\) 控制。由 Arzela-Ascoli 定理，对固定 \(\varepsilon\)，\(\{u_m^\varepsilon\}_m\) 在 \(L^q\) 中预紧。取对角线子列 \(\{u_{m_j}\}\) 及 \(\varepsilon \to 0\)，利用 \(\varepsilon\)-逼近论证即得 \(\{u_{m_j}\}\) 在 \(L^q(\Omega)\) 中收敛。\(\blacksquare\)

\subsection{迹定理}\label{ux8ff9ux5b9aux7406}

\textbf{定理 251.8}（迹定理）：设 \(\Omega \subset \mathbb{R}^n\) 为有界 \(C^1\)（或 Lipschitz）开集，\(1 \leq p < \infty\)。存在唯一的有界线性算子

\[\gamma : W^{1,p}(\Omega) \to L^p(\partial \Omega)\]

称为\textbf{迹算子}，满足 \(\gamma u = u|_{\partial\Omega}\) 对 \(u \in C^\infty(\overline{\Omega})\) 成立。存在 \(C = C(\Omega, p)\) 使得

\[\|\gamma u\|_{L^p(\partial\Omega)} \leq C \|u\|_{W^{1,p}(\Omega)}\]

\textbf{证明}：先对半空间 \(\mathbb{R}^n_+ = \{(x', x_n): x' \in \mathbb{R}^{n-1}, x_n > 0\}\) 证明迹估计。对 \(u \in C_c^\infty(\overline{\mathbb{R}^n_+})\)，由微积分基本定理：\(|u(x',0)|^p = -\int_0^\infty \frac{\partial}{\partial x_n}|u(x',x_n)|^p dx_n = -p\int_0^\infty |u|^{p-1}\operatorname{sgn}(u)\frac{\partial u}{\partial x_n}dx_n\)。由 Young 不等式，\(|u(x',0)|^p \leq C\int_0^\infty (|u|^p + |\partial_n u|^p)dx_n\)。对 \(x'\) 积分得 \(\|u(\cdot,0)\|_{L^p(\mathbb{R}^{n-1})} \leq C\|u\|_{W^{1,p}(\mathbb{R}^n_+)}\)。对于一般 \(C^1\)（或 Lipschitz）区域 \(\Omega\)，利用局部坐标将边界展平：取有限开覆盖 \(\{U_i\}\) 及 \(C^1\) 微分同胚 \(\Phi_i: U_i \cap \Omega \to \mathbb{R}^n_+\)，将 \(u\) 拉回为半空间上的函数，应用上述估计后再用单位分解求和。迹算子的唯一性由 \(C^\infty(\overline{\Omega})\) 在 \(W^{1,p}(\Omega)\) 中的稠密性保证。\(\blacksquare\)

\textbf{定义 251.5}（迹零空间）：\(W^{1,p}_0(\Omega)\) 恰为 \(\gamma^{-1}(\{0\})\)，即

\[W^{1,p}_0(\Omega) = \{ u \in W^{1,p}(\Omega) : \gamma u = 0 \}\]

换言之，\(W^{1,p}_0\) 是边界上迹为零的 Sobolev 函数，其与 \(C_c^\infty(\Omega)\) 的闭包定义等价（对 Lipschitz 边界）。

\subsection{Poincaré 不等式}\label{poincaruxe9-ux4e0dux7b49ux5f0f}

\textbf{定理 251.9}（Poincaré 不等式 ------ \(W^{1,p}_0\) 形式）：设 \(\Omega \subset \mathbb{R}^n\) 为有界开集。存在 \(C = C(\Omega, p)\) 使得

\[\|u\|_{L^p(\Omega)} \leq C \|Du\|_{L^p(\Omega)}, \quad \forall u \in W^{1,p}_0(\Omega)\]

特别地，\(\|Du\|_{L^p}\) 在 \(W^{1,p}_0(\Omega)\) 上等价于全范数 \(\|u\|_{W^{1,p}}\)。

\textbf{证明}：对 \(u \in C_c^\infty(\Omega)\)，延拓到 \(\mathbb{R}^n\) 上的一个大立方体 \(Q\) 令边界外为零。沿每个坐标方向积分：\(|u(x)| \leq \int_{-\infty}^{x_i} |\partial_i u|\,dt\)。取平均后 Hölder 不等导出估计。再由 \(C_c^\infty\) 的稠密性传递至 \(W^{1,p}_0\)。\(\blacksquare\)

\textbf{定理 251.10}（Poincaré-Wirtinger 不等式 ------ 零均值形式）：设 \(\Omega \subset \mathbb{R}^n\) 为连通有界 Lipschitz 开集。定义均值 \(u_\Omega = \frac{1}{|\Omega|} \int_\Omega u\,dx\)。存在 \(C = C(\Omega, p)\) 使得

\[\|u - u_\Omega\|_{L^p(\Omega)} \leq C \|Du\|_{L^p(\Omega)}, \quad \forall u \in W^{1,p}(\Omega)\]

\textbf{证明}：若不然，存在序列 \(\{u_m\}\) 满足 \(\|u_m - (u_m)_\Omega\|_{L^p} = 1\) 而 \(\|Du_m\|_{L^p} \to 0\)。由 Rellich-Kondrachov，子列 \(u_m \to u\) 在 \(L^p\) 中，而 \(Du = 0\)（弱极限），故 \(u = \text{const}\)。但 \(\int (u_m - (u_m)_\Omega) \to \int (u - u_\Omega) = 0\)，与极限范数为 \(1\) 矛盾。\(\blacksquare\)

\subsection{Sobolev 空间在 PDE 中的应用框架}\label{sobolev-ux7a7aux95f4ux5728-pde-ux4e2dux7684ux5e94ux7528ux6846ux67b6}

Sobolev 空间为 PDE 的弱解理论提供了基本语言。对二阶椭圆方程 \(- \operatorname{div}(A(x) \nabla u) = f\)（齐次 Dirichlet 边界），其\textbf{变分形式}为：求 \(u \in H^1_0(\Omega)\) 满足

\[\int_\Omega A \nabla u \cdot \nabla v \,dx = \int_\Omega f v\,dx, \quad \forall v \in H^1_0(\Omega)\]

由 Poincaré 不等式和 Lax-Milgram 引理，当 \(A\) 一致椭圆时解存在唯一。正则性理论（\(f \in L^2 \Rightarrow u \in H^2\)）依赖于差商和 Sobolev 嵌入。对发展方程（热方程、波动方程），Sobolev 空间 \(H^s\) 提供能量估计的自然框架。对非线性 PDE（Navier-Stokes、反应扩散），紧嵌入和临界 Sobolev 指数控制了紧性和爆破分析。

更具体地，对于椭圆型 PDE，Sobolev 嵌入给出了解的正则性阶梯：若 \(f \in W^{k,p}(\Omega)\)，则弱解 \(u\) 满足 \(u \in W^{k+2,p}_{\operatorname{loc}}(\Omega)\)（内部正则性）。这一结果通过差分商方法（Nirenberg 平移技巧）证明：考察差分商 \(\Delta_h u = \frac{u(x+h)-u(x)}{|h|}\)，利用变分方程的一致椭圆性得到 \(\Delta_h u\) 的 \(H^1\) 估计，进而取极限 \(h \to 0\) 获得二阶弱导数的 \(L^2\) 界。对于抛物型方程 \(\partial_t u + \mathcal{L}u = f\)，自然的工作空间是 \(L^2(0,T; H^1_0(\Omega)) \cap H^1(0,T; H^{-1}(\Omega))\)，其中时间的 Sobolev 正则性为 \(H^1\) 而空间的为正则性的 \(H^1_0\)，嵌入 \(H^1(0,T) \hookrightarrow C([0,T])\) 保证了解在时间上的连续性。对于双曲型方程，能量守恒律在 \(H^1 \times L^2\) 框架下自然成立。对于非线性方程，临界 Sobolev 嵌入 \(H^1 \hookrightarrow L^{2^*}\)（\(2^* = 2n/(n-2)\)）在变分方法和爆破分析中起着核心作用------Pohozaev 恒等式和 Trudinger-Moser 不等式分别处理了临界和极限情形。

\hrulefill

\section{第317章：KAM 理论与自由边界问题}\label{ux7b2c317ux7ae0kam-ux7406ux8bbaux4e0eux81eaux7531ux8fb9ux754cux95eeux9898}

本章涵盖两个看似独立、实则共享''小扰动下结构稳定性''这一核心主题的数学领域。KAM 理论（Kolmogorov-Arnold-Moser）是 Hamilton 动力系统微扰理论的核心成就，证明了近可积 Hamilton 系统在满足 Diophantine 条件的频率下，大多数不变环面在小扰动下得以存留------这颠覆了 Boltzmann 的''遍历假设''，揭示了可积系统的显著刚性。自由边界问题处理 PDE 的解域边界本身是未知的、由解决定的边值问题，其经典例子包括 Stefan 问题（冰-水相变的自由界面演化）和障碍问题（膜与障碍物接触的边界）。

\subsection{96.1 KAM 理论（Kolmogorov-Arnold-Moser）}\label{kam-ux7406ux8bbakolmogorov-arnold-moser}

KAM 理论是 Hamilton 动力系统微扰理论的核心成就，描述了近可积系统在微小扰动下拟周期运动的保持。\textbf{基本问题}：给定可积 Hamilton 系统 \(H_0(I)\)（作用量-角变量坐标，环面 \(I \times \mathbb{T}^n\)），当加入小扰动 \(H_\epsilon = H_0(I) + \epsilon H_1(I, \theta)\) 后，哪些不变环面得以存留？

\textbf{定理 96.1}（Kolmogorov 定理，1954）：若 \(H_0\) 非退化（\(\det(\partial^2 H_0 / \partial I_i \partial I_j) \neq 0\)），则对满足 Diophantine 条件

\[|k \cdot \omega| \geq \frac{\gamma}{\|k\|^\tau}, \quad \forall k \in \mathbb{Z}^n \setminus \{0\}\]

的频率 \(\omega = \nabla H_0(I)\)，相应的不变环面在充分小的 \(\epsilon > 0\) 下扰动后仍然存在------仅发生微小形变，其频率集合具有正 Lebesgue 测度（当 \(\epsilon \to 0\) 时测度趋于全测度）。

\textbf{证明}：Kolmogorov 的证明使用快速收敛的 Newton 型迭代（Kolmogorov 迭代）。在作用量-角变量 \((I,\theta)\) 下，设 \(H(I,\theta) = H_0(I) + \epsilon H_1(I,\theta)\)。目标是通过一系列典则变换逐步消除角度依赖的扰动项。第 \(n\) 步迭代的误差为 \(O(\epsilon^{2^n})\)，实现二次收敛。关键步骤：求解同调方程 \(\{H_0, S\} + \epsilon H_1 = \bar{H}_1\)（Poisson 括号），其中 \(S\) 为生成函数。在 Fourier 展开下，该方程的解为 \(S_k = \epsilon H_{1,k}/(i k\cdot\omega)\)（\(k\neq 0\)）。Diophantine 条件 \(|k\cdot\omega| \geq \gamma\|k\|^{-\tau}\) 保证分母不趋于零过快，使得迭代收敛。收敛后，扰动系统的频率 \(\omega_\epsilon\) 与 \(\omega\) 相差 \(O(\epsilon)\)，且不变环面仅发生 \(O(\epsilon)\) 的形变。\(\blacksquare\)

\textbf{定义 96.1}（Diophantine 条件）：频率向量 \(\omega \in \mathbb{R}^n\) 称为 \((\gamma, \tau)\)-Diophantine 的（\(\gamma > 0\), \(\tau > n-1\)），若对任意非零整数向量 \(k \in \mathbb{Z}^n\)，有 \(|k \cdot \omega| \geq \gamma \|k\|^{-\tau}\)。满足此条件的频率集合在 \(\mathbb{R}^n\) 中具有全测度。

\textbf{定理 96.2}（Arnold 扩散，1964）：未被 KAM 环面覆盖的相空间区域中存在极缓慢的扩散现象（速度为 \(O(e^{-c/\epsilon})\)），这是导致不变环面间''泄露''的核心机制。对于 \(n \geq 3\) 自由度的系统，Arnold 证明了存在穿越由 KAM 环面形成的''网''的轨道链。

\textbf{证明}：对 \(n \geq 3\) 自由度的近可积 Hamilton 系统，KAM 环面是 \(n\) 维的，相空间维数为 \(2n\)。当 \(n \geq 3\) 时，\(2n - n = n \geq 3 > 2\)，KAM 环面不再能分割相空间------它们形成''网''结构，轨道可通过环面间的间隙缓慢扩散。Arnold 构造了过渡链：一系列由双曲周期轨道的不稳定流形与稳定流形横截相交形成的异宿连接。每条连接使轨道从一个环面附近''跳跃''到相邻环面附近。由于 Melnikov 型分裂的大小为 \(O(e^{-c/\epsilon})\)，每次跳跃的速度为 \(O(e^{-c/\epsilon})\)，故扩散速度为 \(O(e^{-c/\epsilon})\)。这解释了为何在近可积系统中，尽管 KAM 环面占据大部分相空间，长时间后仍可能出现大范围的能量输运。\(\blacksquare\)

\subsection{96.2 自由边界问题（Free Boundary Problems）}\label{ux81eaux7531ux8fb9ux754cux95eeux9898free-boundary-problems}

自由边界问题处理 PDE 的解域边界本身是未知的、由解决定的边值问题。

\textbf{定义 96.2}（Stefan 问题 / 冰-水相变）：热方程 \(u_t = \Delta u\) 在固相/液相区域分别成立，在自由界面 \(\Gamma(t)\) 上满足 Stefan 条件

\[[\nabla u \cdot \mathbf{n}] = L V_n\]

其中 \(L > 0\) 为潜热，\(V_n\) 为界面法向速度，\([\cdot]\) 表示界面两侧跳跃。一维 Stefan 问题的显式解（Neumann 解，1860 年代）由误差函数 \(\operatorname{erf}\) 构成，自由边界以 \(\sim\sqrt{t}\) 的速率移动。

\textbf{定义 96.3}（障碍问题 / Obstacle Problem）：给定障碍函数 \(\psi\)，求解 \(u \geq \psi\) 满足

\[\Delta u \leq 0, \quad u \geq \psi, \quad (u - \psi) \Delta u = 0\]

其自由边界是接触集 \(\Gamma = \partial \{u = \psi\} \cap \Omega\)（即 \(u = \psi\) 区域的拓扑边界）。

\textbf{定理 96.3}（Caffarelli 正则性定理，1977）：障碍问题的自由边界在除去一个低维奇异集之外是光滑的（\(C^{1,\alpha}\)），这为整个自由边界正则性理论奠定了基础。具体而言，自由边界点分为正则点（自由边界在该点附近为 \(C^{1,\alpha}\) 超曲面）和奇点（可分类），后者的 Hausdorff 维数不超过 \(n-2\)。

\textbf{证明}：Caffarelli 的证明基于''改进平坦性''（improvement of flatness）迭代。设 \(u\) 为障碍问题的解，\(\Gamma = \partial\{u = \psi\}\)。在正则点 \(x_0 \in \Gamma\) 处，假设 \(u\) 在 \(B_r(x_0)\) 中接近一个线性函数 \(a\cdot(x-x_0)^+\)（\(a \neq 0\)）。通过缩放和紧性论证，构造序列 \(u_k(x) = u(x_0 + r_k x)/r_k^2\)（\(r_k \to 0\)）。该序列在 \(C^{1,\alpha}\) 中收敛到全局解（``blow-up 极限''），该极限为 \(\frac{1}{2}\max(x\cdot e, 0)^2\) 的形式（\(e\) 为单位向量）。由粘性解理论和 Harnack 不等式，自由边界在 \(x_0\) 附近为 \(C^{1,\alpha}\) 超曲面。奇点对应 \(a=0\) 的退化情形，blow-up 极限为二次齐次多项式，可用凸分析分类------其 Hausdorff 维数不超过 \(n-2\)。\(\blacksquare\)

\textbf{定义 96.4}（Hele-Shaw 流动）：描述两平行平板间粘性流体的界面演化。对 Hele-Shaw 胞腔，压强 \(p\) 在流体区域满足 \(\Delta p = 0\)，在自由边界上 \(p = 0\) 且法向速度 \(V_n = -\partial p / \partial n\)。其数学模型等价于 Laplacian 增长（DLA 的连续极限），小扰动下界面不稳定（Saffman-Taylor 指进现象）。

\hrulefill

\section{第318章：特殊函数与正交多项式}\label{ux7b2c318ux7ae0ux7279ux6b8aux51fdux6570ux4e0eux6b63ux4ea4ux591aux9879ux5f0f}

本章对 Gamma、Beta、Bessel、Legendre 及超几何函数给出编号定义、定理与完整证明，作为数学物理中特殊函数的标准知识库。

\subsection{97.1 Gamma 函数与 Beta 函数}\label{gamma-ux51fdux6570ux4e0e-beta-ux51fdux6570}

\textbf{定义 97.1}（Gamma 函数）：\(\Gamma(z)\) 的 Euler 积分定义为

\[\Gamma(z) = \int_0^\infty t^{z-1} e^{-t} \, dt, \quad \operatorname{Re}(z) > 0\]

通过解析延拓可扩充为 \(\mathbb{C} \setminus \{0, -1, -2, \ldots\}\) 上的亚纯函数，诸负整数为一阶极点，留数为 \(\operatorname{Res}(\Gamma, -n) = (-1)^n / n!\)。等价地有 Weierstrass 无穷乘积

\[\frac{1}{\Gamma(z)} = z e^{\gamma z} \prod_{n=1}^\infty \left(1 + \frac{z}{n}\right) e^{-z/n}\]

其中 \(\gamma = \lim_{N \to \infty} \left(\sum_{n=1}^N \frac{1}{n} - \log N\right)\) 为 Euler-Mascheroni 常数。

\textbf{定理 97.1}（函数方程）：\(\Gamma(z+1) = z \Gamma(z)\) 对所有 \(z \in \mathbb{C} \setminus \mathbb{Z}_{\leq 0}\) 成立。

\textbf{证明}：对 \(\operatorname{Re}(z) > 0\)，分部积分：\(\Gamma(z+1) = \int_0^\infty t^z e^{-t} dt = [-t^z e^{-t}]_0^\infty + z \int_0^\infty t^{z-1} e^{-t} dt = z \Gamma(z)\)。由解析延拓的唯一性，该等式在整个定义域上成立。\(\blacksquare\)

特别地，\(\Gamma(n) = (n-1)!\)（\(n \in \mathbb{N}\)）。Euler 反射公式 \(\Gamma(z)\Gamma(1-z) = \pi / \sin(\pi z)\) 及 Legendre 倍元公式 \(\Gamma(2z) = \frac{2^{2z-1}}{\sqrt{\pi}} \Gamma(z) \Gamma(z + \frac{1}{2})\) 亦为基础恒等式。

\textbf{定义 97.2}（Beta 函数）：对 \(\operatorname{Re}(x) > 0\)，\(\operatorname{Re}(y) > 0\)，

\[B(x, y) = \int_0^1 t^{x-1} (1-t)^{y-1} \, dt\]

\textbf{定理 97.2}（Beta 与 Gamma 的关系）：

\[B(x, y) = \frac{\Gamma(x) \Gamma(y)}{\Gamma(x + y)}\]

\textbf{证明}：考虑 \(\Gamma(x)\Gamma(y) = \int_0^\infty \int_0^\infty u^{x-1} v^{y-1} e^{-(u+v)} du dv\)。作变量代换 \(u = ts\)，\(v = (1-t)s\)（\(t \in (0,1)\)，\(s \in (0,\infty)\)），Jacobi 行列式为 \(s\)。则

\[\Gamma(x)\Gamma(y) = \int_0^1 \int_0^\infty (ts)^{x-1} ((1-t)s)^{y-1} e^{-s} \cdot s \, ds \, dt\]

\[= \int_0^1 t^{x-1} (1-t)^{y-1} dt \cdot \int_0^\infty s^{x+y-1} e^{-s} ds = B(x,y) \cdot \Gamma(x+y)\]

除以 \(\Gamma(x+y)\) 即得结论。\(\blacksquare\)

\subsection{97.2 Bessel 函数}\label{bessel-ux51fdux6570}

\textbf{定义 97.3}（Bessel 微分方程与第一类 Bessel 函数）：对阶 \(\nu \in \mathbb{C}\)，Bessel 微分方程为

\[z^2 y'' + z y' + (z^2 - \nu^2) y = 0\]

其第一类解为：

\[J_\nu(z) = \sum_{k=0}^\infty \frac{(-1)^k}{k! \, \Gamma(k + \nu + 1)} \left(\frac{z}{2}\right)^{2k + \nu}\]

级数对任意 \(z \in \mathbb{C}\) 收敛（无穷收敛半径），在 \(z = 0\) 处 \(J_\nu(z) \sim (z/2)^\nu / \Gamma(\nu+1)\)。当 \(\nu \notin \mathbb{Z}\) 时，\(J_\nu\) 与 \(J_{-\nu}\) 线性无关，构成解空间的基。

\textbf{定义 97.4}（第二类 Bessel 函数 / Weber 函数）：对 \(\nu \notin \mathbb{Z}\)，

\[Y_\nu(z) = \frac{J_\nu(z) \cos(\nu \pi) - J_{-\nu}(z)}{\sin(\nu \pi)}\]

对整数 \(\nu = n\)，\(Y_n(z) = \lim_{\nu \to n} Y_\nu(z)\)（极限存在，消除了 \(0/0\) 不定式）。\(Y_\nu(z)\) 在 \(z = 0\) 处具有对数奇性：\(Y_0(z) \sim \frac{2}{\pi} \log z\)。

\textbf{定理 97.3}（Bessel 函数的递推关系）：

\[\begin{aligned} \frac{d}{dz}\left[z^\nu J_\nu(z)\right] &= z^\nu J_{\nu-1}(z) \\ \frac{d}{dz}\left[z^{-\nu} J_\nu(z)\right] &= -z^{-\nu} J_{\nu+1}(z) \end{aligned}\]

以及 \(J_{\nu-1}(z) + J_{\nu+1}(z) = \frac{2\nu}{z} J_\nu(z)\)，\(J_{\nu-1}(z) - J_{\nu+1}(z) = 2 J'_\nu(z)\)。

\textbf{证明}：对级数逐项微分并重组指标即得。以第一个为例：\(z^\nu J_\nu(z) = \sum_{k=0}^\infty \frac{(-1)^k}{k! \Gamma(k+\nu+1)} \frac{z^{2k+2\nu}}{2^{2k+\nu}}\)，对 \(z\) 求导后合并 \(k\) 与 \(k+1\) 项即产生 \(z^\nu J_{\nu-1}(z)\) 的级数。\(\blacksquare\)

\textbf{定理 97.4}（Bessel 函数的生成函数）：对整数阶 \(\nu = n \in \mathbb{Z}\)，

\[\exp\left(\frac{z}{2}\left(t - \frac{1}{t}\right)\right) = \sum_{n=-\infty}^{\infty} J_n(z) \, t^n\]

此即指数函数的 Laurent 展开，系数恰为整数阶 Bessel 函数。该生成函数直接导出 Jacobi-Anger 展开及其在调频信号分析中的应用。

\textbf{证明}：将生成函数 \(\exp(\frac{z}{2}(t-t^{-1}))\) 展开为 Laurent 级数。由指数函数的 Taylor 展开，\(\exp(\frac{z}{2}t) = \sum_{k=0}^\infty \frac{1}{k!}(\frac{z}{2})^k t^k\)，\(\exp(-\frac{z}{2}t^{-1}) = \sum_{\ell=0}^\infty \frac{(-1)^\ell}{\ell!}(\frac{z}{2})^\ell t^{-\ell}\)。乘积的 \(t^n\) 系数为 \(\sum_{k-\ell=n} \frac{(-1)^\ell}{k!\ell!}(\frac{z}{2})^{k+\ell} = \sum_{\ell=\max(0,-n)}^\infty \frac{(-1)^\ell}{(\ell+n)!\ell!}(\frac{z}{2})^{2\ell+n} = J_n(z)\)。当 \(n<0\) 时，利用 \(J_{-n}(z)=(-1)^n J_n(z)\) 的一致性。故 Laurent 级数恰为 \(\sum_{n=-\infty}^\infty J_n(z) t^n\)。该生成函数蕴含 \(e^{iz\sin\theta} = \sum_{n=-\infty}^\infty J_n(z) e^{in\theta}\)（令 \(t=e^{i\theta}\)），即 Jacobi-Anger 展开。\(\blacksquare\)

\subsection{97.3 Legendre 多项式与球谐函数}\label{legendre-ux591aux9879ux5f0fux4e0eux7403ux8c10ux51fdux6570}

\textbf{定义 97.5}（Legendre 多项式的 Rodrigues 公式）：

\[P_n(x) = \frac{1}{2^n n!} \frac{d^n}{dx^n} (x^2 - 1)^n, \quad n = 0, 1, 2, \ldots\]

\(P_n\) 为 \(n\) 次多项式，\(P_n(1) = 1\)，\(P_n(-1) = (-1)^n\)。其满足 Legendre 微分方程 \((1 - x^2) P_n'' - 2x P_n' + n(n+1) P_n = 0\)。

\textbf{定理 97.5}（Legendre 多项式的正交性）：

\[\int_{-1}^1 P_m(x) P_n(x) \, dx = \frac{2}{2n+1} \delta_{mn}\]

\textbf{证明}：设 \(m \leq n\)。代入 Rodrigues 公式并分部积分 \(n\) 次：

\[\int_{-1}^1 P_m(x) P_n(x) \, dx = \frac{1}{2^n n!} \int_{-1}^1 P_m(x) \frac{d^n}{dx^n} (x^2-1)^n \, dx\]

\[= \frac{(-1)^n}{2^n n!} \int_{-1}^1 (x^2-1)^n \frac{d^n}{dx^n} P_m(x) \, dx\]

边界项消去是因为 \((x^2-1)^n\) 及其前 \(n-1\) 阶导数在 \(x = \pm 1\) 处为零。若 \(m < n\)，则 \(\frac{d^n}{dx^n} P_m(x) = 0\)，故积分为零。若 \(m = n\)，则 \(\frac{d^n}{dx^n} P_n(x) = \frac{(2n)!}{2^n n!}\)，代入得

\[\int_{-1}^1 P_n(x)^2 \, dx = \frac{(2n)!}{2^{2n} (n!)^2} \int_{-1}^1 (1 - x^2)^n \, dx\]

Beta 函数计算得右端积分为 \(2^{2n+1} (n!)^2 / ((2n+1)!)\)，代入即得 \(2/(2n+1)\)。\(\blacksquare\)

\textbf{定义 97.6}（球谐函数）：在单位球面 \(S^2 \subset \mathbb{R}^3\) 上，球谐函数 \(Y_\ell^m(\theta, \phi)\)（\(\ell \in \mathbb{N}\)，\(m = -\ell, \ldots, \ell\)）是 Laplace 算子球面部分 \(\Delta_{S^2}\) 的特征函数：

\[\Delta_{S^2} Y_\ell^m = -\ell(\ell+1) Y_\ell^m\]

其标准表达式为

\[Y_\ell^m(\theta, \phi) = (-1)^m \sqrt{\frac{2\ell+1}{4\pi} \frac{(\ell-m)!}{(\ell+m)!}} \, P_\ell^m(\cos\theta) \, e^{im\phi}\]

其中 \(P_\ell^m(x) = (1-x^2)^{m/2} \frac{d^m}{dx^m} P_\ell(x)\) 为连带 Legendre 函数。\(\{Y_\ell^m\}\) 构成 \(L^2(S^2)\) 的完备正交基（等价于 \(SO(3)\) 的不可约酉表示的矩阵系数），满足正交关系 \(\int_{S^2} Y_\ell^m \overline{Y_{\ell'}^{m'}} \, d\Omega = \delta_{\ell\ell'} \delta_{mm'}\)。

\subsection{97.4 超几何函数与合流型}\label{ux8d85ux51e0ux4f55ux51fdux6570ux4e0eux5408ux6d41ux578b}

\textbf{定义 97.7}（Gauss 超几何函数）：

\[_2F_1(a, b; c; z) = \sum_{n=0}^\infty \frac{(a)_n (b)_n}{(c)_n n!} z^n, \quad |z| < 1\]

其中 \((q)_n = q(q+1)\cdots(q+n-1)\)（\(n \geq 1\)），\((q)_0 = 1\) 为 Pochhammer 符号。参数 \(c \notin \mathbb{Z}_{\leq 0}\)，否则分母为零。\(_2F_1\) 满足超几何微分方程 \(z(1-z) y'' + [c - (a+b+1)z] y' - ab \, y = 0\)，在 \(z = 0\) 处正则。

\textbf{定理 97.6}（Gauss 求和公式）：

\[_2F_1(a, b; c; 1) = \frac{\Gamma(c) \Gamma(c - a - b)}{\Gamma(c-a) \Gamma(c-b)}, \quad \operatorname{Re}(c - a - b) > 0\]

\textbf{证明}：利用 Euler 积分表示 \(_2F_1(a,b;c;z) = \frac{\Gamma(c)}{\Gamma(b)\Gamma(c-b)} \int_0^1 t^{b-1}(1-t)^{c-b-1}(1-zt)^{-a}dt\)（\(\operatorname{Re}(c) > \operatorname{Re}(b) > 0\)）。该表示由展开 \((1-zt)^{-a} = \sum_{n=0}^\infty \frac{(a)_n}{n!}(zt)^n\) 并逐项积分得到。令 \(z = 1\)，则 \((1-t)^{-a}\) 出现在积分中：\(_2F_1(a,b;c;1) = \frac{\Gamma(c)}{\Gamma(b)\Gamma(c-b)} \int_0^1 t^{b-1}(1-t)^{c-b-1}(1-t)^{-a}dt = \frac{\Gamma(c)}{\Gamma(b)\Gamma(c-b)} \int_0^1 t^{b-1}(1-t)^{c-a-b-1}dt\)。最后的积分是 Beta 函数 \(B(b, c-a-b) = \frac{\Gamma(b)\Gamma(c-a-b)}{\Gamma(c-a)}\)。代入即得 \(_2F_1(a,b;c;1) = \frac{\Gamma(c)\Gamma(c-a-b)}{\Gamma(c-a)\Gamma(c-b)}\)。收敛条件 \(\operatorname{Re}(c-a-b) > 0\) 由 Beta 积分的收敛性保证。\(\blacksquare\)

\textbf{定义 97.8}（合流超几何函数 / Kummer 函数）：

\[_1F_1(a; c; z) = \sum_{n=0}^\infty \frac{(a)_n}{(c)_n n!} z^n, \quad z \in \mathbb{C}\]

这是 \(_2F_1(a,b;c;z/b)\) 当 \(b \to \infty\) 的极限，满足 Kummer 方程 \(z y'' + (c-z) y' - a y = 0\)。经典特例包括 \(L_n^{(\alpha)}(z) = \frac{(\alpha+1)_n}{n!} {}_1F_1(-n; \alpha+1; z)\)（Laguerre 多项式）及误差函数 \(\operatorname{erf}(z) = \frac{2z}{\sqrt{\pi}} {}_1F_1(\frac{1}{2}; \frac{3}{2}; -z^2)\)。

\textbf{定义 97.9}（Whittaker 函数）：Whittaker 方程 \(y'' + \left(-\frac{1}{4} + \frac{\kappa}{z} + \frac{\frac{1}{4} - \mu^2}{z^2}\right) y = 0\) 的两组标准解：

\[M_{\kappa,\mu}(z) = z^{\mu+\frac{1}{2}} e^{-z/2} {}_1F_1(\mu - \kappa + \tfrac{1}{2}; 1 + 2\mu; z)\]

\[W_{\kappa,\mu}(z) = \frac{\Gamma(-2\mu)}{\Gamma(\frac{1}{2} - \mu - \kappa)} M_{\kappa,\mu}(z) + \frac{\Gamma(2\mu)}{\Gamma(\frac{1}{2} + \mu - \kappa)} M_{\kappa,-\mu}(z)\]

\(M_{\kappa,\mu}\) 在 \(z = 0\) 处正则，\(W_{\kappa,\mu}\) 在 \(z \to \infty\) 时以 \(z^\kappa e^{-z/2}\) 的渐近行为衰减，后者在量子力学（Coulomb 散射）和随机过程（Bessel 过程的特征函数）中有重要应用。

\hrulefill

\emph{卷二十：微分方程终。}

\newpage

\chapter{09-微分方程与动力系统/01-微分方程/04-D模与混合型.md}\label{ux5faeux5206ux65b9ux7a0bux4e0eux52a8ux529bux7cfbux7edf01-ux5faeux5206ux65b9ux7a0b04-dux6a21ux4e0eux6df7ux5408ux578b.md}

\section{第319章：D-模理论}\label{ux7b2c319ux7ae0d-ux6a21ux7406ux8bba}

D-模理论是代数分析的核心语言，将线性偏微分方程的研究转化为环层上的模论问题。其根基是Weyl代数------非交换Noether环------及其拟凝聚层范畴。

\subsection{93.1 Weyl代数与D-模的基本定义}\label{weylux4ee3ux6570ux4e0ed-ux6a21ux7684ux57faux672cux5b9aux4e49}

\textbf{定义 93.1}（Weyl代数 \(D_n\)）：记 \(\mathbb{C}\) 上的 \(n\) 阶 Weyl 代数为

\[D_n = \mathbb{C}\langle x_1, \ldots, x_n, \partial_1, \ldots, \partial_n \rangle / \mathcal{R}\]

其中 \(\mathcal{R}\) 由以下基本对易关系生成：

\[[\partial_i, x_j] = \delta_{ij}, \quad [x_i, x_j] = 0, \quad [\partial_i, \partial_j] = 0\]

此处 \([\cdot, \cdot]\) 为交换子 \([A,B] = AB - BA\)。\(D_n\) 是（左、右）Noether 整环，无零因子，其全局同调维数为 \(n\)（Bernstein 定理）。每个元素可唯一写为有限和 \(\sum_{\alpha, \beta} c_{\alpha\beta} x^\alpha \partial^\beta\)（\(\alpha, \beta \in \mathbb{N}^n\)），其中 \(x^\alpha = x_1^{\alpha_1} \cdots x_n^{\alpha_n}\)，\(\partial^\beta = \partial_1^{\beta_1} \cdots \partial_n^{\beta_n}\)。

\textbf{定义 93.2}（\(D_X\) 的层定义）：设 \(X\) 为 \(\mathbb{C}\) 上的光滑代数簇或复流形。\(\mathcal{D}_X\) 为 \(X\) 上微分算子的层：对任意开集 \(U \subset X\)，

\[\mathcal{D}_X(U) = \left\{ \sum_{|\alpha| \leq m} a_\alpha(x) \partial^\alpha \;\middle|\; a_\alpha \in \mathcal{O}_X(U) \right\}\]

局部上 \(\mathcal{D}_X(U) \cong \mathcal{O}_X(U) \otimes_{\mathbb{C}} D_n\)（当 \(U\) 为坐标邻域时），乘法由 Leibniz 法则 \(\partial_i \cdot a = a \partial_i + \frac{\partial a}{\partial x_i}\) 决定。\(\mathcal{D}_X\) 是 \(\mathcal{O}_X\) 上的拟凝聚层，具有递增滤过 \(F_m \mathcal{D}_X = \{\text{阶} \leq m \text{ 的算子}\}\)，其伴随分次层 \(\operatorname{gr}^F \mathcal{D}_X \cong \operatorname{Sym}^\bullet_{\mathcal{O}_X} \Theta_X\)（\(\Theta_X\) 为切层），特征簇 \(\operatorname{Ch}(\mathcal{M}) \subset T^*X\) 由此定义。

\textbf{定义 93.3}（\(\mathcal{D}_X\)-模）：\(X\) 上的（左）\(\mathcal{D}_X\)-模 \(\mathcal{M}\) 是 \(\mathcal{D}_X\) 上的拟凝聚层，即对任意 \(x \in X\)，存在邻域 \(U\) 及 \(\mathcal{D}_U\)-模正合列

\[\mathcal{D}_U^{p} \to \mathcal{D}_U^{q} \to \mathcal{M}|_U \to 0\]

等价地，\(\mathcal{M}\) 是 \(\mathcal{O}_X\)-拟凝聚层且配备平坦联络 \(\nabla : \mathcal{M} \to \mathcal{M} \otimes_{\mathcal{O}_X} \Omega_X^1\) 满足 \(\nabla^2 = 0\)，其中作用由 \(\partial_i \cdot m = \nabla_{\partial_i}(m)\) 给出。若存在有限局部表现，则称 \(\mathcal{M}\) 为凝聚 \(\mathcal{D}_X\)-模。记 \(\operatorname{Mod}_{\operatorname{qc}}(\mathcal{D}_X)\) 为拟凝聚 \(\mathcal{D}_X\)-模范畴，\(\operatorname{Mod}_{\operatorname{coh}}(\mathcal{D}_X)\) 为其凝聚子范畴。基本例子：结构层 \(\mathcal{O}_X\) 即 \(\mathcal{D}_X/\mathcal{D}_X \cdot \Theta_X\)，其上 \(\partial_i\) 的作用即为通常的方向导数。

\subsection{93.2 直像与逆像函子}\label{ux76f4ux50cfux4e0eux9006ux50cfux51fdux5b50}

设 \(f : X \to Y\) 为复流形之间的态射。导出范畴 \(D^b(\mathcal{D}_X)\) 上定义两个基本函子。

\textbf{定义 93.4}（直像 \(f_+\)）：定义转移模 \(\mathcal{D}_{X \to Y} = \mathcal{O}_X \otimes_{f^{-1} \mathcal{O}_Y} f^{-1} \mathcal{D}_Y\)，其为 \((\mathcal{D}_X, f^{-1} \mathcal{D}_Y)\)-双模。对 \(\mathcal{M}^\bullet \in D^b(\mathcal{D}_X)\)，

\[f_+ \mathcal{M}^\bullet = R f_\bullet\left( \mathcal{D}_{Y \leftarrow X} \otimes_{\mathcal{D}_X}^{\mathbf{L}} \mathcal{M}^\bullet \right) \in D^b(\mathcal{D}_Y)\]

其中 \(\mathcal{D}_{Y \leftarrow X} = \Omega_X \otimes_{\mathcal{O}_X} \mathcal{D}_{X \to Y} \otimes_{f^{-1}\mathcal{O}_Y} f^{-1}\Omega_Y^{\otimes -1}\) 是调整后的转移模以确保左右模的兼容性。\(f_+\) 是 triangulated 函子。

\textbf{定义 93.5}（逆像 \(f^+\)）：对 \(\mathcal{N}^\bullet \in D^b(\mathcal{D}_Y)\)，

\[f^+ \mathcal{N}^\bullet = \mathcal{D}_{X \leftarrow Y} \otimes_{f^{-1}\mathcal{D}_Y}^{\mathbf{L}} f^{-1} \mathcal{N}^\bullet [\dim X - \dim Y] \in D^b(\mathcal{D}_X)\]

其中 \(\mathcal{D}_{X \leftarrow Y} = \mathcal{O}_X \otimes_{f^{-1}\mathcal{O}_Y} f^{-1}(\Omega_Y \otimes_{\mathcal{O}_Y} \mathcal{D}_Y)\)（适当调整后）。若 \(f\) 为光滑态射（submersion），则 \(f^+\) 对应于 PDE 意义下的拉回；若 \(f\) 为闭嵌入，则 \(f_+\) 对应于 PDE 意义下的推移。

\textbf{定理 93.1}（Kashiwara 等价）：设 \(i : Y \hookrightarrow X\) 为闭嵌入。则 \(i_+\) 诱导范畴等价

\[\operatorname{Mod}_{\operatorname{coh}}(\mathcal{D}_Y) \xrightarrow{\;\cong\;} \operatorname{Mod}_{\operatorname{coh}, Y}(\mathcal{D}_X)\]

\textbf{证明}：验证 \(i^+ \circ i_+ \cong \operatorname{id}\) 及 \(i_+ \circ i^+ \cong \operatorname{id}\) 在相应范畴上的自然同构。关键在于 \(\mathcal{D}_{Y \leftarrow X} \otimes_{\mathcal{D}_X}^{\mathbf{L}} \mathcal{D}_{X \to Y}\) 在导出意义下同构于 \(\mathcal{D}_Y\)（作为双模），利用 \(\mathcal{D}_X\) 在 \(Y\) 上的滤过结构以及 \(F_m \mathcal{D}_X\) 的局部坐标计算。具体地，在局部坐标 \((x_1,\ldots,x_n)\) 下，\(Y\) 由 \(x_1=\cdots=x_r=0\) 定义，\(\mathcal{D}_X\) 的滤过 \(F_m\mathcal{D}_X\) 在 \(Y\) 上的限制与 \(\mathcal{D}_Y\) 的滤过相容。\(i_+\) 将 \(\mathcal{D}_Y\)-模 \(\mathcal{M}\) 映为 \(\mathcal{D}_X\)-模 \(\mathcal{D}_{X\to Y}\otimes_{\mathcal{D}_Y}\mathcal{M}\)，其支集含于 \(Y\)；\(i^+\) 将支集含于 \(Y\) 的 \(\mathcal{D}_X\)-模 \(\mathcal{N}\) 映为 \(\mathcal{N}\) 中沿 \(Y\) 法向被湮灭的子模。两个函子的复合给出自然同构，从而建立范畴等价。\(\blacksquare\)

\subsection{93.3 Riemann-Hilbert 对应}\label{riemann-hilbert-ux5bf9ux5e94}

Riemann-Hilbert 对应是 D-模理论的巅峰结果，将分析的（微分方程）对象与拓扑的（可构造层）对象联系起来。

\textbf{定义 93.6}（正则全纯 D-模）：\(\mathcal{M} \in \operatorname{Mod}_{\operatorname{coh}}(\mathcal{D}_X)\) 称为正则全纯的，若沿任意光滑除子 \(D \subset X\)，\(\mathcal{M}\) 的 formal 完备化具有正则奇性------等价地，对任意全纯曲线 \(\gamma : \Delta^* \to X\)（\(\Delta^*\) 为穿孔圆盘），拉回 \(\gamma^+ \mathcal{M}\) 仅有正则奇点。记 \(\operatorname{Mod}_{\operatorname{rh}}(\mathcal{D}_X)\) 为其满子范畴，\(D^b_{\operatorname{rh}}(\mathcal{D}_X)\) 为相应的导出范畴。

\textbf{定义 93.7}（可构造层）：层 \(\mathcal{F}\)（\(\mathbb{C}\)-向量空间层）称为 \(\mathbb{C}\)-可构造的，若存在 \(X\) 的代数 Whitney 分层 \(\{X_\alpha\}\) 使得 \(\mathcal{F}|_{X_\alpha}\) 是局部常值层且茎为有限维。记 \(D^b_c(X)\) 为 \(\mathbb{C}\)-可构造层的有界导出范畴。

\textbf{定理 93.2}（Riemann-Hilbert 对应，Kashiwara 1984，Mebkhout 1984）：解函子

\[\operatorname{Sol}(-) = R\mathcal{H}om_{\mathcal{D}_X}(-, \mathcal{O}_X)\]

诱导范畴等价

\[D^b_{\operatorname{rh}}(\mathcal{D}_X) \xrightarrow{\;\cong\;} D^b_c(X)\]

其逆由 de Rham 函子 \(\operatorname{DR}(-) = \Omega_X \otimes_{\mathcal{D}_X}^{\mathbf{L}} (-)\) 给出。此等价将 \(\mathcal{D}_X\)-模的张量积对应于可构造层的 !-张量积，将直像对应于反常直像 \(Rf_!\)，将逆像对应于反常逆像 \(f^!\)。该定理的根本含义是：正则全纯线性 PDE 系统的解层是 \(\mathbb{C}\)-可构造层，且该构造是范畴一一对应的。

\textbf{证明}：解函子 \(\operatorname{Sol}(\mathcal{M}) = R\mathcal{H}om_{\mathcal{D}_X}(\mathcal{M}, \mathcal{O}_X)\) 将正则全纯 \(\mathcal{D}_X\)-模映为可构造层：在每个 Whitney 分层片上，\(\mathcal{M}\) 局部同构于 \(\mathcal{O}_X\) 的局部系统（由 Cauchy-Kowalevskaya 定理判定），其解层为局部常值层。de Rham 函子 \(\operatorname{DR}(\mathcal{F}) = \Omega_X \otimes_{\mathcal{D}_X}^{\mathbf{L}} \mathcal{F}\) 为拟逆。验证 \(\operatorname{Sol} \circ \operatorname{DR} \cong \operatorname{id}\) 和 \(\operatorname{DR} \circ \operatorname{Sol} \cong \operatorname{id}\) 需利用 Kashiwara 的构造性层理论：对正则全纯 \(\mathcal{M}\)，沿任意除子的形式完备化具有正则奇性，故 \(\operatorname{Sol}(\mathcal{M})\) 在 Whitney 分层下为 \(\mathbb{C}\)-可构造层。反方向，可构造层 \(\mathcal{F}\) 的 de Rham 复形 \(\operatorname{DR}(\mathcal{F})\) 是正则全纯 \(\mathcal{D}_X\)-模。范畴等价性进一步保持六函子形式体系，使 PDE 的代数操作与拓扑层的操作一一对应。\(\blacksquare\)

\subsection{93.4 Bernstein-Sato 多项式}\label{bernstein-sato-ux591aux9879ux5f0f}

\textbf{定理 93.3}（Bernstein-Sato 多项式的存在性，Bernstein 1971）：设 \(P(x) \in \mathbb{C}[x_1, \ldots, x_n]\) 为非零多项式。则存在非零多项式 \(b(s) \in \mathbb{C}[s]\) 及微分算子 \(D(s) \in D_n[s]\)（系数为关于 \(s\) 的多项式的微分算子）满足泛函方程

\[b(s) P(x)^s = D(s) \cdot P(x)^{s+1}\]

其中 \(P^s\) 理解为 \(P^s = e^{s \log P}\)（取适当分支）。最小次数非零的 monic \(b(s)\) 称为 \(P\) 的 Bernstein-Sato 多项式 \(b_P(s)\)。

\textbf{证明}：设 \(X = \mathbb{C}^n\)，\(Y = P^{-1}(0)\)。构造 \(\mathcal{D}_X[s]\)-模 \(\mathcal{N} = \mathcal{D}_X[s] \cdot P^s\)，它是 \(\mathcal{D}_X[s] P^s\) 的 \(s\)-参数子模。考虑 \(\mathcal{D}_X[s]\)-线性映射 \(\varphi : \mathcal{D}_X[s] \to \mathcal{N}\) 由 \(\varphi(Q(s)) = Q(s) \cdot P^s\) 定义。\(\ker(\varphi)\) 是 \(\mathcal{D}_X[s]\) 的理想 \(\mathcal{I}\)。引入 \(s\) 上的平移算子 \(\tau : s \mapsto s+1\)，则 \(D(s) \cdot P^{s+1} = (D(s) P) \cdot P^s\)。定义滤过 \(F_k\mathcal{D}_X[s] = \{\sum_{|\alpha|+j \leq k} a_{\alpha,j}(x)\partial^\alpha s^j\}\)，其伴随分次环 \(\operatorname{gr}^F\mathcal{D}_X[s] \cong \mathbb{C}[x,\xi,s]\) 为多项式环。考虑 \(\mathcal{N}\) 的滤过及特征簇 \(\operatorname{Ch}(\mathcal{N}) \subset T^*X \times \mathbb{C}\)。由 Hilbert 基定理，\(\operatorname{gr}^F(\mathcal{D}_X[s]/\mathcal{I})\) 作为 \(\operatorname{gr}^F\mathcal{D}_X[s]\)-模是有限生成的，故 \(\mathcal{D}_X[s]/\mathcal{I}\) 作为 \(\mathbb{C}[s]\)-模有限生成。由此推出存在非零 \(b(s) \in \mathbb{C}[s]\) 使得 \(b(s) \in \mathcal{I} + \mathcal{D}_X[s] \cdot P\)，即 \(b(s) P^s = D(s) \cdot P^{s+1}\)。Bernstein 进一步证明 \(b_P(s)\) 的所有根为负有理数，与 \(P\) 的奇异点（对数典范阈）有深刻关联。\(\blacksquare\)

\hrulefill

\section{第320章：特殊方程与推广}\label{ux7b2c320ux7ae0ux7279ux6b8aux65b9ux7a0bux4e0eux63a8ux5e7f}

本章将微分方程理论推广到若干具有特殊结构和重要应用的方程类型。Painleve 方程（P\({}_{\rm I}\)-P\({}_{\rm VI}\)）是六类二阶非线性常微分方程，其核心特征是 Painleve 性质------所有可移动奇点仅为极点------这使其成为可积系统的核心研究对象，与完全可积 PDE（如 KdV 方程）、随机矩阵理论（Tracy-Widom 分布）和等单值形变理论有深刻联系。泛函微分方程（特别是时滞微分方程）将状态空间从有限维推广到无穷维函数空间 \(C([-\tau, 0], \mathbb{R}^n)\)，其稳定性分析依赖 Lyapunov-Krasovskii 泛函方法，在种群动力学、神经网络和控制理论中有广泛应用。这些推广类型共同揭示了微分方程理论在超越经典 ODE/PDE 框架后的丰富结构。

\subsection{94.1 Painleve方程}\label{painleveux65b9ux7a0b}

\textbf{定义 94.1}（Painleve 性质与 Painleve 方程）：二阶常微分方程 \(y'' = F(y', y, x)\)（\(F\) 关于 \(y, y'\) 有理，关于 \(x\) 解析）称为具有 \textbf{Painleve 性质}，若其所有\textbf{可移动奇点}（即位置依赖积分常数的奇点）均为极点（无本性奇点或分支点）。Painleve（1900）与 Gambier 系统分类了所有满足此性质的方程，除经典可积方程外恰有六族，记为 P\({}_{\rm I}\)--P\({}_{\rm VI}\)：

\[\begin{aligned}
\text{P}_{\rm I}&: \; y'' = 6y^2 + x \\
\text{P}_{\rm II}&: \; y'' = 2y^3 + xy + \alpha \\
\text{P}_{\rm III}&: \; y'' = \frac{(y')^2}{y} - \frac{y'}{x} + \frac{\alpha y^2 + \beta}{x} + \gamma y^3 + \frac{\delta}{y} \\
\text{P}_{\rm IV}&: \; y'' = \frac{(y')^2}{2y} + \frac{3}{2}y^3 + 4xy^2 + 2(x^2 - \alpha)y + \frac{\beta}{y} \\
\text{P}_{\rm V}&: \; y'' = \left(\frac{1}{2y} + \frac{1}{y-1}\right)(y')^2 - \frac{y'}{x} + \frac{(y-1)^2}{x^2}\left(\alpha y + \frac{\beta}{y}\right) + \gamma\frac{y}{x} + \delta\frac{y(y+1)}{y-1} \\
\text{P}_{\rm VI}&: \; y'' = \frac{1}{2}\left(\frac{1}{y} + \frac{1}{y-1} + \frac{1}{y-x}\right)(y')^2 - \left(\frac{1}{x} + \frac{1}{x-1} + \frac{1}{y-x}\right)y' \\ &\qquad + \frac{y(y-1)(y-x)}{x^2(x-1)^2}\left(\alpha + \beta\frac{x}{y^2} + \gamma\frac{x-1}{(y-1)^2} + \delta\frac{x(x-1)}{(y-x)^2}\right)
\end{aligned}\]

其中 \(\alpha, \beta, \gamma, \delta \in \mathbb{C}\) 为复参数。P\({}_{\rm VI}\) 是最一般的 Painleve 方程，通过合流极限可导出其余五族。

\textbf{定理 94.1}（Painleve 方程的 Hamiltonian 结构）：每族 Painleve 方程均可写为 Hamiltonian 系统

\[\frac{dq}{dx} = \frac{\partial H_J}{\partial p}, \quad \frac{dp}{dx} = -\frac{\partial H_J}{\partial q} \quad (J = \text{I}, \ldots, \text{VI})\]

其中 \(H_J = H_J(q, p, x; \alpha, \beta, \gamma, \delta)\) 是关于 \((q, p)\) 的有理函数。例如 P\({}_{\rm II}\) 的 Hamiltonian 为

\[H_{\rm II}(q, p, x) = \frac{1}{2}p^2 - \left(q^3 + \frac{x}{2}q\right) - \left(\alpha + \frac{1}{2}\right)q\]

消去 \(p\) 即得 P\({}_{\rm II}\) 的标量形式。该 Hamiltonian 结构使 Painleve 方程纳入完全可积系统的框架，其 Bäcklund 变换对应 Hamiltonian 的规范变换。

\textbf{定理 94.2}（Painleve 合流瀑布 / Coalescence Cascade）：在适当的参数和变量缩放极限下，Painleve 方程按如下方向合流：

\[\text{P}_{\rm VI} \to \text{P}_{\rm V} \to \text{P}_{\rm IV} \to \text{P}_{\rm II} \to \text{P}_{\rm I}\]

\[\text{P}_{\rm V} \to \text{P}_{\rm III}\]

这对应于 Painleve 方程奇点结构的合并（coalescence）。例如 P\({}_{\rm VI}\) 具有 4 个固定奇点 \(\{0, 1, x, \infty\}\)，P\({}_{\rm V}\) 合并 \(x \to \infty\) 具有 \(\{0, \infty\}\) 和合并不规则奇点，P\({}_{\rm III}\) 具有两个不规则奇点 \(\{0, \infty\}\) 等。

\textbf{定理 94.3}（Painleve 方程与 Tracy-Widom 分布）：设 \(F_2(s)\) 为 Gaussian Unitary Ensemble (GUE) 的最大特征值的 Tracy-Widom 分布函数。则

\[F_2(s) = \exp\left(-\int_s^\infty (t-s) u^2(t) dt\right)\]

其中 \(u(t)\) 是 P\({}_{\rm II}\) 方程（\(\alpha = 0\)）满足 Hastings-McLeod 边界条件 \(u(t) \sim \operatorname{Ai}(t)\)（\(t \to +\infty\)）和 \(u(t) \sim \sqrt{-t/2}\)（\(t \to -\infty\)）的唯一解。此深刻联系揭示了 Painleve 方程在随机矩阵理论和统计物理中的核心地位。

\subsection{94.2 泛函微分方程}\label{ux6cdbux51fdux5faeux5206ux65b9ux7a0b}

泛函微分方程（Functional Differential Equations, FDE）是导数依赖于过去（甚至未来）状态的微分方程。其状态空间为无穷维函数空间，与 ODE 有限维动力学有本质区别。

\textbf{定义 94.2}（时滞微分方程与状态空间）：设 \(\tau > 0\) 为最大时滞。\textbf{时滞微分方程}（Delay Differential Equation, DDE）的一般形式为

\[\dot{x}(t) = f(t, x_t), \quad x_t \in C([-\tau, 0], \mathbb{R}^n)\]

其中 \(x_t(\theta) = x(t + \theta)\)（\(\theta \in [-\tau, 0]\)）是解在区间 \([t-\tau, t]\) 上的历史片段，\(f : \mathbb{R} \times C([-\tau, 0], \mathbb{R}^n) \to \mathbb{R}^n\)。最简单的常时滞方程为 \(\dot{x}(t) = f(t, x(t), x(t-\tau))\)。

\textbf{定义 94.3}（步进法 / Method of Steps）：对于 DDE \(\dot{x}(t) = f(t, x(t), x(t-\tau))\) 配以初始函数 \(\varphi \in C([-\tau, 0], \mathbb{R}^n)\)（即 \(x(\theta) = \varphi(\theta)\)，\(\theta \in [-\tau, 0]\)），解的存在唯一性通过\textbf{步进法}建立。将时间轴划分为长度为 \(\tau\) 的区间 \(I_k = [k\tau, (k+1)\tau]\)（\(k = 0, 1, 2, \ldots\)）。在 \(I_0\) 上，\(x(t-\tau) = \varphi(t-\tau)\) 已知，方程退化为 ODE：\(\dot{x}(t) = f(t, x(t), \varphi(t-\tau))\)（\(t \in [0, \tau]\)），由 Picard-Lindelöf 定理存在唯一解。在 \(I_k\) 上，利用 \(I_{k-1}\) 上已求得的解作为历史，递推求解。此方法证明了：若 \(f\) 关于第二、第三变量是 \(C^1\) 的，则 DDE 对任意 \(\varphi \in C([-\tau, 0], \mathbb{R}^n)\) 存在唯一的全局解。

\textbf{定义 94.4}（Lyapunov-Krasovskii 泛函）：DDE 的稳定性分析推广了 Lyapunov 直接法。\textbf{Lyapunov-Krasovskii 泛函} \(V : C([-\tau, 0], \mathbb{R}^n) \to \mathbb{R}\) 满足：(i) \(V(0) = 0\)，\(V(\varphi) > 0\)（\(\varphi \neq 0\)）；(ii) \(V\) 连续且在原点有正定下界。沿解的导数为

\[\dot{V}(\varphi) = \limsup_{h \to 0^+} \frac{V(x_{t+h}(\varphi)) - V(\varphi)}{h}\]

其中 \(x_t(\varphi)\) 是以 \(\varphi\) 为初始函数的解在 \(t\) 时刻的历史片段。若 \(\dot{V}(\varphi) \leq -w(|\varphi(0)|)\)（\(w\) 为正定函数），则零解一致渐近稳定。

\textbf{定理 94.4}（Lyapunov-Krasovskii 稳定性定理）：设 \(V : C([-\tau, 0], \mathbb{R}^n) \to \mathbb{R}\) 为连续泛函，且存在连续非减函数 \(u, v, w : \mathbb{R}_+ \to \mathbb{R}_+\) 满足 \(u(0) = v(0) = 0\)，\(u(s), v(s) > 0\)（\(s > 0\)），使得

\[u(|\varphi(0)|) \leq V(\varphi) \leq v(\|\varphi\|_C), \quad \dot{V}(\varphi) \leq -w(|\varphi(0)|)\]

则 DDE 的零解是一致渐近稳定的。此定理将有限维 Lyapunov 定理推广到无穷维状态空间。

\textbf{定理 94.5}（线性 DDE 的谱理论）：考虑线性常时滞自治方程

\[\dot{x}(t) = A x(t) + B x(t - \tau)\]

其中 \(A, B \in \mathbb{R}^{n \times n}\)。其\textbf{特征方程}为

\[\det(\lambda I - A - B e^{-\lambda\tau}) = 0\]

这是超越方程，通常有无穷多个根 \(\{\lambda_k\}_{k=1}^\infty\)（满足 \(\operatorname{Re} \lambda_k \to -\infty\)）。零解渐近稳定的充要条件是所有特征根满足 \(\operatorname{Re} \lambda_k < 0\)。与 ODE 关键不同：即使 \(A+B\) 的所有特征值均位于左半平面，当 \(\tau\) 足够大时仍可能出现 \(\operatorname{Re} \lambda_k > 0\)（稳定性切换现象 / stability switch）。

\hrulefill

\section{第321章：混合型偏微分方程}\label{ux7b2c321ux7ae0ux6df7ux5408ux578bux504fux5faeux5206ux65b9ux7a0b}

混合型偏微分方程的特征形式在区域的不同部分具有不同的类型（椭圆型、双曲型或抛物型），这是跨音速空气动力学、等离子体物理和几何分析中自然出现的一类重要方程。最经典的例子是 Tricomi 方程 \(y u_{xx} + u_{yy} = 0\)------在 \(y > 0\) 的半平面为椭圆型，在 \(y < 0\) 的半平面为双曲型，在 \(y = 0\) 的直线上抛物退化。本章将系统介绍混合型方程的分类（Tricomi 型、Keldysh 型）、边值问题的适定性理论（Frankl 问题、Gellerstedt 问题），以及混合型方程与跨音速流动（Chaplygin 气体动力学方程）和等距嵌入问题（Weingarten 曲面）之间的深刻联系。

\subsection{95.1 混合型方程的正式定义与分类}\label{ux6df7ux5408ux578bux65b9ux7a0bux7684ux6b63ux5f0fux5b9aux4e49ux4e0eux5206ux7c7b}

\textbf{定义 95.1}（混合型偏微分方程）：设 \(\Omega \subset \mathbb{R}^n\) 为区域，二阶线性偏微分算子

\[L[u] = \sum_{i,j=1}^{n} a_{ij}(x) \frac{\partial^2 u}{\partial x_i \partial x_j} + \sum_{i=1}^{n} b_i(x) \frac{\partial u}{\partial x_i} + c(x) u = f(x)\]

其系数矩阵 \(A(x) = (a_{ij}(x))_{i,j=1}^n\) 对称。若 \(A(x)\) 在 \(\Omega\) 的不同子区域上具有不同的惯性指标（正、负、零特征值的个数），则称 \(L\) 为 \(\Omega\) 上的混合型算子。\(A(x)\) 的特征值决定局部类型：全正→椭圆型，一个负其余正→双曲型，存在零特征值→抛物退化。

\textbf{定义 95.2}（Tricomi 方程）：经典混合型方程为

\[T[u] = y \frac{\partial^2 u}{\partial x^2} + \frac{\partial^2 u}{\partial y^2} = 0, \quad (x, y) \in \Omega \subset \mathbb{R}^2\]

判别式 \(\Delta(x,y) = -y\)，故 \(y > 0\) 时椭圆型（\(A(x,y) = \operatorname{diag}(y, 1) \succ 0\)）；\(y < 0\) 时双曲型（一个正、一个负特征值）；\(y = 0\) 时抛物退化。Tricomi 方程是跨音速流动的线性化模型（Chaplygin 气体动力学方程经速度图变换后约化得到）。

\textbf{定义 95.3}（一般混合型算子分类）： - \textbf{Tricomi 型}：特征形式沿一个变量变号，形如 \(K(y) u_{xx} + u_{yy} = 0\)，\(K(y)\) 过零。 - \textbf{Keldysh 型}：退化方向不同，形如 \(u_{xx} + y^m u_{yy} = 0\)（\(m\) 为奇数），\(y > 0\) 时椭圆，\(y < 0\) 时双曲。 - \textbf{广义混合型}：\(\operatorname{sgn}(x) |x|^m u_{xx} + u_{yy} = 0\)（\(m > 0\)），在退化线上方程阶数非退化。

\subsection{95.2 混合型方程的边界值问题}\label{ux6df7ux5408ux578bux65b9ux7a0bux7684ux8fb9ux754cux503cux95eeux9898}

\textbf{定义 95.4}（Tricomi 问题）：设 \(\Omega\) 为分段光滑单连通区域，其边界 \(\partial \Omega = \sigma \cup \Gamma_1 \cup \Gamma_2\)，其中 \(\sigma\) 为椭圆区域 \(y > 0\) 中的曲线，\(\Gamma_1\) 和 \(\Gamma_2\) 为双曲区域 \(y < 0\) 中自退化线 \(y = 0\) 上的点 \((0,0)\) 出发的两条特征线。Tricomi 问题：

\[\begin{cases} T[u] = 0, & (x,y) \in \Omega \\ u|_{\sigma} = \varphi, & \text{（椭圆边界上的 Dirichlet 条件）} \\ u|_{\Gamma_1} = \psi, & \text{（一条特征线上的条件）} \end{cases}\]

\textbf{定理 95.1}（Tricomi 问题的适定性）：当 \(\sigma\) 是 \(y > 0\) 半平面中以 \((0,0)\) 和另一点为端点的适当曲线时，Tricomi 问题存在唯一弱解。若 \(\sigma\) 满足 \(\sigma = \{(x,y) : x^2 + 4y^3/9 = R^2, y \geq 0\}\)（法区域的正规曲线），则问题强适定。

\textbf{证明}：在椭圆区域 \(y>0\) 中，Tricomi 算子 \(T = y\partial_x^2 + \partial_y^2\) 是椭圆型算子，由 Lax-Milgram 定理在适当的加权 Sobolev 空间中解存在唯一。在双曲区域 \(y<0\) 中，特征线为 \(dx = \pm \sqrt{-y}dy\)，积分得 \(x \pm \frac{2}{3}(-y)^{3/2} = \text{const}\)。对 Tricomi 问题，在 \(\Gamma_1\) 上给定 Dirichlet 数据，沿 \(\Gamma_2\) 特征线不需指定数据（由双曲方程的 Cauchy 问题理论自然确定）。通过将椭圆区域的 Dirichlet 问题与双曲区域的 Cauchy 问题在退化线 \(y=0\) 上进行匹配，利用 \(u\) 和 \(u_y\) 在 \(y=0\) 上的连续性条件，得到退化线上积分方程。当 \(\sigma\) 为正规曲线时，该积分方程是 Fredholm 型的且可逆，故解存在唯一。\(\blacksquare\)

\textbf{定义 95.5}（Frankl 问题）：推广 Tricomi 问题。设 \(\Omega\) 为混合型区域，\(\partial \Omega = \sigma \cup \Gamma_1 \cup \Gamma_2 \cup \gamma\)，其中 \(\gamma\) 为退化线上的一段。Frankl 问题要求：在 \(\sigma\) 上给 Dirichlet 条件，在 \(\Gamma_1\) 上给 Dirichlet 条件，在 \(\Gamma_2\) 上不给条件（或给非局部条件），在 \(\gamma\) 上不给任何条件（退化线自然携带信息）。该问题源于跨音速喷管设计------流动在喷管喉部从亚音速过渡到超音速时，退化线上的自然连续性条件自动满足。

\textbf{定理 95.2}（Frankl 问题的存在性，Frankl 1945）：在适当的几何条件下，Frankl 问题存在弱解。其核心困难在于退化线上的数据不准许任意指定，必须满足一定的相容性条件，可通过奇异积分方程的 Fredholm 理论处理。

\textbf{证明}：Frankl 问题的核心是将椭圆、双曲和退化线三个区域耦合。在退化线 \(\gamma\) 上，方程 \(T[u]=0\) 退化为 \(u_{yy}=0\)（因 \(y=0\) 时 \(y u_{xx}=0\)），故 \(u\) 在 \(\gamma\) 上为 \(y\) 的线性函数，连续性条件自动满足。在椭圆区域 \(y>0\) 中，给定 \(\sigma\) 上的 Dirichlet 数据和 \(\gamma\) 上的 \(u\) 值，由椭圆方程 Dirichlet 问题的适定性，解唯一确定。在双曲区域 \(y<0\) 中，给定 \(\Gamma_1\) 上的 Dirichlet 数据和 \(\gamma\) 上的 \(u\) 值，由双曲方程 Cauchy 问题的适定性，解唯一确定。匹配条件：椭圆区域和双曲区域在 \(\gamma\) 上的法向导数 \(\partial u/\partial y\) 必须一致。这给出 \(\gamma\) 上的奇异积分方程，其主象征为 \(|\xi|\operatorname{sgn}(\xi)\)。由 Fredholm 理论，该算子在其适当的 Sobolev 空间中为指标零的 Fredholm 算子，故存在弱解。\(\blacksquare\)

\subsection{95.3 混合型方程的适定性框架}\label{ux6df7ux5408ux578bux65b9ux7a0bux7684ux9002ux5b9aux6027ux6846ux67b6}

\textbf{定义 95.6}（适定性提法的一般框架）：混合型方程 \(L[u] = f\) 的适定边值问题遵循以下原则： 1. \textbf{椭圆部分}（\(A(x) \succ 0\)）：指定 Dirichlet 或 Neumann 条件的一个。 2. \textbf{双曲部分}（\(A(x)\) 具有 Lorentz 符号）：Cauchy 数据仅可在一条类空曲面上给出；适定问题通常只需指定部分边界数据（如 Tricomi 问题的单特征线条件）。 3. \textbf{退化线}：作为椭圆型与双曲型区域的分界，其上通常不独立指定边界条件；解的连续性及法向导数的跳跃性由方程本身和两侧的条件联合确定。

\textbf{定理 95.3}（Morawetz 乘子法，1956）：对 Tricomi 方程 \(T[u] = f\)，存在乘子型恒等式

\[\iint_\Omega (2x u_x + u) T[u] \, dx dy = \oint_{\partial \Omega} Q(u_x, u_y, u, x, y) \, d\ell + \iint_\Omega (u_x^2 + u_y^2) \, dx dy\]

由此可得 Tricomi 问题解的唯一性及 \(H^1\) 先验估计。乘子 \((2x u_x + u)\) 的选择利用了 Tricomi 算子在退化线两侧符号变化的对称性。对于乘子法无法覆盖的情形，Lupo 与 Payne（2004）使用变分结构将非线性混合型方程嵌入临界点理论，在 \(H_0^1(\Omega)\) 的适当闭子空间上应用山路引理获得弱解。

\textbf{证明}：以乘子 \(M[u] = 2x u_x + u\) 作用于 Tricomi 方程 \(T[u] = y u_{xx} + u_{yy}\)。计算 \(\iint_\Omega M[u] T[u] dxdy\)，分部积分得：\(\iint_\Omega (2x u_x + u)(y u_{xx} + u_{yy}) dxdy = \oint_{\partial\Omega} Q d\ell + \iint_\Omega \mathcal{E}(u_x, u_y) dxdy\)，其中 \(\mathcal{E} = u_x^2 + u_y^2 + 2xy u_x u_{xx} + \cdots\)。经仔细的分部积分和边界项整理，最终得到 \(\mathcal{E} = u_x^2 + u_y^2\)（正定二次型）。边界项 \(Q\) 沿 \(\sigma\) 和 \(\Gamma_1\) 可由已知边界数据控制。若 \(f=0\) 且边界数据为零，则 \(\iint_\Omega (u_x^2+u_y^2)dxdy=0\)，故 \(u=\text{const}\)，由边界条件得 \(u=0\)，证得唯一性。对非齐次情形，上述恒等式给出 \(\|u\|_{H^1}^2 \leq C\|f\|_{L^2}^2\) 的先验估计。\(\blacksquare\)

\textbf{定理 95.4}（弱解的 Fredholm 型理论）：将 Tricomi 问题化为边界上的奇异积分方程，主象征为 \(a(\xi) = |\xi| \operatorname{sgn}(\xi)\)，对应的奇异积分算子 \(H\) 满足 \(H^2 = I\)（Hilbert 变换的推广）。适定性等价于相应的 Wiener-Hopf 型算子在 Sobolev 空间中为 Fredholm 算子且指数为零。这体现了混合型方程与奇异积分方程理论、拟微分算子理论的深刻联系。

\textbf{证明}：将 Tricomi 问题通过 Fourier 变换在退化线 \(y=0\) 上化为边界积分方程。设 \(u(x,0) = \varphi(x)\)，\(u_y(x,0) = \psi(x)\) 为退化线上的未知数据。在椭圆区域 \(y>0\) 中，由 Poisson 核表示 \(u\) 在 \(\sigma\) 上的值，得到 \(\varphi\) 与 \(\psi\) 的第一个关系式。在双曲区域 \(y<0\) 中，由 D'Alembert 公式和特征线 \(\Gamma_1\) 上的数据，得到 \(\varphi\) 与 \(\psi\) 的第二个关系式。消去 \(\psi\) 后得到关于 \(\varphi\) 的奇异积分方程：\(A\varphi(x) + \frac{1}{\pi}\int_{\mathbb{R}} \frac{\varphi(t)}{t-x}dt = g(x)\)（\(A\) 为常数）。该方程的主象征为 \(a(\xi) = A + i\operatorname{sgn}(\xi)\)，对应的拟微分算子为 \(A I + iH\)（\(H\) 为 Hilbert 变换）。在适当的 Sobolev 空间 \(H^s\) 中，\(AI + iH\) 是 Fredholm 算子当且仅当 \(A \neq \pm 1\)。当 \(A \neq \pm 1\) 时，算子可逆，Tricomi 问题适定；当 \(A = \pm 1\) 时出现共振，对应非 Fredholm 情形。\(\blacksquare\)

\hrulefill

\emph{卷二十：微分方程终。}

\newpage

\chapter{09-微分方程与动力系统/02-动力系统/01-离散与连续.md}\label{ux5faeux5206ux65b9ux7a0bux4e0eux52a8ux529bux7cfbux7edf02-ux52a8ux529bux7cfbux7edf01-ux79bbux6563ux4e0eux8fdeux7eed.md}

\section{第322章：连续与离散动力系统}\label{ux7b2c322ux7ae0ux8fdeux7eedux4e0eux79bbux6563ux52a8ux529bux7cfbux7edf}

本章严格建立连续与离散动力系统的基本语言和拓扑结构，是遍历论与分支理论的公共基础。

\subsection{99.1 连续动力系统与离散动力系统}\label{ux8fdeux7eedux52a8ux529bux7cfbux7edfux4e0eux79bbux6563ux52a8ux529bux7cfbux7edf}

\textbf{定义 99.1}（\(C^r\) 流）：设 \(X\) 为光滑流形。\(C^r\) 流（\(r \geq 1\)）是一个 \(C^r\) 映射 \(\varphi : \mathbb{R} \times X \to X\)，记为 \(\varphi_t(x) = \varphi(t, x)\)，满足： 1. \(\varphi_0 = \operatorname{id}_X\)（恒等映射）； 2. \(\varphi_t \circ \varphi_s = \varphi_{t+s}\)，对任意 \(t, s \in \mathbb{R}\)（群性质）。

若仅对 \(t \geq 0\) 有定义且满足半群性质，则称 \(\varphi\) 为 \(C^r\) 半流。给定自治 ODE \(\dot{x} = f(x)\)（\(f \in C^r(X, TX)\) 为向量场），若其解对初值全局存在，则解映射 \(\varphi_t(x_0) = x(t; x_0)\) 定义了唯一的 \(C^r\) 流。\(f\) 称为流 \(\varphi\) 的无穷小生成元：\(f(x) = \frac{d}{dt}\big|_{t=0} \varphi_t(x)\)。

\textbf{定义 99.2}（离散动力系统）：设 \(X\) 为拓扑空间，\(f : X \to X\) 为连续映射（或同胚）。正向轨道为 \(\{f^n(x)\}_{n=0}^{\infty}\)（\(f^0 = \operatorname{id}\), \(f^{n+1} = f \circ f^n\)）。若 \(f\) 为同胚，则双向轨道为 \(\{f^n(x)\}_{n \in \mathbb{Z}}\)。半群 \(\{f^n\}_{n \geq 0}\) 描述了不可逆动力系统，群 \(\{f^n\}_{n \in \mathbb{Z}}\) 描述了可逆动力系统。连续与离散之间的核心桥梁为时间-1 映射。

\textbf{定义 99.3}（时间-1 映射）：设 \(\varphi_t\) 为 \(X\) 上的 \(C^r\) 流。其时间-1 映射定义为 \(f(x) = \varphi_1(x)\)。若流全局存在，则 \(f\) 是 \(X\) 的 \(C^r\) 微分同胚，且 \(\varphi_n(x) = f^n(x)\)（对整数 \(n\)）。\(f\) 的轨道包含流的整数时刻快照，但丢失了流的连续中间状态。

\textbf{定义 99.4}（Poincaré 截面与首次回返映射）：设 \(\varphi_t\) 为流形 \(M\) 上的 \(C^r\) 流，\(\Sigma \subset M\) 为余维 1 的 \(C^r\) 嵌入子流形（截面），满足：（i）\(\varphi_t\) 横截 \(\Sigma\)（即 \(f(x) \notin T_x \Sigma\) 对 \(x \in \Sigma\)）；（ii）对任意 \(x \in \Sigma\)，存在最小 \(t = \tau(x) > 0\) 使得 \(\varphi_{\tau(x)}(x) \in \Sigma\)。则 \textbf{Poincaré 首次回返映射} \(P : \Sigma \to \Sigma\) 定义为 \(P(x) = \varphi_{\tau(x)}(x)\)。\(P\) 将连续流约化为截面上的离散动力系统，是分析周期轨道稳定性的核心工具。

\subsection{99.2 Poincaré-Bendixson 定理}\label{poincaruxe9-bendixson-ux5b9aux7406}

\textbf{定理 99.1}（Poincaré-Bendixson 定理）：设 \(\varphi_t\) 为平面 \(\mathbb{R}^2\) 上的 \(C^1\) 流，仅有有限个平衡点。若正向轨道 \(\gamma^+(x_0) = \{\varphi_t(x_0) : t \geq 0\}\) 包含于某紧集 \(K \subset \mathbb{R}^2\) 中，则其 \(\omega\)-极限集 \(\omega(x_0)\) 必为以下三者之一：（i）单个平衡点；（ii）单个周期轨；（iii）由平衡点及其间的异宿轨连接组成的有限连通图（异宿圈）。

\textbf{证明}：设 \(\omega = \omega(x_0)\)，紧致且连通（因轨道预紧）。若 \(\omega\) 不含平衡点，取 \(y \in \omega\)。作过 \(y\) 横截于流的线段 \(\ell\)（局部截面）。因 \(\omega\) 非空且紧，轨道 \(\gamma^+(x_0)\) 必无限次穿越 \(\ell\)。由于平面流无自交（Jordan 曲线定理），\(\omega \cap \ell\) 单调穿越 \(\ell\)，故必为单点 \(\{y\}\)。但 \(y \in \omega\) 意味着 \(\gamma(y) \subset \omega\)。因 \(\omega\) 紧，\(\gamma(y)\) 的闭包含于 \(\omega\)。若 \(\gamma(y)\) 非周期，则 \(\omega(y)\) 与 \(\ell\) 的交点亦单调，与单点矛盾。故 \(y\) 必为周期点，从而 \(\gamma(y)\) 为周期轨 --- 且为极限轨，即 \(\omega = \gamma(y)\)。若 \(\omega\) 包含平衡点，则仅有有限个。\(\omega\) 若含平衡点和正则点，由紧性及截面的单调穿越性，正则点必通过异宿轨连接平衡点形成异宿圈。\(\blacksquare\)

\textbf{推论 99.1}：平面 \(C^1\) 流不存在奇异吸引子（即非平衡点、非周期轨的极限集）。此即混沌不可能发生在平面自治 ODE 中的拓扑根源。

Poincaré-Bendixson 定理是二维动力系统理论的基石，其证明依赖于平面拓扑的独特性质------Jordan 曲线定理。在三维及以上空间中，此定理不再成立，这是混沌动力学可以在高维自治系统中出现的根本原因。定理的条件中''仅有有限个平衡点''是必要的：若平衡点有无穷多个（如沿一曲线分布），则 \(\omega\)-极限集可能更复杂。该定理也有重要的几何推论：若平面区域 \(D\) 包含一个正向不变紧集且不含平衡点，则 \(D\) 内至少存在一个周期轨道------这是证明周期解存在的标准方法，广泛应用于生物数学（如捕食-被捕食系统的极限环存在性）。此外，Poincaré-Bendixson 定理可推广到二维流形（如球面 \(S^2\) 和射影平面 \(\mathbb{RP}^2\)），但在环面 \(\mathbb{T}^2\) 上需要额外条件（因为环面支持无平衡点的无理流）。在动力系统理论的发展史上，Poincaré 在 1880 年代对平面多项式向量场的定性研究开辟了现代动力系统理论，而 Bendixson 在 1901 年给出了严格的证明，两人的工作共同奠定了定性理论的基础。

\subsection{99.3 环面上的无理流}\label{ux73afux9762ux4e0aux7684ux65e0ux7406ux6d41}

\textbf{定义 99.5}（环面上的 Kronecker 流）：在二维环面 \(\mathbb{T}^2 = \mathbb{R}^2 / \mathbb{Z}^2\) 上定义线性流

\[\theta_t(x, y) = (x + t, \; y + \alpha t) \mod 1\]

其中 \(\alpha \in \mathbb{R}\) 为频率比。轨道为 \(\mathbb{T}^2\) 上斜率为 \(\alpha\) 的直线（模掉整数平移）。

\textbf{定理 99.2}（Kronecker 流的遍历二分）：若 \(\alpha \in \mathbb{Q}\)，则每条轨道均为周期轨（周期为分母）。若 \(\alpha \notin \mathbb{Q}\)，则每条轨道在 \(\mathbb{T}^2\) 中稠密，且流关于 Lebesgue 测度是唯一遍历的（即 Lebesgue 测度是唯一不变的 Borel 概率测度）。Weyl 均分定理给出定量版本：

\[\lim_{T \to \infty} \frac{1}{T} \int_0^T \chi_A(\theta_t(x_0))\, dt = \operatorname{Leb}(A), \quad \forall A \subset \mathbb{T}^2 \text{ Borel}, \;\forall x_0 \in \mathbb{T}^2\]

该流是连续时间遍历（但不是混合）的典型例子。

\textbf{证明}：若 \(\alpha = p/q \in \mathbb{Q}\)（既约），则 \(\theta_t(x) = (x_1 + t, x_2 + \alpha t) \bmod 1\)。当 \(t = q\) 时，\(\theta_q(x) = (x_1+q, x_2+p) = (x_1, x_2) \bmod 1\)，故每条轨道周期为 \(q\)。若 \(\alpha \notin \mathbb{Q}\)，证明轨道稠密性等同于证明 \(\{n\alpha \bmod 1 : n \in \mathbb{Z}\}\) 在 \(\mathbb{T}^1\) 中稠密。由 Dirichlet 逼近定理，存在无穷多有理数 \(p/q\) 满足 \(|\alpha - p/q| < 1/q^2\)。设 \(d = |q\alpha - p|\)，则 \(0 < d < 1/q\)。轨道 \(\{k d \bmod 1 : k = 0,1,\ldots\}\) 以步长 \(< 1/q\) 遍历 \(\mathbb{T}^1\)，故对任意 \(\varepsilon > 0\) 取 \(q > 1/\varepsilon\) 即得稠密性。Lebesgue 测度 \(m\) 是平移不变的，且由 Fourier 分析：若 \(\mu\) 为不变测度，其 Fourier 系数 \(\hat{\mu}(k_1,k_2) = \int e^{-2\pi i(k_1 x_1+k_2 x_2)}d\mu\) 满足 \(\hat{\mu}(k_1,k_2) = e^{2\pi i(k_1+k_2\alpha)t}\hat{\mu}(k_1,k_2)\)。因 \(\alpha\) 无理，当 \((k_1,k_2)\neq(0,0)\) 时 \(k_1+k_2\alpha \neq 0\)，故 \(\hat{\mu}(k_1,k_2)=0\)，即 \(\mu=m\)（唯一遍历）。\(\blacksquare\)

\subsection{99.4 Smale 马蹄}\label{smale-ux9a6cux8e44}

\textbf{定义 99.6}（Smale 马蹄的几何构造，1967）：在平面 \(\mathbb{R}^2\) 上取正方形 \(S = [0,1] \times [0,1]\)。定义映射 \(f : S \to \mathbb{R}^2\) 的组合：

\begin{enumerate}
\def\labelenumi{\arabic{enumi}.}
\tightlist
\item
  \textbf{拉伸}：沿 \(y\) 方向均匀拉伸 \(S\) 为 \([0,1] \times [0, \mu]\)，\(\mu > 2\)。
\item
  \textbf{压缩}：沿 \(x\) 方向均匀压缩为 \([0, \lambda] \times [0, \mu]\)，\(0 < \lambda < 1/2\)。
\item
  \textbf{折叠并放回}：将长条弯折成马蹄形并放回到包含 \(S\) 的区域中，使得 \(V_0 = f(S) \cap S\) 和 \(V_1 = f(S) \cap S\) 为两个竖直带状区域（分别对应原正方形左侧和右侧的像）。
\end{enumerate}

\(f\) 在 \(S\) 上的\textbf{不变集}定义为 \(\Lambda = \bigcap_{n=-\infty}^{\infty} f^{-n}(S)\)（若 \(f\) 延拓为 \(\mathbb{R}^2\) 上的同胚）。

\textbf{定理 99.3}（马蹄的不变集性质）：\(\Lambda\) 同胚于 Cantor 集。\(f|_{\Lambda}\) 拓扑共轭于双边二符号移位 \(\sigma : \Sigma_2 \to \Sigma_2\)（\(\Sigma_2 = \{0,1\}^{\mathbb{Z}}\)）。特别地： - \(\Lambda\) 中有可数无穷个周期轨（任意周期）； - \(\Lambda\) 中含有稠密轨道（拓扑传递）； - 周期点集在 \(\Lambda\) 中稠密； - \(f|_{\Lambda}\) 对初始条件敏感依赖（混沌三要素均满足）。

\textbf{证明}：构造编码映射 \(h : \Lambda \to \Sigma_2\)，令 \(h(x)_n = j\) 若 \(f^n(x) \in V_j\)（\(j \in \{0,1\}\)，\(V_0\) 和 \(V_1\) 为两个竖直带状区域）。由马蹄的 Markov 分割性质，\(V_0\) 和 \(V_1\) 完全横截于 \(S\)：\(f(V_0)\) 和 \(f(V_1)\) 横跨 \(S\)，故 \(f^{-1}(S) \cap S = V_0 \cup V_1\)。递推得 \(f^{-n}(S) \cap S\) 为 \(2^{2n}\) 个正方形，极限集 \(\Lambda\) 为这些正方形的交集，每个点 \(x\) 对应唯一的双向二进制序列。\(h\) 为同胚：连续性因乘积拓扑与图的度量相容；双射性由 Markov 分割的横截性保证。共轭性 \(h \circ f|_{\Lambda} = \sigma \circ h\) 由构造直接验证。\(\Sigma_2\) 同胚于 Cantor 集，故 \(\Lambda\) 亦然。\(\blacksquare\)

\subsection{99.5 符号动力系统与拓扑共轭}\label{ux7b26ux53f7ux52a8ux529bux7cfbux7edfux4e0eux62d3ux6251ux5171ux8f6d}

\textbf{定义 99.7}（全移位空间）：设 \(\mathcal{A} = \{0, 1, \ldots, n-1\}\) 为有限字母表。\textbf{全 \(n\)-移位}为 \(\Sigma_n = \mathcal{A}^{\mathbb{Z}}\)（双边）或 \(\Sigma_n^+ = \mathcal{A}^{\mathbb{N}}\)（单边），配备乘积拓扑（由离散拓扑诱导，\(\Sigma_n\) 紧致且可度量）。\textbf{移位映射} \(\sigma : \Sigma_n \to \Sigma_n\) 定义为 \((\sigma(x))_k = x_{k+1}\)（向左移动一个坐标）。\((\Sigma_n, \sigma)\) 称为\textbf{全移位动力系统}。其闭的不变子集称为\textbf{子移位}，可由禁止词集合刻画。

\textbf{定理 99.4}（移位空间的基本同胚）：\(\Sigma_n\) 同胚于 Cantor 集。移位映射 \(\sigma\) 关于乘积测度为强混合 Bernoulli 自同构，拓扑熵 \(h_{\operatorname{top}}(\sigma) = \log n\)。

\textbf{证明}：\(\Sigma_n = \mathcal{A}^{\mathbb{Z}}\) 的拓扑由柱集 \([a_{-k},\ldots,a_k] = \{x : x_{-k}=a_{-k},\ldots,x_k=a_k\}\) 生成。每个柱集为闭开集，故 \(\Sigma_n\) 完全不连通。紧致性由 Tychonoff 定理（\(\mathcal{A}\) 为有限离散空间）保证。无孤立点：对任意 \(x\) 和 \(N\)，取 \(y\) 与 \(x\) 在 \([-N,N]\) 外不同，则 \(y \neq x\) 且 \(d(x,y)\) 可任意小。故 \(\Sigma_n\) 为紧致、完全不连通、无孤立点的度量空间，由 Brouwer 定理，它同胚于 Cantor 集。\(\sigma\) 为强混合：对柱集 \(A,B\)，\(\mu(\sigma^{-n}A \cap B) = \mu(A)\mu(B)\) 对 \(n\) 充分大（因 \(A\) 和 \(\sigma^{-n}A\) 的分量不重叠）。拓扑熵：有限剖分 \(\mathcal{P} = \{[0],\ldots,[n-1]\}\) 为生成剖分，\(h_{\operatorname{top}}(\sigma) = \lim_{m\to\infty} \frac{1}{m}\log(\# \text{允许词长} m) = \log n\)。\(\blacksquare\)

\textbf{定义 99.8}（拓扑共轭）：设 \(f : X \to X\) 与 \(g : Y \to Y\) 为两个（离散或连续）动力系统。若存在同胚 \(h : X \to Y\) 使得 \(h \circ f = g \circ h\)，则称 \(f\) 与 \(g\) 拓扑共轭（\(h\) 为共轭映射）。对连续系统 \(\varphi_t\) 与 \(\psi_t\)，要求 \(h \circ \varphi_t = \psi_t \circ h\) 对所有 \(t\) 成立。

\textbf{定理 99.5}（拓扑共轭的不变量）：若 \(f\) 与 \(g\) 拓扑共轭，则以下量相等： 1. \(P_n(f) = |\{x : f^n(x) = x\}|\)（周期 \(n\) 的点数）--- 由 \(h\) 的一一对应推出； 2. 拓扑熵 \(h_{\operatorname{top}}(f) = h_{\operatorname{top}}(g)\)； 3. 非游荡集 \(\Omega(f)\) 的拓扑结构（\(h(\Omega(f)) = \Omega(g)\)）； 4. 链回归集与吸引子/排斥子的格结构。

\textbf{证明}：对于（1），\(f^n(x) = x \iff h(f^n(x)) = h(x) \iff g^n(h(x)) = h(x)\)，故 \(h\) 给出 \(P_n(f)\) 到 \(P_n(g)\) 的双射。对于（2），拓扑熵由开覆盖生成，共轭映射将覆盖及其迭代原样搬运，故熵不变。\(\blacksquare\)

\hrulefill

\hrulefill

\section{第323章：保测变换与遍历性}\label{ux7b2c323ux7ae0ux4fddux6d4bux53d8ux6362ux4e0eux904dux5386ux6027}

本章严格建立遍历论的公理化基石------保测变换与遍历性等价刻画------并为 Birkhoff 与 von Neumann 定理提供完整证明。

\subsection{240.1 保测变换与不变测度}\label{ux4fddux6d4bux53d8ux6362ux4e0eux4e0dux53d8ux6d4bux5ea6}

\textbf{定义 240.1}（保测变换）：设 \((X, \mathcal{B}, \mu)\) 为概率测度空间（\(\mu(X) = 1\)），\(T : X \to X\) 为可测映射。若对任意 \(A \in \mathcal{B}\) 有 \(\mu(T^{-1} A) = \mu(A)\)，则称 \(T\) 为\textbf{保测变换}，\(\mu\) 为 \(T\)-不变测度，\((X, \mathcal{B}, \mu, T)\) 称为\textbf{保测动力系统}。

若 \(T\) 为双射且 \(T^{-1}\) 也可测并保测，则称 \(T\) 为\textbf{可逆保测变换}。对连续时间系统 \(\{T_t\}_{t \in \mathbb{R}}\)，要求对每个 \(t\)，\(T_t\) 保测且 \((t, x) \mapsto T_t(x)\) 可测。

\textbf{定义 240.2}（不变集的 \(\sigma\)-代数）：定义

\[\mathcal{I} = \{A \in \mathcal{B} : T^{-1} A = A \;\; \mu\text{-a.e.}\}\]

\(\mathcal{I}\) 为 \(\mathcal{B}\) 的完备子 \(\sigma\)-代数。若函数 \(f : X \to \mathbb{R}\) 为 \(\mathcal{I}\)-可测，则 \(f \circ T = f\) \(\mu\)-a.e.（\(T\)-不变函数）。

\subsection{240.2 遍历性的等价刻画}\label{ux904dux5386ux6027ux7684ux7b49ux4ef7ux523bux753b}

\textbf{定义 240.3}（遍历性）：保测动力系统 \((X, \mathcal{B}, \mu, T)\) 称为\textbf{遍历的}，若对任意 \(A \in \mathcal{I}\)，要么 \(\mu(A) = 0\) 要么 \(\mu(A) = 1\)（即不变集平凡）。

\textbf{定理 240.1}（遍历性的六条等价条件）：对保测系统 \((X, \mathcal{B}, \mu, T)\)，以下等价： 1. \(T\) 是遍历的（不变集平凡）。 2. 任意 \(T\)-不变可测函数 \(f\)（\(f \circ T = f\) \(\mu\)-a.e.）必几乎处处常数。 3. 对任意 \(A, B \in \mathcal{B}\)，\(\displaystyle \lim_{n \to \infty} \frac{1}{n} \sum_{k=0}^{n-1} \mu(T^{-k} A \cap B) = \mu(A) \mu(B)\)。 4. 对任意 \(f \in L^1(\mu)\)，其在 \(T\) 下的轨道平均依 \(L^1\) 收敛：\(\frac{1}{n} \sum_{k=0}^{n-1} f \circ T^k \xrightarrow{L^1} \int_X f \, d\mu\)。 5. 不存在非平凡的 \(T\)-不变分解：若 \(X = A \cup B\) 且 \(T^{-1}A = A\), \(T^{-1}B = B\), \(A \cap B = \emptyset\)，则 \(\mu(A) = 0\) 或 \(\mu(B) = 0\)。 6. 对任意 \(f \in L^1(\mu)\)，\(\frac{1}{n}\sum_{k=0}^{n-1} f(T^k x)\) 的几乎处处极限为常数 \(\int_X f\, d\mu\)。

\textbf{证明}：(1) \(\Rightarrow\) (2)：若 \(f \circ T = f\) a.e.，则对任意 \(c \in \mathbb{R}\)，\(\{x : f(x) \leq c\} \in \mathcal{I}\)，故其测度为 0 或 1，推出 \(f\) a.e. 常数。(2) \(\Rightarrow\) (5)：令 \(f = \chi_A\)，则 \(f \circ T = f\)，故 \(f\) a.e. 常数，即 \(\mu(A) \in \{0, 1\}\)。(5) \(\Rightarrow\) (1)：平凡。(2) \(\Rightarrow\) (3)：对固定的 \(B\)，构造测度 \(\nu_n(A) = \frac{1}{n} \sum_{k=0}^{n-1} \mu(T^{-k}A \cap B)\)。由 Banach-Alaoglu 定理，\(\{\nu_n\}\) 有弱-\(*\) 聚点 \(\nu\)，且 \(\nu\) 是 \(T\)-不变的。若 \(T\) 遍历，则 \(\nu\) 必与 \(\mu\) 成比例，由 \(\nu(X) = \mu(B)\) 得 \(\nu = \mu(B) \mu\)，从而 (3) 成立。(3) \(\Rightarrow\) (1)：取 \(A \in \mathcal{I}\)，令 \(B = A\)，由 (3) 得 \(\mu(A) = \mu(A)^2\)，故 \(\mu(A) \in \{0, 1\}\)。剩余蕴含关系由 Birkhoff 定理（定理 240.2）建立。\(\blacksquare\)

\subsection{240.3 Birkhoff 逐点遍历定理}\label{birkhoff-ux9010ux70b9ux904dux5386ux5b9aux7406}

\textbf{定理 240.2}（Birkhoff 逐点遍历定理，1931）：设 \((X, \mathcal{B}, \mu, T)\) 为保测动力系统，\(f \in L^1(\mu)\)。则对 \(\mu\)-几乎每个 \(x \in X\)，

\[\lim_{n \to \infty} \frac{1}{n} \sum_{k=0}^{n-1} f(T^k x) = \mathbb{E}[f \mid \mathcal{I}](x)\]

其中 \(\mathbb{E}[f \mid \mathcal{I}]\) 是 \(f\) 关于不变 \(\sigma\)-代数 \(\mathcal{I}\) 的条件期望。特别地，若 \(T\) 遍历，则极限为常数 \(\int_X f \, d\mu\)。极限函数属于 \(L^1(\mu)\) 且收敛亦在 \(L^1\) 意义下成立。

\textbf{证明}：定义 \(A_n f(x) = \frac{1}{n}\sum_{k=0}^{n-1} f(T^k x)\)。对任意 \(N\) 及 \(a < b\)，定义 Bad 集合 \(E_{a,b} = \{x : \limsup A_n f(x) > b > a > \liminf A_n f(x)\}\)。核心为极大遍历引理：对 \(F_N(x) = \max_{1 \leq m \leq N} \sum_{k=0}^{m-1} (f - a)(T^k x)\)，有 \(\int_{\{F_N > 0\}} (f - a) d\mu \geq 0\)。证明：对每个 \(x\)，定义首次极大时间 \(\tau(x) = \min\{1 \leq m \leq N : S_m f(x) > 0\}\)（\(S_m f = \sum_{k=0}^{m-1}(f-a)(T^k x)\)）。对任意 \(m\)，若 \(\tau(x)=m\)，则 \(S_m f(x) > 0\) 且 \(S_m f(Tx) = S_{m+1}f(x) - (f-a)(x)\)。通过重组求和并利用 \(T\) 保测性，得 \(\int_{F_N>0} (f-a)d\mu \geq \sum_{m=1}^N \int_{\tau=m} [S_m f - S_{\tau\circ T}f \circ T]d\mu \geq 0\)。由此可证 \(\mu(E_{a,b}) = 0\)（否则导出矛盾）。对所有有理数 \(a<b\) 取并集得零测集，在其补集上 \(\lim A_n f(x)\) 存在。极限函数 \(\bar{f}(x)\) 满足 \(\bar{f}\circ T = \bar{f}\) a.e.，且 \(\int \bar{f} d\mu = \int f d\mu\) 及 \(\int_A \bar{f} d\mu = \int_A f d\mu\)（\(A \in \mathcal{I}\)），故 \(\bar{f} = \mathbb{E}[f \mid \mathcal{I}]\)。\(\blacksquare\)

\textbf{定理 240.3}（von Neumann 平均遍历定理，1932）：在 Hilbert 空间 \(L^2(\mu)\) 中，保测算子 \(U_T : f \mapsto f \circ T\)（Koopman 算子）是等距算子。轨道平均 \(\frac{1}{n} \sum_{k=0}^{n-1} U_T^k f\) 依 \(L^2\)-范数收敛到 \(f\) 到 \(U_T\)-不变子空间的正交投影 \(P f = \mathbb{E}[f \mid \mathcal{I}]\)。

\textbf{证明}：\(L^2(\mu) = \operatorname{Fix}(U_T) \oplus \overline{\operatorname{Im}(I - U_T)}\)（正交分解）。\(f \in \operatorname{Fix}(U_T)\) 即 \(f \circ T = f\)，此时 \(\frac{1}{n} \sum_{k=0}^{n-1} U_T^k f = f\) 平凡。对 \(f = g - U_T g \in \operatorname{Im}(I - U_T)\)，\(\frac{1}{n} \sum_{k=0}^{n-1} U_T^k f = \frac{1}{n}(g - U_T^n g) \to 0\)（\(L^2\) 有界性）。闭包上的连续性即得一般情形。\(\blacksquare\)

\subsection{240.4 混合性}\label{ux6df7ux5408ux6027}

\textbf{定义 240.4}（强混合与弱混合）：保测系统 \((X, \mathcal{B}, \mu, T)\) 称为： - \textbf{强混合}：对任意 \(A, B \in \mathcal{B}\)，\(\displaystyle \lim_{n \to \infty} \mu(T^{-n} A \cap B) = \mu(A) \mu(B)\)。 - \textbf{弱混合}：对任意 \(A, B \in \mathcal{B}\)，\(\displaystyle \lim_{n \to \infty} \frac{1}{n} \sum_{k=0}^{n-1} |\mu(T^{-k} A \cap B) - \mu(A) \mu(B)| = 0\)。

\textbf{定理 240.4}（Sinai 因子定理的定性表述）：每个正熵保测变换 \(T\) 包含一个 Bernoulli 因子。更精确地说，存在 \(T\) 的因子 \(\sigma\)（即存在可测映射 \(\pi : X \to \Sigma\) 满足 \(\pi \circ T = \sigma \circ \pi\)）使得 \((\Sigma, \mathcal{B}_\Sigma, \nu, \sigma)\) 同构于一个 Bernoulli 移位。

\textbf{证明}：设 \(h_\mu(T) > 0\)。由熵的定义，存在有限可测剖分 \(\mathcal{P}\) 使得 \(h(T,\mathcal{P}) > 0\)。考虑 \(\mathcal{P}\) 的过去和未来演化：令 \(\mathcal{P}_n = \bigvee_{k=-n}^n T^{-k}\mathcal{P}\)，其极限 \(\mathcal{P}_\infty = \bigvee_{k=-\infty}^\infty T^{-k}\mathcal{P}\) 生成一个因子 \(\sigma\)-代数 \(\mathcal{A}\)。由 Sinai 定理，\(\mathcal{A}\) 同构于一个 Bernoulli 移位。证明的关键在于验证 \(\mathcal{A}\) 在给定 \(\mathcal{P}_\infty\) 的条件下具有独立性：利用 Ornstein 的 ``very weak Bernoulli''（VWB）条件，即对任意 \(\varepsilon>0\)，存在 \(N\) 使得 \(d(\mathcal{P}_N, \mathcal{P}_{N+m})\) 充分小，其中 \(d\) 为剖分距离。由正熵性质，\(h(T,\mathcal{P}) > 0\) 蕴含 VWB 条件，从而因子 \(\mathcal{A}\) 是 Bernoulli 的。\(\blacksquare\)

\subsection{240.5 Kolmogorov-Sinai 熵}\label{kolmogorov-sinai-ux71b5}

\textbf{定义 240.5}（剖分与熵）：设 \(\mathcal{P} = \{P_1, \ldots, P_k\}\) 为 \(X\) 的有限可测剖分。其\textbf{信息函数}为 \(I_{\mathcal{P}}(x) = -\sum_{i=1}^{k} \chi_{P_i}(x) \log \mu(P_i)\)。\(\mathcal{P}\) 的\textbf{熵}为

\[H(\mathcal{P}) = \int_X I_{\mathcal{P}} \, d\mu = -\sum_{i=1}^{k} \mu(P_i) \log \mu(P_i)\]

对剖分 \(\mathcal{P}\)，定义 \(T^{-1} \mathcal{P} = \{T^{-1} P_i\}\)，\(\bigvee_{i=0}^{n-1} T^{-i} \mathcal{P}\) 为诸剖分的联合（最细公共加细）。\(T\) 关于 \(\mathcal{P}\) 的\textbf{过程熵}为

\[h(T, \mathcal{P}) = \lim_{n \to \infty} \frac{1}{n} H\left(\bigvee_{i=0}^{n-1} T^{-i} \mathcal{P}\right)\]

极限存在性由 \(H(\mathcal{P} \vee \mathcal{Q}) \leq H(\mathcal{P}) + H(\mathcal{Q})\) 的次可加性保证。\(T\) 的\textbf{Kolmogorov-Sinai 熵}为

\[h(T) = \sup_{\mathcal{P} \text{ 有限剖分}} h(T, \mathcal{P})\]

若 \(\mathcal{P}\) 为生成剖分（\(\bigvee_{n=0}^{\infty} T^{-n} \mathcal{P} = \mathcal{B}\) mod \(\mu\)），则 Kolmogorov-Sinai 定理给出 \(h(T) = h(T, \mathcal{P})\)。

\textbf{定理 240.5}（Shannon-MacMillan-Breiman 定理）：对遍历保测系统 \((X, \mathcal{B}, \mu, T)\) 及有限生成剖分 \(\mathcal{P}\)，

\[-\frac{1}{n} \log \mu\left(\mathcal{P}_n(x)\right) \xrightarrow{\text{a.e. 且 } L^1} h(T)\]

其中 \(\mathcal{P}_n(x)\) 为 \(\bigvee_{i=0}^{n-1} T^{-i} \mathcal{P}\) 中包含 \(x\) 的原子。此定理是信息论中渐近均分性质在动力系统语境下的推广：典型轨道的剖分原子测度以指数速率 \(e^{-n h(T)}\) 衰减。

\textbf{证明}：设 \(f_n(x) = -\frac{1}{n}\log\mu(\mathcal{P}_n(x))\)。定义 \(g_k(x) = -\log\mu(P_{i_k} \mid \bigvee_{j=0}^{k-1} T^{-j}\mathcal{P})(x)\) 为条件信息函数。则 \(f_n = \frac{1}{n}\sum_{k=0}^{n-1} g_k \circ T^k\)。由鞅收敛定理，\(g_k\) 在 \(L^1\) 中收敛到不变函数 \(g\)。由 Birkhoff 遍历定理，\(\frac{1}{n}\sum_{k=0}^{n-1} g \circ T^k \to \int g d\mu\) a.e. 且 \(L^1\)。需验证 \(\int g d\mu = h(T,\mathcal{P})\)。由熵的定义，\(h(T,\mathcal{P}) = \lim_n \frac{1}{n}H(\bigvee_{i=0}^{n-1}T^{-i}\mathcal{P}) = \lim_n \frac{1}{n}\sum_{k=0}^{n-1} H(\mathcal{P} \mid \bigvee_{i=1}^k T^{-i}\mathcal{P}) = \int g d\mu\)（由鞅收敛和条件熵的链式法则）。遍历性保证 \(g\) 为常数 a.e.，故 \(g = h(T,\mathcal{P})\) a.e.。\(\blacksquare\)

\subsection{240.6 Ornstein 同构定理}\label{ornstein-ux540cux6784ux5b9aux7406}

\textbf{定义 240.6}（Bernoulli 移位）：设 \((p_1, \ldots, p_k)\) 为概率向量。Bernoulli 移位 \(B(p_1, \ldots, p_k)\) 是 \(\{1, \ldots, k\}^{\mathbb{Z}}\) 上的乘积测度空间关于移位映射 \(\sigma\) 的动力系统。其为强混合的且 \(h(B) = -\sum_i p_i \log p_i\)。

\textbf{定理 240.6}（Ornstein 同构定理，1970）：两个 Bernoulli 移位是测度论同构的（即存在保测同胚交换移位作用）当且仅当它们具有相同的 Kolmogorov-Sinai 熵。即：熵是 Bernoulli 移位的完全不变量。

此定理是遍历论中最深刻的结构定理之一。其证明过程极为复杂，核心是构造\textbf{精细化匹配}（very weak Bernoulli partitions）的 Ornstein 理论。推论：任何一个有限状态的 Bernoulli 移位同构于其自身的任何平方根熵的 Bernoulli 因子，只要熵足够细分。

\textbf{证明}（Ornstein 理论框架）：设 \(T_1\) 和 \(T_2\) 为两个 Bernoulli 移位，具有相同的熵 \(h\)。需构造保测同构 \(\phi\) 使得 \(\phi \circ T_1 = T_2 \circ \phi\)。(1) 定义「非常弱 Bernoulli」（VWB）性质：分割 \(\mathcal{P}\) 称为 VWB 的，如果其过去和未来的分布以某种度量方式接近独立。Bernoulli 移位的生成分割 \(\mathcal{P}\) 是 VWB 的。(2) 核心引理（Ornstein 匹配定理）：若两个系统有 VWB 的生成分割且熵相同，则存在保测同构。证明通过逐次逼近构造：在每一步构造一个「几乎」与目标分割匹配的有限编码，利用 Sinai 的「弱 Bernoulli」定理处理误差。(3) 对 Bernoulli 移位，生成分割自动为 VWB，且熵 \(h(T) = \sum p_i \log(1/p_i)\)。因此两个 Bernoulli 移位同构当且仅当它们具有相同的熵。\(\blacksquare\)

\hrulefill

\hrulefill

\section{第324章：混沌与分形}\label{ux7b2c324ux7ae0ux6df7ux6c8cux4e0eux5206ux5f62}

混沌理论揭示了确定性非线性系统中的不可预测行为，而分形几何则刻画了混沌动力系统所产生的奇异吸引子的几何结构。本章从一维映射的符号动力学出发，经由 Sharkovsky 序和 Li-Yorke 定理，引入拓扑混沌的严格定义。进而讨论二维连续系统的 Poincaré-Bendixson 定理、极限环理论以及高维系统的奇异吸引子。最后引入分形维数（Hausdorff 维数、盒维数）和迭代函数系统（IFS）的理论。混沌与分形揭示了确定性、随机性和几何复杂性之间的深刻联系。

\subsection{196.1 一维映射的符号动力学}\label{ux4e00ux7ef4ux6620ux5c04ux7684ux7b26ux53f7ux52a8ux529bux5b66}

\textbf{定义 196.1}（符号动力系统）：设 \(\Sigma_2 = \{0, 1\}^{\mathbb{N}}\) 为半无穷二值序列空间，配备移位映射 \(\sigma: \Sigma_2 \to \Sigma_2\)： \[\sigma(s_1 s_2 s_3 \cdots) = s_2 s_3 \cdots\] 赋予度量 \(d(s, t) = \sum_{n=1}^\infty |s_n - t_n| / 2^n\)，\((\Sigma_2, \sigma)\) 为紧致度量空间上的连续映射。符号动力系统是混沌研究中最基本的模型，其动力学性质可精确刻画，且通过拓扑半共轭关系与许多物理系统建立联系。

\textbf{命题 196.1}（移位映射的混沌性质）：\((\Sigma_2, \sigma)\) 具有以下性质： (i) 周期点稠密（周期点集 \(\operatorname{Per}(\sigma)\) 在 \(\Sigma_2\) 中稠密）； (ii) 拓扑传递性（存在轨道稠密的点，即 \(\{\sigma^n(s) : n \geq 0\}\) 在 \(\Sigma_2\) 中稠密）； (iii) 对初值的敏感依赖性（存在 \(\delta > 0\) 使对任意 \(s \in \Sigma_2\) 和 \(\varepsilon > 0\)，存在 \(t\) 与 \(n\) 使 \(d(s,t) < \varepsilon\) 且 \(d(\sigma^n(s), \sigma^n(t)) \geq \delta\)）。

\textbf{证明}：(i) 周期 \(k\) 的点为满足 \(s_{n+k} = s_n\) 的序列，全体周期点可数且在 \(\Sigma_2\) 中稠密（对任意有限前缀，可周期性重复构造周期点）。(ii) 构造 \(s = (0, 1, 00, 01, 10, 11, 000, 001, \ldots)\) 包含所有有限二值块，则 \(\sigma^n(s)\) 遍历所有有限序列，故轨道稠密。(iii) 取 \(\delta = 1/2\)，对任意 \(s\) 和 \(\varepsilon > 0\)，取 \(n\) 使 \(2^{-n} < \varepsilon\)，则 \(t\) 与 \(s\) 在 \(n\) 位后首次不同，且在 \(n+1\) 位后仍不同，则 \(d(\sigma^n(s), \sigma^n(t)) \geq 1/2\)。\(\blacksquare\)

\textbf{定义 196.2}（拓扑 Markov 链 / 子移位）：设 \(A\) 为 \(N \times N\) 的 0-1 转移矩阵（\(A_{ij} \in \{0,1\}\)）。定义\textbf{子移位空间} \[\Sigma_A = \{s \in \{1, \ldots, N\}^{\mathbb{N}} : A_{s_n s_{n+1}} = 1, \; \forall n\}\] \((\Sigma_A, \sigma)\) 称为\textbf{拓扑 Markov 链}。拓扑 Markov 链是符号动力系统的推广，其允许的转移由矩阵 \(A\) 编码，\(A\) 的谱性质（Perron-Frobenius 定理）直接决定了拓扑 Markov 链的动力学性质（周期点计数、拓扑熵等）。

\textbf{定理 196.1}（拓扑 Markov 链的周期点计数）：\((\Sigma_A, \sigma)\) 中周期 \(n\) 的点的个数为 \(\operatorname{tr}(A^n)\)。特别地，Perron-Frobenius 特征值 \(\lambda_1\) 满足 \(\lim_{n \to \infty} \frac{1}{n}\log \operatorname{tr}(A^n) = \log \lambda_1\)，即周期点数的指数增长率等于拓扑熵。

\textbf{证明}：\(\sigma^n(s) = s\) 等价于 \(s_{k+n} = s_k\)（对所有 \(k\)），即序列由 \(n\) 长块 \(s_1 s_2 \cdots s_n\) 周期重复。该块的首尾必须满足 \(A_{s_n s_1} = 1\)，且相邻满足 \(A_{s_k s_{k+1}} = 1\)。故固定点数为 \(s_1\) 到 \(s_n\) 再到 \(s_1\) 的闭路径数，即 \(\sum_{i=1}^N (A^n)_{ii} = \operatorname{tr}(A^n)\)。\(\blacksquare\)

\subsection{196.2 一维映射的混沌理论}\label{ux4e00ux7ef4ux6620ux5c04ux7684ux6df7ux6c8cux7406ux8bba}

\textbf{定义 196.3}（帐篷映射）：设 \(T: [0, 1] \to [0, 1]\) 为 \[T(x) = \begin{cases} 2x, & 0 \leq x \leq 1/2 \\ 2 - 2x, & 1/2 < x \leq 1 \end{cases}\] 帐篷映射是分段线性映射，其斜率的绝对值为 \(2\) 处处成立（除 \(x=1/2\) 外）。帐篷映射是理解混沌的最简单模型之一，其动力学通过二进展开与移位映射精确对应。

\textbf{定理 196.2}（帐篷映射与移位映射的共轭）：帐篷映射 \(T\) 与 \(\Sigma_2\) 上的移位映射 \(\sigma\) 拓扑半共轭：用二进展开 \(x = 0.b_1 b_2 b_3 \cdots\)，\(T\) 的作用相当于移位 \(\sigma\) 或 \(\sigma\) 的补（取决于 \(b_1\)），即 \(T(0.b_1 b_2 b_3 \cdots) = 0.b_2 b_3 \cdots\)（若 \(b_1 = 0\)）或 \(T(0.b_1 b_2 b_3 \cdots) = 0.\bar{b}_2 \bar{b}_3 \cdots\)（若 \(b_1 = 1\)，其中 \(\bar{b} = 1-b\)）。由于二进展开的非唯一性（\(0.0111\cdots = 0.1000\cdots\)），\(T\) 与 \(\sigma\) 不是严格的拓扑共轭，而是半共轭------投影映射 \(\pi: \Sigma_2 \to [0,1]\) 将二进展开映为对应的实数，满足 \(\pi \circ \sigma = T \circ \pi\)，但 \(\pi\) 不是单射。

\textbf{证明}：若 \(x \in [0, 1/2]\)，则 \(b_1 = 0\)，\(T(x) = 2x = 0.b_2 b_3 \cdots\)（小数部分左移一位）。若 \(x \in [1/2, 1]\)，则 \(b_1 = 1\)，\(T(x) = 2-2x = 1 - 0.b_2 b_3 \cdots = 0.\bar{b}_2 \bar{b}_3 \cdots\)。验证 \(T^{n}\) 的迭代与二进展开的对应关系，注意到二进展开的非唯一性仅影响可数多个点，且在这些点上 \(T\) 的行为与移位映射一致。\(\blacksquare\)

\textbf{定义 196.4}（Logistic 映射）：\(f_\mu(x) = \mu x(1 - x)\)，\(x \in [0, 1]\)，\(\mu \in [0, 4]\)。Logistic 映射是种群动力学中离散时间模型的经典例子，也是通向混沌的倍周期分岔路线的最著名实例。

\textbf{定理 196.3}（Logistic 映射的倍周期分岔与 Feigenbaum 普适性）：随着 \(\mu\) 从 \(0\) 增加到 \(4\)，Logistic 映射经历以下动力学变化： - \(0 < \mu \leq 1\)：\(x = 0\) 为全局吸引的不动点 - \(1 < \mu \leq 3\)：\(x = 1 - 1/\mu\) 为稳定不动点 - \(3 < \mu \leq 1 + \sqrt{6} \approx 3.449\)：稳定 2-周期轨 - \(\mu > 3.449\)：稳定 4-周期轨，等等（倍周期分岔级联） - \(\mu \approx 3.569946 \ldots\)：混沌开始（Feigenbaum 点 \(\mu_\infty\)） - \(\mu = 4\)：完全混沌（\(f_4\) 与帐篷映射拓扑共轭）

\textbf{Feigenbaum 普适常数}：倍周期分岔序列满足 \[\delta = \lim_{n \to \infty} \frac{\mu_n - \mu_{n-1}}{\mu_{n+1} - \mu_n} = 4.6692016091\ldots\] \[\alpha = \lim_{n \to \infty} \frac{d_n}{d_{n+1}} = 2.5029078750\ldots\] 其中 \(\mu_n\) 为第 \(n\) 次分岔的 \(\mu\) 值，\(d_n\) 为分岔图中超稳定周期轨的分岔间距。\(\delta\) 和 \(\alpha\) 是\textbf{普适常数}（Feigenbaum 1978），对所有在极大值处具有二次极大值的单峰映射均相同------这是重正化群在动力系统理论中的深刻应用，揭示了倍周期分岔路线通向混沌的普适标度规律。

\textbf{证明}：Feigenbaum 的证明基于重正化群分析。定义加倍算子 \(\mathcal{R}\) 将 \(f\) 映为 \(f\) 的二次迭代的重标度：\(\mathcal{R}[f](x) = -\alpha f(f(-x/\alpha))\)。\(\mathcal{R}\) 在函数空间中有不动点 \(g\)（满足 \(g = \mathcal{R}[g]\)，即 \(g(x) = -\alpha g(g(-x/\alpha))\)），此不动点方程为 Cvitanovic-Feigenbaum 方程。在不动点处线性化，\(\mathcal{R}\) 的线性化算子 \(D\mathcal{R}[g]\) 只有一个特征值 \(\delta > 1\)（相关本征方向），其余特征值模小于 1。\(\delta\) 即为控制分岔参数收敛速率的相关特征值，\(\alpha\) 为函数重标度因子。普适性源于任何具有二次极大值的单峰映射在重正化群流下均收敛到同一不动点 \(g\)。\(\blacksquare\)

\textbf{定理 196.4}（Li-Yorke 定理，``周期三蕴含混沌''，1975）：设 \(f: I \to I\)（\(I\) 为区间）为连续映射。若 \(f\) 有周期 \(3\) 的点，则 \(f\) 具有以下性质： (i) 对每个 \(k \in \mathbb{N}\)，\(f\) 有周期 \(k\) 的点； (ii) 存在不可数集 \(S \subset I\)（不含周期点）使得对任意 \(x, y \in S\)（\(x \neq y\)）： \[\limsup_{n \to \infty} |f^n(x) - f^n(y)| > 0, \quad \liminf_{n \to \infty} |f^n(x) - f^n(y)| = 0\] 且对任意 \(x \in S\) 和任意周期点 \(p\)，\(\limsup_{n \to \infty} |f^n(x) - f^n(p)| > 0\)。

\textbf{证明}：设 \(a < b < c\) 为 \(f\) 的 3-周期轨，不妨设 \(f(a) = b\)，\(f(b) = c\)，\(f(c) = a\)（或 \(b, a, c\) 等，由对称性可归为此情形）。由连续性和介值定理，存在区间 \(J_0 = [a, b]\) 和 \(J_1 = [b, c]\) 满足 \(f(J_0) \supseteq J_1\) 且 \(f(J_1) \supseteq J_0 \cup J_1\)。对任意有限序列 \(\{s_1, \ldots, s_n\} \in \{0, 1\}^n\)，由区间包含关系，存在 \(x \in J_{s_1}\) 使 \(f^k(x) \in J_{s_{k+1}}\)（\(k=1,\ldots,n-1\)），且 \(f^n(x) = x\) 若 \(s_n = s_1\)（周期点）。此构造给出所有周期的周期点和 Li-Yorke 混沌集的不可数性。混沌集的构造：选取无穷序列 \(\{s_n\}\) 使其非最终周期，对应轨道的 \(\limsup\) 和 \(\liminf\) 的分离由区间 \(J_0\) 和 \(J_1\) 的分离性保证。\(\blacksquare\)

\textbf{定理 196.5}（Sharkovsky 定理，1964）：将正整数按 Sharkovsky 序排列： \[3 \triangleright 5 \triangleright 7 \triangleright \cdots \triangleright 2 \cdot 3 \triangleright 2 \cdot 5 \triangleright 2 \cdot 7 \triangleright \cdots \triangleright 2^2 \cdot 3 \triangleright 2^2 \cdot 5 \triangleright \cdots \triangleright 2^3 \triangleright 2^2 \triangleright 2 \triangleright 1\] 若 \(f\) 有周期 \(m\) 的点，且 \(m \triangleright n\)（\(m\) 在 Sharkovsky 序中先于 \(n\)），则 \(f\) 必有周期 \(n\) 的点。特别地，周期 \(3\) 的存在蕴含所有周期的存在（Li-Yorke 定理是 Sharkovsky 定理的特例）。

\textbf{证明思路}：Sharkovsky 的原始证明基于有向图的组合分析和区间映射的介值性质。核心引理：对连续区间映射 \(f\)，若存在周期 \(m\) 的点，则存在一个由 \(m\) 个区间组成的循环覆盖，其转移矩阵蕴含所有 Sharkovsky 序中 \(m\) 之后的周期。通过构造包含 \(m\) 个点的有向图（顶点为周期轨的 \(m\) 个点，边由 \(f\) 的像包含关系决定），并应用区间上的组合引理，证明该图包含所有必要长度的环。\(\blacksquare\)

\subsection{196.3 二维自治系统与 Poincaré-Bendixson 定理}\label{ux4e8cux7ef4ux81eaux6cbbux7cfbux7edfux4e0e-poincaruxe9-bendixson-ux5b9aux7406}

\textbf{定理 196.6}（Poincaré-Bendixson 定理）：设 \(\dot{\mathbf{x}} = \mathbf{f}(\mathbf{x})\) 为平面上的 \(C^1\) 自治系统，\(\gamma^+\) 为有界正向轨道。则 \(\omega\)-极限集 \(\omega(\gamma^+)\) 是以下三种之一： (i) 单一平衡点； (ii) 单一周期轨（极限环）； (iii) 由平衡点和连接它们的轨道组成的连通集（同宿轨或异宿轨）。

\textbf{推论 196.1}：在平面上，有界正向轨道若不趋于平衡点，则必趋于周期轨（或为周期轨）。这解释了为什么平面连续系统不可能出现混沌------混沌需要至少三维的连续时间系统或离散时间系统。

\textbf{证明}：设 \(\Gamma\) 为无切线段（横截线），即 \(\mathbf{f}(\mathbf{x}) \cdot \mathbf{n}(\mathbf{x}) \neq 0\) 对所有 \(\mathbf{x} \in \Gamma\)。关键引理：轨道 \(\gamma^+\) 与 \(\Gamma\) 的交点序列是单调的（沿 \(\Gamma\) 单方向移动）。这是因为平面上的 Jordan 曲线分离性质：轨道不能自交，两条轨道不能于非平衡点处相交。若 \(\omega(\gamma^+)\) 不含平衡点，则 \(\omega(\gamma^+)\) 为紧致不变集，取 \(\mathbf{p} \in \omega(\gamma^+)\) 及过 \(\mathbf{p}\) 的横截线 \(\Gamma\)，\(\gamma^+\) 与 \(\Gamma\) 的交点单调趋于 \(\mathbf{p}\)，故 \(\omega(\gamma^+)\) 为周期轨。若 \(\omega(\gamma^+)\) 含平衡点，则必由平衡点和连接轨道组成。\(\blacksquare\)

\textbf{定理 196.7}（Bendixson-Dulac 判据）：若在单连通区域 \(\Omega\) 上存在 \(C^1\) 函数 \(B(x, y) > 0\)（Dulac 函数）使得 \(\operatorname{div}(B\mathbf{f}) = \frac{\partial}{\partial x}(B f_1) + \frac{\partial}{\partial y}(B f_2)\) 在 \(\Omega\) 上不变号且不恒为零，则 \(\Omega\) 内不存在闭轨（周期轨）。该判据是平面系统极限环排除的基本工具。

\textbf{证明}：若存在周期轨 \(\Gamma \subset \Omega\)，由 Green 定理： \[\iint_{\operatorname{int}(\Gamma)} \operatorname{div}(B\mathbf{f}) \,dx\,dy = \oint_\Gamma B(f_1 dy - f_2 dx) = \oint_\Gamma B(\dot{x} dy - \dot{y} dx) = \oint_\Gamma B(\dot{x} \dot{y} - \dot{y} \dot{x}) dt = 0\] 与 \(\operatorname{div}(B\mathbf{f})\) 不变号且非零矛盾。\(\blacksquare\)

\subsection{196.4 奇异吸引子与 Lorenz 系统}\label{ux5947ux5f02ux5438ux5f15ux5b50ux4e0e-lorenz-ux7cfbux7edf}

\textbf{定义 196.5}（Lorenz 系统，1963）：Lorenz 方程 \[\begin{cases} \dot{x} = \sigma(y - x) \\ \dot{y} = rx - y - xz \\ \dot{z} = xy - bz \end{cases}\] 其中 \(\sigma = 10\)，\(b = 8/3\)，\(r = 28\)（经典参数）。Lorenz 系统是三维连续时间确定性系统，其长时间行为收敛于一个\textbf{奇异吸引子}------一个具有分形结构的紧致不变集，轨道在此吸引子上呈混沌运动。Lorenz 吸引子是第一个被发现的奇异吸引子（1963），其发现开创了混沌理论。

\textbf{Lorenz 系统的性质}： - 三个平衡点：\((0,0,0)\)（对所有 \(r\)）和 \(C^\pm = (\pm\sqrt{b(r-1)}, \pm\sqrt{b(r-1)}, r-1)\)（\(r > 1\) 时） - 当 \(r > 1\) 时，原点为鞍点（一维不稳定流形，二维稳定流形） - 当 \(r > \sigma(\sigma+b+3)/(\sigma-b-1)\) 时，\(C^\pm\) 失去稳定性（Hopf 分岔），产生奇异吸引子 - 能量函数 \(V = rx^2 + \sigma y^2 + \sigma(z-2r)^2\) 满足 \(\dot{V} = -2\sigma(rx^2 + y^2 + b(z-r)^2 - br^2)\)，说明系统最终有界（耗散系统）

\textbf{定理 196.8}（Tucker 定理，2002 --- Lorenz 吸引子的存在性）：对经典参数值 \((\sigma, b, r) = (10, 8/3, 28)\)，Lorenz 系统具有一个\textbf{奇异吸引子}，其上运动为混沌的。Tucker 通过严格的数值验证（区间算术）证明了 Lorenz 吸引子满足几何 Lorenz 模型的条件，从而确证了 Lorenz 吸引子的存在性------这是 Smale 第 14 问题（1998 年提出）的解答。

\textbf{证明概要}：Tucker 使用了严格数值方法（区间算术和正常形式分析）。核心步骤：(1) 在原点附近用正常形式将 Lorenz 流化为简化形式，构造 Poincaré 返回映射 \(P: \Sigma \to \Sigma\)（\(\Sigma\) 为适当选取的 Poincaré 截面）。(2) 证明 \(P\) 可分解为 \(P = G \circ F\)，其中 \(F\) 为奇异压缩（将截面压缩为 Cantor 集），\(G\) 为膨胀和折叠。(3) 利用几何 Lorenz 模型（Guckenheimer-Williams 1979）的理论，证明 \(P\) 在该模型下具有混沌性质：存在不变马蹄映射，蕴含正拓扑熵和敏感依赖性。(4) 通过区间算术严格验证数值条件，确保所有估计在浮点误差范围内成立。\(\blacksquare\)

\subsection{196.5 分形几何}\label{ux5206ux5f62ux51e0ux4f55}

\textbf{定义 196.6}（Hausdorff 维数与 Hausdorff 测度）：设 \(E \subset \mathbb{R}^n\)，\(s \geq 0\)。\(\delta\)-覆盖下的 \(s\) 维 Hausdorff 预测度为 \[\mathcal{H}_\delta^s(E) = \inf\left\{\sum_{i=1}^\infty (\operatorname{diam} U_i)^s : E \subset \bigcup_i U_i, \; \operatorname{diam} U_i \leq \delta\right\}\] \(s\) 维 Hausdorff 测度：\(\mathcal{H}^s(E) = \lim_{\delta \to 0} \mathcal{H}_\delta^s(E)\)。\textbf{Hausdorff 维数}为 \[\dim_H(E) = \inf\{s \geq 0 : \mathcal{H}^s(E) = 0\} = \sup\{s \geq 0 : \mathcal{H}^s(E) = \infty\}\] Hausdorff 维数是最基本的分形维数，对任意集合 \(E\) 有定义，且满足 \(\dim_H(E) \leq \dim_B(E)\)（盒维数）。

\textbf{定义 196.7}（盒维数 / Minkowski 维数）：设 \(N_\delta(E)\) 为覆盖 \(E\) 所需的最小 \(\delta\)-网格立方体数。上盒维数和下盒维数分别为 \[\overline{\dim}_B(E) = \limsup_{\delta \to 0} \frac{\log N_\delta(E)}{-\log \delta}, \quad \underline{\dim}_B(E) = \liminf_{\delta \to 0} \frac{\log N_\delta(E)}{-\log \delta}\] 若两者相等，则称公共值为\textbf{盒维数} \(\dim_B(E)\)。盒维数比 Hausdorff 维数更容易计算，但对可数稠密集（如 \(\mathbb{Q}\)）的盒维数为 \(1\) 而 Hausdorff 维数为 \(0\)，说明盒维数缺乏可数稳定性。

\textbf{定理 196.9}（Cantor 三分集的维数）：Cantor 三分集是 \([0,1]\) 中每次去掉中间三分之一开区间所得集合的极限。其 Hausdorff 维数和盒维数均为 \(\dim_H(C) = \dim_B(C) = \log 2 / \log 3 \approx 0.6309\)。\(C\) 有 \(0\) Lebesgue 测度，不可数，且为完全不连通的完备集。Cantor 集的构造是自相似分形的最简单例子。

\textbf{证明}：Cantor 集 \(C = \bigcap_{n=0}^\infty C_n\)，其中 \(C_n\) 为 \(2^n\) 个长度为 \(3^{-n}\) 的闭区间的并。对 \(s = \log 2 / \log 3\)，上界：\(\mathcal{H}_{3^{-n}}^s(C) \leq 2^n \cdot (3^{-n})^s = 2^n \cdot 3^{-ns} = 2^n \cdot 2^{-n} = 1\)（用 \(C_n\) 的构成区间作为覆盖）。令 \(n \to \infty\) 得 \(\mathcal{H}^s(C) \leq 1\)。下界：对任意覆盖 \(\{U_i\}\)，可证 \(\sum (\operatorname{diam} U_i)^s \geq 1\)（利用质量分布原理：定义 \(C\) 上的自然测度 \(\mu\)，每个 \(C_n\) 的构成区间赋予质量 \(2^{-n}\)，对任意 \(U_i\) 有 \(\mu(U_i) \leq C(\operatorname{diam} U_i)^s\)，故 \(\sum (\operatorname{diam} U_i)^s \geq C^{-1} \sum \mu(U_i) \geq C^{-1}\)）。因此 \(\mathcal{H}^s(C) = 1\)，\(\dim_H(C) = s\)。盒维数：\(N_{3^{-n}}(C) = 2^n\)，故 \(\dim_B(C) = \lim \frac{\log 2^n}{-\log 3^{-n}} = \frac{\log 2}{\log 3}\)。\(\blacksquare\)

\textbf{定义 196.8}（自相似集与迭代函数系统 IFS）：设 \(\{f_i\}_{i=1}^m\) 为 \(\mathbb{R}^n\) 上的压缩映射（\(\|f_i(x) - f_i(y)\| \leq c_i\|x-y\|\)，\(c_i < 1\)）。由 Hutchinson（1981）定理，存在唯一非空紧集 \(K\)（称为\textbf{吸引子}）满足 \[K = \bigcup_{i=1}^m f_i(K)\] \(K\) 称为\textbf{自相似集}（若 \(f_i\) 为相似映射）或\textbf{自仿射集}（若 \(f_i\) 为仿射映射）。Koch 雪花曲线、Sierpinski 三角形、Menger 海绵等均为自相似分形。IFS 的吸引子是分形几何的核心研究对象，其 Hausdorff 维数在自相似且满足开集条件时可由 Moran 方程 \(\sum_{i=1}^m c_i^s = 1\) 唯一确定。

\textbf{定理 196.10}（Hutchinson 定理，1981）：设 \(\{f_i\}_{i=1}^m\) 为完备度量空间 \((X,d)\) 上的压缩映射。定义 Hutchinson 算子 \(F(A) = \bigcup_{i=1}^m f_i(A)\)（\(A \subset X\) 紧致）。则 \(F\) 在 \(X\) 的紧致子集空间（Hausdorff 度量下）上为压缩映射，存在唯一吸引子 \(K\)。此外，对任意紧致集 \(A_0\)，迭代 \(F^n(A_0) \to K\)（Hausdorff 度量下）。

\textbf{证明}：设 \(\mathcal{K}(X)\) 为 \(X\) 的非空紧致子集全体，配备 Hausdorff 度量 \(d_H(A,B) = \max\{\sup_{a \in A} \inf_{b \in B} d(a,b), \sup_{b \in B} \inf_{a \in A} d(a,b)\}\)。\((\mathcal{K}(X), d_H)\) 为完备度量空间。对 \(A, B \in \mathcal{K}(X)\)， \[d_H(F(A), F(B)) = d_H(\bigcup f_i(A), \bigcup f_i(B)) \leq \max_i d_H(f_i(A), f_i(B)) \leq (\max_i c_i) d_H(A,B)\] 由于 \(\max_i c_i < 1\)，\(F\) 为压缩映射。由 Banach 不动点定理，存在唯一不动点 \(K\)。\(\blacksquare\)

\textbf{定义 196.9}（分形维数与开集条件）：设 \(\{f_i\}\) 为相似比为 \(c_i\) 的相似映射族。若存在非空有界开集 \(U\) 使 \(f_i(U) \subset U\) 且 \(f_i(U) \cap f_j(U) = \emptyset\)（\(i \neq j\)），则称满足\textbf{开集条件}（Open Set Condition, OSC）。在 OSC 下，吸引子 \(K\) 的 Hausdorff 维数 \(\dim_H(K)\) 由 Moran 方程 \(\sum_{i=1}^m c_i^s = 1\) 的唯一解 \(s\) 给出，且 \(0 < \mathcal{H}^s(K) < \infty\)。

\textbf{证明}：下界：在 OSC 下，定义 \(K\) 上的自然测度 \(\mu\)（由 \(m\) 个等概率分枝的 Bernoulli 测度通过 \(f_i\) 的像诱导），可证 \(\mu\) 满足 Frostman 条件：\(\mu(U) \leq C(\operatorname{diam} U)^s\) 对所有 \(U\)。上界：由 \(K\) 的 \(n\) 级近似 \(K_n = \bigcup_{i_1,\ldots,i_n} f_{i_1} \circ \cdots \circ f_{i_n}(K)\) 给出覆盖，直径 \(\leq (\max c_i)^n \operatorname{diam} K\)，\(\mathcal{H}^s_{\delta_n}(K) \leq \sum_{i_1,\ldots,i_n} (\prod_{k=1}^n c_{i_k} \operatorname{diam} K)^s = (\operatorname{diam} K)^s\)。\(\blacksquare\)

\subsection{196.6 Julia 集与 Mandelbrot 集}\label{julia-ux96c6ux4e0e-mandelbrot-ux96c6}

\textbf{定义 196.10}（Julia 集与 Mandelbrot 集）：对复二次多项式 \(f_c(z) = z^2 + c\)（\(c \in \mathbb{C}\)），\textbf{填充 Julia 集} \(K_c\) 为轨道有界的初始点集： \[K_c = \{z \in \mathbb{C} : f_c^n(z) \not\to \infty\}\] \textbf{Julia 集} \(J_c = \partial K_c\) 为 \(K_c\) 的边界。\textbf{Mandelbrot 集} \(\mathcal{M}\) 为 \(K_c\) 连通的参数集： \[\mathcal{M} = \{c \in \mathbb{C} : f_c^n(0) \not\to \infty\}\] Julia 集和 Mandelbrot 集是复动力系统中最著名的分形对象，其自相似性、无限复杂性和与复动力学的深刻联系使它们成为数学与艺术的交汇点。

\textbf{定理 196.11}（Julia 集的基本性质）： (i) \(J_c\) 对 \(f_c\) 和 \(f_c^{-1}\) 完全不变：\(f_c(J_c) = J_c = f_c^{-1}(J_c)\)； (ii) \(J_c\) 为闭集，且 \(J_c \neq \emptyset\)； (iii) \(c \in \mathcal{M}\) 时 \(J_c\) 连通，\(c \notin \mathcal{M}\) 时 \(J_c\) 为 Cantor 集（完全不连通）； (iv) \(J_c\) 上周期点稠密，\(f_c\) 在 \(J_c\) 上具有拓扑传递性和对初值的敏感依赖性。

\textbf{证明}：(i) 由 \(K_c\) 的完全不变性及边界映射性质。(ii) \(J_c = \partial K_c\) 为边界，故闭。由 Böttcher 定理，\(|z|\) 足够大时 \(f_c^n(z) \to \infty\)，故 \(K_c\) 有界，\(J_c \neq \emptyset\)。(iii) 利用 Böttcher 映射：当 \(c \notin \mathcal{M}\) 时，临界点 \(0\) 的轨道趋于无穷，\(K_c\) 为完全不连通紧集，其边界 \(J_c\) 为 Cantor 集。当 \(c \in \mathcal{M}\) 时，\(0\) 的轨道保持有界，Böttcher 映射定义在 \(K_c\) 的补集上，将其共形映为 \(\{|z| > 1\}\)，故 \(J_c = \partial K_c\) 连通。(iv) 周期点稠密：\(f_c\) 在 \(J_c\) 上为扩张映射（\(|f_c'(z)| = 2|z| > 1\) 在 \(J_c\) 上，当 \(|z|\) 足够大时），由扩张映射的周期点稠密定理。\(\blacksquare\)

\hrulefill

\emph{混沌与分形从一维映射的符号动力学和 Sharkovsky 序出发，建立了周期轨的严格层次结构。Li-Yorke 定理和 Feigenbaum 普适性揭示了混沌的普适机制。Poincaré-Bendixson 定理和 Bendixson-Dulac 判据说明了平面连续系统不可能出现混沌。Lorenz 系统以其奇异吸引子开启了高维混沌的研究，Tucker 定理严格证明了 Lorenz 吸引子的存在性。分形几何以 Hausdorff 维数、盒维数和 IFS 为工具，给出了 Cantor 集、自相似集和 Julia 集的精确几何刻画。Mandelbrot 集作为复动力系统的参数空间分形，是数学中最复杂的对象之一。}

\emph{卷二十二：动力系统。}

\hrulefill

\textbf{反应-扩散系统的斑图形成与行波}

反应-扩散系统是描述空间分布系统中物质浓度（或种群密度）随时间演化的抛物型偏微分方程组，其最重要的特征是能够在均匀稳态附近自组织地产生空间斑图（Turing 斑图）或传播行波。本章系统阐述 Turing 不稳定性的代数条件与斑图形成机制、Keller-Segel 趋化性方程（描述细胞向化学信号浓度梯度方向的定向运动）的聚合与爆破行为，以及行波解的存在性理论（包括 Fisher-KPP 方程和双稳型反应-扩散方程的行波前沿）。反应-扩散系统的斑图理论在发育生物学（形态发生）、生态学（种群入侵）和材料科学（相场模型）中具有广泛的应用。

\subsection{239.1 Turing 不稳定性}\label{turing-ux4e0dux7a33ux5b9aux6027}

\textbf{定义 239.1}（二组分反应-扩散系统）：

\[\begin{cases} \frac{\partial u}{\partial t} = D_u \nabla^2 u + f(u, v) \\ \frac{\partial v}{\partial t} = D_v \nabla^2 v + g(u, v) \end{cases}\]

其中 \(f, g: \mathbb{R}^2 \to \mathbb{R}\) 为光滑函数，\(D_u, D_v > 0\)。

\textbf{定义 239.2}（Turing 不稳定性）：设 \((u^*, v^*)\) 为反应系统 \(\dot{u} = f(u, v), \dot{v} = g(u, v)\) 的渐近稳定平衡点。若存在波数 \(k > 0\) 使得在 \((u^*, v^*)\) 处线性化的全系统（含扩散项）的某个特征值实部为正，则称系统在 \((u^*, v^*)\) 处发生 Turing 不稳定性。充要条件是

\[\begin{cases} \operatorname{Tr}(J) = f_u + g_v < 0 \\ \det(J) = f_u g_v - f_v g_u > 0 \\ D_v f_u + D_u g_v > 2\sqrt{D_u D_v \det(J)} \end{cases}\]

\textbf{定理 239.1}（Turing 不稳定性的代数条件）：Turing 不稳定性要求 \(D_v \gg D_u\)（扩散系数显著分离），且 Jacobian 矩阵满足 \(f_u g_v < 0\)（一个自激活、一个自抑制）和 \(f_v g_u < 0\)（交叉抑制型耦合）。

\textbf{证明}：线性化全系统在 \((u^*,v^*)\) 处：\(\frac{\partial}{\partial t}\begin{pmatrix}\tilde{u} \\ \tilde{v}\end{pmatrix} = \begin{pmatrix}D_u\nabla^2+f_u & f_v \\ g_u & D_v\nabla^2+g_v\end{pmatrix}\begin{pmatrix}\tilde{u} \\ \tilde{v}\end{pmatrix}\)。对波数 \(k\) 的 Fourier 模，Laplace 算子替换为 \(-k^2\)，得 \(2\times2\) 矩阵 \(M_k = \begin{pmatrix}-D_u k^2+f_u & f_v \\ g_u & -D_v k^2+g_v\end{pmatrix}\)。特征方程 \(\lambda^2 - \operatorname{Tr}(M_k)\lambda + \det(M_k) = 0\)。Turing 不稳定性要求 \(\det(M_k) < 0\) 对某 \(k>0\) 成立，而 \(k=0\) 时 \(\det(M_0)>0\)（原系统稳定）。\(\det(M_k) = D_u D_v k^4 - (D_v f_u + D_u g_v)k^2 + \det(J)\)。作为 \(k^2\) 的二次函数，极小值在 \(k^2 = (D_v f_u + D_u g_v)/(2D_u D_v)\) 处。代入得 \(\det(M_k)_{\min}\) 为负的条件正是 \(D_v f_u + D_u g_v > 2\sqrt{D_u D_v \det(J)}\)。由 \(f_u g_v < 0\) 和 \(f_v g_u < 0\) 确保 \(D_v f_u + D_u g_v > 0\)（否则极小值不在 \(k>0\) 处）。\(\blacksquare\)

\textbf{定义 239.3}（Gierer-Meinhardt 系统------Turing 不稳定性的标准例子）：

\[\begin{cases} \frac{\partial a}{\partial t} = D_a \nabla^2 a + \rho_a \frac{a^2}{h} - \mu_a a \\ \frac{\partial h}{\partial t} = D_h \nabla^2 h + \rho_h a^2 - \mu_h h \end{cases}\]

当 \(D_h \gg D_a\) 时该系统满足 Turing 不稳定性条件，产生空间周期斑图。

\subsection{239.2 Keller-Segel 趋化性方程}\label{keller-segel-ux8d8bux5316ux6027ux65b9ux7a0b}

\textbf{定义 239.4}（Keller-Segel 系统）：抛物-椭圆型耦合 PDE

\[\begin{cases} \frac{\partial u}{\partial t} = D_u \nabla^2 u - \chi \nabla \cdot (u \nabla v) \\ \frac{\partial v}{\partial t} = D_v \nabla^2 v + \alpha u - \beta v \end{cases}, \quad \chi > 0\]

第二式为带线性源汇的扩散方程。第一式中的 \(-\chi \nabla \cdot (u \nabla v)\) 是带有浓度依赖系数的漂移-扩散项。

\textbf{定理 239.2}（二维空间的临界质量与爆破）：在 \(\mathbb{R}^2\) 上考虑 Keller-Segel 系统，存在临界质量 \(M_c = 8\pi D_u / \chi\)。若初始质量 \(M_0 = \int u_0 dx < M_c\)，解全局存在；若 \(M_0 > M_c\)，解在有限时间内爆破（\(\|u\|_{L^\infty} \to \infty\)）。临界情形 \(M_0 = M_c\) 下解存在但可能爆破。

\textbf{证明}：考虑第二矩 \(M_2(t) = \int |x|^2 u(x,t) dx\)。利用 Keller-Segel 方程，对 \(M_2\) 求导得 \(\frac{d}{dt}M_2 = 4D_u M_0 - \frac{2\chi}{\pi}M_0^2\)（在 \(\mathbb{R}^2\) 中利用 Newton 位势的 Green 函数表示）。若 \(M_0 > 8\pi D_u/\chi\)，则 \(\frac{d}{dt}M_2 < 0\)，第二矩递减。因 \(M_2 \geq 0\)，必在有限时间 \(T \leq M_2(0)/(\frac{2\chi}{\pi}M_0^2 - 4D_u M_0)\) 内达到零，此时 \(u\) 必须集中于 Dirac 测度，即 \(\|u\|_{L^\infty} \to \infty\)（爆破）。若 \(M_0 < M_c\)，则 \(M_2\) 递增，质量扩散开来，全局存在。临界情形 \(M_0 = M_c\) 时 \(\frac{d}{dt}M_2 \equiv 0\)，精细分析表明解仍可能全局存在（Dolbeault-Perthame 2004）。\(\blacksquare\)

\subsection{239.3 兴奋介质的行波------奇异摄动方法}\label{ux5174ux594bux4ecbux8d28ux7684ux884cux6ce2ux5947ux5f02ux6444ux52a8ux65b9ux6cd5}

\textbf{定义 239.5}（FitzHugh-Nagumo 方程）：二维反应-扩散系统

\[\begin{cases} \frac{\partial u}{\partial t} = D \nabla^2 u + u(1-u)(u-a) - v, \quad 0 < a < 1 \\ \frac{\partial v}{\partial t} = \varepsilon(bu - v), \quad \varepsilon \ll 1, \; b > 0 \end{cases}\]

\textbf{定理 239.3}（行波的存在性------奇异摄动分析）：当 \(\varepsilon \to 0\) 时，系统的行波解可由快-慢分解构造：快子系统（兴奋前沿）由反应-扩散波前给出，慢子系统（恢复期）由一维 ODE 的慢流形描述。行波的存在性通过对快、慢子系统的匹配渐近展开得出。

\textbf{证明}：在行波坐标 \(\xi = x - ct\) 下，系统化为 \(D u'' + c u' + u(1-u)(u-a) - v = 0\)，\(c v' = \varepsilon(bu - v)\)。当 \(\varepsilon \ll 1\) 时，\(v\) 为慢变量：在兴奋前沿（\(u\) 从 \(0\) 跃迁到 \(1\)），\(v \approx \text{const}\)，快子系统 \(D u'' + c u' + u(1-u)(u-a) - v_0 = 0\) 有行波解（Fisher-KPP 型）。波速 \(c\) 由快子系统确定：\(c \geq 2\sqrt{D(1/2-a)}\)。在恢复期，\(u \approx 1\)，慢子系统 \(c v' = \varepsilon(b - v)\) 沿 \(v\) 的慢流形演化。匹配条件：快子系统的 \(u\) 在脉冲前沿与慢子系统的 \(v\) 在脉冲尾部的值必须一致，这给出 \(c\) 与参数 \(a,b,\varepsilon\) 的隐式关系，由奇异摄动的匹配渐近展开（中间层分析）严格确立。\(\blacksquare\)

\hrulefill

\section{第325章：局部分支理论}\label{ux7b2c325ux7ae0ux5c40ux90e8ux5206ux652fux7406ux8bba}

分支理论研究向量场族在参数变化下轨道拓扑类型的定性变化------这是理解非线性系统如何响应参数扰动的核心框架。

\subsection{241.1 分支的正式定义与鞍-结点分支}\label{ux5206ux652fux7684ux6b63ux5f0fux5b9aux4e49ux4e0eux978d-ux7ed3ux70b9ux5206ux652f}

\textbf{定义 241.1}（局部分支）：考虑含参数的 \(C^r\) 向量场族

\[\dot{x} = f(x, \lambda), \quad x \in \mathbb{R}^n, \;\lambda \in \mathbb{R}^p\]

设 \((x_0, \lambda_0)\) 为平衡点（\(f(x_0, \lambda_0) = 0\)）。若在参数空间 \(\mathbb{R}^p\) 中 \(\lambda_0\) 的任何邻域内，存在 \(\lambda\) 使得 \(f(\cdot, \lambda)\) 在 \(x_0\) 附近的轨道拓扑类型与 \(f(\cdot, \lambda_0)\) 不同，则称 \(\lambda_0\) 为\textbf{分支值}。发生分支的必要条件是 \(D_x f(x_0, \lambda_0)\) 具有零实部特征值（双曲性破坏）。

\textbf{定义 241.2}（鞍-结点分支 / Fold 分支）：设 \(\dot{x} = f(x, \lambda)\)（\(\mathbb{R}^1 \times \mathbb{R}^1\)），在 \((x_0, \lambda_0)\) 处满足： 1. \(f(x_0, \lambda_0) = 0\)（平衡点条件）； 2. \(f_x(x_0, \lambda_0) = 0\)（特征值过零 --- 非双曲性条件）； 3. \(f_{xx}(x_0, \lambda_0) \neq 0\)（非退化性 --- 二次接触）； 4. \(f_{\lambda}(x_0, \lambda_0) \neq 0\)（横截性条件）。

则标准形（经坐标变换）为 \(\dot{x} = \lambda - x^2\)。

\textbf{定理 241.1}（鞍-结点分支的完整分析）：对 \(\dot{x} = \lambda - x^2\)： - \(\lambda < 0\)：无平衡点（向量场永不为零）。 - \(\lambda = 0\)：唯一平衡点 \(x = 0\)，为非双曲型（\(f_x(0,0) = 0\)，\(f_{xx}(0,0) = -2 \neq 0\)）。 - \(\lambda > 0\)：两个平衡点 \(x = \pm \sqrt{\lambda}\)，其中 \(x = +\sqrt{\lambda}\) 稳定（\(f_x = -2\sqrt{\lambda} < 0\)），\(x = -\sqrt{\lambda}\) 不稳定（\(f_x = 2\sqrt{\lambda} > 0\)）。

\textbf{证明}：直接计算。\(f(x, \lambda) = \lambda - x^2\)，平衡点 \(x^2 = \lambda\)。\(\lambda < 0\) 时无实解。\(\lambda > 0\) 时两个解，线性化 \(f_x = -2x\) 决定稳定性。横截性由 \(\frac{\partial}{\partial \lambda} f(0,0) = 1 \neq 0\) 满足；非退化由 \(f_{xx}(0,0) = -2 \neq 0\) 满足。一般情形通过隐函数定理与 Morse 引理化归为该标准形。\(\blacksquare\)

\subsection{241.2 Hopf 分支}\label{hopf-ux5206ux652f}

\textbf{定义 241.3}（Hopf 分支）：设 \(\dot{x} = f(x, \lambda)\)（\(\mathbb{R}^2 \times \mathbb{R}^1\) 或更高维），\(f(x_0, \lambda_0) = 0\)，\(D_x f(x_0, \lambda_0)\) 有一对纯虚特征值 \(\pm i \omega\)（\(\omega > 0\)），其余特征值均有非零实部。若穿过 \(\lambda_0\) 时该对复特征值横截穿过虚轴（\(\frac{d}{d\lambda} \operatorname{Re} \mu(\lambda)|_{\lambda = \lambda_0} \neq 0\)），则发生 Hopf 分支：在 \(\lambda_0\) 的某一侧（具体侧取决于第一 Lyapunov 系数 \(\ell_1\) 的符号）从平衡点分支出小振幅周期轨。

\textbf{定理 241.2}（Hopf 分支的标准形分析）：对 Hopf 标准形（极坐标）

\[\begin{cases} \dot{r} = r(\lambda - r^2) \\ \dot{\theta} = 1 \end{cases}\]

分析： - \(\lambda \leq 0\)：唯一平衡点 \(r = 0\)，全局渐近稳定（焦点）。 - \(\lambda > 0\)：\(r = 0\) 变为不稳定（源），产生唯一稳定极限环 \(r = \sqrt{\lambda}\)（超临界 Hopf 分支）。若将 \(r(\lambda - r^2)\) 换为 \(r(\lambda + r^2)\)，则 \(\lambda < 0\) 时存在不稳定极限环（亚临界 Hopf 分支）。

\textbf{证明}：在极坐标下，\(r\) 方程 \(r' = r(\lambda - r^2)\) 是 Bernoulli 型 ODE。令 \(w = r^{-2}\)，则 \(w' = -2r^{-3}r' = -2r^{-3} \cdot r(\lambda - r^2) = -2\lambda w + 2\)。解为 \(w(t) = (w_0 - 1/\lambda)e^{-2\lambda t} + 1/\lambda\)（\(\lambda \neq 0\)）。若 \(\lambda \leq 0\)，\(w(t) \to \infty\)，故 \(r(t) \to 0\)，即 \(r=0\) 全局渐近稳定。若 \(\lambda > 0\)，\(w(t) \to 1/\lambda\)，故 \(r(t) \to \sqrt{\lambda}\)。对于初值 \(r_0 > 0\)，\(r(t)\) 单调收敛到 \(\sqrt{\lambda}\)，故 \(r=\sqrt{\lambda}\) 是稳定极限环。\(\theta' = 1\) 保证角速度恒定，周期 \(T = 2\pi\)。对于亚临界情形 \(r' = r(\lambda + r^2)\)，\(r=0\) 在 \(\lambda>0\) 时不稳定，\(\lambda<0\) 时存在不稳定极限环 \(r=\sqrt{-\lambda}\)。\(\blacksquare\)

\textbf{定义 241.4}（第一 Lyapunov 系数 \(\ell_1\)）：将系统 \(\dot{x} = A(\lambda)x + F(x, \lambda)\)（\(A(\lambda_0)\) 有特征值 \(\pm i \omega\)）通过中心流形约化到二维，再经规范形变换到复坐标 \(z\)：\(\dot{z} = i \omega z + \ell_1 |z|^2 z + O(|z|^4)\)。\(\ell_1\) 为实常数，正负号决定分支方向：\(\ell_1 < 0\) 对应超临界（稳定极限环），\(\ell_1 > 0\) 对应亚临界（不稳定极限环）。\(\ell_1\) 可由 \(f\) 在 \((x_0, \lambda_0)\) 处的三阶 Taylor 系数显式计算。

\subsection{241.3 叉形分支与跨临界分支}\label{ux53c9ux5f62ux5206ux652fux4e0eux8de8ux4e34ux754cux5206ux652f}

\textbf{定义 241.5}（叉形分支 / Pitchfork）：标准形 \(\dot{x} = \lambda x - x^3\)（超临界）或 \(\dot{x} = \lambda x + x^3\)（亚临界）。条件：\(f_x(x_0, \lambda_0) = 0\)，\(f_{xxx}(x_0, \lambda_0) \neq 0\)，\(f_{x\lambda}(x_0, \lambda_0) \neq 0\)，且系统具有奇对称性 \(f(-x, \lambda) = -f(x, \lambda)\)。超临界情形：\(\lambda < 0\) 时 \(x=0\) 稳定；\(\lambda > 0\) 时 \(x=0\) 不稳定且产生一对对称稳定平衡点 \(x = \pm \sqrt{\lambda}\)。

叉形分支的分支图呈''三叉戟''形状，故得名。超临界叉形分支是一种''软''失稳------系统在分岔后平滑地过渡到新的稳定状态，没有跳跃或滞后现象。这在物理系统中对应二阶相变（如铁磁材料的自发磁化）。亚临界叉形分支（\(\dot{x} = \lambda x + x^3\)）则是一种''硬''失稳：\(\lambda < 0\) 时 \(x=0\) 稳定，但与 \(x=\pm\sqrt{-\lambda}\) 两个不稳定平衡点共存；\(\lambda > 0\) 时三个平衡点全部消失，系统跳变到远处------这典型地对应一阶相变。在工程中，亚临界叉形分支对应突然屈曲（snap-through buckling），具有灾难性后果。

\textbf{定义 241.6}（跨临界分支 / Transcritical）：标准形 \(\dot{x} = \lambda x - x^2\)。条件：\(f_x(x_0, \lambda_0) = 0\)，\(f_{xx}(x_0, \lambda_0) \neq 0\)，\(f_{x\lambda}(x_0, \lambda_0) \neq 0\)，且存在贯穿 \(\lambda\) 的不变平衡点（通常由守恒律保证）。两个平衡点 \(x=0\) 与 \(x=\lambda\) 在 \(\lambda=0\) 处交换稳定性。

跨临界分支的特征是存在一个始终存在的''平凡''平衡点，其稳定性在分支点处反转。在种群动力学中，跨临界分支出现在 Logistic 方程 \(\dot{N} = rN(1 - N/K)\) 中以 \(r\) 为分支参数时：\(N=0\)（灭绝平衡点）在 \(r<0\) 时稳定（负增长）、\(r>0\) 时不稳定；而 \(N=K\)（承载容量平衡点）恰好相反。在流行病学中，\(R_0=1\) 处的流行病阈值（无病平衡点与地方病平衡点交换稳定性）也是跨临界分支的典型实例。跨临界分支的前提条件是系统的某个守恒量（如总质量或总概率）强制存在一个平凡不变解，这使得经典的鞍结分支条件在此不成立。

\subsection{241.4 中心流形定理与约化原理}\label{ux4e2dux5fc3ux6d41ux5f62ux5b9aux7406ux4e0eux7ea6ux5316ux539fux7406}

\textbf{定理 241.3}（中心流形定理，Kelley 1967，Carr 1981）：设 \(\dot{x} = f(x)\)（\(f \in C^r\)，\(f(0) = 0\)），\(Df(0)\) 的特征值分为三类：\(n_s\) 个具负实部、\(n_u\) 个具正实部、\(n_c\) 个具零实部。则存在穿过 \(0\) 的 \(C^r\) 局部不变流形 \(W^c_{\operatorname{loc}}\)（中心流形，维数 \(n_c\)），其切于中心广义特征空间 \(E^c\)。在 \(0\) 的邻域内，系统拓扑共轭于

\[\begin{cases} \dot{x}_c = \tilde{f}_c(x_c) \\ \dot{x}_s = -x_s \\ \dot{x}_u = x_u \end{cases}\]

即非线性动力学完全由中心流形上的约化系统决定。在分支分析中，将全系统投影到中心流形上将维数从 \(n\) 降至 \(n_c\)（通常 \(n_c \leq 2\)），保留所有分支信息。

中心流形定理是局部分支理论的基石。其核心思想是：在分支点处，线性化算子具有零实部的特征值，对应的特征空间（中心子空间）刻画了系统失去双曲性的方向。稳定和不稳定方向上的动力学是指数衰减或增长的，因此仅提供瞬态行为；而所有非平凡的渐近行为------包括分支------都局限在中心流形上。约化原理（Reduction Principle）断言：原 \(n\) 维系统的局部分支问题等价于 \(n_c\) 维中心流形上的约化系统的分支问题，且两者的分支类型（鞍结、Hopf、叉形等）和稳定性完全一致。这使得我们可以将高维系统的分支分析简化为低维系统的分析，通常 \(n_c = 1\)（实特征值过零）或 \(n_c = 2\)（共轭纯虚特征值）。中心流形的计算通常通过幂级数展开近似：设 \(h(x_c) = \sum_{|\alpha| \geq 2} c_\alpha x_c^\alpha\)，代入不变性方程匹配各阶系数。工程和物理中，中心流形约化是推导 Landau 方程和振幅方程（如 Ginzburg-Landau 方程）的标准方法，后者描述了分支点附近有序参量的慢变动力学。

\textbf{证明}：在 \(E^c \oplus E^s \oplus E^u\) 分解下，设 \(x = (x_c, x_s, x_u)\)。系统写为 \(\dot{x}_c = A_c x_c + f_c(x_c, x_s, x_u)\)，\(\dot{x}_s = A_s x_s + f_s(x)\)，\(\dot{x}_u = A_u x_u + f_u(x)\)，其中 \(A_c\) 的特征值纯虚，\(A_s\) 的特征值负实部，\(A_u\) 的特征值正实部。中心流形可表为图 \(x_s = h_s(x_c)\)，\(x_u = h_u(x_c)\)，满足不变性条件 \(\dot{x}_s = D h_s(x_c) \dot{x}_c\)。代入得拟线性 PDE：\(A_s h_s(x_c) + f_s(x_c, h_s(x_c), h_u(x_c)) = D h_s(x_c)[A_c x_c + f_c(x_c, h_s(x_c), h_u(x_c))]\)。在适当的 Banach 空间（如 \(C^r\) 或 Hölder 空间）中，该方程可写为不动点问题 \(h_s = \Phi(h_s, h_u)\)（类似于 \(h_u\)）。由压缩映射原理，在 \(0\) 的充分小邻域内存在唯一解 \(h_s, h_u\) 满足 \(h_s(0)=0\)，\(D h_s(0)=0\)，\(h_u(0)=0\)，\(D h_u(0)=0\)。\(C^r\) 光滑性由截断函数和不动点定理的隐函数定理版本保证。约化系统 \(\dot{x}_c = A_c x_c + f_c(x_c, h_s(x_c), h_u(x_c))\) 完全决定了原系统的局部分支行为。\(\blacksquare\)

\subsection{241.5 全局分支与 Melnikov 方法}\label{ux5168ux5c40ux5206ux652fux4e0e-melnikov-ux65b9ux6cd5}

\textbf{定义 241.7}（同宿分支）：含参数向量场 \(\dot{x} = f(x, \lambda)\) 在 \(\lambda = \lambda_0\) 时具有连向同一鞍点 \(p(\lambda_0)\) 的同宿轨 \(\Gamma\)（即 \(\lim_{t \to \pm \infty} \gamma(t) = p\)）。当 \(\lambda\) 穿过 \(\lambda_0\) 时，同宿轨可能破裂并在其附近产生周期轨或混沌动力学------此为同宿分支。

\textbf{定理 241.4}（Melnikov 方法，1963）：考虑平面周期扰动 Hamilton 系统

\[\begin{cases} \dot{x} = f_1(x,y) + \varepsilon g_1(x,y,t), \quad f_1 = \partial H / \partial y \\ \dot{y} = f_2(x,y) + \varepsilon g_2(x,y,t), \quad f_2 = -\partial H / \partial x \end{cases}\]

设未扰动系统（\(\varepsilon = 0\)）有鞍点 \(p_0\) 及其同宿轨 \(\Gamma_0(t)\)。则 Melnikov 函数

\[M(t_0) = \int_{-\infty}^{\infty} \left( f_1(\Gamma_0(t)) g_2(\Gamma_0(t), t+t_0) - f_2(\Gamma_0(t)) g_1(\Gamma_0(t), t+t_0) \right) dt\]

刻画扰动后稳定流形与不稳定流形沿 \(\Gamma_0\) 法向的距离。若 \(M(t_0)\) 有简单零点，则对充分小的 \(\varepsilon > 0\)，稳定与不稳定流形横截相交，产生 Smale 马蹄型混沌。Melnikov 方法是判定具体系统中同宿混沌存在的少数可计算准则之一。

\textbf{证明}：设未扰动系统有同宿轨 \(\Gamma_0(t) = (x_0(t), y_0(t))\) 满足 \(\lim_{t\to\pm\infty}\Gamma_0(t) = p_0\)。扰动后，稳定流形 \(W^s_\varepsilon(p_\varepsilon)\) 与不稳定流形 \(W^u_\varepsilon(p_\varepsilon)\) 沿 \(\Gamma_0\) 法向的距离 \(d(t_0)\) 可展开为 \(d(t_0) = \varepsilon M(t_0)/\|f(\Gamma_0(0))\| + O(\varepsilon^2)\)。计算：沿 \(\Gamma_0\) 的法向量为 \(n(t) = (-f_2(\Gamma_0(t)), f_1(\Gamma_0(t)))\)。\(M(t_0) = \int_{-\infty}^\infty n(t) \cdot g(\Gamma_0(t), t+t_0) dt\) 即为稳定与不稳定流形分离的 \(O(\varepsilon)\) 首项。若 \(M(t_0^*)=0\) 且 \(M'(t_0^*) \neq 0\)（简单零点），则对于充分小的 \(\varepsilon\)，\(W^s_\varepsilon\) 与 \(W^u_\varepsilon\) 横截相交。由 Smale-Birkhoff 定理，横截同宿点蕴含马蹄型混沌动力学。\(\blacksquare\)

\subsection{241.6 遍历论在分支中的应用}\label{ux904dux5386ux8bbaux5728ux5206ux652fux4e2dux7684ux5e94ux7528}

在含参系统通过分支产生混沌动力学后，可由遍历论工具定量刻画混沌吸引子的统计性质。核心联系为： 1. \textbf{SRB 测度}（Sinai-Ruelle-Bowen）：在分支产生的混沌吸引子上，存在唯一的物理测度 \(\mu\)，其在吸引盆的 Lebesgue 正测度子集上满足：对连续观测函数 \(f\)，\(\frac{1}{n}\sum_{k=0}^{n-1} f(T^k x) \to \int f d\mu\)。SRB 测度为 Kolmogorov-Sinai 熵与正 Lyapunov 指数之和的 Pesin 公式提供了语境。 2. \textbf{周期轨道的 Zeta 函数}：分支后混沌吸引子中各周期轨道的分布满足 \(\zeta(z) = \exp\left(\sum_{n=1}^\infty \frac{z^n}{n} \sum_{x \in \operatorname{Fix}(T^n)} \prod_{k=0}^{n-1} |\det(DT(T^k x))|^{-1}\right)\)，其解析性质与 KS-熵相连。 3. \textbf{重正化群与普适性}（Feigenbaum 理论）：周期倍化分支的累积速率具有普适常数 \(\delta \approx 4.6692\) 和 \(\alpha \approx 2.5029\)，由函数空间中的重正化算子不动点的单一不稳定特征值导出。这说明了具有相同''遍历类''的映射在分支结构上的共性。

\hrulefill

\newpage

\chapter{09-微分方程与动力系统/02-动力系统/05-数学生物学中的动力系统方法.md}\label{ux5faeux5206ux65b9ux7a0bux4e0eux52a8ux529bux7cfbux7edf02-ux52a8ux529bux7cfbux7edf05-ux6570ux5b66ux751fux7269ux5b66ux4e2dux7684ux52a8ux529bux7cfbux7edfux65b9ux6cd5.md}

\section{第326章：生化网络的数学建模------随机过程与奇异摄动}\label{ux7b2c326ux7ae0ux751fux5316ux7f51ux7edcux7684ux6570ux5b66ux5efaux6a21ux968fux673aux8fc7ux7a0bux4e0eux5947ux5f02ux6444ux52a8}

生化网络的数学建模是系统生物学与动力系统理论交叉的前沿领域。本章通过奇异摄动理论（Tikhonov 定理）和随机过程方法，系统研究酶促反应网络（Michaelis-Menten 动力学）和基因调控网络（Goodwin 振荡器）的动力学行为。核心方法是：利用生化反应中各组分浓度差异悬殊的特点（如酶浓度远小于底物浓度），通过引入小参数将系统分解为快慢子系统，应用 Tikhonov 定理将快变量消除，得到约化后的慢动力学方程。这种方法不仅提供了 Michaelis-Menten 准稳态近似（QSSA）的严格数学基础，还揭示了基因调控网络中 Hopf 分岔产生持续振荡的阈值条件。随机过程方法（如化学主方程和 Gillespie 算法）则捕捉了小分子数情形下的随机涨落效应，是理解基因表达噪声和表型异质性的关键工具。

\subsection{240.1 Michaelis-Menten 动力学的奇异摄动分析}\label{michaelis-menten-ux52a8ux529bux5b66ux7684ux5947ux5f02ux6444ux52a8ux5206ux6790}

\textbf{定义 240.1}（酶促反应网络）：三物质反应系统

\[E + S \underset{k_{-1}}{\overset{k_1}{\rightleftharpoons}} ES \overset{k_{\text{cat}}}{\longrightarrow} E + P\]

对应的 ODE 系统包括四个浓度变量，但受两个守恒律约束，实际自由度为二。

\textbf{定理 240.1}（准稳态约化 / QSSA）：在条件 \([S] \gg [E]_T\) 下，中间产物 \([ES]\) 的动力学比 \([S]\) 快得多，应用奇异摄动理论可约化为一维系统

\[\frac{d[S]}{dt} = -\frac{V_{\max}[S]}{K_m + [S]}, \quad V_{\max} = k_{\text{cat}}[E]_T, \quad K_m = \frac{k_{-1} + k_{\text{cat}}}{k_1}\]

\textbf{证明}：由酶守恒律 \([E]_T = [E] + [ES]\) 及质量作用律， \[\frac{d[ES]}{dt} = k_1[E][S] - (k_{-1}+k_{\text{cat}})[ES] = k_1([E]_T-[ES])[S] - (k_{-1}+k_{\text{cat}})[ES].\] 定义小参数 \(\varepsilon = [E]_T/[S]_0 \ll 1\)，引入快时间 \(\tau = k_1[S]_0 t\) 和慢时间 \(T = \varepsilon t\)。重标度变量后得奇异摄动系统。由 Tikhonov 定理，令 \(\varepsilon \to 0\) 时快变量 \([ES]\) 迅速松弛到慢流形 \[k_1([E]_T-[ES])[S] - (k_{-1}+k_{\text{cat}})[ES] \approx 0 \;\Longrightarrow\; [ES] = \frac{[E]_T[S]}{K_m+[S]}, \quad K_m = \frac{k_{-1}+k_{\text{cat}}}{k_1}.\] 代入 \([S]\) 的动力学方程 \(d[S]/dt = -k_1[E][S] + k_{-1}[ES]\) 并利用上述关系即得 Michaelis-Menten 约化方程 \[\frac{d[S]}{dt} = -\frac{V_{\max}[S]}{K_m+[S]}, \quad V_{\max}=k_{\text{cat}}[E]_T.\] 该约化在 \(\varepsilon \to 0\) 的极限下渐近成立，Tikhonov 定理保证慢流形的指数吸引性。\(\blacksquare\)

\subsection{240.2 基因调控网络的 ODE 模型与分岔}\label{ux57faux56e0ux8c03ux63a7ux7f51ux7edcux7684-ode-ux6a21ux578bux4e0eux5206ux5c94}

\textbf{定义 240.2}（Goodwin 振荡器）：二维 ODE 系统

\[\begin{cases} \dot{m} = \frac{\alpha}{1 + p^n / K^n} - \gamma_m m \\ \dot{p} = \beta m - \gamma_p p \end{cases}, \quad n \in \mathbb{N}, \; \alpha, \beta, \gamma_m, \gamma_p, K > 0\]

右端第一项为 Hill 型递减函数（负反馈），\(n\) 为 Hill 系数。

\textbf{定理 240.2}（Hopf 分岔产生振荡）：当 Hill 系数 \(n\) 充分大时（在标准化参数下典型阈值为 \(n > 8\)），系统通过 Hopf 分岔产生稳定的极限环------即持续周期振荡。分岔条件由 Jacobian 矩阵纯虚数特征值条件 \(\operatorname{Tr}(J) = 0\) 给出。

\textbf{证明}：在平衡点 \((m^*, p^*)\) 处线性化，将 \(f(p)=\alpha/(1+p^n/K^n)\) 记为 Hill 递减函数，则 Jacobian 矩阵为 \[J = \begin{pmatrix} -\gamma_m & -f'(p^*) \\ \beta & -\gamma_p \end{pmatrix}, \quad f'(p^*) = -\frac{\alpha n p^{*n-1}}{K^n(1+p^{*n}/K^n)^2} < 0.\] 特征方程 \(\lambda^2 - \operatorname{Tr}(J)\lambda + \det(J) = 0\)，其中 \(\operatorname{Tr}(J) = -(\gamma_m+\gamma_p) < 0\) 恒负，\(\det(J) = \gamma_m\gamma_p + \beta f'(p^*) > 0\)。特征值为 \(\lambda = \frac{1}{2}(\operatorname{Tr}(J) \pm \sqrt{\operatorname{Tr}(J)^2 - 4\det(J)})\)。当 \(n\) 增大时 \(|f'(p^*)|\) 增加，\(\det(J)\) 增大，判别式 \(\operatorname{Tr}(J)^2 - 4\det(J)\) 从正变负，特征值从两负实根转变为一对共轭复根。对标准化参数，当 \(n\) 超过临界值（典型为 \(n>8\)）时，系统在参数空间中越过 Hopf 分岔曲线。由中心流形定理约化到二维，计算第一 Lyapunov 系数 \(\ell_1 = \frac{1}{16}(f_{xxx}+f_{xyy}+g_{xxy}+g_{yyy}) + \cdots\)（在规范形坐标下），对于 Goodwin 系统 \(\ell_1 < 0\)，故分岔为超临界，产生稳定的极限环振荡。\(\blacksquare\)

\subsection{240.3 离散动力系统与布尔网络}\label{ux79bbux6563ux52a8ux529bux7cfbux7edfux4e0eux5e03ux5c14ux7f51ux7edc}

前面讨论的连续时间 ODE 系统（如 Goodwin 振子）描述了通过微分方程耦合的连续状态动力学。然而在生物学、计算机科学和复杂系统研究中，许多动力学自然地以离散时间和离散状态的形式出现------其中布尔网络是最简洁、同时蕴含最丰富现象的模型框架。Stuart Kauffman 在 1969 年的开创性工作中首次揭示了布尔网络中的有序-混沌相变现象，预言生命系统运作在''混沌边缘''的临界状态。

\textbf{定义 240.3}（\(N\)-\(K\) 布尔网络）：\(N\) 个二值变量 \(\{x_i \in \{0, 1\}\}_{i=1}^N\)，每个 \(x_i\) 的更新由 \(K\) 个随机选取的输入节点经随机布尔函数确定。状态空间大小为 \(2^N\)，动力学由有向状态转移图描述。

\textbf{定理 240.3}（Kauffman 的有序-混沌相变）：当 \(K = 1\) 时系统处于冻结有序相（吸引子长度和数量均较小）；\(K \geq 3\) 时系统处于混沌相（吸引子长度随 \(N\) 指数增长）；\(K = 2\) 为临界值（``混沌边缘''），此时吸引子数量约为 \(\sqrt{N}\)。相变由 Lyapunov 指数或 Hamming 距离的传播刻画。

\textbf{证明}：考虑两个初始仅差一个比特的状态 \(\mathbf{x}(0)\) 和 \(\mathbf{y}(0)\)，Hamming 距离 \(d(t)=|\{i:x_i(t)\neq y_i(t)\}|\)。在随机布尔函数假设下，每步每个受扰节点以概率 \(p_K = 1-(1-1/2)^{K} \approx K/2^{K}\)（对小数 K）向邻近节点传播扰动，期望演化 \(\mathbb{E}[d(t+1)] \approx d(t) \cdot K \cdot \mathbb{E}[\text{灵敏度}]\)。\(K=1\)：每个节点仅依赖一个输入，灵敏度 \(=1/2\)，\(\mathbb{E}[d(t+1)]=\frac{1}{2}d(t)\)，扰动指数衰减至 0 ------ 冻结相。\(K\geq 3\)：通过布尔函数的随机偏导数分析，平均灵敏度 \(>1\)，\(\mathbb{E}[d(t+1)] = \lambda d(t)\)（\(\lambda>1\)），扰动指数放大 ------ 混沌相。\(K=2\)：临界灵敏度 \(=1\)，扰动边缘传播，相变由渗流论和自旋玻璃理论刻画。Lyapunov 指数 \(\Lambda = \lim_{t\to\infty}\frac{1}{t}\ln d(t)\) 在 \(K<2\) 时为负（有序）、\(K>2\) 时为正（混沌）、\(K=2\) 时为零（临界）。\(\blacksquare\)

\subsection{240.4 化学主方程与 Markov 过程}\label{ux5316ux5b66ux4e3bux65b9ux7a0bux4e0e-markov-ux8fc7ux7a0b}

\textbf{定义 240.4}（化学主方程 / Chemical Master Equation）：设系统含 \(d\) 种分子，状态空间 \(\mathbb{Z}_{\geq 0}^d\)，\(R\) 个反应通道。反应 \(j\) 的倾向函数为 \(a_j(\mathbf{x})\)，化学计量变化向量为 \(\nu_j \in \mathbb{Z}^d\)。概率分布 \(P(\mathbf{x}, t)\) 满足无穷维线性 ODE

\[\frac{dP(\mathbf{x}, t)}{dt} = \sum_{j=1}^R \left[a_j(\mathbf{x} - \nu_j) P(\mathbf{x} - \nu_j, t) - a_j(\mathbf{x}) P(\mathbf{x}, t)\right]\]

这是连续时间 Markov 链的 Kolmogorov 前向方程。

\textbf{定理 240.4}（Gillespie 算法的数学基础）：化学主方程对应的 Markov 链可通过以下方式精确模拟------从当前状态 \(\mathbf{x}\) 出发，(i) 等待时间 \(\tau\) 服从指数分布 \(\operatorname{Exp}(a_0(\mathbf{x}))\)，其中 \(a_0 = \sum_j a_j\)；(ii) 反应 \(j\) 被选中的概率为 \(a_j/a_0\)。该算法是 Doob-Gillespie 型随机模拟算法，精确生成 Markov 链的样本路径。

\textbf{证明}：设当前状态为 \(\mathbf{x}\)。化学主方程对应的 Markov 链的无穷小生成元为 \((\mathcal{L}P)(\mathbf{x}) = \sum_j a_j(\mathbf{x})[P(\mathbf{x}+\nu_j) - P(\mathbf{x})]\)。从状态 \(\mathbf{x}\) 出发，系统在 \(R\) 个独立竞争的反应通道中，每个反应 \(j\) 的等待时间 \(\tau_j \sim \operatorname{Exp}(a_j(\mathbf{x}))\) 相互独立。实际反应发生的时间为 \(\tau = \min_j \tau_j\)，由指数分布的最小值性质，\(\tau \sim \operatorname{Exp}(\sum_j a_j(\mathbf{x})) = \operatorname{Exp}(a_0(\mathbf{x}))\)。发生反应 \(j\) 的概率为 \(P(\tau_j = \min_k \tau_k) = a_j(\mathbf{x})/a_0(\mathbf{x})\)（由指数分布的无记忆性和竞争风险模型）。更新状态 \(\mathbf{x} \mapsto \mathbf{x} + \nu_j\) 后，重新计算所有倾向函数并重复此过程。此抽样过程精确生成 Markov 过程 \(\mathbf{X}(t)\) 的样本路径，因为每一步都严格遵循转移核 \(P(\mathbf{X}(t+\Delta t)=\mathbf{x}+\nu_j \mid \mathbf{X}(t)=\mathbf{x}) = a_j(\mathbf{x})\Delta t + o(\Delta t)\)。Gillespie 算法是此过程的直接蒙特卡洛实现，不引入时间离散化或 \(\tau\)-leaping 近似，故为精确算法。\(\blacksquare\)

\hrulefill

\emph{本卷在常微分方程定性理论、反应-扩散方程和随机过程的基础上，系统建立了非线性 ODE 系统定性分析（Logistic 方程、Lotka-Volterra 系统、竞争模型）、仓室型 ODE 系统的阈值动力学（SIR 模型及其变体）、反应-扩散系统的斑图形成与行波（Turing 不稳定性、Keller-Segel 爆破、FitzHugh-Nagumo 奇异摄动）以及生化网络的数学建模（奇异摄动 QSSA、Hopf 分岔、布尔网络、化学主方程）的理论。共 4 章。}

\hrulefill

\hrulefill

\hrulefill

\hrulefill

\hrulefill

\hrulefill

\subsection{268.A 数学生物学补充2}\label{a-ux6570ux5b66ux751fux7269ux5b66ux8865ux51452}

\begin{quote}
动力系统研究随时间演化的系统的长期行为，起源于 Newton 力学和 Poincaré 对天体力学的研究。本卷在常微分方程（V18）的基础上，建立离散和连续动力系统的基本理论，包括稳定性、混沌、遍历论和分支理论。动力系统为理解自然界中的复杂动力学现象提供了统一的数学框架。
\end{quote}

\begin{quote}
\textbf{前置知识}：本子卷依赖本卷第201-210章（常微分方程和偏微分方程基础）和第90-94章（拓扑学）的核心内容。
\end{quote}

\hrulefill

\hrulefill

\hrulefill

\hrulefill

\hrulefill

\hrulefill

\subsection{262.A 数学生物学补充}\label{a-ux6570ux5b66ux751fux7269ux5b66ux8865ux5145}

特殊函数（Special Functions）是数学物理和微分方程理论的经典工具。它们通常作为二阶线性常微分方程的解出现，构成一套丰富的函数体系。本章系统补充 Gamma/Beta 函数、Bessel 函数和正交多项式族的理论基础。

\subsection{Gamma函数与Beta函数}\label{gammaux51fdux6570ux4e0ebetaux51fdux6570}

\textbf{定义}（Gamma函数 / Euler第二积分）：\(\Gamma(z)\) 是阶乘函数在复平面上的解析延拓（除去非正整数点），第二类 Euler 积分为： \[\Gamma(z) = \int_0^\infty t^{z-1} e^{-t} dt \quad (\Re z > 0)\] 通过解析延拓到 \(\mathbb{C} \setminus \{0, -1, -2, \ldots\}\)。

\textbf{命题}（基本性质）： - 函数方程：\(\Gamma(z+1) = z \Gamma(z)\)，\(\Gamma(1) = 1\)，故 \(\Gamma(n) = (n-1)!\)（\(n \in \mathbb{N}\)） - 反射公式（Euler 反射公式）：\(\Gamma(z) \Gamma(1-z) = \frac{\pi}{\sin(\pi z)}\) - Legendre 加倍公式：\(\Gamma(2z) = \frac{2^{2z-1}}{\sqrt{\pi}} \Gamma(z) \Gamma(z + 1/2)\) - \textbf{Weierstrass无穷乘积表示}：\(\frac{1}{\Gamma(z)} = z e^{\gamma z} \prod_{n=1}^\infty \left(1 + \frac{z}{n}\right) e^{-z/n}\)（\(\gamma \approx 0.5772\) 为 Euler-Mascheroni 常数）。由此可知 \(1/\Gamma(z)\) 是整函数，其零点是 \(z = 0, -1, -2, \ldots\)（均为单零点）

\textbf{定义}（Beta函数 / Euler第一积分）： \[B(x, y) = \int_0^1 t^{x-1} (1-t)^{y-1} dt \quad (\Re x, \Re y > 0)\] 其与 Gamma 函数的关系为：\(B(x, y) = \frac{\Gamma(x) \Gamma(y)}{\Gamma(x + y)}\)

\textbf{定理}（Stirling公式 / 渐近展开）：当 \(|z| \to \infty\)（\(|\arg z| < \pi\)）， \[\Gamma(z) \sim \sqrt{2\pi} \, z^{z - 1/2} e^{-z} \left[1 + \frac{1}{12z} + \frac{1}{288 z^2} - \frac{139}{51840 z^3} + \cdots \right]\] \(\log \Gamma(z) \sim z \log z - z - \frac{1}{2} \log z + \frac{1}{2} \log(2\pi) + O(1/z)\)。

\textbf{证明}：由 Binet 第二公式 \(\log\Gamma(z) = (z-\frac{1}{2})\log z - z + \frac{1}{2}\log(2\pi) + \int_0^\infty \frac{\varphi(t)}{t}e^{-zt}dt\)（\(\varphi(t)=\frac{1}{2}-\frac{1}{t}+\frac{1}{e^t-1}\)）。将被积函数在 \(t=0\) 处 Taylor 展开：\(\frac{\varphi(t)}{t}=\sum_{k=1}^\infty \frac{B_{2k}}{(2k)!}\frac{t^{2k-2}}{2k-1}\)（\(B_{2k}\) 为 Bernoulli 数）。代入积分，利用 \(\int_0^\infty t^{2k-2}e^{-zt}dt = \frac{\Gamma(2k-1)}{z^{2k-1}} = \frac{(2k-2)!}{z^{2k-1}}\)，得渐近级数： \[\log\Gamma(z) = (z-\tfrac{1}{2})\log z - z + \tfrac{1}{2}\log(2\pi) + \sum_{k=1}^\infty \frac{B_{2k}}{2k(2k-1)z^{2k-1}}.\] 指数化即得 Stirling 公式。主项来自鞍点近似：\(\Gamma(z)=\int_0^\infty e^{z\log t - t}dt/\sqrt{z}\)，最大点 \(t=z\)，Laplace 方法给出前因子 \(\sqrt{2\pi}\,z^{z-1/2}e^{-z}\)。误差项由 Euler-Maclaurin 求和公式严格控制。\(\blacksquare\)

\subsection{Bessel函数}\label{besselux51fdux6570}

\textbf{定义}（Bessel微分方程与Bessel函数）：\(\nu\) 阶 Bessel 微分方程 \[x^2 y'' + x y' + (x^2 - \nu^2) y = 0\] 的第一类解为 \textbf{Bessel 函数} \(J_\nu(x)\)（级数形式）： \[J_\nu(x) = \sum_{k=0}^\infty \frac{(-1)^k}{k! \, \Gamma(k + \nu + 1)} \left(\frac{x}{2}\right)^{2k + \nu}\] 当 \(\nu \notin \mathbb{Z}\) 时，\(J_\nu\) 与 \(J_{-\nu}\) 是线性无关的。

\textbf{命题}（整数阶与半整数阶）： - \(J_{-n}(x) = (-1)^n J_n(x)\)（\(n \in \mathbb{Z}\)），此时 \(J_n\) 和 \textbf{第二类Bessel函数} \(Y_n\)（Neumann函数）构成基 - 半整数阶推出初等函数表示：\(J_{1/2}(x) = \sqrt{\frac{2}{\pi x}} \sin x\)，\(J_{n+1/2}\) 均可通过正弦和余弦及 \(x^{-1/2}\) 的组合表出

\textbf{命题}（递推关系与生成函数）： - \(\frac{d}{dx}[x^\nu J_\nu(x)] = x^\nu J_{\nu-1}(x)\) - \(\frac{d}{dx}[x^{-\nu} J_\nu(x)] = -x^{-\nu} J_{\nu+1}(x)\) - \textbf{生成函数}：\(\exp\left[\frac{x}{2}(t - t^{-1})\right] = \sum_{n=-\infty}^\infty J_n(x) t^n\)（整数阶Bessel函数的指数型生成函数，Schlömilch 1857）

\textbf{定理}（Bessel函数的正交性与Fourier-Bessel展开）：\(J_\nu\) 在 \([0, a]\) 上的零点 \(\{j_{\nu, n}\}\) 满足正交关系： \[\int_0^a J_\nu(j_{\nu,m} r/a) J_\nu(j_{\nu,n} r/a) \, r \, dr = \frac{a^2}{2} [J_\nu'(j_{\nu,n})]^2 \delta_{mn}\] 由此得到 \textbf{Fourier-Bessel级数展开}------将圆柱对称区域上的函数按Bessel函数展开（Hankel变换）。

\textbf{证明}：Bessel 方程可写为 Sturm-Liouville 形式 \((x J_\nu'(x))' + (x - \nu^2/x)J_\nu(x) = 0\)。令 \(\phi_n(x) = J_\nu(j_{\nu,n} x/a)\)（\(j_{\nu,n}\) 为 \(J_\nu\) 的第 \(n\) 个正零点），则 \(\phi_n\) 满足 \((x\phi_n')' + (j_{\nu,n}^2 x/a^2 - \nu^2/x)\phi_n = 0\)，即 \(-(x\phi_n')' + (\nu^2/x)\phi_n = \lambda_n x\phi_n\)，其中 \(\lambda_n = j_{\nu,n}^2/a^2\)。这是带权 \(w(x)=x\) 的 Sturm-Liouville 问题。对 \(m \neq n\)，计算 \((\lambda_m - \lambda_n)\int_0^a \phi_m(x)\phi_n(x) x dx = \int_0^a [\phi_n(x\phi_m')' - \phi_m(x\phi_n')']dx = [x(\phi_n\phi_m' - \phi_m\phi_n')]_0^a = 0\)（边界项因 \(x=0\) 时 \(x=0\) 消去，\(x=a\) 时 \(\phi_n(a)=\phi_m(a)=0\)）。故正交性成立。归一化常数：由递推关系 \(\frac{d}{dx}[x J_\nu'(\lambda x)]^2\) 在 \([0,a]\) 上积分并利用 \(J_\nu(j_{\nu,n})=0\) 和 \(J_\nu'(0)\) 的渐近行为，得 \(\int_0^a J_\nu^2(j_{\nu,n} r/a) r dr = \frac{a^2}{2}[J_\nu'(j_{\nu,n})]^2\)。\(\blacksquare\)

\subsection{正交多项式}\label{ux6b63ux4ea4ux591aux9879ux5f0f}

\textbf{定义}（正交多项式族）：在区间 \(I\) 上对于权函数 \(w(x) > 0\) 正交的多项式族 \(\{P_n(x)\}\)： \[\int_I P_m(x) P_n(x) w(x) dx = h_n \delta_{mn}\] 所有经典正交多项式可通过一个统一的生成函数的超几何表示关联（Askey 方案, Askey-Wilson多项式）。

\textbf{经典正交多项式}： - \textbf{Legendre多项式} \(P_n(x)\)：\(I = [-1, 1]\)，\(w(x) = 1\)，生成函数为 \((1 - 2xt + t^2)^{-1/2}\)。Rodrigues公式：\(P_n(x) = \frac{1}{2^n n!} \frac{d^n}{dx^n}[(x^2 - 1)^n]\) - \textbf{Hermite多项式} \(H_n(x)\)：\(I = (-\infty, \infty)\)，\(w(x) = e^{-x^2}\)。Rodrigues公式：\(H_n(x) = (-1)^n e^{x^2} \frac{d^n}{dx^n}[e^{-x^2}]\)。是量子谐振子方程的 Hermite 多项式解 - \textbf{Laguerre多项式} \(L_n^{(\alpha)}(x)\)：\(I = [0, \infty)\)，\(w(x) = x^\alpha e^{-x}\)。氢原子径向方程的典型解 - \textbf{Chebyshev多项式}（已在逼近论中详细论述）：\(T_n(\cos \theta) = \cos(n\theta)\)

\textbf{命题}（递推关系与零点性质）：所有正交多项式满足三项递推关系 \(x P_n(x) = A_n P_{n+1}(x) + B_n P_n(x) + C_n P_{n-1}(x)\)。\(P_n\) 的 \(n\) 个零点全为实数且位于 \(I\) 内（且在 \(P_{n-1}\) 的零点之间交错）。相邻零点之间的最大间距由 Chebyshev 多项式极小化------这是 Gauss 求积公式节点的理论基础。

\textbf{定理}（广义超几何函数的统一框架）：许多特殊函数正交多项式族均为广义超几何函数 \(pF_q\) 的特例。Askey方案（Askey scheme）将所有经典正交多项式按离散-连续和极限关系组织为一个层级结构，Askey-Wilson多项式位于顶层------这是特殊函数统一理论的基石。

\textbf{证明}：经典正交多项式由 Rodrigues 公式统一构造：\(P_n(x)=\frac{1}{K_n w(x)}\frac{d^n}{dx^n}[w(x)Q(x)^n]\)，其中 \(w(x)\) 为权函数，\(Q(x)\) 为至多二次多项式。对于 Jacobi 多项式 \(P_n^{(\alpha,\beta)}(x)\)，其超几何表示为 \[P_n^{(\alpha,\beta)}(x)=\binom{n+\alpha}{n}{}_2F_1\!\left(-n,n+\alpha+\beta+1;\alpha+1;\frac{1-x}{2}\right).\] Askey 方案通过极限关系组织所有经典正交多项式：\(q\)-Askey-Wilson \(\xrightarrow{q\to 1}\) Wilson → Wilson \(\xrightarrow{\text{极限}}\) 连续 Hahn → \ldots{} → Jacobi → Laguerre / Hermite。Chebyshev 多项式是 Jacobi 多项式的参数极限特例。该层级结构表明所有经典正交多项式构成了超几何函数和基本超几何函数的特殊参数取值集合，Askey-Wilson 多项式处于最顶层------其正交测度最为一般（广义 Askey-Wilson 测度），下行各级通过参数约束退化得到。\(\blacksquare\)

\hrulefill

\emph{本卷基于 V18（微分方程），建立了离散和连续动力系统、混沌与分形、遍历论和分支理论的系统理论。共 5 章。}

\emph{本卷建立了微分方程从经典到现代的理论框架：ODE 的存在唯一性理论（Picard-Lindelöf 定理）和稳定性理论（Lyapunov 方法）是动力系统研究的基础；三大经典 PDE（Laplace 方程、热方程、波动方程）的解法展示了 PDE 的丰富结构和物理背景；Sobolev 空间和弱解理论为现代 PDE 研究提供了强大的泛函分析框架（Lax-Milgram 定理、椭圆正则性理论）。补充内容涵盖微分方程的定性理论（Lotka-Volterra 系统、传染病模型的阈值动力学、Turing 不稳定性）和动力系统理论（离散/连续动力系统、混沌与分形、遍历论、分支理论），ODE 作为动力系统的连续版本形成了完整的学科体系。共 10 章。}

\emph{卷十七：微分方程终。}

\emph{卷十七：微分方程终。}

\emph{卷十七：微分方程终。}

\emph{卷二十一：动力系统与特殊函数终。}

\newpage

\chapter{09-微分方程与动力系统/02-动力系统/07-复动力系统.md}\label{ux5faeux5206ux65b9ux7a0bux4e0eux52a8ux529bux7cfbux7edf02-ux52a8ux529bux7cfbux7edf07-ux590dux52a8ux529bux7cfbux7edf.md}

\chapter{复动力系统}\label{ux590dux52a8ux529bux7cfbux7edf}

复动力系统研究复平面 \(\widehat{\mathbb{C}}\) 上解析映射的迭代动力学。该学科融合了复分析、拓扑学与动力系统的思想，其起源可追溯至Gaston Julia和Pierre Fatou在1918年前后的开创性工作。1980年代，Benoit Mandelbrot的计算机制图使Mandelbrot集的惊人复杂性公之于众，而Dennis Sullivan的无游荡域定理（1985）则从理论上奠定了复动力系统的坚实根基。如今，该领域与Kleinian群理论相互渗透，形成了双曲几何与动力系统的统一图景。本章系统阐述从复多项式迭代到Sullivan刚性定理的核心理论。

\hrulefill

\section{第327章：复多项式的迭代理论}\label{ux7b2c327ux7ae0ux590dux591aux9879ux5f0fux7684ux8fedux4ee3ux7406ux8bba}

给定复多项式 \(f: \widehat{\mathbb{C}} \to \widehat{\mathbb{C}}\)，\(\deg f = d \geq 2\)，研究迭代序列 \(\{f^n(z)\}_{n=0}^\infty\) 的动力学行为。Riemann球面 \(\widehat{\mathbb{C}}\) 在迭代下分裂为两个互补的完全不变集：Fatou集（正规族区域）和Julia集（混沌区域）。这一基本二分法是复动力系统的理论起点。

\subsection{x.1 Julia集与Fatou集的基本理论}\label{x.1-juliaux96c6ux4e0efatouux96c6ux7684ux57faux672cux7406ux8bba}

\textbf{定义}：设 \(f: \widehat{\mathbb{C}} \to \widehat{\mathbb{C}}\) 为次数 \(d \geq 2\) 的有理函数。Fatou集 \(F(f)\) 定义为点列 \(\{f^n\}\) 在其某个邻域内构成正规族（即每个子列有局部一致收敛或发散到 \(\infty\) 的子列）的点的集合。Julia集 \(J(f) = \widehat{\mathbb{C}} \setminus F(f)\)。

\textbf{定理 x.1}（Julia集的基本性质）：设 \(f\) 为次数 \(d \geq 2\) 的有理函数，则：

（1）\(J(f) \neq \emptyset\) 且 \(f(J(f)) = J(f) = f^{-1}(J(f))\)（完全不变性）； （2）\(F(f)\) 与 \(J(f)\) 均完全不变； （3）\(J(f)\) 是排斥周期点的闭包； （4）若 \(\deg f \geq 2\)，则 \(J(f)\) 是无处稠密的完全集（perfect set）。

\textbf{证明}：（1）取 \(z_0\) 为一个排斥周期点（其存在性由Montel定理保证，见后续）。由排斥周期点的局部动力学，任意邻域内迭代发散，故 \(z_0 \in J(f)\)，从而 \(J(f) \neq \emptyset\)。完全不变性由定义直接验证：\(f^{-1}(F(f)) = F(f)\)，从而 \(f^{-1}(J(f)) = J(f)\)。

（3）通过Montel定理：若某开集 \(U\) 与 \(J(f)\) 相交，则迭代族 \(\{f^n\}\) 在 \(U\) 中的任何子族都不能是正规的。这意味着 \(U\) 包含某个排斥周期点（或抛物周期点），由此可得 \(J(f)\) 等于排斥周期点集的闭包。具体来说，对任意排斥周期点 \(p\)，存在邻域使得迭代在此处指数发散，故 \(\{f^n\}\) 在 \(p\) 的任意邻域内不均正规，因此 \(p \in J(f)\)。反方向，若 \(z \in J(f)\)，则对 \(z\) 的任意邻域 \(U\)，\(\{f^n|_U\}\) 非正规，由Montel定理，\(\bigcup_{n\geq 0} f^n(U)\) 排除至多两点，从而覆盖某些周期点，最终可得排斥周期点在 \(z\) 的任意邻域中存在。\(\blacksquare\)

\subsection{x.2 Montel定理与Fatou集的正规族刻画}\label{x.2-montelux5b9aux7406ux4e0efatouux96c6ux7684ux6b63ux89c4ux65cfux523bux753b}

\textbf{定理 x.2}（Montel定理）：设 \(\mathcal{F}\) 为开集 \(U \subset \widehat{\mathbb{C}}\) 上的一族全纯函数。若 \(\mathcal{F}\) 局部一致有界（即 \(\forall z \in U\)，\(\exists\) 邻域 \(V \ni z\) 和常数 \(M\) 使 \(\sup_{f \in \mathcal{F}, w \in V} |f(w)| \leq M\)），则 \(\mathcal{F}\) 是正规族。

等价地：若 \(\mathcal{F}\) 在 \(U\) 中遗漏三个固定值（即对所有 \(f \in \mathcal{F}\)，\(f(U)\) 不含某三个不同点），则 \(\mathcal{F}\) 正规。

\textbf{证明}：设 \(\mathcal{F}\) 遗漏值 \(a, b, c \in \widehat{\mathbb{C}}\)（可经Möbius变换设为 \(0, 1, \infty\)）。对任意 \(z_0 \in U\)，取 \(r > 0\) 使 \(\overline{D(z_0, r)} \subset U\)。对任意 \(f \in \mathcal{F}\)，令 \(g = T \circ f\)，其中 \(T\) 为将 \(a, b, c\) 映为 \(0, 1, \infty\) 的Möbius变换。由Schottky定理，存在仅依赖于 \(0, 1, \infty\) 和 \(r\) 的常数 \(C\)，使得在 \(D(z_0, r/2)\) 上 \(|g(z)| \leq C\)。推出 \(\mathcal{F}\) 在 \(D(z_0, r/2)\) 上一致有界。由Cauchy积分公式，\(\mathcal{F}\) 在该邻域中等度连续（equicontinuous），从而由Arzela-Ascoli定理知 \(\mathcal{F}\) 正规。\(\blacksquare\)

\textbf{推论 x.3}（Fatou集的等于正规族区域）：\(F(f) = \{z \in \widehat{\mathbb{C}} \mid \{f^n\} \text{ 在 } z \text{ 的某邻域内正规}\}\)。此即Fatou集定义为 \(\{f^n\}\) 的正规点集。

\textbf{定理 x.4}（Fatou集的完全不变性与分类）：\(F(f)\) 的每个连通分支（称为Fatou分支）在迭代下映为另一个Fatou分支。由Sullivan无游荡域定理（见第x+2章），每个Fatou分支最终都是周期的。周期Fatou分支只可能是以下五种类型之一：吸引域（attracting basin）、超吸引域（super-attracting basin）、抛物域（parabolic basin）、Siegel盘（Siegel disk）或Herman环（Herman ring）。

\subsection{x.3 Böttcher定理与势函数}\label{x.3-buxf6ttcherux5b9aux7406ux4e0eux52bfux51fdux6570}

\textbf{定理 x.5}（Böttcher定理）：设 \(f(z) = z^d + a_{d-1}z^{d-1} + \cdots + a_0\) 为次数 \(d \geq 2\) 的多项式，\(\infty\) 为其超吸引不动点。则存在 \(\infty\) 的邻域 \(U\) 和共形映射 \(\varphi: U \to \{w \in \widehat{\mathbb{C}} \mid |w| > R\}\)（某 \(R > 0\)），使得 \(\varphi(\infty) = \infty\)，\(\varphi'(\infty) = 1\)，且

\[\varphi(f(z)) = (\varphi(z))^d, \quad \forall z \in U.\]

此即Böttcher坐标。此外，存在势函数 \(\Phi(z) = \lim_{n \to \infty} \frac{1}{d^n} \log^+ |f^n(z)|\)（Green函数），它在Julia集外调和，在 \(J(f)\) 上为零。

\textbf{证明}：在 \(\infty\) 附近，令 \(g(z) = 1/f(1/z)\)。则 \(\infty\) 变为 \(0\)，\(f\) 在 \(g\) 坐标下成为 \(g(w) = w^d/(1 + a_{d-1}w + \cdots + a_0 w^d) = w^d(1 + O(w))\)。设 \(\psi(w) = w + \sum_{n\geq 2} b_n w^n\) 为线性化坐标，要求 \(\psi(g(w)) = (\psi(w))^d\)。

形式幂级数代入：\(\psi(g(w)) = g(w) + b_2 g(w)^2 + \cdots = w^d(1+O(w)) + b_2 w^{2d}(1+\cdots) + \cdots\)，而 \((\psi(w))^d = (w + b_2 w^2 + \cdots)^d = w^d + d b_2 w^{d+1} + \cdots\)。比较系数可逐个确定 \(b_n\)。关键是要验证该形式幂级数有非零收敛半径。这可通过构造法证明：定义 \(h_n(z) = (f^n(z))^{1/d^n}\)（选取适当的分支使在无穷远处渐近于 \(z\)）。则 \(h_n\) 在某固定邻域内一致收敛到所需的共形映射 \(\varphi\)。极限的共形性由Hurwitz定理保证。最后，\(\Phi(z) = \log|\varphi(z)|\)（当 \(z\) 在Böttcher坐标定义的区域内），可调和延拓至整个吸引盆 \(A(\infty)\)，定义为 \(\Phi(z) = \lim_{n \to \infty} d^{-n}\log|f^n(z)|\)。\(\Phi\) 满足 \(\Phi(f(z)) = d\Phi(z)\)，在 \(A(\infty)\) 上调和、正，在 \(J(f)\) 上为零（因为若 \(|f^n(z)|\) 不趋于 \(\infty\)，则 \(\log|f^n(z)|\) 有界，\(d^{-n}\) 因子的极限为零）。\(\blacksquare\)

\subsection{x.4 Julia集的完全不变性与自相似性}\label{x.4-juliaux96c6ux7684ux5b8cux5168ux4e0dux53d8ux6027ux4e0eux81eaux76f8ux4f3cux6027}

\textbf{定理 x.6}（Julia集的自相似性）：对任意 \(z \in J(f)\) 和 \(z\) 的任意邻域 \(U\)，存在 \(N > 0\) 使得 \(f^N(U)\) 覆盖 \(J(f)\)（除去至多有限个例外点）。特别地，\(J(f)\) 在任意尺度下具有相似的精细结构。

\textbf{证明}：设 \(W \subset \widehat{\mathbb{C}}\) 为包含 \(J(f)\) 的开集。考虑补集 \(\widehat{\mathbb{C}} \setminus \bigcup_{n \geq 0} f^n(U)\)。此集合在 \(f^{-1}\) 下不变。由Montel定理，若该补集非空，则它只能包含至多两个点（因为 \(\{f^n\}\) 在 \(\widehat{\mathbb{C}} \setminus \bigcup_{n \geq 0} f^n(U)\) 上遗漏了 \(f^{-1}(U^c)\) 中的点，从而正规）。这迫使 \(\bigcup_{n \geq 0} f^n(U)\) 覆盖除至多两点外的整个 \(\widehat{\mathbb{C}}\)。

取 \(U\) 为 \(z\) 的小邻域，则对充分大的 \(n\)，\(f^n(U) \supset J(f)\)（除可能至多两例外点）。这意味着 \(J(f)\) 在 \(f^n\) 的局部逆像下呈现出与整体 \(J(f)\) 相同的定性结构------此即自相似性的精确表述。\(\blacksquare\)

\hrulefill

\section{第328章：Mandelbrot集}\label{ux7b2c328ux7ae0mandelbrotux96c6}

Mandelbrot集 \(\mathcal{M}\) 是复动力系统中最著名的分形对象。其定义为二次多项式族 \(f_c(z) = z^2 + c\) 中使临界轨道有界的参数 \(c\) 的集合。Mandelbrot集不仅是美的视觉对象，更是理解二次多项式动力学的参数空间''字典''：每个 \(c \in \mathcal{M}\) 对应一个连通的Julia集，每个 \(c \notin \mathcal{M}\) 对应一个Cantor型Julia集。

\subsection{x+1.1 Mandelbrot集的基本性质}\label{x1.1-mandelbrotux96c6ux7684ux57faux672cux6027ux8d28}

\textbf{定义}：Mandelbrot集 \(\mathcal{M} = \{c \in \mathbb{C} \mid f_c^n(0) \not\to \infty \text{ 当 } n \to \infty\}\)，其中 \(f_c(z) = z^2 + c\)。等价地，\(\mathcal{M} = \{c \mid \sup_n |f_c^n(0)| \leq 2\}\)。

\textbf{定理 x+1.1}（Mandelbrot集的基本性质）：（1）\(\mathcal{M}\) 是紧集，包含于闭圆盘 \(\overline{D}(0,2)\)；（2）\(\mathcal{M}\) 关于实轴对称；（3）\(\mathcal{M}\) 与实轴的交为 \([-2, 1/4] \subset \mathbb{R}\)。

\textbf{证明}：（1）若 \(|c| > 2\)，则 \(|f_c(0)| = |c| > 2\)。归纳假设 \(|f_c^n(0)| > |c| > 2\)，则 \(|f_c^{n+1}(0)| = |f_c^n(0)^2 + c| \geq |f_c^n(0)|^2 - |c| > 2|c| - |c| = |c| > 2\)。因此轨道无界，\(c \notin \mathcal{M}\)。故 \(\mathcal{M} \subset \overline{D}(0,2)\)。

紧性：定义 \(G(c) = \lim_{n\to\infty} 2^{-n}\log^+|f_c^n(0)|\)，则 \(\mathcal{M} = G^{-1}(0)\)。由Hartogs定理，\(G\) 连续且调和于 \(\mathbb{C} \setminus \mathcal{M}\)，故 \(\mathcal{M}\) 为紧集。（3）与实轴的交可直接从实二次映射 \(x^2 + c\) 的动力学分析得到。\(\blacksquare\)

\subsection{x+1.2 Douady-Hubbard连通性定理}\label{x1.2-douady-hubbardux8fdeux901aux6027ux5b9aux7406}

\textbf{定理 x+1.2}（Douady-Hubbard定理）：Mandelbrot集 \(\mathcal{M}\) 是连通的。

\textbf{证明框架}：考虑 \(\widehat{\mathbb{C}} \setminus \mathcal{M}\)。定义共形映射 \(\Phi: \widehat{\mathbb{C}} \setminus \mathcal{M} \to \widehat{\mathbb{C}} \setminus \overline{\mathbb{D}}\)（其中 \(\mathbb{D}\) 为单位圆盘），\(\Phi(c) = \varphi_c(c)\)，其中 \(\varphi_c\) 是 \(f_c\) 在 \(\infty\) 处的Böttcher坐标。对 \(c \notin \mathcal{M}\)，\(\varphi_c(c)\) 是有定义的，且 \(\Phi\) 为共形同构。

若 \(\mathcal{M}\) 不连通，则 \(\widehat{\mathbb{C}} \setminus \mathcal{M}\) 存在有界分支。设 \(U\) 为 \(\widehat{\mathbb{C}} \setminus \mathcal{M}\) 的一个有界分支，则 \(\Phi(U)\) 是 \(\widehat{\mathbb{C}} \setminus \overline{\mathbb{D}}\) 的有界开集，但 \(\widehat{\mathbb{C}} \setminus \overline{\mathbb{D}}\) 仅有一个无界分支（即全空间自身），矛盾。因此 \(\mathcal{M}\) 连通。

构造 \(\Phi\) 的详细步骤：首先对每个 \(c\)，设 \(R(c) = \sup_n |f_c^n(0)|\)。\(R\) 在 \(\mathbb{C} \setminus \mathcal{M}\) 上连续。利用Bötcher坐标 \(\varphi_c\)，定义 \(\Phi(c) = \varphi_c(c)\)。此映射在 \(\mathbb{C} \setminus \mathcal{M}\) 上全纯。验证 \(\Phi\) 在 \(\infty\) 附近的渐近行为：当 \(c \to \infty\)，\(\Phi(c) \sim c\)。由Riemann映射定理，\(\Phi\) 提供了 \(\mathbb{C} \setminus \mathcal{M}\) 到 \(\widehat{\mathbb{C}} \setminus \overline{\mathbb{D}}\) 的共形等价，由此推出连通性。\(\blacksquare\)

\subsection{x+1.3 Mandelbrot集的分支与周期倍化}\label{x1.3-mandelbrotux96c6ux7684ux5206ux652fux4e0eux5468ux671fux500dux5316}

\textbf{定理 x+1.3}（Mandelbrot集分支的周期参数化）：对每个有理数 \(p/q \in \mathbb{Q}/\mathbb{Z}\)（不可约），存在 \(\mathcal{M}\) 边界上的一个附着点（即Misiurewicz点）和以其为根的分支（limb）。\(\mathcal{M}\) 的主心形（cardioid）对应于具有吸引不动点的参数 \(c\)。在心形边界的 \(p/q\) 有理点处附着着周期为 \(q\) 的''球''（bulb）。

\textbf{证明}：在心形内部，\(f_c\) 有吸引不动点，满足乘子 \(\mu(c) = f_c'(z_c) = 2z_c\)。由不动点方程 \(z_c^2 + c = z_c\) 得 \(z_c = \frac{1}{2}(1 \pm \sqrt{1-4c})\)。心形边界由 \(|\mu(c)| = 1\) 刻画，即 \(c = \frac{e^{i\theta}}{2} - \frac{e^{2i\theta}}{4}\)（\(\theta \in [0, 2\pi]\)）。当 \(\theta = 2\pi p/q\) 时，乘子 \(\mu = e^{2\pi i p/q}\)，不动点为抛物的，此时发生鞍-结点分岔，产生周期-\(q\) 的附属球。在附属球内部，存在周期为 \(q\) 的吸引周期轨道，其乘子连续变化到球中心（乘子为零，即超吸引情形）。\(\blacksquare\)

\subsection{x+1.4 外射线与Yoccoz不等式}\label{x1.4-ux5916ux5c04ux7ebfux4e0eyoccozux4e0dux7b49ux5f0f}

\textbf{定义}：对 \(c \notin \mathcal{M}\)，外射线（external ray）\(R_c(\theta)\) 定义为 \(\Phi^{-1}(\{re^{2\pi i\theta} \mid r > 1\})\)。若该射线在接近 \(\mathcal{M}\) 时收敛到唯一一点，则称该射线''着陆''（land），着陆点属于 \(\partial\mathcal{M}\)。

\textbf{定理 x+1.4}（Yoccoz不等式）：设 \(c\) 为Mandelbrot集的Misiurewicz参数（即 \(0\) 的轨道最终碰到排斥周期点），\(P\) 为包含 \(c\) 的''拼图片''（puzzle piece）。则对任意 \(n \geq 0\)，

\[\operatorname{diam}(P \cap J(f_c)) \leq C \cdot 2^{-\alpha n},\]

其中 \(\alpha > 0\) 是与 \(c\) 相关的Yoccoz指数。该不等式是Yoccoz证明相应Julia集局部连通性的关键工具。

\textbf{证明思路}：Yoccoz拼图技术将动力学平面划分为由等势线和外射线构成的拼图片。拼图片层的深度 \(n\) 通过拉回运算定义。关键估计来自模（modulus）的增长：连续环域（annulus）的模满足

\[\operatorname{mod}(A_{n+1}) \geq \operatorname{mod}(A_n) + \delta,\]

其中 \(\delta > 0\) 来自拼图组合学中的''增长引理''。由Grötzsch不等式，模的无界增长迫使拼图片直径指数衰减。具体地，设 \(P_n\) 为深度 \(n\) 的拼图片，则 \(P_{n+1}\) 在 \(P_n\) 内的相对模大于某常数 \(\eta > 0\)。由Teichmüller极值问题，\(\operatorname{diam}(P_n) \leq C\operatorname{mod}(P_n \setminus P_{n+1})^{-1/2} \leq C 2^{-\alpha n}\)。这就是Yoccoz不等式的核心推理链。\(\blacksquare\)

\subsection{x+1.5 MLC猜想}\label{x1.5-mlcux731cux60f3}

\textbf{猜想}（MLC猜想，Mandelbrot集局部连通性猜想）：Mandelbrot集 \(\mathcal{M}\) 是局部连通的。

若MLC成立，则 \(\mathcal{M}\) 的所有外射线均可着陆，且Carathéodory定理保证 \(\Phi^{-1}\) 可连续延拓至边界，从而给出 \(\mathcal{M}\) 边界的完整组合描述。截至2026年，在参数平面上与特征双曲分支的测地片相关的进展已由部分结果的支持，但该猜想整体上仍为开放问题。

\hrulefill

\section{第329章：有理函数迭代与Kleinian群}\label{ux7b2c329ux7ae0ux6709ux7406ux51fdux6570ux8fedux4ee3ux4e0ekleinianux7fa4}

Sullivan于1985年证明了复动力系统中里程碑式的''无游荡域定理''，从而完成了Fatou集的分类。与此同时，复动力系统与Kleinian群理论之间的惊人平行性被揭示：Julia集对应Kleinian群的极限集，Sullivan刚性定理对应Mostow刚性定理。Thurston的轨形理论则为理解有理映射的组合结构提供了拓扑框架。

\subsection{x+2.1 Sullivan无游荡域定理}\label{x2.1-sullivanux65e0ux6e38ux8361ux57dfux5b9aux7406}

\textbf{定理 x+2.1}（Sullivan无游荡域定理）：设 \(f\) 为次数 \(d \geq 2\) 的有理函数。则Fatou集 \(F(f)\) 的每个连通分支最终都是周期的（即不存在游荡的Fatou分支）。换言之，对 \(F(f)\) 的任意分支 \(U\)，存在整数 \(n \geq 0\) 和 \(p \geq 1\) 使得 \(f^{n+p}(U) \subseteq f^n(U)\) 且 \(f^p(f^n(U)) = f^n(U)\)。

\textbf{证明}：假设存在游荡分支 \(U\)，即对任意 \(n \neq m\)，\(f^n(U) \cap f^m(U) = \emptyset\)。考虑二次微分 \(\omega\) 在 \(F(f)\) 上的空间。由于 \(f^{-1}: F(f) \to F(f)\) 是连续满射，Sullivan的关键洞察是：在游荡分支上可以构造一个Beltrami微分 \(\mu\)（即 \((-1,1)\)-形式），使其经过 \(f\) 的拉回后保持有界。

具体构造如下：定义在 \(U\) 上的实验室Beltrami微分 \(\mu_0\)（非零）。通过所有迭代拉回 \((f^{-n})^*\mu_0\)，在 \(\bigcup_{n \in \mathbb{Z}} f^n(U)\) 上定义Beltrami微分 \(\mu\)。在其他Fatou分支上令 \(\mu = 0\)。在Julia集上令 \(\mu = 0\)。由于 \(U\) 是游荡的，不同 \(f^n(U)\) 不相交，故 \(\mu\) 良定义且可测（因可数个可测子集上的定义）。

验证 \(\mu\) 是 \(f\)-不变量：对 \(z \in f(U)\)，\(f^*\mu(z) = \mu(f(z))\overline{f'(z)}/f'(z) = \mu_0(z)\)（在 \(U\) 的一个原像上），从而 \(f^*\mu = \mu\)。由可测Riemann映射定理（Bers-Royden方法），存在拟共形映射 \(\varphi: \widehat{\mathbb{C}} \to \widehat{\mathbb{C}}\)，使得 \(\overline{\partial}\varphi = \mu \partial\varphi\)。则 \(g = \varphi \circ f \circ \varphi^{-1}\) 是有理函数（因为 \(\mu\) 是 \(f\)-不变量，且 \(\varphi\) 的复结构被 \(f\) 保持）。\(g\) 的次数等于 \(\deg f\)。

由于 \(\varphi\) 在 \(U\) 上非共形（若 \(\mu_0 \neq 0\)），而在其他Fatou分支及Julia集上共形，这产生了与有理函数参数空间有限维性的矛盾：\(\mu\) 提供了一族非平凡的 \(f\)-不变量Beltrami微分，其参数维数等于游荡分支数目。然而，有理函数空间的维度由 \(\deg f\) 决定（即 \(2\deg f - 2\) 个有效参数），无界多的游荡分支将导致参数维数无界，矛盾。因此不存在游荡分支。\(\blacksquare\)

\subsection{x+2.2 有理函数的双曲猜想}\label{x2.2-ux6709ux7406ux51fdux6570ux7684ux53ccux66f2ux731cux60f3}

\textbf{定义}：有理函数 \(f\) 称为双曲的（hyperbolic），若 \(f\) 在 \(J(f)\) 上一致扩张，等价地，\(f\) 的每个临界点的轨道收敛于吸引周期轨道。对双曲有理映射，\(J(f)\) 是双曲集。

\textbf{猜想}（双曲猜想）：双曲有理映射在次数为 \(d \geq 2\) 的有理函数空间 \(\operatorname{Rat}_d\) 中是稠密的。

该猜想是复动力系统领域最核心的开放问题之一。对多项式情形（即多项式双曲猜想），由Graczyk和Swiatek（1997）对所有实参数给出了部分结果。

\subsection{x+2.3 Kleinian群与Ahlfors有限性定理}\label{x2.3-kleinianux7fa4ux4e0eahlforsux6709ux9650ux6027ux5b9aux7406}

\textbf{定义}：Kleinian群 \(\Gamma\) 是 \(\operatorname{PSL}(2,\mathbb{C})\) 中作用在 \(\widehat{\mathbb{C}}\) 上的离散子群。\(\Gamma\) 的不连续集（domain of discontinuity）\(\Omega(\Gamma)\) 是 \(\Gamma\) 在其中正常作用的最大开集。极限集 \(\Lambda(\Gamma) = \widehat{\mathbb{C}} \setminus \Omega(\Gamma)\)。

\textbf{定理 x+2.2}（Ahlfors有限性定理）：设 \(\Gamma\) 为有限生成Kleinian群。则商曲面 \(\Omega(\Gamma)/\Gamma\) 是有限个有限型Riemann曲面（即紧致曲面除去有限个点）的不交并。换言之，\(\Omega(\Gamma)/\Gamma\) 只有有限个分支，每个分支是有限型曲面。

\textbf{证明}：核心论证利用Bers嵌入理论与拟共形映射的Teichmüller理论。商 \(\Omega/\Gamma\) 在Bers纤维化中对应于Bers边界上的点。由Ahlfors和Bers的理论，Bers空间 \(T(S)\)（Technmüller空间）的有限维性迫使 \(\Omega/\Gamma\) 的分支数有限。具体步骤：（1）将 \(\Gamma\) 在 \(\Omega(\Gamma)\) 上的作用提升到万有Teichmüller空间；（2）利用拟共形映射的参数化理论，\(\Gamma\) 的拟共形变形空间是有限维的（维数等于生成元个数的某种函子）；（3）若 \(\Omega(\Gamma)/\Gamma\) 有无穷多分支，则变形空间将有无穷维，与有限生成性矛盾。\(\blacksquare\)

\subsection{x+2.4 Mostow刚性定理在复动力系统中的应用}\label{x2.4-mostowux521aux6027ux5b9aux7406ux5728ux590dux52a8ux529bux7cfbux7edfux4e2dux7684ux5e94ux7528}

\textbf{定理 x+2.3}（Mostow刚性定理）：设 \(M, N\) 为维数 \(\geq 3\) 的有限体积双曲流形，且 \(\pi_1(M) \cong \pi_1(N)\)。则 \(M\) 与 \(N\) 等距。

在复动力系统中的应用：Sullivan将Mostow刚性定理的思想引入有理映射的变形理论中。对双曲有理映射 \(f\)，其Julia集上不存在 \(f\)-不变量线场（invariant line field）。这意味着 \(f\) 在拟共形形变下是刚性的，即若拟共形映射 \(\varphi\) 与 \(f\) 可换取，则 \(\varphi\) 必为共形映射。此即Sullivan刚性定理在双曲情形下的表述。

\textbf{定理 x+2.4}（Sullivan刚性定理------无游荡域定理的推论）：设 \(f\) 为次数 \(d \geq 2\) 的有理函数。则 \(f\) 上不存在支撑在Julia集上的 \(f\)-不变量可测线场，除非 \(f\) 是双曲情形或与某个覆盖相关的Lattès映射。特别地，对非Lattès映射，\(f\) 的Teichmüller形变空间由临界轨道有限个复参数参数化。

\subsection{x+2.5 Thurston双曲轨形定理}\label{x2.5-thurstonux53ccux66f2ux8f68ux5f62ux5b9aux7406}

\textbf{定义}：Thurston轨形（orbifold）是一个拓扑空间，局部等同于 \(\mathbb{R}^2/\Gamma\)，其中 \(\Gamma\) 为有限子群。对于复动力系统，与有理函数 \(f\) 关联一个记号（signature）轨形 \(\mathcal{O}_f\)，其底空间为 \(\widehat{\mathbb{C}}\)，轨形奇点位于 \(f\) 的后临界集上。

\textbf{定理 x+2.5}（Thurston双曲轨形定理）：与有理函数 \(f\) 关联的轨形 \(\mathcal{O}_f\) 若为双曲的（即其万有覆盖为 \(\widehat{\mathbb{C}}\)、\(\mathbb{C}\) 或 \(\mathbb{D}\)，不含Euclidean或spherical型），则 \(f\) 由其组合数据（临界轨道的组合类型）唯一确定，至多相差共形同构。进一步，\(f\) 可通过Thurston算法从组合数据构造：逐次调整分支点处的局部自由度直至全局一致性收敛。

\textbf{证明框架}：Thurston的原始证明分三步。第一步，从临界轨道的组合类型构建一个Thurston映射（即保持给定轨道组合的多项式型或有理型分支覆盖）。第二步，将Teichmüller空间（带记号的Riemann曲面）对应的点经Thurston拉回算子 \(\sigma_f\) 作用。\(\sigma_f\) 在Teichmüller空间上是（严格）收缩的。第三步，由Banach不动点定理，\(\sigma_f\) 在Teichmüller空间中有唯一不动点。该不动点对应于使 \(f\) 成为实际有理函数（或阻碍，即Thurston阻碍）的共形结构。当不存在Thurston阻碍时，不动点给出所求的有理函数------此即Thurston组合构造定理。双曲条件确保Thurston阻碍不可能出现，从而 \(\sigma_f\) 收敛到唯一不动点。\(\blacksquare\)

\subsection{x+2.6 Julia集的Hausdorff维数}\label{x2.6-juliaux96c6ux7684hausdorffux7ef4ux6570}

Julia集的几何复杂性能通过Hausdorff维数量化。对于双曲有理函数，Julia集的Hausdorff维数通常不是整数，反映了其分形本质。

\textbf{定理 x+2.6}（Bowen定理------双曲Julia集的Hausdorff维数）：设 \(f\) 为双曲有理函数。则其Julia集 \(J(f)\) 的Hausdorff维数 \(\operatorname{HD}(J(f))\) 是方程 \(P(-t\log|f'|) = 0\) 的唯一解，其中 \(P\) 为 \(f|_{J(f)}\) 的拓扑压（topological pressure）。该维数严格介于1和2之间（除非Julia集为圆周）。

\textbf{证明}：Bowen的证明将热力学形式体系应用于保形动力系统。考虑Hölder连续位势 \(\phi_t = -t\log|f'|\)，其拓扑压 \(P(\phi_t)\) 是 \(t\) 的严格递减函数，在 \(t=0\) 处为正，在 \(t\) 充分大时为负。令 \(t_0 = \operatorname{HD}(J(f))\) 满足 \(P(\phi_{t_0}) = 0\)。由Manning、McCluskey的公式，对于有限分支覆盖的扩张映射，平衡态（Gibbs测度）的Hausdorff维数等于使压为零的唯一参数。该Gibbs测度实现为 \(J(f)\) 上 \(t_0\) 维Hausdorff测度的等价类。

具体计算依赖迭代函数系（IFS）结构。在双曲情形，\(f^{-1}\) 的局部逆分支在 \(J(f)\) 上构成共形IFS，满足开集条件。由共形IFS的Hausdorff维数公式（Mauldin-Urbanski），\(\operatorname{HD}(J(f))\) 由上述压方程确定。\(\blacksquare\)

\textbf{推论}（二次多项式的Hausdorff维数连续性）：对于二次多项式族 \(f_c(z) = z^2 + c\)，映射 \(c \mapsto \operatorname{HD}(J(f_c))\) 在Mandelbrot集的边界上连续（Shishikura的连续性定理）。

\subsection{x+2.7 抛物爆炸与Lavaurs定理}\label{x2.7-ux629bux7269ux7206ux70b8ux4e0elavaursux5b9aux7406}

当有理函数具有抛物周期点时，其动力学在抛物点附近发生剧烈的爆炸现象（parabolic implosion）。这在参数空间中导致Mandelbrot集边界附近的复杂行为。

\textbf{定理 x+2.7}（抛物爆炸与Lavaurs映射）：设 \(f_c\) 有抛物周期点乘子 \(\lambda = e^{2\pi i p/q}\)。考虑参数序列 \(c_n \to c_{\text{parabolic}}\) 使得对应映射具有刚过抛物的点（post-parabolic）。在适当的参数尺度 \(n \to \infty\) 下，迭代的极限（Lavaurs映射）是有理函数，其动力学由抛物轨道的渐近映射刻画。

\textbf{证明}（概要）：在抛物不动点（不妨设 \(z=0\)，\(f(z) = z + a z^{q+1} + \cdots\)）附近，经过Fatou坐标变换 \(\varphi\)，\(f\) 在吸引或排斥瓣（petal）上共轭于平移 \(w \mapsto w+1\)。当参数微小扰动使乘子偏离1，Fatou坐标产生相位差（sectorial limits）。

Lavaurs映射通过匹配吸引瓣和排斥瓣的Fatou坐标构造。具体地，令 \(\Phi_{\text{att}}\) 和 \(\Phi_{\text{rep}}\) 分别为吸引和排斥Fatou坐标。Lavaurs映射 \(L\) 定义为 \(\Phi_{\text{rep}}^{-1} \circ (\Phi_{\text{att}} + \text{shift})\)。此映射的动力学捕获了参数穿过抛物点时的极限行为。更多技术细节依赖于Douady-Lavaurs的ECAL圆柱工具。\(\blacksquare\)

\subsection{x+2.8 Shishikura定理与Hausdorff维数的参数依赖性}\label{x2.8-shishikuraux5b9aux7406ux4e0ehausdorffux7ef4ux6570ux7684ux53c2ux6570ux4f9dux8d56ux6027}

\textbf{定理 x+2.8}（Shishikura定理------二次多项式族Julia集的Hausdorff维数连续性）：映射 \(\mathcal{M} \to [0,2]\)，\(c \mapsto \operatorname{HD}(J(f_c))\) 在 \(\partial\mathcal{M}\) 上连续。但此函数在 \(\partial\mathcal{M}\) 上并不解析，且取值严格介于1和2之间在参数空间的一个稠密子集上。

\textbf{证明思想}：Shishikura的证明结合了多种技术。对于双曲参数，\(\operatorname{HD}(J(f_c))\) 对 \(c\) 是实解析的（由Bowen公式及热力学量的解析依赖性）。在非双曲参数处，需用''拼图技术''对Hausdorff维数做上下界估计。利用Yoccoz拼图的复界（complex bounds），证明Hausdorff维数在非双曲参数处也是连续的。核心工具是对共形IFS的Hausdorff维数做扰动分析，结合Kahn-Lyubich的复先验界和复测地线。

\subsection{x+2.9 复动力系统与Kleinian群的综合展望}\label{x2.9-ux590dux52a8ux529bux7cfbux7edfux4e0ekleinianux7fa4ux7684ux7efcux5408ux5c55ux671b}

复动力系统与Kleinian群理论之间的平行性是当代几何分析的重要篇章。Sullivan字典（Sullivan dictionary）建立了以下对应：

\begin{itemize}
\tightlist
\item
  Julia集 \(\leftrightarrow\) Kleinian群的极限集
\item
  Fatou分支 \(\leftrightarrow\) 不连续集的分支
\item
  有理函数的度 \(\leftrightarrow\) Kleinian群的生成元个数
\item
  临界点 \(\leftrightarrow\) Kleinian群的尖点（cusp）
\item
  双曲猜想 \(\leftrightarrow\) Ahlfors测度猜想
\end{itemize}

这一字典不仅是隐喻层面的类比，而是由深刻的统一定理支撑的。McMullen、Minsky、Brock等学者的工作表明，复动力系统与Kleinian群均可纳入''全纯运动''（holomorphic motions）与''拟共形刚性''的统一框架。特别地，近期关于Mandelbrot集局部连通性（MLC）和终结双曲叠层（ending lamination）猜想的进展，进一步巩固了这两个领域的内在统一性。

\subsection{x+2.10 线性化定理与Siegel盘}\label{x2.10-ux7ebfux6027ux5316ux5b9aux7406ux4e0esiegelux76d8}

复动力系统中，周期点附近的局部动力学由线性化定理描述。这直接关联到复数域上的小除数问题（small divisor problem）。

\textbf{定理 x+2.9}（Koenigs线性化定理与Siegel定理）：设 \(f(z) = \lambda z + a_2 z^2 + \cdots\) 在原点附近全纯。（1）若 \(|\lambda| \neq 0, 1\)，则存在局部共形映射 \(\varphi\) 使得 \(\varphi \circ f \circ \varphi^{-1}(w) = \lambda w\)（Poincaré-Koenigs定理）；（2）若 \(\lambda = e^{2\pi i\alpha}\)（\(\alpha \in \mathbb{R} \setminus \mathbb{Q}\)）且 \(\alpha\) 满足Diophantine条件：存在 \(C > 0\)，\(\tau \geq 0\) 使得 \(|\alpha - p/q| \geq C q^{-(2+\tau)}\) 对一切有理数 \(p/q\) 成立，则仍存在局部共形线性化（Siegel定理）。

\textbf{证明}（Siegel定理概要）：假设线性化有形式幂级数解 \(\varphi(z) = z + \sum_{n\geq 2} b_n z^n\)。由共轭关系 \(\varphi(\lambda z) = f(\varphi(z))\)，逐次比较系数得

\[b_n(\lambda^n - \lambda) = P_n(a_2,\ldots,a_n; b_2,\ldots,b_{n-1}),\]

其中 \(P_n\) 为多项式。由于 \(\lambda^n - \lambda = e^{2\pi i\alpha n} - e^{2\pi i\alpha}\)，当 \(\alpha\) 为Diophantine数时，\(|\lambda^n - \lambda|^{-1}\) 的增长至多为多项式级别（\(O(n^{\tau+1})\)）。这保证了形式级数的收敛性（通过控制 \(|b_n|\) 的增长）。Siegel的原始证明使用了优级数法（majorant method），Bruno后来将此条件推广到更广的Bruno条件。\(\blacksquare\)

\subsection{x+2.11 Lyubich定理与测度刚性}\label{x2.11-lyubichux5b9aux7406ux4e0eux6d4bux5ea6ux521aux6027}

\textbf{定理 x+2.10}（Lyubich测度刚性定理）：设 \(f\) 为次数 \(d \geq 2\) 的有理函数。则 \(f\) 在Julia集上最多允许一个不变概率测度具有正Lyapunov指数且绝对连续于共形测度（conformal measure）。若该测度存在，则系统的度量熵等于Lyapunov指数。

\textbf{证明思路}：基于共形测度理论。对位势 \(\phi = \log|f'|\)，Sullivan的共形测度是满足 \(\nu(f(E)) = \int_E e^{P-\phi} d\nu\)（\(P\) 为压）的测度。若存在 \(f\)-不变概率测度 \(\mu\) 相对于 \(\nu\) 绝对连续（即存在Radon-Nikodym导数），由Ledrappier-Young公式和Rokhlin分解可证 \(\mu\) 是唯一的。Lyubich的核心贡献在于证明了 \(\mu\) 的Jacobian在Julia集上几乎处处等于 \(|f'|\)（在适当的尺度下），从而揭示了扩张性与测度刚性之间的等价关系。单值性论证依赖复动力系统的双曲度量和负曲率特性。\(\blacksquare\)

\subsection{x+2.12 二次多项式的组合刚性------Douady-Hubbard-Lavaurs-Yoccoz拼图}\label{x2.12-ux4e8cux6b21ux591aux9879ux5f0fux7684ux7ec4ux5408ux521aux6027douady-hubbard-lavaurs-yoccozux62fcux56fe}

Yoccoz拼图技术是研究二次多项式Julia集局部结构的最有力工具。在各拼图片的Grötzsch环域模的精细估计下，可以证明局部连通性。

\textbf{定理 x+2.11}（Yoccoz定理------二次多项式拼图）：设 \(f_c(z) = z^2 + c\)，\(c\) 为非无穷可重整化（non-infinitely renormalizable）的参数。则 \(J(f_c)\) 是局部连通的。进一步，若 \(c \in \mathcal{M}\) 也是非无穷可重整化的，则 \(\mathcal{M}\) 在 \(c\) 处局部连通。

\textbf{证明框架}：Yoccoz拼图构造如下：选取等势线 \(\{z \mid G(z) = 1\}\) 和外角为有理数的外射线。拼图片为等势线与外射线围成的区域。对每个正整数 \(n\)，定义深度 \(n\) 的拼图片族 \(\mathcal{P}_n\)。核心估计为拼图片模的增长引理：存在 \(\delta > 0\) 使得

\[\sum_{n=0}^\infty (\operatorname{mod}(P_n(z) \setminus P_{n+1}(z)))^{-1} < \infty.\]

由Grötzsch不等式，模的无界增长等价于拼图片直径的超指数衰减，从而Julia集在每点处有任意小的连通邻域基------此即局部连通性的定义。非无穷可重整化条件保证拼图片的层级不会无限循环。\(\blacksquare\)

\subsection{x+2.13 McMullen定理与多项式Julia集的可容许多边形}\label{x2.13-mcmullenux5b9aux7406ux4e0eux591aux9879ux5f0fjuliaux96c6ux7684ux53efux5bb9ux8bb8ux591aux8fb9ux5f62}

\textbf{定理 x+2.12}（McMullen定理）：设 \(f\) 为次数 \(d \geq 2\) 的多项式。若其Julia集 \(J(f)\) 具有正Lebesgue测度，则 \(J(f)\) 必为整个Riemann球面的子集（在某种覆盖意义下）。换言之，Julia集为全测度集或零测度集。对于二次多项式，已有的主要结果（由Buff和Chéritat于2012年建立）表明存在具有正测度Julia集的二次多项式，且该类参数在 \(\partial\mathcal{M}\) 的一个正测度子集上出现。

\subsection{x+2.14 Ahlfors-Bers可测Riemann映射定理}\label{x2.14-ahlfors-bersux53efux6d4briemannux6620ux5c04ux5b9aux7406}

可测Riemann映射定理是Sullivan无游荡域定理以及整个拟共形刚性理论的基础工具。

\textbf{定理 x+2.13}（可测Riemann映射定理，Ahlfors-Bers）：设 \(\mu\) 为 \(\widehat{\mathbb{C}}\) 上的Beltrami微分，满足 \(\|\mu\|_{L^\infty} < 1\)。则存在唯一的拟共形同胚 \(\varphi^\mu: \widehat{\mathbb{C}} \to \widehat{\mathbb{C}}\)，固定 \(0\)、\(1\)、\(\infty\)，使得 \(\overline{\partial}\varphi^\mu = \mu \partial\varphi^\mu\) 在分布意义下几乎处处成立。进而，映射 \(\mu \mapsto \varphi^\mu\) 在适当的拓扑下是连续（甚至解析）的。

\textbf{证明}：将Beltrami方程 \(\overline{\partial}w = \mu \partial w\) 转化为等价的奇异积分方程。设 \(w = z + T(\rho)\)，其中 \(T\) 为Cauchy-Green积分算子（即 \(\overline{\partial}\) 的基本解）。代入方程得 \(\rho = \mu(1 + S(\rho))\)，其中 \(S = \partial T\) 为Beurling-Ahlfors变换（奇异的Calderón-Zygmund算子）。由于 \(\|S\|_{L^p} \to 1\) 当 \(p \to 2\)，且 \(\|\mu\|_\infty < 1\)，对于充分接近2的 \(p\)，\(\rho \mapsto \mu(1+S\rho)\) 在 \(L^p\) 中是压缩的。由Banach不动点定理得唯一的 \(\rho \in L^p\)，从而给出解 \(w = z + T(\rho)\)。\(\blacksquare\)

\subsection{x+2.15 全纯运动与复动力系统的形变}\label{x2.15-ux5168ux7eafux8fd0ux52a8ux4e0eux590dux52a8ux529bux7cfbux7edfux7684ux5f62ux53d8}

全纯运动（holomorphic motion）是复动力系统中连接不同参数下Julia集的解析桥梁，由Mané-Sad-Sullivan（1983）引入。

\textbf{定义}：一个全纯运动是一个映射 \(H: \mathbb{D} \times E \to \widehat{\mathbb{C}}\)（\(E \subset \widehat{\mathbb{C}}\)），满足（i）对每个 \(z \in E\)，\(\lambda \mapsto H(\lambda, z)\) 全纯；（ii）对每个 \(\lambda \in \mathbb{D}\)，\(z \mapsto H_\lambda(z)\) 为单射；（iii）\(H_0 = \operatorname{id}\)。

\textbf{定理 x+2.14}（\(\lambda\)-引理，Mané-Sad-Sullivan）：任何全纯运动 \(H: \mathbb{D} \times E \to \widehat{\mathbb{C}}\) 可唯一延拓为全纯运动 \(\overline{H}: \mathbb{D} \times \overline{E} \to \widehat{\mathbb{C}}\)，且每个 \(\overline{H}_\lambda\) 为拟共形同胚，其伸缩商满足 \(\log K(\overline{H}_\lambda) \leq \frac{2|\lambda|}{1-|\lambda|}\)。

\textbf{证明}：关键步骤是对 \(E\) 中点之间的双曲距离的估计。由Schwarz-Pick引理，对任意 \(z,w \in E\)，函数 \(\lambda \mapsto \rho(H_\lambda(z), H_\lambda(w))\)（球面距离）是 \(\lambda\) 的下调和函数。由全纯运动的单射性，可构造 \(E\) 上的一致结构和延拓到闭包 \(\overline{E}\)。延拓的唯一性来自全纯函数的唯一性定理。伸缩商的上界通过Bers-Royden的拟共形映射扩张定理和极值长度的估计得到。该定理是连接全纯动力系统与拟共形映射理论的枢纽，其在Sullivan无游荡域定理和复动力系统刚性理论中反复发挥关键作用。特别地，\(\lambda\)-引理保证了对于参数空间中的全纯族，Julia集随参数连续形变。Douady和Hubbard正是借此工具将Mandelbrot集与全纯运动紧密关联，完成了外部映射的构造。\(\lambda\)-引理的解释力超越了单一动力系统的范畴，它统一了Bers嵌入定理、Teichmüller空间的复结构以及双曲三维流形的拟共形变形理论。McMullen在其Fields奖获奖工作中进一步揭示了全纯运动与Kleinian群变形空间的深刻同构性。\(\blacksquare\)

\subsection{x+2.16 超越整函数与Eremenko猜想}\label{x2.16-ux8d85ux8d8aux6574ux51fdux6570ux4e0eeremenkoux731cux60f3}

超越整函数的动力系统是多项式情形向无穷远处的自然推广。设 \(f: \mathbb{C} \to \mathbb{C}\) 为超越整函数，其本质奇点位于 \(\infty\) 处。与多项式不同，\(f\) 的Julia集非紧，且 \(\infty\) 属于Julia集。

\textbf{定理 x+2.15}（Baker定理------Julia集的连通性）：对于超越整函数，其Julia集可能具有无界连通分支（Baker域）。若 \(f\) 具有有限多个奇点值（即 \(f\) 为Speiser类函数），则 \(J(f)\) 为连通的，或由无穷多个紧致分支组成。

\textbf{证明}：Speiser类函数的Fatou集分支分类遵循多项式情形的一般图景，但引入Baker域（趋于 \(\infty\) 的Fatou分支）作为新的类型。若 Fatou 集由单个 Baker 域组成则 Julia 集连通。Eremenko猜想（1989年）进一步断言对于此类函数，Julia集的所有连通分支是单点。该猜想与复动力系统的逃逸集（escaping set）结构密切相关，由Rempe-Gillen和Rottenfusser近期工作中对轨道簇的解析延长性质获得了部分实质性进展。\(\blacksquare\)

\hrulefill

\chapter{差分方程与函数方程}\label{ux5deeux5206ux65b9ux7a0bux4e0eux51fdux6570ux65b9ux7a0b}

\hrulefill

\section{第330章：差分方程}\label{ux7b2c330ux7ae0ux5deeux5206ux65b9ux7a0b}

差分方程是离散自变量函数的方程，在数值分析、控制理论、信号处理和数学生物学中具有核心地位。本章系统讨论差分方程的理论体系。

\subsection{线性差分方程的基本理论}\label{ux7ebfux6027ux5deeux5206ux65b9ux7a0bux7684ux57faux672cux7406ux8bba}

\textbf{定义 268.1（差分算子）} 设 \(y_n\) 是定义在整数（或自然数）上的数列。定义前向差分算子 \(\Delta\) 和后向差分算子 \(\nabla\)： \[\Delta y_n = y_{n+1} - y_n, \quad \nabla y_n = y_n - y_{n-1}.\] 高阶差分定义为 \(\Delta^k y_n = \Delta(\Delta^{k-1} y_n)\)。位移算子 \(E\) 定义为 \(E y_n = y_{n+1}\)，从而 \(\Delta = E - I\)。

\textbf{定义 268.2（线性差分方程）} \(k\) 阶齐次线性差分方程的一般形式为 \[y_{n+k} + a_{k-1} y_{n+k-1} + \cdots + a_1 y_{n+1} + a_0 y_n = 0, \quad a_0 \neq 0.\] 其中系数 \(a_0, \ldots, a_{k-1}\) 为常数。

\textbf{定理 268.3（特征方程法）} 方程 \(y_{n+k} + a_{k-1} y_{n+k-1} + \cdots + a_0 y_n = 0\) 的特征方程为 \[\lambda^k + a_{k-1} \lambda^{k-1} + \cdots + a_1 \lambda + a_0 = 0.\]

（1）若特征方程有 \(k\) 个互异实根 \(\lambda_1, \ldots, \lambda_k\)，则通解为 \[y_n = c_1 \lambda_1^n + c_2 \lambda_2^n + \cdots + c_k \lambda_k^n.\]

（2）若 \(\lambda\) 是重数为 \(m\) 的根，则其贡献的解为 \[(c_1 + c_2 n + \cdots + c_m n^{m-1}) \lambda^n.\]

（3）若出现共轭复根 \(\lambda = r e^{\pm i\theta}\)，则对应实解为 \[r^n (A \cos n\theta + B \sin n\theta).\]

\textbf{证明} 设解具有形式 \(y_n = \lambda^n\)，代入差分方程： \[\lambda^{n+k} + a_{k-1} \lambda^{n+k-1} + \cdots + a_0 \lambda^n = 0 \implies \lambda^n (\lambda^k + a_{k-1} \lambda^{k-1} + \cdots + a_0) = 0.\] 故 \(\lambda\) 必须满足特征方程。\(k\) 阶方程的通解空间是 \(k\) 维的（由初始条件唯一确定），故上述 \(k\) 个线性无关解构成基。对于重根情形，验证 \(n \lambda^n, n^2 \lambda^n, \ldots, n^{m-1} \lambda^n\) 是解：由于 \(\lambda\) 同时是特征多项式 \(p(t)\) 及其前 \(m-1\) 阶导数的根，利用位移算子的代数性质即得。\(\blacksquare\)

\textbf{定理 268.4（与线性常微分方程的类比）} 线性差分方程与线性常微分方程之间存在深刻的结构类比：

\begin{longtable}[]{@{}ll@{}}
\toprule
微分方程 & 差分方程 \\
\midrule
\endhead
\bottomrule
\endlastfoot
\(\frac{dy}{dt} = ay\) & \(y_{n+1} = \lambda y_n\) \\
解：\(y(t) = Ce^{at}\) & 解：\(y_n = C\lambda^n\) \\
特征方程由 \(y = e^{rt}\) 代入 & 特征方程由 \(y_n = \lambda^n\) 代入 \\
稳定条件：\(\operatorname{Re}(r) < 0\) & 稳定条件：\(|\lambda| < 1\) \\
\end{longtable}

该类比源于指数映射将加法群 \((\mathbb{Z}, +)\) 的平移作用与乘法群 \((\mathbb{R}_{>0}, \times)\) 的伸缩作用相联系。具体地，若 \(y_{n+1} = \lambda y_n\)，则 \(y_n = \lambda^n y_0 = e^{n \ln \lambda} y_0\)，与微分方程 \(y' = (\ln \lambda) y\) 的离散化对应。

\subsection{常数变易法与待定系数法}\label{ux5e38ux6570ux53d8ux6613ux6cd5ux4e0eux5f85ux5b9aux7cfbux6570ux6cd5}

\textbf{定理 268.5（非齐次方程的通解结构）} \(k\) 阶非齐次线性差分方程 \[y_{n+k} + a_{k-1} y_{n+k-1} + \cdots + a_0 y_n = f(n)\] 的通解为 \(y_n = y_n^{(h)} + y_n^{(p)}\)，其中 \(y_n^{(h)}\) 是对应齐次方程的通解，\(y_n^{(p)}\) 是任一特解。

\textbf{证明} 设 \(y_n^{(p)}\) 是一个特解，\(y_n\) 是任意解。令 \(z_n = y_n - y_n^{(p)}\)。代入原方程： \[z_{n+k} + a_{k-1}z_{n+k-1} + \cdots + a_0 z_n = (y_{n+k} - y_{n+k}^{(p)}) + \cdots + a_0 (y_n - y_n^{(p)}) = f(n) - f(n) = 0.\] 故 \(z_n\) 是齐次解，即 \(z_n = y_n^{(h)}\)。因此 \(y_n = y_n^{(h)} + y_n^{(p)}\)。反之，\(y_n^{(h)} + y_n^{(p)}\) 确实满足原方程（齐次解代入得0，特解得 \(f(n)\)，和为 \(f(n)\)）。故通解恰为齐次通解加任一特解。\(\blacksquare\)

\textbf{定理 268.6（常数变易法）} 设齐次方程的基解为 \(\{y_n^{(1)}, y_n^{(2)}, \ldots, y_n^{(k)}\}\)。非齐次方程的特解可设为 \[y_n^{(p)} = c_1(n) y_n^{(1)} + c_2(n) y_n^{(2)} + \cdots + c_k(n) y_n^{(k)},\] 其中 \(c_i(n)\) 满足方程组 \[\sum_{i=1}^{k} \Delta c_i(n) \cdot y_{n+j}^{(i)} = \begin{cases} 0, & j = 1, \ldots, k-1 \\ f(n), & j = k \end{cases}.\]

\textbf{证明} 将 \(y_n^{(p)}\) 代入原方程并利用Casorati行列式（差分方程的Wronskian类比）非零性求解 \(\Delta c_i(n)\)。Casorati行列式 \(C(n) = \det(y_{n+j-1}^{(i)})_{i,j=1}^{k}\) 满足 \(C(n) = (-1)^k a_0 C(n-1)\)，非零性确保系数唯一确定。\(\blacksquare\)

\textbf{定理 268.7（待定系数法）} 当 \(f(n)\) 具有特殊形式时可直接猜测特解形式：

\begin{itemize}
\item
  若 \(f(n) = P_m(n) \lambda^n\)，其中 \(P_m\) 是 \(m\) 次多项式，则设 \(y_n^{(p)} = n^s Q_m(n) \lambda^n\)，其中 \(s\) 是 \(\lambda\) 作为特征根的重数，\(Q_m\) 是系数待定的 \(m\) 次多项式。
\item
  若 \(f(n) = A \cos(\omega n) + B \sin(\omega n)\)，则设 \(y_n^{(p)} = C \cos(\omega n) + D \sin(\omega n)\)（当 \(e^{\pm i\omega}\) 不是特征根时）。
\end{itemize}

\textbf{证明} 对于 \(f(n) = P_m(n)\lambda^n\)，将 \(y_n^{(p)} = n^s Q_m(n) \lambda^n\) 代入差分方程。利用位移算子 \(E\)：原方程可写为 \(P(E) y_n = f(n)\)，其中 \(P(E) = E^k + a_{k-1}E^{k-1} + \cdots + a_0\)。由于 \(P(E)[\lambda^n g(n)] = \lambda^n P(\lambda E) g(n)\)，代入后得到关于 \(Q_m\) 系数的线性方程组。\(s\) 的选择确保非齐次项与齐次解的共振被消除（\(s\) 为 \(\lambda\) 的重数使得 \(n^s\lambda^n\) 不再属于齐次解空间），此时方程组有唯一解。三角情形同理：利用 Euler 公式 \(\cos(\omega n) = \operatorname{Re}(e^{i\omega n})\) 将三角形式归约为指数形式。\(\blacksquare\)

\subsection{稳定性理论}\label{ux7a33ux5b9aux6027ux7406ux8bba}

\textbf{定义 268.8（稳定性的概念）} 差分方程 \(y_{n+1} = f(y_n)\) 的平衡点 \(y^*\) 满足 \(y^* = f(y^*)\)。

\begin{itemize}
\item
  \(y^*\) 是\textbf{Lyapunov稳定的}，若对任意 \(\varepsilon > 0\)，存在 \(\delta > 0\) 使 \(|y_0 - y^*| < \delta \implies |y_n - y^*| < \varepsilon\) 对所有 \(n\) 成立。
\item
  \(y^*\) 是\textbf{渐近稳定的}，若存在 \(\delta > 0\) 使 \(|y_0 - y^*| < \delta \implies \lim_{n \to \infty} y_n = y^*\)。
\end{itemize}

\textbf{定理 268.9（线性稳定性准则）} 对于线性系统 \(\mathbf{x}_{n+1} = A \mathbf{x}_n\)，零解渐近稳定当且仅当 \(A\) 的所有特征值满足 \(|\lambda| < 1\)（谱半径 \(\rho(A) < 1\)）。

\textbf{证明} （\(\Rightarrow\)）设零解渐近稳定。若存在特征值 \(\lambda\) 满足 \(|\lambda| \geq 1\)，取对应特征向量 \(\mathbf{v} \neq \mathbf{0}\)，则 \(\mathbf{x}_n = A^n \mathbf{v} = \lambda^n \mathbf{v}\)。当 \(|\lambda| > 1\) 时 \(\|\mathbf{x}_n\| \to \infty\)；当 \(|\lambda| = 1\) 时 \(\|\mathbf{x}_n\| = \|\mathbf{v}\|\) 不趋于零，均与渐近稳定性矛盾。

（\(\Leftarrow\)）设 \(\rho(A) < 1\)。由矩阵的 Jordan 分解 \(A = P J P^{-1}\)，\(A^n = P J^n P^{-1}\)。因 \(\rho(A) < 1\)，每个 Jordan 块 \(J_k(\lambda)\) 满足 \(J_k(\lambda)^n \to 0\)（\(|\lambda| < 1\) 时多项式乘以指数衰减）。故 \(A^n \to 0\)。对任意初值 \(\mathbf{x}_0\)，\(\mathbf{x}_n = A^n \mathbf{x}_0 \to \mathbf{0}\)。更精确地，\(\|\mathbf{x}_n\| \leq C \rho^n \|\mathbf{x}_0\|\) 对某 \(\rho \in (\rho(A), 1)\) 和常数 \(C\) 成立。\(\blacksquare\)

\textbf{定理 268.10（Jury判据 - 二阶情形）} 二阶实系数多项式 \(p(z) = z^2 + a_1 z + a_0\) 的所有根满足 \(|z| < 1\) 当且仅当 \[|a_0| < 1, \quad |a_1| < 1 + a_0.\]

\textbf{证明} 利用双线性（Möbius）变换 \(z = \frac{w+1}{w-1}\) 将单位圆盘映射到左半平面，然后应用Routh-Hurwitz判据。具体地，\(p(\frac{w+1}{w-1}) = 0\) 等价于 \((1 + a_1 + a_0) w^2 + 2(1 - a_0) w + (1 - a_1 + a_0) = 0\)。根在左半平面的Hurwitz条件为所有系数同号，即得上述条件。\(\blacksquare\)

\textbf{定理 268.11（Lyapunov函数法）} 考虑系统 \(\mathbf{x}_{n+1} = \mathbf{f}(\mathbf{x}_n)\)，其中 \(\mathbf{f}(\mathbf{0}) = \mathbf{0}\)。若存在连续函数 \(V: \mathbb{R}^d \to \mathbb{R}\) 满足 \(V(\mathbf{0}) = 0\) 且对于 \(\mathbf{x} \neq \mathbf{0}\) 有 \(V(\mathbf{x}) > 0\)，并且 \(\Delta V(\mathbf{x}) = V(\mathbf{f}(\mathbf{x})) - V(\mathbf{x}) < 0\)，则零解是渐近稳定的。

\textbf{证明} 这是连续动力系统Lyapunov定理的离散类比。\(V\) 沿轨道严格递减保证了 \(\lim_{n \to \infty} V(\mathbf{x}_n) = 0\)，由 \(V\) 的正定性得 \(\mathbf{x}_n \to \mathbf{0}\)。线性系统的构造：取 \(V(\mathbf{x}) = \mathbf{x}^T P \mathbf{x}\)，其中 \(P\) 满足离散Lyapunov方程 \(A^T P A - P = -Q\)（\(Q\) 正定）。\(\rho(A) < 1\) 保证正定解 \(P\) 存在。\(\blacksquare\)

\subsection{Z-变换与求解技巧}\label{z-ux53d8ux6362ux4e0eux6c42ux89e3ux6280ux5de7}

\textbf{定义 268.12（Z-变换）} 序列 \(\{x_n\}_{n=0}^{\infty}\) 的\textbf{Z-变换}定义为 \[\mathcal{Z}\{x_n\} = X(z) = \sum_{n=0}^{\infty} x_n z^{-n},\] 其中 \(z\) 是复变量，变换在收敛域 \(|z| > R\) 内有效（\(R\) 为收敛半径）。

\textbf{命题 268.13（Z-变换的基本性质）}

（1）（线性）\(\mathcal{Z}\{a x_n + b y_n\} = a X(z) + b Y(z)\)；

（2）（位移）\(\mathcal{Z}\{x_{n+k}\} = z^k X(z) - z^k x_0 - z^{k-1} x_1 - \cdots - z x_{k-1}\)；

（3）（卷积）\(\mathcal{Z}\{(x * y)_n\} = X(z) Y(z)\)，其中 \((x * y)_n = \sum_{k=0}^{n} x_k y_{n-k}\)；

（4）（尺度变换）\(\mathcal{Z}\{a^n x_n\} = X(z/a)\)。

\textbf{证明（卷积定理）} \[\mathcal{Z}\{(x * y)_n\} = \sum_{n=0}^{\infty} \left(\sum_{k=0}^{n} x_k y_{n-k}\right) z^{-n} = \sum_{k=0}^{\infty} x_k z^{-k} \sum_{n=k}^{\infty} y_{n-k} z^{-(n-k)} = X(z) Y(z).\] 交换求和次序的合法性由级数的绝对收敛性保证。\(\blacksquare\)

\textbf{定理 268.14（Z-反演公式）} \[x_n = \frac{1}{2\pi i} \oint_{C} X(z) z^{n-1} \, dz,\] 其中 \(C\) 是包围 \(X(z) z^{n-1}\) 所有奇点的正向围道。

\textbf{证明} 由 Z-变换定义 \(X(z) = \sum_{k=0}^{\infty} x_k z^{-k}\)，两边乘以 \(z^{n-1}\) 并在围道 \(C\) 上积分： \[\frac{1}{2\pi i} \oint_C X(z) z^{n-1} dz = \frac{1}{2\pi i} \oint_C \sum_{k=0}^{\infty} x_k z^{n-k-1} dz.\] 交换积分与求和（因级数在围道上一致收敛），得 \[\sum_{k=0}^{\infty} x_k \cdot \frac{1}{2\pi i} \oint_C z^{n-k-1} dz.\] 由 Cauchy 积分公式，\(\frac{1}{2\pi i} \oint_C z^{m-1} dz = \delta_{m,0}\)（Kronecker delta）。故仅当 \(n-k-1 = -1\)（即 \(k = n\)）时积分非零，值为 \(1\)。因此积分为 \(x_n\)，证毕。\(\blacksquare\)

\textbf{Z-变换求解差分方程的步骤}：（1）对差分方程两边取Z-变换，利用位移性质将差分方程化为 \(X(z)\) 的代数方程；（2）解出 \(X(z)\)；（3）通过部分分式展开和Z-反变换（或查表）得到 \(x_n\)。这一过程与Laplace变换求解微分方程完全平行。

\subsection{非线性差分方程与混沌}\label{ux975eux7ebfux6027ux5deeux5206ux65b9ux7a0bux4e0eux6df7ux6c8c}

\textbf{定义 268.15（Logistic映射）} \textbf{Logistic映射}是最简单的非线性差分方程之一： \[x_{n+1} = r x_n (1 - x_n), \quad 0 \leq x_n \leq 1, \quad 0 \leq r \leq 4.\]

该映射的动力学行为随参数 \(r\) 变化展现了完整的倍周期分岔路径到混沌的过渡：

\begin{itemize}
\tightlist
\item
  \(0 < r \leq 1\)：\(x^* = 0\) 是全局吸引的平衡点；
\item
  \(1 < r \leq 3\)：\(x^* = 1 - 1/r\) 是稳定的非平凡平衡点；
\item
  \(3 < r \leq 1 + \sqrt{6} \approx 3.449\)：出现稳定的2-周期轨道；
\item
  随后通过倍周期分岔序列（周期 \(2 \to 4 \to 8 \to \cdots\)）通向混沌；
\item
  \(r \approx 3.5699\)（Feigenbaum点）：周期分岔的积累点，此后进入混沌区。
\end{itemize}

\textbf{定理 268.16（Feigenbaum普适性）} 在倍周期分岔序列中，分岔参数值满足 \[\delta = \lim_{k \to \infty} \frac{r_k - r_{k-1}}{r_{k+1} - r_k} \approx 4.6692016\ldots,\] \[\alpha = \lim_{k \to \infty} \frac{d_k}{d_{k+1}} \approx 2.5029078\ldots,\] 其中 \(\delta\) 和 \(\alpha\) 是Feigenbaum常数，对一大类具有二次极大值的单峰映射普遍适用。该普适性源于重整化群算子在函数空间中的鞍点性质。

\hrulefill

\section{第331章：函数方程}\label{ux7b2c331ux7ae0ux51fdux6570ux65b9ux7a0b}

函数方程是以未知函数为变量、函数值之间满足的等式关系。许多基本函数最初即通过函数方程来刻画。

\subsection{Cauchy型函数方程}\label{cauchyux578bux51fdux6570ux65b9ux7a0b}

\textbf{定理 269.1（Cauchy方程的可测解）} 设 \(f: \mathbb{R} \to \mathbb{R}\) 满足 Cauchy 方程 \[f(x + y) = f(x) + f(y), \quad \forall x, y \in \mathbb{R}.\] 若 \(f\) 在一点连续（或在一测度为正的集合上有界，或 Lebesgue 可测），则 \(f(x) = cx\)，其中 \(c = f(1)\)。

\textbf{证明} （1）由加性得 \(f(0) = f(0 + 0) = f(0) + f(0)\)，故 \(f(0) = 0\)。（2）\(f(-x) = -f(x)\) 因 \(0 = f(x + (-x)) = f(x) + f(-x)\)。（3）归纳得 \(f(nx) = n f(x)\) 对所有 \(n \in \mathbb{N}\) 成立，进而对所有 \(n \in \mathbb{Z}\) 成立。（4）令 \(c = f(1)\)，则 \(f(q) = cq\) 对所有 \(q \in \mathbb{Q}\) 成立（因 \(f(p/q) = p f(1/q)\) 且 \(q f(1/q) = f(1) = c\)）。（5）若 \(f\) 在某点 \(x_0\) 处连续，则 \(f\) 处处连续：由 \(f(x) - f(x_0) = f(x - x_0)\) 得连续性延拓。（6）连续性结合有理点上的稠密性得 \(f(x) = cx\) 对所有 \(x \in \mathbb{R}\) 成立。\(\blacksquare\)

\textbf{定理 269.2（Hamel基与非平凡解的存在性）} 若不附加任何正则性条件，Cauchy方程存在非平凡（非线性）解。这等价于将 \(\mathbb{R}\) 视为 \(\mathbb{Q}\) 上的向量空间，任选一组Hamel基并定义 \(f\) 在基上的值。非平凡解的存在性依赖于选择公理（或等价的Hamel基的存在性）。

\textbf{证明} 将 \(\mathbb{R}\) 视为 \(\mathbb{Q}\) 上的向量空间（维数为连续统）。由选择公理，存在Hamel基 \(B = \{b_\alpha\}_{\alpha \in \Lambda} \subset \mathbb{R}\)，即每个 \(x \in \mathbb{R}\) 可唯一表示为 \(x = \sum q_\alpha b_\alpha\)（有限有理线性组合）。定义 \(f\) 在基上任意取值：对每个 \(b_\alpha\)，选取 \(f(b_\alpha) \in \mathbb{R}\)，然后线性延拓 \(f(\sum q_\alpha b_\alpha) = \sum q_\alpha f(b_\alpha)\)。若选取 \(f(b_\alpha)\) 不与 \(b_\alpha\) 成比例，则 \(f\) 满足加性方程但非 \(f(x) = cx\)。这样的 \(f\) 不可测且处处不连续。

这类解具有高度病态的性质：其图像在 \(\mathbb{R}^2\) 中稠密，且在任何区间上均无界（从而处处不连续、不可测）。\(\blacksquare\)

\textbf{定理 269.3（Cauchy方程的推广）}

（1）指数Cauchy方程：若 \(g(x + y) = g(x) g(y)\) 且 \(g\) 可测且不恒为零，则 \(g(x) = e^{cx}\) 或 \(g(x) = 0\)。这通过取对数 \(f(x) = \ln g(x)\) 归约为加法Cauchy方程。

（2）对数Cauchy方程：若 \(h(xy) = h(x) + h(y)\) 对 \(x, y > 0\) 成立且 \(h\) 可测，则 \(h(x) = c \ln x\)。令 \(u = \ln x, v = \ln y\) 即化为加法形式。

（3）乘性Cauchy方程：若 \(\varphi(xy) = \varphi(x) \varphi(y)\) 对 \(x, y > 0\) 成立且 \(\varphi\) 可测，则 \(\varphi(x) = x^c\) 或 \(\varphi(x) = 0\)。

\textbf{证明} （1）首先证 \(g(x) > 0\) 对所有 \(x\)：\(g(x) = g(x/2 + x/2) = g(x/2)^2 \geq 0\)。若存在 \(x_0\) 使 \(g(x_0) = 0\)，则 \(g(x) = g(x - x_0 + x_0) = g(x - x_0)g(x_0) = 0\)，与 \(g\) 不恒为零矛盾。故 \(g(x) > 0\)。令 \(f(x) = \ln g(x)\)，则 \(f(x+y) = f(x) + f(y)\)。由定理 269.1，\(f(x) = cx\)，故 \(g(x) = e^{cx}\)。

（2）令 \(f(u) = h(e^u)\)，\(u, v \in \mathbb{R}\)。则 \(f(u+v) = h(e^{u+v}) = h(e^u e^v) = h(e^u) + h(e^v) = f(u) + f(v)\)。由定理 269.1，\(f(u) = cu\)，故 \(h(x) = c \ln x\)。

（3）令 \(\psi(u) = \ln \varphi(e^u)\)（\(\varphi > 0\) 的情形）。则 \(\psi(u+v) = \ln \varphi(e^{u+v}) = \ln(\varphi(e^u)\varphi(e^v)) = \psi(u) + \psi(v)\)，故 \(\psi(u) = cu\)，\(\varphi(x) = e^{c \ln x} = x^c\)。若 \(\varphi(x_0) = 0\)，类似（1）可证 \(\varphi \equiv 0\)。\(\blacksquare\)

\subsection{Jensen方程与凸性}\label{jensenux65b9ux7a0bux4e0eux51f8ux6027}

\textbf{定理 269.4（Jensen方程）} 设 \(f: \mathbb{R} \to \mathbb{R}\) 满足 Jensen 方程 \[f\left(\frac{x + y}{2}\right) = \frac{f(x) + f(y)}{2}, \quad \forall x, y \in \mathbb{R}.\] 若 \(f\) 在一点连续（或局部有界），则 \(f(x) = ax + b\)，其中 \(a, b \in \mathbb{R}\)。

\textbf{证明} 令 \(g(x) = f(x) - f(0)\)，则 \(g(0) = 0\) 且 \(g(\frac{x+y}{2}) = \frac{g(x) + g(y)}{2}\)。取 \(y = 0\) 得 \(g(x/2) = g(x)/2\)，故 \(g(x+y) = 2g(\frac{x+y}{2}) = g(x) + g(y)\)，即 \(g\) 满足加法Cauchy方程。由连续性和Cauchy方程的正则条件得 \(g(x) = ax\)，从而 \(f(x) = ax + b\)，\(b = f(0)\)。\(\blacksquare\)

\textbf{推论 269.5（中点凸性）} 若 \(f: I \to \mathbb{R}\)（\(I\) 为区间）是Jensen凸的（即 \(f(\frac{x+y}{2}) \leq \frac{f(x)+f(y)}{2}\)）且在一点连续，则 \(f\) 是凸函数（满足 \(f(tx + (1-t)y) \leq t f(x) + (1-t) f(y)\) 对所有 \(t \in [0,1]\) 成立）。

\subsection{D'Alembert方程}\label{dalembertux65b9ux7a0b}

\textbf{定义 269.6（D'Alembert方程）} D'Alembert（或余弦）函数方程为 \[f(x + y) + f(x - y) = 2 f(x) f(y), \quad \forall x, y \in \mathbb{R}.\] 该方程以双曲余弦 \(\cosh(cx)\) 和三角余弦 \(\cos(cx)\) 为其基本解。

\textbf{定理 269.7（D'Alembert方程的正则解）} 设 \(f: \mathbb{R} \to \mathbb{R}\)（或 \(\mathbb{C}\)）是非零函数，满足D'Alembert方程且在某点可微（或连续，或在正测度集上有界）。则 \(f\) 具有形式 \[f(x) = \cosh(\alpha x) = \frac{e^{\alpha x} + e^{-\alpha x}}{2}, \quad \text{或} \quad f(x) = \cos(\beta x) = \frac{e^{i\beta x} + e^{-i\beta x}}{2},\] 其中 \(\alpha, \beta \in \mathbb{C}\)。此外 \(f(0) = 1\) 且 \(f\) 为偶函数。

\textbf{证明概要} （1）取 \(y = 0\) 得 \(2f(x) = 2f(x)f(0)\)，若 \(f\) 不恒为零则 \(f(0) = 1\)。（2）取 \(x = 0\) 得 \(f(y) + f(-y) = 2f(y)\)，故 \(f(-y) = f(y)\)，\(f\) 为偶函数。（3）考虑差分方程并通过对 \(y\) 求导（设 \(f\) 可微）的极限过程导出 \(f''(x) = \gamma f(x)\)，\(\gamma = f''(0)\)。根据 \(\gamma\) 的符号得双曲或三角余弦解。不依赖可微性假设的情形可通过序列逼近和函数方程迭代得出。\(\blacksquare\)

\textbf{定理 269.8（D'Alembert方程与三角恒等式）} 该方程本质上是余弦加法定理的抽象形式： \[\cos(a + b) + \cos(a - b) = 2 \cos a \cos b.\] 同理，双曲余弦满足 \(\cosh(a + b) + \cosh(a - b) = 2 \cosh a \cosh b\)。

\subsection{Pexider型方程及其推广}\label{pexiderux578bux65b9ux7a0bux53caux5176ux63a8ux5e7f}

\textbf{定义 269.9（Pexider方程）} Pexider方程是Cauchy方程的推广，涉及多个未知函数： \[f(x + y) = g(x) + h(y).\]

\textbf{定理 269.10（Pexider方程的解）} 设 \(f, g, h: \mathbb{R} \to \mathbb{R}\) 满足上述方程且在一点连续（或局部有界），则存在常数 \(a, b, c\) 使得 \[f(x) = a x + b + c, \quad g(x) = a x + b, \quad h(y) = a y + c.\]

\textbf{证明} 令 \(\phi(x) = g(x) - g(0)\)，\(\psi(y) = h(y) - h(0)\)。代入原方程得 \(f(x+y) = \phi(x) + \psi(y) + g(0) + h(0)\)。取 \(y = 0\)：\(f(x) = \phi(x) + g(0) + h(0)\)。取 \(x = 0\)：\(f(y) = \psi(y) + g(0) + h(0)\)。从而 \(f(x+y) - f(x) - f(y) = -(g(0) + h(0))\)。令 \(F(x) = f(x) - g(0) - h(0)\)，则 \(F(x+y) = F(x) + F(y)\)，归约为加法Cauchy方程。\(\blacksquare\)

\textbf{定理 269.11（推广的Pexider系统）} 以下均为常见的多函数方程系统：

（1）\(f(x + y) = g(x) h(y)\) \(\implies\) \(f, g, h\) 是指数型函数；

（2）\(f(xy) = g(x) + h(y)\) \(\implies\) \(f, g, h\) 是对数型函数；

（3）\(f(xy) = g(x) h(y)\) \(\implies\) \(f, g, h\) 是幂函数型；

\subsection{迭代函数方程与Schröder方程}\label{ux8fedux4ee3ux51fdux6570ux65b9ux7a0bux4e0eschruxf6derux65b9ux7a0b}

\textbf{定义 269.12（迭代函数方程）} 设 \(\varphi\) 是已知函数。求 \(f\) 使得 \(f\) 与 \(\varphi\) 可交换：\(f(\varphi(x)) = \varphi(f(x))\)，或求 \(f\) 满足 \(f(f(x)) = \varphi(x)\)（Babbage方程，即求 \(\varphi\) 的迭代平方根）。

\textbf{定理 269.13（Schröder方程）} 给定函数 \(\varphi\)，Schröder方程寻求 \(\sigma\) 和常数 \(s\) 满足 \[\sigma(\varphi(x)) = s \cdot \sigma(x).\] 此方程将 \(\varphi\) 的迭代线性化：若 \(\sigma\) 可逆，则 \(\varphi(x) = \sigma^{-1}(s \cdot \sigma(x))\)，从而 \[\varphi^n(x) = \sigma^{-1}(s^n \cdot \sigma(x)).\]

\textbf{证明} 验证线性化公式：由 Schröder 方程，\(\sigma(\varphi(x)) = s \cdot \sigma(x)\)，故 \(\varphi(x) = \sigma^{-1}(s \cdot \sigma(x))\)（若 \(\sigma\) 可逆）。对 \(n\) 归纳：\(n=1\) 时即定义；设 \(\varphi^n(x) = \sigma^{-1}(s^n \cdot \sigma(x))\)，则 \(\varphi^{n+1}(x) = \varphi(\varphi^n(x)) = \sigma^{-1}(s \cdot \sigma(\varphi^n(x))) = \sigma^{-1}(s \cdot s^n \cdot \sigma(x)) = \sigma^{-1}(s^{n+1} \cdot \sigma(x))\)。故公式成立。

例如，对 \(\varphi(x) = x^2\)（在 \(x > 0\) 上），取 \(\sigma(x) = \ln x\)，\(s = 2\)，则 \(\ln(x^2) = 2 \ln x\)，故 \(\varphi^n(x) = x^{2^n}\)。\(\blacksquare\)

\textbf{定理 269.14（Abel方程）} 与Schröder方程密切相关的是Abel方程： \[\alpha(\varphi(x)) = \alpha(x) + 1.\] 这对应于Schröder方程中 \(s = 1\) 的极限情形，或等价于 \(\alpha = \sigma^{-1} \circ \log_s\)。Abel方程的解 \(\alpha\) 提供了迭代的半群结构：\(\varphi^n(x) = \alpha^{-1}(\alpha(x) + n)\)。

\textbf{证明} 与 Schröder 方程的关联：若 \(\sigma\) 满足 \(\sigma(\varphi(x)) = s \cdot \sigma(x)\)（\(s > 0, s \neq 1\)），令 \(\alpha(x) = \log_s(\sigma(x))\)，则 \(\alpha(\varphi(x)) = \log_s(\sigma(\varphi(x))) = \log_s(s \cdot \sigma(x)) = 1 + \log_s(\sigma(x)) = \alpha(x) + 1\)。反之，若 \(\alpha\) 满足 Abel 方程，令 \(\sigma(x) = s^{\alpha(x)}\)（对某 \(s > 0\)），则 \(\sigma(\varphi(x)) = s^{\alpha(\varphi(x))} = s^{\alpha(x) + 1} = s \cdot \sigma(x)\)。迭代公式：\(\varphi^n(x) = \alpha^{-1}(\alpha(x) + n)\)，由归纳直接验证。\(\blacksquare\)

\textbf{定理 269.15（Koenigs定理）} 设 \(\varphi\) 在不动点 \(x_0\)（\(\varphi(x_0) = x_0\)）附近是解析的，且乘子 \(\lambda = \varphi'(x_0)\) 满足 \(0 < |\lambda| < 1\) 或 \(|\lambda| > 1\)（即不动点是\textbf{双曲的}）。则Schröder方程存在局部解析解 \(\sigma\)，满足 \(\sigma(x_0) = 0\)，\(\sigma'(x_0) = 1\)，且 \(s = \lambda\)。

\textbf{证明} 不妨设 \(x_0 = 0\) 且考虑 \(|\lambda| < 1\) 的情形（\(|\lambda| > 1\) 的情形可考虑逆映射）。定义 \(\sigma(x) = \lim_{n \to \infty} \lambda^{-n} \varphi^n(x)\)。由解析性，\(\varphi(x) = \lambda x + a_2 x^2 + \cdots\)。关键估计：当 \(n \to \infty\) 时 \(\varphi^n(x) = \lambda^n x + O(|\lambda|^{2n})\)（因高阶项指数衰减快于主导项），故 \(\lambda^{-n} \varphi^n(x) \to x + \text{解析函数}\)。此极限在 \(0\) 的某邻域内一致收敛（使用 Cauchy 估计控制解析函数的增长），故 \(\sigma\) 解析。验证 Schröder 方程：\(\sigma(\varphi(x)) = \lim_{n \to \infty} \lambda^{-n} \varphi^{n+1}(x) = \lambda \cdot \lim_{n \to \infty} \lambda^{-(n+1)} \varphi^{n+1}(x) = \lambda \cdot \sigma(x)\)。\(\sigma(0) = 0\) 且由展开知 \(\sigma'(0) = 1\)。

该定理由Poincaré（1890）和Koenigs（1884）证明，是现代迭代理论（复动力系统）的基石，将非线性迭代局部化为线性乘法的共轭。\(\blacksquare\)

\newpage

\chapter{10-概率与统计/01-概率论/01-测度论概率基础.md}\label{ux6982ux7387ux4e0eux7edfux8ba101-ux6982ux7387ux8bba01-ux6d4bux5ea6ux8bbaux6982ux7387ux57faux7840.md}

\section{第332章：测度论概率基础}\label{ux7b2c332ux7ae0ux6d4bux5ea6ux8bbaux6982ux7387ux57faux7840}

Kolmogorov 在 1933 年将概率论建立在测度论之上，使概率论成为数学分析的一个分支。本章从概率空间出发，用测度论语言严格定义随机变量、分布、期望和独立性。

\subsection{104.1 概率空间}\label{ux6982ux7387ux7a7aux95f4}

\textbf{定义 104.1}（概率空间）：一个\textbf{概率空间} \((\Omega, \mathcal{F}, \mathbb{P})\) 由以下组成： - \(\Omega\)：样本空间（所有可能结果的集合） - \(\mathcal{F}\)：\(\Omega\) 上的 \(\sigma\)-代数（事件域） - \(\mathbb{P}: \mathcal{F} \to [0, 1]\)：概率测度，满足 \(\mathbb{P}(\Omega) = 1\) 和可数可加性

\textbf{定义 104.2}（\(\sigma\)-代数）：\(\Omega\) 的子集族 \(\mathcal{F}\) 称为 \(\sigma\)\textbf{-代数}，如果： 1. \(\Omega \in \mathcal{F}\) 2. \(A \in \mathcal{F} \Rightarrow A^c \in \mathcal{F}\)（对补集封闭） 3. \(A_1, A_2, \ldots \in \mathcal{F} \Rightarrow \bigcup_{n=1}^\infty A_n \in \mathcal{F}\)（对可数并封闭）

\textbf{定义 104.3}（Borel \(\sigma\)-代数）：\(\mathbb{R}\) 上的 \textbf{Borel \(\sigma\)-代数} \(\mathcal{B}(\mathbb{R})\) 是包含所有开集的最小 \(\sigma\)-代数。\(\mathcal{B}(\mathbb{R}^n)\) 类似定义。

\textbf{命题 104.1}（概率测度的基本性质）： 1. \(\mathbb{P}(\varnothing) = 0\) 2. 有限可加性：若 \(A_1, \ldots, A_n\) 互不相交，则 \(\mathbb{P}(\bigcup_{i=1}^n A_i) = \sum_{i=1}^n \mathbb{P}(A_i)\) 3. 单调性：\(A \subseteq B \Rightarrow \mathbb{P}(A) \leq \mathbb{P}(B)\) 4. 连续性：\(A_n \uparrow A\)（单调解）\(\Rightarrow \mathbb{P}(A_n) \uparrow \mathbb{P}(A)\)；\(A_n \downarrow A \Rightarrow \mathbb{P}(A_n) \downarrow \mathbb{P}(A)\) 5. 次可加性：\(\mathbb{P}(\bigcup_{n=1}^\infty A_n) \leq \sum_{n=1}^\infty \mathbb{P}(A_n)\)

\emph{证明}：直接由测度的一般性质（V9 Ch 65）推导。∎

\subsection{104.2 随机变量与分布}\label{ux968fux673aux53d8ux91cfux4e0eux5206ux5e03}

\textbf{定义 104.4}（随机变量）：\((\Omega, \mathcal{F}, \mathbb{P})\) 上的\textbf{随机变量} \(X\) 是可测函数 \(X: \Omega \to \mathbb{R}\)（即对每个 Borel 集 \(B \in \mathcal{B}(\mathbb{R})\)，\(X^{-1}(B) \in \mathcal{F}\)）。

\(n\) 维\textbf{随机向量} \(\mathbf{X} = (X_1, \ldots, X_n)\) 是可测函数 \(\mathbf{X}: \Omega \to \mathbb{R}^n\)。

\textbf{定义 104.5}（分布）：随机变量 \(X\) 的\textbf{分布}（或\textbf{律}）\(\mu_X\) 是 \(\mathbb{R}\) 上的概率测度，定义为 \(\mu_X(B) = \mathbb{P}(X \in B) = \mathbb{P}(X^{-1}(B))\)。

\textbf{定义 104.6}（分布函数）：\(X\) 的\textbf{累积分布函数}（CDF）为 \(F_X(x) = \mathbb{P}(X \leq x) = \mu_X((-\infty, x])\)。\(F_X\) 是非降右连续函数，满足 \(\lim_{x \to -\infty} F_X(x) = 0\)，\(\lim_{x \to \infty} F_X(x) = 1\)。

\textbf{定义 104.7}（概率密度函数）：若存在非负可测函数 \(f_X\) 使得 \(\mu_X(B) = \int_B f_X(x) dx\)，则 \(f_X\) 称为 \(X\) 的\textbf{概率密度函数}（PDF）。此时 \(X\) 称为\textbf{绝对连续}的。

\subsection{104.3 期望的测度论定义}\label{ux671fux671bux7684ux6d4bux5ea6ux8bbaux5b9aux4e49}

\textbf{定义 104.8}（期望）：随机变量 \(X\) 的\textbf{期望}（或\textbf{均值}）定义为 Lebesgue 积分：

\[\mathbb{E}[X] = \int_\Omega X(\omega) d\mathbb{P}(\omega)\]

若 \(\mathbb{E}[|X|] < \infty\)，则称 \(X\) \textbf{可积}。

\textbf{命题 104.2}（期望的转化公式）：\(\mathbb{E}[X] = \int_{\mathbb{R}} x \, d\mu_X(x)\)。更一般地，对可测函数 \(g\)，\(\mathbb{E}[g(X)] = \int_{\mathbb{R}} g(x) d\mu_X(x)\)。若 \(X\) 有密度 \(f_X\)，则 \(\mathbb{E}[g(X)] = \int_{\mathbb{R}} g(x) f_X(x) dx\)。

\textbf{定义 104.9}（方差与矩）：\(X\) 的\textbf{方差} \(\operatorname{Var}(X) = \mathbb{E}[(X - \mathbb{E}[X])^2] = \mathbb{E}[X^2] - (\mathbb{E}[X])^2\)。\(k\) 阶矩 \(\mathbb{E}[X^k]\)，\(k\) 阶中心矩 \(\mathbb{E}[(X - \mathbb{E}[X])^k]\)。

\textbf{定理 104.3}（Jensen 不等式）：若 \(\varphi\) 是凸函数且 \(X\) 可积，则 \(\varphi(\mathbb{E}[X]) \leq \mathbb{E}[\varphi(X)]\)（假设 \(\mathbb{E}[\varphi(X)]\) 存在）。

\emph{证明}：凸函数 \(\varphi\) 存在支撑线：\(\varphi(x) \geq \varphi(x_0) + a(x - x_0)\)。取 \(x_0 = \mathbb{E}[X]\)，两边求期望。∎

\textbf{定理 104.4}（Hölder 不等式）：\(\mathbb{E}[|XY|] \leq (\mathbb{E}[|X|^p])^{1/p} (\mathbb{E}[|Y|^q])^{1/q}\)，其中 \(1/p + 1/q = 1\)（\(p, q > 1\)）。

\textbf{证明}：当 \(\mathbb{E}[|X|^p] = 0\) 或 \(\mathbb{E}[|Y|^q] = 0\) 时显然。否则令 \(U = |X|/(\mathbb{E}[|X|^p])^{1/p}\)，\(V = |Y|/(\mathbb{E}[|Y|^q])^{1/q}\)。由 Young 不等式 \(UV \leq U^p/p + V^q/q\)。取期望，\(\mathbb{E}[UV] \leq 1\)，化回 \(X,Y\) 即得不等式。等号成立当且仅当 \(|X|^p\) 与 \(|Y|^q\) a.s. 成比例。\(\blacksquare\)

\textbf{推论 104.5}（Cauchy-Schwarz 不等式）：\(\mathbb{E}[|XY|] \leq \sqrt{\mathbb{E}[X^2] \mathbb{E}[Y^2]}\)（\(p = q = 2\)）。

\subsection{104.4 独立性}\label{ux72ecux7acbux6027}

\textbf{定义 104.10}（独立事件）：事件 \(A_1, \ldots, A_n \in \mathcal{F}\) 称为\textbf{独立}的，如果对任意子集 \(\{i_1, \ldots, i_k\}\)，

\[\mathbb{P}(A_{i_1} \cap \cdots \cap A_{i_k}) = \mathbb{P}(A_{i_1}) \cdots \mathbb{P}(A_{i_k})\]

\textbf{定义 104.11}（独立 \(\sigma\)-代数）：\(\sigma\)-子代数 \(\mathcal{G}_1, \ldots, \mathcal{G}_n \subseteq \mathcal{F}\) 称为\textbf{独立}的，如果对任意 \(A_i \in \mathcal{G}_i\)，\(A_1, \ldots, A_n\) 独立。

\textbf{定义 104.12}（独立随机变量）：随机变量 \(X_1, \ldots, X_n\) 称为\textbf{独立}的，如果它们生成的 \(\sigma\)-代数 \(\sigma(X_1), \ldots, \sigma(X_n)\) 独立。等价地，\(\mathbb{P}(X_1 \in B_1, \ldots, X_n \in B_n) = \prod_{i=1}^n \mathbb{P}(X_i \in B_i)\) 对所有 Borel 集 \(B_i\) 成立。

\textbf{命题 104.6}：若 \(X_1, \ldots, X_n\) 独立且可积，则 \(\mathbb{E}[X_1 \cdots X_n] = \mathbb{E}[X_1] \cdots \mathbb{E}[X_n]\)。若两两不相关（\(\operatorname{Cov}(X_i, X_j) = 0\)），则 \(\operatorname{Var}(\sum X_i) = \sum \operatorname{Var}(X_i)\)。

\hrulefill

\hrulefill

\hrulefill

\hrulefill

\hrulefill

\hrulefill

\section{第333章：大数定律与中心极限定理}\label{ux7b2c333ux7ae0ux5927ux6570ux5b9aux5f8bux4e0eux4e2dux5fc3ux6781ux9650ux5b9aux7406}

大数定律（LLN）和中心极限定理（CLT）是概率论的两大基石。LLN 表明样本均值收敛于期望，CLT 描述了围绕均值的波动行为。

\subsection{105.1 大数定律}\label{ux5927ux6570ux5b9aux5f8b}

\textbf{定理 105.1}（弱大数定律，Khinchine）：设 \(X_1, X_2, \ldots\) 是独立同分布（i.i.d.）且 \(\mathbb{E}[|X_1|] < \infty\)。令 \(\mu = \mathbb{E}[X_1]\)。则样本均值 \(\bar{X}_n = \frac{1}{n} \sum_{i=1}^n X_i\) 依概率收敛于 \(\mu\)：

\[\bar{X}_n \xrightarrow{\mathbb{P}} \mu \quad (n \to \infty)\]

\emph{证明}：使用特征函数。\(\varphi_{\bar{X}_n}(t) = \varphi_{X_1}(t/n)^n \to e^{i\mu t}\)（由 \(\varphi_{X_1}(t) = 1 + i\mu t + o(t)\)）。由连续性定理，\(\bar{X}_n \xrightarrow{d} \mu\)（常数），故依概率收敛。∎

\textbf{定理 105.2}（强大数定律，Kolmogorov）：设 \(X_1, X_2, \ldots\) 是 i.i.d. 且 \(\mathbb{E}[|X_1|] < \infty\)。则

\[\bar{X}_n \xrightarrow{\text{a.s.}} \mu \quad \text{几乎必然}\]

即 \(\mathbb{P}(\lim_{n \to \infty} \bar{X}_n = \mu) = 1\)。

\textbf{证明}：设\(X_n\)为独立同分布随机变量，\(\mathbb{E}[X_1]=\mu\)，\(\mathbb{E}[|X_1|]<\infty\)。定义\(S_n=\sum_{i=1}^n X_i\)。由Kolmogorov强大数定律（或Etemadi定理的简化版本），对任意\(\varepsilon>0\)，\(\Pr(\sup_{k\ge n}|S_k/k-\mu|>\varepsilon)\to 0\)（\(n\to\infty\)）。等价地，\(\Pr(|S_n/n-\mu|>\varepsilon)\to 0\)。此由Chebyshev不等式：\(\Pr(|S_n/n-\mu|>\varepsilon)\le\operatorname{Var}(S_n/n)/\varepsilon^2=\sigma^2/(n\varepsilon^2)\to 0\)（假设方差有限；一般情形使用截断方法）。\(\blacksquare\)

\textbf{定理 105.3}（Kolmogorov 0-1 律）：设 \(X_1, X_2, \ldots\) 独立。尾 \(\sigma\)-代数 \(\mathcal{T} = \bigcap_{n=1}^\infty \sigma(X_n, X_{n+1}, \ldots)\) 中的每个事件概率为 \(0\) 或 \(1\)。特别地，\(\limsup \bar{X}_n\) 是常数（几乎必然）。

\textbf{证明}：尾 \(\sigma\)-代数 \(\mathcal{T} = \bigcap_n \sigma(X_k : k \geq n)\) 中的任意事件 \(A\) 与任意有限个 \(X_1,\ldots,X_m\) 独立。因为 \(A \in \sigma(X_{m+1}, X_{m+2}, \ldots)\)，而 \(\sigma(X_1,\ldots,X_m)\) 与 \(\sigma(X_{m+1},\ldots)\) 独立（由独立性）。故 \(A\) 与 \(\bigcup_m \sigma(X_1,\ldots,X_m)\) 生成的代数独立，由 Dynkin \(\pi\)-\(\lambda\) 定理，\(A\) 与整个 \(\sigma(X_1,X_2,\ldots)\) 独立。由于 \(A \in \mathcal{T} \subseteq \sigma(X_1,X_2,\ldots)\)，\(A\) 与自身独立：\(\mathbb{P}(A) = \mathbb{P}(A \cap A) = \mathbb{P}(A)^2\)，故 \(\mathbb{P}(A) \in \{0,1\}\)。\(\blacksquare\)

\subsection{105.2 中心极限定理}\label{ux4e2dux5fc3ux6781ux9650ux5b9aux7406}

\textbf{定理 105.4}（Lindeberg-Lévy 中心极限定理）：设 \(X_1, X_2, \ldots\) 是 i.i.d.，\(\mathbb{E}[X_1] = \mu\)，\(\operatorname{Var}(X_1) = \sigma^2 \in (0, \infty)\)。则

\[\frac{\sqrt{n}(\bar{X}_n - \mu)}{\sigma} \xrightarrow{d} \mathcal{N}(0, 1) \quad (n \to \infty)\]

即标准化样本均值依分布收敛于标准正态分布。

\emph{证明}：使用特征函数。\(\varphi_{Z_n}(t) = \varphi_{X_1-\mu}\left(\frac{t}{\sigma\sqrt{n}}\right)^n\)。由 Taylor 展开，\(\varphi_{X_1-\mu}(s) = 1 - \frac{\sigma^2 s^2}{2} + o(s^2)\)。故 \(\varphi_{Z_n}(t) \to e^{-t^2/2}\)，即 \(\mathcal{N}(0, 1)\) 的特征函数。∎

\textbf{定理 105.5}（Lindeberg-Feller CLT，三角阵列）：设对每个 \(n\)，\(X_{n,1}, \ldots, X_{n, k_n}\) 独立，\(\mathbb{E}[X_{n,i}] = 0\)，\(\sum_{i=1}^{k_n} \operatorname{Var}(X_{n,i}) = 1\)。若 \textbf{Lindeberg 条件}成立：

\[\lim_{n \to \infty} \sum_{i=1}^{k_n} \mathbb{E}[X_{n,i}^2 \mathbf{1}_{\{|X_{n,i}| > \varepsilon\}}] = 0 \quad (\forall \varepsilon > 0)\]

则 \(\sum_{i=1}^{k_n} X_{n,i} \xrightarrow{d} \mathcal{N}(0, 1)\)。Lindeberg 条件保证了没有单个变量主导和的行为。

\textbf{证明}：设 \(\varphi_{n,i}\) 为 \(X_{n,i}\) 的特征函数。截断展开：\(\varphi_{n,i}(t) = 1 - t^2 \sigma_{n,i}^2/2 + R_{n,i}(t)\)，其中 \(\sigma_{n,i}^2 = \mathbb{E}[X_{n,i}^2]\)，余项满足 \(|R_{n,i}(t)| \leq \varepsilon t^2 \sigma_{n,i}^2 + 2 \mathbb{E}[X_{n,i}^2 \mathbf{1}_{\{|X_{n,i}| > \varepsilon\}}]\)。由 Lindeberg 条件，对乘积 \(\prod_i \varphi_{n,i}(t)\) 取对数并展开，主项 \(\sum (-\frac{t^2}{2}\sigma_{n,i}^2) = -\frac{t^2}{2}\)，余项一致趋于 \(0\)。故特征函数逐点收敛于 \(e^{-t^2/2}\)，由 Levy 连续性得定理。\(\blacksquare\)

\textbf{定理 105.6}（Berry-Esseen 界）：在上述 i.i.d. CLT 条件下，若 \(\mathbb{E}[|X_1|^3] < \infty\)，则

\[\sup_{x \in \mathbb{R}} \left| \mathbb{P}\left(\frac{\sqrt{n}(\bar{X}_n - \mu)}{\sigma} \leq x\right) - \Phi(x) \right| \leq \frac{C \mathbb{E}[|X_1 - \mu|^3]}{\sigma^3 \sqrt{n}}\]

其中 \(\Phi\) 是标准正态分布函数，\(C \leq 0.4748\)（已知最优常数）。

\textbf{证明}（Essen 光滑引理法）：设 \(Z_n = \frac{\sqrt{n}(\bar{X}_n - \mu)}{\sigma}\)，\(F_n\) 为分布函数，\(\varphi_n\) 为特征函数。由 Ess en 不等式，\(\sup_x |F_n(x) - \Phi(x)| \leq C \int_{-T}^T |\frac{\varphi_n(t) - e^{-t^2/2}}{t}| dt + C/T\)。利用 \(|e^{itX} - (1 + itX - \frac{t^2 X^2}{2})| \leq \min(|tX|^2, |tX|^3/6)\)，展开 \(\varphi_n(t) = (1 - \frac{t^2}{2n})^n + O(|t|^3 n^{-1/2})\)。取 \(T = C\sqrt{n}\)，积分估计给出 \(O(n^{-1/2})\) 界。常数 \(0.4748\) 来自 Shevtsova 的改进。\(\blacksquare\)

\subsection{105.3 特征函数}\label{ux7279ux5f81ux51fdux6570}

\textbf{定义 105.1}（特征函数）：随机变量 \(X\) 的\textbf{特征函数}为 \(\varphi_X(t) = \mathbb{E}[e^{itX}] = \int_{\mathbb{R}} e^{itx} d\mu_X(x)\)。特征函数是 Fourier 变换（或 Fourier-Stieltjes 变换）。

\textbf{命题 105.7}（特征函数的性质）： 1. \(|\varphi_X(t)| \leq 1\)，\(\varphi_X(0) = 1\) 2. \(\varphi_X\) 一致连续 3. \(\varphi_{aX+b}(t) = e^{itb} \varphi_X(at)\) 4. 若 \(X, Y\) 独立，则 \(\varphi_{X+Y}(t) = \varphi_X(t) \varphi_Y(t)\) 5. 若 \(\mathbb{E}[|X|^k] < \infty\)，则 \(\varphi_X^{(k)}(0) = i^k \mathbb{E}[X^k]\)

\textbf{定理 105.8}（Levy 连续性定理）：设 \(X_n\) 是随机变量序列，\(\varphi_n\) 是其特征函数。若 \(\varphi_n(t) \to \varphi(t)\) 逐点，且 \(\varphi\) 在 \(0\) 连续，则 \(\varphi\) 是某个随机变量的特征函数，且 \(X_n \xrightarrow{d} X\)（\(\varphi_X = \varphi\)）。

\textbf{证明}：由特征函数一致等度连续（\(\varphi_n(0)=1\)，\(\varphi\) 在 0 连续），序列 \(\mu_n\) 是胎紧的：\(\forall \varepsilon > 0 \exists K : \sup_n \mu_n(|X_n| > K) < \varepsilon\)。由 Prokhorov 定理，存在弱收敛子列 \(\mu_{n_k} \xrightarrow{w} \mu\)。相应特征函数 \(\varphi_{n_k}(t) \to \hat{\mu}(t)\)，故 \(\hat{\mu} = \varphi\)。\(\varphi\) 在 0 的连续性确保 \(\mu\) 是概率测度。由于每个收敛子列极限相同（由特征函数唯一性），全序列弱收敛。\(\blacksquare\)

\hrulefill

\section{第334章：概率论极限理论深化}\label{ux7b2c334ux7ae0ux6982ux7387ux8bbaux6781ux9650ux7406ux8bbaux6df1ux5316}

中心极限定理的经典 Lindeberg-Lévy 版本要求独立同分布和有限方差，但在实际统计应用中，常遇到独立但不同分布的数据（如三角阵列）或需要量化正态逼近速率的场景。本章在上一章的基础上深入探讨极限理论的三个进阶方向：Lindeberg-Feller 中心极限定理------给出独立不同分布的三角阵列收敛于正态分布的充分必要条件，证明 Lindeberg 条件的必要性和充分性；Berry-Esseen 定理------量化 i.i.d. 情形下分布函数逼近正态分布的误差速率，其 \(O(n^{-1/2})\) 收敛速率是最优的；以及 Cramér 大偏差定理------研究中心极限定理尾部的精确渐近行为，给出偏离均值较远处的稀有事件概率的对数渐近公式。

\subsection{Lindeberg-Feller中心极限定理}\label{lindeberg-fellerux4e2dux5fc3ux6781ux9650ux5b9aux7406}

\textbf{定理271.1（Lindeberg-Feller CLT------充分必要条件）} 设对每个 \(n\)，\(X_{n,1}, \ldots, X_{n, k_n}\) 是独立随机变量（行内独立，行间可任意），满足 \(\mathbb{E}[X_{n,i}] = 0\) 且 \(0 < \sigma_{n,i}^2 = \operatorname{Var}(X_{n,i}) < \infty\)。令 \(s_n^2 = \sum_{i=1}^{k_n} \sigma_{n,i}^2\)。则以下两条等价：

（i）Lindeberg 条件成立：对任意 \(\varepsilon > 0\)， \[\lim_{n \to \infty} \frac{1}{s_n^2} \sum_{i=1}^{k_n} \mathbb{E}[X_{n,i}^2 \mathbf{1}_{\{|X_{n,i}| > \varepsilon s_n\}}] = 0\]

（ii）\(\frac{1}{s_n} \sum_{i=1}^{k_n} X_{n,i} \xrightarrow{d} \mathcal{N}(0, 1)\) 且 Feller 条件（\(\max_{i} \sigma_{n,i}^2 / s_n^2 \to 0\)）成立。

\textbf{证明}（（i）\(\Rightarrow\)（ii）------充分性）标准化 \(Y_{n,i} = X_{n,i}/s_n\)。由 Lindeberg 条件可推出 Feller 条件（因对 \(\varepsilon > 0\)，\(\sigma_{n,i}^2/s_n^2 = \mathbb{E}[Y_{n,i}^2 \mathbf{1}_{|Y_{n,i}| \leq \varepsilon}] + \mathbb{E}[Y_{n,i}^2 \mathbf{1}_{|Y_{n,i}| > \varepsilon}] \leq \varepsilon^2 + \sum_j \mathbb{E}[Y_{n,j}^2 \mathbf{1}_{|Y_{n,j}| > \varepsilon}]\)，由 Lindeberg 令 \(n \to \infty\)）。对特征函数 \(\varphi_{n,i}(t) = \mathbb{E}[e^{itY_{n,i}}]\) 作 Taylor 展开至二阶：\(\varphi_{n,i}(t) = 1 - \frac{t^2}{2}\sigma_{n,i}^2/s_n^2 + R_{n,i}\)，其中余项满足 \(|R_{n,i}| \leq |t|^3 \mathbb{E}[|Y_{n,i}|^3 \mathbf{1}_{|Y_{n,i}| \leq \varepsilon}]/6 + t^2 \mathbb{E}[Y_{n,i}^2 \mathbf{1}_{|Y_{n,i}| > \varepsilon}]\)。由 Lindeberg 条件和 \(\sum_i \mathbb{E}[|Y_{n,i}|^3 \mathbf{1}_{|Y_{n,i}| \leq \varepsilon}] \leq \varepsilon \sum_i \sigma_{n,i}^2/s_n^2 = \varepsilon\) 得 \(\sum_i |R_{n,i}| \to 0\)。对乘积特征函数取对数展开：\(\ln \prod_i \varphi_{n,i}(t) = \sum_i (\varphi_{n,i}(t) - 1) + o(1) = -t^2/2 + o(1)\)，故特征函数收敛于 \(e^{-t^2/2}\)，由 Lévy 连续性定理得依分布收敛。

（（ii）\(\Rightarrow\)（i）------必要性）设 \(Z_n = \sum_i Y_{n,i} \xrightarrow{d} \mathcal{N}(0, 1)\)。由正态分布的矩存在性，\(\mathbb{E}[Z_n^2] = 1\) 且 \(\sup_n \mathbb{E}[Z_n^4] < \infty\)。由 Lévy 距离和胎紧性，\(\lim_n \sum_i \mathbb{E}[Y_{n,i}^2 \mathbf{1}_{|Y_{n,i}| > \varepsilon}] = \lim_n (\mathbb{E}[Z_n^2] - \sum_i \mathbb{E}[Y_{n,i}^2 \mathbf{1}_{|Y_{n,i}| \leq \varepsilon}]) = 0\)。\(\blacksquare\)

\subsection{Berry-Esseen定理}\label{berry-esseenux5b9aux7406}

\textbf{定理271.2（Berry-Esseen定理------精确陈述）} 设 \(X_1, X_2, \ldots\) 是独立同分布随机变量，\(\mathbb{E}[X_1] = \mu\)，\(\operatorname{Var}(X_1) = \sigma^2 > 0\)，且 \(\mathbb{E}[|X_1|^3] = \rho_3 < \infty\)。令 \(F_n\) 为标准化样本均值 \(Z_n = \frac{\sqrt{n}(\bar{X}_n - \mu)}{\sigma}\) 的分布函数，\(\Phi\) 为标准正态分布函数。则

\[\sup_{x \in \mathbb{R}} |F_n(x) - \Phi(x)| \leq \frac{C \rho_3}{\sigma^3 \sqrt{n}}\]

其中 \(C \leq 0.4748\) 为最优常数（Shevtsova，2013）。该不等式给出了 CLT 收敛速率的最优 \(O(n^{-1/2})\) 阶，且常数 \(C\) 是普适的。

\textbf{证明}（Essen光滑引理方法）

\textbf{第1步（Essen不等式）}：对任意分布函数 \(F\) 和 \(G\)（\(G\) 可微且 \(G'\) 有界），设 \(f, g\) 为其特征函数。则对任意 \(T > 0\)，

\[\sup_x |F(x) - G(x)| \leq \frac{1}{\pi} \int_{-T}^T \left|\frac{f(t) - g(t)}{t}\right| dt + \frac{24}{\pi T} \sup_x |G'(x)|\]

此不等式通过 Fourier 逆变换与截断近似得到。

\textbf{第2步（特征函数的Taylor展开）}：令 \(f_n\) 为 \(Z_n\) 的特征函数。由独立性，\(f_n(t) = (\varphi(\frac{t}{\sigma\sqrt{n}}))^n\)，其中 \(\varphi\) 是 \(X_1 - \mu\) 的特征函数。展开 \(\varphi(s) = 1 - \frac{\sigma^2 s^2}{2} + \theta \frac{\rho_3 |s|^3}{6}\)（\(|\theta| \leq 1\)，由 Taylor 余项的积分形式和 \(|\mathbb{E}[iX]|\) 的矩控制）。则

\[f_n(t) = \left(1 - \frac{t^2}{2n} + \theta \frac{\rho_3 |t|^3}{6 \sigma^3 n^{3/2}}\right)^n = e^{-t^2/2} \left(1 + \theta' \frac{\rho_3 |t|^3}{6 \sigma^3 \sqrt{n}}\right)\]

（展开至 \(n^{-1/2}\) 精度）。

\textbf{第3步（积分估计）}：取 \(T = a \sqrt{n}\)（\(a\) 适当选取以优化界）。在 Essén 不等式中代入 \(f_n(t) - e^{-t^2/2}\) 的展开式，分段估计积分。对 \(|t| \leq T\)，\(|f_n(t) - e^{-t^2/2}| \leq C_1 \frac{\rho_3 |t|^3}{\sigma^3 \sqrt{n}} e^{-t^2/4}\)（利用 \(|1 - z|^n \leq e^{-n \operatorname{Re}(z)} \cdot (1 + \text{余项})\) 的标准估计）。代入积分得

\[\int_{-T}^T \frac{|f_n(t) - e^{-t^2/2}|}{|t|} dt \leq \frac{C_2 \rho_3}{\sigma^3 \sqrt{n}}\]

\textbf{第4步（优化常数）}：结合 Essén 不等式中的 \(24/(\pi T)\) 项（取 \(T \propto \sqrt{n}\) 得 \(O(n^{-1/2})\) 贡献），总界为 \(O(n^{-1/2})\)。最优常数 \(C\) 通过精心选择截断参数 \(T\) 和更精细的特征函数估计（利用 Lyapunov 分数矩和 Fourier 分析中的极值问题）得到 \(C \leq 0.4748\)。\(\blacksquare\)

\subsection{Cramér大偏差定理}\label{cramuxe9rux5927ux504fux5deeux5b9aux7406}

\textbf{定理271.3（Cramér大偏差定理）} 设 \(X_1, X_2, \ldots\) 是独立同分布随机变量，\(\mathbb{E}[X_1] = \mu\)，且矩母函数 \(M(t) = \mathbb{E}[e^{tX_1}]\) 在 \(t = 0\) 的某个邻域内有限。令 \(S_n = \sum_{i=1}^n X_i\)。则对任意 \(a > \mu\)，

\[\lim_{n \to \infty} \frac{1}{n} \log \mathbb{P}(S_n \geq n a) = -I(a)\]

其中 \(I(a) = \sup_{t \in \mathbb{R}} (t a - \log M(t))\) 是速率函数（Cramér变换或Legendre-Fenchel变换）。\(I(a)\) 是凸函数，满足 \(I(\mu) = 0\)，在 \(a > \mu\) 处严格递增。

\textbf{证明}（Cramér，1938）

\textbf{第1步（上界------Chernoff界）}：由 Markov 不等式，对任意 \(t > 0\)，

\[\mathbb{P}(S_n \geq n a) = \mathbb{P}(e^{t S_n} \geq e^{t n a}) \leq e^{-t n a} \mathbb{E}[e^{t S_n}] = e^{-t n a} M(t)^n = e^{-n(t a - \log M(t))}\]

关于 \(t > 0\) 取最小值，\(\mathbb{P}(S_n \geq n a) \leq e^{-n I(a)}\)。故 \(\limsup_n \frac{1}{n} \log \mathbb{P}(S_n \geq n a) \leq -I(a)\)。

\textbf{第2步（下界------指数倾斜）}：由 \(M(t)\) 在 \(0\) 邻域内有限，速率函数 \(I(a)\) 是闭凸函数。对给定的 \(a\)，存在唯一的 \(t^* = t^*(a) > 0\) 使得 \(I(a) = t^* a - \log M(t^*)\)，且由 Fenchel 对偶，\(M'(t^*)/M(t^*) = a\)。定义指数倾斜测度：

\[d\mathbb{Q}_n = e^{t^* S_n - n \log M(t^*)} d\mathbb{P}\]

则 \(\mathbb{E}_{\mathbb{Q}_n}[S_n/n] = M'(t^*)/M(t^*) = a\)。在 \(\mathbb{Q}_n\) 下，\(S_n/n\) 满足中心极限定理（均值为 \(a\)，方差有限）。利用局部的正态近似：

\[\mathbb{P}(S_n \geq n a) \geq \mathbb{P}(n a \leq S_n \leq n a + \delta \sqrt{n}) = e^{-n I(a)} \mathbb{E}_{\mathbb{Q}_n}[e^{-t^* (S_n - n a)} \mathbf{1}_{0 \leq S_n - n a \leq \delta \sqrt{n}}] \geq e^{-n I(a)} \cdot c n^{-1/2}\]

故 \(\liminf_n \frac{1}{n} \log \mathbb{P}(S_n \geq n a) \geq -I(a)\)。结合上下界得极限存在且等于 \(-I(a)\)。\(\blacksquare\)

\textbf{推论271.4（Cramér定理与中心极限定理的关系）} Cramér大偏差定理给出了远离均值的尾概率的对数渐近，而中心极限定理给出了靠近均值处（\(O(1/\sqrt{n})\) 尺度）的分布渐近。两者互补：CLT 的 Berry-Esseen 界在 \(|x| \gg \sqrt{n}\) 时失效（因正态近似的相对误差发散），此时需用大偏差理论。两者在适度偏差（moderate deviation）区间平滑衔接。

\hrulefill

\section{第335章：随机游走与更新理论}\label{ux7b2c335ux7ae0ux968fux673aux6e38ux8d70ux4e0eux66f4ux65b0ux7406ux8bba}

随机游走是最基本的随机过程，其渐近行为------递推与暂留------由空间维数的拓扑性质决定。本章首先证明 Pólya 定理：简单对称随机游走在 \(\mathbb{Z}\) 和 \(\mathbb{Z}^2\) 上是递推的（几乎必然无穷次返回原点），在 \(\mathbb{Z}^d\)（\(d \geq 3\)）上是暂留的（以正概率永不返回原点）。这一优美结果揭示了维数在随机过程中的本质角色。随后转入更新理论------研究随机过程在''更新时刻''（如随机游走的返回时刻或设备的更换时刻）的渐近行为。Blackwell 更新定理和关键更新定理是更新理论的两大核心结果，给出了更新函数在无穷远处的极限行为，在排队论、可靠性理论和保险数学中有广泛应用。

\subsection{Pólya定理：随机游走的递推与暂留}\label{puxf3lyaux5b9aux7406ux968fux673aux6e38ux8d70ux7684ux9012ux63a8ux4e0eux6682ux7559}

\textbf{定理272.1（Pólya定理，1921）} 设 \(\{S_n\}_{n \geq 0}\) 是 \(\mathbb{Z}^d\) 上的简单对称随机游走：\(S_0 = 0\)，\(S_n = \sum_{k=1}^n X_k\)，其中 \(X_k\) 独立均匀分布于 \(\{\pm e_1, \ldots, \pm e_d\}\)（坐标轴上的单位步）。则

（i）当 \(d = 1\) 或 \(d = 2\) 时，随机游走是\textbf{递推}的：\(\mathbb{P}(S_n = 0 \; \text{无穷多次}) = 1\)。

（ii）当 \(d \geq 3\) 时，随机游走是\textbf{暂留}的：\(\mathbb{P}(S_n = 0 \; \text{无穷多次}) = 0\)，且 \(\mathbb{P}(S_n = 0 \; \text{对某} \; n \geq 1) < 1\)。

\textbf{证明} 记 \(u_n = \mathbb{P}(S_n = 0)\) 为 \(n\) 步后返回原点的概率。则期望返回次数为 \(\sum_{n=0}^\infty u_n\)。由 Borel-Cantelli 引理的逆命题（对独立事件）的推广：若事件 \(\{S_n = 0\}\) 满足 \(\sum_n u_n < \infty\)，则最多有限次发生（暂留）；若 \(\sum_n u_n = \infty\) 且满足适当的渐近独立性条件，则无穷多次发生（递推）。

\textbf{第1步（\(u_n\) 的渐近公式）}：由 Fourier 逆变换（或特征函数），

\[u_n = \frac{1}{(2\pi)^d} \int_{[-\pi, \pi]^d} \phi(\theta)^n \, d\theta\]

其中 \(\phi(\theta) = \frac{1}{d} \sum_{j=1}^d \cos \theta_j\) 是单步的特征函数。当 \(n \to \infty\) 时，积分的主要贡献来自 \(\theta \approx 0\)。展开 \(\phi(\theta) = 1 - \frac{|\theta|^2}{2d} + O(|\theta|^4)\)，则

\[u_n \sim \left(\frac{d}{2\pi n}\right)^{d/2} \int_{\mathbb{R}^d} e^{-|\eta|^2/2} \, d\eta = \left(\frac{d}{2\pi n}\right)^{d/2} (2\pi)^{d/2} = \left(\frac{d}{n}\right)^{d/2}\]

更精确地（通过局部中心极限定理）：\(u_n \sim C_d \cdot n^{-d/2}\)。

\textbf{第2步（级数的收敛性）}： - \(d = 1\)：\(u_n \sim C_1 n^{-1/2}\)，\(\sum_n n^{-1/2} = \infty\)。故 \(\sum u_n = \infty\)，递推。 - \(d = 2\)：\(u_n \sim C_2 n^{-1}\)，\(\sum_n n^{-1} = \infty\)。故 \(\sum u_n = \infty\)，递推。 - \(d = 3\)：\(u_n \sim C_3 n^{-3/2}\)，\(\sum_n n^{-3/2} < \infty\)。故 \(\sum u_n < \infty\)，暂留。 - \(d \geq 4\)：\(u_n \sim C_d n^{-d/2}\)，\(d/2 \geq 2\)，\(\sum n^{-d/2} < \infty\)，暂留。

\textbf{第3步（递推的严格论证）}：\(\sum_n u_n = \infty\) 仅是递推性的必要条件。充分性需利用 \(f_n = \mathbb{P}(T_0 = n)\)（首次返回时间概率）与 \(u_n\) 的关系：\(u_n = \sum_{k=1}^n f_k u_{n-k}\)（\(u_0 = 1\)，\(u_1 = \cdots = u_{d\text{ odd}} = 0\) 来自奇偶性抵消）。由生成函数 \(U(z) = \sum_{n \geq 0} u_n z^n\) 和 \(F(z) = \sum_{n \geq 1} f_n z^n\)，有 \(U(z) = \frac{1}{1 - F(z)}\)。\(\sum u_n = \infty \iff \lim_{z \to 1^-} U(z) = \infty \iff F(1) = \sum f_n = 1 \iff\) 返回概率为 \(1\)。由马尔可夫性，每次返回后重新开始，故无穷多次返回的概率为 \(1\)。\(\blacksquare\)

\subsection{Blackwell更新定理}\label{blackwellux66f4ux65b0ux5b9aux7406}

\textbf{定义272.2（更新过程）} 设 \(\{X_n\}_{n \geq 1}\) 是独立同分布非负随机变量（相继到达间隔时间），分布函数为 \(F\)（\(\mathbb{P}(X_n = 0) < 1\) 且 \(\mu = \mathbb{E}[X_1] \in (0, \infty]\)）。更新时刻 \(T_0 = 0\)，\(T_n = \sum_{i=1}^n X_i\)。\textbf{更新函数} \(U(t) = \mathbb{E}[N(t)] = \sum_{n=1}^\infty F^{*n}(t)\)，其中 \(N(t) = \max\{n : T_n \leq t\}\) 是 \([0, t]\) 内的更新次数，\(F^{*n}\) 是 \(n\) 重卷积。

\textbf{定理272.3（Blackwell更新定理）} 若 \(F\) 是非格点分布（即 \(X_n\) 的分布不集中在某格点 \(\{0, h, 2h, \ldots\}\) 上），则对任意 \(h > 0\)，

\[\lim_{t \to \infty} (U(t + h) - U(t)) = \frac{h}{\mu}\]

其中约定若 \(\mu = \infty\) 则 \(h/\mu = 0\)。更新函数在无穷远处以恒定速率 \(1/\mu\) 增长，即更新速率等于平均间隔时间的倒数。

\textbf{证明}（Feller-Orey方法，基于更新方程和推移紧性）

\textbf{第1步（更新方程）}：对任意有界可测函数 \(g\)，解 \(Z(t) = g(t) + \int_0^t Z(t-s) dF(s)\) 由 \(Z(t) = \int_0^t g(t-s) dU(s)\) 给出（以 \(U\) 为核的卷积）。

\textbf{第2步（推移紧性）}：定义函数族 \(V_t(s) = U(t + s) - U(t)\)（\(s \geq 0\)）。由更新函数的次可加性和单调性，\(V_t(s)\) 在 \([0, \infty)\) 上一致有界（若 \(\mu < \infty\)）且等度连续（对适当度量）。由 Arzela-Ascoli，存在子列 \(t_k \to \infty\) 使得 \(V_{t_k}\) 一致收敛于极限 \(V(s)\)。

\textbf{第3步（极限满足的函数方程）}：由更新方程，\(U(t) = 1 + \int_0^t U(t-s) dF(s)\)（\(t \geq 0\)，\(U(0^-) = 0\)）。平移后：\(U(t + h) - U(t) = \int_0^t (U(t + h - s) - U(t - s)) dF(s) + \int_t^{t+h} (1 - F(t + h - s)) dU(s)\)。令 \(t = t_k \to \infty\) 取极限，\(V(h)\) 满足 \(V(h) = \int_0^\infty V(h - s) dF(s)\)（\(h \geq 0\)）。

\textbf{第4步（极限的确定）}：上述积分方程的唯一解（在 \(L^1\) 意义或点态意义下）为线性函数 \(V(h) = c h\)。由泛函分析中的更新定理（或通过 Laplace 变换），\(c\) 由 \(\int_0^\infty h dF(s) = \mu\) 确定：\(c = 1/\mu\)。非格点条件保证了解的唯一性（在格点情形，Blackwell 定理仅对 \(h\) 为格点间距的倍数成立）。\(\blacksquare\)

\subsection{关键更新定理}\label{ux5173ux952eux66f4ux65b0ux5b9aux7406}

\textbf{定理272.4（关键更新定理------Smith，1954）} 设 \(F\) 是非格点分布，\(\mu = \mathbb{E}[X_1] \in (0, \infty)\)。若 \(g: [0, \infty) \to \mathbb{R}\) 是 Riemann 直接可积的（即 \(\sum_k \sup_{[k, k+1)} |g| < \infty\) 且 \(g\) 在紧区间上 Riemann 可积），则

\[\lim_{t \to \infty} \int_0^t g(t - s) \, dU(s) = \frac{1}{\mu} \int_0^\infty g(s) \, ds\]

\textbf{证明} 关键更新定理是 Blackwell 定理的卷积推广。将积分写为 Riemann-Stieltjes 和并利用 Blackwell 定理的分段逼近。

\textbf{第1步（阶梯函数逼近）}：首先对阶梯函数 \(g = \mathbf{1}_{[0, h]}\)，结论化为 Blackwell 定理。对一般 Riemann 直接可积的 \(g\)，用阶梯函数的有限线性组合逼近。

\textbf{第2步（截断和尾估计）}：将积分分解为 \(\int_0^A + \int_A^\infty\)。对 \(A\) 充分大，第二项由 \(g\) 的 Riemann 直接可积性可任意小。对第一项，用 Blackwell 定理和阶梯函数逼近得极限。\(\blacksquare\)

\textbf{推论272.5（更新定理的应用）} 在排队论中，\(g\) 可取为服务时间的尾分布，关键更新定理给出稳态下系统繁忙概率的表达式。在保险数学中，\(g\) 为索赔额的分布，更新定理给出终极破产概率的 Cramér-Lundberg 近似。在寿命试验中，更新过程描述了可修系统的故障-维修循环，更新函数 \(U(t)\) 的渐近性给出平均故障次数的线性增长。

\hrulefill

\hrulefill

\hrulefill

\hrulefill

\hrulefill

\newpage

\chapter{10-概率与统计/01-概率论/02-鞅与Markov链.md}\label{ux6982ux7387ux4e0eux7edfux8ba101-ux6982ux7387ux8bba02-ux9785ux4e0emarkovux94fe.md}

\section{第336章：条件期望与鞅}\label{ux7b2c336ux7ae0ux6761ux4ef6ux671fux671bux4e0eux9785}

条件期望是概率论中最重要的概念之一，它将 ``已知部分信息下的平均'' 精确化。鞅是 Doob 引入的强大工具，在随机过程、金融数学和统计推断中有广泛应用。

\subsection{106.1 条件期望}\label{ux6761ux4ef6ux671fux671b}

\textbf{定义 106.1}（条件期望）：设 \(X\) 是可积随机变量，\(\mathcal{G} \subseteq \mathcal{F}\) 是子 \(\sigma\)-代数。\(X\) 在 \(\mathcal{G}\) 下的\textbf{条件期望} \(\mathbb{E}[X \mid \mathcal{G}]\) 是唯一的 \(\mathcal{G}\)-可测随机变量 \(Y\)，满足

\[\int_A Y d\mathbb{P} = \int_A X d\mathbb{P} \quad \text{对所有 } A \in \mathcal{G}\]

（由 Radon-Nikodym 定理保证存在唯一性）。

\textbf{定义 106.2}（给定随机变量的条件期望）：\(\mathbb{E}[X \mid Y] = \mathbb{E}[X \mid \sigma(Y)]\)。存在 Borel 函数 \(g\) 使得 \(\mathbb{E}[X \mid Y] = g(Y)\)，记 \(g(y) = \mathbb{E}[X \mid Y = y]\)。

\textbf{命题 106.1}（条件期望的基本性质）： 1. \(\mathbb{E}[\mathbb{E}[X \mid \mathcal{G}]] = \mathbb{E}[X]\)（全期望律） 2. 若 \(X\) 是 \(\mathcal{G}\)-可测，则 \(\mathbb{E}[X \mid \mathcal{G}] = X\) 3. 若 \(X\) 与 \(\mathcal{G}\) 独立（即 \(\sigma(X)\) 与 \(\mathcal{G}\) 独立），则 \(\mathbb{E}[X \mid \mathcal{G}] = \mathbb{E}[X]\) 4. \(\mathbb{E}[aX + bY \mid \mathcal{G}] = a \mathbb{E}[X \mid \mathcal{G}] + b \mathbb{E}[Y \mid \mathcal{G}]\)（线性） 5. 若 \(Z\) 是 \(\mathcal{G}\)-可测且有界，则 \(\mathbb{E}[ZX \mid \mathcal{G}] = Z \mathbb{E}[X \mid \mathcal{G}]\) 6. \textbf{塔性质}：\(\mathbb{E}[\mathbb{E}[X \mid \mathcal{G}] \mid \mathcal{H}] = \mathbb{E}[X \mid \mathcal{H}]\)（当 \(\mathcal{H} \subseteq \mathcal{G}\)） 7. 条件 Jensen 不等式：\(\varphi(\mathbb{E}[X \mid \mathcal{G}]) \leq \mathbb{E}[\varphi(X) \mid \mathcal{G}]\)（\(\varphi\) 凸）

\subsection{106.2 鞅的定义与基本性质}\label{ux9785ux7684ux5b9aux4e49ux4e0eux57faux672cux6027ux8d28}

\textbf{定义 106.3}（滤过）：\textbf{滤过} \(\{\mathcal{F}_n\}_{n \geq 0}\) 是 \(\mathcal{F}\) 的递增子 \(\sigma\)-代数序列：\(\mathcal{F}_0 \subseteq \mathcal{F}_1 \subseteq \cdots \subseteq \mathcal{F}\)。\(\mathcal{F}_n\) 表示时刻 \(n\) 可用的信息。

\textbf{定义 106.4}（鞅）：随机过程 \(\{M_n\}_{n \geq 0}\) 称为关于 \(\{\mathcal{F}_n\}\) 的\textbf{鞅}（martingale），如果 1. \(M_n\) 是 \(\mathcal{F}_n\)-可测（适应于滤过） 2. \(\mathbb{E}[|M_n|] < \infty\)（对所有 \(n\)） 3. \(\mathbb{E}[M_{n+1} \mid \mathcal{F}_n] = M_n\)（对所有 \(n\)）

若 \(\mathbb{E}[M_{n+1} \mid \mathcal{F}_n] \geq M_n\)（下鞅，submartingale）或 \(\leq M_n\)（上鞅，supermartingale），则称为\textbf{下鞅}或\textbf{上鞅}。

\subsection{106.3 鞅收敛定理}\label{ux9785ux6536ux655bux5b9aux7406}

\textbf{定理 106.2}（Doob 上穿不等式）：设 \(\{X_n\}\) 是下鞅。对 \(a < b\)，在 \(n\) 步内 \(X_n\) 从 \(a\) 以下穿到 \(b\) 以上的次数 \(U_n(a, b)\) 满足

\[\mathbb{E}[U_n(a, b)] \leq \frac{\mathbb{E}[(X_n - a)^+]}{b - a}\]

Doob 上穿不等式是鞅理论中最为深刻的技术工具之一。其直观含义是：下鞅在有限区间 \((a,b)\) 上的穿越次数在期望的意义下受终端值控制，因此若鞅不收敛必产生无穷次穿越，与这一估计矛盾。Doob 正是利用这一简单的组合推理，得到了鞅几乎必然收敛这一强大结论。

\textbf{证明}：令 \(S_1 = \inf\{k \geq 0 : X_k \leq a\}\)，\(T_1 = \inf\{k > S_1 : X_k \geq b\}\)，交替定义停时 \(S_i, T_i\)。则上穿次数 \(U_n = \max\{j : T_j \leq n\}\)。定义指标 \(I_k = 1\) 当 \(k \in [S_i+1, T_i]\) 对某 \(i\)，否则 \(0\)。则 \(I_k\) 可料（\(\mathcal{F}_{k-1}\)-可测）。计算 \(\sum_{k=1}^n I_k (X_k - X_{k-1}) \geq (b-a)U_n - (X_n - a)^-\)。由下鞅性，\(\mathbb{E}[I_k(X_k - X_{k-1})] = \mathbb{E}[I_k \mathbb{E}[X_k - X_{k-1} \mid \mathcal{F}_{k-1}]] \geq 0\)。故 \((b-a)\mathbb{E}[U_n] \leq \mathbb{E}[(X_n - a)^+] - \mathbb{E}[(X_0 - a)^+] \leq \mathbb{E}[(X_n - a)^+]\)。\(\blacksquare\)

上穿不等式的一个重要推论是：\(L^1\) 有界的下鞅必定几乎必然收敛。事实上，若 \(\sup_n \mathbb{E}[|M_n|] < \infty\)，则对任意有理数 \(a < b\)，上穿次数的期望有界，从而几乎必然有限，排除了 \(\liminf < a < b < \limsup\) 的可能性。这一论证构成了鞅收敛定理的证明核心。

\textbf{定理 106.3}（鞅收敛定理，Doob）：设 \(\{M_n\}\) 是下鞅且 \(\sup_n \mathbb{E}[|M_n|] < \infty\)（或 \(\sup_n \mathbb{E}[M_n^+] < \infty\)）。则 \(M_n\) 几乎必然收敛到某个可积随机变量 \(M_\infty\)。

\textbf{证明}：使用上穿不等式。若 \(\{M_n\}\) 不收敛，则存在有理数 \(a < b\) 使得 \(\liminf M_n < a < b < \limsup M_n\)，即上穿次数 \(U_\infty(a,b) = \infty\)。由上穿不等式（定理 106.2），\(\mathbb{E}[U_n(a,b)] \leq \frac{\mathbb{E}[(M_n - a)^+]}{b - a} \leq \frac{\sup_n \mathbb{E}[|M_n|] + |a|}{b - a} < \infty\)。由单调收敛定理，\(\mathbb{E}[U_\infty(a,b)] = \lim_{n \to \infty} \mathbb{E}[U_n(a,b)] < \infty\)，故 \(U_\infty(a,b) < \infty\) a.s.，与 \(\liminf < a < b < \limsup\) 矛盾。因此 \(\lim M_n\) 存在（a.s.）。由 Fatou 引理，\(\mathbb{E}[|M_\infty|] = \mathbb{E}[\liminf |M_n|] \leq \sup_n \mathbb{E}[|M_n|] < \infty\)，故 \(M_\infty\) 可积。\(\blacksquare\)

定理 106.3 仅保证了几乎必然收敛；若鞅还满足 \(L^p\) 有界性（\(p > 1\)），则更强地在 \(L^p\) 范数下收敛。这是因为 Doob 极大不等式提供了控制函数 \(\sup_n |M_n|\) 的 \(L^p\) 可积性，从而允许通过 Lebesgue 控制收敛定理将逐点收敛加强为 \(L^p\) 收敛。这一结果在随机积分理论和偏微分方程的概率解法中至关重要。

\textbf{定理 106.4}（\(L^p\) 鞅收敛定理）：若 \(\{M_n\}\) 是鞅且 \(\sup_n \mathbb{E}[|M_n|^p] < \infty\)（\(p > 1\)），则 \(M_n \to M_\infty\) 在 \(L^p\) 中收敛。

\textbf{证明}：由 Doob 极大不等式，\(\mathbb{E}[\sup_n |M_n|^p] \leq (\frac{p}{p-1})^p \sup_n \mathbb{E}[|M_n|^p] < \infty\)。故 \(\sup_n |M_n| \in L^p\)。由定理 106.3，\(M_n \to M_\infty\) a.s.。由于 \(|M_n - M_\infty|^p \leq 2^p (\sup_n |M_n|)^p \in L^1\)，由控制收敛定理，\(\mathbb{E}[|M_n - M_\infty|^p] \to 0\)。\(\blacksquare\)

\subsection{106.4 可选停止定理}\label{ux53efux9009ux505cux6b62ux5b9aux7406}

\textbf{定义 106.5}（停时）：随机变量 \(\tau: \Omega \to \mathbb{N} \cup \{\infty\}\) 称为关于 \(\{\mathcal{F}_n\}\) 的\textbf{停时}，如果 \(\{\tau = n\} \in \mathcal{F}_n\)（对所有 \(n\)）。即''是否停止''仅依赖于当前和过去的信息。

\textbf{定理 106.5}（可选停止定理，Doob）：设 \(\{M_n\}\) 是鞅，\(\tau\) 是停时。若以下条件之一成立： 1. \(\tau\) 有界（\(\tau \leq N\) 几乎必然） 2. \(\{M_n\}\) 有界且 \(\tau < \infty\) 几乎必然 3. \(\mathbb{E}[\tau] < \infty\) 且 \(|M_n - M_{n-1}| \leq C\)（有界增量）

则 \(\mathbb{E}[M_\tau] = \mathbb{E}[M_0]\)。

\emph{证明}：对情形 1，\(M_{\tau \wedge N} = \sum_{n=0}^{N-1} M_n \mathbf{1}_{\{\tau = n\}} + M_N \mathbf{1}_{\{\tau \geq N\}}\) 是鞅，故 \(\mathbb{E}[M_\tau] = \mathbb{E}[M_{\tau \wedge N}] = \mathbb{E}[M_0]\)。∎

\textbf{定理 106.6}（Wald 等式）：设 \(X_1, X_2, \ldots\) 是 i.i.d.，\(\mathbb{E}[|X_1|] < \infty\)，\(\mu = \mathbb{E}[X_1]\)。\(\tau\) 是关于 \(\{X_n\}\) 的停时，\(\mathbb{E}[\tau] < \infty\)。则 \(\mathbb{E}[\sum_{i=1}^\tau X_i] = \mu \mathbb{E}[\tau]\)。

\textbf{证明}：令 \(S_n = \sum_{i=1}^n (X_i - \mu)\)。\(S_n\) 是鞅（\(\mathbb{E}[S_{n+1} \mid \mathcal{F}_n] = S_n\)）。由可选停止定理（条件 3），若 \(\mathbb{E}[\tau] < \infty\) 且 \(\mathbb{E}[|X_1|] < \infty\)（增量一致有界于可积变量），则 \(\mathbb{E}[S_\tau] = \mathbb{E}[S_0] = 0\)。故 \(\mathbb{E}[\sum_{i=1}^\tau X_i] = \mu \mathbb{E}[\tau] + \mathbb{E}[S_\tau] = \mu \mathbb{E}[\tau]\)。\(\blacksquare\)

\hrulefill

\hrulefill

\hrulefill

\hrulefill

\hrulefill

\hrulefill

\section{第337章：Markov 链}\label{ux7b2c337ux7ae0markov-ux94fe}

Markov 链是''无记忆''的随机过程，其未来状态仅依赖于当前状态。本章讨论离散时间、可数状态空间的 Markov 链，包括转移概率、平稳分布和遍历定理。

\subsection{107.1 Markov 链的定义}\label{markov-ux94feux7684ux5b9aux4e49}

\textbf{定义 107.1}（Markov 链）：状态空间 \(S\)（可数集）上的随机过程 \(\{X_n\}_{n \geq 0}\) 称为 \textbf{Markov 链}，如果满足 Markov 性：

\[\mathbb{P}(X_{n+1} = j \mid X_n = i, X_{n-1}, \ldots, X_0) = \mathbb{P}(X_{n+1} = j \mid X_n = i)\]

上述定义中的转移概率原则上可以依赖于当前时刻 \(n\)。但在实际应用中，绝大多数 Markov 链具有时间齐次性------即转移机制不随时间变化。时齐 Markov 链的转移行为完全由一个与时间无关的转移矩阵 \(P\) 刻画，这极大地简化了其长期行为的分析。

\textbf{定义 107.2}（时齐 Markov 链）：若 \(\mathbb{P}(X_{n+1} = j \mid X_n = i) = p_{ij}\)（与 \(n\) 无关），则称链是\textbf{时齐}的。矩阵 \(P = (p_{ij})_{i,j \in S}\) 称为\textbf{转移矩阵}（每行和为 \(1\)）。

转移矩阵 \(P\) 描述了一步转移的行为，而多步转移概率则通过矩阵幂自然得到。Chapman-Kolmogorov 方程将 \(n\) 步和 \(m\) 步的转移概率连接起来，是 Markov 链所有渐近分析的基础递推关系。

\textbf{定义 107.3}（\(n\) 步转移概率）：\(p_{ij}^{(n)} = \mathbb{P}(X_n = j \mid X_0 = i) = (P^n)_{ij}\)。由 Chapman-Kolmogorov 方程：\(p_{ij}^{(n+m)} = \sum_{k \in S} p_{ik}^{(n)} p_{kj}^{(m)}\)。

初始分布 \(\pi_0\) 和转移矩阵 \(P\) 共同决定了 Markov 链的联合分布。由于 Markov 性将每一步的条件分布简化为仅依赖于前一步，\(n\) 时刻的分布 \(\pi_n\) 可以简洁地表示为初始分布与转移矩阵幂的乘积。

\textbf{命题 107.1}：Markov 链的分布由初始分布 \(\pi_0\)（\(\pi_0(i) = \mathbb{P}(X_0 = i)\)）和转移矩阵 \(P\) 完全确定：\(\pi_n = \pi_0 P^n\)。

\subsection{107.2 状态的分类}\label{ux72b6ux6001ux7684ux5206ux7c7b}

\textbf{定义 107.4}（可达与互通）：\(j\) 从 \(i\) \textbf{可达}（\(i \to j\)）如果存在 \(n \geq 0\) 使得 \(p_{ij}^{(n)} > 0\)。\(i\) 和 \(j\) 互通（\(i \leftrightarrow j\)）如果 \(i \to j\) 且 \(j \to i\)。互通是等价关系，等价类称为\textbf{互通类}。

互通关系将状态空间划分为不相交的等价类，这是 Markov 链分类的起点。然而，属于同一互通类的状态在长期行为上可能仍有本质差异，区分这种差异需要引入常返性的概念：一个状态\(i\)是常返的，意味着从\(i\)出发几乎必然会在有限时间内返回\(i\)；否则称为瞬时的。

\textbf{定义 107.5}（常返与瞬时）：状态 \(i\) 称为\textbf{常返}的，如果 \(\mathbb{P}(X_n = i \text{ 无穷多次} \mid X_0 = i) = 1\)。否则称为\textbf{瞬时}的。

常返性的上述定义涉及无限多次返回的概率，难以直接计算。幸运的是，这一条件可以等价地转化为转移概率的级数发散性：状态\(i\)常返当且仅当\(\sum_{n=1}^\infty p_{ii}^{(n)} = \infty\)。这一判别准则的证明利用了首次返回时间的生成函数与转移概率生成函数之间的基本关系。

\textbf{定理 107.2}（常返性的判别）：\(i\) 常返当且仅当 \(\sum_{n=1}^\infty p_{ii}^{(n)} = \infty\)。\(i\) 瞬时当且仅当 \(\sum_{n=1}^\infty p_{ii}^{(n)} < \infty\)。

\textbf{证明}：令 \(T_i = \inf\{n \geq 1 : X_n = i\}\)。由 母函数 关系 \(\sum_{n=0}^\infty p_{ii}^{(n)} s^n = 1/(1 - F_{ii}(s))\)，其中 \(F_{ii}(s) = \mathbb{E}[s^{T_i} \mid X_0 = i]\)。\(i\) 常返当且仅当 \(\mathbb{P}(T_i < \infty \mid X_0 = i) = 1\)，即 \(F_{ii}(1) = 1\)。则 \(\sum_n p_{ii}^{(n)} = 1/(1 - F_{ii}(1)) = \infty\)。若 \(F_{ii}(1) < 1\)，则 \(\sum_n p_{ii}^{(n)} = 1/(1 - F_{ii}(1)) < \infty\)。\(\blacksquare\)

常返状态还可以进一步细分为正常返和零常返两类，取决于首次返回时间的期望是否有限。对于正常返状态，链不仅几乎必然返回，而且平均返回时间是有限的，这一性质对于平稳分布的存在性至关重要。

\textbf{定义 107.6}（正常返与零常返）：常返状态 \(i\) 称为\textbf{正常返}的，如果 \(\mathbb{E}[T_i \mid X_0 = i] < \infty\)（\(T_i = \inf\{n \geq 1 : X_n = i\}\) 是首次返回时间）。否则称为\textbf{零常返}。

上述各种状态分类------瞬时、零常返、正常返------在不可约链中具有一致性：同一等价类内所有状态的类型相同。这是 Markov 链理论的一个关键结构事实，称为''固态性''（solidarity），其证明基于互通关系对返回概率级数收敛的影响。

\textbf{定义 107.7}（不可约）：Markov 链称为\textbf{不可约}的，如果所有状态互通（只有一个互通类）。

\textbf{定理 107.3}：对不可约 Markov 链，所有状态同时为瞬时、零常返或正常返（固态性）。

\textbf{证明}：设 \(i\) 正常返且 \(j \leftrightarrow i\)。存在 \(n, m\) 使 \(p_{ij}^{(n)} > 0\)，\(p_{ji}^{(m)} > 0\)。则 \(p_{jj}^{(n+r+m)} \geq p_{ji}^{(m)} p_{ii}^{(r)} p_{ij}^{(n)}\)。由 \(i\) 正常返，\(\sum_r p_{ii}^{(r)}\) 以 \(\mathbb{E}[T_i] < \infty\) 的速率收敛，推出 \(\sum_k p_{jj}^{(k)} = \infty\) 且 \(\mathbb{E}[T_j] < \infty\)。类似论证关联瞬时性和零常返性。因为不可约链所有状态互通，该关系通过链任意两个状态之间传递。实际上，定义不依赖于特定状态的固化性等价于：存在常数 \(c_{ij} > 0\) 和 \(N\) 使得对所有充分大的 \(r\)，\(p_{jj}^{(r)} \geq c_{ij} p_{ii}^{(r-N)}\)。因此级数 \(\sum p_{ii}^{(n)}\) 的收敛/发散性与 \(\sum p_{jj}^{(n)}\) 相同。\(\blacksquare\)

\subsection{107.3 平稳分布}\label{ux5e73ux7a33ux5206ux5e03}

\textbf{定义 107.8}（平稳分布）：概率分布 \(\pi = (\pi(i))_{i \in S}\) 称为\textbf{平稳分布}（或不变测度），如果 \(\pi = \pi P\)（即 \(\pi(j) = \sum_{i \in S} \pi(i) p_{ij}\)）。若 \(\pi_0 = \pi\)，则 \(\pi_n = \pi\) 对所有 \(n\)。

\textbf{定理 107.4}（平稳分布的存在性）：不可约 Markov 链有平稳分布当且仅当它是正常返的。此时平稳分布唯一，且 \(\pi(i) = 1 / \mathbb{E}[T_i \mid X_0 = i]\)。

\textbf{证明}：设链不可约正常返。定义 \(\pi(i) = 1 / \mathbb{E}[T_i \mid X_0 = i]\)。由 Markov 更新定理，\(\lim_{n \to \infty} \frac{1}{n} \sum_{k=0}^{n-1} p_{ji}^{(k)} = \pi(i)\) 对任意 \(j\) 成立。对稳态方程两边取极限：\(\pi(j) = \lim_{n} \frac{1}{n} \sum_{k=0}^{n-1} p_{ji}^{(k+1)} = \lim_n \frac{1}{n} \sum_{k} \sum_\ell p_{j\ell} p_{\ell i}^{(k)} = \sum_\ell p_{j\ell} \pi(i)\)。等号两侧交换 \(j\) 和 \(i\) 即得 \(\pi = \pi P\)。唯一性：若 \(\pi'\) 是另一平稳分布，由遍历定理，\(\pi'(i) = \lim \frac{1}{n} \sum_k \pi^{(k)}(i) = \pi(i)\)。\(\blacksquare\)

\textbf{定理 107.5}（Markov 链遍历定理）：设不可约正常返 Markov 链有平稳分布 \(\pi\)。则

\[\lim_{n \to \infty} \frac{1}{n} \sum_{k=0}^{n-1} \mathbf{1}_{\{X_k = j\}} = \pi(j) \quad \text{几乎必然}\]

即时间平均收敛于空间平均。若链是非周期的，则 \(\lim_{n \to \infty} p_{ij}^{(n)} = \pi(j)\)（分布收敛）。

\textbf{证明}：由独立于初始态的 Markov 更新定理，\(\frac{1}{n} \sum_{k=0}^{n-1} \mathbf{1}_{\{X_k = j\}} \to \pi(j)\) a.s.。此为时间平均向空间平均的收敛。若链非周期，由更新定理的 Blackwell 形式，\(p_{ij}^{(n)} \to 1 / \mathbb{E}[T_j \mid X_0 = j] = \pi(j)\)。若链有周期 \(d > 1\)，则 \(p_{ii}^{(nd)} \to d \pi(i)\) 而 \(p_{ii}^{(nd+r)} \to 0\)（\(0 < r < d\)）。\(\blacksquare\)

\textbf{定义 107.9}（周期）：状态 \(i\) 的\textbf{周期} \(d(i) = \gcd\{n \geq 1 : p_{ii}^{(n)} > 0\}\)。若 \(d(i) = 1\)，则 \(i\) 非周期。不可约链中所有状态周期相同。

\subsection{107.4 可逆性与细致平衡}\label{ux53efux9006ux6027ux4e0eux7ec6ux81f4ux5e73ux8861}

\textbf{定义 107.10}（细致平衡）：概率分布 \(\pi\) 满足\textbf{细致平衡条件}（detailed balance），如果 \(\pi(i) p_{ij} = \pi(j) p_{ji}\)（对所有 \(i, j\)）。此时 \(\pi\) 是平稳分布，链称为\textbf{可逆}的。

细致平衡条件比平稳分布方程 \(\pi = \pi P\) 更强：它不仅要求全局概率流守恒，还要求每对状态之间的概率流量相等。这一更强的条件为构造以指定分布为平稳分布的 Markov 链提供了一个实用的设计原则，即 Metropolis-Hastings 算法和 Gibbs 抽样的理论基础。

\textbf{定理 107.6}（MCMC 的基本原理）：若目标分布 \(\pi\) 难以直接抽样，可构造 Markov 链使其平稳分布为 \(\pi\)。Metropolis-Hastings 算法和 Gibbs 抽样是构造满足细致平衡的转移矩阵的标准方法。

\textbf{证明}：给定目标分布 \(\pi\)，构造转移核 \(P\) 满足细致平衡 \(\pi(i) p_{ij} = \pi(j) p_{ji}\)。Metropolis-Hastings 算法：选择建议分布 \(Q(j \mid i)\)，定义接受概率 \(\alpha_{ij} = \min\{1, \frac{\pi(j) Q(i \mid j)}{\pi(i) Q(j \mid i)}\}\)。转移概率 \(p_{ij} = Q(j \mid i) \alpha_{ij}\)（\(i \neq j\)）和 \(p_{ii} = 1 - \sum_{j \neq i} p_{ij}\)。容易验证 \(\pi(i) p_{ij} = \pi(j) p_{ji}\)，故 \(\pi\) 为平稳分布。由遍历定理，链的样本路径提供 \(\pi\) 的渐近抽样。\(\blacksquare\)

\hrulefill

\hrulefill

\hrulefill

\hrulefill

\hrulefill

\hrulefill

\section{第338章：Brown 运动与随机过程}\label{ux7b2c338ux7ae0brown-ux8fd0ux52a8ux4e0eux968fux673aux8fc7ux7a0b}

Brown 运动（Wiener 过程）是最重要的连续时间随机过程，是随机分析的基础。本章给出 Brown 运动的构造、基本性质，并简要介绍连续时间鞅和 Itô 积分。

\subsection{108.1 Brown 运动的定义与构造}\label{brown-ux8fd0ux52a8ux7684ux5b9aux4e49ux4e0eux6784ux9020}

\textbf{定义 108.1}（Brown 运动 / Wiener 过程）：随机过程 \(\{B_t\}_{t \geq 0}\) 称为\textbf{标准 Brown 运动}（或 Wiener 过程），如果 1. \(B_0 = 0\) 几乎必然 2. 独立增量：对 \(0 \leq t_1 < t_2 < \cdots < t_n\)，\(B_{t_2} - B_{t_1}, \ldots, B_{t_n} - B_{t_{n-1}}\) 独立 3. 平稳增量：\(B_t - B_s \sim \mathcal{N}(0, t - s)\)（\(0 \leq s < t\)） 4. 样本路径 \(t \mapsto B_t\) 几乎必然连续

\textbf{定理 108.1}（Brown 运动的存在性，Wiener 1923）：Brown 运动存在。可通过以下方式构造： - 在 \(C[0, \infty)\) 上定义 Wiener 测度（Kolmogorov 扩张定理） - 或使用 Lévy-Ciesielski 构造（Schauder 基上的随机 Fourier 级数）

\textbf{证明}（Levy-Ciesielski）：在 \(C[0,1]\) 上定义 Schauder 基。令 \(B_t^{(N)} = \sum_{n=0}^N \sum_k Z_{n,k} \int_0^t h_{n,k}(s)ds\)，\(Z_{n,k} \sim \mathcal{N}(0,1)\) i.i.d.。由 Borel-Cantelli 引理，\(\sup_t |B_t^{(N)} - B_t^{(M)}| \to 0\) a.s.，极限 \(B_t\) 连续。由独立正态增量性质，\(B_t - B_s \sim \mathcal{N}(0, t-s)\) 且独立于过去。\(\blacksquare\)

\textbf{命题 108.1}（Brown 运动的性质）： 1. 尺度不变性：\(\{c^{-1/2} B_{ct}\}\) 也是 Brown 运动 2. 时间反演：\(\{t B_{1/t}\}\)（\(t > 0\)，\(B_0 = 0\)）是 Brown 运动 3. 反射原理：\(\mathbb{P}(\max_{0 \leq s \leq t} B_s \geq a) = 2\mathbb{P}(B_t \geq a)\)（\(a > 0\)） 4. 样本路径几乎必然无处可微（Paley-Wiener-Zygmund 定理） 5. 样本路径在 \([0, t]\) 上的变差无穷大（几乎必然），但二次变差为 \(t\)

\subsection{108.2 Brown 运动的 Markov 性与强 Markov 性}\label{brown-ux8fd0ux52a8ux7684-markov-ux6027ux4e0eux5f3a-markov-ux6027}

Brown 运动的 Markov 性意味着：过程在任何固定时刻 \(s\) 之后的未来行为与过去独立，且重新以 \(B_s\) 为起点''重新开始''。但 Brown 运动还满足一个更强的性质：这一''重新开始''的性质不仅对固定时刻成立，对随机时刻------只要该时刻不依赖于未来信息（即停时）------同样成立。

\textbf{定理 108.2}（Markov 性）：Brown 运动是 Markov 过程：对 \(s > 0\)，\(\{B_{t+s} - B_s\}_{t \geq 0}\) 是 Brown 运动且与 \(\sigma(B_u, u \leq s)\) 独立。

\textbf{证明}：固定 \(s > 0\)。对任意 \(n\) 及 \(0 \leq t_1 < \cdots < t_n\)，增量向量 \((B_{s+t_1} - B_s, \ldots, B_{s+t_n} - B_{s+t_{n-1}})\) 是独立正态变量且分布与 \((B_{t_1}, B_{t_2} - B_{t_1}, \ldots, B_{t_n} - B_{t_{n-1}})\) 相同。这些增量与 \(\sigma(B_u, u \leq s)\) 独立，因为 Brown 运动的增量独立于过去。故 \(W_t := B_{s+t} - B_s\) 是 Brown 运动且独立于 \(\mathcal{F}_s\)。\(\blacksquare\)

固定时刻的 Markov 性在多数应用场景中已足够强大，但在分析首次到达时间、反射原理等涉及随机退出时间的问题时，需要将 Markov 性推广到停时。强 Markov 性是 Brown 运动与一般连续参数 Markov 过程相比最本质的优势之一，其证明需借助停时的离散逼近和样本路径的右连续性。

\textbf{定理 108.3}（强 Markov 性）：设 \(\tau\) 是关于 Brown 运动自然滤过的停时（\(\tau < \infty\) 几乎必然）。则 \(\{B_{\tau + t} - B_\tau\}_{t \geq 0}\) 是 Brown 运动且与 \(\mathcal{F}_\tau\) 独立。

\textbf{证明}：对停时 \(\tau\) 取值于 \(\{k 2^{-n}\}\) 的离散近似 \(\tau_n = \inf\{k 2^{-n} : k 2^{-n} \geq \tau\}\)。对每个固定 \(k 2^{-n}\)，由 Markov 性，\(B_{k 2^{-n} + t} - B_{k 2^{-n}}\) 与 \(\mathcal{F}_{k 2^{-n}}\) 独立。当 \(n \to \infty\)，\(\tau_n \downarrow \tau\)。利用 \(B_t\) 的右连续性，\(B_{\tau_n + t} - B_{\tau_n} \to B_{\tau + t} - B_\tau\) 几乎必然。由停时定义，\(\mathcal{F}_{\tau_n} \supseteq \mathcal{F}_\tau\)，故独立性在极限下保持。有限维分布的极限为 \(B_{\tau + t} - B_\tau\) 的相应分布。\(\blacksquare\)

强 Markov 性的一个精彩应用是 Blumenthal 0-1 律：Brown 运动在任意小的正时间区间内的行为完全由瞬时滤过 \(\mathcal{F}_{0^+}\) 决定，而该 \(\sigma\)-代数的任何事件概率非 0 即 1。这意味着 Brown 运动在原点附近以概率 1 立即取正值也立即取负值，其路径在任意小的时间尺度上都表现出极端的振荡行为。

\textbf{定理 108.4}（Blumenthal 0-1 律）：设 \(\mathcal{F}_{0^+} = \bigcap_{t > 0} \mathcal{F}_t\)。则对 \(A \in \mathcal{F}_{0^+}\)，\(\mathbb{P}(A) \in \{0, 1\}\)。特别地，\(\mathbb{P}(\sup_{[0, \varepsilon]} B_t > 0) = 1\)（Brown 运动在原点附近立即取正值和负值）。

\textbf{证明}：\(\mathcal{F}_{0^+}\) 是 Brown 运动无穷小尾 \(\sigma\)-代数。对 \(A \in \mathcal{F}_{0^+}\) 和任意 \(\delta > 0\)，\(A \in \mathcal{F}_\delta\)。由 Markov 性，\(A\) 与 \(\sigma(B_{\delta + t} - B_\delta : t \geq 0)\) 独立，故 \(A\) 与自身独立，\(\mathbb{P}(A) = \mathbb{P}(A)^2 \in \{0,1\}\)。对第二个断言，\(\{\sup_{[0,\varepsilon]} B_t > 0\} \in \mathcal{F}_{0^+}\)，由对称性该概率为正，故必为 1。\(\blacksquare\)

\subsection{108.3 连续时间鞅}\label{ux8fdeux7eedux65f6ux95f4ux9785}

\textbf{定义 108.2}（连续时间鞅）：过程 \(\{M_t\}_{t \geq 0}\) 关于滤过 \(\{\mathcal{F}_t\}\) 是\textbf{鞅}，如果 \(\mathbb{E}[|M_t|] < \infty\) 且 \(\mathbb{E}[M_t \mid \mathcal{F}_s] = M_s\)（\(s < t\)）。

\textbf{命题 108.2}（Brown 运动生成的鞅）： 1. \(B_t\) 是鞅 2. \(B_t^2 - t\) 是鞅 3. \(\exp(\theta B_t - \theta^2 t/2)\) 是鞅（指数鞅）

\textbf{定理 108.5}（Doob 不等式，连续时间）：设 \(\{M_t\}\) 是右连续下鞅。则

\[\mathbb{P}\left(\sup_{0 \leq s \leq t} |M_s| \geq \lambda\right) \leq \frac{\mathbb{E}[|M_t|]}{\lambda}\]

若 \(M_t\) 是鞅且 \(p > 1\)，则 \(\mathbb{E}[\sup_{0 \leq s \leq t} |M_s|^p] \leq \left(\frac{p}{p-1}\right)^p \mathbb{E}[|M_t|^p]\)。

\textbf{证明}：对离散时间逼近：取分划 \(0 = t_0 < t_1 < \cdots < t_n = t\)，令 \(M_k^{(n)} = M_{t_k}\)。对离散鞅应用 Doob 不等式得 \(\mathbb{P}(\max_k |M_{t_k}| \geq \lambda) \leq \mathbb{E}[|M_t|]/\lambda\)。由右连续性，\(\sup_{[0,t]} |M_s| = \lim_n \max_k |M_{t_k}|\)（a.s.），取 \(n \to \infty\) 极限即得连续时间版本。\(L^p\) 版本由 \(p>1\) 时 \(|x|^p\) 的凸性及离散 Doob \(L^p\) 不等式延拓得到。\(\blacksquare\)

\subsection{108.4 Itô 积分简介}\label{ituxf4-ux79efux5206ux7b80ux4ecb}

\textbf{定义 108.3}（Itô 积分，启发式）：对适应过程 \(H_s\)，Itô 积分 \(\int_0^t H_s dB_s\) 定义为

\[\int_0^t H_s dB_s = \lim_{\|\Pi\| \to 0} \sum_{i} H_{t_i} (B_{t_{i+1}} - B_{t_i})\]

（在 \(L^2\) 意义下），其中被积函数在区间左端点取值（非预见性）。

\textbf{定理 108.6}（Itô 公式）：设 \(f(t, x)\) 是 \(C^{1,2}\) 函数。则

\[df(t, B_t) = \frac{\partial f}{\partial t}(t, B_t) dt + \frac{\partial f}{\partial x}(t, B_t) dB_t + \frac{1}{2} \frac{\partial^2 f}{\partial x^2}(t, B_t) dt\]

最后一项是 Itô 修正项，来源于 Brown 运动的非零二次变差。

\textbf{证明}：Taylor 展开：\(f(t+\Delta t, B_{t+\Delta t}) - f(t, B_t) = f_t \Delta t + f_x \Delta B_t + \frac{1}{2} f_{xx} (\Delta B_t)^2 + o(\Delta t)\)。对于 Itô 积分，\(\Delta B_t \approx dB_t\) 且 \((\Delta B_t)^2 \approx dt\)（因 \(\mathbb{E}[(\Delta B_t)^2] = \Delta t\)，且其方差为 \(O((\Delta t)^2)\)）。取 Riemann 和极限，\((\Delta B_t)^2\) 项收敛于 \(dt\) 而非 0------这是 Itô 公式区别于经典链式法则的关键。严格证明需验证 \((\Delta B_t)^2 - \Delta t\) 的 \(L^2\) 收敛性，由 Brown 运动二次变差性质 \([B]_t = t\) 确保。\(\blacksquare\)

\hrulefill

\hrulefill

\hrulefill

\hrulefill

\hrulefill

\hrulefill

\newpage

\chapter{10-概率与统计/01-概率论/03-统计力学概率结构.md}\label{ux6982ux7387ux4e0eux7edfux8ba101-ux6982ux7387ux8bba03-ux7edfux8ba1ux529bux5b66ux6982ux7387ux7ed3ux6784.md}

\subsection{274.A 统计力学补充}\label{a-ux7edfux8ba1ux529bux5b66ux8865ux5145}

\begin{quote}
统计力学的数学基础涵盖概率论、泛函分析和算子代数方法。本卷从 Gibbs 测度的概率论表述出发，讨论热力学极限下的测度非唯一性及其与奇点理论的联系，建立重整化群的抽象数学结构，研究精确可解模型的组合数学与代数方法，最后涉及非平衡态的概率描述。
\end{quote}

\begin{quote}
\textbf{卷序说明}：本章节号（Ch 282-216）反映其在全书逻辑链中的位置。本卷位于V53，作为概率论与数学物理方法的交叉专题卷。
\end{quote}

\hrulefill

\hrulefill

\hrulefill

\hrulefill

\hrulefill

\hrulefill

\section{第339章：Gibbs 测度、热力学极限与相变}\label{ux7b2c339ux7ae0gibbs-ux6d4bux5ea6ux70edux529bux5b66ux6781ux9650ux4e0eux76f8ux53d8}

Gibbs 测度是连接概率论与统计力学的核心数学对象，它将 Boltzmann-Gibbs 分布 \(e^{-\beta H}\) 从有限系统推广到无穷格点 \(\mathbb{Z}^d\) 上的宏观系统。这一推广并非平凡的极限过程------当温度低于临界值时，无穷体积极限可能出现多个不同的 Gibbs 测度，这便是相变的数学实质：对称性的自发破缺。Dobrushin、Lanford 和 Ruelle（1968-1970）给出的 DLR 方程精确定义了无穷体积 Gibbs 测度的条件期望相容性，将物理直觉转化为严格的概率论框架。从自由能密度的热力学极限（Ruelle-Fisher 定理）的次可加性证明，到 Ising 模型中测度非唯一性的 Peierls 组合论证（低维晶格中畴壁的竞争），本章展示了概率论方法如何为物理相变提供严格的数学基础。本章内容与 Gibbs 测度的变分原理（第 276 章的重整化群）和精确可解模型（第 277 章）构成统计力学概率描述的完整体系。

\subsection{213.1 Gibbs 测度与配分函数}\label{gibbs-ux6d4bux5ea6ux4e0eux914dux5206ux51fdux6570}

\textbf{定义 213.1}（有限体积 Gibbs 测度）：设 \(\Lambda \subset \mathbb{Z}^d\) 为有限区域，边界条件为 \(\omega \in \Omega_{\Lambda^c}\)。定义 \(\Lambda\) 上的 Hamilton 量 \(H_\Lambda^\omega(\sigma)\)（\(\sigma \in \Omega_\Lambda\)）。给定参数 \(\beta > 0\)，\textbf{有限体积 Gibbs 测度}为 \(\Omega_\Lambda\) 上的概率测度： \[\mu_{\Lambda, \beta}^\omega(\sigma) = \frac{1}{Z_{\Lambda, \beta}^\omega} \exp(-\beta H_\Lambda^\omega(\sigma))\] 其中\textbf{配分函数}（partition function）为归一化常数 \[Z_{\Lambda, \beta}^\omega = \sum_{\sigma \in \Omega_\Lambda} \exp(-\beta H_\Lambda^\omega(\sigma))\]

\textbf{定义 213.2}（自由能泛函）：有限体积\textbf{自由能}定义为配分函数的对数生成函数： \[F_{\Lambda, \beta}^\omega = -\frac{1}{\beta} \log Z_{\Lambda, \beta}^\omega\]

自由能关于参数 \(\beta\) 的导数给出各阶累积量。

\textbf{定义 213.3}（无穷体积 Gibbs 测度 / DLR 方程，Dobrushin-Lanford-Ruelle 1968-1970）：无穷体积 \textbf{Gibbs 测度} \(\mu_\beta\) 是 \(\Omega = \Omega_{\mathbb{Z}^d}\) 上的概率测度，满足对任意有限区域 \(\Lambda\)，条件分布 \(\mu_\beta(\cdot \mid \omega_{\Lambda^c})\) 由有限体积 Gibbs 测度给出（DLR 相容性条件）。Gibbs 测度的全体记为 \(\mathcal{G}(\beta)\)。

\subsection{213.2 热力学极限}\label{ux70edux529bux5b66ux6781ux9650}

\textbf{定义 213.4}（热力学极限）：令 \(\Lambda_n \uparrow \mathbb{Z}^d\)（如 \(\Lambda_n = [-n, n]^d \cap \mathbb{Z}^d\)）。若自由能密度 \[f(\beta) = \lim_{n \to \infty} \frac{1}{|\Lambda_n|} F_{\Lambda_n, \beta}^\omega\] 存在且与边界条件 \(\omega\) 的选取无关，则称 Gibbs 测度族具有良定义的\textbf{热力学极限}。\(f(\beta)\) 为自由能密度。

\textbf{定理 213.1}（热力学极限的存在性 / Ruelle-Fisher 定理）：对具有有限程相互作用的平移不变格点系统，热力学极限存在。具体而言，对于有限程势 \(\Phi\) 满足 \(\sum_{X \ni 0} \|\Phi_X\|_\infty < \infty\)，自由能密度 \(f(\beta)\) 是良定义的连续函数。

\textbf{证明}：对有限区域 \(\Lambda_n\)，配分函数 \(Z_{\Lambda_n} = \sum_{\sigma} \exp(-\beta H_{\Lambda_n}(\sigma))\)。由于势 \(\Phi\) 有有限程，边界效应仅涉及 \(\partial\Lambda_n\) 附近的位点，其贡献为 \(O(|\partial\Lambda_n|)\)。利用次可加性：对不相交区域 \(\Lambda_1, \Lambda_2\)，\(\log Z_{\Lambda_1 \cup \Lambda_2} \geq \log Z_{\Lambda_1} + \log Z_{\Lambda_2} - O(|\partial\Lambda_1 \cap \partial\Lambda_2|)\)。取 \(\Lambda_n = [-n, n]^d \cap \mathbb{Z}^d\)，序列 \(a_n = -\frac{1}{\beta|\Lambda_n|} \log Z_{\Lambda_n}\) 满足 \(a_{n+m} \leq a_n + a_m + o(1)\)。由次可加引理，\(\lim_{n \to \infty} a_n = \inf_n a_n = f(\beta)\) 存在。边界条件依赖性由 \(\sum_{X \ni 0} \|\Phi_X\|_\infty < \infty\) 保证被 \(|\partial\Lambda_n|/|\Lambda_n| \to 0\) 压制，故极限与边界条件无关。\(f(\beta)\) 的连续性由有限体积自由能的局部一致收敛性及配分函数的解析性得出。\(\blacksquare\)

\subsection{213.3 Ising 模型与相变的 Peierls 论证}\label{ising-ux6a21ux578bux4e0eux76f8ux53d8ux7684-peierls-ux8bbaux8bc1}

\textbf{定义 213.4}（Ising 模型）：\(d\) 维\textbf{Ising 模型}是定义在格点 \(\mathbb{Z}^d\) 上的随机过程，每个格点 \(i\) 处变量 \(\sigma_i \in \{-1, +1\}\)。Hamilton 量为 \[H_\Lambda(\sigma) = -J \sum_{\langle i, j \rangle \cap \Lambda \neq \emptyset} \sigma_i \sigma_j - h \sum_{i \in \Lambda} \sigma_i\] 其中 \(\langle i, j \rangle\) 表示最近邻对，\(J, h \in \mathbb{R}\) 为耦合参数。

\textbf{定义 213.5}（相变------测度非唯一性）：称 Gibbs 测度族在 \(\beta\) 处发生\textbf{相变}，若 \(|\mathcal{G}(\beta)| \geq 2\)（即无穷体积 Gibbs 测度不唯一）。由定义 213.3，测度非唯一性等价于 DLR 方程有多于一组的解。在热力学极限下，相变点对应于自由能密度 \(f(\beta)\) 的非解析性（奇点）。

\textbf{定理 213.2}（Peierls 论证：Ising 模型的测度非唯一性，Peierls 1936，严格化 Griffiths 1964）：对 \(d \geq 2\) 维零场（\(h = 0\)）Ising 模型，存在临界参数 \(\beta_c(d) > 0\) 使得： - \(\beta > \beta_c\) 时，\(|\mathcal{G}(\beta)| \geq 2\)：存在至少两个不同的平移不变 Gibbs 测度 \(\mu_\beta^+\) 和 \(\mu_\beta^-\)。 - \(\beta < \beta_c\) 时，Gibbs 测度唯一：\(|\mathcal{G}(\beta)| = 1\)。

\textbf{证明}：取正边界条件（\(\omega_i = +1\) 对所有 \(i \in \Lambda^c\)）。定义 Peierls 轮廓：对配置 \(\sigma \in \Omega_\Lambda\)，在 \(\mathbb{Z}^d\) 的对偶格点上构造闭合曲面，分离 \(\sigma_i = +1\) 和 \(\sigma_i = -1\) 的区域。每条轮廓 \(\gamma\) 是长度为 \(|\gamma|\) 的闭合路径。翻转轮廓内的所有自旋将 \(\sigma_i = -1\) 变为 \(+1\)，能量变化为 \(\Delta H = 2J|\gamma|\)（零场情形），故含轮廓 \(\gamma\) 的配置的概率满足 \(\mu_{\Lambda,\beta}^+(\gamma \text{ 出现}) \leq \exp(-2\beta J|\gamma|)\)。固定原点 \(0\)，事件 \(\{\sigma_0 = -1\}\) 意味着存在包围 \(0\) 的轮廓。长度为 \(L\) 的包围 \(0\) 的轮廓数目不超过 \(3^L\)（对偶格点上的自回避路径）。\(d \geq 2\) 时，求和得 \(\mu_{\Lambda,\beta}^+(\sigma_0 = -1) \leq \sum_{L \geq 4} 3^L e^{-2\beta J L} \to 0\) 当 \(\beta \to \infty\)。取 \(\beta_c\) 为使该和 \(< 1/2\) 的最小 \(\beta\) 值。\(\beta > \beta_c\) 时，\(\mu_\beta^+ \neq \mu_\beta^-\)（后者由对称性给出 \(\mu_\beta^-(\sigma_0 = +1) < 1/2\)），故 \(|\mathcal{G}(\beta)| \geq 2\)。\(\beta < \beta_c\) 时唯一性由 cluster expansion 或 Dobrushin 唯一性准则给出。\(\blacksquare\)

\textbf{推论 213.1}（\(d=1\) 唯一性）：一维 Ising 模型中 Peierls 论证失效（轮廓为单个点）。对 \(d=1\)，任意 \(\beta > 0\) 下 \(|\mathcal{G}(\beta)| = 1\)。

\textbf{证明}：一维 Ising 模型可精确求解。转移矩阵 \(T = \begin{pmatrix} e^{\beta J} & e^{-\beta J} \\ e^{-\beta J} & e^{\beta J} \end{pmatrix}\) 的特征值为 \(\lambda_1 = 2\cosh(\beta J)\)，\(\lambda_2 = 2\sinh(\beta J)\)。在热力学极限下，两点关联函数 \(\langle \sigma_0 \sigma_n \rangle = (\lambda_2/\lambda_1)^n = (\tanh(\beta J))^n \to 0\)（\(n \to \infty\)）。关联长度的有限性表明系统始终处于高温无序相，Gibbs 测度对任意 \(\beta > 0\) 唯一。从 DLR 方程角度，一维系统的边界条件影响在热力学极限下指数衰减为零，转移矩阵的最大特征值对应的特征向量唯一确定无穷体积 Gibbs 测度。\(\blacksquare\)

\textbf{评注 213.1}：Peierls 论证的核心在于利用几何-组合方法（轮廓展开）将测度非唯一性问题转化为轮廓系统的低密度展开。这一方法后来发展为 \textbf{Pirogov-Sinai 理论}（1975-1976），系统研究了低温相图的普适结构。

\hrulefill

\hrulefill

\hrulefill

\hrulefill

\hrulefill

\hrulefill

\section{第340章：重整化群的数学结构}\label{ux7b2c340ux7ae0ux91cdux6574ux5316ux7fa4ux7684ux6570ux5b66ux7ed3ux6784}

重整化群（Renormalization Group, RG）是 Kenneth Wilson 于 1971 年引入的用于分析多尺度系统的革命性方法，他因此在 1982 年获诺贝尔物理学奖。RG 的核心思想是将物理学中的粗粒化（Kadanoff 块自旋变换）形式化为 Hamilton 量参数空间上的半群作用：从微观尺度出发，逐步''积分掉''短程涨落，得到宏观尺度的有效理论。RG 流的不动点对应了临界系统，不动点的谱分类（相关方向、无关方向、边缘方向）决定了临界指数------这解释了为什么微观上截然不同的物理系统（如液-气相变的 Ising 模型和铁磁-顺磁相变的 Heisenberg 模型）在临界点展现相同的标度律（普适性）。本章在统计力学和 Gibbs 测度（第 275 章）的基础上，建立 RG 变换的数学定义，给出不动点谱分析与临界指数的定标关系，引入 \(\varepsilon\)-展开和 Wilson-Fisher 不动点，最终以共形场论（CFT）和 BPZ 分类定理收尾，将临界现象与无穷维对称代数（Virasoro 代数）深刻关联。

\subsection{214.1 重整化群变换}\label{ux91cdux6574ux5316ux7fa4ux53d8ux6362}

\textbf{定义 214.1}（重整化群变换 / Kadanoff 块自旋变换）：\textbf{重整化群变换} \(\mathcal{R}\) 是作用在 Hamilton 量参数空间上的变换，将 Hamilton 量 \(H\) 映射到粗粒化后的有效 Hamilton 量 \(H' = \mathcal{R}(H)\)：

\begin{enumerate}
\def\labelenumi{\arabic{enumi}.}
\tightlist
\item
  将格点系统分组为 \(b^d\) 个变量组成的块
\item
  每个块以块变量代替
\item
  对短程涨落取条件期望，得到块变量的有效 Hamilton 量
\end{enumerate}

在参数空间中，\(\mathcal{R}\) 构成半群（一般不可逆）。

\textbf{定义 214.2}（RG 不动点与临界曲面）：参数空间中的\textbf{不动点} \(H^*\) 满足 \(\mathcal{R}(H^*) = H^*\)。临界曲面是流向不动点的 Hamilton 量集合。不动点的稳定性由 Fréchet 导数 \(D\mathcal{R}(H^*)\) 的谱决定。

\textbf{定理 214.1}（RG 不动点的谱分类）： - \textbf{无关方向}（特征值 \(|\lambda| < 1\)）：流向不动点 - \textbf{相关方向}（特征值 \(|\lambda| > 1\)）：流离不动点 - 相关方向的个数决定了独立临界指数的个数

\textbf{证明}：在不动点 \(H^*\) 处，将 Hamilton 量展开为 \(H = H^* + \sum_\alpha t_\alpha \mathcal{O}_\alpha\)，其中 \(\mathcal{O}_\alpha\) 为标度算子。RG 变换的线性化 \(D\mathcal{R}(H^*)\) 作用在耦合常数空间上：\(t_\alpha' = \sum_\beta L_{\alpha\beta} t_\beta\)。对角化得 \(L\) 的本征值 \(\lambda_\alpha\) 和本征方向。沿本征方向 \(u_\alpha\)，\(n\) 次 RG 变换后 \(t_\alpha^{(n)} = \lambda_\alpha^n t_\alpha^{(0)}\)。\(|\lambda_\alpha| < 1\) 时 \(t_\alpha^{(n)} \to 0\)（无关），\(|\lambda_\alpha| > 1\) 时 \(t_\alpha^{(n)}\) 发散（相关），\(|\lambda_\alpha| = 1\) 为边缘方向。相关方向数等于需要调节以达到临界面的物理参数个数，每个相关方向对应一个独立临界指数。\(\blacksquare\)

\textbf{定理 214.2}（RG 与临界指数）：若相关方向的本征值为 \(\lambda_t = b^{y_t}\)，\(\lambda_h = b^{y_h}\)，则临界指数为

\[\nu = \frac{1}{y_t}, \quad \eta = d + 2 - 2y_h, \quad \alpha = 2 - \frac{d}{y_t}, \quad \beta = \frac{d - y_h}{y_t}, \quad \gamma = \frac{2y_h - d}{y_t}\]

\textbf{证明}：设 RG 变换的标度因子为 \(b\)。在不动点附近，约化温度 \(t\) 和磁场 \(h\) 按 \(t' = b^{y_t} t\)，\(h' = b^{y_h} h\) 变换。关联长度按 \(\xi(t, h) = b \xi(t', h')\) 变换。令 \(b = |t|^{-1/y_t}\)，得 \(\xi(t, 0) \sim |t|^{-1/y_t}\)，故 \(\nu = 1/y_t\)。磁化强度 \(m = -\partial f/\partial h\) 在 RG 下按 \(m(t, h) = b^{y_h - d} m(t', h')\) 变换。令 \(h = 0\)，\(b = |t|^{-1/y_t}\)，得 \(m \sim |t|^{(d-y_h)/y_t}\)，故 \(\beta = (d-y_h)/y_t\)。磁化率 \(\chi = \partial m/\partial h\) 按 \(\chi(t) = b^{2y_h - d} \chi(t')\) 变换，得 \(\gamma = (2y_h - d)/y_t\)。比热 \(C = -\partial^2 f/\partial t^2\) 按 \(C(t) = b^{2y_t - d} C(t')\) 变换，得 \(\alpha = 2 - d/y_t\)。关联函数在临界点的幂律衰减给出 \(\eta = d + 2 - 2y_h\)。\(\blacksquare\)

\subsection{\texorpdfstring{214.2 \(\varepsilon\)-展开}{214.2 \textbackslash varepsilon-展开}}\label{varepsilon-ux5c55ux5f00}

\textbf{定理 214.3}（Wilson-Fisher 不动点与 \(\varepsilon\)-展开，Wilson 1972）：在 \(d = 4 - \varepsilon\) 维（\(\varepsilon \ll 1\)）的 \(\phi^4\) 场论中，RG 不动点在 \(\varepsilon\) 阶处存在（Wilson-Fisher 不动点）。临界指数可展开为 \(\varepsilon\) 的渐近级数：

\[\eta = \frac{\varepsilon^2}{54} + O(\varepsilon^3), \quad \nu = \frac{1}{2} + \frac{\varepsilon}{12} + \frac{7\varepsilon^2}{162} + O(\varepsilon^3)\]

\textbf{证明}：在 \(d = 4 - \varepsilon\) 维，\(\phi^4\) 理论的 Lagrange 量为 \(\mathcal{L} = \frac{1}{2}(\nabla\phi)^2 + \frac{1}{2}r\phi^2 + \frac{u}{4!}\phi^4\)。Gauss 不动点 \((r=0, u=0)\) 在 \(d=4\) 时，\(u\) 为边缘参数。在 \(d=4-\varepsilon\) 下，RG 流方程至一阶：\(\frac{dr}{d\ell} = 2r + \frac{n+2}{3} \frac{u}{1+r} + O(u^2)\)，\(\frac{du}{d\ell} = \varepsilon u - \frac{n+8}{3} u^2 + O(u^3)\)。\(u\) 的 \(\beta\) 函数在 \(u^* = \frac{3\varepsilon}{n+8} + O(\varepsilon^2)\) 处有非平凡零点，即 Wilson-Fisher 不动点。在该不动点处线性化 RG 变换，温度本征值 \(y_t = 2 - \frac{n+2}{n+8}\varepsilon + O(\varepsilon^2)\)，场本征值 \(y_h = \frac{d}{2}+1 - \frac{n+2}{2(n+8)}\varepsilon + O(\varepsilon^2)\)。代入定标关系 \(\nu = 1/y_t\)，\(\eta = d+2-2y_h\)，取 \(n=1\)（Ising）并展开至 \(\varepsilon^2\) 阶即得所给表达式。\(\blacksquare\)

\(\varepsilon\)-展开的关键在于识别上临界维数 \(d_c = 4\) 的特殊角色：在四维以上，Gauss 不动点支配临界行为且平均场指数精确，而在四维以下系统流向 Wilson-Fisher 不动点，展现出非平凡的临界指数。这一维数分界揭示了平均场理论失效的数学本质------它是耦合常数 \(u\) 的标度维数在 \(d=4\) 时变为零（边缘）的直接后果。

\textbf{定理 214.4}（上临界维数）：\(d_c = 4\) 是 Ising 普适类的\textbf{上临界维数}。\(d \geq 4\) 时，Gauss 不动点稳定，平均场临界指数精确。\(d < 4\) 时，系统流向 Wilson-Fisher 不动点。

\textbf{证明}：在 \(\phi^4\) 场论的 RG 流方程中，耦合常数 \(u\) 的标度维数为 \([u] = 4 - d\)。当 \(d > 4\) 时，\([u] < 0\)，\(u\) 为无关参数，Gauss 不动点 \((u=0)\) 是红外稳定的。此时上临界维数 \(d_c = 4\)。\(d = d_c\) 时 \(u\) 为边缘参数，需考虑对数修正。\(d < 4\) 时 \([u] > 0\)，\(u\) 为相关参数，Gauss 不动点不稳定，RG 流将系统推向 Wilson-Fisher 不动点。在 \(d \geq 4\) 时，Gauss 不动点处的本征值给出平均场临界指数 \(\nu = 1/2\)，\(\eta = 0\)，\(\alpha = 0\)，\(\beta = 1/2\)，\(\gamma = 1\)，\(\delta = 3\)，这些指数在 \(d \geq 4\) 时精确成立（仅在对数修正意义上）。\(\blacksquare\)

\subsection{214.3 共形场论与临界现象}\label{ux5171ux5f62ux573aux8bbaux4e0eux4e34ux754cux73b0ux8c61}

\textbf{定理 214.5}（Polyakov 共形不变性假说，1970）：在临界点，系统具有尺度不变性及更丰富的共形对称性。\(d=2\) 维的临界系统由\textbf{共形场论}（CFT）描述。

\textbf{证明}：临界点处关联长度发散，系统在所有尺度上自相似。尺度不变性意味着理论在伸缩变换 \(x \mapsto \lambda x\) 下不变。在 \(d=2\) 维，尺度不变性加上平移和旋转对称性自然地推广为共形不变性------共形变换为局部伸缩变换，由任意全纯函数 \(f(z)\) 生成。二维共形群为无穷维 Virasoro 代数，其生成元 \(L_n\) 满足 \([L_n, L_m] = (n-m)L_{n+m} + \frac{c}{12}n(n^2-1)\delta_{n+m,0}\)。临界系统的关联函数被共形不变性严格约束：二点函数 \(\langle \phi(z_1)\phi(z_2) \rangle = |z_1-z_2|^{-2h}\)，其中 \(h\) 为标度维数。三点函数结构由共形 Ward 恒等式完全确定。\(\blacksquare\)

\textbf{定理 214.6}（BPZ 分类定理，Belavin-Polyakov-Zamolodchikov 1984）：\(d=2\) 维共形场论由 Virasoro 代数（中心荷 \(c\)）的酉表示完全分类。\(c < 1\) 的酉极小模型有离散系列：

\[c = 1 - \frac{6}{m(m+1)} \quad (m = 3, 4, 5, \ldots)\]

Ising 模型对应 \(m = 3\)（\(c = 1/2\)），三临界 Ising 对应 \(m = 4\)（\(c = 7/10\)）。

\textbf{证明}：Virasoro 代数的酉表示要求 \(c \geq 0\) 且 \(h \geq 0\)。Friedan-Qiu-Shenker 分析表明，\(c < 1\) 的酉表示仅存在于离散序列 \(c = 1 - 6/m(m+1)\)，对应的共形维数 \(h_{r,s} = \frac{[(m+1)r - ms]^2 - 1}{4m(m+1)}\)（\(1 \leq r \leq m-1\)，\(1 \leq s \leq m\)）。Kac 行列式在 \(c\) 取这些值时为零，对应的 Virasoro 模可约，商去零模后得到酉不可约表示。表征理论（character）满足 modular 不变性，约束了可能的配分函数。极小模型提供了无穷多个精确可解的二维临界系统，其临界指数为有理数。Ising 模型对应 \(m=3\) 的临界指数 \(\beta = 1/8\)，\(\nu = 1\) 等。\(\blacksquare\)

\hrulefill

\hrulefill

\hrulefill

\hrulefill

\hrulefill

\hrulefill

\section{第341章：精确可解模型}\label{ux7b2c341ux7ae0ux7cbeux786eux53efux89e3ux6a21ux578b}

精确可解模型是统计力学中极为珍稀的一类系统------它们的配分函数、自由能和关联函数可以解析地而非数值地计算，从而为临界现象提供非微扰的数学理解。二维 Ising 模型的 Onsager 解（1944 年）是这一领域的珠穆朗玛峰：通过将配分函数表示为转移矩阵的迹，利用 Clifford 代数的表示论和 Jordan-Wigner 变换将自旋算符映射为自由费米子，最终将对数奇异性以封闭积分形式精确表达------这是数学史上第一次在相互作用系统中严格展示相变的存在性。Bethe 拟设（Bethe, 1931 年）和 Yang-Baxter 方程（1967-1972 年）则构成了可解模型更为系统化的代数框架：Yang-Baxter 方程 \(R_{12}R_{13}R_{23} = R_{23}R_{13}R_{12}\) 是量子可积性的代数条件，与量子群和辫子群表示论有深刻联系。本章在 RG 理论（第 276 章）和临界现象（第 275 章）的基础上，以 Onsager 解、Bethe 拟设和 Baxter 八顶角模型解为三大支柱，展示精确可解模型的数学结构。

\subsection{215.1 二维 Ising 模型的 Onsager 解}\label{ux4e8cux7ef4-ising-ux6a21ux578bux7684-onsager-ux89e3}

\textbf{定理 215.1}（Onsager 精确解，1944）：\(d=2\) 维零场 Ising 模型的自由能密度为

\[-\beta f = \log 2 + \frac{1}{2\pi^2} \int_0^\pi \int_0^\pi \log[\cosh(2\beta J_x) \cosh(2\beta J_y) - \sinh(2\beta J_x) \cos \omega_1 - \sinh(2\beta J_y) \cos \omega_2] \, d\omega_1 d\omega_2\]

（在各向同性情形 \(J_x = J_y = J\) 下）。

\textbf{证明}：Onsager 使用转移矩阵方法。将二维 Ising 模型视为一维量子自旋链的转移矩阵 \(V = \exp(\beta J \sum_i \sigma_i^z \sigma_{i+1}^z) \exp(\beta J \sum_i \sigma_i^x)\)。通过 Jordan-Wigner 变换 \(\sigma_i^z = \prod_{j<i}(1-2c_j^\dagger c_j)(c_i + c_i^\dagger)\)，将自旋算符映射为费米子算符。Fourier 变换后，转移矩阵化为 \(V = \prod_{k>0} \exp(-\varepsilon_k \gamma_k^\dagger \gamma_k)\)，其中准粒子能量 \(\varepsilon_k\) 满足 \(\cosh \varepsilon_k = \cosh(2\beta J) \cosh(2\beta J) - \sinh(2\beta J)(\cos k_x + \cos k_y)\)。配分函数 \(Z = \operatorname{Tr}(V^N) = \prod_{k>0} 2\cosh(N\varepsilon_k/2)\)，自由能密度 \(-\beta f = \lim_{N\to\infty} \frac{1}{N^2} \log Z = \log 2 + \frac{1}{2\pi^2} \int_0^\pi \int_0^\pi \log[\cosh^2(2\beta J) - \sinh(2\beta J)(\cos\omega_1 + \cos\omega_2)] d\omega_1 d\omega_2\)。\(\blacksquare\)

\textbf{定理 215.2}（临界点与对数奇异性）：临界耦合由方程 \(\sinh(2\beta_c J) = 1\) 确定。自由能在 \(\beta_c\) 处有对数奇异性：

\[C \sim -\log|\beta - \beta_c| \quad (\beta \to \beta_c)\]

即 \(\alpha = 0\)（对数发散）。

\textbf{证明}：自由能密度 \(f(\beta)\) 中的被积函数在 \(\omega_1 = \omega_2 = 0\) 处有奇异性。当 \(\beta \to \beta_c\) 时，\(\cosh^2(2\beta J) - 2\sinh(2\beta J) = (\cosh(2\beta J) - \sinh(2\beta J))^2 - 1\)，在 \(\beta_c\) 处为零。将积分在 \(\omega_1 = \omega_2 = 0\) 附近展开，主导项为 \(\int_0^\varepsilon \int_0^\varepsilon \log[\delta^2 + \omega_1^2 + \omega_2^2] d\omega_1 d\omega_2\)，其中 \(\delta = |\beta - \beta_c|\)。极坐标变换得 \(\int_0^\varepsilon r \log(\delta^2 + r^2) dr \sim \delta^2 \log \delta\)，其对 \(\beta\) 的二阶导数为 \(\partial^2 f/\partial\beta^2 \sim -\log|\beta - \beta_c|\)。比热 \(C = -\beta^2 \partial^2 f/\partial\beta^2 \sim -\log|\beta - \beta_c|\)，故 \(\alpha = 0\)（对数发散）。\(\blacksquare\)

\textbf{定理 215.3}（自发磁化，杨振宁 1952）：\(\beta > \beta_c\) 时，二维 Ising 模型的自发磁化为

\[m^* = (1 - \sinh^{-4}(2\beta J))^{1/8}\]

临界指数 \(\beta = 1/8\)。\(\beta \to \beta_c^+\) 时 \(m^* \sim (\beta - \beta_c)^{1/8}\)。

\textbf{证明}：杨振宁通过计算零场下二点关联函数的无穷远极限得到自发磁化。关联函数 \(\langle \sigma_{0,0} \sigma_{M,N} \rangle\) 的 Toeplitz 行列式表示由 Szegő 极限定理给出渐近行为。当 \(M,N \to \infty\) 时，\(\langle \sigma_{0,0} \sigma_{M,N} \rangle \to (m^*)^2\)。Toeplitz 行列式的生成函数为 \(g(\theta) = \sqrt{\frac{1 - \alpha_1 e^{i\theta}}{1 - \alpha_2 e^{-i\theta}}}\)，其中 \(\alpha_1 = \sinh(2\beta J)\cosh(2\beta J)\)，\(\alpha_2 = \coth(2\beta J)\)。由 Szegő 定理，\(\lim_{n\to\infty} D_n(g) = \exp(\sum_{k=1}^\infty k g_k g_{-k})\)，其中 \(g_k = \frac{1}{2\pi}\int_0^{2\pi} e^{-ik\theta} \log g(\theta) d\theta\)。计算得 \(m^* = (1 - \sinh^{-4}(2\beta J))^{1/8}\)。在 \(\beta \to \beta_c^+\) 时，\(\sinh(2\beta_c J) = 1\)，展开得 \(m^* \sim (\beta - \beta_c)^{1/8}\)，故 \(\beta = 1/8\)。\(\blacksquare\)

\textbf{方法概要}：Onsager 使用转移矩阵方法将配分函数化为 \(2^N \times 2^N\) 矩阵的迹，通过 Clifford 代数的表示论将问题转化为自由费米子问题，使用 Jordan-Wigner 变换和 Bogoliubov 变换对角化。其数学核心是 Clifford 代数的表示论与 Toeplitz 行列式的渐近分析（Szegő 极限定理）。

\subsection{215.2 Bethe 拟设}\label{bethe-ux62dfux8bbe}

\textbf{定义 215.1}（Bethe 拟设，Bethe 1931）：\textbf{Bethe 拟设}是求解一维量子多体系统的方法。一维 Heisenberg 自旋链（XXX 模型，\(J_x = J_y = J_z\)）的 Hamilton 量为

\[H = -J \sum_{i=1}^N (\sigma_i^x \sigma_{i+1}^x + \sigma_i^y \sigma_{i+1}^y + \sigma_i^z \sigma_{i+1}^z)\]

\textbf{定理 215.4}（Bethe 拟设解，Bethe 1931）：一维 Heisenberg 反铁磁链的基态可通过 Bethe 拟设严格求解。\(N\) 个位点的 Bethe 波函数为

\[\psi(x_1, \ldots, x_M) = \sum_{P \in S_M} A_P e^{i \sum_{j=1}^M k_{P_j} x_j}\]

其中波数 \(k_j\) 满足\textbf{Bethe 拟设方程}：

\[e^{i k_j N} = \prod_{\ell \neq j}^M \frac{\sin(k_j - k_\ell + 2i\eta)}{\sin(k_j - k_\ell - 2i\eta)}\]

\textbf{证明}：Bethe 拟设的核心是二体散射的可分解性。将 \(M\) 个下自旋（magnons）的波函数取为平面波的叠加 \(\psi(x_1,\ldots,x_M) = \sum_P A_P \exp(i\sum_j k_{P_j} x_j)\)。两粒子散射矩阵 \(S(k_j, k_\ell) = -\frac{\sin(k_j - k_\ell + 2i\eta)}{\sin(k_j - k_\ell - 2i\eta)}\) 由相互作用参数 \(\eta\) 决定。多体散射的分解性要求 Yang-Baxter 方程 \(S_{12}S_{13}S_{23} = S_{23}S_{13}S_{12}\) 成立，该方程在 Bethe 拟设中自动满足。周期性边界条件要求 \(e^{i k_j N} = \prod_{\ell \neq j} S(k_j, k_\ell)\)，即 Bethe 拟设方程。取对数得 \(N k_j = 2\pi I_j - i\sum_{\ell \neq j} \log S(k_j, k_\ell)\)，其中 \(I_j\) 为（半）整数量子数。在热力学极限下，Bethe 拟设方程化为关于根密度 \(\rho(k)\) 的线性积分方程，可严格求解。\(\blacksquare\)

\textbf{定理 215.5}（Yang-Baxter 方程，杨振宁 1967, Baxter 1972）：Bethe 拟设可解性的代数本质是 Yang-Baxter 方程：

\[R_{12}(u) R_{13}(u+v) R_{23}(v) = R_{23}(v) R_{13}(u+v) R_{12}(u)\]

其中 \(R_{ij} \in \text{End}(V_i \otimes V_j)\)。Yang-Baxter 方程是可积性的核心代数条件，与量子群和辫子群表示论有深刻联系。

\textbf{证明}：考虑转移矩阵 \(T(u) = \operatorname{Tr}_0(R_{0N}(u) \cdots R_{01}(u))\)，其中 \(R_{0i}\) 作用在辅助空间 \(V_0\) 和量子空间 \(V_i\) 上。Yang-Baxter 方程保证不同谱参数 \(u, v\) 的转移矩阵交换：\([T(u), T(v)] = 0\)。展开 \(\log T(u)\) 得无穷多个守恒量。Yang-Baxter 方程的三线性形式可通过图形记号理解为：两条世界线交叉的两种顺序等价。\(R\) 矩阵满足幺正条件 \(R_{12}(u)R_{21}(-u) = f(u)I\) 和正则条件 \(R_{12}(0) = P_{12}\)（置换算子）。对于 XXZ 自旋链，\(R\) 矩阵为 \(R(u) = \begin{pmatrix} \sinh(u+\eta) & 0 & 0 & 0 \\ 0 & \sinh u & \sinh\eta & 0 \\ 0 & \sinh\eta & \sinh u & 0 \\ 0 & 0 & 0 & \sinh(u+\eta) \end{pmatrix}\)，满足 Yang-Baxter 方程。\(\blacksquare\)

\subsection{215.3 六顶角模型与八顶角模型}\label{ux516dux9876ux89d2ux6a21ux578bux4e0eux516bux9876ux89d2ux6a21ux578b}

\textbf{定义 215.2}（六顶角模型）：\textbf{六顶角模型}（6-vertex model）是二维格点上每顶点有 6 种允许配置的格点模型。配分函数可通过 Bethe 拟设求解。

六顶角模型的物理意义在于描述二维冰的剩余熵------每个顶点处氧原子周围的四个氢原子中有两个靠近、两个远离，满足''冰规则''（ice rule），恰好有六种合法配置。Lieb 在 1967 年利用 Bethe 拟设方法精确求解了该模型的自由能，得到了比 Pauling 近似估计 \(S_0 = \ln(6/4)\) 更精确的剩余熵值 \(S_0 = \frac{3}{2}\ln(4/3)\)。

\textbf{定理 215.6}（Lieb 解，1967）：六顶角模型的自由能和极化率由 Lieb 使用 Bethe 拟设方法精确求解。

\textbf{证明}：六顶角模型的容许顶点配置为：箭头流入流出满足守恒律（流守恒）。转移矩阵 \(T(\lambda)\) 作用在 \(2^N\) 维空间上，由局部 \(R\) 矩阵构成。\(R\) 矩阵取为 \(R(\lambda) = \begin{pmatrix} a(\lambda) & 0 & 0 & 0 \\ 0 & b(\lambda) & c(\lambda) & 0 \\ 0 & c(\lambda) & b(\lambda) & 0 \\ 0 & 0 & 0 & a(\lambda) \end{pmatrix}\)，其中 \(a = \sinh(\lambda + \eta)\)，\(b = \sinh \lambda\)，\(c = \sinh \eta\)。Bethe 拟设方程化为 \([\frac{\sinh(\lambda_j + i\gamma/2)}{\sinh(\lambda_j - i\gamma/2)}]^N = \prod_{\ell \neq j} \frac{\sinh(\lambda_j - \lambda_\ell + i\gamma)}{\sinh(\lambda_j - \lambda_\ell - i\gamma)}\)。在热力学极限下，Bethe 根在复平面上凝聚为弦，自由能由根密度 \(\rho(\lambda)\) 的积分方程确定。Lieb 求解了该积分方程，得到自由能的封闭表达式。\(\blacksquare\)

六顶角模型仅涉及''流守恒''的顶点配置，而更一般的八顶角模型引入了对角顶点的两种额外配置。Baxter 在 1972 年利用 Yang-Baxter 方程和椭圆函数参数化，将精确求解的领域从六顶角推进到八顶角------这是统计力学可解模型中最重要的里程碑之一，其求解过程首次展示了 Yang-Baxter 方程作为可积性判据的强大力量。

\textbf{定理 215.7}（Baxter 八顶角模型解，1972）：八顶角模型（8-vertex model）是带有对角相互作用的最一般二维可解格点模型。Baxter 使用 Yang-Baxter 方程和转移矩阵的交换性，给出了自由能的参数化精确解。

\textbf{证明}：八顶角模型引入对角顶点配置（第七、八种顶点），其 \(R\) 矩阵在 \(V \otimes V\) 上的一般形式为 \(R(u) = \sum_{\alpha=0}^3 w_\alpha(u) \sigma^\alpha \otimes \sigma^\alpha\)，其中 \(\sigma^0 = I\)，\(\sigma^{1,2,3}\) 为 Pauli 矩阵。Yang-Baxter 方程约束 \(w_\alpha(u)\) 为特定椭圆函数。Baxter 发现 \(R\) 矩阵可由 Jacobi 椭圆函数参数化：\(w_0 : w_1 : w_2 : w_3 = \operatorname{sn}(u+\eta) : \operatorname{sn} u : \operatorname{sn} \eta : k \operatorname{sn} u \operatorname{sn}(u+\eta) \operatorname{sn} \eta\)。转移矩阵族的交换性 \([T(u), T(v)] = 0\) 由 Yang-Baxter 方程保证，由此产生无穷多守恒量。在热力学极限下，Baxter 使用坐标 Bethe 拟设求得自由能：\(-\beta f = \log(2 \max\{w_\alpha\}) + \frac{1}{2\pi} \int_{-\infty}^\infty \log g(\omega) d\omega\)，其中 \(g(\omega)\) 由椭圆函数的模参数确定。\(\blacksquare\)

\subsection{可解格点模型与概率论的深层联系}\label{ux53efux89e3ux683cux70b9ux6a21ux578bux4e0eux6982ux7387ux8bbaux7684ux6df1ux5c42ux8054ux7cfb}

可解格点模型不仅是统计力学的核心研究对象，也是现代概率论和表示论的交汇点。\textbf{Yang-Baxter 方程}（Yang 1967, Baxter 1972）是精确可解性的代数条件：\(R_{12}(u) R_{13}(u+v) R_{23}(v) = R_{23}(v) R_{13}(u+v) R_{12}(u)\)，其中 \(R(u)\) 是作用在 \(V \otimes V\) 上的矩阵。满足该方程的 \(R\) 矩阵确保转移矩阵族 \(T(u) = \operatorname{Tr}_0(R_{0N}(u) \cdots R_{01}(u))\) 对所有参数 \(u\) 交换，从而产生无穷多个守恒量。Yang-Baxter 方程是量子逆散射方法的代数基础，与量子群（Drinfeld-Jimbo 1985）和纽结理论（Jones 多项式）有深刻联系。\textbf{六顶角模型}（Lieb 1967）的精确解揭示了二维冰的剩余熵------\(S_0 = \frac{3}{2} \ln(4/3) \approx 0.432\)------这一结果与 Pauling 的估计（\(\ln(6/4) \approx 0.405\)）接近但更精确。\textbf{Bethe 拟设}（Bethe ansatz, 1931）是求解一维量子多体系统的核心方法，其基本思想是将多体波函数表示为平面波的叠加，其中系数由两体散射矩阵和 Bethe 方程决定。在热力学极限下，Bethe 根在复平面上凝聚为弦（strings），根密度满足线性积分方程，可通过 Wiener-Hopf 方法或数值方法求解。\textbf{随机 Loewner 演化}（SLE, Schramm 2000）是概率论与统计力学交叉的最新成果：SLE\(_\kappa\) 是一族随机共形映射，由一维 Brown 运动驱动，其随机曲线与二维临界统计力学模型（如渗流、Ising 模型、自规避随机游走）的界面具有相同的统计性质。SLE 的发现（Fields 奖 2006）统一了此前看似无关的多个临界现象，是 21 世纪概率论最重要的成就之一。

\hrulefill

\hrulefill

\section{第342章：非平衡态的概率描述}\label{ux7b2c342ux7ae0ux975eux5e73ux8861ux6001ux7684ux6982ux7387ux63cfux8ff0}

非平衡态统计力学研究的是系统远离热力学平衡时的动力学行为，其数学描述主要依赖三类概率工具：Markov 过程（主方程、细致平衡条件）、随机微分方程（Langevin 方程、Fokker-Planck 方程）和涨落定理（fluctuation theorems）。与平衡态统计力学（第 275-277 章）以 Gibbs 测度为核心的静力学框架不同，非平衡态理论必须面对熵产生、不可逆性和时间箭头等根本性问题。Boltzmann 的 H 定理（1872 年）首次从分子动力学推导出熵的单调非减性质，但其统计解释（Loschmidt 悖论和 Zermelo 回归悖论）引发了持续一个多世纪的哲学争论。现代涨落定理（Jarzynski 等式、Crooks 定理、Gallavotti-Cohen 定理）则在 1990 年代为这些争论提供了严格的概率论解答：熵产生为负的概率虽小但不为零，且满足精确的对称关系，将可逆微观动力学与不可逆宏观热力学统一在概率框架下。

\subsection{216.1 Boltzmann 方程与 H 定理}\label{boltzmann-ux65b9ux7a0bux4e0e-h-ux5b9aux7406}

\textbf{定义 216.1}（Boltzmann 方程，1872）：定义在相空间上的分布函数 \(f(x, v, t)\) 满足

\[\frac{\partial f}{\partial t} + v \cdot \nabla_x f = \left(\frac{\partial f}{\partial t}\right)_{\text{coll}}\]

其中碰撞项为

\[\left(\frac{\partial f}{\partial t}\right)_{\text{coll}} = \int_{\mathbb{R}^3} \int_{S^2} [f(v') f(v_*') - f(v) f(v_*)] B(v - v_*, \omega) \, d\omega \, dv_*\]

（Boltzmann 碰撞算子）。\(B\) 是碰撞核，由二体相互作用势决定。

Boltzmann 方程描述的是稀薄气体中分子速度分布函数的时间演化，其碰撞项 \((\partial f/\partial t)_{\text{coll}}\) 编码了分子二体碰撞的统计效应。Boltzmann 在此基础上定义了 H 泛函------类似于负熵的量------并发现了一个令人震惊的事实：H 泛函在碰撞算子的作用下单调非增，其最小值恰对应于 Maxwell-Boltzmann 平衡分布。这就是著名的 H 定理，它首次从分子动力学的视角为热力学第二定律提供了统计基础。

\textbf{定理 216.1}（Boltzmann H 定理，1872）：定义\textbf{H 泛函} \(H(t) = \int f(x, v, t) \log f(x, v, t) \, dx \, dv\)。则 \(dH/dt \leq 0\)，且 \(dH/dt = 0\) 当且仅当 \(f\) 是局部 Maxwell 分布：

\[f_{\text{eq}}(v) = \frac{n}{(2\pi \theta)^{3/2}} \exp\left(-\frac{|v - u|^2}{2\theta}\right)\]

其中 \(n, u, \theta\) 为分布参数。

\textbf{证明}：由 Boltzmann 方程，\(\frac{dH}{dt} = \int \frac{\partial f}{\partial t} (1 + \log f) dx dv = \int Q(f,f)(1 + \log f) dx dv\)，其中 \(Q(f,f)\) 为碰撞算子。将碰撞算子代入，利用碰撞的微观可逆性（\(B(v-v_*,\omega) = B(v_*-v,\omega)\)）和变量变换 \((v,v_*) \leftrightarrow (v',v_*')\)，得 \(\frac{dH}{dt} = -\frac{1}{4} \int (f' f_*' - f f_*) \log\frac{f' f_*'}{f f_*} B d\omega dv_* dv dx \leq 0\)。被积函数非负（因 \((x-y)\log(x/y) \geq 0\)），故 \(dH/dt \leq 0\)。等号成立当且仅当 \(f' f_*' = f f_*\) 几乎处处成立，即 \(\log f\) 为碰撞不变量：\(\log f = a + b \cdot v + c|v|^2\)。代入归一化条件得局部 Maxwell 分布。\(\blacksquare\)

H 定理确立了热力学第二定律在分子动力学层面的有效性，但 Boltzmann 本人并未解决一个更精细的问题：向平衡态的收敛有多快？这一问题（Cercignani 猜想）直到 2005 年才由 Desvillettes 和 Villani 通过对数 Sobolev 不等式和熵-熵产生方法给出了严格的指数收敛速率证明。

\textbf{定理 216.2}（Cercignani 猜想 / 收敛速率，Desvillettes-Villani 2005）：对 Maxwell 分子（\(B\) 仅依赖于散射角），Boltzmann 方程的解以指数速率收敛到 Maxwell 分布：\(H(t) - H_{\text{eq}} \leq C e^{-t/\tau}\)。

\textbf{证明}：对 Maxwell 分子，碰撞核 \(B\) 仅依赖于散射角，碰撞算子具有卷积结构。Bobylev 利用 Fourier 变换将 Boltzmann 方程化为 \(\partial_t \hat{f}(\xi) = \int_{S^2} [\hat{f}(\xi^+)\hat{f}(\xi^-) - \hat{f}(\xi)\hat{f}(0)] B(\theta) d\omega\)，其中 \(\xi^\pm = (\xi \pm |\xi|\omega)/2\)。定义熵产生泛函 \(D(f) = -\frac{dH}{dt}\)。Cercignani 猜想等价于 \(D(f) \geq \lambda (H(f) - H_{\text{eq}})\) 对某 \(\lambda > 0\) 成立。Desvillettes 和 Villani 通过证明对数 Sobolev 型不等式建立了该熵-熵产生不等式：利用 Landau 方程（Boltzmann 方程的扩散极限）的谱间隙估计和超压缩性，得到 \(\lambda = \inf_{f} D(f)/(H(f)-H_{\text{eq}}) > 0\)。由 Gronwall 不等式，\(H(t) - H_{\text{eq}} \leq (H(0)-H_{\text{eq}}) e^{-\lambda t}\)。\(\blacksquare\)

\subsection{216.2 涨落耗散定理}\label{ux6da8ux843dux8017ux6563ux5b9aux7406}

\textbf{定理 216.3}（涨落耗散定理，Callen-Welton 1951, Kubo 1957）：线性响应函数 \(\chi(t)\) 与平衡涨落的关系：

\[\chi(t) = -\beta \frac{d}{dt} \langle A(t) A(0) \rangle_{\text{eq}} \quad (t > 0)\]

其中 \(\langle A(t) A(0) \rangle\) 是可观测量 \(A\) 的平衡时间关联函数，\(\beta > 0\) 为模型参数。

\textbf{证明}：考虑外加微扰 \(H'(t) = -f(t) A\)，系统的线性响应为 \(\langle A(t) \rangle - \langle A \rangle_{\text{eq}} = \int_{-\infty}^t \chi(t-s) f(s) ds\)。由 Kubo 公式，响应函数 \(\chi(t) = \beta \langle \dot{A}(0) A(t) \rangle_{\text{eq}}\)（\(t>0\)）。利用时间平移不变性 \(\langle A(t) A(0) \rangle = \langle A(0) A(-t) \rangle\)，得 \(\langle \dot{A}(0) A(t) \rangle = -\frac{d}{dt} \langle A(0) A(t) \rangle\)。故 \(\chi(t) = -\beta \frac{d}{dt} \langle A(t) A(0) \rangle_{\text{eq}}\)。在频域中，\(\tilde{\chi}(\omega) = \int_0^\infty e^{i\omega t} \chi(t) dt\)，其虚部给出涨落耗散定理的谱形式：\(\Im \tilde{\chi}(\omega) = \frac{\beta\omega}{2} \tilde{C}_{AA}(\omega)\)，其中 \(\tilde{C}_{AA}\) 为关联函数的 Fourier 变换。\(\blacksquare\)

涨落耗散定理的数学核心在于：线性响应函数与平衡态时间关联函数之间的精确关系不是近似的，而是从 Liouville 方程的线性扰动推导出的严格结果。这一关系在输运理论中的最直接应用便是 Green-Kubo 公式：扩散常数、电导率、黏度等宏观输运系数可以直接从平衡系统的微观动力学关联函数计算出来，从而建立了不可逆过程与可逆微观力学之间的桥梁。

\textbf{定理 216.4}（Green-Kubo 公式，1954, 1957）：输运系数由平衡关联函数的时间积分给出：

\[D = \frac{1}{d} \int_0^\infty \langle v(t) \cdot v(0) \rangle_{\text{eq}} \, dt\]

\textbf{证明}：由 Einstein 关系，扩散系数 \(D = \lim_{t\to\infty} \frac{1}{2dt} \langle |x(t) - x(0)|^2 \rangle\)。将位移写为速度的积分 \(x(t) - x(0) = \int_0^t v(s) ds\)，则 \(\langle |x(t) - x(0)|^2 \rangle = \int_0^t \int_0^t \langle v(s) \cdot v(u) \rangle ds du\)。由平稳性，\(\langle v(s) \cdot v(u) \rangle = \langle v(s-u) \cdot v(0) \rangle\)。变量变换得 \(\langle |x(t)-x(0)|^2 \rangle = 2 \int_0^t (t-\tau) \langle v(\tau) \cdot v(0) \rangle d\tau\)。除以 \(2dt\) 并取 \(t \to \infty\)：若速度关联函数衰减足够快，\(\lim_{t\to\infty} \frac{1}{t} \int_0^t \tau \langle v(\tau) \cdot v(0) \rangle d\tau = 0\)，故 \(D = \frac{1}{d} \int_0^\infty \langle v(t) \cdot v(0) \rangle dt\)。这就是 Green-Kubo 公式，将输运系数表示为平衡涨落的时间关联。\(\blacksquare\)

Green-Kubo 公式直接蕴含了一个更深刻的对称性：Onsager 倒易关系。如果系统中同时存在多种不可逆过程（如热传导与电流），则它们之间的输运系数矩阵是高度对称的。Onsager 在 1931 年指出了这一对称性的根源：它是微观动力学时间反演对称性的宏观遗存，是热力学第二定律在近平衡线性区的精确补充。

\textbf{定理 216.5}（Onsager 倒易关系，1931）：输运系数矩阵 \(L_{ij}\) 对称：\(L_{ij} = L_{ji}\)。这是微观动力学可逆性的宏观推论。

\textbf{证明}：由 Green-Kubo 公式，\(L_{ij} = \int_0^\infty \langle \dot{A}_i(0) \dot{A}_j(t) \rangle_{\text{eq}} dt\)。微观可逆性要求 \(\langle \dot{A}_i(0) \dot{A}_j(t) \rangle_{\text{eq}} = \langle \dot{A}_i(t) \dot{A}_j(0) \rangle_{\text{eq}}\)（时间反演对称性）。由平稳性，\(\langle \dot{A}_i(t) \dot{A}_j(0) \rangle = \langle \dot{A}_i(0) \dot{A}_j(-t) \rangle = \langle \dot{A}_j(0) \dot{A}_i(t) \rangle\)（最后一步用时间反演对称性）。代入得 \(L_{ij} = \int_0^\infty \langle \dot{A}_j(0) \dot{A}_i(t) \rangle dt = L_{ji}\)。若流 \(J_i\) 在时间反演下为奇函数（如热流），则需引入符号因子 \(\varepsilon_i = \pm 1\)，Onsager 倒易关系推广为 \(L_{ij} = \varepsilon_i \varepsilon_j L_{ji}\)。\(\blacksquare\)

\subsection{216.3 涨落定理}\label{ux6da8ux843dux5b9aux7406}

\textbf{定理 216.6}（涨落定理，Evans-Cohen-Morriss 1993, Gallavotti-Cohen 1995）：在非平衡稳态中，熵产生率 \(\Sigma_t\) 满足

\[\frac{\mathbb{P}(\Sigma_t = A)}{\mathbb{P}(\Sigma_t = -A)} = e^{A t}\]

即正熵产生的概率指数地大于负熵产生的概率。

\textbf{证明}：考虑非平衡稳态的相空间轨迹 \(\Gamma_t\)。定义耗散函数 \(\Omega_t(\Gamma_0) = \int_0^t \Sigma(\Gamma_s) ds\)。由 Liouville 方程，相空间体积元满足 \(\frac{d}{dt} \log \frac{d\mu(\Gamma_t)}{d\mu(\Gamma_0)} = -\Lambda(\Gamma_t)\)，其中 \(\Lambda\) 为相空间压缩率。Gallavotti-Cohen 涨落定理基于混沌假设：耗散函数的大偏差速率函数满足对称性 \(I(-A) = I(A) + A\)。据此，\(\frac{1}{t} \log \frac{\mathbb{P}(\Sigma_t = A)}{\mathbb{P}(\Sigma_t = -A)} \to A\)。从微观可逆性出发：设时间反演映射为 \(M: \Gamma \mapsto \Gamma^*\)，则 \(\mathbb{P}(\Sigma_t = -A) = \int \delta(\Sigma_t + A) d\mu = \int \delta(\Sigma_t(\Gamma^*) - A) \exp(-\Omega_t(\Gamma)) d\mu(\Gamma)\)。由 \(\Omega_t(\Gamma) = t\Sigma_t\) 在稳态中成立，得 \(\mathbb{P}(\Sigma_t = -A) = e^{-A t} \mathbb{P}(\Sigma_t = A)\)。\(\blacksquare\)

涨落定理的核心信息在于：熵在短时间内可以自发减少------即存在负熵产生事件------但其概率随观察时间指数衰减。这一结果为理解非平衡系统的不可逆性提供了严格的数学表述。Jarzynski 在此基础上将涨落定理推广到一个极为实用的形式：即便在远离平衡的驱动过程中，非平衡功的指数平均仍然精确等于自由能差的指数函数，从而提供了一种从非平衡测量中提取平衡热力学信息的通用框架。

\textbf{定理 216.7}（Jarzynski 等式，1997）：对非平衡过程，功 \(W\) 的分布满足

\[\langle e^{-\beta W} \rangle = e^{-\beta \Delta F}\]

其中 \(\Delta F\) 是自由能差。该等式允许从非平衡观测中提取平衡态信息。

\textbf{证明}：考虑系统从初始平衡态 \(A\) 经外部驱动到达终态 \(B\) 的非平衡过程。初始分布为 \(\rho_A(\Gamma) = e^{-\beta H_A(\Gamma)}/Z_A\)。功 \(W\) 定义为沿轨迹的总能量变化：\(W = H_B(\Gamma_\tau) - H_A(\Gamma_0)\)。计算 \(\langle e^{-\beta W} \rangle = \int e^{-\beta (H_B(\Gamma_\tau) - H_A(\Gamma_0))} \rho_A(\Gamma_0) d\Gamma_0\)。由 Liouville 定理，相空间体积元守恒 \(d\Gamma_0 = d\Gamma_\tau\)。代入 \(\rho_A\) 得 \(\langle e^{-\beta W} \rangle = \int e^{-\beta H_B(\Gamma_\tau)} \frac{1}{Z_A} d\Gamma_\tau = \frac{Z_B}{Z_A} = e^{-\beta(F_B - F_A)} = e^{-\beta \Delta F}\)。这一推导的关键在于 Liouville 定理：虽然系统在非平衡驱动下演化，相空间体积形式上不变。\(\blacksquare\)

\subsection{216.4 Fokker-Planck 方程与 Langevin 动力学}\label{fokker-planck-ux65b9ux7a0bux4e0e-langevin-ux52a8ux529bux5b66}

\textbf{定义 216.2}（Fokker-Planck 方程）：\textbf{Fokker-Planck 方程}（Smoluchowski 方程）描述扩散过程 \(X_t\) 的概率密度 \(p(x, t)\) 的演化：

\[\frac{\partial p}{\partial t} = -\frac{\partial}{\partial x}[\mu(x) p] + \frac{1}{2} \frac{\partial^2}{\partial x^2}[D(x) p]\]

其中 \(\mu(x)\) 是漂移系数，\(D(x)\) 是扩散系数。

Fokker-Planck 方程是概率密度演化的偏微分方程描述，而 Langevin 方程则从随机微分方程的角度描述单个粒子轨道的随机动力学。两者之间的等价性------即 Ito 或 Stratonovich 随机微分方程的解的边际分布恰好满足 Fokker-Planck 方程------构成了连接随机过程微观描述与宏观概率分布的数学桥梁。

\textbf{定理 216.8}（Fokker-Planck 与 Langevin 的等价性）：Langevin 方程 \(dX_t = \mu(X_t) dt + \sigma(X_t) dW_t\) 的概率密度满足 Fokker-Planck 方程（\(\mu\) 和 \(D = \sigma^2\) 的对应关系由 Stratonovich 或 Itô 积分解释确定）。

\textbf{证明}：考虑 Itô 解释的 Langevin 方程。对任意光滑测试函数 \(\varphi\)，由 Itô 公式，\(d\varphi(X_t) = [\mu(X_t)\varphi'(X_t) + \frac{1}{2}\sigma^2(X_t)\varphi''(X_t)] dt + \sigma(X_t)\varphi'(X_t) dW_t\)。取期望，鞅部分消失：\(\frac{d}{dt}\mathbb{E}[\varphi(X_t)] = \mathbb{E}[\mu(X_t)\varphi'(X_t) + \frac{1}{2}\sigma^2(X_t)\varphi''(X_t)]\)。设 \(p(x,t)\) 为 \(X_t\) 的密度，则左边为 \(\int \varphi(x) \partial_t p(x,t) dx\)，右边为 \(\int [\mu(x)\varphi'(x) + \frac{1}{2}\sigma^2(x)\varphi''(x)] p(x,t) dx\)。分部积分，对任意 \(\varphi\) 成立，故 \(\partial_t p = -\partial_x(\mu p) + \frac{1}{2}\partial_{xx}(\sigma^2 p)\)。若采用 Stratonovich 解释，漂移项替换为 \(\mu \to \mu + \frac{1}{2}\sigma\sigma'\)。\(\blacksquare\)

Fokker-Planck 方程最经典的应用之一是 Kramers 逃逸问题：热噪声驱动下，粒子如何越过势垒从亚稳态逃逸到稳定态。Kramers 在 1940 年给出了平均首次通过时间的渐近公式，揭示了热涨落协助下的势垒跨越行为遵循 Arrhenius 律------平均逃逸时间随势垒高度 \(\Delta E\) 和逆温度 \(\beta\) 指数增长。这一结果不仅是化学反应速率理论的基础，也是统计力学中罕见可解的非平衡动力学问题的典范。

\textbf{定理 216.9}（Kramers 逃逸问题，1940）：Fokker-Planck 方程中，粒子从势阱中逃逸的平均首次通过时间为

\[\tau \sim e^{\beta \Delta E}\]

其中 \(\Delta E\) 是势垒高度。

\textbf{证明}：考虑势函数 \(U(x)\) 中的过阻尼 Langevin 动力学 \(\gamma dX_t = -U'(X_t) dt + \sqrt{2\gamma/\beta} dW_t\)。对应的 Fokker-Planck 方程为 \(\partial_t p = \frac{1}{\gamma}\partial_x(U' p) + \frac{1}{\beta\gamma}\partial_{xx} p\)。稳态分布为 \(p_{\text{eq}}(x) \propto e^{-\beta U(x)}\)。平均首次通过时间 \(\tau\) 满足后向 Kolmogorov 方程 \(\mathcal{L}^\dagger \tau = -1\)，其中 \(\mathcal{L}^\dagger = -\frac{1}{\gamma}U'(x)\partial_x + \frac{1}{\beta\gamma}\partial_{xx}\)。解为 \(\tau(x) = \beta\gamma \int_x^b e^{\beta U(y)} dy \int_a^y e^{-\beta U(z)} dz\)。对于双稳态势 \(U(x)\)，势垒 \(\Delta E = U(x_{\text{max}}) - U(x_{\text{min}})\)。在 \(\beta\Delta E \gg 1\) 时，Laplace 方法给出 \(\tau \sim \frac{2\pi\gamma}{\sqrt{U''(x_{\min})|U''(x_{\max})|}} e^{\beta\Delta E}\)，即 Arrhenius 律。\(\blacksquare\)

\hrulefill

\emph{本卷建立了统计力学的核心数学框架：Gibbs 测度与 DLR 理论将平衡态问题放在严格的概率论基础上，相变表现为热力学极限下 Gibbs 测度的非唯一性；Ising 模型的 Peierls 论证利用组合几何方法严格证明了测度非唯一性；重整化群提供了多尺度系统的统一数学结构（参数空间上的半群作用、不动点的谱分析、\(\varepsilon\)-展开）；精确可解模型（Onsager 解的 Clifford 代数与 Toeplitz 行列式、Bethe 拟设与 Yang-Baxter 方程）给出了强关联系统的非微扰精确结果；非平衡态的概率描述（Boltzmann 方程、涨落耗散定理、涨落定理）将不可逆性纳入概率论框架。统计力学的数学结构是概率论、泛函分析和算子代数的重要交汇领域。}

\hrulefill

\hrulefill

\hrulefill

\hrulefill

\hrulefill

\hrulefill

\hrulefill

\subsection{279.A 统计力学补充2}\label{a-ux7edfux8ba1ux529bux5b66ux8865ux51452}

\begin{quote}
随机矩阵理论（Random Matrix Theory / RMT）研究元素为随机变量的矩阵的谱性质（特征值分布）。它起源于 Wishart（1928）在多元统计中的工作和 Wigner（1955）在核物理中的应用，现已发展成连接数论（Riemann \(\zeta\) 函数零点）、高维统计、机器学习、无线通信、量子混沌和自由概率论的统一数学框架。Wigner 半圆律是「随机矩阵的普遍性」的典范------大量随机矩阵的宏观特征值分布不依赖于微观分布细节。本卷在概率论 II（V21）、线性代数（V7）和泛函分析（V14）的基础上，建立 Wigner 半圆律、自由概率论、Tracy-Widom 分布和随机矩阵现代应用的理论体系。
\end{quote}

\hrulefill

\hrulefill

\hrulefill

\hrulefill

\hrulefill

\newpage

\chapter{10-概率与统计/01-概率论/04-随机矩阵理论.md}\label{ux6982ux7387ux4e0eux7edfux8ba101-ux6982ux7387ux8bba04-ux968fux673aux77e9ux9635ux7406ux8bba.md}

\section{第343章：Wigner 半圆律}\label{ux7b2c343ux7ae0wigner-ux534aux5706ux5f8b}

Wigner 半圆律（1955-1958 年）是随机矩阵理论的奠基性定理，其历史背景植根于核物理学------Wigner 在 1950 年代为描述重原子核的能级统计而引入随机矩阵模型，发现大维度下 Hermitian 随机矩阵的经验特征值分布趋向于半圆形。这一定理的深远意义在于它揭示了随机矩阵的''普遍性''（universality）：无论矩阵元素的具体分布如何，只要满足基本的矩条件，极限谱分布就由半圆律唯一确定。从方法论角度看，Wigner 的矩法证明是组合数学与概率论的经典结合------将矩阵的迹展开为指标配对，Catalan 数自然地作为''非交叉配对''的计数出现，从而将随机矩阵的谱问题转化为组合计数问题。这一方法为后续的自由概率论（第 281 章 Voiculescu）奠定了组合学基础------半圆律在自由概率论中是自由中心极限定理的极限分布。本章还引入 Wishart 矩阵和 Marchenko-Pastur 律，将谱理论拓展至高维样本协方差矩阵的分析。

\subsection{267.1 经典 Wigner 矩阵}\label{ux7ecfux5178-wigner-ux77e9ux9635}

\textbf{定义 267.1}（Wigner 矩阵）：\textbf{Wigner 矩阵} \(H = (h_{ij})_{1 \leq i,j \leq N}\) 是 \(N \times N\) 随机 Hermitian（或实对称）矩阵，满足： - 对角线元素 \(h_{ii}\) 是 i.i.d. 实随机变量（均值 \(0\)，方差 \(\sigma^2\)） - 非对角线元素 \(h_{ij}\)（\(i < j\)）是 i.i.d. 复（或实）随机变量（均值 \(0\)，方差 \(1/N\)，独立于对角线元素） - \(h_{ji} = \overline{h_{ij}}\)

\textbf{定义 267.2}（经验谱分布 / ESD）：矩阵 \(H\) 的\textbf{经验谱分布}（Empirical Spectral Distribution）为经验测度 \[\mu_H = \frac{1}{N} \sum_{i=1}^N \delta_{\lambda_i}\]

其中 \(\lambda_1 \leq \cdots \leq \lambda_N\) 是 \(H\) 的特征值。

\textbf{定理 267.1}（Wigner 半圆律，1955-1958）：设 \(H\) 是 Wigner 矩阵（矩阵元素方差适当缩放）。以概率 \(1\)（几乎必然），当 \(N \to \infty\) 时，经验谱分布 \(\mu_H\) 弱收敛到半圆律： \[d\mu_{sc}(x) = \frac{1}{2\pi} \sqrt{4 - x^2} \, \mathbf{1}_{[-2, 2]}(x) \, dx\]

半圆律的 Stieltjes 变换为 \(s(z) = \frac{-z + \sqrt{z^2 - 4}}{2}\)（\(z \in \mathbb{C}^+\)），满足方程 \(s(z) + \frac{1}{z + s(z)} = 0\)。

\textbf{证明}（矩法）：对 \(k\) 为偶数（奇次矩为零，由对称性），计算 \(\mathbb{E}[\frac{1}{N} \operatorname{Tr}(H^k)] = \frac{1}{N} \sum_{i_1, \ldots, i_k=1}^N \mathbb{E}[h_{i_1 i_2} h_{i_2 i_3} \cdots h_{i_k i_1}]\)。将期望展开，每个项对应一个闭合的指标序列 \((i_1, i_2, \ldots, i_k, i_1)\)（可视作边标记图）。由于矩阵元素独立且均值为零，期望非零仅当每个 \(h_{ij}\) 在乘积中出现偶数次------即指标序列中的每条边出现恰好两次，形成''配对''（pairing）。在 \(N \to \infty\) 极限下，主导贡献来自不会产生交叉的\textbf{非交叉配对}（non-crossing / planar pairings）------因为交叉配对导致独立指标求和时每多一个交叉损失因子 \(1/N\)。\(k\) 个顶点的非交叉配对数目恰由 Catalan 数 \(C_{k/2} = \frac{1}{k/2+1}\binom{k}{k/2}\) 给出。因此 \(\lim_{N \to \infty} \mathbb{E}[\frac{1}{N} \operatorname{Tr}(H^k)] = \begin{cases} C_{k/2}, & k \text{ 为偶数} \\ 0, & k \text{ 为奇数} \end{cases}\)。Catalan 数正是半圆律矩序列。由矩方法（moment method），矩唯一确定紧支撑概率分布，方差估计进一步证明几乎必然收敛到半圆律。\(\blacksquare\)

\textbf{定理 267.2}（半圆律的普遍性 / Universality）：Wigner 半圆律对矩阵元素的分布不敏感------只要元素满足有限的矩条件（如四阶矩有限），极限谱分布就是半圆的。这是随机矩阵理论「普遍性」（universality）的第一个表现。

\textbf{证明}：设 \(H\) 为 Wigner 矩阵，其元素分布满足 \(\mathbb{E}[h_{ij}] = 0\)，\(\mathbb{E}[|h_{ij}|^2] = 1/N\)，且 \(\mathbb{E}[|h_{ij}|^4] \leq C/N^2\)。在矩法计算 \(\mathbb{E}[\operatorname{Tr}(H^k)]/N\) 中，非交叉配对的贡献为 \(C_{k/2}\)（与分布无关）。交叉配对的贡献涉及元素的高阶矩，但每多一个交叉配对，自由指标数减少 \(2\)，导致贡献被 \(O(1/N^2)\) 压制。在 \(N \to \infty\) 下，交叉配对对期望矩的贡献趋于零，极限矩序列仅依赖于非交叉配对数目（Catalan 数），与元素的具体分布无关。方差估计 \(\operatorname{Var}(\operatorname{Tr}(H^k)/N) = O(1/N^2)\) 证明几乎必然收敛。因此只要二阶矩和四阶矩有限，极限谱分布为半圆律。\(\blacksquare\)

\subsection{267.2 矩法与自由独立}\label{ux77e9ux6cd5ux4e0eux81eaux7531ux72ecux7acb}

\textbf{定义 267.3}（Catalan 数与半圆律的矩）：半圆律的 \(2k\) 阶矩（奇阶矩为零）为\textbf{Catalan 数} \(C_k\)： \[\int x^{2k} d\mu_{sc}(x) = C_k = \frac{1}{k+1} \binom{2k}{k}\]

\(k=1\) 时 \(C_1=1\)，\(k=2\) 时 \(C_2=2\)，\(k=3\) 时 \(C_3=5\)。Catalan 数计数了「非交叉配对」的数量。

Catalan 数不仅给出了半圆律的矩序列，更揭示了矩法证明中一个关键的组合事实：在 \(N \to \infty\) 极限下，指标序列的配对必须是非交叉的。交叉配对在独立指标求和时每多一个交叉则损失一个因子 \(1/N\)，从而被渐近压制。这一观察是矩法证明的核心技术环节，也是自由概率论中非交叉分割概念的前身。

\textbf{定理 267.3}（矩法的核心估计）：在 Wigner 矩阵中，\(\mathbb{E}[\operatorname{Tr}(H^k)]/N\) 的主导贡献来自矩阵指标序列的「非交叉配对」（non-crossing / planar pairings），交叉配对的贡献被 \(1/N^2\) 压制。

\textbf{证明}：将 \(\mathbb{E}[\operatorname{Tr}(H^k)]\) 展开为 \(\sum_{i_1,\ldots,i_k} \mathbb{E}[h_{i_1 i_2} \cdots h_{i_k i_1}]\)。每个非零项对应一个配对 \(\pi\)（将 \(k\) 条边配成 \(k/2\) 对）。设 \(V(\pi)\) 为配对 \(\pi\) 中不同顶点的数目。配对的贡献为 \((1/N)^{k/2} \cdot N^{V(\pi)}\)（\(N^{V(\pi)}\) 来自独立指标求和）。对于非交叉配对，\(V(\pi) = k/2 + 1\)，故贡献为 \(N\)。对于有 \(c\) 个交叉的配对，\(V(\pi) = k/2 + 1 - c\)，贡献为 \(N^{1-c}\)。除以 \(N\) 后，非交叉配对贡献 \(O(1)\)，单交叉配对贡献 \(O(1/N)\)，多交叉配对贡献 \(O(1/N^2)\) 或更高阶。在 \(N \to \infty\) 下，仅有非交叉配对的贡献存留。\(\blacksquare\)

\subsection{267.3 Wishart 矩阵与 Marchenko-Pastur 律}\label{wishart-ux77e9ux9635ux4e0e-marchenko-pastur-ux5f8b}

\textbf{定义 267.4}（Wishart / 样本协方差矩阵）：\(X\) 是 \(p \times n\) 矩阵，元素为 i.i.d.（均值 \(0\)，方差 \(1\)）。\textbf{样本协方差矩阵}为 \(S = \frac{1}{n} X X^T\)（\(p \times p\)）。

\textbf{定理 267.4}（Marchenko-Pastur 律，1967）：设 \(p, n \to \infty\) 且 \(p/n \to \gamma \in (0, \infty)\)。\(S\) 的经验谱分布收敛到 \textbf{Marchenko-Pastur 分布}： \[d\mu_{MP}(x) = \frac{\sqrt{(x - a_-)(a_+ - x)}}{2\pi \gamma x} \mathbf{1}_{[a_-, a_+]}(x) dx + (1 - \gamma^{-1})_+ \delta_0(dx)\]

其中 \(a_\pm = (1 \pm \sqrt{\gamma})^2\)。

\textbf{证明}：设 \(S = \frac{1}{n} X X^\top\)，其中 \(X \in \mathbb{R}^{p \times n}\) 元素为 i.i.d. 均值为 \(0\)，方差为 \(1\)。\(S\) 的 Stieltjes 变换 \(s_N(z) = \frac{1}{p} \operatorname{Tr}((S - zI)^{-1})\) 满足自洽方程。利用矩阵恒等式 \((S - zI)^{-1} = -z^{-1}I + z^{-1}S(S - zI)^{-1}\) 和随机矩阵理论中的迹等式，在 \(p, n \to \infty\)，\(p/n \to \gamma\) 下，\(s(z) = \lim s_N(z)\) 满足 \(z s(z)^2 + (z + 1 - \gamma) s(z) + 1 = 0\)。解得 \(s(z) = \frac{-(z+1-\gamma) + \sqrt{(z+1-\gamma)^2 - 4z}}{2z}\)。由 Stieltjes 反演公式，\(\rho(x) = \frac{1}{\pi} \lim_{\epsilon \to 0^+} \Im s(x+i\epsilon)\)。当 \(x \in [a_-, a_+]\) 时，平方根为虚数，给出 \(\rho(x) = \frac{\sqrt{(x-a_-)(a_+-x)}}{2\pi\gamma x}\)。当 \(\gamma > 1\) 时，\(s(z)\) 在 \(z=0\) 处有极点，对应 Dirac 质量 \((1-\gamma^{-1})\delta_0\)。\(\blacksquare\)

\hrulefill

\hrulefill

\hrulefill

\hrulefill

\emph{卷二十二：概率论终。} \hrulefill

\section{第344章：自由概率论}\label{ux7b2c344ux7ae0ux81eaux7531ux6982ux7387ux8bba}

自由概率论（Free Probability Theory）由 Dan Voiculescu 于 1980 年代创立，其核心动机在于解决算子代数中的 von Neumann 因子同构分类问题，但随后被发现是随机矩阵在 \(N \to \infty\) 极限下的自然概率框架。自由概率论将经典概率论中的独立性替换为''自由独立性''（freeness）------一种严格非交换的独立概念，其组合学由非交叉分割（non-crossing partitions）而非经典的所有分割来描述。这一理论的优美之处在于：独立 GUE 随机矩阵族在 \(N \to \infty\) 时渐近自由，其极限谱分布由自由卷积（free convolution）和 \(R\)-变换描述，而半圆律在自由概率论中恰恰扮演了与 Gauss 分布在经典概率论中相同的角色------它是自由中心极限定理的极限。本章在 Wigner 半圆律（第 280 章）的矩法基础上，系统建立非交换概率空间、自由累积量、自由卷积和 \(R\)-变换的理论，并与随机矩阵的渐近自由性质建立桥梁。

\subsection{268.1 非交换概率空间}\label{ux975eux4ea4ux6362ux6982ux7387ux7a7aux95f4-1}

\textbf{定义 268.1}（非交换概率空间 / \(C^*\)-概率空间）：非交换概率空间是 \((\mathcal{A}, \tau)\)，其中 \(\mathcal{A}\) 是单位 \(C^*\)-代数，\(\tau: \mathcal{A} \to \mathbb{C}\) 是单位保迹的线性泛函，满足 \(\tau(1) = 1\) 且 \(\tau(a^*a) \geq 0\)（正性）。

在这个代数框架中，核心的代换是将经典概率论中的''独立性''替换为一种适合非交换算子的新概念。经典独立性要求联合矩分解为边缘矩的乘积，但在算子代数中这种定义不再自然。Voiculescu 的洞察在于用迹的消失条件来刻画一种严格的非交换独立关系。

\textbf{定义 268.2}（自由独立 / Free Independence）：\(\mathcal{A}\) 的子代数族 \(\{\mathcal{A}_i\}_{i \in I}\) 是\textbf{自由独立的}，如果对任意满足 \(\tau(a_k) = 0\) 的非交换多项式 \(a_1 \in \mathcal{A}_{i_1}, \ldots, a_k \in \mathcal{A}_{i_k}\)（相邻指标不同 \(i_1 \neq i_2 \neq \cdots \neq i_k\)），有 \(\tau(a_1 a_2 \cdots a_k) = 0\)。

自由独立的这一定义之所以重要，是因为它与随机矩阵的极限行为之间存在精确的对应关系。独立 Wigner 矩阵族虽然在经典意义下独立，但其非交换结构使得经典独立性不再适用。Voiculescu 在 1991 年的突破性发现指出，在大 \(N\) 极限下，独立随机矩阵族恰好满足自由独立条件。

\textbf{定理 268.1}（自由独立 = 独立矩阵的极限行为，Voiculescu 1991）：若 \(A_N\) 和 \(B_N\) 是独立（经典意义上的独立）的 Wigner 矩阵序列，则它们在 \(N \to \infty\) 极限下是自由独立的（在非交换概率空间中）。

\textbf{证明}：设 \(A_N, B_N\) 为独立 Wigner 矩阵，\(\tau_N(X) = \frac{1}{N}\mathbb{E}[\operatorname{Tr}(X)]\)。要证对任意满足 \(\tau(a_k) = 0\) 的交替乘积，\(\tau(a_1 b_1 a_2 b_2 \cdots) = 0\)。由矩法，展开 \(\operatorname{Tr}(A_N^{p_1} B_N^{q_1} A_N^{p_2} \cdots)\) 为指标序列的求和。由于 \(A_N\) 和 \(B_N\) 独立，期望分解为 \(A_N\) 元素和 \(B_N\) 元素的期望的乘积。在 \(N \to \infty\) 下，主导配对必须是非交叉的。但交替序列的非交叉配对必然将某个 \(A_N\) 元素与 \(B_N\) 元素配对，而它们独立且均值为零，故期望为零。因此交替乘积以零为极限，即自由独立条件成立。\(\blacksquare\)

\subsection{268.2 自由卷积与 R-变换}\label{ux81eaux7531ux5377ux79efux4e0e-r-ux53d8ux6362}

\textbf{定义 268.3}（自由加法卷积 / Free Additive Convolution \textbf{\(\boxplus\)}）：若 \(a, b\) 自由独立，其边缘分布的加法卷积 \(\mu_a \boxplus \mu_b\) 定义为 \(a+b\) 的分布。

\textbf{定义 268.4}（R-变换 / \(R\)-transform，Voiculescu 1986）：对概率测度 \(\mu\)，定义其 Cauchy 变换 \(G_\mu(z) = \int \frac{d\mu(x)}{z - x}\)（\(z \in \mathbb{C}^+\)）。\(K_\mu(z) = G_\mu^{-1}(z)\)（函数逆）。\textbf{R-变换}为 \[R_\mu(z) = K_\mu(z) - \frac{1}{z}\]

\textbf{定理 268.2}（R-变换的线性化性质）：若 \(\mu = \mu_1 \boxplus \mu_2\)，则 \[R_\mu(z) = R_{\mu_1}(z) + R_{\mu_2}(z)\]

即\textbf{自由卷积对应 R-变换的加法}------这是自由概率论的核心结果，类比经典卷积对应特征函数的乘积。对数 Fourier 变换的角色在此被 R-变换取代。

\textbf{证明}：设 \(a, b\) 自由独立，其分布为 \(\mu_a, \mu_b\)。\(a+b\) 的 Cauchy 变换 \(G_{a+b}(z) = \tau((z - a - b)^{-1})\)。由算子值自由概率论中的子ordination性质，\(G_{a+b}(z) = G_a(z - R_b(G_{a+b}(z)))\)。取函数逆：\(K_{a+b}(z) = K_a(z) + K_b(z) - 1/z\)。代入 \(R(z) = K(z) - 1/z\)，得 \(R_{a+b}(z) = R_a(z) + R_b(z)\)。此线性性源于自由独立的矩-累积量关系：在自由独立下，交替累积量消失，故 \(a+b\) 的自由累积量 \(\kappa_n^{a+b} = \kappa_n^a + \kappa_n^b\)。R-变换为自由累积量的生成函数 \(R(z) = \sum_{n \geq 0} \kappa_{n+1} z^n\)，因此线性性自然成立。\(\blacksquare\)

\subsection{268.3 自由乘法卷积与 S-变换}\label{ux81eaux7531ux4e58ux6cd5ux5377ux79efux4e0e-s-ux53d8ux6362}

\textbf{定义 268.5}（自由乘法卷积 / Free Multiplicative Convolution \textbf{\(\boxtimes\)}）：对正算子 \(a, b\) 自由独立，\(\mu_a \boxtimes \mu_b\) 定义为 \(a^{1/2} b a^{1/2}\)（或等价的 \(ab\)）的分布。

\textbf{定义 268.6}（S-变换 / S-transform）：对具有非零均值的正测度 \(\mu\)，\(S_\mu(z)\) 满足 \(\chi_\mu(z S_\mu(z)) = z\)，其中 \(\chi_\mu\) 是矩母函数的相关函数。自由乘法卷积对应 S-变换的乘法：\(S_{\mu_1 \boxtimes \mu_2}(z) = S_{\mu_1}(z) S_{\mu_2}(z)\)。

S-变换的构造涉及矩母函数 \(M_\mu(z) = \sum_{n \geq 1} m_n z^n\) 及其逆函数。定义 \(\psi_\mu(z) = \sum_{n \geq 0} m_{n+1} z^n\)（移位矩母函数），则 \(\chi_\mu(z) = z \psi_\mu^{-1}(z)\)，而 \(S_\mu(z) = \frac{1+z}{z \chi_\mu^{-1}(z)}\)。自由乘法卷积的乘法性质 \(S_{\mu \boxtimes \nu} = S_\mu \cdot S_\nu\) 是自由概率论中与经典独立随机变量乘积的 Mellin 变换可类比的结果。与加法卷积不同，乘法卷积要求算子为正（或至少具有非零均值），这限制了其应用范围。S-变换在随机矩阵乘积的渐近特征值分布中扮演核心角色，例如独立 Wigner 矩阵乘积的极限分布可通过迭代 S-变换计算。自由乘法卷积与自由加法卷积通过指数映射相联系：\(\mu \boxtimes \nu = \exp(\log \mu \boxplus \log \nu)\)（在适当的测度变换下），这提供了两种卷积之间的桥梁。

\subsection{268.4 自由熵与自由 Fisher 信息}\label{ux81eaux7531ux71b5ux4e0eux81eaux7531-fisher-ux4fe1ux606f}

\textbf{定义 268.7}（自由熵 / Free Entropy，Voiculescu 1993-1998）：非交换随机变量 \(X\) 的\textbf{自由熵} \(\chi(X)\) 类比于经典 Shannon 微分熵，定义为近似矩阵模型的「微观态计数」的大偏差极限。自由熵满足自由信息论中的最大化原理：自由半圆变量最大化自由熵（给定方差约束）。

自由熵 \(\chi(X_1, \ldots, X_n)\) 被定义为用矩阵模型逼近非交换随机变量组时，矩阵维数的指数增长率。精确地说，设 \(\Gamma(X_1, \ldots, X_n; N, \varepsilon)\) 为满足 \(|\tau(X_{i_1} \cdots X_{i_k}) - \tau_N(A_{i_1}^{(N)} \cdots A_{i_k}^{(N)})| < \varepsilon\) 的 \(N \times N\) 自伴矩阵组 \((A_1^{(N)}, \ldots, A_n^{(N)})\) 的 Lebesgue 测度，则 \(\chi = \lim_{\varepsilon \to 0} \limsup_{N \to \infty} (\frac{1}{N^2} \log \Gamma + \frac{n}{2} \log N)\)。自由 Fisher 信息 \(\Phi^*(X_1, \ldots, X_n)\) 则通过自由扩散过程的导数定义，满足自由 Cramer-Rao 不等式：\(\Phi^* \cdot \chi \geq n\)（类比经典 Fisher 信息与熵的 Cramer-Rao 界）。自由熵是 Voiculescu 在解决自由群因子同构问题过程中发展的关键工具------它作为 II\(_1\) 因子的不变量，能够区分不同秩的自由群 von Neumann 代数。

\hrulefill

\hrulefill

\hrulefill

\hrulefill

\hrulefill

\hrulefill

\section{第345章：Tracy-Widom 分布与边缘统计}\label{ux7b2c345ux7ae0tracy-widom-ux5206ux5e03ux4e0eux8fb9ux7f18ux7edfux8ba1}

Tracy-Widom 分布（1994 年）是随机矩阵理论在 20 世纪末最重要的发现之一，它在统计物理学中的地位堪比正态分布在经典概率论中的地位。当 Wigner 半圆律描述了随机矩阵特征值的''体''（bulk）统计时，Tracy-Widom 分布 \(F_\beta\)（\(\beta = 1, 2, 4\) 分别对应 GOE、GUE、GSE）描述了最大特征值在缩放极限下的涨落分布。令人惊叹的是，Tracy-Widom 分布的累积分布函数由 Painlevé II 方程的解表达，将随机矩阵的边缘统计与完全可积的非线性微分方程深刻地耦合------这是随机矩阵理论中最优美的分析结果之一。Tracy-Widom 分布的普遍性远超其发源地：它在 KPZ 普适性类（界面生长）、组合学（最长递增子序列的 Baik-Deift-Johansson 定理）和高维统计（主成分分析的 Johnstone 检验）中反复出现。本章从 Gauss 系综（GUE、GOE）的联合特征值分布出发，引出 Tracy-Widom 分布的 Airy 核表征，并证明最大特征值的缩放极限。

\subsection{269.1 Gaussian 系综}\label{gaussian-ux7cfbux7efc}

\textbf{定义 269.1}（Gaussian 系综 / GOE、GUE、GSE）： - \textbf{GOE}（Gaussian Orthogonal Ensemble）：实对称矩阵 \(H = H^T\)，\(h_{ii} \sim \mathcal{N}(0, 2)\)，\(h_{ij} \sim \mathcal{N}(0, 1)\)（\(i < j\)） - \textbf{GUE}（Gaussian Unitary Ensemble）：复 Hermitian 矩阵 \(H = H^\dagger\)，\(\Re h_{ij}, \Im h_{ij} \sim \mathcal{N}(0, 1/2)\) - \textbf{GSE}（Gaussian Symplectic Ensemble）：描述自旋-轨道耦合系统

\textbf{定理 269.1}（GUE 的联合特征值分布）：GUE 矩阵的特征值联合概率密度为

\[p(\lambda_1, \ldots, \lambda_N) = \frac{1}{Z_N} \prod_{i < j} |\lambda_i - \lambda_j|^2 \prod_{i=1}^N e^{-\lambda_i^2 / 2}\]

其中 \(|\lambda_i - \lambda_j|^2\) 是 Vandermonde 行列式的平方------表示特征值之间存在「排斥」效应（level repulsion）。

\textbf{证明}：GUE 矩阵 \(H\) 的概率密度为 \(P(H) \propto \exp(-\frac{1}{2}\operatorname{Tr}(H^2))\)。将 \(H\) 对角化：\(H = U \Lambda U^\dagger\)，其中 \(\Lambda = \operatorname{diag}(\lambda_1,\ldots,\lambda_N)\)，\(U\) 为酉矩阵。变量变换的 Jacobian 为 \(\prod_{i<j} |\lambda_i - \lambda_j|^2\)（由 Weyl 积分公式）。对 \(U\) 积分（Haar 测度上的积分为常数），即得特征值的联合密度 \(p(\lambda_1,\ldots,\lambda_N) \propto \prod_{i<j} |\lambda_i - \lambda_j|^2 \exp(-\frac{1}{2}\sum_i \lambda_i^2)\)。Vandermonde 行列式 \(\Delta(\lambda) = \prod_{i<j} (\lambda_i - \lambda_j)\) 的平方产生特征值排斥------当 \(\lambda_i \approx \lambda_j\) 时密度趋于零，说明特征值不会聚集。\(\blacksquare\)

\subsection{269.2 最大特征值与 Tracy-Widom 分布}\label{ux6700ux5927ux7279ux5f81ux503cux4e0e-tracy-widom-ux5206ux5e03}

\textbf{定理 269.2}（Tracy-Widom 分布，1994-1996）：设 \(\lambda_{\max}\) 是 GUE（或 GOE/GSE）矩阵的最大特征值。经过适当缩放（中心化和标准化）后：

\[\frac{\lambda_{\max} - 2\sqrt{N}}{N^{-1/6}} \xrightarrow{d} F_\beta\]

其中 \(F_\beta\) 是\textbf{Tracy-Widom 分布}（\(\beta = 1\) 为 GOE，\(\beta = 2\) 为 GUE，\(\beta = 4\) 为 GSE）。\(F_2\) 由 Painlevé II 方程的解定义： \[F_2(s) = \exp\left( -\int_s^\infty (x - s) q^2(x) dx \right)\]

其中 \(q(x)\) 满足 \(\ddot{q} = x q + 2 q^3\)（Painlevé II），边界条件 \(q(x) \sim \operatorname{Ai}(x)\)（\(x \to \infty\)）。

\textbf{证明}：设 \(\lambda_{\max}\) 为 GUE 的最大特征值。由行列式点过程结构，\(\mathbb{P}(\lambda_{\max} \leq s) = \det(I - K_N|_{(s,\infty)})\)，其中 \(K_N\) 为 Hermite 核。在边缘缩放 \(x = 2\sqrt{N} + s N^{-1/6}\) 下，Hermite 核收敛到 Airy 核 \(K_{\text{Airy}}(x,y) = \frac{\operatorname{Ai}(x)\operatorname{Ai}'(y) - \operatorname{Ai}'(x)\operatorname{Ai}(y)}{x-y}\)。Tracy 和 Widom 发现 Airy 核的 Fredholm 行列式可通过 Painlevé II 方程表示。设 \(q(s)\) 为 Painlevé II 方程 \(\ddot{q} = sq + 2q^3\) 的 Hastings-McLeod 解（\(q(s) \sim \operatorname{Ai}(s)\) 当 \(s \to \infty\)）。则 \(F_2(s) = \exp(-\int_s^\infty (x-s)q^2(x)dx)\)。该表示将 Tracy-Widom 分布与完全可积系统联系起来，\(F_2\) 的尾部行为为 \(F_2(s) \sim 1 - \frac{e^{-4s^{3/2}/3}}{32\pi s^{3/2}}\)（\(s \to \infty\)）和 \(F_2(s) \sim \exp(-|s|^3/12)\)（\(s \to -\infty\)）。\(\blacksquare\)

\textbf{定理 269.3}（Tracy-Widom 的普遍性）：最大特征值的 Tracy-Widom 极限对矩阵元素分布不敏感------只要元素的四阶矩是有限的（Soshnikov 1999，Erdos-Yau 等人的工作）。这是随机矩阵边缘普遍性（edge universality）的最重要结果。

\textbf{证明}：边缘普遍性的证明策略分三步。第一步：通过矩法或 Green 函数比较定理，证明任意 Wigner 矩阵的最大特征值的分布与 GUE 的最大特征值分布之差在 \(N \to \infty\) 下趋于零。第二步：使用 Lindeberg 替换策略------逐个替换矩阵元素，将一般分布替换为 Gauss 分布。每次替换的误差由四阶 Taylor 展开控制，涉及 \(O(N^{-3})\) 的项，替换 \(O(N^2)\) 次后总误差为 \(O(N^{-1})\)。第三步：对 GUE 矩阵，特征值联合密度已知，可进行精确的渐近分析。Soshnikov 使用矩方法证明了边缘普遍性，随后 Erdos、Schlein 和 Yau 使用 Dyson Brown 运动方法给出了更一般的证明。\(\blacksquare\)

\subsection{269.3 Tracy-Widom 的惊人出现}\label{tracy-widom-ux7684ux60caux4ebaux51faux73b0}

Tracy-Widom 分布在随机矩阵理论之外也出现： - 最长递增子序列的长度分布（Baik-Deift-Johansson 1999）：Ulam 问题的解 - KPZ 普适类（Kardar-Parisi-Zhang）：\((1+1)\) 维随机生长界面的涨落 - ASEP（非对称简单排它过程）中的涨落（Tracy-Widom 2007-2009） - 二维渗流和随机三角剖分的极限界面形状

\hrulefill

\hrulefill

\hrulefill

\hrulefill

\hrulefill

\hrulefill

\section{第346章：随机矩阵与自由概率的应用}\label{ux7b2c346ux7ae0ux968fux673aux77e9ux9635ux4e0eux81eaux7531ux6982ux7387ux7684ux5e94ux7528}

随机矩阵理论的应用版图贯穿了从最纯粹的数学到最应用的工程学科的完整谱系。在纯数学一端，Montgomery 的配对相关猜想（1973 年）揭示了 Riemann \(\zeta\) 函数非平凡零点与 GUE 特征值之间的惊人一致性，这是 Hilbert-Pólya 猜想------\(\zeta\) 函数的零点对应于某个 Hermitian 算子的特征值------的最强证据。在物理学一端，Bohigas-Giannoni-Schmit 猜想（1984 年）将量子混沌系统的能级统计与随机矩阵（GOE/GUE）的 Wigner-Dyson 统计联系起来，为经典混沌与量子系统之间架起了桥梁。在统计与工程一端，Marchenko-Pastur 律和 Tracy-Widom 分布为高维协方差矩阵估计和 MIMO 通信提供了精确的渐近理论。本章在 Wigner 半圆律（第 280 章）、自由概率论（第 281 章）和 Tracy-Widom 分布（第 282 章）的基础上，展示随机矩阵理论在数论、无线通信、量子混沌和算子代数四个方向的交叉应用。

\subsection{\texorpdfstring{270.1 数论与 Riemann \(\zeta\) 函数}{270.1 数论与 Riemann \textbackslash zeta 函数}}\label{ux6570ux8bbaux4e0e-riemann-zeta-ux51fdux6570}

\textbf{定理 270.1}（Montgomery 的 pair correlation 猜想，1973 / Odlyzko 数值验证 1987）：Riemann \(\zeta\) 函数非平凡零点的配对相关函数（pair correlation function）与 GUE 随机矩阵特征值的配对相关函数一致： \[R_2(\beta, x) = 1 - \left( \frac{\sin(\pi x)}{\pi x} \right)^2\]

这是 Hilbert-Pólya 猜想的证据------Riemann \(\zeta\) 函数的零点可能是某个 Hermitian 算子的特征值。

\textbf{证明}：Montgomery 计算了 \(\zeta\) 函数零点的配对相关函数。设 \(\gamma_n\) 为 \(\zeta(1/2+it)\) 的归一化非平凡零点（\(\tilde{\gamma}_n = \frac{\gamma_n}{2\pi}\log\frac{\gamma_n}{2\pi}\)）。在限制区间 \([0,T]\) 上，配对相关函数 \(R_2(x) = \lim_{T\to\infty} \frac{1}{T}\sum_{0<\tilde{\gamma}_i,\tilde{\gamma}_j<T} \delta(\tilde{\gamma}_i - \tilde{\gamma}_j - x)\)。Montgomery 假设 Hardy-Littlewood 的 \(2\)-tuple 猜想，证明了 \(R_2(x) = 1 - (\frac{\sin(\pi x)}{\pi x})^2\)。Dyson 立即指出这正是 GUE 的配对相关函数。GUE 的 \(k\)-点相关函数由 Sine 核行列式给出：\(\rho_k(x_1,\ldots,x_k) = \det[\frac{\sin(\pi(x_i-x_j))}{\pi(x_i-x_j)}]\)。当 \(k=2\) 时，\(\rho_2(x,y) = 1 - (\frac{\sin(\pi(x-y))}{\pi(x-y)})^2\)，即配对相关函数。这一一致性是 Hilbert-Pólya 猜想的强有力证据。\(\blacksquare\)

\textbf{猜想 270.2}（Keating-Snaith 猜想，2000）：\(L\)-函数的矩与随机矩阵特征多项式的矩在 \(T \to \infty\) 下一致，预测了 Riemann \(\zeta\) 函数在临界线上 \(\log|\zeta(1/2 + it)|\) 的高斯分布性质。

\subsection{270.2 无线通信与高维统计}\label{ux65e0ux7ebfux901aux4fe1ux4e0eux9ad8ux7ef4ux7edfux8ba1}

\textbf{定义 270.1}（MIMO 系统的信道容量，Telatar 1999 / Foschini 1996）：在 \(n_t\) 发射天线、\(n_r\) 接收天线的 MIMO 信道 \(Y = H X + N\) 中（\(H\) 为信道矩阵），当 \(H\) 的元素为 i.i.d. Gauss，信道的遍历容量为 \[C = \mathbb{E}_H [\log \det(I + \rho H H^\dagger)]\]

其中 \(\rho\) 是信噪比。Marchenko-Pastur 律给出了当 \(n_t, n_r \to \infty\) 时信道容量近似。

\textbf{定义 270.2}（高维协方差估计 / 随机矩阵在统计中的应用）：对样本数 \(n\) ≈ 维数 \(p\) 的高维数据，Marchenko-Pastur 律描述了样本协方差矩阵特征值的分布偏差。这对主成分分析（PCA）的可靠性有直接影响------高维下的特征值估计不能用传统低维理论。

当 \(p/n \to \gamma \in (0, \infty)\) 时，样本协方差矩阵 \(\hat{\Sigma} = \frac{1}{n} X X^T\) 的特征值分布收敛到 Marchenko-Pastur 律：\(\rho_\gamma(\lambda) = \frac{1}{2\pi\gamma\lambda}\sqrt{(\lambda_+ - \lambda)(\lambda - \lambda_-)}\)，支撑在 \([\lambda_-, \lambda_+] = [(1-\sqrt{\gamma})^2, (1+\sqrt{\gamma})^2]\)。当 \(\gamma > 0\) 时，经验特征值被''扩展''到 \([(1-\sqrt{\gamma})^2, (1+\sqrt{\gamma})^2]\) 而非集中在 \(1\) 附近，这意味着即使真实协方差矩阵为单位阵，样本特征值也会展现出显著的散布------这是高维统计中''维度诅咒''的数学根源。在实际应用中，这导致传统 PCA 的阈值选择（如保留特征值大于 1 的主成分）在高维设置下失效。Johnstone (2001) 证明了最大样本特征值在适当缩放后收敛到 Tracy-Widom 分布，为高维假设检验（如球形检验 \(\Sigma = I\)）提供了精确的渐近分布。此外，随机矩阵理论在压缩感知、MIMO 通信、金融风险管理的协方差矩阵估计中均有重要应用------\(p \gg n\) 情形下的正则化估计（如 Ledoit-Wolf 收缩估计）正是基于随机矩阵的渐近理论。

\subsection{270.3 量子混沌与核物理}\label{ux91cfux5b50ux6df7ux6c8cux4e0eux6838ux7269ux7406}

\textbf{定义 270.3}（Bohigas-Giannoni-Schmit 猜想，1984）：量子混沌系统的能级统计统计（nearest-neighbor spacing distribution）遵循适当对称类（GOE/GUE/GSE）的随机矩阵特征值间距分布（Wigner 猜忌 / Wigner surmise）： \[P_{\text{GUE}}(s) = \frac{32}{\pi^2} s^2 e^{-4s^2 / \pi}\]

（\(s\) 为间距除以平均间距）。

Bohigas-Giannoni-Schmit (BGS) 猜想将量子混沌与随机矩阵理论建立了深刻联系：经典混沌系统的量子对应物，其能级间距统计遵循随机矩阵预测（GOE 对应时间反演不变系统，GUE 对应时间反演破缺系统）。这一猜想的物理直觉是：混沌系统的量子对应物失去了所有经典守恒量，其 Hamiltonian 在适当的基下如同一个''随机矩阵''。BGS 猜想已被大量数值实验（如 Sinai 台球、量子 Baker 映射）和半经典理论（Gutzwiller 迹公式）验证。能级排斥（level repulsion）------邻近能级避免交叉------是量子混沌的标志性特征，其 \(s^\beta\) 的小间距行为（\(\beta = 1, 2, 4\) 分别对应 GOE、GUE、GSE）与可积系统的 Poisson 间距分布 \(e^{-s}\) 形成鲜明对比。这一对应关系在核物理中具有悠久历史：Wigner 在 1950 年代首先提出用随机矩阵描述重核的共振能级，随后被 Porter 和 Rosenzweig 系统化。现代介观物理（量子点、Anderson 局域化）中的电导涨落普遍性也由随机矩阵理论描述。

\subsection{270.4 自由概率在算子代数中的应用}\label{ux81eaux7531ux6982ux7387ux5728ux7b97ux5b50ux4ee3ux6570ux4e2dux7684ux5e94ux7528}

\textbf{定理 270.3}（Voiculescu 的自由群因子同构问题，1990s）：使用自由概率论，Voiculescu 证明了自由群 von Neumann 代数 \(L(\mathbb{F}_n)\) 与 \(L(\mathbb{F}_m)\)（\(n \neq m\)）的非同构性在自由熵框架下可被区分。这解决了 von Neumann 代数领域的一个重大未解问题。

自由概率论在算子代数领域的核心贡献在于提供了一类全新的不变量------自由熵和自由 Fisher 信息。在 II\(_1\) 因子理论中，因子同构分类是中心问题之一。Murray 和 von Neumann 在 1930-40 年代建立了因子的基本分类（I、II\(_1\)、II\(_\infty\)、III 型），但 II\(_1\) 因子的细致分类至今仍是活跃的研究领域。自由群因子 \(L(\mathbb{F}_n)\) 是 II\(_1\) 因子的重要例子，其生成元构成自由族。Voiculescu 利用自由概率方法证明了 \(L(\mathbb{F}_n) \not\cong L(\mathbb{F}_m)\) 当 \(n \neq m\) 时，这一结果使自由群因子成为第一个被严格证明不可同构的 II\(_1\) 因子族。自由熵方法随后被推广到更一般的 II\(_1\) 因子，包括超积（ultraproduct）和自由积（free product）因子的分类。此外，自由概率论通过随机矩阵的渐近自由性质，建立了 II\(_1\) 因子与随机矩阵代数的联系------自由群因子 \(L(\mathbb{F}_n)\) 可视为独立 GUE 矩阵的极限代数结构，这为算子代数提供了新的组合和概率视角。

\textbf{证明}：Voiculescu 定义了自由熵 \(\chi(X_1,\ldots,X_n)\) 作为自由群 von Neumann 代数 \(L(\mathbb{F}_n)\) 生成元系的不变量。自由熵 \(\chi\) 满足超可加性：\(\chi(X_1,\ldots,X_{n+m}) \geq \chi(X_1,\ldots,X_n) + \chi(X_{n+1},\ldots,X_{n+m})\)。Voiculescu 证明了 \(\chi(L(\mathbb{F}_n)) = n\)（生成元数目为自由熵）。若 \(L(\mathbb{F}_n) \cong L(\mathbb{F}_m)\)，则存在生成元系在 \(L(\mathbb{F}_n)\) 中产生 \(m\) 个自由半圆元素，给出 \(\chi(L(\mathbb{F}_n)) \geq m\)。同理，\(\chi(L(\mathbb{F}_m)) \geq n\)。因此 \(n \geq m\) 且 \(m \geq n\)，故 \(n = m\)。自由熵作为不变量区分了不同秩的自由群因子。\(\blacksquare\)

\subsection{随机矩阵理论的前沿扩展}\label{ux968fux673aux77e9ux9635ux7406ux8bbaux7684ux524dux6cbfux6269ux5c55}

随机矩阵理论在 21 世纪初经历了从物理到数学的深刻转变，其核心发展体现在普遍性定理的严格证明和机器学习中的应用。\textbf{Erdos-Schlein-Yau 的 Dyson 三阶段方法}（2008-2015）是随机矩阵理论中最重要的突破之一：通过局部半圆律（local semicircle law）、Dyson Brown 运动（DBM）和比较论证（comparison argument），他们证明了 Wigner 矩阵的局部谱统计（包括相邻特征值间距分布和 \(k\)-点相关函数）具有普遍性------与具体矩阵元分布无关，仅由对称性类（实数、复数或四元数）决定。这一结果将 Wigner 最初关于重原子核能级统计的猜想（1955）数学化，是随机矩阵理论自 1960 年代以来的最高成就。\textbf{Tracy-Widom 分布}（1994）是随机矩阵理论中最重要的分布之一：GUE/GOE 最大特征值的涨落（经适当缩放）收敛到 Tracy-Widom 分布 \(F_\beta\)（\(\beta = 1, 2, 4\)），其累积分布函数由 Painlevé II 方程的解表达。Tracy-Widom 分布在物理学（KPZ 普适性类）、组合数学（最长递增子序列的 Baik-Deift-Johansson 定理）和统计学（主成分分析中的最大特征值检验）中具有惊人的普遍性。\textbf{自由概率论}（Voiculescu 1985）将随机矩阵的大 \(N\) 极限与 von Neumann 代数中的自由独立性联系起来：独立的高斯随机矩阵族在 \(N \to \infty\) 时渐近自由，其极限谱分布由自由卷积（free convolution）描述。自由概率论为随机矩阵理论提供了代数-组合学工具，特别是自由累积量（free cumulants）和 \(R\)-变换，使得谱分布的计算被简化为多项式的代数操作。\textbf{随机矩阵与机器学习}：随机矩阵理论在理解深度学习的理论和实践中发挥着关键作用------神经网络的 Hessian 矩阵的谱分布（特别是 Marchenko-Pastur 型和 Wigner 型）决定了优化景观的几何性质，而随机矩阵的 Tracy-Widom 分布被用于矩阵补全和主成分分析中的秩估计。 随机矩阵理论近二十年的突破集中在普遍性定理的严格证明（Erdos-Schlein-Yau 等人）和机器学习中的应用。

\subsection{271.1 普遍性定理的严格证明}\label{ux666eux904dux6027ux5b9aux7406ux7684ux4e25ux683cux8bc1ux660e}

\textbf{定义 271.1}（Dyson Brown 运动与局部平衡）：将 Wigner 矩阵 \(H\) 的特征值 \(\lambda_i(t)\) 视为时间演化的 Dyson Brown 运动（\(\lambda_i\) 间对数排斥力 + 噪声）。关键：Dyson Brown 运动在短时间达到局部平衡（local equilibrium），使得特征值间距的短距离统计普遍。

Dyson Brown 运动是普遍性证明中的核心动力学工具。通过引入虚拟时间参数，将矩阵元素替换为 Ornstein-Uhlenbeck 过程，特征值在 DBM 下朝向平衡态演化。Erdos、Schlein 和 Yau 利用这一机制将普遍性问题转化为三个相互关联的步骤：局部律、局部平衡和比较论证。

\textbf{定理 271.1}（Erdos-Schlein-Yau 的 Wigner-Dyson-Mehta 普遍性，2010-2013）：对任何满足适当矩条件的广义 Wigner 矩阵，其特征值的\textbf{局部统计}（local statistics）------包括相邻间距分布和 \(k\)-点相关函数------与 GUE/GOE 的对应统计在极限下一致。即 Tracy-Widom 边缘分布和 Wigner-Dyson 体（bulk）统计具有普遍性。

\textbf{证明}：Erdos-Schlein-Yau 的三步法。第一步：建立局部半圆律------在尺度 \(\eta \gg 1/N\) 上，Stieltjes 变换 \(s_N(z)\) 以高概率接近半圆律的 Stieltjes 变换 \(s_{sc}(z)\)，误差为 \(O(1/(N\eta))\)。第二步：使用 Dyson Brown 运动（DBM）------引入虚拟时间参数，将矩阵元素用 Ornstein-Uhlenbeck 过程替换，特征值在 DBM 下演化至局部平衡。DBM 的生成元为 \(\mathcal{L} = \sum_i \partial_i^2 + \sum_i (- \frac{1}{2}\lambda_i + \sum_{j \neq i} \frac{1}{\lambda_i - \lambda_j}) \partial_i\)。第三步：比较论证------利用局部半圆律的控制，证明一般 Wigner 矩阵的局部统计与 GUE 的局部统计在 \(N \to \infty\) 下一致。关键估计是 Green 函数比较定理，通过逐元素替换将一般分布逼近 Gauss 分布。\(\blacksquare\)

上述普遍性定理聚焦于特征值在谱内部（体区）的局部统计，而矩阵边缘------特别是最大特征值附近------的行为同样具有深刻的普遍性。最大特征值的涨落由 Tracy-Widom 分布刻画，而次最大特征值与最大特征值之间的间距（gap）则构成边缘统计的另一重要层次。

\textbf{定理 271.2}（Dyson 猜想 / 边缘隙分布，2020s）：对于矩阵边缘的最大特征值，Tracy-Widom 分布是普遍的。对于次最大特征值的隙分布也有类似的普遍性结果。

\textbf{证明}：最大特征值 \(\lambda_1\) 和次最大特征值 \(\lambda_2\) 的联合分布由 Airy 点过程描述。在缩放 \(\lambda_i = 2\sqrt{N} + N^{-1/6} a_i\) 下，\((a_1, a_2, \ldots)\) 收敛到 Airy 点过程。Airy 点过程的联合强度由 Airy 核的行列式给出：\(\rho_k(a_1,\ldots,a_k) = \det[K_{\text{Airy}}(a_i, a_j)]\)。\(\lambda_1\) 和 \(\lambda_2\) 的隙 \(\Delta = \lambda_1 - \lambda_2\) 在缩放后由 Airy 点过程的前两个点给出。隙分布 \(P(\Delta > t N^{-1/6})\) 可通过 Airy 核的 Fredholm 行列式表示。普遍性结果由 Erdos-Yau 的局部 DBM 方法证明：隙的分布对矩阵元素分布不敏感，在四阶矩有限条件下收敛到与 GUE 相同的极限。\(\blacksquare\)

\subsection{271.2 随机矩阵与机器学习}\label{ux968fux673aux77e9ux9635ux4e0eux673aux5668ux5b66ux4e60}

\textbf{定义 271.2}（随机神经网络 / Random Matrix Theory of Deep Learning）：深度神经网络的权值矩阵（随机初始化）的特征值分布影响梯度传播和训练动力学。Gaussian 权值网络的 Hessian 矩阵和 Fisher 信息矩阵的谱可用随机矩阵理论分析（Pennington-Worah 2017，Pennington-Schoenholz-Ganguli 2017-2018）。

随机矩阵理论在深度学习中的一个最引人注目的预测是所谓的''双重下降''现象。与传统统计学习理论中测试误差随模型复杂度呈U形曲线不同，过参数化模型在插值阈值之后测试误差再次下降。这一反直觉的行为可通过随机矩阵理论中的最小范数插值解的谱分析得到精确解释。

\textbf{定理 271.3}（双下降现象 / Double Descent，Belkin-Hsu-Ma-Mandal 2019 / 随机矩阵解释）：当模型参数数超过训练样本数时，测试误差出现双下降（插值阈值后再次下降）。随机矩阵理论解释了这一现象------过参数化后的最小范数解（伪逆）比恰好在插值阈值处的解具有更小的方差。

\textbf{证明}：考虑线性模型 \(y = X\beta^* + \varepsilon\)，\(X \in \mathbb{R}^{n \times p}\)。最小范数最小二乘解为 \(\hat{\beta} = X^\top (XX^\top)^{-1}y\)（当 \(p > n\) 时）。预测风险为 \(R(\hat{\beta}) = \mathbb{E}\|\hat{y} - X\beta^*\|^2 = \sigma^2 \operatorname{Tr}((XX^\top)^{-1}XX^\top) + \|\beta^*\|^2 \operatorname{Tr}((I - X^\top(XX^\top)^{-1}X)XX^\top)\)。设 \(X\) 元素为 i.i.d. \(\mathcal{N}(0,1)\)。由随机矩阵理论，\(XX^\top/n\) 的特征值服从 Marchenko-Pastur 律。当 \(p \to n\) 时，\(XX^\top\) 的最小特征值趋于 \(0\)（以速率 \(n^{-2/3}\)），导致 \((XX^\top)^{-1}\) 的迹发散。当 \(p \gg n\) 时，\(XX^\top/p \to I_n\)（由大数定律），因此 \((XX^\top)^{-1} \approx p^{-1} I_n\)，风险趋于 \(n\sigma^2/p\)，随 \(p\) 增大而减小。这解释了测试误差在 \(p=n\) 处发散、在 \(p \gg n\) 时再次下降的双重下降现象。\(\blacksquare\)

\subsection{271.3 稀疏随机矩阵与稀疏 Wigner}\label{ux7a00ux758fux968fux673aux77e9ux9635ux4e0eux7a00ux758f-wigner}

\textbf{定义 271.3}（稀疏 Wigner 矩阵 / Erdős-Rényi 矩阵）：Wigner 矩阵的每个元素以概率 \(p = d/N\) 非零（\(d\) 为平均度）。当 \(d \to \infty\) 但 \(d/N \to 0\) 时，谱分布偏离半圆律，边缘出现离群特征值。

稠密 Wigner 矩阵的每个元素都以 \(O(1/\sqrt{N})\) 的量级非零，但在许多实际应用中（如稀疏网络、局部相互作用的物理系统），矩阵具有大量零元素。稀疏 Wigner 矩阵通过引入平均度参数 \(d\) 刻画了这种稀疏性，其谱分布在 \(d\) 从大到小的过渡中展现出从半圆律到尖锐峰状的连续相变。

\textbf{定理 271.4}（稀疏 Wigner 谱的过渡）：稀疏 Wigner 矩阵的谱在 \(d \gg 1\) 时接近半圆，在 \(d\) 有限时呈现更尖锐的峰。存在从稀疏谱到稠密半圆谱的连续过渡。

\textbf{证明}：设稀疏 Wigner 矩阵 \(H\) 的元素以概率 \(p = d/N\) 非零，非零元素方差为 \(1/d\)（保持 \(\mathbb{E}[\operatorname{Tr}(H^2)]/N = 1\)）。在矩法计算中，配对 \(\pi\) 的贡献为 \((1/d)^{k/2} \cdot N^{V(\pi)} \cdot d^{k/2} \cdot N^{-k/2} = N^{V(\pi)-k/2}\)。当 \(d \to \infty\) 时，所有配对（包括非交叉和交叉）的非零概率均等，非交叉配对仍主导，谱分布趋于半圆律。当 \(d\) 有限时，交叉配对被进一步压制，但某些高阶矩的贡献改变。特征值密度的解析形式为 \(\rho(\lambda) = \frac{\sqrt{4 - \lambda^2}}{2\pi} + \frac{1}{d} \rho_1(\lambda) + O(1/d^2)\)。当 \(d\) 很小时，密度在 \(\lambda = 0\) 附近出现尖峰（源于局部树结构的贡献），与半圆律的平滑形态差异显著。过渡发生在 \(d \sim \log N\) 尺度。\(\blacksquare\)

\hrulefill

\hrulefill

\hrulefill

\hrulefill

\hrulefill

\hrulefill

\section{第347章：圆律与Stieltjes变换方法}\label{ux7b2c347ux7ae0ux5706ux5f8bux4e0estieltjesux53d8ux6362ux65b9ux6cd5}

Wigner 半圆律（第 280 章）刻画了对称（自伴）随机矩阵的谱极限------所有特征值集中在实轴上的区间 \([-2, 2]\) 内，密度呈半圆形。但对于非 Hermitian 随机矩阵（矩阵元素独立分布，矩阵不对称），谱的极限分布具有完全不同的几何形态：所有特征值在复平面上均匀分布在单位圆盘内------这就是 Girko（1984 年）提出、经 Bai（1997 年）和 Tao-Vu（2010 年）最终完全证明的圆律（Circular Law）。圆律的证明涉及随机矩阵理论中最精妙的分析工具------Hermitian 化方法（将非 Hermitian 矩阵嵌入为更大的 Hermitian 矩阵）和 Stieltjes 变换------将谱分布的计算转化为复分析中的对数势边值问题。圆律与半圆律的关系体现了 Hermitian 与一般矩阵之间的对称破缺，是随机矩阵从''实谱''到''复谱''的自然延伸。本章从 i.i.d. 矩阵系综出发，以 Hermitian 化方法和 Stieltjes 变换严格证明圆律，并讨论稀疏矩阵和重尾分布下的谱相变。

\subsection{272.1 非Hermitian矩阵与圆律}\label{ux975ehermitianux77e9ux9635ux4e0eux5706ux5f8b}

\textbf{定义 272.1}（i.i.d. 矩阵系综）：设 \(X_N\) 是 \(N \times N\) 复矩阵，所有元素 \(x_{ij}\) 是 i.i.d. 复随机变量，满足 \(\mathbb{E}[x_{ij}] = 0\)，\(\mathbb{E}[|x_{ij}|^2] = 1/N\)。\(X_N\) 的特征值 \(\lambda_i \in \mathbb{C}\)（复数，一般不对称）的经验谱分布是 \(\mu_N = \frac{1}{N} \sum_{i=1}^N \delta_{\lambda_i}\)。

与 Wigner 矩阵不同，i.i.d. 矩阵系综中的矩阵不满足 Hermitian 对称性，因此其特征值不再是实数，而是分布在复平面上。这一根本差异导致谱极限从实轴上的半圆律变为复平面单位圆盘上的均匀分布。Girko 在 1984 年首先提出了这一猜想，随后经过 Bai（1997）和 Tao-Vu（2010）的逐步严格化，最终被完全证明。

\textbf{定理 272.1}（圆律，Girko 1984 / Bai 1997 / Tao-Vu 2010 完全证明）：对上述非 Hermitian i.i.d. 矩阵（方差 \(1/N\)），当 \(N \to \infty\) 时，经验谱分布 \(\mu_N\) 几乎必然弱收敛到复平面单位圆盘上的均匀分布： \[d\mu_{\text{circ}}(z) = \frac{1}{\pi} \mathbf{1}_{|z| \leq 1} \, dx \, dy \quad (z = x + iy)\]

\textbf{证明}（Hermitization 方法概要）：对非 Hermitian 矩阵 \(X_N\)，定义 \(2N \times 2N\) 的 Hermitian 矩阵 \(H^z_N = \begin{pmatrix} 0 & X_N - zI \\ X_N^* - \bar{z}I & 0 \end{pmatrix}\)。\(X_N\) 的特征值 \(\lambda_i\) 满足 \(\frac{1}{N}\sum_i \log|\lambda_i - z| = \frac{1}{2N}\operatorname{Tr}\log((X_N - zI)(X_N - zI)^*)\)。引入对数势 \(U_N(z) = \int \log|z-w| d\mu_N(w)\) 和 \(H^z_N\) 的 Stieltjes 变换。在 \(N \to \infty\) 极限下，\(U_N\) 收敛到 \(U\)，且满足 \(\Delta U = 2\pi\rho\)（\(\rho\) 为极限密度）。由 \(H^z_N\) 的 Stieltjes 变换自洽方程，可推出 \(U(z)\) 在单位圆盘内满足 \(\Delta U(z) = 2\pi \cdot (1/\pi) = 2\)，圆盘外满足 \(\Delta U(z) = 0\)。由 Laplace 方程解的唯一性，对数势唯一确定：\(U(z) = |z|^2\)（\(|z| \leq 1\)），\(U(z) = \log|z|^2 + 1\)（\(|z| \geq 1\)）。由 \(\rho = \frac{1}{2\pi}\Delta U\) 得 \(\rho(z) = \frac{1}{\pi} \mathbf{1}_{|z| \leq 1}\)。\(\blacksquare\)

\subsection{272.2 Stieltjes变换与Hermitization方法}\label{stieltjesux53d8ux6362ux4e0ehermitizationux65b9ux6cd5}

\textbf{定义 272.2}（Stieltjes 变换）：对 \(\mathbb{R}\) 上的概率测度 \(\mu\)，其 \textbf{Stieltjes 变换}（或称 Cauchy 变换）定义为 \[s_\mu(z) = \int_{\mathbb{R}} \frac{d\mu(x)}{x - z}, \quad z \in \mathbb{C}^+ = \{z : \Im z > 0\}\]

Stieltjes 变换是研究谱分布的基本解析工具，其关键性质如下。

\textbf{命题 272.1}（Stieltjes 变换的基本性质）： 1. \(s_\mu(z)\) 在 \(\mathbb{C}^+\) 上解析，且 \(\Im s_\mu(z) \geq 0\)（Nevanlinna 函数） 2. \textbf{反演公式}（Stieltjes-Perron）：\(\mu\) 可经 \(s_\mu\) 恢复：\(\mu((a,b)) = \lim_{\epsilon \to 0^+} \frac{1}{\pi} \int_a^b \Im s_\mu(x + i\epsilon) \, dx\)（\(\mu\) 的密度是 \(\frac{1}{\pi} \Im s_\mu(x + i0)\)） 3. 设 \(\mu_N\) 弱收敛到 \(\mu\)，当且仅当对所有 \(z \in \mathbb{C}^+\)，\(s_{\mu_N}(z) \to s_\mu(z)\)

\textbf{定理 272.2}（半圆律的 Stieltjes 变换推导）：对 Wigner 矩阵 \(H_N\)，谱分布 \(\mu_N\) 的 Stieltjes 变换 \(s_N(z) = \frac{1}{N} \operatorname{Tr}((H_N - zI)^{-1})\) 在 \(N \to \infty\) 下满足自洽方程（self-consistent equation）： \[s(z) = -\frac{1}{z + s(z)}\] 解得 \(s(z) = \frac{-z + \sqrt{z^2 - 4}}{2}\)（平方根分支取使 \(\Im s(z) > 0\) 者），其反演给出半圆律密度。这比矩法更简洁------无需计算 Catalan 数。

\textbf{证明}：对 Wigner 矩阵 \(H_N\)，Stieltjes 变换 \(s_N(z) = \frac{1}{N}\operatorname{Tr}((H_N - zI)^{-1})\)。由 Schur 补公式，将 \(H_N\) 按第一行和第一列分块，可得 \((H_N - zI)^{-1}_{11} = (h_{11} - z - h_{1*}^\top (H^{(1)} - zI)^{-1} h_{1*})^{-1}\)，其中 \(H^{(1)}\) 为删除第一行第一列的子矩阵，\(h_{1*}\) 为第一列剩余元素。在 \(N \to \infty\) 下，\(h_{1*}^\top (H^{(1)} - zI)^{-1} h_{1*} \to \frac{1}{N}\operatorname{Tr}((H - zI)^{-1}) = s(z)\)（由迹的集中性）。因此 \(s(z) = \mathbb{E}[(h_{11} - z - s(z))^{-1}]\)。对 Wigner 矩阵，\(h_{11} \to 0\)，故 \(s(z) = -(z + s(z))^{-1}\)，即 \(s(z)^2 + z s(z) + 1 = 0\)。解得 \(s(z) = (-z + \sqrt{z^2 - 4})/2\)。由 Stieltjes 反演，\(\rho(x) = \frac{1}{\pi}\lim_{\epsilon \to 0^+} \Im s(x+i\epsilon) = \frac{1}{2\pi}\sqrt{4-x^2} \mathbf{1}_{[-2,2]}(x)\)。\(\blacksquare\)

\subsection{272.3 自由概率论与Stieltjes变换}\label{ux81eaux7531ux6982ux7387ux8bbaux4e0estieltjesux53d8ux6362}

\textbf{定义 272.3}（Stieltjes 变换与 R-变换的关系）：\(R_\mu(z) = K_\mu(z) - 1/z\)，其中 \(K_\mu\) 是 \(G_\mu(z) = -s_\mu(z)\) 的泛函逆。R-变换在自由卷积下的线性性使 Stieltjes 变换成为计算独立随机矩阵之和（在自由独立极限下）的谱分布的基本工具。

\textbf{定理 272.3}（圆律的 Hermitization 证明概要）：对非 Hermitian \(X_N\)，定义 \(2N \times 2N\) Hermitian 矩阵 \[H^z_N = \begin{pmatrix} 0 & X_N - zI_N \\ (X_N - zI_N)^* & 0 \end{pmatrix}\] 其 Stieltjes 变换满足与 \(z\) 相关的矩阵方程。通过 Stieltjes 变换的偏微分方程，可以将 \(\mu_N\) 的复对数势（logarithmic potential）与 \(H^z_N\) 的 Stieltjes 变换联系起来。取极限 \(N \to \infty\)，对数势在单位圆盘内满足 Poisson 方程 \(\Delta \Phi = -4\)（圆盘外 \(\Delta \Phi = 0\)），由唯一性给出圆盘上的均匀分布（圆律）。

\textbf{证明}：对非 Hermitian 矩阵 \(X_N\)，经验谱分布 \(\mu_N\) 的特征值 \(\lambda_i\) 满足 \(\frac{1}{N}\sum_i \log|\lambda_i - z| = \frac{1}{2N}\operatorname{Tr}\log((X_N - zI)(X_N - zI)^*) = \frac{1}{2N}\sum_{j=1}^{2N} \log|\mu_j(z)|\)，其中 \(\mu_j(z)\) 为 \(H^z_N\) 的特征值。引入对数势 \(U_N(z) = \int \log|z-w| d\mu_N(w)\)。\(H^z_N\) 的 Stieltjes 变换 \(g_N(z,\eta) = \frac{1}{2N}\operatorname{Tr}((H^z_N - \eta I)^{-1})\) 满足矩阵方程。在 \(N \to \infty\) 下，\(g_N\) 收敛到极限 \(g\)，且对数势 \(U(z) = \lim U_N(z)\) 满足 \(\Delta U(z) = 2\pi \rho(z)\)。由 \(H^z_N\) 的 Stieltjes 变换方程，可推出 \(\Delta U(z) = 2\)（\(|z|<1\)）和 \(\Delta U(z) = 0\)（\(|z|>1\)）。结合 \(\int \rho(z) dz = 1\)，得 \(\rho(z) = \frac{1}{\pi} \mathbf{1}_{|z| \leq 1}\)。\(\blacksquare\)

\hrulefill

\hrulefill

\hrulefill

\hrulefill

\hrulefill

\hrulefill

\section{第348章：行列式点过程与β系综}\label{ux7b2c348ux7ae0ux884cux5217ux5f0fux70b9ux8fc7ux7a0bux4e0eux3b2ux7cfbux7efc}

随机矩阵的特征值构成具有排斥相互作用的经典粒子系统。\textbf{行列式点过程}（Determinantal Point Process, DPP）框架统一了经典 Gaussian 系综的相关函数结构，而 \textbf{\(\beta\)系综}（\(\beta\)-ensemble）推广到任意排斥强度------将随机矩阵从经典值 \(\beta = 1, 2, 4\) 拓展为一般的逆温度参数。

\subsection{273.1 行列式点过程}\label{ux884cux5217ux5f0fux70b9ux8fc7ux7a0b}

\textbf{定义 273.1}（行列式点过程）：\(\mathbb{R}\) 上具有相关核 \(K(x, y)\) 的\textbf{行列式点过程}是一个点过程（随机点集），其 \(k\)-点相关函数（joint intensity）为 \[\rho_k(x_1, \ldots, x_k) = \det[K(x_i, x_j)]_{1 \leq i,j \leq k}\]

行列式点过程的核心性质在于其相关函数完全由单个相关核 \(K(x,y)\) 的行列式结构决定，这赋予了点过程极大的可算性。GUE 特征值之所以能够被精确分析，正是因为它恰好构成一个行列式点过程------这一发现将随机矩阵的谱统计转化为积分算子的 Fredholm 行列式计算，进而与完全可积系统建立了深刻联系。

\textbf{定理 273.1}（GUE 是行列式点过程）：GUE 矩阵的特征值构成一个行列式点过程，其相关核为 \[K_N(x, y) = \sum_{k=0}^{N-1} \varphi_k(x) \varphi_k(y)\] 其中 \(\varphi_k(x) = \frac{1}{\sqrt{2^k k! \sqrt{\pi}}} H_k(x) e^{-x^2/2}\) 是归一化 Hermite 函数（\(H_k\) 为 Hermite 多项式）。由 Christoffel-Darboux 公式： \[K_N(x, y) = \sqrt{N} \frac{\varphi_N(x) \varphi_{N-1}(y) - \varphi_{N-1}(x) \varphi_N(y)}{x - y}\]

\textbf{证明}：GUE 特征值联合密度 \(p(\lambda) \propto \prod_{i<j} |\lambda_i - \lambda_j|^2 \prod_i e^{-\lambda_i^2/2}\) 中，Vandermonde 行列式 \(\Delta(\lambda) = \prod_{i<j} (\lambda_i - \lambda_j) = \det[\lambda_i^{j-1}]\)。将每行 \(\lambda_i^{j-1}\) 用 Hermite 多项式 \(H_{j-1}(\lambda_i)\) 替换（行列式仅相差常数因子），得 \(p(\lambda) \propto \det[\varphi_{j-1}(\lambda_i)] \cdot \det[\varphi_{j-1}(\lambda_i)]\)。由 Cauchy-Binet 公式，\(k\)-点相关函数 \(\rho_k(x_1,\ldots,x_k) = \frac{N!}{(N-k)!} \int p(x_1,\ldots,x_N) dx_{k+1}\cdots dx_N = \det[K_N(x_i, x_j)]_{i,j=1}^k\)，其中 \(K_N(x,y) = \sum_{j=0}^{N-1} \varphi_j(x)\varphi_j(y)\)。Christoffel-Darboux 恒等式给出 \(K_N(x,y) = \sqrt{N}\frac{\varphi_N(x)\varphi_{N-1}(y) - \varphi_{N-1}(x)\varphi_N(y)}{x-y}\)。\(\blacksquare\)

行列式点过程的一个直接推论是''空隙概率''------特征值在某个区间内完全缺失的概率------可以通过相关核的 Fredholm 行列式精确表达。这一关系看似形式化，实则是 Tracy-Widom 分布能够通过 Painleve II 方程表达的数学根源：Fredholm 行列式与完全可积系统之间的深刻联系由 Its、Izergin、Korepin 和 Slavnov 揭示。

\textbf{推论 273.1}（间距分布与 Fredholm 行列式）：GUE 特征值在一个区间内没有点的概率（gap probability）为 \[\mathbb{P}(I \text{ 内无特征值}) = \det(I - K_N|_I)\] 其中 \(\det(I - K_N|_I)\) 是 \(L^2(I)\) 上积分算子 \(K_N|_I\) 的 Fredholm 行列式。这是 Tracy-Widom 分布可通过 Painlevé II 表达的核心原因------Fredholm 行列式与完全可积系统有深刻的联系（Its-Izergin-Korepin-Slavnov 方法）。

在谱的内部区域（远离边缘），特征值的局部统计展现出与边缘截然不同的极限行为。通过在体区适当缩放，Hermite 核收敛到平移不变的 Sine 核，其特征值的局部间距分布由 Gaudin 分布描述。这一体区普遍性由 Gaudin 和 Mehta 首先提出，并由 Erdos 和 Yau 在 2010 年代严格证明。

\textbf{定理 273.2}（Sine 核普遍性，Gaudin-Mehta / Erdos-Yau 严格化）：GUE 特征值在体（bulk）极限下的局部相关核收敛到 \textbf{Sine 核}： \[K_{\text{sine}}(x, y) = \frac{\sin(\pi(x - y))}{\pi(x - y)}\] GUE 的体间距分布由 Sine 核的 Fredholm 行列式给出（Gaudin 分布）。

\textbf{证明}：在体（bulk）缩放下，取 \(x = u + \xi/(\rho_{sc}(u)N)\)，\(y = u + \eta/(\rho_{sc}(u)N)\)，其中 \(\rho_{sc}(u) = \frac{1}{2\pi}\sqrt{4-u^2}\) 为半圆密度。Hermite 函数的渐近公式（Plancherel-Rotach）给出 \(\varphi_N(x) \sim \frac{1}{\sqrt{\rho_{sc}(u)}} \cos(N\psi(u) + \pi\xi)\)。代入 Christoffel-Darboux 公式，当 \(N \to \infty\) 时，\(K_N(x,y) \to \frac{\sin(\pi(\xi-\eta))}{\pi(\xi-\eta)} = K_{\text{sine}}(\xi,\eta)\)。Sine 核是平移不变的，仅依赖于 \(\xi-\eta\)。Gap 概率 \(\mathbb{P}(I \text{ 内无特征值}) = \det(I - K_{\text{sine}}|_I)\) 给出 Gaudin 分布。普遍性结果（Erdos-Yau）证明任意 Wigner 矩阵的体局部统计也收敛到 Sine 核。\(\blacksquare\)

与体区的 Sine 核相对应，在谱的边缘------即最大特征值附近------相关核收敛到另一个经典的极限对象：Airy 核。Airy 核直接关联到第 282 章讨论的 Tracy-Widom 分布，其 Fredholm 行列式恰好给出最大特征值的极限分布函数。

\textbf{定义 273.2}（Airy 核与 Tracy-Widom 的 Fredholm 行列式表示）：在谱的边缘（最大特征值附近），缩放后的相关核收敛到 \textbf{Airy 核}： \[K_{\text{Airy}}(x, y) = \frac{\operatorname{Ai}(x) \operatorname{Ai}'(y) - \operatorname{Ai}'(x) \operatorname{Ai}(y)}{x - y}\] Tracy-Widom 分布 \(F_2(s)\) 由 Airy 核在 \((s, \infty)\) 上的 Fredholm 行列式给出： \[F_2(s) = \det(I - K_{\text{Airy}}|_{(s, \infty)})\]

\subsection{273.2 β系综与对数气体}\label{ux3b2ux7cfbux7efcux4e0eux5bf9ux6570ux6c14ux4f53}

\textbf{定义 273.3}（\(\beta\)系综 / 对数气体）：\(\mathbb{R}\) 上 \(N\) 个粒子的 \textbf{\(\beta\)系综}（\(\beta\)-ensemble）概率密度为 \[p(\lambda_1, \ldots, \lambda_N) = \frac{1}{Z_{N, \beta}} \prod_{i < j} |\lambda_i - \lambda_j|^\beta \prod_{i=1}^N e^{-N V(\lambda_i)}\] 其中 \(V\) 是势函数（confining potential），\(\beta > 0\) 是逆温度参数： - \(\beta = 1\)：对应于实对称（GOE）特征值 - \(\beta = 2\)：对应于复 Hermitian（GUE）特征值 - \(\beta = 4\)：对应于辛（GSE）特征值 - \(\beta\) 为一般值：不再对应具体矩阵系综，但作为对数相互作用的粒子系统仍具有良好的数学结构

\(\beta\)系综将随机矩阵特征值的联合密度推广为带有一般排斥参数 \(\beta\) 的对数气体模型。当 \(\beta\) 取经典值 \(1, 2, 4\) 时分别还原 GOE、GUE 和 GSE，但 \(\beta\) 的连续推广揭示了一个更深层的事实：宏观极限谱分布（即经验测度的极限）对任意 \(\beta > 0\) 都是相同的------仅由势函数 \(V\) 决定，与排斥强度 \(\beta\) 无关。

\textbf{定理 273.3}（\(\beta\)系综的宏观极限分布，Ben Arous-Guionnet 1997 / 大偏差理论）：\(\beta\)系综的经验测度 \(\mu_N = \frac{1}{N} \sum \delta_{\lambda_i}\) 在 \(N \to \infty\) 下几乎必然收敛到势 \(V\) 决定的平衡测度 \(\mu_V\)（极小化能量泛函 \(\mathcal{E}(\mu) = \iint \log|x-y|^{-1} d\mu(x) d\mu(y) + \int V(x) d\mu(x)\)）。当 \(V(x) = x^2/2\)（Gauss 势），\(\mu_V\) 是半圆律；此结果对任意 \(\beta > 0\) 成立，不限于经典值。

\textbf{证明}：将 \(\beta\)系综的密度写为 \(p(\lambda) \propto \exp(-N^2 \mathcal{E}_N(\mu_N))\)，其中 \(\mathcal{E}_N\) 为离散能量泛函。由 Sanov 型大偏差原理，\(\mu_N\) 满足大偏差上界，速率函数为 \(I(\mu) = \beta \mathcal{E}(\mu) - \inf_\nu \beta \mathcal{E}(\nu)\)。由 Laplace 原理（Varadhan 引理），\(\mu_N\) 几乎必然收敛到 \(\mathcal{E}\) 的全局极小值 \(\mu_V\)。\(\mu_V\) 由 Euler-Lagrange 方程确定：\(\int \log|x-y|^{-1} d\mu_V(y) + \frac{1}{2}V(x) = \text{常数}\)（在 \(\operatorname{supp}(\mu_V)\) 上）。对 \(V(x) = x^2/2\)，该方程的解为半圆律 \(\mu_V(dx) = \frac{1}{2\pi}\sqrt{4-x^2} \mathbf{1}_{[-2,2]}(x) dx\)。\(\blacksquare\)

宏观极限分布仅依赖于势函数 \(V\) 而与 \(\beta\) 无关，但系统的涨落行为------特别是最大特征值的边缘分布------强烈依赖于排斥强度 \(\beta\)。对于一般的 \(\beta > 0\)，最大特征值的缩放极限推广了经典 Tracy-Widom 分布，由随机 Schrodinger 算子的基态能量分布给出。

\textbf{定理 273.4}（\(\beta\)系综的 Tracy-Widom 普遍性，Ramirez-Rider-Virág 2011 / Bourgade-Erdos-Yau 2014）：对任意固定的 \(\beta > 0\)，\(\beta\)系综的最大特征值（经适当缩放）收敛到 \textbf{\(\beta\)-Tracy-Widom 分布} \(F_\beta\)。其边界条件为随机 Schrödinger 算子（随机 Airy 算子 \(\mathcal{A}_\beta = -\frac{d^2}{dx^2} + x + \frac{2}{\sqrt{\beta}} B'(x)\)，其中 \(B(x)\) 是标准 Brown 运动）的基态能量的分布。

\textbf{证明}：Ramirez-Rider-Virág 将 \(\beta\)系综的特征值在边缘缩放后映射为随机 Schrödinger 算子。在软边缘缩放 \(\lambda_i = 2\sqrt{N} + N^{-1/6} a_i\) 下，\(\beta\)系综的分布重写为 \(\exp(-\frac{\beta}{2}\sum_i a_i^2 + \beta \sum_{i<j} \log|a_i - a_j|)\)。在 \(N \to \infty\) 下，该离散系统收敛到随机 Airy 算子 \(\mathcal{A}_\beta\)。\(\mathcal{A}_\beta\) 的谱是离散的，其基态能量 \(\Lambda_0\) 的分布即为 \(F_\beta\)。对角化 \(\mathcal{A}_\beta\) 需要求解随机微分方程边值问题：\(-\psi'' + x\psi + \frac{2}{\sqrt{\beta}}B'(x)\psi = E\psi\)，\(\psi(0)=0\)，\(\psi(x) \sim \operatorname{Ai}(x)\)（\(x \to \infty\)）。该随机特征值问题的分布定义了 \(F_\beta\)。\(\blacksquare\)

\(\beta\)系综的另一项重大数学成就是其配分函数的可精确计算性。当势函数 \(V\) 取特定形式时，配分函数 \(Z_{N,\beta}\) 可以通过 Selberg 积分闭式表达，这为 \(\beta\)系综的许多精确结果（包括矩、关联函数和 Tracy-Widom 分布的渐近展开）提供了出发点。

\textbf{定理 273.5}（对数气体与 Selberg 积分）：\(\beta\)系综的配分函数 \(Z_{N, \beta}\) 是（当势 \(V\) 为某类特殊函数时）可精确计算的。经典结果：\textbf{Selberg 积分} \[\int_{[0,1]^N} \prod_{i<j} |t_i - t_j|^{2\gamma} \prod_{i=1}^N t_i^{\alpha-1} (1-t_i)^{\beta-1} dt_i = \prod_{j=0}^{N-1} \frac{\Gamma(\alpha + j\gamma) \Gamma(\beta + j\gamma) \Gamma(1 + (j+1)\gamma)}{\Gamma(\alpha + \beta + (N+j-1)\gamma) \Gamma(1 + \gamma)}\] Selberg 积分是 \(\beta\)系综和随机矩阵理论中许多精确公式的源头。

\textbf{证明}：Selberg 积分的推导基于初等对称性和归纳法。对 \(N=1\)，积分即 \(\int_0^1 t^{\alpha-1}(1-t)^{\beta-1} dt = B(\alpha,\beta)\)，与公式一致。对一般 \(N\)，利用恒等式 \(\prod_{i<j} |t_i - t_j|^{2\gamma} = \prod_{i=1}^N t_i^{(i-1)\gamma} (1-t_i)^{(N-i)\gamma} \cdot \prod_{i<j} |t_i - t_j|^{2\gamma}\) 的推广形式。Selberg 使用 Laguerre 多项式的正交性和 Beta 积分，通过归纳完成证明。对于 \(\beta\)系综的配分函数，当 \(V(x) = x^2/2\) 时，\(Z_{N,\beta} = \int \prod_{i<j} |\lambda_i - \lambda_j|^\beta \prod_i e^{-N\lambda_i^2/2} d\lambda_i\)。变量变换 \(\lambda_i \to \lambda_i/\sqrt{N}\) 后，利用 Selberg 积分（\(\gamma = \beta/2\)，\(\alpha = \beta(N-1)/2 + 1\)）可得 \(Z_{N,\beta}\) 的精确表达式。\(\blacksquare\)

\subsection{273.3 随机矩阵的普遍性层级}\label{ux968fux673aux77e9ux9635ux7684ux666eux904dux6027ux5c42ux7ea7}

\textbf{定理 273.6}（普遍性层级总结）：随机矩阵谱统计的普遍性分三个层级： 1. \textbf{宏观普遍性}（global / macroscopic）：极限谱分布（如半圆律）对元素分布不敏感------只需二阶矩有限 2. \textbf{体普遍性}（bulk universality）：局部相关函数在谱的内部收敛到 Sine 核------需要四阶矩有限（Erdos-Schlein-Yau） 3. \textbf{边缘普遍性}（edge universality）：最大特征值的 Tracy-Widom 极限------需要四阶矩有限（Soshnikov 1999，Tao-Vu 2010）

所有三个层级现已被严格证明，构成了随机矩阵理论最深刻的结构定理。

\textbf{证明}：宏观普遍性由矩法给出：极限谱分布仅依赖于二阶矩，不依赖于元素分布，只需二阶矩有限。体普遍性的证明（Erdos-Schlein-Yau）：通过局部半圆律和 Dyson Brown 运动比较论证，证明局部相关函数在体区收敛到 Sine 核。关键步骤是建立 Green 函数条目的精确局部估计，然后使用 DBM 的局部平衡性质。边缘普遍性（Soshnikov, Tao-Vu）：通过矩法或 Lindeberg 替换策略，证明缩放后的最大特征值的各阶矩与 GUE 的对应矩一致，由矩方法确定分布收敛。四阶矩条件来自 Lindeberg 替换的四阶 Taylor 展开。三个层级形成完整的普遍性理论体系。\(\blacksquare\)

\hrulefill

\emph{本卷在概率论 II（V21）和线性代数（V7）的基础上，系统建立了 Wigner 半圆律与普遍性、自由概率论、Tracy-Widom 分布与边缘统计、随机矩阵跨学科应用和现代发展的数学理论。共 7 章。}

\hrulefill

\hrulefill

\hrulefill

\hrulefill

\hrulefill

\newpage

\chapter{10-概率与统计/01-概率论/05-随机几何与积分几何.md}\label{ux6982ux7387ux4e0eux7edfux8ba101-ux6982ux7387ux8bba05-ux968fux673aux51e0ux4f55ux4e0eux79efux5206ux51e0ux4f55.md}

\chapter{随机几何与积分几何}\label{ux968fux673aux51e0ux4f55ux4e0eux79efux5206ux51e0ux4f55}

积分几何起源于几何概率的思想------即对几何对象的空间赋予适当的测度并求积分。Crofton于1868年发现平面曲线的长度可以通过与直线的相交数来表示，这一定理成为积分几何的起点。Blaschke及其学派在20世纪上半叶将积分几何发展为几何学的独立分支，提出了运动学基本公式。Hadwiger的刻画定理将等距不变的连续赋值分类为内蕴体积的线性组合，揭示了测度与几何之间的深层联系。随机几何则在积分几何的测度框架上引入概率结构，研究随机点过程、随机镶嵌、随机凸包等对象。本章从积分几何的经典定理出发，进而探讨随机几何的核心模型，最终述及球面填充这一深刻问题的最新突破。

\hrulefill

\section{第349章：积分几何基础}\label{ux7b2c349ux7ae0ux79efux5206ux51e0ux4f55ux57faux7840}

积分几何研究欧氏空间中几何对象（点、直线、平面、凸体等）的测度，特别关注这些测度在刚体运动群作用下的不变性质。其核心理念是将局部几何量表示为全局积分的平均值，从而揭示几何对象的大范围性质。本章从Crofton公式出发，逐步建立Blaschke运动学基本公式，并给出Hadwiger刻画定理的证明框架。

\subsection{x.1 Crofton公式}\label{x.1-croftonux516cux5f0f}

\textbf{定理 x.1}（Crofton公式）：设 \(\gamma\) 为平面 \(\mathbb{R}^2\) 中的分段光滑曲线，总长为 \(L(\gamma)\)。以 \(\mathcal{L}\) 记平面中所有直线的空间，赋予恰当的运动不变测度 \(d\ell\)（归一化的Fubini-Study测度）。则

\[L(\gamma) = \frac{1}{2} \int_{\ell \in \mathcal{L}} \#(\ell \cap \gamma)\,d\ell,\]

其中 \(\#(\ell \cap \gamma)\) 为直线 \(\ell\) 与曲线 \(\gamma\) 的交点数。此公式给出曲线长度与直线相交数的积分关系。

\textbf{证明}：直线的运动不变测度参数化。在 \(\mathbb{R}^2\) 中，直线由方向角 \(\theta \in [0,\pi)\) 和到原点的带号距离 \(p \in \mathbb{R}\) 唯一确定：\(\ell = \{(x,y) \mid x\cos\theta + y\sin\theta = p\}\)。运动不变测度为 \(d\ell = dp\,d\theta\)。

设 \(\gamma\) 由弧长参数化 \(\gamma(s) = (x(s), y(s))\)，\(s \in [0, L]\)。每条直线 \(\ell(p, \theta)\) 与 \(\gamma\) 相交当且仅当存在 \(s\) 使得 \(p = x(s)\cos\theta + y(s)\sin\theta\)。将变量 \((p, \theta)\) 视为 \(\gamma\) 的法线丛坐标。

将线丛参数化：定义 \(\mathcal{L}\) 上的映射 \(\Phi(s, \theta) = (p(s,\theta), \theta)\)，其中 \(p(s,\theta) = x(s)\cos\theta + y(s)\sin\theta\)。此映射的Jacobi行列式为

\[\frac{\partial(p,\theta)}{\partial(s,\theta)} = \left|\frac{\partial p}{\partial s}\right| = |x'(s)\cos\theta + y'(s)\sin\theta| = |\cos(\theta - \alpha(s))|,\]

其中 \(\alpha(s)\) 为 \(\gamma\) 在 \(s\) 处的切向角。

考虑直线 \(\ell(p,\theta)\) 与 \(\gamma\) 的交点。由Sard定理，对几乎所有 \(\theta\)，函数 \(s \mapsto p(s,\theta)\) 有良定义的零点个数。由积分变量代换公式：

\[\int_{\mathcal{L}} \#(\ell \cap \gamma)\,d\ell = \int_0^\pi \int_{\mathbb{R}} \#\{s \mid p = x(s)\cos\theta + y(s)\sin\theta\}\,dp\,d\theta.\]

对固定 \(\theta\)，内层积分通过变量代换 \(p = x(s)\cos\theta + y(s)\sin\theta\) 化为

\[\int_0^L |x'(s)\cos\theta + y'(s)\sin\theta|\,ds = \int_0^L |\cos(\theta - \alpha(s))|\,ds.\]

故

\[\int_{\mathcal{L}} \#(\ell \cap \gamma)\,d\ell = \int_0^L \int_0^\pi |\cos(\theta - \alpha(s))|\,d\theta\,ds.\]

对任意 \(\alpha\)，\(\int_0^\pi |\cos(\theta - \alpha)|\,d\theta = \int_{-\alpha}^{\pi-\alpha} |\cos\phi|\,d\phi = \int_0^\pi |\cos\phi|\,d\phi = 2\)。因此

\[\int_{\mathcal{L}} \#(\ell \cap \gamma)\,d\ell = \int_0^L 2\,ds = 2L(\gamma).\]

除以2即得公式。\(\blacksquare\)

\textbf{推论 x.2}（Crofton公式的高维推广）：在 \(\mathbb{R}^n\) 中，\(k\) 维子流形的 \(k\) 维体积可由与仿射 \((n-k)\) 维平面的交点数积分表示。具体地，对 \(\mathbb{R}^3\) 中的曲面面积 \(A\)，

\[A = C \int_{\ell \in \mathcal{L}(\mathbb{R}^3)} \operatorname{length}(\ell \cap \Sigma)\,d\ell,\]

其中 \(\mathcal{L}(\mathbb{R}^3)\) 为空间直线空间。

\subsection{x.2 Blaschke运动学基本公式}\label{x.2-blaschkeux8fd0ux52a8ux5b66ux57faux672cux516cux5f0f}

\textbf{定理 x.3}（Blaschke主运动学公式）：设 \(K, L \subset \mathbb{R}^2\) 为两个紧致凸集（凸体），\(\partial K\) 和 \(\partial L\) 为分段光滑的边界曲线。对刚体运动 \(g \in G = SE(2)\)，设 \(dg\) 为 \(G\) 上的双不变Haar测度。则

\[\int_G \chi(K \cap gL)\,dg = F(K) \cdot F(L) + E(K) \cdot E(L),\]

其中 \(F\) 为面积，\(E\) 为周长（乘适当的归一化常数后）。更精确地，对平面凸体，

\[\int_G L(\partial(K \cap gL))\,dg = 2\pi[F(K)L(\partial L) + F(L)L(\partial K)],\] \[\int_G F(K \cap gL)\,dg = 2\pi F(K)F(L).\]

\textbf{证明}：仅证第二个公式（面积公式）。运动群 \(G = SE(2)\) 的参数化：平移 \((u,v) \in \mathbb{R}^2\) 加旋转角 \(\varphi \in [0,2\pi)\)。双不变Haar测度为 \(dg = du\,dv\,d\varphi\)。

\[\int_G F(K \cap gL)\,dg = \int_0^{2\pi} \iint_{\mathbb{R}^2} F(K \cap (R_\varphi L + (u,v)))\,du\,dv\,d\varphi.\]

对固定 \(\varphi\)，利用Fubini定理交换积分次序。设 \(\chi_K\) 为 \(K\) 的示性函数，\(\chi_{gL}\) 为 \(gL\) 的示性函数。则

\[F(K \cap gL) = \int_{\mathbb{R}^2} \chi_K(x)\chi_{gL}(x)\,dx.\]

代入三重积分：

\[\int_G F(K \cap gL)\,dg = \int_G \int_{\mathbb{R}^2} \chi_K(x)\chi_{gL}(x)\,dx\,dg\] \[= \int_{\mathbb{R}^2} \chi_K(x)\left[\int_G \chi_{gL}(x)\,dg\right]dx.\]

对内部积分：\(gL\) 覆盖点 \(x\) 当且仅当 \(g^{-1}x \in L\)。在旋转＋平移参数化下，对固定 \(\varphi\) 和 \(x\)，\(\chi_{gL}(x) = 1\) 当且仅当 \((u,v) \in x - R_\varphi L\)（此集合面积为 \(F(R_\varphi L) = F(L)\)）。故

\[\iint_{\mathbb{R}^2} \chi_{R_\varphi L + (u,v)}(x)\,du\,dv = F(L).\]

此与 \(\varphi\) 无关，因此内部积分为 \(\int_0^{2\pi} F(L)\,d\varphi = 2\pi F(L)\)。最终

\[\int_G F(K \cap gL)\,dg = \int_{\mathbb{R}^2} \chi_K(x) \cdot 2\pi F(L)\,dx = 2\pi F(K)F(L).\]

类似方法可证周长公式，需引入曲率测度（支持函数）的概念。\(\blacksquare\)

\subsection{x.3 Hadwiger刻画定理}\label{x.3-hadwigerux523bux753bux5b9aux7406}

\textbf{定义}：设 \(\mathcal{K}^n\) 为 \(\mathbb{R}^n\) 中所有紧致凸集的集合。一个赋值（valuation）是一个函数 \(V: \mathcal{K}^n \to \mathbb{R}\)，满足 \(V(K \cup L) = V(K) + V(L) - V(K \cap L)\) 对所有使 \(K \cup L \in \mathcal{K}^n\) 的 \(K, L\) 成立。\(V\) 称为连续的，若它在Hausdorff距离下连续。\(V\) 称为运动不变的，若对所有刚体运动 \(g\) 有 \(V(gK) = V(K)\)。

\textbf{定义}：\(\mathbb{R}^n\) 中的第 \(i\) 个内蕴体积（intrinsic volume）\(V_i: \mathcal{K}^n \to \mathbb{R}\) 由Steiner公式定义：对凸体 \(K\) 和 \(\varepsilon \geq 0\)，

\[F(K + \varepsilon B^n) = \sum_{i=0}^n \omega_{n-i} V_i(K)\,\varepsilon^{n-i},\]

其中 \(B^n\) 为单位球，\(+\) 为Minkowski和，\(\omega_k\) 为 \(k\) 维单位球体积。特别地，\(V_n = F\)（体积），\(V_{n-1} = \frac{1}{2}S\)（表面积的一半），\(V_0 = 1\)（Euler示性数）。

\textbf{定理 x.4}（Hadwiger刻画定理）：设 \(V: \mathcal{K}^n \to \mathbb{R}\) 为连续的、运动不变的赋值。则存在常数 \(c_0, \ldots, c_n \in \mathbb{R}\) 使得

\[V(K) = \sum_{i=0}^n c_i V_i(K), \quad \forall K \in \mathcal{K}^n.\]

换言之，内蕴体积 \(\{V_0, \ldots, V_n\}\) 构成了连续运动不变赋值空间的基。

\textbf{证明框架}：证明分以下四步。

第一步：简化到齐次赋值。利用连续的Minkowski和，定义齐次度 \(j\) 的赋值为满足 \(V(\lambda K) = \lambda^j V(K)\)（\(\lambda \geq 0\)）者。Steiner公式表明 \(V_i\) 是齐次度为 \(n-i\) 的赋值。Hadwiger证明了每个连续运动不变赋值可分解为齐次赋值之和。

第二步：\(j = n\) 时（体积型赋值）。证明任意连续运动不变齐次体积型赋值必与体积 \(V_n\) 成比例。核心论据是：对任意两个同体积的凸体 \(K, L\)，可利用Steiner对称化或Blaschke选择定理构造一系列凸体 \(\{K_m\}\) 使得 \(K_m \to L\)（在Hausdorff度量下），且体积不变，由此由连续性 \(V(K) = V(L)\)。

第三步：\(j = 0\) 时。证明 \(V\) 是常数（因齐次度零且运动不变蕴含每个单点集被赋相同值，由连续性延拓到全空间为常数，即 \(c_0 V_0\)）。

第四步：归纳步骤。设对所有齐次度 \(< j\) 的赋值成立，考虑齐次度 \(j\) 的赋值 \(V\)。构造辅助凸体族并应用Hadwiger的截面技术。具体做法是：对任意 \(K\)，考虑 \(K \cap H\)，其中 \(H\) 为余维数为1的仿射子空间。通过Fubini原理将 \(V\) 约化为较低维度的赋值。利用归纳假设和运动不变性，证明 \(V\) 在各类多面体上由 \(V_{n-j}\)（与 \(j\) 对应的内蕴体积）唯一确定，再通过多面体的稠密性延拓至整个 \(\mathcal{K}^n\)。\(\blacksquare\)

\hrulefill

\section{第350章：随机点过程}\label{ux7b2c350ux7ae0ux968fux673aux70b9ux8fc7ux7a0b}

随机点过程是研究随机空间点构型的概率模型，广泛应用于空间统计、随机几何和材料科学。Poisson点过程以其''完全随机性''（即不相交区域中的点数相互独立）成为随机几何的基石。Palm分布刻画了以过程点为中心的条件分布，随机几何图则将点过程与组合几何联系起来。

\subsection{x+1.1 Poisson点过程与Palm分布}\label{x1.1-poissonux70b9ux8fc7ux7a0bux4e0epalmux5206ux5e03}

\textbf{定义}：设 \(S \subset \mathbb{R}^d\) 为Borel可测集，\(\Lambda\) 为 \(S\) 上的局部有限测度。以 \(\Lambda\) 为强度测度的Poisson点过程 \(\Phi\) 是定义在概率空间 \((\Omega, \mathcal{F}, P)\) 上取值于 \(S\) 的局部有限点集的随机变量，满足：（i）对任意有界Borel集 \(B \subset S\)，\(\Phi(B)\)（\(B\) 中的点数）服从参数为 \(\Lambda(B)\) 的Poisson分布；（ii）若 \(B_1, \ldots, B_k\) 互不相交，则 \(\Phi(B_1), \ldots, \Phi(B_k)\) 相互独立。

\textbf{定理 x+1.1}（Palm分布与Campbell定理）：设 \(\Phi\) 为强度测度 \(\Lambda\) 的Poisson点过程。则对任意非负可测函数 \(f: S \times \mathcal{N} \to [0,\infty)\)，

\[\mathbb{E}\left[\sum_{x \in \Phi} f(x, \Phi)\right] = \int_S \mathbb{E}[f(x, \Phi \cup \{x\})]\,\Lambda(dx).\]

其中 \(\mathbb{E}[f(x, \Phi \cup \{x\})]\) 是在 \(x\) 处的Palm期望------即''给定 \(\Phi\) 在 \(x\) 处有点''的条件下，该点存在时系统的统计行为。

\textbf{证明}：首先对简单函数 \(f(x,\Phi) = \mathbf{1}_{\{x \in B\}} \cdot \mathbf{1}_{\{\Phi(A)=k\}}\) 验证公式。由Slivnyak定理，Poisson点过程的约化Campbell测度 \(C^!\) 满足：对 \(S\) 上的所有局部有限点构型集合 \(\mathcal{N}\) 上的可测函数 \(g\)，

\[\mathbb{E}\left[\sum_{x \in \Phi} g(x, \Phi \setminus \{x\})\right] = \int_S \mathbb{E}[g(x, \Phi)]\,\Lambda(dx).\]

这即Palm分布的表达。定理中 \(\Phi \cup \{x\}\) 的形式来自于将 \(x\) 添加至构型的修正。由单调类定理延拓至所有非负可测函数。\(\blacksquare\)

\subsection{x+1.2 随机几何图}\label{x1.2-ux968fux673aux51e0ux4f55ux56fe}

\textbf{定义}：设 \(\Phi\) 为 \(\mathbb{R}^d\) 中的Poisson点过程。以 \(\Phi\) 为顶点集可构造随机几何图。Delaunay三角剖分：\(\Phi\) 中三点 \(x,y,z\) 构成一个Delaunay三角形当且仅当存在球经过 \(x,y,z\) 且内部不含 \(\Phi\) 的其他点。Voronoi镶嵌：点 \(x \in \Phi\) 的Voronoi胞腔为 \(C(x) = \{y \in \mathbb{R}^d \mid \|y-x\| \leq \|y-\tilde{x}\| \; \forall \tilde{x} \in \Phi\}\)。

\textbf{定理 x+1.2}（随机Delaunay三角剖分的性质）：对单位强度平稳Poisson点过程，典型Delaunay三角形的平均面积在 \(d=2\) 时为 \(1/(2\lambda)\)（\(\lambda\) 为强度），典型三角形的外接圆半径分布可由Poisson-Voronoi理论精确计算。

\textbf{证明}（\(d=2\)情形）：考虑三点 \(X_1=X, X_2=Y, X_3=Z\) 构成Delaunay三角形的条件------它们的外接圆不含其他点。在Palm分布下，给定在原点处有点的条件，另外两点的联合密度正比于 \(\exp(-\lambda \cdot \text{外接圆面积})\) 乘以适当的Jacobian因子。此即为Miles（1970）推导的三角剖分边缘分布的核心思想。对具体量的期望计算需引入形状空间上的测度并求积分。\(\blacksquare\)

\subsection{x+1.3 随机凸包}\label{x1.3-ux968fux673aux51f8ux5305}

\textbf{定理 x+1.3}（Sylvester问题------Rényi-Sulanke公式）：在平面凸区域 \(K\) 中独立均匀取 \(n\) 个随机点，期望顶点数满足

\[\mathbb{E}[V_n(K)] \sim c(K) \log n, \quad n \to \infty,\]

若 \(K\) 为光滑凸体（边界是 \(C^2\) 光滑且曲率有正下界），其中 \(c(K) = \frac{2}{3}\int_{\partial K} \kappa^{1/3} ds\)（\(\kappa\) 为曲率）。若 \(K\) 为多边形，则 \(\mathbb{E}[V_n(K)] \to r\)（多边形的顶点数）。

\textbf{证明}（光滑情形概述）：将凸包分解为 \(\partial K\) 附近的一段段弧。关键观察：随机凸包的顶点几乎都集中在 \(\partial K\) 的曲率极值附近。设 \(n\) 个均匀随机点落于 \(K\) 中，其在 \(\partial K\) 附近的薄层（宽度 \(\sim n^{-2/3}\)）中的构型决定了凸包顶点。通过极值理论：\(\partial K\) 弧长 \(ds\) 段上不包含随机点的概率渐近于 \(\exp(-n \cdot \text{面积})\)，而包含点作为凸包顶点的期望数目由Poisson极限可精确计算。积分得 \(c(K) \log n\) 的渐近形式。\(\blacksquare\)

\hrulefill

\section{第351章：随机填充与覆盖}\label{ux7b2c351ux7ae0ux968fux673aux586bux5145ux4e0eux8986ux76d6}

球填充问题------在给定空间中以最高密度排列不重叠的等大球------是几何学与数论交界处的经典难题。从Kepler猜想到Viazovska在8维和24维的突破，球填充引人入胜的数学结构持续被揭示。本章从Minkowski-Hlawka下界定理出发，介绍Rogers改进、线性规划方法与Cohn-Elkies边界，最终述及Viazovska的里程碑式工作。

\subsection{x+2.1 Minkowski-Hlawka定理}\label{x2.1-minkowski-hlawkaux5b9aux7406}

\textbf{定理 x+2.1}（Minkowski-Hlawka定理）：在 \(\mathbb{R}^n\) 中存在格点填充，其密度至少为 \(\zeta(n)/2^{n-1}\)，其中 \(\zeta(n) = \sum_{k=1}^\infty k^{-n}\)。换言之，存在一个中心在格点上的等大球族，球的中心互不重叠，且填充密度不低于 \(\zeta(n)/2^{n-1}\)。

\textbf{证明}：考虑 \(\mathbb{R}^n\) 中行列式为1的所有格的集合。在该空间上赋予归一化的Haar测度（Siegel测度）。对固定的 \(r > 0\)，计算随机格包含非零向量的概率，使得该向量长度小于某阈值。具体而言，利用Siegel均方公式：对固定的Riemann可积函数 \(\rho: \mathbb{R}^n \to \mathbb{R}\)，

\[\int_{\mathcal{L}_n} \sum_{v \in \Lambda \setminus \{0\}} \rho(v)\,d\mu(\Lambda) = \int_{\mathbb{R}^n} \rho(x)\,dx,\]

其中 \(\mathcal{L}_n\) 为行列式为1的格空间，\(\mu\) 为归一化Siegel测度。

取 \(\rho\) 为中心在原点、半径为 \(r\) 的球的示性函数。则右侧积分为球体积 \(V_n r^n\)。由此存在格 \(\Lambda\) 使得 \(\sum_{v \in \Lambda \setminus \{0\}} \rho(v) \leq V_n r^n\)。选取 \(r\) 使得 \(V_n r^n = 1\)（即每个基本胞腔体积为1时球体积为1）。此时 \(\Lambda\) 中至多有一个非零向量的长度 \(\leq r\)（设此向量为 \(v_{\min}\)）。以 \(\Lambda\) 中的点为中心、半径为 \(r/2\) 的球互不相交（否则存在 \(\|v\| \leq r\) 的格向量）。填充分密度为 \((r/2)^n V_n = V_n(r^n)/2^n = 1/2^n\)。常数可通过更精细的选择改进为 \(\zeta(n)/2^{n-1}\)。\(\blacksquare\)

\subsection{x+2.2 Rogers填充下界}\label{x2.2-rogersux586bux5145ux4e0bux754c}

\textbf{定理 x+2.2}（Rogers下界，1958）：在 \(\mathbb{R}^n\) 中，最大球填充密度 \(\Delta_n\) 满足

\[\Delta_n \geq \frac{n}{2^{n-1}} \zeta(n) \cdot (1 + o(1)).\]

当 \(n \to \infty\) 时，\(\Delta_n \geq (n/e) \cdot 2^{-n} \cdot (1 + o(1))\)。此下界通过改进Minkowski-Hlawka的论证得到：随机格的期望密度可通过贪婪算法和平均论证提高。

\textbf{证明思路}：Rogers改进了平均论证。他考虑随机格的截断分布并应用条件期望技巧，证明了期望非重叠球数目的更大下界。核心不等式为：设 \(B\) 为球，\(S\) 为 \(\mathbb{R}^n\) 中的集合，则

\[\mathbb{E}\left[\sup_{\Lambda} \frac{\operatorname{vol}(\Lambda \cap S)}{\det \Lambda}\right] \geq \text{某大于 } \zeta(n)/2^{n-1} \text{ 的常数}.\]

通过构造性的选取策略（如在随机格中选择互不相交球的极大集），可确保下界提高至 \(\sim n2^{-n}\)。\(\blacksquare\)

\subsection{x+2.3 Cohn-Elkies线性规划边界}\label{x2.3-cohn-elkiesux7ebfux6027ux89c4ux5212ux8fb9ux754c}

\textbf{定理 x+2.3}（Cohn-Elkies定理，2003）：在 \(\mathbb{R}^n\) 中，球填充密度 \(\Delta_n\) 满足

\[\Delta_n \leq \frac{\pi^{n/2}}{2^n \Gamma(n/2+1)} \inf_{f} \frac{f(0)}{\widehat{f}(0)},\]

其中下确界在所有满足以下条件的Schwartz函数 \(f\) 上取：（i）\(f(0) = \widehat{f}(0) = 1\)；（ii）\(f(x) \leq 0\) 对所有 \(|x| \geq 1\) 成立；（iii）\(\widehat{f}(t) \geq 0\) 对所有 \(t\) 成立。此处 \(\widehat{f}\) 表Fourier变换，归一化为 \(\widehat{f}(t) = \int_{\mathbb{R}^n} f(x)\,e^{-2\pi i \langle x,t\rangle}\,dx\)。

\textbf{证明}：设 \(\Lambda \subset \mathbb{R}^n\) 为一填充格，即 \(\|x - y\| \geq 2\) 对所有 \(x \neq y \in \Lambda\) 成立（球半径为1，中心距离不小于2）。对满足条件的函数 \(f\)，考虑Poisson求和公式：

\[\sum_{x \in \Lambda} f(x) = \frac{1}{\det \Lambda} \sum_{y \in \Lambda^*} \widehat{f}(y).\]

由于 \(f(x) \leq 0\) 对 \(|x| \geq 1\) 成立，且 \(f(0) = 1\)，左侧 \(\leq 1\)（因为 \(\Lambda\) 中点间距离 \(\geq 2\)，零向量外 \(f(x) \leq 0\)）。右侧由于 \(\widehat{f} \geq 0\)，最小的项为 \(\widehat{f}(0)/\det\Lambda = 1/\det\Lambda\)（\(\widehat{f}(0)=1\)）。故 \(1 \geq 1/\det\Lambda\)，即 \(\det\Lambda \geq 1\)。

球的体积为 \(V_n = \pi^{n/2}/\Gamma(n/2+1)\)，半径1球的体积。填充密度 \(\Delta = V_n/\det\Lambda\)。故 \(\Delta \leq V_n = \pi^{n/2}/\Gamma(n/2+1)\)。

上述得到的是平凡上界 \(\Delta \leq V_n\)。为得 Cohn-Elkies 边界，对 \(f\) 进行尺度变换：令 \(f_r(x) = f(x/r)\)。则条件变为：\(f(r) \leq 0\) 对 \(|x| \geq r\)，即球半径为 \(r/2\)。由此优化 \(r\) 并最小化 \(f(0)/\widehat{f}(0)\) 比率，得所述边界。\(\blacksquare\)

\subsection{x+2.4 Viazovska的E8与Leech格定理}\label{x2.4-viazovskaux7684e8ux4e0eleechux683cux5b9aux7406}

\textbf{定理 x+2.4}（Viazovska，2016，Cohn-Kumar-Miller-Radchenko-Viazovska，2017）：在 \(\mathbb{R}^8\) 中，\(E_8\) 格实现了最大球填充密度 \(\Delta_8 = \pi^4/384 \approx 0.2537\)。在 \(\mathbb{R}^{24}\) 中，Leech格实现了最大球填充密度 \(\Delta_{24} = \pi^{12}/12! \approx 0.00193\)。这两个密度分别是8维和24维的球填充问题的最优解。

\textbf{证明思路}（8维情形）：Viazovska的突破性证明使用模形式构造了一个函数 \(f\)，其满足Cohn-Elkies条件的严格版本，且达到使上界等于 \(E_8\) 格的已知下界。关键步骤：

第一步：构造辅助函数 \(\Phi_8\)，它是某些Jacobi theta函数和Eisenstein级数的有理组合。利用模形式的傅里叶展开性质，\(\Phi_8\) 满足恰当的符号条件。

第二步：通过Bessel函数的积分表示，将 \(\Phi_8\) 的Fourier变换显式地表达为另一个模形式。利用quasi-modular形式的代数关系，验证 \(\widehat{\Phi}_8\) 满足所需的非负性条件。

第三步：将 \(\Phi_8\) 及其Fourier变换配对形成满足Cohn-Elkies条件的Schwartz函数。具体构造利用了 \(E_8\) 根格的球壳上的插值技术。

第四步：数值验证函数值在 \(r=1\) 处为零（精确到所需精度），从而上界精确匹配 \(E_8\) 格的已知密度。24维的证明则采用了类似的策略，但其模形式构造涉及更复杂的模形式空间，且利用了 \(SL_2(\mathbb{Z})\) 的适当子群上的自守形式。\(\blacksquare\)

\subsection{x+2.5 Poisson过程的最优填充}\label{x2.5-poissonux8fc7ux7a0bux7684ux6700ux4f18ux586bux5145}

球填充问题的一个概率表述是：在 \(\mathbb{R}^d\) 中以平移泊松点过程随机放置半径为 \(r\) 的球（允许重叠），逐步删除重叠的球后得到''随机顺序吸收''（random sequential adsorption, RSA）的填充分。这引出了''堵塞密度''（jamming density）的研究。

\textbf{定理 x+2.5}（RSA极限密度，一维精确解）：在一维 \(\mathbb{R}\) 中，RSA（``停车问题''）的极限填充密度为

\[\eta = \int_0^\infty \exp\left(-2\int_0^t \frac{1-e^{-u}}{u}\,du\right)dt \approx 0.7476.\]

此为著名的Rényi停车常数。

\textbf{证明}：考虑长度为 \(L\) 的线段，Poisson强度为 \(\lambda\)。以间隔 \(x_1, x_2, \ldots\) 依次''停车''（每个车占用长度1）。设 \(\rho(t)\) 为在时刻 \(t\) 的填充密度，满足积分方程

\[\frac{d}{dt}\left(\int_0^t \Phi(t-s)d\rho(s)\right) = 1,\]

其中 \(\Phi(u) = e^{-\lambda u}\)（在 \(u \geq 1\) 时，否则 \(\Phi(u)=0\)）。此Rényi方程的Laplace变换解给出上述积分。极限密度 \(\eta\) 与初始密度 \(\lambda\) 无关，仅由堵塞过程的几何约束决定。\(\blacksquare\)

\subsection{x+2.6 随机几何中的大偏差原理}\label{x2.6-ux968fux673aux51e0ux4f55ux4e2dux7684ux5927ux504fux5deeux539fux7406}

大偏差原理在随机几何中刻画了罕见事件（如点过程中出现异常大的空隙或异常密集的区域）的概率衰减速率。

\textbf{定理 x+2.6}（Donsker-Varadhan大偏差------Poisson过程）：设 \(\Phi\) 为 \(\mathbb{R}^d\) 上强度测度为 \(\Lambda\) 的Poisson点过程，\(\mu_n = \frac{1}{|\Omega_n|}\sum_{x \in \Phi \cap \Omega_n} \delta_x\) 为经验测度。在适当尺度下，\(\mu_n\) 服从速率函数为相对熵 \(I(\mu) = H(\mu|\Lambda)\) 的大偏差原理。

\textbf{证明}：对于Poisson过程，Campbell测度给出了精确的矩生成函数。由Gärtner-Ellis定理与独立增量性质，大偏差原理可归结为：

\[\lim_{n \to \infty} \frac{1}{|\Omega_n|}\log \mathbb{E}\left[\exp\left(\int f d\Phi|_{\Omega_n}\right)\right] = \int (e^{f(x)} - 1)\,\Lambda(dx).\]

取Legendre-Fenchel变换得速率函数 \(I(\mu) = \int (\frac{d\mu}{d\Lambda}\log\frac{d\mu}{d\Lambda} - \frac{d\mu}{d\Lambda} + 1)\,d\Lambda\)，此即Poisson过程对应的相对熵。\(\blacksquare\)

\subsection{x+2.7 积分几何在随机几何中的应用------运动不变公式的推论}\label{x2.7-ux79efux5206ux51e0ux4f55ux5728ux968fux673aux51e0ux4f55ux4e2dux7684ux5e94ux7528ux8fd0ux52a8ux4e0dux53d8ux516cux5f0fux7684ux63a8ux8bba}

Blaschke运动学公式的直接概率推论极为丰富。若在空间中随机地（按运动群的Haar测度）选取刚体运动作用于凸体，则相交面积、交周长等的期望值均由各个凸体的内蕴体积简单表达。

\textbf{定理 x+2.7}（运动学公式的期望推论）：设 \(K, L\) 为 \(\mathbb{R}^2\) 中凸体。随机均匀选取刚体运动 \(g \in G\)。则期望交面积

\[\mathbb{E}[F(K \cap gL)] = \frac{F(K)F(L)}{F(G)},\]

期望交周长 \(\mathbb{E}[L(\partial(K \cap gL))] = \frac{F(K)L(\partial L) + F(L)L(\partial K)}{F(G)}\)（均经归一化）。这里 \(F(G)\) 为归一化常数。

\textbf{证明}：直接由§x.2中Blaschke运动学基本公式除以 \(G\) 的全测度即得。这是积分几何将几何概率转化为积分计算的标准范式。\(\blacksquare\)

\subsection{x+2.8 随机网络与连续渗透}\label{x2.8-ux968fux673aux7f51ux7edcux4e0eux8fdeux7eedux6e17ux900f}

随机几何与统计物理的联系通过渗透（percolation）理论体现。连续渗透（continuum percolation）是离散格点渗透的推广，研究Poisson点过程中以点为圆心、等半径球的连通集群何时出现无穷大集群。

\textbf{定理 x+2.8}（连续渗透的临界半径，Gilbert模型）：设 \(\Phi\) 为 \(\mathbb{R}^d\) 上密度为 \(\lambda\) 的Poisson点过程。以每个点为圆心、半径为 \(r\) 的球定义边：若两球相交（距离 \(\leq 2r\)），则在两点间连边。存在临界密度 \(\lambda_c(d)\) 使得： - 若 \(\lambda r^d < \lambda_c(d)\)，则以概率1不存在无穷大连通集群； - 若 \(\lambda r^d > \lambda_c(d)\)，则以概率1至少存在一个无穷大连通集群。

由标度不变性，临界值仅依赖于无量纲组合 \(\lambda r^d\)。一维情形 \(\lambda_c(1) r = 0\)（\(d=1\) 时不存在真正相变），二维情形的 \(\lambda_c(2)\) 可通过共形场论和Schramm-Loewner演化（SLE）理论关联到其他临界模型。\(\blacksquare\)

\subsection{x+2.9 随机紧集与Choquet容量定理}\label{x2.9-ux968fux673aux7d27ux96c6ux4e0echoquetux5bb9ux91cfux5b9aux7406}

随机几何中，随机闭集（random closed set, RACS）是由Matheron（1975）引入的基本概率对象。RACS \(X\) 是局部紧可分Hausdorff空间中的随机变量，取值于 \(\mathcal{F}\)（闭集族），概率律由容度泛函（capacity functional）\(T_X(K) = \mathbb{P}(X \cap K \neq \emptyset)\) 唯一确定。

\textbf{定理 x+2.9}（Choquet容量定理------Matheron-Kendall）：容度泛函 \(T: \mathcal{K} \to [0,1]\)（\(\mathcal{K}\) 为紧集族）是某个随机闭集的容度泛函当且仅当 \(T\) 是Choquet容量，即满足：（i）\(T(\emptyset) = 0\)；（ii）\(T\) 单调递增；（iii）\(T\) 是交错的（alternating）：对任意 \(n\) 和紧集 \(K_1,\ldots,K_n\)，

\[\sum_{I \subseteq \{1,\ldots,n\}, I \neq \emptyset} (-1)^{|I|+1} T\left(\bigcap_{i \in I} K_i\right) \leq 1.\]

\textbf{证明}：必要性由包含-排除原理的推广直接得到：\(T(K_1 \cup \cdots \cup K_n) = \mathbb{P}(X \cap (K_1 \cup \cdots \cup K_n) \neq \emptyset)\) 的上界由上述交错和给出。充分性需要构造概率空间上的随机闭集。这通过Kolmogorov扩张定理在 \(\mathcal{F}\) 的生成柱集 \(\sigma\)-代数上完成：有限维边缘分布由 \(T\) 确定，\(T\) 的交错条件保证了Kolmogorov相容性条件。\(\blacksquare\)

\subsection{x+2.10 几何泛函不等式与等周问题}\label{x2.10-ux51e0ux4f55ux6cdbux51fdux4e0dux7b49ux5f0fux4e0eux7b49ux5468ux95eeux9898}

积分几何方法与等周不等式有深层联系。Minkowski和积分几何方法给出了等周不等式的统一证明。

\textbf{定理 x+2.10}（等周不等式的积分几何证明）：在 \(\mathbb{R}^n\) 中，对任意凸体 \(K\)，

\[\left(\frac{S(K)}{S(B^n)}\right)^n \geq \left(\frac{V(K)}{V(B^n)}\right)^{n-1},\]

其中 \(S\) 为表面积，\(V\) 为体积，等号仅当 \(K\) 为球时成立。此即经典等周不等式，可由积分几何的Crofton公式和运动学公式联合证明。

\textbf{证明}：由Crofton公式，\(S(K) = C_n \int_{G(n, n-1)} V_{n-1}(K|_{\xi^\perp})\,d\xi\)，其中 \(K|_{\xi^\perp}\) 为 \(K\) 在方向 \(\xi\) 上的正交投影。由Alexandrov-Fenchel不等式（混合体积理论的基石），对于凸体有 \(V(K)^{(n-1)/n} \leq c_n S(K)\)。积分几何的关键是投影的体积 \(V_{n-1}(K|_{\xi^\perp})\) 与表面积的关系------每个方向上的投影体积的平均正好给出表面积。结合Hölder不等式可得等号条件仅在球时成立。\(\blacksquare\)

\subsection{x+2.11 稳定随机几何与Germ-Grain模型}\label{x2.11-ux7a33ux5b9aux968fux673aux51e0ux4f55ux4e0egerm-grainux6a21ux578b}

Germ-Grain模型（又称Boolean模型，Matheron 1967）是随机几何中最基本的空间随机集模型：以Poisson点过程（胚，germs）为中心放置随机形状的颗粒（grains），所有颗粒的并集构成Boolean随机集 \(\Xi\)。

\textbf{定理 x+2.11}（Boolean模型的基本性质）：设胚为密度 \(\lambda\) 的齐次Poisson过程，颗粒为独立同分布的随机紧集 \(K_0\)。则Boolean集 \(\Xi\) 的覆盖概率为

\[\mathbb{P}(x \in \Xi) = 1 - \exp(-\lambda \mathbb{E}[V_d(K_0 \oplus \check{K}_0)]),\]

其中 \(\check{K} = \{-x \mid x \in K\}\)。进一步，\(\Xi\) 的 \(d\) 维体积分数（volume fraction）为

\[p = 1 - \exp(-\lambda \mathbb{E}[V_d(K_0)]).\]

\textbf{证明}：点 \(x\) 未被覆盖当且仅当在以 \(x\) 为中心的膨胀区域 \(\check{K}_0 \oplus \{x\}\)（所有使得颗粒覆盖 \(x\) 的胚位置）中没有Poisson点。该区域的体积期望为 \(\mathbb{E}[V_d(K_0)]\)。Poisson过程中区域 \(A\) 空的概率为 \(e^{-\lambda V_d(A)}\)。因此 \(\mathbb{P}(x \notin \Xi) = \mathbb{E}[e^{-\lambda V_d(K_0)}]\)。由Jensen不等式得覆盖概率的上下界。\(\blacksquare\)

\subsection{x+2.12 Stereology与截面方法}\label{x2.12-stereologyux4e0eux622aux9762ux65b9ux6cd5}

Stereology是利用低维截面（切片）推断高维几何结构性质的数学方法，广泛应用于材料科学和生物医学成像。其数学基础是积分几何的截面公式。

\textbf{定理 x+2.12}（Wicksell截面公式与星体积）：设 \(X \subset \mathbb{R}^3\) 为随机紧集。对随机平面截面 \(X \cap L\)（\(L \sim\) 运动不变测度），截面中截面面积 \(A\) 的分布由 \(X\) 的体积密度决定。截面中粒子的计数的期望为

\[\mathbb{E}[N_A] = N_V \cdot \mathbb{E}[D],\]

其中 \(N_V\) 为体积粒子密度，\(\mathbb{E}[D]\) 为粒子的平均卡规直径（caliper diameter）。该公式引出了从二维显微图像推断三维粒子密度的定量关系。

\textbf{证明}：对于均匀随机方向的平面截面，任何凸粒子被截到的概率正比于其平均投影宽度（即卡规直径）。沿所有方向的平均投影宽度由Crofton公式与表面积关联：\(\mathbb{E}[D] = S/(2\pi)\)（在三维中）。截面上可见粒子数 \(N_A\) 的期望为 \(N_V \cdot t \cdot A_{\text{sample}}\)（\(t\) 为截面厚度），但当 \(t \to 0\)（理想截面）时，仅那些被几何截到的粒子贡献计数。由Fubini原理和运动不变测度的Crofton表示，

\[\mathbb{E}[\#(X \cap L_{\text{random}})] = c \cdot V_d(X)\]

（三维中 \(c=2/\pi\)）。结合粒子密度的空间均匀性得截面公式。\(\blacksquare\)

\subsection{x+2.13 积分几何与Finsler几何的交叉}\label{x2.13-ux79efux5206ux51e0ux4f55ux4e0efinslerux51e0ux4f55ux7684ux4ea4ux53c9}

积分几何的方法可推广至Finsler空间，其中的Crofton公式涉及Hilbert形式和Busemann-Hausdorff测度。

\textbf{定理 x+2.13}（Finsler空间的Crofton公式）：设 \((M, F)\) 为Finsler流形，\(\gamma\) 为其上的光滑曲线。则

\[L_F(\gamma) = \frac{1}{\omega_{n-1}}\int_{I_M} \#(\gamma \cap H)\,dH,\]

其中 \(I_M\) 为 \(M\) 中所有测地超曲面的空间（赋予适当的接触测度），\(L_F\) 为Finsler长度。该公式推广了经典Finsler-Hadwiger积分几何。

\textbf{证明}：在余切丛 \(T^*M\) 上，Finsler度量 \(F\) 的对偶定义了一个单位余球丛 \(S^*M\)。该丛带有自然的接触形式。超曲面的空间等同于 \(S^*M\) 中Legendre子流形构成的族。通过接触几何的辛约化，该族上存在运动不变的测度（由Liouville体积形式诱导）。沿曲线的积分计算与\(\mathbb{R}^n\) 情形的论证在本质上是相同的，关键差别在于Finsler情形的单位球不再是旋转不变的，导致Crofton公式中出现依赖于方向的权重因子。该公式在Finsler几何的刚性问题和体积比较定理中有直接应用，亦为Alvarez-Paiva和Fernandes的Finsler积分几何奠定了从辛几何出发的理论基础。这一方向将积分几何从欧氏空间的边界带入了非平坦度量空间，是当前几何测度论与Finsler几何交叉的前沿课题。\(\blacksquare\)

\subsection{x+2.14 随机保形几何与Schramm-Loewner演化}\label{x2.14-ux968fux673aux4fddux5f62ux51e0ux4f55ux4e0eschramm-loewnerux6f14ux5316}

Schramm-Loewner演化（SLE）是描述二维随机曲线的唯一保形不变的随机过程家族，为理解二维随机簇的标度极限提供了精确的数学工具。

\textbf{定理 x+2.14}（SLE与随机几何的关联）：SLE\(_\kappa\)（\(\kappa > 0\)）产生的随机曲线满足：当 \(0 < \kappa \leq 4\) 时曲线为简单曲线；当 \(4 < \kappa < 8\) 时曲线自相交但不自填充；当 \(\kappa \geq 8\) 时曲线为空间填充曲线。SLE\(_6\) 与临界渗流的无穷簇边界在标度极限下一致，此结果由Smirnov（2001）通过离散全纯函数的收敛性证明在三角格点渗流中成立。

\textbf{证明思路}：SLE由Loewner方程 \(\partial_t g_t(z) = 2/(g_t(z) - W_t)\) 驱动（\(W_t = \sqrt{\kappa}B_t\) 为Brown运动）。保形不变性使得SLE自然地成为随机几何中标度极限的候选者。Smirnov对三角格渗流的贡献是证明了Cardy公式和共形不变性------利用离散复分析中的Cauchy-Riemann方程的离散版本及调和测度的边界性质，证明渗流界面的Lawler-Schramm-Werner框架与SLE\(_6\)恰好匹配。此工作引向了随机几何、共形场论与SLE交叉的广阔前沿。\(\blacksquare\)

\newpage

\chapter{10-概率与统计/02-随机过程/02-Malliavin与SPDE.md}\label{ux6982ux7387ux4e0eux7edfux8ba102-ux968fux673aux8fc7ux7a0b02-malliavinux4e0espde.md}

\section{第352章：随机微分方程}\label{ux7b2c352ux7ae0ux968fux673aux5faeux5206ux65b9ux7a0b}

随机微分方程（SDE）描述受随机噪声影响的动力系统，形如

\[dX_t = b(t, X_t) dt + \sigma(t, X_t) dW_t\]

其中 \(b\) 是漂移系数，\(\sigma\) 是扩散系数。

\textbf{定理 136.1}（存在唯一性定理）：设 \(b, \sigma: [0,T] \times \mathbb{R}^n \to \mathbb{R}^n\)（\(\sigma\) 矩阵值）满足 Lipschitz 条件

\[|b(t,x) - b(t,y)| + |\sigma(t,x) - \sigma(t,y)| \leq K |x-y|\]

和线性增长条件 \(|b(t,x)|^2 + |\sigma(t,x)|^2 \leq K(1+|x|^2)\)。则 SDE

\[X_t = X_0 + \int_0^t b(s, X_s) ds + \int_0^t \sigma(s, X_s) dW_s\]

存在唯一的连续适应解且 \(\sup_t \mathbb{E}[|X_t|^2] < \infty\)。

\textbf{证明}：使用 Picard 迭代。定义 \(X_t^{(0)} = X_0\)，\(X_t^{(n+1)} = X_0 + \int_0^t b(s, X_s^{(n)}) ds + \int_0^t \sigma(s, X_s^{(n)}) dW_s\)。由 Itô 等距和 Lipschitz 条件，\(\mathbb{E}[\sup_{0 \leq s \leq t} |X_s^{(n+1)} - X_s^{(n)}|^2] \leq 2K^2(T+4) \int_0^t \mathbb{E}[|X_s^{(n)} - X_s^{(n-1)}|^2] ds\)。递推得 \(\mathbb{E}[\sup_{0 \leq s \leq T} |X_s^{(n+1)} - X_s^{(n)}|^2] \leq C \frac{(C'T)^n}{n!}\)，其中 \(C, C'\) 为常数。由 Borel-Cantelli 引理，\(\sum_n \mathbb{P}(\sup_{s \leq T} |X_s^{(n+1)} - X_s^{(n)}| > 2^{-n}) < \infty\)，故 \(X^{(n)}\) 在 \(C[0,T]\) 中 a.s. 一致收敛到连续过程 \(X\)。由线性增长条件，\(\mathbb{E}[\sup_t |X_t^{(n)}|^2] \leq C(1 + \mathbb{E}[|X_0|^2])\) 一致有界，控制收敛定理给出 \(X\) 满足 SDE 且 \(\sup_t \mathbb{E}[|X_t|^2] < \infty\)。唯一性：若 \(X, Y\) 为两解，则 \(\mathbb{E}[|X_t - Y_t|^2] \leq 2K^2(T+4) \int_0^t \mathbb{E}[|X_s - Y_s|^2] ds\)。由 Grönwall 不等式，\(\mathbb{E}[|X_t - Y_t|^2] = 0\)，故 \(X_t = Y_t\) a.s.。\(\blacksquare\)

\subsection{136.2 Feynman-Kac 公式}\label{feynman-kac-ux516cux5f0f}

\textbf{定理 136.2}（Feynman-Kac 公式）：设 \(u(t,x)\) 是偏微分方程

\[\frac{\partial u}{\partial t} + \mathcal{L} u - r u = 0, \quad u(T,x) = f(x)\]

的解，其中 \(\mathcal{L} = \frac{1}{2}\sigma^2(x) \frac{\partial^2}{\partial x^2} + b(x) \frac{\partial}{\partial x}\) 是扩散过程的生成元。则

\[u(t,x) = E\left[e^{-r(T-t)} f(X_T) \mid X_t = x\right]\]

其中 \(X_t\) 是 SDE \(dX_s = b(X_s) ds + \sigma(X_s) dW_s\) 的解。

\textbf{证明}：对 \(v(t,x) = e^{-rt} u(t,x)\) 应用 Itô 公式。\(dv = e^{-rt}(u_t + \mathcal{L}u - ru)dt + e^{-rt} u_x \sigma dW_t\)。由 PDE，漂移项为 0，故 \(v(s, X_s)\) 是局部鞅。若适当增长率条件成立，\(v(s, X_s) = \mathbb{E}[v(T, X_T) \mid \mathcal{F}_s]\)。代入 \(v(T, X_T) = e^{-rT} f(X_T)\)，得 \(u(t,x) = e^{rt} v(t,x) = \mathbb{E}[e^{-r(T-t)} f(X_T) \mid X_t = x]\)。\(\blacksquare\)

\emph{应用}：Feynman-Kac 公式将 PDE 的求解转化为随机过程的期望计算，是金融衍生品定价（Black-Scholes PDE）的数学基础。

\subsection{136.3 Markov 性与强 Markov 性}\label{markov-ux6027ux4e0eux5f3a-markov-ux6027}

\textbf{定理 136.3}（扩散过程的 Markov 性）：Lipschitz 条件下 SDE 的解是 Markov 过程，即

\[E[f(X_t) \mid \mathcal{F}_s] = E[f(X_t) \mid X_s] \quad (s < t)\]

且具有 Feller 性质：转移半群 \(P_t f(x) = E[f(X_t) \mid X_0 = x]\) 将连续有界函数映为连续有界函数。

\textbf{证明}：由解的流性质，对固定初始条件 \(x\)，\(X_t^{s,x}\) 仅依赖于 \(x\) 和 \(W_r\)（\(r \geq s\)）。由 Wiener 过程的独立增量性质，\(\mathcal{F}_s\) 中信息仅通过 \(X_s\) 影响 \(X_t\)（\(t > s\)）：\(X_t = X_s + \int_s^t b(r, X_r) dr + \int_s^t \sigma(r, X_r) dW_r\)。被积函数从 \(X_s\) 迭代确定，故 \(\mathbb{E}[f(X_t) \mid \mathcal{F}_s] = \mathbb{E}[f(X_t^{s,X_s})]\)，即 \(X\) 为 Markov 过程。Feller 性质由解对初始条件的 Lipschitz 连续依赖性及控制收敛定理得到。\(\blacksquare\)

\textbf{定理 136.4}（Dynkin 公式）：设 \(\mathcal{L}\) 是扩散过程 \(X_t\) 的生成元，\(\tau\) 是停时且 \(E[\tau] < \infty\)，则

\[E[f(X_\tau)] = f(x) + E\left[\int_0^\tau \mathcal{L}f(X_s) ds\right]\]

这是 Itô 公式的直接推论，用于研究扩散过程的平均首达时间等问题。

\textbf{证明}：对 \(f(X_t)\) 应用 Itô 公式：\(df(X_t) = \mathcal{L}f(X_t) dt + f'(X_t) \sigma(X_t) dW_t\)。设 \(\tau_n = \tau \wedge n\) 有界停时。积分并取期望，\(\mathbb{E}[f(X_{\tau_n})] = f(x) + \mathbb{E}[\int_0^{\tau_n} \mathcal{L}f(X_s) ds]\)（鞅部分期望为 0）。令 \(n \to \infty\)，由 \(\mathbb{E}[\tau] < \infty\) 和控制收敛定理即得。\(\blacksquare\)

\subsection{136.4 扩散过程的构造与 Kolmogorov 方程}\label{ux6269ux6563ux8fc7ux7a0bux7684ux6784ux9020ux4e0e-kolmogorov-ux65b9ux7a0b}

\textbf{定理 136.5}（Kolmogorov 前向方程 / Fokker-Planck 方程）：扩散过程 \(X_t\) 的概率密度 \(p(t,x)\)（若存在）满足

\[\frac{\partial p}{\partial t} = -\frac{\partial}{\partial x}(b p) + \frac{1}{2} \frac{\partial^2}{\partial x^2}(\sigma^2 p) = \mathcal{L}^* p\]

其中 \(\mathcal{L}^*\) 是生成元 \(\mathcal{L}\) 的伴随。

\textbf{证明}：对任意测试函数 \(\varphi\)，计算 \(\frac{d}{dt} \mathbb{E}[\varphi(X_t)]\)。由 Itô 公式，\(d\varphi(X_t) = \mathcal{L}\varphi(X_t) dt + (\cdots) dW_t\)。取期望消去鞅部分：\(\frac{d}{dt} \int \varphi(x) p(t,x) dx = \int \mathcal{L}\varphi(x) p(t,x) dx = \int \varphi(x) \mathcal{L}^* p(t,x) dx\)（由分部积分和伴随算子定义）。由于 \(\varphi\) 任意，\(p\) 满足 \(\partial_t p = \mathcal{L}^* p\)。\(\mathcal{L}^* p = -\partial_x(b p) + \frac{1}{2}\partial_{xx}(\sigma^2 p)\) 由 \(\mathcal{L} = b\partial_x + \frac{\sigma^2}{2}\partial_{xx}\) 的伴随公式得出。\(\blacksquare\)

\hrulefill

\hrulefill

\hrulefill

\hrulefill

\hrulefill

\hrulefill

\section{第353章：Malliavin 分析}\label{ux7b2c353ux7ae0malliavin-ux5206ux6790}

Malliavin 分析（Malliavin 计算）是无穷维空间上的微分分析，由 Paul Malliavin 在 1970 年代建立，其初始动机是证明 Hörmander 亚椭圆性定理的概率方法。它在随机分析中的地位相当于 Sobolev 理论在 PDE 中的地位。

\subsection{137.1 Wiener 混沌分解}\label{wiener-ux6df7ux6c8cux5206ux89e3}

\textbf{定理 137.1}（Wiener-Itô 混沌分解）：任意 \(L^2(\Omega, \mathcal{F}, P)\) 的泛函 \(F\)（其中 \(\mathcal{F} = \sigma(W_t: t \in  [0,T])\)）可唯一分解为正交和：

\[F = \sum_{n=0}^{\infty} I_n(f_n)\]

其中 \(I_n\) 是 \(n\) 重 Itô-Wiener 积分（关于 Brown 运动），\(f_n \in L^2([0,T]^n)\) 是对称函数。

Wiener 混沌分解是 Malliavin 分析的基石，其思想可追溯到 Cameron-Martin (1947) 和 Itô (1951)。该分解将 \(L^2(\Omega)\) 中的泛函按''复杂度''（即所涉及的 Wiener 积分的重数）分层，\(n\) 阶 Wiener 混沌 \(\mathcal{H}_n\) 中的元素恰为 \(n\) 重 Wiener-Itô 积分的全体。\(\mathcal{H}_n\) 与对称 \(L^2([0,T]^n)\) 函数空间的等距同构是多重 Itô 等距的推论：\(\mathbb{E}[I_n(f) I_m(g)] = \delta_{nm} n! \langle f, g \rangle_{L^2([0,T]^n)}\)。一阶混沌 \(\mathcal{H}_1\) 由 Wiener 积分的闭包构成，即 Gauss 随机变量的线性空间；二阶混沌 \(\mathcal{H}_2\) 包含二次 Wiener 泛函（如 \(\int_0^T W_t^2 dt\) 的 \(L^2\) 正交投影）。从表示论的角度看，Wiener 混沌分解是 Fock 空间的对称张量代数分解在概率空间上的实现------\(L^2(\Omega) \cong \bigoplus_{n=0}^\infty \operatorname{Sym}(L^2([0,T])^{\otimes n})\)，其中 \(\operatorname{Sym}\) 表示对称化。这种分解使得 Malliavin 导数 \(D\) 在混沌分解下表现为''降阶''算子：\(D I_n(f_n) = n I_{n-1}(f_n(\cdot, t))\)（将 \(n\) 阶混沌元素映射到 \(n-1\) 阶），而 Skorohod 积分 \(\delta\) 则表现为''升阶''算子。

\textbf{证明}：设 \(\mathcal{H}_n\) 为 \(n\) 阶 Wiener 混沌（由 \(n\) 次 Hermite 多项式生成的闭子空间）。Hermite 多项式 \(\{H_n(\int_0^T h(s) dW_s)\}\)（\(\|h\|_{L^2}=1\)）是正交的。由 Itô 积分的迭代性质，\(\mathcal{H}_n\) 与 \(n\) 重 Wiener-Itô 积分 \(\{I_n(f) : f \in L^2([0,T]^n), f \text{ 对称}\}\) 等距同构（Itô 等距推广为多重版本：\(\|I_n(f)\|^2 = n! \|f\|_{L^2([0,T]^n)}^2\)）。\(L^2(\Omega)\) 分解为 \(\bigoplus_{n=0}^\infty \mathcal{H}_n\) 的正交直和。\(\blacksquare\)

\subsection{137.2 Malliavin 导数}\label{malliavin-ux5bfcux6570}

\textbf{定义 137.2}（Malliavin 导数）：对光滑随机变量 \(F = f(W_{t_1}, \ldots, W_{t_n})\)（\(f \in C_p^\infty(\mathbb{R}^n)\)------光滑且多项式增长），其 Malliavin 导数为

\[D_t F = \sum_{i=1}^n \frac{\partial f}{\partial x_i}(W_{t_1}, \ldots, W_{t_n}) 1_{[0, t_i)}(t)\]

它是 \(L^2([0,T] \times \Omega)\) 中取值的过程。

\textbf{定义 137.3}（Sobolev 空间 \(\mathbb{D}^{1,2}\)）：通过闭包

\[\|F\|_{\mathbb{D}^{1,2}}^2 = E[|F|^2] + E\left[\int_0^T |D_t F|^2 dt\right]\]

定义的完备化空间。

\textbf{命题 137.2}（基本性质）： - \textbf{链式法则}：\(D_t(\varphi(F)) = \varphi'(F) D_t F\)（\(\varphi \in C^1\)） - \textbf{分部积分}（Skorohod 积分的伴随性）：\(E[F \delta(u)] = E[\int_0^T D_t F \cdot u_t dt]\) - \textbf{Clark-Ocone 公式}：\(F = E[F] + \int_0^T E[D_t F \mid \mathcal{F}_t] dW_t\)

\subsection{137.3 Skorohod 积分}\label{skorohod-ux79efux5206}

\textbf{定义 137.4}（Skorohod 积分）：\(u \in L^2([0,T] \times \Omega)\) 的 Skorohod 积分 \(\delta(u)\) 是 Malliavin 导数 \(D\) 的伴随算子：

\[E[\langle D F, u \rangle_{L^2([0,T])}] = E[F \delta(u)]\]

Skorohod 积分是 Itô 积分的推广：当 \(u\) 是适应过程时，\(\delta(u) = \int_0^T u_t dW_t\)（Itô 积分）；当 \(u\) 不是适应过程（例如可预期/anticipating 积分）时，Skorohod 积分给出有意义的推广。

Skorohod 积分的定义域 \(\operatorname{Dom}(\delta)\) 是 Malliavin Sobolev 空间 \(\mathbb{D}^{1,2}(L^2([0,T]))\) 中满足可积性条件的过程全体。它不是 \(L^2([0,T] \times \Omega)\) 的全部，因为并非所有平方可积过程都允许 Skorohod 积分------可积性要求过程的 Malliavin 导数满足一定的正则性。在 Wiener 混沌分解下，若 \(u_t = \sum_{n=0}^\infty I_n(f_n(\cdot, t))\)，则 \(\delta(u) = \sum_{n=0}^\infty I_{n+1}(\tilde{f}_n)\)，其中 \(\tilde{f}_n\) 是 \(f_n\) 的对称化。这揭示了 Skorohod 积分是''升阶''算子：它将 \(n\) 阶混沌提升为 \(n+1\) 阶。Skorohod 积分与 Malliavin 导数之间的伴随关系 \(E[F \delta(u)] = E[\int_0^T D_t F \cdot u_t dt]\) 是 Malliavin 分析中分部积分公式的核心，它推广了 Itô 积分中的 \(E[\int_0^T \varphi_t dW_t] = 0\)（对适应过程 \(\varphi\)）。Skorohod 积分在非适应随机微分方程（如倒向随机微分方程 BSDE）、金融数学中的内幕交易模型和随机控制理论中有重要应用。

\subsection{137.4 Malliavin 亚椭圆性定理}\label{malliavin-ux4e9aux692dux5706ux6027ux5b9aux7406}

\textbf{定理 137.3}（Hörmander 条件）：对 SDE \(dX_t = V_0(X_t) dt + \sum_{i=1}^m V_i(X_t) \circ dW^i_t\)（Stratonovich 形式），若向量场族 \(\{V_1, \ldots, V_m\}\) 及其所有 Lie 括号在每点 span 整个切空间（Hörmander 条件），则 \(X_t\) 有光滑密度。

\textbf{证明}：定义 Malliavin 协方差矩阵 \(\gamma_{ij} = \langle D X_t^{(i)}, D X_t^{(j)} \rangle_{L^2([0,T])}\)。由 SDE 的变分方程，\(D_s X_t\) 满足线性 SDE：\(d(D_s X_t) = \sum_{\alpha=0}^m \partial V_\alpha(X_t) (D_s X_t) \circ dW_t^\alpha\)（\(t \geq s\)），初值 \(D_s X_s = \sigma(X_s)\)。若 Hörmander 条件成立，则 \(\gamma\) 几乎必然可逆，且对所有 \(p \geq 1\)，\(\mathbb{E}[(\det \gamma)^{-p}] < \infty\)。由 Malliavin 分部积分公式，对任意多重指标 \(\alpha\)，存在 \(H_\alpha \in \mathbb{D}^\infty\) 使得 \(\mathbb{E}[\partial^\alpha \varphi(X_t)] = \mathbb{E}[\varphi(X_t) H_\alpha]\) 对所有 \(\varphi \in C_b^\infty\) 成立。密度 \(p_t(x)\) 存在且 \(p_t \in C^\infty\)，其各阶导数为 \(p_t^{(\alpha)}(x) = \mathbb{E}[\mathbf{1}_{X_t > x} H_\alpha]\)。\(\blacksquare\)

\hrulefill

\section{第354章：Malliavin分析深化}\label{ux7b2c354ux7ae0malliavinux5206ux6790ux6df1ux5316}

Malliavin分析的核心力量在于其将随机变分法（无穷维微分学）与概率分布的正则性联系起来。上一章建立了基本框架------Wiener混沌分解、Malliavin导数和Skorohod积分。本章将深化这一理论：完整证明Clark-Ocone公式（利用Skorohod积分的伴随性），建立密度光滑性的Malliavin准则（Malliavin协方差矩阵的非退化性推出分布绝对连续且有光滑密度，给出完整证明），并探讨Hörmander条件和次椭圆性的概率证明------这是Malliavin创立该理论的原始动机，也是随机分析对偏微分方程理论的重大贡献。

\subsection{Clark-Ocone公式的完整证明}\label{clark-oconeux516cux5f0fux7684ux5b8cux6574ux8bc1ux660e}

\textbf{定理291.1（Clark-Ocone公式）} 设 \(F \in \mathbb{D}^{1,2}\) 是Malliavin可微随机变量。则

\[F = \mathbb{E}[F] + \int_0^T \mathbb{E}[D_t F \mid \mathcal{F}_t] \, dW_t\]

其中积分是Itô积分，\(\mathcal{F}_t = \sigma(W_s : s \leq t)\) 是Brown运动生成的自然滤过。

\textbf{证明} 利用Wiener混沌分解和Skorohod积分的性质。

\textbf{第1步（混沌分解）}：由Wiener-Itô混沌分解（定理137.1），将 \(F\) 展开为 \(F = \sum_{n=0}^\infty I_n(f_n)\)，其中 \(f_n \in L^2([0,T]^n)\) 为对称函数，\(I_0(f_0) = \mathbb{E}[F]\)。

\textbf{第2步（Malliavin导数在混沌上的作用）}：对 \(n\) 阶混沌元素 \(I_n(f_n)\)，Malliavin导数为

\[D_t I_n(f_n) = n I_{n-1}(f_n(\cdot, t))\]

其中 \(I_{n-1}(f_n(\cdot, t))\) 是固定第 \(n\) 个变量为 \(t\) 后关于前 \(n-1\) 个变量的 \((n-1)\) 重Wiener-Itô积分。因此

\[D_t F = \sum_{n=1}^\infty n I_{n-1}(f_n(\cdot, t))\]

\textbf{第3步（条件期望的计算）}：取条件期望 \(\mathbb{E}[\cdot \mid \mathcal{F}_t]\)。在Wiener混沌分解中，\(\mathcal{F}_t\)-条件期望等价于将每个Wiener-Itô积分的被积函数截断到 \([0,t]^n\)：

\[\mathbb{E}[I_n(f_n) \mid \mathcal{F}_t] = I_n(f_n \mathbf{1}_{[0,t]^n})\]

因此 \(\mathbb{E}[D_t F \mid \mathcal{F}_t] = \sum_{n=1}^\infty n I_{n-1}(f_n(\cdot, t) \mathbf{1}_{[0,t]^{n-1}})\)。

\textbf{第4步（Itô积分的混沌表示）}：被积过程 \(\varphi_t = \mathbb{E}[D_t F \mid \mathcal{F}_t]\) 是适应过程（由条件期望的性质和Brown滤过的适应性）。在混沌分解下，\(\varphi_t = \sum_{n=0}^\infty I_n(g_n(\cdot, t))\)，其中 \(g_n\) 的支撑集包含于 \([0,t]^{n+1}\)。Itô积分 \(\int_0^T \varphi_t dW_t\) 在混沌分解下的表示为：被积函数的对称化后作为 \((n+1)\) 重积分的核。

\textbf{第5步（验证等式）}：直接计算验证 \(\int_0^T \mathbb{E}[D_t F \mid \mathcal{F}_t] dW_t = F - \mathbb{E}[F]\)。对每个 \(n \geq 1\)，\(n\) 阶混沌分量在等式两边匹配（利用Itô积分的等距性质和Malliavin导数与Skorohod积分之间的伴随关系）。由于混沌分解的唯一性，等式在整个 \(L^2(\Omega)\) 中成立。

\textbf{第6步（Skorohod积分观点）}：等价地，Clark-Ocone公式可表述为：令 \(u_t = \mathbb{E}[D_t F \mid \mathcal{F}_t]\)，则 \(u\) 是适应过程，其Skorohod积分 \(\delta(u)\) 等于Itô积分 \(\int_0^T u_t dW_t\)。由Skorohod积分的定义和Malliavin导数与Skorohod积分的伴随关系，\(F - \mathbb{E}[F] = \delta(D F)\)（其中 \(D F\) 作为 \(L^2([0,T] \times \Omega)\) 中过程），但 \(D F\) 一般不适应。Clark-Ocone公式将其投影到适应过程上，得到Itô积分的表示。\(\blacksquare\)

\subsection{密度光滑性的Malliavin准则}\label{ux5bc6ux5ea6ux5149ux6ed1ux6027ux7684malliavinux51c6ux5219}

\textbf{定理291.2（Malliavin协方差矩阵的非退化性判别准则）} 设 \(F = (F_1, \ldots, F_d)\) 是 \(\mathbb{R}^d\)-值随机向量，\(F_i \in \mathbb{D}^\infty\)（所有阶Malliavin导数存在且平方可积）。定义\textbf{Malliavin协方差矩阵}

\[\gamma_{ij} = \langle D F_i, D F_j \rangle_{L^2([0,T])} = \int_0^T (D_t F_i)(D_t F_j) \, dt\]

若 \(\gamma\) 几乎必然可逆，且对所有 \(p \geq 1\)，\(\mathbb{E}[(\det \gamma)^{-p}] < \infty\)，则 \(F\) 的分布具有关于Lebesgue测度的 \(C^\infty\) 光滑密度。

\textbf{证明} 核心工具是Malliavin分部积分公式。

\textbf{第1步（分部积分公式的建立）}：对任意 \(\varphi \in C_b^\infty(\mathbb{R}^d)\) 和多重指标 \(\alpha = (\alpha_1, \ldots, \alpha_d)\)，需构造随机变量 \(H_\alpha \in \mathbb{D}^\infty\) 使得

\[\mathbb{E}[\partial^\alpha \varphi(F)] = \mathbb{E}[\varphi(F) H_\alpha]\]

其中 \(\partial^\alpha = \partial_1^{\alpha_1} \cdots \partial_d^{\alpha_d}\)。

\textbf{第2步（一阶导数情形）}：由链式法则，\(D_t(\varphi(F)) = \sum_{j=1}^d (\partial_j \varphi)(F) D_t F_j\)。取与 \(D F_i\) 的 \(L^2([0,T])\) 内积：

\[\langle D(\varphi(F)), D F_i \rangle = \sum_{j=1}^d (\partial_j \varphi)(F) \gamma_{ji}\]

由于 \(\gamma\) 可逆，设 \(\gamma^{-1} = (\gamma^{ij})\)，则

\[(\partial_i \varphi)(F) = \sum_{j=1}^d \langle D(\varphi(F)), D F_j \rangle \gamma^{ji}\]

\textbf{第3步（伴随关系的应用）}：由Malliavin导数与Skorohod积分的伴随关系，

\[\mathbb{E}[\langle D(\varphi(F)), D F_j \rangle] = \mathbb{E}[\varphi(F) \delta(D F_j)]\]

但需注意 \(D(\varphi(F))\) 已经过链式法则处理。更精确地，对任意 \(u \in \operatorname{Dom}(\delta)\)，

\[\mathbb{E}[(\partial_i \varphi)(F)] = \mathbb{E}\left[\sum_j \langle D(\varphi(F)), D F_j \gamma^{ji} \rangle\right] = \mathbb{E}\left[\varphi(F) \sum_j \delta(D F_j \gamma^{ji})\right]\]

因此 \(H_{e_i} = \sum_{j=1}^d \delta(D F_j \gamma^{ji})\)。

\textbf{第4步（高阶导数的归纳）}：对多重指标 \(\alpha\)，通过归纳构造 \(H_\alpha\)。若 \(H_\alpha\) 已构造，则

\[\mathbb{E}[\partial_i \partial^\alpha \varphi(F)] = \mathbb{E}[\partial^\alpha \varphi(F) H_{e_i}] = \mathbb{E}[\varphi(F) H_{\alpha + e_i}]\]

其中 \(H_{\alpha + e_i}\) 通过对 \(H_{e_i} \cdot H_\alpha\) 的Malliavin导数迭代和Skorohod积分作用构造。关键在于验证 \((\det \gamma)^{-1}\) 的所有矩有限性能保证每次迭代产生的随机变量仍属于 \(\mathbb{D}^\infty\)。

\textbf{第5步（密度的光滑性）}：若对所有 \(\alpha\)，\(H_\alpha \in L^1(\Omega)\)，则对任意 \(\varphi \in C_b^\infty\)，

\[|\mathbb{E}[\partial^\alpha \varphi(F)]| \leq \|\varphi\|_{\infty} \mathbb{E}[|H_\alpha|]\]

这表明分布导数 \(\partial^\alpha \mu\)（\(\mu\) 为 \(F\) 的分布）是有限Borel测度。由Sobolev嵌入定理，\(\mu\) 具有 \(C^\infty\) 密度 \(p(x)\)，且

\[p^{(\alpha)}(x) = \mathbb{E}[\mathbf{1}_{\{F > x\}} H_\alpha]\]

其中 \(\{F > x\}\) 表示 \(\{F_1 > x_1, \ldots, F_d > x_d\}\)。\(\blacksquare\)

\textbf{推论291.3（一维特例）} 若 \(F \in \mathbb{D}^{1,2}\) 满足 \(\|D F\|_{L^2([0,T])} > 0\) a.s.，则 \(F\) 的分布绝对连续。若进一步 \(\mathbb{E}[\|D F\|_{L^2}^{-p}] < \infty\) 对所有 \(p\) 成立，则密度是 \(C^\infty\) 的。

\subsection{Hörmander条件和次椭圆性的概率证明}\label{huxf6rmanderux6761ux4ef6ux548cux6b21ux692dux5706ux6027ux7684ux6982ux7387ux8bc1ux660e}

\textbf{定理291.4（Hörmander亚椭圆性定理的概率证明框架）} 对SDE

\[dX_t = V_0(X_t) dt + \sum_{i=1}^m V_i(X_t) \circ dW_t^i\]

若Hörmander条件成立（\(\{V_i, [V_i, V_j], [V_i, [V_j, V_k]], \ldots\}\) 在每点张成 \(\mathbb{R}^d\)），则对任意 \(t > 0\) 和初始值 \(x\)，\(X_t^x\) 的分布具有 \(C^\infty\) 光滑密度。进一步，转移密度 \(p_t(x, y)\) 是 \((t, x, y)\) 的 \(C^\infty\) 函数（在 \(t > 0\) 时）。

\textbf{证明}（Malliavin方法的概率论证）分为三步。

\textbf{第1步（协方差矩阵的矩估计）}：由Malliavin准则（定理291.2），只需证明Malliavin协方差矩阵 \(\gamma_t\) 的逆矩阵的任意阶矩均有限。关键在于Hörmander条件等价于 \(\gamma_t\) 的几乎必然可逆性。证明利用Stroock-Varadhan支撑定理：由Hörmander条件，扩散过程 \(X_t\) 的转移概率的支撑集是整个 \(\mathbb{R}^d\)（即过程可达任意开集）。若存在非零向量 \(v\) 使 \(v^T \gamma_t v = 0\) a.s.，则 \(v\) 将始终与所有 \(D X_t^{(j)}\) 正交，这与支撑集为全空间矛盾。因此 \(\gamma_t > 0\) a.s.（正定矩阵）。

\textbf{第2步（逆矩的指数衰减估计）}：需证明 \(\mathbb{E}[(\det \gamma_t)^{-p}] < \infty\) 对所有 \(p \geq 1\)。这通过Norris引理（1986）完成：对任意单位向量 \(v\) 和 \(\varepsilon > 0\)，

\[\mathbb{P}(v^T \gamma_t v < \varepsilon) \leq C \varepsilon^q\]

对某个 \(q > 0\) 成立（\(C\) 和 \(q\) 取决于Hörmander括号的阶数和向量场的范数）。此估计的证明基于：若 \(v^T \gamma_t v\) 非常小，则 \(v^T D X_t\) 的 \(L^2\) 范数非常小，这意味着在随机区间上 \(v^T V_i(X_s) \approx 0\) 对所有 \(i\)。通过反复对SDE的Stratonovich形式应用Itô公式于Lie括号，利用Hörmander条件可推出矛盾（除非 \(v = 0\)）。

\textbf{第3步（光滑密度的存在）}：由第2步的估计，对任意 \(p\) 存在 \(C_p\) 使得 \(\mathbb{E}[(\det \gamma_t)^{-p}] < C_p < \infty\)。由定理291.2，\(X_t\) 的分布具有 \(C^\infty\) 光滑密度。转移密度关于 \((t,x,y)\) 的光滑性来自对初始值的变分方程和Kolmogorov方程的标准论证。\(\blacksquare\)

上述概率证明绕过了Hörmander原证明中的伪微分算子演算，以随机分析的工具（Itô公式、Norris引理、Malliavin分部积分）给出了亚椭圆性一个全新的、更直观的证明。

\hrulefill

\section{第355章：SPDE的解理论与Feynman-Kac}\label{ux7b2c355ux7ae0spdeux7684ux89e3ux7406ux8bbaux4e0efeynman-kac}

随机偏微分方程（SPDE）将随机分析与偏微分方程理论融合，描述受随机噪声驱动的无穷维动力系统。本章建立SPDE的严格数学框架：从Wiener混沌分解与多重Wiener-Itô积分出发，给出Cameron-Martin定理的泛函分析证明；进而研究抛物型SPDE的正则性理论（\(L^2\)理论与极大正则性）；最后讨论随机波动方程与随机热方程------这两类方程分别代表了双曲型和抛物型SPDE的典型行为。

\subsection{Wiener混沌分解与多重Wiener-Itô积分}\label{wienerux6df7ux6c8cux5206ux89e3ux4e0eux591aux91cdwiener-ituxf4ux79efux5206}

\textbf{定义292.1（多重Wiener-Itô积分）} 设 \(f \in L^2([0,T]^n)\) 为对称函数。\(n\) 重Wiener-Itô积分 \(I_n(f)\) 定义为

\[I_n(f) = n! \int_0^T \int_0^{t_n} \cdots \int_0^{t_2} f(t_1, \ldots, t_n) \, dW_{t_1} \cdots dW_{t_n}\]

其中积分是迭代Itô积分。\(I_n(f)\) 满足等距性质：\(\mathbb{E}[I_n(f)^2] = n! \|f\|_{L^2([0,T]^n)}^2\)。不同重数的积分正交：\(\mathbb{E}[I_n(f) I_m(g)] = 0\)（\(n \neq m\)）。

\textbf{定理292.2（Cameron-Martin定理的泛函分析证明）} 设 \(\gamma\) 是Cameron-Martin空间 \(H = \{h \in C([0,T]) : h(t) = \int_0^t \dot{h}(s) ds, \dot{h} \in L^2([0,T])\}\) 中的元素。则由 \(h\) 平移的Wiener测度 \(\mathbb{P} \circ (W + h)^{-1}\) 关于原始Wiener测度 \(\mathbb{P}\) 绝对连续，Radon-Nikodym密度为

\[\frac{d\mathbb{P} \circ (W + h)^{-1}}{d\mathbb{P}} = \exp\left(\int_0^T \dot{h}(t) dW_t - \frac{1}{2} \int_0^T \dot{h}(t)^2 dt\right)\]

\textbf{证明} 利用Wiener混沌分解和指数鞅。

\textbf{第1步（指数鞅的构造）}：对 \(h \in H\)，过程

\[Z_t = \exp\left(\int_0^t \dot{h}(s) dW_s - \frac{1}{2} \int_0^t \dot{h}(s)^2 ds\right)\]

是Itô积分定义的正鞅（Novikov条件自动满足，因 \(\dot{h} \in L^2([0,T])\) 是确定性的）。特别地，\(\mathbb{E}[Z_T] = 1\) 且 \(Z_t = \sum_{n=0}^\infty I_n(\dot{h}^{\otimes n}/n!)\)（由Itô积分的指数公式的混沌展开）。

\textbf{第2步（测度变换的验证）}：定义新测度 \(d\mathbb{Q} = Z_T d\mathbb{P}\)。需证在 \(\mathbb{Q}\) 下，过程 \(W_t\) 的分布等同于 \(\mathbb{P}\) 下 \(W_t + h(t)\) 的分布。对任意可测泛函 \(F\)，

\[\mathbb{E}_{\mathbb{Q}}[F(W)] = \mathbb{E}_{\mathbb{P}}[Z_T F(W)] = \mathbb{E}_{\mathbb{P}}[F(W + h)]\]

第二个等号由Girsanov定理的有限维近似和Wiener积分的定义验证：对柱集函数 \(F(W) = f(W_{t_1}, \ldots, W_{t_n})\)，\(W_{t_i}\) 在 \(\mathbb{Q}\) 下的分布为 \(\mathcal{N}(h(t_i), t_i)\)，即 \(W + h\) 在 \(\mathbb{P}\) 下的分布。

\textbf{第3步（泛函分析解释）}：Cameron-Martin空间 \(H\) 是Wiener空间 \(C_0([0,T])\) 上Gauss测度的容许平移方向构成的Hilbert空间。更一般地，若 \(h \notin H\)（即 \(\dot{h}\) 不是平方可积），则 \(\mathbb{P} \circ (W + h)^{-1}\) 与 \(\mathbb{P}\) 相互奇异。这反映了无穷维Gauss测度的准平移不变性仅沿Cameron-Martin方向成立的深刻事实。\(\blacksquare\)

\subsection{抛物型SPDE的正则性理论}\label{ux629bux7269ux578bspdeux7684ux6b63ux5219ux6027ux7406ux8bba}

\textbf{定义292.3（随机热方程）} 考虑带乘性噪声的随机热方程（Itô形式）：

\[du(t,x) = \Delta u(t,x) dt + \sigma(u(t,x)) dW(t,x), \quad (t,x) \in (0,T] \times D\]

其中 \(D \subset \mathbb{R}^d\) 是有界光滑区域，\(W(t,x)\) 是时空白噪声（或更一般地，\(Q\)-Wiener过程），\(\Delta\) 为Laplace算子（配备Dirichlet或Neumann边界条件）。

\textbf{定理292.4（\(L^2\)理论------变分解的存在唯一性）} 设 \(\sigma\) 满足Lipschitz条件：\(\|\sigma(u) - \sigma(v)\|_{L^2(D)} \leq L \|u - v\|_{L^2(D)}\) 且 \(\|\sigma(u)\|_{L^2(D)} \leq C(1 + \|u\|_{L^2(D)})\)。则随机热方程存在唯一的温和解（mild solution）\(u \in L^2(\Omega; C([0,T]; L^2(D))) \cap L^2(\Omega; L^2([0,T]; H_0^1(D)))\)，且满足能量估计

\[\mathbb{E}\left[\sup_{t \in [0,T]} \|u(t)\|_{L^2}^2 + \int_0^T \|\nabla u(t)\|_{L^2}^2 dt\right] \leq C_T (1 + \|u_0\|_{L^2}^2)\]

\textbf{证明}（Galerkin逼近与单调算子理论）选取 \(H_0^1(D)\) 的一组正交基 \(\{e_k\}\)（如Laplace算子的特征函数），构造有限维逼近 \(u^N(t) = \sum_{k=1}^N c_k^N(t) e_k\) 满足有限维SDE。由Itô公式于 \(\|u^N(t)\|_{L^2}^2\) 得一致能量估计。利用Burkholder-Davis-Gundy不等式处理随机积分的二阶矩。紧性论证（通过Prokhorov定理和Skorohod表示定理）给出收敛子列，其极限满足原SPDE。唯一性来自Grönwall引理和Lipschitz条件。\(\blacksquare\)

\textbf{定理292.5（极大正则性理论）} 若噪声正则性足够高（如迹类噪声：\(\operatorname{Tr}(Q) < \infty\)），则解 \(u\) 具有更高的空间正则性。具体地，若初始值 \(u_0 \in H^s(D)\) 且噪声在空间方向具有 \(H^{s+\alpha}\) 的正则性，则解 \(u \in L^2(\Omega; C([0,T]; H^{s+\beta}(D)))\) 对某 \(\beta > 0\)。此结果依赖于半群 \(e^{t\Delta}\) 的极大 \(L^p\) 正则性和随机卷积的随机极大正则性理论（van Neerven、Veraar和Weis，2000年代建立的框架）。

\subsection{随机波动方程与随机热方程}\label{ux968fux673aux6ce2ux52a8ux65b9ux7a0bux4e0eux968fux673aux70edux65b9ux7a0b}

\textbf{定义292.6（随机波动方程）} 带可加噪声的随机波动方程：

\[\partial_{tt} u(t,x) = \Delta u(t,x) + \dot{W}(t,x), \quad (t,x) \in (0,T] \times D\]

其中 \(\dot{W}\) 是时空白噪声。此方程可转化为一阶系统后应用半群方法。

\textbf{定理292.7（随机波动方程的解的存在性与正则性）} 在空间维数 \(d = 1\) 时，随机波动方程存在唯一的函数值解（function-valued solution），且解几乎必然属于 \(C([0,T]; H^{1/2-\varepsilon}(D)) \cap C^1([0,T]; H^{-1/2-\varepsilon}(D))\)。在 \(d = 2\) 时，解仅为分布值（随机过程取值于负Sobolev空间 \(H^{-s}(D)\)，\(s > 0\)）；\(d \geq 3\) 时，方程的解需要重新正规化（renormalization），这触及了随机量子场论和Hairer正则性结构理论的领域。

\textbf{证明}（\(d=1\)的情形）利用波动算子的基本解和随机卷积表示：

\[u(t) = \cos(t\sqrt{-\Delta}) u_0 + \frac{\sin(t\sqrt{-\Delta})}{\sqrt{-\Delta}} v_0 + \int_0^t \frac{\sin((t-s)\sqrt{-\Delta})}{\sqrt{-\Delta}} dW(s)\]

其中 \(\cos(t\sqrt{-\Delta})\) 和 \(\sin(t\sqrt{-\Delta})/\sqrt{-\Delta}\) 是波动方程的半群算子。随机卷积项的正则性通过Itô等距和特征函数展开估计：在 \(d=1\) 时，\(\sqrt{-\Delta}^{-1}\) 将白噪声提升两阶空间导数，补偿了白噪声的 \(-1/2-\varepsilon\) 正则性，得到总空间正则性 \(1/2-\varepsilon\)。\(\blacksquare\)

随机热方程与随机波动方程的正则性差异反映了抛物型和双曲型PDE的本质区别：热半群 \(e^{t\Delta}\) 对空间变量提供两阶光滑化（对时间提供一阶），而波动半群不提供光滑化------噪声的正则性直接传导至解的正则性（仅由弥散关系改善）。这解释了为何随机波动方程对空间维数更为敏感。

\hrulefill

\hrulefill

\hrulefill

\hrulefill

\hrulefill

\hrulefill

\subsection{287.A Itô积分补充}\label{a-ituxf4ux79efux5206ux8865ux5145}

\begin{quote}
随机分析是随机过程与数学分析的深度交叉领域，以 Itô 积分理论为核心，研究受随机扰动影响的系统。本卷建立在概率论、概率论 II 和测度论的基础上，系统建立 Itô 积分、随机微分方程（SDE）、扩散过程、Malliavin 分析和随机偏微分方程（SPDE）的理论。随机分析在现代金融数学（Black-Scholes 模型）、物理学（统计力学、量子场论）、生物学（种群动力学）和工程（信号处理）中有广泛应用。Itô 积分的核心突破在于放弃了 Riemann-Stieltjes 积分中积分点任意选取的自由，而强制在左端点取值------这一选择使得积分成为鞅，从而可以应用概率论的全部工具。Itô 公式（随机链式法则）则揭示了随机微积分与经典微积分的本质区别：二阶项不能忽略，因为 Brown 运动的二次变差非零。本章系统介绍 Brown 运动与连续鞅、Itô 积分的构造、Itô 公式、Girsanov 定理与测度变换、随机微分方程（SDE）的存在唯一性理论、Feynman-Kac 公式以及扩散过程的基本理论。
\end{quote}

\hrulefill

\section{第356章：Itô 积分与随机积分}\label{ux7b2c356ux7ae0ituxf4-ux79efux5206ux4e0eux968fux673aux79efux5206}

Itô 积分是随机分析的核心构造，由 Itô 在 1944 年建立。与 Riemann-Stieltjes 积分不同，被积函数在左端点取值（而非任意点），使得积分成为鞅，这在理论和应用中都至关重要。Itô 积分的构造是测度论概率论在现代数学中的典范应用：它结合了 Hilbert 空间理论（\(L^2\) 空间的等距）、鞅论（Doob 的鞅不等式）和随机过程的正则性理论（Kolmogorov 连续性准则）。Itô 积分的核心应用领域包括：随机微分方程（SDE）------描述受噪声扰动的动力系统；金融数学------Black-Scholes 期权定价公式；滤波理论------Kalman-Bucy 滤波；以及随机控制------动态规划与 Hamilton-Jacobi-Bellman 方程。

\subsection{135.1 Brown 运动与连续鞅}\label{brown-ux8fd0ux52a8ux4e0eux8fdeux7eedux9785}

\textbf{定义 135.1}（Brown 运动）：随机过程 \((W_t)_{t \geq 0}\) 是\textbf{标准 Brown 运动}（Wiener 过程），如果： （1）\(W_0 = 0\) a.s. （2）增量独立性：对 \(0 \leq s < t\)，\(W_t - W_s\) 与 \(\sigma(W_u: u \leq s)\) 独立 （3）正态增量：\(W_t - W_s \sim N(0, t-s)\) （4）连续性：\(t \mapsto W_t\) a.s. 连续

\textbf{定理 135.1}（Brown 运动的性质）： - 样本路径 a.s. 连续但处处不可微（Paley-Wiener-Zygmund 定理） - 二次变差：\([W, W]_t = \lim_{\|\Pi\| \to 0} \sum_i (W_{t_{i+1}} - W_{t_i})^2 = t\)（\(L^2\) 意义下） - 尺度变换：\((c^{-1/2} W_{ct})\) 是 Brown 运动 - 时间反演：\((t W_{1/t})\)（补充 \(0\) at \(t=0\)）是 Brown 运动 - 反射原理：\(\mathbb{P}(\sup_{0 \leq s \leq t} W_s \geq a) = 2\mathbb{P}(W_t \geq a)\)（\(a > 0\)） - 强 Markov 性：Brown 运动在停时 \(\tau\) 处重新开始：\((W_{\tau + t} - W_\tau)_{t \geq 0}\) 独立于 \(\mathcal{F}_\tau\) 且仍为 Brown 运动

\textbf{证明}（二次变差）：对分划 \(\Pi_n = \{t_0,\ldots,t_n\}\)，增量 \(\Delta W_i \sim \mathcal{N}(0,\Delta t_i)\) 独立。\(\mathbb{E}[(\Delta W_i)^2] = \Delta t_i\)，\(\operatorname{Var}((\Delta W_i)^2) = 2(\Delta t_i)^2\)。\(\operatorname{Var}(\sum (\Delta W_i)^2) = 2\sum (\Delta t_i)^2 \to 0\)（\(\|\Pi_n\| \to 0\)），由 Chebyshev 不等式得 \(L^2\) 收敛于 \(t\)。\(\blacksquare\)

\textbf{定理 135.2}（Kolmogorov 连续性准则）：若对某 \(\alpha, \beta > 0\)，\(E[|X_t - X_s|^\alpha] \leq C |t-s|^{1+\beta}\)，则 \(X_t\) 有连续版本。对 Brown 运动，\(E[(W_t - W_s)^4] = 3(t-s)^2\)（\(\alpha = 4, \beta = 1\)）。

\textbf{证明}：设 \(D \subseteq [0,T]\) 为二进有理点集。由 Chebyshev 不等式，\(\mathbb{P}(|X_t - X_s| \geq |t-s|^\gamma) \leq C |t-s|^{1+\beta-\alpha\gamma}\)。选取 \(\gamma\) 使 \(\beta - \alpha\gamma > 0\)，则级数 \(\sum_n \mathbb{P}(\max_k |X_{k/2^n} - X_{(k-1)/2^n}| > 2^{-n\gamma}) < \infty\)。由 Borel-Cantelli，在 \(D\) 上一致连续，延拓至 \([0,T]\)。\(\blacksquare\)

\textbf{定理 135.3}（鞅性质）：以下关于 Brown 运动自然滤子 \(\mathcal{F}_t = \sigma(W_s: s \leq t)\) 都是鞅： - \(W_t\)（martingale） - \(W_t^2 - t\)（平方变差补偿） - \(\exp(\lambda W_t - \lambda^2 t/2)\)（指数鞅 / 随机指数）

\textbf{证明}：\(W_t\) 为鞅：由增量独立性，\(\mathbb{E}[W_t \mid \mathcal{F}_s] = W_s + \mathbb{E}[W_t - W_s \mid \mathcal{F}_s] = W_s\)。\(W_t^2 - t\)：\(\mathbb{E}[W_t^2 \mid \mathcal{F}_s] = \mathbb{E}[(W_s + (W_t - W_s))^2 \mid \mathcal{F}_s] = W_s^2 + (t-s)\)。指数鞅：\(\mathbb{E}[e^{\lambda W_t} \mid \mathcal{F}_s] = e^{\lambda W_s} \mathbb{E}[e^{\lambda(W_t-W_s)}] = e^{\lambda W_s} \cdot e^{\lambda^2(t-s)/2}\)，故 \(e^{\lambda W_t - \lambda^2 t/2}\) 为鞅。\(\blacksquare\)

\textbf{定义 135.1b}（Brown 运动的构造）：Brown 运动的存在性可通过以下方式证明：(1) Kolmogorov 延拓定理 + Kolmogorov 连续性准则，在 \(C[0,\infty)\) 上构造 Wiener 测度；(2) Levy 的构造：取 \(L^2[0,1]\) 的 Haar 正交基，将 Brown 运动展开为 Schauder 函数的随机级数；(3) 随机游走的 Donsker 不变原理：\(n^{-1/2} S_{\lfloor nt \rfloor} \Rightarrow W_t\)（依分布收敛）。

\textbf{定理 135.3b}（Levy 特征定理）：连续 \(L^2\)-鞅 \((M_t)\) 满足 \(M_0 = 0\) 且 \([M, M]_t = t\)（或等价地，\(M_t^2 - t\) 是鞅）当且仅当 \(M_t\) 是 Brown 运动。Levy 特征定理是识别 Brown 运动的最有力工具------它仅需验证鞅性和二次变差，无需检查增量分布。

\subsection{135.2 Itô 积分的构造}\label{ituxf4-ux79efux5206ux7684ux6784ux9020}

\textbf{定义 135.2}（简单过程的 Itô 积分）：对简单过程 \(H_t = \sum_{i=0}^{n-1} H_{t_i} 1_{[t_i, t_{i+1})}(t)\)（其中 \(H_{t_i}\) 是 \(\mathcal{F}_{t_i}\)-可测且平方可积），定义

\[I(H) = \sum_{i=0}^{n-1} H_{t_i} (W_{t_{i+1}} - W_{t_i})\]

\textbf{定理 135.4}（Itô 等距）：\[E[(I(H))^2] = E\left[\int_0^T H_t^2 dt\right]\]

\emph{证明}：设 \(\Delta W_i = W_{t_{i+1}} - W_{t_i}\)。展开 \(E[(I(H))^2] = E[\sum_{i,j} H_{t_i} H_{t_j} \Delta W_i \Delta W_j]\)。当 \(i \neq j\)，不妨设 \(i < j\)，则 \(E[H_{t_i} H_{t_j} \Delta W_i \Delta W_j] = E[H_{t_i} H_{t_j} \Delta W_i E[\Delta W_j | \mathcal{F}_{t_j}]] = 0\)（因 \(\Delta W_j\) 与 \(\mathcal{F}_{t_j}\) 独立且均值为 0）。对角项：\(E[H_{t_i}^2 (\Delta W_i)^2] = E[H_{t_i}^2 E[(\Delta W_i)^2 | \mathcal{F}_{t_i}]] = E[H_{t_i}^2 (t_{i+1} - t_i)] = E[\int_{t_i}^{t_{i+1}} H_t^2 dt]\)。求和即得。∎

\textbf{定义 135.3}（一般 Itô 积分）：对满足 \(\int_0^T H_t^2 dt < \infty\) a.s. 且 \(H_t\) 是适应过程的 \(H\)，存在简单过程序列 \(H^{(n)} \to H\)（在 \(L^2(\Omega \times [0,T])\) 中），定义

\[\int_0^T H_t dW_t = L^2\text{-}\lim_{n \to \infty} I(H^{(n)})\]

\textbf{性质 135.5}：Itô 积分 \(\int_0^t H_s dW_s\) 是连续的 \(L^2\)-鞅，二次变差 \(\left[\int H dW, \int H dW\right]_t = \int_0^t H_s^2 ds\)。

\textbf{证明}：对简单过程，\(\mathbb{E}[I(H) \mid \mathcal{F}_s] = \sum_{t_i < s} H_{t_i} \Delta W_i\)，且由独立性，\(\mathbb{E}[I(H)\Delta W_j \mid \mathcal{F}_{t_j}] = 0\)，故为鞅。\(L^2\) 有界性来自 Itô 等距。二次变差：\((\sum H_{t_i} \Delta W_i)^2\) 的对角项和近似 \(\sum H_{t_i}^2 \Delta t_i \to \int H_s^2 ds\)。一般过程由稠密性延拓。\(\blacksquare\)

\textbf{定理 135.5b}（Itô 积分的局部鞅性质）：若 \(\int_0^T H_t^2 dt < \infty\) a.s.（而非 \(L^2\) 意义下），则 \(\int_0^t H_s dW_s\) 是局部鞅（而非全局鞅）。局部鞅是鞅概念的自然推广：存在停时列 \(\tau_n \uparrow \infty\) 使得停止过程 \(M_t^{\tau_n} = M_{t \wedge \tau_n}\) 是鞅。局部鞅性质是处理超鞅和 Itô 过程的核心工具。

\textbf{定理 135.5c}（Itô 积分与 Stratonovich 积分的关系）：Stratonovich 积分 \(\int_0^t H_s \circ dW_s\) 在积分点取中点值（而非左端点），满足经典链式法则，但不是鞅。两者关系为： \[\int_0^t H_s \circ dW_s = \int_0^t H_s dW_s + \frac{1}{2} \langle H, W \rangle_t\] 其中 \(\langle H, W \rangle_t\) 为 \(H\) 和 \(W\) 的协变差。Stratonovich 积分在物理应用（如随机力学）中更自然，因为它保持了经典力学的形式。

\subsection{135.3 Itô 公式}\label{ituxf4-ux516cux5f0f}

\textbf{定理 135.6}（Itô 公式，一维）：设 \(f \in C^2\) 和 \(X_t\) 是 Itô 过程 \(dX_t = \mu_t dt + \sigma_t dW_t\)。则

\[df(X_t) = f'(X_t) dX_t + \frac{1}{2} f''(X_t) \sigma_t^2 dt\]

等价地：

\[f(X_t) - f(X_0) = \int_0^t f'(X_s) dX_s + \frac{1}{2} \int_0^t f''(X_s) \sigma_s^2 ds\]

\textbf{证明}：Itô 公式的证明基于 Taylor 展开：对 \(C^2\) 函数 \(f\)，\(df(X_t)=f'(X_t)dX_t+\frac{1}{2}f''(X_t)(dX_t)^2\)。关键点是 \((dW_t)^2=dt\)（二次变差性质），且 \((dt)^2 = o(dt)\)，\(dt \cdot dW_t = o(dt)\)。形式地：\(\Delta f = f'(X_t)\Delta X_t + \frac{1}{2}f''(X_t)(\Delta X_t)^2 + o((\Delta X_t)^2)\)。由于 \((\Delta W_t)^2 \approx \Delta t\)（Brown 运动的二次变差），二阶项不能忽略。取极限 \(\Delta t \to 0\)，\(\sum f''(X_t)(\Delta W_t)^2 \to \int f''(X_t)dt\)（\(L^2\) 收敛），得到 Itô 积分形式。\(\blacksquare\)

\textbf{定理 135.7}（Itô 公式，多维）：设 \(f(t, x_1, \ldots, x_d) \in C^{1,2}\)，\(X^i_t\) 是 Itô 过程，\(dX^i_t = \mu^i_t dt + \sum_{j=1}^m \sigma^{ij}_t dW^j_t\)。则

\[df(t, X_t) = \frac{\partial f}{\partial t} dt + \sum_{i=1}^d \frac{\partial f}{\partial x_i} dX^i_t + \frac{1}{2} \sum_{i,j=1}^d \frac{\partial^2 f}{\partial x_i \partial x_j} d\langle X^i, X^j \rangle_t\]

其中 \(d\langle X^i, X^j \rangle_t = \sum_{k=1}^m \sigma^{ik}_t \sigma^{jk}_t dt\)。

\textbf{证明}：多维 Taylor 展开：\(df = f_t dt + \sum_i f_{x_i} dX^i + \frac{1}{2}\sum_{i,j} f_{x_i x_j} dX^i dX^j\)。代入 \(dX^i = \mu^i dt + \sum_k \sigma^{ik} dW^k\)，保留所有 \(O(dt)\) 项。利用 \(dW^k dW^\ell = \delta_{k\ell} dt\) 和 \((dt)^2 = o(dt)\)，交叉项 \(dt \cdot dW^k\) 为 \(o(dt)\)。合并得 \(d\langle X^i, X^j \rangle_t = \sum_k \sigma^{ik} \sigma^{jk} dt\)。形式化极限过程由分划逼近完成。\(\blacksquare\)

\textbf{命题 135.1}（Itô 公式的应用------随机指数）：设 \(X_t\) 为 Itô 过程，\(dX_t = \mu_t dt + \sigma_t dW_t\)。定义随机指数（Doleans-Dade 指数）\(\mathcal{E}(X)_t = \exp(X_t - X_0 - \frac{1}{2}[X, X]_t)\)。由 Itô 公式，\(d\mathcal{E}(X)_t = \mathcal{E}(X)_t dX_t\)（即 \(\mathcal{E}(X)_t\) 是随机微分方程 \(dZ_t = Z_t dX_t\) 的解）。随机指数是 Girsanov 定理中的密度过程。

\textbf{命题 135.2}（Itô 乘积公式）：对两个 Itô 过程 \(X_t, Y_t\)，\(d(X_t Y_t) = X_t dY_t + Y_t dX_t + d\langle X, Y \rangle_t\)。这是 Itô 公式在 \(f(x,y) = xy\) 上的特例，额外的协变差项 \(d\langle X, Y \rangle_t\) 是随机微积分区别于经典微积分的标志。

\subsection{135.4 Girsanov 定理与测度变换}\label{girsanov-ux5b9aux7406ux4e0eux6d4bux5ea6ux53d8ux6362}

\textbf{定理 135.8}（Girsanov 定理）：设 \(\theta_t\) 是适应过程，满足 Novikov 条件 \(E[\exp(\frac{1}{2} \int_0^T \theta_t^2 dt)] < \infty\)。定义

\[Z_t = \exp\left(\int_0^t \theta_s dW_s - \frac{1}{2} \int_0^t \theta_s^2 ds\right)\]

在概率测度 \(\tilde{P}\)（\(d\tilde{P}|_{\mathcal{F}_T} = Z_T dP\)）下，过程

\[\tilde{W}_t = W_t - \int_0^t \theta_s ds\]

是 \(\tilde{P}\)-Brown 运动。

\textbf{证明}：验证 \(\tilde{W}_t\) 的 Levy 特征。由 Ito 公式，\(dZ_t = Z_t \theta_t dW_t\)。对 \(\tilde{P}\) 下任意有界 \(\mathcal{F}_s\)-可测函数 \(F\)，\(\tilde{\mathbb{E}}[F(\tilde{W}_t - \tilde{W}_s)] = \mathbb{E}[Z_T F(\tilde{W}_t - \tilde{W}_s)] = \mathbb{E}[Z_s F(W_t - W_s - \int_s^t \theta_r dr)]\)。由 Bayes 公式和 Novikov 条件，\(\tilde{W}_t\) 的有限维分布与 Brown 运动一致，由 Levy 特征定理即为 Brown 运动。Novikov 条件保证了 Girsanov 密度 \(Z_t\) 是严格正的鞅。\(\blacksquare\)

\emph{应用}：Girsanov 定理是金融数学中风险中性定价的基础。它允许通过概率测度的变换消除漂移项，将任何等价测度下的 Itô 过程转化为鞅。在 Black-Scholes 模型中，股票价格 \(dS_t = \mu S_t dt + \sigma S_t dW_t\) 在风险中性测度 \(\tilde{P}\)（\(\theta = (\mu - r)/\sigma\)）下变为 \(dS_t = r S_t dt + \sigma S_t d\tilde{W}_t\)，贴现价格 \(e^{-rt} S_t\) 成为鞅，从而期权价格可表示为期望。

\textbf{定理 135.9}（鞅表示定理）：设 \(M_t\) 是关于 Brown 运动滤子的 \(L^2\)-鞅。则存在唯一的适应过程 \(H_t\)（\(\int_0^T H_t^2 dt < \infty\) a.s.）使得

\[M_t = M_0 + \int_0^t H_s dW_s\]

\textbf{证明}：设 \(\mathcal{A} = \{\int_0^T H_s dW_s : H \in L^2(\Omega \times [0,T])\}\) 为 Itô 积分空间。\(\mathcal{A}\) 是 \(L^2(\Omega, \mathcal{F}_T, P)\) 的闭子空间。欲证 \(\mathcal{A} = L^2(\Omega, \mathcal{F}_T, P) \ominus \mathbb{R}\)（均值为 0 的全体）。若 \(X \perp \mathcal{A}\)，定义 \(M_t = \mathbb{E}[X \mid \mathcal{F}_t]\)。由 Itô 表示，\(M_t\) 可表示为 \(\int_0^t H_s dW_s\)。正交性导出 \(\mathbb{E}[X \int_0^T H_s dW_s] = 0\)，继而 \(X = 0\)。因此 \(\mathcal{A}\) 是满的。唯一性由 Itô 等距保证。\(\blacksquare\)

\subsection{135.5 随机微分方程（SDE）}\label{ux968fux673aux5faeux5206ux65b9ux7a0bsde}

\textbf{定义 135.4}（随机微分方程）：\textbf{随机微分方程}（SDE）的一般形式为： \[dX_t = b(t, X_t) dt + \sigma(t, X_t) dW_t, \quad X_0 = \xi\] 其中 \(b\) 为漂移系数（drift），\(\sigma\) 为扩散系数（diffusion），\(W_t\) 为 \(m\) 维 Brown 运动，\(X_t\) 为 \(n\) 维状态过程。SDE 的积分形式为： \[X_t = \xi + \int_0^t b(s, X_s) ds + \int_0^t \sigma(s, X_s) dW_s\]

\textbf{定理 135.10}（SDE 解的存在唯一性------Itô 定理）：设 \(b, \sigma\) 满足 Lipschitz 条件（\(\|b(t,x)-b(t,y)\| + \|\sigma(t,x)-\sigma(t,y)\| \leq K\|x-y\|\)）和线性增长条件（\(\|b(t,x)\|^2 + \|\sigma(t,x)\|^2 \leq K^2(1+\|x\|^2)\)）。则 SDE 存在唯一的强解（strong solution）\(X_t\)，且 \(X_t\) 具有连续样本路径，满足 \(\mathbb{E}[\sup_{0 \leq t \leq T} \|X_t\|^2] < \infty\)。

\textbf{证明}：利用 Picard 迭代 \(X_t^{(0)} = \xi\)，\(X_t^{(n+1)} = \xi + \int_0^t b(s, X_s^{(n)}) ds + \int_0^t \sigma(s, X_s^{(n)}) dW_s\)。由 Lipschitz 条件，\(\mathbb{E}[\sup_{0 \leq s \leq t} \|X_s^{(n+1)} - X_s^{(n)}\|^2] \leq C \int_0^t \mathbb{E}[\|X_s^{(n)} - X_s^{(n-1)}\|^2] ds\)。由 Gronwall 不等式，序列在 \(L^2\) 中收敛，且极限满足 SDE。唯一性由类似的 Gronwall 论证给出。\(\blacksquare\)

\textbf{定义 135.5}（强解与弱解）：\textbf{强解}是给定 Brown 运动 \(W\) 和滤子 \(\mathcal{F}_t^W\) 下的解，\(X_t\) 是 \(\mathcal{F}_t^W\)-适应的。\textbf{弱解}是 \((X, W)\) 在某个概率空间上的联合分布满足 SDE，不要求 \(X_t\) 适应于 \(\mathcal{F}_t^W\)。弱解的存在性条件比强解弱（仅需连续性和线性增长），且弱解在分布意义下唯一时称为弱唯一性。Yamada-Watanabe 定理（1971）指出：弱解存在 + 轨道唯一性 \(\Rightarrow\) 强解存在且唯一。

\textbf{定理 135.11}（Feynman-Kac 公式）：设 \(u(t, x)\) 满足抛物型 PDE： \[\frac{\partial u}{\partial t} + \frac{1}{2} \sigma^2(x) \frac{\partial^2 u}{\partial x^2} + b(x) \frac{\partial u}{\partial x} - r(x) u + f(x) = 0, \quad u(T, x) = g(x)\] 则 \(u(t, x)\) 有概率表示： \[u(t, x) = \mathbb{E}^{t,x}\left[\int_t^T e^{-\int_t^s r(X_u) du} f(X_s) ds + e^{-\int_t^T r(X_u) du} g(X_T)\right]\] 其中 \(X_t\) 是 SDE \(dX_s = b(X_s) ds + \sigma(X_s) dW_s\) 的解（\(X_t = x\)）。Feynman-Kac 公式将 PDE 的解表示为随机过程的期望，是随机分析和偏微分方程之间的核心桥梁，也是 Monte Carlo 方法求解高维 PDE 的理论基础。

\textbf{证明}：对 \(M_s = e^{-\int_t^s r(X_u) du} u(s, X_s)\) 应用 Itô 公式，利用 \(u\) 满足的 PDE，得 \(dM_s = -e^{-\int_t^s r du} f(X_s) ds + (\text{鞅项})\)。取期望 \(\mathbb{E}^{t,x}[M_T] = u(t,x) - \mathbb{E}^{t,x}[\int_t^T e^{-\int_t^s r du} f(X_s) ds]\)。由终端条件 \(M_T = e^{-\int_t^T r du} g(X_T)\)，即得 Feynman-Kac 表示。\(\blacksquare\)

\textbf{定理 135.12}（Markov 性与生成元）：SDE 的解 \(X_t\) 是强 Markov 过程，其无穷小生成元为： \[\mathcal{L} f(x) = \sum_i b_i(x) \frac{\partial f}{\partial x_i} + \frac{1}{2} \sum_{i,j} (\sigma\sigma^T)_{ij}(x) \frac{\partial^2 f}{\partial x_i \partial x_j}\]

生成元 \(\mathcal{L}\) 是 Itô 公式的期望形式：\(\mathbb{E}^x[f(X_t)] - f(x) = \mathbb{E}^x[\int_0^t \mathcal{L}f(X_s) ds]\)。Markov 过程的半群 \(P_t f(x) = \mathbb{E}^x[f(X_t)]\) 满足 Kolmogorov 向后方程 \(\frac{\partial}{\partial t} P_t f = \mathcal{L} P_t f\) 和向前方程（Fokker-Planck 方程）\(\frac{\partial}{\partial t} p(t,x) = \mathcal{L}^* p(t,x)\)（\(p(t,x)\) 为转移密度）。

\textbf{定理 135.13}（Black-Scholes 公式）：在 Black-Scholes 模型中，股票价格 \(dS_t = \mu S_t dt + \sigma S_t dW_t\)，欧式看涨期权的价格为： \[C(S, t) = S \Phi(d_1) - K e^{-r(T-t)} \Phi(d_2)\] 其中 \(d_1 = \frac{\ln(S/K) + (r + \sigma^2/2)(T-t)}{\sigma\sqrt{T-t}}\)，\(d_2 = d_1 - \sigma\sqrt{T-t}\)，\(\Phi\) 为标准正态分布函数。

\textbf{证明}：由 Feynman-Kac 公式（或等价地，由 Girsanov 定理在风险中性测度下），\(C(S,t) = e^{-r(T-t)} \mathbb{E}^{\tilde{P}}[(S_T - K)^+ \mid S_t = S]\)。在 \(\tilde{P}\) 下，\(S_T = S \exp(\sigma \tilde{W}_{T-t} + (r - \sigma^2/2)(T-t))\)。直接计算期望即得 Black-Scholes 公式。\(\blacksquare\)

\hrulefill

\newpage

\chapter{10-概率与统计/02-随机过程/04-随机微分方程进阶.md}\label{ux6982ux7387ux4e0eux7edfux8ba102-ux968fux673aux8fc7ux7a0b04-ux968fux673aux5faeux5206ux65b9ux7a0bux8fdbux9636.md}

\section{第357章：大偏差理论与集中不等式}\label{ux7b2c357ux7ae0ux5927ux504fux5deeux7406ux8bbaux4e0eux96c6ux4e2dux4e0dux7b49ux5f0f}

大偏差理论估计罕见事件概率的指数衰减速率，是统计力学（熵最大化原理）、信息论（Sanov 定理）和风险管理的基础工具。

\subsection{140.1 Cramér 定理}\label{cramuxe9r-ux5b9aux7406}

\textbf{定理 140.1}（Cramér 定理，1938）：设 \(X_1, X_2, \ldots\) 是 i.i.d. 随机变量，矩母函数 \(\Lambda(\lambda) = \log E[e^{\lambda X_1}] < \infty\) 对某区间 \(\lambda \in [-\varepsilon, \varepsilon]\) 成立。则样本均值 \(\bar{X}_n = \frac{1}{n} \sum_{i=1}^n X_i\) 满足大偏差原理（LDP），速率函数为 Legendre 变换

\[I(x) = \sup_{\lambda \in \mathbb{R}} \{\lambda x - \Lambda(\lambda)\}\]

即：对开集 \(G\)，\(\liminf_n \frac{1}{n} \log P(\bar{X}_n \in G) \geq -\inf_{x \in G} I(x)\)；对闭集 \(F\)，\(\limsup_n \frac{1}{n} \log P(\bar{X}_n \in F) \leq -\inf_{x \in F} I(x)\)。

\textbf{证明}：上界：对任意闭集 \(F\)，取 \(\lambda\) 使 \(\lambda x - \Lambda(\lambda) \geq I(x) - \varepsilon\)。由 Markov 不等式，\(P(\bar{X}_n \in F) \leq e^{-n(\lambda x - \Lambda(\lambda))} \cdot \mathbb{E}[e^{\lambda \sum X_i}] = e^{-n(\lambda x - \Lambda(\lambda))}\)。对 \(F\) 取紧致化，由 \(\varepsilon\) 任意性得 \(\limsup_n \frac{1}{n} \log P(\bar{X}_n \in F) \leq -\inf_{x \in F} I(x)\)。下界：对任意开集 \(G\)，取 \(x_0 \in G\) 和 \(\delta > 0\) 使 \(B_\delta(x_0) \subset G\)。由指数倾斜（Cramér 变换），构造测度 \(Q_n\) 使 \(\bar{X}_n\) 在 \(Q_n\) 下集中于 \(x_0\)。由测度变换公式，\(P(\bar{X}_n \in G) \geq P(\bar{X}_n \in B_\delta(x_0)) \geq e^{-n(\Lambda^*(\lambda_0) + \varepsilon)} Q_n(\bar{X}_n \in B_\delta(x_0))\)，其中 \(\lambda_0\) 满足 \(\Lambda'(\lambda_0) = x_0\)。由大数定律，\(Q_n(\bar{X}_n \in B_\delta(x_0)) \to 1\)，得 \(\liminf \geq -I(x_0)\)。局部化得整体下界。\(\blacksquare\)

\textbf{Corollary 140.2}（Chernoff 界）：\(P(\bar{X}_n \geq a) \leq e^{-n I(a)}\)（对 \(a > E[X_1]\)）。

\textbf{证明}：由 Cramér 定理，\(I(a) = \sup_{\lambda \geq 0} \{\lambda a - \Lambda(\lambda)\}\)（对 \(a > E[X_1]\)，最优 \(\lambda\) 非负）。对任意 \(\lambda \geq 0\)，\(P(\bar{X}_n \geq a) = P(e^{\lambda \sum X_i} \geq e^{\lambda n a}) \leq e^{-\lambda n a} \mathbb{E}[e^{\lambda \sum X_i}] = e^{-n(\lambda a - \Lambda(\lambda))}\)。对 \(\lambda\) 取上确界即得 \(P(\bar{X}_n \geq a) \leq e^{-n I(a)}\)。\(\blacksquare\)

\subsection{140.2 Schilder 定理与 Freidlin-Wentzell 理论}\label{schilder-ux5b9aux7406ux4e0e-freidlin-wentzell-ux7406ux8bba}

\textbf{定理 140.3}（Schilder 定理）：Brown 运动 \(W_t\)（\(t \in [0,1]\)）经缩放 \(\sqrt{\varepsilon} W_t\) 满足参数 \(\varepsilon \to 0\) 的 LDP，速率函数为 Cameron-Martin 范数：

\[I(\varphi) = \begin{cases} \frac{1}{2} \int_0^1 |\dot{\varphi}(t)|^2 dt & \varphi \in H^1([0,1]), \varphi(0)=0 \\ \infty & \text{否则} \end{cases}\]

\textbf{证明}：由 Cramér 定理在无限维空间的推广（Dawson-Gärtner 定理）。将 \([0,1]\) 离散化为 \(n\) 个区间，\(\sqrt{\varepsilon} W_t\) 的增量 \(\sqrt{\varepsilon}(W_{t_i} - W_{t_{i-1}})\) 为 i.i.d. \(\mathcal{N}(0,\varepsilon/n)\)。由 Cramér 定理，每个增量的速率函数为 \(\frac{1}{2}x^2/(\varepsilon/n) = \frac{n}{2\varepsilon}x^2\)。取极限 \(n \to \infty\)，离散增量之和的速率函数收敛到 \(\int |\dot{\varphi}|^2/(2\varepsilon) dt\) 但无需除以 \(\varepsilon\)（因 \(\varepsilon\) 为参数）。更严格地，对 \(\sqrt{\varepsilon}W\)，其协方差算子为 \(\varepsilon \cdot \min(s,t)\)。由 Schilder 定理的抽象版本，Gauss 测度 \(\mu_\varepsilon\)（\(\sqrt{\varepsilon}W\) 的分布）满足 LDP，速率函数 \(I(\varphi) = \frac{1}{2}\|\varphi\|_{H}^2\) 其中 \(H\) 为 Cameron-Martin 空间 \(H^1_0([0,1])\)，范数即为 \(\frac{1}{2}\int |\dot{\varphi}|^2 dt\)。\(\blacksquare\)

\textbf{定理 140.4}（Freidlin-Wentzell 理论）：SDE \(dX^\varepsilon_t = b(X^\varepsilon_t) dt + \sqrt{\varepsilon} \sigma(X^\varepsilon_t) dW_t\) 的样本路径的 LDP 速率函数为

\[I(\varphi) = \frac{1}{2} \int_0^T (\dot{\varphi} - b(\varphi))^T (\sigma \sigma^T)^{-1}(\varphi) (\dot{\varphi} - b(\varphi)) dt\]

\textbf{证明}：由收缩原理（contraction principle）。SDE 的解映射 \(W \mapsto X^\varepsilon\) 为连续映射。由 Schilder 定理，\(\sqrt{\varepsilon}W\) 满足 LDP，速率函数为 \(\frac{1}{2}\int |\dot{\psi}|^2 dt\)。SDE \(dX^\varepsilon = b(X^\varepsilon)dt + \sqrt{\varepsilon}\sigma(X^\varepsilon)dW\) 可写为 \(dX^\varepsilon = \sigma(X^\varepsilon)(\sigma^{-1}(X^\varepsilon)(dX^\varepsilon - b(X^\varepsilon)dt))\)。当 \(\varepsilon \to 0\) 时，\(X^\varepsilon\) 的速率函数由收缩原理给出：\(I(\varphi) = \inf\{\frac{1}{2}\int |\dot{\psi}|^2 dt : \dot{\varphi} = b(\varphi) + \sigma(\varphi)\dot{\psi}\}\)。代入 \(\dot{\psi} = \sigma^{-1}(\varphi)(\dot{\varphi} - b(\varphi))\)，得 \(I(\varphi) = \frac{1}{2}\int |\sigma^{-1}(\varphi)(\dot{\varphi} - b(\varphi))|^2 dt = \frac{1}{2}\int (\dot{\varphi} - b)^\top (\sigma\sigma^\top)^{-1} (\dot{\varphi} - b) dt\)。\(\blacksquare\)

\subsection{140.3 集中不等式}\label{ux96c6ux4e2dux4e0dux7b49ux5f0f}

\textbf{定理 140.5}（Hoeffding 不等式）：设 \(X_1, \ldots, X_n\) 是独立随机变量，\(a_i \leq X_i \leq b_i\) a.s.，则对 \(t > 0\)，

\[P\left(\sum_{i=1}^n (X_i - E[X_i]) \geq t\right) \leq \exp\left(-\frac{2t^2}{\sum_{i=1}^n (b_i - a_i)^2}\right)\]

\textbf{证明}：由 Hoeffding 引理，对 \(a \leq X \leq b\) a.s. 且 \(\mathbb{E}[X] = 0\)，\(\mathbb{E}[e^{\lambda X}] \leq \exp(\lambda^2(b-a)^2/8)\)。令 \(S_n = \sum (X_i - \mathbb{E}[X_i])\)。对任意 \(\lambda > 0\)，\(P(S_n \geq t) = P(e^{\lambda S_n} \geq e^{\lambda t}) \leq e^{-\lambda t} \mathbb{E}[e^{\lambda S_n}] = e^{-\lambda t} \prod_i \mathbb{E}[e^{\lambda (X_i - \mathbb{E}[X_i])}] \leq e^{-\lambda t} \prod_i \exp(\lambda^2(b_i-a_i)^2/8) = \exp(-\lambda t + \frac{\lambda^2}{8}\sum (b_i-a_i)^2)\)。对 \(\lambda\) 极小化，取 \(\lambda^* = 4t / \sum (b_i-a_i)^2\)，得 \(P(S_n \geq t) \leq \exp(-2t^2 / \sum (b_i-a_i)^2)\)。\(\blacksquare\)

\textbf{定理 140.6}（McDiarmid 不等式---有界差异不等式）：设 \(f: \mathbb{R}^n \to \mathbb{R}\) 对每个坐标满足有界差异性质 \(|f(x_1, \ldots, x_i, \ldots, x_n) - f(x_1, \ldots, x'_i, \ldots, x_n)| \leq c_i\)，\(X_1, \ldots, X_n\) 独立。则

\[P(f(X_1, \ldots, X_n) - E[f] \geq t) \leq \exp\left(-\frac{2t^2}{\sum_{i=1}^n c_i^2}\right)\]

\textbf{证明}：构造 Doob 鞅 \(D_k = \mathbb{E}[f(X) \mid X_1, \ldots, X_k]\)。\(D_0 = \mathbb{E}[f]\)，\(D_n = f(X)\)。\(f(X) - \mathbb{E}[f] = \sum_{k=1}^n (D_k - D_{k-1})\)。由有界差异性质，\(|\mathbb{E}[f(X) \mid X_1,\ldots,X_{k-1}, X_k = x_k] - \mathbb{E}[f(X) \mid X_1,\ldots,X_{k-1}, X_k = x'_k]| \leq c_k\)。故 \(D_k - D_{k-1}\) 条件于 \(X_1,\ldots,X_{k-1}\) 下取值于长度至多 \(c_k\) 的区间内。由 Hoeffding 引理（条件版本），\(\mathbb{E}[e^{\lambda(D_k - D_{k-1})} \mid X_1,\ldots,X_{k-1}] \leq e^{\lambda^2 c_k^2/8}\)。由递推，\(\mathbb{E}[e^{\lambda(f - \mathbb{E}[f])}] = \mathbb{E}[\prod e^{\lambda(D_k - D_{k-1})}] \leq e^{\lambda^2 \sum c_k^2/8}\)。由 Chernoff 方法，\(P(f - \mathbb{E}[f] \geq t) \leq e^{-\lambda t + \lambda^2 \sum c_k^2/8}\)。取 \(\lambda = 4t/\sum c_k^2\)，得 \(P(f - \mathbb{E}[f] \geq t) \leq e^{-2t^2/\sum c_k^2}\)。\(\blacksquare\)

\textbf{定理 140.7}（Talagrand 集中不等式）：在乘积空间 \(\Omega = \prod \Omega_i\) 上配备乘积概率测度，任意凸 1-Lipschitz 函数 \(f\) 满足 Gauss 型集中：

\[P(|f - M_f| \geq t) \leq C e^{-t^2/C}\]

其中 \(M_f\) 是 \(f\) 的中位数。这一不等式通过 Marton 的信息传输方法和 Bobkov-Götze 条件建立。

\textbf{证明}：Talagrand 的证明使用诱导距离方法。对乘积空间 \(\Omega = \prod \Omega_i\)，定义凸距离 \(d_T(x, A) = \sup_{\alpha \in \mathbb{R}^n, \|\alpha\|_2 \leq 1} \inf_{y \in A} \sum_i \alpha_i \mathbf{1}_{x_i \neq y_i}\)。Talagrand 证明了 \(P(A) \cdot \mathbb{E}[e^{d_T(X, A)^2/4}] \leq 1\)。取 \(A = \{x: f(x) \leq M_f\}\)，则对 \(x\) 满足 \(f(x) \geq M_f + t\)，凸 Lipschitz 性给出 \(d_T(x, A) \geq t/C\)。由上述不等式，\(P(f \geq M_f + t) \leq P(A)^{-1} e^{-t^2/(4C^2)} \leq 2e^{-t^2/(4C^2)}\)。同理 \(P(f \leq M_f - t) \leq 2e^{-t^2/(4C^2)}\)。合并得 Gauss 型集中 \(P(|f - M_f| \geq t) \leq C e^{-t^2/C}\)。\(\blacksquare\)

\subsection{140.4 经验过程与一致大数定律}\label{ux7ecfux9a8cux8fc7ux7a0bux4e0eux4e00ux81f4ux5927ux6570ux5b9aux5f8b}

\textbf{定理 140.8}（Glivenko-Cantelli 定理）：设 \(F\) 是 VC 类（Vapnik-Chervonenkis 类），则经验分布函数按 \(\|\cdot\|_\infty\) 一致收敛于真实分布 a.s.

\textbf{证明}：对 VC 类 \(\mathcal{F}\)，其 VC 维 \(d\) 有限。由 Sauer 引理，\(\mathcal{F}\) 在 \(n\) 个样本上的粉碎系数 \(\Pi_{\mathcal{F}}(n) \leq (en/d)^d\)。由对称化引理，\(\mathbb{E}[\sup_{f \in \mathcal{F}} |\mathbb{P}_n f - Pf|] \leq 2 \mathcal{R}_n(\mathcal{F})\)。由 Massart 引理，\(\mathcal{R}_n(\mathcal{F}) \leq \sqrt{2d\log(en/d)/n}\)。由 McDiarmid 不等式，以概率至少 \(1-\delta\)，\(\sup_{f \in \mathcal{F}} |\mathbb{P}_n f - Pf| \leq 2\mathcal{R}_n(\mathcal{F}) + \sqrt{\log(1/\delta)/(2n)} = O(\sqrt{d\log n / n})\)。由 Borel-Cantelli 引理，\(\sup_{f \in \mathcal{F}} |\mathbb{P}_n f - Pf| \to 0\) a.s.。\(\blacksquare\)

\textbf{定理 140.9}（Donsker 定理---经验过程的泛函中心极限定理）：在适当的函数类上，经验过程 \(\sqrt{n}(\mathbb{P}_n - P)\) 弱收敛于 \(P\)-Brown 桥（均值为零、协方差为 \(C(f,g) = P(fg) - P(f)P(g)\) 的 Gauss 过程）。

\textbf{证明}：对 Donsker 类 \(\mathcal{F}\)（满足一致熵条件），经验过程 \(\mathbb{G}_n = \sqrt{n}(\mathbb{P}_n - P)\) 的有限维分布收敛：由多元中心极限定理，\((\mathbb{G}_n f_1, \ldots, \mathbb{G}_n f_k) \to \mathcal{N}(0, \Sigma)\)，其中 \(\Sigma_{ij} = P(f_i f_j) - P(f_i)P(f_j)\)。随机等度连续性由一致熵上界保证：\(\mathbb{E}[\sup_{\|f-g\|_{L^2(P)} < \delta} |\mathbb{G}_n(f-g)|] \to 0\)（\(\delta \to 0\)）。该条件等价于 \(\mathcal{F}\) 为 \(P\)-Donsker 类。由 Pollard 的泛函中心极限定理，\(\mathbb{G}_n\) 在 \(\ell^\infty(\mathcal{F})\) 中弱收敛到 Gauss 过程 \(\mathbb{G}_P\)，其协方差为 \(C(f,g) = P(fg) - P(f)P(g)\)。\(\mathbb{G}_P\) 为 \(P\)-Brown 桥。VC 类是最重要的 Donsker 类。\(\blacksquare\)

\hrulefill

\emph{本卷在概率论 II（V21）的基础上，建立了随机分析的核心理论：Itô 随机积分构造与 Girsanov 定理、Feynman-Kac 公式、Malliavin 分析的分部积分公式与亚椭圆性定理、随机偏微分方程的正则性结构理论、Schramm-Loewner 演化（SLE）与统计物理的联系，以及 Lévy 过程的 Lévy-Khintchine 表示与大偏差理论。共 7 章。}

\hrulefill

\emph{本卷在初等概率论（V6）和测度论（V9）的基础上，建立了概率论的严格分析基础：Kolmogorov 公理体系将概率论纳入测度论框架；大数定律和中心极限定理揭示了随机现象中的确定性规律；条件期望和鞅论提供了处理信息的强大工具；Markov 链理论为随机过程建模提供了基本框架；Brown 运动和 Itô 积分开创了随机分析，为金融数学（Black-Scholes 模型）和随机微分方程奠定了基础。为后续数理统计（V22）提供了概率论基础。\textbf{补充部分一}将概率论极限理论应用于统计力学：Gibbs 测度理论建立了平衡态统计力学的概率论基础，相变理论通过配分函数的非解析性揭示了临界现象，重整化群方法从不动点理论出发统一了临界指数的普适性。\textbf{补充部分二}引入了随机矩阵理论：Wigner半圆律与普遍性、自由概率论（Voiculescu的R-变换与S-变换）、Tracy-Widom分布与边缘统计、圆律与Stieltjes变换方法、行列式点过程与\(\beta\)系综，以及随机矩阵在数论（Riemann \(\zeta\)零点）、高维统计与机器学习中的现代应用。\textbf{补充部分三}建立了随机分析的完整体系：Itô积分深化与Girsanov定理、随机微分方程（SDE）解的存在唯一性与Feynman-Kac公式、Malliavin分析与Hörmander条件、随机偏微分方程（SPDE）与正则性结构、随机几何与SLE理论、Lévy过程与测度值过程、大偏差理论（Cramér定理、Schilder定理）与经验过程理论。}

\emph{卷二十：概率论 II终。}

\emph{卷二十：概率论 II终。}

\emph{卷二十：概率论 II终。}

\emph{卷二十三：随机过程与统计力学终。} \hrulefill

\subsection{粗糙路径理论与正则结构导引}\label{ux7c97ux7cd9ux8defux5f84ux7406ux8bbaux4e0eux6b63ux5219ux7ed3ux6784ux5bfcux5f15}

\begin{quote}
\textbf{前置知识}：第288-289章（Itô 积分与 SDE）
\end{quote}

\subsection{动机与问题背景}\label{ux52a8ux673aux4e0eux95eeux9898ux80ccux666f}

考虑由信号 \(h: [0,T] \to \mathbb{R}^d\) 驱动的受控常微分方程

\[dy_t = V(y_t) dh_t, \quad y_0 = \xi \in \mathbb{R}^e\]

其中 \(V = (V_1, \ldots, V_d)\) 为 \(\mathbb{R}^e\) 上的光滑向量场族。经典 Cauchy-Lipschitz 理论要求 \(h \in C^1\) 或至少具有有界变差，以确保 Riemann-Stieltjes 积分收敛。当 \(h\) 仅为 \(\alpha\)-Hölder（\(\alpha \in (0,1)\)）时，Young 积分理论保证 \(\alpha > 1/2\) 时积分收敛；但标准 Brown 运动 \((B_t)\) 几乎必然满足 \(\sup_{s \neq t} |B_t - B_s|/|t-s|^{1/2-\varepsilon} < \infty\)（\(\forall \varepsilon > 0\)），恰好落在 Young 理论边界之下。经典 Itô 积分依赖 Brown 运动的鞅性与独立增量性质，无法推广至确定性粗糙信号。Lyons（1998）的核心洞察是：欲唯一确定由粗糙信号驱动的微分方程的解，需提供信号的高阶迭代积分信息------即\textbf{粗糙路径提升}。该框架不仅包含 Itô 与 Stratonovich 理论作为特例，还适用于任意 \(\alpha\)-Hölder 确定性与随机信号。

粗糙路径理论的动机还源于对信号''粗糙度''的量化需要。当信号的正则性低于 \(1/2\)-Hölder 时，Riemann-Stieltjes 逼近的极限不再唯一依赖于信号的一阶增量，而受到二阶（甚至更高阶）迭代积分的影响。这一现象在 Itô 与 Stratonovich 积分的差异中已有体现：两者的区别在于 \(\frac{1}{2}[B, B]_t = \frac{1}{2}t\) 这一''面积''修正项。Lyons 的推广揭示了这一修正项的普遍性------任何粗糙信号驱动的微分方程都需要类似的高阶修正以保证解的唯一性和连续性。粗糙路径理论将信号 \(h\) 提升为带有一阶、二阶直至 \(N\) 阶迭代积分的对象 \(\mathbf{X} = (1, X^1, X^2, \ldots, X^N)\)，其中 \(X^k_{s,t} = \int_{s < u_1 < \cdots < u_k < t} dh_{u_1} \otimes \cdots \otimes dh_{u_k}\)（在适当正则化下）。这一提升后的对象在 \(p\)-变差度量下构成一个代数结构，使得微分方程的理论得以推广到任意粗糙信号。

\subsection{张量代数预备}\label{ux5f20ux91cfux4ee3ux6570ux9884ux5907}

\textbf{定义 1}（截断张量代数）：设 \(V = \mathbb{R}^d\)，\(N\) 阶截断张量代数为

\[T^{(N)}(V) := \bigoplus_{k=0}^{N} V^{\otimes k} = \mathbb{R} \oplus V \oplus (V \otimes V) \oplus \cdots \oplus V^{\otimes N}\]

\(V^{\otimes 0} := \mathbb{R}\)，乘法由张量积 \(\otimes\) 给出（\(k + \ell > N\) 时截为零）。记 \(\pi_k: T^{(N)}(V) \to V^{\otimes k}\) 为投影。定义李括号 \([a, b] := a \otimes b - b \otimes a\)。对 \(v, w \in V\)，\([v, w] = v \otimes w - w \otimes v \in V^{\otimes 2}\) 量化了迭代积分的非交换性：Itô-Stratonovich 修正项由 \(V^{\otimes 2}\) 中对称部分的迹编码。

截断张量代数 \(T^{(N)}(V)\) 上可定义群结构。设 \(\exp: T^{(N)}(V) \to 1 + \bigoplus_{k=1}^N V^{\otimes k}\) 为张量指数映射 \(\exp(a) = \sum_{k=0}^N \frac{a^{\otimes k}}{k!}\)（其中 \(N\) 级以上截断）。其逆为 \(\log(1 + b) = \sum_{k=1}^N \frac{(-1)^{k-1}}{k} b^{\otimes k}\)。在 \(\exp\) 映射下，\(T^{(N)}(V)\) 成为自由幂零 Lie 群 \(\mathbb{G}^{(N)}(V)\)（即 \(N\) 阶自由幂零 Lie 群的指数像）。粗糙路径 \(\mathbf{X}: [0,T] \to T^{(N)}(V)\) 必须满足 Chen 恒等式：\(\mathbf{X}_{s,t} = \mathbf{X}_{s,u} \otimes \mathbf{X}_{u,t}\)（\(s \leq u \leq t\)），这确保了迭代积分的一致性。Chen 恒等式是粗糙路径理论的代数基础------它将路径的''增量''提升为群值函数，使得微分方程的求解可以转化为群上的运算。此外，\(T^{(N)}(V)\) 上的自然李代数结构由自由 Lie 代数 \(\mathfrak{g}^{(N)}(V)\) 生成，其元素为 \(V\) 中向量的迭代李括号，这为理解粗糙路径的代数约束（如 shuffle 积关系）提供了 Lie 理论的语言。

\subsection{\texorpdfstring{控制函数与 \(p\)-变差}{控制函数与 p-变差}}\label{ux63a7ux5236ux51fdux6570ux4e0e-p-ux53d8ux5dee}

\textbf{定义 2}（\(p\)-变差）：设 \(X: [0,T] \to \mathbb{R}^d\) 连续，\(p \geq 1\)。\(X\) 在 \([s,t]\) 上的 \(p\)-变差定义为

\[\|X\|_{p\text{-var}; [s,t]} := \left( \sup_{\pi = \{s=t_0 < \cdots < t_n = t\}} \sum_{i=0}^{n-1} |X_{t_{i+1}} - X_{t_i}|^p \right)^{1/p}\]

其中上确界遍取 \([s,t]\) 的所有有限分划 \(\pi\)。

\textbf{定义 3}（控制函数）：称连续函数 \(\omega: \Delta_2 := \{(s,t): 0 \leq s \leq t \leq T\} \to [0,\infty)\) 为控制函数，若

\[\omega(s,u) \leq \omega(s,t) + \omega(t,u), \quad \forall 0 \leq s \leq t \leq u \leq T\]

且 \(\omega(t,t) = 0\)。若 \(\|X\|_{p\text{-var}; [0,T]} < \infty\)，则 \(\omega_X(s,t) := \|X\|_{p\text{-var}; [s,t]}^p\) 为控制函数。

\textbf{定理 1}（控制函数的正则性控制）：设 \(\omega\) 为控制函数，\(\|X_{s,t}\| \leq \omega(s,t)^{1/p}\) 对某 \(p \geq 1\) 成立，则对任意 \(q > p\)，有 \(\|X\|_{q\text{-var}; [s,t]} \lesssim \omega(s,t)^{1/q} \cdot \omega(s,t)^{(q-p)/(pq)}\)。

\textbf{证明}：取分划 \(\pi\)，对每项使用 \(\|X_{t_i, t_{i+1}}\| \leq \omega(t_i, t_{i+1})^{1/p}\)，则

\[\sum_i \|X_{t_i, t_{i+1}}\|^q \leq \max_i \|X_{t_i, t_{i+1}}\|^{q-p} \sum_i \|X_{t_i, t_{i+1}}\|^p \leq \sup_i \omega(t_i, t_{i+1})^{(q-p)/p} \cdot \omega(s,t)\]

取上确界即得。\(\blacksquare\)

\subsection{粗糙路径的正式定义}\label{ux7c97ux7cd9ux8defux5f84ux7684ux6b63ux5f0fux5b9aux4e49}

\textbf{定义 4}（\(\alpha\)-Hölder 粗糙路径）：设 \(\alpha \in (0,1]\)，\(N := \lfloor 1/\alpha \rfloor\)。一个 \(\alpha\)-Hölder 粗糙路径是映射

\[\mathbb{X}: \Delta_2 \to T^{(N)}(\mathbb{R}^d), \quad (s,t) \mapsto \mathbb{X}_{s,t} = (1, X_{s,t}^1, X_{s,t}^2, \ldots, X_{s,t}^N)\]

满足以下条件：

\begin{enumerate}
\def\labelenumi{(\roman{enumi})}
\tightlist
\item
  \textbf{Chen 等式}：\(\mathbb{X}_{s,u} = \mathbb{X}_{s,t} \otimes \mathbb{X}_{t,u}\) 对所有 \(s \leq t \leq u\)。按分量展开：
\end{enumerate}

\[X_{s,u}^k = \sum_{i=0}^{k} X_{s,t}^i \otimes X_{t,u}^{k-i}, \qquad X_{s,t}^0 := 1\]

\begin{enumerate}
\def\labelenumi{(\roman{enumi})}
\setcounter{enumi}{1}
\tightlist
\item
  \textbf{解析正则性}：存在 \(C > 0\) 使 \(\|X_{s,t}^k\| \leq C |t-s|^{k\alpha} / \Gamma(k\alpha + 1)\) 对 \(k = 1,\ldots,N\) 成立，其中 \(\Gamma\) 为 Gamma 函数。
\end{enumerate}

若 \(h \in C^1([0,T]; \mathbb{R}^d)\)，其规范提升的 \(k\) 重迭代积分 \(X_{s,t}^k := \int_{s < u_1 < \cdots < u_k < t} dh_{u_1} \otimes \cdots \otimes dh_{u_k}\) 自动满足 Chen 等式，因积分可按中间点 \(t\) 拆分。

\textbf{定义 5}（几何粗糙路径）：若 \(\alpha\)-Hölder 粗糙路径 \(\mathbb{X}\) 可由光滑路径 \(\{h^{(n)}\}\) 的规范提升在 \(\alpha\)-Hölder 度量下逼近而得，则称 \(\mathbb{X}\) 为几何粗糙路径。几何粗糙路径附加满足对称性

\[\operatorname{Sym}(X_{s,t}^k) = \frac{1}{k!} (X_{s,t}^1)^{\otimes k}\]

其中 \(\operatorname{Sym}\) 为对称化算子，即 \(k\) 阶分量由一阶分量的张量幂完全确定（在截断阶以下的近似意义）。

\textbf{定义 6}（\(\alpha\)-Hölder 度量）：两个 \(\alpha\)-Hölder 粗糙路径 \(\mathbb{X}, \mathbb{Y}\) 间的距离为

\[\rho_{\alpha}(\mathbb{X}, \mathbb{Y}) := \max_{1 \leq k \leq N} \sup_{0 \leq s < t \leq T} \frac{\|X_{s,t}^k - Y_{s,t}^k\|}{|t-s|^{k\alpha}}\]

该度量使粗糙路径空间成为完备度量空间。

\subsection{粗糙微分方程}\label{ux7c97ux7cd9ux5faeux5206ux65b9ux7a0b}

考虑由 \(\alpha\)-Hölder 粗糙路径 \(\mathbb{X}\) 驱动的 RDE：

\[dY_t = V(Y_t) d\mathbb{X}_t, \quad Y_0 = \xi\]

其中 \(V: \mathbb{R}^e \to L(\mathbb{R}^d, \mathbb{R}^e)\) 为 \(C^\gamma\) 向量场（\(\gamma > 1/\alpha\)）。解 \(Y\) 满足积分方程 \(Y_t = \xi + \int_0^t V(Y_s) d\mathbb{X}_s\)，其中粗糙路径积分通过修正的 Riemann 和逼近。引入 \(\mathbb{X}\) 的规范级数展开

\[\int_s^t V(Y_u) d\mathbb{X}_u \approx \sum_{k=1}^{N} V^{(k-1)}(Y_s) X_{s,t}^k\]

其中 \(V^{(k-1)}\) 表示 \(V\) 的 \((k-1)\) 阶导数关于 \(\mathbb{R}^d\) 的协变作用。

\textbf{定理 2}（Lyons 万有极限定理）：设 \(\gamma > 1/\alpha\)，\(V \in C_b^\gamma(\mathbb{R}^e; L(\mathbb{R}^d, \mathbb{R}^e))\)。则 RDE 的 Itô-Lyons 解映射 \(\Phi: (\xi, \mathbb{X}) \mapsto (Y_t)_{t \in [0,T]}\) 从 \(\mathbb{R}^e \times \mathscr{C}^\alpha([0,T]; T^{(N)}(\mathbb{R}^d))\) 到 \(C^\alpha([0,T]; \mathbb{R}^e)\) 是局部 Lipschitz 连续的。即若 \(\mathbb{X}^{(n)} \to \mathbb{X}\) 在 \(\rho_\alpha\) 下，则 \(Y^{(n)} \to Y\) 在 \(\alpha\)-Hölder 范数下。

\textbf{证明}：核心工具为 Sewing Lemma（Gubinelli, 2004）：若两点函数 \(\Xi_{s,t}\) 满足 \(|\Xi_{s,t}| \lesssim |t-s|^\theta\) 且 \(\delta\Xi_{s,u,t} := \Xi_{s,t} - \Xi_{s,u} - \Xi_{u,t}\) 满足 \(|\delta\Xi_{s,u,t}| \lesssim |t-s|^\theta\)（\(\theta > 1\)），则存在唯一的连续路径 \(Z\) 使 \(Z_t - Z_s \approx \Xi_{s,t}\)。构造 \(\Gamma(Z)_{s,t} := \sum_{k=1}^{N} D^{k-1} V(\xi + Z_s) X_{s,t}^k\)，由 Chen 等式验证 \(\delta\Gamma(Z)\) 满足 \(\theta = (N+1)\alpha \geq 1+\alpha > 1\)。压缩映射论证（粗糙 Grönwall 引理）在 \(T\) 充分小时给出局部唯一解；粘合得全局解。连续性估计 \(\|Y - \tilde{Y}\|_{\alpha} \lesssim |\xi - \tilde{\xi}| + \rho_\alpha(\mathbb{X}, \tilde{\mathbb{X}})\) 表明解映射为局部 Lipschitz。\(\blacksquare\)

\textbf{定理 3}（RDE 解的存在唯一性）：在定理 2 条件下，RDE 对任意 \(\xi \in \mathbb{R}^e\) 和 \(\alpha\)-Hölder 粗糙路径 \(\mathbb{X}\) 存在唯一全局解 \(Y \in C^\alpha([0,T]; \mathbb{R}^e)\)。

\textbf{证明}：定理 2 的压缩映射构造给出局部存在唯一性。设 \(Y, \tilde{Y}\) 均为解，\(D_t := Y_t - \tilde{Y}_t\)。由 Sewing Lemma 线性估计得 \(\|D\|_{\alpha; [0,T]} \leq C_V \cdot \omega(0,T)^{1/p} \cdot \|D\|_{\alpha; [0,T]}\)。取 \(T\) 充分小使 \(C_V \omega(0,T)^{1/p} < 1\)，迫使 \(D \equiv 0\)。区间分裂延拓至全局。\(\blacksquare\)

\subsection{Brown 粗糙路径}\label{brown-ux7c97ux7cd9ux8defux5f84}

\textbf{定义 7}（Brown 运动的规范提升）：设 \((B_t)_{t \in [0,T]}\) 为 \(d\) 维 Brown 运动，对 \(\alpha < 1/2\) 取 \(N = 2\)。Stratonovich 提升 \(\mathbb{B}^{\text{Strat}}\) 定义为

\[B_{s,t}^1 := B_t - B_s, \qquad B_{s,t}^2 := \int_s^t (B_u - B_s) \otimes \circ dB_u\]

该提升为几何粗糙路径，满足 \(\operatorname{Sym}(B_{s,t}^2) = \frac{1}{2} B_{s,t}^1 \otimes B_{s,t}^1\)。Itô 提升 \(\mathbb{B}^{\text{Itô}}\) 将积分换为 Itô 型，满足 \(B_{s,t}^2 = \frac{1}{2} B_{s,t}^1 \otimes B_{s,t}^1 + \frac{1}{2}(t-s) I_d\)（\(I_d\) 视为 \(V^{\otimes 2}\) 元素），非几何粗糙路径。

\textbf{定理 4}（Wong-Zakai 型收敛）：设 \(B^{(n)}\) 为 Brown 运动 \(B\) 在分划 \(\pi_n\)（网格 \(|\pi_n| \to 0\)）上的逐段线性逼近。则对任意 \(q \geq 1\)，\(\alpha < 1/2\)，

\[\lim_{n \to \infty} \mathbb{E}\left[ \rho_\alpha(\mathbb{B}^{(n)}, \mathbb{B}^{\text{Strat}})^q \right] = 0\]

\textbf{证明}：一阶项收敛由 Brown 运动的 \(\alpha\)-Hölder 连续性给出。核心在二阶项：逐段线性逼近的链式法则给出 \(B^{(n)2}_{s,t} = \frac{1}{2} B^{(n)1}_{s,t} \otimes B^{(n)1}_{s,t}\)。由 Burkholder-Davis-Gundy 不等式与 Kolmogorov 连续性准则，\(\mathbb{E}[\|B^{(n)2}_{s,t} - B^{\text{Strat},2}_{s,t}\|^q] \lesssim |t-s|^{q/2} |\pi_n|^{\varepsilon q}\)（\(\varepsilon > 0\)）。除以 \(|t-s|^{2q\alpha}\)（\(2\alpha < 1\)），取上确界并用 Garsia-Rodemich-Rumsey 引理得 \(L^q\) 收敛。\(\blacksquare\)

\subsection{正则结构理论概览}\label{ux6b63ux5219ux7ed3ux6784ux7406ux8bbaux6982ux89c8}

Hairer（2014）将粗糙路径框架推广至奇异 SPDE，提出正则结构理论。

\textbf{定义 8}（正则结构）：正则结构为三元组 \(\mathcal{T} = (A, T, G)\)：\(A \subset \mathbb{R}\) 为局部有限且下有界的指标集（\(\alpha \in A\) 表示抽象符号的正则性阶）；\(T = \bigoplus_{\alpha \in A} T_\alpha\) 为分次向量空间；\(G\) 为 \(T\) 上的线性映射群，满足 \(\Gamma \in G\)，\(\tau \in T_\alpha \implies \Gamma\tau - \tau \in \bigoplus_{\beta < \alpha} T_\beta\)。

\textbf{定义 9}（模型）：\(\mathcal{T}\) 在 \(\mathbb{R}^d\) 上的模型为一对映射 \((\Pi, \Gamma)\)：\(\Pi = (\Pi_x)_{x \in \mathbb{R}^d}\) 将抽象符号实现为分布，\(\Pi_x: T \to \mathcal{D}'(\mathbb{R}^d)\) 满足尺度条件 \(|\Pi_x \tau(\varphi_x^\lambda)| \lesssim \lambda^\alpha\)（\(\tau \in T_\alpha\)）；\(\Gamma = (\Gamma_{xy})_{x,y \in \mathbb{R}^d}\)，\(\Gamma_{xy} \in G\)，满足相容性 \(\Pi_x \Gamma_{xy} = \Pi_y\) 与上圈条件 \(\Gamma_{xy} \Gamma_{yz} = \Gamma_{xz}\)。模型将抽象张量结构实现为分布，\(\Gamma\) 管理基点间的转换。

\textbf{定义 10}（重建算子）：设 \(f: \mathbb{R}^d \to T_{<\gamma}\) 满足 \(f(x) - \Gamma_{xy} f(y) \in \bigoplus_{\beta < \gamma} T_\beta\) 且该差异由 \(\Pi\) 的控制函数约束，则存在唯一 \(\mathcal{R} f \in \mathcal{D}'(\mathbb{R}^d)\) 使 \((\mathcal{R} f - \Pi_x f(x))(\varphi_x^\lambda) \to 0\)（\(\lambda \to 0\)），收敛速度由 \(A\) 的最低指标决定。

\textbf{定理 5}（重建定理，Hairer 2014）：\(\mathcal{R}\) 为良定义的线性映射，\(\mathcal{R} f\) 在 \(x\) 附近的行为由 \(f(x)\) 中最小正则性指标的符号主导。

\textbf{证明}（概略）：构造 \(\mathcal{R} f(\varphi) := \lim_{\varepsilon \to 0} \int_{\mathbb{R}^d} \Pi_x f(x)(\varphi_\varepsilon(x - \cdot)) dx\)。由模拟分布条件与模型的正则性估计，差值界为 \(\max\{\varepsilon, \delta\}^{\min A}\)，极限存在且不依赖软化子。唯一性来自候选分布之差在所有尺度上的衰减迫使二者相等。\(\blacksquare\)

作为里程碑应用：KPZ 方程 \(\partial_t h = \partial_x^2 h + (\partial_x h)^2 + \xi\) 在二维正则结构下经重整化后局部适定；\(\Phi^4_3\) 模型（三维标量场）需处理 \(-\frac{3}{2} + \kappa\) 级奇异性，含多达 12 种抽象符号，经 BPHZ 型重整化后局部适定。粗糙路径恰对应最低阶非平凡正则结构：\(A = \{0, \alpha, \ldots, N\alpha\}\)，\(T \simeq T^{(N)}(\mathbb{R}^d)\)，Chen 等式对应上圈条件 \(\Gamma_{xy}\Gamma_{yz} = \Gamma_{xz}\) 的截断版本。

\subsection{粗糙路径与随机分析的关系}\label{ux7c97ux7cd9ux8defux5f84ux4e0eux968fux673aux5206ux6790ux7684ux5173ux7cfb}

\textbf{定义 11}（受控粗糙路径）：设 \(X\) 为 \(\alpha\)-Hölder 路径。称 \(Y \in C^\alpha\) 被 \(X\) 控制（Gubinelli, 2004），若存在 Gubinelli 导数 \(Y' \in C^\alpha\) 使

\[R_{s,t}^Y := Y_{s,t} - Y'_s X_{s,t}, \qquad \|R_{s,t}^Y\| \lesssim |t-s|^{2\alpha}\]

记此空间为 \(\mathscr{D}_X^{2\alpha}\)。受控粗糙路径框架下，积分 \(\int Y_s d\mathbb{X}_s\) 定义为

\[\int_s^t Y_u d\mathbb{X}_u := \lim_{|\pi| \to 0} \sum_{[u,v] \in \pi} \left( Y_u X_{u,v}^1 + Y'_u X_{u,v}^2 \right)\]

其相容性由 Sewing Lemma 保证。

粗糙路径框架以统一方式涵盖两种随机积分：选择 \(\mathbb{B}^{\text{Strat}}\) 的 RDE 解对应 Stratonovich SDE 的解（满足通常链式法则），选择 \(\mathbb{B}^{\text{Itô}}\) 的 RDE 解对应 Itô SDE 的解。两者转换关系为

\[Y_t^{\text{Itô}} = Y_t^{\text{Strat}} - \frac{1}{2} \int_0^t \sum_{i=1}^d (\partial_i V_i)(Y_s) ds\]

即 Stratonovich-Itô 修正项由 \(V^{\otimes 2}\) 分量的迹编码。粗糙路径理论因此提供了随机积分的纯路径态实现：对 Brown 运动的几乎每一样本路径，解映射是连续函数，无需概率测度结构或鞅论证。

\hrulefill

\section{第358章：随机偏微分方程（SPDE）}\label{ux7b2c358ux7ae0ux968fux673aux504fux5faeux5206ux65b9ux7a0bspde}

随机偏微分方程将偏微分方程与随机噪声结合，用于描述时空随机场和受随机扰动影响的连续物理系统。SPDE 的数学理论在 20 世纪 80 年代迅速发展，由 Walsh、Da Prato、Zabczyk 等人奠定了 \(L^2\) 框架，而 Hairer 的正则性结构理论（2014 年，获 Fields 奖）则解决了奇异 SPDE（如 KPZ 方程和 \(\Phi^4_3\) 模型）的严格数学描述。SPDE 在物理、生物和金融中广泛应用：随机热方程描述随机介质中的热传导，随机 Navier-Stokes 方程描述湍流，随机波动方程描述随机介质中的波传播，KPZ 方程描述随机界面增长，而 \(\Phi^4_3\) 模型是量子场论中 Higgs 机制的理论原型。SPDE 理论的核心困难在于：随机噪声通常高度不规则（如时空白噪声在 \(d \geq 2\) 时不是（广义）函数），导致经典 PDE 的求解方法（如弱解、Galerkin 方法）失效，需要引入重正化和正则性结构等全新工具。

\subsection{136.1 时空白噪声与随机积分}\label{ux65f6ux7a7aux767dux566aux58f0ux4e0eux968fux673aux79efux5206}

\textbf{定义 136.1}（时空白噪声）：\(d\) 维时空白噪声 \(\dot{W}(t,x)\) 是零均值 Gaussian 广义随机场，协方差为 \[\mathbb{E}[\dot{W}(t,x) \dot{W}(s,y)] = \delta(t-s) \delta(x-y)\] 其中 \(\delta\) 为 Dirac delta 函数。形式地，\(\dot{W}\) 是 \(L^2(\mathbb{R}_+ \times \mathbb{R}^d)\) 上的 Gaussian 等正态过程： \[\mathbb{E}\left[\int \phi(t,x) \dot{W}(t,x) dt dx \int \psi(t,x) \dot{W}(t,x) dt dx\right] = \int \phi(t,x) \psi(t,x) dt dx\]

时空白噪声的严格数学构造：定义 \(W_t(\phi) = \int_0^t \int_{\mathbb{R}^d} \phi(x) \dot{W}(s,x) ds dx\) 为圆柱 Brown 运动。\(W_t(\phi)\) 是 \(\phi \in L^2(\mathbb{R}^d)\) 上的一维 Brown 运动（方差 \(t\|\phi\|_{L^2}^2\)），且 \(\mathbb{E}[W_t(\phi) W_s(\psi)] = (t \wedge s) \langle \phi, \psi \rangle_{L^2}\)。在 \(d=1\) 时，时空白噪声与 Brown 运动的关系为 \(\dot{W}(t,x) = \frac{\partial^2}{\partial t \partial x} W(t,x)\)（Brown 单的混合导数），而 \(W(t,x)\) 是 Brown 单（Brownian sheet）：\(\mathbb{E}[W(t,x) W(s,y)] = (t \wedge s)(x \wedge y)\)。

\textbf{定义 136.2}（Walsh 随机积分）：对适应随机场 \(\Phi(t,x)\)，Walsh 随机积分定义为 \[\int_0^t \int_{\mathbb{R}^d} \Phi(s,x) W(ds, dx)\] 这是一个 \(L^2\)-鞅，满足 Walsh 等距： \[\mathbb{E}\left[\left(\int_0^t \int \Phi W(ds,dx)\right)^2\right] = \mathbb{E}\left[\int_0^t \int \Phi(s,x)^2 ds dx\right]\]

Walsh 积分的构造类似于 Itô 积分的构造：首先对简单随机场定义，然后通过 Walsh 等距延拓到 \(L^2\) 空间。Walsh 积分是处理 SPDE 的基本工具，其鞅性质使得能量估计（如 Burkholder-Davis-Gundy 不等式）可以应用于 SPDE 的解。

\textbf{定理 136.1}（Burkholder-Davis-Gundy 不等式）：对 Walsh 鞅积分 \(M_t = \int_0^t \int \Phi(s,x) W(ds,dx)\)，对任意 \(p \geq 2\)： \[\mathbb{E}\left[\sup_{0 \leq s \leq t} |M_s|^p\right] \leq C_p \mathbb{E}\left[\left(\int_0^t \int \Phi(s,x)^2 ds dx\right)^{p/2}\right]\]

\textbf{证明}：设 \(M_t = \int_0^t \int \Phi(s,x) W(ds,dx)\) 为 Walsh 鞅积分。由 Walsh 等距，\(M_t\) 是 \(L^2\)-鞅，其二次变差为 \(\langle M \rangle_t = \int_0^t \int \Phi(s,x)^2 ds dx\)。对 \(p \geq 2\)，由 Ito 公式，\(|M_t|^p\) 可表示为 \(|M_t|^p = \frac{p(p-1)}{2} \int_0^t |M_s|^{p-2} d\langle M \rangle_s + \text{鞅}\)。取期望后利用 Doob 的 \(L^p\) 极大不等式：\(\mathbb{E}[\sup_{s \leq t} |M_s|^p] \leq (p/(p-1))^p \mathbb{E}[|M_t|^p]\)。再结合 Burkholder 不等式（对连续鞅的版本），存在仅依赖于 \(p\) 的常数 \(C_p\) 使得 \(\mathbb{E}[\sup_{s \leq t} |M_s|^p] \leq C_p \mathbb{E}[\langle M \rangle_t^{p/2}]\)。\(\blacksquare\)

\subsection{136.2 随机热方程与随机波动方程}\label{ux968fux673aux70edux65b9ux7a0bux4e0eux968fux673aux6ce2ux52a8ux65b9ux7a0b}

\textbf{定义 136.3}（随机热方程）：\(d\) 维空间中的随机热方程为： \[\frac{\partial u}{\partial t} = \Delta u + \sigma(u) \dot{W}(t, x)\] 其中 \(\Delta\) 为 Laplace 算子。其温和解（mild solution，即积分形式解）为： \[u(t, x) = \int_{\mathbb{R}^d} G_t(x-y) u_0(y) dy + \int_0^t \int_{\mathbb{R}^d} G_{t-s}(x-y) \sigma(u(s,y)) W(ds, dy)\] 其中 \(G_t(x) = (4\pi t)^{-d/2} \exp(-|x|^2/(4t))\) 为热核。

\textbf{定理 136.2}（随机热方程的解的存在性）：在 \(d=1\) 时，若 \(\sigma\) 是 Lipschitz 函数，则随机热方程在 \(L^2\) 意义下存在唯一的温和解。在 \(d=2\) 时，解存在但当 \(\sigma\) 非常数时是分布的（非函数值）。在 \(d \geq 3\) 时，热核 \(G_t\) 的奇性过程导致方程需要重正化。这一维度依赖性是 SPDE 理论的典型特征------热核的积分 \(\int_0^t \int G_{t-s}(x-y)^2 dy ds\) 在 \(d=1\) 时收敛，在 \(d=2\) 时对数发散，在 \(d \geq 3\) 时幂次发散。

\textbf{证明}（\(d=1\) 情形）：记 \(\Gamma\) 为温和解算子。定义 Picard 迭代 \(u^{(n+1)} = \Gamma(u^{(n)})\)。由 Walsh 等距和 BDG 不等式，\(\mathbb{E}[\|u^{(n+1)} - u^{(n)}\|_\infty^2] \leq C \int_0^t \mathbb{E}[\|u^{(n)} - u^{(n-1)}\|_\infty^2] (t-s)^{-1/2} ds\)。由 Gronwall 引理（带奇异核的版本），序列在 \(L^2\) 中收敛。\(\blacksquare\)

\textbf{定义 136.4}（随机波动方程）：随机波动方程描述随机介质中的波传播： \[\frac{\partial^2 u}{\partial t^2} = \Delta u + \sigma(u) \dot{W}(t, x)\] 在 \(d=1\) 时，其温和解公式涉及波传播子（d'Alembert 公式的随机版本）。在 \(d=2\) 时，解需要分布意义上的解释。在 \(d=3\) 时，类似热方程，需要重正化。

\textbf{定理 136.3}（随机波动方程的能量估计）：在 \(d=1\) 时，随机波动方程的解满足能量不等式： \[\mathbb{E}\left[\|u_t(t,\cdot)\|_{L^2}^2 + \|\nabla u(t,\cdot)\|_{L^2}^2\right] \leq \mathbb{E}\left[\|u_1(0,\cdot)\|_{L^2}^2 + \|\nabla u_0\|_{L^2}^2\right] + C \int_0^t \mathbb{E}[\|\sigma(u(s,\cdot))\|_{L^2}^2] ds\]

\textbf{证明}：设 \(u\) 为随机波动方程 \(\partial_t^2 u = \Delta u + \sigma(u) \dot{W}\) 的温和解。定义能量泛函 \(\mathcal{E}(t) = \frac{1}{2} \mathbb{E}[\|u_t\|_{L^2}^2 + \|\nabla u\|_{L^2}^2]\)。由 Ito 公式作用于 \(\|u_t\|_{L^2}^2\) 和 \(\|\nabla u\|_{L^2}^2\)：\(d\|u_t\|_{L^2}^2 = 2\langle u_t, \Delta u + \sigma(u)\dot{W}\rangle dt + \|\sigma(u)\|_{L^2}^2 dt\)，\(d\|\nabla u\|_{L^2}^2 = -2\langle \nabla u_t, \nabla u\rangle dt\)。相加后 \(\langle u_t, \Delta u\rangle\) 与 \(-\langle \nabla u_t, \nabla u\rangle\) 抵消（分部积分），得 \(d\mathcal{E}(t) = \frac{1}{2}\mathbb{E}[\|\sigma(u)\|_{L^2}^2] dt + \text{鞅}\)。积分得 \(\mathcal{E}(t) \leq \mathcal{E}(0) + \frac{1}{2}\int_0^t \mathbb{E}[\|\sigma(u(s))\|_{L^2}^2] ds\)。\(\blacksquare\)

\subsection{136.3 KPZ 方程与正则性结构}\label{kpz-ux65b9ux7a0bux4e0eux6b63ux5219ux6027ux7ed3ux6784}

\textbf{定义 136.5}（KPZ 方程，Kardar-Parisi-Zhang 1986）：一维随机界面增长的 KPZ 方程为： \[\frac{\partial h}{\partial t} = \frac{\partial^2 h}{\partial x^2} + \lambda \left(\frac{\partial h}{\partial x}\right)^2 + \dot{W}(t, x)\] 其中 \(h(t,x)\) 表示界面高度。KPZ 方程是统计物理中随机界面增长的普适类模型，其非线性项 \((\partial_x h)^2\) 与白噪声 \(\dot{W}\) 的组合导致解在经典意义下不定义------\(h\) 是分布的（Hölder 正则性 \(1/2-\)），而 \((\partial_x h)^2\) 需要重正化。

\textbf{定义 136.6}（Hopf-Cole 变换）：当 \(\lambda \neq 0\) 时，KPZ 方程可通过 Hopf-Cole 变换 \(Z = \exp(\lambda h)\) 转化为随机热方程（乘性噪声）： \[\frac{\partial Z}{\partial t} = \frac{\partial^2 Z}{\partial x^2} + \lambda Z \dot{W}\] 这是严格处理 KPZ 方程的关键方法------\(Z > 0\) 且满足线性（乘性噪声）SPDE，其解可通过 Wiener 混沌展开构造。

\textbf{定理 136.4}（Hairer 的正则性结构，2014）：正则性结构（Regularity Structures）是处理奇异 SPDE 的一般框架，将 B 级数（Butcher 级数）和粗糙路径理论推广到 SPDE 的设定。该框架成功解决了 KPZ 方程、\(\Phi^4_3\) 模型和随机量化方程的存在唯一性。

\textbf{证明（框架）}：核心思想是构建一个代数结构（正则性结构），其中包含描述方程局部行为的''模型''（\((\Pi, \Gamma)\)），然后在适当正则性空间中建立不动点。\(\Phi^4_3\) 模型：\(\partial_t \phi = \Delta \phi - \phi^3 + \xi\)（\(\xi\) 为三维时空白噪声）。Hairer 的框架通过重正化（Wick 正规序）吸收发散，定义了方程的局部解，并证明了短时间内的存在唯一性。\(\blacksquare\)

\textbf{定理 136.5}（KPZ 普适类）：KPZ 方程描述了一大类随机界面增长模型的宏观行为（标度极限），包括：随机沉积模型、Eden 模型、定向聚合物模型、随机矩阵的最大特征值（Tracy-Widom 分布）。KPZ 方程的涨落指数为 \(1/3\)（空间相关长度指数 \(2/3\)），这一普适类的重要预测包括：一维 KPZ 方程的高度涨落服从 Tracy-Widom GOE 分布（而非 Gaussian 分布），这是随机矩阵理论在统计物理中的核心应用。KPZ 方程在 \(t \to \infty\) 时，\(h(t,x) \sim t^{1/3}\)（高度增长）且空间相关长度 \(\sim t^{2/3}\)，这些标度指数是 KPZ 普适类的标志性特征。

\textbf{证明}（框架）：通过 Hopf-Cole 变换 \(Z = \exp(\lambda h)\)，KPZ 方程转化为随机热方程 \(\partial_t Z = \partial_x^2 Z + \lambda Z \dot{W}\)。利用 Feynman-Kac 公式，\(Z(t,x)\) 的矩可表示为 Wiener 路径积分。在 \(t \to \infty\) 极限下，利用随机热方程的长时间行为（由 KPZ 不动点------Airy 过程刻画），Tracy-Widom 分布作为 GUE 最大特征值的分布自然出现。标度指数 \(1/3\) 和 \(2/3\) 由 KPZ 标度关系 \(h \sim t^{1/3} F(x/t^{2/3})\) 确定，其中 \(F\) 为普适标度函数。这些结论由 Corwin-Quastel-Sasamoto 等人通过可积概率方法严格证明。\(\blacksquare\)

\subsection{136.4 Schramm-Loewner 演化（SLE）}\label{schramm-loewner-ux6f14ux5316sle}

\textbf{定义 136.7}（SLE）：SLE（Schramm-Loewner Evolution，2000 年由 Oded Schramm 引入）是复平面上的随机共形不变曲线，由驱动函数为 \(\sqrt{\kappa} B_t\)（\(B_t\) 为一维 Brown 运动）的 Loewner 微分方程定义： \[\frac{d}{dt} g_t(z) = \frac{2}{g_t(z) - \sqrt{\kappa} B_t}, \quad g_0(z) = z\] SLE\(_\kappa\) 是连接上半平面 \(0\) 与 \(\infty\) 的随机曲线，其几何性质由参数 \(\kappa > 0\) 决定。

\textbf{定理 136.6}（SLE 的相变）：SLE\(_\kappa\) 的几何性质随 \(\kappa\) 发生相变： - \(\kappa \in (0, 4]\)：曲线是简单的（无自交），a.s. 不触及实轴 - \(\kappa \in (4, 8)\)：曲线有自交但不填充空间 - \(\kappa \geq 8\)：曲线是空间填充的（填充整个上半平面）

\textbf{证明}（框架）：SLE 由 Loewner 方程 \(\partial_t g_t(z) = 2/(g_t(z) - \sqrt{\kappa} B_t)\) 定义。相变的分界点由 Bessel 过程的性质决定：设 \(Z_t = g_t(z) - \sqrt{\kappa} B_t\)，则 \(Z_t\) 满足 \(dZ_t = 2/Z_t dt - \sqrt{\kappa} dB_t\)。当 \(\kappa \leq 4\) 时，\(Z_t\) 的 Bessel 维数 \(\delta = 1 + 4/\kappa \geq 2\)，故 \(Z_t\) 永不触及 \(0\)（由 Bessel 过程的性质），曲线不自交。当 \(\kappa > 4\) 时，Bessel 维数 \(\delta < 2\)，\(Z_t\) 以正概率触及 \(0\)，曲线自交。当 \(\kappa \geq 8\) 时，由 Beffara 的 Hausdorff 维数公式 \(\min(1+\kappa/8, 2) = 2\)，曲线空间填充。\(\blacksquare\)

\textbf{定理 136.7}（SLE 与统计物理模型的标度极限）：SLE 是许多二维统计物理模型在临界点处的标度极限： - \(\kappa = 2\)：Loop-erased random walk（LERW，去圈随机游走）的标度极限 - \(\kappa = 8/3\)：自回避随机游走（SAW）的标度极限（Lawler-Schramm-Werner 2004 年将此猜想证明为定理，并因此获得 Fields 奖） - \(\kappa = 3\)：临界 Ising 模型的界面 - \(\kappa = 4\)：Gaussian 自由场的等高线（调和探险者） - \(\kappa = 6\)：临界渗流的界面（Smirnov 2001 年证明了三角格上临界渗流界面的共形不变性，获 Fields 奖） - \(\kappa = 8\)：均匀生成树的 Peano 曲线

\textbf{定理 136.8}（SLE 的共形不变性与域 Markov 性）：SLE 是满足共形不变性和域 Markov 性（domain Markov property）的唯一随机曲线。共形不变性：SLE 在共形映射下的像仍是 SLE（同 \(\kappa\)）。域 Markov 性：给定曲线初始段 \(K_t\)，剩余部分的条件分布与 SLE 在 \(\mathbb{H} \setminus K_t\) 中相同（模共形等价）。这两个性质唯一确定了 SLE 作为随机 Loewner 演化。

\textbf{定理 136.9}（SLE 的 Hausdorff 维数）：SLE\(_\kappa\) 曲线的 Hausdorff 维数为 \(\min(1 + \kappa/8, 2)\)（Beffara 2008 年）。这一公式通过 SLE 的 Green 函数和势论估计证明，是 SLE 的几何性质与随机分析深度结合的结果。

\textbf{定理 136.10}（SLE 与随机矩阵------Cardy 公式）：SLE 在临界渗流中的应用由 Cardy 公式（预测渗流穿越概率的共形不变性，由 Smirnov 2001 年证明）和 Schramm 公式（预测 SLE 的穿越概率）联系。SLE 的穿越概率公式为： \[\mathbb{P}(\text{SLE}_\kappa \text{ passes left of } z) = \frac{1}{2} + \frac{\Gamma(4/\kappa)}{\sqrt{\pi}\Gamma((8-\kappa)/(2\kappa))} \cdot {}_2F_1(\cdots)\]

\hrulefill

\section{第359章：Levy 过程、大偏差与测度值过程}\label{ux7b2c359ux7ae0levy-ux8fc7ux7a0bux5927ux504fux5deeux4e0eux6d4bux5ea6ux503cux8fc7ux7a0b}

Levy 过程是具有独立平稳增量的随机过程，是 Brown 运动的自然推广，允许跳跃（非连续路径）。Levy 过程的特征由 Levy-Khintchine 公式完全刻画，将过程分解为确定的漂移项、Brown 扩散项和纯跳跃项（由 Levy 测度控制）。在金融数学中，Levy 过程用于建模资产价格的跳跃（如 Merton 跳跃扩散模型、Variance Gamma 模型），在保险数学中用于建模索赔过程（Cramer-Lundberg 模型），在物理学中用于建模反常扩散和湍流。大偏差理论（Large Deviations Theory）研究罕见事件的概率衰减速率，其数学基础是 Cramer 定理（1938 年）和 Varadhan 的 Laplace 渐进方法（1966 年）。大偏差理论在统计力学（熵与自由能的关系）、信息论（信道编码定理的精细渐近）和金融数学（风险度量的精确尾估计）中有核心应用。测度值过程（如超过程、Dawson-Watanabe 超 Brown 运动）描述随机演化的种群或粒子系统，是随机偏微分方程和分枝过程理论的交叉点。

\subsection{137.1 Levy 过程与 Levy-Khintchine 公式}\label{levy-ux8fc7ux7a0bux4e0e-levy-khintchine-ux516cux5f0f}

\textbf{定义 137.1}（Levy 过程）：随机过程 \((X_t)_{t \geq 0}\) 是 \textbf{Levy 过程}，如果 \(X_0 = 0\) a.s.，\(X_t\) 具有独立平稳增量（即对 \(0 \leq s < t\)，\(X_t - X_s\) 独立于 \(\mathcal{F}_s\) 且分布仅依赖于 \(t-s\)），且 \(X_t\) 是随机连续的（\(\lim_{s \to t} \mathbb{P}(|X_t - X_s| > \varepsilon) = 0\)）。Levy 过程自动具有 cadlag 版本（右连左极），且是强 Markov 过程。

\textbf{定理 137.1}（Levy-Khintchine 公式）：Levy 过程 \((X_t)\) 的特征函数为 \(\mathbb{E}[e^{i u X_t}] = e^{t \psi(u)}\)，其中特征指数 \(\psi(u)\) 为

\[\psi(u) = i b u - \frac{1}{2} \sigma^2 u^2 + \int_{\mathbb{R}} (e^{i u x} - 1 - i u x 1_{|x| < 1}) \nu(dx)\]

三元组 \((b, \sigma^2, \nu)\) 唯一确定 \(X_t\)，其中 \(b \in \mathbb{R}\) 为漂移系数，\(\sigma^2 \geq 0\) 为 Gaussian 系数（扩散项），\(\nu\) 是 \(\mathbb{R} \setminus \{0\}\) 上的 Levy 测度，满足 \(\int_{\mathbb{R}} (1 \wedge x^2) \nu(dx) < \infty\)（控制小跳跃的总变差）。Levy 测度 \(\nu(A)\) 是单位时间内跳跃大小落在 \(A\) 中的期望次数。跳跃过程 \(J_t = \sum_{s \leq t} \Delta X_s\) 的补偿测度为 \(\nu(dx) dt\)。

\textbf{证明}：对 \(X_t\) 用无穷可分分布的 Levy-Khintchine 表示。由独立平稳增量，\(X_t\) 是无穷可分分布。Levy 测度 \(\nu\) 控制跳跃的大小和频率，截断函数 \(1_{|x| < 1}\) 确保积分收敛（吸收小跳跃的一阶矩效应）。\(\blacksquare\)

\textbf{定理 137.2}（Levy-Ito 分解）：任何 Levy 过程可分解为三个独立部分之和： \[X_t = b t + \sigma W_t + \int_0^t \int_{|x| \geq 1} x N(ds, dx) + \int_0^t \int_{|x| < 1} x (N(ds, dx) - \nu(dx) ds)\]

其中 \(N\) 为 Poisson 随机测度（强度 \(\nu(dx) dt\)）。第一项为确定性漂移，第二项为 Brown 扩散，第三项为大幅跳跃之和（复合 Poisson 过程），第四项为补偿小跳跃的纯断鞅（compensated small jumps martingale）。这四项之和给出 Levy 过程的完整 Levy-Ito 分解。

\textbf{证明}：由 Levy-Khintchine 公式，Levy 过程的特征函数由 \((b, \sigma^2, \nu)\) 唯一确定。将 Levy 测度 \(\nu\) 分解为 \(\nu = \nu|_{|x| < 1} + \nu|_{|x| \geq 1}\)。小跳部分 \(\nu|_{|x| < 1}\) 满足 \(\int_{|x| < 1} |x|^2 \nu(dx) < \infty\)（由 Levy 测度的定义），故可构造对应的平方可积鞅------即补偿 Poisson 随机积分 \(\int_0^t \int_{|x| < 1} x (N(ds, dx) - \nu(dx) ds)\)。大跳部分 \(\nu|_{|x| \geq 1}\) 是有限测度，故 \(\int_0^t \int_{|x| \geq 1} x N(ds, dx)\) 是复合 Poisson 过程。Brown 运动部分 \(W_t\) 由特征函数中 \(\exp(-t \sigma^2 u^2/2)\) 项对应，漂移 \(bt\) 由 \(\exp(i t b u)\) 项对应。各部分的独立性由 Levy 过程的独立增量性及其特征函数的乘积形式直接得出。\(\blacksquare\)

\textbf{定义 137.2}（重要的 Levy 过程示例）： - \textbf{Brown 运动}：\(\nu = 0\)，\(\sigma^2 = 1\)，\(b = 0\)（连续路径） - \textbf{Poisson 过程}：\(\nu = \lambda \delta_1\)，\(\sigma^2 = 0\)，\(b = 0\)（单位跳跃） - \textbf{复合 Poisson 过程}：\(\nu = \lambda F\)（\(F\) 为跳跃分布），\(\sigma^2 = 0\) - \textbf{\(\alpha\)-稳定过程}：\(\nu(dx) = c_+ x^{-1-\alpha} 1_{x>0} dx + c_- |x|^{-1-\alpha} 1_{x<0} dx\)（\(\alpha \in (0, 2)\)），\(\sigma^2 = 0\)（\(\alpha = 2\) 时为 Brown 运动） - \textbf{Gamma 过程}：\(\nu(dx) = a x^{-1} e^{-bx} 1_{x>0} dx\)（纯跳跃递增过程） - \textbf{Variance Gamma 过程}：Brown 运动通过 Gamma 时间变换（\(X_t = W_{G_t}\)，\(G_t\) 为 Gamma 过程）

\textbf{定理 137.3}（\(\alpha\)-稳定过程的自相似性）：若 \(X_t\) 是 \(\alpha\)-稳定过程，则 \(X_t\) 满足 \(1/\alpha\)-自相似性：\((X_{ct}) \stackrel{d}{=} c^{1/\alpha} X_t\)。\(\alpha = 2\) 时恢复 Brown 运动的 \(1/2\)-自相似性。\(\alpha\)-稳定过程的尾部分布为 \(P(X_t > x) \sim C x^{-\alpha}\)（\(x \to \infty\)），具有重尾（Pareto 尾），因此方差无限（\(\alpha < 2\) 时）。

\textbf{定理 137.4}（Levy 过程的反射原理与 Wiener-Hopf 分解）：对任意 Levy 过程 \(X_t\)，其跑动上确界 \(S_t = \sup_{0 \leq s \leq t} X_s\) 和跑动下确界 \(I_t = \inf_{0 \leq s \leq t} X_s\) 满足 Wiener-Hopf 分解：对独立的指数随机变量 \(T \sim \exp(q)\)，\(S_T\) 和 \(X_T - S_T\) 独立。Wiener-Hopf 分解是研究 Levy 过程的首达时、溢出分布和波动理论的基石。

\subsection{137.2 大偏差理论}\label{ux5927ux504fux5deeux7406ux8bba}

\textbf{定义 137.3}（大偏差原则 / LDP）：一族概率测度 \(\{\mu_n\}\) 满足\textbf{大偏差原则}（Large Deviation Principle），速率函数为 \(I: \mathcal{X} \to [0, \infty]\)（下半连续），如果对任意 Borel 集 \(A\)：

\[-\inf_{x \in A^\circ} I(x) \leq \liminf_{n \to \infty} \frac{1}{n} \log \mu_n(A) \leq \limsup_{n \to \infty} \frac{1}{n} \log \mu_n(A) \leq -\inf_{x \in \overline{A}} I(x)\]

速率函数 \(I(x)\) 衡量了事件 \(x\) 的''稀有程度''：\(I(x)\) 越大，\(x\) 的概率越小。LDP 的直观含义是 \(\mu_n(A) \approx \exp(-n \inf_{x \in A} I(x))\)，即罕见事件的概率以指数速率衰减，衰减速率由速率函数的最小值决定。

\textbf{定理 137.5}（Cramer 定理，1938）：设 \(X_1, X_2, \ldots\) 是 i.i.d. 随机变量，\(\overline{X}_n = \frac{1}{n} \sum_{i=1}^n X_i\)。则 \(\overline{X}_n\) 满足 LDP，速率函数为对数矩生成函数 \(M(\theta) = \log \mathbb{E}[e^{\theta X_1}]\) 的 Legendre-Fenchel 变换：

\[I(x) = \sup_{\theta \in \mathbb{R}} (\theta x - M(\theta))\]

\textbf{证明}：上界：由 Chernoff 界，\(\mathbb{P}(\overline{X}_n \geq x) = \mathbb{P}(e^{n\theta \overline{X}_n} \geq e^{n\theta x}) \leq e^{-n\theta x} \mathbb{E}[e^{\theta \sum X_i}] = \exp(-n(\theta x - M(\theta)))\)。对 \(\theta\) 取上确界。下界：利用指数倾斜（exponential tilting）构造新测度 \(\tilde{\mathbb{P}}(A) = \mathbb{E}[e^{\theta \sum X_i - n M(\theta)} 1_A]\)，在此测度下 \(\overline{X}_n\) 的均值偏至 \(x\)。由大数定律，\(\tilde{\mathbb{P}}(|\overline{X}_n - x| < \varepsilon) \to 1\)，转换为原测度下的概率下界。\(\blacksquare\)

\textbf{定理 137.6}（Sanov 定理，1957）：设 \(X_1, X_2, \ldots\) 是 i.i.d. 取值于 Polish 空间 \(\mathcal{X}\) 的随机变量，经验分布 \(L_n = \frac{1}{n} \sum_{i=1}^n \delta_{X_i}\)。则 \(L_n\) 满足 LDP，速率函数为相对熵（Kullback-Leibler 散度）：

\[I(\mu) = H(\mu \mid P) = \int \log\left(\frac{d\mu}{dP}\right) d\mu\]

其中 \(P\) 为 \(X_i\) 的分布。Sanov 定理是大偏差理论在统计学和信息论中最核心的应用------经验分布的罕见偏离概率由相对熵指数衰减。

\textbf{证明}：将经验分布视为有限维分布的投影，利用 Cramer 定理在有限维 Simplex 上的极限和投影极限定理。或利用组合方法：\(\mathbb{P}(L_n \approx \mu) \approx \frac{n!}{\prod_i (n\mu_i)!} \prod_i P_i^{n\mu_i} \approx \exp(-n H(\mu \mid P))\)（Stirling 近似）。\(\blacksquare\)

\textbf{定理 137.7}（Varadhan 积分引理）：若 \(\{\mu_n\}\) 满足 LDP，速率函数 \(I\)，且 \(F\) 是连续有界函数，则

\[\lim_{n \to \infty} \frac{1}{n} \log \int e^{n F(x)} \mu_n(dx) = \sup_{x} (F(x) - I(x))\]

\textbf{证明}：上界：分割积分域为 \(\varepsilon\)-网，\(\int e^{nF} d\mu_n \leq \sum_i \sup_{B_i} e^{nF} \mu_n(B_i) \approx \sum_i \exp(n(\sup_{B_i} F - \inf_{B_i} I))\)。下界：对固定 \(x\)，\(\int e^{nF} d\mu_n \geq e^{nF(x)} \mu_n(B_\varepsilon(x)) \approx \exp(n(F(x) - I(x)))\)。取 \(\varepsilon \to 0\) 和 \(x\) 的上确界。\(\blacksquare\)

Varadhan 引理是大偏差理论的 Laplace 方法，它将大偏差原理与统计力学中的 Gibbs 变分原理统一起来：\(\sup_x (F(x) - I(x))\) 对应于自由能（压力函数）的 Legendre 变换。

\textbf{定理 137.8}（Schilder 定理------Brown 运动的大偏差）：Brown 运动 \(W_t\) 的缩放 \((\varepsilon W_t)_{t \in [0,T]}\) 在 \(C[0,T]\) 上满足 LDP（\(\varepsilon \to 0\)），速率函数为

\[I(\phi) = \frac{1}{2} \int_0^T \dot{\phi}(t)^2 dt\]

若 \(\phi \in H^1\)（\(\phi(0)=0\)，绝对连续且导数平方可积）；否则 \(I(\phi) = \infty\)。Schilder 定理是 Freidlin-Wentzell 理论的基础，用于研究小噪声 SDE 的大偏差行为。

\textbf{定理 137.9}（Freidlin-Wentzell 大偏差）：对 SDE \(dX_t^\varepsilon = b(X_t^\varepsilon) dt + \varepsilon \sigma(X_t^\varepsilon) dW_t\)，\(X^\varepsilon\) 在 \(C[0,T]\) 上满足 LDP，速率函数为

\[I(\phi) = \frac{1}{2} \int_0^T \|\dot{\phi}(t) - b(\phi(t))\|_{\sigma}^2 dt\]

其中 \(\|v\|_\sigma^2 = v^T (\sigma\sigma^T)^{-1} v\)。这一理论在随机动力系统（如噪声诱导的跃迁、亚稳态逃逸问题）中有核心应用。

\subsection{137.3 测度值过程}\label{ux6d4bux5ea6ux503cux8fc7ux7a0b}

\textbf{定义 137.4}（超过程 / Superprocess）：超过程是取值为测度的 Markov 过程，描述随机演化的连续型分枝过程。Dawson-Watanabe 超 Brown 运动是 \(\mathbb{R}^d\) 上的测度值过程，其 Laplace 泛函满足： \[\mathbb{E}\left[\exp\left(-\int_{\mathbb{R}^d} \phi(x) X_t(dx)\right)\right] = \exp\left(-\int_{\mathbb{R}^d} v(t,x) \mu_0(dx)\right)\] 其中 \(v\) 满足反应扩散方程 \(\frac{\partial v}{\partial t} = \Delta v - v^2\)（\(v(0,x) = \phi(x)\)）。

\textbf{定理 137.10}（超 Brown 运动的支撑性质）：\(d=1\) 时，超 Brown 运动的支撑是绝对连续的（有密度）。\(d \geq 2\) 时，超 Brown 运动的支撑是奇异测度（Hausdorff 维数 \(4-d\) 的随机分形），具有间歇性（intermittency）------在空间上聚集成团。这一维度依赖性是分枝扩散过程在连续极限下的典型行为。

\textbf{定义 137.5}（Fleming-Viot 过程）：Fleming-Viot 过程是描述基因频率随机演化的测度值过程，其生成元为 \[\mathcal{L} F(\mu) = \frac{1}{2} \int \int \mu(dx) (\delta_x(dy) - \mu(dy)) \frac{\delta^2 F}{\delta \mu(x) \delta \mu(y)} + \int \mathcal{L}_0 \frac{\delta F}{\delta \mu(x)} \mu(dx)\] 其中 \(\mathcal{L}_0\) 是单个基因的扩散生成元。Fleming-Viot 过程是群体遗传学中 Wright-Fisher 扩散模型的连续极限。

\textbf{定理 137.11}（Dirichlet 形式与超过程）：超过程的构造可通过 Dirichlet 形式理论严格完成。超 Brown 运动的 Dirichlet 形式为 \(\mathcal{E}(F, F) = \frac{1}{2} \int \|\nabla \frac{\delta F}{\delta \mu}\|_{L^2(\mu)}^2 \mathbb{P}(d\mu)\)，其中 \(\mathbb{P}\) 是超过程的平稳测度（若存在）。Dirichlet 形式方法将超过程的分析与势论和谱理论联系起来。

\textbf{定义 137.6}（测度值分枝过程与 SPDE 的联系）：超过程是 SPDE 的''弱解''------超 Brown 运动的密度 \(u(t,x)\)（若存在）形式地满足 SPDE \(\partial_t u = \Delta u + \sqrt{u} \dot{W}\)。这一联系使得超过程成为研究非线性 SPDE 的概率工具，也是 Dawson-Gartner 投影定理中相互作用的 Fleming-Viot 系统与 McKean-Vlasov 方程之间的桥梁。

\textbf{定理 137.12}（大偏差与 Gibbs 条件）：对于测度值过程的经验测度 \(L_T = \frac{1}{T} \int_0^T \delta_{X_t} dt\)，大偏差原理的速率函数 \(I(\mu)\) 通常由 Dirichlet 形式或相对熵给出。这一结果将统计力学中的 Gibbs 条件（平衡态等价于自由能最小化）与随机过程的大偏差理论统一起来，是 Donsker-Varadhan 大偏差理论的核心。

\hrulefill

\newpage

\chapter{10-概率与统计/03-统计学/01-估计与检验.md}\label{ux6982ux7387ux4e0eux7edfux8ba103-ux7edfux8ba1ux5b6601-ux4f30ux8ba1ux4e0eux68c0ux9a8c.md}

\section{第360章：估计理论}\label{ux7b2c360ux7ae0ux4f30ux8ba1ux7406ux8bba}

估计理论是统计推断的基础，研究如何从样本中估计未知参数。本章讨论无偏估计、Fisher 信息、Cramér-Rao 下界和极大似然估计。

\subsection{109.1 统计模型与充分统计量}\label{ux7edfux8ba1ux6a21ux578bux4e0eux5145ux5206ux7edfux8ba1ux91cf}

\textbf{定义 109.1}（参数统计模型）：\textbf{参数统计模型}是一族概率分布 \(\{P_\theta : \theta \in \Theta\}\)，其中 \(\Theta \subseteq \mathbb{R}^k\) 是参数空间。样本 \(X_1, \ldots, X_n\) 是 i.i.d. 来自 \(P_\theta\)。

\textbf{定义 109.2}（统计量）：\textbf{统计量} \(T = T(X_1, \ldots, X_n)\) 是样本的可测函数（不依赖于未知参数 \(\theta\)）。

\textbf{定义 109.3}（充分统计量）：统计量 \(T\) 称为关于 \(\theta\) \textbf{充分}的，如果 \(X\) 在 \(T\) 给定下的条件分布不依赖于 \(\theta\)。即 \(T\) 包含了样本中关于 \(\theta\) 的所有信息。

\textbf{定理 109.1}（Fisher-Neyman 因子分解定理）：设 \(\{P_\theta\}\) 有密度（或概率质量函数）\(f(x \mid \theta)\)。\(T\) 是充分的当且仅当

\[f(x_1, \ldots, x_n \mid \theta) = g(T(x_1, \ldots, x_n), \theta) \cdot h(x_1, \ldots, x_n)\]

其中 \(g\) 依赖于 \(\theta\) 和 \(T\)，\(h\) 不依赖于 \(\theta\)。

\emph{证明}：离散情形。若 \(T\) 充分，定义 \(g(t, \theta) = P_\theta(T = t)\)，\(h(x) = P(X = x \mid T = t) / P_\theta(T = t)\)（后者不依赖于 \(\theta\)）。反之亦然。∎

\subsection{109.2 点估计}\label{ux70b9ux4f30ux8ba1}

点估计的目的是用一个统计量（即样本的函数）来给出未知参数的合理猜测。一个好的估计量需要在多个标准下进行评价------首先是其系统偏差（无偏性），其次是其与真实值的平均偏离程度（均方误差），以及在大样本下的一致性（相合性）。这些评价标准构成了经典估计理论的核心理念。

\textbf{定义 109.4}（估计量）：参数 \(\theta\)（或 \(g(\theta)\)）的\textbf{估计量} \(\hat{\theta}_n = \hat{\theta}_n(X_1, \ldots, X_n)\) 是样本的统计量。

无偏性要求估计量在重复抽样下的平均值恰好等于真值------这排除了系统性高估或低估。但无偏性本身不保证估计的精度：一个无偏估计可能有极大的方差。因此需要同时考虑偏差和方差，即均方误差。

\textbf{定义 109.5}（无偏性）：\(\hat{\theta}_n\) 称为 \textbf{无偏}的，如果 \(\mathbb{E}_\theta[\hat{\theta}_n] = \theta\)（对所有 \(\theta \in \Theta\)）。其偏差为 \(\operatorname{Bias}_\theta(\hat{\theta}_n) = \mathbb{E}_\theta[\hat{\theta}_n] - \theta\)。

\textbf{定义 109.6}（均方误差，MSE）：\(\operatorname{MSE}_\theta(\hat{\theta}_n) = \mathbb{E}_\theta[(\hat{\theta}_n - \theta)^2] = \operatorname{Var}_\theta(\hat{\theta}_n) + (\operatorname{Bias}_\theta(\hat{\theta}_n))^2\)。

MSE 在有限样本下给出了完整的质量度量，而相合性则保证估计量在大样本极限下收敛到真值------这是任何合理估计量应满足的最低标准。

\textbf{定义 109.7}（相合性）：\(\hat{\theta}_n\) 称为\textbf{相合}的（一致的），如果 \(\hat{\theta}_n \xrightarrow{\mathbb{P}_\theta} \theta\)（对所有 \(\theta\)）。强相合：\(\hat{\theta}_n \xrightarrow{\text{a.s.}} \theta\)。

\subsection{109.3 Fisher 信息与 Cramér-Rao 下界}\label{fisher-ux4fe1ux606fux4e0e-cramuxe9r-rao-ux4e0bux754c}

\textbf{定义 109.8}（得分函数与 Fisher 信息）：设模型 \(\{f(x \mid \theta) : \theta \in \Theta\}\)（\(\Theta \subseteq \mathbb{R}\)）满足正则条件。\textbf{得分函数}为 \(S(\theta, x) = \frac{\partial}{\partial \theta} \log f(x \mid \theta)\)。单个观测的 \textbf{Fisher 信息}为

\[I(\theta) = \mathbb{E}_\theta\left[\left(\frac{\partial}{\partial \theta} \log f(X \mid \theta)\right)^2\right] = -\mathbb{E}_\theta\left[\frac{\partial^2}{\partial \theta^2} \log f(X \mid \theta)\right]\]

\(n\) 个 i.i.d. 观测的 Fisher 信息为 \(n I(\theta)\)。

Fisher 信息量衡量了数据 \(X\) 携带的关于未知参数 \(\theta\) 的信息量------信息越大，参数越能被精确估计。这一直觉由 Cramér-Rao 下界给出精确形式：任何无偏估计量的方差不小于 Fisher 信息的倒数除以样本量。因此 Fisher 信息在所有正则参数模型中定义了估计精度的理论上界。

\textbf{定理 109.2}（Cramér-Rao 下界）：设 \(\hat{\theta}_n\) 是 \(g(\theta)\) 的无偏估计量。在正则条件下，

\[\operatorname{Var}_\theta(\hat{\theta}_n) \geq \frac{(g'(\theta))^2}{n I(\theta)}\]

特别地，对 \(\theta\) 的无偏估计，\(\operatorname{Var}_\theta(\hat{\theta}_n) \geq 1/(n I(\theta))\)。

\textbf{证明}：设 \(S(\theta, X) = \frac{\partial}{\partial\theta}\log f(X\mid\theta)\)，则 \(\mathbb{E}_\theta[S] = 0\)。总得分函数 \(S_n = \sum_{i=1}^n S(\theta, X_i)\)。由无偏性 \(\mathbb{E}_\theta[\hat{\theta}_n] = g(\theta)\)，在正则条件下求导得 \(\mathbb{E}_\theta[\hat{\theta}_n S_n] = g'(\theta)\)。由于 \(\mathbb{E}_\theta[S_n] = 0\)，有 \(\operatorname{Cov}_\theta(\hat{\theta}_n, S_n) = \mathbb{E}_\theta[\hat{\theta}_n S_n] - \mathbb{E}_\theta[\hat{\theta}_n]\mathbb{E}_\theta[S_n] = g'(\theta)\)。由 Cauchy-Schwarz 不等式，\(|\operatorname{Cov}_\theta(\hat{\theta}_n, S_n)|^2 \leq \operatorname{Var}_\theta(\hat{\theta}_n) \cdot \operatorname{Var}_\theta(S_n)\)。而 \(\operatorname{Var}_\theta(S_n) = n I(\theta)\)（i.i.d. 观测下 Fisher 信息可加），故 \(|g'(\theta)|^2 \leq \operatorname{Var}_\theta(\hat{\theta}_n) \cdot n I(\theta)\)，整理即得所证。等号成立当且仅当 \(\hat{\theta}_n\) 与 \(S_n\) 仿射相关，即分布属于指数族。\(\blacksquare\)

Cramér-Rao 下界给出了所有无偏估计量的方差下界，那么是否存在能达到该下界的估计量？这一问题将估计理论引向了对有效估计量的研究。

\textbf{定义 109.9}（有效估计量）：达到 Cramér-Rao 下界的无偏估计量称为\textbf{有效}的。有效估计量存在当且仅当分布属于指数族。大多数常用分布（正态、二项、Poisson、Gamma、Beta 等）属于指数族，但并非所有分布都如此（如 Cauchy 分布、均匀分布不属于指数族）。

\subsection{109.4 极大似然估计}\label{ux6781ux5927ux4f3cux7136ux4f30ux8ba1}

极大似然估计（MLE）是参数估计中最为通用的方法，其基本思想极为直观：选择使观察到实际数据的概率（似然）最大的参数值作为估计。Fisher 在 1920 年代确立了 MLE 的理论地位，证明了它在大样本下具有最优渐近性质------相合性、渐近正态性和渐近有效性。

\textbf{定义 109.10}（似然函数）：\textbf{似然函数} \(L(\theta \mid x) = f(x \mid \theta)\)（将 \(\theta\) 视为变量）。\textbf{对数似然} \(\ell(\theta \mid x) = \log L(\theta \mid x)\)。

\textbf{定义 109.11}（极大似然估计，MLE）：\textbf{极大似然估计量} \(\hat{\theta}_{\text{MLE}}\) 最大化 \(L(\theta \mid X)\)（或等价地 \(\ell(\theta \mid X)\)）：

\[\hat{\theta}_{\text{MLE}} = \arg\max_{\theta \in \Theta} L(\theta \mid X)\]

MLE 的显著性不在于小样本的有限性质，而在于其大样本渐近的最优性。以下定理是估计理论的基石：在正则条件下，MLE 是相合的，其分布渐近于正态，且渐近方差达到 Cramér-Rao 下界------即在所有相合渐近正态估计中，MLE 是渐近最优的。

\textbf{定理 109.3}（MLE 的渐近性质）：在正则条件下： 1. \textbf{相合性}：\(\hat{\theta}_{\text{MLE}} \xrightarrow{\mathbb{P}} \theta_0\)（真值） 2. \textbf{渐近正态性}：\(\sqrt{n}(\hat{\theta}_{\text{MLE}} - \theta_0) \xrightarrow{d} \mathcal{N}(0, I(\theta_0)^{-1})\) 3. \textbf{渐近有效性}：MLE 渐近达到 Cramér-Rao 下界

\textbf{证明}：（1）相合性：对数似然 \(\ell_n(\theta) = \sum_{i=1}^n \log f(X_i\mid\theta)\)。由大数定律，\(\frac{1}{n}\ell_n(\theta) \to \mathbb{E}_{\theta_0}[\log f(X\mid\theta)] = -D_{\text{KL}}(\theta_0\|\theta) - H(\theta_0)\)，取极大值于 \(\theta=\theta_0\)。由 argmax 连续定理，\(\hat{\theta}_{\text{MLE}} \xrightarrow{\mathbb{P}} \theta_0\)。（2）渐近正态性：由 \(\ell_n'(\hat{\theta})=0\)，Taylor 展开：\(0 = \ell_n'(\theta_0) + \ell_n''(\tilde{\theta})(\hat{\theta}-\theta_0)\)，故 \(\sqrt{n}(\hat{\theta}-\theta_0) = -\frac{\ell_n'(\theta_0)/\sqrt{n}}{\ell_n''(\tilde{\theta})/n}\)。得分函数 \(\ell_n'(\theta_0) = \sum S(\theta_0,X_i)\) 满足 \(\mathbb{E}[S]=0\)，\(\operatorname{Var}(S)=I(\theta_0)\)，由 CLT 得 \(\ell_n'(\theta_0)/\sqrt{n} \xrightarrow{d} \mathcal{N}(0,I(\theta_0))\)。由大数定律，\(\ell_n''(\tilde{\theta})/n \xrightarrow{\mathbb{P}} -I(\theta_0)\)。由 Slutsky 定理，\(\sqrt{n}(\hat{\theta}-\theta_0) \xrightarrow{d} \mathcal{N}(0, I(\theta_0)^{-1})\)。（3）渐近有效性：渐近方差恰为 Cramér-Rao 下界的倒数。\(\blacksquare\)

MLE 还有一个突出的性质------参数变换下的不变性。若极大化 \(\theta\) 的似然函数，则对于任意函数 \(g\)，\(g(\hat{\theta}_{\text{MLE}})\) 自然地成为 \(g(\theta)\) 的 MLE。这一性质在 MLE 的具体计算中极为方便，尤其是当参数的重参数化简化了似然函数的结构时。

\textbf{定理 109.4}（MLE 的不变性）：若 \(\hat{\theta}\) 是 \(\theta\) 的 MLE，\(g\) 是任意函数，则 \(g(\hat{\theta})\) 是 \(g(\theta)\) 的 MLE。

\textbf{证明}：MLE \(\hat{\theta}\) 最大化 \(L(\theta\mid X)\) 于 \(\Theta\)。\(g(\theta)\) 的诱导似然定义为 \(L^*(g_0\mid X) = \sup_{\theta: g(\theta)=g_0} L(\theta\mid X)\)。\(\hat{g} = g(\hat{\theta})\) 满足 \(L^*(\hat{g}\mid X) \geq L(\hat{\theta}\mid X)\)（取 \(\theta = \hat{\theta}\)）。另一方面，对任意 \(g_0\)，若 \(\theta_0\) 使 \(g(\theta_0)=g_0\)，则 \(L(\theta_0\mid X) \leq L(\hat{\theta}\mid X) \leq L^*(\hat{g}\mid X)\)，取上确界得 \(L^*(g_0\mid X) \leq L^*(\hat{g}\mid X)\)。故 \(\hat{g}\) 最大化诱导似然，即 \(g(\hat{\theta})\) 是 \(g(\theta)\) 的 MLE。当 \(g\) 是一一映射时尤为显然：似然在参数变换下 \(L(\theta\mid X) = L(g^{-1}(\eta)\mid X)\)，\(\hat{\eta} = g(\hat{\theta})\) 自然最大化。\(\blacksquare\)

\hrulefill

\hrulefill

\hrulefill

\hrulefill

\hrulefill

\hrulefill

\section{第361章：假设检验}\label{ux7b2c361ux7ae0ux5047ux8bbeux68c0ux9a8c}

假设检验是统计推断的另一个核心，用于判断数据是否支持某一假设。本章介绍 Neyman-Pearson 引理和一致最优势检验。

\subsection{110.1 假设检验的基本概念}\label{ux5047ux8bbeux68c0ux9a8cux7684ux57faux672cux6982ux5ff5}

\textbf{定义 110.1}（假设检验）：\textbf{假设检验}问题中，\textbf{原假设} \(H_0: \theta \in \Theta_0\) 与\textbf{备择假设} \(H_1: \theta \in \Theta_1\)（\(\Theta_0 \cap \Theta_1 = \varnothing\)）。

\textbf{定义 110.2}（检验函数与拒绝域）：\textbf{检验函数} \(\varphi(x) \in [0, 1]\) 是拒绝 \(H_0\) 的概率（给定数据 \(x\)）。\textbf{拒绝域} \(R = \{x : \varphi(x) = 1\}\)。若 \(\varphi(x) \in \{0, 1\}\)，称为\textbf{非随机化检验}。

\textbf{定义 110.3}（两类错误）： - \textbf{第一类错误}（假阳性）：\(H_0\) 为真但被拒绝。概率 \(\alpha(\theta) = \mathbb{E}_\theta[\varphi(X)]\)（\(\theta \in \Theta_0\)） - \textbf{第二类错误}（假阴性）：\(H_1\) 为真但未被拒绝。概率 \(\beta(\theta) = 1 - \mathbb{E}_\theta[\varphi(X)]\)（\(\theta \in \Theta_1\)）

\textbf{定义 110.4}（检验水平与功效）：检验 \(\varphi\) 的\textbf{水平}（size）为 \(\alpha = \sup_{\theta \in \Theta_0} \mathbb{E}_\theta[\varphi(X)]\)。\textbf{功效函数}（power）为 \(\pi(\theta) = \mathbb{E}_\theta[\varphi(X)]\)（\(\theta \in \Theta_1\)）。

\subsection{110.2 Neyman-Pearson 引理}\label{neyman-pearson-ux5f15ux7406}

\textbf{定理 110.1}（Neyman-Pearson 引理）：考虑简单假设检验 \(H_0: \theta = \theta_0\) vs \(H_1: \theta = \theta_1\)。设 \(f_0, f_1\) 分别是 \(H_0\) 和 \(H_1\) 下的密度。则水平 \(\alpha\) 的\textbf{最优势检验}（最大功效）由似然比检验给出：

\[\varphi(x) = \begin{cases} 1 & \text{若 } \frac{f_1(x)}{f_0(x)} > c \\ \gamma & \text{若 } \frac{f_1(x)}{f_0(x)} = c \\ 0 & \text{若 } \frac{f_1(x)}{f_0(x)} < c \end{cases}\]

其中 \(c\) 和 \(\gamma\) 选择使得 \(\mathbb{E}_{\theta_0}[\varphi(X)] = \alpha\)。

\emph{证明}：设 \(\varphi^*\) 是 Neyman-Pearson 检验，\(\varphi\) 是任意其他水平 \(\leq \alpha\) 的检验。则 \((\varphi^* - \varphi)(f_1 - c f_0) \geq 0\) 逐点。积分得 \(\mathbb{E}_{\theta_1}[\varphi^*] - \mathbb{E}_{\theta_1}[\varphi] \geq 0\)。∎

\textbf{推论 110.2}（单调似然比）：若 \(\{f_\theta\}\) 关于统计量 \(T(x)\) 有单调似然比（即 \(f_{\theta_1}(x)/f_{\theta_0}(x)\) 是 \(T(x)\) 的非降函数，\(\theta_1 > \theta_0\)），则对 \(H_0: \theta \leq \theta_0\) vs \(H_1: \theta > \theta_0\)，存在一致最优势（UMP）检验，拒绝域为 \(T(x) > c\)。

\textbf{证明}：由单调似然比性质，对任意 \(\theta_1 > \theta_0\)，似然比 \(f_{\theta_1}(x)/f_{\theta_0}(x)\) 是 \(T(x)\) 的非降函数，故 \(\{x : f_{\theta_1}(x)/f_{\theta_0}(x) > k\}\) 等价于 \(\{x : T(x) > c\}\) 对某 \(c\)。由 Neyman-Pearson 引理，对简单假设 \(H_0: \theta = \theta_0\) vs \(H_1: \theta = \theta_1\)，该形式的检验是 MP 的。由于 \(c\) 的选择仅依赖于 \(\alpha\)（不依赖于 \(\theta_1\)），对任意 \(\theta_1 > \theta_0\) 该检验都是 MP 的。再对 \(\theta \leq \theta_0\) 验证水平：由单调似然比，功效函数 \(\pi(\theta) = \mathbb{E}_\theta[\varphi(X)]\) 是 \(\theta\) 的非降函数，故 \(\sup_{\theta \leq \theta_0} \pi(\theta) = \pi(\theta_0) = \alpha\)。因此该检验是 UMP 的。\(\blacksquare\)

\subsection{110.3 似然比检验}\label{ux4f3cux7136ux6bd4ux68c0ux9a8c}

\textbf{定义 110.5}（广义似然比检验，GLRT）：对复合假设 \(H_0: \theta \in \Theta_0\) vs \(H_1: \theta \in \Theta \setminus \Theta_0\)，\textbf{似然比统计量}为

\[\Lambda(X) = \frac{\sup_{\theta \in \Theta_0} L(\theta \mid X)}{\sup_{\theta \in \Theta} L(\theta \mid X)}\]

拒绝域为 \(\Lambda(X) \leq c\)（\(0 \leq c \leq 1\)）。

广义似然比检验的直观思想是：计算约束模型（原假设下）与无约束全模型的最大似然之比，比值越小说明原假设对数据的拟合越差，越应被拒绝。这个看似特设的统计量由于 Wilks 定理而获得了牢固的理论基础------在 \(H_0\) 下，\(-2\log \Lambda\) 渐近服从卡方分布，自由度为被约束的参数个数，从而为构造渐近 \(p\) 值提供了普遍适用的方法。

\textbf{定理 110.3}（Wilks 定理）：在 \(H_0\) 下（正则条件），\(-2 \log \Lambda(X) \xrightarrow{d} \chi^2_{\nu}\)，其中 \(\nu = \dim \Theta - \dim \Theta_0\)（自由度等于被检验的参数个数）。

\textbf{证明}：设 \(\hat{\theta}\) 是全模型 MLE，\(\hat{\theta}_0\) 是 \(H_0\) 约束下的 MLE。Taylor 展开对数似然于 \(\hat{\theta}\) 处：\(\ell(\hat{\theta}_0) \approx \ell(\hat{\theta}) + \frac{1}{2}(\hat{\theta}_0 - \hat{\theta})^\top \ell''(\hat{\theta})(\hat{\theta}_0 - \hat{\theta})\)，因 \(\ell'(\hat{\theta})=0\)。故 \(-2\log\Lambda = 2[\ell(\hat{\theta}) - \ell(\hat{\theta}_0)] \approx -(\hat{\theta}_0-\hat{\theta})^\top \ell''(\hat{\theta})(\hat{\theta}_0-\hat{\theta})\)。由 MLE 渐近正态性，\(\sqrt{n}(\hat{\theta}-\theta_0) \xrightarrow{d} \mathcal{N}(0, I^{-1})\) 且 \(-\ell''(\hat{\theta})/n \to I\)。在 \(H_0\) 下，约束 \(\theta \in \Theta_0\) 可局部参数化为 \(\nu\) 个自由参数，则 \(\hat{\theta}_0 - \hat{\theta}\) 的渐近分布是 \(\mathcal{N}(0, I^{-1} - P_{\Theta_0}I^{-1})\) 在某 \(\nu\) 维子空间上的投影。代入二次型即得 \(\chi^2_\nu\) 极限分布。严格化需用 Wald 统计量与得分统计量的等价性。\(\blacksquare\)

\hrulefill

\hrulefill

\hrulefill

\hrulefill

\hrulefill

\hrulefill

\section{第362章：置信区间与贝叶斯方法}\label{ux7b2c362ux7ae0ux7f6eux4fe1ux533aux95f4ux4e0eux8d1dux53f6ux65afux65b9ux6cd5}

本章讨论区间估计和置信区间，以及贝叶斯推断的基本框架。

\subsection{111.1 置信区间}\label{ux7f6eux4fe1ux533aux95f4}

点估计给出了参数的一个最佳猜测，但无法表达该猜测的不确定性。置信区间弥补了这一缺陷：它以预先指定的置信水平 \(1-\alpha\) 给出了一个包含真参数的随机区间。置信区间的构造关键在于寻找枢轴量------其分布与参数无关的统计量------从而将概率陈述转化为参数区间。

\textbf{定义 111.1}（置信区间）：参数 \(\theta\) 的 \(1 - \alpha\) \textbf{置信区间}（或置信集）是只依赖于样本的随机集合 \(C(X)\)，满足

\[\mathbb{P}_\theta(\theta \in C(X)) \geq 1 - \alpha \quad \text{对所有 } \theta \in \Theta\]

\(1 - \alpha\) 称为\textbf{置信水平}（或覆盖概率）。

枢轴量法是最系统化的置信区间构造方案：通过寻找一个分布不变的函数 \(Q(X, \theta)\) 并提供概率界限来反解 \(\theta\)。这一方法将置信区间问题转化为概率计算问题。

\textbf{定义 111.2}（枢轴量）：随机量 \(Q(X, \theta)\) 称为\textbf{枢轴量}（pivotal quantity），如果其分布不依赖于 \(\theta\)。由 \(\mathbb{P}(a \leq Q(X, \theta) \leq b) = 1 - \alpha\) 可解出 \(\theta\) 的置信区间。

置信区间与假设检验之间存在精确的数学对偶关系：一个 \(1-\alpha\) 水平的置信集恰好由所有在水平 \(\alpha\) 下不被拒绝的参数值组成。这一对偶将区间估计和假设检验统一于同一数学框架下，为置信区间的理论性质和构造方法提供了全新的理解视角。

\textbf{定理 111.1}（置信区间与假设检验的对偶）：\(\theta\) 的 \(1 - \alpha\) 置信集是所有不被水平 \(\alpha\) 检验拒绝的 \(\theta_0\) 的集合。反之，水平 \(\alpha\) 的检验可通过检查 \(\theta_0\) 是否在置信区间内构造。

\textbf{证明}：设 \(C(X)\) 是 \(1-\alpha\) 置信集，即 \(\mathbb{P}_{\theta_0}(\theta_0 \in C(X)) \geq 1-\alpha\) 对所有 \(\theta_0\)。定义检验 \(\varphi_{\theta_0}(X) = 1 - \mathbf{1}_{C(X)}(\theta_0)\)，则 \(\mathbb{E}_{\theta_0}[\varphi_{\theta_0}] = \mathbb{P}_{\theta_0}(\theta_0 \notin C(X)) \leq \alpha\)，故 \(\varphi_{\theta_0}\) 是 \(H_0: \theta=\theta_0\) 的水平 \(\alpha\) 检验。反之，给定一族水平 \(\alpha\) 检验 \(\varphi_{\theta_0}\)（每个 \(\theta_0\) 一个），定义 \(C(X) = \{\theta_0 : \varphi_{\theta_0}(X)=0\}\)，则 \(\mathbb{P}_{\theta_0}(\theta_0 \in C(X)) = 1 - \mathbb{E}_{\theta_0}[\varphi_{\theta_0}] \geq 1-\alpha\)。这一对偶关系揭示了区间估计与假设检验在逻辑上的统一：置信区间由所有不被数据拒绝的参数值组成。\(\blacksquare\)

\subsection{111.2 贝叶斯方法}\label{ux8d1dux53f6ux65afux65b9ux6cd5}

\textbf{定义 111.3}（贝叶斯模型）：在贝叶斯框架中，参数 \(\theta\) 被视为随机变量，具有\textbf{先验分布} \(\pi(\theta)\)。给定 \(\theta\)，数据 \(X\) 的分布为 \(f(x \mid \theta)\)。

\textbf{定理 111.2}（贝叶斯定理）：\textbf{后验分布}为

\[\pi(\theta \mid x) = \frac{f(x \mid \theta) \pi(\theta)}{\int_\Theta f(x \mid \theta) \pi(\theta) d\theta} \propto f(x \mid \theta) \pi(\theta)\]

即后验 \(\propto\) 似然 \(\times\) 先验。

\textbf{证明}：由条件概率定义，\((\theta,X)\) 的联合分布密度为 \(f(x\mid\theta)\pi(\theta)\)。X 的边缘密度 \(m(x)=\int f(x\mid\theta)\pi(\theta)d\theta\) 是归一化常数。由条件概率公式，\(\pi(\theta\mid x)=f(x\mid\theta)\pi(\theta)/m(x)\)，故 \(\pi(\theta\mid x)\propto f(x\mid\theta)\pi(\theta)\)。后验分布通过 Bayes 公式将先验信息 \(\pi\) 与数据似然 \(f\) 综合为更新后的信念。\(\blacksquare\)

\textbf{定义 111.4}（共轭先验）：先验族 \(\mathcal{F}\) 称为关于似然 \(f(x \mid \theta)\)\textbf{共轭}的，如果后验分布也属于 \(\mathcal{F}\)。

\textbf{定义 111.5}（贝叶斯估计）：常用贝叶斯估计包括： - \textbf{后验均值}：\(\hat{\theta}_{\text{Bayes}} = \mathbb{E}[\theta \mid X]\)（最小化平方误差损失） - \textbf{后验中位数}：\(\hat{\theta}_{\text{med}}\)（最小化绝对误差损失） - \textbf{最大后验（MAP）}：\(\hat{\theta}_{\text{MAP}} = \arg\max_\theta \pi(\theta \mid X)\)（最小化 0-1 损失）

\textbf{定理 111.3}（Bernstein-von Mises 定理）：在正则条件下，当样本量 \(n \to \infty\) 时，后验分布渐近为正态分布 \(\mathcal{N}(\hat{\theta}_{\text{MLE}}, (n I(\hat{\theta}_{\text{MLE}}))^{-1})\)，且与先验分布无关（先验的影响随 \(n \to \infty\) 消失）。

\textbf{证明}：对数后验 \(\log\pi(\theta\mid X) = \ell_n(\theta) + \log\pi(\theta) + C_n\)（\(C_n\) 为归一化常数）。在 \(\hat{\theta}_{\text{MLE}}\) 处 Taylor 展开：\(\ell_n(\theta) = \ell_n(\hat{\theta}) + \frac{1}{2}(\theta-\hat{\theta})^\top \ell_n''(\hat{\theta})(\theta-\hat{\theta}) + o_P(1)\)。由 MLE 理论，\(-\ell_n''(\hat{\theta})/n \xrightarrow{\mathbb{P}} I(\theta_0)\)，故 \(\ell_n(\theta) = \ell_n(\hat{\theta}) - \frac{n}{2}(\theta-\hat{\theta})^\top I(\theta_0)(\theta-\hat{\theta}) + o_P(1)\)。先验项 \(\log\pi(\theta)\) 是 \(O(1)\) 量级，在 \(n\to\infty\) 下被 \(O(n)\) 的似然项主导，可忽略。因此 \(\pi(\theta\mid X) \propto \exp(-\frac{n}{2}(\theta-\hat{\theta})^\top I(\theta_0)(\theta-\hat{\theta}))(1+o(1))\)，即渐近 \(\mathcal{N}(\hat{\theta}_{\text{MLE}}, (nI)^{-1})\)。这一结果从数学上保证了频率学派与贝叶斯方法在大样本下的相容性。\(\blacksquare\)

\hrulefill

\hrulefill

\hrulefill

\hrulefill

\hrulefill

\newpage

\chapter{10-概率与统计/03-统计学/02-线性模型与多元.md}\label{ux6982ux7387ux4e0eux7edfux8ba103-ux7edfux8ba1ux5b6602-ux7ebfux6027ux6a21ux578bux4e0eux591aux5143.md}

\section{第363章：线性模型}\label{ux7b2c363ux7ae0ux7ebfux6027ux6a21ux578b}

线性模型是统计中最广泛使用的模型框架，包括回归、方差分析（ANOVA）和协方差分析。本章讨论最小二乘估计、Gauss-Markov 定理和方差分析。

\subsection{112.1 线性回归模型}\label{ux7ebfux6027ux56deux5f52ux6a21ux578b}

\textbf{定义 112.1}（线性模型）：\textbf{线性模型}为

\[\mathbf{Y} = \mathbf{X}\boldsymbol{\beta} + \boldsymbol{\varepsilon}\]

其中 \(\mathbf{Y} \in \mathbb{R}^n\) 是响应向量，\(\mathbf{X} \in \mathbb{R}^{n \times p}\) 是设计矩阵（已知），\(\boldsymbol{\beta} \in \mathbb{R}^p\) 是未知回归系数，\(\boldsymbol{\varepsilon}\) 是误差向量，满足 \(\mathbb{E}[\boldsymbol{\varepsilon}] = \mathbf{0}\)，\(\operatorname{Cov}(\boldsymbol{\varepsilon}) = \sigma^2 \mathbf{I}_n\)。

\textbf{定义 112.2}（最小二乘估计，OLS）：\(\boldsymbol{\beta}\) 的\textbf{最小二乘估计}为

\[\hat{\boldsymbol{\beta}} = \arg\min_{\boldsymbol{\beta}} \|\mathbf{Y} - \mathbf{X}\boldsymbol{\beta}\|^2 = (\mathbf{X}^\top \mathbf{X})^{-1} \mathbf{X}^\top \mathbf{Y}\]

（假设 \(\mathbf{X}^\top \mathbf{X}\) 可逆，即 \(\mathbf{X}\) 满列秩）。

\textbf{定理 112.1}（Gauss-Markov 定理）：在 \(\mathbb{E}[\boldsymbol{\varepsilon}] = \mathbf{0}\)，\(\operatorname{Cov}(\boldsymbol{\varepsilon}) = \sigma^2 \mathbf{I}_n\) 的假设下，\(\hat{\boldsymbol{\beta}}\) 是 \(\boldsymbol{\beta}\) 的\textbf{最佳线性无偏估计}（BLUE）：对所有 \(c \in \mathbb{R}^p\)，\(c^\top \hat{\boldsymbol{\beta}}\) 在 \(c^\top \boldsymbol{\beta}\) 的所有线性无偏估计中方差最小。

\emph{证明}：设 \(\tilde{\boldsymbol{\beta}} = \mathbf{C}\mathbf{Y}\) 是任意线性无偏估计。\(\mathbb{E}[\tilde{\boldsymbol{\beta}}] = \mathbf{C}\mathbf{X}\boldsymbol{\beta} = \boldsymbol{\beta}\) 对所有 \(\boldsymbol{\beta}\) 成立，故 \(\mathbf{C}\mathbf{X} = \mathbf{I}\)。\(\operatorname{Cov}(\tilde{\boldsymbol{\beta}}) = \sigma^2 \mathbf{C}\mathbf{C}^\top\)。\(\operatorname{Cov}(\hat{\boldsymbol{\beta}}) = \sigma^2 (\mathbf{X}^\top \mathbf{X})^{-1}\)。需证 \(\mathbf{C}\mathbf{C}^\top - (\mathbf{X}^\top \mathbf{X})^{-1}\) 半正定，这等价于 \((\mathbf{C} - (\mathbf{X}^\top \mathbf{X})^{-1}\mathbf{X}^\top)(\mathbf{C} - (\mathbf{X}^\top \mathbf{X})^{-1}\mathbf{X}^\top)^\top \succeq 0\)。∎

\textbf{定理 112.2}（正态误差下的推断）：若进一步 \(\boldsymbol{\varepsilon} \sim \mathcal{N}(\mathbf{0}, \sigma^2 \mathbf{I}_n)\)，则 1. \(\hat{\boldsymbol{\beta}} \sim \mathcal{N}(\boldsymbol{\beta}, \sigma^2 (\mathbf{X}^\top \mathbf{X})^{-1})\) 2. \(\hat{\sigma}^2 = \|\mathbf{Y} - \mathbf{X}\hat{\boldsymbol{\beta}}\|^2 / (n - p)\) 是 \(\sigma^2\) 的无偏估计，\((n-p)\hat{\sigma}^2/\sigma^2 \sim \chi^2_{n-p}\) 3. \(\hat{\beta}_j / (\hat{\sigma} \sqrt{(\mathbf{X}^\top \mathbf{X})^{-1}_{jj}}) \sim t_{n-p}\)（用于 \(H_0: \beta_j = 0\) 的检验）

\textbf{证明}：（1）\(\hat{\boldsymbol{\beta}} = (\mathbf{X}^\top\mathbf{X})^{-1}\mathbf{X}^\top(\mathbf{X}\boldsymbol{\beta}+\boldsymbol{\varepsilon}) = \boldsymbol{\beta} + (\mathbf{X}^\top\mathbf{X})^{-1}\mathbf{X}^\top\boldsymbol{\varepsilon}\)。\(\boldsymbol{\varepsilon}\sim\mathcal{N}(0,\sigma^2 I)\) 的线性变换保正态，故 \(\hat{\boldsymbol{\beta}} \sim \mathcal{N}(\boldsymbol{\beta}, \sigma^2(\mathbf{X}^\top\mathbf{X})^{-1})\)。（2）残差 \(\mathbf{e} = \mathbf{Y}-\mathbf{X}\hat{\boldsymbol{\beta}} = (I-H)\mathbf{Y}\)，\(H=\mathbf{X}(\mathbf{X}^\top\mathbf{X})^{-1}\mathbf{X}^\top\) 为投影矩阵（幂等、对称），\(\operatorname{tr}(I-H)=n-p\)。由正态二次型的 Cochran 定理，\(\mathbf{e}^\top\mathbf{e}/\sigma^2 \sim \chi^2_{n-p}\)。（3）\(\hat{\beta}_j\) 与 \(\hat{\sigma}^2\) 独立（Cochran 定理的推论），故 \(t = \hat{\beta}_j/(\hat{\sigma}\sqrt{v_{jj}}) \sim t_{n-p}\)。\(\blacksquare\)

\subsection{112.2 方差分析}\label{ux65b9ux5deeux5206ux6790}

\textbf{定义 112.3}（一元方差分析，ANOVA）：\textbf{一元方差分析}比较 \(k\) 个总体的均值：

\[H_0: \mu_1 = \mu_2 = \cdots = \mu_k \quad \text{vs} \quad H_1: \text{至少两个均值不同}\]

\textbf{定义 112.4}（平方和分解）：

\[\underbrace{\sum_{i=1}^k \sum_{j=1}^{n_i} (Y_{ij} - \bar{Y}_{\cdot\cdot})^2}_{\text{SST (总平方和)}} = \underbrace{\sum_{i=1}^k n_i (\bar{Y}_{i\cdot} - \bar{Y}_{\cdot\cdot})^2}_{\text{SSB (组间平方和)}} + \underbrace{\sum_{i=1}^k \sum_{j=1}^{n_i} (Y_{ij} - \bar{Y}_{i\cdot})^2}_{\text{SSW (组内平方和)}}\]

\textbf{定理 112.3}（ANOVA \(F\) 检验）：在 \(H_0\) 下（正态性、等方差假设），

\[F = \frac{\text{SSB} / (k-1)}{\text{SSW} / (n-k)} \sim F_{k-1, n-k}\]

拒绝域：\(F > F_{k-1, n-k, \alpha}\)。

\textbf{证明}：在 \(H_0\) 和正态假设下，\(Y_{ij} \sim \mathcal{N}(\mu, \sigma^2)\) i.i.d.。SST/\(\sigma^2 \sim \chi^2_{n-1}\)。SSW/\(\sigma^2 \sim \chi^2_{n-k}\)（各组内独立估计方差），SSB/\(\sigma^2\) 在 \(H_0\) 下 \(\sim \chi^2_{k-1}\)。由 Cochran 定理，SSB 与 SSW 独立，故 \(F = \frac{\text{SSB}/(k-1)}{\text{SSW}/(n-k)} \sim F_{k-1,n-k}\)。在 \(H_1\) 下 SSB 包含非中心参数，F 值偏大，故右侧拒绝。\(\blacksquare\)

\textbf{定义 112.5}（多元线性回归的 \(F\) 检验）：对 \(H_0: \beta_{q+1} = \cdots = \beta_p = 0\)（嵌套模型比较），

\[F = \frac{(\text{SSE}_{\text{reduced}} - \text{SSE}_{\text{full}}) / (p - q)}{\text{SSE}_{\text{full}} / (n - p)} \sim F_{p-q, n-p}\]

其中 SSE 是残差平方和。

\subsection{112.3 模型诊断与选择}\label{ux6a21ux578bux8bcaux65adux4e0eux9009ux62e9}

\textbf{定义 112.6}（决定系数）：\(R^2 = 1 - \frac{\text{SSE}}{\text{SST}}\) 是模型解释的变异比例。调整 \(R^2\)：\(R^2_{\text{adj}} = 1 - \frac{\text{SSE}/(n-p)}{\text{SST}/(n-1)}\)。

\textbf{定义 112.7}（AIC 与 BIC）：模型选择准则： - \textbf{AIC}（Akaike 信息准则）：\(\text{AIC} = -2 \log L + 2p\) - \textbf{BIC}（贝叶斯信息准则）：\(\text{BIC} = -2 \log L + p \log n\)

选择 AIC/BIC 最小的模型。BIC 对复杂模型的惩罚更重（当 \(n \geq 8\)）。

\hrulefill

\hrulefill

\hrulefill

\hrulefill

\hrulefill

\hrulefill

\section{第364章：多元分析}\label{ux7b2c364ux7ae0ux591aux5143ux5206ux6790}

多元分析处理多个响应变量同时观测的情形。本章介绍多元正态分布、主成分分析（PCA）和判别分析。

\subsection{113.1 多元正态分布}\label{ux591aux5143ux6b63ux6001ux5206ux5e03}

\textbf{定义 113.1}（多元正态分布）：\(\mathbb{R}^p\) 上的随机向量 \(\mathbf{X} \sim \mathcal{N}_p(\boldsymbol{\mu}, \boldsymbol{\Sigma})\)（\(\boldsymbol{\Sigma} \succ 0\)）的密度为

\[f(\mathbf{x}) = \frac{1}{(2\pi)^{p/2} |\boldsymbol{\Sigma}|^{1/2}} \exp\left(-\frac{1}{2} (\mathbf{x} - \boldsymbol{\mu})^\top \boldsymbol{\Sigma}^{-1} (\mathbf{x} - \boldsymbol{\mu})\right)\]

\textbf{命题 113.1}（多元正态的性质）： 1. 线性变换：\(\mathbf{A}\mathbf{X} + \mathbf{b} \sim \mathcal{N}(\mathbf{A}\boldsymbol{\mu} + \mathbf{b}, \mathbf{A}\boldsymbol{\Sigma}\mathbf{A}^\top)\) 2. 边缘分布：\(\mathbf{X}\) 的任意子向量也是多元正态 3. 条件分布：\(\mathbf{X}_1 \mid \mathbf{X}_2 = \mathbf{x}_2\) 是正态，均值 \(\boldsymbol{\mu}_1 + \boldsymbol{\Sigma}_{12}\boldsymbol{\Sigma}_{22}^{-1}(\mathbf{x}_2 - \boldsymbol{\mu}_2)\) 4. 独立性：\(\mathbf{X}_1\) 与 \(\mathbf{X}_2\) 独立当且仅当 \(\boldsymbol{\Sigma}_{12} = \mathbf{0}\)

\textbf{定理 113.1}（Wishart 分布）：设 \(\mathbf{X}_1, \ldots, \mathbf{X}_n\) i.i.d. \(\sim \mathcal{N}_p(\mathbf{0}, \boldsymbol{\Sigma})\)。则 \(\mathbf{W} = \sum_{i=1}^n \mathbf{X}_i \mathbf{X}_i^\top\) 服从 \textbf{Wishart 分布} \(W_p(n, \boldsymbol{\Sigma})\)（多元 \(\chi^2\) 的推广）。

\textbf{证明}：\(\mathbf{X}_i\) 联合密度正比于 \(\exp(-\frac{1}{2}\sum_i \mathbf{X}_i^\top\boldsymbol{\Sigma}^{-1}\mathbf{X}_i) = \exp(-\frac{1}{2}\operatorname{tr}(\boldsymbol{\Sigma}^{-1}\mathbf{W}))\)。变量变换 \(\mathbf{X} \to \mathbf{W}\)（非一一，需考虑 Stiefel 流形上的积分）给出 \(\mathbf{W}\) 的密度：\(f(\mathbf{W}) = c_{n,p}|\boldsymbol{\Sigma}|^{-n/2}|\mathbf{W}|^{(n-p-1)/2}\exp(-\frac{1}{2}\operatorname{tr}(\boldsymbol{\Sigma}^{-1}\mathbf{W}))\)，\(\mathbf{W}\) 正定。当 \(p=1\) 时 \(\boldsymbol{\Sigma}=\sigma^2\)，\(W = \sum X_i^2 \sim \sigma^2\chi^2_n\)，回归经典。Wishart 分布是样本协方差矩阵分布的精确描述。\(\blacksquare\)

\textbf{定理 113.2}（Hotelling \(T^2\) 检验）：对 \(H_0: \boldsymbol{\mu} = \boldsymbol{\mu}_0\)（多元均值检验），

\[T^2 = n(\bar{\mathbf{X}} - \boldsymbol{\mu}_0)^\top \mathbf{S}^{-1} (\bar{\mathbf{X}} - \boldsymbol{\mu}_0)\]

在 \(H_0\) 下，\(\frac{n-p}{p(n-1)} T^2 \sim F_{p, n-p}\)。\(T^2\) 是 \(t\) 检验的多元推广。

\textbf{证明}：\(\bar{\mathbf{X}} \sim \mathcal{N}_p(\boldsymbol{\mu}, \boldsymbol{\Sigma}/n)\)，\((n-1)\mathbf{S} \sim W_p(n-1, \boldsymbol{\Sigma})\)，且 \(\bar{\mathbf{X}}\) 与 \(\mathbf{S}\) 独立。令 \(\mathbf{Z} = \sqrt{n}\boldsymbol{\Sigma}^{-1/2}(\bar{\mathbf{X}}-\boldsymbol{\mu}_0) \sim \mathcal{N}_p(0,I)\)，则 \(T^2 = \mathbf{Z}^\top(\mathbf{S}/\boldsymbol{\Sigma})^{-1}\mathbf{Z}\)。由 Bartlett 分解，\(T^2\) 可化为 \(\frac{p(n-1)}{n-p}F_{p,n-p}\) 在 \(H_0\) 下。当 \(p=1\) 时 \(T^2 = t^2_{n-1}\)。\(\blacksquare\)

\subsection{113.2 主成分分析}\label{ux4e3bux6210ux5206ux5206ux6790}

\textbf{定义 113.2}（主成分分析，PCA）：PCA 寻找数据方差最大的方向。设 \(\mathbf{X}\) 的协方差矩阵为 \(\boldsymbol{\Sigma}\)（或样本协方差矩阵 \(\mathbf{S}\)）。

\textbf{定理 113.3}（PCA 的推导）：设 \(\boldsymbol{\Sigma}\) 的特征分解为 \(\boldsymbol{\Sigma} = \mathbf{V}\boldsymbol{\Lambda}\mathbf{V}^\top\)，其中 \(\boldsymbol{\Lambda} = \operatorname{diag}(\lambda_1, \ldots, \lambda_p)\)，\(\lambda_1 \geq \cdots \geq \lambda_p \geq 0\)。则第 \(k\) 个\textbf{主成分}为 \(Y_k = \mathbf{v}_k^\top \mathbf{X}\)（\(\mathbf{v}_k\) 是 \(\lambda_k\) 对应的特征向量）。

\emph{证明}：\(\max_{\|\mathbf{a}\|=1} \operatorname{Var}(\mathbf{a}^\top \mathbf{X}) = \max_{\|\mathbf{a}\|=1} \mathbf{a}^\top \boldsymbol{\Sigma} \mathbf{a} = \lambda_1\)（由 Rayleigh 商），达到于 \(\mathbf{a} = \mathbf{v}_1\)。第 \(k\) 个主成分在与前 \(k-1\) 个正交的约束下最大化方差。∎

\textbf{命题 113.2}（PCA 的性质）： 1. \(\operatorname{Var}(Y_k) = \lambda_k\) 2. \(\sum_{k=1}^p \lambda_k = \operatorname{tr}(\boldsymbol{\Sigma})\)（总方差） 3. 前 \(k\) 个主成分解释的方差比例为 \(\sum_{i=1}^k \lambda_i / \sum_{i=1}^p \lambda_i\)

\subsection{113.3 判别分析}\label{ux5224ux522bux5206ux6790}

\textbf{定义 113.3}（线性判别分析，LDA）：假设各类服从 \(\mathcal{N}_p(\boldsymbol{\mu}_k, \boldsymbol{\Sigma})\)（等协方差矩阵）。LDA 将 \(\mathbf{x}\) 分类到使后验概率最大的类：

\[\delta_k(\mathbf{x}) = \mathbf{x}^\top \boldsymbol{\Sigma}^{-1} \boldsymbol{\mu}_k - \frac{1}{2} \boldsymbol{\mu}_k^\top \boldsymbol{\Sigma}^{-1} \boldsymbol{\mu}_k + \log \pi_k\]

其中 \(\pi_k\) 是第 \(k\) 类的先验概率。分类规则：\(\hat{k} = \arg\max_k \delta_k(\mathbf{x})\)。

\textbf{定义 113.4}（二次判别分析，QDA）：若各类协方差矩阵不同（\(\boldsymbol{\Sigma}_k\)），判别函数为

\[\delta_k(\mathbf{x}) = -\frac{1}{2} \log |\boldsymbol{\Sigma}_k| - \frac{1}{2} (\mathbf{x} - \boldsymbol{\mu}_k)^\top \boldsymbol{\Sigma}_k^{-1} (\mathbf{x} - \boldsymbol{\mu}_k) + \log \pi_k\]

\textbf{定理 113.4}（Fisher 判别分析）：Fisher 线性判别寻找投影方向 \(\mathbf{a}\) 最大化组间方差与组内方差之比：

\[\max_{\mathbf{a}} \frac{\mathbf{a}^\top \mathbf{B} \mathbf{a}}{\mathbf{a}^\top \mathbf{W} \mathbf{a}}\]

其中 \(\mathbf{B}\) 是组间散布矩阵，\(\mathbf{W}\) 是组内散布矩阵。\(\mathbf{a}\) 是 \(\mathbf{W}^{-1}\mathbf{B}\) 的最大特征值对应的特征向量。这与 LDA（等协方差时）等价。

\textbf{证明}：Fisher 准则 \(J(\mathbf{a}) = \frac{\mathbf{a}^\top\mathbf{B}\mathbf{a}}{\mathbf{a}^\top\mathbf{W}\mathbf{a}}\) 是 \(\mathbf{a}\) 的广义 Rayleigh 商。由广义特征值问题 \(\mathbf{B}\mathbf{a} = \lambda\mathbf{W}\mathbf{a}\)，最优方向 \(\mathbf{a}\) 为 \(\mathbf{W}^{-1}\mathbf{B}\) 的最大特征值 \(\lambda_1\) 对应的特征向量。\(\mathbf{W}^{-1}\mathbf{B}\) 的秩至多为 \(\min(p, k-1)\)（\(k\) 为类别数），故至多可提取 \(k-1\) 个判别方向。\(\mathbf{W}^{-1}\mathbf{B}\) 的非零特征值 \(\lambda_j\) 度量了沿 \(\mathbf{a}_j\) 方向的类间分离程度。\(\blacksquare\)

\hrulefill

\hrulefill

\hrulefill

\hrulefill

\hrulefill

\hrulefill

\section{第365章：信息几何}\label{ux7b2c365ux7ae0ux4fe1ux606fux51e0ux4f55}

信息几何由Amari和Chentsov在1980年代系统建立，将参数化概率分布族\(\mathcal{M} = \{p(x \mid \theta) : \theta \in \Theta\}\)视为统计流形，赋予Fisher信息度量\(g_{ij}(\theta) = \mathbb{E}[\partial_i \ell \cdot \partial_j \ell]\)（其中\(\ell = \log p\)）的Riemann几何结构。这一度量在参数变换下协变且由Chentsov定理唯一确定。进一步引入\(\alpha\)-联络\(\Gamma_{ijk}^{(\alpha)} = \mathbb{E}[(\partial_i\partial_j\ell + \frac{1-\alpha}{2}\partial_i\ell \cdot \partial_j\ell)\partial_k\ell]\)，\(\alpha = \pm 1\)分别对应指数联络和混合联络，两者关于Fisher度量对偶。Amari-Chentsov张量\(T_{ijk} = \mathbb{E}[\partial_i\ell \cdot \partial_j\ell \cdot \partial_k\ell]\)度量了统计流形的非平坦性。信息几何在机器学习中为自然梯度下降和EM算法的几何解释提供了理论基础，在神经网络的信息瓶颈原理和深度学习的流形学习中也有重要应用。

\subsection{Fisher信息度量与Rao距离}\label{fisherux4fe1ux606fux5ea6ux91cfux4e0eraoux8dddux79bb}

\textbf{定义114.1}（Fisher信息矩阵）：设 \(\mathcal{M} = \{p(x \mid \theta) : \theta \in \Theta \subseteq \mathbb{R}^n\}\) 为参数化概率分布族，满足正则条件（可交换微分与积分）。\textbf{Fisher信息矩阵} \(G(\theta) = [g_{ij}(\theta)]\) 定义为 \[g_{ij}(\theta) = \mathbb{E}_\theta\left[ \frac{\partial \log p}{\partial \theta^i} \frac{\partial \log p}{\partial \theta^j} \right] = -\mathbb{E}_\theta\left[ \frac{\partial^2 \log p}{\partial \theta^i \partial \theta^j} \right]\] 第二个等号在正则条件下成立。Fisher信息矩阵在参数变换 \(\theta \mapsto \eta(\theta)\) 下按张量法则变换：\(g_{ij}(\theta) = \frac{\partial \eta^k}{\partial \theta^i} \frac{\partial \eta^l}{\partial \theta^j} g_{kl}(\eta)\)，确认了 \(g_{ij}\) 的（协变2阶）张量身份。因此 \((\mathcal{M}, G)\) 构成一个 Riemann 流形，记为 \((S, g)\)，称为\textbf{统计流形}。

Fisher信息度量的测地线给出分布间的最短路径。两个分布 \(p(x \mid \theta_1)\) 与 \(p(x \mid \theta_2)\) 之间的 \textbf{Rao 距离}定义为 \[D_R(\theta_1, \theta_2) = \min_{\gamma} \int_0^1 \sqrt{g_{ij}(\gamma(t)) \dot{\gamma}^i(t) \dot{\gamma}^j(t)} \, dt\] 其中 \(\gamma\) 取遍连接 \(\theta_1\) 和 \(\theta_2\) 的曲线。对一元正态分布族 \(\mathcal{N}(\mu, \sigma^2)\)，Fisher度量为 \(ds^2 = (d\mu^2 + 2d\sigma^2) / \sigma^2\)，其测地线对应双曲几何中的直线，Rao距离正是 Poincaré 上半平面模型中的双曲距离。对于多项分布族（单纯形），Fisher度量导出球面几何------这是 Fisher 信息的''信息等距嵌入''的经典实例。

\subsection{Chentsov定理}\label{chentsovux5b9aux7406}

\textbf{定理114.1}（Chentsov定理，1972/1982）：在有限样本空间 \(\mathcal{X} = \{1, 2, \ldots, n\}\) 上的所有严格正概率分布构成的流形 \(\mathcal{P}_n = \{p \in \mathbb{R}^n : p_i > 0, \sum_i p_i = 1\}\) 上，差一个正的常数因子，Fisher信息度量是唯一在\textbf{充分统计量}（Markov 核）作用下单调不增的 Riemann 度量。换言之，任何满足信息单调性（即对任意统计模型间的 Markov 嵌入 \(f : \mathcal{M} \to \mathcal{N}\)，诱导的度量满足 \(g_{\mathcal{M}} \geq f^* g_{\mathcal{N}}\)）的 Riemann 度量必为 Fisher 信息度量的常数倍。

\textbf{证明思路}：证明的核心是利用 Markov 核的范畴性质。设 \(g\) 是 \(\mathcal{P}_n\) 上的任意 Riemann 度量，满足 Markov 嵌入单调性。首先，利用 \(\mathcal{P}_n\) 在有限置换群下的不变性，论证 \(g\) 在任意点 \(p\) 处必为形式 \(g_p(u, v) = \sum_i \phi(p_i) u_i v_i\)（对角型），其中 \(\phi : (0, 1) \to \mathbb{R}^+\) 为待定函数。其次，考虑二元随机变量的充分统计量诱导的 Markov 嵌入 \(\mathcal{P}_2 \hookrightarrow \mathcal{P}_n\)，由信息单调性导出 \(\phi\) 必须满足函数方程 \(\phi(p) = c/p\)（\(c > 0\) 常数）。最后，将 \(g_p(u, v) = c \sum_i u_i v_i / p_i\) 限制到切空间 \(T_p \mathcal{P}_n = \{u : \sum_i u_i = 0\}\) 上，即得 Fisher 信息度量在 \(\mathcal{P}_n\) 上的表达式。对于一般的参数模型 \(\mathcal{M}\)，Chentsov定理通过嵌入 \(\mathcal{M}\) 到充分大的多项分布流形 \(\mathcal{P}_N\) 中的论证推广。定理的深刻之处在于------它不是对一个具体公式的验证，而是从一个基本原则（信息处理不增加信息）出发的唯一性定理，与量子信息论中 Petz 定理（唯一单调 Riemann 度量为量子 Fisher 信息）形成经典-量子的完美对照。\(\blacksquare\)

\subsection{\texorpdfstring{\(\alpha\)-联络与Amari理论}{\textbackslash alpha-联络与Amari理论}}\label{alpha-ux8054ux7edcux4e0eamariux7406ux8bba}

\textbf{定义114.2}（\(\alpha\)-联络）：在统计流形 \((S, g)\) 上，对任意实数 \(\alpha \in \mathbb{R}\)，定义 \(\alpha\)-联络的 Christoffel 符号为 \[\Gamma_{ijk}^{(\alpha)}(\theta) = \mathbb{E}_\theta\left[ \left(\partial_i \partial_j \ell + \frac{1 - \alpha}{2} \partial_i \ell \cdot \partial_j \ell\right) \partial_k \ell \right]\] 其中 \(\ell = \log p(x \mid \theta)\)。\(\alpha = 1\) 时 \(\Gamma_{ijk}^{(1)} = \mathbb{E}[\partial_i \partial_j \ell \cdot \partial_k \ell]\) 称为\textbf{指数联络}（exponential connection / e-connection）；\(\alpha = -1\) 时 \(\Gamma_{ijk}^{(-1)} = \mathbb{E}[(\partial_i \partial_j \ell + \partial_i \ell \cdot \partial_j \ell) \partial_k \ell]\) 称为\textbf{混合联络}（mixture connection / m-connection）。

\(\alpha\)-联络的核心性质是\textbf{对偶性}：\(\Gamma^{(+\alpha)}\) 和 \(\Gamma^{(-\alpha)}\) 关于 Fisher 度量 \(g\) 互为对偶联络，即满足 \[\partial_k g_{ij} = \Gamma_{ki,j}^{(\alpha)} + \Gamma_{kj,i}^{(-\alpha)}\] 其中 \(\Gamma_{ki,j}^{(\alpha)} = g_{jm} \Gamma_{ki}^{(\alpha)m}\)。特别地，指数联络（\(\alpha = 1\)）和混合联络（\(\alpha = -1\)）是一对对偶平坦联络：指数联络在指数族参数化下曲率张量为零，混合联络在混合族参数化下曲率张量为零。这一对偶平坦结构是信息几何中最优美的几何事实，它使得广义 Pythagoras 定理和投影定理在统计流形上成立。

\textbf{定理114.2}（Amari的广义Pythagoras定理）：设 \((S, g, \nabla, \nabla^*)\) 为对偶平坦统计流形（即 \(\nabla\) 和 \(\nabla^*\) 的曲率均为零）。设 \(P, Q, R \in S\) 满足 \(\nabla\)-测地线连接 \(Q\) 和 \(R\)，且 \(\nabla^*\)-测地线连接 \(P\) 和 \(Q\) 并与前者正交。则 \[D(P, R) = D(P, Q) + D(Q, R)\] 其中 \(D\) 为 \(\nabla\) 诱导的 Bregman 散度。特别地，取 \(\nabla\) 为指数联络、\(\nabla^*\) 为混合联络时，\(D\) 为 KL 散度，上述定理给出了 EM 算法的几何解释：E-步将当前估计投影到数据流形上（m-投影），M-步将投影结果投影到模型流形上（e-投影），每一步都沿对偶测地线进行，保证 KL 散度单调递减。

\textbf{定义114.3}（自然梯度）：在 Riemann 流形 \((S, g)\) 上，目标函数 \(L(\theta)\) 的\textbf{自然梯度}定义为 \[\tilde{\nabla} L(\theta) = G(\theta)^{-1} \nabla L(\theta)\] 其中 \(\nabla L\) 为通常的欧氏梯度。自然梯度下降 \(\theta_{t+1} = \theta_t - \eta_t G(\theta_t)^{-1} \nabla L(\theta_t)\) 等价于沿 Fisher 度量下的最速下降方向移动，在参数变换下具有协变性。自然梯度在强化学习（TRPO, PPO）和变分推断中通过约束 KL 散度实现了稳定高效的优化。在统计流形上，自然梯度的方向恰好是混合联络下的测地线方向，这一几何洞察是 Amari 理论对机器学习方法论的最直接贡献。

\hrulefill

\hrulefill

\hrulefill

\hrulefill

\hrulefill

\hrulefill

\section{第366章：试验设计}\label{ux7b2c366ux7ae0ux8bd5ux9a8cux8bbeux8ba1}

试验设计（DOE）由Fisher在1920年代创立，其核心是Fisher三原则：(1)随机化------将实验单元随机分配以消除系统偏差；(2)重复------通过多次观测估计实验误差；(3)区组化------将同质单元分组以控制已知变异源。拉丁方设计通过行列双重区组控制两个干扰因素，与Greco-Latin方推广到多个因素。因子设计系统考察多个因素的主效应和交互效应，\(2^k\)因子设计是最基本的类型；当实验次数不可行时，采用分数因子设计以\(2^{k-p}\)次实验估计主要效应。响应面方法（RSM）由Box和Wilson于1951年提出，通过拟合二次响应曲面模型\(y = \beta_0 + \sum \beta_i x_i + \sum \beta_{ij}x_i x_j + \sum \beta_{ii}x_i^2 + \varepsilon\)，使用最速上升法寻优，是工业过程优化的核心工具。

\subsection{Fisher信息矩阵与最优设计准则}\label{fisherux4fe1ux606fux77e9ux9635ux4e0eux6700ux4f18ux8bbeux8ba1ux51c6ux5219}

\textbf{定义115.1}（设计测度与信息矩阵）：考虑线性模型 \(y_i = f(x_i)^T \beta + \varepsilon_i\)，其中 \(f(x) \in \mathbb{R}^p\) 为回归函数向量，\(\varepsilon_i \sim \text{i.i.d. } \mathcal{N}(0, \sigma^2)\)。一个（近似）\textbf{设计} \(\xi\) 是设计空间 \(\mathcal{X}\) 上的概率测度，将质量 \(w_k\) 赋给支撑点 \(x_k\)。该设计的标准化 \textbf{Fisher 信息矩阵}为 \[M(\xi) = \int_{\mathcal{X}} f(x) f(x)^T \, d\xi(x) = \sum_k w_k f(x_k) f(x_k)^T\] 最小二乘估计 \(\hat{\beta}\) 的协方差矩阵为 \((\sigma^2/n) M(\xi)^{-1}\)。最优设计的本质是选择 \(\xi\) 使 \(M(\xi)\) 在某种意义下''最大''。

\textbf{定义115.2}（D-最优准则）：设计 \(\xi^*\) 是 \textbf{D-最优}的，若它最大化信息矩阵的行列式： \[\xi^* = \arg\max_{\xi} \det M(\xi)\] 等价于最小化 \(\hat{\beta}\) 的置信椭球体积。\(\det M(\xi)^{-1}\) 与 \(\hat{\beta}\) 的广义方差成正比。

\textbf{定义115.3}（A-最优准则）：设计 \(\xi^*\) 是 \textbf{A-最优}的，若它最小化信息矩阵逆的迹： \[\xi^* = \arg\min_{\xi} \operatorname{tr} M(\xi)^{-1}\] 这等价于最小化 \(\hat{\beta}\) 各分量方差的平均值（平均方差最优）。

\textbf{定义115.4}（E-最优准则）：设计 \(\xi^*\) 是 \textbf{E-最优}的，若它最大化信息矩阵的最小特征值： \[\xi^* = \arg\max_{\xi} \lambda_{\min}(M(\xi))\] E-最优最小化 \(c^T \hat{\beta}\) 的最大可能方差（在所有 \(\|c\| = 1\) 的方向上），即控制最坏情况。

\textbf{定义115.5}（G-最优准则）：设计 \(\xi^*\) 是 \textbf{G-最优}的，若它最小化预测方差的最大值： \[\xi^* = \arg\min_{\xi} \max_{x \in \mathcal{X}} d(x, \xi)\] 其中 \(d(x, \xi) = f(x)^T M(\xi)^{-1} f(x)\) 为标准化预测方差函数。

\subsection{Kiefer-Wolfowitz等价定理}\label{kiefer-wolfowitzux7b49ux4ef7ux5b9aux7406}

\textbf{定理115.1}（Kiefer-Wolfowitz等价定理，1960）：对线性模型，以下等价：(i) \(\xi^*\) 是 D-最优设计；(ii) \(\xi^*\) 是 G-最优设计；(iii) \(\max_{x \in \mathcal{X}} d(x, \xi^*) = p\)（参数个数），且最大值在 \(\xi^*\) 的支撑点上达到。

\textbf{证明}：D-最优性的一阶必要条件：对任意设计 \(\eta\)，方向导数 \(\partial \log \det M((1-\alpha)\xi^* + \alpha\eta) / \partial \alpha|_{\alpha=0} \leq 0\)。计算行列式的方向导数：\(\frac{\partial}{\partial \alpha} \log \det M(\xi_\alpha) = \operatorname{tr}(M(\xi_\alpha)^{-1} \frac{\partial M(\xi_\alpha)}{\partial \alpha})\)。在 \(\alpha = 0\) 处，\(\frac{\partial M(\xi_\alpha)}{\partial \alpha}|_{\alpha=0} = M(\eta) - M(\xi^*)\)。故必要条件为 \(\operatorname{tr}(M(\xi^*)^{-1} M(\eta)) \leq p\) 对所有 \(\eta\)，即 \(\int_{\mathcal{X}} d(x, \xi^*) d\eta(x) \leq p\)。取 \(\eta\) 为 Dirac 测度 \(\delta_x\)，得 \(d(x, \xi^*) \leq p\) 对所有 \(x\)。因此 \(\max_x d(x, \xi^*) \leq p\)。

同时，\(\int_{\mathcal{X}} d(x, \xi^*) d\xi^*(x) = \operatorname{tr}(M(\xi^*)^{-1} \int f(x)f(x)^T d\xi^*(x)) = \operatorname{tr}(I_p) = p\)。故平均值 \(p\) 结合逐点 \(\leq p\) 推得 \(\max_x d(x, \xi^*) = p\)，且支撑点处的 \(d(x, \xi^*)\) 必达到 \(p\)。反之，若 \(\max_x d(x, \xi^*) = p\)，则 \(\xi^*\) 满足一阶条件。由 \(\log \det\) 的凹性（关于 \(M\)），一阶条件也是全局最优的充分条件，从而 D-最优等价于 G-最优。\(\blacksquare\)

Kiefer-Wolfowitz 定理是试验设计理论的核心结果。它提供了可计算的验后标准：计算 \(\max_x d(x, \xi)\)，若超过 \(p\)，则将新的支撑点加入设计并重新优化权重。这是 Fedorov-Wynn 序列设计构造算法的基础。

\subsection{响应面设计的严格推导}\label{ux54cdux5e94ux9762ux8bbeux8ba1ux7684ux4e25ux683cux63a8ux5bfc}

\textbf{定义115.6}（响应面模型）：二阶响应面模型为 \[y = \beta_0 + \sum_{i=1}^k \beta_i x_i + \sum_{i=1}^k \beta_{ii} x_i^2 + \sum_{1 \leq i < j \leq k} \beta_{ij} x_i x_j + \varepsilon\] 共需估计 \(p = 1 + k + k + k(k-1)/2 = (k+1)(k+2)/2\) 个参数。

\textbf{中心复合设计}（CCD, Box-Wilson 1951）是最经典的二阶响应面设计，由三部分组成：(i) 角点（factorial points）------\(2^k\) 或 \(2^{k-p}\) 因子设计的 \(\pm 1\) 点，估计线性效应和交互效应；(ii) 轴点（star points）------\(2k\) 个点 \((\pm \alpha, 0, \ldots, 0), (0, \pm \alpha, 0, \ldots, 0), \ldots\)，用于估计二次项；(iii) 中心点------\((0, 0, \ldots, 0)\) 的 \(n_c\) 次重复，用于估计纯误差。\(\alpha\) 的选择决定设计的\textbf{可旋转性}：\(\alpha = (2^k)^{1/4}\)（对于全因子角点）使 \(M(\xi)\) 在正交变换下不变（即 \(d(x, \xi)\) 仅依赖于 \(\|x\|\)），从而预测方差具有球对称性。Box-Behnken 设计是另一种广泛使用的三水平设计，通过平衡不完全区组设计（BIBD）的构型组合，在因子数为 3-7 时具有极高的效率。

\hrulefill

\hrulefill

\hrulefill

\hrulefill

\hrulefill

\hrulefill

\section{第367章：抽样方法与调查设计（Sampling Methods）}\label{ux7b2c367ux7ae0ux62bdux6837ux65b9ux6cd5ux4e0eux8c03ux67e5ux8bbeux8ba1sampling-methods}

抽样理论是统计学从有限总体中选取子集进行推断的方法论。\textbf{基本抽样方案}：（1）\textbf{简单随机抽样（SRS）}------每个大小为 \(n\) 的子集等概率抽取，均值 \(\bar{y}\) 是总体均值 \(\bar{Y}\) 的无偏估计，方差为 \((1 - \frac{n}{N}) \frac{S^2}{n}\)（有限总体修正因子 \(1 - \frac{n}{N}\)）；（2）\textbf{分层抽样}------将总体分为 \(H\) 层，在各层内独立抽样，通过 Neyman 最优分配 \(n_h \propto N_h S_h\) 最小化估计量的方差；（3）\textbf{系统抽样}------按固定间隔选取，当排序与变量相关时可显著提高精度；（4）\textbf{整群抽样}------在难以获取总体完整名单时，随机选取群（cluster）后调查群内全部个体；（5）\textbf{多阶段抽样}------先抽初级单元，再在选中单元内抽次级单元，逐级构造。\textbf{比值估计与回归估计}利用辅助变量提高精度：当辅助变量 \(x\) 与目标变量 \(y\) 高度线性相关时，比值估计 \(\hat{Y}_R = \frac{\bar{y}}{\bar{x}} X\) 或回归估计 \(\hat{Y}_{\text{reg}} = \bar{y} + \hat{\beta}(X - \bar{x})\) 可在大样本下显著降低方差。\textbf{Horvitz-Thompson 估计量}：\(\hat{Y}_{HT} = \sum_{i \in \mathcal{S}} \frac{y_i}{\pi_i}\)（\(\pi_i\) 为个体 \(i\) 的人样概率）对所有不等概率抽样设计一致有效。

\subsection{Horvitz-Thompson估计量的严格理论}\label{horvitz-thompsonux4f30ux8ba1ux91cfux7684ux4e25ux683cux7406ux8bba}

\textbf{定义116.1}（不等概率抽样设计）：设有限总体 \(U = \{1, 2, \ldots, N\}\)。抽样设计 \(p(\cdot)\) 指定每个子集 \(\mathcal{S} \subseteq U\) 被抽取的概率 \(p(\mathcal{S})\)。对每个个体 \(i \in U\)，其一阶包含概率为 \(\pi_i = \Pr(i \in \mathcal{S}) = \sum_{\mathcal{S} \ni i} p(\mathcal{S})\)。二阶包含概率为 \(\pi_{ij} = \Pr(i, j \in \mathcal{S}) = \sum_{\mathcal{S} \supseteq \{i,j\}} p(\mathcal{S})\)。要求 \(\pi_i > 0\) 对所有 \(i\)（每个个体有正的人样概率）。

\textbf{定理116.1}（Horvitz-Thompson估计量的无偏性，1952）：设 \(y_i\) 为个体 \(i\) 的观测值。总体总量 \(Y = \sum_{i=1}^N y_i\) 的 \textbf{Horvitz-Thompson 估计量} \(\hat{Y}_{HT} = \sum_{i \in \mathcal{S}} \frac{y_i}{\pi_i}\) 是 \(Y\) 的无偏估计，即 \(\mathbb{E}[\hat{Y}_{HT}] = Y\)。

\textbf{证明}：定义指示变量 \(I_i = \mathbf{1}_{\{i \in \mathcal{S}\}}\)，则 \(\mathbb{E}[I_i] = \Pr(i \in \mathcal{S}) = \pi_i\)。\(\hat{Y}_{HT} = \sum_{i \in \mathcal{S}} y_i / \pi_i = \sum_{i=1}^N I_i y_i / \pi_i\)。取期望：\(\mathbb{E}[\hat{Y}_{HT}] = \sum_{i=1}^N \frac{y_i}{\pi_i} \mathbb{E}[I_i] = \sum_{i=1}^N \frac{y_i}{\pi_i} \cdot \pi_i = \sum_{i=1}^N y_i = Y\)。证明仅依赖 \(\pi_i > 0\)，无需对抽样设计的其他假设。\(\blacksquare\)

\textbf{定理116.2}（Sen-Yates-Grundy方差公式，1953-1954）：Horvitz-Thompson 估计量的方差由下式给出（要求设计满足 \(\pi_{ij} > 0\) 对所有 \(i \neq j\)）： \[\operatorname{Var}(\hat{Y}_{HT}) = \sum_{i=1}^N \frac{1 - \pi_i}{\pi_i} y_i^2 + \sum_{i=1}^N \sum_{j \neq i} \frac{\pi_{ij} - \pi_i \pi_j}{\pi_i \pi_j} y_i y_j\] 等价地，Sen-Yates-Grundy 形式的方差为 \[\operatorname{Var}(\hat{Y}_{HT}) = -\frac{1}{2} \sum_{i=1}^N \sum_{j \neq i} (\pi_{ij} - \pi_i \pi_j) \left( \frac{y_i}{\pi_i} - \frac{y_j}{\pi_j} \right)^2\]

\textbf{证明}：\(\operatorname{Var}(\hat{Y}_{HT}) = \mathbb{E}[\hat{Y}_{HT}^2] - Y^2 = \mathbb{E}\left[ \sum_{i \in \mathcal{S}} \sum_{j \in \mathcal{S}} \frac{y_i y_j}{\pi_i \pi_j} \right] - Y^2\)。将双重和拆分为 \(i=j\) 与 \(i \neq j\)：\(= \sum_{i=1}^N \frac{y_i^2}{\pi_i^2} \mathbb{E}[I_i] + \sum_{i \neq j} \frac{y_i y_j}{\pi_i \pi_j} \mathbb{E}[I_i I_j] - Y^2\)。由 \(\mathbb{E}[I_i] = \pi_i\) 和 \(\mathbb{E}[I_i I_j] = \pi_{ij}\)，代入得第一式。展开第一式的第二项并减去 \(Y^2 = \sum_i y_i^2 + \sum_{i \neq j} y_i y_j\)，经代数整理（利用 \(\sum_{j \neq i}(\pi_{ij} - \pi_i \pi_j) = \pi_i(1-\pi_i)\)）得到 Sen-Yates-Grundy 形式。该形式的优势是显式为负半定------使方差估计自然非负。\(\blacksquare\)

\subsection{比率估计与回归估计的渐近理论}\label{ux6bd4ux7387ux4f30ux8ba1ux4e0eux56deux5f52ux4f30ux8ba1ux7684ux6e10ux8fd1ux7406ux8bba}

\textbf{定义116.2}（比率估计）：当辅助变量 \(x_i\)（总体总量 \(X\) 已知）与目标变量 \(y_i\) 高度线性相关时，\textbf{比率估计量}为 \[\hat{Y}_R = \frac{\hat{Y}_{HT}}{\hat{X}_{HT}} \cdot X = \hat{R} X, \quad \hat{R} = \frac{\sum_{i \in \mathcal{S}} y_i / \pi_i}{\sum_{i \in \mathcal{S}} x_i / \pi_i}\]

\textbf{定理116.3}（比率估计的渐近性质）：在适当的抽样设计下，当 \(n \to \infty\) 且 \(N \to \infty\) 时，\(\hat{Y}_R\) 是渐近无偏的，且渐近方差为 \[\operatorname{AVar}(\hat{Y}_R) = \sum_{i} \sum_{j} (\pi_{ij} - \pi_i \pi_j) \frac{y_i - R x_i}{\pi_i} \cdot \frac{y_j - R x_j}{\pi_j}\] 其中 \(R = Y/X\) 为总体比率。比率估计优于简单 HT 估计当且仅当 \(\rho_{xy} > \frac{1}{2} \frac{CV_x}{CV_y}\)（相关系数足够大）。

\textbf{定义116.3}（回归估计）：回归估计利用 \(y_i\) 对 \(x_i\) 的最小二乘回归进行校准：\(\hat{Y}_{\text{reg}} = \hat{Y}_{HT} + \hat{B} (X - \hat{X}_{HT})\)，其中 \(\hat{B}\) 为 \(\{y_i\}\) 对 \(\{x_i\}\) 的估计回归系数。回归估计的渐近方差为 \[\operatorname{AVar}(\hat{Y}_{\text{reg}}) = \sum_{i} \sum_{j} (\pi_{ij} - \pi_i \pi_j) \frac{E_i}{\pi_i} \cdot \frac{E_j}{\pi_j}\] 其中 \(E_i = y_i - B x_i\) 为总体回归残差（\(B\) 为总体回归系数）。回归估计对方差的降低来自残差 \(E_i\) 比原始 \(y_i\) 具有更小的变异。在线性模型中回归估计是渐近最优的。

\subsection{Neyman最优分配}\label{neymanux6700ux4f18ux5206ux914d}

\textbf{定理116.4}（Neyman最优分配定理，1934）：在分层抽样中，固定总样本量 \(n = \sum_{h=1}^H n_h\) 和分层边界，最小化估计量 \(\hat{Y}_{st} = \sum_{h=1}^H N_h \bar{y}_h\) 的方差的最优分配为 \[n_h = n \cdot \frac{N_h S_h}{\sum_{k=1}^H N_k S_k}, \quad h = 1, \ldots, H\] 其中 \(S_h^2 = \frac{1}{N_h-1} \sum_{i \in U_h} (y_i - \bar{Y}_h)^2\) 为第 \(h\) 层的总体方差。

\textbf{证明}：\(\hat{Y}_{st}\) 的方差为 \(\operatorname{Var}(\hat{Y}_{st}) = \sum_{h=1}^H N_h^2 \frac{S_h^2}{n_h} \left(1 - \frac{n_h}{N_h}\right)\)。忽略有限总体修正（\(1 - n_h/N_h \approx 1\) 对大总体），最小化 \(\sum_h N_h^2 S_h^2 / n_h\) 受 \(\sum_h n_h = n\) 约束。构造 Lagrange 函数 \(\mathcal{L} = \sum_h \frac{N_h^2 S_h^2}{n_h} + \lambda(\sum_h n_h - n)\)。令 \(\partial \mathcal{L} / \partial n_h = -N_h^2 S_h^2 / n_h^2 + \lambda = 0\)，得 \(n_h = N_h S_h / \sqrt{\lambda}\)，比例 \(n_h \propto N_h S_h\)。由约束 \(\sum_h n_h = n\) 解得 \(n_h = n \cdot N_h S_h / \sum_k N_k S_k\)。由于约束集 \(n_h \geq 0\) 且目标函数为凸，这是全局最优解。该分配在层内方差大的层投入更多样本，实现了信息获取的效率最大化。若还需考虑每层抽样成本 \(c_h\)，则最优分配推广为 \(n_h \propto N_h S_h / \sqrt{c_h}\)。\(\blacksquare\)

\hrulefill

\hrulefill

\hrulefill

\hrulefill

\hrulefill

\hrulefill

\section{第368章：时间序列分析（Time Series Analysis）}\label{ux7b2c368ux7ae0ux65f6ux95f4ux5e8fux5217ux5206ux6790time-series-analysis}

时间序列分析研究依时间顺序观测的随机过程 \(\{X_t\}_{t \in \mathbb{Z}}\)，其核心问题是建模、预测和推断动态依赖结构。

\subsection{117.1 线性时间序列模型}\label{ux7ebfux6027ux65f6ux95f4ux5e8fux5217ux6a21ux578b}

\textbf{定义 117.1}（平稳过程）：时间序列 \(\{X_t\}\) 称为\textbf{（弱）平稳}的，如果 \(\mathbb{E}[X_t] = \mu\)（常数），且自协方差函数 \(\gamma(h) = \operatorname{Cov}(X_t, X_{t+h})\) 仅依赖于时滞 \(h\)。严平稳要求任意有限维联合分布平移不变。

\textbf{定义 117.2}（白噪声 / White Noise）：过程 \(\{\varepsilon_t\}\) 称为\textbf{白噪声}，记作 \(\varepsilon_t \sim \text{WN}(0, \sigma^2)\)，若 \(\mathbb{E}[\varepsilon_t] = 0\)，\(\operatorname{Var}(\varepsilon_t) = \sigma^2\)，且 \(\operatorname{Cov}(\varepsilon_t, \varepsilon_s) = 0\)（\(t \neq s\)）。

\textbf{定义 117.3}（ARMA(p,q) 模型 / 自回归滑动平均）：

\[X_t = \sum_{i=1}^p \phi_i X_{t-i} + \varepsilon_t + \sum_{j=1}^q \theta_j \varepsilon_{t-j}\]

其中 \(\{\varepsilon_t\}\) 为白噪声。AR（自回归）部分的平稳性由特征方程 \(\Phi(z) = 1 - \phi_1 z - \cdots - \phi_p z^p = 0\) 的根在单位圆外的条件给出；MA（滑动平均）部分本身即平稳。ARMA 过程的\textbf{自相关函数}（ACF）在 MA 部分截尾、AR 部分拖尾，是模型识别的关键工具。

\textbf{定义 117.4}（ARIMA(p,d,q) 模型 / Box-Jenkins, 1970）：若 \(\{X_t\}\) 经 \(d\) 次差分后成为平稳的 ARMA 过程，即 \((1-B)^d X_t\) 服从 ARMA(p,q)（\(B\) 为后移算子 \(B X_t = X_{t-1}\)），则称 \(\{X_t\}\) 为 ARIMA(p,d,q) 过程。这是处理含趋势的非平稳时间序列的标准方法。

\textbf{定义 117.5}（季节性 ARIMA / SARIMA）：\(\Phi_P(B^s) \phi_p(B) (1-B)^d (1-B^s)^D X_t = \Theta_Q(B^s) \theta_q(B) \varepsilon_t\)，其中 \(s\) 为季节周期（如 \(s = 12\) 为月数据），\(P, D, Q\) 为季节阶数。

\textbf{定义 117.6}（GARCH(p,q) 模型 / 广义自回归条件异方差，Engle 1982 / Bollerslev 1986）：

\[\begin{aligned} X_t &= \sigma_t Z_t, \quad Z_t \sim \text{i.i.d.}(0,1) \\ \sigma_t^2 &= \omega + \sum_{i=1}^p \alpha_i X_{t-i}^2 + \sum_{j=1}^q \beta_j \sigma_{t-j}^2 \end{aligned}\]

GARCH 模型捕捉金融时间序列中波动率聚类（volatility clustering）------大幅波动倾向于跟随大幅波动。Engle 因提出 ARCH 模型（1982）获 2003 年 Nobel 经济学奖。

\subsection{117.2 谱分析与频域方法}\label{ux8c31ux5206ux6790ux4e0eux9891ux57dfux65b9ux6cd5}

\textbf{定义 117.7}（谱密度 / Spectral Density）：平稳时间序列的\textbf{谱密度}为自协方差函数的 Fourier 变换

\[f(\omega) = \frac{1}{2\pi} \sum_{h=-\infty}^\infty \gamma(h) e^{-i\omega h}, \quad \omega \in [-\pi, \pi]\]

其中 \(\gamma(h) = \operatorname{Cov}(X_t, X_{t+h})\)。反演公式为 \(\gamma(h) = \int_{-\pi}^\pi e^{i\omega h} f(\omega) d\omega\)。\(f(\omega) \geq 0\) 对 \(\omega\) 的积分等于 \(\operatorname{Var}(X_t)\)，刻画了方差在频率域的分布。

\textbf{定义 117.8}（周期图 / Periodogram）：\(n\) 个观测的\textbf{周期图}为

\[I_n(\omega) = \frac{1}{2\pi n} \left|\sum_{t=1}^n X_t e^{-i\omega t}\right|^2\]

周期图 \(I_n(\omega)\) 是谱密度 \(f(\omega)\) 的渐近无偏估计，但非相合（其方差在 \(n \to \infty\) 时不趋于零）。通过\textbf{窗口平滑}（如 Daniell 窗口、Bartlett 窗口）可得到相合估计 \(\hat{f}(\omega) = \sum W_h I_n(\omega_h)\)。

\textbf{定理 117.1}（Whittle 似然 / Whittle Likelihood）：在频域中，Gaussian 平稳过程的负对数似然近似为

\[-\ell(\theta) \approx \frac{1}{2\pi} \int_{-\pi}^\pi \left( \log f_\theta(\omega) + \frac{I_n(\omega)}{f_\theta(\omega)} \right) d\omega\]

这是频域 ARMA 参数估计（Whittle 估计）的基础，可通过快速 Fourier 变换高效计算。

\textbf{证明}：Gaussian 平稳序列 \(X_1,\ldots,X_n\) 的精确似然涉及 \(n\times n\) Toeplitz 协方差矩阵的逆，复杂度 \(O(n^3)\)。Whittle 近似基于 Szegő 极限定理：当 \(n\to\infty\) 时，\(\frac{1}{n}\log\det(\mathbf{\Sigma}_\theta) \to \frac{1}{2\pi}\int_{-\pi}^\pi \log f_\theta(\omega)d\omega\) 且 \(\frac{1}{n}\mathbf{X}^\top\mathbf{\Sigma}_\theta^{-1}\mathbf{X} \to \frac{1}{2\pi}\int_{-\pi}^\pi \frac{I_n(\omega)}{f_\theta(\omega)}d\omega\)。代入高斯对数似然 \(-2\ell(\theta) = \log\det(\mathbf{\Sigma}_\theta) + \mathbf{X}^\top\mathbf{\Sigma}_\theta^{-1}\mathbf{X}\)，忽略 \(O(1/n)\) 误差即得 Whittle 近似似然。频域计算可利用 FFT 降至 \(O(n\log n)\)。\(\blacksquare\)

\hrulefill

\hrulefill

\hrulefill

\hrulefill

\hrulefill

\hrulefill

\newpage

\chapter{10-概率与统计/03-统计学/03-统计学习理论.md}\label{ux6982ux7387ux4e0eux7edfux8ba103-ux7edfux8ba1ux5b6603-ux7edfux8ba1ux5b66ux4e60ux7406ux8bba.md}

\subsection{303.A 统计学习补充}\label{a-ux7edfux8ba1ux5b66ux4e60ux8865ux5145}

\begin{quote}
本卷从概率论、泛函分析和优化理论的视角，系统阐述统计学习理论的核心数学结构------PAC可学性、一致收敛、函数逼近和随机优化。内容涵盖PAC学习与VC维（概率论与组合数学）、神经网络的通用逼近定理（泛函分析中稠密性定理的特例）、随机梯度下降的收敛理论（非凸优化）、以及最优传输（Monge-Kantorovich问题与Brenier定理）。
\end{quote}

\hrulefill

\hrulefill

\hrulefill

\hrulefill

\hrulefill

\hrulefill

\section{第369章：统计学习理论}\label{ux7b2c369ux7ae0ux7edfux8ba1ux5b66ux4e60ux7406ux8bba}

统计学习理论由 Vapnik 和 Chervonenkis 在 1970 年代创立，是理解从有限样本中学习的数学可能性和极限的严格框架。其核心问题看似简单却极为深刻：在什么条件下，有限训练样本足以保证学习算法对未见数据具有泛化能力？VC 维（Vapnik-Chervonenkis 维数）为这一问题提供了精确的组合刻画------假设类 PAC 可学习当且仅当其 VC 维有限，这一结果将学习理论从工程直觉提升为数学定理。统计学习理论不仅是机器学习的数学基石（支持向量机直接源于 VC 理论的结构风险最小化原则），更与概率论（一致大数定律、McDiarmid 不等式）、组合学（Sauer 引理、增长函数）和泛函分析（函数空间的度量熵）深刻交织。本章从 PAC 学习框架（Valiant, 1984）出发，建立 VC 维理论、Rademacher 复杂度和偏差-方差分解的完整体系，并引入过参数化时代的新型泛化理论------双重下降现象。

\subsection{187.1 PAC学习框架}\label{pacux5b66ux4e60ux6846ux67b6}

\textbf{定义 187.1}（PAC学习，Probably Approximately Correct，Valiant 1984）：设\(\mathcal{H}\)是假设类，\(\mathcal{D}\)是未知分布。算法\(\mathcal{A}\)从i.i.d.样本\(S \sim \mathcal{D}^m\)中输出\(h_S \in \mathcal{H}\)。\(\mathcal{H}\)是\textbf{PAC可学习}的，如果存在函数\(m_{\mathcal{H}}: (0,1)^2 \to \mathbb{N}\)和算法\(\mathcal{A}\)，对任意\(\epsilon, \delta \in (0,1)\)和任意分布\(\mathcal{D}\)，当\(m \geq m_{\mathcal{H}}(\epsilon, \delta)\)时，

\[\mathbb{P}_{S \sim \mathcal{D}^m}[L_{\mathcal{D}}(h_S) \leq \min_{h \in \mathcal{H}} L_{\mathcal{D}}(h) + \epsilon] \geq 1 - \delta\]

其中\(L_{\mathcal{D}}(h) = \mathbb{P}_{(x,y) \sim \mathcal{D}}[h(x) \neq y]\)是泛化误差。

\textbf{定义 187.2}（样本复杂度）：\(m_{\mathcal{H}}(\epsilon, \delta)\)是\(\mathcal{H}\)的\textbf{样本复杂度}------达到\((\epsilon, \delta)\)-PAC学习所需的最小样本数。

\textbf{定理 187.1}（有限假设类的PAC可学习性）：若\(\mathcal{H}\)有限，则\(m_{\mathcal{H}}(\epsilon, \delta) \leq \frac{\log |\mathcal{H}| + \log(1/\delta)}{\epsilon^2}\)。即\(\mathcal{H}\)是PAC可学习的，样本复杂度为\(O(\epsilon^{-2} \log |\mathcal{H}|)\)。

\emph{证明}：使用Hoeffding不等式和联合界。对每个\(h \in \mathcal{H}\)，\(|L_{\mathcal{D}}(h) - L_S(h)| \leq \epsilon\)的概率\(\geq 1 - 2e^{-2m\epsilon^2}\)。联合界给出\(\mathbb{P}[\exists h: |L_{\mathcal{D}}(h) - L_S(h)| > \epsilon] \leq 2|\mathcal{H}| e^{-2m\epsilon^2}\)。令此概率\(\leq \delta\)，解出\(m \geq \frac{\log(2|\mathcal{H}|) + \log(1/\delta)}{2\epsilon^2} = O(\epsilon^{-2} \log |\mathcal{H}|)\)。∎

\subsection{187.2 VC维}\label{vcux7ef4}

\textbf{定义 187.3}（VC维，Vapnik-Chervonenkis维数）：二分类假设类\(\mathcal{H}\)的\textbf{VC维}\(\operatorname{VCdim}(\mathcal{H})\)是能被\(\mathcal{H}\)打散（shatter）的最大点集的大小。即存在\(d\)个点，\(\mathcal{H}\)在这\(d\)个点上能实现所有\(2^d\)种标注方式。

\textbf{定理 187.2}（VC维基本定理 / Vapnik-Chervonenkis定理）：假设类\(\mathcal{H}\)是PAC可学习的当且仅当其VC维有限。此时样本复杂度为

\[m_{\mathcal{H}}(\epsilon, \delta) = O\left(\frac{\operatorname{VCdim}(\mathcal{H}) + \log(1/\delta)}{\epsilon^2}\right)\]

\emph{注}：VC维是判断假设类可学习性的精确刻画。有限VC维等价于一致大数定律成立。

\textbf{证明}：分两个方向。（\(\Leftarrow\)）若\(\operatorname{VCdim}(\mathcal{H}) = d < \infty\)，由Sauer引理，增长函数\(\Pi_{\mathcal{H}}(m) \leq (em/d)^d\)。由Vapnik-Chervonenkis不等式，以概率至少\(1-\delta\)，\(\sup_{h \in \mathcal{H}} |L_{\mathcal{D}}(h) - L_S(h)| \leq \sqrt{\frac{8d\log(m+1) + 8\log(4/\delta)}{m}}\)。令此\(\leq \epsilon\)，解出\(m = O(\frac{d + \log(1/\delta)}{\epsilon^2})\)，故\(\mathcal{H}\)是PAC可学习的。（\(\Rightarrow\)）若\(\mathcal{H}\)是PAC可学习的，则存在一致大数定律：\(\sup_{h \in \mathcal{H}} |L_{\mathcal{D}}(h) - L_S(h)| \to 0\)（依概率）。由Vapnik-Chervonenkis的不可学习性定理，这蕴含\(\operatorname{VCdim}(\mathcal{H}) < \infty\)。\(\blacksquare\)

\textbf{定理 187.3}（Sauer引理）：设\(\operatorname{VCdim}(\mathcal{H}) = d\)。则\(\mathcal{H}\)在\(m\)个点上的增长函数满足

\[\Pi_{\mathcal{H}}(m) \leq \sum_{i=0}^d \binom{m}{i} \leq \left(\frac{em}{d}\right)^d\]

Sauer引理是VC维有限\(\Rightarrow\)PAC可学习这一方向的关键组合工具。

\textbf{证明}：对\(m\)归纳。\(m \leq d\)时，\(\Pi_{\mathcal{H}}(m) \leq 2^m = \sum_{i=0}^m \binom{m}{i}\)，成立。设对\(m-1\)成立。取\(m\)个点组成的集合\(S\)，固定一点\(x \in S\)。考虑\(\mathcal{H}\)在\(S\)上的限制\(\mathcal{H}|_S\)。将\(\mathcal{H}|_S\)中的函数按\(x\)处的取值配对：若两个函数在\(S\setminus\{x\}\)上相同但在\(x\)处取值不同，则它们配对。设\(\mathcal{H}'\)为去掉\(x\)后不能配对的函数构成的假设类，则\(\operatorname{VCdim}(\mathcal{H}') \leq d-1\)（否则\(x\)与打散\(d\)个点的集合合并后打散\(d+1\)个点）。因此\(\Pi_{\mathcal{H}}(m) \leq \Pi_{\mathcal{H}}(m-1) + \Pi_{\mathcal{H}'}(m-1) \leq \sum_{i=0}^d \binom{m-1}{i} + \sum_{i=0}^{d-1} \binom{m-1}{i} = \sum_{i=0}^d \binom{m}{i} \leq (em/d)^d\)。最后一个不等式由\(\binom{m}{i} \leq (em/i)^i\)及\(i \leq d\)求和得到。\(\blacksquare\)

\subsection{187.3 Rademacher复杂度}\label{rademacherux590dux6742ux5ea6}

\textbf{定义 187.4}（经验Rademacher复杂度）：给定样本\(S = (z_1, \ldots, z_m)\)，函数类\(\mathcal{F}\)的\textbf{经验Rademacher复杂度}为

\[\widehat{\mathcal{R}}_S(\mathcal{F}) = \mathbb{E}_{\sigma} \left[\sup_{f \in \mathcal{F}} \frac{1}{m} \sum_{i=1}^m \sigma_i f(z_i)\right]\]

其中\(\sigma_i \sim \operatorname{Uniform}(\{-1, +1\})\)是独立Rademacher随机变量。Rademacher复杂度衡量函数类拟合随机噪声的能力。

\textbf{定理 187.4}（Rademacher复杂度泛化界）：对任意\(\delta > 0\)，以概率至少\(1 - \delta\)，

\[\sup_{f \in \mathcal{F}} \left|\mathbb{E}[f(Z)] - \frac{1}{m} \sum_{i=1}^m f(z_i)\right| \leq 2 \mathcal{R}_m(\mathcal{F}) + \sqrt{\frac{\log(2/\delta)}{2m}}\]

其中\(\mathcal{R}_m(\mathcal{F}) = \mathbb{E}_S[\widehat{\mathcal{R}}_S(\mathcal{F})]\)是Rademacher复杂度。Rademacher复杂度给出了比VC维更精细的泛化界，因为它依赖于数据分布。

\textbf{证明}：定义\(\Phi(S) = \sup_{f \in \mathcal{F}} |\mathbb{E}[f] - \frac{1}{m}\sum_i f(z_i)|\)。当替换一个样本\(z_i\)时，\(\Phi(S)\)至多改变\(2\|f\|_\infty/m \leq 2/m\)（假设\(f\)取值于\([0,1]\)）。由McDiarmid不等式，以概率至少\(1-\delta\)，\(\Phi(S) \leq \mathbb{E}_S[\Phi(S)] + \sqrt{\frac{\log(2/\delta)}{2m}}\)。由对称化引理，\(\mathbb{E}_S[\Phi(S)] \leq 2 \mathcal{R}_m(\mathcal{F})\)：引入独立副本\(S'\)，则\(\mathbb{E}_S[\sup_f |\mathbb{E}[f] - \hat{\mathbb{E}}_S[f]|] \leq \mathbb{E}_{S,S'}[\sup_f |\hat{\mathbb{E}}_{S'}[f] - \hat{\mathbb{E}}_S[f]|] = \mathbb{E}_{S,S',\sigma}[\sup_f |\frac{1}{m}\sum_i \sigma_i(f(z_i') - f(z_i))|] \leq 2\mathcal{R}_m(\mathcal{F})\)。\(\blacksquare\)

\textbf{定理 187.5}（Massart有限类引理）：对有限函数类\(\mathcal{F}\)，\(\mathcal{R}_m(\mathcal{F}) \leq \sqrt{2 \log |\mathcal{F}| / m}\)。

\textbf{证明}：设\(\mathcal{F} = \{f_1, \ldots, f_N\}\)。经验Rademacher复杂度为\(\widehat{\mathcal{R}}_S(\mathcal{F}) = \mathbb{E}_\sigma[\max_j \frac{1}{m}\sum_i \sigma_i f_j(z_i)]\)。对每个\(j\)，\(X_j = \frac{1}{m}\sum_i \sigma_i f_j(z_i)\)是亚Gauss随机变量，方差代理为\(\frac{1}{m^2}\sum_i f_j(z_i)^2 \leq 1/m\)（假设\(|f| \leq 1\)）。由亚Gauss随机变量的极大不等式，\(\mathbb{E}_\sigma[\max_j X_j] \leq \sqrt{2\log N / m}\)。取期望\(\mathbb{E}_S\)即得\(\mathcal{R}_m(\mathcal{F}) \leq \sqrt{2\log|\mathcal{F}|/m}\)。\(\blacksquare\)

\subsection{187.4 偏差-方差分解}\label{ux504fux5dee-ux65b9ux5deeux5206ux89e3}

\textbf{定义 187.5}（偏差-方差分解）：对平方损失，期望泛化误差分解为

\[\mathbb{E}[(h_S(x) - y)^2] = \underbrace{(\mathbb{E}[h_S(x)] - \mathbb{E}[y|x])^2}_{\text{偏差}^2} + \underbrace{\operatorname{Var}(h_S(x))}_{\text{方差}} + \underbrace{\operatorname{Var}(y|x)}_{\text{不可约误差}}\]

偏差-方差权衡是理解过拟合与欠拟合的核心数学框架。偏差度量学习算法对真实函数的系统性偏离，方差度量算法对训练样本随机波动的敏感度。

\textbf{定理 187.6}（双重下降现象，Belkin et al.~2019）：在过参数化模型中，测试误差随模型复杂度呈现''双重下降''：先下降（经典偏差-方差权衡），再上升（过拟合），然后再次下降（过参数化区域）。这一现象表明经典偏差-方差权衡在过参数化体制下需要修正，其数学根源在于插值估计量的隐式正则化性质。

\textbf{证明}：考虑线性模型\(y = X\beta^* + \varepsilon\)，其中\(X \in \mathbb{R}^{n \times p}\)。在欠参数化区域（\(p < n\)），最小二乘估计\(\hat{\beta} = (X^\top X)^{-1}X^\top y\)的预测风险为\(\mathbb{E}\|\hat{y} - X\beta^*\|^2 = \sigma^2 p\)。在过参数化区域（\(p > n\)），取最小范数插值解\(\hat{\beta} = X^\top (X X^\top)^{-1} y\)，其预测风险为\(\mathbb{E}\|\hat{y} - X\beta^*\|^2 = \sigma^2 n + \|\beta^*\|^2 \cdot \operatorname{Tr}((I - XX^\top(XX^\top)^{-1})X X^\top)\)。当\(p \to n\)（插值阈值）时，\(X X^\top\)的最小特征值趋于零，风险发散。当\(p \gg n\)时，由随机矩阵理论，\(X X^\top/p \to I_n\)，风险趋于\(\sigma^2 n\)。这解释了测试误差在\(p=n\)附近达到峰值、在\(p \gg n\)时再次下降的双重下降现象。\(\blacksquare\)

\hrulefill

\hrulefill

\hrulefill

\hrulefill

\hrulefill

\hrulefill

\section{第370章：神经网络逼近理论}\label{ux7b2c370ux7ae0ux795eux7ecfux7f51ux7edcux903cux8fd1ux7406ux8bba}

神经网络逼近理论从泛函分析的视角研究一个核心问题：由特定激活函数生成的有限宽、有限深神经网络的函数空间在目标函数空间中的稠密性及逼近速率。Cybenko 的通用逼近定理（1989 年）------单隐层非多项式激活网络在连续函数空间中稠密------虽然奠定了神经网络的理论合法性，但并未回答效率问题。真正深刻的是 Barron 定理（1993 年）：对 Fourier 表示光滑的函数，两层网络的 \(L^2\) 逼近速率为 \(O(N^{-1/2})\) 且与输入维数 \(d\) 无关------这是神经网络克服''维数灾难''的数学根源。深度分离定理（Telgarsky 2016, Eldan-Shamir 2016）则进一步证明深层网络在表达能力上对浅层网络具有超多项式优势。神经正切核（NTK, Jacot et al.~2018）在无限宽度极限下将非凸神经网络训练等价为核回归，架起了深度学习与经典的统计学习理论之间的桥梁。本章在统计学习理论（第 304 章）的基础上，以逼近理论和 NTK 为两大支柱。

\subsection{188.1 通用逼近定理}\label{ux901aux7528ux903cux8fd1ux5b9aux7406}

\textbf{定理 188.1}（Cybenko通用逼近定理，1989）：设\(\sigma\)是连续的非多项式激活函数。则单隐层神经网络\(\sum_{i=1}^N a_i \sigma(w_i^\top x + b_i)\)在\(C([0,1]^d)\)中稠密（在一致范数下）。

\textbf{证明}：设\(\mathcal{M}\)为形如\(\sum_{i=1}^N a_i \sigma(w_i^\top x + b_i)\)的函数全体在\(C([0,1]^d)\)中的闭包。若\(\mathcal{M} \neq C([0,1]^d)\)，则由Hahn-Banach定理，存在非零有界线性泛函\(L \in C([0,1]^d)^*\)使得\(L|_{\mathcal{M}} = 0\)。由Riesz表示定理，存在有限Borel测度\(\mu\)使得\(L(f) = \int f d\mu\)。特别地，对所有\(w, b\)，\(\int \sigma(w^\top x + b) d\mu(x) = 0\)。若\(\sigma\)是非多项式，则\(\sigma\)的有限线性组合的集合在\(L^2\)中稠密，由此推出\(\mu\)的所有矩为零，故\(\mu = 0\)，矛盾。因此\(\mathcal{M} = C([0,1]^d)\)。\(\blacksquare\)

Cybenko 定理确立了神经网络的理论合法性------任何连续函数原则上都可以由一个单隐层网络任意精度地逼近。但它并未回答一个至关重要的效率问题：逼近的速率如何？需要多少神经元才能达到给定精度？早期的实验已经观察到神经网络在克服''维数灾难''方面比传统逼近格式具有经验优势，而 Barron（1993 年）首次为这一优势提供了数学证明------对于 Fourier 表示具有光滑性的函数，两层网络的 \(L^2\) 逼近速率 \(O(N^{-1/2})\) 与输入维数 \(d\) 无关。

\textbf{定理 188.2}（Barron定理，1993）：具有Fourier表示的函数\(f(x) = \int_{\mathbb{R}^d} e^{i\omega \cdot x} \tilde{f}(\omega) d\omega\)，若\(\int \|\omega\| |\tilde{f}(\omega)| d\omega \leq C\)，则可用\(N\)个隐层神经元的网络以\(O(N^{-1/2})\)的速率逼近（在\(L^2\)范数下），与输入维数\(d\)无关。

\textbf{证明}：将Fourier表示写为\(f(x) = \int_{\mathbb{R}^d} \cos(\omega \cdot x + \theta(\omega)) |\tilde{f}(\omega)| d\omega\)，其中\(\theta(\omega)\)为相位。由\(\|f\|_{\mathcal{B}} = \int \|\omega\| |\tilde{f}(\omega)| d\omega \leq C\)，可将积分视为概率测度\(p(\omega) = \|\omega\| |\tilde{f}(\omega)| / C\)下的期望。Monte Carlo抽样：从\(p\)中独立抽取\(\omega_1, \ldots, \omega_N\)，构造\(f_N(x) = \frac{C}{N} \sum_{i=1}^N \frac{\cos(\omega_i \cdot x + \theta(\omega_i))}{\|\omega_i\|}\)。由Jensen不等式和独立性的协方差估计，\(\mathbb{E}\|f - f_N\|_{L^2}^2 \leq \frac{C^2}{N} \int \frac{|\tilde{f}(\omega)|^2}{p(\omega)} d\omega = \frac{C}{N} \|f\|_{\mathcal{B}} \leq \frac{C^2}{N}\)。故\(L^2\)逼近速率为\(O(N^{-1/2})\)，与\(d\)无关。\(\blacksquare\)

Barron 的分析方法自然引出了一个函数空间的精确定义：能够以与维数无关的速率被两层网络逼近的函数构成了 Barron 空间 \(\mathcal{B}\)。该空间的特征化------通过 Fourier 表示的加权 Sobolev 型条件 \(\int_{\mathbb{R}^d} \|\omega\| |\tilde{f}(\omega)| d\omega < \infty\)------为理解神经网络在何种目标函数上能够高效逼近提供了精确的数学语言。

\textbf{定义 188.1}（Barron空间）：\textbf{Barron空间}\(\mathcal{B}\)由满足\(\|f\|_{\mathcal{B}} = \int_{\mathbb{R}^d} \|\omega\| |\tilde{f}(\omega)| d\omega < \infty\)的函数组成。神经网络在Barron空间上具有\(O(N^{-1/2})\)的逼近速率，该速率与输入维数\(d\)无关------这是克服维数灾难的关键数学性质。

\subsection{188.2 深度的优势}\label{ux6df1ux5ea6ux7684ux4f18ux52bf}

上述结果均针对两层（浅层）神经网络。那么增加网络的深度------引入多个隐藏层------是否能从根本上提升表达能力？深度分离定理（Telgarsky 2016 年，Eldan-Shamir 2016 年）从数学上给出了肯定的回答：存在可由深度为 \(L\)、宽度为多项式的 ReLU 网络高效表示的函数，而任何深度为 \(L-1\) 的浅层网络需要指数级宽度才能达到相同精度。

\textbf{定理 188.3}（深度分离定理，Telgarsky 2016, Eldan-Shamir 2016）：存在可由深度为\(L\)、宽度为多项式大小的ReLU网络表示的函数，但任何深度为\(L-1\)的网络需要指数级宽度才能达到相同的逼近精度。这从数学上证明了深层网络在表达能力上对浅层网络具有超多项式优势。

\textbf{证明}：构造锯齿函数族\(\Delta_k(x)\)：\(\Delta_0(x) = x\)，\(\Delta_{k+1}(x) = \Delta_k(2x) - \Delta_k(2x-1)\)。\(\Delta_k\)是周期为\(2^{-k}\)的锯齿波，由\(k\)层ReLU网络计算。定义\(f_L(x) = \Delta_L(x)\)，可由深度\(L\)、宽度\(O(1)\)的ReLU网络表示。但任何深度\(L-1\)的网络表示\(f_L\)需要\(2^{\Omega(L)}\)个神经元。原因：\(f_L\)有\(2^{L-1}\)个线性区域，而深度\(L-1\)、宽度\(W\)的ReLU网络最多有\(O(W^{L-1})\)个线性区域。要逼近\(f_L\)的\(2^{L-1}\)个区域，需\(W = 2^{\Omega(L/(L-1))} = 2^{\Omega(1)}\)在\(L\)固定时，或需要\(W\)随\(L\)指数增长。\(\blacksquare\)

Telgarsky 的构造使用了锯齿函数的迭代复合来证明深度分离，而 Poggio 等人（2017 年）从另一个角度展示了深度的优势：某些简单函数------如 \(d\) 个变量的乘积 \(\prod_{i=1}^d x_i\)------可以通过深度为 \(O(\log d)\) 的二进制树结构用多项式宽度网络高效表示，但恒定深度的浅层网络需要指数级宽度。这一结果直观地揭示了深度网络''组合性逼近''的本质。

\textbf{定理 188.4}（Poggio等，2017）：深度ReLU网络可以指数级更高效地表示某些函数。具体地，函数\(f(x_1, \ldots, x_d) = \prod_{i=1}^d x_i\)可由深度\(O(\log d)\)的ReLU网络表示，但浅层网络需要指数级宽度。

\textbf{证明}：利用二进制树结构。对\(d\)个变量，逐对相乘：\(x_1 x_2, x_3 x_4, \ldots\)，然后对结果再次逐对相乘，\(\log_2 d\)层后得\(\prod_{i=1}^d x_i\)。每层使用ReLU网络的乘法门：\(xy = \frac{1}{4}((x+y)^2 - (x-y)^2)\)，而平方函数\(z \mapsto z^2\)可由两个ReLU神经元在紧集上任意精度逼近：\(z^2 \approx \max(0, z)^2 + \max(0, -z)^2 - 2\max(0,z)\max(0,-z)\)。总深度为\(O(\log d)\)，宽度为\(O(d)\)（多项式）。浅层网络：若深度固定为\(L\)，表示\(d\)次多项式需要宽度\(\Omega(2^{d/L})\)，当\(d\)大时指数级。\(\blacksquare\)

前面两个定理通过具体函数的构造证明了深度的优势是严格存在的。Yarotsky（2017 年）进一步从逼近论的角度给出了深度网络逼近 \(C^s\) 光滑函数的最优速率：深度为 \(L\) 的 ReLU 网络逼近 \(C^s\) 函数的 \(L^2\) 误差下界为 \(O(L^{-2s/d})\)，且该速率是最优的。这一结果将深度网络的逼近能力置于经典的逼近论框架中，为深度学习提供了微积分层面的数学保障。

\textbf{定理 188.5}（Yarotsky定理，2017）：深度\(L\)的ReLU网络逼近\(C^s\)函数的最优速率为\(O(L^{-2s/d})\)。深度网络通过组合性实现高效逼近------函数的层次化复合结构对应网络的逐层表示。

\textbf{证明}：核心是分片多项式逼近。给定\(f \in C^s([0,1]^d)\)，将定义域划分为\(2^{d\ell}\)个边长为\(2^{-\ell}\)的子立方体。在每个子立方体上，用\(f\)的\(s\)阶Taylor展开（\(s\)次多项式）逼近，误差为\(O(2^{-\ell s})\)。使用ReLU网络实现\(s\)次多项式乘法，需要\(O(s)\)层。对于深度\(L\)的网络，可取\(\ell \approx L/s\)层用于划分，总逼近误差为\(O(2^{-L s / d}) = O(L^{-2s/d})\)（在\(L^2\)范数下）。此速率是最优的，由\(C^s\)函数的非线性\(n\)-宽度下界保证：任何具有\(N\)个参数的连续参数化逼近格式的误差下界为\(\Omega(N^{-s/d})\)，而深度\(L\)的网络有\(O(L)\)个参数层次。\(\blacksquare\)

\subsection{188.3 神经正切核}\label{ux795eux7ecfux6b63ux5207ux6838}

\textbf{定义 188.2}（神经正切核，Neural Tangent Kernel，Jacot et al.~2018）：考虑参数为\(\theta\)的神经网络\(f(x; \theta)\)。\textbf{神经正切核}定义为

\[\Theta(x, x') = \mathbb{E}_{\theta \sim \mathcal{N}(0, I)} \left[\nabla_\theta f(x; \theta)^\top \nabla_\theta f(x'; \theta)\right]\]

在无限宽度极限下，NTK在训练过程中保持不变------这是理解神经网络梯度下降动力学的核心数学工具。

\textbf{定理 188.6}（NTK极限下的梯度下降，Jacot et al.~2018）：在无限宽度极限下，梯度下降训练神经网络等价于关于NTK的核回归：

\[f_t(x) \approx \Theta(x, X)^\top \Theta(X, X)^{-1} (I - e^{-t \Theta(X, X)}) Y\]

其中\(\Theta(X, X)_{ij} = \Theta(x_i, x_j)\)。这给出了神经网络训练的精确解析描述，将非凸神经网络优化问题转化为线性核方法。

\textbf{证明}：在无限宽度极限下，由中心极限定理，网络输出\(f(x;\theta)\)的分布收敛到均值为零、协方差为\(\Theta\)的Gauss过程。梯度下降的连续时间极限为\(\frac{d\theta_t}{dt} = -\eta \nabla_\theta \mathcal{L}(\theta_t)\)，其中\(\mathcal{L}\)为经验损失。在无限宽度下，参数的变化是\(O(1/\sqrt{N})\)，故\(\theta_t \approx \theta_0\)。因此\(f_t(x) \approx f_0(x) + \nabla_\theta f(x;\theta_0)^\top (\theta_t - \theta_0)\)。代入\(\theta_t - \theta_0 = -\eta \int_0^t \nabla_\theta \mathcal{L}(\theta_s) ds\)，利用链式法则和NTK的定义，得\(\frac{d f_t(x)}{dt} = -\eta \Theta(x, X) \nabla_f \mathcal{L}(f_t(X))\)。对平方损失\(\mathcal{L}(f) = \frac{1}{2}\|f - Y\|^2\)，解得\(f_t(X) = (I - e^{-\eta t \Theta(X,X)})Y\)。代入插值公式得\(f_t(x) = \Theta(x,X)^\top \Theta(X,X)^{-1}(I - e^{-\eta t \Theta(X,X)})Y\)。\(\blacksquare\)

\textbf{定理 188.7}（NTK的谱性质）：NTK的特征值衰减速率决定了神经网络的学习速度。在均匀分布数据上，NTK的特征值随频率增加而衰减，因此梯度下降倾向于先学习低频函数------这一频谱偏差现象由NTK的Mercer特征分解自然导出。

\textbf{证明}：由Mercer定理，\(\Theta(x,x') = \sum_{k=1}^\infty \lambda_k \phi_k(x) \phi_k(x')\)，其中\(\lambda_k \geq 0\)为特征值，\(\phi_k\)为特征函数。在NTK回归中，\(f_t(x) = \sum_k \phi_k(x) \langle \phi_k, f^* \rangle (1 - e^{-\eta t \lambda_k})\)。对于Laplace核或与ReLU相关的NTK，特征值\(\lambda_k\)按\(\lambda_k \sim k^{-(d+1)/d}\)衰减（在\(d\)维球面上）。因此，与较大\(\lambda_k\)对应的低频成分以速率\(O(e^{-\eta t \lambda_k})\)被学习，而高频成分的学习速率指数级慢。这解释了梯度下降的频谱偏差：网络先拟合数据的光滑（低频）趋势，后拟合高频细节。\(\blacksquare\)

\hrulefill

\hrulefill

\hrulefill

\hrulefill

\hrulefill

\hrulefill

\section{第371章：深度学习的优化理论}\label{ux7b2c371ux7ae0ux6df1ux5ea6ux5b66ux4e60ux7684ux4f18ux5316ux7406ux8bba}

深度学习的核心悖论之一在于：随机梯度下降（SGD）在高达数百万维的非凸参数空间中以惊人的效率收敛到具有卓越泛化能力的解，而经典的非凸优化理论对此几乎无法解释。这一悖论驱动了过去十年深度学习优化理论的快速发展。从收敛性理论角度看，SGD 在非凸光滑损失函数上保证梯度的平均范数以 \(O(1/\sqrt{T})\) 速率收敛到零（即到达驻点），但损失景观的几何分析揭示了一个更深刻的事实：与经典非凸优化中令人恐惧的鞍点问题不同，过参数化网络的损失景观中，高损失局部极小值以指数概率消失，所有局部极小值集中在全局极小值附近且构成连通区域------这是神经网络在实践中高度可训练性的数学根源。隐式正则化理论（Soudry et al., 2018）进一步指出，SGD 本身无须显式正则化就隐式地偏好最大间隔解，将优化算法与泛化理论统一在同一框架下。

\subsection{189.1 随机梯度下降的收敛理论}\label{ux968fux673aux68afux5ea6ux4e0bux964dux7684ux6536ux655bux7406ux8bba}

\textbf{定义 189.1}（随机梯度下降，SGD）：设\(\ell(\theta) = \mathbb{E}_{z \sim \mathcal{D}}[\ell(\theta; z)]\)是期望风险。SGD更新为

\[\theta_{t+1} = \theta_t - \eta_t \nabla_\theta \ell(\theta_t; z_t)\]

其中\(z_t \sim \mathcal{D}\)是随机采样的样本，\(\eta_t\)是步长。

SGD 的收敛性取决于损失函数 \(\ell(\theta)\) 的凸性。当 \(\ell\) 是凸函数时，SGD 在期望意义下以 \(O(\log T / \sqrt{T})\) 的速率收敛到全局最优。这一速率与被加项 \(1/\sqrt{t}\) 的衰减步长策略密切相关------递减步长在收敛速度和噪声平滑之间达成平衡。

\textbf{定理 189.1}（SGD的收敛性，凸情形）：若\(\ell(\theta)\)是凸函数且梯度有界，\(\eta_t = 1/\sqrt{t}\)，则

\[\mathbb{E}\left[\ell\left(\frac{1}{T}\sum_{t=1}^T \theta_t\right)\right] - \ell(\theta^*) \leq O\left(\frac{\log T}{\sqrt{T}}\right)\]

\textbf{证明}：由凸性，\(\ell(\theta_t) - \ell(\theta^*) \leq \nabla\ell(\theta_t)^\top (\theta_t - \theta^*)\)。设\(\|\nabla\ell(\theta_t; z_t)\| \leq G\)。由SGD更新，\(\|\theta_{t+1} - \theta^*\|^2 = \|\theta_t - \theta^*\|^2 - 2\eta_t \nabla\ell(\theta_t; z_t)^\top (\theta_t - \theta^*) + \eta_t^2 \|\nabla\ell(\theta_t; z_t)\|^2\)。取条件期望，\(\mathbb{E}[\|\theta_{t+1} - \theta^*\|^2 \mid \theta_t] \leq \|\theta_t - \theta^*\|^2 - 2\eta_t (\ell(\theta_t) - \ell(\theta^*)) + \eta_t^2 G^2\)。整理并对\(t=1,\ldots,T\)求和，\(\sum_{t=1}^T 2\eta_t \mathbb{E}[\ell(\theta_t) - \ell(\theta^*)] \leq \|\theta_1 - \theta^*\|^2 + G^2 \sum_{t=1}^T \eta_t^2\)。取\(\eta_t = 1/\sqrt{t}\)，则\(\sum \eta_t \sim 2\sqrt{T}\)，\(\sum \eta_t^2 \sim \log T\)。由Jensen不等式，\(\mathbb{E}[\ell(\bar{\theta}_T)] - \ell(\theta^*) \leq \frac{\|\theta_1 - \theta^*\|^2 + G^2 \log T}{2\sqrt{T}} = O(\log T / \sqrt{T})\)。\(\blacksquare\)

凸性的假设极大地简化了收敛分析，但深度神经网络的目标函数是非凸的。在非凸情形下，经典的收敛保证降至驻点收敛------SGD 保证梯度的平均平方范数以 \(O(1/\sqrt{T})\) 的速率衰减。然而，这仅说明算法能找到一个梯度很小的点，但该点可能是鞍点、局部极大值或局部极小值------非凸优化的核心挑战正在于此。

\textbf{定理 189.2}（SGD的收敛性，非凸情形）：若\(\ell(\theta)\)光滑（\(L\)-光滑），则SGD收敛到驻点：

\[\frac{1}{T} \sum_{t=1}^T \mathbb{E}[\|\nabla \ell(\theta_t)\|^2] \leq O\left(\frac{1}{\sqrt{T}}\right)\]

在非凸情形下，SGD保证梯度的平均范数趋于零，但收敛到的驻点可能是鞍点而非局部极小值。

\textbf{证明}：由\(L\)-光滑性，\(\ell(\theta_{t+1}) \leq \ell(\theta_t) + \nabla\ell(\theta_t)^\top (\theta_{t+1} - \theta_t) + \frac{L}{2}\|\theta_{t+1} - \theta_t\|^2\)。代入SGD更新\(\theta_{t+1} - \theta_t = -\eta_t \nabla\ell(\theta_t; z_t)\)并取期望，\(\mathbb{E}[\ell(\theta_{t+1})] \leq \mathbb{E}[\ell(\theta_t)] - \eta_t \mathbb{E}[\|\nabla\ell(\theta_t)\|^2] + \frac{L\eta_t^2}{2} \mathbb{E}[\|\nabla\ell(\theta_t; z_t)\|^2]\)。设\(\mathbb{E}[\|\nabla\ell(\theta_t; z_t)\|^2] \leq \sigma^2 + \|\nabla\ell(\theta_t)\|^2\)（有界方差假设）。取\(\eta_t = \eta = 1/\sqrt{T}\)，求和得\(\eta \sum_{t=1}^T \mathbb{E}[\|\nabla\ell(\theta_t)\|^2] \leq \ell(\theta_1) - \ell^* + \frac{L\eta^2 T}{2}(\sigma^2 + \sup_t \mathbb{E}[\|\nabla\ell(\theta_t)\|^2])\)。整理得\(\frac{1}{T}\sum_t \mathbb{E}[\|\nabla\ell(\theta_t)\|^2] \leq \frac{\ell(\theta_1) - \ell^*}{\eta T} + \frac{L\eta\sigma^2}{2} = O(1/\sqrt{T})\)。\(\blacksquare\)

\subsection{189.2 损失景观的几何结构}\label{ux635fux5931ux666fux89c2ux7684ux51e0ux4f55ux7ed3ux6784}

\textbf{定义 189.2}（损失景观，Loss Landscape）：深度神经网络的损失函数\(\ell(\theta)\)在高维参数空间中的几何结构。关键数学特征包括鞍点的分布、局部极小值的质量、以及全局极小值的连通性。

经典的非凸优化理论将高维非凸景观中的局部极小值和鞍点视为主要障碍。但深度神经网络的经验成功暗示其损失景观具有特殊的几何性质。Choromanska 等人（2015 年）利用随机矩阵理论为这一直觉提供了数学支撑：在大规模网络中，高损失的局部极小值以指数概率消失，所有局部极小值集中在全局极小值附近。

\textbf{定理 189.3}（Choromanska等，2015的随机矩阵模型）：在大规模随机神经网络中，利用随机矩阵理论可以证明：损失景观的局部极小值集中在全局极小值附近，且高损失的局部极小值随网络规模增大以指数概率消失。

\textbf{证明}：将随机神经网络的损失函数写为\(\ell(\theta) = \frac{1}{m}\sum_{i=1}^m \ell_i(\theta)\)，其中\(\ell_i\)对应第\(i\)个训练样本。在随机初始化下，\(\ell_i(\theta)\)可视为关联的随机过程。Hessian矩阵\(H = \nabla^2 \ell(\theta)\)可分解为\(H = H_0 + \Delta\)，其中\(H_0\)为正定部分（由随机矩阵极限给出），\(\Delta\)为波动。由随机矩阵理论，在临界点处\(H\)的特征值分布由Wigner半圆律加上秩1修正给出。当\(m\)（网络宽度）大时，\(\Delta\)的谱半径收缩为\(O(1/\sqrt{m})\)，故\(H\)以高概率正定，即临界点为局部极小值。进一步，所有局部极小值的目标函数值集中在全局极小值附近，高损失极小值的概率按\(e^{-c m}\)衰减。\(\blacksquare\)

Choromanska 等人的随机矩阵分析考虑了非线性网络中关联的随机过程带来的复杂性。一个更简单但富有启发性的起点是深度线性网络------去掉激活函数的网络。Kawaguchi（2016 年）证明了一个优美的结果：深度线性网络的损失景观中，每一个局部极小值都是全局极小值。这一结果暗示非线性网络的''陷阱''来自激活函数引入的非线性。

\textbf{定理 189.4}（Kawaguchi定理，2016）：深度线性网络（无激活函数）的每个局部极小值都是全局极小值。对于非线性网络，该性质在过参数化区域中近似成立------这一结果的证明依赖于对Hessian矩阵谱结构在过参数化极限下的分析。

\textbf{证明}：深度线性网络为\(f(x) = W_L W_{L-1} \cdots W_1 x\)。损失函数\(\ell(W_1,\ldots,W_L) = \|W_L\cdots W_1 X - Y\|_F^2\)。令\(W = W_L \cdots W_1\)，则\(\ell\)仅依赖于乘积\(W\)。设\((W_1,\ldots,W_L)\)为局部极小值。若\(\nabla_{W_i} \ell = 0\)对所有\(i\)成立，则由链式法则，\(W_{i+1}^\top \cdots W_L^\top \nabla_W \ell \cdot W_1^\top \cdots W_{i-1}^\top = 0\)对所有\(i\)。若\(W\)非满秩，可通过构造下降方向得矛盾。若\(W\)满秩，则\(\nabla_W \ell = 0\)，即\(W\)为全局极小值\(W^*\)。因此所有局部极小值对应的\(W\)均满足\((W X - Y)X^\top = 0\)，即\(W\)为最小二乘解，故为全局极小值。\(\blacksquare\)

以上两个结果分别从随机矩阵和代数角度揭示了过参数化网络损失景观的良性结构。Freeman 和 Bruna（2017 年）进一步发现了一个引人注目的几何性质：全局极小值之间并非孤立，而是通过低损失路径相互连通------这一''模式连通性''（mode connectivity）意味着 SGD 即使在高度非线性景观中也有连续路径从一个解游走到另一个解而不遭遇高损失障碍。

\textbf{定理 189.5}（Freeman-Bruna定理，2017）：在过参数化条件下，深度ReLU网络的损失景观在全局极小值附近是连通的（mode connectivity），即任意两个全局极小值可以通过损失值不显著升高的路径相连。

\textbf{证明}：设\(\theta_a, \theta_b\)为两个全局极小值。在过参数化下，网络有冗余参数。构造路径\(\theta(\lambda) = \sqrt{1-\lambda}\theta_a + \sqrt{\lambda}\theta_b\)（\(\lambda \in [0,1]\)）不是直接线性插值（因参数空间非凸），而是通过网络的中间表示插值。利用过参数化性质：存在参数化下的零空间，在其中移动参数不改变网络输出。具体地，对每层加扰动\(\Delta W_i\)满足\(W_{i+1}\Delta W_i + \Delta W_{i+1}W_i = 0\)（至一阶）。沿此零空间方向，可将\(\theta_a\)连续变形为\(\theta_b\)而不显著增加损失。更精确地，构造Bezier曲线\(\theta(t) = (1-t)^2\theta_a + 2t(1-t)\theta_m + t^2\theta_b\)，其中\(\theta_m\)在过参数化空间的零流形上选取，使得沿路径的损失上界为\(O(\max\{\ell(\theta_a),\ell(\theta_b)\} + \varepsilon)\)。\(\blacksquare\)

\subsection{189.3 隐式正则化的数学分析}\label{ux9690ux5f0fux6b63ux5219ux5316ux7684ux6570ux5b66ux5206ux6790}

\textbf{定义 189.3}（隐式正则化，Implicit Regularization）：SGD本身（无显式正则化）倾向于收敛到具有特定数学结构的解------这一现象的数学分析构成了隐式正则化理论。

上述损失景观的分析论证了全局极小值的存在性与连通性，但未解释 SGD 如何在众多可能的全局极小值之间做出选择。隐式正则化理论（Soudry 等人 2018 年）为这一选择提供了数学解释：即便不加显式的 \(\ell_2\) 正则化，梯度下降在可分离数据的线性分类中自动收敛到最大间隔分类器------等价于硬间隔 SVM 的解。

\textbf{定理 189.6}（Soudry等，2018的隐式偏差）：在可分离数据的线性分类中，梯度下降（带指数尾损失）在极限下收敛到最大间隔分类器，等价于硬间隔SVM的解。这从数学上证明了梯度下降隐式地执行了\(\ell_2\)正则化。

\textbf{证明}：考虑线性分类器\(f(x) = w^\top x\)，损失为\(\ell(y f(x))\)，其中\(\ell\)为指数尾损失（如指数损失或logistic损失）。梯度下降更新为\(w_{t+1} = w_t - \eta \nabla_w \mathcal{L}(w_t)\)。在可分离数据下，\(\|w_t\| \to \infty\)且\(w_t/\|w_t\|\)收敛。令\(w_\infty = \lim_{t\to\infty} w_t/\|w_t\|\)。由梯度下降的KKT条件，\(w_\infty\)满足\(\min_i y_i w_\infty^\top x_i = 1\)且\(w_\infty\)为满足\(\|w\|=1\)约束下最大化间隔的方向。具体地，将梯度下降写为\(\frac{dw}{dt} = -\sum_i \ell'(y_i w^\top x_i) y_i x_i\)。在\(t \to \infty\)时，\(\ell'(y_i w^\top x_i)\)对间隔最小的样本贡献主导，迫使\(w\)对齐最大间隔方向。由此\(w_\infty\)与硬间隔SVM的解\(\arg\min_w \|w\|^2\) s.t. \(y_i w^\top x_i \geq 1\)一致。\(\blacksquare\)

Soudry 等人的结果针对的是线性分类器。在矩阵分解中，Gunasekar 等人（2017 年）发现了一个类似的隐式正则化现象：梯度下降在过参数化的矩阵分解中隐式地极小化核范数（即迹范数），从而自动偏向低秩解。这一发现将优化算法的隐式偏差与压缩感知和低秩矩阵补全的核心正则化原则建立了桥梁。

\textbf{定理 189.7}（Gunasekar等，2017）：在矩阵分解问题中，梯度下降隐式地最小化核范数，从而倾向于低秩解。这一结论将梯度下降的隐式偏差与迹范数正则化建立了精确的数学联系。

\textbf{证明}：考虑矩阵分解\(X = UV^\top\)，其中\(U \in \mathbb{R}^{n \times r}\)，\(V \in \mathbb{R}^{m \times r}\)。梯度下降最小化\(\ell(U,V) = \|UV^\top - M\|_F^2\)。在过参数化下（\(r > \operatorname{rank}(M)\)），无穷多\(U,V\)满足\(UV^\top = M\)。梯度下降的连续时间极限为\(\dot{U} = -\nabla_U \ell\)，\(\dot{V} = -\nabla_V \ell\)。观察到\(\frac{d}{dt}(U^\top U - V^\top V) = 0\)，故\(U^\top U - V^\top V\)为运动常数。在适当初始化条件下（\(U_0, V_0\)平衡），梯度下降收敛到最小核范数解：\(\|UV^\top\|_* = \frac{1}{2}(\|U\|_F^2 + \|V\|_F^2)\)。核范数最小化促进低秩性，因为\(\|X\|_* = \sum_i \sigma_i(X)\)在奇异值上施加\(\ell_1\)型惩罚。\(\blacksquare\)

从矩阵分解的隐式低秩偏好自然地引出一个更深的问题：对于更深层网络，梯度下降的隐式偏差如何随层数变化？Arora 等人（2019 年）揭示了层数 \(L\) 在隐式正则化中的关键作用------梯度下降极小化的并非核范数，而是 Schatten \(2/L\)-拟范数，该拟范数随 \(L\) 增大而趋近于秩函数，为深层网络天然的秩缩减现象提供了精确的数学解释。

\textbf{定理 189.8}（Arora等，2019）：在深度线性网络中，梯度下降隐式地倾向于低秩解，且收敛方向与初始化相关。对于深度矩阵分解，梯度下降在参数空间的轨迹可由特定的Riemannian梯度流近似。

\textbf{证明}：考虑深度矩阵分解\(W = W_L W_{L-1} \cdots W_1\)。梯度下降隐式地极小化\(\|W\|_{2/L}^{2/L}\)（Schatten \(2/L\)-拟范数）。当\(L \to \infty\)时，该拟范数趋近于秩函数。具体地，梯度下降的连续时间极限在参数空间上诱导出Riemannian度量。令\(W(t) = W_L(t) \cdots W_1(t)\)。由链式法则，\(\frac{dW}{dt} = \sum_{i=1}^L (W_L \cdots W_{i+1}) \dot{W}_i (W_{i-1} \cdots W_1)\)。代入\(\dot{W}_i = -\nabla_{W_i} \ell\)，得\(\frac{dW}{dt} = -\sum_{i=1}^L P_i P_i^\top \nabla_W \ell\)，其中\(P_i = W_L \cdots W_{i+1}\)。该流在\(W\)的奇异值上诱导出特定的演化速率，使得较小奇异值衰减更快，从而隐式偏向低秩。初始化\(W_i(0) \approx \alpha I\)时，该偏好最显著。\(\blacksquare\)

\textbf{统计学习理论的现代发展与深度学习理论}

统计学习理论在 21 世纪经历了从经典 VC 理论到深度学习理论的范式转变，其核心挑战是理解为什么过参数化的深度神经网络能够泛化------这一现象与经典理论（如奥卡姆剃刀和 VC 维上界）的预测矛盾。\textbf{双重下降现象}（double descent, Belkin et al.~2019）是深度学习理论中最引人注目的发现之一：在经典统计学习理论中，测试误差作为模型复杂度的函数呈 U 形曲线（偏差-方差权衡）；但在深度学习中，测试误差在插值阈值（模型参数数 = 训练样本数）之后再次下降，形成''双重下降''曲线。这一现象表明过参数化（overparameterization）并非有害，反而可能通过隐式正则化（implicit regularization）提升泛化能力。\textbf{神经切线核}（Neural Tangent Kernel, NTK, Jacot-Gabriel-Hongler 2018）提供了理解无限宽神经网络的理论工具：在无限宽极限下，随机初始化的神经网络等效于具有 NTK 的高斯过程，梯度下降在函数空间中的动力学由 NTK 的线性化方程描述，收敛到 NTK 核回归器的解。NTK 理论解释了为什么超宽神经网络可以在训练误差为零的情况下实现良好的泛化，但未能完全解释有限宽度的优越性。\textbf{随机特征模型}（Rahimi-Recht 2007）和\textbf{卷积神经切线核}（CNTK, Arora et al.~2019）将 NTK 理论推广到卷积结构，揭示了卷积归纳偏差（inductive bias）对泛化性能的贡献。\textbf{均值场理论}（mean-field theory, Mei-Montanari-Nguyen 2018, Chizat-Bach 2018）将有限宽神经网络的训练动力学描述为概率测度空间上的 Wasserstein 梯度流，为理解特征学习（feature learning）------NTK 范式无法解释的深度学习的核心优势------提供了数学框架。 最优传输理论起源于Monge（1781）提出的土壤搬运问题，经Kantorovich（1942）赋予线性规划形式，如今已成为连接偏微分方程、概率论和几何的深刻数学理论。本章聚焦于Monge-Kantorovich问题、Brenier定理及其与Monge-Ampère方程的联系。

\subsection{190.1 Monge问题与Kantorovich松弛}\label{mongeux95eeux9898ux4e0ekantorovichux677eux5f1b}

\textbf{定义 190.1}（Monge最优传输问题，1781）：给定概率测度\(\mu, \nu\)在\(\mathbb{R}^d\)上，\textbf{Monge问题}寻求代价最小的传输映射\(T\)：

\[\min_{T: T_\#\mu = \nu} \int_{\mathbb{R}^d} c(x, T(x)) d\mu(x)\]

其中\(T_\#\mu = \nu\)表示\(T\)将\(\mu\)推前（push-forward）为\(\nu\)。对二次代价\(c(x,y) = \|x-y\|^2\)，称为\textbf{二次最优传输问题}。

\textbf{定义 190.2}（Kantorovich松弛，1942）：Monge问题的Kantorovich松弛将确定性的传输映射推广为概率耦合：

\[\min_{\gamma \in \Pi(\mu, \nu)} \int_{\mathbb{R}^d \times \mathbb{R}^d} c(x, y) d\gamma(x, y)\]

其中\(\Pi(\mu, \nu)\)是所有边缘分布分别为\(\mu\)和\(\nu\)的联合分布（耦合）的集合。这是Monge问题的凸松弛------将非凸约束\(T_\#\mu = \nu\)替换为凸约束\(\gamma \in \Pi(\mu, \nu)\)。

\textbf{定理 190.1}（Kantorovich对偶定理）：Kantorovich问题具有以下对偶形式：

\[\min_{\gamma \in \Pi(\mu, \nu)} \int c(x,y) d\gamma = \sup_{\varphi, \psi} \left\{\int \varphi d\mu + \int \psi d\nu : \varphi(x) + \psi(y) \leq c(x,y)\right\}\]

对二次代价\(c(x,y) = \|x-y\|^2\)，对偶问题可化为：

\[W_2^2(\mu, \nu) = \sup_{\varphi \text{ convex}} \left\{\int \varphi^* d\nu - \int \varphi d\mu\right\}\]

其中\(\varphi^*(y) = \sup_x \{\langle x, y \rangle - \varphi(x)\}\)是\(\varphi\)的Legendre-Fenchel共轭。

\textbf{证明}：由Kantorovich对偶的一般形式，\(\min_{\gamma \in \Pi} \int c d\gamma = \sup_{\varphi,\psi} \{\int \varphi d\mu + \int \psi d\nu : \varphi(x) + \psi(y) \leq c(x,y)\}\)。对二次代价\(c(x,y) = \|x-y\|^2\)，展开得\(\|x-y\|^2 = \|x\|^2 + \|y\|^2 - 2\langle x,y\rangle\)。令\(\varphi(x) = \frac{1}{2}\|x\|^2 - u(x)\)，\(\psi(y) = \frac{1}{2}\|y\|^2 - v(y)\)，约束\(\varphi(x)+\psi(y) \leq \|x-y\|^2\)变为\(u(x)+v(y) \geq \langle x,y\rangle\)。取\(v(y) = u^*(y) = \sup_x\{\langle x,y\rangle - u(x)\}\)，则可证明\(u\)为凸函数。对偶目标化为\(\sup_{u \text{ convex}} \{\int u^* d\nu - \int u d\mu\}\)，即\(W_2^2(\mu,\nu) = \sup_{\varphi \text{ convex}} \{\int \varphi^* d\nu - \int \varphi d\mu\}\)。\(\blacksquare\)

\subsection{190.2 Brenier定理与Monge-Ampère方程}\label{brenierux5b9aux7406ux4e0emonge-ampuxe8reux65b9ux7a0b}

\textbf{定理 190.2}（Brenier定理，1991）：设\(\mu, \nu\)是\(\mathbb{R}^d\)上的概率测度，\(\mu\)关于Lebesgue测度绝对连续，且其二阶矩有限。则对二次代价\(c(x,y) = \|x-y\|^2\)：

\begin{enumerate}
\def\labelenumi{\arabic{enumi}.}
\tightlist
\item
  存在唯一的（\(\mu\)-几乎处处）最优传输映射\(T = \nabla\varphi\)，其中\(\varphi: \mathbb{R}^d \to \mathbb{R}\)是凸函数（Brenier势）。
\item
  \(T\)是Monge问题和Kantorovich问题的共同最优解。
\item
  \(\nabla\varphi\)关于\(\mu\)几乎处处是单值的。
\end{enumerate}

Brenier定理将最优传输问题归约为寻找凸函数\(\varphi\)的问题------\(\varphi\)满足\textbf{Monge-Ampère方程}（在适当的意义下）：

\textbf{证明}：设\(\gamma\)为Kantorovich问题的最优耦合。由Kantorovich对偶，存在凸函数\(\varphi\)使得\(\varphi(x) + \varphi^*(y) = \|x-y\|^2/2\)在\(\operatorname{supp}(\gamma)\)上成立。由此\(y = x - \nabla\varphi(x)\)（由Legendre-Fenchel共轭的性质），故\(T(x) = x - \nabla\varphi(x) = \nabla(\|x\|^2/2 - \varphi(x))\)。令\(\tilde{\varphi}(x) = \|x\|^2/2 - \varphi(x)\)，则\(T = \nabla\tilde{\varphi}\)且\(\tilde{\varphi}\)为凸函数（因为\(\varphi\)为凸且\(\|x\|^2/2\)严格凸）。唯一性：若\(T_1, T_2\)均为最优传输映射，则\(\int \|T_1 - T_2\|^2 d\mu = 0\)（由最优性），故\(T_1 = T_2\) \(\mu\)-a.e.。Monge-Ampère方程由变量替换公式\(\nu = T_\#\mu\)导出：对任意测试函数\(\xi\)，\(\int \xi(y) d\nu(y) = \int \xi(T(x)) d\mu(x) = \int \xi(\nabla\varphi(x)) d\mu(x)\)。若\(\mu\)和\(\nu\)有密度，则\(f(x) = g(\nabla\varphi(x)) \det(D^2\varphi(x))\)。\(\blacksquare\)

\[\det(D^2 \varphi(x)) = \frac{d\mu}{d\nu \circ \nabla\varphi}(x)\]

当\(\mu\)和\(\nu\)都有密度\(f\)和\(g\)时，若\(\nabla\varphi\)是\(C^1\)微分同胚，则形式上有：

\[\det(D^2 \varphi(x)) \cdot g(\nabla\varphi(x)) = f(x)\]

这是一类完全非线性椭圆型偏微分方程（Monge-Ampère型），其正则性理论（Caffarelli 1990-1996）是现代偏微分方程研究的重大成果。

\subsection{190.3 Wasserstein距离的几何}\label{wassersteinux8dddux79bbux7684ux51e0ux4f55}

\textbf{定义 190.3}（\(p\)-Wasserstein距离）：对\(p \geq 1\)，概率测度\(\mu, \nu\)之间的\textbf{\(p\)-Wasserstein距离}（地球移动距离）定义为

\[W_p(\mu, \nu) = \left(\min_{\gamma \in \Pi(\mu, \nu)} \int \|x - y\|^p d\gamma(x, y)\right)^{1/p}\]

Wasserstein距离在具有有限\(p\)阶矩的概率测度空间\(\mathcal{P}_p(\mathbb{R}^d)\)上定义了度量，使\((\mathcal{P}_p(\mathbb{R}^d), W_p)\)成为测地空间------其测地线与经典的Riemannian几何测地线不同，由位移插值（McCann插值）给出。

\textbf{定理 190.3}（Kantorovich-Rubinstein对偶，\(p=1\)）：\(W_1\)距离具有以下对偶表示：

\[W_1(\mu, \nu) = \sup_{\|f\|_{\text{Lip}} \leq 1} \left\{\int f d\mu - \int f d\nu\right\}\]

这是Wasserstein-1距离的Kantorovich-Rubinstein对偶定理，将最优传输问题转化为在1-Lipschitz函数空间上的优化问题。

\textbf{证明}：由Kantorovich对偶的一般形式，取\(c(x,y) = \|x-y\|\)。则\(\min_{\gamma \in \Pi} \int \|x-y\| d\gamma = \sup_{\varphi,\psi} \{\int \varphi d\mu + \int \psi d\nu : \varphi(x) + \psi(y) \leq \|x-y\|\}\)。取\(\psi = -\varphi\)，约束变为\(\varphi(x) - \varphi(y) \leq \|x-y\|\)，即\(\varphi\)为1-Lipschitz函数。对偶目标化为\(\sup_{\|\varphi\|_{\text{Lip}} \leq 1} \{\int \varphi d\mu - \int \varphi d\nu\}\)。\(W_1\)距离由此对偶形式给出，其几何意义为：\(W_1\)等于在1-Lipschitz函数空间上\(\mu\)和\(\nu\)的积分差的上确界。这也是积分概率度量（IPM）框架下Wasserstein距离与GAN的判别器网络之间的深层联系。\(\blacksquare\)

\textbf{定理 190.4}（Sinkhorn算法与熵正则化，Cuturi 2013）：加熵正则化的最优传输问题（Sinkhorn距离）：

\[W_\varepsilon(\mu, \nu) = \min_{\gamma \in \Pi(\mu, \nu)} \int \|x - y\|^2 d\gamma + \varepsilon \operatorname{KL}(\gamma \| \mu \otimes \nu)\]

可通过Sinkhorn迭代求解，复杂度\(O(n^2)\)，远优于精确最优传输的\(O(n^3 \log n)\)。熵正则化使问题在计算上可处理，同时保持了Wasserstein度量的核心几何性质。

\textbf{证明}：熵正则化问题为\(\min_{\gamma \in \Pi} \langle C, \gamma \rangle + \varepsilon \sum_{i,j} \gamma_{ij} \log \gamma_{ij}\)，其中\(C_{ij} = \|x_i - y_j\|^2\)。Lagrangian为\(\mathcal{L} = \langle C, \gamma \rangle + \varepsilon \sum \gamma_{ij} \log \gamma_{ij} + \lambda^\top (\gamma \mathbf{1} - \mu) + \eta^\top (\gamma^\top \mathbf{1} - \nu)\)。一阶条件给出\(\gamma_{ij} = \exp(-\frac{C_{ij}}{\varepsilon} + \frac{\lambda_i + \eta_j}{\varepsilon} - 1) = u_i K_{ij} v_j\)，其中\(K_{ij} = e^{-C_{ij}/\varepsilon}\)，\(u_i = e^{\lambda_i/\varepsilon - 1/2}\)，\(v_j = e^{\eta_j/\varepsilon - 1/2}\)。由边缘约束\(u \odot (K v) = \mu\)，\(v \odot (K^\top u) = \nu\)。这是Sinkhorn不动点迭代：\(u^{(k+1)} = \mu \oslash (K v^{(k)})\)，\(v^{(k+1)} = \nu \oslash (K^\top u^{(k+1)})\)。由Sinkhorn定理，此迭代线性收敛。每次迭代\(O(n^2)\)，总复杂度\(O(n^2/\varepsilon)\)。\(\blacksquare\)

\hrulefill

\hrulefill

\hrulefill

\hrulefill

\hrulefill

\hrulefill

\section{第372章：序列决策的数学理论}\label{ux7b2c372ux7ae0ux5e8fux5217ux51b3ux7b56ux7684ux6570ux5b66ux7406ux8bba}

序列决策（sequential decision-making）是强化学习与最优控制共同的数学语言，其公理化框架由 Bellman（1957 年）的马尔可夫决策过程（MDP）奠定。MDP 的深刻之处在于将未来不确定性的折现归约为当前状态的值函数和 Q 函数，这种递归结构催生了 Bellman 最优方程和值迭代/策略迭代等核心算法。序列决策理论与统计学习理论（第 304 章）之间存在本质的联系：在线学习中的遗憾（regret）最小化、多臂老虎机（Lai-Robbins 1985 下界、UCB 算法）和 Thompson 采样，均可视为特殊的 MDP，而 Information Bottleneck 原理（Tishby et al., 1999）提供了一种信息论视角理解探索-利用权衡------最优策略在最大化奖励与最小化与环境的相关信息之间寻求平衡。本章从 MDP 的数学定义出发，系统建立 Bellman 方程、Dynamic Programming 和 Information Bottleneck 方法，为强化学习的数学理论（第 307 章补充）提供基础框架。

\subsection{191.1 马尔可夫决策过程与Bellman方程}\label{ux9a6cux5c14ux53efux592bux51b3ux7b56ux8fc7ux7a0bux4e0ebellmanux65b9ux7a0b}

\textbf{定义 191.1}（马尔可夫决策过程，MDP）：\textbf{MDP}由五元组\((\mathcal{S}, \mathcal{A}, P, R, \gamma)\)定义： - \(\mathcal{S}\)：状态空间（可测空间） - \(\mathcal{A}\)：动作空间 - \(P(s'|s, a)\)：转移概率核 - \(R(s, a)\)：即时奖励函数 - \(\gamma \in [0, 1)\)：折扣因子

MDP是受控Markov过程的数学抽象：在每个时间步，决策者观察状态\(s\)，选择动作\(a\)，获得奖励\(R(s,a)\)，系统按\(P(\cdot|s,a)\)转移到新状态。

\textbf{定义 191.2}（价值函数与Q函数）：策略\(\pi(a|s)\)的\textbf{状态价值函数}和\textbf{动作价值函数}为

\[V^\pi(s) = \mathbb{E}_\pi\left[\sum_{t=0}^\infty \gamma^t R(s_t, a_t) \mid s_0 = s\right]\]

\[Q^\pi(s, a) = \mathbb{E}_\pi\left[\sum_{t=0}^\infty \gamma^t R(s_t, a_t) \mid s_0 = s, a_0 = a\right]\]

\textbf{定理 191.1}（Bellman期望方程）：\(V^\pi\)和\(Q^\pi\)满足以下不动点方程：

\[V^\pi(s) = \mathbb{E}_{a \sim \pi(\cdot|s)}[R(s, a) + \gamma \mathbb{E}_{s' \sim P}[V^\pi(s')]]\]

\[Q^\pi(s, a) = R(s, a) + \gamma \mathbb{E}_{s' \sim P}[\mathbb{E}_{a' \sim \pi}[Q^\pi(s', a')]]\]

Bellman方程表明价值函数是Bellman算子的唯一不动点，折扣因子\(\gamma\)保证了算子的压缩性（Banach不动点定理）。

\textbf{证明}：由\(V^\pi\)的定义，\(V^\pi(s) = \mathbb{E}_\pi[\sum_{t=0}^\infty \gamma^t R(s_t, a_t) \mid s_0 = s] = \mathbb{E}_{a \sim \pi(\cdot|s)}[R(s,a)] + \gamma \mathbb{E}_\pi[\sum_{t=0}^\infty \gamma^t R(s_{t+1}, a_{t+1}) \mid s_0 = s]\)。由Markov性和条件期望的塔性质，第二项\(= \mathbb{E}_{a \sim \pi(\cdot|s)}[\gamma \mathbb{E}_{s' \sim P(\cdot|s,a)}[V^\pi(s')]]\)。同理，\(Q^\pi(s,a) = R(s,a) + \gamma \mathbb{E}_{s' \sim P}[\mathbb{E}_{a' \sim \pi}[Q^\pi(s',a')]]\)。定义Bellman算子\(\mathcal{T}^\pi V(s) = \mathbb{E}_{a \sim \pi}[R(s,a) + \gamma \mathbb{E}_{s'}[V(s')]]\)。对任意\(V_1, V_2\)，\(\|\mathcal{T}^\pi V_1 - \mathcal{T}^\pi V_2\|_\infty \leq \gamma \|V_1 - V_2\|_\infty\)，故\(\mathcal{T}^\pi\)是\(\gamma\)-压缩的。由Banach不动点定理，\(V^\pi\)是唯一不动点。\(\blacksquare\)

\textbf{定理 191.2}（Bellman最优性方程）：最优价值函数\(V^*(s) = \sup_\pi V^\pi(s)\)满足

\[V^*(s) = \max_{a \in \mathcal{A}} \left[R(s, a) + \gamma \mathbb{E}_{s' \sim P}[V^*(s')]\right]\]

最优Bellman算子也是\(\gamma\)-压缩的（在一致范数下），因此\(V^*\)存在且唯一，相应的贪心策略即为最优策略。

\textbf{证明}：定义最优Bellman算子\(\mathcal{T}^* V(s) = \max_{a \in \mathcal{A}} [R(s,a) + \gamma \mathbb{E}_{s' \sim P}[V(s')]]\)。对任意\(V_1, V_2\)，\(|\mathcal{T}^* V_1(s) - \mathcal{T}^* V_2(s)| \leq \max_a |\gamma \mathbb{E}_{s'}[V_1(s') - V_2(s')]| \leq \gamma \|V_1 - V_2\|_\infty\)，故\(\mathcal{T}^*\)是\(\gamma\)-压缩的。由Banach不动点定理，存在唯一不动点\(V^*\)。现证\(V^* = \sup_\pi V^\pi\)。一方面，对任意\(\pi\)，\(V^\pi \leq \mathcal{T}^* V^\pi\)（因\(\pi\)可能非最优），迭代得\(V^\pi \leq V^*\)。另一方面，取贪心策略\(\pi^*(s) = \arg\max_a [R(s,a) + \gamma \mathbb{E}_{s'}[V^*(s')]]\)，则\(V^{\pi^*} = \mathcal{T}^{\pi^*} V^{\pi^*} = \mathcal{T}^* V^{\pi^*}\)，故\(V^{\pi^*} = V^*\)。因此\(\pi^*\)为最优策略。\(\blacksquare\)

\subsection{191.2 策略梯度定理}\label{ux7b56ux7565ux68afux5ea6ux5b9aux7406}

\textbf{定义 191.3}（策略梯度定理，Sutton et al.~1999）：参数化策略\(\pi_\theta\)的性能\(J(\theta) = \mathbb{E}_{s_0 \sim \rho}[V^{\pi_\theta}(s_0)]\)的梯度为

\[\nabla_\theta J(\theta) = \mathbb{E}_{\pi_\theta}\left[\nabla_\theta \log \pi_\theta(a|s) \cdot Q^{\pi_\theta}(s, a)\right]\]

策略梯度定理将期望回报对策略参数的梯度表达为对数策略梯度的加权期望，其数学基础是测度变换下的Score Function恒等式：\(\nabla_\theta \pi_\theta = \pi_\theta \nabla_\theta \log \pi_\theta\)。

策略梯度定理的证明利用了轨迹概率分布的测度变换。设轨迹 \(\tau = (s_0, a_0, s_1, a_1, \ldots)\) 的联合分布为 \(p_\theta(\tau) = \rho(s_0) \prod_{t=0}^\infty \pi_\theta(a_t | s_t) P(s_{t+1} | s_t, a_t)\)，则回报期望 \(J(\theta) = \int R(\tau) p_\theta(\tau) d\tau\)。利用对数导数技巧：\(\nabla_\theta p_\theta = p_\theta \nabla_\theta \log p_\theta\)，得到 \(\nabla_\theta J(\theta) = \mathbb{E}_{\tau \sim p_\theta}[R(\tau) \nabla_\theta \log p_\theta(\tau)]\)。由于转移概率 \(P\) 不依赖于 \(\theta\)，\(\nabla_\theta \log p_\theta(\tau) = \sum_{t=0}^\infty \nabla_\theta \log \pi_\theta(a_t | s_t)\)。代入即得 \(\nabla_\theta J(\theta) = \mathbb{E}_{\pi_\theta}[\sum_{t=0}^\infty \nabla_\theta \log \pi_\theta(a_t | s_t) R(\tau)]\)。利用 \(Q^{\pi_\theta}(s,a) = \mathbb{E}[R(\tau) | s_0 = s, a_0 = a]\) 和因果性（未来动作不影响过去奖励），可化简为 \(\nabla_\theta J(\theta) = \mathbb{E}_{\pi_\theta}[\nabla_\theta \log \pi_\theta(a|s) \cdot Q^{\pi_\theta}(s,a)]\)。REINFORCE 算法（Williams 1992）直接使用此公式进行 Monte Carlo 梯度估计，而 Actor-Critic 方法则通过同时学习策略（Actor）和价值函数（Critic）来降低估计方差。为减少方差，通常引入基线函数 \(b(s)\)：\(\nabla_\theta J(\theta) = \mathbb{E}[\nabla_\theta \log \pi_\theta(a|s) (Q^{\pi_\theta}(s,a) - b(s))]\)，常用的基线是状态值函数 \(V^{\pi_\theta}(s)\)，此时 \(Q^{\pi_\theta} - V^{\pi_\theta}\) 即为优势函数 \(A^{\pi_\theta}(s,a)\)。

\subsection{191.3 信息瓶颈原理}\label{ux4fe1ux606fux74f6ux9888ux539fux7406}

\textbf{定义 191.4}（信息瓶颈，Tishby et al.~1999）：\textbf{信息瓶颈原理}寻求输入随机变量\(X\)的压缩表示\(T\)，使\(T\)保留关于目标变量\(Y\)的最大信息：

\[\min_{p(t|x)} I(X; T) - \beta I(T; Y)\]

其中\(I(\cdot; \cdot)\)是互信息，\(\beta \geq 0\)控制压缩与保留的权衡。当\(\beta \to \infty\)时，\(T\)是\(Y\)的充分统计量；当\(\beta \to 0\)时，\(T\)是\(X\)的极小充分表示。

\textbf{定理 191.4}（信息瓶颈与表示学习，Tishby-Zaslavsky 2015）：深度神经网络的学习过程可被理解为信息瓶颈的迭代优化------隐层逐步压缩输入信息（减少\(I(X; T)\)），同时保留关于标签的信息（维持\(I(T; Y)\)）。这一信息论视角为深度网络的泛化能力提供了基于互信息的数学解释。

\textbf{证明}：将深度网络的第\(k\)层表示\(T_k\)视为Markov链\(Y \to X \to T_1 \to \cdots \to T_L\)中的随机变量。由数据处理不等式，\(I(X; T_{k+1}) \leq I(X; T_k)\)（信息压缩），而\(I(T_k; Y) \leq I(X; Y)\)。在SGD训练中，网络参数\(\theta\)的更新使\(T_k\)逐步逼近信息瓶颈的最优解：\(\min_{p(t|x)} I(X; T) - \beta I(T; Y)\)。训练初期，\(I(T_k; Y)\)快速增加（拟合阶段）；训练后期，\(I(X; T_k)\)逐渐减小（压缩阶段），此即信息瓶颈的两个阶段。互信息在每一层的演化由\(I(X; T_k)\)和\(I(T_k; Y)\)的动力学方程描述，其整体效果等价于在表示空间中执行\(I(X; T)\)的约束优化。\(\blacksquare\)

\textbf{定理 191.5}（互信息的Donsker-Varadhan表示，Belghazi et al.~2018）：互信息可通过以下变分表示来估计：

\[I(X; Z) = \sup_{f: \Omega \to \mathbb{R}} \mathbb{E}_{P_{XZ}}[f(x, z)] - \log \mathbb{E}_{P_X \otimes P_Z}[e^{f(x, z)}]\]

这是Donsker-Varadhan对KL散度的变分表示的直接推论，将互信息的估计转化为可微函数的优化问题。

\textbf{证明}：由Donsker-Varadhan变分公式，KL散度表示为\(\operatorname{KL}(P \| Q) = \sup_{f} \{\mathbb{E}_P[f] - \log \mathbb{E}_Q[e^f]\}\)。取\(P = P_{XZ}\)（联合分布），\(Q = P_X \otimes P_Z\)（边缘乘积），则\(I(X; Z) = \operatorname{KL}(P_{XZ} \| P_X \otimes P_Z) = \sup_f \{\mathbb{E}_{P_{XZ}}[f(x,z)] - \log \mathbb{E}_{P_X \otimes P_Z}[e^{f(x,z)}]\}\)。其中\(f\)可取任意可测函数。该变分表示将互信息估计转化为最大化问题：使用神经网络参数化\(f_\theta\)，通过梯度上升优化\(\theta\)，即可从样本中估计互信息。当\(f\)取最优函数\(f^* = \log(dP_{XZ}/d(P_X \otimes P_Z))\)时达到上确界。\(\blacksquare\)

\hrulefill

\emph{本卷从概率论、泛函分析和优化理论的统一视角，系统阐述了统计学习理论的核心数学结构：PAC可学性与VC维提供了学习的可能性边界（概率论与组合数学）；神经网络逼近理论揭示了有限参数网络的表达能力（泛函分析中的稠密性与逼近速率）；SGD收敛理论和损失景观分析解释了非凸优化的训练动力学；最优传输理论（Monge-Kantorovich-Brenier）建立了概率测度空间上的几何结构；MDP的Bellman方程体系为序列决策提供了不动点理论框架；信息瓶颈原理从信息论角度统一了表示学习的压缩-保留权衡。共5章。}

\emph{本卷在初等统计（V6）和概率论 II（V20）的基础上，建立了数理统计的严格理论框架：估计理论（无偏性、相合性、Cramér-Rao 下界、MLE 的渐近最优性）为参数估计提供了理论基础；Neyman-Pearson 引理和似然比检验为假设检验提供了最优性准则；置信区间与假设检验的对偶关系揭示了统计推断的统一结构；贝叶斯方法通过先验分布融合先验信息，后验分布的渐近正态性（Bernstein-von Mises 定理）保证了与频率学派方法的一致性；线性模型（Gauss-Markov 定理、ANOVA）是应用最广泛的统计工具；多元分析（PCA、LDA）是处理高维数据的标准方法。\textbf{补充部分}将统计推断理论推广至机器学习：PAC学习理论建立了从有限样本到泛化误差的概率界限，VC维理论刻画了假设空间的复杂度，SGD收敛性分析为大规模优化提供了理论保证，Wasserstein距离统一了分布间的最优传输与生成模型，Markov决策过程的Bellman方程与策略梯度定理为强化学习提供了数学基础。}

\emph{卷二十一：统计 II终。}

\emph{卷二十一：统计 II终。}

\emph{卷二十一：统计 II终。}

\emph{卷二十四：统计学终。}

\newpage

\chapter{10-概率与统计/03-统计学/04-现代统计方法.md}\label{ux6982ux7387ux4e0eux7edfux8ba103-ux7edfux8ba1ux5b6604-ux73b0ux4ee3ux7edfux8ba1ux65b9ux6cd5.md}

\chapter{现代统计方法}\label{ux73b0ux4ee3ux7edfux8ba1ux65b9ux6cd5}

\hrulefill

\section{第373章：非参数统计与重抽样}\label{ux7b2c373ux7ae0ux975eux53c2ux6570ux7edfux8ba1ux4e0eux91cdux62bdux6837}

现代统计方法超越了经典参数模型的框架，发展出处理复杂数据结构和高维问题的强大工具。非参数方法和重抽样技术不依赖特定的分布假设，通过数据驱动的方式实现推断，在理论严谨性和实践灵活性之间取得了优雅的平衡。

\subsection{309.1 非参数方法的基本思想}\label{ux975eux53c2ux6570ux65b9ux6cd5ux7684ux57faux672cux601dux60f3}

\textbf{定义 309.1.1（非参数统计）} 非参数统计（nonparametric statistics）是指不假定数据来自有限维参数族分布的一类统计方法。其基本思想是让数据''为自己说话''，通过灵活的模型结构（如无穷维函数空间）自适应地逼近真实的数据生成机制。

\textbf{定义 309.1.2（次序统计量）} 设 \(X_1, X_2, \ldots, X_n\) 为样本。将它们从小到大排序：\(X_{(1)} \leq X_{(2)} \leq \cdots \leq X_{(n)}\)，则 \(X_{(k)}\) 称为第 \(k\) 个次序统计量。次序统计量是非参数推理的基本构件。

\textbf{命题 309.1.3（次序统计量的分布）} 若 \(X_1, \ldots, X_n\) 独立同分布于密度为 \(f\)、分布函数为 \(F\) 的分布，则 \(X_{(k)}\) 的密度为：

\[f_{X_{(k)}}(x) = \frac{n!}{(k-1)! (n-k)!} F(x)^{k-1} (1 - F(x))^{n-k} f(x).\]

特别地，样本中位数 \(X_{(\lceil n/2 \rceil)}\) 的渐近分布为正态分布：\(\sqrt{n}(m_n - m) \xrightarrow{d} \mathcal{N}(0, 1/(4 f(m)^2))\)，其中 \(m = F^{-1}(1/2)\) 为总体中位数。

\textbf{定义 309.1.4（秩与秩检验）} 设 \(X_1, \ldots, X_n\) 为样本。\(X_i\) 的秩 \(R_i\) 是在有序排列中的位置：\(R_i = \sum_{j=1}^n \mathbf{1}_{\{X_j \leq X_i\}}\)。秩统计量是秩的任意函数，它们关于分布单调变换是不变的（分布自由的）。

\textbf{定理 309.1.5（Wilcoxon符号秩检验）} 对于配对样本的对称性检验，Wilcoxon符号秩统计量为 \(W = \sum_{i=1}^n \operatorname{sgn}(Z_i) \cdot R_i^+\)，其中 \(Z_i\) 为配对差值，\(R_i^+\) 为 \(|Z_i|\) 在 \(\{|Z_j|\}\) 中的秩。在零假设（对称分布）下，\(W\) 的精确分布由秩的随机符号配置决定，可通过组合计算求得。对于大样本，\(\frac{W}{\sqrt{n(n+1)(2n+1)/6}}\) 渐近标准正态。

\textbf{证明}：设配对差值 \(Z_1, \ldots, Z_n\) 独立同分布于关于原点对称的连续分布。在零假设下，\(\operatorname{sgn}(Z_i)\) 和 \(|Z_i|\) 是独立的，且 \(\operatorname{sgn}(Z_i)\) 等可能地取 \(\pm 1\)。令 \(R_i^+\) 为 \(|Z_i|\) 的秩。则 \(W = \sum_{i=1}^n \operatorname{sgn}(Z_i) R_i^+\) 在零假设下的分布等价于 \(W = \sum_{i=1}^n \varepsilon_i i\)（重新标号后），其中 \(\varepsilon_i\) 独立取值 \(\pm 1\) 各以概率 \(1/2\)。由此 \(\mathbb{E}[W] = 0\)，\(\operatorname{var}(W) = \sum_{i=1}^n i^2 = n(n+1)(2n+1)/6\)。由独立但不同分布随机变量之和的Lindeberg-Feller中心极限定理可验证Lindeberg条件，从而 \(W / \sqrt{n(n+1)(2n+1)/6} \xrightarrow{d} \mathcal{N}(0,1)\)。\(\blacksquare\)

\textbf{定理 309.1.6（Mann-Whitney U检验）} 对两个独立样本 \(X_1, \ldots, X_m\) 和 \(Y_1, \ldots, Y_n\)，U统计量 \(U = \sum_{i=1}^m \sum_{j=1}^n \mathbf{1}_{\{X_i > Y_j\}}\) 度量了 \(X\) 倾向大于 \(Y\) 的程度。在零假设（两分布相同）下，\(U\) 的均值和方差可精确计算，且 \(U\) 渐近正态分布。

\textbf{证明}：在零假设下 \(X_1,\ldots,X_m\) 和 \(Y_1,\ldots,Y_n\) 独立同分布。U统计量 \(U = \sum_{i=1}^m \sum_{j=1}^n \mathbf{1}_{\{X_i > Y_j\}}\) 的期望为 \(\mathbb{E}[U] = m n P(X > Y) = mn/2\)（由对称性，\(P(X > Y) = P(X < Y) = 1/2\)）。方差可通过U统计量的一般公式计算：\(\operatorname{var}(U) = mn(m+n+1)/12\)。由Hoeffding的U统计量投影定理，\((U - mn/2) / \sqrt{mn(m+n+1)/12} \xrightarrow{d} \mathcal{N}(0,1)\)。\(\blacksquare\)

\subsection{309.2 Bootstrap方法}\label{bootstrapux65b9ux6cd5}

\textbf{定义 309.2.1（Bootstrap原理）} Bootstrap方法由Efron于1979年系统提出，其核心思想是用经验分布 \(\widehat{F}_n\)（将概率质量 \(1/n\) 赋予每个观测值）替代未知的总体分布 \(F\)。Bootstrap样本是从 \(\widehat{F}_n\) 中有放回地抽取的大小为 \(n\) 的样本。

\textbf{定义 309.2.2（Bootstrap估计）} 设 \(\theta = T(F)\) 为感兴趣的参数，\(\widehat{\theta} = T(\widehat{F}_n)\) 为其估计。Bootstrap分布为 \(T(\widehat{F}_n^*)\) 的条件分布（给定原样本），其中 \(\widehat{F}_n^*\) 为Bootstrap样本对应的经验分布。通过Monte Carlo模拟 \(B\) 个Bootstrap样本计算 \(\widehat{\theta}_1^*, \ldots, \widehat{\theta}_B^*\)，以近似Bootstrap分布。

\textbf{定理 309.2.3（Bootstrap的渐近有效性）} 若 \(T\) 在 \(F\) 处是Hadamard可微的，则Bootstrap分布一致地逼近 \(\sqrt{n}(\widehat{\theta} - \theta)\) 的真实抽样分布：

\[\sup_x \left| P^*(\sqrt{n}(\widehat{\theta}^* - \widehat{\theta}) \leq x) - P(\sqrt{n}(\widehat{\theta} - \theta) \leq x) \right| \xrightarrow{P} 0,\]

其中 \(P^*\) 为Bootstrap条件概率，\(P\) 为真实抽样概率。

\textbf{完整证明概要}：由Hadamard可微性，\(\sqrt{n}(T(\widehat{F}_n) - T(F)) \xrightarrow{d} T_F'(\mathbb{G})\)，其中 \(\mathbb{G}\) 为 \(F\)-Brown桥过程。Bootstrap版本满足 \(\sqrt{n}(T(\widehat{F}_n^*) - T(\widehat{F}_n)) \xrightarrow{d^*} T_F'(\mathbb{G})\)（条件于原样本）。若 \(T_F'\) 是线性的或充分光滑，则两者的极限分布相同。核心的论证用到经验过程的Donsker定理和delta方法（泛函形式）。\(\blacksquare\)

\textbf{定理 309.2.4（Bootstrap置信区间的数学原理）} 百分位Bootstrap置信区间（percentile method）和Bootstrap-t置信区间是最常用的构造。

\textbf{百分位Bootstrap区间}：\([\widehat{\theta}_{\alpha/2}^*, \widehat{\theta}_{1-\alpha/2}^*]\)，其中 \(\widehat{\theta}_q^*\) 为Bootstrap分布的 \(q\) 分位数。此方法基于Bootstrap分布的枢轴性质。

\textbf{Bootstrap-t区间}：对每个Bootstrap样本计算 \(t^* = \sqrt{n}(\widehat{\theta}^* - \widehat{\theta}) / \widehat{\operatorname{se}}^*\)，用 \(t^*\) 的经验分位数代替正态分位数。Bootstrap-t区间具有更高阶的精确性：其覆盖误差为 \(O(n^{-1})\) 而非 \(O(n^{-1/2})\)。

\textbf{证明}（Bootstrap-t的高阶精确性）：由Edgeworth展开，真实抽样分布与正态近似的误差为 \(O(n^{-1/2})\)。Bootstrap-t将估计的分布展开式与真实的展开式匹配到 \(O(n^{-1/2})\) 项，使区间端点的覆盖概率误差降至 \(O(n^{-1})\)。确切证明需要验证Edgeworth展开中系数的光滑性和Bootstrap对其的一致性估计。\(\blacksquare\)

\subsection{309.3 Jackknife与交叉验证}\label{jackknifeux4e0eux4ea4ux53c9ux9a8cux8bc1}

\textbf{定义 309.3.1（Jackknife）} Jackknife方法由Quenouille和Tukey发展，系统地从样本中逐一删除观测值来估计偏差和方差。设 \(\widehat{\theta}\) 为基于全部 \(n\) 个观测值的估计，\(\widehat{\theta}_{(-i)}\) 为删除第 \(i\) 个观测值后的估计。则

Jackknife偏差估计：\(\widehat{\operatorname{bias}}_{\text{jack}} = (n - 1)(\overline{\theta}_{(\cdot)} - \widehat{\theta})\)，其中 \(\overline{\theta}_{(\cdot)} = \frac{1}{n} \sum_{i=1}^n \widehat{\theta}_{(-i)}\)。

Jackknife方差估计：\(\widehat{\operatorname{var}}_{\text{jack}} = \frac{n - 1}{n} \sum_{i=1}^n (\widehat{\theta}_{(-i)} - \overline{\theta}_{(\cdot)})^2\)。

\textbf{定义 309.3.2（交叉验证）} \(K\)-折交叉验证将数据随机分为 \(K\) 个等大折。对每折 \(k\)，用其余 \(K-1\) 折数据训练模型，在第 \(k\) 折上计算预测误差。最终估计为各折误差的平均值。取 \(K = n\) 即为留一交叉验证（LOOCV）。

\textbf{定理 309.3.3（交叉验证的模型选择一致性）} 对于线性回归模型，Mallows \(C_p\)、AIC和留一交叉验证在预测误差意义上渐近等价。BIC（贝叶斯信息准则）和删 \(d\) 交叉验证（\(d/n \to 1\)）在模型选择上是一致的（即随 \(n \to \infty\) 以概率 \(1\) 选中真实模型），而AIC和LOOCV是渐近有效的（极小化预测误差）但不一致。

\subsection{309.4 核密度估计与最优窗宽}\label{ux6838ux5bc6ux5ea6ux4f30ux8ba1ux4e0eux6700ux4f18ux7a97ux5bbd}

\textbf{定义 309.4.1（核密度估计）} 设 \(X_1, \ldots, X_n\) 独立同分布于密度 \(f\)。核密度估计为：

\[\widehat{f}_h(x) = \frac{1}{n h} \sum_{i=1}^n K\left(\frac{x - X_i}{h}\right),\]

其中 \(K\) 为核函数（满足 \(\int K = 1\)，\(\int u K(u) du = 0\)），\(h > 0\) 为窗宽（bandwidth）。

\textbf{定理 309.4.2（MISE与渐近性质）} 均方积分误差（MISE）为 \(\operatorname{MISE}(h) = \mathbb{E} \int (\widehat{f}_h(x) - f(x))^2 dx\)。其渐近展开为：

\[\operatorname{AMISE}(h) = \frac{1}{n h} R(K) + \frac{h^4}{4} \mu_2(K)^2 R(f''),\]

其中 \(R(g) = \int g(x)^2 dx\)，\(\mu_2(K) = \int u^2 K(u) du\)。第一项为方差项，第二项为平方偏差项。

\textbf{最优窗宽}：极小化AMISE得 \(h_{\text{opt}} = \left(\frac{R(K)}{\mu_2(K)^2 R(f'')}\right)^{1/5} n^{-1/5}\)。此时达到最优收敛率：

\[\operatorname{MISE}(h_{\text{opt}}) \sim \frac{5}{4} \left(\mu_2(K)^2 R(K)^4 R(f'')\right)^{1/5} n^{-4/5}.\]

\(n^{-4/5}\) 的收敛率是非参数密度估计的极小极大最优率（在Sobolev光滑度为 \(2\) 的情形）。

\textbf{证明}：MISE \(= \int (\mathbb{E}[\widehat{f}_h(x)] - f(x))^2 dx + \int \operatorname{var}(\widehat{f}_h(x)) dx\)。偏差的渐近展开：由Taylor展开，\(\mathbb{E}[\widehat{f}_h(x)] = \int K(u) f(x - hu) du = f(x) + \frac{h^2}{2} \mu_2(K) f''(x) + o(h^2)\)。故 \(\int (\text{bias})^2 dx \sim \frac{h^4}{4} \mu_2(K)^2 R(f'')\)。方差：\(\operatorname{var}(\widehat{f}_h(x)) = \frac{1}{n} \operatorname{var}(K_h(x - X_1)) \sim \frac{1}{nh} f(x) R(K)\)（利用 \(\int K_h(x-y)^2 f(y) dy \sim \frac{1}{h} f(x) R(K)\)）。求和得AMISE公式。对 \(h\) 求导并令导数为零即得最优窗宽 \(h_{\text{opt}}\)。\(\blacksquare\)

\textbf{定理 309.4.3（核密度估计的逐点渐近正态性）} 若 \(f\) 在 \(x\) 处二阶连续可微，\(h = h_n \to 0\) 且 \(n h_n \to \infty\)，则

\[\sqrt{n h} (\widehat{f}_h(x) - f(x) - \frac{h^2}{2} \mu_2(K) f''(x)) \xrightarrow{d} \mathcal{N}(0, f(x) R(K)).\]

\textbf{证明}：\(\widehat{f}_h(x)\) 可写为独立同分布随机变量之和的平均。记 \(Z_{n,i} = \frac{1}{h} K(\frac{x-X_i}{h})\)。则 \(\mathbb{E}[Z_{n,i}] = f(x) + \frac{h^2}{2} \mu_2(K) f''(x) + o(h^2)\)，\(\operatorname{var}(Z_{n,i}) = \frac{1}{h} f(x) R(K) + o(1/h)\)。由Lyapunov中心极限定理（需验证Lyapunov条件：\(n h \to \infty\) 保证三阶矩可忽略），\(\sqrt{nh}(\widehat{f}_h(x) - \mathbb{E}[\widehat{f}_h(x)]) \xrightarrow{d} \mathcal{N}(0, f(x) R(K))\)。加上偏差项即得所述渐近分布。\(\blacksquare\)

\subsection{309.5 经验过程与Dvoretzky-Kiefer-Wolfowitz不等式}\label{ux7ecfux9a8cux8fc7ux7a0bux4e0edvoretzky-kiefer-wolfowitzux4e0dux7b49ux5f0f}

\textbf{定义 309.5.1（经验分布函数与经验过程）} 经验分布函数 \(\mathbb{F}_n(x) = \frac{1}{n} \sum_{i=1}^n \mathbf{1}_{\{X_i \leq x\}}\)。经验过程定义为 \(\sqrt{n}(\mathbb{F}_n(x) - F(x))\)。

\textbf{定理 309.5.2（Donsker定理）} \(\sqrt{n}(\mathbb{F}_n - F)\) 在 \(D[-\infty, \infty]\) 上弱收敛于 \(F\)-Brown桥过程 \(\mathbb{G} \circ F\)，其中 \(\mathbb{G}\) 为标准Brown桥。

\textbf{证明}：经验过程 \(\sqrt{n}(\mathbb{F}_n - F)\) 的有限维分布由多元中心极限定理弱收敛于Gauss过程。紧致性（tightness）由经验过程的经典不等式验证：对任意 \(\varepsilon, \eta > 0\)，存在 \(\delta > 0\) 使得 \(\limsup_{n\to\infty} P(\sup_{|s-t|<\delta} |\sqrt{n}(\mathbb{F}_n(F^{-1}(s)) - s) - \sqrt{n}(\mathbb{F}_n(F^{-1}(t)) - t)| > \varepsilon) < \eta\)。这可由矩不等式或括号熵（bracketing entropy）论证。极限过程 \(\mathbb{G} \circ F\) 的协方差结构为 \(\operatorname{cov}(\mathbb{G}(F(s)), \mathbb{G}(F(t))) = F(s \wedge t) - F(s)F(t)\)，恰为Brown桥。\(\blacksquare\)

\textbf{定理 309.5.3（Dvoretzky-Kiefer-Wolfowitz不等式）} 对任意 \(\varepsilon > 0\)，

\[P\left(\sup_{x \in \mathbb{R}} |\mathbb{F}_n(x) - F(x)| > \varepsilon\right) \leq 2 e^{-2n \varepsilon^2}.\]

此不等式（DKW）是对Glivenko-Cantelli定理的定量强化------不仅给出 \(L^\infty\) 范数的几乎必然收敛，还给出了精确的指数尾概率界。DKW不等式是构造Kolmogorov-Smirnov检验置信带的数学基础。

\textbf{证明概要}：DKW不等式的证明使用了大偏差理论和Massart的精细化。本质上，\(\sup_x |\mathbb{F}_n - F|\) 可化为 \([0,1]\) 上均匀经验过程的 \(L^\infty\) 范数。利用Doob的鞅方法和指数不等式可证其尾概率界。Massart证明了常数 \(2\) 是最优的。\(\blacksquare\)

\hrulefill

\section{第374章：生存分析与可靠性}\label{ux7b2c374ux7ae0ux751fux5b58ux5206ux6790ux4e0eux53efux9760ux6027}

生存分析（Survival Analysis）研究事件发生时间的统计规律，其标志性特征是截断和删失数据的存在。该领域的核心数学工具是计数过程鞅理论，由Aalen于1970年代系统发展。

\subsection{310.1 生存函数与危险率函数}\label{ux751fux5b58ux51fdux6570ux4e0eux5371ux9669ux7387ux51fdux6570}

\textbf{定义 310.1.1（生存函数）} 设 \(T \geq 0\) 为事件发生时间（随机变量），其生存函数定义为：

\[S(t) = P(T > t) = 1 - F(t), \quad t \geq 0.\]

\(S(t)\) 表示个体存活（事件未发生）超过时间 \(t\) 的概率。

\textbf{定义 310.1.2（危险率函数）} 危险率函数（hazard function）或瞬时失效率定义为：

\[\lambda(t) = \lim_{\Delta t \to 0^+} \frac{P(t \leq T < t + \Delta t \mid T \geq t)}{\Delta t} = \frac{f(t)}{S(t)} = -\frac{d}{dt} \log S(t).\]

危险率函数与生存函数之间的关系为 \(S(t) = \exp(-\int_0^t \lambda(u) du)\)。累积危险率 \(\Lambda(t) = \int_0^t \lambda(u) du = -\log S(t)\)。

\textbf{定义 310.1.3（截断与删失）} 右删失（right censoring）：观测到的仅是 \(Y = \min(T, C)\) 和删失指示 \(\delta = \mathbf{1}_{\{T \leq C\}}\)，其中 \(C\) 为删失时间。左截断（left truncation）：仅当 \(T \geq L\) 时个体才进入研究。独立删失假设：\(T \perp C\)（在实际中弱化为条件独立）。

\subsection{310.2 Kaplan-Meier估计与Nelson-Aalen估计}\label{kaplan-meierux4f30ux8ba1ux4e0enelson-aalenux4f30ux8ba1}

\textbf{定义 310.2.1（Kaplan-Meier估计------乘积极限估计）} 设观测到 \(n\) 个个体，以 \(\tilde{t}_1 < \tilde{t}_2 < \cdots < \tilde{t}_k\) 为发生事件的不同时间点。令 \(d_j\) 为在 \(\tilde{t}_j\) 处发生的事件数，\(n_j\) 为在 \(\tilde{t}_j\) 前仍处于风险集中的个体数。Kaplan-Meier生存函数估计为：

\[\widehat{S}(t) = \prod_{j: \tilde{t}_j \leq t} \left(1 - \frac{d_j}{n_j}\right).\]

\textbf{定义 310.2.2（Nelson-Aalen估计）} 累积危险率的Nelson-Aalen估计为：

\[\widehat{\Lambda}(t) = \sum_{j: \tilde{t}_j \leq t} \frac{d_j}{n_j}.\]

\textbf{定理 310.2.3（Kaplan-Meier估计的渐近性质）} 在独立删失假设下，对满足 \(S(t) > 0\) 的 \(t\)，

\[\sqrt{n}(\widehat{S}(t) - S(t)) \xrightarrow{d} \mathcal{N}(0, \sigma^2(t)),\]

其中 \(\sigma^2(t) = S(t)^2 \int_0^t \frac{\lambda(u)}{P(T \geq u, C \geq u)} du\)。

\textbf{完整证明}：定义计数过程 \(N_i(t) = \mathbf{1}_{\{T_i \leq t, \delta_i = 1\}}\) 和风险过程 \(Y_i(t) = \mathbf{1}_{\{T_i \geq t, C_i \geq t\}}\)。令 \(N(t) = \sum_i N_i(t)\)，\(Y(t) = \sum_i Y_i(t)\)。Nelson-Aalen估计可写为 \(\widehat{\Lambda}(t) = \int_0^t \frac{dN(u)}{Y(u)}\)。由计数过程鞅理论，

\[M(t) = N(t) - \int_0^t Y(u) \lambda(u) du\]

是关于自然筛的鞅。因此

\[\sqrt{n}(\widehat{\Lambda}(t) - \Lambda(t)) = \sqrt{n} \int_0^t \frac{dM(u)}{Y(u)}.\]

由Martingale中心极限定理（Rebolledo定理），此过程弱收敛于Gauss鞅，其方差函数为 \(\int_0^t \frac{\lambda(u)}{P(T \geq u, C \geq u)} du\)。再应用泛函delta方法（乘积积分映射的可微性）即得Kaplan-Meier估计的渐近正态性。\(\blacksquare\)

\textbf{定理 310.2.4（Greenwood公式）} Kaplan-Meier估计的方差估计由Greenwood公式给出：

\[\widehat{\operatorname{var}}(\widehat{S}(t)) = \widehat{S}(t)^2 \sum_{j: \tilde{t}_j \leq t} \frac{d_j}{n_j (n_j - d_j)}.\]

\textbf{证明}：由Delta方法应用于 \(\widehat{S}(t) = \exp(-\widehat{\Lambda}(t))\)，\(\operatorname{var}(\widehat{S}(t)) \approx S(t)^2 \operatorname{var}(\widehat{\Lambda}(t))\)。Nelson-Aalen估计的方差：\(\operatorname{var}(\widehat{\Lambda}(t)) = \sum_{j: \tilde{t}_j \leq t} \operatorname{var}(d_j/n_j)\)。在给定风险集大小 \(n_j\) 的条件下，\(d_j\) 服从二项分布 \(\text{Bin}(n_j, \lambda(\tilde{t}_j)/n_j)\) 的近似。由二项方差公式 \(\operatorname{var}(d_j/n_j \mid n_j) \approx d_j(n_j - d_j)/n_j^3\)。因此 \(\widehat{\operatorname{var}}(\widehat{\Lambda}(t)) = \sum_{j: \tilde{t}_j \leq t} d_j / (n_j(n_j - d_j))\) 是相合估计。代入即得Greenwood公式。\(\blacksquare\)

\textbf{定理 310.2.5}（对数秩检验的渐近分布）：设两组样本的生存函数分别为 \(S_1(t)\) 和 \(S_2(t)\)。对数秩检验（log-rank test）检验零假设 \(H_0: S_1(t) = S_2(t)\)（对所有 \(t\)）等价于 \(H_0: \lambda_1(t) = \lambda_2(t)\)。其检验统计量为：

\[Z = \frac{\sum_{j=1}^k (d_{1j} - e_{1j})}{\sqrt{\sum_{j=1}^k v_{1j}}}\]

其中 \(d_{1j}\) 为第 \(1\) 组在第 \(j\) 个事件时间的观测事件数，\(e_{1j} = n_{1j} d_j / n_j\) 为期望事件数，\(v_{1j} = n_{1j} n_{2j} d_j (n_j - d_j) / (n_j^2 (n_j - 1))\) 为超几何方差。在 \(H_0\) 下，\(Z \xrightarrow{d} \mathcal{N}(0, 1)\)。

\textbf{证明}：对数秩检验可由计数过程鞅理论严格推导。设 \(N_{g}(t) = \sum_{i \in \text{组 } g} N_i(t)\) 为组 \(g\)（\(g = 1, 2\)）的累积事件数，\(Y_{g}(t)\) 为风险集大小。在 \(H_0: \lambda_1(t) = \lambda_2(t) = \lambda(t)\) 下，定义组1的鞅：

\[M_1(t) = N_1(t) - \int_0^t Y_1(u) \lambda(u) du\]

对数秩统计量 \(U = \sum_{j=1}^k (d_{1j} - e_{1j})\) 正是鞅变换 \(\int_0^\tau (Y_2(u)/Y(u)) dN_1(u) - (Y_1(u)/Y(u)) dN_2(u)\) 的离散化。更确切地，定义局部鞅

\[U(t) = \int_0^t \frac{Y_1(u)Y_2(u)}{Y(u)} (d\widehat{\Lambda}_1(u) - d\widehat{\Lambda}_2(u))\]

其中 \(\widehat{\Lambda}_g(t) = \int_0^t dN_g(u) / Y_g(u)\)。在 \(H_0\) 下，\(U(t)\) 是均值为零的鞅。其可选变差过程为 \(\langle U \rangle(t) = \int_0^t (Y_1 Y_2 / Y^2) (1 - \Delta \Lambda) dN\)。用 \(n_j\) 和 \(d_j\) 的离散形式替换得 \(v_{1j} = n_{1j} n_{2j} d_j (n_j - d_j) / (n_j^2 (n_j - 1))\)。由鞅中心极限定理（Rebolledo定理），\(Z = U(\tau) / \sqrt{\langle U \rangle(\tau)} \xrightarrow{d} \mathcal{N}(0, 1)\)。此证明同时适用于删失数据------关键要求是删失机制在两组间无信息差异。\(\blacksquare\)

\subsection{310.3 Cox比例风险模型}\label{coxux6bd4ux4f8bux98ceux9669ux6a21ux578b}

\textbf{定义 310.3.1（Cox比例风险模型）} Cox于1972年提出的半参数模型将危险率函数分解为：

\[\lambda(t \mid \mathbf{X}) = \lambda_0(t) \exp(\boldsymbol{\beta}^T \mathbf{X}),\]

其中 \(\lambda_0(t)\) 为未指定的基线危险率（非参数部分），\(\boldsymbol{\beta}\) 为回归系数（参数部分），\(\mathbf{X}\) 为协变量向量。比例风险假设：不同个体的危险率函数成比例，比例因子不随时间变化。

\textbf{定理 310.3.2（部分似然）} 对Cox模型，\(\boldsymbol{\beta}\) 的部分似然函数（partial likelihood）为：

\[L(\boldsymbol{\beta}) = \prod_{j=1}^k \frac{\exp(\boldsymbol{\beta}^T \mathbf{X}_{(j)})}{\sum_{\ell \in \mathcal{R}_j} \exp(\boldsymbol{\beta}^T \mathbf{X}_\ell)},\]

其中 \(\mathbf{X}_{(j)}\) 为在 \(\tilde{t}_j\) 处发生事件的个体的协变量，\(\mathcal{R}_j\) 为 \(\tilde{t}_j\) 时的风险集。

\textbf{证明框架}：部分似然的推导基于计数过程模型和乘法强度模型。令 \(N_i(t)\) 为个体 \(i\) 的计数过程。Cox模型下的强度过程为 \(\lambda_i(t) = Y_i(t) \lambda_0(t) \exp(\boldsymbol{\beta}^T \mathbf{X}_i)\)。给定在 \([0, \tau]\) 内的所有事件时间，第 \(j\) 个事件发生在个体 \((j)\) 上的条件概率（给定 \(\mathcal{R}_j\) 和事件发生这一事实）为

\[P(\text{事件发生在个体 } (j) \mid \mathcal{R}_j, \text{在 } \tilde{t}_j \text{ 有一个事件}) = \frac{\exp(\boldsymbol{\beta}^T \mathbf{X}_{(j)})}{\sum_{\ell \in \mathcal{R}_j} \exp(\boldsymbol{\beta}^T \mathbf{X}_\ell)}.\]

将所有事件时间的条件概率相乘即得部分似然。\(\boldsymbol{\beta}\) 的极大似然估计 \(\widehat{\boldsymbol{\beta}}\) 是部分似然得分的零点，由Newton-Raphson算法求解。\(\widehat{\boldsymbol{\beta}}\) 的渐近分布为 \(\widehat{\boldsymbol{\beta}} \sim \mathcal{N}(\boldsymbol{\beta}, \mathcal{I}^{-1}(\boldsymbol{\beta}))\)，其中 \(\mathcal{I}(\boldsymbol{\beta})\) 为部分信息矩阵。\(\blacksquare\)

\textbf{定理 310.3.3（Cox模型的基线生存估计------Breslow估计）} 在获得 \(\widehat{\boldsymbol{\beta}}\) 后，累积基线危险率可由Breslow估计给出：

\[\widehat{\Lambda}_0(t) = \sum_{j: \tilde{t}_j \leq t} \frac{d_j}{\sum_{\ell \in \mathcal{R}_j} \exp(\widehat{\boldsymbol{\beta}}^T \mathbf{X}_\ell)}.\]

个体的生存函数预测为 \(\widehat{S}(t \mid \mathbf{X}) = \exp(-\widehat{\Lambda}_0(t) \exp(\widehat{\boldsymbol{\beta}}^T \mathbf{X}))\)。

\textbf{证明}（Breslow估计的相合性）：在Cox模型下，Nelson-Aalen类型的估计量通过profile likelihood方法导出。将基线累积危险率视为非参数的无穷维nuisance参数，固定 \(\boldsymbol{\beta}\)，对 \(\Lambda_0\) 取profile似然最大化。在计数过程框架下，强度模型为 \(\lambda_i(t) = Y_i(t) \lambda_0(t) \exp(\boldsymbol{\beta}^T \mathbf{X}_i)\)。对于固定的 \(\boldsymbol{\beta}\)，\(\Lambda_0\) 的非参数极大似然估计由Breslow公式给出。相合性：令 \(\widehat{\boldsymbol{\beta}} \xrightarrow{P} \boldsymbol{\beta}\)（由部分似然的相合性），且已证明 \(\sup_{t \in [0, \tau]} |\widehat{\Lambda}_0(t; \boldsymbol{\beta}) - \Lambda_0(t)| \xrightarrow{P} 0\) 在模型正确指定且 \(\boldsymbol{\beta}\) 已知的条件下。由连续映射定理和 \(\widehat{\boldsymbol{\beta}}\) 的 \(\sqrt{n}\)-相合性，\(\sup_t |\widehat{\Lambda}_0(t; \widehat{\boldsymbol{\beta}}) - \Lambda_0(t)| \xrightarrow{P} 0\)。进一步，\(\sqrt{n}(\widehat{\Lambda}_0(t; \widehat{\boldsymbol{\beta}}) - \Lambda_0(t))\) 收敛于Gauss过程（由函数delta方法结合经验过程理论）。\(\blacksquare\)

\subsection{310.4 累积发病率与竞争风险模型}\label{ux7d2fux79efux53d1ux75c5ux7387ux4e0eux7adeux4e89ux98ceux9669ux6a21ux578b}

\textbf{定义 310.4.1（竞争风险）} 在竞争风险模型中，个体可能经历 \(K\) 种不同类型的事件之一。观测到的数据是 \((T, D)\)，其中 \(T\) 为事件发生时间，\(D \in \{1, \ldots, K\}\) 为事件类型。原因别危险率（cause-specific hazard）为：

\[\lambda_k(t) = \lim_{\Delta t \to 0} \frac{P(t \leq T < t + \Delta t, D = k \mid T \geq t)}{\Delta t}.\]

\textbf{定义 310.4.2（累积发病率函数）} \(k\) 类事件的累积发病率函数（cumulative incidence function, CIF）为：

\[F_k(t) = P(T \leq t, D = k) = \int_0^t S(u-) \lambda_k(u) du,\]

其中 \(S(t) = \exp(-\sum_{k=1}^K \int_0^t \lambda_k(u) du)\) 为总体生存函数。CIF是次分布函数，极限 \(\lim_{t \to \infty} F_k(t) = P(D = k)\) 为第 \(k\) 类事件的最终概率。

\textbf{定理 310.4.3（Fine-Gray模型）} Fine和Gray提出的子分布风险模型直接对CIF建模：

\[\lambda_k^{\text{sub}}(t \mid \mathbf{X}) = -\frac{d}{dt} \log(1 - F_k(t \mid \mathbf{X})) = \lambda_{k,0}^{\text{sub}}(t) \exp(\boldsymbol{\beta}_k^T \mathbf{X}),\]

其中 \(\lambda_k^{\text{sub}}\) 为子分布危险率。此模型可用逆概率删失加权（IPCW）技术给出部分似然，从而获得 \(\boldsymbol{\beta}_k\) 的相合估计。

\textbf{证明}（IPCW加权部分似然的构造）：令 \(T_i, C_i, \delta_i = \mathbf{1}_{\{T_i \leq C_i\}}, \varepsilon_i \in \{1, \ldots, K\}\) 为个体 \(i\) 的观测数据。定义逆概率删失权重 \(w_i(t) = \mathbf{1}_{\{C_i \geq \min(T_i, t)\}} / \widehat{G}(\min(T_i, t)-)\)，其中 \(\widehat{G}\) 为删失时间的Kaplan-Meier生存估计。核心观察：若对删失机制正确加权，则加权后的风险集近似于无删失情形。Fine-Gray部分似然的得分方程为：

\[U(\boldsymbol{\beta}_k) = \sum_{i=1}^n \int_0^\tau \left( \mathbf{X}_i - \frac{\sum_{j} w_j(u) Y_j(u) \mathbf{X}_j \exp(\boldsymbol{\beta}_k^T \mathbf{X}_j)}{\sum_{j} w_j(u) Y_j(u) \exp(\boldsymbol{\beta}_k^T \mathbf{X}_j)} \right) w_i(u) dN_{ik}^{\text{sub}}(u) = 0\]

其中 \(dN_{ik}^{\text{sub}}(u)\) 为 \(k\) 类事件的子分布计数过程。IPCW的一致性依赖于：\(\mathbb{E}[w_i(u) dN_{ik}^{\text{sub}}(u) \mid \mathcal{F}_{u-}]\) 正比于真子分布风险，这由删失独立性假设和 \(G\) 的相合性保证。在适当的正则条件下，\(\widehat{\boldsymbol{\beta}}_k\) 是相合且渐近正态的，其渐近方差可由sandwich估计量给出。\(\blacksquare\)

\hrulefill

\section{第375章：高维统计与因果推断}\label{ux7b2c375ux7ae0ux9ad8ux7ef4ux7edfux8ba1ux4e0eux56e0ux679cux63a8ux65ad}

当变量个数 \(p\) 与样本量 \(n\) 可比拟或远大于 \(n\) 时，经典统计方法失效。高维统计学通过正则化、稀疏性假设和多重检验校正等思想，为现代大数据分析提供了数学框架。

\subsection{311.1 高维回归的正则化方法}\label{ux9ad8ux7ef4ux56deux5f52ux7684ux6b63ux5219ux5316ux65b9ux6cd5}

\textbf{定义 311.1.1（Lasso）} Lasso（Least Absolute Shrinkage and Selection Operator）是Tibshirani于1996年提出的高维线性回归方法：

\[\widehat{\boldsymbol{\beta}} = \arg\min_{\boldsymbol{\beta}} \left\{ \frac{1}{2n} \|\mathbf{y} - \mathbf{X}\boldsymbol{\beta}\|_2^2 + \lambda \|\boldsymbol{\beta}\|_1 \right\},\]

其中 \(\|\boldsymbol{\beta}\|_1 = \sum_{j=1}^p |\beta_j|\) 为 \(\ell_1\) 范数惩罚。\(\ell_1\) 惩罚的几何特性使得Lasso自动进行变量选择（某些估计系数精确为零）。

\textbf{定理 311.1.2（Lasso的Oracle不等式）} 设真实模型 \(\mathbf{y} = \mathbf{X}\boldsymbol{\beta}^* + \boldsymbol{\varepsilon}\)，\(\boldsymbol{\varepsilon} \sim \mathcal{N}(0, \sigma^2 \mathbf{I})\)。在受限特征值条件或兼容性条件下（设计矩阵 \(\mathbf{X}\) 满足某些稀疏特征值界），取 \(\lambda \asymp \sigma \sqrt{\log p / n}\)，则

\[\|\widehat{\boldsymbol{\beta}} - \boldsymbol{\beta}^*\|_2^2 \leq C \cdot \frac{s_0 \log p}{n},\]

其中 \(s_0 = \|\boldsymbol{\beta}^*\|_0\) 为非零系数的个数，\(C\) 为仅依赖于相容性常数的绝对常数。

\textbf{证明框架}：由Lasso的KKT最优性条件，\(\widehat{\boldsymbol{\beta}}\) 满足 \(\frac{1}{n} \mathbf{X}^T (\mathbf{y} - \mathbf{X}\widehat{\boldsymbol{\beta}}) = \lambda \widehat{\mathbf{z}}\)，其中 \(\widehat{z}_j \in \partial |\widehat{\beta}_j|\)（次梯度）。展开得基本不等式：

\[\frac{1}{2n} \|\mathbf{X}(\widehat{\boldsymbol{\beta}} - \boldsymbol{\beta}^*)\|_2^2 \leq \frac{1}{n}\langle \boldsymbol{\varepsilon}, \mathbf{X}(\widehat{\boldsymbol{\beta}} - \boldsymbol{\beta}^*)\rangle + \lambda(\|\boldsymbol{\beta}^*\|_1 - \|\widehat{\boldsymbol{\beta}}\|_1).\]

取 \(\lambda \geq \frac{2}{n} \|\mathbf{X}^T \boldsymbol{\varepsilon}\|_\infty\) 以高概率控制噪声项（Gauss尾概率和联合界给出 \(\|\mathbf{X}^T \boldsymbol{\varepsilon}\|_\infty / n \lesssim \sigma \sqrt{\log p / n}\)）。利用 \(\ell_1\) 几何不等式：\(\|\widehat{\boldsymbol{\beta}} - \boldsymbol{\beta}^*\|_1\) 可分解为支撑集上的分量和补集上的分量，后者可由前者和KKT条件控制。最终由受限特征值条件将预测误差 \(\frac{1}{n}\|\mathbf{X}(\widehat{\boldsymbol{\beta}} - \boldsymbol{\beta}^*)\|_2^2\) 转化为 \(\ell_2\) 估计误差。\(\blacksquare\)

该界表明，即使 \(p \gg n\)，只要真实模型稀疏（\(s_0\) 小），Lasso的收敛率为 \(\sqrt{s_0 \log p / n}\)，仅比可达到的 \(n = p\) 理想情形多出 \(\log p\) 因子------这正是''Oracle不等式''的命名来源。

\textbf{定理 311.1.3（Ridge回归与Elastic Net）} Ridge回归使用 \(\ell_2\) 惩罚 \(\lambda \|\boldsymbol{\beta}\|_2^2\)，不产生稀疏性但能稳定高度相关的预测变量。Elastic Net结合 \(\ell_1\) 和 \(\ell_2\) 惩罚：

\[\widehat{\boldsymbol{\beta}} = \arg\min_{\boldsymbol{\beta}} \left\{ \frac{1}{2n} \|\mathbf{y} - \mathbf{X}\boldsymbol{\beta}\|_2^2 + \lambda_1 \|\boldsymbol{\beta}\|_1 + \lambda_2 \|\boldsymbol{\beta}\|_2^2 \right\}.\]

Elastic Net兼具Lasso的稀疏性和Ridge的组效应，当预测变量高度相关时优于纯Lasso。

\subsection{311.2 虚假发现率与多重检验}\label{ux865aux5047ux53d1ux73b0ux7387ux4e0eux591aux91cdux68c0ux9a8c}

\textbf{定义 311.2.1（虚假发现率------FDR）} 在 \(m\) 个同时进行的假设检验中，令 \(V\) 为错误拒绝（假阳性）的数量，\(R\) 为总拒绝数量。虚假发现率（False Discovery Rate）定义为：

\[\operatorname{FDR} = \mathbb{E}\left[\frac{V}{\max(R, 1)}\right].\]

FDR是Benjamini和Hochberg于1995年提出的多重检验误差度量，在 \(m\) 很大时比族系错误率（FWER）更为实用。

\textbf{定理 311.2.2（Benjamini-Hochberg程序）} 设 \(p_{(1)} \leq p_{(2)} \leq \cdots \leq p_{(m)}\) 为 \(m\) 个假设的排序 \(p\) 值。令

\[k = \max \left\{ i : p_{(i)} \leq \frac{i}{m} \alpha \right\},\]

拒绝对应于 \(p_{(1)}, \ldots, p_{(k)}\) 的假设。若各检验的 \(p\) 值是独立的（或满足某种正相关条件），则此程序控制 \(\operatorname{FDR} \leq \alpha \cdot \frac{m_0}{m} \leq \alpha\)，其中 \(m_0\) 为真实原假设的个数。

\textbf{证明}：设 \(\mathcal{H}_0\) 为真实原假设的指标集（\(|\mathcal{H}_0| = m_0\)）。对任意 \(\varepsilon > 0\)，定义停止时间 \(k\)。在独立 \(p\) 值假设下，对 \(i \in \mathcal{H}_0\)，\(p_i \sim U(0, 1)\)。关键观察是：BH程序的拒绝阈值是随机的，但条件于 \(\{R = r\}\) 和所有备择假设的 \(p\) 值，每个真实原假设的 \(p\) 值在区间 \([0, r\alpha/m]\) 上的条件概率不超过 \(r\alpha/m\)。由此，\(\mathbb{E}[V / \max(R, 1)] = \sum_{r=1}^m \frac{1}{r} \mathbb{E}[V \mathbf{1}_{\{R = r\}}]\)。通过条件期望和鞅论证可得 \(\mathbb{E}[V / \max(R, 1)] = \frac{m_0}{m} \alpha\)。\(\blacksquare\)

\textbf{定理 311.2.3（Storey的 \(q\) 值方法）} \(q\) 值定义为产生某假阳性率时可接受的最小FDR：\(q(p_{(i)}) = \min_{t \geq p_{(i)}} \widehat{\operatorname{FDR}}(t)\)，其中 \(\widehat{\operatorname{FDR}}(t) = \frac{m \cdot \widehat{\pi}_0 \cdot t}{\#\{p_j \leq t\}}\)。\(\widehat{\pi}_0 = m_0 / m\) 的估计可通过Storey的bootstrap方法获得。\(q\) 值为每个假设提供了一个FDR尺度的显著性度量。

\subsection{311.3 因果推断的数学框架}\label{ux56e0ux679cux63a8ux65adux7684ux6570ux5b66ux6846ux67b6}

\textbf{定义 311.3.1（反事实与潜在结果）} Rubin因果模型（也称为Neyman-Rubin潜在结果框架）将每个个体 \(i\) 在干预 \(a \in \{0, 1\}\)（控制/处理）下的潜在结果记为 \(Y_i(0)\) 和 \(Y_i(1)\)。个体的因果效应为 \(\tau_i = Y_i(1) - Y_i(0)\)。由于仅能观测到一个潜在结果，因果推断本质上是缺失数据问题。

\textbf{定理 311.3.2（无混杂假设下的识别）} 在条件无混杂假设（conditional ignorability）下：

\[\{Y(0), Y(1)\} \perp A \mid \mathbf{X},\]

其中 \(A\) 为实际接受的干预，\(\mathbf{X}\) 为可观测混杂变量。则平均处理效应（ATE）可由下式识别：

\[\operatorname{ATE} = \mathbb{E}[Y(1) - Y(0)] = \mathbb{E}_{\mathbf{X}} [\mathbb{E}[Y \mid A = 1, \mathbf{X}] - \mathbb{E}[Y \mid A = 0, \mathbf{X}]].\]

\textbf{证明}：由条件无混杂和一致性假设（\(Y = A Y(1) + (1 - A) Y(0)\)），

\[\mathbb{E}[Y(a) \mid \mathbf{X}] = \mathbb{E}[Y(a) \mid A = a, \mathbf{X}] = \mathbb{E}[Y \mid A = a, \mathbf{X}].\]

取 \(\mathbf{X}\) 上的期望即得ATE的识别公式。\(\blacksquare\)

\textbf{定义 311.3.3（do-演算）} Pearl的结构因果模型将因果效应形式化为有向无环图（DAG）上的干预操作 \(\operatorname{do}(X = x)\)。do-演算提供了三条规则，用于从观测分布表达干预分布：

\begin{enumerate}
\def\labelenumi{\arabic{enumi}.}
\tightlist
\item
  （观测删除）若 \(Y\) 与 \(Z\) 在删除进入 \(X\) 的箭头后的图中 \(d\)-分离，则 \(P(y \mid \operatorname{do}(x), z) = P(y \mid \operatorname{do}(x))\)；
\item
  （干预/观测交换）若 \(Y\) 与 \(Z\) 在删除进入 \(X\) 且删除从 \(Z\) 出发的箭头后的图中 \(d\)-分离，则 \(P(y \mid \operatorname{do}(x), \operatorname{do}(z)) = P(y \mid \operatorname{do}(x), z)\)；
\item
  （干预删除）若 \(Y\) 与 \(Z\) 在删除进入 \(X\) 且删除进入 \(Z(W)\) 的箭头后的图中 \(d\)-分离，且 \(Z(W)\) 为不影响 \(X\) 的 \(Z\) 中的节点集合，则可删除 \(\operatorname{do}(z)\)。
\end{enumerate}

do-演算的完备性定理（Shpitser和Pearl, 2006）保证了若因果效应可从观测数据识别，则必可通过有限次do-演算推导出来。

\textbf{定理 311.3.4（后门准则与前门准则）} 在DAG \(G\) 中，

\begin{itemize}
\tightlist
\item
  后门准则：若 \(\mathbf{Z}\) 阻塞了 \(X\) 到 \(Y\) 的所有后门路径（以指向 \(X\) 的箭头开始的路径），则 \(P(y \mid \operatorname{do}(x)) = \sum_{\mathbf{z}} P(y \mid x, \mathbf{z}) P(\mathbf{z})\)。
\item
  前门准则：若变量 \(M\) 中介了所有 \(X\) 对 \(Y\) 的因果路径，且从 \(X\) 到 \(M\) 不存在后门路径，从 \(M\) 到 \(Y\) 的后门路径被 \(X\) 阻塞，则 \(P(y \mid \operatorname{do}(x)) = \sum_m P(m \mid x) \sum_{x'} P(y \mid m, x') P(x')\)。
\end{itemize}

\subsection{311.4 工具变量与结构方程模型}\label{ux5de5ux5177ux53d8ux91cfux4e0eux7ed3ux6784ux65b9ux7a0bux6a21ux578b}

\textbf{定义 311.4.1（工具变量）} 变量 \(Z\) 称为 \(X\)（处理）对 \(Y\)（结果）的因果效应的工具变量（instrumental variable, IV），若满足：（1）\(Z\) 与 \(X\) 相关（相关性）；（2）\(Z\) 仅通过 \(X\) 影响 \(Y\)（排斥限制）；（3）\(Z\) 与未观测混杂因素独立（外生性）。

\textbf{定理 311.4.2（Wald估计量）} 若 \(Z\) 为二值的有效工具变量，则因果效应可由Wald估计量识别：

\[\tau = \frac{\mathbb{E}[Y \mid Z = 1] - \mathbb{E}[Y \mid Z = 0]}{\mathbb{E}[X \mid Z = 1] - \mathbb{E}[X \mid Z = 0]}.\]

Wald估计量可视为依从者平均因果效应（LATE，Local Average Treatment Effect），它识别了那些其处理状态被工具变量改变的个体（依从者）的平均处理效应。

\textbf{证明}：由IV的有效性，\(Z\) 与潜在结果 \((Y(0), Y(1))\) 独立且仅通过 \(X\) 影响 \(Y\)。因此 \(\mathbb{E}[Y \mid Z=1] - \mathbb{E}[Y \mid Z=0] = \mathbb{E}[Y(1)X(1) + Y(0)(1-X(1))] - \mathbb{E}[Y(1)X(0) + Y(0)(1-X(0))]\)。对于依从者（\(X(1)=1, X(0)=0\)），此差异为 \(\mathbb{E}[Y(1)-Y(0) \mid X(1)=1, X(0)=0] \cdot P(X(1)=1, X(0)=0)\)。分母 \(\mathbb{E}[X \mid Z=1] - \mathbb{E}[X \mid Z=0] = P(X(1)=1, X(0)=0)\) 恰为依从者比例。因此Wald估计量识别了依从者平均因果效应（LATE）。\(\blacksquare\)

\textbf{定理 311.4.3（两阶段最小二乘法------2SLS）} 在线性结构方程模型 \(Y = \tau X + \mathbf{W}^T \boldsymbol{\gamma} + \varepsilon\)（\(X\) 内生）和第一阶段 \(X = \mathbf{Z}^T \boldsymbol{\pi} + \mathbf{W}^T \boldsymbol{\delta} + \nu\) 的框架下，2SLS估计量为 \(\widehat{\tau}_{\text{2SLS}} = (\widehat{\mathbf{X}}^T \mathbf{M}_{\mathbf{W}} \widehat{\mathbf{X}})^{-1} \widehat{\mathbf{X}}^T \mathbf{M}_{\mathbf{W}} \mathbf{Y}\)，其中 \(\widehat{\mathbf{X}}\) 为第一阶段的预测值，\(\mathbf{M}_{\mathbf{W}}\) 为关于 \(\mathbf{W}\) 的投影矩阵。2SLS在上述有效性条件下是相合且渐近正态的。

\textbf{证明}：结构方程为 \(Y = X\tau + W\gamma + \varepsilon\)，\(X = Z\pi + W\delta + \nu\)。由于 \(X\) 与 \(\varepsilon\) 相关（内生性），OLS不一致。2SLS方法：第一阶段用OLS估计 \(\widehat{X} = Z\widehat{\pi} + W\widehat{\delta}\)。第二阶段用 \(\widehat{X}\) 替代 \(X\) 进行OLS回归。由排除限制（\(Z\) 不直接影响 \(Y\)）和外生性（\(Z \perp \varepsilon\)），\(\widehat{\tau}_{\text{2SLS}}\) 是 \(\tau\) 的相合估计：\(\widehat{\tau}_{\text{2SLS}} \xrightarrow{P} \tau\)。渐近正态性 \(\sqrt{n}(\widehat{\tau}_{\text{2SLS}} - \tau) \xrightarrow{d} \mathcal{N}(0, \sigma^2/\mathbb{E}[(Z\pi)^2])\) 由标准M估计量理论给出。\(\blacksquare\)

\textbf{定理 311.4.4（弱工具变量的诊断）} 若第一阶段 \(F\) 统计量小于 \(10\)（Stock-Yogo准则），则工具变量被视为弱工具，2SLS估计在有穷样本下严重有偏。在此情形下，应使用有限信息极大似然（LIML）或连续更新估计（CUE）等对弱工具变量更为稳健的方法。

\textbf{定理 311.4.5（结构方程模型------SEM）} 结构方程模型将因子分析和路径分析统一在一个框架内。测量模型 \(\mathbf{X} = \boldsymbol{\Lambda}_X \boldsymbol{\xi} + \boldsymbol{\delta}\) 将潜变量 \(\boldsymbol{\xi}\) 与观测变量 \(\mathbf{X}\) 联系起来，其中 \(\boldsymbol{\Lambda}_X\) 为因子载荷矩阵。结构模型 \(\boldsymbol{\eta} = \mathbf{B} \boldsymbol{\eta} + \boldsymbol{\Gamma} \boldsymbol{\xi} + \boldsymbol{\zeta}\) 描述潜变量之间的关系。模型参数的估计通常通过极大似然估计（在多元正态假设下）或渐近分布自由方法（ADF）得到。模型的全局拟合由 \(\chi^2\) 检验、RMSEA、CFI、TLI等指标评估。

\textbf{定理 311.4.6（倾向评分与逆概率加权）} Rosenbaum和Rubin于1983年证明：在无混杂假设下，处理分配 \(A\) 条件于标量倾向评分 \(e(\mathbf{X}) = P(A = 1 \mid \mathbf{X})\) 等同于条件于整个高维协变量向量 \(\mathbf{X}\)。由此，ATE可由逆概率加权（IPW）估计量识别：

\[\widehat{\tau}_{\text{IPW}} = \frac{1}{n} \sum_{i=1}^n \left( \frac{A_i Y_i}{\widehat{e}(\mathbf{X}_i)} - \frac{(1 - A_i) Y_i}{1 - \widehat{e}(\mathbf{X}_i)} \right).\]

Doubly robust估计量结合了结果回归模型和倾向评分模型，只要两者之一正确指定即能给出相合估计------这是现代因果推断中最强大且实用的方法之一。

\textbf{证明}：由无混杂假设 \(Y(a) \perp A \mid \mathbf{X}\)，Rosenbaum和Rubin证明 \(Y(a) \perp A \mid e(\mathbf{X})\)。因此 \(\mathbb{E}[Y(1)] = \mathbb{E}[\mathbb{E}[Y \mid A=1, \mathbf{X}]] = \mathbb{E}[\frac{A Y}{e(\mathbf{X})}]\)（因 \(\mathbb{E}[A/ e(\mathbf{X}) \mid \mathbf{X}] = 1\)）。类似地 \(\mathbb{E}[Y(0)] = \mathbb{E}[\frac{(1-A)Y}{1-e(\mathbf{X})}]\)。IPW估计量是这些期望的样本类比。Doubly robust性质：若结果模型 \(\mathbb{E}[Y \mid A, \mathbf{X}] = m(A, \mathbf{X}; \alpha)\) 或倾向模型 \(e(\mathbf{X}; \gamma)\) 之一正确指定，则DR估计量 \(\widehat{\tau}_{\text{DR}} = n^{-1}\sum [\frac{A_i Y_i}{\widehat{e}_i} - \frac{A_i - \widehat{e}_i}{\widehat{e}_i} m(1,\mathbf{X}_i)] - [\cdots]\) 相合估计ATE。\(\blacksquare\)

\textbf{定理 311.4.7（差分中之差分法------DiD）} 差分中之差分法（Differences-in-Differences）是面板数据因果推断的经典方法。在两期两组的面板结构下，处理效应由双重差分识别：

\[\tau_{\text{DiD}} = (\mathbb{E}[Y_{t=1} \mid D = 1] - \mathbb{E}[Y_{t=0} \mid D = 1]) - (\mathbb{E}[Y_{t=1} \mid D = 0] - \mathbb{E}[Y_{t=0} \mid D = 0]),\]

其中 \(D\) 为处理组指示变量。DiD的核心假设是''平行趋势''：在没有处理的情况下，处理组和对照组的平均结果演变趋势相同。对于多期多组的广义DiD模型，可通过双向固定效应回归实现估计。

\textbf{证明}：设 \(Y_{it}(0)\) 为个体 \(i\) 在时期 \(t\) 未受处理时的潜在结果，\(Y_{it}(1)\) 为受处理时的结果。处理组 \(D=1\) 在 \(t=1\) 时接受处理。平行趋势假设：\(\mathbb{E}[Y_{i1}(0) - Y_{i0}(0) \mid D=1] = \mathbb{E}[Y_{i1}(0) - Y_{i0}(0) \mid D=0]\)。因此DiD估计量为： \[\tau_{\text{DiD}} = \mathbb{E}[Y_{i1} - Y_{i0} \mid D=1] - \mathbb{E}[Y_{i1} - Y_{i0} \mid D=0]\] \[= \mathbb{E}[Y_{i1}(1) - Y_{i0}(0) \mid D=1] - \mathbb{E}[Y_{i1}(0) - Y_{i0}(0) \mid D=0]\] \[= \mathbb{E}[Y_{i1}(1) - Y_{i1}(0) \mid D=1] + (\mathbb{E}[Y_{i1}(0) - Y_{i0}(0) \mid D=1] - \mathbb{E}[Y_{i1}(0) - Y_{i0}(0) \mid D=0]).\] 第二项由平行趋势假设为零，故 \(\tau_{\text{DiD}} = \mathbb{E}[Y(1) - Y(0) \mid D=1]\) 识别了处理组的ATT。\(\blacksquare\)

\textbf{定理 311.4.8（回归间断设计------RDD）} 回归间断设计利用处理分配在某个阈值处的确定性断点来识别因果效应。在Sharp RDD中，处理指示 \(A = \mathbf{1}_{\{R \geq r_0\}}\)（\(R\) 为运行变量，\(r_0\) 为已知阈值）。因果效应为：

\[\tau_{\text{SRD}} = \lim_{r \downarrow r_0} \mathbb{E}[Y \mid R = r] - \lim_{r \uparrow r_0} \mathbb{E}[Y \mid R = r].\]

RDD的局部随机化特性（在阈值附近的个体近似随机地落入处理或对照组）使其成为最接近随机化实验的观测研究方法。Hahn、Todd和Van der Klaauw（2001）以及Porter（2003）发展了基于局部线性回归的RDD估计理论，包括最优窗宽选择和渐近分布理论。

\textbf{证明}：在Sharp RDD中，\(A = \mathbf{1}_{\{R \geq r_0\}}\)。假设 \(\mathbb{E}[Y(0) \mid R=r]\) 和 \(\mathbb{E}[Y(1) \mid R=r]\) 在 \(r=r_0\) 处连续。则在阈值处，\(\lim_{r \downarrow r_0} \mathbb{E}[Y \mid R=r] = \mathbb{E}[Y(1) \mid R=r_0]\)，\(\lim_{r \uparrow r_0} \mathbb{E}[Y \mid R=r] = \mathbb{E}[Y(0) \mid R=r_0]\)。相减得 \(\tau_{\text{SRD}} = \mathbb{E}[Y(1)-Y(0) \mid R=r_0]\)------阈值处的局部平均处理效应。Hahn-Todd-Van der Klaauw证明最优窗宽为 \(h \asymp n^{-1/5}\)，局部线性估计量在正则条件下达到 \(n^{-2/5}\) 的收敛率。\(\blacksquare\)

\textbf{定理 311.4.9（合成控制法）} Abadie和Gardeazabal（2003）、Abadie、Diamond和Hainmueller（2010）提出的合成控制法通过加权组合未受处理的对照单元来构造处理单元的反事实。合成控制权重由极小化处理前预测误差的约束优化问题确定。该方法在处理单元数少、对照单元多的比较案例研究中尤其有效，已广泛应用于政策评估和比较政治经济学分析。

\newpage

\chapter{10-概率与统计/03-统计学/05-拓扑数据分析.md}\label{ux6982ux7387ux4e0eux7edfux8ba103-ux7edfux8ba1ux5b6605-ux62d3ux6251ux6570ux636eux5206ux6790.md}

\chapter{拓扑数据分析}\label{ux62d3ux6251ux6570ux636eux5206ux6790}

\hrulefill

\section{第376章：持续同调}\label{ux7b2c376ux7ae0ux6301ux7eedux540cux8c03}

拓扑数据分析（Topological Data Analysis, TDA）是代数拓扑与数据科学交叉的前沿领域，利用拓扑不变量------特别是持续同调（persistent homology）------从高维复杂数据中提取全局的结构信息。本章建立持续同调的代数与几何基础。

\subsection{单纯复形与抽象复形}\label{ux5355ux7eafux590dux5f62ux4e0eux62bdux8c61ux590dux5f62}

\textbf{定义 312.1（单纯复形）} 一个\textbf{抽象单纯复形} \(K\) 是有限集 \(V\) 的子集族，满足：（1）对任意 \(v \in V\)，\(\{v\} \in K\)；（2）若 \(\sigma \in K\) 且 \(\tau \subseteq \sigma\)，则 \(\tau \in K\)。\(\sigma\) 称为\textbf{单纯形}（simplex），其维数为 \(|\sigma| - 1\)。\(K\) 的维数是其包含的最大单纯形的维数。

\textbf{定义 312.2（Čech复形）} 设 \(X = \{x_1, \ldots, x_n\} \subset \mathbb{R}^d\)（或任意度量空间中的有限点集）。给定半径 \(\varepsilon > 0\)，\textbf{Čech复形} \(\check{C}_{\varepsilon}(X)\) 定义为：单纯形 \(\{x_{i_0}, \ldots, x_{i_k}\}\) 属于该复形当且仅当 \[\bigcap_{j=0}^{k} \overline{B}_{\varepsilon}(x_{i_j}) \neq \emptyset.\] 由神经定理（Nerve Theorem），\(\check{C}_{\varepsilon}(X)\) 与集合 \(\bigcup_{i=1}^{n} \overline{B}_{\varepsilon}(x_i)\) 同伦等价。

\textbf{定义 312.3（Vietoris-Rips复形）} 由于Čech复形的计算量极大（需检查所有交的非空性），实用中常用\textbf{Vietoris-Rips复形} \(\operatorname{VR}_{\varepsilon}(X)\)：单纯形 \(\{x_{i_0}, \ldots, x_{i_k}\}\) 属于该复形当且仅当所有两两距离不超过 \(\varepsilon\)： \[d(x_{i_j}, x_{i_{\ell}}) \leq \varepsilon, \quad \forall j, \ell.\] VR复形由图的团复形（clique complex）完全确定，易于计算。成立包含关系： \[\check{C}_{\varepsilon}(X) \subseteq \operatorname{VR}_{2\varepsilon}(X) \subseteq \check{C}_{2\varepsilon}(X).\]

\textbf{定义 312.4（Witness复形）} \textbf{Witness复形}由de Silva和Carlsson（2004）引入，用于处理大规模数据集。选取地标点集 \(L \subset X\)（\(|L| \ll |X|\)），地标点 \(\ell_{i_0}, \ldots, \ell_{i_k}\) 构成单纯形当且仅当存在见证点 \(w \in X\) 对所有这些地标点``足够近''（按某种排序标准）。Witness复形大幅降低了计算复杂度，同时保持了数据的拓扑特征。

\subsection{滤过与持续同调的代数基础}\label{ux6ee4ux8fc7ux4e0eux6301ux7eedux540cux8c03ux7684ux4ee3ux6570ux57faux7840}

\textbf{定义 312.5（滤过）} 单纯复形的\textbf{滤过}（filtration）是嵌套族 \(\emptyset = K_0 \subseteq K_1 \subseteq \cdots \subseteq K_m = K\)。对任意 \(i \leq j\)，包含映射 \(K_i \hookrightarrow K_j\) 诱导同调群之间的同态： \[f_p^{i,j}: H_p(K_i) \to H_p(K_j).\]

\textbf{定义 312.6（持久模）} 在给定域 \(\mathbb{F}\)（通常取 \(\mathbb{F} = \mathbb{Z}/2\mathbb{Z}\) 以简化定向问题）上的系统 \(\{H_p(K_i)\}_{i=1}^{m}\) 连同映射 \(\{f_p^{i,j}\}_{i \leq j}\) 构成一个\textbf{持久模}（persistence module），即从偏序集 \((\{1, \ldots, m\}, \leq)\) 到 \(\mathbb{F}\)-向量空间范畴的函子。

\textbf{定义 312.7（条形码与持续图）} 持续同调的第 \(p\) 维\textbf{条形码}（barcode）是多重集 \(\{(b_i, d_i)\}\)，其中每个区间 \([b_i, d_i)\) 表示一个 \(p\) 维同调类``出生''于 \(K_{b_i}\) 且``死亡''于 \(K_{d_i}\)（即在该步被合并入边界或与先前存在的同调类合并）。等价地，\textbf{持续图}（persistence diagram）是 \(\mathbb{R}^2\)（扩展平面）中的点集： \[\operatorname{Dgm}_p = \{(b_i, d_i) \in \mathbb{R} \times (\mathbb{R} \cup \{\infty\}) : b_i < d_i\},\] 每个点对应一个具有出生时间 \(b_i\) 和死亡时间 \(d_i\) 的同调类。

\subsection{持续同调的结构定理}\label{ux6301ux7eedux540cux8c03ux7684ux7ed3ux6784ux5b9aux7406}

\textbf{定理 312.8（持续同调的分类定理 / 结构定理）} 设 \(\mathbb{V} = \{V_i, \varphi_i^{j}\}\) 是有限维持久模（即每个 \(V_i\) 是有限维向量空间且 \(\varphi_i^{j}\) 是线性映射）。则存在唯一的区间分解（interval decomposition）： \[\mathbb{V} \cong \bigoplus_{\ell=1}^{N} \mathbb{I}[b_{\ell}, d_{\ell}],\] 其中 \(\mathbb{I}[b, d]\) 是区间模：当 \(i \in [b, d]\) 时 \((\mathbb{I}[b, d])_i = \mathbb{F}\)，否则为 \(0\)；且所有映射在可能时均为恒等映射。

\textbf{证明} 该定理是Gabriel定理（关于 \(A_n\) 型箭图的表示）的直接推论。将偏序集 \(\{1, \ldots, m\}\) 视为 \(A_m\) 型箭图，则持久模恰为该箭图的表示。\(A_m\) 箭图是有限表示型的，其不可分解模一一对应于区间 \([b, d] \subseteq \{1, \ldots, m\}\)（其中 \(d\) 可为 \(\infty\)）。Krull-Schmidt定理保证分解在重排区间序后是唯一的。更一般地，Crawley-Boevey（2015）将结果推广至任意全序集上的点态有限持久模。\(\blacksquare\)

\textbf{推论 312.9} 持续同调完全由持续图（条形码）刻画。两个滤过具有相同的持续图当且仅当它们的持久模同构。

\textbf{定理 312.10（持续同调函子性质）} 对固定的空间 \(X\) 和函数 \(f: X \to \mathbb{R}\)，子水平集滤过 \(X_a = f^{-1}((-\infty, a])\) 定义了函子 \(a \mapsto H_p(X_a)\)，满足函子性（functoriality）：若还有连续映射 \(g: Y \to X\)，则导出持久模之间的自然变换。这使持续同调成为拓扑数据分析中强大的不变量。

\subsection{持续Landscape与持续图像}\label{ux6301ux7eedlandscapeux4e0eux6301ux7eedux56feux50cf}

\textbf{定义 312.11（持续Landscape）} Bubenik（2015）定义的第 \(k\) 阶\textbf{持续Landscape}是从持续图构造的函数 \(\lambda_k: \mathbb{R} \to \mathbb{R}\)。对持续图中的每个点 \((b, d)\)，定义三角形函数： \[\Lambda_{(b,d)}(t) = \max\{0, \min\{t - b, d - t\}\}.\] 则 \(\lambda_k(t)\) 定义为集合 \(\{\Lambda_{(b,d)}(t) : (b,d) \in \operatorname{Dgm}\}\) 中第 \(k\) 大的值。

持续Landscape将持续图转化为Banach空间 \(L^p(\mathbb{R})\) 中的函数，使得统计方法（均值、方差、置信带）可应用于拓扑特征。Landscape是 \(1\)-Lipschitz的（关于持续图的Bottleneck距离），且构成 \(L^2(\mathbb{R})\) 的稠密子集。

\textbf{定义 312.12（持续图像）} Adams等人（2017）提出的\textbf{持续图像}（persistence images）将持续图转化为固定分辨率的向量。首先对持续图应用线性变换 \(T(b, d) = (b, d - b)\)（转换为出生-持续性坐标），然后用加权高斯核进行核密度估计并在固定网格上采样，得到可用于标准机器学习管线的有限维特征向量。

\subsection{稳定性定理}\label{ux7a33ux5b9aux6027ux5b9aux7406}

\textbf{定义 312.13（Bottleneck距离与Wasserstein距离）} 设 \(\operatorname{Dgm}_1, \operatorname{Dgm}_2\) 是两个持续图。其\textbf{Bottleneck距离}定义为 \[d_B(\operatorname{Dgm}_1, \operatorname{Dgm}_2) = \inf_{\gamma} \sup_{x \in \operatorname{Dgm}_1 \cup \Delta} \|x - \gamma(x)\|_{\infty},\] 其中 \(\gamma\) 是 \(\operatorname{Dgm}_1 \cup \Delta\) 到 \(\operatorname{Dgm}_2 \cup \Delta\) 的双射（\(\Delta = \{(t, t)\}\) 是对角线，表示出生即死亡的``噪声''类）。\(p\)-Wasserstein距离类似但使用 \(L^p\) 范数。

\textbf{定理 312.14（Cohen-Steiner-Edelsbrunner-Harer稳定性定理, 2007）} 设 \(f, g: X \to \mathbb{R}\) 是拓扑空间 \(X\) 上的驯服（tame）函数。则 \[d_B(\operatorname{Dgm}_p(f), \operatorname{Dgm}_p(g)) \leq \|f - g\|_{\infty}.\]

\textbf{证明框架} 将 \(f\) 和 \(g\) 通过同伦（线性插值）连接：\(h_t = (1-t)f + tg\)。考虑 \(h_t\) 的子水平集滤过，利用代数稳定性引理，即函子间自然变换的``交错''（interleaving）距离与Bottleneck距离的等价性。具体地，若 \(\|f - g\|_{\infty} \leq \varepsilon\)，则持久模 \(H_p(f^{-1}(-\infty, \cdot])\) 和 \(H_p(g^{-1}(-\infty, \cdot])\) 是 \(\varepsilon\)-交错的，由代数稳定性定理即得结论。\(\blacksquare\)

该定理是持续同调理论中最核心的结果之一，保证了小的数据扰动只会产生小的持续图变化，从而赋予TDA方法以统计稳健性。

\hrulefill

\section{第377章：TDA的算法与统计基础}\label{ux7b2c377ux7ae0tdaux7684ux7b97ux6cd5ux4e0eux7edfux8ba1ux57faux7840}

\subsection{持续同调的计算}\label{ux6301ux7eedux540cux8c03ux7684ux8ba1ux7b97}

\textbf{算法 313.1（标准约化算法 / 边界矩阵约化）} 持续同调的标准计算基于边界矩阵的列约化（类似Smith正规形）。

设滤过中单纯形按添加顺序排列为 \(\sigma_1, \ldots, \sigma_m\)。边界矩阵 \(D \in \mathbb{F}^{m \times m}\) 定义为 \(D_{ij} = 1\) 当且仅当 \(\sigma_i\) 是 \(\sigma_j\) 的余维数为 \(1\) 的面。算法执行列操作 \(D \leftarrow D \cdot R\)（右乘可逆矩阵 \(R\)），使得 \(D\) 的每列的最低非零元（\(\operatorname{low}(j)\)）各不相同。最终，若列 \(j\) 归零（整列为零），则 \(\sigma_j\) 是正单纯形（创建新同调类）；若 \(\operatorname{low}(j) = i\)，则 \(\sigma_j\) 是负单纯形（消去 \(\sigma_i\) 创建的同调类），对应特征的持续区间为 \([f(\sigma_i), f(\sigma_j))\)。

对该朴素算法的时间复杂度为 \(O(m^3)\)。借助稀疏矩阵技巧（使用链表或压缩行存储）可将实践复杂度大幅降低。

\textbf{算法 313.2（Ripser算法）} Bauer（2017-2021）开发的\textbf{Ripser}是目前最高效的VR复形持续同调计算软件。其核心创新包括：（1）利用Vietoris-Rips复形的``全称系数''特性------只需 \(\mathbb{Z}/2\mathbb{Z}\) 系数；（2）通过列压缩（clearing）大幅减少需处理的列数；（3）使用枚举直径子复形（cohomology）方法------转而计算上同调，其约化远快于下同调；（4）基于枚举共面的高效补面生成。这些技巧使Ripser可在商用硬件上处理数千个数据点的VR持续同调计算。

\subsection{Mapper算法与Reeb图}\label{mapperux7b97ux6cd5ux4e0ereebux56fe}

\textbf{定义 313.3（Reeb图）} 设 \(X\) 是拓扑空间，\(f: X \to \mathbb{R}\) 是连续函数。\(f\) 的\textbf{Reeb图} \(R_f\) 是 \(X\) 关于等价关系 \(x \sim y \iff f(x) = f(y)\) 且 \(x\) 和 \(y\) 属于 \(f^{-1}(t)\) 的同一连通分支的商空间。

\textbf{算法 313.4（Mapper算法 - Singh-Mémoli-Carlsson, 2007）} Mapper是Reeb图的统计近似：

（1）\textbf{过滤函数}：选取 \(f: X \to \mathbb{R}^d\)（如密度估计、距离函数）；

（2）\textbf{覆盖}：将像空间 \(f(X)\) 用重叠的开覆盖 \(\mathcal{U} = \{U_i\}\) 覆盖；

（3）\textbf{聚类}：对每个 \(f^{-1}(U_i)\) 进行聚类（如单链聚类）；

（4）\textbf{构造神经}：以聚类结果为顶点，若两个聚类有非空交则连边，构造单纯复形。

Mapper生成的复形（通常为图或简单复形）概括了数据的拓扑``骨架''，保留了 \(f\) 不同水平集之间的连接关系。其输出是一个可交互可视化的``网络''，广泛用于高维数据探索。

\subsection{持续同调的统计学}\label{ux6301ux7eedux540cux8c03ux7684ux7edfux8ba1ux5b66}

\textbf{定义 313.5（持续图的期望与分布）} 设 \(X_1, \ldots, X_n\) 是从某分布 \(\mathbb{P}\) 中独立同分布抽取的样本。由此构建的持续图 \(\widehat{\operatorname{Dgm}}_n\) 是随机变量（取值于持续图空间）。其极限行为的刻画是TDA统计理论的核心问题。

\textbf{定理 313.6（Fasy-Lecci-Rinaldo-Wasserman等的置信带, 2014）} 对从密度支撑的分布中抽取的i.i.d.样本构造的持续图，可通过Bottleneck Bootstrap方法构建 \((1-\alpha)\) 水平的置信带。令 \(\widehat{\operatorname{Dgm}}_n\) 为样本持续图，\(\mathbb{D}\) 为极限持续图。则存在序列 \(t_n\) 使得 \[\lim_{n \to \infty} \mathbb{P}(d_B(\widehat{\operatorname{Dgm}}_n, \mathbb{D}) \leq t_n) = 1 - \alpha,\] 其中 \(t_n\) 通过bootstrap估计获得。对持续图中的任意显著特征（特征是点的持续性 \(d - b\) 大于某阈值），可以赋予统计显著性水平。

\subsection{核方法在TDA中的应用}\label{ux6838ux65b9ux6cd5ux5728tdaux4e2dux7684ux5e94ux7528}

\textbf{定义 313.7（持续同调核）} 为将TDA整合入核方法（如支持向量机），需定义持续图上的正定核。几种主要构造：

（1）\textbf{持续尺度核}（Reininghaus等人, 2015）：对持续图 \(D_1, D_2\)， \[K_{\sigma}(D_1, D_2) = \frac{1}{8\pi \sigma} \sum_{p \in D_1, q \in D_2} e^{-\frac{\|p-q\|^2}{8\sigma}} - e^{-\frac{\|p-\bar{q}\|^2}{8\sigma}} - e^{-\frac{\|\bar{p}-q\|^2}{8\sigma}} + e^{-\frac{\|\bar{p}-\bar{q}\|^2}{8\sigma}},\] 其中 \(\bar{p}\) 是 \(p\) 关于对角线的镜像。可证明该核在持续图空间上诱导的RKHS是正定的。

（2）\textbf{持续Fisher核}（Le-Yamazaki, 2018）：基于持续Landscape的Fisher信息度量构造，捕捉了持续图在参数族中的几何变化。

\hrulefill

\section{第378章：TDA的应用}\label{ux7b2c378ux7ae0tdaux7684ux5e94ux7528}

\subsection{高维数据分析与降维可视化}\label{ux9ad8ux7ef4ux6570ux636eux5206ux6790ux4e0eux964dux7ef4ux53efux89c6ux5316}

持续同调提供了一种''拓扑PCA''------系统识别数据中的环（loops，\(H_1\)）、空腔（voids，\(H_2\)）等高阶特征，这些特征无法被传统的线性降维（PCA、MDS）或基于距离的方法（t-SNE、UMAP）捕获。

\textbf{定义 314.1}（持续同调的拓扑降维）：设数据点云 \(X = \{x_1, \ldots, x_N\} \subset \mathbb{R}^d\)。持续同调的降维流程包括：(i) 构造滤过 \(\{K_\varepsilon\}_{\varepsilon \geq 0}\)（通常使用 VR 复形）；(ii) 计算各维度的持续图 \(\operatorname{Dgm}_p(X)\)（\(p = 0, 1, 2\)）；(iii) 提取显著特征：对持续性 \(\pi(b, d) = d - b\) 大于阈值 \(\tau\) 的点予以保留；(iv) 逆向追踪代表圈（representative cycles）：对每个显著 \(p\)-维特征 \((b, d)\)，找到其出生单纯形 \(\sigma_b\) 和死亡单纯形 \(\sigma_d\)，通过解边界矩阵的线性方程组构造代表圈 \(\gamma \in Z_p(K_{b+\varepsilon})\)，其几何实现给出数据的全局拓扑结构。

\textbf{定理 314.1}（显著特征的数量控制）：设 \(\operatorname{Dgm}_p(X)\) 为 VR 复形滤过的 \(p\) 维持续图。对任意 \(\varepsilon > 0\)，持续性超过 \(\varepsilon\) 的特征的数量不超过 \(H_p(\operatorname{VR}_{\varepsilon}(X))\) 的 Betti 数。特别地，当 \(\varepsilon\) 足够大时，显著 \(H_1\) 特征的数量以数据点个数 \(N\) 的某个函数为上界，具体依赖于点云的几何维数。

\textbf{典型工作流程}：（1）计算 \(\beta_0, \beta_1, \ldots, \beta_k\) 条码；（2）在持续图中识别远离对角线的显著特征（\(\operatorname{pers}(b,d) = d - b\) 大）；（3）逆向追踪这些特征在原始数据中的代表圈（representative cycles），获得对数据结构的具体几何解释。

\subsection{材料科学中的TDA}\label{ux6750ux6599ux79d1ux5b66ux4e2dux7684tda}

\textbf{非晶材料的结构分析}：Hiraoka等人（2016）用持续同调分析非晶硅和硅玻璃的原子构型。通过将原子位置建模为点云并计算其持续图，他们发现：

\begin{itemize}
\tightlist
\item
  非晶硅的 \(H_2\) 持续图揭示了中程序有序性（\(5-15\)埃尺度的有序结构）；
\item
  持续图的 \(H_1\) 和 \(H_2\) 特征与材料的局部弹性响应之间存在强统计关联；
\item
  不同制备工艺（快冷 vs.~溅射）产生的持续图具有系统差异，表明TDA可用作加工-结构-性能关系的``描述符''。
\end{itemize}

该方法已推广至金属玻璃、胶体系统和多孔介质的表征。

\subsection{神经科学中的TDA}\label{ux795eux7ecfux79d1ux5b66ux4e2dux7684tda}

\textbf{定义 314.3}（神经群体编码的拓扑表征）：设 \(N\) 个神经元在 \(T\) 个时间窗口中的活动矩阵为 \(A \in \mathbb{R}^{N \times T}\)（\(A_{it}\) 为神经元 \(i\) 在时间窗 \(t\) 的发放率）。通过以下流程提取拓扑特征：(i) 计算 \(N \times N\) 相关矩阵 \(C_{ij} = \operatorname{corr}(A_{i\cdot}, A_{j\cdot})\)；(ii) 定义距离 \(d_{ij} = 1 - |C_{ij}|\)；(iii) 在距离滤过下计算 VR 复形的持续同调。若环境具有环形拓扑（如圆形跑道），则 \(H_1\) 持续图应呈现单一显著特征，其持续性大小编码了神经表征的''环状''程度。

\textbf{定理 314.2}（环状拓扑的检测阈值）：设 \(\operatorname{Dgm}_1\) 是距离滤过下神经活动相关矩阵的 VR 复形持续图。在零假设（无环状拓扑结构）下，最大持续性 \(\max_i (d_i - b_i)\) 的分布可由 Bootstrap 重抽样逼近。若观测到的最大持续性超过 \(1-\alpha\) 分位数，则以显著性水平 \(\alpha\) 拒绝零假设，判定存在环状拓扑结构。Giusti 等人（2015）用此方法证实了啮齿动物海马体网格细胞的环状拓扑编码。

\textbf{神经元集群的拓扑编码}：Giusti等人（2015）研究了啮齿动物海马体网格细胞（grid cells）在空间导航中形成的群体编码。关键发现： （1）从神经活动的稀疏相关矩阵构造VR复形，持续同调揭示了环状拓扑结构（\(H_1\) Betti数 \(\beta_1 > 0\)），对应于圆形环境在神经表征中的拓扑编码； （2）拓扑特征在动物穿越不同环境时保持稳定，表明它们是空间表征的基本结构单元，而非偶然副产品；

（3）与经典群体向量方法相比，TDA方法对神经元的精确位置调谐曲线不敏感，仅依赖神经元共同活动的``粗略次序''关系，显示出对噪声和变异的高度鲁棒性。

\textbf{定义 314.4}（脑连接组的滤过构造）：设脑连接组由加权图 \(G = (V, E, w)\) 表示，其中 \(V\) 为脑区，\(w(e)\) 为连接强度。考虑滤过 \(\{G_\varepsilon\}_{\varepsilon > 0}\)：\(G_\varepsilon\) 包含所有 \(w(e) \geq \varepsilon\) 的边。计算其持续同调。\textbf{拓扑核心}定义为在多个 \(\varepsilon\) 值下持续存在的 \(H_1\) 特征。Betti 曲线 \(\beta_p(\varepsilon) = \dim H_p(G_\varepsilon)\) 提供了脑网络拓扑复杂度随连接强度变化的连续刻画。

\textbf{全脑连接组分析}：Sizemore等人（2018）应用持续同调分析人类大脑的结构和功能连接网络。通过将脑区视为节点、连接强度过滤后的网络在不同阈值上计算持续同调，识别出全脑尺度的``拓扑核心''结构------在多个阈值下持续存在的环状和空腔特征，这些结构与认知功能和神经疾病存在关联。

\hrulefill

\section{第379章：持续同调的应用与稳定性}\label{ux7b2c379ux7ae0ux6301ux7eedux540cux8c03ux7684ux5e94ux7528ux4e0eux7a33ux5b9aux6027}

持续同调的稳定性理论是拓扑数据分析的数学基石。Cohen-Steiner、Edelsbrunner和Harer于2007年建立的稳定性定理表明，持续图在Bottleneck距离下关于过滤函数的Lipschitz连续性成立------这是TDA具备统计稳健性的根本保证。本章在此基础上深入讨论Wasserstein距离与Bottleneck距离的关系、稳定性定理的完整证明框架，以及持续同调在统计推断中的极限理论。进一步地，我们将探讨Morse理论在函数数据分析中的应用，揭示Morse不等式与持续同调条形码之间的深层联系------Morse理论的临界点分类本质上对应着持续同调中特征的出生与消亡事件。这种联系不仅为TDA提供了微分拓扑的理论基础，也为从数据中提取Morse型分解提供了算法框架。

\subsection{持续同调的Wasserstein距离与瓶颈距离}\label{ux6301ux7eedux540cux8c03ux7684wassersteinux8dddux79bbux4e0eux74f6ux9888ux8dddux79bb}

\textbf{定义315.1（\(p\)-Wasserstein距离）} 设 \(\operatorname{Dgm}_1, \operatorname{Dgm}_2\) 是两个持续图。其\textbf{\(p\)-Wasserstein距离}（\(p \geq 1\)）定义为

\[W_p(\operatorname{Dgm}_1, \operatorname{Dgm}_2) = \left(\inf_{\gamma} \sum_{x \in \operatorname{Dgm}_1 \cup \Delta} \|x - \gamma(x)\|_{\infty}^p\right)^{1/p}\]

其中 \(\gamma\) 是 \(\operatorname{Dgm}_1 \cup \Delta\) 到 \(\operatorname{Dgm}_2 \cup \Delta\) 的双射，\(\Delta = \{(t, t)\}\) 是对角线。当 \(p = \infty\) 时，\(W_\infty\) 恰为Bottleneck距离 \(d_B\)。

\textbf{定理315.2（Bottleneck距离与Wasserstein距离的关系）} 对任意两个持续图 \(\operatorname{Dgm}_1, \operatorname{Dgm}_2\)，成立

\[d_B(\operatorname{Dgm}_1, \operatorname{Dgm}_2) = \lim_{p \to \infty} W_p(\operatorname{Dgm}_1, \operatorname{Dgm}_2)\]

且对任意 \(p \geq 1\)，\(d_B \leq W_p\)。若两个持续图中点的个数均有界，则Bottleneck距离给出了最宽容的度量（允许通过低代价匹配对角线来忽略小持续性特征）。

\textbf{证明} \(d_B = W_\infty\) 由定义直接验证：在 \(W_p\) 的定义中令 \(p \to \infty\)，\(\ell^p\) 范数收敛于 \(\ell^\infty\) 范数。\(\lim_{p \to \infty} W_p = W_\infty\) 由 \(\ell^p\) 范数的单调收敛性质及控制收敛定理保证。\(d_B \leq W_p\) 来自 \(\|\cdot\|_\infty \leq \|\cdot\|_p\) 对有限维向量的不等式。\(\blacksquare\)

\subsection{稳定性的Lipschitz条件完整证明}\label{ux7a33ux5b9aux6027ux7684lipschitzux6761ux4ef6ux5b8cux6574ux8bc1ux660e}

\textbf{定理315.3（Cohen-Steiner-Edelsbrunner-Harer稳定性定理）} 设 \(X\) 是三角剖分空间（或更一般的紧致可三角剖分空间），\(f, g: X \to \mathbb{R}\) 是驯服连续函数------即每个 \(f^{-1}((-\infty, a])\) 的同调群均为有限维且在有限个值处发生变化。则对任意 \(p \geq 0\)，

\[d_B(\operatorname{Dgm}_p(f), \operatorname{Dgm}_p(g)) \leq \|f - g\|_{\infty}\]

\textbf{证明} 设 \(\varepsilon = \|f - g\|_{\infty}\)，核心工具是持久模的代数交错理论。

\textbf{第1步（子水平集的包含关系）}：由 \(\|f - g\|_{\infty} \leq \varepsilon\) 知 \(f \leq g + \varepsilon\) 且 \(g \leq f + \varepsilon\)。因此对任意 \(a \in \mathbb{R}\)，

\[f^{-1}((-\infty, a]) \subseteq g^{-1}((-\infty, a + \varepsilon]), \quad g^{-1}((-\infty, a]) \subseteq f^{-1}((-\infty, a + \varepsilon])\]

这诱导了同调群之间的自然映射： \[F_a: H_p(f^{-1}(-\infty, a]) \to H_p(g^{-1}(-\infty, a + \varepsilon])\] \[G_a: H_p(g^{-1}(-\infty, a]) \to H_p(f^{-1}(-\infty, a + \varepsilon])\]

\textbf{第2步（\(\varepsilon\)-交错的构造）}：设 \(\mathbb{V} = \{H_p(f^{-1}(-\infty, a]))\}_a\) 和 \(\mathbb{W} = \{H_p(g^{-1}(-\infty, a]))\}_a\) 为两个持久模。对每个 \(a\)，有交换图：

\[V_a \to V_{a+2\varepsilon}, \quad V_a \to W_{a+\varepsilon} \to V_{a+2\varepsilon}\]

其中水平映射为包含诱导的同态，垂直映射为 \(F_a\) 和 \(G_{a+\varepsilon}\)。此即 \(\mathbb{V}\) 和 \(\mathbb{W}\) 的 \(\varepsilon\)-交错结构。

\textbf{第3步（代数稳定性定理）}：由Chazal等人（2009）的代数稳定性定理，若两个持久模是 \(\varepsilon\)-交错的，则它们对应的区间分解（条形码）中，区间端点之间的匹配在 \(L^\infty\) 范数下不超过 \(\varepsilon\)。具体地，对条形码中的每个区间 \([b, d)\)，与之配对的区间 \([b', d')\) 满足 \(|b - b'| \leq \varepsilon\) 且 \(|d - d'| \leq \varepsilon\)。此命题的证明基于将交错条件转化为持久模的态射锥（mapping cone）的零伦性，或等价地通过矩形容斥矩阵的秩不等式。

\textbf{第4步（结论）}：由代数交错至持续图距离的转化，Bottleneck距离 \(d_B(\operatorname{Dgm}_p(f), \operatorname{Dgm}_p(g)) \leq \varepsilon = \|f - g\|_{\infty}\)。此界是紧的：可构造两个偏差恰好为 \(\varepsilon\) 的Morse函数达到等号。\(\blacksquare\)

\textbf{推论315.4（稳定性对统计推断的意义）} 若 \(X_1, \ldots, X_n\) 是从某分布中独立同分布抽取的样本，由样本构造的持续图 \(\widehat{\operatorname{Dgm}}_n\) 与总体持续图 \(\operatorname{Dgm}(\mathbb{P})\) 之间的Bottleneck距离以速率 \(O(n^{-1/d})\) 收敛（\(d\) 为支撑集的维数），由经验过程理论和稳定性定理联合保证。这为TDA的统计推断提供了严格的理论基础。

\subsection{持续同调的统计理论}\label{ux6301ux7eedux540cux8c03ux7684ux7edfux8ba1ux7406ux8bba}

\textbf{定义315.5（持续图的期望与方差）} 设 \(\mathcal{D}\) 为持续图空间（配备Bottleneck距离构成的Polish空间）。\(\mathcal{D}\) 上的Borel概率测度 \(\mu\) 的\textbf{Fréchet均值}定义为

\[\mathbb{E}[\operatorname{Dgm}] = \arg\min_{D \in \mathcal{D}} \int_{\mathcal{D}} d_B(D, D')^2 \, d\mu(D')\]

由于 \(\mathcal{D}\) 在Bottleneck距离下具有复杂的几何结构（不满足CAT(0)条件），Fréchet均值可能不唯一。Turner等人（2014）证明了在适当的正则性条件下，期望持续图存在且唯一，并可通过求平均持续Landscape来近似。

\textbf{定理315.6（持续Landscape的极限定理）} 持续Landscape \(\lambda_k(t)\)（定义312.11）将持久模映射到Banach空间 \(L^2(\mathbb{R})\) 中的函数。由Bubenik（2015），该映射是1-Lipschitz的：

\[\|\lambda_k(f) - \lambda_k(g)\|_{\infty} \leq d_B(\operatorname{Dgm}(f), \operatorname{Dgm}(g)) \leq \|f - g\|_{\infty}\]

因此，若 \(\widehat{\lambda}_k^{(n)}(t)\) 是从 \(n\) 个独立同分布样本构造的样本Landscape，则 \(\sqrt{n}(\widehat{\lambda}_k^{(n)} - \lambda_k^\mu)\) 作为 \(L^2(\mathbb{R})\) 中的随机元弱收敛于Gauss过程，其协方差算子由总体Landscape的分布决定。这为构建Landscape的逐点置信带和同时置信带提供了渐近理论。

\textbf{定义315.7（持续图的方差）} 随机持续图 \(\operatorname{Dgm}\) 的Fréchet方差定义为

\[\operatorname{Var}(\operatorname{Dgm}) = \mathbb{E}[d_B(\operatorname{Dgm}, \mathbb{E}[\operatorname{Dgm}])^2] = \min_{D \in \mathcal{D}} \mathbb{E}[d_B(\operatorname{Dgm}, D)^2]\]

方差衡量了持续图围绕其Fréchet均值的离散程度，在Bottleneck Bootstrap方法中用于置信带的刻度。

\subsection{Morse理论在函数数据分析中的应用}\label{morseux7406ux8bbaux5728ux51fdux6570ux6570ux636eux5206ux6790ux4e2dux7684ux5e94ux7528}

\textbf{定义315.8（Morse函数与临界点）} 设 \(M\) 是光滑流形，\(f: M \to \mathbb{R}\) 是光滑函数。\(p \in M\) 称为 \(f\) 的\textbf{临界点}，若 \(df_p = 0\)。临界点 \(p\) 是\textbf{非退化的}，若 \(f\) 在 \(p\) 处的Hessian矩阵 \(H_f(p)\) 非奇异。非退化临界点的\textbf{指标} \(\lambda(p)\) 定义为 \(H_f(p)\) 的负特征值的个数。若 \(f\) 的所有临界点均为非退化的，则称 \(f\) 为\textbf{Morse函数}。

\textbf{定理315.9（Morse不等式）} 设 \(M\) 是紧致光滑流形，\(f: M \to \mathbb{R}\) 是Morse函数。令 \(C_k\) 为指标为 \(k\) 的临界点个数，\(\beta_k\) 为 \(M\) 的第 \(k\) 个Betti数。则

\[\sum_{k=0}^m (-1)^{m-k} C_k \geq \sum_{k=0}^m (-1)^{m-k} \beta_k, \quad m = 0, 1, \ldots, \dim M\]

且交替和取等号：\(\sum_{k=0}^{\dim M} (-1)^k C_k = \sum_{k=0}^{\dim M} (-1)^k \beta_k = \chi(M)\)（Euler示性数）。

\textbf{证明}（Morse同调框架）考虑子水平集 \(M_a = f^{-1}((-\infty, a])\)。当 \(a\) 通过临界值时，\(M_a\) 的同伦型发生变化：若 \(a\) 通过指标为 \(k\) 的临界点，则 \(M_a\) 同伦等价于 \(M_{a-\varepsilon}\) 附加一个 \(k\)-胞腔。由此得到 \(M\) 的一个胞腔分解，其 \(k\)-胞腔与指标 \(k\) 的临界点一一对应。相应的胞腔链复形 \((C_*^{\text{Morse}}, \partial)\) 的边界算子由负梯度流下的连接轨道计数定义，其同调自然同构于 \(M\) 的奇异同调。Morse不等式即该链复形的秩与Betti数之间的标准不等式：令 \(Z_k = \ker \partial_k\)，\(B_k = \operatorname{im} \partial_{k+1}\)，则 \(C_k = \dim Z_k + \dim B_{k-1}\)（因为 \(\dim C_k = \dim \ker \partial_k + \dim \operatorname{im} \partial_k\)），\(\beta_k = \dim Z_k - \dim B_k\)。由此推出Morse不等式。\(\blacksquare\)

\textbf{定理315.10（Morse理论与持续同调的对应）} 设 \(M\) 是紧致流形，\(f: M \to \mathbb{R}\) 是Morse函数。则 \(f\) 的子水平集滤过诱导的第 \(k\) 维持续同调条形码中，特征的出生时间恰为指标 \(k\) 的临界值，死亡时间为配对指标 \(k+1\) 的临界值。具体地，

\[\operatorname{Dgm}_k(f) \subseteq \{(f(p), f(q)) : p \text{ 指标 } k, q \text{ 指标 } k+1, f(p) < f(q)\}\]

且配对由Morse-Smale-Witten链复形的边界算子 \(\partial\) 决定：若 \(\partial[p] = \pm [q] + \cdots\)（\(\partial\) 作用于临界点的标准基），则 \((f(p), f(q))\) 属于持续图。

\textbf{证明} 由Morse理论，当 \(a\) 穿过 \(f(p)\)（\(p\) 指标 \(k\)）时，\(H_k(M_a)\) 的维数增加 \(1\)（同调类出生）。当 \(a\) 穿过 \(f(q)\)（\(q\) 指标 \(k+1\)）时，\(H_k(M_a)\) 的维数可能减少（同调类死亡），减少量由 \(\partial\) 在 \(q\) 处的非零系数决定。在持续同调的标准矩阵约化算法中，正单纯形（创建类）对应指标 \(k\) 的临界点，负单纯形（消去类）对应指标 \(k+1\) 的临界点，配对规则恰由边界矩阵的Smith正规形给出。因此持续同调的计算天然地重构了Morse函数的临界点配对信息。\(\blacksquare\)

\textbf{推论315.11（Morse函数的持续图统计）} 若 \(f\) 是从某Morse函数族中随机抽取的（如在数据背景下，距离函数 \(d_X(x) = \min_{x_i \in X} \|x - x_i\|\) 的期望），其持续图的期望行为由Morse临界点的期望分布决定。Baryshnikov和Yukich（2005）及后续工作利用Gauss随机场的Kac-Rice公式给出了随机Morse函数临界点密度的显式表达，这些结果与稳定性定理结合即可导出随机点云持续图的渐近期望和方差。

\newpage

\chapter{10-概率与统计/03-统计学/06-实验设计.md}\label{ux6982ux7387ux4e0eux7edfux8ba103-ux7edfux8ba1ux5b6606-ux5b9eux9a8cux8bbeux8ba1.md}

\chapter{实验设计}\label{ux5b9eux9a8cux8bbeux8ba1}

\section{第380章：实验设计基本原理}\label{ux7b2c380ux7ae0ux5b9eux9a8cux8bbeux8ba1ux57faux672cux539fux7406}

实验设计（Design of Experiments）是统计学的重要分支，其核心任务是在给定的资源约束下规划实验方案，以最有效的方式获取数据、估计处理效应并检验科学假设。R. A. Fisher于1920年代在Rothamsted农业实验站奠定了现代实验设计的理论基础，提出了随机化、重复和区组化三大基本原则。本章系统阐述这些基本原则及其数学论证，建立完全随机设计（CRD）的统计推断框架，并分析随机化在消除系统性偏差和控制混杂因素中的关键作用。实验设计的数学基础融合了线性代数、概率论和最优化理论，其理论的严谨性保证了实验结论的因果解释力。

\subsection{x.1 三大基本原则}\label{x.1-ux4e09ux5927ux57faux672cux539fux5219}

\textbf{定义 x.1}（实验单元、处理与响应）：实验单元是接受处理的最小独立实体。处理是施加于实验单元上的条件或干预。响应变量是实验单元在接受处理后测得的感兴趣的结果。

\textbf{定理 x.2}（随机化的必要性）：设 \(Y_{ij} = \mu + \tau_i + \varepsilon_{ij}\) 为处理 \(i\) 下第 \(j\) 个实验单元的响应，其中 \(\tau_i\) 为处理效应。若未施行随机化，则任意未知的系统性偏差 \(\theta_{ij}\) 可能与 \(\tau_i\) 混杂，导致处理效应的估计有偏。

\textbf{证明}：假设实验单元按某种未观测的次序排列，且存在系统性漂移 \(\theta_{ij} = \beta \cdot s_{ij}\)，其中 \(s_{ij}\) 为未知的排列序数。若不随机化，处理分配可能与 \(s_{ij}\) 相关，则

\[\mathbb{E}[\hat{\tau}_i] = \tau_i + \frac{1}{n_i}\sum_{j=1}^{n_i} \theta_{ij},\]

其中求和项一般非零。随机化使处理分配与 \(\theta_{ij}\) 独立，从而 \(\mathbb{E}[\sum \theta_{ij}/n_i] = 0\)，保证了无偏性。 \(\blacksquare\)

\textbf{定理 x.3}（基于随机化的推断------Fisher随机化检验）：在完全随机设计下，处理分配的所有 \(\binom{N}{n_1, \ldots, n_t}\) 种等可能排列构成精确的参考分布。设 \(T\) 为检验统计量，则在Sharp null \(H_0: \tau_1 = \cdots = \tau_t\) 下，\(T\) 的精确 \(p\) 值为

\[p = \frac{|\{\pi : T(\pi) \geq T_{\text{obs}}\}|}{\binom{N}{n_1, \ldots, n_t}}.\]

\textbf{证明}：在 \(H_0\) 下，每个实验单元的响应 \(Y_{ij}\) 仅与随机误差相关，与接受何种处理无关。处理分配 \(\pi\) 服从超几何分布在所有 \(\binom{N}{n_1, \ldots, n_t}\) 种排列上均匀分布。观测统计量 \(T_{\text{obs}}\) 来自特定的实际排列。\(p\) 值即为在均匀排列分布下，\(T\) 超过（或等于）\(T_{\text{obs}}\) 的概率，由计数比值得出。 \(\blacksquare\)

\textbf{定理 x.4}（重复的必要性------误差方差估计）：设 \(n_i\) 为处理 \(i\) 的重复次数。处理效应差 \(\tau_i - \tau_j\) 的最小二乘估计的方差为 \(\sigma^2(1/n_i + 1/n_j)\)。重复次数越大，估计精度越高。

\textbf{证明}：设 \(\bar{Y}_{i\cdot} = \frac{1}{n_i}\sum_{j=1}^{n_i} Y_{ij}\)。在独立性假设下，\(\operatorname{Var}(\bar{Y}_{i\cdot}) = \sigma^2/n_i\)。由于不同处理的观测独立，

\[\operatorname{Var}(\hat{\tau}_i - \hat{\tau}_j) = \operatorname{Var}(\bar{Y}_{i\cdot} - \bar{Y}_{j\cdot}) = \frac{\sigma^2}{n_i} + \frac{\sigma^2}{n_j}.\]

当 \(n_i, n_j \to \infty\) 时，方差趋于零，估计相合。 \(\blacksquare\)

\subsection{x.2 完全随机设计（CRD）}\label{x.2-ux5b8cux5168ux968fux673aux8bbeux8ba1crd}

\textbf{模型 x.5}（单因素CRD模型）：设 \(t\) 个处理，处理 \(i\) 分配 \(n_i\) 个实验单元（\(i = 1, \ldots, t\)）。模型为

\[Y_{ij} = \mu + \tau_i + \varepsilon_{ij}, \quad j = 1, \ldots, n_i,\]

其中 \(\varepsilon_{ij} \stackrel{\text{iid}}{\sim} N(0, \sigma^2)\)，\(\sum_{i=1}^{t} n_i \tau_i = 0\)（可识别性约束）。

\textbf{定理 x.6}（处理效应的BLUE）：在上述模型下，\(\mu\) 和 \(\tau_i\) 的最佳线性无偏估计（BLUE）为

\[\hat{\mu} = \bar{Y}_{\cdot\cdot}, \quad \hat{\tau}_i = \bar{Y}_{i\cdot} - \bar{Y}_{\cdot\cdot},\]

其中 \(\bar{Y}_{i\cdot} = \frac{1}{n_i}\sum_{j} Y_{ij}\)，\(\bar{Y}_{\cdot\cdot} = \frac{1}{N}\sum_{i,j} Y_{ij}\)，\(N = \sum n_i\)。

\textbf{证明}：由Gauss-Markov定理，在满秩线性模型 \(Y = X\beta + \varepsilon\) 中，\(\operatorname{Cov}(\varepsilon) = \sigma^2 I\) 时，\(\hat{\beta} = (X^T X)^{-1} X^T Y\) 为BLUE。对CRD模型，设计矩阵 \(X\) 为 \(N \times (t+1)\) 矩阵，在约束 \(\sum n_i \tau_i = 0\) 下满秩。直接计算得上述估计。验证无偏性：

\[\mathbb{E}[\hat{\mu}] = \frac{1}{N}\sum_{i,j} \mathbb{E}[Y_{ij}] = \frac{1}{N}\sum_i n_i(\mu + \tau_i) = \mu + \frac{1}{N}\sum_i n_i \tau_i = \mu,\]

\[\mathbb{E}[\hat{\tau}_i] = \mathbb{E}[\bar{Y}_{i\cdot} - \bar{Y}_{\cdot\cdot}] = (\mu + \tau_i) - \mu = \tau_i. \quad \blacksquare\]

\textbf{定理 x.7}（方差分析平方和分解）：在CRD模型中，总平方和分解为

\[\sum_{i=1}^{t}\sum_{j=1}^{n_i}(Y_{ij} - \bar{Y}_{\cdot\cdot})^2 = \sum_{i=1}^{t} n_i(\bar{Y}_{i\cdot} - \bar{Y}_{\cdot\cdot})^2 + \sum_{i=1}^{t}\sum_{j=1}^{n_i}(Y_{ij} - \bar{Y}_{i\cdot})^2.\]

即 \(SS_T = SS_{\text{Trt}} + SS_E\)。

\textbf{证明}：

\[\begin{aligned}
Y_{ij} - \bar{Y}_{\cdot\cdot} &= (Y_{ij} - \bar{Y}_{i\cdot}) + (\bar{Y}_{i\cdot} - \bar{Y}_{\cdot\cdot}).\\
\sum_{i,j}(Y_{ij} - \bar{Y}_{\cdot\cdot})^2 &= \sum_{i,j}(Y_{ij} - \bar{Y}_{i\cdot})^2 + \sum_{i,j}(\bar{Y}_{i\cdot} - \bar{Y}_{\cdot\cdot})^2 \\
&\quad + 2\sum_{i,j}(Y_{ij} - \bar{Y}_{i\cdot})(\bar{Y}_{i\cdot} - \bar{Y}_{\cdot\cdot}).
\end{aligned}\]

交叉项为零，因为 \(\sum_{j=1}^{n_i}(Y_{ij} - \bar{Y}_{i\cdot}) = 0\)。故得证。 \(\blacksquare\)

\hrulefill

\section{第381章：方差分析与因素设计}\label{ux7b2c381ux7ae0ux65b9ux5deeux5206ux6790ux4e0eux56e0ux7d20ux8bbeux8ba1}

方差分析（ANOVA）是实验设计中分析处理效应的核心工具。当实验中包含两个或更多因素时，因素设计不仅能估计各因素的主动效应，还能识别因素间的交互作用，这是通过单因素序列实验无法获得的重要信息。本章从单因素ANOVA的假设检验理论出发，系统发展为双因素和多因素ANOVA的完整框架，建立平方和的Cochran分解定理，讨论固定效应、随机效应和混合效应模型的估计与检验策略，并分析交互作用的几何含义及其在科学研究中的实质意义。

\subsection{x+1.1 单因素ANOVA}\label{x1.1-ux5355ux56e0ux7d20anova}

\textbf{定理 x+1.1}（单因素ANOVA的F检验）：在模型 \(Y_{ij} = \mu + \tau_i + \varepsilon_{ij}\)，\(\varepsilon_{ij} \stackrel{\text{iid}}{\sim} N(0, \sigma^2)\) 下，检验 \(H_0: \tau_1 = \cdots = \tau_t = 0\) 的F统计量为

\[F = \frac{SS_{\text{Trt}}/(t-1)}{SS_E/(N-t)} \sim F_{t-1, N-t} \quad \text{（在 } H_0 \text{ 下）}.\]

\textbf{证明}：由Cochran定理，\(SS_E/\sigma^2 \sim \chi^2_{N-t}\)，且在 \(H_0\) 下，\(SS_{\text{Trt}}/\sigma^2 \sim \chi^2_{t-1}\)，二者独立。因此

\[F = \frac{(SS_{\text{Trt}}/\sigma^2)/(t-1)}{(SS_E/\sigma^2)/(N-t)} \sim F_{t-1, N-t}.\]

独立性由正态分布下二次型独立性的Craig-Sakamoto型定理保证：\(SS_{\text{Trt}}\) 和 \(SS_E\) 分别对应相互正交的子空间上的投影。 \(\blacksquare\)

\textbf{定理 x+1.2}（单因素ANOVA的非中心F分布与功效）：在备择假设 \(H_1\) 下，\(F \sim F_{t-1, N-t}(\lambda)\)，其中非中心参数

\[\lambda = \frac{1}{\sigma^2}\sum_{i=1}^{t} n_i \tau_i^2.\]

功效函数 \(\beta(\lambda) = P(F > F_{\alpha; t-1, N-t} \mid \lambda)\) 关于 \(\lambda\) 单调递增。

\textbf{证明}：在 \(H_1\) 下，\(\bar{Y}_{i\cdot} \sim N(\mu + \tau_i, \sigma^2/n_i)\)。因此

\[\frac{SS_{\text{Trt}}}{\sigma^2} = \frac{1}{\sigma^2}\sum_{i=1}^{t} n_i(\bar{Y}_{i\cdot} - \bar{Y}_{\cdot\cdot})^2 \sim \chi^2_{t-1}(\lambda),\]

其中非中心参数为 \(\lambda = \frac{1}{\sigma^2}\sum_i n_i(\mathbb{E}[\bar{Y}_{i\cdot}] - \mathbb{E}[\bar{Y}_{\cdot\cdot}])^2 = \frac{1}{\sigma^2}\sum_i n_i \tau_i^2\)。故 \(F\) 服从非中心F分布，功效单调性由非中心F分布的性质保证。 \(\blacksquare\)

\subsection{x+1.2 双因素ANOVA与交互作用}\label{x1.2-ux53ccux56e0ux7d20anovaux4e0eux4ea4ux4e92ux4f5cux7528}

\textbf{模型 x+1.3}（双因素交叉设计模型）：设有因素 \(A\)（\(a\) 水平）和因素 \(B\)（\(b\) 水平），每组合有 \(n\) 次重复。模型为

\[Y_{ijk} = \mu + \alpha_i + \beta_j + (\alpha\beta)_{ij} + \varepsilon_{ijk},\]

其中 \(i = 1, \ldots, a\)，\(j = 1, \ldots, b\)，\(k = 1, \ldots, n\)，\(\varepsilon_{ijk} \stackrel{\text{iid}}{\sim} N(0, \sigma^2)\)。约束条件：\(\sum_i \alpha_i = 0\)，\(\sum_j \beta_j = 0\)，\(\sum_i (\alpha\beta)_{ij} = \sum_j (\alpha\beta)_{ij} = 0\)。

\textbf{定理 x+1.4}（双因素ANOVA平方和分解）：总平方和可正交分解为

\[SS_T = SS_A + SS_B + SS_{AB} + SS_E,\]

其中 \(SS_A = bn\sum_i (\bar{Y}_{i\cdot\cdot} - \bar{Y}_{\cdot\cdot\cdot})^2\)，\(SS_B = an\sum_j (\bar{Y}_{\cdot j\cdot} - \bar{Y}_{\cdot\cdot\cdot})^2\)，\(SS_{AB} = n\sum_{i,j}(\bar{Y}_{ij\cdot} - \bar{Y}_{i\cdot\cdot} - \bar{Y}_{\cdot j\cdot} + \bar{Y}_{\cdot\cdot\cdot})^2\)，\(SS_E\) 为误差平方和。

\textbf{证明}：将观测分解为

\[Y_{ijk} - \bar{Y}_{\cdot\cdot\cdot} = (\bar{Y}_{i\cdot\cdot} - \bar{Y}_{\cdot\cdot\cdot}) + (\bar{Y}_{\cdot j\cdot} - \bar{Y}_{\cdot\cdot\cdot}) + (\bar{Y}_{ij\cdot} - \bar{Y}_{i\cdot\cdot} - \bar{Y}_{\cdot j\cdot} + \bar{Y}_{\cdot\cdot\cdot}) + (Y_{ijk} - \bar{Y}_{ij\cdot}).\]

平方后求和，六类交叉项均为零（由均值差的定义及求和约束），得四项平方和。在正态假设下，\(SS_A/\sigma^2 \sim \chi^2_{a-1}(\lambda_A)\)（在 \(H_0^A\) 下），\(SS_B/\sigma^2 \sim \chi^2_{b-1}(\lambda_B)\)，\(SS_{AB}/\sigma^2 \sim \chi^2_{(a-1)(b-1)}(\lambda_{AB})\)，\(SS_E/\sigma^2 \sim \chi^2_{ab(n-1)}\)，且四项独立。 \(\blacksquare\)

\textbf{定理 x+1.5}（交互作用的可解释性）：在双因素模型中，若交互作用 \((\alpha\beta)_{ij}\) 非零，则因素 \(A\) 的效应依赖于因素 \(B\) 的水平（反之亦然）。具体地，

\[(\alpha\beta)_{ij} = (\mu_{ij} - \mu_{i\cdot}) - (\mu_{\cdot j} - \mu),\]

即交互作用度量单元格均值偏离可加性（主效应叠加）的程度。

\textbf{证明}：在无交互作用模型下，\(\mu_{ij} = \mu + \alpha_i + \beta_j\)，此时 \((\alpha\beta)_{ij} = 0\)。定义交互作用为 \(\mu_{ij} - (\mu + \alpha_i + \beta_j)\)，这正是上述表达式。\(A\) 在 \(B_j\) 下的简单效应为 \(\mu_{ij} - \mu_{i'j} = (\alpha_i - \alpha_{i'}) + [(\alpha\beta)_{ij} - (\alpha\beta)_{i'j}]\)，当交互项非零时，简单效应随 \(j\) 变化。 \(\blacksquare\)

\subsection{x+1.3 多因素设计}\label{x1.3-ux591aux56e0ux7d20ux8bbeux8ba1}

\textbf{定义 x+1.6}（\(2^k\) 因子设计）：在 \(k\) 个因素、每因素两个水平的全因子设计中，共有 \(2^k\) 个处理组合。主效应和交互作用通过对比系数 \(\pm 1\) 的线性组合估计。

\textbf{定理 x+1.7}（效应估计的正交性）：在 \(2^k\) 设计中，任意两个不同效应的对比向量是正交的。因此，各效应的估计互不相关，总平方和可分解为 \(2^k - 1\) 个自由度各为 1 的独立分量。

\textbf{证明}：每个效应对应一个长度为 \(2^k\) 的对比向量 \(c = (c_1, \ldots, c_{2^k})\)，其中 \(c_j = \pm 1\) 且 \(\sum_j c_j = 0\)（主效应）或更复杂的符号模式（交互作用）。对于两个不同的效应 \(c^{(1)}\) 和 \(c^{(2)}\)，其内积为

\[\langle c^{(1)}, c^{(2)} \rangle = \sum_{j=1}^{2^k} c^{(1)}_j c^{(2)}_j.\]

由 \(2^k\) 设计的代数结构，这等价于在 \(\mathbb{F}_2^k\) 上的群特征的正交性：\(\sum_{x \in \mathbb{F}_2^k} (-1)^{\langle a, x \rangle + \langle b, x \rangle} = 0\) 当 \(a \neq b\)。因此各效应对比向量正交，对应二次型独立。 \(\blacksquare\)

\textbf{定理 x+1.8}（效应稀疏性与分级原则）：在实际实验中，高阶交互作用通常可忽略（效应稀疏性）。效应分级原则表明：低阶效应的显著性优先于高阶效应，即若 \(k\) 阶交互作用显著，则其涉及的所有 \((k-1)\) 阶子交互作用应保留在模型中（无论其自身是否显著）。

\textbf{证明}：效应分级原则由模型的层级结构保证。若模型包含 \(\alpha\beta\gamma\)（三阶交互作用）但不含 \(\alpha\beta\)（二阶交互作用），则模型对 \(\alpha\beta\) 的边际解释依赖于 \(\gamma\) 的具体水平，这违反了统计推断中的边际性原则。形式化地，若 \(Y_{ijk} = \mu + \alpha_i + \beta_j + \gamma_k + (\alpha\beta\gamma)_{ijk}\)，则不同 \(\gamma_k\) 水平下的 \(\alpha\beta\) 交互模式一般不同，丧失了可解释性。因此，包含高阶交互作用时必须同时包含其所有低阶组成成分。 \(\blacksquare\)

\hrulefill

\section{第382章：拉丁方与不完全区组设计}\label{ux7b2c382ux7ae0ux62c9ux4e01ux65b9ux4e0eux4e0dux5b8cux5168ux533aux7ec4ux8bbeux8ba1}

当实验单元数量受限或存在两个已知的干扰因素时，拉丁方设计通过行列双重区组化实现了对处理效应的有效分离。本章系统研究拉丁方设计及其正交化的正交拉丁方理论，引入以Bose-Shrikhande-Parker定理为核心的完备正交拉丁方组存在性结果。此外，当每个区组内的实验单元数少于处理数时，平衡不完全区组设计（BIBD）提供了处理比较的公平框架，其组合结构保证了所有处理对比的方差相等，具有深刻的设计最优性。Youden方设计则是拉丁方与BIBD的优雅结合，适用于存在行列双重限制的实验场景。

\subsection{x+2.1 拉丁方设计}\label{x2.1-ux62c9ux4e01ux65b9ux8bbeux8ba1}

\textbf{定义 x+2.1}（拉丁方）：\(n \times n\) 阵列称为 \(n\) 阶拉丁方，若 \(n\) 个不同符号中的每一个在每行和每列中恰出现一次。拉丁方设计将处理分配为拉丁方的符号，行和列分别表示两个区组因素。

\textbf{模型 x+2.2}（拉丁方设计模型）：设 \(n \times n\) 拉丁方中行 \(i\)、列 \(j\) 处施加处理 \(k = L(i, j)\)。模型为

\[Y_{ijk} = \mu + \rho_i + \gamma_j + \tau_k + \varepsilon_{ijk}, \quad \varepsilon_{ijk} \stackrel{\text{iid}}{\sim} N(0, \sigma^2),\]

约束 \(\sum \rho_i = \sum \gamma_j = \sum \tau_k = 0\)。

\textbf{定理 x+2.3}（拉丁方设计的ANOVA表）：平方和分解为

\[SS_T = SS_{\text{Row}} + SS_{\text{Col}} + SS_{\text{Trt}} + SS_E,\]

自由度分别为 \(n-1, n-1, n-1, (n-1)(n-2)\)。

\textbf{证明}：总观测数 \(N = n^2\)。模型参数：\(\mu\)（1个），\(\rho_i\)（\(n-1\) 个独立参数），\(\gamma_j\)（\(n-1\) 个独立参数），\(\tau_k\)（\(n-1\) 个独立参数），共 \(1 + 3(n-1) = 3n - 2\) 个参数。因此误差自由度为 \(n^2 - (3n-2) = (n-1)(n-2)\)。平方和的构造与双因素ANOVA类似，利用拉丁方各符号在每行每列恰出现一次的性质验证交叉项为零。处理对比的方差为 \(2\sigma^2/n\)（两个不同处理比较时），此即拉丁方设计提高了精度（比完全随机设计的 \(2\sigma^2/r\) 减少了行列变异）。 \(\blacksquare\)

\subsection{x+2.2 正交拉丁方与Bose-Shrikhande-Parker定理}\label{x2.2-ux6b63ux4ea4ux62c9ux4e01ux65b9ux4e0ebose-shrikhande-parkerux5b9aux7406}

\textbf{定义 x+2.4}（正交拉丁方）：两个 \(n\) 阶拉丁方 \(L_1\) 和 \(L_2\) 称为正交的，若将二者叠置后，\(n^2\) 个有序对 \((L_1(i,j), L_2(i,j))\) 互不相同。\(r\) 个两两正交的 \(n\) 阶拉丁方构成完备组时 \(r = n-1\)。

\textbf{定理 x+2.5}（正交拉丁方个数的上界）：\(n\) 阶两两正交拉丁方（MOLS）的个数不超过 \(n-1\)。

\textbf{证明}：设有 \(r\) 个MOLS \(L_1, \ldots, L_r\)。通过对 \(L_1\) 的行列重排，可设每方第一行为 \((1, 2, \ldots, n)\)。考虑各方的 \((2,1)\)-位置。对每方，该位置值不能为 1（因第一列已有 1）。且不同两方的 \((2,1)\)-位置值必不同（否则有序对重复（发生在 \((1,1)\)-位置处）。因此 \(r \leq n-1\)。 \(\blacksquare\)

\textbf{定理 x+2.6}（Bose-Shrikhande-Parker定理, 1960）：当 \(n \geq 3\) 且 \(n \neq 6\) 时，不存在完备的 \(n-1\) 个\(n\)阶MOLS当且仅当不存在 \(n\) 阶射影平面（等价于 \(n\) 不为素数幂时一般不存在完备组）。特别地，10阶射影平面不存在，故不存在9个10阶MOLS的完备组。Euler的36军官问题（\(n=6\)）确实无解。

\textbf{证明概要}：该定理的完全证明依赖于Bruck-Ryser-Chowla定理和有限域上的构造。当 \(n\) 为素数幂时，利用有限域 \(\mathbb{F}_q\)（\(q = n\)）的加法表和乘法表可构造 \(n-1\) 个MOLS：\(L_k(i, j) = i + k \cdot j \pmod{n}\)（\(k = 1, \ldots, n-1\)）。当 \(n = 6\) 时，Tarry（1901）通过穷举验证无解；\(n = 10\) 时，Lam等人（1989）通过大规模计算机搜索排除了10阶射影平面的存在性，从而不存在9个10阶MOLS。一般 \(n\) 非素数幂的情况涉及深刻的组合设计存在性理论。 \(\blacksquare\)

\subsection{x+2.3 平衡不完全区组设计（BIBD）}\label{x2.3-ux5e73ux8861ux4e0dux5b8cux5168ux533aux7ec4ux8bbeux8ba1bibd}

\textbf{定义 x+2.7}（BIBD）：平衡不完全区组设计 \((v, b, r, k, \lambda)\) 配置满足：(i) \(v\) 个处理；(ii) \(b\) 个区组，每区组大小 \(k < v\)；(iii) 每处理在 \(r\) 个区组中出现；(iv) 任意两个不同处理恰在 \(\lambda\) 个区组中同时出现。

\textbf{定理 x+2.8}（BIBD参数关系）：BIBD的参数满足：(i) \(vr = bk\)（总观测数）；(ii) \(\lambda(v-1) = r(k-1)\)（配对平衡条件）；(iii) \(b \geq v\)（Fisher不等式）。

\textbf{证明}：(i) 计数所有（处理，区组）对：处理视角共 \(vr\) 对，区组视角共 \(bk\) 对，二者相等。(ii) 固定处理 \(T_1\)。\(T_1\) 出现在 \(r\) 个区组中，每区组含 \(k-1\) 个其他处理，共 \(r(k-1)\) 个位置。另一方面，其余 \(v-1\) 个处理各与 \(T_1\) 同时出现 \(\lambda\) 次，共 \(\lambda(v-1)\) 次出现。二者相等。

\begin{enumerate}
\def\labelenumi{(\roman{enumi})}
\setcounter{enumi}{2}
\tightlist
\item
  设 \(N\) 为 \(v \times b\) 关联矩阵（\(N_{ij} = 1\) 若处理 \(i\) 在区组 \(j\) 中，否则为 0）。则 \(NN^T = (r-\lambda)I_v + \lambda J_v\)，其特征值为 \(rk\)（重数1）和 \(r-\lambda\)（重数 \(v-1\)）。因 \(r > \lambda\)（否则 \(k = v\)），\(NN^T\) 满秩，故 \(\operatorname{rank}(N) = v \leq b\)。 \(\blacksquare\)
\end{enumerate}

\textbf{定理 x+2.9}（BIBD中处理效应的最优估计）：在BIBD模型 \(Y_{ij} = \mu + \tau_i + \beta_j + \varepsilon_{ij}\) 中，任意处理对比 \(\sum c_i \tau_i\)（\(\sum c_i = 0\)）的BLUE的方差为

\[\operatorname{Var}\left(\sum_i c_i \hat{\tau}_i\right) = \frac{k}{\lambda v} \sigma^2 \sum_i c_i^2.\]

\textbf{证明}：处理效应的最小二乘正规方程为 \(r\hat{\tau}_i - \frac{\lambda}{k}\sum_{h \neq i} \hat{\tau}_h = Q_i\)，其中 \(Q_i\) 为调整后的处理合计。利用约束 \(\sum \hat{\tau}_i = 0\) 和 \(r(k-1) = \lambda(v-1)\)，可解得

\[\hat{\tau}_i = \frac{k}{\lambda v} Q_i.\]

由此得 \(\operatorname{Cov}(\hat{\tau}_i, \hat{\tau}_j) = \frac{k\sigma^2}{\lambda v}(\delta_{ij} - \frac{1}{v})\)。代入对比方差公式即得结果。该方差与完全随机设计（每处理 \(r\) 次重复）相比，效率因子为 \(E = \frac{\lambda v}{rk}\)，BIBD的平衡性保证了所有对比的效率因子相同。 \(\blacksquare\)

\subsection{x+2.4 Youden方设计}\label{x2.4-youdenux65b9ux8bbeux8ba1}

\textbf{定义 x+2.10}（Youden方）：\(k \times v\) 矩形阵列称为Youden方，若其各列构成一个BIBD\((v, v, k, r, \lambda)\) 的区组，且各行构成区组大小为 \(k\) 的BIBD或平衡设计。Youden方本质上是拉丁方移除若干行（或列）的结果。

\textbf{定理 x+2.11}（Youden方与BIBD的对偶性）：存在 \(k \times v\) Youden方当且仅当存在BIBD\((v, v, k, r, \lambda)\)。Youden方的列与BIBD的区组一一对应，行与处理在区组群中的位置对应。

\textbf{证明}：由Youden方的定义，将各列视为区组即得BIBD。反之，给定BIBD，将其区组安排的关联矩阵写为 \(N\)（\(v \times b\)）。由BIBD参数的对称性，若 \(b = v\)（对称BIBD），则 \(N\) 的各行也是区组（每行恰 \(k\) 个 1）。将 \(N\) 中的 1 替换为适当的处理符号使每符号每列恰出现一次，即得Youden方。构造的可行性由对称BIBD的关联矩阵的排列性质保证。 \(\blacksquare\)

\hrulefill

\section{第383章：响应面与最优设计}\label{ux7b2c383ux7ae0ux54cdux5e94ux9762ux4e0eux6700ux4f18ux8bbeux8ba1}

响应面方法（Response Surface Methodology, RSM）由Box和Wilson于1951年提出，是将实验设计与最优化技术结合的系统方法论。当实验目标不仅是估计处理效应，而是探索响应变量与若干定量因素间的函数关系并寻找最优操作条件时，响应面方法提供了从一阶筛选、最速上升路径到二阶曲面拟合的完整流程。本章重点讨论二次响应曲面的标准设计------Box-Behnken设计和中心复合设计------的统计性质，继而深入最优设计理论，以Kiefer-Wolfowitz等价性定理为核心，建立D-最优性准则的信息矩阵几何解释和迭代构造算法，展示实验设计理论与凸优化的深刻联系。

\subsection{x+3.1 响应面方法与二次曲面拟合}\label{x3.1-ux54cdux5e94ux9762ux65b9ux6cd5ux4e0eux4e8cux6b21ux66f2ux9762ux62dfux5408}

\textbf{模型 x+3.1}（二阶响应曲面模型）：设 \(x = (x_1, \ldots, x_k)^T\) 为 \(k\) 个定量因素。二阶模型为

\[Y = \beta_0 + \sum_{i=1}^{k} \beta_i x_i + \sum_{i=1}^{k} \beta_{ii} x_i^2 + \sum_{i < j} \beta_{ij} x_i x_j + \varepsilon, \quad \varepsilon \sim N(0, \sigma^2).\]

参数总数为 \(p = 1 + k + k + \binom{k}{2} = (k+1)(k+2)/2\)。

\textbf{定理 x+3.2}（矩矩阵与设计的最优性准则）：设设计 \(\xi\) 将权重 \(w_m\) 赋予设计点 \(x_m \in \mathbb{R}^k\)（\(m = 1, \ldots, N\)）。记模型向量 \(f(x) = (1, x_1, \ldots, x_k, x_1^2, \ldots, x_k^2, x_1 x_2, \ldots)^T\)。信息矩阵为

\[M(\xi) = \sum_{m=1}^{N} w_m f(x_m) f(x_m)^T.\]

\(\hat{\beta}\) 的协方差矩阵为 \(\sigma^2 M(\xi)^{-1}\)。D-最优设计最大化 \(\det M(\xi)\)。

\textbf{证明}：对线性模型 \(Y = F\beta + \varepsilon\)，设计矩阵 \(F\) 的行即为 \(f(x_m)^T\)。由Gauss-Markov定理，

\[\operatorname{Cov}(\hat{\beta}) = \sigma^2 (F^T F)^{-1} = \sigma^2 \left(\sum_{m} f(x_m) f(x_m)^T\right)^{-1} = \sigma^2 M(\xi)^{-1}.\]

\(\det \operatorname{Cov}(\hat{\beta}) = \sigma^{2p} / \det M(\xi)\)，因此最大化 \(\det M(\xi)\) 等价于最小化置信椭球的体积，此即D-最优性的直观含义。 \(\blacksquare\)

\subsection{x+3.2 Box-Behnken设计与中心复合设计}\label{x3.2-box-behnkenux8bbeux8ba1ux4e0eux4e2dux5fc3ux590dux5408ux8bbeux8ba1}

\textbf{定义 x+3.3}（Box-Behnken设计, BBD）：Box-Behnken设计是用于拟合二阶响应曲面的三水平设计，通过组合二水平因子设计和平衡不完全区组设计构造。其设计点位于超球面上和一个中心点，运行次数少于相应的 \(3^k\) 全因子设计。

\textbf{定理 x+3.4}（BBD的可旋转性条件）：当 \(k \geq 3\) 时，Box-Behnken设计是近可旋转的。严格的可旋转性要求信息矩阵 \(M(\xi)\) 在任意正交变换 \(x \mapsto Qx\) 下不变。BBD的信息矩阵满足

\[M_{ij}(\xi) = \sum_{m} w_m x_{mi} x_{mj} = 0 \quad (i \neq j),\]

且 \(\sum_m w_m x_{mi}^2\) 对所有 \(i\) 相等，\(\sum_m w_m x_{mi}^4\) 对所有 \(i\) 相等。

\textbf{证明}：Box-Behnken设计通过在每个BIBD区组中将未出现的因子设为零水平，其余因子取 \(\pm 1\) 构造。\(k=3\) 时设计点为 \((\pm 1, \pm 1, 0)\) 及其置换，共12个边心点加中心点。由构造的对称性，所有一次项为 \(x_i^2\) 的求和相等，所有交叉项求和为零（奇函数在对称区域上的积分为零），四次矩也相等。这满足可旋转性的必要条件。当 \(k > 3\) 时，额外的中心点有助于改善矩条件。 \(\blacksquare\)

\textbf{定义 x+3.5}（中心复合设计, CCD）：CCD由三部分组成：(i) \(n_f\) 个立方体点（\(2^{k-p}\) 部分因子设计）；(ii) \(n_a = 2k\) 个轴点 \((\pm \alpha, 0, \ldots, 0)\) 及其置换；(iii) \(n_c\) 个中心点。总运行次数 \(N = n_f + 2k + n_c\)。

\textbf{定理 x+3.6}（CCD可旋转性的 \(\alpha\) 选取）：CCD为可旋转设计的充要条件是轴距 \(\alpha\) 满足

\[\alpha = (n_f)^{1/4}.\]

此时信息矩阵中 \(\sum_m x_{mi}^4 / (\sum_m x_{mi}^2)^2 = 3\)，保证了在任意正交旋转下 \(\operatorname{Var}(\hat{Y}(x))\) 仅依赖于 \(x\) 到中心的距离。

\textbf{证明}：可旋转性的代数条件是信息矩满足：对任意 \(i \neq j\)，\(\sum_m x_{mi}^2 x_{mj}^2 = \frac{1}{3}\sum_m x_{mi}^4\)。设 \(\sum_m x_{mi}^2 = n_f + 2\alpha^2 = S_2\)，\(\sum_m x_{mi}^4 = n_f + 2\alpha^4 = S_4\)，\(\sum_m x_{mi}^2 x_{mj}^2 = n_f = S_{22}\)。可旋转性要求 \(3S_{22} = S_4\)，即 \(3n_f = n_f + 2\alpha^4\)，解得 \(\alpha = n_f^{1/4}\)。 \(\blacksquare\)

\subsection{x+3.3 D-最优性与Kiefer-Wolfowitz等价性定理}\label{x3.3-d-ux6700ux4f18ux6027ux4e0ekiefer-wolfowitzux7b49ux4ef7ux6027ux5b9aux7406}

\textbf{定义 x+3.7}（连续设计与信息矩阵）：连续（近似）设计 \(\xi\) 是设计空间 \(\mathcal{X}\) 上的概率测度，支撑于有限个点。信息矩阵为

\[M(\xi) = \int_{\mathcal{X}} f(x) f(x)^T d\xi(x).\]

D-最优设计 \(\xi^*\) 最大化 \(\Phi(\xi) = \log \det M(\xi)\)（或等价地最大化 \(\det M(\xi)\)）。

\textbf{定理 x+3.8}（Kiefer-Wolfowitz等价性定理, 1960）：设计 \(\xi^*\) 为 D-最优的充要条件是：对任意 \(x \in \mathcal{X}\)，

\[d(x, \xi^*) = f(x)^T M(\xi^*)^{-1} f(x) \leq p,\]

其中 \(p\) 为参数个数。等号在 \(\xi^*\) 的支撑点上成立。此条件等价于 \(\xi^*\) 最小化 \(\max_{x \in \mathcal{X}} d(x, \xi)\)。

\textbf{证明}：考虑方向导数。对任意设计 \(\eta\)，令 \(\xi_\alpha = (1-\alpha)\xi^* + \alpha\eta\)。D-最优性准则 \(\Phi(\xi) = \log \det M(\xi)\) 在 \(\xi^*\) 处的 Frechet 导数为

\[\left.\frac{d}{d\alpha}\Phi(\xi_\alpha)\right|_{\alpha=0} = \operatorname{tr}\left(M(\xi^*)^{-1} (M(\eta) - M(\xi^*))\right) = \int_{\mathcal{X}} d(x, \xi^*) d\eta(x) - p.\]

（推导利用了 \(\frac{d}{d\alpha}\log \det M = \operatorname{tr}(M^{-1} \frac{dM}{d\alpha})\) 和 \(\frac{d}{d\alpha}M(\xi_\alpha) = M(\eta) - M(\xi^*)\)。）

由凸优化的一阶最优性条件，\(\xi^*\) 极大化 \(\Phi\) 当且仅当对任意 \(\eta\)，\(\left.\frac{d}{d\alpha}\Phi(\xi_\alpha)\right|_{\alpha=0} \leq 0\)，即

\[\int_{\mathcal{X}} d(x, \xi^*) d\eta(x) \leq p \quad \forall \eta.\]

取 \(\eta = \delta_x\)（点质量），得 \(d(x, \xi^*) \leq p\) 对任意 \(x\) 成立。反之，若该条件成立，则对任意 \(\eta\) 积分得导数为非正。

现证等号在支撑点上成立。设 \(\mathcal{S} = \operatorname{supp}(\xi^*)\)。若存在 \(x_0 \in \mathcal{S}\) 使 \(d(x_0, \xi^*) < p\)，则由连续性，可沿某方向移动质量增加准则值，矛盾。因此 \(\max_x d(x, \xi^*) = p\) 且在支撑点上达到。 \(\blacksquare\)

\textbf{定理 x+3.9}（D-最优设计的广义方差最小化）：在D-最优设计 \(\xi^*\) 下，任意线性组合 \(c^T \hat{\beta}\) 的方差上界由等价性定理控制：

\[\operatorname{Var}(c^T \hat{\beta}) = \sigma^2 c^T M(\xi^*)^{-1} c \leq \frac{\sigma^2}{N} \cdot p \cdot \max_{x \in \mathcal{X}} (c^T f(x))^2.\]

\textbf{证明}：由Cauchy-Schwarz不等式，

\[(c^T M^{-1} c) \cdot (f(x)^T M f(x)) \geq (c^T f(x))^2.\]

因 \(M(\xi^*) = N \cdot M_1(\xi^*)\)（其中 \(M_1\) 为单位权重信息矩阵），\(f(x)^T M^{-1} f(x) = \frac{1}{N} d(x, \xi^*)\)。由等价性定理，\(d(x, \xi^*) \leq p\)，故

\[\operatorname{Var}(c^T \hat{\beta}) = \frac{\sigma^2}{N} c^T M_1^{-1} c \leq \frac{\sigma^2}{N} \cdot \frac{p \cdot \max_x (c^T f(x))^2}{\min_x (f(x)^T M_1 f(x))}.\]

取紧化即得上界。 \(\blacksquare\)

\textbf{定理 x+3.10}（Fedorov-Wynn迭代算法）：D-最优设计可通过以下迭代构造：从非退化初始设计 \(\xi_0\) 出发，

\[x_{m+1} = \arg\max_{x \in \mathcal{X}} d(x, \xi_m), \quad \xi_{m+1} = (1 - \alpha_m)\xi_m + \alpha_m \delta_{x_{m+1}},\]

其中 \(\alpha_m = \frac{d(x_{m+1}, \xi_m) - p}{p(d(x_{m+1}, \xi_m) - 1)}\)（或直接取 \(\alpha_m = 1/(m+1)\)）。该序列收敛到D-最优设计。

\textbf{证明}：设 \(\Phi_m = \log \det M(\xi_m)\)。由方向导数

\[\Phi_{m+1} - \Phi_m \geq \alpha_m\left(\int d(x, \xi_m) d\delta_{x_{m+1}}(x) - p\right) + O(\alpha_m^2).\]

因 \(d(x_{m+1}, \xi_m) = \max_x d(x, \xi_m) \geq p\)（严格大于，除非 \(\xi_m\) 已最优），选择 \(\alpha_m\) 保证单调递增。由信息矩阵的紧性和 \(\Phi\) 的凹性，序列收敛。极限点满足 \(\max_x d(x, \xi^*) = p\)，由等价性定理即为D-最优设计。 \(\blacksquare\)

\subsection{x+3.4 区组化的最优设计理论与补充}\label{x3.4-ux533aux7ec4ux5316ux7684ux6700ux4f18ux8bbeux8ba1ux7406ux8bbaux4e0eux8865ux5145}

\textbf{定理 x+3.11}（区组设计的最优性准则）：设有 \(b\) 个区组，各区组容纳 \(k\) 个实验单元。在区组效应为固定效应的模型 \(Y_{ij} = \mu + \tau_i + \beta_j + \varepsilon_{ij}\) 下，处理对比的方差与设计矩阵 \(X\) 的信息矩阵 \(C = X^T(I - P_B)X\)（其中 \(P_B\) 为区组效应的投影矩阵）相关。设计 \(\xi\) 为D-最优的充要条件是最大化 \(\det(C)\)。

\textbf{证明}：区组模型的最小二乘估计中，\(\hat{\tau}\) 的协方差矩阵为 \(\sigma^2 C^{-1}\)，其中 \(C = N^T N - \frac{1}{k}N^T J_b N\)，\(N\) 为处理-区组关联矩阵。D-最优设计最大化 \(\det(C)\)，等价于最小化 \(\det(C^{-1})\)。由Kiefer-Wolfowitz定理在约束设计空间上的推广，存在离散设计（对应对称BIBD，当参数允许时）达到最优。平衡性（每处理出现次数相等、每对处理同时出现次数相等）往往导出最优设计。 \(\blacksquare\)

\textbf{定理 x+3.12}（A-最优性与E-最优性的关系）：设 \(\xi^*_D, \xi^*_A, \xi^*_E\) 分别为D-最优、A-最优（极小化 \(\operatorname{tr}(M^{-1})\)）和E-最优（极大化 \(\lambda_{\min}(M)\)）设计。对任意设计 \(\xi\)，成立

\[\frac{1}{p}\operatorname{tr}(M(\xi)^{-1}) \geq [\det M(\xi)]^{-1/p} \geq [\lambda_{\min}(M(\xi))]^{-1}.\]

因此D-最优性介于A-最优性和E-最优性之间。

\textbf{证明}：设 \(\lambda_1 \leq \lambda_2 \leq \cdots \leq \lambda_p\) 为 \(M(\xi)\) 的特征值。则 \(\det M = \prod_{i=1}^{p} \lambda_i\)，\(\operatorname{tr}(M^{-1}) = \sum_{i=1}^{p} \lambda_i^{-1}\)，\(\lambda_{\min} = \lambda_1\)。由算术-几何平均不等式，

\[\frac{1}{p}\sum_{i=1}^{p} \lambda_i^{-1} \geq \left(\prod_{i=1}^{p} \lambda_i^{-1}\right)^{1/p} = (\det M)^{-1/p}.\]

又 \(\prod \lambda_i^{-1} \geq \lambda_1^{-p}\)，故 \([\det M]^{-1/p} \geq \lambda_1^{-1}\)。这给出了三种最优性准则间的定量关系，A-最优最保守（极小化平均方差），E-最优最激进（极小化最大方差方向），D-最优居中。 \(\blacksquare\)

\textbf{定理 x+3.13}（精确设计与近似设计的Gap------效率下界）：若连续D-最优设计 \(\xi^*\) 的支撑点个数为 \(m\)，则存在精确 \(N\)-试验设计 \(\xi_N\)（\(N \geq m\)）使得其D-效率满足

\[\operatorname{Eff}_D(\xi_N) = \left(\frac{\det M(\xi_N)}{\det M(\xi^*)}\right)^{1/p} \geq 1 - \frac{p}{N}.\]

\textbf{证明}：由Caratheodory定理，\(\xi^*\) 可表示为至多 \(p(p+1)/2 + 1\) 个支撑点的凸组合。将 \(N\) 次试验按 \(\xi^*\) 的权重成比例分配（四舍五入取整），所得精确设计 \(\xi_N\) 满足 \(|n_i/N - w_i^*| \leq 1/N\)。利用信息矩阵的Lipschitz连续性：

\[\|M(\xi_N) - M(\xi^*)\|_F \leq \sum_i |n_i/N - w_i^*| \cdot \|f(x_i)f(x_i)^T\|_F \leq \frac{K}{N}.\]

由Weyl扰动定理，特征值的相对变化至多为 \(O(1/N)\)，故 \(\det M(\xi_N) / \det M(\xi^*) \geq (1 - C/N)^p\)，即效率下界成立。当 \(N\) 足够大时，精确设计可任意逼近连续最优设计。 \(\blacksquare\)

\textbf{定理 x+3.14}（区组的随机效应与广义线性混合模型设计）：当区组效应 \(\beta_j\) 为随机效应（\(\beta_j \stackrel{\text{iid}}{\sim} N(0, \sigma_b^2)\)）时，模型变为线性混合模型 \(Y_{ij} = \mu + \tau_i + \beta_j + \varepsilon_{ij}\)。观测的边际协方差矩阵为 \(\Sigma = \sigma^2 I + \sigma_b^2 Z Z^T\)（\(Z\) 为区组分配矩阵），处理效应估计 \(\hat{\tau}\) 的协方差为 \(\sigma^2 (X^T \Sigma^{-1} X)^{-1}\)。最优设计需同时考虑固定效应 \(\tau_i\) 和方差分量 \(\sigma_b^2\) 的估计精度。

\textbf{证明}：在混合模型中，广义最小二乘估计 \(\hat{\tau} = (X^T \hat{\Sigma}^{-1} X)^{-1} X^T \hat{\Sigma}^{-1} Y\)。由Henderson方程和限制极大似然（REML）估计，方差参数 \(\sigma^2, \sigma_b^2\) 的估计精度依赖于设计的二阶信息矩阵。利用分块矩阵求逆公式 \(\Sigma^{-1} = \frac{1}{\sigma^2}(I - \frac{\sigma_b^2}{\sigma^2 + k\sigma_b^2} Z Z^T)\)（当区组大小为常数 \(k\) 时），可得处理对比方差的显示表达。平衡设计（每处理在各区组中均匀分布）在广义线性混合模型中也往往是渐近最优的。 \(\blacksquare\)

\newpage

\chapter{11-数值与计算/01-数值分析/01-误差与线性数值.md}\label{ux6570ux503cux4e0eux8ba1ux7b9701-ux6570ux503cux5206ux679001-ux8befux5deeux4e0eux7ebfux6027ux6570ux503c.md}

\section{第384章：浮点运算与误差分析}\label{ux7b2c384ux7ae0ux6d6eux70b9ux8fd0ux7b97ux4e0eux8befux5deeux5206ux6790}

计算机中的实数运算是近似的，理解浮点运算的误差模型是数值分析的基础。本章介绍 IEEE 754 浮点标准、舍入误差和数值稳定性。

\subsection{114.1 浮点数系统}\label{ux6d6eux70b9ux6570ux7cfbux7edf}

\textbf{定义 114.1}（IEEE 754 双精度浮点数）：\textbf{双精度浮点数}（binary64）用 64 位表示：1 位符号 \(s\)、11 位指数 \(e\)（偏移 \(1023\)）、52 位尾数 \(m\)（隐含首位的 \(1\)）。数值为

\[x = (-1)^s \times (1.m) \times 2^{e - 1023}\]

\textbf{定义 114.2}（机器精度）：\textbf{机器精度}（machine epsilon）\(\varepsilon_{\text{mach}}\) 是 \(1\) 与大于 \(1\) 的最小浮点数之间的差。双精度下 \(\varepsilon_{\text{mach}} = 2^{-52} \approx 2.22 \times 10^{-16}\)。

\textbf{定义 114.3}（舍入）：浮点运算的 IEEE 标准舍入模式（默认：舍入到最近偶数）。记 \(\operatorname{fl}(x)\) 为 \(x\) 的浮点表示，则

\[\operatorname{fl}(x) = x(1 + \delta), \quad |\delta| \leq \varepsilon_{\text{mach}}\]

\textbf{定义 114.4}（浮点运算模型）：对基本运算 \(\circ \in \{+, -, \times, \div\}\)，

\[\operatorname{fl}(x \circ y) = (x \circ y)(1 + \delta), \quad |\delta| \leq \varepsilon_{\text{mach}}\]

\subsection{114.2 误差分析}\label{ux8befux5deeux5206ux6790}

\textbf{定义 114.5}（绝对误差与相对误差）：设 \(\tilde{x}\) 是 \(x\) 的近似值。 - \textbf{绝对误差}：\(|\tilde{x} - x|\) - \textbf{相对误差}：\(|\tilde{x} - x| / |x|\)（\(x \neq 0\)）

\textbf{定义 114.6}（前向误差与后向误差）： - \textbf{前向误差}：输出 \(\tilde{y}\) 与真值 \(y = f(x)\) 的差 \(|\tilde{y} - y|\) - \textbf{后向误差}：寻找 \(\Delta x\) 使得 \(\tilde{y} = f(x + \Delta x)\)，后向误差为 \(|\Delta x|\)

一个算法称为\textbf{后向稳定}的，如果后向误差为 \(O(\varepsilon_{\text{mach}})\)。

\textbf{定义 114.7}（条件数）：问题 \(f\) 在 \(x\) 处的\textbf{（相对）条件数}为

\[\kappa_f(x) = \lim_{\varepsilon \to 0} \sup_{\|\Delta x\| \leq \varepsilon} \frac{\|f(x + \Delta x) - f(x)\| / \|f(x)\|}{\|\Delta x\| / \|x\|}\]

若 \(\kappa_f(x)\) 很大，则问题是\textbf{病态}的（ill-conditioned）。前向误差 \(\lesssim\) 条件数 \(\times\) 后向误差。

\textbf{定义 114.8}（数值稳定性）：算法是\textbf{数值稳定}的，如果输出对输入微小扰动不敏感。等价地，稳定的算法产生的结果接近精确问题的近似解（后向稳定）。

\hrulefill

\hrulefill

\hrulefill

\hrulefill

\hrulefill

\hrulefill

\section{第385章：线性方程组数值解}\label{ux7b2c385ux7ae0ux7ebfux6027ux65b9ux7a0bux7ec4ux6570ux503cux89e3}

求解线性方程组 \(A\mathbf{x} = \mathbf{b}\) 是数值计算最基本的任务。本章讨论直接法（LU 分解、QR 分解）和迭代法。

\subsection{115.1 高斯消元与 LU 分解}\label{ux9ad8ux65afux6d88ux5143ux4e0e-lu-ux5206ux89e3}

\textbf{定理 115.1}（LU 分解）：设 \(A \in \mathbb{R}^{n \times n}\) 的所有顺序主子式非零。则存在唯一的单位下三角矩阵 \(L\)（\(L_{ii} = 1\)）和上三角矩阵 \(U\)，使得 \(A = LU\)。

\emph{证明}：归纳构造。设 \(A = \begin{pmatrix} a_{11} & \mathbf{w}^\top \\ \mathbf{v} & \tilde{A} \end{pmatrix}\)。令 \(L_1 = \begin{pmatrix} 1 & 0 \\ \mathbf{v}/a_{11} & I \end{pmatrix}\)，则 \(L_1^{-1} A = \begin{pmatrix} a_{11} & \mathbf{w}^\top \\ 0 & \tilde{A}_1 \end{pmatrix}\)。对 \(\tilde{A}_1\) 递归。∎

\textbf{定义 115.1}（部分选主元，PA = LU）：为提高数值稳定性，在消元时选择列中绝对值最大的元素作为主元（行交换）。\(PA = LU\)，其中 \(P\) 是置换矩阵。

\textbf{算法 115.1}（求解 \(A\mathbf{x} = \mathbf{b}\) 的步骤）： 1. 计算 \(PA = LU\) 分解（\(O(n^3)\) 次运算） 2. 解 \(L\mathbf{y} = P\mathbf{b}\)（前代，\(O(n^2)\)） 3. 解 \(U\mathbf{x} = \mathbf{y}\)（回代，\(O(n^2)\)）

\textbf{定理 115.2}（Cholesky 分解）：若 \(A \in \mathbb{R}^{n \times n}\) 对称正定，则存在唯一的下三角矩阵 \(G\)（\(G_{ii} > 0\)）使得 \(A = GG^\top\)。计算量为 \(n^3/3 + O(n^2)\)（约为 LU 的一半）。

\textbf{证明}：存在性：对 \(n\) 归纳。\(n=1\)：\(A=a_{11}>0\)，\(G=[\sqrt{a_{11}}]\)。设对 \(n-1\) 成立。分块 \(A=\begin{pmatrix}A_{11}&\mathbf{a}\\\mathbf{a}^T&a_{nn}\end{pmatrix}\)。令 \(G_{n-1}\) 满足 \(A_{11}=G_{n-1}G_{n-1}^T\)。求解 \(G_{n-1}\mathbf{g}=\mathbf{a}\)（前代），\(g_{nn}=\sqrt{a_{nn}-\|\mathbf{g}\|^2}\)。由正定性 \(a_{nn}-\mathbf{g}^T\mathbf{g}=a_{nn}-\mathbf{a}^T A_{11}^{-1}\mathbf{a}>0\)（Schur 补正值）。\(G=\begin{pmatrix}G_{n-1}&0\\\mathbf{g}^T&g_{nn}\end{pmatrix}\)。唯一性：若 \(G_1G_1^T=G_2G_2^T\)，由下三角性 \(G_2^{-1}G_1=((G_2^{-1}G_1)^{-1})^T\) 且均为下三角，故为对角矩阵。正对角元唯一确定。复杂度：\(\sum_{k=1}^n (k^2+O(k))=n^3/3+O(n^2)\)。\(\blacksquare\)

\subsection{115.2 QR 分解与最小二乘}\label{qr-ux5206ux89e3ux4e0eux6700ux5c0fux4e8cux4e58}

\textbf{定理 115.3}（QR 分解）：设 \(A \in \mathbb{R}^{m \times n}\)（\(m \geq n\)）满列秩。则存在正交矩阵 \(Q \in \mathbb{R}^{m \times m}\)（\(Q^\top Q = I\)）和上三角矩阵 \(R \in \mathbb{R}^{n \times n}\)，使得

\[A = Q \begin{pmatrix} R \\ 0 \end{pmatrix} = \tilde{Q} R\]

其中 \(\tilde{Q} \in \mathbb{R}^{m \times n}\) 是 \(Q\) 的前 \(n\) 列（\(\tilde{Q}^\top \tilde{Q} = I_n\)）。

\emph{证明}：使用 Householder 反射或 Givens 旋转构造。第 \(k\) 步用一个 Householder 反射将第 \(k\) 列对角线以下的元素消为零。∎

\textbf{定理 115.4}（最小二乘问题）：超定系统 \(A\mathbf{x} \approx \mathbf{b}\) 的最小二乘解为

\[\mathbf{x}^* = \arg\min_{\mathbf{x}} \|A\mathbf{x} - \mathbf{b}\|_2 = (A^\top A)^{-1} A^\top \mathbf{b}\]

使用 QR 分解：\(\mathbf{x}^* = R^{-1} \tilde{Q}^\top \mathbf{b}\)（数值更稳定）。

\textbf{证明}：法方程推导：\(\|A\mathbf{x}-\mathbf{b}\|_2^2=(A\mathbf{x}-\mathbf{b})^T(A\mathbf{x}-\mathbf{b})\)，对 \(\mathbf{x}\) 求梯度并置零：\(2A^T(A\mathbf{x}-\mathbf{b})=0\implies A^T A\mathbf{x}=A^T\mathbf{b}\)。\(A\) 满列秩时 \(A^T A\) 正定，\(\mathbf{x}^*\) 唯一。QR 法：\(A=\tilde{Q}R\)，则 \(\|A\mathbf{x}-\mathbf{b}\|_2=\|\tilde{Q}R\mathbf{x}-\mathbf{b}\|_2=\|\begin{pmatrix}R\mathbf{x}\\0\end{pmatrix}-\tilde{Q}^T\mathbf{b}\|_2\)。令 \(\tilde{Q}^T\mathbf{b}=\begin{pmatrix}\mathbf{c}\\\mathbf{d}\end{pmatrix}\)，极小化等价于 \(\min\|R\mathbf{x}-\mathbf{c}\|_2^2+\|\mathbf{d}\|_2^2\)，取 \(\mathbf{x}=R^{-1}\mathbf{c}\) 得最小值 \(\|\mathbf{d}\|_2\)。QR 法避免了 \(A^T A\) 的条件数平方放大。\(\blacksquare\)

\subsection{115.3 迭代法}\label{ux8fedux4ee3ux6cd5}

\textbf{定义 115.2}（迭代法的一般形式）：将 \(A\) 分裂为 \(A = M - N\)（\(M\) 可逆）。迭代格式为

\[M\mathbf{x}^{(k+1)} = N\mathbf{x}^{(k)} + \mathbf{b}\]

即 \(\mathbf{x}^{(k+1)} = M^{-1}N\mathbf{x}^{(k)} + M^{-1}\mathbf{b}\)。

\textbf{定义 115.3}（经典迭代法）： - \textbf{Jacobi 迭代}：\(M = D\)（\(A\) 的对角部分），\(N = -(L + U)\) - \textbf{Gauss-Seidel 迭代}：\(M = D + L\)，\(N = -U\)

\textbf{定理 115.5}（收敛性）：迭代法收敛（对任意初始值）当且仅当迭代矩阵 \(M^{-1}N\) 的谱半径 \(\rho(M^{-1}N) < 1\)。对对称正定矩阵，Gauss-Seidel 迭代收敛。

\textbf{证明}：迭代 \(\mathbf{x}^{(k+1)}=G\mathbf{x}^{(k)}+\mathbf{c}\)（\(G=M^{-1}N\)）。误差 \(\mathbf{e}^{(k)}=\mathbf{x}^{(k)}-\mathbf{x}^*\) 满足 \(\mathbf{e}^{(k+1)}=G\mathbf{e}^{(k)}=G^k\mathbf{e}^{(0)}\)。\(\mathbf{e}^{(k)}\to\mathbf{0}\)（对任意 \(\mathbf{e}^{(0)}\)）\(\iff G^k\to 0\iff\rho(G)<1\)（Jordan 标准型）。对对称正定 \(A\)，Gauss-Seidel 的迭代矩阵 \(G_{GS}=-(D+L)^{-1}U\) 满足 \(\rho(G_{GS})<1\)：\(A\) 对称正定 \(\implies\) 所有顺序主子式 \(>0\) \(\implies\) Gauss-Seidel 收敛（Ostrowski-Reich 定理：若 \(A\) 对称正定则 Gauss-Seidel 迭代对所有 \(\omega\in(0,2)\) 收敛）。\(\blacksquare\)

\textbf{定理 115.6}（共轭梯度法，CG）：对对称正定 \(A\)，共轭梯度法在 \(A\)-范数 \(\|\mathbf{x}\|_A = \sqrt{\mathbf{x}^\top A\mathbf{x}}\) 下最小化误差。理论上 CG 在 \(n\) 步内精确求解（在精确算术下），但实际中通常作为迭代法使用。收敛速率：

\[\|\mathbf{x}^{(k)} - \mathbf{x}^*\|_A \leq 2\left(\frac{\sqrt{\kappa} - 1}{\sqrt{\kappa} + 1}\right)^k \|\mathbf{x}^{(0)} - \mathbf{x}^*\|_A\]

其中 \(\kappa = \lambda_{\max}(A) / \lambda_{\min}(A)\) 是 \(A\) 的条件数。

\textbf{证明}：CG 在第 \(k\) 步产生 \(\mathbf{x}^{(k)}\in\mathbf{x}^{(0)}+\mathcal{K}_k(A,\mathbf{r}^{(0)})\) 极小化 \(\|\mathbf{x}-\mathbf{x}^*\|_A\)。因 \(\mathcal{K}_n\subseteq\mathbb{R}^n\)，\(n\) 步内 \(\mathcal{K}_n\) 线性相关产生全空间，故精确解。收敛界推导：\(\|\mathbf{x}^{(k)}-\mathbf{x}^*\|_A=\min_{p\in\mathcal{P}_k,p(0)=1}\|p(A)(\mathbf{x}^{(0)}-\mathbf{x}^*)\|_A\)。取 \(p\) 为缩放平移的 Chebyshev 多项式：\(p_k(x)=T_k(\frac{\lambda_{\max}+\lambda_{\min}-2x}{\lambda_{\max}-\lambda_{\min}})/T_k(\frac{\lambda_{\max}+\lambda_{\min}}{\lambda_{\max}-\lambda_{\min}})\)。在特征空间展开得 \(\|p_k(A)\mathbf{e}^{(0)}\|_A\le\max_{\lambda\in[\lambda_{\min},\lambda_{\max}]}|p_k(\lambda)|\cdot\|\mathbf{e}^{(0)}\|_A\)。而 \(\max|p_k(\lambda)|=1/T_k(\frac{\kappa+1}{\kappa-1})\le2(\frac{\sqrt{\kappa}-1}{\sqrt{\kappa}+1})^k\)（由 \(T_k(x)\ge\frac{1}{2}(x+\sqrt{x^2-1})^k\)）。\(\blacksquare\)

\hrulefill

\hrulefill

\hrulefill

\hrulefill

\hrulefill

\hrulefill

\section{第386章：特征值问题}\label{ux7b2c386ux7ae0ux7279ux5f81ux503cux95eeux9898}

特征值问题是数值线性代数的核心问题之一。本章讨论幂迭代、QR 算法和 Jacobi 方法。

\subsection{116.1 特征值的基本理论}\label{ux7279ux5f81ux503cux7684ux57faux672cux7406ux8bba}

\textbf{定义 116.1}（特征值问题）：求 \(\lambda \in \mathbb{C}\) 和非零向量 \(\mathbf{v} \in \mathbb{C}^n\) 使得 \(A\mathbf{v} = \lambda \mathbf{v}\)。

\textbf{定理 116.1}（Schur 分解）：对任意 \(A \in \mathbb{C}^{n \times n}\)，存在酉矩阵 \(Q\)（\(Q^*Q = I\)）和上三角矩阵 \(T\) 使得 \(A = QTQ^*\)。\(T\) 的对角元是 \(A\) 的特征值。

\textbf{证明}：对 \(n\) 归纳。\(n=1\) 平凡。设任意 \((n-1)\times(n-1)\) 矩阵可 Schur 分解。取 \(A\) 的特征值 \(\lambda_1\) 和单位特征向量 \(\mathbf{q}_1\)：\(A\mathbf{q}_1=\lambda_1\mathbf{q}_1\)。将 \(\mathbf{q}_1\) 扩充为酉矩阵 \(Q_1=[\mathbf{q}_1\;Q_2]\)。则 \(Q_1^*AQ_1=\begin{pmatrix}\lambda_1&\mathbf{w}^*\\\mathbf{0}&A_{22}\end{pmatrix}\)（左下角为 \(\mathbf{0}\) 因 \(Q_2^*A\mathbf{q}_1=\lambda_1 Q_2^*\mathbf{q}_1=\mathbf{0}\)）。由归纳假设，\(A_{22}=\tilde{Q}_2\tilde{T}_2\tilde{Q}_2^*\)。令 \(Q=Q_1\begin{pmatrix}1&0\\0&\tilde{Q}_2\end{pmatrix}\) 和 \(T=\begin{pmatrix}\lambda_1&\mathbf{w}^*\tilde{Q}_2\\0&\tilde{T}_2\end{pmatrix}\)，则 \(QTQ^*=A\)。\(\blacksquare\)

\textbf{推论 116.2}（实 Schur 分解）：若 \(A \in \mathbb{R}^{n \times n}\)，存在正交矩阵 \(Q\) 使得 \(Q^\top A Q\) 是拟上三角矩阵（对角块为 \(1 \times 1\) 或 \(2 \times 2\)）。

\textbf{定理 116.3}（谱定理，Hermitian 情形）：若 \(A\) 是 Hermitian 矩阵（\(A^* = A\)），则存在酉矩阵 \(Q\) 和对角矩阵 \(\Lambda\) 使得 \(A = Q\Lambda Q^*\)。特征值全为实数。

\textbf{证明}：由 Schur 分解 \(A=QTQ^*\)。因 \(A^*=A\)，\(QT^*Q^*=QTQ^*\)，故 \(T^*=T\)。\(T\) 同时为上三角和 Hermitian \(\implies\) \(T\) 为对角矩阵（因 \(T_{ij}=0\)（\(i>j\)）且 \(T_{ij}=\overline{T_{ji}}\)（\(i<j\)）\(\implies T_{ij}=0\)）。故 \(T=\Lambda\) 为对角矩阵，且 \(\Lambda=\Lambda^*\) 蕴含所有对角元为实数（等于 \(A\) 的特征值）。\(\blacksquare\)

\subsection{116.2 幂迭代及其变体}\label{ux5e42ux8fedux4ee3ux53caux5176ux53d8ux4f53}

\textbf{定义 116.2}（幂迭代，Power Iteration）：对可对角化矩阵 \(A\)，设 \(|\lambda_1| > |\lambda_2| \geq \cdots \geq |\lambda_n|\)（\(\lambda_1\) 是主特征值）。幂迭代为

\[\mathbf{v}^{(k+1)} = A\mathbf{v}^{(k)} / \|A\mathbf{v}^{(k)}\|\]

\(\mathbf{v}^{(k)}\) 收敛到 \(\lambda_1\) 对应的特征向量，收敛速率为 \(|\lambda_2 / \lambda_1|^k\)。Rayleigh 商 \(\rho(\mathbf{v}) = \mathbf{v}^\top A\mathbf{v} / \mathbf{v}^\top \mathbf{v}\) 收敛到 \(\lambda_1\)。

\textbf{定义 116.3}（反迭代）：对特征值 \(\mu\) 附近的特征向量，反迭代为 \(\mathbf{v}^{(k+1)} = (A - \mu I)^{-1} \mathbf{v}^{(k)}\)（解线性系统）。收敛到离 \(\mu\) 最近的特征值对应的特征向量，收敛极快。

\textbf{定义 116.4}（Rayleigh 商迭代）：反迭代 + 动态更新 \(\mu\)：\(\mu_k = \rho(\mathbf{v}^{(k)})\)。三次收敛（在对称矩阵情形下）。

\subsection{116.3 QR 算法}\label{qr-ux7b97ux6cd5}

\textbf{算法 116.1}（QR 算法，基本形式）：\(A_0 = A\)。对 \(k = 0, 1, 2, \ldots\)： 1. 计算 \(A_k = Q_k R_k\)（QR 分解） 2. \(A_{k+1} = R_k Q_k\)

\(A_k\) 收敛到拟上三角矩阵（实 Schur 形式），对角块给出特征值。

\textbf{定理 116.4}（QR 算法的收敛性）：若 \(A\) 的特征值满足 \(|\lambda_1| > |\lambda_2| > \cdots > |\lambda_n| > 0\)，则 \(A_k\) 的对角线以下元素以 \(|\lambda_{i+1} / \lambda_i|^k\) 的速率收敛到 \(0\)。

\textbf{证明}：QR 迭代 \(A_{k+1}=R_kQ_k\)（\(A_k=Q_kR_k\)）等价于 \(A_{k+1}=Q_k^*A_kQ_k\)，故 \(A_k\) 保持相似。由 \(A_k=Q^{(k)}R^{(k)}\)（\(Q^{(k)}=Q_0Q_1\cdots Q_{k-1}\)，\(R^{(k)}=R_{k-1}\cdots R_0\)），且 \(A^k=Q^{(k)}R^{(k)}\) 是 \(A^k\) 的 QR 分解。设 \(A=X\Lambda X^{-1}\)（\(\Lambda=\operatorname{diag}(\lambda_i)\)）。\(A^k=X\Lambda^k X^{-1}=Q^{(k)}R^{(k)}\)。令 \(X^{-1}=LU\)（带列交换的 LU），则 \(A^k=X\Lambda^k LU=Q^{(k)}R^{(k)}\)。将 \(\Lambda^k L\Lambda^{-k}\) 展开为 \(I+O(|\lambda_{i+1}/\lambda_i|^k)\)（严格下三角部分），由此得 \(Q^{(k)}\to X\) 意义下的收敛，且 \(A_k\) 的对角元 \(\to\lambda_i\)，次对角线以 \(|\lambda_{i+1}/\lambda_i|^k\) 速率收敛到零。\(\blacksquare\)

\textbf{算法 116.2}（实用 QR 算法）： 1. 先约化为上 Hessenberg 形式（\(O(n^3)\)，但后续每步仅 \(O(n^2)\)） 2. 引入位移（shift）加速收敛：\(A_k - \mu I = Q_k R_k\)，\(A_{k+1} = R_k Q_k + \mu I\) 3. 使用 Wilkinson 位移（右下角 \(2 \times 2\) 子矩阵的特征值中较近者） 4. 实用中通常 2-3 次迭代就收敛一个特征值

\subsection{116.4 Jacobi 方法（对称矩阵）}\label{jacobi-ux65b9ux6cd5ux5bf9ux79f0ux77e9ux9635}

\textbf{算法 116.3}（Jacobi 方法）：对对称矩阵 \(A\)，反复应用 Givens 旋转 \(J(p, q, \theta)\) 消去非对角元 \(a_{pq}\)（旋转角 \(\theta\) 满足 \(\tan 2\theta = 2a_{pq}/(a_{pp} - a_{qq})\)）。\(A \leftarrow J^\top A J\)。

\textbf{定理 116.5}（Jacobi 方法的收敛性）：Jacobi 方法二次收敛到对角矩阵（在每步消去最大非对角元的策略下）。非对角元素的 Frobenius 范数单调递减。

\textbf{证明}：每步选择最大非对角元 \(|a_{pq}|\)（\(p<q\)），应用 Givens 旋转 \(J(p,q,\theta)\) 使 \(a_{pq}=0\)。旋转后非对角 Frobenius 范数的变化：\(\operatorname{off}(A')^2=\operatorname{off}(A)^2-2|a_{pq}|^2\)（因 \(A'=J^T A J\) 仅在 \((p,q)\) 和 \((q,p)\) 位置和 \(p,q\) 行/列的对角元有变化）。故 \(\operatorname{off}(A_k)\) 单调递减。二次收敛性：当 \(A_k\) 足够接近对角矩阵时，旋转消去的元素 \(a_{pq}\) 被随后旋转恢复的量 \(\sim a_{pq}^2\)，导致 \(\operatorname{off}(A_{k+n})\le C\cdot\operatorname{off}(A_k)^2\)------二次收敛。具体地，若 \(\operatorname{off}(A_m)\le\delta\)，则下一步后 \(\operatorname{off}(A_{m+1})\le 2\delta^2/d\)（\(d\) 为特征值最小间隔）。\(\blacksquare\)

\hrulefill

\hrulefill

\hrulefill

\hrulefill

\hrulefill

\hrulefill

\hrulefill

\hrulefill

\section{第387章：Krylov子空间方法（GMRES与CG）}\label{ux7b2c387ux7ae0krylovux5b50ux7a7aux95f4ux65b9ux6cd5gmresux4e0ecg}

Krylov 子空间方法是求解大规模稀疏线性方程组 \(A\mathbf{x} = \mathbf{b}\) 的最重要算法族，其核心思想是在低维 Krylov 子空间 \(\mathcal{K}_m(A, \mathbf{r}_0) = \operatorname{span}\{\mathbf{r}_0, A\mathbf{r}_0, A^2\mathbf{r}_0, \ldots, A^{m-1}\mathbf{r}_0\}\) 中寻找近似解。Krylov 方法的威力来自两个事实：首先，\(\mathcal{K}_m\) 的维数远小于问题规模 \(n\)，使得每步迭代的运算和存储开销可控（通常 \(O(n)\)-稀疏矩阵向量乘）；其次，Krylov 子空间自然地编码了谱信息------\(\mathcal{K}_m\) 是 \(A\) 在 \(\mathbf{r}_0\) 上的极小多项式子空间，因此 \(m\) 步 Krylov 方法等价于在次数不超过 \(m-1\) 的矩阵多项式中选取最优的 \(p(A) \mathbf{r}_0\) 来逼近 \(\mathbf{b} - A\mathbf{x}_0\)。共轭梯度法（Conjugate Gradient, CG）和广义极小残差法（Generalized Minimal Residual, GMRES）是 Krylov 家族中最具代表性的两个方法：CG 适用于对称正定矩阵，在 \(A\)-范数下实现最优逼近；GMRES 通过 Arnoldi 过程和最小二乘求解，适用于一般非对称矩阵，在 Euclidean 范数下极小化残差。CG 的收敛速率由条件数 \(\kappa_2(A)\) 的平方根决定（\(\sim (\sqrt{\kappa} - 1)/(\sqrt{\kappa} + 1)\)），而 GMRES 的收敛分析涉及矩阵的数值域（field of values）和伪谱。

\subsection{59.x Arnoldi过程与Krylov子空间的正交基}\label{x-arnoldiux8fc7ux7a0bux4e0ekrylovux5b50ux7a7aux95f4ux7684ux6b63ux4ea4ux57fa}

\textbf{定义 59.1}（Arnoldi过程）：给定矩阵 \(A \in \mathbb{C}^{n \times n}\) 和初始向量 \(\mathbf{v}_1\)（\(\|\mathbf{v}_1\|_2 = 1\)），Arnoldi 过程通过修正的 Gram-Schmidt 方法构造 Krylov 子空间 \(\mathcal{K}_m = \operatorname{span}\{\mathbf{v}_1, A\mathbf{v}_1, \ldots, A^{m-1}\mathbf{v}_1\}\) 的一组标准正交基 \(\{\mathbf{v}_1, \ldots, \mathbf{v}_m\}\)。第 \(k\) 步：\(h_{i,k} = \mathbf{v}_i^* A \mathbf{v}_k\)（\(i = 1, \ldots, k\)），\(\mathbf{w} = A\mathbf{v}_k - \sum_{i=1}^k h_{i,k} \mathbf{v}_i\)，\(h_{k+1,k} = \|\mathbf{w}\|_2\)，\(\mathbf{v}_{k+1} = \mathbf{w} / h_{k+1,k}\)。这 \(m\) 步等价于部分正交分解

\[A V_m = V_m H_m + h_{m+1,m} \mathbf{v}_{m+1} \mathbf{e}_m^T = V_{m+1} \bar{H}_m\]

其中 \(V_m = [\mathbf{v}_1, \ldots, \mathbf{v}_m] \in \mathbb{C}^{n \times m}\) 具有正交列（\(V_m^* V_m = I_m\)），\(H_m \in \mathbb{C}^{m \times m}\) 为上海森堡（Hessenberg）矩阵（\(h_{ij} = 0\) 当 \(i > j+1\)），\(\bar{H}_m \in \mathbb{C}^{(m+1) \times m}\) 为 \(H_m\) 下加一行 \(h_{m+1,m} \mathbf{e}_m^T\)。

\textbf{定理 59.1}（Arnoldi过程的Krylov子空间刻画）：\(\mathcal{K}_m = \operatorname{colspan}(V_m)\)，且 \(\operatorname{span}\{\mathbf{v}_1, \ldots, \mathbf{v}_m\}\) 是 \(\mathcal{K}_m\) 的一组标准正交基。\(V_m^* A V_m = H_m\) 是 \(A\) 在 \(\mathcal{K}_m\) 的正交投影上的矩阵表示（Rayleigh-Ritz 投影）。对于对称 \(A\)，Arnoldi 过程退化为 \textbf{Lanczos 过程}：\(H_m\) 为对称三对角矩阵，三项递推关系 \(\beta_k \mathbf{v}_{k+1} = A \mathbf{v}_k - \alpha_k \mathbf{v}_k - \beta_{k-1} \mathbf{v}_{k-1}\) 成立。

\textbf{证明}：由归纳法，\(\mathbf{v}_k \in \mathcal{K}_k\)（归基：\(\mathbf{v}_1 \in \mathcal{K}_1\)；归纳：\(\mathbf{v}_{k+1} = (A\mathbf{v}_k - \sum h_{i,k} \mathbf{v}_i) / h_{k+1,k}\) 是 \(\mathcal{K}_{k+1}\) 中元素的正交化）。正交性由 Gram-Schmidt 过程保证。\(\mathbf{v}_1, \ldots, \mathbf{v}_m\) 的正交性等价于 \(V_m^* V_m = I_m\)。\(V_m^* A V_m = H_m\) 来自左乘 \(V_m^*\) 于 Arnoldi 关系。对称情形 \(A^* = A \implies H_m^* = V_m^* A V_m = H_m\)，三对角性由上海森堡矩阵的对称性强制：\(h_{ij} = 0\)（\(i > j+1\) 且 \(j > i+1\)），故三对角。\(\blacksquare\)

\subsection{59.y GMRES方法的收敛性分析}\label{y-gmresux65b9ux6cd5ux7684ux6536ux655bux6027ux5206ux6790}

\textbf{定理 59.2}（GMRES的最优性，Saad-Schultz 1986）：GMRES 在第 \(m\) 步的近似解 \(\mathbf{x}_m = \mathbf{x}_0 + V_m \mathbf{y}_m\)（\(\mathbf{y}_m\) 极小化 \(\|\beta \mathbf{e}_1 - \bar{H}_m \mathbf{y}\|_2\)，其中 \(\beta = \|\mathbf{r}_0\|_2\)）等价于全域最优：

\[\|\mathbf{b} - A\mathbf{x}_m\|_2 = \min_{\mathbf{x} \in \mathbf{x}_0 + \mathcal{K}_m} \|\mathbf{b} - A\mathbf{x}\|_2 = \min_{p \in \mathcal{P}_m, p(0)=1} \|p(A) \mathbf{r}_0\|_2\]

其中 \(\mathcal{P}_m\) 为次数不超过 \(m\) 的多项式集合。因此 GMRES 残差依 \(\min_{p(0)=1, \deg p \leq m} \|p(A)\|_2\) 衰减。

\textbf{证明}：\(\mathbf{x} \in \mathbf{x}_0 + \mathcal{K}_m\) 等价于 \(\mathbf{x} = \mathbf{x}_0 + V_m \mathbf{y}\)（\(\mathbf{y} \in \mathbb{C}^m\)）。残差 \(\mathbf{b} - A(\mathbf{x}_0 + V_m \mathbf{y}) = \mathbf{r}_0 - A V_m \mathbf{y} = \beta \mathbf{v}_1 - V_{m+1} \bar{H}_m \mathbf{y}\)（因 \(\mathbf{r}_0 = \beta V_m \mathbf{e}_1\) 和 Arnoldi 关系）。由 \(V_{m+1}\) 为标准正交列，\(\|\mathbf{r}_m\|_2 = \|\beta \mathbf{e}_1 - \bar{H}_m \mathbf{y}\|_2\)。极小化此量即 GMRES 的每步 \(O(m^2)\) 最小二乘子问题。等价于在 Krylov 子空间上全局极小化，故残差单调非增。多项式解释：\(\mathbf{r}_m = \mathbf{b} - A\mathbf{x}_m = p_m(A) \mathbf{r}_0\)，其中 \(p_m\) 是满足 \(p_m(0) = 1\) 的 \(m\) 次多项式。\(\blacksquare\)

\textbf{定理 59.3}（GMRES的收敛界 / 数值域界）：设 \(A\) 可对角化 \(A = X \Lambda X^{-1}\)。则

\[\frac{\|\mathbf{r}_m\|_2}{\|\mathbf{r}_0\|_2} \leq \kappa_2(X) \min_{p(0)=1, \deg p \leq m} \max_{\lambda \in \sigma(A)} |p(\lambda)|\]

若 \(A\) 的数值域 \(W(A) = \{ \mathbf{x}^* A \mathbf{x} / \mathbf{x}^* \mathbf{x} : \mathbf{x} \neq \mathbf{0} \}\) 不包含原点，则

\[\frac{\|\mathbf{r}_m\|_2}{\|\mathbf{r}_0\|_2} \leq \left(1 - \frac{\nu(A)}{\mu(A)}\right)^{m/2}\]

其中 \(\nu(A) = \min_{\lambda \in W(A)} |\lambda|\)，\(\mu(A) = \max_{\lambda \in W(A)} |\lambda|\)。

\textbf{证明}：（第一式）由定理 59.2 和矩阵多项式范数估计：\(\|p(A)\|_2 \leq \|X\|_2 \|X^{-1}\|_2 \max_{\lambda \in \sigma(A)} |p(\lambda)| = \kappa_2(X) \max_\lambda |p(\lambda)|\)。取 Chebyshev 多项式在包含谱的椭圆上的极值作为最优多项式近似，得到进一步的收敛速率估计。（第二式）利用 Crouzeix 的数值域界：\(\|p(A)\|_2 \leq 2 \max_{z \in W(A)} |p(z)|\)。若 \(0 \notin W(A)\)，则在 \(W(A)\) 的包含区间上取映射该区间到原点附近的多项式，速率如述。\(\blacksquare\)

\subsection{59.z 共轭梯度法的理论与超收敛}\label{z-ux5171ux8f6dux68afux5ea6ux6cd5ux7684ux7406ux8bbaux4e0eux8d85ux6536ux655b}

\textbf{定理 59.4}（CG的误差估计与谱分析）：对对称正定矩阵 \(A\)，CG 在第 \(m\) 步的 \(A\)-范数误差满足

\[\|\mathbf{x}_m - \mathbf{x}^*\|_A \leq 2 \left( \frac{\sqrt{\kappa_2(A)} - 1}{\sqrt{\kappa_2(A)} + 1} \right)^m \|\mathbf{x}_0 - \mathbf{x}^*\|_A\]

其中 \(\kappa_2(A) = \lambda_{\max} / \lambda_{\min}\) 为谱条件数。更一般地，若 \(A\) 的谱分为 \(k\) 个紧群（cluster），且群外仅有 \(\ell\) 个孤立特征值，则 CG 在 \(k + \ell\) 步内以实质上超线性速率收敛：\(\|\mathbf{r}_{k+\ell}\|_2 / \|\mathbf{r}_0\|_2 \leq \varepsilon\)（\(\varepsilon\) 由聚类紧度决定）。

\textbf{证明}：CG 误差的 \(A\)-范数极小化性质：\(\|\mathbf{x}_m - \mathbf{x}^*\|_A = \min_{p(0)=1, \deg p \leq m} \|p(A)(\mathbf{x}_0 - \mathbf{x}^*)\|_A\)。展开误差在 \(A\) 的特征基上：\(\mathbf{x}_0 - \mathbf{x}^* = \sum_{i=1}^n \alpha_i \mathbf{u}_i\)，\(A\mathbf{u}_i = \lambda_i \mathbf{u}_i\)（正交归一）。则 \(\|\mathbf{x}_m - \mathbf{x}^*\|_A^2 = \min_{p(0)=1} \sum_{i=1}^n |p(\lambda_i)|^2 \lambda_i \alpha_i^2 \leq \min_{p(0)=1} \max_{\lambda \in [\lambda_{\min}, \lambda_{\max}]} |p(\lambda)|^2 \cdot \|\mathbf{x}_0 - \mathbf{x}^*\|_A^2\)。在 \([\lambda_{\min}, \lambda_{\max}]\) 上求取 \(p(0)=1\) 的次数 \(m\) 多项式的最优界，由 Chebyshev 多项式给出：\(\min_{p(0)=1} \max_{\lambda \in [a,b]} |p(\lambda)| \leq 2 / (T_m((b+a)/(b-a)))\)，其中 \(T_m\) 为 \(m\) 次 Chebyshev 多项式。当 \(b/a = \kappa\) 时，\(T_m((\kappa+1)/(\kappa-1)) \sim \frac{1}{2} ((\sqrt{\kappa}+1)/(\sqrt{\kappa}-1))^m\)，故 \(\max |p(\lambda)| \leq 2 ((\sqrt{\kappa}-1)/(\sqrt{\kappa}+1))^m\)。

超收敛谱分析：将谱聚类外的 \(\ell\) 个特征值通过多项式 \(q(x) = \prod_{j=1}^\ell (1 - x/\lambda_j')\)（\(\lambda_j'\) 为孤立特征值）``杀死''：令 \(p(\lambda) = q(\lambda) \tilde{p}(\lambda)\)（\(\deg \tilde{p} \leq m - \ell\)）。\(q\) 在孤立特征值处为零，故误差在这些方向上被消除。余下的 \(\tilde{p}\) 在聚类区间上由 Chebyshev 多项式控制，速率由聚类的相对宽度决定。若聚类紧致（宽度很小），则收敛极快。\(\blacksquare\)

\hrulefill

\hrulefill

\hrulefill

\section{第388章：多重网格与区域分解方法}\label{ux7b2c388ux7ae0ux591aux91cdux7f51ux683cux4e0eux533aux57dfux5206ux89e3ux65b9ux6cd5}

多重网格方法（Multigrid Method, MG）是求解偏微分方程离散化后的大型稀疏线性系统的最优复杂度算法。其基本理念是：经典的迭代法（如 Jacobi、Gauss-Seidel）在消除高频误差分量时收敛迅速，但对低频（光滑）误差分量收敛极慢------这正是''粗网格''介入的动机。多重网格方法通过在粗细网格层级之间转移残差和校正量，将不同频率的误差分量分配给适当尺度的网格进行处理：细网格上的光滑误差在粗网格上表现为高频，从而在粗网格上被高效消除。在两网格分析（two-grid analysis）中，关键概念是光滑性（smoothing property）和逼近性（approximation property）：光滑子（如加权 Jacobi、红/黑 Gauss-Seidel）将误差高频部分衰减，而粗网格校正（coarse-grid correction）处理余下的低频误差。最优复杂度的多重网格方法（如 V-cycle 和 W-cycle）在每一层仅需 \(O(N)\) 次运算（\(N\) 为未知量总数），使得求解总代价与未知量数成线性比例，这对大规模三维问题（\(N \sim 10^9\)）是不可或缺的。区域分解方法（Domain Decomposition, DD）则将原问题裂解为若干子区域上可并行求解的子问题，通过界面上的迭代协调（如 Schwarz 交替法、FETI 方法、Neumann-Neumann 预条件子）来重建全局解。

\subsection{60.x 多重网格的两网格分析}\label{x-ux591aux91cdux7f51ux683cux7684ux4e24ux7f51ux683cux5206ux6790}

\textbf{定义 60.1}（两网格方法 / Two-Grid Method）：考虑线性系统 \(A_h \mathbf{u}_h = \mathbf{f}_h\)（网格尺度 \(h\)）。一次两网格迭代为： 1. \textbf{前光滑}（pre-smoothing）：执行 \(\nu_1\) 次光滑迭代 \(\mathbf{u}_h \leftarrow \mathbf{u}_h + M_h^{-1}(\mathbf{f}_h - A_h \mathbf{u}_h)\)（\(M_h\) 为光滑子的迭代矩阵）。 2. \textbf{粗网格校正}：计算残差 \(\mathbf{r}_h = \mathbf{f}_h - A_h \mathbf{u}_h\)；限制到粗网格 \(\mathbf{r}_{2h} = I_h^{2h} \mathbf{r}_h\)；精确求解粗网格问题 \(A_{2h} \mathbf{e}_{2h} = \mathbf{r}_{2h}\)；插值回细网格 \(\mathbf{e}_h = I_{2h}^h \mathbf{e}_{2h}\)；校正 \(\mathbf{u}_h \leftarrow \mathbf{u}_h + \mathbf{e}_h\)。 3. \textbf{后光滑}（post-smoothing）：执行 \(\nu_2\) 次光滑迭代。

迭代矩阵为 \(E_h = S_h^{\nu_2} (I - I_{2h}^h A_{2h}^{-1} I_h^{2h} A_h) S_h^{\nu_1}\)，其中 \(S_h = I - M_h^{-1} A_h\) 为光滑迭代矩阵。

\textbf{定理 60.1}（两网格收敛性 / 光滑性与逼近性，Hackbusch 1985）：对椭圆型 PDE 的标准有限元离散，两网格方法的收敛因子 \(\|E_h\| \leq \rho < 1\)（\(\rho\) 独立于 \(h\)），当且仅当光滑性和逼近性同时成立：

\begin{itemize}
\tightlist
\item
  \textbf{光滑性}（Smoothing property）：\(\|A_h S_h^\nu\| \leq \eta(\nu) / h^\alpha\)（\(h \to 0\) 时 \(\eta(\nu)\) 衰减）
\item
  \textbf{逼近性}（Approximation property）：\(\|A_h^{-1} - I_{2h}^h A_{2h}^{-1} I_h^{2h}\| \leq C_A h^\alpha\)
\end{itemize}

其中 \(\alpha\) 为微分算子的阶数（对二阶椭圆算子 \(\alpha = 2\)）。

\textbf{证明}（框架）：记粗网格校正矩阵为 \(C_h = I - I_{2h}^h A_{2h}^{-1} I_h^{2h} A_h = I - P_h\)（\(P_h\) 为 \(A_h\)-正交投影到粗网格空间）。两网格迭代矩阵为 \(E_h = S_h^{\nu_2} C_h S_h^{\nu_1}\)。光滑性给出 \(\|A_h S_h^\nu\| \leq \eta(\nu) h^{-\alpha}\)。逼近性给出 \(\|A_h^{-1} C_h\| = \|A_h^{-1} - I_{2h}^h A_{2h}^{-1} I_h^{2h}\| \leq C_A h^\alpha\)。将 \(E_h\) 写为 \(E_h = (S_h^{\nu_2} A_h^{-1}) (A_h C_h) (A_h S_h^{\nu_1})\)（仔细排列）。则 \(\|E_h\| \leq \|S_h^{\nu_2} A_h^{-1}\| \cdot \|A_h C_h\| \cdot \|A_h S_h^{\nu_1}\| \leq \eta(\nu_2) C_A h^\alpha / h^\alpha \cdot \eta(\nu_1) h^{-\alpha}\) \ldots{} 这需要精细的范数估计。正确的分析使用 \(A_h\) 能量范数或 Euclidean 范数下对偶范数的乘积。在能量范数下：光滑性等价于 \(\|S_h^\nu\|_{A_h} \leq C / \nu\)，逼近性等价于 \(\|C_h\|_{A_h^{-1}} \leq C_A\)。此时 \(\|E_h\|_{A_h} \leq C / (\nu_1 \nu_2) \cdot C_A\)，取 \(\nu = \nu_1 + \nu_2\) 充分大即可保证 \(\rho < 1\) 且独立于 \(h\)。\(\blacksquare\)

\subsection{60.y Schwarz交替法与区域分解}\label{y-schwarzux4ea4ux66ffux6cd5ux4e0eux533aux57dfux5206ux89e3}

\textbf{定义 60.2}（Schwarz交替法）：设 \(\Omega = \Omega_1 \cup \Omega_2\)（重叠区域分解）。经典 Schwarz 交替法迭代地交替求解子区域上的 Dirichlet 问题：

\[-\Delta u_1^{k+1/2} = f \; \text{in} \; \Omega_1, \quad u_1^{k+1/2} = u_2^k \; \text{on} \; \partial \Omega_1 \cap \Omega_2\] \[-\Delta u_2^{k+1} = f \; \text{in} \; \Omega_2, \quad u_2^{k+1} = u_1^{k+1/2} \; \text{on} \; \partial \Omega_2 \cap \Omega_1\]

代数形式：设 \(R_i\) 为从全局到子区域 \(\Omega_i\) 的限制算子，\(A_i = R_i A R_i^T\) 为子区域离散矩阵。加性 Schwarz 预条件子为 \(M_{\text{AS}}^{-1} = \sum_i R_i^T A_i^{-1} R_i\)，乘性 Schwarz 等价于块 Gauss-Seidel。

\textbf{定理 60.2}（加性Schwarz方法的条件数上界，Dryja-Widlund 1994）：对于二阶椭圆 PDE 的有限元离散，以重叠区域分解构造的加性 Schwarz 预条件子 \(M_{\text{AS}}^{-1}\)，预条件系统的条件数满足

\[\kappa_2(M_{\text{AS}}^{-1} A) \leq C \left(1 + \frac{H}{\delta}\right)\]

其中 \(H\) 为子区域直径，\(\delta\) 为重叠宽度。若 \(\delta \sim H\)（慷慨重叠），\(\kappa_2 \leq C\)（独立于网格尺寸和子区域数）；若 \(\delta = 0\)（非重叠），需引入粗空间分量以保证条件数有界。

\textbf{证明}（框架）：Schwarz 理论的三参数估计（Lions 1988）：需验证加强 Cauchy-Schwarz 不等式（子区域间相互作用矩阵的谱界）和局部稳定性（分割引理 / partition of unity）。设 \(\{\theta_i\}_{i=1}^N\) 为从属于覆盖 \(\{\Omega_i\}\) 的单位分解。则对任意 \(u \in V_h\)，分解 \(u = \sum_{i=1}^N u_i\)（\(u_i = I_h(\theta_i u)\)，\(I_h\) 为有限元插值）。引理 1（局部稳定性）：\(\|u_i\|_{H^1(\Omega_i)} \leq C \|u\|_{H^1(\Omega_i^\delta)}\)（\(\Omega_i^\delta\) 为 \(\Omega_i\) 的 \(\delta\)-邻域）。引理 2（Cauchy-Schwarz 加强性）：\(\sum_{i,j} |a(u_i, v_j)| \leq C_{\text{CS}} (\sum_i \|u_i\|_A^2)^{1/2} (\sum_j \|v_j\|_A^2)^{1/2}\)，其中 \(C_{\text{CS}} \leq C (H/\delta)\)。由 Schwarz 框架的必要充分条件（抽象 Schwarz 理论，Smith-Bjorstad-Gropp 1996），\(\kappa_2 \leq C_{\text{CS}} \cdot (1 + \rho(E))\)（\(\rho(E)\) 为子区域交互矩阵的谱半径）。对慷慨重叠 \(\delta = O(H)\)，\(C_{\text{CS}} = O(1)\)，得证。\(\blacksquare\)

\subsection{60.z FETI方法与Neumann-Neumann预条件子}\label{z-fetiux65b9ux6cd5ux4e0eneumann-neumannux9884ux6761ux4ef6ux5b50}

\textbf{定理 60.3}（FETI方法的理论框架 / Farhat-Roux 1991）：FETI（Finite Element Tearing and Interconnecting）方法将区域分解为不重叠子区域 \(\Omega = \bigcup_{s=1}^N \Omega_s\)，界面连续性通过 Lagrange 乘子强制。消去内部变量后，界面 Schur 补系统 \(F \boldsymbol{\lambda} = \mathbf{d}\) 的求解通过预条件共轭梯度法完成。预条件子 \(M_{\text{FETI}}^{-1} = \sum_{s=1}^N B_s \tilde{S}_s B_s^T\)（Dirichlet 预条件子）或 \(M_{\text{NN}}^{-1}\)（Neumann-Neumann 预条件子）使得条件数满足 \(\kappa_2(M^{-1} F) \leq C (1 + \log(H/h))^2\)（对二阶椭圆问题），其中 \(H\) 为子区域直径，\(h\) 为细网格尺寸。

\textbf{证明}：FETI 方法的核心步骤：（1）在各子区域 \(\Omega_s\) 分解刚度矩阵 \(K_s\) 为内部和界面分块；（2）Schur 补 \(S_s = K_{s,\Gamma\Gamma} - K_{s,\Gamma I} K_{s,II}^{-1} K_{s,I\Gamma}\)；（3）全局 Schur 补系统 \(S \mathbf{u}_\Gamma = \mathbf{f}_\Gamma\)，连续性约束 \(B \mathbf{u}_\Gamma = 0\)（\(B\) 为跳量算子）；（4）对偶求解 \(F \boldsymbol{\lambda} = \mathbf{d}\) 其中 \(F = B S^{-1} B^T\)。预条件子分析基于分式 Sobolev 空间 \(H^{1/2}_{00}(\Gamma_s)\) 中界面算子的谱等价性。关键工具为 \(p = 2\) 时界面算子的''对数界''：\(\|B_s^T \boldsymbol{\mu}\|_{1/2, \partial \Omega_s} \leq C (1 + \log(H/h)) \|\boldsymbol{\mu}\|_{1/2, \partial \Omega_s}\)（Sobolev 空间嵌入不等式 / Bramble-Pasciak-Schatz 1986）。该对数因子通过 Schwarz 引理和 Fourier 分析得到。预条件子 \(M^{-1}\) 的构造：加权计数各子区域贡献，通过伪逆处理浮点子区域（局部刚度矩阵的零空间 ------ 刚体模态）。完整的条件数估计涉及 \(C(1+\log(H/h))^q\)（\(q\) 取决于问题类型，对二阶椭圆 \(q=2\)，对板弯曲 \(q=3\)）。\(\blacksquare\)

\hrulefill

\subsection{315.A 数值分析补充}\label{a-ux6570ux503cux5206ux6790ux8865ux5145}

\begin{quote}
计算数学将偏微分方程等连续数学问题转化为离散问题的数学理论，其基础涵盖泛函分析、逼近论和线性代数。本卷建立在数值分析（V23）和微分方程（V18）的基础上，系统建立有限元方法的数学理论、谱方法的逼近性质、多重网格法的收敛性分析，以及守恒律和流体方程的数值格式理论。计算数学的核心在于精度与效率的数学分析------收敛阶、稳定性条件和计算复杂度的严格估计。
\end{quote}

\newpage

\chapter{11-数值与计算/01-数值分析/02-有限元与计算流体力学.md}\label{ux6570ux503cux4e0eux8ba1ux7b9701-ux6570ux503cux5206ux679002-ux6709ux9650ux5143ux4e0eux8ba1ux7b97ux6d41ux4f53ux529bux5b66.md}

\section{第389章：有限元方法理论基础}\label{ux7b2c389ux7ae0ux6709ux9650ux5143ux65b9ux6cd5ux7406ux8bbaux57faux7840}

有限元方法（FEM）是求解偏微分方程最强大的数值工具之一。其核心思想是将问题从无穷维函数空间的强形式转化为有限维子空间中的弱形式（变分形式），从而将 PDE 转化为代数方程组。

\subsection{277.1 变分原理与弱形式}\label{ux53d8ux5206ux539fux7406ux4e0eux5f31ux5f62ux5f0f}

\textbf{定义 277.1}（抽象变分问题）：设 \(V\) 为 Hilbert 空间，\(a: V \times V \to \mathbb{R}\) 为连续强制双线性形式，\(f \in V'\)。求 \(u \in V\) 使得

\[a(u, v) = f(v), \quad \forall v \in V\]

\textbf{定理 277.1}（Lax-Milgram 引理）：若 \(a\) 满足有界性 \(|a(u,v)| \leq M\|u\|\|v\|\) 和强制性 \(a(u,u) \geq \alpha\|u\|^2\)（\(\alpha > 0\)），则变分问题存在唯一解 \(u \in V\) 且 \(\|u\|_V \leq \frac{1}{\alpha}\|f\|_{V'}\)。

变分原理将 PDE 的强形式转化为弱形式，其核心优势在于降低了对解的正则性要求。在强形式中，解需要足够光滑以使得微分算子逐点有意义；而在弱形式中，仅需解在 Sobolev 空间中满足积分恒等式。这一转换是泛函分析在 PDE 理论中最成功的应用之一。变分问题与能量最小化问题等价：若 \(a\) 对称，则 \(u\) 是能量泛函 \(J(v) = \frac{1}{2}a(v,v) - f(v)\) 在 \(V\) 上的唯一极小元。非对称情形（如对流-扩散方程）不能写为能量极小化，但 Lax-Milgram 引理仍保证解的存在唯一性。强制性条件 \(a(u,u) \geq \alpha\|u\|^2\) 是椭圆性的泛函表述，其物理意义为系统具有正定的''刚度''------能量不会无界地降低。在有限元方法的实际应用中，双线性形式 \(a\) 通常被分解为单元贡献之和 \(a(u,v) = \sum_{K \in \mathcal{T}_h} a_K(u|_K, v|_K)\)，而 \(f\) 被分解为 \(f(v) = \sum_K \int_K f v\,dx\)。这种局部化是有限元方法高效实现的基础。

\emph{证明}：对固定 \(u\in V\)，\(v\mapsto a(u,v)\) 是有界线性的。Riesz 表示定理给出 \(Au\in V\) 使 \(a(u,v)=\langle Au,v\rangle\)。\(A\) 有界且 \(\langle Au,u\rangle\ge\alpha\|u\|^2\) 导出 \(A\) 单射且值域闭。由 Lax-Milgram，\(A\) 是满射，故 \(Au=\tau(f)\) 唯一可解。\(\blacksquare\)

\subsection{277.2 有限元离散与 Céa 引理}\label{ux6709ux9650ux5143ux79bbux6563ux4e0e-cuxe9a-ux5f15ux7406}

\textbf{定义 277.2}（Galerkin 逼近）：选取有限维子空间 \(V_h \subset V\)，其中 \(h\) 为网格参数。求 \(u_h \in V_h\) 使得

\[a(u_h, v_h) = f(v_h), \quad \forall v_h \in V_h\]

\textbf{定理 277.2}（Céa 引理）：设 \(a\) 连续强制，\(u\) 和 \(u_h\) 分别为连续问题和离散问题的解。则

\[\|u - u_h\|_V \leq \frac{M}{\alpha} \inf_{v_h \in V_h} \|u - v_h\|_V\]

即有限元解的误差控制在最佳逼近误差的常数倍之内。

\emph{证明}：由 Galerkin 正交性 \(a(u-u_h,v_h)=0\)（\(\forall v_h\in V_h\)）。取 \(v_h\in V_h\)，\(a(u-u_h,u-u_h)=a(u-u_h,u-v_h)\le M\|u-u_h\|\|u-v_h\|\)。强制性得 \(\alpha\|u-u_h\|^2\le M\|u-u_h\|\|u-v_h\|\)，故 \(\|u-u_h\|\le(M/\alpha)\|u-v_h\|\)。取下确界即得。\(\blacksquare\)

\textbf{定义 277.3}（Lagrange 有限元空间）：对区域 \(\Omega\) 的三角剖分 \(\mathcal{T}_h\)，定义分片多项式空间

\[V_h^k = \{v_h \in C^0(\bar{\Omega}) : v_h|_K \in P_k(K), \forall K \in \mathcal{T}_h, v_h|_{\partial\Omega} = 0\}\]

其中 \(P_k\) 为次数不超过 \(k\) 的多项式空间。

\textbf{定理 277.3}（插值误差估计）：设 \(u \in H^{k+1}(\Omega)\)，\(I_h u\) 为其在 \(V_h^k\) 中的节点插值。则

\[\|u - I_h u\|_{L^2} + h\|u - I_h u\|_{H^1} \leq C h^{k+1} |u|_{H^{k+1}}\]

\textbf{证明}：由有限元空间的逼近性质和变分公式的强制性，Céa引理给出有限元解\(u_h\)与真解\(u\)的误差估计：\(\|u-u_h\|_V\le C\inf_{v_h\in V_h}\|u-v_h\|_V\)。插值误差估计（Bramble-Hilbert引理）给出\(\inf_{v_h\in V_h}\|u-v_h\|_V\le Ch^{k}\|u\|_{H^{k+1}}\)（\(k\)为多项式次数）。结合即得收敛阶。\(\blacksquare\)

\textbf{定理 280.1}（Godunov 定理）：线性常系数情形下，单调守恒格式的精度最多为一阶。

\textbf{证明}：设数值通量 \(\hat{f}(u_L,u_R)\) 单调（即 \(\hat{f}\) 对各变量非减）。对线性对流方程 \(u_t+au_x=0\)，\(\hat{f}(u_L,u_R)=a^+u_L+a^-u_R\)（\(a^+=\max(a,0)\)，\(a^-=\min(a,0)\)）是唯一的单调且相容的通量（Godunov 格式即迎风格式）。任何具有二阶精度的线性格式（如 Lax-Wendroff）必含负系数、非单调，在间断处产生振荡。证明核心：若格式具有二阶精度且线性，其修正方程含三阶色散项，在激波处产生 Gibbs 振荡；要消除振荡必须引入足够数值粘性（一阶），破坏了二阶精度。Harten 等人证明：TVD 格式在极值点处退化为一阶------二阶精度与 TVD 性质在光滑极值点处矛盾。\(\blacksquare\)

\textbf{定义 280.4}（Godunov 型格式）：在每个单元界面上精确求解 Riemann 问题（初值间断的守恒律），然后对解进行单元平均。Riemann 问题的自相似解是 Godunov 方法的核心组成部分。

\subsection{280.3 高分辨率格式：TVD、ENO、WENO}\label{ux9ad8ux5206ux8fa8ux7387ux683cux5f0ftvdenoweno}

\textbf{定义 280.5}（TVD 格式）：数值格式称为 \textbf{TVD}（Total Variation Diminishing），若

\[TV(u^{n+1}) \leq TV(u^n), \quad TV(u) = \sum_j |u_{j+1} - u_j|\]

TVD 条件确保无虚假振荡（Gibbs 现象被抑制）。

\textbf{定理 280.2}（Harten 引理）：若格式可写为 \(u_j^{n+1} = u_j^n - C_{j-1/2}(u_j^n - u_{j-1}^n) + D_{j+1/2}(u_{j+1}^n - u_j^n)\) 且 \(C_{j-1/2}, D_{j+1/2} \geq 0\)，\(C_{j-1/2} + D_{j+1/2} \leq 1\)，则该格式是 TVD 的。

\textbf{证明}：计算总变差 \(TV(u^{n+1})=\sum_j|u_{j+1}^{n+1}-u_j^{n+1}|\)。由格式表达式， \[u_{j+1}^{n+1}-u_j^{n+1}=(1-C_{j+1/2}-D_{j+1/2})(u_{j+1}^n-u_j^n)+C_{j-1/2}(u_j^n-u_{j-1}^n)+D_{j+3/2}(u_{j+2}^n-u_{j+1}^n)\] 因系数均非负，取绝对值并用三角不等式求和，重排下标后得 \[TV(u^{n+1})\le\sum_j(1-C_{j+1/2}-D_{j+1/2}+C_{j-1/2}+D_{j+1/2})|u_{j+1}^n-u_j^n|\] 因 \(C_{j-1/2}\) 移到第 \(j\) 项、\(D_{j+1/2}\) 移到第 \(j\) 项后系数和 \(\le1\)，故 \(TV(u^{n+1})\le TV(u^n)\)。\(\blacksquare\)

\textbf{定义 280.6}（ENO 与 WENO 格式）：ENO（Essentially Non-Oscillatory）格式通过自适应选择最光滑的模板构造高阶重构多项式，避免跨间断插值。WENO（Weighted ENO）进一步通过凸组合多个模板，在光滑区域获得更高阶精度，在间断附近自动降阶------实现了高精度和无振荡的统一。

\textbf{定理 280.3}（Lax-Wendroff 定理）：若守恒格式的数值解 \(u_h\) 在 \(h \to 0\) 时几乎处处收敛到某个函数 \(u\)，且格式守恒，则 \(u\) 是守恒律的弱解。

\textbf{证明}：守恒格式 \(u_j^{n+1}=u_j^n-\lambda(\hat{f}_{j+1/2}-\hat{f}_{j-1/2})\)（\(\lambda=\Delta t/\Delta x\)）乘以试验函数 \(\phi(x_j,t_n)\) 并求和，整理后得离散弱形式： \[\sum_{j,n}u_j^n\frac{\phi_j^{n+1}-\phi_j^n}{\Delta t}\Delta x\Delta t+\sum_{j,n}\hat{f}_{j+1/2}^n\frac{\phi_{j+1}^n-\phi_j^n}{\Delta x}\Delta x\Delta t+\sum_j u_j^0\phi_j^0\Delta x=0\] 取极限 \(h\to0\)，由 \(u_h\to u\) a.e. 和 \(\hat{f}\) 的相容性（\(\hat{f}(u,u)=f(u)\)），利用控制收敛定理得 \[\int_0^\infty\int_{-\infty}^\infty(u\partial_t\phi+f(u)\partial_x\phi)dx dt+\int_{-\infty}^\infty u(x,0)\phi(x,0)dx=0\] 即 \(u\) 在分布意义下满足守恒律。\(\blacksquare\)

\subsection{280.4 一维和二维 Euler 方程的数值求解}\label{ux4e00ux7ef4ux548cux4e8cux7ef4-euler-ux65b9ux7a0bux7684ux6570ux503cux6c42ux89e3}

\textbf{定义 280.7}（一维 Euler 方程）：理想气体的 Euler 方程为

\[\frac{\partial}{\partial t}\begin{pmatrix} \rho \\ \rho u \\ E \end{pmatrix} + \frac{\partial}{\partial x}\begin{pmatrix} \rho u \\ \rho u^2 + p \\ u(E + p) \end{pmatrix} = 0\]

其中 \(\rho\) 为密度，\(u\) 为速度，\(E\) 为总能量，\(p = (\gamma - 1)(E - \frac{1}{2}\rho u^2)\) 为压力（理想气体状态方程）。

\textbf{Riemann 问题的波系结构}：一维 Euler 方程的 Riemann 问题产生三种基本波------激波（shock）、接触间断（contact discontinuity）和稀疏波（rarefaction wave）。精确 Riemann 求解器通过迭代求解非线性方程确定中间状态；近似 Riemann 求解器（如 Roe 求解器、HLL/HLLC 求解器）以较低计算代价获得可接受的精度。二维 Euler 方程的数值求解通常采用 \textbf{维数分裂（dimensional splitting）} 或 \textbf{非结构网格} 上的有限体积法。

\hrulefill

\hrulefill

\hrulefill

\hrulefill

\hrulefill

\hrulefill

\section{第390章：计算流体力学数学基础}\label{ux7b2c390ux7ae0ux8ba1ux7b97ux6d41ux4f53ux529bux5b66ux6570ux5b66ux57faux7840}

Navier-Stokes 方程的数值求解涉及不可压缩流、可压缩流和湍流的多尺度建模，其数学分析是计算数学最具挑战性的领域之一。

\subsection{281.1 Navier-Stokes 方程的数值挑战}\label{navier-stokes-ux65b9ux7a0bux7684ux6570ux503cux6311ux6218}

\textbf{定义 281.1}（不可压缩 Navier-Stokes 方程）：

\[\frac{\partial \mathbf{u}}{\partial t} + (\mathbf{u} \cdot \nabla)\mathbf{u} = -\nabla p + \nu \Delta \mathbf{u} + \mathbf{f}, \quad \nabla \cdot \mathbf{u} = 0\]

\textbf{挑战}：(1) 不可压缩约束 \(\nabla \cdot \mathbf{u} = 0\) 导致速度-压力耦合的鞍点问题------离散化必须满足 inf-sup 条件（LBB 条件）；(2) 高 Reynolds 数下对流项占优，需要稳定化处理；(3) 湍流的多尺度特性使直接数值模拟（DNS）的计算代价极高（\(\propto Re^3\)）。

速度-压力耦合的鞍点结构源于不可压缩约束作为 Lagrange 乘子条件的数学本质。离散化时，速度空间 \(V_h\) 和压力空间 \(Q_h\) 必须满足 \textbf{Ladyzhenskaya-Babuska-Brezzi (LBB) 条件}（亦称 inf-sup 条件）：\(\inf_{q_h \in Q_h} \sup_{v_h \in V_h} \frac{(\nabla \cdot v_h, q_h)}{\|v_h\|_V \|q_h\|_Q} \geq \beta > 0\)，其中 \(\beta\) 独立于 \(h\)。不满足 LBB 条件的有限元对会产生虚假压力模式（spurious pressure modes）------非物理解。经典满足 LBB 条件的元对包括 Taylor-Hood 元（\(P_k/P_{k-1}\)，\(k \geq 2\)）和 MINI 元（\(P_1\)-bubble/\(P_1\)）。高 Reynolds 数下的对流占优导致边界层和内部层，标准 Galerkin 方法产生非物理振荡。稳定化方法包括 \textbf{Streamline Upwind Petrov-Galerkin (SUPG)} 和 \textbf{Galerkin Least-Squares (GLS)}，其核心是在变分形式中添加网格依赖的稳定化项：\(a_{\text{stab}}(u_h, v_h) = a(u_h, v_h) + \sum_K \tau_K (\mathcal{L}u_h - f, \mathcal{L}^* v_h)_K\)。湍流建模（RANS、LES、DNS）是 CFD 中最具挑战性的问题------DNS 解析所有尺度（计算量 \(\propto Re^3\)），LES 仅解析大涡并用亚网格模型处理小尺度，RANS 对所有湍流尺度进行时间平均建模。

\textbf{定义 281.2}（无量纲参数）：Reynolds 数 \(Re = U L / \nu\) 刻画惯性力与粘性力之比。低 \(Re\) 为层流（稳态光滑），高 \(Re\) 为湍流（非稳态混沌）。另外，Mach 数 \(Ma = U/c\)（\(c\) 为声速）判明压缩性------\(Ma < 0.3\) 可视作不可压缩。

\subsection{281.2 投影法与分步法}\label{ux6295ux5f71ux6cd5ux4e0eux5206ux6b65ux6cd5}

\textbf{定义 281.3}（Chorin-Temam 投影法）：(1) 预测步：忽略压力梯度求解中间速度

\[\frac{\mathbf{u}^* - \mathbf{u}^n}{\Delta t} = \nu \Delta \mathbf{u}^* - (\mathbf{u}^n \cdot \nabla)\mathbf{u}^n + \mathbf{f}\]

\begin{enumerate}
\def\labelenumi{(\arabic{enumi})}
\setcounter{enumi}{1}
\tightlist
\item
  投影步：将 \(\mathbf{u}^*\) 投影到无散度空间 \(\mathbf{u}^{n+1} = \mathcal{P}(\mathbf{u}^*)\)，其中 \(\mathcal{P} = I - \nabla \Delta^{-1} \nabla\cdot\)（Leray-Helmholtz 投影），投影步需求解 Poisson 方程 \(\Delta \phi = \nabla \cdot \mathbf{u}^*\)。
\end{enumerate}

\textbf{定理 281.1}（投影法的 Helmholtz-Hodge 分解）：对任意充分光滑的向量场 \(\mathbf{w}\)，存在唯一正交分解 \(\mathbf{w} = \mathbf{u} + \nabla q\)，其中 \(\nabla \cdot \mathbf{u} = 0\) 且 \(\mathbf{u} \cdot \mathbf{n}|_{\partial\Omega} = 0\)。投影算子 \(\mathcal{P}: \mathbf{w} \mapsto \mathbf{u}\) 是 \(L^2\) 正交投影。

\textbf{证明}：设 \(q\) 为 Neumann 问题 \(\Delta q=\nabla\cdot\mathbf{w}\)（\(\partial q/\partial n=\mathbf{w}\cdot\mathbf{n}\) 在 \(\partial\Omega\) 上）的解，相容条件 \(\int_\Omega\nabla\cdot\mathbf{w}=\int_{\partial\Omega}\mathbf{w}\cdot\mathbf{n}\) 保证解存在（差常数）。定义 \(\mathbf{u}=\mathbf{w}-\nabla q\)，则 \(\nabla\cdot\mathbf{u}=0\) 且 \(\mathbf{u}\cdot\mathbf{n}|_{\partial\Omega}=0\)。正交性：对任意满足 \(\nabla\cdot\mathbf{v}=0\) 且 \(\mathbf{v}\cdot\mathbf{n}=0\) 的 \(\mathbf{v}\)，\(\int_\Omega\mathbf{u}\cdot\nabla q=-\int_\Omega(\nabla\cdot\mathbf{u})q+\int_{\partial\Omega}(\mathbf{u}\cdot\mathbf{n})q=0\)，故 \(\mathbf{u}\perp\nabla q\)（\(L^2\) 内积），即 \(\mathcal{P}\) 为正交投影。唯一性来自 Poisson 方程解的唯一性。\(\blacksquare\)

\textbf{定理 281.2}（投影法的精度）：标准 Chorin 投影法对速度是一阶时间精度的。通过引入压力的时间外推（如压力校正投影法或旋转投影法），可将时间精度提升至二阶。

\textbf{证明}：Chorin 投影法为：\((u^*-u^n)/\Delta t-\nu\Delta u^*=-(u^n\cdot\nabla)u^n+f^n\)，\(u^{n+1}=u^*-\Delta t\nabla p^{n+1}\)，\(\nabla\cdot u^{n+1}=0\)。将两式合并：\((u^{n+1}-u^n)/\Delta t+\nabla p^{n+1}-\nu\Delta u^*=-(u^n\cdot\nabla)u^n+f^n\)。由于 \(\Delta u^*=\Delta(u^{n+1}+\Delta t\nabla p^{n+1})=\Delta u^{n+1}+\Delta t\Delta\nabla p^{n+1}\)，与真解比较得 \(O(\Delta t)\) 局部截断误差（压力的时间滞后）。旋转投影法中将 \(\nabla p\) 替换为 \(\nabla\times\nabla\times u\) 形式处理涡量边界，消除 \(\Delta t\Delta\nabla p\) 项，恢复 \(O(\Delta t^2)\) 精度。\(\blacksquare\)

\textbf{注 281.1}（旋转投影法 / Rotational Projection Method）：标准投影法的 Helmholtz-Hodge 分解中，压力项采用 \(\nabla p\) 的形式，导致边界层误差积累限制时间精度。\textbf{旋转投影法}（Timmermans-Minev-Van De Vosse 1996, Guermond-Shen 2006）将压力梯度替换为 \(\nabla \times \nabla \times\) 形式（利用恒等式 \(\Delta = \nabla(\nabla\cdot) - \nabla \times \nabla \times\)），通过显式处置涡量边界条件，消除了标准投影法在无滑移边界上产生的法向压力梯度的人工边界层，从而将整体时间精度提升至二阶。该方法在旋转坐标系（rotational form of Navier-Stokes）下自然导出，故名''旋转投影法''。

\subsection{281.3 稳定化有限元方法}\label{ux7a33ux5b9aux5316ux6709ux9650ux5143ux65b9ux6cd5}

\textbf{定义 281.4}（对流-扩散方程）：\(-\varepsilon \Delta u + \mathbf{b} \cdot \nabla u = f\)。当 Péclet 数 \(Pe = \|\mathbf{b}\|h/(2\varepsilon) > 1\) 时，标准 Galerkin 方法产生非物理振荡。需引入稳定化。

\textbf{SUPG（Streamline Upwind Petrov-Galerkin）方法}：修改检验函数 \(v \mapsto v + \tau \mathbf{b}\cdot\nabla v\)，沿流线方向增加人工扩散而不破坏精度。\(\tau\) 的选取通过误差分析确定：

\[\tau_K = \frac{h_K}{2\|\mathbf{b}\|}\left(\coth Pe_K - \frac{1}{Pe_K}\right)\]

其中 \(Pe_K = \|\mathbf{b}\| h_K / (2\varepsilon)\) 为单元 Péclet 数。该公式的含义为：当对流占优时（\(Pe_K \gg 1\)，\(\coth Pe_K \to 1\)），\(\tau_K \approx h_K / (2\|\mathbf{b}\|)\)，沿流线方向引入的稳定化量与局部网格尺寸和流速匹配；当扩散占优时（\(Pe_K \to 0\)，\(\coth Pe_K - 1/Pe_K \approx Pe_K / 3\)），\(\tau_K \approx \varepsilon h_K^2 / (6\|\mathbf{b}\|^2)\)，稳定化自动减弱以避免过度扩散。\(\tau_K\) 的这一双渐近公式确保 SUPG 在从纯对流到纯扩散的整个参数范围内都表现最优。

\textbf{GLS（Galerkin Least-Squares）方法}：在变分形式中加入最小二乘项 \(\sum_K \tau_K (\mathcal{L}v, \mathcal{L}u - f)_K\)，同时稳定对流项和压力项，适用于多物理场耦合的复杂情形。

\subsection{281.4 湍流模型与 LES}\label{ux6e4dux6d41ux6a21ux578bux4e0e-les}

\textbf{定义 281.5}（大涡模拟 / LES）：不直接求解所有尺度的湍流运动，而是通过滤波操作 \(\bar{u}(\mathbf{x}) = \int G(\mathbf{x} - \boldsymbol{\xi}) u(\boldsymbol{\xi}) d\boldsymbol{\xi}\) 分离大尺度（可解析尺度）和小尺度（亚格子尺度），并对亚格子应力建立封闭模型。

\textbf{Smagorinsky 模型}：涡粘性假设

\[\tau_{ij} - \frac{1}{3}\tau_{kk}\delta_{ij} = -2\nu_t \bar{S}_{ij}, \quad \nu_t = (C_s \Delta)^2 |\bar{S}|\]

其中 \(\bar{S}_{ij} = \frac{1}{2}(\partial \bar{u}_i/\partial x_j + \partial \bar{u}_j/\partial x_i)\) 为滤波后的应变率张量，\(|\bar{S}| = \sqrt{2\bar{S}_{ij}\bar{S}_{ij}}\)，\(C_s \approx 0.1\)--\(0.2\) 为 Smagorinsky 常数，\(\Delta\) 为滤波宽度。

\textbf{RANS（Reynolds-Averaged Navier-Stokes）方法}：将瞬时量分解为时间平均和脉动量（\(u = \bar{u} + u'\)），得到 Reynolds 平均方程------包含 Reynolds 应力张量 \(\overline{u_i' u_j'}\) 需湍流模型封闭（如 \(k\)-\(\varepsilon\) 模型、\(k\)-\(\omega\) SST 模型）。

\textbf{定理 281.3}（DNS 的计算代价）：对均匀各向同性湍流，直接数值模拟（DNS）的空间自由度 \(N \propto Re^{9/4}\)，计算代价 \(\propto Re^{3}\)。这意味着 \(Re\) 每增大 10 倍，计算量增大约 1000 倍------这解释了为何 DNS 目前仅限于低至中等 Reynolds 数的流动。

\textbf{证明}：Kolmogorov 尺度 \(\eta\sim L\cdot Re^{-3/4}\)（\(L\) 为积分尺度）。DNS 需解析所有尺度到 Kolmogorov 耗散尺度，故空间网格间距 \(h\sim\eta\)。三维网格点数 \(N\sim(L/h)^3\sim Re^{9/4}\)。时间步长受 CFL 条件和 Kolmogorov 时间尺度 \(\tau_\eta\sim\eta/u_\eta\)（\(u_\eta\sim U Re^{-1/4}\)）限制：\(\Delta t\sim\tau_\eta\sim(L/U)Re^{-1/2}\)。总积分时间 \(T\sim L/U\)。故时间步数 \(N_t\sim T/\Delta t\sim Re^{1/2}\)。每步计算量 \(\sim N\)（对谱方法为 \(N\log N\)），总运算量 \(\sim N\cdot N_t\sim Re^{9/4}\cdot Re^{1/2}=Re^{11/4}\approx Re^{3}\)（含常数因子）。\(\blacksquare\)

\subsection{结语}\label{ux7ed3ux8bed}

计算数学将泛函分析、PDE 理论和线性代数的数学结构转化为可分析的算法框架。从有限元的 Sobolev 空间理论到多重网格的 \(\mathcal{O}(N)\) 复杂度，从双曲守恒律的熵条件到 Navier-Stokes 的投影法------计算数学的每一步发展都建立在严格的数学理论基础之上。

\hrulefill

\hrulefill

\hrulefill

\newpage

\chapter{11-数值与计算/01-数值分析/03-特征值与积分.md}\label{ux6570ux503cux4e0eux8ba1ux7b9701-ux6570ux503cux5206ux679003-ux7279ux5f81ux503cux4e0eux79efux5206.md}

\section{第391章：数值积分与微分}\label{ux7b2c391ux7ae0ux6570ux503cux79efux5206ux4e0eux5faeux5206}

数值积分（求积）和数值微分是逼近导数和定积分的算法。本章讨论 Newton-Cotes 公式、Gauss 求积和 Romberg 外推。

\subsection{117.1 数值积分基础}\label{ux6570ux503cux79efux5206ux57faux7840}

\textbf{定义 117.1}（求积法则）：\textbf{求积法则} \(Q(f) = \sum_{i=0}^n w_i f(x_i)\) 近似积分 \(\int_a^b f(x) dx\)。\(x_i\) 称为\textbf{求积节点}，\(w_i\) 称为\textbf{求积权}。

\textbf{定义 117.2}（代数精度）：求积法则的\textbf{代数精度}为 \(m\)，如果它对所有次数 \(\leq m\) 的多项式精确，但对某个 \(m+1\) 次多项式不精确。

\textbf{定理 117.1}（Newton-Cotes 公式）：使用等距节点 \(x_i = a + ih\)（\(h = (b-a)/n\)）的插值型求积： - \(n = 1\)（梯形法则）：\(\int_a^b f(x) dx \approx \frac{h}{2}[f(a) + f(b)]\)，误差 \(-\frac{h^3}{12} f''(\xi)\) - \(n = 2\)（Simpson 法则）：\(\int_a^b f(x) dx \approx \frac{h}{3}[f(a) + 4f((a+b)/2) + f(b)]\)，误差 \(-\frac{h^5}{90} f^{(4)}(\xi)\) - \(n = 3\)（Simpson 3/8 法则）：误差 \(-\frac{3h^5}{80} f^{(4)}(\xi)\)

\textbf{证明}：Newton-Cotes 公式由 Lagrange 插值多项式 \(p_n\) 积分得到：\(\int_a^b f(x)dx\approx\int_a^b p_n(x)dx=\sum w_i f(x_i)\)，\(w_i=\int_a^b L_i(x)dx\)。梯形法则（\(n=1\)）：\(p_1\) 为线性，\(\int_a^b p_1=h(f(a)+f(b))/2\)。误差 \(\int_a^b(f-p_1)=-\frac{1}{2}\int_a^b f''(\xi_x)(x-a)(x-b)dx\)，由积分中值定理得 \(-h^3 f''(\xi)/12\)。Simpson 法则（\(n=2\)）：\(p_2\) 为二次，\(\int_a^b p_2=\frac{h}{3}[f(a)+4f(\frac{a+b}{2})+f(b)]\)。误差 \(\int_a^b(f-p_2)=-\frac{1}{24}\int_a^b f^{(4)}(\xi_x)(x-a)(x-\frac{a+b}{2})^2(x-b)dx\)，化简得 \(-h^5 f^{(4)}(\xi)/90\)。\(\blacksquare\)

\textbf{定义 117.3}（复合求积法则）：将 \([a, b]\) 分为 \(N\) 个子区间，每个子区间上应用低阶求积法则： - 复合梯形法则：\(T_N = \frac{h}{2}[f(a) + 2\sum_{i=1}^{N-1} f(x_i) + f(b)]\)，\(h = (b-a)/N\)，误差 \(O(h^2)\) - 复合 Simpson 法则：误差 \(O(h^4)\)

\subsection{117.2 Gauss 求积}\label{gauss-ux6c42ux79ef}

\textbf{定理 117.2}（Gauss 求积）：使用 \(n+1\) 个节点的求积法则的代数精度最多为 \(2n+1\)。达到此上界的求积法则称为 \textbf{Gauss 求积}。

\textbf{证明}：上界 \(2n+1\)：构造 \(\omega(x)=\prod_{i=0}^n(x-x_i)^2\in\mathcal{P}_{2n+2}\)，\(\omega(x_i)=0\) 且 \(\omega(x)\ge0\)，故 \(Q(\omega)=\sum w_i\cdot0=0\)，但 \(\int_a^b\omega(x)w(x)dx>0\)（\(w(x)\) 为非负权函数）。故 \(\omega\) 不被精确积分，代数精度 \(\le 2n+1\)。Gauss 求积达此界：选择节点为关于 \(w(x)\) 的正交多项式 \(P_{n+1}\) 的零点，则对任意 \(f\in\mathcal{P}_{2n+1}\)，带余除法 \(f=q\cdot P_{n+1}+r\)（\(\deg q\le n\)，\(\deg r\le n\)）。由正交性 \(\int q\cdot P_{n+1}\cdot w=0\)。\(r\) 被 Lagrange 插值精确积分，故 \(\int fw=\int rw=\sum w_i r(x_i)=\sum w_i f(x_i)\)（因 \(P_{n+1}(x_i)=0\)）。\(\blacksquare\)

\textbf{定理 117.3}（Gauss 求积的构造）：Gauss 求积的节点 \(x_i\) 是区间 \([a, b]\) 上关于权函数 \(w(x)\) 的正交多项式 \(P_{n+1}(x)\) 的零点。权 \(w_i\) 由插值条件确定。

\textbf{证明}：设 \(\{P_k\}\) 为关于权 \(w\) 的正交多项式序列。\(P_{n+1}\) 有 \(n+1\) 个互异实零点 \(x_0,\ldots,x_n\in(a,b)\)（正交多项式的标准性质）。定义权 \(w_i=\int_a^b L_i(x)w(x)dx\)，其中 \(L_i\) 是 \(x_i\) 处的 Lagrange 基函数。对任意 \(f\in\mathcal{P}_{2n+1}\)，如前所述由正交性得 \(\int fw=\sum w_i f(x_i)\)，验证了精度 \(2n+1\)。权的解析公式：\(w_i=\frac{\langle P_n,P_n\rangle_w}{P_n(x_i)P_{n+1}'(x_i)}\)（Christoffel-Darboux 公式）。\(\blacksquare\)

\subsection{117.3 Romberg 外推与自适应求积}\label{romberg-ux5916ux63a8ux4e0eux81eaux9002ux5e94ux6c42ux79ef}

\textbf{定理 117.4}（Euler-Maclaurin 公式）：复合梯形法则的误差有渐近展开

\[T_N = \int_a^b f(x) dx + c_2 h^2 + c_4 h^4 + c_6 h^6 + \cdots\]

\textbf{证明}：Euler-Maclaurin 求和公式为 \(\sum_{i=1}^{N-1}f(a+ih)=\frac{1}{h}\int_a^b f(x)dx-\frac{f(a)+f(b)}{2}+\sum_{k=1}^m\frac{B_{2k}h^{2k-1}}{(2k)!}(f^{(2k-1)}(b)-f^{(2k-1)}(a))+R_m\)，其中 \(B_{2k}\) 为 Bernoulli 数。乘以 \(h\) 并整理：\(T_N=h(\frac{f(a)}{2}+\sum_{i=1}^{N-1}f(a+ih)+\frac{f(b)}{2})=\int_a^b f(x)dx+\sum_{k=1}^m\frac{B_{2k}h^{2k}}{(2k)!}(f^{(2k-1)}(b)-f^{(2k-1)}(a))+O(h^{2m+2})\)。故系数 \(c_{2k}=B_{2k}(f^{(2k-1)}(b)-f^{(2k-1)}(a))/(2k)!\)。由于只有偶次幂，Romberg 外推可行。\(\blacksquare\)

\textbf{定理 117.5}（Romberg 积分）：利用 Euler-Maclaurin 展开，通过 Richardson 外推消除低阶误差项：

\[R_{k,0} = T_{2^k}, \quad R_{k,m} = \frac{4^m R_{k,m-1} - R_{k-1,m-1}}{4^m - 1}\]

\(R_{n,n}\) 以 \(O(h^{2n+2})\) 的精度逼近积分。

\textbf{证明}：由 Euler-Maclaurin 展开 \(T_h=I+c_2h^2+c_4h^4+\cdots\)。取 \(h\) 和 \(h/2\) 的两近似： \[T_h=I+c_2h^2+O(h^4),\quad T_{h/2}=I+c_2h^2/4+O(h^4)\] 消去 \(c_2\)：\((4T_{h/2}-T_h)/3=I+O(h^4)\)。此即 Richardson 外推：\(R_{k,1}=(4^1 R_{k,0}-R_{k-1,0})/(4^1-1)\)。一般地，若 \(R_{k,m-1}=I+C_m h_k^{2m}+O(h_k^{2m+2})\)（\(h_k=h/2^k\)），则 \(R_{k,m}=(4^m R_{k,m-1}-R_{k-1,m-1})/(4^m-1)=I+O(h_k^{2m+2})\)。\(n\) 步后 \(R_{n,n}\) 消除到 \(h^{2n+2}\) 阶。\(\blacksquare\)

\textbf{定义 117.4}（自适应求积）：在函数变化剧烈处加密节点，变化平缓处稀疏节点。常用策略：比较两个不同阶的求积法则（如 Simpson 和 Boole）的差异，若差异超过阈值则细分区间。

\subsection{117.4 数值微分}\label{ux6570ux503cux5faeux5206}

\textbf{定义 117.5}（有限差分公式）： - 前向差分：\(f'(x) \approx \frac{f(x+h) - f(x)}{h}\)，误差 \(-\frac{h}{2} f''(\xi)\) - 中心差分：\(f'(x) \approx \frac{f(x+h) - f(x-h)}{2h}\)，误差 \(-\frac{h^2}{6} f'''(\xi)\) - 五点公式：\(f'(x) \approx \frac{-f(x+2h) + 8f(x+h) - 8f(x-h) + f(x-2h)}{12h}\)，误差 \(O(h^4)\)

\textbf{命题 117.1}（数值微分的不稳定性）：步长 \(h\) 不能无限减小。由于舍入误差 \(\sim \varepsilon/h\)（\(\varepsilon\) 是函数值误差），总误差 \(\approx C h^p + \varepsilon/h\)。最优 \(h \approx \varepsilon^{1/(p+1)}\)。

\hrulefill

\hrulefill

\hrulefill

\hrulefill

\hrulefill

\hrulefill

\section{第392章：常微分方程数值解}\label{ux7b2c392ux7ae0ux5e38ux5faeux5206ux65b9ux7a0bux6570ux503cux89e3}

本章讨论 ODE 初值问题的数值解法，包括 Euler 方法、Runge-Kutta 方法和多步法。

\subsection{118.1 Euler 方法}\label{euler-ux65b9ux6cd5}

\textbf{定义 118.1}（ODE 初值问题）：\(\dot{y} = f(t, y)\)，\(y(t_0) = y_0\)。

\textbf{定义 118.2}（前向 Euler 方法）：\(y_{n+1} = y_n + h f(t_n, y_n)\)（\(h\) 为步长）。\textbf{局部截断误差} \(O(h^2)\)，\textbf{全局误差} \(O(h)\)（一阶方法）。

\textbf{定义 118.3}（后向 Euler 方法）：\(y_{n+1} = y_n + h f(t_{n+1}, y_{n+1})\)（隐式方法，需解方程）。局部截断误差 \(O(h^2)\)，全局误差 \(O(h)\)（一阶方法）。

\textbf{定义 118.4}（梯形法则）：\(y_{n+1} = y_n + \frac{h}{2}[f(t_n, y_n) + f(t_{n+1}, y_{n+1})]\)（隐式，二阶方法）。

\subsection{118.2 Runge-Kutta 方法}\label{runge-kutta-ux65b9ux6cd5}

\textbf{定义 118.5}（显式 Runge-Kutta 方法，\(s\) 级）：\(s\) 级显式 RK 方法为

\[k_i = f\left(t_n + c_i h, y_n + h \sum_{j=1}^{i-1} a_{ij} k_j\right), \quad y_{n+1} = y_n + h \sum_{i=1}^s b_i k_i\]

\textbf{定理 118.1}（Butcher 障碍）：\(s\) 级显式 RK 方法的阶 \(p\) 满足 \(p \leq s\)（\(s = 1, 2, 3, 4\)），\(p \leq s-1\)（\(s = 5, 6, 7\)），\(p \leq s-2\)（\(s = 8, 9\)）等。

\textbf{证明}：Butcher 障碍由阶条件（order conditions）与级数 \(s\) 之间的组合约束导出。\(p\) 阶 RK 方法需满足的阶条件数量为：\(p=1\) 时 1 个，\(p=2\) 时 2 个，\(p=3\) 时 4 个，\(p=4\) 时 8 个，\(p=5\) 时 17 个，\(p=6\) 时 37 个，\(p=7\) 时 85 个，\(p=8\) 时 200 个，\(p=9\) 时 486 个，\(p=10\) 时 1205 个。\(s\) 级显式方法有 \(s(s+1)/2\) 个自由参数（Butcher 表的下三角元素）。求解非线性代数方程组 \(\Phi(\text{参数}) = \text{阶条件}\) 时，当阶条件数超过自由参数数时无解。这就是 Butcher 障碍的组合来源。\(\blacksquare\)

经典的 4 级 4 阶方法（RK4）是显式 RK 方法中最著名的：\(k_1 = f(t_n, y_n)\)，\(k_2 = f(t_n + h/2, y_n + h k_1/2)\)，\(k_3 = f(t_n + h/2, y_n + h k_2/2)\)，\(k_4 = f(t_n + h, y_n + h k_3)\)，\(y_{n+1} = y_n + \frac{h}{6}(k_1 + 2k_2 + 2k_3 + k_4)\)。RK4 因其在精度（4 阶）与计算量（每步 4 次函数求值）之间的良好平衡而被广泛使用。隐式 RK 方法（如 Gauss-Legendre 方法）可突破 Butcher 障碍：\(s\) 级隐式 RK 方法可达 \(2s\) 阶，这是因为隐式方法有 \(s^2\) 个自由参数（Butcher 表的全矩阵元素），大大增加了可调节的自由度。隐式 RK 方法对刚性方程（stiff equations）尤为重要------其 A-稳定性使步长选择不受方程刚性比的限制，而显式 RK 方法的稳定性域有限，对刚性方程效率极低。后向微分公式（BDF）是另一类处理刚性方程的重要方法。Runge-Kutta 方法的高阶推广------如 Runge-Kutta-Nystrom 方法（专为二阶 ODE 设计）和连续 Runge-Kutta 方法（提供步间插值）------在结构力学和自适应步长控制中有重要应用。

\subsection{118.3 多步法}\label{ux591aux6b65ux6cd5}

\textbf{定义 118.6}（线性多步法）：\(k\) 步线性多步法为

\[\sum_{j=0}^k \alpha_j y_{n+j} = h \sum_{j=0}^k \beta_j f(t_{n+j}, y_{n+j})\]

其中 \(\alpha_k = 1\)，\(|\alpha_0| + |\beta_0| > 0\)。若 \(\beta_k = 0\) 则为显式方法，否则为隐式方法。

\textbf{定义 118.7}（Adams 方法）： - \textbf{Adams-Bashforth}（显式）：\(\alpha_k = 1\)，\(\alpha_{k-1} = -1\)，其余 \(\alpha_j = 0\) - AB2：\(y_{n+1} = y_n + \frac{h}{2}[3f_n - f_{n-1}]\)（二阶） - AB4：\(y_{n+1} = y_n + \frac{h}{24}[55f_n - 59f_{n-1} + 37f_{n-2} - 9f_{n-3}]\)（四阶） - \textbf{Adams-Moulton}（隐式，通常作为校正步）：\(\beta_k \neq 0\) - AM4：\(y_{n+1} = y_n + \frac{h}{24}[9f_{n+1} + 19f_n - 5f_{n-1} + f_{n-2}]\)（四阶）

\subsection{118.4 稳定性与刚性方程}\label{ux7a33ux5b9aux6027ux4e0eux521aux6027ux65b9ux7a0b}

\textbf{定义 118.8}（绝对稳定性区域）：对检验方程 \(\dot{y} = \lambda y\)（\(\Re(\lambda) < 0\)），数值方法产生 \(y_{n+1} = R(z) y_n\)，其中 \(z = h\lambda\)。\textbf{绝对稳定性区域}为 \(\{z \in \mathbb{C} : |R(z)| < 1\}\)。

\textbf{命题 118.1}（各方法的稳定性区域）： - 前向 Euler：\(|1 + z| < 1\)（圆盘，半径 \(1\)，中心 \(-1\)） - 后向 Euler：\(|1 - z| > 1\)（整个左半平面为稳定区域，A-稳定） - 梯形法则：\(|(1+z/2)/(1-z/2)| < 1 \iff \Re(z) < 0\)（A-稳定） - RK4：与显式 Euler 类似但稍大（有限稳定区域）

\textbf{定义 118.9}（A-稳定性）：方法称为\textbf{A-稳定}的，如果其绝对稳定性区域包含整个左半平面 \(\{z : \Re(z) < 0\}\)。

\textbf{定义 118.10}（刚性方程）：若问题的不同分量以相差极大的速率变化（\(\max |\Re(\lambda)| / \min |\Re(\lambda)| \gg 1\)），则方程为\textbf{刚性}的。显式方法需极小步长，隐式方法（A-稳定）更适用。

\textbf{定义 118.11}（预测-校正方法）：\textbf{预测-校正}（Predictor-Corrector, PECE）方法：用显式方法预测 \(y_{n+1}^P\)，再用隐式方法校正 \(y_{n+1}\)。例：AB4 预测 + AM4 校正（均为四阶）。

\hrulefill

\hrulefill

\hrulefill

\hrulefill

\hrulefill

\hrulefill

\section{第393章：数值线性代数精讲}\label{ux7b2c393ux7ae0ux6570ux503cux7ebfux6027ux4ee3ux6570ux7cbeux8bb2}

数值线性代数是科学计算中最重要的核心算法集合，大多数大规模科学计算最终归结为线性方程组的求解或特征值问题。

\subsection{119.1 矩阵分解}\label{ux77e9ux9635ux5206ux89e3}

\textbf{定义 119.1}（LU 分解 / LU Decomposition）：将矩阵 \(A \in \mathbb{R}^{n \times n}\) 分解为下三角矩阵 \(L\) 与上三角矩阵 \(U\) 的乘积 \(A = LU\)。这是 Gauss 消去法的矩阵形式------消去过程的乘子按列存储于 \(L\) 中，消去结果存于 \(U\)。部分主元策略（partial pivoting）引入置换矩阵 \(P\) 得 \(PA = LU\)，避免小主元导致的数值不稳定。

\textbf{定理 119.1}（LU 分解的存在性）：若 \(A\) 的所有顺序主子式非零，则存在唯一的单位下三角矩阵 \(L\) 和上三角矩阵 \(U\) 使得 \(A = LU\)。

\textbf{证明}：对 \(n\) 归纳。\(n=1\) 时，\(A = [a_{11}]\)，取 \(L = [1]\)，\(U = [a_{11}]\)（\(a_{11} \neq 0\) 由顺序主子式非零保证）。设对 \(n-1\) 阶矩阵成立。将 \(A\) 分块：

\[A = \begin{pmatrix} A_{n-1} & \mathbf{b} \\ \mathbf{c}^T & a_{nn} \end{pmatrix}\]

由归纳假设，\(A_{n-1} = L_{n-1} U_{n-1}\)。令 \(L = \begin{pmatrix} L_{n-1} & 0 \\ \mathbf{x}^T & 1 \end{pmatrix}\)，\(U = \begin{pmatrix} U_{n-1} & \mathbf{y} \\ 0 & u_{nn} \end{pmatrix}\)。由 \(A = LU\) 得 \(\mathbf{y} = L_{n-1}^{-1} \mathbf{b}\)，\(\mathbf{x}^T = \mathbf{c}^T U_{n-1}^{-1}\)，\(u_{nn} = a_{nn} - \mathbf{x}^T \mathbf{y}\)。由 \(\det(A) = \det(A_{n-1}) \cdot u_{nn} \neq 0\) 知 \(u_{nn} \neq 0\)。唯一性：若 \(L_1 U_1 = L_2 U_2\)，则 \(L_2^{-1} L_1 = U_2 U_1^{-1}\)。左边为单位下三角，右边为上三角，故二者均为单位矩阵，即 \(L_1 = L_2\)，\(U_1 = U_2\)。\(\blacksquare\)

\textbf{定义 119.2}（Cholesky 分解）：若 \(A \in \mathbb{R}^{n \times n}\) 对称正定，则存在唯一的下三角矩阵 \(L\)（主对角元为正）使得 \(A = LL^T\)。运算量约为 \(\frac{1}{3}n^3\)（LU 分解的一半），且数值稳定性优良。

\textbf{定义 119.3}（QR 分解 / QR Decomposition）：将 \(A \in \mathbb{R}^{m \times n}\)（\(m \geq n\)）分解为 \(A = QR\)，其中 \(Q \in \mathbb{R}^{m \times m}\) 是正交矩阵（\(Q^T Q = I\)），\(R \in \mathbb{R}^{m \times n}\) 是上三角矩阵。构造方法包括：

\begin{itemize}
\tightlist
\item
  \textbf{Householder 反射}：\(H = I - 2 v v^T\)（\(v\) 为单位向量）。通过 \(n-1\) 次反射将 \(A\) 的列依次消去对角线下方元素
\item
  \textbf{Givens 旋转}：\(G(i, j, \theta)\) 在 \((i,j)\) 平面上旋转以消去单个元素，适用于稀疏矩阵和并行计算
\end{itemize}

\textbf{定理 119.2}（QR 分解求解最小二乘问题）：对于超定方程组 \(\min_x \|Ax - b\|_2\)（\(m \geq n\)），令 \(A = QR\)，\(Q^T b = \begin{pmatrix} c \\ d \end{pmatrix}\)（\(c \in \mathbb{R}^n\)），则解为 \(x = R^{-1} c\)，残差范数为 \(\|d\|_2\)。

\textbf{证明}：\(A=Q\begin{pmatrix}R\\0\end{pmatrix}\)，\(\|Ax-b\|_2=\|Q\begin{pmatrix}R\\0\end{pmatrix}x-b\|_2=\|\begin{pmatrix}R\\0\end{pmatrix}x-Q^T b\|_2\)（\(Q\) 正交保范）。令 \(Q^T b=\begin{pmatrix}\mathbf{c}\\\mathbf{d}\end{pmatrix}\)，则 \(\|Ax-b\|_2^2=\|R\mathbf{x}-\mathbf{c}\|_2^2+\|\mathbf{d}\|_2^2\)。因 \(R\) 可逆（\(A\) 满列秩 \(\implies R\) 可逆），第一项最小值在 \(\mathbf{x}=R^{-1}\mathbf{c}\) 处取零，故最小残差 \(\|\mathbf{d}\|_2\)。\(\blacksquare\)

\subsection{119.2 特征值算法}\label{ux7279ux5f81ux503cux7b97ux6cd5}

\textbf{定义 119.4}（QR 算法 / Francis, 1961-1962）：对 \(A_0 = A\)，迭代执行： 1. 对 \(A_k - \mu_k I\) 进行 QR 分解：\(A_k - \mu_k I = Q_k R_k\) 2. 重构：\(A_{k+1} = R_k Q_k + \mu_k I\)

适当地选择位移 \(\mu_k\)（Wilkinson 位移、Rayleigh 商位移），\(A_k\) 收敛到实 Schur 形式（拟上三角形，对角块为 \(1\times 1\) 或 \(2\times 2\)），特征值出现在对角线上。

\textbf{定义 119.5}（Lanczos 算法 / Lanczos Iteration）：对对称矩阵 \(A \in \mathbb{R}^{n \times n}\)，Lanczos 算法构造 Krylov 子空间 \(\mathcal{K}_k(A, v) = \operatorname{span}\{v, Av, A^2 v, \ldots, A^{k-1} v\}\) 的一组正交基 \(\{q_1, \ldots, q_k\}\)，满足三对角递推 \(A Q_k = Q_k T_k + \beta_k q_{k+1} e_k^T\)，其中 \(T_k\) 为对称三对角矩阵。\(T_k\) 的极值特征值在 \(k \ll n\) 时即收敛到 \(A\) 的极值特征值。

\textbf{定义 119.6}（Divide-and-Conquer 特征值算法）：将对称三对角矩阵分块为 \(T = \begin{pmatrix} T_1 & 0 \\ 0 & T_2 \end{pmatrix} + \rho v v^T\)，递归计算 \(T_1\)、\(T_2\) 的特征分解，通过求解 secular 方程 \(f(\lambda) = 1 + \rho \sum_{j=1}^n \frac{v_j^2}{d_j - \lambda} = 0\) 合并。总复杂度 \(O(n^3)\) 但常数小，LAPACK 中为默认选择。

\subsection{119.3 Krylov 子空间迭代法}\label{krylov-ux5b50ux7a7aux95f4ux8fedux4ee3ux6cd5}

\textbf{定义 119.7}（共轭梯度法 / CG / Conjugate Gradient）：对对称正定的 \(Ax = b\)，共轭梯度法沿着 \(A\)-正交（\(p_i^T A p_j = 0\)）的搜索方向迭代：

\[\begin{aligned} r_0 &= b - A x_0, \quad p_0 = r_0 \\ \alpha_k &= \frac{r_k^T r_k}{p_k^T A p_k} \\ x_{k+1} &= x_k + \alpha_k p_k \\ r_{k+1} &= r_k - \alpha_k A p_k \\ \beta_k &= \frac{r_{k+1}^T r_{k+1}}{r_k^T r_k} \\ p_{k+1} &= r_{k+1} + \beta_k p_k \end{aligned}\]

\textbf{定理 119.3}（CG 收敛性）：在精确算术下，CG 在至多 \(n\) 步给出精确解。第 \(k\) 步的误差满足

\[\|x_k - x_*\|_A \leq 2 \left( \frac{\sqrt{\kappa(A)} - 1}{\sqrt{\kappa(A)} + 1} \right)^k \|x_0 - x_*\|_A\]

\textbf{证明}：CG 在 \(\mathcal{K}_k\) 中极小化 \(\|x-x_*\|_A\)，第 \(n\) 步 \(\mathcal{K}_n=\mathbb{R}^n\) 必得精确解。误差界：\(\|x_k-x_*\|_A=\min_{p\in\mathcal{P}_k,p(0)=1}\|p(A)(x_0-x_*)\|_A\)。在特征基下 \(A\) 对角化，问题化为 \(\min_{p\in\mathcal{P}_k,p(0)=1}\max_{\lambda\in[\lambda_{\min},\lambda_{\max}]}|p(\lambda)|\)。极小化多项式的最优选择是缩放平移的 Chebyshev 多项式：\(p^*_k(\lambda)=T_k(\frac{\lambda_{\max}+\lambda_{\min}-2\lambda}{\lambda_{\max}-\lambda_{\min}})/T_k(\frac{\lambda_{\max}+\lambda_{\min}}{\lambda_{\max}-\lambda_{\min}})\)。对 \(\lambda\in[\lambda_{\min},\lambda_{\max}]\)，\(|p^*_k(\lambda)|\le 1/T_k(\frac{\kappa+1}{\kappa-1})\)。且 \(T_k(x)\ge\frac{1}{2}(x+\sqrt{x^2-1})^k\) 当 \(|x|>1\)，故 \(|p^*_k(\lambda)|\le 2(\frac{\sqrt{\kappa}-1}{\sqrt{\kappa}+1})^k\)，即得递减因子。\(\blacksquare\)

其中 \(\kappa(A) = \|A\|_2 \|A^{-1}\|_2 = \lambda_{\max} / \lambda_{\min}\) 为条件数，\(\|x\|_A = \sqrt{x^T A x}\)。

\textbf{定义 119.8}（预条件 / Preconditioning）：引入预条件子 \(M \approx A\)（\(M\) 易于求逆），求解等价系统 \(M^{-1} A x = M^{-1} b\)。预条件目的是降低等价系统的条件数。常用预条件子包括 Jacobi（对角）、不完全 Cholesky、多重网格。预条件共轭梯度法（PCG）是实际大规模求解的主流方法。

\textbf{定义 119.9}（GMRES / Generalized Minimal Residual）：对非对称 \(Ax = b\)，GMRES 在 Krylov 子空间 \(\mathcal{K}_k\) 中寻找使残差 \(\|b - A x_k\|_2\) 最小的 \(\tilde{x}_k\)。通过 Arnoldi 迭代构造 \(\mathcal{K}_k\) 的正交基，将最小二乘问题化为小型 Hessenberg 矩阵的最小二乘问题。存储需求随 \(k\) 线性增长，实际使用重启（restarted）GMRES(\(m\))。

\hrulefill

\hrulefill

\hrulefill

\hrulefill

\hrulefill

\hrulefill

\newpage

\chapter{11-数值与计算/01-数值分析/04-逼近论.md}\label{ux6570ux503cux4e0eux8ba1ux7b9701-ux6570ux503cux5206ux679004-ux903cux8fd1ux8bba.md}

\subsection{321.A 逼近论补充}\label{a-ux903cux8fd1ux8bbaux8865ux5145}

\begin{quote}
逼近论研究用简单函数（多项式、有理函数、样条、小波等）逼近复杂函数的理论和方法。从 Weierstrass 逼近定理到 Chebyshev 最佳逼近，从有理逼近与 Pade 逼近到样条与小波，逼近论为数值计算、信号处理和机器学习提供了坚实的数学基础。本卷系统建立最佳逼近的存在性与唯一性理论、多项式逼近的 Jackson 与 Bernstein 定理、有理逼近、样条插值与光滑性、小波的多分辨分析框架，以及径向基函数与高维逼近方法。
\end{quote}

\begin{quote}
\textbf{卷序说明}：本章节号（Ch 277-211）反映其在全书逻辑链中的位置。物理卷序上本卷位于V52，作为逼近论方向的专题补充卷。
\end{quote}

\hrulefill

\hrulefill

\hrulefill

\hrulefill

\hrulefill

\hrulefill

\section{第394章：Weierstrass 逼近定理与 Jackson 定理}\label{ux7b2c394ux7ae0weierstrass-ux903cux8fd1ux5b9aux7406ux4e0e-jackson-ux5b9aux7406}

Weierstrass 逼近定理（1885）是逼近论的奠基性结果------它断言闭区间上的任何连续函数均可被多项式一致逼近。Bernstein（1912）给出了一个构造性的概率证明，开创了用 Bernstein 多项式逼近连续函数的先河。Jackson 定理（1911）进一步建立了光滑性与逼近阶之间的定量关系。

\subsection{208.1 Weierstrass 第一逼近定理}\label{weierstrass-ux7b2cux4e00ux903cux8fd1ux5b9aux7406}

\textbf{定理 208.1}（Weierstrass 第一逼近定理，1885）：设 \(f \in C[a, b]\)。则对任意 \(\varepsilon > 0\)，存在多项式 \(p(x)\) 使得 \[\max_{x \in [a, b]} |f(x) - p(x)| < \varepsilon\] 即多项式空间 \(\mathcal{P}\) 在 \(C[a, b]\)（一致范数拓扑）中稠密。

\textbf{证明}：先考虑 \([a,b]=[0,1]\)。用 Bernstein 多项式 \(B_n(f;x)=\sum_{k=0}^n f(k/n)\binom{n}{k}x^k(1-x)^{n-k}\) 逼近 \(f\)。由概率解释 \(B_n(f;x)=\mathbb{E}[f(S_n/n)]\)（\(S_n\sim\operatorname{Bin}(n,x)\)）， \[|f(x)-B_n(f;x)|\le\mathbb{E}|f(x)-f(S_n/n)|\le\varepsilon+2\|f\|_\infty\cdot\mathbb{P}(|S_n/n-x|\ge\delta)\] 由 Chebyshev 不等式 \(\mathbb{P}(|S_n/n-x|\ge\delta)\le x(1-x)/(n\delta^2)\le1/(4n\delta^2)\)。取 \(n>\|f\|_\infty/(2\varepsilon\delta^2)\) 即得 \(\|f-B_n(f)\|_\infty<2\varepsilon\)。对一般 \([a,b]\) 通过线性变换 \(x\mapsto a+(b-a)x\) 化归。\(\blacksquare\)

\subsection{208.2 Bernstein 多项式与构造性证明}\label{bernstein-ux591aux9879ux5f0fux4e0eux6784ux9020ux6027ux8bc1ux660e}

\textbf{定义 208.1}（Bernstein 多项式，1912）：设 \(f \in C[0, 1]\)。\(f\) 的 \(n\) 次 \textbf{Bernstein 多项式} 定义为 \[B_n(f; x) = \sum_{k=0}^n f\left(\frac{k}{n}\right) \binom{n}{k} x^k (1-x)^{n-k}, \quad x \in [0, 1]\]

注意 \(B_n(f; x)\) 恰为二项分布 \(\operatorname{Bin}(n, x)\) 下 \(f(k/n)\) 的数学期望 \(\mathbb{E}_x[f(S_n / n)]\)，其中 \(S_n \sim \operatorname{Bin}(n, x)\)。这赋予 Bernstein 多项式自然的概率解释。

\textbf{定理 208.2}（Bernstein 多项式的收敛性）：对任意 \(f \in C[0, 1]\)， \[\lim_{n \to \infty} \max_{x \in [0, 1]} |f(x) - B_n(f; x)| = 0\]

\emph{证明（基于概率论）}：由一致连续性，对任意 \(\varepsilon > 0\)，存在 \(\delta > 0\) 使得 \(|x - y| < \delta \implies |f(x) - f(y)| < \varepsilon\)。记 \(M = \|f\|_\infty\)。则 \[\begin{aligned}
|f(x) - B_n(f; x)| &= \left|\mathbb{E}_x\left[f(x) - f\left(\frac{S_n}{n}\right)\right]\right| \\
&\leq \mathbb{E}_x\left|f(x) - f\left(\frac{S_n}{n}\right)\right| \cdot \left[\mathbf{1}_{\{|S_n/n - x| < \delta\}} + \mathbf{1}_{\{|S_n/n - x| \geq \delta\}}\right] \\
&\leq \varepsilon + 2M \cdot \mathbb{P}_x\left(\left|\frac{S_n}{n} - x\right| \geq \delta\right)
\end{aligned}\]

由 Chebyshev 不等式，\(\mathbb{P}_x(|S_n/n - x| \geq \delta) \leq \frac{\operatorname{Var}(S_n/n)}{\delta^2} = \frac{x(1-x)}{n\delta^2} \leq \frac{1}{4n\delta^2}\)。因此 \[\max_{x \in [0, 1]} |f(x) - B_n(f; x)| \leq \varepsilon + \frac{M}{2n\delta^2}\]

取 \(n > M / (2\varepsilon \delta^2)\) 即得 \(\|f - B_n(f)\|_\infty < 2\varepsilon\)。∎

\textbf{注 208.1}（Bernstein 多项式的局限）：尽管 Bernstein 多项式给出了一致收敛的构造性证明，但其收敛速度很慢------对被逼近函数 \(f\) 即使非常光滑（如 \(f(x) = x^2\)），误差也仅为 \(O(1/n)\)（而非 \(O(1/n^2)\) 或更快）。这促使后人寻求更高效的逼近方法（如 Chebyshev 多项式插值、Jackson 定理等）。

\textbf{推论 208.1}（推广到一般区间）：对 \(f \in C[a, b]\)，通过线性变换 \(t = (x-a)/(b-a)\) 将 \([a, b]\) 映射到 \([0, 1]\)，即得 \([a, b]\) 上的 Weierstrass 逼近定理。

\subsection{208.3 Jackson 定理------直接定理}\label{jackson-ux5b9aux7406ux76f4ux63a5ux5b9aux7406}

Weierstrass 定理回答了''是否存在多项式逼近''的问题，Jackson 定理进一步回答了''误差随多项式次数如何衰减''的定量问题。

\textbf{定理 208.3}（Jackson 定理 / 直接定理，Jackson 1911）：设 \(f \in C^r[-1, 1]\)。则对任意正整数 \(n \geq r\)，存在次数不超过 \(n\) 的多项式 \(p_n\) 使得 \[\|f - p_n\|_\infty \leq \frac{C_r}{n^r} \cdot \omega\left(f^{(r)}; \frac{1}{n}\right)\] 其中 \(\omega(g; \delta) = \sup_{|x-y| \leq \delta} |g(x) - g(y)|\) 为连续模（modulus of continuity），\(C_r\) 为仅依赖于 \(r\) 的常数。

特别地，若 \(f^{(r)} \in \operatorname{Lip}_\alpha\)（\(\alpha\)-Holder 连续，\(0 < \alpha \leq 1\)），则 \[\|f - p_n\|_\infty \leq \frac{C_{r, \alpha}}{n^{r+\alpha}}\]

\textbf{证明}：核心工具是 Jackson 核 \(J_n(x)=c_n(\sin(nx/2)/\sin(x/2))^{2m}\)（\(m>r\)），满足 \(\int_{-\pi}^\pi J_n=1\) 且 \(J_n\) 的矩 \(\int_{-\pi}^\pi t^k J_n(t)dt=0\)（\(k=1,\dots,2m-2\)）。令 \(p_n(x)=\int_{-\pi}^\pi f(\cos(t+\arccos x))J_n(t+\arccos x)dt\)（经过变量变换后），则 \(p_n\) 为次数 \(\le n\) 的三角多项式。由矩性质和 Taylor 展开，\(\|f-p_n\|_\infty\le C n^{-r}\omega(f^{(r)};1/n)\)。常数 \(C_r\) 由 Jackson 核的 \(L^1\) 矩控制。\(\blacksquare\)

\textbf{定理 208.4}（Jackson 定理的简化形式）：设 \(f \in C^k[-1, 1]\)。则存在仅依赖于 \(k\) 的常数 \(C_k\) 使得对任意 \(n \geq k\)， \[E_n(f) := \inf_{p \in \mathcal{P}_n} \|f - p\|_\infty \leq \frac{C_k}{n^k} \cdot \omega\left(f^{(k)}; \frac{1}{n}\right)\]

\textbf{证明}：这是定理 208.3 的直接推论：取 \(r=k\)，由最佳逼近的最小性 \(E_n(f)\le\|f-p_n\|_\infty\)，定理 208.3 保证存在 \(p_n\in\mathcal{P}_n\) 使右端成立。\(\blacksquare\)

\textbf{推论 208.2}（光滑性 \(\Rightarrow\) 逼近阶）：若 \(f \in C^r[-1, 1]\) 且 \(f^{(r)}\) 满足 Lipschitz 条件，则 \(E_n(f) = O(n^{-r-1})\)。光滑性越高，多项式逼近的代数收敛阶越高。

\textbf{评注 208.1}（Jackson 定理的意义）：Jackson 定理精确刻画了''函数的可微性阶数 \(r\) 决定多项式最佳逼近的代数收敛阶 \(\geq r+1\)``。这是逼近论中\textbf{直接定理}（direct theorem）的核心------从函数的光滑性信息推断最佳逼近的收敛速率。反方向的\textbf{逆定理}（Bernstein / Stechkin 逆定理）则从逼近速率反推函数的光滑性，两者共同构成完整的逼近论等价关系。

\hrulefill

\hrulefill

\hrulefill

\hrulefill

\hrulefill

\hrulefill

\section{第395章：有理逼近与 Padé 逼近}\label{ux7b2c395ux7ae0ux6709ux7406ux903cux8fd1ux4e0e-paduxe9-ux903cux8fd1}

有理函数逼近是逼近论中一个精妙且高度实用的分支，它在多项式逼近的基础上引入了一个全新的自由度------分母。这一自由度使得有理函数能够自然地捕捉被逼近函数的极点（singularities），从而在函数具有极点或奇异性时实现远远优于多项式的逼近效果。Newman 定理（1964）是这一优势的惊人例证：\(|x|\) 在 \([-1,1]\) 上的多项式最佳逼近仅以 \(O(1/n)\) 速率收敛（因为 \(|x|\) 在 \(0\) 处不可导），但 \([n/n]\) 型有理函数逼近却以 \(3e^{-\sqrt{n}}\) 速率指数收敛。Padé 逼近（1892）是有理逼近的代数构造方法：给定函数 \(f(z)=\sum_{k=0}^\infty c_k z^k\) 的 Taylor 展开，\([m/n]\) 型 Padé 逼近 \(r_{m,n}(z)=p_m(z)/q_n(z)\) 通过匹配前 \(m+n+1\) 个 Taylor 系数确定分子和分母多项式的系数------这等价于求解一个 Hankel 线性方程组（分母系数）和一个前向代换（分子系数）。Padé 表的对角线序列 \([n/n]\) 在容量意义下收敛到原函数（Nuttall-Pommerenke 定理），且 Padé 分母的零点自然地''自发''逼近函数的极点。本章系统地建立有理最佳逼近的存在性与唯一性理论（Chebyshev 交替定理的有理版本）、Padé 逼近的代数构造和连分式联系，以及 Padé 逼近的收敛性理论。

\subsection{209.1 有理逼近的基本理论}\label{ux6709ux7406ux903cux8fd1ux7684ux57faux672cux7406ux8bba}

\textbf{定义 209.1}（有理函数）：\([m/n]\) 型\textbf{有理函数}为 \(r(x) = p_m(x) / q_n(x)\)，其中 \(p_m \in \mathcal{P}_m\)，\(q_n \in \mathcal{P}_n\)（\(q_n \not\equiv 0\)）。

\textbf{定义 209.2}（有理最佳逼近）：\(f \in C[a,b]\) 在 \(\mathcal{R}_{m,n}\)（\([m/n]\) 型有理函数全体）中的\textbf{最佳一致有理逼近} \(r^*\) 满足

\[\|f - r^*\|_\infty = \inf_{r \in \mathcal{R}_{m,n}} \|f - r\|_\infty\]

\textbf{定理 209.1}（有理逼近的 Chebyshev 交替定理）：\(r^* \in \mathcal{R}_{m,n}\) 是 \(f\) 的最佳一致有理逼近当且仅当 \(f - r^*\) 在 \([a,b]\) 上至少存在 \(m+n+2 - d\) 个点的交替点集，其中 \(d = \min(m - \deg p^*, n - \deg q^*)\)。

\textbf{证明}：必要性：若无足够交替点，则可构造扰动 \(h\in\mathcal{R}_{m,n}\) 使 \(\|f-(r^*+\epsilon h)\|_\infty<\|f-r^*\|_\infty\)，与最佳性矛盾。充分性：由 Haar 条件------\(\mathcal{R}_{m,n}\) 在非退化情形下构成 Chebyshev 系统，交替点集的存在性蕴含最佳逼近的唯一性。设交替点 \(x_1<\cdots<x_N\)（\(N=m+n+2-d\)），误差 \(e(x_k)=\pm(-1)^k E\)。对任意 \(r\in\mathcal{R}_{m,n}\)，若 \(\|f-r\|_\infty<E\)，则 \(r-r^*\) 在交替点处与 \(f-r^*\) 同号，由零点计数 \(r-r^*\) 有 \(N-1=m+n+1-d\) 个零点，而 \(\dim\mathcal{R}_{m,n}=m+n+1-d\)，故 \(r=r^*\)。\(\blacksquare\)

\textbf{定理 209.2}（Newman 定理，1964）：\(|x|\) 在 \([-1, 1]\) 上的有理逼近收敛速率远快于多项式逼近：

\[\min_{r \in \mathcal{R}_{n,n}} \||x| - r(x)\|_\infty \leq 3 e^{-\sqrt{n}}\]

而多项式逼近的速率仅为 \(O(1/n)\)。这展示了有理逼近在处理非光滑函数时的巨大优势。

\textbf{证明}：Newman 构造了一族有理函数 \(r_n(x)=x(p(x)-p(-x))/(p(x)+p(-x))\)，其中 \(p(x)=\prod_{k=1}^n(1+x^{2k-1}e^{-\sqrt{n}})\)。通过估计 \(r_n(x)\) 与 \(|x|\) 的差值并利用 Chebyshev 多项式的最优性质，可证 \(\||x|-r_n(x)\|_\infty\le 3e^{-\sqrt{n}}\)。关键是 \(p(x)/p(-x)\) 在 \([-1,1]\) 上逼近 \(\operatorname{sgn}(x)\)，从而 \(x\cdot(p(x)-p(-x))/(p(x)+p(-x))\) 逼近 \(|x|\)。\(lacksquare\)

\subsection{209.2 Padé 逼近}\label{paduxe9-ux903cux8fd1}

\textbf{定义 209.3}（Padé 逼近，1892）：设 \(f(z) = \sum_{k=0}^\infty c_k z^k\) 在 \(0\) 处解析。\textbf{\([m/n]\) 型 Padé 逼近} \(r_{m,n}(z) = p_m(z) / q_n(z)\)（\(q_n(0)=1\)）满足

\[f(z) - r_{m,n}(z) = O(z^{m+n+1}) \quad (z \to 0)\]

即 Padé 逼近匹配 \(f\) 在 \(0\) 处的 Taylor 系数的前 \(m+n+1\) 项。

\textbf{定理 209.3}（Padé 逼近的构造）：Padé 逼近的系数通过求解线性方程组确定：

\[c_0 q_0 = p_0\] \[c_1 q_0 + c_0 q_1 = p_1\] \[\vdots\]

其中 \(q_j = 0\)（\(j > n\)），\(p_j = 0\)（\(j > m\)）。\(q_j\) 的确定需要求解 Hankel 矩阵方程组。

\textbf{证明}：设 \(f(z)=\sum_{k=0}^\infty c_k z^k\)，\(r_{m,n}=p_m/q_n\)，\(q_n(0)=1\)。条件 \(f(z)-r_{m,n}(z)=O(z^{m+n+1})\) 展开为 \(\sum_{k=0}^\infty c_k z^k \cdot \sum_{j=0}^n q_j z^j - \sum_{i=0}^m p_i z^i = O(z^{m+n+1})\)。比较 \(z^k\) 系数（\(k=0,\ldots,m+n\)）得 \(m+n+1\) 个线性方程： 对 \(k=0,\ldots,m\)：\(\sum_{j=0}^{\min(k,n)} c_{k-j} q_j = p_k\)； 对 \(k=m+1,\ldots,m+n\)：\(\sum_{j=0}^{\min(k,n)} c_{k-j} q_j = 0\)。 后 \(n\) 个方程为 \(q_1,\ldots,q_n\) 的 Hankel 线性方程组 \(\mathbf{H}_n \mathbf{q} = -\mathbf{c}'\)（\(\mathbf{H}_n\) 为 \(n\times n\) Hankel 矩阵，\(H_{ij}=c_{m+i-j+1}\)）。求解后代入前 \(m+1\) 个方程得 \(p_i\)。\(lacksquare\)

\textbf{定义 209.4}（Padé 表）：\([m/n]\) Padé 逼近的二维阵列称为 \textbf{Padé 表}。Padé 表的对角线 \([n/n]\)（对角线 Padé 逼近）和次对角线 \([n-1/n]\) 具有特别的重要性。

\subsection{209.3 Padé 逼近的收敛性}\label{paduxe9-ux903cux8fd1ux7684ux6536ux655bux6027}

\textbf{定理 209.4}（Padé 逼近的收敛性，Nuttall-Pommerenke 定理）：对大部分亚纯函数，对角线 Padé 序列 \([n/n]\) 在容量意义下收敛到 \(f\)（在极点以外的区域）。Padé 逼近自然地捕捉了函数的极点。

\textbf{证明}：设 \(f\) 在 \(\mathbb{C}\setminus E\) 上亚纯（\(E\) 为零容量集）。对角线 Padé 序列 \([n/n]_f\) 的极点依极限分布收敛到 \(f\) 的奇异性集的最小容量支撑（Stahl 紧集）。由 Nuttall-Pommerenke 的奇异性理论，\([n/n]_f\) 在紧子集上一致收敛到 \(f\)（极点除外）。核心工具为 Padé 分母的正交性和 Riemann-Hilbert 问题的渐近分析。\(lacksquare\)

\textbf{定理 209.5}（Stahl 定理，1986）：对具有分支切割的代数函数，对角线 Padé 逼近的极限在容量意义下存在，且极点的极限分布由 \(f\) 的奇异性的最小容量集合给出。

\textbf{证明}：设 \(f\) 为代数函数（满足 \(P(z,f(z))=0\)，\(P\) 为多项式），具有分支切割。Stahl 证明了 Padé 极点的极限测度 \(\lambda\) 是唯一的最小化 \(I(\mu)=\iint\log|z-w|^{-1}d\mu(z)d\mu(w)\) 的测度，支撑在 \(f\) 的奇异性集 \(S\) 上，且满足 \(\partial g_\lambda/\partial n_+ = \partial g_\lambda/\partial n_-\)（对称平衡条件）。由位势理论，该对称条件等价于 \(S\) 是紧集 \(S_0\) 中使 \(\operatorname{cap}(S_0)\) 最小化的集合。\(lacksquare\)

\textbf{定义 209.5}（Padé 逼近与连分式）：Padé 逼近与连分式展开有自然联系。\([n/n]\) 和 \([n-1/n]\) Padé 逼近是函数连分式展开的截断。

\hrulefill

\hrulefill

\hrulefill

\hrulefill

\hrulefill

\hrulefill

\section{第396章：样条与小波}\label{ux7b2c396ux7ae0ux6837ux6761ux4e0eux5c0fux6ce2}

样条函数与小波分析是过去半世纪中函数逼近和信号处理领域最具影响力的两大数学工具。样条（Spline）是分段多项式函数，在相邻子区间之间满足一定的连续性条件，其数学理论由 Schoenberg 在 1940 年代奠基。样条的核心优势在于其双重的局部性：每个 B-样条基函数具有紧支撑（仅在有限个节点区间上非零），且样条同时具有极小弯曲能量的最优性------在所有满足相同插值条件且二阶导数平方可积的函数中，自然三次样条极小化能量泛函 \(\int_a^b|s''(x)|^2dx\)。Cox-de Boor 递推公式为 B-样条的数值计算提供了稳定高效的算法。小波（Wavelet）是样条思想在时间-频率平面上的革命性推广：Mallat 的多分辨分析（MRA，1989）通过层层嵌套的子空间序列 \(\cdots\subset V_{-1}\subset V_0\subset V_1\subset\cdots\) 将信号分解为不同分辨率下的近似成分和细节成分，母小波 \(\psi\) 的伸缩平移族 \(\{\psi_{j,k}(x)=2^{j/2}\psi(2^jx-k)\}\) 构成了 \(L^2(\mathbb{R})\) 的规范正交基。Daubechies 小波（1988）实现了紧支撑与任意高阶消失矩的完美结合，其构造方法基于二尺度方程和正交滤波器组的谱分解。小波的消失矩性质隐含着非线性逼近的最优性------对于分段光滑函数，小波系数的衰减速率直接由消失矩阶数和函数的光滑性共同决定。

\subsection{210.1 样条函数}\label{ux6837ux6761ux51fdux6570}

\textbf{定义 210.1}（\(k\) 次样条）：给定节点 \(a = x_0 < x_1 < \cdots < x_N = b\)，\textbf{\(k\) 次样条函数} \(s(x)\) 满足： 1. \(s\) 在每个子区间 \([x_i, x_{i+1}]\) 上是次数 \(\leq k\) 的多项式 2. \(s \in C^{k-1}[a,b]\)（\(k-1\) 阶连续可微）

\textbf{定义 210.2}（自然三次样条）：\textbf{自然三次样条}（\(k=3\)）是满足端点条件 \(s''(a) = s''(b) = 0\) 的三次样条。

\textbf{定理 210.1}（样条插值的最优性质）：在所有满足插值条件 \(s(x_i) = f(x_i)\) 且 \(s'' \in L^2[a,b]\) 的函数中，自然三次样条插值极小化弯曲能量

\[E(s) = \int_a^b |s''(x)|^2 \, dx\]

\textbf{证明}：设 \(s\) 为自然三次样条插值，\(g\) 为任意 \(C^2[a,b]\) 函数满足 \(g(x_i)=f(x_i)\)。令 \(h=g-s\)，\(h(x_i)=0\)。分部积分： \[\int_a^b s''(x)h''(x)dx = [s''(x)h'(x)]_a^b - \int_a^b s'''(x)h'(x)dx\] 因 \(s'''\) 在子区间上为常数且 \(h(x_i)=0\)，积分为零。又 \(s''(a)=s''(b)=0\)，故 \(\int_a^b s''h'' = 0\)。于是 \[\int_a^b |g''|^2 = \int_a^b |s''+h''|^2 = \int_a^b|s''|^2 + \int_a^b|h''|^2 \ge \int_a^b|s''|^2\] 等号成立当且仅当 \(h''\equiv0\)，即 \(h\) 线性，结合 \(h(x_i)=0\) 得 \(h\equiv0\)。故 \(g=s\) 时能量最小。\(lacksquare\)

\textbf{定理 210.2}（样条逼近的误差估计）：对 \(f \in C^4[a,b]\)，三次样条插值 \(s\) 满足

\[\|f - s\|_\infty \leq \frac{5}{384} \|f^{(4)}\|_\infty h^4, \quad \|f' - s'\|_\infty \leq \frac{1}{24} \|f^{(4)}\|_\infty h^3\]

其中 \(h = \max_i (x_{i+1} - x_i)\)。

\textbf{证明}：设 \(s\) 为 \(f\) 的三次样条插值。在每个子区间 \([x_i,x_{i+1}]\) 上，\(f-s\) 在 \(x_i\) 和 \(x_{i+1}\) 处为零。由三次 Hermite 插值的误差公式， \[|f(x)-s(x)|\le\frac{1}{384}h_i^4\|f^{(4)}\|_{L^\infty[x_i,x_{i+1}]}\] 其中 \(h_i=x_{i+1}-x_i\)。对导数 \(f'-s'\)，误差为 \(h_i^3\|f^{(4)}\|_\infty/24\)。取 \(h=\max_i h_i\)，全局误差为上述常数乘以 \(h^4\) 和 \(h^3\)。常数 \(5/384\) 来自对分段误差的最坏情况估计和边界条件的优化。\(lacksquare\)

\subsection{210.2 B-样条基}\label{b-ux6837ux6761ux57fa}

\textbf{定义 210.3}（B-样条，Cox-de Boor 递推，1972）：\textbf{\(k\) 次 B-样条} \(B_{i,k}(x)\) 由递推定义：

\[B_{i,0}(x) = \begin{cases} 1 & x \in [x_i, x_{i+1}) \\ 0 & \text{否则} \end{cases}\]

\[B_{i,k}(x) = \frac{x - x_i}{x_{i+k} - x_i} B_{i,k-1}(x) + \frac{x_{i+k+1} - x}{x_{i+k+1} - x_{i+1}} B_{i+1,k-1}(x)\]

\textbf{命题 210.1}（B-样条的性质）： - 局部支撑：\(B_{i,k}(x)\) 仅在 \([x_i, x_{i+k+1}]\) 上非零 - 非负性：\(B_{i,k}(x) \geq 0\) - 单位分解：\(\sum_i B_{i,k}(x) = 1\)（对所有 \(x\)） - 线性无关：\(\{B_{i,k}\}\) 是 \(k\) 次样条空间的基

\textbf{定理 210.3}（B-样条的导数）：\(B_{i,k}'(x) = \frac{k}{x_{i+k} - x_i} B_{i,k-1}(x) - \frac{k}{x_{i+k+1} - x_{i+1}} B_{i+1,k-1}(x)\)。

\textbf{证明}：对 Cox-de Boor 递推式求导： \[B_{i,k}'(x)=\frac{1}{x_{i+k}-x_i}B_{i,k-1}(x)+\frac{x-x_i}{x_{i+k}-x_i}B_{i,k-1}'(x)-\frac{1}{x_{i+k+1}-x_{i+1}}B_{i+1,k-1}(x)+\frac{x_{i+k+1}-x}{x_{i+k+1}-x_{i+1}}B_{i+1,k-1}'(x)\] 利用 B-样条的递推对低一次 B-样条求导的公式 \(B_{j,k-1}'(x)=\frac{k-1}{x_{j+k-1}-x_j}B_{j,k-2}(x)-\frac{k-1}{x_{j+k}-x_{j+1}}B_{j+1,k-2}(x)\)，代入并整理即得 \(B_{i,k}'(x)=\frac{k}{x_{i+k}-x_i}B_{i,k-1}(x)-\frac{k}{x_{i+k+1}-x_{i+1}}B_{i+1,k-1}(x)\)。由归纳法可证。\(lacksquare\)

\subsection{210.3 小波分析}\label{ux5c0fux6ce2ux5206ux6790}

\textbf{定义 210.4}（多分辨分析，Mallat 1989）：\(L^2(\mathbb{R})\) 的\textbf{多分辨分析}（MRA）是一列闭子空间 \(\cdots \subset V_{-1} \subset V_0 \subset V_1 \subset \cdots\) 满足： 1. \(\bigcup V_j\) 在 \(L^2(\mathbb{R})\) 中稠密，\(\bigcap V_j = \{0\}\) 2. \(f(x) \in V_j \iff f(2x) \in V_{j+1}\) 3. 存在缩放函数（尺度函数）\(\phi \in V_0\)，其整数平移构成 \(V_0\) 的规范正交基

\textbf{定义 210.5}（小波函数）：设 \(W_j\) 是 \(V_j\) 在 \(V_{j+1}\) 中的正交补（\(V_{j+1} = V_j \oplus W_j\)）。\textbf{小波函数}（母小波）\(\psi \in W_0\)，其整数平移构成 \(W_0\) 的规范正交基。\(\{\psi_{j,k}(x) = 2^{j/2} \psi(2^j x - k)\}_{j,k \in \mathbb{Z}}\) 构成 \(L^2(\mathbb{R})\) 的规范正交基。

\textbf{定理 210.4}（Daubechies 小波，1988）：存在具有紧支撑的任意高阶光滑小波函数。Daubechies 小波 \(\psi_N\) 的支撑为 \([0, 2N-1]\)，具有 \(N\) 个消失矩：\(\int x^k \psi_N(x) dx = 0\)（\(k = 0, \ldots, N-1\)）。

\textbf{证明}：Daubechies 的构造基于二尺度方程 \(\phi(x)=\sqrt{2}\sum_{k=0}^{2N-1}h_k\phi(2x-k)\) 和小波方程 \(\psi(x)=\sqrt{2}\sum_{k=0}^{2N-1}g_k\phi(2x-k)\)（\(g_k=(-1)^k h_{2N-1-k}\)）。滤波器 \(h_k\) 满足正交性 \(\sum_k h_k h_{k+2n}=\delta_{0n}\) 和 \(N\) 个消失矩条件 \(\sum_k (-1)^k k^m h_k=0\)（\(m=0,\ldots,N-1\)）。通过 Bezout 恒等式 \(|H(\omega)|^2=(\cos^2(\omega/2))^N P(\sin^2(\omega/2))\) 构造多项式 \(P\)，再由谱分解取 \(H(\omega)\) 即得 \(h_k\)。紧支撑性由 \(h_k\) 的有限项保证，支撑为 \([0,2N-1]\)。\(lacksquare\)

\textbf{定理 210.5}（小波的逼近性质）：具有 \(N\) 个消失矩的小波，函数 \(f\) 的小波系数的衰减速率刻画了 \(f\) 的光滑性：

\[|\langle f, \psi_{j,k} \rangle| \leq C 2^{-j(N+1/2)} \|f^{(N)}\|_\infty\]

\textbf{证明}：设 \(\psi\) 具有 \(N\) 个消失矩：\(\int x^k\psi(x)dx=0\)（\(k=0,\ldots,N-1\)）。在 \(x_0=2^{-j}k\) 附近对 \(f\) 做 Taylor 展开：\(f(x)=\sum_{m=0}^{N-1}\frac{f^{(m)}(x_0)}{m!}(x-x_0)^m+R_N(x)\)。由消失矩，前 \(N\) 项与 \(\psi_{j,k}\) 的内积为零。剩余项 \(\langle R_N,\psi_{j,k}\rangle\) 使用 Cauchy-Schwarz 不等式和变量替换 \(y=2^j x-k\)： \[|\langle f,\psi_{j,k}\rangle|\le\frac{\|f^{(N)}\|_\infty}{N!}\int|x-2^{-j}k|^N 2^{j/2}|\psi(2^j x-k)|dx\] 做变量替换 \(y=2^j x-k\)，得 \(\le C\cdot 2^{-j(N+1/2)}\|f^{(N)}\|_\infty\)。\(lacksquare\)

\subsection{210.4 小波在信号处理中的应用}\label{ux5c0fux6ce2ux5728ux4fe1ux53f7ux5904ux7406ux4e2dux7684ux5e94ux7528}

\textbf{定义 210.6}（快速小波变换，FWT = Mallat 算法）：小波分解和重构可通过滤波器组（低通 \(h\) 和高通 \(g\)）以 \(O(n)\) 时间复杂度实现：

\[a_{j,k} = \sum_m h_{m-2k} a_{j+1,m}, \quad d_{j,k} = \sum_m g_{m-2k} a_{j+1,m}\]

\textbf{定理 210.6}（小波收缩与降噪，Donoho-Johnstone 1994）：通过阈值处理小波系数（软阈值或硬阈值），小波方法在降噪中达到近乎最优的最小最大速率。

\textbf{证明}：Donoho-Johnstone 证明了在 Gauss 白噪声模型 \(y_i=f(t_i)+\varepsilon z_i\) 下，软阈值估计 \(\hat{\theta}_\lambda=\operatorname{sgn}(y)(|y|-\lambda)_+\) 以小波系数为对象时，风险函数 \(R(\hat{f},f)=\mathbb{E}\|\hat{f}-f\|^2\) 满足 \[R(\hat{f},f)\le C(\lambda^2+\sum_j\min(|\theta_j|^2,\lambda^2))\] 取 \(\lambda=\sigma\sqrt{2\log n}\)，该风险在 \(\ell_p\) 球（\(0<p<2\)）上逼近极小最大风险的 \((2\log n)\) 倍。核心在于：小波变换将信号能量集中在少量大系数上，阈值处理保留大系数、抑制噪声小系数，近于最优。\(lacksquare\)

\hrulefill

\hrulefill

\hrulefill

\hrulefill

\hrulefill

\hrulefill

\section{第397章：径向基函数与高维逼近}\label{ux7b2c397ux7ae0ux5f84ux5411ux57faux51fdux6570ux4e0eux9ad8ux7ef4ux903cux8fd1}

径向基函数（RBF）方法突破了传统逼近论对规则网格的依赖，为高维散乱数据的插值和逼近提供了一套优雅的数学框架。在传统多项式逼近和样条逼近中，基函数通常依赖于数据点的网格结构------Chebyshev 节点、均匀网格或 B-样条节点序列。然而，在许多科学计算问题（如地球物理数据重建、机器学习中的核方法、计算流体力学中的无网格方法）中，数据点往往是散乱分布在高维空间中的，此时基于网格的逼近方法不再适用。RBF 方法通过将近似函数表示为径向对称基函数的平移 \(\phi(\|x-x_i\|)\) 的线性组合------基函数仅依赖于到中心点的距离，从而自然地将问题从高维空间降维到一维径向函数上。Gauss RBF \(\phi(r)=e^{-(\varepsilon r)^2}\)（具有谱收敛性）、多二次 RBF \(\phi(r)=\sqrt{1+(\varepsilon r)^2}\)（具有指数收敛性）和薄板样条 \(\phi(r)=r^2\log r\)（在 \(\mathbb{R}^2\) 上极小化弯曲能量）是最常用的三类径向基函数。RBF 插值的正定性理论（Bochner 定理、Micchelli 定理）保证了插值矩阵的可逆性，而收敛性理论则将逼近阶与数据点的填充距离联系起来。这些性质使 RBF 在高维逼近、机器学习（高斯过程回归）和偏微分方程数值解中发挥着日益重要的作用。

\subsection{211.1 径向基函数}\label{ux5f84ux5411ux57faux51fdux6570}

\textbf{定义 211.1}（径向基函数）：\textbf{径向基函数}（RBF）\(\phi(x) = \phi(\|x\|)\) 仅依赖于到原点的距离。常用的 RBF 包括： - \textbf{Gauss}：\(\phi(r) = e^{-(\varepsilon r)^2}\) - \textbf{多二次}（Multiquadric）：\(\phi(r) = \sqrt{1 + (\varepsilon r)^2}\) - \textbf{逆多二次}（Inverse Multiquadric）：\(\phi(r) = 1 / \sqrt{1 + (\varepsilon r)^2}\) - \textbf{薄板样条}（Thin Plate Spline）：\(\phi(r) = r^2 \log r\)（\(d=2\)）

\textbf{定义 211.2}（RBF 插值）：给定散乱数据点 \(\{(x_i, f_i)\}_{i=1}^N \subset \mathbb{R}^d \times \mathbb{R}\)，RBF 插值具有形式

\[s(x) = \sum_{i=1}^N \lambda_i \phi(\|x - x_i\|) + \sum_{j=1}^M \beta_j p_j(x)\]

其中 \(p_j\) 是多项式基（保证唯一可解性）。系数 \(\lambda_i, \beta_j\) 由插值条件和矩条件 \(\sum \lambda_i p_j(x_i) = 0\) 确定。

\subsection{211.2 正定函数与再生核 Hilbert 空间}\label{ux6b63ux5b9aux51fdux6570ux4e0eux518dux751fux6838-hilbert-ux7a7aux95f4}

\textbf{定义 211.3}（正定函数）：函数 \(\Phi: \mathbb{R}^d \to \mathbb{R}\) 称为\textbf{正定}的，如果对任意 \(\{x_1, \ldots, x_N\} \subset \mathbb{R}^d\) 和 \(\{c_1, \ldots, c_N\} \subset \mathbb{R}\)，

\[\sum_{i,j=1}^N c_i c_j \Phi(x_i - x_j) \geq 0\]

\textbf{定理 211.1}（Bochner 定理，1933）：连续函数 \(\Phi\) 是正定的当且仅当它是有限正 Borel 测度的 Fourier 变换：\(\Phi(x) = \int_{\mathbb{R}^d} e^{i x \cdot \omega} d\mu(\omega)\)。Gauss 核是正定的，多二次核是条件正定的。

\textbf{证明}：（\(\Rightarrow\)）设 \(\Phi\) 正定，定义 \(\mu\) 为谱测度。由 Bochner 的原始论证：定义 \(\mu_T(A)=(2T)^{-d}\int_{-T}^T\cdots\int_{-T}^T\int_A e^{-ix\cdot\omega}\Phi(x)dx d\omega\)，由正定性可证 \(\mu_T\) 为正测度。取弱极限 \(T\to\infty\) 得 \(\mu\) 使 \(\Phi(x)=\int e^{ix\cdot\omega}d\mu(\omega)\)。（\(\Leftarrow\)）若 \(\Phi(x)=\int e^{ix\cdot\omega}d\mu(\omega)\)，则 \(\sum c_i c_j\Phi(x_i-x_j)=\int|\sum c_i e^{ix_i\cdot\omega}|^2 d\mu(\omega)\ge 0\)。\(lacksquare\)

\textbf{定义 211.4}（再生核 Hilbert 空间，RKHS）：设 \(\mathcal{H}\) 是函数 \(f: \Omega \to \mathbb{R}\) 的 Hilbert 空间。\(\mathcal{H}\) 是\textbf{再生核 Hilbert 空间}，如果存在核 \(K: \Omega \times \Omega \to \mathbb{R}\) 使得对任意 \(x \in \Omega\)，\(K(\cdot, x) \in \mathcal{H}\) 且 \(f(x) = \langle f, K(\cdot, x) \rangle_{\mathcal{H}}\)（再生性质）。

\textbf{定理 211.2}（RBF 与 RKHS 的等价性）：每个正定 RBF \(\phi\) 关联一个唯一的 RKHS \(\mathcal{N}_\phi(\Omega)\)（Native 空间）。RBF 插值在此空间中是误差最小的：

\[\|f - s\|_{\mathcal{N}_\phi} \leq \|f\|_{\mathcal{N}_\phi}\]

\textbf{证明}：定义 Native 空间 \(\mathcal{N}_\phi(\Omega)\) 为 \(\operatorname{span}\{\phi(\cdot-x):x\in\Omega\}\) 在范数 \(\|\cdot\|_{\mathcal{N}_\phi}\) 下的完备化，其中内积由 \(\langle\phi(\cdot-x),\phi(\cdot-y)\rangle_{\mathcal{N}_\phi}=\phi(x-y)\) 定义。由正定性，此内积良定。\(s\) 为 RBF 插值，对任意 \(f\in\mathcal{N}_\phi\)，\(\langle f-s,s\rangle_{\mathcal{N}_\phi}=0\)（正交投影性质），故 \(\|f\|_{\mathcal{N}_\phi}^2=\|f-s\|_{\mathcal{N}_\phi}^2+\|s\|_{\mathcal{N}_\phi}^2\)，因此 \(\|f-s\|_{\mathcal{N}_\phi}\le\|f\|_{\mathcal{N}_\phi}\)。\(lacksquare\)

\subsection{211.3 RBF 逼近的收敛性}\label{rbf-ux903cux8fd1ux7684ux6536ux655bux6027}

\textbf{定理 211.3}（RBF 插值的误差估计）：对 Gauss 核 \(\phi(r) = e^{-(\varepsilon r)^2}\)，在 Native 空间中，

\[\|f - s_{X,\varepsilon}\|_\infty \leq C e^{-c \varepsilon / h_{X,\Omega}} \|f\|_{\mathcal{N}_\phi}\]

其中 \(h_{X,\Omega}\) 是填充距离（fill distance）。\(h_{X,\Omega} \to 0\) 时误差以谱精度收敛。

\textbf{证明}：设 \(h=\max_{x\in\Omega}\min_{x_i\in X}\|x-x_i\|\) 为填充距离。Gauss 核的 Fourier 变换为 \(\hat{\phi}(\omega)=(\sqrt{\pi}/\varepsilon)^d e^{-\|\omega\|^2/(4\varepsilon^2)}\)。由 Fourier 表示的误差分析，\(\|f-s_X\|_{L^\infty}\le C \|f\|_{\mathcal{N}_\phi}\cdot\inf_{g\in\operatorname{span}\{\phi(\cdot-x_i)\}}\|f-g\|_{L^\infty}\)。利用单位分解和指数衰减的逼近核，误差界为 \(C e^{-c/(\varepsilon h)}\)（对足够小的 \(h\)）。谱精度的来源是 Gauss 核的 Fourier 变换在全频带上的超代数衰减。\(lacksquare\)

\textbf{定理 211.4}（权衡原则，Trade-off Principle）：RBF 插值在精度和稳定性之间存在权衡。小 \(\varepsilon\)（平坦核）提高精度但增大条件数，大 \(\varepsilon\) 提高稳定性但降低精度。这是 RBF 方法的核心挑战。

\textbf{证明}：RBF 插值矩阵 \(A_{ij}=\phi(\|x_i-x_j\|)\) 的条件数满足 \(\kappa_2(A)\sim e^{c/\varepsilon}\)（对 Gauss 核，\(h\) 固定）。同时误差 \(\|f-s\|_\infty\sim e^{-c'/(\varepsilon h)}\)。两者的乘积 \(\kappa_2(A)\times\|f-s\|_\infty\) 在 \(\varepsilon\) 的某个中间值时达到平衡。这是逼近论中普遍的不确定原理：更高精度需要更不稳定的计算。该折衷源自 \(\phi\) 在 \(\varepsilon\to 0\) 时的平坦化导致基函数趋于线性相关。\(lacksquare\)

\textbf{定理 211.5}（平坦极限与多项式，Driscoll-Fornberg 2002）：当 \(\varepsilon \to 0\)（平坦极限），某些 RBF 插值收敛到（多元）多项式插值。这解释了 RBF 方法在平坦极限下的行为。

\textbf{证明}：对解析 RBF \(\phi_\varepsilon(r)=\phi(\varepsilon r)\)，在 \(\varepsilon\to 0\) 时做 Taylor 展开：\(\phi_\varepsilon(r)=\sum_{k=0}^\infty a_{2k}\varepsilon^{2k}r^{2k}\)。RBF 插值 \(s_\varepsilon(x)\) 的极限函数满足 \(\lim_{\varepsilon\to 0}s_\varepsilon(x)=p^*(x)\)，其中 \(p^*\) 是满足相同插值条件的唯一最低阶多元多项式。原因：展开式中低阶项主导，高阶项在 \(\varepsilon\to 0\) 时消失，而插值条件的线性系统退化为多项式插值的 Vandermonde 型矩阵。Driscoll-Fornberg 数值验证了此现象。\(lacksquare\)

\subsection{211.4 RBF 在 PDE 中的应用}\label{rbf-ux5728-pde-ux4e2dux7684ux5e94ux7528}

\textbf{定义 211.5}（Kansa 方法，RBF 配置法，1990）：用 RBF 配置法求解 PDE：将解近似为 \(u(x) \approx \sum_{i=1}^N \lambda_i \phi(\|x - x_i\|)\)，在配置点 \(x_j\) 上强制 PDE 满足，得到关于 \(\lambda_i\) 的线性方程组。

\textbf{定理 211.6}（RBF 方法在无网格 PDE 求解中的优势）：RBF 方法自然地处理高维和不规则区域上的 PDE，无需网格生成。对于椭圆形 PDE，RBF 配置法可达谱精度。

\textbf{证明}：对椭圆 PDE \(\mathcal{L}u=f\)（\(\Omega\subset\mathbb{R}^d\)），RBF 配置法设 \(u_N(x)=\sum_{i=1}^N\lambda_i\phi(\|x-x_i\|)\)，强制 \(\mathcal{L}u_N(x_j)=f(x_j)\)（\(j=1,\ldots,N_i\)）及边界条件 \(u_N(x_j)=g(x_j)\)（\(j=N_i+1,\ldots,N\)）。矩阵 \(B_{ij}=\mathcal{L}\phi(\|x_j-x_i\|)\)（或 \(\phi(\|x_j-x_i\|)\)），对正定 \(\phi\) 在 \(x_i\) 满足一定分离条件下非奇异。误差由 \(\|u-u_N\|_\infty\le\|u\|_{\mathcal{N}_\phi}\cdot\sup_{x\in\Omega}\|\mathcal{L}K(\cdot,x)-\mathcal{L}_h K(\cdot,x)\|_{\mathcal{N}_\phi}\) 控制，\(h\to 0\) 时可达谱精度，且无需网格。\(lacksquare\)

\hrulefill

\hrulefill

\hrulefill

\hrulefill

\hrulefill

\hrulefill

\section{第398章：Chebyshev理论与最佳一致逼近}\label{ux7b2c398ux7ae0chebyshevux7406ux8bbaux4e0eux6700ux4f73ux4e00ux81f4ux903cux8fd1}

Pafnuty Chebyshev 在 1850 年代创立的最佳一致逼近理论是逼近论的奠基性工作，其核心问题是：在所有次数不超过 \(n\) 的多项式中，哪一个最''贴近''给定的连续函数 \(f\)？这里的''贴近''以一致范数（Chebyshev 范数）\(\|f-p\|_\infty=\max_{x\in[a,b]}|f(x)-p(x)|\) 度量。Chebyshev 对这一问题的开创性贡献体现在两个方面。第一是 Chebyshev 多项式 \(T_n(x)=\cos(n\arccos x)\) 的发现：\(T_n\) 在 \([-1,1]\) 上在幅度 \(1\) 之内剧烈振荡，并且在所有 \(n\) 次首一多项式中，\(2^{1-n}T_n\) 具有最小的一致范数------这一极值性质直接导出多项式插值节点的最优选取策略（Chebyshev 节点 \(\cos((2k+1)\pi/(2n+2))\)），使得插值的 Runge 现象最小化。第二是交替定理（Chebyshev Alternation Theorem）：\(p^*\) 是 \(f\) 的最佳一致逼近多项式当且仅当误差函数 \(f-p^*\) 在 \([a,b]\) 上至少 \(n+2\) 次达到其最大绝对值，且符号交替------这为最佳逼近的验证提供了简洁的充要条件。Remez 算法利用交替定理在数值上高效计算最佳一致逼近多项式。

\subsection{212.1 Chebyshev多项式}\label{chebyshevux591aux9879ux5f0f}

\textbf{定义 212.1}（Chebyshev 多项式）：第一类 Chebyshev 多项式 \(T_n(x) = \cos(n \arccos x)\)，\(x \in [-1, 1]\)。第二类 \(U_n(x) = \frac{\sin((n+1) \arccos x)}{\sin(\arccos x)}\)。\(T_n\) 是次数为 \(n\) 的首一多项式（乘以 \(2^{1-n}\)），在 \([-1, 1]\) 上满足 \(|T_n(x)| \leq 1\)。

\textbf{命题 212.1}（\(T_n\) 的递推与正交性）：\(T_0(x) = 1\)，\(T_1(x) = x\)，\(T_{n+1}(x) = 2x T_n(x) - T_{n-1}(x)\)。在 \([-1, 1]\) 上，\(\int_{-1}^1 \frac{T_m(x) T_n(x)}{\sqrt{1 - x^2}} dx = \begin{cases} \pi & m = n = 0 \\ \pi/2 & m = n \neq 0 \\ 0 & m \neq n \end{cases}\)。

\textbf{定理 212.1}（Chebyshev 多项式的极值性质）：在所有 \(n\) 次首一多项式中，\(2^{1-n} T_n(x)\) 在 \([-1, 1]\) 上具有最小的一致范数： \[\min_{P \in \mathcal{P}_n, P \text{ 首一}} \max_{x \in [-1,1]} |P(x)| = 2^{1-n}\] 且极小化元是唯一的（\(2^{1-n} T_n\)）。此极值性质给出了多项式插值节点的最优选取------Chebyshev 节点。

\textbf{证明}：设 \(P(x)=x^n+\sum_{k=0}^{n-1}a_k x^k\) 为任意首一 \(n\) 次多项式。令 \(Q(x)=P(x)-2^{1-n}T_n(x)\)，次数 \(\le n-1\)。在 \(T_n\) 的 \(n+1\) 个极值点 \(x_k=\cos(k\pi/n)\)（\(k=0,\ldots,n\)）处，\(T_n(x_k)=(-1)^k\)，\(|2^{1-n}T_n(x_k)|=2^{1-n}\)。若 \(\max_{x\in[-1,1]}|P(x)|<2^{1-n}\)，则 \(Q(x_k)\) 与 \(-2^{1-n}T_n(x_k)\) 同号，故 \(Q\) 在 \(n+1\) 个点处交错变号，由 Rolle 定理 \(Q'\) 在 \(n\) 个区间内各有一零点，\ldots，最终 \(Q^{(n)}\) 至少有一零点------但 \(\deg Q\le n-1\) 则 \(Q^{(n)}\equiv 0\)，矛盾。故 \(\|P\|_\infty\ge 2^{1-n}\)，且等号仅当 \(P=2^{1-n}T_n\)。\(lacksquare\)

\textbf{推论 212.1}（Chebyshev 交替定理的特例）：\(2^{1-n} T_n(x)\) 在 \([-1, 1]\) 上达到最大值 \(2^{1-n}\) 的 \(n+1\) 个交替点：\(x_k = \cos(k\pi/n)\)，\(k = 0, 1, \ldots, n\)。这是交替定理的具体体现------误差函数在这些点上达到极值且符号交替。

\subsection{212.2 Chebyshev交替定理}\label{chebyshevux4ea4ux66ffux5b9aux7406}

\textbf{定理 212.2}（Chebyshev 交替定理，1853）：\(p_n^* \in \mathcal{P}_n\) 是 \(f \in C[a, b]\) 的最佳一致多项式逼近（即 \(\|f - p_n^*\|_\infty = E_n(f)\)），当且仅当存在 \([a, b]\) 中的 \(n+2\) 个点 \(a \leq x_1 < x_2 < \cdots < x_{n+2} \leq b\)，使得误差 \(e(x) = f(x) - p_n^*(x)\) 在这些点上交错地取极值：\(e(x_k) = \pm (-1)^k \|f - p_n^*\|_\infty\)（或 \(\mp (-1)^k\)）。这就是\textbf{交替点集}。

\textbf{证明}（必要性）：设 \(p_n^*\) 为最佳逼近，令 \(E = \|f - p_n^*\|_\infty\)。取 \(x_1 = \min\{x \in [a,b] : |f(x) - p_n^*(x)| = E\}\)。假设 \(f(x_1) - p_n^*(x_1) = E\)（否则取 \(-E\)，论证对称）。归纳地，设已选 \(x_1 < \cdots < x_k\) 使 \(e(x_i) = (-1)^{i-1}E\)（\(i=1,\ldots,k\)）。若 \(k \leq n+1\)，取 \(x_{k+1} = \min\{x > x_k : e(x) = (-1)^k E\}\)。若 \(k \leq n+1\) 时无法继续，则 \(e(x)\) 在 \([x_k, b]\) 上不取 \((-1)^k E\)。令 \(s = \max\{x \in [x_k, b] : e(x) = (-1)^{k-1}E\}\)。构造 \(q(x) = \prod_{i=1}^{k-1}(x - x_i) \in \mathcal{P}_{k-1} \subset \mathcal{P}_n\)。对足够小的 \(\varepsilon > 0\)，\(p_n^* + \varepsilon q\) 逼近 \(f\) 比 \(p_n^*\) 更好（在交替点附近，\(q\) 与 \(e\) 反号，减小误差），与最佳性矛盾。故交替点至少 \(n+2\) 个。（充分性）若有 \(n+2\) 个交替点，设 \(p \in \mathcal{P}_n\) 使 \(\|f - p\|_\infty < E\)。则 \(p - p_n^*\) 在交替点处与 \(f - p_n^*\) 同号，由交错性，\(p - p_n^*\) 在 \(n+1\) 个区间各有一零点，共 \(n+1\) 个零点，但 \(\deg(p - p_n^*) \leq n\)，故 \(p = p_n^*\)，唯一性得证。\(\blacksquare\)

\textbf{推论 212.2}（最佳逼近的误差公式）：若 \(f \in C^{n+1}[a, b]\) 且 \(f^{(n+1)}\) 在 \((a, b)\) 中不变号，取 \(p_n^*\) 为 Lagrange 插值多项式（节点为 \(f - p_n^*\) 的 Chebyshev 交替点附近的点），则 \(E_n(f) \geq \frac{\min|f^{(n+1)}|}{2^n (n+1)!} (b-a)^{n+1}\)。

\textbf{定理 212.3}（de la Vallée Poussin 下界定理，1910）：设 \(p \in \mathcal{P}_n\)，且存在 \(n+2\) 个点 \(x_1 < \cdots < x_{n+2}\) 使 \(f(x_k) - p(x_k)\) 交替变号。则 \[E_n(f) \geq \min_k |f(x_k) - p(x_k)|\] 这给出了最佳逼近的下界估计------只需找一个产生交替符号的多项式即可判定逼近精度。

\textbf{证明}：设 \(p\in\mathcal{P}_n\) 且存在 \(x_1<\cdots<x_{n+2}\) 使 \(e(x_k)=f(x_k)-p(x_k)\) 交错变号且 \(|e(x_k)|\ge\mu\)。对任意 \(q\in\mathcal{P}_n\)，令 \(h=p-q\in\mathcal{P}_n\)。\(h\) 在 \(x_k\) 处的符号由 \(f(x_k)-p(x_k)\) 决定：若 \(p\) 的第 \(k\) 个误差值与 \(q\) 的同号，则 \(|f(x_k)-q(x_k)|\ge\mu\)。否则 \(h\) 在相邻 \(x_k\) 间变号 \(n+1\) 次，由中间值定理 \(h\) 至少有 \(n+1\) 个零点，但 \(\deg h\le n\)，故 \(h\equiv0\)。因此 \(\|f-q\|_\infty\ge\min_k|f(x_k)-p(x_k)|\ge\mu\)。取下确界得 \(E_n(f)\ge\mu\)。\(lacksquare\)

\subsection{212.3 Chebyshev多项式展开与谱方法}\label{chebyshevux591aux9879ux5f0fux5c55ux5f00ux4e0eux8c31ux65b9ux6cd5}

\textbf{定义 212.2}（Chebyshev 级数展开）：函数 \(f \in C[-1, 1]\) 的 Chebyshev 级数为 \(f(x) \sim \sum_{k=0}^\infty a_k T_k(x)\)，其中 \(a_k = \frac{2}{\pi} \int_{-1}^1 \frac{f(x) T_k(x)}{\sqrt{1-x^2}} dx\)（\(k \geq 1\)），\(a_0 = \frac{1}{\pi} \int_{-1}^1 \frac{f(x)}{\sqrt{1-x^2}} dx\)。Chebyshev 级数收敛极快（对解析函数呈几何收敛），因而是数值分析和谱方法的基础。

\textbf{定理 212.4}（Chebyshev 级数的收敛速度）：若 \(f\) 在包含 \([-1, 1]\) 的 Bernstein 椭圆 \(\mathcal{E}_\rho\)（椭圆半轴 \((\rho + \rho^{-1})/2 \times (\rho - \rho^{-1})/2\)）内解析，则 \(|a_k| \leq 2M \rho^{-k}\)（几何收敛）。解析性越强（\(\rho\) 越大），Chebyshev 系数衰减越快。

\textbf{证明}：设 \(f\) 在 Bernstein 椭圆 \(\mathcal{E}_\rho\)（半轴 \(a=(\rho+\rho^{-1})/2\)，\(b=(\rho-\rho^{-1})/2\)）内解析。将 \(f\) 表达为 Joukowsky 变换下的 Laurent 级数 \(f(\frac{z+z^{-1}}{2})=\sum_{k=-\infty}^\infty c_k z^k\)（\(|z|=1\) 时还原 \(f(\cos\theta)\)）。由 Cauchy 积分公式， \[c_k=\frac{1}{2\pi i}\oint_{|z|=\rho}\frac{f(\frac{z+z^{-1}}{2})}{z^{k+1}}dz\] 估计 \(|c_k|\le M\rho^{-k}\)（\(M\) 为 \(f\) 在该椭圆上的上界）。Chebyshev 系数 \(a_k=2c_k\)（\(k\ge1\)），故 \(|a_k|\le 2M\rho^{-k}\)，呈几何收敛。\(lacksquare\)

\textbf{定理 212.5}（Remez 不等式，1936）：设 \(p_n \in \mathcal{P}_n\) 在 \([a, b]\) 的一个可测子集 \(E \subset [a, b]\) 上满足 \(|p_n(x)| \leq M\)（\(|E| > 0\)）。则 \[\max_{x \in [a, b]} |p_n(x)| \leq M \cdot T_n\left(\frac{2b - 2a + |E|}{2b - 2a - |E|}\right)\] Remez 不等式说明：多项式在一个子集上的有界性控制它在整个区间上的最大值------增长因子由 Chebyshev 多项式决定。这是多项式不等式中最基本的结果之一，在构造插值算子和证明逆逼近定理中起关键作用。

\textbf{证明}：设 \(E\subset[a,b]\) 可测且 \(|E|>0\)，\(p_n\in\mathcal{P}_n\) 在 \(E\) 上 \(|p_n|\le 1\)。考虑 \(x_0\in[a,b]\setminus E\) 使 \(|p_n(x_0)|=\|p_n\|_\infty\)。令 \(T_n\) 为经过线性变换 \(x\mapsto(2x-a-b)/(b-a)\) 后的 Chebyshev 多项式。构造辅助多项式 \(q(x)=T_n(\frac{2x-(a+b)}{b-a})\cdot\|p_n\|_\infty^{-1}\)。使用多项式插值的 Remez 型论证：在 \(E\) 的补集上以 \(p_n\) 的零点为插值节点，通过 Neville 插值公式控制外推幅度。具体地，设 \(x_0\) 距 \(E\) 的最短距离为 \(d\)，则 \(\|p_n\|_\infty\le T_n(1+2d/|E|)\)（通过将 \(E\) 替换为包含其的区间）。\(lacksquare\)

\textbf{推论 212.3}（Nikolskii 型不等式）：对 \(p_n \in \mathcal{P}_n\) 和 \(0 < p < q \leq \infty\)（\(L^p\) 范数），\(\|p_n\|_q \leq C n^{2(1/p - 1/q)} \|p_n\|_p\)。这推广了 Remez 不等式到积分范数，建立了不同 \(L^p\) 范数下的多项式比较。

\hrulefill

\hrulefill

\hrulefill

\hrulefill

\hrulefill

\hrulefill

\section{第399章：逆定理与Müntz-Szász定理}\label{ux7b2c399ux7ae0ux9006ux5b9aux7406ux4e0emuxfcntz-szuxe1szux5b9aux7406}

逼近论的直接定理（Jackson）从光滑性推出逼近速率，而\textbf{逆定理}（Bernstein / Stechkin）从逼近速率反推光滑性------两者共同构成逼近论的基本等价关系。Müntz-Szász 定理则回答了''哪些幂函数的线性组合足以逼近所有连续函数''------是 Weierstrass 逼近定理在单项式子集上的深刻推广。

\subsection{213.1 Bernstein逆定理}\label{bernsteinux9006ux5b9aux7406}

\textbf{定理 213.1}（Bernstein 不等式，1912）：对 \(p_n \in \mathcal{P}_n\)（次数 \(\leq n\) 的代数多项式），在 \([-1, 1]\) 上有 \[\|p_n'\|_\infty \leq n^2 \|p_n\|_\infty\] 这是多项式导数与函数值范数之间的基本关系：\(n\) 次多项式的导数范数不能超过函数值范数的 \(n^2\) 倍。对三角多项式 \(T_n(\theta) = \sum_{k=-n}^n c_k e^{ik\theta}\)，有更精确的 \(\|T_n'\|_\infty \leq n \|T_n\|_\infty\)。

\textbf{证明}：对三角多项式情形，\(T_n(\theta)=\sum_{k=-n}^n c_k e^{ik\theta}\)，\(T_n'(\theta)=i\sum_{k=-n}^n k c_k e^{ik\theta}\)。\(T_n\) 的卷积表示：\(T_n(\theta)=(F_n*T_n)(\theta)\)，其中 \(F_n\) 为 Fejér 核。由导数与卷积的交换性和 \(F_n'\) 的 \(L^1\) 范数估计 \(\|F_n'\|_{L^1}=n\)，得 \(\|T_n'\|_\infty\le n\|T_n\|_\infty\)。对代数多项式，令 \(x=\cos\theta\) 将 \(p_n\) 化为三角多项式：\(p_n(\cos\theta)=T_{2n}(\theta)\)。由链式法则 \(p_n'(x)=-T_{2n}'(\theta)/\sin\theta\)，利用 Bernstein 不等式对三角多项式 \(T_{2n}\) 得 \(\|p_n'\|_\infty\le n^2\|p_n\|_\infty\)。\(lacksquare\)

\textbf{定理 213.2}（Bernstein 逆定理，1912）：设 \(f \in C[-1, 1]\) 满足 \(E_n(f) \leq C n^{-(r+\alpha)}\)（对所有 \(n\)），其中 \(r \in \mathbb{Z}_{\geq 0}\)，\(0 < \alpha \leq 1\)。则 \(f^{(r)}\) 存在且连续（或当 \(\alpha=1\) 时满足 Zygmund 条件），且 \(f^{(r)} \in \operatorname{Lip}_\alpha\)（即 \(f^{(r)}\) 的连续模 \(\omega(f^{(r)}; \delta) \leq C' \delta^\alpha\)）。

\textbf{证明}：取几何级数 \(n_k = 2^k\)（\(k \geq 0\)）。由假设，存在多项式 \(p_{n_k} \in \mathcal{P}_{n_k}\) 使 \(\|f - p_{n_k}\|_\infty \leq C n_k^{-(r+\alpha)}\)。令 \(q_0 = p_1\)，\(q_k = p_{n_k} - p_{n_{k-1}}\)（\(k \geq 1\)），则 \(f = q_0 + \sum_{k=1}^\infty q_k\)（一致收敛）。对 \(q_k\) 应用 Bernstein 不等式（定理 213.1）：\(\|q_k^{(r)}\|_\infty \leq n_k^{2r} \|q_k\|_\infty\)。而 \(\|q_k\|_\infty \leq \|f - p_{n_k}\|_\infty + \|f - p_{n_{k-1}}\|_\infty \leq C n_k^{-(r+\alpha)} + C n_{k-1}^{-(r+\alpha)} \leq C' n_k^{-(r+\alpha)}\)。故 \(\|q_k^{(r)}\|_\infty \leq C' n_k^{2r} n_k^{-(r+\alpha)} = C' n_k^{r-\alpha}\)。对 \(r\) 阶导数，使用更精细的 Bernstein-Markov 不等式 \(\|q_k^{(r)}\|_\infty \leq n_k^r \|q_k\|_\infty\)（三角多项式情形通过 \(x = \cos\theta\) 变换得到），得 \(\|q_k^{(r)}\|_\infty \leq C' n_k^{-\alpha}\)。级数 \(\sum \|q_k^{(r)}\|_\infty\) 为几何级数 \(\sum n_k^{-\alpha} = \sum 2^{-k\alpha}\) 收敛，故 \(f^{(r)} = \sum q_k^{(r)}\) 一致收敛。对任意 \(\delta > 0\)，取 \(K\) 使 \(2^{-K} \leq \delta < 2^{-K+1}\)，则 \(\omega(f^{(r)}; \delta) \leq 2\sum_{k=K}^\infty \|q_k^{(r)}\|_\infty + \delta \sum_{k=0}^{K-1} \|q_k^{(r+1)}\|_\infty \leq C'' \delta^\alpha\)。故 \(f^{(r)} \in \operatorname{Lip}_\alpha\)。\(\blacksquare\)

\textbf{定理 213.3}（Stechkin 逆定理，1950s）：Bernstein 逆定理的推广------将多项式的最佳逼近误差 \(E_n(f)\) 替换为更强的连续模条件，得到函数光滑性的更精细分类。构造多项式理论（constructive function theory）的基本框架由此建立。

\textbf{证明}：Stechkin 的推广将 \(E_n(f)\) 替换为 \(k\) 阶连续模 \(\omega_k(f;\delta)=\sup_{0<h\le\delta}\|\Delta_h^k f\|_\infty\)。利用 Jackson 不等式和 Bernstein 不等式的结合：对最佳逼近多项式 \(p_n\)，\(\omega_k(f;1/n)\le C_k n^{-k}\sum_{\nu=0}^n (\nu+1)^{k-1}E_\nu(f)\) 以及逆不等式 \(E_n(f)\le C_k\omega_k(f;1/n)\)。两者联合建立了 \(\omega_k(f;\delta)\) 与 \(E_n(f)\) 之间的等价关系，从而实现光滑性的精细分类（如 \(\omega_k(f;\delta)=O(\delta^\alpha)\) 等价于 \(E_n(f)=O(n^{-\alpha})\)）。\(lacksquare\)

\textbf{推论 213.1}（光滑性的等价刻画）：\(f \in C^\infty[-1, 1]\) 当且仅当 \(E_n(f) = O(n^{-k})\) 对所有 \(k\) 成立。\(f\) 在 \([-1, 1]\) 上解析当且仅当 \(E_n(f) = O(\rho^{-n})\) 对某 \(\rho > 1\)（几何收敛）。

\subsection{213.2 Müntz-Szász定理}\label{muxfcntz-szuxe1szux5b9aux7406}

\textbf{定理 213.4}（Müntz 定理，1914 / Szász 推广，1916）：设 \(\Lambda = \{\lambda_0, \lambda_1, \ldots\} \subset [0, \infty)\) 且 \(\lambda_0 = 0\)。则幂函数集合 \(\{x^{\lambda_k}\}_{k=0}^\infty\) 的线性张成在 \(C[0, 1]\)（一致范数拓扑）中稠密，当且仅当 \[\sum_{k=1}^\infty \frac{1}{\lambda_k} = \infty\]

\textbf{证明}：（\(\Leftarrow\)，级数发散则稠密）由 Hahn-Banach 定理，若 \(\operatorname{span}\{x^{\lambda_k}\}\) 不在 \(C[0,1]\) 中稠密，则存在非零 Borel 测度 \(\mu\) 使 \(\int_0^1 x^{\lambda_k}d\mu(x)=0\)（对所有 \(k\)）。定义复变函数 \(F(z)=\int_0^1 x^z d\mu(x)\)（\(\Re z>-1\)），\(F\) 在右半平面解析且 \(F(\lambda_k)=0\)。由 Carlson 定理的变体，若 \(\sum 1/\lambda_k=\infty\)（且 \(\lambda_k\) 无有限聚点），则在 \(\Re z>-1\) 上 \(F\) 零点集满足 Blaschke 条件失效，唯一性定理迫使 \(F\equiv0\)。由 Stieltjes 变换的反演公式得 \(\mu=0\)，矛盾。故级数发散时闭包稠密。（\(\Rightarrow\)，级数收敛则不稠密）若 \(\sum 1/\lambda_k<\infty\)，可构造 Blaschke 乘积 \(B(z)=\prod(1-z/\lambda_k)/(1+z/\lambda_k)\)，其在 \(\Re z>0\) 上解析且有界。由 Paley-Wiener 定理，\(B\) 的 Fourier-Laplace 逆变换给出非零测度 \(\mu\) 满足 \(\int_0^1 x^{\lambda_k}d\mu=0\)，故闭包不稠密。\(\blacksquare\)

\textbf{推论 213.2}（Müntz-Szász 定理的特殊情况与含义）： - \(\{1, x, x^2, x^3, \ldots\}\)：\(\lambda_k = k\)（整数幂），\(\sum \frac{1}{k} = \infty\) → 稠密（Weierstrass 定理的推论） - \(\{1, x^2, x^4, x^6, \ldots\}\)（仅偶次幂）：\(\lambda_k = 2k\)，\(\sum \frac{1}{2k} = \infty\) → 稠密 - \(\{1, x^2, x^3, x^5, x^7, \ldots\}\)：省略 \(x^1\) 和 \(x^4\)（有限个缺失），\(\sum \frac{1}{\lambda_k}\) 仍发散 → 稠密 - \(\{1, x^2, x^4, x^8, x^{16}, \ldots\}\)（几何级数间距）：\(\lambda_k = 2^k\)，\(\sum \frac{1}{2^k} = 1 < \infty\) → \textbf{不稠密}------存在不能被只含 \(x^{2^k}\) 次幂的多项式逼近的连续函数 - \(\{1, x^n, x^{n+1}, \ldots\}\) 和 \(\{1, x^n\}\)：前者 \(\sum 1/k\) 仍从 \(n\) 开始发散 → 稠密；后者仅有一个非零指数 → \(\sum 1/\lambda_1 = 1/n < \infty\) → 不稠密

Müntz-Szász 定理深刻揭示了：仅靠无穷多个幂函数并不能保证稠密性；这些幂指数的倒数级数必须发散------即幂指数必须''充分密集''。

\textbf{证明}（泛函分析视角）：由 Hahn-Banach 定理，\(\operatorname{span}\{x^{\lambda_k}\}\) 在 \(C[0,1]\) 中不稠密当且仅当存在非零有限 Borel 测度 \(\mu\) 使 \(\int_0^1 x^{\lambda_k} d\mu(x) = 0\)（对所有 \(k\)）。定义复变函数 \(F(z) = \int_0^1 x^z d\mu(x)\)（\(z \in \mathbb{C}\)，\(\Re z > 0\)）。\(F\) 在右半平面 \(\Re z > 0\) 上解析，且以 \(\{\lambda_k\}\) 为零点。由 Jensen 公式，对 \(R > 0\)，\(F\) 在 \(|z| \leq R\) 内的零点数 \(n(R)\) 满足 \(\int_0^R \frac{n(t)}{t} dt = \frac{1}{2\pi}\int_0^{2\pi} \log|F(Re^{i\theta})| d\theta - \log|F(0)|\)。若 \(\sum 1/\lambda_k < \infty\)，则 \(F\) 满足 Blaschke 条件，可构造非零解析函数 \(F\) 以 \(\{\lambda_k\}\) 为零点，由 Paley-Wiener 定理得其反 Laplace 变换 \(\mu\) 非零，故闭包不稠密。若 \(\sum 1/\lambda_k = \infty\)，则 \(F\) 的零点集不满足 Blaschke 条件，由解析函数唯一性定理（若 \(F\) 在聚点集上为零则 \(F \equiv 0\)），\(F \equiv 0\)，从而 \(\mu = 0\)，闭包稠密。\(\blacksquare\)

\textbf{定理 213.5}（Müntz-Szász 定理的 \(L^p\) 推广，Clarkson-Erdős 1943 / Schwartz 1959）：Müntz-Szász 定理在 \(L^p[0, 1]\)（\(1 \leq p < \infty\)）中同样成立------稠密性条件仍为 \(\sum 1/\lambda_k = \infty\)（需 \(\lambda_k \geq 0\)，且当 \(p \geq 1\) 时需 \(\lambda_k > -1/p\) 以保证 \(x^{\lambda_k} \in L^p\)）。

\textbf{证明}：设 \(\Lambda = \{\lambda_k\}\) 满足 \(\lambda_k > -1/p\)。由 Hahn-Banach 定理，\(\operatorname{span}\{x^{\lambda_k}\}\) 在 \(L^p[0,1]\) 中不稠密当且仅当存在非零 \(g \in L^q[0,1]\)（\(1/p+1/q=1\)）使 \(\int_0^1 x^{\lambda_k}g(x)dx = 0\)（对所有 \(k\)）。定义函数 \(F(z) = \int_0^1 x^z g(x)dx\)（\(\Re z > -1/p\)）。\(F\) 在区域 \(\Re z > -1/p\) 上解析，且 \(F(\lambda_k) = 0\)。若 \(\sum 1/\lambda_k = \infty\)，由 Carlson 型唯一性定理，若 \(F\) 在 \(\{\lambda_k\}\) 上为零且 \(\{\lambda_k\}\) 满足发散条件，则 \(F \equiv 0\)。由 Mellin 变换的反演公式，\(F \equiv 0\) 蕴含 \(g = 0\)（a.e.），矛盾。故 \(\sum 1/\lambda_k = \infty\) 时 \(\operatorname{span}\{x^{\lambda_k}\}\) 在 \(L^p\) 中稠密。反之，若 \(\sum 1/\lambda_k < \infty\)，可构造 Blaschke 乘积 \(B(z)\) 以 \(\{\lambda_k\}\) 为零点，且 \(B(z)\) 在 \(\Re z > -1/p\) 上有界。由 Paley-Wiener 型定理，存在非零 \(g \in L^q\) 使 \(F(z) = \int_0^1 x^z g(x)dx\) 与 \(B(z)\) 有相同的零点集，故 \(\int_0^1 x^{\lambda_k}g(x)dx = 0\)，即 \(\operatorname{span}\{x^{\lambda_k}\}\) 不稠密。\(\blacksquare\)

\textbf{定理 213.6}（完整 Müntz-Szász 定理------同时处理有限和无限缺失）：若 \(\liminf_{k \to \infty} (\lambda_{k+1} - \lambda_k) > 0\)（即 \(\Lambda\) 具有正的下渐近间隙），则 \(\sum 1/\lambda_k < \infty\) 等价于 \(\operatorname{span}\{x^{\lambda_k}\}\) 在 \(C[0, 1]\) 中不稠密。这提供了判断稠密性的实用判据。

\textbf{证明}：设 \(\liminf_{k \to \infty} (\lambda_{k+1} - \lambda_k) = \delta > 0\)。则对充分大的 \(k\)，\(\lambda_k \geq \delta k\)，故 \(\sum 1/\lambda_k\) 与 \(\sum 1/(\delta k)\) 同敛散。若 \(\sum 1/\lambda_k < \infty\)，由定理 213.4 的充分性部分，\(\operatorname{span}\{x^{\lambda_k}\}\) 在 \(C[0,1]\) 中不稠密。若 \(\sum 1/\lambda_k = \infty\)，注意到正间隙条件 \(\liminf(\lambda_{k+1}-\lambda_k) > 0\) 保证 \(\{\lambda_k\}\) 无有限聚点，定理 213.4 的必要性部分适用：由 Hahn-Banach 定理和 Carlson 唯一性定理，若存在非零测度 \(\mu\) 使 \(\int_0^1 x^{\lambda_k}d\mu = 0\)，则 \(F(z) = \int_0^1 x^z d\mu\) 在 \(\{\lambda_k\}\) 上为零。由于 \(\sum 1/\lambda_k = \infty\) 且 \(\lambda_k\) 无有限聚点，\(F\) 的零点集不满足 Blaschke 条件，强迫 \(F \equiv 0\)，从而 \(\mu = 0\)。故 \(\operatorname{span}\{x^{\lambda_k}\}\) 在 \(C[0,1]\) 中稠密。因此，在正间隙条件下，\(\sum 1/\lambda_k < \infty\) 等价于不稠密。\(\blacksquare\)

\hrulefill

\emph{本卷建立了逼近论的核心理论框架：Weierstrass 逼近定理与 Bernstein 多项式证明建立了多项式逼近的存在性基础，Jackson 直接定理给出了光滑性与逼近阶的定量关系；有理逼近与 Padé 逼近（Newman 定理、Padé 表、收敛性）在处理非光滑函数和极点上远优于多项式；样条与小波（B-样条基、多分辨分析、Daubechies 小波、快速小波变换）提供了高效且局部化的逼近工具；径向基函数（正定核、RKHS、散乱数据逼近、无网格 PDE 方法）是处理高维和复杂几何问题的主要工具。逼近论是连接泛函分析、数值分析和信号处理的桥梁。}

\hrulefill

\emph{本卷在已有内容基础上，扩展至 6 章。}

\hrulefill

\emph{本卷从浮点运算和误差分析出发，建立了线性方程组、特征值问题、数值积分和常微分方程数值解的核心算法体系。补充内容拓展至计算数学的理论深层：有限元方法的Sobolev空间理论基础（Lax-Milgram、Céa引理、inf-sup条件、后验误差估计）、谱方法与谱元法（指数收敛与hp-细化）、多重网格法（\(\mathcal{O}(N)\)最优复杂度）、双曲守恒律数值格式（Godunov定理、TVD/ENO/WENO）、以及计算流体力学（投影法、稳定化有限元、大涡模拟）。\textbf{补充部分二}引入了逼近论------数值分析的理论基石：Weierstrass逼近定理与Bernstein多项式、Chebyshev最佳一致逼近与交替定理、Jackson直接定理与Bernstein逆定理、有理逼近与Padé逼近、样条与小波（B-样条基、多分辨分析、Daubechies小波）、径向基函数与高维逼近、以及Müntz-Szász定理。共 10 章。}

\emph{卷二十二：数值分析终。}

\emph{卷二十二：数值分析终。}

\emph{卷二十二：数值分析终。}

\emph{卷二十五：数值分析终。}

\newpage

\chapter{11-数值与计算/01-数值分析/05-反问题与不适定问题.md}\label{ux6570ux503cux4e0eux8ba1ux7b9701-ux6570ux503cux5206ux679005-ux53cdux95eeux9898ux4e0eux4e0dux9002ux5b9aux95eeux9898.md}

\chapter{反问题与不适定问题}\label{ux53cdux95eeux9898ux4e0eux4e0dux9002ux5b9aux95eeux9898}

反问题（Inverse Problem）是根据观测结果推断模型参数或系统输入的数学问题，广泛存在于地球物理、医学成像、无损检测和遥感等领域。与正问题（已知原因求结果）相反，反问题通常是''不适定的''（ill-posed）------Hadamard在1902年关于偏微分方程适定性的经典定义恰好刻画了反问题的核心困难：解的存在性、唯一性或连续依赖性中至少一项不成立。正则化理论为处理不适定反问题提供了系统的数学框架，其基本思想是用一簇适定的''邻近''问题逼近原不适定问题。本章从不适定性的一般理论出发，系统阐述线性与非线性正则化方法，并涵盖反散射、断层成像和Bayes反问题框架。

\hrulefill

\section{第400章：不适定问题与正则化理论}\label{ux7b2c400ux7ae0ux4e0dux9002ux5b9aux95eeux9898ux4e0eux6b63ux5219ux5316ux7406ux8bba}

Hadamard提出的适定性概念是数学物理问题的基本准则。适定性包括三个条件：解的存在性、唯一性和对数据的连续依赖性。反问题往往违反第三条------即解不连续依赖于数据，微小的观测噪声可导致解的剧烈震荡。正则化理论通过引入附加信息（如光滑性先验）来稳定求解过程，其核心是平衡数据拟合与模型复杂度。

\subsection{x.1 Hadamard适定性与不适定性}\label{x.1-hadamardux9002ux5b9aux6027ux4e0eux4e0dux9002ux5b9aux6027}

\textbf{定义}（Hadamard适定性）：设 \(X, Y\) 为赋范线性空间，\(K: X \to Y\) 为线性算子。方程 \(Kx = y\) 称为适定的（well-posed），若：（1）\(\forall y \in Y\)，\(\exists x \in X\) 使 \(Kx = y\)（存在性）；（2）解唯一（\(K\) 单射）；（3）解连续依赖于数据 \(y\)，即 \(K^{-1}: \operatorname{ran}(K) \to X\) 连续。若任一条件不成立，则问题称为不适定的（ill-posed）。

\textbf{经典示例}：第一类积分方程（Fredholm积分方程）

\[(Kx)(s) = \int_a^b k(s,t)x(t)\,dt = y(s), \quad s \in [c,d],\]

当 \(k \in C([c,d] \times [a,b])\) 时，算子 \(K: L^2([a,b]) \to L^2([c,d])\) 是紧算子。若 \(K\) 的无限维值域（必有此事），则 \(K^{-1}\) 无界------即反问题不适定。

\textbf{定理 x.1}（紧算子逆的无界性）：设 \(K: X \to Y\) 为无限维赋范空间之间的紧线性算子。若 \(K\) 的单射且有无限维值域，则 \(K^{-1}: \operatorname{ran}(K) \to X\) 无界（从而不连续）。

\textbf{证明}：假设 \(K^{-1}\) 有界。则对任意 \(x \in X\)，\(\|x\| \leq \|K^{-1}\|\,\|Kx\|\)。取 \(X\) 中任一有界序列 \(\{x_n\}\)，由 \(K\) 的紧性，\(\{Kx_n\}\) 有收敛子列。由 \(K^{-1}\) 的有界性假设，\(\{x_n\}\) 也有收敛子列。但这意味着 \(X\) 的单位球是紧的------仅当 \(\dim X < \infty\) 时成立。矛盾。故 \(K^{-1}\) 无界。\(\blacksquare\)

\subsection{x.2 紧算子的奇值分解与Picard准则}\label{x.2-ux7d27ux7b97ux5b50ux7684ux5947ux503cux5206ux89e3ux4e0epicardux51c6ux5219}

\textbf{定理 x.2}（紧算子的奇值分解，SVD）：设 \(K: X \to Y\) 为Hilbert空间之间的紧线性算子。则存在奇值系统 \(\{(\sigma_n, u_n, v_n)\}_{n=1}^\infty\)，其中 \(\sigma_1 \geq \sigma_2 \geq \cdots > 0\)，\(\{u_n\} \subset X\)，\(\{v_n\} \subset Y\) 为标准正交系，使得

\[Ku_n = \sigma_n v_n, \quad K^*v_n = \sigma_n u_n, \quad \forall n.\]

且对任意 \(x \in X\)，\(Kx = \sum_{n=1}^\infty \sigma_n \langle x, u_n\rangle_X \, v_n\)（在 \(Y\) 中收敛）。

\textbf{证明}：考虑自伴紧正算子 \(K^*K: X \to X\)。由紧自伴算子的谱定理，存在标准正交基 \(\{u_n\}\) 和正特征值 \(\{\sigma_n^2\}\) 使得 \(K^*Ku_n = \sigma_n^2 u_n\)（\(\sigma_n > 0\)）。定义 \(v_n = \frac{1}{\sigma_n}Ku_n \in Y\)。验证 \(\{v_n\}\) 为标准正交系：

\[\langle v_n, v_m\rangle_Y = \frac{1}{\sigma_n\sigma_m}\langle Ku_n, Ku_m\rangle = \frac{1}{\sigma_n\sigma_m}\langle K^*Ku_n, u_m\rangle = \frac{\sigma_n}{\sigma_m}\langle u_n, u_m\rangle = \delta_{nm}.\]

由 \(\operatorname{ran}(K)\) 中元素的稠密性论证可得 \(\{v_n\}\) 为 \(\overline{\operatorname{ran}(K)}\) 的标准正交基。SVD的收敛级数表示由Parseval恒等式推出。\(\blacksquare\)

\textbf{定理 x.3}（Picard准则）：设 \(K\) 有奇值系统 \(\{(\sigma_n, u_n, v_n)\}\)。方程 \(Kx = y\) 有解 \(x \in X\) 当且仅当

\[\sum_{n=1}^\infty \frac{|\langle y, v_n\rangle_Y|^2}{\sigma_n^2} < \infty,\]

且此时解为 \(x = \sum_{n=1}^\infty \frac{\langle y, v_n\rangle}{\sigma_n} u_n\)。

\textbf{证明}：（\(\Rightarrow\)）若 \(Kx = y\)，展开 \(x = \sum \langle x, u_n\rangle u_n\)（加上可能的核空间分量），则

\[y = Kx = \sum_{n=1}^\infty \sigma_n\langle x, u_n\rangle v_n,\]

故 \(\langle y, v_n\rangle = \sigma_n\langle x, u_n\rangle\)。由Parseval等式，\(\sum |\langle x, u_n\rangle|^2 = \sum |\langle y, v_n\rangle|^2/\sigma_n^2 \leq \|x\|^2 < \infty\)。

（\(\Leftarrow\)）若上述级数收敛，定义 \(x = \sum_{n=1}^\infty \frac{\langle y, v_n\rangle}{\sigma_n} u_n\)（加上核空间的正交补项，可不失一般性设 \(y \perp \ker K^*\)）。由收敛性，\(x \in X\) 且 \(Kx = y\)（因 \(K\) 作用于该级数收敛于 \(y\)）。\(\blacksquare\)

Picard准则揭示反问题不适定的本质：\(y\) 的系数 \(\langle y, v_n\rangle\) 的衰减必须快于 \(\sigma_n\) 的衰减才能保证解的存在性。但实际中，含噪声数据 \(y^\delta = y + \delta \eta\) 的Fourier系数 \(\langle y^\delta, v_n\rangle\) 仅以 \(O(1/n)\) 或更慢的速率衰减（噪声白化效应），而 \(\sigma_n \to 0\)，因此 \(\sum |\langle y^\delta, v_n\rangle|^2/\sigma_n^2 = \infty\) 以概率1成立------此即反问题对噪声极度敏感的数学根源。

\subsection{x.3 Tikhonov正则化的收敛性理论}\label{x.3-tikhonovux6b63ux5219ux5316ux7684ux6536ux655bux6027ux7406ux8bba}

\textbf{定理 x.4}（Tikhonov正则化的收敛性------先验参数选取）：设 \(K: X \to Y\) 为紧单射线性算子，\(y^\delta\) 满足 \(\|y^\delta - y\| \leq \delta\)，\(y = Kx^\dagger\)（\(x^\dagger\) 为真解）。Tikhonov正则化解

\[x_\alpha^\delta = \arg\min_{x \in X} \{\|Kx - y^\delta\|^2 + \alpha\|x\|^2\} = (K^*K + \alpha I)^{-1}K^*y^\delta.\]

若选取 \(\alpha = \alpha(\delta)\) 满足 \(\alpha(\delta) \to 0\) 且 \(\delta^2/\alpha(\delta) \to 0\) 当 \(\delta \to 0\)，则 \(x_{\alpha(\delta)}^\delta \to x^\dagger\) 在 \(X\) 中。

\textbf{证明}：分解误差：

\[\|x_\alpha^\delta - x^\dagger\| \leq \|x_\alpha^\delta - x_\alpha\| + \|x_\alpha - x^\dagger\|,\]

其中 \(x_\alpha = (K^*K + \alpha I)^{-1}K^*y\) 为无噪声Tikhonov解。

第一项（噪声传播）：由 \((K^*K + \alpha I)^{-1}K^*\) 的算子范数 \(\leq 1/(2\sqrt{\alpha})\)，

\[\|x_\alpha^\delta - x_\alpha\| = \|(K^*K + \alpha I)^{-1}K^*(y^\delta - y)\| \leq \frac{\delta}{2\sqrt{\alpha}}.\]

第二项（正则化偏差）：利用SVD，

\[x_\alpha - x^\dagger = [(K^*K + \alpha I)^{-1}K^*K - I]x^\dagger = \sum_{n=1}^\infty \left(\frac{\sigma_n^2}{\sigma_n^2 + \alpha} - 1\right)\langle x^\dagger, u_n\rangle u_n = -\sum_{n=1}^\infty \frac{\alpha}{\sigma_n^2 + \alpha}\langle x^\dagger, u_n\rangle u_n.\]

故 \(\|x_\alpha - x^\dagger\|^2 = \sum_{n=1}^\infty \frac{\alpha^2}{(\sigma_n^2 + \alpha)^2}|\langle x^\dagger, u_n\rangle|^2\)。由控制收敛定理，\(\lim_{\alpha \to 0} \|x_\alpha - x^\dagger\| = 0\)（因 \(x^\dagger \in X\) 且 \(\|x^\dagger\|^2 = \sum |\langle x^\dagger, u_n\rangle|^2 < \infty\)）。

由假设 \(\alpha \to 0\) 且 \(\delta/\sqrt{\alpha} \to 0\)，总误差 \(\|x_\alpha^\delta - x^\dagger\| \leq \delta/(2\sqrt{\alpha}) + \|x_\alpha - x^\dagger\| \to 0\)。\(\blacksquare\)

\textbf{定理 x.5}（Morozov偏差原理------后验参数选取）：设 \(\tau > 1\) 为固定常数。选取 \(\alpha = \alpha(\delta)\) 为满足偏差方程

\[\|Kx_\alpha^\delta - y^\delta\| = \tau\delta\]

的解（若存在）。则 \(\|x_{\alpha(\delta)}^\delta - x^\dagger\| \to 0\) 当 \(\delta \to 0\)，且该收敛关于方法自身具有最优阶（在源条件下）。

\textbf{证明}：首先验证偏差方程的解的存在性。定义函数 \(\varphi(\alpha) = \|Kx_\alpha^\delta - y^\delta\|^2 - \tau^2\delta^2\)。由SVD：

\[\varphi(\alpha) = \sum_{n=1}^\infty \frac{\alpha^2}{(\sigma_n^2 + \alpha)^2}|\langle y^\delta, v_n\rangle|^2 - \tau^2\delta^2.\]

当 \(\alpha \to 0^+\)，\(\varphi(\alpha) \to \sum |\langle y^\delta, v_n\rangle|^2 - \tau^2\delta^2\)不等。当 \(\alpha \to \infty\)，\(\varphi(\alpha) \to \|y^\delta\|^2 - \tau^2\delta^2\)。若 \(\|y^\delta\| > \tau\delta\) 且存在适当的中间条件，则 \(\varphi\) 有零点 \(\alpha(\delta)\)。

收敛性：设 \(x_{\alpha(\delta)}^\delta\) 为偏差原理选取的解。利用偏差方程 \(\|Kx_{\alpha(\delta)}^\delta - y^\delta\| = \tau\delta\)，结合

\[\|K(x_{\alpha(\delta)}^\delta - x^\dagger)\| \leq \|Kx_{\alpha(\delta)}^\delta - y^\delta\| + \|y^\delta - y\| \leq (\tau+1)\delta.\]

由正则化参数的性质和Tikhonov正则化的一般收敛性理论，\(x_{\alpha(\delta)}^\delta \to x^\dagger\)。在源条件 \(x^\dagger = (K^*K)^\nu w\)（\(0 < \nu \leq 1\)，\(\|w\| < \infty\)）下，可得最优收敛阶 \(\|x_{\alpha(\delta)}^\delta - x^\dagger\| = O(\delta^{2\nu/(2\nu+1)})\)。\(\blacksquare\)

\hrulefill

\section{第401章：线性反问题与非迭代方法}\label{ux7b2c401ux7ae0ux7ebfux6027ux53cdux95eeux9898ux4e0eux975eux8fedux4ee3ux65b9ux6cd5}

线性反问题构成了反问题理论的重要组成部分。反散射问题由Helmholtz方程在无穷远处的渐近行为（远场模式）刻画，Radon变换是计算机断层扫描（CT）的核心数学工具，而电阻抗成像的Calderón问题则联系起反问题与几何分析。本章展开这些经典线性反问题的数学理论。

\subsection{x+1.1 反散射问题}\label{x1.1-ux53cdux6563ux5c04ux95eeux9898}

\textbf{定义}：声波反散射问题考虑Helmholtz方程在无界域中的解。总场 \(u\) 满足

\[\Delta u + k^2 n(x)u = 0, \quad x \in \mathbb{R}^d,\]

其中 \(k > 0\) 为波数，折射率 \(n(x) - 1\) 具有紧支撑（散射体）。\(u = u^i + u^s\)，\(u^i\) 为入射平面波 \(e^{ikx\cdot\theta}\)（\(\theta \in S^{d-1}\)），\(u^s\) 满足Sommerfeld辐射条件：

\[\lim_{r \to \infty} r^{(d-1)/2}\left(\frac{\partial u^s}{\partial r} - iku^s\right) = 0, \quad r = |x|.\]

此时远场模式 \(u^\infty(\hat{x}, \theta)\)（\(\hat{x} = x/|x|\)）由渐近式 \(u^s(x) = \frac{e^{ikr}}{r^{(d-1)/2}}u^\infty(\hat{x}, \theta) + o(r^{-(d-1)/2})\) 定义。反散射问题即：给定 \(u^\infty\)，求 \(n(x)\)。

\textbf{定理 x+1.1}（远场算子的紧性）：在适当条件下，远场算子 \(F: n \mapsto u^\infty\) 是紧的，从而反问题不适定。进一步，\(F\) 的Fréchet导数为紧算子，其SVD提供了线性化反散射的正则化框架。

\subsection{x+1.2 Radon变换与断层扫描}\label{x1.2-radonux53d8ux6362ux4e0eux65adux5c42ux626bux63cf}

\textbf{定义}：二维Radon变换将函数 \(f: \mathbb{R}^2 \to \mathbb{R}\) 映为其沿直线的线积分：

\[(\mathcal{R}f)(\theta, s) = \int_{\{x \mid x\cdot\theta = s\}} f(x)\,d\ell(x) = \int_{-\infty}^\infty f(s\theta + t\theta^\perp)\,dt,\]

其中 \(\theta = (\cos\varphi, \sin\varphi)\) 为方向向量，\(s \in \mathbb{R}\) 为到原点的带号距离。

\textbf{定理 x+1.2}（Central Slice定理）：Radon变换与Fourier变换的关系为

\[\widehat{(\mathcal{R}f)}(\theta, \sigma) = \widehat{f}(\sigma\theta),\]

其中左侧为一维Fourier变换（对第二变量 \(s\)），右侧为二维Fourier变换。即Radon变换在方向 \(\theta\) 上的一维Fourier变换等于 \(f\) 的二维Fourier变换在射线 \(\sigma\theta\) 上的限制。

\textbf{证明}：由定义，

\[\widehat{(\mathcal{R}f)}(\theta, \sigma) = \int_{-\infty}^\infty (\mathcal{R}f)(\theta, s)\,e^{-2\pi i s\sigma}\,ds = \int_{-\infty}^\infty \int_{-\infty}^\infty f(s\theta + t\theta^\perp)\,e^{-2\pi i s\sigma}\,dt\,ds.\]

引入变量代换 \(x = s\theta + t\theta^\perp\)，即 \((s, t) = (x\cdot\theta, x\cdot\theta^\perp)\)。Jacobi行列式为1（旋转坐标）。则

\[\widehat{(\mathcal{R}f)}(\theta, \sigma) = \int_{\mathbb{R}^2} f(x)\,e^{-2\pi i (x\cdot\theta)\sigma}\,dx = \int_{\mathbb{R}^2} f(x)\,e^{-2\pi i x\cdot(\sigma\theta)}\,dx = \widehat{f}(\sigma\theta).\]

这个简洁的恒等式将Radon反演转化为Fourier反演。\(\blacksquare\)

\textbf{定理 x+1.3}（滤波反投影FBP公式）：\(f\) 可由Radon变换通过滤波反投影重建：

\[f(x) = \int_0^\pi (\mathcal{R}f(\theta, \cdot) * h)(x\cdot\theta)\,d\theta,\]

其中 \(h\) 是满足 \(\widehat{h}(\sigma) = |\sigma|\) 的斜坡滤波器（ramp filter），\(*\) 表示一维卷积。

\textbf{证明}：由Fourier反演公式和Central Slice定理：

\[f(x) = \int_{\mathbb{R}^2} \widehat{f}(\xi)\,e^{2\pi i x\cdot\xi}\,d\xi.\]

引入极坐标 \((\sigma, \theta)\)（即 \(\xi = \sigma\theta\)），\(d\xi = |\sigma|\,d\sigma\,d\theta\)（\(\theta \in [0,\pi)\)，\(\sigma \in \mathbb{R}\)）。则

\[f(x) = \int_0^\pi \int_{-\infty}^\infty \widehat{f}(\sigma\theta)\,e^{2\pi i \sigma(x\cdot\theta)}\,|\sigma|\,d\sigma\,d\theta\] \[= \int_0^\pi \int_{-\infty}^\infty \widehat{(\mathcal{R}f)}(\theta, \sigma)\,|\sigma|\,e^{2\pi i \sigma(x\cdot\theta)}\,d\sigma\,d\theta.\]

内积分恰为 \((\mathcal{R}f(\theta,\cdot) * h)(x\cdot\theta)\)（其中 \(\widehat{h}(\sigma) = |\sigma|\)）。在实践中，\(h\) 的频域截断给出实际的滤波反投影重建算法。此即现代CT扫描仪的数学基础。\(\blacksquare\)

\subsection{x+1.3 电阻抗成像的Calderón问题}\label{x1.3-ux7535ux963bux6297ux6210ux50cfux7684calderuxf3nux95eeux9898}

\textbf{定义}：Calderón问题（电阻抗成像EIT的数学模型）考虑边界测量确定内部电导率的问题。设 \(\Omega \subset \mathbb{R}^n\) 为有界区域，电导率 \(\gamma(x) > 0\)。电势 \(u\) 满足

\[\nabla \cdot (\gamma \nabla u) = 0, \quad x \in \Omega, \quad u|_{\partial\Omega} = f.\]

Dirichlet-to-Neumann (DtN)映射 \(\Lambda_\gamma: f \mapsto \gamma\frac{\partial u}{\partial\nu}\big|_{\partial\Omega}\)。Calderón问题为：由 \(\Lambda_\gamma\) 确定 \(\gamma\)。

\textbf{定理 x+1.4}（Calderón唯一性问题）：在 \(\dim \geq 3\) 时，若 \(\gamma \in C^2(\overline{\Omega})\)，则 \(\Lambda_\gamma\) 唯一确定 \(\gamma\)。在 \(\dim = 2\) 时，唯一性在共形等价意义下成立。

\textbf{证明要点}：（\(\dim \geq 3\)情形）Sylvester和Uhlmann（1987）的证明利用了复几何光学（Complex Geometrical Optics, CGO）解。构造指数增长解 \(u = e^{\rho\cdot x}(1 + \psi_\rho)\)，其中 \(\rho \in \mathbb{C}^n\) 满足 \(\rho \cdot \rho = 0\)（注意 \(\rho\) 为复向量）。代入方程得 \(\psi_\rho\) 满足带奇异势的方程，可用适当的函数空间理论求解。通过选取合适的 \(\rho\)，由 \(\Lambda_{\gamma_1} = \Lambda_{\gamma_2}\) 可推得 \(\gamma_1\) 和 \(\gamma_2\) 的Fourier变换在适当的加权空间中相等，从而唯一性。二维情形因复结构而不同：DtN映射在共形变换下不变，故仅能确定共形等价类中的电导率。\(\blacksquare\)

\hrulefill

\section{第402章：非线性反问题与迭代正则化}\label{ux7b2c402ux7ae0ux975eux7ebfux6027ux53cdux95eeux9898ux4e0eux8fedux4ee3ux6b63ux5219ux5316}

许多实际反问题是非线性的，这给正则化带来了额外的挑战。迭代方法通过在达到合适的数据拟合水平时提前终止来隐式地实现正则化。Landweber迭代是最经典的迭代正则化方法，Gauss-Newton法适用于非线性算子方程，而Bayes框架则为反问题提供了概率解释。

\subsection{x+2.1 Landweber迭代的正则化性质}\label{x2.1-landweberux8fedux4ee3ux7684ux6b63ux5219ux5316ux6027ux8d28}

\textbf{定义}：对于线性方程 \(Kx = y\)，Landweber迭代定义为

\[x_{k+1} = x_k + \omega K^*(y - Kx_k), \quad x_0 = 0,\]

其中 \(0 < \omega < 2/\|K\|^2\) 为松弛参数。迭代次数 \(k\) 充当正则化参数的角色。

\textbf{定理 x+2.1}（Landweber迭代的收敛性）：设 \(y \in \operatorname{ran}(K)\)，以 \(x^\dagger = K^\dagger y\) 记最小范数解。则 \(x_k \to x^\dagger\) 当 \(k \to \infty\)。

\textbf{证明}：引入误差 \(e_k = x_k - x^\dagger\)。迭代方程为 \(e_{k+1} = (I - \omega K^*K)e_k\)。由归纳法，\(e_k = (I - \omega K^*K)^k e_0 = -(I - \omega K^*K)^k x^\dagger\)（因 \(x_0 = 0\)）。

由SVD分析：

\[\|e_k\|^2 = \sum_{n=1}^\infty (1 - \omega\sigma_n^2)^{2k} |\langle x^\dagger, u_n\rangle|^2.\]

由于 \(0 < \omega < 2/\sigma_1^2\)，各项 \((1 - \omega\sigma_n^2) \in (-1, 1)\)，且当 \(\sigma_n \to 0\) 时该项趋于1，即高频分量衰减缓慢------这正是Landweber迭代的半收敛（semi-convergence）特性。

由控制收敛定理，\(\lim_{k \to \infty} \|e_k\| = 0\)（对每个固定 \(n\)，\((1 - \omega\sigma_n^2)^{2k} \to 0\) 当 \(k \to \infty\)）。在实际含噪情形，需在适当的 \(k(\delta)\) 处停止，平衡迭代误差与噪声放大。\(\blacksquare\)

\textbf{定理 x+2.2}（Landweber迭代的正则化------噪声情形）：设 \(\|y^\delta - y\| \leq \delta\)。若选取停止指标 \(k = k(\delta)\) 使得 \(\|Kx_{k(\delta)}^\delta - y^\delta\| \leq \tau\delta < \|Kx_{k(\delta)-1}^\delta - y^\delta\|\)（\(\tau > 1\)，偏差型停止准则），则 \(x_{k(\delta)}^\delta \to x^\dagger\) 当 \(\delta \to 0\)。

\textbf{证明}：关键在于证明 \(k(\delta)\) 随 \(\delta \to 0\) 趋于无穷，且误差在停止时被控制。利用残差的单调递减性质：\(\|Kx_{k+1}^\delta - y^\delta\| < \|Kx_k^\delta - y^\delta\|\)。在停止时：

\[\|Kx_{k(\delta)}^\delta - y^\delta\| \leq \tau\delta.\]

同时由停止准则的下界，\(\tau\delta < \|Kx_{k(\delta)-1}^\delta - y^\delta\|\)，这意味着 \(k(\delta)\) 是最小的满足 \(\leq \tau\delta\) 的指标，从而随 \(\delta \to 0\) 发散。总误差：

\[\|x_{k(\delta)}^\delta - x^\dagger\| \leq \|x_{k(\delta)}^\delta - x_{k(\delta)}\| + \|x_{k(\delta)} - x^\dagger\|.\]

第二项由无噪声情形的收敛性趋于零。第一项由Landweber迭代的Lipschitz稳定性：\(\|x_k^\delta - x_k\| \leq \text{const} \cdot \sqrt{k}\,\delta\)。由 \(k(\delta)\) 的发散速度（通常 \(k(\delta) = O(\delta^{-2/(2\nu+1)})\) 在源条件下），可证第一项也趋于零。\(\blacksquare\)

\subsection{x+2.2 非线性反问题中的迭代方法}\label{x2.2-ux975eux7ebfux6027ux53cdux95eeux9898ux4e2dux7684ux8fedux4ee3ux65b9ux6cd5}

对于非线性算子方程 \(F(x) = y\)（\(F: X \to Y\) 为Fréchet可微），常用的迭代正则化方法包括：

\textbf{正则化的Gauss-Newton法（IRGNM）}：第 \(k\) 步迭代：

\[x_{k+1} = x_k - (F'(x_k)^*F'(x_k) + \alpha_k I)^{-1}F'(x_k)^*(F(x_k) - y^\delta).\]

每一步求解一个线性化的Tikhonov问题。\(\alpha_k \to 0\) 为正则化参数序列。

\textbf{Levenberg-Marquardt方法}：

\[x_{k+1} = x_k - (F'(x_k)^*F'(x_k) + \lambda_k I)^{-1}F'(x_k)^*(F(x_k) - y^\delta),\]

其中 \(\lambda_k\) 的选取满足使 \(\|F(x_{k+1}) - y^\delta\| < \|F(x_k) - y^\delta\|\)（充分下降条件）且 \(\|x_{k+1} - x_k\| \leq C\lambda_k^{-1/2}\|F(x_k) - y^\delta\|\)。

\textbf{定理 x+2.3}（IRGNM的收敛性）：设 \(F\) 满足非线性条件（切锥条件 \(F(x+h) - F(x) - F'(x)h = O(\|h\|^2)\) 或适当的光滑性与弱闭性），\(\{\alpha_k\}\) 适当递减（如 \(\alpha_{k+1} = q\alpha_k\)，\(0 < q < 1\)）。则含噪声的迭代序列 \(\{x_k^\delta\}\) 在适当停止准则下收敛到真解。

\textbf{证明要点}：核心是建立每个迭代步的误差减少不等式。在适当的源条件和非线性条件下可证：

\[\|x_{k+1}^\delta - x^\dagger\| \leq c(\|x_k^\delta - x^\dagger\|^2 + \alpha_k + \delta/\sqrt{\alpha_k}).\]

通过精细的参数选取，确保迭代的收敛性和正则化性质。\(\blacksquare\)

\subsection{x+2.3 Bayes反问题框架}\label{x2.3-bayesux53cdux95eeux9898ux6846ux67b6}

\textbf{定义}：在Bayes框架下，反问题表述为：给定观测数据 \(y\) 和模型 \(y = F(x) + \eta\)（\(\eta\) 为随机噪声），利用Bayes定理从先验分布 \(\pi(x)\) 和似然 \(\pi(y|x)\) 计算后验分布 \(\pi(x|y)\)。

\textbf{定理 x+2.4}（Bayes定理在无限维中的推广）：设 \(\pi_0\) 为可分Hilbert空间 \(X\) 上的Gaussian测度（先验），观测模型 \(y = Kx + \eta\)，\(\eta \sim \mathcal{N}(0, \Gamma)\) 为Gaussian噪声，\(\Gamma\) 为对称正定迹类算子。则后验测度 \(\pi^y\) 在 \(X\) 上绝对连续于先验 \(\pi_0\)，Radon-Nikodym导数为

\[\frac{d\pi^y}{d\pi_0}(x) \propto \exp\left(-\frac{1}{2}\|\Gamma^{-1/2}(y - Kx)\|^2\right).\]

后验测度的最大值（MAP估计）对应于Tikhonov正则化解：\(x_{\text{MAP}} = \arg\min_x\{\|y - Kx\|_\Gamma^2 + \|x\|_{C_0}^2\}\)，其中 \(C_0\) 为先验的协方差算子。

\textbf{证明}：由Bayes定理 \(\pi^y(dx) \propto \pi(y|x)\pi_0(dx)\)。Gaussian似然密度（关于 \(\pi_0\)）为 \(\pi(y|x) \propto \exp(-\frac{1}{2}\|\Gamma^{-1/2}(y - Kx)\|^2)\)。由此Radon-Nikodym导数即为该指数函数（乘常数因子）。

\(\pi^y\) 的绝对连续性来自Cameron-Martin定理的推广（Feldman-Hajek定理）：两个Gaussian测度在可分Hilbert空间上等价当且仅当它们的协方差算子差异是Hilbert-Schmidt的。先验 \(\pi_0 = \mathcal{N}(0, C_0)\) 和后验在前述条件下绝对连续。MAP估计的等价性通过计算后验的对数密度并求极大得到。\(\blacksquare\)

Bayes框架的优势在于：它天然提供了解的不确定性量化（通过后验分布的宽度或可信区间），且可用MCMC（马尔可夫链蒙特卡洛）等采样方法数值实现，适用于高度非线性和高维反问题。

\subsection{x+2.4 梯度型迭代法的正则化半收敛性质}\label{x2.4-ux68afux5ea6ux578bux8fedux4ee3ux6cd5ux7684ux6b63ux5219ux5316ux534aux6536ux655bux6027ux8d28}

Landweber迭代和共轭梯度法（CG）在半收敛（semi-convergence）方面展现出相似的行为：迭代误差先减小后增大，存在最优停止时间。

\textbf{定理 x+2.4}（CG方法的半收敛性）：对于线性不适定问题 \(Kx = y\)，共轭梯度法（CGLS）应用于法方程 \(K^*Kx = K^*y\) 产生的迭代序列 \(\{x_k\}\) 满足：\(\|x_k - K^\dagger y\|\) 先减小后增大（当 \(k \to \infty\) 时发散，若数据含噪声）。最优停止步数由偏差原理或广义交叉验证（GCV）确定。

\textbf{证明}：CGLS的 Krylov 子空间性质：\(x_k\) 在 \(\mathcal{K}_k(K^*K, K^*y)\) 中极小化残差。由 \(K\) 的紧性和 Picard 准则，\(K^\dagger y\) 可由奇向量 \(\{u_n\}\) 刻画。CGLS 有效地滤除高频噪声分量，但迭代步数过多时，高频噪声被放大。利用CGLS的多项式滤波特性------\(x_k = P_k(K^*K)K^*y\)，其中 \(P_k\) 为第 \(k\) 次 Ritz 多项式------可解析地证明 \(\|x_k - K^\dagger y\|\) 的U形行为。最优停止的偏差原理保证 \(\|Kx_k^\delta - y^\delta\| \leq \tau\delta\) 时的 \(k\) 提供了适当的正则化。\(\blacksquare\)

\subsection{x+2.5 全变差正则化与反问题}\label{x2.5-ux5168ux53d8ux5deeux6b63ux5219ux5316ux4e0eux53cdux95eeux9898}

Tikhonov正则化光滑解、模糊边缘。全变差（TV）正则化克服此缺陷，特别适用于分段常数的解（如医学成像中的组织边界）。

\textbf{定理 x+2.5}（TV正则化的适定性与收敛性）：设 \(K: L^2(\Omega) \to Y\) 为有界线性算子。TV正则化解为

\[x_\alpha^\delta = \arg\min_{x \in BV(\Omega)} \left\{\|Kx - y^\delta\|_Y^2 + \alpha \operatorname{TV}(x)\right\}.\]

若 \(K\) 单射且 \(\operatorname{ran}(K)\) 闭（对带限算子需适当修改），则 \(x_\alpha^\delta\) 存在。在源条件 \(x^\dagger \in BV(\Omega)\) 下，取 \(\alpha(\delta) \sim \delta\) 时有收敛阶 \(\|x_{\alpha(\delta)}^\delta - x^\dagger\|_{L^1} = O(\delta)\)。

\textbf{证明}：存在性由BV空间的紧嵌入性质和TV泛函的下半连续性给出。Convergence rates use Bregman distances associated with the TV functional. The key estimate:

\[\|x_\alpha^\delta - x^\dagger\|_{L^1} \leq C\left(\frac{\delta}{\sqrt{\alpha}} + \sqrt{\alpha} \operatorname{TV}(x^\dagger)\right).\]

Balancing \(\alpha \sim \delta\) gives \(O(\delta)\) rate. More precisely, one uses the Fenchel dual and the source condition in the subgradient of TV.\(\blacksquare\)

\subsection{\texorpdfstring{x+2.6 稀疏约束反问题与\(\ell_1\)正则化}{x+2.6 稀疏约束反问题与\textbackslash ell\_1正则化}}\label{x2.6-ux7a00ux758fux7ea6ux675fux53cdux95eeux9898ux4e0eell_1ux6b63ux5219ux5316}

\(L^1\) 正则化在解具有稀疏表示的反问题中特别有效。考虑模型 \(y = Kx + \eta\)，其中 \(x\) 在适当的基（小波、框架等）下是稀疏的。

\textbf{定理 x+2.6}（\(\ell_1\) 正则化反问题的误差估计）：设真解 \(x^\dagger\) 满足源条件（存在 \(w \in Y\) 使 \(K^*w \in \partial\|x^\dagger\|_1\)），噪声水平 \(\|y^\delta - y\| \leq \delta\)。则 \(\ell_1\)-Tikhonov解 \(\hat{x}_\alpha^\delta = \arg\min_x\{\|Kx - y^\delta\|^2 + \alpha\|x\|_1\}\)（\(\alpha \sim \delta\)）满足

\[\|\hat{x}_\alpha^\delta - x^\dagger\| \leq C\sqrt{\delta} \quad (\text{在 } \ell_2 \text{ 范数下}).\]

\textbf{证明}：利用Moreau包络和邻近算子的非扩张性质。在源条件下，\(\ell_1\) 正则化等价于Bregman迭代的极限，其收敛率由邻近梯度法的理论得出。关键步骤是证明近似解与真解在支撑集上的一致性（oracle property），继而误差仅由噪声主导。\(\blacksquare\)

\subsection{x+2.7 分数阶反问题与扩散光学层析}\label{x2.7-ux5206ux6570ux9636ux53cdux95eeux9898ux4e0eux6269ux6563ux5149ux5b66ux5c42ux6790}

许多应用中的反问题涉及分数阶微分算子，这类问题的不适定性通常更加严重。以分数扩散方程的反源问题为例：已知终值 \(u(T, x)\)，求初值 \(u(0, x)\) 或源项。

\textbf{定理 x+2.7}（分数扩散反问题的对数型稳定性）：对时间分数扩散方程 \(\partial_t^\alpha u = \Delta u\)（\(0 < \alpha < 1\)，\(\partial_t^\alpha\) 为 Caputo 分数导数），反初值问题的稳定性估计为对数型：

\[\|u(0,\cdot) - v(0,\cdot)\|_{L^2} \leq C\left(\log\frac{1}{\|u(T,\cdot) - v(T,\cdot)\|}\right)^{-\gamma},\]

指数 \(\gamma = \alpha/(1-\alpha)\) 依赖于分数阶数。与经典热传导（\(\alpha=1\)）相比，分数阶情形的不适定性显著加剧。

\textbf{证明}：利用 Mittag-Leffler 函数表示的分离变量解。在特征函数基下，算子半群作用为 \(u(T) = \sum a_n E_\alpha(-\lambda_n T^\alpha)\phi_n\)。演化因子的衰减为 \(E_\alpha(-\lambda T^\alpha) \sim (\lambda T^\alpha)^{-1}/\Gamma(1-\alpha)\)（当 \(\lambda\) 大时），远慢于经典情形的指数衰减 \(e^{-\lambda T}\)。这导致高频分量的衰减极慢，反问题因此仅有对数型稳定性。\(\blacksquare\)

\subsection{x+2.8 反问题的唯一性和识别性}\label{x2.8-ux53cdux95eeux9898ux7684ux552fux4e00ux6027ux548cux8bc6ux522bux6027}

反散射和Calderón型问题属于识别性问题，唯一性是核心。

\textbf{定义}：Dirichlet-to-Neumann映射 \(\Lambda_\gamma: H^{1/2}(\partial\Omega) \to H^{-1/2}(\partial\Omega)\) 定义为 \(\Lambda_\gamma f = \gamma\frac{\partial u}{\partial\nu}|_{\partial\Omega}\)，其中 \(u\) 满足 \(\nabla\cdot(\gamma\nabla u) = 0\) 且 \(u|_{\partial\Omega} = f\)。

\textbf{定理 x+2.8}（电导率确定的边界唯一性------Kohn-Vogelius不等式）：设 \(\gamma_1, \gamma_2\) 为 \(\Omega\) 上两个各向同性电导率。若 \(\Lambda_{\gamma_1} = \Lambda_{\gamma_2}\)，则 \(\gamma_1 = \gamma_2\) 在 \(\partial\Omega\) 的邻域内。

\textbf{证明}（Kohn-Vogelius）：对任意边界数据 \(f\)，两电导率对应的解分别记为 \(u_1, u_2\)。由 \(\Lambda_{\gamma_1} = \Lambda_{\gamma_2}\) 得 \(\gamma_1\nabla u_1\cdot\nu = \gamma_2\nabla u_2\cdot\nu\) 在 \(\partial\Omega\) 上。令 \(w = u_1 - u_2\)。则 \(\nabla\cdot(\gamma_1\nabla w) = \nabla\cdot((\gamma_2 - \gamma_1)\nabla u_2)\)。运用能量估计，由边界条件的匹配性可递推地唯一确定 \(\gamma_1\) 和 \(\gamma_2\) 在边界处的各阶法向导数。当边界为解析时，由Cauchy-Kovalevskaya定理得全局唯一性。对于一般光滑情形，利用Calderón的唯一延拓性质（由复几何光学解的稠密性）可得全局唯一性。\(\blacksquare\)

\subsection{x+2.9 反源问题与Green函数的辨识}\label{x2.9-ux53cdux6e90ux95eeux9898ux4e0egreenux51fdux6570ux7684ux8fa8ux8bc6}

反源问题要求根据边界或内部的有限观测恢复未知源项，其在环境监测和生物医学成像中广泛应用。

\textbf{定理 x+2.9}（椭圆型反源问题的对数稳定性）：考虑方程 \(-\Delta u + qu = f\)（\(x \in \Omega\)），\(u|_{\partial\Omega} = 0\)。已知边界的Neumann数据 \(\partial u/\partial\nu|_{\partial\Omega}\)，求源 \(f\)。此反问题仅有对数型稳定性：当 \(f \in H^s\) 有界时，

\[\|f_1 - f_2\|_{L^2} \leq C\left(\log\frac{1}{\|\partial u_1/\partial\nu - \partial u_2/\partial\nu\|_{H^{-1/2}}}\right)^{-2s/(2s+n)}.\]

\textbf{证明}：利用谱分解和正则化解析。设 \(\{\lambda_k, \phi_k\}\) 为 \(-\Delta + qI\) 在Dirichlet条件下的特征对。观测数据为 \(b_k = \langle\partial u/\partial\nu, \phi_k|_{\partial\Omega}\rangle\)（通过Green恒等式与 \(f\) 的系数 \(a_k\) 关联）。由特征值的Weyl渐近律和源的正则性条件，高频分量的信息衰减为代数型：\(|b_k| \sim \lambda_k^{-1}|a_k|\)。正则化解的截断误差为 \(\|f^{\text{true}} - f^{\text{reg}}\|^2 = \sum_{k>K} a_k^2 + \sum_{k \leq K} (b_k^\delta - b_k)^2/\lambda_k^2\)。选取截断 \(K\) 平衡截断误差和噪声传播，由特征值的增长 \(\lambda_k \sim k^{2/n}\) 导出对数型稳定性。\(\blacksquare\)

\newpage

\chapter{11-数值与计算/02-计算理论/01-数理逻辑.md}\label{ux6570ux503cux4e0eux8ba1ux7b9702-ux8ba1ux7b97ux7406ux8bba01-ux6570ux7406ux903bux8f91.md}

\subsection{328.A 数理逻辑补充}\label{a-ux6570ux7406ux903bux8f91ux8865ux5145}

数理逻辑（Mathematical Logic）是研究数学推理本身的形式化科学------它不满足于''做数学''，而是要回答''什么是数学证明''这一基础问题。这一学科的思想发源于 Leibniz 的''通用表意文字''（Characteristica Universalis）梦想（17 世纪），经 Boole 的命题代数（1854）和 Frege 的谓词逻辑（1879）的奠基，在 20 世纪达到了辉煌的巅峰：Gödel 完备性定理（1929）证明了一阶逻辑的语义与语法等价，Gödel 不完全性定理（1931）粉碎了 Hilbert 纲领中数学完全性与一致性的梦想，Turing 的可计算性理论（1936）为计算机科学提供了理论基石，Cohen 的力迫法（1963）证明连续统假设独立于 ZFC。本卷从一阶逻辑的语法和语义出发，系统构建数理逻辑的核心支柱：完备性定理的 Henkin 构造、模型论的量词消去与稳定性理论、可计算性的 Turing 机与 Church-Turing 论题、不完全性定理的 Gödel 编码与对角线引理，并深入集合论的序数-基数算术和力迫法。这些看似高度抽象的结果------一个公式的存在与否、一个集合的大小------对 20 世纪整个数学的面貌产生了深远的影响，定义了''数学基础''本身的边界。

\hrulefill

\hrulefill

\hrulefill

\hrulefill

\hrulefill

\hrulefill

\section{第403章：一阶逻辑完备性}\label{ux7b2c403ux7ae0ux4e00ux9636ux903bux8f91ux5b8cux5907ux6027}

一阶逻辑的完备性定理是 20 世纪数理逻辑最璀璨的基石之一。一阶逻辑（Predicate Logic of First Order）由 Gottlob Frege 在《概念文字》（1879）中首次引入，经 Peano、Russell-Whitehead（《数学原理》1910-1913）的完善，最终在 David Hilbert 和 Wilhelm Ackermann（1928）的教材《理论逻辑基础》中获得了现代形式化。完备性定理断言：如果一个一阶逻辑公式在每一个模型中都为真（语义上有效），则它可以由逻辑公理（命题重言式、等词公理、量词规则）形式地推导出来（语法上可证）。Kurt Gödel 在他 1929 年的博士论文中首次证明了这一等价性，其证明的核心思想是''模型存在引理''：任何一致的公式集都有一个模型。Gödel 的构造过程------用 Henkin 常数消除存在量词、构造极大一致集、由等价物项（term）的等价类构造项模型（term model）------被称为 Henkin-Hasenjaeger 方法（Leon Henkin 1949 年给出了更简练的证明）。完备性定理的推论与逆否命题同等重要：紧致性定理（若 \(\Gamma\) 的每个有限子集有模型，则 \(\Gamma\) 有模型）是完备性定理的简单且强大的推论，为 Robinson 的非标准分析（1961）和代数闭域的可公理化理论提供了严格的逻辑基础。本章从语法（形式语言、证明系统）和语义（结构、真值）两方面建立一阶逻辑，并证明 Gödel 完备性定理：语义真蕴含语法可证。

\subsection{99.1 一阶语言}\label{ux4e00ux9636ux8bedux8a00}

\textbf{定义 99.1}（一阶语言）：一个\textbf{一阶语言} \(\mathcal{L}\) 由以下符号组成： - \textbf{逻辑符号}：变元 \(x_1, x_2, \ldots\)；连接词 \(\neg, \land, \lor, \to\)；量词 \(\forall, \exists\)；等号 \(=\)（通常包含） - \textbf{非逻辑符号}：常数符号 \(\{c_i\}\)；函数符号 \(\{f_j\}\)（带元数）；关系符号 \(\{R_k\}\)（带元数）

\textbf{定义 99.2}（项与公式）：\(\mathcal{L}\) 的\textbf{项}归纳定义：变元和常数是项；若 \(f\) 是 \(n\) 元函数符号，\(t_1, \ldots, t_n\) 是项，则 \(f(t_1, \ldots, t_n)\) 是项。

\textbf{原子公式}为 \(R(t_1, \ldots, t_n)\)（\(R\) 是 \(n\) 元关系符号）或 \(t_1 = t_2\)。\textbf{公式}由原子公式通过连接词和量词递归构造。

\textbf{定义 99.3}（句子与自由变元）：不含自由变元的公式称为\textbf{句子}。一个\textbf{理论} \(T\) 是一组 \(\mathcal{L}\)-句子的集合。

\subsection{99.2 语义}\label{ux8bedux4e49}

\textbf{定义 99.4}（\(\mathcal{L}\)-结构）：\(\mathcal{L}\) 的一个\textbf{结构}（或模型）\(\mathfrak{A}\) 由以下组成： - 一个非空集合 \(A\)（论域） - 对每个常数符号 \(c\)，\(\mathfrak{A}\) 中一个元素 \(c^{\mathfrak{A}} \in A\) - 对每个 \(n\) 元函数符号 \(f\)，一个函数 \(f^{\mathfrak{A}}: A^n \to A\) - 对每个 \(n\) 元关系符号 \(R\)，一个关系 \(R^{\mathfrak{A}} \subseteq A^n\)

\textbf{定义 99.5}（真值，Tarski 真定义）：对 \(\mathcal{L}\)-结构 \(\mathfrak{A}\) 和句子 \(\varphi\)，\(\mathfrak{A} \models \varphi\)（\(\varphi\) 在 \(\mathfrak{A}\) 中为真）由 Tarski 真定义递归定义： - \(\mathfrak{A} \models R(t_1, \ldots, t_n)\) 当且仅当 \((t_1^{\mathfrak{A}}, \ldots, t_n^{\mathfrak{A}}) \in R^{\mathfrak{A}}\) - \(\mathfrak{A} \models \neg \varphi\) 当且仅当 \(\mathfrak{A} \not\models \varphi\) - \(\mathfrak{A} \models \varphi \land \psi\) 当且仅当 \(\mathfrak{A} \models \varphi\) 且 \(\mathfrak{A} \models \psi\) - \(\mathfrak{A} \models \forall x \varphi(x)\) 当且仅当对每个 \(a \in A\)，\(\mathfrak{A} \models \varphi(a)\)

\textbf{定义 99.6}（逻辑蕴含）：\(T \models \varphi\) 表示对每个 \(\mathcal{L}\)-结构 \(\mathfrak{A}\)，若 \(\mathfrak{A} \models T\)（即 \(\mathfrak{A}\) 满足 \(T\) 中所有句子），则 \(\mathfrak{A} \models \varphi\)。

\subsection{99.3 证明系统}\label{ux8bc1ux660eux7cfbux7edf}

\textbf{定义 99.7}（Hilbert 式证明系统）：一阶逻辑的公理包括： 1. 命题逻辑公理（重言式） 2. \(\forall x \varphi(x) \to \varphi(t)\)（\(t\) 对 \(\varphi\) 中 \(x\) 可代入） 3. \(\forall x (\varphi \to \psi) \to (\forall x \varphi \to \forall x \psi)\) 4. \(\varphi \to \forall x \varphi\)（\(x\) 不在 \(\varphi\) 中自由） 5. 等号公理：\(x = x\)，\(x = y \to y = x\)，\((x = y \land y = z) \to x = z\)，以及函数和关系的代换公理

推理规则：\textbf{Modus Ponens}（从 \(\varphi\) 和 \(\varphi \to \psi\) 推出 \(\psi\)）和 \textbf{全称概括}（Generalization：从 \(\varphi\) 推出 \(\forall x \varphi\)，要求 \(x\) 不在前提中自由出现）。

\textbf{定义 99.8}（可证性）：\(T \vdash \varphi\) 表示存在从 \(T\) 中公理出发、使用推理规则得到的 \(\varphi\) 的有限证明序列。

\subsection{99.4 Gödel 完备性定理}\label{guxf6del-ux5b8cux5907ux6027ux5b9aux7406}

\textbf{定理 99.1}（可靠性定理）：若 \(T \vdash \varphi\)，则 \(T \models \varphi\)。

\emph{证明}：对证明长度归纳。公理在任意结构中为真，推理规则保持真值。∎

\textbf{定理 99.2}（Gödel 完备性定理，1929）：若 \(T \models \varphi\)，则 \(T \vdash \varphi\)。等价地，每个一致的理论（\(T \nvdash \bot\)）有模型。

\textbf{证明}：由一阶逻辑的完备性定理（Gödel 1930），Henkin构造法：对一致理论\(T\)，通过添加Henkin常量和极大化，构造项模型\(\mathcal{M}\)。\(\mathcal{M}\)的元素为闭项，谓词解释由\(T\)中的原子公式决定。归纳证明\(\mathcal{M}\models\varphi\iff T\vdash\varphi\)，从而\(T\)一致必可满足。\(\blacksquare\)

\textbf{推论 99.3}（紧致性定理）：若 \(T\) 的每个有限子集有模型，则 \(T\) 有模型。

\emph{证明}：若 \(T\) 无模型，由完备性定理，\(T \vdash \bot\)。证明是有限的，仅用到 \(T\) 的有限子集 \(T_0\)。故 \(T_0 \vdash \bot\)，由可靠性，\(T_0\) 无模型，矛盾。∎

\textbf{推论 99.4}（Löwenheim-Skolem 定理，向下）：若可数语言 \(\mathcal{L}\) 的理论 \(T\) 有无限模型，则 \(T\) 有任意无限基数 \(\kappa \geq \aleph_0\) 的模型（向上 Löwenheim-Skolem：若 \(T\) 有无限模型，则 \(T\) 有任意基数 \(\kappa \geq |\mathcal{L}|\) 的模型）。

Löwenheim-Skolem 定理是模型论中最令人惊讶的元定理之一，它揭示了形式语言的表达能力与其语义之间的深刻张力。向下 Löwenheim-Skolem 定理（Löwenheim 1915 / Skolem 1920）表明：任何可数语言中的一致理论都有可数模型。向上 Löwenheim-Skolem 定理（Tarski 1920s / Mal'tsev 1936）则表明：若理论有无限模型，则对任意基数 \(\kappa \geq |\mathcal{L}|\)，理论都有基数至少为 \(\kappa\) 的模型。这两个定理合在一起意味着：一阶逻辑无法刻画''可数性''或''不可数性''------任何一阶理论（只要它有无穷模型）都既有可数模型又有任意大的不可数模型。这导致了著名的 \textbf{Skolem 悖论}：ZFC 集合论（可数语言中的一阶理论）若一致，则有一个可数模型，但在此可数模型中却可以证明''存在不可数集合''（Cantor 定理在该模型内部成立）。悖论得以消解的关键在于：在可数模型中，``不可数''的含义是''不存在模型内部的函数将该集合映射到 \(\mathbb{N}\)``，而非''该集合在外部意义下不可数''。从外部看，该集合确实是可数的。这一现象揭示了形式化数学中''内部''与''外部''视角的根本区别，也是 Gödel 不完全性定理和 Cohen 力迫法的重要前奏。

\emph{证明}（向上）：添加 \(\kappa\) 个新常元 \(\{c_\alpha\}_{\alpha < \kappa}\) 和公理 \(c_\alpha \neq c_\beta\)（\(\alpha \neq \beta\)）。每个有限子集有模型（原无限模型可解释有限个新常元为不同元素）。由紧致性，整个理论有模型，其基数 \(\geq \kappa\)。∎

\hrulefill

\hrulefill

\hrulefill

\hrulefill

\hrulefill

\hrulefill

\section{第404章：模型论基础}\label{ux7b2c404ux7ae0ux6a21ux578bux8bbaux57faux7840}

模型论（Model Theory）研究数学结构与其形式化理论（一阶逻辑公式的集合）之间的关系，它将逻辑从''文字游戏''提升为真正能够研究代数、几何和数论具体对象的数学分支。这一转变的里程碑是 Abraham Robinson 关于代数闭域的可公理化理论和非标准分析的发现。模型论的核心概念------初等等价（两个结构满足完全相同的一阶语句）、初等嵌入（保持所有一阶公式真值不变的单射）和完全理论（理论的每个语句或其否定被理论唯一确定）------构成了连接纯逻辑与具体数学对象的桥梁。模型论的代数视角体现在量词消去（Quantifier Elimination）概念上：一个理论具有量词消去当且仅当每个一阶公式都等价于一个无量词公式------这使得模型的可定义集合具有简单的几何描述。Tarski（1931）首次证明了实闭域 \((\mathbb{R};+,\cdot,<)\) 具有量词消去------所有可定义集合都是半代数集（多项式等式和不等式的有限布尔组合）。进而 Tarski 由此得出实闭域的一阶理论是完全的、可判定的，且 \(\mathbb{R}\) 和 \(\mathbb{C}\) 的真初等性质都具有量词消去------这一系列结果是模型论应用于代数几何的经典范例。模型论的稳定性理论（Shelah 1969-1978）通过精细地分类一阶理论中的''可定义集''的复杂程度，产生了一套令人惊叹的普遍分类纲领，在代数几何（Zariski 几何）和数论（差分域）中都产生了深远的影响。

\subsection{100.1 初等等价与初等扩张}\label{ux521dux7b49ux7b49ux4ef7ux4e0eux521dux7b49ux6269ux5f20}

\textbf{定义 100.1}（初等等价）：两个 \(\mathcal{L}\)-结构 \(\mathfrak{A}\) 和 \(\mathfrak{B}\) 称为\textbf{初等等价}的（记作 \(\mathfrak{A} \equiv \mathfrak{B}\)），如果它们满足相同的 \(\mathcal{L}\)-句子。

\textbf{定义 100.2}（初等子模型）：\(\mathfrak{A} \preceq \mathfrak{B}\)（\(\mathfrak{A}\) 是 \(\mathfrak{B}\) 的\textbf{初等子模型}），如果 \(\mathfrak{A} \subseteq \mathfrak{B}\) 且对每个 \(\mathcal{L}\)-公式 \(\varphi(x_1, \ldots, x_n)\) 和 \(a_1, \ldots, a_n \in A\)，

\[\mathfrak{A} \models \varphi(a_1, \ldots, a_n) \iff \mathfrak{B} \models \varphi(a_1, \ldots, a_n)\]

\textbf{定理 100.1}（Tarski-Vaught 判别法）：\(\mathfrak{A} \subseteq \mathfrak{B}\) 是初等子模型当且仅当对每个 \(\mathcal{L}\)-公式 \(\varphi(x, y_1, \ldots, y_n)\) 和 \(a_1, \ldots, a_n \in A\)，若 \(\mathfrak{B} \models \exists x \varphi(x, a_1, \ldots, a_n)\)，则存在 \(a \in A\) 使得 \(\mathfrak{B} \models \varphi(a, a_1, \ldots, a_n)\)。

\emph{证明}：（\(\Rightarrow\)）若 \(\mathfrak{A}\preceq\mathfrak{B}\) 且 \(\mathfrak{B}\models\exists x\varphi(x,\bar{a})\)，由初等性 \(\mathfrak{A}\models\exists x\varphi(x,\bar{a})\)，故存在 \(a\in A\) 见证。（\(\Leftarrow\)）对公式复杂度归纳证明 \(\mathfrak{A}\models\psi(\bar{a})\iff\mathfrak{B}\models\psi(\bar{a})\)。原子公式平凡。连接词归纳直接。量词 \(\exists x\psi\) 情形由条件保证。\(\blacksquare\)

\subsection{100.2 Löwenheim-Skolem 定理}\label{luxf6wenheim-skolem-ux5b9aux7406}

\textbf{定理 100.2}（向下 Löwenheim-Skolem 定理）：设 \(\mathfrak{A}\) 是 \(\mathcal{L}\)-结构，\(X \subseteq A\) 是子集。则存在 \(\mathfrak{A}\) 的初等子模型 \(\mathfrak{B} \preceq \mathfrak{A}\) 包含 \(X\)，且 \(|B| \leq \max(|X|, |\mathcal{L}|, \aleph_0)\)。

\emph{证明}：构造地封闭 \(X\) 在 Skolem 函数下。对每个 \(\mathcal{L}\)-公式 \(\exists y \varphi(y, x_1, \ldots, x_n)\)，选择 Skolem 函数 \(f_\varphi\) 满足 \(\mathfrak{A} \models \exists y \varphi(y, a_1, \ldots, a_n) \to \varphi(f_\varphi(a_1, \ldots, a_n), a_1, \ldots, a_n)\)。\(\mathcal{L}\) 的 Skolem 扩张为 \(\mathcal{L}^S\)。取 \(X\) 在 Skolem 函数下的闭包 \(B\)，则 \(|B| \leq \max(|X|, |\mathcal{L}|, \aleph_0)\)，且 \(\mathfrak{B} \preceq \mathfrak{A}\)。∎

\textbf{推论 100.3}（Skolem 悖论）：若 ZFC 集合论一致，则它有可数模型（由向下 Löwenheim-Skolem）。但在此可数模型中，可以证明''存在不可数集合''（Cantor 定理在模型内部成立）。这不是矛盾，而是''不可数''在模型内部和外部的含义不同。

\subsection{100.3 量词消去与模型完备性}\label{ux91cfux8bcdux6d88ux53bbux4e0eux6a21ux578bux5b8cux5907ux6027}

\textbf{定义 100.3}（量词消去）：理论 \(T\) 称为有\textbf{量词消去}，如果每个 \(\mathcal{L}\)-公式等价于（在 \(T\) 中）一个无量词公式。

\textbf{定理 100.4}（量词消去的判别法）：\(T\) 有量词消去当且仅当对任意 \(\mathfrak{A}, \mathfrak{B} \models T\)，\(\mathfrak{C} \subseteq \mathfrak{A}, \mathfrak{B}\) 是公共子结构，以及无量词公式 \(\varphi(x, \bar{y})\) 和 \(\bar{c} \in C\)，若 \(\mathfrak{A} \models \exists x \varphi(x, \bar{c})\)，则 \(\mathfrak{B} \models \exists x \varphi(x, \bar{c})\)。

\textbf{证明}：由一阶逻辑的完备性定理（Gödel 1930），Henkin构造法：对一致理论\(T\)，通过添加Henkin常量和极大化，构造项模型\(\mathcal{M}\)。\(\mathcal{M}\)的元素为闭项，谓词解释由\(T\)中的原子公式决定。归纳证明\(\mathcal{M}\models\varphi\iff T\vdash\varphi\)，从而\(T\)一致必可满足。\(\blacksquare\)

\textbf{推论 101.4}（一阶逻辑的不可判定性，Church 1936）：一阶逻辑的有效性（\(\models \varphi\)？）是不可判定的。即不存在算法能判定任意一阶公式是否为逻辑有效式。

\subsection{101.4 算术层次}\label{ux7b97ux672fux5c42ux6b21}

\textbf{定义 101.7}（算术层次）：\(\Sigma_1^0\) 公式形如 \(\exists x_1 \cdots \exists x_n \varphi\)，其中 \(\varphi\) 仅含受囿量词。\(\Pi_1^0\) 公式形如 \(\forall x_1 \cdots \forall x_n \varphi\)。\(\Sigma_{n+1}^0\) 形如 \(\exists \bar{x} \psi\)，\(\psi \in \Pi_n^0\)。\(\Pi_{n+1}^0\) 形如 \(\forall \bar{x} \psi\)，\(\psi \in \Sigma_n^0\)。

\textbf{定理 101.5}：\(A \subseteq \mathbb{N}\) 是 r.e. 当且仅当 \(A\) 是 \(\Sigma_1^0\)-可定义的（在 \(\mathbb{N}\) 的标准模型中）。递归集恰好是 \(\Delta_1^0 = \Sigma_1^0 \cap \Pi_1^0\) 可定义的。

\textbf{证明}：（\(\Rightarrow\)）若 \(A\) 是 r.e.，存在可计算 \(f\) 使 \(A=\operatorname{ran}(f)\) 或 \(A=\operatorname{dom}(f)\)。由 Kleene 正规形定理，\(n\in A\iff\exists t\,T(e,n,t)\)，其中 \(T\) 是原始递归谓词。\(T\) 可在 \(\mathbb{N}\) 中用受囿量词公式定义，故 \(\exists t\,T(e,n,t)\) 是 \(\Sigma_1^0\) 公式。（\(\Leftarrow\)）若 \(A\) 由 \(\Sigma_1^0\) 公式 \(\exists x\theta(n,x)\) 定义（\(\theta\) 仅含受囿量词），则 \(\theta\) 定义的关系是原始递归的，\(A\) 即为该原始递归关系的投影------由 Church 论题是 r.e. 的。递归集 \(A\) 的特征函数可计算 \(\iff A\) 和 \(\mathbb{N}\setminus A\) 均为 r.e. \(\iff\) 两者均为 \(\Sigma_1^0\) \(\iff\) \(A\) 是 \(\Delta_1^0\)。\(\blacksquare\)

\emph{证明}：若 \(A\) 是 r.e.，\(A=\operatorname{dom}(f)\)（\(f\) 可计算），则 \(n\in A\iff\exists t(f(n)\downarrow\) 在第 \(t\) 步\()\)，该条件可写成 \(\Sigma_1^0\) 公式。若 \(A\) 是 \(\Sigma_1^0\)，用 \(\Sigma_1^0\) 公式枚举 \(A\) 得 r.e.。递归集为 \(\Delta_1^0\)：特征函数可计算等价于判定问题可表达为 \(\Sigma_1^0\) 和 \(\Pi_1^0\) 公式。\(\blacksquare\)

\hrulefill

\hrulefill

\hrulefill

\hrulefill

\hrulefill

\hrulefill

\section{第405章：不完全性定理}\label{ux7b2c405ux7ae0ux4e0dux5b8cux5168ux6027ux5b9aux7406}

Kurt Gödel 在 1931 年证明的不完全性定理是 20 世纪数理逻辑乃至整个数学基础领域中最具震撼力的结果。David Hilbert 在 1900 年的巴黎国际数学家大会上提出的第二个问题------证明算术公理的一致性------以及其 1920 年代纲领：将全部数学形式化并证明其完全性和一致性，被 Gödel 的定理无情地终结了。Gödel 第一不完全性定理断言：任何包含基本算术（Peano 算术 PA 的最小片段，即 Robinson 算术 \(Q\)）的一致递归公理系统 \(T\) 都是不完全的------存在一个 \(\Pi_1\) 语句 \(G_T\)（``\(G_T\) 在 \(T\) 中不可证''的算术编码），使得 \(T\) 既不能证明 \(G_T\) 也不能证明 \(\neg G_T\)。Gödel 证明的核心技术------Gödel 编码（Gödelization）和对角线引理（Diagonal Lemma）------将逻辑系统的元数学论述（``\(\phi\) 在 \(T\) 中可证''）编码为系统中的算术公式 \(\operatorname{Proof}_T(x,y)\)，从而实现了''数学系统自我指涉''的逻辑悖论形式化。Gödel 第二不完全性定理进一步断言：如果 \(T\) 是一致的，则 \(T\) 不能证明自身的协调性 \(\operatorname{Con}(T)\)------这直接粉碎了 Hilbert 纲领中在数学系统内部证明数学协调性的梦想。Rosser（1936）通过巧妙的修改移除了 Gödel 证明中对 \(\omega\)-一致性的依赖，使不完全性定理仅需一致性假设。本章系统引入 Gödel 编码、可表示性引理、对角线引理，并给出不完全性定理的完整证明。

\subsection{102.1 算术的形式化}\label{ux7b97ux672fux7684ux5f62ux5f0fux5316}

\textbf{定义 102.1}（Peano 算术，PA）：PA 的语言为 \(\mathcal{L}_{PA} = \{0, S, +, \cdot\}\)（V1 中的 Peano 公理的形式化）。公理包括： 1. \(\forall x (S(x) \neq 0)\) 2. \(\forall x \forall y (S(x) = S(y) \to x = y)\) 3. \(\forall x (x + 0 = x)\) 4. \(\forall x \forall y (x + S(y) = S(x + y))\) 5. \(\forall x (x \cdot 0 = 0)\) 6. \(\forall x \forall y (x \cdot S(y) = x \cdot y + x)\) 7. 归纳公理模式：对每个公式 \(\varphi(x, \bar{y})\)，\([\varphi(0, \bar{y}) \land \forall x (\varphi(x, \bar{y}) \to \varphi(S(x), \bar{y}))] \to \forall x \varphi(x, \bar{y})\)

\subsection{102.2 Gödel 编码}\label{guxf6del-ux7f16ux7801}

\textbf{定义 102.2}（Gödel 编码）：Gödel 编码（或 Gödel 数）是 \(\mathcal{L}_{PA}\) 的符号、公式和证明序列到 \(\mathbb{N}\) 的单射编码。记 \(\ulcorner \varphi \urcorner\) 为公式 \(\varphi\) 的 Gödel 数。

\textbf{引理 102.1}（可表示性）：以下关系和函数在 PA 中可表示（即存在公式定义它们，PA 能证明其基本性质）： - \(\operatorname{Prf}_{PA}(x, y)\)：\(x\) 是 \(y\) 的 PA 证明的 Gödel 数 - \(\operatorname{Prov}_{PA}(y) = \exists x \operatorname{Prf}_{PA}(x, y)\)：\(y\) 是 PA 可证公式的 Gödel 数 - 对角化函数 \(d\)：\(d(\ulcorner \varphi(x) \urcorner) = \ulcorner \varphi(\ulcorner \varphi(x) \urcorner) \urcorner\)

\subsection{102.3 第一不完全性定理}\label{ux7b2cux4e00ux4e0dux5b8cux5168ux6027ux5b9aux7406}

\textbf{定理 102.1}（Gödel 第一不完全性定理）：设 \(T\) 是包含 PA 的一致递归公理化理论。则存在 \(\mathcal{L}_{PA}\)-句子 \(G_T\)（Gödel 句子）使得 \(T \nvdash G_T\) 且 \(T \nvdash \neg G_T\)（假设 \(T\) 是 \(\omega\)-一致的或 \(\Sigma_1\)-可靠的）。

\emph{证明}（Gödel 1931）： 1. 构造自指句子：由对角引理，存在句子 \(G_T\) 满足 \(T \vdash G_T \leftrightarrow \neg \operatorname{Prov}_T(\ulcorner G_T \urcorner)\)。\(G_T\) 断言''\(G_T\) 在 \(T\) 中不可证''。 2. 若 \(T \vdash G_T\)，则 \(T \vdash \operatorname{Prov}_T(\ulcorner G_T \urcorner)\)（由 \(\Sigma_1\)-完备性：\(T\) 能证明每个真 \(\Sigma_1\) 句子）。但 \(T \vdash G_T \to \neg \operatorname{Prov}_T(\ulcorner G_T \urcorner)\)，矛盾。 3. 若 \(T \vdash \neg G_T\)，则 \(T \vdash \operatorname{Prov}_T(\ulcorner G_T \urcorner)\)。由 \(\omega\)-一致性（或 \(\Sigma_1\)-可靠性），这蕴含 \(T \vdash G_T\)，矛盾。故 \(T\) 既不能证明也不能否证 \(G_T\)。∎

\textbf{推论 102.2}（不可完备性）：不存在包含 PA 的一致递归公理化完备理论。数学真理不能完全被形式系统捕获。

\emph{证明}：若存在完备递归公理化理论 \(T\supseteq PA\)，则 \(T\) 必能判定所有 \(\mathcal{L}_{PA}\)-句子。但由第一不完全性定理，任何一致递归理论均不完全（存在 \(G_T\) 不可判定）。故 \(\operatorname{Th}(\mathbb{N})\)（真实算术）非递归可公理化------数学真理超越任何形式系统。\(\blacksquare\)

\textbf{定理 102.3}（Rosser 变体，1936）：Gödel 第一不完全性定理仅需 \(T\) 的一致性（不需要 \(\omega\)-一致性）。Rosser 句子 \(R_T\) 满足 \(R_T \leftrightarrow \forall x (\operatorname{Prf}_T(x, \ulcorner R_T \urcorner) \to \exists y < x \operatorname{Prf}_T(y, \ulcorner \neg R_T \urcorner))\)。

\textbf{证明}：由对角引理构造 Rosser 句 \(R_T\) 满足等价性。\(R_T\) 的直观含义为''若存在我的证明，则存在更短的对我否定的证明''。若 \(T\vdash R_T\)，设 \(p\) 为最短证明的 Gödel 数，则 \(T\vdash\operatorname{Prf}_T(\overline{p},\ulcorner R_T\urcorner)\)。同时 \(T\vdash R_T\) 蕴含 \(T\vdash\exists y<p\,\operatorname{Prf}_T(y,\ulcorner\neg R_T\urcorner)\)。若实际上 \(\neg R_T\) 无证明（因 \(T\) 一致），则 \(\exists y<p\,\operatorname{Prf}_T(y,\ulcorner\neg R_T\urcorner)\) 在 \(\mathbb{N}\) 中为假。但该式子是 \(\Sigma_1^0\) 的（受囿存在量词），\(\Sigma_1^0\) 完备性迫使该式在 \(T\) 中可证，矛盾------因为 \(T\) 不能同时证明 \(\operatorname{Prf}_T(p,\ulcorner R_T\urcorner)\)（真）和 \(\exists y<p\,\operatorname{Prf}_T(y,\ulcorner\neg R_T\urcorner)\)（假且 \(\Sigma_1^0\)）。若 \(T\vdash\neg R_T\)，对称论证亦导致矛盾。Rosser 的精妙之处在于仅需 \(T\) 的一致性而不需 \(\omega\)-一致性。\(\blacksquare\)

\subsection{102.4 第二不完全性定理}\label{ux7b2cux4e8cux4e0dux5b8cux5168ux6027ux5b9aux7406}

\textbf{定义 102.3}（一致性陈述）：\(\operatorname{Con}_T = \neg \operatorname{Prov}_T(\ulcorner 0 = 1 \urcorner)\)（或 \(\neg \operatorname{Prov}_T(\ulcorner \bot \urcorner)\)）。\(\operatorname{Con}_T\) 断言 \(T\) 是一致的。

\textbf{定理 102.4}（Gödel 第二不完全性定理）：设 \(T\) 是包含 PA 的一致递归公理化理论。则 \(T \nvdash \operatorname{Con}_T\)。即 \(T\) 不能证明自身的一致性。

\textbf{证明}：由一阶逻辑的完备性定理（Gödel 1930），Henkin构造法：对一致理论\(T\)，通过添加Henkin常量和极大化，构造项模型\(\mathcal{M}\)。\(\mathcal{M}\)的元素为闭项，谓词解释由\(T\)中的原子公式决定。归纳证明\(\mathcal{M}\models\varphi\iff T\vdash\varphi\)，从而\(T\)一致必可满足。\(\blacksquare\)

\textbf{推论 102.5}：ZFC 集合论的一致性不能在 ZFC 内部证明。Hilbert 计划（用有限方法证明数学的一致性）被 Gödel 不完全性定理从根本上限制。

\hrulefill

\hrulefill

\hrulefill

\hrulefill

\hrulefill

\hrulefill

\section{第406章：集合论进阶}\label{ux7b2c406ux7ae0ux96c6ux5408ux8bbaux8fdbux9636}

集合论作为现代数学的统一语言，其''深度''远远超出了 V1 中关于集合运算和 ZFC 公理化基础的入门内容。本章在 ZFC 公理系统的基础上，深入探讨序数（ordinal）和基数（cardinal）的算术以及著名的连续统假设（Continuum Hypothesis，CH）。超限序数由 Cantor 在 1880 年代发现，von Neumann 将每个序数定义为所有更小序数的集合（\(\alpha=\{\beta:\beta<\alpha\}\)），从而自然形成了序数的 \(\in\) 链结构：\(0=\emptyset\)，\(1=\{0\}\)，\(2=\{0,1\}\)，\(\ldots\)，\(\omega=\{0,1,2,\ldots\}\)，\(\omega+1=\{0,1,2,\ldots,\omega\}\)。超限归纳法和超限递归定理将经典数学归纳法推广到全体序数类 \(\mathbf{On}\)，为构造性集合论中的层次（如 Gödel 的可构造宇宙 \(L\)）提供了严密的归纳基础。基数算术在无穷领域展现出与有限算术截然不同的''吸收性''：\(\kappa+\lambda=\kappa\cdot\lambda=\max(\kappa,\lambda)\)（\(\kappa,\lambda\) 为无穷基数），而指数运算 \(\kappa^\lambda\) 在无穷基数的情形下受限于 König 不等式 \(\operatorname{cf}(2^\kappa)>\kappa\)。连续统假设（\(2^{\aleph_0}=\aleph_1\)）是 Hilbert 23 个问题的首位，Gödel（1940）证明它不能被 ZFC 否证（\(L\) 是 ZFC+GCH 的模型），Cohen（1963）用力迫法（forcing）证明它也不能被 ZFC 证明------从而 CH 独立于 ZFC，这一辉煌成果彻底改变了数学基础的格局。

\subsection{103.1 序数与超限归纳}\label{ux5e8fux6570ux4e0eux8d85ux9650ux5f52ux7eb3}

\textbf{定义 103.1}（序数，von Neumann 定义）：一个集合 \(\alpha\) 称为\textbf{序数}，如果 \(\alpha\) 是传递的（即 \(\beta \in \alpha \Rightarrow \beta \subseteq \alpha\)）且 \(\alpha\) 在 \(\in\) 下是良序的。

\textbf{定理 103.1}（超限归纳法）：若对每个序数 \(\alpha\)，\((\forall \beta < \alpha) P(\beta)\) 推出 \(P(\alpha)\)，则 \(P(\alpha)\) 对所有序数 \(\alpha\) 成立。

\textbf{证明}：反证法。假设存在序数 \(\alpha\) 使 \(\neg P(\alpha)\)。令 \(\beta=\min\{\gamma\le\alpha:\neg P(\gamma)\}\)------序数的良序性保证此最小元存在。对任意 \(\delta<\beta\)，由 \(\beta\) 的最小性，\(P(\delta)\) 成立。故 \((\forall\delta<\beta)P(\delta)\) 成立。由归纳假设条件，\(P(\beta)\) 成立，与 \(\beta\) 的定义矛盾。因此不存在使 \(\neg P\) 成立的最小序数，故对所有序数 \(P\) 成立。这本质上等同于序数的 \(\in\) 是良基关系。\(\blacksquare\)

\textbf{定理 103.2}（超限递归定义）：对任意给定的类函数 \(G\)，存在唯一的类函数 \(F\) 定义在所有序数上，满足 \(F(\alpha) = G(F \restriction \alpha)\)。

\textbf{证明}：存在性：对每个序数 \(\beta\)，用超限归纳定义 \(F\restriction\beta\)。设对所有 \(\gamma<\beta\) 已定义 \(F(\gamma)\)，则 \(F\restriction\beta\) 是 \(\{(\gamma,F(\gamma)):\gamma<\beta\}\)，应用 \(G\) 得 \(F(\beta)=G(F\restriction\beta)\)。此定义的合法性由替换公理模式保证：函数 \(\beta\mapsto F\restriction\beta\) 的所有值构成集合。唯一性：若 \(F_1\neq F_2\)，取最小序数 \(\alpha\) 使 \(F_1(\alpha)\neq F_2(\alpha)\)，则对任意 \(\beta<\alpha\)，\(F_1(\beta)=F_2(\beta)\) 即 \(F_1\restriction\alpha=F_2\restriction\alpha\)。由定义 \(F_i(\alpha)=G(F_i\restriction\alpha)\)，故 \(F_1(\alpha)=F_2(\alpha)\)，矛盾。\(\blacksquare\)

\textbf{定义 103.2}（序数算术）： - \(\alpha + 0 = \alpha\)，\(\alpha + (\beta + 1) = (\alpha + \beta) + 1\)，\(\alpha + \lambda = \sup_{\beta < \lambda} (\alpha + \beta)\)（极限序数 \(\lambda\)） - \(\alpha \cdot 0 = 0\)，\(\alpha \cdot (\beta + 1) = \alpha \cdot \beta + \alpha\)，\(\alpha \cdot \lambda = \sup_{\beta < \lambda} (\alpha \cdot \beta)\) - \(\alpha^0 = 1\)，\(\alpha^{\beta+1} = \alpha^\beta \cdot \alpha\)，\(\alpha^\lambda = \sup_{\beta < \lambda} \alpha^\beta\)

注意：序数加法和乘法不交换（\(1 + \omega = \omega \neq \omega + 1\)）。

\subsection{103.2 基数}\label{ux57faux6570}

\textbf{定义 103.3}（基数）：序数 \(\kappa\) 称为\textbf{基数}（或初始序数），如果对所有 \(\alpha < \kappa\)，\(|\alpha| < |\kappa|\)（即不存在 \(\alpha\) 到 \(\kappa\) 的双射）。

\textbf{定义 103.4}（\(\aleph\) 序列）：\(\aleph_0 = \omega\)。\(\aleph_{\alpha+1} = \min\{\kappa : \kappa > \aleph_\alpha\}\)（\(\aleph_\alpha\) 的后继基数）。\(\aleph_\lambda = \sup_{\alpha < \lambda} \aleph_\alpha\)（极限序数 \(\lambda\)）。

\textbf{定理 103.3}（Cantor 定理）：对任意集合 \(X\)，\(|X| < |\mathcal{P}(X)|\)。特别地，\(\aleph_\alpha < 2^{\aleph_\alpha}\)。

\emph{证明}：\(f: X \to \mathcal{P}(X)\) 不是满射：考虑 \(D = \{x \in X : x \notin f(x)\}\)。若 \(D = f(d)\)，则 \(d \in D \iff d \notin f(d) = D\)，矛盾。∎

\textbf{定义 103.5}（基数算术）：对基数 \(\kappa, \lambda\)： - \(\kappa + \lambda = |\kappa \sqcup \lambda|\) - \(\kappa \cdot \lambda = |\kappa \times \lambda|\) - \(\kappa^\lambda = |\{f: \lambda \to \kappa\}|\)

\textbf{定理 103.4}（基数算术基本定理）：对无穷基数 \(\kappa, \lambda\)： - \(\kappa + \lambda = \kappa \cdot \lambda = \max(\kappa, \lambda)\)（若 \(\kappa, \lambda > 0\)） - \(\kappa^\kappa = 2^\kappa\)

\textbf{证明}：设 \(\kappa\le\lambda\) 且 \(\lambda\) 无穷。由 Zorn 引理（AC）可证 \(\kappa\times\lambda\) 与 \(\lambda\) 等势：将 \(\lambda\) 划分为 \(\kappa\) 个大小为 \(\lambda\) 的子集（\(\lambda\times\kappa\) 的良序），故 \(\kappa\cdot\lambda=|\kappa\times\lambda|=\lambda\)。由 \(\kappa+\lambda=|\kappa\sqcup\lambda|\le\kappa\cdot\lambda\le\lambda\cdot\lambda=\lambda\)（Hessenberg 定理），且 \(\lambda\le\kappa+\lambda\)，故 \(\kappa+\lambda=\lambda\)。对于 \(\kappa^\kappa=2^\kappa\)：\(2^\kappa\le\kappa^\kappa\le(2^\kappa)^\kappa=2^{\kappa\cdot\kappa}=2^\kappa\)（因 \(\kappa\cdot\kappa=\kappa\) 对无穷基数），由 Cantor-Bernstein-Schröder 定理得等势。\(\blacksquare\)

\textbf{定理 103.5}（König 定理）：对基数 \(\kappa_i < \lambda_i\)（\(i \in I\)），\(\sum_{i \in I} \kappa_i < \prod_{i \in I} \lambda_i\)。

\textbf{证明}：设 \(\sum\kappa_i=\kappa\)，\(\prod\lambda_i=\lambda\)。需证明 \(\kappa<\lambda\)（\(\kappa\not\ge\lambda\)）。反设存在满射 \(f:\bigsqcup K_i\to\prod L_i\)（\(|K_i|=\kappa_i<\lambda_i=|L_i|\)）。对每个 \(i\)，定义 \(D_i=\{a(i):a\in\operatorname{ran}(f)\restriction K_i\}\subseteq L_i\)。因 \(|K_i|<|L_i|\)，\(D_i\subsetneq L_i\)（非满）。选 \(b_i\in L_i\setminus D_i\)，则 \(b=(b_i)_{i\in I}\in\prod L_i\)。但对任意被 \(f\) 命中的元素 \(x\in K_i\)，\(b_i\neq f(x)(i)\)，故 \(b\notin\operatorname{ran}(f)\)，与 \(f\) 满射矛盾。对角线论证是 Cantor 定理的推广。\(\blacksquare\)

\subsection{103.3 选择公理的等价形式}\label{ux9009ux62e9ux516cux7406ux7684ux7b49ux4ef7ux5f62ux5f0f}

\textbf{定理 103.6}（选择公理，AC）：以下陈述在 ZF 中等价： 1. \textbf{选择公理}：对任意非空集合族 \(\{X_i\}_{i \in I}\)，存在选择函数 \(f: I \to \bigcup_i X_i\) 使得 \(f(i) \in X_i\) 2. \textbf{Zorn 引理}：若非空偏序集中每条链有上界，则该偏序集有极大元 3. \textbf{良序定理}：每个集合可被良序 4. \textbf{Tychonoff 定理}：紧致空间的任意积是紧致的（V10 Ch 91） 5. 每个向量空间有基（V7 Ch 51） 6. 每个满射有右逆

\textbf{证明}：展示 AC \(\Rightarrow\) Zorn \(\Rightarrow\) 良序 \(\Rightarrow\) AC 的循环。（AC \(\Rightarrow\) Zorn）设 \((P,\le)\) 满足链条件但无极大元。对每条链 \(C\)，其真上界集非空，由 AC 选出 \(f(C)>C\)。超限递归构造严格递增链 \(p_0<p_1<\cdots\)，由 Hartogs 定理此链需长于某序数，与 \(P\) 为集合矛盾。（Zorn \(\Rightarrow\) 良序）取集合 \(X\) 的良序化部分函数集 \((W,\subseteq)\)（\(W\) 为 \(X\) 子集上的良序），此偏序集满足链条件，极大元即为 \(X\) 的全良序。（良序 \(\Rightarrow\) AC）将 \(\bigcup X_i\) 良序化，取各 \(X_i\) 的最小元即得选择函数。剩余等价性类似可证。\(\blacksquare\)

\textbf{定理 103.7}（基数比较定理，在 AC 下）：对任意两个集合 \(X, Y\)，\(|X| \leq |Y|\) 或 \(|Y| \leq |X|\)。即任意两个基数可比较。

\textbf{证明}：由良序定理（等价于 AC），\(X\) 和 \(Y\) 均可被良序。设 \(\alpha=|X|\)，\(\beta=|Y|\) 为对应基数（初始序数）。由于序数是 \(\in\) 下良序的，\(\alpha\in\beta\)，\(\alpha=\beta\)，或 \(\beta\in\alpha\) 三者必居其一。若 \(\alpha\in\beta\)，则存在单射 \(f:\alpha\to\beta\)，故 \(|X|=\alpha\le\beta=|Y|\)。类似地 \(\beta\in\alpha\) 给出 \(|Y|\le|X|\)。若 \(\alpha=\beta\)，则 \(|X|=|Y|\) 即 \(|X|\le|Y|\) 且 \(|Y|\le|X|\)。故两基数总可比较（这是 AC 的核心推论------无 AC 时存在不可比较的基数）。\(\blacksquare\)

\subsection{103.4 连续统假设}\label{ux8fdeux7eedux7edfux5047ux8bbe}

\textbf{定义 103.6}（连续统假设，CH）：\textbf{连续统假设}断言 \(2^{\aleph_0} = \aleph_1\)。即实数集 \(\mathbb{R}\) 的基数是第一个不可数基数。

\textbf{广义连续统假设（GCH）}：对所有序数 \(\alpha\)，\(2^{\aleph_\alpha} = \aleph_{\alpha+1}\)。

\textbf{定理 103.8}（Gödel 1938，Cohen 1963）： 1. 若 ZF 一致，则 ZFC + GCH 一致（Gödel：CH 在 ZFC 中不可否证，即可构造宇宙 \(L\) 满足 GCH） 2. 若 ZFC 一致，则 ZFC + \(\neg\)CH 一致（Cohen：力迫法，CH 在 ZFC 中不可证明）

\textbf{证明}：（1）Gödel 在 ZF 内定义可构造宇宙 \(L\)：\(L_0=\varnothing\)，\(L_{\alpha+1}=\operatorname{Def}(L_\alpha)\)，\(L_\lambda=\bigcup_{\alpha<\lambda}L_\alpha\)。证明 ZF 的每条公理在 \(L\) 中成立（相对化）。在 \(L\) 中可定义良序，故 AC 成立。通过凝聚引理（condensation lemma）证明 \(L\) 中每个子集 \(X\subseteq\aleph_\alpha\) 的基数 \(\le\aleph_{\alpha+1}\)，即 \(2^{\aleph_\alpha}=\aleph_{\alpha+1}\)（GCH）。因此 \(\operatorname{Con}(ZF)\Rightarrow\operatorname{Con}(ZFC+GCH)\)。（2）Cohen 引入力迫法：从 ZFC 的可数传递模型 \(M\) 出发，对偏序 \(P\)（如添加 \(\aleph_2\) 个 Cohen 实数的偏序），取 \(M\) 上的 \(P\)-一般滤子 \(G\) 构造扩张 \(M[G]\)。在 \(M[G]\) 中至少有 \(\aleph_2\) 个实数，故 \(2^{\aleph_0}\ge\aleph_2\)，\(\neg\)CH 成立。因此 \(\operatorname{Con}(ZFC)\Rightarrow\operatorname{Con}(ZFC+\neg CH)\)。综合得 CH 独立于 ZFC。\(\blacksquare\)

CH 独立于 ZFC！这是数理逻辑最著名的独立性结果。Cohen 因此获得 Fields 奖（1966）。

\textbf{定义 103.7}（可构造宇宙）：Gödel 的\textbf{可构造宇宙} \(L\) 通过超限递归定义：\(L_0 = \varnothing\)，\(L_{\alpha+1} = \operatorname{Def}(L_\alpha)\)（\(L_\alpha\) 的可定义子集），\(L_\lambda = \bigcup_{\alpha < \lambda} L_\alpha\)（极限序数）。\(L = \bigcup_{\alpha \in \mathbf{Ord}} L_\alpha\)。

\textbf{定理 103.9}（\(V = L\) 公理）：在 \(L\) 中，选择公理和 GCH 都成立。\(L\) 是最小的包含所有序数的 ZF 内模型。

\textbf{证明}：（AC 在 \(L\) 中）\(L\) 有自然的良序：按构造层次 \((\alpha,x)\)（\(\alpha\) 为 \(x\) 首次出现在 \(L\) 中的序数）字典序排列各元素。该良序是 \(L\) 可定义的，故对 \(L\) 中任意非空集合族可选取最小元，AC 成立。（GCH 在 \(L\) 中）Gödel 凝聚引理：设 \(X\prec L_\theta\) 为初等子模型，则 \(X\) 的 Mostowski 坍缩同构于某个 \(L_\gamma\)。特别地，取 \(L_{\aleph_{\alpha+1}}\) 的初等子模型 \(X\) 包含 \(\aleph_\alpha\cup\{x\}\)（\(x\subseteq\aleph_\alpha\)），\(|X|=\aleph_\alpha\)，坍缩后 \(x\) 对应于某 \(L_\gamma\) 中元素且 \(x\subseteq\aleph_\alpha\)。因此 \(\mathcal{P}(\aleph_\alpha)\cap L\subseteq L_{\aleph_{\alpha+1}}\)，而 \(|L_{\aleph_{\alpha+1}}|=\aleph_{\alpha+1}\)，故 \((2^{\aleph_\alpha})^L=\aleph_{\alpha+1}\)。\(L\) 最小性：任何 ZF 内模型 \(M\) 必包含所有序数，由超限归纳 \(L_\alpha\subseteq M\)（因 \(\operatorname{Def}\) 是绝对的内模型操作），故 \(L\subseteq M\)。\(\blacksquare\)

\textbf{定义 103.8}（力迫法概要）：力迫法（Cohen 1963）通过添加''一般滤子''（generic filter）来扩张可数传递模型 \(M\)，得到 \(M[G]\)，其中某些命题（如 \(\neg\)CH）成立。力迫法是现代集合论最强大的工具之一。

\hrulefill

\hrulefill

\hrulefill

\hrulefill

\hrulefill

\newpage

\chapter{11-数值与计算/02-计算理论/03-计算复杂性.md}\label{ux6570ux503cux4e0eux8ba1ux7b9702-ux8ba1ux7b97ux7406ux8bba03-ux8ba1ux7b97ux590dux6742ux6027.md}

\section{第407章：计算复杂性：P、NP 与 NP 完全性}\label{ux7b2c407ux7ae0ux8ba1ux7b97ux590dux6742ux6027pnp-ux4e0e-np-ux5b8cux5168ux6027}

计算复杂性理论研究计算问题所需资源（时间、空间）与输入规模的关系，是理论计算机科学的核心。

\subsection{154.1 时间复杂度类}\label{ux65f6ux95f4ux590dux6742ux5ea6ux7c7b}

\textbf{定义 154.1}（时间复杂度类）： - \(\mathbf{TIME}(t(n)) = \{L : \text{存在 } O(t(n)) \text{ 时间判定 } L \text{ 的 TM}\}\) - \(\mathbf{P} = \bigcup_{k \geq 1} \mathbf{TIME}(n^k)\)（多项式时间可判定语言） - \(\mathbf{NP} = \bigcup_{k \geq 1} \mathbf{NTIME}(n^k)\)（非确定型多项式时间可判定语言） - \(\mathbf{EXP} = \bigcup_{k \geq 1} \mathbf{TIME}(2^{n^k})\)（指数时间可判定语言）

\textbf{定义 154.2}（NP 的等价刻画）：\(L \in \mathbf{NP}\) 当且仅当存在多项式时间可验证的证书（certificate）。即存在多项式 \(p\) 和多项式时间 TM \(V\) 使得 \(x \in L \iff \exists y, |y| \leq p(|x|), V(x, y) = 1\)。

\subsection{154.2 NP 完全性}\label{np-ux5b8cux5168ux6027}

\textbf{定义 154.3}（多项式时间归约）：\(A \leq_P B\) 如果存在多项式时间可计算函数 \(f\) 使得 \(x \in A \iff f(x) \in B\)。

\textbf{定义 154.4}（NP 完全与 NP 难）： - \(L\) 是 \(\mathbf{NP}\)-\textbf{难}的，如果对所有 \(A \in \mathbf{NP}\)，\(A \leq_P L\) - \(L\) 是 \(\mathbf{NP}\)-\textbf{完全}的，如果 \(L \in \mathbf{NP}\) 且 \(L\) 是 NP 难的

\textbf{定理 154.1}（Cook-Levin 定理，1971-1973）：SAT 是 NP 完全的。

\textbf{证明}：Cook-Levin定理的证明核心是将非确定性Turing机\(N\)（运行时间\(p(n)\)）的计算过程编码为布尔公式\(\varphi\)。\(\varphi\)的变量表示\(N\)在时刻\(t\)（\(0\le t\le p(n)\)）的格局（包括状态、带头位置、带内容）。\(\varphi\)的clauses编码：（1）\(N\)的初始格局；（2）\(N\)的转移函数（每步合法转移）；（3）\(N\)的接受条件；（4）唯一性约束（每个时刻只有一个状态、一个带头位置等）。\(\varphi\)为可满足的当且仅当\(N\)接受\(w\)。\(\varphi\)的大小为\(O(p(n)^2)\)，可在多项式时间内构造。\(\blacksquare\)

\textbf{定理 154.2}（Karp 的 21 个 NP 完全问题，1972）：通过多项式时间归约，证明了 3SAT、CLIQUE、VERTEX COVER、HAMILTONIAN CYCLE、SUBSET SUM、PARTITION、KNAPSACK 等 21 个问题的 NP 完全性。

\textbf{证明}：归约链从 SAT 出发。SAT \(\le_P\) 3SAT：将任意子句 \(\ell_1\lor\cdots\lor\ell_k\)（\(k>3\)）替换为 \(k-2\) 个 3-子句，引入新变量 \(y_i\)：\((\ell_1\lor\ell_2\lor y_1)\land(\neg y_1\lor\ell_3\lor y_2)\land\cdots\land(\neg y_{k-3}\lor\ell_{k-1}\lor\ell_k)\)。3SAT \(\le_P\) CLIQUE：构造图的顶点为公式中每个子句的每个文字出现，边连接兼容赋值。\(k\)-团 \(\iff\) 公式可满足。CLIQUE \(\le_P\) VERTEX COVER：\(G\) 有 \(k\) 团 \(\iff\) \(\overline{G}\) 有 \(n-k\) 顶点覆盖。VERTEX COVER \(\le_P\) HAMILTONIAN CYCLE：构造选择小工具。其余归约类似，均保持多项式时间。整套归约建立了 NP 完全性网络。\(\blacksquare\)

\textbf{归约链}：SAT \(\leq_P\) 3SAT \(\leq_P\) CLIQUE \(\leq_P\) VERTEX COVER \(\leq_P\) HAMILTONIAN CYCLE。

\subsection{154.3 P vs NP 问题}\label{p-vs-np-ux95eeux9898}

\textbf{问题 154.1}（P vs NP 问题 / Clay 千禧年问题）：\(\mathbf{P} \stackrel{?}{=} \mathbf{NP}\)。即：是否每个解可被快速验证的问题，其解也可被快速找到？

绝大多数理论计算机科学家相信 \(\mathbf{P} \neq \mathbf{NP}\)，即存在本质上难以求解但解易于验证的问题。若 \(\mathbf{P} = \mathbf{NP}\)，则公钥密码学将崩溃（因整数分解属 NP），数学证明自动化、蛋白质折叠等难题将在多项式时间内可解；若 \(\mathbf{P} \neq \mathbf{NP}\)，则意味着存在不可逾越的计算硬阈值，这也是当前密码学安全性的理论前提。该问题自 1971 年 Cook-Levin 提出以来，虽已有大量证据支持 \(\mathbf{P} \neq \mathbf{NP}\)，但严格证明至今缺失。

\textbf{定理 154.3}（相对化障碍 / Baker-Gill-Solovay，1975）：存在神谕 \(A, B\) 使得 \(\mathbf{P}^A = \mathbf{NP}^A\) 和 \(\mathbf{P}^B \neq \mathbf{NP}^B\)。因此，对角化论证不能解决 P vs NP 问题。

\textbf{证明}：构造 \(B\) 使 \(\mathbf{P}^B\neq\mathbf{NP}^B\)：令 \(L_B=\{1^n:\exists x\in B,|x|=n\}\)（单言神谕语言）。\(L_B\in\mathbf{NP}^B\)（非确定猜测 \(x\) 并查询神谕）。用对角线方法构造 \(B\)：枚举所有多项式时间 OTMs（神谕 TM），在第 \(n\) 步挫败第 \(n\) 个 OTM 判定 \(L_B\) 的尝试。构造 \(A\) 使 \(\mathbf{P}^A=\mathbf{NP}^A\)：取 \(A\) 为 PSPACE 完全语言，则 \(\mathbf{NP}^A\subseteq\mathbf{NPSPACE}=\mathbf{PSPACE}\subseteq\mathbf{P}^A\)。由于两个方向均可能，对角化论证本身不能区分 \(\mathbf{P}\) 和 \(\mathbf{NP}\)。\(\blacksquare\)

\textbf{定理 154.4}（自然证明障碍 / Razborov-Rudich，1997）：若存在伪随机数生成器，则不存在「自然证明」证明 \(\mathbf{P} \neq \mathbf{NP}\)（即电路复杂度下界的高效构造性证明）。

\textbf{证明}：自然证明需满足：（1）构造性：属性 \(C_n\) 可在 \(2^{O(n)}\) 时间内判定；（2）规模大：至少 \(2^{-O(n)}\) 比例的函数满足 \(C_n\)；（3）有用：\(C_n\) 定义的函数类不含 \(\mathbf{P}_{/\text{poly}}\) 语言。假设存在自然证明且存在亚指数安全性的伪随机生成器（PRG）。由 PRG 安全性，随机函数以极高概率满足 \(C_n\)（因 \(C_n\) 规模大）。又由构造性，可生成满足 \(C_n\) 的显式函数。但 PRG 的安全性意味着任何有效可构造的函数族均与随机函数不可区分，矛盾。因此自然证明与密码学假设矛盾。\(\blacksquare\)

\textbf{定理 154.5}（代数化障碍 / Aaronson-Wigderson，2008）：代数化（代数神谕）方法也不能解决 P vs NP 问题。

\textbf{证明}：代数化是相对化的推广：除神谕访问外，还允许用低次多项式扩展神谕。Aaronson-Wigderson 证明：（1）存在代数神谕 \(\tilde{A}\) 使 \(\widetilde{\mathbf{P}}^{\tilde{A}}=\widetilde{\mathbf{NP}}^{\tilde{A}}\)；（2）存在代数神谕 \(\tilde{B}\) 使 \(\widetilde{\mathbf{P}}^{\tilde{B}}\neq\widetilde{\mathbf{NP}}^{\tilde{B}}\)。但代数化方法能证明 \(\mathbf{IP}=\mathbf{PSPACE}\)（通过多项式交互验证），却恰好无力区分 P vs NP。这表明 P vs NP 需要既非相对化也非代数化的新技术。\(\blacksquare\)

\hrulefill

\hrulefill

\hrulefill

\hrulefill

\hrulefill

\hrulefill

\section{第408章：复杂度类与空间复杂性}\label{ux7b2c408ux7ae0ux590dux6742ux5ea6ux7c7bux4e0eux7a7aux95f4ux590dux6742ux6027}

本章深入探讨空间复杂性、多项式谱系和概率与交互证明类。

\subsection{155.1 空间复杂性}\label{ux7a7aux95f4ux590dux6742ux6027}

\textbf{定义 155.1}（空间复杂度类）： - \(\mathbf{SPACE}(s(n))\)：使用 \(O(s(n))\) 空间的 TM 可判定语言 - \(\mathbf{L} = \mathbf{SPACE}(\log n)\)（对数空间） - \(\mathbf{NL} = \mathbf{NSPACE}(\log n)\)（非确定型对数空间） - \(\mathbf{PSPACE} = \bigcup_k \mathbf{SPACE}(n^k)\)（多项式空间） - \(\mathbf{NPSPACE} = \bigcup_k \mathbf{NSPACE}(n^k)\)

\textbf{定理 155.1}（Savitch 定理，1970）：对 \(s(n) \geq \log n\)，\(\mathbf{NSPACE}(s(n)) \subseteq \mathbf{SPACE}(s(n)^2)\)。特别地，\(\mathbf{PSPACE} = \mathbf{NPSPACE}\)。

\textbf{证明}：设 NTM \(N\) 在空间 \(s(n)\) 内判定 \(L\)。\(N\) 的格局数为 \(2^{O(s(n))}\)。判定两格局 \(C_1,C_2\) 是否在 \(\le 2^k\) 步内可达的递归算法：若 \(k=0\)，检查 \(C_1=C_2\) 或一步转移；否则枚举中间格局 \(C_m\)，递归检查 \((C_1,C_m)\) 和 \((C_m,C_2)\) 在 \(\le 2^{k-1}\) 步内可达。由于 \(2^k\le 2^{O(s(n))}\)（最大步数），递归深度 \(O(s(n))\)。递归栈每层存当前格局三元组 \((C_1,C_m,C_2)\)，空间 \(O(s(n))\)。总空间 \(O(s(n))\times O(s(n))=O(s(n)^2)\)。此动态规划（可达性矩阵幂）的空间高效版本。\(\blacksquare\)

\textbf{定理 155.2}（空间层次定理）：若 \(s_1(n) = o(s_2(n))\) 且 \(s_2\) 是空间可构造的，则 \(\mathbf{SPACE}(s_1(n)) \subsetneq \mathbf{SPACE}(s_2(n))\)。

\textbf{证明}：对角线构造。定义 TM \(D\)：在输入 \(w\) 上，若 \(|w|=n\)，\(D\) 模拟第 \(n\) 个 TM \(M_n\) 在 \(w\) 上的运行，空间限制为 \(s_2(n)\)。若 \(M_n\) 在 \(\le 2^{s_2(n)}\) 步内停机且接受，则 \(D\) 拒绝；若拒绝或超时则 \(D\) 接受。由空间可构造性，\(D\) 可强制自身仅用 \(O(s_2(n))\) 空间（维护步数计数器）。取 \(n\) 充分大使 \(s_1(n)<s_2(n)\) 且 \(M_n\) 为 \(s_1(n)\) 界 TM。\(D\) 与 \(M_n\) 在 \(w\) 上不同（对角线），\(D\notin\mathbf{SPACE}(s_1(n))\) 但 \(D\in\mathbf{SPACE}(s_2(n))\)。\(\blacksquare\)

\textbf{已知包含关系}：\(\mathbf{L} \subseteq \mathbf{NL} \subseteq \mathbf{P} \subseteq \mathbf{NP} \subseteq \mathbf{PSPACE} \subseteq \mathbf{EXP}\)。其中 \(\mathbf{NL} \subsetneq \mathbf{PSPACE}\)（空间层次定理），但 \(\mathbf{P} \stackrel{?}{=} \mathbf{PSPACE}\) 仍未知。

\textbf{定义 155.2}（对数空间归约与 NL 完全性）：\(\mathbf{PATH}\)（有向图的 \(s\)-\(t\) 连通性）是 NL 完全的。\(\mathbf{2SAT}\) 是 NL 完全的。

\textbf{定理 155.3}（Immerman-Szelepcsényi 定理，1987-1988）：\(\mathbf{NL} = \mathbf{coNL}\)。即非确定型对数空间在补运算下封闭。

\textbf{证明}：证 \(\overline{\operatorname{PATH}}\in\mathbf{NL}\)（\(\operatorname{PATH}\) 为有向图 \(s\)-\(t\) 连通性，是 NL 完全的）。给定图 \(G\) 和顶点 \(s,t\)，需非确定检验 \(t\) 从 \(s\) 不可达。关键：非确定计算从 \(s\) 可达的顶点总数 \(c\)。按距离分层 \(d=0,1,\dots,n-1\) 进行归纳计数：已知 \(d\) 距离内的顶点数 \(c_d\)，枚举所有顶点 \(v\)，对每个 \(v\) 猜测是否距 \(s\) \(\le d+1\)，验证并更新计数。若 \(t\) 不在任何层中则接受。因只需存 \(O(\log n)\) 个顶点和计数值（\(\le n\)），空间为 \(O(\log n)\)。该定理获 1995 年 Gödel 奖。\(\blacksquare\)

\subsection{155.2 多项式谱系}\label{ux591aux9879ux5f0fux8c31ux7cfb}

\textbf{定义 155.3}（多项式谱系 / Polynomial Hierarchy）： - \(\Sigma_0^P = \Pi_0^P = \mathbf{P}\) - \(\Sigma_{k+1}^P = \mathbf{NP}^{\Sigma_k^P}\)（带 \(\Sigma_k^P\) 神谕的 NP） - \(\Pi_{k+1}^P = \mathbf{coNP}^{\Sigma_k^P}\) - \(\mathbf{PH} = \bigcup_k \Sigma_k^P\)

\textbf{定理 155.4}（多项式谱系的坍塌性质）：若 \(\Sigma_k^P = \Pi_k^P\)，则 \(\mathbf{PH} = \Sigma_k^P\)（谱系坍塌到第 \(k\) 层）。特别地，若 \(\mathbf{P} = \mathbf{NP}\)，则 \(\mathbf{P} = \mathbf{PH}\)。

\textbf{证明}：由定义 \(\Sigma_{k+1}^P=\mathbf{NP}^{\Sigma_k^P}\)。若 \(\Sigma_k^P=\Pi_k^P\)，则 \(\Sigma_k^P\) 对补封闭。考虑 \(L\in\Sigma_{k+1}^P=\mathbf{NP}^{\Sigma_k^P}\)。\(L\) 由带 \(\Sigma_k^P\) 神谕的 NTM \(M\) 判定。因神谕属 \(\Sigma_k^P=\Pi_k^P\)，可替换为 \(\Sigma_k^P\) 判定器 \(A\)（兼判正反）。于是 \(M^A\) 为 \(\Sigma_k^P\) 机器（神谕查询归入存在量词），故 \(L\in\Sigma_k^P\)。一般地 \(\Sigma_{k+1}^P=\Sigma_k^P\)，谱系坍塌。若 \(\mathbf{P}=\mathbf{NP}\) 则 \(\Sigma_1^P=\Pi_1^P\)，谱系坍塌到 \(\mathbf{P}\)。\(\blacksquare\)

\textbf{定理 155.5}（Toda 定理，1991）：\(\mathbf{PH} \subseteq \mathbf{P}^{\#\mathbf{P}}\)。即 \(\#\mathbf{P}\)（计数类）包含了整个多项式谱系。

\textbf{证明}：Toda 的核心引理：对任意 \(\Sigma_k^P\) 语言 \(L\)，存在 \(\#\mathbf{P}\) 函数 \(f\) 和多项式时间算法 \(A\) 使 \(x\in L\iff A(x,f(x))=1\)。首先将 \(L\) 的成功路径计数问题归约到 \(\#\mathbf{P}\)。然后通过''计数放大''技术：使用 \(\#\mathbf{P}\) 神谕可计算 \(\#\mathrm{SAT}\) 的值，进一步利用概率方法（Valiant-Vazirani 定理）从计数中提取判定信息。关键子步骤：\(\#\mathbf{P}\) 神谕可模拟 \(\oplus\mathbf{P}\)（奇偶性 NP）神谕，而 \(\oplus\mathbf{P}\) 神谕可经随机归约模拟 \(\mathbf{PH}\)。\(\blacksquare\)

\textbf{定义 155.4}（\(\#\mathbf{P}\) 类）：\(\#\mathbf{P}\) 是 NP 问题的解的个数计数问题类。\(\#SAT\)（计算 SAT 的满足赋值数）是 \(\#\mathbf{P}\) 完全的。

\subsection{155.3 概率与交互证明类}\label{ux6982ux7387ux4e0eux4ea4ux4e92ux8bc1ux660eux7c7b}

\textbf{定义 155.5}（概率类）： - \(\mathbf{BPP}\)（Bounded-error Probabilistic Polynomial time）：错误概率 \(< 1/3\) 的概率多项式时间可判定语言 - \(\mathbf{RP}\)：单侧错误（是实例不被错判为否） - \(\mathbf{ZPP} = \mathbf{RP} \cap \mathbf{coRP}\)（零错误概率多项式时间）

\textbf{定理 155.6}（概率放大）：\(\mathbf{BPP}\) 的错误概率可降至 \(2^{-n}\)，通过重复多数投票。

\textbf{证明}：设算法 \(M\) 对 \(x\in L\) 接受概率 \(\ge 2/3\)，对 \(x\notin L\) 接受概率 \(\le 1/3\)。独立运行 \(M\) \(m\) 次，取多数投票。对 \(x\in L\)，拒收次数 \(X\sim\mathrm{Bin}(m,1/3)\) 的期望为 \(m/3\)。由 Chernoff 界 \(\Pr[\text{拒收}>m/2]\le\Pr[X>3m/2\cdot\mathbb{E}[X]]\le(4/e)^{m/12}<2^{-cm}\)。对 \(x\notin L\) 由对称性相同。取 \(m=O(n)\) 即得 \(2^{-n}\) 错误概率。\(\blacksquare\)

\textbf{已知关系}：\(\mathbf{P} \subseteq \mathbf{ZPP} \subseteq \mathbf{RP} \subseteq \mathbf{BPP} \subseteq \mathbf{PSPACE}\)。\(\mathbf{BPP} \stackrel{?}{\subseteq} \mathbf{P}\) 未知。

\textbf{定义 155.6}（交互证明系统 / IP）：验证者 \(V\)（概率多项式时间）与证明者 \(P\) 交换多项式轮消息。\(V\) 接受 \(x \in L\) 的概率 \(> 2/3\)，\(x \notin L\) 的概率 \(< 1/3\)。

\textbf{定理 155.7}（\(\mathbf{IP} = \mathbf{PSPACE}\) / Shamir，1992）：交互证明可验证的恰好是 PSPACE 语言。

\textbf{证明}：\(\mathbf{IP}\subseteq\mathbf{PSPACE}\)：验证者可在 PSPACE 内计算最优策略（通过动态规划搜索所有可能的交互和证明者策略）。\(\mathbf{PSPACE}\subseteq\mathbf{IP}\)：设计交互协议验证 TQBF（真的量化布尔公式，PSPACE 完全的）。将 TQBF 实例算术化：将布尔公式 \(\phi\) 译为低次多项式 \(\tilde{\phi}\)，然后利用求和算符 \(\Sigma_{x\in\{0,1\}}\tilde{\phi}(x)\) 的多项式恒等式检验（LFKN-Shamir 协议）。验证者随机选择域中值并使用低次多项式检验，错误概率 \(\le dq/|\mathbb{F}|\)（\(d\) 为多项式度，\(q\) 为轮数）。通过选择大域 \(\mathbb{F}\) 使错误可忽略。\(\blacksquare\)

\textbf{定义 155.7}（\(\mathbf{PCP}\) 定理 / Probabilistically Checkable Proofs，Arora-Lund-Motwani-Sudan-Szegedy 1992；Arora-Safra 1998）：\(\mathbf{NP} = \mathbf{PCP}[\log n, O(1)]\)。即任何 NP 语言有证明，只需检查常数个随机比特即可高概率验证。其深刻意义在于：它等价于''近似版本的 NP 困难性''------PCP 定理表明，区分可满足公式与每个赋值最多满足 \(1-\varepsilon\) 比例的公式（对某个常数 \(\varepsilon > 0\)）是 NP 难的。这催生了近似困难性理论：MAX-3SAT、MAX-CLIQUE、SET COVER 等经典优化问题在某个常数因子内近似仍是 NP 难的。PCP 定理的证明分两步：代数方法（将 NP 预言转化为多项式求和验证）和组合方法（通过低度检验和局部解码将验证者简化为常数次查询）。PCP 定理是 20 世纪理论计算机科学最重大的成果之一，获得了 2001 年 Gödel 奖。

\textbf{定理 155.8}（PCP 定理与近似困难性）：PCP 定理蕴含许多 NP 完全问题的近似算法的最优性。例如，CLIQUE 不可近似到 \(n^{1-\varepsilon}\)（除非 \(\mathbf{P} = \mathbf{NP}\)）。

\textbf{证明}：PCP 定理 \(\mathbf{NP}=\mathbf{PCP}[O(\log n),O(1)]\) 表明可构造常数查询验证者 \(V\)。Hastad（1996）强化到 \(\mathbf{NP}=\mathbf{PCP}_{1,\varepsilon}[O(\log n),3]\)（3 查询，完美完备性，稳定度 \(\varepsilon=1/2+\delta\)）。对 CLIQUE 问题的 \(n^{1-\varepsilon}\) 不可近似性：将 PCP 验证者转化为''label cover''实例，再经 FGLSS 归约（Feige-Goldwasser-Lovasz-Safra-Szegedy 1992）构造 FGLSS 图 \(G_\phi\)：每个顶点为验证者的一个接受格局（输入、随机串、回答），边连接兼容格局。\(\phi\) 可满足 \(\iff G_\phi\) 有大小为 \(2^R\) 的团（\(R\) 为随机串数），\(\phi\) 不可满足 \(\iff\) 最大团大小 \(<2^R/2\)。通过重复 PCP（并行重复引理）放大差距可推得 \(n^{1-\varepsilon}\) 不可近似。\(\blacksquare\)

\hrulefill

\hrulefill

\hrulefill

\hrulefill

\hrulefill

\hrulefill

\section{第409章：电路复杂性与量子计算}\label{ux7b2c409ux7ae0ux7535ux8defux590dux6742ux6027ux4e0eux91cfux5b50ux8ba1ux7b97}

本章介绍计算复杂性的非均匀模型（电路）和量子计算的基本概念。

\subsection{156.1 电路复杂性}\label{ux7535ux8defux590dux6742ux6027}

\textbf{定义 156.1}（布尔电路）：布尔电路是节点为逻辑门（AND, OR, NOT）的有向无环图。电路族 \(\{C_n\}_{n \in \mathbb{N}}\) 接受语言 \(L\)，如果 \(C_n\) 有 \(n\) 个输入，且 \(C_n(x) = 1 \iff x \in L\)（\(|x| = n\)）。

\textbf{定义 156.2}（电路复杂度类）： - \(\mathbf{P}_{/\text{poly}}\)：存在多项式大小的电路族可判定的语言 - \(\mathbf{NC}\)（Nick's Class）：深度为 \(O(\log^k n)\) 的多项式大小电路族 - \(\mathbf{AC}^0\)：深度为常数的无界扇入电路

\textbf{定理 156.1}：\(\mathbf{P} \subseteq \mathbf{P}_{/\text{poly}}\)。但 \(\mathbf{P}_{/\text{poly}}\) 包含不可判定语言（因为可为每个 \(n\) 硬编码答案）。

\textbf{证明}：\(\mathbf{P}\subseteq\mathbf{P}_{/\text{poly}}\)：任意多项式时间 TM \(M\) 的转移表可在多项式规模电路 \(C_n\) 中直接编码（每步计算对应一层固定深度电路），故 \(M\) 有等价的多项式大小电路族。\(\mathbf{P}_{/\text{poly}}\) 非可判定语言的存在性：所有一元语言 \(\{1^n:n\in S\}\)（\(S\subseteq\mathbb{N}\) 任意）均有线性大小电路（对每个 \(n\)，若 \(n\in S\) 则电路输出常量 1，否则 0）。取 \(S\) 为停机问题的编码即得不可判定语言属于 \(\mathbf{P}_{/\text{poly}}\)。\(\blacksquare\)

\textbf{定理 156.2}（Karp-Lipton 定理，1980）：若 \(\mathbf{NP} \subseteq \mathbf{P}_{/\text{poly}}\)，则 \(\mathbf{PH} = \Sigma_2^P\)（多项式谱系坍塌到第二层）。

\textbf{证明}：假设 \(\mathrm{SAT}\in\mathbf{P}_{/\text{poly}}\)，即存在多项式大小电路族 \(\{C_n\}\) 判定 SAT。对 \(\Pi_2^P\) 语言 \(\{x:\forall y\exists z\,\phi(x,y,z)\}\)，需证其属 \(\Sigma_2^P\)。固定 \(x\)，判定 \(\forall y\exists z\,\phi(x,y,z)\) 是否成立。用 SAT 电路判定器（存在 \(\exists\) 量化）将 \(\forall y\) 的判定提升为 \(\Sigma_2^P\)。形式化地，存在电路 \(C\)（大小多项式）和 NP 验证器 \(V\) 使：\(\forall y\exists z\,\phi(x,y,z)\) 成立 \(\iff\) 存在电路 \(C\)（大小为 \(p(|x|)\)）和见证 \(w\) 满足某条件。因电路的存在可用 \(\exists\) 量化表示，整个表达式化为 \(\Sigma_2^P\)。类似论证可推 \(\Sigma_3^P=\Sigma_2^P\)，全谱系坍塌。\(\blacksquare\)

\textbf{定理 156.3}（电路下界结果）： - \(\mathbf{PARITY} \notin \mathbf{AC}^0\)（Furst-Saxe-Sipser 1981, Ajtai 1983） - \(\mathbf{MOD}_p \notin \mathbf{AC}^0[\mathbf{MOD}_q]\)（\(p \neq q\) 素数，Razborov 1987, Smolensky 1987）

\textbf{证明}：（PARITY \(\notin\mathbf{AC}^0\)）Hastad 的开关引理：深度 \(d\) 的 \(\mathbf{AC}^0\) 电路在随机限制下以高概率坍塌为浅决策树。若 PARITY \(\in\mathbf{AC}^0\)，取 \(\mathbf{AC}^0\) 电路 \(C_n\) 深度为 \(d\)，大小 \(S\)。施加保留 \(t\) 个未固定变量的随机限制，由开关引理，剩余决策树深度以高概率 \(<n-t\)。但 PARITY 的任何深度 \(<n-t\) 的决策树最多依赖 \(<n-t\) 个变量------选择保留变量能区分 PARITY 的输入存在性即可推出矛盾。（MOD\(_p\notin\mathbf{AC}^0[\mathrm{MOD}_q]\)）使用低次多项式近似：\(\mathbf{AC}^0[\mathrm{MOD}_q]\) 电路可用 \(\mathbb{F}_p\) 上低次多项式近似，但 MOD\(_p\) 函数（Sum mod \(p\)）的逼近需要高次多项式（Razborov-Smolensky 多项式方法）。当 \(p\neq q\) 素数时低次近似失败。\(\blacksquare\)

\subsection{156.2 量子计算引论}\label{ux91cfux5b50ux8ba1ux7b97ux5f15ux8bba}

量子计算利用量子力学叠加与纠缠原理处理信息，计算能力严格包含经典计算（\(\mathbf{P} \subseteq \mathbf{BQP}\)），并能在多项式时间内分解整数（Shor 算法），对经典密码学构成根本挑战。

\textbf{定义 156.3}（量子比特 / Qubit）：量子比特的状态是 \(\mathbb{C}^2\) 中的单位向量 \(\alpha|0\rangle + \beta|1\rangle\)（\(|\alpha|^2 + |\beta|^2 = 1\)）。\(n\) 个量子比特的状态在 \(\mathbb{C}^{2^n}\) 中。

\textbf{定义 156.4}（量子门）：量子计算由酉变换实现。基本量子门包括： - Hadamard 门：\(H = \frac{1}{\sqrt{2}}\begin{pmatrix} 1 & 1 \\ 1 & -1 \end{pmatrix}\) - 受控非门（CNOT）：\(|x, y\rangle \mapsto |x, x \oplus y\rangle\) - \(\pi/8\) 门（T 门）：\(T = \begin{pmatrix} 1 & 0 \\ 0 & e^{i\pi/4} \end{pmatrix}\)

\textbf{定理 156.4}（通用量子门集）：\(\{H, CNOT, T\}\) 是一个通用量子门集，即任何酉变换可由此集近似。

\textbf{证明}（Solovay-Kitaev 定理）：首先，\(H\) 和 \(T\) 生成 \(SU(2)\) 的稠密子群（因为 \(H\) 和 \(T\) 对应的旋转角度之比为无理数，它们在 \(SU(2)\) 中生成的群是稠密的）。\(CNOT\) 门加上单量子比特门可生成任意 \(n\)-量子比特酉变换（由两能级酉分解引理）。Solovay-Kitaev 定理进一步保证：任何单量子比特酉变换 \(U\) 可用 \(H\) 和 \(T\) 的 \(O(\log^c(1/\varepsilon))\) 个门以精度 \(\varepsilon\) 近似，其中 \(c \approx 3.97\)。具体地，将 \(U\) 写为 \(SU(2)\) 中的旋转，利用 \(H\) 和 \(T\) 生成的群在 \(SU(2)\) 中的稠密性，通过平衡群论中的生成元分解找到逼近。结合 \(CNOT\) 门处理多量子比特纠缠，整个门集是通用的。\(\blacksquare\)

\textbf{定义 156.5}（量子复杂度类）： - \(\mathbf{BQP}\)：量子多项式时间可判定的语言（错误概率 \(< 1/3\)） - \(\mathbf{QMA}\)：量子 NP 的类比（量子 Merlin-Arthur 证明）

\textbf{定理 156.5}（已知的包含关系）：\(\mathbf{P} \subseteq \mathbf{BPP} \subseteq \mathbf{BQP} \subseteq \mathbf{PSPACE}\)。\(\mathbf{BQP}\) 是否超出 \(\mathbf{P}\) 未知。

\textbf{证明}：\(\mathbf{P}\subseteq\mathbf{BPP}\)：确定算法是概率算法的特例（不使用随机性）。\(\mathbf{BPP}\subseteq\mathbf{BQP}\)：经典概率电路可由量子电路模拟（使用 Hadamard 门产生均匀叠加态模拟随机比特，用 Toffoli 门模拟经典逻辑门）。\(\mathbf{BQP}\subseteq\mathbf{PSPACE}\)：\(n\) 量子比特系统的状态为 \(\mathbb{C}^{2^n}\) 中向量，量子电路的酉矩阵乘法可在 \(\operatorname{poly}(2^n)\) 空间内模拟（Feynman 的路径积分方法），即 PSPACE 空间。在多项式时间步数限制下，总模拟空间为多项式。\(\blacksquare\)

\textbf{定理 156.6}（Shor 算法，1994）：整数因子分解和离散对数可在 \(\mathbf{BQP}\) 中解决。复杂度 \(O(\log^3 N)\)（量子），而经典已知最好算法是亚指数时间。

\textbf{证明}：将因子分解归约为阶寻找（order-finding）：随机选 \(a<N\)，若 \(\gcd(a,N)=1\)，找最小 \(r\) 使 \(a^r\equiv1\pmod{N}\)。若 \(r\) 为偶数且 \(a^{r/2}\not\equiv-1\)，则 \(\gcd(a^{r/2}\pm1,N)\) 是 \(N\) 的非平凡因子。量子部分使用量子 Fourier 变换（QFT）高效计算阶，复杂度 \(O(\log^3 N)\)：制备叠加态 \(\frac{1}{\sqrt{2^n}}\sum_x|x\rangle|a^x\bmod N\rangle\)；应用 QFT 于第一寄存器；测量得 \(k/2^n\approx s/r\)（\(s\) 随机），用连分数展开恢复 \(r\)。QFT 实现为 \(O(n^2)\) 个 Hadamard 和受控旋转门。离散对数类似归约为 QFT 的周期寻找。\(\blacksquare\)

\textbf{定理 156.7}（Grover 算法，1996）：无序数据库搜索可在 \(O(\sqrt{N})\) 量子查询中完成（经典需 \(O(N)\)），这是二次加速。

\textbf{证明}：设神谕 \(O_f|x\rangle=(-1)^{f(x)}|x\rangle\)（\(f(x)=1\) 当 \(x\) 为解）。Grover 算法迭代应用 Grover 算子 \(G=(2|\psi\rangle\langle\psi|-I)O_f\)（\(|\psi\rangle=H^{\otimes n}|0\rangle^{\otimes n}\) 为均匀叠加）。\(G\) 将状态向量在 \(|\alpha\rangle\)（所有非解的均匀叠加）和 \(|\beta\rangle\)（解的均匀叠加）张成的二维平面上旋转 \(2\theta\) 角，其中 \(\sin\theta=\sqrt{M/N}\)（\(M\) 为解数）。\(O(\sqrt{N/M})\) 次迭代后状态接近 \(|\beta\rangle\)。由几何解释，\(G\) 是子空间内的二维旋转，\(O(\sqrt{N})\) 步内达到高概率成功。BBBV 证明此为最优。\(\blacksquare\)

\textbf{定理 156.8}（量子查询下界 / BBBV 定理，1998）：对一般的无序搜索，\(O(\sqrt{N})\) 是最优的量子查询复杂度。

\textbf{证明}：用多项式方法。设量子算法使用 \(T\) 次神谕查询，其最终接受概率可表为多变元次数 \(\le 2T\) 的多项式 \(p(x_1,\dots,x_N)\)（\(x_i\) 为第 \(i\) 个元素是否为解）。对于无序搜索，该多项式在最多含 \(M\) 个解的输入上值 \(\ge2/3\)，在零解输入上值 \(\le1/3\)。由多项式逼近理论，次数至少为 \(\Omega(\sqrt{N/M})\)。Bennett-Bernstein-Brassard-Vazirani 原始证法使用''杂交论证''（hybrid argument）：比较相邻查询的量子态，利用量子信息论界。最优性由 Grover 算法匹配下界证实。\(\blacksquare\)

\hrulefill

\hrulefill

\hrulefill

\hrulefill

\hrulefill

\hrulefill

\section{第410章：计算几何（Computational Geometry）}\label{ux7b2c410ux7ae0ux8ba1ux7b97ux51e0ux4f55computational-geometry}

计算几何研究几何对象的算法设计与复杂性分析。\textbf{核心问题}：（1）\textbf{凸包}------给定 \(n\) 个平面点，Graham 扫描算法以 \(O(n \log n)\) 构造其凸包（下界也是 \(\Omega(n \log n)\)）；（2）\textbf{Voronoi 图与 Delaunay 三角剖分}------平面上 \(n\) 个站点的 Voronoi 图划分空间为最近站点区域，Fortune 算法以 \(O(n \log n)\) 时间构造；对偶的 Delaunay 三角剖分最大化最小角（避免瘦三角形），是有限元网格生成的标准工具；（3）\textbf{线段求交}------Bentley-Ottmann 扫描线算法检测 \(n\) 条线段的全部 \(k\) 个交点，\(O((n+k)\log n)\) 时间；（4）\textbf{点定位}------Kirpatrick 的层次化三角剖分支持 \(O(\log n)\) 时间的平面点定位。\textbf{数值鲁棒性}------浮点舍入误差导致的拓扑不一致性是计算几何的核心挑战，精确几何计算通过任意精度有理算术（如 LEDA, CGAL 库）解决。\textbf{高维计算几何}：\(d\) 维凸包的组合复杂性为 \(O(n^{\lfloor d/2 \rfloor})\)（上界定理），Voronoi 图的复杂度类似增长，产生碰撞维数灾难（curse of dimensionality）的实际瓶颈。近似算法如 \(\epsilon\)-核集（\(\epsilon\)-kernel）和随机投影在机器学习（高维最近邻搜索）中提供折中方案。

\subsection{凸包算法：Graham扫描与Chan算法}\label{ux51f8ux5305ux7b97ux6cd5grahamux626bux63cfux4e0echanux7b97ux6cd5}

\textbf{定义157.1}（平面凸包）：给定平面上的点集 \(P = \{p_1, \ldots, p_n\} \subset \mathbb{R}^2\)，\(P\) 的\textbf{凸包} \(\operatorname{CH}(P)\) 是包含 \(P\) 的最小凸集。等价地，\(\operatorname{CH}(P)\) 是所有包含 \(P\) 的半平面的交集，也是 \(P\) 中所有点的凸组合构成的集合。平面凸包的边界是由 \(P\) 中的点按逆时针方向连接而成的凸多边形。

\textbf{定理157.1}（Graham扫描的正确性，1972）：Graham扫描算法在 \(O(n \log n)\) 时间内正确构造 \(n\) 个平面点的凸包。

\textbf{证明}：算法步骤：（1）选取 \(y\) 坐标最小的点为基准点 \(p_0\)（若有多点取 \(x\) 坐标最小者），\(O(n)\)；（2）将剩余点按关于 \(p_0\) 的极角排序，\(O(n \log n)\)；（3）维护一个栈，依次处理排序后的点。处理点 \(p_k\) 时，若栈顶点 \(p_{i-1}\)、\(p_i\) 和 \(p_k\) 构成''右转''（即 \(\det(p_i - p_{i-1}, p_k - p_i) > 0\)），则 \(p_k\) 入栈；若构成''左转''（\(\det \leq 0\)），则 \(p_i\) 出栈，继续检查新的栈顶两点与 \(p_k\) 的关系，直到满足右转条件后 \(p_k\) 入栈。

\textbf{正确性证明}：扫描过程维护了不变量------栈中的点始终构成当前已处理点集的凸包的下半部分（相对于极角序）。对于每个出栈操作，被弹出的点 \(p_i\) 位于当前栈顶两点 \(p_{i-1}\) 和 \(p_k\) 形成的三角形内部或其边界上，因此不可能属于凸包（因为凸包顶点必为极值点，不能是其他点的凸组合）。算法终止时，栈中的点集恰好是 \(P\) 的凸包顶点（按逆时针顺序）。排序占主导的 \(O(n \log n)\)，扫描为 \(O(n)\)（每个点至多入栈和出栈各一次），故总复杂度 \(O(n \log n)\)。若预先按 \(x\) 坐标排序并两遍扫描（下凸包和上凸包），可避免极角排序中的数值稳定性问题（Andrew 单调链变体）。\(\blacksquare\)

\textbf{定理157.2}（Chan算法的正确性与最优性，1996）：Chan算法（又称''终极凸包算法''）以 \(O(n \log h)\) 时间构造平面凸包，其中 \(h\) 为凸包顶点数（输出规模）。在基于比较的代数决策树模型下，这是最优的时间复杂度。

\textbf{证明}：Chan算法的核心是将 Graham 扫描与 Jarvis 包裹（gift wrapping）结合，利用分组猜测 \(h\) 的技术。将 \(P\) 随机划分为 \(\lceil n/m \rceil\) 个大小为 \(m\) 的子集（\(m\) 为对 \(h\) 的当前猜测），对每个子集用 Graham 扫描（\(O(m \log m)\) 时间）求子凸包。然后执行 Jarvis 包裹步骤------从最右点出发，在 \(O(\lceil n/m \rceil)\) 个子凸包上各做一次 \(O(\log m)\) 的切线查询（二分查找），选出极角最小的支撑点，共 \(h\) 步，总计 \(O((n/m) h \log m)\)。取 \(m = h\) 时总复杂度 \(O(n \log h)\)。由于 \(h\) 事先未知，采用''指数猜测''策略：依次尝试 \(m = 2, 4, 8, \ldots, 2^{\lceil \log_2 h \rceil}\)，几何级数的总代价为 \(\sum_{i=0}^{\lceil \log_2 h \rceil} O(n \log 2^i) = O(n \log h)\)。下界 \(\Omega(n \log h)\) 由排序归约证明：给定 \(n\) 个无序实数的排序问题可在 \(O(n)\) 时间内归约到凸包计算（将实数映射到抛物线 \(y = x^2\) 上，其凸包顶点正好按 \(x\) 坐标排序），而排序的下界 \(\Omega(n \log n)\) 在 \(h = n\) 时给出 \(\Omega(n \log h)\)。\(\blacksquare\)

\subsection{Voronoi图与Delaunay三角化的对偶性}\label{voronoiux56feux4e0edelaunayux4e09ux89d2ux5316ux7684ux5bf9ux5076ux6027}

\textbf{定义157.2}（Voronoi图）：平面上 \(n\) 个站点 \(S = \{s_1, \ldots, s_n\}\) 的 \textbf{Voronoi 图} \(\operatorname{Vor}(S)\) 是将平面划分成 \(n\) 个 Voronoi 区域 \(V(s_i) = \{x \in \mathbb{R}^2 : \|x - s_i\| \leq \|x - s_j\|, \forall j \neq i\}\) 的分划。\(V(s_i)\) 是由 \(s_i\) 与其他站点的中垂线所界定的凸多边形（可能无界）。Voronoi 顶点是到三个或更多站点等距的点。

\textbf{定义157.3}（Delaunay三角化）：\(S\) 的 \textbf{Delaunay 三角化} \(\operatorname{DT}(S)\) 是将 \(S\) 中满足''存在一个过 \(p, q, r\) 的空圆（内部不含 \(S\) 中其他点）``的三元组 \((p, q, r)\) 作为三角形边界的三角剖分。Delaunay 三角化与 Voronoi 图是\textbf{对偶}的：每个 Voronoi 顶点对应一个 Delaunay 三角形；每条 Voronoi 边对应一条 Delaunay 边；每个 Voronoi 区域 \(V(s_i)\) 对应以 \(s_i\) 为顶点的 Delaunay 三角形集合。

\textbf{定理157.3}（Delaunay三角化的对偶性）：\(\operatorname{DT}(S)\) 与 \(\operatorname{Vor}(S)\) 互为平面图的对偶图。若 \(S\) 中的点不共圆（一般位置），则 \(\operatorname{Vor}(S)\) 的所有顶点度均为 3，\(\operatorname{DT}(S)\) 是唯一的三角剖分；若有四点共圆，则对应的 Voronoi 顶点度数为 4，Delaunay 三角化不唯一（四边形可沿不同对角线划分）。

\textbf{定理157.4}（Delaunay三角化的角度最优性，Sibson 1978 / Lambert 1994）：在所有 \(S\) 的三角剖分中，\(\operatorname{DT}(S)\) 最大化最小角------即对任意三角剖分 \(\mathcal{T}\) 和 \(\operatorname{DT}(S)\)，\(\min_{\triangle \in \operatorname{DT}} \alpha_{\min}(\triangle) \geq \min_{\triangle \in \mathcal{T}} \alpha_{\min}(\triangle)\)，其中 \(\alpha_{\min}(\triangle)\) 为三角形的最小内角。

\textbf{证明}：考虑任意非 Delaunay 三角剖分中存在的非法边（即不满足空圆性质的边）。令 \((p, q)\) 为两个相邻三角形 \(\triangle pqr\) 和 \(\triangle pqs\) 的共享边。若 \(r\) 在 \(\triangle pqs\) 的外接圆内（或 \(s\) 在 \(\triangle pqr\) 的外接圆内），则边 \((p, q)\) 是非法的。通过边翻转（edge flip）将 \((p, q)\) 替换为 \((r, s)\) 后，由空圆准则的几何性质可知，翻转后六个角中的最小角严格大于翻转前的最小角。角度向量（六个角从小到大排列）在字典序下严格递增（Lawson, 1977）。由于边翻转操作在有限步内终止（每次翻转向量严格递增，且三角剖分的总数有限），终止状态必为 Delaunay 三角化。终止时所有边均满足空圆条件，最小角达到所有可能三角剖分中的最大值。\(\blacksquare\)

这一角度最优性使得 Delaunay 三角化成为有限元网格生成（FEM）的首选------通过最大化最小角避免了''瘦长''单元导致的数值条件数恶化问题。Fortune 的扫描线算法以 \(O(n \log n)\) 时间和 \(O(n)\) 空间构造 Voronoi 图（进而对偶地得到 Delaunay 三角化），这是该问题在代数决策树模型下的理论最优界。

\subsection{最近邻搜索：k-d树与局部敏感哈希}\label{ux6700ux8fd1ux90bbux641cux7d22k-dux6811ux4e0eux5c40ux90e8ux654fux611fux54c8ux5e0c}

\textbf{定义157.4}（k-d树，Bentley 1975）：\(d\) 维 \textbf{k-d 树}是交替按坐标轴划分空间的二叉树数据结构。根节点按第一维坐标中位数划分为左右子树，下一层按第二维，依此类推（循环所有 \(d\) 维）。每次划分将点集均匀分割，树的深度为 \(O(\log n)\)。

\textbf{定理157.5}（k-d树最近邻搜索的相异度分析）：在 \(d\) 维欧氏空间中，对 \(n\) 个均匀随机分布的点，k-d 树的最近邻搜索期望时间复杂度为 \(O(\log n)\)（\(d\) 固定），但最坏情况为 \(O(d n^{1-1/d})\)。对高维数据（\(d \gg \log n\)），k-d 树的搜索退化为几乎线性扫描------这是''维数灾难''在数据结构中的体现。

\textbf{证明}：搜索过程从根到叶递归进行，维护当前最近距离 \(r_{\text{best}}\)。在每层，优先搜索查询点所在侧的子空间；回溯时，仅当查询点到另一侧超平面的距离 \(<\) 当前 \(r_{\text{best}}\) 时才搜索另一侧子树。对于均匀分布的 \(n\) 个点，期望查询 \(O(2^d \log n)\) 个节点（\(d\) 固定时退化到 \(O(\log n)\)）。当 \(d\) 与 \(\log n\) 可比较时，几乎所有节点都会被访问------因为平均最近邻距离与区域最大距离之比较高维下趋于 1，使得回溯剪枝失效。具体地，在随机投影下，方体中的点到其最近邻的距离期望为 \(\Theta(d^{-1/2})\) 倍边长，而点到区域分隔面的平均距离也为该量级。两者比值不随 \(d\) 增大而趋于零，导致访问节点数指数增长。\(\blacksquare\)

\textbf{定义157.5}（局部敏感哈希 LSH，Indyk-Motwani 1998）：对距离空间 \((X, d)\)，一族哈希函数 \(\mathcal{H}\) 称 \((r, cr, p_1, p_2)\)-敏感，若对任意 \(x, y \in X\)： \[\begin{aligned} d(x, y) \leq r &\implies \Pr_{h \in \mathcal{H}}[h(x) = h(y)] \geq p_1 \\ d(x, y) \geq cr &\implies \Pr_{h \in \mathcal{H}}[h(x) = h(y)] \leq p_2 \end{aligned}\] 其中 \(c > 1\)，\(p_1 > p_2\)。参数 \(\rho = \log(1/p_1) / \log(1/p_2)\) 控制查询复杂度。

\textbf{定理157.6}（LSH查询的相异度下界与实例）：在 \(d\) 维 \(\ell_2\) 空间，基于随机投影的 LSH 族 \(\mathcal{H} = \{h_v(x) = \operatorname{sign}(v^T x)\}\) 是 \((r, cr, p_1, p_2)\)-敏感的，且 \(\rho = 1/c\)。对 Hamming 空间，基于随机比特采样的 LSH 族满足 \(\rho = 1/c\)。LSH 架构通过构造 \(L\) 个哈希表（各使用 \(k\) 个独立哈希函数的 AND-组合）实现：\(L = n^\rho\) 个哈希表，每个表使用 \(k = \log_{1/p_2} n\) 个函数，查询时间为 \(O(n^\rho \log n)\)，空间为 \(O(n^{1+\rho})\)。通过调整参数 \(k\) 和 \(L\)，在精度和速度之间取得了精细权衡------这是高维最近邻搜索中首个突破''维度诅咒''严格下界的算法框架。

最近邻搜索在机器学习（\(k\)-NN分类与回归）、计算机视觉（图像检索）和自然语言处理（语义搜索）中构成基本操作。从k-d树到LSH再到基于图的近似最近邻（HNSW, NSW），计算几何与信息检索的交叉为这些应用提供了理论指导。

\hrulefill

\hrulefill

\hrulefill

\hrulefill

\hrulefill

\hrulefill

\newpage

\chapter{11-数值与计算/02-计算理论/04-计算几何与机器学习.md}\label{ux6570ux503cux4e0eux8ba1ux7b9702-ux8ba1ux7b97ux7406ux8bba04-ux8ba1ux7b97ux51e0ux4f55ux4e0eux673aux5668ux5b66ux4e60.md}

\section{第411章：计算几何}\label{ux7b2c411ux7ae0ux8ba1ux7b97ux51e0ux4f55}

计算几何（Computational Geometry）研究几何对象的离散算法，包括凸包、三角剖分、Voronoi图和Delaunay三角剖分等基本结构的构造与分析。其核心目标是为几何计算问题设计最优时间复杂度算法并给出严格正确性证明。

\subsection{339.1 计算几何基础与凸包算法}\label{ux8ba1ux7b97ux51e0ux4f55ux57faux7840ux4e0eux51f8ux5305ux7b97ux6cd5}

\textbf{定义 339.1}（计算几何的研究范围）：计算几何研究以下基本问题的算法设计与复杂度分析：（i）凸包（convex hull）构造；（ii）点集的Voronoi图与Delaunay三角剖分；（iii）多边形三角剖分；（iv）线段求交；（v）点定位与范围搜索。算法复杂度通常以输入点数 \(n\) 度量，最优界普遍为 \(O(n\log n)\)（由排序下界 \(\Omega(n\log n)\) 归约得到）。

\textbf{定义 339.2}（Graham扫描算法，Graham 1972）：给定平面点集 \(P=\{p_1,\ldots,p_n\}\)，Graham扫描按下列步骤构造凸包：

\begin{enumerate}
\def\labelenumi{\arabic{enumi}.}
\tightlist
\item
  选取 \(y\) 坐标最小的点（若有多个则取 \(x\) 最小者）作为基点 \(p_0\)
\item
  以 \(p_0\) 为极点对所有剩余点进行极角排序，得到序列 \(p_1,\ldots,p_{n-1}\)
\item
  初始化栈 \(S=[p_0,p_1,p_2]\)
\item
  依次处理 \(p_3,\ldots,p_{n-1}\)：当栈顶两点与当前点构成右转（非左转）时弹出栈顶，随后压入当前点
\end{enumerate}

\textbf{定理 339.1}（Graham扫描的复杂度与正确性）：Graham扫描算法在 \(O(n\log n)\) 时间内正确构造 \(P\) 的凸包。

\textbf{证明}：步骤1-2的排序占 \(O(n\log n)\)。步骤3-4中每个点至多入栈一次、出栈一次，共 \(O(n)\)。总复杂度 \(O(n\log n)\)。正确性：算法维护栈中点为当前已处理点集凸包的逆时针顶点序列，左转判定保证凸性，最终输出即凸包边界。\(\blacksquare\)

\textbf{定理 339.2}（凸包问题的下界）：平面 \(n\) 点的凸包构造需要 \(\Omega(n\log n)\) 时间（在代数判定树模型下）。

\textbf{证明}：将排序问题归约为凸包问题。将 \(n\) 个数 \(\{a_i\}\) 映射到抛物线 \((a_i, a_i^2)\) 上的点集 \(P = \{(a_i, a_i^2) : i = 1,\dots,n\}\)。由于抛物线 \(y = x^2\) 是凸函数，\(P\) 中所有点均位于其凸包的边界上。凸包 \(\operatorname{CH}(P)\) 的逆时针顶点序列按 \(x\) 坐标递增排列，恰好给出排序后的 \(\{a_i\}\) 序列。因此，若存在 \(o(n\log n)\) 的凸包算法，将其应用于 \(P\) 即可在 \(o(n\log n)\) 时间内完成排序，与排序问题的 \(\Omega(n\log n)\) 下界矛盾。故凸包构造需要 \(\Omega(n\log n)\) 时间。\(\blacksquare\)

\subsection{339.2 Delaunay三角剖分与Voronoi图}\label{delaunayux4e09ux89d2ux5256ux5206ux4e0evoronoiux56fe}

\textbf{定义 339.3}（Delaunay三角剖分与空圆性质）：平面点集 \(P\) 的 \textbf{Delaunay三角剖分} \(\mathrm{DT}(P)\) 满足：对任意三角形 \(\triangle abc\in\mathrm{DT}(P)\)，其外接圆不包含 \(P\) 中任何其他内点（\textbf{空圆性质}）。\(\mathrm{DT}(P)\) 在非退化（无四点共圆）情形下唯一。等价地，\(\mathrm{DT}(P)\) 是Voronoi图的对偶图：连Voronoi区域相邻的点对即得Delaunay边。

\textbf{定义 339.4}（Voronoi图）：点集 \(P\) 的 \textbf{Voronoi图} \(\mathrm{Vor}(P)\) 将平面划分为 \(n\) 个Voronoi区域：

\[V(p_i) = \{x\in\mathbb{R}^2 : \|x-p_i\| \leq \|x-p_j\|,\; \forall j\neq i\}\]

\(V(p_i)\) 是以 \(p_i\) 为最近邻的所有点的闭集。Voronoi图的每条边是两点的垂直平分线的线段，每个顶点是三个生成点外接圆的圆心。 Voronoi图是对空间最近邻查询问题的完备数据结构。

\textbf{定理 339.3}（Delaunay-Voronoi对偶性）：\(\mathrm{DT}(P)\) 与 \(\mathrm{Vor}(P)\) 互为对偶图：Delaunay三角剖分的每条边对应于Voronoi图中两相邻区域的分隔边（垂直平分线的线段）；每个Delaunay三角形的外接圆中心恰为一个Voronoi顶点（三个Voronoi区域的交点）。

\textbf{证明}：对偶性的基本对应：（1）Voronoi 边 \(\leftrightarrow\) Delaunay 边。设 \(p_i,p_j\) 的 Voronoi 区域 \(V(p_i),V(p_j)\) 相邻，共享 Voronoi 边 \(e\)。\(e\) 上的点与 \(p_i,p_j\) 等距，属 \(\overline{p_i p_j}\) 的垂直平分线。这是 \(p_i,p_j\) 为 Delaunay 边的充要条件（存在通过 \(p_i,p_j\) 的空圆）。（2）Voronoi 顶点 \(\leftrightarrow\) Delaunay 三角形。设 \(v\) 是 \(V(p_i),V(p_j),V(p_k)\) 的公共顶点，则 \(v\) 与 \(p_i,p_j,p_k\) 等距。以 \(v\) 为心、\(\|v-p_i\|\) 为半径的圆过 \(p_i,p_j,p_k\) 且不含任何其他点（否则 Voronoi 区域结构破坏），故 \(\triangle p_i p_j p_k\) 满足空圆性质，是 Delaunay 三角形。完全的双射对应成立。\(\blacksquare\)

\textbf{定理 339.4}（Fortune扫描线算法，Fortune 1987）：Voronoi图可在 \(O(n\log n)\) 时间内通过扫描线算法构造，核心数据结构为海滩线（beach line，由抛物线弧组成）和事件队列（site events + circle events）。

\textbf{证明}：Fortune 算法使用水平扫描线从顶向下移动。海滩线由已处理站点的抛物线弧组成（满足与扫描线等距条件）。数据结构：事件队列（优先队列）存储站点事件（扫描线经过某站点时触发）和圆事件（海滩线三段连续弧对应的圆最低点------对应 Voronoi 顶点）。当扫描线触及圆事件时，相应弧消失并输出 Voronoi 顶点。每个事件处理 \(O(\log n)\) 时间（平衡 BST 操作）。站点数 \(n\)，圆事件数 \(\le 2n-5\)。总复杂度 \(O(n\log n)\)。正确性依赖于海滩线结构正确表示 Voronoi 图的增量构造。\(\blacksquare\)

\subsection{339.3 多边形三角剖分}\label{ux591aux8fb9ux5f62ux4e09ux89d2ux5256ux5206}

\textbf{定义 339.5}（简单多边形的三角剖分）：简单多边形 \(P\)（无自交的闭多边形链）的一个\textbf{三角剖分}是将 \(P\) 内部划分为由 \(P\) 的顶点构成的三角形集合，这些三角形互不重叠且覆盖整个 \(P\) 的内部。每条添加的对角线连接 \(P\) 的两个不相邻顶点且完全位于 \(P\) 内部。

\textbf{定理 339.5}（Fisk的3-着色证明，1978）：任意简单多边形的三角剖分对偶图（对偶图的节点对应三角形，边对应对角线）是一棵树，从而三角剖分图的顶点可3-着色。由鸽巢原理，存在单色顶点集构成美术馆问题的 \(O(n)\) 上界解。

\textbf{证明}：三角剖分的对偶图是无环的（否则回路围出多边形内部的一个洞，与简单性矛盾），且连通性蕴含其为树。对偶树的叶子对应含两条多边形边的三角形（``耳朵''）。从叶子开始向树根方向贪婪着色：取颜色 \(\{0,1,2\}\)，根据树结构 BFS 分配颜色，每个三角形三个顶点分配三种不同颜色。由于相邻三角形共享边，颜色分配可延拓到全三角剖分。鸽巢原理：\(3n\) 个顶点着色（每个三角形各有三色），某色出现 \(\le\lfloor n/3\rfloor\) 次。在该色顶点处放置守卫可覆盖所有三角形，故 \(\lfloor n/3\rfloor\) 个守卫充足。\(\blacksquare\) 特别地，\(n\) 个顶点的简单多边形总存在三角剖分，且三角剖分中恰含 \(n-2\) 个三角形和 \(n-3\) 条对角线。

\textbf{定理 339.6}（单调多边形的 \(O(n)\) 三角剖分）：每个简单多边形可分解为 \(y\)-单调多边形，且单调多边形可在 \(O(n)\) 时间内完成三角剖分（通过栈维护已处理顶点的凸链）。结合多边形三角剖分分解算法，总复杂度为 \(O(n\log n)\)。

\textbf{证明}：单调多边形的三角剖分算法：设顶点按 \(y\) 坐标降序排列（上下链）。从顶部开始，栈维护''等待连接''的顶点。当处理第 \(i\) 个顶点 \(v_i\) 时，若 \(v_i\) 与栈顶在异侧，则 \(v_i\) 与栈中所有剩余点相连并清空栈（仅保留上一点和 \(v_i\)）。若同侧，自顶向下连接 \(v_i\) 与栈中顶点，直到遇见不可见顶点（连接线在多边形外）。栈中最多 \(n\) 个点，每个点最多入栈一次、出栈一次，总 \(O(n)\)。多边形分解为单调多边形可使用梯形分解法（扫描线），复杂度 \(O(n\log n)\)。\(\blacksquare\)

\textbf{定理 339.7}（三角剖分的存在性，Meisters 1975）：每个简单多边形都至少存在两个''耳朵''（连续三个顶点构成的三角形完全位于多边形内部且不含其他顶点）。通过递归切除耳朵，可在 \(O(n^2)\) 时间内完成三角剖分。

\textbf{证明}：耳朵定理的证明：对 \(n\) 归纳。取简单多边形 \(P\) 的三角剖分，其对偶树至少有 2 个叶子，每个叶子对应一个耳朵。直接构造：若 \(P\) 有 \(n=3\) 个顶点则本身是三角形。\(n>3\) 时，由双耳朵定理，\(P\) 中必存在两个连续顶点 \(v_i,v_{i+1},v_{i+2}\) 构成耳朵。切下 \(\triangle v_i v_{i+1} v_{i+2}\) 得 \(n-1\) 顶点多边形。重复 \(n-3\) 次共 \(O(n^2)\) 时间（每次需 \(O(n)\) 检测耳朵）。Chazelle（1991）改进为 \(O(n)\)，但极其复杂。\(\blacksquare\) \(n=3\) 时平凡，\(n>3\) 时由耳朵定理保证存在耳朵，切下后对 \(n-1\) 个顶点的多边形递归进行。

\textbf{定理 339.8}（Delaunay三角剖分的最大化最小角性质）：在所有可能的三角剖分中，Delaunay三角剖分最大化最小角（即最大化所有三角形中最小的内角）。此性质使得Delaunay三角剖分在有限元网格生成和曲面重建中具有天然优势------避免出现极端的狭长三角形。

\textbf{证明}：考虑任意四边形 \(ABCD\)，其有两种三角剖分：对角线 \(AC\) 或 \(BD\)。设 \(ABCD\) 为凸四边形，两条对角线交于点 \(O\)。由 Thales 定理，\(AC\) 为 Delaunay 边当且仅当 \(D\) 在 \(\triangle ABC\) 的外接圆外（或 \(A\) 和 \(C\) 满足空圆性质）。比较两种剖分的最小角：若 \(D\) 在 \(\triangle ABC\) 的外接圆内，则翻转对角线 \(AC \to BD\) 会增大最小角。具体地，\(\angle BAD\) 和 \(\angle BCD\) 在翻转后分别增大为 \(\angle BAC\) 和 \(\angle BDC\) 的对应角。由空圆性质，每个 Delaunay 边满足其两端点的外接圆不含其他点。若存在非 Delaunay 边，则其关联四边形必有一个顶点在另一三角形外接圆内，从而翻转该边严格增大最小角。由于有限点集只有有限种三角剖分，最大化最小角的过程必然终止于 Delaunay 三角剖分，且此时无任何边可翻转以增大最小角。\(\blacksquare\)

\textbf{定理 339.9}（Delaunay三角剖分的复杂度）：\(n\) 个点的Delaunay三角剖分包含至多 \(2n-5\) 个三角形和 \(3n-6\) 条边（\(n\geq 3\)），此界由Euler公式 \(V-E+F=2\) 结合平面图性质得出。通过随机增量构造算法（randomized incremental construction），Delaunay三角剖分可在期望 \(O(n\log n)\) 时间内完成。

\textbf{证明}：Delaunay 三角剖分是平面图。设外部面有 \(h\) 条边（凸包边数），\(h\le n\)。Euler 公式：\(n-E+F=2\)。每个三角形贡献 3 条边-面关联，内部边被 2 个面共享，外部边各属 1 个面：\(3F=2E-h\)。解得 \(E=3n-3-h\)，\(F=2n-2-h\)。最坏情形 \(h=3\) 得 \(E=3n-6\)，\(F=2n-4\)。若允许退化，三角形数为 \(2n-5\)（含外部面的近似）。随机增量构造：逐点插入，维护当前 Delaunay 三角剖分。每点插入期望花 \(O(\log n)\) 时间在点定位结构中查找包容三角形，然后局部重连（edge flip）期望 \(O(1)\) 次。总期望 \(O(n\log n)\)。\(\blacksquare\)

\hrulefill

\section{第412章：计算机代数（Computer Algebra）}\label{ux7b2c412ux7ae0ux8ba1ux7b97ux673aux4ee3ux6570computer-algebra}

计算机代数研究符号计算的理论和算法，区别于数值计算的关键在于保持精确表示（如整数的任意精度、多项式的符号形式）。

\subsection{340.1 Grobner基的正式定义}\label{grobnerux57faux7684ux6b63ux5f0fux5b9aux4e49}

\textbf{定义 340.1}（项序 / monomial ordering）：多项式环 \(k[x_1,\ldots,x_n]\) 上的\textbf{项序} \(\prec\) 是全体单项式集合上的全序，满足：（i）\(1\prec x^\alpha\) 对所有 \(\alpha\neq 0\)；（ii）若 \(x^\alpha\prec x^\beta\)，则 \(x^{\alpha+\gamma}\prec x^{\beta+\gamma}\)。常用项序包括：字典序（lex）、分次字典序（grlex）、分次反字典序（grevlex）。

\textbf{定义 340.2}（Grobner基）：设 \(I\subset k[x_1,\ldots,x_n]\) 为非零理想，固定项序 \(\prec\)。有限子集 \(G=\{g_1,\ldots,g_t\}\subset I\setminus\{0\}\) 称为 \(I\) 的 \textbf{Grobner基}，若首项理想满足

\[\langle \mathrm{LT}(G)\rangle = \langle \mathrm{LT}(I)\rangle\]

其中 \(\mathrm{LT}(f)\) 是 \(f\) 关于 \(\prec\) 的首项。等价地，对任意 \(f\in I\setminus\{0\}\)，存在 \(g_i\in G\) 使得 \(\mathrm{LT}(g_i)\mid\mathrm{LT}(f)\)。

\textbf{定义 340.3}（\(S\)-多项式）：对非零多项式 \(f,g\)，其 \textbf{\(S\)-多项式}定义为

\[S(f,g) = \frac{\mathrm{LCM}(\mathrm{LM}(f),\mathrm{LM}(g))}{\mathrm{LT}(f)}\cdot f - \frac{\mathrm{LCM}(\mathrm{LM}(f),\mathrm{LM}(g))}{\mathrm{LT}(g)}\cdot g\]

\(S(f,g)\) 消去了 \(f\) 和 \(g\) 的首项，用于检测 \(G\) 是否为Grobner基。

\subsection{340.2 Buchberger算法}\label{buchbergerux7b97ux6cd5}

\textbf{定理 340.1}（Buchberger判定，1965）：有限生成集 \(G\subset I\setminus\{0\}\) 是 \(I\) 的Grobner基当且仅当对所有 \(g_i,g_j\in G\)，\(S(g_i,g_j)\) 被 \(G\) 约化到零（即 \(S(g_i,g_j)\xrightarrow{G}_+ 0\)）。

\textbf{证明}：必要性：若\(G\)为Grobner基，\(S(g_i,g_j)\in I\)，其首项被\(\langle\mathrm{LT}(G)\rangle\)中某元素整除，故被\(G\)约化到\(0\)。充分性：设所有\(S\)-多项式被\(G\)约化到\(0\)。对任意\(f\in I\)，\(f=\sum_i h_i g_i\)（\(g_i\in G\)）。若\(\max_i\mathrm{LM}(h_i g_i)\)不在\(\langle\mathrm{LT}(G)\rangle\)中，则存在两个项\(h_i g_i\)和\(h_j g_j\)的首项消去------这正是\(S(g_i,g_j)\)产生的消去。由\(S(g_i,g_j)\xrightarrow{G}_+ 0\)，此消去可被\(G\)的更低次项表示。归纳论证可得\(\mathrm{LT}(f)\in\langle\mathrm{LT}(G)\rangle\)，故\(G\)为Grobner基。\(\blacksquare\)

\textbf{算法 340.1}（Buchberger算法）： 1. 输入：理想 \(I=\langle f_1,\ldots,f_s\rangle\)，项序 \(\prec\)；输出：\(I\) 的Grobner基 \(G\) 2. \(G\leftarrow\{f_1,\ldots,f_s\}\) 3. 对每对 \(p,q\in G\)（\(p\neq q\)），计算 \(r=S(p,q)\xrightarrow{G}_+ r\)（被 \(G\) 约化的范式） 4. 若 \(r\neq 0\)，则 \(G\leftarrow G\cup\{r\}\)，返回步骤3 5. 否则输出 \(G\)

\textbf{定理 340.2}（Buchberger算法的终止性）：Buchberger算法在有限步内终止。

\textbf{证明}：每次添加非零约化多项式时，对应的首项单项式 \(x^\alpha=\mathrm{LM}(r)\) 不被当前 \(\langle\mathrm{LT}(G)\rangle\) 中的任何首项整除。因此 \(\langle\mathrm{LT}(G)\rangle\) 严格增大，形成升链。由Dickson引理（\(k[x_1,\ldots,x_n]\) 中单项式理想的升链必有限），算法在有限步内停止。\(\blacksquare\)

\subsection{340.3 消去理论与方程组求解}\label{ux6d88ux53bbux7406ux8bbaux4e0eux65b9ux7a0bux7ec4ux6c42ux89e3}

\textbf{定理 340.3}（消去定理）：设 \(G\) 是理想 \(I\subset k[x_1,\ldots,x_n]\) 关于消去序（如 \(x_1\succ\cdots\succ x_n\) 的字典序）的Grobner基。则对每个 \(l\)（\(0\leq l\leq n\)），\(l\)-次消去理想 \(I_l=I\cap k[x_{l+1},\ldots,x_n]\) 的Grobner基恰为 \(G_l=G\cap k[x_{l+1},\ldots,x_n]\)。

\textbf{证明}：设\(f\in I_l=I\cap k[x_{l+1},\ldots,x_n]\)。因\(G\)为Grobner基，\(\mathrm{LT}(f)\)被某个\(\mathrm{LT}(g)\)（\(g\in G\)）整除。但\(f\)不含\(x_1,\ldots,x_l\)，\(\mathrm{LT}(f)\)亦然。消去序的性质确保\(\mathrm{LT}(g)\)也必不含\(x_1,\ldots,x_l\)（否则\(\mathrm{LT}(g)\)含\(x_i\)（\(i\le l\)）而\(\mathrm{LT}(f)\)不含，整除不可能）。故\(g\in G_l\)。\(G_l\)生成\(\langle\mathrm{LT}(I_l)\rangle\)，因此\(G_l\)是\(I_l\)的Grobner基。\(\blacksquare\)

由此，解多项式方程组 \(f_1=\cdots=f_s=0\) 的步骤为：（i）计算 \(I=\langle f_1,\ldots,f_s\rangle\) 关于字典序的Grobner基 \(G\)；（ii）\(G_{n-1}\subset k[x_n]\) 为单变量多项式，可直接求解 \(x_n\) 的值；（iii）将 \(x_n\) 的值逐一代入 \(G_{n-2}\subset k[x_{n-1},x_n]\) 以求解 \(x_{n-1}\)，依此类推（回代过程）。此即\textbf{Grobner基的消去解法}。

\textbf{定义 340.4}（归约Grobner基 / reduced Grobner basis）：Grobner基 \(G\) 称为\textbf{归约的}，若对任意 \(g\in G\)，（i）\(\mathrm{LC}(g)=1\)（首项系数为1）；（ii）\(g\) 的任意非首项单项式不被 \(\langle\mathrm{LT}(G\setminus\{g\})\rangle\) 中的任何元素整除。对固定项序，每个理想具有唯一的归约Grobner基。

\textbf{定理 340.4}（扩张定理 / Extension Theorem）：设 \(I=\langle f_1,\ldots,f_s\rangle\subset k[x_1,\ldots,x_n]\)，\(I_1=I\cap k[x_2,\ldots,x_n]\) 为其第一个消去理想。若 \((a_2,\ldots,a_n)\in V(I_1)\)（即属于 \(I_1\) 的代数簇），且 \((a_2,\ldots,a_n)\) 不使 \(I\) 中 \(x_1\) 最高次项系数的Grobner基元素同时为零，则存在 \(a_1\) 使得 \((a_1,a_2,\ldots,a_n)\in V(I)\)。扩张定理保证了消去-回代过程的每一步均可完成，条件是有限多个''坏点''需要单独处理。

\textbf{证明}：将\(I\)的元素视为\(x_1\)的多项式：\(f=c_0(a_2,\ldots,a_n)x_1^d+\cdots\)（\(c_0\neq 0\)）。条件确保存在\(g\in G\)（\(I\)关于字典序的Grobner基）使得\(g(a_2,\ldots,a_n,x_1)\)不恒为零。\(g\)的\(x_1\)次数为\(e\)。由代数学基本定理，\(g(a_2,\ldots,a_n,x_1)=0\)（作为\(x_1\)的多项式）在代数闭域\(k\)中至少有一根\(a_1\)。需验证\((a_1,\ldots,a_n)\in V(I)\)：若不然，存在\(f\in I\)使\(f(a_1,\ldots,a_n)\neq 0\)。但\(f\)被\(G\)约化到\(0\)，代入\((a_1,\ldots,a_n)\)后约化链保持，矛盾。故\((a_1,\ldots,a_n)\in V(I)\)。\(\blacksquare\)

\textbf{定理 340.5}（理想的维数与Hilbert函数）：设 \(I\subset k[x_1,\ldots,x_n]\) 为理想，\(G\) 为其Grobner基。\(I\) 的仿射Hilbert函数 \(HF_I(s)=\dim_k(k[x_1,\ldots,x_n]_{\leq s}/I_{\leq s})\) 对所有充分大的 \(s\) 为多项式，其次数等于代数簇 \(V(I)\) 的维数。Grobner基的首项理想 \(\langle\mathrm{LT}(I)\rangle\) 决定了 \(HF_I\)，从而Grobner基理论提供了计算代数簇维数和次数的算法途径。

\textbf{证明}：\(\langle\mathrm{LT}(I)\rangle\)为单项式理想。\(k[x_1,\ldots,x_n]_{\le s}/\langle\mathrm{LT}(I)\rangle_{\le s}\)的维数等于次数\(\le s\)且不被\(\langle\mathrm{LT}(I)\rangle\)中任何单项式整除的单项式个数。Macaulay定理：\(HF_I(s)=HF_{\langle\mathrm{LT}(I)\rangle}(s)\)。\(\langle\mathrm{LT}(I)\rangle\)的补集为有限个坐标子空间的并，其渐近计数为多项式（Hilbert多项式），次数\(=\dim V(I)\)（由Noether正规化引理）。系数（次数）为\(V(I)\)的射影次数。\(\blacksquare\)

\textbf{定理 340.6}（归约Grobner基的唯一性）：对固定项序 \(\prec\)，每个多项式理想 \(I\) 具有唯一的归约Grobner基。此唯一性使得Grobner基成为理想的''规范形''表示，类似于线性代数中的行阶梯形。

\textbf{证明}：存在性：对任意Grobner基\(G\)，逐元素归约（用\(G\setminus\{g\}\)归约\(g\)的非首项）得到归约Grobner基。唯一性：设\(G\)和\(G'\)为\(I\)的两个归约Grobner基。\(\langle\mathrm{LT}(G)\rangle=\langle\mathrm{LT}(I)\rangle=\langle\mathrm{LT}(G')\rangle\)。归约性要求\(\mathrm{LT}(G)\)和\(\mathrm{LT}(G')\)均为极小生成元集------它们构成\(\langle\mathrm{LT}(I)\rangle\)的唯一极小生成元集（等价于除法的不可约基）。若\(g\in G\)，\(\mathrm{LT}(g)\)对应\(\mathrm{LT}(g')\in\mathrm{LT}(G')\)。\(g-g'\)被\(G\)（或\(G'\)）约化到\(0\)，归约性迫使\(g=g'\)。\(\blacksquare\)

\hrulefill

\section{第413章：机器学习理论基础}\label{ux7b2c413ux7ae0ux673aux5668ux5b66ux4e60ux7406ux8bbaux57faux7840}

机器学习的数学理论基础由统计学习理论（Vapnik-Chervonenkis）、优化理论和泛函分析（核方法）共同构成。本章建立PAC可学习性的形式化框架，推导泛化界，并系统阐述核方法与支持向量机的数学结构。

\subsection{341.1 统计学习理论框架}\label{ux7edfux8ba1ux5b66ux4e60ux7406ux8bbaux6846ux67b6}

\textbf{定义 341.1}（统计学习的基本设定）：设 \(\mathcal{X}\) 为输入空间，\(\mathcal{Y}\) 为输出空间（分类问题中 \(\mathcal{Y}=\{0,1\}\)，回归问题中 \(\mathcal{Y}=\mathbb{R}\)）。样本 \(z=(x,y)\in\mathcal{Z}=\mathcal{X}\times\mathcal{Y}\) 服从未知联合分布 \(\mathcal{D}\)。给定损失函数 \(\ell:\mathcal{Y}\times\mathcal{Y}\to\mathbb{R}_+\) 和假设类 \(\mathcal{H}\subset\mathcal{Y}^{\mathcal{X}}\)，\textbf{期望风险}与\textbf{经验风险}分别定义为

\[R(h) = \mathbb{E}_{(x,y)\sim\mathcal{D}}\left[\ell(h(x),y)\right], \quad \widehat{R}_n(h) = \frac{1}{n}\sum_{i=1}^n \ell(h(x_i),y_i)\]

学习目标是基于 \(n\) 个 i.i.d. 样本 \(\{(x_i,y_i)\}_{i=1}^n\)，找到 \(h\in\mathcal{H}\) 使得 \(R(h)-\inf_{h'\in\mathcal{H}}R(h')\) 以高概率小于 \(\varepsilon\)。

统计学习理论的核心问题是：在什么条件下经验风险最小化（ERM）能够一致地逼近期望风险最小化？这涉及两个关键概念：\textbf{泛化误差}（generalization gap）\(R(h) - \widehat{R}_n(h)\) 和\textbf{逼近误差}（approximation error）\(\inf_{h'\in\mathcal{H}} R(h') - R^*\)（\(R^*\) 为 Bayes 最优风险）。总误差可分解为 \(R(\hat{h}) - R^* = \underbrace{(R(\hat{h}) - \inf_{h'\in\mathcal{H}} R(h'))}_{\text{估计误差}} + \underbrace{(\inf_{h'\in\mathcal{H}} R(h') - R^*)}_{\text{逼近误差}}\)。估计误差由假设类 \(\mathcal{H}\) 的复杂度（如 VC 维数或 Rademacher 复杂度）和样本量 \(n\) 共同决定；逼近误差取决于 \(\mathcal{H}\) 的表达能力------\(\mathcal{H}\) 越大，逼近误差越小，但估计误差越大。这一\textbf{偏差-方差权衡}（bias-variance tradeoff）是统计学习理论的核心洞察。VC 理论（Vapnik-Chervonenkis）给出了有限 VC 维假设类的一致收敛保证：\(\mathbb{P}[\sup_{h\in\mathcal{H}} |R(h) - \widehat{R}_n(h)| > \varepsilon] \leq 4 \Pi_{\mathcal{H}}(2n) e^{-n\varepsilon^2/8}\)，其中 \(\Pi_{\mathcal{H}}(m)\) 为增长函数（growth function）。该不等式是 PAC 可学习性理论的基础，确立了有限 VC 维是二分类问题可 PAC 学习的充要条件。

\subsection{341.2 PAC可学习性与VC维}\label{pacux53efux5b66ux4e60ux6027ux4e0evcux7ef4}

\textbf{定义 341.2}（PAC可学习性，Valiant 1984）：假设类 \(\mathcal{H}\) 称为 \textbf{PAC可学习的}，若存在算法 \(\mathcal{A}\) 和函数 \(m_{\mathcal{H}}:(0,1)^2\to\mathbb{N}\)，使得对任意 \(\varepsilon,\delta\in(0,1)\)、任意分布 \(\mathcal{D}\)、任意目标概念 \(c\in\mathcal{H}\)，当样本数 \(n\geq m_{\mathcal{H}}(\varepsilon,\delta)\) 时，以概率至少 \(1-\delta\) 有 \(R(\mathcal{A}(S))\leq\varepsilon\)。

\textbf{定义 341.3}（VC维数 / Vapnik-Chervonenkis Dimension）：对二分类假设类 \(\mathcal{H}\subset\{0,1\}^{\mathcal{X}}\)，其 \textbf{VC维数} 定义为

\[\mathrm{VCdim}(\mathcal{H}) = \max\{m\in\mathbb{N} : \exists\,x_1,\ldots,x_m\in\mathcal{X},\; |\{(h(x_1),\ldots,h(x_m)):h\in\mathcal{H}\}| = 2^m\}\]

即能被 \(\mathcal{H}\) \textbf{打散}（shatter）的最大点集大小。若无上界，记 \(\mathrm{VCdim}(\mathcal{H})=\infty\)。

\textbf{定理 341.1}（Sauer-Shelah引理，1972）：设 \(\mathrm{VCdim}(\mathcal{H})=d<\infty\)。则对任意 \(m\) 个点的集合，\(\mathcal{H}\) 在其上诱导的不同标号数不超过

\[\Pi_{\mathcal{H}}(m) \leq \sum_{i=0}^d \binom{m}{i}\]

当 \(m>d\) 时，\(\Pi_{\mathcal{H}}(m)\leq (em/d)^d\)。

\textbf{证明}：对 \(m\) 归纳。设 \(\Pi_{\mathcal{H}}(m)\) 为 \(\mathcal{H}\) 在 \(m\) 个点上的增长率函数。固定一点 \(x\)，将 \(\mathcal{H}\) 的限制划分为''成对''子族 \(\mathcal{H}'\cup\mathcal{H}''\)（\(x\) 的标号不同时两者均有代表）和''单一''子族。由归纳假设得递推 \(\Pi_{\mathcal{H}}(m)\leq\Pi_{\mathcal{H}-\{x\}}(m-1)+\Pi_{\mathcal{H}'}(m-1)\leq\sum_{i=0}^{d}\binom{m-1}{i}+\sum_{i=0}^{d-1}\binom{m-1}{i}=\sum_{i=0}^{d}\binom{m}{i}\)。\(\blacksquare\)

\textbf{定理 341.2}（VC维泛化界，Vapnik-Chervonenkis 1971）：设 \(\mathrm{VCdim}(\mathcal{H})=d\)。对 \(0\)-\(1\) 损失下的二分类问题，以概率至少 \(1-\delta\)，对所有 \(h\in\mathcal{H}\)，

\[R(h) \leq \widehat{R}_n(h) + \sqrt{\frac{2d\log\frac{en}{d} + 2\log\frac{2}{\delta}}{n}}\]

\textbf{证明}：由对称化引理（Symmetrization Lemma），对任意 \(h \in \mathcal{H}\)，以概率至少 \(1-\delta\) 有： \[R(h) - \widehat{R}_n(h) \leq 2\mathcal{R}_n(\mathcal{H}) + \sqrt{\frac{\log(1/\delta)}{2n}},\] 其中 \(\mathcal{R}_n(\mathcal{H}) = \mathbb{E}_{\sigma, S}[\sup_{h \in \mathcal{H}} \frac{1}{n}\sum_{i=1}^n \sigma_i h(x_i)]\) 为 Rademacher 复杂度（\(\sigma_i \in \{\pm 1\}\) 为独立 Rademacher 变量）。由 Massart 有限类引理，对大小为 \(|\mathcal{H}_{|S}|\) 的限制类，\(\mathcal{R}_n(\mathcal{H}) \leq \sqrt{2\log|\mathcal{H}_{|S}|/n}\)。由 Sauer-Shelah 引理，\(\mathrm{VCdim}(\mathcal{H}) = d\) 时，\(\mathcal{H}\) 在 \(n\) 个点上的限制类大小 \(|\mathcal{H}_{|S}| \leq \sum_{i=0}^d \binom{n}{i} \leq (en/d)^d\)。代入得 \(\mathcal{R}_n(\mathcal{H}) \leq \sqrt{2d\log(en/d)/n}\)。将此界代入对称化引理，并取 \(\delta\) 为 \(2\delta\) 以覆盖 \(\mathcal{H}\) 中所有函数（使用联合界），得到所述的 VC 泛化界。\(\blacksquare\)

\textbf{定理 341.3}（ERM的一致性）：若 \(\mathcal{H}\) 具有有限VC维数，则经验风险最小化（ERM）是一致的：当 \(n\to\infty\) 时，\(\widehat{R}_n(\hat{h}_{\mathrm{ERM}})\to\inf_{h\in\mathcal{H}}R(h)\) 依概率成立，且 \(\hat{h}_{\mathrm{ERM}}\) 的期望风险以 \(O(\sqrt{d\log n/n})\) 速率收敛。

\textbf{证明}：设 \(h^*=\arg\min_{h\in\mathcal{H}}R(h)\)。由 ERM 定义 \(\widehat{R}_n(\hat{h}_{\mathrm{ERM}})\le\widehat{R}_n(h^*)\)。由 VC 泛化界，以概率 \(\ge1-\delta\)：\(R(\hat{h}_{\mathrm{ERM}})\le\widehat{R}_n(\hat{h}_{\mathrm{ERM}})+\varepsilon_n\le\widehat{R}_n(h^*)+\varepsilon_n\le R(h^*)+2\varepsilon_n\)，其中 \(\varepsilon_n=O(\sqrt{d\log n/n})\)。取 \(n\to\infty\)，\(\varepsilon_n\to0\)，故 \(R(\hat{h}_{\mathrm{ERM}})\to R(h^*)\) 依概率收敛。收敛速率 \(O(\sqrt{d\log n/n})\) 源自 VC 泛化界的平方根依赖。\(\blacksquare\)

\subsection{341.3 核方法与RKHS}\label{ux6838ux65b9ux6cd5ux4e0erkhs}

\textbf{定义 341.4}（正定核与再生核Hilbert空间 / RKHS）：设 \(\mathcal{X}\) 为非空集。函数 \(K:\mathcal{X}\times\mathcal{X}\to\mathbb{R}\) 称为\textbf{正定核}，若对任意 \(n\in\mathbb{N}\)、\(\{x_i\}_{i=1}^n\subset\mathcal{X}\) 和 \(\{c_i\}_{i=1}^n\subset\mathbb{R}\)，有 \(\sum_{i,j=1}^n c_i c_j K(x_i,x_j)\geq 0\)。相应的 \textbf{RKHS} \(\mathcal{H}_K\) 是满足再生性的函数Hilbert空间：对任意 \(f\in\mathcal{H}_K\) 和 \(x\in\mathcal{X}\)，\(f(x)=\langle f, K(\cdot,x)\rangle_{\mathcal{H}_K}\)。

\textbf{定理 341.4}（表示定理，Kimeldorf-Wahba 1970；Scholkopf-Herbrich-Smola 2001）：设 \(\mathcal{H}_K\) 为核 \(K\) 对应的RKHS，\(\Omega:[0,\infty)\to\mathbb{R}\) 严格单调增。则正则化经验风险最小化问题

\[\min_{f\in\mathcal{H}_K} \frac{1}{n}\sum_{i=1}^n \ell(f(x_i),y_i) + \lambda\,\Omega(\|f\|_{\mathcal{H}_K})\]

的最优解形如 \(f^*(x)=\sum_{i=1}^n \alpha_i K(x_i,x)\)（\(\alpha_i\in\mathbb{R}\)），即解落在训练数据的有限维子空间中。

\textbf{证明}：将任意 \(f\in\mathcal{H}_K\) 正交分解为 \(f=f_\parallel+f_\perp\)，其中 \(f_\parallel\in\mathrm{span}\{K(\cdot,x_i)\}\)。由再生性 \(f(x_i)=\langle f_\parallel+f_\perp,K(\cdot,x_i)\rangle=\langle f_\parallel,K(\cdot,x_i)\rangle\)，故损失项仅依赖 \(f_\parallel\)。又 \(\|f\|^2=\|f_\parallel\|^2+\|f_\perp\|^2\)，正则项 \(\Omega(\|f\|)\) 在 \(\|f_\perp\|>0\) 时严格增大。因此 \(f^*\) 必满足 \(f_\perp=0\)。\(\blacksquare\)

\subsection{341.4 Mercer定理与SVM}\label{mercerux5b9aux7406ux4e0esvm}

\textbf{定理 341.5}（Mercer定理，Mercer 1909）：设 \(\mathcal{X}\subset\mathbb{R}^d\) 为紧集，\(K:\mathcal{X}\times\mathcal{X}\to\mathbb{R}\) 为连续对称函数。\(K\) 是正定核当且仅当存在 \(L^2(\mathcal{X})\) 的标准正交基 \(\{\phi_j\}_{j=1}^\infty\) 和非负特征值 \(\lambda_j\geq 0\)（\(\sum\lambda_j<\infty\)），使得

\textbf{证明}：SGD更新\(\mathbf{w}_{t+1}=\mathbf{w}_t-\eta_t\nabla f(\mathbf{w}_t;\xi_t)\)（\(\xi_t\)为随机样本）。设\(f\)为\(\mu\)-强凸且\(\nabla f\)为\(L\)-Lipschitz。\(\mathbb{E}[\|\mathbf{w}_{t+1}-\mathbf{w}^*\|^2|\mathbf{w}_t]\le\|\mathbf{w}_t-\mathbf{w}^*\|^2-2\eta_t\mu\|\mathbf{w}_t-\mathbf{w}^*\|^2+\eta_t^2\sigma^2\)。取\(\eta_t=1/(\mu t)\)，递推得\(\mathbb{E}[\|\mathbf{w}_t-\mathbf{w}^*\|^2]=O(1/t)\)。\(\blacksquare\)

\[K(x,y) = \sum_{j=1}^\infty \lambda_j\,\phi_j(x)\phi_j(y)\]

级数对 \((x,y)\) 绝对且一致收敛。此特征分解将核方法嵌入到可能无限维的特征空间 \(\Phi(x)=(\sqrt{\lambda_j}\phi_j(x))_{j=1}^\infty\) 中，\(K(x,y)=\langle\Phi(x),\Phi(y)\rangle\)。

\textbf{定义 341.5}（SVM原问题与对偶问题）：给定训练集 \(\{(x_i,y_i)\}_{i=1}^n\)（\(y_i\in\{\pm 1\}\)），软间隔SVM的\textbf{原问题}为

\[\min_{\mathbf{w},b,\boldsymbol{\xi}} \frac{1}{2}\|\mathbf{w}\|^2 + C\sum_{i=1}^n \xi_i \quad \text{s.t.} \quad y_i(\langle\mathbf{w},\Phi(x_i)\rangle+b)\geq 1-\xi_i,\;\xi_i\geq 0\]

引入Lagrange乘子 \(\alpha_i,\mu_i\geq 0\)，消去 \(\mathbf{w},b,\xi_i\) 得\textbf{对偶问题}：

\[\max_{\boldsymbol{\alpha}} \sum_{i=1}^n \alpha_i - \frac{1}{2}\sum_{i,j=1}^n \alpha_i\alpha_j y_i y_j K(x_i,x_j) \quad \text{s.t.} \quad \sum_{i=1}^n \alpha_i y_i = 0,\; 0\leq\alpha_i\leq C\]

对偶问题的解 \(\boldsymbol{\alpha}^*\) 中 \(\alpha_i^*>0\) 的样本即为支持向量，决策函数 \(h(x)=\mathrm{sign}(\sum_{i}\alpha_i^* y_i K(x_i,x)+b^*)\)。

\subsection{341.5 神经网络的通用逼近定理}\label{ux795eux7ecfux7f51ux7edcux7684ux901aux7528ux903cux8fd1ux5b9aux7406}

\textbf{定理 341.6}（通用逼近定理，Cybenko 1989；Hornik-Stinchcombe-White 1989）：设 \(\sigma:\mathbb{R}\to\mathbb{R}\) 为非常数、有界、连续且单调递增的Sigmoid激活函数。则形如

\[G(x) = \sum_{j=1}^N \alpha_j\,\sigma(\mathbf{w}_j^T x + b_j)\]

的单隐藏层前馈神经网络在 \(C([0,1]^n)\)（紧集上连续函数空间）中关于一致范数 \(\|\cdot\|_\infty\) 稠密。

*\textbf{证明}：考虑假设类\(\mathcal{H}\)的VC维\(d\)。对任意分布\(\mathcal{D}\)和\(\varepsilon>0\)，以概率\(\ge 1-\delta\)，对所有\(h\in\mathcal{H}\)，\(|\operatorname{err}(h)-\widehat{\operatorname{err}}(h)|\le O\left(\sqrt{\frac{d+\log(1/\delta)}{m}}\right)\)（Vapnik-Chervonenkis界）。\(\mathcal{H}\)为PAC可学习的当且仅当\(d<\infty\)。充分性：取\(m\ge O((d+\log(1/\delta))/\varepsilon^2)\)样本，ERM算法输出\(\hat{h}\)满足\(\operatorname{err}(\hat{h})\le\varepsilon\)。必要性：若\(d=\infty\)，存在分布使任意\(m\)样本的ERM失败。\(\blacksquare\)

lacksquare\$

\hrulefill

\subsection{342.A 计算几何补充}\label{a-ux8ba1ux7b97ux51e0ux4f55ux8865ux5145}

非标准分析（Nonstandard Analysis）由 Abraham Robinson 于 1961 年创立，使用模型论中的超幂构造（ultrapower construction）将实数域 \(\mathbb{R}\) 扩充为包含''无限小''和''无限大''元素的非标准模型 \({}^*\mathbb{R}\)，使 Leibniz 的无限小演算获得了严格的数学基础。

\subsection{超幂构造与紧致性定理}\label{ux8d85ux5e42ux6784ux9020ux4e0eux7d27ux81f4ux6027ux5b9aux7406}

\textbf{定义}（超滤子 / ultrafilter）：集合 \(I\) 上的\textbf{超滤子} \(\mathcal{U}\) 是 \(I\) 的幂集 \(\mathcal{P}(I)\) 的子集，满足：(1) \(\emptyset \notin \mathcal{U}\)；(2) 若 \(A, B \in \mathcal{U}\)，则 \(A \cap B \in \mathcal{U}\)；(3) 若 \(A \in \mathcal{U}\) 且 \(A \subset B\)，则 \(B \in \mathcal{U}\)；(4) 对任意 \(A \subset I\)，\(A \in \mathcal{U}\) 或 \(I \setminus A \in \mathcal{U}\)（极大性）。

不包含任何有限集（即仅包含无限集）的非主超滤子（nonprincipal ultrafilter）的存在性信赖于选择公理（经 Zorn 引理）。

\textbf{定义}（超幂 / ultrapower）：设 \(\mathcal{U}\) 是 \(\mathbb{N}\) 上的非主超滤子。\(\mathbb{R}\) 的超幂是积 \(\mathbb{R}^{\mathbb{N}}\) 模去如下等价关系的商： \[(a_n) \sim (b_n) \iff \{n : a_n = b_n\} \in \mathcal{U}\] 所得的域记为 \({}^*\mathbb{R} = \mathbb{R}^{\mathbb{N}} / \mathcal{U}\)，称为\textbf{非标准实数域}（超实数域）。

\textbf{定理}（Łoś定理 / 超积基本定理，1955）：对任意语言 \(\mathcal{L}\) 中的一阶公式 \(\varphi(x_1, \ldots, x_k)\) 和任意元素 \([a_n] \in {}^*\mathbb{R}\)（\(a_n\) 是第 \(n\) 个分量）， \[{}^*\mathbb{R} \models \varphi([a_n^{(1)}], \ldots, [a_n^{(k)}]) \iff \{n : \mathbb{R} \models \varphi(a_n^{(1)}, \ldots, a_n^{(k)})\} \in \mathcal{U}\] 特别地，标准实数 \(\mathbb{R}\) 上所有真的一阶语句在 \({}^*\mathbb{R}\) 中也真（传递原理 transfer principle）。

\subsection{超实数与非标准模型}\label{ux8d85ux5b9eux6570ux4e0eux975eux6807ux51c6ux6a21ux578b}

\textbf{定义}（超实数的代数结构）：\({}^*\mathbb{R}\) 是包含 \(\mathbb{R}\) 作为子域的有序域。嵌入通过将 \(r \in \mathbb{R}\) 映为常数列 \((r, r, r, \ldots)\) 的等价类实现。

\textbf{定义}（无限小 / 无限大 / 有限超实数）： - \(x \in {}^*\mathbb{R}\) 是\textbf{无限小}（infinitesimal），若对所有 \(n \in \mathbb{N}\)，\(|x| < 1/n\)。记作 \(x \approx 0\) - \(x \in {}^*\mathbb{R}\) 是\textbf{无限大}（infinite / unlimited），若对所有 \(n \in \mathbb{N}\)，\(|x| > n\) - \(x \in {}^*\mathbb{R}\) 是\textbf{有限}（finite），若存在标准 \(n \in \mathbb{N}\) 使 \(|x| < n\)。所有有限超实数构成 \({}^*\mathbb{R}\) 的子环 \(\mathcal{O}\)

\textbf{定理}（标准部分定理 / standard part）：每个有限超实数 \(x \in {}^*\mathbb{R}\) 有唯一的标准部分 \(\operatorname{st}(x) \in \mathbb{R}\)（即与 \(x\) 相差一个无限小的唯一标准实数）：\(x - \operatorname{st}(x) \approx 0\)。从有限超实数到 \(\mathbb{R}\) 的映射 \(\operatorname{st}: \mathcal{O} \to \mathbb{R}\) 是满环同态，核为无限超实数理想 \(\mathfrak{m} \subset \mathcal{O}\)。

\subsection{非标准分析与非标准模型}\label{ux975eux6807ux51c6ux5206ux6790ux4e0eux975eux6807ux51c6ux6a21ux578b}

\textbf{定义}（超幂扩张 / enlargement）：设 \(\mathcal{U}\) 是足够大索引集 \(I\) 上的非主超滤子（包含所有 \(\mathbb{N}\) 中的有限余子集）。则\textbf{超结构}的扩张 \(\bigcup_{n} V_n(X) \hookrightarrow \bigcup_{n} {}^*V_n(X)\)（其中 \(V_0(X) = X\)，\(V_{n+1}(X) = V_n(X) \cup \mathcal{P}(V_n(X))\)）给出一个满足传递原理和\textbf{饱和原理}（\(\kappa\)-saturated for sufficiently large \(\kappa\)）的非标准模型。

\textbf{定义}（饱和性 / saturation）：非标准扩张是 \textbf{\(\kappa\)-饱和的}，如果对索引集基数 \(< \kappa\) 的任意一族内部集（internal sets）\(\{A_i\}\)，若其具有有限交性质（任意有限个子族的交非空），则全族有非空交。饱和性保证了非标准扩张''足够大''------包含充分多的理想元素。

\textbf{定理}（非标准分析的微积分 / Leibniz 无穷小演算的严格化）： - \(f\) 在 \(x_0\) 连续 \(\iff\) 对所有 \(x \approx x_0\)（\(x\) 是 \(x_0\) 的无限小邻域内的超实数），\(f(x) \approx f(x_0)\) - \(f\) 在 \(x_0\) 可微且导数为 \(L\) \(\iff\) 对所有 \(dx \approx 0, dx \neq 0\)，\(\frac{f(x_0 + dx) - f(x_0)}{dx} \approx L\)（且此关系与 \(dx\) 的选取无关） - Riemann积分 \(\int_a^b f(x) dx\) 由无限细分的 Riemann 和的公共标准部分给出（Keisler 1960s）

\subsection{非标准分析的应用}\label{ux975eux6807ux51c6ux5206ux6790ux7684ux5e94ux7528}

\textbf{定义}（非标准Hilbert空间与Loeb测度）：非标准分析在泛函分析中给出紧算子和Fredholm理论的简化处理。\textbf{Loeb测度}（1975）将标准Borel测度构造为非标准计数测度的标准部分，为构造Brown运动等提供了新路径。

\textbf{定理}（Bernstein-Robinson 多项式紧算子不变子空间定理，1966）：用非标准分析给出了 Hilbert 空间上多项式紧算子存在非平凡不变子空间的第一个证明------经标准部分翻译后得到经典分析的结果。这证明了非标准方法在泛函分析中的威力。

\textbf{推论}（非标准测度论）：\(\sigma\)-加性测度空间的非标准构造（Loeb空间）为随机分析、鞅论和Brown运动理论提供了统一的框架：标准随机过程可视为内部过程的''投影''。

\hrulefill

\emph{本卷完成了数理逻辑的核心内容：一阶逻辑的完备性（Gödel 1929）建立了语法与语义的桥梁；紧致性定理和 Löwenheim-Skolem 定理揭示了形式语言的局限性；Turing 机和停机问题为可计算性划定了边界；Gödel 不完全性定理（1931）是 20 世纪数学最深刻的成果之一，它证明了任何足够强的形式系统都不可避免地存在不可判定的命题；序数和基数理论为无穷提供了精确的度量；连续统假设的独立性（Gödel 1938 / Cohen 1963）展示了形式系统的根本局限性。数理逻辑不仅是数学的基础，也是计算机科学（计算理论、形式化验证）的基石。}

\emph{本卷在计算理论（自动机、形式语言、可计算性、计算复杂性、量子计算）的基础上，整合了数理逻辑（命题逻辑、一阶逻辑、模型论、证明论、Gödel 不完备定理）和计算几何与计算机代数。从哥德尔到图灵，形成了数学基础与计算理论的完整链条。共 12 章。}

\emph{卷三十一：计算理论与逻辑终。}

\newpage

\chapter{11-数值与计算/02-计算理论/04-计算理论.md}\label{ux6570ux503cux4e0eux8ba1ux7b9702-ux8ba1ux7b97ux7406ux8bba04-ux8ba1ux7b97ux7406ux8bba.md}

\section{第414章：自动机与形式语言}\label{ux7b2c414ux7ae0ux81eaux52a8ux673aux4e0eux5f62ux5f0fux8bedux8a00}

自动机理论是计算模型的最简形式，与形式语言（Chomsky 谱系）之间有精确的对应关系。

\subsection{152.1 有限自动机}\label{ux6709ux9650ux81eaux52a8ux673a}

\textbf{定义 152.1}（确定型有限自动机 / DFA）：一个\textbf{确定型有限自动机}是五元组 \(M = (Q, \Sigma, \delta, q_0, F)\)，其中 \(Q\) 为有限状态集，\(\Sigma\) 为输入字母表，\(\delta: Q \times \Sigma \to Q\) 为转移函数，\(q_0 \in Q\) 为初始状态，\(F \subseteq Q\) 为接受状态集。

\(M\) 接受字符串 \(w = w_1 w_2 \cdots w_n\) 如果存在状态序列 \(r_0, r_1, \ldots, r_n\) 满足 \(r_0 = q_0\)，\(r_{i+1} = \delta(r_i, w_{i+1})\)，且 \(r_n \in F\)。

\textbf{定义 152.2}（正则语言）：语言 \(L \subseteq \Sigma^*\) 是\textbf{正则语言}，如果存在 DFA 使得 \(L = L(M)\)（\(M\) 接受的语言）。

\textbf{定义 152.3}（非确定型有限自动机 / NFA）：转移函数 \(\delta: Q \times \Sigma_\varepsilon \to \mathcal{P}(Q)\)（\(\Sigma_\varepsilon = \Sigma \cup \{\varepsilon\}\)）。NFA 接受字符串如果存在某条接受路径。

\textbf{定理 152.1}（DFA 与 NFA 的等价性 / Rabin-Scott，1959）：任何 NFA 可转化为等价的 DFA（子集构造法）。正则语言类在并、交、补、连接、Kleene 星号下封闭。

\textbf{证明}：设 NFA \(N=(Q,\Sigma,\delta,q_0,F)\)。构造 DFA \(M=(\mathcal{P}(Q),\Sigma,\delta',\{q_0\},F')\)，\(\delta'(S,a)=\bigcup_{q\in S}\delta(q,a)\)（含 \(\varepsilon\)-闭包），\(F'=\{S\subseteq Q:S\cap F\neq\emptyset\}\)。归纳可证 \(\hat{\delta}'(\{q_0\},w)\) 即 \(N\) 读 \(w\) 后可能状态集，故 \(L(M)=L(N)\)。封闭性：补由交换 DFA 的接受/非接受状态；交由乘积自动机 \(Q_1\times Q_2\)；并由 \(\overline{\overline{L_1}\cap\overline{L_2}}\) 或 NFA 构造；连接由 \(\varepsilon\)-转移串联 NFA；Kleene 星号由 \(\varepsilon\)-反馈。\(\blacksquare\)

\textbf{定理 152.2}（正则语言的泵引理）：若 \(L\) 是正则语言，则存在泵长度 \(p\) 使得任何 \(w \in L\)（\(|w| \geq p\)）可分解为 \(w = xyz\)，满足 \(|xy| \leq p\)，\(|y| > 0\)，且 \(xy^iz \in L\) 对所有 \(i \geq 0\)。

\textbf{证明}：设 DFA \(M\) 有 \(p\) 个状态接受 \(L\)。取 \(w\in L\) 且 \(|w|\ge p\)。\(M\) 读 \(w\) 经过 \(|w|+1\) 个状态 \(r_0,\dots,r_{|w|}\)。由鸽巢原理，前 \(p+1\) 个状态中必有重复：\(r_i=r_j\)（\(0\le i<j\le p\)）。令 \(x=w_1\cdots w_i\)，\(y=w_{i+1}\cdots w_j\)，\(z=w_{j+1}\cdots w_{|w|}\)。\(|xy|\le p\)，\(|y|>0\)。读 \(y\) 使 \(M\) 从 \(r_i\) 回到 \(r_i\)，故读 \(y^i\) 亦循环，\(xy^iz\) 均被接受。\(\blacksquare\)

\textbf{定理 152.3}（Myhill-Nerode 定理）：语言 \(L\) 是正则的当且仅当 \(L\) 的右不变等价关系 \(\equiv_L\) 的指数（等价类个数）有限。这给出了正则语言的最小 DFA 构造方法。

\textbf{证明}：定义 \(x\equiv_L y\iff(\forall z\in\Sigma^*)(xz\in L\leftrightarrow yz\in L)\)。若 \(\equiv_L\) 指数有限，构造 DFA：状态为等价类 \([x]\)，转移 \(\delta([x],a)=[xa]\)，接受 \(F=\{[x]:x\in L\}\)。态数=指数且为最小 DFA。反之若 \(L\) 被 \(p\) 态 DFA \(M\) 接受，定义 \(x\sim y\) 若 \(\hat{\delta}(q_0,x)=\hat{\delta}(q_0,y)\)。\(\sim\) 细化 \(\equiv_L\)（因状态同则未来行为同），指数 \(\le p\) 有限。\(\blacksquare\)

\subsection{152.2 下推自动机与上下文无关语言}\label{ux4e0bux63a8ux81eaux52a8ux673aux4e0eux4e0aux4e0bux6587ux65e0ux5173ux8bedux8a00}

\textbf{定义 152.4}（上下文无关文法 / CFG）：文法 \(G = (V, \Sigma, R, S)\)，其中产生式规则形如 \(A \to \alpha\)（\(A \in V\)，\(\alpha \in (V \cup \Sigma)^*\)）。\(L(G)\) 是能从 \(S\) 推导出的终结符串集合。

\textbf{定义 152.5}（下推自动机 / PDA）：PDA 是有限自动机加上一个栈（stack）。转移函数 \(\delta: Q \times \Sigma_\varepsilon \times \Gamma_\varepsilon \to \mathcal{P}(Q \times \Gamma_\varepsilon)\)，其中 \(\Gamma\) 为栈字母表。

\textbf{定理 152.4}（PDA 与 CFG 的等价性）：语言 \(L\) 是上下文无关的当且仅当存在 PDA 接受 \(L\)。

\textbf{证明}：（CFG\(\Rightarrow\)PDA）对 CFG \(G\)，构造单状态 PDA：栈初始为 \(S\)。对每个产生式 \(A\to\alpha\)，可弹出栈顶 \(A\) 并压入 \(\alpha^R\)（逆序压入）。对终结符 \(a\)，若栈顶为 \(a\) 则弹出。PDA 以空栈接受。（PDA\(\Rightarrow\)CFG）对 PDA \(P\)，构造变元 \(A_{pq}\)（从状态 \(p\) 到 \(q\) 且空栈）。产生式：\(A_{pp}\to\varepsilon\)；\(A_{pq}\to A_{pr}A_{rq}\)（路径分割）；\(A_{pq}\to a A_{rs} b\)（读 \(a\) 压 \(X\)，最终弹出读 \(b\)）。归纳可证 \(A_{pq}\) 生成恰好从 \(p\) 到 \(q\) 空栈可接受的字符串。\(\blacksquare\)

\textbf{定理 152.5}（CFL 的泵引理）：若 \(L\) 是 CFL，则存在 \(p\) 使得任何 \(w \in L\)（\(|w| \geq p\)）可分解为 \(w = uvxyz\)，满足 \(|vxy| \leq p\)，\(|vy| > 0\)，且 \(uv^ixy^iz \in L\) 对所有 \(i \geq 0\)。

\textbf{证明}：设 CFG 的 Chomsky 范式中变元数为 \(v\)。取 \(w\in L\) 且 \(|w|\ge 2^{v+1}\)。\(w\) 的语法树深度 \(>v\)，故某路径上必有变元重复 \(A\)。则 \(S\Rightarrow^* uAz\Rightarrow^* uvAyz\Rightarrow^* uvxyz=w\)。由重复结构，\(|vxy|\le 2^{v+1}\)（子树不高于 \(v+1\) 的产物），\(|vy|>0\)（否则 \(A\Rightarrow^*A\) 可消去）。因 \(A\Rightarrow^* vAy\) 且 \(A\Rightarrow^* x\)，可重复 \(vAy\) 部分：\(uv^ixy^iz\) 均可生成。取 \(p=2^{v+1}\)。\(\blacksquare\)

\textbf{定理 152.6}（Chomsky 范式）：任何不含 \(\varepsilon\) 的 CFG 可转化为等价的 Chomsky 范式格式（产生式仅形如 \(A \to BC\) 或 \(A \to a\)）。

\textbf{证明}：逐步转换。（1）添加新始符 \(S_0\to S\)。（2）消除 \(\varepsilon\)-产生式 \(A\to\varepsilon\)：对每个含 \(A\) 的产生式 \(B\to\alpha A\beta\)，添加 \(B\to\alpha\beta\)，最后删除 \(A\to\varepsilon\)。（3）消除单产生式 \(A\to B\)：若 \(B\to\gamma\) 非单产生式，添加 \(A\to\gamma\)。（4）拆分长产生式 \(A\to X_1X_2\cdots X_k\)（\(k>2\)）：引入新变元 \(C_1,\dots,C_{k-2}\)，替换为 \(A\to X_1 C_1\)，\(C_1\to X_2 C_2\)，\ldots，\(C_{k-2}\to X_{k-1}X_k\)。（5）处理混合产生式中的终结符：对 \(A\to aB\)（\(a\) 终结符），引入 \(T_a\to a\) 替换。最终仅含 \(A\to BC\) 和 \(A\to a\)，文法等价。\(\blacksquare\)

\subsection{152.3 Turing 机与 Chomsky 谱系}\label{turing-ux673aux4e0e-chomsky-ux8c31ux7cfb}

\textbf{定义 152.6}（Turing 机 / TM）：Turing 机是七元组 \(M = (Q, \Sigma, \Gamma, \delta, q_0, q_{\text{accept}}, q_{\text{reject}})\)，其中 \(\Gamma\) 是带字母表，\(\delta: Q \times \Gamma \to Q \times \Gamma \times \{L, R\}\)（转移函数）。

\textbf{定理 152.7}（Chomsky 谱系）： - 类型 3（正则文法）\(\leftrightarrow\) DFA/NFA \(\leftrightarrow\) 正则语言 - 类型 2（上下文无关文法）\(\leftrightarrow\) PDA \(\leftrightarrow\) CFL - 类型 1（上下文有关文法）\(\leftrightarrow\) 线性有界自动机（LBA） - 类型 0（无限制文法）\(\leftrightarrow\) Turing 机 \(\leftrightarrow\) 递归可枚举语言（r.e.）

\textbf{证明}：各对应关系的严格性由泵引理和不可判定性保证。正则 \(\subsetneq\) CFL：\(\{a^n b^n\}\) 是 CFL（\(S\to aSb|\varepsilon\)）但泵引理证其非正则。CFL \(\subsetneq\) CSL：\(\{a^n b^n c^n\}\) 由 CSG 生成（非收缩规则）但 CFL 泵引理证其非 CFL。CSL \(\subsetneq\) r.e.：停机问题是 r.e. 但非 CSL（CSL 对补封闭且停机问题不可判定，而 r.e. 对补不封闭）。所有包含关系恰为 Chomsky 层级。\(\blacksquare\)

\textbf{定理 152.8}（Turing 机的变体等价性）：多带 TM、非确定型 TM（NTM）、双向无限带 TM 在计算能力上与单带 TM 等价。

\textbf{证明}：（多带到单带）\(k\) 带 TM \(M\) 的模拟：单带 TM \(S\) 将 \(k\) 条带交错存储，带头位置用特殊符号标记。\(S\) 扫描整带收集 \(k\) 个带头下符号，按 \(M\) 转移函数更新各带。（NTM\(\to\)DTM）使用 dovetailing 宽度优先遍历 NTM 计算树，语言接受能力等价（时间可能指数膨胀）。（双向无限带）通过折叠（fold）映射将双向无限带编码为单向无限带：\(\dots b_{-2}b_{-1}b_0 b_1 b_2 \dots\) 映射到带的两轨存储。所有变体接受相同语言类（r.e.）。\(\blacksquare\)

\hrulefill

\hrulefill

\hrulefill

\hrulefill

\hrulefill

\hrulefill

\section{第415章：可计算性理论深化}\label{ux7b2c415ux7ae0ux53efux8ba1ux7b97ux6027ux7406ux8bbaux6df1ux5316}

本章在 V20 Ch 325 的基础上，深入探讨可计算性的核心定理和不可判定问题。

\subsection{153.1 可计算函数的精确定义}\label{ux53efux8ba1ux7b97ux51fdux6570ux7684ux7cbeux786eux5b9aux4e49}

\textbf{定义 153.1}（部分递归函数）：\textbf{部分递归函数}是从初始函数（零函数 \(Z(x)=0\)、后继函数 \(S(x)=x+1\)、投影函数 \(P_i^n(x_1,\ldots,x_n)=x_i\)）通过复合、原始递归和极小化（\(\mu\)-算子）得到的函数类。\textbf{全递归函数}是处处有定义的部分递归函数。

\textbf{定理 153.1}（Church-Turing 论题）：以下计算模型定义相同的函数类： - Turing 机可计算函数 - \(\lambda\)-可定义函数（Church 的 \(\lambda\) 演算，1936） - 部分递归函数（Gödel-Kleene，1936） - 随机存取机（RAM）可计算函数 - 所有「合理」的计算模型

\textbf{证明}：等价性证明分多方向。（TM \(\to\) 部分递归）用 Gödel 编码将 TM 格局编码为自然数，停机格局的 \(\mu\) 极小化搜索得到部分递归函数。（部分递归 \(\to\) TM）原始递归函数可由带计数器的 TM 实现；\(\mu\) 算子对应穷举搜索。（\(\lambda\)-演算 \(\to\) TM）用 SECD 机或 Krivine 机器的 TM 模拟实现 \(\beta\)-归约。（TM \(\to\) \(\lambda\)-演算）用 Church 编码表示自然数和带状态，归约模拟转移。论题本身不可形式证明（涉及''直观可计算''概念），但所有已知''有效可计算''过程均可归约为 Turing 机。\(\blacksquare\)

\textbf{定理 153.2}（Kleene 正规形定理）：存在原始递归谓词 \(T(e, x, y)\)（Kleene T-谓词）和原始递归函数 \(U(y)\) 使得任何部分递归函数可表为 \(f(x) = U(\mu y \, T(e, x, y))\)。

\textbf{证明}：\(T(e,x,y)\) 的直觉：\(e\) 编码 TM 的索引，\(y\) 编码完整的停机计算历史。判定 \(T(e,x,y)\) 仅需局部验证：初始格局正确、每步转移合法、最终格局为停机状态（且输出位置包含有效值）。这些全部是原始递归的（有限状态转移表查表）。\(U(y)\) 从计算历史的编码中提取输出值（最后一个格局中约定的输出位置）。故 \(f(x)\simeq U(\mu y\,T(e,x,y))\) 将部分递归函数表为单次 \(\mu\) 算子作用于原始递归谓词，达到''一个 \(\mu\) 够用''的极小性。\(\blacksquare\)

\textbf{定理 153.3}（\(s\)-\(m\)-\(n\) 定理 / 参数化定理）：存在原始递归函数 \(S_n^m\) 使得 \(\varphi_{S_n^m(e, x_1, \ldots, x_n)}^{(m)}(y_1, \ldots, y_m) = \varphi_e^{(m+n)}(x_1, \ldots, x_n, y_1, \ldots, y_m)\)。

\textbf{证明}：\(S_n^m\) 输入程序索引 \(e\) 和固定参数 \(\vec{x}\)，输出新程序索引 \(e'\)。新程序 \(e'\) 在输入 \(\vec{y}\) 时的行为：将硬编码的 \(\vec{x}\) 和输入 \(\vec{y}\) 拼合为 \((\vec{x},\vec{y})\)，然后模拟第 \(e\) 个程序。这类似于部分求值（partial evaluation）。在 Turing 机模型中，\(S_n^m\) 通过修改 \(e\) 的代码实现：将读输入的指令替换为直接产生 \(\vec{x}\) 和 \(\vec{y}\) 拼合的指令。此修改是原始递归的（有限符号替换）。Kleene 给出了基于 Gödel 编码的显式构造。\(\blacksquare\)

\subsection{153.2 不可判定性}\label{ux4e0dux53efux5224ux5b9aux6027}

\textbf{定理 153.4}（停机问题 / Halting Problem，Turing 1936）：语言 \(HALT = \{\langle M, w \rangle : M\) 在输入 \(w\) 上停机\(\}\) 是递归可枚举的，但不是可判定的（递归的）。

\emph{证明}：假设存在判定器 \(H\)。构造 \(D\)：在输入 \(\langle M \rangle\) 时，\(D\) 以 \(\langle M, \langle M \rangle \rangle\) 运行 \(H\)，若 \(H\) 接受则拒绝，反之接受。则 \(D(\langle D \rangle)\) 接受当且仅当 \(D(\langle D \rangle)\) 不接受，矛盾。∎

\textbf{定理 153.5}（Rice 定理，1953）：递归可枚举语言的任何非平凡性质（即存在具有和不具有该性质的 r.e. 语言）都是不可判定的。即 \(\{\langle M \rangle : L(M) \in \mathcal{P}\}\) 不可判定，只要 \(\mathcal{P}\) 非平凡。

\textbf{证明}：构造接受\(L\)的Turing机\(M\)。\(L\)为正则语言，存在DFA \(D\)识别\(L\)。\(M\)模拟\(D\)的状态转移：读入符号后更新状态，输入结束时检查是否处于接受状态。Turing机可模拟任意DFA（DFA是Turing机的特例，读写头仅向右移动且不写入）。因此\(L\)可被Turing机判定，\(L\in\mathbf{R}\)。\(\blacksquare\)

\textbf{推论 153.6}（不可判定问题的例子）：以下问题均不可判定： - 判定两个 CFG 是否生成相同的语言 - 判定一个 CFG 是否是歧义的 - 判定一个 TM 是否接受空语言 - Hilbert 第十问题：判定 Diophantine 方程是否有整数解（Matiyasevich，1970，基于 Davis-Putnam-Robinson 的工作）

\subsection{153.3 相对计算与 Turing 度}\label{ux76f8ux5bf9ux8ba1ux7b97ux4e0e-turing-ux5ea6}

\textbf{定义 153.2}（神谕 Turing 机 / Oracle TM）：带有神谕 \(A \subseteq \Sigma^*\) 的 Turing 机 \(M^A\) 可以查询「\(w \in A\)？」。相对可计算性：\(B \leq_T A\) 如果 \(B\) 在 \(M^A\) 下可判定。

\textbf{定义 153.3}（Turing 度）：\(\leq_T\) 的等价类构成 Turing 度。Turing 度构成上半格。\(\mathbf{0}\) 是可计算集的度，\(\mathbf{0}'\) 是停机问题的度。

\textbf{定理 153.7}（Post 问题，1944 / Friedberg-Muchnik 1956-1957）：存在递归可枚举的 Turing 度严格介于 \(\mathbf{0}\) 和 \(\mathbf{0}'\) 之间。这通过优先方法（priority method）证明，是递归论的核心技术。

\textbf{证明}：构造 r.e. 集 \(A\) 满足 \(\mathbf{0}<_T A<_T\mathbf{0}'\)。需求：\(R_{2e}\)：\(\Phi_e\neq A\)（\(A\) 非递归）；\(R_{2e+1}\)：\(\Phi_e^A\) 非全或 \(\Phi_e^A\neq K\)（\(A\) 不计算停机问题 \(K\)）。有限损伤优先方法：需求按 \(R_0,R_1,R_2,\dots\) 优先级排序。对 \(R_{2e}\)：选靶 \(x\) 暂不放入 \(A\)，若某步发现 \(\Phi_e(x)\downarrow=0\)，则将 \(x\) 放入 \(A\) 使 \(\Phi_e(x)\neq A(x)\)。放入 \(x\) 可能损伤低优先需求（它们依赖''\(A\) 不变''的假设），但每个需求只被更高优先需求有限次损伤。对 \(R_{2e+1}\)：监控 \(\Phi_e^A\) 是否在某点与 \(K\) 不同，否则保持 \(A\) 不动。归纳可证每个需求最终永久满足（不再被损伤），故 \(0<_T A\) 且 \(K\not\le_T A\)（即 \(A<_T\mathbf{0}'\)）。\(\blacksquare\)

\textbf{定理 153.8}（Kleene-Post 定理）：存在不可比较的 Turing 度（即 \(\mathbf{a} \nleq_T \mathbf{b}\) 且 \(\mathbf{b} \nleq_T \mathbf{a}\)）。

\textbf{证明}：构造 \(\Delta_2^0\) 集 \(A,B\subseteq\mathbb{N}\) 满足对每个 \(e\)：（1）\(\Phi_e^A\neq B\)；（2）\(\Phi_e^B\neq A\)。构造分步进行：设 \(A_s,B_s\) 为前 \(s\) 步的有限近似（以 \(\{0,1\}\) 值序列表示）。第 \(e\) 步处理需求 \(\Phi_e^A\neq B\)：找扩展 \(A\supseteq A_s\) 和 \(x\)，使在有限信息下 \(\Phi_e^A(x)\) 已收敛，定义 \(B(x)\neq\Phi_e^A(x)\)。因 \(A\) 的有限初始段不能强制神谕 \(\Phi_e^A\) 对所有输入收敛，总存在适当扩展实现分离。需求 \(\Phi_e^B\neq A\) 对称处理。对偶数需求优先处理 \(A\)，奇数优先处理 \(B\)，无冲突。最终 \(A,B\) 均为 \(\Delta_2^0\)（极限可计算），且相互不可 Turing 归约。\(\blacksquare\)

\hrulefill

\section{第416章：递归论深化}\label{ux7b2c416ux7ae0ux9012ux5f52ux8bbaux6df1ux5316}

递归论（可计算性理论）的深层结构远远超出了停机问题的基本不可判定性。本章将深入探讨递归论的三个核心主题：Kleene递归定理------揭示了计算系统自我引用的能力，是不动点定理在可计算性中的体现；算术阶层与Post定理------在可定义性与图灵度之间建立了精确的对应关系，是描述集合论在自然数上的特化版本；以及Friedberg-Muchnik的有限损害优先方法------这是递归论中最精巧的构造技术之一，给出了Post问题的正面回答，证明了存在严格介于可计算度和停机问题度之间的递归可枚举度。

\subsection{Kleene递归定理}\label{kleeneux9012ux5f52ux5b9aux7406}

\textbf{定理336.1（Kleene第一递归定理------不动点定理）} 设 \(f: \mathbb{N} \to \mathbb{N}\) 是全递归函数。则存在 \(e \in \mathbb{N}\)（称为 \(f\) 的不动点）使得 \(\varphi_e = \varphi_{f(e)}\)，即第 \(e\) 个部分递归函数与第 \(f(e)\) 个部分递归函数是同一个函数。换言之，任何可计算地变换程序索引的操作都有一个''不动点程序''。

\textbf{证明} 考虑函数 \(g(x) = \varphi_x(x)\)（即第 \(x\) 个程序在输入 \(x\) 上的行为）。\(g\) 是部分递归的。定义部分递归函数

\[h(x) = \begin{cases} \varphi_{\varphi_x(x)}(y), & \text{若 } \varphi_x(x) \downarrow \\ \uparrow, & \text{否则} \end{cases}\]

即 \(h(x)\) 是索引为 \(\varphi_x(x)\) 的程序（如果定义的话）。\(h\) 是部分递归的，设其索引为 \(e_0\)，即 \(h = \varphi_{e_0}\)。则 \(\varphi_{e_0}(x) = \varphi_{\varphi_x(x)}\)。

令 \(d(y) = f(\varphi_{e_0}(y))\)（若 \(\varphi_{e_0}(y)\) 定义；否则发散）。\(d\) 是部分递归的。现在应用 \(h\) 的定义于某个特殊输入：取 \(x\) 使 \(\varphi_x = d\)。则

\[\varphi_{\varphi_x(x)}(y) = h(x)(y) = \varphi_{e_0}(x)(y) = \varphi_{\varphi_x(x)}(y)\]

所以 \(\varphi_{f(\varphi_{e_0}(x))} = \varphi_{\varphi_{e_0}(x)}\)。令 \(e = \varphi_{e_0}(x)\)，则 \(f(e)\) 和 \(e\) 索引同一函数，即 \(\varphi_{f(e)} = \varphi_e\)。\(\blacksquare\)

\textbf{定理336.2（Kleene第二递归定理------自引用定理）} 对任意部分递归函数 \(f(x, y)\)，存在 \(e\) 使得 \(\varphi_e(y) = f(e, y)\) 对所有 \(y\) 成立。即存在一个程序，其行为等价于''将该程序自身的索引作为第一个参数的 \(f\)``。

\textbf{证明} 由 \(s\)-\(m\)-\(n\) 定理（定理153.3），存在全递归函数 \(s\) 使得 \(\varphi_{s(z)}(y) = f(z, y)\)。对 \(s\) 应用第一递归定理（定理336.1），存在 \(e\) 使得 \(\varphi_e = \varphi_{s(e)}\)。则 \(\varphi_e(y) = \varphi_{s(e)}(y) = f(e, y)\)，证毕。\(\blacksquare\)

\textbf{推论336.3（递归定理的应用）} 存在程序输出自身的源代码（Quine程序）；存在程序其定义域恰为 \(\{e\}\)（单元素集）；对任意 \(n\)，存在索引 \(e\) 使得 \(\varphi_e\) 恰在输入 \(e, e+1, \ldots, e+n-1\) 处定义。以上均为第二递归定理的直接应用。

\subsection{算术阶层与Post定理}\label{ux7b97ux672fux9636ux5c42ux4e0epostux5b9aux7406}

\textbf{定义336.4（算术阶层）} 自然数上的关系按定义复杂性分为：

\begin{itemize}
\tightlist
\item
  \(\Sigma_0^0 = \Pi_0^0\)：递归关系（原始递归可判定）。
\item
  \(\Sigma_{n+1}^0\)：形如 \(\exists y_1 \forall y_2 \cdots Q y_{n+1} R(x, y_1, \ldots, y_{n+1})\) 的关系，其中 \(R \in \Sigma_0^0\)，量词以 \(\exists\) 开头。
\item
  \(\Pi_{n+1}^0\)：形如 \(\forall y_1 \exists y_2 \cdots Q y_{n+1} R(x, y_1, \ldots, y_{n+1})\) 的关系，量词以 \(\forall\) 开头。
\item
  \(\Delta_n^0 = \Sigma_n^0 \cap \Pi_n^0\)：既可通过 \(\Sigma_n^0\) 又可通过 \(\Pi_n^0\) 公式定义的关系。
\end{itemize}

\textbf{定理336.5（算术层级的严格性）} \(\Sigma_n^0 \subsetneq \Sigma_{n+1}^0\) 且 \(\Pi_n^0 \subsetneq \Pi_{n+1}^0\) 对所有 \(n \geq 1\)。特别地，停机问题 \(K = \{e : \varphi_e(e) \downarrow\}\) 是 \(\Sigma_1^0\)-完备的（即每个 \(\Sigma_1^0\) 集可 \(1\)-归约于 \(K\)）。

\textbf{证明} 利用对角线法。定义 \(\Sigma_n^0\)-完备集 \(K_n\) 为 \(\{e : \varphi_e^{\emptyset^{(n-1)}}(e) \downarrow\}\)（其中 \(\emptyset^{(n-1)}\) 是第 \(n-1\) 个图灵跳跃）。\(K_n\) 是 \(\Sigma_n^0\) 但非 \(\Pi_n^0\)：若 \(K_n \in \Pi_n^0\)，则 \(K_n \in \Delta_n^0\)，由此可推出 \(\emptyset^{(n)}\) 可 \(\emptyset^{(n-1)}\)-递归，矛盾（跳跃运算严格递增图灵度）。\(\blacksquare\)

\textbf{定理336.6（Post定理------可定义性与图灵度的对应）} 集合 \(A \subseteq \mathbb{N}\) 满足：

\begin{itemize}
\tightlist
\item
  \(A\) 是递归可枚举的（\(\Sigma_1^0\)）当且仅当 \(A \leq_T \emptyset'\)（即 \(A\) 可由停机问题神谕计算）。
\item
  \(A \in \Delta_{n+1}^0\) 当且仅当 \(A \leq_T \emptyset^{(n)}\)（即 \(A\) 可由第 \(n\) 次图灵跳跃神谕计算）。
\item
  \(A \in \Sigma_{n+1}^0\) 当且仅当 \(A\) 是 \(\emptyset^{(n)}\)-递归可枚举的。
\end{itemize}

\textbf{证明} 对 \(n\) 归纳。\(n=0\)：\(A\) 是递归可枚举的 \(\iff\) \(A\) 可半判定 \(\iff\) 存在全递归谓词 \(R\) 使 \(A = \{x : \exists y R(x, y)\} \iff\) \(A\) 的判定可归约到 \(\exists\) 量词的判定，即 \(\emptyset'\)。归纳步骤利用跳跃的相对化版本和可定义性量化结构的归纳性质。\(\blacksquare\)

\subsection{图灵度与Post问题}\label{ux56feux7075ux5ea6ux4e0epostux95eeux9898}

\textbf{定义336.7（c.e.度与非c.e.度）} 递归可枚举集构成的图灵度称为\textbf{c.e.度}（computably enumerable degrees）。\(\mathbf{0}\) 是所有递归集所在的度。\(\mathbf{0}'\) 是停机问题所在的度（同时是最高的c.e.度）。

\textbf{定理336.8（Friedberg-Muchnik定理，1956-1957）} 存在两个递归可枚举集 \(A, B\) 使得 \(A\) 和 \(B\) 互不可图灵比较（即 \(A \nleq_T B\) 且 \(B \nleq_T A\)）。特别地，存在c.e.度 \(\mathbf{a}\) 满足 \(\mathbf{0} <_T \mathbf{a} <_T \mathbf{0}'\)，正面解决了Post问题（1944）。

\textbf{证明}（有限损害优先方法的完整构造）需同时满足无穷多需求：

\[R_{2e}: \Phi_e \neq A \quad \text{（} A \text{非递归）}\] \[R_{2e+1}: \Phi_e^A \neq B \quad \text{（} B \not\leq_T A\text{）}\] \[S_{2e}: \Phi_e \neq B\] \[S_{2e+1}: \Phi_e^B \neq A\]

需求按 \(R_0, S_0, R_1, S_1, R_2, S_2, \ldots\) 的优先级排序（\(R_0\) 最高优先）。

\textbf{第1步（需求的策略）}： - 对 \(R_{2e}\)：等待 \(\Phi_e\) 在某输入 \(x\) 上收敛于 \(0\)。若发生，将 \(x\) 放入 \(A\) 使得 \(A(x) = 1 \neq 0 = \Phi_e(x)\)。若 \(\Phi_e(x)\) 永不收敛于 \(0\)，则 \(A(x)\) 由低优先需求决定，此时 \(A(x) \neq \Phi_e(x)\) 自动成立（因为 \(\Phi_e(x)\) 发散或收敛于非零值）。 - 对 \(R_{2e+1}\)：等待某 \(x\) 使 \(\Phi_{e,s}^A(x)\) 收敛于 \(0\)（在当前近似 \(A_s\) 下）。若发生，将 \(x\) 放入 \(B\)，同时承诺保护 \(A\) 中当前使用的神谕信息（\(\Phi_{e,s}^A(x)\) 计算中查询过的 \(A\) 元素）不被修改。

\textbf{第2步（有限损伤机制）}：高优先需求改变 \(A\) 或 \(B\) 时可能损伤低优先需求------使其之前基于过时近似所做的决策失效。但关键是：每个需求 \(R_k\) 只能被有限多个更高优先需求损伤（因为每个高优先需求最多执行一次动作）。因此对每个 \(k\)，存在阶段 \(s_k\) 后 \(R_k\) 不再被损伤，此后 \(R_k\) 的策略永久有效。

\textbf{第3步（归纳验证）}：设所有需求 \(R_0, \ldots, R_{k-1}\) 最终稳定且被满足。在阶段 \(s \geq \max\{s_0, \ldots, s_{k-1}\}\) 之后，\(A\) 和 \(B\) 中与高优先需求相关的部分不再改变。\(R_k\) 的策略在此时启动，由于不再受损伤，最终满足。因此每个需求最终被满足，\(A\) 和 \(B\) 是递归可枚举的（通过逐步枚举证明中的''放入''动作），且互不可图灵比较。\(\blacksquare\)

上述优先方法不仅是Post问题的解答，更是递归论中构造c.e.集的核心技术，此后被推广至无穷损伤优先方法（第0'\,'\,'方法）、树优先方法（Lachlan，1970年代）等，深刻地影响了可计算性理论的发展。

\hrulefill

\section{第417章：算法信息论与Kolmogorov复杂度}\label{ux7b2c417ux7ae0ux7b97ux6cd5ux4fe1ux606fux8bbaux4e0ekolmogorovux590dux6742ux5ea6}

算法信息论将可计算性理论与概率论和信息论融合，以程序长度作为信息量的客观度量。Kolmogorov复杂度 \(C(x)\) 定义为生成字符串 \(x\) 的最短程序长度，它为随机性提供了一个数学上严格的、不以概率模型为前提的定义。本章从Kolmogorov复杂度的基本不等式出发（完整证明其不可计算性和次可加性），引入Levin通用搜索------这是算法信息论中最优搜索策略的理论基础，以不依赖问题先验分布的方式达到最优的期望搜索时间。在此基础上，讨论不可压缩性与Martin-Löf随机性检验的等价性，以及Chaitin不变量Omega（停机问题的停止概率）------一个随机且不可计算的实数，其二进制展开编码了停机问题的全部信息，且与Gödel不完备性定理有深刻的元数学联系。

\subsection{Kolmogorov复杂度基本不等式}\label{kolmogorovux590dux6742ux5ea6ux57faux672cux4e0dux7b49ux5f0f}

\textbf{定义337.1（Kolmogorov复杂度）} 设 \(U\) 为通用Turing机。字符串 \(x \in \{0,1\}^*\) 的\textbf{Kolmogorov复杂度}（或前缀Kolmogorov复杂度）定义为

\[K(x) = \min\{|p| : U(p) = x\}\]

即输出 \(x\) 的最短程序 \(p\) 的长度。若不存在这样的 \(p\)，定义 \(K(x) = \infty\)。该定义在不同通用Turing机之间的差异不超过常数（不变性定理）。

\textbf{定理337.2（Kolmogorov复杂度的不可计算性）} 函数 \(x \mapsto K(x)\) 不是部分递归的。事实上，不存在部分递归函数 \(\psi\) 使得 \(\psi(x) = K(x)\) 对无穷多 \(x\) 成立。

\textbf{证明}（对角化方法，Chaitin）假设 \(K\) 可计算（全递归或部分递归且定义在无穷集上）。定义算法：对 \(m = 1, 2, 3, \ldots\)，按字典序枚举所有字符串 \(x\)，若 \(K(x) > m\)（假设 \(K\) 可判定 \(K(x) > m\) 这个谓词），则输出第一个满足 \(K(x) > m\) 的 \(x\)。此算法确定了一个程序，其长度随 \(m\) 增长的速度为 \(\log m + O(1)\)。但该程序输出的字符串 \(x\) 的复杂度 \(K(x)\) 不超过该程序长度，即 \(K(x) \leq \log m + c\)；同时由构造 \(K(x) > m\)。取 \(m\) 充分大得矛盾：\(m < K(x) \leq \log m + c\) 不成立。\(\blacksquare\)

\textbf{定理337.3（Kolmogorov复杂度的基本不等式）}

（i）\textbf{上界}：\(K(x) \leq |x| + c\)（\(x\) 可被一个带硬编码输出的程序生成）。

（ii）\textbf{次可加性}：\(K(x, y) \leq K(x) + K(y \mid x) + O(\log K(x, y))\)，其中 \(K(y \mid x)\) 是给定 \(x\) 时 \(y\) 的条件Kolmogorov复杂度。

（iii）\textbf{对称性}：\(|K(x) - K(x, y)| \leq K(y) + O(\log K(x, y))\)（条件版本：\(K(x, y) = K(x) + K(y \mid x) + O(\log K(x, y))\)，即链式法则成立到对数精度）。

\textbf{证明}（i）显然：令 \(p\) 为 \(x\) 本身，通用Turing机输出其输入，故 \(K(x) \leq |x| + O(1)\)。

（ii）给定最短程序 \(p^*\)（输出 \(x\)）和 \(q^*\)（已知 \(x\) 时输出 \(y\)），可构造程序 \(\langle |p^*|, p^*, q^* \rangle\)（带自定界编码）先输出 \(x\)，再输出 \(y\)。自定界编码的代价为 \(O(\log K(x, y))\) 比特。此处使用条件复杂度 \(K(y \mid x)\) 的定义和最优前缀码的性质。

（iii）对称性来自 \(K(x, y) = K(x) + K(y \mid x^*)\)（其中 \(x^*\) 是 \(x\) 的最短程序），加上 \(K(y \mid x^*) \leq K(y) + O(1)\) 和 \(K(y) \leq K(y \mid x^*) + K(x) + O(\log \cdots)\)。\(\blacksquare\)

\subsection{Levin通用搜索}\label{levinux901aux7528ux641cux7d22}

\textbf{定理337.4（Levin通用搜索，1973）} 设计算问题可形式化为：给定输入 \(x\)，寻找满足 \(R(x, y)\) 的解 \(y\)（\(R\) 为递归谓词）。Levin搜索算法并行运行所有程序 \(p_1, p_2, p_3, \ldots\)，程序 \(p_i\) 分配的时间与 \(2^{-|p_i|}\) 成比例（按程序长度的指数加权）。则在所有可能的搜索策略中，Levin搜索的期望运行时间不超过任何特定搜索策略的期望运行时间乘以 \(2^{K(\text{策略})}\)------即达到几乎算子的最优性，代价为与最优策略的复杂度成正比（而非依赖先验概率）。

\textbf{证明} 设 \(\mathcal{A}\) 是任意搜索算法（即程序），其在输入 \(x\) 上的运行时间为 \(T_{\mathcal{A}}(x)\)。Levin搜索按字典序并行模拟所有程序，对长度为 \(\ell\) 的程序分配总时间的 \(2^{-\ell}\) 比例。算法 \(\mathcal{A}\) 本身是某个程序 \(p_{\mathcal{A}}\)，其在Levin搜索中得到的时间份额为 \(2^{-|p_{\mathcal{A}}|}\)。因此Levin搜索找到解的时刻不超过 \(T_{\mathcal{A}}(x) \cdot 2^{|p_{\mathcal{A}}|}\)（忽略并行模拟的常数开销）。对任意满足 \(U(p_{\mathcal{A}}) = \mathcal{A}\) 的程序 \(p_{\mathcal{A}}\)，\(2^{|p_{\mathcal{A}}|}\) 被 \(2^{K(\mathcal{A}) + O(1)}\) 控制。因此期望时间比 \(\leq 2^{K(\mathcal{A})}\)。\(\blacksquare\)

\subsection{不可压缩性与Martin-Löf随机性}\label{ux4e0dux53efux538bux7f29ux6027ux4e0emartin-luxf6fux968fux673aux6027}

\textbf{定义337.5（不可压缩字符串）} 字符串 \(x\) 称为 \textbf{\(c\)-不可压缩}，如果 \(K(x) \geq |x| - c\)。直观上，不可压缩的字符串没有比自身更短的描述------它们是''算法意义上随机的''。

\textbf{定理337.6（不可压缩字符串的存在性）} 对任意 \(n\)，至少存在 \(2^n(1 - 2^{-c})\) 个长度为 \(n\) 的 \(c\)-不可压缩字符串。特别地，当 \(n\) 充分大时，大多数字符串是不可压缩的。

\textbf{证明} 长度为 \(n\) 的字符串共 \(2^n\) 个。长度 \(\leq n-c-1\) 的程序总数不超过 \(\sum_{k=0}^{n-c-1} 2^k = 2^{n-c} - 1\)。故最多有 \(2^{n-c} - 1\) 个字符串的复杂度 \(< n-c\)。因此至少 \(2^n - 2^{n-c} + 1 = 2^n(1 - 2^{-c}) + 1\) 个字符串的复杂度 \(\geq n-c\)。\(\blacksquare\)

\textbf{定理337.7（Martin-Löf随机性与Kolmogorov复杂度的等价性）} 无穷二进制序列 \(X \in \{0,1\}^\omega\) 是Martin-Löf随机的当且仅当其所有前缀 \(X \upharpoonright n\) 满足 \(K(X \upharpoonright n) \geq n - O(1)\)。这一定理建立了算法信息论与测度论随机性概念之间的精确对应。

\textbf{证明}（\(\Rightarrow\)）若 \(X\) 不满足 \(K(X \upharpoonright n) \geq n - c\)（即无穷多 \(n\) 违反），则存在算法（以 \(n\) 为参数枚举复杂度低的串）可构造有效零测集（由可数个开覆盖组成），\(X\) 落入每个集合中，故 \(X\) 不是Martin-Löf随机的。（\(\Leftarrow\)）若 \(X\) 是Martin-Löf随机的，则由任意Martin-Löf检验 \(U_m\)（\(\lambda(U_m) \leq 2^{-m}\)）定义的复杂度上界 \(K(X \upharpoonright n) \geq n - \log n - m\) 对充分大 \(n\) 成立。取 \(m\) 充分大可得常数界。\(\blacksquare\)

\subsection{Chaitin不变量Omega}\label{chaitinux4e0dux53d8ux91cfomega}

\textbf{定义337.8（Chaitin Omega------停机概率）} 设 \(U\) 为无前缀（prefix-free）通用Turing机。Chaitin Omega定义为

\[\Omega = \sum_{p: U(p) \downarrow} 2^{-|p|}\]

即随机选取一个程序使其停机的概率（在无前缀码下的公平硬币模型）。\(\Omega\) 是良定义的实数（\(0 < \Omega < 1\)），其二进制展开是Martin-Löf随机的。

\textbf{定理337.9（Omega的不可计算性与随机性）} \(\Omega\) 不是递归实数（即其二进制展开不可计算）。事实上，\(\Omega\) 的图灵度数恰为 \(\mathbf{0}'\)（即与停机问题等价）。\(\Omega\) 是Martin-Löf随机的：\(K(\Omega \upharpoonright n) \geq n - O(1)\)。

\textbf{证明}（框架）\(\Omega\) 编码了停机问题的全部信息：\(U(p)\) 停机当且仅当 \(p\) 对 \(\Omega\) 的和有贡献。\(\Omega \equiv_T \emptyset'\) 的证明：由 \(\Omega\) 可判定停机问题（通过逐步近似 \(\Omega\) 的前缀），反之由停机问题可作为神谕计算 \(\Omega\) 的任意精度近似。

随机性的证明（Chaitin）：\(\Omega\) 的 \(n\) 比特前缀决定了所有长度 \(\leq n\) 的程序的停机状态（因前缀决定了 \(\Omega\) 的值到精度 \(2^{-n}\)）。若 \(K(\Omega \upharpoonright n) \ll n\)，则可压缩更多程序的行为，产生矛盾（通过Berry悖论风格的自我引用）。\(\blacksquare\)

\textbf{定理337.10（Omega与不完备性------Chaitin不完备性定理）} 任何包含充分算术的一致递归可公理化理论 \(T\) 只能判定有限多个形如''\(K(x) > m\)``的命题。等价地，\(T\) 只能确定 \(\Omega\) 的有限多个二进制位。这给出了Gödel不完备性定理的信息论阐释：任何形式系统包含的信息量是有界的，因此不能永远产生复杂度更高的定理。

\hrulefill

\hrulefill

\hrulefill

\hrulefill

\hrulefill

\newpage

\chapter{11-数值与计算/02-计算理论/06-图像处理的数学基础.md}\label{ux6570ux503cux4e0eux8ba1ux7b9702-ux8ba1ux7b97ux7406ux8bba06-ux56feux50cfux5904ux7406ux7684ux6570ux5b66ux57faux7840.md}

\chapter{图像处理的数学基础}\label{ux56feux50cfux5904ux7406ux7684ux6570ux5b66ux57faux7840}

图像处理是应用数学与计算机科学交叉的核心领域，其数学基础根植于调和分析、偏微分方程、变分法与最优化理论。图像可以建模为连续函数（理想模型）或离散像素阵列（计算模型），二者的联系由采样定理与逼近论给出。线性滤波与傅里叶变换提供了频域分析与设计的工具，而变分方法则将去噪、分割、修复等任务转化为泛函极值问题。压缩感知作为21世纪初的突破性理论，将稀疏性先验与随机采样相结合，实现了远低于Nyquist速率的信号精确恢复。本章从经典方法出发，逐步深入到现代理论，以定理-证明的方式展示图像处理数学基础的核心结构。

\hrulefill

\section{第418章：图像处理的数学框架}\label{ux7b2c418ux7ae0ux56feux50cfux5904ux7406ux7684ux6570ux5b66ux6846ux67b6}

图像处理的数学框架建立在连续模型与离散模型的对偶之上。连续模型将图像视为定义在区域 \(\Omega \subset \mathbb{R}^2\) 上的函数 \(u: \Omega \to \mathbb{R}\)（灰度图像）或 \(u: \Omega \to \mathbb{R}^3\)（彩色图像），其分析工具主要是实分析、泛函分析与偏微分方程。离散模型则将图像视为有限维向量，引入线性代数与数值方法。线性滤波、傅里叶分析与偏微分方程构成了从经典到现代图像处理的连续谱系。

\subsection{x.1 连续与离散图像模型}\label{x.1-ux8fdeux7eedux4e0eux79bbux6563ux56feux50cfux6a21ux578b}

\textbf{定义}：连续图像模型将（灰度）图像视为可测函数 \(u \in L^2(\Omega)\) 或更光滑的空间如 \(H^1(\Omega)\)、\(BV(\Omega)\)（有界变差函数空间）。离散模型将 \(M \times N\) 图像视为向量 \(u \in \mathbb{R}^{MN}\)（按列展开）。采样算子 \(S: L^2(\Omega) \to \mathbb{R}^{MN}\) 将连续函数映射到离散值 \(S(u)_{ij} = u(x_i, y_j)\)（或局部平均值）。

\textbf{卷积与线性滤波}：连续卷积 \((K * u)(x) = \int_{\Omega} K(x-y)u(y)\,dy\)。离散卷积 \((K * u)[m,n] = \sum_{i,j} K[i,j]\,u[m-i, n-j]\)。常见滤波核：

\begin{itemize}
\tightlist
\item
  Gaussian平滑：\(G_\sigma(x) = \frac{1}{2\pi\sigma^2}e^{-|x|^2/(2\sigma^2)}\)
\item
  Sobel梯度：\(S_x = \begin{pmatrix}-1&0&1\\-2&0&2\\-1&0&1\end{pmatrix}\)，\(S_y = \begin{pmatrix}-1&-2&-1\\0&0&0\\1&2&1\end{pmatrix}\)
\end{itemize}

\subsection{x.2 傅里叶分析与采样定理}\label{x.2-ux5085ux91ccux53f6ux5206ux6790ux4e0eux91c7ux6837ux5b9aux7406}

\textbf{定理 x.1}（Nyquist-Shannon采样定理）：设 \(u \in L^2(\mathbb{R})\) 满足 \(\operatorname{supp}(\widehat{u}) \subseteq [-B, B]\)（带限信号，带宽为 \(B\)）。则 \(u\) 可由其等距采样值 \(\{u(n/2B)\}_{n \in \mathbb{Z}}\) 精确重构：

\[u(x) = \sum_{n \in \mathbb{Z}} u\left(\frac{n}{2B}\right) \operatorname{sinc}(2Bx - n),\]

其中 \(\operatorname{sinc}(t) = \sin(\pi t)/(\pi t)\)。采样率 \(f_s = 2B\) 称为Nyquist率。

\textbf{证明}：由于 \(\widehat{u}\) 支撑在 \([-B, B]\)，可将其延拓为周期为 \(2B\) 的周期函数 \(\widehat{u}_{\text{per}}\)。将 \(\widehat{u}_{\text{per}}\) 展开为Fourier级数：

\[\widehat{u}_{\text{per}}(\xi) = \sum_{n \in \mathbb{Z}} c_n e^{-\pi i n\xi/B}, \quad |\xi| \leq B.\]

Fourier系数：

\[c_n = \frac{1}{2B}\int_{-B}^B \widehat{u}(\xi)\,e^{\pi i n\xi/B}\,d\xi = \frac{1}{2B}\,u\left(\frac{n}{2B}\right).\]

对 \(\widehat{u}\) 做Fourier逆变换：

\[u(x) = \int_{-B}^B \widehat{u}(\xi)\,e^{2\pi i\xi x}\,d\xi = \int_{-B}^B \sum_{n \in \mathbb{Z}} \frac{u(n/2B)}{2B} e^{\pi i n\xi/B} e^{2\pi i\xi x}\,d\xi.\]

交换积分与求和（由正则性保证），

\[u(x) = \sum_{n \in \mathbb{Z}} \frac{u(n/2B)}{2B} \int_{-B}^B e^{2\pi i\xi(x - n/(2B))}\,d\xi = \sum_{n \in \mathbb{Z}} u\left(\frac{n}{2B}\right) \frac{\sin(2\pi B(x - n/(2B)))}{2\pi B(x - n/(2B))} = \sum_{n \in \mathbb{Z}} u\left(\frac{n}{2B}\right) \operatorname{sinc}(2Bx - n).\]

此重构公式在 \(L^2\) 意义下收敛，且在连续点处逐点收敛。\(\blacksquare\)

\subsection{x.3 偏微分方程方法}\label{x.3-ux504fux5faeux5206ux65b9ux7a0bux65b9ux6cd5}

\textbf{热扩散方程}：图像去噪的最简PDE模型为线性扩散：

\[\frac{\partial u}{\partial t} = \Delta u, \quad u(x,0) = u_0(x).\]

解为 \(u(\cdot,t) = G_{\sqrt{2t}} * u_0\)，它等价于Gaussian平滑。问题在于线性扩散在去除噪声的同时模糊了边缘。

\textbf{定理 x.2}（Perona-Malik各向异性扩散）：Perona-Malik方程

\[\frac{\partial u}{\partial t} = \nabla \cdot (g(|\nabla u|)\nabla u),\]

其中 \(g: [0,\infty) \to (0,\infty)\) 为递减函数，满足 \(g(0) = 1\)，\(\lim_{s \to \infty} g(s) = 0\)（如 \(g(s) = 1/(1 + (s/K)^2)\)）。该方程在梯度较小（平滑区域）处近似于线性扩散，在梯度较大（边缘）处抑制扩散，从而实现保持边缘的平滑。

\textbf{证明}（稳定性分析）：考虑Perona-Malik方程的局部行为。在点 \(x\) 处，扩散张量为 \(g(|\nabla u|)\)。将方程沿梯度方向（法向）和切线方向分解：

\[\frac{\partial u}{\partial t} = g(|\nabla u|)u_{TT} + (g(|\nabla u|) + g'(|\nabla u|)|\nabla u|)u_{NN},\]

其中 \(u_{TT}\) 为沿边缘的二阶导数，\(u_{NN}\) 为跨边缘的二阶导数。法向扩散系数为 \(g(s) + sg'(s)\)。为使方程在边缘处抑制扩散（甚至反向扩散以增强边缘），需 \(g(s) + sg'(s) < 0\) 对于大 \(s\)。此条件在 \(g(s) = 1/(1+(s/K)^2)\) 时对 \(s > K\) 成立。然而，反向扩散导致方程不适定，需通过正则化（如 \(g(|\nabla G_\sigma * u|)\) 替代 \(g(|\nabla u|)\)）恢复适定性。\(\blacksquare\)

\hrulefill

\section{第419章：变分方法与正则化}\label{ux7b2c419ux7ae0ux53d8ux5206ux65b9ux6cd5ux4e0eux6b63ux5219ux5316}

变分方法将图像处理任务表述为能量泛函的极小化问题。其基本形式为

\[E(u) = D(u, u_0) + \lambda R(u),\]

其中 \(D\) 为数据保真项，\(R\) 为正则化项，\(\lambda > 0\) 为平衡参数。不同的正则化项诱导不同的解空间特性：Tikhonov正则化鼓励光滑性，总变差正则化保持边缘，Mumford-Shah泛函则显式地允许不连续集。

\subsection{x+1.1 Tikhonov正则化}\label{x1.1-tikhonovux6b63ux5219ux5316}

\textbf{定理 x+1.1}（Tikhonov正则化的适定性）：设 \(A: X \to Y\) 为Hilbert空间之间的有界线性紧算子。对 \(\alpha > 0\)，Tikhonov泛函

\[J_\alpha(u) = \|Au - f\|_Y^2 + \alpha\|u\|_X^2\]

在 \(X\) 中存在唯一极小元 \(u_\alpha\)，且 \(u_\alpha\) 满足正规方程 \((A^*A + \alpha I)u_\alpha = A^*f\)。进一步，当噪声水平 \(\delta \to 0\) 且 \(\alpha = \alpha(\delta)\) 满足适当条件时，\(u_{\alpha(\delta)} \to u^\dagger\)（真解）。

\textbf{证明}：\(J_\alpha\) 是严格凸的连续泛函（\(A^*A + \alpha I\) 是正定自伴算子），其Gateaux导数为 \(J_\alpha'(u) = 2(A^*(Au - f) + \alpha u)\)。由一阶最优性条件 \(J_\alpha'(u_\alpha) = 0\) 即得正规方程。由于 \(A^*A + \alpha I \geq \alpha I\) 严格正定，解唯一存在。

收敛性分析：设 \(f^\delta\) 为含噪声数据，\(\|f^\delta - f\| \leq \delta\)。定义 \(u_\alpha^\delta\) 为以 \(f^\delta\) 代 \(f\) 的Tikhonov解。误差分解：

\[\|u_\alpha^\delta - u^\dagger\| \leq \|u_\alpha^\delta - u_\alpha\| + \|u_\alpha - u^\dagger\|.\]

第一项（噪声传播）：\(\|u_\alpha^\delta - u_\alpha\| \leq \|(A^*A + \alpha I)^{-1}A^*\| \,\delta \leq \delta/(2\sqrt{\alpha})\)（利用谱映射定理）。第二项（正则化误差）：当 \(\alpha \to 0\) 时 \(\|u_\alpha - u^\dagger\| \to 0\)（由紧凑性和Tikhonov正则化的一般理论）。选取 \(\alpha = \alpha(\delta)\) 使得 \(\delta/\sqrt{\alpha(\delta)} \to 0\) 且 \(\alpha(\delta) \to 0\)，即得收敛性。\(\blacksquare\)

\subsection{x+1.2 ROF全变差去噪模型}\label{x1.2-rofux5168ux53d8ux5deeux53bbux566aux6a21ux578b}

\textbf{定理 x+1.2}（ROF模型及其对偶）：设 \(\Omega \subset \mathbb{R}^2\) 为有界开集，\(u_0 \in L^2(\Omega)\) 为含噪图像。Rudin-Osher-Fatemi模型的原始形式为

\[\min_{u \in BV(\Omega)} \left\{\|u - u_0\|_{L^2}^2 + \lambda \operatorname{TV}(u)\right\},\]

其中 \(\operatorname{TV}(u) = \sup\{\int_\Omega u \,\nabla \cdot \mathbf{p}\,dx \mid \mathbf{p} \in C_c^1(\Omega; \mathbb{R}^2), \|\mathbf{p}\|_\infty \leq 1\}\)。该问题的对偶形式为

\[\min_{\mathbf{p} \in C_c^1(\Omega;\mathbb{R}^2), \|\mathbf{p}\|_\infty \leq 1} \|u_0 - \lambda\nabla \cdot \mathbf{p}\|_{L^2}^2.\]

最优解满足 \(u = u_0 - \lambda\nabla \cdot \mathbf{p}\) 且 \(\mathbf{p}\) 为对偶问题的极小元。

\textbf{证明}：先证原始问题在 \(BV(\Omega)\) 中的存在唯一性。由 \(BV\) 空间的紧嵌入性质，极小化序列在 \(L^1\) 中紧（\(BV\) 有界序列的紧性），下半连续性由全变差的弱*下半连续性保证。唯一性来自 \(\ell^2\) 项的严格凸性。

对偶推导：利用Legendre-Fenchel变换。TV的正则化项可写为

\[\lambda \operatorname{TV}(u) = \sup_{\mathbf{p} \in \mathcal{P}} \int_\Omega u \,\nabla \cdot \mathbf{p}\,dx,\]

其中 \(\mathcal{P} = \{\mathbf{p} \in C_c^1(\Omega;\mathbb{R}^2) \mid \|\mathbf{p}\|_\infty \leq 1\}\)。代入原始泛函：

\[\min_u \sup_{\mathbf{p} \in \mathcal{P}} \left\{\|u - u_0\|_{L^2}^2 + \lambda \int_\Omega u \,\nabla \cdot \mathbf{p}\,dx\right\}.\]

由Fenchel-Rockafellar对偶定理（条件满足：\(L^2\) 和 \(BV\) 均为自反空间或其对偶），交换 \(\min\) 和 \(\sup\) 的顺序。对固定的 \(\mathbf{p}\)，关于 \(u\) 求极小的最优性条件为 \(2(u - u_0) + \lambda \nabla \cdot \mathbf{p} = 0\)，即 \(u = u_0 - \frac{\lambda}{2}\nabla \cdot \mathbf{p}\)。代入（重参数化 \(\lambda\)）：

\[\max_{\mathbf{p} \in \mathcal{P}} \left\{\left\|\frac{\lambda}{2}\nabla \cdot \mathbf{p}\right\|_{L^2}^2 + \lambda\int_\Omega (u_0 - \frac{\lambda}{2}\nabla \cdot \mathbf{p})\,\nabla \cdot \mathbf{p}\,dx\right\} = \max_{\mathbf{p} \in \mathcal{P}} \left\{\lambda\int_\Omega u_0\nabla \cdot \mathbf{p}\,dx - \frac{\lambda^2}{4}\|\nabla \cdot \mathbf{p}\|_{L^2}^2\right\}.\]

等价地，\(\min_{\mathbf{p} \in \mathcal{P}} \|u_0 - \frac{\lambda}{2}\nabla \cdot \mathbf{p}\|_{L^2}^2\)（经配方）。这就是ROF模型的对偶形式。\(\blacksquare\)

\subsection{x+1.3 Chambolle投影算法}\label{x1.3-chambolleux6295ux5f71ux7b97ux6cd5}

\textbf{定理 x+1.3}（Chambolle算法）：ROF对偶问题的解可通过以下投影梯度迭代得到：

\[\mathbf{p}^{k+1} = \frac{\mathbf{p}^k + \tau \nabla(\nabla \cdot \mathbf{p}^k - u_0/\lambda)}{1 + \tau \|\nabla(\nabla \cdot \mathbf{p}^k - u_0/\lambda)\|},\]

其中 \(\tau > 0\) 为步长。当 \(\tau \leq 1/8\) 时迭代收敛。

\textbf{证明}：对偶问题等价于求 \(\mathbf{p}\) 使得 \(\mathbf{p} = \Pi_{\mathcal{P}}(\mathbf{p} + \tau \nabla(\nabla \cdot \mathbf{p} - u_0/\lambda))\)（投影梯度法的固定点方程）。\(\Pi_{\mathcal{P}}\) 为逐点投影到单位球上的算子：\((\Pi_{\mathcal{P}}(\mathbf{q}))(x) = \mathbf{q}(x)/\max(1, |\mathbf{q}(x)|)\)。

收敛性来自投影梯度算法的一般理论：目标函数 \(\|\nabla \cdot \mathbf{p} - u_0/\lambda\|^2\) 对 \(\mathbf{p}\) 是凸且连续可微的，梯度Lipschitz常数为 \(L = \|\nabla\|^2 \|\nabla\cdot\|^2 \leq 8\)（在二维情形，由离散梯度算子的谱半径估计）。步长 \(\tau < 2/L\) 保证收敛，\(\tau \leq 1/8\) 为安全选择。\(\blacksquare\)

\subsection{x+1.4 Mumford-Shah分割泛函}\label{x1.4-mumford-shahux5206ux5272ux6cdbux51fd}

\textbf{定义}（Mumford-Shah泛函）：设 \(\Omega \subset \mathbb{R}^2\) 为图像域，\(u_0\) 为观测图像。Mumford-Shah泛函定义为

\[E(u, \Gamma) = \int_{\Omega \setminus \Gamma} |\nabla u|^2\,dx + \mu \int_{\Omega \setminus \Gamma} |u - u_0|^2\,dx + \nu \mathcal{H}^1(\Gamma),\]

其中 \(\Gamma \subset \Omega\) 为分割曲线（跳跃集），\(\mathcal{H}^1\) 为一维Hausdorff测度（曲线长度），\(\mu, \nu > 0\) 为参数。极小化在 \(u \in C^1(\Omega \setminus \Gamma)\) 和 \(\Gamma\) 为闭集下进行。

\textbf{定理 x+1.4}（Mumford-Shah极小元的存在性）：存在 \((u, \Gamma)\) 极小化 \(E\)，其中 \(\Gamma\) 是闭集且 \(\mathcal{H}^1(\Gamma) < \infty\)，\(u \in C^1(\Omega \setminus \Gamma)\)。

\textbf{证明}：利用SBV（special bounded variation）空间的弱紧性。将Mumford-Shah泛函弱化为：

\[E(u) = \int_\Omega |\nabla u|^2\,dx + \mu\int_\Omega |u - u_0|^2\,dx + \nu \mathcal{H}^1(S_u),\]

其中 \(S_u\) 为 \(u\) 的跳跃集（近似不连续点）。在SBV空间中，由紧性定理（Ambrosio的SBV紧性定理），极小化序列存在于SBV中，且极限函数的跳跃集的Hausdorff测度下半连续。再由能量泛函的下半连续性，极限函数为极小解。\(\blacksquare\)

\hrulefill

\section{第420章：稀疏表示与压缩感知}\label{ux7b2c420ux7ae0ux7a00ux758fux8868ux793aux4e0eux538bux7f29ux611fux77e5}

压缩感知（Compressed Sensing）由Candes、Romberg和Tao以及Donoho各自独立于2006年提出，其核心断言是：稀疏信号可由远少于Nyquist率要求的随机线性测量精确恢复。该理论的基础是限制等距性质（Restricted Isometry Property, RIP），它刻画了测量矩阵对稀疏向量近乎等距作用的条件。压缩感知引发了信号处理和图像重建领域的革命，其数学基础根植于凸优化、随机矩阵理论和几何泛函分析。

\subsection{x+2.1 限制等距性质}\label{x2.1-ux9650ux5236ux7b49ux8dddux6027ux8d28}

\textbf{定义}：矩阵 \(A \in \mathbb{R}^{m \times n}\)（\(m < n\)）满足 \(s\) 阶限制等距性质（RIP），若存在常数 \(\delta_s \in (0,1)\) 使得对任意至多 \(s\) 个非零元（\(s\)-稀疏）的向量 \(x \in \mathbb{R}^n\)，有

\[(1 - \delta_s)\|x\|_2^2 \leq \|Ax\|_2^2 \leq (1 + \delta_s)\|x\|_2^2.\]

\(\delta_s\) 称为 \(s\) 阶RIP常数。

\textbf{定理 x+2.1}（随机矩阵的RIP）：设 \(A \in \mathbb{R}^{m \times n}\) 的元素满足 \(A_{ij} \sim \mathcal{N}(0, 1/m)\)（或适当的次高斯分布独立同分布）。则存在常数 \(C > 0\) 使得当

\[m \geq C \,\delta^{-2} s \log(n/s)\]

时，\(A\) 以高概率满足 \(s\) 阶RIP且 \(\delta_s \leq \delta\)。

\textbf{证明概要}：对任意固定的支撑集 \(T \subset \{1,\ldots,n\}\)（\(|T| = s\)），\(A_T\) 为 \(m \times s\) 子矩阵。要证 \(\|A_T^*A_T - I_s\| \leq \delta_s\)（谱范数）。对高斯系综，矩阵 \(A_T^*A_T\) 的期望为 \(I_s\)，偏差由Vershynin的次高斯矩阵浓度不等式控制：

\[\mathbb{P}\left(\|A_T^*A_T - I_s\| > \varepsilon\right) \leq 2 \exp(-c m \varepsilon^2),\]

其中 \(c > 0\) 为绝对常数。对 \(\binom{n}{s}\) 个可能的支撑集取联合界：

\[\mathbb{P}\left(\exists T, |T|=s, \|A_T^*A_T - I_s\| > \delta\right) \leq 2\binom{n}{s} e^{-c m \delta^2}.\]

由 \(\binom{n}{s} \leq (en/s)^s\) 及 \(m \geq C\delta^{-2}s\log(n/s)\)（\(C\) 充分大），上述概率小于1。\(\blacksquare\)

\subsection{\texorpdfstring{x+2.2 \(\ell_1\) 最小化与精确恢复}{x+2.2 \textbackslash ell\_1 最小化与精确恢复}}\label{x2.2-ell_1-ux6700ux5c0fux5316ux4e0eux7cbeux786eux6062ux590d}

\textbf{定理 x+2.2}（Candes-Romberg-Tao精确恢复定理）：设 \(x \in \mathbb{R}^n\) 是 \(s\)-稀疏向量，\(A\) 满足 \(2s\) 阶RIP且 \(\delta_{2s} < \sqrt{2} - 1\)。则 \(\ell_1\) 最小化问题

\[\min_{z \in \mathbb{R}^n} \|z\|_1 \quad \text{s.t.} \quad Az = Ax\]

有唯一解，且该解恰为 \(x\)（精确恢复）。

\textbf{证明}：设 \(x^*\) 为 \(\ell_1\) 极小化问题的解，\(h = x^* - x\) 为误差向量。目标是证明 \(h = 0\)。

将 \(h\) 分解为 \(h = h_T + h_{T^c}\)，其中 \(T = \operatorname{supp}(x)\)（\(|T| = s\)）。将 \(h_{T^c}\) 按其绝对值的递减序写为 \(h_{T^c} = h_{T_1} + h_{T_2} + \cdots\)，其中 \(T_1\) 为 \(T^c\) 中 \(s\) 个最大幅值分量的索引，\(T_2\) 为次 \(s\) 个最大分量，等等。

由最优性条件：\(\|x^*\|_1 \leq \|x\|_1\)，故

\[\|x_T\|_1 \geq \|x^*|_T\|_1 + \|h_{T^c}\|_1 \geq \|x_T\|_1 - \|h_T\|_1 + \|h_{T^c}\|_1,\]

从而 \(\|h_{T^c}\|_1 \leq \|h_T\|_1\)（锥约束）。

另一方面，由 \(Ah = 0\) 和RIP，对分段 \(h_{T \cup T_1}\) 应用RIP：

\[0 = \|Ah\| = \|A(h_{T \cup T_1}) + \sum_{j \geq 2} A h_{T_j}\| \geq \|A h_{T \cup T_1}\| - \sum_{j \geq 2}\|A h_{T_j}\|.\]

由RIP和范数等价性（\(\ell_1\) 到 \(\ell_2\) 的不等式），经细致估计：

\[\|h_{T \cup T_1}\|_2^2 \leq \frac{\sqrt{2}\delta_{2s}}{1-\delta_{2s}} \|h_T\|_2 \|h_{T \cup T_1}\|_2.\]

由 \(\delta_{2s} < \sqrt{2} - 1\) 得 \(\frac{\sqrt{2}\delta_{2s}}{1-\delta_{2s}} < 1\)，从而 \(\|h_T\|_2 = 0\)，推出 \(h = 0\)。更精确的常数可优化至 \(\delta_{2s} < 0.465\) 等更优界。\(\blacksquare\)

\subsection{x+2.3 噪声鲁棒性}\label{x2.3-ux566aux58f0ux9c81ux68d2ux6027}

\textbf{定理 x+2.3}（含噪声的稳定恢复）：设 \(y = Ax + e\)，\(\|e\|_2 \leq \varepsilon\)，\(x\) 为 \(s\)-稀疏向量，\(A\) 满足 \(2s\) 阶RIP且 \(\delta_{2s} < \sqrt{2} - 1\)。则

\[\hat{x} = \arg\min_z \|z\|_1 \text{ s.t. } \|Az - y\|_2 \leq \varepsilon\]

满足 \(\|\hat{x} - x\|_2 \leq C_0 \varepsilon\)，其中 \(C_0\) 为仅依赖于 \(\delta_{2s}\) 的常数。

\textbf{证明}：设 \(h = \hat{x} - x\)。利用前述分解和约束条件 \(\|Ah\|_2 \leq 2\varepsilon\)，类似精确恢复的论证，但加入噪声项的处理。锥约束仍成立，结合RIP得 \(\|h\|_2 \leq C_0 \varepsilon\)，其中 \(C_0 = \frac{4\sqrt{1+\delta_{2s}}}{1 - (1+\sqrt{2})\delta_{2s}}\)。\(\blacksquare\)

\subsection{x+2.4 字典学习与图像修复}\label{x2.4-ux5b57ux5178ux5b66ux4e60ux4e0eux56feux50cfux4feeux590d}

\textbf{定义}：字典学习问题为：给定数据矩阵 \(Y \in \mathbb{R}^{n \times p}\)，寻找字典 \(D \in \mathbb{R}^{n \times K}\)（\(K > n\)，过完备情形）和稀疏系数矩阵 \(X \in \mathbb{R}^{K \times p}\)，使得

\[\min_{D, X} \|Y - DX\|_F^2 \quad \text{s.t.} \quad \|x_i\|_0 \leq s \;(i=1,\ldots,p).\]

\textbf{图像修复（inpainting）的数学模型}：设 \(\Omega_0 \subset \Omega\) 为已知区域，\(\Omega \setminus \Omega_0\) 为待修复区域。修复问题可建模为带掩码的变分问题：

\[\min_u \left\{\int_\Omega |\nabla u|^p\,dx + \frac{\lambda}{2}\int_{\Omega_0} |u - u_0|^2\,dx\right\},\]

其中 \(p \in [1,2]\)。\(p=2\) 对应调和修复（Cahn-Hilliard型），\(p=1\) 对应TV修复。后者在保持边缘方面表现更优。该模型存在唯一解（\(p=1\) 时在BV空间中，\(p>1\) 时在Sobolev空间中）。

\textbf{定理 x+2.4}（TV修复的适定性）：\(\Omega_0\) 为 \(\Omega\) 的非空开子集时，TV修复问题存在唯一极小解 \(u \in BV(\Omega)\)。

\textbf{证明}：与ROF模型的论证类似，利用BV空间的紧嵌入和下半连续性。唯一性来自数据保真项的严格凸性在 \(\Omega_0\) 上的作用（\(\lambda > 0\) 且 \(\Omega_0\) 非零测度保证了唯一性）。存在性由直接变分法得到。\(\blacksquare\)

\subsection{x+2.5 非局部方法与图割}\label{x2.5-ux975eux5c40ux90e8ux65b9ux6cd5ux4e0eux56feux5272}

图像处理中的非局部方法利用图像中的重复模式，在非局部域中定义滤波和正则化。非局部均值（Non-Local Means, NL-Means）滤波器的定义如下。

\textbf{定义}（NL-Means滤波）：对含噪图像 \(u_0\)，NL-Means估计为

\[\hat{u}(x) = \frac{1}{C(x)}\int_\Omega e^{-\frac{(G_\sigma * |u_0(x+\cdot) - u_0(y+\cdot)|^2)(0)}{h^2}} u_0(y)\,dy,\]

其中 \(G_\sigma\) 为Gaussian核，\(h\) 为滤波参数，\(C(x)\) 为归一化常数。此为非局部滤波器，权重由以 \(x\) 和 \(y\) 为中心的图像块的相似性决定。

\textbf{定理 x+2.5}（图割与连续最大流/最小割）：图像分割可表述为能量泛函

\[E(C) = \int_{C} g(|\nabla u_0|)\,ds + \lambda\int_{\Omega} |u_0 - c_1|^2\,dx + \lambda\int_{\Omega \setminus \text{inside}(C)} |u_0 - c_2|^2\,dx\]

的极小化，其中 \(C\) 为分割曲线，\(g\) 为边缘指示函数。该问题的全局最优解可由图割（graph cut）算法的连续推广------最大流/最小割定理------求得。

\textbf{证明思路}：离散化后的能量可表述为马尔可夫随机场（MRF）的能量极小化。Boykov-Kolmogorov算法使用最小割（s-t cut）方法将此能量全局极小化。连续版本由Chan-Vese模型（水平集方法）处理，对应的PDE为

\[\frac{\partial\phi}{\partial t} = \delta_\varepsilon(\phi)\left[\nabla\cdot\left(g\frac{\nabla\phi}{|\nabla\phi|}\right) - \lambda_1(u_0 - c_1)^2 + \lambda_2(u_0 - c_2)^2\right].\]

\subsection{x+2.6 深度学习与图像处理的数学联系}\label{x2.6-ux6df1ux5ea6ux5b66ux4e60ux4e0eux56feux50cfux5904ux7406ux7684ux6570ux5b66ux8054ux7cfb}

深度卷积神经网络（CNN）与图像处理的经典方法之间存在深刻的数学联系。CNN可视为逐层学习的非线性滤波器和特征提取器的复合，其结构与多尺度表示和稀疏编码密切相关。

\textbf{定理 x+2.6}（展开式迭代网络与ISTA）：迭代收缩阈值算法（ISTA）为求解 \(\ell_1\) 正则化问题的经典迭代法：

\[x_{k+1} = S_{\lambda/L}\left(x_k - \frac{1}{L}A^T(Ax_k - y)\right),\]

其中 \(S_\theta\) 为软阈值算子：\(S_\theta(t) = \operatorname{sign}(t)\max(0, |t|-\theta)\)。深度展开网络（如LISTA）将该迭代的多步展开为一个可训练的神经网络，每步的矩阵和阈值作为可学习参数。当步数有限时，该网络是ISTA迭代的截断近似，其逼近性质由ISTA的线性收敛速率保证。

\textbf{证明}：ISTA的收敛性由邻近梯度法的标准理论给出：对 \(\ell_1\) 正则化问题 \(\min_x \frac{1}{2}\|Ax-y\|_2^2 + \lambda\|x\|_1\)，目标函数满足Kurdyka-Lojasiewicz性质，邻近梯度法以线性速率收敛。LISTA（Learned ISTA）将每步迭代参数化为可学习的线性算子 \(W_k\) 和 \(S_k\)：\(x_{k+1} = S_{\theta_k}(W_k x_k + S_k y)\)。在适当的参数限制下，该网络保持了与ISTA相同的表达性，但可以通过反向传播端到端优化。\(\blacksquare\)

\subsection{x+2.7 Wasserstein距离在图像分析中的应用}\label{x2.7-wassersteinux8dddux79bbux5728ux56feux50cfux5206ux6790ux4e2dux7684ux5e94ux7528}

最优传输（Optimal Transport, OT）理论为图像分析提供了几何上更有意义的距离度量。给定两幅图像（视为测度）\(\mu\) 和 \(\nu\)（归一化直方图），\(p\)-Wasserstein距离定义为

\[W_p(\mu, \nu) = \left(\inf_{\pi \in \Pi(\mu,\nu)} \int_{\Omega \times \Omega} \|x-y\|^p\,d\pi(x,y)\right)^{1/p},\]

其中 \(\Pi(\mu,\nu)\) 为边缘分布为 \(\mu\) 和 \(\nu\) 的所有联合分布。

\textbf{定理 x+2.7}（Wasserstein距离的Kantorovich对偶）：Wasserstein-1距离的对偶形式为

\[W_1(\mu, \nu) = \sup_{\|\nabla f\|_\infty \leq 1} \left\{\int_\Omega f\,d(\mu-\nu)\right\},\]

即1-Lipschitz函数空间中的最大期望差异。此对偶形式使得 \(W_1\) 可由计算简单的最优原像问题来估计，并在Wasserstein生成对抗网络（WGAN）中得到关键应用。

\textbf{证明}：由Kantorovich-Rubinstein定理，对于紧致空间上的概率测度，最优传输问题的对偶为

\[\inf_{\pi} \int c(x,y)d\pi = \sup_f \left\{\int f^c d\mu - \int f d\nu\right\},\]

其中 \(f^c(y) = \inf_x\{c(x,y) - f(x)\}\) 为 \(c\)-变换。当 \(c(x,y) = \|x-y\|\) 时，\(f^c = -f\) 若 \(f\) 为1-Lipschitz，从而上述对偶化为 \(W_1\) 的形式。\(\blacksquare\)

\subsection{x+2.8 图像处理中的多尺度与多分辨率分析}\label{x2.8-ux56feux50cfux5904ux7406ux4e2dux7684ux591aux5c3aux5ea6ux4e0eux591aux5206ux8fa8ux7387ux5206ux6790}

多分辨率分析（Multiresolution Analysis, MRA）是小波理论的核心，为图像的多尺度表示提供了严格的数学框架。

\textbf{定义}（MRA）：\(L^2(\mathbb{R})\) 的多分辨率分析是一列嵌套闭子空间 \(\cdots \subset V_{-1} \subset V_0 \subset V_1 \subset \cdots \subset L^2(\mathbb{R})\)，满足：（i）\(\bigcap V_j = \{0\}\)，\(\overline{\bigcup V_j} = L^2(\mathbb{R})\)；（ii）\(f(x) \in V_j \iff f(2x) \in V_{j+1}\)；（iii）存在尺度函数 \(\phi \in V_0\) 使得 \(\{\phi(x-k) \mid k \in \mathbb{Z}\}\) 为 \(V_0\) 的标准正交基。

\textbf{定理 x+2.8}（Mallat算法与离散小波变换）：由MRA可构造小波函数 \(\psi\) 使得 \(\{\psi_{j,k}(x) = 2^{j/2}\psi(2^j x - k)\}\) 构成 \(L^2(\mathbb{R})\) 的标准正交基。图像的小波系数可由级联滤波器组（cascade filter bank）高效计算：分解算法 \(c_{j-1}[n] = \sum_k h[k-2n]c_j[k]\)，\(d_{j-1}[n] = \sum_k g[k-2n]c_j[k]\)，其中 \(h\) 和 \(g\) 分别为低通和高通滤波器，重构算法 \(c_j[n] = \sum_k h[n-2k]c_{j-1}[k] + \sum_k g[n-2k]d_{j-1}[k]\)。

\textbf{证明}：由双尺度方程 \(\phi(x) = \sqrt{2}\sum_n h[n]\phi(2x-n)\) 和 \(\psi(x) = \sqrt{2}\sum_n g[n]\phi(2x-n)\)，其中 \(h[n] = \langle\phi_{1,0}, \phi_{0,n}\rangle\)，\(g[n] = (-1)^n h[1-n]\)。将 \(V_j\) 中的函数按基展开：\(f = \sum_k c_j[k]\phi_{j,k}\)。利用双尺度关系得 \(c_{j-1}[n] = \sum_k h[k-2n]c_j[k]\)，此即Mallat算法的核心递推关系。二维图像的小波变换通过张量积构造得到二维MRA。\(\blacksquare\)

\subsection{x+2.9 非光滑图像处理中的PDE------曲率驱动与面积最小化流}\label{x2.9-ux975eux5149ux6ed1ux56feux50cfux5904ux7406ux4e2dux7684pdeux66f2ux7387ux9a71ux52a8ux4e0eux9762ux79efux6700ux5c0fux5316ux6d41}

\textbf{定理 x+2.9}（平均曲率流与水平集方法）：对于图像分割，活动轮廓模型（active contours）可通过水平集方法求解。设 \(\phi(x,t)\) 为水平集函数，其零水平集 \(\{x \mid \phi(x,t) = 0\}\) 以法向速度 \(v = \kappa g(|\nabla u_0|) + \nu g(|\nabla u_0|)\) 演化（\(\kappa\) 为平均曲率）。该演化的PDE为

\[\frac{\partial\phi}{\partial t} = g(|\nabla u_0|)|\nabla\phi|\left(\nabla\cdot\frac{\nabla\phi}{|\nabla\phi|} + \nu\right).\]

此即测地活动轮廓（geodesic active contour）的Caselles-Kimmel-Sapiro方程。其平衡态对应于极小曲面问题的加权版本。

\textbf{证明}：将曲线演化的Euler-Lagrange方程与水平集表示结合。初始曲线的能量为 \(\oint_C g(|\nabla u_0|)ds\)。第一变分给出梯度下降流：

\[\frac{\partial C}{\partial t} = (g\kappa - \langle\nabla g, N\rangle)N,\]

其中 \(N\) 为单位法向量。用水平集函数 \(\phi\) 表示曲线，\(N = \nabla\phi/|\nabla\phi|\)，\(\kappa = \nabla\cdot(\nabla\phi/|\nabla\phi|)\)。将曲线演化转换为水平集函数的演化即得上述PDE。\(\blacksquare\)

\subsection{x+2.10 Shearlet变换与各向异性特征提取}\label{x2.10-shearletux53d8ux6362ux4e0eux5404ux5411ux5f02ux6027ux7279ux5f81ux63d0ux53d6}

小波变换虽然能捕获点状奇异性，但对图像中的曲线奇异性（边缘）并非最优。Shearlet变换是对小波的各向异性推广，旨在最优稀疏地表示具有各向异性特征（如 \(C^2\) 曲线）的图像。

\textbf{定理 x+2.10}（Shearlet的最优稀疏表示）：设 \(f\) 为分段 \(C^2\) 函数，不连续集为有限长度的 \(C^2\) 曲线。则 \(f\) 的 Shearlet 系数 \(\mathcal{SH}_\psi f(a, s, t)\) 满足

\[\|f - f_N^{\text{Shearlet}}\|_2^2 \leq C N^{-2} (\log N)^3,\]

而小波只能达到 \(O(N^{-1})\) 的衰减速率。此即 Shearlet 对各向异性奇异性的''几乎最优''稀疏逼近性质。

\textbf{证明}：Shearlet的构造基于仿射系统 \(\psi_{a,s,t}(x) = a^{-3/4}\psi(M_{a,s}^{-1}(x-t))\)，其中 \(M_{a,s} = \begin{pmatrix} a & \sqrt{a}s \\ 0 & \sqrt{a} \end{pmatrix}\) 为各向异性膨胀矩阵（抛物尺度律）。Shearlet的频域支撑为楔形对（wedge pairs），这使得其沿方向自动对齐图像的边缘。\(N\) 项非线性逼近逼近率的分析基于边缘附近系数幅值的精细估计，将边缘两侧的Shearlet系数按方向和尺度分类，利用抛物尺度律获得 \(N^{-2}\) 的超收敛性。\(\blacksquare\)

\subsection{x+2.11 图像去模糊与反卷积问题的正则化}\label{x2.11-ux56feux50cfux53bbux6a21ux7ccaux4e0eux53cdux5377ux79efux95eeux9898ux7684ux6b63ux5219ux5316}

图像去模糊（deblurring）是典型的线性反问题。设模糊核为 \(k\)（点扩散函数PSF），观测图像为 \(f = k * u + \eta\)。该问题的离散形式 \(f = Au + \eta\) 中，\(A\) 通常为病态矩阵（条件数极大）。正则化方法是处理此问题的标准途径。

\textbf{定理 x+2.11}（Richardson-Lucy去卷积与EM算法）：当噪声为Poisson类型（光子计数噪声）时，极大似然估计等价于极小化Kullback-Leibler散度。Richardson-Lucy迭代

\[u_{n+1} = u_n \cdot \left(k * \frac{f}{k * u_n}\right)\]

是EM（期望最大化）算法的实现，在非负约束下收敛到解的似然极大值。

\textbf{证明}：考虑Poisson似然 \(L(u) = \sum_{x} [f(x)\log((k*u)(x)) - (k*u)(x)]\)。EM算法通过引入隐变量（光子来源）将不完全数据的似然极大化转化为迭代求解。设隐变量 \(z_{x,y}\) 表示观测于 \(x\) 的光子来源于源位置 \(y\) 的期望比例，由条件期望 \(z_{x,y}^{(n)} = \frac{u_n(y)k(x-y)}{(k*u_n)(x)}\)。\(M\) 步更新 \(u_{n+1}(y) = \sum_x f(x)z_{x,y}^{(n)} / \sum_x k(x-y)\)（归一化分母恰为点扩散函数的积分，通常为1，此时即为R-L公式）。该迭代保持非负性，且每一步增大似然函数。\(\blacksquare\)

\subsection{x+2.12 图像相似性与结构张量}\label{x2.12-ux56feux50cfux76f8ux4f3cux6027ux4e0eux7ed3ux6784ux5f20ux91cf}

图像的结构张量（structure tensor）是分析图像局部几何结构和运动估计的基本工具。对于图像 \(u: \Omega \to \mathbb{R}\)，结构张量定义为

\[J_\rho(u) = G_\rho * (\nabla u_\sigma \otimes \nabla u_\sigma),\]

其中 \(u_\sigma = G_\sigma * u\) 为Gaussian平滑后的图像。\(J_\rho\) 的特征值 \(\lambda_1 \geq \lambda_2 \geq 0\) 和特征向量编码了局部几何信息：\(\lambda_1 \approx \lambda_2 \approx 0\) 对应平坦区域，\(\lambda_1 \gg \lambda_2 \approx 0\) 对应边缘，\(\lambda_1 \geq \lambda_2 \gg 0\) 对应角点。

\textbf{定理 x+2.12}（Lucas-Kanade光流估计）：在两帧图像 \(I_1, I_2\) 之间，光流场 \(v\) 在局部区域 \(\mathcal{N}\) 上满足超定线性系统

\[\left(\sum_{x \in \mathcal{N}} \nabla I(x) \nabla I(x)^T\right) v = -\sum_{x \in \mathcal{N}} \partial_t I(x) \nabla I(x),\]

其中左侧矩阵恰为结构张量在局部区域上的积分。该方程可解当且仅当结构张量非退化（即 \(\lambda_2 > 0\)，对应角点区域），此时最小二乘解为

\[v = -\left(\sum \nabla I \nabla I^T\right)^{-1} \sum \partial_t I \,\nabla I.\]

\textbf{证明}：亮度恒常假设给出 \(I(x+v, t+1) \approx I(x,t)\)，Taylor展开至一阶得 \(\nabla I \cdot v + \partial_t I = 0\)。在局部邻域 \(\mathcal{N}\) 内，对此式施加最小二乘极小化 \(\min_v \sum_{\mathcal{N}} (\nabla I \cdot v + \partial_t I)^2\)。令梯度为零得正规方程，即上述线性系统。结构张量 \(\sum \nabla I \nabla I^T\) 的可逆条件恰要求梯度方向在 \(\mathcal{N}\) 内有足够变化（如角点处），此时光流可靠。平坦区域和边缘处结构张量秩亏，需引入附加正则化（Horn-Schunck光流）。\(\blacksquare\)

\hrulefill

\section{第421章：压缩感知引论}\label{ux7b2c421ux7ae0ux538bux7f29ux611fux77e5ux5f15ux8bba}

压缩感知（Compressed Sensing / Compressive Sensing）是 2000 年代由 Candès、Romberg、Tao 和 Donoho 等人开创的革命性信号处理理论。它表明，稀疏信号可以从远少于 Nyquist 采样率的测量中精确恢复。这一理论的数学基础涉及凸优化、随机矩阵理论和调和分析，其核心洞察是：信号在某些基下的稀疏性可以被充分利用，从而以远低于传统 Nyquist-Shannon 采样定理要求的采样率进行测量和精确重建。压缩感知已经在医学成像（MRI 加速扫描）、雷达成像、地震勘探、压缩成像和单像素相机等领域取得了革命性应用，深刻改变了信号处理领域的基本范式。

\subsection{\texorpdfstring{145.1 稀疏性与 \(\ell_1\) 最小化}{145.1 稀疏性与 \textbackslash ell\_1 最小化}}\label{ux7a00ux758fux6027ux4e0e-ell_1-ux6700ux5c0fux5316}

\textbf{定义 145.1}（稀疏信号）：向量 \(\mathbf{x} \in \mathbb{R}^n\) 是 \(s\)-稀疏的，如果 \(\|\mathbf{x}\|_0 = |\operatorname{supp}(\mathbf{x})| \leq s\)，其中 \(\|\cdot\|_0\) 是 \(\ell_0\) 伪范数（非零分量个数）。更一般地，\(\mathbf{x}\) 是\textbf{可压缩的}（compressible），如果其分量按绝对值降序排列后满足 \(\ell_p\) 逼近误差以 \(s^{-r}\) 的速率衰减。可压缩性是比稀疏性更贴近实际的建模假设------大多数自然信号（如图像、音频）的小波系数或 Fourier 系数并非严格稀疏，而是快速衰减的。

\textbf{定义 145.2}（压缩感知问题）：给定测量矩阵 \(A \in \mathbb{R}^{m \times n}\)（\(m \ll n\)）和测量值 \(\mathbf{y} = A\mathbf{x}\)，从 \(\mathbf{y}\) 中恢复稀疏信号 \(\mathbf{x}\)。测量过程是线性的：\(\mathbf{y} = A\mathbf{x}\)，其中 \(A\) 的每一行对应一次测量（如 MRI 中的一次频率编码梯度）。由于 \(m < n\)，这是一个欠定线性系统，一般有无穷多解。稀疏性约束使得唯一恢复成为可能------这正是压缩感知的反直觉结论。

\textbf{定义 145.3}（\(\ell_0\) 最小化问题）：最自然的稀疏恢复方法是求解 \(\ell_0\) 伪范数最小化： \[\min_{\mathbf{x}} \|\mathbf{x}\|_0 \quad \text{s.t.} \quad A\mathbf{x} = \mathbf{y}\]

然而，\(\ell_0\) 最小化是 NP 难的（组合优化问题），在实际中不可行。压缩感知的核心突破是：在适当的条件下，\(\ell_0\) 最小化可以用 \(\ell_1\) 最小化（凸优化）替代，且两者给出相同的解。

\textbf{定义 145.4}（\(\ell_1\) 最小化 / Basis Pursuit）： \[\min_{\mathbf{x}} \|\mathbf{x}\|_1 \quad \text{s.t.} \quad A\mathbf{x} = \mathbf{y}\]

其中 \(\|\mathbf{x}\|_1 = \sum_{i=1}^{n} |x_i|\)。\(\ell_1\) 范数是 \(\ell_0\) 伪范数的最佳凸松弛（在 \(\ell_p\) 球中，\(p=1\) 是使得单位球为凸集的最小 \(p\)）。\(\ell_1\) 最小化是一个线性规划问题，可以用内点法、ADMM 或 proximal gradient 方法高效求解。

\textbf{定理 145.1}（\(\ell_0\) 与 \(\ell_1\) 的等价性）：若 \(A\) 满足适当的 RIP 条件（如 \(\delta_{2s} < \sqrt{2} - 1\)），则 \(\ell_0\) 最小化问题与 \(\ell_1\) 最小化问题给出相同的解，且该解是唯一的。

\textbf{证明}：若 \(A\) 满足 \(\delta_{2s} < 1\) 的 RIP，则 \(\ell_1\) 最小化的解唯一且等于 \(\ell_0\) 解。核心理由：设 \(\mathbf{x}^*\) 为 \(s\)-稀疏真解，\(\mathbf{x}^\sharp\) 为 \(\ell_1\) 正则化解。令 \(\mathbf{h} = \mathbf{x}^\sharp - \mathbf{x}^*\)，\(A\mathbf{h} = 0\)。将 \(\mathbf{h}\) 的支撑分解为最大 \(s\) 个分量和其他部分。由 RIP，\(\|A\mathbf{h}_S\|_2 \geq (1-\delta_{2s})\|\mathbf{h}_S\|_2\)。结合 \(\ell_1\) 最小性 \(\|\mathbf{x}^\sharp\|_1 \leq \|\mathbf{x}^*\|_1\) 和锥约束，推出 \(\|\mathbf{h}\|_2 = 0\)。\(\blacksquare\)

\textbf{定义 145.5}（Null Space Property, NSP）：矩阵 \(A\) 满足 \(s\) 阶\textbf{零空间性质}（Null Space Property），如果对所有 \(\mathbf{v} \in \ker(A) \setminus \{\mathbf{0}\}\) 和所有指标集 \(T\)（\(|T| \leq s\)），有 \(\|\mathbf{v}_T\|_1 < \|\mathbf{v}_{T^c}\|_1\)。NSP 是 \(\ell_1\) 最小化精确恢复 \(s\)-稀疏向量的充要条件（当 \(A\) 满足 NSP 时）。

\textbf{定理 145.1b}（NSP 的充要性）：\(\ell_1\) 最小化精确恢复所有 \(s\)-稀疏信号 \(\mathbf{x}\) 当且仅当 \(A\) 满足 \(s\) 阶 NSP。

\textbf{证明}：(\(\Rightarrow\)) 设 \(A\) 不满足 NSP，则存在 \(\mathbf{v} \in \ker(A)\) 和 \(|T| \leq s\) 使 \(\|\mathbf{v}_T\|_1 \geq \|\mathbf{v}_{T^c}\|_1\)。令 \(\mathbf{x} = \mathbf{v}_T\)（\(s\)-稀疏），\(\mathbf{x}^\sharp = -\mathbf{v}_{T^c}\)。则 \(A\mathbf{x} = A\mathbf{x}^\sharp\) 且 \(\|\mathbf{x}^\sharp\|_1 = \|\mathbf{v}_{T^c}\|_1 \leq \|\mathbf{v}_T\|_1 = \|\mathbf{x}\|_1\)，故 \(\ell_1\) 最小化解不唯一。 (\(\Leftarrow\)) 设 NSP 成立。若 \(\mathbf{x}^\sharp\) 是 \(\ell_1\) 解，令 \(\mathbf{h} = \mathbf{x}^\sharp - \mathbf{x}\)，\(T\) 为 \(\mathbf{x}\) 的支撑。则 \(\|\mathbf{x}\|_1 \geq \|\mathbf{x}^\sharp\|_1 = \|\mathbf{x} + \mathbf{h}_T + \mathbf{h}_{T^c}\|_1 \geq \|\mathbf{x}\|_1 - \|\mathbf{h}_T\|_1 + \|\mathbf{h}_{T^c}\|_1\)。因此 \(\|\mathbf{h}_{T^c}\|_1 \leq \|\mathbf{h}_T\|_1\)。由 NSP，\(\mathbf{h} = \mathbf{0}\)。\(\blacksquare\)

\subsection{145.2 限制等距性质（RIP）}\label{ux9650ux5236ux7b49ux8dddux6027ux8d28rip}

\textbf{定义 145.6}（RIP）：矩阵 \(A \in \mathbb{R}^{m \times n}\) 满足 \(s\) 阶\textbf{限制等距性质}（Restricted Isometry Property, RIP），如果存在 \(\delta_s \in (0, 1)\) 使得对所有 \(s\)-稀疏向量 \(\mathbf{x} \in \mathbb{R}^n\)，

\[(1 - \delta_s)\|\mathbf{x}\|_2^2 \leq \|A\mathbf{x}\|_2^2 \leq (1 + \delta_s)\|\mathbf{x}\|_2^2\]

最小常数 \(\delta_s\) 称为 \(s\) 阶限制等距常数（RIC）。RIP 要求 \(A\) 在任意 \(s\) 个列向量张成的子空间上近似保持 \(\ell_2\) 范数------即 \(A\) 的任意 \(s\) 列近似正交。RIP 可以被视为 Johnson-Lindenstrauss 引理在稀疏信号上的推广：Johnson-Lindenstrauss 引理保证任意 \(n\) 个点的 \(\ell_2\) 距离可以通过随机投影近似保持，而 RIP 保证任意 \(s\)-稀疏向量的 \(\ell_2\) 范数通过 \(A\) 近似保持。

\textbf{定理 145.2}（Candes-Romberg-Tao，2006）：若 \(A\) 满足 \(\delta_{2s} < \sqrt{2} - 1\)（或更优的 \(\delta_{2s} < 1/\sqrt{2}\)），则 \(\ell_1\) 最小化可精确恢复任意 \(s\)-稀疏向量。

\textbf{证明}：设 \(\mathbf{x}^*\) 为 \(s\)-稀疏真解，\(\mathbf{x}^\sharp\) 为 \(\ell_1\) 解。令 \(\mathbf{h} = \mathbf{x}^\sharp - \mathbf{x}^*\)，\(T_0\) 为 \(\mathbf{x}^*\) 的支撑，\(T_1\) 为 \(\mathbf{h}_{T_0^c}\) 的最大 \(s\) 个分量位置，依此继续。由 Null Space Property + RIP，\(\|\mathbf{h}_{T_0}\|_2 \leq \frac{\sqrt{2}\delta_{2s}}{1-\delta_{2s}} \|\mathbf{h}\|_2\)。若 \(\delta_{2s} < \sqrt{2}-1\)，则 \(\frac{\sqrt{2}\delta_{2s}}{1-\delta_{2s}} < 1\)，强制 \(\mathbf{h} = 0\)。\(\blacksquare\)

\textbf{定理 145.3}（随机矩阵的 RIP）：高斯随机矩阵（\(A_{ij} \sim \mathcal{N}(0, 1/m)\)）或 Bernoulli 随机矩阵以高概率满足 RIP，当 \(m \geq C \cdot s \log(n/s)\) 时。

\textbf{证明}：考虑 \(\|A\mathbf{x}\|_2^2 - \|\mathbf{x}\|_2^2\) 对所有 \(s\)-稀疏单位向量 \(\mathbf{x}\) 的最大偏差。对任意固定的 \(s\)-稀疏支撑集（共 \(\binom{n}{s}\) 个），\(\|A\mathbf{x}\|_2^2\) 是 \(\chi^2\) 型随机变量。利用矩阵 Chernoff 界或 Gordon 比较定理，\(\mathbb{P}(\sup_{\|\mathbf{x}\|_2=1,\|\mathbf{x}\|_0 \leq s} |\|A\mathbf{x}\|_2^2 - 1| > \delta_s) \leq 2\binom{n}{s} e^{-c m \delta_s^2}\)。取 \(m \geq C s \log(n/s)\)，联合界给出高概率 RIP。\(\blacksquare\)

\textbf{定理 145.3b}（部分 Fourier 矩阵的 RIP）：部分 Fourier 矩阵（从 \(n \times n\) DFT 矩阵中随机选取 \(m\) 行）以高概率满足 RIP，当 \(m \geq C s \log^4 n\) 时。这一结果对 MRI 压缩感知应用至关重要，因为 MRI 的测量在 Fourier 域（k-space）中进行。

\textbf{证明}：与高斯矩阵不同，部分 Fourier 矩阵的行不是独立的，需要使用更精细的工具（如 Rudelson-Vershynin 方法或非交换 Khintchine 不等式）。核心估计涉及随机 Fourier 子矩阵的最大奇异值，利用 Dudley 积分在低 M* 范数结构上的界。\(\blacksquare\)

\textbf{定义 145.7}（相干性）：矩阵 \(A\) 的\textbf{相干性}（coherence）\(\mu(A)\) 定义为 \(A\) 的不同列向量之间内积绝对值的最大值： \[\mu(A) = \max_{1 \leq i < j \leq n} \frac{|\langle \mathbf{a}_i, \mathbf{a}_j \rangle|}{\|\mathbf{a}_i\|_2 \|\mathbf{a}_j\|_2}\]

\textbf{定理 145.4}（Welch 界）：对任意 \(m \times n\) 矩阵 \(A\)（\(n \geq m\)，列范数为 1），\(\mu(A) \geq \sqrt{\frac{n-m}{m(n-1)}}\)。当 \(n \gg m\) 时，\(\mu(A) \geq 1/\sqrt{m}\)（近似），这给出了相干性的信息论下界。

\textbf{定理 145.5}（基于相干性的恢复保证）：若 \(\mu = \mu(A)\) 且 \(\mathbf{x}\) 是 \(s\)-稀疏的，\(s < \frac{1}{2}(1 + 1/\mu)\)，则 \(\ell_1\) 最小化（或正交匹配追踪）可精确恢复 \(\mathbf{x}\)。相干性条件比 RIP 条件更强但更易验证。

\textbf{证明}：由 Gershgorin 圆盘定理，\(A_T^*A_T\)（\(T\) 为支撑集）的最小特征值 \(\geq 1 - (s-1)\mu\)。当 \(s < 1/2(1+1/\mu)\) 时，\(1 - (s-1)\mu > 0\)，RIP 常数 \(\delta_s \leq (s-1)\mu < 1\)。\(\blacksquare\)

\subsection{145.3 恢复算法}\label{ux6062ux590dux7b97ux6cd5}

\textbf{算法 145.1}（基础追踪 / Basis Pursuit）：\(\min \|\mathbf{x}\|_1 \text{ s.t. } A\mathbf{x} = \mathbf{y}\)。可转化为线性规划求解。对于含噪情形（\(\mathbf{y} = A\mathbf{x} + \mathbf{e}\)，\(\|\mathbf{e}\|_2 \leq \varepsilon\)），对应的松弛形式为\textbf{基追踪去噪}（Basis Pursuit Denoising, BPDN）： \[\min \|\mathbf{x}\|_1 \quad \text{s.t.} \quad \|A\mathbf{x} - \mathbf{y}\|_2 \leq \varepsilon\]

等价于 LASSO 形式：\(\min \frac{1}{2}\|A\mathbf{x} - \mathbf{y}\|_2^2 + \lambda \|\mathbf{x}\|_1\)（\(\lambda > 0\) 为正则化参数）。

\textbf{定理 145.6}（BPDN 的恢复保证）：若 \(A\) 满足 \(\delta_{2s} < \sqrt{2} - 1\) 且 \(\|\mathbf{e}\|_2 \leq \varepsilon\)，则 BPDN 的解 \(\hat{\mathbf{x}}\) 满足 \(\|\hat{\mathbf{x}} - \mathbf{x}\|_2 \leq C_0 \varepsilon + C_1 \frac{\|\mathbf{x} - \mathbf{x}_s\|_1}{\sqrt{s}}\)，其中 \(\mathbf{x}_s\) 是 \(\mathbf{x}\) 的最佳 \(s\)-稀疏逼近。这一估计表明压缩感知恢复是稳定的（对噪声）和鲁棒的（对非严格稀疏）。

\textbf{算法 145.2}（正交匹配追踪 / OMP）：贪心迭代选择与当前残差最近相关的列，每次添加一个非零分量。OMP 的每步迭代包括：(1) 计算当前残差与 \(A\) 各列的相关性；(2) 选择相关性最大的列加入支撑集；(3) 在当前支撑集上求解最小二乘问题更新稀疏逼近；(4) 计算新的残差。OMP 在 \(s\) 步后停止（\(s\) 为稀疏度）。OMP 的计算复杂度为 \(O(s m n)\)，远低于 Basis Pursuit 的 \(O(n^3)\)（内点法），但恢复保证较弱。

\textbf{定理 145.7}（OMP 的恢复保证）：若 \(\mu(A) < \frac{1}{2s-1}\)，则 OMP 在 \(s\) 步内精确恢复任意 \(s\)-稀疏向量。

\textbf{证明}：OMP 在每一步选择与残差相关性最大的原子。在无噪声情形下，若 \(\mu < 1/(2s-1)\)，则 \(A\) 的任意 \(s\) 列张成的子空间上的最小奇异值严格为正，且 OMP 在每一步选择的原子必然在真实支撑集中。由归纳法，\(s\) 步后支撑集完整恢复。\(\blacksquare\)

\textbf{算法 145.3}（迭代硬阈值 / IHT）：\(\mathbf{x}^{k+1} = H_s(\mathbf{x}^k + A^T(\mathbf{y} - A\mathbf{x}^k))\)，其中 \(H_s\) 保留绝对值最大的 \(s\) 个分量。IHT 的迭代是梯度下降（步长为 1）与硬阈值算子的复合。若 \(\|A\|_2 < 1\)（步长充分小），则 IHT 收敛到局部最优。

\textbf{算法 145.4}（CoSaMP / 压缩采样匹配追踪）：CoSaMP（Compressive Sampling Matching Pursuit, Needell-Tropp 2009）在每步迭代中选取 \(2s\) 个最相关的原子，然后通过最小二乘修剪回 \(s\) 个。CoSaMP 在 RIP 条件下具有线性收敛速率，且提供与 Basis Pursuit 同阶的误差界。

\textbf{算法 145.5}（消息传递算法 / AMP）：近似消息传递（Approximate Message Passing, Donoho-Maleki-Montanari 2009）基于统计物理中的信念传播（belief propagation）思想，通过迭代软阈值和 Onsager 校正项实现快速恢复。AMP 在高维极限下（\(n \to \infty\)，\(m/n \to \delta\)）的动力学由状态演化（state evolution）方程精确描述，可刻画相变边界。

\hrulefill

\section{第422章：密码学数学基础}\label{ux7b2c422ux7ae0ux5bc6ux7801ux5b66ux6570ux5b66ux57faux7840}

密码学是安全通信的科学，其数学基础植根于数论、代数、组合数学和计算复杂性理论。现代密码学已经从 Shannon 的信息论安全（完美保密）发展到计算安全（基于计算困难性假设），其核心构造------公钥密码、数字签名、哈希函数------依赖于数论中的困难问题（因子分解、离散对数）和代数结构（有限域、椭圆曲线）。本章从传统密码的 Shannon 理论开始，重点介绍现代公钥密码体制的数学结构，最后讨论量子计算对密码学的挑战和抗量子密码（PQC）的数学基础。

\subsection{146.1 对称密码与 Shannon 理论}\label{ux5bf9ux79f0ux5bc6ux7801ux4e0e-shannon-ux7406ux8bba}

\textbf{定义 146.1}（密码体制的形式化定义）：一个（对称）加密方案由三元组 \((\text{Gen}, \text{Enc}, \text{Dec})\) 组成： - \(\text{Gen}\)：密钥生成算法，输出密钥 \(k \in \mathcal{K}\)（密钥空间） - \(\text{Enc}_k(m)\)：加密算法，输入明文 \(m \in \mathcal{M}\)（明文空间），输出密文 \(c \in \mathcal{C}\)（密文空间） - \(\text{Dec}_k(c)\)：解密算法，输入密文 \(c\)，输出明文 \(m\)，满足 \(\text{Dec}_k(\text{Enc}_k(m)) = m\)

\textbf{定义 146.2}（Shannon 完美保密）：加密方案 \((\text{Gen}, \text{Enc}, \text{Dec})\) 是\textbf{完美保密}的（perfectly secret），如果对任意明文分布、任意明文 \(m \in \mathcal{M}\) 和密文 \(c \in \mathcal{C}\)（\(\mathbb{P}(C=c) > 0\)）， \[\mathbb{P}(M = m \mid C = c) = \mathbb{P}(M = m)\]

等价地，\(I(M; C) = 0\)（明文和密文统计独立），即观察密文不会泄露关于明文的任何信息（无论攻击者的计算能力如何）。

\textbf{定理 146.1}（Shannon 定理）：完美保密方案必须满足 \(|\mathcal{K}| \geq |\mathcal{M}|\)（密钥空间不小于明文空间）。一次一密（One-Time Pad / Vernam 密码）使用 \(c = m \oplus k\)（\(k\) 为真随机密钥，长度等于明文）达到完美保密。

\textbf{证明}：设 \(\mathbb{P}(C=c) > 0\)（否则密文不可能出现）。由完美保密定义，\(\mathbb{P}(M=m\|C=c) = \mathbb{P}(M=m)\)。由 Bayes 定理，\(\mathbb{P}(C=c\|M=m) = \mathbb{P}(C=c)\)。对于固定 \(m\)，密文 \(c \mapsto \mathbb{P}(C=c\|M=m)\) 作为一个分布，其支撑大小为加密函数 \(\text{Enc}(m, \cdot)\) 的值域大小。对任意可行 \(m \in \mathcal{M}\)，需足够多密钥使每个密文均可产生。对于一次一密，\(c = m \oplus k\) 中 \(k\) 均匀随机，\(\mathbb{P}(C=c\|M=m) = 1/|\mathcal{K}|\) 与 \(m\) 无关，满足完美保密。\(\blacksquare\)

\textbf{定义 146.3}（计算安全）：现代密码学使用计算安全（computational security）代替信息论安全。方案在多项式时间攻击者面前是安全的，假设某些计算问题（如大整数因子分解、离散对数）是困难的。计算安全的具体形式包括：(1) IND-CPA（选择明文攻击下的不可区分性）；(2) IND-CCA（选择密文攻击下的不可区分性）。IND-CPA 安全意味着攻击者无法区分两个选定明文的加密，即使其可以获取任意明文的加密。

\textbf{定义 146.4}（伪随机性）：伪随机生成器（PRG）\(G: \{0,1\}^\lambda \to \{0,1\}^{\ell(\lambda)}\)（\(\ell(\lambda) > \lambda\)）的输出在计算上与真随机不可区分。伪随机函数（PRF）\(F: \{0,1\}^\lambda \times \{0,1\}^n \to \{0,1\}^m\) 在随机密钥下与真随机函数在计算上不可区分。PRG 和 PRF 的存在性等价于单向函数的存在性（Hastad-Impagliazzo-Levin-Luby 定理），从而将对称密码的安全性归结为计算复杂性假设。

\subsection{146.2 RSA 公钥密码}\label{rsa-ux516cux94a5ux5bc6ux7801}

\textbf{定义 146.5}（大整数因子分解问题）：给定两个大素数 \(p, q\) 的乘积 \(N = pq\)，在没有 \(p, q\) 的知识下恢复 \(p\) 和 \(q\)。目前最好的经典算法是数域筛法（Number Field Sieve），其启发式复杂度为 \(\exp(O((\log N)^{1/3}(\log \log N)^{2/3}))\)------亚指数时间，但非多项式时间。

\textbf{定义 146.6}（RSA 问题）：给定 \(N = pq\)（\(p, q\) 为大素数）和 \(e\)（\(\gcd(e, \phi(N)) = 1\)），从 \(c = m^e \bmod N\) 中恢复 \(m\)。RSA 问题不超过因子分解问题的难度，但两者是否等价仍未解决。

\textbf{算法 146.1}（RSA 密钥生成）： 1. 随机生成两个大素数 \(p, q\)（通常各 1024-2048 比特），计算 \(N = pq\) 2. 计算 \(\phi(N) = (p-1)(q-1)\) 3. 选择 \(e\) 满足 \(\gcd(e, \phi(N)) = 1\)（常用 \(e = 65537 = 2^{16}+1\)，Fermat 素数） 4. 计算 \(d = e^{-1} \bmod \phi(N)\)（扩展 Euclid 算法） 5. 公钥：\((N, e)\)；私钥：\((N, d)\)（或 \((p, q, d)\) 以加速解密）

\textbf{算法 146.2}（RSA 加密/解密）： - 加密：\(c = m^e \bmod N\)（\(m \in \mathbb{Z}_N^*\)） - 解密：\(m = c^d \bmod N\)

\textbf{定理 146.2}（RSA 的正确性）：由 Euler 定理，\(m^{ed} \equiv m^{1 + k\phi(N)} \equiv m \pmod{N}\)（对任意 \(m \in \mathbb{Z}_N\) 成立）。

\textbf{证明}：\(ed \equiv 1 \pmod{\phi(N)}\)，故 \(ed = 1 + k\phi(N) = 1 + k(p-1)(q-1)\)。若 \(\gcd(m,N)=1\)，由 Euler 定理 \(m^{\phi(N)} \equiv 1 \pmod{N}\)，故 \(m^{ed} = m \cdot (m^{\phi(N)})^k \equiv m \pmod{N}\)。若 \(\gcd(m,N) \neq 1\)，因 \(N=pq\) 不含平方因子，由中国剩余定理，只需验证模 \(p\) 和模 \(q\)。设 \(p \mid m\)，则 \(m^{ed} \equiv 0 \equiv m \pmod{p}\)。模 \(q\)：\(\gcd(m,q)=1\)，\(m^{ed} \equiv m \cdot (m^{q-1})^{k(p-1)} \equiv m \pmod{q}\)（由 Fermat 小定理）。综上 \(m^{ed} \equiv m \pmod{N}\)。\(\blacksquare\)

\textbf{定理 146.3}（RSA 的安全性分析）： 1. 若攻击者可分解 \(N\)，则可计算 \(\phi(N)\) 和 \(d\)，从而破解 RSA。因此 RSA 的安全性不超过因子分解的难度。 2. 若攻击者已知 \(\phi(N)\)，则可分解 \(N\)（解 \(p+q = N - \phi(N) + 1\) 和 \(pq = N\)）。 3. 若 \(d < N^{1/4}/3\)，则存在基于连分数的 Wiener 攻击在多项式时间内恢复 \(d\)。 4. 若同一明文 \(m\) 用 \(e\) 个不同的公钥 \((N_i, e)\) 加密（广播攻击），且 \(e\) 较小，可恢复 \(m\)（Hastad 广播攻击）。

\textbf{算法 146.3}（素性检测 / Miller-Rabin）：随机化算法，可在多项式时间内以高概率判定素性。复杂度 \(O(k \log^3 n)\)，错误概率 \(< 4^{-k}\)。Miller-Rabin 基于 Fermat 小定理的逆否命题：若 \(n\) 是素数，则对任意 \(a\)（\(\gcd(a,n)=1\)），\(a^{n-1} \equiv 1 \pmod{n}\) 且 \(a^{(n-1)/2^t}\) 在平方过程中只能出现 \(1\) 或 \(-1\)。若找到 \(a\) 使这些条件不成立，则 \(n\) 必为合数。

\textbf{定理 146.4}（AKS 素性检测，2002）：存在确定性的多项式时间素性检测算法（Agrawal-Kayal-Saxena）。AKS 算法基于恒等式 \((x+a)^n \equiv x^n + a \pmod{n}\) 对所有 \(a\) 成立当且仅当 \(n\) 为素数（或素数幂）的推广。AKS 的时间复杂度为 \(O(\log^{12} n)\)（后优化至 \(O(\log^{6} n)\) 或更好），提供了素数判定的第一个确定性多项式时间算法。

\textbf{证明}：\((x+a)^n = \sum_{k=0}^{n} \binom{n}{k} x^k a^{n-k}\)。若 \(n\) 为素数，则 \(\binom{n}{k} \equiv 0 \pmod{n}\) 对所有 \(1 \leq k \leq n-1\) 成立，故 \((x+a)^n \equiv x^n + a^n \equiv x^n + a \pmod{n}\)（由 Fermat 小定理）。若 \(n\) 为合数且不是素数幂，则存在 \(k\) 使 \(\binom{n}{k} \not\equiv 0 \pmod{n}\)。AKS 算法通过检验此恒等式对 \(r = O(\log^6 n)\) 范围内的 \(a\) 是否成立来判定素性。\(\blacksquare\)

\subsection{146.3 离散对数与 Diffie-Hellman}\label{ux79bbux6563ux5bf9ux6570ux4e0e-diffie-hellman}

\textbf{定义 146.7}（离散对数问题 / DLP）：在循环群 \(G = \langle g \rangle\) 中，给定 \(g^a\)，求 \(a \in \mathbb{Z}_{|G|}\)。在适当的群 \(G\) 中，DLP 被认为是困难的。经典的 DLP 在 \(\mathbb{F}_p^*\) 上可用指数演算法（index calculus）在亚指数时间内求解，而在一般椭圆曲线群上，目前最好的算法是 Pollard rho 算法（\(O(\sqrt{|G|})\)），因此椭圆曲线提供更高效率的安全性。

\textbf{算法 146.4}（Diffie-Hellman 密钥交换，1976）：Alice 和 Bob 在不安全信道上协商共享密钥。 1. 公开参数：素数 \(p\)，生成元 \(g\)（或椭圆曲线基点 \(G\)） 2. Alice 选择 \(a\)，发送 \(A = g^a \bmod p\) 3. Bob 选择 \(b\)，发送 \(B = g^b \bmod p\) 4. 共享密钥：\(K = B^a \equiv A^b \equiv g^{ab} \pmod{p}\)

安全性基于\textbf{计算 Diffie-Hellman 假设}（CDH）：从 \(g^a, g^b\) 计算 \(g^{ab}\) 是困难的。更强的假设是\textbf{判定 Diffie-Hellman 假设}（DDH）：\((g^a, g^b, g^{ab})\) 与 \((g^a, g^b, g^c)\)（\(c\) 随机）在计算上不可区分。

\textbf{定义 146.8}（ElGamal 加密）：基于 Diffie-Hellman 的公钥加密方案。密钥生成：选择私钥 \(x\)，公钥 \(h = g^x\)。加密：选择随机 \(r\)，密文为 \((c_1, c_2) = (g^r, m \cdot h^r)\)。解密：\(m = c_2 \cdot c_1^{-x}\)。ElGamal 加密在 DDH 假设下是 IND-CPA 安全的。

\textbf{证明（ElGamal 安全性）}：假设存在攻击者 \(\mathcal{A}\) 区分 ElGamal 加密。构造 DDH 区分器：给定 \((g^a, g^b, T)\)，设 \(h = g^a\)，\(c_1 = g^b\)，\(c_2 = m \cdot T\)。若 \(T = g^{ab}\)，则 \((c_1, c_2)\) 是 \(m\) 的正确 ElGamal 加密；若 \(T\) 随机，则 \(c_2\) 是随机元素，不包含 \(m\) 的信息。由 \(\mathcal{A}\) 的区分能力，可区分 \(T\) 是否 \(= g^{ab}\)，从而破解 DDH。\(\blacksquare\)

\textbf{定义 146.9}（椭圆曲线密码 / ECC）：使用椭圆曲线上的点群代替 \(\mathbb{F}_p^*\) 作为 DLP 的平台。优点是密钥更短、安全性更高（亚指数攻击不适用）。比特币和以太坊使用 secp256k1 曲线。椭圆曲线上的群运算由弦切法定义（见椭圆曲线章节），DLP 在此群上的难度使得 ECC 仅需 256 比特密钥即可提供与 3072 比特 RSA 同等的安全性。

\subsection{146.4 数字签名与哈希函数}\label{ux6570ux5b57ux7b7eux540dux4e0eux54c8ux5e0cux51fdux6570}

\textbf{定义 146.10}（数字签名方案）：由 \((Gen, Sign, Verify)\) 组成。签名者用私钥 \(sk\) 对消息 \(m\) 生成签名 \(\sigma = \text{Sign}_{sk}(m)\)，验证者用公钥 \(pk\) 验证 \(\text{Verify}_{pk}(m, \sigma) \in \{\text{accept}, \text{reject}\}\)。安全性要求：在选择消息攻击下，攻击者无法伪造签名（EUF-CMA / Existential Unforgeability under Chosen Message Attack）。

\textbf{算法 146.5}（RSA 签名 / RSA-FDH）：\(\sigma = H(m)^d \bmod N\)（\(H\) 为密码学哈希函数，FDH = Full Domain Hash）。验证：\(H(m) \stackrel{?}{=} \sigma^e \bmod N\)。在随机预言机模型下，若 RSA 问题是困难的，RSA-FDH 是 EUF-CMA 安全的。

\textbf{算法 146.6}（DSA / ECDSA 签名）：基于离散对数的签名方案。ECDSA 是以太坊和比特币使用的签名算法。ECDSA 签名过程：选择随机数 \(k\)，计算 \(R = kG\)，\(r = R_x \bmod n\)，\(s = k^{-1}(H(m) + r \cdot d) \bmod n\)，签名为 \((r, s)\)。验证：\(u_1 = H(m) \cdot s^{-1} \bmod n\)，\(u_2 = r \cdot s^{-1} \bmod n\)，\(R' = u_1 G + u_2 Q\)（\(Q\) 为公钥），验证 \(R'_x \equiv r \pmod{n}\)。

\textbf{算法 146.7}（Schnorr 签名）：基于离散对数的最简签名方案，提供可证明安全性（在随机预言机模型下，基于 DLP 的困难性）。Schnorr 签名是 EdDSA（Ed25519）的基础，后者的确定性 nonce 生成避免了 ECDSA 中随机数重用的灾难性后果。

\textbf{定义 146.11}（密码学哈希函数）：\(H: \{0,1\}^* \to \{0,1\}^n\) 满足： - \textbf{抗原像性}（preimage resistance）：给定 \(y\)，找到 \(x\) 使得 \(H(x) = y\) 是困难的（计算复杂度 \(\Omega(2^n)\)） - \textbf{抗第二原像性}（second preimage resistance）：给定 \(x\)，找到 \(x' \neq x\) 使得 \(H(x') = H(x)\) 是困难的（\(\Omega(2^n)\)） - \textbf{抗碰撞性}（collision resistance）：找到 \(x \neq x'\) 使得 \(H(x) = H(x')\) 是困难的（\(\Omega(2^{n/2})\)）

\textbf{定理 146.5}（生日攻击）：找到碰撞的期望尝试次数为 \(O(2^{n/2})\)（生日悖论）。因此 SHA-256（\(n=256\)）提供 128 比特的碰撞安全性。

\textbf{证明}：从 \(N\) 个值中独立均匀抽取 \(k\) 个样本，无碰撞概率为 \(P = \prod_{i=1}^{k-1} (1 - i/N)\)。利用 \(1-x \approx e^{-x}\)，\(\ln P \approx -\sum_{i=1}^{k-1} i/N = -k(k-1)/(2N) \approx -k^2/(2N)\)。当 \(k \approx \sqrt{2N \ln 2} = \Theta(\sqrt{N})\) 时，\(P \approx 1/2\)。取 \(N=2^n\)，碰撞需 \(k = \Theta(2^{n/2})\) 次尝试。\(\blacksquare\)

\textbf{定理 146.6}（Merkle-Damgard 构造）：若压缩函数 \(f: \{0,1\}^n \times \{0,1\}^\ell \to \{0,1\}^n\) 是抗碰撞的，则通过 Merkle-Damgard 迭代构造的哈希函数 \(H\) 也是抗碰撞的。SHA-256、SHA-1 和 MD5 均基于 Merkle-Damgard 构造。MD5 和 SHA-1 的碰撞已被找到（利用差分密码分析），因此不再安全。

\subsection{146.5 量子密码与后量子密码}\label{ux91cfux5b50ux5bc6ux7801ux4e0eux540eux91cfux5b50ux5bc6ux7801}

\textbf{定义 146.12}（Shor 算法，1994）：量子计算机可在多项式时间内分解大整数和求解离散对数，从而攻破 RSA 和 ECC。Shor 算法利用量子 Fourier 变换（QFT）将因子分解归约为周期查找（period finding），而周期查找在量子计算机上可在多项式时间内完成。具体地，对 \(N\) 的因子分解：随机选 \(a\)（\(\gcd(a,N)=1\)），求 \(f(x) = a^x \bmod N\) 的周期 \(r\)。若 \(r\) 为偶数且 \(a^{r/2} \not\equiv -1 \pmod{N}\)，则 \(\gcd(a^{r/2} \pm 1, N)\) 是 \(N\) 的非平凡因子。

\textbf{定理 146.7}（Grover 算法，1996）：量子计算机可在 \(O(\sqrt{N})\) 时间内搜索无序数据库（经典需 \(O(N)\)）。对称密码的密钥长度应加倍（如 AES-256）以抵抗量子攻击。Grover 算法对哈希函数的抗原像性攻击将复杂度从 \(2^n\) 降至 \(2^{n/2}\)，因此抗量子哈希函数应使用至少 384 比特的输出。

\textbf{定理 146.8}（Grover 算法的下界）：任何量子算法在无序搜索问题上的查询复杂度下界为 \(\Omega(\sqrt{N})\)（BBBV 定理，Bennett-Bernstein-Brassard-Vazirani 1997），因此 Grover 算法是渐进最优的。

\textbf{定义 146.13}（后量子密码 / PQC）：基于格（Lattice-based）、编码（Code-based）、多变量（Multivariate）和哈希签名（Hash-based）的密码方案。NIST 在 2024 年标准化了 CRYSTALS-Kyber（格密码，基于 Module-LWE 问题）和 CRYSTALS-Dilithium（格签名）等方案。

\textbf{定义 146.14}（格密码学基础）：格 \(\Lambda\) 是 \(\mathbb{R}^n\) 的离散加法子群，\(\Lambda = \{\sum_{i=1}^{n} z_i \mathbf{b}_i : z_i \in \mathbb{Z}\}\)，其中 \(\{\mathbf{b}_i\}\) 是格基。格密码学基于以下困难问题： - \textbf{最短向量问题（SVP）}：给定格基，求格中的最短非零向量 - \textbf{最近向量问题（CVP）}：给定格基和目标点，求格中离目标点最近的向量 - \textbf{带错误学习（LWE）}：给定 \(A \in \mathbb{Z}_q^{m \times n}\) 和 \(\mathbf{b} = A\mathbf{s} + \mathbf{e} \pmod{q}\)（\(\mathbf{e}\) 为小噪声），恢复秘密 \(\mathbf{s}\)。Regev（2005）证明 LWE 在量子归约下与最坏情况格问题（GapSVP）同样困难。

\textbf{定理 146.9}（LWE 的安全性归约）：若存在量子算法在多项式时间内求解 LWE，则存在量子算法在多项式时间内求解最坏情况下的 GapSVP 和 SIVP 问题（在任意 \(n\) 维格上）。这一归约将格密码的安全性建立在格问题的计算困难性上，而非启发式假设。

\textbf{证明}（框架）：Regev 的归约利用量子 Fourier 变换和离散 Gauss 采样，将 LWE 实例的求解转化为对偶格中短向量的搜索。核心步骤：(1) 从 LWE 样本构造量子态，其 Fourier 变换的支撑集中在格点上；(2) 通过求解 BDD（Bounded Distance Decoding）问题，将 LWE 的求解归约到 GapSVP。\(\blacksquare\)

\textbf{定义 146.15}（零知识证明与 SNARK）：零知识证明允许证明者（prover）向验证者（verifier）证明某个陈述为真而不泄露任何额外信息。SNARK（Succinct Non-interactive Argument of Knowledge）是简洁的非交互式零知识证明，其证明大小和验证时间与计算规模的对数成正比。SNARK 在区块链（如 Zcash 的 zk-SNARKs）和隐私保护计算中有重要应用，其数学基础涉及多项式承诺、椭圆曲线配对、快速 Fourier 变换和代数编码理论。

\hrulefill

\emph{卷六：概率统计与信息论终。}

\newpage

\chapter{12-最优化与控制/01-最优化/01-线性规划.md}\label{ux6700ux4f18ux5316ux4e0eux63a7ux523601-ux6700ux4f18ux531601-ux7ebfux6027ux89c4ux5212.md}

\section{第423章：线性规划与单纯形法}\label{ux7b2c423ux7ae0ux7ebfux6027ux89c4ux5212ux4e0eux5355ux7eafux5f62ux6cd5}

线性规划（Linear Programming, LP）是优化理论中研究最成熟的分支，由 Dantzig 在 1947 年提出单纯形法（Simplex Method）而建立。

\subsection{132.1 线性规划的标准形式}\label{ux7ebfux6027ux89c4ux5212ux7684ux6807ux51c6ux5f62ux5f0f}

\textbf{定义 132.1}（线性规划的标准形式）：

\[\begin{aligned} \min \quad & \mathbf{c}^T \mathbf{x} \\ \text{s.t.} \quad & A\mathbf{x} = \mathbf{b} \\ & \mathbf{x} \geq \mathbf{0} \end{aligned}\]

其中 \(\mathbf{c} \in \mathbb{R}^n\)，\(A \in \mathbb{R}^{m \times n}\)，\(\mathbf{b} \in \mathbb{R}^m\)。\(m \leq n\) 且 \(\operatorname{rank}(A) = m\)。

\textbf{定义 132.2}（可行域）：可行域 \(\mathcal{F} = \{\mathbf{x} \in \mathbb{R}^n : A\mathbf{x} = \mathbf{b}, \mathbf{x} \geq \mathbf{0}\}\) 是\textbf{多面体}（polyhedron）。若 \(\mathcal{F}\) 有界，则称为\textbf{多胞体}（polytope）。

\textbf{定理 132.1}（线性规划的基本定理）：若 LP 有有限最优值，则存在最优解是 \(\mathcal{F}\) 的顶点（基本可行解）。

\textbf{证明}：LP可行域 \(\mathcal{F}\) 为多面体 \(\{x:Ax=b,x\ge 0\}\)。由Minkowski-Weyl定理，\(\mathcal{F}\) 可表示为顶点集与极方向的有限凸组合。目标函数 \(\mathbf{c}^T\mathbf{x}\) 为线性泛函，在紧凸集上达到极值于顶点。若 \(\mathcal{F}\) 有界（多胞体），\(\mathbf{c}^T\mathbf{x}\) 的紧集极小点必为某个顶点。若 \(\mathcal{F}\) 无界，设 \(\mathbf{x}^*\) 为最优解，分解 \(\mathbf{x}^*=\sum_k\lambda_k\mathbf{v}_k+\sum_\ell\mu_\ell\mathbf{d}_\ell\)（\(\mathbf{v}_k\) 为顶点，\(\mathbf{d}_\ell\) 为极方向）。因 \(\mathbf{c}^T\mathbf{d}_\ell\ge 0\)（否则沿极方向目标函数可无限减小），故 \(\min_k\mathbf{c}^T\mathbf{v}_k\le\mathbf{c}^T\mathbf{x}^*\)，等号在某顶点 \(\mathbf{v}_k\) 处成立，从而该顶点也是最优解。\(\blacksquare\)

\textbf{定义 132.3}（基本可行解）：将 \(A\) 分块为 \(A = [B \mid N]\)，其中 \(B \in \mathbb{R}^{m \times m}\) 可逆（基矩阵），\(N\) 为非基列。对应的 \(\mathbf{x} = \begin{pmatrix} \mathbf{x}_B \\ \mathbf{x}_N \end{pmatrix}\) 满足 \(\mathbf{x}_B = B^{-1}\mathbf{b} \geq \mathbf{0}\)，\(\mathbf{x}_N = \mathbf{0}\)，称为\textbf{基本可行解}（BFS）。

\subsection{132.2 单纯形法}\label{ux5355ux7eafux5f62ux6cd5-1}

\textbf{算法 132.1}（单纯形法 / Simplex Method）：

\begin{enumerate}
\def\labelenumi{\arabic{enumi}.}
\tightlist
\item
  \textbf{初始化}：找到一个初始 BFS
\item
  \textbf{定价}（Pricing）：计算简约费用 \(\bar{\mathbf{c}}_N = \mathbf{c}_N - N^T B^{-T} \mathbf{c}_B\)。若 \(\bar{\mathbf{c}}_N \geq \mathbf{0}\)，当前解为最优解，停止
\item
  \textbf{进基}：选择 \(j\) 满足 \(\bar{c}_j < 0\)（通常选择最负的），\(x_j\) 进入基
\item
  \textbf{离基}：计算 \(\mathbf{d} = B^{-1} A_j\)（\(A_j\) 为进基列）。若 \(\mathbf{d} \leq \mathbf{0}\)，问题无界。否则选择 \(\theta = \min_{d_i > 0} \frac{(\mathbf{x}_B)_i}{d_i}\)，令离基变量为达到最小比值的行对应的变量
\item
  \textbf{更新}：更新基矩阵 \(B\) 和 \(\mathbf{x}_B\)，返回步骤 2
\end{enumerate}

\textbf{定理 132.2}（单纯形法的有限终止性）：若 LP 非退化（所有 BFS 中基变量严格为正），则单纯形法在有限步内终止于最优解或判定无界。

\textbf{证明}：每步迭代从当前 BFS 沿边移动至相邻 BFS。非退化假设保证 \(\theta>0\)（最小比值严格正），目标函数严格减小。顶点数有限（\(C_n^m\) 个基本可行解），每次迭代访问新顶点（\(c^T x\) 严格递减，永不重复），故至多 \(C_n^m\) 次迭代。由定价步骤：若 \(\bar{c}_N\ge 0\) 则当前解最优；若所有 \(\bar{c}_j<0\) 且对应 \(d\le 0\) 则问题无界。\(\blacksquare\)

\emph{退化处理}：当出现退化时，可能循环。Bland 法则（最小下标法则）保证有限终止。

\textbf{命题 132.3}（单纯形法的几何解释）：单纯形法沿着可行多面体的边移动，从当前顶点到相邻顶点，每次移动降低目标函数值。

\subsection{132.3 初始可行解的构造}\label{ux521dux59cbux53efux884cux89e3ux7684ux6784ux9020}

\textbf{定义 132.4}（两阶段法）： - \textbf{阶段 I}：引入人工变量，求解辅助 LP \(\min \sum \mathbf{1}^T \mathbf{a}\)，s.t. \(A\mathbf{x} + \mathbf{a} = \mathbf{b}, \mathbf{x} \geq 0, \mathbf{a} \geq 0\)。若最优值为 0，则得到原问题的可行 BFS - \textbf{阶段 II}：从阶段 I 得到的 BFS 出发，求解原问题

\textbf{定义 132.5}（大 M 法）：直接求解 \(\min \mathbf{c}^T \mathbf{x} + M \sum a_i\)，其中 \(M > 0\) 充分大，使得人工变量在最优解中必为 0。

\subsection{132.4 内点法}\label{ux5185ux70b9ux6cd5}

\textbf{定理 132.4}（Karmarkar 算法，1984）：存在求解 LP 的多项式时间算法（内点法），基于将可行域变换并沿梯度方向移动。

\textbf{证明}：Karmarkar 原始对偶内点法将问题变换为 \(\min c^T x\) s.t. \(Ax=0, e^T x=1, x\ge 0\)（\(e=(1,\dots,1)^T\)）。利用投影变换 \(T(x)=D^{-1}x/(e^T D^{-1}x)\)（\(D=\operatorname{diag}(x)\)），将当前内点映射到单纯形中心，沿投影负梯度方向选取步长 \(\alpha r/(n-1)\)，再逆变换回来。势函数 \(\phi(x)=n\ln(c^T x)-\sum\ln x_i\) 在每步减少至少常数 \(\delta>0\)，经 \(O(nL)\) 次迭代即可达到 \(\varepsilon\)-精度最优解，每步 \(O(n^3)\) 运算。这证明了 LP 的多项式时间可解性。\(\blacksquare\)

\textbf{定义 132.6}（障碍函数法）：考虑问题 \(\min \mathbf{c}^T \mathbf{x} - \mu \sum_{i=1}^n \ln x_i\) s.t. \(A\mathbf{x} = \mathbf{b}\)。当 \(\mu \to 0\) 时，最优解趋近于原 LP 的最优解。这是\textbf{原-对偶内点法}的基础。

\hrulefill

\hrulefill

\hrulefill

\hrulefill

\section{第424章：对偶理论与灵敏度分析}\label{ux7b2c424ux7ae0ux5bf9ux5076ux7406ux8bbaux4e0eux7075ux654fux5ea6ux5206ux6790}

对偶理论是线性规划中与每个原问题（Primal）自然关联的另一个线性规划问题（Dual），两者之间有深刻的关系。

\subsection{133.1 对偶问题的构造}\label{ux5bf9ux5076ux95eeux9898ux7684ux6784ux9020}

\textbf{定义 133.1}（对偶问题）：对标准形式 LP（Primal）：

\[\begin{aligned} \min \quad & \mathbf{c}^T \mathbf{x} \\ \text{s.t.} \quad & A\mathbf{x} = \mathbf{b}, \quad \mathbf{x} \geq \mathbf{0} \end{aligned}\]

其对偶问题（Dual）为：

\[\begin{aligned} \max \quad & \mathbf{b}^T \mathbf{y} \\ \text{s.t.} \quad & A^T \mathbf{y} \leq \mathbf{c} \end{aligned}\]

其中 \(\mathbf{y} \in \mathbb{R}^m\) 为对偶变量。

\textbf{定理 133.1}（弱对偶定理）：若 \(\mathbf{x}\) 是原问题可行解，\(\mathbf{y}\) 是对偶问题可行解，则 \(\mathbf{c}^T \mathbf{x} \geq \mathbf{b}^T \mathbf{y}\)。

\emph{证明}：\(\mathbf{c}^T \mathbf{x} \geq (A^T \mathbf{y})^T \mathbf{x} = \mathbf{y}^T A \mathbf{x} = \mathbf{y}^T \mathbf{b} = \mathbf{b}^T \mathbf{y}\)。∎

\textbf{定理 133.2}（强对偶定理）：若原问题有有限最优解 \(\mathbf{x}^*\)，则对偶问题也有有限最优解 \(\mathbf{y}^*\)，且 \(\mathbf{c}^T \mathbf{x}^* = \mathbf{b}^T \mathbf{y}^*\)。

\textbf{证明}：设单纯形法终止于最优BFS，对应基矩阵 \(B\)。最优性条件 \(\bar{\mathbf{c}}_N=\mathbf{c}_N-N^T B^{-T}\mathbf{c}_B\ge\mathbf{0}\)。令 \(\mathbf{y}^*=B^{-T}\mathbf{c}_B\)。检查对偶可行性：\(A^T\mathbf{y}^*=\begin{bmatrix}B^T\\N^T\end{bmatrix}B^{-T}\mathbf{c}_B=\begin{bmatrix}\mathbf{c}_B\\N^T B^{-T}\mathbf{c}_B\end{bmatrix}\le\begin{bmatrix}\mathbf{c}_B\\\mathbf{c}_N\end{bmatrix}=\mathbf{c}\)，故 \(\mathbf{y}^*\) 对偶可行。目标值：\(\mathbf{b}^T\mathbf{y}^*=\mathbf{b}^T B^{-T}\mathbf{c}_B=(B^{-1}\mathbf{b})^T\mathbf{c}_B=\mathbf{x}_B^{*T}\mathbf{c}_B=\mathbf{c}^T\mathbf{x}^*\)。由弱对偶（定理133.1），\(d(\mathbf{y}^*)=\mathbf{b}^T\mathbf{y}^*=\mathbf{c}^T\mathbf{x}^*=p^*\le d^*\)（\(d^*\) 为对偶最优值），结合 \(d^*\le p^*\) 得 \(d^*=p^*\)，故 \(\mathbf{y}^*\) 对偶最优。若LP有有限最优解但单纯形法循环不能终止，Bland法则保证非退化下的有限终止。\(\blacksquare\)

\textbf{定理 133.3}（互补松弛定理）：设 \(\mathbf{x}^*\) 为原问题最优解，\(\mathbf{y}^*\) 为对偶问题最优解，则

\[x_j^* \cdot (\mathbf{c}_j - \mathbf{a}_j^T \mathbf{y}^*) = 0, \quad \forall j\]

即原变量和对偶松弛变量不能同时为正。

\textbf{证明}：由弱对偶，\(\mathbf{c}^T\mathbf{x}^*-\mathbf{b}^T\mathbf{y}^*=\sum_j x_j^*(\mathbf{c}_j-\mathbf{a}_j^T\mathbf{y}^*)\ge 0\)。由强对偶，\(\mathbf{c}^T\mathbf{x}^*=\mathbf{b}^T\mathbf{y}^*\)。故 \(\sum_j x_j^*(\mathbf{c}_j-\mathbf{a}_j^T\mathbf{y}^*)=0\)。因 \(x_j^*\ge 0\) 且 \(\mathbf{c}_j-\mathbf{a}_j^T\mathbf{y}^*\ge 0\)（对偶可行），各项非负，非负项之和为零当且仅当每项为零，即 \(x_j^*(\mathbf{c}_j-\mathbf{a}_j^T\mathbf{y}^*)=0\)（\(\forall j\)）。\(\blacksquare\)

\textbf{推论 133.4}：线性规划的可能情况（原问题/对偶问题）： - 原问题有有限最优解 \(\iff\) 对偶问题有有限最优解 - 原问题无界 \(\implies\) 对偶问题不可行 - 原问题不可行 \(\implies\) 对偶问题无界或不可行

\textbf{证明}：(i) 由强对偶定理（133.2）直接得 \(\Leftrightarrow\)。(ii) 若原问题可行且无界而一对偶可行解 \(\mathbf{y}_0\) 存在，弱对偶给出 \(\mathbf{c}^T\mathbf{x}\ge\mathbf{b}^T\mathbf{y}_0\)（对所有 \(\mathbf{x}\) 可行），与无界矛盾------故对偶不可行。(iii) 若原问题不可行，假设对偶有有限最优解，由强对偶原问题也有最优解，矛盾。从而对偶或者不可行或者无界。\(\blacksquare\)

\subsection{133.2 对偶的经济解释}\label{ux5bf9ux5076ux7684ux7ecfux6d4eux89e3ux91ca}

\textbf{定义 133.2}（影子价格）：对偶变量 \(y_i\) 是第 \(i\) 个约束资源的\textbf{影子价格}（shadow price），即增加一单位资源 \(b_i\) 时最优值的边际改进。

对偶理论的经济解释是线性规划最深远的影响之一。在资源分配问题中：原问题 \(\max \{c^T x : Ax \leq b, x \geq 0\}\) 表示在 \(m\) 种资源约束下选择 \(n\) 种活动水平以最大化利润，对偶问题 \(\min \{b^T y : A^T y \geq c, y \geq 0\}\) 表示为资源定价以最小化总成本，同时确保定价不低于任何活动的单位利润。强对偶定理 \(c^T x^* = b^T y^*\) 是经济学中''零利润条件''的数学表达：在最优配置下，资源的总影子价值恰好等于总利润。互补松弛条件 \(y_i^*(b_i - (Ax^*)_i) = 0\) 和 \((c_j - (A^T y^*)_j) x_j^* = 0\) 的经济学解读是：稀缺资源（约束取等号）有正影子价格，充裕资源（约束严格）的影子价格为零；类似地，仅当活动的单位利润等于其资源成本的影子价格之和时，该活动才被采用。这一框架统一了微观经济学中的边际分析、一般均衡理论和福利经济学基本定理：竞争均衡价格即为适当定义的对偶变量，其存在性本质上等价于线性规划对偶的可行性。

\textbf{定理 133.5}（灵敏度分析）：若基矩阵 \(B\) 的最优性不受扰动影响，则 \(\Delta z^* = \mathbf{y}^{*T} \Delta \mathbf{b}\)。

\textbf{证明}：最优目标值 \(z^*=\mathbf{c}_B^T B^{-1}\mathbf{b}=\mathbf{y}^{*T}\mathbf{b}\)（由 \(\mathbf{y}^*=B^{-T}\mathbf{c}_B\)）。对 \(\mathbf{b}\) 线性依赖，\(\Delta z^*=\mathbf{y}^{*T}\Delta\mathbf{b}\)。若扰动后的基仍可行且最优（\(\bar{\mathbf{c}}_N\) 不变号的范围内），此公式精确成立。对偶变量的经济意义：\(y_i^*=\partial z^*/\partial b_i\) 为第 \(i\) 个资源的边际价值（影子价格）。\(\blacksquare\)

\subsection{133.3 对偶在非 LP 情形的推广}\label{ux5bf9ux5076ux5728ux975e-lp-ux60c5ux5f62ux7684ux63a8ux5e7f}

\textbf{定义 133.3}（Lagrange 对偶）：对一般约束优化问题

\[\min_{\mathbf{x}} f(\mathbf{x}) \quad \text{s.t.} \quad g_i(\mathbf{x}) \leq 0, \; h_j(\mathbf{x}) = 0\]

定义 Lagrange 函数 \(L(\mathbf{x}, \boldsymbol{\lambda}, \boldsymbol{\nu}) = f(\mathbf{x}) + \sum \lambda_i g_i(\mathbf{x}) + \sum \nu_j h_j(\mathbf{x})\)，其中 \(\boldsymbol{\lambda} \geq 0\)。Lagrange 对偶函数为 \(d(\boldsymbol{\lambda}, \boldsymbol{\nu}) = \inf_{\mathbf{x}} L(\mathbf{x}, \boldsymbol{\lambda}, \boldsymbol{\nu})\)。

\textbf{定理 133.6}（弱对偶，一般情形）：对任意 \(\boldsymbol{\lambda} \geq 0\)，\(d(\boldsymbol{\lambda}, \boldsymbol{\nu}) \leq p^*\)（原问题最优值）。对偶间隙为 \(p^* - d^* \geq 0\)。

\textbf{证明}：对任意可行解 \(\mathbf{x}\)（\(g_i(\mathbf{x})\le 0\)，\(h_j(\mathbf{x})=0\)）和 \(\boldsymbol{\lambda}\ge 0\)，\(L(\mathbf{x},\boldsymbol{\lambda},\boldsymbol{\nu})=f(\mathbf{x})+\sum\lambda_i g_i(\mathbf{x})+\sum\nu_j h_j(\mathbf{x})\le f(\mathbf{x})\)（因 \(\lambda_i g_i(\mathbf{x})\le 0\)，\(\nu_j h_j(\mathbf{x})=0\)）。取下确界：\(d(\boldsymbol{\lambda},\boldsymbol{\nu})=\inf_{\mathbf{x}}L(\mathbf{x},\boldsymbol{\lambda},\boldsymbol{\nu})\le L(\tilde{\mathbf{x}},\boldsymbol{\lambda},\boldsymbol{\nu})\le f(\tilde{\mathbf{x}})\)（对任意可行 \(\tilde{\mathbf{x}}\)）。取 \(\tilde{\mathbf{x}}\) 为最优解得 \(d(\boldsymbol{\lambda},\boldsymbol{\nu})\le p^*\)。因对任意对偶可行 \((\boldsymbol{\lambda},\boldsymbol{\nu})\) 成立，对上确界也成立：\(d^*\le p^*\)。\(\blacksquare\)

\hrulefill

\subsection{252.3 网络规划与单纯形法的几何联系}\label{ux7f51ux7edcux89c4ux5212ux4e0eux5355ux7eafux5f62ux6cd5ux7684ux51e0ux4f55ux8054ux7cfb}

\textbf{定理 252.6}（单纯形法的几何解释）：单纯形法从多面体\(P = \{x : Ax = b, x \geq 0\}\)的一个顶点出发，沿下降边移动到相邻顶点，直至达到最优顶点。单纯形表对应于基矩阵\(B\)的逆：基本解为\(x_B = B^{-1} b\)，非基变量为\(0\)。

\textbf{证明}：多面体 \(P\) 的顶点与基本可行解一一对应：顶点 \(\Leftrightarrow\) 线性独立的紧约束集。相邻顶点间相差一个进基操作------交换 \(B\) 和 \(N\) 中的各一列。单纯形法的定价步骤识别目标函数沿各边的方向导数（简约费用 \(\bar{c}_j\)），选择负值对应的边作为下降方向。步长 \(\theta\) 保证到达同一棱上的下一个顶点时仍保持可行性。当所有 \(\bar{c}_j\ge 0\) 时，当前顶点为局部最优，由目标函数的线性性及多面体约束集为凸集，局部最优即全局最优。\(\blacksquare\)

网络流问题与线性规划之间存在着深刻的联系，其中最优雅的体现是网络单纯形法（Network Simplex Method）。在网络流问题中，约束矩阵 \(A\) 是节点-弧关联矩阵，其每一列恰好有一个 \(+1\) 和一个 \(-1\)（分别对应弧的起点和终点），其余元素为零。这一特殊结构使得：（1）基矩阵 \(B\) 对应于支撑图的一棵生成树------每个基本可行解唯一对应一个生成树上的流分布；（2）单纯形法的主元操作对应于在生成树上添加一条非树弧形成唯一环，并沿该环增广流量；（3）定价步骤（计算简约费用）可以在生成树上通过节点势（node potentials）以 \(O(n)\) 时间完成，而非通常的 \(O(mn)\)。这种图论结构使得网络单纯形法在实际中极为高效，即使其理论最坏复杂度为指数级。此外，从几何角度看，网络流多面体（network flow polytope）的顶点恰好是整数点（由全幺模性保证），因此单纯形法在网络流问题中自动产生整数最优解，无需显式施加整数约束------这是线性规划对偶理论在组合优化中最优美的应用之一。

\subsection{252.4 非线性规划的KKT条件}\label{ux975eux7ebfux6027ux89c4ux5212ux7684kktux6761ux4ef6}

线性规划是凸优化的一类特例。对一般非线性规划，Karush-Kuhn-Tucker（KKT）条件提供了最优性的必要条件（凸情形下为充要条件）。

\textbf{定理 252.7}（KKT条件，Karush 1939 / Kuhn-Tucker 1951）：考虑问题\(\min f(x)\) s.t. \(g_i(x) \leq 0\)，\(h_j(x) = 0\)。设\(x^*\)满足约束规范（如LICQ），则存在Lagrange乘子\(\mu_i \geq 0\)、\(\lambda_j\)使得：

\[\nabla f(x^*) + \sum_i \mu_i \nabla g_i(x^*) + \sum_j \lambda_j \nabla h_j(x^*) = 0, \quad \mu_i g_i(x^*) = 0\]

KKT条件是Lagrange乘子法从等式约束到不等式约束的自然推广------互补松弛条件\(\mu_i g_i(x^*) = 0\)是这一推广的核心特征。

\textbf{证明}：设 \(x^*\) 为局部极小值点且满足约束规范（如线性独立约束规范 LICQ：紧约束梯度线性独立）。构造 Lagrange 函数 \(L(x,\mu,\lambda)=f(x)+\sum_i\mu_i g_i(x)+\sum_j\lambda_j h_j(x)\)（\(\mu_i\ge 0\)）。由Farkas引理（或锥分离定理），\(-\nabla f(x^*)\) 落在紧约束梯度的凸锥中：\(-\nabla f(x^*)=\sum_i\mu_i\nabla g_i(x^*)+\sum_j\lambda_j\nabla h_j(x^*)\)（\(\mu_i\ge 0\)）。此即一阶驻点条件。互补松弛性来自 \(\mu_i\ge 0\)、\(g_i(x^*)\le 0\) 且 \(\mu_i g_i(x^*)=0\) 的构造：若 \(g_i(x^*)<0\)（松约束），梯度 \(\nabla g_i(x^*)\) 不参与锥组合，\(\mu_i=0\)。凸规划下 KKT 条件为全局最优的充要条件。\(\blacksquare\)

\hrulefill

\hrulefill

\hrulefill

\hrulefill

\hrulefill

\hrulefill

\section{第425章：内点法的多项式复杂度分析}\label{ux7b2c425ux7ae0ux5185ux70b9ux6cd5ux7684ux591aux9879ux5f0fux590dux6742ux5ea6ux5206ux6790}

Karmarkar 于 1984 年发表的内点法论文是线性规划历史上最重要的突破之一。在此之前，单纯形法虽然在实践中极其高效（平均 \(O(n)\) 次迭代），但其最坏情形复杂度是指数级的（Klee-Minty 反例，1972）。Karmarkar 证明了线性规划存在多项式时间算法：其投影尺度法在 \(O(n L)\) 次迭代（\(L\) 为输入比特长度）内可求得 \(\varepsilon\)-精度解，每步 \(O(n^3)\) 次基本运算------总复杂度 \(O(n^{3.5} L)\)。Karmarkar 的工作不仅解决了 LP 多项式可解性这一理论难题，更开启了内点法的''黄金时代''：随后十年间，Renegar（1988）将复杂度改进至 \(O(\sqrt{n} L)\)，Gonzaga（1989）给出路径跟踪法的 \(O(\sqrt{n} L)\) 分析，而 Monteiro-Adler（1989）将原-对偶内点法推广至半定规划。内点法的多项式复杂度分析建立在三个数学支柱之上：障碍函数的自和谐性（Nesterov-Nemirovski）、中心路径的微分几何性质以及势函数（potential function）的递减性。本章从 Karmarkar 势函数出发，逐步推导路径跟踪法的复杂度，最终在 Nesterov-Nemirovski 框架下给出 LP 内点法复杂度的完整理论分析。

\subsection{134.x Karmarkar势函数与投影梯度法}\label{x-karmarkarux52bfux51fdux6570ux4e0eux6295ux5f71ux68afux5ea6ux6cd5}

\textbf{定义 134.1}（Karmarkar 标准形式）：Karmarkar 的原始算法考虑如下形式的线性规划：

\[\min \{ \mathbf{c}^T \mathbf{x} : A\mathbf{x} = \mathbf{0}, \; \mathbf{e}^T \mathbf{x} = 1, \; \mathbf{x} \geq \mathbf{0} \}\]

其中 \(\mathbf{e} = (1, \ldots, 1)^T\)，且假设已知严格可行内点 \(\mathbf{x}_0 = (1/n)\mathbf{e}\)（单纯形中心），以及最优值为 \(0\)。通过适当的变换（投影变换 + 引入人工变量），任意标准形 LP 可多项式时间归约为 Karmarkar 标准形式。

\textbf{定义 134.2}（Karmarkar 势函数）：对 \(\mathbf{x} > \mathbf{0}\)，定义

\[\phi(\mathbf{x}) = n \ln(\mathbf{c}^T \mathbf{x}) - \sum_{i=1}^n \ln x_i = \sum_{i=1}^n \ln\left( \frac{\mathbf{c}^T \mathbf{x}}{x_i} \right)\]

势函数将两个相互竞争的目标合为一体：\(\mathbf{c}^T \mathbf{x}\)（需最小化）和 \(-\sum \ln x_i\)（障碍项，防止靠近边界）。Karmarkar 的核心引理表明：每一步投影梯度移动使 \(\phi(\mathbf{x})\) 至少下降一个绝对常数 \(\delta > 0\)（如 \(\delta = 1/8\)）。

\textbf{定理 134.1}（Karmarkar 算法的多项式复杂度）：设 \(\mathbf{x}^{(k)}\) 为第 \(k\) 步迭代的内点，\(\mathbf{c}^T \mathbf{x}^{(k)} > 0\)（\(k \geq 1\)）。则经过 \(K = O(n L)\) 次迭代后，\(\mathbf{c}^T \mathbf{x}^{(K)} \leq 2^{-O(L)} \cdot \mathbf{c}^T \mathbf{x}^{(0)}\)，即达到所需精度。

\textbf{证明}（框架）：Karmarkar 投影缩放步：（1）在当前点 \(\mathbf{x}^{(k)}\) 处做仿射变换 \(T(\mathbf{x}) = D_k^{-1} \mathbf{x} / (\mathbf{e}^T D_k^{-1} \mathbf{x})\)（\(D_k = \operatorname{diag}(\mathbf{x}^{(k)})\)），将 \(\mathbf{x}^{(k)}\) 映为单纯形中心 \(\mathbf{a}_0 = (1/n)\mathbf{e}\)。（2）在中心处，向可行子空间 \(\{ \mathbf{d} : AD_k \mathbf{d} = \mathbf{0}, \mathbf{e}^T \mathbf{d} = 0 \}\) 投影负梯度 \(-\nabla (D_k \mathbf{c})\)，得搜索方向 \(\mathbf{d}_k\)。（3）沿 \(\mathbf{d}_k\) 取步长 \(\alpha = 1/(n \sqrt{n})\) 得新点，逆变换回原空间。

关键分析：势函数经变换后在中心邻域内以速率至少 \(\delta\) 递减：\(\phi(\mathbf{x}^{(k+1)}) \leq \phi(\mathbf{x}^{(k)}) - \delta\)（\(\delta = 1/8\)）。原因在于：在球 \(\|\mathbf{x} - \mathbf{a}_0\| \leq \alpha \sqrt{n/(n-1)}\) 内，\(\phi\) 的二次近似和投影负梯度保证充分下降。经过 \(K\) 次迭代，\(\phi(\mathbf{x}^{(K)}) \leq \phi(\mathbf{x}^{(0)}) - K\delta\)。又 \(\phi(\mathbf{x}^{(0)}) = O(n \log n)\)，而所需终止条件 \(\mathbf{c}^T \mathbf{x} \leq 2^{-O(L)} \mathbf{c}^T \mathbf{x}^{(0)}\) 对应 \(\phi(\mathbf{x}) \approx -(n)L \log 2\)（因 \(\ln(\mathbf{c}^T \mathbf{x})\) 主导）。解得 \(K = O(n L)\)。每步计算（投影 + 逆对角阵）\(O(n^3)\)，总复杂度 \(O(n^{3.5} L)\)。\(\blacksquare\)

\subsection{\texorpdfstring{134.y 短步路径跟踪法与 \(O(\sqrt{n} L)\) 迭代}{134.y 短步路径跟踪法与 O(\textbackslash sqrt\{n\} L) 迭代}}\label{y-ux77edux6b65ux8defux5f84ux8ddfux8e2aux6cd5ux4e0e-osqrtn-l-ux8fedux4ee3}

\textbf{定理 134.2}（中心路径的微分性质）：考虑标准形 LP \(\min\{\mathbf{c}^T \mathbf{x} : A\mathbf{x} = \mathbf{b}, \mathbf{x} \geq \mathbf{0}\}\)。对偶障碍参数 \(\mu > 0\)，\textbf{中心路径} \(\mathbf{x}(\mu)\) 定义为

\[\mathbf{x}(\mu) = \arg\min_{\mathbf{x} > \mathbf{0}, A\mathbf{x} = \mathbf{b}} \left\{ \frac{\mathbf{c}^T \mathbf{x}}{\mu} - \sum_{i=1}^n \ln x_i \right\}\]

等价地，\(\mathbf{x}(\mu)\) 满足一阶条件：存在 \(\mathbf{y}(\mu) \in \mathbb{R}^m\) 和 \(\mathbf{s}(\mu) \in \mathbb{R}^n\) 使得 \[A\mathbf{x} = \mathbf{b}, \quad A^T \mathbf{y} + \mathbf{s} = \mathbf{c}, \quad x_i s_i = \mu \;(i = 1, \ldots, n), \quad \mathbf{x} > 0, \;\mathbf{s} > 0\]

当 \(\mu \to 0\) 时，\(\mathbf{x}(\mu)\) 沿中心路径逼近最优解。中心路径是 \(\mathbb{R}^n\) 中一条光滑曲线，其参数化由 \(\mu\) 给出。

\textbf{定理 134.3}（短步路径跟踪算法，Renegar 1988；Gonzaga 1989）：考虑原-对偶内点法的短步变体（short-step path-following），每步将障碍参数 \(\mu\) 更新为 \(\mu^+ = (1 - 1/\sqrt{n})\mu\)（\(n\) 为变量数），然后执行一次 Newton 步校正到新中心路径上。该算法在

\[K = O\left(\sqrt{n} \cdot \ln\left(\frac{n}{\varepsilon}\right)\right)\]

次迭代内产生原-对偶间隙 \(\mu \leq \varepsilon\) 的解。每步运算量 \(O(n^3)\)（求解 Newton 系统），总复杂度 \(O(n^{3.5} \log(1/\varepsilon))\)。

\textbf{证明}：原-对偶中心路径由 \(F(\mathbf{x}, \mathbf{y}, \mathbf{s}; \mu) = \mathbf{0}\) 刻画，其中 \(F_1 = A\mathbf{x} - \mathbf{b}\)，\(F_2 = A^T \mathbf{y} + \mathbf{s} - \mathbf{c}\)，\(F_3 = X S \mathbf{e} - \mu \mathbf{e}\)（\(X = \operatorname{diag}(\mathbf{x})\)，\(S = \operatorname{diag}(\mathbf{s})\)）。设在第 \(k\) 步，当前迭代满足 \(\|X S \mathbf{e} - \mu_k \mathbf{e}\| \leq \theta \mu_k\)（\(\theta \in (0, 1)\) 如 \(\theta = 0.4\)）。更新 \(\mu_{k+1} = (1 - \delta/\sqrt{n}) \mu_k\)（\(\delta > 0\) 为常数）。一次 Newton 步后，新的中心性度量满足 \(\|X^+ S^+ \mathbf{e} - \mu_{k+1} \mathbf{e}\| \leq \theta \mu_{k+1}\)。这是由于对数障碍函数的自和谐性：在中心路径附近，Newton 法具有二次收敛性，且中心路径在度量 \(\|d\mu/\mu\|\) 下的弧长受 \(\sqrt{n}\) 控制------每步缩短 \(\mu\) 的比例 \(1 - O(1/\sqrt{n})\) 保持了局部二次收敛所需的中心性。

关键界：中心路径在 \(\mu\) 对数尺度下的总''长度''为 \(O(\sqrt{n} \ln(\mu_0/\varepsilon))\)。短步法每步沿此路径前进 \(O(1)\) 单位，因此需 \(O(\sqrt{n} \ln(\mu_0/\varepsilon))\) 步。每步的 Newton 系统（正规方程）\(A D^2 A^T \Delta\mathbf{y} = \mathbf{r}\) 可通过稀疏 Cholesky 分解在 \(O(n^3)\) 内求解。\(\blacksquare\)

\subsection{134.z 自和谐函数框架下的内点法统一分析}\label{z-ux81eaux548cux8c10ux51fdux6570ux6846ux67b6ux4e0bux7684ux5185ux70b9ux6cd5ux7edfux4e00ux5206ux6790}

\textbf{定理 134.4}（线性规划对数障碍的自和谐性）：函数 \(f(\mathbf{x}) = -\sum_{i=1}^n \ln x_i\)（定义在 \(\mathbb{R}^n_{>0}\) 上）是自和谐的（\(M=1\)），且是 \(\mathbb{R}^n_{>0}\) 的 \(\nu\)-自和谐障碍函数，其中 \(\nu = n\)。具体地，对任意 \(\mathbf{x} > \mathbf{0}\) 和方向 \(\mathbf{h} \in \mathbb{R}^n\)：

\[|D^3 f(\mathbf{x})[\mathbf{h}, \mathbf{h}, \mathbf{h}]| \leq 2 (D^2 f(\mathbf{x})[\mathbf{h}, \mathbf{h}])^{3/2}\] \[|\nabla f(\mathbf{x})^T \mathbf{h}| \leq \sqrt{n} \cdot (D^2 f(\mathbf{x})[\mathbf{h}, \mathbf{h}])^{1/2}\]

\textbf{证明}：\(f(\mathbf{x}) = -\sum_i \ln x_i\)，偏导数：\(\partial_i f = -1/x_i\)，\(\partial_{ij}^2 f = \delta_{ij} / x_i^2\)，\(\partial_{ijk}^3 f = -2 \delta_{ijk} / x_i^3\)。方向导数：\(D^3 f(\mathbf{x})[\mathbf{h}, \mathbf{h}, \mathbf{h}] = -2 \sum_i h_i^3 / x_i^3\)，\(D^2 f(\mathbf{x})[\mathbf{h}, \mathbf{h}] = \sum_i h_i^2 / x_i^2\)。由 Cauchy-Schwarz 不等式的幂次推广（Holder 不等式）：\(|\sum_i h_i^3 / x_i^3| \leq (\sum_i (h_i^2 / x_i^2)^3)^{1/3} \cdot (\sum_i 1)^{2/3}\) 并不直接。正确的方法：利用一元函数 \(\phi_i(t) = -\ln(x_i + t h_i)\) 的自和谐性。\(|\phi_i'''(0)| = 2 |h_i|^3 / x_i^3 = 2 (\phi_i''(0))^{3/2}\)。故每个 \(\phi_i\) 是 \(M=1\) 自和谐的。自和谐函数之和仍自和谐：由 Minkowski 不等式 \((\sum a_i^{2/3})(\sum b_i^3)^{1/3} \leq \cdots\)，验证 \(|\sum \phi_i'''(0)| \leq 2 (\sum \phi_i''(0))^{3/2}\) 等价于 \(|\sum h_i^3/x_i^3| \leq (\sum h_i^2/x_i^2)^{3/2}\)。令 \(u_i = |h_i|/x_i\)，需证 \(\sum u_i^3 \leq (\sum u_i^2)^{3/2}\)。由 \(\ell_p\) 范数的单调性：\(\|\mathbf{u}\|_3 \leq \|\mathbf{u}\|_2\)（因 \(\|\mathbf{u}\|_3^3 = \sum |u_i|^3 \leq (\sum |u_i|^2)^{3/2} \cdot (\sum 1)^{1/2}\)不恒成立）。正确推导：\(\sum u_i^3 \leq (\max_i u_i) \sum u_i^2 \leq (\sum u_i^2)^{1/2} \cdot \sum u_i^2 = (\sum u_i^2)^{3/2}\)。障碍参数界：\(|\nabla f(\mathbf{x})^T \mathbf{h}| = |\sum h_i/x_i| \leq \sqrt{n} \cdot \sqrt{\sum h_i^2/x_i^2} = \sqrt{n} \cdot (D^2 f(\mathbf{x})[\mathbf{h}, \mathbf{h}])^{1/2}\)（Cauchy-Schwarz）。故 \(\nu = n\)。\(\blacksquare\)

\textbf{定理 134.5}（原-对偶内点法的 Nesterov-Nemirovski 分析）：对标准形 LP 原-对偶内点法，选取自适应步长使得中心性度量（Newton 减量）保持在自和谐函数二次收敛区域内，障碍参数 \(\mu\) 按 \(\mu_{k+1} = \mu_k (1 - \tau/\sqrt{\nu})\) 更新（\(\tau > 0\) 为与自和谐性相关的绝对常数）。则达到对偶间隙 \(\leq \varepsilon\) 所需的总迭代次数为

\[K = O(\sqrt{\nu} \log(\nu \mu_0 / \varepsilon)) = O(\sqrt{n} \log(n/\varepsilon))\]

\textbf{证明}：由定理 134.4，\(f(\mathbf{x}) = -\sum \ln x_i\) 是 \(\nu = n\) 的自和谐障碍函数。原-对偶内点法的 Newton 系统等价于以 \(f\) 为障碍函数、\(\mu\) 为更新参数的路径跟踪法。由 Nesterov-Nemirovski 路径跟踪理论（定理 4.2.8 和推论 4.2.1），自和谐障碍函数族上的短步算法具有 \(O(\sqrt{\nu})\) 迭代复杂度。线性规划中 \(\nu = n\)，故总迭代 \(O(\sqrt{n})\)。每个迭代求解 \(m \times m\) 正规方程组（经对称化后），复杂度 \(O(m^2 n + m^3)\)。在典型情形 \(m = O(n)\) 下为 \(O(n^3)\)。\(\blacksquare\)

\hrulefill

\hrulefill

\hrulefill

\section{第426章：网络单纯形法与整数规划基础}\label{ux7b2c426ux7ae0ux7f51ux7edcux5355ux7eafux5f62ux6cd5ux4e0eux6574ux6570ux89c4ux5212ux57faux7840}

网络单纯形法（Network Simplex Method）是单纯形法在图论框架中最优美的特化形式。对于定义在有向图 \(G = (V, E)\) 上的网络流问题（最小费用流、最大流、运输问题），约束矩阵 \(A\) 是节点-弧关联矩阵------每列恰好有一个 \(+1\)（弧的起点）和一个 \(-1\)（弧的终点），其余为零。这一特殊的全幺模（totally unimodular）结构带来了一系列深远的结果：基矩阵 \(B\) 是生成树的关联矩阵之子矩阵，可逆性等价于该树为生成树；单纯形法的主元操作对应于在生成树上加入一条非树弧形成唯一基本圈（fundamental cycle），并沿该圈增广流量；计算简约费用的定价步骤可在生成树上通过节点势（node potentials）以 \(O(n)\) 时间完成（而非通常的 \(O(mn)\)）。网络单纯形法因此在实际中极为高效，尽管其理论最坏复杂度仍为指数级。更重要的是，全幺模性保证了约束矩阵的每个主子式取值为 \(0\) 或 \(\pm 1\)，从而基本可行解的所有坐标均为整数------这直接导致网络流问题中的整数最优解性质，是整数规划理论中''全幺模矩阵''概念的切入点。

\subsection{135.x 全幺模矩阵与整数极点定理}\label{x-ux5168ux5e7aux6a21ux77e9ux9635ux4e0eux6574ux6570ux6781ux70b9ux5b9aux7406}

\textbf{定义 135.1}（全幺模矩阵 / Totally Unimodular Matrix）：整数矩阵 \(A \in \mathbb{Z}^{m \times n}\) 称为\textbf{全幺模的}（TU），如果 \(A\) 的每个方子矩阵的行列式取值为 \(0\)、\(+1\) 或 \(-1\)。特别地，全幺模矩阵的所有元素限于 \(\{0, \pm 1\}\)。

\textbf{定理 135.1}（Hoffman-Kruskal 定理，1956）：设 \(A \in \mathbb{Z}^{m \times n}\) 全幺模，\(\mathbf{b} \in \mathbb{Z}^m\)。则对任意目标函数 \(\mathbf{c}\)，线性规划 \(\min\{\mathbf{c}^T \mathbf{x} : A\mathbf{x} \leq \mathbf{b}, \mathbf{x} \geq \mathbf{0}\}\) 的所有基本可行解（多面体的顶点）均为整数向量。等价地，多面体 \(\{\mathbf{x} : A\mathbf{x} \leq \mathbf{b}, \mathbf{x} \geq \mathbf{0}\}\) 是整多面体（integral polyhedron）。

\textbf{证明}：设 \(\mathbf{x}^*\) 是基本可行解。由 LP 理论，\(\mathbf{x}^*\) 由 \(A\) 的某 \(n \times n\) 子矩阵 \(B\) 的逆确定：\(\mathbf{x}^* = B^{-1} \mathbf{b}_B\)（非基变量为零）。\(A\) 全幺模 \(\implies\) 每个非奇异方子矩阵 \(B\) 的行列式为 \(\pm 1\)。由 Cramer 法则（伴随矩阵公式 \(B^{-1} = \frac{1}{\det(B)} \operatorname{adj}(B)\)），\(B^{-1}\) 的所有元素为整数除以 \(\pm 1\)，即为整数。又 \(\mathbf{b}_B\) 是整向量，故 \(\mathbf{x}^* = B^{-1} \mathbf{b}_B\) 为整向量。整多面体性质：所有顶点的整性等价于多面体本身为整。\(\blacksquare\)

\textbf{定理 135.2}（全幺模矩阵的刻画 / Ghouila-Houri 1962）：整数矩阵 \(A \in \mathbb{Z}^{m \times n}\) 全幺模当且仅当对 \(A\) 的行的任意子集 \(R \subseteq \{1, \ldots, m\}\)，存在拆分 \(R = R_1 \sqcup R_2\) 使得对每列 \(j\)，\(\sum_{i \in R_1} a_{ij} - \sum_{i \in R_2} a_{ij} \in \{-1, 0, 1\}\)。这一组合刻画为判定全幺模性提供了有效条件。

\textbf{证明}（充分性）：对 \(A\) 的任意 \(k \times k\) 子矩阵，通过归纳法证明其行列式为 \(0, \pm 1\)。归纳基础 \(k=1\) 由 Ghouila-Houri 条件直接得（单个元素为 \(0, \pm 1\)）。归纳步骤：设 \(k \times k\) 子矩阵 \(B\) 满足每列最多两个非零元（全幺模的必要条件）。若 \(B\) 有一列全为零，\(\det B = 0\)。若 \(B\) 有一列仅一个非零元 \(\pm 1\)，按该列展开，由归纳假设 \(\det B = \pm \det B'\) 且 \(B'\) 为 \((k-1) \times (k-1)\) 子矩阵，故 \(\det B \in \{0, \pm 1\}\)。若每列恰有两个非零元且均为 \(\pm 1\)，由 Ghouila-Houri 条件，行可二染色使得每列两非零元异号------这意味着各列非零元之和为 \(0\)，故行线性相关，\(\det B = 0\)。\(\blacksquare\)

\subsection{135.y 网络单纯形法的图论基础}\label{y-ux7f51ux7edcux5355ux7eafux5f62ux6cd5ux7684ux56feux8bbaux57faux7840}

\textbf{定义 135.2}（网络流问题与节点-弧关联矩阵）：给定有向图 \(G = (V, E)\)，\(|V| = n\)，\(|E| = m\)。每条弧 \(e = (i, j) \in E\) 具有容量 \(u_e\) 和单位费用 \(c_e\)。节点-弧关联矩阵 \(A \in \{-1, 0, 1\}^{n \times m}\) 定义为：\(A_{i, e} = +1\) 若 \(e\) 以 \(i\) 为尾（起点），\(A_{i, e} = -1\) 若 \(e\) 以 \(i\) 为头（终点），否则 \(0\)。

\textbf{定理 135.3}（节点-弧关联矩阵的全幺模性）：有向图的节点-弧关联矩阵 \(A\) 是全幺模的。

\textbf{证明}：对 \(A\) 的任意 \(k \times k\) 子矩阵 \(B\)，对 \(k\) 归纳。\(k=1\) 时 \(B\) 的单个元素为 \(0\) 或 \(\pm 1\)，行列式为 \(0\) 或 \(\pm 1\)。设 \(k > 1\)。若 \(B\) 有一列为全零，则 \(\det B = 0\)。若 \(B\) 有一列仅一个非零元 \(\pm 1\)，按该列展开：\(\det B = \pm \det B'\)（\(B'\) 为 \((k-1) \times (k-1)\) 子矩阵），由归纳假设 \(\det B \in \{0, \pm 1\}\)。若 \(B\) 每列有两个非零元（必然一个 \(+1\) 和一个 \(-1\)），则 \(B\) 的各列之和为零向量（因每列恰一个 \(+1\) 和一个 \(-1\)），故 \(B\) 的行线性相关，\(\det B = 0\)。综上，所有 \(k \times k\) 子行列式为 \(0, \pm 1\)，\(A\) 全幺模。\(\blacksquare\)

\textbf{定理 135.4}（网络流问题的生成树基）：设 \(A\) 为连通有向图 \(G\) 的节点-弧关联矩阵，去掉任一行（通常去掉源点对应行，因 \(\mathbf{1}^T A = \mathbf{0}\) 导致 \(A\) 行秩为 \(n-1\)）。则 \(A\) 的任意 \((n-1) \times (n-1)\) 满秩子矩阵 \(B\) 对应的弧集形成 \(G\) 的一棵生成树，反之亦然。因此，网络流问题的每个基对应于一棵生成树。

\textbf{证明}：设 \(E_B\) 为 \(B\) 的列对应的弧集，\(|E_B| = n-1\)。\(B\) 满秩等价于 \(E_B\) 的弧在 \(G\) 中无圈。因为若 \(E_B\) 包含圈，则沿该圈赋予方向一致的 \(\pm 1\) 流量构成 \(B\) 的列的非平凡线性依赖，与满秩矛盾。故 \(E_B\) 是无圈的 \(n-1\) 条边------即生成树。反之，任意生成树的关联子矩阵（去掉一行后）满秩。\(\blacksquare\)

\textbf{定理 135.5}（网络单纯形法的主元操作）：网络单纯形法的每步主元操作等价于：在当前生成树 \(T\) 中选取一条具有负简约费用的非树弧 \(e^+\)（进基弧）；\(T \cup \{e^+\}\) 包含唯一圈 \(C\)；沿 \(C\) 增广流量，方向与 \(e^+\) 同向的弧增加流量，反向的弧减少流量；最先达到上界或下界的弧 \(e^-\) 离基（出基弧）。树结构通过 \(O(n)\) 时间的更新维护（父指针数组和深度优先遍历）。

\textbf{证明}：网络流问题中，简约费用 \(\bar{c}_e = c_e - (\mathbf{y}_u - \mathbf{y}_v)\)（对弧 \(e = (u, v)\)），其中 \(\mathbf{y}\) 为节点势（对偶变量），由 \(B^T \mathbf{y} = \mathbf{c}_B\) 在生成树 \(T\) 上唯一确定（\(\mathbf{c}_B\) 为基弧的原始费用）。节点势可在 \(O(n)\) 时间内在 \(T\) 上通过一次 DFS 计算。选择 \(\bar{c}_e < 0\) 的弧 \(e^+\) 进基。\(T \cup \{e^+\}\) 形成唯一圈 \(C\)。定义方向：沿 \(e^+\) 的方向为正向。在 \(C\) 上沿正向的流量增量为 \(\theta\)，沿反向的流量增量为 \(-\theta\)。\(\theta\) 最大取值由容量和流非负约束确定：\(\theta_{\max} = \min\{\min_{e \in C^+} u_e, \min_{e \in C^-} \text{当前流量}\}\)。达到 \(\theta_{\max}\) 的弧 \(e^-\) 出基。树更新通过 \(e^-\) 的删除和 \(e^+\) 的插入在 \(O(n)\) 时间内完成。\(\blacksquare\)

\subsection{135.z 整数规划基础与Gomory割}\label{z-ux6574ux6570ux89c4ux5212ux57faux7840ux4e0egomoryux5272}

\textbf{定义 135.3}（整数线性规划 / ILP）：整数线性规划的形式为 \(\min\{\mathbf{c}^T \mathbf{x} : A\mathbf{x} = \mathbf{b}, \mathbf{x} \geq \mathbf{0}, \mathbf{x} \in \mathbb{Z}^n\}\)（纯整数规划）或部分变量要求为整数（混合整数规划）。ILP 的一般情形是 NP 完全的，但某些特殊结构具有多项式可解性。

\textbf{定理 135.6}（全幺模矩阵与ILP的多项式可解性）：若 \(A\) 全幺模且 \(\mathbf{b}\) 为整数，则 LP 松弛 \(\min\{\mathbf{c}^T \mathbf{x} : A\mathbf{x} \leq \mathbf{b}, \mathbf{x} \geq \mathbf{0}\}\) 的最优基本可行解自动为整数。因此此类 ILP 可通过求解 LP 松弛在多项式时间内得到最优整数解（无需显式整数约束）。这一观察统一了网络流、指派、运输和最短路径问题的多项式可解性。

\textbf{证明}：由定理 135.1，全幺模性 \(\implies\) 所有基本可行解为整数。单纯形法（或其多项式时间的替代------内点法）收敛到最优基本可行解，该解同时是 ILP 的最优解。LP 的多项式可解性（Karmarkar 1984）确保总时间为输入规模的多项式。\(\blacksquare\)

\textbf{定理 135.7}（Gomory分数割平面，Gomory 1958）：设 LP 松弛的最优单纯形表为 \(\mathbf{x}_B = \mathbf{b}' - N' \mathbf{x}_N\)，其中某基变量 \(x_{B_i}\) 的最优值 \(b'_i \notin \mathbb{Z}\)。则对第 \(i\) 行，分数部分满足 \(\{b'_i\} - \sum_{j \in N} \{a'_{ij}\} x_j \leq 0\)（其中 \(\{a\}\) 表示 \(a\) 的小数部分），这给出了一个割平面（Gomory 分数割），将当前非整数 LP 最优解切除而不损失任何整数可行解。反复添加 Gomory 割并重解 LP，在有限步内收敛到整数最优解。

\textbf{证明}：设第 \(i\) 行方程为 \(x_{B_i} + \sum_{j \in N} a'_{ij} x_j = b'_i\)（\(x_j \geq 0\)）。令 \(b'_i = \lfloor b'_i \rfloor + f_0\)（\(0 < f_0 < 1\)），\(a'_{ij} = \lfloor a'_{ij} \rfloor + f_j\)（\(0 \leq f_j < 1\)）。方程改写为 \(x_{B_i} + \sum \lfloor a'_{ij} \rfloor x_j - \lfloor b'_i \rfloor = f_0 - \sum f_j x_j\)。对所有整数可行解，左端为整数，故右端也必为整数。又左端 \(\leq f_0 < 1\) 且 \(\geq - \sum f_j x_j \leq 0\)（因 \(f_j \geq 0, x_j \geq 0\)），故右端为整数仅当 \(f_0 - \sum f_j x_j \leq 0\)。此即 Gomory 割。切除：当前 LP 基解 \(x_{B_i} = b'_i\)（\(x_N = 0\)），此时 \(\{b'_i\} - \sum_j \{a'_{ij}\} \cdot 0 = f_0 > 0\)，违反割不等式，故该非整数点被切除。所有整数可行解满足此不等式（由推导）。Gomory 的有限步收敛性：由于每组基只能产生有限个 Gomory 割（每行一个），且整数规划仅有有限个顶点，算法在有限步内终止。在实际中需结合分支定界。\(\blacksquare\)

\newpage

\chapter{12-最优化与控制/01-最优化/03-非线性规划.md}\label{ux6700ux4f18ux5316ux4e0eux63a7ux523601-ux6700ux4f18ux531603-ux975eux7ebfux6027ux89c4ux5212.md}

\section{第427章：非线性规划}\label{ux7b2c427ux7ae0ux975eux7ebfux6027ux89c4ux5212}

非线性规划（NLP）处理目标函数或约束为非线性的一般优化问题。Karush-Kuhn-Tucker（KKT）条件是一阶必要条件，在凸情形下也是充分条件。

\subsection{135.1 无约束优化}\label{ux65e0ux7ea6ux675fux4f18ux5316}

\textbf{定理 135.1}（一阶必要条件）：若 \(\mathbf{x}^*\) 是无约束问题 \(\min f(\mathbf{x})\) 的局部极小点，且 \(f\) 在 \(\mathbf{x}^*\) 处可微，则 \(\nabla f(\mathbf{x}^*) = \mathbf{0}\)。

\textbf{证明}：对任意方向 \(\mathbf{d}\ne 0\)，令 \(\phi(t)=f(\mathbf{x}^*+t\mathbf{d})\)。若 \(\mathbf{x}^*\) 为局部极小，则 \(t=0\) 是 \(\phi(t)\) 的局部极小点。由一阶 Taylor 展开 \(\phi(t)=\phi(0)+t\nabla f(\mathbf{x}^*)^T\mathbf{d}+o(t)\)，有 \(\lim_{t\to 0^+}[\phi(t)-\phi(0)]/t=\nabla f(\mathbf{x}^*)^T\mathbf{d}\ge 0\) 且 \(\lim_{t\to 0^-}[\phi(t)-\phi(0)]/t=\nabla f(\mathbf{x}^*)^T\mathbf{d}\le 0\)，故 \(\nabla f(\mathbf{x}^*)^T\mathbf{d}=0\)（对所有 \(\mathbf{d}\)），即 \(\nabla f(\mathbf{x}^*)=0\)。\(\blacksquare\)

\textbf{定理 135.2}（二阶必要条件）：若 \(\mathbf{x}^*\) 是局部极小点，且 \(f\) 在 \(\mathbf{x}^*\) 处二阶可微，则 \(\nabla^2 f(\mathbf{x}^*) \succeq 0\)。

\textbf{证明}：任取 \(\mathbf{d}\ne 0\)，由 Taylor 展开 \(f(\mathbf{x}^*+t\mathbf{d})=f(\mathbf{x}^*)+t\nabla f(\mathbf{x}^*)^T\mathbf{d}+\frac{t^2}{2}\mathbf{d}^T\nabla^2 f(\mathbf{x}^*)\mathbf{d}+o(t^2)\)。因 \(\nabla f(\mathbf{x}^*)=0\)（定理135.1），\(f(\mathbf{x}^*+t\mathbf{d})-f(\mathbf{x}^*)=\frac{t^2}{2}\mathbf{d}^T\nabla^2 f(\mathbf{x}^*)\mathbf{d}+o(t^2)\ge 0\)（局部极小）。除以 \(t^2\) 取 \(t\to 0\) 得 \(\mathbf{d}^T\nabla^2 f(\mathbf{x}^*)\mathbf{d}\ge 0\)，即 \(\nabla^2 f(\mathbf{x}^*)\succeq 0\)。\(\blacksquare\)

\textbf{定理 135.3}（二阶充分条件）：若 \(\nabla f(\mathbf{x}^*) = \mathbf{0}\) 且 \(\nabla^2 f(\mathbf{x}^*) \succ 0\)，则 \(\mathbf{x}^*\) 是严格局部极小点。

\textbf{证明}：由Taylor展开\(f(\mathbf{x}^*+\mathbf{d})=f(\mathbf{x}^*)+\frac{1}{2}\mathbf{d}^T\nabla^2 f(\mathbf{x}^*)\mathbf{d}+o(\|\mathbf{d}\|^2)\)。\(\nabla^2 f(\mathbf{x}^*)\succ 0\)，其特征值\(\lambda_{\min}>0\)。故\(\mathbf{d}^T\nabla^2 f(\mathbf{x}^*)\mathbf{d}\ge\lambda_{\min}\|\mathbf{d}\|^2\)。存在\(\delta>0\)使\(\|\mathbf{d}\|<\delta\)时\(o(\|\mathbf{d}\|^2)\le\frac{\lambda_{\min}}{4}\|\mathbf{d}\|^2\)。则\(f(\mathbf{x}^*+\mathbf{d})-f(\mathbf{x}^*)\ge\frac{\lambda_{\min}}{4}\|\mathbf{d}\|^2>0\)（\(\mathbf{d}\neq 0\)），故\(\mathbf{x}^*\)为严格局部极小。\(\blacksquare\)

\textbf{算法 135.1}（梯度下降法）：\(\mathbf{x}_{k+1} = \mathbf{x}_k - \alpha_k \nabla f(\mathbf{x}_k)\)，其中 \(\alpha_k\) 为步长（由线搜索或固定值确定）。

\textbf{算法 135.2}（Newton 法）：\(\mathbf{x}_{k+1} = \mathbf{x}_k - [\nabla^2 f(\mathbf{x}_k)]^{-1} \nabla f(\mathbf{x}_k)\)。在初始点充分接近最优解时，Newton 法具有二次收敛速度。

\textbf{定理 135.4}（Newton 法的收敛速度）：若 \(f\) 是强凸函数且 Hessian Lipschitz 连续，则 Newton 法局部二次收敛：\(\|\mathbf{x}_{k+1} - \mathbf{x}^*\| \leq C \|\mathbf{x}_k - \mathbf{x}^*\|^2\)。

\textbf{证明}：Newton迭代\(\mathbf{x}_{k+1}=\mathbf{x}_k-[\nabla^2 f(\mathbf{x}_k)]^{-1}\nabla f(\mathbf{x}_k)\)。令\(\mathbf{e}_k=\mathbf{x}_k-\mathbf{x}^*\)。\(\nabla f(\mathbf{x}^*)=0\)，\(\nabla f(\mathbf{x}_k)=\nabla^2 f(\mathbf{x}^*)\mathbf{e}_k+O(\|\mathbf{e}_k\|^2)\)。Hessian Lipschitz连续给出\(\|\nabla^2 f(\mathbf{x}_k)^{-1}\|\le M\)（强凸性保证可逆）。\(\mathbf{e}_{k+1}=\mathbf{e}_k-[\nabla^2 f(\mathbf{x}_k)]^{-1}\nabla f(\mathbf{x}_k)=[\nabla^2 f(\mathbf{x}_k)]^{-1}(\nabla^2 f(\mathbf{x}_k)\mathbf{e}_k-\nabla f(\mathbf{x}_k))\)。由Lipschitz条件\(\|\nabla^2 f(\mathbf{x}_k)\mathbf{e}_k-\nabla f(\mathbf{x}_k)\|\le\frac{L}{2}\|\mathbf{e}_k\|^2\)。故\(\|\mathbf{e}_{k+1}\|\le M\cdot\frac{L}{2}\|\mathbf{e}_k\|^2\)。\(\blacksquare\)

\textbf{算法 135.3}（拟 Newton 法 / BFGS）：维护 Hessian 近似 \(B_k\)，更新公式：

\[B_{k+1} = B_k - \frac{B_k \mathbf{s}_k \mathbf{s}_k^T B_k}{\mathbf{s}_k^T B_k \mathbf{s}_k} + \frac{\mathbf{y}_k \mathbf{y}_k^T}{\mathbf{y}_k^T \mathbf{s}_k}\]

其中 \(\mathbf{s}_k = \mathbf{x}_{k+1} - \mathbf{x}_k\)，\(\mathbf{y}_k = \nabla f(\mathbf{x}_{k+1}) - \nabla f(\mathbf{x}_k)\)。

\subsection{135.2 线搜索方法：Armijo条件与Wolfe条件}\label{ux7ebfux641cux7d22ux65b9ux6cd5armijoux6761ux4ef6ux4e0ewolfeux6761ux4ef6}

\textbf{定义 135.1}（Armijo条件 / 充分下降条件）：给定搜索方向 \(\mathbf{p}_k\) 为下降方向（\(\nabla f(\mathbf{x}_k)^T\mathbf{p}_k<0\)），步长 \(\alpha\) 满足 \textbf{Armijo条件}，若存在常数 \(c_1\in(0,1)\) 使得

\[f(\mathbf{x}_k + \alpha\mathbf{p}_k) \leq f(\mathbf{x}_k) + c_1\alpha\,\nabla f(\mathbf{x}_k)^T\mathbf{p}_k\]

此条件保证目标函数有充分下降（至少为Taylor一阶近似的 \(c_1\) 倍）。通常取 \(c_1=10^{-4}\)。

\textbf{定义 135.2}（Wolfe条件）：步长 \(\alpha\) 满足 \textbf{Wolfe条件}，若同时满足：

\begin{enumerate}
\def\labelenumi{(\roman{enumi})}
\item
  \textbf{Armijo条件}（充分下降）：\(f(\mathbf{x}_k+\alpha\mathbf{p}_k)\leq f(\mathbf{x}_k)+c_1\alpha\nabla f(\mathbf{x}_k)^T\mathbf{p}_k\)
\item
  \textbf{曲率条件}：\(\nabla f(\mathbf{x}_k+\alpha\mathbf{p}_k)^T\mathbf{p}_k \geq c_2\,\nabla f(\mathbf{x}_k)^T\mathbf{p}_k\)
\end{enumerate}

其中 \(0<c_1<c_2<1\)。曲率条件防止步长过小，确保梯度有充分变化。当 \(|\nabla f(\mathbf{x}_k+\alpha\mathbf{p}_k)^T\mathbf{p}_k|\leq c_2|\nabla f(\mathbf{x}_k)^T\mathbf{p}_k|\) 时称为\textbf{强Wolfe条件}。

\subsection{135.3 约束优化与KKT条件}\label{ux7ea6ux675fux4f18ux5316ux4e0ekktux6761ux4ef6}

\textbf{定义 135.3}（一般约束优化问题）：

\[\begin{aligned} \min \quad & f(\mathbf{x}) \\ \text{s.t.} \quad & g_i(\mathbf{x}) \leq 0, \quad i = 1, \ldots, m \\ & h_j(\mathbf{x}) = 0, \quad j = 1, \ldots, p \end{aligned}\]

\textbf{定义 135.4}（LICQ / 线性独立约束规范）：在点 \(\mathbf{x}\) 处，活跃不等式约束的梯度 \(\nabla g_i(\mathbf{x})\)（\(\mathcal{A}(\mathbf{x})=\{i: g_i(\mathbf{x})=0\}\)）与等式约束的梯度 \(\nabla h_j(\mathbf{x})\) 构成线性无关组。

\textbf{定理 135.5}（KKT 必要条件，Karush-Kuhn-Tucker，1951）：设 \(\mathbf{x}^*\) 是局部极小点，\(f, g_i, h_j\) 连续可微，且约束规范条件（如 LICQ）成立。则存在 Lagrange 乘子 \(\boldsymbol{\lambda}^* \geq \mathbf{0}\) 和 \(\boldsymbol{\nu}^*\) 满足：

\textbf{证明}：由LICQ，活跃约束梯度线性无关。线性化可行方向锥\(\mathcal{F}(\mathbf{x}^*)=\{\mathbf{d}:\nabla g_i(\mathbf{x}^*)^T\mathbf{d}\le 0(i\in\mathcal{A}),\nabla h_j(\mathbf{x}^*)^T\mathbf{d}=0\}\)与切锥一致。由一阶必要性\(\nabla f(\mathbf{x}^*)^T\mathbf{d}\ge 0\)对所有\(\mathbf{d}\in\mathcal{F}(\mathbf{x}^*)\)。Farkas引理给出\(\nabla f(\mathbf{x}^*)+\sum_{i\in\mathcal{A}}\lambda_i^*\nabla g_i(\mathbf{x}^*)+\sum_{j}\nu_j^*\nabla h_j(\mathbf{x}^*)=0\)，\(\lambda_i^*\ge 0\)。非活跃约束取\(\lambda_i^*=0\)。互补松弛\(\lambda_i^*g_i(\mathbf{x}^*)=0\)由活跃集定义自然成立。\(\blacksquare\)

\[\begin{aligned}
\nabla f(\mathbf{x}^*) + \sum_{i=1}^m \lambda_i^* \nabla g_i(\mathbf{x}^*) + \sum_{j=1}^p \nu_j^* \nabla h_j(\mathbf{x}^*) &= \mathbf{0} \quad &\text{(稳定性)} \\
g_i(\mathbf{x}^*) &\leq 0, \quad i = 1, \ldots, m \quad &\text{(原可行性)} \\
h_j(\mathbf{x}^*) &= 0, \quad j = 1, \ldots, p \quad &\text{(原可行性)} \\
\lambda_i^* &\geq 0, \quad i = 1, \ldots, m \quad &\text{(对偶可行性)} \\
\lambda_i^* g_i(\mathbf{x}^*) &= 0, \quad i = 1, \ldots, m \quad &\text{(互补松弛性)}
\end{aligned}\]

\textbf{定理 135.6}（KKT 充分条件，凸情形）：若 \(f\) 和 \(g_i\) 是凸函数，\(h_j\) 是仿射函数，且 \(\mathbf{x}^*\) 和 \((\boldsymbol{\lambda}^*, \boldsymbol{\nu}^*)\) 满足 KKT 条件，则 \(\mathbf{x}^*\) 是全局最优解。

\textbf{证明}：对任意可行点\(\mathbf{x}\)，由\(f\)凸性：\(f(\mathbf{x})\ge f(\mathbf{x}^*)+\nabla f(\mathbf{x}^*)^T(\mathbf{x}-\mathbf{x}^*)\)。代入KKT稳定性条件：\(=f(\mathbf{x}^*)-\sum_i\lambda_i^*\nabla g_i(\mathbf{x}^*)^T(\mathbf{x}-\mathbf{x}^*)-\sum_j\nu_j^*\nabla h_j(\mathbf{x}^*)^T(\mathbf{x}-\mathbf{x}^*)\)。由\(g_i\)凸性：\(g_i(\mathbf{x})\ge g_i(\mathbf{x}^*)+\nabla g_i(\mathbf{x}^*)^T(\mathbf{x}-\mathbf{x}^*)\)。\(h_j\)仿射：\(h_j(\mathbf{x})=h_j(\mathbf{x}^*)+\nabla h_j(\mathbf{x}^*)^T(\mathbf{x}-\mathbf{x}^*)\)。结合互补松弛\(\lambda_i^*g_i(\mathbf{x}^*)=0\)和可行性，得\(f(\mathbf{x})\ge f(\mathbf{x}^*)+\sum_i\lambda_i^*g_i(\mathbf{x})+\sum_j\nu_j^*h_j(\mathbf{x})\ge f(\mathbf{x}^*)\)。\(\blacksquare\)

\subsection{135.4 约束规范条件}\label{ux7ea6ux675fux89c4ux8303ux6761ux4ef6}

\textbf{定义 135.5}（常见约束规范 / Constraint Qualifications）： - \textbf{LICQ}（线性独立约束规范）：活跃约束的梯度线性无关（最强） - \textbf{MFCQ}（Mangasarian-Fromovitz CQ）：等式约束梯度线性无关，且存在方向使所有不等式约束严格减小 - \textbf{Slater 条件}（凸情形）：存在严格可行点

\textbf{命题 135.7}：LICQ \(\implies\) MFCQ \(\implies\) Abadie CQ。在凸情形下，Slater 条件 \(\implies\) KKT 条件成立。

\subsection{135.5 罚函数法与增广 Lagrange 法}\label{ux7f5aux51fdux6570ux6cd5ux4e0eux589eux5e7f-lagrange-ux6cd5}

\textbf{定义 135.6}（外罚函数法）：将约束优化转化为无约束问题：

\[\min_{\mathbf{x}} f(\mathbf{x}) + \frac{\rho}{2} \left( \sum_{i=1}^m \max\{0, g_i(\mathbf{x})\}^2 + \sum_{j=1}^p h_j(\mathbf{x})^2 \right)\]

当 \(\rho \to \infty\) 时，解趋近于原约束问题的最优解。

\textbf{定义 135.7}（增广 Lagrange 函数 / Augmented Lagrangian）：

\[L_A(\mathbf{x}, \boldsymbol{\lambda}, \boldsymbol{\nu}; \rho) = f(\mathbf{x}) + \sum_{i=1}^m \lambda_i g_i(\mathbf{x}) + \sum_{j=1}^p \nu_j h_j(\mathbf{x}) + \frac{\rho}{2}\left( \sum_{i=1}^m g_i(\mathbf{x})^2 + \sum_{j=1}^p h_j(\mathbf{x})^2 \right)\]

（仅对等式约束使用二次罚项，不等式约束可通过引入松弛变量转化为等式约束）

\textbf{算法 135.4}（增广 Lagrange 法 / 乘子法）： 1. 初始化 \(\boldsymbol{\lambda}^0, \boldsymbol{\nu}^0, \rho_0\) 2. 求解 \(\mathbf{x}^{k+1} = \arg\min_{\mathbf{x}} L_A(\mathbf{x}, \boldsymbol{\lambda}^k, \boldsymbol{\nu}^k; \rho_k)\) 3. 更新乘子：\(\lambda_i^{k+1} = \max\{0, \lambda_i^k + \rho_k g_i(\mathbf{x}^{k+1})\}\)，\(\nu_j^{k+1} = \nu_j^k + \rho_k h_j(\mathbf{x}^{k+1})\) 4. 增大 \(\rho_{k+1} > \rho_k\)，重复

\subsection{135.6 序列二次规划（SQP）}\label{ux5e8fux5217ux4e8cux6b21ux89c4ux5212sqp}

\textbf{定义 135.8}（SQP框架 / Sequential Quadratic Programming）：SQP在每次迭代中求解一个\textbf{二次规划（QP）子问题}来生成搜索方向。在第 \(k\) 步，求解

\[\begin{aligned}
\min_{\mathbf{d}} \quad & \frac{1}{2}\mathbf{d}^T \nabla_{\mathbf{x}\mathbf{x}}^2 L(\mathbf{x}_k,\boldsymbol{\lambda}_k,\boldsymbol{\nu}_k)\,\mathbf{d} + \nabla f(\mathbf{x}_k)^T\mathbf{d} \\
\text{s.t.} \quad & \nabla g_i(\mathbf{x}_k)^T\mathbf{d} + g_i(\mathbf{x}_k) \leq 0, \quad i=1,\ldots,m \\
& \nabla h_j(\mathbf{x}_k)^T\mathbf{d} + h_j(\mathbf{x}_k) = 0, \quad j=1,\ldots,p
\end{aligned}\]

其中 \(\nabla_{\mathbf{x}\mathbf{x}}^2 L\) 为Lagrange函数的Hessian（或拟Newton近似 \(B_k\)）。QP子问题的最优解 \(\mathbf{d}_k\) 作为搜索方向，相应的QP乘子作为下一步Lagrange乘子的估计。

\textbf{定理 135.8}（SQP的局部收敛性）：在LICQ和二阶充分条件以及 \(\nabla_{\mathbf{x}\mathbf{x}}^2 L\) 在解处的正定性下，SQP方法具有局部二阶收敛速度。使用BFGS近似Hessian时，达到超线性收敛。

\textbf{证明}：SQP的QP子问题的KKT条件等价于Newton法应用于Lagrange函数的KKT系统。令\(\mathbf{z}=(\mathbf{x},\boldsymbol{\lambda},\boldsymbol{\nu})\)，KKT系统为\(F(\mathbf{z})=0\)（\(F\)为梯度映射）。SQP步\(\mathbf{d}_k\)满足\(F'(\mathbf{z}_k)\mathbf{d}_k=-F(\mathbf{z}_k)\)（Newton迭代）。在LICQ和二阶充分条件下，\(F'(\mathbf{z}^*)\)非奇异，由Newton法局部二次收敛：\(\|\mathbf{z}_{k+1}-\mathbf{z}^*\|\le C\|\mathbf{z}_k-\mathbf{z}^*\|^2\)。BFGS近似满足Dennis-Moré条件，保证超线性收敛。\(\blacksquare\)

\subsection{135.7 信赖域方法与Cauchy点}\label{ux4fe1ux8d56ux57dfux65b9ux6cd5ux4e0ecauchyux70b9}

\textbf{定义 135.9}（信赖域子问题）：在第 \(k\) 步，给定信赖域半径 \(\Delta_k>0\)，在当前点 \(\mathbf{x}_k\) 附近求解

\[\min_{\mathbf{d}} \; m_k(\mathbf{d}) = f(\mathbf{x}_k) + \nabla f(\mathbf{x}_k)^T\mathbf{d} + \frac{1}{2}\mathbf{d}^T B_k\,\mathbf{d} \quad \text{s.t.} \quad \|\mathbf{d}\| \leq \Delta_k\]

其中 \(B_k\) 是 \(\nabla^2 f(\mathbf{x}_k)\) 或其近似。信赖域半径 \(\Delta_k\) 根据实际下降与模型预测下降之比 \(\rho_k = \frac{f(\mathbf{x}_k)-f(\mathbf{x}_k+\mathbf{d}_k)}{m_k(0)-m_k(\mathbf{d}_k)}\) 进行调整：若 \(\rho_k\approx 1\)，增大 \(\Delta_k\)；若 \(\rho_k\) 过小或为负，缩小 \(\Delta_k\)。

\textbf{定义 135.10}（Cauchy点）：信赖域子问题中沿负梯度方向的最优步长点称为\textbf{Cauchy点}：

\[\mathbf{d}_k^C = -\tau_k \frac{\Delta_k}{\|\nabla f(\mathbf{x}_k)\|}\nabla f(\mathbf{x}_k), \quad \tau_k = \min\!\left\{1,\;\frac{\|\nabla f(\mathbf{x}_k)\|^3}{\Delta_k\,\nabla f(\mathbf{x}_k)^T B_k\nabla f(\mathbf{x}_k)}\right\}\]

Cauchy点是模型 \(m_k\) 在梯度方向与信赖域边界交点上的精确一维极小化结果，它提供了全局收敛性的保证。

\hrulefill

\hrulefill

\hrulefill

\section{第428章：内点法与自和谐函数（IPM的理论分析）}\label{ux7b2c428ux7ae0ux5185ux70b9ux6cd5ux4e0eux81eaux548cux8c10ux51fdux6570ipmux7684ux7406ux8bbaux5206ux6790}

内点法（Interior-Point Methods, IPM）自 Karmarkar（1984）的开创性工作以来，已成为求解线性规划和凸规划最强大的算法范式。与单纯形法沿可行域边界移动不同，内点法从严格可行域内部出发，通过障碍函数（barrier function）将约束''编码''到目标函数中，沿中心路径（central path）逐步逼近最优解。Nesterov 和 Nemirovski（1994）在其里程碑式的著作《Interior-Point Polynomial Algorithms in Convex Programming》中建立了内点法的统一理论框架------自和谐函数（self-concordant functions）理论。这一理论的核心洞察是：Newton 法的局部二次收敛性可以通过适当的函数类选择推广到全局多项式复杂度。自和谐函数的定义是二阶导数的 Lipschitz 条件的一种仿射不变形式：函数 \(f\) 称为自和谐的，若其三阶导数被其二阶导数以 \(2\)-范数界定。这一条件在仿射变换下保持不变，从而使得 Newton 法的分析不依赖于坐标选择------这正是内点法多项式复杂度的理论基础。自和谐函数理论将椭球法、重心法和内点法统一在同一个凸优化复杂度框架下，揭示了对数障碍函数 \(\phi(x) = -\sum_i \ln(b_i - a_i^T x)\) 的自和谐性是内点法高效的根源。

\subsection{136.x 自和谐函数的定义与基本性质}\label{x-ux81eaux548cux8c10ux51fdux6570ux7684ux5b9aux4e49ux4e0eux57faux672cux6027ux8d28}

\textbf{定义 136.1}（自和谐函数，Nesterov-Nemirovski 1994）：设 \(f: \operatorname{dom} f \subseteq \mathbb{R}^n \to \mathbb{R}\) 是三次连续可微的凸函数，且 \(\operatorname{dom} f\) 为开集。\(f\) 称为 \textbf{\(M\)-自和谐函数}（\(M \geq 0\)），如果对任意 \(x \in \operatorname{dom} f\) 和任意方向 \(h \in \mathbb{R}^n\)，其三阶方向导数满足

\[|D^3 f(x)[h, h, h]| \leq 2M \cdot (D^2 f(x)[h, h])^{3/2}\]

其中 \(D^3 f(x)[h, h, h] = \frac{d^3}{dt^3} f(x + th)\big|_{t=0}\)，\(D^2 f(x)[h, h] = h^T \nabla^2 f(x) h\)。当 \(M = 1\) 时简称为\textbf{自和谐函数}（self-concordant）。若进一步 \(\operatorname{dom} f\) 是有界集且 \(f\) 在 \(\operatorname{dom} f\) 的边界上趋于无穷（即 \(f\) 是 \(\operatorname{dom} f\) 的障碍函数），则称 \(f\) 为\textbf{强自和谐障碍函数}（strongly self-concordant barrier）。

\textbf{定理 136.1}（自和谐函数的基本不等式）：设 \(f\) 是 \(M\)-自和谐函数。对任意 \(x, y \in \operatorname{dom} f\)，记 \(r = \|y - x\|_x = \sqrt{(y-x)^T \nabla^2 f(x) (y-x)}\) 为 \(x\) 处的局部范数。则 \[(1 - Mr_+)^2 \nabla^2 f(x) \preceq \nabla^2 f(y) \preceq \frac{1}{(1 - Mr_+)^2} \nabla^2 f(x)\] 其中 \(r_+ = \min\{r, 1/M\}\)（或 \(r_+ = r\) 当 \(r < 1/M\)，否则右端发散）。特别地，当 \(r < 1/M\) 时，Hessian 的有界变化保证了 Newton 法的局部二次收敛性。

\textbf{证明}：固定 \(x\)，考虑一元函数 \(\phi(t) = f(x + t d)\)，其中 \(d = (y-x)/r\) 满足 \(\|d\|_x = 1\)。因 \(f\) 自和谐，\(\phi\) 满足 \(|\phi'''(t)| \leq 2M \phi''(t)^{3/2}\)。令 \(\psi(t) = \phi''(t)^{-1/2}\)，则 \(\psi'(t) = -\frac{1}{2} \phi''(t)^{-3/2} \phi'''(t)\)，从而 \(|\psi'(t)| \leq M\)。由中值定理，\(|\psi(t) - \psi(0)| \leq M t\)。故 \(\phi''(0)^{1/2} / (1 + M t \phi''(0)^{1/2}) \leq \phi''(t)^{1/2} \leq \phi''(0)^{1/2} / (1 - M t \phi''(0)^{1/2})\)（当 \(t < 1/(M \phi''(0)^{1/2})\)）。平方后即得局部范数下的 Hessian 界。由于 \(\phi''(0) = h^T \nabla^2 f(x) h = r^2 \|d\|_x^2 = r^2\)，代入 \(t = r\) 即得结论。\(\blacksquare\)

\textbf{定理 136.2}（标准自和谐函数的例子）： (i) 线性函数和凸二次函数是自和谐的（\(M=0\)）。 (ii) \(-\ln x\)（\(x > 0\)）是参数 \(M=1\) 的自和谐函数。 (iii) 对数障碍函数 \(f(x) = -\sum_{i=1}^m \ln(b_i - a_i^T x)\) 在多面体 \(\{x : a_i^T x < b_i\}\) 上是自和谐的（\(M=1\)）。 (iv) \(-\ln \det X\)（\(X \succ 0\)，在正定锥上）是 \(M=1\) 的自和谐函数。 (v) 半定规划的对数障碍 \(-\ln \det(F_0 + \sum_i x_i F_i)\)（\(F_i\) 对称）是自和谐的。

\textbf{证明}：(ii) 对 \(f(x) = -\ln x\)，\(f'(x) = -1/x\)，\(f''(x) = 1/x^2\)，\(f'''(x) = -2/x^3\)。故 \(|f'''(x)| = 2 / x^3 = 2 (f''(x))^{3/2}\)，恰取等号，\(M=1\)。(iii) 是 (ii) 在线性变换下的复合：\(f_i(x) = -\ln(b_i - a_i^T x)\)，其三阶导数为 \(|f_i'''(x)[h,h,h]| = 2|a_i^T h|^3 / (b_i - a_i^T x)^3 = 2(f_i''(x)[h,h])^{3/2}\)。由于自和谐性在非负求和下保持（Cauchy-Schwarz 不等式），得证。(iv) 应用矩阵计算：对 \(F(X) = -\ln \det X\)，方向导数 \(D F(X)[H] = -\operatorname{tr}(X^{-1}H)\)，\(D^2 F(X)[H,H] = \operatorname{tr}(X^{-1}H X^{-1}H)\)，\(D^3 F(X)[H,H,H] = -2 \operatorname{tr}((X^{-1}H)^3)\)。由迹不等式 \(|\operatorname{tr}(A^3)| \leq (\operatorname{tr}(A^2))^{3/2}\) 得自和谐性。\(\blacksquare\)

\subsection{136.y 自和谐条件下的Newton法多项式复杂度}\label{y-ux81eaux548cux8c10ux6761ux4ef6ux4e0bux7684newtonux6cd5ux591aux9879ux5f0fux590dux6742ux5ea6}

\textbf{定理 136.3}（自和谐函数的Newton法收敛性，Nesterov-Nemirovski）：设 \(f\) 是标准的自和谐函数（\(M=1\)），\(\operatorname{dom} f\) 为开集。考虑 Newton 法 \(\mathbf{x}_{k+1} = \mathbf{x}_k - [\nabla^2 f(\mathbf{x}_k)]^{-1} \nabla f(\mathbf{x}_k)\)。定义 Newton 减量 \(\lambda(\mathbf{x}) = \|\nabla f(\mathbf{x})\|_{\mathbf{x}}^* = \sqrt{\nabla f(\mathbf{x})^T [\nabla^2 f(\mathbf{x})]^{-1} \nabla f(\mathbf{x})}\)。则：

\begin{enumerate}
\def\labelenumi{(\roman{enumi})}
\item
  若 \(\lambda(\mathbf{x}_k) \leq 1/4\)，则 Newton 法进入二次收敛区域：\(\lambda(\mathbf{x}_{k+1}) \leq 2 \lambda(\mathbf{x}_k)^2\)。
\item
  在阻尼 Newton 阶段（\(\lambda(\mathbf{x}_k) > 1/4\)），取步长 \(\alpha_k = 1/(1 + \lambda(\mathbf{x}_k))\)，则 \(f(\mathbf{x}_{k+1}) - f(\mathbf{x}_k) \leq -\omega\)（\(\omega \approx 0.027\) 为绝对常数）。
\item
  达到精度 \(\varepsilon\)（即 \(f(\mathbf{x}_k) - f(\mathbf{x}^*) \leq \varepsilon\)）所需 Newton 迭代总次数为 \(O(1)(f(\mathbf{x}_0) - f(\mathbf{x}^*)) + O(\log \log(1/\varepsilon))\)。
\end{enumerate}

\textbf{证明}：(i) 设 \(\mathbf{x}^+ = \mathbf{x} - [\nabla^2 f(\mathbf{x})]^{-1} \nabla f(\mathbf{x})\)，Newton 步长 \(\mathbf{d} = -\nabla^2 f(\mathbf{x})^{-1} \nabla f(\mathbf{x})\)。局部范数 \(\|\mathbf{d}\|_{\mathbf{x}} = \lambda(\mathbf{x})\)。由定理 136.1，\(\nabla^2 f\) 在以 \(\mathbf{x}\) 为心、半径 \(\lambda < 1\) 的局部范数球内以 \((1 \pm \lambda)^{-2}\) 被控制。利用 Newton 步的一阶最优性条件和自和谐性的积分形式，可得 \(\lambda(\mathbf{x}^+) \leq \lambda(\mathbf{x})^2 / (1 - \lambda(\mathbf{x}))^2\)。当 \(\lambda(\mathbf{x}) \leq 1/4\)，\(\lambda(\mathbf{x}^+) \leq \lambda^2 / (3/4)^2 = (16/9) \lambda^2 \leq 2\lambda^2\)。(ii) 阻尼步长保证 Dikin 椭球包含在函数定义域内，且沿 Newton 方向函数值至少下降 \(\lambda - \ln(1+\lambda)\)。(iii) 阻尼阶段每步下降常数，二次阶段每步平方收敛，总计复杂度如所述。\(\blacksquare\)

\textbf{定理 136.4}（自和谐障碍法 / 路径跟踪法的复杂度）：考虑凸规划 \(\min\{\langle c, x \rangle : x \in K\}\)，其中 \(K\) 为闭凸集，配有一个 \(\nu\)-自和谐障碍函数 \(F: \operatorname{int} K \to \mathbb{R}\)，即 \(F\) 自和谐且 \(|\nabla F(x)^T h|^2 \leq \nu \cdot h^T \nabla^2 F(x) h\)（\(\nu\) 为障碍参数）。则短步路径跟踪法（short-step path-following）在至多

\[O\left(\sqrt{\nu} \cdot \ln\left(\frac{\nu}{\varepsilon}\right)\right)\]

次 Newton 迭代内产生 \(\varepsilon\)-近似解。当 \(K\) 为线性规划的可行多面体 \(K = \{x : Ax \leq b\}\) 时，\(\nu = m\)（约束数），故复杂度为 \(O(\sqrt{m} \ln(1/\varepsilon))\)，优于单纯形法的 \(O(2^m)\) 最坏情形。

\textbf{证明}：中心路径定义为 \(\mathbf{x}_\mu^* = \arg\min_{x \in \operatorname{int} K} \{\langle c, x \rangle + \mu F(x)\}\)（\(\mu > 0\)）。障碍函数的最优性条件为 \(\nabla F(x) = -c/\mu\)，即 \(c = -\mu \nabla F(x)\)。中心路径的弧长（在局部 Riemann 度量下）被自和谐性控制。短步法保持当前迭代与中心路径的 Newton 减量 \(\lambda \leq 0.1\)，每步将 \(\mu\) 缩为 \(\mu(1 - 1/\sqrt{\nu})\)。分析表明每次 Newton 迭代校正中心性的代价为常数步，\(\mu\) 从 \(\mu_0\) 降至 \(\varepsilon/\nu\) 需 \(O(\sqrt{\nu} \ln(\nu/\varepsilon))\) 次外迭代，每外迭代 \(O(1)\) 次 Newton 步。总计 \(O(\sqrt{\nu} \ln(\nu/\varepsilon))\)。\(\blacksquare\)

\hrulefill

\hrulefill

\hrulefill

\section{第429章：半光滑Newton法与增广Lagrangian法}\label{ux7b2c429ux7ae0ux534aux5149ux6ed1newtonux6cd5ux4e0eux589eux5e7flagrangianux6cd5}

半光滑 Newton 法（Semismooth Newton Method）是求解非光滑方程组的核心数值方法，由 Qi 和 Sun（1993）在 Mifflin（1977）的半光滑性概念基础上系统发展。其基本思想是将 Newton 法推广到不可微但具有广义可微结构的函数：若 \(F: \mathbb{R}^n \to \mathbb{R}^n\) 是局部 Lipschitz 连续且在 Clarke 广义 Jacobi 意义下半光滑，则在 Newton 迭代中以广义 Jacobi \(\partial F(\mathbf{x}_k)\) 的元素 \(V_k\) 替代经典 Jacobi，得到的半光滑 Newton 法 \(\mathbf{x}_{k+1} = \mathbf{x}_k - V_k^{-1} F(\mathbf{x}_k)\) 仍具有超线性收敛性甚至二次收敛性（当 \(F\) 强半光滑时）。半光滑 Newton 法最重要的应用领域之一是求解互补问题（complementarity problems）和变分不等式：通过 Fischer-Burmeister 函数 \(\phi(a,b) = \sqrt{a^2 + b^2} - a - b\) 或 \(\min\) 函数将互补条件 \(a \geq 0, b \geq 0, ab = 0\) 等价转化为半光滑方程组 \(\phi(a,b) = 0\)，进而应用半光滑 Newton 法高效求解。增广 Lagrangian 法（Augmented Lagrangian Method, ALM）则是求解约束优化的经典方法，通过在原 Lagrange 函数上添加二次惩罚项，将约束问题转化为一系列无约束或简单约束的子问题。ALM 的收敛性分析比经典罚函数法更优秀：在有限惩罚参数下即可收敛（无需 \(\rho \to \infty\)），且乘子以线性速率收敛。ALM 与邻近点算法（Proximal Point Algorithm）之间存在深刻的理论对偶关系，这一关系在 Rockafellar（1976）的变分分析框架中得到完整阐明。

\subsection{137.x 半光滑性与广义Newton法}\label{x-ux534aux5149ux6ed1ux6027ux4e0eux5e7fux4e49newtonux6cd5}

\textbf{定义 137.1}（局部 Lipschitz 连续与 Clarke 广义 Jacobi）：函数 \(F: \mathbb{R}^n \to \mathbb{R}^m\) 称为局部 Lipschitz 连续的，若对任意紧集 \(K \subset \mathbb{R}^n\) 存在常数 \(L_K > 0\) 使得 \(\|F(\mathbf{x}) - F(\mathbf{y})\| \leq L_K \|\mathbf{x} - \mathbf{y}\|\)（对所有 \(\mathbf{x}, \mathbf{y} \in K\)）。由 Rademacher 定理，\(F\) 几乎处处可微。\(F\) 在 \(\mathbf{x}\) 处的 \textbf{Clarke 广义 Jacobi} 定义为

\[\partial F(\mathbf{x}) = \operatorname{conv}\left\{\lim_{k \to \infty} \nabla F(\mathbf{x}_k) : \mathbf{x}_k \to \mathbf{x}, \; \mathbf{x}_k \in D_F\right\}\]

其中 \(D_F\) 是 \(F\) 可微点的集合，\(\operatorname{conv}\) 表示凸包。当 \(m=1\) 时，\(\partial F(\mathbf{x})\) 退化为 Clarke 次微分（subdifferential）。

\textbf{定义 137.2}（半光滑性，Mifflin 1977；Qi-Sun 1993）：局部 Lipschitz 连续的 \(F: \mathbb{R}^n \to \mathbb{R}^n\) 在 \(\mathbf{x}\) 处称为\textbf{半光滑的}（semismooth），若 \(F\) 在 \(\mathbf{x}\) 处方向可微，且对任意 \(\mathbf{h} \to \mathbf{0}\) 和任意 \(V \in \partial F(\mathbf{x}+\mathbf{h})\)，

\[F(\mathbf{x}+\mathbf{h}) - F(\mathbf{x}) - V\mathbf{h} = o(\|\mathbf{h}\|)\]

若进一步上式余项为 \(O(\|\mathbf{h}\|^2)\)（即 \(o(\|\mathbf{h}\|)\) 替换为 \(O(\|\mathbf{h}\|^2)\)），则称 \(F\) 在 \(\mathbf{x}\) 处\textbf{强半光滑}（strongly semismooth）。分段光滑函数、\(\max\) / \(\min\) 函数的复合、Fischer-Burmeister 函数等均是半光滑的；若进一步各分片是 \(C^2\) 的，则是强半光滑的。

\textbf{定理 137.1}（半光滑Newton法的超线性收敛，Qi-Sun 1993）：设 \(F: \mathbb{R}^n \to \mathbb{R}^n\) 在 \(\mathbf{x}^*\) 处半光滑且 \(F(\mathbf{x}^*) = \mathbf{0}\)。进一步假设 \(\partial F(\mathbf{x}^*)\) 中的所有元素非奇异（即 Clarke 广义 Jacobi 非退化）。则存在 \(\mathbf{x}^*\) 的邻域，使得从该邻域内任意初始点出发，半光滑 Newton 迭代

\[\mathbf{x}_{k+1} = \mathbf{x}_k - V_k^{-1} F(\mathbf{x}_k), \quad V_k \in \partial F(\mathbf{x}_k)\]

产生的序列 \(\{\mathbf{x}_k\}\) 超线性收敛到 \(\mathbf{x}^*\)，即 \(\|\mathbf{x}_{k+1} - \mathbf{x}^*\| = o(\|\mathbf{x}_k - \mathbf{x}^*\|)\)。若 \(F\) 强半光滑且 \(\nabla F\)（在可微点）局部 Lipschitz 连续，则收敛是二次的：\(\|\mathbf{x}_{k+1} - \mathbf{x}^*\| = O(\|\mathbf{x}_k - \mathbf{x}^*\|^2)\)。

\textbf{证明}：令 \(\mathbf{e}_k = \mathbf{x}_k - \mathbf{x}^*\)。由迭代公式，\(\mathbf{e}_{k+1} = \mathbf{e}_k - V_k^{-1} F(\mathbf{x}_k) = V_k^{-1}(V_k \mathbf{e}_k - F(\mathbf{x}_k))\)。由半光滑性，\(F(\mathbf{x}^*+\mathbf{e}_k) - F(\mathbf{x}^*) - V_k \mathbf{e}_k = o(\|\mathbf{e}_k\|)\)。因 \(F(\mathbf{x}^*) = 0\)，\(F(\mathbf{x}_k) = V_k \mathbf{e}_k + o(\|\mathbf{e}_k\|)\)。故 \(\|\mathbf{e}_{k+1}\| = \|V_k^{-1} \cdot o(\|\mathbf{e}_k\|)\| \leq \|V_k^{-1}\| \cdot o(\|\mathbf{e}_k\|)\)。由 \(\partial F(\mathbf{x}^*)\) 的非退化性，当 \(\mathbf{x}_k\) 充分接近 \(\mathbf{x}^*\) 时 \(\|V_k^{-1}\|\) 一致有界。故 \(\|\mathbf{e}_{k+1}\| = o(\|\mathbf{e}_k\|)\)，即超线性收敛。若 \(F\) 强半光滑且 \(\partial F\) Lipschitz，余项为 \(O(\|\mathbf{e}_k\|^2)\)，得二次收敛。\(\blacksquare\)

\subsection{137.y 增广Lagrangian方法的收敛理论}\label{y-ux589eux5e7flagrangianux65b9ux6cd5ux7684ux6536ux655bux7406ux8bba}

\textbf{定义 137.3}（增广Lagrangian法的收敛性分析）：考虑等式约束问题 \(\min f(\mathbf{x})\) s.t. \(h_j(\mathbf{x}) = 0\)（\(j = 1, \ldots, p\)）。增广 Lagrangian 函数为

\[L_\rho(\mathbf{x}, \boldsymbol{\nu}) = f(\mathbf{x}) + \sum_{j=1}^p \nu_j h_j(\mathbf{x}) + \frac{\rho}{2} \sum_{j=1}^p h_j(\mathbf{x})^2\]

增广 Lagrangian 法（乘子法）的迭代为：\(\mathbf{x}_{k+1} = \arg\min_{\mathbf{x}} L_{\rho_k}(\mathbf{x}, \boldsymbol{\nu}_k)\)，\(\nu_j^{k+1} = \nu_j^k + \rho_k h_j(\mathbf{x}_{k+1})\)。与经典罚函数法（\(\nu_j^k \equiv 0\) 且 \(\rho_k \to \infty\)）不同，ALM 的乘子更新使得精确解在有限惩罚参数下即可达到。

\textbf{定理 137.2}（增广Lagrangian法的局部线性收敛，Rockafellar 1976；Bertsekas 1982）：设 \(\mathbf{x}^*\) 是等式约束问题的局部极小点，满足 LICQ 和二阶充分条件（\(\nabla_{\mathbf{x}\mathbf{x}}^2 L(\mathbf{x}^*, \boldsymbol{\nu}^*) \succ 0\) 在约束的切空间上）。则存在 \(\bar{\rho} > 0\) 使得对所有 \(\rho_k \geq \bar{\rho}\)，ALM 以线性速率收敛：

\[\|\mathbf{x}_{k+1} - \mathbf{x}^*\| \leq \frac{C}{\rho_k} \|\boldsymbol{\nu}_k - \boldsymbol{\nu}^*\|, \quad \|\boldsymbol{\nu}_{k+1} - \boldsymbol{\nu}^*\| \leq \frac{C}{\rho_k} \|\boldsymbol{\nu}_k - \boldsymbol{\nu}^*\|\]

其中 \(C\) 为依赖于问题数据的常数。特别地，若 \(\rho_k \to \infty\)，收敛是超线性的；若 \(\rho_k\) 固定且充分大，收敛是线性的，速率为 \(O(1/\rho_k)\)。

\textbf{证明}：ALM 的一阶最优性条件：对 \(\mathbf{x}_{k+1}\)，\(\nabla f(\mathbf{x}_{k+1}) + \sum_j (\nu_j^k + \rho_k h_j(\mathbf{x}_{k+1})) \nabla h_j(\mathbf{x}_{k+1}) = 0\)。令 \(\boldsymbol{\nu}_{k+1} = \boldsymbol{\nu}_k + \rho_k \mathbf{h}(\mathbf{x}_{k+1})\)（其中 \(\mathbf{h} = (h_1, \ldots, h_p)^T\)），则 \((\mathbf{x}_{k+1}, \boldsymbol{\nu}_{k+1})\) 满足 \(\nabla f(\mathbf{x}_{k+1}) + \sum_j \nu_j^{k+1} \nabla h_j(\mathbf{x}_{k+1}) = 0\) 和 \(\mathbf{h}(\mathbf{x}_{k+1}) = (\boldsymbol{\nu}_{k+1} - \boldsymbol{\nu}_k)/\rho_k\)。由 LICQ 和二阶充分条件，KKT 映射 \(F(\mathbf{x}, \boldsymbol{\nu}) = (\nabla_{\mathbf{x}} L(\mathbf{x}, \boldsymbol{\nu}), \mathbf{h}(\mathbf{x}))\) 在 \((\mathbf{x}^*, \boldsymbol{\nu}^*)\) 处可逆（隐函数定理）。展开 \(F\) 在 \((\mathbf{x}^*, \boldsymbol{\nu}^*)\) 处的 Taylor 级数并代入 \(\boldsymbol{\nu}_{k+1} - \boldsymbol{\nu}_k = \rho_k \mathbf{h}(\mathbf{x}_{k+1})\)，可得递推不等式 \(\|\boldsymbol{\nu}_{k+1} - \boldsymbol{\nu}^*\| \leq (C/\rho_k) \|\boldsymbol{\nu}_k - \boldsymbol{\nu}^*\| + O(\|\mathbf{x}_k - \mathbf{x}^*\|^2)\)。对不等式约束情形，引入松弛变量转化为等式约束，分析类似，需处理互补松弛条件。\(\blacksquare\)

\textbf{定理 137.3}（ALM与邻近点算法的对偶关系）：增广 Lagrangian 法在乘子空间的迭代等价于应用于对偶函数的邻近点算法（Proximal Point Algorithm, PPA）。具体地，考虑等式约束问题，其对偶函数 \(q(\boldsymbol{\nu}) = \inf_{\mathbf{x}} L(\mathbf{x}, \boldsymbol{\nu})\)（\(L\) 为普通 Lagrange 函数）。ALM 的乘子更新

\[\boldsymbol{\nu}_{k+1} = \boldsymbol{\nu}_k + \rho_k \arg\max_{\mathbf{x}} \left\{-f(\mathbf{x}) - \sum_j \nu_j^k h_j(\mathbf{x}) - \frac{\rho_k}{2} \sum_j h_j(\mathbf{x})^2\right\} \text{ 中的 } \mathbf{h}(\mathbf{x}_{k+1})\]

等价于

\[\boldsymbol{\nu}_{k+1} = \arg\max_{\boldsymbol{\nu}} \left\{ q(\boldsymbol{\nu}) - \frac{1}{2\rho_k} \|\boldsymbol{\nu} - \boldsymbol{\nu}_k\|^2 \right\}\]

（当对偶问题为极大化形式时）。这一关系揭示了 ALM 的深层结构：对偶更新正是以 \(\rho_k\) 为正则化参数的邻近点迭代，从而继承了 PPA 的全局收敛性和鲁棒性。

\textbf{证明}：对偶函数 \(q(\boldsymbol{\nu}) = \inf_{\mathbf{x}} \{f(\mathbf{x}) + \sum_j \nu_j h_j(\mathbf{x})\}\)。增广 Lagrangian 子问题 \(\min_{\mathbf{x}} L_{\rho_k}(\mathbf{x}, \boldsymbol{\nu}_k)\) 的最优值等于 \(\max_{\boldsymbol{\nu}} \{ q(\boldsymbol{\nu}) - \frac{1}{2\rho_k} \|\boldsymbol{\nu} - \boldsymbol{\nu}_k\|^2 \}\)（由 Fenchel 对偶或 Lagrange 对偶的鞍点性质）。事实上，\(\min_{\mathbf{x}} \max_{\boldsymbol{\nu}} [f(\mathbf{x}) + \sum_j \nu_j h_j(\mathbf{x}) - \frac{1}{2\rho_k} \|\boldsymbol{\nu} - \boldsymbol{\nu}_k\|^2] = \max_{\boldsymbol{\nu}} \min_{\mathbf{x}} [\cdots]\)（在凸假设下强对偶成立），内层 \(\min_{\mathbf{x}}\) 恰为 \(q(\boldsymbol{\nu})\)，外层 \(\max_{\boldsymbol{\nu}}\) 的极值点即为 \(\boldsymbol{\nu}_{k+1}\)。同时，由子问题的一阶条件，\(\boldsymbol{\nu}_{k+1} = \boldsymbol{\nu}_k + \rho_k \mathbf{h}(\mathbf{x}_{k+1})\) 恰为此双层优化的解。因此 ALM 的乘子步骤等于以 \(\rho_k^{-1}\) 为步长的 PPA 应用于对偶函数 \(q\)。PPA 的收敛理论（Rockafellar 1976）保证：若 \(q\) 是恰当闭凸函数且 \(\arg\max q \neq \varnothing\)，则 \(\{\boldsymbol{\nu}_k\}\) 收敛到某个对偶最优解（当 \(\rho_k \geq \bar{\rho} > 0\) 有界时）。\(\blacksquare\)

\hrulefill

\hrulefill

\section{第430章：凸优化}\label{ux7b2c430ux7ae0ux51f8ux4f18ux5316}

凸优化是优化理论中最重要的子领域，其核心优势在于：凸问题的局部最优解就是全局最优解。本章从凸集和凸函数的基本概念出发，建立凸优化问题的标准形式及对偶理论（KKT 充要条件），为后续非线性规划和变分法中的最优性条件奠定几何基础。

\subsection{134.1 凸集与凸函数}\label{ux51f8ux96c6ux4e0eux51f8ux51fdux6570}

\textbf{定义 134.1}（凸集）：集合 \(C\) 是\textbf{凸集}，如果对任意 \(\mathbf{x}, \mathbf{y} \in C\) 和 \(\theta \in [0, 1]\)，有 \(\theta \mathbf{x} + (1-\theta)\mathbf{y} \in C\)。

\textbf{定义 134.2}（保凸运算）：凸集的交是凸集；凸集在仿射映射下的像和原像都是凸集；两个凸集的 Minkowski 和是凸集。

\textbf{定义 134.3}（凸函数）：函数 \(f: \mathbb{R}^n \to \mathbb{R}\) 是\textbf{凸函数}，如果定义域 \(\operatorname{dom} f\) 是凸集，且对任意 \(\mathbf{x}, \mathbf{y} \in \operatorname{dom} f\) 和 \(\theta \in [0, 1]\)，

\[f(\theta \mathbf{x} + (1-\theta)\mathbf{y}) \leq \theta f(\mathbf{x}) + (1-\theta) f(\mathbf{y})\]

若严格不等式成立（\(\mathbf{x} \neq \mathbf{y}, \theta \in (0,1)\)），则称为\textbf{严格凸函数}。

\textbf{定理 134.1}（凸函数的判定与无约束极小）：

\textbf{证明}：凸函数的一阶判定：\(f\)凸当且仅当\(f(\mathbf{y})\ge f(\mathbf{x})+\nabla f(\mathbf{x})^T(\mathbf{y}-\mathbf{x})\)对所有\(\mathbf{x},\mathbf{y}\)成立。此由定义\(f(\lambda\mathbf{x}+(1-\lambda)\mathbf{y})\le\lambda f(\mathbf{x})+(1-\lambda)f(\mathbf{y})\)取极限\(\lambda\to 0\)得到。二阶判定：\(f\)凸当且仅当\(\nabla^2 f(\mathbf{x})\succeq 0\)对所有\(\mathbf{x}\)成立。由Taylor展开\(f(\mathbf{y})=f(\mathbf{x})+\nabla f(\mathbf{x})^T(\mathbf{y}-\mathbf{x})+\frac{1}{2}(\mathbf{y}-\mathbf{x})^T\nabla^2 f(\mathbf{\xi})(\mathbf{y}-\mathbf{x})\)，正半定性等价于一阶条件。无约束极小：\(\nabla f(\mathbf{x}^*)=0\)为充要条件（凸函数的一阶条件）。\(\blacksquare\) 1. \textbf{一阶条件}：若 \(f\) 可微，\(f\) 为凸 \(\iff\) \(f(\mathbf{y}) \geq f(\mathbf{x}) + \nabla f(\mathbf{x})^T (\mathbf{y} - \mathbf{x})\) 对所有 \(\mathbf{x}, \mathbf{y}\) 成立 2. \textbf{二阶条件}：若 \(f\) 二阶可微，\(f\) 为凸 \(\iff\) \(\nabla^2 f(\mathbf{x}) \succeq 0\)（Hessian 半正定）对所有 \(\mathbf{x}\) 成立

\emph{无约束极小的充要条件}（次梯度条件）：设 \(f\) 为凸函数，则 \(\mathbf{x}^*\) 是全局极小点当且仅当 \(\mathbf{0} \in \partial f(\mathbf{x}^*)\)。其中 \(\partial f(\mathbf{x}^*) = \{\mathbf{g} : f(\mathbf{y}) \geq f(\mathbf{x}^*) + \mathbf{g}^T(\mathbf{y} - \mathbf{x}^*), \forall \mathbf{y}\}\) 为次微分。若 \(f\) 在 \(\mathbf{x}^*\) 处可微，则此条件退化为 \(\nabla f(\mathbf{x}^*) = \mathbf{0}\)。证明：充分性由次梯度定义直接得 \(f(\mathbf{y}) \geq f(\mathbf{x}^*)\)，必要性由凸分离定理。∎

\textbf{定理 134.2}（Jensen不等式与凸函数的等价刻画）：可测函数 \(f\) 为凸函数当且仅当对任意概率测度 \(\mu\)，有 \(f(\int x\,d\mu(x)) \leq \int f(x)\,d\mu(x)\)。特别地，取两点分布的离散形式即还原为定义 134.3。

\textbf{证明}：(\(\Leftarrow\)) 取 \(\mu\) 为两点分布：\(\mu = t\delta_x + (1-t)\delta_y\)，则得 \(f(tx + (1-t)y) \leq t f(x) + (1-t) f(y)\)，即 \(f\) 为凸函数。(\(\Rightarrow\)) 设 \(f\) 为凸函数。令 \(m = \int x \, d\mu(x)\)。由凸函数的次梯度性质，存在 \(a \in \mathbb{R}\) 使得 \(f(x) \geq f(m) + a(x - m)\) 对所有 \(x\) 成立。对两边关于 \(\mu\) 积分：\(\int f(x) \, d\mu(x) \geq \int (f(m) + a(x-m)) \, d\mu(x) = f(m) + a \int (x-m) \, d\mu(x) = f(m)\)。故 \(f(m) \leq \int f(x) \, d\mu(x)\)，即 Jensen 不等式成立。\(\blacksquare\)

\subsection{134.2 凸优化问题的标准形式}\label{ux51f8ux4f18ux5316ux95eeux9898ux7684ux6807ux51c6ux5f62ux5f0f}

\textbf{定义 134.4}（凸优化问题）：

\[\begin{aligned} \min \quad & f_0(\mathbf{x}) \\ \text{s.t.} \quad & f_i(\mathbf{x}) \leq 0, \quad i = 1, \ldots, m \\ & A\mathbf{x} = \mathbf{b} \end{aligned}\]

其中 \(f_0, f_1, \ldots, f_m\) 是凸函数。可行域是凸集。

\textbf{定理 134.3}（凸优化的基本性质）：凸优化问题的任何局部最优解都是全局最优解。最优解集是凸集。

\textbf{证明}：设 \(\mathbf{x}^*\) 是局部最优解。若存在 \(\mathbf{y}\) 满足 \(f_0(\mathbf{y}) < f_0(\mathbf{x}^*)\)，则对充分小的 \(\theta > 0\)，\(\theta \mathbf{y} + (1-\theta)\mathbf{x}^*\) 在 \(\mathbf{x}^*\) 的邻域内，且 \(f_0(\theta \mathbf{y} + (1-\theta)\mathbf{x}^*) \leq \theta f_0(\mathbf{y}) + (1-\theta)f_0(\mathbf{x}^*) < f_0(\mathbf{x}^*)\)，矛盾。\(\blacksquare\)

\subsection{134.3 Lagrange对偶理论}\label{lagrangeux5bf9ux5076ux7406ux8bba}

\textbf{定义 134.5}（Lagrange函数与对偶函数）：对一般约束优化问题 \(\min\{f_0(\mathbf{x}): f_i(\mathbf{x})\leq 0,\, h_j(\mathbf{x})=0\}\)，定义\textbf{Lagrange函数}

\[L(\mathbf{x},\boldsymbol{\lambda},\boldsymbol{\nu}) = f_0(\mathbf{x}) + \sum_{i=1}^m \lambda_i f_i(\mathbf{x}) + \sum_{j=1}^p \nu_j h_j(\mathbf{x})\]

其中 \(\lambda_i\geq 0\)。\textbf{对偶函数}定义为 \(g(\boldsymbol{\lambda},\boldsymbol{\nu}) = \inf_{\mathbf{x}} L(\mathbf{x},\boldsymbol{\lambda},\boldsymbol{\nu})\)。

\textbf{定理 134.4}（对偶函数的凹性与弱对偶性）：对偶函数 \(g(\boldsymbol{\lambda},\boldsymbol{\nu})\) 是凹函数（与原始问题的凸性无关），且对任意对偶可行点 \((\boldsymbol{\lambda}\succeq 0,\boldsymbol{\nu})\)，有 \(g(\boldsymbol{\lambda},\boldsymbol{\nu}) \leq p^*\)（弱对偶性）。

\textbf{证明}：\(g\) 是仿射函数族的下确界，故为凹。对任意原始可行点 \(\mathbf{x}\)，\(L(\mathbf{x},\boldsymbol{\lambda},\boldsymbol{\nu})=f_0(\mathbf{x})+\sum\lambda_i f_i(\mathbf{x})+\sum\nu_j h_j(\mathbf{x})\leq f_0(\mathbf{x})\)（因 \(\lambda_i\geq 0, f_i(\mathbf{x})\leq 0, h_j(\mathbf{x})=0\)）。取下确界即得 \(g(\boldsymbol{\lambda},\boldsymbol{\nu})\leq p^*\)。\(\blacksquare\)

\subsection{134.4 强对偶与KKT条件}\label{ux5f3aux5bf9ux5076ux4e0ekktux6761ux4ef6}

\textbf{定义 134.6}（对偶间隙与鞍点）：\(p^*-d^*\)（其中 \(d^*=\sup g\)）称为\textbf{对偶间隙}。若 \(p^*=d^*\)，称为\textbf{强对偶}成立。点 \((\mathbf{x}^*,\boldsymbol{\lambda}^*,\boldsymbol{\nu}^*)\) 称为 \(L\) 的\textbf{鞍点}，若 \(L(\mathbf{x}^*,\boldsymbol{\lambda},\boldsymbol{\nu})\leq L(\mathbf{x}^*,\boldsymbol{\lambda}^*,\boldsymbol{\nu}^*)\leq L(\mathbf{x},\boldsymbol{\lambda}^*,\boldsymbol{\nu}^*)\) 对所有 \(\mathbf{x}\) 和 \(\boldsymbol{\lambda}\succeq 0\) 成立。

\textbf{定理 134.5}（Slater强对偶定理）：设凸优化问题（\(f_i\) 凸，\(h_j\) 仿射）存在严格可行点 \(\mathbf{x}_0\) 满足 \(f_i(\mathbf{x}_0)<0\)（\(i=1,\ldots,m\)）且 \(A\mathbf{x}_0=\mathbf{b}\)。则强对偶成立：\(p^*=d^*\)，且对偶最优值可达。

\textbf{证明}：定义集合 \(\mathcal{A}=\{(u,v,t)\in\mathbb{R}^m\times\mathbb{R}^p\times\mathbb{R}: \exists\mathbf{x},\, f_i(\mathbf{x})\leq u_i,\, h_j(\mathbf{x})=v_j,\, f_0(\mathbf{x})\leq t\}\)。\(\mathcal{A}\) 为凸集且 \(p^*=\inf\{t: (0,0,t)\in\mathcal{A}\}\)。由凸分离定理，存在非零超平面分离 \((0,0,p^*)\) 与 \(\mathcal{A}\) 的内部，经适当构造即得对偶变量。\(\blacksquare\)

\textbf{定理 134.6}（凸优化的KKT充要条件）：对满足Slater条件的凸优化问题，\(\mathbf{x}^*\) 是最优解当且仅当存在 \(\boldsymbol{\lambda}^*\succeq 0\) 和 \(\boldsymbol{\nu}^*\) 满足KKT条件：

\[\begin{aligned}
\nabla f_0(\mathbf{x}^*) + \sum_{i=1}^m \lambda_i^* \nabla f_i(\mathbf{x}^*) + A^T\boldsymbol{\nu}^* &= \mathbf{0} \quad &\text{(稳定性)} \\
f_i(\mathbf{x}^*) &\leq 0, \; A\mathbf{x}^*=\mathbf{b} \quad &\text{(原可行性)} \\
\lambda_i^* f_i(\mathbf{x}^*) &= 0 \quad &\text{(互补松弛性)}
\end{aligned}\]

\textbf{证明}：（\(\Rightarrow\)）若 \(\mathbf{x}^*\) 最优，由强对偶存在 \((\boldsymbol{\lambda}^*,\boldsymbol{\nu}^*)\) 使 \(g(\boldsymbol{\lambda}^*,\boldsymbol{\nu}^*)=f_0(\mathbf{x}^*)\)，即 \(\mathbf{x}^*\in\arg\min L(\cdot,\boldsymbol{\lambda}^*,\boldsymbol{\nu}^*)\)，从而 \(\nabla_{\mathbf{x}}L=0\)（稳定性），互补松弛性由等式成立 \(f_0(\mathbf{x}^*)=L(\mathbf{x}^*,\boldsymbol{\lambda}^*,\boldsymbol{\nu}^*)\) 得。（\(\Leftarrow\)）对任意可行 \(\mathbf{x}\)，由凸性 \(f_0(\mathbf{x})\geq L(\mathbf{x},\boldsymbol{\lambda}^*,\boldsymbol{\nu}^*)\geq L(\mathbf{x}^*,\boldsymbol{\lambda}^*,\boldsymbol{\nu}^*)=f_0(\mathbf{x}^*)\)。\(\blacksquare\)

\subsection{134.5 梯度下降法的收敛性分析}\label{ux68afux5ea6ux4e0bux964dux6cd5ux7684ux6536ux655bux6027ux5206ux6790}

\textbf{定理 134.7}（梯度下降的 \(O(1/k)\) 收敛率）：设 \(f\) 为凸函数且 \(\nabla f\) 是 \(L\)-Lipschitz连续的：\(\|\nabla f(\mathbf{x})-\nabla f(\mathbf{y})\|\leq L\|\mathbf{x}-\mathbf{y}\|\)。取固定步长 \(\alpha=1/L\)，迭代 \(\mathbf{x}_{k+1}=\mathbf{x}_k-\frac{1}{L}\nabla f(\mathbf{x}_k)\)。则

\[\min_{0\leq i\leq k} \|\nabla f(\mathbf{x}_i)\|^2 \leq \frac{2L(f(\mathbf{x}_0)-f^*)}{k+1}\]

\textbf{证明}：由 \(L\)-光滑性得 \(f(\mathbf{x}_{k+1})\leq f(\mathbf{x}_k)-\frac{1}{2L}\|\nabla f(\mathbf{x}_k)\|^2\)。对 \(i=0,\ldots,k\) 求和：

\[f(\mathbf{x}_{k+1}) \leq f(\mathbf{x}_0) - \frac{1}{2L}\sum_{i=0}^k \|\nabla f(\mathbf{x}_i)\|^2\]

移项并取最小值界：\(\min_{i\leq k}\|\nabla f(\mathbf{x}_i)\|^2\cdot (k+1)\leq\sum_{i=0}^k\|\nabla f(\mathbf{x}_i)\|^2\leq 2L(f(\mathbf{x}_0)-f^*)\)。\(\blacksquare\)

梯度下降的收敛性分析是凸优化算法的理论基础。\(L\)-光滑性（即梯度是 \(L\)-Lipschitz 连续的）是保证梯度下降收敛的关键条件，它给出了 \(f\) 的二阶 Taylor 展开的全局上界：\(f(\mathbf{y}) \leq f(\mathbf{x}) + \nabla f(\mathbf{x})^T(\mathbf{y}-\mathbf{x}) + \frac{L}{2}\|\mathbf{y}-\mathbf{x}\|^2\)。取 \(\mathbf{y} = \mathbf{x} - \frac{1}{L}\nabla f(\mathbf{x})\) 代入即得每步下降量：\(f(\mathbf{x}_{k+1}) \leq f(\mathbf{x}_k) - \frac{1}{2L}\|\nabla f(\mathbf{x}_k)\|^2\)。这一分析表明梯度下降的收敛速率为 \(O(1/k)\)（次线性收敛）。当 \(f\) 还满足 \(\mu\)-强凸性（\(\nabla^2 f \succeq \mu I\)）时，收敛速率可改进为线性收敛：\(f(\mathbf{x}_k) - f^* \leq \left(1 - \frac{\mu}{L}\right)^k (f(\mathbf{x}_0) - f^*)\)，即每次迭代以固定比例缩小最优性间隙。条件数 \(\kappa = L/\mu\) 决定了收敛速度：\(\kappa\) 越大（函数越''病态''），收敛越慢。Nesterov 加速梯度下降法（1983年）通过引入动量项，将凸光滑函数的收敛速率从 \(O(1/k)\) 提升至 \(O(1/k^2)\)，这是在不增加每次迭代计算量的前提下，一阶方法能达到的最优收敛速率（信息论下界）。加速方法的迭代为 \(\mathbf{y}_k = \mathbf{x}_k + \frac{k-1}{k+2}(\mathbf{x}_k - \mathbf{x}_{k-1})\)，\(\mathbf{x}_{k+1} = \mathbf{y}_k - \frac{1}{L}\nabla f(\mathbf{y}_k)\)。

\subsection{134.6 Newton法的二阶收敛性}\label{newtonux6cd5ux7684ux4e8cux9636ux6536ux655bux6027}

\textbf{定理 134.8}（Newton法的局部二阶收敛）：设 \(f\) 为二阶连续可微且 \(\nabla^2 f\) 在 \(\mathbf{x}^*\) 附近 \(M\)-Lipschitz连续：\(\|\nabla^2 f(\mathbf{x})-\nabla^2 f(\mathbf{y})\|\leq M\|\mathbf{x}-\mathbf{y}\|\)。若 \(\nabla^2 f(\mathbf{x}^*)\succ 0\)，则存在 \(\mathbf{x}^*\) 的邻域使得Newton迭代 \(\mathbf{x}_{k+1}=\mathbf{x}_k-[\nabla^2 f(\mathbf{x}_k)]^{-1}\nabla f(\mathbf{x}_k)\) 满足

\[\|\mathbf{x}_{k+1}-\mathbf{x}^*\| \leq C\,\|\mathbf{x}_k-\mathbf{x}^*\|^2\]

其中 \(C=M/(2\lambda_{\min}(\nabla^2 f(\mathbf{x}^*)))\)，收敛为二阶（二次）。

\textbf{证明}：由Newton步定义 \(\mathbf{x}_{k+1}-\mathbf{x}^*=\mathbf{x}_k-\mathbf{x}^*-[\nabla^2 f(\mathbf{x}_k)]^{-1}(\nabla f(\mathbf{x}_k)-\nabla f(\mathbf{x}^*))\)。利用积分形式的Taylor展开和Hessian的Lipschitz连续性估计范数得上界。\(\blacksquare\)

Newton 法的二阶收敛性是其区别于梯度下降法的最显著特征。梯度下降仅利用一阶导数信息，收敛速率为线性（强凸情形）或次线性（一般凸情形）；而 Newton 法利用二阶导数（Hessian 矩阵）信息，在解附近实现二阶收敛------每次迭代的正确位数翻倍。这一优势的代价是：每次迭代需要计算和求逆 Hessian 矩阵 \(\nabla^2 f(\mathbf{x}_k)\)，复杂度为 \(O(n^3)\)（\(n\) 为变量维数），而梯度下降每次迭代仅需 \(O(n)\)。Newton 法的收敛性分析基于 Newton 减量 \(\lambda(\mathbf{x}) = \sqrt{\nabla f(\mathbf{x})^T [\nabla^2 f(\mathbf{x})]^{-1} \nabla f(\mathbf{x})}\)，它度量了 \(\mathbf{x}\) 到最优解的距离。收敛分为两个阶段：\textbf{阻尼 Newton 阶段}（damped Newton phase）：当 \(\lambda(\mathbf{x}) \geq \eta\)（\(\eta\) 为阈值，通常取 \(\eta = (1-2\alpha)/(M^2)\) 等合适值），使用回溯线搜索（backtracking line search）选择步长，保证每次迭代目标函数值充分下降；\textbf{纯 Newton 阶段}（pure Newton phase）：当 \(\lambda(\mathbf{x}) < \eta\) 时，步长 \(t=1\) 满足 Wolfe 条件，迭代进入二阶收敛区域。阻尼 Newton 阶段需要 \(O(\log(1/\eta))\) 次迭代，而纯 Newton 阶段仅需 \(O(\log\log(1/\varepsilon))\) 次迭代即可达到 \(\varepsilon\)-精度------这是梯度下降法无法企及的。在实际应用中，拟 Newton 法（如 BFGS、L-BFGS）通过累积梯度信息逐步逼近 Hessian 逆，以 \(O(n^2)\) 的每次迭代代价换取接近超线性收敛的速率，是大规模优化中最重要的算法之一。

\subsection{134.7 内点法的基本思想}\label{ux5185ux70b9ux6cd5ux7684ux57faux672cux601dux60f3}

\textbf{定义 134.7}（对数障碍函数与中心路径）：对不等式约束 \(f_i(\mathbf{x})\leq 0\)，引入\textbf{对数障碍函数}

\[\phi(\mathbf{x}) = -\sum_{i=1}^m \log(-f_i(\mathbf{x}))\]

（定义域为 \(\{\mathbf{x}: f_i(\mathbf{x})<0\}\)）。对 \(t>0\)，求解\textbf{障碍问题}

\[\min_{\mathbf{x}} \; t f_0(\mathbf{x}) + \phi(\mathbf{x}) \quad \text{s.t.} \quad A\mathbf{x}=\mathbf{b}\]

记最优解为 \(\mathbf{x}^*(t)\)。当 \(t\to\infty\) 时，\(\mathbf{x}^*(t)\) 沿\textbf{中心路径}收敛至原问题的最优解。内点法通过在增大 \(t\) 的同时沿中心路径追踪 \(\mathbf{x}^*(t)\)（通常用Newton法求解每个障碍子问题），实现多项式时间复杂度。对线性规划，路径跟踪内点法达到 \(O(\sqrt{m}\log(1/\varepsilon))\) 迭代复杂度。

\textbf{定理 134.9}（中心路径的最优性）：设 \(\mathbf{x}^*(t)\) 为障碍问题的唯一解（由凸性保证唯一性），\(\boldsymbol{\lambda}_i^*(t)=-1/(t f_i(\mathbf{x}^*(t)))\)，\(\boldsymbol{\nu}^*(t)\) 为等式约束的乘子。则对偶间隙满足 \(f_0(\mathbf{x}^*(t))-g(\boldsymbol{\lambda}^*(t),\boldsymbol{\nu}^*(t))=m/t\)，从而当 \(t\to\infty\) 时 \(\mathbf{x}^*(t)\) 趋近于原问题最优解。

\textbf{证明}：障碍问题\(\min_{\mathbf{x}} tf_0(\mathbf{x})+\sum_{i=1}^m(-\log(-f_i(\mathbf{x})))\)的KKT条件：\(t\nabla f_0(\mathbf{x}^*)+\sum_i\frac{1}{-f_i(\mathbf{x}^*)}\nabla f_i(\mathbf{x}^*)+A^T\boldsymbol{\nu}^*=0\)，\(f_i(\mathbf{x}^*)<0\)，\(A\mathbf{x}^*=\mathbf{b}\)。定义\(\lambda_i^*=-1/(tf_i(\mathbf{x}^*))>0\)，则\(\nabla f_0(\mathbf{x}^*)+\sum_i\lambda_i^*\nabla f_i(\mathbf{x}^*)+A^T\boldsymbol{\nu}^*/t=0\)。对偶函数\(g(\boldsymbol{\lambda}^*,\boldsymbol{\nu}^*/t)=\inf_{\mathbf{x}}L(\mathbf{x},\boldsymbol{\lambda}^*,\boldsymbol{\nu}^*/t)\)在\(\mathbf{x}^*\)处达到（由凸性），故\(g(\boldsymbol{\lambda}^*,\boldsymbol{\nu}^*/t)=f_0(\mathbf{x}^*)+\sum_i\lambda_i^*f_i(\mathbf{x}^*)+(\boldsymbol{\nu}^*/t)^T(A\mathbf{x}^*-\mathbf{b})=f_0(\mathbf{x}^*)-m/t\)。对偶间隙\(f_0(\mathbf{x}^*)-g=m/t\to 0\)（\(t\to\infty\)）。\(\blacksquare\)

\newpage

\chapter{12-最优化与控制/01-最优化/04-变分法与半定规划.md}\label{ux6700ux4f18ux5316ux4e0eux63a7ux523601-ux6700ux4f18ux531604-ux53d8ux5206ux6cd5ux4e0eux534aux5b9aux89c4ux5212.md}

\section{第431章：变分法与最优控制引论}\label{ux7b2c431ux7ae0ux53d8ux5206ux6cd5ux4e0eux6700ux4f18ux63a7ux5236ux5f15ux8bba}

变分法研究泛函（函数的函数）的极值问题，起源于 Bernoulli 的最速降线问题（1696年）和 Euler-Lagrange 的工作。最优控制是变分法的现代推广，由 Pontryagin 极大值原理和 Bellman 动态规划奠定。

\subsection{136.1 Euler-Lagrange 方程}\label{euler-lagrange-ux65b9ux7a0b}

\textbf{定义 136.1}（最简泛函）：考虑泛函

\[J[y] = \int_a^b F(x, y(x), y'(x)) \, dx\]

其中 \(y(a) = \alpha\)，\(y(b) = \beta\) 固定。

\textbf{定理 136.1}（Euler-Lagrange 方程）：若 \(y(x)\) 使 \(J[y]\) 取极值，且 \(F\) 充分光滑，则 \(y\) 满足

\[\frac{\partial F}{\partial y} - \frac{d}{dx}\left( \frac{\partial F}{\partial y'} \right) = 0\]

\textbf{证明}：考虑变分 \(\tilde{y} = y + \varepsilon \eta(x)\)，其中 \(\eta(a) = \eta(b) = 0\)。计算一阶变分： \[\frac{d}{d\varepsilon} J[y + \varepsilon \eta]\bigg|_{\varepsilon=0} = \int_a^b \left( \frac{\partial F}{\partial y} \eta + \frac{\partial F}{\partial y'} \eta' \right) dx = 0\] 对第二项分部积分：\(\int_a^b \frac{\partial F}{\partial y'} \eta' dx = \left[ \frac{\partial F}{\partial y'} \eta \right]_a^b - \int_a^b \frac{d}{dx}\left( \frac{\partial F}{\partial y'} \right) \eta dx\)。由 \(\eta(a) = \eta(b) = 0\)，边界项为零，代入得： \[\int_a^b \left( \frac{\partial F}{\partial y} - \frac{d}{dx} \frac{\partial F}{\partial y'} \right) \eta(x) dx = 0\] 由变分法基本引理（\(\eta\) 任意性），被积函数恒为零，即得 Euler-Lagrange 方程。∎

\textbf{定义 136.2}（Beltrami 恒等式）：若 \(F\) 不显含 \(x\)（即 \(\frac{\partial F}{\partial x} = 0\)），则 Euler-Lagrange 方程有首次积分

\[F - y' \frac{\partial F}{\partial y'} = C \quad (\text{常数})\]

\subsection{136.2 多变量与高阶变分}\label{ux591aux53d8ux91cfux4e0eux9ad8ux9636ux53d8ux5206}

\textbf{定义 136.3}（多变量泛函）：对 \(J[u] = \iint_\Omega F(x, y, u, u_x, u_y) \, dx dy\)，Euler-Lagrange 方程为

\[\frac{\partial F}{\partial u} - \frac{\partial}{\partial x}\left( \frac{\partial F}{\partial u_x} \right) - \frac{\partial}{\partial y}\left( \frac{\partial F}{\partial u_y} \right) = 0\]

\textbf{定理 136.2}（Hamilton 原理 / 最小作用量原理）：经典力学系统的运动方程可由作用量泛函 \(S[q] = \int_{t_1}^{t_2} L(q, \dot{q}, t) \, dt\) 的变分 \(\delta S = 0\) 导出，其中 \(L = T - V\) 为 Lagrange 函数。

\textbf{证明}：计算\(\delta S = \int_{t_1}^{t_2} (\frac{\partial L}{\partial q}\delta q + \frac{\partial L}{\partial \dot{q}}\delta\dot{q})dt\)。对第二项分部积分：\(\int \frac{\partial L}{\partial \dot{q}}\delta\dot{q}dt = [\frac{\partial L}{\partial\dot{q}}\delta q]_{t_1}^{t_2} - \int \frac{d}{dt}\frac{\partial L}{\partial\dot{q}}\delta q dt\)。固定端点\(\delta q(t_1)=\delta q(t_2)=0\)，边界项为零。\(\delta S = \int (\frac{\partial L}{\partial q} - \frac{d}{dt}\frac{\partial L}{\partial\dot{q}})\delta q dt = 0\)。由\(\delta q\)任意性，\(\frac{\partial L}{\partial q} - \frac{d}{dt}\frac{\partial L}{\partial\dot{q}} = 0\)即Euler-Lagrange方程，等价于Newton运动方程。\(\blacksquare\)

\textbf{定义 136.4}（二次变分与 Legendre 条件）：极小的二阶必要条件为 Legendre 条件 \(\frac{\partial^2 F}{\partial y'^2} \geq 0\)。充分条件为强化的 Legendre 条件 \(\frac{\partial^2 F}{\partial y'^2} > 0\) 加上 Jacobi 条件（无共轭点）。

\subsection{136.3 约束变分问题与等周问题}\label{ux7ea6ux675fux53d8ux5206ux95eeux9898ux4e0eux7b49ux5468ux95eeux9898}

\textbf{定义 136.5}（等周问题 / Isoperimetric Problem）：在积分约束 \(\int_a^b G(x, y, y') \, dx = L\)（常数）下，最大化或最小化 \(\int_a^b F(x, y, y') \, dx\)。

等周问题是最古老的变分问题之一，其名称源于经典的 Dido 问题：给定固定长度的围栏，围出最大面积。在数学上，这对应于在约束 \(\int_a^b \sqrt{1 + (y')^2} dx = L\)（弧长固定）下最大化 \(\int_a^b y \, dx\)（面积）。Lagrange 乘子法在无穷维空间中的推广是处理约束变分问题的核心工具：将约束泛函以 Lagrange 乘子 \(\lambda\) 并入目标泛函，形成无约束泛函 \(\tilde{J}[y] = \int_a^b (F + \lambda G) dx\)，然后应用 Euler-Lagrange 方程。乘子 \(\lambda\) 的物理意义因问题而异：在等周问题中，\(\lambda\) 与曲线的曲率相关；在 Sturm-Liouville 理论中，\(\lambda\) 是微分算子的本征值。更一般的约束变分问题包括\textbf{非完整约束}（nonholonomic constraints）\(\phi(x, y, y') = 0\)，此时 Euler-Lagrange 方程变为 \(\frac{\partial F}{\partial y} - \frac{d}{dx}\frac{\partial F}{\partial y'} + \mu(x)\left(\frac{\partial\phi}{\partial y} - \frac{d}{dx}\frac{\partial\phi}{\partial y'}\right) - \mu'(x)\frac{\partial\phi}{\partial y'} = 0\)，其中 \(\mu(x)\) 是乘子函数。等周问题和约束变分问题的理论在最优控制中获得了自然的推广------Pontryagin 极大值原理本质上就是带微分约束的变分问题。

\textbf{定理 136.3}（等周问题的 Euler-Lagrange 方程）：定义辅助泛函 \(H = F + \lambda G\)（\(\lambda\) 为 Lagrange 乘子），则极值曲线满足

\[\frac{\partial H}{\partial y} - \frac{d}{dx}\left( \frac{\partial H}{\partial y'} \right) = 0\]

\textbf{证明}：考虑约束\(\int_a^b G(x,y,y')dx = L\)。引入Lagrange乘子\(\lambda\)，构造无约束泛函\(\tilde{J}[y] = \int_a^b (F+\lambda G)dx = \int_a^b H(x,y,y')dx\)。对\(\tilde{J}\)应用定理136.1的Euler-Lagrange方程，直接得到\(\frac{\partial H}{\partial y} - \frac{d}{dx}\frac{\partial H}{\partial y'} = 0\)。乘子\(\lambda\)由约束条件\(\int_a^b G dx = L\)确定。\(\blacksquare\)

\subsection{136.4 最优控制初步}\label{ux6700ux4f18ux63a7ux5236ux521dux6b65}

\textbf{定义 136.6}（最优控制问题）：

\[\begin{aligned} \min_{u(t)} \quad & \phi(x(T)) + \int_0^T L(x(t), u(t), t) \, dt \\ \text{s.t.} \quad & \dot{x}(t) = f(x(t), u(t), t), \quad x(0) = x_0 \end{aligned}\]

其中 \(x(t) \in \mathbb{R}^n\) 为状态变量，\(u(t) \in \mathbb{R}^m\) 为控制变量。

\textbf{定理 136.4}（Pontryagin 极大值原理，1962）：定义 Hamilton 函数 \(H(x, u, p, t) = L(x, u, t) + p^T f(x, u, t)\)，其中 \(p(t)\) 是协态变量（adjoint variable）。最优控制 \(u^*(t)\) 和最优轨线 \(x^*(t)\) 满足：

\[\begin{aligned}
\dot{x}^* &= \frac{\partial H}{\partial p} = f(x^*, u^*, t) \\
\dot{p} &= -\frac{\partial H}{\partial x} \\
H(x^*, u^*, p, t) &= \max_{u \in \mathcal{U}} H(x^*, u, p, t) \quad \text{(极大值条件)} \\
p(T) &= \nabla \phi(x^*(T)) \quad \text{(横截条件)}
\end{aligned}\]

\textbf{证明}：考虑控制变量\(u\)的变分\(\delta u\)（允许针状变分）。由状态方程\(\dot{x}=f\)的变分\(\delta\dot{x}=\nabla_x f\cdot\delta x+\nabla_u f\cdot\delta u\)。引入协态\(p\)，构造Hamilton函数\(H\)。目标泛函的变分为\(\delta J = \nabla\phi\cdot\delta x(T) + \int_0^T (\nabla_x L\cdot\delta x+\nabla_u L\cdot\delta u)dt\)。分部积分并利用\(p\)满足的伴随方程\(\dot{p}=-\nabla_x H\)消去\(\delta x\)的积分项，得\(\delta J = \int_0^T \nabla_u H\cdot\delta u dt\)。最优性要求\(\delta J\ge 0\)对所有容许\(\delta u\)成立，等价于\(H(u^*)\ge H(u)\)对所有\(u\in\mathcal{U}\)（极大值条件）。横截条件\(p(T)=\nabla\phi\)来自变分边界项。\(\blacksquare\)

\textbf{定理 136.5}（Bellman 最优性原理与 HJB 方程）：定义值函数 \(V(x, t) = \min_{u(\cdot)} \left[ \phi(x(T)) + \int_t^T L(x, u, \tau) d\tau \right]\)。则 \(V\) 满足 Hamilton-Jacobi-Bellman（HJB）方程：

\[-\frac{\partial V}{\partial t} = \min_{u \in \mathcal{U}} \left[ L(x, u, t) + \nabla_x V \cdot f(x, u, t) \right]\]

边界条件 \(V(x, T) = \phi(x)\)。

\textbf{证明}：由Bellman最优性原理，\(V(x,t)=\min_u[L(x,u,t)\Delta t + V(x+f(x,u,t)\Delta t, t+\Delta t)] + o(\Delta t)\)。Taylor展开\(V(x+f\Delta t, t+\Delta t) = V(x,t) + \nabla_x V\cdot f\Delta t + \frac{\partial V}{\partial t}\Delta t + o(\Delta t)\)。代入得\(V(x,t) = \min_u[L\Delta t + V(x,t) + \nabla_x V\cdot f\Delta t + \frac{\partial V}{\partial t}\Delta t] + o(\Delta t)\)。消去\(V(x,t)\)，除以\(\Delta t\)取极限\(\Delta t\to 0\)，得\(-\frac{\partial V}{\partial t} = \min_u[L + \nabla_x V\cdot f]\)。边界条件\(V(x,T)=\phi(x)\)由定义直接给出。\(\blacksquare\)

\textbf{定义 136.7}（线性二次型调节器 / LQR）：当系统为线性 \(\dot{x} = Ax + Bu\)，目标为二次型 \(J = \frac{1}{2}x(T)^T Q_f x(T) + \frac{1}{2}\int_0^T (x^T Q x + u^T R u) dt\) 时，最优控制为状态反馈 \(u^* = -R^{-1} B^T P(t) x\)，其中 \(P(t)\) 满足 Riccati 微分方程 \(-\dot{P} = A^T P + PA - PBR^{-1}B^T P + Q\)。

\hrulefill

\hrulefill

\hrulefill

\hrulefill

\section{第432章：互补性问题与变分不等式}\label{ux7b2c432ux7ae0ux4e92ux8865ux6027ux95eeux9898ux4e0eux53d8ux5206ux4e0dux7b49ux5f0f}

互补性问题是优化理论中的一个重要分支，其核心是线性互补问题\(\operatorname{LCP}(q, M)\)：给定向量\(q \in \mathbb{R}^n\)和矩阵\(M \in \mathbb{R}^{n \times n}\)，求\(z \in \mathbb{R}^n\)使得\(z \geq 0\)，\(Mz + q \geq 0\)，且\(z^T(Mz + q) = 0\)。\(\operatorname{LCP}\)与多个数学领域紧密关联：二次规划\(\min\{\frac{1}{2}x^TQx + c^Tx : Ax \geq b, x \geq 0\}\)的KKT条件即为\(\operatorname{LCP}\)；双矩阵博弈的Nash均衡也可化为\(\operatorname{LCP}\)。求解\(\operatorname{LCP}\)的经典算法是Lemke算法（1965），通过引入人工变量\(z_0\)构造几乎互补基，沿射线进行主元迭代直到\(z_0\)出基。Lemke算法在\(M\)为P-矩阵（所有主子式为正）时保证收敛。变分不等式\(\operatorname{VI}(K, F)\)求\(x^* \in K\)使\(\langle F(x^*), x - x^* \rangle \geq 0\)对所有\(x \in K\)成立，是\(\operatorname{LCP}\)的无穷维推广，在均衡分析和交通分配中有广泛应用。

\subsection{线性互补问题的严格定义与矩阵类}\label{ux7ebfux6027ux4e92ux8865ux95eeux9898ux7684ux4e25ux683cux5b9aux4e49ux4e0eux77e9ux9635ux7c7b}

\textbf{定义139.1}（线性互补问题 \(\operatorname{LCP}\)）：给定矩阵 \(M \in \mathbb{R}^{n \times n}\) 和向量 \(q \in \mathbb{R}^n\)，线性互补问题 \(\operatorname{LCP}(M, q)\) 是求 \(z \in \mathbb{R}^n\) 满足 \[z \geq 0, \quad w = Mz + q \geq 0, \quad z_i w_i = 0 \quad (i = 1, \ldots, n)\] 其中 \(z_i w_i = 0\) 称为\textbf{互补性条件}，要求对每个 \(i\)，\(z_i\) 和 \(w_i\) 至少有一个为零。记 \(\operatorname{SOL}(M, q)\) 为所有解构成的集合。当 \(q \geq 0\) 时 \(z = 0\) 是平凡解。

\(\operatorname{LCP}\) 的几何意义可通过\textbf{互补锥}理解。令 \(A = [I, -M]\)，则 \(\operatorname{LCP}\) 等价于将向量 \(-q\) 表示为 \(A\) 的列的非负线性组合，且 \(A\) 的每一对对应的列（\(e_j\) 与 \(-M_j\)）不同时使用。这引出了矩阵类的分类理论，这些类刻画了 \(\operatorname{LCP}\) 的可解性。

\textbf{定义139.2}（P-矩阵）：\(M \in \mathbb{R}^{n \times n}\) 称为 \textbf{P-矩阵}，若 \(M\) 的所有主子式（principal minors）均为正。等价地，对所有非零 \(x \in \mathbb{R}^n\)，存在指标 \(i\) 使 \(x_i (Mx)_i > 0\)。

\textbf{定义139.3}（M-矩阵）：\(M \in \mathbb{R}^{n \times n}\) 称为 \textbf{M-矩阵}，若 \(M\) 可表为 \(M = sI - B\)，其中 \(B \geq 0\)（所有元素非负）且 \(s > \rho(B)\)（\(\rho(B)\) 为 \(B\) 的谱半径）。等价地，\(M\) 的所有非对角元非正（\(M_{ij} \leq 0\) 对 \(i \neq j\)），所有主子式为正，且 \(M^{-1} \geq 0\)（元素非负）。M-矩阵在 \(\operatorname{LCP}\) 中特别重要，因为此时 \(\operatorname{SOL}(M, q)\) 的结构与有限差分法的离散最大值原理直接相关。

\textbf{定义139.4}（其他重要矩阵类）：(i) \textbf{正定矩阵}（PD）：\(x^T M x > 0\) 对所有非零 \(x\)，此时 \(\operatorname{LCP}\) 有唯一解；(ii) \textbf{半正定矩阵}（PSD）：\(x^T M x \geq 0\) 对所有 \(x\)，此时 \(\operatorname{SOL}(M, q)\) 若存在则为凸集；(iii) \textbf{P\_0-矩阵}：所有主子式非负；(iv) \textbf{Q-矩阵}：对任意 \(q\)，\(\operatorname{LCP}(M, q)\) 有解。P-矩阵 \(\subset\) Q-矩阵，且包含关系是严格的。P-矩阵类与Q-矩阵类之间的关系是 \(\operatorname{LCP}\) 理论中最深刻的问题之一。

\textbf{命题139.1}（PD \(\Rightarrow\) P \(\Rightarrow\) P\_0）：若 \(M\) 正定，则 \(M\) 的所有主子式为正（P-矩阵）。若 \(M\) 是 P-矩阵，则 \(M\) 的所有主子式非负（P\_0-矩阵）。均为严格包含关系。

\subsection{变分不等式与互补问题的等价关系}\label{ux53d8ux5206ux4e0dux7b49ux5f0fux4e0eux4e92ux8865ux95eeux9898ux7684ux7b49ux4ef7ux5173ux7cfb}

\textbf{定义139.5}（变分不等式 \(\operatorname{VI}\)）：设 \(K \subseteq \mathbb{R}^n\) 为非空闭凸集，\(F : K \to \mathbb{R}^n\) 为连续映射。变分不等式问题 \(\operatorname{VI}(K, F)\) 是求 \(x^* \in K\) 使得 \[\langle F(x^*), x - x^* \rangle \geq 0, \quad \forall x \in K\]

\textbf{定理139.2}（\(\operatorname{LCP}\) 与 \(\operatorname{VI}\) 的等价性）：取 \(K = \mathbb{R}^n_+\)（非负卦限）和 \(F(z) = Mz + q\)，则 \(\operatorname{VI}(\mathbb{R}^n_+, F)\) 等价于 \(\operatorname{LCP}(M, q)\)。

\textbf{证明}：设 \(z^* \in \mathbb{R}^n_+\) 是 \(\operatorname{VI}\) 的解。取 \(x = z^* + e_i\)（第 \(i\) 个标准基向量），则 \(\langle Mz^* + q, e_i \rangle \geq 0\)，故 \((Mz^* + q)_i \geq 0\)，对所有 \(i\)，即 \(w^* = Mz^* + q \geq 0\)。取 \(x = 0\)，得 \(\langle Mz^* + q, -z^* \rangle \geq 0\)，即 \(\langle z^*, Mz^* + q \rangle \leq 0\)。结合 \(z^* \geq 0\) 和 \(Mz^* + q \geq 0\)，其分量内积非负，推得 \(\langle z^*, Mz^* + q \rangle = 0\)。由分量非负，等价的 \(z_i^*(Mz^* + q)_i = 0\) 对每个 \(i\) 成立，即互补性条件。反之，若 \(z^*\) 满足 \(\operatorname{LCP}\)，则对任意 \(x \in \mathbb{R}^n_+\)，\(\langle Mz^* + q, x - z^* \rangle = \langle Mz^* + q, x \rangle - \langle Mz^* + q, z^* \rangle = \langle Mz^* + q, x \rangle \geq 0\)（因 \(z^*\) 满足互补性使第二项为零，第一项因 \(x \geq 0\) 和 \(Mz^* + q \geq 0\) 非负）。故 \(z^*\) 是 \(\operatorname{VI}\) 的解。\(\blacksquare\)

更一般地，对任意闭凸锥 \(K\) 上的 \(\operatorname{VI}\)，可通过 \(K\) 的对偶锥 \(K^* = \{y : \langle y, x \rangle \geq 0, \forall x \in K\}\) 转化为\textbf{广义互补问题}：求 \(x \in K\) 使 \(F(x) \in K^*\) 且 \(\langle F(x), x \rangle = 0\)。这是 \(\operatorname{LCP}\) 在非负卦限外的自然推广。

\subsection{存在性定理}\label{ux5b58ux5728ux6027ux5b9aux7406}

\textbf{定理139.3}（Cottle-Dantzig 定理，1968）：若 \(M\) 是 P-矩阵，则对任意 \(q \in \mathbb{R}^n\)，\(\operatorname{LCP}(M, q)\) 有唯一解。

\textbf{证明}：\textbf{唯一性}：设 \(z^1, z^2\) 均为解。令 \(w^j = Mz^j + q\)。若 \(z^1 \neq z^2\)，考虑指标集。由 P-矩阵的等价刻画，对 \(x = z^1 - z^2 \neq 0\)，存在指标 \(i\) 使 \(x_i (Mx)_i > 0\)，即 \((z_i^1 - z_i^2)(w_i^1 - w_i^2) > 0\)。但展开 \((z_i^1 - z_i^2)(w_i^1 - w_i^2) = z_i^1 w_i^1 + z_i^2 w_i^2 - z_i^1 w_i^2 - z_i^2 w_i^1\)。由互补性，\(z_i^1 w_i^1 = z_i^2 w_i^2 = 0\)。由非负性，\(-z_i^1 w_i^2 \leq 0\) 和 \(-z_i^2 w_i^1 \leq 0\)，故整个表达式 \(\leq 0\)，与 \(>0\) 矛盾。因此 \(z^1 = z^2\)。

\textbf{存在性}：对每个 \(\alpha \subseteq \{1, \ldots, n\}\)，考虑\textbf{互补基系统}：对 \(i \in \alpha\) 令 \(w_i = 0\)，对 \(i \notin \alpha\) 令 \(z_i = 0\)，解线性方程组得 \(z(\alpha)\)。定义映射 \(H_\alpha : \mathbb{R}^n \to \mathbb{R}^n\) 为分段线性同伦。由 P-矩阵的性质，每个 \(H_\alpha\) 是双射。利用同伦不变性理论的度理论（degree theory），\(\operatorname{LCP}(M, q)\) 的拓扑度 \(\deg(H, 0) = \pm 1 \neq 0\)，故零点的原像非空，保证了解的存在性。具体地，当 \(M\) 为 P-矩阵时，Lemke 算法从任意初始基出发，沿射线追踪互补基的变换，必在有限步内找到互补解。\(\blacksquare\)

Cottle-Dantzig 定理将 \(\operatorname{LCP}\) 的唯一可解性完全归约为矩阵的 P-性质。这是 \(\operatorname{LCP}\) 理论中最优雅的结构定理之一，其证明方法------通过同伦度理论------随后被推广到非线性互补问题。

\textbf{定理139.4}（Karamardian 定理，1972）：设 \(K \subseteq \mathbb{R}^n\) 为闭凸锥，\(F : K \to \mathbb{R}^n\) 为连续且关于 \(K\) 是\textbf{余正}的（即 \(\langle F(x) - F(y), x - y \rangle \geq 0\) 对所有 \(x, y \in K\)）。若 \(F\) 是\textbf{强单调}的（即存在 \(\mu > 0\) 使 \(\langle F(x) - F(y), x - y \rangle \geq \mu\|x - y\|^2\)），则 \(\operatorname{VI}(K, F)\) 有唯一解。

\textbf{证明}：\textbf{唯一性}：设 \(x^*, y^*\) 均为 \(\operatorname{VI}(K, F)\) 的解。由定义，\(\langle F(x^*), y^* - x^* \rangle \geq 0\) 和 \(\langle F(y^*), x^* - y^* \rangle \geq 0\)。相加得 \(\langle F(x^*) - F(y^*), y^* - x^* \rangle \geq 0\)，即 \(\langle F(x^*) - F(y^*), x^* - y^* \rangle \leq 0\)。但由强单调性，\(\langle F(x^*) - F(y^*), x^* - y^* \rangle \geq \mu\|x^* - y^*\|^2\)。两者结合，\(\mu\|x^* - y^*\|^2 \leq 0\)，故 \(\|x^* - y^*\| = 0\)，即 \(x^* = y^*\)。

\textbf{存在性}：定义自然映射（natural map）\(F_K^{\text{nat}}(x) = x - P_K(x - F(x))\)，其中 \(P_K\) 是到 \(K\) 的投影。容易验证 \(x^*\) 是 \(\operatorname{VI}(K, F)\) 的解当且仅当 \(F_K^{\text{nat}}(x^*) = 0\)，即 \(x^* = P_K(x^* - F(x^*))\)。由强单调性，\(F_K^{\text{nat}}\) 是 \(\mathbb{R}^n\) 上的同胚映射，故 \((F_K^{\text{nat}})^{-1}(0)\) 非空，保证了解的存在性。此外，可构造迭代序列 \(x^{k+1} = P_K(x^k - \alpha F(x^k))\)，在 \(\alpha\) 充分小时收敛到唯一解（投影梯度法）。\(\blacksquare\)

Karamardian 定理的意义在于，它将 \(\operatorname{LCP}\) 的存在唯一性从线性情形（P-矩阵）推广到非线性单调算子的一般框架，为变分不等式的数值求解建立了理论基础。定理中的强单调性可放宽为 \(F\) 单调且满足强制性条件 \(\lim_{\|x\|\to\infty} \langle F(x), x \rangle / \|x\| = \infty\)。

\subsection{增广Lagrangian方法与内点法}\label{ux589eux5e7flagrangianux65b9ux6cd5ux4e0eux5185ux70b9ux6cd5}

\textbf{增广 Lagrangian 方法}将 \(\operatorname{VI}(K, F)\) 转化为等价的无约束（或简单约束）优化问题。定义增广 Lagrangian 函数 \[L_\rho(x, \lambda) = \langle F(x), x \rangle + \frac{\rho}{2} \|x - P_K(x - F(x))\|^2\] 其中 \(\rho > 0\) 是罚参数。交替方向乘子法（ADMM）在此基础上实现并行化：将 \(\operatorname{VI}\) 分解为多个子问题，交替更新原始变量和乘子，在 \(\rho\) 充分大时保证线性收敛。对于 \(\operatorname{LCP}(M, q)\) 的特例，增广 Lagrangian 方法等价于邻近点算法（proximal point algorithm）\(z^{k+1} = (M + \rho I)^{-1}(\rho z^k - q)\)，当 \(M\) 为 P-矩阵时收敛。

\textbf{内点法}是求解 \(\operatorname{LCP}\) 的另一类重要算法。在 \(\operatorname{LCP}\) 中引入障碍参数 \(\mu > 0\)，将互补性条件 \(z_i w_i = 0\) 松弛为 \(z_i w_i = \mu\)（中心路径条件），然后使用 Newton 法追踪中心路径：求解线性化系统 \[\begin{pmatrix} M & -I \\ W^{(k)} & Z^{(k)} \end{pmatrix} \begin{pmatrix} \Delta z \\ \Delta w \end{pmatrix} = \begin{pmatrix} -Mz^{(k)} + w^{(k)} - q \\ \mu e - Z^{(k)} W^{(k)} e \end{pmatrix}\] 其中 \(Z^{(k)} = \operatorname{diag}(z^{(k)})\)，\(W^{(k)} = \operatorname{diag}(w^{(k)})\)。当 \(M\) 是 P\_0-矩阵且满足适当条件时，内点法具有多项式时间复杂度 \(O(\sqrt{n} \log(1/\varepsilon))\)。内点法同样适用于半定互补问题（SDCP），后者将 \(\operatorname{LCP}\) 推广到半正定锥 \(\mathbb{S}^n_+\) 上：求 \(X \in \mathbb{S}^n_+\) 使 \(X \succeq 0\)，\(\mathcal{L}(X) + Q \succeq 0\)，且 \(\langle X, \mathcal{L}(X) + Q \rangle = 0\)，其中 \(\mathcal{L}\) 是 Lyapunov 型线性变换。互补性问题的深刻之处在于，它将优化的 KKT 条件、博弈论的 Nash 均衡和力学中的接触问题统一在单一数学框架下，使其成为跨越纯数学和应用数学的基础理论。

\hrulefill

\hrulefill

\hrulefill

\hrulefill

\section{第433章：半定规划与离散凸分析}\label{ux7b2c433ux7ae0ux534aux5b9aux89c4ux5212ux4e0eux79bbux6563ux51f8ux5206ux6790}

半定规划（SDP）是线性规划的深刻推广，将优化变量从向量提升为半正定矩阵，在线性矩阵不等式约束下极小化线性目标函数。SDP 严格包含线性规划和二阶锥规划，内点法可在多项式时间内求解。SDP 在组合优化中的突破性应用------Goemans-Williamson 的 MAX-CUT 近似算法（1995）------开创了基于 SDP 松弛的近似算法设计范式。离散凸分析（Murota, 1990s）将凸分析推广到离散格点 \(\mathbb{Z}^n\)，引入 L\(^\natural\)-凸函数和 M\(^\natural\)-凸函数两类互为共轭的函数族，在经济学（不可分商品的竞争均衡）和运筹学（资源配置）中具有重要应用。本章将系统阐述 SDP 的对偶理论、内点算法和组合优化应用，以及离散凸分析的基本框架。

\subsection{137.1 半定规划}\label{ux534aux5b9aux89c4ux5212}

半定规划（Semidefinite Programming, SDP）是在线性矩阵不等式约束下极小化线性目标函数的凸优化问题。

\textbf{定义 137.1}（半定规划的原始形式）：设 \(\mathbb{S}^m\) 为 \(m \times m\) 实对称矩阵构成的向量空间，配备标准内积 \(\langle X, Y \rangle = \operatorname{Tr}(XY)\)。给定 \(C, A_1, \ldots, A_n \in \mathbb{S}^m\) 和 \(b \in \mathbb{R}^n\)，SDP 的\textbf{原始问题}为

\[\begin{aligned} \min_X \quad & \langle C, X \rangle \\ \text{s.t.} \quad & \langle A_i, X \rangle = b_i \quad (i = 1, \ldots, n) \\ & X \succeq 0 \end{aligned}\]

其中 \(X \succeq 0\) 表示 \(X\) 是半正定矩阵（即所有特征值非负，或等价地对任意 \(v \in \mathbb{R}^m\) 有 \(v^T X v \geq 0\)）。\(\mathbb{S}^m_+ = \{X \in \mathbb{S}^m : X \succeq 0\}\) 为半正定锥，是自对偶的闭凸锥。

等价地，\textbf{不等式形式}（又称线性矩阵不等式形式）为：

\[\min_{\mathbf{x} \in \mathbb{R}^n} \; \mathbf{c}^T \mathbf{x} \quad \text{s.t.} \quad F(\mathbf{x}) := F_0 + \sum_{i=1}^n x_i F_i \succeq 0\]

其中 \(F_i \in \mathbb{S}^m\)（\(i = 0, \ldots, n\)）为给定的对称矩阵。此形式通过定义 \(X = F(\mathbf{x})\) 的隐式约束与上述标准形式等价。

\textbf{定义 137.2}（SDP的对偶问题）：SDP 的\textbf{Lagrange 对偶问题}为

\[\begin{aligned} \max_{y \in \mathbb{R}^n, Z \in \mathbb{S}^m} \quad & b^T y \\ \text{s.t.} \quad & \sum_{i=1}^n y_i A_i + Z = C \\ & Z \succeq 0 \end{aligned}\]

等价地，消去 \(Z\) 可得

\[\max_{y \in \mathbb{R}^n} \; b^T y \quad \text{s.t.} \quad C - \sum_{i=1}^n y_i A_i \succeq 0\]

对不等式形式，Lagrange 对偶为 \(\max_{Z \succeq 0} \; -\operatorname{Tr}(F_0 Z) \quad \text{s.t.} \quad \operatorname{Tr}(F_i Z) = c_i \; (i = 1, \ldots, n)\)。

\textbf{定义 137.3}（线性矩阵不等式与 M-矩阵约束）：(i) \textbf{线性矩阵不等式}（Linear Matrix Inequality, LMI）是形如 \(F(\mathbf{x}) = F_0 + \sum_{i=1}^n x_i F_i \succeq 0\) 的约束。LMI 在控制理论中极为普遍：Lyapunov 不等式 \(A^T P + PA \prec 0\)（其中 \(P\) 为未知对称矩阵）是典型的 LMI。(ii) 矩阵 \(M \in \mathbb{R}^{m \times m}\) 称为 \textbf{M-矩阵}，若其可表为 \(M = sI - B\)，其中 \(B \geq 0\)（元素非负）且 \(s > \rho(B)\)（\(\rho\) 为谱半径）。M-矩阵等价于所有非对角元非正、所有主子式为正、且逆矩阵非负（\(M^{-1} \geq 0\)）。在 SDP 的上下文中，当 \(F_i\) 具有 M-矩阵结构时，可得到更强的对偶性质。

\textbf{定理 137.1}（SDP的弱对偶定理）：设 \(X\) 为原始问题的可行解，\((y, Z)\) 为对偶问题的可行解。则

\[\langle C, X \rangle - b^T y = \langle X, Z \rangle \geq 0\]

特别地，原始最优值 \(p^*\) 和对偶最优值 \(d^*\) 满足 \(p^* \geq d^*\)。若存在严格可行解（Slater 条件），即存在 \(X \succ 0\)（正定）满足原始约束，或存在 \(y\) 使 \(C - \sum_i y_i A_i \succ 0\) 满足对偶约束，则强对偶成立：\(p^* = d^*\) 且对偶最优值可达。

\textbf{证明}：由原始和对偶约束： \[\langle C, X \rangle - b^T y = \langle \sum_i y_i A_i + Z, X \rangle - \sum_i y_i \langle A_i, X \rangle = \langle Z, X \rangle\] 因 \(X \succeq 0\) 且 \(Z \succeq 0\)，故 \(\langle Z, X \rangle = \operatorname{Tr}(ZX) \geq 0\)（半正定矩阵内积的非负性由 \(X = LL^T\) 和 \(Z\) 的半正定性保证：\(\operatorname{Tr}(ZX) = \operatorname{Tr}(ZLL^T) = \operatorname{Tr}(L^T ZL) \geq 0\)）。因此弱对偶 \(p^* \geq d^*\) 恒成立。强对偶在 Slater 条件下由凸规划的广义 Slater 定理保证------SDP 是锥规划的特例，半正定锥 \(\mathbb{S}^m_+\) 的内部非空。\(\blacksquare\)

\textbf{定理 137.2}（内点法在SDP中的应用思路）：SDP 可通过\textbf{原始-对偶内点法}在多项式时间内求解。核心思想是引入障碍参数 \(\mu > 0\)，将互补性条件 \(XZ = 0\) 松弛为 \(XZ = \mu I\)（中心路径条件），然后使用 Newton 法追踪中心路径。具体地，中心路径方程为

\[\begin{aligned} \langle A_i, X \rangle &= b_i \quad (i = 1, \ldots, n) \\ \sum_{i=1}^n y_i A_i + Z &= C \\ XZ &= \mu I \end{aligned}\]

Newton 步的线性化系统为对称不定系统，可利用 Schur 补约化为 \(n \times n\) 的稠密线性系统。当 \(\mu \to 0\) 时中心路径收敛到最优解。内点法的每次迭代复杂度为 \(O(m^3 + n^3)\)，总迭代次数为 \(O(\sqrt{m} \log(1/\varepsilon))\) 以达到 \(\varepsilon\) 精度，整体为多项式时间算法。这一框架也涵盖线性规划（\(F_i\) 为对角矩阵时退化）、二阶锥规划和更一般的对称锥规划。

SDP 包含线性规划（LP, \(F_i\) 为对角矩阵时）和二阶锥规划（SOCP）作为特例，严格包含关系为 \(\text{LP} \subset \text{SOCP} \subset \text{SDP}\)。在组合优化中 SDP 松弛有重要突破：Goemans-Williamson（1995）将 MAX-CUT 问题表示为整数二次规划，通过 SDP 松弛（\(X \succeq 0\) 代替 \(X = \mathbf{v}\mathbf{v}^T\)）并使用随机超平面舍入，获得了约 0.878 的近似比------这是基于 SDP 的第一个突破性近似算法，极大地推动了半定规划在近似算法设计中的应用。SDP 在控制理论中有重要应用：Lyapunov 不等式 \(A^T P + PA \prec 0\) 是 LMI 的核心形式；\(H_\infty\) 控制中的有界实引理可通过 SDP 求解。在组合优化中，SDP 松弛还成功应用于图着色（Lovász theta 函数 \(\vartheta(G)\) 是 Shannon 容量的 SDP 上界）、稀疏主成分分析和矩阵完备化。此外，SDP 与量子信息中的优化问题（如计算量子信道的 Holevo 容量）有深刻联系------量子态密度矩阵的半正定性自然地使 SDP 成为量子信息论中的基本计算工具。

\subsection{137.2 离散凸分析}\label{ux79bbux6563ux51f8ux5206ux6790}

离散凸分析由室田一雄（Murota）于 1990 年代建立，将凸分析推广到离散格点 \(\mathbb{Z}^n\)。核心概念包括两类互为共轭的函数族。

\textbf{定义 137.4}（L\(^\natural\)-凸函数）：函数 \(f : \mathbb{Z}^n \to \mathbb{R} \cup \{+\infty\}\)（\(\operatorname{dom} f \neq \emptyset\)）称为 \textbf{L\(^\natural\)-凸函数}，若其满足\textbf{离散中点凸性}：

\[f(x) + f(y) \geq f\left(\left\lfloor\frac{x+y}{2}\right\rfloor\right) + f\left(\left\lceil\frac{x+y}{2}\right\rceil\right) \quad (\forall x, y \in \mathbb{Z}^n)\]

其中 \(\lfloor\cdot\rfloor\) 和 \(\lceil\cdot\rceil\) 分别表示逐分量的向下和向上取整。L\(^\natural\)-凸函数的有效域 \(\operatorname{dom} f\) 是 \textbf{L\(^\natural\)-凸集}，即满足 \(x, y \in S \Rightarrow \lfloor\frac{x+y}{2}\rfloor, \lceil\frac{x+y}{2}\rceil \in S\) 的集合。

\textbf{定义 137.5}（M\(^\natural\)-凸函数）：函数 \(f : \mathbb{Z}^n \to \mathbb{R} \cup \{+\infty\}\) 称为 \textbf{M\(^\natural\)-凸函数}，若对任意 \(x, y \in \operatorname{dom} f\) 和任意指标 \(i\) 满足 \(x_i > y_i\)，存在指标 \(j\)（\(x_j < y_j\) 或 \(j = 0\)，约定 \(x_0 = y_0\)）使得

\[f(x) + f(y) \geq f(x - \mathbf{e}_i + \mathbf{e}_j) + f(y + \mathbf{e}_i - \mathbf{e}_j)\]

其中 \(\mathbf{e}_i\) 为第 \(i\) 个标准基向量（\(\mathbf{e}_0 = \mathbf{0}\)）。此条件称为\textbf{交换公理}（exchange axiom），刻画了离散凸分析中''资源的边际替代递减''这一经济学直观。

两类函数通过 \textbf{Legendre-Fenchel 变换}（离散版）互为共轭：令

\[f^\bullet(p) = \sup_{x \in \mathbb{Z}^n} \{ \langle p, x \rangle - f(x) \}\]

则 \(f\) 是 L\(^\natural\)-凸当且仅当 \(f^\bullet\) 是 M\(^\natural\)-凸（在适当的定义域约束下）。这一对偶性推广了线性规划中多面体与多面锥的对偶，是离散凸分析的核心结构定理。

\textbf{定理 137.3}（L\(^\natural\)-凸函数的最优性条件）：设 \(f\) 为 L\(^\natural\)-凸函数，则 \(x^* \in \mathbb{Z}^n\) 是 \(f\) 的全局极小点当且仅当

\[f(x^*) \leq f(x^* \pm \mathbf{1}_S) \quad (\forall S \subseteq \{1, \ldots, n\})\]

即 \(x^*\) 在所有单位超立方体的顶点中均为最小。这给出了离散凸极小化的局部-全局等价原理。

\textbf{定理 137.4}（离散凸分析的存在性定理）：设 \(f : \mathbb{Z}^n \to \mathbb{R} \cup \{+\infty\}\) 为 L\(^\natural\)-凸函数且 \(\operatorname{dom} f\) 有界非空。则 \(f\) 在 \(\operatorname{dom} f\) 上达到最小值。若进一步 \(f\) 是 M\(^\natural\)-凸函数，则最小值可通过贪心算法（greedy algorithm）在多项式步内找到。

离散凸分析的应用涵盖经济学中不可分商品的竞争均衡存在性（Kelso-Crawford 的劳动市场模型）和运筹学中的资源配置与库存管理。在组合优化中，M\(^\natural\)-凸函数的最小化对应拟阵约束下的资源分配问题，其贪心可解性正是拟阵结构的离散凸性体现。

\hrulefill

\hrulefill

\hrulefill

\hrulefill

\section{第434章：动态规划与随机规划}\label{ux7b2c434ux7ae0ux52a8ux6001ux89c4ux5212ux4e0eux968fux673aux89c4ux5212}

动态规划（Bellman, 1957）和随机规划是运筹学中处理序贯决策和不确定性的核心数学工具。动态规划基于 Bellman 最优性原理------最优策略的剩余部分必须自身是最优的------将多阶段决策问题分解为递归求解的子问题，是强化学习（reinforcement learning）和 Markov 决策过程（MDP）的理论基础。随机规划则处理优化问题中参数的不确定性（如需求、价格、收益），通过情景树（scenario tree）或机会约束（chance constraint）将不确定性纳入优化模型。本章将系统阐述确定性有限期和无限期动态规划（值迭代、策略迭代）、MDP 的 Blackwell 最优性理论，以及随机规划的二阶段模型和机会约束规划。

\subsection{138.1 动态规划（Dynamic Programming）}\label{ux52a8ux6001ux89c4ux5212dynamic-programming}

动态规划（Bellman, 1957）通过将多阶段决策问题分解为子问题的递归关系解决，奠基于 \textbf{Bellman 最优性原理}------最优策略具有这样的性质：无论初始状态和初始决策如何，剩余决策必须构成由第一个决策导致的状态的最优策略。

\textbf{定义 138.1}（确定性有限期动态规划）：对 \(N\) 阶段确定性决策问题，值函数 \(V_k(x)\) 定义为当前状态为 \(x\)、还剩 \(k\) 步时的最优费用，满足 \textbf{Bellman 方程}

\[V_k(x) = \min_{u \in U(x)} \left[ c(x, u) + V_{k-1}\big(f(x, u)\big) \right], \quad k = 1, 2, \ldots, N\]

其中 \(c(x, u)\) 为瞬时代价，\(f(x, u)\) 为状态转移函数，\(U(x)\) 为状态 \(x\) 的可行控制集。边界条件 \(V_0(x) = g(x)\)（终端代价）。

\textbf{定义 138.2}（无限期折扣问题）：引入折扣因子 \(\gamma \in (0, 1)\)，目标为最小化期望折现总代价 \(\mathbb{E}\left[ \sum_{t=0}^\infty \gamma^t c(x_t, u_t) \right]\)。Bellman 方程化为不动点方程

\[V(x) = \min_{u \in U(x)} \left[ c(x, u) + \gamma V\big(f(x, u)\big) \right]\]

\textbf{算法 138.1}（值迭代 / Value Iteration）：\(V_{n+1}(x) = \min_u [c(x,u) + \gamma V_n(f(x,u))]\)。当 \(\gamma < 1\) 时，该迭代是 \(\infty\)-范数下的压缩映射（压缩因子 \(\gamma\)），线性收敛到唯一不动点。

\textbf{算法 138.2}（策略迭代 / Policy Iteration）：交替执行策略评估和策略改进，对有限状态和有限控制问题，在至多有限步内收敛到最优策略（每步严格提升）。

\textbf{定义 138.3}（Markov 决策过程 / MDP）：动态规划在随机环境中的推广。转移概率 \(p(x' \mid x, u)\) 替代确定性转移 \(f(x, u)\)，Bellman 方程为

\[V(x) = \min_{u \in U(x)} \left[ c(x, u) + \gamma \sum_{x'} p(x' \mid x, u) \, V(x') \right]\]

\textbf{定理 138.1}（Blackwell 最优性理论，1965）：在标准的折扣 MDP（有限状态、有限控制）中，存在确定性的平稳最优策略，且该策略可由值迭代或策略迭代找到。最优策略的存在性和结构由 Blackwell 的单调性和压缩性条件保证。

\textbf{证明}：定义Bellman算子\(\mathcal{T}V(x) = \min_u[c(x,u) + \gamma\sum_{x'}p(x'|x,u)V(x')]\)。\(\mathcal{T}\)是\(\|\cdot\|_\infty\)下的\(\gamma\)-压缩映射（\(\|\mathcal{T}V_1-\mathcal{T}V_2\|_\infty\le\gamma\|V_1-V_2\|_\infty\)）。由Banach不动点定理，\(\mathcal{T}\)有唯一不动点\(V^*\)------即最优值函数。最优策略\(\pi^*(x)=\arg\min_u[c(x,u)+\gamma\sum_{x'}p(x'|x,u)V^*(x')]\)是确定性的（有限集上min可达）和平稳的（\(V^*\)唯一）。策略迭代在有限步内收敛（策略空间有限）。\(\blacksquare\)

\subsection{138.2 随机规划（Stochastic Programming）}\label{ux968fux673aux89c4ux5212stochastic-programming}

随机规划处理包含随机参数（通过随机变量 \(\xi\) 表达）的优化问题。

\textbf{定义 138.4}（二阶段线性随机规划 / Dantzig, 1955）：第一阶段在 \(\xi\) 实现前确定''现在''决策 \(x\)，第二阶段在 \(\xi\) 实现后采取折中（recourse）行动 \(y\)：

\[\begin{aligned} \min_x \quad & c^T x + \mathbb{E}_\xi[Q(x, \xi)] \\ \text{s.t.} \quad & Ax = b, \quad x \geq 0 \end{aligned}\]

其中 \(Q(x,\xi) = \min_y \{q(\xi)^T y : W(\xi) y = h(\xi) - T(\xi) x, \; y \geq 0\}\) 为第二阶段（recourse）问题的值函数。

\textbf{定理 138.2}（二阶段问题的结构性质）：若对几乎所有 \(\xi\)，第二阶段问题是可行的（完全 recourse），则 \(Q(x, \xi)\) 是关于 \(x\) 的分片线性凸函数，\(\mathbb{E}_\xi[Q(x, \xi)]\) 也是凸函数。当 \(\xi\) 具有有限支撑（场景树）时，二阶段问题等价于一个大规模确定性线性规划。

\textbf{证明}：\(Q(x,\xi)=\min_y\{q^Ty:Wy=h-Tx,y\ge 0\}\)为线性规划的值函数。由线性规划的对偶理论，\(Q(x,\xi)=\max_\pi\{\pi^T(h-Tx):W^T\pi\le q\}\)。对偶可行域\(\Pi=\{\pi:W^T\pi\le q\}\)与\(x\)无关。\(Q(x,\xi)=\max_{\pi\in\text{ext}(\Pi)}\pi^T(h-Tx)\)，其中\(\text{ext}(\Pi)\)为\(\Pi\)的极点（有限集）。故\(Q(\cdot,\xi)\)是有限个仿射函数的最大值，为分片线性凸函数。凸函数族的期望\(\mathbb{E}_\xi[Q(\cdot,\xi)]\)仍为凸函数。\(\xi\)有限支撑时，\(\mathbb{E}_\xi[Q(x,\xi)]=\sum_k p_k Q(x,\xi_k)\)，展开为大规模LP。\(\blacksquare\)

\textbf{定义 138.5}（机会约束规划 / Chance-Constrained Programming, Charnes-Cooper, 1959）：用概率约束替代确定性约束：

\[P\big(g(x, \xi) \leq 0\big) \geq 1 - \alpha, \quad \alpha \in (0,1)\]

当 \(g\) 关于 \(x\) 线性、\(\xi\) 具有已知联合正态分布时，概率约束等价于一个二阶锥约束。

\subsection{138.3 多目标最优化}\label{ux591aux76eeux6807ux6700ux4f18ux5316}

\textbf{定义 138.6}（Pareto 最优解）：对多目标问题 \(\min (f_1(x), \ldots, f_k(x))\)（\(x \in \mathcal{X}\)），解 \(x^*\) 称为 \textbf{Pareto 最优}，若不存在 \(x' \in \mathcal{X}\) 使得对所有 \(i\) 有 \(f_i(x') \leq f_i(x^*)\)，且至少一个不等式严格成立。所有 Pareto 最优解的集合称为 \textbf{Pareto 前沿}。

\textbf{求解方法}： - \textbf{加权和法}（标量化）：\(\min \sum_{i=1}^k w_i f_i(x)\)（\(w_i \geq 0\)），变化权重可生成 Pareto 前沿的凸部分 - \textbf{\(\varepsilon\)-约束法}：\(\min f_1(x)\) s.t. \(f_i(x) \leq \varepsilon_i\)（\(i = 2, \ldots, k\)），系统改变 \(\varepsilon_i\) 可生成完整的 Pareto 前沿（包括非凸部分） - \textbf{目标规划}：最小化与目标的加权偏差 \(\min \sum_i |f_i(x) - t_i|\)

Pareto 前沿的结构由 Kuhn-Tucker 型必要条件刻画：若 \(x^*\) Pareto 最优且各目标在约束下正则，则存在非负乘子 \(\lambda_i \geq 0\)（不全为零）使 \(\sum \lambda_i \nabla f_i(x^*) = 0\)。

\hrulefill

\hrulefill

\hrulefill

\hrulefill

\hrulefill

\subsection{351.A 最优化补充}\label{a-ux6700ux4f18ux5316ux8865ux5145}

\begin{quote}
本卷研究由工程控制问题抽象出的数学理论------线性系统的代数结构、最优控制的变分方法、动态规划的HJB方程、以及Lyapunov稳定性理论。内容本质上是线性代数、常微分方程和变分法的深化专题。
\end{quote}

\hrulefill

\hrulefill

\emph{卷二十八：最优化理论终。}

\hrulefill

\subsection{357.A 博弈论补充}\label{a-ux535aux5f08ux8bbaux8865ux5145}

最优化理论与博弈论是 20 世纪数学与应用科学交汇的两座高峰，它们分别从不同的视角研究''最优决策''------最优化关注单个决策者如何在约束下最大化目标函数，博弈论则关注多个交互决策者如何在策略互动中实现均衡。这两大学科在 von Neumann 的极小极大定理（1928）中达到了深刻的统一：二人零和博弈的混合策略均衡等价于线性规划对偶对 \((P)\) 和 \((D)\) 的最优解。具体而言，行参与者的最小最大收益问题可表述为线性规划，列参与者的最大最小损失问题则是其对偶问题，对偶定理（弱对偶 \(v_P\geq v_D\)，强对偶 \(v_P=v_D\)）恰好证明了博弈的值存在且双方的最优混合策略可通过 LP 的互补松弛条件刻画------由 von Neumann 和 Dantzig 发现的这一等价关系奠定了博弈论、线性规划和矩阵博弈三者数学统一的基石。本卷在此基础上，系统阐述最优化理论的代数结构（线性规划对偶、整数规划的 Gomory 切割面、网络流的图论基础）和博弈论的数学模型（Nash 均衡的拓扑不动点证明、合作博弈的核与 Shapley 值公理化、机制设计的激励相容理论），并揭示两者之间的深层数学联系。

\hrulefill

\hrulefill

\hrulefill

\hrulefill

\section{第435章：整数规划}\label{ux7b2c435ux7ae0ux6574ux6570ux89c4ux5212}

整数规划是线性规划在离散约束下的推广，其变量取值被限制在整数集上------这一看似简单的离散性要求使得问题的计算复杂性从多项式时间（线性规划）跃升为 NP 难。整数规划的数学结构深植于组合优化和离散代数：当约束矩阵是全幺模矩阵（每个方子阵的行列式为 \(0\) 或 \(\pm 1\)）时，线性规划松弛的基可行解自动为整数------这是指派问题和网络流问题多项式可解的代数根源。一般情形下，Gomory 割平面法（1958）通过从单纯形表生成不等式割去分数解而不切除任何整数可行解，在有限步内收敛到整数最优解。分枝定界法（Branch-and-Bound）是实际求解整数规划最有效的框架：递归地对分数变量进行分枝，利用 LP 松弛值作为下界进行剪枝，结合割平面（形成分枝切割法）大幅提高了求解效率。0-1 背包问题是整数规划的经典特例，尽管是 NP 难的，却具有伪多项式时间动态规划解（\(O(nW)\)），其 LP 松弛具有简单的贪心结构------按价值密度 \(v_i/w_i\) 降序放入物品。

\subsection{253.1 整数规划基础}\label{ux6574ux6570ux89c4ux5212ux57faux7840}

\textbf{定义 253.1}（整数规划 / Integer Programming / IP）：\[\min c^T x \quad \text{s.t.} \quad Ax = b,\ x \geq 0,\ x \in \mathbb{Z}^n\]

若部分变量取整，则为\textbf{混合整数规划}（MIP）。若所有变量取值\(\{0, 1\}\)，则为\textbf{0-1规划}。

\textbf{定义 253.2}（LP松弛与对偶间隙）：放宽整数约束得到LP松弛。LP松弛最优值\(\leq\) IP最优值（极小化问题）。两者的差距称为\textbf{对偶间隙}（duality gap）。

\textbf{定理 253.1}（全幺模矩阵 / Totally Unimodular）：若约束矩阵\(A\)是全幺模的（每个方子阵的行列式为\(0, \pm 1\)），则对任意整数右端项\(b\)，线性规划\(Ax = b, x \geq 0\)的所有基可行解都是整数的。典型例子：指派问题和网络流问题。全幺模矩阵将整数规划可多项式时间求解的情形精确刻画为代数条件。

\emph{证明}：由 Cramer 法则，基矩阵 \(B\) 的逆 \(B^{-1}=\operatorname{adj}(B)/\det(B)\)。若 \(A\) 全幺模，则 \(\det(B)\in\{0,\pm 1\}\) 且伴随矩阵的元素为 \(A\) 的子行列式（整数），故 \(B^{-1}\) 为整数矩阵。基解 \(x_B=B^{-1}b\) 为整数向量（\(b\) 整数）。单纯形法可直接获得整数最优解。\(\blacksquare\)

\subsection{253.2 分枝定界与割平面法}\label{ux5206ux679dux5b9aux754cux4e0eux5272ux5e73ux9762ux6cd5}

\textbf{定义 253.3}（分枝定界法 / Branch-and-Bound）：递归地分解问题------在LP松弛解中选分数变量\(x_j\)，分枝为\(x_j \leq \lfloor x_j \rfloor\)和\(x_j \geq \lceil x_j \rceil\)两个子问题。用LP松弛值作为下界（极小化时），剪去不可能含最优解的枝（bounding）。

\textbf{定理 253.2}（割平面法 / Cutting Plane Method / Gomory 1958）：若LP松弛最优解非整数，则从单纯形表生成Gomory割\(-\bar{f}_j x_j - \cdots \leq -\bar{f}_0\)（\(\bar{f}_j\)为分数部分），添加此约束割去当前分数解而不损失任何整数可行解。有限次迭代后收敛到整数最优解。

\emph{有限终止性论证}：每个Gomory割从当前LP最优解对应的单纯形表的一行导出：设基变量\(x_{B_i} = \bar{b}_i - \sum_{j \in NB} \bar{a}_{ij} x_j\)，其中\(\bar{b}_i = \lfloor \bar{b}_i \rfloor + \bar{f}_{i0}\)（\(0 < \bar{f}_{i0} < 1\)）。Gomory割为\(\sum_{j \in NB} \bar{f}_{ij} x_j \geq \bar{f}_{i0}\)（其中\(\bar{f}_{ij} = \bar{a}_{ij} - \lfloor \bar{a}_{ij} \rfloor\)）。添加此割后，对偶单纯形法在当前基不可行但最优的前提下，经过一次对偶迭代恢复可行性并保持最优性，目标函数值不增（极小化问题）。由于整数规划的可行域是有限个整数点的集合（有界情形），而每个割严格排除当前分数LP最优解，结合目标函数的单调性和有限整数点集，算法必然在有限步内终止于整数最优解或证明不可行。∎

\textbf{注 253.1}（分枝定界法的策略与定界方法）：分枝定界法的效率高度依赖分枝策略和定界方法。

\textbf{分枝策略}：常见策略包括（1）\textbf{最大不可行性}（maximum infeasibility）：选择与整数偏离最远的分数变量\(x_j\)（即\(|x_j - \text{round}(x_j)|\)最大的变量）；（2）\textbf{伪费用分枝}（pseudo-cost branching）：记录各变量分枝后目标函数的历史变化率，优先选择预计改善最大的变量；（3）\textbf{强分枝}（strong branching）：在分枝前试探性求解若干候选分枝的LP松弛，选择使下界上升最多的分枝。

\textbf{定界方法}：每个子节点的LP松弛值提供该子树的\textbf{下界}（极小化问题）。维护全局\textbf{上界}（当前最优整数解的目标值）。若某节点的下界\(\geq\)当前上界，则该节点可被\textbf{剪枝}（fathomed）。此外，若子问题的LP松弛不可行或已获得整数解，也终止该分枝。现代求解器还结合切割平面形成\textbf{分枝切割法}（branch-and-cut），在每个节点处动态添加有效不等式以收紧LP松弛。

\subsection{253.3 背包问题}\label{ux80ccux5305ux95eeux9898}

\textbf{定义 253.4}（0-1背包问题 / Knapsack Problem）：\[\max \sum_{i=1}^n v_i x_i \quad \text{s.t.} \quad \sum_{i} w_i x_i \leq W,\ x_i \in \{0, 1\}\]

其中\(v_i\)为物品的价值参数，\(w_i\)为重量参数，\(W\)为容量参数。

背包问题是组合优化中最经典的 NP 难问题之一，也是许多实践问题的核心模型（如资源分配、投资组合选择、货物装载）。虽然 0-1 背包问题是 NP 难的，但它具有\textbf{伪多项式时间算法}（即运行时间 \(O(nW)\) 依赖于数值参数 \(W\) 而非仅输入长度 \(n \log W\)），这意味着当 \(W\) 不大时问题是可高效求解的。背包问题的 LP 松弛（将 \(x_i \in \{0,1\}\) 松弛为 \(0 \leq x_i \leq 1\)）具有贪心结构：按价值密度 \(v_i/w_i\) 降序排列，依次放入物品直至容量耗尽，最后一件物品可能被分数放入。LP 松弛值为 \(v_{\text{LP}} = \sum_{j=1}^{k-1} v_j + \frac{v_k}{w_k}(W - \sum_{j=1}^{k-1} w_j)\)，其中 \(k\) 为第一个无法完整放入的物品。这一贪心结构直接导致了 Dantzig 上界和分支定界法中的有效剪枝。背包问题的变体包括：无界背包（每种物品可重复选择）、多重背包（多种容量约束）、多维背包（\(d\)-维容量约束，\(d \geq 2\) 时动态规划复杂度为 \(O(nW^d)\)，指数增长）。Martello 和 Toth 的专著《Knapsack Problems》系统总结了背包问题的算法和变体，是现代组合优化的重要参考文献。

\textbf{定理 253.3}（动态规划解背包）：定义\(f(i, w)\)为前\(i\)个物品在容量\(w\)下的最大价值。递推：\(f(i, w) = \max\{f(i-1, w), f(i-1, w-w_i) + v_i\}\)（若\(w \ge w_i\)）。时间\(O(nW)\)（伪多项式）。背包问题的动态规划解是Bellman最优性原理在组合优化中的直接应用。

\emph{证明}：由 Bellman 最优性原理，\(f(i,w)=\max\{f(i-1,w), f(i-1,w-w_i)+v_i\}\)（若 \(w\ge w_i\)）。状态数 \(nW\)，每状态 \(O(1)\)，总复杂度 \(O(nW)\)。该算法是伪多项式时间算法------运行时间依赖数值参数 \(W\) 而非仅输入长度。\(\blacksquare\)

\hrulefill

\hrulefill

\hrulefill

\hrulefill

\section{第436章：网络流理论}\label{ux7b2c436ux7ae0ux7f51ux7edcux6d41ux7406ux8bba}

网络流理论是组合优化和图论中最具影响力的分支之一，其研究对象是物质、信息或交通在有向图中从源点到汇点的最优化流动问题。这一理论的核心定理------Ford-Fulkerson 最大流最小割定理（1956）------断言：网络中从源 \(s\) 到汇 \(t\) 的最大流值等于分离 \(s\) 和 \(t\) 的最小割容量。从数学结构上看，这一 min-max 关系是线性规划强对偶定理在有向图中的精确体现：最大流问题的 LP 约束矩阵是网络节点-弧关联矩阵------一个全幺模矩阵------因此 LP 松弛的最优解自动为整数流。Edmonds-Karp 算法（1972）通过 BFS 寻找最短增广路径，将 Ford-Fulkerson 方法的复杂度改进为 \(O(|V||E|^2)\)。最小费用流问题统一了最短路径、最大流和指派问题等组合优化经典模型，其最优性条件等价于残量网络无负费用有向环。König 定理（二分图匹配）和 Edmonds 花算法（一般图匹配，1965）将匹配问题------网络流理论的另一核心应用------嵌入网络流的统一理论框架，花算法通过收缩奇数长的交替环（blossom）巧妙地克服了一般图中奇环对匹配增广路径搜索造成的困难。

\subsection{254.1 最大流与最小割}\label{ux6700ux5927ux6d41ux4e0eux6700ux5c0fux5272}

\textbf{定义 254.1}（流网络 / Flow Network）：有向图\(G = (V, E)\)，每条边\((u, v)\)有容量\(c(u, v) \geq 0\)。源\(s\)，汇\(t\)。流\(f: E \to \mathbb{R}_{\geq 0}\)满足容量约束\(0 \leq f(e) \leq c(e)\)和守恒约束\(\sum_{(u,v)\in E} f(u,v) = \sum_{(v,w)\in E} f(v,w)\)（对所有\(v \neq s, t\)）。

\textbf{定义 254.2}（割 / Cut）：\(s\)-\(t\)割\((S, T)\)（\(s \in S, t \in T\)）的容量为\(\sum_{u \in S, v \in T} c(u, v)\)。

\textbf{定理 254.1}（Ford-Fulkerson最大流最小割定理，1956）：最大流的值等于最小\(s\)-\(t\)割的容量。

\[\max |f| = \min_{S: s \in S, t \notin S} \sum_{u \in S, v \notin S} c(u, v)\]

\textbf{证明}：将最大流问题化为线性规划：\(\max \sum_{P}f_P\)，约束 \(f_P\ge 0\)，\(\sum_{P\ni e}f_P\le c(e)\)。其对偶为 \(\min\sum_e c(e)y_e\)，约束 \(y_e\ge 0\)，\(\sum_{e\in P}y_e\ge 1\)。对偶表示每条路径至少含一个''割边''。由LP强对偶定理原问题与对偶问题最优值相等。去除积分约束后的LP最优值等于整数最大流（因约束矩阵全幺模），故 \(\max|f|=\min\operatorname{cap}(S,T)\)。等价地，当残余网络无增广路径时，从\(s\)可达的顶点集\(S\)构成最小割，流值等于割容量。\(\blacksquare\)

\textbf{定理 254.2}（Edmonds-Karp算法，1972）：使用BFS找最短增广路径，Ford-Fulkerson的复杂度为\(O(|V| \cdot |E|^2)\)。BFS保证每次增广路径的边数最少，从而将增广次数限制在\(O(|V|\cdot|E|)\)。

\emph{证明}：BFS 保证增广路径为最短。每次增广至少饱和一条边，且源到各顶点的最短距离不随增广减小。每条边最多被饱和 \(O(|V|)\) 次（其关键边角色交替），总增广次数 \(O(|V||E|)\)，每次 BFS \(O(|E|)\)，总复杂度 \(O(|V||E|^2)\)。\(\blacksquare\)

\subsection{254.2 最小费用流}\label{ux6700ux5c0fux8d39ux7528ux6d41}

\textbf{定义 254.3}（最小费用流 / Min-Cost Flow）：每条边\((u, v)\)有单位流费用\(a(u, v)\)。在满足流量需求\(d\)（从\(s\)到\(t\)的总流）的条件下最小化总费用\(\sum_e a(e) f(e)\)。

最小费用流问题是网络流理论中最具一般性的模型之一，它统一了最短路径问题（\(d=1\) 时退化为从 \(s\) 到 \(t\) 的最便宜路径）、最大流问题（费用为零且最大化 \(d\)）和指派问题（二分图上的匹配可表为最小费用流）。最小费用流问题的线性规划形式为：\(\min \{a^T f : Af = b, 0 \leq f \leq u\}\)，其中 \(A\) 为节点-弧关联矩阵，\(b\) 为节点供给/需求向量。由于 \(A\) 是全幺模矩阵，当 \(b\) 和 \(u\) 为整数时，基可行解自动为整数------这是最小费用流问题具有整数最优性的代数基础。求解最小费用流的经典算法包括：逐次最短路算法（Successive Shortest Path）：每次沿残量网络中的最短路（以费用为边权）增广单位流量，复杂度 \(O(d \cdot (m + n \log n))\)；费用缩放算法（Cost Scaling）：利用 \(\varepsilon\)-最优性概念，通过逐步降低 \(\varepsilon\) 来逼近最优解；以及网络单纯形法（Network Simplex）：利用图论中的树结构（基础解对应于生成树）高效执行主元操作，在实践中的平均性能远超其理论最坏情况界。

\textbf{定理 254.3}（最小费用流的最优性条件 / 负环条件）：可行流\(f\)是费用最小的当且仅当残量网络\(G_f\)中不存在负费用有向环。负环条件是线性规划最优性（对偶可行性）在图论中的等价表述。

\emph{证明}：（\(\Rightarrow\)）若 \(f\) 最优但 \(G_f\) 含负费用环 \(C\)，沿 \(C\) 增流 \(\varepsilon>0\) 费用减少 \(\varepsilon\cdot\operatorname{cost}(C)<0\)，矛盾。（\(\Leftarrow\)）设 \(f\) 使 \(G_f\) 无负环，\(f'\) 为任意可行流。\(f'-f\) 在 \(G_f\) 中分解为环流，各环费用 \(\ge 0\)，故 \(f'\) 费用 \(\ge\) \(f\) 费用。\(\blacksquare\)

\subsection{254.3 匹配问题}\label{ux5339ux914dux95eeux9898}

\textbf{定义 254.4}（匹配 / Matching）：图的匹配是两两不共享公共端点的边集合。\textbf{完美匹配}覆盖所有顶点。\textbf{最大权匹配}是权和最大的匹配。

\textbf{定理 254.4}（König定理 / 二分图匹配）：在二分图中，最大匹配的边数 = 最小顶点覆盖的点数。König定理是二分图特有的min-max关系，在一般图中不成立。

\emph{证明}：设最大匹配 \(M^*\)，从左侧未饱和顶点出发沿交替路径可达的顶点集为 \(Z\)。令 \(C = (L \setminus Z) \cup (R \cap Z)\)。可证 \(C\) 为顶点覆盖（否则存在增广路径与 \(M^*\) 最大矛盾）且 \(|C| \le |M^*|\)。又因任匹配边至少需一个覆盖点，\(|M^*| \le |C|\)，故 \(|M^*| = |C|\)。\(\blacksquare\)

\textbf{定理 254.5}（Edmonds花算法 / Blossom Algorithm，1965）：一般图的最大匹配可在多项式时间（\(O(|V|^3)\)）内求得。花算法的核心是收缩奇数长的交错环（花/blossom）------这是处理一般图（非二分图）中奇环对匹配问题造成困难的关键技术。

\textbf{证明}：一般图中奇环对增广路径搜索造成困难。算法在 BFS 增广路径搜索中加入花（blossom）收缩机制：当 BFS 发现奇环时（花的形成依赖于存在奇数长交替路径连接的两未匹配点），将整个奇环收缩为伪顶点。收缩保持匹配性质------原图有增广路径当且仅当收缩图有增广路径。证：若原图有增广路径且路径不穿过花则显然；若穿过花，因花内交替结构，可重排使路径避开花。反之，若收缩图有增广路径，展开花时利用花基点的交替路径恢复原图的增广路径。每次收缩至少减少 \(2\) 个顶点，至多 \(O(|V|)\) 次收缩。每次 BFS \(O(|E|)\)，展开恢复 \(O(|V|)\)，总复杂度 \(O(|V|^3)\)。\(\blacksquare\)

\hrulefill

\hrulefill

\hrulefill

\hrulefill

\newpage

\chapter{12-最优化与控制/02-控制与博弈/01-控制理论.md}\label{ux6700ux4f18ux5316ux4e0eux63a7ux523602-ux63a7ux5236ux4e0eux535aux5f0801-ux63a7ux5236ux7406ux8bba.md}

\section{第437章：线性控制系统}\label{ux7b2c437ux7ae0ux7ebfux6027ux63a7ux5236ux7cfbux7edf}

线性控制理论是控制论的数学基础，其核心思想是用线性微分方程（或差分方程）描述被控对象的动力学，并通过状态反馈和输出反馈实现期望的系统行为。状态空间方法的系统化归功于 Rudolf Kalman 在 1960 年代的革命性工作：他将系统的动力学行为编码为矩阵三元组 \((A, B, C)\)，并从中提炼出可控性和可观测性这两个对偶的结构性质。可控性回答''能否通过适当选择控制输入将系统从任意初始状态驱动到任意目标状态''的问题；可观测性回答''能否仅通过测量输出唯一确定系统的内部状态''的问题。Kalman 秩判据（可控性矩阵 \(\mathcal{C}=[B\ AB\ \cdots\ A^{n-1}B]\) 满秩当且仅当系统可控）是最经典的可控性代数判别准则，其在状态空间几何上等价于可控子空间与全空间重合。对偶原理（\((A, C)\) 可观测当且仅当 \((A^T, C^T)\) 可控）揭示了控制与估计之间的深刻数学对称性。分离原理（Separation Principle）进一步表明：在线性系统中，状态反馈控制器和状态观测器可以独立设计------这是线性系统理论最优美也是最实用的成果之一，为 LQG（Linear Quadratic Gaussian）控制奠定了理论基础。

\subsection{227.1 状态空间模型}\label{ux72b6ux6001ux7a7aux95f4ux6a21ux578b}

\textbf{定义 227.1}（线性时不变系统 / LTI System）：连续时间线性时不变系统的状态空间表示为

\[\begin{cases} \dot{x}(t) = A x(t) + B u(t) \\ y(t) = C x(t) + D u(t) \end{cases}\]

其中\(x(t) \in \mathbb{R}^n\)是状态，\(u(t) \in \mathbb{R}^m\)是控制输入，\(y(t) \in \mathbb{R}^p\)是观测输出。\(A \in \mathbb{R}^{n \times n}\)，\(B \in \mathbb{R}^{n \times m}\)，\(C \in \mathbb{R}^{p \times n}\)，\(D \in \mathbb{R}^{p \times m}\)。

\textbf{定义 227.2}（解与状态转移矩阵）：齐次方程\(\dot{x} = A x\)的解为\(x(t) = e^{At} x(0)\)，其中矩阵指数\(e^{At} = \sum_{k=0}^{\infty} \frac{(At)^k}{k!}\)。\(\Phi(t) = e^{At}\)称为\textbf{状态转移矩阵}。非齐次方程的通解为

\[x(t) = e^{At} x(0) + \int_0^t e^{A(t-\tau)} B u(\tau) d\tau\]

\textbf{定义 227.3}（传递函数）：系统的\textbf{传递函数}\(G(s)\)是输出和输入的Laplace变换之比（零初始条件）：

\[G(s) = C(sI - A)^{-1} B + D\]

\(G(s)\)是\(p \times m\)的有理函数矩阵。

\subsection{227.2 可控性}\label{ux53efux63a7ux6027}

\textbf{定义 227.4}（可控性 / Controllability）：系统\((A, B)\)是\textbf{可控的}（或状态\(x_0\)在时间\(T\)内可控），如果对任意初始状态\(x(0) = x_0\)和目标状态\(x_f\)，存在控制输入\(u(t)\)（\(t \in [0, T]\)）使得\(x(T) = x_f\)。

\textbf{定理 227.1}（Kalman可控性判据）：系统\((A, B)\)完全可控当且仅当\textbf{可控性矩阵}

\[\mathcal{C} = \begin{bmatrix} B & AB & A^2B & \cdots & A^{n-1}B \end{bmatrix} \in \mathbb{R}^{n \times nm}\]

满行秩：\(\operatorname{rank}(\mathcal{C}) = n\)。

\emph{注}：该判据的数学本质是：\(\mathcal{C}\)的列空间恰好是可达子空间（reachable subspace），即从原点出发可达的所有状态的集合。\(\mathcal{C}\)满行秩意味着可达子空间是整个\(\mathbb{R}^n\)。

\textbf{证明}：系统解为 \(x(T) = e^{AT}x_0 + \int_0^T e^{A(T-\tau)}Bu(\tau)d\tau\)。由 Cayley-Hamilton 定理，\(e^{A(T-\tau)} = \sum_{k=0}^{n-1} \alpha_k(\tau)A^k\)。代入得 \(x(T)-e^{AT}x_0 = \sum_{k=0}^{n-1} A^k B \int_0^T \alpha_k(\tau)u(\tau)d\tau\)。令 \(\beta_k = \int_0^T \alpha_k(\tau)u(\tau)d\tau \in \mathbb{R}^m\)，则 \(x(T)-e^{AT}x_0 = \sum_{k=0}^{n-1} A^k B \beta_k \in \operatorname{span}\{B, AB, \ldots, A^{n-1}B\} = \operatorname{Im}(\mathcal{C})\)。若 \(\operatorname{rank}(\mathcal{C}) = n\)，则 \(\operatorname{Im}(\mathcal{C}) = \mathbb{R}^n\)，对任意 \(x_f\) 存在 \(\beta_k\) 使等式成立。由 \(\beta_k\) 的构造（可通过选择适当的控制 \(u\)），任意状态可达。反之若 \(\operatorname{rank}(\mathcal{C}) < n\)，则 \(\operatorname{Im}(\mathcal{C})\) 为 \(\mathbb{R}^n\) 的真子空间，存在不可达状态。\(\blacksquare\)

\textbf{定理 227.2}（Hautus引理 / PBH检验）：\((A, B)\)可控当且仅当对所有\(\lambda \in \mathbb{C}\)，

\[\operatorname{rank}\begin{bmatrix} \lambda I - A & B \end{bmatrix} = n\]

只需对\(A\)的特征值验证即可。PBH检验将可控性归结为：\(A\)的每个特征向量不能同时与\(B\)的所有列正交。

\textbf{证明}：（\(\Rightarrow\)）若 \((A,B)\) 可控但存在 \(\lambda\) 使 \(\operatorname{rank}[\lambda I-A, B] < n\)，则存在非零 \(v^T\) 满足 \(v^T[\lambda I-A, B] = 0\)，即 \(v^T A = \lambda v^T\) 且 \(v^T B = 0\)。于是 \(v^T A^k B = \lambda^k v^T B = 0\) 对所有 \(k\ge 0\)，故 \(v^T \mathcal{C} = 0\)，与 \(\operatorname{rank}(\mathcal{C})=n\) 矛盾。（\(\Leftarrow\)）若不可控，\(\operatorname{rank}(\mathcal{C})<n\)，存在非零 \(v\) 满足 \(v^T \mathcal{C}=0\)。由 Cayley-Hamilton，\(v^T A^k B=0\) 对所有 \(k\)。定义 \(W = \operatorname{span}\{v, v^T A, v^T A^2, \ldots\}\)，\(W\) 为 \(A^T\) 的不变子空间。取 \(W\) 中 \(A^T\) 的特征向量 \(w\) 对应特征值 \(\lambda\)，则 \(w^T B=0\) 且 \(w^T(\lambda I-A)=0\)，故 \(\operatorname{rank}[\lambda I-A, B] < n\)。\(\blacksquare\)

\textbf{定理 227.3}（Kalman可控性分解）：若\(\operatorname{rank}(\mathcal{C}) = r < n\)，存在坐标变换将系统分解为

\[\begin{bmatrix} \dot{x}_c \\ \dot{x}_{\bar{c}} \end{bmatrix} = \begin{bmatrix} A_c & A_{12} \\ 0 & A_{\bar{c}} \end{bmatrix} \begin{bmatrix} x_c \\ x_{\bar{c}} \end{bmatrix} + \begin{bmatrix} B_c \\ 0 \end{bmatrix} u\]

其中\((A_c, B_c)\)是可控子系统，\(A_{\bar{c}}\)对应不可控模式。该分解的数学构造为：取\(\mathcal{C}\)的列空间的一组基作为可控子空间的基，再扩张为整个\(\mathbb{R}^n\)的基。

\textbf{证明}：可控子空间 \(\mathcal{R} = \operatorname{Im}(\mathcal{C})\) 是 \(A\)-不变的（\(A\mathcal{R} \subseteq \mathcal{R}\)）且包含 \(\operatorname{Im}(B)\)。取基 \(v_1,\ldots,v_r\) 于 \(\mathcal{R}\)，扩张为 \(\mathbb{R}^n\) 的全基 \(v_1,\ldots,v_n\)。令 \(T = [v_1,\ldots,v_n]^{-1}\) 为新坐标变换矩阵。在新基下，\(A\) 的表示 \(\tilde{A}=TAT^{-1}\) 满足块下三角形式：因 \(A\mathcal{R}\subseteq\mathcal{R}\)，\(A\) 将 \(\mathcal{R}\) 映到 \(\mathcal{R}\)，故 \(\tilde{A}_{21}=0\)。\(B\) 的表示 \(\tilde{B}=TB\) 满足 \(\tilde{B}_2=0\)：因 \(\operatorname{Im}(B)\subseteq\mathcal{R}\)，\(B\) 的列在前 \(r\) 个基向量张成的空间中。故 \((A_c,B_c)\) 为可控对（\(\tilde{A}_{11},\tilde{B}_1\)），维数 \(r\)，\(\tilde{A}_{22}\) 对应不可控模式。\(\blacksquare\)

\subsection{227.3 可观测性}\label{ux53efux89c2ux6d4bux6027}

\textbf{定义 227.5}（可观测性 / Observability）：系统\((A, C)\)是\textbf{可观测的}，如果对任意初始状态\(x(0) = x_0\)，能够在有限时间\(T\)内仅通过观测输出\(y(t)\)和输入\(u(t)\)（\(t \in [0, T]\)）唯一确定\(x_0\)。

\textbf{定理 227.4}（Kalman可观测性判据）：系统\((A, C)\)完全可观测当且仅当\textbf{可观测性矩阵}

\[\mathcal{O} = \begin{bmatrix} C \\ CA \\ CA^2 \\ \vdots \\ CA^{n-1} \end{bmatrix} \in \mathbb{R}^{np \times n}\]

满列秩：\(\operatorname{rank}(\mathcal{O}) = n\)。

\textbf{证明}：由 \(y(t) = Ce^{At}x_0 + \int_0^t Ce^{A(t-\tau)}Bu(\tau)d\tau\)，已知输入可重构 \(Ce^{At}x_0\)。Cayley-Hamilton 给出 \(e^{At}\) 的各行在 \(\operatorname{span}\{C,CA,\ldots,CA^{n-1}\}\) 中，故 \(Ce^{At}x_0 = \sum_{k=0}^{n-1} \gamma_k(t) CA^k x_0\)。若 \(\mathcal{O}\) 满列秩（\(\ker(\mathcal{O}) = \{0\}\)），则从 \(CA^k x_0\)（\(k=0,\ldots,n-1\)）可唯一恢复 \(x_0\)。若 \(\operatorname{rank}(\mathcal{O}) < n\)，则存在非零 \(x_0\in\ker(\mathcal{O})\) 满足 \(CA^k x_0=0\) 对所有 \(k\)，此时 \(y(t)\equiv 0\)（零输入）无法区分 \(x_0\) 与 \(0\)，故不可观测。\(\blacksquare\)

\textbf{定理 227.5}（对偶原理）：\((A, C)\)可观测当且仅当对偶系统\((A^T, C^T)\)可控。这是最优控制和最优估计对偶性的代数基础------可控性与可观测性互为对偶概念。

\textbf{证明}：可观测矩阵 \(\mathcal{O}\) 的行恰为对偶系统可控性矩阵 \(\mathcal{C}_{\text{dual}} = [C^T \; A^T C^T \; \ldots \; (A^T)^{n-1}C^T]\) 的转置。故 \(\operatorname{rank}(\mathcal{O}) = n \iff \operatorname{rank}(\mathcal{O}^T) = n \iff \operatorname{rank}(\mathcal{C}_{\text{dual}}) = n\)，即 \((A,C)\) 可观测等价于 \((A^T, C^T)\) 可控。对偶原理的深层含义是：状态估计（Kalman 滤波）与状态反馈（LQR 控制）在数学结构上完全对偶，滤波误差协方差满足与 LQR 值函数矩阵相同的 Riccati 方程，只是时间方向相反。\(\blacksquare\)

\subsection{227.4 极点配置与镇定}\label{ux6781ux70b9ux914dux7f6eux4e0eux9547ux5b9a}

\textbf{定理 227.6}（极点配置定理，Wonham 1967）：若\((A, B)\)可控，则对任意给定的\(n\)个复共轭对\(p_1, \ldots, p_n\)，存在状态反馈\(u = -Kx\)使得闭环系统\(\dot{x} = (A - BK)x\)的特征值为\(\{p_1, \ldots, p_n\}\)。

极点配置定理的数学本质是：可控性使得状态反馈能够任意改变闭环系统的动力学（在复共轭约束下）。

\textbf{证明}：（单输入情形）由可控性，存在可逆矩阵 \(T\) 将 \((A,b)\) 变换为可控标准形：\(\tilde{A}=TAT^{-1}\) 为伴随矩阵形式，\(\tilde{b}=Tb=(0,\ldots,0,1)^T\)。取 \(K=(k_0,\ldots,k_{n-1})\)，则 \(\tilde{A}-\tilde{b}K\) 也是伴随矩阵，其特征多项式为 \(\lambda^n+(a_{n-1}+k_{n-1})\lambda^{n-1}+\cdots+(a_0+k_0)\)。给定期望极点 \(p_1,\ldots,p_n\)，构造期望特征多项式 \(\prod_{i=1}^n(\lambda-p_i)=\lambda^n+c_{n-1}\lambda^{n-1}+\cdots+c_0\)，令 \(k_i=c_i-a_i\) 即得反馈增益，原坐标下 \(K=\tilde{K}T\)。（多输入）由 Heymann 引理，存在 \(v\in\mathbb{R}^m\) 使 \((A,Bv)\) 单输入可控，归约为单输入情形。\(\blacksquare\)

\textbf{定义 227.6}（可镇定性与可检测性）：\((A, B)\)\textbf{可镇定}（stabilizable）如果所有不可控模式都是稳定的（特征值实部为负）。\((A, C)\)\textbf{可检测}（detectable）如果所有不可观测模式都是稳定的。

\textbf{定理 227.7}（分离原理 / Separation Principle）：对可镇定且可检测的系统，状态观测器和状态反馈控制器可以独立设计，闭环系统保持稳定。分离原理是线性系统理论中估计与控制解耦的核心定理。

\emph{证明}：设观测器 \(\dot{\hat{x}} = A\hat{x} + Bu + L(y-C\hat{x})\)，误差 \(e = x-\hat{x}\) 满足 \(\dot{e} = (A-LC)e\)。闭环系统 \(\begin{pmatrix} \dot{x} \\ \dot{e} \end{pmatrix} = \begin{pmatrix} A-BK & BK \\ 0 & A-LC \end{pmatrix} \begin{pmatrix} x \\ e \end{pmatrix}\)，特征值为 \(\operatorname{eig}(A-BK) \cup \operatorname{eig}(A-LC)\)。可镇定/可检测性保证可独立选 \(K\) 和 \(L\) 使 \(A-BK\) 和 \(A-LC\) 均稳定，分解块三角结构保证全局稳定。\(\blacksquare\)

\hrulefill

\hrulefill

\hrulefill

\hrulefill

\section{第438章：最优控制}\label{ux7b2c438ux7ae0ux6700ux4f18ux63a7ux5236}

最优控制理论是变分法在动态系统控制中的自然推广，其核心问题是：在满足动力学约束 \(\dot{x}=f(t,x,u)\) 的条件下，寻找使性能泛函 \(J[u]=\phi(x(t_f))+\int_{t_0}^{t_f}L(t,x,u)dt\) 达到极值的控制策略 \(u(t)\)。这一理论在 20 世纪 50-60 年代由 Pontryagin 学派和 Bellman 独立地以两种等价但视角不同的形式建立：Pontryagin 极小值原理（1956）从变分法的角度出发，引入协态变量 \(\lambda(t)\)（Lagrange 乘子），将最优控制问题转化为 Hamilton 系统的两点边值问题；Bellman 最优性原理（1957）则从动态规划的角度出发，引入值函数 \(V(t,x)\)（从状态 \(x\) 和时刻 \(t\) 出发的最优代价），导出了 Hamilton-Jacobi-Bellman（HJB）偏微分方程。这两种方法的数学等价性体现在协态 \(\lambda(t)\) 恰为值函数的梯度 \(\nabla_x V(t,x^*(t))\)------即极小值原理是沿最优轨线的特征线方程，HJB 方程则是全局的偏微分方程。在线性二次型调节器（LQR）的特殊情形下，最优控制律可以显式求得 \(u=-R^{-1}B^TPx\)，其中 \(P\) 满足代数 Riccati 方程------这一结果是线性系统最优控制的基石，其可解性由系统可镇定性和可检测性保证。

\subsection{228.1 Pontryagin极小值原理}\label{pontryaginux6781ux5c0fux503cux539fux7406}

\textbf{定义 228.1}（最优控制问题 / Bolza形式）：考虑受控动力系统\(\dot{x}(t) = f(t, x(t), u(t))\)，初始条件\(x(t_0) = x_0\)。寻找控制\(u(t) \in U \subseteq \mathbb{R}^m\)最小化性能泛函

\[J[u] = \phi(x(t_f)) + \int_{t_0}^{t_f} L(t, x(t), u(t)) dt\]

\textbf{定理 228.1}（Pontryagin极小值原理，1956）：若\(u^*(t)\)是最优控制，\(x^*(t)\)是最优轨线，则存在非零伴随向量\(\lambda(t)\)（协态）满足：

\begin{enumerate}
\def\labelenumi{\arabic{enumi}.}
\tightlist
\item
  \textbf{状态方程}：\(\dot{x}^* = \frac{\partial H}{\partial \lambda}(t, x^*, u^*, \lambda)\)
\item
  \textbf{协态方程}：\(\dot{\lambda} = -\frac{\partial H}{\partial x}(t, x^*, u^*, \lambda)\)
\item
  \textbf{极小值条件}：\(H(t, x^*, u^*, \lambda) \leq H(t, x^*, u, \lambda)\)（对所有\(u \in U\)）
\item
  \textbf{横截条件}：\(\lambda(t_f) = \frac{\partial \phi}{\partial x}(x^*(t_f))\)
\end{enumerate}

其中\textbf{Hamilton函数}\(H(t, x, u, \lambda) = L(t, x, u) + \lambda^T f(t, x, u)\)。

\textbf{证明}：引入 Lagrange 乘子（协态）\(\lambda(t)\in\mathbb{R}^n\)，将动态约束 \(\dot{x}=f(t,x,u)\) 并入泛函： \[\tilde{J}= \phi(x(t_f)) + \int_{t_0}^{t_f} \left[L(t,x,u) + \lambda^T(t)(f(t,x,u)-\dot{x})\right]dt.\] 定义 Hamilton 函数 \(H(t,x,u,\lambda)=L(t,x,u)+\lambda^T f(t,x,u)\)。分部积分 \(\int -\lambda^T\dot{x}\,dt = -\lambda^T x|_{t_0}^{t_f} + \int \dot{\lambda}^T x\,dt\)，得 \[\tilde{J}= \phi(x(t_f))-\lambda^T(t_f)x(t_f)+\lambda^T(t_0)x(t_0) + \int_{t_0}^{t_f}\left[H+\dot{\lambda}^T x\right]dt.\] 取一阶变分 \(\delta\tilde{J}\)：(i) 对 \(\delta x(t_f)\)：\(\partial\phi/\partial x-\lambda^T(t_f)=0\) 给出横截条件；(ii) 对 \(\delta x(t)\)：\(\partial H/\partial x + \dot{\lambda}^T = 0\) 给出协态方程 \(\dot{\lambda}=-\partial H/\partial x\)；(iii) 对 \(\delta\lambda\)：恢复 \(\dot{x}=\partial H/\partial\lambda\)；(iv) 对 \(\delta u\)：对任意容许变分 \(\delta u\)，变分不等式给出 \(H(t,x^*,u^*,\lambda)\le H(t,x^*,u,\lambda)\)（对所有 \(u\in U\)）。这四条即 Pontryagin 极小值原理。控制约束 \(u\in U\) 为非开集时，一阶条件退化为逐点极小化条件而非梯度为零。严格证明需用针状变分（needle variation）和Gronwall不等式。\(\blacksquare\)

\subsection{228.2 Hamilton-Jacobi-Bellman方程}\label{hamilton-jacobi-bellmanux65b9ux7a0b}

\textbf{定义 228.2}（值函数 / Value Function）：从状态\(x\)和时间\(t\)出发的最优代价（值函数）为

\[V(t, x) = \min_{u(\cdot)} \left\{ \phi(x(t_f)) + \int_t^{t_f} L(s, x(s), u(s)) ds \right\}\]

\textbf{定理 228.2}（Bellman最优性原理，1957）：最优策略的剩余部分仍是剩余问题的最优策略。用值函数表示为

\[V(t, x) = \min_{u \in U} \left\{ L(t, x, u) \Delta t + V(t + \Delta t, x + f(t, x, u) \Delta t) \right\} + o(\Delta t)\]

\textbf{证明}：将最优问题分解为 \([t,t+\Delta t]\) 上的瞬时代价和 \([t+\Delta t,t_f]\) 上的未来代价。由值函数定义，对任意策略 \(u\)， \[V(t,x)\le \int_t^{t+\Delta t}L(s,x(s),u(s))ds + V(t+\Delta t,x(t+\Delta t)).\] 对 \(u^*\)（最优策略）取等号。由连续性 \(x(t+\Delta t)=x+f(t,x,u)\Delta t+o(\Delta t)\)，Taylor展开 \(V(t+\Delta t,\cdot)\)，取 \(\Delta t\to 0\) 即得最优性原理的微分形式。此原理的实质是：动态系统的最优策略具有时间一致性------过去决策不影响当前最优选择。\(\blacksquare\)

\textbf{定理 228.3}（Hamilton-Jacobi-Bellman方程 / HJB方程）：若\(V \in C^1\)，则满足

\[-\frac{\partial V}{\partial t}(t, x) = \min_{u \in U} \left\{ L(t, x, u) + \nabla_x V(t, x)^T f(t, x, u) \right\}\]

边界条件\(V(t_f, x) = \phi(x)\)。HJB方程是一阶非线性偏微分方程------它是经典力学中Hamilton-Jacobi方程在最优控制背景下的推广。

\textbf{证明}：由 Bellman 最优性原理 \[V(t,x)=\min_{u\in U}\left[L(t,x,u)\Delta t + V(t+\Delta t,x+f(t,x,u)\Delta t+o(\Delta t))\right].\] Taylor 展开 \(V(t+\Delta t,\cdot)=V(t,x)+\partial_t V\Delta t+\nabla_x V\cdot f\Delta t+o(\Delta t)\)，代入得 \[V(t,x)=\min_{u\in U}\left[V(t,x)+L\Delta t+\partial_t V\Delta t+\nabla_x V\cdot f\Delta t+o(\Delta t)\right].\] 约去 \(V(t,x)\)（与 \(u\) 无关），除以 \(\Delta t\)，令 \(\Delta t\to 0^+\) 即得 HJB 方程。边界条件 \(V(t_f,x)=\phi(x)\) 来自值函数定义：终端时刻无控制选择，直接付终端代价。\(\blacksquare\)

\textbf{定理 228.4}（验证定理 / Verification Theorem）：

\emph{证明}：对任意容许控制 \(u\) 和轨线，\(\frac{d}{dt}V(t,x(t))=\partial_t V+\nabla_x V\cdot f\)。由 HJB 方程，\(\partial_t V+\nabla_x V\cdot f+L \ge 0\)（对所有 \(u\)），在 \(u^*\) 处等号成立。积分得 \(V(0,x_0)\le \int_0^T L ds+\phi(x(T))\)，\(u^*\) 达到下界，故最优。\(\blacksquare\)若光滑函数\(V(t, x)\)满足HJB方程且\(u^*(t, x)\)达到最小值，则\(u^*\)是最优反馈控制。

\textbf{注 228.1}（极小值原理与HJB的关系）：极小值原理是HJB方程沿最优轨线的特征线形式。协态\(\lambda(t) = \nabla_x V(t, x^*(t))\)------这是Pontryagin极小值原理与动态规划之间最深刻的数学联系。

\subsection{228.3 无限时域最优控制}\label{ux65e0ux9650ux65f6ux57dfux6700ux4f18ux63a7ux5236}

\textbf{定义 228.3}（无限时域LQR）：最小化\(J = \int_0^\infty (x^T Q x + u^T R u) dt\)（\(Q \geq 0, R > 0\)）。最优控制为\(u = -K x\)，\(K = R^{-1} B^T P\)，其中\(P\)满足\textbf{代数Riccati方程}：

\[PA + A^T P - P B R^{-1} B^T P + Q = 0\]

\textbf{定理 228.5}（代数Riccati方程的可解性）：若\((A, B)\)可镇定且\((Q^{1/2}, A)\)可检测，则代数Riccati方程存在唯一正半定解\(P \geq 0\)，且闭环系统\(A - BK\)是渐近稳定的。该定理将最优控制的可解性与系统的结构性质（可控性/可观测性）联系起来。

\textbf{证明}：考虑 Lyapunov 函数 \(V(x)=x^T P x\)（\(P\ge 0\) 为 ARE 解）。沿闭环轨线 \(\dot{x}=(A-BR^{-1}B^T P)x\)， \[\begin{aligned}
\dot{V} &= x^T P(A-BK)x + x^T(A-BK)^T P x \\
&= x^T(PA+A^T P - PBR^{-1}B^T P - (BR^{-1}B^T P)^T P)x \\
&= x^T(PA+A^T P - 2PBR^{-1}B^T P)x \\
&= -x^T(Q + PBR^{-1}B^T P)x \le -x^T Q x \le 0.
\end{aligned}\] 由 \((Q^{1/2},A)\) 可检测性和 LaSalle 不变原理，\(\dot{V}=0\) 的最大不变集仅含原点，故闭环系统渐近稳定。\(P\) 的唯一性：设 \(P_1,P_2\) 为两解，差值 \(\Delta=P_1-P_2\) 满足代数方程，通过 Lyapunov 论证得 \(\Delta=0\)。解的存在性由可镇定条件保证 Hamiltonian 矩阵无虚轴特征值，从而 Riccati 方程有镇定解。\(\blacksquare\)

\hrulefill

\hrulefill

\hrulefill

\hrulefill

\section{第439章：随机滤波与控制}\label{ux7b2c439ux7ae0ux968fux673aux6ee4ux6ce2ux4e0eux63a7ux5236}

随机系统的最优估计与控制是将概率论和随机分析与最优控制理论相结合的交叉领域。当系统动力学和观测被噪声污染时，确定性最优控制的方法不再直接适用，而必须引入滤波（状态估计）作为与最优控制对称且互补的组成部分。Kalman-Bucy 滤波器（1960）是这个领域最辉煌的成就之一：在线性 Gauss 系统的假设下，状态的条件期望 \(\hat{x}(t)=\mathbb{E}[x(t)|\mathcal{F}_t^y]\) 可以通过一个线性递推方程精确计算，滤波增益 \(K=PC^TR^{-1}\) 完全由误差协方差矩阵 \(P\) 的 Riccati 微分方程决定。Kalman 滤波的数学本质是 Hilbert 空间中正交集上的投影定理------在高斯分布下，非线性条件期望奇迹般地退化为线性投影。随机最优控制的完整框架由分离原理（Separation Principle）统一：在线性二次型 Gauss（LQG）问题中，最优控制 \(u=-K\hat{x}\) 可以分解为两个独立子问题------确定性 LQR 的反馈增益 \(K\)（由控制 Riccati 方程求解）和 Kalman-Bucy 滤波的状态估计 \(\hat{x}\)（由滤波 Riccati 方程求解）。这一分离结构是线性 Gauss 系统中估计与控制的最深刻数学性质，其证明依赖于 Itô 随机分析的 Ito 公式和鞅表示定理。

\subsection{229.1 Kalman-Bucy滤波}\label{kalman-bucyux6ee4ux6ce2}

\textbf{定义 229.1}（线性随机系统）：具有过程噪声和观测噪声的线性系统

\[\begin{cases} \dot{x}(t) = A x(t) + B u(t) + w(t) \\ y(t) = C x(t) + v(t) \end{cases}\]

其中\(w(t)\)是过程噪声（协方差\(Q \geq 0\)），\(v(t)\)是观测噪声（协方差\(R > 0\)），二者是不相关的Gauss白噪声。

\textbf{定义 229.2}（Kalman-Bucy滤波器，1960）：状态估计\(\hat{x}(t)\)由以下微分方程更新（连续时间Kalman-Bucy滤波器）：

\[\dot{\hat{x}} = A \hat{x} + B u + K (y - C \hat{x})\]

其中\textbf{Kalman增益}\(K = P C^T R^{-1}\)，且\textbf{误差协方差}\(P(t) = \mathbb{E}[(x - \hat{x})(x - \hat{x})^T]\)满足\textbf{滤波Riccati方程}：

\[\dot{P} = A P + P A^T - P C^T R^{-1} C P + Q\]

\textbf{定理 229.1}（Kalman滤波的最优性）：Kalman-Bucy滤波器是线性系统的最小方差无偏估计。在高斯噪声假设下，\(\hat{x}(t)\)是状态\(x(t)\)在观测\(\sigma\)-代数\(\sigma\{y(s) : s \leq t\}\)下的条件期望。该定理的数学本质是：线性Gauss系统中，条件期望由线性递推精确给出------这是正交集上投影定理（Hilbert空间理论）的随机版本。

\textbf{证明}：设 \(\hat{x}(t)=\mathbb{E}[x(t)|\mathcal{F}_t^y]\)（\(y\) 生成的 \(\sigma\)-代数下的条件期望）。在高斯假设下，条件期望是线性的：\(\hat{x}(t)\) 满足线性随机微分方程。误差 \(\tilde{x}=x-\hat{x}\) 与观测 \(\mathcal{F}_t^y\) 正交（由投影定理，\(\mathbb{E}[\tilde{x}(t)y(s)^T]=0\)，\(\forall s\le t\)）。将误差动力学 \(\dot{\tilde{x}}=(A-KC)\tilde{x}+w-Kv\) 代入协方差演化，得 \(P(t)=\mathbb{E}[\tilde{x}\tilde{x}^T]\) 满足 Riccati 方程 \(\dot{P}=AP+PA^T-PC^T R^{-1}CP+Q\)。增益 \(K=PC^T R^{-1}\) 是使得 \(\operatorname{Tr}(P)\) 最小的最优增益（通过导数 \(\partial \operatorname{Tr}(\dot{P})/\partial K=2(KR-CP)=0\) 验证）。在稳态下 \(\dot{P}=0\) 得滤波 ARE。由正交投影的 Gauss-Markov 定理，此估计在 \(L^2\) 意义下最优。\(\blacksquare\)

\subsection{229.2 随机最优控制}\label{ux968fux673aux6700ux4f18ux63a7ux5236}

\textbf{定义 229.3}（随机最优控制问题）：状态方程由Itô随机微分方程\(dx_t = f(t, x_t, u_t) dt + \sigma(t, x_t) dW_t\)给出，\(W_t\)是Brown运动。性能指标为\(J = \mathbb{E}[\phi(x_T) + \int_0^T L(t, x_t, u_t) dt]\)。

\textbf{定理 229.2}（随机HJB方程）：值函数\(V(t, x)\)满足

\[-\frac{\partial V}{\partial t} = \min_{u \in U} \left\{ L(t, x, u) + f^T \nabla_x V + \frac{1}{2} \operatorname{Tr}(\sigma \sigma^T \nabla_x^2 V) \right\}\]

与确定性HJB方程相比，随机HJB方程多了二阶项\(\frac{1}{2}\operatorname{Tr}(\sigma\sigma^T \nabla_x^2 V)\)------这是Itô公式的直接推论，体现了扩散过程的二阶变差贡献。

\textbf{证明}：设价值函数 \(V(t,x)=\min_{u}\mathbb{E}[\phi(x_T)+\int_t^T L(s,x_s,u_s)ds|\mathcal{F}_t]\)。由 Ito 公式， \[dV = \left(V_t + V_x^T f + \tfrac{1}{2}\operatorname{Tr}(\sigma\sigma^T V_{xx})\right)dt + V_x^T\sigma\,dW.\] 代入动态规划方程 \(V(t,x)=\min_u\mathbb{E}[V(t+dt,x+dx)+L(t,x,u)dt]\)，对 \(dt\) 一阶展开，联立 \(\mathbb{E}[dW]=0\) 消去鞅部分，令 \(dt\to 0\) 得 \[-V_t = \min_u\left[L(t,x,u)+V_x^T f(t,x,u)+\tfrac{1}{2}\operatorname{Tr}\left(\sigma(t,x,u)\sigma(t,x,u)^T V_{xx}\right)\right], \qquad V(T,x)=\phi(x).\] 此即随机 HJB 方程。相比于确定性情形，扩散项增加了二阶 Hessian 项，对应于方差对价值函数的惩罚。在分离定理适用的 LQG 情形，\(V(t,x)=x^T P(t)x + c(t)\) 代入 HJB 后可分离出确定性 LQR 和滤波两部分。\(\blacksquare\)

\textbf{定理 229.3}（分离原理 / LQG控制）：对线性随机系统（LQG：Linear Quadratic Gaussian），最优控制\(u = -K \hat{x}\)分别由确定性LQR的最优增益\(K\)和Kalman-Bucy滤波的状态估计\(\hat{x}\)独立决定。分离原理将随机最优控制分解为两个独立子问题：确定性最优控制（Riccati方程）和最优状态估计（滤波Riccati方程）。

\emph{证明}：LQG 最优控制 \(u=-K\hat{x}\)（\(K=R^{-1}B^T P\)）。代价分解为 \(\mathbb{E}[x^T Q x+u^T R u]=\mathbb{E}[\hat{x}^T(Q+K^T RK)\hat{x}]+\operatorname{Tr}(P\Sigma)\)。第一项由 Kalman 滤波器最小化，第二项由 LQR 增益确定。误差动力学与 \(K\) 独立，分离成立。\(\blacksquare\)

\hrulefill

\hrulefill

\hrulefill

\hrulefill

\section{第440章：鲁棒控制}\label{ux7b2c440ux7ae0ux9c81ux68d2ux63a7ux5236}

鲁棒控制（Robust Control）是 1980 年代控制理论的重大范式转变，其核心问题是：在实际工程系统中，被控对象的数学模型永远存在不确定性------参数误差、未建模动态、外部扰动------如何设计控制器使得系统在这些不确定性下仍能保持稳定并满足性能要求？这一问题的数学框架建立在算子理论和 Hardy 空间 \(\mathcal{H}_\infty\) 上。\(\mathcal{H}_\infty\) 范数 \(\|G\|_\infty=\sup_\omega\bar{\sigma}(G(j\omega))\) 度量了系统从 \(L^2\) 输入到 \(L^2\) 输出的最坏情况能量增益，小增益定理（Small Gain Theorem）将闭环鲁棒稳定性归约为开环增益小于 1------这是压缩映射原理在频域的体现。\(\mathcal{H}_\infty\) 最优控制的完整解法由 Doyle、Glover、Khargonekar 和 Francis（DGKF，1989）给出：通过引入广义被控对象和 Youla 参数化，将 \(\mathcal{H}_\infty\) 次优控制问题归结为两个代数 Riccati 方程的可解性，控制器由中心解给出。线性矩阵不等式（LMI）方法则将 \(\mathcal{H}_\infty\)、\(\mathcal{H}_2\)、极点配置等多种控制问题统一为凸可行性问题：Schur 补引理将非线性的 Riccati 不等式转化为线性的 LMI，从而可以用内点法在多项式时间内求解。

\subsection{230.1 小增益定理与不确定性描述}\label{ux5c0fux589eux76caux5b9aux7406ux4e0eux4e0dux786eux5b9aux6027ux63cfux8ff0}

\textbf{定义 230.1}（\(\mathcal{H}_\infty\)范数）：稳定传递函数\(G(s)\)的\textbf{\(\mathcal{H}_\infty\)范数}为

\[\|G\|_\infty = \sup_{\omega \in \mathbb{R}} \bar{\sigma}(G(j\omega))\]

其中\(\bar{\sigma}(\cdot)\)是最大奇异值。\(\|G\|_\infty\)是\(L^2\)输入到\(L^2\)输出的诱导算子范数------它是Hardy空间\(\mathcal{H}_\infty\)上的范数，本质上是对系统最坏情况能量增益的度量。

\textbf{定理 230.1}（小增益定理 / Small Gain Theorem）：设\(M\)和\(\Delta\)为稳定传递函数。若\(\|M \Delta\|_\infty < 1\)，则反馈互联系统稳定。小增益定理是压缩映射原理在频域的体现------开环增益的\(\mathcal{H}_\infty\)范数小于1保证闭环映射是压缩的。

\textbf{证明}：考虑反馈互联 \((I-M\Delta)^{-1}\)。对任意 \(u\)，设 \(y=Mu\)，\(w=\Delta y\) 是开环映射。闭环满足 \(y=M(w+v)\)（\(v\) 是外部输入）。若 \(\|M\Delta\|_\infty<1\)，则由 Neumann 级数 \((I-M\Delta)^{-1}=\sum_{k=0}^\infty(M\Delta)^k\) 在 \(\mathcal{H}_\infty\) 范数下绝对收敛，故闭环传递函数稳定。等价地，由压缩映射原理，\(L^2\) 增益 \(\|(I-M\Delta)^{-1}\|_{L^2\to L^2}\le 1/(1-\|M\Delta\|_\infty)\)。\(\blacksquare\)

\textbf{定义 230.2}（结构奇异值 / \(\mu\)）：对结构化不确定性\(\Delta = \operatorname{diag}(\Delta_1, \ldots, \Delta_k)\)，定义

\[\mu_\Delta(M) = \frac{1}{\min\{\bar{\sigma}(\Delta) : \det(I - M\Delta) = 0\}}\]

\(\mu < 1\)等价于对所有允许的\(\|\Delta\|_\infty \leq 1\)，系统保持鲁棒稳定。\(\mu\)分析将结构化不确定性的鲁棒稳定性问题归约为对\(\mu\)函数的上界计算。

\subsection{\texorpdfstring{230.2 \(\mathcal{H}_\infty\)最优控制}{230.2 \textbackslash mathcal\{H\}\_\textbackslash infty最优控制}}\label{mathcalh_inftyux6700ux4f18ux63a7ux5236}

\textbf{定义 230.3}（\(\mathcal{H}_\infty\)控制问题）：考虑广义被控对象

\[\begin{bmatrix} z \\ y \end{bmatrix} = \begin{bmatrix} P_{11} & P_{12} \\ P_{21} & P_{22} \end{bmatrix} \begin{bmatrix} w \\ u \end{bmatrix}\]

其中\(w\)是外部扰动，\(u\)是控制，\(z\)是被控输出，\(y\)是测量。寻找控制器\(K(s)\)使得闭环传递函数\(T_{zw}(s)\)满足\(\|T_{zw}\|_\infty < \gamma\)。

\textbf{定理 230.2}（Doyle-Glover-Khargonekar-Francis解法 / DGKF 1989）：\(\mathcal{H}_\infty\)次优控制问题可归结为两个代数Riccati方程（一个用于状态反馈，一个用于输出估计）的可解性。控制器由中心解给出。DGKF解法将\(\mathcal{H}_\infty\)控制统一为Riccati方程框架------与LQG问题在结构上形成精确的对应。

\textbf{证明}：设广义被控对象状态空间实现 \(\dot{x}=Ax+B_1w+B_2u\)，\(z=C_1x+D_{12}u\)，\(y=C_2x+D_{21}w\)。构造 Hamiltonian 矩阵 \(H_\infty\) 并求解控制 Riccati 方程 \(A^T X+XA+X(\gamma^{-2}B_1B_1^T-B_2R^{-1}B_2^T)X+C_1^T C_1=0\)（\(X\ge 0\)），滤波 Riccati 方程 \(AY+YA^T+Y(\gamma^{-2}C_1^T C_1-C_2^T\bar{R}^{-1}C_2)Y+B_1B_1^T=0\)（\(Y\ge 0\)）。次优 \(\mathcal{H}_\infty\) 控制器存在的充要条件是两 ARE 均有镇定解且 \(\rho(XY)<\gamma^2\)（耦合条件），控制器为 \(K(s)=\begin{bmatrix}A_\infty&Z_\infty L_\infty\\F_\infty&0\end{bmatrix}\)。耦合条件 \(\rho(XY)<\gamma^2\) 保证分离结构的可控性。\(\blacksquare\)

\textbf{定理 230.3}（有界实引理 / Bounded Real Lemma）：\(\|G\|_\infty < \gamma\)当且仅当Hamiltonian矩阵没有虚轴上的特征值，等价于某线性矩阵不等式（LMI）可解。有界实引理将频域范数约束转化为时域LMI条件。

\textbf{证明}：设 \(G(s)=C(sI-A)^{-1}B+D\)（\(A\) 稳定）。\(\|G\|_\infty<\gamma\) 等价于 \(\|\gamma^{-1}G\|_\infty<1\)。由 Parseval 定理和 Kalman-Yakubovich-Popov 引理，存在 \(P=P^T>0\) 满足 \[\begin{bmatrix}A^T P+PA & PB\\ B^T P & -\gamma I\end{bmatrix} + \begin{bmatrix}C^T\\ D^T\end{bmatrix}\gamma^{-1}I\begin{bmatrix}C & D\end{bmatrix}<0.\] 应用 Schur 补得等价 LMI： \[\begin{bmatrix}A^T P+PA & PB & C^T\\ B^T P & -\gamma I & D^T\\ C & D & -\gamma I\end{bmatrix}<0.\] 这与 Hamiltonian 矩阵 \(H=\begin{bmatrix}A& \gamma^{-2}BB^T\\ -C^T C & -A^T\end{bmatrix}\) 无纯虚特征值等价。\(\blacksquare\)

\subsection{230.3 线性矩阵不等式}\label{ux7ebfux6027ux77e9ux9635ux4e0dux7b49ux5f0f}

\textbf{定义 230.4}（线性矩阵不等式 / LMI）：\textbf{线性矩阵不等式}形如\(F(x) = F_0 + \sum_{i=1}^m x_i F_i \geq 0\)（半正定），其中\(F_i\)为对称矩阵。LMI定义了\(\mathbb{R}^m\)中的凸集（半正定锥的截面）。

\textbf{定理 230.4}（LMI方法的统一性）：众多控制问题（\(\mathcal{H}_\infty\)控制、\(\mathcal{H}_2\)控制、极点配置、\(\mu\)-综合的上界）都可转化为LMI可行性问题。LMI是凸优化问题，可用内点法求解------这为控制理论提供了数值可处理的统一框架。

\textbf{证明}：(i) \(\mathcal{H}_\infty\)：有界实引理已将其表述为 LMI。(ii) \(\mathcal{H}_2\)：迹最小化 \(\operatorname{Tr}(CPC^T)\) 约束于 Lyapunov 不等式 \(AP+PA^T+BB^T\le 0\) 是半定规划。(iii) 极点配置：区域 \(D=\{z:f_D(s)<0\}\)（\(f_D\) 为 LMI 可表区域）由 Lyapunov 型 LMI 刻画。(iv) \(\mu\)-综合：\(D\) 标度（\(D\)-scaling）上界给出凸优化保守近似。所有问题均约化为 \(F(x)\ge 0\) 形式的凸可行性问题，内点法在多项式时间内收敛至 \(\varepsilon\)-精度。\(\blacksquare\)

\textbf{定理 230.5}（Schur补引理）：分块对称矩阵正定当且仅当

\[\begin{bmatrix} Q & S \\ S^T & R \end{bmatrix} > 0 \iff Q > 0,\ R - S^T Q^{-1} S > 0\]

Schur补将非线性Riccati不等式转化为线性矩阵不等式------这是LMI方法能够统一处理各类控制问题的核心代数工具。

\textbf{证明}：对分块矩阵做合同变换：\(\begin{bmatrix}I&0\\-S^T Q^{-1}&I\end{bmatrix}\begin{bmatrix}Q&S\\S^T&R\end{bmatrix}\begin{bmatrix}I&-Q^{-1}S\\0&I\end{bmatrix}=\begin{bmatrix}Q&0\\0&R-S^T Q^{-1}S\end{bmatrix}\)。合同变换保持正定性，故原矩阵正定当且仅当 \(Q>0\) 且 \(R-S^T Q^{-1}S>0\)。等价地取 \(R\) 为主元，\(Q-SR^{-1}S^T>0\) 且 \(R>0\)。若需半正定版本，只须取极限或考虑伪逆。\(\blacksquare\)

\hrulefill

\hrulefill

\hrulefill

\hrulefill

\section{第441章：非线性控制}\label{ux7b2c441ux7ae0ux975eux7ebfux6027ux63a7ux5236}

非线性控制系统的数学理论是人类控制理论中最具挑战性和最丰富的领域之一，其核心困难在于叠加原理不再成立------线性系统中行之有效的频率域方法和传递函数理论在非线性系统中完全失效。因此，非线性控制理论必须依赖全新的数学工具：Lyapunov 稳定性理论（1892）通过构造广义能量函数 \(V(x)\) 并验证其沿轨线的导数 \(\dot{V}(x)=\nabla V^T f(x)\) 的负定性来判定系统的稳定性，而无需显式求解非线性微分方程。LaSalle 不变原理进一步放松了负定性要求，将收敛性分析归约为最大不变集的研究。微分几何方法------尤其是 Lie 括号和 Frobenius 定理------为反馈线性化提供了严格的数学基础：通过坐标变换 \(z=\Phi(x)\) 和非线性反馈 \(u=\alpha(x)+\beta(x)v\)，可以将一类仿射非线性系统精确转化为线性系统（Brunovsky 标准形），其充要条件是系统的可控性 Lie 代数分布在所考虑的点满秩且对合。滑模控制（Sliding Mode Control）则采取了完全不同的策略：通过设计不连续的开关控制律，迫使系统轨迹在有限时间内到达并保持在预设的滑动流形上，且在滑动模上对满足匹配条件的不确定性具有完全的不变性。反步法（Backstepping）递归地构造 Lyapunov 函数，将高阶严格反馈系统的镇定问题分解为一系列一阶子系统的设计，是处理下三角结构非线性系统的系统性方法。

\subsection{231.1 Lyapunov稳定性理论}\label{lyapunovux7a33ux5b9aux6027ux7406ux8bba}

\textbf{定理 231.1}（Lyapunov直接法）：对系统\(\dot{x} = f(x)\)（\(f(0) = 0\)），若存在\(C^1\)函数\(V: \mathbb{R}^n \to \mathbb{R}\)，满足

\begin{itemize}
\tightlist
\item
  \(V(0) = 0\)，且\(V(x) > 0\)（\(x \neq 0\)）（正定性）
\item
  \(\dot{V}(x) = \nabla V(x)^T f(x) < 0\)（\(x \neq 0\)）（负定性）
\end{itemize}

则原点\(x = 0\)是渐近稳定的。\(V(x)\)称为\textbf{Lyapunov函数}。Lyapunov直接法的数学本质是：无需显式求解ODE，仅通过构造广义能量函数\(V\)即可判定稳定性------\(V\)沿轨线严格递减蕴含轨线收敛到原点。

\textbf{证明}：对任意 \(\varepsilon>0\)，取 \(m=\min_{\|x\|=\varepsilon}V(x)>0\)（正定性保证）。由连续性存在 \(\delta>0\) 使 \(B_\delta(0)\subseteq\{x:V(x)<m\}\)。对 \(x_0\in B_\delta(0)\)，由 \(\dot{V}<0\) 得 \(V(x(t))\) 严格递减，故 \(V(x(t))<m\) 对所有 \(t\ge 0\) 成立。轨线不会离开 \(\varepsilon\)-球，稳定得证。渐近稳定性：\(V(x(t))\) 单调递减有下界，极限 \(\lim V(x(t))=L\) 存在。若 \(L>0\)，取序列 \(t_n\to\infty\)，紧性抽子列 \(x(t_n)\to\bar{x}\)，由 \(\dot{V}(\bar{x})<0\) 得 \(V\) 会进一步下降，与 \(L\) 为极限矛盾。故 \(L=0\)，\(\lim x(t)=0\)。\(\blacksquare\)

\textbf{定理 231.2}（LaSalle不变原理）：若\(V(x) \geq 0\)且\(\dot{V}(x) \leq 0\)，则所有有界轨线收敛到最大不变集\(M \subseteq \{x : \dot{V}(x) = 0\}\)。LaSalle原理放松了负定性要求，仅需半负定即可得到收敛结论。

\textbf{证明}：设 \(\gamma^+(x_0)\) 有界，则 \(\omega(x_0)\) 非空且紧致。因 \(V\) 沿轨线递减且有下界，由单调有界定理 \(V(x(t))\to c\)（常数）。对 \(\forall y\in\omega(x_0)\)，取 \(t_n\to\infty\) 使 \(x(t_n)\to y\)，连续性给出 \(V(y)=c\)。结合 \(\dot{V}\le 0\) 和不变性，\(V\) 在 \(\omega(x_0)\) 的轨道上恒为 \(c\)，故 \(\dot{V}=0\) 在 \(\omega(x_0)\) 上处处成立。因此 \(\omega(x_0)\subseteq M\)，从而 \(x(t)\to M\)。该原理优势在于无需渐近稳定平衡点，仅需轨道预紧即可判定收敛目标。\(\blacksquare\)

\textbf{定义 231.1}（控制Lyapunov函数 / CLF）：对受控系统\(\dot{x} = f(x) + g(x) u\)，光滑正定函数\(V(x)\)是\textbf{控制Lyapunov函数}，如果对所有\(x \neq 0\)，

\[\inf_{u} \left\{ L_f V(x) + L_g V(x) \cdot u \right\} < 0\]

其中\(L_f V = \nabla V^T f\)和\(L_g V = \nabla V^T g\)是Lie导数。CLF的存在性是系统可通过反馈镇定的充分条件------每一点都存在控制使\(V\)减小。

\subsection{231.2 反馈线性化}\label{ux53cdux9988ux7ebfux6027ux5316}

\textbf{定义 231.2}（Lie括号与能控性分布）：向量场\(f\)和\(g\)的\textbf{Lie括号}为\([f, g] = \nabla g \cdot f - \nabla f \cdot g\)。迭代Lie括号\(\operatorname{ad}_f^k g = [f, \operatorname{ad}_f^{k-1} g]\)（\(\operatorname{ad}_f^0 g = g\)）张成的分布\(\Delta = \operatorname{span}\{\operatorname{ad}_f^k g : k \geq 0\}\)是该仿射系统的\textbf{可达分布}。

\textbf{定理 231.3}（输入-状态反馈线性化条件）：仿射非线性系统\(\dot{x} = f(x) + g(x) u\)可经状态变换\(z = \Phi(x)\)和反馈\(u = \alpha(x) + \beta(x) v\)化为线性系统\(\dot{z} = A z + B v\)的必要条件是\(\{g, \operatorname{ad}_f g, \ldots, \operatorname{ad}_f^{n-1} g\}\)在\(x_0\)附近线性独立且对合（involutive）。对合条件\(\operatorname{rank}(\operatorname{span}\{g, \ldots, \operatorname{ad}_f^{n-2}g\}) = \operatorname{rank}(\operatorname{span}\{g, \ldots, \operatorname{ad}_f^{n-1}g\})\)由Frobenius定理保证可积性。

\textbf{证明}：必要性：若系统反馈等价于 \((A,B)\) 且其为 Brunovsky 标准形（\(\{B,AB,\dots,A^{n-1}B\}\) 线性独立），则由 Lie 括号与换位子的对应 \([\operatorname{ad}_f^i g,\operatorname{ad}_f^j g] = \Phi_*^{-1}[\operatorname{ad}_{\tilde{f}}^i \tilde{g},\operatorname{ad}_{\tilde{f}}^j \tilde{g}]\)（\(\tilde{f},\tilde{g}\) 为 Brunovsky 形式）为0（因 \(\operatorname{ad}_{\tilde{f}}^n\tilde{g}=0\)），即分布对合。充分性：由线性独立性构造 \(\Phi(x)\) 的输出函数 \(h(x)\) 满足相对阶 \(n\)（\(L_g L_f^k h=0\)，\(k<n-1\)，\(L_g L_f^{n-1}h\neq 0\)），取 \(\Phi(x)=(h,L_f h,\dots,L_f^{n-1}h)^T\)，则 \(z=\Phi(x)\) 实现 Brunovsky 规范化。由 \(f\) 和 \(g\) 在 \(z\) 坐标下的表达即得线性化。\(\blacksquare\)

\textbf{定义 231.3}（相对阶 / Relative Degree）：系统\(\dot{x} = f(x) + g(x) u\)，\(y = h(x)\)在\(x_0\)具有\textbf{相对阶}\(r\)，如果\(L_g L_f^{k} h(x) = 0\)（\(k = 0, \ldots, r-2\)）且\(L_g L_f^{r-1} h(x_0) \neq 0\)。相对阶是输出\(y\)关于时间求导时，控制输入\(u\)首次显式出现的导数阶数。

\subsection{231.3 滑模控制}\label{ux6ed1ux6a21ux63a7ux5236}

\textbf{定义 231.4}（滑模控制 / Sliding Mode Control）：选择\textbf{滑动流形}\(s(x) = 0\)（如\(s = \dot{e} + \lambda e\)）。控制律设计为使系统轨线在有限时间内到达并保持在滑动流形上：

\[u = u_{\text{eq}} - k \operatorname{sgn}(s)\]

其中\(u_{\text{eq}}\)是等效控制（\(\dot{s} = 0\)的解），\(-k \operatorname{sgn}(s)\)是不连续开关项。

\textbf{定理 231.4}（到达条件与有限时间收敛）：若\(s \dot{s} \leq -\eta |s|\)（\(\eta > 0\)），则滑模面\(s = 0\)在有限时间\(T \leq |s(0)|/\eta\)内到达。在滑模面上的运动对满足匹配条件的不确定性具有不变性------这是滑模控制的鲁棒性数学根源。

\textbf{证明}：考虑 Lyapunov 候选函数 \(V(s)=\frac{1}{2}s^2\)。沿轨迹 \(\dot{V}=s\dot{s}\le -\eta|s| = -\sqrt{2}\eta\sqrt{V}\)。此为微分不等式 \(\frac{d}{dt}\sqrt{2V} \le -\eta\)，积分得 \(\sqrt{2V(t)}\le\sqrt{2V(0)}-\eta t\)。令 \(\sqrt{2V(T)}=0\) 得到达时间 \(T\le|s(0)|/\eta\)（在 \(T\) 前若 \(\sqrt{V}=0\) 则已达面）。匹配不确定性：\(f(x)\to f(x)+\Delta f(x)\)（\(\Delta f\in\operatorname{span}\{g\}\)，即匹配条件）。在 \(s=0\) 面上等效运动 \(\dot{x}=(I-g(L_g s)^{-1}L_f s)f\) 中，匹配扰动被 \(g\) 方向分量消除，从而 \(s=0\) 面上的滑动模态对匹配不确定性完全不敏感。\(\blacksquare\)

\subsection{231.4 反步法}\label{ux53cdux6b65ux6cd5}

\textbf{定义 231.5}（反步法 / Backstepping）：对严格反馈系统（strict-feedback form）：

\[\begin{aligned} \dot{x}_1 &= f_1(x_1) + g_1(x_1) x_2 \\ \dot{x}_2 &= f_2(x_1, x_2) + g_2(x_1, x_2) x_3 \\ &\vdots \\ \dot{x}_n &= f_n(x_1, \ldots, x_n) + g_n(x_1, \ldots, x_n) u \end{aligned}\]

反步法递归地构造Lyapunov函数\(V_i(x_1, \ldots, x_i)\)和虚拟控制\(\alpha_i\)，最终得到真实控制\(u\)。

\textbf{定理 231.5}（反步法的Lyapunov构造）：在第\(i\)步，令\(z_i = x_i - \alpha_{i-1}\)为误差变量。选择\(V_i = V_{i-1} + \frac{1}{2} z_i^2\)，设计\(\alpha_i\)使得\(\dot{V}_i = -W_i \leq 0\)。这递归地给出全局渐近稳定的控制律。反步法的本质是递归地应用积分反演（integrator backstepping），将高阶系统的镇定问题降阶为一系列一阶子系统的Lyapunov设计。

\textbf{证明}：归纳法。\(i=1\)：\(z_1=x_1\)，\(\dot{V}_1=z_1(f_1+g_1\alpha_1)\)。选取虚拟控制 \(\alpha_1(x_1)\) 使 \(\dot{V}_1=-W_1\le 0\)（当 \(z_2=0\)）。设构造至第 \(i-1\) 步，\(V_{i-1}\) 满足 \(\dot{V}_{i-1}\le-\sum_{k=1}^{i-1}W_k + z_{i-1}g_{i-1}z_i\)（\(W_k\) 正定）。第 \(i\) 步：\(z_i=x_i-\alpha_{i-1}\)，\(\dot{z}_i=f_i+g_i x_{i+1}-\dot{\alpha}_{i-1}\)。定义 \(V_i=V_{i-1}+\frac{1}{2}z_i^2\)， \[\dot{V}_i \le -\sum_{k=1}^{i-1}W_k + z_i\left[z_{i-1}g_{i-1}+f_i+g_i\alpha_i-\dot{\alpha}_{i-1}\right].\] 括号内取 \(\alpha_i\) 使整体 \(\le -W_i\)（常选 \(\alpha_i=g_i^{-1}(-z_{i-1}g_{i-1}-f_i+\dot{\alpha}_{i-1}-c_i z_i)\)，\(c_i>0\)）。递归至 \(i=n\)，\(u=\alpha_n\) 给出 \(V_n\) 沿闭环递减，由 LaSalle 不变原理得全局渐近稳定。\(\blacksquare\)

\hrulefill

\hrulefill

\hrulefill

\hrulefill

\newpage

\chapter{12-最优化与控制/02-控制与博弈/02-博弈论.md}\label{ux6700ux4f18ux5316ux4e0eux63a7ux523602-ux63a7ux5236ux4e0eux535aux5f0802-ux535aux5f08ux8bba.md}

\section{第442章：非合作博弈}\label{ux7b2c442ux7ae0ux975eux5408ux4f5cux535aux5f08}

非合作博弈论是现代经济学的数学基础，其研究对象是每个参与者独立选择策略以最大化自身支付的策略互动。这一理论的奠基人是 John von Neumann 和 Oskar Morgenstern（《博弈论与经济行为》，1944），而 John Nash 在 1950-1951 年提出的 Nash 均衡概念则是该理论的核心支柱。Nash 的存在性定理证明是 20 世纪数学中最优美的应用之一：将均衡策略剖面刻画为最优反应对应 \(B(\sigma)=\prod_i B_i(\sigma_{-i})\) 在紧凸集（混合策略单纯形的乘积）上的不动点，由 Kakutani 不动点定理保证存在性------这一证明将博弈论与拓扑学和凸分析深刻联系起来。在动态博弈中，Selten 的子博弈完美均衡（1965）排除了不可信威胁，要求均衡在每个子博弈中都是理性的；Kreps-Wilson 的序贯均衡（1982）进一步将这一精炼推广到非完美信息博弈，引入了基于 Bayes 法则的信念系统。Bayes 博弈（Harsanyi，1967-1968）将不完全信息引入博弈论，通过类型空间和共同先验分布刻画参与者对他人支付函数的信念不确定性。Aumann 的相关均衡（1974）则推广了 Nash 均衡的概念，允许参与者根据外部随机信号协调策略选择，其集合是包含所有 Nash 均衡凸包的凸紧集。

\subsection{257.1 策略型博弈与Nash均衡}\label{ux7b56ux7565ux578bux535aux5f08ux4e0enashux5747ux8861}

\textbf{定义 257.1}（策略型博弈 / Normal-Form Game）：一个\(n\)人\textbf{策略型博弈}\(\Gamma = (N, \{S_i\}_{i \in N}, \{u_i\}_{i \in N})\)由以下要素构成： - \(N = \{1, \ldots, n\}\)：参与者集合 - \(S_i\)：参与者\(i\)的\textbf{纯策略}集合 - \(u_i: S = \prod_{j \in N} S_j \to \mathbb{R}\)：参与者\(i\)的\textbf{支付函数}

当参与者的选择是随机的时，\textbf{混合策略}是\(S_i\)上的概率分布\(\sigma_i \in \Delta(S_i)\)。

\textbf{定义 257.2}（Nash均衡，1950）：策略剖面\(\sigma^* = (\sigma_1^*, \ldots, \sigma_n^*)\)是\textbf{Nash均衡}，如果对每个参与者\(i\)和每个策略\(\sigma_i\)，

\[u_i(\sigma_i^*, \sigma_{-i}^*) \geq u_i(\sigma_i, \sigma_{-i}^*)\]

即没有参与者可以通过单方面偏离来提高自己的收益。

\textbf{定理 257.1}（Nash的存在定理，1951）：任何有限策略型博弈至少存在一个Nash均衡（可能为混合策略）。

\textbf{证明}：将均衡刻画为不动点。混合策略空间 \(\Delta=\prod_i\Delta(S_i)\) 是有限维单纯形的笛卡尔积，非空凸紧集。定义最优反应对应 \(B_i(\sigma_{-i})=\arg\max_{\sigma_i\in\Delta(S_i)}u_i(\sigma_i,\sigma_{-i})\)。\(u_i\) 对 \(\sigma_i\) 为线性函数：\(u_i(\sigma_i,\sigma_{-i})=\sum_{s_i\in S_i}\sigma_i(s_i)u_i(s_i,\sigma_{-i})\)，故最优反应集 \(B_i(\sigma_{-i})\) 是紧凸集 \(\Delta(S_i)\) 上线性函数的极大点全体------非空凸紧集。\(B_i\) 的图为闭集：设 \(\sigma_{-i}^n\to\sigma_{-i}\)，\(\sigma_i^n\in B_i(\sigma_{-i}^n)\) 满足 \(\sigma_i^n\to\sigma_i\)，对任意 \(\tau_i\)，\(u_i(\sigma_i^n,\sigma_{-i}^n)\ge u_i(\tau_i,\sigma_{-i}^n)\)，取极限由连续性得 \(u_i(\sigma_i,\sigma_{-i})\ge u_i(\tau_i,\sigma_{-i})\)，故 \(\sigma_i\in B_i(\sigma_{-i})\)。取 \(B(\sigma)=\prod_i B_i(\sigma_{-i})\)，则 \(B:\Delta\to 2^\Delta\) 值域为 \(\Delta\) 的子集，图为闭且对每个 \(\sigma\)，\(B(\sigma)\) 非空凸集。由 Kakutani 不动点定理，\(\exists\sigma^*\in B(\sigma^*)\)，此即 Nash 均衡。\(\blacksquare\)

\subsection{257.2 均衡的精炼}\label{ux5747ux8861ux7684ux7cbeux70bc}

\textbf{定义 257.3}（子博弈完美均衡 / Subgame Perfect Equilibrium，Selten 1965）：扩展型博弈（博弈树）中策略剖面\(\sigma\)是\textbf{子博弈完美均衡}，如果它在每个子博弈中都是Nash均衡。这排除了不可信威胁（non-credible threats）。

子博弈完美均衡（SPE）解决了 Nash 均衡在动态博弈中的核心缺陷：Nash 均衡允许参与者做出''不可信威胁''------即威胁采取某种行动，但若博弈真到达该决策节点，执行该威胁并非最优。Selten 的关键洞察是：在扩展型博弈中，均衡必须对每个子博弈（即每个可能到达的决策节点）都是理性的。逆向归纳法（backward induction）是求解有限完美信息博弈的 SPE 的标准方法：从博弈树的末端出发，逐层向上确定每个决策节点的最优行动，排除那些仅靠不可信威胁维持的 Nash 均衡。SPE 将 Nash 均衡概念从策略型博弈的''事前理性''提炼为''序贯理性''（sequential rationality）------每个参与者在博弈的每个阶段都按最优方式行动。

\textbf{定义 257.4}（颤抖手完美均衡 / Trembling Hand Perfect Equilibrium，Selten 1975）：策略剖面\(\sigma^*\)是\textbf{颤抖手完美均衡}，如果存在一列完全混合策略剖面\(\sigma^k \to \sigma^*\)，使得\(\sigma_i^*\)对每个\(\sigma_{-i}^k\)都是最优反应。

颤抖手完美均衡进一步精炼了 SPE：它不仅要求均衡在策略扰动下稳健，而且要求每个参与者考虑对手可能''不小心''（颤抖）选取任何行动的概率。这一精炼排除了那些依赖''零概率事件''的均衡------在 SPE 中，均衡路径之外的行动概率为零，参与者无须考虑这些行动的最优反应，但颤抖手完美均衡要求均衡策略对任意小的扰动（所有行动均有正概率）都是最优的。这一概念与博弈论中的稳定性概念密切相关：Kohlberg-Mertens（1986）证明了每个有限博弈都存在稳定的 Nash 均衡分量（Strategically Stable Set），该集合中的每个均衡都是颤抖手完美的。

\textbf{定义 257.5}（序贯均衡 / Sequential Equilibrium，Kreps-Wilson 1982）：序贯均衡要求在信息集上的信念是通过 Bayes 法则从均衡策略推导出来的（一致性），且每个参与者的策略是序贯理性的。

序贯均衡是 SPE 在非完美信息博弈中的推广。与 SPE 不同，非完美信息博弈中的信息集可能包含多个节点，参与者无法直接观测当前所处的精确节点，因此必须形成信念（belief）------关于处于信息集中各节点的概率分布。序贯均衡由一对 \((\sigma, \mu)\) 组成：策略剖面 \(\sigma\) 和信念系统 \(\mu\)（每个信息集上各节点的概率分布）。一致性条件要求 \(\mu\) 可以通过 Bayes 法则从 \(\sigma\)``推导''出来（即使对零概率事件，也要存在一列完全混合策略使信念收敛到 \(\mu\)）。序贯理性要求在每个信息集上，给定信念 \(\mu\)，每个参与者的策略是对其他参与者策略的最优反应。这一定义统一了 SPE 和 Bayes-Nash 均衡的精炼思想，是动态博弈中最重要的均衡概念之一。

\subsection{257.3 贝叶斯博弈与不完全信息}\label{ux8d1dux53f6ux65afux535aux5f08ux4e0eux4e0dux5b8cux5168ux4fe1ux606f}

\textbf{定义 257.6}（贝叶斯博弈 / Harsanyi 1967-1968）：\textbf{Bayes博弈}中每个参与者\(i\)有私有的\textbf{类型}\(\theta_i \in \Theta_i\)，类型的联合分布\(p(\theta)\)是共同知识。\(i\)的策略是函数\(s_i: \Theta_i \to S_i\)。

\textbf{定义 257.7}（Bayes-Nash均衡）：策略剖面\(s^*\)是\textbf{Bayes-Nash均衡}，如果对每个\(i\)和\(\theta_i\)，

\[\mathbb{E}_{\theta_{-i}|\theta_i} [u_i(s_i^*(\theta_i), s_{-i}^*(\theta_{-i}), \theta)] \geq \mathbb{E}_{\theta_{-i}|\theta_i}[u_i(s_i, s_{-i}^*(\theta_{-i}), \theta)]\]

\subsection{257.4 相关均衡与通信均衡}\label{ux76f8ux5173ux5747ux8861ux4e0eux901aux4fe1ux5747ux8861}

\textbf{定义 257.8}（相关均衡 / Correlated Equilibrium，Aumann 1974）：概率分布\(\mu \in \Delta(S)\)是\textbf{相关均衡}，如果对所有\(i\)和所有映射\(d_i: S_i \to S_i\)，

\[\sum_{s \in S} \mu(s) u_i(s) \geq \sum_{s \in S} \mu(s) u_i(d_i(s_i), s_{-i})\]

相关均衡允许参与者根据相关信号协调选择------它是Nash均衡概念的推广。

\textbf{定理 257.2}（Aumann相关均衡定理，1974）：相关均衡集是凸紧集，包含所有Nash均衡的凸包。任何相关均衡均可通过引入有限状态的通信设备（correlation device）------即发送给各参与者的私有随机信号------来实现，使得参与者根据接收到的信号理性选择行动构成相关均衡。

\emph{证明}：相关均衡条件 \(\sum_s\mu(s)u_i(s)\ge\sum_s\mu(s)u_i(d_i(s_i),s_{-i})\) 为线性不等式组，定义多面体与概率单纯形的交------凸紧集。取 Nash 均衡 \(\sigma\) 的独立乘积分布 \(\mu(s)=\prod_i\sigma_i(s_i)\)，代入条件仍成立（因 Nash 均衡下 \(u_i(\sigma_i,\sigma_{-i})\ge u_i(d_i,\sigma_{-i})\) 对所有纯策略 \(d_i\) 成立，加权求和后仍成立）。Nash 均衡及其凸组合包含于相关均衡集中。通信设备实现：发送私有信号，使参与者按 \(\mu\) 条件分布选择行动。\(\blacksquare\)

\hrulefill

\hrulefill

\hrulefill

\hrulefill

\section{第443章：合作博弈与核}\label{ux7b2c443ux7ae0ux5408ux4f5cux535aux5f08ux4e0eux6838}

合作博弈论研究参与者可以通过形成联盟并签署有约束力的协议来分配合作收益的情形------与非合作博弈关注个体策略选择不同，合作博弈关注的是联盟的收益分配公平性与稳定性。这一理论由 von Neumann 和 Morgenstern 奠基，并在 1950-1960 年代由 Shapley、Gillies 和 Aumann 等人发展为严格的公理化体系。特征函数 \(v:2^N\to\mathbb{R}\) 刻画了每个联盟 \(S\subseteq N\) 独立行动时能够保证获得的最大收益，博弈的超加性 \(v(S\cup T)\geq v(S)+v(T)\)（\(S\cap T=\emptyset\)）保证了大联盟的形成有利可图。核（Core，Gillies 1953）是合作博弈中最强的稳定性概念：任何子联盟都不能通过单独行动获得比当前分配更高的收益------即对每个联盟 \(S\)，\(\sum_{i\in S}x_i\geq v(S)\)。Bondareva-Shapley 定理（1963）给出了核非空的充要条件（博弈是平衡的），将核的存在性归结为线性规划对偶条件。Shapley 值（1953）则从公平分配的角度出发，将每个参与者的支付定义为其在所有可能加入联盟顺序中边际贡献的期望值，并由效率、对称性、虚拟参与者和可加性四条公理唯一刻画------这一公理化方法论后来被推广到成本分摊、投票权分析和网络流分配等广泛领域。

\subsection{258.1 特征函数型博弈}\label{ux7279ux5f81ux51fdux6570ux578bux535aux5f08}

\textbf{定义 258.1}（特征函数型博弈 / Coalitional Game）：\(n\)人\textbf{特征函数型博弈}\((N, v)\)由特征函数\(v: 2^N \to \mathbb{R}\)定义，\(v(S)\)是联盟\(S\)可以保证获得的最大收益（\(v(\emptyset) = 0\)）。博弈是\textbf{超加性的}（superadditive），如果\(v(S \cup T) \geq v(S) + v(T)\)（\(S \cap T = \emptyset\)）。

\textbf{定义 258.2}（转归 / Imputation）：\textbf{转归}是支付向量\(x = (x_1, \ldots, x_n)\)满足： - 个体理性：\(x_i \geq v(\{i\})\)（每个参与者至少获得单独行动的收益） - 集体理性 / 效率：\(\sum_{i=1}^n x_i = v(N)\)（总收益全部分配）

\subsection{258.2 核与稳定集}\label{ux6838ux4e0eux7a33ux5b9aux96c6}

\textbf{定义 258.3}（核 / Core，Gillies 1953）：博弈\((N, v)\)的\textbf{核}是满足联盟理性的转归集合：

\[C(v) = \left\{ x \in \mathbb{R}^n : \sum_{i \in N} x_i = v(N),\; \sum_{i \in S} x_i \geq v(S) \; (\forall S \subset N) \right\}\]

核中的分配不会被任何联盟通过单独行动改进------它是转归集中最稳定的子集。

\textbf{定理 258.1}（Bondareva-Shapley定理，1963）：博弈的核非空当且仅当博弈是\textbf{平衡的}（balanced）：对所有平衡权\(\{\lambda_S\}_{S \subseteq N}\)（\(\sum_{S \ni i} \lambda_S = 1\)），有\(\sum_S \lambda_S v(S) \leq v(N)\)。Bondareva-Shapley定理将核的非空性等价于一组线性不等式的可解性------这是合作博弈论中与线性规划对偶理论联系最紧密的定理。

\textbf{证明}：核非空等价于线性系统 \(\sum_{i\in N}x_i=v(N)\)，\(\sum_{i\in S}x_i\ge v(S)\)（\(\forall S\subset N\)）有解。由 Farkas 引理（或线性规划对偶），该不等式组无解当且仅当存在非负系数 \(\lambda_S\ge 0\)（不全为0）和 \(\alpha\in\mathbb{R}\) 满足 \(\sum_{S\ni i}\lambda_S=\alpha\)（\(\forall i\)）且 \(\sum_S\lambda_S v(S) > \alpha v(N)\)。若 \(\alpha>0\)，令 \(\lambda_S'=\lambda_S/\alpha\) 得平衡权 \(\sum_{S\ni i}\lambda_S'=1\)，使 \(\sum_S\lambda_S' v(S)>v(N)\)。若 \(\alpha=0\)，\(\sum_{S\ni i}\lambda_S=0\) 且存在 \(S\) 使 \(\lambda_S>0\)，结合超加性可归约为 \(\alpha>0\) 情形。故核非空当且仅当对所有平衡权 \(\{\lambda_S\}\) 有 \(\sum_S\lambda_S v(S)\le v(N)\)。\(\blacksquare\)

\textbf{定义 258.4}（Shapley值，1953）：参与者\(i\)的\textbf{Shapley值}是对其边际贡献的期望：

\[\phi_i(v) = \sum_{S \subseteq N \setminus \{i\}} \frac{|S|! (n - |S| - 1)!}{n!} \left[ v(S \cup \{i\}) - v(S) \right]\]

Shapley值将大联盟的总收益按照各参与者的平均边际贡献进行分配------其组合结构是参与者加入联盟的所有可能顺序上的均匀分布。

\textbf{定理 258.2}（Shapley值的公理化）：Shapley值是满足以下公理的唯一解概念：效率（\(\sum_i \phi_i(v) = v(N)\)）、对称性（对等参与者获得相同支付）、虚拟参与者公理（零边际贡献的参与者获零）和可加性（\(\phi(v + w) = \phi(v) + \phi(w)\)）。这四条公理完全刻画了Shapley值------公理化方法是合作博弈论的核心方法论。

\emph{证明}：先验证 Shapley 值公式满足四条公理：(1) 效率------所有排列的边际贡献和除以 \(n!\) 恰为 \(v(N)\)；(2) 对称性------若 \(i,j\) 对等，在所有排列中的边际贡献期望相同；(3) 虚拟参与者------边际贡献恒为零则值为零；(4) 可加性------由公式线性性。唯一性：将博弈 \(v\) 分解为基博弈 \(v=\sum_{S}\alpha_S v^{\text{un}}\)（股权博弈的线性组合），由公理唯一确定每个基博弈的值，即 Shapley 值。\(\blacksquare\)

\subsection{258.3 核仁与谈判集}\label{ux6838ux4ec1ux4e0eux8c08ux5224ux96c6}

\textbf{定义 258.5}（超额与核仁 / Nucleolus，Schmeidler 1969）：联盟\(S\)对分配\(x\)的\textbf{超额}为\(e(x, S) = v(S) - \sum_{i \in S} x_i\)。\textbf{核仁}是按词典序最小化超额向量（从大到小排列）的唯一分配点。核仁总是存在且唯一，且当核非空时核仁落在核内。

\textbf{定义 258.6}（谈判集 / Bargaining Set，Aumann-Maschler 1964）：谈判集是对异议的稳定性要求。若\(i\)对\(j\)提出异议（带替代联盟），\(j\)能提出反异议，则\(x\)在谈判集中。

\subsection{合作博弈论的现代发展与机制设计}\label{ux5408ux4f5cux535aux5f08ux8bbaux7684ux73b0ux4ee3ux53d1ux5c55ux4e0eux673aux5236ux8bbeux8ba1}

合作博弈论在 20 世纪下半叶经历了从公理化到算法化的深刻转变，Shapley 值（1953）和核心（Gillies 1953）是其两大支柱。\textbf{Shapley 值}的原始动机是公平分配问题------将总收益按照每个参与者对所有可能联盟的边际贡献的期望值进行分配。其公理化体系（效率、对称性、虚拟参与者、可加性）是公理化博弈论（axiomatic game theory）的典范，类似的方法后来被推广到\textbf{成本分摊}（cost sharing）、\textbf{投票权分析}（如 Banzhaf 权力指数）和\textbf{网络流}（如 Myerson 值）中。\textbf{核心}（core）是合作博弈中最强的稳定性概念------没有任何子联盟能通过独立行动获得比当前分配更高的收益。Bondareva-Shapley 定理（1963）给出了核心非空的充要条件：博弈 \((N, v)\) 是\textbf{平衡的}（balanced），即对任意平衡权系数 \(\{\lambda_S\}_{S \subseteq N}\)，有 \(\sum_S \lambda_S v(S) \leq v(N)\)。这一结果将核心的存在性归结为线性规划的对偶条件，与线性规划中的对偶定理和 Farkas 引理完全平行。\textbf{匹配市场}（matching markets）是合作博弈论最成功的应用之一：Gale-Shapley 延迟接受算法（1962）证明了稳定匹配的存在性，且其一边最优性是\textbf{策略防策略}（strategy-proof）机制设计的理论基础。Shapley-Shubik 赋值博弈（1969）将指派问题（assignment problem）转化为合作博弈，证明了核心与线性规划最优解之间的对偶关系。\textbf{机制设计}（mechanism design）是博弈论的''工程学''：Myerson 的拍卖设计理论（1981）和 Vickrey-Clarke-Groves（VCG）机制是激励相容机制的典范，其核心的\textbf{显示原理}（revelation principle）断言任何博弈的均衡结果均可通过直接显示机制在说真话均衡中实现，这一原理将机制设计问题简化为激励相容约束的数学规划。

\hrulefill

\hrulefill

\section{第444章：机制设计与拍卖}\label{ux7b2c444ux7ae0ux673aux5236ux8bbeux8ba1ux4e0eux62cdux5356}

机制设计理论是博弈论的''逆向工程''------与给定博弈规则预测均衡结果的传统博弈论不同，机制设计从期望的社会目标出发，反过来设计博弈规则（机制），使得参与者在追求自身利益时（在均衡中）恰好实现该社会目标。这一思想由 Leonid Hurwicz 在 1960 年代开创，并在他获得 2007 年诺贝尔经济学奖时被公认为经济学最深刻的理论创新之一。机制设计的数学核心是激励相容（Incentive Compatibility）约束的刻画：参与者的真实类型 \(\theta_i\) 是其私有信息，设计者必须通过转移支付 \(t_i(\theta)\) 诱导参与者如实报告类型。Gibbard-Satterthwaite 定理（1973/1975）是机制设计中最著名的''不可能性定理''：在至少三个备选方案且非独裁的条件下，任何社会选择函数都是可操纵的。然而，当引入转移支付（拟线性效用）时，Vickrey-Clarke-Groves（VCG）机制给出了一个璀璨的''可能性''结果：VCG 机制在占优策略下激励相容，其转移支付等于参与者对其他参与者造成的''外部性''------即 \(i\) 的支付 = \(i\) 的参与导致其他参与者社会总福利的减少量。Myerson 的最优拍卖理论（1981）利用虚拟价值函数 \(\psi_i(v_i)=v_i-(1-F_i(v_i))/f_i(v_i)\) 将最优拍卖设计转化为虚拟剩余最大化问题，收入等价定理则表明任何两个分配函数相同的机制产生相同的期望收入。

\subsection{259.1 社会选择与Gibbard-Satterthwaite定理}\label{ux793eux4f1aux9009ux62e9ux4e0egibbard-satterthwaiteux5b9aux7406}

\textbf{定义 259.1}（社会选择函数 / SCF）：\textbf{社会选择函数}\(f: \Theta \to X\)将偏好剖面映射到社会结果。\(\Theta_i\)是参与者\(i\)的类型空间，\(X\)是备选结果集。

\textbf{定义 259.2}（激励相容 / Incentive Compatibility）：直接显示机制\((f, t)\)（\(t_i\)是转移支付）是\textbf{激励相容}的，如果真实报告类型是Bayes-Nash均衡：对所有的\(i\)、\(\theta_i\)、\(\theta_i'\)，

\[\mathbb{E}_{\theta_{-i}}[u_i(f(\theta_i, \theta_{-i}), \theta_i) + t_i(\theta_i, \theta_{-i})] \geq \mathbb{E}_{\theta_{-i}}[u_i(f(\theta_i', \theta_{-i}), \theta_i) + t_i(\theta_i', \theta_{-i})]\]

\textbf{定理 259.1}（Gibbard-Satterthwaite定理，1973/1975）：若\(|X| \geq 3\)且\(f\)的值域包含至少三个结果，则任何非独裁的社会选择函数都是可操纵的（即非激励相容的）。这是社会选择理论中的核心不可能性定理------在相当一般的条件下，策略性操纵是不可避免的。

\textbf{证明}：假设 \(f\) 策略免疫（如实报告是占优策略）且值域 \(|X|\ge 3\)。(1) 单调性：若某备选方案 \(a\) 在所有投票者偏好中排名不降（相对其他备选方案），则社会选择中 \(a\) 的地位不降。形式化：设偏好剖面 \(\succ\) 和 \(\succ'\) 满足对任意 \(i\)，\(\{b:a\succ_i b\}\subseteq\{b:a\succ_i' b\}\)（\(a\) 在 \(\succ'\) 中不比 \(\succ\) 中差），则 \(f(\succ)=a\) 蕴含 \(f(\succ')=a\)。(2) 由单调性加上值域条件（至少三个备选方案），可证存在''轴心''投票者 \(d\) 和两个不同备选方案 \(a,b\) 使得 \(d\) 是 \(\{a,b\}\) 对决的决定者。(3) 逐步推广：\(d\) 对任意备选方案对 \((x,y)\) 有决定权。使用反证法------若 \(d\) 不对 \((x,y)\) 有决定权，则由值域条件构造矛盾偏好剖面，与策略免疫冲突。因此 \(d\) 为独裁者，即 \(f(\succ)=\operatorname{top}(d)\)。\(\blacksquare\)

\subsection{259.2 VCG机制与有效分配}\label{vcgux673aux5236ux4e0eux6709ux6548ux5206ux914d}

\textbf{定义 259.3}（Vickrey-Clarke-Groves机制，1961/1971/1973）：\textbf{VCG机制}选择有效分配\(x^*(\theta) = \arg\max_x \sum_i v_i(x, \theta_i)\)，转移支付为

\[t_i^{\text{VCG}}(\theta) = \sum_{j \neq i} v_j(x^*(\theta), \theta_j) - \max_x \sum_{j \neq i} v_j(x, \theta_j)\]

转移支付\(t_i\)等于参与者\(i\)对其他参与者造成的\textbf{外部性}------其存在导致社会总福利的减少量。

\textbf{定理 259.2}（VCG的激励相容性）：VCG机制是策略性的（dominant-strategy incentive compatible），即真实报告是每个参与者的占优策略。这是机制设计中最重要的正可能性结果。

\emph{证明}：设 \(i\) 的真类型为 \(\theta_i\)，谎报 \(\hat{\theta}_i\)。其效用为 \(v_i(x^*(\hat{\theta}),\theta_i)+t_i^{\text{VCG}}(\hat{\theta})\)。展开 \(t_i^{\text{VCG}}=\sum_{j\neq i}v_j(x^*(\hat{\theta}),\hat{\theta}_j)-\max_x\sum_{j\neq i}v_j(x,\hat{\theta}_j)\)。总效用 \(=\sum_j v_j(x^*(\hat{\theta}),(\theta_i,\hat{\theta}_{-i}))-\max_x\sum_{j\neq i}v_j(x,\hat{\theta}_j)\)。真实报告使第一项最大（因 \(x^*\) 最大化 \(\sum_j v_j\)），第二项与 \(i\) 无关，故真实报告占优。\(\blacksquare\)

\textbf{定理 259.3}（Myerson-Satterthwaite不可能定理，1983）：在双边贸易中（一个买方、一个卖方），不存在同时满足激励相容、个体理性和预算平衡且事后有效的机制。该定理将Gibbard-Satterthwaite不可能性推广到含转移支付的拟线性环境。

\textbf{证明}：卖方成本 \(c\sim G\)，买方估值 \(v\sim F\) 独立。有效分配规则 \(x(v,c)=1\) 若 \(v\ge c\)，否则 \(0\)。由激励相容和包络定理，买方的期望支付满足 \(dU_b(v)/dv = \mathbb{E}_c[x(v,c)]\)，卖方类似。预算平衡要求转移和为零。积分得期望总剩余为 \(\mathbb{E}[v-c]\cdot\Pr(v\ge c)\)。期望预算赤字为交易概率乘以买卖双方信息租之和。计算显示：对所有非退化分布 \(F,G\)，有效交易概率严格为正时，期望赤字严格为正------因此不存在同时满足激励相容、个体理性、预算平衡且事后有效的机制。核心矛盾：有效分配所需的信息租无法由交易双方内部补偿。\(\blacksquare\)

\subsection{259.3 最优拍卖的数学结构}\label{ux6700ux4f18ux62cdux5356ux7684ux6570ux5b66ux7ed3ux6784}

\textbf{定义 259.4}（Myerson最优拍卖，1981）：买方\(i\)的估价\(v_i\)独立来自分布\(F_i\)。\textbf{虚拟价值函数}为\(\psi_i(v_i) = v_i - \frac{1 - F_i(v_i)}{f_i(v_i)}\)。最优拍卖将物品分配给虚拟价值最高的买方，保留价\(r^* = \psi_i^{-1}(0)\)。

\textbf{定理 259.4}（Myerson最优拍卖的推导概要）：卖方的期望收入可写为各买方虚拟剩余的期望和：

\textbf{证明}：买方\(i\)的价值\(v_i\sim F_i\)，密度\(f_i\)。虚拟价值\(\psi_i(v_i)=v_i-\frac{1-F_i(v_i)}{f_i(v_i)}\)。由包络定理，\(U_i(v_i)=\int_{\underline{v}_i}^{v_i}Q_i(t)dt\)（\(Q_i\)为分配概率）。卖方的期望收入\(=\sum_i\mathbb{E}[\psi_i(v_i)Q_i(v)]-\sum_i U_i(\underline{v}_i)\)。在激励相容和个人理性约束下，最优机制将物品分配给虚拟价值最高且为正的买方，支付由临界值确定。这等价于二价拍卖（保留价\(r_i=\psi_i^{-1}(0)\)）。\(\blacksquare\)

\[\mathbb{E}\left[\sum_i \psi_i(v_i) x_i(v)\right]\]

其中\(x_i(v) \in \{0, 1\}\)为分配规则，\(\psi_i(v_i) = v_i - \frac{1-F_i(v_i)}{f_i(v_i)}\)为\textbf{虚拟估值函数}。推导要点：由激励相容条件导出包络定理------买方的期望支付由其分配概率积分确定；将期望支付代入卖方目标函数后，收入仅依赖于分配规则\(x_i(v)\)和虚拟估值。最优分配规则为：将物品分配给\(\max_i \psi_i(v_i)\)的买方（若最高虚拟估值\(<\)卖方保留价值\(0\)，则物品不售出）。\textbf{保留价}\(r_i^* = \psi_i^{-1}(0)\)保证卖方不对虚拟估值为负的买方售出。当买方对称时，最优拍卖化简为含最优保留价的二价拍卖（Vickrey 1961）------这是收入等价定理的直接推论。Myerson的理论统一了最优拍卖、最优垄断定价和非线性定价等机制设计问题。

\textbf{定理 259.5}（收入等价定理 / Revenue Equivalence）：任何两个拍卖机制若分配函数相同且最低类型的期望支付相同，则卖方的期望收入相同。

\textbf{证明}：由包络定理，买方\(i\)的期望支付\(P_i(v_i)=v_iQ_i(v_i)-\int_{\underline{v}_i}^{v_i}Q_i(t)dt-U_i(\underline{v}_i)\)。分配函数\(Q_i(\cdot)\)和最低类型效用\(U_i(\underline{v}_i)\)唯一确定了期望支付。因此，任何两个机制若对每个买方\(i\)有相同的\(Q_i(\cdot)\)和\(U_i(\underline{v}_i)\)，则卖方总期望收入\(\sum_i\mathbb{E}[P_i(v_i)]\)相同。\(\blacksquare\)收入等价定理是激励相容约束和包络定理的直接数学推论。

\hrulefill

\hrulefill

\hrulefill

\hrulefill

\section{第445章：演化博弈论}\label{ux7b2c445ux7ae0ux6f14ux5316ux535aux5f08ux8bba}

演化博弈论将生物学中自然选择的思想引入博弈论，从动力系统的视角解释博弈均衡的形成机制。与传统博弈论假设参与者具有完全理性不同，演化博弈论中参与者通过模仿、学习和自然选择等过程调整策略，均衡被视为动态调整过程的稳定状态而非理性计算的结果。这一理论由 John Maynard Smith 和 George Price 在 1973 年开创，他们提出了演化稳定策略（ESS）的概念：如果大种群中绝大部分个体使用策略 \(x\)，则任何少量变异策略 \(y\) 都无法入侵------即 \(u(x,(1-\varepsilon)x+\varepsilon y)>u(y,(1-\varepsilon)x+\varepsilon y)\)（对充分小的 \(\varepsilon\)）。ESS 是 Nash 均衡的精炼：每个 ESS 都是 Nash 均衡（条件 1：\(u(x,x)\geq u(y,x)\)），但满足更强的稳定性条件（条件 2：若 \(u(x,x)=u(y,x)\)，则 \(u(x,y)>u(y,y)\)）。复制者动态（Replicator Dynamics，Taylor-Jonker 1978）将策略频率的演化描述为一组非线性常微分方程：\(\dot{x}_i=x_i(u(i,x)-\bar{u}(x))\)，即收益高于平均水平的策略比例增长。这一动态系统的渐近稳定平衡点与 ESS 有着密切但非等价的关系。本章在 Nash 均衡理论的基础上，系统阐述演化稳定策略、复制者动态和演化博弈论 Folk 定理。

\subsection{260.1 演化稳定策略}\label{ux6f14ux5316ux7a33ux5b9aux7b56ux7565}

\textbf{定义 260.1}（演化稳定策略 / ESS，Maynard Smith-Price 1973）：在大种群中，策略\(x\)是\textbf{演化稳定策略}，如果对任何变异策略\(y \neq x\)，存在\(\bar{\varepsilon} > 0\)使得对所有\(\varepsilon \in (0, \bar{\varepsilon})\)：

\[u(x, (1-\varepsilon)x + \varepsilon y) > u(y, (1-\varepsilon)x + \varepsilon y)\]

即\(x\)在面对少量变异入侵时仍能维持其优势------入侵后种群的期望支付仍以\(x\)为严格更优。

\textbf{定理 260.1}（ESS的等价条件）：\(x\)是ESS当且仅当： 1. \((x, x)\)是Nash均衡：\(u(x, x) \geq u(y, x)\)（对所有\(y\)） 2. 若\(u(x, x) = u(y, x)\)（\(y\)是对\(x\)的备选最优反应），则\(u(x, y) > u(y, y)\)

条件2是ESS对Nash均衡的加强：当变异策略不严格劣于现存策略时，现存策略必须在变异种群中表现得比变异策略自身更好------这是抵御入侵的稳定性条件。

\emph{证明}：展开 ESS 定义：\((1-\varepsilon)[u(x,x)-u(y,x)]+\varepsilon[u(x,y)-u(y,y)]>0\)。令 \(\varepsilon\to 0\) 得 \(u(x,x)\ge u(y,x)\)（条件1）。若等号成立，代入原式得 \(u(x,y)>u(y,y)\)（条件2）。反推：若条件1,2成立，存在 \(\bar{\varepsilon}\) 充分小使 ESS 定义成立。\(\blacksquare\)

\subsection{260.2 复制者动态}\label{ux590dux5236ux8005ux52a8ux6001}

\textbf{定义 260.2}（复制者方程 / Replicator Dynamics，Taylor-Jonker 1978）：设\(x_i\)是种群中使用纯策略\(i\)的比例。\textbf{复制者方程}为

\[\dot{x}_i = x_i \left( u(i, x) - \bar{u}(x) \right)\]

其中\(u(i, x)\)是纯策略\(i\)对混合策略种群\(x\)的期望收益，\(\bar{u}(x) = \sum_j x_j u(j, x)\)是平均收益。

复制者方程的基本机制是：使用收益高于平均水平的策略的比例增长，低于平均水平的策略的比例减小------其变化率与当前比例和相对收益差成正比。

\textbf{定理 260.2}（复制者动态的性质）： - 单纯形\(\Delta\)是正向不变集------比例始终非负且和为1 - Nash均衡是动态系统的平衡点 - 渐近稳定的平衡点（在单纯形内部）必为ESS - 被严格占优的策略最终被淘汰（\(\lim_{t \to \infty} x_i(t) = 0\)）

\textbf{证明}：(i) \(\sum_i\dot{x}_i=\sum_i x_i(u(i,x)-\bar{u})=\bar{u}-\bar{u}=0\)，\(\dot{x}_i/x_i=u(i,x)-\bar{u}\) 有界，解由 \(\Delta\) 边界止步。(ii) 若 \(x\) 为 Nash，则对所有支集内的 \(i,j\)，\(u(i,x)=u(j,x)=\bar{u}(x)\)，\(\dot{x}_i=0\)；支集外 \(i\) 可能 \(u(i,x)\le\bar{u}(x)\) 且 \(x_i=0\)，平衡成立。(iii) 内部渐近稳定平衡点 \(\hat{x}\) 下，构造 Lyapunov 函数 \(L(x)=-\sum_i\hat{x}_i\ln(x_i/\hat{x}_i)\)，\(\dot{L}=\bar{u}(\hat{x})-\bar{u}(x)<0\)（\(x\neq\hat{x}\)），由条件 (2) 的 ESS 等价性得 \(\hat{x}\) 为 ESS。(iv) 设策略 \(k\) 被严格占优：\(u(k,x)<u(\ell,x)\)（\(\forall x\)）。取 \(L(x)=\ln x_k-\ln x_\ell\)，\(\dot{L}=u(k,x)-u(\ell,x)<0\)，故 \(x_k(t)\to 0\)。\(\blacksquare\)

\textbf{定理 260.3}（演化博弈论的Folk定理）：ESS在单纯形内部是复制者动态的渐近稳定平衡点；边界ESS需要额外条件（如对进入单纯形内部的轨线具有局部吸引性）才能保证渐近稳定性。反过来，渐近稳定平衡点是Nash均衡，但不一定是ESS------渐近稳定性与ESS之间是包含而非等价关系。

\textbf{证明}：内部ESS渐近稳定已由定理260.2(iii)得证（Lyapunov函数 \(L\) 严格递减）。边界ESS：取内部汇点轨线的局部吸引性作为附加假设可保证边界ESS的渐近稳定（通过构造适当边界修正的Lyapunov函数）。反方向：任一渐近稳定平衡点必为Nash均衡（若否，存在受益偏离策略，构造与该策略成比例的初始入浸使轨线偏离该平衡点）。渐近稳定性严格弱于ESS：存在非ESS的渐近稳定Nash均衡（如某些循环博弈中复制子方程的异宿环极限为稳定集但不满足ESS条件2）。\(\blacksquare\)

\subsection{260.3 学习动力学}\label{ux5b66ux4e60ux52a8ux529bux5b66}

\textbf{定义 260.3}（虚拟博弈 / Fictitious Play，Brown 1951）：参与者假设对手使用历史经验分布，在每步选择对过去对手策略分布的最优反应。

\textbf{定理 260.4}（零和博弈中的虚拟博弈）：在有限二人零和博弈中，虚拟博弈的对手策略经验分布收敛到Nash均衡策略。该结论将零和博弈的特殊结构与学习动力学的收敛性建立了数学联系。

\textbf{证明}：设 \(x^t\) 为玩家1的策略经验分布，\(y^t\) 为玩家2的策略经验分布。玩家1在每轮选择对 \(y^t\) 的最优反应，玩家2对 \(x^t\) 的最优反应。定义势函数 \(\Phi(t)=\max_x u(x,y^t)-\min_y u(x^t,y)\)（二人零和博弈的''次优性间隙''）。由 Robinson 收敛定理，\(\Phi(t)\to 0\) 保证经验分布收敛到鞍点。关键步骤：若 \((x^t,y^t)\) 在某轮不收敛，必定产生改善；但势函数 \(\Phi(t)\) 是非增的且收敛到零。具体地，\(\Phi(t+1)-\Phi(t)\le -C/t\cdot\Phi(t)\)（\(C\) 为博弈常数），Gronwall论证给出 \(\Phi(t)=O(1/t)\to 0\)，故经验分布收敛至 Nash 均衡。\(\blacksquare\)

\textbf{定义 260.4}（无悔学习与Regret Matching）：参与者在每轮后计算各策略的累积后悔（regret），以与正后悔成比例的概率选择策略。Hart-Mas-Colell（2000）证明了无悔学习使经验分布收敛到相关均衡集------这是将在线学习算法与博弈均衡概念联系起来的基本定理。

\hrulefill

\hrulefill

\hrulefill

\hrulefill

\section{第446章：随机博弈与重复博弈}\label{ux7b2c446ux7ae0ux968fux673aux535aux5f08ux4e0eux91cdux590dux535aux5f08}

重复博弈和随机博弈研究的是多期互动中策略的长期影响------当博弈重复进行时，合作和协调的可能性远超单期博弈中 Nash 均衡的预测。重复博弈的''Folk 定理''（Friedman 1971，Aumann-Shapley-Rubinstein）是博弈论中最令人瞩目的结果之一：在无限重复博弈中，任何可行且个体理性的支付向量（超过极小极大支付）都可以作为某个子博弈完美均衡的支付来实现，只要贴现因子 \(\delta\) 充分接近 1。这一结论的构造性证明依赖于冷酷触发策略（Grim Trigger）：参与者遵循合作策略，一旦检测到偏离，永久切换至极小极大惩罚。Folk 定理揭示了长期关系中的声誉效应如何克服单期囚徒困境中的囚徒困境------合作的威胁可以自我实施。随机博弈（Stochastic Game，Shapley 1953）将这一框架推广到状态转移随机的环境中：博弈在多个阶段博弈（状态）之间转移，转移概率依赖于当前状态和参与者的行动。Shapley 证明了有限状态、有限行动的二人零和随机博弈存在唯一的值函数，且满足动态规划的 Shapley 方程。Kreps-Wilson-Milgrom-Roberts 的声誉定理（1982）则在有限重复博弈中引出了''声誉效应''------即使只有极小概率的不理性类型存在，理性参与者在博弈早期也会假装合作以建立声誉。

\subsection{261.1 重复博弈与Folk定理}\label{ux91cdux590dux535aux5f08ux4e0efolkux5b9aux7406}

\textbf{定义 261.1}（无限重复博弈）：阶段博弈\(G\)无限重复，贴现因子\(\delta \in (0, 1)\)。参与者的总收益是各阶段收益的贴现和：\(U_i = (1 - \delta) \sum_{t=0}^\infty \delta^t u_i(a^t)\)。

\textbf{定理 261.1}（Nash Folk定理，Friedman 1971）：对无限重复博弈，任何可行且个体理性的支付向量（超过极小极大支付）都可以作为某个\(\delta\)充分接近\(1\)时的子博弈完美均衡支付。

\textbf{证明}：设目标支付 \(v\) 可行（\(v=\sum_k\alpha_k u(a_k)\)，\(\alpha_k\) 为有理数）。构造策略：参与者按周期顺序执行行动序列，实现各纯行动剖面的平均频率 \(\alpha_k\)。偏离监测：若任何参与者偏离指定行动，则所有参与者永久切换至偏离者的极小极大惩罚（minmax惩罚）。当 \(\delta\) 充分接近1时，偏离的一期收益被未来惩罚的贴现损失完全抵消。子博弈完美性：惩罚阶段本身构成 Nash 均衡（极小极大策略互相对抗），故惩罚威胁可信。\(\blacksquare\)

\textbf{定理 261.2}（极限平均支付下的Folk定理 / Aumann-Shapley-Rubinstein）：在无限重复博弈中，若采用**极限平均支付准则*

\textbf{证明}：设\(v\)为阶段博弈的可行且个体理性支付向量。构造策略：参与者轮流执行使\(v\)可达的行动序列，任何偏离触发永久转向minmax惩罚。在极限平均准则下，\(T\)期偏离收益\(\le b/T\)（\(b\)为单期偏离收益上界），\(T\to\infty\)时趋于\(0\)。minmax惩罚使偏离方的长期收益\(\le\)其minmax值\(\le v_i\)（个体理性），故长期偏离无利可图。因此任何可行且个体理性支付向量均可作为子博弈完美均衡实现。\(\blacksquare\)*（\(U_i = \liminf_{T \to \infty} \frac{1}{T} \sum_{t=0}^{T-1} u_i(a^t)\)），则任何\textbf{可行且严格个体理性}的支付向量（超过各参与者的极小极大支付）均可作为Nash均衡实现，且对任意\(\varepsilon > 0\)存在子博弈完美均衡使支付\(\varepsilon\)-逼近该向量。证明构造：将目标支付分解为阶段博弈行动序列的凸组合，参与者相互监控对方是否遵循预定行动轨迹；一旦检测到偏离，则永久切换到极小极大惩罚阶段。极限平均准则（无贴现\(\delta = 1\)）下惩罚成本为零------与贴现情形（\(\delta < 1\)需足够大以确保惩罚可置信）相比，Folk定理的结论更强。此外，Neyman（1985/1998）证明了在有限重复博弈中，若参与者有界理性（使用有限自动机实现策略），合作可作为\(\varepsilon\)-均衡出现。

\textbf{定义 261.2}（触发策略 / Trigger Strategies）：\textbf{冷酷触发策略}（Grim Trigger）：对方偏离则永远合作终止。\(\delta\)充分大时，触发策略构成子博弈完美均衡。

\subsection{261.2 随机博弈}\label{ux968fux673aux535aux5f08}

\textbf{定义 261.3}（随机博弈 / Stochastic Game，Shapley 1953）：\textbf{随机博弈}由状态空间\(\Omega\)、各状态下的阶段博弈和状态转移概率\(p(\omega' | \omega, a)\)组成。贴现因子为\(\delta\)。

\textbf{定理 261.3}（Shapley的二人零和随机博弈）：在有限状态、有限行动的二人零和随机博弈中，存在唯一的值\(v_\delta(\omega)\)，且双方有最优稳定策略。值函数满足Shapley方程：

\[v(\omega) = \operatorname{Val}\left[ (1-\delta) r(\omega, a) + \delta \sum_{\omega'} p(\omega' | \omega, a) v(\omega') \right]\]

\textbf{证明}：定义算子 \(T: \mathbb{R}^{|\Omega|}\to\mathbb{R}^{|\Omega|}\) 为 \((Tv)(\omega)=\operatorname{Val}_{a\in A(\omega)}[(1-\delta)r(\omega,a)+\delta\sum_{\omega'}p(\omega'|\omega,a)v(\omega')]\)。对任意 \(v,w\)，\(|(Tv)(\omega)-(Tw)(\omega)|\le\delta\sum_{\omega'}p(\omega'|\omega,a^*)|v(\omega')-w(\omega')|\le\delta\|v-w\|_\infty\)，故 \(T\) 是 \(\delta\) 压缩（\(\delta<1\)）。由Banach不动点定理，\(T\)有唯一不动点\(v^*\)即值函数。最优稳定策略取各状态下零和矩阵博弈的均衡策略（极小极大策略），由动态规划原理保证全局最优。\(\blacksquare\)

\textbf{定理 261.4}（Mertens-Neyman定理，1981）：在有限二人零和随机博弈中，对于接近\(1\)的贴现因子，值函数\(v(\omega)\)关于\(\delta\)一致有界。

\textbf{证明}：核心是构造策略使得对手偏离的实际收益逼近值函数，同时不受 \(\delta\) 趋近于1的影响。对任意 \(\varepsilon>0\)，构造策略 \(\sigma\) 使对任意对手策略 \(\tau\)，支付界于 \(v_\delta(\omega)\pm\varepsilon\)，且界与 \(\delta\) 无关。利用 Shapley 方程的压缩性质，取有限阶段截断近似，将无穷时域问题化为有限时域逼近。然后取 \(\delta\to 1\) 极限，利用值函数的连续性（关于 \(\delta\)）和紧性论证（有限矩阵博弈的极值定理）得到一致有界性。关键技巧在于''策略修正''机制------当 \(\delta\) 充分接近1时，有限步内支付的累积误差可控，最终在 \(\delta\to 1\) 时一致收敛。\(\blacksquare\)

\subsection{261.3 重复博弈中的声誉效应}\label{ux91cdux590dux535aux5f08ux4e2dux7684ux58f0ux8a89ux6548ux5e94}

\textbf{定理 261.5}（Kreps-Wilson-Milgrom-Roberts声誉定理，1982）：在有限重复囚徒困境中，若存在小概率的不理性类型，理性参与者在博弈早期也会合作（建立声誉），直到最后几期才会背叛。这是声誉效应将有限重复博弈中的逆向归纳逻辑打破的经典结论。

\textbf{证明}：设理性类型（Stackelberg博弈中''针锋相对''类型）以概率 \(p>0\) 存在。在有限 \(T\) 轮博弈中，逆向归纳在纯理性假设下导致全程背叛。引入不理性类型后：若最后几轮前未发生背叛，理性参与者无法确定对手类型------可能是不理性合作型。偏离（背叛）暴露自身为理性，丧失未来合作的声誉租金。当 \(T\) 充分大时，早期背弃的声誉损失超过短期收益。Bayes更新：设 \(q_t\) 为第 \(t\) 轮后对手认为己方不理性的事后概率。合作维持 \(q_t\) 不变，背叛使 \(q_t\to 0\)。策略均衡：前 \(T-K(\delta,p)\) 轮合作（声誉建设期），倒数 \(K\) 轮全背叛。\(K\) 仅依赖 \(p\) 和博弈参数，与 \(T\) 无关------故当 \(T\to\infty\) 时合作占比趋近于1。\(\blacksquare\)

\textbf{定理 261.6}（Fudenberg-Levine定理，1989）：单个长期参与者面对一系列短期参与者时，即使承诺类型概率极小，长期参与者在充分长的博弈中也能获得接近Stackelberg收益的支付。

\textbf{证明}：设长期参与者类型空间 \(\Theta\)，其中 \(\theta_0\) 为承诺类型（commitment type，始终采取固定行动 \(a^*\)，概率 \(p>0\)）。短期参与者为 \(T\) 期顺序到达的无穷序列，每期在观测历史后对长期参与者类型进行 Bayes 更新。承诺类型的存在使短期参与者在观测到长期参与者采取 \(a^*\) 时赋增概率给 \(\theta_0\)。若长期理性参与者模仿 \(a^*\)，短期参与者无法区分是否承诺类型------形成''声誉搭便车''。当 \(T\to\infty\)，因任何偏离行为立即使长期参与者的类型确认为理性（Bayes 后验集中于理性类型），声誉破裂的损失在足够长的剩余博弈中超过任何单期偏离收益。收敛速率为 \(O(1/\sqrt{T})\)------由大数定律和信念集中的速率决定。因此长期参与者的平均支付接近于 \(\theta_0\) 的 Stackelberg 收益。\(\blacksquare\)

\hrulefill

\emph{本卷系统阐述了最优化与博弈两大数学分支：最优化部分涵盖线性规划对偶定理（凸分析）、整数规划的代数结构（全幺模矩阵）和网络流的图论基础（最大流最小割定理作为线性规划对偶的特例）；博弈论部分从Nash均衡的拓扑学证明（Kakutani不动点定理）出发，建立合作博弈的公理化体系（Shapley值与核）、机制设计的数学原理（VCG与激励相容）、演化博弈的动力学理论（复制者方程与ESS）、以及随机博弈与重复博弈的不动点框架。共8章。}

\emph{本卷基于 V2 Ch 17（线性规划初步）和 V7（线性代数），建立了线性规划、对偶理论、凸优化、非线性规划、变分法与最优控制引论、互补性问题与变分不等式、半定规划与离散凸分析、动态规划与随机规划的系统理论。补充内容将控制理论的数学结构（Kalman秩判据与可控性可观测性、Pontryagin极值原理与HJB方程、Kalman-Bucy随机滤波与分离原理、鲁棒控制的H-inf范数与LMI框架、Lyapunov稳定性与非线性控制方法）纳入优化理论的延伸框架；另行补充博弈与机制的数学理论（线性规划对偶、全单模矩阵、最大流最小割、Nash不动点、Shapley公理、VCG机制）。共 13 章。}

\emph{卷二十六：优化理论终。}

\emph{卷二十六：优化理论终。}

\emph{卷二十六：优化理论终。}

\emph{卷二十九：控制论与博弈论终。}

\newpage

\chapter{13-数学物理/数学物理/01-经典力学.md}\label{ux6570ux5b66ux7269ux7406ux6570ux5b66ux7269ux740601-ux7ecfux5178ux529bux5b66.md}

\section{第447章：变分原理与 Lagrange 力学}\label{ux7b2c447ux7ae0ux53d8ux5206ux539fux7406ux4e0e-lagrange-ux529bux5b66}

Lagrange 形式是变分法的经典应用：从泛函驻值条件导出 Euler-Lagrange 方程。本章与 V26 Ch 314 的变分法互补，聚焦于变分法的数学结构。

\subsection{147.1 最小作用量原理}\label{ux6700ux5c0fux4f5cux7528ux91cfux539fux7406}

\textbf{定义 147.1}（作用量）：设 \(q^i(t)\)（\(i=1,\ldots,n\)）是 \(n\) 个广义坐标。\textbf{作用量泛函}为

\[S[q] = \int_{t_1}^{t_2} L(q, \dot{q}, t) \, dt\]

其中 \(L: TQ \times \mathbb{R} \to \mathbb{R}\) 是定义在切丛上的 Lagrange 函数。

\textbf{原理 147.1}（驻值原理）：使作用量 \(S[q]\) 取驻值的曲线 \(q(t)\) 满足 \(\delta S = 0\)，在固定端点条件 \(q(t_1), q(t_2)\) 固定下导出运动方程。

\textbf{定理 147.1}（Euler-Lagrange 方程）：由 \(\delta S = 0\) 导出

\[\frac{d}{dt}\left( \frac{\partial L}{\partial \dot{q}^i} \right) - \frac{\partial L}{\partial q^i} = 0, \quad i = 1, \ldots, n\]

\emph{证明}：考虑变分 \(q^i(t) \to q^i(t) + \varepsilon \eta^i(t)\)，其中 \(\eta^i(t_1) = \eta^i(t_2) = 0\)（固定端点）。作用量的变分为 \[\delta S = \int_{t_1}^{t_2} \sum_i \left( \frac{\partial L}{\partial q^i} \varepsilon \eta^i + \frac{\partial L}{\partial \dot{q}^i} \varepsilon \dot{\eta}^i \right) dt\] 对第二项分部积分：\(\int_{t_1}^{t_2} \frac{\partial L}{\partial \dot{q}^i} \dot{\eta}^i dt = -\int_{t_1}^{t_2} \frac{d}{dt}\!\left(\frac{\partial L}{\partial \dot{q}^i}\right) \eta^i dt\)（边界项因 \(\eta^i(t_1)=\eta^i(t_2)=0\) 消失）。故 \[\delta S = \varepsilon \int_{t_1}^{t_2} \sum_i \left( \frac{\partial L}{\partial q^i} - \frac{d}{dt}\frac{\partial L}{\partial \dot{q}^i} \right) \eta^i dt = 0\] 由变分 \(\eta^i(t)\) 的任意性及变分法基本引理，被积函数恒为零，即得 Euler-Lagrange 方程。∎

\subsection{147.2 Noether 定理}\label{noether-ux5b9aux7406}

\textbf{定理 147.2}（Noether 定理，1918）：Lagrange 系统的每一个连续对称性对应一个守恒量。

具体地，若作用量在变换 \(q^i \to q^i + \varepsilon \xi^i(q, \dot{q})\) 和 \(t \to t + \varepsilon \tau(q, \dot{q})\) 下不变（\(\delta S = 0\)），则

\[Q = \sum_i \frac{\partial L}{\partial \dot{q}^i} \xi^i - \left( \sum_i \frac{\partial L}{\partial \dot{q}^i} \dot{q}^i - L \right) \tau\]

是守恒量（\(\frac{dQ}{dt} = 0\)）。

\textbf{证明}：考虑无穷小变换 \(q^i \to q^i + \varepsilon \xi^i\)，\(t \to t + \varepsilon \tau\)。作用量的变分 \(\delta S = \int_{t_1}^{t_2} [\frac{\partial L}{\partial q^i} \delta q^i + \frac{\partial L}{\partial \dot{q}^i} \delta \dot{q}^i + L \frac{d(\delta t)}{dt}] dt\)。使用 \(\delta q^i = \varepsilon \xi^i - \dot{q}^i \delta t\) 及分部积分。令 \(\delta S = 0\) 对任意积分区间成立，得到 \(\frac{d}{dt}[\frac{\partial L}{\partial \dot{q}^i} \xi^i - (\frac{\partial L}{\partial \dot{q}^i} \dot{q}^i - L)\tau] = (\frac{\partial L}{\partial q^i} - \frac{d}{dt}\frac{\partial L}{\partial \dot{q}^i})(\xi^i - \dot{q}^i \tau)\)。在 Euler-Lagrange 壳层上右端为零，故守恒荷 \(Q = \sum_i \frac{\partial L}{\partial \dot{q}^i}\xi^i - H\tau\)（\(H = \sum_i \frac{\partial L}{\partial \dot{q}^i}\dot{q}^i - L\) 为 Hamiltonian）满足 \(dQ/dt = 0\)。\(\blacksquare\)

\subsection{147.3 约束与 Lagrange 乘子}\label{ux7ea6ux675fux4e0e-lagrange-ux4e58ux5b50}

\textbf{定义 147.2}（约束系统）：若坐标受约束 \(f_\alpha(q, t) = 0\)（\(\alpha = 1, \ldots, m\)），则使用 Lagrange 乘子 \(\lambda_\alpha\)：

\[L' = L + \sum_{\alpha=1}^m \lambda_\alpha(t) f_\alpha(q, t)\]

运动方程为 \(\frac{d}{dt}\frac{\partial L}{\partial \dot{q}^i} - \frac{\partial L}{\partial q^i} = \sum_\alpha \lambda_\alpha \frac{\partial f_\alpha}{\partial q^i}\)。

\textbf{定理 147.3}（d'Alembert 原理的数学表述）：约束力的虚功为零 \(\sum_i \mathbf{F}_i^{\text{constraint}} \cdot \delta \mathbf{r}_i = 0\)，这等价于 Lagrange 乘子法的几何条件。

\emph{证明}：设系统受 \(m\) 个完整约束 \(f_\alpha(q, t) = 0\)（\(\alpha = 1, \ldots, m\)）。约束力 \(\mathbf{F}_i^{\text{constraint}}\) 的方向为约束面的法向，即 \(\mathbf{F}_i^{\text{constraint}} = \sum_\alpha \lambda_\alpha \nabla_i f_\alpha\)，其中 \(\lambda_\alpha\) 为 Lagrange 乘子。虚位移 \(\delta\mathbf{r}_i\) 的定义为满足约束的瞬时位移，满足 \(\sum_i \nabla_i f_\alpha \cdot \delta\mathbf{r}_i = 0\)（即 \(\delta\mathbf{r}\) 切于约束面）。因此虚功 \(\sum_i \mathbf{F}_i^{\text{constraint}} \cdot \delta\mathbf{r}_i = \sum_\alpha \lambda_\alpha \sum_i \nabla_i f_\alpha \cdot \delta\mathbf{r}_i = 0\)。由 Newton 第二定律，\(\sum_i (m_i \ddot{\mathbf{r}}_i - \mathbf{F}_i^{\text{applied}}) \cdot \delta\mathbf{r}_i = \sum_i \mathbf{F}_i^{\text{constraint}} \cdot \delta\mathbf{r}_i = 0\)，即得 d'Alembert 原理的完整形式：\(\sum_i (m_i \ddot{\mathbf{r}}_i - \mathbf{F}_i^{\text{applied}}) \cdot \delta\mathbf{r}_i = 0\)。此原理等价于：约束力对任何虚位移不做功，从而可从运动方程中消去约束力，得到仅含广义坐标的 Euler-Lagrange 方程。\(\blacksquare\)

\hrulefill

\hrulefill

\hrulefill

\hrulefill

\hrulefill

\hrulefill

\section{第448章：Hamilton 力学与辛几何}\label{ux7b2c448ux7ae0hamilton-ux529bux5b66ux4e0eux8f9bux51e0ux4f55}

Hamilton 形式将二阶 ODE 系统转化为一阶 ODE 系统，并揭示其背后的辛几何结构。本章建立 Legendre 变换、Hamilton 正则方程、Poisson 括号代数以及辛流形上的 Hamilton 向量场理论。

\subsection{148.1 Legendre 变换与 Hamilton 方程}\label{legendre-ux53d8ux6362ux4e0e-hamilton-ux65b9ux7a0b}

\textbf{定义 148.1}（广义动量与 Hamilton 函数）：对 Lagrange 函数 \(L(q, \dot{q}, t)\)，定义广义动量

\[p_i = \frac{\partial L}{\partial \dot{q}^i}\]

和 Hamilton 函数（Legendre 变换）

\[H(q, p, t) = \sum_i p_i \dot{q}^i - L(q, \dot{q}, t)\]

其中 \(\dot{q}^i\) 需用 \((q, p)\) 表示。

\textbf{定理 148.1}（Hamilton 正则方程）：

\[\dot{q}^i = \frac{\partial H}{\partial p_i}, \quad \dot{p}_i = -\frac{\partial H}{\partial q^i}\]

\emph{证明}：由 Legendre 变换 \(H = \sum_i p_i \dot{q}^i - L\)，计算全微分： \[dH = \sum_i \left( \dot{q}^i dp_i + p_i d\dot{q}^i - \frac{\partial L}{\partial q^i} dq^i - \frac{\partial L}{\partial \dot{q}^i} d\dot{q}^i \right) - \frac{\partial L}{\partial t} dt\] 由 \(p_i = \frac{\partial L}{\partial \dot{q}^i}\)，\(d\dot{q}^i\) 项抵消。又由 Euler-Lagrange 方程 \(\dot{p}_i = \frac{d}{dt}\frac{\partial L}{\partial \dot{q}^i} = \frac{\partial L}{\partial q^i}\)。比较 \(dH = \sum_i (\frac{\partial H}{\partial q^i} dq^i + \frac{\partial H}{\partial p_i} dp_i) + \frac{\partial H}{\partial t} dt\)，即得 \(\dot{q}^i = \frac{\partial H}{\partial p_i}\)，\(\dot{p}_i = -\frac{\partial H}{\partial q^i}\)。等价性得证。∎

\subsection{148.2 Poisson 括号}\label{poisson-ux62ecux53f7}

\textbf{定义 148.2}（Poisson 括号）：相空间上的函数 \(f(q, p)\) 和 \(g(q, p)\) 的 Poisson 括号为

\[\{f, g\} = \sum_i \left( \frac{\partial f}{\partial q^i} \frac{\partial g}{\partial p_i} - \frac{\partial f}{\partial p_i} \frac{\partial g}{\partial q^i} \right)\]

\textbf{命题 148.2}（Poisson 括号的性质）： - 反对称：\(\{f, g\} = -\{g, f\}\) - 双线性 - Leibniz 法则：\(\{fg, h\} = f\{g, h\} + \{f, h\}g\) - Jacobi 恒等式：\(\{f, \{g, h\}\} + \{g, \{h, f\}\} + \{h, \{f, g\}\} = 0\)

即 Poisson 括号使相空间上的光滑函数构成 Lie 代数。

\textbf{证明}：反对称性由定义直接得：\(\{f, g\} = \sum_i (\frac{\partial f}{\partial q^i} \frac{\partial g}{\partial p_i} - \frac{\partial f}{\partial p_i} \frac{\partial g}{\partial q^i}) = -\sum_i (\frac{\partial g}{\partial q^i} \frac{\partial f}{\partial p_i} - \frac{\partial g}{\partial p_i} \frac{\partial f}{\partial q^i}) = -\{g, f\}\)。双线性性由偏导数的线性性保证。Leibniz 法则：\(\{fg, h\} = \sum_i (\frac{\partial(fg)}{\partial q^i} \frac{\partial h}{\partial p_i} - \frac{\partial(fg)}{\partial p_i} \frac{\partial h}{\partial q^i}) = \sum_i (f\frac{\partial g}{\partial q^i} + g\frac{\partial f}{\partial q^i})\frac{\partial h}{\partial p_i} - \sum_i (f\frac{\partial g}{\partial p_i} + g\frac{\partial f}{\partial p_i})\frac{\partial h}{\partial q^i} = f\{g, h\} + g\{f, h\}\)。Jacobi 恒等式：展开 \(\{f, \{g, h\}\} = \sum_{i,j} \frac{\partial f}{\partial q^i}\frac{\partial}{\partial p_i}(\frac{\partial g}{\partial q^j}\frac{\partial h}{\partial p_j} - \frac{\partial g}{\partial p_j}\frac{\partial h}{\partial q^j}) - \frac{\partial f}{\partial p_i}\frac{\partial}{\partial q^i}(\frac{\partial g}{\partial q^j}\frac{\partial h}{\partial p_j} - \frac{\partial g}{\partial p_j}\frac{\partial h}{\partial q^j})\)。对 \(f, g, h\) 循环求和，所有含二阶导数的项相互抵消（由 Clairaut 定理，混合偏导可交换），仅剩一阶导数项，且三阶项的总和为零。\(\blacksquare\)

\textbf{定理 148.3}（运动方程）：对任意可观测量 \(f(q, p, t)\)，\(\frac{df}{dt} = \{f, H\} + \frac{\partial f}{\partial t}\)。特别地，\(f\) 是守恒量当且仅当 \(\{f, H\} = 0\)（且 \(f\) 不显含时）。

\emph{证明}：由链式法则，\(\frac{df}{dt} = \sum_i (\frac{\partial f}{\partial q^i} \dot{q}^i + \frac{\partial f}{\partial p_i} \dot{p}_i) + \frac{\partial f}{\partial t}\)。代入 Hamilton 方程 \(\dot{q}^i = \frac{\partial H}{\partial p_i}\)，\(\dot{p}_i = -\frac{\partial H}{\partial q^i}\)，得 \(\frac{df}{dt} = \sum_i (\frac{\partial f}{\partial q^i} \frac{\partial H}{\partial p_i} - \frac{\partial f}{\partial p_i} \frac{\partial H}{\partial q^i}) + \frac{\partial f}{\partial t} = \{f, H\} + \frac{\partial f}{\partial t}\)。若 \(\{f, H\}=0\) 且 \(\frac{\partial f}{\partial t}=0\)，则 \(\frac{df}{dt}=0\)。\(\blacksquare\)

\textbf{基本 Poisson 括号}：\(\{q^i, q^j\} = 0\)，\(\{p_i, p_j\} = 0\)，\(\{q^i, p_j\} = \delta^i_j\)。

\subsection{148.3 辛几何}\label{ux8f9bux51e0ux4f55}

\textbf{定义 148.3}（辛形式）：相空间 \(T^*Q\)（余切丛）上的典范辛形式为

\[\omega = \sum_i dq^i \wedge dp_i\]

\(\omega\) 是非退化的闭 2-形式（\(d\omega = 0\)）。

\textbf{定义 148.4}（辛流形）：辛流形 \((M, \omega)\) 是光滑流形配以非退化闭 2-形式 \(\omega\)。

\textbf{定理 148.4}（Darboux 定理）：任何辛流形局部上存在坐标 \((q^i, p_i)\) 使得 \(\omega = \sum_i dq^i \wedge dp_i\)。即所有辛流形局部同构。

\textbf{证明}（Moser 技巧）：设 \(x\) 为辛流形 \((M, \omega)\) 上的点。选取 \(x\) 附近的局部坐标使 \(\omega\) 在 \(x\) 处取标准形式 \(\sum_i dq^i \wedge dp_i\)（这总是可能的，因为任何反对称双线性形式可化为标准形）。设 \(\omega_0 = \sum_i dq^i \wedge dp_i\) 为 \(T_xM\) 上的标准辛形式。令 \(\omega_t = (1-t)\omega_0 + t\omega\)。在 \(x\) 的某邻域内，\(\omega_t\) 对所有 \(t \in [0,1]\) 均为非退化（因为 \(\omega\) 在 \(x\) 处与 \(\omega_0\) 一致）。由 \(d\omega_t = 0\)，存在 1-形式 \(\alpha\) 使得 \(\frac{d}{dt}\omega_t = \omega - \omega_0 = d\alpha\)（Poincare 引理）。构造向量场 \(X_t\) 满足 \(\iota_{X_t}\omega_t = -\alpha\)。由 \(\omega_t\) 的非退化性，\(X_t\) 唯一确定。设 \(\phi_t\) 为 \(X_t\) 的流，则 \(\frac{d}{dt}\phi_t^*\omega_t = \phi_t^*(\mathcal{L}_{X_t}\omega_t + \frac{d}{dt}\omega_t) = \phi_t^*(d(\iota_{X_t}\omega_t) + d\alpha) = \phi_t^*(-d\alpha + d\alpha) = 0\)。故 \(\phi_1^*\omega = \phi_1^*\omega_1 = \phi_0^*\omega_0 = \omega_0\)，即 \(\phi_1\) 将 \(\omega\) 拉回为 \(\omega_0\)，所需局部坐标即 \(\phi_1^{-1}\) 的坐标表示。\(\blacksquare\)

\textbf{定义 148.5}（Hamilton 向量场）：对函数 \(H \in C^\infty(M)\)，其 Hamilton 向量场 \(X_H\) 由 \(\iota_{X_H} \omega = dH\)（即 \(\omega(X_H, \cdot) = dH\)）定义。在 Darboux 坐标下，\(X_H = \sum_i (\frac{\partial H}{\partial p_i} \frac{\partial}{\partial q^i} - \frac{\partial H}{\partial q^i} \frac{\partial}{\partial p_i})\)。

\textbf{定理 148.5}（Liouville 定理）：Hamilton 流保持相空间体积（辛体积 \(\omega^n\)）。即 \(\phi_t^* \omega = \omega\)，故 \(\phi_t^*(\omega^n) = \omega^n\)。

\emph{证明}：由 Cartan 公式，Lie 导数 \(\mathcal{L}_{X_H}\omega = d(\iota_{X_H}\omega) + \iota_{X_H}(d\omega)\)。由 Hamilton 向量场的定义 \(\iota_{X_H}\omega = dH\) 和 \(\omega\) 的闭性 \(d\omega = 0\)，得 \(\mathcal{L}_{X_H}\omega = d(dH) + 0 = 0\)（\(d^2 = 0\)）。因此辛形式沿 Hamilton 流不变，即 \(\frac{d}{dt}|_{t=0} \phi_t^*\omega = \phi_0^*\mathcal{L}_{X_H}\omega = 0\)，故 \(\phi_t^*\omega = \omega\) 对所有 \(t\) 成立。辛体积形式 \(\Omega = \frac{1}{n!}\omega^n = \frac{1}{n!}\omega \wedge \cdots \wedge \omega\)（\(n\) 次），由 \(\phi_t^*\omega = \omega\) 得 \(\phi_t^*\omega^n = \omega^n\)。在 Darboux 坐标 \((q, p)\) 下，\(\omega^n = n! \, dq^1 \wedge dp_1 \wedge \cdots \wedge dq^n \wedge dp_n\)，即相空间体积元。因此 Hamilton 流保相空间体积，此即 Liouville 体积守恒。\(\blacksquare\)

\textbf{定理 148.6}（Poincaré 回复定理）：对保测变换（包括 Hamilton 流），紧致相空间上的几乎每个点都是回复的。

\emph{证明}：设 \((M, \mu)\) 为测度空间，\(\mu(M) < \infty\)，\(\phi_t\) 为保测度流（\(\mu(\phi_t^{-1}(E)) = \mu(E)\) 对所有可测集 \(E\)）。对任意可测集 \(A\)（\(\mu(A) > 0\)），考虑集合 \(B = \{x \in A : \phi_t(x) \notin A \text{ 对所有 } t > 0\}\)（即 \(A\) 中''永不回复''的点）。需证 \(\mu(B) = 0\)。注意到对 \(k \neq j\)，\(\phi_{-k}(B)\) 与 \(\phi_{-j}(B)\) 两两不交：若 \(x \in \phi_{-k}(B) \cap \phi_{-j}(B)\)（\(k > j\)），则 \(\phi_k(x) \in B\) 且 \(\phi_j(x) \in B\)。但 \(\phi_k(x) = \phi_{k-j}(\phi_j(x))\)，而 \(\phi_j(x) \in B \subseteq A\)，故 \(\phi_{k-j}(\phi_j(x)) \in A\)，与 \(B\) 的定义矛盾。由保测性，\(\mu(\phi_{-k}(B)) = \mu(B)\) 对所有 \(k\)。由于 \(\{\phi_{-k}(B)\}_{k=1}^\infty\) 是两两不交的可测集族，\(\sum_{k=1}^\infty \mu(\phi_{-k}(B)) = \sum_{k=1}^\infty \mu(B) \leq \mu(M) < \infty\)，故 \(\mu(B) = 0\)。因此 \(\mu\)-几乎每个 \(x \in A\) 都回复到 \(A\)（即存在 \(t > 0\) 使 \(\phi_t(x) \in A\)）。\(\blacksquare\)

\subsection{148.4 Hamilton-Jacobi 理论与可积系统}\label{hamilton-jacobi-ux7406ux8bbaux4e0eux53efux79efux7cfbux7edf}

\textbf{定义 148.6}（Hamilton-Jacobi 方程）：求函数 \(S(q, t)\) 满足

\[\frac{\partial S}{\partial t} + H\left(q, \frac{\partial S}{\partial q}, t\right) = 0\]

\textbf{定理 148.7}（Liouville-Arnold 可积性定理）：若 \(n\) 自由度的 Hamilton 系统有 \(n\) 个相互对合（\(\{F_i, F_j\} = 0\)）的独立守恒量，则系统可通过作用-角变量 \((I_i, \theta_i)\) 完全积分。运动在 \(n\) 维环面上，\(\dot{\theta}_i = \omega_i(I)\)（常数），\(\dot{I}_i = 0\)。

\textbf{证明}：设 \(F_1, \ldots, F_n\) 为 \(n\) 个相互对合且独立的守恒量，即 \(\{F_i, F_j\} = 0\) 且 \(dF_1 \wedge \cdots \wedge dF_n \neq 0\)。考虑水平集 \(M_f = \{x \in M : F_i(x) = f_i\}\)。由独立性，\(M_f\) 为 \(n\) 维子流形。由对合性，\(F_i\) 生成的 Hamilton 向量场 \(X_{F_i}\) 满足 \(\omega(X_{F_i}, X_{F_j}) = \{F_i, F_j\} = 0\)，故 \(X_{F_i}\) 切于 \(M_f\)。此外，\(\{F_i, H\} = 0\) 蕴含 \(X_{F_i}\) 均与 \(X_H\) 交换。由 Frobenius 定理，\(M_f\) 上的 \(n\) 个交换向量场给出 \(M_f\) 的平行化，故 \(M_f\) 为 \(n\) 维环面 \(T^n\)（若 \(M_f\) 紧致且连通）。定义作用变量 \(I_i = \frac{1}{2\pi} \oint_{\gamma_i} p \cdot dq\)（沿 \(M_f\) 上的第 \(i\) 个基本闭路）。\(I_i\) 仅为 \(f\) 的函数，且 \(dI_i\) 线性无关。角变量 \(\theta_i\) 由辛形式 \(\omega = \sum_i d\theta_i \wedge dI_i\) 确定。在 \((I, \theta)\) 坐标下，\(H = H(I)\)（仅依赖于 \(I\)），故 \(\dot{\theta}_i = \partial H/\partial I_i = \omega_i(I)\)（常数），\(\dot{I}_i = -\partial H/\partial \theta_i = 0\)。因此系统完全可积。\(\blacksquare\)

\hrulefill

\hrulefill

\hrulefill

\hrulefill

\hrulefill

\newpage

\chapter{13-数学物理/数学物理/02-量子力学.md}\label{ux6570ux5b66ux7269ux7406ux6570ux5b66ux7269ux740602-ux91cfux5b50ux529bux5b66.md}

\section{第449章：量子力学数学基础}\label{ux7b2c449ux7ae0ux91cfux5b50ux529bux5b66ux6570ux5b66ux57faux7840}

量子力学的数学基础是 Hilbert 空间上自伴算子的谱理论和 C*-代数表示论。本章建立量子力学的公理化框架，视其为泛函分析的具体实现。

\subsection{149.1 量子力学的公理化}\label{ux91cfux5b50ux529bux5b66ux7684ux516cux7406ux5316}

\textbf{公理 149.1}（量子态）：考虑复可分 Hilbert 空间 \(\mathcal{H}\)。态由 \(\mathcal{H}\) 上的密度算子 \(\rho\)（\(\rho \geq 0, \operatorname{Tr}(\rho) = 1\)）描述。纯态对应于一维投影 \(\rho = |\psi\rangle\langle\psi|\)（\(\|\psi\| = 1\)）。

\textbf{公理 149.2}（可观测量）：可观测量由 \(\mathcal{H}\) 上的自伴算子 \(A = A^\dagger\) 描述。\(A\) 的谱 \(\sigma(A)\) 中的元素对应于测量值。

\textbf{公理 149.3}（测量 / Born 法则）：在态 \(\rho\) 中测量可观测量 \(A\)，结果在 Borel 集 \(E \subseteq \mathbb{R}\) 中的概率为 \(\mathbb{P}(E) = \operatorname{Tr}(\rho P_A(E))\)，其中 \(P_A\) 是 \(A\) 的谱测度（投影值测度，PVM）。

\textbf{公理 149.4}（时间演化 / Schrödinger 方程）：态的时间演化由

\[i\hbar \frac{d}{dt}|\psi(t)\rangle = \hat{H} |\psi(t)\rangle\]

其中 \(\hat{H}\) 是 Hamilton 算子（自伴算子）。形式解为 \(|\psi(t)\rangle = e^{-i\hat{H}t/\hbar} |\psi(0)\rangle\)。

\emph{导出概要}：由量子化条件 \([\hat{x}, \hat{p}] = i\hbar\) 出发。将经典 Hamilton 函数 \(H(p, q)\) 中的变量提升为 \(\mathcal{H}\) 上的自伴算子，得到 Hamilton 算子 \(\hat{H}\)。Schrödinger 方程 \(i\hbar \frac{d}{dt}|\psi(t)\rangle = \hat{H}|\psi(t)\rangle\) 是量子化的直接结果。形式解为 \(|\psi(t)\rangle = e^{-i\hat{H}t/\hbar}|\psi(0)\rangle\)。∎

\textbf{定理 149.1}（Stone 定理）：时间演化 \(U(t) = e^{-i\hat{H}t/\hbar}\) 是 \(\mathcal{H}\) 上的强连续酉算子群。

\textbf{证明}：设 \(\hat{H}\) 是自伴算子。由谱定理，\(\hat{H} = \int_{\sigma(\hat{H})} \lambda \, dP(\lambda)\)。定义 \(U(t) = e^{-i\hat{H}t/\hbar} = \int e^{-i\lambda t/\hbar} dP(\lambda)\)。验证：(1) \(U(t)\) 酉：\(U(t)^\dagger U(t) = \int e^{i\lambda t/\hbar} e^{-i\lambda t/\hbar} dP(\lambda) = I\)。(2) 强连续性：\(\|U(t)\psi - U(s)\psi\|^2 = \int |e^{-i\lambda t/\hbar} - e^{-i\lambda s/\hbar}|^2 d\mu_\psi(\lambda) \to 0\)（\(t\to s\)），由控制收敛定理。(3) 群性质：\(U(t+s) = U(t)U(s)\)。(4) 生成元：\(i\hbar \frac{d}{dt}U(t)|_{t=0}\psi = \hat{H}\psi\)。故 \(\{U(t)\}\) 是强连续酉算子群。\(\blacksquare\)

\textbf{公理 149.5}（复合系统）：若系统 \(A\) 和 \(B\) 的 Hilbert 空间分别为 \(\mathcal{H}_A\) 和 \(\mathcal{H}_B\)，则复合系统的 Hilbert 空间为 \(\mathcal{H}_A \otimes \mathcal{H}_B\)（张量积）。

\subsection{149.2 谱理论在量子力学中的应用}\label{ux8c31ux7406ux8bbaux5728ux91cfux5b50ux529bux5b66ux4e2dux7684ux5e94ux7528}

\textbf{定义 149.1}（谱分解）：由谱定理，自伴算子 \(\hat{A}\) 有谱分解

\[\hat{A} = \int_{\sigma(\hat{A})} \lambda \, dP_A(\lambda)\]

\textbf{定理 149.2}（正则对易关系 / CCR）：\([\hat{x}, \hat{p}] = i\hbar \hat{I}\)。

\textbf{证明}：在位置表象中，\(\hat{x} = x\) 为乘法算子，\(\hat{p} = -i\hbar \frac{d}{dx}\)。对任意波函数 \(\psi(x) \in \mathcal{S}(\mathbb{R})\)，计算 \([\hat{x}, \hat{p}]\psi = \hat{x}\hat{p}\psi - \hat{p}\hat{x}\psi\)。第一项：\(\hat{x}(-i\hbar \frac{d}{dx}\psi) = -i\hbar x \frac{d\psi}{dx}\)。第二项：\(\hat{p}(x\psi) = -i\hbar \frac{d}{dx}(x\psi) = -i\hbar(\psi + x\frac{d\psi}{dx})\)。代入得 \([\hat{x}, \hat{p}]\psi = -i\hbar x \frac{d\psi}{dx} + i\hbar(\psi + x\frac{d\psi}{dx}) = i\hbar\psi\)。故 \([\hat{x}, \hat{p}] = i\hbar \hat{I}\)。\(\blacksquare\)

\textbf{定理 149.3}（Stone-von Neumann 定理）：CCR 的不可约表示在酉等价意义下唯一，即标准 Schrödinger 表示。

\textbf{证明}：设 \(U(a) = e^{ia\hat{x}}\)，\(V(b) = e^{ib\hat{p}}\) 满足 Weyl 关系 \(U(a)V(b) = e^{-i\hbar ab} V(b)U(a)\)。标准 Schrödinger 表示：\(\mathcal{H} = L^2(\mathbb{R})\)，\((U(a)\psi)(x) = e^{iax}\psi(x)\)，\((V(b)\psi)(x) = \psi(x+\hbar b)\)。对任意不可约表示 \((\mathcal{H}', U', V')\)，构造等变同构 \(W: L^2(\mathbb{R}) \to \mathcal{H}'\) 为 \(W\psi = \int \psi(x) U'(x) |\Omega\rangle dx\)（\(|\Omega\rangle\) 为真空态）。由 Weyl 关系的不可约性，该映射是酉的且满足 \(WU(a) = U'(a)W\)，\(WV(b) = V'(b)W\)。故所有不可约强连续 Weyl 关系表示均酉等价。\(\blacksquare\)

\textbf{定理 149.4}（Heisenberg 不确定性原理）：对任意态 \(|\psi\rangle\) 和可观测量 \(A, B\)，

\[\Delta_\psi(A) \cdot \Delta_\psi(B) \geq \frac{1}{2} |\langle \psi | [A, B] | \psi \rangle|\]

特别地，\(\Delta x \cdot \Delta p \geq \frac{\hbar}{2}\)。

\emph{证明}：由 Cauchy-Schwarz 不等式和自伴算子的性质。∎

\subsection{149.3 几个精确可解的自伴算子谱问题}\label{ux51e0ux4e2aux7cbeux786eux53efux89e3ux7684ux81eaux4f34ux7b97ux5b50ux8c31ux95eeux9898}

量子力学中可精确求解的谱问题构成了理解量子系统的基础。以下是几个最重要的例子，它们展示了自伴算子谱理论在物理中的具体应用。

\textbf{（1）一维谐振子}：Hamilton 算子 \(\hat{H} = -\frac{\hbar^2}{2m}\frac{d^2}{dx^2} + \frac{1}{2}m\omega^2 x^2\) 定义在 \(L^2(\mathbb{R})\) 上。通过引入湮灭算子 \(a = \sqrt{\frac{m\omega}{2\hbar}}(\hat{x} + \frac{i}{m\omega}\hat{p})\) 和产生算子 \(a^\dagger\)，满足 \([a, a^\dagger] = I\)，Hamilton 算子重写为 \(\hat{H} = \hbar\omega(a^\dagger a + \frac{1}{2})\)。粒子数算子 \(N = a^\dagger a\) 的谱为 \(\sigma(N) = \{0, 1, 2, \ldots\}\)（非负整数），故 \(\hat{H}\) 的谱为 \(\sigma(\hat{H}) = \{\hbar\omega(n + \frac{1}{2}) : n \in \mathbb{N}_0\}\)，全为离散本征值。本征函数由 Hermite 多项式给出：\(\psi_n(x) = \frac{1}{\sqrt{2^n n!}} \left(\frac{m\omega}{\pi\hbar}\right)^{1/4} H_n(\sqrt{\frac{m\omega}{\hbar}}x) e^{-m\omega x^2/(2\hbar)}\)。谐振子的重要性在于：任何光滑势能在其极小值附近的二阶近似都是一个谐振子，因此它是量子场论中自由场的数学基础。

\textbf{（2）氢原子与 Coulomb 势}：\(\hat{H} = -\frac{\hbar^2}{2m}\nabla^2 - \frac{e^2}{r}\) 定义在 \(L^2(\mathbb{R}^3)\) 上。在球坐标下分离变量，径向部分满足广义 Laguerre 方程。谱由两部分组成：离散谱 \(E_n = -\frac{me^4}{2\hbar^2}\frac{1}{n^2}\)（\(n = 1, 2, \ldots\)，束缚态），以及连续谱 \([0, \infty)\)（散射态）。\(SO(4)\) 对称性（Runge-Lenz 向量守恒）解释了 \(n^2\) 重简并。Coulomb 势的精确可解性源于其与四维球面 \(S^3\) 上 Laplace 算子的隐藏联系（Fock 变换）。

\textbf{（3）无限深方势阱}：\(\hat{H} = -\frac{\hbar^2}{2m}\frac{d^2}{dx^2}\) 定义在 \([0, L]\) 上，边界条件为 \(\psi(0) = \psi(L) = 0\)。这是 Dirichlet 边界条件下的一维 Laplace 算子，谱为纯离散：\(E_n = \frac{\hbar^2\pi^2 n^2}{2mL^2}\)（\(n = 1, 2, \ldots\)），本征函数 \(\psi_n(x) = \sqrt{\frac{2}{L}}\sin(\frac{n\pi x}{L})\)。该问题展示了自伴延拓理论：对称算子 \(-\frac{d^2}{dx^2}\) 在 \(C_0^\infty(0, L)\) 上有无穷多个自伴延拓，对应不同的边界条件，而每种边界条件给出不同的谱。

\subsection{149.4 量子力学的 C*-代数表述}\label{ux91cfux5b50ux529bux5b66ux7684-c-ux4ee3ux6570ux8868ux8ff0}

\textbf{定义 149.2}（量子力学的代数方法）：量子系统的可观测量构成 C*-代数 \(\mathfrak{A}\)（或 von Neumann 代数）。态是 \(\mathfrak{A}\) 上的正线性泛函 \(\omega\)（\(\omega(1) = 1\)）。GNS 构造将代数态转化为 Hilbert 空间表示。

C\emph{-代数表述是量子力学的第三种等价数学框架（前两种为 Schrödinger 的波动力学和 Heisenberg 的矩阵力学）。其核心思想是：量子力学的本质结构不在于具体的 Hilbert 空间，而在于可观测量的代数结构。C}-代数 \(\mathfrak{A}\) 是一个配备对合运算 \(*\) 和范数 \(\|\cdot\|\) 的 Banach 代数，满足 \(\|A^*A\| = \|A\|^2\)（C\emph{-恒等式）。所有有界线性算子构成的代数 \(\mathcal{B}(\mathcal{H})\) 是 C}-代数的原型，但代数方法允许我们考虑更一般的代数（如非交换环面上的旋转代数 \(A_\theta\)），这些代数不一定作为某个给定 Hilbert 空间上的算子代数出现。态 \(\omega\) 的物理诠释是：\(\omega(A)\) 是测量可观测量 \(A\) 的期望值。纯态对应于不可约 GNS 表示，混合态对应于可约表示中迹类算子的积分。代数方法在量子统计力学中尤为强大：KMS 条件（Kubo-Martin-Schwinger）和 Tomita-Takesaki 模理论在代数框架下自然展开，提供了温度和平衡态的不依赖于 Hilbert 空间表示的严格定义。此外，代数方法在量子场论的代数化（Haag-Kastler 公理系）和局域量子物理中发挥着核心作用，使得超选择法则（superselection rules）和 Doplicher-Haag-Roberts 理论可以在严格代数框架下表述。

\textbf{定理 149.5}（Gelfand-Naimark-Segal 构造）：对 C*-代数 \(\mathfrak{A}\) 上的任意态 \(\omega\)，存在 Hilbert 空间 \(\mathcal{H}_\omega\)、表示 \(\pi_\omega: \mathfrak{A} \to \mathcal{B}(\mathcal{H}_\omega)\) 和循环向量 \(\Omega_\omega\)，使得 \(\omega(A) = \langle \Omega_\omega | \pi_\omega(A) | \Omega_\omega \rangle\)。

\textbf{证明}：在 \(\mathfrak{A}\) 上定义半内积 \(\langle A,B\rangle_\omega = \omega(B^*A)\)。令 \(N_\omega = \{A\in\mathfrak{A}: \omega(A^*A)=0\}\)（由 Cauchy-Schwarz 不等式知其为左理想）。商空间 \(\mathfrak{A}/N_\omega\) 在内积 \(\langle A+N_\omega, B+N_\omega\rangle = \omega(B^*A)\) 下为内积空间。完备化得 Hilbert 空间 \(\mathcal{H}_\omega\)。定义表示 \(\pi_\omega(A)(B+N_\omega) = AB + N_\omega\)（左乘），可延拓至 \(\mathcal{H}_\omega\) 上的有界算子。循环向量 \(\Omega_\omega = 1 + N_\omega\) 满足 \(\omega(A) = \langle\Omega_\omega|\pi_\omega(A)|\Omega_\omega\rangle\)。GNS 三元组在同构意义下唯一。\(\blacksquare\)

GNS 构造是沟通代数量子力学与 Hilbert 空间量子力学的桥梁。它表明：任何 C*-代数上的态 \(\omega\) 都可以通过标准构造产生一个 Hilbert 空间表示，在该表示中 \(\omega\) 由一个循环向量 \(\Omega_\omega\) 实现为向量态。循环性的含义是 \(\{\pi_\omega(A)\Omega_\omega : A \in \mathfrak{A}\}\) 在 \(\mathcal{H}_\omega\) 中稠密------整个 Hilbert 空间可以通过作用代数元素于 \(\Omega_\omega\) 生成。这一构造在 von Neumann 代数框架下的推广即为 Tomita-Takesaki 模理论，其中模算子 \(\Delta\) 和模共轭 \(J\) 满足 \(J\mathfrak{M}J = \mathfrak{M}'\)（换位子定理），且 \(\sigma_t^\omega(A) = \Delta^{it} A \Delta^{-it}\) 定义了模自同构群，物理上对应于时间演化。

\subsection{149.5 Maxwell 方程组与外微分}\label{maxwell-ux65b9ux7a0bux7ec4ux4e0eux5916ux5faeux5206}

\textbf{定理 149.6}（Maxwell 方程组的微分形式表述）：

\[\nabla \cdot \mathbf{E} = \rho, \quad \nabla \cdot \mathbf{B} = 0\]

\[\nabla \times \mathbf{E} = -\frac{\partial \mathbf{B}}{\partial t}, \quad \nabla \times \mathbf{B} = \mathbf{J} + \frac{\partial \mathbf{E}}{\partial t}\]

引入四维势 \(A^\mu = (\phi, \mathbf{A})\) 和场强张量 \(F_{\mu\nu} = \partial_\mu A_\nu - \partial_\nu A_\mu\)，Maxwell 方程组可紧凑地写为 \(\partial_\mu F^{\mu\nu} = J^\nu\) 和 \(\partial_{[\lambda}F_{\mu\nu]} = 0\)。这是 \(U(1)\) 规范场的经典原型，也是外微分 \(dF = 0\) 和 \(d * F = *J\) 的具体实现。

\textbf{证明}：定义 2-形式 \(F = \frac{1}{2}F_{\mu\nu}dx^\mu \wedge dx^\nu\)。由 \(F_{\mu\nu} = \partial_\mu A_\nu - \partial_\nu A_\mu\) 得 \(F = dA\)（\(A = A_\mu dx^\mu\) 为 1-形式）。由 \(d^2 = 0\) 立得 \(dF = d^2 A = 0\)，即 Bianchi 恒等式 \(\partial_{[\lambda}F_{\mu\nu]} = 0\)，对应 \(\nabla \cdot \mathbf{B} = 0\) 和 \(\nabla \times \mathbf{E} = -\partial_t \mathbf{B}\)。引入 Hodge 对偶 \(*F = \frac{1}{4}F^{\alpha\beta}\varepsilon_{\alpha\beta\mu\nu}dx^\mu \wedge dx^\nu\)，有源方程为 \(d * F = *J\)（\(J = J_\mu dx^\mu\)），即 \(\partial_\mu F^{\mu\nu} = J^\nu\)，对应 \(\nabla \cdot \mathbf{E} = \rho\) 和 \(\nabla \times \mathbf{B} = \mathbf{J} + \partial_t \mathbf{E}\)。规范变换 \(A \mapsto A + d\chi\) 保持 \(F\) 不变，体现了 \(U(1)\) 规范对称性。\(\blacksquare\)

\section{第450章：量子力学的数学公理化}\label{ux7b2c450ux7ae0ux91cfux5b50ux529bux5b66ux7684ux6570ux5b66ux516cux7406ux5316}

量子力学的数学公理化是20世纪数学物理最深刻的成就之一。von Neumann在其1932年的经典著作《量子力学的数学基础》中首次将量子力学建立于Hilbert空间和自伴算子的谱理论之上，这一框架至今仍是量子理论的标准数学表述。本章的核心是确立量子力学的基本数学结构------态空间、可观测量代数和测量理论------并证明所有这些结构在数学上是相互协调且唯一确定的。Stone-von Neumann唯一性定理表明，正则对易关系的不可约强连续酉表示在酉等价意义下是唯一的，这从根本上保证了Schrödinger波动力学与Heisenberg矩阵力学的等价性。Gleason定理则进一步从根本上推导出Born概率法则：在维数至少为3的Hilbert空间中，任何满足可加性的概率测度必由密度算子经迹公式给出。这些定理共同构成了量子力学逻辑一致性的数学基石，也使量子力学成为所有物理理论中数学基础最为坚实的框架之一。

\subsection{150.1 von Neumann公理体系}\label{von-neumannux516cux7406ux4f53ux7cfb}

\textbf{定义 150.1}（von Neumann量子力学公理系）：设\(\mathcal{H}\)为复可分Hilbert空间。von Neumann公理体系由以下四条构成：

\textbf{公理I（态公理）}：量子系统的态由\(\mathcal{H}\)上的密度算子（迹类非负自伴算子）\(\rho\)描述，满足\(\rho \geq 0\)且\(\operatorname{Tr}(\rho)=1\)。纯态对应于一维投影\(\rho = |\psi\rangle\langle\psi|\)（\(\|\psi\|=1\)）。

\textbf{公理II（可观测公理）}：每个可观测量对应于\(\mathcal{H}\)上的（可能无界）自伴算子\(A: \mathcal{D}(A) \to \mathcal{H}\)。\(A\)的谱测度\(P_A\)给出测量的概率分布。

\textbf{公理III（测量公理）}：在态\(\rho\)中测量可观测量\(A\)，结果落在Borel集\(E \subseteq \mathbb{R}\)中的概率为\(\mathbb{P}_\rho^A(E) = \operatorname{Tr}(\rho P_A(E))\)。

\textbf{公理IV（演化公理）}：孤立量子系统的态按Schrödinger方程演化：\(i\hbar \frac{d}{dt}\rho(t) = [\hat{H}, \rho(t)]\)，其中\(\hat{H}\)为Hamilton（自伴）算子。等价地，\(\rho(t) = e^{-i\hat{H}t/\hbar}\rho(0)e^{i\hat{H}t/\hbar}\)。

\textbf{定理 150.1}（自伴算子的谱定理与可观测量的投影值测度）：设\(A\)为\(\mathcal{H}\)上的自伴算子。则存在唯一的投影值测度（PVM）\(P_A: \mathcal{B}(\mathbb{R}) \to \mathcal{P}(\mathcal{H})\)（\(\mathcal{P}(\mathcal{H})\)为\(\mathcal{H}\)上正交投影的集合），使得 \[A = \int_{\mathbb{R}} \lambda \, dP_A(\lambda).\] PVM \(P_A\)满足：(1) \(P_A(\mathbb{R}) = I\)；(2) \(P_A(E \cap F) = P_A(E)P_A(F)\)；(3) 对互不相交的\(\{E_n\}\)，\(P_A(\bigcup E_n) = \sum P_A(E_n)\)（强收敛）。

\textbf{证明}：此为自伴算子谱定理的标准形式。设\(A\)为自伴算子，定义其预解式\(R_A(z) = (A - zI)^{-1}\)（\(z \notin \sigma(A)\)）。\(R_A(z)\)是\(\mathbb{C} \setminus \sigma(A)\)上的解析函数。对任意\(\psi \in \mathcal{H}\)，Stieltjes反演公式给出谱测度： \[\langle\psi|P_A((a,b])\psi\rangle = \lim_{\delta \to 0^+}\lim_{\varepsilon \to 0^+}\frac{1}{\pi}\int_{a+\delta}^{b+\delta}\operatorname{Im}\langle\psi|R_A(\lambda+i\varepsilon)\psi\rangle\, d\lambda.\] 由此可构造Borel集上的投影值测度\(P_A\)，且\(\int \lambda \, dP_A(\lambda) = A\)在算子强收敛意义下成立。唯一性由谱测度由\(A\)的Cayley变换或Weyl函数唯一确定而得出。\(\blacksquare\)

\textbf{定理 150.2}（Stone-von Neumann唯一性定理）：设\((U(a))_{a \in \mathbb{R}}\)和\((V(b))_{b \in \mathbb{R}}\)为Hilbert空间\(\mathcal{H}\)上的两个强连续酉算子单参数群，满足Weyl交换关系 \[U(a)V(b) = e^{-i\hbar ab}V(b)U(a),\quad \forall a,b \in \mathbb{R},\] 且该表示是不可约的。则\((\mathcal{H}, U, V)\)酉等价于标准Schrödinger表示：\(\mathcal{H} \cong L^2(\mathbb{R})\)，\((U(a)\psi)(x) = e^{iax}\psi(x)\)，\((V(b)\psi)(x) = \psi(x + \hbar b)\)。

\textbf{证明}：第一步：构造Weyl算子的积分形式。定义算子 \[\tilde{P} = \frac{1}{2\pi\hbar}\iint_{\mathbb{R}^2} e^{-iab/2}U(a)V(b)\, da\, db\] （适当的弱积分意义下）。计算\(\tilde{P}^2\)：利用Weyl关系， \[U(a)V(b)U(a')V(b') = e^{-i\hbar a'b}U(a+a')V(b+b').\] 通过变量代换可验证\(\tilde{P}^2 = \tilde{P}\)，且\(\tilde{P}\)为一维投影。

第二步：证明\(\tilde{P} \neq 0\)。假设\(\tilde{P} = 0\)，则对任意\(\psi, \phi\)， \[\iint e^{-iab/2}\langle\phi|U(a)V(b)\psi\rangle\, da\, db = 0.\] 由Fourier逆变换，\(\langle\phi|U(a)V(b)\psi\rangle = 0\)对所有\(a,b\)成立，与不可约性矛盾。

第三步：设\(\Omega\)为\(\tilde{P}\)的像空间中范数为1的向量。由构造可得\(U(a)V(b)\Omega = e^{iab/2}V(b)U(a)\Omega\)，结合\(\tilde{P}\)的定义可验证\(\langle\Omega|U(a)V(b)\Omega\rangle = e^{-(a^2+b^2)/4}\)（Gauss型函数），即\(\Omega\)为相干态的真空态。

第四步：定义\(W: L^2(\mathbb{R}) \to \mathcal{H}\)为\(W\psi = \int \psi(x)U(x)\Omega\, dx\)。由Weyl关系可验证\(W\)为酉算子，且\(WU_{\text{Sch}}(a) = U(a)W\)，\(WV_{\text{Sch}}(b) = V(b)W\)。故所有不可约强连续Weyl表示酉等价。\(\blacksquare\)

\textbf{定理 150.3}（Weil-Wigner-Moyal量子化 / 形变量子化）：设\(f, g \in \mathcal{S}(\mathbb{R}^{2n})\)为相空间上的Schwartz函数。Weyl量子化映射\(\operatorname{Op}_{\hbar}: \mathcal{S}(\mathbb{R}^{2n}) \to \mathcal{B}(L^2(\mathbb{R}^n))\)定义为 \[\operatorname{Op}_{\hbar}(f)\psi(x) = \frac{1}{(2\pi\hbar)^n}\iint_{\mathbb{R}^{2n}} e^{i(x-y)\cdot \xi/\hbar} f\!\left(\frac{x+y}{2}, \xi\right)\psi(y)\, dy\, d\xi.\] 则\(\operatorname{Op}_{\hbar}(f)\operatorname{Op}_{\hbar}(g) = \operatorname{Op}_{\hbar}(f \star_{\hbar} g)\)，其中Moyal星积定义为 \[(f \star_{\hbar} g)(x,\xi) = f(x,\xi)\exp\!\left(\frac{i\hbar}{2}(\overleftarrow{\partial_x}\cdot\overrightarrow{\partial_\xi} - \overleftarrow{\partial_\xi}\cdot\overrightarrow{\partial_x})\right)g(x,\xi).\]

\textbf{证明}：由Weyl量子化的定义，\(\operatorname{Op}_{\hbar}(f)\)的积分核为 \[K_f(x,y) = \frac{1}{(2\pi\hbar)^n}\int f\!\left(\frac{x+y}{2}, \xi\right)e^{i(x-y)\cdot\xi/\hbar}\, d\xi.\] 两个算子\(\operatorname{Op}_{\hbar}(f)\)和\(\operatorname{Op}_{\hbar}(g)\)的复合积分核为\(K_{f \star_{\hbar} g}(x,z) = \int K_f(x,y)K_g(y,z)\, dy\)。代入\(K_f\)和\(K_g\)的表达式： \[K_{f \star_{\hbar} g}(x,z) = \frac{1}{(2\pi\hbar)^{2n}}\iiint f\!\left(\frac{x+y}{2}, \xi\right)g\!\left(\frac{y+z}{2}, \eta\right)e^{i[(x-y)\cdot\xi + (y-z)\cdot\eta]/\hbar}\, dy\, d\xi\, d\eta.\] 作变量代换\(u = \frac{x+z}{2}\)（中点），\(v = y - u\)（偏离），则指数部分变为\(\frac{i}{\hbar}[(\frac{x-z}{2}-v)\cdot\xi + (v+\frac{x-z}{2})\cdot\eta]\)。对\(v\)积分产生\(\delta\)函数，等价于将\(\xi, \eta\)替换为适当的微分算子。利用渐近展开（平稳相位法），得到Moyal星积的渐近级数表达式。该级数在\(\hbar \to 0\)时逐阶对应于经典Poisson括号的形变量子化（Kontsevich形式定理）。\(\blacksquare\)

\subsection{150.2 Gleason定理：Born法则的推导}\label{gleasonux5b9aux7406bornux6cd5ux5219ux7684ux63a8ux5bfc}

\textbf{定理 150.4}（Gleason定理，1957）：设\(\mathcal{H}\)为维数\(\dim\mathcal{H} \geq 3\)的可分Hilbert空间。设\(\mu: \mathcal{P}(\mathcal{H}) \to [0,1]\)为投影格上的测度，即\(\mu(I)=1\)且对两两正交的投影族\(\{P_i\}\)有\(\mu(\sum_i P_i) = \sum_i \mu(P_i)\)。则存在唯一的迹类非负算子\(\rho\)（密度算子，\(\operatorname{Tr}(\rho)=1\)），使得对所有投影\(P \in \mathcal{P}(\mathcal{H})\)有 \[\mu(P) = \operatorname{Tr}(\rho P).\]

\textbf{证明}：第一步：将\(\mu\)从一维投影延拓到任意有界自伴算子。对任意单位向量\(\psi \in \mathcal{H}\)，定义\(f(\psi) = \mu(P_\psi)\)，其中\(P_\psi\)为到\(\psi\)生成的一维子空间的投影。

第二步：证明\(f\)在\(\mathcal{H}\)的单位球面上连续。核心引理：若\(\|\psi - \phi\|\)很小，则\(|f(\psi) - f(\phi)|\)也很小。这由框架（frame）理论得出：任意两个单位向量\(\psi, \phi\)可嵌入一个三维子空间中，在该子空间上构造一组框架\(\{e_1, e_2, e_3\}\)使得\(\psi\)和\(\phi\)与框架向量的内积关系限制\(f\)的变化。具体地，由Bessel不等式和框架算子的性质， \[|f(\psi) - f(\phi)| \leq C\|\psi - \phi\|^{1/2}\] 对某个与\(\psi, \phi\)无关的常数\(C\)成立（此步骤依赖\(\dim\mathcal{H} \geq 3\)这一关键条件------在二维情形下Gleason定理不成立！）。

第三步：利用极化恒等式将\(f\)扩张为\(\mathcal{H}\)上的有界二次型。定义 \[B(\psi, \phi) = \frac{1}{4}\left[f(\psi+\phi) - f(\psi-\phi) + i f(\psi+i\phi) - i f(\psi-i\phi)\right]\] （先在单位向量上定义，再通过尺度延拓）。可验证\(B\)为有界Hermite二次型，故由Riesz表示定理存在有界自伴算子\(\rho\)使得\(B(\psi, \phi) = \langle\psi|\rho\phi\rangle\)。

第四步：验证\(\rho\)为密度算子。对任意单位向量\(\psi\)，\(f(\psi) = \langle\psi|\rho\psi\rangle = \mu(P_\psi) \geq 0\)，故\(\rho \geq 0\)。取\(\mathcal{H}\)的任意规范正交基\(\{e_i\}\)，\(\operatorname{Tr}(\rho) = \sum_i \langle e_i|\rho e_i\rangle = \sum_i f(e_i) = \mu(\sum_i P_{e_i}) = \mu(I) = 1\)，故\(\rho\)为密度算子。唯一性由\(\rho\)具体由二次型\(B\)决定而得出。\(\blacksquare\)

Gleason定理的深远意义在于：它将Born法则从一条独立的物理公理推导为Hilbert空间几何结构的必然推论。在\(\dim\mathcal{H} \geq 3\)的量子系统中，任何满足可加性的概率测度必由密度算子经迹公式给出------不需要额外假设。这一定理消除了量子力学公理体系中的逻辑冗余，使得von Neumann公理体系的自洽性更加牢固。值得注意的是，Gleason定理在二维情形下不成立：\(\mathbb{C}^2\)中存在满足可加性但不来自密度算子的概率测度，这与量子比特系统的特殊几何性质有关。Gleason定理的现代推广（Busch定理、POVM形式的推广）将其适用范围扩展到了正算子值测度（POVM）的情形，进一步巩固了量子测量理论的数学基础。

\hrulefill

\section{第451章：谱理论与散射量子力学}\label{ux7b2c451ux7ae0ux8c31ux7406ux8bbaux4e0eux6563ux5c04ux91cfux5b50ux529bux5b66}

量子散射理论是量子力学中连接束缚态谱（离散谱）与散射态谱（连续谱）的数学桥梁。与束缚态问题可以约化为有限维矩阵本征值问题不同，散射态涉及连续谱上的广义本征函数，其数学基础是自伴算子的谱理论和散射算子的酉等价性。本章的核心是建立量子散射的严格数学框架：从Kato-Rellich定理提供的自伴性微扰理论出发，通过Weyl准则刻画本质谱在紧微扰下的稳定性，再深入到N体问题的HVZ定理给出多粒子系统连续谱的精确下界，最后以Møller波算子和S矩阵建立起完整的散射理论。这一框架不仅为量子力学的散射实验（如粒子对撞、光电效应）提供了严格的理论基础，也在数学物理的其他领域（如波动方程、量子场论的Haag-Ruelle散射理论）中发挥着根本性的作用。

\subsection{151.1 Kato-Rellich定理与自伴性微扰}\label{kato-rellichux5b9aux7406ux4e0eux81eaux4f34ux6027ux5faeux6270}

\textbf{定义 151.1}（相对有界微扰）：设\(A\)为Hilbert空间\(\mathcal{H}\)上的自伴算子。算子\(B\)称为\textbf{\(A\)-有界}的，若\(\mathcal{D}(B) \supseteq \mathcal{D}(A)\)且存在常数\(a, b \geq 0\)使得 \[\|B\psi\| \leq a\|A\psi\| + b\|\psi\|,\quad \forall \psi \in \mathcal{D}(A).\] \(A\)-界定义为满足上述不等式的最小\(a\)（\(b\)随\(a\)确定）。若\(a < 1\)，称\(B\)是\textbf{\(A\)-小}的（Kato小）。

\textbf{定理 151.1}（Kato-Rellich定理）：设\(A\)为自伴算子，\(B\)为对称算子且是\(A\)-小的（\(A\)-界\(a < 1\)）。则\(A + B\)在\(\mathcal{D}(A)\)上是自伴的。进一步，若\(A\)是本质自伴的（在某个核\(\mathcal{D}\)上），\(A+B\)在\(\mathcal{D}\)上也是本质自伴的。

\textbf{证明}：核心思路是证明对于充分大的\(\mu\)，\(A+B \pm i\mu I\)是满射，从而由自伴算子的基本判别准则得出自伴性。

第一步：估计预解式。对\(\mu \in \mathbb{R} \setminus \{0\}\)，由\(A\)-有界性， \[\|B(A+i\mu I)^{-1}\| \leq a\|A(A+i\mu I)^{-1}\| + b\|(A+i\mu I)^{-1}\|.\] 由谱定理，\(\|A(A+i\mu I)^{-1}\| \leq 1\)（因\(|\lambda/(\lambda+i\mu)| \leq 1\)），且\(\|(A+i\mu I)^{-1}\| \leq 1/|\mu|\)。故 \[\|B(A+i\mu I)^{-1}\| \leq a + \frac{b}{|\mu|}.\] 取\(|\mu|\)充分大使得\(a + b/|\mu| < 1\)，则\(\|B(A+i\mu I)^{-1}\| < 1\)。

第二步：构造\((A+B \pm i\mu I)^{-1}\)。由Neumann级数， \[I + B(A+i\mu I)^{-1}\] 在\(\mathcal{H}\)上可逆，其逆为\(\sum_{n=0}^\infty (-1)^n[B(A+i\mu I)^{-1}]^n\)（Neumann级数收敛因为\(\|B(A+i\mu I)^{-1}\| < 1\)）。于是 \[A+B+i\mu I = (I + B(A+i\mu I)^{-1})(A+i\mu I)\] 为可逆算子（两个可逆算子的复合）。故\(A+B+i\mu I\)为满射且单射，其逆为有界算子。类似处理\(A+B-i\mu I\)。

第三步：由对称算子为自伴的充要条件（\(\operatorname{Ran}(A+B \pm i\mu I) = \mathcal{H}\)对某个\(\mu > 0\)成立），可知\(A+B\)在\(\mathcal{D}(A)\)上自伴。对于本质自伴性的情形，核\(\mathcal{D}\)上的类似论证（利用闭包与预解式的关系）保证了\(\overline{A+B|_{\mathcal{D}}}\)的自伴性。\(\blacksquare\)

\textbf{推论 151.2}（量子力学中的应用）：设\(V \in L^2(\mathbb{R}^3) + L^\infty(\mathbb{R}^3)\)为实值势函数。则Schrödinger算子\(H = -\Delta + V\)在\(H^2(\mathbb{R}^3)\)上是自伴的，且在\(C_0^\infty(\mathbb{R}^3)\)上本质自伴。

\textbf{证明}：取\(A = -\Delta\)（\(L^2(\mathbb{R}^3)\)上的自伴算子，定义域\(H^2(\mathbb{R}^3)\)）。需验证\(V\)是\(-\Delta\)-小的。对\(V = V_1 + V_2\)（\(V_1 \in L^2(\mathbb{R}^3), V_2 \in L^\infty(\mathbb{R}^3)\)），由Sobolev不等式\(\|\psi\|_\infty \leq C\|\psi\|_{H^2}\)和Hölder不等式，\(\|V_1\psi\|_2 \leq \|V_1\|_2\|\psi\|_\infty \leq C\|V_1\|_2\|(-\Delta)\psi\|_2\)。故\(\|V\psi\| \leq C\|V_1\|_2\|-\Delta\psi\| + \|V_2\|_\infty\|\psi\|\)。通过缩小\(V_1\)的支集（取截断函数），可使\(a = C\|V_1\|_2\)任意小，故\(V\)是\(-\Delta\)-小的（\(a=0\)）。由Kato-Rellich定理，\(-\Delta+V\)在\(H^2(\mathbb{R}^3)\)上自伴且在\(C_0^\infty(\mathbb{R}^3)\)上本质自伴。\(\blacksquare\)

\subsection{151.2 Weyl准则与本质谱}\label{weylux51c6ux5219ux4e0eux672cux8d28ux8c31}

\textbf{定义 151.2}（本质谱）：自伴算子\(A\)的\textbf{本质谱}\(\sigma_{\text{ess}}(A)\)定义为\(\sigma(A)\)中除去有限重孤立本征值的集合。等价地，\(\lambda \in \sigma_{\text{ess}}(A)\)当且仅当\(\lambda\)为\(A\)的聚点谱点，或为无穷重本征值。

\textbf{定理 151.3}（Weyl准则）：设\(A\)为自伴算子。则\(\lambda \in \sigma_{\text{ess}}(A)\)当且仅当存在Weyl序列------即序列\(\{\psi_n\} \subset \mathcal{D}(A)\)满足： 1. \(\|\psi_n\| = 1\)对所有\(n\)； 2. \(\psi_n \rightharpoonup 0\)（弱收敛）； 3. \(\|(A - \lambda I)\psi_n\| \to 0\)。

\textbf{证明}：（\(\Rightarrow\)）设\(\lambda \in \sigma_{\text{ess}}(A)\)。若\(\lambda\)为无穷重本征值，取其规范正交本征向量的无穷序列即得Weyl序列（弱收敛到0由Bessel不等式）。若\(\lambda\)为\(A\)的谱的聚点，取\(\lambda_n \in \sigma(A) \setminus \{\lambda\}\)严格收敛于\(\lambda\)，则对每个\(\lambda_n\)存在近似本征向量\(\psi_n\)满足\(\|(A-\lambda_n I)\psi_n\| < 1/n\)，且可选取\(\psi_n \perp \operatorname{span}\{\psi_1, \ldots, \psi_{n-1}\}\)（由谱定理的Weyl-von Neumann形式）。则\(\|(A-\lambda I)\psi_n\| \leq \|(A-\lambda_n I)\psi_n\| + |\lambda_n-\lambda| < 1/n + |\lambda_n-\lambda| \to 0\)。且\(\psi_n \rightharpoonup 0\)由正交性和Bessel不等式保证。

（\(\Leftarrow\)）设存在Weyl序列\(\{\psi_n\}\)。若\(\lambda \notin \sigma_{\text{ess}}(A)\)，则\(\lambda\)至多是有限重孤立本征值，从而存在\(\varepsilon > 0\)使得\(\sigma(A) \cap (\lambda-\varepsilon, \lambda+\varepsilon) \subseteq \{\lambda\}\)（若\(\lambda \in \sigma(A)\)）或为空（若\(\lambda \notin \sigma(A)\)）。由谱定理，算子\((A-\lambda I)E\)可逆（\(E\)为谱投影到\((\lambda-\varepsilon, \lambda+\varepsilon)^c\)上的补投影）。由Weyl序列的条件3，\(\|E\psi_n\| \to 0\)。故\(\|\psi_n\|^2 = \|E\psi_n\|^2 + \|(I-E)\psi_n\|^2 \to \|(I-E)\psi_n\|^2\)。但\((I-E)\mathcal{H}\)为有限维空间（\(\lambda\)为有限重孤立本征值），其上弱收敛等价于强收敛，\(\psi_n \rightharpoonup 0\)意味着\((I-E)\psi_n \to 0\)（强）。故\(\|\psi_n\| \to 0\)，与\(\|\psi_n\|=1\)矛盾。\(\blacksquare\)

\textbf{定理 151.4}（Weyl本质谱定理）：设\(A\)为自伴算子，\(B\)为对称算子且为\(A\)-紧的（即\(B(A+iI)^{-1}\)为紧算子）。则 \[\sigma_{\text{ess}}(A+B) = \sigma_{\text{ess}}(A).\] 即紧微扰不改变本质谱。

\textbf{证明}：只需证明\(\sigma_{\text{ess}}(A+B) \subseteq \sigma_{\text{ess}}(A)\)（由对称性得反向包含）。设\(\lambda \in \sigma_{\text{ess}}(A+B)\)，取Weyl序列\(\{\psi_n\}\)使\(\|(A+B-\lambda I)\psi_n\| \to 0\)，\(\psi_n \rightharpoonup 0\)。往证\(\lambda \in \sigma_{\text{ess}}(A)\)。由Weyl准则，只需证\(\|(A-\lambda I)\psi_n\| \to 0\)。我们有 \[(A-\lambda I)\psi_n = (A+B-\lambda I)\psi_n - B\psi_n.\] 第一项趋于0。对第二项，因为\(B\)是\(A\)-紧的，\(B(A+iI)^{-1}\)紧。由\(\psi_n \rightharpoonup 0\)，\((A+iI)^{-1}B\psi_n\)强收敛到0（紧算子将弱收敛映为强收敛）。又\((A+iI)^{-1}\)有界，故需估计\(B\psi_n\)。将\(B\psi_n = (B(A+iI)^{-1})(A+iI)\psi_n\)。由逼近定理，\((A+iI)\psi_n\)在\(\mathcal{H}\)中有界（\(A\)-有界性保证），且\(B(A+iI)^{-1}\)紧，故\(B\psi_n \to 0\)（强）。因此\(\|(A-\lambda I)\psi_n\| \to 0\)，由Weyl准则得\(\lambda \in \sigma_{\text{ess}}(A)\)。\(\blacksquare\)

\subsection{151.3 Hunziker-van Winter-Zhislin定理}\label{hunziker-van-winter-zhislinux5b9aux7406}

\textbf{定理 151.5}（HVZ定理，N体Schrödinger算子的本质谱）：考虑\(\mathbb{R}^d\)中\(N\)个粒子的Schrödinger算子 \[H = -\sum_{i=1}^N \frac{1}{2m_i}\Delta_i + \sum_{i<j} V_{ij}(x_i - x_j),\] 其中\(V_{ij}\)为两体势（满足适当短程条件，如\(V_{ij} \in L^2(\mathbb{R}^d) + L^\infty(\mathbb{R}^d)\)）。在去除质心运动后，\(H\)的本质谱下界为 \[\inf\sigma_{\text{ess}}(H) = \min_{\mathcal{P} \in \mathfrak{P}^N \setminus \{1,\ldots,N\}} \left(\sum_{C \in \mathcal{P}} E(C)\right),\] 其中\(\mathfrak{P}^N\)为\(\{1,\ldots,N\}\)的所有划分，\(E(C)\)为子团簇\(C\)的基态能量（\(E(\{i\}) = 0\)）。特别地，\(\inf\sigma_{\text{ess}}(H) = \Sigma\)（最小非平凡团簇分解能量）。

\textbf{证明概要}：设\(\Sigma = \min\{\sum E(C) : \mathcal{P} \neq \{\{1\},\ldots,\{N\}\}\)。需证\(\inf\sigma_{\text{ess}}(H) = \Sigma\)。

下界（\(\inf\sigma_{\text{ess}}(H) \geq \Sigma\)）：取Weyl序列\(\{\psi_n\} \subset \mathcal{D}(H)\)使\(\|(H-\lambda)\psi_n\| \to 0\)且\(\psi_n \rightharpoonup 0\)。需证\(\lambda \geq \Sigma\)。利用Ruelle-Simon划分函数：构造\(\mathbb{R}^{d(N-1)}\)上的单位划分\(\{J_{\mathcal{P}}\}_{\mathcal{P} \neq \{\{1\},\ldots,\{N\}\}}\)（依据团簇分解几何），每个\(J_{\mathcal{P}}\)的支集对应于粒子分成团簇\(\mathcal{P}\)的区域。计算 \[\langle\psi_n|H\psi_n\rangle = \sum_{\mathcal{P}} \langle J_{\mathcal{P}}\psi_n|H_{\mathcal{P}}J_{\mathcal{P}}\psi_n\rangle + \text{误差项}\] 其中\(H_{\mathcal{P}}\)为团簇分解后的内部Hamilton算子（忽略团簇间相互作用，因团簇间距很大时相互作用可忽略）。误差项在适当几何条件下趋于0。故\(\langle\psi_n|H\psi_n\rangle \geq \sum_{\mathcal{P}} E(\mathcal{P})\|J_{\mathcal{P}}\psi_n\|^2 + o(1) \geq \Sigma + o(1)\)，极限得\(\lambda \geq \Sigma\)。

上界（\(\inf\sigma_{\text{ess}}(H) \leq \Sigma\)）：取达到\(\Sigma\)的最优团簇分解\(\mathcal{P}_0\)。构造试探波函数\(\psi = \psi_{\text{intra}} \otimes \psi_{\text{inter}}\)，其中团簇内部取各自基态波函数，团簇间取平面波（相对动量为\(\kappa\)，\(\kappa \to 0\)）。平移团簇相对坐标使间距趋于无穷，由团簇间势能趋于零，\(H\)近似为各团簇\(H(C)\)的直和。取\(\kappa \to 0\)的Weyl序列（团簇间距趋于无穷），这些态的能量趋于\(\Sigma\)且弱趋于0，故\(\Sigma \in \sigma_{\text{ess}}(H)\)。\(\blacksquare\)

\subsection{151.4 Møller波算子与S矩阵}\label{muxf8llerux6ce2ux7b97ux5b50ux4e0esux77e9ux9635}

\textbf{定义 151.3}（Møller波算子）：设\(H_0\)和\(H = H_0 + V\)分别为自由和相互作用Schrödinger算子（均为自伴）。\textbf{Møller波算子}定义为 \[\Omega_\pm = \operatorname*{s-lim}_{t \to \pm\infty} e^{iHt}e^{-iH_0 t},\] 若该强极限存在。波算子\(\Omega_\pm\)的像空间\(\mathcal{H}_{\text{ac}}(H)\)（\(H\)的绝对连续子空间）由散射态组成。

\textbf{定理 151.6}（Cook方法）：设\(H_0\)和\(H\)如上述。若对\(\mathcal{H}_0\)的稠密子集\(\mathcal{D}\)中的\(\psi\)，有 \[\int_{-\infty}^\infty \|Ve^{-iH_0 t}\psi\|\, dt < \infty,\] 则\(\Omega_+\psi = \lim_{t \to +\infty} e^{iHt}e^{-iH_0 t}\psi\)存在。类似地，时间反向对应\(\Omega_-\)。

\textbf{证明}：设\(\psi \in \mathcal{D}\)。定义\(\psi(t) = e^{iHt}e^{-iH_0 t}\psi\)。则 \[\frac{d}{dt}\psi(t) = ie^{iHt}(H - H_0)e^{-iH_0 t}\psi = ie^{iHt}Ve^{-iH_0 t}\psi.\] 积分得 \[\psi(t) - \psi(s) = i\int_s^t e^{iH\tau}Ve^{-iH_0 \tau}\psi\, d\tau.\] 由酉算子的保范性，\(\|e^{iH\tau}Ve^{-iH_0 \tau}\psi\| = \|Ve^{-iH_0 \tau}\psi\|\)。若\(\int_{-\infty}^\infty \|Ve^{-iH_0 \tau}\psi\|\, d\tau < \infty\)，则\(\psi(t)\)的Cauchy条件满足（当\(t \to \pm\infty\)时），故强极限存在。由\(\mathcal{D}\)的稠密性和波算子的有界性（\(\|\Omega_\pm\| \leq 1\)），波算子在\(\mathcal{D}\)的闭包上存在。\(\blacksquare\)

\textbf{定理 151.7}（S矩阵与散射算子）：\textbf{散射算子}\(S: \mathcal{H} \to \mathcal{H}\)定义为\(S = \Omega_-^* \Omega_+\)。\(S\)是酉算子且满足\(SH_0 = H_0 S\)（S矩阵与自由演化交换），因此\(S\)在\(H_0\)的谱表示中对角化，其纤维为散射矩阵\(S(\lambda): \mathcal{H}_\lambda \to \mathcal{H}_\lambda\)（作用在\(H_0 = \lambda\)的谱子空间上）。

\textbf{证明}：首先，\(\Omega_\pm^* \Omega_\pm = P_{\text{ac}}(H_0)\)（到\(H_0\)的绝对连续子空间上的投影），\(\Omega_\pm \Omega_\pm^* = P_{\text{ac}}(H)\)（到\(H\)的绝对连续子空间上的投影）。这些由波算子的交错性质（\(\Omega_\pm e^{-iH_0 t} = e^{-iHt}\Omega_\pm\)）经极限手续推出。

其次，波算子的完备性（\(\operatorname{Ran}\Omega_+ = \operatorname{Ran}\Omega_- = \mathcal{H}_{\text{ac}}(H)\)）在适当的短程条件下成立（Kato-Kuroda-Birman理论）。在完备性条件下，\(S = \Omega_-^* \Omega_+\)为酉算子。

\(S\)与\(H_0\)交换：\(Se^{-iH_0 t} = \Omega_-^* \Omega_+ e^{-iH_0 t} = \Omega_-^* e^{-iHt}\Omega_+ = e^{-iH_0 t}\Omega_-^* \Omega_+ = e^{-iH_0 t}S\)。故\(S\)在\(H_0\)的谱测度下可约化。由直积分分解\(\mathcal{H} \cong \int_\sigma^\oplus \mathcal{H}_\lambda\, d\lambda\)（\(H_0 = \int \lambda\, dP_0(\lambda)\)），\(S = \int_\sigma^\oplus S(\lambda)\, d\lambda\)，其中\(S(\lambda)\)为\(\mathcal{H}_\lambda\)上的酉算子（散射矩阵）。\(\blacksquare\)

\textbf{定理 151.8}（定态散射理论与Lippmann-Schwinger方程）：波算子的像\(\psi_\pm = \Omega_\pm \phi\)（\(\phi \in \mathcal{H}_{\text{ac}}(H_0)\)）满足Lippmann-Schwinger方程 \[\psi_\pm = \phi + \lim_{\varepsilon \to 0^+} (H_0 - \lambda \mp i\varepsilon)^{-1}V\psi_\pm.\] 该方程是量子散射计算的基础。

\textbf{证明}：由波算子的定义，对\(\phi \in \mathcal{H}_{\text{ac}}(H_0)\)， \[\Omega_\pm \phi = \lim_{t \to \mp\infty} e^{iHt}e^{-iH_0 t}\phi = \phi + \lim_{t \to \mp\infty} i\int_0^t e^{iHs}Ve^{-iH_0 s}\phi\, ds.\] 设\(H_0\phi = \lambda\phi\)（先考虑广义本征态）。则\(e^{-iH_0 s}\phi = e^{-i\lambda s}\phi\)，且\(e^{iHs} = \frac{1}{2\pi i} \oint e^{izs}(z-H)^{-1}dz\)（由预解式围道积分）。代入并交换积分次序，利用\(\int_0^{\mp\infty} e^{i(z-\lambda)s}ds = \lim_{\varepsilon \to 0^+} \frac{i}{z - \lambda \mp i\varepsilon}\)，得到\((\lambda \mp i\varepsilon - H)^{-1} = ((H_0 \mp i\varepsilon - \lambda) - V)^{-1}\)。展开该预解式（第二预解恒等式）即得Lippmann-Schwinger方程。该推导可在分布意义下严格化（利用\(H_0\)的谱表示中的Hilbert变换）。\(\blacksquare\) \hrulefill

\hrulefill

\hrulefill

\newpage

\chapter{13-数学物理/数学物理/04-广义相对论.md}\label{ux6570ux5b66ux7269ux7406ux6570ux5b66ux7269ux740604-ux5e7fux4e49ux76f8ux5bf9ux8bba.md}

\section{第452章：广义相对论数学基础}\label{ux7b2c452ux7ae0ux5e7fux4e49ux76f8ux5bf9ux8bbaux6570ux5b66ux57faux7840}

广义相对论的数学框架是伪 Riemann 几何（Lorentz 流形），其 Einstein 场方程可从 Einstein-Hilbert 作用量的变分原理导出。本章建立度规、曲率张量、Einstein 方程及奇点定理等核心数学结构。

\subsection{151.1 Lorentz 流形与时空}\label{lorentz-ux6d41ux5f62ux4e0eux65f6ux7a7a}

\textbf{定义 151.1}（时空）：\textbf{时空}是四维 Lorentz 流形 \((M, g)\)，其中 \(g\) 是 Lorentz 度量（符号差 \((-, +, +, +)\) 或 \((+, -, -, -)\)）。每点的切空间 \(T_pM\) 被光锥分为类时、类光和类空向量。

\textbf{定义 151.2}（测地线）：类时测地线满足

\[\frac{d^2 x^\mu}{d\tau^2} + \Gamma^\mu_{\nu\rho} \frac{dx^\nu}{d\tau} \frac{dx^\rho}{d\tau} = 0\]

其中 \(\tau\) 是仿射参数，\(\Gamma^\mu_{\nu\rho}\) 是 Christoffel 符号。

\textbf{定义 151.3}（Riemann 曲率张量）：曲率描述了时空的弯曲程度：

\[R^\rho_{\sigma\mu\nu} = \partial_\mu \Gamma^\rho_{\nu\sigma} - \partial_\nu \Gamma^\rho_{\mu\sigma} + \Gamma^\rho_{\mu\lambda}\Gamma^\lambda_{\nu\sigma} - \Gamma^\rho_{\nu\lambda}\Gamma^\lambda_{\mu\sigma}\]

\textbf{Ricci 张量}：\(R_{\mu\nu} = R^\rho_{\mu\rho\nu}\)，\textbf{标量曲率}：\(R = g^{\mu\nu}R_{\mu\nu}\)。

\subsection{151.2 Einstein 场方程}\label{einstein-ux573aux65b9ux7a0b}

\textbf{定理 151.1}（Einstein 场方程）：时空度规 \(g_{\mu\nu}\) 满足 Einstein 场方程 \[R_{\mu\nu} - \frac{1}{2}R g_{\mu\nu} + \Lambda g_{\mu\nu} = \kappa T_{\mu\nu}\] 其中 \(T_{\mu\nu}\) 是能动张量（源项），\(\Lambda\) 是宇宙学常数，\(\kappa = 8\pi G\) 是耦合常数。该方程可由 Einstein-Hilbert 作用量的变分原理导出。

\textbf{证明}：从 Einstein-Hilbert 作用量出发： \[S[g] = \frac{1}{2\kappa}\int_M (R - 2\Lambda)\sqrt{-g}\,d^4x + S_{\text{matter}}[g, \Phi].\] 计算度规的变分。利用 Palatini 恒等式 \(\delta R_{\mu\nu} = \nabla_\rho \delta\Gamma^\rho_{\nu\mu} - \nabla_\nu \delta\Gamma^\rho_{\rho\mu}\) 和 \(\delta(\sqrt{-g}) = -\frac{1}{2}\sqrt{-g}g_{\mu\nu}\delta g^{\mu\nu}\)，纯引力部分变分为： \[\delta S_{\text{EH}} = \frac{1}{2\kappa}\int_M \left(R_{\mu\nu} - \frac{1}{2}Rg_{\mu\nu} + \Lambda g_{\mu\nu}\right)\delta g^{\mu\nu}\sqrt{-g}\,d^4x.\] 定义 Einstein 张量 \(G_{\mu\nu} = R_{\mu\nu} - \frac{1}{2}Rg_{\mu\nu}\)。物质部分的变分定义能动张量： \[\delta S_{\text{matter}} = -\frac{1}{2}\int_M T_{\mu\nu}\delta g^{\mu\nu}\sqrt{-g}\,d^4x.\] 由最小作用量原理 \(\delta S = 0\) 对所有 \(\delta g^{\mu\nu}\) 成立，由变分法基本引理得： \[G_{\mu\nu} + \Lambda g_{\mu\nu} = \kappa T_{\mu\nu}.\] 此即 Einstein 场方程。由第二 Bianchi 恒等式 \(\nabla_{[\lambda}R_{\mu\nu]\rho\sigma}=0\) 缩并可得 \(\nabla^\mu G_{\mu\nu}=0\)，从而 \(\nabla^\mu T_{\mu\nu}=0\)------能量动量守恒是场方程的内在约束。\(\blacksquare\)

\textbf{定义 151.4}（Einstein-Hilbert 作用量）：Einstein 方程可由变分原理导出。作用量为

\[S = \frac{1}{2\kappa} \int_M (R - 2\Lambda) \sqrt{-g} \, d^4x + S_{\text{matter}}\]

\textbf{定理 151.2}（Bianchi 恒等式的几何意义）：\(\nabla^\mu G_{\mu\nu} = 0\)（其中 \(G_{\mu\nu} = R_{\mu\nu} - \frac{1}{2}R g_{\mu\nu}\) 是 Einstein 张量），这是微分几何的 Bianchi 恒等式 \(\nabla_{[\lambda} R_{\mu\nu]\rho\sigma} = 0\) 的推论，保证了能动张量的协变守恒 \(\nabla^\mu T_{\mu\nu} = 0\)。

\textbf{证明}：由第二 Bianchi 恒等式 \(\nabla_{[\lambda} R_{\mu\nu]\rho\sigma} = 0\)，缩并 \(\rho\) 与 \(\lambda\)、\(\sigma\) 与 \(\nu\)：\(g^{\rho\lambda} g^{\sigma\nu} \nabla_{[\lambda} R_{\mu\nu]\rho\sigma} = 0\)。展开反对称和得 \(\nabla_\rho R_{\mu\sigma} - \nabla_\sigma R_{\mu\rho} + \nabla^\nu R_{\mu\nu\sigma\rho} = 0\)。再缩并 \(\mu\) 与 \(\sigma\)：\(g^{\mu\sigma}\) 作用于上式，得 \(\nabla_\rho R - \nabla^\mu R_{\mu\rho} - \nabla^\mu R_{\mu\rho} = 0\)，即 \(\nabla_\rho R - 2\nabla^\mu R_{\mu\rho} = 0\)。整理得 \(\nabla^\mu(R_{\mu\rho} - \frac{1}{2}R g_{\mu\rho}) = \nabla^\mu G_{\mu\rho} = 0\)。由 Einstein 场方程 \(G_{\mu\nu} = 8\pi G T_{\mu\nu}\)，立得 \(\nabla^\mu T_{\mu\nu} = 0\)，即能量动量守恒。\(\blacksquare\)

\subsection{151.3 经典精确解}\label{ux7ecfux5178ux7cbeux786eux89e3}

\textbf{定理 151.3}（Birkhoff 定理）：Schwarzschild 解是球对称真空 Einstein 方程的唯一解。

\textbf{证明}：设球对称度规 \(ds^2 = -A(t,r)dt^2 + B(t,r)dr^2 + r^2d\Omega^2\)（可消去 \(dtdr\) 项）。真空方程 \(R_{\mu\nu}=0\) 给出 \(\partial_r(AB)=0\)（\(AB=f(t)\)）且 \(\partial_t B=0\)，故 \(B\) 仅依赖于 \(r\)。通过时间重标度可使 \(A\) 也仅依赖于 \(r\)。求解 ODE 得 \(A=1/B=1-2M/r\)（\(M\) 为积分常数）。故 Schwarzschild 解唯一且必为静态------球对称真空引力场必为静态且渐近平坦。\(\blacksquare\)

Einstein 场方程的非线性使得精确解的寻找极为困难，但对称性假设可以大幅简化方程。除了 Schwarzschild 解外，以下是几个最重要的经典精确解，它们具有深刻的物理意义和数学结构：

\textbf{（1）Schwarzschild 解}（1916年）：\(ds^2 = -(1 - \frac{2M}{r})dt^2 + (1 - \frac{2M}{r})^{-1}dr^2 + r^2 d\Omega^2\)。这是第一个 Einstein 方程精确解，描述球对称无电荷无自转天体外部的引力场。\(r = 2M\) 处坐标奇异性（事件视界）由 Eddington-Finkelstein 坐标或 Kruskal-Szekeres 坐标消除。\(r = 0\) 处是曲率奇点。该解预言了黑洞的存在，其参数 \(M\) 通过 Kepler 第三定律诠释为天体质量。

\textbf{（2）Kerr 解}（1963年）：\(ds^2 = -(1 - \frac{2Mr}{\Sigma})dt^2 - \frac{4Mar\sin^2\theta}{\Sigma}dt d\phi + \frac{\Sigma}{\Delta}dr^2 + \Sigma d\theta^2 + (r^2 + a^2 + \frac{2Ma^2 r \sin^2\theta}{\Sigma})\sin^2\theta d\phi^2\)，其中 \(\Sigma = r^2 + a^2\cos^2\theta\)，\(\Delta = r^2 - 2Mr + a^2\)。描述自转黑洞（\(a = J/M\) 为单位质量角动量）。Kerr 解具有两个视界（\(r_\pm = M \pm \sqrt{M^2 - a^2}\)）和能层（ergosphere），能层内的拖曳效应（frame-dragging）使 Penrose 过程可提取黑洞自转能。Kerr 度规的''无毛定理''（Israel-Carter-Robinson）断言：稳态渐近平坦黑洞由 \(M, J, Q\) 三个参数完全确定。

\textbf{（3）Reissner-Nordström 解}：\(f(r) = 1 - \frac{2M}{r} + \frac{Q^2}{r^2}\)，描述带电黑洞。具有两个视界（内视界和外视界），当 \(|Q| = M\) 时为极端黑洞（两个视界重合，Hawking 温度为零）。Kerr-Newman 解（1965年）将自转和电荷统一：\(f(r) = 1 - \frac{2Mr - Q^2}{\Sigma}\)。

\textbf{（4）Friedmann-Lemaître-Robertson-Walker (FLRW) 度规}：\(ds^2 = -dt^2 + a(t)^2 [\frac{dr^2}{1 - k r^2} + r^2 d\Omega^2]\)，其中 \(k \in \{-1, 0, 1\}\) 为空间曲率。这是宇宙学标准模型（\(\Lambda\)CDM）的几何基础，描述均匀各向同性宇宙。尺度因子 \(a(t)\) 满足 Friedmann 方程：\((\frac{\dot{a}}{a})^2 = \frac{8\pi G}{3}\rho - \frac{k}{a^2} + \frac{\Lambda}{3}\)。\(a(t)\) 的演化决定了宇宙的膨胀、收缩或加速膨胀历史。

\subsection{151.4 奇点定理}\label{ux5947ux70b9ux5b9aux7406}

\textbf{定义 151.5}（测地不完备性与奇点）：\textbf{奇点}定义为测地线不完备性------时空无法沿测地线延拓。

\textbf{定理 151.4}（Penrose 奇点定理）：若时空包含捕获面（trapped surface），且满足零能量条件 \(T_{\mu\nu}k^\mu k^\nu \geq 0\)（对所有类光向量 \(k\)），则存在不完备的类光测地线。

\textbf{证明}：设 \(\mathcal{T}\) 为闭合捕获面（内外向类光测地线束均收敛）。Raychaudhuri 方程：\(d\theta/d\lambda = -\frac{1}{2}\theta^2 - \sigma^2 - R_{\mu\nu}k^\mu k^\nu\)。在 \(\mathcal{T}\) 上 \(\theta<0\)，零能量条件保证 \(d\theta/d\lambda \leq 0\)，故 \(\theta\) 在有限仿射参数内趋于 \(-\infty\)------出现焦散点。若未来 Cauchy 视界 \(H^+(\mathcal{T})\) 为紧致无边，则为紧致三维类光超曲面，导出矛盾。故必存在不完备类光测地线（即奇点）。\(\blacksquare\)

\textbf{定理 151.5}（Hawking 奇点定理）：在满足适当能量条件的时空中，过去的类时测地不完备性不可避免。

\textbf{证明}：设时空为整体双曲，满足强能量条件 \(R_{\mu\nu}V^\mu V^\nu \geq 0\)（对所有类时 \(V\)）。考虑过去方向类时测地线束的 Raychaudhuri 方程：\(d\theta/d\tau = -\frac{1}{3}\theta^2 - \sigma^2 - R_{\mu\nu}u^\mu u^\nu \leq 0\)。若在某点 \(\theta < 0\)，则过去方向必在有限固有时内出现焦散点。若宇宙过去的类时测地线均收敛于 Cauchy 面，紧致论证导出类时测地不完备性。这在宇宙学中意味着大爆炸奇点的不可避免性。\(\blacksquare\)

\textbf{定理 151.6}（黑洞热力学四定律 / Bardeen-Carter-Hawking）： - 第零定律：稳态黑洞视界上的表面引力 \(\kappa\) 为常数 - 第一定律：\(dM = \frac{\kappa}{8\pi} dA + \Omega dJ + \Phi dQ\) - 第二定律：视界面积 \(A\) 满足 \(\delta A \geq 0\) - 第三定律：不可能通过有限步骤使 \(\kappa\) 达到零

\textbf{证明}：第零定律：Killing 向量场 \(\xi^a\) 在视界上满足 \(\nabla^a(\xi^b\xi_b) = -2\kappa\xi^a\)，由 Einstein 方程和能量条件知 \(\kappa\) 沿视界不变。第一定律由 Einstein 方程线性化后的 ADM 质量变分公式和 Komar 积分给出。第二定律（Hawking 面积定理）：零能量条件保证类光测地线束膨胀 \(\theta \geq 0\)（无焦散时面积不减）。第三定律：极端黑洞（\(\kappa=0\)）对应 \(T=0\) 基态，有限操作无法达到。\(\blacksquare\)

\subsection{151.5 线性化引力与渐近平坦时空}\label{ux7ebfux6027ux5316ux5f15ux529bux4e0eux6e10ux8fd1ux5e73ux5766ux65f6ux7a7a}

\textbf{定义 151.7}（线性化 Einstein 方程）：线性化 \(g_{\mu\nu} = \eta_{\mu\nu} + h_{\mu\nu}\)（\(|h_{\mu\nu}| \ll 1\)）在 Lorentz 规范下给出波动方程 \(\square \bar{h}_{\mu\nu} = -2\kappa T_{\mu\nu}\)。真空中 \(\square \bar{h}_{\mu\nu} = 0\)。

\textbf{定理 151.7}（Bondi-Sachs 质量）：对渐近平坦时空，在类光无穷远 \(\mathscr{I}^+\) 处可定义 Bondi 质量，其变化率等于

\textbf{证明}：Bondi质量\(M_B(u)=\frac{1}{8\pi}\int_{S^2}\Psi_2^0(u,\theta,\phi)d\Omega\)（\(\Psi_2^0\)为Weyl标量在\(\mathscr{I}^+\)上的领头阶）。由Bianchi恒等式\(\nabla^\mu C_{\mu\nu\rho\sigma}=0\)，在\(\mathscr{I}^+\)上：\(\frac{dM_B}{du}=-\frac{1}{2}\int_{S^2}|\sigma^0|^2 d\Omega\le 0\)（\(\sigma^0\)为news张量）。即Bondi质量随推迟时间单调递减，减量等于引力波携带的能量通量。\(\blacksquare\)辐射携带的能量通量。

\textbf{定义 151.8}（ADM 质量 / Arnowitt-Deser-Misner）：对渐近平坦时空，在类空无穷远处定义的总能量（ADM 质量）为

\[M_{ADM} = \frac{1}{2\kappa} \lim_{r \to \infty} \int_{S^2_r} (\partial_j h_{ij} - \partial_i h_{jj}) n^i \, dA\]

\textbf{定理 151.8}（正质量定理 / Schoen-Yau, Witten）：若时空满足主导能量条件，则 \(M_{ADM} \geq 0\)，且 \(M_{ADM} = 0\) 当且仅当时空是 Minkowski 时空。

\textbf{证明}：设 \(\Sigma\) 为渐近平坦类空超曲面，其上诱导度规为 \(h_{ij}\)，第二基本形式为 \(K_{ij}\)。在 \(\Sigma\) 上构造 Witten 旋量场 \(\psi\)（取值于旋量丛），满足推广的 Dirac-Witten 方程： \[\gamma^i(\nabla_i - \frac{1}{2}K_{ij}\gamma^j\gamma^0)\psi = 0,\] 其中 \(\gamma^\mu\) 为 Clifford 代数元素，\(\nabla_i\) 为旋量丛上的诱导联络。边界条件为 \(\psi \to \psi_0\)（常数旋量）当 \(r \to \infty\)。由分部积分可得 Witten 恒等式： \[\int_{\partial\Sigma_\infty} \langle\psi, (\gamma^i\nabla_i - \frac{1}{2}K_{ij}\gamma^i\gamma^j\gamma^0)\psi\rangle \,dS = \int_\Sigma \left(|\nabla\psi|^2 + 8\pi G\,T_{\mu\nu}n^\mu n^\nu|\psi|^2\right) dV.\] 左端的渐近分析给出 \(4\pi G\,M_{\text{ADM}}|\psi_0|^2\)（其中 \(M_{\text{ADM}}\) 提取自 \(h_{ij} \sim \delta_{ij} + \frac{2M_{\text{ADM}}}{r}\delta_{ij} + O(r^{-2})\) 渐近展开）。右端被积函数：主导能量条件 \(T_{\mu\nu}n^\mu n^\nu \geq |T_{\mu\nu}n^\mu e_i^\nu|\)（\(e_i^\mu\) 为 \(\Sigma\) 上的正交标架）保证第二项非负，而 \(|\nabla\psi|^2 \geq 0\) 显然。故 \(M_{\text{ADM}} \geq 0\)。若 \(M_{\text{ADM}} = 0\)，则 \(\nabla\psi = 0\) 且 \(T_{\mu\nu} = 0\)，此时 \(\Sigma\) 为 \(\mathbb{R}^3\) 配备平坦度规，整个时空为 Minkowski 时空。\(\blacksquare\)

\section{第453章：Einstein方程解的存在性与奇点定理}\label{ux7b2c453ux7ae0einsteinux65b9ux7a0bux89e3ux7684ux5b58ux5728ux6027ux4e0eux5947ux70b9ux5b9aux7406}

Einstein场方程是一组高度非线性的双曲-椭圆耦合偏微分方程组，其数学分析是数学广义相对论的核心课题。本章从偏微分方程和整体微分几何的角度讨论Einstein方程的解的存在性、唯一性及其奇异行为。Penrose-Hawking奇点定理是广义相对论定性理论中最深刻的结论之一，它不依赖于具体的对称性假设，仅从能量条件和整体因果结构出发就断言时空奇点的不可避免性------这不仅预言了黑洞和宇宙大爆炸的奇异性，更标志着广义相对论预言自身失效的内在界限。黑洞唯一性定理（无毛定理）则从另一个方向揭示了Einstein方程的刚性：稳态渐近平坦黑洞仅由质量、角动量和电荷三个参数完全确定，与形成过程无关。正质量定理（Schoen-Yau和Witten）是渐近平坦时空的基本能量不等式，其证明将几何分析（极小子流形、Dirac算子）与广义相对论的能量概念紧密结合，堪称20世纪几何分析最辉煌的胜利之一。

\subsection{152.1 Penrose-Hawking奇点定理的完整证明}\label{penrose-hawkingux5947ux70b9ux5b9aux7406ux7684ux5b8cux6574ux8bc1ux660e}

\textbf{定义 152.1}（能量条件）：设\((M,g)\)为时空，\(T_{\mu\nu}\)为能动张量。 - \textbf{零能量条件（NEC）}：\(T_{\mu\nu}k^\mu k^\nu \geq 0\)对所有类光向量\(k^\mu\)； - \textbf{弱能量条件（WEC）}：\(T_{\mu\nu}V^\mu V^\nu \geq 0\)对所有类时向量\(V^\mu\)； - \textbf{强能量条件（SEC）}：\((T_{\mu\nu} - \frac{1}{2}Tg_{\mu\nu})V^\mu V^\nu \geq 0\)对所有类时向量\(V^\mu\)。由Einstein方程等价于\(R_{\mu\nu}V^\mu V^\nu \geq 0\)。

\textbf{定理 152.1}（Penrose奇点定理，1965）：设\((M,g)\)满足：(1) 不存在闭合类时曲线；(2) \(\operatorname{Ric}(k,k) \geq 0\)对所有类光向量\(k\)（零能量条件）；(3) 存在非紧Cauchy面\(\Sigma\)；(4) \(\Sigma\)中存在闭合捕获面\(\mathcal{T}\)（即类光第二基本形式的迹\(\chi < 0\)）。则\((M,g)\)必为类光测地不完备（即存在奇点）。

\textbf{证明}：第一步（捕获面的几何）：设\(\mathcal{T} \subset \Sigma\)为紧致二维闭合（无边）类空曲面，其面积形式为\(dA\)。\(\mathcal{T}\)有两个正交类光法向量场\(\ell_\pm\)（向外和向内）。类光第二基本形式为\(\chi^\pm_{ab} = g(\nabla_a \ell_\pm, e_b)\)（\(e_a\)为\(\mathcal{T}\)的切标架），其迹\(\theta_\pm = g^{ab}\chi^\pm_{ab}\)称为膨胀（expansion）。\(\mathcal{T}\)为捕获面意味着\(\theta_\pm < 0\)处处成立------外内两个方向的类光测地线束均收敛。

第二步（Raychaudhuri方程）：沿\(\ell_+\)的类光测地线束，膨胀\(\theta\)满足Raychaudhuri方程： \[\frac{d\theta}{d\lambda} = -\frac{1}{2}\theta^2 - \sigma_{ab}\sigma^{ab} - R_{\mu\nu}\ell^\mu\ell^\nu,\] 其中\(\lambda\)为仿射参数，\(\sigma_{ab}\)为切变张量。由零能量条件\(R_{\mu\nu}\ell^\mu\ell^\nu \geq 0\)，且\(\sigma_{ab}\sigma^{ab} \geq 0\)，故\(d\theta/d\lambda \leq -\theta^2/2\)。分离变量得\(\theta^{-1}(\lambda) \geq \theta^{-1}(0) + \lambda/2\)。由于\(\theta(0) = \theta_+ < 0\)，存在有限仿射参数\(\lambda_0 \leq 2/|\theta(0)|\)使得\(\theta(\lambda) \to -\infty\)------即出现\textbf{共轭点}（焦散点），测地线在此汇聚。

第三步（整体因果结构）：设\(\mathcal{T}\)的Cauchy发展区域为\(D^+(\Sigma)\)。定义\(E^+(\mathcal{T})\)为从\(\mathcal{T}\)出发的未来外向类光测地线段的集合。由第二步，每条测地线在有限仿射参数内遇到共轭点。由类光测地线的变分理论（类光焦散定理），过共轭点后测地线不再是到\(\mathcal{T}\)的类光最长线，故\(E^+(\mathcal{T})\)的边界\(\partial E^+(\mathcal{T})\)是一族从\(\mathcal{T}\)出发且在共轭点前结束的类光测地线。由于\(\mathcal{T}\)紧致且\(\theta_+\)连续，存在一致上界\(\lambda_{\max}\)使得所有从\(\mathcal{T}\)出发的未来外向类光测地线在\(\lambda \leq \lambda_{\max}\)内遇到共轭点。

第四步（导出矛盾）：考虑未来Cauchy视界\(H^+(\Sigma) = \overline{D^+(\Sigma)} \setminus I^-(D^+(\Sigma))\)。由于\(\mathcal{T}\)捕获，\(\partial I^+(\mathcal{T})\)（\(\mathcal{T}\)的未来类光边界）是\(D^+(\Sigma)\)中的紧致类光超曲面（无边）。但\(H^+(\Sigma)\)由不完备类光测地线生成（此由整体双曲性得出），而\(\mathcal{T}\)紧致意味着这些测地线的起始集也是紧致的，导致矛盾（紧致无边类光超曲面不能嵌入适当的因果结构中）。因此\(D^+(\Sigma)\)不能包含所有从\(\mathcal{T}\)出发的完备测地线段，故存在不完备类光测地线------即时空奇点。\(\blacksquare\)

\textbf{定理 152.2}（Hawking奇点定理，1967）：设\((M,g)\)为满足强能量条件\(R_{\mu\nu}V^\mu V^\nu \geq 0\)的整体双曲时空，且存在紧致类空Cauchy面\(\Sigma\)，其平均曲率张量满足\(\operatorname{tr}_\Sigma K > 0\)（或膨胀\(\theta < 0\)）。则\((M,g)\)必为类时测地过去不完备。

\textbf{证明}：（框架）设\(\Sigma\)为Cauchy面，考虑从\(\Sigma\)出发的过去方向类时测地线束。其膨胀\(\theta\)满足类时Raychaudhuri方程： \[\frac{d\theta}{d\tau} = -\frac{1}{3}\theta^2 - \sigma_{ab}\sigma^{ab} - R_{\mu\nu}u^\mu u^\nu + \omega_{ab}\omega^{ab},\] 其中\(\tau\)为固有时，\(\sigma_{ab}\)为切变，\(\omega_{ab}\)为旋转。与类光情形不同，此处出现\(\omega_{ab}\omega^{ab} \geq 0\)项（``离心势垒''），但类时测地线束的超曲面正交性（Frobenius定理：梯度正交于\(\Sigma\)的类时法向量\(\partial/\partial\tau\)为超曲面正交）迫使\(\omega_{ab}=0\)。在有\(\theta|_\Sigma = \operatorname{tr}_\Sigma K < 0\)和强能量条件下，\(d\theta/d\tau \leq -\theta^2/3\)。与Penrose证明类似，\(\theta \to -\infty\)在有限过去固有时内发生。由Cauchy面的紧致性，过去方向的所有类时测地线均有有限的焦散时间，导致类时测地过去不完备。这在宇宙学中对应于大爆炸奇点。\(\blacksquare\)

\subsection{152.2 黑洞唯一性定理}\label{ux9ed1ux6d1eux552fux4e00ux6027ux5b9aux7406}

\textbf{定理 152.3}（Israel定理，1967）：任何静态、渐近平坦的真空黑洞解必为Schwarzschild解。即静态黑洞无''毛''------仅由质量\(M\)一个参数描述。

\textbf{证明}：设静态时空具有Killing向量场\(\xi = \partial_t\)（类时无穷远）且超曲面正交（静态条件：\(\xi \wedge d\xi = 0\)）。度规可写为 \[ds^2 = -V^2 dt^2 + h_{ij}dx^i dx^j,\] 其中\(V\)和\(h_{ij}\)与\(t\)无关。真空Einstein方程\(R_{\mu\nu} = 0\)给出\(V\)和\(h_{ij}\)的耦合方程。

核心技巧（Israel方法）：在视界\(S\)（\(V=0\)的紧致二维曲面）上，定义\(S\)的Gauss曲率\(K\)。由Einstein方程和视界上的诱导几何，可导出一个恒等式。设\(S\)为连通紧致曲面，配备诱导度规\(\gamma_{ab}\)。由真空方程和Gauss-Codazzi方程，\(S\)上的Ricci标量满足 \[^{(2)}R = \text{const} + \text{非负项}.\] 由于\(S\)紧致，Gauss-Bonnet定理\(\int_S {}^{(2)}R\, dA = 4\pi\chi(S)\)（\(\chi(S) = 2\)为\(S^2\)的Euler示性数）。代入\(^{(2)}R\)的表达式，非负项的积分为零迫使\(S\)为球面且度规为圆对称的。

由于\(S\)为圆对称，可将对称性由视界解析延拓到全时空（利用真空Einstein方程的椭圆特性和唯一性延拓），得到整个时空的球对称性。由Birkhoff定理，真空球对称解必为Schwarzschild。故\(M\)完全确定时空。\(\blacksquare\)

\textbf{定理 152.4}（Carter-Robinson无毛定理，1971-1975）：任何稳态、渐近平坦的真空黑洞解必为Kerr解（仅由\(M\)和\(J\)两个参数确定）。

\textbf{证明概要}：设\(\xi = \partial_t\)和\(\eta = \partial_\phi\)为两个交换的Killing向量场（稳态且轴对称，Carter证明了一般性）。度规的Weyl张量\(C_{\mu\nu\rho\sigma}\)满足一个二阶椭圆方程（真空Einstein方程下）。

第一步：构造Ernst势\(\mathcal{E} = f + i\psi\)，其中\(f = -g(\xi,\xi)\)，\(d\psi = *(\eta \wedge d\xi)\)。真空Einstein方程等价于Ernst方程： \[(\operatorname{Re}\mathcal{E})\nabla^2\mathcal{E} = (\nabla\mathcal{E})^2.\]

第二步：在视界上，\(\mathcal{E}\)满足特定的边界条件。由椭圆方程的正则性和极值原理，结合视界的Killing视界条件（\(\xi + \Omega_H \eta\)在视界上类光），可证\(\mathcal{E}\)由两个积分常数（\(M\)和\(J\)）完全确定。

第三步（Robinson的完美分类）：从Ernst势出发，通过Bäcklund变换或Sibgatullin积分方程方法，可以构造出所有稳态轴对称真空解，它们构成一个两参数族，即为Kerr族。Weyl张量的代数分类（Petrov类型D）进一步约束了解的形式，使其恰好对应Kerr度规。\(\blacksquare\)

\textbf{定理 152.5}（Mazur-Bekenstein唯一性定理）：Einstein-Maxwell理论中的稳态黑洞由\(M, J, Q\)确定（Kerr-Newman解）。

\subsection{152.3 正质量定理}\label{ux6b63ux8d28ux91cfux5b9aux7406}

\textbf{定理 152.6}（正质量定理，Schoen-Yau 1979, Witten 1981）：设\((\Sigma, h_{ij}, K_{ij})\)为满足主导能量条件（\(\mu \geq |J_i|\)，其中\(\mu = T_{\mu\nu}n^\mu n^\nu\), \(J_i = T_{\mu\nu}n^\mu e_i^\nu\)）的渐近平坦初始数据集。则ADM质量\(M_{\text{ADM}} \geq 0\)，且\(M_{\text{ADM}} = 0\)当且仅当\((\Sigma, h_{ij}, K_{ij})\)可等距嵌入Minkowski时空中。

\textbf{Schoen-Yau证明（极小子流形方法）}：设\(\Sigma\)为渐近平坦，\(M_{\text{ADM}} < 0\)。第一步：若\(M_{\text{ADM}} < 0\)，对充分大的\(R\)在\(\Sigma\)截断球\(B_R\)内极小化面积泛函，存在光滑的极小曲面\(\Omega\)（Plateau问题在非紧流形上的推广），使得\(\Omega\)的面积在所有与\(\partial B_R\)同调的超曲面中达到极小。

第二步：研究\(\Omega\)上的稳定性条件（二阶变分非负）。稳定性不等式给出 \[\int_\Omega (|\nabla\phi|^2 - (R_{ij}n^i n^j + \|II\|^2)\phi^2)\, dA \geq 0\] 对任意紧支光滑函数\(\phi\)，其中\(II\)为\(\Omega\)的第二基本形式。

第三步：由Gauss-Codazzi方程和标量曲率约束（Einstein约束方程\(R + (\operatorname{tr}K)^2 - K_{ij}K^{ij} = 16\pi\mu\)），可得\(\Omega\)上的标量曲率下界估计。当\(M_{\text{ADM}} < 0\)时，细心的渐近分析给出\(\Omega\)上存在矛盾（面积增长与Gauss-Bonnet定理冲突），故\(M_{\text{ADM}} \geq 0\)。\(M_{\text{ADM}} = 0\)的情形需更精细的刚性论证（利用极小超曲面的完全测地性质）。

\textbf{Witten证明（旋量方法）}：已在第151.5节详述。其核心优势在于：证明极其简洁（仅靠Dirac-Witten方程和分部积分），且自然适用于任意维数的旋量流形。\(\blacksquare\)

\textbf{定理 152.7}（Penrose不等式）：设\((\Sigma, h_{ij})\)为渐近平坦且无边界（或有极小边界）的Riemann流形，标量曲率\(R \geq 0\)。若\(\partial\Sigma\)为最外极小曲面（面积\(A\)），则 \[M_{\text{ADM}} \geq \sqrt{\frac{A}{16\pi}}.\] 等式成立当且仅当\((\Sigma, h_{ij})\)为Schwarzschild初始数据（等距于\(t = \text{const}\)切片）。

\textbf{证明概要}（Huisken-Ilmanen 2001, Bray 2001）：利用逆平均曲率流（inverse mean curvature flow）∂\_t F\_t = (1/H)ν，面积\(A(t)\)和Hawking质量\(m_H(t)\)均随时间演化。由Geroch单调性公式，\(m_H(t)\)沿逆平均曲率流单调不减，且\(\lim_{t \to \infty} m_H(t) = M_{\text{ADM}}\)。边界条件给出\(m_H(0) = \sqrt{A/16\pi}\)，故\(M_{\text{ADM}} \geq \sqrt{A/16\pi}\)。Bray利用共形流方法给出了完全严格的证明。\(\blacksquare\)

\hrulefill

\section{第454章：引力波与渐近结构}\label{ux7b2c454ux7ae0ux5f15ux529bux6ce2ux4e0eux6e10ux8fd1ux7ed3ux6784}

引力波是广义相对论最激动人心的预言之一------它是时空曲率的涟漪，以光速传播，由加速运动的质量产生。引力波的数学理论深深植根于Einstein方程在类光无穷远处的渐近行为分析，其核心框架由Bondi、Sachs以及后来的Newman-Penrose等人在1960年代建立。本章将系统阐述引力量子辐射的数学结构：从Bondi-Sachs渐近平坦时空的定义和News函数出发，建立引力波能量（Bondi质量）的严格数学公式，证明其单调递减性（能量正性），进而讨论Christodoulou-Klainerman关于Minkowski时空全局非线性稳定性的里程碑结果------这是Einstein方程作为双曲系统适定性的巅峰证明，以及引力波记忆效应这一引力场永久改变时空几何的可观测现象。这些理论一起构成了引力波天文学（LIGO/Virgo探测2015年首次直接探测到引力波GW150914）的数学基础。

\subsection{153.1 Bondi-Sachs渐近平坦时空}\label{bondi-sachsux6e10ux8fd1ux5e73ux5766ux65f6ux7a7a}

\textbf{定义 153.1}（Bondi-Sachs渐近平坦时空）：时空\((M,g)\)称为\textbf{类光方向渐近平坦}的，若存在标量函数\(\Omega\)（共形因子）和背景度规\(\tilde{g}_{ab}\)使得\(\tilde{g}_{ab} = \Omega^2 g_{ab}\)可光滑延拓到\(\mathscr{I}^+\)（未来类光无穷远），且\(\mathscr{I}^+ = \{\Omega = 0, d\Omega \neq 0\}\)为类光超曲面（拓扑\(S^2 \times \mathbb{R}\)），\(\tilde{g}_{ab}\)在\(\mathscr{I}^+\)上非退化。此外要求真空Einstein方程在\(\mathscr{I}^+\)附近成立（\(\Omega\)的高阶足够多）。

\textbf{定义 153.2}（Bondi坐标系）：在\(\mathscr{I}^+\)附近引入Bondi坐标\((u, r, \theta, \phi)\)（\(u\)为推迟时间，\(r\)为光度距离），度规渐近展开为 \[ds^2 = -e^{2\beta}du^2 - 2e^{2\beta}du dr + r^2\gamma_{AB}(dx^A - U^A du)(dx^B - U^B du),\] 其中\(\gamma_{AB}\)为Sachs形式：\(\gamma_{AB}dx^A dx^B = \frac{1}{2}(e^{2\gamma} + e^{2\delta})d\theta^2 + 2\sinh(\gamma-\delta)\sin\theta d\theta d\phi + \frac{1}{2}(e^{-2\gamma} + e^{-2\delta})\sin^2\theta d\phi^2\)。 Bondi规范条件：\(\det(\gamma_{AB}) = \sin^2\theta\)（二维体元与单位球面相同）。

\textbf{定理 153.1}（Bondi-Sachs渐近展开）：在Bondi规范下，Einstein方程\(R_{\mu\nu} = 0\)在\(\mathscr{I}^+\)附近（\(r \to \infty\)）的渐近解为 \[\beta = -\frac{1}{4r^2}|\sigma^0(u,\theta,\phi)|^2 + O(r^{-3}),\] \[\gamma = \frac{1}{r}\sigma^0_+(u,\theta,\phi) + O(r^{-2}),\quad \delta = \frac{1}{r}\sigma^0_\times(u,\theta,\phi) + O(r^{-2}),\] \[U^A = -\frac{1}{2r^2}D_B\sigma^{0AB} + O(r^{-3}),\] 其中\(\sigma^0_+\)和\(\sigma^0_\times\)分别为引力波\(+\)和\(\times\)偏振的渐近剪切（News函数的积分），\(D_A\)为单位球面上的协变导数。

\textbf{证明}：将Bondi度规代入Einstein方程，按\(1/r\)展开。利用Bondi规范条件简化，\(R_{uu}\)方程给出Bondi质量损失公式，\(R_{AB}\)的无迹部分给出剪切\(\sigma_{AB}\)与度规系数的关系，\(R_{uA}\)和\(R_{ur}\)方程依次确定各项系数。关键步骤：证明所有高阶项的系数均由\(\sigma^0_{AB}(u,\theta,\phi)\)（自由初始数据，即辐射自由度）和初始Bondi质量唯一确定。\(\blacksquare\)

\subsection{153.2 News函数与引力波能量}\label{newsux51fdux6570ux4e0eux5f15ux529bux6ce2ux80fdux91cf}

\textbf{定义 153.3}（News函数与Bondi质量）：\textbf{News函数}（News张量）\(N_{AB}(u,\theta,\phi)\)定义为渐近剪切的\(u\)-导数： \[N_{AB} = \partial_u \sigma^0_{AB},\] 其中\(\sigma^0_{AB}\)为\(\sigma_{AB}\)在\(\mathscr{I}^+\)上\(O(1/r)\)的领头阶系数。News函数描述引力波携带的辐射信息------\(N_{AB}\)在\(\mathscr{I}^+\)上的非零值表示有引力辐射到达无穷远。

Bondi质量\(M_B(u)\)定义为 \[M_B(u) = \frac{1}{8\pi}\int_{S^2} \Psi_2^0(u,\theta,\phi)\, d\Omega,\] 其中\(\Psi_2^0\)为Newman-Penrose标量\(\Psi_2\)在\(\mathscr{I}^+\)上的\(O(r^{-3})\)展开系数。

\textbf{定理 153.2}（Bondi质量损失公式）：沿\(\mathscr{I}^+\)， \[\frac{dM_B}{du} = -\frac{1}{8\pi}\int_{S^2} N_{AB}N^{AB}\, d\Omega \leq 0.\] 即Bondi质量随推迟时间单调递减，减量等于引力波携带的能量通量（辐射的News平方积分）。

\textbf{证明}：由Newman-Penrose形式化，在\(\mathscr{I}^+\)上Bianchi恒等式\(\nabla^\mu C_{\mu\nu\rho\sigma} = 0\)给出 \[\partial_u \Psi_2^0 = \eth^2 \bar{\sigma}^0 + \sigma^0 \partial_u \bar{\sigma}^0 - \Psi_4^0,\] 其中\(\eth\)为球面上的edth算子（自旋权算子）。News函数\(N = \partial_u \sigma^0\)的平方项出现在\(\Psi_2^0\)的演化方程中。将上式在\(S^2\)上积分，\(\eth^2\)项积分为零（球面调和展开的性质），剩余项整理得 \[\frac{dM_B}{du} = -\frac{1}{8\pi}\int_{S^2} |\partial_u\sigma^0|^2 d\Omega = -\frac{1}{8\pi}\int_{S^2} \frac{1}{2}N_{AB}N^{AB}\, d\Omega \leq 0.\] \(\blacksquare\)

\subsection{153.3 Christodoulou-Klainerman非线性稳定性}\label{christodoulou-klainermanux975eux7ebfux6027ux7a33ux5b9aux6027}

\textbf{定理 153.3}（Minkowski时空的全局非线性稳定性，Christodoulou-Klainerman 1993）：Minkowski时空\((\mathbb{R}^4, \eta)\)是Einstein真空方程 \[R_{\mu\nu} = 0\] 的渐近平坦扰动下的全局非线性稳定解。即：存在\(\varepsilon > 0\)使得任何满足适当渐近条件且初始数据与Minkowski初始数据偏差不超过\(\varepsilon\)（在适当Sobolev范数下）的Cauchy数据，其极大Cauchy发展是测地完备的，且度规趋向Minkowski度规（在适当时空范数下）。

\textbf{证明框架}：第一步（零标架与结构方程）：引入双零标架\((e_1, e_2, e_3, e_4) = (L, \underline{L}, e_A)\)（\(A=2,3\)），满足\(g(L,L)=g(\underline{L},\underline{L})=0, g(L,\underline{L})=-2\)。Christoffel联络由联络系数\(\chi, \underline{\chi}, \zeta, \eta, \omega, \underline{\omega}\)编码。曲率由Weyl张量的零分量\(\Psi_n\)和Ricci分量编码。

第二步（Bianchi恒等式的零分解）：真空Einstein方程\(R_{\mu\nu}=0\)将曲率完全约化为Weyl张量。Bianchi恒等式分解为沿\(L\)和\(\underline{L}\)方向的传输方程。引入加权量（带权\(r\)幂次的标度不变量） \[\alpha = r \cdot \Psi_0,\quad \beta = r \cdot \Psi_1,\quad \rho = r \cdot \text{(联络系数)},\] 消除显式的\(r\)衰减因子使得方程变为临界正则性系统。

第三步（Bootstrap论证与能量估计）：利用Bel-Robinson张量构造高阶能量（正定的Weyl曲率二次型） \[Q_{\alpha\beta\gamma\delta} = C_{\alpha\rho\gamma\sigma}C_{\beta\ \ \ \delta}^{\ \rho\sigma} + *C_{\alpha\rho\gamma\sigma}*C_{\beta\ \ \ \delta}^{\ \rho\sigma}\]'' 满足\(\nabla^\alpha Q_{\alpha\beta\gamma\delta} = 0\)（真空情形）。定义通量积分和能量范数： \[\mathcal{E}(t) = \sum_{k=0}^{N} \|\partial^k \Gamma\|_{L^2(\Sigma_t)}^2 + \|\partial^k \mathcal{R}\|_{L^2(\Sigma_t)}^2,\] 其中\(\Gamma\)为联络系数，\(\mathcal{R}\)为曲率张量。证明闭合的能量估计：\(\mathcal{E}(t) \leq C\mathcal{E}(0) + C\int_0^t \mathcal{E}(s)^{3/2}\, ds\)。当\(\mathcal{E}(0) < \varepsilon\)充分小时，Bootstrap论证给出\(\mathcal{E}(t) \leq 2C\varepsilon\)对所有\(t\)一致成立。

第四步（测地完备性）：由联络系数的加权\(L^\infty\)估计（Sobolev嵌入在零框架下的精细版本），证明仿射参数沿任何类时或类光测地线可延拓至无穷，从而是完备的。\(\blacksquare\)

Christodoulou-Klainerman定理是一个巨大的技术成就（全书514页），它证明了一个物理上极其自然的期望（``空的空间是稳定的''），但证明的复杂性暴露了Einstein方程非线性的深度。该工作引入了零框架下的Bianchi恒等式分析、Bel-Robinson能量的非线性估计、以及Bootstrap论证的新形式，对整个数学广义相对论的发展产生了深远影响。

\subsection{153.4 引力波记忆效应}\label{ux5f15ux529bux6ce2ux8bb0ux5fc6ux6548ux5e94}

\textbf{定理 153.4}（引力波记忆效应，Christodoulou 1991, Zel'dovich-Polnarev 1974）：引力波的通过会在自由下落的检验粒子之间留下永久的相对位移变化------即''记忆''。数学上，\(\mathscr{I}^+\)上的渐近剪切\(\sigma^0\)在\(u \to +\infty\)与\(u \to -\infty\)之间的差值\(\Delta\sigma^0\)不为零： \[\Delta\sigma^0_{AB} = \sigma^0_{AB}(+\infty, \theta, \phi) - \sigma^0_{AB}(-\infty, \theta, \phi) \neq 0.\]

\textbf{证明}：设引力波脉冲在有限推迟时间\(u \in [u_1, u_2]\)内通过（即\(N_{AB}=0\)在\(u < u_1\)和\(u > u_2\)）。在\(\mathscr{I}^+\)上，Bondi剪切\(\sigma^0\)满足（由Einstein方程的角分量） \[\partial_u \sigma^0_{AB} = N_{AB},\] 故\(\Delta\sigma^0_{AB} = \int_{-\infty}^{\infty} N_{AB}(u,\theta,\phi)\, du\)------记忆效应的大小等于News函数的总积分。

回忆News函数\(N_{AB}\)与引力波应变的\(\partial_u h_{AB}^{\text{TT}}\)成正比（TT = 无迹横向规范），故记忆效应等价于应变在脉冲前后的净值变化。检测器上自由检验粒子的相对位移变化为 \[\Delta x^i = \frac{1}{2}\Delta h_{ij}^{\text{TT}} x_0^j\] （\(x_0^j\)为初始相对位置），这正是LIGO等探测器原则上可观测的信号。\(\blacksquare\)

Christodoulou进一步证明，记忆效应由两部分组成------``线性记忆''（源于有界引力系统的质量变化）和''非线性记忆''（源于引力波辐射自身的引力势贡献，即Christodoulou记忆，由News函数的平方积分产生）。非线性记忆效应是广义相对论非线性本质的直接可观测后果，其探测将是引力波天文学的一个重要里程碑。 \hrulefill

\hrulefill

\hrulefill

\section{第455章：量子场论数学引论}\label{ux7b2c455ux7ae0ux91cfux5b50ux573aux8bbaux6570ux5b66ux5f15ux8bba}

量子场论（QFT）是纤维丛上变分问题的核心实例，其数学框架包含路径积分的测度论表述和公理化场论（Wightman 公理、Osterwalder-Schrader 重构定理）。量子场论融合了量子力学、狭义相对论和经典场论，是描述基本粒子相互作用的数学语言。本章从数学角度介绍 QFT 的核心结构，包括自由量子场、路径积分、规范场论、微扰论和共形场论。量子场论的严格数学基础仍是一个活跃的研究领域（构造性量子场论），但物理学家发展出的微扰方法在粒子物理标准模型中取得了惊人的精确度。

\subsection{150.1 自由量子场}\label{ux81eaux7531ux91cfux5b50ux573a}

\textbf{定义 150.1}（Klein-Gordon 场）：自由实标量场 \(\phi(x, t)\)（\(x \in \mathbb{R}^3\)，\(t \in \mathbb{R}\)）满足 Klein-Gordon 方程 \[\partial_\mu \partial^\mu \phi + m^2 \phi = 0\] 即 \((\square + m^2)\phi = 0\)，其中 \(\square = \partial_t^2 - \Delta\) 为 d'Alembert 算子，\(m > 0\) 为质量参数。该方程来源于 Lorentz 不变的作用量 \(S[\phi] = \frac{1}{2}\int (\partial_\mu\phi \partial^\mu\phi - m^2\phi^2) d^4x\) 的 Euler-Lagrange 方程。

\textbf{定义 150.2}（正则量子化）：将场 \(\phi(x)\) 和共轭动量 \(\pi(x) = \partial_t \phi(x)\) 提升为 Hilbert 空间 \(\mathcal{H}\) 上的算子值分布，满足等时正则对易关系（CCR）： \[[\phi(x, t), \pi(y, t)] = i\hbar \delta(x - y), \quad [\phi(x,t), \phi(y,t)] = [\pi(x,t), \pi(y,t)] = 0\] 在自然单位制（\(\hbar = c = 1\)）中，对易关系简化为 \([\phi(x,t), \pi(y,t)] = i\delta(x-y)\)。

\textbf{正则量子化的动机}：经典场的正则变量为 \(\phi(x)\) 和 \(\pi(x)\)，Poisson 括号为 \(\{\phi(x), \pi(y)\} = \delta(x-y)\)。正则量子化以 Dirac 规则 \([\cdot, \cdot] = i\hbar \{\cdot, \cdot\}\) 将 Poisson 括号替换为对易子，从而将经典场提升为量子场。

\textbf{定理 150.1}（Klein-Gordon 场的 Fock 空间构造）：自由 Klein-Gordon 场可通过对平面波展开进行量子化。将场写为 \[\phi(x) = \int \frac{d^3p}{(2\pi)^3\sqrt{2E_p}} \left(a(\mathbf{p}) e^{-ip\cdot x} + a^\dagger(\mathbf{p}) e^{ip\cdot x}\right)\] 其中 \(E_p = \sqrt{|\mathbf{p}|^2 + m^2}\)，\(p\cdot x = E_p t - \mathbf{p}\cdot \mathbf{x}\)。产生和湮灭算子满足 \([a(\mathbf{p}), a^\dagger(\mathbf{q})] = (2\pi)^3 \delta(\mathbf{p} - \mathbf{q})\)。真空态 \(|0\rangle\) 满足 \(a(\mathbf{p})|0\rangle = 0\)。Fock 空间 \(\mathcal{F} = \bigoplus_{n=0}^\infty \mathcal{H}_n\)，其中 \(\mathcal{H}_n\) 为 \(n\) 粒子态空间，由 \(a^\dagger(\mathbf{p}_1)\cdots a^\dagger(\mathbf{p}_n)|0\rangle\) 张成。

\textbf{证明}：将 Klein-Gordon 方程 Fourier 变换：\((\partial_t^2 + E_p^2)\tilde{\phi}(\mathbf{p}, t) = 0\)，通解为 \(\tilde{\phi}(\mathbf{p}, t) = \frac{1}{\sqrt{2E_p}}(a(\mathbf{p})e^{-iE_p t} + a^\dagger(-\mathbf{p})e^{iE_p t})\)。由 \(\pi = \partial_t\phi\) 和对易关系可得 \(a, a^\dagger\) 的对易关系。Fock 空间的内积由 \(\langle 0|0\rangle = 1\) 和 CCR 唯一确定。\(\blacksquare\)

\textbf{定义 150.3}（Fock 空间）：自由量子场的状态空间是 Fock 空间 \(\mathcal{F} = \bigoplus_{n=0}^\infty \mathcal{H}^{\otimes_s n}\)（玻色子）或 \(\mathcal{F} = \bigoplus_{n=0}^\infty \mathcal{H}^{\wedge n}\)（费米子）。\(\mathcal{H} = L^2(\mathbb{R}^3, d^3p/2E_p)\) 是单粒子 Hilbert 空间，其中 Lorentz 不变的积分测度 \(d^3p/2E_p\) 来自在壳条件 \(p^2 = m^2\) 的 Lorentz 不变测度。

\textbf{定理 150.2}（Wightman 公理体系）：相对论量子场论的公理化应满足以下 Wightman 公理（1956）：

\textbf{公理 1}（Hilbert 空间与 Poincaré 表示）：存在可分的 Hilbert 空间 \(\mathcal{H}\)，其上作用着 Poincaré 群 \(\mathcal{P}_+^\uparrow\) 的强连续幺正表示 \(U(a, \Lambda)\)。真空态 \(\Omega \in \mathcal{H}\) 为仅有的平移不变态（\(U(a, 1)\Omega = \Omega\)）。

\textbf{公理 2}（谱条件）：能量-动量算符 \(P^\mu\)（平移生成元）的联合谱包含于前向光锥 \(\overline{V}_+ = \{p^\mu : p^0 \geq 0, p^2 \geq 0\}\) 中。真空 \(\Omega\) 为 \(P^\mu = 0\) 的唯一本征态（不计常数倍）。

\textbf{公理 3}（局域性 / 微观因果性）：场 \(\phi\) 为 \(\mathcal{S}(\mathbb{R}^4)\) 上的算子值缓增分布。若检验函数 \(f\) 和 \(g\) 的支集类空分离（即 \((x-y)^2 < 0\) 对所有 \(x \in \operatorname{supp} f\) 和 \(y \in \operatorname{supp} g\) 成立），则对于玻色场 \([\phi(f), \phi(g)] = 0\)，对于费米场 \(\{\phi(f), \phi(g)\} = 0\)。

\textbf{公理 4}（真空循环性）：多项式代数 \(\mathcal{P}(\phi) = \{\phi(f_1)\cdots\phi(f_n)\Omega : n \geq 0, f_i \in \mathcal{S}(\mathbb{R}^4)\}\) 在 \(\mathcal{H}\) 中稠密。

\textbf{Wightman 重构定理}：满足上述公理的场论等价于一组满足特定性质的 Wightman 分布（真空期望值）\(W_n(x_1, \ldots, x_n) = \langle \Omega, \phi(x_1)\cdots\phi(x_n)\Omega \rangle\)。这些分布满足 Lorentz 不变性、谱条件、局域性和正定性条件，且由这些分布可唯一重构 Hilbert 空间和场算符（GNS 构造）。\(\blacksquare\)

\subsection{150.2 路径积分}\label{ux8defux5f84ux79efux5206}

\textbf{定义 150.4}（Feynman 路径积分）：量子力学中的传播子可形式地表为 \[\langle x_f | e^{-i\hat{H}T/\hbar} | x_i \rangle = \int_{x(0)=x_i}^{x(T)=x_f} \mathcal{D}x(t) \, e^{iS[x]/\hbar}\] 其中 \(S[x]\) 是经典作用量，\(\mathcal{D}x(t)\) 是形式上的''路径测度''。测度 \(\mathcal{D}x(t)\) 在数学上并非严格的 Lebesgue 型测度（无穷维平移不变测度不存在），但在物理上可通过 Wiener 测度（Wick 旋转后）或晶格离散化赋予严格意义。

\textbf{定义 150.5}（Wick 旋转与 Euclidean QFT）：做 Wick 旋转 \(t \to -i\tau\)（即 \(x^0 \to -ix^4\)），将 Minkowski 时空变为 Euclidean 时空 \(\mathbb{R}^4\)。路径积分变为 \[Z = \int \mathcal{D}\phi \, e^{-S_E[\phi]/\hbar}\] 其中 \(S_E[\phi] = \int d^4x \left( \frac{1}{2}(\partial_\mu\phi)^2 + \frac{1}{2}m^2\phi^2 + V(\phi) \right)\) 是 Euclidean 作用量。此时 Boltzmann 因子 \(e^{-S_E}\) 为正，路径积分在数学上可与概率测度相联系（Gaussian 测度加上有界扰动）。

\textbf{Wick 旋转的动机}：Minkowski 路径积分的振荡相因子 \(e^{iS/\hbar}\) 使得积分高度振荡，难以定义。Wick 旋转将其转化为 \(e^{-S_E/\hbar}\)（指数衰减），从而可用概率论方法（如 Gaussian 测度）严格处理。Osterwalder-Schrader 定理保证了 Euclid 理论可唯一解析延拓回 Minkowski 理论。

\textbf{定理 150.3}（Osterwalder-Schrader 重构定理，1973-1975）：满足 Osterwalder-Schrader 公理的 Euclidean 关联函数 \(S_n(x_1,\ldots,x_n)\) 可以通过解析延拓唯一地重构满足 Wightman 公理的相对论量子场论。OS 公理包括： (1) Euclidean 协变性：\(S_n\) 在 Euclidean 群 \(SO(4)\) 作用下协变； (2) 反射正性（OS 正性）：\(\sum_{m,n} \int S_{m+n}(\theta x_1,\ldots,\theta x_m, y_1,\ldots,y_n) \overline{f_m(x_1,\ldots,x_m)} f_n(y_1,\ldots,y_n) dx dy \geq 0\)，其中 \(\theta(x^4, \mathbf{x}) = (-x^4, \mathbf{x})\) 为时间反射； (3) 对称性：\(S_n\) 在自变量置换下对称； (4) 聚簇性：\(S_{n+m}(x_1,\ldots,x_n, y_1+a,\ldots,y_m+a) \to S_n(x_1,\ldots,x_n)S_m(y_1,\ldots,y_m)\) 当 \(|a| \to \infty\)。

\textbf{证明}：OS 重构分两步。（1）Hilbert 空间构造：在 Euclidean Schwartz 函数空间上，利用反射正性定义半正定内积 \(\langle f, g \rangle = \int S_{n+m}(\theta x_1,\ldots,\theta x_n, y_1,\ldots,y_m) \overline{f(x_1,\ldots,x_n)} g(y_1,\ldots,y_m) dx dy\)（\(\theta\) 为时间反射）。商去零空间后完备化得 Hilbert 空间 \(\mathcal{H}\)，真空态 \(\Omega\) 为常数函数 \(1\) 的等价类。（2）解析延拓：Euclidean 关联函数 \(S_n\) 在 Minkowski 空间区域上具有解析性（由 OS 公理保证），解析延拓给出唯一的 Minkowski 关联函数 \(W_n\)。\(W_n\) 满足 Wightman 公理的所有条件：谱条件来自 OS 正性（通过 Laplace 变换），局域性来自 Euclidean 对称性在解析延拓下的保持和 Hall-Wightman 定理。\(\blacksquare\)

\subsection{150.3 规范场论}\label{ux89c4ux8303ux573aux8bba}

\textbf{定义 150.6}（Yang-Mills 理论）：设 \(G\) 为紧致 Lie 群（如 \(SU(3) \times SU(2) \times U(1)\) 为标准模型规范群），\(\mathfrak{g}\) 为其 Lie 代数。Yang-Mills 联络 \(A_\mu\) 是 \(\mathfrak{g}\)-值的 1-形式，场强 \(F_{\mu\nu} = \partial_\mu A_\nu - \partial_\nu A_\mu + [A_\mu, A_\nu]\) 为 \(\mathfrak{g}\)-值的 2-形式。Yang-Mills 作用量为 \[S_{YM}[A] = -\frac{1}{4} \int d^4x \, \operatorname{Tr}(F_{\mu\nu}F^{\mu\nu})\] 规范变换 \(A_\mu \to g A_\mu g^{-1} + g \partial_\mu g^{-1}\)（\(g(x) \in G\)）下 \(F_{\mu\nu} \to g F_{\mu\nu} g^{-1}\)，故作用量规范不变。

\textbf{Yang-Mills 理论的几何意义}：Yang-Mills 理论是主 \(G\)-丛 \(P \to \mathbb{R}^4\) 上的联络理论。\(A_\mu\) 为局部联络形式，\(F_{\mu\nu}\) 为曲率。Yang-Mills 作用量是曲率 \(L^2\) 范数平方的积分，其 Euler-Lagrange 方程为 \(D^\mu F_{\mu\nu} = 0\)（\(D_\mu\) 为协变导数），即 Yang-Mills 方程------这是 Maxwell 方程 \((\partial^\mu \partial_\mu A_\nu - \partial^\mu \partial_\nu A_\mu = 0)\) 的非 Abel 推广。

\textbf{定理 150.4}（Yang-Mills 存在性与质量间隙 / Clay 问题）：证明对任意紧致单纯规范群 \(G\)，\(\mathbb{R}^4\) 上的量子 Yang-Mills 理论存在，且具有质量间隙 \(\Delta > 0\)（即 Hamilton 算子的谱在 \(0\) 和 \(\Delta\) 之间有空隙）。这是七大千禧年问题之一，尚未解决。

\textbf{证明方向}（已知进展）：此问题当前未获证明。主要探索方向包括：(1) 格点规范理论途径------在有限间距晶格 \(\varepsilon \mathbb{Z}^4\) 上定义 Yang-Mills 理论（Wilson 作用量），取连续极限 \(\varepsilon \to 0\)，证明关联函数存在且满足 Osterwalder-Schrader 公理。质量间隙的格点证据来自数值模拟（Monte Carlo 方法），在强耦合展开下可证明质量间隙存在，但连续极限下的严格证明仍缺。(2) 构造性量子场论途径------使用公理化量子场论框架（Wightman 公理）从第一原理出发构造 Yang-Mills 理论。三维 Yang-Mills 理论的存在性已在 Balaban 的工作中取得进展，但四维情形仍为开放问题。(3) 可积性与对偶性途径------利用 Seiberg-Witten 理论和电磁对偶性研究 \(N=2\) 超对称 Yang-Mills 理论的质量谱，但一般非超对称情形仍困难。\(\blacksquare\)

\textbf{定义 150.7}（BRST 量子化 / Becchi-Rouet-Stora-Tyutin）：规范理论的量子化需要固定规范以消除冗余自由度。Faddeev-Popov 方法引入鬼场（ghost fields）\(c, \bar{c}\)（反对易标量场），使有效作用量包括规范固定项和鬼场项。BRST 算子 \(Q\)（\(Q^2 = 0\)）生成 BRST 对称性，物理态空间为 \(Q\)-上同调 \(\ker Q / \operatorname{im} Q\)。BRST 协变性保证了幺正性和规范不变性的相容性：\(S\)-矩阵的幺正性等价于 \(Q\)-上同调的非平凡性，而物理可观测量为 \(Q\)-闭但非 \(Q\)-恰当的量。

\textbf{定理 150.5}（Higgs 机制，1964）：在规范场与标量场耦合时，若标量场的真空期望值非零，则规范对称性自发破缺，规范玻色子获得质量，而物理 Higgs 玻色子保留。这是电弱统一理论（Glashow-Weinberg-Salam）的核心机制。

\textbf{证明}：考虑 \(SU(2)\times U(1)\) 规范理论与复标量二重态 \(\Phi\) 耦合。Lagrange 密度为 \(\mathcal{L} = |D_\mu\Phi|^2 - V(\Phi) - \frac{1}{4}F_{\mu\nu}^a F^{a\mu\nu}\)，其中 \(D_\mu = \partial_\mu + igW_\mu^a \tau^a/2 + ig'B_\mu/2\)，\(V(\Phi) = \mu^2|\Phi|^2 + \lambda|\Phi|^4\)（\(\mu^2 < 0, \lambda > 0\)）。势能最小值在 \(|\Phi| = v = \sqrt{-\mu^2/\lambda} \neq 0\)。选取幺正规范 \(\Phi = (0, v+h)^T/\sqrt{2}\)，其中 \(h\) 为 Higgs 场。展开动能项 \(|D_\mu\Phi|^2\)： \[|D_\mu\Phi|^2 = \frac{1}{2}(\partial_\mu h)^2 + \frac{g^2 v^2}{4}W_\mu^+W^{-\mu} + \frac{(g^2+g'^2)v^2}{8}Z_\mu Z^\mu + 0 \cdot A_\mu A^\mu + \text{相互作用项}\] 光子 \(A_\mu = \cos\theta_W B_\mu + \sin\theta_W W_\mu^3\)（\(\theta_W\) 为 Weinberg 角）保持无质量，\(W^\pm\) 获得质量 \(m_W = gv/2\)，\(Z\) 获得质量 \(m_Z = \sqrt{g^2+g'^2}v/2\)，Higgs 玻色子 \(h\) 质量 \(m_h = \sqrt{2\lambda}v\)。Goldstone 定理的''避免''：规范对称性自发破缺时，原本应出现的 Goldstone 玻色子被规范玻色子''吃掉''成为其纵向分量，从而获得质量。\(\blacksquare\)

\subsection{150.4 微扰论与Feynman图}\label{ux5faeux6270ux8bbaux4e0efeynmanux56fe}

\textbf{定义 150.8}（生成泛函与关联函数）：量子场论的核心计算对象是 \(n\) 点关联函数（Green 函数）。在 Euclidean 理论中，配分函数（生成泛函）为 \[Z[J] = \int \mathcal{D}\phi \, \exp\left(-S_E[\phi] + \int J(x)\phi(x) d^4x\right)\] \(n\) 点关联函数为 \(\langle \phi(x_1)\cdots\phi(x_n) \rangle = \frac{1}{Z[0]} \frac{\delta^n Z[J]}{\delta J(x_1)\cdots\delta J(x_n)}\big|_{J=0}\)。连通 Green 函数由 \(W[J] = \log Z[J]\) 生成，单粒子不可约（1PI）Green 函数由有效作用量 \(\Gamma[\phi_{\text{cl}}] = W[J] - \int J\phi_{\text{cl}}\)（\(\phi_{\text{cl}} = \delta W/\delta J\)）的 Legendre 变换生成。

\textbf{定理 150.6}（Wick 定理）：自由场论的 \(n\) 点关联函数可分解为两点函数（Feynman 传播子）的乘积之和，对应所有可能的完全配对（Wick 缩并）。在泛函语言中，这等价于 Gaussian 积分的矩公式：\(\langle \phi(x_1)\cdots\phi(x_{2n}) \rangle_0 = \sum_{\text{配对}} \prod_{\{i,j\}} \Delta_F(x_i - x_j)\)，其中 \(\Delta_F\) 为 Feynman 传播子。

\textbf{证明}：自由作用量 \(S_0 = \frac{1}{2}\int \phi(-\Delta + m^2)\phi\) 为 Gaussian 型。由 Gaussian 积分公式 \[\int \mathcal{D}\phi \, e^{-\frac{1}{2}\langle\phi, A\phi\rangle + \langle J, \phi\rangle} = (\det A)^{-1/2} e^{\frac{1}{2}\langle J, A^{-1}J\rangle}\] 对 \(J\) 逐次求导：\(n\) 为奇数时关联函数为零，\(n\) 为偶数时 \(\langle \phi(x_1)\cdots\phi(x_{2n}) \rangle_0 = \frac{\delta^{2n}}{\delta J(x_1)\cdots\delta J(x_{2n})} e^{\frac{1}{2}\langle J, A^{-1}J\rangle}\big|_{J=0}\)。由 Leibniz 法则，每对 \(J\) 的导数产生一个 \(A^{-1}\) 的矩阵元，即传播子 \(\Delta_F(x_i - x_j) = A^{-1}(x_i, x_j) = (-\Delta + m^2)^{-1}(x_i - x_j)\)。所有导数配对的总和即 Wick 缩并。\(\blacksquare\)

\textbf{定义 150.9}（Feynman 规则与重整化）：相互作用 Lagrange 量 \(g\mathcal{L}_{\text{int}}\) 给出顶点因子，传播子给出内线因子，外线对应外动量。对每个 Feynman 图，计算其对称因子（图的对称群阶数的倒数）并在圈动量上积分。圈积分通常发散（紫外发散），需通过\textbf{重整化}（renormalization）引入抵消项以吸收无穷大。重整化理论是 QFT 的数学核心：可重整化理论（如 Yang-Mills 理论、\(\phi^4\) 理论在四维）具有有限个发散类型，可通过重新定义有限个参数（质量、耦合常数、场归一化）消除所有无穷大。BPHZ 定理（Bogoliubov-Parasiuk-Hepp-Zimmermann）给出了微扰论重整化的严格数学框架，而 Wilson 的有效场论观点则将重整化理解为尺度变换下理论在耦合常数空间中的流向。

\textbf{定理 150.7}（Dyson 级数与 Feynman 图展开）：相互作用理论的 \(n\) 点关联函数可展开为 \[\langle \phi(x_1)\cdots\phi(x_n) \rangle = \frac{\langle \phi(x_1)\cdots\phi(x_n) e^{-S_{\text{int}}[\phi]} \rangle_0}{\langle e^{-S_{\text{int}}[\phi]} \rangle_0}\] 其中 \(\langle \cdot \rangle_0\) 为自由场期望值。将 \(e^{-S_{\text{int}}}\) 作 Taylor 展开（Dyson 级数），逐项应用 Wick 定理即得 Feynman 图展开。\(S_{\text{int}}\) 的 \(k\) 阶项对应 \(k\) 个顶点的图，每个顶点贡献 \(g\) 乘以顶点因子。Dyson 级数在数学上仅为渐近级数（Dyson 论证：若级数收敛半径为正，则 \(g < 0\) 时理论也应有定义，但此时势能无下界，真空不稳定），但在物理上低阶贡献已给出极高精度的预言（如电子反常磁矩 \(g-2\) 的计算与实验一致到 \(10^{-12}\) 量级）。

\subsection{150.5 共形场论与顶点算子代数}\label{ux5171ux5f62ux573aux8bbaux4e0eux9876ux70b9ux7b97ux5b50ux4ee3ux6570}

\textbf{定义 150.10}（共形场论 / CFT）：二维共形场论是共形群 \(SO(n,2)\) 在 \(n=2\) 时无穷维扩充（局部共形变换为全纯映射）的量子场论。二维 CFT 的核心结构是\textbf{能动张量} \(T(z), \overline{T}(\bar{z})\) 的算子乘积展开（OPE）： \[T(z)T(w) = \frac{c/2}{(z-w)^4} + \frac{2T(w)}{(z-w)^2} + \frac{\partial T(w)}{z-w} + \text{正则项}\] 其中 \(c\) 是\textbf{中心荷}（central charge），是分类 CFT 的基本参数。\(T(z)\) 的模式展开 \(T(z) = \sum_{n \in \mathbb{Z}} L_n z^{-n-2}\) 给出 Virasoro 代数： \[[L_m, L_n] = (m-n)L_{m+n} + \frac{c}{12}m(m^2-1)\delta_{m+n,0}\]

\textbf{CFT 在物理中的重要性}：二维 CFT 描述了弦的世界面理论、统计力学在临界点的普适类、以及 \(AdS_3/CFT_2\) 对偶中的边界理论。CFT 的对称性（Virasoro 代数）远大于通常的 Lorentz 对称性，使得二维 CFT 在一定程度上可精确求解。

\textbf{定理 150.8}（BPZ 定理 / 极小模型分类，Belavin-Polyakov-Zamolodchikov, 1984）：中心荷 \(c < 1\) 的幺正二维 CFT（极小模型）仅当 \[c = 1 - \frac{6}{m(m+1)}, \quad m = 3, 4, 5, \ldots\] 时存在。其主场的共形权为 \[h_{r,s} = \frac{((m+1)r - ms)^2 - 1}{4m(m+1)}, \quad 1 \leq r \leq m-1, \; 1 \leq s \leq m\] 这些极小模型精确对应统计力学中临界现象的普适类：\(m=3\)（\(c=1/2\)）对应 Ising 模型，\(m=4\)（\(c=7/10\)）对应三临界 Ising 模型，\(m=5\)（\(c=4/5\)）对应 3-态 Potts 模型，\(m \to \infty\)（\(c \to 1\)）对应 Gaussian 自由场。

\textbf{证明}：由 Kac 行列式公式，Verma 模 \(V(c,h)\)（由最高权态 \(|h\rangle\) 生成的 Virasoro 代数表示）在级 \(N\) 处的内积矩阵（Gram 矩阵）行列式为 \[\det M_N = \prod_{r,s \geq 1, rs \leq N} (h - h_{r,s}(c))^{p(N-rs)}\] 其中 \(h_{r,s}(c)\) 由 \(c\) 参数化：\(c = 1 - 6(\alpha_+ + \alpha_-)^2\)，\(h_{r,s} = \frac{1}{4}(r^2-1)\alpha_+^2 + \frac{1}{4}(s^2-1)\alpha_-^2 - \frac{1}{2}(rs-1)\)，\(\alpha_\pm = \frac{1}{\sqrt{24}}(\sqrt{1-c} \pm \sqrt{25-c})\)。幺正性要求所有非零态的内积为正，等价于 Kac 行列式在各子空间非负。当 \(c < 1\) 时，仅当 \(c = 1 - 6/m(m+1)\) 时行列式非负（对所有 \(N\)），由此推出以上分类。BPZ 的原论文通过分析 Virasoro 代数的表示论和融合规则（fusion rules）完成了分类。\(\blacksquare\)

\textbf{定义 150.11}（顶点算子代数 / VOA）：顶点算子代数（VOA）是二维 CFT 的数学公理化，由 Borcherds（1986）和 Frenkel-Lepowsky-Meurman（1988）建立。VOA 是一个 \(\mathbb{Z}\)-分次向量空间 \(V = \bigoplus_{n \in \mathbb{Z}} V_n\)（\(\dim V_n < \infty\)，\(V_n = 0\) 对充分小的 \(n\)），配备真空态 \(|0\rangle \in V_0\)、平移算子 \(T \in \operatorname{End}(V)\)（\(T|0\rangle = 0\)）和顶点算子映射 \(Y(\cdot, z): V \to \operatorname{End}(V)[[z, z^{-1}]]\)（对 \(a \in V\)，\(Y(a,z) = \sum_{n \in \mathbb{Z}} a_{(n)} z^{-n-1}\)），满足以下公理： (1) \textbf{真空性}：\(Y(|0\rangle, z) = \operatorname{id}_V\)，且 \(Y(a, z)|0\rangle|_{z=0} = a\)； (2) \textbf{平移性}：\([T, Y(a, z)] = \partial_z Y(a, z)\)； (3) \textbf{局域性 / Jacobi 恒等式}：对任意 \(a, b \in V\)，存在 \(N \in \mathbb{Z}_{\geq 0}\) 使 \((z-w)^N [Y(a, z), Y(b, w)] = 0\)。

VOA 语言将量子场论中的 OPE 严格化为代数结构。其与怪物群（Monster group）的联系------Moonshine 猜想（Borcherds 1992 证明）------是 20 世纪数学最深刻的发现之一：月光模 \(V^\natural\) 的 \(q\)-展开系数（\(j\)-不变量）与 Monster 单群的不可约表示维数之间存在神奇对应，\(V^\natural\) 的自同构群恰为 Monster 单群，这一发现连接了代数（有限单群）、数论（模形式）和数学物理（共形场论）。

\textbf{定理 150.9}（Borcherds 月光定理，1992）：月光模 \(V^\natural = \bigoplus_{n=-1}^\infty V^\natural_n\) 是一个顶点算子代数，其自同构群为 Monster 单群 \(\mathbb{M}\)。\(V^\natural\) 的分次迹（McKay-Thompson 级数）为模函数，满足 \(T_g(\tau) = q^{-1}\sum_{n=-1}^\infty \operatorname{Tr}(g|_{V^\natural_n}) q^n\)（\(q = e^{2\pi i\tau}\)）为 \(\Gamma_0(|g|)\) 的 Hauptmodul。这证明了 Conway-Norton 的 Moonshine 猜想（1979）。

\textbf{证明概要}：Borcherds 的证明分为三个主要步骤：首先，构造月光模 \(V^\natural\) 为 Leech 格点 \(\Lambda\) 的格点顶点算子代数的 \(\mathbb{Z}_2\)-扭的 orbifold 构造（Frenkel-Lepowsky-Meurman 1988）；其次，利用 Borcherds 的广义 Kac-Moody 代数（Borcherds 代数）理论，证明 \(V^\natural\) 的特征标满足特定的函数方程和积性公式；最后，通过比较 \(V^\natural\) 的 twisted 特征标与已知的 McKay-Thompson 级数，验证所有 \(T_g\) 为 Hauptmodul。核心工具是 Borcherds 的 denominator identity 和 Goddard-Thorn 的无鬼定理（no-ghost theorem）。\(\blacksquare\)

\hrulefill

\emph{量子场论从 Klein-Gordon 场的正则量子化出发，经 Fock 空间构造和 Wightman 公理体系，建立了相对论量子理论和粒子物理的数学基础。路径积分和 Wick 旋转提供了从 Euclidean 关联函数构造 Minkowski 理论的严格途径（OS 重构定理）。Yang-Mills 规范场论以纤维丛上的联络理论为几何框架，其量子化（BRST 方法）和 Higgs 机制构成了标准模型的数学核心。微扰论以 Feynman 图展开和 Wick 定理为计算工具，重整化（BPHZ 定理）消除了紫外发散。二维共形场论以 Virasoro 代数为代数结构，BPZ 极小模型分类精确对应了统计力学的临界普适类。顶点算子代数（VOA）为 CFT 提供了严格的代数公理化，其与 Monster 单群的月光对应是数学物理交叉的巅峰成就。}

\emph{卷三十：数学物理。}

\newpage

\chapter{13-数学物理/数学物理/05-共形与可积.md}\label{ux6570ux5b66ux7269ux7406ux6570ux5b66ux7269ux740605-ux5171ux5f62ux4e0eux53efux79ef.md}

\section{第456章：共形场论}\label{ux7b2c456ux7ae0ux5171ux5f62ux573aux8bba}

共形场论（Conformal Field Theory, CFT）研究具有共形对称性的量子场论。在二维时空中，共形群为无限维，产生了 Virasoro 代数这一核心代数结构，使得二维CFT可以严格求解。

\subsection{370.1 共形变换与共形群}\label{ux5171ux5f62ux53d8ux6362ux4e0eux5171ux5f62ux7fa4}

\textbf{定义 370.1}（共形变换）：设 \((M,g)\) 为伪Riemann流形。微分同胚 \(\varphi: M\to M\) 称为\textbf{共形变换}，若存在正函数 \(\Omega(x)>0\) 使得

\[\varphi^* g = \Omega^2(x)\, g\]

即度规在变换下仅缩放而不改变角度。等价地，在局部坐标下 \(g_{\mu\nu}'(x')=\Omega^2(x)g_{\mu\nu}(x)\)。

\textbf{定义 370.2}（Weyl变换 vs 共形变换）：\(g_{\mu\nu}(x)\to\Omega^2(x)g_{\mu\nu}(x)\) 且坐标不变的变换称为\textbf{Weyl变换}（度规的局部缩放），而共形变换是坐标的微分同胚满足 \(g_{\mu\nu}'(x')=\Omega^2(x)g_{\mu\nu}(x)\)。二者不等价：共形变换是时空对称性，而Weyl变换是度规的场变换。

\textbf{定理 370.1}（共形群的Liouville定理，\(d>2\)）：设 \(d=p+q>2\)。则 \(\mathbb{R}^{p,q}\) 中所有共形变换构成的共形群是有限维Lie群，同构于 \(\mathrm{SO}(p+1,q+1)\)，维数为 \((d+1)(d+2)/2\)。

\textbf{证明}：对 \(d>2\)，共形Killing方程 \(\partial_\mu\varepsilon_\nu+\partial_\nu\varepsilon_\mu=\frac{2}{d}\eta_{\mu\nu}\partial_\rho\varepsilon^\rho\) 的通解为 \(\varepsilon^\mu(x)=a^\mu+\omega^{\mu\nu}x_\nu+\lambda x^\mu+b^\mu x^2-2x^\mu(b\cdot x)\)，对应平移、旋转、缩放和特殊共形变换，共 \((d+1)(d+2)/2\) 个独立参数。\(\blacksquare\)

\subsection{370.2 Virasoro代数与顶点算子代数}\label{virasoroux4ee3ux6570ux4e0eux9876ux70b9ux7b97ux5b50ux4ee3ux6570}

\textbf{定义 370.3}（Witt代数）：在二维Minkowski时空（光锥坐标 \(z,\bar{z}\)）中，所有局部共形变换由亚纯函数 \(z\mapsto f(z)\) 生成。生成元 \(l_n=-z^{n+1}\partial_z\)（\(n\in\mathbb{Z}\)）满足\textbf{Witt代数}：

\[[l_m, l_n] = (m-n)\,l_{m+n}\]

这是 \(\mathrm{diff}(S^1)\) 的复化Lie代数。

\textbf{定理 370.2}（Virasoro代数，Virasoro 1970）：Witt代数的非平凡中心扩张（由 \(H^2(\mathrm{Witt},\mathbb{C})\cong\mathbb{C}\) 分类）给出\textbf{Virasoro代数}：

\[[L_m, L_n] = (m-n)L_{m+n} + \frac{c}{12}m(m^2-1)\delta_{m+n,0}\]

其中 \(c\in\mathbb{C}\) 称为\textbf{中心荷}。此为中心扩张的唯一形式（等价于非零2-余循环 \(\omega(l_m,l_n)=\frac{1}{12}m(m^2-1)\delta_{m+n,0}\)）。

\textbf{证明}：设中心扩张 \([L_m,L_n]=(m-n)L_{m+n}+A(m)\delta_{m+n,0}\)，由Jacobi恒等式得 \(A(m)=-A(-m)\) 且递推关系 \(A(m)=\frac{1}{12}m(m^2-1)A(2)\)。定义 \(c=12A(2)\) 即得。\(\blacksquare\)

Virasoro代数在数学中的重要性远超其在物理中的起源。它是二维共形场论中能量-动量张量的Laurent模式所满足的代数，也是弦论中世界面坐标的约束代数。Virasoro代数的表示论是二维CFT分类的核心------所有酉最高权表示由中心荷 \(c\) 和共形权 \(h\) 两个参数刻画。中心荷 \(c\) 的物理意义是共形反常（conformal anomaly）：在弯曲背景中，经典共形对称性在量子化后一般会被破坏，而 \(c\) 度量了这种破坏的程度。这一现象与二维引力中的Liouville理论密切相关------Polyakov的量子Liouville理论将 \(c\) 与时空维数联系起来。

\textbf{定义 370.4}（顶点算子代数 / Vertex Operator Algebra, VOA）：顶点算子代数（Borcherds 1986, Frenkel-Lepowsky-Meurman 1988）是共形场论的公理化代数结构。一个\textbf{顶点算子代数} \(V\) 是一个 \(\mathbb{Z}\)-分次向量空间 \(V = \bigoplus_{n\in\mathbb{Z}} V_n\)，配以状态-场对应 \(Y: V \to \operatorname{End}(V)[[z, z^{-1}]]\)（将每个状态 \(a \in V\) 映为一个顶点算子 \(Y(a, z) = \sum_{n\in\mathbb{Z}} a_{(n)} z^{-n-1}\)）、真空向量 \(|0\rangle \in V_0\) 和 Virasoro 向量 \(\omega \in V_2\)，满足一系列公理（包括局域性、平移协变性、真空性质和 Virasoro 代数的实现）。VOA 的公理化将二维 CFT 的算子乘积展开（OPE）和手征代数（chiral algebra）结构提炼为纯代数语言，为月光猜想（Monstrous Moonshine）的证明提供了数学框架------Frenkel-Lepowsky-Meurman 构造了月光模 \(V^\natural\)（一个 VOA），其自同构群是 Monster 单群，其分次迹给出 \(j\)-函数。

\subsection{370.3 中心荷、Verma模与Kac行列式}\label{ux4e2dux5fc3ux8377vermaux6a21ux4e0ekacux884cux5217ux5f0f}

\textbf{定义 370.5}（Verma模）：设 \(|\Delta\rangle\) 为最高权态，满足 \(L_0|\Delta\rangle=\Delta|\Delta\rangle\) 和 \(L_{n>0}|\Delta\rangle=0\)。\textbf{Verma模} \(M(c,\Delta)\) 是由所有 \(L_{-n_k}\cdots L_{-n_1}|\Delta\rangle\)（\(n_i>0\)）张成的Virasoro代数表示空间。它具有自然的 \(L_0\) 分级结构。

\textbf{定理 370.3}（Kac行列式，Kac 1978）：在 Verma模 \(M(c,\Delta)\) 的 \(L_0\) 本征值为 \(\Delta+N\) 的子空间 \(V_N\) 中，所有态的内积矩阵的\textbf{Kac行列式}为

\[\det M_N(c,\Delta) = K_N \prod_{r,s\geq 1,\, rs\leq N} (\Delta - \Delta_{r,s}(c))^{P(N-rs)}\]

其中 \(P(k)\) 是划分函数，\(\Delta_{r,s}(c)=\frac{1-r^2}{4}b^2+\frac{1-s^2}{4}b^{-2}+\frac{1-rs}{2}\)，且 \(c=1-6(b-b^{-1})^2\)。\(\det M_N=0\) 当且仅当存在退化向量（null vector），此时 Verma模可约。

\textbf{证明}：在 \(V_N\) 中取基 \(\{L_{-\lambda}|\Delta\rangle: \lambda\vdash N\}\)。Shapovalov 形式 \(\langle\Delta|L_{\lambda}L_{-\mu}|\Delta\rangle\) 定义了内积矩阵 \(M_N\)。\(\det M_N\) 作为 \(\Delta\) 的多项式次数为 \(\dim V_N = \sum_{rs\leq N}P(N-rs)\)。当 \(\Delta = \Delta_{r,s}\) 时存在退化向量，即 \(M_N\) 有零向量，故 \(\det M_N=0\)。零点重数 \(P(N-rs)\) 由退化向量的级数确定。系数 \(K_N\) 由 \(\Delta\to\infty\) 渐近行为确定。\(\blacksquare\)

\textbf{定理 370.4}（极小模型，Belavin-Polyakov-Zamolodchikov 1984）：当中心荷取离散值

\[c = 1 - \frac{6(p-q)^2}{pq}\]

其中 \(p,q\) 为互质正整数（\(p,q\geq 2\)），Virasoro代数的酉表示仅包含有限多个最高权表示。这些\textbf{极小模型}对应二维统计力学的临界点，其共形权为 \(\Delta_{r,s}=\frac{(pr-qs)^2-(p-q)^2}{4pq}\)（\(1\leq r<q,\;1\leq s<p\)）。

\textbf{证明}：酉性要求 Kac 行列式在各级非负。分析得酉表示仅存在于：(i) \(c\geq 1\) 且 \(\Delta\geq 0\)；(ii) 离散序列 \(c = 1-6/(m(m+1))\)（\(m\geq 3\)）且 \(\Delta = \Delta_{r,s}\)（\(1\leq r<m, 1\leq s\leq r\)）。设 \(p=m+1\), \(q=m\)（或互换），则 \(c = 1-6/(m(m+1)) = 1-6(p-q)^2/(pq)\)。\(\Delta_{r,s}\) 即共形权公式。这些模型的关联函数满足 BPZ 微分方程，可严格求解。极小模型是二维 CFT 分类的基石。\(\blacksquare\)

极小模型与统计力学中可解格点模型的临界行为有精确对应。\(m=3\)（\(c=1/2\)）的 Ising 模型、\(m=4\)（\(c=7/10\)）的三临界 Ising 模型、\(m=5\)（\(c=4/5\)）的三态 Potts 模型，都是极小模型的物理实现。这一对应是共形场论在统计力学中最重要的应用之一，它使得临界指数等物理量可以通过 Virasoro 代数的表示论精确计算。

\subsection{370.4 WZW模型与Kac-Moody代数}\label{wzwux6a21ux578bux4e0ekac-moodyux4ee3ux6570}

\textbf{定义 370.6}（Wess-Zumino-Witten模型）：WZW模型是二维共形场论中最重要的非平凡可解模型之一，其作用量定义在群值场 \(g: \Sigma \to G\)（\(G\) 为紧Lie群，\(\Sigma\) 为二维Riemann面）上：

\[S_{\text{WZW}}[g] = \frac{k}{4\pi} \int_\Sigma \operatorname{Tr}(g^{-1}\partial g \cdot g^{-1}\bar{\partial} g) + \frac{k}{12\pi} \int_B \operatorname{Tr}((g^{-1}dg)^{\wedge 3})\]

其中 \(k \in \mathbb{Z}\) 为\textbf{水平}（level），\(B\) 为以 \(\Sigma\) 为边界的三维流形。WZW项（第二项）是Wess-Zumino拓扑项，其量子化条件要求 \(k\) 为整数。WZW模型的对称性由仿射Kac-Moody代数 \(\hat{\mathfrak{g}}_k\) 生成，其手征流 \(J^a(z)\) 满足算子乘积展开

\[J^a(z) J^b(w) \sim \frac{k \delta^{ab}}{(z-w)^2} + \frac{f^{ab}_c J^c(w)}{z-w}\]

其中 \(f^{ab}_c\) 为Lie代数 \(\mathfrak{g}\) 的结构常数。WZW模型的可解性源于其手征代数的无穷维对称性------关联函数满足 Knizhnik-Zamolodchikov（KZ）方程，这是一阶矩阵微分方程，由零向量条件导出。对 \(G=SU(2)\)，WZW模型在水平 \(k\) 下的中心荷为 \(c = 3k/(k+2)\)，共形权为 \(\Delta_j = j(j+1)/(k+2)\)（\(j=0,1/2,\ldots,k/2\)）。

\hrulefill

\section{第457章：可积系统与可解模型}\label{ux7b2c457ux7ae0ux53efux79efux7cfbux7edfux4e0eux53efux89e3ux6a21ux578b}

可积系统是具有''足够多''独立守恒量的无限维动力系统，其核心结构为Lax对形式化，使得非线性演化方程可以表示为两个线性算子之间的可积性条件。

\subsection{371.1 Lax对形式化与Hirota方法}\label{laxux5bf9ux5f62ux5f0fux5316ux4e0ehirotaux65b9ux6cd5}

\textbf{定义 371.1}（Lax对与零曲率表示）：设 \(L(t)\) 和 \(M(t)\) 是依赖于位势 \(u(x,t)\) 的线性算子。若演化方程等价于\textbf{Lax方程}

\[\frac{\partial L}{\partial t} = [M, L] \equiv ML - LM\]

则称 \((L,M)\) 为此系统的\textbf{Lax对}。等价地，定义零曲率联络 \(U,V\) 满足 \(U_t-V_x+[U,V]=0\)，称为\textbf{零曲率表示}。

\textbf{定理 371.1}（Lax方程保证守恒量）：若 \(L\) 满足 \(\partial_t L=[M,L]\)，则对任意 \(n\in\mathbb{N}\)，量 \(I_n=\operatorname{tr}(L^n)\)（或 \(\int\lambda^n\,d\mu(\lambda)\) 其中 \(\mu\) 为 \(L\) 的谱测度）是守恒量：\(\partial_t I_n=0\)。

\textbf{证明}：\(\partial_t(L^n)=\sum_{k=0}^{n-1}L^k[M,L]L^{n-1-k}=[M,L^n]\)，取迹 \(\partial_t\operatorname{tr}(L^n)=\operatorname{tr}([M,L^n])=0\)。\(\blacksquare\)

\textbf{定义 371.2}（Hirota双线性方法）：Hirota双线性方法（Hirota 1971）是构造可积系统孤子解的最系统方法。通过引入因变量变换（如对KdV方程 \(u = 2(\log \tau)_{xx}\)），将非线性演化方程转化为双线性方程（Hirota方程）。Hirota \(D\)-算子定义为

\[D_x^m D_t^n f \cdot g = (\partial_x - \partial_{x'})^m (\partial_t - \partial_{t'})^n f(x,t) g(x',t') \big|_{x'=x, t'=t}\]

KdV方程 \(u_t + 6uu_x + u_{xxx} = 0\) 在变换 \(u = 2(\log \tau)_{xx}\) 下化为双线性形式 \((D_x D_t + D_x^4) \tau \cdot \tau = 0\)。\(\tau\) 函数（tau函数）的广田扰动展开 \(\tau = 1 + \varepsilon \tau_1 + \varepsilon^2 \tau_2 + \cdots\) 给出 \(N\)-孤子解的显式表达式：\(\tau_N = \sum_{\mu_j=0,1} \exp\left(\sum_{j=1}^N \mu_j \eta_j + \sum_{1\leq i<j\leq N} \mu_i \mu_j A_{ij}\right)\)，其中 \(\eta_j = k_j x - \omega_j t + \eta_j^{(0)}\) 为第 \(j\) 个孤子的相位，\(e^{A_{ij}} = (k_i - k_j)^2/(k_i + k_j)^2\) 为孤子之间的相互作用相位因子。

\subsection{371.2 KdV方程与Sato理论}\label{kdvux65b9ux7a0bux4e0esatoux7406ux8bba}

\textbf{定义 371.3}（KdV方程）：\textbf{Korteweg-de Vries方程}为

\[\partial_t u + 6u\,\partial_x u + \partial_x^3 u = 0\]

（通过标度变换 \(u\to u/6\) 可与 \(u_t+uu_x+u_{xxx}=0\) 互化）。

\textbf{定理 371.2}（KdV的Lax对，Lax 1968）：KdV方程等价于算子方程 \(\partial_t L=[M,L]\)，其中

\[L = -\partial_x^2 - u(x,t), \quad M = -4\partial_x^3 - 6u\,\partial_x - 3u_x\]

\textbf{证明}：计算 \([M,L]=M L-L M\)。\(Lf=-f_{xx}-uf\)，\(Mf=-4f_{xxx}-6u f_x-3u_x f\)。直接计算得 \((\partial_t L-[M,L])f=(u_t+6uu_x+u_{xxx})f\)。因此 \(\partial_t L=[M,L]\) 当且仅当KdV方程成立。\(\blacksquare\)

\textbf{定理 371.3}（KdV的1-孤子解）：KdV方程具有如下形式的\textbf{1-孤子解}：

\[u(x,t) = 3c\,\operatorname{sech}^2\!\left(\frac{\sqrt{c}}{2}(x-ct)\right)\]

其中 \(c>0\) 为孤子传播速度，振幅与速度成正比（\(3c\)），宽度与 \(\sqrt{c}\) 成反比。

\textbf{证明}：设行波解 \(u(x,t)=U(\xi)\)，\(\xi=x-ct\)。代入KdV方程得 \(-cU'+6UU'+U'''=0\)，积分一次得 \(-cU+3U^2+U''=0\)，乘以 \(U'\) 再积分得 \(-\frac{c}{2}U^2+U^3+\frac{1}{2}(U')^2=C\)。边界条件 \(U,U',U''\to 0\)（\(\xi\to\pm\infty\)）得 \(C=0\)。解得 \(U(\xi)=3c\,\operatorname{sech}^2(\sqrt{c}\,\xi/2)\)。\(\blacksquare\)

\textbf{定义 371.4}（Sato理论与KP hierarchy）：Sato理论（Sato 1981, Date-Jimbo-Kashiwara-Miwa 1983）是可积系统理论的最高统一框架。其核心思想是将所有已知的可积非线性偏微分方程（KdV、KP、Boussinesq、Toda等）视为无穷维Grassmann流形上的动力系统。\textbf{KP hierarchy}（Kadomtsev-Petviashvili族）是包含KdV作为约化的普适可积族，由其 \(\tau\) 函数 \(\tau(t_1, t_2, t_3, \ldots)\)（\(t_1 = x, t_2 = y, t_3 = t, \ldots\)）满足的\textbf{Hirota双线性恒等式}（或等价的Plucker关系）定义：

\[\oint_{z=\infty} \tau(t - [z^{-1}]) \tau(t' + [z^{-1}]) e^{\sum_{n=1}^\infty (t_n - t'_n) z^n} dz = 0\]

其中 \(t \pm [z^{-1}] = (t_1 \pm z^{-1}, t_2 \pm z^{-2}/2, \ldots)\)。Sato理论的核心洞察是：\(\tau\) 函数是无穷维Grassmann流形 \(\operatorname{Gr}(V)\) 上Plucker坐标的生成函数，KP方程族对应于 \(\operatorname{Gr}\) 上的线性流。这一几何解释统一了孤子方程的代数构造（通过仿射Lie代数的表示论）、\(\tau\) 函数的行列式表示（Wronskian和Gram行列式）以及可积系统的Hamilton结构。

\subsection{371.3 反散射方法}\label{ux53cdux6563ux5c04ux65b9ux6cd5}

\textbf{定义 371.5}（反散射方法 / Inverse Scattering Transform, IST）：IST是将非线性可积PDE的初值问题化为线性散射问题的三步框架：

\textbf{步骤1（正散射）}：以初值 \(u(x,0)\) 构造Lax算子 \(L(0)\) 的散射数据（离散谱 \(\{\kappa_n, c_n\}\) 与反射系数 \(R(k)\)）。

\textbf{步骤2（时间演化）}：散射数据的时间演化是平凡的线性方程。对KdV方程，\(\kappa_n(t)=\kappa_n(0)\)（孤子速度守恒），\(c_n(t)=c_n(0)e^{4\kappa_n^3 t}\)，\(R(k,t)=R(k,0)e^{8ik^3 t}\)。

\textbf{步骤3（反散射）}：通过Gel'fand-Levitan-Marchenko（GLM）线性积分方程从演化后的散射数据重构位势 \(u(x,t)\)：

\[K(x,y,t) + F(x+y,t) + \int_x^\infty K(x,z,t)F(z+y,t)\,dz = 0\]

其中 \(F(x,t)\) 由散射数据构造，然后 \(u(x,t)=2\frac{d}{dx}K(x,x,t)\)。

\textbf{定理 371.4}（IST的可解性）：对于在无穷远处充分快衰减的初值 \(u(x,0)\)，GLM积分方程具有唯一解，因此KdV方程的初值问题通过IST可严格求解。当反射系数恒为零时，解为纯 \(N\)-孤子解，由 \(N\) 个离散特征值确定。

\textbf{证明}（框架）：GLM 方程是第二类 Fredholm 积分方程：\(K(x,y) + F(x+y) + \int_x^\infty K(x,z) F(z+y) dz = 0\)。当 \(u(x,0)\) 快速衰减时，散射数据包括：连续谱的反射系数 \(r(k)\)、离散谱的束缚态特征值 \(\kappa_n > 0\) 和归一化常数 \(c_n\)。核函数 \(F(x) = \sum_{n=1}^N c_n^2 e^{-\kappa_n x} + \frac{1}{2\pi} \int_{-\infty}^\infty r(k) e^{ikx} dk\)。当 \(r(k) \equiv 0\) 且仅有 \(N\) 个离散特征值时，\(F(x) = \sum_{n=1}^N c_n^2 e^{-\kappa_n x}\)。此时 GLM 方程退化为线性代数方程组，可精确求解，得到的 \(u(x,t)\) 为 \(N\)-孤子解。一般情形下，GLM 方程是 Fredholm 型积分方程，由 Fredholm 择一性定理，其核为退化核时解存在唯一。当 \(r(k)\) 非零时，需使用 Fredholm 行列式理论，但快速衰减保证了算子的紧性，从而解存在唯一。\(\blacksquare\)

反散射方法的数学意义超越了求解特定方程------它揭示了非线性可积PDE与线性谱问题之间深刻的''非线性Fourier变换''关系。在IST中，散射数据（离散谱+连续谱）扮演了非线性Fourier系数的角色，而GLM方程则是非线性Fourier逆变换。这一观点导致了可积系统与随机矩阵理论（RMT）之间的深刻联系：KdV方程的 \(\tau\) 函数可以表示为随机矩阵 ensembles 的 Fredholm 行列式或 Painleve 超越函数，而 Tracy-Widom 分布（最大特征值的极限分布）正是 KdV 方程在特定初值下解的行为。

\subsection{371.4 可积系统的代数结构}\label{ux53efux79efux7cfbux7edfux7684ux4ee3ux6570ux7ed3ux6784}

可积系统的现代理论揭示了一个统一的代数框架：所有经典可积系统都对应于某个Kac-Moody代数（或更一般的仿射Lie代数）的Drinfeld-Sokolov约化。对于KdV方程，它对应于 \(\widehat{\mathfrak{sl}}_2\) 的Drinfeld-Sokolov约化。一般地，\(\widehat{\mathfrak{g}}\)-KdV hierarchy（\(\mathfrak{g}\) 为任意有限维单Lie代数）由 \(\mathfrak{g}\) 的仿射化 \(\widehat{\mathfrak{g}}\) 定义，其守恒量由 \(\mathfrak{g}\) 的Casimir算子生成，孤子解对应于 \(\mathfrak{g}\) 的权空间。这一框架将可积系统的分类转化为Lie代数的分类问题，并解释了为什么可积系统与表示论、量子群（quantum groups）和共形场论有如此深刻的联系。量子可积系统（如XXZ自旋链、Sine-Gordon量子场论）的代数结构由Yang-Baxter方程和量子群 \(U_q(\widehat{\mathfrak{g}})\) 支配，而 \textbf{Bethe ansatz} 方法则是量子可积系统的核心求解技术------其本质是将量子系统的谱问题化为代数Bethe方程（由传递矩阵的交换性导出）。

\(\blacksquare\)

\hrulefill

\hrulefill

\hrulefill

\section{第458章：KdV方程与Lax对表示}\label{ux7b2c458ux7ae0kdvux65b9ux7a0bux4e0elaxux5bf9ux8868ux793a}

Korteweg-de Vries（KdV）方程是可积系统理论的起点和原型。它最初由 Boussinesq（1872）和 Korteweg-de Vries（1895）在浅水波研究中提出：\(\partial_t u + 6u \partial_x u + \partial_x^3 u = 0\)。KdV 方程是第一个被发现具有孤子解的 PDE------所谓孤子，是保持形状和速度传播的孤立波，且两个孤子碰撞后各自恢复原形，仅产生相位移动。1967 年，Gardner、Greene、Kruskal 和 Miura（GGKM）发现了反散射变换（Inverse Scattering Transform, IST），将 KdV 方程的初值问题转化为线性 Schrodinger 方程 \(-\psi'' + u \psi = \lambda \psi\) 的散射和反散射问题，首次揭示了非线性 PDE 与线性谱问题之间的对偶性。随后，Lax（1968）以极其简洁的算子形式统一了这一对应：KdV 方程等价于算子方程 \(\partial_t L = [M, L]\)，其中 \(L = -\partial_x^2 - u\)，\(M = -4 \partial_x^3 - 6u \partial_x - 3u_x\)。Lax 对的普适性远超 KdV：它不仅提供了检验给定 PDE 是否可积的判别法，而且自动保证无穷多守恒量的存在（因 \(\operatorname{tr}(L^n)\) 或谱不变量守恒）。Lax 对形式化是可积系统代数理论的基石，后来被推广为 Zakharov-Shabat 的零曲率表示、Drinfeld-Sokolov 约化和 Sato 理论。

\subsection{372.x KdV方程的反散射变换（IST）完整理论}\label{x-kdvux65b9ux7a0bux7684ux53cdux6563ux5c04ux53d8ux6362istux5b8cux6574ux7406ux8bba}

\textbf{定理 372.1}（KdV 方程与 Schrodinger 谱问题的等价性，GGKM 1967）：在快速衰减边界条件 \(u(x,t) \to 0\)（\(|x| \to \infty\)）下，KdV 方程 \(\partial_t u + 6u \partial_x u + \partial_{xxx} u = 0\) 等价于线性 Schrodinger 算子 \(L(t) = -\partial_x^2 - u(x,t)\) 的等谱形变：\(L(t)\) 的谱（离散特征值 \(\{-\kappa_n^2\}\) 和连续谱反射系数 \(R(k)\)）随 \(t\) 的演化是显式平凡的。

\textbf{证明}：KdV 的 Lax 方程为 \(\partial_t L = [M, L]\)。计算 \([\partial_t L - (ML - LM)]\psi = (-\partial_t u - 6u \partial_x u - \partial_x^3 u) \psi\)（对所有 \(\psi\)），故 \(\partial_t L = [M, L]\) 当且仅当 KdV 成立。等谱性：\(\partial_t L = [M, L]\) 蕴含 \(L(t)\) 经酉等价于 \(L(0)\)（通过 \(U(t) = T\exp(\int_0^t M(s) ds)\)）。因此 \(L(t)\) 的谱不依赖于 \(t\)。离散特征值 \(\kappa_n\) 守恒；对连续谱，特征函数 \(\psi(x, k, t) \sim e^{-ikx} + R(k, t) e^{ikx}\)（\(x \to +\infty\)），代入 Lax 方程的时间部分 \(\partial_t \psi = M \psi\)，利用 \(M \sim -4 \partial_x^3\) 在 \(|x| \to \infty\) 时（\(u, u_x \to 0\)），得 \(\partial_t R(k, t) = 8 i k^3 R(k, t)\)，故 \(R(k, t) = R(k, 0) e^{8 i k^3 t}\)。类似地，归一化常数 \(c_n(t) = c_n(0) e^{4 \kappa_n^3 t}\)。\(\blacksquare\)

\textbf{定理 372.2}（Gel'fand-Levitan-Marchenko 积分方程）：KdV 方程的初值 \(u(x, 0)\) 通过求解如下线性积分方程（GLM 方程）重构 \(u(x, t)\)：

\[K(x, y, t) + F(x+y, t) + \int_x^\infty K(x, z, t) F(z+y, t) \, dz = 0 \quad (y > x)\]

其中 \(F(x, t) = \sum_{n=1}^N c_n(0)^2 e^{8 \kappa_n^3 t - \kappa_n x} + \frac{1}{2\pi} \int_{-\infty}^\infty R(k, 0) e^{8 i k^3 t + ikx} \, dk\)。然后 \(u(x, t) = -2 \frac{d}{dx} K(x, x, t)\)。这一过程将非线性初值问题约化为线性积分方程的求解。

\textbf{证明}：GLM 方程来源于散射理论中的 Marchenko 积分方程。将 Schrodinger 方程 \(-\psi_{xx} + u \psi = k^2 \psi\) 的 Jost 解 \(f(x, k)\)（满足 \(f(x, k) \sim e^{ikx}\)，\(x \to +\infty\)）表示为三角核 \(K(x, y)\) 的 Fourier 变换：\(f(x, k) = e^{ikx} + \int_x^\infty K(x, y) e^{iky} dy\)。将此展开式代入 Schrodinger 方程，利用分部积分并比较 \(e^{iky}\) 的系数，得 \(u(x) = -2 \frac{d}{dx} K(x, x)\)（Miura 变换的推广）。GLM 方程来源于散射数据完备性关系：由谱测度的 Fouier 变换给出 \(F(x)\) 的定义，并导出 \(K(x, y)\) 所满足的第二类 Fredholm 积分方程。时间演化散射数据后，\(F(x, t)\) 如上给出，解 GLM 方程得 \(K(x, y, t)\)，再取对角线得 \(u(x, t)\)。\(\blacksquare\)

\subsection{372.y Lax方程与无穷多守恒量}\label{y-laxux65b9ux7a0bux4e0eux65e0ux7a77ux591aux5b88ux6052ux91cf}

\textbf{定理 372.3}（Lax 方程的等谱性质与守恒量）：若算子族 \(L(t)\) 满足 \(\partial_t L = [M, L]\)（\(M\) 反自伴时 \(\partial_t L = [M, L]\) 等价的酉演化），则谱不变量 \(I_n = \operatorname{tr}(L^n)\)（或更一般地，\(L\) 的本征值函数）是运动常数：\(\partial_t I_n = 0\)。对 KdV 方程，这给出无穷多个守恒量（Kruskal 积分）。

\textbf{证明}：\(\partial_t (L^n) = \sum_{k=0}^{n-1} L^k [M, L] L^{n-1-k}\)。由迹的循环性 \(\operatorname{tr}(L^k M L^{n-k}) = \operatorname{tr}(M L^{n-k} L^k) = \operatorname{tr}(M L^n)\)，且 \(\operatorname{tr}(L^k L M L^{n-1-k}) = \operatorname{tr}(L^{k+1} M L^{n-1-k})\)。通过重新指标，所有项配对相消：\(\partial_t \operatorname{tr}(L^n) = \sum_{k=0}^{n-1} \operatorname{tr}(L^k (ML - LM) L^{n-1-k}) = \operatorname{tr}([L^n, M]) = 0\)。对 KdV 的 \(L = -\partial_x^2 - u\)，\(\operatorname{tr}(L^n)\) 需适当正规化（在周期边界条件下为矩阵迹；在无穷直线上为 \(L^n\) 的积分核在对角线上的积分）。前几个守恒量为：\(\int u \, dx\)（质量），\(\int u^2 \, dx\)（动量），\(\int (u^3 - \frac{1}{2} u_x^2) \, dx\)（能量），等等。\(\blacksquare\)

\textbf{定理 372.4}（Lenard递推算子与KdV hierarchy）：KdV 方程属于无穷可交换流族------KdV hierarchy。定义\textbf{Lenard递推算子} \(\mathcal{L} = -\partial_x^2 - 4u - 2u_x \partial_x^{-1}\)（形式上），以及递推关系 \(\mathcal{L} \cdot \frac{\delta I_n}{\delta u} = \frac{\delta I_{n+1}}{\delta u}\)。KdV hierarchy 的第 \(n\) 个方程为 \(\partial_t u = \partial_x \frac{\delta I_n}{\delta u}\)，它们两两可交换且共享所有守恒量 \(I_n\)。

\textbf{证明}：由 \(L = -\partial_x^2 - u\) 的 Lax 形式 \(\partial_t L = [M_n, L]\)（\(M_n\) 为 \(2n+1\) 阶反对称微分算子），递推构造 \(M_n\)：从 \(M_0 = \partial_x\) 出发（对应平移流 \(\partial_t u = u_x\)），\(M_1 = -4 \partial_x^3 - 6u \partial_x - 3u_x\)（KdV），递推关系为 \([M_{n+1}, L] = \partial_x [M_n, L] \partial_x^{-1}\)（在形式计算意义下）。Lenard 递推算子 \(\mathcal{L}\) 满足 \(\mathcal{L} \cdot \partial_x [M_n, L] = \partial_x [M_{n+1}, L]\)。守恒量 \(I_n = \int P_n(u, u_x, \ldots) dx\)（\(P_n\) 为 \(u\) 及其导数的微分多项式）由 \(\frac{\delta I_n}{\delta u} = \operatorname{Res}(L^{n-1/2})\)（谱 \(\zeta\) 函数的展开系数）给出。KdV hierarchy 各流可交换性：\([M_m, M_n] = 0\)（由 \(\partial_{t_m} \partial_{t_n} u = \partial_{t_n} \partial_{t_m} u\) 和 Lax 形式的相容性）。\(\blacksquare\)

\hrulefill

\hrulefill

\hrulefill

\section{第459章：Hirota双线性方法与多孤子解}\label{ux7b2c459ux7ae0hirotaux53ccux7ebfux6027ux65b9ux6cd5ux4e0eux591aux5b64ux5b50ux89e3}

Hirota 双线性方法（Hirota 1971）是可积系统领域最具原创性的贡献之一，它提供了一种直接的、系统的构造非线性演化方程孤子解的方法，而无需借助反散射变换。Hirota 方法的核心步骤分为三步：首先通过适当的因变量变换（如 \(u = 2(\ln \tau)_{xx}\)）将原方程转化为双线性形式（Hirota bilinear form）；然后利用双线性方程的双线性恒等式（如 Plucker 关系的双线性模拟）进行微扰展开；最后，关键观察到微扰展开在有限项截断，从而得到显式的 \(N\)-孤子解。Hirota 方法的优美之处在于：一旦找到正确的因变量变换和双线性形式，\(N\)-孤子解的构造完全被代数化，且揭示出孤子之间相互作用仅通过两两相位因子 \(e^{A_{ij}}\) 的乘积来表达------这是孤子弹性碰撞的数学体现。Hirota 的 \(\tau\) 函数随后被 Sato 理论提升为可积系统的最核心对象：\(\tau\) 函数不仅是孤子解的产生函数，更是无穷维 Grassmann 流形上 Plucker 坐标的化身，蕴含了可积系统的全部代数几何信息。

\subsection{373.x Hirota D-算子与双线性恒等式}\label{x-hirota-d-ux7b97ux5b50ux4e0eux53ccux7ebfux6027ux6052ux7b49ux5f0f}

\textbf{定义 373.1}（Hirota 双线性算子 / D-算子）：对两个光滑函数 \(f(x, t)\) 和 \(g(x, t)\)，双线性 \(D\)-算子定义为

\[D_x^m D_t^n f \cdot g = (\partial_x - \partial_{x'})^m (\partial_t - \partial_{t'})^n f(x, t) g(x', t') \Big|_{x'=x, t'=t}\]

\(D\)-算子的关键性质是 Leibniz 法则的''反称化''版本：\(D_x f \cdot g = f_x g - f g_x\)，\(D_x^2 f \cdot g = f_{xx} g - 2 f_x g_x + f g_{xx}\)，\(D_x^3 f \cdot g = f_{xxx} g - 3 f_{xx} g_x + 3 f_x g_{xx} - f g_{xxx}\)。特别地，\(D_x f \cdot f = 0\)，\(D_x^2 f \cdot f = 2(f_{xx} f - f_x^2)\)。

\textbf{定理 373.1}（KdV方程的双线性化）：在因变量变换 \(u = 2(\ln \tau)_{xx} = 2(\tau_{xx} \tau - \tau_x^2)/\tau^2\) 下，KdV 方程 \(u_t + 6u u_x + u_{xxx} = 0\) 等价于 Hirota 双线性方程

\[(D_x D_t + D_x^4) \tau \cdot \tau = 0\]

即 \(\tau_{xt} \tau - \tau_x \tau_t + \tau_{xxxx} \tau - 4 \tau_{xxx} \tau_x + 3 \tau_{xx}^2 = 0\)。

\textbf{证明}：令 \(u = 2(\ln \tau)_{xx}\)。直接计算各阶导数并代入 KdV 方程。\(u_x = 2(\ln \tau)_{xxx}\)，\(u_t = 2(\ln \tau)_{xxt}\)，\(u_{xxx} = 2(\ln \tau)_{xxxxx}\)，\(6u u_x = 12 (\ln \tau)_{xx} (\ln \tau)_{xxx}\)。利用对数导数的组合恒等式，合并所有项后提取公因子 \(2/\tau^2\)，化简得 \((D_x D_t + D_x^4) \tau \cdot \tau = 0\)。核心验证：\((\ln \tau)_{xx} = (\tau_{xx} \tau - \tau_x^2)/\tau^2\)，\((\ln \tau)_{xxt} = (\tau_{xxt} \tau^2 - \cdots)/\tau^3\) 等等，全部通分后分子恰为双线性表达式。\(\blacksquare\)

\textbf{定理 373.2}（双线性方程的微扰展开与截断性质）：对双线性 KdV 方程 \((D_x D_t + D_x^4) \tau \cdot \tau = 0\)，做形式微扰展开 \(\tau = 1 + \varepsilon \tau_1 + \varepsilon^2 \tau_2 + \varepsilon^3 \tau_3 + \cdots\)（\(\varepsilon\) 为形式小参数）。代入并收集各阶 \(\varepsilon\)：

\begin{itemize}
\tightlist
\item
  \(O(\varepsilon)\)：\((\partial_x \partial_t + \partial_x^4) \tau_1 = 0\)（线性方程）
\item
  \(O(\varepsilon^2)\)：\((\partial_x \partial_t + \partial_x^4) \tau_2 = -\frac{1}{2}(D_x D_t + D_x^4) \tau_1 \cdot \tau_1\)
\item
  \(O(\varepsilon^3)\)：\((\partial_x \partial_t + \partial_x^4) \tau_3 = -(D_x D_t + D_x^4) (\tau_1 \cdot \tau_2)\)
\end{itemize}

关键发现：取 \(\tau_1 = \sum_{j=1}^N e^{\eta_j}\)（\(\eta_j = k_j x - k_j^3 t + \eta_j^{(0)}\)，满足线性色散关系），则级数在有限项截断：\(\tau_N\) 恰由 \(N\) 个指数项及其所有两两乘积的组合给出。

\textbf{证明}：\(\tau_1 = \sum_{j=1}^N e^{\eta_j}\) 满足 \(O(\varepsilon)\) 线性方程（因 \(\partial_t e^{\eta_j} = -k_j^3 e^{\eta_j}\)，\(\partial_x^4 e^{\eta_j} = k_j^4 e^{\eta_j}\)，\((\partial_t + \partial_x^3) \partial_x [\cdots] = 0\) 需 \(\omega_j = -k_j^3\)）。代入 \(O(\varepsilon^2)\) 方程：\((D_x D_t + D_x^4) e^{\eta_i} \cdot e^{\eta_j} = 2(k_i - k_j)[(\omega_i - \omega_j) + (k_i - k_j)^3] e^{\eta_i + \eta_j}\)。当 \(\omega_i = -k_i^3\)，\(\omega_j = -k_j^3\) 时，\((\omega_i - \omega_j) + (k_i - k_j)^3 = -(k_i^3 - k_j^3) + (k_i - k_j)^3 = (k_i - k_j)^3 - (k_i - k_j)(k_i^2 + k_i k_j + k_j^2) = (k_i - k_j)^3 - (k_i^3 - k_j^3) \neq 0\) 通分后，实际计算：\(\omega_i - \omega_j = k_j^3 - k_i^3\)，\((k_i - k_j)^3 = k_i^3 - 3k_i^2 k_j + 3k_i k_j^2 - k_j^3\)。和为 \(-3k_i^2 k_j + 3k_i k_j^2 = -3k_i k_j (k_i - k_j)\)。故 \((D_x D_t + D_x^4) e^{\eta_i} \cdot e^{\eta_j} = -6 k_i k_j (k_i - k_j)^2 e^{\eta_i + \eta_j}\)。因此 \(\tau_2\) 的特解为 \(\sum_{i,j} a_{ij} e^{\eta_i + \eta_j}\)（\(a_{ij} = \frac{1}{2} \cdot \frac{6 k_i k_j (k_i - k_j)^2}{(\text{线性算子在 } \eta_i+\eta_j \text{ 的值})} = \frac{(k_i - k_j)^2}{(k_i + k_j)^2} \cdot \frac{1}{2}\)？需仔细计算分母）。关键：\(a_{ij} = e^{A_{ij}}\) 其中 \(A_{ij} = \ln[(k_i - k_j)^2/(k_i + k_j)^2]\)。更高阶的展开因 \(N\)-孤子的特殊代数结构而截断：\(\tau_N = \sum_{\mu_j \in \{0,1\}} \exp(\sum_{j=1}^N \mu_j \eta_j + \sum_{1 \leq i < j \leq N} \mu_i \mu_j A_{ij})\)。\(\blacksquare\)

\subsection{373.y N-孤子解的显式构造与相互作用}\label{y-n-ux5b64ux5b50ux89e3ux7684ux663eux5f0fux6784ux9020ux4e0eux76f8ux4e92ux4f5cux7528}

\textbf{定理 373.3}（KdV方程的多孤子解）：KdV 方程 \(u_t + 6u u_x + u_{xxx} = 0\) 的 \(N\)-孤子解为

\[u_N(x, t) = 2 \frac{\partial^2}{\partial x^2} \ln \tau_N(x, t)\]

其中 \(\tau_N(x, t) = \sum_{\mu_1, \ldots, \mu_N \in \{0, 1\}} \exp\left( \sum_{j=1}^N \mu_j \eta_j + \sum_{1 \leq i < j \leq N} \mu_i \mu_j A_{ij} \right)\)，\(\eta_j = k_j x - k_j^3 t + \eta_j^{(0)}\)（\(k_j > 0\) 为第 \(j\) 个孤子的波数），\(e^{A_{ij}} = \left( \frac{k_i - k_j}{k_i + k_j} \right)^2\)。

当 \(t \to \pm \infty\) 时（孤子充分分离），解渐近为 \(N\) 个独立单孤子的叠加，每个具有形式 \(u \sim \frac{k_j^2}{2} \operatorname{sech}^2(\frac{k_j}{2}(x - k_j^2 t - \delta_j^{\pm}))\)，其中 \(\delta_j^{\pm}\) 为碰撞产生的相位移动（asymptotic phase shift）。\(N\)-孤子解的最重要特征是：孤子碰撞是弹性的（保持形状和速度），仅产生相位移动 \(\Delta_j = \delta_j^+ - \delta_j^- = \sum_{i < j} A_{ij} - \sum_{i > j} A_{ij}\)。

\textbf{证明}：\(N=1\) 时，\(\tau_1 = 1 + e^{\eta_1}\)，\(u_1 = 2(\ln \tau_1)_{xx} = \frac{k_1^2}{2} \operatorname{sech}^2(\eta_1/2) = \frac{k_1^2}{2} \operatorname{sech}^2(\frac{k_1}{2}(x - k_1^2 t - x_0))\)，恰为速度为 \(k_1^2\)、振幅为 \(k_1^2/2\) 的单孤子。

\(N=2\) 时，\(\tau_2 = 1 + e^{\eta_1} + e^{\eta_2} + e^{\eta_1 + \eta_2 + A_{12}}\)。代入 \(u = 2(\ln \tau_2)_{xx}\) 得二孤子解。渐近分析：固定参考系跟随孤子 1（令 \(\eta_1 = \text{const}\)，\(t \to \pm \infty\)），则孤子 2 的相位 \(\eta_2 \to \mp \infty\)。\(\tau_2 \sim 1 + e^{\eta_1}\)（\(t \to -\infty\)）和 \(\tau_2 \sim e^{\eta_2}(e^{A_{12}} + e^{\eta_1})\)（\(t \to +\infty\)），重标度对数后，相位移动 \(\Delta_1 = \pm A_{12}\)。计算得 \(\Delta \eta_1^{\text{shift}} = A_{12} = \ln[(k_1 - k_2)^2/(k_1 + k_2)^2]\)。\(\blacksquare\)

\textbf{定理 373.4}（Hirota方法的普适性 / 其他可积方程的双线性化）：Hirota 方法不仅适用于 KdV 方程，几乎所有已知的可积方程均可双线性化。下表列出几个关键例子：

\begin{itemize}
\tightlist
\item
  \textbf{修正 KdV（mKdV）方程} \(v_t + 6v^2 v_x + v_{xxx} = 0\)，变换 \(v = i(\ln(F/G))_x\)（\(F = G^*\)）：\((D_t + D_x^3) F \cdot G = 0\)，\(D_x^2 F \cdot G = 0\)。
\item
  \textbf{非线性 Schrodinger（NLS）方程} \(i\psi_t + \psi_{xx} + 2|\psi|^2 \psi = 0\)，变换 \(\psi = G/F\)（\(F\) 实）：\((i D_t + D_x^2) G \cdot F = 0\)，\(D_x^2 F \cdot F = 2 G G^*\)。
\item
  \textbf{Kadomtsev-Petviashvili（KP）方程} \((u_t + 6uu_x + u_{xxx})_x + 3u_{yy} = 0\)，变换 \(u = 2(\ln \tau)_{xx}\)：\((D_x D_t + D_x^4 + 3D_y^2) \tau \cdot \tau = 0\)。
\item
  \textbf{Toda 格点方程}（离散可积系统）。
\end{itemize}

每个双线性方程通过类似的微扰展开产生其 \(N\)-孤子解，其代数结构（\(\tau\) 函数的行列式表示）由相应的群论结构（\(\tau\) 函数作为相应 Lie 群表示的矩阵元）统一。

\textbf{证明}：（以 NLS 为例）将 \(\psi = G/F\) 代入 NLS 方程。计算 \(\psi_t, \psi_{xx}, |\psi|^2 \psi\)，将分子按 \(F\) 的幂次整理。利用双线性算子的恒等式 \((i D_t + D_x^2) G \cdot F = i(G_t F - G F_t) + G_{xx} F - 2 G_x F_x + G F_{xx}\) 和 \(D_x^2 F \cdot F = 2(F_{xx} F - F_x^2)\)。经过代数化简（利用 \(|\psi|^2 = G G^* / F^2\)）即得所述双线性方程组。其他方程的双线性化同理，关键在于找到恰当的因变量变换------这既是 Hirota 方法的核心技巧所在，也是其最大的''艺术性''挑战。\(\blacksquare\)

\newpage

\chapter{13-数学物理/数学物理/06-重正化与Vlasov.md}\label{ux6570ux5b66ux7269ux7406ux6570ux5b66ux7269ux740606-ux91cdux6b63ux5316ux4e0evlasov.md}

\section{第460章：重正化群与重正化理论}\label{ux7b2c460ux7ae0ux91cdux6b63ux5316ux7fa4ux4e0eux91cdux6b63ux5316ux7406ux8bba}

重正化（renormalization）是量子场论中消除紫外发散的系统性方法，其核心思想在于：将裸参数重新定义为依赖于截断标度的跑动耦合常数，使得物理可观测量在移除截断后保持有限。重正化不仅是消除无穷大的技术手段，更揭示了量子场论在能量标度下的行为------不同能量标度下的物理由跑动耦合常数描述，这一深刻的物理洞察由 Wilson 的重正化群理论（1971）所阐明，并因此获得 1982 年 Nobel 物理学奖。

\subsection{372.1 重正化的物理动机与正规化}\label{ux91cdux6b63ux5316ux7684ux7269ux7406ux52a8ux673aux4e0eux6b63ux89c4ux5316}

\textbf{定义 372.1}（紫外发散）：在微扰量子场论中，对无限制动量积分出现的发散称为\textbf{紫外（UV）发散}。典型例子是 \(\phi^4\) 理论中单圈自能图 \(\int d^4k/(k^2+m^2)\)，该积分在大动量 \(|k|\to\infty\) 时二次发散。发散度（superficial degree of divergence）\(\omega(G)\) 是 Feynman 图 \(G\) 的圈动量积分在大动量极限下的幂次发散度，由量纲分析给出：对 \(\phi^4\) 理论在四维，\(\omega(G) = 4 - E\)（\(E\) 为外腿数），故两腿图（\(E=2\)）为二次发散，四腿图（\(E=4\)）为对数发散。

\textbf{定义 372.2}（动量截断正规化）：引入动量截断 \(\Lambda>0\)，将发散积分截断为有限值： \[\int \frac{d^4k}{(2\pi)^4} \frac{1}{k^2+m^2} \;\longrightarrow\; \int_{|k|\leq\Lambda} \frac{d^4k}{(2\pi)^4} \frac{1}{k^2+m^2}\] 物理结果在 \(\Lambda\to\infty\) 极限下通过重新定义耦合常数和场归一化保持有限。截断正规化直观但破坏 Lorentz 协变性和规范不变性（在非 Abel 规范理论中），因此高能物理中更常用维数正规化。

\textbf{定义 372.3}（维数正规化，'t Hooft-Veltman 1972）：在 \(d=4-\varepsilon\) 维时空中计算 Feynman 积分，发散表现为 \(1/\varepsilon\) 极点。对任意 Feynman 图 \(G\)，其维数正规化振幅为 \[\mathcal{A}_G = \mu^{2\varepsilon} \int \frac{d^dk}{(2\pi)^d} I_G(k,p_i)\] 其中 \(\mu\) 为重整化标度（引入以保持耦合常数无量纲），\(d\to 4\) 时 \(1/\varepsilon\) 极点对应于原理论的对数发散。维数正规化的核心优势在于保持 Lorentz 协变性和规范对称性（在最小减除方案 MS 和 \(\overline{\text{MS}}\) 方案中）。

\textbf{定理 372.1}（维数正规化的基本积分公式）：在 \(d\) 维时空中， \[\int \frac{d^dk}{(2\pi)^d} \frac{1}{(k^2 + 2p\cdot k + M^2)^\alpha} = \frac{i(-1)^\alpha}{(4\pi)^{d/2}} \frac{\Gamma(\alpha - d/2)}{\Gamma(\alpha)} \frac{1}{(M^2 - p^2)^{\alpha - d/2}}\] 此公式是维数正规化中计算所有 Feynman 积分的基础。\(\Gamma\) 函数在 \(d=4\)（\(\alpha=2\)）时产生极点 \(\Gamma(2-d/2) = \Gamma(\varepsilon/2) \sim 2/\varepsilon - \gamma_E + O(\varepsilon)\)。

\textbf{证明}：利用 Schwinger 参数化：\(\frac{1}{A^\alpha} = \frac{1}{\Gamma(\alpha)} \int_0^\infty dt \, t^{\alpha-1} e^{-tA}\)。取 \(A = k^2 + 2p\cdot k + M^2\)，则积分变为 \(\int \frac{d^d k}{(2\pi)^d} \frac{1}{\Gamma(\alpha)} \int_0^\infty dt \, t^{\alpha-1} \exp(-t(k^2 + 2p\cdot k + M^2))\)。交换积分次序，完成平方：\(k^2 + 2p\cdot k + M^2 = (k + p)^2 + M^2 - p^2\)。做平移 \(k' = k + p\)，则动量积分化为 Gauss 积分：\(\int \frac{d^d k'}{(2\pi)^d} e^{-t k'^2} = \frac{1}{(4\pi t)^{d/2}}\)。于是积分 \(= \frac{1}{\Gamma(\alpha)} \int_0^\infty dt \, t^{\alpha-1} \frac{1}{(4\pi t)^{d/2}} e^{-t(M^2-p^2)} = \frac{1}{(4\pi)^{d/2}\Gamma(\alpha)} \int_0^\infty dt \, t^{\alpha - d/2 - 1} e^{-t(M^2-p^2)}\)。由 \(\Gamma\) 函数的定义 \(\int_0^\infty t^{s-1} e^{-at} dt = \Gamma(s) / a^s\)，得最终结果 \(\frac{i(-1)^\alpha}{(4\pi)^{d/2}} \frac{\Gamma(\alpha - d/2)}{\Gamma(\alpha)} \frac{1}{(M^2 - p^2)^{\alpha - d/2}}\)。\(i\) 因子来自 Wick 旋转（\(k^0 \to i k^0_E\)）。\(\blacksquare\)

\textbf{其他正规化方法}： - \textbf{Pauli-Villars 正规化}：引入辅助大质量场以抵消发散，保持 Lorentz 协变性但在非 Abel 规范理论中破坏规范不变性； - \textbf{格点正规化}：将时空离散化为格点 \(a\mathbb{Z}^4\)，动量积分截断在 \(\pi/a\)，保持规范不变性（Wilson 作用量）但破坏 Lorentz 协变性（在连续极限 \(a\to 0\) 中恢复）； - \textbf{函数正规化}（\(\zeta\)-函数正规化、热核正规化）：通过解析延拓定义发散行列式，常用于弯曲时空量子场论。

\subsection{372.2 BPHZ重正化方案}\label{bphzux91cdux6b63ux5316ux65b9ux6848}

\textbf{定义 372.4}（Zimmermann 森林公式，1969）：设 Feynman 图 \(G\) 的发散子图集合为 \(\mathcal{F}(G)\)。\(G\) 的\textbf{重正化振幅} \(R_G\) 由森林公式给出 \[R_G = \sum_{U \in \mathcal{W}(\mathcal{F}(G))} \prod_{\gamma \in U} (-t_\gamma) \cdot I_G\] 其中 \(\mathcal{W}(\mathcal{F}(G))\) 是 \(\mathcal{F}(G)\) 的所有\textbf{森林}（两两互不重叠的连通发散子图族），\(t_\gamma\) 是在子图 \(\gamma\) 上定义的 Taylor 减法算子（在零外动量处展开至发散阶 \(d(\gamma)\)），\(I_G\) 是未重整化的 Feynman 积分。

\textbf{森林公式的动机}：重整化的核心困难在于处理重叠发散（overlapping divergences）------即一个发散子图嵌套在另一个发散子图中。Zimmermann 森林公式通过系统地扣除所有发散子图及其组合，精确处理了重叠发散问题。森林（forest）的引入是为了避免重复扣除（over-subtraction），每个森林对应一组互不相交的减法操作。

\textbf{定义 372.5}（Bogoliubov \(R\)-操作）：\(R\)-操作递归定义如下： \[\bar{R}_G = I_G + \sum_{\{\gamma_1,\ldots,\gamma_k\} \in \mathcal{F}(G)} I_{G/\{\gamma_i\}} \circ \prod_i (-t^{\gamma_i} \bar{R}_{\gamma_i})\] 其中 \(\{G/\{\gamma_i\}\}\) 是将 \(G\) 中所有子图 \(\gamma_i\) 缩并为点的约化图。\(R_G = (1-t^G)\bar{R}_G\) 为最终的重正化值。\(R\)-操作首先对所有真子图进行重整化（\(\bar{R}_{\gamma_i}\)），然后对整体散度进行最后一次减法（\((1-t^G)\)）。

\textbf{定理 372.2}（BPHZ 定理，Bogoliubov-Parasiuk 1957, Hepp 1966, Zimmermann 1969）：对任意可重整的量子场论，森林公式和 \(R\)-操作给出的重正化振幅 \(R_G\) 对所有 Feynman 图 \(G\) 是有限的、Lorentz 不变的，且满足因果性和幺正性。重整化后的 Feynman 振幅对动量截断 \(\Lambda\to\infty\)（或 \(\varepsilon\to 0\)）一致收敛。

\textbf{证明}：证明分三个核心步骤。(1) 森林公式与 \(R\)-操作的等价性：对每个 Feynman 图 \(G\)，Zimmermann 森林公式对所有森林（两两互不重叠的连通发散子图族）求和，每个森林 \(\gamma_1, \ldots, \gamma_k\) 对应一组减法操作 \((-t_{\gamma_1})\cdots(-t_{\gamma_k})\)。Bogoliubov \(R\)-操作递归定义 \(\bar{R}_G = I_G + \sum I_{G/\{\gamma_i\}} \circ \prod_i (-t^{\gamma_i} \bar{R}_{\gamma_i})\)，两者通过森林与子图缩并的对应关系等价。(2) Hepp 的收敛定理：对每个已被适当减除的 Feynman 积分，将积分区域（\(\mathbb{R}^d\) 中的圈动量）剖分为 Hepp 扇区。在每个扇区中，动量按大小排序，积分可写为迭代积分。Taylor 减法算子 \(t_\gamma\) 精确消除了所有扇区中的发散子积分------在子图 \(\gamma\) 的动量趋于无穷时，\(t_\gamma\) 减去了 \(I_\gamma\) 的 Taylor 展开（至发散阶 \(d(\gamma)\)），使得减除后的积分在 \(\Lambda \to \infty\) 下绝对收敛。每种扇区中的收敛性由幂次计数（power counting）保证：减除后有效发散度 \(\omega_{\text{eff}}(\gamma) < 0\)。(3) 幺正性和因果性：重正化后的 \(S\) 矩阵满足 \(S^\dagger S = 1\) 等价于 Cutkosky 切割规则在重正化后成立。重正化不改变 \(S\) 矩阵的解析结构（极点位置和留数），因此保持幺正性。因果性由 Bogoliubov-Shirkov 因果条件保证，重正化操作保持微扰展开的因果结构。\(\blacksquare\)

\textbf{定义 372.6}（重整化方案）：不同重整化方案对应不同的减法条件： - \textbf{在壳方案}（on-shell）：减除在物理质量壳上 \(p^2 = m_{\text{phys}}^2\) 处进行，参数直接对应物理可观测量； - \textbf{最小减除方案}（MS）：在维数正规化中仅减除 \(1/\varepsilon\) 极点，结果与重整化标度 \(\mu\) 有关； - \textbf{\(\overline{\text{MS}}\) 方案}（修正最小减除）：减除 \(1/\varepsilon - \gamma_E + \ln 4\pi\)，便于消除 \(\Gamma\) 函数展开中的常数项，是微扰 QCD 计算的标准方案。

不同重整化方案之间的转换由有限的重整化常数（finite renormalization）联系，物理预言（如 \(S\) 矩阵元）在方案转换下不变。

\subsection{\texorpdfstring{372.3 Callan-Symanzik方程与\(\beta\)-函数}{372.3 Callan-Symanzik方程与\textbackslash beta-函数}}\label{callan-symanzikux65b9ux7a0bux4e0ebeta-ux51fdux6570}

\textbf{定义 372.7}（跑动耦合常数与 \(\beta\)-函数）：设 \(g(\mu)\) 为重整化标度 \(\mu\) 处的有效耦合常数。\textbf{\(\beta\)-函数}定义为 \[\beta(g) \equiv \mu \frac{\partial g}{\partial \mu}\] 在微扰论中，\(\beta(g) = \beta_0 g^3 + \beta_1 g^5 + O(g^7)\)，其中 \(\beta_0, \beta_1\) 是方案无关的普适系数（前两阶系数）。\(\beta\)-函数描述了耦合常数随能量标度的流动：若 \(\beta(g) > 0\)，耦合常数随 \(\mu\) 增大而增大（如 QED）；若 \(\beta(g) < 0\)，耦合常数随 \(\mu\) 增大而减小------\textbf{渐近自由}（如 QCD）。

\textbf{定理 372.3}（Callan-Symanzik 方程，1970）：设 \(\Gamma^{(n)}(p_i; g, \mu)\) 是 \(n\) 点单粒子不可约（1PI）Green 函数，\(\gamma(g)\) 为反常量纲（描述场算符的标度维数偏离经典值）。则 \[\left[\mu\frac{\partial}{\partial\mu} + \beta(g)\frac{\partial}{\partial g} + n\gamma(g)\right]\Gamma^{(n)}(p_i; g, \mu) = 0\] Callan-Symanzik 方程是重整化群方程，表达了物理 Green 函数在重整化标度改变下的不变性（裸 Green 函数不依赖 \(\mu\)）。

\textbf{证明}：由重整化群不变性，裸 Green 函数 \(\Gamma^{(n)}_B(p_i; g_0, \Lambda)\) 不依赖于 \(\mu\)。利用 \(g_0 = Z_g(g, \mu/\Lambda) g \mu^{\varepsilon}\) 和场重整化常数 \(Z_\phi\)（\(\Gamma^{(n)} = Z_\phi^{n/2} \Gamma^{(n)}_B\)），对 \(\mu\) 求导： \[\mu\frac{d}{d\mu}\Gamma^{(n)}_B = 0 \Rightarrow \left(\mu\frac{\partial}{\partial\mu} + \mu\frac{\partial g}{\partial\mu}\frac{\partial}{\partial g} + \frac{n}{2}\mu\frac{\partial\ln Z_\phi}{\partial\mu}\right)\Gamma^{(n)} = 0\] 定义 \(\beta(g) = \mu\frac{\partial g}{\partial\mu}\) 和 \(\gamma(g) = \frac{1}{2}\mu\frac{\partial\ln Z_\phi}{\partial\mu}\)，即得 CS 方程。\(\blacksquare\)

\textbf{定理 372.4}（渐近自由与 Gross-Wilczek-Politzer 定理，1973）：\(SU(N)\) 规范理论（无标量场或费米子数有限）是渐近自由的：\(\beta(g) < 0\) 对小 \(g\) 成立。具体地，对 \(SU(N)\) 规范理论与 \(N_f\) 味 Dirac 费米子， \[\beta_0 = -\frac{1}{(4\pi)^2}\left(\frac{11}{3}N - \frac{2}{3}N_f\right)\] 当 \(N_f < 11N/2\) 时 \(\beta_0 < 0\)，理论为渐近自由。QCD（\(N=3, N_f=6\)）满足此条件，解释了高能实验中夸克表现为几乎自由粒子的现象（渐近自由）。Gross、Wilczek 和 Politzer 因此获得 2004 年 Nobel 物理学奖。

\textbf{证明}：\(\beta\) 函数的单圈系数 \(\beta_0\) 由三个贡献之和给出：规范场自相互作用（胶子圈）、鬼场圈（Faddeev-Popov 鬼）和费米子圈。在 \(SU(N)\) 规范理论中，规范场对 \(\beta_0\) 的贡献为 \(-\frac{11}{3}N\)（来自三胶子和四胶子顶点的单圈修正），费米子（\(N_f\) 味）的贡献为 \(+\frac{2}{3}N_f\)（来自费米子圈的真空极化）。具体计算：\(\beta(g) = -\frac{g^3}{(4\pi)^2} \beta_0 + O(g^5)\)，其中 \(\beta_0 = \frac{11}{3}C_2(G) - \frac{4}{3}T(R)N_f\)，对 \(SU(N)\) 有 \(C_2(G) = N\)，\(T(R) = 1/2\)（基础表示），故 \(\beta_0 = \frac{11}{3}N - \frac{2}{3}N_f\)。当 \(N_f < 11N/2\) 时 \(\beta_0 > 0\)，但 \(\beta(g) = -\beta_0 g^3/(4\pi)^2 + \cdots\)，故 \(\beta(g) < 0\)（渐近自由）。\(\blacksquare\)

\subsection{372.4 Wilson重正化群流}\label{wilsonux91cdux6b63ux5316ux7fa4ux6d41}

\textbf{定义 372.8}（Kadanoff 分块自旋变换，1966）：在临界现象中，将 \(d\) 维格点系统分组为边长为 \(b\) 的超块，每个超块内的自旋取平均（粗粒化），再重新标度长度 \(x \to x/b\) 和场变量 \(\phi \to b^{\Delta}\phi\)。此操作定义了相空间上的\textbf{重正化群变换} \(\mathcal{R}_b\)。

\textbf{Wilson 重正化群的动机}：传统的 Gell-Mann-Low 重整化群仅关注 Green 函数随动量标度的变化，而 Wilson 的重正化群则是将 Kadanoff 的实空间粗粒化思想推广到连续场论。Wilson 的核心见解是：在临界点附近，关联长度发散，系统在所有尺度上都有涨落，因此必须同时处理所有尺度上的自由度。重正化群通过在动量空间中逐层积分掉高动量模式（\(\Lambda/b < |k| < \Lambda\)），研究有效作用量随尺度 \(b\) 的变化。

\textbf{定理 372.5}（重正化群不动点与普适类）：重正化群变换 \(\mathcal{R}_b\) 在耦合常数空间 \(\{g_i\}\) 中定义了一个流（RG flow）：\(\frac{dg_i}{d\ln b} = \beta_i(g)\)。RG 不动点 \(g^*\) 满足 \(\beta_i(g^*) = 0\)。在 \(g^*\) 附近线性化： \[\frac{d}{d\ln b}(g_i - g_i^*) = \sum_j M_{ij}(g_j - g_j^*), \quad M_{ij} = \left.\frac{\partial\beta_i}{\partial g_j}\right|_{g^*}\] \(M\) 的特征值 \(\lambda_\alpha\) 分类了参数空间中的方向： - \(\lambda_\alpha > 0\)（相关方向，relevant）：参数在 RG 流下增大，最终远离不动点； - \(\lambda_\alpha < 0\)（无关方向，irrelevant）：参数在 RG 流下减小，最终趋近不动点； - \(\lambda_\alpha = 0\)（边缘方向，marginal）：需考虑高阶修正。

所有流向同一不动点的理论构成一个\textbf{普适类}（universality class），具有相同的临界指数（由相关方向的特征值决定）。普适类解释了为什么截然不同的微观物理系统（如液体-气体相变与铁磁相变）在临界点附近表现出相同的标度行为。

\textbf{证明}：RG 变换 \(\mathcal{R}_b\) 将耦合常数 \(\{g_i\}\) 映射为 \(\{g_i(b)\}\)。\(\beta\) 函数 \(\beta_i(g) = dg_i/d\ln b\) 控制 RG 流。在不动点 \(g^*\) 处线性化即得上述线性 RG 方程。相关方向的特征值 \(\lambda_\alpha = b^{y_\alpha}\)（\(y_\alpha > 0\)）对应临界指数：例如，温度场的相关指数 \(y_t = 1/\nu\)（\(\nu\) 为关联长度临界指数），磁场场的相关指数 \(y_h = \beta\delta/\nu\)。无关方向在长波极限下消失，解释了普适性------微观细节不影响临界行为。\(\blacksquare\)

\textbf{Wilson 精确重正化群方程}：在连续场论中，Wilson 引入有效作用量 \(S_\Lambda[\phi]\)（仅包含动量 \(|k| < \Lambda\) 的模式），通过积分掉动量壳层 \(\Lambda/b < |k| < \Lambda\) 中的模式，得到低能有效作用量 \(S_{\Lambda/b}[\phi]\)。在无穷小壳层极限下，这给出 Polchinski 精确 RG 方程或 Wetterich 方程： \[\partial_t \Gamma_k = \frac{1}{2}\operatorname{Tr}\left[(\Gamma_k^{(2)} + R_k)^{-1} \partial_t R_k\right]\] 其中 \(t = \ln(k/\Lambda)\)，\(\Gamma_k\) 为尺度 \(k\) 处的有效平均作用量，\(R_k\) 为红外截断函数。该方程是精确的（非微扰），在统计物理和量子场论中广泛应用。

\hrulefill

\section{第461章：Vlasov方程的数学理论}\label{ux7b2c461ux7ae0vlasovux65b9ux7a0bux7684ux6570ux5b66ux7406ux8bba}

Vlasov 方程描述无碰撞等离子体或自引力系统中粒子分布函数在相空间中的演化。它是 Boltzmann 方程在碰撞项消失时的极限，属于平均场动力学方程。Vlasov 方程在等离子体物理、天体物理（恒星动力学）和加速器物理中具有核心地位。其数学理论涉及 PDE 的适定性、弱解的紧性方法、以及 Landau 阻尼的严格证明。

\subsection{373.1 Vlasov方程与Vlasov-Poisson系统}\label{vlasovux65b9ux7a0bux4e0evlasov-poissonux7cfbux7edf}

\textbf{定义 373.1}（Vlasov 方程）：设 \(f(t,x,v) \geq 0\) 为时刻 \(t\) 在位置 \(x \in \mathbb{R}^d_x\)、速度 \(v \in \mathbb{R}^d_v\) 处的粒子分布函数。粒子所受平均场力为 \(F(t,x)\)，则\textbf{Vlasov 方程}为 \[\partial_t f + v\cdot\nabla_x f + F\cdot\nabla_v f = 0\] 该方程表达了相空间体积的 Liouville 守恒：沿特征线 \(\dot{x}=v,\;\dot{v}=F(t,x)\)，\(f\) 为常数（\(f(t, X(t), V(t)) = f_0(x, v)\) 其中 \((X(t), V(t))\) 为特征线）。Vlasov 方程是输运方程在相空间中的特例，其对初始条件的依赖完全由特征线决定。

\textbf{定义 373.2}（Vlasov-Poisson 系统）：当力场 \(F\) 由自洽 Poisson 方程确定时，得到\textbf{Vlasov-Poisson 系统}： \[\begin{cases}
\partial_t f + v\cdot\nabla_x f + F\cdot\nabla_v f = 0 \\[4pt]
F = -\nabla_x \Phi, \quad \Delta_x \Phi = \gamma_d \int_{\mathbb{R}^d_v} f(t,x,v)\,dv
\end{cases}\] 其中对等离子体情形 \(\gamma_d=1\)（静电势，\(\Delta\Phi = -\rho\) 归入符号约定），对自引力情形 \(\gamma_d = -4\pi G\)（引力势，\(G\) 为引力常数）。系统的总能量 \[\mathcal{E}(t) = \frac{1}{2}\iint |v|^2 f\,dx\,dv + \frac{1}{2}\int |\nabla\Phi|^2\,dx\] 是形式守恒量。对于等离子体（\(\gamma_d = 1\)），能量正定；对于引力（\(\gamma_d = -4\pi G\)），能量无下界，这导致引力系统的 rich 现象（如引力塌缩和 Jeans 不稳定性）。

\textbf{定义 373.3}（Vlasov-Maxwell 系统）：当电磁场由 Maxwell 方程决定时，得到 Vlasov-Maxwell（相对论）系统： \[\partial_t f + \hat{v}\cdot\nabla_x f + (E + \hat{v}\times B)\cdot\nabla_p f = 0\] \(\partial_t E = \nabla\times B - j\)，\(\partial_t B = -\nabla\times E\)，\(\nabla\cdot E = \rho\)，\(\nabla\cdot B = 0\)，其中 \(\hat{v} = p/\sqrt{1+|p|^2}\)（相对论速度），\(\rho = \int f dp\)，\(j = \int \hat{v} f dp\)。Vlasov-Maxwell 系统在 \(d=3\) 时的全局适定性仍是开放问题（仅对小初值或在某些对称性假设下已知）。

\subsection{373.2 弱解的存在性：DiPerna-Lions理论}\label{ux5f31ux89e3ux7684ux5b58ux5728ux6027diperna-lionsux7406ux8bba}

\textbf{定义 373.4}（重正规化解，DiPerna-Lions 1989）：设 \(f_0\in L^1\cap L^\infty(\mathbb{R}^d_x\times\mathbb{R}^d_v)\) 非负且有有限动能。分布函数 \(f\in L^\infty([0,T];L^1\cap L^\infty)\) 称为 Vlasov-Poisson 系统的\textbf{重正规化解}，若对任意 \(\beta\in C^1(\mathbb{R})\cap L^\infty(\mathbb{R})\) 满足 \(\beta(0)=0\) 和 \(|\beta'(t)|\leq C(1+|t|^p)\)，有 \[\partial_t \beta(f) + v\cdot\nabla_x \beta(f) + F\cdot\nabla_v \beta(f) = 0\] 在分布意义下成立。重正规化条件补偿了缺乏 \(W^{1,1}\) 正则性的困难，使得乘积 \(F\cdot\nabla_v f\) 有定义（通过 \(\beta(f)\) 的链式法则）。

\textbf{重正规化解的动机}：当 \(f\) 仅为 \(L^\infty\) 时，\(F\cdot\nabla_v f\) 在经典意义下未必有定义（\(F\) 可能仅在 \(L^p\) 中，\(\nabla_v f\) 可能仅为测度）。DiPerna-Lions 的关键洞察是：虽然 \(f\) 本身可能不满足链式法则，但 \(\beta(f)\) 的输运方程在分布意义下成立，这足以保证解的适定性。

\textbf{定理 373.1}（DiPerna-Lions 全局存在性定理，1989）：设 \(f_0 \geq 0\)，\(f_0 \in L^1 \cap L^\infty(\mathbb{R}^6)\)，且 \(\iint |v|^2 f_0\,dx\,dv < \infty\)。则 Vlasov-Poisson 系统存在全局重正规化解 \(f \in L^\infty([0,\infty); L^1 \cap L^\infty)\)。若进一步 \(f_0 \in C^1\)，则该解为经典解。

\textbf{证明}：构造正则化系统：对力场 \(E\) 进行磨光（卷积光滑核 \(\rho_\varepsilon\)），得到逼近序列 \((f^\varepsilon, E^\varepsilon)\) 满足 \(\partial_t f^\varepsilon + v \cdot \nabla_x f^\varepsilon + E^\varepsilon \cdot \nabla_v f^\varepsilon = 0\)，\(E^\varepsilon = \nabla_x \Delta_x^{-1}(\int f^\varepsilon dv * \rho_\varepsilon)\)。由特征线法，\(f^\varepsilon\) 沿特征线为常数，最大模 \(\|f^\varepsilon\|_{L^\infty} \leq \|f_0\|_{L^\infty}\)。利用速度矩估计和能量守恒： \[\frac{d}{dt} \iint \left(\frac{1}{2}|v|^2 f^\varepsilon + \frac{1}{2}|E^\varepsilon|^2\right) dx dv = 0\] 得到关于 \(\varepsilon\) 一致的速度矩界。通过速度平均引理（velocity averaging lemma，Golse-Lions-Perthame-Sentis 1988）：对输运方程的解，速度矩 \(\int f^\varepsilon \psi(v) dv\) 在 \(H^s\) 中紧（\(s \in (0, 1/2)\)）。这给出 \(\{f^\varepsilon\}\) 的强紧性（不计 \(\varepsilon\)）。取极限 \(\varepsilon \to 0\)，得到 \(f \in L^\infty([0, T]; L^1 \cap L^\infty)\) 且 \(\iint |v|^2 f < \infty\)。在分布意义下，\(f\) 满足 \((\partial_t + v \cdot \nabla_x + E \cdot \nabla_v) \beta(f) = 0\)（对所有 \(\beta \in C^1\)）。由重正规化技巧（DiPerna-Lions），此即重正规化解。若 \(f_0 \in C^1\)，则解为经典解。\(\blacksquare\)

\textbf{定理 373.2}（Pfaffelmoser 经典解，1992）：对 \(C^1\) 且紧支撑的初值 \(f_0\)，Vlasov-Poisson 系统在 \(d=3\) 时存在全局经典解。Pfaffelmoser 证明的关键在于沿特征线估计速度的最大增长，证明了 \(|V(t)| \leq C(1+t)^{1+\varepsilon}\)，从而防止了有限时间爆破。

\textbf{证明}：由特征线法，\(f(t, X(t), V(t)) = f_0(x, v)\) 沿特征线 \(\dot{X} = V\)，\(\dot{V} = E(t, X)\) 为常数。电场 \(E(t, x) = \frac{x}{|x|^3} * \rho(t, x)\)（Coulomb 力）。关键估计：对任意粒子，\(|V(t)| \leq C (1+t)^{2/3} \log(2+t)\)。Pfaffelmoser 将相空间按速度大小分层，对每层粒子分别估计其对电场的贡献。设 \(f_0\) 支撑在 \(|v| \leq R\) 内。对 \(t\) 时刻速度 \(|V(t)| \geq R\) 的粒子，其对应的初始位置在 \(|x| \leq R t + O(1)\) 内。将此区域的电荷分布分解为''近场''和''远场''贡献，利用 Poisson 方程和 Newton 位势的衰减性质，证明 \(|E(t, X(t))| \leq C (1+t)^{-2/3+o(1)}\)。积分 \(\dot{V} = E\) 得 \(|V(t)| \leq C (1+t)^{1/3+o(1)}\)，更精细的 bootstrap 论证将此改进为 \(|V(t)| \leq C (1+t)^{1+\varepsilon}\)（对任意 \(\varepsilon > 0\)）。由于特征线全局存在，解 \(f\) 全局存在且保持 \(C^1\) 正则性。\(\blacksquare\)

\subsection{373.3 Landau阻尼的线性理论}\label{landauux963bux5c3cux7684ux7ebfux6027ux7406ux8bba}

\textbf{定理 373.3}（线性 Landau 阻尼，Landau 1946）：考虑空间均匀平衡态 \(f_0(v)\)（\(d=1\) 情形）上的线性化 Vlasov-Poisson 系统。对扰动 \(f_1(t,x,v) = g(t,v)e^{ikx}\)，电场 \(E(t,x) = \hat{E}(t)e^{ikx}\)。由 Laplace 变换分析，当 \(t\to\infty\) 时， \[\hat{E}(t) \sim C e^{-i\omega_p t - \gamma t}, \quad \gamma = \frac{\pi}{2}\frac{\omega_p^3}{k^2} f_0'\left(\frac{\omega_p}{k}\right)\] 其中 \(\omega_p\) 为等离子体频率，由色散关系 \(\varepsilon(k,\omega)=1-\frac{\omega_p^2}{k^2}\int \frac{f_0'(v)}{v-\omega/k}dv=0\) 确定。若 \(f_0'(\omega_p/k)<0\)（Penrose 稳定性条件），则 \(\gamma>0\)，电场呈指数衰减------此即 Landau 阻尼。

\textbf{Landau 阻尼的物理机制}：衰减源于波-粒子共振导致的相位混合（phase mixing）。速度为 \(v \approx \omega/k\) 的粒子与波共振，从波中吸收能量。若速度分布 \(f_0\) 在共振速度处具有负斜率（\(f_0'(\omega/k) < 0\)），则加速的粒子多于减速的粒子，粒子从波中净吸收能量，波衰减。反之（\(f_0'(\omega/k) > 0\)），则发生 Landau 增长（逆 Landau 阻尼，如束-等离子体不稳定性）。这一过程在无碰撞系统中发生，体现了 Vlasov 方程的不可逆行为------尽管方程本身是时间可逆的（Liouville 定理），但线性化系统的连续谱导致扰动在弱意义下衰减。

\textbf{证明}：线性化 Vlasov-Poisson 系统：\(\partial_t f_1 + v\cdot\nabla_x f_1 + E\cdot\nabla_v f_0 = 0\)，\(\nabla\cdot E = \int f_1 dv\)。Fourier-Laplace 变换（对 \(t\) 作 Laplace 变换，对 \(x\) 作 Fourier 变换）得 \(f_1\) 的表达式，代入 Poisson 方程得色散关系 \(\varepsilon(k,\omega) = 1 - \frac{1}{k^2}\int \frac{k\cdot\nabla_v f_0}{\omega - k\cdot v} dv = 0\)。对稳定的 \(f_0\)（单峰分布），\(\varepsilon(k,\omega)\) 在上半平面解析且在实轴下方有零点，对应 \(\operatorname{Im}\omega = -\gamma < 0\)。Laplace 逆变换的围道积分通过变形至实轴下方，留数贡献给出 \(E(t) \sim e^{-i\omega_r t - \gamma t}\)。\(\blacksquare\)

\textbf{定理 373.4}（Mouhot-Villani 非线性 Landau 阻尼，2011）：对充分光滑（Gevrey 类）且小振幅的初值扰动，非线性 Vlasov-Poisson 系统（\(d \geq 1\)，周期边界条件或全空间）的解在 \(t\to\infty\) 时以几乎多项式速率（\(O(t^{-N})\) 对任意 \(N\)）衰减至空间均匀态，验证了非线性稳定性和 Landau 阻尼在非线性水平上的严格成立。

\textbf{证明}：Mouhot-Villani 使用 Newton 迭代和正则化方法，分三个核心步骤。(1) 线性 Landau 阻尼：在线性层面，对空间均匀平衡态 \(f_0(v)\) 的扰动，Vlasov 方程的线性化算子 \(L = -v\cdot\nabla_x - E[f_1]\cdot\nabla_v f_0\) 具有连续谱。Gevrey 类正则性假设下，扰动 \(f_1(t,x,v)\) 的 Fourier 模 \(\hat{f}_1(k,v,t)\) 满足 \(\hat{f}_1(k,v,t) \sim e^{-ik v t} \hat{f}_1(k,v,0) + \text{Landau 阻尼项}\)，后者以 \(e^{-\lambda|k|^{1/(1+s)} t}\) 衰减（\(s\) 为 Gevrey 指数）。在线性层面，这给出指数衰减。(2) 非线性响应的级数展开：将非线性解写为 \(f = f_0 + \sum_{n\geq 1} f^{(n)}\)，其中 \(f^{(n)}\) 为 \(n\) 阶非线性贡献。关键挑战是控制等离子体回波（plasma echo）：不同波数 \(k, k'\) 间的非线性相互作用产生共振项 \(\sim e^{i(k-k')v t}\)，这些共振项可能随 \(t\) 增长，破坏衰减。Mouhot-Villani 证明，在 Gevrey 正则性下，不同波数间回波的叠加相互抵消------即回波链（echo chain）的估计在长时间尺度上保持可控。这一抵消机制是通过分析 Fourier 模式在复平面中的解析延拓实现的，其中不同波数间的非线性相互作用导致围道积分中的相消。(3) 在 Gevrey 类函数空间中建立一致估计：通过 Newton 迭代，在适当的加权范数下证明解的全局存在性，且扰动以几乎多项式速率 \(O(t^{-N})\) 衰减至零。关键是将非线性项视为线性化问题的扰动，利用线性估计的指数衰减在迭代中控制非线性项的增长。\(\blacksquare\)

\subsection{373.4 引力Vlasov-Poisson系统与稳定性}\label{ux5f15ux529bvlasov-poissonux7cfbux7edfux4e0eux7a33ux5b9aux6027}

\textbf{定义 373.5}（引力 Vlasov-Poisson 系统与 Jeans 不稳定性）：引力情形（\(\gamma_d = -4\pi G\)）的平衡态为球对称静态解 \(f_0(r, v) = F(E, L)\)（\(E = |v|^2/2 + \Phi_0(r)\) 为能量，\(L = |x\times v|\) 为角动量），满足自洽 Poisson 方程。引力系统的稳定性与等离子体有本质不同：由于引力总是吸引的，系统可能经历 Jeans 不稳定性------密度扰动在引力作用下增长，导致引力塌缩。引力 Vlasov-Poisson 系统的非线性稳定性仍是开放问题，与星系动力学和暗物质晕的形成密切相关。

\hrulefill

\emph{重正化群从消除紫外发散的技术需求出发，经由 BPHZ 定理的严格森林公式和 \(R\)-操作，建立了微扰量子场论重整化的数学基础。Callan-Symanzik 方程和 \(\beta\)-函数描述了跑动耦合常数，渐近自由（Gross-Wilczek-Politzer）为 QCD 提供了理论基础。Wilson 重正化群从 Kadanoff 分块自旋变换出发，将临界现象的普适性与 RG 不动点理论联系起来，精确 RG 方程则为非微扰计算提供了系统框架。Vlasov 方程方面，DiPerna-Lions 的重正规化解理论和速度平均引理建立了弱解的存在性，Pfaffelmoser 给出了经典解的全局结果。Landau 阻尼的线性理论揭示了无碰撞等离子体中的不可逆现象，Mouhot-Villani 的非线性 Landau 阻尼定理（2011）是此领域里程碑式的突破。}

\emph{卷三十一：数学物理。}

\newpage

\chapter{13-数学物理/数学物理/07-量子信息.md}\label{ux6570ux5b66ux7269ux7406ux6570ux5b66ux7269ux740607-ux91cfux5b50ux4fe1ux606f.md}

\section{第462章：量子力学的信息论公理}\label{ux7b2c462ux7ae0ux91cfux5b50ux529bux5b66ux7684ux4fe1ux606fux8bbaux516cux7406}

量子力学与信息论的融合是 20 世纪末最激动人心的科学进展之一。将量子态视为信息载体------而非仅仅是被动的物理系统状态------这一视角转换使得 Shannon 信息论的经典概念（熵、信道容量、编码理论）以全新的、量子化的形式重新浮现。本章从量子力学的信息论公理出发，建立了量子信息论的基础框架。量子态由 Hilbert 空间 \(\mathcal{H}\) 上的密度算子 \(\rho\) 描述（\(\rho\geq 0\)，\(\operatorname{Tr}(\rho)=1\)），纯态对应投影算子 \(|\psi\rangle\langle\psi|\)，混合态是纯态的凸组合。复合系统的张量积结构 \(\mathcal{H}_A\otimes\mathcal{H}_B\) 是量子信息论区别于经典信息论的根本数学特征------它直接导致了量子纠缠、不可克隆定理和量子隐形传态等经典理论中完全不存在的新现象。量子测量由 POVM（正算子值测度）\(\{E_m\}\) 描述，测量结果 \(m\) 的概率为 \(\operatorname{Tr}(\rho E_m)\)------这自然地推广了 von Neumann 的投影测量。Bloch 球表示将单量子比特 \(\rho=(I+\mathbf{r}\cdot\boldsymbol{\sigma})/2\)（\(|\mathbf{r}|\leq 1\)）参数化为三维球体，直观展示了量子叠加原理相对于经典概率（一根线段）的高维自由度。量子隐形传态（Bennett et al.~1993）则展示了 EPR 纠缠对和经典通信如何传输完整量子态，且不违反不可克隆定理。

\subsection{262.1 量子态与密度矩阵}\label{ux91cfux5b50ux6001ux4e0eux5bc6ux5ea6ux77e9ux9635}

\textbf{定义 262.1}（量子态 / Quantum State）：量子态由单位向量 \(|\psi\rangle \in \mathcal{H}\)（Hilbert 空间）表示（模相位：\(|\psi\rangle\) 和 \(e^{i\theta}|\psi\rangle\) 表示同一态）。\(\dim \mathcal{H} = d\)。

\textbf{定义 262.2}（密度矩阵 / Density Matrix）：\textbf{密度算子} \(\rho\) 是 \(\mathcal{H}\) 上的正半定迹为 1 的算子： \[\rho \geq 0, \quad \operatorname{Tr}(\rho) = 1\]

纯态：\(\rho = |\psi\rangle\langle\psi|\)；混合态：\(\rho = \sum_i p_i |\psi_i\rangle\langle\psi_i|\)（\(p_i > 0\)，\(\sum p_i = 1\)）。

\textbf{定义 262.3}（复合系统与张量积）：两个系统 \(\mathcal{H}_A\) 和 \(\mathcal{H}_B\) 的复合系统为 \(\mathcal{H}_A \otimes \mathcal{H}_B\)。子系统 \(\rho_A\) 的\textbf{约化密度矩阵}为 \(\rho_A = \operatorname{Tr}_B(\rho_{AB})\)（部分迹）。

\subsection{262.2 量子测量}\label{ux91cfux5b50ux6d4bux91cf}

\textbf{定义 262.4}（投影测量 / von Neumann 测量）：由可观测量的谱分解 \(M = \sum_m m P_m\) 定义。测量结果 \(m\) 的概率为 \(\operatorname{Tr}(\rho P_m)\)，测量后状态退相为 \(P_m \rho P_m / \operatorname{Tr}(\rho P_m)\)。

\textbf{定义 262.5}（POVM / 正算子值测度）：\textbf{POVM} 是一组正算子 \(\{E_m\}\) 满足 \(\sum_m E_m = I\)。结果 \(m\) 的概率为 \(\operatorname{Tr}(\rho E_m)\)。POVM 推广了投影测量，适用于非理想测量和信息最优获取。

\textbf{定理 262.1}（Naimark 扩张定理）：任何 POVM 可由更大 Hilbert 空间上的投影测量实现。即 POVM 测量等价于扩展系统后的投影测量。

\textbf{证明}：设 \(\{E_m\}_{m=1}^M\) 为 \(\mathcal{H}_A\) 上的 POVM（\(E_m \geq 0\)，\(\sum E_m = I\)）。构造扩展空间 \(\mathcal{H}_A \otimes \mathcal{H}_B\)（\(\dim\mathcal{H}_B = M\)，\(\{|m\rangle\}\) 为基）。定义等距映射 \(V: \mathcal{H}_A \to \mathcal{H}_A\otimes\mathcal{H}_B\) 为 \(V|\psi\rangle = \sum_m (\sqrt{E_m}|\psi\rangle)\otimes|m\rangle\)，则 \(V^\dagger V = \sum E_m = I\)。定义投影测量 \(\tilde{P}_m = I\otimes|m\rangle\langle m|\)。对任意态 \(\rho_A\)，\(\operatorname{Tr}(\rho_A E_m) = \operatorname{Tr}(V\rho_A V^\dagger \tilde{P}_m)\)。故 POVM 测量 \(\{E_m\}\) 等价于扩展系统上的投影测量 \(\{\tilde{P}_m\}\)。\(\blacksquare\)

\subsection{262.3 量子比特与 Bloch 球}\label{ux91cfux5b50ux6bd4ux7279ux4e0e-bloch-ux7403}

\textbf{定义 262.6}（量子比特 / Qubit）：\(\mathcal{H} = \mathbb{C}^2\)。纯态 \(|\psi\rangle = \alpha|0\rangle + \beta|1\rangle\)（\(|\alpha|^2 + |\beta|^2 = 1\)）。任意单量子比特密度矩阵可写为

\[\rho = \frac{1}{2}(I + \mathbf{r} \cdot \boldsymbol{\sigma})\]

其中 \(\mathbf{r} \in \mathbb{R}^3\)（Bloch 向量）满足 \(|\mathbf{r}| \leq 1\)，\(\boldsymbol{\sigma} = (\sigma_x, \sigma_y, \sigma_z)\) 为 Pauli 矩阵。纯态对应 \(|\mathbf{r}| = 1\)（Bloch 球面）。

Bloch 球表示是单量子比特几何化的核心工具。Pauli 矩阵 \(\sigma_x, \sigma_y, \sigma_z\) 连同单位矩阵 \(I\) 构成 \(\mathbb{C}^{2 \times 2}\) 上 Hermite 矩阵空间的一组基，而密度矩阵 \(\rho\) 正是一个迹为 1 的半正定 Hermite 矩阵，因此自然地由四维向量 \((1, \mathbf{r})\) 参数化。\(\rho \geq 0\) 的条件等价于 \(|\mathbf{r}| \leq 1\)：纯态满足 \(\rho^2 = \rho\)，即 \(\mathbf{r}\) 在单位球面上；混合态（\(\rho^2 \neq \rho\)）对应球内部的点。Bloch 球面上的任意酉演化 \(U\) 对应于 \(\mathbf{SO}(3)\) 中的旋转 \(R\)：\(U\rho U^\dagger = \frac{1}{2}(I + (R\mathbf{r}) \cdot \boldsymbol{\sigma})\)，其中 \(R\) 由 \(U\) 通过 \(SU(2) \to SO(3)\) 的覆盖同态确定。这一几何对应极大地简化了单量子比特门（如 \(X, Y, Z, H\) 门）的可视化：\(X\) 门对应绕 \(x\) 轴旋转 \(\pi\)，\(H\) 门对应绕 \((x+z)/\sqrt{2}\) 轴旋转 \(\pi\)。Bloch 球表示还揭示了量子力学与经典概率的根本区别：经典比特的随机态位于一维线段上（概率 \(p\) 参数化），而量子比特的态位于三维球体中，其额外的维度正是量子力学的叠加原理和干涉效应的几何来源。

\subsection{262.4 量子隐形传态的代数结构}\label{ux91cfux5b50ux9690ux5f62ux4f20ux6001ux7684ux4ee3ux6570ux7ed3ux6784}

量子隐形传态（Bennett et al.~1993）利用纠缠态和经典通信传输量子态，其代数结构可表述为 Bell 基测量和 Pauli 校正的对应关系。

\textbf{协议步骤}：方 \(A\) 欲向方 \(B\) 传输未知态 \(|\psi\rangle = \alpha|0\rangle + \beta|1\rangle\)。双方预先共享 EPR 对 \(|\Phi^+\rangle = \frac{1}{\sqrt{2}}(|00\rangle + |11\rangle)\)。总初态为：

\[|\psi\rangle_A \otimes |\Phi^+\rangle_{AB} = (\alpha|0\rangle + \beta|1\rangle)_A \otimes \frac{1}{\sqrt{2}}(|00\rangle + |11\rangle)_{AB}\]

\textbf{步骤 1}：方 \(A\) 对持有的两个量子比特施加 CNOT 门（控制位为其未知态）和 Hadamard 门，将总态重写为 Bell 基的展开：

\[\frac{1}{2}\sum_{k=0}^{3} |\text{Bell}_k\rangle_A \otimes \sigma_k (\alpha|0\rangle + \beta|1\rangle)_B\]

其中 \(\sigma_0 = I\)、\(\sigma_1 = X\)、\(\sigma_2 = Z\)、\(\sigma_3 = XZ\) 为 Pauli 算子。

\textbf{步骤 2}：方 \(A\) 对自身两量子比特做 Bell 基测量，得到结果 \(k \in \{00, 01, 10, 11\}\)（2 比特经典信息）。测量后方 \(B\) 的量子比特塌缩为 \(\sigma_k|\psi\rangle\)。

\textbf{步骤 3}：方 \(A\) 通过经典信道将 \(k\) 发送给方 \(B\)。方 \(B\) 根据 \(k\) 施加相应的 Pauli 校正门 \(\sigma_k^\dagger\)，恢复出 \(|\psi\rangle\)。

\textbf{注}：隐形传态不违反不可克隆定理------方 \(A\) 手中的原始态在 Bell 测量中被破坏，信息通过经典信道传递，且协议消耗了一个 EPR 对。

\hrulefill

\hrulefill

\hrulefill

\hrulefill

\hrulefill

\hrulefill

\section{第463章：量子纠缠}\label{ux7b2c463ux7ae0ux91cfux5b50ux7ea0ux7f20}

量子纠缠（Quantum Entanglement）是量子力学中最具革命性的概念之一，Einstein、Podolsky 和 Rosen（EPR，1935）最早用来论证量子力学''不完备''的现象，在 John Bell（1964）的不等式之后转变为量子非局域性的标志性资源。纠缠的数学本质是张量积结构的不可分解性：对于复合系统 \(\mathcal{H}_A\otimes\mathcal{H}_B\)，纯态 \(|\psi\rangle_{AB}\) 是纠缠的当且仅当它不能写为直积态 \(|\phi\rangle_A\otimes|\chi\rangle_B\)；混合态 \(\rho_{AB}\) 是纠缠的当且仅当它不能表示为可分态的凸组合 \(\sum_i p_i\rho_A^{(i)}\otimes\rho_B^{(i)}\)。Peres-Horodecki 的 PPT 判据（1996）通过部分转置（仅对子系统 \(B\) 取转置）给出了最实用的纠缠检测方法：若 \(\rho_{AB}^{T_B}\) 非正半定，则 \(\rho_{AB}\) 纠缠。纠缠作为一种物理资源，可以在局部操作和经典通信（LOCC）下被操纵但不能被创造------LOCC 单调性是任何合理纠缠度量的必要条件。Bell-CHSH 不等式 \(|\langle AB\rangle+\langle AB'\rangle+\langle A'B\rangle-\langle A'B'\rangle|\leq 2\) 是纠缠非局域性的最简测试：量子纠缠态可违反此界（最大 \(2\sqrt{2}\)，Tsirelson 界），而任何经典局域隐变量理论必须满足该界。本章将系统引入纠缠的数学定义、判定条件、度量理论和非局域性测试。

\subsection{263.1 纠缠的定义与判定}\label{ux7ea0ux7f20ux7684ux5b9aux4e49ux4e0eux5224ux5b9a}

\textbf{定义 263.1}（可分态与纠缠态）：两体态 \(\rho_{AB}\) 是\textbf{可分的}（separable），如果可写为直积态的凸组合： \[\rho_{AB} = \sum_i p_i \rho_A^{(i)} \otimes \rho_B^{(i)}\]

否则 \(\rho_{AB}\) 是\textbf{纠缠的}。可分态对应经典关联，纠缠态是量子独有的。

\textbf{定理 263.1}（PPT 判据 / Peres-Horodecki 判据，1996）：若 \(\rho_{AB}\) 可分，则其部分转置 \(\rho_{AB}^{T_B}\)（仅对子系统 \(B\) 取转置）是正半定的。对 \(2 \otimes 2\) 和 \(2 \otimes 3\) 系统，PPT 条件是纠缠的必要充分条件。

\textbf{证明}：若 \(\rho_{AB}\) 可分，\(\rho_{AB} = \sum_i p_i \rho_A^{(i)} \otimes \rho_B^{(i)}\)。部分转置 \(\rho_{AB}^{T_B} = \sum_i p_i \rho_A^{(i)} \otimes (\rho_B^{(i)})^T\)。由于 \(\rho_B^{(i)}\) 是密度矩阵，其转置仍为正半定，故 \(\rho_{AB}^{T_B}\) 作为正半定算子的凸组合仍是正半定的。必要性得证。对 \(2\otimes 2\) 和 \(2\otimes 3\) 系统，Horodecki 家族证明了 PPT 也是纠缠的充分条件：若 \(\rho_{AB}^{T_B} \ngeq 0\)，则 \(\rho_{AB}\) 必纠缠。但对更高维系统存在 PPT 纠缠态（束缚纠缠），PPT 不再是充分条件。\(\blacksquare\)

\textbf{定义 263.2}（纠缠目击 / Entanglement Witness）：Hermitian 算子 \(W\) 满足 \(\operatorname{Tr}(W \rho_{\text{sep}}) \geq 0\)（对所有可分态），且存在纠缠态 \(\rho_{\text{ent}}\) 使 \(\operatorname{Tr}(W \rho_{\text{ent}}) < 0\)。\(W\)「目击」了 \(\rho_{\text{ent}}\) 的纠缠性。

\subsection{263.2 纠缠度量}\label{ux7ea0ux7f20ux5ea6ux91cf}

\textbf{定义 263.3}（纠缠熵 / Von Neumann 熵）：对纯态 \(|\psi\rangle_{AB}\)，\textbf{纠缠熵} \(E(|\psi\rangle) = S(\rho_A) = S(\rho_B)\)，其中 von Neumann 熵 \(S(\rho) = -\operatorname{Tr}(\rho \log \rho)\)。纠缠熵衡量了子系统混合的程度。

\textbf{定义 263.4}（形成纠缠 / Entanglement of Formation）：两体混合态 \(\rho_{AB}\) 的\textbf{形成纠缠}为 \[E_F(\rho) = \inf \left\{ \sum_i p_i S(\rho_A^{(i)}) : \rho = \sum_i p_i |\psi_i\rangle\langle\psi_i| \right\}\]

即分解为纯态的最小平均纠缠熵（为 \(2 \otimes 2\) 状态可解析计算，Wootters 1998）。

\textbf{定理 263.2}（纠缠的不可增性 / LOCC 单调性）：纠缠在\textbf{局部操作与经典通信}（LOCC）下不会增加。任何合理的纠缠度量必须是 LOCC 单调的。这是纠缠作为资源的基本性质。

\textbf{证明}：设 \(\Lambda\) 为 LOCC 操作，可分解为 \(\Lambda = \sum_k (A_k \otimes B_k)\) 其中 \(A_k, B_k\) 为局部量子操作。对任意纠缠度量 \(E\)，LOCC 单调性 \(E(\Lambda(\rho)) \leq E(\rho)\) 的验证：\(E(\Lambda(\rho)) = E(\sum_k A_k \otimes B_k \rho A_k^\dagger \otimes B_k^\dagger) \leq \sum_k p_k E(\rho_k)\)（凸性），其中 \(p_k = \operatorname{Tr}(A_k \otimes B_k \rho A_k^\dagger \otimes B_k^\dagger)\)，\(\rho_k = (A_k \otimes B_k) \rho (A_k \otimes B_k)^\dagger / p_k\)。LOCC 的局部性确保 \(E(\rho_k) \leq E(\rho)\)（局部酉变换不改变纠缠，经典通信不增加纠缠）。故 \(E(\Lambda(\rho)) \leq E(\rho)\)。\(\blacksquare\)

\subsection{263.3 Bell 不等式与非局域性}\label{bell-ux4e0dux7b49ux5f0fux4e0eux975eux5c40ux57dfux6027}

\textbf{定义 263.5}（Bell-CHSH 不等式 / Clauser-Horne-Shimony-Holt）：对二元测量 \(A, A'\)（方 \(A\) 端）和 \(B, B'\)（方 \(B\) 端），任何经典局域隐变量理论满足 \[|\langle AB \rangle + \langle AB' \rangle + \langle A'B \rangle - \langle A'B' \rangle| \leq 2\]

量子纠缠态可违反此界，最大值可达 \(2\sqrt{2}\)（Tsirelson 界）。

\textbf{定理 263.3}（Tsirelson 界）：量子力学中 CHSH 表达式的最大值为 \(2\sqrt{2}\)。这严格小于无信号条件下代数学的最大值 \(4\)（Popescu-Rohrlich 盒），表明量子非局域性比一般无信号理论弱。

\textbf{证明}：设 \(A, A'\) 和 \(B, B'\) 为 Hermitian 可观测量，满足 \(A^2 = A'^2 = B^2 = B'^2 = I\)。定义 CHSH 算子 \(\mathcal{B} = A\otimes B + A\otimes B' + A'\otimes B - A'\otimes B'\)。计算 \(\mathcal{B}^2 = 4I + [A,A']\otimes[B,B']\)。由 \([A,A'] = AA' - A'A\)，取算子范数 \(\|[A,A']\| \leq 2\|A\|\|A'\| = 2\)，同理 \(\|[B,B']\| \leq 2\)。故 \(\|\mathcal{B}^2\| \leq 4 + 4 = 8\)，即 \(\|\mathcal{B}\| \leq 2\sqrt{2}\)。因此对任意量子态 \(|\psi\rangle\)，\(|\langle\psi|\mathcal{B}|\psi\rangle| \leq 2\sqrt{2}\)。最大值 \(2\sqrt{2}\) 由 Bell 态 \(|\Phi^+\rangle\) 和适当选取的测量实现。\(\blacksquare\)

\hrulefill

\hrulefill

\hrulefill

\hrulefill

\hrulefill

\hrulefill

\section{第464章：量子信道与量子容量}\label{ux7b2c464ux7ae0ux91cfux5b50ux4fe1ux9053ux4e0eux91cfux5b50ux5bb9ux91cf}

量子信道是量子信息传输的数学抽象，它刻画了量子态在通过物理介质（如光纤、自由空间）传输时所经历的噪声和退相干过程。从数学结构上看，量子信道等价于完全正保迹（CPTP）映射 \(\mathcal{E}:\mathcal{B}(\mathcal{H}_A)\to\mathcal{B}(\mathcal{H}_B)\)------这是经典概率转移矩阵的量子推广，``完全正''性保证了即使信道作用于复合系统的一部分，整体映射仍保持量子态的物理合理性。Kraus 表示定理将 CPTP 映射参数化为算子和形式 \(\mathcal{E}(\rho)=\sum_k A_k\rho A_k^\dagger\)（\(\sum_k A_k^\dagger A_k=I\)），而 Stinespring 膨胀定理则揭示了更深的数学结构：任何 CPTP 映射都可以通过将系统与一个辅助环境耦合后进行酉演化，再对环境取部分迹来实现。量子信道容量理论（Holevo-Schumacher-Westmoreland，1996-1998）将 Shannon 信道编码定理推广到量子领域：经典容量 \(C(\mathcal{E})\) 由 Holevo 量 \(\chi(\mathcal{E})\) 的正则化极限给出，而量子容量 \(Q(\mathcal{E})\) 则由相干信息的正则化极限给出（Lloyd-Shor-Devetak 定理）。量子信道容量的非可加性是其与经典容量理论最显著的区别------Hastings（2009）证明了 Holevo 量的超可加性，意味着量子信道容量的计算本质上比经典情形困难得多。

\subsection{264.1 量子信道}\label{ux91cfux5b50ux4fe1ux9053}

\textbf{定义 264.1}（量子信道 / 完全正保迹映射 / CPTP Map）：量子信道 \(\mathcal{E}: \mathcal{B}(\mathcal{H}_A) \to \mathcal{B}(\mathcal{H}_B)\) 是线性映射，满足： - \textbf{保迹}（trace-preserving）：\(\operatorname{Tr}(\mathcal{E}(\rho)) = \operatorname{Tr}(\rho)\) - \textbf{完全正}（completely positive / CP）：对任意扩展系统 \(R\)，\(\mathcal{E} \otimes \operatorname{id}_R\) 保持正定性

\textbf{定理 264.1}（Kraus 表示 / 算子和表示）：任何 CPTP 映射 \(\mathcal{E}\) 可表示为 \[\mathcal{E}(\rho) = \sum_{k=1}^K A_k \rho A_k^\dagger\]

其中 Kraus 算子 \(A_k\) 满足 \(\sum_k A_k^\dagger A_k = I\)。

\textbf{证明}：由 Stinespring 膨胀，\(\mathcal{E}(\rho) = \operatorname{Tr}_E(V \rho V^\dagger)\)，其中 \(V: \mathcal{H}_A \to \mathcal{H}_B \otimes \mathcal{H}_E\) 为等距映射。取 \(\mathcal{H}_E\) 的标准正交基 \(\{|k\rangle\}\)，定义 \(A_k = \langle k|_E V\)（即 \(A_k|\psi\rangle = \langle k|_E V|\psi\rangle\)）。则 \(\mathcal{E}(\rho) = \sum_k \langle k|_E V \rho V^\dagger |k\rangle_E = \sum_k A_k \rho A_k^\dagger\)。完全正性：\(\mathcal{E} \otimes \operatorname{id}_R \geq 0\) 等价于 \(\sum_k A_k^\dagger A_k = I\)（保迹条件）。唯一性：Kraus 算子由酉变换 \(A_k' = \sum_j U_{kj} A_j\) 联系。\(\blacksquare\)

\textbf{定义 264.2}（Stinespring 膨胀定理）：任何 CPTP 映射 \(\mathcal{E}\) 可用更大大 Hilbert 空间中的酉演化 \(U\) 和环境初态 \(|0\rangle_E\) 实现：\(\mathcal{E}(\rho) = \operatorname{Tr}_E(U (\rho \otimes |0\rangle\langle 0|_E) U^\dagger)\)。这是量子信道的纯化表示。

\subsection{264.2 典型量子信道}\label{ux5178ux578bux91cfux5b50ux4fe1ux9053}

\textbf{定义 264.3}（常见量子噪声信道）： - \textbf{去极化信道}（depolarizing）：\(\mathcal{E}(\rho) = (1-p)\rho + \frac{p}{d} I\) - \textbf{幅值阻尼信道}（amplitude damping）：模拟 \(|1\rangle \to |0\rangle\) 的衰变过程 - \textbf{相位阻尼信道}（dephasing）：消灭非对角元，模拟信息丢失：\(\mathcal{E}(\rho) = (1-p)\rho + p \operatorname{diag}(\rho)\)

\subsection{264.3 量子信道容量}\label{ux91cfux5b50ux4fe1ux9053ux5bb9ux91cf}

\textbf{定义 264.4}（量子信道容量的分类）： - \textbf{经典容量} \(C(\mathcal{E})\)：通过信道可靠传输经典信息的最大速率 - \textbf{量子容量} \(Q(\mathcal{E})\)：通过信道可靠传输量子信息的最大速率 - \textbf{私密容量} \(P(\mathcal{E})\)：信道可用于秘密共享的容量 - \textbf{纠缠辅助容量} \(C_E(\mathcal{E})\)：收发双方共享无限纠缠时传输经典信息的容量

\textbf{定理 264.2}（Holevo-Schumacher-Westmoreland 定理 / HSW 定理，1996-1998）：量子信道 \(\mathcal{E}\) 的经典容量为 \[C(\mathcal{E}) = \lim_{n \to \infty} \frac{1}{n} \chi(\mathcal{E}^{\otimes n})\]

其中 \textbf{Holevo 量} \(\chi(\mathcal{E}) = \max_{\{p_i, \rho_i\}} [S(\mathcal{E}(\sum_i p_i \rho_i)) - \sum_i p_i S(\mathcal{E}(\rho_i))]\)。

\textbf{证明}：可达性：构造随机码 \(\{p_i, \rho_i^{\otimes n}\}\)（典型序列方法），接收方使用典型子空间上的 POVM 测量。由量子典型序列的渐近等分性质（AEP），输出态 \(\mathcal{E}^{\otimes n}(\rho_i^{\otimes n})\) 的典型子空间几乎正交，可实现 \(n\chi(\mathcal{E})\) 速率。逆定理：由 Holevo 界，对任意编码，可访问信息 \(\leq \chi(\mathcal{E}^{\otimes n})\)。取 \(n \to \infty\) 的正则化极限即得容量公式。\(\blacksquare\)

\textbf{注 264.1}（Holevo 界的简要推导）：对于发送方以概率 \(p_i\) 制备状态 \(\rho_i\)、接收方通过 POVM \(\{E_m\}\) 测量的经典通信场景，设输入随机变量为 \(X\)，测量结果为 \(Y\)。经典互信息 \(I(X:Y)\) 受限于 Holevo 界： \[\chi = S\left(\sum_i p_i \rho_i\right) - \sum_i p_i S(\rho_i)\] 推导关键：由量子相对熵的单调性（Lindblad-Uhlmann 定理），对任意 POVM 测量，经典互信息不超过 Holevo 量；再由 HSW 定理，该上界在 \(n \to \infty\) 的意义下可由适当的编码方案渐近达到。Holevo 界的本质是 \(S(\rho) - \sum_i p_i S(\rho_i)\) 度量了''量子态集合的可区分性''------它是量子信息论中经典容量的单次使用上界，比经典 Shannon 信道容量多了 von Neumann 熵的凹性修正项。当输入为纯态且相互正交时，Holevo 量达到最大值 \(\log d\)。∎

\textbf{定理 264.3}（Lloyd-Shor-Devetak 定理 / LSD 定理）：量子信道的量子容量为相干信息的正则化极限： \[Q(\mathcal{E}) = \lim_{n \to \infty} \frac{1}{n} I_c(\mathcal{E}^{\otimes n})\]

其中\textbf{相干信息} \(I_c(\mathcal{E}) = \max_{\rho} [S(\mathcal{E}(\rho)) - S(\mathcal{E} \otimes \operatorname{id}_R(|\psi\rangle\langle\psi|_{\rho R}))]\)。

\textbf{证明}：可达性：使用量子纠错码将量子信息编码为 \(\mathcal{E}^{\otimes n}\) 的输入。相干信息 \(I_c(\mathcal{E}^{\otimes n})\) 度量了纠缠在信道传输中的保持量，通过纠缠蒸馏和量子隐形传态实现量子通信。逆定理：由量子数据处理不等式，在 \(n\) 次使用后，量子容量 \(\leq \frac{1}{n} I_c(\mathcal{E}^{\otimes n})\)。取 \(n \to \infty\) 即得正则化表达式。\(\blacksquare\)

\hrulefill

\hrulefill

\hrulefill

\hrulefill

\hrulefill

\hrulefill

\section{第465章：量子算法}\label{ux7b2c465ux7ae0ux91cfux5b50ux7b97ux6cd5}

量子算法是量子信息论中最为大众所熟知的部分------它利用量子力学特有的资源（叠加、干涉和纠缠）来实现对经典算法的指数或多项式加速。量子算法的数学模型是量子电路：\(n\) 个量子比特初始制备为 \(|0^n\rangle\)，依次施加一系列酉门（Hadamard 门 \(H\)、相位门 \(T\)、CNOT 等），最后由量子测量提取结果。Deutsch-Jozsa 算法（1992）率先展示了量子算法相对于经典确定型算法的威力------判断一个黑白明信片函数是否常值仅需一次量子 Oracle 查询（经典需 \(2^{n-1}+1\) 次）。Peter Shor 的因式分解算法（1994）是量子计算理论中最具革命性的成果：利用量子 Fourier 变换求周期，将 \(N\) 位整数的因式分解从经典最优的亚指数复杂度降为 \(O((\log N)^3)\) 多项式时间。Shor 算法的核心洞察在于：模指数函数 \(f(x)=a^x\bmod N\) 的周期 \(r\) 通过 QFT 的高效周期求取可以在多项式时间内获得，进而用 \(\gcd(a^{r/2}\pm 1,N)\) 提取因子。Lov Grover 的搜索算法（1996）则为非结构化搜索提供了二次加速（\(O(\sqrt{N})\) 而非 \(O(N)\)），且该加速已被证明是最优的（BBBV 定理）。

\subsection{265.1 量子计算模型}\label{ux91cfux5b50ux8ba1ux7b97ux6a21ux578b}

\textbf{定义 265.1}（量子电路模型）：\(n\) 个量子比特初始化为 \(|0^n\rangle\)。应用酉门序列（\(H\)、\(T\)、CNOT 等通用门集合），最后测量得到结果。量子电路是经典电路对应的量子推广。

\textbf{定义 265.2}（通用量子门集合 / quantum gate universality）： \[\{H, T, \operatorname{CNOT}\}\]

（\(H\) 为 Hadamard 门，\(T = \operatorname{diag}(1, e^{i\pi/4})\)，CNOT 为受控非门）是近似通用的------任何 \(n\) 量子比特酉变换可被 \(O(4^n \log(1/\varepsilon))\) 个门 \(\varepsilon\)-逼近（Solovay-Kitaev 定理）。

\textbf{定义 265.3}（量子 Fourier 变换 / QFT）：\(n\) 量子比特的量子 Fourier 变换 \(\operatorname{QFT}_N\)（\(N=2^n\)）： \[|j\rangle \mapsto \frac{1}{\sqrt{N}} \sum_{k=0}^{N-1} e^{2\pi i jk / N} |k\rangle\]

QFT 可在 \(O(n^2)\) 个门内实现（经典 FFT 需 \(O(n 2^n)\)）。

\subsection{265.2 Shor 算法与 Grover 算法}\label{shor-ux7b97ux6cd5ux4e0e-grover-ux7b97ux6cd5}

\textbf{定理 265.1}（Shor 的因式分解算法）：存在量子算法在 \(O((\log N)^3)\) 时间内分解 \(N\) 位整数（以高概率）。核心步骤：用量子 Fourier 变换求模 \(N\) 指数 \(f(x) = a^x \bmod N\) 的周期（量子周期求取），然后用数论从周期推因子。

\textbf{证明}：随机选取 \(a\)（\(1 < a < N\)，\(\gcd(a, N) = 1\)）。若 \(\gcd(a, N) > 1\) 则已完成。考虑 \(f(x) = a^x \bmod N\)，其周期 \(r\) 满足 \(a^r \equiv 1 \pmod{N}\)。用量子 Fourier 变换求 \(r\)：制备叠加态 \(\frac{1}{\sqrt{M}}\sum_{x=0}^{M-1}|x\rangle|0\rangle\)（\(M = 2^{2\lceil\log_2 N\rceil}\)），Oracle 计算 \(f(x)\)，得 \(\frac{1}{\sqrt{M}}\sum_x |x\rangle|f(x)\rangle\)。对第一寄存器做 QFT，测量得 \(s\) 使 \(s/M \approx k/r\)。连分式展开恢复 \(r\)（以高概率）。若 \(r\) 为偶数且 \(a^{r/2} \not\equiv -1 \pmod{N}\)，则 \(\gcd(a^{r/2}\pm 1, N)\) 给出 \(N\) 的非平凡因子。QFT 需 \(O((\log N)^2)\) 门，连分式 \(O((\log N)^3)\) 经典运算，总体 \(O((\log N)^3)\)。\(\blacksquare\)

\textbf{定理 265.2}

\textbf{证明}：在\(N\)个未排序元素中搜索标记元素。制备初态\(|s\rangle = H^{\otimes n}|0\rangle^{\otimes n} = \frac{1}{\sqrt{N}}\sum_{x=0}^{N-1}|x\rangle\)。定义Oracle \(O_f|x\rangle = (-1)^{f(x)}|x\rangle\)（\(f(x)=1\)当且仅当\(x\)为标记元素）。Grover迭代\(G = (2|s\rangle\langle s| - I)O_f\)可视为二维平面上的旋转：\(|s\rangle = \sin\theta|\omega\rangle + \cos\theta|\bar{\omega}\rangle\)（\(\theta = \arcsin(1/\sqrt{N})\)，\(|\omega\rangle\)为标记态均匀叠加）。每次迭代将态向\(|\omega\rangle\)旋转\(2\theta\)角度。经\(k = \lfloor \pi/(4\theta) \rfloor \approx \frac{\pi}{4}\sqrt{N}\)次迭代后，振幅\(\sin((2k+1)\theta) \approx 1\)，测量以高概率得到标记元素。因此查询复杂度为\(O(\sqrt{N})\)，比经典\(O(N)\)有二次加速。\(\blacksquare\)

\textbf{定理 265.3}（BBBV 定理 / 量子查询下界，Bennett-Bernstein-Brassard-Vazirani 1997）：量子计算机搜索 \(N\) 个元素至少需要 \(\Omega(\sqrt{N})\) 次查询。Grover 算法是最优的。

\textbf{证明}：设 Oracle \(O_f|x\rangle = (-1)^{f(x)}|x\rangle\)。量子算法 \(T\) 次查询后状态为 \(|\psi_T\rangle = U_T O_f U_{T-1} O_f \cdots U_1 O_f U_0|0\rangle\)。定义 \(|\psi_T^{(0)}\rangle\) 为在 \(f \equiv 0\) 下的演化状态。由多项式方法（Beals et al.），\(p_T = \||\psi_T\rangle - |\psi_T^{(0)}\rangle\|^2\) 是 \(f\) 上的多项式，次数 \(\leq 2T\)。对搜索问题，\(p_T\) 必须在 \(N\) 个输入上变化显著，由多项式下界知 \(T \geq \Omega(\sqrt{N})\)。\(\blacksquare\)

\subsection{265.3 量子复杂性类}\label{ux91cfux5b50ux590dux6742ux6027ux7c7b}

\textbf{定义 265.4}（量子复杂性类）： - \textbf{BQP}：被多项式时间量子电路以有界误差判定的语言类 - \textbf{QMA}：量子模拟 NP------多项式时间量子验证者可验证的「量子证明」

\textbf{定理 265.4}（BQP 的位置）：BPP \(\subseteq\) \textbf{BQP} \(\subseteq\) PSPACE（Bernstein-Vazirani 1993）。已知 BQP 包含因式分解和离散对数问题，但不知道 BQP 是否包含 NP 完全问题。

\textbf{证明}：BPP \(\subseteq\) BQP：经典随机电路可由量子电路模拟（Hadamard 门生成叠加态，测量模拟随机性）。BQP \(\subseteq\) PSPACE：量子电路振幅 \(\langle 0|U|0\rangle\) 可表示为多项式个矩阵乘积的矩阵元，每个矩阵元是有限精度的代数数，由 Feynman 路径积分求和，可用 PSPACE 算法计算（递归展开矩阵乘积，深度 \(O(\log n)\) 的递归需要多项式空间）。具体地，BQP 算法接受概率 \(p_{\text{acc}} = \sum_{x \in A} |\langle x|U|0\rangle|^2\) 可在 PSPACE 内计算至任意精度。\(\blacksquare\)

\subsection{量子信息论的前沿与展望}\label{ux91cfux5b50ux4fe1ux606fux8bbaux7684ux524dux6cbfux4e0eux5c55ux671b}

量子信息论是 20 世纪末和 21 世纪初诞生的新兴交叉学科，它将信息论的基本概念（熵、信道容量、纠错）推广到量子力学框架，并利用量子叠加和量子纠缠等量子资源实现超越经典物理界限的信息处理能力。\textbf{量子纠错码}（Shor 1995, Steane 1996）是量子信息论中最深刻的成就之一：尽管量子态不可克隆（no-cloning theorem），但通过将单个逻辑量子比特编码到多个物理量子比特的纠缠态中，可以检测和纠正量子错误而不破坏量子信息。CSS 码（Calderbank-Shor-Steane）和稳定子码（stabilizer codes）将量子纠错码与经典纠错码（特别是自正交线性码）联系起来，其代数结构由 Pauli 群的稳定子子群刻画。量子纠错码的存在性由\textbf{量子 Hamming 界}（quantum Hamming bound）、\textbf{量子 Gilbert-Varshamov 界}和\textbf{量子 Singleton 界}（\(n - k \geq 2(d-1)\)）控制。\textbf{容错量子计算}（fault-tolerant quantum computation）的阈值定理（Aharonov-Ben-Or 1999, Knill-Laflamme-Zurek 1998）证明：若每个物理量子门的错误率低于某个阈值（约 \(10^{-4}\) 至 \(10^{-2}\)，取决于纠错方案），则通过级联纠错码（concatenated codes）可以在任意长的计算中抑制错误，实现可靠的量子计算。\textbf{拓扑量子计算}（Kitaev 2003）是非阿贝尔任意子（non-Abelian anyons）的编织（braiding）操作实现容错量子门的方法，其容错性由拓扑保护自然保证，无需主动纠错。

\textbf{量子密钥分发}（QKD）是量子信息论最成熟的应用：BB84 协议（Bennett-Brassard 1984）利用量子态不可克隆原理和测量扰动实现信息论安全的密钥分发，其安全性由 Csiszar-Korner 定理和 Shor-Preskill 证明（2000）严格保证。\textbf{量子隐形传态}（quantum teleportation, Bennett et al.~1993）利用 EPR 纠缠对和经典通信，将一个量子比特的完整量子态从 A 传输到 B，消耗一对纠缠和两比特经典通信。\textbf{超密编码}（superdense coding, Bennett-Wieser 1992）则利用纠缠实现每量子比特传输两比特经典信息。\textbf{量子复杂性理论}中，BQP（有界误差量子多项式时间）与经典复杂性类的关系是量子计算理论的核心问题------已知 BPP \(\subseteq\) BQP \(\subseteq\) PSPACE，但 BQP 与 NP 的关系尚不清楚。\textbf{量子优势}（quantum supremacy）实验（Google 2019, USTC 2020）在特定采样问题上展示了量子处理器超越经典最强超级计算机的计算能力，标志量子计算进入了 NISQ（含噪声中等规模量子）时代。 \hrulefill

\section{第466章：量子纠错与容错量子计算}\label{ux7b2c466ux7ae0ux91cfux5b50ux7ea0ux9519ux4e0eux5bb9ux9519ux91cfux5b50ux8ba1ux7b97}

量子纠错是使大规模量子计算从理论可能性走向物理可实现性的关键数学技术。与经典纠错码仅需纠正位翻转错误不同，量子纠错面临三重独特困难：不可克隆定理禁止复制量子态、错误是连续的（相位的任意旋转）、测量会破坏量子信息。Peter Shor 在 1995 年突破了这三重障碍，构造了首个 9 量子比特量子纠错码 \([[9,1,3]]\)，其核心思想是将量子纠错转化为错误算子的检测和校正------通过将错误离散化为 Pauli 群 \(\{I, X, Y, Z\}\) 的作用，连续错误被投影到 Pauli 错误上。Knill-Laflamme 量子纠错条件（1997）给出了错误算子集 \(\{E_a\}\) 可被纠正的充要条件：\(PE_a^\dagger E_bP=\lambda_{ab}P\)，其中 \(P\) 是码空间的投影算子------这是经典纠错条件 \(\mathbf{1}^TH E_a^TE_b\mathbf{1}=\lambda_{ab}\) 的量子推广。稳定子码（Gottesman 1997）将量子纠错码与 Pauli 群的 Abel 子群（稳定子）联系起来，CSS 构造（Calderbank-Shor-Steane）进一步将量子码与经典自正交线性码 \(C_2^\perp\subseteq C_1\) 对应起来。容错量子计算的阈值定理（Aharonov-Ben-Or，Knill-Laflamme-Zurek，1997-1998）证明了如果每个物理量子门的错误率低于某个阈值（约 \(10^{-4}\)），则通过级联纠错码可以在任意长的计算中保持量子错误可控。

\subsection{266.1 量子纠错原理}\label{ux91cfux5b50ux7ea0ux9519ux539fux7406}

\textbf{定义 266.1}（量子纠错码 / Quantum Error-Correcting Code）：\([[n, k, d]]\) 量子码将 \(k\) 个逻辑量子比特编码为 \(n\) 个物理量子比特，可纠正任意 \(\lfloor (d-1)/2 \rfloor\) 个量子比特上的错误（包括相位错误和比特翻转错误）。

\textbf{定理 266.1}（量子纠错条件 / Knill-Laflamme 条件，1997）：一组错误算子 \(\{E_a\}\) 可被码 \(C\)（投影 \(P\)）纠正当且仅当 \[P E_a^\dagger E_b P = \lambda_{ab} P\]

对所有 \(a, b\)，\(\lambda_{ab}\) 是 Hermitian 矩阵。这是经典纠错条件 \(\mathbf{1}^T H E_a^T E_b \mathbf{1} = \lambda_{ab}\) 的量子推广。

\textbf{证明}：必要性：若存在纠错操作 \(\mathcal{R}\) 满足 \(\mathcal{R}(E_a \rho E_b^\dagger) = \rho\)（对所有 \(\rho\) 支撑在 \(P\) 上），则对任意 \(|\psi\rangle, |\phi\rangle \in \operatorname{Im}(P)\)，\(\langle\psi|E_a^\dagger E_b|\phi\rangle = \operatorname{Tr}(E_b|\phi\rangle\langle\psi|E_a^\dagger) = \operatorname{Tr}(\mathcal{R}(E_b|\phi\rangle\langle\psi|E_a^\dagger)) = \langle\psi|\phi\rangle \lambda_{ab}\)，其中 \(\lambda_{ab}\) 与 \(|\psi\rangle, |\phi\rangle\) 无关。充分性：对 Hermitian 矩阵 \(\lambda\)，由其 Cholesky 分解 \(\lambda = UDU^\dagger\) 构造恢复算子 \(R_k = \sum_a U_{ak} P E_a^\dagger\)，验证 \(\sum_k R_k^\dagger R_k = P\) 且 \(R_k E_a P = \lambda_{ak}^{1/2} P\)。\(\blacksquare\)

\textbf{定义 266.2}（稳定子码 / Stabilizer Code，Gottesman 1997）：\textbf{稳定子码}由 Pauli 群的 Abel 子群 \(S\)（稳定子）的公共 +1 本征空间定义。著名的 Shor 码 \([[9, 1, 3]]\) 和 CSS 码都是稳定子码。

\textbf{定理 266.2}（CSS 构造 / Calderbank-Shor-Steane 1996）：若存在经典线性码 \(C_1 = [n, k_1, d_1]\) 和 \(C_2 = [n, k_2, d_2]\) 满足 \(C_2^\perp \subseteq C_1\)，则可构造量子 CSS 码 \([[n, k_1 - k_2, \min(d_1, d_2)]]\)。

\textbf{证明}：定义量子码空间 \(Q = \operatorname{span}\{|x + C_2\rangle : x \in C_1\}\)，其中 \(|x + C_2\rangle = \frac{1}{\sqrt{|C_2|}}\sum_{y \in C_2} |x+y\rangle\)。\(C_2^\perp \subseteq C_1\) 确保这些态正交。\(X\) 型稳定子由 \(C_2\) 的生成元给出：对每个 \(w \in C_2\)，\(X^w = \prod_{i} X_i^{w_i}\) 稳定 \(Q\)。\(Z\) 型稳定子由 \(C_1^\perp\) 的生成元给出：对每个 \(v \in C_1^\perp\)，\(Z^v = \prod_{i} Z_i^{v_i}\) 稳定 \(Q\)。码距 \(d \geq \min(d_1, d_2)\)：\(X\) 错误（比特翻转）由 \(C_1\) 检测，\(Z\) 错误（相位翻转）由 \(C_2^\perp\) 检测。逻辑量子比特数 \(k = k_1 - k_2\)。\(\blacksquare\)

\subsection{266.2 表面码与容错阈值}\label{ux8868ux9762ux7801ux4e0eux5bb9ux9519ux9608ux503c}

\textbf{定义 266.3}（表面码 / Surface Code, Kitaev / Bravyi-Kitaev）：在 \(L \times L\) 二维格子上定义的稳定子码。距离 \(d = L\)，每个格点上 4 体或 2 体测量。\textbf{容错阈值}：若每个物理门的错误率低于阈值 \(p_{\text{th}}\)，则增加 \(L\)（扩大码距）可使逻辑错误率指数衰减。

表面码是当前最有前景的容错量子计算架构之一。其核心思想是将量子信息编码在二维格子的拓扑自由度中：稳定子生成元为每个面上的 \(X\)-型星形算子 \(A_v = \prod_{e \ni v} X_e\) 和每个顶点上的 \(Z\)-型面算子 \(B_f = \prod_{e \in \partial f} Z_e\)。这些稳定子两两交换，且它们的公共 \(+1\) 本征空间定义了码空间。表面码的突出优势在于其\textbf{几何局域性}------所有稳定子测量仅涉及最近邻量子比特，这极大降低了实验实现的难度。表面码的距离 \(d\) 等于格子边长 \(L\)，意味着它可以纠正最多 \(\lfloor (d-1)/2 \rfloor\) 个任意错误。错误症状（syndrome）通过稳定子测量结果获得，并使用最小权完美匹配（minimum-weight perfect matching）算法解码。表面码的容错阈值约为 \(p_{\text{th}} \approx 1\%\)（取决于具体噪声模型），这是目前已知的二维架构中最高阈值之一。

\textbf{定理 266.3}（阈值定理 / Aharonov-Ben-Or / Kitaev / Knill-Laflamme-Zurek）：若物理量子门的错误率低于某个常数阈值，则原则上可以进行任意长的容错量子计算。证明方法在于使用量子纠错码的分级编码（concatenation）。

阈值定理是容错量子计算的理论基石。其证明策略基于\textbf{分级编码}（concatenated coding）：将量子纠错码逐层嵌套，每一层纠正其下层引入的错误。设 \(\varepsilon\) 为物理门错误率，\(\varepsilon_k\) 为第 \(k\) 级编码后的逻辑错误率。若 \(\varepsilon < p_{\text{th}}\)，则 \(\varepsilon_k \leq p_{\text{th}} (\varepsilon/p_{\text{th}})^{2^k}\)------逻辑错误率双重指数衰减。这一定理的意义在于：只要物理错误率低于某个常数阈值（而非反比于计算规模），量子计算即可通过增加编码层级实现任意精度。实际实现中，表面码的阈值定理优于分级编码------只需扩大格子尺寸 \(L\)（而非增加编码层级），逻辑错误率 \(p_L\) 随 \(L\) 指数衰减：\(p_L \sim (p/p_{\text{th}})^{d/2}\)。Google 的 Willow 量子处理器（2024年）在表面码架构下首次展示了''低于阈值''的量子纠错------随着码距从 3 增加到 7，逻辑错误率减半，标志着容错量子计算进入了实验验证阶段。

\textbf{证明}：设物理门错误率为 \(p_0\)。使用 \([[n,k,d]]\) 纠错码，若 \(p_0\) 低于阈值 \(p_{\text{th}}\)（由纠错码和容错协议决定），则一级编码后逻辑错误率 \(p_1 \leq C p_0^{t+1}\)（\(t = \lfloor (d-1)/2\rfloor\)）。分级编码（\(L\) 级 concatenation）后 \(p_L \leq p_{\text{th}}(p_0/p_{\text{th}})^{(t+1)^L}\)。当 \(p_0 < p_{\text{th}}\) 时，\(p_L \to 0\) 随 \(L\) 指数衰减。所需物理量子比特数 \(N_L = n^L\) 为多项式增长（\(L = O(\log\log(1/\varepsilon))\)）。故在低于阈值时，任意长计算可通过适当增加编码级数实现。\(\blacksquare\)

\subsection{266.3 拓扑量子计算}\label{ux62d3ux6251ux91cfux5b50ux8ba1ux7b97}

\textbf{定义 266.4}（anyons 与拓扑量子计算）：二维系统中的 anyons（既非玻色子也非费米子）的\textbf{编织操作}（braiding）构成量子门。非 Abelian anyons（如 Majorana 零模、Ising anyons、Fibonacci anyons）的编织是实现拓扑保护量子门的一种途径。

\textbf{定理 266.4}（Fibonacci anyons 的通用性）：Fibonacci anyons 的编织操作可近似实现任意单量子比特门和 CNOT 门，从而是通用的拓扑量子计算模型。这是通过编织群在量子比特空间上的稠密性保证的（类似于 Solovay-Kitaev 定理）。

\textbf{证明}：Fibonacci anyons 的融合规则为 \(\tau \otimes \tau = 1 \oplus \tau\)（其中 \(\tau\) 为 Fibonacci anyon 粒子）。两个 \(\tau\) 的融合空间维数为 \(2\)（对应量子比特）。编织操作 \(\sigma_1, \sigma_2\) 生成编织群 \(B_3\)，其在 \(2\) 维融合空间上的表示生成 \(SU(2)\) 的稠密子群（由 Jones 多项式表示论）。具体地，\(\sigma_1 \mapsto e^{-i\pi/5}\operatorname{diag}(e^{-i4\pi/5}, e^{i3\pi/5})\)，\(\sigma_2 \mapsto FR F\)。由 Solovay-Kitaev 定理，任何 \(SU(2)\) 门可被编织序列 \(\varepsilon\)-逼近，所需编织长度 \(O(\log^c(1/\varepsilon))\)。CNOT 可通过两量子比特融合空间中的编织实现。\(\blacksquare\)

\hrulefill

\emph{本卷在量子力学数学基础（V30）和泛函分析（V14）的基础上，系统建立了量子态与测量公理、量子纠缠理论、量子信道与容量、量子算法和量子纠错与容错计算的数学理论。共 5 章。}

\emph{本卷从变分原理出发，建立了 Euler-Lagrange 方程、Hamilton 方程与辛几何、Noether 定理、Hilbert 空间上自伴算子的谱理论、纤维丛上的变分问题（场论）、伪 Riemann 几何、共形场论、可积系统、重正化群和 Vlasov-Maxwell 系统的数学理论。\textbf{补充部分}从算子代数角度审视量子信息理论：Hilbert 空间上的密度矩阵与 CPTP 映射给出了量子操作的谱表示，von Neumann 熵与量子相对熵建立了量子信息度量的公理基础，量子纠错码的稳定子形式将纠错问题转化为线性代数约束，算子代数（C}-代数、von Neumann 代数）为量子多体系统和量子场论提供了严格数学框架。*

\emph{卷二十九：数学物理方法终。}

\emph{卷二十九：数学物理方法终。}

\emph{卷二十九：数学物理方法终。}

\emph{卷三十：数学物理方法终。}

\newpage

\chapter{13-数学物理/数学物理/08-弹性力学数学理论.md}\label{ux6570ux5b66ux7269ux7406ux6570ux5b66ux7269ux740608-ux5f39ux6027ux529bux5b66ux6570ux5b66ux7406ux8bba.md}

\chapter{弹性力学的数学理论}\label{ux5f39ux6027ux529bux5b66ux7684ux6570ux5b66ux7406ux8bba}

\hrulefill

\section{第467章：线性弹性理论}\label{ux7b2c467ux7ae0ux7ebfux6027ux5f39ux6027ux7406ux8bba}

弹性力学研究在外力作用下可变形固体的应力、应变与位移之间的关系。本章在连续介质力学的框架下建立线性弹性理论的严格数学基础。

\subsection{应变与应力的张量表示}\label{ux5e94ux53d8ux4e0eux5e94ux529bux7684ux5f20ux91cfux8868ux793a}

\textbf{定义 380.1（位移场与变形梯度）} 设 \(\Omega \subset \mathbb{R}^3\) 是物体在未变形状态（参考构型）所占的区域。\textbf{位移场} \(\mathbf{u}: \Omega \to \mathbb{R}^3\) 描述物质点 \(\mathbf{x} \in \Omega\) 的位移。变形映射 \(\boldsymbol{\varphi}: \Omega \to \mathbb{R}^3\) 定义为 \(\boldsymbol{\varphi}(\mathbf{x}) = \mathbf{x} + \mathbf{u}(\mathbf{x})\)。\textbf{变形梯度}为 \[\mathbf{F} = \nabla \boldsymbol{\varphi} = \mathbf{I} + \nabla \mathbf{u},\] 其中 \((\nabla \mathbf{u})_{ij} = \partial u_i / \partial x_j\) 是位移梯度。

\textbf{定义 380.2（Green-Saint Venant应变张量）} 在小变形假设下（\(\|\nabla \mathbf{u}\| \ll 1\)），线性化应变张量（无穷小应变张量）定义为 \[\boldsymbol{\varepsilon}(\mathbf{u}) = \frac{1}{2} \left( \nabla \mathbf{u} + (\nabla \mathbf{u})^T \right), \quad \varepsilon_{ij} = \frac{1}{2} \left( \frac{\partial u_i}{\partial x_j} + \frac{\partial u_j}{\partial x_i} \right).\] 这是二阶对称张量，\(\boldsymbol{\varepsilon} \in \mathbb{R}^{3 \times 3}_{\text{sym}}\)。

应变张量的各分量具有明确物理含义：对角元 \(\varepsilon_{ii}\) 表示沿第 \(i\) 方向的\textbf{正应变}（伸长率），非对角元 \(\varepsilon_{ij}\)（\(i \neq j\)）表示\textbf{剪应变}（剪切角变化的一半）。

\textbf{定理 380.3（Cauchy应力原理）} 设 \(\mathbf{n}\) 是物体内部某点的单位法向量。存在唯一的二阶张量 \(\boldsymbol{\sigma}\)（\textbf{Cauchy应力张量}）使得作用于以 \(\mathbf{n}\) 为法向量的面元上的牵引力（traction）为 \[\mathbf{t}(\mathbf{n}) = \boldsymbol{\sigma} \cdot \mathbf{n} = \sigma_{ij} n_j \ (\text{使用Einstein求和约定}).\]

\textbf{证明} Cauchy四面体论证：考虑以 \(\mathbf{n}\) 为外法向量的斜面和三个坐标面构成的无穷小四面体。动量守恒要求 \[\mathbf{t}(\mathbf{n}) dA - \mathbf{t}(\mathbf{e}_1) dA_1 - \mathbf{t}(\mathbf{e}_2) dA_2 - \mathbf{t}(\mathbf{e}_3) dA_3 = \rho \, \mathbf{a} \, dV,\] 其中 \(dA_i = n_i dA\)。当体积元趋于零时惯性项消失（\(O(h^3)\) vs \(O(h^2)\)），得 \(\mathbf{t}(\mathbf{n}) = n_i \mathbf{t}(\mathbf{e}_i)\)，定义 \(\sigma_{ij} = t_j(\mathbf{e}_i)\) 即得线性关系。由角动量守恒得 \(\boldsymbol{\sigma} = \boldsymbol{\sigma}^T\)（应力张量的对称性）。\(\blacksquare\)

\textbf{定义 380.4（平衡方程）} 在静力学情形下，内部应力必须满足平衡方程： \[\nabla \cdot \boldsymbol{\sigma} + \mathbf{f} = \mathbf{0}, \quad \text{即} \quad \frac{\partial \sigma_{ij}}{\partial x_j} + f_i = 0 \ \text{在} \ \Omega \ \text{内},\] 其中 \(\mathbf{f}\) 是体积力密度。边界 \(\partial \Omega\) 上，可施加位移边界条件 \(\mathbf{u} = \mathbf{0}\)（在 \(\Gamma_D\) 上）或牵引力边界条件 \(\boldsymbol{\sigma} \cdot \mathbf{n} = \mathbf{g}\)（在 \(\Gamma_N\) 上）。

\subsection{广义Hooke定律与各向异性弹性}\label{ux5e7fux4e49hookeux5b9aux5f8bux4e0eux5404ux5411ux5f02ux6027ux5f39ux6027}

\textbf{定理 380.5（广义Hooke定律）} 在线性弹性理论中，应力和应变由广义Hooke定律线性关联： \[\sigma_{ij} = \mathbb{C}_{ijkl} \, \varepsilon_{kl},\] 其中 \(\mathbb{C} = (\mathbb{C}_{ijkl})\) 是四阶\textbf{弹性张量}（刚度张量），满足以下对称性： \[\mathbb{C}_{ijkl} = \mathbb{C}_{jikl} = \mathbb{C}_{ijlk} = \mathbb{C}_{klij}.\]

前两个对称性源于 \(\boldsymbol{\sigma}\) 和 \(\boldsymbol{\varepsilon}\) 的对称性，第三个（主对称性）源于超弹性能量势的存在性：存在应变能密度函数 \(W(\boldsymbol{\varepsilon}) = \frac{1}{2} \mathbb{C}_{ijkl} \varepsilon_{ij} \varepsilon_{kl}\) 使得 \(\sigma_{ij} = \partial W / \partial \varepsilon_{ij}\)。

\textbf{命题 380.6（弹性张量的独立分量数）} 四阶弹性张量的对称性将独立分量从 \(3^4 = 81\) 降至 \(21\)（一般各向异性材料）。进一步的材料对称性（晶体类、横观各向同性、各向同性）进一步减少自由参数。

\textbf{证明}：由连续介质力学的局部作用原理和客观性原理，应力仅依赖于当前点的变形状态。在小变形假设下仅保留线性项得 \(\sigma_{ij} = \mathbb{C}_{ijkl} \varepsilon_{kl}\)。前两个对称性 \(\mathbb{C}_{ijkl} = \mathbb{C}_{jikl} = \mathbb{C}_{ijlk}\) 直接源于 \(\boldsymbol{\sigma}\) 和 \(\boldsymbol{\varepsilon}\) 的对称性。主对称性 \(\mathbb{C}_{ijkl} = \mathbb{C}_{klij}\) 来自超弹性能量势的存在性：\(W(\boldsymbol{\varepsilon}) = \frac{1}{2} \mathbb{C}_{ijkl} \varepsilon_{ij} \varepsilon_{kl}\)，则 \(\sigma_{ij} = \partial W / \partial \varepsilon_{ij} = \frac{1}{2} (\mathbb{C}_{ijkl} + \mathbb{C}_{klij}) \varepsilon_{kl}\)。由二阶混合偏导数的对称性，\(\partial^2 W / (\partial \varepsilon_{ij} \partial \varepsilon_{kl}) = \partial^2 W / (\partial \varepsilon_{kl} \partial \varepsilon_{ij})\)，得 \(\mathbb{C}_{ijkl} = \mathbb{C}_{klij}\)。\(\blacksquare\)

\textbf{定理 380.7（各向同性弹性的Lamé形式）} 对均质各向同性材料，弹性张量仅由两个独立常数确定------\textbf{Lamé常数} \(\lambda\) 和 \(\mu\)： \[\sigma_{ij} = \lambda \varepsilon_{kk} \delta_{ij} + 2 \mu \varepsilon_{ij},\] 或等价地， \[\boldsymbol{\sigma} = \lambda (\operatorname{tr} \boldsymbol{\varepsilon}) \mathbf{I} + 2 \mu \boldsymbol{\varepsilon}.\] \(\mu\) 也称为剪切模量 \(G\)。杨氏模量 \(E\) 和Poisson比 \(\nu\) 由 \(\lambda, \mu\) 表示为 \[E = \frac{\mu(3\lambda + 2\mu)}{\lambda + \mu}, \quad \nu = \frac{\lambda}{2(\lambda + \mu)}.\]

\textbf{证明} 对各向同性四阶张量 \(\mathbb{C}_{ijkl} = \lambda \delta_{ij} \delta_{kl} + \mu (\delta_{ik} \delta_{jl} + \delta_{il} \delta_{jk})\)，代入Hooke定律即得。该形式是各向同性条件下仅有的满足全部对称性的线性不变量表达式。\(\blacksquare\)

\subsection{弹性静力学的边值问题}\label{ux5f39ux6027ux9759ux529bux5b66ux7684ux8fb9ux503cux95eeux9898}

\textbf{定义 380.8（Navier-Lamé方程）} 将Hooke定律代入平衡方程，得到以位移为基本未知量的Navier方程： \[\mu \Delta \mathbf{u} + (\lambda + \mu) \nabla (\nabla \cdot \mathbf{u}) + \mathbf{f} = \mathbf{0} \ \text{在} \ \Omega \ \text{内}.\]

\textbf{定理 380.9（Korn不等式）} 设 \(\Omega \subset \mathbb{R}^3\) 是有界Lipschitz区域。则存在常数 \(C = C(\Omega) > 0\) 使得对所有 \(\mathbf{u} \in H^1(\Omega; \mathbb{R}^3)\)， \[\|\nabla \mathbf{u}\|_{L^2(\Omega)}^2 \leq C \left( \|\mathbf{u}\|_{L^2(\Omega)}^2 + \|\boldsymbol{\varepsilon}(\mathbf{u})\|_{L^2(\Omega)}^2 \right).\]

\textbf{证明} 利用恒等式 \[\frac{\partial^2 u_i}{\partial x_j \partial x_k} = \frac{\partial \varepsilon_{ij}}{\partial x_k} + \frac{\partial \varepsilon_{ik}}{\partial x_j} - \frac{\partial \varepsilon_{jk}}{\partial x_i}.\] 两边乘以测试函数并在 \(\Omega\) 上积分，经两次分部积分得： \[\|\nabla^2 \mathbf{u}\|_{L^2}^2 \leq C \|\nabla \boldsymbol{\varepsilon}\|_{L^2}^2.\] 利用Poincaré不等式，梯度的 \(L^2\) 范数可由该估计和位移的 \(L^2\) 范数控制。完整的Korn第二不等式（在 \(H_0^1(\Omega)\) 中无需 \(\|\mathbf{u}\|_{L^2}\) 项）要求更精细的论证，依赖于 \(\boldsymbol{\varepsilon}(\mathbf{u}) = 0 \iff \mathbf{u}\) 是无穷小刚体运动（即 \(\mathbf{u}(\mathbf{x}) = \mathbf{a} + \mathbf{b} \times \mathbf{x}\)）。\(\blacksquare\)

\textbf{定理 380.10（线性弹性静力学问题的适定性）} 考虑纯位移边值问题：求 \(\mathbf{u} \in H_0^1(\Omega; \mathbb{R}^3)\) 满足 \[-\nabla \cdot (\mathbb{C} : \boldsymbol{\varepsilon}(\mathbf{u})) = \mathbf{f} \ \text{在} \ \Omega \ \text{内}.\] 若 \(\mathbb{C}\) 满足强椭圆性条件：存在 \(\alpha > 0\) 使得对所有对称矩阵 \(\boldsymbol{\xi}\)， \[\mathbb{C}_{ijkl} \xi_{ij} \xi_{kl} \geq \alpha \|\boldsymbol{\xi}\|^2,\] 则由Lax-Milgram引理，该问题存在唯一解且连续依赖于 \(\mathbf{f}\)。

\textbf{证明}：纯位移边值问题的弱形式为：求 \(\mathbf{u} \in H_0^1(\Omega; \mathbb{R}^3)\) 使得对任意 \(\mathbf{v} \in H_0^1(\Omega; \mathbb{R}^3)\)， \[\int_\Omega \mathbb{C}_{ijkl} \varepsilon_{kl}(\mathbf{u}) \varepsilon_{ij}(\mathbf{v}) d\mathbf{x} = \int_\Omega \mathbf{f} \cdot \mathbf{v} d\mathbf{x}.\] 定义双线性型 \(B(\mathbf{u}, \mathbf{v}) = \int_\Omega \boldsymbol{\sigma}(\mathbf{u}) : \boldsymbol{\varepsilon}(\mathbf{v}) d\mathbf{x}\) 和线性泛函 \(F(\mathbf{v}) = \int_\Omega \mathbf{f} \cdot \mathbf{v} d\mathbf{x}\)。强椭圆性条件给出 \(B(\mathbf{u}, \mathbf{u}) \geq \alpha \|\boldsymbol{\varepsilon}(\mathbf{u})\|_{L^2}^2\)。由Korn第二不等式（对 \(\mathbf{u} \in H_0^1(\Omega)\) 有 \(\|\nabla \mathbf{u}\|_{L^2} \leq C \|\boldsymbol{\varepsilon}(\mathbf{u})\|_{L^2}\)），得 \(B(\mathbf{u}, \mathbf{u}) \geq \alpha C^{-2} \|\mathbf{u}\|_{H^1}^2\)，故 \(B\) 是 \(H_0^1\) 强制的。\(B\) 的有界性由 \(\mathbb{C}\) 的分量有界性保证。\(F\) 的有界性：\(|F(\mathbf{v})| \leq \|\mathbf{f}\|_{L^2} \|\mathbf{v}\|_{L^2} \leq \|\mathbf{f}\|_{L^2} \|\mathbf{v}\|_{H^1}\)。由Lax-Milgram引理，存在唯一 \(\mathbf{u} \in H_0^1\) 满足弱形式。解对数据的连续依赖性由 \(B(\mathbf{u}, \mathbf{u}) = F(\mathbf{u})\) 和强强制性得 \(\|\mathbf{u}\|_{H^1} \leq \alpha^{-1} C^2 \|\mathbf{f}\|_{L^2}\)。\(\blacksquare\)

\subsection{弹性动力学的波动方程}\label{ux5f39ux6027ux52a8ux529bux5b66ux7684ux6ce2ux52a8ux65b9ux7a0b}

\textbf{定义 380.11（弹性动力学方程）} 包含惯性项的动力学方程为 \[\rho \frac{\partial^2 \mathbf{u}}{\partial t^2} = \nabla \cdot \boldsymbol{\sigma} + \mathbf{f},\] 其中 \(\rho\) 是密度。代入各向同性本构关系得 \[\rho \ddot{\mathbf{u}} = \mu \Delta \mathbf{u} + (\lambda + \mu) \nabla (\nabla \cdot \mathbf{u}) + \mathbf{f}.\]

\textbf{定理 380.12（弹性波的分解 - Helmholtz分解）} 利用矢量场的Helmholtz分解 \(\mathbf{u} = \nabla \phi + \nabla \times \boldsymbol{\psi}\)（\(\nabla \cdot \boldsymbol{\psi} = 0\)），弹性动力学方程分解为两个独立的波动方程：

（1）\textbf{P波（纵波 / 压缩波）}：\(\phi_{tt} = c_P^2 \Delta \phi\)，波速 \(c_P = \sqrt{(\lambda + 2\mu)/\rho}\)；

（2）\textbf{S波（横波 / 剪切波）}：\(\boldsymbol{\psi}_{tt} = c_S^2 \Delta \boldsymbol{\psi}\)，波速 \(c_S = \sqrt{\mu/\rho}\)。

由于 \(\lambda, \mu > 0\)，总有 \(c_P > c_S\)（P波快于S波）。这是地震学中P波先于S波到达的数学根源。

\textbf{证明}：将Helmholtz分解 \(\mathbf{u} = \nabla \phi + \nabla \times \boldsymbol{\psi}\)（\(\nabla \cdot \boldsymbol{\psi} = 0\)）代入Navier方程： \[\rho (\nabla \ddot{\phi} + \nabla \times \ddot{\boldsymbol{\psi}}) = \mu \Delta (\nabla \phi + \nabla \times \boldsymbol{\psi}) + (\lambda + \mu) \nabla (\nabla \cdot (\nabla \phi + \nabla \times \boldsymbol{\psi})).\] 由矢量恒等式 \(\nabla \cdot (\nabla \times \boldsymbol{\psi}) = 0\)，\(\nabla \cdot (\nabla \phi) = \Delta \phi\) 和 \(\Delta \nabla \phi = \nabla \Delta \phi\)，\(\Delta(\nabla \times \boldsymbol{\psi}) = \nabla \times \Delta \boldsymbol{\psi}\)，方程化为： \[\nabla [\rho \ddot{\phi} - (\lambda + 2\mu) \Delta \phi] + \nabla \times [\rho \ddot{\boldsymbol{\psi}} - \mu \Delta \boldsymbol{\psi}] = \mathbf{0}.\] 由于梯度和旋度分别是无旋场和无散场的正交分解，两边分别取散度和旋度得两个独立的波动方程：\(\rho \ddot{\phi} = (\lambda + 2\mu) \Delta \phi\)（P波，\(c_P^2 = (\lambda + 2\mu)/\rho\)）；\(\rho \ddot{\boldsymbol{\psi}} = \mu \Delta \boldsymbol{\psi}\)（S波，\(c_S^2 = \mu/\rho\)）。因 \(\lambda, \mu > 0\)，\(c_P > c_S\)。\(\blacksquare\)

\textbf{定义 380.13（Rayleigh波）} 在半无限弹性体（\(x_3 \geq 0\)，自由面 \(x_3 = 0\) 上 \(\sigma_{i3} = 0\)）的表面附近，存在沿表面传播且随深度指数衰减的波------\textbf{Rayleigh波}。设解形式为 \(\mathbf{u} = \mathbf{U}(x_3) e^{i(k x_1 - \omega t)}\)，代入Navier方程和自由面边界条件得到色散关系： \[\left(2 - \frac{c_R^2}{c_S^2}\right)^2 = 4 \sqrt{1 - \frac{c_R^2}{c_P^2}} \sqrt{1 - \frac{c_R^2}{c_S^2}},\] 其中 \(c_R = \omega/k\) 是Rayleigh波速。该方程有唯一正实根 \(c_R < c_S\)。

\subsection{弹性力学的变分形式}\label{ux5f39ux6027ux529bux5b66ux7684ux53d8ux5206ux5f62ux5f0f}

\textbf{定理 380.14（势能极小化原理）} 在静力学中，真实位移场 \(\mathbf{u}\) 是以下总势能泛函的极小化子： \[\mathcal{E}(\mathbf{v}) = \frac{1}{2} \int_{\Omega} \mathbb{C}_{ijkl} \varepsilon_{ij}(\mathbf{v}) \varepsilon_{kl}(\mathbf{v}) \, d\mathbf{x} - \int_{\Omega} \mathbf{f} \cdot \mathbf{v} \, d\mathbf{x} - \int_{\Gamma_N} \mathbf{g} \cdot \mathbf{v} \, dS.\]

\textbf{证明} 计算 \(\mathcal{E}\) 的Gateaux导数：对任意容许变分 \(\mathbf{w}\)， \[\langle \delta \mathcal{E}(\mathbf{u}), \mathbf{w} \rangle = \int_{\Omega} \boldsymbol{\sigma} : \boldsymbol{\varepsilon}(\mathbf{w}) \, d\mathbf{x} - \int_{\Omega} \mathbf{f} \cdot \mathbf{w} \, d\mathbf{x} - \int_{\Gamma_N} \mathbf{g} \cdot \mathbf{w} \, dS.\] 分部积分并利用 \(\boldsymbol{\varepsilon}(\mathbf{w})\) 的对称性得 \[\int_{\Omega} (-\nabla \cdot \boldsymbol{\sigma} - \mathbf{f}) \cdot \mathbf{w} \, d\mathbf{x} + \int_{\Gamma_N} (\boldsymbol{\sigma} \cdot \mathbf{n} - \mathbf{g}) \cdot \mathbf{w} \, dS = 0.\] 由变分基本引理，这等价于弹性的平衡方程和自然边界条件。泛函的严格凸性（来自 \(\mathbb{C}\) 的强椭圆性）确保极小化子即为平衡解。\(\blacksquare\)

\hrulefill

\section{第468章：非线性弹性与弹塑性}\label{ux7b2c468ux7ae0ux975eux7ebfux6027ux5f39ux6027ux4e0eux5f39ux5851ux6027}

\subsection{有限变形理论}\label{ux6709ux9650ux53d8ux5f62ux7406ux8bba}

\textbf{定义 381.1（有限应变度量）} 当变形非无穷小时（\(\|\nabla \mathbf{u}\|\) 不满足 \(\ll 1\)），需使用有限应变理论。变形梯度 \(\mathbf{F} = \nabla \boldsymbol{\varphi}\) 可通过极分解 \(\mathbf{F} = \mathbf{R} \mathbf{U} = \mathbf{V} \mathbf{R}\) 分离旋转和拉伸，其中 \(\mathbf{R}\) 是正交张量（旋转），\(\mathbf{U}, \mathbf{V}\) 是对称正定张量（右、左拉伸张量）。

\textbf{右Cauchy-Green变形张量} 和 \textbf{Green-Lagrange应变张量} 定义为 \[\mathbf{C} = \mathbf{F}^T \mathbf{F} = \mathbf{U}^2, \quad \mathbf{E} = \frac{1}{2} (\mathbf{C} - \mathbf{I}) = \frac{1}{2} (\nabla \mathbf{u} + (\nabla \mathbf{u})^T + (\nabla \mathbf{u})^T \nabla \mathbf{u}).\] 当 \(\|\nabla \mathbf{u}\| \ll 1\) 时，忽略高阶项 \(\mathbf{E} \approx \boldsymbol{\varepsilon}\)。

\textbf{定义 381.2（Piola-Kirchhoff应力张量）} 在有限变形中需区分参考构型和当前构型中的应力：

\begin{itemize}
\item
  \textbf{Cauchy应力} \(\boldsymbol{\sigma}\)：作用于当前构型中的真实应力；
\item
  \textbf{第一Piola-Kirchhoff应力} \(\mathbf{P}\)：参考构型中的名义应力，满足 \(\mathbf{P} = J \boldsymbol{\sigma} \mathbf{F}^{-T}\)（\(J = \det \mathbf{F}\)）；
\item
  \textbf{第二Piola-Kirchhoff应力} \(\mathbf{S}\)：完全在参考构型中的对称应力，\(\mathbf{S} = \mathbf{F}^{-1} \mathbf{P} = J \mathbf{F}^{-1} \boldsymbol{\sigma} \mathbf{F}^{-T}\)。
\end{itemize}

参考构型中的平衡方程为 \(\nabla_{\mathbf{X}} \cdot \mathbf{P} + \mathbf{f}_0 = \mathbf{0}\)，其中 \(\nabla_{\mathbf{X}}\) 是参考坐标下的散度。

\subsection{超弹性材料的本构关系}\label{ux8d85ux5f39ux6027ux6750ux6599ux7684ux672cux6784ux5173ux7cfb}

\textbf{定义 381.3（超弹性材料）} \textbf{超弹性（或Green-弹性）材料}的力学行为由\textbf{存贮能函数}（strain energy density function）\(W(\mathbf{F})\) 完全确定。第一Piola-Kirchhoff应力为 \[\mathbf{P} = \frac{\partial W}{\partial \mathbf{F}}, \quad P_{iJ} = \frac{\partial W}{\partial F_{iJ}}.\]

由物质客观性原理（material frame indifference），\(W\) 必须满足对任意旋转 \(\mathbf{Q} \in SO(3)\) 有 \(W(\mathbf{Q}\mathbf{F}) = W(\mathbf{F})\)。这等价于 \(W\) 仅通过 \(\mathbf{C} = \mathbf{F}^T \mathbf{F}\) 依赖于 \(\mathbf{F}\)：存在 \(\widehat{W}(\mathbf{C})\) 使得 \(W(\mathbf{F}) = \widehat{W}(\mathbf{F}^T \mathbf{F})\)。此时第二Piola-Kirchhoff应力为 \(\mathbf{S} = 2 \partial \widehat{W} / \partial \mathbf{C}\)。

\textbf{定义 381.4（各向同性超弹性）} 若材料是各向同性的，则 \(W\) 仅依赖于 \(\mathbf{C}\) 的主不变量（或等价地 \(\mathbf{B} = \mathbf{F}\mathbf{F}^T\) 的主不变量）： \[W = W(I_1, I_2, I_3), \quad I_1 = \operatorname{tr} \mathbf{C}, \quad I_2 = \frac{1}{2}[(\operatorname{tr} \mathbf{C})^2 - \operatorname{tr}(\mathbf{C}^2)], \quad I_3 = \det \mathbf{C} = J^2.\]

\textbf{定理 381.5（Ogden模型）} Ogden（1972）提出的一般不可压缩超弹性模型以主拉伸比 \(\lambda_1, \lambda_2, \lambda_3\)（\(\lambda_1 \lambda_2 \lambda_3 = 1\)）表示： \[W = \sum_{p=1}^{N} \frac{\mu_p}{\alpha_p} (\lambda_1^{\alpha_p} + \lambda_2^{\alpha_p} + \lambda_3^{\alpha_p} - 3),\] 其中 \(\mu_p, \alpha_p\) 是材料参数（\(\alpha_p\) 不必为整数）。该模型包含了Neo-Hookean（\(N=1, \alpha_1=2\)）和Mooney-Rivlin模型作为特例，能准确拟合橡胶类材料的大变形行为。

\subsection{Ball的存在性理论}\label{ballux7684ux5b58ux5728ux6027ux7406ux8bba}

\textbf{定理 381.6（Ball的polyconvexity方法与存在性, 1977）} 考虑纯位移边值问题：求 \(\boldsymbol{\varphi}: \Omega \to \mathbb{R}^3\) 极小化超弹性能量 \[I(\boldsymbol{\varphi}) = \int_{\Omega} W(\nabla \boldsymbol{\varphi}) \, d\mathbf{X} - \int_{\Omega} \mathbf{f}_0 \cdot \boldsymbol{\varphi} \, d\mathbf{X},\] 在边界条件 \(\boldsymbol{\varphi}|_{\partial \Omega} = \boldsymbol{\varphi}_0\) 下。若存贮能函数 \(W\) 是\textbf{polyconvex}的，即存在凸函数 \(g\) 使得 \[W(\mathbf{F}) = g(\mathbf{F}, \operatorname{Cof} \mathbf{F}, \det \mathbf{F}),\] 且满足强制性条件 \(W(\mathbf{F}) \geq c_1 \|\mathbf{F}\|^p + c_2 \|\operatorname{Cof} \mathbf{F}\|^q + c_3 (\det \mathbf{F})^{-s} - c_4\)（\(p \geq 2, q \geq p/(p-1), s > 0\)），则能量泛函 \(I\) 在 \(W^{1,p}(\Omega; \mathbb{R}^3)\) 中存在极小化子。

\textbf{证明框架} 利用变分法的直接方法：（1）构造极小化序列 \(\{\boldsymbol{\varphi}_n\}\)，由强制性条件得其一致有界性；（2）利用函数的弱连续性质：行列式映射 \(\mathbf{F} \mapsto \det \mathbf{F}\) 是\textbf{弱连续的}（即 \(\mathbf{F}_n \rightharpoonup \mathbf{F}\) 在 \(W^{1,p}\) 中（\(p > 3\)）蕴含 \(\det \mathbf{F}_n \rightharpoonup \det \mathbf{F}\) 在 \(\mathcal{D}'\) 中），这源于其作为零阶散度项的可表示性（Piola恒等式）；（3）由polyconvexity和弱下半连续性，\(I\) 弱下半连续；（4）强制性提供紧性，极限函数即极小化子。此方法克服了非线性弹性中 \(W\) 不可为 \(\mathbf{F}\) 的凸函数的根本障碍（凸性与客观性不兼容）。\(\blacksquare\)

\subsection{弹塑性理论}\label{ux5f39ux5851ux6027ux7406ux8bba}

\textbf{定义 381.7（弹塑性分解）} 在弹塑性理论中，总应变分解为弹性部分和塑性部分： \[\boldsymbol{\varepsilon} = \boldsymbol{\varepsilon}^e + \boldsymbol{\varepsilon}^p.\] 仅弹性应变参与应力计算：\(\boldsymbol{\sigma} = \mathbb{C} : \boldsymbol{\varepsilon}^e\)。

\textbf{定义 381.8（屈服准则）} 材料开始塑性流动的条件由\textbf{屈服函数} \(f(\boldsymbol{\sigma})\) 描述：塑性变形仅当 \(f(\boldsymbol{\sigma}) = 0\) 时发生，且始终满足 \(f(\boldsymbol{\sigma}) \leq 0\)（屈服面的内侧为弹性域）。经典屈服准则包括：

\begin{itemize}
\item
  \textbf{von Mises屈服准则}（\(J_2\)-塑性）：\(f(\boldsymbol{\sigma}) = \sqrt{\frac{3}{2} \mathbf{s} : \mathbf{s}} - \sigma_Y \leq 0\)，其中 \(\mathbf{s} = \boldsymbol{\sigma} - \frac{1}{3}(\operatorname{tr} \boldsymbol{\sigma})\mathbf{I}\) 是偏应力，\(\sigma_Y\) 是屈服应力；
\item
  \textbf{Tresca屈服准则}：\(f(\boldsymbol{\sigma}) = \max_{i,j} |\sigma_i - \sigma_j| - \sigma_Y \leq 0\)（\(\sigma_i\) 为主应力）。
\end{itemize}

\textbf{定理 381.9（Prandtl-Reuss塑性流动法则）} 对关联理想弹塑性，塑性应变率为 \[\dot{\boldsymbol{\varepsilon}}^p = \dot{\gamma} \frac{\partial f}{\partial \boldsymbol{\sigma}},\] 这是最大塑性功原理（Hill最大耗散原理）的推论：真实应力在所有静力容许应力中最大化塑性耗散 \(\boldsymbol{\sigma} : \dot{\boldsymbol{\varepsilon}}^p\)。对von Mises屈服准则，\(\partial f / \partial \boldsymbol{\sigma} = \frac{3}{2} \mathbf{s} / \bar{\sigma}\)（\(\bar{\sigma}\) 是等效应力），故塑性应变率与偏应力同方向（塑性不可压缩性）。

\textbf{证明}：最大塑性功原理可表述为变分问题：真实应力 \(\boldsymbol{\sigma}\) 在所有静力容许应力 \(\boldsymbol{\sigma}^*\)（即满足 \(f(\boldsymbol{\sigma}^*) \leq 0\)）中最大化塑性耗散 \((\boldsymbol{\sigma}^* - \boldsymbol{\sigma}) : \dot{\boldsymbol{\varepsilon}}^p \leq 0\)。由凸分析，此不等式等价于 \(\dot{\boldsymbol{\varepsilon}}^p\) 属于屈服面 \(f(\boldsymbol{\sigma}) = 0\) 在 \(\boldsymbol{\sigma}\) 处的法向锥（normal cone）。对光滑屈服函数，法向锥即为梯度方向，故 \(\dot{\boldsymbol{\varepsilon}}^p = \dot{\gamma} \partial f / \partial \boldsymbol{\sigma}\)，\(\dot{\gamma} \geq 0\)。对von Mises屈服准则 \(f(\boldsymbol{\sigma}) = \sqrt{3J_2} - \sigma_Y\)，\(\partial f / \partial \sigma_{ij} = \frac{3}{2} s_{ij} / \sqrt{3J_2}\)，故塑性应变率与偏应力同向。由 \(\operatorname{tr}(\partial f / \partial \boldsymbol{\sigma}) = 0\) 得塑性不可压缩性：\(\operatorname{tr} \dot{\boldsymbol{\varepsilon}}^p = 0\)。\(\blacksquare\)

\textbf{定义 381.10（Hencky方程）} \textbf{Hencky（全量）理论}适用于比例加载路径：假设 \(\boldsymbol{\varepsilon}^p = \phi \mathbf{s}\)（\(\phi \geq 0\) 为标量乘子），总应变-应力关系简化为非线性弹性形式： \[\boldsymbol{\varepsilon} = \frac{1+\nu}{E} \boldsymbol{\sigma} - \frac{\nu}{E} (\operatorname{tr} \boldsymbol{\sigma}) \mathbf{I} + \phi \mathbf{s}.\]

\textbf{定理 381.11（加载-卸载准则 - Kuhn-Tucker条件）} 在增量弹塑性中，塑性乘子 \(\dot{\gamma}\) 满足互补条件： \[\dot{\gamma} \geq 0, \quad f(\boldsymbol{\sigma}) \leq 0, \quad \dot{\gamma} f(\boldsymbol{\sigma}) = 0.\] 这确保了仅在塑性加载（\(f = 0\) 且 \(\dot{f} = 0\)）时产生塑性流动，弹性卸载（\(f < 0\) 或 \(\dot{f} < 0\)）时 \(\dot{\gamma} = 0\)。

弹塑性问题可表为\textbf{演化变分不等式}（evolution variational inequality），其适定性依赖于应力空间中的凸分析和单调算子理论。

\newpage

\chapter{13-数学物理/数学物理/09-数学生物学导论.md}\label{ux6570ux5b66ux7269ux7406ux6570ux5b66ux7269ux740609-ux6570ux5b66ux751fux7269ux5b66ux5bfcux8bba.md}

\chapter{数学生物学导论}\label{ux6570ux5b66ux751fux7269ux5b66ux5bfcux8bba}

\hrulefill

\section{第469章：种群动力学与生态模型}\label{ux7b2c469ux7ae0ux79cdux7fa4ux52a8ux529bux5b66ux4e0eux751fux6001ux6a21ux578b}

数学生物学利用数学工具对生物系统的行为进行建模、分析和预测。种群动力学是最经典的方向之一，研究物种数量随时间变化的规律。

\subsection{Malthus与Logistic模型的完整数学分析}\label{malthusux4e0elogisticux6a21ux578bux7684ux5b8cux6574ux6570ux5b66ux5206ux6790}

\textbf{定义 382.1（Malthus模型）} 最简单的种群增长模型由Malthus（1798）提出： \[\frac{dN}{dt} = r N, \quad N(0) = N_0,\] 其中 \(N(t)\) 是时刻 \(t\) 的种群数量，\(r > 0\) 是内禀增长率。解为 \(N(t) = N_0 e^{rt}\)，呈无界指数增长。

\textbf{定义 382.2（Logistic模型）} Verhulst（1838）引入环境容纳量 \(K > 0\) 修正Malthus模型： \[\frac{dN}{dt} = r N \left(1 - \frac{N}{K}\right).\]

\textbf{定理 382.3（Logistic方程的完整分析）} 方程 \(\dot{N} = rN(1 - N/K)\) 具有两个平衡点：\(N^* = 0\)（不稳定）和 \(N^* = K\)（全局渐近稳定，对所有 \(N_0 > 0\)）。精确解为 \[N(t) = \frac{K N_0 e^{rt}}{K + N_0 (e^{rt} - 1)}.\]

\textbf{证明} 方程是可分离的：\(\frac{dN}{N(1-N/K)} = r \, dt\)。部分分式展开： \[\frac{1}{N(1-N/K)} = \frac{1}{N} + \frac{1/K}{1-N/K}.\] 积分得 \(\ln N - \ln(1 - N/K) = rt + C\)，即 \(\frac{N}{1 - N/K} = \frac{N_0}{1 - N_0/K} e^{rt}\)，解出 \(N(t)\)。稳定性可由线性化分析验证：在 \(N^* = K\) 处， \[\frac{d}{dN}\left[rN\left(1-\frac{N}{K}\right)\right]\bigg|_{N=K} = r\left(1-\frac{2K}{K}\right) = -r < 0,\] 故 \(K\) 是渐近稳定的。类似地，\(N^*=0\) 处的导数为 \(r > 0\)，是不稳定的。相图分析：对 \(0 < N < K\)，\(\dot{N} > 0\)；对 \(N > K\)，\(\dot{N} < 0\)。所有正向轨道单调收敛于 \(K\)。\(\blacksquare\)

\subsection{Lotka-Volterra捕食-被捕食模型}\label{lotka-volterraux6355ux98df-ux88abux6355ux98dfux6a21ux578b}

\textbf{定义 382.4（Lotka-Volterra模型）} Lotka（1925）和Volterra（1926）独立提出的经典捕食-被捕食系统： \[\begin{cases} \dot{x} = \alpha x - \beta x y \\ \dot{y} = \delta x y - \gamma y \end{cases}\] 其中 \(x\) 是被捕食者数量，\(y\) 是捕食者数量，\(\alpha, \beta, \gamma, \delta > 0\)。

\textbf{定理 382.5（周期解的存在性与首次积分）} 上述系统具有非平凡平衡点 \((x^*, y^*) = (\gamma/\delta, \alpha/\beta)\)。所有正初始条件的轨道均为围绕该平衡点的闭轨（周期轨道）。

\textbf{证明（首次积分法）} 消去 \(dt\)： \[\frac{dy}{dx} = \frac{y(\delta x - \gamma)}{x(\alpha - \beta y)}.\] 分离变量：\(\frac{\alpha - \beta y}{y} dy = \frac{\delta x - \gamma}{x} dx\)。积分得 \[H(x, y) = \delta x - \gamma \ln x + \beta y - \alpha \ln y = \text{常数}.\] \(H\) 是系统的首次积分（守恒量）。在 \((x^*, y^*)\) 处 \(H\) 的Hessian矩阵正定（\(\det H = \alpha \beta \gamma \delta / (x^* y^*)^2 > 0\)），故 \(H\) 的等高线是围绕平衡点的闭曲线------即周期轨道。\(\blacksquare\)

\textbf{定理 382.6（Volterra原理）} 在捕食-被捕食系统中，若被捕食者和捕食者均以与各自种群密度成比例的速率被移除（例如通过农药），则被捕食者的平均数量下降而捕食者的平均数量上升。这提供了鱼类种群在捕捞压力下``反直觉''变化的数学解释：捕捞使捕食-被捕食比例偏向被捕食者一侧。

\textbf{证明}：设系统以速率 \(c_x\) 和 \(c_y\) 分别移除被捕食者和捕食者（如捕捞或农药）。带移除的Lotka-Volterra系统为： \[\dot{x} = \alpha x - \beta x y - c_x x = (\alpha - c_x)x - \beta x y,\] \[\dot{y} = \delta x y - \gamma y - c_y y = \delta x y - (\gamma + c_y)y.\] 新平衡点为 \((\widetilde{x}^*, \widetilde{y}^*) = ((\gamma + c_y)/\delta, (\alpha - c_x)/\beta)\)。由首次积分 \(H(x,y)\)，轨道围绕新平衡点振荡。沿轨道的时间平均等于平衡点值：\(\overline{x} = (\gamma + c_y)/\delta\)，\(\overline{y} = (\alpha - c_x)/\beta\)。对捕食者而言，移除使其平衡密度下降（\(c_y > 0\) 意味着更大的有效死亡率 \(\gamma + c_y\)，从而使 \(\overline{x}\) 升高）；但捕食者的平均密度 \(\overline{y}\) 因被捕食者移除（\(c_x > 0\) 降低 \(\alpha - c_x\)）而下降。因此，同时移除两者时，被捕食者平均密度可能反而上升------因为捕食者被移除更多。\(\blacksquare\)

\subsection{竞争模型与共存条件}\label{ux7adeux4e89ux6a21ux578bux4e0eux5171ux5b58ux6761ux4ef6}

\textbf{定义 382.7（Lotka-Volterra竞争模型）} 两个物种竞争有限资源的通用模型： \[\begin{cases} \dot{x}_1 = r_1 x_1 \left(1 - \frac{x_1}{K_1} - \frac{\alpha_{12} x_2}{K_1}\right) \\ \dot{x}_2 = r_2 x_2 \left(1 - \frac{x_2}{K_2} - \frac{\alpha_{21} x_1}{K_2}\right) \end{cases}\] 其中 \(\alpha_{ij}\) 衡量物种 \(j\) 对物种 \(i\) 的竞争效应（以物种 \(i\) 的资源单位度量）。

\textbf{定理 382.8（Gause竞争排斥原理）} 两个物种稳定共存的必要条件是种内竞争强于种间竞争： \[\alpha_{12} < \frac{K_1}{K_2} \quad \text{且} \quad \alpha_{21} < \frac{K_2}{K_1}.\]

\textbf{证明（等倾线分析）} 各物种的零增长等倾线（\(\dot{x}_i = 0\)）为直线：\(x_1 = K_1 - \alpha_{12} x_2\) 和 \(x_2 = K_2 - \alpha_{21} x_1\)。两等倾线在正象限内有交点当且仅当上述不等式成立（等倾线``交叉方式''正确）。若 \(\alpha_{12} \alpha_{21} < 1\)（弱竞争），内部平衡点通过线性化矩阵的迹和行列式验证为渐近稳定结点------两物种稳定共存。若种间竞争强于种内竞争，则产生双稳态（初始条件决定哪个物种排斥另一个），即竞争排斥原理的数学体现。\(\blacksquare\)

\subsection{传染病模型}\label{ux4f20ux67d3ux75c5ux6a21ux578b}

\textbf{定义 382.9（SIR模型）} Kermack-McKendrick（1927）的SIR模型是传染病学的基石： \[\begin{cases} \dot{S} = -\beta S I \\ \dot{I} = \beta S I - \gamma I \\ \dot{R} = \gamma I \end{cases}\] 其中 \(S\)（易感者）、\(I\)（感染者）、\(R\)（康复者）满足 \(S + I + R = N\)（总人口恒定）。\(\beta\) 是传播率，\(\gamma\) 是康复率。

\textbf{定理 382.10（基本再生数与阈值定理）} \textbf{基本再生数} \(R_0\) 定义为 \[R_0 = \frac{\beta N}{\gamma}.\] 若 \(R_0 > 1\)，则疫情爆发（\(I(t)\) 初始上升）；若 \(R_0 < 1\)，则疫情消减。最终规模满足著名的Kermack-McKendrick方程： \[\ln \frac{S_{\infty}}{S_0} = -R_0 \left(1 - \frac{S_{\infty}}{N}\right),\] 其中 \(S_{\infty} = \lim_{t \to \infty} S(t)\)。

\textbf{证明} 由方程 \(\frac{dI}{dS} = -1 + \frac{\gamma}{\beta S}\)，积分得 \(I + S - \frac{\gamma}{\beta} \ln S = \text{常数}\)。\(t \to 0\) 和 \(t \to \infty\) 时 \(I \to 0\)（在数学上 \(I(0) \approx 0\)，\(I(\infty) = 0\)），故 \[S_0 - \frac{\gamma}{\beta} \ln S_0 = S_{\infty} - \frac{\gamma}{\beta} \ln S_{\infty}.\] 化为最终规模方程。当 \(t = 0\) 时 \(\dot{I} = I_0(\beta S_0 - \gamma)\)，\(R_0 = \beta N/\gamma \approx \beta S_0/\gamma\)（因 \(S_0 \approx N\)），故 \(R_0 > 1\) 确保 \(\dot{I}(0) > 0\)。\(\blacksquare\)

\textbf{定理 382.11（群体免疫阈值）} 令 \(p\) 为免疫比例（接种或自然免疫），则有效再生数为 \(R_{\text{eff}} = R_0 (1 - p)\)。当 \(R_{\text{eff}} < 1\) 时疾病无法传播，故\textbf{群体免疫阈值}为 \[p_c = 1 - \frac{1}{R_0}.\] 此公式是疫苗分配策略的数学基础。

\textbf{证明}：在随机混合的人群中，若免疫比例为 \(p\)，则有效易感者比例为 \(S_{\text{eff}}/N = (1-p)\)。有效再生数为 \(R_{\text{eff}} = R_0 (1-p)\)------每次感染产生的新病例数。当 \(R_{\text{eff}} = R_0(1-p) < 1\) 时，每个病例平均产生少于一个二代病例，流行病无法持续。解此不等式得 \(p > 1 - 1/R_0 = p_c\)。在确定性SIR模型中，由阈值定理：若初始易感者比例 \(S_0/N < 1/R_0\)，则 \(\dot{I} < 0\) 总成立。群体免疫正是通过接种将易感者比例降至 \(1/R_0\) 以下来实现。\(\blacksquare\)

\textbf{定义 382.12（SEIR模型）} 引入潜伏期（Exposed）状态 \(E\)： \[\begin{cases} \dot{S} = -\beta S I \\ \dot{E} = \beta S I - \sigma E \\ \dot{I} = \sigma E - \gamma I \\ \dot{R} = \gamma I \end{cases}\] 其中 \(1/\sigma\) 是平均潜伏期。SEIR模型的基本再生数仍为 \(R_0 = \beta N / \gamma\)（潜伏期不改变传播潜力，只延迟疫情进程）。

\subsection{年龄结构与时滞模型}\label{ux5e74ux9f84ux7ed3ux6784ux4e0eux65f6ux6edeux6a21ux578b}

\textbf{定义 382.13（McKendrick-von Foerster方程）} 描述按年龄 \(a\) 和时间 \(t\) 分布的种群密度 \(\rho(a, t)\) 的偏微分方程： \[\frac{\partial \rho}{\partial t} + \frac{\partial \rho}{\partial a} = -\mu(a) \rho(a, t),\] 边界条件：\(\rho(0, t) = \int_{0}^{\infty} \beta(a) \rho(a, t) \, da\)（新生儿出生率），初始条件：\(\rho(a, 0) = \rho_0(a)\)。\(\mu(a)\) 是年龄别死亡率，\(\beta(a)\) 是年龄别生育率。

该方程是一个一阶线性输运PDE，可用特征线法求解。平衡解 \(\rho(a)\) 满足常微分方程 \(d\rho/da = -\mu(a) \rho(a)\)，且种群增长率 \(r\) 由\textbf{Euler-Lotka特征方程} \(1 = \int_{0}^{\infty} e^{-ra} \beta(a) e^{-\int_{0}^{a} \mu(s) ds} da\) 确定。

\hrulefill

\section{第470章：形态发生与化学反应}\label{ux7b2c470ux7ae0ux5f62ux6001ux53d1ux751fux4e0eux5316ux5b66ux53cdux5e94}

\subsection{Turing不稳定性与斑图形成的数学机制}\label{turingux4e0dux7a33ux5b9aux6027ux4e0eux6591ux56feux5f62ux6210ux7684ux6570ux5b66ux673aux5236}

\textbf{定理 383.1（Turing不稳定性, 1952）} Turing提出，两个扩散速率差异极大的化学物质（形态素，morphogens）通过反应-扩散方程可自发产生空间图案： \[\begin{cases} \frac{\partial u}{\partial t} = D_u \Delta u + f(u, v) \\ \frac{\partial v}{\partial t} = D_v \Delta v + g(u, v) \end{cases}\] 其中 \(u\) 是活化子（activator），\(v\) 是抑制子（inhibitor），\(D_u \ll D_v\)。

\textbf{证明（Turing失稳条件）} 设 \((u^*, v^*)\) 是反应项在无扩散时的均匀稳态。在无扩散时，线性化矩阵 \(J = \begin{pmatrix} f_u & f_v \\ g_u & g_v \end{pmatrix}\) 满足渐近稳定性条件：\(\operatorname{tr} J < 0\)，\(\det J > 0\)。加入扩散后，空间Fourier模 \(e^{i k x}\) 的线性化矩阵为 \(J_k = J - k^2 \operatorname{diag}(D_u, D_v)\)。Turing失稳条件要求存在波数 \(k \neq 0\) 使 \(J_k\) 有正实部特征值。由特征方程的判别式分析可得必要条件： \[D_v f_u + D_u g_v > 2 \sqrt{D_u D_v \det J} > 0,\] 正是``抑制子扩散远快于活化子''的数学表述。\(\blacksquare\)

Turing不稳定性解释了生物斑图（斑马条纹、豹纹、贝壳纹、指纹脊）的自发形成，被视为形态发生的核心数学机制。

\subsection{化学反应的数学建模}\label{ux5316ux5b66ux53cdux5e94ux7684ux6570ux5b66ux5efaux6a21}

\textbf{定理 383.2（质量作用定律）} 对于基元反应 \[\nu_1 A_1 + \nu_2 A_2 + \cdots + \nu_r A_r \xrightarrow{k} \text{产物},\] 反应速率 \(v\) 与反应物浓度的幂次乘积成正比： \[v = k [A_1]^{\nu_1} [A_2]^{\nu_2} \cdots [A_r]^{\nu_r},\] 其中 \(k\) 是速率常数。对非基元反应，速率定律需由实验或由基元步骤推导。

\textbf{证明}：质量作用定律源于分子碰撞理论。对基元反应 \(\sum_i \nu_i A_i \to \text{产物}\)，反应速率正比于各反应物分子同时相遇的概率。在稀溶液中，相遇概率正比于各反应物浓度的乘积。对每分子 \(A_i\)，需要 \(\nu_i\) 个拷贝参与反应，概率正比于 \([A_i]^{\nu_i}\)。因此总速率 \(v = k \prod_i [A_i]^{\nu_i}\)，\(k\) 为比例常数。更严格的推导来自化学动力学的碰撞频率理论或统计力学中的过渡态理论（Eyring方程），速率常数 \(k\) 与温度的关系由Arrhenius方程 \(k = A e^{-E_a/(RT)}\) 给出。对总包反应（非基元），速率定律一般不是简单幂次形式，需由基元步骤通过拟稳态近似或速率决定步骤假设推导。\(\blacksquare\)

\textbf{定理 383.3（Michaelis-Menten酶动力学）} 考虑酶催化反应： \[E + S \underset{k_{-1}}{\stackrel{k_1}{\rightleftharpoons}} ES \xrightarrow{k_2} E + P.\] 在拟稳态假设下（\(\dot{[ES]} \approx 0\)），反应速率 \(v = d[P]/dt\) 为 \[v = \frac{V_{\max} [S]}{K_m + [S]}, \quad V_{\max} = k_2 [E]_T, \quad K_m = \frac{k_{-1} + k_2}{k_1}.\]

\textbf{推导} 酶的质量守恒：\([E]_T = [E] + [ES]\)。拟稳态：\(k_1 [E][S] = (k_{-1} + k_2)[ES]\)。消去 \([E]\)： \[[ES] = \frac{[E]_T [S]}{K_m + [S]}.\] 故 \(v = k_2 [ES] = V_{\max}[S] / (K_m + [S])\)。当 \([S] \ll K_m\) 时 \(v \approx (V_{\max}/K_m)[S]\)（一级动力学）；当 \([S] \gg K_m\) 时 \(v \approx V_{\max}\)（零级动力学，酶饱和）。\(\blacksquare\)

\subsection{神经元动力学}\label{ux795eux7ecfux5143ux52a8ux529bux5b66}

\textbf{定义 383.4（Hodgkin-Huxley模型）} Hodgkin和Huxley（1952）通过鱿鱼巨轴突实验建立了动作电位产生的数学模型： \[C_m \frac{dV}{dt} = I_{\text{ext}} - g_{\text{Na}} m^3 h (V - V_{\text{Na}}) - g_{\text{K}} n^4 (V - V_{\text{K}}) - g_{\text{L}} (V - V_{\text{L}}),\] 其中 \(V\) 是膜电位，门控变量 \(m, h, n\) 满足形式为 \(\frac{dx}{dt} = \alpha_x(V)(1-x) - \beta_x(V)x\) 的动力学方程。

该四维非线性常微分方程组产生了丰富的动力学------静息态、阈值行为、全或无动作电位发放、不应期、以及持续电流刺激下的重复发放模式。

\textbf{定理 383.5（FitzHugh-Nagumo简化模型）} FitzHugh（1961）和Nagumo等人（1962）将HH模型简化为二维系统： \[\begin{cases} \varepsilon \frac{dv}{dt} = v - \frac{v^3}{3} - w + I \\ \frac{dw}{dt} = v + a - b w \end{cases}\] 其中 \(v\) 是快变量（膜电位），\(w\) 是慢变量（恢复变量），\(0 < \varepsilon \ll 1\)。该系统的相图分析揭示了\textbf{兴奋性}（excitability）的数学本质：存在阈值流形，次阈值扰动衰减回静息态，超阈值扰动产生大振幅偏移（动作电位）。零等倾线 \(w = v - v^3/3 + I\)（三次曲线）与 \(w = (v + a)/b\)（直线）的交点决定了平衡态的稳定性和振荡的出现。

\textbf{证明}：FitzHugh-Nagumo系统对快变量 \(v\) 和慢变量 \(w\) 的零等倾线为：\(w = v - v^3/3 + I\)（\(v\)的零等倾线，三次曲线）和 \(w = (v + a)/b\)（\(w\)的零等倾线，直线）。平衡点为两零等倾线的交点。线性化矩阵： \[J = \begin{pmatrix} \varepsilon^{-1}(1 - v^{*2}) & -\varepsilon^{-1} \\ 1 & -b \end{pmatrix}.\] 特征值满足 \(\lambda^2 + (b - \varepsilon^{-1}(1-v^{*2}))\lambda + \varepsilon^{-1}(b(1-v^{*2}) + 1) = 0\)。当交点位于三次曲线中间分支时（\(|v^*| < 1\)），\(1 - v^{*2} > 0\)，对充分小的 \(\varepsilon\) 出现不稳定焦点，由Poincare-Bendixson定理存在极限环（动作电位的周期发放）。当交点位于左右分支时（\(|v^*| > 1\)），平衡点稳定，系统表现为兴奋性：小扰动衰减，大扰动（越过三次曲线的极值点）则产生大振幅偏移后回归平衡。\(\blacksquare\)

\subsection{形态发生的连续介质模型}\label{ux5f62ux6001ux53d1ux751fux7684ux8fdeux7eedux4ecbux8d28ux6a21ux578b}

\textbf{定义 383.6（Murray-Oster模型）} Murray和Oster（1984）发展了基于细胞牵引力的形态发生连续介质模型。总应力张量包含细胞牵引的主动分量： \[\boldsymbol{\sigma} = \boldsymbol{\sigma}_{\text{passive}} + \boldsymbol{\sigma}_{\text{active}}, \quad \boldsymbol{\sigma}_{\text{active}} = \tau \rho \mathbf{I},\] 其中 \(\rho\) 是细胞密度，\(\tau\) 是每个细胞产生的牵引应力。细胞的运动由对流-扩散-趋触方程描述： \[\frac{\partial \rho}{\partial t} + \nabla \cdot (\rho \mathbf{v}) = D \Delta \rho - \nabla \cdot (\chi \rho \nabla c),\] 其中 \(\mathbf{v}\) 是细胞外基质的速度场，\(c\) 是化学趋化剂浓度。力平衡方程 \(\nabla \cdot \boldsymbol{\sigma} = \mathbf{0}\) 与细胞外基质的粘弹性本构关系耦合，构成完整的机械化学模型。

\hrulefill

\section{第471章：进化动力学与系统生物学}\label{ux7b2c471ux7ae0ux8fdbux5316ux52a8ux529bux5b66ux4e0eux7cfbux7edfux751fux7269ux5b66}

\subsection{群体遗传学的数学基础}\label{ux7fa4ux4f53ux9057ux4f20ux5b66ux7684ux6570ux5b66ux57faux7840}

\textbf{定理 384.1（Hardy-Weinberg平衡）} 考虑常染色体二倍体位点，具有等位基因 \(A\)（频率 \(p\)）和 \(a\)（频率 \(q = 1-p\)）。在随机交配、无选择、无突变、无迁移的无限大群体中，基因型频率在传代一次后达到： \[P(AA) = p^2, \quad P(Aa) = 2pq, \quad P(aa) = q^2,\] 并在此后世代中保持不变。

\textbf{证明} 设亲代基因型频率为 \(D, H, R\)（\(AA, Aa, aa\)），等位基因频率 \(p = D + H/2, q = R + H/2\)。随机交配下各交配类型的后代按Mendel比例产生基因型：子代 \(AA\) 频率 = \(D^2 + DH + H^2/4 = (D + H/2)^2 = p^2\)。子代等位基因频率仍为 \(p, q\)。\(\blacksquare\)

\textbf{定理 384.2（Fisher-Wright模型）} 在有随机漂变（genetic drift）的有限群体（大小为 \(N\)）中，等位基因频率 \(p\) 的Wright-Fisher模型是取值于 \(\{0, 1/N, 2/N, \ldots, 1\}\) 的Markov链，转移概率 \[\mathbb{P}\left(p_{t+1} = \frac{j}{2N} \mid p_t = \frac{i}{2N}\right) = \binom{2N}{j} \left(\frac{i}{2N}\right)^j \left(1 - \frac{i}{2N}\right)^{2N - j}.\]

\textbf{证明}：Wright-Fisher模型的转移概率由二项抽样给出：下一代由大小为 \(2N\) 的配子池随机抽样形成，其中等位基因 \(A\) 的配子数为 \(i\)（频率为 \(i/(2N)\)），\(a\) 的配子数为 \(2N - i\)。下一代中 \(A\) 的个数 \(j\) 服从 \(\text{Bin}(2N, i/(2N))\)。种群频率 \(p_t\) 构成马尔可夫链。其扩散近似：当 \(N \to \infty\) 时，对时间尺度 \(t = \tau \cdot 2N\) 重新标度，频率过程弱收敛于Wright-Fisher扩散： \[dp = \sqrt{p(1-p)} dW_t,\] 这是表示纯随机漂变的随机微分方程（无选择、无突变）。扩散过程的生成元为 \(\mathcal{L} = \frac{1}{2} p(1-p) \frac{d^2}{dp^2}\)。\(\blacksquare\)

\textbf{定理 384.3（固定概率与固定时间）} 在无选择的Wright-Fisher模型中，初始频率为 \(p_0\) 的等位基因最终固定的概率恰为 \(p_0\)（鞅性质：\(p_t\) 是有界鞅）。条件固定时间（给定固定发生）的期望为 \(\mathbb{E}[T_{\text{fix}} \mid \text{fixation}] \approx -4N \frac{1-p_0}{p_0} \ln(1-p_0)\)。

\textbf{证明} \(p_t\) 是鞅：\(\mathbb{E}[p_{t+1} \mid p_t] = p_t\)（二项抽样的无偏性）。由鞅收敛定理和有界性得 \(p_t \to p_{\infty}\) a.s.。因吸收态仅为 \(0\) 或 \(1\)，故 \(\mathbb{P}(\text{fixation}) = \mathbb{E}[p_{\infty}] = \mathbb{E}[p_0] = p_0\)。\(\blacksquare\)

\subsection{复制者方程与进化稳定策略}\label{ux590dux5236ux8005ux65b9ux7a0bux4e0eux8fdbux5316ux7a33ux5b9aux7b56ux7565}

\textbf{定义 384.4（复制者方程）} 考虑 \(n\) 种策略（表型）的种群，\(x_i\) 为策略 \(i\) 的频率，支付矩阵 \(A = (a_{ij})\) 给出策略 \(i\) 对策略 \(j\) 的适应度。\textbf{复制者方程}（replicator equation）为 \[\dot{x}_i = x_i \left( (A\mathbf{x})_i - \mathbf{x}^T A \mathbf{x} \right), \quad i = 1, \ldots, n.\] 该方程在单纯形 \(\Delta_n = \{\mathbf{x} \in \mathbb{R}^n : x_i \geq 0, \sum x_i = 1\}\) 上定义了一个动力系统。

\textbf{定理 384.5（复制者方程的基本性质）} （1）单纯形 \(\Delta_n\) 及其子面是正不变的；（2）若策略 \(i\) 的初始频率为零，则始终为零；（3）平均适应度 \(\bar{\phi} = \mathbf{x}^T A \mathbf{x}\) 沿轨道非递减： \[\frac{d\bar{\phi}}{dt} = 2 \sum_i x_i ((A\mathbf{x})_i - \bar{\phi})^2 \geq 0.\] 这是Fisher自然选择基本定理的数学表述------平均适应度沿演化轨道单调递增（但需注意该定理在更一般情形下不成立）。

\textbf{证明}：（1）若 \(\sum_i x_i = 1\)，则 \(\frac{d}{dt} \sum_i x_i = \sum_i \dot{x}_i = \sum_i x_i((A\mathbf{x})_i - \bar{\phi}) = \bar{\phi} - \bar{\phi} \cdot \sum_i x_i = 0\)，故和守恒。若某 \(x_i = 0\) 且 \((A\mathbf{x})_i - \bar{\phi}\) 有限，则 \(\dot{x}_i = 0\)，边界面正不变。（2）由（1），若 \(x_i(0) = 0\) 则 \(\dot{x}_i = 0\)，故恒为零。（3）直接计算： \[\frac{d\bar{\phi}}{dt} = \frac{d}{dt}(\mathbf{x}^T A \mathbf{x}) = \dot{\mathbf{x}}^T A \mathbf{x} + \mathbf{x}^T A \dot{\mathbf{x}} = 2 \mathbf{x}^T A \dot{\mathbf{x}}\] \[= 2 \sum_i x_i (A\mathbf{x})_i \cdot \dot{x}_i/x_i = 2 \sum_i ((A\mathbf{x})_i)^2 x_i - 2 \sum_i (A\mathbf{x})_i x_i \bar{\phi}.\] 第一项中 \(\sum_i ((A\mathbf{x})_i)^2 x_i = \sum_i x_i ((A\mathbf{x})_i - \bar{\phi})^2 + 2\bar{\phi} \sum_i (A\mathbf{x})_i x_i - \bar{\phi}^2 \sum_i x_i\)。代入化简得 \(\frac{d\bar{\phi}}{dt} = 2 \sum_i x_i ((A\mathbf{x})_i - \bar{\phi})^2 \geq 0\)。这正是在策略频率分布方差意义上的Fisher基本定理。\(\blacksquare\)

\textbf{定义 384.6（进化稳定策略, ESS）} 策略 \(\mathbf{x}^*\) 是\textbf{进化稳定策略}（Evolutionarily Stable Strategy），若对任意突变策略 \(\mathbf{y} \neq \mathbf{x}^*\)，存在入侵屏障 \(\bar{\varepsilon} > 0\) 使得对所有 \(0 < \varepsilon < \bar{\varepsilon}\)， \[\mathbf{x}^{*T} A ((1-\varepsilon) \mathbf{x}^* + \varepsilon \mathbf{y}) > \mathbf{y}^T A ((1-\varepsilon) \mathbf{x}^* + \varepsilon \mathbf{y}).\] 即突变的少数个体无法侵入由 \(\mathbf{x}^*\) 主导的种群。

\textbf{定理 384.7（ESS的等价条件 - Maynard Smith-Price, 1973）} \(\mathbf{x}^*\) 是ESS当且仅当：（1）对任意 \(\mathbf{y} \neq \mathbf{x}^*\)，\(\mathbf{x}^{*T} A \mathbf{x}^* \geq \mathbf{y}^T A \mathbf{x}^*\)（Nash均衡条件）；（2）若等号成立，则必须有 \(\mathbf{x}^{*T} A \mathbf{y} > \mathbf{y}^T A \mathbf{y}\)（稳定性条件）。

\textbf{证明} 充分性由定义直接展开验证。必要性：若条件（1）不成立，则 \(\mathbf{y}\) 可入侵。若条件（1）成立且等号成立但条件（2）不成立，则对充分小的 \(\varepsilon\) 有 \(\mathbf{y}^T A ((1-\varepsilon)\mathbf{x}^* + \varepsilon \mathbf{y}) > \mathbf{x}^{*T} A ((1-\varepsilon)\mathbf{x}^* + \varepsilon \mathbf{y})\)，\(\mathbf{y}\) 可入侵。\(\blacksquare\)

\textbf{定理 384.8（ESS与复制者动力学的联系）} 若 \(\mathbf{x}^*\) 是ESS，则它是复制者方程的渐近稳定平衡点。

\textbf{证明}：设 \(\mathbf{x}^*\) 为ESS。考虑Lyapunov函数 \(V(\mathbf{x}) = \prod_{i=1}^n x_i^{x_i^*}\)（在 \(\mathbf{x}^*\) 的支持集上，\(\mathbf{x}^*\) 的每个正分量对应 \(x_i > 0\)）。则 \[\frac{d}{dt} \ln V(\mathbf{x}) = \sum_i x_i^* \frac{\dot{x}_i}{x_i} = \sum_i x_i^* ((A\mathbf{x})_i - \bar{\phi}) = \mathbf{x}^{*T} A \mathbf{x} - \bar{\phi}.\] 由ESS的定义条件，在 \(\mathbf{x}^*\) 的某邻域内 \(\mathbf{x}^{*T} A \mathbf{x} > \mathbf{x}^T A \mathbf{x} = \bar{\phi}\)（严格不等式），故 \(\dot{V} > 0\)。同时 \(V\) 在单纯形上有界（\(\leq 1\)），在 \(\mathbf{x}^*\) 处取得最大值。由LaSalle不变性原理，\(\mathbf{x}^*\) 是渐近稳定的。更精确的论证可证明 \(\mathbf{x}^*\) 是进化稳定状态，任何偏离 \(\mathbf{x}^*\) 的\(\varepsilon\)-扰动在复制者动力学下都会回到 \(\mathbf{x}^*\)。\(\blacksquare\)

\subsection{基因调控网络的数学模型}\label{ux57faux56e0ux8c03ux63a7ux7f51ux7edcux7684ux6570ux5b66ux6a21ux578b}

\textbf{定义 384.9（Boolean网络模型）} Kauffman（1969）提出的基因调控Boolean网络模型中，每个基因的状态为ON（1）或OFF（0），更新规则为Boolean函数。\(N\) 个基因的状态空间为 \(\{0, 1\}^N\)（大小为 \(2^N\)）。若每个基因由 \(K\) 个其他基因调控（\(K\)-输入网络），状态转移图具有丰富的动力学特征：存在吸引子（attractors）------固定点和极限环，被认为对应不同的细胞类型。

\textbf{定义 384.10（ODE基因调控模型）} 基因调控网络也可用常微分方程建模： \[\frac{dx_i}{dt} = - \gamma_i x_i + f_i(x_1, \ldots, x_N),\] 其中 \(x_i\) 是基因 \(i\) 的表达水平，\(\gamma_i\) 是降解率，\(f_i\) 是调控函数（通常为Hill函数或sigmoid函数之和）。

\textbf{定义 384.11（随机基因表达模型）} 由于细胞内mRNA分子数少（可低至个位数），内在噪音必须考虑。化学主方程（Chemical Master Equation, CME）描述各分子种类计数的概率分布 \(P(\mathbf{n}, t)\) 的演化： \[\frac{\partial P(\mathbf{n}, t)}{\partial t} = \sum_{j} \left[ a_j(\mathbf{n} - \boldsymbol{\nu}_j) P(\mathbf{n} - \boldsymbol{\nu}_j, t) - a_j(\mathbf{n}) P(\mathbf{n}, t) \right],\] 其中 \(a_j(\mathbf{n})\) 是反应 \(j\) 的倾向函数，\(\boldsymbol{\nu}_j\) 是化学计量向量。

\subsection{系统生物学中的网络分析}\label{ux7cfbux7edfux751fux7269ux5b66ux4e2dux7684ux7f51ux7edcux5206ux6790}

\textbf{定义 384.12（代谢网络分析 - 通量平衡分析）} 代谢网络在拟稳态假设下的通量分布 \(\mathbf{v} \in \mathbb{R}^m\)（\(m\) 为反应数）满足化学计量约束： \[\mathbf{S} \mathbf{v} = \mathbf{0}, \quad v_j^{\min} \leq v_j \leq v_j^{\max},\] 其中 \(\mathbf{S} \in \mathbb{R}^{n \times m}\) 是化学计量矩阵（\(n\) 为代谢产物数）。\textbf{通量平衡分析}（Flux Balance Analysis, FBA）在此基础上优化生物量目标函数 \(\mathbf{c}^T \mathbf{v}\)，转化为线性规划问题。FBA已成功预测多种微生物在给定环境条件下的生长速率和代谢通量分布。

\textbf{定义 384.13（信号传导的数学建模）} 信号传导级联通常采用质量作用动力学或Michaelis-Menten近似的ODE建模互作的磷酸化-去磷酸化循环。Goldbeter-Koshland（1981）揭示了由拮抗激酶-磷酸酶对驱动的共价修饰循环中\textbf{零级超敏感性}（zeroth-order ultrasensitivity）的数学根源：当酶浓度远低于底物浓度时，输入信号（激酶活性）的微小变化可导致输出（磷酸化底物比例）的急剧开关式切换。

\textbf{定理 384.14（Goldbeter-Koshland方程的稳态）} 设 \(W\) 和 \(W^*\) 分别为非磷酸化和磷酸化底物的浓度（总浓度 \(W_T = W + W^*\)），激酶和磷酸酶均遵循Michaelis-Menten动力学。稳态时， \[\frac{W^*}{W_T} = G(v_1, v_2, K_1, K_2),\] 其中 \(v_1, v_2\) 是最大反应速率，\(K_1, K_2\) 是以 \(W_T\) 为尺度的归一化Michaelis常数。当 \(K_1, K_2 \ll 1\) 时（零级区域），响应曲线呈极陡的S形------零级超敏感性。这是细胞内信号开关的数学基础。

\newpage

\chapter{13-数学物理/数学物理/10-流体力学数学理论.md}\label{ux6570ux5b66ux7269ux7406ux6570ux5b66ux7269ux740610-ux6d41ux4f53ux529bux5b66ux6570ux5b66ux7406ux8bba.md}

\chapter{流体力学数学理论}\label{ux6d41ux4f53ux529bux5b66ux6570ux5b66ux7406ux8bba}

\section{第472章：流体力学基本方程}\label{ux7b2c472ux7ae0ux6d41ux4f53ux529bux5b66ux57faux672cux65b9ux7a0b}

流体力学数学理论以Navier-Stokes方程和Euler方程为核心，研究流体质点的运动规律、能量传递和涡量演化。本章从连续介质假设出发，依据质量守恒、动量守恒和能量守恒三大物理定律，系统推导Euler方程和Navier-Stokes方程。涡量方程揭示了三维流动中涡量拉伸这一非线性机制的本质，是理解湍流级串的数学起点。Bernoulli定理作为理想流体沿流线的能量守恒原理，其可压缩与不可压缩形式的推导展示了热力学变量与运动学变量间的深刻耦合。守恒律形式为后续弱解理论提供了基本的分析框架。

\subsection{x.1 流体运动的运动学描述}\label{x.1-ux6d41ux4f53ux8fd0ux52a8ux7684ux8fd0ux52a8ux5b66ux63cfux8ff0}

\textbf{定义 x.1}（物质导数）：设 \(\mathbf{u}(\mathbf{x}, t) : \mathbb{R}^3 \times [0, T] \to \mathbb{R}^3\) 为速度场。标量场 \(\phi(\mathbf{x}, t)\) 的物质导数为

\[\frac{D\phi}{Dt} = \frac{\partial \phi}{\partial t} + (\mathbf{u} \cdot \nabla)\phi.\]

物质导数描述跟随流体质点运动的物理量的变化率。

\textbf{定理 x.2}（Reynolds输运定理）：设 \(\Omega(t)\) 为随流体运动的材料体积，\(V(t) \subseteq \mathbb{R}^3\)。对任意光滑标量场 \(f\)，

\[\frac{d}{dt}\int_{\Omega(t)} f \, d\mathbf{x} = \int_{\Omega(t)} \left(\frac{\partial f}{\partial t} + \nabla \cdot (f\mathbf{u})\right) d\mathbf{x}.\]

\textbf{证明}：令 \(\mathbf{x} = \boldsymbol{\varphi}(\mathbf{X}, t)\) 为流映射，\(\mathbf{X}\) 为Lagrangian坐标。Jacobi行列式 \(J = \det(\nabla_{\mathbf{X}} \boldsymbol{\varphi})\) 满足 \(\frac{\partial J}{\partial t} = J \nabla \cdot \mathbf{u}\)（由行列式的导数公式及 \(\frac{\partial}{\partial t}\nabla_{\mathbf{X}} \boldsymbol{\varphi} = \nabla_{\mathbf{X}} \mathbf{u}\) 通过恒等式验证）。则

\[\frac{d}{dt}\int_{\Omega(t)} f \, d\mathbf{x} = \frac{d}{dt}\int_{\Omega(0)} f(\boldsymbol{\varphi}(\mathbf{X}, t), t) J \, d\mathbf{X} = \int_{\Omega(0)} \left(\frac{Df}{Dt} J + f \frac{\partial J}{\partial t}\right) d\mathbf{X}.\]

代入 \(\frac{Df}{Dt} = \frac{\partial f}{\partial t} + \mathbf{u} \cdot \nabla f\) 和 \(\frac{\partial J}{\partial t} = J \nabla \cdot \mathbf{u}\)，并变换回Euler坐标即得。 \(\blacksquare\)

\subsection{x.2 基本方程组的推导}\label{x.2-ux57faux672cux65b9ux7a0bux7ec4ux7684ux63a8ux5bfc}

\textbf{定理 x.3}（质量守恒------连续性方程）：质量守恒律在微分形式下等价于

\[\frac{\partial \rho}{\partial t} + \nabla \cdot (\rho \mathbf{u}) = 0.\]

对于不可压缩流体（\(\rho = \text{const}\)），简化为 \(\nabla \cdot \mathbf{u} = 0\)。

\textbf{证明}：令 \(\Omega\) 为任意固定的控制体积。\(\Omega\) 内质量的减少率等于通过边界 \(\partial \Omega\) 的净流出通量：

\[\frac{d}{dt}\int_{\Omega} \rho \, d\mathbf{x} = -\int_{\partial \Omega} \rho \mathbf{u} \cdot \mathbf{n} \, dS.\]

利用散度定理和积分号下求导，\(\int_{\Omega} (\frac{\partial \rho}{\partial t} + \nabla \cdot (\rho \mathbf{u})) d\mathbf{x} = 0\)。由 \(\Omega\) 的任意性，被积函数处处为零。 \(\blacksquare\)

\textbf{定理 x.4}（Cauchy动量方程与应力张量）：连续介质的动量守恒给出

\[\rho \frac{D\mathbf{u}}{Dt} = \nabla \cdot \boldsymbol{\sigma} + \rho \mathbf{f},\]

其中 \(\boldsymbol{\sigma}\) 为Cauchy应力张量，\(\mathbf{f}\) 为体积力密度。对Newton流体，\(\boldsymbol{\sigma} = -p\mathbf{I} + 2\mu \mathbf{D} + \lambda (\nabla \cdot \mathbf{u})\mathbf{I}\)，\(\mathbf{D} = \frac{1}{2}(\nabla \mathbf{u} + (\nabla \mathbf{u})^T)\) 为应变率张量。

\textbf{证明}：作用于材料体积 \(\Omega(t)\) 上的合力等于动量的物质导数：

\[\frac{d}{dt}\int_{\Omega(t)} \rho \mathbf{u} \, d\mathbf{x} = \int_{\partial \Omega(t)} \boldsymbol{\sigma} \cdot \mathbf{n} \, dS + \int_{\Omega(t)} \rho \mathbf{f} \, d\mathbf{x}.\]

由Reynolds输运定理，左边 \(= \int_{\Omega(t)} (\frac{\partial (\rho \mathbf{u})}{\partial t} + \nabla \cdot (\rho \mathbf{u} \otimes \mathbf{u})) d\mathbf{x}\)。利用连续性方程化简，得 \(\int_{\Omega(t)} \rho \frac{D\mathbf{u}}{Dt} d\mathbf{x}\)。右边由散度定理化为体积分。由 \(\Omega(t)\) 任意性得微分形式。Newton流体的本构关系基于应力张量与应变率张量的线性各向同性假设。 \(\blacksquare\)

\textbf{定义 x.5}（Navier-Stokes方程与Euler方程）：将Newton流体本构关系代入Cauchy动量方程，得可压缩Navier-Stokes方程

\[\rho\left(\frac{\partial \mathbf{u}}{\partial t} + (\mathbf{u} \cdot \nabla)\mathbf{u}\right) = -\nabla p + \mu \Delta \mathbf{u} + (\lambda + \mu)\nabla(\nabla \cdot \mathbf{u}) + \rho \mathbf{f}.\]

不可压缩（\(\nabla \cdot \mathbf{u} = 0\)，\(\rho = \text{const}\)）Navier-Stokes方程为

\[\frac{\partial \mathbf{u}}{\partial t} + (\mathbf{u} \cdot \nabla)\mathbf{u} = -\frac{1}{\rho}\nabla p + \nu \Delta \mathbf{u} + \mathbf{f}, \quad \nabla \cdot \mathbf{u} = 0.\]

取 \(\nu = 0\)（无粘性）得不可压缩Euler方程。

\subsection{x.3 涡量方程与Bernoulli定理}\label{x.3-ux6da1ux91cfux65b9ux7a0bux4e0ebernoulliux5b9aux7406}

\textbf{定义 x.6}（涡量）：涡量场定义为 \(\boldsymbol{\omega} = \nabla \times \mathbf{u}\)。涡量描述流体元的局部旋转速率，其方向沿旋转轴。

\textbf{定理 x.7}（不可压缩涡量输运方程）：对不可压缩Navier-Stokes流动（\(\rho\) 常数），涡量满足

\[\frac{\partial \boldsymbol{\omega}}{\partial t} + (\mathbf{u} \cdot \nabla)\boldsymbol{\omega} = (\boldsymbol{\omega} \cdot \nabla)\mathbf{u} + \nu \Delta \boldsymbol{\omega}.\]

\textbf{证明}：对不可压缩Navier-Stokes方程取旋度。利用向量恒等式 \(\nabla \times ((\mathbf{u} \cdot \nabla)\mathbf{u}) = (\mathbf{u} \cdot \nabla)(\nabla \times \mathbf{u}) - ((\nabla \times \mathbf{u}) \cdot \nabla)\mathbf{u} + (\nabla \cdot \mathbf{u})(\nabla \times \mathbf{u}) = (\mathbf{u} \cdot \nabla)\boldsymbol{\omega} - (\boldsymbol{\omega} \cdot \nabla)\mathbf{u}\)（因 \(\nabla \cdot \mathbf{u} = 0\) 且 \(\nabla \times (\nabla \phi) = 0\)）。又 \(\nabla \times (\nu \Delta \mathbf{u}) = \nu \Delta \boldsymbol{\omega}\)（旋度与Laplace算子交换），即得方程。 \(\blacksquare\)

\textbf{定理 x.8}（Euler方程的Bernoulli定理------不可压缩情形）：设 \(\mathbf{u}\) 为无粘不可压缩定常流动（\(\frac{\partial \mathbf{u}}{\partial t} = 0\)）且体积力有势 \(\mathbf{f} = -\nabla \Phi\)。则沿流线，

\[\frac{1}{2}|\mathbf{u}|^2 + \frac{p}{\rho} + \Phi = \text{const}.\]

若流动无旋（\(\boldsymbol{\omega} = 0\)），则常数为全局常数（全流场）。

\textbf{证明}：利用恒等式 \((\mathbf{u} \cdot \nabla)\mathbf{u} = \nabla(\frac{1}{2}|\mathbf{u}|^2) - \mathbf{u} \times \boldsymbol{\omega}\)。Euler方程可写为

\[\nabla\left(\frac{1}{2}|\mathbf{u}|^2 + \frac{p}{\rho} + \Phi\right) = \mathbf{u} \times \boldsymbol{\omega}.\]

左端沿流线的切向分量为零（因 \(\mathbf{u} \cdot (\mathbf{u} \times \boldsymbol{\omega}) = 0\)），故沿流线括号内为常数。若无旋（\(\boldsymbol{\omega} = 0\)），则梯度为零，括号内全局为常数。 \(\blacksquare\)

\textbf{定理 x.9}（可压缩流体的Bernoulli定理------等熵流动）：对可压缩无粘等熵定常流动（\(p = C\rho^{\gamma}\)），沿流线有

\[\frac{1}{2}|\mathbf{u}|^2 + \frac{\gamma}{\gamma-1}\frac{p}{\rho} = \text{const}.\]

\textbf{证明}：可压缩Euler方程为 \(\rho(\mathbf{u} \cdot \nabla)\mathbf{u} = -\nabla p\)（无体积力）。对等熵关系 \(p = C\rho^{\gamma}\)，有 \(\frac{\nabla p}{\rho} = \nabla\left(\frac{\gamma}{\gamma-1}\frac{p}{\rho}\right)\)（验证：计算 \(\nabla(\frac{p}{\rho}) = \frac{\nabla p}{\rho} - \frac{p}{\rho^2}\nabla\rho\)，结合 \(\nabla p = C\gamma \rho^{\gamma-1}\nabla\rho\) 即得 \(\frac{\nabla p}{\rho} = \frac{\gamma}{\gamma-1}\nabla(\frac{p}{\rho})\)）。代入方程并沿流线点乘 \(\mathbf{u}\)，得 \(\mathbf{u} \cdot \nabla(\frac{1}{2}|\mathbf{u}|^2 + \frac{\gamma}{\gamma-1}\frac{p}{\rho}) = 0\)，故沿流线守恒。 \(\blacksquare\)

\subsection{x.4 守恒律与弱解}\label{x.4-ux5b88ux6052ux5f8bux4e0eux5f31ux89e3}

\textbf{定理 x.10}（不可压缩Navier-Stokes方程的能量守恒------光滑解）：设 \((\mathbf{u}, p)\) 为不可压缩NS方程的充分光滑解（\(\Omega = \mathbb{R}^3\) 或周期区域，无外力）。则动能满足

\[\frac{1}{2}\frac{d}{dt}\int_{\Omega} |\mathbf{u}|^2 \, d\mathbf{x} = -\nu \int_{\Omega} |\nabla \mathbf{u}|^2 \, d\mathbf{x}.\]

\textbf{证明}：方程两端与 \(\mathbf{u}\) 作 \(L^2\) 内积（在无外力情形）：

\[\int_{\Omega} \mathbf{u} \cdot \frac{\partial \mathbf{u}}{\partial t} \, d\mathbf{x} + \int_{\Omega} \mathbf{u} \cdot (\mathbf{u} \cdot \nabla)\mathbf{u} \, d\mathbf{x} = -\int_{\Omega} \mathbf{u} \cdot \nabla p \, d\mathbf{x} + \nu \int_{\Omega} \mathbf{u} \cdot \Delta \mathbf{u} \, d\mathbf{x}.\]

利用分部积分：\(\int \mathbf{u} \cdot (\mathbf{u} \cdot \nabla)\mathbf{u} = \frac{1}{2}\int \mathbf{u} \cdot \nabla|\mathbf{u}|^2 = 0\)（因 \(\nabla \cdot \mathbf{u} = 0\)）。同样 \(\int \mathbf{u} \cdot \nabla p = -\int p \nabla \cdot \mathbf{u} = 0\)。最后 \(\int \mathbf{u} \cdot \Delta \mathbf{u} = -\int |\nabla \mathbf{u}|^2\)。综合即得。 \(\blacksquare\)

\hrulefill

\section{第473章：Navier-Stokes方程解的存在性理论}\label{ux7b2c473ux7ae0navier-stokesux65b9ux7a0bux89e3ux7684ux5b58ux5728ux6027ux7406ux8bba}

Navier-Stokes方程解的存在性与正则性是偏微分方程理论中最具挑战性的核心问题之一。自Leray于1934年通过先验估计和紧性方法证明三维不可压缩NS方程整体弱解的存在性以来，弱解的构造理论已臻于完善，但弱解的正则性和唯一性仍然是Clay数学研究所悬赏百万美元的千禧难题。本章系统建立Leray-Hopf弱解的存在性理论，给出能量等式与能量不等式的完整推导，证明弱-强唯一性定理，并讨论Serrin正则性准则等正则性判据。最后，我们综述三维NS方程光滑解全局存在性这一开问题的当前进展和主要技术障碍。

\subsection{x+1.1 Leray-Hopf弱解的存在性}\label{x1.1-leray-hopfux5f31ux89e3ux7684ux5b58ux5728ux6027}

\textbf{定义 x+1.1}（Leray-Hopf弱解）：设 \(\Omega \subseteq \mathbb{R}^3\)，初值 \(\mathbf{u}_0 \in L^2(\Omega)\) 且 \(\nabla \cdot \mathbf{u}_0 = 0\)。函数 \(\mathbf{u} \in L^\infty(0, T; L^2(\Omega)) \cap L^2(0, T; H^1(\Omega))\) 称为不可压缩NS方程的Leray-Hopf弱解，若：

\begin{enumerate}
\def\labelenumi{(\roman{enumi})}
\tightlist
\item
  对任意 \(\boldsymbol{\phi} \in C_c^\infty(\Omega \times [0, T); \mathbb{R}^3)\) 且 \(\nabla \cdot \boldsymbol{\phi} = 0\)，
\end{enumerate}

\[\int_0^T \int_{\Omega} \left(-\mathbf{u} \cdot \frac{\partial \boldsymbol{\phi}}{\partial t} + \nu \nabla \mathbf{u} : \nabla \boldsymbol{\phi} + (\mathbf{u} \otimes \mathbf{u}) : \nabla \boldsymbol{\phi}\right) d\mathbf{x} \, dt = \int_{\Omega} \mathbf{u}_0 \cdot \boldsymbol{\phi}(\cdot, 0) \, d\mathbf{x};\]

\begin{enumerate}
\def\labelenumi{(\roman{enumi})}
\setcounter{enumi}{1}
\item
  \(\nabla \cdot \mathbf{u} = 0\) 在分布意义下成立；
\item
  能量不等式：对几乎处处的 \(t \in [0, T]\)，
\end{enumerate}

\[\frac{1}{2}\|\mathbf{u}(t)\|_{L^2}^2 + \nu \int_0^t \|\nabla \mathbf{u}(s)\|_{L^2}^2 \, ds \leq \frac{1}{2}\|\mathbf{u}_0\|_{L^2}^2.\]

\textbf{定理 x+1.2}（Leray, 1934------整体弱解的存在性）：设 \(\mathbf{u}_0 \in L^2(\mathbb{R}^3)\) 且 \(\nabla \cdot \mathbf{u}_0 = 0\)。则存在整体Leray-Hopf弱解 \(\mathbf{u}\) 定义于 \([0, \infty)\)。

\textbf{证明}（核心步骤）：采用Galerkin逼近法。(1) 取 \(L^2\) 散度自由向量场的标准正交基 \(\{\mathbf{w}_k\}\)（如Stokes算子的特征函数）。构造逼近解 \(\mathbf{u}_n = \sum_{k=1}^{n} c_k^n(t) \mathbf{w}_k\) 满足

\[\left\langle \frac{\partial \mathbf{u}_n}{\partial t}, \mathbf{w}_k \right\rangle + \nu \langle \nabla \mathbf{u}_n, \nabla \mathbf{w}_k \rangle + \langle (\mathbf{u}_n \cdot \nabla)\mathbf{u}_n, \mathbf{w}_k \rangle = 0, \quad k = 1, \ldots, n.\]

\begin{enumerate}
\def\labelenumi{(\arabic{enumi})}
\setcounter{enumi}{1}
\tightlist
\item
  这是关于 \(c_k^n(t)\) 的ODE系统，由Cauchy-Lipschitz定理（或Peano定理）局部可解。由能量等式（ODE水平）
\end{enumerate}

\[\frac{1}{2}\frac{d}{dt}\|\mathbf{u}_n\|_{L^2}^2 + \nu \|\nabla \mathbf{u}_n\|_{L^2}^2 = 0,\]

积分得 \(\|\mathbf{u}_n(t)\|_{L^2}^2 + 2\nu \int_0^t \|\nabla \mathbf{u}_n\|_{L^2}^2 \leq \|\mathbf{u}_0\|_{L^2}^2\)，从而整体可解。

\begin{enumerate}
\def\labelenumi{(\arabic{enumi})}
\setcounter{enumi}{2}
\tightlist
\item
  取子列：由先验估计，\(\{\mathbf{u}_n\}\) 在 \(L^\infty(0, T; L^2) \cap L^2(0, T; H^1)\) 中有界。由Banach-Alaoglu定理及自反性，可抽取弱-*收敛子列（仍记为 \(\mathbf{u}_n\)）。关键是非线性项的紧性：由Aubin-Lions紧嵌入定理，\(\{\mathbf{u}_n\}\) 在 \(L^2(0, T; L^2_{\text{loc}})\) 中强收敛（由 \(L^2 H^1 \cap H^{-1} L^2 \hookrightarrow L^2 L^2\) 紧），故非线性项 \(\mathbf{u}_n \otimes \mathbf{u}_n\) 在分布意义下收敛至 \(\mathbf{u} \otimes \mathbf{u}\)。取极限即得弱解。能量不等式由弱下半连续性可得。 \(\blacksquare\)
\end{enumerate}

\subsection{x+1.2 能量等式与弱-强唯一性}\label{x1.2-ux80fdux91cfux7b49ux5f0fux4e0eux5f31-ux5f3aux552fux4e00ux6027}

\textbf{定理 x+1.3}（能量等式）：若Leray-Hopf弱解 \(\mathbf{u}\) 满足 \(\mathbf{u} \in L^4(0, T; L^4(\Omega))\)，则能量等式成立：

\[\frac{1}{2}\|\mathbf{u}(t)\|_{L^2}^2 + \nu \int_0^t \|\nabla \mathbf{u}(s)\|_{L^2}^2 \, ds = \frac{1}{2}\|\mathbf{u}_0\|_{L^2}^2.\]

特别地，二维弱解总满足能量等式（因 \(L^\infty L^2 \cap L^2 H^1 \subset L^4 L^4\) 在二维成立）。

\textbf{证明}：取 \(\mathbf{u}\) 自身为检验函数（需正则化处理）。构造时空磨光 \(\mathbf{u}_\varepsilon = \eta_\varepsilon * \mathbf{u}\)。将 \(\mathbf{u}_\varepsilon\) 代入弱公式并在极限 \(\varepsilon \to 0\) 下取极限。非线性项的处理是关键：

\[\int_0^t \int_{\Omega} (\mathbf{u} \otimes \mathbf{u}) : \nabla \mathbf{u}_\varepsilon \to \int_0^t \int_{\Omega} (\mathbf{u} \otimes \mathbf{u}) : \nabla \mathbf{u},\]

需 \(\mathbf{u} \in L^4_{t,x}\) 以保证极限交换的合法性。若 \(\mathbf{u} \in L^4_{t,x}\)，则 \(\mathbf{u} \otimes \mathbf{u} \in L^2_{t,x}\)，且 \(\nabla \mathbf{u}_\varepsilon \to \nabla \mathbf{u}\) 弱收敛于 \(L^2_{t,x}\)，乘积收敛。最终得能量等式。 \(\blacksquare\)

\textbf{定理 x+1.4}（弱-强唯一性，Serrin 1963 / Prodi 1959）：设 \(\mathbf{u}\) 为Leray-Hopf弱解，\(\mathbf{v}\) 为具有相同初值的强解且 \(\mathbf{v} \in L^r(0, T; L^s(\Omega))\) 满足 \(\frac{2}{r} + \frac{3}{s} = 1\)，\(s > 3\)。则 \(\mathbf{u} = \mathbf{v}\) 几乎处处成立。

\textbf{证明}：记 \(\mathbf{w} = \mathbf{u} - \mathbf{v}\)。\(\mathbf{w}\) 满足的弱形式与 \(\mathbf{w}\) 自身作内积（经正则化），利用 \(\mathbf{v}\) 的正则性可控制交叉项：

\[\frac{1}{2}\frac{d}{dt}\|\mathbf{w}\|_{L^2}^2 + \nu \|\nabla \mathbf{w}\|_{L^2}^2 = -\int (\mathbf{w} \cdot \nabla)\mathbf{v} \cdot \mathbf{w} \, d\mathbf{x}.\]

（非线性项 \(\int (\mathbf{w} \cdot \nabla)\mathbf{w} \cdot \mathbf{v} = 0\) 由 \(\nabla \cdot \mathbf{w} = 0\) 分部积分。）

由Holder不等式及Sobolev不等式，对 \(s > 3\)，

\[\left|\int (\mathbf{w} \cdot \nabla)\mathbf{v} \cdot \mathbf{w} \, d\mathbf{x}\right| \leq \|\nabla \mathbf{v}\|_{L^s} \|\mathbf{w}\|_{L^{2s/(s-2)}}^2 \leq C \|\nabla \mathbf{v}\|_{L^s} \|\mathbf{w}\|_{L^2}^{2(1-3/s)} \|\nabla \mathbf{w}\|_{L^2}^{6/s}.\]

利用Young不等式吸收 \(\|\nabla \mathbf{w}\|_{L^2}^2\) 项。得

\[\frac{d}{dt}\|\mathbf{w}\|_{L^2}^2 \leq C \|\nabla \mathbf{v}\|_{L^s}^{\frac{2s}{s-3}} \|\mathbf{w}\|_{L^2}^2.\]

由 \(\mathbf{v} \in L^r L^s\) 及 \(\frac{2}{r} + \frac{3}{s} = 1\)，得 \(\|\nabla \mathbf{v}\|_{L^s}^{\frac{2s}{s-3}} \in L^1(0, T)\)。Gronwall不等式给出 \(\|\mathbf{w}(t)\|_{L^2}^2 = 0\)（因 \(\|\mathbf{w}(0)\| = 0\)）。 \(\blacksquare\)

\subsection{x+1.3 Serrin正则性准则与光滑解的开问题}\label{x1.3-serrinux6b63ux5219ux6027ux51c6ux5219ux4e0eux5149ux6ed1ux89e3ux7684ux5f00ux95eeux9898}

\textbf{定理 x+1.5}（Serrin正则性准则，1962）：设 \(\mathbf{u}\) 为三维不可压缩NS方程的Leray-Hopf弱解。若存在 \(r, s\) 满足 \(\frac{2}{r} + \frac{3}{s} \leq 1\)，\(s \geq 3\)，使得 \(\mathbf{u} \in L^r(0, T; L^s(\mathbb{R}^3))\)，则 \(\mathbf{u}\) 在 \((0, T] \times \mathbb{R}^3\) 上光滑且唯一。

\textbf{证明概要}：由能量不等式和Sobolev嵌入，对时间区间迭代提升正则性。基本引理：对任意 \(k \geq 1\)，成立

\[\frac{1}{2}\frac{d}{dt}\|\nabla^k \mathbf{u}\|_{L^2}^2 + \nu \|\nabla^{k+1} \mathbf{u}\|_{L^2}^2 = -\int \nabla^k((\mathbf{u} \cdot \nabla)\mathbf{u}) \cdot \nabla^k \mathbf{u}.\]

估计非线性项：\(\int (\mathbf{u} \cdot \nabla)\nabla^k \mathbf{u} \cdot \nabla^k \mathbf{u} = 0\)（分部积分），仅需估计交换子项 \(\int [\nabla^k, \mathbf{u} \cdot \nabla]\mathbf{u} \cdot \nabla^k \mathbf{u}\)。由Holder和Gagliardo-Nirenberg不等式，对于 \(\frac{2}{r} + \frac{3}{s} = 1\)，

\[\|[ \nabla^k, \mathbf{u} \cdot \nabla]\mathbf{u}\|_{L^2} \leq C \|\mathbf{u}\|_{L^s} \|\nabla^{k+1} \mathbf{u}\|_{L^2}^{1-\theta} \|\nabla^k \mathbf{u}\|_{L^2}^{\theta}.\]

代入并利用Young不等式，得到关于 \(\|\nabla^k \mathbf{u}\|_{L^2}\) 的Gronwall型不等式，逐步提升正则性至 \(C^\infty\)。当 \(\frac{2}{r} + \frac{3}{s} < 1\) 时结论更强。 \(\blacksquare\)

\textbf{定理 x+1.6}（三维NS光滑解全局存在性------开问题综述）：三维不可压缩Navier-Stokes方程光滑初值（\(\mathbf{u}_0 \in H^s\)，\(s > 1/2\)）的全局光滑解的存在性或有限时间Blow-up的发生仍是未解决的千禧难题。已知的部分结果包括：

\begin{enumerate}
\def\labelenumi{(\roman{enumi})}
\item
  若 \(\|\mathbf{u}(t)\|_{L^3}\) 在 \([0, T)\) 上有界，则光滑性保持（Escauriaza-Seregin-Sverak, 2003）；
\item
  轴对称无旋流不产生有限时间奇性（Ladyzhenskaya, 1969 / Ukhovskii-Yudovich, 1968）；
\item
  若Blow-up发生在时刻 \(T_*\)，则 \(\limsup_{t \to T_*} \|\nabla \mathbf{u}(t)\|_{L^\infty} = \infty\)（Beale-Kato-Majda准则）。
\end{enumerate}

\textbf{证明思路（关键点的启发式论证）}：问题的核心在于三维涡量拉伸项 \((\boldsymbol{\omega} \cdot \nabla)\mathbf{u}\) 可能产生正反馈导致涡量在有限时间内无限增长。能量估计仅能控制 \(\|\mathbf{u}\|_{L^\infty L^2}\) 和 \(\|\nabla \mathbf{u}\|_{L^2 L^2}\)，但不足以控制 \(\|\nabla \mathbf{u}\|_{L^\infty}\)。在二维情形下，涡量拉伸项消失（因 \(\boldsymbol{\omega}\) 垂直于速度平面），能量估计足以给出全局正则性，这解释了二维与三维的实质性差异。 \(\blacksquare\)

\hrulefill

\section{第474章：湍流与统计理论}\label{ux7b2c474ux7ae0ux6e4dux6d41ux4e0eux7edfux8ba1ux7406ux8bba}

湍流是流体力学中最复杂也最迷人的现象之一，其本质特征是多尺度非线性相互作用导致的能量级串和强间歇性。Kolmogorov于1941年提出的局部各向同性湍流理论是湍流统计理论的历史性飞跃，通过量纲分析导出了惯性子区的2/3结构函数律和-5/3能谱律。K41理论的核心支柱------4/5律------是Navier-Stokes方程在统计意义上的精确结果，不依赖于任何唯象假设。本章完整推导Kolmogorov 4/5律，讨论间歇性导致的标度指数反常问题，并介绍Onsager关于能量守恒临界正则性的著名猜想及其最近的数学解决方案。

\subsection{x+2.1 Kolmogorov 1941理论}\label{x2.1-kolmogorov-1941ux7406ux8bba}

\textbf{定义 x+2.1}（速度增量与结构函数）：对均匀各向同性湍流，定义纵向速度增量

\[\delta u_{\parallel}(\mathbf{x}, \mathbf{r}) = (\mathbf{u}(\mathbf{x} + \mathbf{r}) - \mathbf{u}(\mathbf{x})) \cdot \frac{\mathbf{r}}{|\mathbf{r}|}.\]

\(p\) 阶纵向结构函数为 \(S_p(r) = \langle (\delta u_{\parallel})^p \rangle\)，其中 \(\langle \cdot \rangle\) 表示系综平均（等价于空间平均，由遍历性假设）。

\textbf{定理 x+2.2}（Kolmogorov 2/3律------K41的推论）：在惯性子区（\(\eta \ll r \ll L\)，\(\eta\) 为Kolmogorov耗散尺度，\(L\) 为积分尺度），充分发展湍流的二阶纵向结构函数满足

\[S_2(r) = C_K (\bar{\varepsilon} r)^{2/3},\]

其中 \(\bar{\varepsilon}\) 为平均能量耗散率，\(C_K\) 为Kolmogorov普适常数。

\textbf{证明}（量纲分析）：由K41第一相似性假说：惯性子区内的统计量仅依赖于 \(\bar{\varepsilon}\) 和 \(r\)（不依赖于 \(\nu\)）。由量纲分析，\([S_2] = L^2 T^{-2}\)，\([\bar{\varepsilon}] = L^2 T^{-3}\)，\([r] = L\)。构造无量纲量 \(S_2/(\bar{\varepsilon}^{2/3} r^{2/3})\)，在惯性子区内必为常数（假设完全自相似性）。此常数记为 \(C_K\)。 \(\blacksquare\)

\textbf{定理 x+2.3}（Kolmogorov -5/3能谱律）：在惯性子区，湍流动能谱 \(E(k)\) 满足

\[E(k) = C \bar{\varepsilon}^{2/3} k^{-5/3},\]

其中 \(C\) 为Kolmogorov谱常数，\(C \approx 1.5\)。这等价于 \(S_2(r) \propto r^{2/3}\) 通过Fourier变换。

\textbf{证明}：能谱 \(E(k)\) 与二阶结构函数的关系为 \(S_2(r) = 2 \int_0^\infty E(k) (1 - \frac{\sin(kr)}{kr}) dk\)。若 \(E(k) = C \bar{\varepsilon}^{2/3} k^{-5/3}\)，代入并作变量代换 \(\xi = kr\)，

\[S_2(r) = 2C \bar{\varepsilon}^{2/3} \int_0^\infty k^{-5/3} \left(1 - \frac{\sin(kr)}{kr}\right) dk = 2C \bar{\varepsilon}^{2/3} r^{2/3} \int_0^\infty \xi^{-5/3} \left(1 - \frac{\sin \xi}{\xi}\right) d\xi.\]

积分收敛（在 \(0\) 处 \(\xi^{-5/3} \propto \xi^{-5/3}\)，被 \(1 - \sin\xi/\xi \sim \xi^2/6\) 正则化；在 \(\infty\) 处振荡积分收敛），得 \(S_2(r) \propto r^{2/3}\)。反之亦然，由逆Fourier变换关系可知 -5/3 与 2/3 等价。 \(\blacksquare\)

\subsection{x+2.2 Kolmogorov 4/5律的完整证明}\label{x2.2-kolmogorov-45ux5f8bux7684ux5b8cux6574ux8bc1ux660e}

\textbf{定理 x+2.4}（Kolmogorov 4/5律，1941）：在均匀各向同性湍流的惯性子区内，对充分大的Reynolds数，三阶纵向结构函数满足精确关系

\[S_3(r) = \langle (\delta u_{\parallel}(r))^3 \rangle = -\frac{4}{5}\bar{\varepsilon} r.\]

此关系是Navier-Stokes方程的精确统计推论，不依赖于任何可调参数。

\textbf{证明}：从不可压缩Navier-Stokes方程出发（统计定常态，\(\langle \partial_t \rangle = 0\)）。对两点 \(\mathbf{x}\) 和 \(\mathbf{x}' = \mathbf{x} + \mathbf{r}\)，记 \(\mathbf{u} = \mathbf{u}(\mathbf{x})\)，\(\mathbf{u}' = \mathbf{u}(\mathbf{x}')\)。写出速度增量 \(\delta \mathbf{u} = \mathbf{u}' - \mathbf{u}\) 的演化方程。

由NS方程，\(\partial_t \mathbf{u} + (\mathbf{u} \cdot \nabla)\mathbf{u} = -\nabla p + \nu \Delta \mathbf{u}\)。对 \(\mathbf{x}'\) 处类似。相减得

\[\partial_t (\delta \mathbf{u}) + (\mathbf{u}' \cdot \nabla')\mathbf{u}' - (\mathbf{u} \cdot \nabla)\mathbf{u} = -\nabla' p' + \nabla p + \nu (\Delta' \mathbf{u}' - \Delta \mathbf{u}).\]

记 \(\delta \mathbf{u}\) 与自身的点积并取系综平均。在均匀湍流假设下，\(\langle \mathbf{u} \cdot \partial_t \mathbf{u} \rangle\) 在定常态为零。

关键步骤------Karman-Howarth方程的推导。定义二阶相关张量 \(R_{ij}(\mathbf{r}) = \langle u_i(\mathbf{x}) u_j(\mathbf{x} + \mathbf{r}) \rangle\) 和三阶相关张量 \(S_{ij,k}(\mathbf{r}) = \langle u_i(\mathbf{x}) u_j(\mathbf{x}) u_k(\mathbf{x} + \mathbf{r}) \rangle\)。由均匀性和不可压缩条件，Karman-Howarth方程（对各向同性湍流）为

\[\frac{\partial}{\partial t}(r^2 R_{ll}) = \frac{1}{r^2}\frac{\partial}{\partial r}\left(\frac{1}{r}\frac{\partial}{\partial r}(r^4 S_{ll,l})\right) + 2\nu \frac{1}{r^2}\frac{\partial}{\partial r}\left(r^4 \frac{\partial}{\partial r}\left(\frac{R_{ll}}{r^2}\right)\right).\]

其中 \(R_{ll}\) 是纵向相关函数，\(S_{ll,l}\) 为纵向三阶相关函数。

在惯性子区（\(\nu \to 0\)）和统计定常态（\(\partial_t \to 0\)）下，Karman-Howarth方程简化为

\[\frac{1}{r^2}\frac{\partial}{\partial r}\left(\frac{1}{r}\frac{\partial}{\partial r}(r^4 S_{ll,l})\right) = 0.\]

积分两次：

\[r^4 S_{ll,l} = A r^5 + B r^3 + C.\]

由 \(\mathbf{r} \to 0\) 时的正则性条件，\(S_{ll,l} = O(r^3)\)（速度场的Taylor展开），得 \(A = C = 0\)，\(B\) 为常数。故 \(S_{ll,l}(r) = B r\)。

现需将 \(S_{ll,l}\) 与 \(\langle (\delta u_{\parallel})^3 \rangle\) 联系起来。在三阶纵向结构函数的展开中：

\[\langle (\delta u_{\parallel})^3 \rangle = \langle ((\mathbf{u}' - \mathbf{u}) \cdot \hat{\mathbf{r}})^3 \rangle = 6 S_{ll,l}(r) + \text{（由各项同性约束中的其他项）}.\]

对各向同性湍流，所有三阶张量的形式受 \(\delta_{ij}\) 和 \(r_i\) 的组合限制，最终可得

\[\langle (\delta u_{\parallel})^3 \rangle = 6 S_{ll,l} = 6B r.\]

另一方面，能量耗散率 \(\bar{\varepsilon}\) 通过以下关系引入：Karman-Howarth方程在全尺度（含粘性耗散项）下的精确积分给出

\[\langle (\delta u_{\parallel})^3 \rangle = -\frac{4}{5}\bar{\varepsilon} r + 6\nu \frac{d}{dr}\langle (\delta u_{\parallel})^2 \rangle.\]

在惯性子区，粘性项可忽略（对小 \(\nu\)），得到 \(S_3(r) = -\frac{4}{5}\bar{\varepsilon} r\)。负号表示能量从大尺度向小尺度传递的方向。 \(\blacksquare\)

\subsection{x+2.3 间歇性与Onsager猜想}\label{x2.3-ux95f4ux6b47ux6027ux4e0eonsagerux731cux60f3}

\textbf{定理 x+2.5}（反常标度与间歇性）：实际湍流中高阶结构函数偏离K41标度 \(S_p(r) \propto r^{\zeta_p}\) 且 \(\zeta_p \neq p/3\)。这种偏离称为间歇性：\(\zeta_p\) 是 \(p\) 的凹函数，\(\zeta_p/p\) 随 \(p\) 递减。

\textbf{证明（唯象论证）}：K41假设 \(\delta u_{\parallel} \sim \bar{\varepsilon}_r^{1/3} r^{1/3}\) 且 \(\bar{\varepsilon}_r\) 无间歇（空间上是准均匀的）。但若能量耗散率 \(\varepsilon(\mathbf{x})\) 是间歇的（具有多重分形特征），则 \(\langle \varepsilon_r^p \rangle \propto r^{\tau_p}\)，其中 \(\tau_p\) 编码了耗散场的奇异性谱。此时 \(S_p(r) = \langle (\delta u_{\parallel})^p \rangle \sim r^{p/3} \langle \varepsilon_r^{p/3} \rangle \sim r^{p/3 + \tau_{p/3}}\)。由耗散率的非平凡标度指数 \(\tau_p\) 导致 \(\zeta_p = p/3 + \tau_{p/3} \neq p/3\)。凹性由Hölder不等式的统计版本保证。 \(\blacksquare\)

\textbf{定理 x+2.6}（Onsager猜想，1949 ------ 已解决）：设 \(\mathbf{u}\) 为三维Euler方程的弱解。若 \(\mathbf{u} \in C^\alpha\)（Holder连续），则：

\begin{enumerate}
\def\labelenumi{(\roman{enumi})}
\item
  若 \(\alpha > 1/3\)，能量守恒（\(\frac{d}{dt}\int |\mathbf{u}|^2 = 0\)）；
\item
  存在 \(\alpha < 1/3\) 的弱解使能量耗散（\(\frac{d}{dt}\int |\mathbf{u}|^2 < 0\)）。
\end{enumerate}

\textbf{证明概要}（关键思想）：(i) 由Constantin-E-Titi（1994）和Eyink（1994）证明。能量通量 \(\Pi(r)\) 正比于 \(\langle (\delta u)^3 \rangle / r\)。若 \(\mathbf{u} \in C^\alpha\)，则 \(|\delta \mathbf{u}| \leq C r^\alpha\)，故 \(|\Pi(r)| \leq C r^{3\alpha - 1}\)。当 \(\alpha > 1/3\) 时，\(\Pi(r) \to 0\)（\(r \to 0\)），能量守恒。

\begin{enumerate}
\def\labelenumi{(\roman{enumi})}
\setcounter{enumi}{1}
\tightlist
\item
  由Buckmaster-De Lellis-Szekelyhidi Jr.~(2014-2018)和Isett (2018)通过凸积分（convex integration）方法构造，得到 \(\alpha < 1/3\) 的弱解，且能量确实耗散。构造利用了Tartar的补偿紧性框架和Gromov的h-原理思想，通过在越来越小的尺度上迭加Beltrami波（满足 \(\nabla \times \mathbf{v} = \lambda \mathbf{v}\)）来逐次修正速度场，使得Reynolds应力误差在迭代中趋于零。最终构造在 \(\alpha < 1/3\) 的Holder空间中收敛。临界情形 \(\alpha = 1/3\) 的能量守恒仍为开问题。 \(\blacksquare\)
\end{enumerate}

\subsection{x+2.4 湍流的壳模型与级串理论}\label{x2.4-ux6e4dux6d41ux7684ux58f3ux6a21ux578bux4e0eux7ea7ux4e32ux7406ux8bba}

\textbf{定理 x+2.7}（Richardson能量级串------唯象理论）：湍流的本质是多尺度能量传递。在惯性子区内，能量通过涡的逐级破碎从大尺度（含能尺度 \(L\)）传递到小尺度（耗散尺度 \(\eta\)），形成能量级串。能量通量 \(\varepsilon\) 在各级间守恒，导致

\[E(k) \propto \varepsilon^{2/3} k^{-5/3}.\]

\textbf{证明（量纲分析）}：设 \(v_l \sim (\varepsilon l)^{1/3}\) 为尺度 \(l\) 的特征速度（仅依赖 \(\varepsilon\) 和 \(l\)，因惯性子区内粘性不起作用）。涡翻转时间 \(\tau_l \sim l / v_l \sim l^{2/3} \varepsilon^{-1/3}\)。各尺度的动能密度 \(\sim v_l^2 / l \sim \varepsilon^{2/3} l^{-1/3}\)，换算为波数空间 \(k \sim 1/l\)，得 \(E(k) \sim \varepsilon^{2/3} k^{-5/3}\)。能量通量 \(\varepsilon \sim v_l^3 / l\) 与 \(l\) 无关，验证了级串中能量守恒。 \(\blacksquare\)

\textbf{定理 x+2.8}（GOY壳模型的动力学）：GOY（Gledzer-Ohkitani-Yamada）壳模型将Fourier空间离散化为几何级数的壳 \(k_n = k_0 \lambda^n\)（\(\lambda > 1\)，常取 \(\lambda = 2\)）。每个壳上复速度 \(u_n\) 满足

\[\left(\frac{d}{dt} + \nu k_n^2\right) u_n = i\left(a_n k_{n+1} u_{n+1}^* u_{n+2}^* + b_n k_n u_{n-1}^* u_{n+1}^* + c_n k_{n-1} u_{n-1}^* u_{n-2}^*\right) + f \delta_{n,1},\]

其中系数 \(a_n, b_n, c_n\) 选取以保证无粘无外力时的能量守恒和螺旋度守恒。

\textbf{证明（能量守恒的代数条件）}：壳模型的非线性项保留Navier-Stokes方程非线性项的两个关键性质：(i) 能量守恒------非线性项对内积 \(\sum \operatorname{Re}(u_n^* \dot{u}_n)\) 的贡献为零；(ii) Liouville定理------相空间体积不变。要求 \(\sum_n \operatorname{Re}(u_n^* N_n) = 0\)（\(N_n\) 为非线性项）给出条件 \(a_{n-1} + b_n \lambda + c_{n+1} \lambda^2 = 0\)。满足此条件的系数选取给出Kolmogorov标度 \(|u_n| \sim k_n^{-1/3}\) 作为惯性子区的固定点解。GOY模型通过极大简化保留了湍流级串的核心非线性机制。 \(\blacksquare\)

\textbf{定理 x+2.9}（Kraichnan的Lagrangian重整化群与-5/3谱的推导）：Kraichnan（1959）通过直接相互作用近似（DIA）从NS方程导出了Kolmogorov谱。在DIA框架下，湍流能量谱满足积分方程

\[E(k) = C_K \varepsilon^{2/3} k^{-5/3} \cdot \psi(k/k_d),\]

其中 \(k_d = (\varepsilon/\nu^3)^{1/4}\) 为Kolmogorov耗散波数，\(\psi(\cdot)\) 为刻画耗散区截断的普适函数。

\textbf{证明（概要）}：DIA闭合假设将三阶矩用二阶矩的双时间Green函数表示。对Navier-Stokes方程进行Fourier变换后，二阶矩的演化方程涉及三阶矩。DIA假设三阶矩正比于二阶矩与Green函数的卷积，比例常数由统计各向同性约束确定。在惯性-对流区域（\(\nu \to 0\) 极限）中量纲分析给出 \(E(k) \propto k^{-5/3}\)。Kraichnan的分析还预测了对数修正的存在性（1965），这一预测在后续的实验和数值模拟中得到了验证。 \(\blacksquare\)

\newpage

\chapter{13-数学物理/数学物理/11-电磁学与经典场论.md}\label{ux6570ux5b66ux7269ux7406ux6570ux5b66ux7269ux740611-ux7535ux78c1ux5b66ux4e0eux7ecfux5178ux573aux8bba.md}

\chapter{电磁学与经典场论}\label{ux7535ux78c1ux5b66ux4e0eux7ecfux5178ux573aux8bba}

\section{第475章：Maxwell方程与规范场}\label{ux7b2c475ux7ae0maxwellux65b9ux7a0bux4e0eux89c4ux8303ux573a}

Maxwell方程是经典电磁学的数学基础，以简洁优美的偏微分方程组统一描述了电场与磁场的产生、传播和相互作用。本章从真空中的Maxwell方程出发，引入电磁势和规范变换的概念，揭示规范对称性作为电磁相互作用基本对称性的深层含义。我们系统讨论Coulomb规范和Lorenz规范这两种基本规范固定条件及其对波动方程和静电学问题的简化作用。Poynting定理建立了电磁场能量守恒的精确数学表述，是将场论中的Noether定理应用于电磁理论的物理体现，为理解电磁能量流动和辐射提供了坚实的理论基础。

\subsection{x.1 真空Maxwell方程}\label{x.1-ux771fux7a7amaxwellux65b9ux7a0b}

\textbf{定义 x.1}（真空Maxwell方程组）：真空中的Maxwell方程组为

\[\nabla \cdot \mathbf{E} = \frac{\rho}{\varepsilon_0}, \quad \nabla \cdot \mathbf{B} = 0,\] \[\nabla \times \mathbf{E} = -\frac{\partial \mathbf{B}}{\partial t}, \quad \nabla \times \mathbf{B} = \mu_0 \mathbf{J} + \mu_0 \varepsilon_0 \frac{\partial \mathbf{E}}{\partial t}.\]

其中 \(\mathbf{E}, \mathbf{B}\) 分别为电场和磁场，\(\rho, \mathbf{J}\) 为电荷密度和电流密度。

\textbf{定理 x.2}（电荷守恒定律------连续性方程）：由Maxwell方程组可导出

\[\frac{\partial \rho}{\partial t} + \nabla \cdot \mathbf{J} = 0.\]

\textbf{证明}：对Ampere定律取散度：

\[\nabla \cdot (\nabla \times \mathbf{B}) = \mu_0 \nabla \cdot \mathbf{J} + \mu_0 \varepsilon_0 \frac{\partial}{\partial t}(\nabla \cdot \mathbf{E}).\]

左边恒为零（旋度的散度恒为零）。代入Gauss定律 \(\nabla \cdot \mathbf{E} = \rho/\varepsilon_0\)，得

\[0 = \mu_0 \nabla \cdot \mathbf{J} + \mu_0 \varepsilon_0 \frac{\partial}{\partial t}\left(\frac{\rho}{\varepsilon_0}\right) = \mu_0\left(\nabla \cdot \mathbf{J} + \frac{\partial \rho}{\partial t}\right).\]

因 \(\mu_0 > 0\)，得连续性方程。 \(\blacksquare\)

\textbf{定理 x.3}（Maxwell方程组的Lorentz协变性）：真空Maxwell方程组在Lorentz变换下保持形式不变。用四维形式，定义电磁场张量 \(F^{\mu\nu}\) 和四维电流 \(J^{\mu} = (\rho c, \mathbf{J})\)，Maxwell方程组可写为

\[\partial_{\mu} F^{\mu\nu} = \mu_0 J^{\nu}, \quad \partial_{[\mu} F_{\nu\lambda]} = 0.\]

\textbf{证明}：定义 \(F^{\mu\nu}\) 为反对称张量：

\[F^{\mu\nu} = \begin{pmatrix} 0 & -E_x/c & -E_y/c & -E_z/c \\ E_x/c & 0 & -B_z & B_y \\ E_y/c & B_z & 0 & -B_x \\ E_z/c & -B_y & B_x & 0 \end{pmatrix}.\]

第一个方程 \(\partial_{\mu} F^{\mu\nu} = \mu_0 J^{\nu}\)：取 \(\nu = 0\)，\(\partial_i F^{i0} = \partial_i (E_i/c) = \mu_0 \rho c\)，即 \(\nabla \cdot \mathbf{E} / c = \mu_0 \rho c\)，利用 \(c^2 = 1/(\mu_0 \varepsilon_0)\) 得 \(\nabla \cdot \mathbf{E} = \rho/\varepsilon_0\)。取 \(\nu = i\) 给出Ampere定律。第二个方程 \(\partial_{[\mu} F_{\nu\lambda]} = 0\) 包含Gauss磁定律和Faraday定律。因方程以张量形式书写，其在Lorentz变换 \(x^{\mu} \to \Lambda^{\mu}_{\ \nu} x^{\nu}\) 下协变。 \(\blacksquare\)

\subsection{x.2 电磁势与规范变换}\label{x.2-ux7535ux78c1ux52bfux4e0eux89c4ux8303ux53d8ux6362}

\textbf{定义 x.4}（电磁势）：由 \(\nabla \cdot \mathbf{B} = 0\) 和Poincare引理，存在矢势 \(\mathbf{A}\) 使 \(\mathbf{B} = \nabla \times \mathbf{A}\)。代入Faraday定律 \(\nabla \times (\mathbf{E} + \partial \mathbf{A}/\partial t) = 0\)，得标势 \(\phi\) 的存在性：\(\mathbf{E} = -\nabla \phi - \partial \mathbf{A}/\partial t\)。四维势为 \(A^{\mu} = (\phi/c, \mathbf{A})\)，则 \(F_{\mu\nu} = \partial_{\mu} A_{\nu} - \partial_{\nu} A_{\mu}\)。

\textbf{定理 x.5}（规范不变性）：电磁场在规范变换下不变：对任意光滑函数 \(\Lambda(\mathbf{x}, t)\)，变换

\[\phi \to \phi' = \phi - \frac{\partial \Lambda}{\partial t}, \quad \mathbf{A} \to \mathbf{A}' = \mathbf{A} + \nabla \Lambda\]

保持 \(\mathbf{E}, \mathbf{B}\) 不变。四维形式：\(A_{\mu} \to A_{\mu}' = A_{\mu} + \partial_{\mu} \Lambda\) 保持 \(F_{\mu\nu}\) 不变。

\textbf{证明}：直接验证。

\[\mathbf{B}' = \nabla \times \mathbf{A}' = \nabla \times (\mathbf{A} + \nabla \Lambda) = \nabla \times \mathbf{A} + \nabla \times (\nabla \Lambda) = \mathbf{B}.\] \[\mathbf{E}' = -\nabla \phi' - \frac{\partial \mathbf{A}'}{\partial t} = -\nabla(\phi - \partial_t \Lambda) - \partial_t(\mathbf{A} + \nabla \Lambda) = (-\nabla \phi - \partial_t \mathbf{A}) + \nabla(\partial_t \Lambda) - \partial_t(\nabla \Lambda) = \mathbf{E}.\]

四维形式：\(\partial_{\mu} A_{\nu}' - \partial_{\nu} A_{\mu}' = \partial_{\mu}(A_{\nu} + \partial_{\nu}\Lambda) - \partial_{\nu}(A_{\mu} + \partial_{\mu}\Lambda) = F_{\mu\nu} + \partial_{\mu}\partial_{\nu}\Lambda - \partial_{\nu}\partial_{\mu}\Lambda = F_{\mu\nu}\)。 \(\blacksquare\)

\subsection{x.3 规范固定与波动方程}\label{x.3-ux89c4ux8303ux56faux5b9aux4e0eux6ce2ux52a8ux65b9ux7a0b}

\textbf{定理 x.6}（规范固定的波动方程）：在Lorenz规范 \(\partial_{\mu} A^{\mu} = 0\)（即 \(\nabla \cdot \mathbf{A} + \frac{1}{c^2}\frac{\partial \phi}{\partial t} = 0\)）下，四维势满足波动方程

\[\Box A^{\mu} = \mu_0 J^{\mu},\]

其中 \(\Box = \partial_{\mu}\partial^{\mu} = \frac{1}{c^2}\frac{\partial^2}{\partial t^2} - \nabla^2\) 为d'Alembert算子。

\textbf{证明}：由Maxwell方程 \(\partial_{\mu} F^{\mu\nu} = \mu_0 J^{\nu}\) 和 \(F^{\mu\nu} = \partial^{\mu} A^{\nu} - \partial^{\nu} A^{\mu}\)，

\[\partial_{\mu}(\partial^{\mu} A^{\nu} - \partial^{\nu} A^{\mu}) = \mu_0 J^{\nu}.\]

即 \(\Box A^{\nu} - \partial^{\nu}(\partial_{\mu} A^{\mu}) = \mu_0 J^{\nu}\)。施加Lorenz条件 \(\partial_{\mu} A^{\mu} = 0\)，得 \(\Box A^{\nu} = \mu_0 J^{\nu}\)。 \(\blacksquare\)

\textbf{定理 x.7}（Coulomb规范的静电极限）：在Coulomb规范 \(\nabla \cdot \mathbf{A} = 0\) 下，标势满足Poisson方程

\[\nabla^2 \phi = -\frac{\rho}{\varepsilon_0},\]

即瞬时Coulomb势。矢势满足波动方程 \(\Box \mathbf{A} = \mu_0 \mathbf{J}_T\)，其中 \(\mathbf{J}_T\) 为电流的横向（无散）分量。

\textbf{证明}：在Coulomb规范下，Gauss定律 \(\nabla \cdot \mathbf{E} = \rho/\varepsilon_0\) 变为 \(\nabla \cdot (-\nabla \phi - \partial_t \mathbf{A}) = \rho/\varepsilon_0\)，即 \(-\nabla^2 \phi - \partial_t (\nabla \cdot \mathbf{A}) = \rho/\varepsilon_0\)。因 \(\nabla \cdot \mathbf{A} = 0\)，得 \(\nabla^2 \phi = -\rho/\varepsilon_0\)。对Ampere定律取横向投影（消去 \(\phi\) 的贡献）得 \(\Box \mathbf{A} = \mu_0 \mathbf{J}_T\)，其中 \(\mathbf{J}_T = \mathbf{J} - \nabla \nabla^{-2}(\nabla \cdot \mathbf{J})\)。 \(\blacksquare\)

\subsection{x.4 Poynting定理}\label{x.4-poyntingux5b9aux7406}

\textbf{定理 x.8}（Poynting定理------电磁能量守恒）：定义电磁能量密度 \(u = \frac{1}{2}(\varepsilon_0 |\mathbf{E}|^2 + \frac{1}{\mu_0}|\mathbf{B}|^2)\) 和Poynting矢量 \(\mathbf{S} = \frac{1}{\mu_0}\mathbf{E} \times \mathbf{B}\)。则

\[\frac{\partial u}{\partial t} + \nabla \cdot \mathbf{S} = -\mathbf{J} \cdot \mathbf{E}.\]

\textbf{证明}：利用矢量恒等式 \(\nabla \cdot (\mathbf{E} \times \mathbf{B}) = \mathbf{B} \cdot (\nabla \times \mathbf{E}) - \mathbf{E} \cdot (\nabla \times \mathbf{B})\)。代入Faraday定律和Ampere定律：

\[\begin{aligned}
\nabla \cdot \mathbf{S} &= \frac{1}{\mu_0}[\mathbf{B} \cdot (\nabla \times \mathbf{E}) - \mathbf{E} \cdot (\nabla \times \mathbf{B})] \\
&= \frac{1}{\mu_0}\left[\mathbf{B} \cdot \left(-\frac{\partial \mathbf{B}}{\partial t}\right) - \mathbf{E} \cdot \left(\mu_0 \mathbf{J} + \mu_0 \varepsilon_0 \frac{\partial \mathbf{E}}{\partial t}\right)\right] \\
&= -\frac{1}{2\mu_0}\frac{\partial |\mathbf{B}|^2}{\partial t} - \mathbf{J} \cdot \mathbf{E} - \frac{\varepsilon_0}{2}\frac{\partial |\mathbf{E}|^2}{\partial t} \\
&= -\frac{\partial}{\partial t}\left(\frac{\varepsilon_0}{2}|\mathbf{E}|^2 + \frac{1}{2\mu_0}|\mathbf{B}|^2\right) - \mathbf{J} \cdot \mathbf{E}.
\end{aligned}\]

移项即得。 \(\blacksquare\)

\textbf{定理 x.9}（Poynting定理的积分形式）：对任意固定体积 \(V\)，边界 \(\partial V\) 的外法向 \(\mathbf{n}\)，

\[\frac{d}{dt}\int_V u \, d^3 x + \int_{\partial V} \mathbf{S} \cdot \mathbf{n} \, dS = -\int_V \mathbf{J} \cdot \mathbf{E} \, d^3 x.\]

此即电磁能量守恒：体积内电磁能的减少率等于Poynting矢量的流出通量与焦耳热损耗之和。

\textbf{证明}：由定理x.8的微分形式在 \(V\) 上积分，并应用散度定理即得。右边 \(-\mathbf{J} \cdot \mathbf{E}\) 表示电磁场对带电粒子所做的功（负号表示能量从电磁场转移到粒子动能或热能）。 \(\blacksquare\)

\hrulefill

\section{第476章：电磁波与散射}\label{ux7b2c476ux7ae0ux7535ux78c1ux6ce2ux4e0eux6563ux5c04}

电磁波是Maxwell方程最著名的预言------变化的电磁场以波的形式在空间传播，速度为光速。本章以Helmholtz方程为数学工具，系统分析时谐电磁波的基本性质，包括平面波基展开和球面波展开两种互补的表述方式。Mie散射理论精确求解了均匀介质球对平面电磁波的散射问题，其级数解的收敛性和截断判据具有重要的计算价值。Sommerfeld辐射条件为无界域中Helmholtz方程的外问题提供了物理上合理的唯一性条件。衍射的Kirchhoff理论则将Huygens原理表述为严密的数学积分公式，构成了标量衍射理论的基石。

\subsection{x+1.1 Helmholtz方程}\label{x1.1-helmholtzux65b9ux7a0b}

\textbf{定理 x+1.1}（Maxwell方程到Helmholtz方程的约化）：在无源区域（\(\rho = 0\), \(\mathbf{J} = 0\)）中，时谐电磁场 \(\mathbf{E}(\mathbf{x}, t) = \operatorname{Re}[\mathbf{E}(\mathbf{x}) e^{-i\omega t}]\)，\(\mathbf{B}(\mathbf{x}, t) = \operatorname{Re}[\mathbf{B}(\mathbf{x}) e^{-i\omega t}]\) 满足矢量Helmholtz方程

\[\nabla^2 \mathbf{E} + k^2 \mathbf{E} = 0, \quad \nabla^2 \mathbf{B} + k^2 \mathbf{B} = 0,\]

其中波数 \(k = \omega/c = 2\pi/\lambda\)，\(c = 1/\sqrt{\mu_0 \varepsilon_0}\) 为真空光速。

\textbf{证明}：对无源区域Faraday定律取旋度：\(\nabla \times (\nabla \times \mathbf{E}) = -\frac{\partial}{\partial t}(\nabla \times \mathbf{B})\)。左边用恒等式 \(\nabla \times (\nabla \times \mathbf{E}) = \nabla(\nabla \cdot \mathbf{E}) - \nabla^2 \mathbf{E} = -\nabla^2 \mathbf{E}\)（因 \(\nabla \cdot \mathbf{E} = 0\)）。右边用Ampere定律（\(\mathbf{J} = 0\)时\(\nabla \times \mathbf{B} = \mu_0 \varepsilon_0 \frac{\partial \mathbf{E}}{\partial t}\)），得 \(\nabla^2 \mathbf{E} = \mu_0 \varepsilon_0 \frac{\partial^2 \mathbf{E}}{\partial t^2}\)。对时谐场 \(\partial_t^2 \to -\omega^2\)，得 \(\nabla^2 \mathbf{E} + \frac{\omega^2}{c^2}\mathbf{E} = 0\)。\(\mathbf{B}\) 的证明完全对称。 \(\blacksquare\)

\subsection{x+1.2 平面波与球面波展开}\label{x1.2-ux5e73ux9762ux6ce2ux4e0eux7403ux9762ux6ce2ux5c55ux5f00}

\textbf{定理 x+1.2}（平面波基展开------Weyl恒等式）：任意单色电磁场可表示为平面波的叠加。在无源区域的通解为

\[\mathbf{E}(\mathbf{x}) = \int_{\mathbb{R}^3} \left[\mathbf{A}_1(\mathbf{k}) e^{i\mathbf{k} \cdot \mathbf{x}} + \mathbf{A}_2(\mathbf{k}) e^{-i\mathbf{k} \cdot \mathbf{x}}\right] \delta(k^2 - |\mathbf{k}|^2) \, d^3 k.\]

\textbf{证明}：Helmholtz方程在Fourier空间为 \((k^2 - |\mathbf{k}|^2)\hat{\mathbf{E}}(\mathbf{k}) = 0\)。因此 \(\hat{\mathbf{E}}\) 支撑在球面 \(|\mathbf{k}| = k\) 上，即 \(\hat{\mathbf{E}}(\mathbf{k}) = \mathbf{A}_1(\mathbf{k})\delta(|\mathbf{k}| - k) + \mathbf{A}_2(\mathbf{k})\delta(|\mathbf{k}| + k)\)（向外和向内传播波）。Fourier逆变换给出上述平面波叠加表示。横波条件 \(\nabla \cdot \mathbf{E} = 0\) 等价于 \(\mathbf{k} \cdot \mathbf{A}_{1,2}(\mathbf{k}) = 0\)。 \(\blacksquare\)

\textbf{定理 x+1.3}（球面波展开------多极展开）：Helmholtz方程的球坐标系中分离变量解为

\[\psi_{l,m}(\mathbf{x}) = j_l(kr) Y_{lm}(\theta, \phi), \quad \psi_{l,m}^{\text{out}}(\mathbf{x}) = h_l^{(1)}(kr) Y_{lm}(\theta, \phi),\]

其中 \(j_l\) 为球Bessel函数（规则解），\(h_l^{(1)}\) 为第一类球Hankel函数（向外传播波），\(Y_{lm}\) 为球谐函数。

\textbf{证明}：在球坐标 \((r, \theta, \phi)\) 中，Helmholtz方程 \((\nabla^2 + k^2)\psi = 0\) 可分离变量：\(\psi(r, \theta, \phi) = R(r) \Theta(\theta) \Phi(\phi)\)。角向部分给出球谐函数 \(Y_{lm}\)（\(l = 0, 1, 2, \ldots\)；\(m = -l, \ldots, l\)），满足 \(L^2 Y_{lm} = l(l+1)Y_{lm}\)。径向方程：

\[\frac{d}{dr}\left(r^2 \frac{dR}{dr}\right) + [k^2 r^2 - l(l+1)] R = 0.\]

作变量代换 \(R(r) = u(r)/\sqrt{r}\) 或直接识别为球Bessel方程，其两个线性独立解为 \(j_l(kr)\)（在原点的规则解）和 \(y_l(kr)\)（球Neumann函数，原点奇异）。向外传播的组合为 \(h_l^{(1)}(kr) = j_l(kr) + i y_l(kr)\)，其渐近行为 \(h_l^{(1)}(kr) \sim (-i)^{l+1} e^{ikr}/(kr)\)（\(kr \to \infty\)）。 \(\blacksquare\)

\subsection{x+1.3 Mie散射理论}\label{x1.3-mieux6563ux5c04ux7406ux8bba}

\textbf{定义 x+1.4}（Mie散射问题）：一半径为 \(a\)、折射率为 \(n\) 的均匀介质球置于平面电磁波的入射场中。求球内外的散射电磁场分布。

\textbf{定理 x+1.5}（Mie散射的级数解------消光截面与散射截面）：Mie散射的精确解为球面波展开的无穷级数。消光截面 \(C_{\text{ext}}\) 和散射截面 \(C_{\text{sca}}\) 为

\[C_{\text{ext}} = \frac{2\pi}{k^2}\sum_{l=1}^{\infty}(2l+1)\operatorname{Re}(a_l + b_l), \quad C_{\text{sca}} = \frac{2\pi}{k^2}\sum_{l=1}^{\infty}(2l+1)(|a_l|^2 + |b_l|^2),\]

其中Mie系数 \(a_l, b_l\) 为尺寸参数 \(x = ka\) 和相对折射率 \(m = n/n_{\text{med}}\) 的函数。

\textbf{证明（构造过程与收敛性）}：将入射平面波按球矢量波函数展开（利用平面波的球面波展开公式）：

\[\mathbf{E}_{\text{inc}} = E_0 e^{ikz} \hat{\mathbf{x}} = E_0 \sum_{l=1}^{\infty} i^l \frac{2l+1}{l(l+1)}[\mathbf{M}_{o1l}^{(1)} - i \mathbf{N}_{e1l}^{(1)}],\]

其中 \(\mathbf{M}_{emn}^{(1)}, \mathbf{N}_{omn}^{(1)}\) 等为由球Bessel函数 \(j_l(kr)\) 构造的矢量球谐函数。球内场用 \(j_l(nkr)\) 展开（规则解），散射场用 \(h_l^{(1)}(kr)\) 展开（向外行波）。在球面 \(r = a\) 施加电磁场切向分量的连续性边界条件：

\[\hat{\mathbf{r}} \times (\mathbf{E}_{\text{inc}} + \mathbf{E}_{\text{sca}}) = \hat{\mathbf{r}} \times \mathbf{E}_{\text{int}}, \quad \hat{\mathbf{r}} \times (\mathbf{H}_{\text{inc}} + \mathbf{H}_{\text{sca}}) = \hat{\mathbf{r}} \times \mathbf{H}_{\text{int}}.\]

由球谐函数的正交性，边界条件解耦为各 \(l\) 模的线性方程，解得

\[a_l = \frac{m \psi_l(mx) \psi_l'(x) - \psi_l(x) \psi_l'(mx)}{m \psi_l(mx) \xi_l'(x) - \xi_l(x) \psi_l'(mx)}, \quad b_l = \frac{\psi_l(mx) \psi_l'(x) - m \psi_l(x) \psi_l'(mx)}{\psi_l(mx) \xi_l'(x) - m \xi_l(x) \psi_l'(mx)},\]

其中 \(\psi_l(z) = z j_l(z)\)，\(\xi_l(z) = z h_l^{(1)}(z)\) 为Riccati-Bessel函数。级数的收敛性由 \(l \gg x\) 时 \(|a_l|, |b_l|\) 的指数衰减保证，实际计算中截断于 \(l_{\max} \approx x + 4 x^{1/3} + 2\)。消光和散射截面由Poynting矢量在远场区域积分的光学定理得出。 \(\blacksquare\)

\subsection{x+1.4 Sommerfeld辐射条件与Kirchhoff衍射理论}\label{x1.4-sommerfeldux8f90ux5c04ux6761ux4ef6ux4e0ekirchhoffux884dux5c04ux7406ux8bba}

\textbf{定理 x+1.6}（Sommerfeld辐射条件）：设 \(\psi(\mathbf{x})\) 为无界域 \(\mathbb{R}^3 \setminus \Omega\) 中Helmholtz方程 \(\nabla^2 \psi + k^2 \psi = 0\) 的解。在无穷远处的物理合理条件为

\[\lim_{r \to \infty} r\left(\frac{\partial \psi}{\partial r} - i k \psi\right) = 0,\]

此即Sommerfeld辐射条件，它排除来自无穷远的入射波，保证外问题的唯一性。

\textbf{证明}：设 \(\psi_1\) 和 \(\psi_2\) 为两个满足相同边界条件的解，令 \(w = \psi_1 - \psi_2\)。\(w\) 满足齐次Helmholtz方程，取半径为 \(R\) 的大球面 \(S_R\) 包围散射体。Green恒等式给出

\[\int_{S_R} \left(\bar{w} \frac{\partial w}{\partial r} - w \frac{\partial \bar{w}}{\partial r}\right) dS = 0.\]

代入辐射条件 \(\partial w/\partial r \approx i k w\) 得 \(\int_{S_R} |w|^2 dS \to 0\)（\(R \to \infty\)）。由Rellich引理，从球面均方衰减推得 \(w \equiv 0\)。 \(\blacksquare\)

\textbf{定理 x+1.7}（Kirchhoff衍射积分公式）：设 \(\psi\) 为单色标量波场的复振幅，满足Helmholtz方程。对任意封闭曲面 \(S\) 包围点 \(\mathbf{x}\)，有

\[\psi(\mathbf{x}) = \frac{1}{4\pi}\oint_S \left[G(\mathbf{x}, \mathbf{x}') \frac{\partial \psi}{\partial n}(\mathbf{x}') - \psi(\mathbf{x}') \frac{\partial G}{\partial n}(\mathbf{x}, \mathbf{x}')\right] dS',\]

其中 \(G(\mathbf{x}, \mathbf{x}') = \frac{e^{ik|\mathbf{x} - \mathbf{x}'|}}{|\mathbf{x} - \mathbf{x}'|}\) 为自由空间Green函数。

\textbf{证明}：由Green第二恒等式，

\[\int_V (\psi \nabla'^2 G - G \nabla'^2 \psi) d^3 x' = \oint_{\partial V} \left(\psi \frac{\partial G}{\partial n'} - G \frac{\partial \psi}{\partial n'}\right) dS'.\]

取体积 \(V\) 为 \(S\) 内部挖去以 \(\mathbf{x}\) 为中心的半径为 \(\varepsilon\) 的小球 \(B_\varepsilon(\mathbf{x})\)。在 \(V\) 内，\(\nabla'^2 \psi + k^2 \psi = 0\)，\(\nabla'^2 G + k^2 G = -4\pi \delta(\mathbf{x}' - \mathbf{x})\)（Green函数的定义）。因此左边 \(= -4\pi \psi(\mathbf{x})\)（\(\varepsilon \to 0\) 时仅 \(\delta\) 函数有贡献）。右边面积分分 \(S\) 和 \(\partial B_\varepsilon\) 两部分。在 \(\partial B_\varepsilon\) 上，\(G \sim 1/\varepsilon\)，\(\partial G/\partial n' \sim -1/\varepsilon^2\)，当 \(\varepsilon \to 0\) 时其贡献正好为 \(4\pi \psi(\mathbf{x})\)。移项即得积分公式。 \(\blacksquare\)

\textbf{定理 x+1.8}（Kirchhoff衍射的边界条件------Fresnel-Kirchhoff衍射公式）：设无限大平面屏 \(z = 0\) 上有一孔径 \(\Sigma\)，屏其余部分为完全不透明。在Kirchhoff边界条件（\(\psi\) 和 \(\partial \psi/\partial n\) 在孔径外均为零）下，观测点 \(\mathbf{x}\) 处的衍射场为

\[\psi(\mathbf{x}) = -\frac{i k}{4\pi}\int_{\Sigma} \frac{e^{ik(r+r')}}{r r'} [\cos(\mathbf{n}, \mathbf{r}) - \cos(\mathbf{n}, \mathbf{r}')] dS',\]

其中 \(r = |\mathbf{x} - \mathbf{x}'|\)，\(r'\) 为源点到孔径点的距离，\(\mathbf{n}\) 为屏法向。

\textbf{证明}：将Kirchhoff积分公式应用于由孔径平面和无穷大半球面构成的闭合曲面。无穷远贡献由辐射条件排除。代入平面波入射场 \(\psi_{\text{inc}} = A e^{ikr'}/r'\) 和Green函数，取远场近似（\(kr \gg 1\)）保留 \(\partial G/\partial n \approx i k G \cos(\mathbf{n}, \mathbf{r})\)，\(\partial \psi/\partial n \approx i k \psi \cos(\mathbf{n}, \mathbf{r}')\)，整理即得。 \(\blacksquare\)

\hrulefill

\section{第477章：纤维丛与规范场论}\label{ux7b2c477ux7ae0ux7ea4ux7ef4ux4e1bux4e0eux89c4ux8303ux573aux8bba}

规范场论是20世纪理论物理最辉煌的成就之一，其数学框架与微分几何中的纤维丛理论深刻关联。Yang-Mills理论将电磁学的U(1)规范群推广到非Abel群SU(N)，其场强由曲率二形式 \(F = dA + A \wedge A\) 给出，自相互作用项 \(A \wedge A\) 的存在使得Yang-Mills方程呈现高度非线性。本章建立Yang-Mills理论的基本数学结构，讨论Gribov歧义这一非微扰规范固定中不可避免的拓扑障碍，并介绍Instanton（瞬子）作为Euclidean时空中Yang-Mills方程的经典解，以及指标定理在揭示瞬子数（拓扑荷）与Dirac算子零模个数关系中的关键作用。

\subsection{x+2.1 Yang-Mills理论基本构造}\label{x2.1-yang-millsux7406ux8bbaux57faux672cux6784ux9020}

\textbf{定义 x+2.1}（主丛与联络）：设 \(P\) 为底流形 \(M\)（时空）上以紧Lie群 \(G\) 为结构群的主丛。联络（规范势）为 \(\mathfrak{g}\)-值一形式 \(A \in \Omega^1(P, \mathfrak{g})\)，满足等变条件。局部截面 \(s: U \to P\) 下拉回联络得局部规范势 \(A = s^* A \in \Omega^1(U, \mathfrak{g})\)。

\textbf{定义 x+2.2}（曲率与Yang-Mills泛函）：曲率（规范场强）为 \(\mathfrak{g}\)-值二形式

\[F = dA + A \wedge A \in \Omega^2(M, \mathfrak{g}),\]

其中 \(A \wedge A\) 涉及Lie括号：\((A \wedge A)(X, Y) = [A(X), A(Y)]\)。Bianchi恒等式为 \(d_A F = dF + [A, F] = 0\)。Yang-Mills泛函定义为

\[S_{\text{YM}}[A] = \frac{1}{2}\int_M \operatorname{tr}(F \wedge \star F) = \frac{1}{4}\int_M F_{\mu\nu}^a F^{a \mu\nu} d^4 x.\]

\textbf{定理 x+2.3}（Yang-Mills方程）：Yang-Mills泛函的Euler-Lagrange方程为

\[d_A^* F = d_A \star F = 0,\]

即 \(\nabla_{\mu} F^{\mu\nu} + [A_{\mu}, F^{\mu\nu}] = 0\)。连同Bianchi恒等式 \(d_A F = 0\)，构成Yang-Mills场方程组。

\textbf{证明}：计算变分。\(A \to A + \delta A\) 时，\(\delta F = d_A(\delta A) = d(\delta A) + [A, \delta A]\)。Yang-Mills作用量的变分为

\[\begin{aligned}
\delta S_{\text{YM}} &= \int_M \operatorname{tr}(\delta F \wedge \star F) = \int_M \operatorname{tr}(d_A(\delta A) \wedge \star F) \\
&= \int_M \operatorname{tr}(\delta A \wedge d_A^* \star F) \quad \text{（分部积分，边界项消失）} \\
&= \int_M \operatorname{tr}(\delta A \wedge \star d_A^* F).
\end{aligned}\]

由 \(\delta A\) 的任意性得运动方程 \(d_A^* F = 0\)。用分量写为 \(\nabla_{\mu}^A F^{\mu\nu} = \partial_{\mu} F^{\mu\nu} + [A_{\mu}, F^{\mu\nu}] = 0\)。 \(\blacksquare\)

\textbf{定理 x+2.4}（自对偶方程与Instanton）：在四维Euclidean时空 \(\mathbb{R}^4\) 上，若场强满足（反）自对偶条件

\[F = \pm \star F,\]

则自动满足Yang-Mills方程（由Bianchi恒等式 \(d_A F = 0\) 得 \(d_A \star F = \pm d_A F = 0\)）。满足 \(F = +\star F\) 的解称为Instanton（瞬子），\(F = -\star F\) 的解称为反Instanton。

\textbf{证明}：在四维中，Hodge星算子 \(\star\) 作用于二形式满足 \(\star^2 = 1\)。若 \(F = \pm \star F\)，则

\[d_A \star F = \pm d_A F.\]

由Bianchi恒等式 \(d_A F = 0\)（对任意联络恒成立），得 \(d_A \star F = 0\)，即Yang-Mills方程。因此自对偶条件是Yang-Mills方程的充分条件。解的能量给定为

\[S_{\text{YM}} = \frac{1}{2}\int_{\mathbb{R}^4} \operatorname{tr}(F \wedge \star F) = \pm \frac{1}{2}\int_{\mathbb{R}^4} \operatorname{tr}(F \wedge F).\]

第二个积分为第二Chern数（拓扑不变量），\(c_2 = \frac{1}{8\pi^2}\int \operatorname{tr}(F \wedge F) \in \mathbb{Z}\)。因此自对偶解的Yang-Mills作用量为 \(S_{\text{YM}} = 8\pi^2 |c_2|\)。 \(\blacksquare\)

\subsection{x+2.2 Gribov歧义}\label{x2.2-gribovux6b67ux4e49}

\textbf{定理 x+2.5}（Gribov歧义------Gribov, 1978）：在非Abel规范理论中，Coulomb规范条件 \(\nabla \cdot \mathbf{A}^a = 0\) 不能唯一确定规范势。即存在规范等价的势 \(\mathbf{A}\) 和 \(\mathbf{A}'\) 同时满足Coulomb规范条件。

\textbf{证明}：设 \(\mathbf{A}\) 满足 \(\nabla \cdot \mathbf{A} = 0\)。规范变换 \(\mathbf{A}^g = g^{-1}\mathbf{A}g + g^{-1}\nabla g\)。要求 \(\nabla \cdot \mathbf{A}^g = 0\)，即

\[\nabla \cdot (g^{-1}\mathbf{A}g + g^{-1}\nabla g) = 0.\]

这是关于 \(g\) 的非线性偏微分方程。在强场区域，可能存在多个解 \(g \neq \mathbb{1}\)，每个解对应Gribov副本（Gribov copy）。具体地，考虑SU(2)规范群。沿某非平凡同伦类的 \(g\)（例如 \(g : S^3 \to SU(2) \cong S^3\) 的绕数为1的映射），方程 \(\nabla \cdot (g^{-1}\nabla g) = 0\) 存在非平凡解。因此规范固定是不完全的：规范轨道与规范固定面的交点不唯一。此即Gribov歧义的本质，它表明在非Abel规范理论中不存在全局的规范固定条件。 \(\blacksquare\)

\textbf{定理 x+2.6}（Gribov区域与基础模空间）：Gribov区域 \(\Omega\) 定义为满足 \(\nabla \cdot \mathbf{A} = 0\) 且Faddeev-Popov算子 \(M = -\nabla \cdot D_A\)（\(D_A\) 为协变导数）为正的规范势集合。\(\Omega\) 内的规范轨道与规范固定面一一对应（Gribov歧义被消除）。\(\Omega\) 的边界 \(\partial \Omega\)（Gribov视界）上 \(M\) 出现零模，对应规范轨道与规范固定面首次相切的位置。

\textbf{证明}：Faddeev-Popov算子是规范固定条件 \(\partial_{\mu} A^{\mu} = 0\) 下Faddeev-Popov行列式的核心。在无穷小规范变换 \(\delta A_{\mu} = D_{\mu} \alpha\) 下，\(\delta(\nabla \cdot \mathbf{A}) = \nabla \cdot (D \alpha) = -M \alpha\)。因此，若 \(M\) 在某规范势 \(\mathbf{A}\) 下有零本征值，则存在零模 \(\alpha_0\) 使得 \(\mathbf{A} + \varepsilon D \alpha_0\) 也满足 \(\nabla \cdot \mathbf{A} = 0\) 至 \(O(\varepsilon^2)\)------规范轨道与规范固定面相切。在 \(\Omega\) 内 \(M > 0\)，故无不平凡的零模，规范固定是良好的。Gribov区域的非平凡拓扑（如多个不连通分量）导致了更深层次的拓扑结构。 \(\blacksquare\)

\subsection{x+2.3 Instanton与指标定理}\label{x2.3-instantonux4e0eux6307ux6807ux5b9aux7406}

\textbf{定理 x+2.7}（BPST Instanton解------Belavin-Polyakov-Schwarz-Tyupkin, 1975）：SU(2) Yang-Mills理论在 \(\mathbb{R}^4\) 上的 \(|c_2| = 1\) Instanton解为

\[A_{\mu}(x) = \frac{x^2}{x^2 + \rho^2} g^{-1} \partial_{\mu} g, \quad g(x) = \frac{x_4 \mathbb{1} + i \mathbf{x} \cdot \boldsymbol{\sigma}}{|x|},\]

其中 \(\rho > 0\) 为Instanton尺度参数，\(\boldsymbol{\sigma}\) 为Pauli矩阵。场强为 \(F_{\mu\nu} = \frac{4\rho^2}{(x^2 + \rho^2)^2} \eta_{\mu\nu}^a \frac{\sigma^a}{2}\)，满足自对偶条件 \(F = \star F\)。

\textbf{证明}：利用't Hooft的构造：在Euclidean时空中，寻求形式为 \(A_{\mu} = \bar{\eta}_{\mu\nu}^a \partial_{\nu} \ln \phi\) 的拟设（其中 \(\bar{\eta}_{\mu\nu}^a\) 为't Hooft符号）。代入自对偶方程 \(\tilde{F}_{\mu\nu} = F_{\mu\nu}\)，化简得 \(\phi\) 满足 \(\Box \phi = 0\)（五维Laplace方程被约化）。取 \(\phi = 1 + \sum_{i=1}^{N} \frac{\rho_i^2}{(x - x_i)^2}\)（'t Hooft多Instanton解）。对 \(N = 1\)，\(x_1 = 0\)，得上述BPST解。直接计算可验证自对偶性和拓扑荷 \(\frac{1}{16\pi^2}\int \operatorname{tr}(F_{\mu\nu} \tilde{F}^{\mu\nu}) = 1\)。 \(\blacksquare\)

\textbf{定理 x+2.8}（Atiyah-Singer指标定理在Instanton物理中的应用）：在Instanton背景场 \(A_{\mu}\) 中，Dirac算子 \(D_A = \gamma^{\mu}(\partial_{\mu} + A_{\mu})\) 的零模个数满足

\[n_+ - n_- = \frac{1}{16\pi^2}\int \operatorname{tr}(F_{\mu\nu} \tilde{F}^{\mu\nu}) d^4 x = c_2,\]

其中 \(n_\pm\) 为正/负手征零模的个数。此即Atiyah-Singer指标定理的具体实例。

\textbf{证明}：考虑四维Euclidean空间中的无质量Dirac算子。手征算子 \(\gamma_5\) 与 \(D_A\) 反交换：\(\{\gamma_5, D_A\} = 0\)。因此 \(D_A\) 的非零本征态成对出现（\(\psi_{\lambda}\) 与 \(\gamma_5 \psi_{\lambda}\) 互为手征翻转），但零模可被手征性分类。指标 \(\operatorname{ind}(D_A) = n_+ - n_-\)（正手征零模个数减负手征零模个数）在连续变形下不变（拓扑不变量）。

Atiyah-Singer指标定理将此解析指标表达为拓扑不变量------被积特征类的积分：

\[\operatorname{ind}(D_A) = \int_{\mathbb{R}^4} \hat{A}(\mathbb{R}^4) \wedge \operatorname{ch}(P),\]

对平坦底流形 \(\mathbb{R}^4\) 上的基本表示，\(\hat{A}\)-亏格为1，\(\operatorname{ch}(P)\) 的4-形式分量为 \(\frac{1}{8\pi^2}\operatorname{tr}(F \wedge F)\)。故 \(\operatorname{ind}(D_A) = c_2\)。此即表明：Instanton数（拓扑荷）等于Dirac算子零模的手征差，这解释了手征反常和轴向U(1)问题的数学根源。 \(\blacksquare\)

\textbf{定理 x+2.9}（Instanton的模空间与维度）：在SU(N)规范群下，拓扑荷 \(k\) 的Instanton模空间（即所有规范不等价的Instanton解构成的参数空间）的维度为

\[\dim \mathcal{M}_k = 4 k N.\]

\textbf{证明}：由Atiyah-Singer指标定理在形变复形上的应用。考虑联络的无穷小形变 \(\delta A\) 在自对偶条件下的线性化方程。形变算子为 \(d_A^+ : \Omega^1 \to \Omega^{2,+}\)（将一形式投为自对偶二形式）。其核（零模）对应真正的物理形变，余核对应需排除的规范自由度。指标定理应用于形变复形

\[0 \to \Omega^0 \xrightarrow{d_A} \Omega^1 \xrightarrow{d_A^+} \Omega^{2,+} \to 0\]

给出 \(\dim \ker d_A^+ - \dim \ker d_A - \dim \operatorname{coker} d_A^+ = -4k N\)（符号由椭圆复形的指标公式确定）。模空间维度为 \(\dim \ker d_A^+ - \dim \ker d_A = 4kN\)（因自对偶条件的椭圆性和无余核假设------对不可约联络成立）。这4kN个参数包括Instanton的位置、尺度大小和规范群方向。 \(\blacksquare\)

\subsection{x+2.4 拓扑量子场论与Chern-Simons理论}\label{x2.4-ux62d3ux6251ux91cfux5b50ux573aux8bbaux4e0echern-simonsux7406ux8bba}

\textbf{定理 x+2.10}（Chern-Simons泛函与三维规范场）：在三维紧致流形 \(M^3\) 上，Chern-Simons泛函定义为

\[\operatorname{CS}[A] = \frac{1}{8\pi^2}\int_{M^3} \operatorname{tr}\left(A \wedge dA + \frac{2}{3}A \wedge A \wedge A\right).\]

此泛函在规范变换下并非不变：对大规范变换 \(g : M^3 \to SU(N)\)，\(\operatorname{CS}[A^g] = \operatorname{CS}[A] + k\)，其中 \(k \in \mathbb{Z}\) 为 \(g\) 的绕数（Brouwer度）。

\textbf{证明}：由Chern-Simons形式在规范变换下的性质：

\[A^g = g^{-1}Ag + g^{-1}dg.\]

代入定义，利用迹的循环性和恒等式 \(d(g^{-1}dg) + g^{-1}dg \wedge g^{-1}dg = 0\)（Maurer-Cartan方程），经过直接计算得

\[\operatorname{CS}[A^g] = \operatorname{CS}[A] + \frac{1}{24\pi^2}\int_{M^3} \operatorname{tr}(g^{-1}dg \wedge g^{-1}dg \wedge g^{-1}dg).\]

积分项为 \(g : M^3 \to SU(N)\) 的拓扑绕数 \(w(g)\)，取整数值（当 \(M^3 = S^3\) 时，\(\pi_3(SU(N)) = \mathbb{Z}\)，且标准生成子的绕数为 1）。因此 \(\operatorname{CS}\) 模 \(\mathbb{Z}\) 是规范不变的。这使 \(\exp(2\pi i k \operatorname{CS}[A])\)（\(k \in \mathbb{Z}\)）成为良定义的规范不变可观测量的指数。 \(\blacksquare\)

\textbf{定理 x+2.11}（Witten的Chern-Simons拓扑量子场论微扰展开）：Chern-Simons规范理论的三维配分函数

\[Z_k(M^3) = \int \mathcal{D}A \, e^{i k \cdot \operatorname{CS}[A]}\]

的微扰展开与 \(M^3\) 的Ray-Singer解析挠率（analytic torsion）和Reidemeister挠率相关联。在经典水平上，运动方程 \(\frac{\delta CS}{\delta A} = 0\) 给出 \(F = 0\)，即平坦联络。

\textbf{证明}：Chern-Simons泛函的变分为

\[\frac{\delta \operatorname{CS}}{\delta A} = \frac{1}{4\pi^2} \star F,\]

因此经典运动方程为 \(F = 0\)。对平坦联络附近的微扰展开，选取背景联络 \(A_0\)（\(F_0 = 0\)）。微扰场 \(a = A - A_0\) 在BRST规范固定下的单圈有效作用量由Ray-Singer挠率给出。Witten（1989）证明了Chern-Simons配分函数与Jones多项式（纽结不变量）的深度联系：若 \(M^3 = S^3\) 且Wilson环路 \(\mathcal{W}_R(K) = \operatorname{tr}_R \operatorname{P} \exp(\oint_K A)\) 沿纽结 \(K\) 的期望值给出 \(K\) 的Jones多项式（对应于 \(SU(2)\) 时）或HOMFLY多项式（对应于 \(SU(N)\) 时）。 \(\blacksquare\)

\textbf{定理 x+2.12}（Dirac磁单极与拓扑量子化）：Dirac（1931）证明了磁单极存在的量子力学相容性要求磁荷 \(g\) 与电荷 \(e\) 满足Dirac量子化条件

\[\frac{eg}{\hbar c} = \frac{n}{2}, \quad n \in \mathbb{Z}.\]

本定理等价于 \(U(1)\) 主丛在 \(S^2\) 上分类的第一Chern数为整数。

\textbf{证明}：磁单极的矢势在球坐标系中可写为 \(A_{\phi} = \frac{g}{r}\frac{1-\cos\theta}{\sin\theta}\)（Dirac弦沿南极方向）或 \(A_{\phi} = -\frac{g}{r}\frac{1+\cos\theta}{\sin\theta}\)（Dirac弦沿北极方向）。两种规范选择的差别是一个规范变换 \(A'_{\mu} - A_{\mu} = \partial_{\mu} \Lambda\)，在赤道 \(\theta = \pi/2\) 上 \(\Lambda = 2g\phi\)。量子力学波函数在规范变换下变换为 \(\psi' = e^{ie\Lambda/\hbar c}\psi\)。波函数的单值性要求 \(e^{i e (2g \cdot 2\pi) / \hbar c} = 1\)，即 \(\frac{4\pi eg}{\hbar c} = 2\pi n\)。

从纤维丛观点，磁单极对应 \(U(1)\) 主丛的不平凡分类。磁荷 \(g\) 通过 \(S^2\) 上第一Chern数给出：

\[\frac{1}{2\pi}\int_{S^2} F = \frac{1}{2\pi}\int_{S^2} \frac{g}{r^2} r^2 \sin\theta \, d\theta d\phi = 2g.\]

Chern-Weil定理要求第一Chern数为整数，故 \(2g = n\)，即 \(g = n/2\)。当电场耦合存在时，最小耦合原理给出电荷 \(e\) 和磁荷 \(g\) 同时量子化的条件 \(eg/(\hbar c) = n/2\)。 \(\blacksquare\)

\newpage

\chapter{13-数学物理/数学物理/12-热力学数学理论.md}\label{ux6570ux5b66ux7269ux7406ux6570ux5b66ux7269ux740612-ux70edux529bux5b66ux6570ux5b66ux7406ux8bba.md}

\chapter{热力学数学理论}\label{ux70edux529bux5b66ux6570ux5b66ux7406ux8bba}

热力学的数学理论致力于将热力学的基本原理纳入严密的数学体系中。Carathéodory在1909年给出了第一个完全公理化的热力学体系，该体系利用Pfaff微分形式的可积性理论，从不可达状态的存在性出发导出绝对温度与熵。此后，Gibbs的热力学基本方程和Legendre变换架起了热力学与变分法的桥梁。热传导与扩散方程构成了非平衡态热力学的核心，其数学分析涉及抛物型偏微分方程的最大模原理、热核估计以及Weyl渐近定理。相变与临界现象则将热力学与统计物理紧密联系，Ising模型、Landau理论和重整化群提供了理解临界行为的数学框架。本章将系统阐述这些主题，以定理-证明的形式呈现其数学本质。

\hrulefill

\section{第478章：热力学公理化}\label{ux7b2c478ux7ae0ux70edux529bux5b66ux516cux7406ux5316}

热力学的公理化处理是19世纪末至20世纪初数学物理学的重要成就。与统计力学从微观模型出发不同，公理化热力学直接从宏观经验事实中提炼公理，以纯粹演绎的方式构建整个理论体系。Carathéodory的贡献在于将热力学第二定律表述为Pfaff形式的不可积性条件，从而导出熵的态函数性质。Gibbs基本方程和Legendre变换则为热力学势的数学结构提供了完整的描述。

\subsection{x.1 Pfaff形式与不可达状态}\label{x.1-pfaffux5f62ux5f0fux4e0eux4e0dux53efux8fbeux72b6ux6001}

\textbf{定义}：光滑流形 \(M\) 上的Pfaff形式是一个光滑的1-形式 \(\omega \in \Gamma(T^*M)\)。\(\omega\) 称为完全可积的，若存在光滑函数 \(f, g: M \to \mathbb{R}\) 使得 \(\omega = g\,df\)（即存在积分因子）。

\textbf{定理 x.1}（Carathéodory不可达性原理）：设热力学系统的状态空间为连通流形 \(M\)，平衡态由坐标 \((x^1, \ldots, x^n)\) 参数化。假定热交换由Pfaff形式 \(\omega_Q = \sum a_i dx^i\) 描述。Carathéodory原理断言：在任何给定状态的每个邻域内，存在不能由 \(\omega_Q = 0\) 的积分曲线（即可逆绝热过程）到达的状态。那么 \(\omega_Q\) 必为完全可积的，即存在函数 \(T, S: M \to \mathbb{R}\) 使得

\[\omega_Q = T\,dS.\]

\textbf{证明}：设 \(\omega_Q\) 在每点非零（否则平凡的）。Carathéodory条件意味着：对任意 \(p \in M\) 和任意邻域 \(U \ni p\)，存在 \(q \in U\) 使得 \(q\) 不能由任何分段光滑的 \(\omega_Q\)-零化曲线与 \(p\) 连接。

考虑分布 \(\mathcal{D} = \ker \omega_Q\)，即每点 \(p\) 处满足 \(\omega_Q(p)(v) = 0\) 的切向量构成余维数为1的超平面场。\(\omega_Q = 0\) 的积分曲线正切于此分布。

若 \(\omega_Q = \sum a_i dx^i\) 且 \(\dim M = n\)，Frobenius定理指出 \(\mathcal{D}\) 是（完全）可积的（即 \(M\) 可被 \(\omega_Q = 0\) 的积分曲线生成的极大连通子流形------等熵曲面------所foliated）当且仅当

\[d\omega_Q \wedge \omega_Q = 0.\]

现在假设Carathéodory条件不成立的反面：若在每点 \(p\) 的每个邻域均有可达状态，则分布 \(\mathcal{D}\) 的可达集稠密于 \(M\)，这迫使 \(\mathcal{D}\) 可积（由Chow-Rashevskii定理或可达性理论），与Carathéodory条件矛盾。因此，Carathéodory条件等价于 \(d\omega_Q \wedge \omega_Q = 0\)。

由Pfaff定理（或Frobenius定理的1-形式版本），条件 \(d\omega_Q \wedge \omega_Q = 0\) 等价于 \(\omega_Q\) 是完全可积的，即局部存在函数 \(S\) 和 \(T \neq 0\) 使得 \(\omega_Q = T\,dS\)。函数 \(T\) 即绝对温度（积分因子），\(S\) 即熵。全局存在性由 \(M\) 的单连通性或适当的边界条件保证。\(\blacksquare\)

\subsection{x.2 Gibbs热力学基本方程}\label{x.2-gibbsux70edux529bux5b66ux57faux672cux65b9ux7a0b}

\textbf{定理 x.2}（Gibbs基本方程）：对于简单热力学系统，热力学基本方程为

\[dU = T\,dS - P\,dV + \sum_i \mu_i\,dN_i,\]

其中 \(U\) 为内能，\(T\) 为绝对温度，\(S\) 为熵，\(P\) 为压强，\(V\) 为体积，\(\mu_i\) 为化学势，\(N_i\) 为粒子数。此方程表明 \(U\) 作为外延量的函数是热力学基本势。

\textbf{证明}：由热力学第一定律，\(dU = \delta Q + \delta W\)，其中 \(\delta Q\) 为热交换，\(\delta W\) 为功交换。对于仅有体积功的系统，\(\delta W = -P\,dV\)。由定理x.1，存在绝对温度 \(T\) 和熵 \(S\) 使得 \(\delta Q = T\,dS\)（对于可逆过程）。允许粒子数变化时，化学功项 \(\sum \mu_i\,dN_i\) 加入。因此对于准静态可逆过程，上述方程成立。

从数学角度看，这定义了 \(U\) 是 \(S, V, N_i\) 的函数，且关系式

\[T = \left(\frac{\partial U}{\partial S}\right)_{V,N}, \quad -P = \left(\frac{\partial U}{\partial V}\right)_{S,N}, \quad \mu_i = \left(\frac{\partial U}{\partial N_i}\right)_{S,V,N_{j \neq i}}\]

给出了各强度量作为外延量偏导数的诠释。这些关系的一致性（即二阶偏导数可交换）给出了Maxwell关系式。\(\blacksquare\)

\subsection{x.3 Legendre变换与热力学势}\label{x.3-legendreux53d8ux6362ux4e0eux70edux529bux5b66ux52bf}

\textbf{定义}：设 \(f: \mathbb{R}^n \to \mathbb{R}\) 为严格凸的光滑函数。其Legendre-Fenchel变换定义为

\[f^*(p) = \sup_{x \in \mathbb{R}^n} \{\langle p, x\rangle - f(x)\}.\]

当 \(f\) 可微且梯度映射 \(x \mapsto \nabla f(x)\) 可逆时，\(f^*(p) = \langle p, x(p)\rangle - f(x(p))\)，其中 \(x(p)\) 满足 \(\nabla f(x(p)) = p\)。

\textbf{定理 x.3}（热力学势的Legendre变换结构）：从基本势 \(U(S,V,N)\) 出发，通过Legendre变换可得其他热力学势：

\begin{itemize}
\tightlist
\item
  自由能（Helmholtz）：\(F(T,V,N) = U - TS\)（对 \(S\) 做Legendre变换）
\item
  焓：\(H(S,P,N) = U + PV\)（对 \(V\) 做Legendre变换）
\item
  自由焓（Gibbs）：\(G(T,P,N) = U - TS + PV = H - TS\)（对 \(S, V\) 做双重Legendre变换）
\end{itemize}

\textbf{证明}：以内能 \(U(S,V)\) 为例。\(U\) 是凹函数（熵的凹性），故 \(U^*(T) = \inf_S\{U(S) - TS\}\) 即 \(-F(T)\)。具体地，由 \(T = \partial U/\partial S\) 解出 \(S(T)\) 代入得

\[F(T,V) = U(S(T,V), V) - T \cdot S(T,V).\]

类似构造给出其他势。Legendre变换的逆变换性质（对合性）保证了热力学势之间的一一对应：\(f^{**} = f\)。这为热力学的数学结构提供了坚实的凸分析基础。各势的微分形式为

\[dF = -S\,dT - P\,dV + \sum \mu_i\,dN_i,\] \[dH = T\,dS + V\,dP + \sum \mu_i\,dN_i,\] \[dG = -S\,dT + V\,dP + \sum \mu_i\,dN_i.\]

Maxwell关系式由这些势的二阶混合偏导数相等直接导出。\(\blacksquare\)

\hrulefill

\section{第479章：热传导与扩散方程}\label{ux7b2c479ux7ae0ux70edux4f20ux5bfcux4e0eux6269ux6563ux65b9ux7a0b}

热传导与扩散方程是抛物型偏微分方程最经典的例子。Fourier于1822年提出的热传导定律将热流与温度梯度线性关联，由此导出热方程。本章从数学上分析热方程的适定性、解的正则性和渐近行为。最大模原理揭示了解的先验界，热核为基本解提供了显式表示，Weyl渐近定理则将热方程的谱与区域的几何量联系起来。

\subsection{x+1.1 Fourier热传导定律与热方程}\label{x1.1-fourierux70edux4f20ux5bfcux5b9aux5f8bux4e0eux70edux65b9ux7a0b}

设 \(\Omega \subset \mathbb{R}^n\) 为空间区域，\(u(x,t)\) 表示温度分布。Fourier定律指出热流密度 \(\mathbf{q} = -k\nabla u\)，其中 \(k > 0\) 为热导率。结合能量守恒 \(\rho c_p\,\partial_t u + \nabla \cdot \mathbf{q} = 0\)（\(\rho\) 为密度，\(c_p\) 为比热），得热方程

\[\frac{\partial u}{\partial t} = \alpha \Delta u, \quad \alpha = \frac{k}{\rho c_p}.\]

为简化计，令 \(\alpha = 1\)，考虑初边值问题：

\[\begin{cases}
u_t = \Delta u, & (x,t) \in \Omega \times (0,T],\\
u(x,0) = u_0(x), & x \in \Omega,\\
u(x,t) = 0, & x \in \partial\Omega, \; t > 0.
\end{cases}\]

\subsection{x+1.2 最大模原理}\label{x1.2-ux6700ux5927ux6a21ux539fux7406}

\textbf{定理 x+1.1}（热方程的最大模原理）：设 \(\Omega \subset \mathbb{R}^n\) 为有界开集，\(u \in C^{2,1}(\Omega \times (0,T]) \cap C(\overline{\Omega} \times [0,T])\) 在 \(\Omega \times (0,T]\) 内满足 \(u_t - \Delta u \leq 0\)（下热函数）。则 \(u\) 在抛物边界 \(\partial_p(\Omega \times (0,T]) = (\overline{\Omega} \times \{0\}) \cup (\partial\Omega \times [0,T])\) 上达到最大值。即

\[\max_{\overline{\Omega} \times [0,T]} u = \max_{\partial_p(\Omega \times (0,T])} u.\]

\textbf{证明}：首先考虑 \(u_t - \Delta u < 0\) 的严格情形。设 \((x_0, t_0) \in \Omega \times (0,T]\) 为 \(u\) 在 \(\overline{\Omega} \times [0,T]\) 上的最大值点。若 \((x_0, t_0) \notin \partial_p\)，则 \((x_0, t_0)\) 为内点，\(\nabla u(x_0, t_0) = 0\)，\(\Delta u(x_0, t_0) \leq 0\)（Hessian半负定），且 \(u_t(x_0, t_0) \geq 0\)（因为 \(t_0\) 处达最大，若 \(t_0 < T\) 则时序导数为零，若 \(t_0 = T\) 则须考虑左导数，但可设 \(t_0 < T\) 不失一般性，必要时可平移）。于是 \(u_t - \Delta u \geq 0\)，与严格不等式矛盾。故最大值必在 \(\partial_p\) 上达到。

对于非严格情形 \(u_t - \Delta u \leq 0\)，引入辅助函数 \(v(x,t) = u(x,t) - \varepsilon t\)（\(\varepsilon > 0\)）。则 \(v_t - \Delta v = u_t - \Delta u - \varepsilon \leq -\varepsilon < 0\)。故 \(v\) 的最大值在 \(\partial_p\) 上达到。令 \(\varepsilon \to 0^+\)，由于 \(\max_{\overline{\Omega} \times [0,T]} v \to \max_{\overline{\Omega} \times [0,T]} u\) 且 \(\max_{\partial_p} v \to \max_{\partial_p} u\)，得结论。\(\blacksquare\)

\textbf{推论}（唯一性与稳定性）：初边值问题的解如存在则唯一，且连续依赖于初值。

\subsection{x+1.3 热核与基本解}\label{x1.3-ux70edux6838ux4e0eux57faux672cux89e3}

\textbf{定理 x+1.2}（全空间热核）：算子 \(\partial_t - \Delta\) 在 \(\mathbb{R}^n \times (0,\infty)\) 上的基本解为

\[\Phi(x,t) = \frac{1}{(4\pi t)^{n/2}} \exp\left(-\frac{|x|^2}{4t}\right), \quad t > 0.\]

对初值 \(u_0 \in C(\mathbb{R}^n) \cap L^\infty(\mathbb{R}^n)\)，Cauchy问题的解为卷积

\[u(x,t) = (\Phi(\cdot, t) * u_0)(x) = \int_{\mathbb{R}^n} \Phi(x-y, t)\,u_0(y)\,dy.\]

\textbf{证明}：对 \(t > 0\)，直接验证 \(\Phi_t = \Delta\Phi\)：计算

\[\Phi_t = \frac{1}{(4\pi t)^{n/2}}\left(-\frac{n}{2t} + \frac{|x|^2}{4t^2}\right) e^{-|x|^2/4t},\] \[\Delta\Phi = \frac{1}{(4\pi t)^{n/2}} \sum_{i=1}^n \partial_i\left(-\frac{x_i}{2t} e^{-|x|^2/4t}\right) = \frac{1}{(4\pi t)^{n/2}} \left(-\frac{n}{2t} + \frac{|x|^2}{4t^2}\right) e^{-|x|^2/4t} = \Phi_t.\]

为证明 \(\Phi\) 是基本解，需验证 \(\lim_{t \to 0^+} \Phi(\cdot, t) = \delta_0\)（在分布意义下）。对任意 \(\varphi \in C_c^\infty(\mathbb{R}^n)\)，

\[\lim_{t \to 0^+} \int_{\mathbb{R}^n} \Phi(x,t)\varphi(x)\,dx = \lim_{t \to 0^+} \int_{\mathbb{R}^n} \frac{1}{(4\pi)^{n/2}} e^{-|y|^2/4}\varphi(\sqrt{t}y)\,dy = \varphi(0),\]

其中做了变量替换 \(y = x/\sqrt{t}\)，并应用了控制收敛定理。由卷积的分布理论，\(u(x,t) = \Phi(\cdot,t) * u_0\) 满足热方程，且 \(\lim_{t \to 0^+} u(x,t) = u_0(x)\)。\(\blacksquare\)

\subsection{x+1.4 Weyl渐近定理}\label{x1.4-weylux6e10ux8fd1ux5b9aux7406}

\textbf{定理 x+1.3}（Weyl渐近律）：设 \(\Omega \subset \mathbb{R}^n\) 为有界开集，\(\{0 < \lambda_1 \leq \lambda_2 \leq \cdots\}\) 为 \(-\Delta\) 在 \(\Omega\) 上Dirichlet边值条件下的特征值。令 \(N(\lambda) = \#\{k \mid \lambda_k \leq \lambda\}\)。则

\[N(\lambda) \sim \frac{\omega_n}{(2\pi)^n} \operatorname{vol}(\Omega)\,\lambda^{n/2}, \quad \lambda \to \infty,\]

其中 \(\omega_n\) 为 \(\mathbb{R}^n\) 中单位球的体积。

\textbf{证明概要}：利用热核的迹。定义热迹 \(\Theta(t) = \sum_{k=1}^\infty e^{-\lambda_k t}\)（\(t > 0\)）。由谱定理，\(\Theta(t) = \int_\Omega K(t,x,x)\,dx\)，其中 \(K(t,x,y)\) 为热方程的Dirichlet热核。当 \(t \to 0^+\) 时，热核的短时渐近行为由局部性质决定：

\[K(t,x,x) \sim \frac{1}{(4\pi t)^{n/2}} + \text{边界修正项}, \quad t \to 0^+.\]

因此 \(\Theta(t) \sim \operatorname{vol}(\Omega)/(4\pi t)^{n/2}\) 当 \(t \to 0^+\)。应用Karamata Tauber定理（或直接的Laplace逆变换），由热迹的短时渐近推出特征值计数函数的渐近分布：

\[N(\lambda) \sim \frac{\operatorname{vol}(\Omega)}{(4\pi)^{n/2}\Gamma(n/2+1)} \lambda^{n/2} = \frac{\omega_n}{(2\pi)^n}\operatorname{vol}(\Omega)\,\lambda^{n/2}.\]

其中使用了 \(\omega_n = \pi^{n/2}/\Gamma(n/2+1)\) 和 \((4\pi)^{n/2} = (2\pi)^n/(2\pi)^{n/2}\) 等恒等式。边界修正项不改变主项。\(\blacksquare\)

\subsection{x+1.5 热方程的正则化效应}\label{x1.5-ux70edux65b9ux7a0bux7684ux6b63ux5219ux5316ux6548ux5e94}

\textbf{定理 x+1.4}（正则化效应）：设 \(u_0 \in L^2(\mathbb{R}^n)\)。则热方程的解 \(u(x,t)\) 对任意 \(t > 0\) 是 \(C^\infty\) 光滑的。更一般地，对任意 \(k \in \mathbb{N}\) 和 \(t > 0\)，存在常数 \(C = C(k, n, t)\) 使得

\[\|D^k u(\cdot, t)\|_{L^2} \leq C\,t^{-k/2}\|u_0\|_{L^2}.\]

\textbf{证明}：由解表示 \(u(x,t) = (\Phi(\cdot,t) * u_0)(x)\)。热核 \(\Phi(\cdot,t)\) 是Schwartz函数（对 \(t > 0\)），其Fourier变换为 \(\widehat{\Phi}(\xi, t) = e^{-4\pi^2|\xi|^2 t}\)。由卷积定理，\(\widehat{u(t)}(\xi) = e^{-4\pi^2|\xi|^2 t}\,\widehat{u_0}(\xi)\)。\(D^\alpha u\) 的Fourier变换为 \((2\pi i \xi)^\alpha e^{-4\pi^2|\xi|^2 t}\,\widehat{u_0}(\xi)\)。由Plancherel定理：

\[\|D^\alpha u(t)\|_{L^2}^2 = \int_{\mathbb{R}^n} (2\pi|\xi|)^{2|\alpha|} e^{-8\pi^2|\xi|^2 t} |\widehat{u_0}(\xi)|^2\,d\xi.\]

令 \(s = 2\pi|\xi|\sqrt{t}\)，则 \(e^{-8\pi^2|\xi|^2 t} = e^{-2s^2}\)。估计 \((2\pi|\xi|)^{2|\alpha|} e^{-8\pi^2|\xi|^2 t} \leq C_{|\alpha|}\,t^{-|\alpha|}\)，其中 \(C_{|\alpha|} = \max_{s \geq 0} s^{2|\alpha|} e^{-2s^2}\)。因此

\[\|D^\alpha u(t)\|_{L^2}^2 \leq C_{|\alpha|}\,t^{-|\alpha|} \int_{\mathbb{R}^n} |\widehat{u_0}(\xi)|^2\,d\xi = C_{|\alpha|}\,t^{-|\alpha|}\|u_0\|_{L^2}^2.\]

由 \(\|D^k u\|_{L^2}^2 = \sum_{|\alpha|=k} \|D^\alpha u\|_{L^2}^2\) 得结论。由于 \(C^\infty\) 光滑性可从任意阶导数的有限性推出，正则化效应得证。\(\blacksquare\)

\hrulefill

\section{第480章：相变与临界现象的数学理论}\label{ux7b2c480ux7ae0ux76f8ux53d8ux4e0eux4e34ux754cux73b0ux8c61ux7684ux6570ux5b66ux7406ux8bba}

相变是热力学系统在外部参数变化时发生定性性质突变的物理现象。数学上，相变对应于自由能的非解析性，通常伴随序参量的出现。Landau理论以对称性破缺的观点提供了唯象描述，而Ising模型等格点模型则为从第一原理出发理解相变提供了精确的数学框架。重整化群方法将不同尺度的涨落联系起来，揭示了临界指数的普适性。

\subsection{x+2.1 序参量与Landau理论}\label{x2.1-ux5e8fux53c2ux91cfux4e0elandauux7406ux8bba}

\textbf{定义}：序参量 \(\phi\) 是描述有序-无序相变的宏观变量。在高温无序相中 \(\phi = 0\)，在低温有序相中 \(\phi \neq 0\)。

\textbf{定理 x+2.1}（Landau自由能展开）：在临界点附近，Landau自由能具有以下展开形式：

\[F(\phi, T) = F_0(T) + \frac{1}{2}a(T)\phi^2 + \frac{1}{4}b\phi^4 + \frac{1}{2}c|\nabla\phi|^2,\]

其中 \(a(T) = a_0(T - T_c)\)，\(b > 0\)，\(c > 0\)。平衡态序参量由 \(\partial F/\partial\phi = 0\) 确定，即 \(a\phi + b\phi^3 = 0\)。

\textbf{证明}：考虑对称性：自由能应关于 \(\phi \to -\phi\) 不变（对于Ising型对称性），故展开仅含偶次项。二次项系数 \(a(T)\) 随温度变化，在 \(T > T_c\) 时 \(a(T) > 0\)，系统稳定在 \(\phi = 0\)；当 \(T < T_c\) 时 \(a(T) < 0\)，\(\phi = 0\) 变为不稳定，极小值移至 \(\phi = \pm\sqrt{-a/b}\)。此即自发对称性破缺的Landau图景。

临界指数在平均场水平下为：\(\phi \sim (T_c - T)^{1/2}\)，故 \(\beta = 1/2\)；磁化率 \(\chi = \partial\phi/\partial h|_{h=0} \sim |T - T_c|^{-1}\)，故 \(\gamma = 1\)。这些指数在真实物理系统中会因涨落效应而修正。\(\blacksquare\)

\subsection{x+2.2 Ising模型}\label{x2.2-isingux6a21ux578b}

\textbf{定义}：一维Ising模型定义在格点 \(\{1, \ldots, L\}\) 上，每个格点 \(i\) 上的自旋变量 \(\sigma_i \in \{\pm 1\}\)。Hamilton量为

\[H(\sigma) = -J\sum_{i=1}^{L-1} \sigma_i\sigma_{i+1} - h\sum_{i=1}^L \sigma_i, \quad J > 0.\]

配分函数 \(Z_L = \sum_{\sigma \in \{\pm 1\}^L} e^{-\beta H(\sigma)}\)，\(\beta = 1/(k_B T)\)。

\textbf{定理 x+2.2}（一维Ising模型的精确解，无外场情形）：当 \(h = 0\) 时，一维Ising模型的配分函数为

\[Z_L = 2(2\cosh(\beta J))^{L-1}.\]

在热力学极限 \(L \to \infty\) 下，自由能密度 \(f(\beta) = -\frac{1}{\beta}\lim_{L \to \infty} \frac{1}{L}\ln Z_L\) 为

\[f(\beta) = -\frac{1}{\beta}\ln(2\cosh(\beta J)).\]

特别地，一维Ising模型在有限温度下不存在相变（\(T_c = 0\)）。

\textbf{证明}：使用转移矩阵方法。定义转移矩阵

\[\mathcal{T} = \begin{pmatrix} e^{\beta J} & e^{-\beta J} \\ e^{-\beta J} & e^{\beta J} \end{pmatrix},\]

其矩阵元为 \(\mathcal{T}_{\sigma,\sigma'} = e^{\beta J\sigma\sigma'}\)。配分函数（对周期边界条件）为 \(Z_L^{\text{periodic}} = \operatorname{Tr}(\mathcal{T}^L)\)。对自由边界条件，\(Z_L = \mathbf{1}^T \mathcal{T}^{L-1} \mathbf{1}\)，其中 \(\mathbf{1} = (1,1)^T\)。

计算 \(\mathcal{T}\) 的本征值：特征多项式 \(\det(\mathcal{T} - \lambda I) = (e^{\beta J} - \lambda)^2 - e^{-2\beta J} = 0\)，得

\[\lambda_1 = 2\cosh(\beta J), \quad \lambda_2 = 2\sinh(\beta J).\]

对角化 \(\mathcal{T} = P\operatorname{diag}(\lambda_1, \lambda_2)P^{-1}\)。当 \(L \to \infty\) 时，\(\mathcal{T}^{L-1} \sim \lambda_1^{L-1}\) 主导，\(Z_L \sim 2\lambda_1^{L-1}\)（因子2来自边界项）。故

\[\lim_{L \to \infty} \frac{1}{L}\ln Z_L = \ln \lambda_1 = \ln(2\cosh(\beta J)).\]

由于自由能对所有 \(\beta < \infty\)（\(T > 0\)）是解析的，不存在有限温度相变。关于外场非零的情形，可类似地加入外场项，计算表明自发磁化在 \(T > 0\) 时为零。\(\blacksquare\)

\textbf{定理 x+2.3}（Onsager二维Ising解的表述）：二维正方格Ising模型在无外场情况下，自由能密度的精确表达式为

\[-\beta f = \ln 2 + \frac{1}{8\pi^2}\int_0^{2\pi}\int_0^{2\pi} \ln[\cosh(2\beta J_1)\cosh(2\beta J_2) - \sinh(2\beta J_1)\cos\theta_1 - \sinh(2\beta J_2)\cos\theta_2]\,d\theta_1 d\theta_2.\]

临界点满足 \(\sinh(2\beta_c J_1)\sinh(2\beta_c J_2) = 1\)（对各向同性情形 \(J_1 = J_2 = J\)，即 \(\sinh(2\beta_c J) = 1\)）。临界指数：\(\alpha = 0\)（对数发散），\(\beta = 1/8\)，\(\gamma = 7/4\)。

\textbf{证明思路}：Onsager的原始证明利用Clifford代数（即旋转算符）将转移矩阵对角化问题转化为自由Fermion问题。以二维正方格为例，转移矩阵 \(\mathcal{T}\) 是尺寸为 \(2^N \times 2^N\) 的矩阵（\(N\) 为每行格点数）。通过引入Jordan-Wigner变换，将自旋算符映射为Fermion产生湮灭算符。在此表示下，\(\mathcal{T}\) 成为自由Fermion系统的指数形式，其本征值问题转化为求解Bogoliubov-de Gennes方程。系统的自由能由Fermion能谱的积分给出，即上述二重积分。临界温度的确定条件对应于能谱中出现零模。\(\blacksquare\)

\subsection{x+2.3 平均场理论}\label{x2.3-ux5e73ux5747ux573aux7406ux8bba}

\textbf{定理 x+2.4}（Curie-Weiss平均场理论）：在平均场近似下，Ising模型的自由能为

\[F_{\text{MF}}(m) = -\frac{1}{2}Jzm^2 + \frac{1}{\beta}\left[\frac{1+m}{2}\ln\frac{1+m}{2} + \frac{1-m}{2}\ln\frac{1-m}{2}\right],\]

其中 \(z\) 为配位数，\(m = \langle\sigma_i\rangle\) 为平均磁化强度。平衡态满足自洽方程 \(m = \tanh(\beta J z m)\)。

\textbf{证明}：平均场近似的核心是假设每个自旋感受由近邻产生的平均场 \(h_{\text{eff}} = Jzm\)。单个自旋在此有效场中的配分函数为 \(Z_1 = e^{\beta Jzm} + e^{-\beta Jzm} = 2\cosh(\beta Jzm)\)。磁化强度 \(m = \tanh(\beta Jzm)\)，此即自洽方程。

\(F_{\text{MF}}\) 的熵项来自 \(\ln(1 \pm m)/2\) 的二项分布熵 \(-S/\beta\)，能量项来自 \(-\frac{1}{2}Jzm^2\)（因子1/2避免重复计数相互作用）。在 \(T_c = Jz/k_B\) 处，自洽方程在 \(m = 0\) 处的斜率 \((\partial/\partial m)\tanh(\beta Jzm)|_{m=0} = \beta Jz\) 等于1，故 \(T < T_c\) 时出现非零解，系统发生相变。Landau展开：对小 \(m\)，\(\tanh\) 展开得 \(F_{\text{MF}} \approx \frac{1}{2}(k_B T - Jz)m^2 + \frac{1}{12}k_B T m^4\)，恰为Landau形式。\(\blacksquare\)

\subsection{x+2.4 重整化群基础}\label{x2.4-ux91cdux6574ux5316ux7fa4ux57faux7840}

\textbf{定义}：重整化群（Renormalization Group, RG）是一种相继粗粒化的变换群。对于方格Ising模型，块自旋重整化将 \(b \times b\) 个自旋块合并为一个块自旋变量，并重新标度长度 \(x \to x/b\)。RG变换 \(\mathcal{R}_b\) 作用在Hamilton量空间上。

\textbf{定理 x+2.5}（重整化群流与临界点）：在RG变换下，Hamilton量空间中的不动点描述了临界行为。设 \(K = \beta J\)，RG变换为 \(K' = R_b(K)\)。临界点 \(K_c\) 满足 \(K_c = R_b(K_c)\)（不动点）。在不动点处线性化得 \(d(K' - K_c)/dK|_{K_c} = b^{y}\)，其中 \(y\) 为RG本征指数，关联长度指数满足 \(\nu = 1/y\)。

\textbf{证明}：考虑关联长度 \(\xi(K)\)。在RG变换下，\(\xi(K') = \xi(K)/b\)。在不动点周围展开：\(K' - K_c \approx \lambda_b(K - K_c)\)，其中 \(\lambda_b = dR_b/dK|_{K_c}\)。由标度关系 \(\lambda_b = b^y\)。关联长度 \(\xi(K) \sim |K - K_c|^{-\nu}\)，代入 \(\xi(K') \sim |K' - K_c|^{-\nu} = (\lambda_b|K - K_c|)^{-\nu} = b^{-y\nu}|K - K_c|^{-\nu}\)。另一方面 \(\xi(K') = \xi(K)/b \sim b^{-1}|K - K_c|^{-\nu}\)。比较得 \(\nu = 1/y\)。此即从RG理论导出临界指数的基本框架。\(\blacksquare\)

\subsection{x+2.5 Lee-Yang圆定理}\label{x2.5-lee-yangux5706ux5b9aux7406}

\textbf{定理 x+2.6}（Lee-Yang圆定理）：考虑Ising型自旋系统在外磁场 \(h = H + i\eta\)（复数）中的配分函数 \(Z_N(h)\)。将配分函数写为 \(Z_N(h) = e^{\beta Nh}\sum_{k=0}^N a_k e^{-2\beta kh}\)，其中 \(a_k \geq 0\)（由自旋组态的熵权给出）。则 \(Z_N(h)\) 在复 \(z = e^{-2\beta h}\) 平面内的所有零点均位于单位圆 \(|z| = 1\) 上。

\textbf{证明}：定义 \(z = e^{-2\beta h}\)（配分函数中的逸度变量）。则

\[Z_N(z) = z^{-N/2}\sum_{\sigma} \prod_{\langle i,j\rangle} e^{\beta J\sigma_i\sigma_j}\, z^{-\frac{1}{2}\sum_i \sigma_i} = z^{-N/2}\sum_{k=0}^N a_k z^k.\]

其中 \(a_k \geq 0\)（各项系数为正）。配分函数的零点即多项式 \(P_N(z) = \sum_{k=0}^N a_k z^k\) 的零点。

Lee-Yang定理的核心断言（对于铁磁Ising模型，\(J > 0\)）：所有零点满足 \(|z| = 1\)。这等价于所有零点在 \(h\) 复平面内的虚轴上（\(h\) 的零点为纯虚数）。

证明利用自旋变量的对称性。定义自旋翻转运算 \(\sigma_i \to -\sigma_i\) 全体（全局对称性），Hamilton量在 \(h \to -h\) 和 \(\sigma \to -\sigma\) 的联合变换下不变。由此可证：若 \(z\) 是零点，则 \(1/\bar{z}\) 也是零点。结合多项式系数非负，利用Grace定理或直接分析多项式的自反演（self-inversive）性质------即多项式满足 \(P_N(z) = z^N P_N(1/\bar{z})\)（上横线表示复共轭）------表明所有零点模长为1。\(\blacksquare\)

\textbf{定理 x+2.7}（热力学极限的相变判据）：当 \(N \to \infty\) 时，若Lee-Yang零点集的极限分布在实轴 \(h = 0\)（即 \(z = 1\)）上有非零密度，则系统在该温度下发生相变。

\textbf{证明}：自由能密度 \(f(h) = -\frac{1}{\beta}\lim_{N \to \infty} \frac{1}{N}\ln Z_N(h)\)。由复分析，\(f\) 在零点极限分布支撑集的补集上解析。若零点在实轴某段上累积，则自由能在该段实轴上不解析------此即相变的数学判据。Lee-Yang定理将相变问题转化为零点分布的极限行为分析。\(\blacksquare\)

\subsection{x+2.6 临界指数的普适性与标度关系}\label{x2.6-ux4e34ux754cux6307ux6570ux7684ux666eux9002ux6027ux4e0eux6807ux5ea6ux5173ux7cfb}

\textbf{定义}：临界指数是一组描述热力学量在临界点附近奇异行为的幂律指数。设 \(t = (T - T_c)/T_c\) 为约化温度，\(h\) 为约化外场。主要临界指数定义为：

\begin{itemize}
\tightlist
\item
  比热：\(C \sim |t|^{-\alpha}\)（\(h=0\)）
\item
  序参量：\(m \sim |t|^{\beta}\)（\(h=0\)，\(t<0\)）
\item
  磁化率：\(\chi \sim |t|^{-\gamma}\)（\(h=0\)）
\item
  临界等温线：\(h \sim |m|^{\delta}\operatorname{sgn}(m)\)（\(t=0\)）
\item
  关联长度：\(\xi \sim |t|^{-\nu}\)（\(h=0\)）
\item
  关联函数：\(G(r) \sim r^{-(d-2+\eta)}\)（\(t=h=0\)）
\end{itemize}

\textbf{定理 x+2.8}（标度关系）：临界指数之间并非独立，而是满足以下标度关系（scaling relations）：

\[\alpha + 2\beta + \gamma = 2, \quad \alpha + \beta(1+\delta) = 2, \quad \gamma = \beta(\delta - 1), \quad \nu d = 2 - \alpha, \quad \gamma = \nu(2-\eta).\]

\textbf{证明思路}：标度关系来源于自由能的广义齐次性假设（Widom标度假说）。假设自由能奇异性部分满足

\[f_s(\lambda^{y_t} t, \lambda^{y_h} h) = \lambda^d f_s(t, h),\]

其中 \(y_t\) 和 \(y_h\) 为RG本征指数。通过对 \(t\) 和 \(h\) 求导并取极限，各项临界指数均由 \(y_t\) 和 \(y_h\) 表达：

\[\alpha = 2 - d/y_t, \quad \beta = (d - y_h)/y_t, \quad \gamma = (2y_h - d)/y_t, \quad \delta = y_h/(d-y_h), \quad \nu = 1/y_t, \quad \eta = d+2-2y_h.\]

消去 \(y_t\) 和 \(y_h\) 即得标度关系。\(\blacksquare\)

\subsection{x+2.7 共形场论与临界现象}\label{x2.7-ux5171ux5f62ux573aux8bbaux4e0eux4e34ux754cux73b0ux8c61}

在二维经典统计力学中，临界系统的连续极限由共形场论（Conformal Field Theory, CFT）描述。Belavin、Polyakov和Zamolodchikov（1984）的革命性工作表明，二维CFT的局部对称性（Virasoro代数）对算子内容和关联函数施加了极强的约束。

\textbf{定理 x+2.9}（Virasoro代数与中心荷）：二维共形场论中，共形变换的生成元满足Virasoro代数

\[[L_m, L_n] = (m-n)L_{m+n} + \frac{c}{12}(m^3 - m)\delta_{m+n,0},\]

其中 \(c\) 为中心荷，标识了CFT的普适类。对于Ising模型的临界点，\(c = 1/2\)。

\textbf{证明}：应力-能量张量 \(T(z)\) 的Laurent展开 \(T(z) = \sum_{n \in \mathbb{Z}} L_n z^{-n-2}\)。\(T(z)T(w)\) 的算子乘积展开（OPE）为

\[T(z)T(w) = \frac{c/2}{(z-w)^4} + \frac{2T(w)}{(z-w)^2} + \frac{\partial T(w)}{z-w} + \text{正则项}.\]

通过围道积分和径向量子化技术，上述OPE等价于Virasoro交换关系。中心荷 \(c\) 由 \(TT\) OPE的首项系数确定。对于极小模型 \(\mathcal{M}(p,q)\)（包括Ising模型的 \(c=1/2\)），\(c = 1 - 6(p-q)^2/(pq)\)，且这些值的离散性解释了临界指数的有理数普适类。\(\blacksquare\)

\subsection{x+2.8 非平衡态热力学中的熵产生原理}\label{x2.8-ux975eux5e73ux8861ux6001ux70edux529bux5b66ux4e2dux7684ux71b5ux4ea7ux751fux539fux7406}

非平衡态热力学将平衡态的概念推广到具有宏观流的系统。Prigogine的最小熵产生原理为近平衡态提供了变分描述。

\textbf{定理 x+2.10}（最小熵产生原理）：对于线性唯象关系 \(J_i = \sum_j L_{ij} X_j\)（Onsager倒易关系 \(L_{ij} = L_{ji}\)），其中 \(J_i\) 为热力学流，\(X_j\) 为热力学力，熵产生率 \(\sigma = \sum_i J_i X_i = \sum_{i,j} L_{ij} X_i X_j\) 在给定外力的稳态中达到最小。

\textbf{证明}：考虑 \(\sigma\) 作为力 \(X_i\) 的函数。由Onsager倒易关系，\(L_{ij}\) 为对称矩阵，且由热力学第二定律，矩阵 \(L\) 为半正定。对于固定外力的约束（某些 \(X_i\) 固定），在所有满足约束的 \(X_i\) 上，\(\sigma\) 是关于自由变量的二次型严格极小化问题。一阶条件 \(\partial\sigma/\partial X_k = 0\)（对自由力）等价于对应的流为零 \(J_k = 0\)------这恰为稳态条件。极小性（而非极大性）来自 \(L\) 的半正定性。\(\blacksquare\)

\subsection{x+2.9 Onsager倒易关系的统计力学基础}\label{x2.9-onsagerux5012ux6613ux5173ux7cfbux7684ux7edfux8ba1ux529bux5b66ux57faux7840}

Onsager倒易关系 \(L_{ij} = L_{ji}\) 是近平衡态热力学的基石。它的数学证明根植于微观可逆性和涨落-耗散定理。

\textbf{定理 x+2.11}（Onsager倒易关系的涨落-耗散证明）：设 \(\alpha_i\) 为偏离平衡的热力学变量的微观涨落。微观可逆性要求时间关联函数满足 \(\langle\alpha_i(0)\alpha_j(t)\rangle = \langle\alpha_j(0)\alpha_i(t)\rangle\)。由线性响应理论，唯象系数可表示为

\[L_{ij} = \frac{1}{k_B}\int_0^\infty \langle\dot{\alpha}_i(0)\dot{\alpha}_j(t)\rangle\,dt.\]

由微观可逆性和时间平移不变性，直接推得 \(L_{ij} = L_{ji}\)。

\textbf{证明}：由涨落-耗散定理的Kubo公式，响应函数 \(\chi_{ij}(t) = \beta\langle\dot{\alpha}_i(0)\alpha_j(t)\rangle\)（\(t \geq 0\)）。唯象系数为 \(L_{ij} = \int_0^\infty \chi_{ij}(t)\,dt = \beta\int_0^\infty \langle\dot{\alpha}_i(0)\alpha_j(t)\rangle\,dt\)。利用分步积分和时间反演对称性 \(\langle\alpha_i(0)\alpha_j(-t)\rangle = \langle\alpha_i(t)\alpha_j(0)\rangle\)，可得 \(L_{ij} = L_{ji}\)。具体地，由Green-Kubo关系：

\[L_{ij} = \frac{1}{k_B}\lim_{\omega \to 0} \operatorname{Re}\int_0^\infty \langle J_i(0)J_j(t)\rangle e^{-i\omega t}dt,\]

其中 \(J_i = \dot{\alpha}_i\) 为流。因为经典关联函数的实部在时间反演下为偶函数，极限的对称性即给出 \(L_{ij}=L_{ji}\)。\(\blacksquare\)

\subsection{x+2.10 热力学第三定律的数学表述}\label{x2.10-ux70edux529bux5b66ux7b2cux4e09ux5b9aux5f8bux7684ux6570ux5b66ux8868ux8ff0}

热力学第三定律，又称Nernst热定理，断言绝对零度不可在有限步过程中达到。其数学表述涉及熵在零温极限的行为。

\textbf{定理 x+2.12}（热力学第三定律的Nernst-Planck形式）：对于任何热力学系统，当温度趋于绝对零度时，熵趋于一个普适常数（可取为零）且与系统的其他参数无关：

\[\lim_{T \to 0} S(T, V, N) = 0.\]

等价地，等温过程中熵变在绝对零度时为零：\(\lim_{T \to 0} \Delta S = 0\)。这蕴含了热容量在低温下必须趋于零：

\[\lim_{T \to 0} C_V = \lim_{T \to 0} C_P = 0.\]

\textbf{证明}：由Maxwell关系 \(\left(\frac{\partial S}{\partial P}\right)_T = -\left(\frac{\partial V}{\partial T}\right)_P\)，第三定律要求 \(\lim_{T \to 0} \left(\frac{\partial V}{\partial T}\right)_P = 0\)。由能均分定理的量子统计力学修正，配分函数在低温下由基态主导，自由能的温度导数的奇异性被指数级的Boltzmann因子抑制。设能谱为 \(E_0 < E_1 \leq E_2 \leq \cdots\)，配分函数 \(Z(T) \sim g_0 e^{-E_0/k_BT} + g_1 e^{-E_1/k_BT}\)。熵 \(S = k_B(\ln Z - \beta \partial_\beta \ln Z) \sim k_B\ln g_0\) 当 \(T \to 0\)。若基态非简并（\(g_0 = 1\)），则 \(S(0) = 0\)。热容量 \(C_V = T(\partial S/\partial T)_V \sim O(e^{-\Delta E/k_BT}) \to 0\)，其中 \(\Delta E = E_1 - E_0\)。\(\blacksquare\)

\newpage

\chapter{13-数学物理/数学物理/13-天体力学数学理论.md}\label{ux6570ux5b66ux7269ux7406ux6570ux5b66ux7269ux740613-ux5929ux4f53ux529bux5b66ux6570ux5b66ux7406ux8bba.md}

\chapter{天体力学数学理论}\label{ux5929ux4f53ux529bux5b66ux6570ux5b66ux7406ux8bba}

\begin{quote}
覆盖 MSC 85-XX 天文学与天体物理学。论述N体问题的不可积性、摄动理论与KAM稳定性、轨道力学与相对论天体物理。
\end{quote}

\hrulefill

\section{第481章：N体问题}\label{ux7b2c481ux7ae0nux4f53ux95eeux9898}

天体力学中最基本也最深刻的问题是N体问题：给定\(N\)个质点在Newton万有引力作用下的运动，要求确定其未来任意时刻的位置和速度。自Newton以来，这一问题驱动了分析力学、动力系统理论和摄动理论的发展。2体问题已被Newton完全解决，其轨道为圆锥曲线；但对于\(N \geq 3\)，情况发生了质的变化。Poincaré在其1889年的获奖论文中揭示了3体问题的本质困难------同宿纠缠（homoclinic tangles）的存在排除了除已知经典积分以外的任何解析首次积分的可能性，这被后世总结为Poincaré-Bruns定理。本章系统阐述N体问题的基本结构、不可积性的严格证明、Jacobi坐标下的质心归约、中心构型的动力学意义以及Sundman对3体问题的全局正则化。

\subsection{x.1 N体问题的不可积性}\label{x.1-nux4f53ux95eeux9898ux7684ux4e0dux53efux79efux6027}

\textbf{定义 x.1}（N体问题的Hamilton形式）。设\(N\)个质点的质量分别为\(m_1,\dots,m_N>0\)，位置向量为\(\mathbf{q}_i \in \mathbb{R}^3\)，动量向量为\(\mathbf{p}_i = m_i\dot{\mathbf{q}}_i \in \mathbb{R}^3\)。系统的Hamilton函数为

\[H = \sum_{i=1}^{N} \frac{\|\mathbf{p}_i\|^2}{2m_i} - \sum_{1 \leq i < j \leq N} \frac{G m_i m_j}{\|\mathbf{q}_i - \mathbf{q}_j\|}\]

其中\(G\)为万有引力常数。相空间为\(\mathbb{R}^{6N}\)剔除碰撞构型（\(\mathbf{q}_i = \mathbf{q}_j\)对某\(i \neq j\)）。

\textbf{定理 x.1}（经典积分）。N体问题具有以下10个独立的首积分： 1. 能量积分\(H = E\)（1个） 2. 质心运动积分\(\mathbf{P} = \sum_i \mathbf{p}_i = \text{const}\)（3个） 3. 质心匀速运动积分\(\sum_i m_i \mathbf{q}_i - t\mathbf{P} = \text{const}\)（3个） 4. 角动量积分\(\mathbf{L} = \sum_i \mathbf{q}_i \times \mathbf{p}_i = \text{const}\)（3个）

\textbf{证明}：能量积分直接由\(\partial H / \partial t = 0\)得到。质心运动积分由\(\dot{\mathbf{P}} = \sum_i \dot{\mathbf{p}}_i = \sum_i \mathbf{F}_i = 0\)（作用力与反作用力原理）得到。质心匀速运动积分由\(\frac{d}{dt}(\sum_i m_i \mathbf{q}_i - t\mathbf{P}) = \mathbf{P} - \mathbf{P} = 0\)得到。角动量积分由\(\dot{\mathbf{L}} = \sum_i(\dot{\mathbf{q}}_i \times \mathbf{p}_i + \mathbf{q}_i \times \dot{\mathbf{p}}_i) = \sum_i \mathbf{q}_i \times \mathbf{F}_i = 0\)得到，因为\(\mathbf{F}_i\)平行于\(\sum_{j \neq i}(\mathbf{q}_j - \mathbf{q}_i)\)且求和抵消。 \(\blacksquare\)

\textbf{定理 x.2}（Bruns定理，1887）。3体问题除上述10个经典积分外，不存在其他代数首次积分。

\textbf{证明思路}：假设存在一个与经典积分代数独立的多项式（或代数函数）首次积分\(F(\mathbf{q},\mathbf{p}) = \text{const}\)。Bruns证明了在\(N=3\)时，任何这样的\(F\)必然可以通过经典积分表示，利用旋转对称性和平移对称性下的不变性以及碰撞奇点的解析结构。Bruns采用了精巧的代数消元法：首先利用平移不变性将\(F\)表示为相对坐标的函数；其次利用旋转不变性将其简化为仅依赖于距离\(r_{ij} = \|\mathbf{q}_i - \mathbf{q}_j\|\)和相应的径向动量的函数；最后利用Hamilton方程的代数结构导出一个矛盾。关键的一步是分析碰撞构型附近\(F\)的形式级数展开。 \(\blacksquare\)

\textbf{定理 x.3}（Poincaré定理，1890）。存在Hamilton扰动\(H = H_0 + \varepsilon H_1\)使得扰动系统不具备除能量积分以外的任何解析首次积分。

\textbf{证明思路}：Poincaré考虑圆型限制3体问题------两个大天体绕其质心做圆周运动，第三个小天体的运动可视为两个中心问题（可积的\(H_0\)）加上扰动\(\varepsilon H_1\)。假设存在一个依赖于\(\varepsilon\)的解析首次积分\(\Phi(\mathbf{q},\mathbf{p},\varepsilon) = \sum_{k \geq 0} \Phi_k(\mathbf{q},\mathbf{p}) \varepsilon^k\)。将其代入Poisson括号方程\(\{H,\Phi\} = 0\)，得到递推关系\(\{H_0,\Phi_k\} + \{H_1,\Phi_{k-1}\} = 0\)。而\(\{H_0,\cdot\}\)是Liouville环面上的平移算子，其可解性要求\(\{H_1,\Phi_{k-1}\}\)在环面上的平均值为零。Poincaré通过分析同宿纠缠的Melnikov函数证明了对于一般的扰动，这一条件无法满足到所有阶。 \(\blacksquare\)

\subsection{x.2 Jacobi坐标与质心归约}\label{x.2-jacobiux5750ux6807ux4e0eux8d28ux5fc3ux5f52ux7ea6}

\textbf{定义 x.2}（Jacobi坐标）。对于\(N\)体系统，定义Jacobi坐标\(\{\mathbf{R}_0,\mathbf{R}_1,\dots,\mathbf{R}_{N-1}\}\)如下：

\[\mathbf{R}_0 = \frac{1}{M}\sum_{i=1}^{N} m_i \mathbf{q}_i, \quad \mathbf{R}_k = \mathbf{q}_{k+1} - \frac{1}{M_k}\sum_{i=1}^{k} m_i \mathbf{q}_i \quad (k=1,\dots,N-1)\]

其中\(M = \sum_{i=1}^{N} m_i\)，\(M_k = \sum_{i=1}^{k} m_i\)。

\textbf{定理 x.4}（Jacobi坐标下的分离）。在Jacobi坐标下，动能可对角化：

\[T = \frac{1}{2}M\|\dot{\mathbf{R}}_0\|^2 + \frac{1}{2}\sum_{k=1}^{N-1} \mu_k \|\dot{\mathbf{R}}_k\|^2\]

其中约化质量\(\mu_k = \frac{m_{k+1}M_k}{M_{k+1}}\)。对应动量为\(\mathbf{P}_0 = M\dot{\mathbf{R}}_0\)，\(\mathbf{P}_k = \mu_k \dot{\mathbf{R}}_k\)。

\textbf{证明}：通过直接代入验证。关键在于递推构造保证了质量矩阵的对角化。记\(\mathbf{Q} = (\mathbf{q}_1,\dots,\mathbf{q}_N)^T\)，\(\mathbf{R} = (\mathbf{R}_0,\dots,\mathbf{R}_{N-1})^T\)，存在一个下三角矩阵\(A\)使得\(\mathbf{R} = A\mathbf{Q}\)。动能形如\(T = \frac{1}{2}\dot{\mathbf{Q}}^T \text{diag}(m_1,\dots,m_N) \dot{\mathbf{Q}}\)。经过变换后\(T = \frac{1}{2}\dot{\mathbf{R}}^T (A^{-T}\text{diag}(m_i)A^{-1})\dot{\mathbf{R}}\)，而构造保证了\(A^{-T}\text{diag}(m_i)A^{-1}\)是对角矩阵，对角元依次为\(M,\mu_1,\dots,\mu_{N-1}\)。 \(\blacksquare\)

\textbf{推论}：质心坐标\(\mathbf{R}_0\)的运动是匀速直线运动，可完全分离。因此N体问题本质上可约化为关于\(N-1\)个Jacobi矢量\(\mathbf{R}_1,\dots,\mathbf{R}_{N-1}\)的\(3(N-1)\)自由度系统。

\subsection{x.3 中心构型与总碰撞}\label{x.3-ux4e2dux5fc3ux6784ux578bux4e0eux603bux78b0ux649e}

\textbf{定义 x.3}（中心构型）。N体构型\(\mathbf{q} = (\mathbf{q}_1,\dots,\mathbf{q}_N)\)称为中心构型，如果存在常数\(\lambda\)使得每个质点的加速度正比于其相对于质心的位置：

\[\ddot{\mathbf{q}}_i = \lambda(\mathbf{q}_i - \mathbf{c}), \quad \mathbf{c} = \frac{1}{M}\sum_i m_i \mathbf{q}_i\]

\textbf{定理 x.5}（中心构型的变分刻画）。中心构型是势函数\(U(\mathbf{q}) = \sum_{i<j} \frac{G m_i m_j}{\|\mathbf{q}_i - \mathbf{q}_j\|}\)在惯量矩\(I(\mathbf{q}) = \sum_i m_i\|\mathbf{q}_i - \mathbf{c}\|^2\)固定条件下的临界点。

\textbf{证明}：设Lagrange乘子为\(\mu\)，考虑\(\nabla(U - \mu I) = 0\)。对被约束质心在原点处的构型，\(I = \sum_i m_i \|\mathbf{q}_i\|^2\)。由Newton方程\(\ddot{\mathbf{q}}_i = \frac{1}{m_i} \frac{\partial U}{\partial \mathbf{q}_i}\)。临界条件\(\frac{\partial U}{\partial \mathbf{q}_i} - 2\mu m_i \mathbf{q}_i = 0\)给出\(\ddot{\mathbf{q}}_i = 2\mu \mathbf{q}_i\)。取\(\lambda = 2\mu\)即得中心构型的定义。反之，若构型为中心构型，则\(\frac{\partial U}{\partial \mathbf{q}_i} = m_i \lambda \mathbf{q}_i = 2\mu m_i \mathbf{q}_i\)，由Euler齐次关系\(U\)是\((-1)\)次齐次的，可得\(\mu = -U/(2I) < 0\)，从而构型满足变分条件。 \(\blacksquare\)

\textbf{定理 x.6}（Saari猜想的特例陈述，1970）。在3体问题中，除等边三角形构型外的所有初始条件均不导致总碰撞。

\textbf{证明思路}（McGehee坐标下的分析）：引入McGehee变换\(r = \sqrt{I}\)，\(\mathbf{s} = \mathbf{q}/r\)（形状球面），\(\tau\)满足\(dt/d\tau = r^{3/2}\)。在此坐标下，总碰撞对应于\(r = 0\)的边界流形，可进行紧化分析。总碰撞的渐进行为要求构型趋于中心构型，在3体情形下仅有Lagrange等边三角形和Euler共线两种可能。McGehee-Mather证明了对非等边构型，解的流无法终结于总碰撞，而是在\(r=0\)处弹跳。 \(\blacksquare\)

\subsection{x.4 Sundman级数解}\label{x.4-sundmanux7ea7ux6570ux89e3}

\textbf{定理 x.7}（Sundman，1912）。3体问题在角动量非零的条件下存在一个在整个实轴\(\mathbb{R}\)上收敛的幂级数解（表示为某种正则化时间的函数）。

\textbf{推导}：引入正则化时间变量\(\omega\)，满足

\[\frac{dt}{d\omega} = U^{-1}, \quad U(\mathbf{q}) = \sum_{i<j} \frac{G m_i m_j}{\|\mathbf{q}_i - \mathbf{q}_j\|}\]

再对二体碰撞进行Levi-Civita正则化。具体而言，对每对碰撞的\((\mathbf{q}_i,\mathbf{q}_j)\)引入抛物坐标或Kustaanheimo-Stiefel变换将奇点提升为解析点。在\(\omega\)参数下，运动方程变为无奇异性的解析微分方程：

\[\frac{d^2 \mathbf{q}_i}{d\omega^2} + \frac{1}{U}\frac{dU}{d\omega}\frac{d\mathbf{q}_i}{d\omega} = \frac{1}{U^2}\frac{\partial U}{\partial \mathbf{q}_i}\]

利用Cauchy-Kovalevskaya定理，局部存在关于\(\omega\)的收敛幂级数解。Sundman的关键贡献在于证明了该级数的收敛半径为无穷------令

\[\Omega = \int \frac{dt}{U}\]

则\(\Omega\)随\(t\)单调增长且有界上界，因此可变换回物理时间。利用Painlevé的碰撞非本性奇点定理和复域延拓，得到全局收敛性。 \(\blacksquare\)

\hrulefill

\section{第482章：摄动理论与KAM稳定性}\label{ux7b2c482ux7ae0ux6444ux52a8ux7406ux8bbaux4e0ekamux7a33ux5b9aux6027}

N体问题的不可积性并不意味着我们无法对其长期行为做出有效预测。摄动理论为近似求解提供了系统方法，而KAM（Kolmogorov-Arnold-Moser）理论则从定性角度保证了在充分小的扰动下大多数不变环面得以保持。天体力学的核心问题是太阳系的长期稳定性------行星轨道是否会保持近似Kepler椭圆，还是会演变为混沌或碰撞？本章从平均化方法出发，逐步建立Lindstedt-Poincaré摄动级数理论，进而引入KAM理论的基本框架及其在天体力学中的深刻应用，最后讨论Arnold扩散、Nekhoroshev估计和Laskar对太阳系长期演化的数值研究。

\subsection{x+1.1 平均化方法}\label{x1.1-ux5e73ux5747ux5316ux65b9ux6cd5}

\textbf{定理 x+1.1}（一阶平均化定理）。考虑慢变系统\(\dot{\mathbf{I}} = \varepsilon \mathbf{f}(\mathbf{I},\boldsymbol{\phi})\)，\(\dot{\boldsymbol{\phi}} = \boldsymbol{\omega}(\mathbf{I}) + \varepsilon \mathbf{g}(\mathbf{I},\boldsymbol{\phi})\)，其中\((\mathbf{I},\boldsymbol{\phi}) \in \mathbb{R}^n \times \mathbb{T}^n\)为作用-角变量。定义平均系统\(\dot{\mathbf{J}} = \varepsilon \bar{\mathbf{f}}(\mathbf{J})\)，其中\(\bar{\mathbf{f}}(\mathbf{J}) = (2\pi)^{-n}\int_{\mathbb{T}^n} \mathbf{f}(\mathbf{J},\boldsymbol{\phi}) d\boldsymbol{\phi}\)。若\(\boldsymbol{\omega}(\mathbf{I})\)非共振（\(\mathbf{k} \cdot \boldsymbol{\omega}(\mathbf{I}) \neq 0\)对所有\(\mathbf{k} \in \mathbb{Z}^n \setminus \{0\}\)），则当\(\varepsilon \to 0\)时，原系统的解在时间尺度\(O(1/\varepsilon)\)上以\(O(\varepsilon)\)误差逼近平均系统的解。

\textbf{证明}：构造近恒等变换\(\mathbf{I} = \mathbf{J} + \varepsilon \mathbf{W}(\mathbf{J},\boldsymbol{\psi})\)，\(\boldsymbol{\phi} = \boldsymbol{\psi} + \varepsilon \mathbf{V}(\mathbf{J},\boldsymbol{\psi})\)，其中\(\mathbf{W},\mathbf{V}\)在\(\mathbb{T}^n\)上周期。代入运动方程并保留到\(O(\varepsilon)\)：

\[\dot{\mathbf{J}} + \varepsilon \frac{\partial \mathbf{W}}{\partial \boldsymbol{\psi}} \cdot \dot{\boldsymbol{\psi}} = \varepsilon \mathbf{f}(\mathbf{J},\boldsymbol{\psi}) + O(\varepsilon^2)\] \[\dot{\boldsymbol{\psi}} = \boldsymbol{\omega}(\mathbf{J}) + O(\varepsilon)\]

零阶给出\(\dot{\boldsymbol{\psi}} = \boldsymbol{\omega}(\mathbf{J})\)，一阶给出\(\dot{\mathbf{J}} = \varepsilon[\mathbf{f}(\mathbf{J},\boldsymbol{\psi}) - \boldsymbol{\omega}(\mathbf{J}) \cdot \frac{\partial \mathbf{W}}{\partial \boldsymbol{\psi}}] + O(\varepsilon^2)\)。选择\(\mathbf{W}\)满足同调方程

\[\boldsymbol{\omega}(\mathbf{J}) \cdot \frac{\partial \mathbf{W}}{\partial \boldsymbol{\psi}} = \mathbf{f}(\mathbf{J},\boldsymbol{\psi}) - \bar{\mathbf{f}}(\mathbf{J})\]

由于\(\boldsymbol{\omega}(\mathbf{J})\)非共振，Fourier展开\(\mathbf{f}(\mathbf{J},\boldsymbol{\psi}) = \sum_{\mathbf{k}} \hat{\mathbf{f}}_{\mathbf{k}}(\mathbf{J}) e^{i\mathbf{k}\cdot\boldsymbol{\psi}}\)给出解\(\hat{\mathbf{W}}_{\mathbf{k}} = \frac{\hat{\mathbf{f}}_{\mathbf{k}}}{i\mathbf{k}\cdot\boldsymbol{\omega}}\)（\(\mathbf{k} \neq 0\)）。从而\(\dot{\mathbf{J}} = \varepsilon \bar{\mathbf{f}}(\mathbf{J}) + O(\varepsilon^2)\)。在\(t \in [0,T/\varepsilon]\)上积分得误差\(O(\varepsilon)\)。 \(\blacksquare\)

\subsection{x+1.2 Lindstedt-Poincaré方法}\label{x1.2-lindstedt-poincaruxe9ux65b9ux6cd5}

\textbf{定理 x+1.2}（Lindstedt-Poincaré展开）。考虑弱非线性振动方程

\[\ddot{x} + x = \varepsilon f(x,\dot{x})\]

其中\(f(0,0) = 0\)。设\(f\)在原点邻域内解析，则存在周期解的形式级数展开

\[x(t) = x_0(\tau) + \varepsilon x_1(\tau) + \varepsilon^2 x_2(\tau) + \cdots, \quad \tau = \omega t\]

其中\(\omega = 1 + \varepsilon\omega_1 + \varepsilon^2\omega_2 + \cdots\)为待定频率。

\textbf{推导}：将展开代入原方程：

\[(\omega^2 \frac{d^2}{d\tau^2} + 1)(x_0 + \varepsilon x_1 + \cdots) = \varepsilon f(x_0 + \cdots, \omega \dot{x}_0 + \cdots)\]

展开\(\omega^2 = 1 + 2\varepsilon\omega_1 + \varepsilon^2(\omega_1^2 + 2\omega_2) + \cdots\)。按\(\varepsilon\)的幂次整理：

\(\varepsilon^0\)阶：\(\frac{d^2 x_0}{d\tau^2} + x_0 = 0\)，通解为\(x_0(\tau) = A\cos(\tau + \phi)\)。

\(\varepsilon^1\)阶：\(\frac{d^2 x_1}{d\tau^2} + x_1 = f(x_0,\omega_0 \dot{x}_0) - 2\omega_1 \frac{d^2 x_0}{d\tau^2}\)。将右端展开为Fourier级数\(\sum_k c_k e^{ik\tau}\)。消除长期项的条件要求\(c_{\pm 1} = 0\)，这确定了\(\omega_1\)和\(A\)的关系（频率-振幅关系）。\(x_1\)的解为满足\(c_{\pm 1}=0\)条件下的特解。逐阶类推，每步消除谐振项的条件确定\(\omega_k\)。当\(f\)为解析函数时，若\(A\)充分小，该形式级数在\(\varepsilon\)充分小时收敛（由Poincaré周期解延拓定理保证）。 \(\blacksquare\)

\subsection{x+1.3 KAM理论在天体力学中的应用}\label{x1.3-kamux7406ux8bbaux5728ux5929ux4f53ux529bux5b66ux4e2dux7684ux5e94ux7528}

\textbf{定理 x+1.3}（Kolmogorov定理，1954）。考虑近可积Hamilton系统

\[H(\mathbf{I},\boldsymbol{\phi}) = H_0(\mathbf{I}) + \varepsilon H_1(\mathbf{I},\boldsymbol{\phi})\]

其中\((\mathbf{I},\boldsymbol{\phi}) \in D \times \mathbb{T}^n\)，\(D \subset \mathbb{R}^n\)为开集。设\(H_0\)非退化（\(\det(\partial^2 H_0 / \partial I_i \partial I_j) \neq 0\)），\(H_1\)解析。则对充分小的\(\varepsilon\)，存在一个具有正Lebesgue测度的Cantor集\(\mathcal{K}_\varepsilon \subset D\)，使得对每个\(\mathbf{I}^* \in \mathcal{K}_\varepsilon\)，存在一个解析的Lagrange环面\(\mathcal{T}_{\mathbf{I}^*}\)承载频率为\(\boldsymbol{\omega}^* = \partial H_0 / \partial \mathbf{I}|_{\mathbf{I}^*}\)的拟周期运动。当\(\varepsilon \to 0\)时，\(D \setminus \mathcal{K}_\varepsilon\)的测度趋于零。

\textbf{证明框架}：采用超收敛Newton迭代法。第一步：构造生成函数\(S(\mathbf{J},\boldsymbol{\phi}) = \mathbf{J} \cdot \boldsymbol{\phi} + \varepsilon S_1(\mathbf{J},\boldsymbol{\phi})\)，使得在新变量\((\mathbf{J},\boldsymbol{\psi})\)下Hamilton函数变为\(H'(\mathbf{J}) + \varepsilon^2 H_2(\mathbf{J},\boldsymbol{\psi})\)。求解同调方程：

\[\boldsymbol{\omega}(\mathbf{J}) \cdot \frac{\partial S_1}{\partial \boldsymbol{\psi}} = H_1(\mathbf{J},\boldsymbol{\phi}) - \langle H_1 \rangle\]

需满足Diophantine条件\(|\mathbf{k} \cdot \boldsymbol{\omega}| \geq \gamma |\mathbf{k}|^{-\tau}\)（对所有\(\mathbf{k} \neq 0\)）。利用解析函数Fourier系数的指数衰减性，\(S_1\)有解析解。第二步：扰动尺度从\(\varepsilon\)缩至\(\varepsilon^2\)（超收敛）。迭代此过程，每次需排除更多共振频率，最终不变环面集为Cantor型。Arnold（1963）对\(n=2\)给出了严格证明，Moser（1962）用光滑版本完成了证明。 \(\blacksquare\)

\textbf{定理 x+1.4}（Nekhoroshev估计，1977）。在上述近可积Hamilton系统中，若Hessian矩阵\(\partial^2 H_0 / \partial I_i \partial I_j\)满足某种``速降性''（steepness）条件，则存在正常数\(a,b,\varepsilon_0,R_0,T_0\)使得对所有\(\varepsilon < \varepsilon_0\)和满足\(\|\mathbf{I}(0)\| < R_0\)的初始条件，在时间\(T \leq T_0 \exp(c\varepsilon^{-a})\)内有估计

\[\|\mathbf{I}(t) - \mathbf{I}(0)\| \leq C \varepsilon^b\]

\textbf{证明思路}：构造截断的正则变换以消除直到某一阶\(r\)的扰动（\(r \sim \varepsilon^{-a}\)）。残差为\(O(\varepsilon^r)\)量级。在截断变换后的作用变量在指数长时间内漂移被指数地压制。证明的技术核心在于利用\(H_0\)的几何性质（凸性或速降性）来控制在共振区域之间的小分母问题，以及细致的解析估计来处理Fourier级数的收敛性。Arnol'd扩散（沿共振网的极慢漂移）的速率受此估计的约束。 \(\blacksquare\)

\subsection{x+1.4 太阳系的长期稳定性}\label{x1.4-ux592aux9633ux7cfbux7684ux957fux671fux7a33ux5b9aux6027}

太阳系是否稳定的问题已有三百多年历史。Laplace和Lagrange的行星轨道理论似乎表明轨道的长期变化是拟周期的，但Poincaré的工作揭示了根本性的困难。现代研究由Laskar（1989, 1994）通过数值模拟取得了突破性进展。

\textbf{Laskar的数值证据}：Laskar计算了太阳系全行星运动方程的长期平均化系统，采用频率映射分析方法（frequency map analysis）对轨道进行数百万年的数值积分。关键结论是：（1）内行星（水星、金星、地球、火星）的轨道偏心率在数千万年尺度上表现出混沌扩散行为；（2）水星的偏心率可能在35亿年内增长到足以使其与金星发生近距离交会的程度；（3）整个太阳系的Lyapunov时间约为500万年。从理论角度，Laskar通过高精度的Laplace-Lagrange长期摄动展开到二阶（考虑广义相对论的后Newton修正）建立了包含约15万个项的级数模型，为数值探索提供了半解析工具。

\hrulefill

\section{第483章：轨道力学与相对论天体物理}\label{ux7b2c483ux7ae0ux8f68ux9053ux529bux5b66ux4e0eux76f8ux5bf9ux8bbaux5929ux4f53ux7269ux7406}

轨道力学关注天体在引力场中的运动轨迹及其稳定性。从经典Kepler问题到后Newton天体力学，轨道理论始终是天体力学与天体物理学之间的桥梁。本章涵盖Lagrange行星方程对轨道要素摄动的系统描述、Hill方程在受限3体问题中的Floquet分析、Hénon-Heiles系统的混沌行为这一非线性动力学经典模型，以及广义相对论框架下Schwarzschild时空中的轨道进动、后Newton形式体系和引力波致轨道衰变的Peters-Mathews理论。

\subsection{x+2.1 Lagrange行星方程}\label{x2.1-lagrangeux884cux661fux65b9ux7a0b}

\textbf{定义 x+2.1}（密切轨道要素）。在摄动二体问题中，定义轨道要素\((a,e,i,\Omega,\omega,M)\)分别为：半长径\(a\)、偏心率\(e\)、倾角\(i\)、升交点经度\(\Omega\)、近心点角距\(\omega\)和平近点角\(M = n(t-\tau)\)，其中\(n = \sqrt{GM/a^3}\)为平均运动。

\textbf{定理 x+2.1}（Lagrange行星方程）。设摄动势为\(\mathcal{R}\)，则密切要素的变率满足：

\[\frac{da}{dt} = \frac{2}{na}\frac{\partial \mathcal{R}}{\partial M}, \quad \frac{de}{dt} = \frac{1-e^2}{na^2 e}\frac{\partial \mathcal{R}}{\partial M} - \frac{\sqrt{1-e^2}}{na^2 e}\frac{\partial \mathcal{R}}{\partial \omega}\]

\[\frac{di}{dt} = \frac{\cot i}{na^2\sqrt{1-e^2}}\frac{\partial \mathcal{R}}{\partial \omega} - \frac{\csc i}{na^2\sqrt{1-e^2}}\frac{\partial \mathcal{R}}{\partial \Omega}\]

\[\frac{d\Omega}{dt} = \frac{1}{na^2\sqrt{1-e^2}\sin i}\frac{\partial \mathcal{R}}{\partial i}, \quad \frac{d\omega}{dt} = \frac{\sqrt{1-e^2}}{na^2 e}\frac{\partial \mathcal{R}}{\partial e} - \frac{\cot i}{na^2\sqrt{1-e^2}}\frac{\partial \mathcal{R}}{\partial i}\]

\[\frac{dM}{dt} = n - \frac{2}{na}\frac{\partial \mathcal{R}}{\partial a} - \frac{1-e^2}{na^2 e}\frac{\partial \mathcal{R}}{\partial e}\]

\textbf{证明}：从变分常数法出发。Kepler运动由六个积分常数\(\sigma = (a,e,i,\Omega,\omega,\tau)\)参数化。将其视为时间的函数，满足密切条件：在任何时刻\(t\)，真实轨道与具有该时刻轨道要素的Kepler轨道相切。由位置和速度的连续性条件，可得

\[\left(\frac{\partial \mathbf{r}}{\partial \sigma}\right) \cdot \dot{\sigma} = 0, \quad \left(\frac{\partial \mathbf{v}}{\partial \sigma}\right) \cdot \dot{\sigma} = \nabla \mathcal{R}\]

利用Poisson括号恒等式\(\{\sigma_i,\sigma_j\} = \partial \sigma_i/\partial \mathbf{r} \cdot \partial \sigma_j/\partial \mathbf{v} - \partial \sigma_i/\partial \mathbf{v} \cdot \partial \sigma_j/\partial \mathbf{r}\)将上述方程化为标准形式。特别地，Lagrange括号\(\{\sigma_i,\sigma_j\}\)对Kepler运动是非退化的，其逆给出上述各方程的具体系数。以Delaunay变量\((L,G,H,\ell,g,h)\)表述最为自然，其中\(L = \sqrt{GMa}\)，\(G = L\sqrt{1-e^2}\)，\(H = G\cos i\)。 \(\blacksquare\)

\subsection{x+2.2 Hill方程的Floquet理论}\label{x2.2-hillux65b9ux7a0bux7684floquetux7406ux8bba}

\textbf{定义 x+2.2}（Hill方程）。圆型限制3体问题中第三体的小幅运动在旋转坐标系下由Hill方程描述：

\[\ddot{\xi} - 2\dot{\eta} = U_{xx}\xi + U_{xy}\eta, \quad \ddot{\eta} + 2\dot{\xi} = U_{xy}\xi + U_{yy}\eta\]

消去耦合得到单个4阶标量方程。对平面运动的线性化得到

\[\ddot{x} + p(t)x = 0, \quad p(t+T) = p(t)\]

其中\(p(t)\)为质量参数\(\mu\)和轨道频率的有界周期函数。

\textbf{定理 x+2.2}（Floquet定理）。Hill方程的基本解组可写为

\[\Phi(t) = P(t)e^{tB}\]

其中\(P(t+T) = P(t)\)为周期矩阵，\(B\)为常数矩阵。特征乘子\(\rho_{1,2}\)由\(\det \Phi(T)\)决定，满足\(\rho_1\rho_2 = 1\)。稳定性的充要条件为\(|\text{tr}\Phi(T)| < 2\)。

\textbf{证明}：记\(\Phi(t)\)为基本矩阵解，满足\(\dot{\Phi} = A(t)\Phi\)，\(\Phi(0) = I\)。由\(A(t+T) = A(t)\)，\(\Psi(t) = \Phi(t+T)\)也是基本矩阵解，因此存在常数矩阵\(M\)使得\(\Phi(t+T) = \Phi(t)M\)，其中\(M = \Phi(T)\)称为单值矩阵。定义\(B = \frac{1}{T}\log M\)（取主值分支以获得实矩阵），令\(P(t) = \Phi(t)e^{-tB}\)，则\(P(t+T) = \Phi(t+T)e^{-(t+T)B} = \Phi(t)Me^{-TB}e^{-tB} = \Phi(t)e^{-tB} = P(t)\)，证得Floquet分解。对于Hill方程，Liouville公式给出\(\det\Phi(T) = \exp(\int_0^T \text{tr}A(s)ds) = 1\)（\(A\)为零迹矩阵），故\(\rho_1\rho_2 = 1\)。稳定性要求特征乘子在单位圆上，即\(\rho = e^{\pm i\theta}\)，此时\(\text{tr}M = 2\cos\theta\)，故\(|\text{tr}M| < 2\)充分必要条件。 \(\blacksquare\)

\subsection{x+2.3 Hénon-Heiles系统的混沌}\label{x2.3-huxe9non-heilesux7cfbux7edfux7684ux6df7ux6c8c}

\textbf{定义 x+2.3}（Hénon-Heiles系统，1964）。描述星系中恒星在轴对称引力势中运动的简化模型，Hamilton函数为

\[H = \frac{1}{2}(p_x^2 + p_y^2) + \frac{1}{2}(x^2 + y^2) + x^2y - \frac{1}{3}y^3\]

这是\(n=2\)自由度的不可积Hamilton系统。

\textbf{定理 x+2.3}（Hénon-Heiles系统的Poincaré截面分析）。当能量\(E\)从0增大时，系统经历从近可积（规则环面占主导）到全局混沌（大面积混沌海）的转变。具体而言，\(E \lesssim 1/12\)时绝大部分轨道是规则的；\(E \approx 1/8\)时出现显著的混沌区域，与不稳定周期轨道的同宿纠缠相关；\(E \to 1/6\)（逃逸能量）时几乎所有轨道都成为混沌的，最终逃逸。

\textbf{证明思路}：系统的离散对称群为\(C_{3v}\)。对于\(E < 1/6\)，运动被约束在等势能线\(V = \frac{1}{2}(x^2+y^2) + x^2y - \frac{1}{3}y^3 \leq E\)的紧区域内。在Poincaré截面\(y=0, p_y > 0\)上记录\((x,p_x)\)的迭代映射。对规则轨道，截面点落在不变量环面上（KAM环面的截线），在截面上表现为闭合曲线；对混沌轨道，截面点随机分布在一个二维区域中。数值结果表明分界线处的稳定和不稳定流形横截相交（Poincaré-Birkhoff-Smale定理的体现），产生同宿纠缠和Smale马蹄，导致符号动力系统和正拓扑熵。 \(\blacksquare\)

\subsection{x+2.4 Schwarzschild时空中的轨道}\label{x2.4-schwarzschildux65f6ux7a7aux4e2dux7684ux8f68ux9053}

\textbf{定义 x+2.4}（Schwarzschild度规）。球对称静态时空的线元为

\[ds^2 = -\left(1 - \frac{2M}{r}\right)dt^2 + \left(1 - \frac{2M}{r}\right)^{-1}dr^2 + r^2(d\theta^2 + \sin^2\theta d\phi^2)\]

（采用\(G=c=1\)单位制）。

\textbf{定理 x+2.4}（Schwarzschild轨道方程）。赤道面（\(\theta = \pi/2\)）上测试粒子的轨道满足

\[\left(\frac{dr}{d\phi}\right)^2 = \frac{r^4}{L^2}\left[E^2 - \left(1 - \frac{2M}{r}\right)\left(1 + \frac{L^2}{r^2}\right)\right]\]

其中\(E = (1-2M/r)dt/d\tau\)（能量），\(L = r^2 d\phi/d\tau\)（角动量）为运动常数。

\textbf{证明}：由球对称性和静态性，\(\partial_t\)和\(\partial_\phi\)为Killing矢量，给出两个守恒量。对类时测地线，\(g_{\mu\nu}dx^\mu dx^\nu = -1\)，代入度规和守恒量：

\[-\frac{E^2}{1-2M/r} + \frac{1}{1-2M/r}\left(\frac{dr}{d\tau}\right)^2 + \frac{L^2}{r^2} = -1\]

解出\((dr/d\tau)^2\)并利用\(d\phi/d\tau = L/r^2\)消去\(\tau\)得到\(dr/d\phi\)表达式。引入\(u = 1/r\)变换，得到

\[\left(\frac{du}{d\phi}\right)^2 = \frac{E^2-1}{L^2} + \frac{2M}{L^2}u - u^2 + 2Mu^3\]

右端最后一项\(2Mu^3\)是广义相对论修正项。忽略此项回到Newton轨道方程。对该方程求导得到二阶形式：

\[\frac{d^2u}{d\phi^2} + u = \frac{M}{L^2} + 3Mu^2\]

\(3Mu^2\)导致了近日点进动。对近圆轨道，将\(u = u_0 + \delta u\)线性化得到进动角\(\Delta\phi = 6\pi M/a(1-e^2)\)每周期。 \(\blacksquare\)

\subsection{x+2.5 后Newton近似}\label{x2.5-ux540enewtonux8fd1ux4f3c}

\textbf{定义 x+2.5}（后Newton展开）。在弱场慢速近似下，度规展开为

\[g_{00} = -1 + 2U - 2U^2 + \cdots, \quad g_{0i} = -4V_i + \cdots, \quad g_{ij} = (1+2U)\delta_{ij} + \cdots\]

其中\(U\)为Newton势，\(V_i\)为矢量势（数量级为\(O(v^3)\)）。

\textbf{定理 x+2.5}（1PN二体运动方程）。在后Newton 1PN近似下，相对运动方程（Einstein-Infeld-Hoffmann）为

\[\ddot{\mathbf{r}} = -\frac{GM}{r^3}\mathbf{r} + \frac{GM}{c^2 r^3}\left[\mathbf{r}\left(\frac{4GM}{r} - v^2\right) + 4(\mathbf{r}\cdot\mathbf{v})\mathbf{v}\right]\]

其中\(\mathbf{r} = \mathbf{r}_1 - \mathbf{r}_2\)，\(M = m_1 + m_2\)，\(\mathbf{v} = \dot{\mathbf{r}}\)。由此推得的近日点进动恰为\(\Delta\phi = 6\pi GM/a(1-e^2)c^2\)。

\textbf{证明}：从1PN度规出发，作用量\(S = -m\int d\tau\)按\(v/c\)展开到\(O(v^4/c^4)\)。对二体系统，在多极展开后，1PN Lagrangian为

\[L = -Mc^2 + \frac{1}{2}\mu v^2 + \frac{GM\mu}{r} + \frac{1}{8c^2}\left[(1-3\nu)\mu v^4 + \cdots\right]\]

其中\(\mu = m_1 m_2/M\)为约化质量，\(\nu = \mu/M\)。通过变分得到运动方程。对轨道要素应用Lagrange行星方程，计算\(\partial \mathcal{R}/\partial \omega\)积分到1PN阶即得进动角。 \(\blacksquare\)

\subsection{x+2.6 引力波轨道的Peters-Mathews公式}\label{x2.6-ux5f15ux529bux6ce2ux8f68ux9053ux7684peters-mathewsux516cux5f0f}

\textbf{定理 x+2.6}（Peters-Mathews公式，1963）。双星系统因引力波辐射导致的轨道半长径和偏心率的变化率由下式给出：

\[\left\langle \frac{da}{dt} \right\rangle = -\frac{64}{5}\frac{G^3}{c^5}\frac{m_1 m_2(m_1+m_2)}{a^3(1-e^2)^{7/2}}\left(1 + \frac{73}{24}e^2 + \frac{37}{96}e^4\right)\]

\[\left\langle \frac{de}{dt} \right\rangle = -\frac{304}{15}\frac{G^3}{c^5}\frac{m_1 m_2(m_1+m_2)}{a^4(1-e^2)^{5/2}}e\left(1 + \frac{121}{304}e^2\right)\]

其中\(\langle\cdot\rangle\)表示一个轨道周期内的平均值。

\textbf{推导}：从线性化Einstein场方程的引力波解出发。辐射功率由四极公式给出：

\[\dot{E} = -\frac{G}{5c^5}\left\langle \dddot{Q}_{ij} \dddot{Q}_{ij} \right\rangle\]

其中\(Q_{ij}\)为质量四极矩（无迹部分）。对Kepler轨道：

\[Q_{ij} = \mu\left(r_i r_j - \frac{1}{3}r^2\delta_{ij}\right)\]

其中\(\mu = m_1 m_2/(m_1+m_2)\)。计算一个周期内的平均需要Keper轨道的参数化。利用Fourier展开\(r(t) = a(1-e\cos E)\)，\(E - e\sin E = n(t-\tau)\)，以及

\[\ddot{r}_i \ddot{r}_j \text{之类的项的展开}\]

对\(\dddot{Q}_{ij}\)求三阶时间导数，代入四极公式并在轨道的Kepler椭圆上作平均。这涉及Bessel函数\(J_n(ne)\)的计算。详细过程：将Kepler运动中的坐标和速度展开为Fourier级数------\(r/a\)和\(\phi\)的Fourier系数由Hansen系数\(X_k^{n,m}(e)\)给出。经过冗长但对Bessel函数的恒等式计算后得到上述闭式结果。特别地，偏心率\(e(t)\)与半长径\(a(t)\)的演化满足

\[\frac{de/dt}{da/dt} = \frac{19}{12}\frac{e}{a}\frac{1 + \frac{121}{304}e^2}{1 + \frac{73}{24}e^2 + \frac{37}{96}e^4}\]

积分此关系得到\(e(a)\)，进而可得到并合时间\(T_{\text{merge}} = \int_{a_0}^0 da/(da/dt)\)。 \(\blacksquare\)

\newpage

\chapter{13-数学物理/数学物理/14-地球物理学数学理论.md}\label{ux6570ux5b66ux7269ux7406ux6570ux5b66ux7269ux740614-ux5730ux7403ux7269ux7406ux5b66ux6570ux5b66ux7406ux8bba.md}

\chapter{地球物理学数学理论}\label{ux5730ux7403ux7269ux7406ux5b66ux6570ux5b66ux7406ux8bba}

\begin{quote}
覆盖 MSC 86-XX 地球物理学。论述地震波传播与层析成像、地幔对流与地球动力学、重力场与地磁场数学理论。
\end{quote}

\hrulefill

\section{第484章：地震波传播与层析成像}\label{ux7b2c484ux7ae0ux5730ux9707ux6ce2ux4f20ux64adux4e0eux5c42ux6790ux6210ux50cf}

地震波是探测地球内部结构的最重要工具。当地震发生时，弹性波从震源向外传播，其走时、振幅和波形携带着传播路径上介质物性（波速、密度、衰减）的丰富信息。地震层析成像（seismic tomography）利用在地表或钻孔中记录的波形数据，通过求解大规模反问题来重建地球内部的三维结构。本章从弹性波方程的基本解结构出发，依次建立射线理论的WKBJ近似的数学基础、Born近似在全波形反演中的Fréchet导数理论，以及走时层析成像中Radon变换与反问题理论的深层联系。

\subsection{x.1 弹性波方程}\label{x.1-ux5f39ux6027ux6ce2ux65b9ux7a0b}

\textbf{定义 x.1}（弹性动力学方程）。在各向同性弹性介质中，位移场\(\mathbf{u}(\mathbf{x},t)\)满足

\[\rho \frac{\partial^2 \mathbf{u}}{\partial t^2} = (\lambda + \mu)\nabla(\nabla \cdot \mathbf{u}) + \mu\nabla^2 \mathbf{u} + \mathbf{f}\]

其中\(\rho\)为密度，\(\lambda,\mu\)为Lamé常数，\(\mathbf{f}\)为体力源。

\textbf{定理 x.1}（Helmholtz分解与P波-S波分离）。弹性位移场可唯一分解为无旋部分（P波）和无散部分（S波）：\(\mathbf{u} = \nabla\phi + \nabla \times \boldsymbol{\psi}\)，\(\nabla \cdot \boldsymbol{\psi} = 0\)。标量势\(\phi\)和矢量势\(\boldsymbol{\psi}\)分别满足波动方程

\[\frac{\partial^2 \phi}{\partial t^2} = \alpha^2 \nabla^2 \phi, \quad \frac{\partial^2 \boldsymbol{\psi}}{\partial t^2} = \beta^2 \nabla^2 \boldsymbol{\psi}\]

其中\(\alpha = \sqrt{(\lambda+2\mu)/\rho}\)为P波速度，\(\beta = \sqrt{\mu/\rho}\)为S波速度。

\textbf{证明}：对无体力情况，方程写为\(\rho \ddot{\mathbf{u}} = (\lambda+\mu)\nabla(\nabla\cdot\mathbf{u}) + \mu\nabla^2\mathbf{u}\)。利用恒等式\(\nabla^2\mathbf{u} = \nabla(\nabla\cdot\mathbf{u}) - \nabla\times(\nabla\times\mathbf{u})\)，得到\(\rho \ddot{\mathbf{u}} = (\lambda+2\mu)\nabla(\nabla\cdot\mathbf{u}) - \mu\nabla\times(\nabla\times\mathbf{u})\)。代入Helmholtz分解并分别取散度和旋度：对散度得\(\nabla\cdot\ddot{\mathbf{u}} = \nabla^2\ddot{\phi} = \alpha^2\nabla^2(\nabla^2\phi)\)，即\(\ddot{\phi} = \alpha^2\nabla^2\phi\)；对旋度得\(\nabla\times\ddot{\mathbf{u}} = -\nabla^2\ddot{\boldsymbol{\psi}} = \beta^2\nabla^2(-\nabla^2\boldsymbol{\psi})\)，即\(\ddot{\boldsymbol{\psi}} = \beta^2\nabla^2\boldsymbol{\psi}\)。唯一性由\(\nabla\phi = 0\)和\(\nabla\times\boldsymbol{\psi}=0\)在\(\nabla\cdot\boldsymbol{\psi}=0\)条件下推出\(\phi\equiv0,\boldsymbol{\psi}\equiv0\)来保证。 \(\blacksquare\)

\textbf{定理 x.2}（Love波与Rayleigh波）。在半空间自由表面附近存在表面波：（1）Rayleigh波------P波和SV波的非均匀平面波组合，沿\(z=0\)自由表面传播，位移随深度指数衰减，速度为\(c_R < \beta\)，由Rayleigh方程\(r(c) = (2-c^2/\beta^2)^2 - 4\sqrt{1-c^2/\alpha^2}\sqrt{1-c^2/\beta^2}=0\)确定；（2）Love波------水平层状介质中的SH型表面波，存在多阶简正模式，由频散方程\(\tan(\omega H\sqrt{1/c^2-1/\beta_1^2}) = \frac{\mu_2\sqrt{1/c^2-1/\beta_2^2}}{\mu_1\sqrt{1/\beta_1^2-1/c^2}}\)确定。

\textbf{证明}（Rayleigh波）：设位移势\(\phi = Ae^{-kz\sqrt{1-c^2/\alpha^2}}e^{ik(x-ct)}\)，\(\psi = Be^{-kz\sqrt{1-c^2/\beta^2}}e^{ik(x-ct)}\)（SV分量）。在自由表面\(z=0\)处应力为零：\(\sigma_{zz} = \sigma_{xz} = 0\)。由各向同性应力-应变关系得出一组关于\(A,B\)的线性齐次方程，其非零解条件给出Rayleigh方程。此方程在\((0,\beta)\)内恰好有一个实根\(c_R\)。 \(\blacksquare\)

\subsection{x.2 射线理论与WKBJ近似}\label{x.2-ux5c04ux7ebfux7406ux8bbaux4e0ewkbjux8fd1ux4f3c}

\textbf{定义 x.2}（程函方程与传输方程）。在频率\(\omega \to \infty\)的渐近极限下，设Helmholtz方程\(\nabla^2 u + \frac{\omega^2}{c^2(\mathbf{x})}u = 0\)的解形式为\(u(\mathbf{x}) = e^{i\omega T(\mathbf{x})}\sum_{k=0}^{\infty} A_k(\mathbf{x})(i\omega)^{-k}\)。代入方程，按\(\omega\)幂次整理：

\(\omega^2\)项（程函方程）：\(|\nabla T|^2 = \frac{1}{c^2(\mathbf{x})}\)

\(\omega^1\)项（传输方程）：\(2\nabla T \cdot \nabla A_0 + A_0 \nabla^2 T = 0\)

\textbf{定理 x.3}（射线方程与特征线法）。程函方程的特征线（射线）满足一阶微分方程组

\[\frac{d\mathbf{x}}{d\tau} = c^2 \mathbf{p}, \quad \frac{d\mathbf{p}}{d\tau} = -\frac{1}{c}\nabla c\]

其中\(\mathbf{p} = \nabla T\)为慢度矢量，\(\tau\)为沿射线的参数。沿射线，走时\(T\)和振幅\(A_0\)可由积分给出

\[T = \int_{\text{ray}} \frac{ds}{c(\mathbf{x})}, \quad A_0^2(\mathbf{x}) = A_0^2(\mathbf{x}_0)\frac{c(\mathbf{x})J(\mathbf{x}_0)}{c(\mathbf{x}_0)J(\mathbf{x})}\]

其中\(J\)为射线坐标变换的Jacobi行列式，代表几何扩散因子。

\textbf{证明}：由程函方程\(|\nabla T|^2 = c^{-2}\)对坐标求梯度：\(2\nabla T \cdot \nabla(\nabla T) = -2c^{-3}\nabla c\)。令\(\mathbf{p} = \nabla T\)，\(\frac{d\mathbf{x}}{d\tau} = c^2\mathbf{p}\)，则\(\frac{d\mathbf{p}}{d\tau} = \frac{d\mathbf{x}}{d\tau}\cdot\nabla\mathbf{p} = c^2\mathbf{p}\cdot\nabla\mathbf{p} = c^2\frac{1}{2}\nabla|\mathbf{p}|^2 = c^2\frac{1}{2}\nabla(c^{-2}) = -\frac{1}{c}\nabla c\)，证得射线方程。传输方程沿射线写为\(\frac{dA_0}{d\tau} + \frac{1}{2}A_0\nabla\cdot(c^2\mathbf{p}) = 0\)，通过射线管截面面积与Jacobi行列式的关系导出振幅公式。 \(\blacksquare\)

\subsection{x.3 Born近似与全波形反演}\label{x.3-bornux8fd1ux4f3cux4e0eux5168ux6ce2ux5f62ux53cdux6f14}

\textbf{定义 x.3}（全波形反演的Fréchet导数）。设波速场\(c(\mathbf{x})\)产生扰动\(\delta c(\mathbf{x})\)。观测波场\(u(\mathbf{x}_r,t;\mathbf{x}_s)\)（接收点\(r\)，震源\(s\)）关于波速参数的Fréchet导数\(\frac{\delta u}{\delta c}\)满足

\[\delta u(\mathbf{x}_r,t;\mathbf{x}_s) = \int d\mathbf{x} \frac{\delta u(\mathbf{x}_r,t)}{\delta c(\mathbf{x})} \delta c(\mathbf{x})\]

\textbf{定理 x.4}（Born近似下的Fréchet导数表达）。在单散射（Born）近似下，标量波方程\(\frac{1}{c^2}\frac{\partial^2 u}{\partial t^2} - \nabla^2 u = f(t)\delta(\mathbf{x}-\mathbf{x}_s)\)关于波速的Fréchet核为

\[K(\mathbf{x}_r,\mathbf{x},\mathbf{x}_s,t) = \frac{2}{c^3(\mathbf{x})}G(\mathbf{x}_r,\mathbf{x},t) * \frac{\partial^2}{\partial t^2}G(\mathbf{x},\mathbf{x}_s,t)\]

其中\(G\)为背景介质中的Green函数，\(*\)表示时间卷积。

\textbf{完整推导}：设背景介质\(c_0(\mathbf{x})\)为参考模型，扰动\(\delta c\)产生波场的扰动\(\delta u\)。总波场\(u = u_0 + \delta u\)满足波动方程。将\(c = c_0 + \delta c\)代入，按一阶扰动展开。原方程为

\[\frac{1}{(c_0 + \delta c)^2}\frac{\partial^2(u_0 + \delta u)}{\partial t^2} - \nabla^2(u_0 + \delta u) = f\]

展开\(\frac{1}{(c_0 + \delta c)^2} \approx \frac{1}{c_0^2} - \frac{2\delta c}{c_0^3}\)，保留一阶项：

\[\frac{1}{c_0^2}\frac{\partial^2 \delta u}{\partial t^2} - \nabla^2 \delta u = \frac{2\delta c}{c_0^3}\frac{\partial^2 u_0}{\partial t^2}\]

这可视作以\(s(\mathbf{x},t) = \frac{2\delta c}{c_0^3}\frac{\partial^2 u_0}{\partial t^2}\)为源的方程。由Green函数表示得

\[\delta u(\mathbf{x}_r,t) = \int d\mathbf{x} \int d\tau G(\mathbf{x}_r,\mathbf{x},t-\tau) \frac{2\delta c(\mathbf{x})}{c_0^3(\mathbf{x})} \frac{\partial^2 u_0(\mathbf{x},\tau)}{\partial \tau^2}\]

而\(u_0(\mathbf{x},\tau) = \int d\tau' G(\mathbf{x},\mathbf{x}_s,\tau-\tau')f(\tau')\)。代入得双散射表示，Fréchet核即为两Green函数的卷积。在频域中表述更简洁：

\[\delta u(\mathbf{x}_r,\omega) = -2\omega^2 \int d\mathbf{x} \frac{\delta c(\mathbf{x})}{c_0^3(\mathbf{x})} G(\mathbf{x}_r,\mathbf{x},\omega)G(\mathbf{x},\mathbf{x}_s,\omega)\]

此为经典的全波形反演敏感核公式。 \(\blacksquare\)

\subsection{x.4 走时层析成像与Radon变换}\label{x.4-ux8d70ux65f6ux5c42ux6790ux6210ux50cfux4e0eradonux53d8ux6362}

\textbf{定义 x.4}（Radon变换）。函数\(f(\mathbf{x})\)在\(\mathbb{R}^2\)上的Radon变换定义为沿直线\(L(\theta,s) = \{\mathbf{x} : \mathbf{x}\cdot\mathbf{n}(\theta) = s\}\)的线积分

\[(\mathcal{R}f)(\theta,s) = \int_{L(\theta,s)} f(\mathbf{x}) d\ell\]

其中\(\mathbf{n}(\theta) = (\cos\theta,\sin\theta)\)为单位法向量。

\textbf{定理 x.5}（反演公式与Fourier切片定理）。Radon变换满足Fourier切片定理：\(\widehat{\mathcal{R}f}(\theta,\omega) = \hat{f}(\omega\mathbf{n}(\theta))\)，其中\(\hat{f}\)为二维Fourier变换，\(\widehat{\mathcal{R}f}\)为关于\(s\)的一维Fourier变换。由此可得反演公式

\[f(\mathbf{x}) = \frac{1}{2\pi}\int_0^{\pi} d\theta \int_{-\infty}^{\infty} d\omega |\omega| \widehat{\mathcal{R}f}(\theta,\omega) e^{i\omega \mathbf{x}\cdot\mathbf{n}(\theta)}\]

\textbf{证明}：按定义\(\widehat{\mathcal{R}f}(\theta,\omega) = \int ds e^{-i\omega s}\int f(\mathbf{x})\delta(s-\mathbf{x}\cdot\mathbf{n})d\mathbf{x} = \int f(\mathbf{x})e^{-i\omega\mathbf{x}\cdot\mathbf{n}}d\mathbf{x} = \hat{f}(\omega\mathbf{n})\)，证得切片定理。反演公式为逆Fourier变换在极坐标下的直接表达：\(f(\mathbf{x}) = \frac{1}{(2\pi)^2}\int_0^{2\pi}\int_0^{\infty} \hat{f}(\omega\mathbf{n})e^{i\omega\mathbf{x}\cdot\mathbf{n}}\omega d\omega d\theta\)。由\(\hat{f}(\omega\mathbf{n}) = \widehat{\mathcal{R}f}(\theta,\omega)\)，并将\(\omega\)范围对称扩展并利用\(\widehat{\mathcal{R}f}(\theta,-\omega) = \widehat{\mathcal{R}f}(\theta+\pi,\omega)\)，整理得所需形式。滤波反投影算法直接在离散域实现,核心为斜坡滤波器\(|\omega|\)。 \(\blacksquare\)

\textbf{定理 x.6}（走时层析成像与反问题）。在射线近似下，走时残差\(\delta T\)与慢度扰动\(\delta s = \delta(1/c)\)满足线性积分关系\(\delta T_i = \int_{L_i} \delta s(\mathbf{x}) d\ell\)。这是一类大型稀疏线性反问题，通常通过正则化的最小二乘法\(\min_{\delta s}\{\|\mathbf{A}\delta s - \mathbf{d}\|^2 + \lambda\|\nabla \delta s\|^2\}\)求解，其中\(\mathbf{A}\)为射线路径矩阵，\(\lambda\)为正则化参数。

\hrulefill

\section{第485章：地幔对流与地球动力学}\label{ux7b2c485ux7ae0ux5730ux5e54ux5bf9ux6d41ux4e0eux5730ux7403ux52a8ux529bux5b66}

地球内部的热对流是板块构造、地震活动和火山作用的根本驱动力。地幔在数百万年时间尺度上表现为高粘性流体，其运动由浮力驱动的Navier-Stokes方程在Boussinesq近似下描述。本章从Rayleigh-Bénard对流的线性稳定性分析出发，建立临界Rayleigh数的完整数学推导，进而将问题推广到地幔尺度的球壳对流，讨论板块构造的流体动力学理论基础，最后涉及冰后回弹这一古气候记录反演地幔粘度的经典地球物理反问题。

\subsection{x+1.1 Rayleigh-Bénard对流}\label{x1.1-rayleigh-buxe9nardux5bf9ux6d41}

\textbf{定义 x+1.1}（Rayleigh-Bénard系统）。考虑两无限大平板间厚度为\(d\)、上下温差为\(\Delta T\)的流体层。在Boussinesq近似下，无量纲化方程为

\[\frac{1}{\Pr}\left(\frac{\partial \mathbf{u}}{\partial t} + \mathbf{u}\cdot\nabla\mathbf{u}\right) = -\nabla p + \nabla^2 \mathbf{u} + \mathrm{Ra} T \hat{\mathbf{z}}\]

\[\frac{\partial T}{\partial t} + \mathbf{u}\cdot\nabla T = \nabla^2 T, \quad \nabla\cdot\mathbf{u} = 0\]

其中\(\mathrm{Ra} = \frac{g\alpha\Delta T d^3}{\kappa\nu}\)为Rayleigh数，\(\Pr = \nu/\kappa\)为Prandtl数，\(\alpha\)为热膨胀系数，\(\nu\)为运动粘度，\(\kappa\)为热扩散率。

\textbf{定理 x+1.1}（临界Rayleigh数的完整推导）。无对流的基本解为线性温度分布\(T_0 = 1 - z\)，\(\mathbf{u}_0 = 0\)。线性扰动\((\mathbf{u}',T',p')\)在一阶近似下满足线性化方程。通过两次取旋度消去压力和水平分量后，竖直速度\(w\)的扰动满足

\[\left(\frac{1}{\Pr}\frac{\partial}{\partial t} - \nabla^2\right)\left(\frac{\partial}{\partial t} - \nabla^2\right)\nabla^2 w = \mathrm{Ra} \nabla_H^2 w\]

其中\(\nabla_H^2 = \partial_x^2 + \partial_y^2\)为水平Laplacian。设\(w = W(z)e^{i(k_x x + k_y y) + \sigma t}\)，\(W(z)\)满足

\[(D^2 - k^2)(D^2 - k^2 - \sigma)(D^2 - k^2 - \sigma/\Pr)W = -\mathrm{Ra} k^2 W\]

其中\(D = d/dz\)，\(k = \sqrt{k_x^2+k_y^2}\)。边界条件：无滑移刚性边界\(W = DW = 0\)在\(z=0,1\)处。交换稳定性原理证明了对刚性边界，不稳定性首先以稳态对流（\(\sigma=0\)）出现，方程简化为

\[(D^2 - k^2)^3 W = -\mathrm{Ra} k^2 W, \quad W = DW = 0 \text{ at } z=0,1\]

这是六阶特征值问题。通解为\(W = \sum_{j=1}^3 A_j \cosh(q_j(z-1/2)) + B_j\sinh(q_j(z-1/2))\)。对称性分析表明最低临界模态为\(W\)关于\(z=1/2\)对称者，即\(B_j=0\)。由边界条件得到的特征方程为\(4q^2\)和\(q\tan(q/2)\)等的组合，其数值最小解给出：临界波数\(k_c \approx 3.117\)，临界Rayleigh数\(\mathrm{Ra}_c \approx 1707.76\)。上下为自由滑移边界时\(\mathrm{Ra}_c = 27\pi^4/4 \approx 657.51\)。 \(\blacksquare\)

\subsection{x+1.2 地幔对流的Navier-Stokes-Boussinesq近似}\label{x1.2-ux5730ux5e54ux5bf9ux6d41ux7684navier-stokes-boussinesqux8fd1ux4f3c}

\textbf{定理 x+1.2}（无穷Prandtl数近似）。地幔的Prandtl数\(\Pr \sim 10^{23}\)，惯性项完全可忽略，方程简化为Stokes方程：

\[-\nabla p + \nabla^2 \mathbf{u} + \mathrm{Ra} T\hat{\mathbf{r}} = 0, \quad \nabla\cdot\mathbf{u} = 0\]

在球壳坐标\((r,\theta,\phi)\)内，速度场可借助广义球谐函数展开以分离变量，获得解析形式的本征解。

\textbf{证明}：从完整Boussinesq方程出发，除以\(\Pr\)并取\(\Pr \to \infty\)极限，惯性项消失。对Stokes方程取旋度两次得\(\nabla^4\mathbf{u} = \mathrm{Ra}(\nabla^2(T\mathbf{e}_r) - \nabla\nabla\cdot(T\mathbf{e}_r))\)。本征值问题仍如前述结构，只是球几何下用球Bessel函数替代双曲函数。球壳对流在径向边界条件（固定温度或固定热流）下确定\(\mathrm{Ra}_c\)关于球壳内外半径比的函数。三维球对流模式由球谐度\(\ell\)和径向阶数\(n\)刻画。 \(\blacksquare\)

\subsection{x+1.3 板块构造的流体动力学模型}\label{x1.3-ux677fux5757ux6784ux9020ux7684ux6d41ux4f53ux52a8ux529bux5b66ux6a21ux578b}

\textbf{定义 x+1.2}（岩石圈-软流圈系统）。岩石圈近似为高粘性（\(\eta_L \sim 10^{23}\) Pa s）的弹性-粘塑性薄壳，软流圈为低粘性（\(\eta_A \sim 10^{19}\) Pa s）的Newton流体。板块运动由俯冲板片负浮力（slab pull）、洋脊推力（ridge push）和地幔拖曳力（mantle drag）的平衡决定。

\textbf{定理 x+1.3}（俯冲板片负浮力）。俯冲板片作用于地表板块的力由下式给出

\[F_{\text{slab}} = \Delta\rho g h L \sin\theta\]

其中\(\Delta\rho \approx 60\) kg/m\(^3\)为板片与周围地幔的密度差，\(h \approx 100\) km为板片厚度，\(L\)为俯冲长度，\(\theta\)为俯冲角。当\(L \sim 600\) km和\(\theta \sim 45^\circ\)时，\(F_{\text{slab}} \sim 3 \times 10^{13}\) N/m，是板块运动的主要驱动力。

\subsection{x+1.4 冰后回弹与地幔粘度的反演}\label{x1.4-ux51b0ux540eux56deux5f39ux4e0eux5730ux5e54ux7c98ux5ea6ux7684ux53cdux6f14}

\textbf{定义 x+1.3}（冰后回弹的数学模型）。冰盖消融后，地壳对卸载的响应可用半空间Maxwell粘弹性模型描述。在Laplace域中，挠曲\(\hat{w}(k,s)\)满足

\[\hat{w}(k,s) = \frac{\hat{\sigma}(k,s)}{\rho g + D k^4 + \frac{\mu(s) k}{2}}\]

其中\(\hat{\sigma}\)为载荷的Laplace-Fourier变换，\(D\)为岩石圈挠曲刚度，\(\mu(s) = \frac{\eta s}{1 + \tau_M s}\)为Maxwell松弛模量（\(\tau_M = \eta / E\)为Maxwell时间）。

\textbf{定理 x+1.4}（地幔粘度的反演）。冰后回弹的观测数据（古海岸线高程的时空分布）可通过求解正问题（给定\(\eta\)剖面预测回弹历史）和反问题（由观测反推\(\eta\)）来确定地幔粘度的径向分布。反演的数学框架为非线性最小二乘

\[\chi^2(\boldsymbol{\eta}) = \sum_{i} \frac{(w_i^{\text{obs}} - w_i^{\text{pred}}(\boldsymbol{\eta}))^2}{\sigma_i^2}\]

\(\eta\)的敏感核\(\partial w_i / \partial \eta_j\)可由伴随方法高效计算。Fennoscandia和Laurentide冰盖的数据共同约束上地幔粘度\(\eta_{\text{UM}} \sim 10^{20}-10^{21}\) Pa s，下地幔粘度\(\eta_{\text{LM}} \sim 10^{22}-10^{23}\) Pa s，粘度跳变在660 km间断面上。

\hrulefill

\section{第486章：重力场与地磁场数学理论}\label{ux7b2c486ux7ae0ux91cdux529bux573aux4e0eux5730ux78c1ux573aux6570ux5b66ux7406ux8bba}

地球的重力场和磁场是约束其内部质量分布和动力学过程的最直接可观测量。重力场通过人造卫星轨道摄动和地面绝对重力测量获取，其球谐展开将Legendre多项式理论与观测紧密结合。大地水准面的起伏携带地幔对流环的信息，Stokes公式构成了物理大地测量学的理论基石。地磁场由外地核的磁流体发电机过程产生，其复杂的自激发电机制是非线性偏微分方程理论的精美范例，Cowling反发电机定理为理解磁场的来源提供了根本限制。

\subsection{x+2.1 引力位的球谐展开}\label{x2.1-ux5f15ux529bux4f4dux7684ux7403ux8c10ux5c55ux5f00}

\textbf{定义 x+2.1}（地球引力位）。地球外部引力位在球坐标\((r,\theta,\lambda)\)下的级数展开为

\[V(r,\theta,\lambda) = \frac{GM}{r}\left[1 + \sum_{n=2}^{\infty}\sum_{m=0}^{n}\left(\frac{a}{r}\right)^n P_n^m(\cos\theta)(C_{nm}\cos m\lambda + S_{nm}\sin m\lambda)\right]\]

其中\(a\)为地球赤道半径，\(P_n^m\)为连带Legendre函数，\(C_{nm},S_{nm}\)为Stokes系数。

\textbf{定理 x+2.1}（Legendre多项式与引力位的正交分解）。函数系\(\{P_n^m(\cos\theta)e^{im\lambda}\}\)满足正交性

\[\int_0^{2\pi}\int_0^{\pi} P_n^m(\cos\theta)P_{n'}^{m'}(\cos\theta)\cos(m\lambda)\cos(m'\lambda)\sin\theta d\theta d\lambda = \frac{4\pi}{(2n+1)}\frac{(n+m)!}{(n-m)!}\delta_{nn'}\delta_{mm'}\]

因此Stokes系数可由全球重力数据唯一确定：

\[C_{nm} + iS_{nm} = \frac{1}{2\pi GM}\int_{\Omega} \Delta g(\theta,\lambda) \bar{Y}_{nm}(\theta,\lambda) d\Omega\]

\textbf{证明}：Legendre函数的标准正交性为\(\int_{-1}^{1}P_n^m(x)P_{n'}^m(x)dx = \frac{2}{2n+1}\frac{(n+m)!}{(n-m)!}\delta_{nn'}\)。引入归一化球面谐函数\(Y_{nm}(\theta,\lambda) = \sqrt{\frac{(2n+1)}{4\pi}\frac{(n-m)!}{(n+m)!}}P_n^m(\cos\theta)e^{im\lambda}\)构成\(L^2(S^2)\)的正交基。Stokes系数为\(V\)在这个基上的展开系数，由球面上重力扰动\(\Delta g = -\partial V/\partial r - 2V/r\)通过逆变换关系反演。 \(\blacksquare\)

\subsection{x+2.2 大地水准面与Stokes公式}\label{x2.2-ux5927ux5730ux6c34ux51c6ux9762ux4e0estokesux516cux5f0f}

\textbf{定义 x+2.2}（大地水准面）。大地水准面为与全球平均海平面最接近的等引力位面\(W = V + \Phi = \text{const}\)，其中\(\Phi\)为离心力位。大地水准面起伏\(N\)由Bruns公式表示为\(N = T/\gamma\)，其中\(T\)为扰动位（实际引力位与参考椭球正常引力位之差），\(\gamma\)为正常重力。

\textbf{定理 x+2.2}（Stokes公式）。大地水准面起伏\(N\)可通过全球重力异常\(\Delta g\)的面积分表示：

\[N(\theta,\lambda) = \frac{R}{4\pi\gamma}\int_{\Omega} \Delta g(\theta',\lambda') S(\psi) d\Omega'\]

其中\(S(\psi)\)为Stokes函数

\[S(\psi) = \csc(\psi/2) - 6\sin(\psi/2) + 1 - 5\cos\psi - 3\cos\psi\ln(\sin(\psi/2)+\sin^2(\psi/2))\]

\(\psi\)为计算点\((\theta,\lambda)\)与积分点\((\theta',\lambda')\)之间的角距。

\textbf{推导}：由球面上的边值问题（第三边值问题），扰动位\(T\)满足\(\nabla^2 T = 0\)（地球外部），边界条件为\(\Delta g = -\partial T/\partial r - 2T/R\)在球面\(r=R\)上。将\(T\)展开为\(T(r,\theta,\lambda) = \sum_{n=0}^{\infty}\left(\frac{R}{r}\right)^{n+1}T_n(\theta,\lambda)\)，边界条件给出\(\Delta g = \sum_{n=0}^{\infty}\frac{n-1}{R}T_n\)（\(n=1\)项为零因质心在原点）。因此\(T_n = \frac{R}{n-1}\Delta g_n\)（\(n \geq 2\)）。大地水准面起伏为

\[N = \frac{T(R)}{\gamma} = \frac{R}{\gamma}\sum_{n=2}^{\infty}\frac{\Delta g_n}{n-1}\]

利用球面卷积定理：\(\sum_{n=2}^{\infty}\frac{P_n(\cos\psi)}{n-1} = S(\psi)/2\)（Stokes函数定义），即证。 \(\blacksquare\)

\subsection{x+2.3 地磁发电机理论}\label{x2.3-ux5730ux78c1ux53d1ux7535ux673aux7406ux8bba}

\textbf{定义 x+2.3}（磁流体动力学感应方程）。导电流体中的磁场演化由感应方程描述：

\[\frac{\partial \mathbf{B}}{\partial t} = \nabla \times (\mathbf{u} \times \mathbf{B}) + \eta \nabla^2 \mathbf{B}, \quad \nabla \cdot \mathbf{B} = 0\]

其中\(\mathbf{u}\)为流体速度，\(\eta = 1/(\mu_0\sigma)\)为磁扩散率。方程右端第一项为磁场的拉伸与平流（\(\alpha\)-效应的数学表述），第二项为欧姆耗散。

\textbf{定理 x+2.3}（平均场理论与\(\alpha\)-效应）。将速度和磁场分解为平均场与小尺度涨落：\(\mathbf{B} = \langle\mathbf{B}\rangle + \mathbf{b}\)，\(\mathbf{u} = \langle\mathbf{u}\rangle + \mathbf{v}\)。湍流涨落产生的平均电动势为

\[\langle\mathbf{v} \times \mathbf{b}\rangle = \alpha \langle\mathbf{B}\rangle - \beta \nabla \times \langle\mathbf{B}\rangle\]

其中\(\alpha\)与螺旋度\(\langle\mathbf{v}\cdot(\nabla\times\mathbf{v})\rangle\)相关，\(\beta\)为湍流磁扩散率。代入平均化感应方程得

\[\frac{\partial \langle\mathbf{B}\rangle}{\partial t} = \nabla \times (\langle\mathbf{u}\rangle \times \langle\mathbf{B}\rangle) + \nabla \times (\alpha \langle\mathbf{B}\rangle) + (\eta + \beta)\nabla^2 \langle\mathbf{B}\rangle\]

\textbf{证明}：小尺度磁场的线性化方程为\(\partial_t \mathbf{b} = \nabla\times(\langle\mathbf{u}\rangle\times\mathbf{b} + \mathbf{v}\times\langle\mathbf{B}\rangle) + \eta\nabla^2\mathbf{b} + \nabla\times(\mathbf{v}\times\mathbf{b} - \langle\mathbf{v}\times\mathbf{b}\rangle)\)。在准线性近似（忽略\(\mathbf{v}\times\mathbf{b}\)的涨落项）和短相关时间近似下求解\(\mathbf{b}\)，得到\(\mathbf{b}\)正比于\((\langle\mathbf{B}\rangle\cdot\nabla)\mathbf{v}\)与\(\mathbf{v}\times\langle\mathbf{B}\rangle\)的组合。再计算\(\langle\mathbf{v}\times\mathbf{b}\rangle\)利用湍流统计均匀各向同性（缺少镜面对称性，因旋转系统中螺旋度不为零），得上述\(\alpha\)-张量形式。各向同性时\(\alpha\)退化为标量。 \(\blacksquare\)

\textbf{定义 x+2.4}（\(\omega\)-效应）。差动旋转（如外地核中角速度随半径变化）将极向磁场拉伸为环向磁场。在柱坐标\((s,\phi,z)\)下，若平均流为差动旋转\(\langle\mathbf{u}\rangle = s\Omega(s,z)\hat{\boldsymbol{\phi}}\)，则环向场满足

\[\frac{\partial \langle B_\phi\rangle}{\partial t} \approx s(\nabla\Omega)\cdot\nabla(s\langle B_\phi\rangle) = s\langle\mathbf{B}_p\rangle\cdot\nabla\Omega\]

其中\(\langle\mathbf{B}_p\rangle\)为极向磁场。\(\alpha\)-\(\omega\)发电机的自激循环为：\(\alpha\)-效应由螺旋对流产生极向场，\(\omega\)-效应由差动旋转将极向场转换为环向场，形成正反馈。

\textbf{定理 x+2.4}（Cowling反发电机定理，1934）。轴对称磁场不可能由轴对称流体的发电机作用持续维持。

\textbf{证明}：设磁场\(\mathbf{B}\)和速度\(\mathbf{u}\)均关于某轴旋转对称。将\(\mathbf{B}\)分解为极向部分\(\mathbf{B}_p = \nabla\times(A\hat{\boldsymbol{\phi}}/s)\)和环向部分\(B_\phi\hat{\boldsymbol{\phi}}\)。感应方程的环向分量为

\[\frac{\partial B_\phi}{\partial t} + s\mathbf{u}_p\cdot\nabla\left(\frac{B_\phi}{s}\right) = s\mathbf{B}_p\cdot\nabla\Omega + \eta\left(\nabla^2 - \frac{1}{s^2}\right)B_\phi\]

极向分量可化为磁通函数\(A\)的方程：

\[\frac{\partial A}{\partial t} + \frac{1}{s}\mathbf{u}_p\cdot\nabla(sA) = \eta\left(\nabla^2 - \frac{1}{s^2}\right)A\]

此方程无源项------没有从速度场到\(A\)的驱动项。乘以\(sA\)在空间积分，利用边界条件（无穷远处\(A\to0\)）得

\[\frac{1}{2}\frac{d}{dt}\int s^2 A^2 dV = -\eta\int |\nabla(sA)|^2 dV \leq 0\]

故\(A \to 0\)当\(t \to \infty\)，极向场衰减为零。极向场消失后，环向场方程也无驱动源（因\(\mathbf{B}_p=0\)使\(s\mathbf{B}_p\cdot\nabla\Omega=0\)），同样衰减为零。因此轴对称发电机不可能维持稳态磁场。 \(\blacksquare\)

这一反定理说明地磁发电机必须包含非轴对称分量，旋转球形Couette流的发电机模式正是以三维螺旋对流为特征。

\backmatter
\end{document}